\documentclass[11pt,letterpaper,reqno]{amsart}
\makeatletter
\NeedsTeXFormat{LaTeX2e}[2020/10/01]

\RequirePackage[T1]{fontenc}
\RequirePackage{lmodern}
\RequirePackage{microtype}
\RequirePackage[letterpaper,margin=1in,headheight=15pt]{geometry}
\RequirePackage{amsmath,amsthm,amssymb,mathtools}
\RequirePackage{enumitem}
\RequirePackage{perpage}
\RequirePackage{xcolor}
\definecolor{manuscriptred}{rgb}{0.6,0.1,0.1}
\definecolor{manuscriptblue}{rgb}{0.1,0.3,0.5}
\definecolor{manuscriptmagenta}{HTML}{60445C}
\RequirePackage[breaklinks=true,colorlinks=true,
  linkcolor=manuscriptred,citecolor=manuscriptblue,urlcolor=manuscriptmagenta]{hyperref}
\RequirePackage{cleveref}
\hypersetup{pdfauthor={}}

\numberwithin{equation}{section}
\renewcommand{\theequation}{\thesection.\roman{equation}}
\renewcommand{\eqref}[1]{\textup{eq.~\ref{#1}}}
\crefname{equation}{eq.}{eqs.}
\Crefname{equation}{eq.}{eqs.}
\creflabelformat{equation}{#2#1#3}
\setcounter{tocdepth}{1}
\renewcommand{\thefigure}{\Roman{figure}}
\renewcommand{\thetable}{\Roman{table}}
\MakePerPage{footnote}
\renewcommand{\thefootnote}{%
  \ifnum\value{footnote}>9\arabic{footnote}\else\fnsymbol{footnote}\fi}

\theoremstyle{definition}
\newtheorem{ms@theorem}{Theorem}
\renewcommand{\thems@theorem}{\Alph{ms@theorem}}
\newtheorem*{ms@theoremstar}{Theorem}
\newtheorem*{ms@theoremzero}{Theorem 0}
\newtheorem{ms@proposition}{Proposition}[section]
\newtheorem{ms@lemma}[ms@proposition]{Lemma}
\newtheorem{ms@definition}[ms@proposition]{Definition}
\newtheorem{ms@corollary}[ms@proposition]{Corollary}
\newtheorem*{ms@corollarystar}{Corollary}

\newcommand{\ms@majorstart}{%
  \pushQED{\qed}\renewcommand{\qedsymbol}{$\blacksquare\square$}}
\newcommand{\ms@minorstart}{%
  \pushQED{\qed}\renewcommand{\qedsymbol}{$\blacksquare$}}
\newenvironment{theorem}[1][]{\ms@majorstart\ms@theorem[#1]}
  {\popQED\endms@theorem}
\newenvironment{theorem*}[1][]{\ms@majorstart\ms@theoremstar[#1]}
  {\popQED\endms@theoremstar}
\newenvironment{theoremzero}[1][]{\ms@majorstart\ms@theoremzero[#1]}
  {\popQED\endms@theoremzero}
\newenvironment{proposition}[1][]{\ms@minorstart\ms@proposition[#1]}
  {\popQED\endms@proposition}
\newenvironment{lemma}[1][]{\ms@minorstart\ms@lemma[#1]}
  {\popQED\endms@lemma}
\newenvironment{definition}[1][]{\ms@minorstart\ms@definition[#1]}
  {\popQED\endms@definition}
\newenvironment{corollary}[1][]{\ms@majorstart\ms@corollary[#1]}
  {\popQED\endms@corollary}
\newenvironment{corollary*}[1][]{\ms@majorstart\ms@corollarystar[#1]}
  {\popQED\endms@corollarystar}

\theoremstyle{remark}
\newtheorem{remark}[ms@proposition]{Remark}
\newtheorem*{remark*}{Remark}

\crefname{ms@theorem}{theorem}{theorems}
\Crefname{ms@theorem}{Theorem}{Theorems}
\crefname{ms@proposition}{proposition}{propositions}
\Crefname{ms@proposition}{Proposition}{Propositions}
\crefname{ms@lemma}{lemma}{lemmas}
\Crefname{ms@lemma}{Lemma}{Lemmas}
\crefname{ms@definition}{definition}{definitions}
\Crefname{ms@definition}{Definition}{Definitions}
\crefname{ms@corollary}{corollary}{corollaries}
\Crefname{ms@corollary}{Corollary}{Corollaries}
\crefname{remark}{remark}{remarks}
\Crefname{remark}{Remark}{Remarks}

\makeatother
\usepackage{mathrsfs,array,longtable,graphicx}
\allowdisplaybreaks
\setlength{\emergencystretch}{3em}
\DeclareOldFontCommand{\cal}{\relax}{\mathcal}
\newcommand{\R}{\mathbb R}
\newcommand{\C}{\mathbb C}
\newcommand{\Z}{\mathbb Z}
\newcommand{\N}{\mathbb N}
\newcommand{\E}{\mathbb E}
\newcommand{\Sspace}{\mathcal S}
\newcommand{\cS}{\mathcal S}
\newcommand{\dd}{\mathrm d}
\DeclareMathOperator{\supp}{supp}
\DeclareMathOperator{\tr}{tr}
\DeclareMathOperator{\diag}{diag}
\newsavebox{\MathBox}
\newcommand{\FitMath}[1]{\sbox{\MathBox}{$\displaystyle #1$}\ifdim\wd\MathBox>\linewidth\typeout{MATHFIT line=\the\inputlineno\space width=\the\wd\MathBox}\resizebox{\linewidth}{!}{\usebox{\MathBox}}\else\usebox{\MathBox}\fi}
\setlist[enumerate]{label=(\roman*)}
\hypersetup{hypertexnames=false,pdftitle={Four-dimensional SU(2) Yang-Mills theory: a proof outline for tr(F squared)},pdfauthor={}}
\crefname{ms@theorem}{Thm.}{Thms.}\Crefname{ms@theorem}{Thm.}{Thms.}
\crefname{ms@proposition}{Prop.}{Props.}\Crefname{ms@proposition}{Prop.}{Props.}
\crefname{ms@lemma}{Lem.}{Lems.}\Crefname{ms@lemma}{Lem.}{Lems.}
\crefname{ms@definition}{Def.}{Defs.}\Crefname{ms@definition}{Def.}{Defs.}
\crefname{ms@corollary}{Cor.}{Cors.}\Crefname{ms@corollary}{Cor.}{Cors.}
\crefname{remark}{Rem.}{Rems.}\Crefname{remark}{Rem.}{Rems.}
\crefname{section}{\S}{\S\S}\Crefname{section}{\S}{\S\S}
\crefname{subsection}{\S}{\S\S}\Crefname{subsection}{\S}{\S\S}
\crefname{subsubsection}{\S}{\S\S}\Crefname{subsubsection}{\S}{\S\S}
\crefname{paragraph}{\S}{\S\S}\Crefname{paragraph}{\S}{\S\S}
\crefformat{section}{\S#2#1#3}
\Crefformat{section}{\S#2#1#3}
\crefrangeformat{section}{\S\S#3#1#4--#5#2#6}
\Crefrangeformat{section}{\S\S#3#1#4--#5#2#6}
\crefmultiformat{section}{\S\S#2#1#3}{ and~#2#1#3}{, #2#1#3}{ and~#2#1#3}
\Crefmultiformat{section}{\S\S#2#1#3}{ and~#2#1#3}{, #2#1#3}{ and~#2#1#3}
\crefformat{subsection}{\S#2#1#3}
\Crefformat{subsection}{\S#2#1#3}
\crefrangeformat{subsection}{\S\S#3#1#4--#5#2#6}
\Crefrangeformat{subsection}{\S\S#3#1#4--#5#2#6}
\crefmultiformat{subsection}{\S\S#2#1#3}{ and~#2#1#3}{, #2#1#3}{ and~#2#1#3}
\Crefmultiformat{subsection}{\S\S#2#1#3}{ and~#2#1#3}{, #2#1#3}{ and~#2#1#3}
\crefformat{subsubsection}{\S#2#1#3}
\Crefformat{subsubsection}{\S#2#1#3}
\crefrangeformat{subsubsection}{\S\S#3#1#4--#5#2#6}
\Crefrangeformat{subsubsection}{\S\S#3#1#4--#5#2#6}
\crefmultiformat{subsubsection}{\S\S#2#1#3}{ and~#2#1#3}{, #2#1#3}{ and~#2#1#3}
\Crefmultiformat{subsubsection}{\S\S#2#1#3}{ and~#2#1#3}{, #2#1#3}{ and~#2#1#3}
\crefformat{paragraph}{\S#2#1#3}
\Crefformat{paragraph}{\S#2#1#3}
\crefrangeformat{paragraph}{\S\S#3#1#4--#5#2#6}
\Crefrangeformat{paragraph}{\S\S#3#1#4--#5#2#6}
\crefmultiformat{paragraph}{\S\S#2#1#3}{ and~#2#1#3}{, #2#1#3}{ and~#2#1#3}
\Crefmultiformat{paragraph}{\S\S#2#1#3}{ and~#2#1#3}{, #2#1#3}{ and~#2#1#3}
\title[Continuum correlations of $\operatorname{tr}(F^2)$]{Four-dimensional $\operatorname{SU}(2)$ Yang--Mills theory\newline A proof outline for $\operatorname{tr}(F^2)$}
\date{}
\setcounter{secnumdepth}{4}
\begin{document}
\pagenumbering{Alph}
\begin{abstract}
Under the stated uniform renormalization and short-distance hypotheses,
we outline the continuum construction of the renormalized
$\operatorname{tr}(F^2)$ observable in four-dimensional Wilson
$\operatorname{SU}(2)$ theory at sufficiently small reference coupling.
The proof follows the source-dependent multiscale integration through
uniform correlation estimates, volume-first limits and canonical
decomposition to scalar Osterwalder--Schrader reconstruction.
\end{abstract}
\maketitle
\tableofcontents

\section{Introduction}\label{sec:introduction}
\phantomsection\label{page:main-theorem}
\begin{theoremzero}
Assume the uniform renormalization and short-distance estimates
stated below. For sufficiently small reference coupling, the correlations of
$\operatorname{tr}(F^2)$ in four-dimensional Wilson $\operatorname{SU}(2)$
theory have nontrivial volume-first subsequential continuum limits.
In every canonical component, the Schwinger functions satisfy the
Osterwalder--Schrader axioms and reconstruct a local Hermitian scalar
Wightman theory with positive energy and a unique vacuum.
\end{theoremzero}
The precise statements are \cref{thm:continuum-main,thm:phase-main,thm:reconstruction-main}.

Here $\operatorname{tr}(F^2)$ is the scalar curvature-squared observable
with Wilson normalization. At each lattice cutoff we subtract the
expectation of its smeared version $\mathcal O(f)$. A volume-first limit
takes infinite volume at each fixed spacing before sending the spacing
to zero along the weak-coupling sequence defined in \cref{candidate:orientation}.
The Euclidean functionals are first constructed on tests with separated
supports and then extended to Schwartz tests flat on all coincidence
diagonals. The reconstructed field generates the scalar sector associated
with this observable.

The main task is uniform control of the source-dependent multiscale
integration. The background field, its constrained inverse, the polymer
sum and the large-field history sum must remain compatible on one
analytic domain. In particular, a derivative of the normalizer and a
change of exterior configuration must satisfy the same estimates used
to bound local correlations. The proof develops these estimates in
\cref{sec:background,sec:interactions} before passing to limits.

\Cref{sec:continuum} controls the number of insertions and extends the
limits to diagonal-flat tests. \Cref{sec:short-distance} gives the
short-distance law and covariance. \Cref{sec:reconstruction} treats
spectral representations, canonical components and scalar reconstruction.
\Cref{app:calculations} contains the auxiliary lattice estimates.

\clearpage
\pagenumbering{arabic}

\section{Lattice fields and the main results}\label{sec:lattice}
This section fixes the Wilson measure, the normalization of $\operatorname{tr}(F^2)$, the admissible limiting sequences and the test spaces. The two main continuum statements follow the setting. The final subsection states the analytic inputs used in the proof, with their domains and uniformity requirements.

% BEGIN MODULE chapters/foundations.tex
\subsection{Finite-cutoff identities and compactness}
\label{ch:foundations}

This section establishes the finite-dimensional identities and compactness
arguments used throughout the construction. None of its statements assumes
a continuum Yang--Mills measure. Uniformity in the cutoff requires constants independent of the number
of plaquettes, even though each integral is finite.

\phantomsection\label{notdef:bond-fields}
\subsubsection{The configuration space and its probability measure}

Let $G=\operatorname{SU}(2)$, the compact group of complex $2\times2$ matrices
$U$ with $U^*U=1$ and $\det U=1$. Its normalized Haar probability measure
is denoted by $dg$. Fix a finite oriented graph obtained from a cubic lattice
on a four-dimensional torus. Choose one orientation of each geometric bond;
write $B_+$ for this set and impose $U_{\bar b}=U_b^{-1}$ for the reverse
orientation. The configuration space is the compact space
$\mathcal V=G^{B_+}$, with product probability measure $dU$.

A path is an ordered list of consecutive oriented bonds. Its transport is
the ordered product of their matrices. For each geometrical elementary
square choose one boundary orientation and one starting vertex, and denote
its transport by $U_p$. Define

\begin{gather*}
\FitMath{a_p(U)=1-\tfrac12\operatorname{Re}\operatorname{tr}U_p,} \\[3pt]
\FitMath{A(U)=\sum_{p\in P}a_p(U),} \\[3pt]
\FitMath{N_P=|P|.}
\end{gather*}

Changing the starting vertex conjugates $U_p$, and reversing orientation
inverts it. Both changes preserve $a_p$.

\begin{proposition}[Finite-cutoff existence and elementary bounds]
For every real $\beta\ge0$,

\begin{gather*}
\FitMath{e^{-2\beta N_P}\le Z(\beta):=\int e^{-\beta A(U)}dU\le1.}
\end{gather*}

Consequently $Z(\beta)^{-1}e^{-\beta A(U)}dU$ is a well-defined probability
measure. Every bounded Borel observable has a finite expectation.
\end{proposition}
\begin{proof}
The eigenvalues of an element of $G$ are $e^{it},e^{-it}$ for a real $t$.
Thus $\frac12\operatorname{Re}\operatorname{tr}U_p=\cos t$ lies in $[-1,1]$,
so $0\le a_p\le2$ and $0\le A\le2N_P$. Exponentiate these inequalities and
integrate against a probability measure. The lower bound is strictly
positive and supplies the normalization.
\end{proof}

The proof uses only compactness and positivity. Its constants degenerate
when $N_P\to\infty$. In particular it supplies neither tightness of
renormalized local fields nor a useful estimate for a continuum limit.

A gauge transformation is a family $(h_x)$ of group elements indexed by
vertices. It acts by $U_b\mapsto h_{b_-}U_bh_{b_+}^{-1}$. Successive factors
cancel along a path, so a closed transport based at $x$ is conjugated by
$h_x$. Haar measure on each bond is unchanged by its independent left and
right multiplications. This proves gauge invariance of the action and of
the probability measure. A permutation of vertices and bonds coming from
a lattice isometry also preserves the product measure and permutes the
plaquettes. Therefore the finite torus measure has every isometry of that
torus. This conclusion concerns these finite torus measures; it is not a
proof that every infinite-volume periodic Gibbs state has every reflection.

\subsubsection{The continuum normalization is a local expansion}

Represent the Lie algebra by Hermitian traceless matrices and set
$|F|^2=\frac12\operatorname{tr}(F^2)$. For a fixed such matrix,

\begin{gather*}
\FitMath{1-\tfrac12\operatorname{Re}\operatorname{tr}e^{itF}
   =\tfrac{t^2}{4}\operatorname{tr}(F^2)+O(t^4|F|^4)
   =\tfrac12t^2|F|^2+O(t^4|F|^4).}
\end{gather*}

The expansion follows by the matrix exponential series: the linear trace
vanishes and the odd powers have zero real trace. If a smooth connection
has elementary holonomy $U_p=\exp(i a^2F_{\mu\nu}(x_p)+O(a^3))$, with all
required derivatives uniformly bounded, it follows that

\begin{gather*}
\FitMath{a_p=\tfrac12a^4|F_{\mu\nu}(x_p)|^2+O(a^5).}
\end{gather*}

Hence, on a fixed compact region, summing $f(x_p)a_p$ over plaquettes gives
the corresponding Riemann integral of
$\frac12 f\sum_{\mu<\nu}|F_{\mu\nu}|^2$ as $a\downarrow0$.

This is a deterministic smooth-field calculation. Typical configurations
under a quantum measure have not been shown to satisfy its hypotheses.
The calculation fixes the intended normalization of the observable; it
does not prove convergence of that observable under the lattice measures.

\subsubsection{Sources, cumulants, and centering}

For bounded functions $f_i$ on the plaquette barycentres put
$O_i=\sum_p f_i(x_p)a_p$ and

\begin{gather*}
\FitMath{Z(z)=\int\exp\{-\beta A(U)+\sum_{i=1}^n z_iO_i(U)\}\,dU,} \\[3pt]
\FitMath{z\in\mathbb C^n.}
\end{gather*}

The estimate $|O_i|\le2\sum_p|f_i(x_p)|$ gives an integrable uniform majorant
on every compact source set. Termwise differentiation of the exponential
series is therefore valid, and $Z$ is entire. Since $Z(0)>0$, continuity
gives a sufficiently small polydisc about zero where $Z$ has no zeros. On that polydisc
there is a unique holomorphic logarithm $W$ with $W(0)=\log Z(0)$.

The joint cumulant is

\begin{gather*}
\FitMath{\kappa(O_1,\ldots,O_n)=
 \left.\partial_{z_1}\cdots\partial_{z_n}W(z)\right|_{z=0}.}
\end{gather*}

Its expression in moments is

\begin{gather*}
\FitMath{\kappa(O_1,\ldots,O_n)=
 \sum_{\pi\in\operatorname{Part}(n)}
 (|\pi|-1)!(-1)^{|\pi|-1}
 \prod_{J\in\pi}\mathbb E\prod_{i\in J}O_i.}
\end{gather*}

Indeed, write $Z(z)/Z(0)=1+u(z)$ and use
$\log(1+u)=\sum_{m\ge1}(-1)^{m-1}u^m/m$ near zero. The coefficient of
$z_1\cdots z_n$ in $u^m$ sums over ordered partitions into $m$ nonempty
blocks. Dividing by $m$ turns their $m!$ orderings into $(m-1)!$.
Conversely, expanding $e^{W(z)-W(0)}$ proves the moment--cumulant relation

\begin{gather*}
\FitMath{\mathbb E\prod_iO_i=\sum_{\pi\in\operatorname{Part}(n)}
 \prod_{J\in\pi}\kappa((O_i)_{i\in J}).}
\end{gather*}


Replacing $O_i$ by $O_i-c_i$ changes the logarithm only by
$-\sum_i z_ic_i$. Thus every cumulant of order at least two is unchanged.
If $c_i=\mathbb E O_i$, every singleton cumulant vanishes. The centered
moment functional is consequently

\begin{gather*}
\FitMath{\mathbb E\prod_i(O_i-\mathbb E O_i)
 =\sum_{\pi\in\operatorname{Part}_{\ge2}(n)}
 \prod_{J\in\pi}\kappa((O_i)_{i\in J}).}
\end{gather*}

There are finitely many partitions at fixed $n$. Convergence of every
separated cumulant of order at least two therefore implies convergence of
the centered moments on the same tuples. Positivity must be applied to
these moments, not to the cumulants separately.

The source is exactly a local coupling shift:

\begin{gather*}
\FitMath{-\beta A+\sum_i z_iO_i=-\sum_p(\beta-w(x_p))a_p,} \\[3pt]
\FitMath{w=\sum_i z_if_i.}
\end{gather*}

For a constant source $c$, $Z_\beta(c+w)=Z_{\beta-c}(w)$. This algebraic
identity holds before any renormalization procedure is chosen.

\subsubsection{What a mass-preserving transformation preserves}

Let $T$ be a fixed linear map on finite measures on a measurable space,
with $(T\nu)(1)=\nu(1)$ for all measures in its domain. If
$\nu_z=e^{\sum_i z_iO_i}\nu$ belongs to that domain, then

\begin{gather*}
\FitMath{(T\nu_z)(1)=\nu_z(1).}
\end{gather*}

This proves exact equality of sourced partition functions, and hence of
their derivatives when the analytic dependence can be justified. It does
not identify $T\nu$ as the distribution of a predetermined block variable.
For that stronger conclusion one needs a measurable map $M$ and the
identity $(T\nu)(F)=\nu(F\circ M)$ for every bounded measurable $F$.

For example, on the two-point set $\{0,1\}$ let
$T\nu=\nu(\{0,1\})\delta_0$. This positive linear map preserves total mass
but changes the expectation of the coordinate from $\nu(\{1\})$ to zero
when $\nu$ is a probability measure. Thus mass preservation alone contains
no statement about observable expectations. Sources must be transported
through the exact integral, with the choices of the transformation fixed
throughout the source family.

\subsubsection{A fixed-spacing infinite-volume limit}

Consider the countable cubic lattice at a fixed positive spacing. Its
configuration space is $\mathcal V_\infty=G^{B_+}$ with the product
topology. This is a compact metrizable space: enumerate the bonds and use,
for example, the sum of $2^{-j}$ times bounded compatible distances on
their group coordinates. The space of Borel probability measures is
sequentially compact in the weak topology.

For a finite set of bonds $D$, define the conditional specification

\begin{gather*}
\FitMath{(\gamma_DF)(U)=
 \frac{\int F(V_D,U_{D^c})\exp[-\beta A_D(V_D,U_{D^c})]\,dV_D}
 {\int\exp[-\beta A_D(V_D,U_{D^c})]\,dV_D},} \\[3pt]
\FitMath{A_D=\sum_{p:p\cap D\ne\varnothing}a_p.}
\end{gather*}

The sum is finite and the denominator is bounded below by
$e^{-2\beta N(D)}$, where $N(D)$ is the number of plaquettes meeting $D$.
For continuous cylinder $F$, numerator and denominator depend continuously
on only finitely many boundary bonds; dominated convergence proves that
$\gamma_DF$ is another continuous cylinder function.

\begin{proposition}[DLR existence at fixed lattice parameters]
Take torus Wilson measures at fixed spacing and fixed $\beta$, and extend
their bond configurations periodically to the infinite lattice. Every
sequence of torus side lengths tending to infinity has a weakly convergent
subsequence. Each limit $\nu$ satisfies $\nu(\gamma_DF)=\nu(F)$ for every
finite $D$ and bounded Borel cylinder function $F$. It is invariant under
all lattice translations and all lattice isometries preserved by the
finite torus measures.
\end{proposition}
\begin{proof}
Compactness gives the subsequence. Fix $D$ and a continuous cylinder $F$.
For sufficiently large tori the finite set consisting of $D$, all adjacent
plaquettes and the support of $F$ embeds without identifications. The
finite-volume conditional expectation identity then reads
$\mu_\ell(\gamma_DF)=\mu_\ell(F)$ after periodic extension. Both integrands
are continuous cylinder functions, so the identity passes to the weak
limit. Continuous functions on the finite-dimensional compact coordinate
spaces determine measures. The identity therefore extends first to bounded
Borel cylinder functions, and, if desired, to all bounded Borel functions
by the monotone class theorem. Invariance passes to the limit by the same
continuous-cylinder argument, since a fixed lattice isometry acts
continuously on the product space.
\end{proof}

This proposition does not change the spacing or the coupling. It provides
no bounds uniform as the spacing tends to zero and gives no uniqueness or
extremality of the limiting state. A theorem about all periodic Gibbs
states is stronger than a theorem about these torus limit points.

\subsubsection{Compactness of separated correlation forms}

Let $\mathbb E$ be $\mathbb R^4$ or a fixed torus. A separated tuple consists
of finitely many compactly supported $C^1$ functions with pairwise disjoint
supports. Their compactness makes the minimum distance between nonempty
supports positive. Fix the norm
$\|f\|_1=\max(\|f\|_\infty,\|\nabla f\|_\infty)$.

\begin{proposition}[A diagonal extraction principle]
Suppose $T_{n,j}$ are symmetric multilinear forms on separated tuples and
for every fixed $n$, $R<\infty$ and $d>0$ there is a finite $C(n,R,d)$,
independent of all sufficiently large $j$, such that

\begin{gather*}
\FitMath{|T_{n,j}(f_1,\ldots,f_n)|\le C(n,R,d)\prod_i\|f_i\|_1}
\end{gather*}

whenever the supports lie in $B(0,R)$ and are separated by at least $d$.
Then a subsequence converges at every separated tuple, simultaneously for
all $n$. The limits obey the same bounds.
\end{proposition}
\begin{proof}
Choose all finite families of finite unions of rational closed boxes with positive mutual
distance. The $C^1$ functions supported in each union form a separable normed
space, being a subspace of a finite product of spaces of continuous
functions on a compact set. Choose countable dense subsets. There are
countably many finite tuples drawn from these subsets. Each sequence of
values is eventually bounded by the assumed estimate. Bolzano--Weierstrass
and a successive diagonal extraction give a common subsequence on which
all these values converge.

Given any separated tuple, enclose its supports in pairwise separated
finite unions of rational boxes, still inside a fixed bounded region.
Approximate each function in $C^1$ by elements of the chosen dense sets.
The telescoping multilinear identity

\begin{gather*}
\FitMath{T(f_1,\ldots,f_n)-T(g_1,\ldots,g_n)
 =\sum_{i=1}^n T(g_1,\ldots,g_{i-1},f_i-g_i,f_{i+1},\ldots,f_n)}
\end{gather*}

and the uniform bound make the approximation error uniformly small along
the extracted subsequence. This proves the Cauchy property for the
original tuple. Multilinearity, symmetry and the bound pass to the limit
on each product of spaces with separated supports.
\end{proof}

The difficult Yang--Mills input is the hypothesis supplying $C(n,R,d)$.
Finite-cutoff boundedness with a constant involving the number of
plaquettes does not meet it. An exponentially large dependence on the
physical torus volume also does not supply a bound uniform on
$\mathbb R^4$. These are different compactness problems.

\subsubsection{Reflection positivity and closed limits}

Let $\theta$ be a geometric reflection, and let $\overline X$ denote
componentwise complex conjugation of a test tuple. Real reflection-invariant centering preserves
reflection positivity: replacing each $O_j(f)$ by $O_j(f)-m_j(f)$ maps every
positive-half polynomial to another such polynomial and commutes with reflection
when $m_j(\theta f)=m_j(f)$. Suppose for each cutoff a centered moment
functional $\sigma_j$ satisfies

\begin{gather*}
\FitMath{\sum_{r,s=1}^m\overline{c_r}c_s\,
 \sigma_j(\theta\overline X_r\cdot X_s)\ge0}
\end{gather*}

for every finite collection of tuples $X_r$ in the positive half-space.
Within each $X_r$ the supports are separated; the reflected negative tuple
and the positive tuple are automatically separated from one another.
Overlap between supports from $X_r$ and $X_s$ on the same side causes no
problem, because these supports do not occur together in the same
concatenated pair $\theta X_r\cdot X_s$.

If all the entries converge, the finite sum converges and remains
nonnegative. Thus the limiting moment functional is reflection positive
on this vector space of separated monomials. This elementary closedness statement is the
precise sense in which positivity passes to limits. Passing from that
domain to a Schwartz test space requires a density result and a continuous
bound for the limiting distributions; it cannot be justified by closedness
alone.

The algebraic form $\langle F,G\rangle=\sigma(\theta\overline F\cdot G)$
is positive semidefinite. Positivity of
$\langle F+zG,F+zG\rangle$ for every complex $z$ gives
$|\langle F,G\rangle|^2\le\langle F,F\rangle\langle G,G\rangle$.
Consequently a null vector is orthogonal to every vector. The quotient by
the null space is an inner-product space, and its completion is a Hilbert
space. Defining translations on this quotient additionally requires that
they preserve the null space and satisfy the appropriate bounds. Strong
continuity and a nonnegative generator are later conclusions with their
own hypotheses.

\subsubsection{Two limits of what scalar correlations prove}

Let $X$ be a real standard Gaussian variable and put $Y=X^2-1$. For $t<1/2$
near zero,

\begin{gather*}
\FitMath{\log\mathbb E e^{tY}=-t-\tfrac12\log(1-2t)
  =\sum_{n\ge2}\frac{2^{n-1}}n t^n.}
\end{gather*}

It follows that $\kappa_n(Y)=2^{n-1}(n-1)!$ for every $n\ge2$.
In particular a nonzero connected four-point function of a quadratic
observable is compatible with a Gaussian underlying theory. For distinct
centered jointly Gaussian variables the same conclusion is expressed by
the connected cycle contractions of their Wick squares.

Accordingly a non-Gaussian $\operatorname{tr}(F^2)$ field is an important property
of that field, but is not by itself a characterization of an interacting
Yang--Mills gauge theory. Reconstruction from the $\operatorname{tr}(F^2)$
correlations produces the cyclic sector generated by that observable.
Identifying a richer gauge-invariant algebra requires additional
observables and compatibility arguments.

A polynomial cluster estimate likewise gives no positive mass gap. In a
Hilbert space with a nonnegative self-adjoint $H$, a connected time
correlation has the spectral form
$\langle v,e^{-tH}v\rangle=\int_0^\infty e^{-tE}\,d\mu_v(E)$ for a vector
orthogonal to the vacuum. If $d\mu_v(E)=E e^{-E}\,dE$, the integral is
$(1+t)^{-2}$. It decays, yet its spectral measure meets every interval
$(0,\varepsilon)$. The support condition required for a gap would instead
exclude all such intervals below one fixed positive threshold.

These examples delimit the conclusions of the conditional construction.
They do not contradict the proposed ultraviolet estimates, and they do
not replace the work of proving those estimates.


% END MODULE chapters/foundations.tex

% BEGIN MODULE chapters/rg_objects_v2.tex
\subsection{Renormalisation objects and norms}
\label{ch:rg-objects-v2}

The renormalisation step must represent the same integral after one
scale is integrated, with local interactions whose bounds remain uniform
in the number of scales. We describe those interactions by their local
functions and measures, their field domains, and their norms. This gives
a precise meaning to preservation of the class under a step. Constructing
backgrounds and keeping every substituted field inside its stated domain
require analytic estimates beyond these definitions.

\subsubsection{Lattices, currents, and auxiliary geometric data}

\phantomsection\label{notdef:gauge-convention}
The normalization in this section is that of $G=\operatorname{SU}(2)$.
Its Lie algebra is represented by traceless Hermitian matrices, with
$\langle A,B\rangle=\tfrac12\operatorname{tr}(AB)$ and complexification
$\mathfrak g^{\mathbb C}$. Fix an odd integer $L>1$. The level-$j$
lattice, denoted $\mathbb T^{(j)}$, consists of the sites retained after
$j$ blockings of a periodic bare lattice. Its bare spacing, expressed in
level-$j$ units, is $\eta_j=L^{-j}$. A bond configuration satisfies
$U(\bar b)=U(b)^{-1}$. Products along oriented paths are ordered in path
order. Positive coordinate bonds suffice to specify a configuration.
All suprema over bonds below use these bonds. In the left multiplication
chart, reversal is $J(\bar b)=-\operatorname{Ad}(U_b^{-1})J(b)$.
The complex configuration group is
$\operatorname{SL}(2,\mathbb C)$.

For a positively oriented plaquette $p$, based at its chosen corner,
write $U_p$ for its holonomy and set

\begin{gather*}
\FitMath{a_p(U)=1-\tfrac12\operatorname{tr}U_p,} \\[3pt]
\FitMath{A(U)=\sum_pa_p(U),} \\[3pt]
\FitMath{A(c,U)=\sum_pc(x(p))a_p(U).}
\end{gather*}

Here $x(p)$ is the plaquette midpoint. On real configurations this is
the real Wilson action; the same formula is its holomorphic extension.
For definiteness the current on a lattice of spacing $\eta$ is specified
by its variational normalization:

\refstepcounter{equation}
\begin{gather*}
\FitMath{\left.\frac{d}{dt}A(e^{i\eta tB}U)\right|_{t=0}
   =\eta^4\sum_b\langle J_\eta(U)(b),B(b)\rangle.}\tag{\theequation}\label{eq:rg-objects-R1}
\end{gather*}

The sum is over positive bonds and $B$ is arbitrary. This identity
defines every component of $J_\eta$, also for complex $U$ by analytic
continuation. It fixes the transports and factors implicit in the
notation $D_U^*F_U$. One must use the plaquette divergence associated
with \eqref{eq:rg-objects-R1} in that notation; the adjoint of the site gradient alone has
a different domain. At $U=\mathbf1$, differentiation in the potential coordinate
$U(t)=e^{i\eta tB}$ gives
$J'_\eta(\mathbf1)B=\mathrm{d}_\eta^*\mathrm{d}_\eta B$, where $\mathrm{d}_\eta$ here is the
oriented plaquette difference on one-forms divided by $\eta$.
If the numerical mesh changes from $\eta$ to $\eta'=L^m\eta$ while the group
configuration stays fixed, potentials are multiplied by $L^{-m}$ and currents
by $L^{-3m}$. This follows directly from \eqref{eq:rg-objects-R1} and the chart $e^{i\eta A}$,
and specifies the direction of the unit change.
The current variable of a complex pair, introduced below, is independent
of this particular function of $U$.

\phantomsection\label{notdef:localization-cubes}
Fix powers of $L$ called $\mathfrak M,M_1',M$, a gap
$\rho=R_1\mathfrak M$, and constants $C_{\rm box},C_{\rm pot}$.
They are part of the auxiliary data, selected before domains or norms
are formed. The first two sizes concern backgrounds and walks; $M$
is the localization cube size. They are not interchangeable. The
partition $\pi_j$ consists of the centered, half-open $M$-cubes in
level-$j$ units. Boundary conventions are fixed once on every lattice.
For a cubical region $Y$, $Y^{\sim n}$ means addition of $n$ layers
in the partition currently named; $Y^{\sim-n}$ means removal of those
layers. A hull is the smallest union of cubes of the named partition
containing the set. These operations include their periodic versions.

The two erosion assignments are fixed as follows. Take $L\ge13$
odd, $R_1\ge L$, $M_1'\ge R_1\mathfrak M$ and
$M\ge L^4M_1'$, in addition to the background floors. For
$Z\in\mathcal D_{p-1}$ put $\Omega_0^{(1)}=\Omega_0^{(2)}=Z$.
In bare units their successive members are
\begin{equation}\label{htc:erosion-rules}
\begin{aligned}
 \Omega_{q+1}^{(1)}
 &=\bigcup_{\substack{C\text{ of side }\mathfrak M L^{q+1}\\
 C\subset\Omega_q^{(1)},\
 \operatorname{dist}_\infty(C,(\Omega_q^{(1)})^c)>R_1\mathfrak M L^q}}C,\\
 \Omega_{q+1}^{(2)}
 &=\bigcup_{\substack{C\text{ of side }M_1'L^{q+1}\\
 C\subset\Omega_q^{(2)},\
 \operatorname{dist}_\infty(C,(\Omega_q^{(2)})^c)\ge M_1'L^{q+1}}}C,
 \qquad 0\le q<p.
\end{aligned}
\end{equation}
Each union contains every cube of the named aligned partition
satisfying its test. A tower of depth $n\le p$ is the truncation
of this construction at $n$. These definitions commute with the
allowed translations and common lifts. Empty members remain empty;
for a full torus the complement distance is infinite and the
members remain full. The maximal-cube rule is essential: it gives
the upper bound on total erosion proved in \cref{htc:geometry}.

Here is the determining set of a nested sequence. Work temporarily in
level-$k$ units, so the spacing of level $j$ is $\xi_j=L^{j-k}$.
Let $\Omega_0\supset\cdots\supset\Omega_k$ and
$\Omega_{k+1}=\varnothing$. The set $\Gamma_j$ consists of bonds of
$\mathbb T^{(j)}$ whose endpoints lie in $\Omega_0$, at least one
endpoint lies in $\Omega_j$, and neither lies in $\Omega_{j+1}$.
Thus inward bonds crossing into the next layer are excluded. Put

\refstepcounter{equation}
\begin{gather*}
\FitMath{\mathbf B=\bigcup_j\Gamma_j,} \\[3pt]
\FitMath{M_{\mathbf B}(U)
       =(\mathcal M^jU|_{\Gamma_j})_j.}\tag{\theequation}\label{eq:rg-objects-R2}
\end{gather*}

The endpoints of these bonds are the nodes. They are the tops of the
blocks, or pillars, partitioning the determining region. A residual
gauge function equals $\mathbf1$ at all nodes. In particular, imposing
additional inward-bond constraints changes this problem.

The averaging map $\mathcal M$ used in \eqref{eq:rg-objects-R2} is the symmetric contour
average. For group elements $V_i$ of small diameter their mean is the
unique nearby $V$ satisfying $\sum_i\log(V_iV^{-1})=0$. The contour
transport $U(y,x)$ is the mean over the shortest paths from $y$ to $x$
in its centered block that traverse the coordinate directions in one
of their orders. Set $U(x,y)=U(y,x)^{-1}$. For a coarse bond $c$,
$x'=x+(c_+-c_-)$, and straight transports $U[c]$ and $U([x,x'])$,

\refstepcounter{equation}
\begin{gather*}
\FitMath{\mathcal M(U)(c)=
 \exp\!\left(L^{-4}\sum_{x\in B(c_-)}
 \log\{U(c_-,x)U([x,x'])U(x',c_+)U[c]^{-1}\}\right)U[c].}\tag{\theequation}\label{eq:rg-objects-R3}
\end{gather*}

All logarithms in regular neighborhoods are their near-identity
branches. Away from those neighborhoods a measurable extension of
the contour means and logarithms is auxiliary input. Exact identities
using that extension require its gauge covariance; small-field
analyticity concerns only the displayed branches. Existence of a
globally symmetric extension with every desired property is not
asserted merely by writing \eqref{eq:rg-objects-R3}.

\phantomsection\label{notdef:background-orbit}
For every determining sequence and real regular datum $\mathbf V$,
the background input is a regular critical orbit for
$\sum_{p\subset\Omega_0}a_p$ on
$\{U:M_{\mathbf B}(U)=\mathbf V\}$, together with a chosen
analytic continuation $\Phi_{\mathbf V}(B)$ at datum
$e^{iB}\mathbf V$. The representatives, gauge, analytic domain and
transport under local charts are included in this input. Uniqueness,
existence and uniform estimates for these maps are separate background
theorems. The word ``background'' does not substitute for them.

\subsubsection{Hereditary domains and their shifts}

A triple of level $j$ on $X$ is defined on the bare mesh $\eta_j$
when lengths are expressed in level-$j$ units. It is
$P=(U,A',\mathbf J)$, with $U$ real and $A',\mathbf J$ complex
Lie-algebra-valued bond forms. Its pair is
$\mathbf p(P)=(\mathbf U,\mathbf J)$, where
$\mathbf U=e^{i\eta_jA'}U$. A triple is a splitting of its pair.
Fix positive radii $a_0,a_1,a_2$. The simple conditions are

\refstepcounter{equation}
\begin{gather*}
\FitMath{\begin{aligned}
 &|\partial U-\mathbf1|,|\partial\mathbf U-\mathbf1|
       <a_0\eta_j^2,\qquad
   |A'|,|\nabla_UA'|<a_1,\qquad |\mathbf J|<a_2;\\
 &|A|,|\nabla A|<C_{\rm pot}Ma_0
       \quad\text{in the real gauges described below.}
 \end{aligned}}\tag{\theequation}\label{eq:rg-objects-R4}
\end{gather*}

For each cubical box contained in $X$ with side at most
$C_{\rm box}LM$, such a real gauge must exist, with
$U=e^{i\eta_jA}$. The norms are maximum matrix norms in the named local units.
For a one-form component $A_\mu$ based at its initial site, fix the forward convention

\begin{gather*}
\FitMath{(\nabla_\nu^U A)_\mu(x)=\eta_j^{-1}
 \{\operatorname{Ad}(U_\nu(x))A_\mu(x+\eta_je_\nu)-A_\mu(x)\}.}
\end{gather*}

The maximum of these component norms defines $|\nabla_U A|$.
In the usual two-radius notation $a_2=a_0$.

For every subdomain $Z\subset X$ at every level $p\le j$, for both
erosion assignments, and every truncation depth $1\le n\le p$, form
its determining set $\mathbf B'$. Define its tower configuration by

\refstepcounter{equation}
\begin{gather*}
\FitMath{U_n(Z;P)=\Phi_{M_{\mathbf B'}(U)}
 \left(\frac1i\log\{M_{\mathbf B'}(\mathbf U)
                           M_{\mathbf B'}(U)^{-1}\}\right).}\tag{\theequation}\label{eq:rg-objects-R5}
\end{gather*}

This expression is required to exist in the chosen continuation
domain. On $Z$ with two layers of the $M$-cubes of level $p-1$
removed, require

\refstepcounter{equation}
\begin{gather*}
\FitMath{\begin{aligned}
 |\partial U_n(Z;P)-\mathbf1|
      &<a_0L^{-(j-p)}\eta_p^2,\\
 |J(U_n(Z;P))|
      &<a_0L^{-(j-p)}L^{2(n-p)}.
 \end{aligned}}\tag{\theequation}\label{eq:rg-objects-R6}
\end{gather*}

The current in the second line is formed in level-$n$ units.
Impose the same conditions with $\mathbf U$ replaced by $U$.
These are quantifiers over every allowed $Z,p,n$, not just the domain
on which an activity is eventually evaluated. The hereditary core
$\mathcal U_j^\circ(X;a_0,a_1,a_2)$ is the set of triples satisfying
\eqref{eq:rg-objects-R4}--\eqref{eq:rg-objects-R6}. Its pair image is the analytic pair domain; holomorphy
means holomorphy in $(A',\mathbf J)$ on every open slice with $U$
fixed. The inequalities are strict. Requiring the tower maps to exist
in their open continuation domains is part of membership.

An activity on this core must have the tie: its values at any two
splittings of the same complex pair agree. Gauge invariance means
invariance under simultaneous transformation of both entries of the
pair; the current at a bond is based at its initial site and transforms
by conjugation there. The core itself is not saturated by complex
gauge transformations. Extension of a function to the union of complex
orbits requires equality on all intersections and holomorphic descent;
it is an additional statement when used.

Restriction $\rho_X$ discards bonds and variables outside $X$ without
extending them across the boundary. If an evaluation changes its numerical
mesh from $\eta_j$ to $\eta'=L^m\eta_j$, the coordinate change is explicitly

\refstepcounter{equation}
\begin{gather*}
\FitMath{(U,A',\mathbf J)\longmapsto
 (U,L^{-m}A',L^{-3m}\mathbf J),} \\[3pt]
\FitMath{m\ge0.}\tag{\theequation}\label{eq:rg-objects-R7}
\end{gather*}

For a shift $(\mathbf H,\mathsf j)$ on the same lattice define

\refstepcounter{equation}
\begin{gather*}
\FitMath{P\dotplus(\mathbf H,\mathsf j)
  =(U,A'\boxplus\mathbf H,\mathbf J+\mathsf j),} \\[3pt]
\FitMath{e^{i\eta_j(A'\boxplus\mathbf H)}
       =e^{i\eta_j\mathbf H}e^{i\eta_jA'}.}\tag{\theequation}\label{eq:rg-objects-R8}
\end{gather*}

The logarithm defining $\boxplus$ is the small branch; inclusion in
its domain is required. The independent current is shifted additively.
Replacing it by $J(e^{i\eta_j\mathbf H}\mathbf U)$ at every complex
argument would change the construction.

The decoupled response on a cubical region $Y$ is the pair of maps
$\mathfrak h_Y(s;P,B')=(\mathbf H,\mathsf j)$, holomorphic in
complex cube parameters $|s(\Delta)|<e^{\bar\kappa}$, the complex
triple variables, and the datum $B'$. It vanishes at $B'=0$.
Its first component is the solution of the actual background equations
with the propagators decoupled by these parameters:

\refstepcounter{equation}
\begin{gather*}
\FitMath{\begin{gathered}
 \mathbf H=\mathbf a-\mathsf H(s)\mathsf D,
 \qquad \mathsf D=C_k(\mathbf a-\mathsf H(s)\mathsf D),\\
 \mathbf a=\mathsf H_1(s)B'+\mathbf A_1,
 \qquad \mathbf A_1+\mathfrak G(s)\mathsf W_s(\mathbf a)=0.
 \end{gathered}}\tag{\theequation}\label{eq:rg-objects-R9}
\end{gather*}

Here $C_k$ is the nonlinear part of the straightened constraint map;
$\mathsf H,\mathsf H_1$ are the constrained linear minimizers;
$\mathfrak G$ is the inverse on the simultaneous Landau and
zero-average slice; and $\mathsf W_s$ is the nonlinear gradient
remainder. The actual raw-coordinate force and complete
carrier prescription are given in \cref{dnr:linear-families,dnr:parameter-families}
and \eqref{dnr:force}. Their first-order regular-pair domain,
independent current radius, and parameter discs are part of that construction.
The reduced current is the multiplier expression
\eqref{dnr:multiplier}, with the uniform bound
\eqref{dnr:current-bound}. At one change of units its value is
$L^{-3}\mathbf J+\mathsf j$.

The response has realness, gauge covariance,
compatibility under restriction to wall-components, and transport by
common lifts. A wall-component is a component of cubes sharing
three-dimensional faces. When $s$ vanishes on separating cubes, the
equations split over the remaining components, and the response on a
component with no datum is zero. In local first-order norms its size
is at most $C e^{C\bar\kappa}\sup|B'|$; a response viewed $k-j$
levels below has the additional potential factor $L^{-(k-j)}$.
At the whole-sequence all-one point over a physical critical
real reference, its shifted pair
is the corresponding background and its own current, with the
recorded node gauge factor, by \cref{dnr:physical}.
Existence, uniqueness and the stated properties are proved in
\cref{dnr:response,dnr:walls}.

For a term $F$ on $X$ let $X_0$ be its hull in the cubes of the
decoupling scale. The interpolation is

\refstepcounter{equation}
\begin{gather*}
\FitMath{\Phi_Y[F](s,\nu)=F\big((\rho_XP)
                          \dotplus\nu\mathfrak h_Y(s;P,B')\big).}\tag{\theequation}\label{eq:rg-objects-R10}
\end{gather*}

Every occurrence of \eqref{eq:rg-objects-R10} requires a domain-inclusion hypothesis:
for all the indicated $s$ and $|\nu|<\varrho_F$, its triple lies in
the core of $F$. One sufficient quantitative hypothesis, with a fixed
constant $c_F>0$, is inclusion throughout

\refstepcounter{equation}
\begin{gather*}
\FitMath{|\nu|<c_F L^{k-j}a_{0,j}\delta(\theta_k)^{-1}
                   e^{-C_F\bar\kappa}=:\varrho_F>1.}\tag{\theequation}\label{eq:rg-objects-R11}
\end{gather*}

In inner and removed substitutions separate inclusion bounds and
their own radii are required. \eqref{eq:rg-objects-R11} is not a consequence of
the formal shift law in \eqref{eq:rg-objects-R8}.

For disjoint finite cube sets $\mathcal K,\mathcal N$ define

\refstepcounter{equation}
\begin{gather*}
\FitMath{F[\mathcal K\mid\mathcal N]
 =\sum_{\mathcal K'\subset\mathcal K}
   (-1)^{|\mathcal K|-|\mathcal K'|}
      F(\mathbf1_{\mathcal K'\cup\mathcal N}).}\tag{\theequation}\label{eq:rg-objects-R12}
\end{gather*}

This equals the integral of $\prod_{\Delta\in\mathcal K}
\partial_{s(\Delta)}F(s+\mathbf1_{\mathcal N})$ over the unit cube,
by iterating the one-variable fundamental theorem of calculus.
Summing \eqref{eq:rg-objects-R12} over subsets gives its exact resummation identity.
For the rooting assertion, use the wall-components of
$\mathcal K\cup\mathcal N\cup X_0$. If such a component contains a
cube of the differenced set $\mathcal K$ and misses the base $X_0$,
then locality and splitting make the function independent of that component's
parameters, so the alternating bracket is zero. An additional component made
only of fixed $\mathcal N$ cubes need not make a bracket zero. For
$\Phi_Y[F](s,\nu)-\Phi_Y[F](s,0)$ the expression vanishes if the
component of the base has no datum for the shifted response, since that
response is then zero on the base. These conclusions require the stated
splitting and support properties of the actual response. If $|\Phi_Y[F]|\le M_F$ on the stated polydisc,
Schwarz's lemma gives $|\Phi_Y[F](s,1)-\Phi_Y[F](s,0)|
\le2M_F/\varrho_F$. Cauchy's formula on the $s$ discs then bounds
the bracket of this excess by
$2M_F\varrho_F^{-1}(e^{\bar\kappa}-1)^{-|\mathcal K|}$.
For $\bar\kappa\ge1$ this is at most
$2M_F\varrho_F^{-1}e^{-(\bar\kappa-1)|\mathcal K|}$.
When $\mathcal K$ is nonempty, this also bounds the bracket of
$\Phi_Y[F](s,1)$ itself: its value at $\nu=0$ is the $s$-independent
constant $F(\rho_XP)$, which the bracket annihilates. For an empty differenced
set, an un-subtracted constant does not acquire $\varrho_F^{-1}$.
Thus the algebraic bracket bound is proved once the response and
inclusion hypotheses have been supplied.

\subsubsection{Pointed families, marginal coefficients, and list norms}

\phantomsection\label{notdef:tree-length}
Let $\mathcal D_j$ be the connected unions of cubes of $\pi_j$.
Set $d_j(X)=M^{-1}$ times the length of a shortest tree in $X$
meeting each cube; it is zero for one cube. A pointed family is the
collection $E(X,\mathbf p,z)$ for $X\in\mathcal D_j$ and unit sites
$z\in X\cap\mathbb T^{(j)}$, with each term analytic on its core,
gauge invariant and tied. Its norm is

\refstepcounter{equation}
\begin{gather*}
\FitMath{\|E\|_{j,\kappa}=
 \sup_{\text{admissible tori and their companion}}
 \sup_{X,z\in X} e^{\kappa d_j(X)}
 \sup_{P\in\mathcal U_j^\circ(X)}|E(X,\mathbf p(P),z)|.}\tag{\theequation}\label{eq:rg-objects-R13}
\end{gather*}

The supremum in $z$ is not a sum over base points. For an unpointed
family omit $z$ from this formula. Sums over domains require the
separate lattice-animal summability bound at the chosen decay rate.

A companion is a single family $\widehat E$ on $\mathbb Z^4$,
agreeing with every torus term on every common lift of a non-wrapping
domain with the necessary collar. For a fixed level, fixed frame prefix and fixed initial weight, the same
companion serves every cutoff on which those input data coincide. Its
existence is compatibility between tori, not a property
that can be tested on just one torus. Individual terms are required
to be covariant under the translations preserving $\pi_j$.
The pointed total has the stronger property

\refstepcounter{equation}
\begin{gather*}
\FitMath{\mathcal P_z^{\mathsf T}(W)=
    \sum_{X\ni z}E(X,\mathbf p_W,z),} \\[3pt]
\FitMath{\mathcal P_{rz}^{\mathsf T}(rW)=\mathcal P_z^{\mathsf T}(W)}\tag{\theequation}\label{eq:rg-objects-R14}
\end{gather*}

for every lattice isometry $r$, including unit translations. Here

\begin{gather*}
\FitMath{\mathbf p_W=(U_j(W),J(U_j(W))).}
\end{gather*}

The covariance class consists
of the $W$ for which these backgrounds exist, every restriction lies
in the required core, and

\refstepcounter{equation}
\begin{gather*}
\FitMath{L^{2j}\max_p|\partial U_j(W)(p)-\mathbf1|
       <c\min\{\varepsilon(\theta_{j-1}),
                         \varepsilon(\theta_j)\}.}\tag{\theequation}\label{eq:rg-objects-R15}
\end{gather*}

Both $W$ and $rW$ must be in this class. Invariance of the sum over
$z$ alone is weaker than \eqref{eq:rg-objects-R14}.

The site tent is
$h_z(x)=\prod_{\mu=1}^4\max\{1-|x_\mu-z_\mu|,0\}$.
Evaluate it at bond and plaquette midpoints. Then
$\sum_zh_z(x)=1$ and $h_{rz}(rx)=h_z(x)$. A bond-rooted term is
pointed by $E(X,\cdot,z)=\sum_bh_z(x_b)E(X,b;\cdot)$; its collar
must contain the supports of these tents. The pointed Wilson family
has terms $A(h_z,\mathbf U)$ on the cube of $z$ enlarged by one
layer, and zero elsewhere. This convention distributes a coupling
counterterm by the same base points as the regular family.

To define the coefficient, use the chart
$\mathsf e_X(A,J)=(e^{i\eta_jA},J)$ at $(\mathbf1,0)$.
Directions are pairs of complex bond forms, with norm
$\max\{\sup|A|,\sup|\nabla A|,\sup|J|\}$.
Let $\mathsf H_\infty$ be the derivative of the infinite-lattice
background map at its trivial datum, in its fixed gauge, and put

\begin{gather*}
\FitMath{\mathfrak h_\infty B=
       (\mathsf H_\infty B,\mathrm{d}^*\mathrm{d}\mathsf H_\infty B).}
\end{gather*}

This linear response is an input from the background theory; it must
include its domain and decay estimates. Choose an orthonormal color
basis $t_1,t_2,t_3$, and let $\delta_c^a$ be $t_a$ on the unit bond
$c$ and zero elsewhere. With Hessians taken in the displayed chart,

\refstepcounter{equation}
\begin{gather*}
\FitMath{\begin{aligned}
 \Pi^E(c,c';z)&=\tfrac13\sum_{a=1}^3\sum_{\widehat X\ni z}
   \operatorname{Hess}\widehat E(\widehat X,\cdot,z)
       [\mathfrak h_\infty\delta_c^a,
        \mathfrak h_\infty\delta_{c'}^a],\\
 \Pi^E(c,c')&=\sum_z\Pi^E(c,c';z)
                    =:\Pi^E_{\mu\nu}(c'_--c_-),\\
 b[\widehat E]&=\sum_{u\in\mathbb Z^4}
                   \Pi^E_{\mu\nu}(u)u_\mu u_\nu,
                    \qquad\mu\ne\nu.
 \end{aligned}}\tag{\theequation}\label{eq:rg-objects-R16}
\end{gather*}

The coefficient is defined only when all these sums converge
absolutely, the middle expression depends just on directions and
displacement, and the last value is independent of the distinct
directions chosen. Gauge independence, the Ward identities and the
reduction of the second-moment tensor to this scalar require proof.
They are not incorporated by declaring an arbitrary Hessian to have
a coefficient. In particular, a claim that the Wilson tent family has
coefficient one requires the linear-response identities with the
normalization in \eqref{eq:rg-objects-R1}.

A member is a triple $(E,\widehat E,e)$ satisfying \eqref{eq:rg-objects-R13}--\eqref{eq:rg-objects-R16},
with $e=b[\widehat E]$. Activity bounds and coefficient bounds use their
declared physical radii, including the inverse-square readout factor of
\cref{ch:extraction-radius-normalization}. A list at step
$k+1$ contains such members for $j\le k$ and unpointed slot
families $H^{(j)}$ with no coupling counterterm. Its action is

\refstepcounter{equation}
\begin{gather*}
\FitMath{\begin{aligned}
 S[\mathfrak f](U)=&\sum_{j\le k}\left[
   \sum_z\sum_{X\ni z}
       \{E^{(j)}(X,U,z)-E^{(j)}(X,\mathbf1,z)\}
                         -e_jA(U)\right]\\
 &+\sum_{j,X}\{H^{(j)}(X,U)-H^{(j)}(X,\mathbf1)\}.
 \end{aligned}}\tag{\theequation}\label{eq:rg-objects-R17}
\end{gather*}

Every evaluation uses the background and units appropriate to its
level. Slot families have companions, term covariance under the
localization translations, and invariant totals on the stated
background class, but no requirement of a pointed coefficient.
For $0<\alpha_{\rm H}<1$ their list size is

\refstepcounter{equation}
\begin{gather*}
\FitMath{N_{\mathfrak a}(\mathfrak f)=
 \sum_{j\le k}(k+2-j)^{3/2}L^{-\alpha_{\rm H}(k-j)}
                      \|E^{(j)}\|_{j,\kappa}
       +\sum_j\varrho_j^2\|H^{(j)}\|_{j,\kappa},} \\[3pt]
\FitMath{\varrho_j=\max_\iota
       \frac{(1+\beta_{\rm s})a_{\iota,k+1}}{a_{\iota,j}}.}\tag{\theequation}\label{eq:rg-objects-R18}
\end{gather*}

All three radii, including the independent-current radius, enter this
maximum. Admissible radius systems require
$CMa_{0,j}<a_{1,j}$, $a_{1,j}^2\le Ca_{0,j}$,
$a_{\iota,j}\le(1+\beta_{\rm s})a_{\iota,k+1}
\le C(k+2-j)^{1/2}a_{\iota,j}$, and a fixed small upper bound on
the absolute radii. They also require
$\mathsf s_{\mathfrak a}:=\delta(\theta_k)/a_{0,k+1}\le c$.
The constants and the admissible systems are specified before taking
\eqref{eq:rg-objects-R18}. For the running system use
$a_{\iota,j}=\check\alpha_{\iota,j}
:=\min\{\alpha_{\iota,j-1},\alpha_{\iota,j}\}$ at $j\le k$ and
$a_{\iota,k+1}=\alpha_{\iota,k}$. For the absolute system use fixed
$\alpha_\iota$ at $j\le k$ and
$a_{\iota,k+1}=\alpha_\iota/(1+\beta_{\rm s})$.

\phantomsection\label{notdef:frames}
\subsubsection{Frames, histories, structures, and stored variables}

A frame is a real sequence $\theta_0,\ldots,\theta_k\ge\gamma^{-2}$,
with $0<\gamma<e^{-1/2}$, such that, for $i\le j$,

\refstepcounter{equation}
\begin{gather*}
\FitMath{\underline\lambda(j-i)-c_5\le\theta_i-\theta_j
                \le\bar B(j-i),} \\[3pt]
\FitMath{\check R_j=R(\min_{i\le j}\theta_i),} \\[3pt]
\FitMath{\check R_j\le L\check R_{j+1},}\tag{\theequation}\label{eq:rg-objects-R19}
\end{gather*}

where $R(y)$ is the smallest power of $L$ not below $(\log y)^r$.
Set $\bar g_j=\theta_j^{-1/2}$,
$p_0(g)=A_0(\log g^{-2})^{\mathfrak p_0}$,
$p_1(g)=(\log g^{-2})^{\mathfrak p_1}$,
$\varepsilon_j=\bar g_jp_0(\bar g_j)$,
$\delta_j=(A_1/A_0)\varepsilon_j$ and
$\alpha_{\iota,j}=C_\iota\theta_j^{-1/2}(\log\theta_j)^q$.
Here $\mathfrak p_0>\mathfrak p_1$, $A_1<A_0$, and all constants
are fixed. The waiting depth is $\mathsf N_j=\check R_j$.
The barred couplings belong only to the auxiliary frame. The actual running
couplings $x_j^{-1/2}$ are separate. The weights $x_j$ are separate entries, with
$x_0=\beta$ and $x_j=x_{j-1}-b_j$. A tube requires

\refstepcounter{equation}
\begin{gather*}
\FitMath{|\theta_j-x_j|\le\mathsf t_j,} \\[3pt]
\FitMath{\mathsf t_j\le c\theta_{j+1}
              (\log\theta_{j+1})^{-2q_{\rm t}},} \\[3pt]
\FitMath{q_{\rm t}\ge q,} \\[3pt]
\FitMath{2q_{\rm t}>q+\mathfrak p_0.}\tag{\theequation}\label{eq:rg-objects-R20}
\end{gather*}

Its radii are non-increasing. At a terminal frame with no next
member the radius must instead be supplied explicitly; for a
one-member frame require
$\mathsf t_0\le2c\theta_0(\log\theta_0)^{-2q_{\rm t}}$.
All tests, regions and kernels are formed from $\theta$, not from
the unknown entries $x_j$.

\phantomsection\label{notdef:histories}
A history is a finite record of every outcome of the tests of the
steps and of the operations inside large-field components, including
coordinate frames in the signed-permutation group $W_4$. The
formation data assign to each record its regions, stored bonds and
integration variables. In particular they give total tests
$\chi,1-\chi$ as Borel functions on the actual variable space,
with each value in $\{0,1\}$. A plaquette test is
$\mathbf1\{|\partial U_{j,\square}-\mathbf1|
<\varepsilon_j\eta_j^2\}$; a fluctuation test is
$\mathbf1\{|\mathcal A_j|<\delta_j\}$.
The local background map, the neighborhoods of each test, and the
collar and rim tests with their sharper thresholds are part of the
formation data. Listing the names of tests does not define them.

\phantomsection\label{notdef:history-regions}
The label is $\mathcal S=(\Omega_j,\Lambda_j,S_j)_{j\le k}$ with
$T_\ell\supset\Omega_1\supset\Lambda_1\supset\cdots
\supset\Lambda_k$; $S_j\subset\Omega_j\setminus\Lambda_j$
records strips and failed fluctuation regions. Write
$Z_j=\Lambda_j^c$. The step creates unions of
$M\check R_j$-cubes, with gaps at least $cM\check R_j$ in local
units. Nests inside an operation may have gaps of only one $M$-cube
of the later level. The formation data must state the actual
enlargements and test neighborhoods; bounds described by a fixed
number of layers leave that number as a required datum.

An operation at level $t$ acts on the prescribed class of components
which are small, quiet for $\mathsf N_t$ steps, and satisfy the
waiting and quarantine conditions. It first integrates stored fields
in $\mathsf N_t$ stages, then tests extension of the exterior field,
integrates the inner box in an axial gauge with frame in $W_4$,
restores block variables through normalized reference laws, localizes
the exponent, and performs its last test. A healed component leaves
the label; an island passes the earlier tests and fails the last;
a drop-out fails after $n\le\mathsf N_t$ stages and keeps those
stages. These alternatives and every operation kernel must be given
as measurable maps and positive measures. Their exactness is a
separate equality of represented measures.

An operation can use its waiting window only when
$t_\flat=t-\mathsf N_t\ge0$. Before then it is not eligible; no negative-level
history is introduced. A drop-out $C$ records its former
regions $\Omega_l^\circ\cap C$ and its stage boxes $Z''_l(C)$.
For each integer $0\le\nu(C)\le n(C)$ define

\refstepcounter{equation}
\begin{gather*}
\FitMath{\Omega_l^{\boldsymbol\nu}\cap C=
 \begin{cases}
 C\setminus Z''_l(C),&t_\flat<l\le t_\flat+\nu(C),\\
 \Omega_l^\circ\cap C,&t_\flat+\nu(C)<l
                                      \le t_\flat+n(C).
 \end{cases}}\tag{\theequation}\label{eq:rg-objects-R21}
\end{gather*}

Outside the affected levels retain the label regions. The choice
$\boldsymbol\nu=n$ is the determining sequence; it governs
backgrounds. The choice $\boldsymbol\nu=0$ is the action sequence;
it governs the sums and the coupling profiles. Standing marks record
the analogous $\Lambda_l,S_l$ data. Their waiting windows are
required to be disjoint at a point, and old marks disappear when
swallowed by a healed component or an island. These geometric
facts require verification for the actual operation.

A structure $\sigma=(\mathcal S,\boldsymbol\nu;t,\phi)$ also
records a phase. Its phase weight is $0$ at output, $1$ at the
step input, and $1-n/(2\mathsf N_t)$ after $n$ stages. A structure
precedes another if it counts no more completed stages and has
lower level, or the same level with larger phase weight. At an index
$0\le m\le t$ the radius coefficient is

\refstepcounter{equation}
\begin{gather*}
\FitMath{\begin{gathered}
 c_\sigma(x\mid m)=1-\beta_{\rm s}(1-2^{-(t-m)})
       +\varpi_\phi\mathsf g_t(m)
       -\beta_{\rm R}\sum_{l\in\operatorname{Re}_\sigma(x)}
                                      2^{-|m-l|},\\
 \mathsf g_t(m)=
 \begin{cases}\theta_{\rm s}\beta_{\rm s}2^{-(t-m)},&m<t,\\
 \beta_{\rm s},&m=t.
 \end{cases}
 \end{gathered}}\tag{\theequation}\label{eq:rg-objects-R22}
\end{gather*}

Here $0<\theta_{\rm s}<1$ and
$\beta_{\rm R}=c_{\rm R}\beta_{\rm s}$ are fixed fractions;
$\operatorname{Re}_\sigma(x)$ consists exactly of the levels of
the first $\nu(C)$ stages of each drop-out containing $x$.
The layer-radius formation data give the original regularity layer
and the layer after each stage. Their nonnegative difference is
$e_\sigma(x)$. The lift is $L_\flat^{2e_\sigma(x)}$ for a fixed
$L_\flat\ge1$. The layer-$m$ radii are multiplied by this lift and,
except for the real-potential gauge bounds, by \eqref{eq:rg-objects-R22}. Unfinished structures
here have $t\ge1$: at layer zero use level-one radii; at the unfinished
top use level-$t-1$ radii. The initial $t=0$ base structure is separate
input with its own specified radii. The data must ensure $c\le c_\sigma\le1+\beta_{\rm s}$.

The space $\widetilde U^c_\sigma(X)$ imposes the simple, tower
and small-field-core clauses with those layer radii, with one
splitting, on every subdomain and every preceding structure. The
layer assignments, erosion rules, continuation maps and recorded
remade levels are thus indispensable parts of its definition.
The claims that restriction or a stage substitution maps one such
space into another are analytic inclusion hypotheses. They require
the specified radii and constants; a formal partial order by itself
does not prove them.

\paragraph{Pointwise and local gauge clauses.}
For a structure with first member $\Omega_0$, let $\mathsf i(Y)$
consist of the bare bonds with both endpoints in $Y\cap\Omega_0$.
The pointwise part of the structured space imposes the real and
complex curvature bounds on every plaquette formed by these bonds.
The complex potential, its covariant differences on their retained
stencils and the independent current have their layer bounds on the
same carrier. In particular the pointwise bounds extend through
$\Omega_0\setminus\Omega_1$, with the layer-zero radii above.

In addition to the running gauge bounds, for each address
$(\sigma'',i,Z)$ of a preceding or current structure with $Z\subset Y$,
the structured space requires good real gauges on a fixed family
of cubes in two consecutive layers of that address. Each cube lies
in $Z\cap\Omega_0$ for the corresponding first member. In using
the strip response choose this fixed family to include all its gauge
cubes of side at most $12\mathfrak M\xi_n$ at index $n$, including
cubes contained in layer $n+1$. The gauge radius $a_\flat$ is below
the response ceiling. The width and radius are fixed before the
coupling ceiling, which is chosen for that family. In original bare
units the good gauge is $U^v=e^{i\mathsf v}$ with
$L^n|\mathsf v|,L^{2n}|\delta\mathsf v|\le a_\flat$.
The same condition in level units uses $\xi_n|A|$ and
$\xi_n^2|\nabla^\xi A|$, with $\mathsf v=\xi A$.
\label{rg:structured-local-gauge}

Use an address whose domain covers the full carrier of the response,
then restrict its gauges and stencils to the active union of parameter
cubes. When applying a substitution or taking a restriction, we must
verify that the resulting triple has these clauses on that address.
The layered bounds on $Y\cap\Omega_1$ alone give only that part
of this verification. In particular an independent current at a
layer-zero bond is unconstrained by those bounds, even if the
configuration and the tower backgrounds are flat.

Stored fluctuation variables are real bond forms
$\mathcal A_l^{\rm s}(b)$ belonging to the kernels outside the
exponent. They transform by $\operatorname{Ad}u(b_-)$.
Near bonds are the bonds in collars, rims and the nonanalytic cutoff
zone; all other stored bonds are far. Fix
$\lambda_l^{\rm st}=(\log\theta_l)^{6r+1}>1$.
A stored-field function is Borel in the near variables and is the
value at $\mu_b=1$ of a function jointly analytic in the pair, the
far variables, and switches $|\mu_b|<\lambda_l^{\rm st}$.
It is independent of a near variable when its switch is zero, and
independent of that switch when the variable is zero. For far
variables the analytic ball has radius $\lambda_l^{\rm st}$ times
their specified real range. That range is part of the kernel data.
The pair tie is required at every switch value.

\subsubsection{Presentations, standing operations, and their bounds}

\phantomsection\label{notdef:activity-families}
A real presentation consists of the frame and formation data above,
weights $x_j$, regular members
$(E^{(j),\mathrm u},E^{(j),\mathrm m},\widehat E^{(j)},
b_j^{\mathrm u},b_j^{\mathrm m})$, remainder families
$(R^{(j)},\widehat R^{(j)})$, boundary functions
$B_{\mathcal S}(j,X,o)$, and the standing-operation exponents below.
Here $E=E^{\mathrm u}+E^{\mathrm m}$,
$b_j=b_j^{\mathrm u}+b_j^{\mathrm m}=x_{j-1}-x_j
=b[\widehat E^{(j)}]$. Boundary origins $o$ range over five
types: the step, the last layer of a regular family, a step remainder
beside an island, a stage, and exponentiation. A slot is the tuple
$(j,X,o)$ together with the structure-domain on which the function
was formed. It retains that domain at later steps.

For a finite exempt family $\mathsf H$ of nonempty regions, define
$d^\natural_{j,\mathsf H}(X)$ by minimizing $M^{-1}$ times tree
length also over subfamilies $\mathsf F\subset\mathsf H$ contained
in $X$: the tree meets each cube outside $\bigcup\mathsf F$ and
each member of $\mathsf F$, but need not fill those members.
For a label use the level-$j$ hulls of its standing islands as
$\mathsf H$. Denote the result by $d_j^\natural(\mathcal S;X)$.
It is at most $d_j(X)$ and asserts no decay across an exempt island.
The real class requires

\refstepcounter{equation}
\begin{gather*}
\FitMath{\|E^{(j)}\|_{j,\kappa}\le E_0,} \\[3pt]
\FitMath{\|R^{(j)}\|_{j,\kappa}\le \bar g_j^{\kappa_0},} \\[3pt]
\FitMath{|B_{\mathcal S}(j,X,o)|\le\mathsf B
                e^{-\kappa d_j^\natural(\mathcal S;X)}.}\tag{\theequation}\label{eq:rg-objects-R23}
\end{gather*}

Boundary functions have the stored-field property. All entries are
real on real arguments and invariant under the fixed simultaneous
action on pairs and stored data. For regular families in \eqref{eq:rg-objects-R23} is
evaluated at $\check\alpha_{\iota,j}$; for the unmarked family
also require a member of size at most $E_0$ on fixed absolute radii.
The whole-lattice step is additional formed data: its positive tame
kernel is $\mathbf T^{\rm sf}_{j-1}[\theta_{j-1}]$, with its
non-pivot bond coordinates and sharp fluctuation box. For a list
action $S$ its step functional is

\begin{gather*}
\FitMath{\begin{split}
 \mathcal E^{(j)}[\theta;x;S](W)
   ={}&\log\big(\mathbf T^{\rm sf}_{j-1}[\theta]
                   e^{-xA(U_{j-1})+S(U_{j-1})}\big)(W)\\
    &+xA(U_j(W))-S(U_j(W)).
 \end{split}}
\end{gather*}

The kernel, the fibre coordinates, the normalization and the
near-real logarithm branch must be given for this definition.
The localization rule producing a pointed family from this
functional is supplied once, for every torus and list; it is not
determined by its total. The identification clause requires
$E^{(j)}$ to be exactly this localized output from the earlier
entries. Its unmarked version is the same output using only their
unmarked entries.

The formed data also contain pointed Gaussian families
$\mathbf N_j^\circ$ and $\mathbf G_{j,\tau}$ on those absolute
cores, the first localizing the Gaussian logarithmic normalizer
and its trivial-field local numbers, and the second localizing
the change of its quadratic form when $x/\theta=\tau$.
Their finite-dimensional covariance matrices, localization maps and
logarithm branches are typed input data, so the following conditions
have definite meanings. Set
$b_j^0=b[\mathbf N_j^\circ]$ and
$\tau_{j-1}=x_{j-1}/\theta_{j-1}$.
Require, for fixed constants $C_{\rm m},C_{\rm u},C_b$,

\begin{gather*}
\FitMath{\begin{split}
 \|E^{(j),\mathrm u}-\mathbf N_j^\circ
                 -\mathbf G_{j,\tau_{j-1}}\|_{j,\kappa}
       &\le C_{\rm u}\bar g_{j-1}(\log\theta_{j-1})^{\mathfrak p_0},\\
 \|E^{(j),\mathrm m}\|_{j,\kappa}
       &\le C_{\rm m}\bar g_{j-1}^{\kappa_0-C_{\rm m}},\\
 |b_j^{\mathrm m}|&\le C_b \bar g_{j-1}^{\kappa_0-4},\qquad
 |b_j^{\mathrm u}-b_j^0|
       \le C_b \bar g_{j-1}(\log\theta_{j-1})^{\mathfrak p_0}.
 \end{split}}
\end{gather*}

No universality formula or asymptotic-freedom estimate for $b_j^0$
is implied by calling these families Gaussian. Their covariance
construction and coefficient estimates are separate analytic input.

\phantomsection\label{notdef:standing-history}
A standing operation at an island $(t,Y)$ has a countable or finite
local-history index $\omega$, a Borel subset $S_\omega$ of a finite
product of group variables, real bond variables and finite frame
sets, and a positive Borel measure $\mu_\omega$ on that set.
These sets and measures are formed independently of the entries.
Its exponents are Borel in $\xi\in S_\omega$, holomorphic in the
pair, real at real arguments, and have the form

\refstepcounter{equation}
\begin{gather*}
\FitMath{\mathbf T'_t(Y)F=\frac1{|W_4|}\sum_\omega
       \int_{S_\omega}F(\xi)
            e^{\Phi_\omega(\xi;\mathbf p)}\mu_\omega(d\xi),} \\[3pt]
\FitMath{\Phi_\omega=\Phi_\omega^\circ(\xi)
          +\sum_f\Phi_{\omega,f}(\xi;\mathbf p|_{R_f}).}\tag{\theequation}\label{eq:rg-objects-R24}
\end{gather*}

Each carrier $R_f\subset Y^\sim$ supports both dependencies of
its factor. The finite factor index set, the measurable coordinate
projections onto stored variables in each carrier, and any nested
island variables are part of the data. No variable is integrated
twice: swallowed-operation variables retain their earlier labels
inside the product assigned to the swallowing component.

If $(s,D)$ ranges over the forest of nested stored components, set
$v(\omega)=\sum_{(s,D)}\log\theta_s\,\#_s(D)$, where
$\#_s$ counts sites in their own level. For a fixed creation-time
complex domain $\mathcal K_{t,Y}$ require

\refstepcounter{equation}
\begin{gather*}
\FitMath{\begin{aligned}
 &\frac1{|W_4|}\sum_\omega e^{c_{\rm r}v(\omega)}
    \int_{S_\omega}^{*}\sup_{\mathbf p\in\mathcal K_{t,Y}}
       |e^{\Phi_\omega(\xi;\mathbf p)}|\mu_\omega(d\xi)
       \le C e^{-2p_0(\bar g_t)/(1+\beta_0)},\\
 &\sum_{f:R_f=R}\operatorname{osc}\Phi_{\omega,f}
       \le\mathsf B_\Phi(\log\theta_t)^{2\mathfrak p_0}
                  e^{-\kappa d_t^\natural(\mathcal S;R)}.
 \end{aligned}}\tag{\theequation}\label{eq:rg-objects-R25}
\end{gather*}

The first integral is an upper integral, allowing a nonmeasurable
supremum. The oscillation is the supremum of
$|\Phi_{\omega,f}(\xi;p)-\Phi_{\omega,f}(\xi;p')|$ over $\xi$
and both pairs in that fixed domain. The domain is supplied from
the data at creation and must contain every later evaluation;
this last inclusion is an independent requirement.
$c_{\rm r}\ge0$ is fixed before the coupling ceiling, and
$\beta_0>0$ satisfies
$p_0(\bar g_m)\ge(1+\beta_0)^{-1}p_0(\bar g_n)$ for $m<n$.
The supremum remains inside the integral and sum. There is no
asserted small factor in the size of $Y$. A drop-out has no new free
entry of type \eqref{eq:rg-objects-R24}: it keeps its partially formed operation, and
its carrying Wilson factor is included in its separate majorant.
Such a majorant must be supplied to use a drop-out estimate.

The effective exponent additionally requires finite index maps and
vacuum numbers, which are formed data. At label $\mathcal S$ the
map $\mathcal I_j^E(\mathcal S)$ is a subset of
$\{(z,X):z\in\Lambda_j\cap\mathbb T^{(j)},\ X\in\mathcal D_j,
z\in X\}$; it gives the actual enlarged-region rule for the regular
terms. The enlargement width and treatment of boundary-crossing
domains must be included in that map. The remainder index map is
$\mathcal I_j^R(\mathcal S)=
\{X\in\mathcal D_j:X\subset\Lambda_j^{\sim-1}\}$.
The boundary index map selects the slots meeting the one-layer
enlargement of $Z_j$. These maps use the action sequence in \eqref{eq:rg-objects-R21}.
The profiles use the small-field regions of the action sequence.
Fix a nonnegative smooth function \(\psi\) on \(\mathbb R^4\),
invariant under signed coordinate permutations, supported in
\(\{|u|_\infty\le1/4\}\), and normalized by \(\int\psi=1\).
For example, normalize the product of the four functions
\(b(u_\nu)\), where
\[
 b(v)=
 \begin{cases}
  \exp[-1/(1-16v^2)],&|v|<1/4,\\
  0,&|v|\ge1/4.
 \end{cases}
\]
At level \(j\), put \(\mathsf s_j=M\check R_j\) when
\(\Lambda_j\) is a union of those cubes, and \(\mathsf s_j=M\)
otherwise. The latter region is a union of \(M\)-cubes.
Measure distances and translations in level-\(j\) units and set
\(D_j(x)=\operatorname{dist}_\infty(x,\Lambda_j^c)\), with
\(D_j=+\infty\) when \(\Lambda_j=T_\ell\). Define on the torus
\begin{equation}\label{eq:rg-profile-construction}
 \varphi_j(x)=\int_{\mathbb R^4}\psi(u)
       \mathbf1\{D_j(x-\mathsf s_ju)\ge\mathsf s_j/2\}\,du,
\end{equation}
and evaluate it at plaquette midpoints in the action.
The geometry and this rule are held fixed as a source varies.

The distance function is one-Lipschitz when the complement is
nonempty. Thus \(0\le\varphi_j\le1\), it vanishes where
\(D_j\le\mathsf s_j/4\), and it equals one where
\(D_j\ge3\mathsf s_j/4\). Indeed, wherever \(\psi(u)\ne0\)
one has \(|u|_\infty<1/4\), so the argument of \(D_j\) moves
its value by strictly less than \(\mathsf s_j/4\).
In particular \(\varphi_j=1\) one layer inside \(\Lambda_j\),
with that layer taken in the \(\mathsf s_j\)-cube partition.
Its support is contained in \(\{D_j\ge\mathsf s_j/4\}\),
which is a compact subset of the interior of \(\Lambda_j\).
For the whole torus the formula gives \(\varphi_j=1\).
On the periodic lift, changing variables in the convolution and
differentiating its smooth kernel gives, for every multi-index,
\[
 \|\partial^\alpha\varphi_j\|_\infty
 \le\mathsf s_j^{-|\alpha|}
                  \|\partial^\alpha\psi\|_{L^1(\mathbb R^4)}.
\]
This also proves smoothness, since the indicator is bounded and
each kernel derivative has compact support. Cross-level ordering,
where used below, is an additional condition. It follows, for
example, if \(\Lambda_{j+1}\subset
\{D_j\ge3\mathsf s_j/4\}\): then \(\varphi_j=1\) throughout
the support of \(\varphi_{j+1}\).
With $x_k(p)=\beta-\sum_{j\le k}b_j\varphi_j(p)$, the exponent is

\begin{gather*}
\FitMath{\begin{split}
 A_{k,\mathcal S}(U)={}&-A(x_k(\cdot),U)
    +\sum_{j,(z,X)\in\mathcal I_j^E(\mathcal S)}
       \{E^{(j)}(X,U,z)-E^{(j)}(X,\mathbf1,z)\}\\
 &-\sum_{j\le k}b_jA(\varphi_j,U)
    +\sum_{j,X\in\mathcal I_j^R(\mathcal S)}
       \{R^{(j)}(X,U)-R^{(j)}(X,\mathbf1)\}\\
 &+\sum_{(j,X,o)\text{ in boundary index map}}
          B_{\mathcal S}(j,X,o;U,\text{stored data})
                  -\mathcal E_{k,\mathcal S}.
 \end{split}}
\end{gather*}

The real scalar $\mathcal E_{k,\mathcal S}$ is supplied with its
vacuum bookkeeping rule, including the stage numbers retained by
drop-outs. The pointwise identity
$x_k(p)+\sum_{j\le k}b_j\varphi_j(p)=\beta$ is exact. It does not
authorize changing the index maps or the vacuum rule. These definitions specify the objects; their construction and uniform
estimates remain hypotheses where no proof is given.

Erase every healed component together with its ancestral cluster
from a history $h$; the result is its standing part $h^\flat$.
Erasure being an admissible history with the same label, and the
hard-core decomposition used to sum healed clusters, are structural
requirements on the formation data. They are not assumed merely
from the word ``erase.'' The never-large standing part has
$\Omega_j=\Lambda_j=T_\ell$. The presentation gives finite positive
measures $\widehat\tau_{h^\flat}$ on $V_k$, stored fields and
coordinate frames, with marginals $\tau_{h^\flat}$ on $V_k$.
It also retains the individual-history measures $\widehat\tau_h$
produced by the exact kernels. Their identities are

\refstepcounter{equation}
\begin{gather*}
\FitMath{\rho_k(V_k)dV_k=\sum_{h^\flat}\tau_{h^\flat}
             =\sum_h\tau_h,} \\[3pt]
\FitMath{\sum_h\tau_h(\mathcal V_k)=Z.}\tag{\theequation}\label{eq:rg-objects-R26}
\end{gather*}

The first equality uses the summed healed operations; the exponent
of an individual $h$ need not have the effective-action form.

The reading at $X$ records the standing regions, outcomes, marks,
stored test bonds and frames within the hull of $X$ enlarged by one
level-$l$ cell at every scale $l$. A chart is a translation by
top-level cells, a common lift between tori, or their composite.
Coherence requires a chart matching readings to carry the terms,
carriers, Borel sets and measures to their counterparts. Factorwise,

\refstepcounter{equation}
\begin{gather*}
\FitMath{\Phi_{\iota\omega,\iota f}(\iota\xi;\iota p)
                       =\Phi_{\omega,f}(\xi;p).}\tag{\theequation}\label{eq:rg-objects-R27}
\end{gather*}

Unit translations not preserving the partitions are covered by the
pointed-total requirement \eqref{eq:rg-objects-R14}, not by this chart definition.
The simultaneous lattice symmetries preserving level-$j$ cells
must also transport the measures with their variables.

For a datum at $a$ just below a pillar top $y$, let
$\tau[U](a)=(\mathcal M^{m-1}U)(y,a)$ be its contour transport.
An adapted background satisfies $\tau[U](a)=\mathbf1$ at every
stored datum site. To define $\widehat\tau$ require that all
backgrounds exist on its support, every argument belongs to its
entry's domain, and the exponential has one value at all adapted
representatives. The last clause is independent of simultaneous
gauge invariance: an orbit transformation need not fix every datum
site. These are part of class membership, not formal consequences
of the norm bounds.

Finally let $\mathcal Z_k(h)$ be the live region after healing.
For a live component $X$, define its lineage downwards by
$\mathsf Z_k(X)=X$ and by taking at level $j$ all live components
contained in $\mathsf Z_{j+1}(X)$. Write $|C|_j$ for its number of
level-$j$ cells and $d'_j(W)$ for tree length in those cells.
The created sets in a lineage are the components $W$ of the step's
creation set there; operation marks do not enter its deposit.
With one fixed deposit constant $c>0$, set

\refstepcounter{equation}
\begin{gather*}
\FitMath{\begin{aligned}
 \operatorname{Vol}_k(X;h)&=
       \sum_{(j,C)\text{ in its lineage}}
                       2M^4\check R_j^5|C|_j,\\
 \operatorname{Dep}_k(X;h)&=
       \sum_{(j,W)\text{ created in its lineage}}
                       c p_0(\bar g_j)^2(d'_j(W)+1),\\
 \Theta^{(k)}(h)&=\sum_{X\text{ live at }k}
       \{w\operatorname{Dep}_k(X;h)
                          +c_a\operatorname{Vol}_k(X;h)\},
                 \quad 0\le w<1,\quad c_a\ge0.
 \end{aligned}}\tag{\theequation}\label{eq:rg-objects-R28}
\end{gather*}

The end mark is $\Theta^{(K)}$. Label marks and this numerical
weight are different objects. Prepaid stability would assert
$\sum_he^{-E}\tau_h(\mathcal V_k)e^{\Theta^{(k)}(h)}
\le e^{E_+\ell^4}$, with the stipulated normalization and dials;
it requires a proof from creation deposits and the operation
majorants. Nothing in \eqref{eq:rg-objects-R28} establishes that assertion.

\phantomsection\label{notdef:presentation-class}
Thus the class $\mathfrak C_k(\theta)$ is the set of these complete
typed presentations satisfying \eqref{eq:rg-objects-R13}--\eqref{eq:rg-objects-R28}, all their stated
structural and positivity clauses, and the specified analytic
domains. A step-closure theorem must provide the next presentation,
prove those clauses with the same constants and margins, and prove
equality of the represented measures under its kernels. The
background, response, substitution, island-mass, drop-out-majorant
and closure hypotheses remain the substantive unresolved inputs
where independent proofs have not been supplied.


% END MODULE chapters/rg_objects_v2.tex

% BEGIN MODULE chapters/01_orientation.tex
\subsection{The weak-coupling sequence and continuum statements}
\label{candidate:orientation}
\subsubsection{The object of the construction and the source of its difficulty}

We state the continuum and phase theorems. The coupling flow fixes the physical scale, and source derivatives define the connected correlations.

\phantomsection\label{notdef:observables}
\textbf{Objects.} We take $G=\operatorname{SU}(2)$ and $N=2$. Put \(T_\ell=(\mathbb R/\ell\mathbb Z)^4\) and
\(\Lambda_{K,\ell}=a_K\mathbb Z^4/\ell\mathbb Z^4\), where
\(a_K=L^{-K}\). The compatible sides are
\(\mathfrak L=\{2L^m:m\in\mathbb Z_{\ge0}\}\).
At reference coupling \(g\) we use
\(\mathfrak L(g)=\{\ell\in\mathfrak L:\ell\ge\ell_\star(g)\}\);
the floor in \eqref{eq:flow-RF35} makes room for the localization
cubes at every level. On such a torus we consider the following objects: the lattice $\Lambda_{K,\ell}$, whose spacing is $a_K=L^{-K}$; the configurations $U$ on it; the quantity $a_p(U)=1-\tfrac12\operatorname{Re}\operatorname{tr}U_p$; Wilson's measure; the smeared curvature density; and the connected correlation functions of that density:

\begin{gather*}
\FitMath{\mu_{K,\ell,\beta}(dU)=Z_{K,\ell,\beta}^{-1}\,e^{-\beta\sum_pa_p(U)}\,dU,} \\[3pt]
\FitMath{\mathcal{O}(f)=\sum_pf(x(p))\,a_p(U),} \\[3pt]
\FitMath{S^T_{n;K,\ell,\beta}(f_1,\dots,f_n)=\langle\mathcal{O}(f_1);\dots;\mathcal{O}(f_n)\rangle_{K,\ell,\beta}.}
\end{gather*}

In view of $a_p\approx\tfrac12a_K^4|F_U(p)|^2$, the sum $\mathcal{O}(f)$, which carries no prefactor involving the spacing, is a Riemann-sum approximation of $\tfrac12\int f\sum_{\mu<\nu}|F_{\mu\nu}|^2$. It is therefore the $\operatorname{tr}(F^2)$ observable. The test functions are restricted to separated tuples
\((f_i)\in\mathfrak S_n(\mathbb E)\): tuples of real \(C_c^1\)
functions with pairwise disjoint supports. We use
\(\|f\|_\sharp=\max\{\|f\|_\infty,c_M\operatorname{Lip}f\}\),
where \(c_M>0\) is the fixed multiple of the localization size
\(M\) used in \cref{candidate:sources}. The reason is given under ``The frame'' in \cref{thm:continuum-main}.

\phantomsection\label{notdef:first-passage}
\textbf{The cutoff is tied to the coupling.} The renormalisation group of \cref{candidate:averaging,candidate:large_field} assigns to $\beta=g_0^{-2}$ a sequence of running inverse couplings $x_0=\beta$, $x_k=x_{k-1}-b_k$, whose increments satisfy $b_k\le B_\sharp$. These increments are independent of the torus \cref{sec:flow-a-common-companion-and-the-torus-floor}, and also of the cutoff $K'\ge k$ used to carry out the steps \cref{sec:flow-units-common-prefixes-and-change-of-reference}. We fix a small reference coupling $g$. As the unit of length we take the scale at which the flow reaches, for the first time, the reference window $g^{-2}\le x\le g^{-2}+B_\sharp$ belonging to $g$:

\begin{gather*}
\FitMath{K(\beta,g)=\min\{k:\,x_k(\beta)\le g^{-2}+B_\sharp\},} \\[3pt]
\FitMath{\mathcal{W}_K(g)=\{\beta\ge g^{-2}+B_\sharp:\,K(\beta,g)=K\};}
\end{gather*}

a pair $(K,\beta)$ such that $\beta\in\mathcal{W}_K(g)$ is called a first-passage datum.

The remaining pairs with $\beta\ge g^{-2}+B_\sharp$ are accounted for as follows. Suppose that $K\ge K'=K(\beta,g)$, put $b=K-K'$, and let $f^{[b]}(x)=f(L^{-b}x)$ on $T_{L^b\ell}$. Then

\begin{gather*}
\FitMath{\langle\mathcal{O}(f_1);\dots;\mathcal{O}(f_n)\rangle_{K,\ell,\beta}=\langle\mathcal{O}(f_1^{[b]});\dots;\mathcal{O}(f_n^{[b]})\rangle_{K',\,L^b\ell,\,\beta}}
\end{gather*}

holds exactly, for all bounded $f_i$. To see this, note that the dilation $x\mapsto L^bx$ maps $\Lambda_{K,\ell}$ onto $\Lambda_{K',L^b\ell}$, plaquettes to plaquettes and barycentres to barycentres. Wilson's weight depends only on the plaquette variables and on Haar measures, and no power of the spacing enters $\mathcal{O}$. The same identity holds between a periodic Gibbs state at $\beta$ on $\mathbb{L}_K$ and its image on $\mathbb{L}_{K'}$ \cref{sec:flow-units-common-prefixes-and-change-of-reference}. Consequently a pair with $K>K(\beta,g)$ amounts to the first-passage datum of $\beta$ on a dilated torus, on which $\|f^{[b]}\|_\sharp\le\|f\|_\sharp$ and every distance is multiplied by $L^b$. The restriction to first passage therefore loses none of these lattice measures; its only effect is to fix the unit of length. (When $K<K(\beta,g)$ the identity is used in the opposite direction, on a contracted torus. This requires a lower bound on $\ell$, and the norms grow by at most the factor $L^{K(\beta,g)-K}$ \cref{sec:flow-units-common-prefixes-and-change-of-reference}.)

\phantomsection\label{notdef:frozen-coupling}
The final value $x_K(\beta)$ lies in the window but is not fixed at $g^{-2}$. This is why (a)(ii) passes to subsequences, and it is the purpose of clause (d). On each torus, provided the conditions of (d)(1) are satisfied, that clause dispenses with the subsequences as soon as the renormalised coupling, defined below, converges along the sequence. On $\mathbb{R}^4$, in (a)(iii), a subsequence is taken in every case, including the case in which the renormalised coupling converges (as stated in \cref{thm:continuum-main}(a)(iii)). The renormalised inverse coupling $x^\natural_K(\beta)$ is the final value of a particular run of $K$ steps starting at $\beta$. In that run, the auxiliary choices of level $k$ that depend on the coupling are made with $\theta_k=g^{-2}+c_0(K-k)$ instead of $x_k$, where $c_0=\tfrac{11N^2}{12\pi^2}\log L$. Here and in all steps, $\log$ denotes the natural logarithm. A second kind of run is needed for the following reason. Continuity and holomorphy in the bare coupling are established by setting the constructions at two bare couplings side by side and comparing them entry by entry, and by letting a single fixed linear operation act on a whole family of bare densities. Each of these two devices needs the choices to coincide at the two couplings. In a run made with its own couplings, however, the choices vary with $\beta$. Some of them, namely the cube sizes $\check{R}_k$, change through integer values \cref{sec:flow-frozen-choices-slopes-and-tuned-bare-couplings}. Freezing the choices leaves $Z_{K,\ell,\beta}$ and $S^T_n$ unchanged, since these are quantities of the measure. At first-passage data, and under the conditions of (d)(1), the shift that freezing produces in the couplings is at most $g$ times a power of $\log g^{-2}$ \cref{sec:flow-frozen-choices-slopes-and-tuned-bare-couplings}, \cref{candidate:flow,ch:comparison-capsule-v7}. By a continuum limit at $g$ we mean a family $S=(S_n)_{n\ge2}$ with

\begin{gather*}
\FitMath{S_n(X)=\lim_{j\to\infty}S^T_{n;K_j,\ell_j,\beta_j}(X),} \\[3pt]
\FitMath{\beta_j\in\mathcal{W}_{K_j}(g),} \\[3pt]
\FitMath{K_j\to\infty,}
\end{gather*}

taken either on $T_\ell$ (with all $\ell_j=\ell$) or on $\mathbb{R}^4$ (with $\ell_j\to\infty$). Limits on a torus along the tuned couplings of (a)(i) are also counted as continuum limits. \phantomsection\label{notdef:torus-limit-states}
A volume-first limit at $g$ is a limit $\lim_jS^T_n[\nu_j]$, taken along first-passage data $(K_j,\beta_j)$ with $K_j\to\infty$. In it each $\nu_j\in\mathcal{G}^{\mathrm{tor}}_{K_j,\beta_j}$ is a limit point of the periodically extended torus measures as $\ell\to\infty$, with spacing and coupling held fixed. Thus the volume is removed first and the spacing only afterwards.

\textbf{The two theorems in brief.} The exponents are fixed first. Then $L$ is taken large, then the remaining constants are chosen in their order, and finally $g$ is taken small. \Cref{thm:continuum-main} states the following. The two assertions below hold uniformly in the first-passage datum, and the left-hand one holds also uniformly in the torus as soon as $\ell\ge Cn(R+1)$:

\begin{gather*}
\FitMath{\big|S^T_{n;K,\ell,\beta}(f_1,\dots,f_n)\big|\le C^\infty_n(\mathsf{d};R)\prod_i\|f_i\|_\sharp,} \\[3pt]
\FitMath{S_2(f,h)=\frac{3(N^2-1)}{\pi^4}\Big(\frac{12\pi^2}{11N^2}\Big)^2\frac{\mathcal{K}_{\mathrm{free}}(f,h)}{(\log(1/\mathsf{d}))^2}\,(1+\eta)>0.}
\end{gather*}

In the left-hand bound, $n\ge2$, and the supports lie in balls of radius $R$ and are pairwise at distance $\ge\mathsf{d}$. In the right-hand formula, $S$ is a limit along first-passage data, $f,h\ge0$ are bumps whose size and separation are of the order of a small $\mathsf{d}$, and $\mathcal{K}_{\mathrm{free}}(f,h)=\iint f(x)h(y)|x-y|^{-8}\,dx\,dy$. The factor $1+\eta$ is bounded above and below by positive constants, and on each torus $\eta\to0$ as $\mathsf{d}\to0$; on $\mathbb{R}^4$ the clause consists of the bounds by the two constants and nothing more. The five clauses of this theorem are the following:
\begin{itemize}
\item limits exist on the tori and on $\mathbb{R}^4$ (a);
\item they inherit translation and hypercubic invariance, reflection positivity, symmetry, regularity and growth (b);
\item they are non-trivial and have, at short distance, the normalisation of asymptotic freedom (c);
\item on a torus they depend only on the renormalised coupling (d);
\item on $\mathbb{R}^4$ they are rotation invariant if and only if a flux through far shells vanishes (e).
\end{itemize}

\phantomsection\label{notdef:centered-functional}
\cref{thm:phase-main} states that a volume-first limit $S$ is $\operatorname{O}(4)$-invariant. Its centred functional $\sigma_S=\sum_{\pi\in\operatorname{Part}_{\ge2}}\prod_{J\in\pi}S_{|J|}$ is precisely the system of moments obtained when the one-point function is set to zero, which is to say the moments of the centred density. This functional admits a canonical decomposition, which consists of a probability measure $\lambda$, carried by a compact set $\mathsf{X}=\operatorname{supp}\lambda$, together with a continuous injective map $\xi\mapsto\sigma_\xi$ such that

\begin{gather*}
\FitMath{\int_{\mathsf{X}}\prod_{i=1}^q\sigma_\xi(X^i)\,\lambda(d\xi)=\lim_{|v_q|\to\infty}\cdots\lim_{|v_2|\to\infty}\sigma_S\big(X^1\cdot\tau_{v_2}X^2\cdots\tau_{v_q}X^q\big)} \\[3pt]
\FitMath{(q\ge1),} \\[3pt]
\FitMath{\text{so that}} \\[3pt]
\FitMath{\sigma_S=\int_{\mathsf{X}}\sigma_\xi\,\lambda(d\xi).}
\end{gather*}

Each ergodic component $\sigma_\xi$ satisfies the full set of axioms of Osterwalder and Schrader.


\paragraph{The theorem, stated clause by clause}
\label{sec:orientation-the-theorem-stated-clause-by-clause}

There exist a shape constant $\vartheta_0\in(0,1]$ and, for each source radius $\varrho>0$ (\cref{candidate:sources}) and each shape $\vartheta\in(0,1]$, thresholds $0<\gamma_{U'}\le\gamma_U\le\gamma_J$ and $0<\gamma_W(\vartheta)\le\gamma_J$ such that \cref{thm:phase-main} and the clauses below hold. Throughout, $0<g\le\gamma_J$, $\ell\in\mathfrak{L}(g)$, and $(K,\beta)$ is a first-passage datum. Apart from dependences that are displayed, the constants do not vary with $K,\beta,\ell$ or with positions. The ceilings $\gamma_{U'},\gamma_U,\gamma_J$ depend only on the constants fixed before them, among which is $\varrho$, and $\gamma_W(\vartheta)$ depends only on these and on $\vartheta$. Each constant, threshold and function introduced by a clause of the two theorems may in addition depend on $g$ and on the quantities displayed with it. 

\textbf{The frame is adapted to the tools.} When the side is $\ell=2L^m$ with $m\in\mathbb{N}$, each axis carries an even number $2L^{m+K}$ of sites. Each mirror pair $\{y_\mu=c\}\cup\{y_\mu=c+\ell/2\}$ therefore cuts the torus into two halves, and the chessboard cells come in pairs \cref{candidate:sources,ch:all-state-source-transfer}. Because $L$ is odd, the mirrors that bisect bonds at successive cutoffs are nested: $a_K(\mathbb{Z}+\tfrac12)\subset a_{K+1}(\mathbb{Z}+\tfrac12)$. Because barycentres are used, the identity $\mathcal{O}(f)\circ\mathsf{r}^*=\mathcal{O}(\mathsf{r}f)$ is exact for each lattice isometry $\mathsf{r}$, where $\mathsf{r}^*$ denotes its action on configurations. If the support of $f$ has distance greater than $Ca_K$ from the mirror, every contributing plaquette and all its bonds stay on that same side. A barycentre alone does not guarantee this for a support touching the mirror; the continuum reflection-positivity passage uses fixed positive support gaps. In the Gaussian curvature model of
\cref{prop:averaging-1a-2-10}, let \(V_{\rm lat}=(\ell/a_K)^4\)
be the number of sites. Its exact two-cochain space has dimension
\(3(V_{\rm lat}-1)\) per colour: the kernel of the one-cochain
coboundary consists of the \(V_{\rm lat}-1\) gradients and four
harmonic one-cochains. Translation and hypercubic symmetry make
the diagonal of its orthogonal projector constant on the
\(6V_{\rm lat}\) plaquettes. Thus the quadratic observable has
expectation
\[
 \big\langle\mathcal O_{\rm quad}(f)\big\rangle_{\rm G}
 =\frac{N^2-1}{4\beta}
       \left(1-\frac1{V_{\rm lat}}\right)
          \sum_p f(x(p)).
\]
For fixed \(\ell\) and smooth \(f\), its leading Riemann-sum term is
\(\frac{3(N^2-1)}{2\beta}a_K^{-4}\int f\).
Along the weak-coupling sequence \(\beta\asymp K\), this term
diverges when \(\int f\ne0\). This is a calculation in the Gaussian
model; the interacting continuum construction uses the exact
lattice expectation for centering. Moreover $|x-y|^{-8}$ is not locally integrable. For these reasons only $n\ge2$ and only separated tuples are considered, and positivity is stated for the centred $\sigma_S$. The axioms of Osterwalder and Schrader require nothing more, because their test functions vanish to infinite order wherever points coincide. By a mirror $\Sigma$ of $\mathbb{R}^4$ we mean a hyperplane $\{u\cdot x=c\}$ with either $u=e_\mu$, the unit vector of the $\mu$-th coordinate (coordinate), or $u=(e_\mu\pm e_\nu)/\sqrt2$ for some $1\le\mu<\nu\le4$ (diagonal). Its reflection is $\theta_\Sigma$ and its open side is $\mathbb{E}_+(\Sigma)$. A torus has only the coordinate pairs as mirrors, and there the open side is a slab.

\begin{theorem}[Continuum correlations]\label{thm:continuum-main}
Assume the uniform background, interaction, history, flow and ordinary-source
estimates of \cref{ch:dependency-catalogue}, with their stated
domains and normalization. For the short-distance clauses assume the
interacting comparison and running-normalization estimates there; for
the matched-cutoff and Ward clauses assume their additional estimates.
Then the following conclusions hold.
\textbf{(a) Existence.} Each run started from $\beta\ge g^{-2}+B_\sharp$ reaches first passage. For such a run, $c(\beta-g^{-2})-C\le K(\beta,g)\le C(\beta-g^{-2})+C$, and since $b_k\le B_\sharp$ also $g^{-2}<x_K(\beta)\le g^{-2}+B_\sharp$. Hence $K(\beta,g)\to\infty$ as $\beta\to\infty$, and so $\mathcal{W}_K(g)\ne\varnothing$ for infinitely many $K$. The values $x_k(\beta)$ are the same whatever torus the steps are carried out on, and $\beta/\log(1/a_K)\to\tfrac{11N^2}{12\pi^2}$ uniformly on $\mathcal{W}_K(g)$. (i) Let $g\le\gamma_U$. Then for each $K$ there is a tuned $\beta^\natural_K$, that is, one with $x^\natural_K(\beta^\natural_K)=g^{-2}$, and it satisfies $\tfrac12c_0K-c_5\le\beta^\natural_K-g^{-2}\le\tfrac32c_0K+c_5$. Moreover $S^T_{n;K,\ell,\beta^\natural_K}\to S^{\natural,\ell}_n$ as $K\to\infty$ along the whole sequence. This convergence holds on each torus of the family and on tuples of Lipschitz functions whose supports are pairwise disjoint. (ii) On a fixed torus, from every sequence of first-passage data with $K_j\to\infty$ one can extract a subsequence that converges on $\mathfrak{S}(T_\ell)$. (iii) The same holds on $\mathfrak{S}(\mathbb{R}^4)$ when in addition $\ell_j\to\infty$. The limits obtained in this way along tori, with defining sequence $(K_j,\ell_j,\beta_j)$, are what the phrase ``a limit of (a)(iii)'' means. The same also holds for $S^T_n[\nu_j]$ with $\nu_j\in\mathcal{G}^{\mathrm{tor}}_{K_j,\beta_j}$, so volume-first limits exist.

\textbf{(b) Symmetry, positivity and growth.} Each continuum limit has the following properties.
(i) It is invariant under $W_4$ and under the translations of $\mathbb{E}$.
(ii) It is reflection positive at each coordinate mirror: $\sum_{i,j}\bar c_ic_j\,\sigma_S(\theta_\Sigma X_i\cdot X_j)\ge0$ whenever $X_i\in\mathfrak{S}(\mathbb{E}_+(\Sigma))$.
(iii) It is symmetric.
(iv) It is regular: there is a unique distribution of locally finite order on $\mathbb{E}^n_{\ne}$ of which $S_n$ is the restriction. On $C^1$ functions, and with the same constants, it satisfies (c)(1) on $T_\ell$, and it satisfies both bounds (c)(2) when it is a limit of (a)(ii)--(iii).
(v) On $\mathbb{R}^4$ it satisfies the growth condition 

\begin{gather*}
\FitMath{|\mathsf{S}^{\sigma_S}_n(\varphi)|\le c_{\mathrm{E}}C_{\mathrm{E}}^n(n!)^{\varkappa}|\varphi|_{k_{\mathrm{E}}n}}
\end{gather*}

 on $\mathcal{S}_{\ne}(\mathbb{R}^{4n})$. Take the form $\langle F,G\rangle=\sigma_S(\theta\bar F\cdot G)$ on $\mathbb{C}\mathfrak{S}(\{x_4>0\})$. The separable space $\mathcal{H}(S)$ then carries three structures:
\begin{itemize}
\item $\mathcal{T}_t[F]=[\tau_{te_4}F]$, $\mathcal{T}_t=e^{-tH}$, which form a strongly continuous semigroup of self-adjoint contractions, where $H\ge0$ and $H\Omega=0$;
\item a strongly continuous unitary representation $e^{iv\cdot P}$ of $\mathbb{R}^3$ that leaves $\Omega$ fixed, the $P_j$ commuting strongly with $H$;
\item a unitary representation of $W_3$ that leaves $\Omega$ fixed and commutes with $H$.
\end{itemize}

Already for $g\le\gamma_J$, each volume-first limit is a limit of (a)(iii) and is, in addition, reflection positive at each diagonal mirror.

\textbf{(c) Correlation bounds and nontriviality.} Let $n\ge2$, let the $f_i$ be Lipschitz with supports pairwise at distance $\ge\mathsf{d}>0$, and in (2) let the supports also lie in balls of radius $R$. Then

\begin{gather*}
\FitMath{(1)\,|S^T_n|\le C_n(\mathsf{d})\,e^{nc_V\ell^4}\prod_i\|f_i\|_\sharp;} \\[3pt]
\FitMath{(2)\,|S^T_n|\le C^\infty_n(\mathsf{d};R)\prod_i\|f_i\|_\sharp\,\,(\ell\ge Cn(R+1)),} \\[3pt]
\FitMath{|S^T_n|\le A(\mathsf{d};R)B(R)^n(n!)^{\varkappa_\star}\prod_i\|f_i\|_\sharp\,\,(\ell\ge\ell_G);}
\end{gather*}

Here $C^\infty_n(\mathsf{d};R)$ does not increase as $\mathsf{d}$ grows. The bounds (2) also hold, with no condition on a volume, for $S^T_n[\nu]$, where $\nu$ is an arbitrary periodic Gibbs state on $\mathbb{L}_K$ of Wilson's weight at $\beta$, and where $(K,\beta)$ is a first-passage datum.

\phantomsection\label{notdef:positive-test-pairs}
(3) The class $\mathfrak{T}_{\mathbb{E}}(\mathsf{d};\vartheta)$ consists of the pairs of Lipschitz functions $f,h\ge0$, neither of which is identically zero, that satisfy the following conditions. Each is supported in a ball of radius $\mathsf{d}$, and $\mathsf{d}\le\operatorname{dist}(\operatorname{supp}f,\operatorname{supp}h)\le6\mathsf{d}$. In addition, $\int f\ge\vartheta\|f\|_\infty\mathsf{d}^4$ and $\operatorname{Lip}f\le\|f\|_\infty/(\vartheta\mathsf{d})$, and the same two conditions hold for $h$. Assume $g\le\gamma_W(\vartheta)$, $\mathsf{d}\le\mathsf{d}_0$ and $(f,h)\in\mathfrak{T}_{T_\ell}(\mathsf{d};\vartheta)$. Assume also $K-s(\mathsf{d})\ge n_0$, so that at least $n_0$ steps lie between the scale $\mathsf{d}$ and the cutoff. Then

\begin{gather*}
\FitMath{\langle\mathcal{O}(f);\mathcal{O}(h)\rangle_{K,\ell,\beta}=\zeta_K(\mathsf{d})^2\,\tfrac{3(N^2-1)}{\pi^4}\,g(\mathsf{d})^4\,\mathcal{K}_{\mathrm{free}}(f,h)\,(1+\varepsilon),} \\[3pt]
\FitMath{|\varepsilon|\le\min\{\tfrac12,\bar\varepsilon_\ell(\mathsf{d})\}.}
\end{gather*}

In this formula $g(\mathsf{d})=g_{K-s(\mathsf{d})}(\beta)$, where $L^{-s(\mathsf{d})}\approx\mathsf{d}$ up to logarithmic factors. Also $\underline x_s\le g(\mathsf{d})^{-2}\le\bar x_s$, and both bounds are affine in $s=s(\mathsf{d})$. The factor $\zeta_K(\mathsf{d})=\zeta'_{K-s(\mathsf{d})}$ is the running normalisation, that is, it measures how $x_{K-s(\mathsf{d})}$ responds when $\beta$ is shifted by a constant \cref{sec:sources-analytic-dependence-on-the-source-locality-and-decay-in-tame}. One has $|\zeta_K(\mathsf{d})-1|\le\bar\eta(\mathsf{d})$ and $|\tfrac{11N^2}{12\pi^2}g(\mathsf{d})^2\log\tfrac1{\mathsf{d}}-1|\le\bar\eta(\mathsf{d})$. Here $\mathsf{d}_0\in(0,1)$, and the two functions $\bar\eta:(0,\mathsf{d}_0]\to[0,\frac13]$ and $\bar\varepsilon_\ell:(0,\mathsf{d}_0]\to[0,\frac12]$ are non-decreasing; moreover $\bar\eta,\bar\varepsilon_\ell\to0$ as $\mathsf{d}\to0$. Hence, on each torus, $\varepsilon\to0$ uniformly with respect to the cutoff, to the bare coupling as it ranges over its window, to the pair, and to where the pair is placed. For periodic Gibbs states the identity remains true, but only with $|\varepsilon|\le\tfrac12$.

(4) Consequently the limits of (a)(ii)--(iii) obey the law stated at the opening of the text, for those pairs in $\mathfrak{T}_{\mathbb{E}}(\mathsf{d};\vartheta)$ in which $f,h$ are of class $C^1$; the tuned $S^{\natural,\ell}$ are not included. The error satisfies $(1-\bar\varepsilon)(1-\bar\eta)^4\le1+\eta\le(1+\bar\varepsilon)(1+\bar\eta)^4$, with $\bar\varepsilon=\bar\varepsilon_\ell(\mathsf{d})$ on $T_\ell$ and $\bar\varepsilon=\tfrac12$ on $\mathbb{R}^4$, where in addition $\dim\mathcal{H}(S)\ge2$.

For clause (5), assume in addition the comparison
\eqref{eq:four-point-unresolved-interacting-target} for the actual normalized
history/block-field law, with $x=g(\mathsf d)^{-2}$ and
$\zeta=\zeta_K(\mathsf d)$.  On the whole prescribed four-test class,
its remainder must satisfy
\[
 |\mathcal R_4(\underline f)|
 \le\epsilon_{4,\ell}(\mathsf d)\,g(\mathsf d)^8
                         \prod_i\|f_i\|_\infty,
 \qquad\epsilon_{4,\ell}(\mathsf d)\longrightarrow0
            \quad(\mathsf d\downarrow0),
\]
for each fixed admissible torus, uniformly in the permitted cutoff,
bare coupling and support positions.  The comparison retains the actual
constraint, gauge-normalization and closed-mode identifications accompanying
\eqref{eq:four-point-background-projections}, the coarse-complement estimate
\eqref{eq:four-point-needed-complement-tail}, and the weighted lineage-age
estimate of \cref{prop:source-domains-4-14} at a rate strictly greater than
$16$, or a stronger direct estimate for these four insertions.  All the
nonquadratic, source, boundary and excluded-history terms listed after
\eqref{eq:four-point-unresolved-interacting-target} belong in
$\mathcal R_4$.  This additional input has the torus scope of clause (5).

\phantomsection\label{notdef:four-point-tests}
(5) The class $\mathfrak{T}^{(4)}_{\mathbb{E}}(\mathsf{d};\vartheta)$ consists of the quadruples whose pairs all belong to $\mathfrak{T}_{\mathbb{E}}(\mathsf{d};\vartheta)$ and for which $\operatorname{supp}f_i\subset B(c_i,\mathsf{d}/16)$, the four centres $c_i$ lying on one line. On the tori there exists $\vartheta_\square\in(0,1]$ with the following property. Let $\vartheta\le\vartheta_\square$ and $g\le\gamma_W(\vartheta)$. Then there are constants $0<c_\square\le C_\square$, which may depend on $g$ and $\vartheta$, such that the bound below holds whenever $K-s(\mathsf{d})\ge n_4$, where $n_4\ge n_0$ is an integer, $\mathsf{d}$ is smaller than a threshold, which may depend on $\ell$, and the quadruple belongs to $\mathfrak{T}^{(4)}_{T_\ell}(\mathsf{d};\vartheta)$:

\begin{gather*}
\FitMath{c_\square(N^2-1)\,g(\mathsf{d})^8\prod_i\|f_i\|_\infty\le S^T_{4;K,\ell,\beta}(f_1,\dots,f_4)\le C_\square(N^2-1)\,g(\mathsf{d})^8\prod_i\|f_i\|_\infty:}
\end{gather*}

on these quadruples the connected four-point function is therefore strictly positive. On rescaled quadruples it decreases as the fourth power of the reciprocal of the logarithm of the scale.

\textbf{(d) Fixed-torus uniqueness.} Assume $g\le\gamma_U$. Take $(\beta^\natural_K)$ and $S^{\natural,\ell}$ as in (a)(i).
(1) Let $\varrho$ be large and $g\le\gamma_{U'}$. Then the first-passage windows lie close to the tuned couplings: $\mathcal{W}_K(g)\subset[\beta^\natural_K-\tfrac14\varrho,\ \beta^\natural_K+\tfrac14\varrho]$ for each $K$.
(2) For $|\beta-\beta^\natural_K|\le\tfrac14\varrho$ the reference runs of length $K$ are well defined, and $\beta\mapsto x^\natural_K(\beta)$ is bi-Lipschitz uniformly in $K$. As $K\to\infty$, the functions $t\mapsto x^\natural_K(\beta^\natural_K-t)$ converge uniformly on $|t|\le\tfrac14\varrho$. The limit function has a compact interval $I$ as its range, and this interval contains $[g^{-2}-c'\varrho,g^{-2}+c'\varrho]$.
(3) For each torus of the family, one has a family $(S^{(\ell,x)})_{x\in I}$. Each of its members is a family of connected functions, defined on tuples of Lipschitz functions whose supports are pairwise disjoint. It satisfies $S^{(\ell,g^{-2})}=S^{\natural,\ell}$, and $x\mapsto S^{(\ell,x)}_n$ is continuous on $I$. Moreover, whenever $K_j\to\infty$, $|\beta_j-\beta^\natural_{K_j}|\le\tfrac14\varrho$ and $x^\natural_{K_j}(\beta_j)\to x$, it follows that $x\in I$ and $S^T_{n;K_j,\ell,\beta_j}\to S^{(\ell,x)}_n$, with no passage to a subsequence.

Consequently, take two sequences of bare data close to the tuned couplings whose renormalised couplings converge to the same value. They have the same limit on $T_\ell$. Under the conditions of (1) this applies to every sequence of first-passage data, and every limit of (a)(ii) is equal to some $S^{(\ell,x)}$. On each torus, the renormalised coupling by itself determines the limit theory.

\textbf{(e) Rotations.} Let $\omega\in\mathfrak{so}(4)$, and let $\chi$ be a cutoff whose support lies in $(-\ell/2+2a_K,\ell/2-2a_K)^4$. Consider the rotational Ward identity satisfied by $\mu_{K,\ell,\beta}$ with respect to the deformation field $\chi(y)\,\omega y$. The shell flux is the boundary term of this identity:

\begin{gather*}
\FitMath{\mathfrak{B}_{\omega,\chi}=\beta\sum_pa_p\,(\nabla^{\mathrm{s}}_\omega\chi)(x(p))-\beta\sum_p\mathfrak{h}_{\omega,\chi}(p)-\tfrac32\sum_pa_p\sum_{\rho\in\{\mu,\nu\}}(\omega x(p))_\rho\,(\nabla^{\mathrm{c}}_\rho\chi)(x(p)),}
\end{gather*}

where $\mu<\nu$ are the two directions spanning $p$ and $\nabla^{\mathrm{s}}_\omega,\nabla^{\mathrm{c}}_\rho$ are lattice differences, and where 

\begin{gather*}
\FitMath{\mathfrak{h}_{\omega,\chi}(p)=a_K^{-1}\sum_{b\subset\partial p}\sum_\rho[\chi(x_b)-\chi(x(p))]\,(\omega x_b)_\rho\,\mathfrak{d}_{b,p,\rho}}
\end{gather*}

, the $\mathfrak{d}_{b,p,\rho}$ being bounded, gauge-invariant local densities. Each term contains a difference of $\chi$.
For $\mathfrak{B}_{\omega,\mathcal{R}}:=\mathfrak{B}_{\omega,\chi_{\mathcal{R}}}$,
only terms whose plaquette anchors lie in
$\{\mathcal{R}-a\le|x(p)|\le2\mathcal{R}+a\}$ can contribute,
where $a=a_K$. The local densities also read neighboring plaquettes.
Their full field dependence lies in the enlarged shell
$\{\mathcal{R}-3a\le|x|\le2\mathcal{R}+3a\}$, as in
\eqref{wd:flux-energy-bound}. These shells surround the transition region
$\mathcal{R}\le|y|\le2\mathcal{R}$, where $\chi_{\mathcal{R}}$ falls from $1$ to $0$. A limit of (a)(iii) is $\operatorname{O}(4)$-invariant if and only if the following holds for one $\omega\ne0$ (equivalently, for all $\omega$), for all $n\ge2$ and for all separated tuples with smooth entries:

\begin{gather*}
\FitMath{\liminf_{\mathcal{R}\to\infty}\limsup_{j\to\infty}\big|\langle\mathfrak{B}_{\omega,\mathcal{R}};\mathcal{O}(f_1);\dots;\mathcal{O}(f_n)\rangle_{K_j,\ell_j,\beta_j}\big|=0.}
\end{gather*}

In that case $\mathcal{H}(S)$ carries a strongly continuous unitary representation of $\operatorname{O}(3)$. This representation leaves $\Omega$ fixed, commutes with $H$, and extends the representation of $W_3$. This equivalence singles out the one infrared input, and for volume-first limits \cref{thm:phase-main} settles it.
\end{theorem}

\emph{What the equivalence uses.} When $\chi=1$ near the supports, the identity takes the following form (the exact lattice Ward hypothesis in \eqref{eq:comparison-capsule-CP10}):

\begin{gather*}
\FitMath{\sum_{i=1}^n\big\langle\mathcal{O}(f_1);\dots;\mathcal{O}(\nabla^{\mathrm{s}}_\omega f_i);\dots;\mathcal{O}(f_n)\big\rangle_{K,\ell,\beta}=\mathfrak{U}_{\omega,\chi}(f_1,\dots,f_n)} \\[3pt]
\FitMath{+\big\langle\mathfrak{B}_{\omega,\chi};\mathcal{O}(f_1);\dots;\mathcal{O}(f_n)\big\rangle_{K,\ell,\beta}:}
\end{gather*}

besides the flux, it contains a bulk term $\mathfrak{U}_{\omega,\chi}$. This term is the correlation between the $\mathcal{O}(f_i)$ and lattice artefacts of dimension six, distributed over $\operatorname{supp}\chi$. Once $\mathfrak{U}_{\omega,\chi_{\mathcal{R}}}\to0$ along every defining sequence, the rotation defect of $S$, namely the derivative $\tfrac{d}{dt}S_n(e^{t\omega}f_1,\dots,e^{t\omega}f_n)|_{t=0}$, equals minus the limit of the correlation between the flux and the $\mathcal{O}(f_i)$, and clause (e) follows.

This vanishing requires the typed first-order insertion estimates, the joint-insertion moment bound, and the volume-transfer estimate in \hyperref[hyp:ward-bulk]{the Ward hypotheses}. The first-order estimates are needed on fixed tori and slowly growing volumes; the joint-insertion bound and volume transfer are needed for arbitrary sequences with $\ell_j\to\infty$. On this route the anisotropy estimate also requires the exponent stated next.

One marginal operator remains in $\mathfrak{U}$: a traceless diagonal anisotropy of Wilson's weight. Its coefficient is of the size $g_0^{-1}\approx K^{1/2}$, and it disappears only through its own decay $K^{-c_1}$ along the flow. What is needed is $c_1>\tfrac12$.

This number is expressed through one-loop data. Put the weight $1+\mathsf{e}_{\pi(p)}$ on each $a_p$, where $\pi(p)$ is the coordinate plane of $p$. Differentiating the one-loop coefficient of step $l$ plane by plane at $\mathsf{e}=0$ gives a derivative that commutes with $W_4$. Let $\mathsf{m}(l)$ denote the eigenvalue of this derivative on the traceless diagonal weights $\mathsf{e}_{\{\mu,\nu\}}=\mathsf{e}_\mu+\mathsf{e}_\nu$, $\sum_\mu\mathsf{e}_\mu=0$. Then $c_1\in(\tfrac12,1]$ is an exponent such that $\sum_{i<l\le j}\mathsf{m}(l)\ge c_1c_0(j-i)-C$. The one-loop law in this channel, $|\sum_{l\le k}\mathsf{m}(l)-c_0k|\le C$, gives $c_1=1$.

The law comes from a telescoping sum. By the law of total covariance, the sum of the one-loop responses of the steps $l\le k$ is a single trace on the bare lattice. The logarithmic part of this trace is $\int_1^{L^{2k}}d\sigma/\sigma=2k\log L$, in proper time $\sigma$, multiplied by the heat coefficient on a constant abelian background. It equals $kc_0$, with the same $c_0$ as the coupling itself. In the continuum, an anisotropy of the weight amounts, to first order, to a change of coordinates, and it does not run (see \hyperref[hyp:ward-bulk]{the Ward comparison hypotheses}). Rotation invariance in \cref{thm:phase-main} (1) follows from \cref{thm:all-orders-angular-closure}, whose proof uses no Ward identity and does not depend on $c_1$.

\begin{theorem}[Canonical components]\label{thm:phase-main}
Assume the background, interaction, history and flow inputs of
\cref{ch:dependency-catalogue}, the ordinary-source hypothesis D1 in
its stated rooted-domain geometry, and the interacting two-point and
running-normalization inputs F2--F3 there.  The latter inputs are required
uniformly over the periodic Gibbs states used to realize the limit and
its components, on the full test class specified in
\cref{thm:angular-two-point-elimination}, including
$\vartheta_*=1/12000$.
Let $g\le\min\{\gamma_W(\vartheta),\gamma_W(\vartheta_0)\}$, and let $S$ be a volume-first limit with defining sequence $(K_j,\beta_j,\nu_j)$.
(1) Every $\nu_j$ is a periodic Gibbs state and is invariant under all translations of the lattice; moreover $S$ is a limit of (a)(iii). In addition to (b) and (c)(4), it satisfies $H\ge|P_j|$ for $j=1,2,3$, is $\operatorname{O}(4)$-invariant, and is reflection positive at the diagonal mirrors.
(2) Its canonical decomposition exists and is unique up to isomorphism. Here an isomorphism is a homeomorphism of the spaces $\mathsf{X}$ that carries $\lambda$ and the $\sigma_\xi$ along with it. Let $P_0$ be the projection onto $\ker H$. The map $P_0[F]\mapsto(\xi\mapsto\sigma_\xi(F))$ extends to a unitary map $\ker H\to L^2(\mathsf{X},\lambda)$ with $\Omega\mapsto1$, by the zero-energy representation theorem of \cref{ch:phase-vacuum}. If $\mathsf X$ has exactly $m<\infty$ points, then $\dim\ker H=m$; if $\mathsf X$ is infinite, then $\ker H$ has countably infinite Hilbert dimension. This does not identify the set-theoretic cardinality of an infinite compact phase space with its Hilbert dimension. The components are permuted, and $\lambda$ is preserved, by the isometries that leave $S$ invariant (translations, $W_4$, rotations).
(3) Each $\sigma_\xi$ has the following properties.
(i) It has the same invariances and positivity properties, and it is symmetric and regular.
(ii) Its one-point function is $S^{\sigma_\xi}_1(f)=\mathfrak{c}_\xi\int f$, where $\mathfrak{c}_\xi$ is real and continuous in $\xi$ and $\int\mathfrak{c}_\xi\,\lambda(d\xi)=0$. With the same constants, it satisfies the growth condition and the bounds (c)(2); it also satisfies the law (c)(4) in its form on $\mathbb{R}^4$.
(iii) On $\mathcal{H}(\sigma_\xi)$ it has the operators of (b)(v) and a unique vacuum, $\ker H_{\sigma_\xi}=\mathbb{C}\Omega_{\sigma_\xi}$. Its joint spectrum lies in the closed forward cone $\{E\ge|q|\}$. It also has a strongly continuous unitary representation of $\operatorname{O}(3)$ that leaves $\Omega_{\sigma_\xi}$ fixed and commutes with $H_{\sigma_\xi}$.
(iv) It has the cluster property:

\begin{gather*}
\FitMath{|\sigma_\xi(X\cdot\tau_vY)-\sigma_\xi(X)\,\sigma_\xi(Y)|\le C(X,Y)\,|v|^{-2}} \\[3pt]
\FitMath{(|v|\ge r(X,Y)),}
\end{gather*}

where $C(X,Y)$ and $r(X,Y)$ are the same for all $\xi$ and for all volume-first limits at $g$.
(4) Consequently every Osterwalder--Schrader axiom holds for $\sigma_\xi$: growth, invariance under $\mathbb{R}^4\rtimes\operatorname{O}(4)$, reflection positivity, symmetry and clustering. Its Schwinger functions coincide with the Schwinger functions of a Wightman theory. This theory has a single Hermitian scalar field on a separable Hilbert space and a unique vacuum. It is covariant under the proper orthochronous Poincar{\'e} group, its spectrum lies in the closed forward cone, and it is local.
\end{theorem}

% END MODULE chapters/01_orientation.tex

% BEGIN MODULE chapters/dependency_catalogue.tex
\subsection{Technical lemmas and proof dependencies}
\label{ch:dependency-catalogue}
\subsubsection{Lemmas and dependencies of the proof outline}
This catalogue gives the technical statements and dependencies of the
author's proof outline. It fixes the full domains, quantifiers and uniform
constants of each lemma invoked in the argument. The accompanying sections
expand some proofs; other technical statements are quoted here without
their detailed proofs. A reference to one of these statements always
retains its complete hypotheses and only the conclusions stated below.
The catalogue identifies where the principal multiscale estimates enter.
Those estimates require proofs beyond the finite-dimensional and reference
calculations given here.

Different routes use different groups of hypotheses. On an existing sequence of arbitrarily
fine admissible data, the scalar route uses D1 with its actual rooted
geometry, G1, and the ultraviolet inputs F2--F3. The bridge sections
derive compactness, D2, H1, compact Hilbert evolution and G2 from those
data. The joint induction outlined in A--C and the flow section is responsible
for constructing these data with the stated uniformity. Weighted and
matched-run comparisons and typed Ward estimates remain requirements for
the stronger statements that invoke them; the scalar route requires only its ordinary D1 and ultraviolet premises,
together with the consequences derived from them.

\phantomsection\label{notdef:weighted-action}
Here are conventions shared by the packages. For numerical coefficients written in terms of $N$, the group is $\operatorname{SU}(N)$ and here $N=2$. Fix an odd integer $L>1$, a compact
simple matrix group $G$, and dimension four. The principal example is
$G=\operatorname{SU}(2)$. The lattice spacing is $a_K=L^{-K}$, and a finite torus has
physical side $\ell$. There is one group variable $U_b$ per positively oriented
bond, $U_{b^{-1}}=U_b^{-1}$, and normalized product Haar measure $dU$.
Write

\begin{gather*}
\FitMath{a_p(U)=1-\tfrac12\operatorname{Re}\operatorname{tr}U_p,} \\[3pt]
\FitMath{A(w,U)=\sum_p w(x(p))a_p(U),} \\[3pt]
\FitMath{Z_{K,\ell}(w)=\int e^{-\beta A(1,U)+A(w,U)}dU.}
\end{gather*}

The trace convention in this display is for $\operatorname{SU}(2)$; another group
requires its own normalization. Let $\mathcal O_K(f)=A(f,U)$. The connected
functions are the mixed derivatives at zero of
$\log Z_{K,\ell}(\sum_i z_if_i)$. Test functions used together have pairwise
disjoint compact supports, separation at least $d>0$, and individual supports in
balls of radius $R$. Set
$\|f\|_\sharp=\max(\|f\|_\infty,c_M\|\nabla f\|_\infty)$, with $c_M$ fixed.
All uniform constants below must be independent of $K$, the torus, the history,
the positions of the tests, and the relevant Gibbs state. Dependence on
$g,n,d,R$ must be displayed. A lower bound on admissible $\ell$ may depend on
$g$ and on constants fixed earlier. Replacing this uniformity by a bound at one
fixed cutoff does not satisfy an input.

\paragraph*{A. Background fields and propagators: four inputs}

\noindent\textbf{Technical statement A1: the constrained critical orbit.}
For every allowed nested family of small-field regions and every regular coarse
datum $V$, there must exist a regular critical orbit $[U(V)]$ for the Wilson
action with the prescribed block averages. It is unique among regular critical
orbits in a specified larger neighborhood. A chosen gauge gives a representative
which is Borel in all real data and holomorphic in a complex neighborhood of each
regular real datum. Both neighborhoods have stated radii on the local scale;
their radii and constants are uniform in the number of scales. Existence of an
unrestricted global minimizer is insufficient: the needed assertion concerns
the constrained regular orbit and its gauge representative. Compatibility with
restriction to subregions, with seam constraints, and with changing the torus
must be part of the statement. This is the background-orbit input behind the
local assertions numbered \cref{prop:averaging-1a-3-6}, \cref{prop:averaging-1a-3-8} and the background theorems.
For the pointed window deduction, the actual restrictions and substitutions
must remain in their prescribed pair cores. These cores contain the usable
Taylor chart ball and obey the curvature and own tent-action bounds on the
whole admitted core, including the finitely many young ages.
\Cref{ch:block-average-pivot-chart} proves a local chart for
the actual one-step average near a pure-gauge double block, including
the pivot derivative, exact straightening and measure factors. Its
explicit chart radius and nonwrapping premise do not imply existence,
uniformity or all-substitution room for the many-level critical orbit.


\paragraph{The regular proper background.}
\Cref{ch:first-step-background-expanded,ch:first-step-nonflat-precision}
realize the first physical blocking step at fixed $L$, with arbitrary
admitted real flat holonomies. For the separated determining sets
\eqref{mb:determining-rows}, which retain interior and outward rows
and omit inward rows, \cref{bgc:proper-background} proves existence
and uniqueness among all regular critical orbits. Its proof uses the
actual proper strong inverse, forced trial correction, datum descent,
and a reset of curvature and current after every ascent.
\Cref{bgc:global-comparison} places arbitrary regular competitors in
one actual Landau chart; \cref{bgc:strong-comparison} obtains the
required strong norm from their actual currents.
\Cref{bgc:uniform-physical-chart} then supplies a fixed physical
holomorphic continuation ball at every such real datum, uniformly
in depth and volume. The exact comb construction in
\cref{ch:multilevel-backgrounds} selects its Borel real representative.
The larger history input A1 also specifies its actual restriction,
substitution and core requirements; those clauses retain their
separate hypotheses and proofs.

\medskip\noindent\textbf{Technical statement A2: coercivity on a fluctuation fiber.}
Let $Q$ be the derivative of the chosen block average, and let $\mathcal E$ be
the gauge-fixed Hessian on the fluctuation variables at a regular background.
The chosen coordinates must remove all residual gauge directions. On each
conditioning set and with its specified boundary conditions require

\begin{gather*}
\FitMath{\operatorname{Re}\langle B,\mathcal E B\rangle
 \ge c\|B\|_2^2,} \\[3pt]
\FitMath{Q B=0,}
\end{gather*}

with the norm in the units of that fluctuation scale. The same estimate, with
$c/2$, must hold throughout the claimed complex neighborhoods. Zero modes,
torus holonomies, and node constraints must either be absent from this space or
treated separately. A mass for this constrained operator is not a mass gap for
the physical continuum Hamiltonian. This is the quantitative content required by
the conditioning and covariance assertions numbered 1c.6 and related claims.
For the actual unweighted classical flat reference,
\cref{ch:flat-contour-precision-gap,ch:flat-whitney-comparison} prove
\(c(L)I\le\mathsf A_k^{\rm flat}\le CL^8I\), uniformly in the
number of previous classical blocking steps and in compatible tori,
including real flat holonomies. They use the actual plain constraint
and the soft internal-contour gauge. The curl-energy comparison is
\(\|\mathrm{d}a\|^2\le E_k(a)\le256\|\mathrm{d}a\|^2\).
These are the flat classical estimates used in the nonflat comparisons below;
their statements retain that reference-case scope.


\paragraph{The physical first step and a specified current extension.}
The nonflat first-step precision theorem supplies positivity
for the actual Wilson action with its soft gauge term on the
curved constraint fibre, at an admitted real block-axial trial
background.  Its radius and current ceilings are retained.
\Cref{ch:explicit-relative-current-extension} then defines
one full relative auxiliary-current extension, identifies its
physical action at a critical physical pair, and proves the
complete first-step half-disc estimate with the gauge part
included.  This chosen definition must be used consistently
when invoking that result.  Neither zero coarse readout nor a
Ward identity identifies a separately unspecified completion.
For the canonical proper precision \eqref{pgc:actual-precision},
\cref{pgc:proper-covariance} proves uniform real coercivity and
holomorphic invertibility on a common weighted $C^1$ complex
background tube and independent-current ball. The estimate includes
the outer second-jet correction and relative gauge, is uniform in
the proper inner hierarchy and conditioning set, and also covers
the stated strict tree-zero subspace with the gauge term omitted.
The local inverse, logarithm and normalized Gaussian families are
constructed in \cref{pgc:section}. Invoking these results requires
the actual precision, readouts and conditioning coordinates to be
those specified there, with their stated regular-pair ceilings.
Signed profiles, counterterm-modified precisions and interacting Hessians
require their own comparisons on the joint history cores.

\paragraph{Layered site and projected-form reductions.}
\Cref{ch:layered-current-precision} proves the site intertwining,
site inverse/Green/Gram and projected physical form/readout reductions
within its separated geometry and regular-pair class, with field-row/Feynman
and depth-uniform density-coefficient comparisons. Verifying that every
actual background and history substitution lies in this class remains
necessary. The fine Landau coercivity and its complex moving-slice
continuation are proved in
\cref{thm:npl-proper-gap,thm:cml-moving-slice}; the specified current
completion is identified in \cref{prop:acc-current-identification}.
The named propagator and decoupled families are supplied by
\cref{prop:ncw-named-families,dnr:linear-families,dnr:parameter-families}.
Each result retains its own geometry, complete-read and regular-pair
hypotheses. The six-propagator normalizations are recorded in
\cref{candidate:propagators}.

\medskip\noindent\textbf{Technical statement A3: local inverse and response bounds.}
The inverses and response operators actually used, including conditional
covariances and derivatives of $U(V)$, must obey their declared pointwise
and source-norm bounds with exponential decay in scale units. The shorthand
$C e^{-c\operatorname{dist}(x,y)}$ applies after the local source and output
normalizations have been specified. In \cref{candidate:propagators} the
kernels use counting measure: the fine Green source estimates retain
the factors $[\xi_y^2,\xi_y,\xi_y,1]$, the inverse site Gram entries
retain $\xi_y^{-4}$, and the inverse coarse field-row Schur entries
retain $\xi_c^{-2}$. The stated short-distance, first-difference and
H\"older factors must also be kept. For a fine-to-fine Green map, its
density against physical fine volume is $\eta^{-4}$ times its counting
kernel; a counting-entry bound is not the same density bound. Coarse row
operators instead use their own row weights and rescalings, as in
\eqref{ncw:normalized-field-row}. Weighted cell bounds in norms such as
\eqref{scf:force-cell-norm} are operator estimates, not pointwise density
estimates. All required bounds must hold throughout the full complex
datum polydisc, not only at the identity or at real data.
For the pointed window this is the product bond-log domain \eqref{eq:pointed-window-remainder-PW1}, with
the scale-weighted potential, gradient and current response in \eqref{eq:pointed-window-remainder-PW2}--\eqref{eq:pointed-window-remainder-PW3}.
For a region with
several local scales, use a precisely defined weighted path distance and the
corresponding powers of the local spacings. Holomorphic derivatives in each
decoupling parameter must satisfy Cauchy bounds on a fixed enlarged complex
disc. A derivative joining distant sets must retain the decay needed by the
polymer norm. The mesh-uniform scalar estimate and its Dirichlet
compressions are proved in the block-resolvent section.
On the admitted separated regular-pair domains,
\cref{dnr:linear-families,dnr:parameter-families,dnr:response,dnr:walls}
prove the non-Abelian decoupled response, its screened field and current
bounds, holomorphy, and wall decomposition. The all-one physical
identification and coupled response are proved in
\cref{dnr:physical,dnr:coupled}; \cref{dnr:interpolation} identifies the
prescribed response maps. This supplies the assertion numbered 1c.10
on those domains. Each application must still verify that its actual
history data, complete determining reads and substitution lie in the
stated domains with the required radius margins.
\Cref{ch:flat-exponential-kernels} proves exponential kernels
for the actual unweighted classical flat precision and its constrained
covariance, uniformly in depth and compatible volume. It derives the
exact alias symbol, cancels its apparent gauge denominator on the full
complex zero cone, proves a uniform strip, and then inverts using the
flat gap. Periodization and unitary flat holonomies preserve the bounds.
This establishes the stated flat reference kernels, not the nonlinear
response or decoupling bounds required on the full interacting domains.
The finite-box affine comparison in
\eqref{eq:pointed-window-remainder-PW12} prescribes the
affine test field's exact averages on the same determining
set, including the seam rows. The weighted first response
controls the effect of discarding its remote rows.


\paragraph{Primitive inversion versus nonlinear response.}
\Cref{ch:nonlinear-response-reductions} gives the exact
coupled field--multiplier inverse, with the true datum
differential, weighted transposes and constrained Hessian.
It proves the response and complex-ball deductions from stated
two-sided primitive inverse bounds.  Boundary forcing is kept
in its own norm.  \Cref{ch:discrete-holder-estimate}
proves mesh-uniform local scalar estimates and their
flat-unitary covariant version.  Applying them to a nonlinear
response requires the actual connection, forcing and covariant
Laplacian point bounds recorded there.  The nonlinear application uses those primitive equations and bounds
for the actual response on every allowed core. Once they hold, the
H\"older conclusion follows from the displayed argument; no further
H\"older lemma is invoked.

\medskip\noindent\textbf{Technical statement A4: averaging schemes and locality.}
Any transfer from one averaging scheme to another must provide a uniform
intertwining map, including its coarse gauge factor. On fluctuation fibers the
transferred Hessians, inverses and constraints must obey A1--A3 with the same
allowed losses, for every number of blocking steps. Restrictions to a region
must depend only on an explicitly bounded enlargement of that region. Equality
of two linearized averages only at the identity is insufficient. A proof must
bound the accumulation of their differences at nonflat backgrounds and must
state the gauge in which the comparison is made.


\paragraph{Exact restriction with its boundary price.}
The finite factorization in
\cref{ch:multilevel-backgrounds} shows how a new
constrained variation fixes the old data and why restriction
preserves criticality.  Cutting the plaquette action can
change its boundary current.  The changed bonds are in the
new bare boundary layer, where the local scale is \(\eta\);
the at-most-six-plaquette estimate pays a fixed current
constant.  That constant and the datum/seam losses are
included in the actual uniqueness-input ceiling.  Under this
admission and regular critical uniqueness, the restriction
identity is a deduction.  More general changes of averaging
scheme still need the stated physical and complex-domain
intertwining estimates; the restriction proof alone does not
supply such a comparison.

\paragraph*{B. Cluster expansion and power counting: four requirements}

\noindent\textbf{Requirement B1: a conditional polymer majorant.}
Let polymers be finite connected sets of scale cubes, with incompatibility
meaning intersecting supports. A specified positive size $a(X)$ and activities
$z(X)$ must satisfy, uniformly in all external data,

\begin{gather*}
\FitMath{\sum_{Y:Y\not\sim X}|z(Y)|e^{a(Y)}\le a(X).}
\end{gather*}

The same inequality, with strict room, must hold for complex sources and
decoupling parameters on the domains used for differentiation. Require the
absolutely convergent connected-cluster formula for the logarithm, including
termwise differentiation and a uniform majorant for each derivative used later.
The finite-volume identity must be exact before taking limits. This is the
content needed from the convergence statement numbered 1d.4.1; a formal graph
expansion is not sufficient.
The unmarked polymer part of \cref{ch:local-formation-program}
derives the local unmarked activity and union-filed logarithm bounds
from its actual enlarged
Gaussian matrix/core hypotheses, sharp-cutoff structure and local
potential budgets. The Mayer-parameter discs are established at the
Gaussian integral, before zero-freeness is inferred. Gaussian volume
costs, failed bonds, actual neighbor-layer closures and the spare
potential slot are all paid for in the rate allocation. This also
supplies \eqref{eq:local-formation-SI5} under those primitive hypotheses. Their verification for
the actual complex interacting matrices and all allowed cores remains
necessary; convergence of an unrelated model gas would not supply it.

\medskip\noindent\textbf{Requirement B2: conditional pointed marginal extraction.}
\cref{ch:pointed-hessian-cancellation,ch:pointed-window-remainder} prove the following implication
for an actual regular member with no stored color spectators.
Its pointed total and differentiated sums must converge normally;
it is locally gauge invariant, and the individual member's pointed
totals are covariant under unit translations and the full hypercubic
group. Use the actual free response, with gauge covariance and affine
constant-curvature reproduction modulo gauge, with zero linear current.
The usable pair core, localized response and entropy gap must give
the absolute Hessian moments in \eqref{ph:exponential-moment}.
Then \eqref{ph:extractor} defines its coefficient and
\eqref{ph:own-counterterm} proves cancellation by that same member's
own midpoint tent action. Different members may be chosen at different
sites; their coefficients cannot be interchanged.

On physical radius \(a\), the readout retains
\(|b(E)|+\|T[E]\|\le Ca^{-2}\|E\|\).
\Cref{ch:extraction-radius-normalization} supplies its stated
near-zero inclusion and normalization. For the nonlinear remainder
also require A1--A3 throughout the complex datum domain, the actual
core inclusions, spatial H\"older regularity, compatible boundary
equations and the whole-core bound on the own tent action.
Put \(\sigma=L^{-(n-j)}\), \(\rho=\max(1,a_{\rm out}/a)\) and
\(0<\alpha_H\le1\). Under the window domain conditions the
near-domain sum at one base point, after its vacuum and own
counterterm subtraction, has bound

\begin{gather*}
\FitMath{C\|E\|\left[
 \rho^2\sigma^{4+\alpha_H}+\rho^3\sigma^6
 +(1+\rho^2\sigma^4)e^{-c/\sigma}\right].}
\end{gather*}

The final exponential term includes remote domains without a
nonwrapping chart. No factor proportional to the actual curvature
is asserted for that tail. Under the ordinary age bound
\(\rho^2\le C(1+n-j)\), the finite young-age core fallback and the
\(O(\sigma^{-4})\) site count give the summable weight
\(C\|E\|(1+n-j)^{3/2}L^{-\alpha_H(n-j)}\).

These are conditional deductions, not verification that all actual
RG formed members have the required hypotheses. They do not establish
arbitrary torus-sensitive or field-sensitive tails, source derivatives
of the formation map, or the fluctuation disintegration laws in F2.
Stronger assertions in 1e.1.13, 1e.1.15 and 1e.1.16 require
separate estimates.

\medskip\noindent\textbf{Technical statement B3: the one-loop law with a summable lattice correction and controlled nonlinear errors.}
For the chosen normalization let
$c_0=11N^2\log L/(12\pi^2)$. The Gaussian contribution $b_j^0$ must be
independent of the bare coupling and obey
$|b_j^0-c_0|\le C L^{-\sigma j}$, or a stated summable substitute. The full
increment must have a decomposition

\begin{gather*}
\FitMath{b_j=b_j^0+f_j+q_j,} \\[3pt]
\FitMath{|f_j|\le Cx_{j-1}^{-1/2}(\log x_{j-1})^p,} \\[3pt]
\FitMath{|q_j|\le Cx_{j-1}^{-(\kappa_0-4)/2},} \\[3pt]
\FitMath{\kappa_0>4.}
\end{gather*}

If a later deduction needs a stronger error, that stronger estimate must be
proved separately. The coefficient $c_0$ requires an actual lattice trace
calculation with the stated zero-mode convention. A continuum perturbative
coefficient quoted without this comparison does not prove this input.

\phantomsection\label{hyp:B4}
\medskip\noindent\textbf{Technical statement B4: the ordinary step and its finer derivatives.}
The ordinary source box allows a separate Euclidean regular member at each base
point, its own scalar counterterm, a variable Wilson weight, and scale-small
noncovariant defects. The step must be holomorphic on this product box with a
majorant independent of the number of labels. It is not necessary that the
whole spatial array be translation invariant. The ordinary regular-output
estimate does not by itself prove every marked or large-field derivative.

For the stronger comparisons actually invoked, require the potential derivative
bound in \eqref{eq:wext-potential-test} and the one-zone complex-mass bound in \eqref{eq:wext-zone-test}, on the domains stated in that section, with constants
independent of the levels and source positions. The regular irrelevance weight,
squared-radius weights for defects, and large-field survival weights have
different roles. In particular the survival factor belongs in the bound on a
zone's output; it does not enlarge its input radius by its reciprocal. Joint
marked derivatives and typed tensor counter-insertions must satisfy their
corresponding estimates. For the stronger routes these finer conditions are stated technical
lemmas; entire dependence of the bare partition function alone does
not imply them.
They are required for the retained stronger comparison and Ward routes.
The ordinary-source insertion and angular proofs below do not obtain their
operator estimates by assuming these typed Ward derivatives.
\Cref{ch:ordinary-source-formation-audit} proves ordinary
old-level summations, scalar action-load and reserve implications.
The source-insertion and renewed-pointing parts of
\cref{ch:local-formation-program} derive the local first response
under their local Gaussian and bond-split hypotheses; its unmarked
polymer part derives the needed majorant from the stronger primitive
matrix/core and potential package. Its window-formation part derives
the ordinary formation bound using one near window member at each old
base point before its bracket is applied. The linear bootstrap retains
the source norm and gradient gain under these inputs, the constant-source
slopes and the old-remainder bounds.
Actual nonlinear background/core inclusions and history formation
remain to be verified; these implications do not supply them.

\paragraph*{C. Large-field histories and closed induction: four inputs}

\phantomsection\label{hyp:C1}
\noindent\textbf{Technical statement C1: exact measurable operations.}
Every move must be a specified Borel kernel, change of variables, or partition
of unity on an explicitly defined space. The tests are total, including at
configurations where logarithms or minimizers fail to exist. Every stored
variable is integrated exactly once or retained exactly once. For every bounded
gauge-invariant insertion $F$, the sourced construction must preserve the total
integral of $F e^{-\beta A}$. An operation preserving only the integral without
$F$ cannot be used to transport that observable. For real inputs every summand
used as a probability weight must be nonnegative.

\phantomsection\label{hyp:C2}
\medskip\noindent\textbf{Technical statement C2: geometric margins and determination.}
The allowed histories must have a constructive definition. At level $j$, tests,
local backgrounds and retained terms depend on at most a stated enlargement of
the determining regions. Every enlargement lies inside the margin supplied by
the next region. Healing, restriction, joining and continuation must preserve
these properties, including near holes, seams, and waiting windows. On an
admissible torus each cell, each box on which the operation acts, and
its prescribed bounded neighborhood admit a common lift. Use the
companion identities on non-wrapping domains in matching step units.
Other carriers enter the torus estimates with their winding retained.
The number of histories with a specified size must have a proved entropy bound,
with the same size function used in the rarity estimate.

\phantomsection\label{hyp:C3}
\medskip\noindent\textbf{Technical statement C3: rarity with entropy and insertions paid for.}
For each created large-field component $X$ at level $j$, require a mass bound
$\mathfrak M_j(X)\le \exp[-c p_0(g_j)-\kappa |X|_j]$ with the stated
supremum over external complex data inside its internal integral when
component majorants are multiplied.
Here $p_0(g)=(\log g^{-2})^p$, and $\kappa$ exceeds the entropy rate of C2.
The bound must include the Wilson factor actually available to pay for that
component, all stored Jacobians and gauge factors, and the designated insertion
weights. A failure of existence of a background, as opposed to a quantitative
large-field inequality, must also be covered. Drop-outs without their own
Wilson factor need a separate bound after attachment to their carrying term.

\medskip\noindent\textbf{Technical statement C4: invariance of the inductive class.}
There must be one nonempty class $\mathfrak C_j$ of complete presentations, with
the same constants at every level, containing the bare measure. Under the
increment and tube conditions, the exact step maps
$\mathfrak C_j$ into $\mathfrak C_{j+1}$. Its conclusions include the analytic
polymer bounds, C1--C3, boundary bounds, covariance, and the required stability
floor. All constants must admit one simultaneous order of choice, without a
parameter needing to be both smaller and larger than quantities chosen later.
The step assertion 2.11 requires all these properties; each must be proved.

\paragraph*{D. Sources and local correlation bounds: four requirements}

\noindent\textbf{Technical statement D1: the ordinary localized source bound and stability.}
Require the precise torus hypotheses in \eqref{eq:all-state-source-transfer-ST5}--\eqref{eq:all-state-source-transfer-ST6} of
\cref{ch:all-state-source-transfer}: holomorphic point counterterms
and localized source energies, a strict entropy gap, source dependence only
on the specified regions, and independence of the representation of a
source function by a list of tests. On the ordinary source ball, require
the normalized upper stability bound and unsourced partition-function
floor there, with constants uniform in cutoff and allowed torus. The
source class must include the global mirror-periodized Lipschitz profiles
used in the chessboard construction. Fixed-cutoff bounds for point
counterterms may depend on the cutoff but must be uniform in the torus.

For the graded deduction, this input also specifies the actual domain
indexing of \cref{ch:graded-rooted-source}. Take $M=L^{m_0}$ with
fixed integer $m_0\ge0$, and the cofinal compatible tori $\ell=2L^m$,
$m\ge m_0$. At scale $a_s=L^{-s}$ the localization cubes have side
$Ma_s$; their counts in every direction are even integers. Every fixed
power-of-$L$ chessboard tiles all sufficiently large members of this family.
All actual scales $0\le s\le K$ are retained. A domain $X$ is indexed by
its union of these cubes, with $m(X)$ the minimum number of edges of a
finite connected cubical edge tree meeting every cube. Its source
dependence is contained in a fixed number of neighboring cube layers.
The source-energy bound uses $e^{-\kappa m(X)}$, with $\kappa$ above the
rooted entropy rate for this actual indexing. If extra domain types are
present, their multiplicity must be included in that stated entropy gap.

The closed-walk count in \eqref{eq:graded-source-bound-GS1}, the ball-root count in \eqref{eq:graded-source-bound-GS2}, and the global
weighted count in \eqref{eq:graded-source-bound-GS3} are then proved geometric consequences. In particular,
the sharp local count is $C(1+u/a_s)^4e^{-ct}$, and the global count is
$C\ell^4a_s^{-4}$. The weaker abstract estimate in \eqref{eq:all-state-source-transfer-ST6} alone does not imply
these sharper bounds and is insufficient for the graded conclusion below.
Checking that the localized RG energies use this specified geometry is
part of D1, not supplied by writing an abstract domain sum.

These are raw partition-function bounds. A zero-free logarithm on the
whole global source ball is not assumed. The fixed-volume normalized
zero-free lemma gives a suitable neighborhood when needed; the all-state
transfer theorem instead constructs common holomorphic cell normalizers
and obtains the derivatives at zero directly. The elementary transfer does
not prove the ordinary RG source estimates in this input.
\Cref{ch:ordinary-history-stability} proves explicit sufficient
criteria for the upper bound and floor from the normalized positive
label/history representation: total-variation bounds with a weighted
label sum, and a regular all-small-field label with its own action
bound. Those criteria, including the reserve and allocation of Wilson factors,
remain to be verified for the construction. A pointwise
polymer logarithm cannot replace that history representation without
an exact common integral formula.

% Revision4 promotion follows the completed independent QF review.

\paragraph{The actual history norm and the remaining local checks.}
\Cref{ch:history-mass-program} proves the weighted
genealogy count, reserve arithmetic, gauge-normalizer
half-disc bound, strong-mass induction and loaded label
sum from their explicit finite-operation hypotheses.
Its norm integrates the supremum over all admitted reads;
a bound at each fixed read is insufficient.  To apply the
result one must identify the literal stored-coordinate
operation and its fixed reference measure, admit every
later read to the same core, prove the complete quadratic
domination and localized energy envelope, and construct
owned witnesses for the failed tests.  The original inner
Wilson weight is retained.  The source extension has two
proved routes: an actual pointwise excess bound, or an ordinary kernel
holomorphic and bounded in its Banach norm on a complex source ball
over a fixed reference measure, together with a real estimate with a sufficiently large prepaid
mass-weight coefficient on a full real source ball and one common actual
field/read core for both estimates. Both routes retain their specified
formation and domain premises.  The lower
partition floor also needs its admitted positive regular
label and common normalization; an upper history-mass
estimate does not prove that floor.

\medskip\noindent\textbf{Requirement D2: graded all-state bounds derived from D1 and G1.}
\cref{thm:graded-rooted-source} uses D1 in its actual rooted setting
above and the proved chessboard domination from G1. It yields one fixed
finite integer $k$ and $B(R)=C(R+1)^4$ such that, for periodic Gibbs states
$\rho,\nu$ at the same datum,

\begin{gather*}
\FitMath{\left|\rho\prod_{i=1}^n
    [\mathcal O_K(f_i)-\nu(\mathcal O_K(f_i))]\right|
 \le B(R)^n(n!)^k\max\{1,d^{-5n}\}\prod_i\|f_i\|_\sharp.}
\end{gather*}

For $n\ge2$ the cumulants under $\rho$ satisfy the same estimate with
$[2B(R)]^n$ in place of $B(R)^n$. Each support lies in its own radius-$R$
ball, placed anywhere, and different supports have distance at least $d$.
The first-order case uses the one-point difference bound. All constants
are uniform in cutoff, states and support positions.

The proof gives a polynomial cell-normalizer loss $Q_i\le Cn^p$ and
$k=\lceil p\rceil+5$. It uses a genuine product source polydisc from
separate chessboard factors and common univariate normalizers. The
separation singularity is explicitly $d^{-5n}$, not a single $d^{-5}$
prefactor independent of $n$. D2 is therefore a conditional deduction
of the ordinary analytic input D1 with its specified geometry and G1;
it is not an additional graded-source hypothesis. D1's analytic estimates
and actual RG realization are not proved by this deduction.

\medskip\noindent\textbf{Technical statement D3: age-resolved uniform rarity.}
The probabilities or positive weights of live regions persisting from a given
level to a later level must have a summable bound uniform in volume and in the
location of a test. Witness selection, replication by reflections, and pattern
counting must preserve the insertion and source weights. State the resulting
tail function $r(m)$ with $\sum_m r(m)<\infty$, together with every additional
moment in $m$ used later. A rarity estimate for one newly created component is
insufficient for this persistence assertion. This input represents the age and
survival arguments wherever actual history persistence or weighted comparisons
are used. The ordinary graded all-state estimate D2 is obtained by its GS
route above and does not require this separate history-persistence hypothesis.

For deductions from a fixed-bump Riemann sum, the cutoff is required to be sufficiently fine.

\medskip\noindent\textbf{Requirement D4: one-point transfer derived from D1.}
\cref{thm:all-state-source-transfer} proves, under D1, that periodic
Gibbs states $\nu,\rho$ at the same datum satisfy

\begin{gather*}
\FitMath{|\nu(\mathcal O_K(f))-\rho(\mathcal O_K(f))|
       \le C_1(R)\|f\|_\sharp.}
\end{gather*}

If both states are invariant under all plaquette isometries, the bound implies
$|\nu(a_p)-\rho(a_p)|\le C a_K^4$ by testing a fixed nonnegative bump of
nonzero integral. The uniform bound on differences follows from a common holomorphic
normalizer; neither
individual uncentered mean is required to remain bounded. This is essential
when a pure component is centered with respect to a mixture.

\paragraph*{E. Flow, tuning and comparison of cutoffs: four inputs}

\noindent\textbf{Technical statement E1: a consistent flow.}
The recursion $x_j=\beta-\sum_{i\le j}b_i$ must be defined up to first passage,
with $b_i\le B_\sharp$ and
$\lambda(j-i)-C\le x_i-x_j\le B_\sharp(j-i)$, $\lambda>0$. Its increments
must be independent of the admissible torus and of the unused finer cutoff.
These are hypotheses for the elementary first-passage argument; they cannot be
deduced from a recursion whose next step is itself undefined outside the
invariant tube specified by the flow lemmas. Together with B3 they imply the leading logarithmic relation
between $\beta$ and the cutoff.

\medskip\noindent\textbf{Technical statement E2: a frozen calibration and tuning map.}
For auxiliary choices fixed independently of $\beta$, the terminal map
$\beta\mapsto x_K^{\natural}(\beta)$ must be holomorphic on a stated complex
neighborhood of a real interval. On that interval require
$0<c\le\partial_\beta x_K^{\natural}\le C$, and endpoint values bracketing
the prescribed renormalized inverse coupling. All intermediate values must
remain in the allowed tubes. This supplies existence and local uniqueness of a
tuned bare coupling by elementary analysis. The holomorphy assertion is not
valid merely because the original finite-lattice partition function is entire.
The flow section proves the slope implication both from the stronger
independent coupling polydiscs and from the weaker marked bare-disc
route. In the latter, the marked forcing is of size
\(C\epsilon'_i/\varrho\), not \(C\epsilon'_i/x_{i-1}\), and
is retained in the discrete Gronwall estimate. Its summability needs
the stated exponent restriction. Canonical tuning uses the admitted
real interval; broader presentation gluing additionally needs an
admitted real overlap. Overlapping geometric complex discs alone
do not establish that admission.

\medskip\noindent\textbf{Technical statement E3: a summable matched-cutoff comparison.}
At matched terminal coupling, require for each fixed torus and separated tuple

\begin{gather*}
\FitMath{|S^T_{n;K+1}(X)-S^T_{n;K}(X)|
     \le C(n,d,g,\ell)\epsilon_K\prod_i\|f_i\|_\sharp,} \\[3pt]
\FitMath{\sum_K\epsilon_K<\infty.}
\end{gather*}

The floor in $K$ may depend on the fixed torus. To compare two nearby terminal
couplings, add a uniform Lipschitz term proportional to their difference.
Neither assertion implies infinite-volume uniqueness. This is the actual
analytic input behind the matched-run claims numbered 8a and the tuned Cauchy
sequence; a bounded difference at each cutoff does not suffice.

\medskip\noindent\textbf{Technical statement E4: compatibility of calibrations.}
When moving between a self-calibrated first-passage run and a frozen run, the
maps must represent the same bare integral and satisfy a stated uniform bound
on their coupling discrepancy. That discrepancy must be small enough for every
comparison domain and source radius used above. Exact dilation identifies the
same finite lattice in two units, but supplies no estimate on this calibration
discrepancy. Uniqueness statements must state which terminal coupling and which
calibration they use.

\paragraph*{F. Short-distance normalization and nontriviality: four requirements}

\noindent\textbf{Requirement F1: the proved free-kernel comparison.}
The exact continuum coefficient is proved in \cref{ch:free-curvature-exact}.
For the remaining lattice comparison, let $a=\ell/n$, with integer $n$, and use
plain coordinate inner products on site, bond and positively oriented plaquette
cochains of the periodic lattice. Let $\mathrm{d}_1$ be its unscaled circulation matrix
and $P_{\rm ex}$ the orthogonal projection onto $\operatorname{ran}\mathrm{d}_1$.
A standard Gaussian on that range has covariance $P_{\rm ex}$. Its physical
curvature kernel is the $6\times6$ matrix of components

\begin{gather*}
\FitMath{K_{a,\ell;P,Q}(p,q)=a^{-4}P_{\rm ex}(p_P,q_Q),}
\end{gather*}

with each component sampled at its own plaquette barycentre. \Cref{ch:lattice-curvature-comparison} proves the following
comparison with the continuum matrix $K$, with $\gamma=1/2$:

\begin{gather*}
\FitMath{|K_{a,\ell;P,Q}(p,q)-K_{P,Q}(x(p)-x(q))|
 \le C r^{-4}\big[(a/r)^\gamma+(r/\ell)^4\big],} \\[3pt]
\FitMath{a\le r/10,} \\[3pt]
\FitMath{r\le\ell/10,} \\[3pt]
\FitMath{r=|x(p)-x(q)|,}
\end{gather*}

for some fixed $\gamma>0$, uniformly in $a,\ell$, plaquette positions and
orientations. Periodic differences are represented by the short displacement
in this range. The same section proves separated smearing convergence and the resulting
quadratic-density coefficient. This closes the free exact-cochain Gaussian
requirement; the interacting extraction F2 remains an independent input.

The projection convention in this input is distinct from physical weighted
block means. When passing to unit physical blocks use the adjoints and physical
difference $\mathrm{d}_\eta$ established in the block-normalization section. Under plain
fine and coarse inner products an arithmetic block mean does not have an
isometric adjoint; no projection identity may silently interchange these
conventions. Neither normalization proves the interacting extraction F2.

\phantomsection\label{hyp:F2}
\medskip\noindent\textbf{Technical statement F2: extraction of the leading two-point term.}
For the prescribed short-distance test pairs require the finite-cutoff identity

\begin{gather*}
\FitMath{S^T_{2;K,\ell}(f,h)=C_G\zeta_K(d)^2g_K(d)^4
                      \mathcal K(f,h)(1+r_{K,\ell}(f,h)),}
\end{gather*}

where $C_G>0$, the real-data error satisfies $|r|\le1/2$, and all lower
cutoff requirements are explicit. For the angular deduction this input
must cover the full class $\mathfrak T(d;\vartheta_*)$ in \eqref{eq:angular-two-point-elimination-TU9}--\eqref{eq:angular-two-point-elimination-TU13}
of \cref{ch:angular-two-point-elimination}, with
$\vartheta_*=1/12000$, or a broader class. These are nonnegative $C^1$
tests in radius-$d$ balls, with support distance between $d$ and $6d$,
the stated lower integral and upper Lipschitz bounds, and arbitrary
positions. A law proved only for a narrower shape class is insufficient
without a further inclusion argument.

The same bound is required for $S^T_2[\rho]$ in every periodic Gibbs state
used to approximate the limit or any component. At fixed $d$ its
eventual cutoff threshold and constants must be uniform in those states
and support positions. This quantifier is part of the stated F2
lemma, not a consequence of a torus-only law. On each fixed torus any
asserted $r\to0$ needs a separate uniform estimate; it must not be silently
promoted to a statement uniform in volume.
This is the substantive Gaussian-extraction input behind the short-distance
law, including all interaction and large-field remainders.
The finite source insertion and renewed-pointing proofs retain the
\(\alpha^2s\) potential norm under the actual unmarked-majorant,
leg-path and complex-domain hypotheses. The window formation and
linear bootstrap propagate it under their explicit ordinary inputs;
the resulting near-field estimate requires the complete field to
have size \(\delta\), not merely a \(\delta\)-sized increment on an
arbitrary \(\alpha\)-background. Remote and holonomy tails retain
their separate bounds. These deductions do not verify the background
and history inputs for the interacting construction.

The actual clean-history normal form must still control every old
activity's localized Hessian, the conditional exterior force and
pointwise deletion with the same exterior field and correct flag mass.
That route also requires the nonlinear averaging-fibre comparison; gauge fixing the actual run
provides a route that avoids it. The corresponding normal-form estimates are invoked in F2.
An arbitrary resolved coupling supplies neither law.
\Cref{ch:two-point-extraction-tests} proves the covariance
arithmetic and common-normalizer transfer once their inputs hold.
The old-activity part of \cref{ch:local-formation-program}
derives the localized Hessian from the actual all-substitution window room,
window-first value bounds and holomorphy on a Gaussian-coordinate
ball of radius \(c\alpha/g\). The resulting scale-small bound has
factor \((g/\alpha)^2\). Establishing that room, and the conditional
exterior-force and pointwise clean-class laws, remains necessary.
The local one-step chart or a value bound without a complex domain
does not supply those conditions.

The transfer of a Gaussian reference law has two different scopes.
D1's common cell normalizer compares any two periodic infinite-volume
states; an infinite reference obtained with a Gaussian estimate uniform
along arbitrarily large tori therefore needs no finite-volume vacuum
comparison. A law proved only on one prescribed finite reference torus
instead requires a comparison with those states. Compatible local
counterterms, the floor and vacuum-density comparison give one
sufficient route. The local companion criterion in \eqref{eq:two-point-extraction-partial-ET8}--\eqref{eq:two-point-extraction-partial-ET10} proves
that vacuum comparison from explicit nonwrapping and normalization
data; it is not required in the infinite-reference route. Per-torus
estimates whose distance thresholds decrease with the volume do not
establish the latter reference law.


\paragraph{The conditional law cannot be replaced by a formal Gaussian.}
\Cref{ch:clean-conditional-law} gives the actual
coarea/conditional-density bookkeeping and the exact
Stein identities, including cutoff boundary fluxes.
It retains the exterior force, the conditional mean,
the third and fourth cumulants, and the change of
conditional means when clean and bad classes are mixed.
A small Hessian for old activities bounds none of
these other quantities by itself.  The remaining
comparison must identify one actual positive
conditional measure and its normalizer, and bound
these terms uniformly on the specified exterior
fields before applying the covariance arithmetic.
The localization of both the observable and its
deletion must be the one in that identity.  The test
transfer additionally uses a uniform scaled \(C^1\)
shape bound and the stated lattice-to-support ratio,
so a positive continuum mass cannot disappear at
every sampling point.  These are explicit remaining
law and boundary estimates for the F2 route.

\paragraph{The Gaussian cutoff reference.}
\Cref{ch:gaussian-cutoff-clean-law} proves uniform normalized
boundary-flux and quadratic Wick-error estimates for centered Gaussian
laws with the stated convex physical cutoffs, including dependent
blocks. The actual flat first-step elimination is proved at \(x=\theta\),
with curvature covariance \(P_{\mathrm{d}_1\ker P_L}\). Comparison with the full
\(P_{\mathrm{ex}}\) reference also requires its coarse complement.
The interacting conditional law, complete Hessian, exterior force,
normalized force/boundary cumulants and common-reference means remain
model inputs.

\medskip\noindent\textbf{Technical statement F3: the running normalization at the witness scale.}
The witness depth must be defined without reference to a desired conclusion,
and its lattice separation must lie within the permitted comparison range.
Require

\begin{gather*}
\FitMath{g_K(d)^{-2}=c_*\log(1/d)(1+o(1)),} \\[3pt]
\FitMath{\zeta_K(d)=1+o(1),} \\[3pt]
\FitMath{c_*>0,}
\end{gather*}

with the precise uniformity needed by F2. A derivative of a running coupling
with respect to a constant source must be shown to equal this normalization;
it is not an independent choice of field renormalization. F2--F3 give a
nonzero two-point limit by a closed lower bound.
\Cref{ch:comparison-capsule-v7} proves summability of the
normalization slopes from the actual increment estimates on
coupling-sized independent complex polydiscs and control of the
large-field remainder loss. The flow section preserves the alternative
with an actual marked remainder on a fixed bare-coupling disc and its
separate summable forcing. A sharp \(O(g\log^p g^{-2})\) slope rate
on this weaker route needs \(p_R\ge3/2\); summability alone needs
\(p_R>1\). Analyticity on an ordinary source disc alone does not
provide the required marked size and history bounds.

Here is their sufficient ultraviolet consequence, with its size estimate
included. Use

\begin{gather*}
\FitMath{\mathcal K(f,h)=\iint f(x)h(y)|x-y|^{-8}\,dx\,dy .}
\end{gather*}

If the nonnegative supports lie in radius-$d$ balls at distance at least
$d$, and $v_4$ is the volume of the unit ball in $\mathbb R^4$, then

\begin{gather*}
\FitMath{0\le\mathcal K(f,h)\le d^{-8}\|f\|_{L^1}\|h\|_{L^1}
       \le v_4^2\|f\|_\infty\|h\|_\infty .}
\end{gather*}

The uniform F3 bounds imply $|\zeta_K(d)|\le Z$ and
$g_K(d)^4\le C[\log(1/d)]^{-2}$ for small $d$ and all permitted sufficiently
fine cutoffs. Thus F2 gives

\begin{gather*}
\FitMath{0\le S^T_2[\rho](f,h)
     \le C'[\log(1/d)]^{-2}\|f\|_\infty\|h\|_\infty .}
\end{gather*}

The lower inequality uses $1+r\ge1/2$; the upper inequality uses
$1+r\le3/2$. At each fixed pair these closed inequalities pass to every
limit and to the all-state phase realizations. Taking
$u(d)=C'[\log(1/d)]^{-2}$ proves precisely \eqref{eq:angular-two-point-elimination-TU10}, with $u(d)\to0$.
This deduction does not verify the state-uniform extraction or running
normalization required in F2--F3.


\paragraph{Flow bookkeeping and its actual coefficients.}
\Cref{candidate:flow} proves the arithmetic
first-passage, forced-slope and tuning implications
under the stated coefficient, source-disc and
history-error bounds.  The reference frame, initial
physical weight and vacuum normalization are kept
fixed in each comparison.  Tensor reduction uses
the period average where only periodic covariance
has been established; a pointwise own counterterm
requires its full pointed covariance.  What remains
for the interacting flow is to identify and estimate
the literal formed coefficients and their admitted
overlap domains.  An alternative auxiliary pointing
or current convention cannot supply those coefficient
identities without its own comparison.

\medskip\noindent\textbf{Technical statement F4: interpretation of the observable sector.}
The conclusions of F1--F3 concern the scalar $\operatorname{tr}(F^2)$ sector. To
identify a reconstructed theory as the full Yang--Mills theory, specify a
sufficient gauge-invariant observable family, its consistency relations,
renormalization, and its relation to this scalar sector. Supply the needed
continuum gauge identities or equivalent identification. A nonzero scalar
Wightman field alone does not provide this identification. If no such input is
verified, the conclusion must be stated as existence of the conditional scalar
observable theory, with full gauge-theory identification left open.


\paragraph{A complete cutoff family and a precise extension problem.}
\Cref{ch:gauge-observable-algebra,ch:su2-wilson-orbit-coordinates} construct the full
cutoff invariant algebra.  Finite Wilson trace words
separate every finite \(\operatorname{SU}(2)\) gauge orbit,
including degenerate ranks, and their polynomials are
uniformly dense in the invariant continuous functions.
\Cref{ch:scalar-observable-information} exhibits
configurations with arbitrarily small plaquette curvature
and identical action data but distinct Wilson traces.
Its example states are not asserted to be Wilson Gibbs
or continuum states.  The conclusion is the precise
information distinction between the two algebras.
\Cref{ch:gauge-observable-extension-criteria}
proves the mixed positivity, relation, consistency,
support, graph-domain and representation consequences
of specified joint readouts.  The remaining F4 task
is to construct a sufficient family of such readouts
from this same lattice theory, prove its compatible
joint continuum estimates and reconstruction, and
justify sufficiency for the intended gauge theory.
Cutoff density alone is not a uniform continuum
approximation theorem.

\paragraph*{G. Symmetries, rotations and phases: four requirements}

\noindent\textbf{Requirement G1: proved finite-lattice symmetries and positivity.}
\cref{ch:periodic-link-symmetry} proves link-reflection invariance
and the common period $2\mathbb Z^4$ for every periodic $\operatorname{SU}(2)$
Wilson DLR state at fixed $\beta>0$. \Cref{ch:all-mirror-periodic}
then proves invariance and reflection positivity at every coordinate site,
coordinate link, and diagonal lattice mirror, on the full bounded Borel
half-space bond algebra. No restriction to torus-selected states is made.
The compositions give all unit translations and signed permutations.
At $\beta=0$ the unique product Haar state has the same properties.
Adapted split tori are used for diagonal mirrors; an invalid diagonal
split of an ordinary cubic torus is not assumed.

\medskip\noindent\textbf{Requirement G2: scalar rotations derived from ordinary sources and the ultraviolet law.}
Consider limits through periodic Gibbs states, including the canonical
components with their all-state realizations. G1 and the compactness
inheritance argument provide their coordinate and diagonal reflection
positivity; \cref{ch:euclidean-hilbert-bridge} provides the exact
coordinate cube cone. A component with mean $c\int f$ is first centered
at its own mean. Its self-centered cutoff moments converge to this
centered functional by the finite subset expansion and the convergent
one-point differences; the same expansion restores the invariant mean
after the proof. Under D1 in its actual rooted geometry,
\cref{thm:uniform-one-field-insertion} proves a Hilbert insertion
bound uniform in the number of surrounding fields. This operator estimate
does not follow from D2 alone.

\Cref{thm:plate-density-from-insertion} converts that bound to an
analytic $L^1$ plate density, with the explicit loss of five powers of the
time margin. \Cref{ch:angular-kernel-construction} constructs the
holomorphic product kernels and vectors, glues their frames, and proves the
finite angular band and continued Gram positivity.
\Cref{thm:angular-two-point-elimination} uses
\eqref{eq:angular-two-point-elimination-TU18} to remove every
nonconstant two-point harmonic, including the critical frequency.
Its diagonal measure and cone support are constructed in the proof.
The shape bound \eqref{eq:angular-two-point-elimination-TU10} gives
the required limit. Alternatively \cref{sec:angular-mirror-budget}
deduces it from the ordinary source estimates at the running stops,
with their stated logarithmic dependence on the inverse coupling. \Cref{thm:all-orders-angular-closure} proves polygon
divisibility, integer convexity, removal of exceptional distributional
support by common cuts, and full $\operatorname{O}(4)$ covariance. QF transfers
this covariance to the diagonal-flat Schwartz domain.
For the canonical components, the
\hyperlink{phase-isometries}{symmetry of the decomposition}
also transfers covariance once it has been established for the
mixture. The separate component two-point law uses the
state-uniform estimates on their realizing states.

This proves G2 conditionally for this state class, using either
supplier of \eqref{eq:angular-two-point-elimination-TU18}.
The ordinary-source route uses the quantitative stopped estimates.
Typed Ward counter-insertions and distant-shell flux belong to the
separate route below. The \hyperlink{growing-torus-recognition}{density comparison}
also applies the route, with the same hypotheses, to limits of
first-passage sequences above the cutoff-dependent side threshold.
For a general growing cubic-torus limit one must establish the
mirror and cone interface separately.

The stronger Ward route requires the following: an exact
lattice identity must pass to
$\delta_\omega S_n(X)=-\lim_{R\to\infty}\lim_j
\langle\mathfrak B_{\omega,R};X\rangle_j$, with every bulk and
discretization error controlled and the claimed boundary limit existing
and vanishing. The order of those limits and the typed estimates needed
for them remain assumptions wherever that route is invoked. The angular
proof does not establish those separate comparison estimates.

\medskip
\noindent\textbf{Requirement G3: the proved finite-spacing plateau law.}
For a spacing $a>0$, let $\Omega_a=\operatorname{SU}(2)^{\mathcal B_a}$, where
$\mathcal B_a$ is the countable set of positively oriented bonds of
$a\mathbb Z^4$. Give this space its product topology. Let $\mathcal G_{a,\beta}$
be the translation-invariant probability measures satisfying the Wilson DLR
equations at inverse coupling $\beta$. The finite-bond conditional kernel
integrates the bonds in $W$ against the normalized density
$\exp[-\beta\sum_{p:p\cap W\ne\varnothing}a_p]$ and holds the other bonds
fixed. These kernels take continuous local functions to continuous local
functions. Thus $\mathcal G_{a,\beta}$ is a compact convex set. Its reflection invariance and reflection positivity are proved under G1.

\phantomsection\label{notdef:invariant-plateau}
Let $\mathcal I$ be the completed $\sigma$-field of events invariant, modulo
the measure, under all translations in $a\mathbb Z^4$. For bounded local
functions $A,B$ define

\begin{gather*}
\FitMath{\operatorname{Pl}_\nu(A,B)
 =\nu\bigl(\overline{E_\nu[A\mid\mathcal I]}
                         E_\nu[B\mid\mathcal I]\bigr).}
\end{gather*}

\Cref{thm:lattice-plateau} proves the following law, with one numerical
constant $C$ independent of $a,\beta,\nu,r,n,A,B$. If $r\ge1$ is an integer,
$A,B$ depend only on bonds with both endpoints in $a[-r,r]^4$, and
$y\in a\mathbb Z^4$ has $n=|y|_\infty/a-2r\ge1$, then

\begin{gather*}
\FitMath{|\nu(\overline A\,\tau_yB)-\operatorname{Pl}_\nu(A,B)|
 \le C\left[(r/n)^2+(r/n)^3\right]
                \mathcal N_{\nu,r}(A)^{1/2}
                \mathcal N_{\nu,r}(B)^{1/2},}
\end{gather*}

where

\begin{gather*}
\FitMath{\mathcal N_{\nu,r}(A)
 =\max_{\mu=1,\ldots,4;\,\gamma\in\Gamma_{16}}
 \nu\bigl(\overline{\theta_\mu\gamma A}
                         \tau_{2rae_\mu}\gamma A\bigr)
 \le\|A\|_\infty^2.}
\end{gather*}

Here $\theta_\mu$ reflects $x_\mu$ to $-x_\mu$, and $\Gamma_{16}$ is the group
of coordinate sign changes. Nonnegativity of these reflection sizes follows from G1. The law holds for every $\nu\in\mathcal G_{a,\beta}$, including
mixtures. It is a statement before any continuum component is constructed.

\Cref{thm:compact-conditional-law} constructs the required conditional
probabilities on this compact metric space, including Borel versions and
simultaneous probability-valued uniqueness. Applied to $\mathcal I$, this
theorem, the DLR equations, and the law above yield a unique decomposition
into ergodic DLR states. The argument and the required closedness of the pure
state set are proved in the phase-interface section. Compactness of that pure
set and existence of continuum components are conclusions, not assumptions.

\medskip\noindent\textbf{Requirement G4: all-state moments derived from ordinary sources.}
Choose $a_j=L^{-K_j}\downarrow0$, couplings $\beta_j$, and centering states
$\nu_j\in\mathcal G_{a_j,\beta_j}$. Write
$\mathcal O_j(f)=\sum_p f(x(p))a_p$. Under the precise ordinary source hypothesis D1,
\cref{thm:all-state-source-transfer} proves for every
$\rho\in\mathcal G_{a_j,\beta_j}$ that

\begin{gather*}
\FitMath{|\rho(\mathcal O_j(f))-\nu_j(\mathcal O_j(f))|
       \le C_1(R)\|f\|_\sharp,} \\[3pt]
\FitMath{|\kappa^\rho_n(\mathcal O_j(f_1),\ldots,\mathcal O_j(f_n))|
       \le C_n(d;R)\prod_i\|f_i\|_\sharp} \\[3pt]
\FitMath{(n\ge2).}
\end{gather*}

The tests are real $C^1$ functions; each support is contained in a ball of
radius $R$, and different supports have distance at least $d>0$ in the second
bound. Here $\|f\|_\sharp=\max(\|f\|_\infty,c_M\|\nabla f\|_\infty)$.
The constants are independent of $j,\rho$ and support positions. They also
apply to the separated reflected copies used in G3. For the chosen sequence,
assume that the moments centered at $\nu_j$ converge on every separated tuple;
alternatively obtain such a sequence from the preceding compactness theorem.
The per-order transfer alone would not establish factorial growth. In the
actual rooted setting of D1, \cref{thm:graded-rooted-source} also
provides D2, and \cref{thm:quantitative-flat-extension} gives H1.
The bounds displayed here and G3 suffice for the canonical decomposition constructed below.
They do not assume that decomposition or clustering of its future components.


\paragraph*{H. Test spaces and reconstruction: four requirements}

\noindent\textbf{Requirement H1: quantitative diagonal-flat extension derived from D2.}
Let $\mathcal S_{\ne}(\mathbb R^{4n})$ be the Schwartz functions flat, to all
orders, whenever two arguments coincide. Under the D2 moment bound, passed
to a separated-tuple limit, \cref{thm:quantitative-flat-extension}
constructs its unique continuous extension to this space and proves

\begin{gather*}
\FitMath{|\mathsf S_n(\varphi)|\le C_{\rm E}^n(n!)^{\kappa_{\rm E}}
 \sup_{|\alpha|\le k_{\rm E} n,x}(1+|x|^2)^{k_{\rm E}n/2}
                          |\partial^\alpha\varphi(x)|,} \\[3pt]
\FitMath{k_{\rm E}=20,} \\[3pt]
\FitMath{\kappa_{\rm E}=k+32.}
\end{gather*}

The constants are independent of $n$. The same uniform statement applies
to component limits obtained from the all-state bounds. Only the
position-uniform radius-one moment bound is used; a unit partition already
gives an effective $B(R)\le CB(1)(1+R)^4$. The explicit dyadic/product-box
proof includes tensor-relation independence, all diagonal cutoffs,
infinity cutoffs, and the positive-time boundary. Its constrained tensor
density transfers reflection positivity to finite positive-time Schwartz
sequences. Thus H1 follows from D1 in its GS geometry and from G1.

A distribution only on the complement of the diagonals with unspecified
constants would not establish this result. Diagonal-flat functions are
not dense in the full Schwartz space; an extension across diagonals or
contact-term identities still need a separate renormalization statement.

\medskip\noindent\textbf{Requirement H2: positivity and continuous Euclidean evolution.}
\cref{ch:euclidean-hilbert-bridge} proves the compact-tuple Hilbert
construction from invariant separated moments, full reflection positivity
and position-independent per-ball bounds. It proves strong continuity, time
contraction and the cube spectral cone for the stated lattice approximation.
For the full positive-time test sequences in the precise reconstruction
theorem, the required finite quadratic inequalities are
$\sum_{p,q}\mathsf S_{p+q}(\Theta\overline\varphi_p\otimes\varphi_q)\ge0$.
All tensors in this expression must belong to the declared domain. Under H1 and the constrained tensor-density lemma, those full sequence
inequalities and the same Hilbert evolution follow from the compact-tuple
construction. The proof bounds all translated matrix elements before
iterating Cauchy--Schwarz. Its position-independent per-ball hypotheses
are supplied by D1 for the ordinary phase moment bounds. The compact
argument does not by itself give quantitative Schwartz growth; that
additional conclusion follows from the GS/QF estimates in H1.

\medskip\noindent\textbf{Requirement H3: the checked reconstruction theorem import.}
\cref{ch:quantitative-reconstruction-matching}
states the corrected Osterwalder--Schrader theorem in its
diagonal-flat domain and verifies its quantitative hypotheses.
The norm comparison tracks the dimension-dependent constants
as the field number grows, so both derivative order and
factorial coefficient growth have the required form.
H1--H2 supply their stated flat test continuity and positive
ordered form.  G2 supplies full Euclidean covariance and
permutation symmetry on the declared ordinary-source and
ultraviolet route.  Clustering is used for the unique-vacuum
component; the cluster-free version retains the stated
possibly nonunique vacuum.

The energy-band and test-space sections derive the actual
fields, positive unrestricted tempered moments and joint
boundary values from the damped insertions, and derive
Lorentz covariance from G2.  Their identification with the
imported realization is proved before locality is transferred.
The resulting field is on \(\mathcal K_{\rm ord}\).
Equality with the whole Euclidean Hilbert space is not
required.  The representation and phase sections give the
remaining domain and scalar direct-integral identifications.

The corrected theorem applies by the hypothesis checks above.
We cite its continuation and locality conclusions; the
model estimates A--F require their own proofs.  Contact terms invisible to
diagonal-flat Euclidean data are not assigned arbitrary
values.  Actual D1, F2--F3 and the existence of the
admissible data supplying these hypotheses remain the
model requirements.

% Keep the final requirement together instead of leaving its last sentence alone.
\clearpage

\paragraph{What is reconstructed directly.}
The energy-band construction gives the actual fields
and their common invariant energy domain.
The test-space estimates supply full non-tensor
Schwartz extensions, joint boundary limits and the
needed graph approximations.  The representation
section identifies their canonical closures and
generators.  \Cref{ch:scalar-phase-hilbert-bridge}
constructs the measurable direct integral of the
ordered component sectors and proves that its
mixed cyclic range is onto that integral.  The
remaining named reconstruction import supplies
locality and its stated stronger output order
after the verified hypothesis matching.  Its proof
is an identified functional-analysis import,
not an additional unknown model estimate.

\noindent\textbf{Technical statement H4: scope of the final conditional assertion.}
The conditional existence statement may use only verified deductions on its
declared route. For periodic-state limits the ordinary D1 hypothesis in its
actual geometry, the ultraviolet F2--F3 hypotheses, and the proved bridges
give the scalar conditions on a sequence of admissible data;
H3 is supplied by the checked theorem import just stated. Actual verification of
D1 and of the interacting short-distance law is not supplied here.
The additional weighted, matched-cutoff and Ward conclusions require
their stated hypotheses. Full gauge-theory
identification additionally needs F4. A mass-gap conclusion would require a
new estimate, for example a bound yielding spectrum outside zero contained in
$[m,\infty)$ for some $m>0$. None of the packages above supplies that estimate.
A finite lattice measure, a subsequential separated-point limit,
an axiomatic scalar theory, and its Yang--Mills identification are
successive claims, each requiring its own proof.

\paragraph{The exact spectral question left open.}
The scalar direct-integral theorem identifies
the mixed zero-energy space with \(L^2(X,\lambda)\).
It also proves that the gap above that entire
space is the essential infimum of the component
gaps.  Positivity of a component gap at each phase
does not imply a positive uniform lower bound.
An equivalent sufficient target is a common
positive exponential decay rate for the positive
heat-semigroup matrix elements on a dense core,
after subtracting the full zero-energy projection.
No such rate is proved here.  Enlarging the
observable family requires the corresponding
joint spectral and vacuum assertions on the
larger Hilbert space, rather than only on its
scalar subspace.



% END MODULE chapters/dependency_catalogue.tex

% BEGIN LAYERED EXPANSION analytic_construction_map.tex
\subsection{Domains and dependencies of the analytic construction}
\label{map:analytic-construction}

The inverse, background and integration results use
different analytic domains. We specify their order
here, including the two changes of domain which
enter the regular trajectory. This also fixes the
meaning of the physical identifications used below.
All constants are uniform in depth and compatible
volume after the displayed choices have been made.
They need not be uniform as an exponent approaches
an endpoint or as the blocking integer is changed.

\subsubsection{The reference for the inverse}

An admitted proper geometry specifies the field
and site rows, their assigned scales, bare
components and complete determining collars.
It is given before an inverse is constructed.
Choose a real current-regular reference \(U_*\)
in that geometry. A nearby pair is
\begin{equation}\label{map:linear-domain}
 U=e^{i\eta E}U_*,
 \qquad \|E\|_{\mathcal A_\alpha}<r_*,
 \qquad \|J\|_{\infty,3}<\gamma_* .
\end{equation}
Here \(r_*,\gamma_*\) are the smaller
strong-family radii chosen after
\(\alpha,p_{\rm ell},\kappa_+\) and
the final separation and reblocking
parameters. The admitted geometry
includes these separation choices.
The energy current ceiling alone
does not give these strong families.
The frames and norms are fixed at \(U_*\).
The present field \(U\), its restoration and its
actual constraint are allowed to vary.
The current \(J\) is an independent parameter.
It is not defined by differentiating the action
of \(U\).

The order of the operators is consequential.
Positive local Dirichlet columns first give
the field and site primitives. Positive coarse
inversion then gives the projected zero-current
inverse. The first current correction is
performed before its row Schur inverse and
first lift are formed. That lift supplies
the multiplier readout used in the second
current correction. In the notation of
\cref{prop:ncw-named-families}, this order is
\begin{equation}\label{map:inverse-order}
 G_0,\quad G,\quad H,\quad
            \omega_*=H^\sharp J,\quad G_1,\quad H_1 .
\end{equation}
The second lift is not substituted for \(H\)
in the definition of \(\omega_*\).
The physical transposes include the row
lengths and fine counting densities before
any supremum norm is estimated.

The fourth-order construction in
\cref{prop:bhr-strong-restoration-family}
provides a new coefficient family for the
site restoration. It supplies two
H\"older derivatives of the gauge parameter.
The strong current construction in
\cref{prop:scf-strong-current-families}
then provides the full field--multiplier
saddle in the required finite-\(p\)
force classes. These are estimates on
coefficients and their complete carriers.
An endpoint decay bound for the completed
operator would not suffice for this step.

There is a second distinction in the norm.
For joint Taylor expansions of local base
data the required positive majorant is
\begin{equation}\label{map:local-sup-order}
 \sup_y\sum_X e^{\kappa_+\sigma_c(X)}
       \sup_{P_X\in\mathfrak D_X^+}\|K_y[X](P_X)\|.
\end{equation}
The supremum over each local coefficient
domain precedes the carrier sum. A bound
with one global supremum outside that
sum does not imply this bound.
\Cref{lem:ngc-local-data-majorants} proves
the stronger order of operations.
In a coarse short-power stage it uses
the uniform entry bound and the full
endpoint count, giving
\(C(1+B_m)^4(q')^m\).
The exponential contraction pays that
polynomial after the fixed construction
parameters are enlarged.

\subsubsection{The nonlinear germ and its physical restriction}

The auxiliary scalar action is specified
explicitly in
\eqref{cng:auxiliary-action}. Its current
convention is part of its definition.
The nonlinear equation is solved on
the domain \eqref{map:linear-domain}
with a smaller strong radius.
At zero reduced datum its auxiliary
potential is zero. Put
\(C(B)=\Pi_U(A(b_U(B)))\) for
the completed field before the final
literal node gauge. The derivatives are
\begin{equation}\label{map:germ-jets}
 D_bA(0)=H_1,\qquad
 D_b\Pi_U(A(b))\big|_{b=0}=SH_1,\qquad
 D_BC(0)=SH_1Db_U(0).
\end{equation}
The last expression includes the actual
literal-to-reduced datum map. A further
literal node gauge contributes its
specified fine gauge vector to the
literal field derivative.

The remainders are formed before they
are assigned to carriers. In particular,
\(QS=Q\) is used to write the constraint
remainder as \(Q\Xi+N_q(\Pi)\), so each
of its terms has quadratic field degree.
The current and multiplier remainders
are rewritten in their actual point-force
classes before restoration adjoints are
applied. Thus the tree series has no
linear term that must cancel between
different carriers.
\Cref{prop:ngc-datum-connected-family,prop:ngc-joint-connected-family}
prove convergence over every datum
degree and every local base degree.
The base variables in that statement
are supported strong Banach slots,
with bounded-overlap reconstruction.
They are not independent discs for
every raw fine link.

The physical restriction is an exact
scalar identity. At
\begin{equation}\label{map:physical-pair}
                   J=j_U=Q^\sharp\omega_{\rm phys},
\end{equation}
the selected auxiliary action equals
the physical action on the finite
gauge completion, with its constraint
multiplier and explicit penalties.
The gauge tests remove the gauge penalty.
The artificial row mass is constant
on the prescribed constraint and is
removed from the physical multiplier
by the shift \(\nu_{\rm phys}=\nu-a_fb\).
Uniqueness then identifies the solution
with the physical constrained orbit.
This is the content of
\cref{prop:cng-physical-matching}.
The off-current construction remains
the selected convention just defined;
its existence does not identify any
other prescribed auxiliary action.

\subsubsection{Producing the homogeneous reference}

The reference in
\eqref{map:linear-domain} is an input
for a proper geometry. In the homogeneous
case it can instead be constructed from
coarse curvature alone. For a real
coarse field \(V\) with sufficiently
small plaquettes,
\cref{prop:cta-uniform-trial}
constructs a smooth quaternionic
bundle connection whose original
coarse-edge transports are exactly
those of \(V\). Its fine samples
have physical curvature and Wilson
current \(O(\kappa)\).
Their actual hierarchical means
differ from \(V\) by \(O(\kappa)\),
uniformly in the number of means.

The current estimate pairs opposite
plaquettes before division by the
third power of the fine spacing.
The mean estimate uses the full
plain-mean row majorant. These two
operations prevent, respectively,
a loss of one inverse fine spacing
and a factor growing exponentially
with the number of blocking steps.
The forced trial equation retains
the trial current as a source and
corrects both defects.
\Cref{prop:cta-curvature-critical-atlas}
therefore produces the physical
homogeneous orbit without assuming
a global gauge close to a flat
connection.

The bundle construction also supplies
the comparison of nearby real
references. Its Kato map aligns the
connections before fine edges are
sampled. A change of a coarse datum
then gives a relative fine potential
of size \(O(\|B\|)\), rather than a
frame jump divided by the fine spacing.
The resulting strong comparison,
with the larger uniqueness margin,
identifies complex germs on overlaps.
There is no need to assume that an
arbitrary complex overlap contains
a real point.

For constant diagonal plaquette weights,
\cref{prop:wca-weighted-curvature-atlas}
uses the same smooth trial.
Its principal perturbation is estimated
from \(W^{2,p}\) to \(L^p\).
The full joint inverse is first proved
in that stronger field norm.
It is not inferred from a
\(C^{1,\alpha}\) output estimate.
The weighted current of the corrected
physical orbit is small because it is
the actual row transpose of the bounded
physical multiplier. No point bound
for its unweighted current is inserted
into this argument.

These statements produce homogeneous
references, including large holonomies
and nontrivial bundle topology.
The definition and construction of
mixed-scale proper histories still
use their stated geometry and
admission hypotheses.

\subsubsection{The physical Gaussian and the integral domain}

At a physical critical pair the
mass-added second inverse has a
positive energy gap. Its effective
coarse action Hessian is
\begin{equation}\label{map:physical-Schur}
 D^2S=(QG_1Q^\sharp)^{-1}-a_f .
\end{equation}
The artificial row mass is subtracted.
The actual outer precision also
contains the first action gradient
paired with the second derivative
of the outer chart:
\begin{equation}\label{map:outer-precision}
 K_\tau=\tau\{\mathsf C_V^\sharp(D^2S)\mathsf C_V+
                       \langle\omega,\Phi''(0)\rangle\}+K_G .
\end{equation}
Here \(\mathsf C_V=D\Phi_V(0)\)
is the linear outer chart map.
The lower bound for this matrix
is proved by flux transfer and
the geometric curl--gauge estimate.
It does not follow merely from
positivity of the mass-added
coarse Schur matrix.
\Cref{neg:outer-Gaussian-gap,prop:wca-outer-Gaussian-gap}
prove the real physical gaps, and
\cref{acg:canonical-Gaussian-admission}
constructs the local Gaussian
families with their sector and
normalization bounds on its declared
unweighted canonical domain near a
physical pair with small unweighted
current. The weighted physical
continuation has the sector proved
in \cref{prop:wca-complex-outer-sector}.
That sector statement does not by
itself supply independent weighted
Gaussian cube families.

For the regular trajectory the old
interaction is defined on a fixed
raw field ball of radius \(r\).
It must be evaluated there throughout
the integration. Let \(R>r\) be the
larger output radius. The literal
intermediate field and the outer
graph satisfy the two separate bounds
\begin{equation}\label{map:two-domain-buffers}
 \|\log\mathsf p_j(W)\|_\infty
           \le C_pR/L<r/64,\qquad
 \|\log Z_V(\gamma B)-\log V\|_\infty
           \le C_{\rm dom}(L)|\gamma|b<r/64 .
\end{equation}
Thus an output argument on the
larger ball is sent well inside
the old ball before the fluctuation
is added.
The complete old functions are
evaluated there. Gauge identities
on admitted overlaps are used
after that inclusion has been
proved.

Here is a useful explicit order
for these choices. First fix
\(0<\alpha<1\) and
\(p_{\rm ell}>4/(1-\alpha)\)
for the strong estimates. Fix the desired
transport shrink \(t\), final spatial
and period rates, and a larger
intrinsic rate. Choose the odd
blocking integer \(L\) from
\eqref{lcf:blocking-buffer};
the literal amplitude constant
in that choice is uniform in \(L\).
Choose the Gaussian and carrier
localization parameters for those
fixed rates. Choose \(R\) within
the resulting common analytic
radii and put \(r=tR\).
The constants \(C_{\rm dom}(L)\)
and the positive Gaussian margins
are now fixed. Finally enlarge the
coupling floor until
\eqref{map:two-domain-buffers}
holds at the physical
\(\gamma=\theta^{-1/2}\) with
\(b=A_1(\log\theta)^{p_{\rm cut}}\),
where \(p_{\rm cut}>1/2\).
The last choice is possible since
\(\theta^{-1/2}(\log\theta)^{p_{\rm cut}}\)
tends to zero. It changes no
previous geometric parameter.

The finite exponent used in the
elliptic estimates is a different
parameter, fixed before the strong
inverse constants and the blocking choice.
The logarithmic exponent
\(p_{\rm cut}\) belongs to the
real fluctuation cutoff. Keeping
these roles distinct prevents a
false endpoint conclusion for
the elliptic estimates.

\subsubsection{Formation, tuning and differentiated cutoffs}

For a fixed profile ceiling the
formation bounds give
\begin{equation}\label{map:formation-slopes}
 \|\mathbf F_j\|+\|D_h\mathbf F_j\|
       \le C_H\epsilon(\theta_j),\qquad
 \|\partial_x\mathbf F_j\|
       \le C_H\epsilon(\theta_j)/\theta_j,\qquad
 \epsilon(t)=(\log t)^m/t .
\end{equation}
The reference \(\theta_j\), its
cutoff and the Gaussian coordinate
scaling are frozen. Differentiation
in \(x=\theta_j\tau\) means
\(\theta_j^{-1}\partial_\tau\).
The common-radius transport has
an exact semigroup and exponential
age decay. Its adapted norm gives
a one-step contraction. These are
the inputs for the opposite-end
conditions \(h_0=0\), \(x_N=X\).

The scale sums needed for tuning
have different sizes:
\begin{equation}\label{map:tuning-sums}
 \sum_{j<N}\epsilon(a_j)
       \le C\{1+(\log a_0)^{m+1}\},\qquad
 \sum_{j<N}\frac{\epsilon(a_j)}{a_j}
       \le C(\log X)^m/X ,
 \qquad a_j\asymp X+N-j .
\end{equation}
The first sum need not be small.
It is a sublinear correction to
the bare coupling. The second
sum controls the contraction and
can be made small independently
of \(N\).
\Cref{prop:hrt-regular-trajectory}
uses these two sums separately.
\Cref{prop:rts-bounded-UV-tuning,prop:rts-exponential-universality}
allow a bounded irrelevant
initial profile and estimate its
effect after retuning to the
same endpoint.

Changing a plaquette weight is
a different derivative.
The moving weighted background
changes the centre of the free
colour balls. For its first
weight derivative,
\cref{prop:fcw-frozen-weight-jet}
compares it with the fixed
isotropic kernel on the same
Haar-disintegrated fibre.
The difference is a normalized
rooted sphere functional with
bound
\begin{equation}\label{map:cutoff-tail}
 C_H\theta^{1/2}(1+b)^K e^{-cb^2}.
\end{equation}
This is smaller than every
fixed inverse power of \(\theta\).
Its six curvature readouts
vanish. It must nevertheless
be included in the profile
forcing before the diagonal
stress recurrence is estimated.
\Cref{prop:rds-regular-stress-decay}
uses this actual fixed-kernel
first jet. It makes no claim
about a uniform fixed-centre
complex weight disc.

\subsubsection{The interface with the continuum statements}

The analytic constructions above
give the specified inverses,
physical homogeneous backgrounds,
regular integration maps and
tuned regular trajectories.
Their identities retain their
reference domains and the
complete functions to which
pointings were assigned.
They provide concrete inputs
to the multiscale source argument.

The main continuum statements
also require the stated estimates
for complementary histories,
exterior changes, source normalizers
and the short-distance comparison.
The regular trajectory does not
replace those operations by a
bound for a single positive
quadratic form. Nor does a
classical or Gaussian descent
identity identify quantum
partition functions on different
volumes. The hypotheses in
\cref{thm:continuum-main,thm:phase-main,thm:reconstruction-main}
retain those specified model
estimates. Each use below refers
to the corresponding quantitative
statement, with its original
source and geometry class.

% END LAYERED EXPANSION analytic_construction_map.tex


\section{Background fields and constrained inverses}\label{sec:background}
The renormalization step integrates fluctuations around a constrained background field. We construct the averaging maps and their coordinate charts, compare the flat and background-dependent quadratic forms, and state the uniform inverse estimates on the required current spaces. The nonlinear response estimates then control differentiation of the background and of its propagator. These are the inputs to the local integration in \cref{sec:interactions}.

% BEGIN MODULE chapters/02_averaging.tex
\subsection{Block averages and gauge fixing}
\label{candidate:averaging}
\subsubsection{Block averages, gauge fixing, and the renormalisation group of this construction}
\label{sec:averaging-block-averages-gauge-fixing-and-the-renormalisation-group-of-this}

A block average must turn a fine gauge field into a coarse one while
preserving gauge covariance and the coarse-lattice symmetries. The
associated change of variables must preserve the integral and provide
local fluctuation coordinates. We first specify the averaging map and
its gauge choices, which are the geometric ingredients of the
renormalisation step.

\textbf{Setting and objects.} \emph{The numbers.} Each statement of this part bears a number belonging to one of three groups. The group 1a.1 concerns the block average of a gauge field; the group 1a.2 concerns the renormalisation transformation, namely its exact identities, its invariant form and the anatomy of a single step; the group 1a.3 concerns gauge fixing and regular configurations. A label ``Theorem, Proposition, Lemma 1a.n.m'' denotes statement m in group 1a.n, while ``1a.n'' standing alone denotes the group itself. \emph{Requirements on an average.} We work in level-$k$ units, with $L$ odd and centred blocks $B(y)=\{x:|x-y|_\infty\le\frac12(L-1)\}$. We put $\varpi(W)=\sup_p|\partial W(p)-\mathbf{1}|$, and $\varpi_c$ denotes the same supremum taken over the double block $B(c_-)\cup B(c_+)$. For kernels we use the norms $\|K\|_{\mathrm{row}}=\sup_c\sum_b|K(c,b)|$ and $\|K\|_{\mathrm{col}}=\sup_b\sum_c|K(c,b)|$. Elements of the group cannot be added, and transports joining different pairs of points cannot even be compared. A block average $\mathcal{M}:\mathcal{V}_k\to\mathcal{V}_{k+1}$ is required to have the following properties: gauge covariance; locality; linearity at first order; analyticity near flat fields, on $G^c$ as well; symmetry under the isometries of the coarse lattice; and Borel measurability on the whole of $\mathcal{V}_k$, because $\mathbf{T}$ integrates over all configurations. For the mean of group elements Ba{\l}aban gives axioms (109 (1987), (0.5)--(0.9)). Such a mean is any analytic $G^c$-valued function $M$ on families of small diameter that satisfies

\begin{gather*}
\FitMath{M(\{\mathbf{U}_j^{-1}\})=M(\{\mathbf{U}_j\})^{-1},} \\[3pt]
\FitMath{M(\{u\mathbf{U}_jv\})=u\,M(\{\mathbf{U}_j\})\,v,} \\[3pt]
\FitMath{M(\pi\{\mathbf{U}_j\})=M(\{\mathbf{U}_j\}),}
\end{gather*}


\begin{gather*}
\FitMath{\tfrac1i\log M(\{e^{iA_j}\})=\tfrac1n\sum_jA_j+(\text{higher orders}),} \\[3pt]
\FitMath{\mathbf{U}_j\in G\text{ for all }j\,\Rightarrow\,M(\{\mathbf{U}_j\})\in G.}
\end{gather*}

According to Ba{\l}aban, the results of his 1987 paper and of the earlier papers hold for every mean of this kind. For the map $\mathcal{M}$, the second of the two symmetric block averages that he introduces, he states that every result of the paper on averaging operations (Comm. Math. Phys. 98) holds, with proofs that are for the most part unchanged. This statement concerns that single paper; the broader sentence, which refers to ``this, and previous papers'', is the one about means. Of the two symmetric averages together he states that they are equally good, and that many other averages can also be used. He also states that the block average admits an axiomatisation as well. Already in 1985 he states that the analysis of the paper on averaging operations, together with its propositions, extends easily to other averages that satisfy the two requirements set out in its introduction, namely gauge covariance and agreement with the linear block average at first order. No proof accompanies any of these statements. For the theorems of 1984--85 on averaging, gauge fixing, propagators and backgrounds, the proofs are carried out for an average that is not symmetric. \Cref{prop:averaging-1a-1-13} states which properties of an average the proofs use, and it compares the linear parts (1a.1).

\emph{The mean, the contours, the map.} The \emph{mean} $\operatorname{av}(V_\bullet)$ (Federbush) is the solution, close to the family, of $\sum_i\log(V_iV^{-1})=0$. By $\mathcal{G}(y,x)$ we denote the family of shortest paths from $y$ to $x\in B(y)$ that complete the displacement in one direction before passing to the next. There is one such path for each ordering of the directions, so the family has at most $4!$ members and does not contain every shortest path. The \emph{comb} $\Gamma_{y,x}$ is the member with the order $4,3,2,1$. The averaged contour variable is $U(y,x):=\operatorname{av}(U(\Gamma))_{\Gamma\in\mathcal{G}(y,x)}$, and we set $U(x,y)=U(y,x)^{-1}$ and $U(y,y)=\mathbf{1}$. With $x'=x+(c_+-c_-)$, one defines

\begin{gather*}
\FitMath{W_x=U(c_-,x)\,U([x,x'])\,U(x',c_+)\,U[c]^{-1},} \\[3pt]
\FitMath{\mathcal{M}(U)(c)=\exp\Big(L^{-4}\!\!\sum_{x\in B(c_-)}\!\!\log W_x\Big)U[c]} \\[3pt]
\FitMath{\text{on }\mathrm{Dom}(c)=\{\varpi_c<C^{-1}L^{-2}\}:}
\end{gather*}

the $W_x$ are closed transports based at $c_-$ and lying near $\mathbf{1}$, so their logarithms can be added. Where $U(y,x)=\mathbf{1}$, only straight segments are left, so the form is adapted to the block axial gauge. Off $\mathrm{Dom}(c)$ the map is given by two cruder rules, both covariant. When none of the elements $Y_{x,\Gamma,\Gamma'}=U(\Gamma)U([x,x'])U(\Gamma')^{-1}U[c]^{-1}$ equals $-\mathbf{1}$, one replaces $\log W_x$ by the arithmetic mean, over $(\Gamma,\Gamma')$, of the intrinsic logarithms of the $Y_{x,\Gamma,\Gamma'}$; this is Ba{\l}aban's first symmetrised average. In the remaining case one sets $\mathcal{M}(U)(c)=U[c]$. This third rule is required because Ba{\l}aban defines his average on all configurations by means of the spectral logarithm, which satisfies $\log(-\mathbf{1})=i\pi\mathbf{1}\notin i\mathfrak{g}$ and so leads out of $\operatorname{SU}(2)$ (1a.1). By $\mathcal{M}^k$ we mean the iterate only, and not a mean taken in a single stroke over $B^k$. In Ba{\l}aban's paper of 1987 the lattices consist of centres of cubes and have coordinates in $a_j(\mathbb{Z}+\frac12)$. The single translation $x\mapsto x+\frac12L^m\sum_\mu e_\mu$ of $T_\ell$, $\ell=2L^m$, maps them onto the $\Lambda_{j,\ell}$ at every level simultaneously, since $\frac12L^m=\frac12L^{m+j}a_j$ and $L^{m+j}$ is odd. This is the second place where the oddness of $L$ is used, the first being the centring of the blocks (1a.1). Ba{\l}aban also remarks that all the results were obtained first for the two-dimensional chiral models, in which the fields sit on the sites and the action is $\sum_b[1-\frac12\operatorname{Re}\operatorname{tr}U(b_-)U(b_+)^{-1}]$. In those models block averages are means of group elements, and neither contours nor a gauge group enter.

\emph{Linear averages} of bond functions $\mathsf{a}$, which enter in the form $e^{i\mathsf{a}}W$ without any length factor, are defined as follows, where $\bar{\mathsf{a}}(y,x)$ denotes the arithmetic mean of $\mathsf{a}(\Gamma)$ over $\mathcal{G}(y,x)$:

\begin{gather*}
\FitMath{(\mathsf{Q}^{\mathrm{s}}\mathsf{a})(c)=L^{-4}\!\!\sum_{x\in B(c_-)}\!\!\mathsf{a}([x,x']),} \\[3pt]
\FitMath{\sigma(y)=L^{-4}\!\!\sum_{x\in B(y)}\!\!\bar{\mathsf{a}}(y,x),} \\[3pt]
\FitMath{(\mathsf{Q}^{\mathrm{pl}}\mathsf{a})(c)=\sigma(c_-)+(\mathsf{Q}^{\mathrm{s}}\mathsf{a})(c)-\sigma(c_+),} \\[3pt]
\FitMath{(q\lambda)(y)=L^{-4}\!\!\sum_{x\in B(y)}\!\!\lambda(x);}
\end{gather*}

here $\mathsf{Q}^{\mathrm{pl}}$ is the derivative of $\mathcal{M}$ at $\mathbf{1}$. The operator $\mathsf{S}(W_0)$ is $\mathsf{Q}^{\mathrm{s}}$ formed with parallel transport to $c_-$, namely 

\begin{gather*}
\FitMath{(\mathsf{S}(W_0)\mathsf{a})(c)=L^{-4}\sum_{x\in B(c_-)}\sum_{b\in[x,x']}\operatorname{Ad}\big(W_0(c_-,x)W_0([x,b_-])\big)\mathsf{a}(b)}
\end{gather*}

.

\emph{The straightened average} of $\mathbf{W}=e^{i\mathsf{a}}W_0$ relative to $W_0$ and its \emph{block factor} $\mathsf{u}$ are given by

\begin{gather*}
\FitMath{\mathsf{u}(y)=\exp\Big(L^{-4}\!\!\sum_{x\in B(y)}\!\!\log\big[\mathbf{W}(y,x)W_0(y,x)^{-1}\big]\Big),} \\[3pt]
\FitMath{e^{i\mathcal{Q}^1(W_0,\mathsf{a})(c)}\mathcal{M}(W_0)(c)=\mathsf{u}(c_-)^{-1}\mathcal{M}(\mathbf{W})(c)\,\mathsf{u}(c_+).}
\end{gather*}

The right-hand side equals $\mathcal{M}(\mathbf{W}^v)$, where $v$ is the gauge transformation with $\mathbf{W}^v(y,x)=W_0(y,x)$ (the block axial gauge relative to $W_0$) and with trivial average $\bar{R}_{W_0}v=\mathbf{1}$; here $\bar{R}_W$ is the averaging of gauge transformations that is defined below under \emph{Gauge transformations}. Together these two conditions imply $v=\mathsf{u}^{-1}$ at the coarse sites. Ba{\l}aban gives two reasons for the straightening. First, the variational problems of \cref{ch:multilevel-backgrounds} involve composing the average with the passage from the axial gauge to the Landau gauge. Second, only the straight linear average admits the explicit representations of the propagators \cref{sec:propagators-propagators-in-a-background-field-exponential-decay-produced-by-the}; this average is compatible with composition, has an explicit Fourier symbol and does not change under $q\lambda=0$. The change of gauge has to preserve this form, and the non-linear counterpart of $q\lambda=0$ is $\bar{R}_{W_0}v=\mathbf{1}$. A third reason appears in the proof of \cref{prop:averaging-1a-1-6}.

\phantomsection\label{notdef:covariant-averages}
\emph{Order $k$; the maps denoted by $Q$.} Along the sequence $W_{0,n}=\mathcal{M}^nU_0$ we set $\mathsf{a}_0=\eta_kA$, $\mathsf{a}_{n+1}=\mathcal{Q}^1(W_{0,n},\mathsf{a}_n)$, $v_0=\mathbf{1}$, $v_{n+1}=v_n\,\mathsf{u}(W_{0,n},\mathsf{a}_n)$, and define

\begin{gather*}
\FitMath{\mathcal{Q}^{\sharp}_k(U_0;A):=\mathsf{a}_k=Q^{\sharp}_k(U_0)A+C_k(U_0,A),} \\[3pt]
\FitMath{\mathcal{Q}_k(U_0;A)=\tfrac1i\log\big[\mathcal{M}^k(e^{i\eta_kA}U_0)\mathcal{M}^k(U_0)^{-1}\big]=Q_k(U_0)A+C^{\mathrm{pl}}_k(U_0,A).}
\end{gather*}

A sharp indicates the straightened objects, which are the ones that enter the propagators; Ba{\l}aban's $Q_k(U_0,\eta A)$ is the same object built with combs. Without the sharp, $\mathcal{Q}_k$ and $Q_k$ refer to the plain average, whose remainder is denoted $C^{\mathrm{pl}}_k$, while $C_k$ denotes the remainder of the straightened average. The product $v_k=v_k[A]$ of the $k$ block factors is the \emph{node factor}, and $\dot{v}_k$ denotes the linear part of $\frac1i\log v_k$. The letter $Q$ is used for maps of four kinds (1a.1). (1) The non-linear constraint map $\mathcal{Q}_k$ of the plain iterate $\mathcal{M}^k$, and next to it its straightened form $\mathcal{Q}^{\sharp}_k$. (2) $Q_k(U)=\partial_A\mathcal{Q}_k(U;A)|_{A=0}$, together with $Q=Q_1(\mathbf{1})$ and $Q_V=Q_1(V)$; this is the operator in the delta function of a step. (3) $Q^{\sharp}_k(U)$, the operator that enters the propagators and the Landau gauge. Its leading term is the \emph{straight} average with parallel transport, $Q^{\mathrm{st}}_k(U)=\eta_k\mathsf{S}(W_{0,k-1})\cdots\mathsf{S}(W_{0,0})$. The two differ by $O(\alpha_0)$ in row norm, and they are equal at $U=\mathbf{1}$. There, after the translation of \cref{prop:averaging-1a-1-13}, each of them is Ba{\l}aban's $Q_k$ of 1984. (4) $Q'_k(U)$, which averages site functions. In a background, the linear average in Ba{\l}aban's papers (99 (1985) 389, (3.15)) is $Q^{\mathrm{B}}_j(U):=Q(\bar{U}^{j-1})\cdots Q(\bar{U})Q(U)$. It is the linear part of his gauge-corrected average built with combs, and each of its factors equals the straight covariant average up to $O(L^2\alpha_0)$. \Cref{prop:averaging-1a-1-6,prop:averaging-1a-1-13} provide two comparisons, that of (2) with (3) and that of $Q^{\mathrm{B}}_k$ with $Q^{\sharp}_k$, and \cref{sec:propagators-propagators-in-a-background-field-exponential-decay-produced-by-the} relies on both.

\emph{Gauge transformations} are averaged by means of the same contours,

\begin{gather*}
\FitMath{(\bar{R}_Wv)(y)=v(y)\exp\Big(L^{-4}\!\!\sum_{x\in B(y)}\!\!\log\big[v(y)^{-1}\operatorname{Ad}(W(y,x))v(x)\big]\Big),}
\end{gather*}

and the iterates of this map give $\bar{R}^{(j)}_U$. The linear part, $Q'_j(U)\lambda$, is the mean over $B^j(y)$ of $\lambda$ after transport to $y$.

\emph{Isometries on $1$-forms.} An isometry $\mathsf{r}$ acts by $(\mathsf{r}U)(b)=U(\mathsf{r}^{-1}b)$. Bond functions are specified on the positively oriented bonds, whereas a reflection reverses orientation. Since $(e^{i\mathsf{a}}W_0)(\bar{b})=\exp\big(-i\operatorname{Ad}(W_0(b)^{-1})\mathsf{a}(b)\big)W_0(\bar{b})$, the action relative to $W_0$ must include a twist:

\begin{gather*}
\FitMath{(\mathsf{r}_{W_0}\mathsf{a})(b')=\mathsf{a}(b)\,\text{ if }\mathsf{r}^{-1}b'=b,} \\[3pt]
\FitMath{(\mathsf{r}_{W_0}\mathsf{a})(b')=-\operatorname{Ad}\big(W_0(b)^{-1}\big)\mathsf{a}(b)\,\text{ if }\mathsf{r}^{-1}b'=\bar{b};}
\end{gather*}

on pairs the action is $\mathsf{r}(\mathbf{U},\mathbf{J})=(\mathsf{r}\mathbf{U},\mathsf{r}_{\mathbf{U}}\mathbf{J})$ (1a.1).

\emph{Lift, regions, determining sets.} The \emph{lift} is defined by $(Q^{s*}_kV)(b)=V(c)$ when $b_\pm\in B^k(c_\pm)$; on every bond contained in a single block of order $k$ its value is $\mathbf{1}$. For a bond function $B$ on $\mathbb{T}^{(k)}$, its linear lift is the $1$-form that takes the value $\eta_k^{-1}B(c)$ on these crossing bonds and the value $0$ on bonds lying within a block. The value $\mathcal{M}^q(W)(c)$, and also which rule applies, depend on $W$ only in $B^q(c_-)\cup B^q(c_+)$. Averages are therefore defined for a field given on part of a lattice; they remain covariant and obey the semigroup law $\mathcal{M}^{q+1}W(c')=\mathcal{M}(\mathcal{M}^qW)(c')$. Understood in this way, $M_{\mathbf{B}}(U):=(\mathcal{M}^jU|_{\Gamma_j})_j$ is defined on configurations given on a region, and the problems that define the localised backgrounds $U_{k,\square}$ are set up near a cube without reference to the torus (1a.1).

\emph{The transformation.} Here $\delta$ denotes the Dirac distribution at $\mathbf{1}\in G$ with respect to Haar measure. It is a class function, and it is the limit $\mathfrak{a}\to\infty$ of the normalised weights $\propto e^{-\mathfrak{a}(1-\frac12\operatorname{Re}\operatorname{tr}W)}$. The transformation sends a complex Borel measure $\nu$ to its image $\mathbf{T}\nu=\mathcal{M}_*\nu$, which is written with the kernel $\delta(\mathcal{M}(U)V^{-1})$. Thus a coarse configuration receives the total weight of the fine configurations whose average it is. Wilson and Kadanoff introduced the block-spin transformations $(\mathbf{T}_t\rho)(V)=\int dU\,t(V,U)\rho(U)$, with $t\ge0$, $\int dV\,t(V,U)=1$ and $t(V^{u|_{\mathrm{coarse}}},U^u)=t(V,U)$; the transformation here is the member of this family with a hard constraint. Decompositions of unity and the factor $\mathfrak{G}_k$ are inserted into the integral. They are built with the auxiliary choices $\varepsilon_k,\delta_k,\alpha_{i,k},\check{R}_k$ and with the weight of $\mathfrak{G}_k$. All of these depend on a single number $\mathsf{t}$, which is $x_k$ in the run and $\theta_k$ in a reference run. Only $\check{R}_k$ is formed with the minimum of these numbers taken over all levels $i\le k$, which makes it non-increasing in $k$. A history $h$ records outcomes, and these yield the term $\tau_h$, or the term $\tau^{\mathfrak{w}}_h$ when $\rho_0$ is multiplied by a bounded complex weight $\mathfrak{w}(U)$. In addition to $\chi_k$ and $\chi^{(k)}$ there is the approximate fluctuation test $\chi'_k(\square)=\mathbf{1}\{|V_k(V^{\mathrm{bg}}_{k,\square})^{-1}-\mathbf{1}|<2\delta_k\text{ near }\square\}$, which is built with the local averaged background $V^{\mathrm{bg}}_{k,\square}=\mathcal{M}^kU_{k+1,\square}$. Braces denote a difference between backgrounds: $\{F\}:=F(U_k(V_k))-F(U_{k+1})$.

\emph{The block algebra.} The block field of depth $s$, $V^{(s)}:=\mathcal{M}^{K-s}U$, is defined on $\Lambda_{s,\ell}$, a lattice that does not depend on the cutoff $K\ge s$. The \emph{block algebra} $\mathfrak{A}$ is the set of gauge-invariant functions $F$, bounded and continuous, that depend on only finitely many block variables $V^{(s)}(b)$. The composite $F\circ\mathcal{M}^{K-s}$ is defined for every $K\ge s$, so these observables exist at all cutoffs simultaneously (1a.2). The block algebra is a notion separate from \cref{thm:continuum-main,thm:phase-main}, whose observables are the smeared curvature evaluated at the bare scale; neither theorem refers to $\mathfrak{A}$.

\emph{Two operations} (1a.2). Let $\Delta=\mathsf{C}^{\mathsf{T}}\Delta^{(k)}\mathsf{C}$, let $z^{(k)}$ be a product of Jacobians, and let $\chi(S_{k+1})$ denote the tests on $\Omega_{k+1}\setminus\Lambda_{k+1}$. To obtain $\zeta(\Omega^c_{k+1})$, form for each outcome that determines this same $\Omega_{k+1}$ the product of the tests and gauge-fixing factors localised in $\Omega_k\setminus\Omega_{k+1}$, and add these products over all such outcomes. The one-step operation collects everything that step $k+1$ does not carry out:

\begin{gather*}
\FitMath{\mathbf{T}^{(k)}(Z_{k+1})F=\int dV_k\big|_{\Omega_{k+1}^c}\delta(\mathcal{M}(V_k)V_{k+1}^{-1})\,\zeta(\Omega^c_{k+1})\;z^{(k)}\!\int d\mathcal{A}^{\mathrm{s}}_k\big|_{\Omega_{k+1}\setminus\Lambda_{k+1}}\chi(S_{k+1})\,e^{-\frac12\langle\mathcal{A}^{\mathrm{s}}_k,[\Delta-\Delta C^{(k)}(\Lambda_{k+1})\Delta]\mathcal{A}^{\mathrm{s}}_k\rangle}F,}
\end{gather*}

where, more precisely, the Gaussian exponent of the operation is a part of this quadratic form, namely a part that couples no two components of $\Omega_{k+1}\setminus\Lambda_{k+1}$; the remainder of the form is assigned to the boundary terms $B^{(k+1)}$ \cref{sec:large-field-stored-variables-and-integration-in-a-large-field-patch}. As long as $\mathbf{R}$ does not intervene, $\mathbf{T}'_{k+1}=\mathbf{T}^{(k)}(Z_{k+1})\mathbf{T}'_k$. For $x\in\mathbb{C}$ and a function $S$ of the bare fields, we define the pure small-field integral and then the regular part of step $k+1$ by

\begin{gather*}
\FitMath{(\mathbf{T}^{\mathrm{sf}}_k[\mathsf{t}]F)(W)=\int dV\,\delta(\mathcal{M}(V)W^{-1})\,\mathfrak{G}_k(V)\,\chi^{(k)}(T_\ell)\,F(V),}
\end{gather*}


\begin{gather*}
\FitMath{\mathfrak{E}_{k+1}[\mathsf{t};x;S](W)=\log\big(\mathbf{T}^{\mathrm{sf}}_k[\mathsf{t}]\,e^{-xA(U_k(\cdot))+S(U_k(\cdot))}\big)(W)+xA(U_{k+1}(W))-S(U_{k+1}(W)).}
\end{gather*}


\emph{Regular configurations.} We consider sequences $\Omega_0\supset\dots\supset\Omega_k$, with $\Omega_{k+1}=\emptyset$, in which each $\Omega_j$ is a union of $M_1$-cubes of level $j$ and $\operatorname{dist}(\Omega_j^c,\Omega_{j+1})\ge CM_1L^j\eta_k$. The small-field case is the constant sequence $\Omega_j=T_\ell$. Layer $j$ is the set $\Omega_j\setminus\Omega_{j+1}$. It has the local scale $\xi_j=L^j\eta_k$ and the tops $\mathbb{Y}_j=\Lambda_j(\mathbf{B})$, which are its sites in $\mathbb{T}^{(j)}$. Its pillars are the $B^j(y)$, $y\in\mathbb{Y}_j$, and over all $j\le k$ the pillars partition the bare sites of $\Omega_0$. We define

\begin{gather*}
\FitMath{\mathfrak{U}_k(\Omega_\bullet,\alpha)=\big\{U:\,|\partial U-\mathbf{1}|<\alpha\,\eta_k^2\xi_j^{-2},\,\,|J(U)|<\alpha\,\xi_j^{-3}\,\text{ on layer }j\big\},} \\[3pt]
\FitMath{|f|_{(-s)}=\sup_j\sup_{\text{layer }j}\xi_j^s|f|,}
\end{gather*}

together with 

\begin{gather*}
\FitMath{\|A\|_1=\max\{|A|_{(-1)},|\nabla_{U_0}A|_{(-2)}\}}
\end{gather*}

 and 

\begin{gather*}
\FitMath{\|\lambda\|_0=\max\{\sup|\lambda|,|D_{U_0}\lambda|_{(-1)}\}}
\end{gather*}

. Both conditions are invariant under gauge transformations. Combined with the Bianchi identity, the second condition is an elliptic system for the field strength $F_U$, and all bounds on derivatives come from this system; minimisers satisfy the second condition. The complex neighbourhood of radius $\alpha_1$ is the set of $\mathbf{U}=e^{i\eta_kA'}U$ with $U$ real and regular and $\|A'\|_1<\alpha_1$. The norm $\|A'\|_1$ controls only first differences, whereas the current of $\mathbf{U}$ is a second difference. Thus the current enters the domains $U^c_j(X,\alpha_0,\alpha_1)$ as an independent variable $\mathbf{J}$ with $|\mathbf{J}|<\alpha_0$. The real centre $U$ is required to have small curvature and small local potentials, but not a small current. The backgrounds of the averages of $\mathbf{U}$ are required to satisfy what the condition on the current provided, and $U^c_j$ is the set of all such pairs. This set is not saturated with respect to the $G^c$-valued $u$, which act by 

\begin{gather*}
\FitMath{(\mathbf{U},\mathbf{J})\to(\mathbf{U}^u,\operatorname{Ad}(u)\mathbf{J})}
\end{gather*}

. The extent to which a term is invariant under these transformations is given in \cref{prop:averaging-1a-3-9} and in \cref{ch:local-formation-program}.

\emph{Gauges.} Because $G$ is compact, the gauge is not fixed in order to obtain finiteness. The reason for fixing it is that the integrand of $\mathbf{T}$ does not change under the block group $\mathcal{G}_{\mathrm{bl}}=\{u:u=\mathbf{1}\text{ on }\mathbb{T}^{(k+1)}\}$. Consequently, in each block there are $3(L^4-1)$ directions $D_V\lambda$ along which the Hessian at the minimum vanishes, and there is then no Gaussian measure, no covariance and no decay. Similarly, the residual group $\mathcal{G}^{\mathrm{res}}$ leaves the constraints $\mathcal{M}^jU=V_j$ on $\Gamma_j$ unchanged. This group consists of the bare $u$ that equal $\mathbf{1}$ at the \emph{nodes}, which are the end points of the bonds of $\mathbf{B}[\Omega_\bullet]$. Hence a background is an orbit and not a point. No gauge has all of the following properties at once: locality (on a torus, no global gauge turns small curvature into a small potential), compatibility with the averages, covariance under the remaining gauge group, symmetry under the hypercubic group $W_4$, and regularity uniform in $k$. Axial gauges are algebraic and strictly local and have Jacobian one; they appear in the integrals and serve to select representatives. \emph{Box axial:} the condition $V=\mathbf{1}$ on the tree of combs that issues from the base point $y_P$ of a box. \emph{Block axial:} the condition $V(y,x)=\mathbf{1}$ on the averaged contour variables. It is imposed softly on a set $P''$ of blocks, at those pairs $(y,x)$ for which the mean $V(y,x)$ is defined (no factor is inserted at the other pairs), by means of

\begin{gather*}
\FitMath{\mathfrak{G}_k(V;P'')=\prod_{y\in P''}\prod_{x\in B(y)\setminus y}z_{\mathfrak{G}}^{-1}\,\mathbf{1}\{|V(y,x)-\mathbf{1}|<\varepsilon_k\}\,e^{-x_k(1-\frac12\operatorname{Re}\operatorname{tr}V(y,x))},} \\[3pt]
\FitMath{\mathrm{GF}_k(V)=\sum_{y,x}\big[1-\tfrac12\operatorname{Re}\operatorname{tr}V(y,x)\big]\ge0.}
\end{gather*}

\emph{Tower axial} relative to $U_0$, written $\operatorname{Ax}(\Omega_\bullet,U_0)$: in every pillar and at every level $n<j$, the contour variables of $\mathcal{M}^nU$ coincide with those of $\mathcal{M}^nU_0$. When the tops are the only nodes, solving this condition from the top down selects one representative in each orbit of $\mathcal{G}^{\mathrm{res}}$. An example, for $U_0=\mathbf{1}$, is the representative of the orbit $U_k(\mathbf{V})$ of the regular critical configurations of \cref{ch:multilevel-backgrounds}. Suppose that an inward bond of layer $j-1$, that is, a bond of $\mathbb{T}^{(j-1)}$ joining a point of layer $j-1$ to a point of $\Omega_j$, is retained as an additional constraint (see \cref{ch:multilevel-backgrounds}, \emph{Inward bonds}; by the Notation, \ensuremath{\S}4, the determining set contains no such bond). Its end point then lies inside a pillar of layer $j$. That end point is a node but not a top, and it is called a \emph{junction node}; in \cref{ch:multilevel-backgrounds} it is the seam node. The transport from the top to a junction node is invariant under $\mathcal{G}^{\mathrm{res}}$, so the condition cannot be achieved there. Leaving out these conditions gives the gauge $\operatorname{Ax}^\circ(\Omega_\bullet,U_0)$, which meets in exactly one point every orbit where the contour variables are defined. Representatives are covariant. Given a gauge transformation $v$ on the nodes, its lift is the bare transformation $\tilde{v}$ whose value at a bare site is the value of $v$ at the first node reached when one ascends from that site through the centres of its blocks. The map $W\mapsto W^{\tilde{v}}$ then sends the constraint set of $\mathbf{V}$ onto the constraint set of $\mathbf{V}^v$; it also sends orbits to orbits and maps $\operatorname{Ax}^\circ(\Omega_\bullet,\mathbf{1})$ onto itself. Hence, for the representative taken at the reference $U_0=\mathbf{1}$, $U_k(\mathbf{V}^v)=U_k(\mathbf{V})^{\tilde{v}}$. For nests of boxes whose lattices are of different levels \cref{sec:large-field-stored-variables-and-integration-in-a-large-field-patch}, the tree consists of shell trees, connected by one entrance bond in each box. Trees depend on the choice of a frame in $W_4$, and averaging over the $2^4\,4!$ frames restores the symmetry (1a.3).

The estimates are carried out in the elliptic Landau gauge, and this gauge comes with a determinant. We put $Q^{\sharp}=(Q^{\sharp}_j|_{\Gamma_j})_j$, with each component normalised to total weight $1+O(\alpha_0)$, so that $\mathcal{Q}^{\sharp}_j(U_0;A)=\xi_jQ^{\sharp}_jA+C_j(A)$. Further, $Q'=(Q'_j|_{\mathbb{Y}_j})_j$, and its null space is denoted $N(Q')$. On layer $j$ we set $\xi=\xi_j$ and $a=\xi_j^{-2}$. We also put $G'=(\Delta_{U_0}+Q'^*aQ')^{-1}$, and $R$ is the orthogonal projection onto $\Delta_{U_0}N(Q')$. Relative to $U_0$, a \emph{Landau pair} $(A,u)$ of $U$ is a pair satisfying

\begin{gather*}
\FitMath{U^{u^{-1}}=e^{i\eta_kA}U_0,} \\[3pt]
\FitMath{R\,D^*_{U_0}A=0,} \\[3pt]
\FitMath{\bar{R}^{(j)}_{U_0}u=\mathbf{1}\,\text{ on }\mathbb{Y}_j,}
\end{gather*}

in which conditions on averages take the place of the condition $u=\mathbf{1}$ at the nodes, since that condition cannot be handled analytically. This has a cost. Transformations $u$ of this kind do not form a group: $\bar{R}u_1=\bar{R}u_2=\mathbf{1}$ does not imply $\bar{R}(u_1u_2)=\mathbf{1}$. The word ``orbit'' therefore has no meaning, and the set that such transformations generate depends on its base point. For this reason the operations come in a fixed order. First, the tower axial gauge $\operatorname{Ax}^\circ$ selects a base point in each orbit of $\mathcal{G}^{\mathrm{res}}$. Next, the set is generated from that base point. Finally, the set is intersected with the slice. The average that belongs to this gauge is the \emph{reduced average} 

\begin{gather*}
\FitMath{\mathcal{M}_{\mathrm{red}}(W;U_0)=\mathcal{M}(W)^{\mathsf{n}_W^{-1}}\mathcal{M}(U_0)^{-1}=e^{i\mathcal{Q}^1(U_0,\mathsf{a})}}
\end{gather*}

; here $\mathsf{n}_W$ stands for the block factor $\mathsf{u}(U_0,\mathsf{a})$ taken at $W=e^{i\mathsf{a}}U_0$. The reduced average is the straightened average, under the name it bears in the theory of gauges. The contour variable is required to be the same in all four objects $\mathcal{M}$, $\mathsf{n}_W$, $\bar{R}$ and $\operatorname{Ax}$: it must be the mean in each of them, or the comb in each of them.

\textbf{Statements.} \begin{lemma}[{the mean}]\label{lem:averaging-1a-1-1} Consider families $V_\bullet=(V_1,\dots,V_n)$ in $G^c$ of small diameter $\operatorname{diam}V_\bullet:=\max_{i,i'}|V_iV_{i'}^{-1}-\mathbf{1}|$, where the radii do not depend on $n$. For such families the mean near the family is unique. It is holomorphic, with bounds independent of $n$ (by the contraction principle), and it is symmetric. On families in $G$ it takes values in $G$, and there it is the only mean lying in $G$. Here a mean of $V_\bullet$ is an element $V$ with $\max_i|V_iV^{-1}-\mathbf{1}|<1$ and $\sum_i\log(V_iV^{-1})=0$; the logarithm is the series logarithm, which has a meaning because of the first condition. Next, $\operatorname{av}(TV_\bullet T')=T\operatorname{av}(V_\bullet)T'$ holds for $T,T'\in G$. It also holds for $T\in G^c$ as long as $\operatorname{cn}(T)\operatorname{diam}V_\bullet$ remains below the radius, where $\operatorname{cn}(T):=|T|\,|T^{-1}|$ is the condition number. The identity $\operatorname{av}(V_\bullet^{-1})=\operatorname{av}(V_\bullet)^{-1}$ holds in $G$. It holds in $G^c$ when $\operatorname{diam}V_\bullet$ is below half the radius and $\operatorname{cn}(V_i)\le2$ for all $i$. Finally, $\frac1i\log\operatorname{av}(e^{iA_\bullet})=\frac1n\sum_jA_j+O(\max_j|A_j|^2)$.
\end{lemma}

\begin{lemma}[{small curvature: nearby transports, comb gauge}]\label{lem:averaging-1a-1-2} Consider a box that contains at most $s$ sites along each direction and has plaquette size $\varpi$. Suppose the $m$ steps of a path are rearranged so that the rearranged paths stay inside the box. Then the transport along the path changes by at most $(1+\varpi)^{m(m-1)/2}-1$. Moreover, with $h(x)=W(\Gamma_{y,x})$,

\begin{gather*}
\FitMath{|W^h(b)-\mathbf{1}|\le|b_--y|_1\,\varpi,} \\[3pt]
\FitMath{W^h=e^{i\mathsf{a}^0},} \\[3pt]
\FitMath{|\mathsf{a}^0|\le Cs\varpi.}
\end{gather*}

Consequently the contour families on $\mathrm{Dom}(c)$ are of diameter $\le CL^2\varpi_c$, and one has $\mathcal{M}(W)(c)=e^{iY}W[c]$ where $|Y|\le CL^2\varpi_c$.
\end{lemma}

\begin{proposition}[{the averaging map}]\label{prop:averaging-1a-1-3} Let $U$, $u$ be arbitrary, and let $\mathsf{r}$ be any isometry that preserves $\mathbb{T}^{(k+1)}$ (translations, $W_4$, coordinate and diagonal reflections). Then

\begin{gather*}
\FitMath{\mathcal{M}(U^u)=\mathcal{M}(U)^{u|_{\mathrm{coarse}}},} \\[3pt]
\FitMath{\mathcal{M}(\mathsf{r}U)=\mathsf{r}\,\mathcal{M}(U);}
\end{gather*}

and $\mathcal{M}^n(U)(c)$ depends on $U$ only in $B^n(c_-)\cup B^n(c_+)$. The averaging map has four further properties.
\begin{itemize}
\item It is Borel. Each of the three rules applies on a Borel set, and $\mathcal{M}$ is continuous on each of these sets but not across their boundaries. Outside $\mathrm{Dom}(c)$ no further regularity is required.
\item It is holomorphic at $e^{i\mathsf{a}}W_0$ for $\varpi_c(W_0),|\mathsf{a}|<c_L$.
\item It is onto: $\mathcal{M}^k(Q^{s*}_kV)=V$.
\item On small fields of one colour it is linear: $\mathcal{M}(e^{ifT})=e^{i(\mathsf{Q}^{\mathrm{pl}}f)T}$ whenever $T\in\mathfrak{g}$, $|T|=1$, and $f$ is real with $\sup|f|<c_L$.
\end{itemize}

For a $G^c$-valued $g$, covariance is obtained at the expense of room in the domain, and the amount of room lost is $\operatorname{cn}(g)$. Precisely, consider the graded domains $\mathrm{Dom}^c_\theta(c)$, $0<\theta\le1$, on which the contour families and the $W_x$ lie within $\theta$ times the radii. If $\operatorname{cn}(g(c_\pm))\le\varkappa$ and $\varkappa\theta\le1$, then $\mathbf{W}^g\in\mathrm{Dom}^c_{\varkappa\theta}(c)$ and

\begin{gather*}
\FitMath{\mathcal{M}(\mathbf{W}^g)(c)=g(c_-)\,\mathcal{M}(\mathbf{W})(c)\,g(c_+)^{-1}.}\qedhere
\end{gather*}



\end{proposition}

\begin{proposition}[{one averaging multiplies curvature by $L^2$}]\label{prop:averaging-1a-1-4}

\begin{gather*}
\FitMath{|\partial\mathcal{M}(W)(p')-\mathbf{1}|\le L^2\varpi+C(L^2\varpi)^2;} \\[3pt]
\FitMath{|\partial U-\mathbf{1}|\le\alpha_0\eta_k^2,\,\alpha_0\le c} \\[3pt]
\FitMath{\Rightarrow} \\[3pt]
\FitMath{\varpi(\mathcal{M}^nU)\le2\alpha_0(L^n\eta_k)^2} \\[3pt]
\FitMath{(n\le k),}
\end{gather*}

and $\mathcal{M}^nU(c)$ differs from the straight bare transport by at most $C\alpha_0(L^n\eta_k)^2$. Consequently the first rule is applicable on every level, and on $\{|\partial U-\mathbf{1}|<c\,\eta_k^2\}$ the map $\mathcal{M}^k$ is real-analytic.
\end{proposition}

\begin{proposition}[{one step relative to a reference}]\label{prop:averaging-1a-1-5} Let $\varpi=\varpi(W_0)\le c_L$, and consider complex $|\mathsf{a}|<c_L$. Then $\mathcal{Q}^1$ and $\frac1i\log\mathsf{u}$ depend holomorphically on $\mathsf{a}$, are covariant and do not exceed $C|\mathsf{a}|$ in size; moreover

\begin{gather*}
\FitMath{\mathcal{Q}^1(W_0,\mathsf{a})=\mathsf{Q}(W_0)\mathsf{a}+C^1(W_0,\mathsf{a}),} \\[3pt]
\FitMath{\mathsf{Q}(\mathbf{1})=\mathsf{Q}^{\mathrm{s}},} \\[3pt]
\FitMath{\|\mathsf{Q}(W_0)-\mathsf{S}(W_0)\|_{\mathrm{row},\,\mathrm{col}}\le C\varpi,}
\end{gather*}


\begin{gather*}
\FitMath{\|\mathsf{S}(W_0)\|_{\mathrm{row}}\le L,} \\[3pt]
\FitMath{\|\mathsf{S}(W_0)\|_{\mathrm{col}}\le L^{-3},} \\[3pt]
\FitMath{|C^1|\le C|\mathsf{a}|^2.}\qedhere
\end{gather*}



\end{proposition}

\begin{proposition}[{$k$ steps, uniformly in $k$}]\label{prop:averaging-1a-1-6} Assume $|\partial U_0-\mathbf{1}|\le\alpha_0\eta_k^2$ with $\alpha_0\le c$, and let $A$ be complex with $\sup|A|<c$. Here $c,C$ depend only on $L$. The following hold holomorphically in $A$ and in the complex shifts $\mathbf{U}=e^{i\eta_kA'}U_0$ of $U_0$ with $\sup|A'|<\alpha_1\le c$. The bounds are uniform in $A'$; in the bounds on the row norm and on the functional derivative, $\alpha_0+\alpha_1$ then replaces $\alpha_0$.

\begin{gather*}
\FitMath{\mathcal{M}^n(e^{i\eta_kA}U_0)=\big(e^{i\mathsf{a}_n}\mathcal{M}^nU_0\big)^{v_n},} \\[3pt]
\FitMath{|\mathsf{a}_n|\le2L^n\eta_k\sup|A|,} \\[3pt]
\FitMath{|v_n-\mathbf{1}|\le CL^n\eta_k\sup|A|,}
\end{gather*}


\begin{gather*}
\FitMath{Q^{\sharp}_k(U_0)=\eta_k\,\mathsf{Q}(W_{0,k-1})\cdots\mathsf{Q}(W_{0,0}),} \\[3pt]
\FitMath{\|Q^{\sharp}_k(U_0)\|_{\mathrm{row}}\le1+C\alpha_0,} \\[3pt]
\FitMath{|C_k|\le C\sup|A|^2,}
\end{gather*}


\begin{gather*}
\FitMath{Q_k(U_0)A=Q^{\sharp}_k(U_0)A-D_{\mathcal{M}^kU_0}\dot{v}_k(A),} \\[3pt]
\FitMath{|Q_kA|\le C\sup|A|,} \\[3pt]
\FitMath{|C^{\mathrm{pl}}_k|\le C\sup|A|^2:}
\end{gather*}

thus a coarse gradient separates the two linear parts, and this gradient does not become small as the curvature does. More generally, let $h$ be $G^c$-valued with $\operatorname{cn}(h)\operatorname{cn}(v_n)\le2$. Then $(e^{i\eta_kA}U_0)^h$ has well-defined averages, and $\mathcal{M}^n\big((e^{i\eta_kA}U_0)^h\big)=\big(e^{i\mathsf{a}_n}\mathcal{M}^nU_0\big)^{hv_n}$. This ball is much larger than the regular class, because the plaquettes of $e^{i\eta_kA}U_0$ are $O(\eta_k)$ rather than $O(\eta_k^2)$. Accordingly, the validity of the first rule on every level and the analyticity are consequences of the factorisation, not of \cref{prop:averaging-1a-1-4}. The remainder decomposes as $C_k=C_k^{(2)}+C_k^{(3)}+\cdots$, where

\begin{gather*}
\FitMath{C_k^{(2)}(U_0,A)=\sum_{n<k}\mathsf{Q}(W_{0,k-1})\cdots\mathsf{Q}(W_{0,n+1})\,C^{1,(2)}\big(W_{0,n},\mathsf{a}^{\mathrm{lin}}_n\big),} \\[3pt]
\FitMath{|n\text{-th term}|\le CL^{n-k}\sup|A|^2:}
\end{gather*}

the non-linearity consists of commutators and is concentrated at the coarsest of the $k$ scales. The derivative $K_A(c,b):=\partial\mathcal{Q}^{\sharp}_k(c)/\partial A(b)$ satisfies $\|K_A\|_{\mathrm{col}}\le2\eta_k^4$ and the Lipschitz bound $\|K_A-K_{A'}\|_{\mathrm{col}}\le C\eta_k^4\sup|A-A'|$. Consequently $1+C\alpha_0+C|A|$ is a bound for the functional derivative $\eta_k^{-4}K_A$, the adjoint $K_A^*$ is a bounded operator for the sup norm, and $K_A$ acts boundedly from $L^2$ to $\ell^2$. The Lagrange equations of \cref{ch:multilevel-backgrounds} involve $K_A^*$. Complex regular fields $\mathbf{U}=e^{i\eta_kA'}U_0$ with $\sup|A'|<\alpha_1\le c$ have averages that are regular, each on its own scale:
\begin{gather*}
\FitMath{|\partial\mathbf{U}-\mathbf{1}|\le\alpha_0'\eta_k^2,\quad\alpha_0'\le c}
\\[3pt]
\FitMath{\Longrightarrow\quad
 |\partial\mathcal{M}^n\mathbf{U}-\mathbf{1}|
       \le2\alpha_0'(L^n\eta_k)^2.}
\end{gather*}
Here the complex curvature bound is proved directly. Fix a level-$n$
plaquette $p_n$ and let $Y$ be the union of the four complete
backward blocks of its corners, with all bonds and plaquettes whose
vertices lie there. If this read wraps, lift its literal averaging formulas to a periodic
cover on which the read embeds. Put $\ell=L^n$, $t=\ell\eta_k$ and
$r=\alpha_0't^2$. The case $n=0$ is the hypothesis. For $n\ge1$,
use the coordinate comb in this box. Its paths have at most $C\ell$
bare links. The real factors are unitary, and each complex increment
and its inverse have norm at most $e^{\eta_k\alpha_1}$. The resulting
gauge $g$ therefore satisfies
\[
 \sup_x\operatorname{cn}(g(x))\le K:=e^{C\alpha_1t}.
\]
The strip between neighboring combs contains at most $C\ell$
plaquettes. Their conjugating paths have the same bound on their
condition numbers. The product and logarithm estimates give
\[
 \mathbf U^g=e^{ia},\qquad
       \|a\|_\infty\le C\ell\alpha_0'\eta_k^2
                         =Cr/\ell.
\]
All constants here depend only on $L$; the fixed ceiling is decreased
so that the products stay in their logarithm charts.

Extend $a$ by zero outside $Y$ only to use the flat supremum-ball
Taylor estimate; its supremum does not increase, and the four output
rows do not change. No curvature bound across this new boundary is
asserted. Let
$P_n(a)=i^{-1}\log\mathcal M^n(e^{ia})$ and let $L_n$ be its
linear part at zero. The flat recursion
\eqref{hb:actual-recursion} and its node factors give the uniform
holomorphic ball of \cref{wmc:full-plain-Lipschitz}. Cauchy on a
fixed larger ball therefore gives
\[
 \|P_n(a)\|_\infty\le Cr,\qquad
             \|P_n(a)-L_na\|_\infty\le Cr^2.
\]
Indeed its normalized input is $\ell a$, of size at most $Cr$, on
a fixed larger ball. Let $\mathrm d_0$ and $\mathrm d_n$ denote
oriented sums around fine and level-$n$ plaquettes, respectively.
In the one-step linear formula the two contour terms at each corner
cancel. Its curl is the average of the sums over the $L^2$ fine
plaquettes filling the corresponding square. Iterating this exact
identity gives a plaquette average with nonnegative weights and
total weight $\ell^2$. For the selected plaquette every fine
plaquette in this identity lies in $Y$, so
\[
 |(\mathrm d_nL_na)(p_n)|
       \le\ell^2\sup_{p\subset Y}|(\mathrm d_0a)(p)|.
\]
On precisely this stock the finite plaquette product gives
$\sup_{p\subset Y}|(\mathrm d_0a)(p)|
 \le K\alpha_0'\eta_k^2+C\|a\|_\infty^2$.
Apply the same product estimate to the four coarse logarithms.
Complex covariance then conjugates the result back by the gauge
at the base point. This covariance is valid throughout the retained
chart. The flat raw-angle recursion bounds the intermediate plain
logarithms by $CrL^{s-n}$, so every one-step contour family has a
fixed reserve. Conjugation by $g$ or $g^{-1}$ spends at most $K$
of that reserve, and $K$ is as close to one as required by decreasing
$c$. This admission uses the flat recursion, independently of the
complex curvature estimate being proved.
Consequently
\[
 |\partial\mathcal M^n\mathbf U-1|
       \le K^2r+Cr^2
       \le r\bigl(1+C\alpha_1t+Cr\bigr)\le2r.
\]
The choice is uniform because $0<t\le1$. The formulas commute with
pullback, so the same bound descends from the cover.
From $|\nabla_{U_0}A'|<\alpha_1$ one obtains only
$\alpha_0'=\alpha_0+C\alpha_1$. Thus the domains $U^c_j$ require,
as a separate condition, that the complex plaquettes be small;
the estimate just proved preserves their prescribed curvature scale.

All of this holds on determining sets, layer by layer. At the bare bonds of layer $j$ the local scale is $\xi(b)=\xi_j=L^j\eta_k$; it is the local scale of layer $j$ in the sense of ``Regular configurations''. On determining sets the ceilings on $\alpha_0$ and on $\sup_b\xi(b)|A(b)|$ are lowered by the factors $L^{-2}$ and $L^{-1}$. At a bond of $\Gamma_j$ that lies astride a seam, one factor is lost. Writing $K_A(c,b)=\partial\mathsf{a}_j(c)/\partial A(b)$ for $c\in\Gamma_j$, one has

\begin{gather*}
\FitMath{|\partial\mathcal{M}^jU_0-\mathbf{1}|\le2L^2\alpha_0,} \\[3pt]
\FitMath{|\mathsf{a}_j|\le2L\sup_b\xi(b)|A(b)|,} \\[3pt]
\FitMath{\sup_c\sum_b\frac{|K_0(c,b)|}{\xi(b)}\le2L,} \\[3pt]
\FitMath{\sum_j\sum_{c\in\Gamma_j}|K_A(c,b)|\le\frac{C\eta_k^4}{\xi(b)^3},}
\end{gather*}

and this factor is lost once, without compounding; it cannot be removed (\cref{prop:averaging-1a-1-14}). The last bound is a sum over the levels at which a given bare bond enters.
\end{proposition}

\begin{proposition}[{linear averages; gauge transformations}]\label{prop:averaging-1a-1-7} The operator 

\begin{gather*}
\FitMath{Q^{\mathrm{st}}_k:=Q^{\mathrm{st}}_k(\mathbf{1})=Q^{\sharp}_k(\mathbf{1})=\eta_k(\mathsf{Q}^{\mathrm{s}})^k}
\end{gather*}

 is the straight average taken over the blocks of order $k$, and

\begin{gather*}
\FitMath{(Q^{\mathrm{st}}_kA)(c)=\eta_k^5\sum_{x\in B^k(c_-)}A([x,x']),} \\[3pt]
\FitMath{u_k(p)=\prod_\mu\frac{\eta_k(e^{ip_\mu}-1)}{e^{i\eta_kp_\mu}-1},} \\[3pt]
\FitMath{c\le Q^{\mathrm{st}}_kQ^{\mathrm{st}*}_k\le1,}
\end{gather*}

where $u_k$ coincides with the Fourier symbol of $Q'_k$ up to a phase factor. The normalisation is as follows. Constants are mapped to the same constants, and $\|Q^{\mathrm{st}}_k\|_{\mathrm{row}}=1$. Composing $Q^{\mathrm{st}}_k$ with the linear lift gives the identity, and $Q'_kQ_k'^*=1$. The adjoint does not coincide with the lift. The operator $Q^{\mathrm{st}}_kQ^{\mathrm{st}*}_k$ is not a multiple of the identity: it is a convolution that couples each coarse bond to its two neighbours on the same line. This is why the display gives only a two-sided bound. In a background one has $|Q^{\sharp}_k(U_0)A|\le Q^{\mathrm{st}}_k|A|+C\alpha_0Q''_k|A|$, where $Q''_k$ takes the mean of $|A|$ over every bond in the double block of order $k$. With gradients $(\mathrm{d}\lambda)(b)=\lambda(b_+)-\lambda(b_-)$, the two averages intertwine gradients in different ways:

\begin{gather*}
\FitMath{\mathsf{Q}^{\mathrm{pl}}\mathrm{d}\lambda=\mathrm{d}(\lambda|_{\mathrm{coarse}}),} \\[3pt]
\FitMath{\mathsf{Q}^{\mathrm{s}}\mathrm{d}\lambda=\mathrm{d}(q\lambda),} \\[3pt]
\FitMath{Q^{\sharp}_k(U)D_U\lambda=D_{\mathcal{M}^kU}\,Q'_k(U)\lambda,}
\end{gather*}

the last of these relations being exact and valid for the straightened average only. In a background the plain average satisfies, likewise exactly, $Q_k(U)D_U\lambda=D_{\mathcal{M}^kU}(\lambda|_{\mathbb{T}^{(k)}})$, where the restriction stands in place of $Q'_k$. For the main term $Q^{\mathrm{st}}_k(U)$, by contrast, the relation involving $Q'_k$ is valid only up to $O(\alpha_0)$.

It follows that there are two residual invariances. The constraint $\mathcal{M}^kU=V$ is preserved by every transformation with $u=\mathbf{1}$ on $\mathbb{T}^{(k)}$. The straightened constraint is preserved by the $w$ with $r_k=\mathbf{1}$, where $r_k$ is defined below. In linearised form this condition is $Q'_k\lambda=0$. When $A=0$, and likewise at a field tower axial relative to $U_0$, these transformations are precisely the $w$ whose averages are trivial. Of the two sets, only the first forms a group. The set generated by the second depends on its base point, and the tower axial gauge selects that base point in each orbit of the first. The two invariances correspond to each other through $v_k[A]$.

Now let $A$ be real with $\sup|A|\le c$. Put $w=e^{i\lambda}$ and $e^{i\eta_kA^w}U_0:=(e^{i\eta_kA}U_0)^w$, and for a site function assume $\|\lambda\|_1:=\max\{\sup|\lambda|,\sup|D_{U_0}\lambda|\}\le c$. This norm is taken in level-$k$ units, without local scales, and it differs from the norm $\|\cdot\|_1$ of $1$-forms. It equals $\|\lambda\|_0$ for the constant sequence $\Omega_j=T_\ell$, whose only layer is $k$, with $\xi_k=1$. The straightened averages change by $r_n$, where $r_0=w$ and $r_{n+1}=\mathsf{u}(W_{0,n},\mathsf{a}_n(A^w))^{-1}\,r_n\,\mathsf{u}(W_{0,n},\mathsf{a}_n(A))$:

\begin{gather*}
\FitMath{e^{i\mathsf{a}_n(A^w)}W_{0,n}=\big(e^{i\mathsf{a}_n(A)}W_{0,n}\big)^{r_n},} \\[3pt]
\FitMath{w|_{\mathbb{T}^{(n)}}\,v_n[A]=v_n[A^w]\,r_n,} \\[3pt]
\FitMath{r_n|_{A=0}=\bar{R}^{(n)}_{U_0}w,} \\[3pt]
\FitMath{\tfrac1i\log r_k=Q'_k(U_0)\lambda+O\big((\sup|A|+\|\lambda\|_1)\|\lambda\|_1\big).}
\end{gather*}

We say that $u$ is of class $\alpha$ when each $\bar{R}^{(j)}u$ is within $\alpha$ of $\mathbf{1}$ and its variation inside any block is less than $\alpha L^{j+1}\eta_k$. If $u_1,u_2$ are of class $\alpha\le c$, the averages are multiplicative up to second order,

\begin{gather*}
\FitMath{\bar{R}^{(j)}(u_1u_2)=\bar{R}^{(j)}u_1\,\bar{R}^{(j)}u_2\,e^{i\mathsf{z}_j},} \\[3pt]
\FitMath{|\mathsf{z}_j|\le C(\alpha L^j\eta_k)^2,}
\end{gather*}

and $u_1u_2$ is of class $2\alpha+C\alpha^2$. The averages of products depend holomorphically on complex $\lambda$. Precisely, let $u'=e^{i\lambda}$ with $\|\lambda\|_1\le\alpha'\le c$, and let $u_1$ be of class $\alpha$. Then the quotients $\tilde{u}^{(j)}:=\bar{R}^{(j)}(u'u_1)(\bar{R}^{(j)}u_1)^{-1}$ have four properties: they are holomorphic in $\lambda$; they lie within $C\alpha'$ of $\mathbf{1}$; their covariant differences are $\le C\alpha'L^j\eta_k$; and $\frac1i\log\tilde{u}^{(k)}=Q'_k(U_0)\lambda+O((\alpha+\alpha')\alpha')$. 1a.3 uses this to show that the Landau gauge depends analytically on the fields.
\end{proposition}

\begin{lemma}[{rigidity of one step}]\label{lem:averaging-1a-1-9} Suppose that $W,W_0$ have plaquette size $\le\varpi$, that their contour variables agree up to a defect $\delta$, and that $|\mathcal{M}(W)\mathcal{M}(W_0)^{-1}-\mathbf{1}|\le\alpha_1$. Then

\begin{gather*}
\FitMath{|W(b)W_0(b)^{-1}-\mathbf{1}|\le\alpha_1+C\varpi+2\delta}
\end{gather*}

at every bond of the double block. Put differently, modulo the curvature, each bond variable is a function of the contour variables together with the average:

\begin{gather*}
\FitMath{\big|W(b)-W(y,b_-)^{-1}W(y,b_+)\big|\le C\varpi,} \\[3pt]
\FitMath{\big|W(b')-W(c_-,b'_-)^{-1}\,\mathcal{M}(W)(c)\,W(c_+,b'_+)\big|\le C\varpi,}
\end{gather*}

for $b$ inside $B(y)$ with $y=c_\pm$, and for $b'$ in the corridor of $c$, that is, among the bonds parallel to $c$ joining $B(c_-)$ to $B(c_+)$. The defect has two sources: soft gauge fixing gives $|V(y,x)-\mathbf{1}|<\varepsilon_k$ instead of equality, and seams contribute the defect of \cref{lem:averaging-1a-3-11}.
\end{lemma}

\begin{proposition}[{the fibre is curved; the current enters}]\label{prop:averaging-1a-1-10} Let $|\partial U-\mathbf{1}|\le\alpha_0\eta_k^2$ with $\alpha_0\le c$ (the range of \cref{prop:averaging-1a-1-6}). Let moreover $U\in\mathfrak{U}_k(\Omega_\bullet,\alpha_0)$ at the constant sequence $\Omega_j=T_\ell$; this membership brings in the further condition $|J(U)|<\alpha_0$. Suppose that $M_1$, $M\alpha_0$ and the number of sites of $\mathbb{T}^{(k)}$ satisfy the hypotheses of \cref{lem:averaging-1a-3-4} [stated technical lemma; Group A], which provides $G$ and $\mathsf{H}_k$. Finally, let $A$ be small. On $Q^{\sharp}_k(U)A+C_k(U,A)=0$, in the Landau gauge, one has $A=A_1-\mathsf{H}_kC_k^{(2)}(A_1)+\dots$ with $Q^{\sharp}_k(U)A_1=0$, and

\begin{gather*}
\FitMath{A(e^{i\eta_kA}U)=A(U)+\langle J(U),A_1\rangle+\tfrac12\langle A_1,\Delta^{\mathrm{W}}_UA_1\rangle-\langle\mathsf{H}_kC_k^{(2)}(A_1),J(U)\rangle+O(|A_1|^3).}
\end{gather*}

In this expansion $\mathsf{H}_k$ denotes the right inverse $\mathsf{H}=GQ^*(QGQ^*)^{-1}$ of $Q=Q^{\sharp}_k(U)$, built from the $G$ of \cref{sec:propagators-propagators-in-a-background-field-exponential-decay-produced-by-the}; with this choice the relation $RD^*GQ^*=0$ is exact. For the derivative of $V\mapsto U_k(V)$ the corresponding objects are $G_1$, $\mathsf{H}_1$. The term $O(|A_1|^3)$ collects all terms of order $\ge3$ in $A_1$. Indeed, the constraint is solved up to second order, and the solution is substituted into the Taylor expansion of the action.
\end{proposition}

\begin{proposition}[{isometries and linearised objects}]\label{prop:averaging-1a-1-11} Suppose that $\mathsf{r}$ preserves every lattice involved. Then $\mathsf{r}(e^{i\mathsf{a}}W_0)=e^{i\mathsf{r}_{W_0}\mathsf{a}}\,\mathsf{r}W_0$, and

\begin{gather*}
\FitMath{\mathcal{Q}^1(\mathsf{r}W_0,\mathsf{r}_{W_0}\mathsf{a})=\mathsf{r}_{\mathcal{M}(W_0)}\mathcal{Q}^1(W_0,\mathsf{a}),} \\[3pt]
\FitMath{Q^{\sharp}_k(\mathsf{r}U)\,\mathsf{r}_U=\mathsf{r}_{\mathcal{M}^kU}\,Q^{\sharp}_k(U),} \\[3pt]
\FitMath{J(\mathsf{r}U)=\mathsf{r}_UJ(U),} \\[3pt]
\FitMath{\Delta^{\mathrm{W}}_{\mathsf{r}U}\mathsf{r}_U=\mathsf{r}_U\Delta^{\mathrm{W}}_U,}
\end{gather*}

and the corresponding relations hold for $Q_k$, $Q'_k$, $\bar{R}$, $\mathsf{u}$. For real $W_0$, the map $\mathsf{r}_{W_0}$ is an isometry both for the sup norm and for $\langle\cdot,\cdot\rangle$. These maps compose according to $(\mathsf{r}\mathsf{r}')_{W_0}=\mathsf{r}_{\mathsf{r}'W_0}\mathsf{r}'_{W_0}$. If the twisted action is replaced by the plain one (the action at $W_0=\mathbf{1}$), these relations in general fail for a reflection reversing bonds. The plain and the twisted action are equal where $\operatorname{Ad}(W_0(b))$ acts as the identity on the reversed bonds $b$. Consequently, propagators, $\mathsf{H}_k$, $\Delta^{(k)}$ and $C^{(k)}$ are covariant under all of $W_4$ only in the twisted sense, and \cref{sec:propagators-propagators-in-a-background-field-exponential-decay-produced-by-the}, \cref{ch:local-formation-program} and \cref{ch:pointed-hessian-cancellation} use the symmetry in this form.
\end{proposition}

\begin{proposition}[{the fundamental equality; the transformation to the tower axial gauge}]\label{prop:averaging-1a-1-12} Let $U_0$, $A$ be as in \cref{prop:averaging-1a-1-6}. There exists exactly one $u$, defined on the bare sites, for which $(e^{i\eta_kA}U_0)^u$ is tower axial relative to $U_0$ and $\bar{R}^{(k)}_{U_0}u=\mathbf{1}$. Its averages are given by $\bar{R}^{(n)}_{U_0}u=u\,v_n$, so that $u=v_k^{-1}$ on $\mathbb{T}^{(k)}$, and

\begin{gather*}
\FitMath{\mathcal{M}^k\big((e^{i\eta_kA}U_0)^u\big)=e^{i\mathcal{Q}^{\sharp}_k(U_0;A)}\,\mathcal{M}^kU_0:}
\end{gather*}

once the gauge has been changed in this way, averaging of order $k$ amounts to composing the straightened one-step averages $k$ times. The transformation can be written explicitly. Let $x=x_0,\dots,x_k=y$ be the tower of block centres. Let $\mathsf{w}^{\mathrm{rel}}_j(z,x)=(e^{i\mathsf{a}_j}W_{0,j})(z,x)W_{0,j}(z,x)^{-1}$ be the relative contour variables of the straightened averages, let $\mathsf{u}_j$ be the block factors, and let $\mathsf{T}_{U_0}$ be the transport down the tower. Then

\begin{gather*}
\FitMath{u(x)=\prod_{j=k-1}^{0}\operatorname{Ad}\big(\mathsf{T}_{U_0}(x_{j+1},x)\big)^{-1}\Big[\mathsf{u}_j(x_{j+1})^{-1}\,\mathsf{w}^{\mathrm{rel}}_j(x_{j+1},x_j)\Big],} \\[3pt]
\FitMath{|u-\mathbf{1}|\le C\sup|A|.}
\end{gather*}

Consequently $u$ depends locally and holomorphically on $A$ and is of class $C\sup|A|$. At the bare scale, however, it fails to be regular: it jumps by $O(L^j\eta_k\sup|A|)$ between neighbouring bare sites belonging to different blocks of order $j$.
\end{proposition}

\begin{proposition}[{two schemes; the linear parts compared; what carries over}]\label{prop:averaging-1a-1-13} Ba{\l}aban's theorems of 1984--85 are proved for blocks that have their site at a corner and for single combs; this is scheme $\mathsf{c}$. The map $\mathcal{M}$ uses centred blocks and means; this is scheme $\mathsf{m}$. Since $L$ is odd, one translation of his entire setting, by $\mathsf{e}_j=\frac12(L^j-1)a_K(1,1,1,1)$ at level $j$, carries his blocks onto the centred blocks. Of the contour variable, the proofs of \cref{lem:averaging-1a-1-2}--\cref{lem:averaging-1a-1-9} use only the following properties:
\begin{enumerate}
\item there is one contour variable for each pair $(y,x)$, and it is covariant;
\item its paths stay in the block and make fewer than $L$ steps in each direction;
\item coarse translations carry the contour variables into one another;
\item it is $G$-valued on $G$;
\item it is holomorphic near $\mathbf{1}$, with a mean of the $\mathsf{a}(\Gamma)$ as its linear part;
\item where the bond variables commute with one another, it is the exponential of that mean.
\end{enumerate}

This list is one way of carrying out the axiomatisation that Ba{\l}aban mentions. The results of \cref{lem:averaging-1a-1-2}--\cref{lem:averaging-1a-1-9} therefore hold in both schemes, with two exceptions. The first exception is covariance under isometries other than coarse translations. It fails for combs, since under an element of $W_4$ a comb goes over into a comb of a different order, and a lowest corner into a different corner; this failure is precisely why the mean is used. The second exception is the three rules, which are formulated for $\mathcal{M}$ only. When the background is trivial, the two schemes have the same straightened linear part, $\mathsf{Q}_\star(\mathbf{1})=\mathsf{Q}^{\mathrm{s}}$, but different plain linear parts. In a background, let $t$ denote a transport of $U_0$ joining the two base sites of a block. Then

\begin{gather*}
\FitMath{Q^{\sharp}_k(U_0)=\operatorname{Ad}(t)\big(Q^{\mathrm{B}}_k(U_0)+\mathsf{D}_k\big),} \\[3pt]
\FitMath{\|\mathsf{D}_k\|_{\mathrm{row}}\le C\alpha_0,} \\[3pt]
\FitMath{\|\mathsf{D}_k\|_{\mathrm{col}}\le C\alpha_0\eta_k^4,}
\end{gather*}

and the same holds for $Q'_k$. Moreover, $\mathsf{D}_k(c,b)=0$ unless $b$ lies under $c$. In each scheme, the linear part and the straight average of that scheme differ by a kernel of row norm $C\alpha_0$. In this leading part, a bare bond $b$ parallel to $c$ receives as its weight $\eta_k^5$ multiplied by the count of segments $[x,x']$, $x\in B^k(c_-)$, passing through $b$ (a tent along the direction of $c$). Every such segment enters with the transport of $U_0$ taken down a chain of contours and segments. The exact identity $Q_k=Q^{\sharp}_k-D\dot{v}_k$ is supplied in the present text. The identity involving $\mathsf{D}_k$ takes the place of the assertion of 1987 (109, p. 254) that the new average satisfies all the results of the paper on averaging operations. For real $U_0$ the factor $\operatorname{Ad}(t)$ is orthogonal, so it cancels from $Q^{\sharp}_k(U_0)^*\,a\,Q^{\sharp}_k(U_0)$. Consequently, in the two schemes the forms defining $G'$, $R$, $\mathsf{G}_k$, $\mathsf{H}_k$ differ from each other by local forms whose relative size is $C\alpha_0$. Through this comparison, the operators built from $Q^{\mathrm{B}}_k(U_0)$ are used in \cref{ch:multilevel-backgrounds}. The operators are built afresh from $Q^{\sharp}_k(U_0)$, and their bounds are derived from the case $U_0=\mathbf{1}$, in which, after the blocks have been aligned, the two linear parts are equal \cref{sec:propagators-propagators-in-a-background-field-exponential-decay-produced-by-the}.
\end{proposition}

\begin{proposition}[{sharpness: fields of one colour}]\label{prop:averaging-1a-1-14} Let $T\in\mathfrak{g}$ with $|T|=1$. For $U=e^{ifT}$, with $f$ real and small (\cref{prop:averaging-1a-1-3}), everything is linear and explicitly computable. A constant curvature $\varphi$ gives $\partial\mathcal{M}U(p')=e^{iL^2\varphi T}$, so the factor $L^2$ of \cref{prop:averaging-1a-1-4} is attained.

Now consider a seam. Let $\Omega_1$ be, locally, a slab $\{x_1\ge\text{const}\}$. Let $z\in\mathbb{T}^{(1)}$ be the last coarse site outside it, and let $p'$ be the coarse plaquette with corners $z,z+Le_1,z+Le_1+Le_2,z+Le_2$. Set $f(\langle x,x+e_2\rangle)=\mathsf{f}^\Sigma(x_1)$ and $f=0$ on the other bonds, where $\mathsf{f}^\Sigma(t):=\sum_{t'<t}\mathsf{f}(t')$. Here $\mathsf{f}=\varphi$ on the $L-1$ bonds in direction $1$ that lie in the block of $z$; on the $L-1$ bonds preceding them the value is $-\varphi$, and everywhere else it is $0$. The curvature has size at most $\varphi$ and vanishes in $\Omega_1$, while $\partial\mathcal{M}U(p')=e^{i\Theta T}$ with $\Theta=\frac12L(L-1)\varphi$. Thus no bound $C\alpha_0$ with $C$ independent of $L$ can take the place of the bound $2L^2\alpha_0$ of \cref{prop:averaging-1a-1-6}.

At $U_0=\mathbf{1}$, take a bond $c\in\Gamma_j$ astride a seam with outer end $c_-$. Such a bond has one end, its outer end, outside $\Omega_j$; in \cref{ch:multilevel-backgrounds} it is called an outward bond. The weighted row sum equals $\sum_b\frac{|K_0(c,b)|}{\xi(b)}=\tfrac12L(1-L^{-j})+\tfrac12(1+L^{-j})$, and its limit as $j\to\infty$ is $\frac12(L+1)$. Indeed, along a segment $[x,x']$, a bare bond contained in $B^j(c_-)$ has $\xi(b)=L^{j-1}\eta_k$, while the remaining ones extend into $B^j(c_+)$, where the scale equals $L^j\eta_k$.

For a bond of $\Gamma_j$ astride a seam, an additional block factor arises at its outer end. Moreover, the datum entering the straightened constraint (the reduced datum of \cref{ch:multilevel-backgrounds}) is not the same as the plain datum. The factor is genuinely present. At $U_0=\mathbf{1}$, take $\eta_kA=tT$ on the single bare bond $\langle z,z+e_\nu\rangle$, where $z$ is the outer end of a bond of $\Gamma_1$, and $A=0$ elsewhere. Then $\mathsf{u}(z)=\exp\big(it\frac{L^4-1}{8L^4}T\big)\ne\mathbf{1}$. Finally, $\|\mathsf{Q}^{\mathrm{pl}}\|_{\mathrm{row}}>2L$.
\end{proposition}

\begin{proposition}[{the exact identities}]\label{prop:averaging-1a-2-1} The transformation $\mathbf{T}$ preserves mass, positivity ($|\mathbf{T}\nu|\le\mathbf{T}|\nu|$) and gauge invariance. It commutes with the isometries that map $\mathbb{T}^{(k+1)}$ onto itself, and

\begin{gather*}
\FitMath{\mathbf{T}^n\nu=(\mathcal{M}^n)_*\nu,} \\[3pt]
\FitMath{\langle F(\mathcal{M}^kU)\rangle_{K,\ell,\beta}=\int F\,d(\mathbf{T}^k\rho_0)\Big/\int d(\mathbf{T}^k\rho_0).}
\end{gather*}

Consequently, whenever $F\in\mathfrak{A}$ has largest depth $s$, each of the numbers $\langle F\circ\mathcal{M}^{K-s}\rangle_{K,\ell,\beta}$, $K\ge s$, is the integral of one fixed function with respect to the normalised measures $\mathbf{T}^{K-s}\rho_0$, all on the same compact space $\mathcal{V}(\Lambda_{s,\ell})$.
\end{proposition}

\begin{proof}
The Borel property of the actual average follows from its three
defining rules. Fix a coarse bond \(c\), and put
\[
 \mathcal N_c=
 \bigcup_{x,\Gamma,\Gamma'}
 \{U:Y_{x,\Gamma,\Gamma'}(U)=-\mathbf1\},
\]
with the finite index sets used in the definition of the average.
Each holonomy is continuous, so \(\mathcal N_c\) is closed.
The small-curvature construction makes the first rule valid on the
open set \(\mathrm{Dom}(c)\). The other two rules apply respectively
on
\[
 \mathrm{Dom}(c)^c\setminus\mathcal N_c,\qquad
 \mathrm{Dom}(c)^c\cap\mathcal N_c.
\]
The first rule is analytic; the second is continuous on its domain
because the intrinsic logarithm is continuous on
\(G\setminus\{-\mathbf1\}\); the last is the continuous straight-path
transport \(U[c]\). Thus the bond average is Borel on this finite
partition. The configuration space has finitely many independent
positive-bond coordinates, so \(\mathcal M\) is a Borel map on the
whole configuration space.

In particular \(K(U,\cdot)=\delta_{\mathcal M(U)}\) is a probability
kernel: \(K(U,B)=\mathbf1_B(\mathcal M(U))\) is Borel for each Borel
set \(B\). Its defining identity is
\[
 \int F(V)\,K(U,\mathrm{d}V)=F(\mathcal M(U))
\]
for every bounded Borel \(F\). Integration against the input measure
gives the asserted pushforward identity. Kernel normalization proves
mass preservation, and positivity gives
\(|\mathbf T\nu|\le\mathbf T|\nu|\). Iterating the same identity gives
\(\mathbf T^n\nu=(\mathcal M^n)_*\nu\).
Gauge covariance gives gauge invariance of the pushforward, since
every coarse gauge transformation extends to a fine one.
The same argument with lattice isometries gives their commutation
with \(\mathbf T\). Apply the iterated identity to
\(\nu=\rho_0\,\mathrm{d}U\) and divide by its positive total mass to
obtain the expectation formula.
\end{proof}

\begin{proposition}[{soft and hard constraints}]\label{prop:averaging-1a-2-2} Each kernel $t$ of the kind described above preserves mass, positivity and gauge invariance. In the case of linear site fields and of block means $Q$, the Gaussian kernels $t_{\mathfrak{a}}\propto e^{-\frac{\mathfrak{a}}2|\Psi-Q\Phi|^2}$, taken after the canonical rescaling by $L^{(d-2)/2}$, form a family that is closed under composition. In these kernels $\Phi$, $\Psi$ are site functions in dimension $d$, $Q$ is the rescaled block mean, and $Q_k$ is the mean over $B^k$, rescaled in the same way; the composition law reads

\begin{gather*}
\FitMath{t_{\mathfrak{a}}*t_{\mathfrak{a}_k}=t_{\mathfrak{a}_{k+1}},} \\[3pt]
\FitMath{\mathfrak{a}_{k+1}^{-1}=\mathfrak{a}^{-1}+L^{-2}\mathfrak{a}_k^{-1},} \\[3pt]
\FitMath{\mathfrak{a}_k\to\mathfrak{a}(1-L^{-2}),}
\end{gather*}

and these kernels give rise to the propagators $(\Delta+\mathfrak{a}_kQ_k^*Q_k)^{-1}$, where $\Delta=D^*D$ acts on site functions. For the delta kernel the fibres are nested, and over regular data, up to gauge,

\begin{gather*}
\FitMath{\{\mathcal{M}^{k+1}U=W\}=\bigsqcup_{\mathcal{M}(V)=W}\{\mathcal{M}^kU=V\}} \\[3pt]
\FitMath{\Rightarrow} \\[3pt]
\FitMath{U_k(V^{\mathrm{bg}}_k)=U_{k+1}(W),} \\[3pt]
\FitMath{V^{\mathrm{bg}}_k=\mathcal{M}^kU_{k+1}(W),} \\[3pt]
\FitMath{U_j(\mathcal{M}^jU_k(V))=U_k(V)\,\,(j\le k):}
\end{gather*}

every background therefore coincides with the background of its own averages, and $V^{\mathrm{bg}}_k$, with $\mathcal{M}^kU$ taken in the block axial gauge, is in that gauge the critical point of $V\mapsto A(U_k(V))$ on $\{\mathcal{M}(V)=W\}$. Suppose that $\Delta\ge0$ is positive on $\ker Q$ and that $Q$ is surjective. Then, under the constraint $QA=B$, the form $\frac12\langle A,\Delta A\rangle$ attains its minimum at $\mathsf{H}B$, where $\mathsf{H}=\mathsf{G}Q^*(Q\mathsf{G}Q^*)^{-1}$ and $\mathsf{G}=(\Delta+aQ^*Q)^{-1}$. This holds for every $a>0$. \emph{Remark.} Let $W$ be real and lie in a neighbourhood of $\mathbf{1}$; this neighbourhood may depend on $k$ and on the torus. Then $V^{\mathrm{bg}}_k$ is, in that gauge, the unique point of minimum of $V\mapsto A(U_k(V))$ on $\{\mathcal{M}(V)=W\}$.
\end{proposition}

\begin{proposition}[{the steps with their insertions are exact}]\label{prop:averaging-1a-2-3} Each operation of $\mathbf{T}$ and of $\mathbf{R}$ consists of finitely many positive kernels, one for each outcome of the operation, whose sum is a probability kernel. These kernels depend on the auxiliary choices. Once the choices are fixed, the kernels depend neither on $\beta$ nor on the density. When the run makes its own choices, the operations themselves vary with $\beta$. For a bounded, gauge-invariant, complex-valued $\mathfrak{w}$,

\begin{gather*}
\FitMath{\tau^{\mathfrak{w}}_h(dV)=\int dU\,\mathfrak{w}(U)\rho_0(U)\,\mathsf{K}_h(U,dV),} \\[3pt]
\FitMath{\sum_h\mathsf{K}_h(U,\mathcal{V}_k)=1.}
\end{gather*}

Consequently $\tau_h\ge0$ and $|\tau^{\mathfrak{w}}_h|\le\sup|\mathfrak{w}|\,\tau_h$. Moreover $\int\rho_k=Z_{K,\ell,\beta}$ and $\int\rho^w_k=Z(w)$ at every level, on a torus that is large compared with the cubes of side $M\check{R}_j$ of every level. This holds irrespective of the gauges and representatives chosen. It also holds for every real sequence of numbers $\mathsf{t}$ among those admitted in \cref{sec:large-field-the-induction-closed-the-order-of-the-constants-the-sum}. These are the sequences whose terms are all $\ge\gamma^{-2}$, whose differences obey two-sided bounds of the same type as the flow inequalities of \cref{prop:averaging-1a-2-7}, and for which $\check{R}_j\le L\check{R}_{j+1}$. The sum over outcomes of the operations of $\mathbf{T}$ is $\mathcal{M}_*$. By contrast, $\mathbf{R}$ preserves only those integrals of $F$ that do not depend on the components on which it acts; in general, therefore, $\rho_k\ne\mathbf{T}^k\rho_0$. \emph{Exact insertion:} for every gauge-invariant $F$,

\begin{gather*}
\FitMath{\langle F(\mathcal{M}^kU)\rangle_{K,\ell,\beta}=\sum_h\int F^{(h)}d\tau_h\Big/\sum_h\int d\tau_h,} \\[3pt]
\FitMath{|F^{(h)}|\le\sup|F|,}
\end{gather*}

where $F^{(h)}$ denotes the conditional expectation, under the joint law constructed from the kernels, of $F(\mathcal{M}^kU)$ given $(h,V_k)$. If $F$ depends only on fields away from those components, then $F^{(h)}=F$. \emph{Covariance:} $\mathsf{r}_*\tau_h=\tau_{\mathsf{r}h}$ for every isometry that preserves all the partitions involved. An individual $E^{(j)}(X,\cdot)$ is covariant under the translations that preserve $\pi_j$. An individual $R^{(j)}(X,\cdot)$ is covariant under the translations that preserve the $M\check{R}_j$-cubes. Covariance under all isometries $\mathsf{r}$ of $\mathbb{T}^{(j)}$ holds only for the unrestricted pointed totals $\sum_{X\ni z}E^{(j)}(X,U_j,z)$. These totals are taken at $U_j=U_j(V)$, where $V$ and $\mathsf{r}V$ both belong to a class of real block fields of level $j$ lying near the regular ones \cref{ch:local-formation-program}. Wherever the covariance of pointed totals is invoked below, it has this meaning.
\end{proposition}

\begin{lemma}[{units; the delta function resolved; regrouping}]\label{lem:averaging-1a-2-4} When the plaquette variables are close to $\mathbf{1}$, $\xi^{d-4}a_p=\frac12\xi^d\langle F_U,F_U\rangle(1+O(\xi^4|F_U|^2))$. Hence, for $d=4$, Wilson's action is a Riemann sum for $\frac12\int|F|^2$ whatever the unit of length. A change of unit leaves $A$, $\rho_k$ and $x_k$ unchanged. On small fields, $\mathcal{M}(e^{i\mathcal{A}}V^{\mathrm{bg}}_k)V_{k+1}^{-1}=\exp i\mathcal{Q}_1(L\mathcal{A})$ and $\mathcal{Q}_1(L\mathcal{A})=LQ\mathcal{A}+C^{\mathrm{pl},(2)}(\mathcal{A})+\dots$. Here $\mathcal{A}$ carries no factor of length. Up to the scaling of its argument by $L$, the map is $\mathcal{Q}_1(\cdot)=\mathcal{Q}_1(V^{\mathrm{bg}}_k;\cdot)$, and this scaling accounts for the factor $L$ in front of its linear part. The average here is the plain $Q=Q_{V^{\mathrm{bg}}_k}$, and $C^{\mathrm{pl},(2)}$ is the second-order part of the plain remainder $C^{\mathrm{pl}}_1(V^{\mathrm{bg}}_k,L\mathcal{A})$. The symbol $C^{(2)}_k$ is reserved for the straightened average.

An analytic substitution $\mathcal{A}=\psi^{\mathrm{lin}}(\mathcal{A}')=\mathcal{A}'-\mathsf{h}\tilde{D}(\mathcal{A}')$ changes only the variables on the pivot bonds. It turns the delta function into the linear one, $\delta(LQ\mathcal{A}')$, which is solved for a single pivot bond $b_0(c)$ for each coarse bond. The substitution comes from the implicit function theorem. In it, $\mathsf{h}$ inverts $LQ$ on the pivot bonds, $LQ\mathsf{h}=1$. The linearising correction $\tilde{D}$ is of second order and is written $\widetilde{\mathsf{D}}$ in \cref{ch:local-formation-program}. It is the unique small solution of $\mathcal{Q}_1\big(L(\mathcal{A}'-\mathsf{h}\tilde{D}(\mathcal{A}'))\big)=LQ\mathcal{A}'$, obtained as the fixed point of a contraction. For $F$ supported in this regular chart, and with the free integral restricted to its inverse image, the local coarea identity is

\begin{gather*}
\FitMath{\int dV_k\,\delta(\mathcal{M}(V_k)V_{k+1}^{-1})F(V_k)=z^{(k)}\int\prod_{b\ne b_0(c)}d\mathcal{A}'(b)\;e^{\operatorname{Tr}\log(1-\mathsf{h}\,\partial\tilde{D})+\sum_b\log\sigma_{\mathrm{H}}(\mathcal{A}(b))}\,F\big(e^{i\psi^{\mathrm{lin}}(\mathsf{C}\mathcal{A}')}V^{\mathrm{bg}}_k\big),}
\end{gather*}

where $\sigma_{\mathrm{H}}$ is the density of Haar measure in exponential coordinates. Once this density is divided by its value $\sigma_{\mathrm{H}}(0)$ at $\mathbf{1}$, the Jacobians are $1+O(\delta_k)$ for each bond. They are local and analytic, and they are included in $\mathbb{P}^{(k)}$. The factor $\sigma_{\mathrm{H}}(0)$, taken once for each non-pivot bond, is absorbed into the constant of the step. For every number $b$,

\begin{gather*}
\FitMath{-x_kA+\sum_XE^{(k+1)}(X,\cdot)=-(x_k-b)A+\Big(\sum_X\big[E^{(k+1)}(X,\cdot)-E^{(k+1)}(X,\mathbf{1})\big]-bA\Big)+\sum_XE^{(k+1)}(X,\mathbf{1}),}
\end{gather*}

and, summed over the levels, $-A(x_k(\cdot),U)-\sum_{j\le k}b_jA(\varphi_j,U)=-\beta A(U)$. The constant of a step depends on the regions of that step. Let $|\Omega^{(k)*}|$ be the number of non-pivot bonds of $\mathbb{T}^{(k)}$ in $\Omega$, and $\#_{k+1}(\Omega)$ the corresponding number of sites of $\mathbb{T}^{(k+1)}$. Let $z_{\mathfrak{G}}=\mathfrak{z}(x_k,\varepsilon_k)$ be the normaliser of $\mathfrak{G}_k$ for each block site (1a.3), and $\nu_{\mathfrak{g}}=\log\int_{\mathfrak{g}}e^{-|Y|^2/2}dY$. Then

\begin{gather*}
\FitMath{\mathsf{E}^{(k)}(\Omega,\Lambda)=\big(\log\sigma_{\mathrm{H}}(0)-\tfrac12(N^2-1)\log x_k\big)|\Omega^{(k)*}|-(L^4-1)\log z_{\mathfrak{G}}\,\#_{k+1}(\Omega)+\nu_{\mathfrak{g}}|\Lambda^{(k)*}|+\sum_{X\subset\Lambda}\big[E^{(k+1)}+R^{(k+1)}\big](X,\mathbf{1}),}
\end{gather*}

together with $-\mathcal{E}_k=\sum_{j<k}\mathsf{E}^{(j)}(\Omega_{j+1},\Lambda_{j+1})$ and $E=\sum_{j<K}\mathsf{E}^{(j)}(T_\ell,T_\ell)$. Here $\mathcal{E}_k$ depends on the label, and at the all-small-field label the vacuum constant $E$ equals $-\mathcal{E}_K$.
\end{lemma}

\paragraph{Primitive hypotheses for the joint step}
\label{rg:primitive-closure-premises}
Write \(\mathcal H_{\mathrm{step}}\) for the following inputs on every
actual determining region, stored-field fibre and allowed parameter
domain. Retain the background and substitution admission in A1--A4
of \cref{ch:dependency-catalogue}, and the actual matrix/core, potential
and window estimates used in \cref{ch:local-formation-program}.
Their operations and local families must obey the declared covariance
and common-lift compatibility rules.

For the history operations retain the common-core integrable envelope,
kept Gaussian integral and owned failure-energy certificates in
\eqref{hi:pair-sup-envelope}, \eqref{hi:reference-gaussian-integral}
and \eqref{hi:event-energy} of \cref{ch:history-mass-program}.
Use the stated label representation, common vacuum normalization and
finite read ownership. Drop-out evaluations must preserve their failure
certificates under the supremum inside the integral. In the sourced
application also retain \eqref{hi:source-excess-budget} on the same
fixed representation, or the ordinary complex/prepaid real extension
hypotheses stated there that replace it.

The formation, geometric counting and integration arguments give their
conditional successor bounds under these inputs. Any primitive estimate
not yet verified for all actual histories remains a technical hypothesis
of the closure assertions below. Class invariance C4 is their target,
not an assumption used to prove them. Exact transport and mass
preservation follow separately from the normalized measurable moves
in C1. A homogeneous regular-sector construction alone does not verify
the all-history premises in \(\mathcal H_{\mathrm{step}}\).

\begin{proposition}[{the invariant form}]\label{prop:averaging-1a-2-7}
Assume \(\mathcal H_{\mathrm{step}}\).
Let \(\mathfrak F_k\) denote the complete presentations in
\(\mathfrak C_k(\theta)\), written through their represented measures
of the form below. Membership includes the specified domains, stored
variables, measures, marks and reading rules. The step acts on this
presentation; \(\mathbf R\mathbf T\mathfrak F_k\) below denotes its
successor presentations. In it the $E^{(j)}$ are pointed families, the pointed totals of which are covariant with respect to every isometry of $\mathbb{T}^{(j)}$, and $b_j$ is their coefficient \cref{ch:pointed-hessian-cancellation}. A step does not preserve the form and the bounds alone \cref{sec:large-field-the-induction-closed-the-order-of-the-constants-the-sum}. The measures of the class, with their bounds, are

\begin{gather*}
\FitMath{\sum_{\mathcal{S}}\mathbf{T}'_k(\mathcal{S})\big[\chi_{k,\mathcal{S}}\,e^{A_{k,\mathcal{S}}(U_k(\mathbf{V}))}\big],} \\[3pt]
\FitMath{\|E^{(j)}\|_{j,\kappa}\le E_0,} \\[3pt]
\FitMath{\|R^{(j)}\|_{j,\kappa}\le g_j^{\kappa_0},} \\[3pt]
\FitMath{|B^{(j)}(X,\cdot)|\le\mathsf{B}\,e^{-\kappa d^\natural_j(\mathcal{S};X)},} \\[3pt]
\FitMath{\mathsf{M}(i,X)\le e^{-(2-c')\,p_0(g_i)},}
\end{gather*}

where $\mathsf{B}$ is a constant chosen before $\gamma$ (\cref{candidate:large_field}, \cref{sec:large-field-the-induction-closed-the-order-of-the-constants-the-sum}), and where the last inequality, the \emph{mass bound}, concerns those components $X$ that $\mathbf{R}$ has acted on at some level $i\le k$. Such a component was healed, or left standing as an island, or else it dropped out at a failed test, either of a stage or of the main step of $\mathbf{R}$. No assumption is made on a drop-out: its bound follows from the bound for the islands of the lower levels, the Wilson weight of the carrying term over $X$ being counted together with its operation. Let the constants be chosen in order, and let the torus be large compared with the $M\check{R}_j$-cubes at each level. If, in addition, $g_j\le\gamma$ and

\begin{gather*}
\FitMath{\lambda_\sharp(j-i)-c_5\le x_i-x_j\le B_\sharp(j-i)} \\[3pt]
\FitMath{(i\le j\le k+1),}
\end{gather*}

hold, then $\rho_0\in\mathfrak{F}_0$ and $\mathbf{R}\mathbf{T}\mathfrak{F}_k\subset\mathfrak{F}_{k+1}$ \emph{with the same constants}. No family of a level $j\le k$ is altered; only the background at which they are evaluated changes, $U_k\to U_{k+1}$. The new families come out of $\mathbf{T}$ with a surplus, and $\mathbf{R}$ uses it up: radii $(1+\beta_{\mathrm{s}})\alpha_{i,k}\ge\alpha_{i,k+1}$ and rate $(1+4\beta_{\mathrm{s}})\kappa$, where $\beta_{\mathrm{s}}$ is a fixed fraction. For given $(x_0,\dots,x_j)$, the activities $E^{(j)}(X,\cdot)$, $R^{(j)}(X,\cdot)$ do not depend on the label. When $X$ does not wrap around the torus, they depend neither on the torus nor on the cutoff. Further, $x_{k+1}=x_k-b_{k+1}(x_0,\dots,x_k)$ with $|b_{k+1}|\le B_\sharp$. When the label has no large fields, $|\mathbf{E}_k(U_k)|+|\mathbf{R}_k(U_k)|\le C|\mathbb{T}^{(k)}|$. For a general label, every layer is measured in its own scale, and for each $j\le k$ one adds the sites of level $j$ that lie in a rim around the large-field region $Z_j$ (\cref{ch:pointed-hessian-cancellation}, \cref{sec:flow-the-run-its-coefficient-and-the-analytic-requirements}). Finally, stability holds if the flow inequalities hold for all $i\le j\le K$:

\begin{gather*}
\FitMath{\chi_k\,e^{-x_kA(U_k)-E_-|\mathbb{T}^{(k)}|}\le\bar\rho_k,} \\[3pt]
\FitMath{\int\bar\rho_k\le e^{E_+|\mathbb{T}^{(k)}|},} \\[3pt]
\FitMath{E_+\le\tfrac32(N^2-1)\log x_k+O(1).}
\end{gather*}

\emph{The large-field factor} (part \cref{sec:large-field-rare-patches-gains-per-creation-and-costs-per-site}). The operation $\mathbf{T}_k(Z)$ contains delta functions, tests and conditional Gaussian weights. Apart from the core factors of the islands in $Z$, it contains no Wilson weight. Hence, for a component produced by failed plaquette tests, $\mathbf{T}_k(Z)1$ carries no small factor. If no $\mathbf{R}$ has acted on $Z$, then under the tests of the term it obeys only the bound $\le \exp[C\sum_{t\le k}\#_t(Z_t\cap Z)]$. The small factor lies in three places. They depend on pairwise disjoint sets of plaquettes, factors and bonds, so no part of the small factor is used twice. The first is the Wilson weight: near a failed plaquette test of step $j$, $x_jA\ge c\,p_0(g_j)^2$. The second is the kept gauge-fixing exponents; together with the Wilson sum, they pay for a failed approximate fluctuation test. The third is the Gaussian forms of the fluctuation fields, which pay for a failed fluctuation test.

(1) Fix a history, and a component $X$ of $Z_k$ that has not been acted on by any $\mathbf{R}$. Let $(j_i,\mathfrak{c}_i)$ be the creation cells of the history in $X$, namely those cells of the creation set of step $j_i$, which are $M\check{R}_{j_i}$-cubes, that meet a failed test \cref{sec:large-field-rare-patches-gains-per-creation-and-costs-per-site}. Then, for all values of the exterior variables, and with the norm taken as the total variation of that part which the term of the history carries,

\begin{gather*}
\FitMath{\big\|\mathbf{T}_k(X)\big[\chi_{k,\mathcal{S}}\,e^{-A(x_k(\cdot),U_k)}\big]\big\|\le\prod_ie^{-c\,p_0(g_{j_i})^2}\,\exp\Big[C\sum_{t\le k}\#_t(Z_t\cap X)\Big]:}
\end{gather*}

there is one factor for each creation, taken at the coupling of its own step, and one constant for each site.

(2) Bounds of this kind do not multiply over the components, because the background couples the components to one another.

(3) What does multiply, piece by piece, is the pointwise majorant on which (1) rests. A piece $(t,D)$ consists of a component $D$ of the large-field region formed at step $t$, whose enlargement $D^\sim$ shares no cube with any component that $\mathbf{R}$ has acted on at a level in $[t,k]$, together with the stored variables $\xi_{t,D}$ belonging to it and the weight $\mathsf{w}_{t,D}$ of those variables. Throughout the support of the history,

\begin{gather*}
\FitMath{A(x_k(\cdot),U_k)\ge A_{\mathrm{cr}},} \\[3pt]
\FitMath{e^{-A_{\mathrm{cr}}}\prod_{(t,D)}\mathsf{w}_{t,D}\le\prod_{\mathrm{creations}}e^{-c\,p_0(g_t)^2}\prod_{(t,D)}\bar{\mathsf{w}}_{t,D}(\xi_{t,D}),} \\[3pt]
\FitMath{\int\bar{\mathsf{w}}_{t,D}\le e^{C\#_t(D)},}
\end{gather*}

where $A_{\mathrm{cr}}$ denotes the Wilson sums near the creations in these pieces. Only those creations are counted on the right-hand side. Each $\bar{\mathsf{w}}_{t,D}$ is a function of the variables of its own piece only.

(4) After integration, take a whole label together with the list of those outermost components $(i,X)$ that $\mathbf{R}$ has acted on and that still stand in the label. For some history having this label and this list, the bound carries one factor for each of its creation cells lying outside the components that $\mathbf{R}$ has acted on and that still stand in the label (islands and drop-outs; the creations inside such a component are spent in its mass), and one mass for each outermost one among these components:

\begin{gather*}
\FitMath{\int\mathbf{T}'_k(\mathcal{S})\big[\chi_{k,\mathcal{S}}\,e^{-A(x_k(\cdot),U_k)}\big]dV_k\le\exp\Big[C\sum_j(1+\log x_j)\#_j(Z_j)\Big]\prod_{\text{creation cells }(t,\mathfrak{c})}e^{-c\,p_0(g_t)}\prod_{(i,X)}\mathsf{M}(i,X),}
\end{gather*}


\begin{gather*}
\FitMath{\mathsf{M}(i,X)=\frac1{|W_4|}\sum_\omega e^{c\,\mathsf{v}(\omega)}\int\sup_{\mathbf{p}}\big|e^{\Phi_\omega(\xi;\mathbf{p})}\big|\,\mu_\omega(d\xi),}
\end{gather*}

where $\omega$ ranges over the local histories of $X$, each equipped with a hypercubic frame in $W_4$ of the gauge tree, and $\xi$ denotes the stored variables. The variable $\mathbf{p}=(\mathbf{U},\mathbf{J})$ ranges over the complex pairs on $X^\sim$ at which a term evaluates its factor $e^{\Phi_\omega}$. The exponent has the form $\Phi_\omega=\Phi^\circ_\omega(\xi)+\sum_f\Phi_{\omega,f}(\xi;\mathbf{p}|_{R_f})$ with $R_f\subset X^\sim$, and $\mathsf{v}(\omega)$ is a space-time volume. The supremum is taken inside the integral and the sum. This is the form of the mass for a healed component or an island. In the case of a drop-out, the Wilson weight of the carrying term over $X$ appears inside the integral, and the supremum extends over the evaluations at $\xi$. These are the real configurations that are compatible with the stored variables and regular layer by layer, and at which every test listed by $\omega$ takes its recorded value, with the exception of the test whose failure made $X$ a drop-out. They form a set fibred over $\xi$ \cref{sec:large-field-rare-patches-gains-per-creation-and-costs-per-site}. The induction carries the mass bound at the coupling of level $i$ of the action. In this bound the $c$ of $2-c$ is small, whatever the $c$ in the weight $e^{c\,\mathsf{v}(\omega)}$. The mass evaluated at a later level does not exceed the mass at level $i$.

(5) The integral operation $\mathbf{T}'_k(X)$ on a component on which $\mathbf{R}$ acts contains, except in the case of a drop-out, a Wilson factor of its own, namely that of the core. When the component is healed or remains as an island, the total mass that is bounded is that of this operation, not that of $\mathbf{T}_k(Z)$:

\begin{gather*}
\FitMath{\mathbf{T}'_k(X)\,1\le\sup\,\exp\Big\{\sum_{j\le k}CM^4\check{R}_j^5\,d'_j(Z_j)\Big\}\prod_{j\le k}\prod_i\exp\Big(-c\,p_0(g_j)^2\big(d'_j(Z^{(i)}_j)+1\big)-2p_0(g_j)\Big),}
\end{gather*}

where the supremum extends over every admissible choice of the domains $Z_j$, $Z^{(i)}_j$ (Ba{\l}aban, Comm. Math. Phys. 122 (1989) 355, (1.79)), and the sums over the admissible sequences are counted among the costs. The inner product is taken over the regions $Z^{(i)}_j$ created at step $j$. For a component $Z$ of $Z_j$ which has been neither healed nor marked, the factor is written $\exp(-\kappa_j(Z)-2p_0(g_{j(Z)}))$; here $j(Z)$ is the level at which the first creation in $Z$ occurs. The \emph{credit} satisfies $\kappa_j(Z)\ge\sum_{n=j+1}^{j+\mathsf{k}}CM^4\check{R}_n^5\,d'_n(S^{n-j}Z)$, where the sum runs over the $\mathsf{k}$ further steps at which nothing is created. Here $S$ denotes the hull in the next partition with a margin added. The number of $M\check{R}_n$-cubes that meet the region takes the place of the linear size $d'_n$, because a single cube has $d'=0$ and still carries a cost. The credit covers the volume costs of every step until $\mathbf{R}$ acts on $Z$. Once no further step is owed, $\mathbf{T}'_k(X)1\le e^{-(2-c)\,p_0(g_k)}$. Moreover, in the weighted sum over the labels, every component $C$ of $Z_k$ keeps a factor $\exp[-c\,p_0(g_k)-\kappa_1|C|_k]$, where $|C|_k$ is its number of cells.
\end{proposition}

\begin{proposition}[{the small-field step as a map on actions}]\label{prop:averaging-1a-2-8} Suppose that the tests and $\mathfrak{G}_k$ are formed at $\mathsf{t}$, where $\mathsf{t}$ is large, and let $x\in\mathbb{C}$ with $|x-\mathsf{t}|\le c\,\mathsf{t}$. Suppose that the analyticity radii $(\alpha_0,\alpha_1)$ of the levels $j\le k+1$ are small, with $CM\alpha_0<\alpha_1$ and $\alpha_1^2\le C\alpha_0$ at each level. Suppose also that they are comparable: the quotient of the radii of level $k+1$ by those of any level $j\le k$ lies between $(1+\beta_{\mathrm{s}})^{-1}$ and $(1+\beta_{\mathrm{s}})(k+2-j)^{1/2}$. These are the admissible systems of \cref{ch:local-formation-program}. Let $S$ be local, analytic and \emph{renormalised} in the same way as $\mathbf{E}_k+\mathbf{R}_k$:

\begin{gather*}
\FitMath{S=\sum_{j\le k}\Big[\sum_z\sum_{X\ni z}\big(G^{(j)}(X,\cdot,z)-G^{(j)}(X,\mathbf{1},z)\big)-e_jA\Big]+\sum_{j\in\mathcal{J}}\sum_X\big(H^{(j)}(X,\cdot)-H^{(j)}(X,\mathbf{1})\big),}
\end{gather*}


\begin{gather*}
\FitMath{N(S)=\sum_{j\le k}\omega_{k+1-j}\|G^{(j)}\|_{j,\kappa}+\sum_{j\in\mathcal{J}}\mathsf{q}_j^2\|H^{(j)}\|_{j,\kappa},} \\[3pt]
\FitMath{\omega_n\le Cn^CL^{-\alpha_{\mathrm{H}}n},}
\end{gather*}

where the pointed totals of $G^{(j)}$ are covariant under all isometries of $\mathbb{T}^{(j)}$, and $e_j$ is the coefficient \cref{ch:pointed-hessian-cancellation} of $G^{(j)}$. The totals of $H^{(j)}$ are invariant under the isometries of $\mathbb{T}^{(k+1)}$. Finally, $\mathsf{q}_j$ is the ratio of the analyticity radii of level $k+1$ to those of level $j$. Geometric weights occur only in the first sum. Both symmetries are part of the definition of ``renormalised''. The covariance of the pointed totals is what makes the bracket left after the subtraction of $e_jA$ irrelevant \cref{ch:pointed-hessian-cancellation}. Consider $W$ in a class of real block fields, namely those at which every background formed in the step exists and yields pairs in the domains of the terms. Then, on the enlarged radii and uniformly in $k$, 

\begin{gather*}
\FitMath{\mathfrak{E}_{k+1}[\mathsf{t};x;S]=\log Z^{(k)}+\sum_{X\in\mathcal{D}_{k+1}}\mathsf{F}(X,U_{k+1})+\mathrm{const}}
\end{gather*}

, where

\begin{gather*}
\FitMath{|\mathsf{F}(X,\cdot)|\le C\big(|1-x/\mathsf{t}|+\delta(\mathsf{t})+\mathsf{s}\,N(S)+e^{-c\,\mathsf{t}\delta(\mathsf{t})^2}\big)\,e^{-(1+4\beta_{\mathrm{s}})\kappa d_{k+1}(X)},}
\end{gather*}

with holomorphic dependence on $x$ and on the parameters of $S$. The map $S\mapsto\mathsf{F}$ satisfies a Lipschitz bound with constant $C\mathsf{s}$; here $\mathsf{s}\ll1$ denotes the quotient of the fluctuation threshold by the analyticity radius. Everything above holds under the condition $\mathsf{s}\,(1+N(S))\le c$. Along the run, $G^{(j)}=E^{(j)}$ and $e_j=b_j$, while the $H^{(j)}$ are the aged $R^{(j)}$ of \cref{prop:averaging-1a-2-9}.
\end{proposition}

\begin{proposition}[{anatomy of step $k+1$}]\label{prop:averaging-1a-2-9} The insertions are made in the following order. First, the plaquettes of the local backgrounds of $V_{k+1}$ are tested. Next, on the blocks $P''$ that lie away from failures and from $Z_k$, where the $V(y,x)$ are defined, the factor $\mathfrak{G}_k(\cdot;P'')$ is inserted. Then the approximate fluctuation $V_k(V^{\mathrm{bg}}_{k,\square})^{-1}$ is tested, and this determines $\Omega_{k+1}$. On that region $V^{\mathrm{bg}}_k$ is defined and $V_k=e^{i\mathcal{A}_k}V^{\mathrm{bg}}_k$ is set. Finally, $\mathcal{A}_k$ is tested, and this determines $\Lambda_{k+1}$. The insertions and tests just described, together with the regions that they determine, are those of Ba{\l}aban's form, whereas the sharper tests on collar and rim, together with their outcomes, are given in \cref{sec:large-field-the-tests-the-nested-regions-and-the-content-of-a}. As a formula,

\begin{gather*}
\FitMath{(\mathbf{T}\rho_k)(V_{k+1})=\sum_{P,Q,R}\int dV_k\,\delta(\mathcal{M}(V_k)V_{k+1}^{-1})\;\chi_{k+1}\chi^c_{k+1}(P)\;\mathfrak{G}_k(V_k)\;\chi'_k\chi'^c_k(Q)\;\chi^{(k)}\chi^{(k)c}(R)\;\rho_k(V_k).}
\end{gather*}

After \cref{lem:averaging-1a-2-4} has been applied and the field has been scaled, $\mathcal{A}_k=g_kB$, a Taylor expansion about the critical point yields

\begin{gather*}
\FitMath{x_k\big[A(U_k(V_k))+\mathrm{GF}_k(V_k)\big]=x_kA(U_{k+1})+\tfrac12\langle B,\Delta^{(k)}B\rangle+\sum_bO(g_k|B|^3).}
\end{gather*}

Completing the square in $B|_{\Lambda_{k+1}}$ produces $C^{(k)}(\Lambda_{k+1})$, together with the conditional Gaussian that is kept in $\mathbf{T}^{(k)}(Z_{k+1})$:

\begin{gather*}
\FitMath{(\mathbf{T}\rho_k)(V_{k+1})=\sum_{\mathcal{S}'}\chi_{k+1}(\Omega_{k+1})\,\mathbf{T}^{(k)}(Z_{k+1})\mathbf{T}'_k(\mathcal{S})\Big[e^{A_{k,\mathcal{S}}(U_{k+1})+E^{(k)}_0(\Omega_{k+1})}\,\mathsf{I}\Big],} \\[3pt]
\FitMath{\mathsf{I}=Z^{(k)}(\Lambda_{k+1})\!\int d\mu_{C^{(k)}(\Lambda_{k+1})}(B)\,\chi^{(k)}(\Lambda_{k+1})\,e^{\mathbb{P}^{(k)}+\{A_{k,\mathcal{S}}\}'+\mathsf{W}^{(k)}},}
\end{gather*}

where $E^{(k)}_0(\Omega)$ consists of the first two terms of $\mathsf{E}^{(k)}(\Omega,\Lambda)$. The measure $\mu_{C^{(k)}(\Lambda_{k+1})}$ is the Gaussian measure conditioned on the exterior fluctuation field, and its mean is linear in that field. The prime indicates that the quadratic part of $-x_kA$ has been absorbed into the measure. The term $\mathsf{W}^{(k)}$ is left over from a change of variables on the strip, that is, on the part of $S_{k+1}$ lying along $\partial\Lambda_{k+1}$. Before any expansion, the logarithm is divided into three parts by interpolation:

\begin{gather*}
\FitMath{\log\mathsf{I}=\int_0^1\!dt_1\,\big\langle\{\mathbf{B}_k-A(x_k(\cdot)-x_k,\cdot)\}+\mathsf{W}^{(k)}\big\rangle_{t_1,1}+\int_0^1\!dt_2\,\big\langle\{\mathbf{R}''_k\}\big\rangle_{0,t_2}+\log\mathsf{I}\big|_{t_1=t_2=0},}
\end{gather*}

where $\mathbf{R}''_k$ contains the young $R^{(j)}$, namely those whose cubes are larger than a block. The logarithm is expanded \cref{ch:local-formation-program} and regrouped (\cref{lem:averaging-1a-2-4}, with $b_{k+1}$, $\varphi_{k+1}$ and $\mathsf{E}^{(k)}(\Omega_{k+1},\Lambda_{k+1})$). This gives the form of level $k+1$:

\begin{gather*}
\FitMath{(\mathbf{T}\rho_k)(V_{k+1})=\sum_{\mathcal{S}'}\chi_{k+1}\,\mathbf{T}'_{k+1}(\mathcal{S}')\exp\Big[A_{k,\mathcal{S}}(U_{k+1})+\!\!\sum_{X\subset\Lambda_{k+1}}\!\!\big[E^{(k+1)}+R^{(k+1)}\big](X,U_{k+1})} \\[3pt]
\FitMath{+\!\!\sum_{X\cap Z^\sim_{k+1}\ne\emptyset}\!\!B^{(k+1)}(X;U_{k+1},\underline{\mathfrak{f}})+E^{(k)}_0(\Omega_{k+1})+\nu_{\mathfrak{g}}|\Lambda_{k+1}^{(k)*}|\Big]}
\end{gather*}


\begin{gather*}
\FitMath{=\sum_{\mathcal{S}'}\chi_{k+1}\,\mathbf{T}'_{k+1}(\mathcal{S}')\exp A_{k+1,\mathcal{S}'}(U_{k+1});}
\end{gather*}

after this, $\mathbf{R}$ acts. The terms with $X\subset\Lambda_{k+1}$ coincide with the terms of the pure small-field integral (among the $R^{(k+1)}$, here and in $\mathsf{E}^{(k)}$, a label keeps those whose $X$ lies one cell inside $\Lambda_{k+1}$, and the others are counted among the $B^{(k+1)}$); this coincidence is the locality stated in \cref{prop:averaging-1a-2-7}.
\end{proposition}

\begin{proposition}[{a model case: the free field; the renormalisation group as nested orthogonal projections}]\label{prop:averaging-1a-2-10} Let the \emph{free curvature} $F$ be the standard Gaussian vector in the Euclidean space $\Omega^2_{\mathrm{ex}}=\operatorname{ran}\mathrm{d}$ of exact $2$-cochains on the bare lattice. Here $\mathrm{d}$ is the coboundary taking real $1$-cochains to real $2$-cochains on the bare lattice, and $\mathrm{d}_{(j)}$ is the coboundary of $\mathbb{T}^{(j)}$. This vector stands for one colour of the gauge field, taken in the quadratic approximation and up to closed forms. The plain linear average and the straight linear average induce the same average of curvatures, called the \emph{plaquette average}:

\begin{gather*}
\FitMath{\mathrm{d}_{(j)}(\mathsf{Q}^{\mathrm{pl}})^j=\mathrm{d}_{(j)}(\mathsf{Q}^{\mathrm{s}})^j=\mathsf{A}^{(j)}\mathrm{d},} \\[3pt]
\FitMath{(\mathsf{A}^{(j)}F)(P)=L^{-4j}\sum_{z\in B^j(y)}\,\sum_{p\in\mathrm{Sq}^{(j)}(z;P)}F(p),}
\end{gather*}

where $\mathrm{Sq}^{(j)}(z;P)$ is the set of $L^{2j}$ bare plaquettes which together cover the plaquette $P$ of $\mathbb{T}^{(j)}$ with corner $y$, once it is translated to $z$. Put $\mathsf{W}_j:=\Omega^2_{\mathrm{ex}}\ominus(\Omega^2_{\mathrm{ex}}\cap\ker\mathsf{A}^{(j)})$. These are subspaces of $\Omega^2_{\mathrm{ex}}$ whose level is written as a subscript, and they are distinct from the term $\mathsf{W}^{(k)}$ of \cref{prop:averaging-1a-2-9}, whose level is written as a superscript. Let $P_{\mathsf{W}}$ denote the orthogonal projections. Then the spaces are nested, $\mathsf{W}_0\supset\mathsf{W}_1\supset\dots\supset\mathsf{W}_K$. Given the block field of level $k$, the conditional expectation of $F$ is $P_{\mathsf{W}_k}F$. This is the curvature of the linear background, namely the exact cochain of smallest norm that has the prescribed plaquette averages. The covariances of the background and of the fluctuation are two complementary orthogonal projections, $\mathbb{E}\big[P_{\mathsf{W}_k}F(p)\,P_{\mathsf{W}_k}F(q)\big]=P_{\mathsf{W}_k}(p,q)$ and $P_{\mathrm{ex}}=P_{\mathsf{W}_k}+P_{\mathsf{W}_k^\perp}$, where $P_{\mathrm{ex}}$ is the projection onto $\Omega^2_{\mathrm{ex}}$. Finally,

\begin{gather*}
\FitMath{F=P_{\mathsf{W}_K}F+\sum_{j=1}^K\mathcal{C}_{(j)}F,} \\[3pt]
\FitMath{\mathcal{C}_{(j)}=P_{\mathsf{W}_{j-1}}-P_{\mathsf{W}_j}=\mathcal{H}_{j-1}C^{\mathrm{tr}}_j\mathcal{H}_{j-1}^*,} \\[3pt]
\FitMath{|\mathcal{C}_{(j)}(p,q)|\le C\eta_{j-1}^4e^{-\delta|x(p)-x(q)|/a_{K-j+1}},}
\end{gather*}

where the summands are independent and the covariance of each summand is its own projection. Here, for a block field $B$ of level $j-1$, $\mathcal{H}_{j-1}B$ denotes the curvature of its linear background, namely the exact $2$-cochain on the bare lattice of smallest norm satisfying $\mathsf{A}^{(j-1)}\mathcal{H}_{j-1}B=\mathrm{d}_{(j-1)}B$; and $C^{\mathrm{tr}}_j$ is the fluctuation covariance, in a tree gauge, of step $j$. Thus step $j$ integrates out precisely $\mathcal{C}_{(j)}F$, and the amplitude and range of each slice are those of its own scale. In this model the semigroup law is the nesting of the spaces, and the delta constraint is conditioning. The nesting of backgrounds takes the form $P_{\mathsf{W}_j}P_{\mathsf{W}_k}=P_{\mathsf{W}_k}$ ($j\le k$). Three features of the full theory are absent from the model: the fibre is flat, no coupling runs, and no field is large. The current of a linear background is not zero: $\mathrm{d}^*P_{\mathsf{W}_k}F=((\mathsf{Q}^{\mathrm{s}})^k)^*\mathrm{d}_{(k)}^*\phi$, which is the image of a Lagrange multiplier. What is missing is the term $\langle\mathsf{H}_kC^{(2)}_k,J\rangle$ of \cref{prop:averaging-1a-1-10}, because $C^{(2)}_k=0$ on a flat fibre. In the non-linear theory the current and the multiplier do vanish, at $U=\mathbf{1}$. For this reason the linearisation at that point of the minimiser of $A$ subject to $\mathcal{M}^kU=V$ is the minimiser of $\frac12\|\mathrm{d}\mathsf{a}\|^2$ subject to the linearised constraint, and the second-order part of the averages makes no contribution. Hence the model is the linearisation at $\mathbf{1}$. It serves as the template for \cref{sec:short-distance-the-short-distance-two-point-law-and-its-asymptotic-freedom}, \cref{sec:short-distance-isolating-a-single-scale-coupling-and-normalisation-at-the-scale}.
\end{proposition}

\begin{lemma}[{axial gauges}]\label{lem:averaging-1a-3-1} Consider a box of side $D$ on which $|\partial V-\mathbf{1}|\le\epsilon$. Among the representatives $V^{\psi'}$ with $\psi'(y_P)=\mathbf{1}$, the configuration $V^{\mathrm{ax}}=V^\psi$, $\psi(w)=V(\Gamma_{y_P,w})$, is the only one that equals $\mathbf{1}$ on the tree. The non-abelian Stokes theorem expresses $V^{\mathrm{ax}}(b)$ as a product of conjugated plaquette variables, at most $|b_--y_P|_1$ in number:

\begin{gather*}
\FitMath{V^{\mathrm{ax}}(b)=\prod_{i=1}^m\psi(w_i)\,\partial V(p_i)^{\pm1}\psi(w_i)^{-1},} \\[3pt]
\FitMath{|V^{\mathrm{ax}}(b)-\mathbf{1}|\le|b_--y_P|_1\,\epsilon\le4D\epsilon;}
\end{gather*}

Suppose that $D\epsilon\le c$ and that $F=\frac1i\log\partial V$ is H{\"o}lder along combs; it is distinct from the field strength $F_V$ introduced in the Notation, \ensuremath{\S}2, being the logarithm of each plaquette variable. Then $\frac1i\log V^{\mathrm{ax}}_\mu(x)=\sum_{\rho<\mu}(x_\rho-y_\rho)F_{\rho\mu}(y_P)+O(|x-y_P|^2)\times(\text{H{\"o}lder constant}+\epsilon^2)$, which is the expansion used in \cref{ch:pointed-hessian-cancellation} and \cref{sec:large-field-stored-variables-and-integration-in-a-large-field-patch}. On a separated sequence suppose that $U_0,U\in\mathfrak{U}_k(\Omega_\bullet,\alpha_0)$, with $\alpha_0>0$ below the fixed geometric ceiling, and that their raw averaged data differ by at most $\alpha_1\ge0$ on every interior and outward determining row. Let $\hat g$ be the real gauge, equal to one at every layer top, whose descending construction aligns the averaged contour variables of $U^{\hat g}$ with those of $U_0$. The construction following \eqref{mb:rigidity-sum} gives $|\mathcal M^j(U^{\hat g})(b)\mathcal M^jU_0(b)^{-1}-\mathbf1|\le C_L(\alpha_0+\alpha_1)$ for $0\le j\le k$ and every carried level-$j$ bond $b$ inner to $\Omega_0$ with an endpoint in $\Omega_j$, uniformly in $k$. This is a raw group-element bound, before division by the local spacing. On a single level (compare the rigidity \cref{lem:averaging-1a-1-9}), $\mathcal{A}_k$ is defined, and is $O(\varepsilon_k)$, only after the gauge has been fixed.
\end{lemma}

\begin{proposition}[{Faddeev--Popov identities with Jacobian one}]\label{prop:averaging-1a-3-3} Let a compact group $H$ act, and let $\mathsf{k}\ge0$ satisfy $\int_H\mathsf{k}(h\cdot x)\,dh=1$. Then for every finite complex measure $\nu$ and every bounded invariant $F$,

\begin{gather*}
\FitMath{\int\nu(dx)\int_Hdh\,\mathsf{k}(h\cdot x)\,F(h\cdot x)=\int F\,d\nu;}
\end{gather*}

The measure obtained in this way is constructed as follows: one chooses $h$ with law $\mathsf{k}(h\cdot x)dh$ and replaces $x$ by $h\cdot x$. This measure equals $\mathsf{k}\,\bar\nu$, where $\bar\nu$ is the gauge average of $\nu$. When $\nu$ is invariant, it is the product $\mathsf{k}\,\nu$. If $u=\mathbf{1}$ on the coarse sites, then $V^u(y,x)=V(y,x)u(x)^{-1}$. Where $V(y,x)$ is undefined, no factor is inserted at the pair $(y,x)$, and with this convention

\begin{gather*}
\FitMath{\int\prod_xdu(x)\,\mathfrak{G}_k(V^u;P'')=1,} \\[3pt]
\FitMath{(\mathbf{T}\rho_k)(V_{k+1})=\int dV\,\delta(\mathcal{M}(V)V_{k+1}^{-1})\,\mathfrak{G}_k(V;P'')\,\rho_k(V).}
\end{gather*}

In addition, $\mathfrak{G}_k(V^v)=\mathfrak{G}_k(V)$ whenever $v$ is constant on blocks. The same holds for delta functions on spanning trees. Let $T$ be a tree whose root lies outside the sites of the group, and let $V_0$ be any configuration on $T$. Then a unique $u$ satisfies $(u\cdot V)|_T=V_0$, and the measure $dV$ splits into factors along the tree:

\begin{gather*}
\FitMath{\int f\,dV=\int du\int\prod_{b\notin T}dS(b)\,f\big(u\cdot(S\sqcup V_0)\big),} \\[3pt]
\FitMath{\text{so}} \\[3pt]
\FitMath{\int dV\,\delta_T(VV_0^{-1})\,\rho F=\int dV\,\rho F}
\end{gather*}

for invariant $\rho,F$. Finally, the delta function can be removed: the gauge average of $\Psi\,\delta_T(VV_0^{-1})dV$ is $(\Psi\circ\pi_T)\,dV$, where $\pi_T$ is the projection onto the slice along the orbits.
\end{proposition}

\begin{lemma}[{[stated technical lemma; Group A] (the Landau gauge)}]\label{lem:averaging-1a-3-4} Let $U$ be tower axial and let $W=U^{u^{-1}}$. Then $\mathsf{n}_W=u^{-1}\bar{R}_{U_0}u$, so a $u$ with trivial averages leaves the straightened averages exactly unchanged. In linearised form this reads $Q^{\sharp}_j(U_0)D\lambda=D_{\mathcal{M}^jU_0}Q'_j(U_0)\lambda$. Hence $Q^{\sharp}A$ does not change under $A\to A-D\lambda$, $\lambda\in N(Q')$. Up to terms that contain $J(U_0)$, gradients are null directions of $\Delta^{\mathrm{W}}$. Indeed $D^*J(U_0)=0$, and, writing $\lambda_+$ for $\lambda(b_+)$ transported to $b_-$,

\begin{gather*}
\FitMath{\langle A-D\lambda,\Delta^{\mathrm{W}}(A-D\lambda)\rangle=\langle A,\Delta^{\mathrm{W}}A\rangle-\big\langle i[\lambda,A]-i[A,\lambda_+]-i[\lambda,D\lambda],\,J(U_0)\big\rangle.}
\end{gather*}

The projection is given by

\begin{gather*}
\FitMath{R=1-G'Q'^*(Q'G'^2Q'^*)^{-1}Q'G',} \\[3pt]
\FitMath{\operatorname{ran}R=\{f:Q'G'f=0\},} \\[3pt]
\FitMath{\inf_{\lambda\in N(Q')}\|D^*A-\Delta\lambda\|^2=\|(1-R)D^*A\|^2,}
\end{gather*}

and $R\Delta\lambda=\Delta\lambda=G'^{-1}\lambda$ for $\lambda\in N(Q')$. The map $\lambda\mapsto(R\Delta\lambda,Q'\lambda)$ is a bijection onto $\operatorname{ran}R\times\{\text{functions on tops}\}$. Its inverse is $(f,\nu)\mapsto G'f+H'\nu$, where $\nu$ is a function on the tops and $H'$ is a right inverse of $Q'$, constructed from $G'$, with $R\Delta H'=0$. Let $U\in\mathfrak{U}_k(\Omega_\bullet,\alpha_0)$ and let $M_1$ be large. Assume that $M\alpha_0$ (and therefore $M_1\alpha_0$) lies below a number that depends only on $L$, and that $\mathbb{T}^{(k)}$ has at least $CM_1$ sites in each coordinate direction. Then the form

\begin{gather*}
\FitMath{\langle A,\Delta^{\mathrm{W}}_UA\rangle+\|RD^*_UA\|^2+\langle Q^{\sharp}A,aQ^{\sharp}A\rangle}
\end{gather*}

is positive definite. Its inverse $\mathsf{G}_k(U)$ maps $|\cdot|_{(-3)}$ into $\|\cdot\|_1$, with a norm that depends only on $L$. On the set $\{RD^*A=0,\ \xi Q^{\sharp}A=B\}$ the form differs from the Hessian by a constant. Consequently $\mathsf{H}_kB=\mathsf{G}Q^{\sharp*}(Q^{\sharp}\mathsf{G}Q^{\sharp*})^{-1}\xi^{-1}B$ minimises the Hessian itself under the constraint, and $RD^*\mathsf{G}Q^{\sharp*}=0$. Here $\mathsf{G}$ stands for Ba{\l}aban's $G$, and not for the inverse of the displayed form, which is his $G_0$. The formula for $\mathsf{H}_kB$ and the identity just stated hold exactly only when the Hessian in the form is replaced by its invariant extension $\Delta_\pi$. This extension is defined by $\langle A,\Delta_\pi A\rangle=\langle A-D\lambda,\Delta^{\mathrm{W}}(A-D\lambda)\rangle$, where $\lambda\in N(Q')$ is the one with $RD^*(A-D\lambda)=0$. The replacement is a perturbation of size $CM\alpha_0$, which is small together with the current because $\alpha_0$ is chosen after $M$. The bounds persist under this perturbation, and it vanishes where $J=0$. In the symbols of the Notation, \ensuremath{\S}5, this operator equals $\mathsf{H}\,\xi^{-1}$. There $\mathsf{H}$ denotes the linear minimiser constructed from $G$, and it satisfies $Q^{\sharp}\mathsf{H}=1$. Where $J=0$ and $\xi=1$, the operator coincides with the $\mathsf{H}_k$ of the Notation. It carries the normalisation for the constraint $\xi Q^{\sharp}A=B$. With his $Q^{\mathrm{B}}$, Ba{\l}aban uses both normalisations: the minimiser under $QA=B$ and the minimiser under $L^j\eta\,Q_jA=B$. On layer $j$ the two differ by the factor $\xi^{-1}=(L^j\eta_k)^{-1}$. The normalisation with $\xi$ is the one taken in his paper on the variational problem.
\end{lemma}

\begin{proposition}[{from the axial to the Landau gauge}]\label{prop:averaging-1a-3-6} Work in the scheme $\mathsf{c}$, in which Ba{\l}aban's average of 1985 replaces $\mathcal{M}$ in every object, and fix the block and gap numbers $\mathfrak M,R_1$. The constants $c_1,B_1>0$ may depend on $L,\mathfrak M,R_1$, but are independent of depth and compatible volume. Assume the hypotheses of \cref{lem:averaging-1a-3-1} [stated technical lemma; Group A], with the full gauge $\operatorname{Ax}(\Omega_\bullet,U_0)$ and not $\operatorname{Ax}^\circ$, and assume $\alpha_0+\alpha_1\le c_1$. Then $U$ has exactly one Landau pair that satisfies

\begin{gather*}
\FitMath{|A|_{(-1)},\,|\nabla_{U_0}A|_{(-2)},\,|D^*DA|_{(-3)},\,|\Delta_{U_0}A|_{(-3)}\le B_1(\alpha_0+\alpha_1),} \\[3pt]
\FitMath{\sup|u-\mathbf{1}|\le CB_1(\alpha_0+\alpha_1),}
\end{gather*}

that satisfies H{\"o}lder bounds on $\nabla A$ (these bounds give rise to the $\alpha_{\mathrm{H}}$ below), and that has the straightened averages of $U$,

\begin{gather*}
\FitMath{\mathcal{Q}^{\sharp}_j(U_0;A)=B\,\text{ on the bonds at }\mathbb{Y}_j,} \\[3pt]
\FitMath{B=\tfrac1i\log\big[\mathcal{M}^jU(\mathcal{M}^jU_0)^{-1}\big],} \\[3pt]
\FitMath{|B|\le C\alpha_1}
\end{gather*}

(with an explicit correction entering on bonds that join two layers). For $\mathcal{M}$ itself, under the same hypotheses with $\alpha_0+\alpha_1\le c$, there exists a pair $(A,u)$ in the slice of \cref{ch:multilevel-backgrounds} that has these averages and satisfies $\|A\|_1,\ |D^*DA|_{(-3)}\le C(\alpha_0+\alpha_1)$. That slice has the same form as the Landau gauge but a different projection, whose site averages run along combs from the block corners.
\end{proposition}

\begin{proposition}[{a priori bounds on the slice}]\label{prop:averaging-1a-3-7} Assume that $U_0,e^{i\eta_kA}U_0\in\mathfrak{U}_k(\Omega_\bullet,\alpha_0)$, that $RD^*A=0$, and that the averages are as above. Assume further that the only bound given is $|A|_{(-1)}<\alpha_2$, with $\alpha_0,\alpha_1,\alpha_2\le c$ ($c$ chosen for the fixed block and gap geometry of \cref{prop:averaging-1a-3-6}) and $C\alpha_2^2\le\alpha_0+\alpha_1$. Then, in the scheme $\mathsf{c}$, all the bounds of the theorem on $A$ hold, with $B_1=CL\times$ the norms of $\mathsf{G}_k(U_0)$, $\mathsf{H}_k(U_0)$. For $\mathcal{M}$, the bounds on $\|A\|_1$ and $|D^*DA|_{(-3)}$ hold, in the slice of \cref{ch:multilevel-backgrounds}.
\end{proposition}

\begin{proposition}[{regular means small potential, locally}]\label{prop:averaging-1a-3-8} Let $U\in\mathfrak{U}_k(\Omega_\bullet,\alpha_0)$ and impose the fixed background geometry floors. Take $M$ to be a positive integer multiple of $\mathfrak r_{\mathrm p}$ and take a cube of side $M\xi_j$ aligned on the $\mathfrak r_{\mathrm p}\xi_j$ grid whose concentric enlargement of side $(M+4\rho)\xi_j$ lies in $\Omega_j$ and is shorter than half the torus. The curvature and full-current bounds are required on every complete stencil meeting this enlargement, as in \eqref{mb:local-regularity-conclusion}. Assume $CM\alpha_0\le c_1$. Then on the original cube, for some $u$

\begin{gather*}
\FitMath{U^{u^{-1}}=e^{i\eta_kA},} \\[3pt]
\FitMath{\xi_j|A|,\,\xi_j^2|\nabla A|,\,\xi_j^3|D_{\mathbf{1}}^*D_{\mathbf{1}}A|,\,\xi_j^3|\Delta A|\le CM\alpha_0,}
\end{gather*}

For the additional normalization in scheme $\mathsf{c}$, choose the gauge by \cref{prop:averaging-1a-3-6}. First put the top box datum in its box axial gauge, then impose the lower full tower gauge while fixing the top nodes, and finally take the normalized Landau pair relative to $\mathbf{1}$. The lower tower changes preserve the top datum. The Landau change has trivial nonlinear site averages, so the reduced-average identity of \cref{lem:averaging-1a-3-4} preserves its straightened value. On the homogeneous top rows this gives exactly $\mathcal{Q}^{\sharp}_j(\mathbf{1};A)=\frac1i\log(\mathcal{M}^jU)^{\mathrm{ax}}$. This identity uses the stated axial-to-Landau theorem and its choice of gauge. Conversely, let $U_0\in\mathfrak{U}_k(\Omega_\bullet,\alpha_0)$ and $\|A\|_1,|D^*DA|_{(-3)}<\alpha_2$, with $\alpha_0,\alpha_2$ small. Then $e^{i\eta_kA}U_0\in\mathfrak{U}_k(\Omega_\bullet,\alpha_0+3\alpha_2)$. Now take a complex $\mathbf{U}=e^{i\eta_kA'}U$ with $\xi_j|A'|<\alpha_1$ whose own plaquettes are small, $|\partial\mathbf{U}-\mathbf{1}|<\alpha_0\eta_k^2\xi_j^{-2}$ on the cube. (From $\|A'\|_1<\alpha_1$ this follows only with $\alpha_0+C\alpha_1$ in place of $\alpha_0$, \cref{prop:averaging-1a-1-6}.) Assume also $CM(\alpha_0+\alpha_1)\le c$. Then the averages are $V^v$. Here $V$ is what the averages of $\mathbf{U}$ become in the comb gauge of the cube, and $v=v'v_{\mathrm{r}}$ is a $G^c$-valued function on the nodes. The factor $v_{\mathrm{r}}$ is $G$-valued and does not depend on $A'$, the factor $v'$ is holomorphic in $A'$, and

\begin{gather*}
\FitMath{|V-\mathbf{1}|\le CM\alpha_0,} \\[3pt]
\FitMath{\operatorname{cn}(v)=|v|\,|v^{-1}|\le e^{CM\alpha_1},} \\[3pt]
\FitMath{|v'-\mathbf{1}|\le CM\alpha_1.}
\end{gather*}

The third bound of Ba{\l}aban, $|v-\mathbf{1}|\le CM\alpha_1$, holds for $v'$. For $v$ it holds only at pairs normalised by $U(\Gamma_{y_\square,x})=\mathbf{1}$. Indeed, for $U=\mathbf{1}^g$, $A'=0$ one has $v=g(\cdot)g(y_\square)^{-1}$, which is arbitrary, whereas $\operatorname{cn}(v)=1$.
\end{proposition}

\begin{proposition}[{transversality, locality, node factor; the choices disappear}]\label{prop:averaging-1a-3-9} The slice is transversal. Let $\tilde{A}$ be real with $\|\tilde{A}\|_1\le c$, and let a $G$-valued $\mathsf{y}$ near $\mathbf{1}$ be prescribed on the tops. Then some $w=e^{i\lambda}$ brings $\tilde{A}$ into the slice with $r_j=\mathsf{y}$ on $\mathbb{Y}_j$, and $\|\lambda\|_0\le C(|D^*\tilde{A}|_{(-2)}+\sup|\mathsf{y}-\mathbf{1}|)$. Here $r_j$ is defined as in \cref{prop:averaging-1a-1-7} and formed at $\tilde{A}$; at $\tilde{A}=0$ it is $\bar{R}^{(j)}_{U_0}w$. The pair $(A,u)$ depends locally on the data $(J(U),B)$, with tails $Ce^{-\delta d(y,y')}$ and $\delta<\delta_0$. Distances are measured in local scales, $J(U)$ in $|\cdot|_{(-3)}$, $B$ in the sup norm and $A$ in $\|\cdot\|_1$. Along differentiable families of $U$, this locality follows from the relation

\begin{gather*}
\FitMath{\delta A=\mathsf{G}_k(U_0)\big[\operatorname{Ad}(u^{-1})\delta J(U)+\mathcal{L}\delta A\big]+\mathsf{G}_k(U_0)Q^{\sharp*}\xi^{-3}\big[\delta B-\partial_AC(A)\delta A\big].}
\end{gather*}

For backgrounds over a real datum $\mathbf{V}$, assume that the datum is regular in the sense of \cref{ch:multilevel-backgrounds}: its test loops lie within a small $\varepsilon$ of $\mathbf{1}$, and, where inward bonds are kept, so do the two-scale loops at the seams. For the retained-inward variant also assume the original and peeled admissions of \cref{lem:mb-retained-edge-Schur}. Under these hypotheses the pair depends holomorphically on the datum within a ball of radius $c$, the representative at $e^{iB}\mathbf{V}$ being $(e^{i\eta_kA(B)}U_k(\mathbf{V}))^{h(B)}$. None of these statements depends on $k$. The plain averages of $e^{i\eta_kA}U_0$ are the datum transformed by the \emph{node factor}: $\mathcal{M}^j(e^{i\eta_kA}U_0)(c)=\mathsf{n}(c_-)V_j(c)\mathsf{n}(c_+)^{-1}$, where $\mathsf{n}=u^{-1}|_{\mathrm{nodes}}$. At the tops the node factor is the iterated $\mathsf{n}_W$, namely $v_j[A]$. At a junction node it contains one further factor, the quotient of the datum by the reduced average on one bond ending at that node. Every ingredient of a step is gauge covariant, and the map $(A,u)\to(\operatorname{Ad}(v)A,vuv^{-1})$ takes the set of Landau pairs of $U$ with respect to $U_0$ onto the set of Landau pairs of $U^v$ with respect to $U_0^v$. Consequently $E^{(k+1)}(X,\cdot)$ is invariant under 

\begin{gather*}
\FitMath{(\mathbf{U},\mathbf{J})\to(\mathbf{U}^u,\operatorname{Ad}(u)\mathbf{J})}
\end{gather*}

 for $G$-valued $u$. It is therefore also invariant for the $G^c$-valued $u$ that are connected to these within the domain, and it can be evaluated in any gauge. A function $F$ constructed on a slice is invariant under the subgroup that preserves the slice, and it has a unique invariant extension, $F^\natural(\mathfrak{z})=F(\mathsf{c}[\mathfrak{z}]\cdot\mathfrak{z})$. Here $\mathsf{c}[\mathfrak{z}]$ is the transformation that takes $\mathfrak{z}=(U,A',\mathbf{J})$ to the slice; this symbol, which carries an argument, is distinct from the scheme $\mathsf{c}$ of \cref{prop:averaging-1a-1-13}. The extension retains locality, the bounds, and holomorphy in $(A',\mathbf{J})$ at fixed real $U$.
\end{proposition}

\begin{lemma}[{the block axial gauge outside the nodes}]\label{lem:averaging-1a-3-11} Near a seam, the block of a coarse site contains nodes, and every element of $\mathcal{G}^{\mathrm{res}}$ is pinned at these nodes. In that case the gauge condition is $(\mathcal{M}^iU)(y,x)=\mathbf{1}$ for $y\in P''$ and for those $x\in B(y)\setminus y$ other than nodes. For a pair of nodes, the value $(\mathcal{M}^iU)(y,x)$ is invariant on the orbit, and the datum does not contain it. If $P''$ consists of nodes, the gauge is met once by every orbit of $\mathcal{G}^{\mathrm{res}}$, meaning that the averages $\mathcal{M}^iU$ on the bonds within $B(P'')$ and the nodes depend only on the orbit. Hence $V^{\mathrm{bg}}_k=\mathcal{M}^kU$, with $U$ in the orbit $U_{k+1}(\mathbf{V})$ and in the gauge, is a function of the datum; this is the case $i=k$, with $k+1$ in place of $k$. In this situation

\begin{gather*}
\FitMath{\mathrm{GF}_i(\mathcal{M}^iU;P'')=\sum_{y\in P''}\,\sum_{x\in B(y)\setminus y,\,x\text{ a node}}\big[1-\tfrac12\operatorname{Re}\operatorname{tr}(\mathcal{M}^iU)(y,x)\big],}
\end{gather*}

which is constant on the orbit. Near a seam, this is what is left of the identity $\mathrm{GF}_k(V^{\mathrm{bg}}_k)=0$.
\end{lemma}

\begin{proposition}[{a generalisation of \cref{prop:averaging-1a-3-6}, in the scheme $\mathsf{c}$}]\label{prop:averaging-1a-3-12} Fix $\gamma'>0$. Where the constant $B_1$ of the theorem enters, $\gamma'$ is not large; $\gamma'=1$ is an example. Assume the hypotheses of the theorem, with $c_1$ depending on $\gamma'$ as well; for arbitrary $\gamma'$, the constant $B_1$ must also depend on $\gamma'$. Let $f\in\operatorname{ran}R$ with $|f|_{(-2)}<\gamma'(\alpha_0+\alpha_1)$. Then the contraction in $\|\lambda\|_0$ from the proof of the theorem produces, for $\lambda$ in a ball of fixed radius, a unique $u$ whose averages are trivial and for which $R\,D^*_{U_0}A=f$. The consequence below uses only this conclusion. Ba{\l}aban's Theorem 8 [stated technical lemma; Group E] also asserts, without proof, all the bounds of the theorem, with the same $B_1$. The second-order and H{\"o}lder bounds require, in addition, a bound on $Df$ in $|\cdot|_{(-3)}$. As a consequence, the set $\{A:\mathcal{Q}^{\sharp}(U_0;A)=B,\ RD^*A=f\text{ for some small }f\}$ consists of pieces of gauge orbits, and each piece meets the slice $f=0$ in exactly one point. On each piece the action takes a single value. Hence a point of the slice that is critical there is critical on the thickened set as well, so that the Lagrange equations no longer involve the gauge condition. Ba{\l}aban makes use of this a single time, to obtain the second-order bounds on the minimiser. In \cref{sec:propagators-propagators-in-a-background-field-exponential-decay-produced-by-the} the corresponding step, which passes from the slice to a multiplier attached to the averages only, is carried out with gauge directions the images of which under $RD^*$ fill all of $\operatorname{ran}R$. The domains $U^c_j$ have real centres that are regular, though not nearly critical, and at these centres the construction demands small local potentials outright.
\end{proposition}

\begin{proposition}[{a seam in dimension one}]\label{prop:averaging-1a-3-13} Take all definitions on $\mathbb{Z}\supset L\mathbb{Z}$, and perform one step on the weight $1$ on two blocks, the second of which is centred at $y$; the third block, $B(y')$, remains fine. The seam bond $s$ is kept in the datum, $\mathsf{p}$ is the transport from $s$ to $y'$, and $f_{\mathfrak{G}}$ is the one-site factor of $\mathfrak{G}_k$. Let $V^+(c)$ be the coarse variable on $c=\langle y,y'\rangle$. Then the integral of the step can be computed exactly:

\begin{gather*}
\FitMath{\Phi=f_{\mathfrak{G}}\big(V^+(c)\,\mathsf{p}^{-1}V(s)^{-1}\big).}
\end{gather*}

Under the node-wise gauge group of the datum, this integral fails to be invariant. Let a transformation be $g$ at the centre $y$ and $\mathbf{1}$ at all remaining end points of bonds of the datum, the end of $s$ among them. It changes the integral into $f_{\mathfrak{G}}(gV^+(c)\mathsf{p}^{-1}V(s)^{-1})$. The integral is invariant under the action in which the end of $s$ moves with the centre of its block, and with respect to this action the tree insertion is a valid identity. Under the node-wise action, $\Phi$ has gauge average equal to the constant $1$. The procedure of \cref{prop:averaging-1a-3-3} therefore produces the bare slice measure, not $\Phi$ times that measure, and in this case the insertion is not an identity. If $s$ is integrated together with its block, then $\int dV(s)\Phi=1$, and nothing remains to vary.
\end{proposition}

\textbf{Proofs.} \emph{Averaging.} Every covariance property follows from the identity $\log(TXT^{-1})=T(\log X)T^{-1}$. Each expression under a logarithm is conjugated at a single site ($c_-$ for $\log W_x$). The norms are gauge invariant, and so therefore is the choice of rule. For \cref{lem:averaging-1a-1-1}, in the chart $V=e^YV_1$ the means are exactly the fixed points of the map $Y\mapsto Y+\frac1n\sum_i\log(V_iV_1^{-1}e^{-Y})$. 

Denote this fixed-point map by $\Phi$. For explicit radii, put
$d=\operatorname{diam}V_\bullet\le1/40$
and $T_i=V_iV_1^{-1}$. On $|Y|\le1/8$ one has
$|T_ie^{-Y}-1|\le(1+d)e^{1/8}-1<1/6$.
The exponential series and the logarithmic resolvent identity give
\[
 \operatorname{Lip}(\Phi)\le
 d+(1+d)(e^{1/8}-1)+\tfrac15(1+d)e^{1/8}<\tfrac25,
 \qquad |\Phi(0)|\le\frac{d}{1-d}.
\]
Thus $\Phi$ maps this closed ball into itself and its fixed point
satisfies $|Y_*|\le(5/3)d/(1-d)\le2d$. Its iterates from zero
converge uniformly, which proves holomorphy in the family on the
open diameter domain. For determinant-one inputs the iterates are
traceless: every logarithmic argument has determinant one and lies
within $1/6$ of one, so the trace of its series logarithm is zero.
For unitary inputs they are also anti-Hermitian. The constructed
mean therefore belongs to $G^c$, and to $G$ on real inputs.

To prove uniqueness among all admissible real means, let $V'\in G$
be any mean and put $Z_i=V_iV'^{-1}$. Then
$|Z_i-1|<1$, $\sum_i\log Z_i=0$, and $|Z_i-Z_1|\le d$.
For unitary $A,B$ with $|A-1|,|B-1|<1$, the identity
\[
 \log A-\log B
 =\int_0^1(1+t(A-1))^{-1}(A-B)(1+t(B-1))^{-1}\,dt
\]
implies $|\log A-\log B|\le2|A-B|$. Indeed each resolvent
has norm at most $(1-t/2)^{-1}$, since the real parts of the
eigenvalues of $A$ and $B$ exceed $1/2$, and
$\int_0^1(1-t/2)^{-2}\,dt=2$.
Consequently
\[
 |\log Z_1|
 =\left|\frac1n\sum_i(\log Z_1-\log Z_i)\right|
 \le2d<\tfrac18.
\]
With $Y=-\log Z_1$ we have $V'=e^YV_1$; its mean equation
says precisely $\Phi(Y)=Y$. Hence $Y=Y_*$ and
$V'=\operatorname{av}(V_\bullet)$. All constants above are
independent of the number of family members.

Under a conjugation that is not unitary, diameters change by at most $\operatorname{cn}(T)$. The isometry $\mathsf{r}$ permutes the centred blocks and the families $\mathcal{G}$, but it carries a comb to a comb in a different order. This is why Ba{\l}aban's first average, built on cornered blocks and combs, needed to be symmetrised. \Cref{lem:averaging-1a-1-2} (lattice homotopy): each interchange of two steps in a path costs one plaquette. The contours sweep out $O(L^2)$ plaquettes, and this determines $\mathrm{Dom}(c)$. In the comb gauge the bound increases with the size of the box, so small curvature yields a small potential only on boxes of bounded size. In the comb gauge one has $\mathcal{M}(e^{i\mathsf{a}})=\exp(i\mathsf{Q}^{\mathrm{pl}}\mathsf{a}+O(|\mathsf{a}|^2))$. Since the contour terms $\sigma$ cancel at the corners of $p'$, the lattice Stokes formula yields

\begin{gather*}
\FitMath{\sum_{c\in\partial p'}\pm(\mathsf{Q}^{\mathrm{pl}}\mathsf{a})(c)=L^{-4}\sum_{x\in B(y)}\sum_{p\subset S_x}\mathsf{a}(\partial p),} \\[3pt]
\FitMath{S_x=p'\text{, with corner }y\text{, translated to }x,\,L^2\text{ plaquettes}.}
\end{gather*}

Consequently $\hat\varpi_n=\varpi(\mathcal{M}^nU)/\alpha_0(L^n\eta_k)^2$ satisfies $\hat\varpi_{n+1}\le\hat\varpi_n(1+C\alpha_0L^{-2(k-n)})$. Measured in its own scale, the curvature of level $n$ is summable, and this outweighs the accumulation of $k$ errors.

A single step (\cref{prop:averaging-1a-1-5}) is proved by universality and Cauchy bounds, whereas Ba{\l}aban computes $\mathsf{Q}(W_0)$ in closed form from the Campbell--Hausdorff formula. In the comb gauge everything reduces to finitely many holomorphic functions that depend on $L$ alone, and hence $\mathcal{Q}^1=\mathsf{Q}(W_0)\mathsf{a}+O(|\mathsf{a}|^2)$. At $W_0=\mathbf{1}$ one has $\frac1i\log\mathsf{u}=\sigma+O(|\mathsf{a}|^2)$. Thus the block factors take away precisely the contour terms of $\mathsf{Q}^{\mathrm{pl}}$, and $\mathsf{Q}(\mathbf{1})=\mathsf{Q}^{\mathrm{s}}$. The kernels of $\mathsf{Q}(W_0)$ and $\mathsf{S}(W_0)$ are holomorphic and coincide at $\mathbf{1}$, so they differ by $O(\varpi)$. The norms of $\mathsf{S}$ reduce to counting: there are $L^4\cdot L$ pairs $(x,b)$ for each coarse bond, each fine bond is contained in $L$ segments, and every pair carries the weight $L^{-4}$. For $k$ steps, the factorisation follows from covariance by induction. At a complex $1$-form $A$, covariance under the $G^c$-valued $v_n$ uses up part of the room in the domain, and it is applied with $\operatorname{cn}(v_n)=|v_n|\,|v_n^{-1}|\le2$. Setting $t_n=\sup|\mathsf{a}_n|/L^n\eta_k$ and $\varpi_n=\varpi(W_{0,n})$, one has

\begin{gather*}
\FitMath{t_{n+1}\le t_n\big(1+C\varpi_n+CL^n\eta_kt_n\big),} \\[3pt]
\FitMath{\sum_n\varpi_n\le C\alpha_0,} \\[3pt]
\FitMath{\sum_{n<k}L^n\eta_k\le2L^{-1}.}
\end{gather*}

The recursion closes because $\|\mathsf{S}(W_{0,n})\|_{\mathrm{row}}$ does not exceed $L$, the ratio of the spacings, with no loss. This is the third gain of straightening. Since $\|\mathsf{Q}^{\mathrm{pl}}\|_{\mathrm{row}}\ge L+\frac{3(L^2-1)}{2L}>2L$, a product of one-step bounds for the plain average grows geometrically in $k$. The derivative is bounded through column sums and Schur's test. When the plaquettes are complex, the real comb gauge does not work: at every level its error is as large as the bound itself, which gives only $1+Ck\alpha_1$. The $G^c$-valued comb gauge of $\mathcal{M}^n\mathbf{U}$, whose condition number is $1+C\alpha_1L^n\eta_k$, gives $\varpi^c_{n+1}\le L^2\varpi^c_n(1+C\alpha_1L^n\eta_k)(1+CL^2\varpi^c_n)$. The price of the non-unitary conjugation is therefore one factor per level, and these factors are summable. The determining sets are obtained by folding with reflections. The double block of order $j$ beneath a bond of $\Gamma_j$ lies in $\Omega_{j-1}$. There the fields are, in the worst case, regular only in the next finer scale, and this costs one factor, a single time.

Intertwining: at a pure gauge $e^{i\mathsf{a}^w}W_0=W_0^w$ one has $\mathsf{u}(W_0,\mathsf{a}^w)=w\,(\bar{R}_{W_0}w)^{-1}$. It is for the sake of this identity that $\bar{R}$ is defined as it is. \Cref{prop:averaging-1a-1-7} rests on a coupled recursion for the size of $r_n$ and for its covariant difference, which the same geometric sums close. For products one has $|\mathsf{z}_{j+1}|\le|\mathsf{z}_j|+C(|\mathsf{z}_j|+\alpha L^{j+1}\eta_k)^2$, because the defect added at each step is quadratic in the variation within a block. This is the origin of the second condition in the definition of a class. \Cref{lem:averaging-1a-1-9}: up to an error $C\varpi$, a bond inside a block is the quotient of two contour transports, and the average then determines the crossing bonds. The comparison is made factor by factor. \Cref{prop:averaging-1a-1-10}: the constraint is solved to second order by means of the right inverse $\mathsf{H}_k$. At a critical point under the constraint, $J(U)=Q^{\sharp}_k(U)^*\omega$ is in general not zero. It is the image of a Lagrange multiplier $\omega$, whose size is that of the curvature of $V$. The extra term does not occur in abelian theories, and it gives rise to $G_1$, $\mathsf{H}_1$ in \cref{sec:propagators-propagators-in-a-background-field-exponential-decay-produced-by-the}. \Cref{prop:averaging-1a-1-11} follows from the first identity together with $\mathcal{M}(\mathsf{r}U)=\mathsf{r}\mathcal{M}(U)$ and $(\mathsf{r}U)(\mathsf{r}y,\mathsf{r}x)=U(y,x)$. \Cref{prop:averaging-1a-1-12} is purely algebraic. Writing $\hat{u}_n=u\,v_n$, the axial condition at level $n$ takes the form 

\begin{gather*}
\FitMath{\hat{u}_n(x)=\operatorname{Ad}(W_{0,n}(z,x))^{-1}[\hat{u}_{n+1}(z)\mathsf{u}_n(z)^{-1}\mathsf{w}^{\mathrm{rel}}_n(z,x)]}
\end{gather*}

. This fixes $u$ level by level, starting at the top and going down, and it yields $\bar{R}_{W_{0,n}}\hat{u}_n=\hat{u}_{n+1}$. \Cref{prop:averaging-1a-1-13}: two chains of transports with common end points differ by the curvature that they enclose. At level $n$ this is $O(\alpha_0L^{-2(k-n)})$, so the differences form a geometric series. The transport $t$ is unavoidable, because in each block the two schemes see a gauge transformation at two different sites. \Cref{prop:averaging-1a-1-14} follows from linearity and the Stokes formula, which give $\Theta=L^{-4}\sum_{x\in B(z)}L[\mathsf{f}^\Sigma(x_1+L)-\mathsf{f}^\Sigma(x_1)]$, together with a count.

\emph{The transformation.} \cref{prop:averaging-1a-2-1} combines the definition of an image measure with a property of $\mathcal{M}$. The renormalisation transformations obey the semigroup law $\mathbf{T}^n\nu=(\mathcal{M}^n)_*\nu$, which relates different spaces. In physical units the symmetry group becomes smaller as $k$ grows; in the units of the level it is the symmetry group of the unit lattice. Finite range does not survive, since $\mathbf{T}\rho_0$ couples all bonds. Reflection positivity is required of the bare measure only \cref{candidate:sources,ch:all-state-source-transfer}, \cref{ch:all-mirror-periodic}, \cref{ch:continuum-compactness,ch:all-mirror-periodic}. \Cref{prop:averaging-1a-2-2}: the variances are additive (taking a block mean divides by $L^d$, the rescaling multiplies by $L^{d-2}$, and the new kernel contributes $\mathfrak{a}^{-1}$). Nested backgrounds: backgrounds exist only over regular data, so minimising in two stages is not enough. The argument is instead as follows. A point that is critical on the larger fibre is critical on the smaller one, regularity improves in finer units, and uniqueness \cref{ch:multilevel-backgrounds} completes the proof. Because the argument involves critical points alone, it extends to complex data. The minimiser $\mathsf{H}B$ is the one given by the method of Lagrange multipliers. The delta kernel secures four things:
\begin{itemize}
\item closure under composition. Since $\mathcal{M}$ is non-linear, the composite of two soft kernels on $G$ is not a function of $W$ and $\mathcal{M}^2U$ alone.
\item a single large term. The critical point does not depend on the coupling, the leading term is exactly $x_kA(U_{k+1})$, and terms produced at $U_j$ can be evaluated at $U_k$ at every later level.
\item exact averages. Since $V_k=\mathcal{M}^kU$, a test on the block fields is a test on the bare field.
\item no analytic loss. The mass $aQ^{\sharp*}Q^{\sharp}$ remains in $\mathsf{G}_k$ as a device, with $a>0$ arbitrary.
\end{itemize}

The cost is that the fibre is curved. For the gauge fixing the opposite choice is made: it enters through a soft weight.

\Cref{prop:averaging-1a-2-3} (stochastic kernels, balance of masses). The mass, summed over the outcomes, is preserved by operations of three kinds. \emph{Tests:} at any given state exactly one outcome takes the value $1$. A test is counted as failed whenever the object it examines (a mean, a logarithm, a minimal orbit, a minimiser) is undefined, so every test is defined on all configurations. The domains and values of the real tests for the block average, regular critical orbit and unique attained minimum are Borel by the finite formulas for those objects in the \hyperlink{real-background-tests}{measurability argument below}. \emph{Draws:} the kernel $\kappa_{\mathfrak{G}}(V,\cdot)=\int du\,\mathfrak{G}_k(V^u)\delta_{V^u}$. \emph{Maps:} adjoining $\mathcal{M}(V_k)$, and forgetting a variable. No smallness is involved in any of this. The operation $\mathbf{R}$ draws the variables of a component afresh using the reference kernel specified below. Because of the partitions, single terms have only restricted symmetry. Only the totals at a base point are covariant under the full group $W_4$, which \cref{ch:pointed-hessian-cancellation} requires. Base points are required because the sum $\sum_{X\subset\Lambda_j}$ is not invariant and the cancellation against $b_jA(\varphi_j,\cdot)$ must be made at each point separately. When $z$ lies one layer of $M\check{R}_j$-cubes inside $\Lambda_j$, the restricted pointed total differs from the unrestricted one by $O(e^{-\frac12\kappa\check{R}_j})$.

\Cref{lem:averaging-1a-2-4}: a compact group admits no dilation. Hence $x_k$ \emph{is} the normalisation of the kinetic term. Gauge invariance links the quadratic, cubic and quartic parts of the marginal term, so the renormalisations of the field strength and of the coupling are given by a single number. The substitution $\mathcal{A}_k=g_k\mathcal{A}^{\mathrm{s}}_k$ takes the place of the rescaling. Consequently there is no scaling transformation, no fixed point and no eigen-direction. The word ``irrelevant'' can then mean only a bound $O(\eta_{k-j}^{4+\alpha_{\mathrm{H}}})$ at regular backgrounds, where $\alpha_{\mathrm{H}}<1$ is a H{\"o}lder exponent \cref{ch:pointed-hessian-cancellation}. Level $k$ is a problem on a unit lattice that carries a fine lattice of spacing $\eta_k\to0$. Every estimate in \cref{ch:multilevel-backgrounds}, \cref{sec:propagators-propagators-in-a-background-field-exponential-decay-produced-by-the} must therefore hold uniformly in $\eta_k$. \Cref{lem:averaging-1a-2-4}: the density itself does not change. Its exponent is always $-\beta A(U_k)$ together with the terms in their unrenormalised form. The point is which grouping of the terms admits uniform bounds:
\begin{itemize}
\item without grouping, the bound $\|E^{(j)}\|\le E_0$ yields no more than $CE_0\eta_k^{-4}|\mathbb{T}^{(k)}|$;
\item the number $b=b_{k+1}$ by itself leaves a geometric sum over $j$;
\item if the marginal terms are left among the corrections, they add up to $c_0kA$ to leading order;
\item once absorbed into $x_k$, they fix the next thresholds and $g_{k+1}$, the width of the following Gaussian.
\end{itemize}

The profiles $\varphi_j$ keep track of the place at which step $j$ was performed. In the constant, the contribution $-\frac12(N^2-1)\log x_k$ for each free bond is the Jacobian of $\mathcal{A}_k=g_kB$. Laplace's method gives $z_{\mathfrak{G}}\approx Cx_k^{-(N^2-1)/2}$. With four bonds per site, less one gauge direction, this gives $\mathsf{E}^{(j)}(T_\ell,T_\ell)=-(\frac32(N^2-1)\log x_j+O(1))(|\mathbb{T}^{(j)}|-|\mathbb{T}^{(j+1)}|)$. The term $\log x_k$ in $E_+$ arises from the constants $\mathsf{E}^{(j)}(T_\ell,T_\ell)$, $k\le j<K$, which belong to the steps $k+1,\dots,K$ and remain owed in $\bar\rho_k=e^{-E}\rho_k$ (when $k=K$ no constant remains owed, and $E_+=O(1)$). Over a large-field region, the constants differ by a volume factor $e^{C\log x_j\#_j(Z_{j+1})}$, and the small factors of \cref{sec:large-field-rare-patches-gains-per-creation-and-costs-per-site} must outweigh it.

\Cref{prop:averaging-1a-2-7} (the division of the work among the parts). Each part supplies the following:
\begin{itemize}
\item \cref{ch:multilevel-backgrounds}: the backgrounds;
\item \cref{sec:propagators-propagators-in-a-background-field-exponential-decay-produced-by-the}: the propagators;
\item \cref{ch:local-formation-program}: the new $E^{(k+1)}$, by a polymer expansion under a criterion of Koteck{\'y}--Preiss type;
\item \cref{ch:pointed-hessian-cancellation}: the coupling and the bounds on the sums;
\item \cref{sec:large-field-the-tests-the-nested-regions-and-the-content-of-a}: tests, regions and histories;
\item \cref{sec:large-field-stored-variables-and-integration-in-a-large-field-patch}: $\mathbf{R}$;
\item \cref{sec:large-field-rare-patches-gains-per-creation-and-costs-per-site}: the large-field factor and stability;
\item \cref{candidate:large_field}: the terms that $\mathbf{R}$ leaves behind;
\item \cref{sec:large-field-the-induction-closed-the-order-of-the-constants-the-sum}: the order of the constants and the inductive form;
\item \cref{sec:flow-the-run-its-coefficient-and-the-analytic-requirements}: the flow;
\item \cref{sec:flow-a-common-companion-and-the-torus-floor}: independence from the label, the torus and the cutoff.
\end{itemize}

The last item rests on locality. A term localised in $X$ is assembled from walks, potentials and tests within $X$. It therefore does not depend on the label, and, when $X$ does not wrap around the torus, it depends on neither the torus nor the cutoff. The small-field half requires $g_j\le\gamma$. It also requires $-C\le x_i-x_j\le C(j-i)$ for $i\le j\le k$, because the analyticity radii of level $j$ run with $g_j$ and the radii of two levels must be comparable. For radii that do not run, no relation between the couplings is required; Ba{\l}aban's remark that asymptotic freedom is not assumed refers to this case. For the complete model, what is used is the lower half of asymptotic freedom. At the levels of creation $j\le k$ the small factors were deposited, and since $x_j\ge x_k-C$ one has $p_0(g_j)\ge c\,p_0(g_k)$. The terms $R^{(j)}$ carry the vacuum subtraction but no counterterm for the coupling, and they are marginal. On the support of the tests, the following bound holds for every $X\in\mathcal{D}_j$ that contains a site $z$ of $\Omega_k$ and remains within $cML^{k-j}$ of $z$ in level-$j$ units (a longer $X$ satisfies $d_j(X)\ge cL^{k-j}-C$ and has only the bound $2g_j^{\kappa_0}e^{-\kappa d_j(X)}$; around a site of $\Omega_n\setminus\Omega_{n+1}$, $j\le n<k$, the right-hand side holds with $n$ in place of $k$):

\begin{gather*}
\FitMath{|R^{(j)}(X,U_k)-R^{(j)}(X,\mathbf{1})|\le C(d_j(X)+3)^2(x_j/x_k)^2\eta_{k-j}^4\,g_j^{\kappa_0}e^{-\kappa d_j(X)},} \\[3pt]
\FitMath{\sum_{j\le k}(x_j/x_k)^2g_j^{\kappa_0}\le Cg_k^{\kappa_0-2}:}
\end{gather*}

these terms do not decrease as $k-j$ grows. Only the linear growth of $x_j$ makes them summable, and it does so for $\kappa_0>6$. If one assumes no more than $x_j\ge x_k-C$, the sum can grow to order $kg_k^{\kappa_0}$. It is this summability that makes a counterterm for them unnecessary. It is also why aged $R$-terms change $b_j$ by no more than $C\gamma^{\kappa_0-C}$, so that $\lambda_\sharp>0$ is preserved \cref{sec:flow-the-run-its-coefficient-and-the-analytic-requirements}. The upper half ensures that the thresholds of adjacent levels are comparable: $\varepsilon_{k+1}\le C\varepsilon_k$.

\emph{The large-field factor: the reasons for these forms.} The fields at which a plaquette test fails do not form a small set for the measures contained in $\mathbf{T}_k(Z)$. They do form a small set for the Wilson weight, since $e^{-x_jA}\le e^{-c\,x_j\varepsilon_j^2}=e^{-c\,p_0(g_j)^2}$. Therefore, if no $\mathbf{R}$ has acted on a component $X$, the small factor that a failed plaquette test carries there is obtained from $\mathbf{T}_k(X)$ applied to $\chi\,e^{-A(x_k(\cdot),U_k)}$, and never from $\mathbf{T}_k(X)1$. This weight is evaluated at the background, which depends on the variables of all the components simultaneously. For this reason, bounds that are uniform in the exterior variables do not combine multiplicatively. The phenomenon already occurs in measure theory. Take Lebesgue measure on $[0,\mathsf{l}]$ in each of two variables, $0<\epsilon\ll\mathsf{l}$, and $F=\mathbf{1}\{|\xi_1-\xi_2|<\epsilon\}$. Then

\begin{gather*}
\FitMath{\sup_{\xi_2}\int F\,d\xi_1=2\epsilon,} \\[3pt]
\FitMath{\iint F\,d\xi_1d\xi_2=2\epsilon\,\mathsf{l}-\epsilon^2\gg(2\epsilon)^2.}
\end{gather*}

The remedy is a pointwise bound by a product of functions, each depending only on its own group of variables, the groups being disjoint, combined with Tonelli's theorem. The small factors come out as numbers from the Wilson sums near the creations, joined with a fixed fraction both of the Gaussian forms and of the kept gauge-fixing exponents. What remains of the weights is bounded by functions that no longer depend on the background. The supremum in $\mathsf{M}$ is placed under the integral for the same reason. The factor of $X$ then becomes a function of $\xi$ only, and the masses multiply. When a component has dropped out, the supremum is taken over the evaluations at $\xi$. These are the real configurations that are compatible with the stored variables and regular layer by layer, and at which every test listed by $\omega$ takes its recorded value, with one exception: the test whose failure turned $X$ into a drop-out. They form a set fibred over $\xi$. A product set contains a flat $\mathbf{U}$, at which $e^{-A}=1$, so taking the supremum over it eliminates the Wilson weight. For a live component, a bound $e^{-c\,p_0(g_k)}$ does not reproduce itself, since a step without a creation costs $e^{+CM^4\check{R}_k^5(d'_k(Z)+1)}$. What is carried along is the credit. As long as nothing is created, linear sizes counted in cells decrease geometrically. At a drop of $\check{R}$ the two partitions coincide and nothing shrinks, and no two consecutive steps are both drops. The waiting lasts at most $\check{R}_j$ steps, unless two components join or a mark is set. In that case the waiting begins anew. In the count of \cref{sec:large-field-rare-patches-gains-per-creation-and-costs-per-site}, every joining is charged to a separate created cell belonging to it alone, and a mark sets one factor $e^{-p_0(g_t)}$ of its mass against the same $\check{R}_t^6$. Hence the costs still owed add up to at most $CM^4L^5\check{R}_j^6(d'_j(Z)+1)$. The $p_0(g_j)^2$ per cube deposited at the creation covers $\check{R}_j^6$ per cube, and the powers of $\log g^{-2}$ are chosen in this order. After $\mathbf{R}$ has acted, the mass bound is what is carried forward.

Three losses are made good:
\begin{itemize}
\item The expansions use up part of $\kappa$ and of the radii. Cubes that are $L$ times larger give back a factor of order $L$ in the rate, and the curvature decreases by $L^{-2}$ at each level. $\mathbf{T}$ reserves a surplus of this gain for $\mathbf{R}$, whose re-expansions consume that surplus and restore exactly $\kappa$, $\alpha_i$.
\item The marginal term is absorbed into $x_k$.
\item $e^{-p_0(g_j)}$ is smaller than any power of $g_j$, but it is not a power of the spacing when $d=4$, so it cannot pay for all later steps. This is the reason for $\mathbf{R}$.
\end{itemize}

In the presence of large fields, $\log\rho_k$ is not a sum of local terms, and this is why the construction uses densities and sums over labels. Analyticity serves as a device. A single supremum bound on a complex domain stands in for bounds on all derivatives, and it is preserved under substitution, integration and taking logarithms.

\Cref{prop:averaging-1a-2-8} is \cref{ch:local-formation-program} carried out with parameters. The size of $\mathbb{P}^{(k)}$ is $g|B|^3$, and under the test $g|B|<\delta(\mathsf{t})$ this is $\le\delta(\mathsf{t})|B|^2$. The bracket of level $j$ in $\{S\}$ is $O(\mathsf{q}_j^3\,\eta_{k-j}^{4+\alpha_{\mathrm{H}}})$. Here, as in \cref{prop:averaging-1a-2-8}, $\mathsf{q}_j$ is the quotient of the analyticity radii at level $k+1$ by those at level $j$, and it is $\le C(k+2-j)^{1/2}$ thanks to the comparability assumed there. In \cref{ch:local-formation-program} the quotient is taken with the output radii, which are larger than those of level $k+1$ by a fixed factor, and in that part it goes by the name of the ratio. What the remainders contribute to the increment $b_j$ is a different quantity, even though \cref{ch:pointed-hessian-cancellation}, \cref{sec:large-field-the-induction-closed-the-order-of-the-constants-the-sum}, \cref{sec:flow-the-run-its-coefficient-and-the-analytic-requirements} and \cref{sec:flow-frozen-choices-slopes-and-tuned-bare-couplings} write it with the same letter. Each Cauchy estimate taken at a later background costs one power of $\mathsf{q}_j$, up to the third, and on radii that do not run there is no such loss. Applied to a shift of size $\delta$, Cauchy's estimate produces $\mathsf{s}$, and a cube of level $k$ contains $\eta_{k-j}^{-4}$ cubes of $\pi_j$. What is left at the age $n=k+1-j$ is $\mathsf{q}_j^3\,\eta_{k-j}^{\alpha_{\mathrm{H}}}\le C(n+1)^{3/2}L^{-\alpha_{\mathrm{H}}(n-1)}$, and these are the weights $\omega_n$. The $H^{(j)}$ are marginal, and their sum is controlled only through their norms. If the raw local terms are taken without the counterterms $b_jA$, then $\{S\}$ contains $\sum_{j\le k}b_j\{A\}\approx c_0k\{A\}$. In that case no bound holds uniformly in $k$, and the statement does not hold. Keeping $\mathsf{t}$ (which governs the tests, the gauge-fixing weight and the scaling) separate from $x$, the coefficient of $A$, comes at the price of a quadratic term proportional to $1-x/\mathsf{t}$. In return it allows frozen choices \cref{sec:flow-frozen-choices-slopes-and-tuned-bare-couplings} and complex weights \cref{sec:sources-the-curvature-source-with-fixed-reference-data}, \cref{ch:comparison-capsule-v7,ch:weighted-matched-domains}. Once the choices are fixed, the steps are fixed operations that are linear in the measure, and $\beta$ enters only through the weight of the action.

\Cref{prop:averaging-1a-2-9}: $\mathrm{GF}_k(V^{\mathrm{bg}}_k)=0$. The linear term $x_k\langle\mathsf{H}_k\mathcal{A}',J(U_{k+1})\rangle$ vanishes on $Q_{V^{\mathrm{bg}}_k}\mathcal{A}'=0$, since the current is the image of a Lagrange multiplier. Because the fibre is curved, the term $\langle\mathsf{H}_k\mathsf{h}C^{\mathrm{pl},(2)}(B),J\rangle$ remains, and it is included in $\Delta^{(k)}$. Hence $x_kA(U_{k+1})$ is the only place where a negative power of $g_k$ occurs. The choices are made in the following order. $b_{k+1}$ is obtained on the infinite lattice from the choices of level $k$, so $x_{k+1}$, $\varepsilon_{k+1}$, $\check{R}_{k+1}$ are fixed in advance of testing $V_{k+1}$. \emph{The new tests imply the old:} on $\Omega_{k+1}$, wherever the new functions are supported, the old ones take the value $1$. By the nesting, $U_{k,\square}=(e^{i\eta_k\mathbf{H}}U_{k+1,\square'})^u$. Here $\mathbf{H}$ denotes the response of \cref{ch:multilevel-backgrounds} when the datum is shifted, the shift being $O(\delta_k)$ everywhere but in a boundary layer at a distance of order $M_2\check{R}_k$. Therefore

\begin{gather*}
\FitMath{|\mathbf{H}|,|\nabla\mathbf{H}|\le C\big(\delta_k+e^{-\check{R}_k}\varepsilon_k\big)\le c\,\varepsilon_k,} \\[3pt]
\FitMath{|\partial U_{k,\square}-\mathbf{1}|\le2\varepsilon_{k+1}\eta_{k+1}^2+c\,\varepsilon_k\eta_k^2<\varepsilon_k\eta_k^2:}
\end{gather*}

so $\chi_k$ passes into $\chi_{k+1}$. This holds because the new threshold is smaller by the factor $L^{-2}\varepsilon_{k+1}/\varepsilon_k\le CL^{-2}$ (by the upper half of the flow inequalities, thresholds of neighbouring levels remain comparable), and because $\delta_k=(A_1/A_0)\varepsilon_k\ll\varepsilon_k$ (both carry the same power of the logarithm, and their ratio is small and fixed). Likewise every indicator in $\mathfrak{G}_k$ is equal to $1$, and $\chi'_k$ follows from $\chi^{(k)}$. Since the local backgrounds are so close to the global ones, the tests can be carried out on local objects, which factorise, while the expansions are made around global ones. \emph{The split:} an aged $R^{(j)}$ has each of its cubes inside a single block, its total transforms covariantly under the isometries of $\mathbb{T}^{(k+1)}$, and it is counted together with $\mathbf{E}_k$. For a young one this is not so, and a young one may not be admitted to the family from which $b_{k+1}$ is computed. The $t_1$-integral, which comprises boundary terms, the part of the coupling that depends on position, and the strip term $\mathsf{W}^{(k)}$, contributes only to $B^{(k+1)}$. The $t_2$-integral takes the expectation of the young terms in a measure which itself contains them. Consequently at least one factor $g_j^{\kappa_0}$ appears in every term of its expansion, and this integral contributes to both $R^{(k+1)}$ and $B^{(k+1)}$. The third term contains only regular terms. A further interpolation makes it independent of the exterior fluctuation field, and the resulting contribution goes to $B^{(k+1)}$. Its cluster expansion then produces, for $X\subset\Lambda_{k+1}$, the pure small-field integral's terms, $E^{(k+1)}$, and in addition boundary terms. $b_{k+1}$ is computed from this symmetric remainder. Thus a contribution to $R^{(k+1)}$ comes from $\mathbf{T}$ already, ahead of the one that $\mathbf{R}$ adds. \Cref{prop:averaging-1a-2-10}: iterating the Stokes formula of \cref{prop:averaging-1a-1-4} gives the first display. When a standard Gaussian vector is conditioned on the value of a linear map, its conditional expectation projects orthogonally onto the complement of the kernel of that map, and its components in mutually orthogonal subspaces are independent. The slice is a Schur complement, and $\mathcal{H}_{j-1}$ transports it. The only property of it that has to be analysed is its decay, and this is the analysis of \cref{sec:propagators-propagators-in-a-background-field-exponential-decay-produced-by-the} at $\mathbf{1}$.

\emph{Gauges.} \cref{lem:averaging-1a-3-1} is Stokes' formula on a tree. Because of the factor $|b_--y_P|$, boxes have side $O(M)$ and $M\varepsilon_k$ has to be small; that is, the coupling is taken small after $M$. \Cref{lem:averaging-1a-3-1} is a descent through the levels, with increments $C\alpha_0\sum_iL^{-2i}$. \Cref{prop:averaging-1a-3-3} follows from Fubini's theorem and the invariance of Haar measure. No determinant arises, because the gauge function is a group variable translated by $u(x)$. The mean is used because single combs break $W_4$. The constraint is taken soft because fixing the mean by a delta function imposes a non-linear restriction which couples all the bonds of the paths to one another and has no explicit solution. The weight is $x_k$, and with this weight

\begin{gather*}
\FitMath{x_k\,\mathrm{GF}_k(e^{ig_kB}V^{\mathrm{bg}}_k)=\tfrac12\sum_{y,x}\big[|B(y,x)|^2+O(g_k|B|^3)\big]}
\end{gather*}

is of order one. Here $g_kB(y,x)$ denotes the linear part in $B$ of $\frac1i\log$ applied to the mean $(e^{ig_kB}V^{\mathrm{bg}}_k)(y,x)$, and when $V^{\mathrm{bg}}_k=\mathbf{1}$, $B(y,x)$ coincides with the mean of $B(\Gamma)$ taken over $\mathcal{G}(y,x)$. This quantity is strictly positive on each direction $D\lambda$ with $\lambda|_{\mathrm{coarse}}=0$ (on these $|B(y,x)|=|\lambda(x)|$), and these directions exhaust the null directions of the Hessian. The weight itself is auxiliary: the first identity of \cref{prop:averaging-1a-3-3} yields an identity for every choice of weight, so the weight may be frozen \cref{sec:flow-frozen-choices-slopes-and-tuned-bare-couplings} and need not depend on the source \cref{sec:sources-the-curvature-source-with-fixed-reference-data}. \emph{Procedure and product:} the insertion, taken as a draw, is exact on any measure, and on invariant measures it agrees with the written product. A term with stored fields is invariant. Indeed, the distance from $P''$ to $Z_k$ is $\ge cM\check{R}_k$. Hence every transformation over $P''$ can be lifted to a bare transformation which is constant on pillars and equal to $\mathbf{1}$ near $Z_k$, the region where all stored quantities enter. In the case of the trees of \cref{sec:large-field-stored-variables-and-integration-in-a-large-field-patch}, the stored fields are acted on by the group via the tops of their pillars. The exact model \cref{prop:averaging-1a-3-13} shows that this verification is not idle: the group that acts is decided by whether or not the seam bonds lie in the determining set. The removal identity shows that dropping $\delta_T$ from the density gives an equivalent density, not an equal one: the two coincide when integrated against invariant functions.

\Cref{lem:averaging-1a-3-4}: in the pillars, $U$ and $U_0$ have the same contour variables. Hence $W(y,x)U_0(y,x)^{-1}=u(y)^{-1}\operatorname{Ad}(U_0(y,x))u(x)$, and the block factor $\mathsf{n}_W$ becomes the definition of $u^{-1}\bar{R}_{U_0}u$. This is Ba{\l}aban's deeper reason for imposing the axial gauge first and for keeping a single contour variable throughout (with two of them, circulations of $A$ survive in the linear part). \emph{Origin of $R$:} the term sought is Feynman's $\|\partial^*A\|^2$, but only the $\lambda\in N(Q')$ leave $\delta(B-Q^{\sharp}A)$ unchanged. Put \(N=\ker Q'\), equipped with the induced physical inner product and its Lebesgue measure \(d\lambda_N\). In the two gauge integrals below we integrate with this measure on \(N\). The Faddeev--Popov quotient with gauge parameter $\alpha_{\mathrm{g}}$ is

\begin{gather*}
\FitMath{e^{-\frac1{2\alpha_{\mathrm{g}}}\|\partial^*A\|^2}\Big/\int_N d\lambda_N\,e^{-\frac1{2\alpha_{\mathrm{g}}}\|\partial^*A-\Delta\lambda\|^2}=\mathrm{const}\cdot e^{-\frac1{2\alpha_{\mathrm{g}}}\|R\,\partial^*A\|^2}.}
\end{gather*}

In a real background, let \(E_R=\operatorname{ran}R\) carry its induced physical inner product, and let \(\delta_{E_R}\) be its unit delta distribution. The restriction \(T=\Delta_{U_0}|_N:N\longrightarrow E_R\) is a bijection; its absolute determinant is \(J_T=\det(T^*T)^{1/2}\), with the determinant taken on \(N\). Then

\begin{gather*}
\FitMath{e^{-\frac12\langle A,\Delta_\pi A\rangle}=J_T\int_N d\lambda_N\,\delta_{E_R}(RD^*A-T\lambda)\,e^{-\frac12\langle A-D\lambda,\Delta^{\mathrm{W}}(A-D\lambda)\rangle}.}
\end{gather*}

The change of variables \(y=T\lambda\) evaluates the integral at \(\lambda_*=T^{-1}RD^*A\), with factor \(J_T^{-1}\), and the definition of \(\Delta_\pi\) gives the identity. The determinant here is the Landau gauge's Faddeev--Popov determinant, and axial gauges have no such determinant. In normalised expressions such as $\mathsf{H}_kB$ it cancels; in Gaussian normalisers it does not. In the one-loop coefficient \cref{ch:pointed-hessian-cancellation}, \cref{sec:flow-the-run-its-coefficient-and-the-analytic-requirements} it enters together with the determinant of the gauge-fixed Hessian. Positivity (\cref{lem:averaging-1a-3-4}): consider $U=\mathbf{1}$ (the remaining cases are perturbations of it). A gradient $\partial\lambda$ ($\partial=D_{\mathbf{1}}$) is detected by $Q^{\sharp}$ if $\partial Q'\lambda\ne0$, and by $R\,\partial^*$ otherwise, since in that case $\lambda\in N(Q')$ once a constant is subtracted. The full $DD^*$ is positive as well, but admissible transformations cannot reach $D^*A=0$, and with the full operator $\mathsf{G}Q^{\sharp*}(Q^{\sharp}\mathsf{G}Q^{\sharp*})^{-1}$ is no longer the minimiser of the Hessian. $RD^*\mathsf{G}Q^{\sharp*}=0$ requires the invariant extension $\Delta_\pi$, which is a perturbation of size $CM\alpha_0$. (Ba{\l}aban's $G_0$ is defined by the form of the statement, and his $G$ goes with the identity. Adding the term of \cref{prop:averaging-1a-1-10} gives $G_1$. The three satisfy the same bounds, with the exception of the one with the covariant Laplacian, and they coincide where $J=0$.)

\Cref{prop:averaging-1a-3-7} (elliptic regularity through $\mathsf{G}_k(U_0)$): both currents, each smaller than $\alpha_0$, give $D^*DA=\mathcal{J}$ with $|\mathcal{J}|_{(-3)}\le2\alpha_0+C\alpha_2|\nabla A|_{(-2)}+\dots$. The averages give $\xi Q^{\sharp}A=B-C(A)=:\mathsf{b}$ with $|C(A)|\le C\alpha_2^2$ (\cref{prop:averaging-1a-1-6} leaves a remainder that is of second order), and this accounts for the $\alpha_1$ in the result and for the condition on $\alpha_2^2$. Put $\Delta'=\Delta^{\mathrm W}_{U_0}-D^*D$; the plaquette Hessian comparison gives $|\Delta'A|_{(-3)}\le C_L\alpha_0|A|_{(-1)}$. Use the fixed auxiliary row weights $a=\xi^{-2}$ in $\mathsf G_k$ for this estimate. On the slice,

\begin{gather*}
\FitMath{A=\mathsf{G}_k(U_0)(\mathcal{J}+\Delta'A)+\mathsf{G}_k(U_0)Q^{\sharp*}a\xi^{-1}\mathsf{b},} \\[3pt]
\FitMath{\|A\|_1\le B_0\big(|\mathcal{J}|_{(-3)}+C_L\alpha_0|A|_{(-1)}+\sup|\mathsf{b}|\big),}
\end{gather*}

where $a\xi^{-1}=\xi^{-3}$ for the chosen row weights. Since $|A|_{(-1)}\le\|A\|_1$, absorption follows once $B_0C_L(\alpha_0+\alpha_2)\le\frac12$. \Cref{prop:averaging-1a-3-6} (Ba{\l}aban's argument, written out for the scheme $\mathsf{c}$) is proved by induction on $k$. On $\Omega_k$ the pair $(A,u_1)$ of depth $k-1$ is a factor $L$ worse; here $u_1$ denotes the explicit product from \cref{prop:averaging-1a-1-12}. One seeks $u=u_1e^{i\lambda}$. By Campbell--Hausdorff, once the non-linear average has been straightened to $Q'\lambda=0$,

\begin{gather*}
\FitMath{\lambda=-G'RD^*A+G'R\,\mathfrak{N}_4(\lambda,D\lambda,A,D^*A),} \\[3pt]
\FitMath{|R\mathfrak{N}_4|_{(-2)}\le CB_1(\alpha_0+\alpha_1)\|\lambda\|_0,}
\end{gather*}

and this is a contraction on complex $\lambda$ in $\|\lambda\|_0$. (The bound is valid on the ball of radius $CB_1(\alpha_0+\alpha_1)$, which contains the solution; within a ball of fixed radius the solution is unique.) Curvature and mismatch thus have to beat fixed norms for the chosen block and gap geometry. These norms do not grow with the depth. The new constant is of order $LB_1$, and \cref{prop:averaging-1a-3-7} resets it to $B_1$; without this reset its growth in $k$ is geometric. Uniqueness holds among small $A$, so no Gribov ambiguity arises. In \cref{ch:multilevel-backgrounds} the passage to the slice is redone for $\mathcal{M}$. There Brouwer's theorem replaces the contraction. It provides a pair, though not its uniqueness, and with it the bounds on $\|A\|_1$ and $|D^*DA|_{(-3)}$. For a sequence with layers, the induction is replaced by a single preliminary gauge followed by an ascent through the truncations $(\Omega_0,\dots,\Omega_j)$. \Cref{prop:averaging-1a-3-8}: the cube is surrounded by a tower that descends to the bare scale. After passing to tower and box axial gauges, one applies the theorem to $(\mathbf{1},U)$ within the scheme $\mathsf{c}$. Only propagators at $\mathbf{1}$ enter, and this breaks the circularity with \cref{sec:propagators-propagators-in-a-background-field-exponential-decay-produced-by-the}. At complex fields we have $v=w^{-1}$ and $w(x)=\mathbf{U}(\Gamma_{y_\square,x})=\Pi\,v_{\mathrm{r}}^{-1}$. Here $\Pi$ is a product of factors $\exp(i\eta_k\,tA'(b)t^{-1})$, where each $t$ is a unitary transport of $U$, so $\operatorname{cn}(v)=\operatorname{cn}(\Pi)$. The potential of the background inside a large-field box is obtained in \cref{sec:large-field-stored-variables-and-integration-in-a-large-field-patch} by the same sequence of steps: the box axial gauge, used for the average, beneath it the tower axial gauge, and then Landau. Every change of its gauge is compensated on the stored variables. Local gauges $\hat{k}_y$ on overlapping boxes are glued together by $k(x)=\mathcal{N}\big(\sum_y\chi_y(x)\hat{k}_y(x)\big)$, $\mathcal{N}(X)=X/\sqrt{\det X}$. In $\operatorname{SU}(2)$, a real combination of group elements close to one another is a positive number times a group element, and the map $\mathcal{N}$ is holomorphic.

\Cref{prop:averaging-1a-3-9}: with weighted norms, or on complex balls, the induction incurs a loss at every level, and there is no reset. Instead one differentiates the identity of \cref{prop:averaging-1a-3-7}, with $\mathsf{G}_k$ bounded in the sup norms that carry the weight $e^{\delta d(\cdot,y')}$. This is a single linear estimate and needs no induction. That the averages are trivial does not force $u$ to be trivial at the nodes, and this gives rise to the node factor. Invariant functions do not detect it. At sites that carry stored fields but have ceased to be tops, $\mathsf{n}^{-1}$ is continued into each pillar by means of a smooth bump, and in this way the stored data are not acted on. \Cref{prop:averaging-1a-3-9}: the objects are words in transports and logarithms, each typed at a site, and fixed points are covariant by uniqueness. \Cref{prop:averaging-1a-3-6} gives the uniqueness of its Landau pairs only in the scheme $\mathsf{c}$; the set of these pairs is nevertheless covariant, since every one of the three defining conditions of a pair is covariant. The case of complex $u$ follows by the identity theorem, because the $G$-valued $u$ form a totally real submanifold of full dimension. \Cref{lem:averaging-1a-3-11} is a matter of covariance only: $(\mathcal{M}^iU^w)(y,x)=w(y)(\mathcal{M}^iU)(y,x)w(x)^{-1}$ with $w=\mathbf{1}$ at the nodes. \Cref{prop:averaging-1a-3-12}: the equation becomes $\lambda=G'(f-RD^*A)+G'R\,\mathfrak{N}_4$, whose inhomogeneity is of the size of the other terms. Ba{\l}aban states that the remaining proofs carry over with only minor changes. However, $\Delta_{U_0}A-D^*DA$ equals $Df+D(1-R)D^*A$ up to curvature terms, so the bounds of the theorem force $|Df|_{(-3)}\le C(\alpha_0+\alpha_1)$. On the other hand, $|Df|$ is of order $\eta_k^{-1}|f|$ for $f=\Delta\lambda$ with $\lambda\in N(Q')$ oscillating at the bare spacing. \Cref{prop:averaging-1a-3-13}: the contour variables of a block are coordinates that preserve Haar measure. All the integrals except one give $\int f_{\mathfrak{G}}=1$, and the remaining one, $\int dt\,f_{\mathfrak{G}}(t)\,\delta(tV(s)\mathsf{p}V^+(c)^{-1})$, is the display. In four dimensions the same occurs at each coarse site whose block contains an end of a seam bond.

\emph{Routine and delicate.} Routine are the algebra of means, of logarithms and of covariance, the exact identities, Stokes' formula and the Haar identities. The delicate points are the following:
\begin{itemize}
\item uniformity in $k$, which is reached solely via the straightened average, and only because curvature and potentials are measured on the scale of each level;
\item complex arguments, where each use of covariance costs room in the domain;
\item the two averages, together with the invariances of the constraints they impose, which have to be kept apart;
\item normalisation by averages instead of pointwise normalisation, the group property that is lost, and the node factor that is left over;
\item the junctions of layers;
\item the invariance of whatever an insertion encounters after fields have been stored.
\end{itemize}

For the construction as a whole, the difficulty lies in the words ``with the same constants''.


% END MODULE chapters/02_averaging.tex

% BEGIN PROOF EXPANSION formation/nonlinear_averaging_energy.tex
\subsection{Nonlinear averaging and the Wilson energy}
\label{candidate:nonlinear-averaging-energy}

The four-dimensional normalization of the contour average gives an
almost-contraction of the Wilson energy. Its leading constant is one.
This matters when the average is repeated: a fixed larger constant at
each step would grow with the number of levels. We prove the estimate
for the symmetric contour map in \eqref{eq:rg-objects-R3}, using its
near-identity means as in \cite{B85a}.

Work first on a periodic lattice of spacing one. Its side lengths
are divisible by a fixed odd integer \(L\ge3\). The coarse sites form
\(L\mathbb Z^4\), and \(B(y)=y+\mathcal B\) is the translated
centred block of \(L^4\) fine sites belonging to \(y\). A translated
choice of the block origins gives the same proof. Each independently
oriented fine plaquette is counted once. For a real \(\operatorname{SU}(2)\)
configuration \(W\), write
\begin{equation}
 F_p=-i\log\partial W(p),\qquad
 \varepsilon(W)=\max_p|F_p|,\qquad
 A(W)=\sum_p\left(1-\tfrac12\operatorname{Tr}\partial W(p)\right).
 \label{ne:curvature-energy}
\end{equation}
The logarithms are the small branches. Lie algebra norms use
\(\langle X,Y\rangle=\tfrac12\operatorname{Tr}(XY)\), so
\(1-\tfrac12\operatorname{Tr}e^{iF}=1-\cos|F|\).
All constants below can depend on \(L\), and none depends on the
fine spacing, number of blocks or number of averaging steps.

\begin{proposition}[One averaging step]
\label{prop:nonlinear-average-energy}
There are \(\varepsilon_L>0\) and \(C_L<\infty\) such that
\(\varepsilon(W)\le\varepsilon_L\) admits every local contour
mean and logarithm defining \(\mathcal M(W)\), and
\begin{align}
 \varepsilon(\mathcal M(W))
     &\le L^2\varepsilon(W)+C_L\varepsilon(W)^2,
 \label{ne:curvature-step}\\
 A(\mathcal M(W))^{1/2}
     &\le(1+C_L\varepsilon(W))A(W)^{1/2}.
 \label{ne:energy-step}
\end{align}
These estimates hold for real fields. No critical-point equation or
inverse Hessian estimate is required.
\end{proposition}

\subsubsection{The linear flux weights}

Let \(a\) be a Lie algebra one-cochain. If \(t_a(y,x)\)
denotes the mean of its line integrals over the prescribed shortest
contour paths, differentiation of the actual contour formula at the
identity gives, for \(c=(y,y+Le_\mu)\),
\begin{align}
 (Qa)_\mu(y)={}&L^{-4}\sum_{x\in B(y)}
       \sum_{i=0}^{L-1}a_\mu(x+ie_\mu)
                    +h_a(y)-h_a(y+Le_\mu),\notag\\
 h_a(y)={}&L^{-4}\sum_{x\in B(y)}t_a(y,x).
 \label{ne:linear-contour-average}
\end{align}
Indeed the logarithm of a product has, at the identity, the sum of
the logarithmic increments. The derivative of a symmetric mean of
nearby transports is their arithmetic mean. The central straight
transport in the contour formula cancels its copy inside the
averaged logarithm. The two endpoint transports give the displayed
coarse gradient of \(h_a\).

The curl of this gradient vanishes. The remaining straight
transports bound an \(L\)-by-\(L\) plaquette, so exact discrete
Stokes gives
\begin{equation}
 (\mathrm{d}_cQa)_{\mu\nu}(y)
 =L^{-4}\sum_{x\in B(y)}\sum_{i,j=0}^{L-1}
      (\mathrm{d}a)_{\mu\nu}(x+ie_\mu+je_\nu)
 =\sum_pT_{C,p}(\mathrm{d}a)_p.
 \label{ne:positive-flux}
\end{equation}
Here \(C=(y,\mu,\nu)\) is a coarse plaquette, and \(T_{C,p}\)
is zero for other fine orientations. Every coefficient is
nonnegative. Its row and column sums are exactly
\begin{equation}
 \sum_pT_{C,p}=L^2,\qquad \sum_CT_{C,p}=L^{-2}.
 \label{ne:row-column-sums}
\end{equation}
The row sum counts \(L^4\) starting sites and \(L^2\)
plaquettes per surface, each with weight \(L^{-4}\). For the
column sum fix a fine plaquette and its orientation. In each of
its two coordinate directions, each of the \(L\) choices of
surface offset determines one block residue and one coarse
origin. In the transverse directions the block residue determines
the coarse origin uniquely. There are therefore \(L^2\)
representations, with total weight \(L^2L^{-4}\). This count
also holds modulo the periodic side lengths, with multiplicity.

More generally let \(O_{C,p}\) be arbitrary orthogonal color
matrices. Cauchy--Schwarz and \eqref{ne:row-column-sums} imply
\begin{align}
 \sum_C\left|\sum_pT_{C,p}O_{C,p}f_p\right|^2
 &\le L^2\sum_{C,p}T_{C,p}|f_p|^2
  =\sum_p|f_p|^2.
 \label{ne:transported-flux-contraction}
\end{align}
Thus independent local gauge frames do not change the leading
constant. In dimension \(d\) the same argument gives norm
\(L^{2-d/2}\); it equals one in dimension four.

\subsubsection{Local gauges and the nonlinear remainder}

For each coarse plaquette \(C\), choose a rectangular fine box
\(D_C\) containing all paths in its four contour averages, with
a fixed additional layer. Its side lengths are at most \(cL\).
The family of these boxes has bounded overlap. If a box crosses a
periodic seam, perform this construction on a lift to the covering
lattice. Path identities and the following overlap counts include
the resulting periodic multiplicities.

Choose a coordinate comb tree in \(D_C\), rooted at the initial
vertex of \(C\), and gauge its edge transports to one. Write the
resulting local field as \(W^{u_C}_b=e^{ia_b}\). There are
constants depending only on \(L\) such that
\begin{equation}
 \max_{b\subset D_C}|a_b|\le C_L\varepsilon(W),
 \qquad
 \left(\sum_{b\subset D_C}|a_b|^2\right)^{1/2}
     \le C_L\left(\sum_{p\subset D_C}|F_p|^2\right)^{1/2}.
 \label{ne:local-comb-estimate}
\end{equation}
Enlarge \(D_C\) by another fixed layer if needed for the
plaquettes in the right side.

Here are finite proofs of these bounds. A remaining edge, together
with its two comb paths, forms a loop. Move the coordinate steps
of one path past those of the other, cancelling immediate
backtracks. Each interchange inserts one elementary plaquette,
conjugated by a prefix transport. There are at most \(C_L\)
such interchanges. This gives an exact product of conjugated
plaquette holonomies for \(W^{u_C}_b\). All factors and conjugating
transports are unitary, and
\(\|1-\prod_iV_i\|_{\rm op}\le\sum_i\|1-V_i\|_{\rm op}\).
The small logarithm has a uniformly bounded derivative near one.
Consequently
\(|a_b|\le C\sum_pn_{b,p}|F_p|\), where the nonnegative
integers \(n_{b,p}\) belong to a finite matrix of bounded row
and column sums depending only on \(L\). Its Schur bound gives
the second inequality in \eqref{ne:local-comb-estimate}; its row
bound gives the first.

These same local charts admit the actual means. For finitely many
nearby transports \(V_i\), the defining equation
\(\sum_i\log(V_iV^{-1})=0\) has derivative minus the number
of transports times the identity in the logarithmic coordinate of
\(V\), at the identity. The analytic implicit function theorem
therefore gives a unique small mean on a fixed ball. There are
only finitely many local patterns for fixed \(L\). Decrease
\(\varepsilon_L\) so that \eqref{ne:local-comb-estimate}
places all their transports in these balls. Gauge covariance of
the mean follows from uniqueness and conjugation of the small
logarithm. Means constructed in different local gauges therefore
give the same original coarse transport.

Let \(F_p^{(C)}=\operatorname{Ad}(u_C(p_-))F_p\). Taylor's
formula for the fine plaquette logarithm on this local chart gives
\begin{equation}
 F_p^{(C)}=(\mathrm{d}a)_p+r_p(a),\qquad
 \left(\sum_{p\subset D_C}|r_p(a)|^2\right)^{1/2}
 \le C_L\varepsilon(W)
          \left(\sum_{p\subset D_C}|F_p|^2\right)^{1/2}.
 \label{ne:fine-Taylor}
\end{equation}
For clarity, the Taylor remainder of one plaquette is bounded by
\(C\max_b|a_b|\sum_{b\subset\partial p}|a_b|\) on this
fixed matrix ball. Sum its square, use the bounded incidence of
edges and plaquettes, and apply \eqref{ne:local-comb-estimate}.
This proves the displayed estimate, including its dependence on
the curvature rather than merely on the number of edges.

The coarse plaquette logarithm is an analytic function of the
same finite list \(a\), vanishes at zero, and has derivative
\((\mathrm{d}_cQa)_C\) by \eqref{ne:linear-contour-average}. Its Taylor
remainder has the bound
\[
 |r_C^{\rm av}(a)|\le C_L\max_b|a_b|
                       \left(\sum_{b\subset D_C}|a_b|^2\right)^{1/2}.
\]
Indeed the second derivative is bounded on a smaller fixed
analytic ball, and the equivalence of the finite-dimensional
norms costs only a constant depending on \(L\). Substitute
\eqref{ne:fine-Taylor} in \eqref{ne:positive-flux}. After
conjugating the coarse logarithm back, this gives
\begin{equation}
 F_C^{\rm av}=\sum_pT_{C,p}O_{C,p}F_p+R_C,
 \qquad
 |R_C|\le C_L\varepsilon(W)
              \left(\sum_{p\subset D_C}|F_p|^2\right)^{1/2},
 \label{ne:actual-nonlinear-flux}
\end{equation}
Here the sum keeps separate the fine plaquette occurrences in
the chosen lift of \(D_C\). Each \(O_{C,p}\) is an adjoint rotation
and hence orthogonal. If several lifted occurrences project to the
same physical plaquette, their rotations need not agree. In the
global column sum group their scalar weights, retaining the separate
rotations in the flux sum. The resulting physical column weight is
still \(L^{-2}\), and every row weight is still \(L^2\).
Thus the coefficients are exactly the weights of
\eqref{ne:positive-flux}, with their periodic multiplicities;
no nonlinear coefficient replaces either leading sum.

\begin{proof}[Proof of \cref{prop:nonlinear-average-energy}]
The row bound in \eqref{ne:row-column-sums} and the finite size
of \(D_C\) give \eqref{ne:curvature-step} from
\eqref{ne:actual-nonlinear-flux}. Bounded overlap and
\eqref{ne:transported-flux-contraction} give
\begin{equation}
 \left(\sum_C|F_C^{\rm av}|^2\right)^{1/2}
 \le(1+C_L\varepsilon(W))
                    \left(\sum_p|F_p|^2\right)^{1/2}.
 \label{ne:log-curvature-contraction}
\end{equation}
The scalar Taylor estimate
\[
 \tfrac12t^2(1-Ct^2)\le1-\cos t\le\tfrac12t^2
\]
holds for small real \(t\). Apply the upper estimate to the
coarse field and the lower estimate to the fine field. Absorbing
\((1-C\varepsilon(W)^2)^{-1/2}\) in
\(1+C_L\varepsilon(W)\) proves \eqref{ne:energy-step}.
\end{proof}

\subsubsection{Iteration without a loss in the number of levels}

\begin{corollary}[An all-depth energy bound]
\label{cor:iterated-average-energy}
There is \(\alpha_L>0\) such that, for every integer \(k\ge0\)
and every lattice admitting \(k\) block averages, the condition
\begin{equation}
 \varepsilon(U)\le\alpha L^{-2k},\qquad0<\alpha\le\alpha_L,
 \label{ne:bare-regularity}
\end{equation}
admits all the near-identity branches used in those averages and
implies
\par\noindent
\begin{minipage}{\linewidth}
\begin{equation}
 \varepsilon(\mathcal M^jU)\le2\alpha L^{-2(k-j)}
 \quad(0\le j\le k),\qquad
 A(\mathcal M^kU)\le2A(U).
 \label{ne:iterated-energy}
\end{equation}
\popQED\par
\end{minipage}
\end{corollary}

\begin{proof}
Write \(\varepsilon_j=\varepsilon(\mathcal M^jU)\).
The first recurrence of \cref{prop:nonlinear-average-energy} gives
\[
 \varepsilon_{j+1}\le L^2\varepsilon_j
                    (1+C_L\varepsilon_j/L^2).
\]
Suppose the curvature bound in \eqref{ne:iterated-energy} holds
up to a given level. Up to that level,
\[
 \sum_{j<k}\varepsilon_j
 \le2\alpha\sum_{s\ge1}L^{-2s}
 =\frac{2\alpha}{L^2-1}.
\]
Iterating the recurrence therefore bounds
\(\varepsilon_j/(\alpha L^{-2(k-j)})\) by
\(\exp[2C_L\alpha/(L^2(L^2-1))]\).
Choose \(\alpha_L\) so that this is strictly less than two
and \(2\alpha_L<\varepsilon_L\). The first possible failure
of the curvature bound is then impossible, and every mean used
at the next level is admitted.

Iterate \eqref{ne:energy-step}, retaining its leading constant
one. It gives
\[
 A(\mathcal M^kU)^{1/2}
 \le\exp\left(C_L\sum_{j<k}\varepsilon_j\right)A(U)^{1/2}
 \le\exp\left(\frac{2C_L\alpha}{L^2-1}\right)A(U)^{1/2}.
\]
Decrease \(\alpha_L\) so that the last exponential is at most
\(\sqrt2\). Squaring proves the energy assertion.
\end{proof}

For example, if \(\nu(p)\ge\nu_*>0\) and one coarse plaquette
obeys \(\|1-\partial\mathcal M^kU(C)\|_{\rm op}\ge a\), then
\begin{equation}
 A(\nu,U)\ge\nu_*A(U)
       \ge\tfrac14\nu_*a^2.
 \label{ne:average-event-energy}
\end{equation}
The last inequality uses
\(1-\tfrac12\operatorname{Tr}V=\tfrac12\|1-V\|_{\rm op}^2\)
for \(V\in \operatorname{SU}(2)\), and \eqref{ne:iterated-energy}.
This is a global one-event certificate under
\eqref{ne:bare-regularity}. Several history events require disjoint
ownership or a bounded-overlap allocation of the action. A field
containing an unregularized old core does not satisfy the hypotheses
of this corollary.

% END PROOF EXPANSION formation/nonlinear_averaging_energy.tex


% BEGIN MODULE chapters/block_average_pivot_chart_v8.tex
\subsection{The block average and its pivot chart}
\label{ch:block-average-pivot-chart}

The finite calculation here concerns the actual centered one-step
average.  Its definition uses an implicit group mean for each internal
contour transport, followed by the exponential of an arithmetic mean
of closed logarithms for the block bond. We use both operations.  The
calculation proves the local pivot and quadratic-remainder estimates;
it does not prove a many-level background theorem or an exterior-force
estimate.

\subsubsection*{The exact finite map}

Work on a unit-spacing lattice in dimension four, with odd \(L\ge3\).
The block at a coarse site \(y\) is
\(B(y)=\{x:|x-y|_\infty\le(L-1)/2\}\).
For a positive coarse bond \(c=(y,y+Le_\mu)\), put
\(y'=y+Le_\mu\) and \(D_c=B(y)\cup B(y')\).
On a torus we assume that this double block and its one-bond
neighborhood have an embedded rectangular lift preserving their
unit bond graph, with no periodic edge joining opposite faces of
the patch.  A sufficient condition for a cubic torus is that its
period exceeds \(4L+4\).  All bonds used below have both endpoints
in \(D_c\), with the rectangular bond graph of this lift.

For \(x\in B(y)\), let \(\mathcal G(y,x)\) be the shortest paths
which complete one coordinate displacement before starting the next.
There are at most \(24\) such paths, each of length at most \(2L\).
Repeated paths, if retained, are counted with their prescribed equal
weights.  The transport along an oriented path is the ordered product
of its bond matrices; reversed bonds use the inverse matrix.
For nearby matrices \(T_1,\ldots,T_n\), define their contour mean
as the solution near them of
\begin{equation}
 \frac1n\sum_{j=1}^n\log(T_jT^{-1})=0.
 \label{eq:block-average-pivot-chart-QC1}
\end{equation}
All logarithms below are the local analytic logarithm near the
identity.  Set
\begin{equation}
 U(y,x)=\operatorname{av}\big(U(\Gamma):\Gamma\in\mathcal G(y,x)\big),
 \qquad U(x,y)=U(y,x)^{-1}.
 \label{eq:block-average-pivot-chart-QC2}
\end{equation}
With \(x'=x+Le_\mu\), the actual block map is
\begin{equation}
 \begin{split}
 Z_x(U)&=U(y,x)U([x,x'])U(x',y')U[c]^{-1},\\
 \mathcal M(U)(c)
   &=\exp\left[L^{-4}\sum_{x\in B(y)}\log Z_x(U)\right]U[c].
 \end{split}\label{eq:block-average-pivot-chart-QC3}
\end{equation}
Here \([x,x']\) and \([c]\) are straight segments of length \(L\).
In particular, the outer operation in \eqref{eq:block-average-pivot-chart-QC3} is not the implicit
mean \eqref{eq:block-average-pivot-chart-QC1}.  The cruder Borel definitions of the block map away
from this analytic neighborhood do not enter the calculation.

For a background \(V\), let \(W=\mathcal M(V)\).  In physical bond
logarithms, with no scale factor in the exponential, put
\begin{equation}
 q_{V,c}(a)=\frac1i\log\left[
       \mathcal M(e^{ia}V)(c)\mathcal M(V)(c)^{-1}\right].
 \label{eq:block-average-pivot-chart-QC4}
\end{equation}
The matrices \(a(b)\) are complex traceless \(2\) by \(2\) matrices,
and their norm is the supremum of the operator norms over positive
bonds.  On real variables they are Hermitian.  All constants denoted
\(C,c\) below are universal; their values may be enlarged or decreased
at successive occurrences.

\subsubsection*{The contour mean, with quantitative bounds}

The elementary matrix inequalities
\begin{equation}
 \|e^X-I-X\|\le\tfrac12e^{\|X\|}\|X\|^2,
 \qquad
 \|\log(I+H)-H\|\le\frac{\|H\|^2}{2(1-\|H\|)}
 \label{eq:block-average-pivot-chart-QC5}
\end{equation}
hold for \(\|H\|<1\), by the power series and submultiplicativity.
They show, on a fixed small ball, that
\begin{equation}
 \begin{aligned}
 \log(e^Xe^{-Z})&=X-Z+E(X,Z),\\
 \|E(X,Z)\|&\le C(\|X\|+\|Z\|)^2,\\
 \|D_ZE(X,Z)\|&\le C(\|X\|+\|Z\|).
 \end{aligned}
 \label{eq:block-average-pivot-chart-QC6}
\end{equation}
For the derivative estimate one can differentiate the convergent
logarithm series term by term and use
\(D e^Z[H]=\int_0^1e^{(1-t)Z}He^{tZ}\,dt\).
The derivative of each remainder in \eqref{eq:block-average-pivot-chart-QC5} is bounded by a constant
times the radius.  Thus the constants in \eqref{eq:block-average-pivot-chart-QC6} depend on no number
of contour paths.

Write \(T_j=e^{X_j}\), with \(\max_j\|X_j\|\le t\), and seek
\(T=e^Z\).  \eqref{eq:block-average-pivot-chart-QC1} is equivalent to
\begin{equation}
 Z=\overline X+\frac1n\sum_jE(X_j,Z),
 \qquad \overline X=\frac1n\sum_jX_j.
 \label{eq:block-average-pivot-chart-QC7}
\end{equation}
For a fixed sufficiently small \(t\), the right side maps
\(\|Z\|\le2t\) into itself and has Lipschitz constant at most
\(1/2\).  Iteration from \(\overline X\) proves existence and
uniqueness there.  It also proves
\(Z=\overline X+O(t^2)\), uniformly in \(n\).
The iterates are holomorphic in all \(X_j\), and convergence is
locally uniform on a slightly smaller fixed ball.  Their limit is
therefore holomorphic.  Trace zero is preserved in \eqref{eq:block-average-pivot-chart-QC7}.
For anti-Hermitian \(X_j\), anti-Hermiticity is preserved as well,
so the mean lies in \(\operatorname{SU}(2)\) on real group inputs.

Uniqueness proves permutation symmetry and left-right covariance.
Indeed under \(T_j\mapsto uT_jv\), with \(T\mapsto uTv\), every
logarithm in \eqref{eq:block-average-pivot-chart-QC1} is conjugated by \(u\), whenever the transformed
families remain in their analytic domains.  The inverse property
follows from
\(\log(T_j^{-1}T)=-\operatorname{Ad}(T^{-1})\log(T_jT^{-1})\).
Thus the locally constructed mean is the mean in \eqref{eq:block-average-pivot-chart-QC2}, including
families centered at a unitary matrix other than \(I\).

\subsubsection*{Local gauges do not constrain coarse holonomies}

For a real background \(V_0\) on \(D_c\), put
\(\varpi_c=\sup_{p\subset D_c}\|\partial V_0(p)-I\|\).
There is a unitary gauge \(s_c\) on this rectangle such that
\begin{equation}
 \sup_{b\subset D_c}\|V_0^{s_c}(b)-I\|
                                      \le C L\varpi_c.
 \label{eq:block-average-pivot-chart-QC8}
\end{equation}
Here \(V^s(b)=s(b_-)V(b)s(b_+)^{-1}\).
To prove \eqref{eq:block-average-pivot-chart-QC8}, fix a coordinate-order comb from \(y\) to every
site of the rectangle and set \(s_c(x)=V_0(\Gamma_{y,x})\).
For any bond \(b\), the product
\(s_c(b_-)V_0(b)s_c(b_+)^{-1}\) is its comb loop.  Slide the
additional bond past the later coordinate segments of the comb.
Each slide exchanges two adjacent perpendicular steps and inserts
one plaquette, conjugated by a partial transport.  There are at
most \(CL\) slides, since every coordinate interval in the rectangle
has length at most \(2L\).  All these conjugations are unitary.
The norm of a product of such plaquette factors minus \(I\) is
at most the sum of their deviations, proving \eqref{eq:block-average-pivot-chart-QC8}.

The same gauge, fixed from the real background, works for
\(V=e^{iE}V_0\): its distance from \(I\) is at most
\(C(L\varpi_c+\|E\|_\infty)\), when these quantities are small.
More generally we can start directly with a unitary local gauge
in which
\begin{equation}
 \sup_{b\subset D_c}\|V^{s_c}(b)-I\|\le\epsilon .
 \label{eq:block-average-pivot-chart-QC9}
\end{equation}
Different double blocks may use different gauges.  We never require
one gauge close to \(I\) on an entire observation window.

The group mean and \eqref{eq:block-average-pivot-chart-QC3} are gauge covariant.  The internal transport
transforms with its two endpoint factors; each \(Z_x\) transforms
by conjugation at \(y\).  Consequently
\begin{equation}
 q_{V^{s_c},c}(a^{s_c})
    =\operatorname{Ad}(s_c(y))q_{V,c}(a),
 \qquad a^{s_c}(b)=\operatorname{Ad}(s_c(b_-))a(b).
 \label{eq:block-average-pivot-chart-QC10}
\end{equation}
All these adjoint actions have norm one.  A coarse bond value
\(\mathcal M(V)(c)\) in the original gauge may be anywhere in
\(\operatorname{SU}(2)\).  It is the ratio in \eqref{eq:block-average-pivot-chart-QC4}, rather than that
coarse bond itself, whose logarithm is taken near \(I\).
Nontrivial coarse holonomies therefore cause no loss in these
local estimates.

\subsubsection*{Uniform Taylor estimates for the actual average}

We first work in the gauge \eqref{eq:block-average-pivot-chart-QC9}, so \(V=e^{iv}\) with
\(\|v\|_\infty\le C\epsilon\).  If
\(\|v\|_\infty,\|a\|_\infty<c_0/L\), every path product entering
\eqref{eq:block-average-pivot-chart-QC2} stays in a fixed small neighborhood of \(I\).
For example, a product of \(N\) factors, each within \(t\) of \(I\),
is within \(e^{CNt}-1\) of \(I\).  Its reversed factors satisfy
the same bound.  Here the internal paths have \(N\le2L\).
The contour-mean estimate just proved applies, and the closed
products in \eqref{eq:block-average-pivot-chart-QC3}, of total path length at most \(6L\), are also
in a fixed logarithm neighborhood.  The outer logarithm average
is normalized by \(L^{-4}\); its size is bounded by its largest
summand, without a factor \(L^4\).

It follows that
\begin{equation}
 F_c(v,a)=\frac1i\log\left[
   \mathcal M(e^{ia}e^{iv})(c)\mathcal M(e^{iv})(c)^{-1}\right]
 \label{eq:block-average-pivot-chart-QC11}
\end{equation}
is holomorphic and bounded by a universal constant on the joint
supremum balls of radius \(c_0/L\).  It vanishes at \(a=0\).
This establishes the domain before any derivative estimate.
On balls a fixed fraction smaller, two scalar Cauchy formulas
give
\begin{equation}
 \|D_a F_c\|_{\infty\to\mathrm{op}}\le CL,
 \qquad
 \|D_a^2F_c\|+\|D_vD_aF_c\|\le CL^2.
 \label{eq:block-average-pivot-chart-QC12}
\end{equation}
For example, for unit supremum directions \(u,w\), apply Cauchy
to \(F_c(v,a+zu+tw)\), with both scalar circles of radius
\(c/L\).  This proves the mixed bilinear bound without dividing
the radius by the number of bond coordinates.

Let \(K_V=Dq_V(0)\) and \(R_V(a)=q_V(a)-K_Va\).
Taylor's formula with integral remainder gives
\begin{equation}
 \begin{aligned}
 \|K_V\|_{\infty\to\infty}&\le CL,\\
 \|R_V(a)\|_\infty&\le CL^2\|a\|_\infty^2,\\
 \|DR_V(a)\|_{\infty\to\infty}&\le CL^2\|a\|_\infty .
 \end{aligned}
 \label{qc:quadratic-residual}
\end{equation}
for \(\|a\|_\infty<c/L\) and \(\epsilon<c/L\).
The estimates transfer to the original frame by \eqref{eq:block-average-pivot-chart-QC10}.
They are uniform over all coarse bonds, hence over the volume.
In addition, varying the background in \eqref{eq:block-average-pivot-chart-QC12} gives
\begin{equation}
 \|Dq_{e^{iv}}(0)-Dq_I(0)\|_{\infty\to\infty}
                                      \le CL^2\|v\|_\infty.
 \label{eq:block-average-pivot-chart-QC14}
\end{equation}
These are estimates for the explicit nonlinear average in \eqref{eq:block-average-pivot-chart-QC3}, not
for a model linear average.

\subsubsection*{The pivot coefficient and finite elimination}

Write \(h=(L-1)/2\).  The pivot of \(c=(y,y+Le_\mu)\) is
\(b_0(c)=(y+he_\mu,y+(h+1)e_\mu)\).
It crosses from \(B(y)\) to \(B(y')\).  No internal contour path
contains it.  Among all coarse pivots, only \(b_0(c)\) occurs
in the support of \(q_{V,c}\).  Indeed internal contours stay
inside one block; each straight segment in \eqref{eq:block-average-pivot-chart-QC3} crosses only
the face between these two blocks and is parallel to \(e_\mu\).
The only coarse pivot on that face is \(b_0(c)\).  This support
assertion, rather than a claim about all bonds joining two
vertices of a periodic double block, is what is used below.

At the identity background define
\(\overline a(y,x)\) to be the arithmetic mean of the oriented
path sums of \(a\) over \(\mathcal G(y,x)\), and
\(\sigma_a(y)=L^{-4}\sum_{x\in B(y)}\overline a(y,x)\).
The first derivative of the contour mean is this arithmetic mean.
Differentiating \eqref{eq:block-average-pivot-chart-QC3}, the negative derivative of \(U[c]\) in
each closed logarithm cancels its final positive derivative.
Thus
\begin{equation}
 (K_Ia)(c)=\sigma_a(y)
       +L^{-4}\sum_{x\in B(y)}a([x,x'])-\sigma_a(y').
 \label{eq:block-average-pivot-chart-QC15}
\end{equation}
The two \(\sigma_a\) terms contain no pivot.  Exactly \(L\)
segments in the remaining sum contain \(b_0(c)\): their transverse
coordinates equal those of \(y\), and their longitudinal starting
coordinate can be any of the \(L\) coordinates in the block.
All occurrences have positive orientation.  Therefore
\begin{equation}
 J_c(I):=D_{a(b_0(c))}q_{I,c}(0)=L^{-3}I_3.
 \label{eq:block-average-pivot-chart-QC16}
\end{equation}
For a pure gauge \(V^s=I\), the coefficient in its original frame is
\begin{equation}
 J_c(V)=L^{-3}\operatorname{Ad}(s(y)^{-1}s(b_0(c)_-)).
 \label{eq:block-average-pivot-chart-QC17}
\end{equation}
The unitary matrix in this formula is the prefix transport from
\(y\) to the start of the pivot.  Replacing it by \(I\) in the
original gauge would be incorrect.

From \eqref{eq:block-average-pivot-chart-QC14}, after the unitary local change of frame, we have
\(\|J_c(V^{s_c})-L^{-3}I_3\|\le CL^2\epsilon\).
Consequently, under the conservative condition
\begin{equation}
 \epsilon\le cL^{-5},
 \qquad\text{in particular if }\quad
       \|E\|_\infty+L\varpi_c\le cL^{-5},
 \label{eq:block-average-pivot-chart-QC18}
\end{equation}
every \(J_c\) is invertible and \(\|J_c^{-1}\|\le2L^3\).
There is no off-diagonal pivot coefficient.  Set
\begin{equation}
 (h_V\nu)(b_0(c))=J_c^{-1}\nu(c),\qquad
 (h_V\nu)(b)=0\ (b\text{ nonpivot}),
 \qquad
 \mathsf C_V=\iota-h_VK_V\iota,
 \label{eq:block-average-pivot-chart-QC19}
\end{equation}
where \(\iota\) extends the nonpivot coordinates by zero.
Then
\begin{equation}
 K_Vh_V=1,\quad K_V\mathsf C_V=0,\quad
 (\mathsf C_Vu)|_{\rm free}=u,
 \qquad
 \|h_V\|\le2L^3,\quad\|\mathsf C_V\|\le CL^4 .
 \label{qc:pivot-linear-maps}
\end{equation}
These are supremum operator norms.  The same powers bound column
norms.  Indeed a fine bond belongs to at most a fixed number of
double blocks, and \eqref{eq:block-average-pivot-chart-QC12} bounds the coefficient in each of their
rows.  The maps are local: a pivot row uses only its own double
block.  They are holomorphic in the background, by the Neumann
series for \(J_c^{-1}\).

For comparison with a convention which inserts \(a/L\) in the
one-step exponential, its linear average is \(Q_V=K_V/L\), and
its pivot coefficient is
\begin{equation}
 \mathsf m_c=J_c/L,\qquad
 \mathsf m_c(I)=L^{-4}I_3,
 \qquad \det J_c=L^3\det\mathsf m_c.
 \label{eq:block-average-pivot-chart-QC21}
\end{equation}
This accounts for the factor \(L^{-3}\) in a real elimination
Jacobian; it is not a change in the number of color coordinates.

\subsubsection*{The nonlinear graph, constructed rather than assumed}

For a full physical bond variable \(t\), solve
\begin{equation}
 d=R_V(t-h_Vd),\qquad \Psi_V(t)=t-h_Vd(t).
 \label{eq:block-average-pivot-chart-QC22}
\end{equation}
On \(\|t\|_\infty<r_1=cL^{-5}\), \eqref{qc:quadratic-residual} maps the ball
\(\|d\|_\infty\le CL^2\|t\|_\infty^2\) into itself, after
choosing its fixed prefactor large enough.  Its Lipschitz constant
in \(d\) is at most \(CL^5r_1<1/2\).  Iteration proves the unique
solution, its holomorphic dependence, and
\begin{equation}
 \begin{aligned}
 q_V(\Psi_V(t))&=K_Vt,&
 \|d(t)\|&\le CL^2\|t\|^2,\\
 \|D_td(t)\|&\le CL^2\|t\|,&
 \|D\Psi_V(t)-1\|&\le CL^5\|t\|.
 \end{aligned}
 \label{eq:block-average-pivot-chart-QC23}
\end{equation}
The derivative estimate follows by differentiating the contraction
and inverting \(1+DR_Vh_V\) by its convergent Neumann series.
Only the pivot values are changed.  There is even an explicit
inverse, on a smaller neighborhood:
\(\Psi_V^{-1}(a)=a+h_VR_V(a)\).  Substitution in \eqref{eq:block-average-pivot-chart-QC22} verifies
this formula by uniqueness.

Each equation for \(d(c)\) uses only its own pivot value and
the nonpivot data in \(D_c\).  The equations for different
coarse bonds are independent once the nonpivot data are fixed.
Thus the graph is exactly finite range, not merely exponentially
local.  In particular
\begin{equation}
 a(u)=\Psi_V(\mathsf C_Vu),\quad q_V(a(u))=0,\quad
 a(u)|_{\rm free}=u,
 \qquad \|u\|_\infty<cL^{-9}
 \label{eq:block-average-pivot-chart-QC24}
\end{equation}
is the actual small fibre chart.  It obeys
\(\|a(u)\|_\infty+\|D_ua(u)\|\|u\|_\infty
 \le CL^4\|u\|_\infty\).
Conversely, any sufficiently small solution of \(q_V(a)=0\)
has this form: \(t=a+h_VR_V(a)\) satisfies \(K_Vt=0\), hence
\(t=\mathsf C_V(a|_{\rm free})\), and \eqref{eq:block-average-pivot-chart-QC23} gives uniqueness.

For the canonical Gaussian coordinate \(u=gy\), the physical
chart therefore contains \(\|y\|_\infty<cL^{-9}/g\).
At weak coupling it contains any radius \(c_*\alpha/g\) with
fixed \(c_*\) and sufficiently small \(\alpha\).
This establishes the one-step primitive radius; inclusion in
all the older background domains remains a different estimate.

\subsubsection*{Real coarea normalization and complex continuation}

We record the normalization only on the real positive chart first.
Let \(\mathcal E,\mathcal C,\mathcal F\) be the fine positive bonds,
coarse positive bonds and nonpivots, respectively.  Thus
\(|\mathcal F|=|\mathcal E|-|\mathcal C|\).
In real Lie algebra coordinates Haar measure is
\(d(e^{ia})=\sigma(a)d^3a\), where
\(\sigma(a)=\sigma_0(\sin|a|/|a|)^2\) and \(\sigma_0>0\).
Group delta in the same chart satisfies
\(\delta_G(e^{iq})=\sigma_0^{-1}\delta_{\mathbb R^3}(q)\).

In the linear variables \(t=\mathsf C_Vu+h_Vv\), free coordinates
are \(u\) and \(K_Vt=v\).  Its determinant is
\(\prod_c\det J_c^{-1}\).  The full straightening \(a=\Psi_V(t)\)
has derivative close to the identity by \eqref{eq:block-average-pivot-chart-QC23}.  At real
backgrounds, the determinants \(\det J_c\) are positive: in the
local unitary frame \(J_c\) is within half of the positive scalar
matrix \(L^{-3}I_3\), and the real adjoint frame transformations
have determinant one.  Also \(\det D\Psi_V>0\), by continuity
from its value one at zero and the uniform derivative bound.
The ordinary real change of variables therefore gives, on the fibre,
\begin{equation}
 \begin{split}
 &\prod_{b\in\mathcal E}dU_b
       \prod_{c\in\mathcal C}
            \delta_G(\mathcal M(U)(c)W(c)^{-1})\\
 &\quad=\sigma_0^{|\mathcal F|}
       \prod_c(\det J_c)^{-1}
       \det D\Psi_V(\mathsf C_Vu)
       \prod_{b\in\mathcal E}\frac{\sigma(a_b(u))}{\sigma_0}
       \prod_{b\in\mathcal F}d^3u_b .
 \end{split}\label{eq:block-average-pivot-chart-QC25}
\end{equation}
This is an identity of local fibre measures, with the group
delta interpreted by this change of coordinates.
The straightening derivative is block triangular: its free block
is the identity, and its pivot diagonal has one three by three
block per coarse bond.  Hence its log determinant is a sum of
local trace logarithms.  Substituting \(u=gy\) in \eqref{eq:block-average-pivot-chart-QC25} supplies
\((\sigma_0g^3)^{|\mathcal F|}\), while
\((\det J_c)^{-1}=L^{-3}(\det\mathsf m_c)^{-1}\).

For complex backgrounds the coefficient functions in \eqref{eq:block-average-pivot-chart-QC25}
have their coherent holomorphic continuations on the chart just
constructed.  In particular \((\det J_c)^{-1}\) and the branch
of \(\log\det J_c\) continuing its real value are used, not an
absolute value of a complex determinant.  The Haar density is
continued by the entire power series
\(\sigma_0[\sin\sqrt z/\sqrt z]^2\),
\(z=\tfrac12\operatorname{tr}a^2\), rather than by a complex
operator norm.  No complex measure or complex delta function is
being asserted by this continuation.

\subsubsection*{Dependence on the remaining scales}

The average, its raw Taylor constant \(CL^2\), and the local
pivot construction contain no \(M\), observation radius or
number of free variables.  If the supplied local near-pure-gauge
bound is only \(\|V^s-I\|\le C_bMa\), \eqref{eq:block-average-pivot-chart-QC18} requires
\(a\le c/(C_bML^5)\).  Alternatively the direct real condition
\(\varpi_c\le cL^{-6}\), with a small complex perturbation, suffices.
The input-to-chart constants are conservative fixed powers of \(L\);
no claim of optimal powers is needed.

If only a complex gauge with condition numbers at most \(K\ge1\)
is available, conjugation charges one factor \(K\) per input
direction and one for the output.  Thus safe bounds are
\(C K^3L^2\) for the quadratic residual,
\(2K^2L^3\) for \(h_V\), and \(CK^4L^4\) for \(\mathsf C_V\).
The conservative physical nonpivot radius becomes
\(cK^{-9}L^{-9}\).  A condition number growing with a larger
window must therefore be charged; it cannot be discarded by
calling the one-step chart local.

The calculation discharges the one-step pivot and Taylor primitives
in these actual physical coordinates.  It does not bound the
many-level response or a further coordinate transformation on a
large run window.  In particular, a polynomial loss in that window
radius in the background/core inclusion must still be retained
in a later Cauchy radius.  The free precision gap, the all-window
core and tower estimates, the conditional exterior forces and the
clean-class law remain separate assertions.


% END MODULE chapters/block_average_pivot_chart_v8.tex

% BEGIN MODULE chapters/flat_contour_precision_gap_v8.tex
\subsection{A gap for the classical flat precision}
\label{ch:flat-contour-precision-gap}

The gauge form in this section is the actual soft contour gauge.
It is not a divergence projector.  At the flat field its quadratic
form is the sum of the squares of the arithmetic internal contour
path sums.  The constraint is the plain derivative of the nonlinear
block average, rather than the straight average alone.  With these
operators the classical flat precision has a gap uniform in the
torus volume and in the number of finer averaging levels.

The statement concerns the unweighted Wilson action at a flat
background, with its prescribed positive soft gauge coefficient.
Signed counterterms, nonconstant layer weights, nonzero-current
constraint corrections and complex decoupling matrices are not
included in the conclusion.

\subsubsection*{The actual quadratic forms and their normalization}

Use raw bond logarithms on a unit-spacing four-dimensional lattice.
For a scalar site function \(f\) and a positively oriented bond
field \(a\), set
\begin{equation}
 \begin{aligned}
 (\mathrm{d}f)_\mu(x)&=f(x+e_\mu)-f(x),\\
 (\mathrm{d}a)_{\mu\nu}(x)&=a_\mu(x)+a_\nu(x+e_\mu)
                         -a_\mu(x+e_\nu)-a_\nu(x).
 \end{aligned}
 \label{eq:flat-contour-precision-gap-CF1}
\end{equation}
The same convention, with step \(n\), is written \(\mathrm{d}_n\) on a
lattice of spacing \(n\).  Norms are counting \(\ell^2\) norms of
sites, positive bonds or plaquettes, with each unordered direction
pair counted once.  One color suffices for the proof; the flat
\(\operatorname{SU}(2)\) forms are three identical copies.

The Wilson convention
\(A(U)=\sum_p[1-\tfrac12\operatorname{Re}\operatorname{tr}U_p]\)
gives
\(A(e^{ita})=\tfrac12t^2\|\mathrm{d}a\|^2+O(t^3)\).
For the actual soft gauge,
\begin{equation}
 \mathrm{GF}(V)=\sum_y\sum_{x\in B(y)\setminus\{y\}}
          [1-\tfrac12\operatorname{Re}\operatorname{tr}V(y,x)].
 \label{eq:flat-contour-precision-gap-CF2}
\end{equation}
Here \(V(y,x)\) is the implicit mean of the transports along the
coordinate-order internal paths.  Its first derivative at the
identity is the arithmetic mean of the corresponding oriented
path sums.  Denote this linear functional by
\(\gamma_{y,x}(a)\), and put \((Ga)_{y,x}=\gamma_{y,x}(a)\).
Consequently
\begin{equation}
 \mathrm{GF}(e^{ita})=\tfrac12t^2\|Ga\|^2+O(t^3).
 \label{eq:flat-contour-precision-gap-CF3}
\end{equation}
Internal paths contain no crossing pivot.  Thus the gauge form
on nonpivot variables is local to each block and has no pivot
coordinate.  For a gradient,
\(\gamma_{y,x}(\mathrm{d}f)=f(x)-f(y)\).

The physical step uses the scale \(g=\theta^{-1/2}\) and gauge
weight \(\theta\).  For the real unweighted step the action
coefficient is also \(\theta\).  Hence \(\theta g^2=1\), and
the two quadratic forms above have coefficient one in the
canonical Gaussian variables.  Normalizing constants and the
Haar Jacobian belong to other parts of the step; they do not
change this definition of its free precision.

\subsubsection*{Plain and straight averaging}

For odd \(n\), let \(B_n(y)\) be the centered block of \(n^4\)
sites.  The raw straight bond average is
\begin{equation}
 (S_na)_\mu(y)=n^{-4}\sum_{x\in B_n(y)}
                         \sum_{t=0}^{n-1}a_\mu(x+te_\mu).
 \label{eq:flat-contour-precision-gap-CF4}
\end{equation}
At one step define
\(\sigma(a)(y)=L^{-4}\sum_{x\in B_L(y)}\gamma_{y,x}(a)\).
Direct differentiation of the actual outer exponential of closed
logarithms gives the plain derivative
\begin{equation}
 P_La=S_La-\mathrm{d}_L\sigma(a).
 \label{eq:flat-contour-precision-gap-CF5}
\end{equation}
The derivative of its reference straight segment cancels its
negative occurrence in every closed logarithm.  This explains
why no extra reference-segment term is present in \eqref{eq:flat-contour-precision-gap-CF5}.

Let \(q_Lf\) be the arithmetic site average over \(B_L(y)\),
and \(r_Lf=f|_{\text{block centers}}\).  Telescoping along paths
proves
\begin{equation}
 S_L\mathrm{d}f=\mathrm{d}_Lq_Lf,\qquad
 \sigma(\mathrm{d}f)=q_Lf-r_Lf,\qquad
 P_L\mathrm{d}f=\mathrm{d}_Lr_Lf.
 \label{eq:flat-contour-precision-gap-CF6}
\end{equation}
In particular it is restriction at the centers, not a site
average, that intertwines with the plain constraint.

Write \(P^{[k]}\) for the derivative at the identity of the
plain \(k\)-fold nonlinear average, with raw logarithms at both
ends.  For \(n=L^k\),
\begin{equation}
 P^{[k]}a=S_na+\mathrm{d}_n\psi_k(a)
 \label{eq:flat-contour-precision-gap-CF7}
\end{equation}
for a linear site functional \(\psi_k\).
Here is a direct verification.  Composition of the straight
averages equals \(S_{nL}\): in each coordinate the centered
digits \(nr+u\), with
\(r\in\{-(L-1)/2,\ldots,(L-1)/2\}\) and
\(u\in\{-(n-1)/2,\ldots,(n-1)/2\}\),
tile the larger centered block, and the segment indices
\(nt+j\), \(0\le t<L\), \(0\le j<n\), tile the longer straight
segment.  All weights multiply
to \((nL)^{-4}\).  Applying the next plain step to the gradient
in \eqref{eq:flat-contour-precision-gap-CF7} gives another coarse gradient by \eqref{eq:flat-contour-precision-gap-CF6}.  Its application
to \(S_na\) differs from the next straight average by a gradient,
by \eqref{eq:flat-contour-precision-gap-CF5}.  This proves \eqref{eq:flat-contour-precision-gap-CF7} inductively, starting with \eqref{eq:flat-contour-precision-gap-CF5}.

\subsubsection*{Curvature averaging is a contraction in dimension four}

Expanding the boundary of an \(n\) by \(n\) square gives
\begin{equation}
 (\mathrm{d}_nS_na)_{\mu\nu}(y)
    =n^{-4}\sum_{x\in B_n(y)}\sum_{r,s=0}^{n-1}
             (\mathrm{d}a)_{\mu\nu}(x+re_\mu+se_\nu).
 \label{eq:flat-contour-precision-gap-CF8}
\end{equation}
Every interior bond cancels in this expansion; its four remaining
sides are exactly the four straight averages in \(\mathrm{d}_nS_n\).
Call the positive averaging operator on plaquettes on the right
\(T_n\).  Each row has sum \(n^2\).  Each fine plaquette has
total column weight \(n^{-2}\): for each pair \((r,s)\) its
starting point determines exactly one block \(B_n(y)\), since
those blocks partition the fine sites.  There are \(n^2\) such
pairs, each with weight \(n^{-4}\).

For nonnegative weights \(t_{Pp}\), weighted Cauchy gives
\begin{equation}
 \sum_P\left|\sum_p t_{Pp}F_p\right|^2
 \le\sum_P\left(\sum_pt_{Pp}\right)
                   \left(\sum_pt_{Pp}|F_p|^2\right)
 \le n^2n^{-2}\sum_p|F_p|^2.
 \label{eq:flat-contour-precision-gap-CF9}
\end{equation}
Thus \(\|T_n\|_{2\to2}\le1\).  Counting with multiplicities
makes the same proof valid on a periodic lattice.  The product
of row and column sums in dimension \(d\) would be \(n^{4-d}\);
its cancellation here uses dimension four.
Taking curls of \eqref{eq:flat-contour-precision-gap-CF7} therefore proves the exact uniform bound
\begin{equation}
 \|\mathrm{d}_nP^{[k]}w\|\le\|\mathrm{d}w\|.
 \label{cf:curvature-contraction}
\end{equation}
No bound on a background propagator was used.

\subsubsection*{The classical constrained pullback}

For current bond data \(a\), on the level-\(k\) lattice, define
the classical flat quadratic value
\begin{equation}
 E_k(a)=\min\{\|\mathrm{d}w\|^2:P^{[k]}w=a\}.
 \label{eq:flat-contour-precision-gap-CF11}
\end{equation}
The affine set is nonempty.  For instance place the value of a
coarse bond on every fine bond crossing its block face and put
zero on fine bonds internal to a block of side \(n=L^k\).
The nonlinear group lift with internal bonds \(I\) has exactly
the prescribed average: all parallel segments cross the same
face value and all closed logarithms vanish.  Applying this
observation successively through the centered subdivisions gives
the same assertion at \(k\) levels.  Its derivative is a right
inverse for \(P^{[k]}\).

Existence of the minimum in \eqref{eq:flat-contour-precision-gap-CF11} needs no uniform inverse
theorem.  The image under \(d\) of the finite-dimensional affine
constraint set is a closed affine subspace; its point nearest
zero exists.  Equivalently, quotient
\(\ker P^{[k]}\cap\ker \mathrm{d}\), restrict to its orthogonal complement,
and solve the positive finite-dimensional normal equations there.
This also shows that the minimum is a quadratic form in \(a\).
Its construction at this stage has no claimed uniform inverse bound.
\eqref{cf:curvature-contraction}, applied to every admissible lift, gives
\begin{equation}
 E_k(a)\ge\|\mathrm{d}_na\|^2.
 \label{eq:flat-contour-precision-gap-CF12}
\end{equation}
We henceforth express the current lattice in units of its own
spacing, so the curl on the right is written \(\mathrm{d}a\).

This is the classical action part of the actual flat precision.
Indeed the first variation of the bare action at the identity is
zero.  A differentiable exact background lift has a tangent \(w\)
with \(P^{[k]}w=a\); its action Hessian is \(\|\mathrm{d}w\|^2\), because
its second derivative is multiplied by that zero first variation.
Linearized constrained stationarity makes this tangent minimize
the quadratic form in \eqref{eq:flat-contour-precision-gap-CF11}, modulo its gauge null space.
The same observation eliminates the second-order one-step
constraint correction at a flat point.  In particular, no
nonzero-current correction has been dropped in establishing
\eqref{eq:flat-contour-precision-gap-CF11} at the flat point.  The finite Gaussian form in \eqref{eq:flat-contour-precision-gap-CF11} is
defined without assuming a nonlinear background theorem; when
that background is used, this argument identifies its flat
quadratic pullback.  Even an arbitrary exact differentiable lift
would suffice for the lower bound in \eqref{eq:flat-contour-precision-gap-CF12}.

Let \(C\) be the actual flat pivot inclusion, with identity on
nonpivot coordinates.  The physical flat free precision is
characterized by
\begin{equation}
 \langle u,A_k^{\rm flat}u\rangle
       =E_k(Cu)+\|GCu\|^2
       \ge\|\mathrm{d}Cu\|^2+\|GCu\|^2.
 \label{eq:flat-contour-precision-gap-CF13}
\end{equation}
This formula is for the actual contour gauge and the actual
plain fibre.  It is not obtained by substituting a Landau
projection for \(G\).

\subsubsection*{A fixed-cell Fourier calculation}

Group the current lattice into blocks of side \(L\).  In one
cell let \(R=\{-(L-1)/2,\ldots,(L-1)/2\}^4\).
A bond is indexed by its starting site \(r\in R\) and direction.
There are \(4L^4\) bond coordinates and four pivot coordinates.
The free space consequently has dimension \(4L^4-4\), independently
of the volume.  For a coarse Bloch parameter
\(p\in(\mathbb R/2\pi\mathbb Z)^4\), extend a cell field by
\begin{equation}
 a(Lm+r,\mu)=e^{ip\cdot m}a(r,\mu).
 \label{eq:flat-contour-precision-gap-CF14}
\end{equation}
The curl matrix \(\mathrm{d}_p\) has entries which are constants or
\(e^{\pm ip_\mu}\).  The internal-contour matrix \(G\) uses
only edges within its cell, hence has no \(p\) dependence.
The plain constraint matrix \(P(p)\), obtained from \eqref{eq:flat-contour-precision-gap-CF5}, is
also a finite matrix with such phase entries.

The middle pivot of direction \(\mu\) is the crossing bond
from \(he_\mu\) to \((h+1)e_\mu\), \(h=(L-1)/2\).
At the identity its coefficient in the raw constraint is
\(L^{-3}\), since exactly \(L\) of the \(L^4\) parallel
segments contain it.  There is no other pivot in that constraint.
It is the own pivot in the cell of the coarse bond's start,
so this coefficient has no Bloch phase.  Thus the pivot block
of \(P(p)\) is \(L^{-3}I_4\) for every \(p\).
If \(h_0\) inserts \(L^3\) on those pivots and \(\iota\)
inserts free coordinates with zero pivots, then
\begin{equation}
 C(p)=\iota-h_0P(p)\iota,
 \quad P(p)C(p)=0,\quad C(p)|_{\rm free}=1,
 \quad \operatorname{Ran}C(p)=\ker P(p).
 \label{eq:flat-contour-precision-gap-CF15}
\end{equation}
The last equality follows by solving each pivot equation.
It also proves the rank \(4L^4-4\) for all \(p\).
All entries of \(C(p)\) are continuous Laurent polynomials in
the Bloch phases.  This is the same finite-range pivot map as
in the physical chart, evaluated at the flat field.

Define the Hermitian finite matrix
\begin{equation}
 B(p)=C(p)^*\big(\mathrm{d}_p^*\mathrm{d}_p+G^*G\big)C(p).
 \label{eq:flat-contour-precision-gap-CF16}
\end{equation}
We prove that it has no zero vector for every \(p\), including
\(p=0\).  Suppose \(a=C(p)u\) and
\(\mathrm{d}_pa=0\), \(Ga=0\).  Extend it by \eqref{eq:flat-contour-precision-gap-CF14}.
Within a block, two coordinate-order paths to the same site
differ by plaquette swaps.  Zero curl makes their sums equal.
The same swaps express a comb loop with one additional internal
bond as a sum of internal plaquette curls.  Consequently zero
\(G\) forces every internal path sum, and then every internal
bond, to vanish.

Across a face perpendicular to \(e_\mu\), a plaquette with
one transverse step compares two neighboring crossing values.
Its other two bonds are internal to the two blocks, hence zero.
The crossing values are therefore equal.  The transverse face
grid is connected, so all crossing bonds on that face have a
common value \(\beta_{y,\mu}\).  Every straight segment from
\(x\in B(y)\) to \(x+Le_\mu\) then has sum
\(\beta_{y,\mu}\), and all internal contour sums are zero.
The actual constraint gives
\((P_La)(y,\mu)=\beta_{y,\mu}\).
It vanishes by \eqref{eq:flat-contour-precision-gap-CF15}, killing every crossing bond as well.
Thus \(a=0\), and \(u=0\) by its free restriction.
The argument is valid for every real Bloch parameter, with complex
fiber vectors.

One can also describe the possible closed fields explicitly.
On the cell, the \(L^4\) functions
\(L^{-2}\exp(i\xi_\ell\cdot r)\),
\(\xi_\ell=(p+2\pi\ell)/L\), form an orthonormal basis,
where each component of \(\ell\) runs from \(0\) to \(L-1\).
Orthogonality is the finite geometric-series identity.
Fine differences multiply a mode by
\(t_\mu=e^{i\xi_{\ell,\mu}}-1\).
The closedness equations are
\(t_\mu a_\nu-t_\nu a_\mu=0\).  If one \(t_\mu\ne0\),
they say exactly \(a=t\phi\), a gradient.
The only remaining case is \(p=0,\ell=0\), giving a constant
harmonic vector.  The local argument above kills that case
as well; no zero-frequency mode has been silently omitted.

The entries of \(B(p)\) are continuous on the compact Bloch
torus.  Its least Rayleigh quotient therefore has a positive
uniform minimum \(c(L)>0\).  Explicitly, if the infimum were
zero, a sequence of parameters and unit free vectors would
have a convergent subsequence.  Their limit would be a unit
zero vector for the limiting \(B(p)\), contradicting the
previous paragraph.  Thus
\begin{equation}
 B(p)\ge c(L)1\qquad\hbox{for every }p.
 \label{eq:flat-contour-precision-gap-CF17}
\end{equation}

On a torus with \(N_\mu\) coarse blocks in each direction,
the allowed parameters are \(p_\mu=2\pi j_\mu/N_\mu\).
The normalized discrete Fourier transform in the block index
is unitary, again by the same geometric-series identity.
It converts the comparison form on the right of \eqref{eq:flat-contour-precision-gap-CF13} into
the sum of \eqref{eq:flat-contour-precision-gap-CF16} over these parameters.  Parseval and \eqref{eq:flat-contour-precision-gap-CF17}
therefore prove
\begin{equation}
 A_k^{\rm flat}\ge c(L)1
 \label{cf:classical-flat-gap}
\end{equation}
on the actual free nonpivot coordinates, uniformly in the torus
volume and in \(k\).  Local formulas can first be used on the
admissible large tori; no small-torus identification of different
pivots is needed.  The three color copies have the same constant.

\subsubsection*{Flat holonomies and the limit of the conclusion}

The same conclusion holds for a real flat background in its
prescribed block axial representative.  Flatness makes internal
transports path independent.  Since their means in this
representative are \(I\), every internal bond is \(I\).
Flat plaquettes make the crossing value constant on each block
face.  The background is therefore the lift of a flat coarse
connection.  A gauge constant on each block flattens it on
the universal cover and preserves the soft contour gauge norm.
Only twisted periodicity remains.

The four torus holonomies commute: their commutator is the
product of the plaquettes of a fundamental rectangle, all equal
to \(I\).  Commuting unitary matrices have a common orthonormal
eigenbasis; choose an eigenspace of any noncentral holonomy and
use its invariance under the others, or any basis if all are
central.  Their adjoint actions are consequently commuting
unitaries on the complexified color space.  Decomposing into
their common eigenspaces gives characters \(e^{i\chi_\mu}\).
The allowed block momenta in that character are
\((2\pi j_\mu+\chi_\mu)/N_\mu\), up to the sign convention
for the boundary twist.  The local matrices and the curvature
averaging proof are unchanged: the boundary phases have modulus
one, so the absolute row and column counts in \eqref{eq:flat-contour-precision-gap-CF9} still apply.
Since \eqref{eq:flat-contour-precision-gap-CF17} holds for every
real Bloch parameter, all these shifted sets satisfy the same
gap.  No smallness of global holonomies is required.

The result proves the classical flat gap for the actual contour
gauge and plain constraint.  It is a lower bound, and does not
prove a uniform upper bound or a decay estimate for a dressed
precision.  Nonzero-current terms from curved constraints,
signed layer profiles, counterterm arrays, non-Abelian
perturbations and enlarged complex decoupling parameters still
need their own quantitative comparison with this positive
flat form.  No full background or interacting Gaussian-law
assertion follows from \eqref{cf:classical-flat-gap} alone.


% END MODULE chapters/flat_contour_precision_gap_v8.tex

% BEGIN MODULE chapters/flat_whitney_comparison_v8.tex
\subsection{An upper bound for the classical flat action}
\label{ch:flat-whitney-comparison}

This section concerns the same unweighted classical Gaussian value
as the flat-gap calculation.  Let \(n=L^k\), with \(L\) odd, and
let \(P^{[k]}\) be the raw derivative of the actual plain iterated
average.  For a coarse bond field \(a\), put
\begin{equation}
 E_k(a)=\min\{\|\mathrm{d}w\|^2:P^{[k]}w=a\}.
 \label{eq:flat-whitney-comparison-WI1}
\end{equation}
We construct an exact lift with energy at most
\(256\|\mathrm{d}_na\|^2\), uniformly in \(k\) and in the volume.
No bound on a non-Abelian background response is used.

\subsubsection*{Discrete interpolation on the actual centered grids}

Fine sites are integers and coarse sites are \(nm\),
\(m\in\mathbb Z^4\).  Write
\begin{equation}
 \varphi_n(t)=\max\{1-|t|/n,0\},\qquad
 \beta_n(t)=\mathbf1_{\{0\le t<n\}},\quad t\in\mathbb Z.
 \label{eq:flat-whitney-comparison-WI2}
\end{equation}
Periodize these functions on a finite torus.  The identities
\(\sum_m\varphi_n(x-nm)=1\),
\(\sum_m\beta_n(x-nm)=1\), and
\(\sum_{x\in\mathbb Z}\varphi_n(x)=n\) follow by summing the
two linear pieces.  Also
\begin{equation}
 \varphi_n(t+1)-\varphi_n(t)
       =n^{-1}\big[\beta_n(t+n)-\beta_n(t)\big].
 \label{eq:flat-whitney-comparison-WI3}
\end{equation}
Define interpolation of sites, bonds and plaquettes by
\begin{equation}
 \begin{aligned}
 (I_0f)(x)&=\sum_m f(m)\prod_\lambda\varphi_n(x_\lambda-nm_\lambda),\\
 (I_1a)_\mu(x)&=n^{-1}\sum_m a_\mu(m)
       \beta_n(x_\mu-nm_\mu)
       \prod_{\lambda\ne\mu}\varphi_n(x_\lambda-nm_\lambda),\\
 (I_2F)_{\mu\nu}(x)&=n^{-2}\sum_m F_{\mu\nu}(m)
       \beta_n(x_\mu-nm_\mu)\beta_n(x_\nu-nm_\nu)\\
   &\qquad\times\prod_{\lambda\notin\{\mu,\nu\}}
                          \varphi_n(x_\lambda-nm_\lambda).
 \end{aligned}\label{eq:flat-whitney-comparison-WI4}
\end{equation}
The coarse arrays in these formulas are indexed by their starting
node \(m\).  Their values are raw bond or plaquette values; no
length factor has been included in the coarse difference \(\mathrm{d}_n\).
\eqref{eq:flat-whitney-comparison-WI3}, followed by reindexing \(m\), proves
\begin{equation}
 \mathrm{d}I_0=I_1\mathrm{d}_n,\qquad \mathrm{d}I_1=I_2\mathrm{d}_n.
 \label{eq:flat-whitney-comparison-WI5}
\end{equation}
For example, differentiating a transverse nodal hat in \(I_1\)
replaces it by the difference of two interval indicators divided
by \(n\); their coefficients are precisely the two terms of
the corresponding coarse plaquette boundary.  This proves the
second identity with the same orientation as
\((\mathrm{d}a)_{\mu\nu}=a_\mu+\nabla_\mu a_\nu-a_\mu(\cdot+e_\nu)\).

For a fixed plaquette direction, \(I_2\) has nonnegative weights.
Their row sum is \(n^{-2}\), by the two partition identities.
Their column sum is \(n^2\): each interval contributes \(n\)
fine positions, each of the remaining two nodal hats has sum
\(n\), and the prefactor is \(n^{-2}\).
Weighted Cauchy therefore gives
\begin{equation}
 \|I_2F\|^2\le n^{-2}n^2\|F\|^2=\|F\|^2.
 \label{eq:flat-whitney-comparison-WI6}
\end{equation}
These counting norms are those of the four-dimensional Wilson
quadratic energy.  As for curvature averaging, dimension four
is used in the cancellation of the powers of \(n\).

\subsubsection*{The exact averaging symbol}

With centered blocks, the raw straight average is
\[
(S_nw)_\mu(y)=n^{-4}\sum_{x\in B_n(y)}
\sum_{j=0}^{n-1}w_\mu(x+je_\mu).
\]
We compute \(S_nI_1\) exactly to obtain a right inverse.  Put \(h=(n-1)/2\).  In a transverse direction
the coefficient of the neighboring coarse index \(m\) is
\(n^{-1}\sum_{x=-h}^h\varphi_n(x-nm)\).
In the component direction it is
\(n^{-2}\#\{(x,j):-h\le x\le h,\ 0\le j<n,
\ nm\le x+j<n(m+1)\}\).
Both expressions are zero unless \(m=0,\pm1\), and both have
the same three values:
\begin{equation}
 a_0=1-\frac{h(h+1)}{n^2}=\frac34+\frac1{4n^2},
 \qquad
 a_1=\frac{h(h+1)}{2n^2}=\frac{1-n^{-2}}8
       \quad(m=\pm1).
 \label{eq:flat-whitney-comparison-WI7}
\end{equation}
For the longitudinal count, the negative and positive overshoots
each number \(1+\cdots+h=h(h+1)/2\); the remaining
\(n^2-h(h+1)\) pairs give the central coefficient.
For the transverse hat, its positive-side and negative-side
linear sums give the same two triangular numbers.

Let \(T_\lambda\) translate a coarse array by one node in
direction \(\lambda\).  The tensor product in \eqref{eq:flat-whitney-comparison-WI4} and the
identical one-dimensional coefficients just computed show
\begin{equation}
 S_nI_1=\mathcal A_n,
 \qquad \mathcal A_n=\prod_{\lambda=1}^4
       \big[a_0I+a_1(T_\lambda+T_\lambda^{-1})\big].
 \label{eq:flat-whitney-comparison-WI8}
\end{equation}
Crucially, this is the same scalar convolution on all four
bond components.  Each one-dimensional symbol is
\begin{equation}
 a_n(t)=\frac34+\frac1{4n^2}
                   +\frac{1-n^{-2}}4\cos t,
 \qquad \frac12\le a_n(t)\le1.
 \label{eq:flat-whitney-comparison-WI9}
\end{equation}
Fourier transform on the coarse torus, or the quadratic-form
bound \(a_0-2a_1\ge1/2\) for a unitary shift, proves
\begin{equation}
 \|\mathcal A_n^{-1}\|_{2\to2}\le16.
 \label{eq:flat-whitney-comparison-WI10}
\end{equation}
All shifts commute.  Their scalar convolution commutes with
every coarse difference, on sites, bonds and plaquettes.
Its inverse has the same property.

\subsubsection*{An exact lift for the plain constraint}

For prescribed coarse \(a\), take
\(w_0=I_1\mathcal A_n^{-1}a\).
\eqref{eq:flat-whitney-comparison-WI5}, \eqref{eq:flat-whitney-comparison-WI6}, \eqref{eq:flat-whitney-comparison-WI8} and \eqref{eq:flat-whitney-comparison-WI10} imply
\begin{equation}
 S_nw_0=a,\qquad
 \|\mathrm{d}w_0\|\le\|\mathrm{d}_n\mathcal A_n^{-1}a\|
                       \le16\|\mathrm{d}_na\|.
 \label{eq:flat-whitney-comparison-WI11}
\end{equation}
This is an exact straight lift.  The actual constraint is plain.
For completeness its two relevant identities are
\begin{equation}
 P^{[k]}w=S_nw+\mathrm{d}_n\psi_k(w),\qquad
 P^{[k]}\mathrm{d}f=\mathrm{d}_n(r_nf),
 \label{eq:flat-whitney-comparison-WI12}
\end{equation}
where \(r_n\) restricts a fine site function to the centers
\(n\mathbb Z^4\).  At one step these follow by differentiating
the actual closed-logarithm formula: \(P_L=S_L-\mathrm{d}_L\sigma\),
and on gradients an internal contour is
\(f(x)-f(y)\), so \(P_L\mathrm{d}f=\mathrm{d}_Lr_Lf\).
Iterating proves the second identity.  Composition of straight
averages is \(S_n\); applying a plain average to a gradient
again gives a gradient, proving the first by induction.

Nodal interpolation satisfies \(r_nI_0f=f\).  Set
\begin{equation}
 f=-\psi_k(w_0),\qquad w=w_0+\mathrm{d}I_0f.
 \label{eq:flat-whitney-comparison-WI13}
\end{equation}
Then \eqref{eq:flat-whitney-comparison-WI11}--\eqref{eq:flat-whitney-comparison-WI12} give \(P^{[k]}w=a\) exactly, while
\(\mathrm{d}w=\mathrm{d}w_0\).  The gradient correction has no classical
curvature cost.  The energy estimate itself needs no amplitude
bound; such a bound is proved below.  We have proved
\begin{equation}
 E_k(a)\le256\|\mathrm{d}_na\|^2.
 \label{eq:flat-whitney-comparison-WI14}
\end{equation}
The complementary lower bound follows from the exact curvature
average: \(\mathrm{d}_nP^{[k]}w=T_n\mathrm{d}w\), where \(T_n\) has row sum
\(n^2\) and total column weight \(n^{-2}\).  The same weighted
Cauchy argument gives \(\|T_n\|\le1\).  Thus
\begin{equation}
 \boxed{\ \|\mathrm{d}_na\|^2\le E_k(a)\le256\|\mathrm{d}_na\|^2\ }.
 \label{wi:energy-comparison}
\end{equation}
All statements include \(n=1\), where the interpolations and
\(\mathcal A_1\) are identities.  No zero-frequency inverse
is taken: \(\mathcal A_n\) itself is strictly positive even
on constants.

\subsubsection*{Consequences and scope}

For the actual flat free nonpivot precision,
\(\langle u,A_k^{\rm flat}u\rangle
=E_k(Cu)+\|GCu\|^2\).
The local pivot basis has counting \(\ell^2\) norm at most
\(CL^4\), from its row and column bounds.  Also
\(\|\mathrm{d}\|^2\le24\) on one-forms in four dimensions.
The internal contour matrix has row sum at most \(CL\)
and column sum at most \(CL^4\), hence \(\|G\|^2\le CL^5\).
It contains no pivot, so \(GC=G\iota\), where \(\iota\)
inserts free coordinates.  \eqref{eq:flat-whitney-comparison-WI14} and these bounds
give the conservative uniform upper bound
\begin{equation}
 A_k^{\rm flat}\le CL^8 I.
 \label{eq:flat-whitney-comparison-WI16}
\end{equation}
Together with the separate flat-gap theorem this bounds the
same classical precision above and below independently of
the number of finer levels and the volume.

\subsubsection*{Supremum norms of the constructed linear lift}

The same scalar inverse is bounded in the supremum norm.
For one direction the off-diagonal part of
\(a_0I+a_1(T+T^{-1})\), divided by \(a_0\), has supremum
operator norm at most \(2a_1/a_0<1\).  Its inverse Neumann
series therefore has norm at most
\((a_0-2a_1)^{-1}\le2\).  Multiplying the four commuting
direction inverses gives
\begin{equation}
 \|\mathcal A_n^{-1}\|_{\infty\to\infty}\le16,
 \qquad \|w_0\|_\infty\le\frac{16}{n}\|a\|_\infty .
 \label{eq:flat-whitney-comparison-WI17}
\end{equation}
The second inequality uses the row sum \(1/n\) of \(I_1\).

There is an explicit choice of the scalar in \eqref{eq:flat-whitney-comparison-WI12}.
Let \(\psi_0=0\).  When the current spacing is \(L^j\),
apply \(P_L=S_L-\mathrm{d}_L\sigma_L\) to
\(P^{[j]}w_0=S_{L^j}w_0+\mathrm{d}_{L^j}\psi_j\).
The next plain step sends the last gradient to the gradient
of the center restriction.  Thus
\begin{equation}
 \psi_{j+1}=r_L\psi_j-\sigma_L(S_{L^j}w_0).
 \label{eq:flat-whitney-comparison-WI18}
\end{equation}
All fields in this recurrence are raw bond or site values
on their current lattices.  There is no missing spacing
factor.  Internal paths have at most \(2(L-1)\) steps,
so \(\|\sigma_L\|_{\infty\to\infty}\le2(L-1)\).
Also \(\|S_{L^j}\|_{\infty\to\infty}=L^j\), whereas
restriction has norm one.  The geometric sum yields
\(\|\psi_k\|_\infty
 \le2(L-1)\sum_{j<k}L^j\|w_0\|_\infty
 =2(n-1)\|w_0\|_\infty\le32\|a\|_\infty\).

Since \(\mathrm{d}I_0f=I_1\mathrm{d}_nf\), the correcting gradient in \eqref{eq:flat-whitney-comparison-WI13}
has norm at most \(2\|f\|_\infty/n\).  Consequently the
particular exact lift constructed above obeys
\begin{equation}
 \|\mathrm{d}I_0f\|_\infty\le\frac{64}{n}\|a\|_\infty,
 \qquad \|w\|_\infty\le\frac{80}{n}\|a\|_\infty .
 \label{wi:linear-lift-sup}
\end{equation}
The scalar convolution and its inverse commute with curl.
The row sum \(n^{-2}\) of \(I_2\) and \eqref{eq:flat-whitney-comparison-WI17} also give
\begin{equation}
 \|\mathrm{d}w\|_\infty=\|\mathrm{d}w_0\|_\infty
              \le\frac{16}{n^2}\|\mathrm{d}_na\|_\infty.
 \label{eq:flat-whitney-comparison-WI20}
\end{equation}
These are norms of one explicitly defined linear right
inverse.  They are not estimates for a minimizing response.

The interpolation identities also hold with commuting unitary
boundary twists.  Work on the universal cover and periodize
with the corresponding character; absolute row and column
sums are unchanged.  Twisted shifts remain unitary and
commuting, so \eqref{eq:flat-whitney-comparison-WI9}--\eqref{eq:flat-whitney-comparison-WI10} retain their lower bound for every
shifted Bloch parameter.  Thus the same energy comparison
applies to real flat holonomies in the block axial gauge.

The construction gives energy and supremum estimates for
a particular linear lift.  It does not bound a minimizing
or nonlinear background map, or its complex core and tower
substitutions.  It does not prove exponential decay of the
effective Hessian.  Nor does it
control signed layer profiles, nonzero-current corrections,
counterterm perturbations or complex decoupling arrays.
Those estimates must compare their actual operators with
this unweighted classical flat form.


% END MODULE chapters/flat_whitney_comparison_v8.tex

% BEGIN MODULE chapters/flat_exponential_kernels_v8.tex
\subsection{Exponential kernels of the classical flat precision}
\label{ch:flat-exponential-kernels}

We compute the classical flat quadratic form of the actual plain
iterated average, and then the covariance on its next-step nonpivot
fiber with the actual soft contour gauge.  A spectral gap by itself
does not imply kernel decay.  The proof below supplies a uniform
complex strip for the exact symbols before using the flat gap.
All fields are in raw bond coordinates, indexed by their starting
sites.  No midpoint phase is suppressed.

Let \(n=L^k\), with \(L\) odd.  The current lattice is in units
of its own spacing; its finer lattice has \(n\) subdivisions per
direction.  We first use one scalar color.  At the identity the
three \(\operatorname{SU}(2)\) colors give identical copies.  The
normalization of the action is
\(A(e^{ita})=\tfrac12t^2\|\mathrm{d}a\|^2+O(t^3)\).
After the canonical scaling \(g=\theta^{-1/2}\), the real
unweighted step has \(\theta g^2=1\).  Thus the matrices calculated
here are precisely its classical quadratic precision matrices.

\subsubsection*{The two constrained minima are equal}

Write \(P^{[k]}\) for the raw derivative of the actual plain
\(k\)-fold average.  The straight average is
\begin{equation}
 (S_nw)_\mu(ny)=n^{-4}\sum_{r\in B_n(0)}
                         \sum_{j=0}^{n-1}w_\mu(ny+r+je_\mu).
 \label{eq:flat-exponential-kernels-FX1}
\end{equation}
Here \(B_n(0)=\{-(n-1)/2,\ldots,(n-1)/2\}^4\).
The actual one-step map is that of
\cref{ch:block-average-pivot-chart}.  Let
\(\mathcal G_L(y,x)\) be its coordinate-order shortest-path
multiset from \(y\) to \(x\in B_L(y)\), with its prescribed
multiplicities and equal weights.  In positive bond coordinates set
\begin{equation}
 \begin{aligned}
 \gamma_{y,x}(a)&=\frac1{\#\mathcal G_L(y,x)}
       \sum_{\Gamma\in\mathcal G_L(y,x)}
                \sum_{b\in\Gamma}\varepsilon_\Gamma(b)a(b),\\
 (Ga)_{y,x}&=\gamma_{y,x}(a)\quad(x\ne y),\qquad
 \sigma(a)(y)=L^{-4}\sum_{x\in B_L(y)}\gamma_{y,x}(a).
 \end{aligned}\label{eq:flat-exponential-kernels-FX1a}
\end{equation}
Here the signs record the path orientation and
\(\gamma_{y,y}=0\).  Differentiation gives
\(P_L=S_L-\mathrm{d}_L\sigma\).  On gradients the plain map intertwines
with restriction at the centers, whereas the straight map
intertwines with the site average.  Iteration gives
\begin{equation}
 P^{[k]}w=S_nw+\mathrm{d}_n\psi_k(w),\qquad
 P^{[k]}\mathrm{d}f=\mathrm{d}_n r_nf,\qquad S_n\mathrm{d}f=\mathrm{d}_nq_nf.
 \label{eq:flat-exponential-kernels-FX2}
\end{equation}
Both site maps are onto: for \(r_n\) assign the values at the
centers, and for \(q_n\) use a function constant on each centered
block.  These are statements about finite linear cochains; no
regularity bound for those corrections is required.

Suppose \(P^{[k]}w=a\).  Then \(S_nw=a-\mathrm{d}_n\psi_k(w)\).
Choose \(q_nf=\psi_k(w)\).  The field \(w+\mathrm{d}f\) has straight
average \(a\) and the same curl.  Conversely, for \(S_nw=a\),
choose \(r_nf=-\psi_k(w)\) to correct the plain average.
Consequently the actual classical quadratic value is
\begin{equation}
 E_n(a):=\min_{P^{[k]}w=a}\|\mathrm{d}w\|^2
          =\min_{S_nw=a}\|\mathrm{d}w\|^2.
 \label{eq:flat-exponential-kernels-FX3}
\end{equation}
Nonemptiness also follows from the filter calculation below:
its central alias is nonzero in every component at \(p\ne0\),
and a constant fine field lifts a constant coarse field.
The two gradient corrections transfer feasibility between the
constraints.  Finite dimensional projection after quotienting
the curl-null directions then gives both minima.
The equality does not change the actual soft gauge, which is
added later on the current one-step fiber.

\subsubsection*{Fourier normalization and the auxiliary covariance}

Take a coarse torus with \(\mathsf s\) sites per direction and
a fine torus with \(n\mathsf s\).  Use the unitary transforms
\begin{equation}
 \widehat w(\xi)=(n\mathsf s)^{-2}\sum_xe^{-i\xi\cdot x}w(x),
 \qquad
 \widehat a(p)=\mathsf s^{-2}\sum_ye^{-ip\cdot y}a(ny).
 \label{eq:flat-exponential-kernels-FX4}
\end{equation}
Fine aliases of a coarse momentum \(p\) are
\(\xi_m=(p+2\pi m)/n\), with
\(m\in I_n^4\), \(I_n=\{-(n-1)/2,\ldots,(n-1)/2\}\).
Put
\begin{equation}
 b_n(z)=\frac{\sin(nz/2)}{n\sin(z/2)},\qquad
 B_n(\xi)=\prod_\nu b_n(\xi_\nu),\qquad h=(n-1)/2.
 \label{eq:flat-exponential-kernels-FX5}
\end{equation}
The centered block sum is \(B_n\), and the straight segment
sum is \(ne^{ih\xi_\mu}b_n(\xi_\mu)\).  Thus the exact filter
of component \(\mu\) in \eqref{eq:flat-exponential-kernels-FX1} is
\begin{equation}
 s_{\mu,n}(\xi)=ne^{ih\xi_\mu}B_n(\xi)b_n(\xi_\mu),
 \qquad
 \widehat a_\mu(p)=n^{-2}\sum_{m\in I_n^4}
                      s_{\mu,n}(\xi_m)\widehat w_\mu(\xi_m).
 \label{eq:flat-exponential-kernels-FX6}
\end{equation}
The factor \(n^{-2}\) is the square root of the ratio of the
coarse and fine numbers of sites.  In particular its square,
not an unnormalized alias sum, enters a covariance.

The auxiliary fine Feynman form is
\(\|\mathrm{d}w\|^2+\|\mathrm{d}^*w\|^2\).  Its symbol is component-diagonal,
with
\(\lambda(\xi)=4\sum_\nu\sin^2(\xi_\nu/2)\).
For \(p\ne0\), every alias has positive \(\lambda\).
The coarse covariance of this auxiliary form is diagonal:
\begin{equation}
 D_{\mu,n}(p)=n^{-4}\sum_m
                  \frac{|s_{\mu,n}(\xi_m)|^2}{\lambda(\xi_m)}.
 \label{eq:flat-exponential-kernels-FX7}
\end{equation}
This formula can equally be obtained by minimizing the positive
fine Feynman form under its linear constraint; no probability
measure on the fine constant mode is needed.

For complex continuation define
\begin{equation}
 \begin{aligned}
 q_{\nu,m}(p)&=
 \frac{\sin((p_\nu+2\pi m_\nu)/2)}
      {n\sin((p_\nu+2\pi m_\nu)/(2n))},\\
 f_{\mu,m}(p)&=\left(\prod_\nu q_{\nu,m}(p)^2\right)
                                              q_{\mu,m}(p)^2,\\
 s_{n,m}(p)&=4n^2\sum_\nu
                    \sin^2((p_\nu+2\pi m_\nu)/(2n)).
 \end{aligned}\label{eq:flat-exponential-kernels-FX8}
\end{equation}
Removable ratios are filled in at zero.  On real momenta the
phase \(e^{ih\xi_\mu}\) in \eqref{eq:flat-exponential-kernels-FX6} cancels with its conjugate.
Its analytic version is the product of the filters at \(\xi\)
and \(-\xi\), giving exactly
\begin{equation}
 D_{\mu,n}(p)=\sum_{m\in I_n^4}
                           \frac{f_{\mu,m}(p)}{s_{n,m}(p)}.
 \label{fx:alias-covariance}
\end{equation}
Indeed \(s_{\mu,n}(\xi)s_{\mu,n}(-\xi)
=n^2B_n(\xi)^2b_n(\xi_\mu)^2\), and
\(s_{n,m}=n^2\lambda(\xi_m)\).  This checks every power of
\(n\).  Squared absolute values in \eqref{eq:flat-exponential-kernels-FX7} are not used as
holomorphic functions.

\subsubsection*{Recovering the gauge-invariant energy}

For real \(p\), let
\(t_\mu^+(p)=e^{ip_\mu}-1\) and
\(t_\mu^-(p)=e^{-ip_\mu}-1\).
The latter is the conjugate of the former only on real momenta.
In analytic formulas below \((t^-)^{\mathsf T}\) is a holomorphic
row, not a conjugation of a complex variable.

At nonzero fine frequency, the elementary Hodge identity is
\(\|\mathrm{d}\widehat w\|^2+\|\mathrm{d}^*\widehat w\|^2
=\lambda\|\widehat w\|^2\).
Subtracting a fine gradient makes a field divergence free
without changing its curl.  Its straight average changes by
a coarse gradient, by \eqref{eq:flat-exponential-kernels-FX2}.  Conversely any coarse gradient
can be corrected by a fine gradient, since \(q_n\) is onto.
It follows that the curl minimum in \eqref{eq:flat-exponential-kernels-FX3} equals the minimum
of the Feynman quadratic value over coarse gauge changes:
\begin{equation}
 E_n(p;a)=\min_{z\in\mathbb C}
            (a+t^+z)^*D_n(p)^{-1}(a+t^+z).
 \label{eq:flat-exponential-kernels-FX10}
\end{equation}
For clarity, the two inequalities use different simple steps.
A divergence-free representative of a curl-minimizing field
has Feynman energy equal to its curl energy, with coarse value
\(a+t^+z\).  In the other direction a Feynman minimizer at
that coarse value can be gradient-corrected back to \(a\),
and its curl energy is no greater than its Feynman energy.

Solving the scalar minimum gives the precision symbol
\begin{equation}
 \mathcal E_n=D_n^{-1}
   -\frac{D_n^{-1}t^+(t^-)^{\mathsf T}D_n^{-1}}
                 {(t^-)^{\mathsf T}D_n^{-1}t^+}.
 \label{eq:flat-exponential-kernels-FX11}
\end{equation}
Here \(D_n=\operatorname{diag}(D_{\mu,n})\).
At the real zero momentum the classical value is zero on every
constant coarse one-form: its constant fine lift is \(a/n\).
We accordingly set \(\mathcal E_n(0)=0\).  This zero mode is
retained, not removed by a mass term.  We next prove that this
definition is the value of a uniformly holomorphic symbol.

\subsubsection*{Uniform noncentral alias bounds}

Fix a small \(\rho>0\) and allow
\(|\operatorname{Re}p_\nu|<\pi+\rho\),
\(|\operatorname{Im}p_\nu|<\rho\).
For \(m_\nu\ne0\),
\begin{equation}
 |n\sin((p_\nu+2\pi m_\nu)/(2n))|
                      \ge c(1+|m_\nu|).
 \label{eq:flat-exponential-kernels-FX12}
\end{equation}
To see this, the real part of the sine argument stays in a
fixed subinterval of \((-\pi,\pi)\), and its absolute value
is at least \(c|m_\nu|/n\).  On that interval
\(|\sin x|\ge c|x|\), and
\(|\sin(x+iy)|^2=\sin^2x+\sinh^2y\).
The numerator in \eqref{eq:flat-exponential-kernels-FX8} is bounded.  For \(m_\nu=0\) the ratio
is a quotient of two sinc functions on a fixed compact
zero-free neighborhood.  Hence, uniformly in \(n\),
\begin{equation}
 |q_{\nu,m}(p)|\le\frac C{1+|m_\nu|}.
 \label{eq:flat-exponential-kernels-FX13}
\end{equation}

For \(m\ne0\) the denominator in \eqref{eq:flat-exponential-kernels-FX8} has a uniform positive
real part.  In fact
\(\operatorname{Re}\sin^2(x+iy)
\ge\sin^2x-\sinh^2y\).
After multiplication by \(4n^2\), the imaginary displacement
costs at most \(C\rho^2\), whereas the real sum is at least
\(c|m|^2\).  Decrease \(\rho\) to obtain
\begin{equation}
 \operatorname{Re}s_{n,m}(p)\ge c(1+|m|^2)
                                \qquad(m\ne0).
 \label{eq:flat-exponential-kernels-FX14}
\end{equation}
The constants do not grow as the alias set grows.

Write \(s_n=s_{n,0}\), \(q_\nu=q_{\nu,0}\),
\(Q_n=\prod_\nu q_\nu^2\), and \(f_\mu=Q_nq_\mu^2\).
Then
\begin{equation}
 D_{\mu,n}=\frac{f_\mu}{s_n}+R_{\mu,n},\qquad
 R_{\mu,n}=\sum_{m\ne0}\frac{f_{\mu,m}}{s_{n,m}}.
 \label{eq:flat-exponential-kernels-FX15}
\end{equation}
Each summand of \(R\) is bounded by the absolutely summable
majorant
\begin{equation}
 \frac{C}{1+|m|^2}
          \prod_{\nu=1}^4(1+|m_\nu|)^{-2}.
 \label{eq:flat-exponential-kernels-FX16}
\end{equation}
The extra factor \((1+|m_\mu|)^{-2}\) may be discarded.
Thus \(R\) is holomorphic and uniformly bounded on the stated
box, independently of \(n\).  Cauchy on a smaller box bounds
all its fixed-order derivatives uniformly as well.  For \(n=1\)
the sum is empty.  Every noncentral summand vanishes at \(p=0\),
so \(R_{\mu,n}(0)=0\); only boundedness is needed below.

\subsubsection*{Cancellation of the apparent gauge denominator}

Near \(p=0\), uniformly in \(n\),
\(q_\nu=1+O(|p|^2)\), \(s_n=O(|p|^2)\), and
\(f_\mu=1+O(|p|^2)\).  Define
\begin{equation}
 h_\mu=(f_\mu+s_nR_{\mu,n})^{-1},\qquad
 H_n=\operatorname{diag}(h_\mu).
 \label{eq:flat-exponential-kernels-FX17}
\end{equation}
On a sufficiently small fixed complex polydisc these are
holomorphic and uniformly bounded.  The exact filter identity
\begin{equation}
 \sum_\mu\frac{t_\mu^+t_\mu^-}{f_\mu}
       =\frac{4n^2\sum_\mu\sin^2(p_\mu/(2n))}{Q_n}
       =\frac{s_n}{Q_n}
 \label{eq:flat-exponential-kernels-FX18}
\end{equation}
is essential.  It follows directly by canceling the numerator
\(\sin^2(p_\mu/2)\) against \(q_\mu^2\), including removable
values.  It is not an approximation to a gauge denominator.

Since \(h_\mu=f_\mu^{-1}-s_nR_{\mu,n}h_\mu/f_\mu\),
\eqref{eq:flat-exponential-kernels-FX18} gives
\begin{equation}
 \begin{aligned}
 F_n&:=\sum_\mu t_\mu^+t_\mu^-h_\mu=s_nJ_n,\\
 J_n&=Q_n^{-1}
       -\sum_\mu\frac{t_\mu^+t_\mu^-R_{\mu,n}h_\mu}{f_\mu}.
 \end{aligned}\label{eq:flat-exponential-kernels-FX19}
\end{equation}
Here \(J_n(0)=1\) and \(|J_n-1|\le C|p|^2\) uniformly.
It is therefore nonzero on a smaller fixed polydisc.
On its meromorphic domain \(D_n^{-1}=s_nH_n\), and the
denominator in \eqref{eq:flat-exponential-kernels-FX11} equals \(s_nF_n=s_n^2J_n\).
The actual precision consequently has the holomorphic form
\begin{equation}
 \mathcal E_n(p)=s_nH_n
       -\frac{H_nt^+(t^-)^{\mathsf T}H_n}{J_n}.
 \label{fx:origin-cancellation}
\end{equation}
This removes the whole complex zero cone of \(s_n\), not merely
the value at the real origin.  The expression is uniformly
bounded there, is \(O(|p|^2)\), and gives
\(\mathcal E_n(0)=0\).  It also obeys the exact Ward identity
\(\mathcal E_nt^+=0\), by \eqref{eq:flat-exponential-kernels-FX19}.
For \(n=1\), it reduces to
\(\lambda(p)I-t^+(t^-)^{\mathsf T}\), the ordinary curl
quadratic form.  This checks the normalization and phases.

\subsubsection*{A uniform strip and its Fourier kernels}

We establish a complex domain away from zero explicitly.  On the real
Brillouin cube the central \(q_\nu\) are positive and bounded
above and below uniformly in \(n\).  Also
\(c|p|^2\le s_n(p)\le C|p|^2\).
For real \(|p|\ge r>0\), the central term in \eqref{eq:flat-exponential-kernels-FX15} and the
nonnegative alias terms give
\begin{equation}
 c_D\le D_{\mu,n}(p)\le C_r,
 \qquad
 (t^-)^{\mathsf T}D_n^{-1}t^+\ge c_r>0.
 \label{eq:flat-exponential-kernels-FX21}
\end{equation}
The second inequality uses
\(\sum_\mu|t_\mu^+|^2\ge c|p|^2\) and the upper bound on \(D\).

On a sufficiently thin complex neighborhood of this real
compact set, \(\operatorname{Re}s_n\ge cr^2-C|\operatorname{Im}p|^2\).
Thus the central term is holomorphic and bounded there;
\eqref{eq:flat-exponential-kernels-FX12}--\eqref{eq:flat-exponential-kernels-FX16} do the same for the aliases.  Cauchy bounds
the first derivatives uniformly in \(n\).  Shrinking the
neighborhood by a fixed amount depending on \(r\), \eqref{eq:flat-exponential-kernels-FX21}
therefore keeps every \(D_\mu\) and its scalar gauge denominator
nonzero.  The formula in \eqref{eq:flat-exponential-kernels-FX11} is holomorphic and uniformly bounded
on that neighborhood.  This is a quantitative perturbation
of uniformly separated real denominators.

Choose \(r\) inside the origin polydisc of \eqref{fx:origin-cancellation}, then reduce
the common width once.  The two formulas agree wherever their
denominators are nonzero, by their exact algebraic derivation,
and hence on their overlaps by continuity.  Together they give
constants \(\eta>0,M_E<\infty\), independent of \(n\), such that
\begin{equation}
 \mathcal E_n\text{ is holomorphic and }
 \sup_{|\operatorname{Im}p_\nu|<\eta}\|\mathcal E_n(p)\|\le M_E.
 \label{eq:flat-exponential-kernels-FX22}
\end{equation}
The symbol is \(2\pi\)-periodic.  In \eqref{fx:alias-covariance}, changing a coarse
momentum by \(2\pi\) reindexes the aliases; an endpoint alias
changes its fine frequency by \(2\pi\), under which the actual
filter and denominator are periodic.  The same is true of
\eqref{eq:flat-exponential-kernels-FX11}, and its analytic continuation has that periodicity.

Define its infinite-lattice kernel by
\begin{equation}
 \mathcal K_n(z)=(2\pi)^{-4}\int_{[-\pi,\pi]^4}
                         e^{ip\cdot z}\mathcal E_n(p)\,dp.
 \label{eq:flat-exponential-kernels-FX23}
\end{equation}
For \(0<\eta'<\eta\), shift the \(\nu\)-th contour to
\(\operatorname{Im}p_\nu=\eta'\operatorname{sgn}z_\nu\).
Periodicity cancels the vertical faces.  The exponential then
gives
\begin{equation}
 \|\mathcal K_n(z)\|\le M_Ee^{-\eta'|z|_1}.
 \label{eq:flat-exponential-kernels-FX24}
\end{equation}
All diagonal terms are retained in this Fourier definition.
The kernel and its weighted absolute row sums are bounded
uniformly in \(k\).  Its zero-frequency value and Ward identity
are those already proved; no contact term is added afterward.

\subsubsection*{The actual next-step precision and its inverse}

Now group the current lattice into the fixed blocks of side
\(L\) used by the next RG step.  For block Bloch momentum
\(P\), let \(r\in B_L(0)\), \(j\in I_L^4\), and
\(\zeta_j=(P+2\pi j)/L\).  The finite cell transform and its
analytic inverse are
\begin{equation}
 U(P)_{rj}=L^{-2}e^{i\zeta_j\cdot r},\qquad
 V(P)_{jr}=L^{-2}e^{-i\zeta_j\cdot r}.
 \label{eq:flat-exponential-kernels-FX25}
\end{equation}
Finite geometric-series orthogonality proves \(VU=UV=1\).
For real \(P\), \(V=U^*\).  For complex \(P\) the displayed
inverse, not literal conjugation, is used.  Their norms are
bounded on a fixed strip, since their extra diagonal factors
are \(e^{\pm iP\cdot r/L}\) and \(|r_\nu|/L<1/2\).
The folded classical precision is
\begin{equation}
 \mathcal E_n^{\rm cell}(P)
     =U(P)\operatorname{diag}_j(\mathcal E_n(\zeta_j))V(P),
 \label{eq:flat-exponential-kernels-FX26}
\end{equation}
with the identity on component indices understood.  It is
periodic and holomorphic on a fixed block strip, uniformly
in \(n\).  Periodicity again follows by the alias permutation.

Let \(C(P)\) be the actual flat middle-pivot basis for the
plain one-step constraint.  Its free restriction is the
identity and its pivot block solves the coefficient
\(L^{-3}I_4\); hence its entries are fixed Laurent polynomials
in the block phases.  Let \(G\) be the internal contour
matrix defined in \eqref{eq:flat-exponential-kernels-FX1a}, with rows \((y,x)\), \(x\ne y\).
Define
\(C^\sharp(P)=C(-P)^{\mathsf T}\); on real \(P\) this is \(C(P)^*\).
The exact classical free nonpivot precision is
\begin{equation}
 \mathsf A_n(P)=C^\sharp(P)
        [\mathcal E_n^{\rm cell}(P)+G^{\mathsf T}G]C(P).
 \label{eq:flat-exponential-kernels-FX27}
\end{equation}
This is the action pullback plus the actual soft gauge.  The
Feynman form was used only to compute the gauge-invariant
minimum in \eqref{eq:flat-exponential-kernels-FX3}; it has not replaced the soft gauge in \eqref{eq:flat-exponential-kernels-FX27}.

For clarity the uniform real gap used next can be obtained
independently of this symbol calculation.  Expanding the curl
of \eqref{eq:flat-exponential-kernels-FX1} gives an average of \(n\) by \(n\) fine plaquette
sheets.  Its nonnegative row weights sum to \(n^2\) and its
total column weights to \(n^{-2}\).  Weighted Cauchy gives
\(\|\mathrm{d}_nS_nw\|\le\|\mathrm{d}w\|\), hence
\(E_n(a)\ge\|\mathrm{d}a\|^2\) in current units.
Consequently \eqref{eq:flat-exponential-kernels-FX27} dominates the fixed-cell form
\(C(P)^*[d_P^*d_P+G^{\mathsf T}G]C(P)\).
This form has no zero vector: zero curl makes internal path
sums independent of the path, so zero \(G\) kills every
internal bond.  A plaquette crossing a block face then makes
its crossing values equal tangentially.  The plain constraint
kills that common face value.  The free restriction of \(C\)
is the identity, so the original free vector is zero.
The matrices in this comparison are continuous on the compact
block Bloch torus, with fixed dimension and fixed rank.  Their
least Rayleigh quotient on the unit sphere therefore has a
positive minimum \(c(L)\).  This proves
\(\mathsf A_n(P)\ge c(L)1\) uniformly in \(n\), including \(P=0\).
\eqref{eq:flat-exponential-kernels-FX22}, the finite transform and the fixed polynomial
matrices in \eqref{eq:flat-exponential-kernels-FX27} give a uniform bound for \(\mathsf A_n\)
on a complex strip.  Cauchy on a smaller strip bounds its
first derivatives by a constant \(M_A(L)\).
Choose its width so that
\begin{equation}
 \|\mathsf A_n(P+iY)-\mathsf A_n(P)\|
          \le4M_A(L)|Y|_\infty<c(L)/2.
 \label{eq:flat-exponential-kernels-FX28}
\end{equation}
The Neumann series around the real positive matrix then proves
\begin{equation}
 \mathsf A_n^{-1}\text{ is holomorphic},\qquad
          \|\mathsf A_n(P+iY)^{-1}\|\le2/c(L)
 \label{eq:flat-exponential-kernels-FX29}
\end{equation}
on one strip independent of \(n\).  This step uses both the
proved analytic bounds and the gap.

Contour shifting as in \eqref{eq:flat-exponential-kernels-FX23}--\eqref{eq:flat-exponential-kernels-FX24} now bounds the Fourier
coefficients of the block precision and its inverse by
\(C(L)e^{-c(L)|m|_1}\).  The cell has fixed size, so converting
its indices back to current nonpivot bonds only changes the
constants.  We obtain
\begin{equation}
 |\mathsf A_k^{\rm flat}(b,b')|
       +|\mathsf C_k^{\rm flat}(b,b')|
       \le C_L e^{-\mu_L d(b,b')},
 \qquad \mathsf C_k^{\rm flat}=(\mathsf A_k^{\rm flat})^{-1}.
 \label{fx:flat-kernel-bounds}
\end{equation}
Here \(d(b,b')\) is bond distance in current lattice units;
using starting-site distance or midpoint distance only changes
the fixed constants.
The symbols \(\mathsf C_k^{\rm flat}\) in this conclusion are
covariances; the undecorated pivot map \(C(P)\) is a different
rectangular matrix.  Exponential weighted row and column
sums follow by summing the fixed-dimensional lattice shells.

\subsubsection*{Finite tori, holonomies, and scope}

For a periodic lattice the exact quadratic minimization is
diagonal on its allowed Fourier momenta.  Its symbols are
the samples calculated above, with no volume correction.
If \(K_n(m)\) is any infinite block kernel just obtained and
the block period is \(\mathsf s_{\rm b}\), discrete Fourier
orthogonality and absolute convergence give
\begin{equation}
 K_{n,\mathsf s_{\rm b}}(m)
          =\sum_{v\in\mathbb Z^4}K_n(m+\mathsf s_{\rm b}v).
 \label{eq:flat-exponential-kernels-FX31}
\end{equation}
The image sum preserves an exponential bound in periodic
distance, uniformly over the admissible periods.  It identifies
the infinite-lattice kernel determined by the same symbol,
with exponentially small image errors on fixed nonwrapping
neighborhoods.  It does not assert exact equality of a full
torus kernel and an infinite kernel on a lift.  The same
argument applies to the inverse symbol, since its sampled
matrix is the inverse of each sampled positive precision.

For real flat holonomies in the prescribed block axial
representative, a unitary gauge constant on each block
flattens the connection on the universal cover and preserves
the contour gauge norm.  The commuting adjoint holonomies
split into unitary characters.  Each character shifts the
allowed block momenta by a real amount; the infinite symbols
are unchanged.  Equivalently the image sum in \eqref{eq:flat-exponential-kernels-FX31} receives
unit-modulus factors.  All absolute kernel bounds remain
valid, without requiring small global holonomies.

The conclusion is the exponential precision and covariance
bound for the actual unweighted classical flat quadratic
problem.  The auxiliary \(D_n\) itself is massless and is not
asserted to have an exponential kernel.  The gap and decay
in \eqref{fx:flat-kernel-bounds} concern the next-step constrained nonpivot variables
with the soft gauge, not a physical Yang--Mills mass gap.
No signed profile, counterterm, nonzero-current, non-Abelian
or enlarged complex-decoupling comparison has been proved
by this calculation.  Those require their actual operators
and uniform perturbation estimates.


% END MODULE chapters/flat_exponential_kernels_v8.tex

% BEGIN MODULE chapters/first_step_nonflat_precision.tex
\subsection{First-step precision at small curvature}
\label{ch:first-step-nonflat-precision}

This section concerns one unweighted step on a unit-spacing lattice.
It starts with a given real field and the actual nonlinear block
average.  It proves a gap and exponential covariance bounds for the
physical curved-fibre Hessian at a given constrained-critical
background.  Small plaquettes suffice; a small global potential or
small torus holonomies are not required.  The argument does not
construct the background, iterate the conclusion through lower
levels, or control signed profiles and complex decoupling parameters.

A fine current and a coarse multiplier have different types.  We
keep them distinct throughout.  An auxiliary extension with an
independently prescribed fine current will be useful in the proof.
Its agreement with any separately named off-shell completion is an
additional identification, stated precisely at the end.

\subsubsection*{The actual map and the two currents}

Fix odd \(L\ge3\).  Positive fine bonds carry Hermitian traceless
matrices with colour pairing
\(\langle X,Y\rangle=\tfrac12\operatorname{tr}(XY)\).
Complexification uses this bilinear pairing and its transpose.
The \(\ell^2\) norm on real fields uses the associated Euclidean
colour norm.  Operator estimates on complex fields use its Hermitian
extension.  No bare-spacing factor occurs in the exponential.

Use the centered blocks, contour-path multiset and middle pivots of
\cref{ch:block-average-pivot-chart}.  Denote the actual contour
mean from \(y\) to \(x\in B(y)\) by \(T_{y,x}(U)\).  Thus it solves
\(\sum_{\Gamma}\log(U(\Gamma)T_{y,x}^{-1})=0\), with the prescribed
equal weights.  The block bond \(c=(y,y')\), \(y'=y+Le_\mu\), is
\begin{equation}
 \begin{aligned}
 Z_x(U)&=T_{y,x}(U)U([x,x+Le_\mu])
 T_{y',x+Le_\mu}(U)^{-1}U[c]^{-1},\\
 \mathcal M(U)(c)&=
 \exp\left\{L^{-4}\sum_{x\in B(y)}\log Z_x(U)\right\}U[c].
 \end{aligned}
 \label{eq:first-step-nonflat-precision-NFP1}
\end{equation}
The internal implicit mean and the outer explicit exponential are
different operations.  The analytic branches are the ones fixed by
the local pivot construction.  On a torus every relevant double
block and its bond neighbourhood must have the embedded rectangular
lift used there.  Periods divisible by \(L\) and exceeding \(4L+4\)
are sufficient.  Work on the periodic cover whenever a larger local
patch is needed.

The unweighted action and actual soft gauge are
\begin{equation}
 A(U)=\sum_p\left[1-\tfrac12\operatorname{tr}U(\partial p)\right],
 \qquad
 \mathrm{GF}(U)=\sum_y\sum_{x\in B(y)\setminus\{y\}}
             \left[1-\tfrac12\operatorname{tr}T_{y,x}(U)\right].
 \label{eq:first-step-nonflat-precision-NFP2}
\end{equation}
These expressions equal their usual real-trace versions on
\(\operatorname{SU}(2)\).  Their complex continuations use the trace in
\eqref{eq:first-step-nonflat-precision-NFP2}, without taking a real part.  All internal contour paths avoid
every middle crossing pivot.  Hence GF depends on no pivot variable.
On the infinite lattice these sums specify local interactions:
derivatives are taken of their finite differences under compactly
supported variations.  Only finitely many local terms then change.
The coefficient bounds below extend the resulting matrices to
\(\ell^2\); no finite value for the whole background action is assumed.

Let \(\mathfrak X\) be the positive nonpivot bonds and \(\iota\) their
inclusion into all fine bonds.  For a background \(U\), define
\begin{equation}
 \begin{gathered}
 q_U(a)=\frac1i\log\big[\mathcal M(e^{ia}U)\mathcal M(U)^{-1}\big],\\
 P_U=Dq_U(0),\qquad P_Uh_U=1,\qquad
 C_U=\iota-h_UP_U\iota.
 \end{gathered}
 \label{eq:first-step-nonflat-precision-NFP3}
\end{equation}
Here \(h_U\) is supported only on the pivot of each coarse bond.
The actual fibre chart \(\Phi_U\) has free restriction \(u\), satisfies
\(q_U(\Phi_U(u))=0\), and obeys
\begin{equation}
 \Phi_U(0)=0,\quad D\Phi_U(0)=C_U,\quad
 D^2\Phi_U(0)[u,v]
 =-h_UD^2q_U(0)[C_Uu,C_Uv].
 \label{eq:first-step-nonflat-precision-NFP4}
\end{equation}
Indeed differentiating the constraint gives its displayed linear and
quadratic identities.  The second derivative of the chart is
supported on pivots, so applying \(h_UP_U\) to it is the identity.

Write \(J_A(U)\) for the Wilson gradient in the left exponential
coordinates.  For an arbitrary fine current \(J\), set
\(\mu_U(J)=h_U^{\mathsf T}J\); this is a coarse row covector.
Define the auxiliary matrix by its bilinear form
\begin{equation}
 \begin{split}
 \mathscr K_U[J](u,v)
 &=D^2 A_U[C_Uu,C_Uv]
   +D^2\mathrm{GF}_U[\iota u,\iota v]\\
 &\quad-\big\langle\mu_U(J),
           D^2q_U(0)[C_Uu,C_Uv]\big\rangle .
 \end{split}\label{eq:first-step-nonflat-precision-NFP5}
\end{equation}
The equality \(D^2\mathrm{GF}[C_Uu,C_Uv]
=D^2\mathrm{GF}[\iota u,\iota v]\) holds because GF has no pivot
coordinate.  The same fact gives
\(h_U^{\mathsf T}J_{\mathrm{GF}}(U)=0\) identically, even if the
contour means are not the identity.

The chain rule and \eqref{eq:first-step-nonflat-precision-NFP4} now give the exact physical identity
\begin{equation}
 \begin{split}
 K_{\rm phys}(U)
 &:=D_u^2\{A(e^{i\Phi_U(u)}U)
                 +\mathrm{GF}(e^{i\Phi_U(u)}U)\}_{u=0}
       ,\\
 K_{\rm phys}(U)&=\mathscr K_U[J_A(U)],\\
 \mathscr K_U[J]-K_{\rm phys}(U)
 &=\big\langle J-J_A(U),D^2\Phi_U(0)\big\rangle .
 \end{split}\label{eq:first-step-nonflat-precision-NFP6}
\end{equation}
No criticality is needed for this coordinate identity.  Its Gaussian
interpretation uses a given constrained-critical background in the
block-axial representative.  All internal contour means then equal
one.  The GF gradient vanishes and its Hessian is the sum of the squares of
the derivatives of the actual contour means.  Constrained criticality
also gives
\begin{equation}
 J_A(U)=P_U^{\mathsf T}\mu_U(J_A(U)).
 \label{eq:first-step-nonflat-precision-NFP7}
\end{equation}
To see this, the gradient annihilates \(\ker P_U\), and \(P_U\) is
onto by its pivot right inverse.  Thus it is \(P_U^{\mathsf T}\mu\);
applying \(h_U^{\mathsf T}\) identifies \(\mu\).  An independent
\(J\) is not asserted to satisfy \eqref{eq:first-step-nonflat-precision-NFP7}.

At this first step there is no inner background minimization: the
zero-fold averaging and site maps are identities.  In a formulation
using an intermediate Landau projector, that projector is zero,
and the zero-fold constraint has no quadratic part.  These facts
do not by themselves identify or annihilate a separately defined
invariant-Hessian correction such as \(\Delta'_\pi[J]\), even after
projection.  The physical identity above follows directly from the
actual chart chain rule, independently of that normal-form check.
In canonical Gaussian coordinates
\(a=gu\), an action coefficient \(x\) with \(xg^2=1\) gives precisely
this same unweighted Hessian.

\subsubsection*{A joint analytic domain and explicit support counts}

We record the finite bounds used below.  Constants \(c,C\) in this
section are universal, with their value allowed to change.  There
are radii
\begin{equation}
 r_b=c_bL^{-5},\qquad r_u=c_uL^{-9}
 \label{eq:first-step-nonflat-precision-NFP8}
\end{equation}
with universal positive \(c_b,c_u\), the latter fixed after the
former, such that, for \(U=e^{iX}\), the actual chart is jointly holomorphic
on \(\|X\|_\infty<8r_b\), \(\|u\|_\infty<8r_u\), and
\begin{equation}
 \begin{aligned}
 \|h_U\|_{\infty\to\infty},\ \|h_U\|_{1\to1}&\le CL^3,\\
 \|C_U\|_{\infty\to\infty},\ \|C_U\|_{1\to1}&\le CL^4,\\
 \|\Phi_U(u)\|_\infty&\le CL^4\|u\|_\infty.
 \end{aligned}
 \label{eq:first-step-nonflat-precision-NFP9}
\end{equation}
The notation for \(\ell^1\) and \(\ell^\infty\) here includes the
fixed three-dimensional colour space; changing among its coordinate
norms changes only \(C\).

Here are the domain and locality details.  Products along a contour
of length \(O(L)\), the implicit internal mean, and the logarithms in
\eqref{eq:first-step-nonflat-precision-NFP1} are jointly holomorphic for \(\|X\|,\|a\|<c/L\).
The product and logarithm estimates of the pivot construction give
\(\|P_U\|_{\infty\to\infty}\le CL\) and a quadratic constraint
remainder at most \(CL^2\|a\|^2\).  A background derivative of
\(P_U\), obtained on this larger joint domain by Cauchy's formula,
is at most \(CL^2\).  Consequently the pivot coefficient differs
from \(L^{-3}1\) by at most \(CL^2\|X\|\).
For the first radius in \eqref{eq:first-step-nonflat-precision-NFP8}, its inverse has norm at most \(CL^3\).
The pivot matrix is block diagonal over coarse bonds: no other pivot
occurs in the support of a given constraint row.  This proves the
bound for \(h_U\).

A constraint row uses at most \(CL^4\) fine bonds.  A fine bond lies
in at most eight relevant double blocks.  These facts and the same
directional Cauchy estimate give the column bound \(CL\) for
\(P_U\), as well as its row bound.  The formula in \eqref{eq:first-step-nonflat-precision-NFP3} proves both bounds
for \(C_U\).  Solving the nonlinear remainder with the pivot inverse
has contraction constant at most
\(CL^3L^2(L^4\|u\|)=CL^9\|u\|\).
Taking the second radius in \eqref{eq:first-step-nonflat-precision-NFP8} small secures a contraction on
the whole indicated joint domain.  Its iterates are holomorphic in
both parameters and converge uniformly there on smaller closed
polydiscs.  Passing their Cauchy integrals to the limit proves joint holomorphy.

For locality, the pivot value \(\Phi_U(u)(b_0(c))\) depends only on
the background and free variables in its own double block.  The
component contraction for this pivot does not involve another pivot.
All nonpivot values are exactly \(u\).  In particular:

\begin{enumerate}
\item A given free coordinate affects its own full bond and at most
eight pivot components.  Each of these bonds belongs to six
plaquettes.  Thus at most \(54\) Wilson plaquette terms can depend
on that free coordinate through the chart.
\item GF has no pivot.  A free coordinate can occur only in the
\(L^4-1\) internal means of its own block.  Counting contour-path
multiplicities does not increase the number of these terms.
\item Each such composed term depends on at most \(CL^4\) scalar
coordinates.  For a fixed free bond, the union of all possible other
free coordinates coupled to it still lies in a fixed number of
neighbouring blocks and has size \(CL^4\).  This union bound does
not multiply the two preceding counts together.
\end{enumerate}
All matrix coefficients are therefore finite range.  A double block
has \(\ell^1\) diameter at most \(5L\).  Two pivot bonds entering
one plaquette have starting sites at distance at most two.  Thus the
union of their chart supports has diameter at most \(10L+2\le12L\).
The safe range \(12L\), measured by starting sites of fine bonds,
also accommodates the internal contour terms.
The same bound applies to the background variables on which a
coefficient depends.

Each composed local action or gauge term is bounded by a universal
constant on the joint domain above.  Indeed all the products have
\(O(L)\) factors, with \(X\) and \(\Phi\) of size \(O(L^{-5})\),
and the contour mean remains in its fixed analytic neighbourhood.
Cauchy's formula in two free scalar directions bounds its Hessian
coefficient by \(Cr_u^{-2}\le CL^{18}\).
The \(CL^4\) possible terms give \(CL^{22}\) per matrix entry, and
the \(CL^4\) possible coupled coordinates give the Schur bound
\(CL^{26}\).  One more background-direction Cauchy estimate costs
\(r_b^{-1}\le CL^5\).  Here and below the Schur norm is the maximum
of the absolute row and column sums of colour block norms.  We have
proved
\begin{equation}
 \begin{aligned}
 \|K_{\rm phys}(e^{iX})\|_{\rm Sch}&\le CL^{26},\\
 \|K_{\rm phys}(e^{iX_1})-K_{\rm phys}(e^{iX_0})\|_{\rm Sch}
       &\le CL^{31}\|X_1-X_0\|_\infty.
 \end{aligned}
 \label{eq:first-step-nonflat-precision-NFP10}
\end{equation}
on a fixed smaller background domain.  Only local terms were
estimated; no extensive action bound independent of the volume was
used.

For the current term, Cauchy's formula and \eqref{eq:first-step-nonflat-precision-NFP9} give
\(|\partial_{u_b}\partial_{u_{b'}}\Phi(b_0(c))|
\le CL^4/r_u\le CL^{13}\).
A fixed free coordinate occurs in at most eight such pivot
components and couples to \(CL^4\) others.  Hence
\begin{equation}
 \|\langle J,D^2\Phi_U(0)\rangle\|_{\rm Sch}
       \le CL^{17}\|J\|_\infty.
 \label{eq:first-step-nonflat-precision-NFP11}
\end{equation}
The corresponding background derivative costs another \(CL^5\).
The Wilson gradient near the identity obeys
\(\|J_A(e^{iX})\|_\infty\le C\|X\|_\infty\), by differentiating
the six adjacent plaquette products.  Combining \eqref{eq:first-step-nonflat-precision-NFP6}, \eqref{eq:first-step-nonflat-precision-NFP10} and
\eqref{eq:first-step-nonflat-precision-NFP11} proves, for \(\|J_i\|\le1\),
\begin{equation}
 \begin{split}
 \|\mathscr K_{e^{iX}}[J]-K_0\|_{\rm Sch}
       &\le CL^{32}(\|X\|_\infty+\|J\|_\infty),\\
 \|\mathscr K_{e^{iX_1}}[J_1]-\mathscr K_{e^{iX_0}}[J_0]\|_{\rm Sch}
       &\le CL^{32}
          (\|X_1-X_0\|_\infty+\|J_1-J_0\|_\infty).
 \end{split}\label{eq:first-step-nonflat-precision-NFP12}
\end{equation}
Here \(K_0=\mathscr K_1[0]\) is the actual first flat precision.
These are derived coefficient estimates, not assumptions about a
matrix perturbation.  All counts and domains are independent of
the torus volume.

\subsubsection*{Local frames which preserve the actual soft gauge}

Suppose now that \(U\) is real and block axial, and
\(\sup_p\|U(\partial p)-1\|\le\epsilon\).
The relevant symmetry of GF is a unitary gauge transformation
constant on each centered block.  Its values may vary from block
to block.  If \(v\) is such a transformation and
\((V_v a)(b)=\operatorname{Ad}(v(b_-))a(b)\), then
\begin{equation}
 \Phi_{U^v}(V_vu)=V_v\Phi_U(u),\qquad
 \mathscr K_{U^v}[V_vJ]=V_v\mathscr K_U[J]V_v^{-1}
 \label{eq:first-step-nonflat-precision-NFP13}
\end{equation}
on free coordinates.  In this formula \(V_v\) is restricted to the
appropriate full or free space.  To verify it, internal means
transform by conjugation with the block value, so GF is unchanged.
The relative coarse logarithm in \eqref{eq:first-step-nonflat-precision-NFP3} transforms by that same
adjoint action at the coarse bond's starting site.  The pivot's
starting site belongs to that block.  The pivot right inverse and
free-identity basis therefore intertwine these actions.  Uniqueness
of the actual chart proves its first identity, and differentiating
the invariant local functions proves the second.  The transformations
are unitary in the colour norm.  No covariance of GF under an
arbitrary site gauge is used.

Consider a rectangular patch of diameter \(O(LR)\), with a fixed
additional margin of complete blocks.  A unitary comb gauge \(g\)
on that patch makes
\begin{equation}
 \|U^g(b)-1\|\le CLR\epsilon.
 \label{eq:first-step-nonflat-precision-NFP14}
\end{equation}
For completeness, compare the two comb paths to the endpoints of
a bond.  Moving that bond successively past the later coordinate
segments expresses the resulting closed transport as a product
of at most \(CLR\) conjugated elementary plaquettes.  Unitarity and
the telescoping identity for a product of unitary matrices bound
its distance from the identity by the sum of their distances.
This proves \eqref{eq:first-step-nonflat-precision-NFP14}, with no bound on \(g-1\).

Although \(g\) need not be block constant, block axiality supplies
the missing adjustment.  The contour means transform exactly as
\begin{equation}
 T_{y,x}(U^g)=g(y)g(x)^{-1}.
 \label{eq:first-step-nonflat-precision-NFP15}
\end{equation}
Every path in such a mean has length \(O(L)\).  By \eqref{eq:first-step-nonflat-precision-NFP14} its
transport, and therefore its implicit mean, is within
\(CL^2R\epsilon\) of the identity.  This last statement follows
directly from the contraction defining the mean: inputs within
\(t\) of the identity have mean within \(Ct\).
Let \(v(x)=g(y)\) on \(B(y)\).  Then
\(v(x)g(x)^{-1}=T_{y,x}(U^g)\), so \eqref{eq:first-step-nonflat-precision-NFP14}--\eqref{eq:first-step-nonflat-precision-NFP15} imply
\begin{equation}
 \|U^v(b)-1\|\le CL^2R\epsilon,
 \qquad v\text{ is block constant}.
 \label{eq:first-step-nonflat-precision-NFP16}
\end{equation}
We can take its logarithm with the same bound.  Thus the local frame
preserves the actual soft gauge and gives a small potential on the
patch.  Arbitrary coarse holonomies have been retained outside it.

If a real field is initially given without the axial condition,
small plaquettes make the contour-path transports mutually close:
the same plaquette sliding uses at most \(CL^2\) plaquettes.
The actual implicit means therefore exist.  Taking the site gauge
\(u(x)=T_{y,x}(U)\), \(u(y)=1\), makes all internal means equal to
one and leaves the coarse datum unchanged.  This supplies the
representative used here.  It does not assert existence of a
constrained-critical orbit.

\subsubsection*{The gap from a covariant partition of unity}

Let \(c_0=c_0(L)>0\) be the proved flat first-step gap
\(K_0\ge c_0\), as in \eqref{cf:classical-flat-gap}.
Inequalities for real symmetric matrices in this section use their
Hermitian quadratic forms on the complexified \(\ell^2\) space.
Increase a universal constant once so that
\(K_L=CL^{40}\) dominates the coefficient, support and partition
constants just obtained.  Choose
\begin{equation}
 \begin{gathered}
 R\ge40,\qquad K_L/R^2\le c_0/8,\\
 CL^2R\epsilon\le r_b/8,\qquad
 K_L(R\epsilon+\gamma)\le c_0/8.
 \end{gathered}
 \label{eq:first-step-nonflat-precision-NFP17}
\end{equation}
Here \(\gamma\) bounds the real fine current \(J\), and it is also
at most one.  These are explicit smallness requirements in terms
of \(L\), the finite flat gap and the already proved chart radii.

Construct a scalar squared partition on the coarse lattice.
Let \(h(t)=\max(1-|t|,0)\), put
\(h_z(x)=\prod_\nu h(x_\nu/R-z_\nu)\), and divide by
\((\sum_z h_z(x)^2)^{1/2}\).
At least one factor product is \(\ge2^{-4}\), and at most
\(2^4\) are nonzero.  The resulting \(\chi_z\) satisfy
\(\sum_z\chi_z^2=1\), have support of radius \(R\), and obey
\begin{equation}
 \sum_z|\chi_z(x)-\chi_z(x')|^2
       \le C|x-x'|_1^2/R^2
 \quad\text{if }|x-x'|_1\le20.
 \label{eq:first-step-nonflat-precision-NFP18}
\end{equation}
The estimate follows by differentiating the normalized hats along
the segment; only a bounded number of them occur there.  Evaluate
them at the block containing a free bond's starting site.

For a compactly supported free field \(f\), the exact matrix identity
and its bound are
\begin{equation}
 \begin{split}
 \langle f,\mathscr K_U[J]f\rangle
 &=\sum_z\langle\chi_zf,\mathscr K_U[J]\chi_zf\rangle+\mathcal E,\\
 |\mathcal E|&\le (K_L/R^2)\|f\|_2^2.
 \end{split}\label{eq:first-step-nonflat-precision-NFP19}
\end{equation}
Indeed \(1-\sum_z\chi_z(b)\chi_z(b')
=\tfrac12\sum_z(\chi_z(b)-\chi_z(b'))^2\).
The matrix has range \(12L\), so the block indices differ by at
most a fixed number.  Apply \eqref{eq:first-step-nonflat-precision-NFP18}, the row and column bounds,
and \(2|f_b||f_{b'}|\le|f_b|^2+|f_{b'}|^2\).
The localization is on free variables.  Its full tangent field is
\(C_U(\chi_zf)\); it is not replaced by \(\chi_zC_Uf\).

Use \eqref{eq:first-step-nonflat-precision-NFP16} on a larger patch around each support.  The margin
contains every background variable of its matrix coefficients.
\eqref{eq:first-step-nonflat-precision-NFP13} carries the localized form into that block-constant
frame.  Extend the small background logarithms by zero outside the
larger patch, and the transformed current by zero.  The coefficients
in the localized form are unchanged, by their strict locality.
The near-identity bound in \eqref{eq:first-step-nonflat-precision-NFP12} and the flat gap therefore give
\begin{equation}
 \langle\chi_zf,\mathscr K_U[J]\chi_zf\rangle
 \ge\{c_0-K_L(R\epsilon+\gamma)\}\|\chi_zf\|_2^2.
 \label{eq:first-step-nonflat-precision-NFP20}
\end{equation}
The harmless \(L^2\) in \eqref{eq:first-step-nonflat-precision-NFP16} has already been absorbed in
\(K_L=CL^{40}\).  On a double block the same local-frame argument
also proves the uniform coefficient bounds and existence of its
pivot inverse before \eqref{eq:first-step-nonflat-precision-NFP19} is used; this involves only the smaller
\(R=O(1)\) patch and has no dependence on the desired global gap.
Thus no matrix bound needed in the IMS error has been assumed.
Summing \eqref{eq:first-step-nonflat-precision-NFP20} and using \eqref{eq:first-step-nonflat-precision-NFP17}--\eqref{eq:first-step-nonflat-precision-NFP19} proves
\begin{equation}
 \mathscr K_U[J]\ge\tfrac12 c_0\,1.
 \label{nfp:real-gap}
\end{equation}

This proof takes place on the periodic cover.  It also gives the
torus inequality without requiring one local frame to descend to
the torus.  Lift a fixed periodic vector to the cover and restrict it
to a box of \(N^4\) period cells.  Its squared norm is \(N^4\) times
the period norm.  Finite range makes its quadratic form equal
\(N^4\) times the torus form, up to a boundary error of order
\(N^3\), for that fixed vector and period.  Divide the cover
inequality by \(N^4\) and let \(N\) tend to infinity.  The constant
in \eqref{nfp:real-gap} is independent of the period.  This also explains why
no assumption on noncontractible holonomies entered the proof.

For the physical current, each first variation of one plaquette
vanishes when its plaquette transport is the identity.  Differentiating
its four factors and using unitarity bounds its contribution by
\(C\|U(\partial p)-1\|\).  There are six adjacent plaquettes, hence
\begin{equation}
 \|J_A(U)\|_\infty\le C\epsilon.
 \label{eq:first-step-nonflat-precision-NFP22}
\end{equation}
Consequently a sufficiently small \(\epsilon\), depending only on
\(L\), satisfies \eqref{eq:first-step-nonflat-precision-NFP17} with \(J=J_A(U)\).  \eqref{eq:first-step-nonflat-precision-NFP6} and
\eqref{nfp:real-gap} prove the actual physical first-step gap at every given
real block-axial constrained-critical field of this curvature size.

\subsubsection*{Complex neighbourhood and exponential kernels}

The same local frames yield a quantitative complex statement.
Fix a real pair satisfying \eqref{eq:first-step-nonflat-precision-NFP17}, with room in those inequalities,
and put
\(\mathbf U=e^{iE}U\), \(\mathbf J=J+Z\), for complex fine fields.
The local unitary gauges preserve \(\|E\|_\infty\) and
\(\|Z\|_\infty\).  The logarithm in a small local frame has bounded
derivative, so the joint-domain estimate in \eqref{eq:first-step-nonflat-precision-NFP12}, now applied to
the background difference, proves
\begin{equation}
 \|\mathscr K_{\mathbf U}[\mathbf J]-\mathscr K_U[J]\|_{\rm Sch}
       \le K_L(\|E\|_\infty+\|Z\|_\infty)
 \label{eq:first-step-nonflat-precision-NFP23}
\end{equation}
provided \(\|E\|_\infty\le r_b/16\) and the total current remains
below one.  In particular require
\begin{equation}
 K_L(\|E\|_\infty+\|Z\|_\infty)\le c_0/8.
 \label{eq:first-step-nonflat-precision-NFP24}
\end{equation}
Schur's inequality bounds the \(\ell^2\) perturbation by the same
quantity.  The inverse follows by a Neumann series about the real
matrix and has norm at most \(4/c_0\).
The coefficients are holomorphic local functions, and their uniform
Cauchy bounds make that Neumann inverse holomorphic in every
admitted finite complex parameter family.

For physical analytic continuation take
\(\mathbf J=J_A(\mathbf U)\).
The adjacent-plaquette formula also gives
\(\|J_A(e^{iE}U)-J_A(U)\|_\infty\le C\|E\|_\infty\)
on the stated unitary-centred complex neighbourhood.  Thus a
smaller explicit bound \(\|E\|\le c_0/(CK_L)\), together with the
chart bound, implies \eqref{eq:first-step-nonflat-precision-NFP24}.  If such a field is a given block-axial complex
critical background, its physical Gaussian is the analytic matrix
in \eqref{eq:first-step-nonflat-precision-NFP6}.  Existence of that background is still a separate question.

Write \(r_L=12L\), increase \(K_L\) to bound the complex row and
column norms by \(2K_L\), and choose
\begin{equation}
 \tau_L=r_L^{-1}\log\big(1+c_0/(16K_L)\big)>0.
 \label{eq:first-step-nonflat-precision-NFP25}
\end{equation}
For a bounded real function \(\varphi\) which is one-Lipschitz in
starting-site distance, multiplication by \(e^{\tau_L\varphi}\)
changes the matrix by an operator of norm at most
\begin{equation}
 2K_L(e^{\tau_Lr_L}-1)\le c_0/8.
 \label{eq:first-step-nonflat-precision-NFP26}
\end{equation}
This follows entry by entry, from finite range, and then from the
row and column bounds.  The inverse of the conjugated matrix has
norm at most \(8/c_0\).  Take
\(\varphi(b)=\min\{d(b,b'),N\}\), apply this estimate between
coordinate vectors and then let \(N\) increase.  On the infinite
lattice the weights used before this limit are bounded; no
unbounded conjugation identity is assumed.  We obtain
\begin{equation}
 \|\mathscr K_{\mathbf U}[\mathbf J]^{-1}(b,b')\|
       \le (8/c_0)e^{-\tau_Ld(b,b')}.
 \label{nfp:kernel-decay}
\end{equation}
The same proof on a torus uses its periodic distance.  Fixed-matrix
Dirichlet compressions retain the real lower bound, coefficient
estimates and conjugation bound, so \eqref{nfp:kernel-decay} also holds for them.
This does not identify a matrix rebuilt from a different local
background or a different sequence of regions with that compression.
Finite range gives an exponential bound for the precision itself.
Summing the kernel bound with the smaller weight
\(e^{\tau_Ld/2}\) gives absolute weighted row and column bounds,
since a four-dimensional lattice has at most \(C(m+1)^3\) bonds
at distance in \([m,m+1)\).  All these constants depend only on
\(L\) and the fixed flat gap, not on volume or global holonomies.

\subsubsection*{The remaining identification and the exact scope}

The physical conclusion is fully determined by the actual chart
and the actual Wilson and contour-gauge functionals.  An extension
to independent-current cores must specify its complete projected
current and gauge form.  If its isolated one-step constraint
subtraction is called \(E^{(2)}_U[J]\), the necessary constraint-readout
condition for agreement with \eqref{eq:first-step-nonflat-precision-NFP5} is
\begin{equation}
 \langle C_Uu,E^{(2)}_U[J]C_Uv\rangle
   =\langle h_U^{\mathsf T}J,
                   D^2q_U(0)[C_Uu,C_Uv]\rangle.
 \label{eq:first-step-nonflat-precision-NFP28}
\end{equation}
Equivalently, if \(q_U^{(2)}\) denotes the homogeneous quadratic
part, its diagonal readout is
\(2\langle h_U^{\mathsf T}J,q_U^{(2)}(C_Uu)\rangle\).
This fixes the factor of two and the pivot-right-inverse composition.
The \hyperlink{first-step-named-form}{first-step calculation below}
verifies \eqref{eq:first-step-nonflat-precision-NFP28} for the
named form and checks its remaining terms on the axial locus.
The condition alone does not identify a general all-core
completion.  Its projected gauge block and every remaining
current or gauge-restoration contribution, including a term denoted
\(\Delta'_\pi[J]\), must also match.  The Hessian of the absolute
contour gauge used here agrees with the relative-contour gauge
Hessian on the axial locus, but need not agree away from it.
Agreement at physical critical pairs alone proves none of these
off-core identities.  We have defined and bounded \eqref{eq:first-step-nonflat-precision-NFP5} explicitly,
without assuming the identification being tested.

Thus the actual on-shell first unweighted Gaussian, and its admitted
physical analytic continuation, have the proved gap and kernels.
The auxiliary independent-current extension has the same bounds
under its explicit current ceiling.  No assertion about all lower
levels, signed layer coefficients, nonzero-current many-level
Schur complements, enlarged complex interpolation domains, or the
conditional exterior-force law follows from this one-step result.


% END MODULE chapters/first_step_nonflat_precision.tex

% BEGIN MODULE chapters/first_step_background.tex
\subsection{The first-step background and its regular minimum}
\label{ch:first-step-background-expanded}

We construct a background for one raw, unweighted step.  The input is
a real coarse connection with sufficiently small plaquettes.  Its
potential and its noncontractible holonomies need not be small.  The
output is an actual block-axial fine connection, with the prescribed
nonlinear average, which is constrained-critical for the Wilson action
and minimizes it within the explicit regular class proved below.
The construction gives exponential first responses, product-topology
continuity, and compatible local complex continuations.  It uses the
actual contour soft gauge as an auxiliary function in the construction.

The constants below depend on the fixed odd integer \(L\ge3\), but
not on volume.  We use the actual average, middle pivots and physical
left bond logarithms of
\cref{ch:first-step-nonflat-precision}.  On a periodic lattice
we retain its embedded-double-block condition and use the periodic
cover for larger patches.  The infinite-lattice action is interpreted
through finite local variations.  No finite total action is assumed.

\subsubsection*{An exact initial field}

For a coarse connection \(V\), define its face lift \(U_0(V)\) as
follows.  A fine bond internal to a centered block has value \(I\).
Every positive fine bond crossing from \(B(y)\) to \(B(y+Le_\mu)\)
has value \(V(y,\mu)\).  The value on a reversed bond is the inverse.
Every internal contour transport, and therefore every actual implicit
contour mean, is \(I\).  Every straight length-\(L\) segment used in
the average crosses its designated coarse face once and has transport
\(V(y,\mu)\), as does its reference segment.  Consequently
\begin{equation}
 T_{y,x}(U_0(V))=I,\qquad
 Z_x(U_0(V))=I,\qquad
 \mathcal M_L U_0(V)=V.
 \label{eq:first-step-background-BGE1}
\end{equation}
These are exact nonlinear identities.  A fine plaquette whose two
directions are not both crossing directions is \(I\).  If both
directions cross, its four factors are precisely the four factors of
the corresponding coarse plaquette.  In particular, for real \(V\),
\begin{equation}
 \sup_p\|U_0(V)(\partial p)-I\|
       \le\sup_P\|V(\partial P)-I\|=:\epsilon.
 \label{eq:first-step-background-BGE2}
\end{equation}
There is no holonomy restriction in this construction.

Write \(\Phi_V:=\Phi_{U_0(V)}\) for the actual fibre chart, with
identity free restriction.  Its free bond set \(\mathfrak X\) depends
only on the fixed block partition.  Define
\begin{equation}
 U_V(u)=e^{i\Phi_V(u)}U_0(V),\qquad
 F_V(u)=D_u\{A(U_V(u))+\mathrm{GF}(U_V(u))\}.
 \label{eq:first-step-background-BGE3}
\end{equation}
We identify a row derivative with a field using the fixed bilinear
colour pairing.  On the infinite lattice, a component of \(F_V\) is
the derivative under its compactly supported free coordinate variation.
Only finitely many local terms enter that component.  The estimates
below give a bounded analytic map on the stated \(\ell^\infty\) ball;
\eqref{eq:first-step-background-BGE3} does not assign a value to an infinite sum of action terms.

Let \(c_0=c_0(L)>0\) be the actual flat gap.  Fix a physical curvature
ceiling \(\alpha_L>0\) so that every real block-axial field of curvature
at most \(\alpha_L\), together with its Wilson current, satisfies the
smallness conditions of the preceding section with room.  At the face
lift, the actual GF gradient vanishes.  Thus
\begin{equation}
 \begin{aligned}
 F_V(0)&=C_{U_0(V)}^{\mathsf T}J_A(U_0(V)),&
 \|F_V(0)\|_\infty&\le A_L\epsilon,\\
 K_V&:=DF_V(0)=K_{\rm phys}(U_0(V)),&
 K_V&\ge\tfrac12c_0 I,
 \qquad A_L=C_*L^4.
 \end{aligned}
 \label{eq:first-step-background-BGE4}
\end{equation}
Here and below \(C_*\) is a sufficiently large fixed universal
constant.  The gap is legitimate before criticality is established:
the exact chain rule \eqref{eq:first-step-nonflat-precision-NFP6} and the proof of \eqref{nfp:real-gap} require a real
block-axial small-curvature field and its small physical current.
Criticality was required there only to interpret the resulting matrix
as a background Gaussian.  No such interpretation is used to obtain
\eqref{eq:first-step-background-BGE4}.

\subsubsection*{The bounded-field inverse and the nonlinear estimates}

The matrix \(K_V\) is real symmetric, has range \(12L\), and has
uniform absolute row and column bounds.  Its \(\ell^2\) inverse has
the actual exponential kernels proved in \eqref{nfp:kernel-decay}.  If \(\tau_L>0\)
is the rate there, put
\begin{equation}
 S_4(\tau_L)=\sum_{m\ge0}C_4(m+1)^3e^{-\tau_Lm},\qquad
 B_L=1+\frac8{c_0}S_4(\tau_L).
 \label{eq:first-step-background-BGE5}
\end{equation}
The universal \(C_4\) bounds the number of positive fine bonds in a
distance shell; it is valid also for periodic distance.  We may enlarge
it so \(B_L\ge8/c_0\).  Absolute kernel rows define an operator
\(H_V\) on bounded free fields, with
\begin{equation}
 \|H_V\|_{\infty\to\infty}\le B_L,\qquad
 K_VH_V=H_VK_V=I\quad\hbox{on }\ell^\infty.
 \label{eq:first-step-background-BGE6}
\end{equation}
For clarity, both identities are proved here.  Each matrix entry of
\(K_VH_V\) or \(H_VK_V\) equals the identity entry by the
\(\ell^2\) inverse identities.  Applying the former to a bounded field
involves finitely many \(K_V\) entries in each row and absolute
\(H_V\) rows, so the sums interchange.  For the latter, the absolute
double sum is bounded by
\[
 \sum_j\|H_V(b,j)\|\sum_{b'}\|K_V(j,b')\|\,\|f\|_\infty
 \le B_L\|K_V\|_{\rm Sch}\|f\|_\infty.
\]
It therefore interchanges as well.  This proves \eqref{eq:first-step-background-BGE6} on
\(\ell^\infty\), rather than merely extending an \(\ell^2\)
notation to that space.

We need nonlinear estimates for the actual local functions.  The
joint background and free radii \(r_b=c_bL^{-5}\),
\(r_u=c_uL^{-9}\) of \eqref{eq:first-step-nonflat-precision-NFP8} are retained.  In each local real
block-constant frame the face lift has small potential.  The actual
joint chart therefore exists on a fixed multiple of these radii,
with uniform local bounds.  Its component dependence is the same
as in the preceding section.  A fixed free coordinate enters at most
\(54\) composed Wilson terms and \(L^4-1\) actual GF terms, and its
coupled free coordinates lie in a fixed number of blocks.  Two scalar
free Cauchy formulas cost \(r_u^{-2}\), the term count costs
\(CL^4\), and the coupled-coordinate count costs \(CL^4\).  Hence,
on a fixed smaller ball,
\begin{equation}
 \begin{gathered}
 \|DF_V(u)\|_{\rm Sch}\le M_L,\qquad
 \|DF_V(u)-K_V\|_{\rm Sch}\le N_L\|u\|_\infty,\\
 M_L=C_*L^{26},\qquad N_L=C_*L^{35}.
 \end{gathered}
 \label{eq:first-step-background-BGE7}
\end{equation}
To obtain the second estimate, apply one more scalar Cauchy formula
in an arbitrary unit supremum free direction.  It costs
\(r_u^{-1}=CL^9\).  Integrate that derivative along the free
segment from zero to \(u\).  No count of all coordinates in the
direction is inserted.  All these Hessians have range \(12L\),
and their coefficients depend on background variables within a fixed
neighbourhood of the same size.  These are derived bounds for the
composed actual action and actual contour GF; a bound on their desired
matrix perturbation is not an input.

We will also vary the coarse datum.  Fix a real \(V\) as above, and
write
\begin{equation}
 \begin{gathered}
 V_B(c)=e^{iB(c)}V(c),\qquad
 U_0(V_B)=e^{iE_0(B)}U_0(V),\\
 E_0(B)(b)=
 \begin{cases}B(c),&b\hbox{ crosses face }c,\\0,&b\hbox{ is internal}.
 \end{cases}
 \end{gathered}
 \label{eq:first-step-background-BGE8}
\end{equation}
In particular \(\|E_0(B)\|_\infty\le\|B\|_\infty\).
The real local frames are fixed while \(B\) varies.  Write
\(F(B,u)=F_{V_B}(u)\).  On \(\|B\|_\infty<r_b/C_*\) and the
retained free ball, the actual chart and all component functions are
jointly holomorphic, and
\begin{equation}
 \begin{aligned}
 \|F(B,0)-F(0,0)\|_\infty&\le Q_L\|B\|_\infty,
       &Q_L&=C_*L^{18},\\
 \|D_uF(B,u)-K_V\|_{\rm Sch}
       &\le D_L\|B\|_\infty+N_L\|u\|_\infty,
       &D_L&=C_*L^{31}.
 \end{aligned}
 \label{eq:first-step-background-BGE9}
\end{equation}
For the first estimate the local scalar Cauchy costs are \(L^9\)
for one free derivative and \(L^5\) for one background direction,
followed by the \(L^4\) term count.  For the second, apply the
background Cauchy loss \(L^5\) to the \(L^{26}\) Schur estimate
and combine with \eqref{eq:first-step-background-BGE7}.  The face replication in \eqref{eq:first-step-background-BGE8} has supremum
norm one, so no number of crossing bonds is charged to the directional
background Cauchy radius.  Local real unitary frames preserve all
input norms.  These estimates are consequently uniform even when the
coarse holonomies are large.

\subsubsection*{Explicit radii and the real fixed point}

Enlarge \(S_L=C_*L^4\) so that on the retained chart
\begin{equation}
 \|\Phi_V(u)\|_\infty\le S_L\|u\|_\infty,
 \qquad
 \sup_p\|U_V(u)(\partial p)-I\|
           \le\epsilon+S_L\|u\|_\infty
 \label{eq:first-step-background-BGE10}
\end{equation}
for real \(u\).  The second estimate follows by telescoping four
unitary bond factors and using the first.  Set
\begin{equation}
 \begin{gathered}
 r_N=\min\left\{r_u/16,\ (16B_LN_L)^{-1},\
                         \alpha_L/(16S_L)\right\},\\
 \eta_{\rm reg}=\min\left\{\alpha_L/4,\ c_*L^{-5},\
                             r_N/(16C_{\rm cmp}L^3)\right\},\\
 P_L=1+2B_LS_LA_L,\qquad
 \epsilon_*=\min\left\{\alpha_L/16,\
                 r_N/(16B_LA_L),\eta_{\rm reg}/(16P_L)\right\}.
 \end{gathered}
 \label{eq:first-step-background-BGE11}
\end{equation}
Here \(C_{\rm cmp}\) is the universal comparison constant proved
below, and \(c_*\) is small enough for its local logarithm branches.
This is a sequential choice: \(r_N\) is fixed first, then
\(\eta_{\rm reg}\), then \(\epsilon_*\).  All are positive and
depend only on \(L\).

Suppose \(\epsilon\le\epsilon_*\).  On the closed bounded-field
ball \(\|u\|_\infty\le r_N\), let
\begin{equation}
 \mathcal T_V(u)=u-H_VF_V(u).
 \label{eq:first-step-background-BGE12}
\end{equation}
By \eqref{eq:first-step-background-BGE6}--\eqref{eq:first-step-background-BGE7},
\begin{equation}
 \|D\mathcal T_V(u)\|_{\infty\to\infty}
       \le B_LN_Lr_N\le1/16,
 \qquad
 \|\mathcal T_V(0)\|_\infty
       \le B_LA_L\epsilon\le r_N/16.
 \label{eq:first-step-background-BGE13}
\end{equation}
Thus \(\mathcal T_V\) maps the ball strictly into itself and is a
contraction.  Its iterates from zero are real and converge to a unique
root \(u(V)\) in that ball.  The two-sided identity in \eqref{eq:first-step-background-BGE6} makes
the fixed-point equation equivalent to \(F_V(u(V))=0\).  Moreover
\begin{equation}
 \|u(V)\|_\infty\le2B_LA_L\epsilon,
 \qquad
 \|\Phi_V(u(V))\|_\infty\le2S_LB_LA_L\epsilon.
 \label{eq:first-step-background-BGE14}
\end{equation}
The resulting real field has curvature at most
\(P_L\epsilon\le\eta_{\rm reg}/16\).  In particular it remains
inside the physical NFP curvature domain.  No critical point was
assumed to start this iteration.

\subsubsection*{Why the root is a Wilson critical background}

We first check the tangent space used by the stationarity equation.
On the retained joint domain, the pivot derivative of the actual
constraint stays invertible.  The chart has identity free restriction.
Differentiating its constraint and solving the pivot part therefore
shows that its derivative spans the entire tangent kernel of that
constraint, at every point in the ball.  The component derivative is
strictly local.  In particular a compactly supported full tangent
variation is represented by its compactly supported free restriction;
no tangent directions are omitted in \eqref{eq:first-step-background-BGE3}.

Let \(x\in B(y)\setminus\{y\}\).  Apply a site gauge supported at
\(x\), with \(s_t(x)=e^{itX}\), and identity at every other site.
It fixes all coarse centers, so the actual nonlinear average is
unchanged.  It is an admitted fibre variation for sufficiently small
\(t\).  The Wilson derivative along it is zero by gauge invariance.
Endpoint covariance of the actual contour mean gives
\begin{equation}
 T_{y,x}(U^{s_t})=T_{y,x}(U)e^{-itX}.
 \label{eq:first-step-background-BGE15}
\end{equation}
Every other internal mean is unchanged: any occurrence of \(x\)
inside a contour cancels between its two adjacent factors.  Consequently
stationarity of \eqref{eq:first-step-background-BGE3} at the root implies
\begin{equation}
 0=\frac i2\operatorname{tr}\big(T_{y,x}(U)X\big)
       \quad\hbox{for every traceless Hermitian }X.
 \label{eq:first-step-background-BGE16}
\end{equation}
The trace pairing then says that \(T_{y,x}(U)\) is scalar.  Its
determinant is one, so it is \(I\) or \(-I\).  The chart keeps all
internal bonds, their \(O(L)\)-length contour transports and their
means in the retained neighbourhood of \(I\); the choice \(r_N\le
r_u/16\) secures this uniformly.  The minus sign is excluded.  Thus
\begin{equation}
 T_{y,x}(U(V))=I\quad\hbox{for every internal pair},\qquad
 D\mathrm{GF}(U(V))=0.
 \label{eq:first-step-background-BGE17}
\end{equation}
The second identity holds on the full fine-bond space: the trace
derivative of an \(\operatorname{SL}(2)\) curve through \(I\) is zero.
The root equation is therefore Wilson stationarity along the entire
actual fibre.  We have constructed
\begin{equation}
 \boxed{\quad
 \mathcal M_LU(V)=V,\qquad U(V)\hbox{ is block axial},\qquad
 J_A(U(V))\perp\ker P_{U(V)}.
 \quad}
 \label{eq:first-step-background-BGE18}
\end{equation}
All arguments on the infinite lattice used compactly supported free
or site variations; they require no value for the background action.
Here the perpendicular notation means annihilation of compactly
supported tangent fields.  The bounded pointwise multiplier identity,
which implies no \(\ell^2\) membership, is proved in \eqref{eq:first-step-background-BGE34}--\eqref{eq:first-step-background-BGE35}.
At this actual critical background the physical gap and covariance
bounds of the preceding section now have their stated Gaussian
interpretation.

\subsubsection*{Uniqueness and the comparison with the face lift}

An axial real Wilson critical field represented by \(U_V(u)\) in
the closed \(r_N\)-ball has zero GF gradient and hence satisfies
\(F_V(u)=0\).  The contraction proves that it equals \(U(V)\).
The comparison in \eqref{eq:first-step-background-BGE11} identifies a geometric uniqueness class
without assuming chart membership.

Let \(W\) be a real block-axial field with fine plaquette deviation
at most \(\eta\).  On one block a unitary comb gauge \(g\), anchored
at its center, makes its internal bonds \(CL\eta\)-close to \(I\).
Every transformed internal path has length \(O(L)\), so its actual
mean is \(CL^2\eta\)-close to \(I\).  But axiality and covariance
give that mean as \(g(y)g(x)^{-1}=g(x)^{-1}\).  Thus the original
internal bonds obey
\begin{equation}
 \|W(b)-I\|\le CL^2\eta
       \quad(b\hbox{ internal}).
 \label{eq:first-step-background-BGE19}
\end{equation}

For a coarse face \(c\), write \(S_x=W([x,x+Le_\mu])\).
Coordinate-path swaps express each contour/reference loop in the
double block as a product of at most \(CL^2\) conjugated elementary
plaquettes.  Unitarity makes its deviation at most \(CL^2\eta\).
Two internal coordinate-order paths differ by the same number of
swaps, so their transports differ by that amount.  The contraction
for the implicit mean, recentered at any one unitary input, shows
that every such path transport differs from its actual mean by
\(CL^2\eta\).  Since those means are \(I\), replacing the two
contour transports in the closed loop by their means gives
\begin{equation}
 \|S_xW[c]^{-1}-I\|\le CL^2\eta.
 \label{eq:first-step-background-BGE20}
\end{equation}
The actual outer exponential of the arithmetic logarithm mean has
the same bound relative to \(W[c]\).  Hence
\begin{equation}
 \|S_x\mathcal M_L(W)(c)^{-1}-I\|\le CL^2\eta.
 \label{eq:first-step-background-BGE21}
\end{equation}
Each \(S_x\) contains one crossing bond and at most \(L-1\)
internal bonds.  Stripping those unitary internal factors with
\eqref{eq:first-step-background-BGE19} proves the safe bound
\begin{equation}
 \sup_b\|W(b)-U_0(\mathcal M_LW)(b)\|\le C_{\rm cmp}L^3\eta.
 \label{eq:first-step-background-BGE22}
\end{equation}
Enlarge \(C_{\rm cmp}\) to give the same estimate for the physical
left logarithm of the displayed ratio on its near-identity branch.
No global small gauge or small coarse potential was used.

If \(\mathcal M_LW=V\) and \(\eta\le\eta_{\rm reg}\), that
logarithm is within the ordinary local graph uniqueness radius and
its free restriction is within \(r_N/16\).  The actual fibre chart
therefore represents \(W\); if \(W\) is constrained-critical it
equals \(U(V)\).  For a real field not initially axial, small
plaquettes first make its internal contour families mutually close.
Their actual means exist, even if centered at a large unitary
transport.  The site gauge
\(s(x)=T_{y,x}(W)\), \(s(y)=I\), sets every mean to \(I\) and
preserves the coarse datum.  Curvature and constrained criticality
are preserved.  Applying \eqref{eq:first-step-background-BGE22} to this representative proves
uniqueness of the center-fixing gauge orbit among real regular
fields of curvature at most \(\eta_{\rm reg}\).
The constructed field belongs to this same class by \eqref{eq:first-step-background-BGE14} and
the final smallness choice in \eqref{eq:first-step-background-BGE11}.  Two axial representatives
in a center-fixing orbit are equal: covariance of their means forces
the site gauge itself to be identity.

\subsubsection*{Complex continuation and first response kernels}

Keep the real \(V\) fixed.  Enlarge a universal constant in
\(P_L'=C_*+2B_LS_LQ_L\), and set
\begin{equation}
 b_N=\min\left\{r_b/(16C_*),\ (16B_LD_L)^{-1},\
              r_N/(16B_LQ_L),\eta_{\rm reg}/(16P_L')\right\}.
 \label{eq:first-step-background-BGE23}
\end{equation}
For \(\|B\|_\infty<b_N\), use the same fixed real inverse
\(H_V\) in
\[
 \mathcal T_B(u)=u-H_VF(B,u).
\]
\eqref{eq:first-step-background-BGE9}--\eqref{eq:first-step-background-BGE13} give
\begin{equation}
 \|D\mathcal T_B(u)\|_{\infty\to\infty}\le1/8,
 \qquad
 \|\mathcal T_B(0)\|_\infty\le r_N/8.
 \label{eq:first-step-background-BGE24}
\end{equation}
The contraction is uniform on the complex ball and maps it strictly
inside.  Its iterates are holomorphic and converge uniformly on
smaller closed parameter domains.  Thus \(u(B)\), and
\(U(V_B)=e^{i\Phi_{V_B}(u(B))}U_0(V_B)\), are holomorphic in every
admitted finite complex parameter family.  Reality is preserved
when \(B\) is real.  The root satisfies
\begin{equation}
 \|u(B)\|_\infty\le2B_L(A_L\epsilon+Q_L\|B\|_\infty).
 \label{eq:first-step-background-BGE25}
\end{equation}
The site variation argument in \eqref{eq:first-step-background-BGE15}--\eqref{eq:first-step-background-BGE17} applies to a complex
root as well.  The equation pairs with every complexified traceless
direction; a scalar mean in \(\operatorname{SL}(2,\mathbb C)\) is
\(\pm I\), and the retained branch selects \(I\).  Hence this is
an actual complex block-axial Wilson critical background, rather
than a critical point only of \(A+\mathrm{GF}\).  No independent
fine current is introduced.

For completeness we derive response locality from actual finite
matrices.  At the root put
\(\mathcal H_B=D_uF(B,u(B))\).  The preceding estimates give
\begin{equation}
 \|\mathcal H_B-K_V\|_{\rm Sch}
       \le(8B_L)^{-1}\le c_0/64,
 \qquad \|\mathcal H_B\|_{\rm Sch}\le M_L.
 \label{eq:first-step-background-BGE26}
\end{equation}
They also retain range \(r_0=12L\).  Choose
\begin{equation}
 \sigma_L=r_0^{-1}\log\big(1+c_0/(64M_L)\big)>0.
 \label{eq:first-step-background-BGE27}
\end{equation}
Multiplication by \(e^{\sigma_L\varphi}\), for a bounded real
one-Lipschitz \(\varphi\), changes \(\mathcal H_B\) by at most
\(M_L(e^{\sigma_Lr_0}-1)=c_0/64\) in Schur norm.  A Neumann
argument about \(K_V\ge c_0/2\), first for the complex perturbation
and then for this conjugation, gives an inverse norm at most
\(8/c_0\).  Bounded weights
\(\varphi(b)=\min\{d(b,b'),N\}\), followed by \(N\to\infty\),
give the actual kernel bound
\begin{equation}
 \|\mathcal H_B^{-1}(b,b')\|
        \le(8/c_0)e^{-\sigma_Ld(b,b')}.
 \label{eq:first-step-background-BGE28}
\end{equation}
Absolute rows again give both inverse identities on bounded fields.
At a real root this matrix is a positive coordinate Hessian.  At
any root the physical free coordinates centered at that root differ
from the original free coordinates by pointwise invertible
three-by-three exponential-coordinate Jacobians.  Their norms and
inverse norms are uniformly bounded on the retained ball.  Since
the first derivative vanishes, the Hessians are related by ordinary
congruence with these pointwise matrices.  Thus the physical
centered Hessian also has the corresponding inverse kernel bounds.

Define the source derivative by first differentiating the first argument
of \(F(B,u)\) at fixed \(u\), and only then substituting the root:
\(\mathcal R_B=(\partial_1F)(B,u(B))\).  Thus the total derivative
of the root equation is
\(0=\mathcal R_B+\mathcal H_BD_Bu\).
This partial source derivative is strictly local.  A free row sees
only a fixed number of coarse faces in its
dependence neighbourhood.  The Cauchy count giving \(Q_L\) bounds
each coefficient by \(CL^{18}\), while a coarse face can reach at
most \(CL^4\) free coordinates.  Therefore
\begin{equation}
 \|\mathcal R_B\|_{\rm Sch}\le C_*L^{22},\qquad
 D_Bu=-\mathcal H_B^{-1}\mathcal R_B.
 \label{eq:first-step-background-BGE29}
\end{equation}
The support distance between a free bond and a contributing coarse
face, represented by its pivot, is at most a fixed multiple of \(L\).
Combining this fact with \eqref{eq:first-step-background-BGE28} gives
\begin{equation}
 \|D_{B(c)}u(b)\|\le C_L
                     e^{-\sigma_Ld(b,b_0(c))}.
 \label{eq:first-step-background-BGE30}
\end{equation}
The derivative \(D_u\Phi\) has the same strict locality and
absolute row and column bounds \(CL^4\) on the shrunken domain,
by one free Cauchy formula.  Its direct background derivative is
local and bounded by \(CL^9\|u\|_\infty\), using a background
Cauchy formula and \(\Phi_B(0)=0\).  The exponential map and the
direct face lift have bounded local derivatives.  Consequently the
actual fine-field derivative in left trivialization satisfies
\begin{equation}
 \left\|\frac1i(D_{B(c)}U(V_B)(b))U(V_B)(b)^{-1}\right\|
       \le C_L e^{-\sigma_Ld(b,b_0(c))}.
 \label{eq:first-step-background-BGE31}
\end{equation}
The constants in these response bounds are independent of volume
and holonomies.  The response formulas concern this first-step
solution; they are not many-level background derivative estimates.

\subsubsection*{Interior locality, periodicity and exact scope}

There is a useful finite comparison without interpolating between
two distant coarse connections.  Let \(V,V'\) be admitted real
coarse fields which agree on a coarse patch, and let \(u,u'\) be
their constructed free solutions.  Both arrays have the same free
coordinate labels and lie in the uniform \(r_N\)-ball.  The function
\(F_V(u')\) is defined even where the two face lifts differ.  Strict
locality gives \(F_V(u')=F_{V'}(u')=0\) on every row whose entire
coarse dependence neighbourhood lies in the common patch.  Call this
set of rows \(I\).  On the convex free segment define
\begin{equation}
 \overline{\mathcal H}
    =\int_0^1DF_V(u'+t(u-u'))\,dt.
 \quad\hbox{Then}\quad
 \overline{\mathcal H}(u-u')=-F_V(u').
 \label{eq:first-step-background-BGE32}
\end{equation}
Every integrand is a real symmetric range-\(12L\) matrix with the
same Schur bound and the same positive gap: it differs from \(K_V\)
by at most \(N_Lr_N\le c_0/128\).  Its average has those properties
as well, so the bounded-conjugation argument gives exponential
inverse kernels.  The residual in \eqref{eq:first-step-background-BGE32} is supported in \(I^c\)
and has norm at most \(2M_Lr_N\), by the derivative bound along the
same segment.  Summing the inverse with the smaller exponential
weight gives
\begin{equation}
 |u(b)-u'(b)|\le C_Lr_N
                  e^{-\sigma_Ld(b,I^c)/2}.
 \label{eq:first-step-background-BGE33}
\end{equation}
The local chart derivative gives the analogous bound for fine fields
whose background dependence neighbourhoods are in the common patch;
the fixed additional range is absorbed in \(C_L\).  No path from
\(V\) to \(V'\) through small-curvature coarse fields has been
assumed.  This is a deterministic locality statement for these
canonical first-step axial representatives.

The face lift, chart and actual GF intertwine every real unitary
gauge transformation constant on blocks.  The unique fixed point
therefore has the corresponding coarse gauge covariance.  The same
uniqueness proves compatibility with translations preserving the
block partition.  If the coarse input is periodic, translating the
cover root by a period gives another root in the same ball; it is
therefore periodic.  Conversely a finite-torus root lifts to such a
cover root, so the two constructions agree.  This uses locality and
uniqueness of the bounded-field equation, rather than silently
periodizing an \(\ell^2\) inverse.  All local frames and kernel
estimates are valid on the cover, and their constants do not depend
on the noncontractible holonomies.

\subsubsection*{Bounded currents and pointwise constrained stationarity}

Let \(U=U(V)\), and use the actual chart centered at \(U\) for this
paragraph.  Write \(P=P_U\), \(h=h_U\),
\(C=C_U=\iota-hP\iota\), and \(J=J_A(U)\).
Every component of \(J\) is bounded by the physical plaquette estimate.
The finite-range maps \(P,h,C\), and their transposes, act on bounded
fields by finite component sums.  In particular the coarse covector
\(\mu=h^{\mathsf T}J\) is pointwise well-defined and obeys
\begin{equation}
 \mu(c)=J_c(U)^{-\mathsf T}J(b_0(c)),\qquad
 \|\mu\|_\infty\le CL^3\|J\|_\infty,
 \quad J_c(U)=D_{a(b_0(c))}q_{U,c}(0).
 \label{eq:first-step-background-BGE34}
\end{equation}
The symbol \(J_c(U)\) here is the three-by-three pivot coefficient,
whereas \(J(b)\) is a fine Wilson current.

For each compactly supported free field \(v\), \(Cv\) is a compactly
supported tangent field.  Wilson stationarity gives
\(\langle J,Cv\rangle=0\); testing one free coordinate at a time
therefore gives \(C^{\mathsf T}J=0\) pointwise.  Set
\(R=J-P^{\mathsf T}h^{\mathsf T}J\).
Its free restriction is exactly \(C^{\mathsf T}J\), hence zero.
On a pivot \(b_0(c)\), no other constraint row contains that pivot,
and \eqref{eq:first-step-background-BGE34} gives
\[
 R(b_0(c))=J(b_0(c))-J_c(U)^{\mathsf T}
                              J_c(U)^{-\mathsf T}J(b_0(c))=0.
\]
Both restrictions vanish.  Thus the exact bounded-current statement is
\begin{equation}
 \boxed{\quad J_A(U)=P_U^{\mathsf T}h_U^{\mathsf T}J_A(U)
                   \quad\hbox{pointwise on fine bonds}.\quad}
 \label{eq:first-step-background-BGE35}
\end{equation}
Conversely, this pointwise identity annihilates every compactly supported
\(a\in\ker P_U\): both \(a\) and \(P_Ua\) have finite support,
so transposition is a finite-sum identity and
\(\langle J_A(U),a\rangle=\langle h_U^{\mathsf T}J_A(U),P_Ua\rangle=0\).
The multiplier is unique since \(Ph=I\).  These equivalences are
statements about bounded currents and compact variations; their proof
does not place \(J_A(U)\) or \(\mu\) in \(\ell^2\).

\subsubsection*{Local interactions and bounded-field differentiability}

We specify the derivatives used in the iteration, source response and
comparison arguments.  For a finite-support increment of free variables,
each changed pivot is in a fixed neighbouring double block.  Only finitely
many Wilson and GF terms change, so the relative value of
\(A+\mathrm{GF}\), its first variation, and its second variation are
ordinary finite sums.  The \(b\)-component of \(F(B,u)\) is the
derivative under a variation of the single free coordinate \(b\).
It is a function of finitely many local free and coarse coordinates.
This defines the same component even when \(u\) and the background
do not have finite support.  The row/column estimates previously proved
extend its derivative matrices to bounded fields.

Here is a uniform remainder check that proves actual Fr\'echet
differentiability on \(\ell^\infty\).  Put
\begin{equation}
 \rho_b=c_{\rm loc}r_b,\qquad \rho_u=r_u/4,\qquad
 G_L=C_*L^{13}.
 \label{eq:first-step-background-BGE36}
\end{equation}
Here \(c_{\rm loc}>0\) is a universal constant fixed by the
matrix-logarithm and local-frame domain estimates before \(C_*\)
is chosen.  Take \(C_*\) large enough that \(b_N\le\rho_b/4\)
and to cover all derivative constants below.  Each
component of \(F\) is holomorphic and bounded by \(G_L\) on the
joint local domain of radii \(2\rho_b,2\rho_u\).  Indeed one free
Cauchy derivative of a bounded local interaction costs \(CL^9\),
and at most \(CL^4\) terms enter that component.  The estimate is
independent of the number of coordinates elsewhere on the lattice.
The radii in \eqref{eq:first-step-background-BGE11} and \eqref{eq:first-step-background-BGE23}, with its reduced first entry, lie
strictly inside half these local radii.

At such an interior point define
\(\nu=\|\Delta B\|_\infty/\rho_b+
       \|\Delta u\|_\infty/\rho_u\le1/4\).
For each row apply scalar Cauchy to
\(t\mapsto F_b(B+t\Delta B,u+t\Delta u)\) on a circle of radius
\((2\nu)^{-1}\); the case \(\nu=0\) is immediate.  The circle
stays in the stated local domain.  Summing the scalar Taylor remainder
of order at least two, uniformly in the row, gives
\begin{equation}
 \begin{split}
 &\|F(B+\Delta B,u+\Delta u)-F(B,u)
       -(\partial_1F)(B,u)\Delta B
       -(\partial_2F)(B,u)\Delta u\|_\infty\\
 &\hspace{35mm}\le C G_L\nu^2.
 \end{split}
 \label{eq:first-step-background-BGE37}
\end{equation}
The first Taylor coefficient is a linear map of the increments because
each component is an ordinary finite-dimensional derivative.  Its
uniform Cauchy bound makes that map bounded on \(\ell^\infty\).
\eqref{eq:first-step-background-BGE37} is the required uniform Fr\'echet remainder.
Applying the same argument to the derivatives proves their continuity
and higher local analytic remainders.  It also applies to the chart
components, using their uniform bound \(S_L\|u\|_\infty\).
Consequently the iteration is a genuine bounded-field equation.
Its holomorphic iterates converge uniformly; their component Cauchy
integrals and uniform remainders pass to the limit.  The resulting
bounded-field maps are locally holomorphic, as are the finite parameter
families already stated.  No extensive action estimate or formal
derivative of an infinite action value is needed.

\subsubsection*{The regular minimum and its quantitative strictness}

For every real \(u\) in the whole closed \(r_N\)-ball, \eqref{eq:first-step-background-BGE7}
and the radius choice give the Hermitian form bound
\begin{equation}
 DF_V(u)\ge\left(c_0/2-N_Lr_N\right)I
           \ge c_0/4\,I.
 \label{eq:first-step-background-BGE38}
\end{equation}
In fact \(N_Lr_N\le(16B_L)^{-1}\le c_0/128\).
On a finite torus this proves uniform strict convexity of the actual
\(A+\mathrm{GF}\) in free coordinates.  On the infinite lattice
the same quadratic inequality applies to compact variations.

On a finite torus let \(W\) be any real fibre field of curvature
at most \(\eta_{\rm reg}\), with \(\mathcal M_LW=V\).
Its canonical gauge fixing, equal to identity at the coarse centers,
gives an axial representative \(W_{\rm ax}\).
By \eqref{eq:first-step-background-BGE22} it has actual coordinates \(u_W\) in the same ball.
Let \(u_*=u(V)\) and \(d=u_W-u_*\).
The full free segment \(u_*+sd\), \(0\le s\le1\), remains in
the convex ball.  The segment is used for \(A+\mathrm{GF}\);
axiality is needed at its two endpoints.  Since \(F_V(u_*)=0\),
Taylor's integral identity
and \eqref{eq:first-step-background-BGE38} give
\begin{equation}
 \begin{split}
 &[A+\mathrm{GF}](U_V(u_W))
                  -[A+\mathrm{GF}](U_V(u_*))\\
 &\quad=\int_0^1(1-s)
            \langle d,DF_V(u_*+sd)d\rangle\,ds
       \ge(c_0/8)\|d\|_2^2.
 \end{split}
 \label{eq:first-step-background-BGE39}
\end{equation}
Both endpoints are axial, so their actual GF values vanish.
Wilson gauge invariance gives \(A(W)=A(W_{\rm ax})\).  Therefore
\begin{equation}
 \boxed{\quad A(W)-A(U(V))\ge(c_0/8)\|u_W-u(V)\|_2^2.\quad}
 \label{eq:first-step-background-BGE40}
\end{equation}
Equality forces the canonical axial representatives to coincide,
so equality occurs precisely in the constructed center-fixing gauge
orbit.  The constructed field belongs to this regular class by
\eqref{eq:first-step-background-BGE11}--\eqref{eq:first-step-background-BGE14}.  Hence the construction gives its minimizing orbit
among the stated regular fibre fields; no minimum over fields outside
that class is claimed.

On the infinite lattice take instead a compact perturbation \(W\)
of \(U(V)\), still real, regular and in the same fibre.
Only finitely many blocks can have changed internal means.  Its
canonical axializing site gauge is therefore identity outside finitely
many blocks, because \(U(V)\) is already axial.  Thus
\(W_{\rm ax}-U(V)\), and hence \(u_W-u_*\), have finite support.
Chart locality makes the entire segment in \eqref{eq:first-step-background-BGE39} a variation in a
fixed finite fine neighbourhood.  All relative action values and
derivatives in that equation are finite interaction differences.
The same Taylor argument proves \eqref{eq:first-step-background-BGE40} with exactly that meaning of
the left-hand side.  It does not compare infinite extensive action
values.

\subsubsection*{Dependence on coarse data in the product topology}

For a free row \(b\), let \(\mathsf D_b\) be its finite set of
coarse face variables on which \(F_V(u)_b\) depends.  This set is
contained in a fixed number of neighbouring coarse blocks, for every
admitted \(u\).  For a finite coarse patch \(\mathcal P\), put
\begin{equation}
 I_{\mathcal P}=\{b:\mathsf D_b\subset\mathcal P\},\qquad
 \delta_{\mathcal P}=\sup_{c\in\mathcal P}\|V(c)-V'(c)\|,
 \qquad \delta_L=\min\{1/8,r_b/(8C_*)\}.
 \label{eq:first-step-background-BGE41}
\end{equation}
Suppose \(V,V'\) are both admitted real data and
\(\delta_{\mathcal P}\le\delta_L\).  Write \(u=u(V)\),
\(u'=u(V')\), and evaluate the latter array in \(F_V\).
It is in the common free ball even though the data can differ
arbitrarily outside \(\mathcal P\).

We justify a local source estimate without a patch-volume loss.
For \(c\in\mathcal P\) take the principal left logarithm
\(B(c)=i^{-1}\log(V'(c)V(c)^{-1})\); it has norm at most
\(C\delta_{\mathcal P}\).  Extend it by zero outside the patch.
On a row in \(I_{\mathcal P}\), strict locality identifies the
endpoint data with \(V'\).  Uniqueness of each actual local pivot
graph identifies that endpoint with the chart used for \(V'\),
independently of data outside its dependence set.
Fixed local real frames for \(V\)
keep the background interpolation \(e^{itB}V\) inside the actual
joint analytic chart.  This interpolation need not stay in the
admitted small-curvature class.  One background-direction and one
free-direction Cauchy formula, with \(u'\) held fixed, give
\begin{equation}
 \sup_{b\in I_{\mathcal P}}
       |F_V(u')_b-F_{V'}(u')_b|
          \le Q'_L\delta_{\mathcal P},\qquad Q'_L=C_*L^{18}.
 \label{eq:first-step-background-BGE42}
\end{equation}
The count is \(L^5\) for the background radius, \(L^9\) for the
free radius and \(L^4\) local terms.  The background direction has
supremum norm at most \(C\delta_{\mathcal P}\), so its number of
nonzero faces is absent from this estimate.

The same real averaged Hessian as in \eqref{eq:first-step-background-BGE32}, with the data \(V\)
fixed, satisfies
\begin{equation}
 \overline{\mathcal H}(u-u')=-F_V(u'),\qquad
 |F_V(u')_b|\le
 \begin{cases}
 Q'_L\delta_{\mathcal P},&b\in I_{\mathcal P},\\
 2M_Lr_N,&b\notin I_{\mathcal P}.
 \end{cases}
 \label{eq:first-step-background-BGE43}
\end{equation}
The interior estimate uses \(F_{V'}(u')=0\); the exterior estimate
uses \(F_V(u)=0\) and the uniform derivative bound along the free
segment.  No comparison of the distant coarse data is used there.
Define the finite weighted kernel-row constant
\begin{equation}
 \mathcal B_L=\frac8{c_0}
       \sum_{m\ge0}C_4(m+1)^3e^{-\sigma_Lm/2},\qquad
 C_{{\rm cont},L}=\mathcal B_L\max\{Q'_L,2M_L\}.
 \label{eq:first-step-background-BGE44}
\end{equation}
The conjugation proof applies to \(\overline{\mathcal H}\) by
\eqref{eq:first-step-background-BGE38}.  Its exponentially weighted absolute rows, multiplied by
the two parts of \eqref{eq:first-step-background-BGE43}, yield
\begin{equation}
 |u(b)-u'(b)|\le C_{{\rm cont},L}\delta_{\mathcal P}
      +C_{{\rm cont},L}r_N
               e^{-\sigma_Ld(b,I_{\mathcal P}^c)/2}.
 \label{eq:first-step-background-BGE45}
\end{equation}
If the complement is empty, its distance is infinity and the second
term is zero.  The estimate contains neither the patch cardinality
nor the torus volume.  Local chart derivatives give the analogous
fine-field estimate, with a fixed additional neighbourhood absorbed
in the constant and in the distance margin.

For clarity the admitted infinite-volume real data set is
\begin{equation}
 \mathcal A_L=\{V\in\operatorname{SU}(2)^{\mathcal E_{\rm coarse}}:
             \|V(\partial P)-I\|\le\epsilon_*
                   \text{ for every coarse plaquette }P\}.
 \label{eq:first-step-background-BGE46}
\end{equation}
The bond set is countable.  Each displayed condition is continuous
in finitely many product coordinates, so \(\mathcal A_L\) is closed
in the compact metrizable product configuration space.
Given a fixed free bond and an error tolerance, first choose a finite
patch with distance to its exterior large enough to make the second
term of \eqref{eq:first-step-background-BGE45} small.  Then constrain its finitely many face variables
so \(\delta_{\mathcal P}\) makes the first term small.  This proves
coordinatewise continuity of \(V\mapsto u(V)\) on \(\mathcal A_L\)
in its relative product topology.  Every fine output component depends
continuously on finitely many data and chart coordinates, so
\(V\mapsto U(V)\) is continuous in the fine product topology.
It is consequently Borel.  The Borel structures here are those of
the two countable product configuration spaces.

\subsubsection*{Compatibility of the complex germs}

For each admitted real \(V\), use the complex neighbourhood with
principal coarse coordinates
\(W=e^{iB}V\), \(\|B\|_\infty<b_N\).
The fixed-real-base construction gives a root \(u_V(W)\); denote
its associated fine field by \(\mathcal U_V(W)\).
Its improved bound follows directly from \eqref{eq:first-step-background-BGE24}:
\begin{equation}
 \|u_V(W)\|_\infty
      \le\frac{\|\mathcal T_B(0)\|_\infty}{1-1/8}
      \le r_N/7<r_N/4.
 \label{eq:first-step-background-BGE47}
\end{equation}
Consider a common complex datum
\(W=e^{iB}V=e^{iB'}V'\) in two such neighbourhoods.
The face lift \(U_0(W)\) is the same exact bond configuration in
both presentations.  Its internal bonds and their means are \(I\),
and its exact average is \(W\).  The free coordinates are the
principal small left bond logarithms relative to this same lift.

We check that the actual chart is shared.  At zero, all closed
products in its average are \(I\).  For the retained small full
chart variables, the internal path transports, the closed products,
and the relative coarse products stay in their same near-identity
logarithm neighbourhoods.  The face values of \(W\) and their
inverses are uniformly bounded, because \(W\) is in a retained
complex neighbourhood of a real unitary datum.  Internal means use
the unique near-identity implicit mean.  The closed and relative
logarithms use the unique principal branch there.  Unitary local
frames conjugate these branches; they do not select new ones.
It follows that both presentations have the same actual function
\begin{equation}
 \begin{gathered}
 q_W(a)=i^{-1}\log[\mathcal M_L(e^{ia}U_0(W))W^{-1}],\\
 P_W=Dq_W(0),\qquad R_W(a)=q_W(a)-P_Wa.
 \end{gathered}
 \label{eq:first-step-background-BGE48}
\end{equation}
Here \(h_W\) is the inverse of the same diagonal pivot coefficients,
inserted at the same pivots.  For a fixed free \(u\), either chart
satisfies the identical local equation
\begin{equation}
 a=C_Wu-h_WR_W(a),\qquad a|_{\rm free}=u.
 \label{eq:first-step-background-BGE49}
\end{equation}
The existing joint pivot contraction has Lipschitz constant less than
one on its retained full-variable neighbourhood.  Both chart outputs
are in that neighbourhood, since
\(\|a\|_\infty\le S_Lr_N\) and \(r_N\le r_u/16\).
Its uniqueness therefore proves equality of the two graphs for
\(\|u\|_\infty\le r_N\), component by component.  The resulting
equation \(F_W(u)=0\) is identical in the two constructions.

Both roots in \eqref{eq:first-step-background-BGE47} are in the entire \(r_N\)-ball on which either
fixed-real-base map is a contraction.  Either inverse \(H_V\) or
\(H_{V'}\) is two-sided on bounded fields, so its fixed points in
that ball are exactly the roots of this common equation.
Uniqueness at the common complex datum gives
\begin{equation}
 u_V(W)=u_{V'}(W),\qquad \mathcal U_V(W)=\mathcal U_{V'}(W).
 \label{eq:first-step-background-BGE50}
\end{equation}
The uniform local Fr\'echet estimates and holomorphic contractions
then glue these continuations on overlaps.  This uses the common
complex datum and the actual pivot-graph uniqueness; it does not
assume that the overlap contains an admitted real point or extend
the argument to unrelated average maps or logarithm branches.

We have proved existence, a quantitative regular minimizing orbit,
bounded pointwise current stationarity, product-topology continuity,
compatible complex continuation and exponential first responses for
an actual first-step background. The
\hyperlink{first-step-named-form}{named comparison below}
checks independent currents on this first-step axial locus.
Nothing here constructs a many-level background, treats signed
layer coefficients or complex decoupling arrays, or proves the
conditional exterior-force, history-class or interacting
Gaussian-law assertions.


% END MODULE chapters/first_step_background.tex

% BEGIN MODULE chapters/domain_core_ledger.tex
\subsection{Domains of local substitutions}
\label{ch:domain-core-ledger}

The ordinary step uses three coordinate systems at once. A bond logarithm
is a group coordinate, a potential is a group coordinate divided by a
length, and a Gaussian coordinate is a physical logarithm divided by a
coupling. Its analytic domains also have two distinct parts: bounds on a
field and bounds on every background obtained from its restricted
averages. We fix these conventions and give quantitative admission tests.
The tests apply to the nonlinear local maps themselves.

The finite averaging chart, the classical flat comparison, the first-step
physical Hessian and the first-step background construction are used in
their stated scopes. The last construction includes its regular minimum,
product continuity and complex gluing conclusions. The many-level
background and response estimates are not conclusions of this section.
They occur below as explicitly named domain estimates. Once those estimates
hold, all the inclusions, changes of units and Cauchy estimates used here
are proved with constants independent of depth and compatible volume.

\subsubsection{Lengths, covectors and quadratic coefficients}

Work in the units of a fixed level. The bare spacing is \(\eta\), with
\(0<\eta\le1\); a bond configuration is written
\begin{equation}
 \mathbf U(b)=e^{ia(b)}U(b)=e^{i\eta A'(b)}U(b),
 \qquad a=\eta A'.
 \label{dc:bond-potential}
\end{equation}
Here \(U\) is real. The independent analytic current will be denoted by
\(\mathbf J\). The Wilson gradient is denoted by \(J_A(U)\) when its
normalization has been specified. These are different variables.
The color pairing is \((X,Y)=\frac12\operatorname{tr}(XY)\).
It is complex bilinear. Operator norms on complex coordinate spaces use
the corresponding Hermitian norm, and therefore use adjoints rather than
bilinear transposes in positivity arguments.

Suppose that the current density is normalized by
\begin{equation}
 DA_W(U)[\eta A]
   =\eta^4\sum_b(J_A(U)(b),A(b)).
 \label{dc:current-density}
\end{equation}
The covector for the raw logarithm \(a\), with counting pairings, is
\(\jmath_A(U)=\eta^3J_A(U)\). Indeed substitution of \(A=a/\eta\)
in \eqref{dc:current-density} gives \(DA_W(U)[a]=\sum_b(\eta^3J_A(U)(b),a(b))\).
At unit spacing, these two current conventions coincide. The formula in \eqref{dc:current-density}
also explains the dimension three weight on a current: the action has
dimension zero, while a potential has dimension one in four dimensions.

More generally let a linear response \(H\) map datum coordinates to field
coordinates with positive diagonal pairing matrices \(W_r,W_f\). Its
bilinear transpose and its transpose after length normalization are
\begin{equation}
 H^{\mathsf T}=W_r^{-1}H^{\mathrm t}W_f,
 \qquad P=\mathsf m Q,
 \qquad P^{\mathsf T}=Q^{\mathsf T}\mathsf m^{\mathsf T}.
 \label{dc:weighted-transpose}
\end{equation}
The matrix \(\mathsf m\) multiplies a row by its length \(\xi_c\).
This formula is valid for arbitrary specified row weights; no row-volume
factor is inferred from the letter \(Q\). To verify it, write
\(u^{\mathrm t}W_fHv=(H^{\mathrm t}W_fu)^{\mathrm t}v\) and compare
with the row pairing. Thus every transpose in a field-current identity
must refer to its two declared pairings. Counting row pairings and
volume field pairings can be used together, provided \eqref{dc:weighted-transpose} is used.

Let \(h\) be a right inverse for the raw constraint derivative \(P\).
The right inverse for the per-length derivative \(Q\) is \(h_Q=h\mathsf m\).
For the corresponding raw and per-length multipliers one obtains
\begin{equation}
 Ph=1,\quad Qh_Q=1,\qquad
 \mu=h^{\mathsf T}J,\quad
 \mu_Q=h_Q^{\mathsf T}J=\mathsf m^{\mathsf T}\mu.
 \label{dc:multiplier-length}
\end{equation}
If \(q_Q=\mathsf m^{-1}q\), then
\(\langle\mu_Q,D^2q_Q\rangle=\langle\mu,D^2q\rangle\).
Consequently changing a constraint from raw to per-length coordinates
changes its multiplier at the same time. It does not change the physical
curvature pairing.

For a change of length unit, let the new unit be \(L^{-n}\) times the old
one. Numerical lengths multiply by \(L^n\). The same configuration and
current are represented by
\begin{equation}
 \eta_{\rm new}=L^n\eta_{\rm old},\quad
 A_{\rm new}=L^{-n}A_{\rm old},\quad
 J_{\rm new}=L^{-3n}J_{\rm old},\quad
 \nabla_{\rm new}A_{\rm new}=L^{-2n}\nabla_{\rm old}A_{\rm old}.
 \label{dc:unit-change}
\end{equation}
The first identity and \(a=\eta A\) prove the second. A difference
quotient contains one inverse numerical length, which proves the fourth.
The third follows either from \eqref{dc:current-density}, or from the dimension three of its
covector. In particular
\(\eta^4\langle J,A\rangle\) is unchanged. If a bond lies in a layer of
local scale \(\xi\), the quantities
\begin{equation}
 \xi|A|,\qquad \xi^2|\nabla_UA|,\qquad
 \xi^3|J|,\qquad (\xi/\eta)^2|\partial U-1|
 \label{dc:dimensionless-tests}
\end{equation}
are unchanged as well. These are the quantities that enter a layered
regularity test. Applying the same radius to their unweighted numerical
coordinates would describe a different domain.

At a step whose fixed reference inverse coupling is \(\mathsf t\), put
\(g=\mathsf t^{-1/2}\). If \(y\) is the Gaussian free coordinate and
\(\Phi_U\) the actual physical fibre chart, the configuration is
\begin{equation}
 U_y=e^{i\Phi_U(gy)}U,
 \qquad |y|<\mathfrak b,\qquad
 \delta=g\mathfrak b.
 \label{dc:gaussian-coordinate}
\end{equation}
Thus a physical free-coordinate radius \(r\) becomes \(r/g\), before
allowances for the norm of the chart and the background response. The
chart is applied to \(gy\), not to \(y\).

For a holomorphic scalar functional \(F(B)\) of a physical datum, let
\(F_s(B_s)=F(gB_s)\). The chain rule gives
\begin{equation}
 D^2F_s(0)=g^2D^2F(0),\qquad
 |b(F_s)|\le CM(g/r)^2
 \ \Longleftrightarrow\ |b(F)|\le CMr^{-2}.
 \label{dc:quadratic-conversion}
\end{equation}
Here the last equivalence means the corresponding bounds under the same
linear coefficient readout. It is not an assertion that a general
coefficient is saturated by its Cauchy bound. The free-response and
spatial-moment sums in that readout have their own fixed constants.
In particular a radius proportional to \(g(\log g^{-2})^q\) is a growing
logarithmic radius in the Gaussian variable. Its physical extraction
factor remains \(g^{-2}(\log g^{-2})^{-2q}\).

\subsubsection{The actual average and its three linearizations}

The internal transport \(T_{y,x}(U)\), \(x\in B(y)\), is the mean of
the transports along coordinate-order shortest paths from \(y\) to \(x\).
There is one path for each ordering of the four directions, with
multiplicity; there are at most \(4!\) such paths. It is the solution near
that family of
\(\sum_\Gamma\log(U(\Gamma)T_{y,x}^{-1})=0\).
It is not an arithmetic average of matrices. With \(x'=x+c_+-c_-\), set
\begin{equation}
 \begin{aligned}
 W_x(U)&=T_{c_-,x}(U)U([x,x'])T_{c_+,x'}(U)^{-1}U[c]^{-1},\\
 \mathcal M(U)(c)&=
 \exp\left(L^{-4}\sum_{x\in B(c_-)}\log W_x(U)\right)U[c].
 \end{aligned}
 \label{dc:actual-average}
\end{equation}
This formula uses the regular local logarithm. The total measurable
average also has its prescribed fallback rules. Those rules make real
measure identities meaningful on every configuration. Cauchy estimates
below use only the analytic extension of \eqref{dc:actual-average}, with all physical
substitution values remaining in its regular branch.

At the identity let \(\overline a(y,x)\) be the arithmetic mean of the
signed path sums. Define the raw straight and plain derivatives by
\begin{equation}
 \begin{aligned}
 (S_La)(c)&=L^{-4}\sum_{x\in B(c_-)}a([x,x']),\\
 (\Sigma_La)(y)&=L^{-4}\sum_{x\in B(y)}\overline a(y,x),\\
 P_La&=S_La-\mathrm{d}_c\Sigma_La,\qquad
 (q_L\lambda)(y)=L^{-4}\sum_{x\in B(y)}\lambda(x),
 \end{aligned}
 \label{dc:linear-averages}
\end{equation}
where \((\mathrm{d}_cf)(c)=f(c_+)-f(c_-)\).
The sign in \eqref{dc:linear-averages} follows by differentiating the four transports in
\eqref{dc:actual-average}. For a gradient \(a=\mathrm{d}\lambda\), every path sum telescopes.
Consequently
\begin{equation}
 \Sigma_L\mathrm{d}\lambda=q_L\lambda-\lambda|_{\rm centers},
 \qquad S_L\mathrm{d}\lambda=\mathrm{d}_cq_L\lambda,
 \qquad P_L\mathrm{d}\lambda=\mathrm{d}_c(\lambda|_{\rm centers}).
 \label{dc:two-gradient-identities}
\end{equation}
This proves both gradient identities directly. The scalar site average
and restriction to centers have different kernels.

The straightening relative to a real background \(U_0\) removes a block
factor. In raw coordinates its one-step definition is
\begin{equation}
 \begin{aligned}
 u_y(a)&=\exp\left(L^{-4}\sum_{x\in B(y)}
       \log[T_{y,x}(e^{ia}U_0)T_{y,x}(U_0)^{-1}]\right),\\
 e^{i\mathcal Q^1(U_0,a)(c)}\mathcal M(U_0)(c)
   &=u_{c_-}(a)^{-1}\mathcal M(e^{ia}U_0)(c)u_{c_+}(a).
 \end{aligned}
 \label{dc:straightening}
\end{equation}
For several steps this definition is iterated, including the accumulated
node factors. Differentiating at a flat background removes precisely the
\(\mathrm{d}_c\Sigma_L\) term and gives the straight derivative. For a general
background, the resulting derivative is the straightened derivative
\(Q^\sharp\). Its leading transported straight derivative is a third
operator, \(Q^{\rm st}\). The plain derivative \(Q^{\rm pl}\),
\(Q^\sharp\), and \(Q^{\rm st}\) must therefore have separate names.

In physical potential coordinates, the derivative also contains the input
length factor. After \(k\) straight steps at the identity it is
\(\eta_k S_L^k\), where \(\eta_k=L^{-k}\). Its row sum is one, because
the raw straight row sum is \(L\) at every step. With covariant difference
quotients and the appropriate coarse units, differentiation of gauge
covariance gives
\begin{equation}
 Q^{\rm pl}_k(U)D_U\lambda
      =D_{\mathcal M^kU}(\lambda|_{\rm centers}),
 \qquad
 Q^\sharp_k(U)D_U\lambda
      =D_{\mathcal M^kU}Q'_k(U)\lambda.
 \label{dc:covariant-intertwining}
\end{equation}
The first formula follows from the exact bi-equivariance of
\(\mathcal M\), including its node endpoints. The second follows from
the same calculation after \eqref{dc:straightening}; its site factor is the derivative of
the averaged gauge transformation. A transported straight main term
does not satisfy the second identity exactly unless the omitted
straightening correction is zero.

There is a quantitative reason to keep these distinctions even at a flat
field. For the raw straight map, every fixed fine bond occurs in exactly
\(L\) of the intervals \([x,x+Le_\mu]\), counted over all block starts.
Its unsigned column sum is \(L^{-3}\). Its row sum is \(L\).
In contrast choose the first internal bond from a block center in direction
\(\mu\). Among the path orderings, at least a fraction \(1/4\) start in
that direction. There are \((L-1)L^3/2\) endpoints with positive
\(\mu\)-coordinate. Therefore its coefficient in \(\Sigma_La(y)\) is at
least
\begin{equation}
 L^{-4}\frac14\frac{L-1}{2}L^3
      =\frac{L-1}{8L}\ge\frac1{12}\qquad(L\ge3).
 \label{dc:plain-column-example}
\end{equation}
Take an outgoing coarse bond in another direction. Its straight part
does not see this fine component, and its plain derivative contains this
\(\Sigma_L\) coefficient. Hence the plain column bound is not the
straight column bound. Zero displacement in some path directions causes
no problem, since the coordinate orders form a multiset. This calculation
uses a nonwrapping local lift. Periodic tori large enough for the same
double block have the identical calculation.

On curvature quotients the difference in \eqref{dc:linear-averages} is a gradient, so the
classical energy comparison can transfer after an exact gradient
correction. That quotient statement does not identify the full gauge,
current or node-factor operators. In particular a column estimate for
one of those operators is an estimate in its declared coordinates, not
an estimate for an unnamed map denoted by \(Q\).

\subsubsection{The fibre chart, current subtraction and contour gauges}

On the regular logarithm branch, define
\[
q_U(a)=i^{-1}\log[\mathcal M(e^{ia}U)\mathcal M(U)^{-1}].
\]
Put \(P_U=Dq_U(0)\). The pivot construction
gives a right inverse \(h_U\), supported on the crossing pivots. If
\(\iota\) inserts the free coordinates and leaves pivots zero, set
\begin{equation}
 C_U=\iota-h_UP_U\iota,\qquad
 q_U(\Phi_U(u))=0,\qquad
 \Phi_U(0)=0,\quad D\Phi_U(0)=C_U.
 \label{dc:chart-identity}
\end{equation}
The chart has identity restriction to the free coordinates; its nonlinear
correction is on the pivots. Two differentiations of the constraint give
\begin{equation}
 D^2\Phi_U(0)[u,v]
      =-h_UD^2q_U(0)[C_Uu,C_Uv].
 \label{dc:chart-second}
\end{equation}
Indeed its image is on the pivots, where \(P_U\) has inverse \(h_U\).
This also proves the sign without appealing to a convention for a
``quadratic part.'' If \(q_U^{(2)}\) is the homogeneous quadratic
polynomial, then \(D^2q_U(0)[a,a]=2q_U^{(2)}(a)\).

Let \(\mathscr A_U\) denote the raw Wilson Hessian, and let \(J\) be a
raw fine covector in this paragraph. The explicitly defined auxiliary
form and its physical value are
\begin{equation}
 \begin{aligned}
 K_U^{\rm aux}[J](u,v)
   &=\langle C_Uu,\mathscr A_UC_Uv\rangle
      +G_U^{\rm abs}(u,v)
      -\langle h_U^{\mathsf T}J,D^2q_U[C_Uu,C_Uv]\rangle,\\
 K_U^{\rm phys}&=K_U^{\rm aux}[\jmath_A(U)].
 \end{aligned}
 \label{dc:auxiliary-form}
\end{equation}
Here \(G_U^{\rm abs}\) is the Hessian of the absolute contour gauge
composed with the chart; pivot avoidance also identifies it with
\(C_U^{\mathsf T}(\operatorname{Hess}{\rm GF})C_U\).
The equality is the ordinary chain rule, since
the pivots occur in no internal contour path and hence
\(h_U^{\mathsf T}\jmath_{\rm GF}=0\). It does not require the starting
field to be a critical point. Criticality is used for the stationary Gaussian interpretation.
The proved small-curvature physical gap also applies to the noncritical
starting face lift used to construct the first-step background.

For clarity the two contour functionals are
\begin{equation}
 \begin{aligned}
 \mathrm{GF}(V)&=\sum_{y,x\ne y}\left(1-\tfrac12\operatorname{tr}T_{y,x}(V)\right),\\
 \mathrm{GF}_{{\rm rel},U}(a)&=\sum_{y,x\ne y}
 \left(1-\tfrac12\operatorname{tr}
       [T_{y,x}(e^{ia}U)T_{y,x}(U)^{-1}]\right).
 \end{aligned}
 \label{dc:two-gauges}
\end{equation}
The second has zero first derivative at every \(U\). If \(G_Ua\)
is the first derivative of the relative contour logarithm, its Hessian
is \(G_U^{\mathsf T}G_U\), with the declared color pairings.
This follows by expanding \(1-\frac12\operatorname{tr}e^{iX}
=\frac12|X|^2+O(|X|^4)\). The first functional has that Hessian only
on the axial locus \(T_{y,x}(U)=1\). For example the second derivative
of \(1-\frac12\operatorname{tr}(e^{it\sigma_3}e^{i\vartheta\sigma_3})\)
is \(\cos\vartheta\), whereas the relative expression has second
derivative one. Absolute gauge is invariant under block-constant gauges.
The relative function is jointly invariant under arbitrary site gauges:
the product of its two contour factors transforms by conjugation at
\(y\). These are different covariance statements.

The gauge-restoration term in a named propagator requires an additional
check. Let \(Q'\) be the assigned average of site functions,
\(\Delta_{\rm site}=D^{\mathsf T}D\), and
\(G'=(\Delta_{\rm site}+Q'^{\mathsf T}aQ')^{-1}\), with the
specified positive site weight \(a\). The real gauge projection \(R\)
projects onto \(\Delta_{\rm site}\ker Q'\); on the admitted complex
domain the same operators use their bilinear analytic formulas.
Their existence on that domain is part of the primitive operator
requirement. Put \(S=1-DG'RD^{\mathsf T}\). On physical currents its
projected Hessian is expressed by
\begin{equation}
 \Delta_\pi=S^{\mathsf T}\Delta^{\rm W}S,\qquad
 \Delta'_\pi=\Delta^{\rm W}-\Delta_\pi.
 \label{dc:physical-gauge-restoration}
\end{equation}
Thus \(R=0\) makes this physical correction zero. For an
independent-current completion \(\Delta'_\pi[J]\), it is necessary to
inspect its actual formula. If it is obtained by expanding \eqref{dc:physical-gauge-restoration} and
extending its current factors without removing the factors of \(R\),
then it is zero when \(R=0\) for every \(J\): every term of
\(\Delta^{\rm W}-(1-T)^{\mathsf T}\Delta^{\rm W}(1-T)\),
\(T=DG'RD^{\mathsf T}\), contains a \(T\). This is a sufficient
algebraic criterion, not a conclusion from agreement on the physical
graph alone.

Normalize every quadratic matrix by the exponent
\(-\frac12\langle u,Ku\rangle\). At the first step, where the
zero-fold average is the identity, the difference between a named
completion and \eqref{dc:auxiliary-form} has the complete form
\begin{equation}
 \begin{split}
 \mathcal D_U[J]={}&C_U^{\mathsf T}
   (\mathscr G_U^{\rm named}-\mathscr G_U^{\rm abs}
                         -\Delta'_\pi[J])C_U\\
 &-C_U^{\mathsf T}E_U^{(2)}[J]C_U+\mathcal E_U[J],\\
 \mathcal E_U[J](u,v)&=
   \langle h_U^{\mathsf T}J,D^2q_U[C_Uu,C_Uv]\rangle.
 \end{split}
 \label{dc:complete-comparison}
\end{equation}
The notation for the current readout inside \(E^{(2)}\) must first be
converted by \eqref{dc:weighted-transpose}--\eqref{dc:multiplier-length}. The formula in \eqref{dc:complete-comparison} is obtained by subtracting
the two quadratic forms term by term. Matching just its last line is
the constraint-readout condition; it is not the complete comparison.
In particular a relative gauge prescription at a seam-free stage must
be reconciled with the named matrix used at that stage.

\paragraph{Domain estimate I: the named completion.}
For the actual regular pair on each required region, compute all terms
of \eqref{dc:complete-comparison} with one declared pairing and Hessian normalization. Prove
either \(\mathcal D_U[J]=0\), or an explicit local coefficient estimate
small enough for the available gap. The quantifiers include independent
complex \(J\) on its own radius, the chosen gauge representatives and
all admitted local restrictions. The physical identity in \eqref{dc:auxiliary-form} and the
first-step physical gap already hold independently of this estimate.
The calculation at zero inner depth can be done directly.

%
%
\paragraph{The named form at the first step.}
\hypertarget{first-step-named-form}{}
Take the unweighted bare torus with no lower layers, and use
bare units. In the lower sequence all bond and site rows
are the zero-fold rows \(Q=1\) and \(Q'=1\). The Landau
projection \(R\) is zero. In the definition of the
gauge-restoration term, the gauge read is
\(\mathsf V A=G'RD^{\mathsf T}A\). Its vanishing gives
\(\Delta'_\pi[J]=0\) directly from the defining quadratic
form, for independent \(J\) as well as for the physical
current. The lower straightened constraint
\(\mathcal Q_0(U;A)=A\) is linear. Its quadratic correction
\(\Delta^{(2)}\) is therefore zero. With the bare row weight
one, the ordered inverse formulas give
\[
 \mathsf G_1^{-1}=\mathscr A_U+1,\qquad
 \mathsf H_1=1,\qquad
 (Q\mathsf G_1Q^{\mathsf T})^{-1}-1=\mathscr A_U .
\]
These identities are on the domain of the named inverses.

The outer row is still the genuine one-step average. The
source's per-length derivative is \(L^{-1}Dq_U(0)\), so at
a pivot its coefficient is \(L^{-1}\) times that of the raw
derivative. The pivot operator with coefficient
\(L^{-1}\kappa_U(c)^{-1}\) is exactly \(h_U\) in raw
constraint coordinates. The polarized quadratic coefficient
of the raw logarithm is \(\tfrac12 D^2q_U(0)\).
In the bare pairing the lower current read
\(\mathsf H_1^{\mathsf T}J\) is \(J\) itself, restricted to
the bonds of the next step. The named correction pairs twice
that polarized coefficient, after \(h_U\), with this current.
Consequently
\[
 \langle C_Uu,\mathsf E^{(2)}[\mathsf H_1^{\mathsf T}J]C_Uv\rangle
  =\langle h_U^{\mathsf T}J,D^2q_U(0)[C_Uu,C_Uv]\rangle.
\]
The named gauge term is the relative contour Hessian in
\eqref{dc:two-gauges}. On the axial locus the base contour
values equal one, so the relative and absolute
functionals agree there as functions of the perturbation.
With the common free chart, all terms of the named form
therefore agree with \eqref{dc:auxiliary-form} on this
first-step locus for independent admitted currents. Away
from that locus the calculation gives the named relative
completion, with the absolute-to-relative difference
explicitly displayed in \eqref{dc:complete-comparison}.

\subsubsection{Triples, simple bounds and hereditary cores}

Write \(\eta_j=L^{-j}\) in level-\(j\) units. A core is initially a
domain of triples \((U,A',\mathbf J)\), with \(U\) real and
\(\mathbf U=e^{i\eta_jA'}U\). For fixed \(U\), analyticity means
holomorphy in \((A',\mathbf J)\). On a single layer its simple conditions
at radii \((b_0,b_1)\) are
\begin{equation}
 \begin{gathered}
 |\partial U-1|,\ |\partial\mathbf U-1|<b_0\eta_j^2,
 \qquad |A'|,\ |\nabla_UA'|<b_1,\qquad |\mathbf J|<b_0,\\
 U=e^{i\eta_jA}\text{ in real gauges on the specified }O(LM)
 \text{ cubes},\qquad |A|,|\nabla A|<C_gM b_0.
 \end{gathered}
 \label{dc:simple-core}
\end{equation}
A convention with \(LM\) instead of \(M\) in the last constant is
absorbed in \(C_g(L)\), not in a level-dependent radius. On a layered
domain the inequalities are written in the dimensionless quantities
\eqref{dc:dimensionless-tests}, with the assigned radii of that layer.

For every subdomain \(Z\) of level \(p\le j\), form each of the two
assigned eroded towers. The first uses its prescribed cube size and
outer-unit gap, the second its prescribed minimal determining size and
inner-unit gap. If \(\mathbf B'\) is the determining set of a truncated
tower, its configuration is the background continued from the real
datum, evaluated at
\begin{equation}
 B_Z=\frac1i\log\left[
 M_{\mathbf B'}(e^{i\eta_jA'}U)M_{\mathbf B'}(U)^{-1}\right].
 \label{dc:tower-datum}
\end{equation}
This definition keeps the real base \(U\). Existence of a continuation
from one base does not by itself identify it with a continuation from
another base representing the same complex pair.

Let \(Z^{\ominus2}\) mean removal of the prescribed two cube layers.
The hereditary clauses, for every truncation \(n\le p\) of both rules,
are existence and the bounds
\begin{equation}
 \begin{aligned}
 |\partial U_n(Z;\mathbf U)-1|
       &<b_0L^{-(j-p)}\eta_p^2,\\
 |J_A(U_n(Z;\mathbf U))|
       &<b_0L^{-(j-p)}L^{2(n-p)}
       \qquad\text{on }Z^{\ominus2}.
 \end{aligned}
 \label{dc:tower-clauses}
\end{equation}
The current in the second line is formed in level-\(n\) units.
The same clauses are imposed on the real factor \(U\). This current
is not the independent argument \(\mathbf J\) of an activity. The
two kinds of current enter different tests.

We denote the resulting hereditary core by
\(\mathcal C_j(X;b_0,b_1)\). A local function descends from triples to
pairs only when its actual definition gives the same value on two
admitted splittings of that pair. In collar expressions this requirement
may belong to a specified difference of terms rather than each term.
Restriction to a smaller set is permitted only with its required collars
and determining neighborhoods. In particular it is not an operation that
discards some of the quantifiers in \eqref{dc:tower-clauses}.

The linear regular-pair class has a related but different convention.
On a layer of scale \(\xi\) it asks for curvature
\(\alpha_0\eta^2\xi^{-2}\), imaginary potential and first difference
\(\alpha_1\xi^{-1},\alpha_1\xi^{-2}\), independent current
\(\gamma_0\xi^{-3}\), and real gauges on \(O(M_1\xi)\) cubes with
\(\xi|A|,\xi^2|\nabla A|\le\alpha_0\). If those cubes fit into the
core's gauge cubes, the simple part of \eqref{dc:simple-core} enters this linear class
whenever
\begin{equation}
 b_0\le\min\{\alpha_0,\gamma_0\},\qquad
 b_1\le\alpha_1,\qquad C_gM b_0\le\alpha_0.
 \label{dc:core-linear-class}
\end{equation}
These inequalities follow by substitution in the four tests. The last
one is essential: the linear class's gauge bound is \(\alpha_0\) on
the entire cube, not \(M\alpha_0\). The required containment of gauge
cubes, and their nonwrapping torus lift, are geometric conditions that
must be checked before this inclusion is used.

\subsubsection{Quantitative room inside a core}

Here is a precise meaning of spending a fraction of a core. Divide each
left side of \eqref{dc:simple-core} and \eqref{dc:tower-clauses} by its stated radius and scale factor.
Include the corresponding real-gauge bounds and the datum-radius tests
for each atomic continuation used in \eqref{dc:tower-datum}. Call the resulting
nonnegative tests \(N_\beta\). A central core of fraction \(\vartheta\)
means \(N_\beta\le\vartheta<1\) for every test, on the same declared
atomic domains. At fixed finite world there are finitely many tests,
although their number need not have a useful bound independent of depth.

\textbf{Margin lemma.} Suppose a proposed substitution \(p_z\) has all
its atomic maps defined on a joint parameter domain and
\begin{equation}
 N_\beta(p_0)\le\vartheta,\qquad
 |N_\beta(p_z)-N_\beta(p_0)|\le\epsilon
       \quad\text{for every }\beta,\qquad \vartheta+\epsilon<1.
 \label{dc:margin-lemma}
\end{equation}
Then the substitution belongs to the core throughout that domain. This
is immediate for each test, and the same bound holds after taking their
supremum. No factor counting the tests occurs. The hard part of an
application is proving the uniform test difference, including every
normalized tower test. Simple-condition bounds alone are not that proof.

Near the identity a useful inclusion follows from a precise tower
estimate. Put
\(\|v\|_1=\max(\|A'\|_\infty,\|\nabla A'\|_\infty,
\|\mathbf J\|_\infty)\). The plaquette expansion gives
\begin{equation}
 |\partial e^{i\eta A'}-1|
   \le C\eta^2(\|\nabla A'\|_\infty+\|A'\|_\infty^2).
 \label{dc:plaquette-bound}
\end{equation}
To prove it, the terms linear in the four exponentials add to the
discrete curl times \(i\eta^2\). The remaining series has at least two
factors \(\eta A'\) and is bounded by
\(C\eta^2\|A'\|_\infty^2\) on a fixed small chart. This proof is
uniform as \(\eta\downarrow0\). A supremum bound on \(A'\) without
its difference quotient would instead allow a linear term of order
\(\eta\|A'\|_\infty\).

\paragraph{Domain estimate II: uniform atomic towers.}
\Cref{htc:all-towers} proves that simple bounds at sufficiently small
radii $(e_0,e_1)$ give every continued tower in
\eqref{dc:tower-datum}, for both erosion rules and every truncation.
The tests in \eqref{dc:tower-clauses}, divided by their displayed
length and scale factors, are at most $K_T(e_0+e_1^2)$.
The constant is independent of $j,p,n,Z$, compatible volume and
the real base. The common complex ball and the physical-current
estimate are proved in \cref{htc:common-ball,htc:tower-bound}.

Consequently there is $c>0$ such that
\begin{equation}
 \|v\|_1<c\min\{b_0,b_1,1\}
 \quad\Longrightarrow\quad
 (1,A',\mathbf J)\in\mathcal C_j(X;b_0,b_1).
 \label{dc:zero-core-ball}
\end{equation}
For a proof choose
\(e_0=\min\{b_0/(4K_T),a_T/4\}\) and
\(e_1=\min\{b_1/2,\sqrt{b_0/(4K_T)},a_T/4\}\).
Then $e_0+e_1<a_T$, as required by \cref{htc:all-towers}. The right side of
\eqref{dc:plaquette-bound} divided by \(\eta^2\), as well as \(|\mathbf J|\), is less
than \(e_0\) when \(c\) is small. The potential tests are less than
\(e_1\). The real field is the identity, so its gauge tests vanish.
Finally \(K_T(e_0+e_1^2)\le b_0/2\), which proves every tower clause.
This argument also proves a fixed central-core margin by reducing \(c\).

If a continued background pair \(p(B)\), on the physical datum ball
\(\|B\|_\infty<r_B\), satisfies
\(\|p(B)-(1,0)\|_1\le C_B\|B\|_\infty\), \eqref{dc:zero-core-ball} proves admission on
\begin{equation}
 \|B\|_\infty<
 \min\{r_B/2,\ c\min(b_0,b_1,1)/(2C_B)\}.
 \label{dc:datum-admission}
\end{equation}
The radius is in the physical datum. In a Gaussian datum it is divided
by \(g\). Thus the background estimate and the tower estimate supply a
near-zero ball directly; a further radius loss is not inferred just
because the physical radius runs with the coupling.

For two local holomorphic representatives, equality on a real open subset
of a common connected datum domain proves equality throughout that
domain. One proof expands at a real interior point: every real partial
derivative of their difference vanishes, hence every complex Taylor
coefficient vanishes. Analytic continuation along overlapping balls in
the connected domain finishes the proof. This applies to the first-step
background's admitted common complex charts. It applies to functions of
independent \((A',J)\) only when the real agreement is open in both
variables. Equality just at \(J=J_A(U)\) is not such an agreement;
the holomorphic function \(J-J_A(U)\) already demonstrates the distinction.

\subsubsection{Absolute and running systems of radii}

Set \(f_q(t)=t^{-1/2}(\log t)^q\). The two running field radii have
the form \(\alpha_{i,j}=C_i f_q(\theta_j)\), \(i=0,1\), while
\(\delta_j=C_\delta\theta_j^{-1/2}(\log\theta_j)^{p_0}\), with
\(q>p_0\). The constants \(C_i\) are fixed after the geometric
parameters. They are not changed when a complex source is introduced.
For a step of depth \(k\), a system \(\mathfrak a\) is admissible when
\begin{equation}
 \begin{gathered}
 C_sMa_{0,j}<a_{1,j},\qquad a_{1,j}^2\le C_s a_{0,j},\\
 a_{i,j}\le\bar a_i:=(1+c_o)a_{i,k+1}
       \le C_o\sqrt{k+2-j}\,a_{i,j},\\
 s_{\mathfrak a}:=\delta_k/a_{0,k+1}\le s_*,
 \qquad \bar a_0+\bar a_1\le a_*.
 \end{gathered}
 \label{dc:admissible-system}
\end{equation}
The first two conditions govern the shape of the core; the next line
compares every input radius with the output radius. The last line gives
fluctuation room and an absolute ceiling. An output radius larger than
an input radius is permitted; its ratio must be recorded in estimates.

The absolute system has entries \(a_{i,j}=\alpha_i\) for \(j\le k\)
and \(a_{i,k+1}=\alpha_i/(1+c_o)\). Its output radii are
\(\bar a_i=\alpha_i\). The running system has
\begin{equation}
 a_{i,j}=\min\{\alpha_{i,j-1},\alpha_{i,j}\}\quad(1\le j\le k),
 \qquad a_{i,k+1}=\alpha_{i,k},\qquad
 \varrho_j=\max_i\frac{\bar a_i}{a_{i,j}}.
 \label{dc:running-system}
\end{equation}
If a level-zero entry is needed, set \(a_{i,0}=\alpha_{i,0}\);
there is no negative-index radius. For \(k=0\) the old-member list
is empty and its output entry is \(a_{i,1}=\alpha_{i,0}\).
The adjacent minimum admits the same old member at the two adjacent
steps. It is not replaced by a radius chosen after that member is read.

\textbf{Radii lemma.} Suppose the real frame satisfies
\begin{equation}
 \theta_j\ge T,\qquad
 -c_5\le\theta_i-\theta_j\le B(j-i)\quad(i\le j).
 \label{dc:weak-frame}
\end{equation}
Choose \(C_1>C_sMC_0\). For sufficiently large \(T\), the absolute
system with sufficiently small fixed \(\alpha_i\), and the running
system below it, satisfy \eqref{dc:admissible-system}. Moreover
\begin{equation}
 \varrho_j\le C\sqrt{k+2-j},\qquad
 s_{\mathfrak a}\le C(\log\theta_k)^{p_0-q}.
 \label{dc:radius-ratios}
\end{equation}
For the proof, \((\log f_q)'(t)=(-\frac12+q/\log t)/t\), so
\(f_q\) is decreasing above \(e^{2q}\). Its logarithmic derivative
has absolute value at most \(1/t\) there. Thus bounded displacements of
\(t\) change \(f_q(t)\) by \(1+O(T^{-1})\), uniformly for
\(t\ge T\). The lower inequality in \eqref{dc:weak-frame} then proves
\(a_{i,j}\le(1+c_o)\alpha_{i,k}\), after increasing \(T\).
The adjacent minimum is at least \(C_i f_q(\theta_j+B)\).
For \(u\ge\theta_k\),
\(f_q(\theta_k)/f_q(u)\le\sqrt{u/\theta_k}\); for
\(\theta_k-c_5\le u<\theta_k\) its ratio is uniformly bounded.
Since \(\theta_j+B\le\theta_k+B(k+1-j)\), this proves the first
bound in \eqref{dc:radius-ratios}.

All entries of the running system are a common scalar times \(C_i\).
Its first shape condition follows from the choice of \(C_1/C_0\).
Its second follows from
\(a_{1,j}^2/a_{0,j}=(C_1^2/C_0)\min f_q(\theta)\to0\)
uniformly as \(T\to\infty\). The remaining ceiling follows the same
way, and division of \(\delta_k\) by \(\alpha_{0,k}\) proves the
second bound in \eqref{dc:radius-ratios}. For the absolute system choose the fixed
\(\alpha_i\) with the same shape, then increase \(T\) until the
running system lies below it and \(\delta_k/a_{0,k+1}\) is small.
Every choice is independent of \(k\).

There are two useful radius ratios for an old member. Its general
unmarked bound may use \(\max\{1,\delta_k/\alpha_j\}\). A homogeneous
near-field Taylor bound uses the actual ratio
\(\rho_j=\delta_k/\alpha_j\), which can be less than one. With
\(a=k-j\), \(\sigma_a=L^{-a}\), the frame bounds imply
\begin{equation}
 \frac{C_{\rm ax}\delta_k\sigma_a}{r_{B,j}}
    \le C s_{\mathfrak a}\sqrt{a+2}\,L^{-a},
 \qquad\text{provided }\alpha_j/r_{B,j}\le C.
 \label{dc:half-datum-disc}
\end{equation}
The supremum of \(\sqrt{a+2}L^{-a}\) is finite: it is bounded by
\(\sum_{a\ge0}(a+2)L^{-a}\). Consequently a single small choice of
\(s_*\) puts the central axial rows in a half datum disc at every
age. The remote rows have no such age gain. Their admission uses
the fixed absolute whole-datum ceiling
\eqref{pw:absolute-datum-ceiling} for the actual window maps.
For a different datum chart this inclusion must be verified
separately; the displayed radius ratio does not imply it.

The corresponding quadratic extraction weights obey the exact identity
\begin{equation}
 w_{n,j}=(\alpha_n/\alpha_j)^2,\qquad
 \alpha_n^{-2}w_{n,j}=\alpha_j^{-2}.
 \label{dc:radius-readout-weight}
\end{equation}
This identifies where the physical readout is paid. For example, if
\(r_j\le C\theta_j^{-p}\), \(p>2\), and the stronger frame bound
\(\theta_j\ge\theta_n+\lambda(n-j)-c_5\) holds, then
\begin{equation}
 \alpha_n^{-2}\sum_{j\le n}w_{n,j}r_j
 \le C\theta_n^{2-p}(\log\theta_n)^{-2q}.
 \label{dc:old-age-readout}
\end{equation}
Indeed the left side is bounded by
\(C\sum_{a\ge0}(\theta_n+ca)^{1-p}
(\log(\theta_n+ca))^{-2q}\), with bounded changes of the lower
floor absorbed in \(C\). Pull out \((\log\theta_n)^{-2q}\) and
compare the decreasing power with its integral. The integral is a
constant times \(\theta_n^{2-p}\). Thus summing ages and extracting a
quadratic coefficient are two specified operations, not two additional
copies of the same radius loss.

\subsubsection{Admission on a joint parameter domain}

Several distinct parameters occur in a local expression. A fine
decoupling parameter \(s_\Delta\) varies on \(|s_\Delta|<R\),
\(R=e^{\bar\kappa}>1\). A multiplier \(\nu\) amplifies an actual
response. A reference ratio \(\tau=x/\mathsf t\) lies in a fixed disc
\(|\tau-1|<c_\tau\). Physical sources vary the exponent members on
their ordinary source domain. An auxiliary potential parameter can have
a larger domain. The real frame \(\theta\), numerical cutoffs,
Gaussian scaling, reference operations and geometry are frozen before
any of these source derivatives are taken.

One cannot choose these domains independently by checking one real
parameter at a time. A useful exact calculus is the following.
Consider a finite ordered collection of atomic holomorphic maps
\(F_v\), each already defined on the ball
\(\|d-d_v^0\|<r_v\). The input \(d_v\) of a later map is a
specified holomorphic expression in the external parameters and earlier
outputs. Suppose normalized displacements satisfy majorants
\begin{equation}
 u_v\le e_v+\sum_{w<v}A_{vw}u_w,\qquad A_{vw}\ge0,
 \qquad u_v=\|d_v-d_v^0\|/r_v.
 \label{dc:admission-majorant}
\end{equation}
These inequalities must be verified on the proposed joint domain, using
primitive derivative estimates inside the atomic balls. If the
recursively computed right sides are at most \(\vartheta<1\), then
every composition is holomorphic there. This follows by induction:
the previous compositions exist, their values put the next input in its
already assigned ball, and composition of holomorphic maps is
holomorphic. There is no step of continuing an unknown map past its
domain.

Writing \(A\) for the strictly lower triangular matrix, the exact
finite bound is
\begin{equation}
 u\le(1+A+\cdots+A^{N-1})e.
 \label{dc:admission-paths}
\end{equation}
Each matrix product counts ordered paths through the dependency graph.
If its row norm is at most \(\zeta<1\), then
\(\sup u_v\le\sup e_v/(1-\zeta)\), independently of the number of
maps. Alternatively actual scale weights can bound the path sum even
when the unweighted row norm is not less than one. This bound must be
proved. For example the scalar recursion \(d_v=d_{v-1}+z\) has input
\(Nz\) at its last node; an individual bound \(|z|<r\) supplies no
depth-independent admission disc.

For a localized expression at old level \(j\le k\), the actual
substitution is of the form
\begin{equation}
 p_{j,Y}(s,\nu;B)
   =\rho_Xp\dotplus\nu\,\mathfrak h_Y(s;B),
 \qquad |s_\Delta|<e^{\bar\kappa}.
 \label{dc:local-substitution}
\end{equation}
The operation \(\dotplus\) changes both the potential and current by
the specified nonlinear local response; \(\rho_X\) is the restriction
appropriate to that term. It is not addition only to its potential
component. For an inner window the background is first rewritten on the
assigned box. Its datum is the actual nonlinear average
\begin{equation}
 M_{\mathbf B_j}\left(
 e^{i\eta_k(t+t_\square\zeta_\square)\mathbf H_k(B')}
 U_+\right),\qquad 0\le t\le1,
 \label{dc:window-actual-datum}
\end{equation}
with its declared continuation in \(t_\square\). The background of
this datum is inserted before the bracket interpolation. This is where
the pointed marginal cancellation acts. Its datum is not obtained by
replacing \(M_{\mathbf B_j}\) by its linearization.

\paragraph{Domain estimate III: the full local substitution.}
For each inner, outer and removed origin, each allowed localization
domain \(Y\), all its restricted towers and the entire joint polydisc
in \eqref{dc:local-substitution}, prove admission to the core of the term being evaluated.
In the normalized tests of \eqref{dc:margin-lemma}, a sufficient quantitative bound is
\begin{equation}
 \sup_\beta|N_\beta(p_{j,Y}(s,\nu;B))-
                       N_\beta(\rho_Xp)|
 \le C_Re^{C_R\bar\kappa}
       |\nu|\,\delta_k L^{-(k-j)}/a_{0,j}.
 \label{dc:substitution-core-target}
\end{equation}
The domain includes every admitted external complex pair, all real free
fields on their prescribed window box and every auxiliary parameter of
the expression. Atomic datum distances must satisfy their own radii as
in \eqref{dc:admission-majorant}. If the actual bound has a further geometric loss, that loss
must be included in the right side before choosing the parameters.

The numerical consequence of \eqref{dc:substitution-core-target} is complete. On a central core of
fraction \(\vartheta\), choose
\begin{equation}
 \rho_{j,k}=c_R e^{-C_R\bar\kappa}
       \frac{a_{0,j}}{\delta_k L^{-(k-j)}}.
 \label{dc:nu-radius}
\end{equation}
For \(|\nu|<\rho_{j,k}\), a sufficiently small \(c_R\) spends less
than half the initial margin. \eqref{dc:running-system}--\eqref{dc:radius-ratios} give
\begin{equation}
 \rho_{j,k}^{-1}\le
 C e^{C_R\bar\kappa}s_{\mathfrak a}
           \sqrt{k+2-j}\,L^{-(k-j)}.
 \label{dc:nu-small-factor}
\end{equation}
Hence one small, depth-independent choice of \(s_*\) makes every
\(\rho_{j,k}>2\), so \(\nu=1\) is an interior point. The gain
\(\rho_{j,k}^{-1}\) is obtained by a Schwarz estimate after admission;
it is not a device to create admission. The scale factor in \eqref{dc:nu-radius} is
what prevents a bound growing with the number of old levels.

\subsubsection{Cauchy estimates and the count of parameters}

Let \(F\) be holomorphic and bounded by \(M_F\) on a supremum ball
of radius \(r\). For two arbitrary directions \(u,v\), apply Cauchy
to \(F(z u/\|u\|_\infty+w v/\|v\|_\infty)\) on
\(|z|,|w|<r/3\). It proves
\begin{equation}
 |D^2F(0)[u,v]|
       \le 9M_Fr^{-2}\|u\|_\infty\|v\|_\infty.
 \label{dc:directional-cauchy}
\end{equation}
The same proof on a smaller centered ball works uniformly away from the
origin. It has no dimension factor. If instead only entrywise Hessian
bounds \(|\partial_i\partial_jF|\le M_2\) are known for \(m\)
coordinates, the elementary bound on the Euclidean operator norm is
\(mM_2\). A better bound needs information about the full scalar
supremum or the sum of the local entries. In a sum of local energies,
the global scalar supremum usually grows with volume; \eqref{dc:directional-cauchy} does not
turn that extensive bound into a volume-independent local one.

For independent decoupling coordinates the calculation is different.
If \(F(s,\nu)\) is bounded by \(M_F\) on
\(|s_\Delta|<R\), \(|\nu|<\rho\) with \(\rho>1\), and
\(F(s,0)\) is subtracted, Schwarz followed by coordinate Cauchy gives
\begin{equation}
 \left|\int_{[0,1]^K}\partial_K
      [F(s+1_N,1)-F(s+1_N,0)]\,ds\right|
       \le\frac{2M_F}{\rho}(R-1)^{-|K|}.
 \label{dc:bracket-cauchy}
\end{equation}
Here \(K,N\) are disjoint sets of coordinates and \(\partial_K\)
differentiates each coordinate in \(K\) once. A circle of radius
arbitrarily close to \(R-1\) fits at every real center in \([0,1]\).
Successive one-variable Cauchy formulae prove \eqref{dc:bracket-cauchy}, and the limit in
the radii is harmless. There is no \(|K|!\), because no coordinate is
differentiated repeatedly. By contrast, \(m\) derivatives in one
coordinate produce \(m!r^{-m}\). With \(R=e^{\bar\kappa}\) and
\(\bar\kappa\ge1\),
\begin{equation}
 (R-1)^{-|K|}\le e^{-(\bar\kappa-1)|K|}.
 \label{dc:disc-rate}
\end{equation}
This follows from \(1-e^{-\bar\kappa}\ge e^{-1}\).

If local domains of tree length between \(m\) and \(m+1\), rooted
at a fixed cube, number at most \(C_De^{\kappa_Dm}\), then every
fixed polynomial multiplicity is summable with a strict rate surplus:
\begin{equation}
 \sum_{m\ge0}(m+1)^A e^{-(\kappa-\kappa_D)m}<\infty
       \qquad(\kappa>\kappa_D).
 \label{dc:entropy-surplus}
\end{equation}
For completeness, put \(h=(\kappa-\kappa_D)/2\). The continuous
function \((t+1)^Ae^{-ht}\) is bounded on \([0,\infty)\), so the
series is bounded by a constant times \(\sum e^{-hm}\).
Using \eqref{dc:disc-rate} in a rooted bracket expansion requires the geometric
comparison between its number of parameter cubes and its output tree
length. It also requires the actual rooting identity: a component
disjoint from the term's base contributes zero. Analyticity by itself
does not give this identity.

For a family with potential norm \(\varepsilon\), one spare finite
potential slot of size \(w_*\) admits an auxiliary parameter on
\(|z|<w_*/\varepsilon\). If the uninserted local logarithm has norm
at most \(C_*\) on that slot, Cauchy proves
\begin{equation}
 \|DE[V]v\|_{\rm out}\le
          (2C_*/w_*)\|v\|_{\rm pot}.
 \label{dc:spare-slot}
\end{equation}
There is no dependence on the number of field coordinates or the total
volume in this estimate when the local majorant has no such dependence.
For two directions of norms \(\varepsilon_1,\varepsilon_2\), use
\(|z_i|<w_*/(4\varepsilon_i)\). The sum of the two potentials fits
in half the slot, and
\begin{equation}
 \|D^2E[V](v_1,v_2)\|_{\rm out}
       \le16C_*w_*^{-2}\varepsilon_1\varepsilon_2.
 \label{dc:two-spare-directions}
\end{equation}
For \(n\) unrelated simultaneous directions the analogous sum-budget
disc has radius proportional to \(w_*/(n\varepsilon_i)\) in each
coordinate. Reusing the two-direction radii for arbitrary \(n\)
would exceed the slot. This count concerns potential allowance, not the
distinct-cube calculation in \eqref{dc:bracket-cauchy}.

The auxiliary slot need not be an enlarged physical source domain for
histories. Suppose the actual source curve has
\(V(w)=V(0)+wv+O(w^2)\) in potential norm on its ordinary fixed disc.
The norm Taylor expansion of the holomorphic local map gives
\(\partial_wE(V(w))|_0=DE[V(0)]v\). Its quadratic remainder is bounded
by \(C_*\|V(w)-V(0)\|_{\rm pot}^2/w_*^2\), by one-variable Cauchy
on the line through that difference. Thus the large auxiliary disc
computes the actual first jet, even if the nonlinear physical curve is
defined only on its smaller source disc.

\subsubsection{Support, Schur norms, and finite coordinate counts}

For a matrix of color blocks, normalize the input and output coordinates
by their declared positive weights. Write \(B=D_oAD_i^{-1}\) for the
normalized matrix and put
\begin{equation}
 R_B=\sup_i\sum_j\|B_{ij}\|,\qquad
 C_B=\sup_j\sum_i\|B_{ij}\|.
 \label{dc:schur-norms}
\end{equation}
Then \(\|B\|_{2\to2}\le\sqrt{R_BC_B}\). To prove this, set
\(b_{ij}=\|B_{ij}\|\). For finitely supported vectors,
\(\|(Bx)_i\|^2\le(\sum_jb_{ij})(\sum_jb_{ij}\|x_j\|^2)\).
Sum over \(i\), use the column bound, and then complete. This also
proves the assertion on a countable lattice. A row estimate alone is not
the required Schur estimate. For example the matrix with a single
nonzero column containing \(m\) ones has row sum one and norm
\(\sqrt m\).

Suppose one scalar local energy has modulus at most \(M_0\) on a joint
coordinate domain with free radius \(h_u\) and background radius \(h_b\).
Suppose at most \(T\) such energies contribute to any fixed row and at
most \(N\) columns couple to that row, with the same bounds for columns.
Cauchy, first in two free directions and then in a background direction,
proves
\begin{equation}
 \|K\|_{\rm Schur}\le C NT M_0h_u^{-2},\qquad
 \|K(U')-K(U)\|_{\rm Schur}
       \le C NT M_0h_u^{-2}h_b^{-1}\|U'-U\|.
 \label{dc:local-cauchy-count}
\end{equation}
The backgrounds lie in a smaller common analytic chart, and the segment
between their coordinate values stays there. The first bound counts the
terms in an entry, then entries in a row. The second integrates the
third derivative along that segment. This proof specifies where
coordinate counts enter; they are not powers of the number of parameters
in \eqref{dc:bracket-cauchy}.

For the actual first-step chart, the verified finite geometry gives
\begin{equation}
 h_u=cL^{-9},\quad h_b=c'L^{-5},\quad
 T,N\le CL^4.
 \label{dc:first-step-ledger}
\end{equation}
A free coordinate changes its own bond and a fixed number of nearby
crossing pivots. Wilson terms have fixed plaquette incidence, and a
gauge block contains at most \(L^4-1\) internal contour rows. Thus
the number of scalar terms affecting one row and the number of coupled
free coordinates are each \(CL^4\), not their product before taking
the row sum. With \(M_0\le C\), \eqref{dc:local-cauchy-count} gives \(CL^{26}\) and
\(CL^{31}\|U'-U\|\). Also
\(\|\Phi_U(u)\|_\infty\le CL^4\|u\|_\infty\) on its free ball
implies, by two-direction Cauchy,
\(\|D^2\Phi_U(0)[e_i,e_j]\|\le CL^{13}\).
Only a fixed number of pivot outputs occur for a given row. After the
\(CL^4\) column count, the independent-current part is bounded by
\(CL^{17}\|J\|_\infty\). These are local first-step estimates.
They provide no hidden bound on a many-level window chart with a number
of coordinates growing with its side.

The actual constraint basis must also be used when localizing fields.
For a free-coordinate cutoff \(\chi\), the tangent field is
\(C_U(\chi u)\). It need not equal \(\chi C_Uu\), since pivot
values depend on neighboring free coordinates. A quadratic partition
argument on free variables therefore produces the commutator of the
actual compressed free-coordinate matrix. It is not obtained by clipping
the fine constraint rows. If \(H_0\ge c_0\) and a derived local
coefficient estimate gives \(\|H-H_0\|\le\varepsilon<c_0\), then
\begin{equation}
 \operatorname{Re}(u,Hu)_{\rm H}
        \ge(c_0-\varepsilon)\|u\|^2,\qquad
 \|H^{-1}\|\le(c_0-\varepsilon)^{-1}.
 \label{dc:gap-transfer}
\end{equation}
The inverse bound follows from the Neumann series with \(H_0^{-1}\).
A compression of this same matrix inherits the form inequality. An
operator rebuilt with different block means or boundary constraints
requires its own coefficient comparison; it is not that compression.

\paragraph{Domain estimate IV: coordinate counts across scales}
On the actual run window of side parameter \(r\), determine its free
coordinate set, its constraint inverse and its joint background domain.
Prove explicitly the bounds replacing \(h_u,h_b,T,N\) in \eqref{dc:local-cauchy-count},
including their dependence on \(r,L,M\), the region and the local layer.
Polynomial losses in \(r\) are allowed. The normalization matrices
\(D_i,D_o\), the compatible boundary rows and the nonwrapping lift are
part of this estimate. The finite one-step constants in \eqref{dc:first-step-ledger} are not
evidence that these window constants are independent of \(r\).

\subsubsection{Gauge changes, pivots and the real coarea factor}

Real unitary changes of local gauge preserve all the norm tests that are
written covariantly. A complex change has a cost. If
\(K=\sup_x\|g(x)\|\|g(x)^{-1}\|\), then
\(\|\operatorname{Ad}(g(x))X\|\le K\|X\|\).
For the actual average,
\(W_x(U^g)=g(c_-)W_x(U)g(c_-)^{-1}\). Its logarithm has the same
covariance as long as both points belong to the regular analytic branch.
The internal mean has the analogous bi-equivariance, since its defining
logarithmic equation is conjugated at its initial endpoint. Therefore
\begin{equation}
 \max_x\|W_x(U)-1\|<r_{\log}/(2K)
 \quad\Longrightarrow\quad
 \max_x\|W_x(U^g)-1\|<r_{\log}/2.
 \label{dc:complex-gauge-room}
\end{equation}
This requires a smaller analytic domain; it does not imply that a
complex gauge orbit saturates the original core. It also avoids multiplying a condition
number separately at every bond of a closed path: covariance reduces
that product to one conjugation. A real flat holonomy need not be close
to one for this calculation, since the closed comparison transports are
the arguments of the logarithms.

At a flat block-axial reference, the crossing pivot appears once in each
of \(L\) longitudinal segments and in no internal contour path. Its
raw constraint derivative is \(L^{-3}\operatorname{Ad}(V_c)\), with
the actual prefix transport \(V_c\). No other pivot occurs in that
constraint's support, under the nonwrapping double-block geometry. Thus
\begin{equation}
 J_c:=\partial_{a_{p(c)}}q_U(c)=L^{-3}\operatorname{Ad}(V_c),
 \qquad \det_{\mathbb R^3}J_c=L^{-9}
       \quad\text{at a real flat reference}.
 \label{dc:pivot-determinant}
\end{equation}
The determinant of a real adjoint transport is one: conjugation preserves
the color metric and is connected to the identity in \(\operatorname{SU}(2)\).
For the per-length pivot matrix \(\mathsf m_c=J_c/L\),
\(\det J_c=L^3\det\mathsf m_c\). These two determinants cannot
be interchanged in a coarea factor.

Let \(\sigma_H(a)d^3a\) be Haar measure in its regular exponential
chart. In the same chart
\(\delta_G(e^{iq})=\delta_{\mathbb R^3}(q)/\sigma_H(0)\).
Integrating one real pivot from a constrained integral therefore leaves
\begin{equation}
 \frac{\sigma_H(a_{p(c)})}
      {\sigma_H(0)\,|\det\partial_{a_{p(c)}}q_U(c)|}.
 \label{dc:real-coarea}
\end{equation}
At the flat reference this is \(L^9\). The remaining free bonds still
carry their Haar densities, and \(a_b=gy_b\) contributes
\(g^{3N_f}\) for \(N_f\) free bonds. The positivity and absolute
determinant in \eqref{dc:real-coarea} belong to real coarea. On a connected small real
chart its oriented determinant stays positive. The analytic factor used
later is its holomorphic determinant inverse, times the continued Haar
density, on the specified complex chart where the determinant is
nonzero. One does not analytically continue an absolute value.

\subsubsection{Finite potentials, cutoffs and renewed pointing}

The local potential norm used by the ordinary finite integral is
\begin{equation}
 \|\mathbf V\|_\lambda
  =\sup_{Y,P,B}
  \frac{e^{\lambda d(Y)}|V(Y;P,B)|}{Q_Y(B)},
 \qquad Q_Y(B)=1+M^{-4}\sum_{b\subset Y}|B_b|^2.
 \label{dc:potential-norm}
\end{equation}
The supremum uses its specified complex local cores and field box.
The coordinate \(B\) in this paragraph is the Gaussian variable of
\eqref{dc:gaussian-coordinate}. It is neither the physical datum of \eqref{dc:quadratic-conversion} nor the pivot basis
\(C_U\). An origin tag distinguishes several potential members with
the same geometric carrier. The tag changes neither \(d(Y)\) nor the
geometric support, but its norm is charged in the potential allowance.

Use nested boxes \(\mathcal B_0\subset\mathcal B_1\) with bounds
\(\mathfrak b\) and \((1+\epsilon_b)\mathfrak b\),
\(\mathfrak b\ge1\). Potentials to be differentiated are
holomorphic on the larger box; the finite slot majorant accepts the
resulting potentials on the smaller box. The real cutoff is kept on
that smaller box. For \(t\in[0,1]\), a coordinate Cauchy circle at
\(tB\) of radius
\begin{equation}
 r_{b,a}=c_b(1+|B_b^a|),\qquad c_b<\epsilon_b/4,
 \label{dc:adaptive-circle}
\end{equation}
fits in the larger box after reducing the fixed constant if needed.
The Euclidean color norm differs from the coordinate supremum by a
fixed factor. On that circle \(Q_Y\) is at most \(C_bQ_Y(B)\),
because the changed squared coordinate is at most
\(C(|B_b^a|^2+1)\). Consequently
\begin{equation}
 |B_b^a\partial_{b,a}v(Y;tB)|
       \le C_b\|\mathbf v\|_\lambda e^{-\lambda d(Y)}Q_Y(B).
 \label{dc:coordinate-source-leg}
\end{equation}
This estimate is uniform in \(\mathfrak b\). A fixed circle of radius
proportional to \(\mathfrak b\), estimated with the quadratic weight
at its largest point, would give an unnecessary \(\mathfrak b^2\)
when the original coordinate is small. The formula in \eqref{dc:adaptive-circle} retains the
quadratic norm of the actual field.

An older beta-source member follows the path \(v(tB)\). The intrinsic
action and gauge-fixing remainders, which vanish through order two,
instead have the normalization \(t^{-2}V(tB)\). This distinction also
has a direct analytic check. If \(V(tB)=\sum_{m\ge3}t^mV_m(B)\)
is bounded on a fixed larger \(t\)-disc \(|t|<1+\epsilon_t\), then
\begin{equation}
 \int_0^1\left|\frac d{dt}\{t^{-2}V(tB)\}\right|dt
       \le\sum_{m\ge3}|V_m(B)|
       \le C_{\epsilon_t}\sup_{|z|<1+\epsilon_t}|V(zB)|.
 \label{dc:radial-integrability}
\end{equation}
The first inequality integrates
\((m-2)t^{m-3}\), whose integral is one. Cauchy bounds the final
series geometrically. Applying \(t^{-2}\) to a general older source
with a linear or quadratic part would not give this integrability.
The actual bond-indexed split of the three interpolation legs must still
be used: a bound on their total does not bound the integral of the
absolute values of an arbitrarily chosen split.

\paragraph{Domain estimate V: the actual root slot.}
For each of the three assigned interpolation legs, prove the integrated
potential norm of its actual bond-root density, on \(\mathcal B_0\),
uniformly in the external core, carrier and interpolation time. The
larger box \(\mathcal B_1\) must provide the coordinate circles in
\eqref{dc:adaptive-circle}. Intrinsic legs use their actual third-order cancellation;
older source legs use \eqref{dc:coordinate-source-leg}. This is the precise root input to renewed
pointing, not a replacement of that pointing by an Euler split.

There is a simple exact interpretation of the connected term. On a
complete physical pair let the fixed reference integral be
\begin{equation}
 \Xi(z)=\int\chi(B)e^{V(B)+zv(B)}\,d\mu(B).
 \label{dc:physical-reference-integral}
\end{equation}
The reference Gaussian and cutoffs have no \(z\)-derivative. Wherever
the local expansion supplies its nonzero denominator,
\begin{equation}
 \frac d{dz}\langle a_z\rangle_z\bigg|_0
    =\langle\dot a_0\rangle_0
      +\langle a_0v\rangle_0-\langle a_0\rangle_0\langle v\rangle_0.
 \label{dc:connected-pointing}
\end{equation}
This follows by differentiating numerator and denominator, with the
same normalizer. The double-direction estimate in \eqref{dc:two-spare-directions}, union filing
and the rooted cluster bound estimate this root-plus-source term on
the actual local cores. On a general independent-current core the
localized logarithmic terms are still defined and differentiated by
their local finite formulas. They need not be marginals of one global
\(\Xi\); \eqref{dc:connected-pointing} is the physical sum interpretation, not an off-graph
identity inferred from holomorphy.

The number of scalar Gaussian coordinates in a local set is
\[
 N(Z)=3\#\{\text{free bonds in }Z\}
\]
and is of order \((LM)^4N_c(Z)\) when \(N_c\) counts coarse cubes. Thus a Gaussian
majorant \(e^{CuN(Z)}\) equals a bound of the form
\(e^{C'(LM)^4uN_c(Z)}\). The dimension factor is counted once.
Neither holomorphy of an integrand nor smallness of \(u\) alone
gives domination: the real reference precision, the fixed integration
sets and the available Gaussian quadratic reserve must be specified.

\subsubsection{The real weight, the reference Gaussian and domination}

The reference inverse coupling \(\theta\) and the real action weight \(x\)
have different roles. Put \(\tau=x/\theta\). In Gaussian coordinates the
actual quadratic form consists of \(\tau\) times the physical action
Hessian and the gauge Hessian at the reference weight. Write these two
free-coordinate matrices as \(K_A,K_G\). Then
\begin{equation}
 K_\tau=\tau K_A+K_G,\qquad K=K_A+K_G,\qquad
 e^{-\frac12 B^{\mathsf T}K_\tau B}
 =e^{-\frac{\tau}{2}B^{\mathsf T}KB}
       e^{\frac{\tau-1}{2}B^{\mathsf T}K_GB}.
 \label{dc:reference-representation}
\end{equation}
Thus using the normalized Gaussian with covariance \(\tau^{-1}K^{-1}\)
is a presentation choice with a compensating free potential. It does not
multiply the actual gauge weight by \(\tau\). Its normalizing determinant
contributes the factor \(\tau^{-N/2}\) for \(N\) real scalar coordinates.
The logarithm is chosen continuously from \(\tau=1\). Physical sources
are not inserted in \(x\) in this presentation; their Wilson terms occur
as potential members. Consequently neither this determinant nor \(g\)
has a physical-source derivative.

Here is the finite domination calculation with all its premises. Let
\(A_0\) be a real positive precision on \(N\) coordinates, \(A_0\ge m_0\).
Suppose the actual complex symmetric matrix satisfies
\begin{equation}
 \|A_0^{-1/2}(A-A_0)A_0^{-1/2}\|\le\epsilon,\qquad
 \rho\ge0,\qquad \epsilon+2\rho/m_0<\tfrac12 .
 \label{dc:relative-gaussian-premise}
\end{equation}
If a measurable integrand has modulus at most \(F_0e^{\rho|B|^2}\),
then the coherently normalized complex Gaussian gives
\begin{equation}
 \left|\int F(B)\,d\mu_A(B)\right|
 \le F_0\left(\frac{1+\epsilon}
                    {1-\epsilon-2\rho/m_0}\right)^{N/2}
 \le F_0e^{C(\epsilon+\rho/m_0)N}.
 \label{dc:gaussian-domination}
\end{equation}
Indeed the determinant ratio \(|\det A|/\det A_0\) is at most
\((1+\epsilon)^N\), while
\(\operatorname{Re}A-2\rho I\ge
(1-\epsilon-2\rho/m_0)A_0\). Integrating that real Gaussian proves
the first bound. The logarithms of numerator and denominator prove
the second. The same computation works with an integrand restricted
to a fixed Borel set. On smaller parameter polydiscs its right side
is an integrable bound independent of the parameters. If the densities
and integrands are holomorphic there, dominated Cauchy formulae or
Morera's theorem prove holomorphy of the integral.

The hypothesis on parameter-independent sets is substantive. A moving
sharp cutoff is not covered by the preceding proof. It must either be
frozen, as it is for the ordinary source construction, or be transformed
by an explicitly controlled change of variables, including its boundary
and Jacobian. Likewise holomorphy of a normalized integral uses the
coherent determinant branch and its nonvanishing matrix domain.

There is an explicit coefficient route to \eqref{dc:relative-gaussian-premise}. If the actual entries
of \(A-A_0\) satisfy row and column sums at most \(\epsilon m_0\),
\eqref{dc:schur-norms} gives \eqref{dc:relative-gaussian-premise}. For a ratio perturbation in \eqref{dc:reference-representation}, the additional
relative norm is at most
\begin{equation}
 |\tau-1|\,\|K^{-1/2}K_AK^{-1/2}\|.
 \label{dc:tau-room}
\end{equation}
At a flat reference \(0\le K_A\le K\), so this factor is at most
\(|\tau-1|\). Away from that reference its bound follows from the
actual local action coefficients and the proved gap, not from the name
of the coupling ratio. Here \(K\) is the real reference precision;
an additional complex-pair perturbation is tested by \eqref{dc:relative-gaussian-premise}.
Thus the allowed \(\tau\)-disc is chosen after
these constants and the unused quadratic potential reserve.

\paragraph{Domain estimate VI: primitive Gaussian and potential room.}
For every local conditioned matrix used by the assigned decoupling,
derive the normalized entry sums needed for \eqref{dc:relative-gaussian-premise}, on its enlarged
parameter discs and its external core. Determine the scalar coordinate
count, the reference precision and the total positive quadratic growth
of the admitted potentials. Verify that the sum leaves the strict
inequality in \eqref{dc:relative-gaussian-premise}, with the small per-cube cost required by the
ordinary polymer criterion. This is one matrix/coefficient and
potential-budget estimate. Once it holds, \eqref{dc:gaussian-domination}, the spare-slot
derivatives and domination of parameter integrals are conclusions.
It is not necessary to identify all these local matrices with marginals
of one Gaussian at a general independent-current pair.

\subsubsection{The site-frozen complex source box}

The ordinary complex data allow each site to carry its own Euclidean
member. Such a member consists of a pointed local family, its infinite
companion and its own scalar coefficient, with covariance of the
pointed total under the required site symmetries. Members chosen at
different sites need not be related. For a label \(\mathcal S\), let
\(h_z^{(j)}\) be its fixed nonnegative tents, with
\(\sum_z h_z^{(j)}=1\), and let \(0\le\varphi_j\le1\) be the nested
profiles. Write \(b_j\) for the real unsourced increments. Define
\begin{equation}
 \begin{aligned}
 c_{\mathcal S}(p)
   &=u(p)+\sum_{j\le k}\varphi_j(p)
             \sum_z h_z^{(j)}(p)[e_j(z)-b_j],\\
 c_{j,\mathcal S}(p)
   &=\varphi_j(p)\sum_z h_z^{(j)}(p)e_j(z),\\
 x_k(p)&=\beta-\sum_{j\le k}b_j\varphi_j(p).
 \end{aligned}
 \label{dc:site-frozen-weights}
\end{equation}
Each site uses the coefficient \(e_j(z)\) of its own full Euclidean
member, even though the exponent reads just that member's terms
pointed at \(z\). A defect does not inherit this symmetry by being
placed at that site.

The exact cancellation of the scalar action weights is
\begin{equation}
 -A(x_k-c_{\mathcal S},U)-\sum_{j\le k}A(c_{j,\mathcal S},U)
       =-A(\beta-u,U).
 \label{dc:weight-cancellation}
\end{equation}
Expand \eqref{dc:site-frozen-weights}, use \(\sum_zh_z^{(j)}=1\), and use linearity of the
Wilson action in its weight. This proves the identity label by label.
It does not depend on the number of sites. It also shows why replacing
an own counterterm by a coefficient from another member changes the
presentation.

The scalar part of the box is measured by
\begin{equation}
 \mathsf u=\|u\|_\infty+
         \sum_{j\le k}\sup_z|e_j(z)-b_j|.
 \qquad \|c_{\mathcal S}\|_\infty\le\mathsf u .
 \label{dc:scalar-box-bound}
\end{equation}
The inequality uses only the tents and \(0\le\varphi_j\le1\).
For ordered profiles, \(x_k(p)\) is a convex combination of
\(\beta,x_1,\ldots,x_k\): its coefficients are
\(1-\varphi_1,\varphi_1-\varphi_2,\ldots,\varphi_k\).
Thus on real field integrals
\(\operatorname{Re}(x_k(p)-c_{\mathcal S}(p))
\ge\min_i x_i-\mathsf u\). A fixed scalar box \(\mathsf u\le\bar c\)
preserves a positive Wilson coefficient when the real floor exceeds
\(\bar c\). This is a weight calculation; it is not a positivity
assertion for arbitrary complex local activities.

Let \(d_j(X)\) be the localization tree length and \(d_j^\natural(X)\)
the hull length with the specified standing islands exempted. The
remaining coordinates of the ordinary box have sizes
\begin{equation}
 \begin{gathered}
 \mathsf g_j=\sup_{z_0,X,z,P}e^{\kappa d_j(X)}|G_{z_0}^{(j)}(X,P,z)|,
 \quad
 \mathsf a_j=\sup_{X,z,P} e^{\kappa d_j(X)}|D^{(j)}(X,P,z)|,\\
 \mathsf r_j=\sup_{X,P}e^{\kappa d_j(X)}|H^{(j)}(X,P)|,\qquad
 \mathsf h=\sup\frac{|B_{\mathcal S}^{(j)}(X)|}
                       {\bar B_{\mathcal S}^{(j)}(X)},\qquad
 \mathsf f=\sup_\omega\frac{|\Psi_\omega|}{n(\omega)} .
 \end{gathered}
 \label{dc:box-sizes}
\end{equation}
The external \(P\) ranges over the prescribed hereditary core.
The boundary denominator is its actual unsourced majorant, proportional
to \(e^{-\kappa d_j^\natural}\). If
\(\Psi_\omega=\Psi_\omega^\circ+\sum_f\psi_f(P|_{R_f})\), a separate
oscillation size is
\begin{equation}
 \mathsf f_\Phi=
 \sup_{\omega,R}
 \frac{\sum_{f:R_f=R}\operatorname{osc}_P\psi_f}
      {(\log\theta_i)^{2p_0}e^{-\kappa d_i^\natural(R)}} .
 \label{dc:history-oscillation}
\end{equation}
Amplitude and oscillation are different coordinates: a constant exponent
has zero oscillation and may have arbitrarily large amplitude.

At input or output index \(m\in\{k,k+1\}\), the actual box bounds are
\begin{equation}
 \begin{gathered}
 \mathsf u\le\bar c,\quad \mathsf g_j\le E_0+1,\quad
 \mathsf a_j\le L^{-(m-j)},\quad
 \mathsf r_j\le C\theta_j^{-\kappa_0/2},\\
 \mathsf h\le C,\qquad \mathsf f\le\bar f,\qquad
 \mathsf f_\Phi\le Cc_\times(\bar c).
 \end{gathered}
 \label{dc:actual-source-box}
\end{equation}
The label spaces, tests, real backgrounds, reference operations, measures
and all numerical Gaussian constants are the real ones. The complex
variables are exponent members and their own coefficients. This is why
local functions can be differentiated without differentiating a sharp
test or an adaptively chosen threshold.

For a physical source \(w\), the site's member is the one for the
constant source \(\zeta=w(p_z)\). The defect is
\(E^{(j),w}(X,P,z)-E^{(j),w(p_z)}(X,P,z)\), not an independent
Euclidean member. A localization that sees only \(w\) within a domain
of diameter \(D\) gives the elementary difference bound
\(\sup|w-w(p_z)|\le D\,\operatorname{Lip}w\). Cauchy in this localized
source direction then gives the source-gradient gain when its actual
coefficient norm is known. The spatial inequality supplies the gain in
\(\operatorname{Lip}w\); it does not supply a missing \(\alpha_j^2\)
in the coefficient norm.

The fixed scalar ceiling in \eqref{dc:actual-source-box} also explains the two source discs.
The potential \(uA\) has size \(O(\alpha_j^2\|u\|_\infty)\) on a
small-field chart, and therefore fits an enlarged auxiliary finite slot.
The same \(\|u\|_\infty\) occurs without this factor in \eqref{dc:scalar-box-bound}.
It can also occur repeatedly in a history exponent. Hence finite-slot
admission cannot enlarge the full source box by \(\alpha_j^{-2}\).
The first-jet calculation in \eqref{dc:spare-slot} remains valid.

There is an exact elementary reserve calculation for history amplitudes.
Write \(n(\omega)=\sum_j n_j(\omega)\), including each swallowed history
exactly once, and \(v(\omega)=\sum_j\log\theta_j\,n_j(\omega)\).
If \(\log\theta_j\ge2\bar f\), then
\begin{equation}
 |\exp\Psi_\omega|\le e^{\bar f n(\omega)}
                \le e^{v(\omega)/2}.
 \label{dc:history-reserve}
\end{equation}
It can spend half of a previously established \(e^{-v(\omega)}\)
reserve, while keeping the remaining half. The retained inner Wilson
factor and component rarity must still be present in the original
mass estimate. The formula in \eqref{dc:history-reserve} does not manufacture that estimate.

\paragraph{Domain estimate VII: source incidence and excess exponents.}
For one actual step or one large-field component, list every use of a
given complex member, its own coefficient and its scalar anchor.
Prove the aggregate site bound and the bounds in \eqref{dc:box-sizes}--\eqref{dc:actual-source-box} after
formation, using a single unused action reserve. For a component mass
the supremum over its admitted external pair is placed inside its
stored-variable integral. For drop-outs it is over the actual fibre
of stored variables and recorded tests, not a product set containing an
unrelated flat field. These specified incidence and fibre estimates
are the remaining source-domain calculation; independent one-factor
disc checks do not give the joint box.

\subsubsection{Large run windows and old activities}

The actual free-coordinate window may have
\(m_I\le C(r+D)^4\) scalar coordinates. A useful larger analytic chart
is constructed in physical coordinates and then divided by \(g_k\);
it is not obtained by continuing the cutoff indicator. Suppose the
actual physical pivot map on this window and every atomic background
composition have radius at least
\begin{equation}
 h_W\ge c\,\alpha_k r^{-A_c},\qquad
 R_y=h_W/g_k,
 \label{dc:window-gaussian-radius}
\end{equation}
and their normalized displacements obey the admitted graph bounds in \eqref{dc:admission-majorant}--\eqref{dc:admission-paths}. This defines a precise chart target, including the
polynomial loss \(A_c\). It is not a uniform one-block constant.
For old level \(i\), the output displacement appearing in that graph
has the scale ratio \(L^{-(k-i)}\alpha_k/\alpha_i\); by \eqref{dc:radius-ratios} its
supremum is bounded. Amplified window directions use the same estimate
with their actual multiplier, not a separate analytic continuation
argument for each real direction.

If the cutoff is \(\mathfrak b_k\asymp(\log\theta_k)^{p_0}\) and
\(r\le C(\log\theta_k)^{p_r}\), then the cutoff and two fixed-fraction
coordinate circles fit in \eqref{dc:window-gaussian-radius} whenever
\begin{equation}
 q>p_0+A_cp_r,\qquad \theta_k\ge T_*.
 \label{dc:window-box-room}
\end{equation}
Indeed \(R_y\ge c(\log\theta_k)^{q-A_cp_r}\), so
\(\mathfrak b_k/R_y\to0\). A finite increase of \(T_*\) gives any
prescribed strict fractional room. If the complete old exponent has
scalar supremum at most \(Cr^{A_0}\) on this larger chart, its
directional Hessian is bounded by \eqref{dc:directional-cauchy}; an entrywise and row-count
proof can instead yield a bound of the form
\begin{equation}
 \|\operatorname{Hess}_{y_I}P_{\rm old}\|
       \le Cr^{A_{\rm tot}}(\log\theta_k)^{-2q}.
 \label{dc:old-hessian-budget}
\end{equation}
Here \(A_{\rm tot}\) must include the two losses \(2A_c\), the actual
amplitude and coordinate-count losses, and any response losses.
For \(r\le C(\log\theta_k)^{p_r}\), this tends to zero if
\(2q>A_{\rm tot}p_r\). If the free precision has a fixed gap \(m_*\),
increasing \(T_*\) makes the bound less than \(m_*/4\).

\paragraph{Domain estimate VIII: the actual window chart and exponent.}
Prove \eqref{dc:window-gaussian-radius} for the actual constrained window, exterior coordinates
and gauge quotient; prove the graph bounds for all local substitutions;
and derive its complete old-exponent amplitude or localized entry
bounds before using \eqref{dc:old-hessian-budget}. The count \(m_I\), the old-age summation
and the normalization of each old activity are part of this estimate.
The localized old-activity theorem supplies the Hessian implication
from these primitive data and the pointed remainder. It does not
identify a conditional exterior law.

For the near, nonwrapping part there is a sharper estimate. Let a regular
source derivative born at level \(j\) have norm
\(A_j\le C\alpha_j^2s_j\), and put
\(\rho_j=\delta_n/\alpha_j\), \(a=n-j\).
The actual pointed window theorem, after its fine-site count, gives
\begin{equation}
 C A_j\left(\rho_j^2L^{-\alpha_Ha}
                      +\rho_j^3L^{-2a}\right).
 \label{dc:true-near-ratio}
\end{equation}
The needed hypothesis is a bound \(C\delta_n\) on the complete
background-plus-fluctuation field used in the near Taylor estimate.
A bound only on its increment does not give it. Under that hypothesis,
\eqref{dc:half-datum-disc} admits the central half-disc at every age,
even when \(\rho_j<1\); the remote datum uses the separate absolute
ceiling \eqref{pw:absolute-datum-ceiling}. On the complete homogeneous
window class, \eqref{wf:whole-window-amplitude} and the full-profile
calculation preceding \eqref{eq:local-formation-LB17} supply the
proportional \(C\delta_n\) remote-datum bound as well. For a class
without that bound its response is retained as a separate tail.
The frame gives \(A_j/\alpha_j^2\le Cs_n\) and
\(\rho_j\le Cs_n\sqrt{a+2}\). Therefore the sum of \eqref{dc:true-near-ratio} is at most
\(C\delta_n^2s_n\), by geometric summation. A defect with an additional
factor \(h_j\operatorname{Lip}w\) sums the same way since
\(h_j=h_nL^{-a}\). A field-independent remote or winding term has no
such homogeneous radius factor and is assigned its own decaying class.

The exterior-force estimate is a separate bounded target. In the
assigned conditional coordinates \(I,O\), it must control
\begin{equation}
 Q^{\rm sc}y,\quad Q^e_{IO}y_O(e),\quad
       (Q^e_{II}-Q^{e_0}_{II})y,\quad
       \nabla_I(P_e-P_{e_0}),\qquad
 Q^{\rm sc}=Q_{IO}Q_{OO}^{-1}Q_{OI}.
 \label{dc:exterior-force-targets}
\end{equation}
For example the third term needs the actual bound
\(\left|((Q^e_{II}-Q^{e_0}_{II})y)_b\right|
\le C\mathfrak b\,e^{-c\,d(b,O)}\), uniformly over the admitted
exterior data and the real window; the fourth needs the corresponding
exponential derivative bound without an uncontrolled exterior factor.
They require control of the actual exterior field and conditional
exponent. Neither \eqref{dc:old-hessian-budget} nor decay of one covariance proves them.

\subsubsection{Carriers, world floors and the source norm}

A world is a torus with physical side \(\ell\) and cutoff \(K\).
Its numerical side at level \(j\) is
\(\mathfrak s_j=\ell L^{K-j}\). A comparison of local formulas requires
one injective rectangular lift containing the whole bond/path
neighborhood being compared. A separate lift for each endpoint is not
enough. In particular the pivot determinant calculation requires the
support of the relevant constraint to contain no other crossing pivot.
That support, including its contour paths, is what is lifted.

For the actual finite average the determining neighborhood grows by a
fixed number of blocks at each scale. If \(R_j\) is its radius in bare
numerical units after \(j\) steps, the elementary support recursion is
\begin{equation}
 R_{j+1}\le R_j+C_LL^{j+1}\eta,\qquad
 L^{-k}\eta^{-1}R_k
       \le C_L\sum_{a=0}^{k-1}L^{-a}
       \le \frac{C_L}{1-L^{-1}}.
 \label{dc:average-support}
\end{equation}
Every internal path lies in its assigned block and every straight
segment in the assigned double block, which proves the first
inequality. Composing supports proves the recursion; the displayed
geometric sum proves the depth-independent radius in final units.
The differentiated map has the same determining neighborhood.
Taking a local plaquette or current adds its explicit finite range.

This proof concerns the average, not the full nonlinear minimizing
background. The latter has inverse tails. A localized walk term has
its own finite surroundings; the complete inverse is an absolutely
convergent sum of such terms. Carrier identities for a localized
substitution use its decoupling and rooting rule. They cannot be
deduced by applying \eqref{dc:average-support} to an unlocalized inverse.
For instance a scale-dependent collar of width \(e_j\) contributes
\(\sum_{j\le k}e_jL^{j-k}\) in final units. If \(e_j\) is bounded,
this is a geometric sum. The weaker relation \(e_j\le Le_{j+1}\)
alone does not give a depth-independent sum: \(e_j=L^{k-j}\)
makes every summand one. Thus changes of localization size must be
checked in the actual carrier construction.

The nonwrapping conditions have three separate uses. The double blocks
of an average and the cubes of local gauges must be shorter than half
the torus. A gauge cube enlarged by its comparison collar must still
have that property. Finally each tower is formed only to a depth at
which its prescribed inner cube size fits in half the torus. If the
largest required numerical size in the relevant level is \(D_j\), the
test is
\begin{equation}
 2D_j<\ell L^{K-j}.
 \label{dc:world-floor}
\end{equation}
This condition is unchanged by moving the lifted patch around the torus.
It asserts no smallness of holonomy around a noncontractible cycle.
The first-step curved physical gap handles arbitrary real holonomies
by its local gauge and partition argument; it is not obtained by
choosing one globally small potential.

\paragraph{Domain estimate IX: compatible carriers and floors.}
For each average, tower, conditioned inverse term and inner, outer or
removed substitution, give its actual determining neighborhood and
collars. Prove that its required lifts satisfy \eqref{dc:world-floor}, and that each
restriction or transport carries all its defining rows and parameter
cubes. For full inverses, give the separately summed boundary/image
tail. This finite-list geometric estimate must accompany the analytic
estimate III; replacing a region operator by a clipped average does
not preserve its determining neighborhood.

Let \(h_j=L^{j-K}\) be the physical spacing of level \(j\), and let
\(f\) be a Lipschitz test on the physical torus. Use
\begin{equation}
 \|f\|_{\sharp,j}=
       \max\{\|f\|_\infty,c_Mh_j\operatorname{Lip}f\},
       \qquad c_M\ge C_{\rm geo}M.
 \label{dc:sharp-test-norm}
\end{equation}
The constant \(c_M\) contains the fixed block-size factor. If a domain
of tree length \(d_j(X)\) sees plaquettes within physical distance
\(C_{\rm geo}Mh_j(d_j(X)+1)\) of its anchor, then
\begin{equation}
 |f(p)-f(p_z)|
  \le C_{\rm geo}Mh_j(d_j(X)+1)\operatorname{Lip}f
  \le(d_j(X)+1)\|f\|_{\sharp,j}.
 \label{dc:source-anchor-bound}
\end{equation}
This follows from the Lipschitz definition and the diameter of a
connected tree of cubes. The polynomial \(d_j(X)+1\) is absorbed by
the strict entropy surplus in \eqref{dc:entropy-surplus}. Moreover
\(\|f\|_{\sharp,j}\le\|f\|_{\sharp,n}\) for \(j\le n\), and
\(h_j/h_n=L^{-(n-j)}\). These are exactly the scale factors of a
source-gradient defect in a later window. No second block factor
is removed from \eqref{dc:source-anchor-bound}; \(M/c_M\), rather than \(M\), is the fixed
ratio being bounded.

\subsubsection{Tubes of frozen couplings and real admission}

The frozen reference choices are formed from a real frame \(\theta\).
Let the real recursion have weights
\(x_j^\natural=\beta-\sum_{i\le j}b_i\), deviations
\(z_j=\theta_j-x_j^\natural\), and nonincreasing tube radii \(t_j\).
The actual tube requires
\(t_j(\log\theta_{j+1})^{C_t}\le c_t\theta_{j+1}\).
Room \(2r\) means \(|z_j|+2r\le t_j\) at every intermediate level.
It is a simultaneous condition, not just a condition at the endpoint.

Suppose the constant-source coefficients are holomorphic on
\(|\zeta|<2r\), and their difference quotients have bounds \(d_i\)
with \(\sum_i d_i\le B_0\). Their exact running shift is
\begin{equation}
 \zeta_j(\zeta)=\zeta+\sum_{i\le j}[b_i(\zeta)-b_i(0)],
 \qquad
 |\zeta_j(\zeta)|\le(1+B_0)|\zeta|.
 \label{dc:constant-source-tube}
\end{equation}
Thus \(|\zeta|<r/(1+B_0)\) leaves room \(r\) inside every tube.
The estimate is just the triangle inequality for the displayed
recursion and is independent of the number of levels. The substantive
input is the summable quotient bound on its actual disc, including the
marked bare-disc forcing if that is how the bound is obtained.

A constant source may be read either in the fixed source presentation
of \eqref{dc:site-frozen-weights}, or as an alternative frozen-weight presentation. In the
former, it remains absent from the Gaussian structures; in the latter,
the compensating presentation identity in \eqref{dc:reference-representation} must be used with the
new weights. Equality of the finite integrals includes all these
normalizing and interaction factors. It is not justified by changing
\(\tau\) in only the Gaussian determinant.

Real admission also matters when comparing two holomorphic
presentations. On a connected complex overlap that contains a real
open set where both are admitted and represent the same bare integral,
the real identity extends holomorphically. The proof is the Taylor
argument following \eqref{dc:datum-admission}. Geometric overlap of two complex discs
without a common admitted real set gives no such conclusion for the
actual presentations. Endpoint bracketing and slope bounds on a
canonical admitted interval can therefore be used for tuning on that
interval; broader calibration gluing requires its own real-admission
estimate.

\paragraph{Domain estimate X: the actual coupling tube and overlap.}
Prove the coefficient quotient bounds used in \eqref{dc:constant-source-tube} on the declared
frozen discs, preservation of the intermediate tube with its prescribed
slack, and the actual real admission of any overlap on which gluing is
claimed. The objects are the unchanged frozen choices and the actual
constant-source coefficients. This estimate is separate from the
many-level field core: a bound on a physical field radius does not
give a bound on an analytic coupling disc.

\subsubsection{A simultaneous choice of the remaining small parameters}

The elementary choice problems can now be separated from their model
estimates. First determine the polynomial degrees and their dependencies.
The logarithmic exponents are chosen before geometric parameters when
the actual construction requires that order. The fixed block size,
its finite-map constants, the propagator rates, tower and collar sizes,
remaining analytic radii and localization sizes are then chosen in their
declared order; the weak-coupling floor is last. For the numerical
calculation below these geometric data and a strict entropy surplus
have already been fixed.
Choose the enlarged decoupling radius \(e^{\bar\kappa}\) so that the
remaining rate in \eqref{dc:disc-rate}, after geometric and entropy charges, is still
positive. Choose the shape constants \(C_0,C_1\) and the absolute
ceiling to satisfy \eqref{dc:core-linear-class} and \eqref{dc:admissible-system}. The quantities \(L,M,\bar\kappa\)
are then fixed; a constant such as \(e^{C_R\bar\kappa}\) is not allowed
to depend on the depth.

The analytic logarithmic exponent \(q\) has to satisfy every required
finite lower bound, including
\begin{equation}
 q>p_0,\qquad q>p_0+A_cp_r,\qquad
       2q>A_{\rm tot}p_r.
 \label{dc:logarithmic-choice}
\end{equation}
Any additional power of the window radius or history reserve used by
the actual estimates must be added before that exponent is chosen.
The elementary feasibility conclusion assumes that the displayed
polynomial degrees do not themselves increase with the new choice
of \(q\). When \(q\) must be chosen before later geometry, these
degree bounds must also be uniform over the later admitted choices,
or a simultaneous feasibility list must be solved in the required
order. Fixed-geometry logarithmic smallness alone does not prove that
earlier choice. If the degrees do depend on \(q\), the resulting
simultaneous inequalities must
first be solved with that dependence; increasing \(q\) alone is no
proof of feasibility. For finitely
many conditions \(C_\nu(\log T)^{-b_\nu}<\epsilon_\nu\), with
\(b_\nu>0\), a single sufficiently large \(T\) satisfies them all:
take \(\log T>\max_\nu(C_\nu/\epsilon_\nu)^{1/b_\nu}\).
The conditions \(CM f_q(T)<c\) hold after a further increase, since
every power of \(\log T\) is \(o(T^{1/2})\). This proves simultaneous
smallness of the radius ceilings, \(s_{\mathfrak a}\), the factors in
\eqref{dc:nu-small-factor}, the Gaussian ratio perturbation and every explicitly listed
logarithmic loss. It is a finite feasibility calculation after the
exponents of the actual estimates have been obtained.

Finally fix the source scalar ceiling and excess-exponent ceiling, and
increase the real floor so that
\(\log T\ge2\bar f\) and the Wilson weights in \eqref{dc:scalar-box-bound} remain positive
on the required real integrals. The spare potential allowance is chosen
inside the already established Gaussian/polymer reserve; the nested
field boxes spend their fixed fractional slack once. None of these
choices alters an earlier real cutoff or reference operation under a
source derivative.

At this point the algebraic domain reductions are closed. Domain
estimates I--X specify what remains to be proved for the actual
many-level operators, towers, substitutions, roots, Gaussian matrices,
source arrays and run windows. They have explicit objects, radii and
uniformities. The finite chart and physical background results realize only their
stated first-step chart, background, gap and radius instances.
Independent-current normal-form matching, actual integrated roots,
complete local Gaussian budgets and source/history estimates remain
separate wherever they have not been supplied, even at a first step. At further levels the
action is the actual composed action \(A(U_k(V))\), with its constraint
and current terms. Reusing the numerical first-step constants without
the corresponding many-level estimate would not be an induction.





% END MODULE chapters/domain_core_ledger.tex

% BEGIN MODULE chapters/explicit_relative_current_extension.tex
\subsection{A relative-gauge current extension at the first step}
\label{ch:explicit-relative-current-extension}

An independent-current operator is not determined by its values on the
physical current graph.  We therefore give a concrete alternative
definition at one step and prove what it supplies.  Its complete
nonlinear action, quadratic form, gauge covariance, real gap and complex
inverse are specified.  At a constrained-critical axial background it
has the actual physical action as its value, throughout the admitted
fibre chart.  Agreement with any other current completion requires a separate
identification.  Using it at later scales
requires the further consistency and closure checks stated below.

\subsubsection*{The raw fibre and the relative contour functional}

Use fine lattice spacing one in this section.  Currents are raw
covectors for the bond logarithms with the counting pairing.  The
conversion to a potential or a different length unit is the domain
section's conversion, including its dual cell weights.
For the actual one-step group average put
\begin{equation}
 q_U(a)=\frac1i\log[\mathcal M(e^{ia}U)\mathcal M(U)^{-1}],
 \qquad P_U=Dq_U(0).
 \label{eq:explicit-relative-current-extension-NX1}
\end{equation}
All logarithms and contour means use their admitted regular branches.
The geometry is the embedded double-block geometry of the pivot chart,
on the infinite lattice or a compatible torus.  In particular no second
crossing pivot enters a row.

Let $h_U$ be the pivot right inverse, let $\iota$ insert free
coordinates with zero pivots, and set
\begin{equation}
 C_U=\iota-h_UP_U\iota,\qquad
 q_U(\Phi_U(u))=0,\qquad
 \Phi_U(0)=0,\qquad D\Phi_U(0)=C_U .
 \label{eq:explicit-relative-current-extension-NX2}
\end{equation}
The free restriction of $\Phi_U(u)$ is exactly $u$.  The
quantitative chart gives a joint background and free-variable
neighbourhood, with free radius $r_u=cL^{-9}$ and background
radius $r_b=c'L^{-5}$ after its specified fixed shrinkages.
Two differentiations give
\begin{equation}
 D^2\Phi_U(0)[u,v]
           =-h_UD^2q_U(0)[C_Uu,C_Uv].
 \label{eq:explicit-relative-current-extension-NX3}
\end{equation}
The second derivative has only pivot components, so the
right inverse on those components proves the identity and its sign.

Write $f_U(a)=A_W(e^{ia}U)$ and let
$j_U=Df_U(0)$ be its raw gradient.  On a finite torus these
are ordinary finite-dimensional functions.  On the infinite
lattice every scalar action difference below is first evaluated
for finitely supported free variations.  The locally finite
decomposition and its derivatives then define the kernels.
No extensive scalar action is assigned a finite supremum on
all bounded fields.

For every internal contour row define its relative logarithm by
\begin{equation}
 \ell_{U;y,x}(a)=\frac1i\log[
        T_{y,x}(e^{ia}U)T_{y,x}(U)^{-1}],\qquad
 G_Ua=(D\ell_{U;y,x}(0)a)_{y,x\ne y}.
 \label{eq:explicit-relative-current-extension-NX4}
\end{equation}
The relative gauge function is
\begin{equation}
 {\cal G}_U(a)=\sum_{y,x\ne y}
   \left(1-\tfrac12\operatorname{tr}
       [T_{y,x}(e^{ia}U)T_{y,x}(U)^{-1}]\right).
 \label{eq:explicit-relative-current-extension-NX5}
\end{equation}
It has value and first derivative zero at $a=0$.  Since
$1-\frac12\operatorname{tr}e^{iX}
=\frac12|X|^2+O(|X|^4)$, its Hessian is $G_U^{\mathsf T}G_U$.
Here $G_U$ acts on full bond variations; on free coordinates
the contour differential is $G_UC_U$.  This distinction fixes
the placement of both copies of $C_U$.

\subsubsection*{A stationary extension for every independent current}

For a declared independent raw current $J$ define
\begin{equation}
 \begin{split}
 S^A_U[J](u)={}&f_U(\Phi_U(u))-f_U(0)
                         -\langle j_U,\Phi_U(u)\rangle\\
             &+\langle J,\Phi_U(u)-C_Uu\rangle,\\
 S^{\rm rel}_U[J](u)={}&S^A_U[J](u)+{\cal G}_U(\Phi_U(u)).
 \end{split}\label{eq:explicit-relative-current-extension-NX6}
\end{equation}
The pairings are complex bilinear.  These formulas are definitions
on the actual pivot chart, rather than an extension by agreement
on a lower-dimensional graph.

Every term in the first derivative cancels:
$Df_U(0)C_U=j_UC_U$, the next term subtracts it, and
$D(\Phi_U-C_Uu)(0)=0$.  The relative gauge also has zero
first derivative.  Thus
\begin{equation}
 S^{\rm rel}_U[J](0)=0,\qquad DS^{\rm rel}_U[J](0)=0
                         \quad\hbox{for every admitted }(U,J).
 \label{eq:explicit-relative-current-extension-NX7}
\end{equation}
The cancellation uses the actual $j_U$ in the subtraction.
Replacing it by $J$ there would give a different linear term.

Put $\mathscr A_U=D^2f_U(0)$.  The chain rule and \eqref{eq:explicit-relative-current-extension-NX3} give
\begin{equation}
 \begin{split}
 K^A_U[J](u,v)
   ={}&\langle C_Uu,\mathscr A_U C_Uv\rangle
       -\langle h_U^{\mathsf T}J,
                         D^2q_U(0)[C_Uu,C_Uv]\rangle,\\
 K^{\rm rel}_U[J]
   ={}&K^A_U[J]+(G_UC_U)^{\mathsf T}(G_UC_U).
 \end{split}\label{eq:explicit-relative-current-extension-NX8}
\end{equation}
Indeed the $j_U D^2\Phi_U$ term from the first composition
is cancelled by $-\langle j_U,\Phi_U\rangle$.  The $J$
term supplies $\langle J,D^2\Phi_U\rangle$.  There is no
constraint-curvature term from the relative gauge because
its full-bond first derivative is zero.
All matrices are normalized as Hessians, with quadratic
exponent $-\frac12\langle u,Ku\rangle$.

For a fixed reference inverse coupling $\theta=g^{-2}$
and physical action weight $x$ put $\tau=x/\theta$.
The gauge part keeps its reference weight.  Consequently
the corresponding free precision is
\begin{equation}
 K_{U,\tau}^{\rm rel}[J]=\tau K^A_U[J]
                              +(G_UC_U)^{\mathsf T}(G_UC_U).
 \label{eq:explicit-relative-current-extension-NX9}
\end{equation}
The factor $\tau$ has not multiplied the gauge block.
This is the actual distinction needed in a complex-weight
comparison at a frozen frame.

\subsubsection*{Exact physical agreement and the noncritical linear term}

If $U$ is constrained-critical, $C_U$ parametrizes its
fibre tangent and therefore
\begin{equation}
 C_U^{\mathsf T}j_U=0 .
 \label{eq:explicit-relative-current-extension-NX10}
\end{equation}
At the physical current $J=j_U$, \eqref{eq:explicit-relative-current-extension-NX6} then gives, for
every free $u$ in the chart,
\begin{equation}
 S^A_U[j_U](u)=A_W(e^{i\Phi_U(u)}U)-A_W(U).
 \label{eq:explicit-relative-current-extension-NX11}
\end{equation}
This is an equality of the whole action difference.
It is stronger than equality of Hessians.

If $U$ is also contour-axial, $T_{y,x}(U)=1$.  Hence
${\cal G}_U(a)=\mathrm{GF}(e^{ia}U)$ identically on the
chart, and
\begin{equation}
 \begin{split}
 g^{-2}\{\tau S^A_U[j_U](gy)+{\cal G}_U(\Phi_U(gy))\}
 ={}&x\{A_W(e^{i\Phi_U(gy)}U)-A_W(U)\}\\
 &+\theta\,\mathrm{GF}(e^{i\Phi_U(gy)}U).
 \end{split}\label{eq:explicit-relative-current-extension-NX12}
\end{equation}
Thus the extension produces the actual weighted constrained
action and the reference gauge weight at this physical background.

For a noncritical $U$ the exact identity instead has
\begin{equation}
 S^A_U[j_U](u)
      =A_W(e^{i\Phi_U(u)}U)-A_W(U)
                           -\langle C_U^{\mathsf T}j_U,u\rangle.
 \label{eq:explicit-relative-current-extension-NX13}
\end{equation}
The missing physical linear force must be restored.
It has zero Hessian, so at $J=j_U$ the quadratic comparison
is still the physical chain-rule comparison.
But its contribution after Gaussian rescaling is of order
$g^{-1}$ unless that tangent force vanishes.  No
coupling-independent holomorphic potential claim for the
whole noncritical physical action follows from \eqref{eq:explicit-relative-current-extension-NX7}.

\subsubsection*{Gauge covariance from the actual graph}

For a site gauge $s$ the left logarithm transforms as
$a^s(b)=\operatorname{Ad}(s(b_-))a(b)$.  Denote this
full-bond linear map by ${\cal U}_s$ and its free restriction
by ${\cal U}_s^f$.  The actual average is bi-equivariant.
On common regular branches this gives
\begin{equation}
 q_{U^s}({\cal U}_sa)
        ={\cal U}_s^{\rm coarse}q_U(a).
 \label{eq:explicit-relative-current-extension-NX14}
\end{equation}
Free restriction and pivot insertion are stable under this
bondwise transformation.  The vector
${\cal U}_s\Phi_U(u)$ solves the transformed constraint
and has free restriction ${\cal U}_s^fu$.  Uniqueness of
the pivot graph proves
\begin{equation}
 \Phi_{U^s}({\cal U}_s^fu)
      ={\cal U}_s\Phi_U(u),\qquad
 C_{U^s}{\cal U}_s^f={\cal U}_sC_U .
 \label{eq:explicit-relative-current-extension-NX15}
\end{equation}
For complex gauges this assertion is made only on the
common admitted graph branch.  It does not say that a
complex orbit saturates a core.

Wilson gauge invariance gives
$j_{U^s}={\cal U}_sj_U$ under the invariant bilinear
colour pairing.  The product in \eqref{eq:explicit-relative-current-extension-NX5} is conjugated
at its contour's initial node.  Combining these
facts proves
\begin{equation}
 S^{\rm rel}_{U^s}[{\cal U}_sJ]({\cal U}_s^fu)
                            =S^{\rm rel}_U[J](u).
 \label{eq:explicit-relative-current-extension-NX16}
\end{equation}
The same holds separately for the action and gauge
parts, and hence for every fixed $\tau$ in \eqref{eq:explicit-relative-current-extension-NX9}.
In an orthonormal colour basis ${\cal U}_s^f$ is
bilinearly orthogonal even for admitted complex gauges.
For real unitary gauges it is also unitary in the
Hermitian Hilbert norm.  Therefore
\begin{equation}
 K_{U^s,\tau}^{\rm rel}[{\cal U}_sJ]
        ={\cal U}_s^f K_{U,\tau}^{\rm rel}[J]
                                      ({\cal U}_s^f)^{-1}.
 \label{eq:explicit-relative-current-extension-NX17}
\end{equation}
This proves matrix covariance directly.  It does not
use an integral change by a complex gauge through a
sharp real cutoff.

\subsubsection*{Axialization and the independent-current real gap}

For a real regular field, define within each block
$s(y)=1$ at its centre and $s(x)=T_{y,x}(U)$.
The internal contour means are unitary.  Their
bi-equivariance gives
\begin{equation}
 T_{y,x}(U^s)=1,\qquad \mathcal M(U^s)=\mathcal M(U).
 \label{eq:explicit-relative-current-extension-NX18}
\end{equation}
The second equality uses $s=1$ at the coarse centres.
The regular logarithm branches remain the same:
after this transformation all path transports are
near identity by the small-plaquette comparison, and
uniqueness of the internal mean identifies the transformed
mean.  This is the finite averaging chart's real
gauge-covariance construction.

Plaquette norms and the norm of the transformed current
are unchanged.  The gauge is periodic on a compatible
torus and is defined on the full cover as well.
No noncontractible holonomy is made small.

Let $c_0(L)$ be the proved flat gap.  Choose $R,\epsilon,\gamma$
as in the first-step gap theorem, in particular
\begin{equation}
 R\ge40,\quad K_L/R^2\le c_0/8,\quad
 CL^2R\epsilon\le r_b/8,\quad
 K_L(R\epsilon+\gamma)\le c_0/8,\quad
 K_L=CL^{40},\quad \gamma\le1 .
 \label{eq:explicit-relative-current-extension-NX19}
\end{equation}
Suppose all fine plaquettes have size at most $\epsilon$
and the real raw current has norm at most $\gamma$.
At the axialized field the relative and absolute
contour Hessians coincide.  The formula in \eqref{eq:explicit-relative-current-extension-NX8} is then
exactly the explicit axial auxiliary form of the
first-step gap theorem.  That theorem gives
$K_{U^s}^{\rm rel}[{\cal U}_sJ]\ge c_0/2$.
Unitary conjugation in \eqref{eq:explicit-relative-current-extension-NX17} proves
\begin{equation}
 K_U^{\rm rel}[J]\ge c_0/2
 \label{eq:explicit-relative-current-extension-NX20}
\end{equation}
for every such real regular $U$, without assuming that
$U$ was axial or critical.  The inequality uses the
Hermitian quadratic form of the real symmetric matrix
on its complexification.  The current is the declared
raw current of this one-step problem.

This transfer is possible because the relative gauge
is jointly site-gauge covariant.  Applying it to an
arbitrary absolute-gauge operator off the axial locus
would not be valid.

\subsubsection*{The full half-disc with the gauge weight retained}

There is a stronger weight statement than a small
operator-norm neighbourhood of $\tau=1$.  Choose
the universal constant in $K_L$ in \eqref{eq:explicit-relative-current-extension-NX19} to
control the action and gauge blocks separately
as well as their sum.  The same free-coordinate
partition proof gives
\begin{equation}
 K^A_U[J]\ge
       -K_L\{R^{-2}+R\epsilon+\gamma\}
                         \ge-c_0/4 .
 \label{eq:explicit-relative-current-extension-NX20a}
\end{equation}
Here the flat action block is nonnegative:
it is $C_0^{\mathsf T}\mathscr A_{\mathbf1}C_0$,
the compression of the free Wilson curl form.
In each localized patch the action block differs
from it by at most $K_L(R\epsilon+\gamma)$,
by the same actual chart coefficient comparison
used for the complete form.  Localizing the
action block on free variables costs at most
$K_L/R^2$.  Summing the squared partition
therefore proves \eqref{eq:explicit-relative-current-extension-NX20a}.  No positive gap
for the action block alone has been assumed.
Real gauge covariance transfers this statement
from the axial representative just as before.

Write $K=K^A_U[J]+K^G_U$, where
$K^G_U=(G_UC_U)^{\mathsf T}(G_UC_U)\ge0$.
\eqref{eq:explicit-relative-current-extension-NX20} and \eqref{eq:explicit-relative-current-extension-NX20a} imply
\begin{equation}
 B=K^{-1/2}K^A_U[J]K^{-1/2},\qquad
                    -\tfrac12I\le B\le I.
 \label{eq:explicit-relative-current-extension-NX20b}
\end{equation}
The upper bound follows from $K^G_U\ge0$.
The lower bound uses $K^{-1}\le2/c_0$.
In particular $B$ is a bounded self-adjoint
operator with norm at most one, and
\begin{equation}
 K_{U,\tau}^{\rm rel}[J]
   =K^{1/2}[I+(\tau-1)B]K^{1/2}.
 \label{eq:explicit-relative-current-extension-NX20c}
\end{equation}
For $|\tau-1|\le1/2$ this proves
$\operatorname{Re}(v,K_{U,\tau}^{\rm rel}[J]v)_{\rm H}
\ge\frac12(v,Kv)_{\rm H}\ge(c_0/4)\|v\|_2^2$
and inverse norm at most $4/c_0$.
The inverse is holomorphic on $|\tau-1|<1$
by the Neumann series for its middle factor.
The gauge block has stayed at its original weight.

In finite dimension let $\beta_a\in[-1/2,1]$
be the eigenvalues of $B$.  Then
$\det K_\tau=\det K\prod_a[1+(\tau-1)\beta_a]$.
Every factor has positive real part on that
disc.  Its logarithm continued from one
defines the coherent Gaussian determinant
branch.  This is a matrix normalizer
statement, not a zero-free theorem for an
interacting partition function with cutoffs.

One also obtains a useful absolute bound.
For $N$ real coordinates, $0\le\rho<1/4$,
and $|F(y)|\le F_0e^{\rho(y,Ky)}$,
\begin{equation}
 \left|\int F\,d\mu_{K_\tau}\right|
 \le F_0\left({3/2\over1/2-2\rho}\right)^{N/2},
                      \qquad |\tau-1|\le1/2 .
 \label{eq:explicit-relative-current-extension-NX20d}
\end{equation}
The determinant modulus is at most
$(3/2)^N\det K$, and
$\operatorname{Re}K_\tau-2\rho K
\ge(1/2-2\rho)K$.  Integrating this
real Gaussian proves the inequality.
Its cost is exponential in the actual
number of coordinates.  A component
mass estimate must pay that cost from
its stated reserve; it has not become
a volume-independent scalar bound.

\subsubsection*{A complex neighbourhood and exponentially local inverse}

The matrix has the fixed finite range of the actual
one-step chart.  Relative contour factors add no
new geometric range: both factors use the same
internal paths.  On the joint chart their analytic
values and inverses are bounded in a local unitary
frame.  The local Cauchy estimates and coordinate counts give
\begin{equation}
 \|K^{\rm rel}_{e^{iE}U,\tau}[J+Z]-K_U^{\rm rel}[J]\|_{\rm Sch}
       \le \widehat K_L
                   (\|E\|_\infty+\|Z\|_\infty+|\tau-1|),
       \qquad \widehat K_L=C'L^{40}.
 \label{eq:explicit-relative-current-extension-NX21}
\end{equation}
Here use $\|E\|_\infty\le r_b/16$, retained chart
room and total current at most one.  To justify
the same safe power, each scalar relative contour
term is bounded on the same free radius $cL^{-9}$
and background radius $c'L^{-5}$.  At most $CL^4$
terms affect one row and at most $CL^4$ columns
couple to it.  Its free Hessian costs $CL^{26}$
and one background difference $CL^{31}$.
The current curvature term has the previously
counted $CL^{17}$ current cost.  The action block
in the $\tau$ derivative has a bounded row sum
of this same safe order.  Increasing $C'$ gives
\eqref{eq:explicit-relative-current-extension-NX21}, independently of volume.

If the right side is at most $c_0/8$, the Neumann
series about \eqref{eq:explicit-relative-current-extension-NX20} gives inverse norm at most
$4/c_0$.  Finite range and its uniform row and
column sums also give a local exponential inverse.
For clarity let $d$ be the fine lattice distance
between free coordinates and let $\psi$ be any
bounded one-Lipschitz function.  For small $m>0$,
\begin{equation}
 \|e^{m\psi}Ke^{-m\psi}-K\|_{\rm Sch}
           \le C_L(e^{m r_0}-1),\qquad r_0\le12L .
 \label{eq:explicit-relative-current-extension-NX22}
\end{equation}
This follows entry by entry from the range bound.
Choose $m=m(L)>0$ so this perturbation is at most
$c_0/8$.  The conjugated complex matrix is still
invertible by comparison with the same real
matrix.  Apply its inverse bound with
$\psi(b)=\min\{d(b,b'),N\}$ and let $N$ increase.
One obtains
\begin{equation}
 \|(K^{-1})_{bb'}\|\le(4/c_0)e^{-m d(b,b')}.
 \label{eq:explicit-relative-current-extension-NX23}
\end{equation}
The distance is periodic distance on a torus.
Counting its shells in the four-dimensional
cover shows that row and column sums of this
kernel are bounded uniformly in compatible volume.
Consequently it is a two-sided inverse on bounded
fields as well.  The left identity uses the finite
range of $K$; the right identity follows by
absolute summation using that row bound.
This also proves uniqueness of a bounded solution.


The half-disc has a uniform complex-field
version as well.  Bound the action-block
and gauge-block differences separately
by the same local Cauchy estimates.
Since $|\tau|\le3/2$, their sum has
norm at most $C_L(\|E\|_\infty+\|Z\|_\infty)$,
uniformly on $|\tau-1|\le1/2$.
Make this norm at most $c_0/16$.
Comparison with the inverse in \eqref{eq:explicit-relative-current-extension-NX20c}
gives the inverse for the complex pair.
For the weighted conjugation choose
$m(L)>0$ so its additional row/column
cost is at most $c_0/16$.
The total perturbation of the same
real-pair half-disc operator is then
at most $c_0/8$.  Its inverse norm
and kernel satisfy
\begin{equation}
 \|K_{\mathbf U,\tau}^{\rm rel}[\mathbf J]^{-1}\|
       \le8/c_0,\qquad
 \|(K_{\mathbf U,\tau}^{\rm rel}[\mathbf J]^{-1})_{bb'}\|
       \le(8/c_0)e^{-m(L)d(b,b')}.
 \label{eq:explicit-relative-current-extension-NX23a}
\end{equation}
This argument is uniform in $\tau$
on the full closed half-disc.  The
smaller near-one comparison in \eqref{eq:explicit-relative-current-extension-NX21}
is a sufficient alternative; it is
not needed to restrict $\tau$ in
this stronger version.

\subsubsection*{Local intrinsic remainders and an admitted coupling disc}

There is a locally finite decomposition of \eqref{eq:explicit-relative-current-extension-NX6}
in which each scalar term has value and derivative
zero at zero.  For every Wilson plaquette take its
action difference minus its own raw gradient
paired with $\Phi_U$.  For every bond take
$\langle J_b,(\Phi_U-C_Uu)_b\rangle$.
Add each individual relative contour term.
The plaquette gradient terms sum to $j_U$,
so this is exactly \eqref{eq:explicit-relative-current-extension-NX6}.

Each term depends on a fixed neighbourhood of its
plaquette, bond or contour.  On the common
free ball of a fixed fraction of $r_u$ its
holomorphic supremum is bounded by a constant
$M_L$.  The unitary local frames and the
actual chart give this bound uniformly in
the real regular base, with the retained complex
perturbation and current ceilings.
Subtract half its local Hessian and denote
the result by $R_\alpha(u)$.  It vanishes
through degree two.  One-variable Cauchy
on the line through $u$ gives
\begin{equation}
 |R_\alpha(u)|\le
          2M_L(\|u\|_\infty/r)^3,\qquad
                         \|u\|_\infty\le r/2 ,
 \label{eq:explicit-relative-current-extension-NX24}
\end{equation}
for a fixed smaller radius $r=r(L)>0$.
Indeed the Taylor coefficients of order $n$
are bounded by $M_Lr^{-n}$ and the sum
from $n=3$ is geometric.

Let the fixed Gaussian field box have
$\|y\|_\infty\le b$, $b\ge1$.  For an auxiliary
coupling $q$ with $|q|b\le r/2$, the singularity
at zero in the following expression is removable:
\begin{equation}
 q^{-2}R_\alpha(qy),\qquad
 |q^{-2}R_\alpha(qy)|
                           \le 2M_L|q|b^3/r^3 .
 \label{eq:explicit-relative-current-extension-NX25}
\end{equation}
The gauge and action parts may have the separate
fixed coefficients $1,\tau$ of \eqref{eq:explicit-relative-current-extension-NX9}.
All constants are uniform for $\tau$ on its
retained disc.

The actual Haar and pivot-coarea factors
also give local analytic potentials.  Factor
their value at $u=0$ and take the holomorphic
logarithm of the ratio on a smaller common
chart.  Each factor and its inverse are
bounded there by the quantitative pivot
inverse and Haar chart.  Cauchy bounds
its logarithmic derivative on a half ball;
integration along the radial segment proves
$|\log(J_\alpha(u)/J_\alpha(0))|
\le C_L\|u\|_\infty/r$.
This does not analytically continue an absolute
value: the real coarea orientation is positive,
and the complex factor is its coherent determinant
inverse.

Group these fixed-range terms into their actual
local carriers at a fixed block size $M$.
Only $C_{L,M}$ terms enter a carrier or
a rooted carrier sum.  The potential weight
$Q_Y(y)\ge1$ and the bounded tree length of
these carriers then give
\begin{equation}
 \|V_{q,\tau}\|_\lambda
                  \le C_{L,M,\lambda}|q|b^3/r^3 .
 \label{eq:explicit-relative-current-extension-NX26}
\end{equation}
This is a potential-norm estimate, independent of
the number of carriers.  For any fixed allowed
potential budget $w_*$ one can therefore take
an admitted disc $|q|<c_{L,M,\lambda,w_*}b^{-3}$,
also shrunk to ensure $|q|b<r/2$.  This is an
actual lower bound on an admissible radius for
these first-step intrinsic and measure potentials.
It is not just an upper restriction on a
radius that has not been constructed.

\subsubsection*{The finite integral and the factors that remain outside}

Here $I$ denotes the constrained integral with this specified
soft gauge weight and regular field box.  Its identification
with a regular part of the original Wilson measure uses
the actual gauge-fixing identity and its prefactor;
neither is created by the pivot change of variables.

At a physical critical axial pair, let
$\mathcal J_U(u)$ be the whole real Haar/coarea
density in the free chart and let $N$ be the
number of real free coordinates.  On its
fixed Gaussian box the exact integral is
\begin{equation}
 \begin{split}
 I={}&e^{-xA_W(U)}g^N\mathcal J_U(0)
             (2\pi)^{N/2}(\det K_{U,\tau}^{\rm rel}[j_U])^{-1/2}\\
 &\quad\cdot\int\chi(y)e^{V_{g,\tau}(y)}
                              d\mu_{K_{U,\tau}^{\rm rel}[j_U]}(y).
 \end{split}\label{eq:explicit-relative-current-extension-NX27}
\end{equation}
This follows from \eqref{eq:explicit-relative-current-extension-NX12}, the exact chart
coarea factor and subtraction of its quadratic
exponent.  The nonconstant potential consists
of the negative intrinsic remainders divided
by $g^2$ and the logarithm of
$\mathcal J_U(gy)/\mathcal J_U(0)$.
The Gaussian determinant uses its coherent
branch from the real positive matrix.

The prefactor is independent of the integrated
field $y$, but generally depends on the coarse
datum $U$ and on the specified current and
weight.  It is not automatically a numerical
vacuum constant.  Its field-dependent part
still needs the assigned localization; its
value at the reference pair is counted once
in the scalar normalization.  The powers
of $g$ have been factored explicitly before
introducing the auxiliary analytic coupling
$q$ in \eqref{eq:explicit-relative-current-extension-NX25}.

The fixed box may be chosen as a product of
colour-norm balls, so real site gauges preserve
it.  Complex gauge covariance of matrices,
already proved in \eqref{eq:explicit-relative-current-extension-NX17}, is not by itself
a contour-deformation theorem for a sharp
real integral.  Such an integral's local
holomorphic gauge identity follows on an
admitted connected continuation from its
real gauge identity; arbitrary full-core
orbit descent requires the additional
domain comparison.

We have supplied an explicit consistent
first-step extension, the physical action
identity, its real gap, complex inverse,
and the intrinsic potential radius.
The enlarged decoupling matrices, integrated
bond-root bounds, all-substitution core
admission, localization of the prefactor,
and the many-level effective action and
history operations remain separate.
If this extension is used beyond the
first step, each of them must use the
same definition and prove its physical
sum identity.  It is not legitimate
to import estimates for a differently
defined current completion without
the complete comparison.  Conversely,
the named-completion comparison need
not be imposed in addition when the
entire alternate construction has
actually been rebuilt from \eqref{eq:explicit-relative-current-extension-NX6}.


% END MODULE chapters/explicit_relative_current_extension.tex

% BEGIN MODULE chapters/multilevel_background_reductions.tex
\subsection{Constrained backgrounds at several scales}
\label{candidate:background}
\label{ch:multilevel-backgrounds}

A background is a regular critical orbit of the fine Wilson action with
its assigned block averages, not a new independent coarse action. Its
construction has three separate ingredients. The determining rows and
their residual gauge group are fixed geometrically. The critical equations
are then solved on a slice by a coupled field and multiplier argument.
Finally a local regularity improvement identifies the solution among all
regular orbits and gives the bounds needed for restriction and analytic
continuation.

This section develops the geometric and nonlinear construction using
the actual finite averaging map. Its proper-region inverse and local
coefficient estimates are proved subsequently; their application to
all sufficiently small regular mixed data is completed in
\cref{bgc:proper-background}. The construction includes arbitrary
regular competitors, the action-boundary current, and the outward
node factors. The physical constrained Hessian agrees with an
independently defined current operator only when their complete
projected forms agree.

The six background statements concern the regular critical
orbit, local regularity, complex continuation, response, hereditary
complex domains, and restriction with boundary comparison. A proved
reduction to an explicitly stated primitive estimate is distinguished from
a proof of that estimate.

\subsubsection{Determining rows and the geometry of a seam}

Fix an odd integer \(L\ge3\) and a depth \(k\). Work in level-\(k\)
units, so the bare spacing and intermediate spacings are
\begin{equation}
 \eta=L^{-k},\qquad \xi_j=L^{j-k},\qquad 0\le j\le k.
 \label{mb:scales}
\end{equation}
The regions \(\Omega_0\supset\cdots\supset\Omega_k\) are unions of
the assigned aligned cubes of side \(\mathfrak M\xi_j\), where
\(\mathfrak M\) is an odd power of \(L\). Set \(\Omega_{k+1}=\varnothing\).
Half-open cells specify boundary membership. The separation condition is
\begin{equation}
 \operatorname{dist}_\infty(\Omega_j^c,\Omega_{j+1})
                     >\rho\xi_j,\qquad \rho\ge C_{\rm geom}L.
 \label{mb:separation}
\end{equation}
The number \(C_{\rm geom}\) is chosen large enough for the fixed
one-step path neighborhoods and collars used below. This condition
does not require a region to be connected or simply connected.

Because all cubes are aligned and their sizes are powers of \(L\),
each \(\Omega_j\) is a union of the blocks \(B^j(y)\) of side
\(\xi_j\), and each \(\Omega_{j+1}\) is also a union of those blocks.
A top is a site \(y\in\mathbb T^{(j)}\cap
(\Omega_j\setminus\Omega_{j+1})\); its pillar is \(B^j(y)\).
These pillars partition \(\Omega_0\). To prove this, take a bare
cell and the largest \(j\) for which it lies in \(\Omega_j\).
Its level-\(j\) block is contained in \(\Omega_j\), while
\(\Omega_{j+1}\), being a union of level-\(j\) blocks, contains
none of that block. Its center is therefore a top of level \(j\).
Two such blocks at one level are disjoint. Blocks at different levels
are nested if they meet, which is incompatible with the corresponding
disjoint layers. This proves existence and uniqueness of its pillar.

For a bare bond or plaquette \(e\), let \(\mathsf L(e)\) be the largest
top scale attached to its corners. The numerical action and regularity
tests use
\begin{equation}
 A_{\Omega_0}(U)=\sum_{p\subset\Omega_0}
       \left(1-\tfrac12\operatorname{tr}\partial U(p)\right),\qquad
 \eta^{-2}\mathsf L(p)^2|\partial U(p)-1|,
 \quad \mathsf L(b)^3|J_A(U)(b)|.
 \label{mb:action-regularity}
\end{equation}
Only inner plaquettes contribute to this action. Its boundary current
therefore need not contain all six plaquettes of a bond. The current
is its actual derivative with the length and pairing normalization of
\cref{ch:domain-core-ledger}.

The determining bonds of level \(j\) are
\begin{equation}
 \begin{gathered}
 \Gamma_j=\left\{c\subset\Omega_0:
   \begin{array}{l}
     c\text{ is a level-}j\text{ bond},\
        c\text{ touches }\Omega_j,\\
     c\text{ has no endpoint in }\Omega_{j+1}
   \end{array}\right\},\\
 \mathbf B=\bigcup_{j=0}^k\Gamma_j .
 \end{gathered}
 \label{mb:determining-rows}
\end{equation}
No inward row is kept. A row whose two endpoints lie in \(\Omega_j\)
is interior. An outward row has one endpoint outside \(\Omega_j\).
Every endpoint is a top, possibly of the preceding level.

Here is the precise local geometry behind the last assertion. Let
\(c=(y,x)\in\Gamma_j\), \(j\ge1\), with
\(y\in\Omega_j\) and \(x\notin\Omega_j\).
Their distance is \(\xi_j=L\xi_{j-1}\). By \eqref{mb:separation}, neither endpoint
nor its fixed double-block neighborhood can leave \(\Omega_{j-1}\).
The endpoint \(x\) consequently belongs to
\(\Omega_{j-1}\setminus\Omega_j\), and it is a site of
\(\mathbb T^{(j-1)}\). It is a top of level \(j-1\), even when it
also happens to be a site of a coarser lattice. Its centered
level-\(j\) block is entirely outside \(\Omega_j\), because that
region is a union of level-\(j\) blocks. The separation puts this
block entirely in \(\Omega_{j-1}\). Hence every level-\((j-1)\)
bond internal to that block is a prescribed interior row of
\(\Gamma_{j-1}\).

The same argument shows why a row's double block reads no third
layer. Each endpoint block misses \(\Omega_{j+1}\), by the no-inward
condition and block alignment, and the neighborhood is inside
\(\Omega_{j-1}\) by separation. A one-step seam calculation therefore
contains the levels \(j-1,j\), whose length ratio is the fixed number
\(L\). It does not contain a product of ratios over all earlier
levels. Larger constructions must explicitly preserve these determining
neighborhoods; the assertion is not true for an arbitrary enlarged
box at a seam.

The datum and residual group are
\begin{equation}
 \begin{aligned}
 M_{\mathbf B}(U)&=(\mathcal M^jU|_{\Gamma_j})_{j=0}^k,
 &\mathcal F(\mathbf V)&=\{U:M_{\mathbf B}(U)=\mathbf V\},\\
 \mathcal G_{\rm res}&=\{g:g(y)=1\text{ at every node}\}.
 \end{aligned}
 \label{mb:fibre-group}
\end{equation}
Exact gauge covariance of the iterated average proves that this
group preserves the fibre. Its definition involves all nodes, not
only the centers of the last lattice.

\paragraph{Test loops and the missing inward rows}

At an ordinary coarse plaquette the datum loop uses \(V_j\).
Where the loop crosses into the preceding layer it uses the actual
one-step averages of \(V_{j-1}\) on those edges. These are the mixed
test loops. Their values transform by conjugation at their starting
node; their distance from one is invariant under real unitary gauges.
This is the local compatibility information that a regular datum
supplies. It is not a bound on all individual \(V_j(c)-1\).

If inward rows are retained, their endpoint in \(\Omega_{j+1}\)
becomes an additional seam node. There is then a genuine distinction
between fixing all nodes and imposing a full block-axial gauge in the
inner block. The following example is exact. Take an inner top \(y\)
and an additional seam node \(x\ne y\) in its block. Choose a site
gauge with \(g(y)=1\) and \(g(x)\ne1\), as close to one as desired,
and set \(U=1^g\). All plaquettes and the action current are zero.
All closed datum tests are also one. But
\begin{equation}
 T_{y,x}(\mathcal M^jU)=g(y)g(x)^{-1}\ne1 .
 \label{mb:seam-node-obstruction}
\end{equation}
Every residual transformation fixing both nodes leaves this transport
unchanged. Thus this pure-gauge orbit cannot be made fully block axial
there. The example does not obstruct an axial gauge \emph{outside}
the nodes; that gauge is proved next. It explains why a step that
requires a full inner block-axial representative omits the inward rows,
or else states its exceptional-node gauge explicitly.

With two kept inward edges meeting the same seam node, compatibility
also records their relative frames. In an adapted real datum frame one
can write these as products of a coarse edge, the internal fine contour
transport and the kept inward edge. The product of one such route with
the inverse of another is a closed two-scale loop. It is a datum test
independent of the fine representative. Omitting these tests while
retaining the edges changes the regular-datum class; elementary
plaquettes in each separate layer do not record their relative frame.

\subsubsection{An exact representative outside the nodes}

The residual orbit has a canonical comb representative independently
of any elliptic equation. A comb is the fixed coordinate-order path
inside a block. The only properties needed here are locality,
bi-equivariance and the identity transport at a zero-length path.
The construction works as well with an admitted mean transport
having those properties.

For a nonnode \(x\), let \(n(x)\) be its largest lattice index:
\(x\in\mathbb T^{(n)}\setminus\mathbb T^{(n+1)}\).
Its parent \(z(x)\) is the center of its level-\((n+1)\) block.
Its ascent terminates at its pillar top. We seek a residual gauge
such that, for every such site,
\begin{equation}
 (\mathcal M^{n(x)}U)(\Gamma_{z(x),x})=1.
 \label{mb:comb-axial}
\end{equation}
Process the indices in descending order. At index \(n\), leave the
gauge one at all nodes and at \(\mathbb T^{(n+1)}\). At a nonnode
of exact index \(n\), assign it the current transport
\((\mathcal M^nU)(\Gamma_{z(x),x})\). Extend this gauge by one at
the sites not being processed. Bi-equivariance makes that transport
one. Gauge covariance of \(\mathcal M^m\), \(m>n\), sees only
values at \(\mathbb T^{(m)}\), which have not changed. Thus previously
imposed conditions remain true. After finitely many stages every
condition in \eqref{mb:comb-axial} holds, and every prescribed node value has stayed one.

Uniqueness is equally direct. Suppose two representatives in one
residual orbit satisfy \eqref{mb:comb-axial}, and let \(h\) take one to the other.
It is one at the tops. Descending an ascent, if \(h(z)=1\), covariance
of a transport already equal to one gives
\begin{equation}
 1=h(z)\,1\,h(x)^{-1},\qquad h(x)=1.
 \label{mb:axial-uniqueness}
\end{equation}
Every nonnode is reached, so \(h=1\). The representative is therefore
a function of the orbit. On real configurations the construction is
Borel: the total average is Borel, the comb transports are finite
products, and all inversions and multiplications are continuous
group operations. On an admitted regular complex chart it is
holomorphic, since there the assigned averages are holomorphic.

This argument gives no depth-independent bound for the derivative
of the gauge. A real pure gauge may have arbitrary values at fine
sites while fixing the tops. Norm bounds on a slice require the
separate estimates below. Likewise a complex representative is
used only when its values remain in the admitted analytic branches
and its condition numbers have the required bounds. The Borel orbit
construction alone does not establish that admission.

For a sequence of depth \(k+1\) with no inward rows, an inner
level-\((k+1)\) block has no node except its center. Indeed every
lower-level site of the block lies in \(\Omega_{k+1}\), so cannot
be a top of a preceding layer, and alignment supplies no other
coarse center in that block. A residual gauge can consequently
set all its level-\(k\) internal contour variables to one while
fixing the datum. This is the exact node-counting fact used in the
definition of the next step's physical background.

\subsubsection{The reduced datum has an explicit local inverse}

Fix a real reference field \(U_0\) whose used averages lie on the
regular branches of the actual finite mean. For an interior row
the reduced datum is \(b(B)(c)=B(c)\). On an outward row oriented
toward its outer endpoint \(z\), it is
\begin{equation}
 e^{ib(B)(c)}=e^{iB(c)}Z_c(B_{\rm int}),\qquad
 Z_c(B_{\rm int})
       =\operatorname{Ad}(V_0(c))\,v_z(B_{\rm int}).
 \label{mb:reduced-datum}
\end{equation}
The factor \(v_z\) is the actual one-step block factor made from
the interior data in the finer block centered at \(z\).
The geometry following \eqref{mb:determining-rows} proves that these are prescribed
interior rows, not another set of unknown outward rows.
For the reverse orientation use the inverse group identity.

There is therefore a direct inverse. Recover \(B_{\rm int}=b_{\rm int}\),
then set
\begin{equation}
 B(c)=\frac1i\log\big(e^{ib(c)}
                         Z_c(b_{\rm int})^{-1}\big)
 \quad\text{on outward rows}.
 \label{mb:reduced-inverse}
\end{equation}
All logarithms are on the common small datum chart. Conjugation by
the real \(V_0(c)\) has norm one; it imposes no smallness of that
coarse edge or of a global holonomy. The finite mean construction
gives a radius \(r_0(L)>0\) on its admitted local reference class.
On a fixed smaller ball, its block factors and their first two
derivatives are bounded by constants depending on that finite
construction. Thus \eqref{mb:reduced-datum}--\eqref{mb:reduced-inverse} are admitted on balls
\(\|B\|_\infty<r_0/C_L\), respectively
\(\|b\|_\infty<r_0/C'_L\), and satisfy
\begin{equation}
 \|b(B)\|_\infty\le C_L\|B\|_\infty,\qquad
 \|b^{-1}(b)\|_\infty\le C'_L\|b\|_\infty.
 \label{mb:reduced-bounds}
\end{equation}
The uniform statement is over local reference patches on which that
same finite chart and its radius are admitted. Establishing those
patches for a multilevel background is part of its regular-domain
estimate; it is not inferred just from the existence of \(U_0\).

It is useful to see exactly why the inverse does not contain a
long product of triangular matrices. Define the right logarithmic
differential
\(\mathcal E_X(H)=(1/i)(d/dt)e^{i(X+tH)}|_{t=0}e^{-iX}\).
Differentiating \eqref{mb:reduced-datum} gives the outward diagonal
\(D_c=\mathcal E_{b(c)}^{-1}\mathcal E_{B(c)}\), not the identity
in general. Let \(D\) be this block diagonal matrix, with identity
on interior rows, and let \(T\) contain all derivatives of block
factors. Its only nonzero blocks have interior columns and outward
rows. Consequently
\begin{equation}
 Db=D+T,\qquad TD^{-1}T=0,\qquad
 (Db)^{-1}=D^{-1}-D^{-1}TD^{-1}.
 \label{mb:seam-datum-nilpotence}
\end{equation}
Multiplying the last two matrices verifies both inverse identities.
This nilpotence is a property of the actual reduced datum. An
arbitrary lower triangular coupling between successive levels would
not have it.

For the row count, one outward factor reads the \(O(L^4)\) interior
data of its single finer block. For the column count, an interior
bond belongs to one assigned finer block, and that block can be
the outer endpoint block of at most eight coarse incident rows.
Thus both row and column counts are bounded at fixed \(L\),
independently of the number of layers. After inserting the actual
cell weights and amplitudes, an interaction changes scale only
from \(j-1\) to \(j\); its weight ratios are fixed powers of \(L\).
\eqref{mb:reduced-bounds}--\eqref{mb:seam-datum-nilpotence} therefore also give two-sided weighted
Schur bounds in those declared norms. Higher derivatives follow
from the same finite group formula and its fixed datum disc.

This local datum calculation is different from constructing a
right inverse of all the field constraint rows. The latter must
also lift fields and control their energy and gauge slice at a
seam. That primitive linear task is retained in the precision
section. The exact derivative \(Db\), including its nonidentity
outward diagonal, is the datum forcing in the response equations.

\subsubsection{Real reconstruction on one block}

The next estimates are finite group calculations. They concern real
unitary fields on a rectangular lift of one or two blocks. They do not
require those fields to be close to one on every bond. Let
\(\varpi(U)=\sup_p|\partial U(p)-1|\) on that neighborhood.

A coordinate-order path can be changed to another such path by
interchanging adjacent steps in different directions and canceling
backtracks. Each interchange inserts one plaquette, conjugated by the
transport of its preceding path. On a box of side \(O(L)\), the
number of interchanges is bounded in terms of \(L\), and more specifically
at most a dimensional constant times \(L^2\) for the fixed path
families used here. Thus a contractible comparison loop \(\gamma\)
in this family is a product of transported plaquette factors and
\begin{equation}
 |U(\gamma)-1|\le C_L\varpi(U).
 \label{mb:finite-stokes}
\end{equation}
Indeed for unitary matrices \(|AB-1|\le|A-1|+|B-1|\), and
conjugation by a unitary transport preserves the norm. Inducting
over the inserted factors proves \eqref{mb:finite-stokes}. The constant is not
multiplied by the size of an individual open transport. This is
why arbitrary real holonomies cause no loss in the local argument.

Fix the comb from the block center. The relative transports of
the coordinate-order path family differ from that comb by \eqref{mb:finite-stokes}.
The local mean lemma, on its fixed regular radius, then implies
\begin{equation}
 |T_{y,x}(U)-U(\Gamma_{y,x})|\le C_L\varpi(U).
 \label{mb:mean-versus-comb}
\end{equation}
To see that the mean lemma applies without a small open transport,
multiply every member on the right by the inverse of the comb
transport. The resulting family is near one. The implicit
logarithmic mean has a uniformly bounded derivative on a smaller
fixed ball, by its finite-dimensional inverse construction.
Its value therefore differs from one by \(C_L\varpi(U)\).
Right equivariance restores the original transport and proves
\eqref{mb:mean-versus-comb}.

For a fine bond \(b\), let \(y_-,y_+\) be the centers of its endpoint
blocks and use the same center when both endpoints lie in one block.
Define its transported value by
\begin{equation}
 \widetilde U(b)=T_{y_-,b_-}(U)\,U(b)\,
                            T_{y_+,b_+}(U)^{-1}.
 \label{mb:transported-bond}
\end{equation}
If the centers coincide, this is a comparison loop, so
it is within \(C_L\varpi(U)\) of one. Otherwise the centers are
joined by a coarse bond \(c\). Replace the means by combs using
\eqref{mb:mean-versus-comb}. The resulting open route and the central straight route
from \(y_-\) to \(y_+\) differ by a loop of the type in \eqref{mb:finite-stokes}.
The closed comparison transports in the definition of
\(\mathcal M(U)(c)\) satisfy the same bound. The finite exponential
and logarithm estimate consequently gives
\begin{equation}
 |\widetilde U(b)-\mathcal M(U)(c)|\le C_L\varpi(U)
       \quad (y_-\ne y_+).
 \label{mb:bond-reconstruction}
\end{equation}
All paths used are in the fixed double block. The regularity
threshold makes the closed logarithms well defined; the final
coarse bond itself may be far from one.

For two such fields \(U,U_0\), put
\(d_T(x)=|T_{y,x}(U)-T_{y,x}(U_0)|\), using the endpoint's block
center. Combining \eqref{mb:transported-bond}--\eqref{mb:bond-reconstruction} and using unitarity gives
\begin{equation}
 \begin{split}
 |U(b)-U_0(b)|\le{}&
    |\mathcal M(U)(c)-\mathcal M(U_0)(c)|
       +d_T(b_-)+d_T(b_+)\\
 &+C_L\{\varpi(U)+\varpi(U_0)\}.
 \end{split}
 \label{mb:rigidity-defect}
\end{equation}
For a bond internal to one block, omit the coarse-bond difference.
The coefficient of the coarse difference is exactly one, since
it is multiplied only by unitary transports. In a comparison where
the contour means agree the two \(d_T\) terms vanish. If only the
combs agree, \eqref{mb:mean-versus-comb} absorbs them into the curvature error.

An iteration with this coefficient is qualitatively different
from an iteration with an unspecified constant in front of the
coarse error. If a bond is in a pillar of height \(j\), and the
regularity at intermediate index \(n<j\) is
\(\varpi_n\le\varepsilon L^{2(n-j)}\), \eqref{mb:rigidity-defect} yields, along an
aligned descent,
\begin{equation}
 e_n\le e_{n+1}+C_L\varepsilon L^{2(n-j)},\qquad
 e_n\le e_j+\frac{C_L\varepsilon}{L^2-1}.
 \label{mb:rigidity-sum}
\end{equation}
The second inequality is obtained by summing the first. An
uncontrolled endpoint contour defect would have to be added at
each step; it cannot be omitted. For the preliminary gauge on a separated real sequence, let
\(W_n=\mathcal M^nU\) and \(W_{0,n}=\mathcal M^nU_0\).
Begin with \(g_k=1\) and descend. Set \(g_n=1\) on the
tops \(\Lambda_n\) of that layer. On the block of an interior
next-level center \(y\), set
\[
 g_n(x)=W_{0,n}(y,x)^{-1}g_{n+1}(y)W_n(y,x).
\]
These blocks and the layer tops partition the level-\(n\)
sites of \(\Omega_n\). Since the contour from \(y\) to itself
is one, the formula extends \(g_{n+1}\) and gives a single
gauge on the bare sites. Covariance of the mean then makes
the two contour values equal in every such block. The
interior and outward determining bonds have both endpoints
at tops, so their original datum mismatch is unchanged.

At an inward bond orient it towards the seam node. Its
corridor belongs to the outward bond one level above,
whose outer endpoint block lies wholly in \(\Lambda_n\).
The inward end uses the aligned contour. At the outer end
each contour reads at most \(2(L-1)\) interior bonds of
the recorded block. With a mismatch at most \(\alpha_1\)
on the proper determining rows, telescoping along the
path and replacing the mean by the path costs at most
\(2(L-1)\alpha_1+C_L\alpha_0\). Here \(\alpha_0\) bounds
the regularity radius and the double block lies in
\(\Omega_n\). Including the outward bond one level above
gives the inward bound \((2L-1)\alpha_1+C_L\alpha_0\).
This applies whether the inward row is retained or omitted.
For a bond with both ends in a deeper layer, the endpoint
blocks follow the same largest common layer as the bond,
by alignment. Its double-block curvature at level \(n\)
is bounded by \(C\alpha_0L^{2(n-j)}\) when \(j\) is that
largest layer. Applying \eqref{mb:rigidity-sum} from the
proper or inward bound at \(j\) then gives a uniform
bound \(C_L(\alpha_0+\alpha_1)\) at every carried link.

The estimate in \eqref{mb:rigidity-sum} is a bound on raw relative group elements.
It does not divide those differences by the fine spacing and
does not prove a bound on a first difference quotient. The
elliptic slice estimate used later is precisely what improves
this preliminary closeness to the natural physical norm.

\paragraph{The sharp linear curvature factor}

A rough application of \eqref{mb:finite-stokes} to each coarse edge gives
\(C_L\varpi\), which is not enough for a uniform iteration in
physical units. The leading coefficient can be computed exactly.
On the same finite rectangular patch take a comb tree gauge.
Each remaining bond is a comparison loop with the tree, hence
is within \(C_L\varpi\) of one. For small \(\varpi\) write it as
\(e^{ia_b}\) with \(\|a\|_\infty\le C_L\varpi\).
Taylor expansion of the finite products gives
\begin{equation}
 |\mathrm{d}a(p)|\le\varpi+C_L\varpi^2,\qquad
 \log\partial\mathcal M(e^{ia})(p')/i
         =\mathrm{d}_c(P_La)(p')+O(C_L\varpi^2).
 \label{mb:curvature-linearization}
\end{equation}
The first linear term is the discrete curl; every other product
term has at least two \(a\)'s. The same is true of the finite
implicit mean and logarithm on their fixed chart. These statements
are estimates for a fixed number of coordinates depending on \(L\),
so the quadratic constants are finite and independent of depth.

The flat identity \(P_L=S_L-\mathrm{d}_c\Sigma_L\) implies
\(\mathrm{d}_cP_L=\mathrm{d}_cS_L\). For a \(\mu\nu\) coarse plaquette,
direct telescoping of the straight sums gives the exact flux formula
\begin{equation}
 (\mathrm{d}_cS_La)_{\mu\nu}(y)
   =L^{-4}\sum_{x\in B(y)}
       \sum_{r,s=0}^{L-1}
       (\mathrm{d}a)_{\mu\nu}(x+re_\mu+se_\nu).
 \label{mb:sheet-flux}
\end{equation}
For example the difference of the two \(\nu\)-edge sums separated
by \(Le_\mu\) telescopes in the \(\mu\) direction; the other two
edge sums telescope in \(\nu\) with the opposite sign. Their
combination is the displayed sheet curl. There are \(L^4L^2\)
terms with weight \(L^{-4}\), so its supremum norm is at most
\(L^2\|\mathrm{d}a\|_\infty\). Finally the exponential changes a small
coarse logarithmic curvature into its group value with a quadratic
error. Together with \eqref{mb:curvature-linearization} this proves the real one-step estimate
\begin{equation}
 \varpi(\mathcal M U)\le L^2\varpi(U)+C_L\varpi(U)^2.
 \label{mb:one-step-curvature}
\end{equation}
The tree gauge was local and unitary, so the conclusion holds
in the original frame.

For a homogeneous regular patch let
\(\varpi(U)\le\varepsilon L^{-2k}\). Suppose all determining
neighborhoods stay in that patch. Put
\(v_n=L^{2(k-n)}\varpi(\mathcal M^nU)\). \eqref{mb:one-step-curvature} gives
\begin{equation}
 v_{n+1}\le v_n+
      C_LL^{-2(k-n)-2}v_n^2.
 \label{mb:curvature-iteration}
\end{equation}
The coefficients in this nonlinear recursion have a sum bounded
independently of \(k\). Choose \(\varepsilon\) small enough that
four times this sum times \(\varepsilon\) is less than one.
Inductively, if \(v_q\le2\varepsilon\) up to \(n\), summing the
recursion gives \(v_{n+1}<2\varepsilon\). Thus
\begin{equation}
 |\partial\mathcal M^nU-1|
                  \le2\varepsilon L^{2(n-k)}
                      \quad(0\le n\le k).
 \label{mb:uniform-real-average}
\end{equation}
The same induction establishes admission to every one-step
regular mean branch, by reducing \(\varepsilon\) below its fixed
threshold. On a proper region it is used only at points whose
full determining neighborhood remains in the region.

Across a seam the comparison uses the finer of the two adjacent
length scales. A curvature measured in the coarser scale can
therefore lose a fixed factor \(L^2\). The geometry of \eqref{mb:determining-rows}
permits that local loss and no chain of unexamined intermediate
layers in one row. Proving an entire tower or joined-region
statement still requires its determining neighborhoods to have
the advertised alignment and collars. The formula in \eqref{mb:uniform-real-average} is the
averaging estimate; it does not solve a constrained variational
problem.

\subsubsection{A domain specified before the unknown background}

An inverse estimate cannot be used to define its own range of
application. We therefore first specify a class of possible reference
fields without referring to the unknown solution.
On each fixed determining patch use a real unitary local gauge and
ask that its elementary plaquettes satisfy the bound in \eqref{mb:action-regularity},
with a small specified constant \(\epsilon_0\).
The bound is a test of the given field. Its coarse tests are obtained
from the actual averages. For real homogeneous subpatches \eqref{mb:uniform-real-average}
proves admission to their regular branches. At a seam the preceding
two-scale neighborhood is tested, and its finite averaging factors
are tested in the same regular chart. The latter tests are explicit
finite group inequalities, not assertions about the critical field.

Let \(\mathcal R_{\rm lin}(\epsilon_0,\gamma_0)\) consist of
such real fields and independent currents with
\begin{equation}
 \sup_{p\subset\Omega_0}
      \eta^{-2}\mathsf L(p)^2|\partial U(p)-1|<\epsilon_0,
 \qquad
 \sup_b\mathsf L(b)^3|J(b)|<\gamma_0 .
 \label{mb:primitive-real-domain}
\end{equation}
Include the actual finite-chart and row-rank inequalities needed
by the chosen slice, and name their radii. These conditions can
be checked on \((U,J)\) before stationarity is imposed.
For a physical trial field the current used in a physical
curved-fibre Hessian is \(J_A(U)\). A separately prescribed \(J\)
in \eqref{mb:primitive-real-domain} is not set equal to it. The two domains may coincide
on a physical submanifold while their off-submanifold
completions differ.

Complex fields are introduced relative to one such real reference
by \(U'=e^{i\eta A'}U\). Their domain records, separately, bounds
on the physical logarithm and on every actual finite averaging
or gauge formula used:
\begin{equation}
 \|\mathsf L A'\|_\infty<r_A,\qquad
 \|\mathsf L^2 D_UA'\|_\infty<r_D,\qquad
 \|\mathsf L^3 J'\|_\infty<r_J,\qquad
 \mathcal Q_\beta(A',J')<r_\beta .
 \label{mb:primitive-complex-domain}
\end{equation}
Here \(\mathcal Q_\beta\) ranges over a declared finite list of
local readouts and atomic branch tests, including the outward
datum factors, slice maps, current readouts and inverses actually
used. The tests are normalized by length and scale, not by their
target radii. The separate inequalities and their margins are
essential: a point satisfying the curvature inequality need not
satisfy a complex logarithm or multiplier inequality.

The linearized operator has to be defined on this primitive
domain. If its definition uses a nonlinear background at a
smaller scale, that background must already have been constructed
in the induction, with its own admitted domain and margins.
It cannot use the background presently being solved for. The
induction is a directed finite list: an action or readout at
level \(j\) reads only completed levels below \(j\), and its
new constraint rows are the specified level-\(j\) rows.
This is a condition on the actual construction, not on notation.

There are two useful consequences of this arrangement. A primitive
matrix theorem on \eqref{mb:primitive-real-domain}--\eqref{mb:primitive-complex-domain} can be applied at a noncritical
trial field. Also the complex inverse can be continued on an
entire admitted ball by the two-sided Schur argument, rather
than being asserted only along the eventual solution. Both
consequences are needed in a first-exit proof of admission.

\paragraph{A regular section and what it must prove}

A regular datum is defined by its elementary and two-scale
test loops. Before the variational construction one needs a
configuration realizing those data sufficiently accurately, or
exactly, in the primitive domain.
There are four concrete tasks in such a construction.

First, on the finite coarse patches, choose local unitary frames
from the datum transports. Their transition matrices must obey
the actual node cocycle. Products around contractible test loops
are small, while noncontractible holonomies are retained as
unitary transition matrices. Replacing them by the identity
would change the fibre.

Second, extend the local coarse connection through a block with
smooth fields on its free coordinates. The required estimates
are on the physical potential, its first differences and its
current or covariant Laplacian, at that block's scale. A raw
face-supported lift need not have these estimates. For instance
a jump in its plaquette field on a coarse face can produce a
fine current of order \(\eta^{-1}\) times that jump.

Third, evaluate the \emph{actual} finite mean of this trial field.
The discrepancy from the datum is a named local function
\(\mathfrak e(B)\), not the discrepancy of a linear surrogate.
The finite chart gives its derivative and quadratic remainder.
The linear row right inverse must map this discrepancy into the
same potential and current norms.

Fourth, the local corrections must agree on overlaps and leave
previously frozen rows unchanged. At a seam this includes the
outward factor in \eqref{mb:reduced-datum} and its interior dependence, not merely
the interior average.

The nonlinear correction in the third task is elementary once
its linear and finite coefficient estimates have been proved.
In declared field and datum norms let
\begin{equation}
 \begin{aligned}
 q(A)&=Q_0A+N_q(A),\qquad Q_0E=1,\qquad
 \|E b\|_X\le C_E\|b\|_Y,\\
 \|N_q(A)-N_q(A')\|_Y
 &\le C_q(\|A\|_X+\|A'\|_X)\|A-A'\|_X .
 \end{aligned}
 \label{mb:section-input}
\end{equation}
All maps are the actual constraint and its actual derivative.
Writing \(A=Eu\) reduces \(q(A)=b\) to
\begin{equation}
 u=b-N_q(Eu).
 \label{mb:section-contraction}
\end{equation}
On \(\|u\|_Y\le2\|b\|_Y\) the right side is a self-map and
has Lipschitz constant at most \(4C_qC_E^2\|b\|_Y\).
If that number is less than one half, successive substitution
converges geometrically, uniformly in the finite lattice.
The resulting correction is at most \(2C_E\|b\|_Y\) in \(X\).
For complex data, uniform convergence on a smaller open ball
makes it holomorphic: each iterate is holomorphic, and applying
continuous linear functionals and Cauchy's integral formula
proves holomorphy of the limit.

This proves the constraint correction from \eqref{mb:section-input}. It does not
prove the scale-uniform right-inverse estimate, the smoothed
trial estimate, or the overlap cocycle just listed.
Those are bounded geometric and operator tasks for the actual
map. The target is that their constants are independent of
depth, volume, region shape and real holonomy, for the fixed
allowed local geometry. A proof for only one block realizes
the corresponding first-step instance.

\subsubsection{The actual critical equations}

Use fixed reference coordinates
\(U(A)=e^{i\eta A}U_0\), with the prescribed slice imposed.
Write the constraint in the coordinates just specified as
\(q(A)=b(B)\), and set \(F(A)=A_{\Omega_0}(U(A))\).
Here \(q\) is the chosen actual plain constraint, or the
straightened constraint with the accompanying node factors and
reduced datum identified below. These choices have different
derivatives; every use of a linear inverse specifies which is used.
All transposes below are bilinear cell-weighted transposes.
For real positivity the pairing is Hermitian instead.
In particular a change from raw bond logarithms to physical
potentials also changes the dual current weights.

Let \(J_A(U)\) be the physical left-variation current:
\begin{equation}
 \left.\frac d{dt}A_{\Omega_0}(e^{i\eta th}U)\right|_{t=0}
       =\eta^4\sum_b (J_A(U)(b),h(b)).
 \label{mb:physical-current-definition}
\end{equation}
For the reference-coordinate derivative this gives
\begin{equation}
 DF(A)[h]
   =\eta^4\sum_b
       (J_A(U(A))(b),\mathcal E_{\eta A(b)}h(b)).
 \label{mb:reference-gradient}
\end{equation}
Indeed
\((d/dt)e^{i\eta(A+th)}|_0
       =i\eta\mathcal E_{\eta A}(h)e^{i\eta A}\).
Thus the reference covector at nonzero \(A\) includes a
right-logarithmic differential. Dropping this factor from a
second derivative changes its current terms.

The coupled critical equations are
\begin{equation}
 DF(A)-Dq(A)^\sharp\omega=0,\qquad q(A)=b(B).
 \label{mb:stationarity}
\end{equation}
Here \(\omega\) is a coarse covector in the normalization
dual to \(q\). For raw constraints \(q=mQ\), one has
\(\omega_Q=m^\sharp\omega\). The scalar
\(\langle\omega,D^2q\rangle
 =\langle\omega_Q,D^2Q\rangle\) is unchanged.

These equations are equivalent to constrained criticality
when \(Dq\) is onto on the chosen slice. One direction follows
by testing \eqref{mb:stationarity} on \(\ker Dq\). Conversely choose any
right inverse \(E_A\) and suppose \(DF(A)\) annihilates
that kernel. Since
\(h-E_ADq\,h\in\ker Dq\),
\begin{equation}
 DF(A)[h]=DF(A)[E_ADq\,h]
   =\langle E_A^\sharp DF(A),Dq\,h\rangle.
 \label{mb:multiplier-existence}
\end{equation}
This proves the assertion and the uniqueness of the multiplier:
if \(Dq^\sharp\omega=0\), testing on \(E_Av\) gives
\(\langle\omega,v\rangle=0\) for every \(v\).
Surjectivity and its quantitative bound are therefore genuine
parts of the background theorem.

The Lagrangian Hessian is
\begin{equation}
 K(A,\omega)=D^2F(A)-\langle\omega,D^2q(A)\rangle .
 \label{mb:lagrangian}
\end{equation}
It is its restriction to tangent fields, not \(D^2F\) alone,
that controls the critical fibre. This follows either by
differentiating the multiplier equation or directly from a
curve \(q(A(t))=b\):
\(Dq\,\ddot A=-D^2q[\dot A,\dot A]\) and hence
\begin{equation}
 \frac{d^2}{dt^2}F(A(t))\big|_0
       =\langle\dot A,K(A,\omega)\dot A\rangle
 \quad\text{at a critical point}.
 \label{mb:curved-fibre-hessian}
\end{equation}

At a noncritical point the same fibre coordinate has an
additional term involving its acceleration and \(DF\).
For the exact pivot chart \(\Phi\), differentiating
\(q(\Phi(u))=q(0)\) gives
\(D^2\Phi=-h\,D^2q[D\Phi,D\Phi]\).
Consequently the Hessian of \(F\circ\Phi\) is the projected
action Hessian minus the contraction with
\(h^\sharp DF\), with its correct physical current and chart
differentials.  This identity follows from the finite
chart differentiation.  The first-step positivity actually
proved by the finite precision results is for
\(D^2[(F+{\rm GF}_{\rm actual})\circ\Phi]\) at an
admitted real block-axial trial background, under the
stated first-step radius and current ceilings.  The actual
soft gauge term and its identification are retained;
at the physical current this is the corresponding
action-plus-gauge fibre Hessian.  A positive action-only
Hessian on a separately chosen slice would require its
own slice estimate and is not imported from that result.
Criticality is needed to identify the multiplier
interpretation in \eqref{mb:lagrangian}, not for the stated
action-plus-gauge trial-field estimate.

\paragraph{Absolute and relative soft gauges}

The internal transport is a mean of coordinate-order paths.
The absolute gauge term is
\begin{equation}
 {\rm GF}_{\rm abs}(U)=
   \sum_{y,x}\left(1-\tfrac12\operatorname{tr}T_{y,x}(U)\right).
 \label{mb:absolute-gauge}
\end{equation}
Its full Hessian is the one in a physical first-step
calculation on an axial field. If a construction instead uses
the relative gauge with reference \(U\),
\begin{equation}
 {\rm GF}_{{\rm rel},U}(a)=
   \sum_{y,x}\left(1-\tfrac12\operatorname{tr}
          [T_{y,x}(e^{ia}U)T_{y,x}(U)^{-1}]\right),
 \label{mb:relative-gauge}
\end{equation}
its gradient at \(a=0\) vanishes for every background. If
\(G_Ua\) is the right-logarithmic first derivative of its
transport, its quadratic form is the corresponding positive
sum of squares, \(G_U^\sharp G_U\), with the actual trace
normalization. The second transport derivative has zero
contribution because it is paired with the derivative of
the scalar function at the identity.

The absolute Hessian has the same form only where
\(T_{y,x}(U)=1\). For example along
\(T(t)=e^{i(\theta+t)\sigma_3}\) the absolute second
derivative of \(1-\cos(\theta+t)\) is \(\cos\theta\),
whereas the relative second derivative is one.
They differ at arbitrarily small nonzero \(\theta\).
Both absolute and relative terms are actual finite formulas;
one cannot select between them by positivity alone.

The relative expression is covariant under a joint site
gauge of its reference and variable field: both transports
transform bi-equivariantly, and their product transforms by
conjugation at the center. The absolute expression has the
corresponding invariance only when the endpoint and center
gauge values agree, or when the prescribed gauge convention
otherwise makes its transport a closed quantity.
At a seam, removing bonds from a determining set can make
the relative convention the appropriate one. A retained
boundary datum may instead require the absolute convention.
The local row definition must decide which formula is
present.

Thus an independent-current precision theorem must compare
its entire projected form with the physical or prescribed
auxiliary one. Write \(\mathcal J=DF\) for the physical comparison,
and let \(\mathcal J\) be the declared independent fine covector,
after the same dual-coordinate conversion, for an auxiliary comparison:

\begin{equation}
 C^\sharp\big[
     \Delta_{\rm named}(U,J)
       -D^2F(U)+\langle h^\sharp\mathcal J,D^2q\rangle
       -K_{\rm gauge,actual}(U,J)\big]C=0,
 \label{mb:complete-normal-form}
\end{equation}
after every term has been expressed in the same raw or
physical coordinates. Equivalently the difference may be
estimated with an explicit error and charged to the gap.
The displayed subtraction is schematic only in its
\emph{name}: each named term must have its literal local
definition, including every projected extra-current term.
The comparison is an equality of fully typed bilinear
forms. Constraint-curvature matching alone checks only
one summand. Zero-fold averaging maps do not imply that a
separate projected correction \(\Delta'_\pi[J]\) vanishes.
The precision section retains this complete matching task.

\subsubsection{The action induced on the coarse datum}

Suppose the critical solution of \eqref{mb:stationarity} has been constructed
on an open real datum ball and admits differentiation there.
The induced action is
\(S(B)=F(A(B))\). Let \(Q=Dq(A)\), and let \(T\) denote
the energetic tangent lift: \(QT=1\), while
\(KTv\) annihilates \(\ker Q\). The response equations give
\begin{equation}
 \dot A=T\,Db\,\dot B,\qquad
 \dot\omega=\mathcal M\,Db\,\dot B,\qquad
 KT=Q^\sharp\mathcal M .
 \label{mb:energetic-lift}
\end{equation}
The letter \(\mathcal M\) here denotes the multiplier response,
not the group averaging map. It is determined by these
equations and is symmetric in the real bilinear pairing.
Indeed
\(\langle v,\mathcal M w\rangle
 =\langle Tv,KTw\rangle\) by \(QT=1\).

Differentiating \(S\) and using \eqref{mb:stationarity} proves
\begin{equation}
 \begin{split}
 DS(B)[v]&=\langle\omega,Db\,v\rangle,\\
 D^2S(B)[v,w]&=
    \langle Db\,v,\mathcal M\,Db\,w\rangle
               +\langle\omega,D^2b[v,w]\rangle .
 \end{split}
 \label{mb:coarse-action-hessian}
\end{equation}
In particular an outward seam has the second term even
when its interior datum coordinates are linear. An
argument using only \(\mathcal M\) would lose this actual
datum-curvature term. There is no such term in a chart
where the datum itself is the constraint coordinate;
converting that chart back to \(B\) restores it.

There is also an exact formula from an admitted soft
completion. Let \(a\) be the coarse gauge or penalty
matrix in that completion and assume
\begin{equation}
 G=(K+Q^\sharp aQ)^{-1},\qquad
 W=QGQ^\sharp
 \quad\text{are defined, with }W\text{ invertible}.
 \label{mb:soft-completion}
\end{equation}
Set \(H=GQ^\sharp W^{-1}\). Then \(QH=1\), and
\begin{equation}
 KH=Q^\sharp(W^{-1}-a),\qquad
 H=T,\qquad \mathcal M=W^{-1}-a.
 \label{mb:soft-readout}
\end{equation}
The first identity follows by multiplying
\((K+Q^\sharp aQ)G=1\) by \(Q^\sharp W^{-1}\).
It proves the energetic property of \(H\). Uniqueness
follows from invertibility of the constrained tangent
operator: the difference of two energetic lifts lies
in \(\ker Q\) and is orthogonal under \(K\) to that
kernel. This proves \eqref{mb:soft-readout}.

These are algebraic identities on the admitted domain.
They supply no bound for \(W^{-1}\), no positivity for
a completion not actually present, and no identification
of an independent-current operator. Once those primitive
facts have been proved, however, the coarse readout and
its correct subtraction are forced by \eqref{mb:coarse-action-hessian}--\eqref{mb:soft-readout}.
On the next scale the action is this \(S\), that is the
original fine Wilson action evaluated at its constrained
critical background. It must not be replaced by the
bare Wilson formula on the coarse variables.

\subsubsection{Straightening and the named current completion}

Plain, straight and straightened averages must remain distinct.
Write \(q_n^{\rm st}(A)\) for the nonlinear straightened raw
constraint at level \(n\). Its defining node factors satisfy
\begin{equation}
 \mathcal M^n(e^{i\eta A}U_0)
      =\big(e^{iq_n^{\rm st}(A)}\mathcal M^nU_0\big)^{v_n[A]} .
 \label{mb:node-diagram}
\end{equation}
For each top of level \(n\), prescribe
\(u[A](y)v_n[A](y)=1\). Interior rows then have physical datum
\(e^{iq_n^{\rm st}(A)}V_0\). At an outward endpoint of level
\(n-1\) the value prescribed for \(u[A]\) uses \(v_{n-1}\),
not \(v_n\). With the orientation of \eqref{mb:reduced-datum}, covariance gives
\begin{equation}
 e^{iq_n^{\rm st}(A)(c)}
   =e^{iB(c)}
       \operatorname{Ad}(V_0(c))
          [v_{n-1}[A](z)^{-1}v_n[A](z)] .
 \label{mb:outward-node-identity}
\end{equation}
On the fully prescribed outer block the bracket is exactly
the one-step factor \(v_z(B_{\rm int})\) in \eqref{mb:reduced-datum}.
This follows by evaluating its defining one-step averaging
identity on the already prescribed level-\((n-1)\) interior
data. Therefore the straightened equation is
\(q^{\rm st}(A)=b(B)\) with the reduced datum, including its
outward factors. \eqref{mb:outward-node-identity} also checks orientation
and multiplication order.

After an admitted Landau slice has been constructed, the
physical representative is
\begin{equation}
 \Phi(B)=(e^{i\eta\mathcal H(B)}U_0)^{u_B},\qquad
 RD^\sharp\mathcal H(B)=0,\qquad
 u_Bv_n[\mathcal H(B)]=1\text{ at level-}n\text{ tops}.
 \label{mb:landau-pair}
\end{equation}
An independent comb slice gives a different potential and
a different gauge factor. It has the same physical orbit
and the same reduced datum, but need not coincide with
the pair in \eqref{mb:landau-pair}. Equality requires that its potential
already obey the Landau condition. A bound on \(u_B-1\)
only at the tops cannot be extended to all bare sites
without the separate comb gauge estimate.

Gauge invariance of the fine action makes \(F(\mathcal H(B))\)
the action of \(\Phi(B)\). Constrained criticality is
therefore described by \eqref{mb:stationarity} on this slice, provided that
the slice and node construction parametrize the local
physical fibre modulo the residual group. The primitive
gauge estimate must prove that parametrization, its
uniqueness and its admitted ball. The node identities
above prove the datum correspondence; they alone do not
give the quantitative gauge estimate.

For reference the actual named propagator family uses the
linear part \(Q\) per unit row length of this straightened
constraint. Let \(D=D_U\) be the covariant site difference,
let \(Q'\) be the corresponding site average to the tops,
and let \(a\) have the actual row or site weight. On its
admitted domain use the zero-extension site Laplacian
\(\Delta_{\rm s}\), and set
\begin{equation}
 \begin{split}
 G'&=(\Delta_{\rm s}+Q'^\sharp aQ')^{-1},\\
 R&=1-G'Q'^\sharp
          (Q'G'^2Q'^\sharp)^{-1}Q'G',\\
 \mathsf S&=1-DG'RD^\sharp,\qquad
 \mathsf G_0^{-1}
       =\Delta_U^{\rm W}+DRD^\sharp+Q^\sharp aQ .
 \end{split}
 \label{mb:named-gauge-operators}
\end{equation}
At a real field \(R\) is the declared orthogonal site
projector; on complex fields the formula uses bilinear
transposes and analytic continuation on its admitted
domain. Its inverses and its relation to the chosen
slice are estimates to be established, not consequences
of displaying the formula.

Let \(\Delta'_\pi[J]\) be the specified linear-current
correction and let \(\Delta^{(2)}[\mathsf j]\) be the
constraint quadratic form, normalized by
\(\langle A,\Delta^{(2)}[\mathsf j]A\rangle
       =2\langle C^{(2)}(A),\mathsf j\rangle\).
The subsequent definitions are ordered:
\begin{equation}
 \begin{aligned}
 \mathsf G^{-1}&=\mathsf G_0^{-1}-\Delta'_\pi[J],\\
 \mathsf H&=\mathsf GQ^\sharp(Q\mathsf GQ^\sharp)^{-1},\\
 \mathsf G_1^{-1}
   &=\mathsf G^{-1}
       -\mathsf S^\sharp\Delta^{(2)}[\mathsf H^\sharp J]\mathsf S,\\
 \mathsf H_1&=\mathsf G_1Q^\sharp(Q\mathsf G_1Q^\sharp)^{-1}.
 \end{aligned}
 \label{mb:named-current-operators}
\end{equation}
In particular the \(\mathsf H\) inside the definition of
\(\mathsf G_1\) is formed first. Replacing it by
\(\mathsf H_1\) there gives a different implicit operator.
Likewise \(Q\) in these formulas is not the plain derivative
of \(\mathcal M^n\), and the gauge term is not replaced by
the contour soft gauge of the next step.

The next-step coarse form has the readout
\begin{equation}
 \begin{aligned}
 \Delta_k&=(Q\mathsf G_1Q^\sharp)^{-1}-a
                 -\mathsf E^{(2)}[\mathsf H_1^\sharp J],\\
 \mathsf A&=\mathsf C_V^\sharp
          (\mathsf G_V^{(2)}+\Delta_k)\mathsf C_V,
 \qquad V=\mathcal M^kU .
 \end{aligned}
 \label{mb:next-step-form}
\end{equation}
Here \(\mathsf C_V\) eliminates the actual plain-average
pivots, \(\mathsf E^{(2)}\) is that one-step plain constraint's
quadratic form, and \(\mathsf G_V^{(2)}\) is the prescribed
soft contour gauge there. The last current contraction
is the datum-curvature subtraction obtained when the
coarse datum is itself restricted to the next-step fibre.
It is not the earlier straightened \(\Delta^{(2)}\).

On a physical critical pair, identifying
\(\mathsf G_1^{-1}\) with the actual Lagrangian Hessian
plus its actual gauge and row terms makes \eqref{mb:coarse-action-hessian}--\eqref{mb:soft-readout}
the derivation of this coarse readout. If an alternative
expression subtracts an additional
\(\mathsf H_1^\sharp DRD^\sharp\mathsf H_1\), it agrees
only after that term has been shown to vanish, for example
by the proved physical slice relation \(RD^\sharp\mathsf H_1=0\).
There is no corresponding automatic identity for an
independent \(J\).

Thus the precision problem contains two different tasks:
prove bounds for the literal ordered family in \eqref{mb:named-gauge-operators}--\eqref{mb:next-step-form}
on its primitive independent-current domain, and identify
its physical readout with the Hessian of the induced action.
The latter uses the full projected-current and gauge comparison,
not merely the constraint quadratic term. A proof of the
first-step Wilson curved-fibre gap does not by itself
perform either multilevel task.

\subsubsection{A quantitative coupled construction}

Here the primitive inverse and coefficient estimates are
stated with their actual objects. This makes the existence
argument a finite, checkable nonlinear reduction rather
than an invocation of an unspecified implicit-function
theorem.

Fix a real trial field \(U_0\), not necessarily critical,
in \eqref{mb:primitive-real-domain}, with admitted constraint and slice charts.
Let \(Q_0=Dq(0)\). Choose a coarse trial covector
\(\omega_0\) and put
\begin{equation}
 \begin{aligned}
 f_0&=-DF(0)+Q_0^\sharp\omega_0,\qquad
 K_0=D^2F(0)-\langle\omega_0,D^2q(0)\rangle,\\
 \mathcal S_0&=
   \begin{pmatrix}K_0&-Q_0^\sharp\\Q_0&0\end{pmatrix}.
 \end{aligned}
 \label{mb:trial-system}
\end{equation}
The field is restricted to the fixed slice, and the
force covectors are restricted to its dual. If the
gauge kernel has not been quotiented, its gauge equation
and its actual site inverse have to be included in
this system. There is no inverse on an unremoved gauge
kernel.

Let \(X,Z,Y,X^*\) be the declared field, multiplier,
datum and force spaces with their physical weights.
The primitive inverse estimate to be proved is a
two-sided kernel estimate for
\(\mathcal S_0^{-1}:X^*\oplus Y\to X\oplus Z\).
Its consequence used in this section is
\begin{equation}
 \|\mathcal S_0^{-1}(f,b)\|_{X\oplus Z}
                    \le C_0(\|f\|_{X^*}+\|b\|_Y),
 \label{mb:primitive-inverse}
\end{equation}
with constants uniform over volume, depth, admissible
regions, local scales and real holonomies. The precision
section separates its gauge, tangent, row and seam
estimates. It must also prove that the complete matrix
in \eqref{mb:trial-system}, rather than a convenient surrogate, has this
bound on the primitive trial class.

Write \(A\) for the field and \(\nu=\omega-\omega_0\).
Taylor's formula gives the exact equations
\begin{equation}
 \begin{split}
 \mathcal S_0(A,\nu)
  =\big(&f_0-R_F(A)+R_{q,\omega_0}(A)
               +(Dq(A)-Q_0)^\sharp\nu,\quad
                    b-N_q(A)\big),\\
 R_F(A)&=DF(A)-DF(0)-D^2F(0)A,\\
 R_{q,\omega_0}(A)
    &=[Dq(A)-Q_0-D^2q(0)A]^\sharp\omega_0,\\
 N_q(A)&=q(A)-q(0)-Q_0A .
 \end{split}
 \label{mb:exact-fixed-point}
\end{equation}
The right datum \(b\) here is measured from \(q(0)\).
The term linear in \(A\) and in \(\omega_0\) was included
in \(K_0\). Retaining it again on the right would
double count the constraint curvature.

For clarity, the finite coefficient task is to verify,
on the entire declared complex ball, the inequalities
\begin{equation}
 \begin{split}
 \|R_F(A)-R_F(A')\|_{X^*}
       &\le c_F(r)\|A-A'\|_X,\\
 \|R_{q,\omega_0}(A)-R_{q,\omega_0}(A')\|_{X^*}
       &\le c_{q\omega}(r)\|A-A'\|_X,\\
 \|N_q(A)-N_q(A')\|_Y
       &\le c_q(r)\|A-A'\|_X,\\
 \|(Dq(A)-Q_0)^\sharp\nu
       -(Dq(A')-Q_0)^\sharp\nu'\|_{X^*}
       &\le c_{\rm mix}(r)
               (\|A-A'\|_X+\|\nu-\nu'\|_Z).
 \end{split}
 \label{mb:coefficient-list}
\end{equation}
Here \(\|A\|_X,\|A'\|_X,\|\nu\|_Z,\|\nu'\|_Z\le r\).
At a fixed finite lattice, Taylor integrals prove the
inequalities with finite constants. For example
\begin{equation}
 R_F(A)=\int_0^1(1-t)\,D^3F(tA)[A,A]\,dt .
 \label{mb:taylor-force}
\end{equation}
The notation contracts two slots, leaving a covector.
The analogous integral for \(Dq-Q_0\) has one slot
contracted, and that for \(R_{q,\omega_0}\) has two.
A uniform estimate must count the actual local
coordinates in these tensors and use the force and
field weights before summing. A scalar supremum bound
times the total number of fine variables is not
uniform in this problem.

Set \(c(r)=c_F(r)+c_{q\omega}(r)+c_q(r)+c_{\rm mix}(r)\).
Assume the verified inequalities give
\begin{equation}
 C_0c(r)\le\tfrac12,\qquad
 C_0(\|f_0\|_{X^*}+\|b\|_Y)\le r/2.
 \label{mb:contraction-conditions}
\end{equation}
Applying \(\mathcal S_0^{-1}\) to \eqref{mb:exact-fixed-point} defines a
contraction of the closed radius-\(r\) ball into itself.
Starting at zero, its successive iterates satisfy
\(\|x_{n+1}-x_n\|\le2^{-n}\|x_1\|\); summing this
geometric series proves convergence in \(X\oplus Z\).
Continuity of the map proves the equation. The same
estimate on the difference of two fixed points proves
uniqueness in that ball.

For two forcings whose actual maps share this ball,
subtracting their fixed-point equations gives
\begin{equation}
 \|A-\widetilde A\|_X+\|\nu-\widetilde\nu\|_Z
 \le2C_0\big(\|f_0-\widetilde f_0\|_{X^*}
                         +\|b-\widetilde b\|_Y\big)
 \label{mb:fixed-point-response}
\end{equation}
when their nonlinear coefficient maps agree.
If those maps also change, add their evaluated
difference on the common ball to the right-hand side.
Its smallness must be proved from the actual coefficient
comparison, not from a difference of the desired solutions.

The argument permits a noncritical trial field through
the explicit \(f_0\). It does not establish that an
arbitrary regular critical field lies in this ball.
That additional uniqueness statement needs a
preliminary slice-closeness and local regularity
argument. The raw group closeness \eqref{mb:rigidity-sum} is one
ingredient, but does not by itself give \(X\)-closeness.

\paragraph{Whole-ball holomorphy and reality}

Assume the finite products, means, logs, slice readouts
and operator inverse are admitted on the whole complex
ball used in \eqref{mb:exact-fixed-point}, with a fixed margin. If the datum
parameter varies over an open disc for which \eqref{mb:contraction-conditions}
holds uniformly, all iterates of the fixed-point map
are holomorphic and converge uniformly on smaller
closed discs. Cauchy's formula on each finite collection
of parameter directions then passes to the limit.
The solution is holomorphic and has the same bound \(r\).
This proof does not require a volume-dependent radius.

For real reference data, conjugating the equations by
\(A\mapsto A^*\), \(\omega\mapsto\omega^*\) and
\(g\mapsto(g^*)^{-1}\) preserves their finite formulas.
The conjugate solution solves the equation at the
conjugate datum and remains in the same ball.
Uniqueness proves
\begin{equation}
 A(B^{\mathfrak c})=A(B)^{\mathfrak c},\qquad
 \omega(B^{\mathfrak c})=\omega(B)^{\mathfrak c},\qquad
 \Phi(B^{\mathfrak c})=\Phi(B)^{\mathfrak c}.
 \label{mb:reality}
\end{equation}
At real \(B\) this is the required real solution.
Gauge covariance is proved in the same way, provided
the gauge maps both primitive domains and their slice
normalizations to one another. The uniqueness theorem
then identifies the transformed solution. For a
complex gauge this proviso includes its condition
number and every branch inequality; algebraic
covariance alone does not make the complex core
gauge saturated.

\subsubsection{Local gauges can be glued without a large holonomy assumption}

We next prove the gluing calculation used by a local regularity
argument. Its hypotheses are local potential and first-difference
estimates on patches of fixed size. Its conclusion does not
include a componentwise Laplacian estimate; that estimate has
a separate elliptic role.

Let a rectangular cube have side \(m\xi\), with
\(\eta\le\xi\), and enlarge it by a fixed collar.
Cover it by overlapping patches \(W_a\) of side \(C\rho\xi\),
whose centers lie on a rectangular grid of spacing
comparable to \(\rho\xi\). Here \(\rho\) is fixed.
Assume bounded overlap and a partition of unity with
\begin{equation}
 \sum_a\chi_a=1,\qquad \chi_a\ge0,\qquad
 |\nabla^r\chi_a|\le C_r\xi^{-r}\quad(r=1,2).
 \label{mb:gauge-partition}
\end{equation}
Such a partition is obtained by sampling fixed smooth
tensor bumps and dividing by their positive sum.
The finite-difference product formula, or the mean
value theorem applied to the smooth bumps, proves
the bounds uniformly in \(\eta\).

Suppose the real field has
\(|\partial U-1|\le\epsilon\eta^2\xi^{-2}\)
and on each \(W_a\) there is a real unitary gauge \(k_a\)
in which
\begin{equation}
 U^{k_a}=e^{i\eta A_a},\qquad
 \xi|A_a|+\xi^2|\nabla A_a|\le C\epsilon .
 \label{mb:local-gauge-input}
\end{equation}
On the real cubes used below,
\hyperlink{mb:real-cube-gauge}{the elliptic section construction}
supplies \eqref{mb:local-gauge-input} from curvature and full-current
bounds. The actual current transforms by conjugation and can be
carried along in the patch comparison.

We first align the constant frames. Choose a corner
of the large rectangular lift as root. For each center
\(y_a\), choose the coordinate-order path from the
root to \(y_a\), and multiply \(k_a\) by the constant
\begin{equation}
 c_a=U(\Gamma_{{\rm root},y_a})\,k_a(y_a)^{-1},
 \qquad \widehat k_a=c_ak_a.
 \label{mb:constant-alignment}
\end{equation}
This does not change either norm in \eqref{mb:local-gauge-input}, since it
conjugates \(A_a\) by a constant unitary matrix.
If \(x\in W_a\), the transport of the locally gauged
field from \(y_a\) to \(x\) along a path of length
\(O(\rho\xi)\) differs from one by \(C\epsilon\).
The unitary product bound proves this by summing
\(\eta|A_a|\) on that path. Therefore
\(\widehat k_a(x)\) differs by \(C\epsilon\) from the
original transport along the root path followed by
that local path.

For overlapping \(W_a,W_b\), those two rooted paths
differ by a strip, not an arbitrary surface spanning
the whole large cube. To verify this, change their
center coordinates one at a time. The difference
between the centers is \(O(\rho\xi)\) in each
coordinate. Commuting this short displacement through
the four ordered root segments sweeps at most
\(C m\rho\xi^2\) area. The extra local paths sweep
only \(C\rho^2\xi^2\). The lattice interchange proof
of \eqref{mb:finite-stokes}, now counting fine plaquettes in these
strips, gives
\begin{equation}
 |\widehat k_a(x)\widehat k_b(x)^{-1}-1|
                         \le C_\rho m\epsilon
                           \quad(x\in W_a\cap W_b).
 \label{mb:transition-value}
\end{equation}
The estimate is valid on the rectangular lift that
contains those paths. A periodic cube wider than
half the torus is excluded from this local assertion.
No global Wilson line was replaced by one; the
possibly large root transports cancel in a comparison
loop and conjugate its plaquettes unitarily.

The transition also has controlled derivatives.
Put \(r_{ab}=\widehat k_a\widehat k_b^{-1}\).
For an oriented fine bond, covariance of \(U\) says
\begin{equation}
 r_{ab}(x+\eta e_\nu)
       =e^{-i\eta\widehat A_{a,\nu}(x)}
             r_{ab}(x)e^{i\eta\widehat A_{b,\nu}(x)} .
 \label{mb:transition-equation}
\end{equation}
Subtract \(r_{ab}(x)\) and divide by \(\eta\).
Unitarity and \eqref{mb:local-gauge-input} give
\(|\nabla r_{ab}|\le C\epsilon/\xi\).
Taking another difference of this exact formula uses
only \(\nabla A_a,\nabla A_b,\nabla r_{ab}\).
The derivative bound for the exponential on a fixed
small raw ball gives
\begin{equation}
 |\nabla r_{ab}|\le C\epsilon\xi^{-1},\qquad
 |\nabla^2r_{ab}|\le C(\epsilon+\epsilon^2)\xi^{-2}.
 \label{mb:transition-jets}
\end{equation}
This proves the two derivative bounds without
differentiating the long root paths.

Form \(X(x)=\sum_a\chi_a(x)\widehat k_a(x)\).
A real linear combination of \({\rm SU}(2)\)
matrices is a quaternionic matrix. Its determinant
is the squared Euclidean norm of its four real
quaternion coordinates, so, when it is nonzero,
\begin{equation}
 \mathcal N(X)=X/\sqrt{\det X}\in{\rm SU}(2).
 \label{mb:quaternion-normalization}
\end{equation}
At a point select any active anchor \(a_0\), and
write \(X=Y\widehat k_{a_0}\), where
\(Y=\sum_a\chi_a r_{a a_0}\).
\eqref{mb:transition-value} gives
\(\|Y-1\|\le C_\rho m\epsilon\).
If this is less than a sufficiently small fixed
number, then \(\det Y\) stays in a neighborhood
of one not meeting zero. Thus \(\mathcal N(X)\)
is defined, independently of the selected anchor,
and
\begin{equation}
 k=\mathcal N(X)=\mathcal N(Y)\widehat k_{a_0},
 \qquad
 |\mathcal N(Y)-1|+
     \xi|\nabla\mathcal N(Y)|+
     \xi^2|\nabla^2\mathcal N(Y)|
                        \le C_\rho m\epsilon .
 \label{mb:glued-relative-jets}
\end{equation}
The derivative estimates follow from \eqref{mb:gauge-partition},
\eqref{mb:transition-value}--\eqref{mb:transition-jets} and the uniformly bounded first two
derivatives of \(Y\mapsto Y/\sqrt{\det Y}\) near
one. When taking differences one keeps the anchor
on a neighborhood containing those differences;
the bounded enlarged overlap makes this possible.
The resulting \(k\) itself does not depend on
that choice.

In the anchor frame, \(U^k\) is
\begin{equation}
 \mathcal N(Y)(x)\,e^{i\eta\widehat A_{a_0}(b)}
                            \mathcal N(Y)(x+\eta e_\nu)^{-1}.
 \label{mb:glued-field}
\end{equation}
The relative endpoint product is \(1+O(\eta m\epsilon/\xi)\).
The admitted small logarithm therefore defines \(A\).
Taking one difference, using \eqref{mb:glued-relative-jets}, gives
\begin{equation}
 U^k=e^{i\eta A},\qquad
       \xi|A|+\xi^2|\nabla A|\le C_\rho m\epsilon .
 \label{mb:glued-potential}
\end{equation}
These bounds can be verified directly by differentiating
the logarithm on its small fixed raw ball: the endpoint
increment divided by \(\eta\) is \(O(m\epsilon/\xi)\);
its first difference is \(O(m\epsilon/\xi^2)\).
All remainders have an extra factor
\(\eta m\epsilon/\xi\le m\epsilon\), which is small.
If \(\xi^3|J_A(U)|\le\epsilon\), the same bound holds
after \(k\), since the current is conjugated at its
fine endpoint. In particular no factor \(m\) is
needed for that covariant current bound.

The normalization has an admitted holomorphic version.
If the local gauges depend holomorphically on
parameters and their anchored \(Y\)'s remain in the
same small complex neighborhood of one, the branch
of \(\sqrt{\det Y}\) is the principal branch there.
The formula is holomorphic and the relative
normalizations agree on overlaps by the exact
identity \(\det\widehat k_a=1\).
For real gauges it is the positive square root.
No absolute value of a complex determinant occurs.

The overlap argument gives the potential, first difference
and covariant current bounds. For the componentwise Laplacian
of this particular patched potential one would need a further
estimate. The real cubes in the layered profile instead use
the single gauge constructed below. Given the componentwise
Laplacian bound, the discrete interior estimates give
\begin{equation}
 [\nabla A]_{\alpha;\square}
   \le C_{\alpha,\rho}m\epsilon\,\xi^{-2-\alpha},
                  \qquad 0<\alpha<1 .
 \label{mb:holder-reduction}
\end{equation}
Here the input is the bounded potential and
componentwise Laplacian on the larger cube, with
their displayed physical scales. The constant
may diverge as \(\alpha\uparrow1\). Thus this
Hölder conclusion follows from those two bounds.

%
%
\paragraph{A single real gauge with second-order bounds.}
\hypertarget{mb:real-cube-gauge}{}
Work first in units of the cube scale, with bare spacing
\(\eta\le1/13\). On a lifted box of side \(s\ge5\), assume
\(|\partial U-1|\le\alpha\eta^2\) on its plaquettes and
\(|D_U^{\mathsf T}\partial U|\le\alpha\eta^2\) at its full bonds.
Put \(e=s\alpha\) and take it below a fixed small constant.
The comb section \(\phi_0\) is unitary. Comparing the two
comb paths at the endpoints of a bond inserts at most
\(3s/\eta\) plaquettes, so \(\|D_0\phi_0\|_\infty\le3e\).
Here \(D_0\) is the covariant difference of a section, formed
only on bonds of the box. Smooth \(\phi_0\) by solving
\[
 (D_0^*D_0+1)\phi=\phi_0,
 \qquad \psi=D_0\phi .
\]
The operator is positive and invertible in the real quaternionic
space. At a site of maximum section norm the inner product of
\(D_0^*D_0\phi\) with \(\phi\) is nonnegative, by unitarity of
each transport. It follows that \(\|\phi\|_\infty\le1\).

Differentiate the equation. Moving a transported difference
past a perpendicular one inserts the corresponding plaquette
holonomy. In the equation for \(\psi_\mu\) the difference of
the opposite plaquette contributions gives the covariant
plaquette divergence. The remaining commutators have one
curvature factor and one first difference of \(\phi\).
At a missing perpendicular neighbor there is also an unpaired
plaquette contribution of size \(5\alpha/\eta\). More precisely,
with \(\Sigma=\|\psi\|_\infty\), the modulus is a subsolution of
\[
 (\mathsf L_{0,\mu}+1)|\psi_\mu|(x)
 \le \|D_0\phi_0\|_\infty+2\alpha+12\alpha\Sigma
                       +5\alpha\eta^{-1} n_\mu(x),
\]
where \(n_\mu\) counts the missing perpendicular neighbors.
The scalar bond Laplacian \(\mathsf L_{0,\mu}\) uses only the
available perpendicular neighbors and zero extension along
direction \(\mu\). This inequality follows by taking the
inner product with the unit direction of \(\psi_\mu\), using
unitarity on the transported terms. At a zero the scalar
left-hand side is nonpositive. To pay the boundary term,
for each perpendicular direction take the one-dimensional
function on \(N\) sites
\[
 \beta(l)=\frac{\cosh(\kappa\eta(l-(N-1)/2))}
                    {\sinh(\kappa\eta N/2)},\qquad
 2\sinh(\kappa\eta/2)=\eta .
\]
It is bounded by two, solves the massive homogeneous
difference equation in the interior and gives \(\eta^{-1}\)
at each endpoint with the absent neighbor omitted. The sum
of the perpendicular barriers is a supersolution for
\(\eta^{-1}n_\mu\), because the zero extension in direction
\(\mu\) adds a nonnegative term. Scalar comparison yields
\(\Sigma\le\|D_0\phi_0\|_\infty+32\alpha+12\alpha\Sigma\).
In particular \(\Sigma\le19e\) for the admitted smallness.

In the comb frame, \(\phi_0^{-1}(\phi-\phi_0)\) is the flat
divergence of \(\phi_0^{-1}\psi\), up to an error of size
\(4\|D_0\phi_0\|_\infty\Sigma\). Average on a subbox of unit
side up to one mesh spacing. The divergence telescopes in
each coordinate, so its mean is at most \(8\Sigma\). Across
one bond the section difference changes by at most
\(\eta(\Sigma+\|D_0\phi_0\|_\infty)\). Comparing the mean
with the value at an included site gives
\[
 \|\phi-\phi_0\|_\infty\le
 12\Sigma+4\|D_0\phi_0\|_\infty+4\|D_0\phi_0\|_\infty\Sigma .
\]
Thus \(\rho_\phi=|\phi|\ge1/2\) and \(u=\phi/\rho_\phi\)
is a real gauge. On the bond \(b=(x,x+\eta e_\mu)\) put
\(\omega_\mu(x)=\rho_\phi(x)^{-1}u(x)^{-1}\psi_\mu(x)\).
The transformed link is
\((1+\eta\omega_\mu)/|1+\eta\omega_\mu|\). In the interior
the product rule in the differentiated section equation gives
\[
 |\Delta_\eta\omega_\mu(x)|\le
 102e+304e\max_{x'\sim x}\eta^{-1}
                                |\omega_\mu(x')-\omega_\mu(x)| .
\]
Indeed the transported Laplacian of \(\psi\) is bounded by
\(\Sigma+\|D_0\phi_0\|_\infty+2\alpha+12\alpha\Sigma\)
there. The derivative of the normalizing section contributes
one term bounded by \(|\phi-\phi_0||\omega|\) and an edge
sum bounded by \(8\Sigma\max_{x'\sim x}|\omega(x')-\omega(x)|/\eta\)
before division by \(\rho_\phi\). To absorb the first-order
term, maximize the size of a first difference multiplied by
the smaller of one and the distance to the boundary. When
this distance is at most \(64\eta\), use the sup bound for
\(\omega\). Otherwise use the interior gradient estimate in
a cube of radius one quarter of the weighted distance. The
distances there and on the neighboring stencil are at least
half the maximizing distance. If \(H\) is the maximum, the
gradient estimate gives \(H\le Ce+CeH\). Decrease the
ceiling for \(e\) to absorb the last term.
At distance at least one this gives \(|\nabla\omega|\le Ce\),
and the equation bounds \(|\Delta_\eta\omega|\) on the
next interior stencil by \(Ce\).

Take the small logarithm of the normalized link, divided by
\(i\eta\), to obtain \(A\). Its first difference is controlled
by \(\nabla\omega\), and the Laplacian has one additional
chain-rule term bounded by \(C\eta|\nabla\omega|^2\).
The flat curl divergence is obtained separately. Expand the
four factors of a plaquette to first order and take the
transported differences of the three opposite plaquette pairs
of a full bond. The transported linear parts give
\(\partial^\sharp\partial A\) up to products of \(A\) and its
first differences. In the nonlinear parts the difference of
opposite remainders is bounded by \(C A_0A_1\eta^3\), and
transport of a quadratic remainder costs \(C A_0^3\eta^3\).
Here \(A_0,A_1\) are stencil suprema of \(|A|,|\nabla A|\).
The resulting bound is
\[
 |\partial^\sharp\partial A|(b)\le
 \eta^{-2}|(D_U^{\mathsf T}\partial U)(b)|+C(A_0A_1+A_0^3).
\]
Both second-order bounds are therefore \(Ce\) on the
prescribed interior. This supplies the real gauge in
\eqref{mb:local-gauge-input} with the full current hypothesis.

\paragraph{The cubes used in the layered profile.}
Let the target cube have side \(2m'\xi_j\), with \(m'\) a
multiple of the separation radius \(\rho\). Its concentric
enlargement of side \((2m'+4\rho)\xi_j\) lies in \(\Omega_j\)
and is shorter than half the torus. Suppose the curvature
and full-current radii on elements meeting that enlargement
are at most \(\epsilon\). In level-\(j\) units the radius bound
gives \(|\partial U-1|\le\epsilon\eta^2\) and \(|J(U)|\le\epsilon\).
Writing \(\partial U(p)=c(p)1+i\eta^2 F(p)\), the scalar
coefficient satisfies \(1-c(p)=|\partial U(p)-1|^2/2\).
In the transported divergence the imaginary coefficient is
\(i\eta^2J(U)\). The scalar difference of each opposite pair
is bounded by \(\epsilon^2\eta^4/2\), so at a full bond
\[
 |D_U^{\mathsf T}\partial U|
       \le (\epsilon+\tfrac32\epsilon^2)\eta^2 .
\]
For \(j\ge1\), use a cube of side \(2m'+2\mathfrak r_{\mathrm p}\)
aligned to the \(\mathfrak r_{\mathrm p}\) grid. Its left face
in each coordinate is the last grid face at least four units
to the left of the target. The cube then contains the full
four-unit neighborhood of the target. Its further concentric
enlargement by \(2L\mathfrak r_{\mathrm p}\) on each side and
the full plaquette stencils fit in the original \(2\rho\) margin.
These inclusions use the separation floor and
\(\mathfrak r_{\mathrm p}\ge13\). Apply the real cube gauge
with \(\alpha=2\epsilon\). The side of the aligned cube is
at most \(3m'\), so \(m'\epsilon\le c\) gives all four bounds
of size \(Cm'\epsilon\) on that cube. At level zero the mesh
is one in these units. The comb gauge already bounds the
potential by \(Cm'\epsilon\), and each difference in the
displayed operators is a bounded stencil of these values.

For the gradient Hölder bound compare sites within distance
one quarter in the scale units. Multiply the relevant
component \(A_\mu\) by a cutoff equal to one within distance
one and supported within distance two of the first site.
The cutoff and its neighbor stencils lie in the aligned cube.
Write the result as \(f=G'h\) for the constant-layer site inverse
\(G'=(\Delta+\varkappa_j Q'^*Q')^{-1}\), with \(\Delta\) the
positive flat Laplacian. The product rule and the four bounds
give \(\|h\|_\infty\le Cm'\epsilon\). The support of \(h\)
meets at most \(8^4\) unit blocks. Apply the flat gradient
Hölder bound to each supported part, using the same local
gradient at both sites since the cutoff is one there and on
their forward stencils. At larger distances the sup bound
on the gradient suffices. Rescaling gives
\eqref{mb:holder-reduction} with \(m=2m'\). On a critical
field the radius \(\epsilon\) can be taken from
\eqref{mb:weighted-maximum} on the entire enlargement.

\subsubsection{The site projector and the slice estimate}

Some of the gauge-slice algebra can be completed before
the quantitative kernel bounds. Put
\(\Delta=\Delta_{\rm s}\), and suppose
\(G'=(\Delta+Q'^\sharp aQ')^{-1}\) exists, \(Q'\)
is onto, and the site Gram matrix
\(B'=Q'G'^2Q'^\sharp\) is invertible. At real
data \(G'\) is symmetric. The formula in \eqref{mb:named-gauge-operators} then
makes \(R\) the orthogonal projection onto
\begin{equation}
 \ker(Q'G')=\Delta\ker Q'.
 \label{mb:site-range}
\end{equation}
Indeed \(G'Q'^\sharp(B')^{-1}Q'G'\) is the
orthogonal projection onto \({\rm ran}(G'Q'^\sharp)\).
If \(Q'\lambda=0\), then \(G'\Delta\lambda=\lambda\),
so \(Q'G'\Delta\lambda=0\). Conversely
\(\Delta\) is injective on \(\ker Q'\): if it
annihilates such a \(\lambda\), the invertible
\(\Delta+Q'^\sharp aQ'\) also annihilates it.
The two subspaces in \eqref{mb:site-range} have the same dimension.
This proves equality. On an admitted complex domain
the identities remain bilinear matrix identities.

In particular
\begin{equation}
 Q'G'R=0,\qquad \Delta G'R=R,\qquad
                 R\Delta G'R=R.
 \label{mb:site-inverse-identities}
\end{equation}
The first follows from the projector formula, and
the second from
\(\Delta G'=1-Q'^\sharp aQ'G'\).
For the actual regional sequence the site and inner-bond
Laplacians satisfy \(\Delta-D^\sharp D=\eta^{-2}n\), where
\(n(x)\) counts missing bare neighbors at the outer face.
The separation puts these sites in the bare layer, with all
their inner neighbors also in that layer. If \(f\) is supported
on these sites, \(G'^{-1}f\) is supported on bare tops.
The site row there is evaluation, so \(G'^{-1}f\) lies in
\({\rm ran}\,Q'^\sharp\). Hence \(Rf=0\) and
\(RD^\sharp D=R\Delta\). Also \(Q'\lambda=0\) forces \(\lambda\)
to vanish on the bare tops. These facts give
\begin{equation}
 \mathsf S^2=\mathsf S,\qquad
 RD^\sharp\mathsf S=0,\qquad
 \mathsf S D\lambda=0\quad(Q'\lambda=0).
 \label{mb:field-slice-projector}
\end{equation}
For example the quadratic term in the square of
\(1-DG'RD^\sharp\) reduces using
\(R\Delta G'R=R\). The last identity uses
\(\Delta\lambda\in{\rm ran}R\) and
\(G'\Delta\lambda=\lambda\).
Thus \(\mathsf S\) projects onto the fixed Landau
slice along the infinitesimal residual gauges.

If the actual straightened rows and site readout
satisfy \(QD=\bar DQ'\), then also
\begin{equation}
 Q\mathsf S=Q .
 \label{mb:slice-preserves-rows}
\end{equation}
This is immediate from \(Q'G'R=0\).
The intertwining has to be checked with the actual
outer-endpoint readout of every outward row.
Scale-tagged endpoint values cannot be silently
identified when they are averages at different
levels. For plain rows the corresponding site
readout is different even at a flat field.
\eqref{mb:slice-preserves-rows} is therefore a consequence of the
specified intertwining, not a reason to assume it.

There is a complete linear inverse for the site
equation. Let \(E_0\) be any admitted right inverse
of \(Q'\). For \(f\in{\rm ran}R\) the unique solution
of \(R\Delta\lambda=f,\ Q'\lambda=b\) is
\begin{equation}
 \lambda=G'f+(1-G'R\Delta)E_0b.
 \label{mb:site-pair-inverse}
\end{equation}
The first constraint follows from \(Q'G'R=0\).
For the other, \eqref{mb:site-inverse-identities} cancels the two terms
\(R\Delta E_0b\). A homogeneous solution has
\(\lambda\in\ker Q'\) and
\(\Delta\lambda\in{\rm ran}R\), hence
\(\Delta\lambda=0\), and is zero by the injectivity
already proved.
The formula is independent of the chosen \(E_0\),
since uniqueness identifies all choices.

For a real small field, the actual nonlinear gauge
equations are the Landau condition on the logarithm
of \(U^gU_0^{-1}\) and the pillar-top normalization
\(g(y)=v_n[A](y)\) at the top of a pillar of level \(n\).
The retained seam nodes are handled by the outward factors
in \eqref{mb:outward-node-identity}. Their derivative consists of
the site pair in \eqref{mb:site-pair-inverse}, with the signs determined
by the chosen infinitesimal gauge convention.
The infinitesimal gauge is defined by
\((d/dt)U^{e^{it\lambda}}|_0\), so the same sign
is used in \(D\), in the node derivative and in
\(\bar D\). The derivative of the node equation
is the declared \(Q'\lambda\), as checked from
the straightening identity in \eqref{mb:node-diagram}.

%
%
\paragraph{The real site correction.}
Use the weighted site inverse on the separated real sequence.
Write \(|f|_{(-2)}\) for the site sup norm with weight \(\xi^2\).
For \(f\in\operatorname{ran}R\) and node data \(\nu\), let
\(\lambda(f,\nu)\) be the solution of
\(R\Delta\lambda=f\), \(Q'\lambda=\nu\) given by
\eqref{mb:site-pair-inverse}. In the regional norms this solution
and \(R\) satisfy
\[
 \|\lambda(f,\nu)\|_{Y_{\rm site}}\le C(|f|_{(-2)}+\|\nu\|_\infty),
 \qquad |Rf|_{(-2)}\le C|f|_{(-2)}.
\]
The site norm includes the weighted first difference and the
zero-extension Laplacian. It controls the corresponding norm
\(Y_{\rm site}^{\circ}\) with the inner-bond Laplacian
\(D^\sharp D\). Indeed their difference is \(\eta^{-2}n\)
on the bare layer, with \(n\le8\), so the comparison constant
is at most nine.

Let \(A_0\) be a preliminary real chart potential, and write
\(A_0^w\) for the logarithm after applying \(w=e^{i\lambda}\),
with the reference \(U_0\) kept fixed. For a top of level \(j\)
the node factor of this change is
\(r_j=v_j[A_0^w]^{-1}w v_j[A_0]\). Let \(\mathsf T\) be the
node transport rotation in the straightened linearization.
The nonlinear divergence and node equations use
\[
 D^\sharp A_0^w=D^\sharp A_0-D^\sharp D\lambda+\mathcal N(A_0,\lambda),
 \qquad
 \mathsf T^{-1}\tfrac1i\log r_j=Q'\lambda+q(A_0,\lambda).
\]
For the first remainder the logarithmic gauge change gives
\[
 |\mathcal N(A_0,\lambda)|_{(-2)}
 \le C(\|A_0\|_X+\|\lambda\|_{Y_{\rm site}^{\circ}})
                                    \|\lambda\|_{Y_{\rm site}^{\circ}}.
\]
It also gives \(|A_0^w-A_0|_{(-1)}\le C\|\lambda\|_{Y_{\rm site}}\).
In the second remainder the straightened node derivative differs
from \(Q'\) by the curvature of the reference and the variations
of the averaging factors. If \(\alpha_0\) is the reference radius,
the resulting bound is
\[
 \|q(A_0,\lambda)\|_\infty
 \le C(\alpha_0+|A_0|_{(-1)}+\|\lambda\|_{Y_{\rm site}})
                                    \|\lambda\|_{Y_{\rm site}} .
\]
The factor at a bare top is exactly \(w\), so the node remainder
there vanishes. These bounds use the logarithms on the full
admitted ball and the literal recursive node factors.

Prescribe \(r_j=\mathsf r\) on the corresponding tops and put
\(\nu_*=\mathsf T^{-1}(\log\mathsf r)/i\). Let
\(m_0=|D^\sharp A_0|_{(-2)}+\|\mathsf r-1\|_\infty\).
On the closed real ball of radius \(C_0m_0\) in the pair
\((f,\nu)\), use the map
\[
 (f,\nu)\longmapsto
 \big(R[D^\sharp A_0+\mathcal N(A_0,\lambda(f,\nu))],
                   \nu_*-q(A_0,\lambda(f,\nu))\big).
\]
Choose \(C_0\) first to pay for the two linear forcing terms.
Then decrease the ceiling for \(\alpha_0,\|A_0\|_X\) and
\(\|\mathsf r-1\|_\infty\) so the remainders are at most half
the ball radius. The continuous map sends the ball into
itself, so the finite-dimensional fixed point gives the
real correction. Use \(RD^\sharp D=R\Delta\) to get
\(RD^\sharp A_0^w=0\), and the second equation gives the
prescribed top factors. The two estimates are
\(\|\lambda\|_{Y_{\rm site}}\le Cm_0\) and
\(|A_0^w-A_0|_{(-1)}\le Cm_0\). When \(m_0=0\) take
\(\lambda=0\). The construction specifies top factors. The
retained seam factors follow from the row equations.

The a priori improvement in that slice is a
different statement. It concerns \emph{any}
two regular configurations \(U_0,e^{i\eta A}U_0\)
already in the slice and with \(\|\mathsf L A\|_\infty\)
below a fixed admitted ceiling. Its target is
\begin{equation}
 \|A\|_X+\|\mathsf L^3D^\sharp DA\|_\infty
      \le B_1\big(\epsilon_0+
                       \|q^{\rm st}(A)\|_\infty\big),
 \label{mb:slice-reset}
\end{equation}
where \(B_1\) is independent of the depth and
does not depend on the larger preliminary
constant used to put \(A\) under the ceiling.
No criticality of these two configurations is
assumed. The currents of both configurations
control the covariant second-order forcing.
For the nonlinear group formula the forcing
must be computed with the reference logarithmic
differential \eqref{mb:reference-gradient}.

%
Put \(a_2=\|\mathsf L A\|_\infty\), and let \(\alpha_0\) be
the common regularity bound of the two configurations, in the
normalization of the slice. Write \(\beta_0\) for the maximum of
the logarithmic row coordinates \(|a_0(A)|\) on interior level-zero
rows, and \(\beta_1\) for the corresponding maximum of \(|a_j(A)|\)
on all interior and outward rows. Write
\(|f|_{(-3)}=\|\mathsf L^3 f\|_\infty\) for the bond norm.
On a full bond the Wilson currents at the two configurations and
the Taylor comparison of the logarithmic action bound the linear
second difference by \(3\alpha_0+C_L a_2\|A\|_X\) in this norm.
At an inner bond that is not full, every bond of its inner
plaquettes is a level-zero interior bond. There
\(a_0(A)(c)=\eta A(c)\), so the stencil gives the separate
bound \(24\beta_0\). Thus, with \(J'=D^\sharp DA\),
\[
 |J'|_{(-3)}\le
 \max\{3\alpha_0+C_L a_2\|A\|_X,24\beta_0\}.
\]
In the real slice \(RD^\sharp A=0\). Straightening the proper
rows gives \(QA=\widehat a-C^+(A)\), where \(\widehat a\)
is the conjugated row coordinate, with the same norm, and
\(\|C^+(A)\|_\infty\le C_L(\alpha_0+a_2)a_2\).
For data on these rows, extended by zero to all other rows,
\(|Q^*\omega|_{(-3)}\le 2(1+L^{-3})\|\omega\|_\infty\).
Indeed the finer row that could read a bond in the next layer
would be inward, so it is absent. The same-layer and next-coarser
row counts give the two displayed weights.

Use the actual inverse \(G_0\) of the regional Wilson Hessian
completed by \(DRD^\sharp+Q^*Q\). If \(\Delta'\) is the
Wilson Hessian minus \(D^\sharp D\), the slice equation is
\[
 G_0^{-1}A=J'+\Delta'A+Q^*(\widehat a-C^+(A)),
 \qquad |\Delta'A|_{(-3)}\le C_L\alpha_0a_2.
\]
The weighted regional inverse bound and \(a_2\le\|A\|_X\)
therefore give
\begin{equation}
 \|A\|_X\le C_1(\alpha_0+\beta_0+\beta_1)
                +C_2(\alpha_0+a_2)\|A\|_X.
 \label{mb:slice-absorption}
\end{equation}
For an interior row its logarithmic coordinate is the reduced
datum \(q^{\rm st}\). At an outward row the additional seam
factor is the block average at the outer end. It reads only
interior rows one level below, so the group angle inequality
gives \(\beta_1\le(1+\kappa_6)\|q^{\rm st}(A)\|_\infty\).
In particular no bound on inward or edge rows is needed.
With \(\alpha_0,a_2\) below a fixed ceiling for which the last
coefficient is at most one half, absorption followed by the
bound on \(J'\) proves \eqref{mb:slice-reset}.

To reach that slice from a real regular pair, first align
the mean contours level by level with the gauge equal to one
on the final pillar tops. The contour comparison gives a
relative link bound \(C_{\rm prod}(\alpha_0+\alpha_1)\) on
the carried links at each level, where \(\alpha_1\) bounds
the mismatch on the retained datum. Use the truncated
sequences up to each level in turn. At level zero the
link logarithm in this preliminary gauge has the same
bound in the field norm, since the mesh is one in its units.

At ascent to level \(j\) replace the previous logarithm
by \(A_0=L A^{(j-1)}\) in the new units. The products
\(\eta_j A_0\) and \(\eta_{j-1}A^{(j-1)}\) are equal,
so the old node factors are unchanged. The field norm
costs at most \(L^2\). At the old tops the correction
target \(\mathsf r\) is one. At the newly promoted tops
the old representative uses \(v_{j-1}[A_0]\) in place of
\(v_j[A_0]\). The correction target is the last one-step
node factor \(v_{j-1}[A_0]^{-1}v_j[A_0]\), with defect
bounded by \(C_L|A_0|_{(-1)}\) by the full block estimate.
With this target the correction cancels the initial node
mismatch. Its potential stays under the fixed ceiling when
\(C_LL^2\) times the previous field bound fits that ceiling.

The reduced proper rows now compare the two averaged
data directly. The edge row is a quotient of two links,
so the mismatch there costs at most twice the link bound.
At intermediate levels use the preliminary bound
\(C_{\rm prod}(\alpha_0+\alpha_1)\). At the final level use
\(\alpha_1\) on the actual retained rows: the preliminary gauge
fixes interior and outward endpoints and cancels at the
common seam node of an edge row. Then \eqref{mb:slice-reset}
resets the norm at every ascent to a fixed multiple of
\(\alpha_0+\alpha_1\). At the final seam the remaining
node factor is fixed by its inward-link equation, with
the retained endpoint and the appropriate orientation.
In particular, when the retained data are the same and the
coupled system identifies \(A=0\), the normalized gauge
fixes the tops and that link equation fixes every retained
seam node. This yields the residual orbit on the full
retained datum.

\subsubsection{The local improvement and its weighted maximum argument}

Define the distance \(d_{\mathbf B}\) by assigning a path
the length \(\int ds/\xi(s)\), where \(\xi(s)\) is the
local pillar scale, and taking the infimum over paths.
One may equivalently use the corresponding weighted
lattice path distance, with a fixed comparison constant.
At a seam the adjacent scales differ by \(L\); the
comparison is made using the same convention in every
kernel estimate. Test plaquettes are placed at a fixed
corner or center, at bounded distance in this metric.

Let \(r(x)\) be the maximum local dimensionless curvature
and current radius of a regular critical field, as
in \eqref{mb:action-regularity}, and let \(v(\tau)=|\partial\mathbf V(\tau)-1|\)
be its assigned datum tests. For a radius \(R\) in
this metric write \(r_R(x)\) for the supremum over
the \(R\)-ball and \(v_R(x)\) for the supremum of
tests in that ball. The actual local improvement
estimate to be established is
\begin{equation}
 r(x)\le C v_R(x)+
       \{C\sigma^2e^{-c\delta_*\sigma}
                              +C\sigma^4r_R(x)\}\,r_R(x),
 \qquad R\le C_RL\sigma .
 \label{mb:one-patch-improvement}
\end{equation}
The constants \(C,c,C_R,\delta_*>0\) must be independent
of \(\sigma\), depth, volume and region shape.
The patch must contain the actual determining
neighborhood of its cutoff fields and rows.
This is a critical-field estimate. It is not an
inequality for every field with small plaquettes.

%
%
The local proof is given under
\hyperlink{hi:real-datum-profile}{the recovery of the regularity radius}.
It uses an auxiliary separated sequence whose positive-level
members are inside the trivialized cube. Extend the local potential
by zero on the auxiliary torus. Constrained variations of this
extension patch to constrained variations of the original field,
which is how the critical equation is obtained in the auxiliary
flat slice. The current created by the extension is paid in the
bare-scale norm outside the positive-level members.

The datum gauge costs one further factor of \(\sigma\).
The datum and preliminary slice bounds are both at most
\(C\sigma^2r_R(x)\). Split the flat kernel at the inner cube, then
make the same split in the multiplier equation. The central
datum rows grow only linearly with distance from the observation,
while the outer rows and outer differentiated fields are at
distance at least \(c\sigma\). The field and row remainders have
two small factors in the preliminary slice. Thus the contribution
of the central datum, the near remainders and the far reads to the
plaquette and current estimates is bounded, respectively, by
\begin{equation}
 C v_R(x),\qquad C\sigma^4r_R(x)^2,\qquad
 C\sigma^2 e^{-c\delta_*\sigma}r_R(x).
 \label{mb:patch-forcing-list}
\end{equation}
This calculation is on the real separated geometry with the
auxiliary flat-slice and local gauge estimates. Here the distance
to a tested plaquette is the minimum over its corners. When an
inward row is kept one also counts the pinch tests, with distance
measured to their seam node. On a determining set in force all
inward rows have been removed, so only tested plaquettes remain.

We now give the entire maximum argument from
\eqref{mb:one-patch-improvement}. Choose \(\sigma\) so large that
\(C\sigma^2e^{-c\delta_*\sigma}\le1/8\).
Then choose the regularity ceiling \(r_*\)
so small that \(C\sigma^4r_*\le1/8\).
For \(r\le r_*\), \eqref{mb:one-patch-improvement} becomes
\begin{equation}
 r(x)\le C v_R(x)+\theta r_R(x),\qquad \theta\le\tfrac14.
 \label{mb:local-halving}
\end{equation}
Choose \(\delta>0\) with
\(\theta e^{\delta R}\le1/2\), and define
\begin{equation}
 \widetilde v_\delta(x)
   =\sup_\tau e^{-\delta d_{\mathbf B}(x,\tau)}v(\tau).
 \label{mb:datum-profile}
\end{equation}
The triangle inequality gives
\(\widetilde v_\delta(y)\le
 e^{\delta d_{\mathbf B}(x,y)}\widetilde v_\delta(x)\).
In particular \(v_R(x)\le e^{\delta R}\widetilde v_\delta(x)\).

For \(s>0\) let
\(M_s=\sup_x r(x)/(\widetilde v_\delta(x)+s)\).
It is finite since \(r\le r_*\). \eqref{mb:local-halving} gives
\begin{equation}
 M_s\le C e^{\delta R}
                       +\theta e^{\delta R}M_s,\qquad
 r(x)\le2C e^{\delta R}\widetilde v_\delta(x).
 \label{mb:weighted-maximum}
\end{equation}
The second inequality follows by absorption
and \(s\downarrow0\). No maximum has to be
attained, so the proof works on infinite
volume as well when the input radius is bounded.
If the datum tests in layer \(q\) are at most
\(\epsilon^{(q)}\), then
\begin{equation}
 r(x)\le C'\sup_q
       \epsilon^{(q)}
           e^{-\delta d_{\mathbf B}(x,{\rm layer}\ q)} .
 \label{mb:layer-profile}
\end{equation}
This proves the profile assertion from the
local improvement, with its exact choice order.

The unweighted estimate \eqref{mb:patch-forcing-list} closes
the radius in the admitted real induction. If a provisional
construction is regular below the input ceiling, its output
radius is $C'\|v\|_\infty$, independently of the larger
provisional constant. The conditions
for applying that improvement are checked
before it is used. One cannot use the smaller
output radius to justify a construction
already outside its input ceiling.

If \(v=0\), \eqref{mb:weighted-maximum} proves that curvature and
current vanish. This alone does not identify global holonomies.
If, in addition, the entire retained datum is trivial,
\(M_{\mathbf B}(U)=1\), the separated determining rows
\eqref{mb:determining-rows} imply \(U=1^g\) for
\(g\in\mathcal G_{\rm res}\), including when \(\Omega_0\)
has holes or several components.

Here is the geometric argument. Extend \(U\) by one outside
the inner bonds of \(\Omega_0\). A plaquette meeting its
complement has no corner in \(\Omega_1\), by separation;
each of its bonds is therefore either an original trivial
level-zero row or an added bond with value one. The extension
is flat and has trivial datum for the sequence with first
member replaced by the torus. If a periodic coordinate is too
short for a required pillar or double block to embed, make
this local construction on a compatible periodic cover.
The rows and their complete averaging reads lift with the
same values. The pillar transport rule below commutes with
the deck translations, so the resulting gauge descends.
On each pillar with top \(y\),
define \(t(x)\) as the transport from \(y\) to \(x\)
inside that pillar. Flatness makes this transport independent
of the path. Also every retained average equals its straight
path transport: on its complete double block a flat field
is a pure gauge, and exact covariance of the actual average
reduces this assertion to \(\mathcal M^q1=1\).

For a bond \(b=(x,x')\) inside one pillar, flatness gives
\(t(x)U(b)t(x')^{-1}=1\). For adjacent pillars of the same
order, close this path with their retained interior row
inside its double block; that row has transport one.
Adjacent pillars of different orders have orders \(q,q+1\).
The seam geometry gives a retained outward row from the upper
top to a site \(w\) whose level-\((q+1)\) block contains the
lower top. Join \(w\) to the lower top by level-\(q\)
interior rows in that block. These rows and the outward row
all have transport one, and they close the path through
\(b\) inside the outward row's complete double block.
Thus \(t(x)U(b)t(x')^{-1}=1\) in this last case as well;
no omitted inward row is used. Every node is a pillar top,
so \(t=1\) at the nodes. Restricting to \(\Omega_0\) and
setting \(g=t^{-1}\) yields
\(U(b)=t(x)^{-1}t(x')=1^g(b)\) with
\(g\in\mathcal G_{\rm res}\). The argument uses the
separated determining geometry, not a topological condition
imposed separately on the region.

\paragraph{Second-order and Hölder readouts}

For a local field \(h\), exponential decay of
the inverse gives a bound on \(h\) in its
declared field norm. It does not automatically
give a bound on a second derivative or on
the ordinary Laplacian. The actual field
equation must be used. In a real local
gauge the covariant Laplacian and the
componentwise Laplacian differ by terms
containing \(A\), a first difference of
\(A\), and at most one difference of \(h\).
On an admitted potential ball these satisfy
the local numerical comparison
\begin{equation}
 |(L_U+\Delta_\eta)h|
 \le C\big(
      |A|\,|\nabla h|+
      (|\nabla A|+|A|^2)|h|\big),
 \qquad L_U=\sum_\nu\nabla_{U,\nu}^\sharp\nabla_{U,\nu}.
 \label{mb:laplacian-comparison}
\end{equation}
Here \(\Delta_\eta\) is the nonpositive scalar
heat generator. The signs in \eqref{mb:laplacian-comparison} are
therefore a plus on its left. The estimate
follows by expanding each transported
difference and its transpose and using
\(|e^{i\eta A}-1|\le C\eta|A|\).
For the second-order coefficient difference,
subtract the two endpoint potentials and
use \(|\nabla A|\); the remaining product
has two \(A\)'s. Both contributions have
the displayed dimensions.

Given bounds on \(h,\nabla_Uh,L_Uh\) and
on the actual gauge coefficients, \eqref{mb:laplacian-comparison}
supplies a componentwise Laplacian bound.
The discrete interior theorem then gives
the first-difference Hölder bound.
It also supplies the coarse-mesh fallback
when the number of sites in the interior
ball is bounded. This removes a separate
Hölder estimate from the primitive list,
but not the pointwise gradient bound,
the covariant Laplacian bound or the local
gauge construction.

For a differentiated background at a
critical pair, the slice can give the
separate second-order member
\(D^\sharp Dh\). At a noncritical
independent-current pair the actual
linear equation instead controls the
combinations
\begin{equation}
 (D^\sharp D+DRD^\sharp)h,\qquad
                  D(1-R)D^\sharp h .
 \label{mb:off-current-readouts}
\end{equation}
Their sum has to be compared to the
covariant componentwise Laplacian with
the actual curvature and current terms.
The critical slice identity cannot be
inserted at an independent-current pair.
Near the outer boundary all these claims
are restricted to bonds with the required
neighbors; the missing-plaquette terms of
the action are separate boundary forcing.

\subsubsection{The quadratic real-to-complex reduction}

The square of the complex radius is forced
by the actual plaquette product. Consider
a real unitary \(U\), a complex physical
potential \(A=A_{\rm r}+iA_{\rm i}\), and
\(U'=e^{i\eta A}U\). On a patch of scale
\(\xi\), assume
\begin{equation}
 \xi|A|\le e_1,\qquad
 |\partial U-1|+|\partial U'-1|
                       \le e_0\eta^2\xi^{-2},
 \qquad e_0,e_1\le c .
 \label{mb:complex-curvature-input}
\end{equation}
Transport the four oriented increments to
the starting corner of a plaquette, using
the real base field. With raw increments
\(a_b=\eta A_b\), reversing an edge uses
\(a_{\bar b}=-\operatorname{Ad}(U_b^{-1})a_b\).
The exact product identity is
\begin{equation}
 \partial U'\,\partial U^{-1}
                   =\prod_{r=1}^4e^{iX_r}.
 \label{mb:plaquette-increments}
\end{equation}
Each \(X_r\) is an oriented \(a_b\) conjugated
by the preceding real base transport.
Consequently \(|X_r|\le\eta e_1/\xi\).
The finite product expansion gives
\begin{equation}
 \left|\prod_{r=1}^4e^{iX_r}-1-i\sum_{r=1}^4X_r\right|
                   \le C\eta^2e_1^2\xi^{-2}.
 \label{mb:plaquette-quadratic}
\end{equation}
One can prove this without a series rearrangement:
the second derivative of the product with
all \(X_r\) multiplied by one scalar is
bounded by \(C(\sum|X_r|)^2\) on the unit
disc, and the integral Taylor remainder
has the same bound.

The sum of transported increments is
\(\eta^2\mathrm{d}_UA\), up to the terms changing
the closing base transport to the identity.
Those terms are bounded by
\(C\eta|A|\,|\partial U-1|\).
Dividing \eqref{mb:plaquette-increments}--\eqref{mb:plaquette-quadratic} by \(\eta^2\)
and using \eqref{mb:complex-curvature-input} proves
\begin{equation}
 \xi^2|\mathrm{d}_UA|\le C(e_0+e_1^2).
 \label{mb:complex-linear-curvature}
\end{equation}
Indeed the closing error is at most
\(Ce_0e_1\eta\xi^{-3}\), bounded by
\(Ce_0\xi^{-2}\) since \(\eta\le\xi\).
The real base transports make \(\mathrm{d}_U\)
real linear on the Hermitian Lie algebra.
The same estimate therefore holds for
\(A_{\rm r}\) and \(A_{\rm i}\) separately.
Applying the product expansion again shows
that the real unitary configurations
\(e^{i\eta A_{\rm r}}U\) and
\(e^{i\eta A_{\rm i}}U\) have curvature
at most \(C(e_0+e_1^2)\eta^2\xi^{-2}\).
This is the real comparison used by the
continuation argument.

Here is the precise analytic step. Let
\(\mathcal F(z_1,z_2)\) be any one of the
fixed holomorphic background readouts
obtained from
\(e^{i\eta(z_1A_{\rm r}+z_2A_{\rm i})}U\),
with all its determining maps and gauges
admitted. Suppose
\begin{equation}
 \|\mathcal F(z_1,z_2)\|_{\rm sc}\le M
       \quad\text{for }|z_1|,|z_2|<r_0/(2e_1),
 \qquad e_1\le r_0/16 .
 \label{mb:complex-readout-disc}
\end{equation}
The subscript sc means that the output
is normalized by its physical scale.
It does not divide by the desired
smaller target radius. A two-variable
Cauchy estimate gives
\begin{equation}
 \|\mathcal F(1,i)-\mathcal F(1,0)
                  -i\mathcal F(0,1)+i\mathcal F(0,0)\|_{\rm sc}
                  \le C M e_1^2/r_0^2 .
 \label{mb:complex-polarization-remainder}
\end{equation}
For completeness, expand on a polydisc
of radius \(R=r_0/(4e_1)\ge4\).
Each coefficient of total degree \(n\)
is at most \(MR^{-n}\). The displayed
linear combination annihilates degrees
zero and one. There are \(n+1\) monomials
of degree \(n\); the sum over \(n\ge2\)
is bounded by \(CMR^{-2}\).
This proves \eqref{mb:complex-polarization-remainder} with no dependence on
the number of lattice variables.

The three real configurations on the
right of \eqref{mb:complex-polarization-remainder} are treated by the
real background estimate, using \eqref{mb:complex-linear-curvature}.
If the actual crude analytic estimate
has \(M/r_0^2\le C_{\rm geom}\), this gives
\begin{equation}
 \|\mathcal F(1,i)\|_{\rm sc}
                  \le C_{\rm geom}(e_0+e_1^2).
 \label{mb:complex-heredity}
\end{equation}
Thus the hereditary complex estimate is
reduced to a fixed-radius crude analytic
bound and the real estimate. It is not
an independent assumption of the
desired improved radius.
If \(r_0\) depends on a running field
normalization, its actual factor in
\(M/r_0^2\) must be retained. A bound
on a shrinking disc without this
ratio does not prove \eqref{mb:complex-heredity}.
The physical and Gaussian-coordinate
conversion in the domain estimates determines
which radius enters that ratio.

\subsubsection{Restriction is an exact constrained-variation argument}

The iterated average has a stronger locality than an
arbitrary composition of finite-range maps. A level-\(j\)
row reads the two level-\(j\) endpoint blocks and their
joining bonds. To prove this, proceed by induction.
At one step this is its definition. At the next
step, a block is the disjoint union of its \(L^4\)
subblocks. An internal contour reads only subblocks
inside that block. The straight segment between
corresponding points of two adjacent blocks reads
subblocks in their union and their joining bonds.
The lower-step locality closes the induction.
The fixed local lift includes all these paths.

For a determining row these two blocks miss the
next region, and they lie in the preceding region,
by the seam geometry. Their diameter in
\(d_{\mathbf B}\) is bounded by a constant depending
only on the fixed local geometry. The number of
bare variables can be large; the bound is on
the physical determining neighborhood, not on
that number.

Consider two aligned separated sequences with
\(\Omega'_q\subset\Omega_q\), and let
\(\mathbf B'\) be the new determining set.
The smaller regions give a refinement of the
local top scales. An old row's value can be
evaluated from the new determining values and
the fixed field outside \(\Omega'_0\).
Here is a terminating algorithm that proves
the assertion.

Start with the averaging expression for an
old row of level \(j\). No endpoint block lies
in \(\Omega'_{j+1}\), since it misses
\(\Omega_{j+1}\). If the row touches
\(\Omega'_j\), it is a new determining row
and its value is prescribed. Its full
neighborhood is inside \(\Omega'_0\) by
separation when \(j\ge1\). If neither
endpoint lies in \(\Omega'_j\), both
endpoint blocks are entirely outside that
region, by alignment. Expand its one-step
formula into level-\((j-1)\) row values.
Their endpoint blocks have no point in
\(\Omega'_j\). A lower row touching
\(\Omega'_{j-1}\) is prescribed; otherwise
expand it again. The index strictly
decreases. At level zero the bond is
either a prescribed new row or a bond
of the fixed exterior.

Rows whose determining neighborhood meets
the boundary of \(\Omega'_0\) cause no
exception. If a level-\(r\) endpoint were
in \(\Omega'_r\), the sum of the gaps
below that level and the chosen
\(\rho\gg L\) would put its two-block
neighborhood inside \(\Omega'_0\).
Thus such a row is expanded until its
bare exterior bonds are fixed. This
proves the finite factorization
\begin{equation}
 M_{\mathbf B}(U)
   =\mathcal F_{\mathbf B,\mathbf B'}
          (M_{\mathbf B'}(U|_{\Omega'_0}),
                             U|_{\rm exterior})
 \label{mb:refinement-factorization}
\end{equation}
for those rows affected by an interior
variation. The function \(\mathcal F\)
is built from the actual products,
means and averages just evaluated.
No linearization is involved.

Let \(U_t\) be a differentiable local
variation on the new fibre, and extend
it by the unchanged field outside
\(\Omega'_0\). The level-zero new
datum fixes every boundary-layer bond:
its free bonds can touch only
\(\Omega'_1\), which is separated
from \((\Omega'_0)^c\) by more than
two bare spacings. Consequently every
old plaquette changed by the variation
is an inner plaquette of \(\Omega'_0\).
The remaining old action terms are
constant, and \eqref{mb:refinement-factorization} fixes every old
constraint. Therefore
\begin{equation}
 \left.\frac d{dt}A_{\Omega'_0}(U_t)\right|_0
  =\left.\frac d{dt}A_{\Omega_0}(\widetilde U_t)\right|_0=0
 \label{mb:restriction-criticality}
\end{equation}
whenever \(U\) is critical for the old
problem. This also explains why an
arbitrary cut with freely varying
boundary bonds would not prove the
same statement.

The new local scales are no larger
than the old scales at their elements,
so the normalized fine curvature
bound improves.  The current needs
a separate boundary check, because
the new action omits some old
plaquettes.  Write \(j_{\rm old}\)
and \(j_{\rm new}\) for the two
left-variation currents, in the
physical normalization used above.
If a bond has the same incident
retained plaquettes in both actions,
these currents agree and its new
current weight improves.

Every remaining changed-current bond
is within a fixed bare plaquette
distance of the new boundary.
The gap between \(\Omega'_1\) and
\((\Omega'_0)^c\) therefore puts
it in the bare boundary layer,
where its new local scale is
\(\eta\).  The left-variation
formula has at most six retained
plaquette contributions.  Unitarity
and the fixed trace normalization give
\begin{equation}
 \eta^3|j_{\rm new}(b)|
 \le C_\partial
       \max_{\substack{p\subset\Omega'_0\\p\ni b}}|U_p-I|
 \le C_\partial\epsilon_{\rm old}.
 \label{eq:multilevel-background-MB89a}
\end{equation}
The maximum is zero when there is
no retained incident plaquette.
Here \(\epsilon_{\rm old}\) is the
old regularity parameter; the
normalized fine curvature estimate
bounds each displayed bare plaquette
by that parameter.  This proves current regularity
with a fixed boundary constant.
It does not require monotonicity
of the current under restriction
of the action.
The new datum is regular by the
actual averaging estimate. On its
two-block neighborhoods the smallest
old scale is at least its new scale
divided by \(L\), giving the fixed
seam loss \(C_LL^2\). Mixed tests are
the same physical averages by the
finite factorization.  Let
\(C_{\rm res}\) dominate \(1\),
\(C_\partial\), and the stated
datum and seam constant \(C_LL^2\).
The restriction identity is asserted
when \(C_{\rm res}\epsilon_{\rm old}\)
lies below all the corresponding
regularity input ceilings of the
real uniqueness theorem.  Under
this actual admission condition,
\begin{equation}
 U_{\mathbf B'}\!\left(
       M_{\mathbf B'}(U_{\mathbf B}(\mathbf V)|_{\Omega'_0})
                      \right)
       =U_{\mathbf B}(\mathbf V)|_{\Omega'_0}
       \pmod{\mathcal G_{\rm res}(\mathbf B')}.
 \label{mb:restriction-identity}
\end{equation}
This conclusion uses regular critical
uniqueness, not global minimization.
It is a complete reduction of the
restriction identity to that theorem.

The same argument proves the background
nesting used for the next step.
In a depth-\((k+1)\) problem the
single inner node of each last block
permits an axial mean representative
there. Its \(k\)-fold average on
the bonds touching \(\Omega_{k+1}\)
is independent of the representative
chosen: the gauge relating two such
representatives is one at the block
centers and, by their axial means,
at every level-\(k\) site in the
block. Covariance fixes those
coarse bonds. Freeze that averaged
field as the new last datum and
apply \eqref{mb:restriction-identity} to the truncation.
Thus the shorter background and
the longer one are the same orbit
on their common region. No new
effective action is introduced by
this identity.

\subsubsection{Collars and joined sequences}

Take \(\mu\ge\max\{\rho,\mathfrak M,L^2\}\) to be a power of \(L\),
with \(\rho\ge L^2\).
A part \(A_a\) is a union of cubes of side \(s_a=\mu\xi_{j_a}\)
at level \(j_a\). Write \(A_a^{(-r)}\) for the result of
removing \(r\) layers of these cubes, and use cube-wise meets
throughout. For distances use the closed representatives of the
cubes. Suppose that the four-layer region
\(F_a=A_a\setminus A_a^{(-4)}\) lies in the ambient layer \(j_a\)
and that for \(a\ne b\),
\[
       \operatorname{dist}_\infty(A_a^c,A_b^c)
                  \ge 10\max(s_a,s_b).
\]
Put \(W=\mathop{\sqcap}\limits_a A_a\),
\(W^{(-r)}=\mathop{\sqcap}\limits_a A_a^{(-r)}\), and
\(Q_a=A_a\setminus A_a^{(-3)}\).
If \(x\in F_a\), its distance from \(A_a^c\) is at most \(4s_a\),
so its distance from \(A_b^c\) is at least \(6s_b\).
The entire \(s_b\)-cube through \(x\) is then in \(A_b^{(-4)}\).
In particular \(Q_a\subset F_a\subset W\), and the wide collars
\(Q_a\) are pairwise disjoint.

The \(\mu\)-tower of \(A_a\) has \(A_{a,0}=A_a\). For
\(0\le q<j_a\), take
\begin{equation}\label{mb:eroded-tower}
 A_{a,q+1}
 =\bigcup_{\substack{C\text{ of side }\mu\xi_{q+1},\
                       C\subset A_{a,q}\\
               \operatorname{dist}_\infty(C,A_{a,q}^c)
                                      \ge\mu\xi_{q+1}}} C .
\end{equation}
Here the indicated aligned partitions exist whenever the
level-\(j_a\) partition of side \(s_a\) exists. The construction gives
\begin{equation}\label{mb:tower-gap}
 \operatorname{dist}_\infty(A_{a,q}^c,A_{a,q+1})
                   \ge\mu\xi_{q+1}>\rho\xi_q .
\end{equation}
At depth \(q\), every cube of side \(\mu\xi_q\) at distance at
least \(2\mu\xi_q\) from \(A_a^c\) lies in the tower. Indeed, let
\(C\) be such a cube at depth \(q+1\). Every depth-\(q\) cube
meeting its open \(\mu\xi_{q+1}\)-neighborhood has distance at
least \(\mu\xi_{q+1}-\mu\xi_q\ge2\mu\xi_q\) from \(A_a^c\).
By induction these cubes lie in \(A_{a,q}\), which proves that
\(C\) passes the next erosion. Thus \(A_a^{(-2)}\) lies in every
member of the tower of \(A_a\) up to depth \(j_a\).

Define
\[
 \Theta^a_q=
 \begin{cases}
 A_{a,q},&q\le j_a,\\
 A_a^{(-2)},&q>j_a,
 \end{cases}
 \qquad
 \Omega_q^W=\Omega_q\sqcap\mathop{\sqcap}\limits_a\Theta^a_q .
\]
The sequence decreases and has first member \(\Omega_0\sqcap W\).
For \(q\le j_a\) the mask is a union of ambient level-\(q\) cubes
by grid refinement. For \(q>j_a\), the ambient region misses \(F_a\).
Intersecting it with \(A_a^{(-2)}\) therefore selects exactly those
ambient cubes on the interior side of the collar. This proves
alignment.

To check separation at \(q\), first use the ambient gap for the
complement of \(\Omega_q\). For a mask with \(q<j_a\) use
\eqref{mb:tower-gap}. For \(q\ge j_a\), the points of
\(\Omega^W_{q+1}\) lie in \(A_a^{(-4)}\), while the complement
of \(\Theta_q^a\) lies at distance at most \(2s_a\) from \(A_a^c\).
At \(q=j_a\) this gives a gap of at least \(2s_a\).
At \(q>j_a\) a shortest segment from \(\Omega^W_{q+1}\) to
the complement of \(\Theta_q^a\) crosses the surface at distance
\(3s_a\) from \(A_a^c\). That surface lies in \(F_a\subset\Omega_q^c\).
The ambient gap then bounds the segment's length from below.
The same argument applies part by part, so the joined sequence
has the prescribed separation. A smaller first member can be
used when \(\Omega^W_1\) is a union of \(\mu\xi_1\)-cubes and
is at distance at least \(\mu\xi_1\) from
\((\Omega_0\sqcap W)^c\). In that case enlarge \(\Omega^W_1\)
by one layer of those cubes to obtain the collared first member.
It stays inside the full first member and keeps a gap of at least
\(\mu\xi_1\) from its own complement to \(\Omega^W_1\).

On \(W^{(-2)}\) the positive-index masks leave the ambient
regions unchanged. In particular, on bonds of the joined first
member with both ends in \(W^{(-2)}\), the determining rows
are the ambient rows. A new row of level \(q\) can arise only
when an end that was in \(\Omega_{q+1}\) is cut out of the
new region at index \(q+1\). That cut must occur in a tower
with \(q<j_a\), since the level-\(j_a\) boundary collar misses
\(\Omega_{j_a+1}\). The removed end is at distance at most
\(2\mu\xi_{q+1}\le2s_a\) from \(A_a^c\), by the interior
inclusion just proved. It lies in \(A_a\setminus A_a^{(-2)}\).
Its other end is within one level-\(q\) spacing and lies in
\(Q_a\). Both ends are therefore in the unique \(Q_a\) that
contains the row. These are the collar rows.

The localized datum agrees
with the original datum on
old rows. On a new collar
row it is the exact face
lift of the ambient
level-\(j_a\) datum:
\begin{equation}
 V^{\rm lift}(b)=
 \begin{cases}
  1,&b\text{ stays in one ambient }j_a\text{-block},\\
  V_{j_a}(c),&
     b\text{ crosses from the block of }c_-
                         \text{ to that of }c_+ .
 \end{cases}
 \label{mb:collar-face-lift}
\end{equation}
The orientation of a reverse
bond uses the group inverse.
On each successive grid
the central straight route
crosses a fixed coarse
face once. An internal
contour is one. The one-step
average of this lift is
therefore exactly the
preceding lifted datum,
by its defining exponential
average. This proves
constraint compatibility.
The lift can concentrate
curvature on its faces;
it is not asserted to
be a regular fine
configuration. Its
role is to prescribe
the collar data for
the localized background.

For a tested plaquette \(p\) of level \(q\), let
\(\mathcal R(p)\) be the union of the spacing-\(\xi_q\) cubes
of its corners. Its diameter is at most \(2\xi_q\) in the
sup distance. It is contained in \(W\), by separation at
positive levels and by the inner-bond condition at level zero.
Every joined test is of one of two kinds. Either
\(\mathcal R(p)\subset W^{(-2)}\), or
\(\mathcal R(p)\subset Q_a\) for a part with \(q\le j_a\).
To check this, fix \(a\). If \(q\le j_a\) and some point of
\(\mathcal R(p)\) has distance at most \(2s_a\) from \(A_a^c\),
the whole test lies at distance at most
\(2s_a+2\xi_q<3s_a\), hence in \(Q_a\). If no such point
exists, the test lies in \(A_a^{(-2)}\). If \(q>j_a\), a
corner in \(\Omega_q^W\) has distance at least \(4s_a\) from
\(A_a^c\). When \(q=j_a+1\), the test diameter keeps the
whole test at distance greater than \(2s_a\). When \(q>j_a+1\),
a segment to a point within distance \(2s_a\) would meet
\(F_a\subset\Omega_{q-1}^c\) within \(2\xi_q\) of the corner.
The ambient gap excludes this since \(\rho>2L\).

In the interior case all determining rows used by the test
are ambient rows. The holonomy, including its one-step
averages on mixed rows, is the original holonomy. In the
collar case the joined datum on the test is the lift from
level \(j_a\). A common row in that collar is necessarily
of level \(j_a\), so its prescribed value agrees with the
lift as well. The exact one-step averaging identity for
\eqref{mb:collar-face-lift} reduces a mixed test to a test
of this same lift. Its holonomy is one unless the plaquette
crosses two ambient faces, in which case it is the holonomy
of the corresponding level-\(j_a\) plaquette of \(Q_a\).
All its ambient defining rows lie in \(W\). Thus the
joined datum is determined by the ambient rows with both
ends in \(W\). It inherits their test bound \(\varepsilon_1\)
with no increase, provided the bound is imposed on the
ambient tests whose corner cubes are contained in \(W\).
The real background theorem applies to this regular datum
on the joined sequence when its geometric floors and
regularity radii hold. The reference field is the
normalized representative of the resulting regular
critical orbit.

\subsubsection{Response, collars and a common admitted path}

The source direction in a datum derivative includes the
reduced-datum map. At a critical solution, differentiating
\eqref{mb:stationarity} gives
\begin{equation}
 \begin{pmatrix}K&-Q^\sharp\\Q&0\end{pmatrix}
       \binom{\dot A}{\dot\omega}
       =\binom{0}{Db\,\dot B}.
 \label{mb:datum-response-system}
\end{equation}
No derivative of \(F\) with respect to \(B\) occurs
in fixed reference coordinates. If the reference
or an independent current also changes, its actual
derivative is a force in the first row and a
datum force in the second. The fixed-reference
response calculation supplies those terms.

Assume the actual coupled inverse has weighted
pointwise kernels for both columns and all
declared readouts. In normalized coordinates
write its field--datum entries as \(T_{yd}\).
For a point source at row \(c\), \eqref{mb:seam-datum-nilpotence} gives
\begin{equation}
 \dot A(y)=\sum_d T_{yd}(Db)_{dc}e,\qquad
 \|\dot A(y)\|_{X(y)}
                      \le C e^{-m d_{\mathbf B}(y,y_c)}|e|.
 \label{mb:point-response}
\end{equation}
The sum contains only \(d=c\) and outward
rows whose one-step outer block contains
the interior row \(c\). Its number and
distance from \(c\) are bounded at fixed
geometry; its length weights change by
only one factor of \(L\).
The triangle inequality proves the
exponential estimate from the inverse's
entry bound. The same calculation
applies to the multiplier and to each
proved differentiated field readout.

For an extended source, the required
bound is a row \emph{sum}, not just a
bound on each entry. If the normalized
matrix satisfies
\(\sup_y\sum_c e^{m d(y,y_c)}\|T_{yc}\|\le C\),
then a source supported on \(\mathcal C\)
obeys
\begin{equation}
 \|\dot A(y)\|_{X(y)}
      \le C e^{-m d_{\mathbf B}(y,\mathcal C)}
                              \|\dot B\|_\infty .
 \label{mb:extended-source}
\end{equation}
If kernels are written with cell measures,
those measures are included in this
Schur row sum. Mixing a measure-weighted
kernel convention with unweighted point
entries would give the wrong source
normalization.

The current derivative must be read
from the actual local Wilson current
at \(e^{i\eta A}U_0\), including its
logarithmic differential. Its stencil
contains second-order field terms.
The same is true of the covariant
Laplacian member. Once these readouts
and the pointwise gradient have the
weighted bounds, the interior theorem
gives their gradient Hölder member
with any fixed exponent below one.
The node derivative follows from
the actual differentiated node
factor. Its weighted bound is
another readout of the primitive
operator package. The mere algebraic
existence of \(v_n[A]\) does not
bound the product of its derivatives
over many levels.

Suppose two real data sets on a joined
sequence differ only on its collar
rows and have been put into a common
node frame. Write them there as
\(V_1=e^{iB}V_0\), where
\(\|B\|_\infty\le C_\mu\epsilon\).
The frame comparison is made by
the exact finite contour
reconstruction \eqref{mb:transported-bond}--\eqref{mb:rigidity-sum}.
At every collar it involves one
ambient level and its finer
tower; it must respect the
recorded old edges. This is
the datum alignment step.
The smallness condition on
\(\epsilon\) depends on \(\mu\).

For this comparison, the output norm \(X(y)\) uses the local
scale of the joined sequence, whereas \(d_{\mathbf B}\) is the
ambient metric. Joined scales are no larger than ambient scales,
so the native joined path distance dominates \(d_{\mathbf B}\).
The native weighted Schur bound therefore implies the same bound
with this ambient exponential weight and unchanged joined
source and output normalizations.

Take the actual path \(V_t=e^{itB}V_0\).
Assume it stays in the common admitted datum ball.
Integrating \eqref{mb:extended-source} gives a field comparison
whose forcing is supported only on the collar.
In the ambient metric each collar point is within \(C\mu\)
of \(W^c\), because the supplied collar lies in one ambient layer.
Hence
\begin{equation}
 \|f(y)\|_{X(y)}
    \le C_\mu\epsilon
           e^{-m d_{\mathbf B}(y,W^c)},\qquad
                  C_\mu\le C_\mu^{(0)}e^{Cm\mu}.
 \label{mb:boundary-comparison}
\end{equation}
Indeed \(d_{\mathbf B}(y,\mathcal C)\ge
d_{\mathbf B}(y,W^c)-C\mu\); \eqref{mb:extended-source}
and the fundamental theorem of calculus prove the assertion.
The factor \(C_\mu^{(0)}\) is the explicit finite alignment
and path constant. It is not claimed uniform as \(\mu\)
tends to infinity. Where the joined and ambient scales agree,
this is also an ambient output-norm bound. At a finer tower
element, stronger ambient readout weights require the separate
scale factors and lag or small-source estimate in
\eqref{mb:ambient-tower-weights}--\eqref{mb:settled-distance-gain};
they are not absorbed into \(C_\mu\).
\par{}
The physical comparison has the
form \(U|_W=(e^{i\eta f}U_W^\circ)^w\).
At tops away from the collar
the node identity prescribes
\(w=v_q[f]^{-1}\).
At other points it has no
specified normalization.
In particular \eqref{mb:boundary-comparison} is a
bound on the comparison
potential and its declared
readouts, not a bound on
that gauge at every fine
site.

\paragraph{Comparison with the real joined background.}
Let \(V\) be a real datum with all test holonomies within
\(\varepsilon_1\) of one. Fix the collar width \(\mu\).
Take the no-inward determining sets, the separation and
weighted-response floors of the real background theorem,
and at least \(6\rho\) sites in each direction at the last
level, with \(\rho\) taken above those floors. Decrease
\(\varepsilon_1\) so that the radii used below are admitted.
Let \(U\) be a critical configuration in the regular orbit
of \(V\), and let \(U^\circ_W\) be the normalized
representative of the regular orbit of \(V^W\).
The orbit radius \(\varepsilon_0\) obeys
\(B_3\varepsilon_1\le\varepsilon_0\le a_0\).
Regularity improvement puts \(U\) at radius
\(\epsilon=B_3\varepsilon_1\). By
\eqref{mb:restriction-identity}, \(U|_{\Omega_0^W}\) is
critical at its induced datum \(M_{\mathbf B^W}
(U|_{\Omega_0^W})\), which agrees with \(V^W\) on common
rows. We compare the induced datum with \(V^W\) in the
same node frame.

In \(Q_a\) with \(j=j_a\ge1\), each level-\(q\) average
\(W_q=\mathcal M^qU\), \(0\le q\le j\), obeys
\[
 \max_{p\text{ inner in }Q_a}|\partial W_q(p)-1|
                      \le2\epsilon L^{2(q-j)} .
\]
This is the averaging curvature bound in the ambient
layer \(j\). At \(q=j\) the inner bonds of \(Q_a\)
are determining rows, so \(W_j=V_j\) there.
With \(C_{\rm lift}\) the constant of the finite alignment estimate
\eqref{mb:transported-bond}--\eqref{mb:rigidity-sum}
there is a real gauge on \(Q_a\), equal to one on level-\(j\)
sites, whose level-\(q\) averages for \(q<j\) differ
from the lift of \(W_j\) by at most
\[
                 2C_{\rm lift}\epsilon/(L^2-1).
\]
The denominator comes from summing the smaller curvature
bounds over \(q<j\). The gauges on the wide collars
combine with the identity elsewhere to give a gauge
\(g\) on the joined first member. It leaves all common
rows unchanged. Indeed, if a common row has an end in
\(Q_a\), its layer end is either in \(F_a\), in which case
the row has level \(j_a\), or outside \(F_a\). In the latter
case the row would have length at least \(s_a\), hence
level at least \(j_a+2\) since \(\mu\ge L^2\). Separation
would put its end in \(Q_a\) inside \(\Omega_{j_a+1}\),
which is impossible. Thus on every common row meeting
a gauge support, the gauge equals one at the relevant
endpoints. There are no new rows of a part of level zero.

Consequently the induced datum of \((U|_{\Omega_0^W})^g\)
is \(e^{iB}V^W\), with \(B=0\) away from the collar rows.
For the discrepancy on a collar row use the intrinsic
logarithm. Once \(2C_{\rm lift}\epsilon/(L^2-1)<1\),
the logarithm is bounded by \(\pi/3\) times the group
discrepancy. In particular
\begin{equation}\label{mb:real-collar-datum}
 \|B\|_\infty\le
       \frac{2\pi C_{\rm lift}B_3}{3(L^2-1)}\varepsilon_1 .
\end{equation}
Take the right-hand side below the radius \(r_{\mathrm B}\)
of the joined response polydisc. The line \(tB\) now lies
in that polydisc. The real endpoints represent the
regular critical orbits at \(V^W\) and \(e^{iB}V^W\).
This uses gauge covariance of restriction and uniqueness
at the enlarged regularity radius of the induced datum.
Let \(\mathcal H_W\) denote the response in its normalized
slice. The comparison field is \(f=\mathcal H_W(B)\), and
restoring the gauges gives
\[
 U|_{\Omega_0^W}=(e^{i\eta f}U^\circ_W)^w .
\]
Outside the wide collars \(g=1\). The residual gauge
fixes the nodes, so at the level-\(q\) tops there
\(w=v_q[U^\circ_W;f]^{-1}\). If all parts have level zero,
\(B=f=0\) and the comparison is by a residual gauge.

Write \(\mathsf N_x(f)\) for the maximum of the potential,
first covariant difference and covariant \(D^{\mathsf T}D\)
readouts at a joined bond \(x\), multiplied by the first,
second and third powers of the joined top spacing,
respectively. The point response at a collar row \(c\)
with top \(y_c\) is bounded in this norm at the top \(y\)
of \(x\) by
\(B_5 e^{-\delta_0 d_{\mathbf B^W}(y,y_c)}\) times the
size of the row increment. Use the part of the rate
\(\delta_0/2\) that gives a uniform sum of row weights,
as in \eqref{mb:extended-source}. For a collar point,
the distance to \(W^c\) in the ambient metric is at most
\(12\mu\). Indeed the path to the corresponding
\(A_a^c\) stays in the single layer \(j_a\) until it reaches
the boundary and has bare length at most \(12\mu L^{j_a}\).
For a path from \(y\) to a collar row, stop the path
when it first reaches the boundary of \(W^{(-3)}\). Up to that point
the ambient and joined layer weights agree. The stopping
point is in a wide collar, and the preceding bound and
the triangle inequality therefore give
\[
 \inf_c d_{\mathbf B^W}(y,y_c)\ \ge\
                         d_{\mathbf B}(y,W^c)-12\mu .
\]
The inequality also holds if \(y\) already lies in a
wide collar. Integrating the derivative of
\(\mathcal H_W(tB)\) over \(0\le t\le1\), using
\eqref{mb:real-collar-datum}, gives
\begin{equation}\label{mb:real-collar-norm}
 \mathsf N_x(f)\le C_\mu\varepsilon_1
                      e^{-\delta_0 d_{\mathbf B}(y,W^c)/2}.
\end{equation}
Here the row sum and the alignment constant are independent
of the depth and volume, and the conversion of distances
costs at most \(e^{6\mu\delta_0}\).

In this gauge the difference of plaquette variables is
bounded by the first covariant differences and the
quadratic remainder in \(f\). The current difference
also uses \(D^{\mathsf T}Df\). On these fixed stencils
the assigned scales differ by at most one level, and
changing the observation point within the stencil costs
a fixed factor in the exponential. Hence the plaquette
difference at \(p\) and the bare-unit current difference
at a full bond \(b\) have, in addition to the right-hand
side of \eqref{mb:real-collar-norm} measured at the corner
or endpoint, the factors \(L^{-2\lambda^W(p)}\) and
\(L^{-3\lambda^W(b)}\), respectively. Here \(\lambda^W\)
is the assigned level in the joined sequence. This
last power uses \(J^{\rm bare}=L^{-3i}J^{(i)}\)
when the current is calculated in level-\(i\) units.

Two separated sequences which agree on \(W\), with real
data agreeing on all determining bonds with both ends
in \(W\), have the same joined datum and normalized
joined background. The distance from an interior point
to \(W^c\) is the same for the two sequences, since the
initial part of any path up to its first exit has the
same layer weights. Apply the estimate to each
background against this common joined background.
For plaquettes and full-bond current stencils in
\(W^{(-2)}\sqcap\Omega^W_0\) the assigned levels there
are the ambient ones. The triangle inequality gives
the comparison of the two readouts, with twice the
constant, in the gauges just specified.

The real comparison supplies the endpoint estimate on
these joined domains. To use it for complex retained
forms, we must first admit the full and restricted
histories on one domain of independent complex
potentials and currents, with the actual collar
substitutions and their remaining radius margins.
The real datum path above varies neither an independent
current nor the history's stored variables. The
fixed-tower complex response also keeps its determining
sequence fixed. A comparison of the complete forms
on that common domain, including the error from
discarded chains, is needed before transferring a
positive lower bound to a retained form.

\paragraph{The real base of a layered domain.}
Let \(e(x)\) be a regularity profile for a real critical
configuration, so it bounds the scaled curvature and the
current at full bonds based at \(x\). Use a separated
sequence with gaps greater than \(8\rho L^q\) between
levels \(q\) and \(q+1\), and with at least \(6\rho\) top
sites in each torus direction. On a nonempty aligned
localization \(X'\) of \(M_1L^i\)-cubes, intersect the
sequence with the \(M_1\)-tower of \(X'\).
For each depth \(p\le i\), restriction of the critical
configuration is critical at the induced datum by the
restriction argument for \eqref{mb:restriction-identity}.
The profile bound persists for the plaquettes and full
bonds of the localized first member. Their assigned
level can only decrease, and the current at those full
bonds has the same incident plaquettes. A current of the
restricted action at a changed stencil instead uses
\eqref{eq:multilevel-background-MB89a} and its bare-layer
bound. If \(e'\) is the profile supremum
on the localized first member, the tested plaquettes of
the induced datum have defects at most \(C_L e'\), by
averaging on their full backward cells. Use the same
datum ceiling for the tested seam holonomies if present.
Take \(e'\) below
the test and orbit ceilings, including \(L^2e'\) below the
curvature radius of one-step averaging and \(C_Le'\) below
the datum radius of the tower. The same restriction then
lies in the regular orbit of the induced datum. This
identifies the normalized tower configuration with the
restriction up to its residual gauge.

For a point of original level \(\lambda(x)\), the layer
of the localization is \(r=\min\{\lambda(x),i\}\).
A sufficient profile inequality for layer radius \(b_r\) is
\begin{equation}\label{mb:real-layer-profile}
         e(x)L^{-2(\lambda(x)-i)_+}\le b_r .
\end{equation}
In level-\(i\) units let \(\xi=L^{-i}\) and
\(\xi_r=L^{r-i}\). The curvature bound on this layer is
\(b_r\xi^2\xi_r^{-2}\) at plaquettes in the field's domain.
In converting the full-bond current from bare units,
the factor \(L^{3i}\) multiplies its bound
\(e(x)L^{-3\lambda(x)}\). The extra suppression at
\(\lambda(x)>i\) has exponent three, so the same
condition gives \(b_r\xi_r^{-3}\).
At a smaller observation level \(i''\), the radii on
layers \(r<i''\) can be kept, and the new top radius
can be chosen to dominate
\[
                    \max_{i''\le r\le i} b_rL^{-2(r-i'')}.
\]
This is the profile comparison for all nested
subdomains whose tower grids are aligned.

The gauge clauses require a separate local estimate.
For a cube of side \(C_{\mathrm g}ML^r\) in bare units,
take the profile supremum \(e_\square\) on its
\(8\rho L^{r-1}\) neighborhood. In the nonwrapping
geometry the local transport and critical background
estimate supplies a real gauge with both
\(\xi_r|A|\) and \(\xi_r^2|\nabla A|\) bounded by
\(C_L C_{\mathrm g}M e_\square\). In using it require
that the profile is below the one-step curvature
ceiling and that \(2\mathsf k_L C_{\mathrm g}ML
e_\square\le \rho a_1\), where \(\mathsf k_L\) is the
transport constant and \(a_1\) the regular datum
ceiling. Require also that this gauge bound fits the
gauge radius of the layer. These checks, together
with \eqref{mb:real-layer-profile}, give the real
initial triple \((U,0,J(U))\) in the layered domain
with the plaquette and full-bond tests just described.
If the domain also tests the restricted current at
changed stencils, its bare-layer radius must allow the
boundary bound above. The potential slot is zero here.
Admission of the varied potential and independent
current along the complex path requires the margins
in the following argument.

\paragraph{Admission is checked for every object}

The assumption on the common path
can be reduced to a finite
first-exit calculation. List
all the actual core quantities
\(\mathcal R_\beta(t)\), their
strict bounds \(r_\beta\),
and initial margins
\(s_\beta=r_\beta-\mathcal R_\beta(0)>0\).
Include small logarithm
branches, mean Jacobians,
current ceilings, gauge
condition numbers, site
Gram inverses, field
inverses and tower clauses.
If an inverse is an
atomic object, its norm
or its distance from
singularity is one of
these tests.

While the path is admitted,
suppose its actual response
and coefficient estimates
give
\begin{equation}
 |\mathcal R_\beta(t)-\mathcal R_\beta(0)|
                 \le K_\beta t\,\|B\|,\qquad
 K_\beta\|B\|\le s_\beta/2
                        \quad\text{for every }\beta.
 \label{mb:all-clause-path}
\end{equation}
At a proposed first exit
every inequality still
has half its margin.
On a finite lattice
the bounded coordinates
have a limiting point.
The branch and inverse
margins extend every
atomic formula there,
and the uniform coupled
inverse and coefficient
bound extend the solution
by its local contraction.
This contradicts a first
exit before \(t=1\).
The constants must be
uniform in the lattice,
but compactness here is
only the finite-dimensional
continuation argument.

Checking only curvature in
\eqref{mb:all-clause-path} would not suffice.
Nor does a response bound
at already admitted points
establish that every
substitution is admitted.
The actual inner substitution
also includes the fluctuation,
its amplified window datum,
the coarse background,
the cutoff, and the
independent-current readout.
Each has its own test
in this finite list.
The domain estimates give
the physical/Gaussian
radius conversion and
the slow-radius
inequalities used to
verify those numerical
margins.

\subsubsection{Gluing real constructions and their analytic maps}

There are two different gluing
arguments. For real data,
uniqueness among all regular
critical orbits identifies
constructions made in
different neighborhoods.
Applying the unique comb
normalization outside the
nodes identifies the
representatives themselves.
For complex data one uses
either a common contraction
neighborhood or analytic
uniqueness on the relevant
connected overlap.

The analytic uniqueness fact
can be proved directly.
If two holomorphic maps
on a connected complex
datum ball agree on a
real open ball, their
difference vanishes on
a complex polydisc.
Fix all coordinates
but one to be real and
apply the one-variable
identity theorem. Then
allow that coordinate
to be complex and
repeat with the next.
This proves vanishing
on the polydisc.
Continuation along
a finite chain of
intersecting polydiscs
in the connected
domain proves equality
everywhere.

For maps with different real
centers this argument
applies to each connected
component of their
complex overlap that
contains the required
real open overlap.
It does not identify
an unrelated overlap
component without such
an argument. An
alternative is to
put both solutions
into a single admitted
reference chart and
the single uniqueness
ball of \eqref{mb:contraction-conditions}.
That gives equality
pointwise on the
common chart,
without an overlap
connectedness claim.
The chart transition
includes its node
factors and gauge,
not just a difference
of coarse logarithms.

Let \(\mathcal D\) be the
real admitted datum
domain covered by
these local constructions.
It is open in the
finite product of
real group coordinates.
Every local fixed
point is continuous
in its datum by
the uniform
contraction estimate.
The comb representative
is continuous on the
admitted mean branches.
On overlaps the
representatives agree,
so these give a
continuous, hence
Borel, real background
map on \(\mathcal D\).
Equivalently choose
a countable chart
cover from the
countable basis of
this finite-dimensional
manifold and use
the first chart
containing each point;
overlap agreement
makes this a
Borel map.

If an integral uses this map
only under a cutoff
supported in \(\mathcal D\),
one can assign the
identity outside
\(\mathcal D\) to
obtain a Borel
extension. This is
only a definition
outside the regular
problem. It makes
no claim about its
fibre or criticality.
Holomorphic integrals
use the actual common
complex domain and
its parameter-independent
reference measure;
the Borel extension
is not an analytic
substitute for that
domain.

\paragraph{Measurability of the real background tests.}\hypertarget{real-background-tests}{}
Fix the finite lattice, determining regions, and numerical parameters.
In real matrix coordinates, $G$ is algebraic. The intrinsic logarithm
has the graph
$X\in\mathfrak g$, $|X|<\pi$, $e^{iX}=H$; the exponential on this
bounded set is restricted analytic. The finite equations defining a
mean, together with the requirement that their solution be unique,
therefore define its graph and its domain in $\mathbb R_{\mathrm{an}}$.
The actual three branches of $\mathcal M$, its finite iterates,
the determining-data map, and the prescribed axial gauge conditions
are definable as well. We use closure under real quantification and
cell decomposition in this structure \cite{Edmundo2000}.

It remains to check the criticality condition, which quantifies over
all differentiable curves in a fibre. Write
$C_d=\{U:M_{\mathbf B}(U)=d\}$ and regard it as a definable subset
of a finite-dimensional real matrix space. At $U\in C_d$ define
the contingent cone by the finite formula
\[
 \begin{split}
 v\in T_d(U)\quad\Longleftrightarrow\quad
 &\forall\epsilon>0\ \forall\delta>0\ \exists t,z:\\[-2pt]
 &0<t<\delta,\quad |z-v|<\epsilon,\quad U+tz\in C_d.
 \end{split}
\]
The velocities of two-sided differentiable curves in $C_d$ through
$U$ are exactly $T_d(U)\cap(-T_d(U))$. One inclusion follows from
the difference quotients. For the converse, if $v\in T_d(U)$,
then $(0,v)$ belongs to the closure of the definable set
$\{(t,z):t>0,\ U+tz\in C_d\}$. Definable curve selection and
one-variable monotonicity \cite[Lemma~2 and Proposition~1]{Kurdyka1998}
give, after reparametrizing by the positive coordinate $t$, a
differentiable curve $z(t)$ tending to $v$. Thus $U+tz(t)$ lies in
$C_d$ and has right derivative $v$ at $0$. If also $-v\in T_d(U)$,
apply the same argument to $-v$ and reverse time on that branch.
The two branches joined at $U$ are differentiable on an open interval
and have derivative $v$ at $0$.

The Wilson action has a polynomial extension in real matrix
coordinates, using the conjugate transpose for inverse link
variables. Hence the source criticality condition is exactly
\[
 \forall v\in T_d(U)\cap(-T_d(U)),\qquad
             DA_{\Omega_0}(U)v=0.
\]
It is therefore definable even at a singular fibre. Adding the
prescribed regularity inequalities defines the relation
$\mathscr C(d,U)$ of all regular critical configurations. The
minimal-orbit domain is given by
\[
 \begin{split}
 &\exists U\ \mathscr C(d,U),\\
 &\forall U_1,U_2\quad
  \mathscr C(d,U_1)\wedge\mathscr C(d,U_2)
  \Longrightarrow
  \exists g\in\mathcal G^{\mathrm{res}}(\mathbf B)\quad
                     U_2=U_1^g.
 \end{split}
\]
The residual group fixes every prescribed node, including seam nodes;
it is a finite-dimensional compact algebraic group. Gauge invariance
of $\mathscr C$ makes this formula precisely the condition that the
nonempty set of regular critical configurations be one orbit.

The gauge-invariant action and plaquette readouts have definable
graphs obtained by quantifying over this orbit. The prescribed
representative in $\operatorname{Ax}^{\circ}$ has a definable graph
on its exact domain of existence and uniqueness. For a block axial
readout on $Y$, retain the conditions that $Y$ consists of nodes,
that $B(Y)\setminus Y$ contains no nodes, and that the required means
exist. Gauge covariance supplies a representative in this gauge;
two such representatives differ by a residual gauge transformation
equal to the identity on $B(Y)$. Their averaged fields therefore
agree on the interior bonds used by the test. Existential
quantification over these representatives defines that same
single-valued readout, without choosing a different representative.

Finally, in the averaged-background problem let $\mathsf W_d^{\mathrm{reg}}$
be the prescribed constrained regular set, and write
$a_d(w)=A_K(U_k(w))$. The preceding orbit formula and gauge invariance
make both this set and this action readout definable. The graph of
its unique attained minimum is exactly
\[
 \begin{split}
 w\in\mathsf W_d^{\mathrm{reg}},\qquad
 &\forall z\in\mathsf W_d^{\mathrm{reg}},\quad a_d(w)\le a_d(z),\\
 &\forall z\in\mathsf W_d^{\mathrm{reg}},\quad
       a_d(z)=a_d(w)\Longrightarrow z=w.
 \end{split}
\]
Its projection is the exact existence set used in the definition.
Cell decomposition makes these domains and their single-valued maps
Borel. Consequently the real background tests, with their prescribed
failure values when an object is undefined, are Borel. Existence of the required orbit and attainment of the minimum on
the intended supports use the corresponding real estimates.

\subsubsection{Two finite mechanisms that prevent growth with depth}

\paragraph{Eliminating the multiplier}

It is unnecessary to assume an a priori small
multiplier in a uniqueness statement once its
actual row lift is bounded. Fix
\(Q_0E_0=1\), and write
\(Dq(A)=Q_0+C'(A)\). At a critical
field, \eqref{mb:multiplier-existence} gives
\begin{equation}
 [1+(C'(A)E_0)^\sharp]\omega=E_0^\sharp DF(A).
 \label{mb:eliminate-multiplier}
\end{equation}
If \(\|(C'(A)E_0)^\sharp\|_{Z\to Z}<1/2\),
the bracket has its convergent Neumann inverse,
and the multiplier is the explicit function
\begin{equation}
 \omega(A)=
    [1+(C'(A)E_0)^\sharp]^{-1}E_0^\sharp DF(A).
 \label{mb:multiplier-function}
\end{equation}
The required transpose bound is a consequence of
the column as well as row estimate of the
actual lift, with the declared force and
multiplier amplitudes. It does not follow
from an unweighted supremum bound on \(E_0\).

Suppose on a field ball
\(\|DF(A)\|_{X^*}\le\gamma\),
\(\|DF(A)-DF(A')\|_{X^*}\le C_F\|A-A'\|_X\),
\(\|E_0^\sharp\|\le C_E\), and the bracket's
perturbation has Lipschitz constant \(C_N\).
The inverse difference identity gives
\begin{equation}
 \begin{split}
 \|\omega(A)\|_Z&\le2C_E\gamma,\\
 \|\omega(A)-\omega(A')\|_Z
       &\le(2C_EC_F+4C_E\gamma C_N)\|A-A'\|_X .
 \end{split}
 \label{mb:multiplier-bounds}
\end{equation}
The second constant need not be small.
However a field fixed-point equation reads
it through the small factor \(C'(A)^\sharp\).
If \(\|C'(A)^\sharp\|\le cr\) and its
Lipschitz constant is \(c\), then
\begin{equation}
 \|C'(A)^\sharp\omega(A)
                 -C'(A')^\sharp\omega(A')\|_{X^*}
 \le \{cr\,{\rm Lip}(\omega)+2cC_E\gamma\}
                                      \|A-A'\|_X .
 \label{mb:small-multiplier-entry}
\end{equation}
The coefficient is small with \(r+\gamma\).
This is the field-only form of the coupled
contraction. It shows why a non-small
derivative of the multiplier is not an
obstruction.

Alternatively suppose the two-component
iteration has a Lipschitz matrix bounded by
\(\left(\begin{smallmatrix}\epsilon&\epsilon\\C&\epsilon\end{smallmatrix}\right)\).
In the norm
\(\max(\|A\|_X,\Theta^{-1}\|\omega\|_Z)\)
its row sums are at most
\(\epsilon+\epsilon\Theta\) and
\(C/\Theta+\epsilon\).
Choose a fixed \(\Theta\ge4C\) first,
then \(\epsilon(1+\Theta)\le1/4\).
Both row sums are less than one half.
The admitted multiplier radius must be
at least \(\Theta\) times the field
radius used by this ball. These choices
are made with the actual physical
norms, not by suppressing the multiplier
dimension in notation.

\paragraph{Separation certificates for the region operations}

For sets, write the gap condition as
\(N_{\rho\xi_q}(\Omega_{q+1})\subset\Omega_q\),
where \(N_r\) is the metric neighborhood.
Unions preserve it, since
\(N_r(\bigcup A)=\bigcup N_r(A)\).
Intersections also preserve it, since
\(N_r(\bigcap A)\subset\bigcap N_r(A)\).
Thus levelwise cube-wise unions and
meets of separated aligned sequences
are separated and aligned.

Enlargement and rounding spend a
computable part of the gap.
If the next region is enlarged by
at most \(a\xi_{q+1}\), its distance
to the old complement decreases by
at most \(aL\xi_q\). Rounding to
aligned \(\mathfrak M\xi_{q+1}\)
cubes adds at most
\(C\mathfrak M L\xi_q\).
Hence
\begin{equation}
 \operatorname{gap}_{\rm new}(q)
   \ge \{\rho-L(a+C\mathfrak M)\}\xi_q .
 \label{mb:gap-spending}
\end{equation}
An enlargement of the current region
can only improve this lower bound.
If the total finite list of
enlargements spends at most half
the initial coefficient, all
resulting gaps remain at least
\(\rho\xi_q/2\).
The list must be counted for the
actual construction; repeating an
unrecorded erosion arbitrarily many
times is not covered by this claim.

This gives a mathematical certificate
for each sequence used in the
induction: its levels are unions
and meets of the declared cube
families, its finite enlargements
have the total cost \eqref{mb:gap-spending},
and the only change of local
scale in a collar is the
single-layer change proved in
\eqref{mb:eroded-tower}--\eqref{mb:tower-gap}.
The original label sequence,
the extension by the next-step
domain, healed sequences,
prepared sequences and stage
sequences, towers, joined
strip and glued regions,
and cube-wise localizations
all require such a
certificate from their actual
definitions. The history
construction supplies the
first cube families; the
elementary closure facts
above supply the geometric
conclusion once that finite
certificate is verified.

\paragraph{The two-radius datum descent}

A section may lose a constant
\(C_\downarrow\ge1\) on tests
near its seam and be exact
away from that seam. It
would be incorrect to
iterate the global bound
\(\epsilon\mapsto
C_\downarrow\epsilon\).
The relevant object is
the finite dependency
graph of tests.

At every descent stage assign
an input vertex to each
actual test read by the
section and an output
vertex to each resulting
test. A good edge carries
coefficient one. A bad
edge carries at most
\(C_\downarrow\) and has
its input in the seam
determining collar.
The finite section
estimate has the form
\begin{equation}
 t_{\rm out}(\tau)\le
       \max_{\tau'\to\tau}
               c(\tau',\tau)t_{\rm in}(\tau'),
 \qquad c\in\{1,C_\downarrow\}.
 \label{mb:test-dependency}
\end{equation}
A sum over finitely many
local tests is absorbed
in \(C_\downarrow\) at
the bad edge, not added
at each later exact
transfer.

Assume that a seam collar's
full test-determining
neighborhood is disjoint
from the output footprints
of the previously spoiled
collars. This is the
specific collar condition
to check using \eqref{mb:gap-spending};
disjoint geometric faces
without their determining
neighborhoods would not
suffice. Then a directed
dependency path meets
at most one bad edge.
Iterating \eqref{mb:test-dependency} gives
\begin{equation}
 t_{\rm final}(\tau)
     \le \max_{\pi\ {\rm ending\ at}\ \tau}
           \prod_{e\in\pi}c(e)\,
                   \max t_{\rm initial}
     \le C_\downarrow\epsilon .
 \label{mb:datum-no-growth}
\end{equation}
This proves the uniform
descent bound.
At the input to each
fresh collar the radius
is still \(\epsilon\).
Away from those collars
the carried radius may
be \(C_\downarrow\epsilon\),
but the next exact
transfer does not
multiply it.

The estimate is a complete
finite graph argument.
Its model inputs are
the actual one-step
section estimate and
the collar-neighborhood
certificate. A pointwise
regularity estimate
for a section with
no record of where
it loses its constant
does not imply this
two-radius conclusion.
The same distinction
is needed when sections
are folded onto boxes:
every reflected or
translated determining
neighborhood must
remain in its recorded
collar.

\subsubsection{Criticality through two actual averaging levels}

The induction changes the number of constraints.
Its chain rule must use the induced fine action,
not a fresh coarse Wilson action.
Let \(m\) be the actual datum map for a truncation,
let \(r\) be its next-step average, including
the actual outward reduced datum, and suppose
the weaker constraint is \(r\circ m\).
Let \(\Psi(V)\) be a differentiable section
with \(m(\Psi(V))=V\), critical for the
fine action \(F\) on each \(m\)-fibre.
Define \(S(V)=F(\Psi(V))\).

Every fine tangent vector \(h\) at \(\Psi(V)\)
has the exact decomposition
\begin{equation}
 h=D\Psi(V)\,Dm(h)+h_{\rm vert},\qquad
                           Dm(h_{\rm vert})=0.
 \label{mb:vertical-decomposition}
\end{equation}
This follows by differentiating \(m\circ\Psi=1\).
The derivative of \(F\) annihilates \(h_{\rm vert}\).
Consequently
\begin{equation}
 DF[h]=DS[Dm(h)] .
 \label{mb:transitive-variation}
\end{equation}
If \(Dr\,Dm(h)=0\) and \(V\) is critical
for \(S\) on its \(r\)-fibre, this is zero.
Conversely, if \(\Psi(V)\) is critical
on the \((r\circ m)\)-fibre, test with
\(h=D\Psi\,v\), \(Dr\,v=0\), to see
that \(V\) is critical for \(S\).
Thus
\begin{equation}
 \Psi(V)\text{ is critical on }(r\circ m)^{-1}(r(V))
 \ \Longleftrightarrow\
 V\text{ is critical for }S\text{ on }r^{-1}(r(V)).
 \label{mb:criticality-transitivity}
\end{equation}
In orbit coordinates, choose the admitted
representatives and carry the node gauges
through both maps. The residual gauge
directions contribute zero by invariance.
The same proof then applies to the
quotient tangent vectors. Its geometric
premise, \(r\circ m\) equal to the
desired weak constraint, is the
finite averaging identity and the
outward factor in \eqref{mb:outward-node-identity}; it is not a
formal convention about the names
of the two data sets.

This supplies a precise noncircular
order for constructing depth \(k+1\).
Assume the real backgrounds through
depth \(k\) have already been
constructed on their admitted
domains, together with the needed
actual readouts.
Complete the last coarse datum
to a datum for the depth-\(k\)
truncation by the one-step
section, leaving all earlier
rows unchanged. At an outward
last row its pivot crosses
from an inner block to the
outer block. That pivot is
not an earlier kept inward
row, precisely because
those rows were omitted.
The finite pivot chart
therefore supplies a
local constraint correction.
Its regular-test and
collar estimates are the
section task in \eqref{mb:section-input}
and \eqref{mb:test-dependency}, not an
assumption of a critical
field.

The available background at the
completed datum is a trial
field for the weaker problem.
It is already critical
in its vertical directions,
but need not be critical
for the weak constraint.
Its remaining force is
the derivative of the
induced action \(S_k\)
on the last coarse
datum. Its Hessian is
the actual readout
\eqref{mb:coarse-action-hessian}--\eqref{mb:next-step-form}.
The primitive inverse
and finite coefficient
estimates for that
readout are established
on the real primitive
class before solving
the weaker problem.
They may use the
backgrounds of depths
at most \(k\); they
must not use the
unknown depth-\((k+1)\)
critical field.

The lower background is evaluated
on the admitted analytic
datum neighborhood already
constructed at depth \(k\).
Exploring this coarse
variation does not
restart the real datum
descent at every point.
The marked two-radius
class is needed to
construct its real
center; the actual
whole-disc continuation
and its fixed margins
control the subsequent
coarse variation.

Apply the one-step coupled
construction to \(S_k\)
and its actual plain
constraint chart, with
its prescribed soft
gauge. The soft gauge
has to be removed or
identified by the
actual gauge equation
on the physical
solution; its mere
positivity does not
make a critical point
of the soft action
critical for \(S_k\).
Once that identity
is proved, \eqref{mb:criticality-transitivity}
makes the lifted fine
field critical for
the desired depth-\((k+1)\)
problem.

Suppose its provisional
regular radius is
\(C_{\rm up}C_\downarrow\epsilon\).
All constants here are
the fixed section,
inverse and coefficient
constants, independent
of \(k\). Require
\begin{equation}
 C_{\rm up}C_\downarrow\epsilon<r_*,
 \qquad
 C_{\rm prof}\epsilon<r_{\rm induct},
 \label{mb:induction-radius}
\end{equation}
where \(r_*\) is the
input ceiling of
\eqref{mb:one-patch-improvement}, and
\(r_{\rm induct}\)
is the radius promised
to the next induction
step. The profile
estimate then resets
the output to
\(C_{\rm prof}\epsilon\).
The two-radius datum
certificate prevents
\(C_\downarrow\) from
compounding during
the completion.
This proves closure
of the radius
induction from its
listed estimates.

At depth zero all the
fine bonds are datum:
there are no free
fibre directions.
The finite plaquette
derivative gives
\(\xi_0^3|J_A|\le
C\sup|\partial V-1|\),
because a bond meets
at most six inner
plaquettes and
\(\xi_0=\eta\).
This is the base
regularity assertion.
The first nontrivial
constraint chart and
physical background
step are the finite
results already
proved. They realize
those instances of
the induction, not
the uniform
multilevel estimates
that its later
steps require.

\paragraph{Uniqueness among regular orbits}

Take two regular critical
fields for the same
datum. Put the first
in the reference role.
The preliminary contour
alignment and \eqref{mb:rigidity-sum}
give a raw closeness
below the fixed gauge
ceiling after choosing
the common regularity
radius small enough.
The actual site
correction puts the
second field in the
reference slice.
Its a priori estimate in \eqref{mb:slice-reset} gives
\(\|A\|_X\le B_1\epsilon\),
independent of the
preliminary constants.
Both configurations
then lie in the
single uniqueness
ball of the coupled
construction.

There is no independent
multiplier hypothesis:
\eqref{mb:multiplier-function}--\eqref{mb:multiplier-bounds} bound
the multiplier from
the actual current
and row lift.
Thus the two solutions
of \eqref{mb:stationarity}, with the
same reduced datum,
are equal in that
slice. Their relating
gauge is one at the
nodes by the node
normalization.
They belong to one
residual orbit.
On the admitted real regional sequence the site correction
and \eqref{mb:slice-reset} give the field in the ball used here.
The uniqueness assertion also uses the admitted coupled system
for the full set of retained rows.

\subsubsection{Why the physical soft gauge disappears from stationarity}

There is an exact local argument for the gauge
step used in the physical construction. Let \(S(V)\)
be a gauge-invariant coarse action and fix its
plain one-step averages. Assume each interior
block has no prescribed node except its center.
For each noncenter \(x\) let \(T_{y,x}(V)\) be
the admitted mean transport from its center.
Let \(T^0_{y,x}\) be a fixed reference transport,
and put
\begin{equation}
 G_{\rm soft}(V)=
   \sum_{y,x\ne y}\kappa_{y,x}
       \left(1-\tfrac12\operatorname{tr}
             [T_{y,x}(V)(T^0_{y,x})^{-1}]\right),
 \qquad \kappa_{y,x}>0 .
 \label{mb:physical-soft-gauge}
\end{equation}
The coefficients are fixed on the
variation under consideration.
The absolute gauge is the case
\(T^0_{y,x}=1\).
All transports and reference ratios
lie on the regular branch near
the identity.

A gauge \(h\) equal to one at
the block centers preserves the
plain average by covariance.
Vary only \(h(x)=e^{itH}\)
at one noncenter site.
Every other mean transport
is unchanged, and
\(T_{y,x}\) changes to
\(T_{y,x}e^{-itH}\).
Gauge invariance gives
no derivative of \(S\).
If \(S+G_{\rm soft}\) is
critical on this fibre,
\begin{equation}
 \operatorname{tr}
       [(T^0_{y,x})^{-1}T_{y,x}H]=0
                  \quad\text{for all traceless }H .
 \label{mb:soft-gauge-stationarity}
\end{equation}
The orthogonal complement of the
traceless \(2\) by \(2\) matrices
under the trace pairing is the
scalar matrices. The determinant
is one, so
\((T^0)^{-1}T\) equals \(1\)
or \(-1\). The admitted branch
excludes \(-1\). Thus
\begin{equation}
 T_{y,x}(V)=T^0_{y,x}
                      \quad\text{at every such site}.
 \label{mb:soft-gauge-vanishes}
\end{equation}
At these values the differential
of \(G_{\rm soft}\) is zero
for every field variation,
not only a gauge variation:
the derivative of the trace
at the identity annihilates
the traceless tangent.
Consequently the critical
point is critical for
\(S\) on the fibre.

Conversely a physical critical
point in this prescribed
transport gauge is critical
for \(S+G_{\rm soft}\).
The proof holds for real
unitary fields and on an
admitted complex branch
using the bilinear
holomorphic derivative.
The positive real coefficient
may be replaced in the
complex statement by an
admitted nonzero coefficient.

This argument explains all its
scope requirements. A retained
seam node cannot be varied,
so \eqref{mb:soft-gauge-stationarity} need not hold
there. The no-inward geometry
removes that obstruction.
A sharp regular-box cutoff
must be constant on the
stationary neighborhood;
otherwise its derivative
is an extra term.
An independent-current
auxiliary action containing
linear source subtractions
is not the gauge-invariant
physical \(S\) in this
argument. Its physical
matching must first
identify the action
and gauge terms.
Finally the choice of
relative or absolute
gauge still matters
off the axial locus,
as \eqref{mb:absolute-gauge}--\eqref{mb:relative-gauge}
show.

\subsubsection{Ambient towers and their scale factors}

The tower of a localized field may
sit inside a deeper separated sequence.
At an element write \(l\) for the
ambient scale and \(\xi'\le l\)
for the tower scale. The natural
tower curvature and current
readouts are multiplied by the
same squared scale ratio:
\begin{equation}
 \left(\frac l{\xi'}\right)^2
     \big(\eta^{-2}\xi'^2|\partial U-1|,
                            \xi'^3|J|\big)
       =\big(\eta^{-2}l^2|\partial U-1|,
                            l^2\xi'|J|\big).
 \label{mb:ambient-tower-weights}
\end{equation}
The current factor is \(l^2\xi'\),
not \(l^3\). This is the
two-power improvement asserted
by the tower estimate, not
the three-power dimensional
weight of a fully ambient
current.
Under a unit change by \(t\),
\(\eta,l,\xi'\) multiply by
\(t\), \(A\) by \(t^{-1}\),
and \(J\) by \(t^{-3}\).
Both entries of \eqref{mb:ambient-tower-weights}
are invariant. This
calculation fixes the
weights in a joined
sequence independently
of an index convention.

The ambient estimate is proved in \cref{atc:theorem} for every
fixed pair of contained separated sequences. At a settled element
write $s(x)=l(x)/\xi'(x)$. The near input is cut off on a ball
of radius $l(x)$; its weighted tower norm gains $s(x)^{-1}$,
so the quadratic Cauchy remainder pays the factor $s(x)^2$.
For a far input bond $b$, the actual geometry gives
\begin{equation}\label{mb:settled-distance-input}
 d'(x,x_b)\ge d^\flat(x,x_b),\qquad
 s(x)\le1+d'(x,x_b),\qquad
 \frac{s(x)}{s(x_b)}\le L^{2+2d'(x,x_b)/\rho}.
\end{equation}
The complete row read has bounded diameter in $d'$, and its
raw logarithmic derivative contributes $s(x_b)^{-1}$.
Thus, with $D=d'(x,x_b)$ and the fixed rate reserve,
\begin{equation}\label{mb:settled-distance-gain}
 s(x)^2e^{-m d'(x,y_c)}s(x_b)^{-1}
       \le C(1+D)e^{-(m-2\log L/\rho)D}
       \le Ce^{-2\delta D}.
\end{equation}
This is the far-input estimate; a positive distance gain is
not needed for the near input. The physical current follows
from the differentiated constrained Euler equation, with all
six plaquettes and the actual row-transpose densities retained.
The output weights are precisely \eqref{mb:ambient-tower-weights}.
The real reference, determining tower and ambient sequence remain
fixed in the comparison. Constants are uniform over those
geometries. The smaller rate in \cref{atc:theorem} includes the
scale-ratio and carrier losses. Fixed higher derivatives follow
on smaller analytic domains with their own Cauchy constants.
\paragraph{Outer and inner substitutions}

For an outer substitution
\(\mathbf U\mapsto
e^{i\eta\mathbf A}\mathbf U\),
the finite product formula
gives, in normalized
first-order norms,
\begin{equation}
 e_{0,\rm new}\le
  C(e_{0,\rm old}+\|\mathbf A\|_1
                              +\|\mathbf A\|_1^2),\qquad
 e_{1,\rm new}\le
              e_{1,\rm old}+C\|\mathbf A\|_1 .
 \label{mb:outer-substitution}
\end{equation}
Here \(\|\mathbf A\|_1\)
uses the local length
and derivative factors,
and \(e_0\) includes
all actual tower
curvature readouts.
For a tower entry the
same inequality uses
its applicable admitted
readout estimate: the
settled estimate at a
settled element and
the boundary estimate
otherwise. It does
not follow merely
from the elementary
plaquette formula.
Current and atomic
branch tests have
their own finite
Lipschitz estimates.
If the old tests
use half the target
radii, the inequalities
\(\alpha_0\le\alpha_1\),
\(\alpha_1^2\le C\alpha_0\),
and
\(\|\mathbf A\|_1\le c\alpha_0\)
leave the needed
margin once \(c\)
is chosen from
these actual
constants.

For the inner substitution,
the datum is the actual
window expression
\(\varsigma B+B'\),
\(0\le\varsigma\le1\).
The first term comes
from the averages
of the sum of the
coarse background
and the fluctuation
in the specified
order of multiplication.
The second is its
boundary correction.
It is not an
independent scalar
source.
The entire nonlinear
datum must first
be placed in its
absolute physical
disc. The free
response and its
amplification are
then applied, with
the current computed
from the same
composed field.

The numerical age reduction
can be stated without
suppressing its source
of smallness. Write
\(\delta_n=\alpha_ns_n\)
for the final physical
fluctuation amplitude,
\(\rho_j=\delta_n/\alpha_j\),
\(a=n-j\), and now
\(\sigma=L^{-a}\).
The ordinary slow-radius
estimate is
\(\alpha_n/\alpha_j
\le C\sqrt{1+a}\).
The independent-coordinate
datum radius satisfies
\(\alpha_j/r_{B,j}\le C\)
on the actual near
window class. If
the actual fluctuation
amplitude and window
response obey
\begin{equation}
 \frac{\delta_n}{r_{B,j}}\le C\rho_j,\qquad
 \frac{C_{\rm ax}\delta_n\sigma}{r_{B,j}}
       \le C s_n\sqrt{1+a}\,L^{-a},
              \qquad a=n-j\ge0 ,
 \label{mb:inner-age-margin}
\end{equation}
then the second quantity
is uniformly small
at all ages by
choosing the common
small fluctuation
constant first.
Indeed
\(\sup_{a\ge0}\sqrt{1+a}L^{-a}<\infty\).
Young ages cause
no exception once
that constant is
chosen with the
actual response
and radius factors.

This verifies the central datum-disc
margin. The full datum also needs the separate absolute
ceiling \eqref{pw:absolute-datum-ceiling}. To infer that
the complete inner
substitution belongs
to every running core,
the actual local
window map must put
its background-plus-
fluctuation pair in
the primitive class,
including the entire
outgoing physical
amplitude, its
gradient, the
independent current,
the tower and the
atomic branch tests.
An increment bound
alone does not
control that
outgoing amplitude.
The first-exit
list \eqref{mb:all-clause-path}
and the domain
estimates isolate
those remaining
finite inequalities.

\subsubsection{A real one-step section at a seam}

The finite section used in the descent
can be made explicit. Work for this
paragraph on the level-\(k\) lattice,
so a new block contains \(L^4\)
sites. The old level-\(k\) data
are prescribed outside the next
region and on its exterior bonds;
the newly prescribed data are
on the level-\((k+1)\) rows
touching the next region.
Assume all ordinary and mixed
plaquette tests are at most
\(\epsilon\), and use an embedded
two-block neighborhood for every
new pivot.

For a block outside the next
region, all its internal
level-\(k\) bonds are old
prescribed data. Compute
their real mean transports
and use the gauge
\(h(x)=T_{y,x}\), \(h(y)=1\),
inside that block.
The average covariance fixes
all level-\((k+1)\) endpoint
values. The internal
transport means become
one, exactly.
\eqref{mb:mean-versus-comb}--\eqref{mb:bond-reconstruction}
then put every internal
old bond within
\(C_L\epsilon\) of one.
This intermediate gauge
is permitted to change
the old datum; it
will be undone at
the end.
It is not a residual
gauge of the original
fixed datum.

Complete the unassigned
level-\(k\) field in
this frame as follows.
Inside each new block
put its internal bonds
equal to one. On every
fine bond crossing a
new coarse face put
the given coarse
value \(V(c)\).
Leave all old
prescribed bonds
unchanged. The
no-inward convention
ensures that these
instructions do
not prescribe an
old fixed bond
twice: every
crossing bond
with an endpoint
in the new region
was omitted from
the old determining
set.

The curvature of this trial
completion is at most
\(C_L\epsilon\).
Here is the full local
check. A fine plaquette
crossing no coarse
face is either an
old plaquette or
has value one.
A plaquette crossing
one face has equal
coarse transports
on its two parallel
crossing bonds;
they cancel, up
to internal old
transports already
within \(C_L\epsilon\)
of one.
A plaquette crossing
two faces corresponds
to one coarse corner.
If all of its
coarse edges are
new, its product
is the new coarse
plaquette.
If some are old
exterior edges,
\eqref{mb:bond-reconstruction} replaces
each old crossing
bond by its actual
coarse mean with
error \(C_L\epsilon\).
The resulting product
is exactly the
assigned mixed
test.
These alternatives
exhaust the \(0\),
\(1\) and \(2\)
coordinate faces
that a fine
plaquette can
cross.
All prefactors
are unitary.

For a new coarse row,
its inner internal
transports are one
and its outer
internal transports
are one or within
\(C_L\epsilon\) of
one. Each straight
route crosses its
coarse face once.
The defining closed
logarithmic factors
are therefore
\(O(C_L\epsilon)\).
Its actual average
differs from the
desired \(V(c)\)
by that amount.
There is no
smallness assumption
on \(V(c)-1\).

Correct this discrepancy
on the assigned
crossing pivot.
The finite pivot
chart gives a
uniformly invertible
real derivative
and a fixed
admitted radius
depending on \(L\).
For \(\epsilon\)
below its fixed
threshold it gives
a pivot logarithm
\(t_c\) with
\begin{equation}
 \mathcal M(V^{\rm trial}[t_c])(c)=V(c),\qquad
                  |t_c|\le C_L\epsilon .
 \label{mb:section-pivot-correction}
\end{equation}
The internal transport
paths contain no
crossing pivot,
and a straight
segment of another
coarse row contains
no such pivot:
the segment either
has a different
orientation or
crosses a different
coarse face.
Thus its average
does not depend
on \(t_c\).
All these corrections
can be performed
independently.
No norm cost is
proportional to
the number of
new rows.

A fine plaquette has at
most four corrected
edges. Its unitary
product changes by
at most
\(4\sup_c|t_c|\)
times a fixed
small-chart constant.
The completed
datum therefore
satisfies
\begin{equation}
 \sup|\partial V^\downarrow-1|
       \le C_\downarrow\epsilon,\qquad
 \mathcal M(V^\downarrow)|_{\Gamma_{k+1}}=V_{k+1},
 \qquad
 V^\downarrow|_{\Gamma_k}=V_k .
 \label{mb:real-seam-section}
\end{equation}
Undoing the intermediate
unitary gauge restores
the old prescribed
values and preserves
these test bounds.
Mixed tests for the
completed truncation
are checked by the
same corner argument.
The one-step coarse
constraint is exact.

This is a real finite
section with a fixed
constant \(C_\downarrow(L)\).
Its hypotheses include
the actual ordinary
and mixed tests,
the regular finite
mean branch and
the pivot support
premise. Its output
is a coarse datum,
not a regular fine
critical field.
In a refinement
where no seam is
present, the face
lift is already
exact before the
pivot corrections.
To use a coefficient
one in the
two-radius graph
\eqref{mb:test-dependency}, one records
only the test
coordinates on
which that exact
transfer preserves
the original test;
the remaining
corner and pivot
coordinates belong
to the spoiled
collar footprint.
Their future
determining
neighborhoods
must satisfy the
disjointness
certificate of
that graph.

The finite construction does
not prove a smoothed
trial bound for the
bare current at
arbitrary depth.
It supplies the
one-step datum
completion used
before the lower
background is
evaluated. The
regular fine
trial then comes
from that already
constructed lower
background, as
in \eqref{mb:vertical-decomposition}--\eqref{mb:induction-radius}.
This avoids
mistaking the
face-supported
datum for the
desired continuum
scale potential.

\subsubsection{Descent of the datum and ascent of the critical field}
\label{mb:datum-descent-critical-ascent}

The induction on the number of levels need not preserve an exact
one-step average during the descent of the datum.  It must preserve
the older rows exactly and leave a uniformly small error on the new
rows.  The coupled critical equation removes this error during the
ascent.  The following argument keeps the error from acquiring a
factor at each level.

We use the no-inward determining sets of
\eqref{mb:determining-rows}.  In addition to the geometry already
specified, fix powers \(\mathfrak M,R_1,\rho\) of \(L\) with
\[
 L\ge13,\qquad \mathfrak M\ge\mathfrak m^{\mathrm s},\qquad
 R_1\ge L,\qquad R_1\mathfrak M\ge L^2\rho,
 \qquad \rho\ge\mathfrak r^{\mathrm s}.
\]
Every level-\(k\) coordinate circle has at least \(6\rho\)
sites.  The gap in \eqref{mb:separation} is greater than
\(R_1\mathfrak M\xi_j\).  The fixed floors
\(\mathfrak m^{\mathrm s},\mathfrak r^{\mathrm s}\)
depend on \(L\).  Empty regions and regions equal to the whole
torus are allowed.  These are quantitative restrictions on the
geometry; separation with an unspecified constant is not a
substitute for them.

Let \(\mathbf B^{(n)}\) be the determining set of the truncated
sequence \((\Omega_0,\ldots,\Omega_n)\), with all its inward
rows removed. Its last member contains the level-\(n\) bonds
inner to \(\Omega_0\) that touch \(\Omega_n\). Its earlier
members agree with the
corresponding retained members of \(\mathbf B\).

\paragraph{The two radii for a datum.}
Fix the constant \(C_\downarrow\ge1\) of the finite seam
construction.  A datum belongs to \(\mathcal D(\varepsilon)\)
if all its ordinary and mixed test holonomies are defined and
within \(C_\downarrow\varepsilon\) of one, and the tested plaquettes of
levels \(j-1,j\) within \(6\xi_j\) of
\(\partial\Omega_j\) are within \(\varepsilon\) of one.
All distances in this paragraph use level-\(k\) units.  For a
determining set retaining inward rows, one must also impose the
\(\varepsilon\) bound on each pinch holonomy.  Such holonomies
compare the two retained transports ending at the same seam node;
they are additional compatibility tests.  They are absent for
the determining sets used here.

For \(C_\downarrow\varepsilon\le\varepsilon_\downarrow\)
the finite seam construction gives data \(V^{(n)}\) on
\(\mathbf B^{(n)}\), with \(V^{(k)}=V\), such that
\begin{align}
 V^{(n)}_j&=V_j &&\text{on the original rows of level }j\le n,
                                                        \label{mb:descent-old-rows}\\
 |\mathcal M(V^{(n-1)}_{n-1})(c)(V^{(n)}_n(c))^{-1}-1|
      &\le C_\downarrow\varepsilon
       &&(c\in\Gamma_n^{(n)}).                         \label{mb:descent-defect}
\end{align}
The descended datum satisfies \(V^{(n)}\in\mathcal D(\varepsilon)\).

Here is the reason for using two radii.  In a last-level cell the
unassigned fine bonds are filled by the face lift: internal bonds
are one, and bonds crossing a coarse face carry its coarse value,
in the local transport frame.  A fine plaquette deep inside the
new region has value one or the corresponding coarse plaquette.
It therefore inherits the larger bound without another factor.
Near the seam the comparison with the old prescribed bonds costs
\(C_\downarrow\).  All tests used in that comparison lie within
three coarse spacings of the row or plaquette under consideration.
They are within six spacings of the current boundary and hence
obey the smaller bound.  Separation keeps these tests outside the
region altered in the preceding descent step.  Older tests remain
unchanged.  This proves the assertions by downward induction;
the bound is \(C_\downarrow\varepsilon\), not
\(C_\downarrow^{k-n}\varepsilon\).

This descent only asserts \eqref{mb:descent-defect} at a general
seam row. Its average is exact when both row endpoints
and their sup-norm neighbor shells of one level-\(n\)
spacing lie in the new region. The independent pivot corrections of
\eqref{mb:section-pivot-correction} produce a different finite
section under their stated support premise.  The present
induction uses neither equality of these sections nor exact
seam averages during descent.

\paragraph{The one-level analytic input.}
Write \(\mathfrak U'_n(r)\) for the regularity class with
curvature and full-bond current at radius \(r\). In bare
units (\(\eta=1\)), if \(j_n(e)\) is the largest \(j\le n\) for which
\(e\) has a corner in \(\Omega_j\), its tests are
\[
 |\partial U(p)-1|<rL^{-2j_n(p)},\qquad
 |\eta D_U^{\mathsf T}\partial U(b)|<rL^{-3j_n(b)}.
\]
The second expression is the covariant plaquette divergence on
full bonds; the truncated boundary action has its separate
boundary current.  This convention is needed in applying the
local regularity estimate.

The analytic step needed in this induction has the following
precise form. Let \(\varepsilon'>0\) be small enough that
\(L^3B_3^{\mathrm s}\varepsilon'\le c_{10}\) and
\(K_9L^3B_3^{\mathrm s}\varepsilon'\le a_5(\rho)\), where
\(K_9\) is the provisional regularity constant below.
Suppose the new datum has all test holonomies defined
and strictly within \(\varepsilon'\) of one, and a regular
field \(U_0\in\mathfrak U'_{n-1}(B_3^{\mathrm s}\varepsilon')\)
has the prescribed averages on all older rows and differs from
the desired level-\(n\) average by at most \(\varepsilon'\).
The coupled equation, followed by regularity improvement, must
produce a critical field for the new datum in
\(\mathfrak U'_n(B_3^{\mathrm s}\varepsilon')\).
The same constants and smallness thresholds must apply at every
truncation.  This is an assertion about the actual row equations,
their node factors and the Wilson current.

The quantitative use of the coupled equation is as follows.
In level-\(n\) units put
\(\alpha_0=L^3B_3^{\mathrm s}\varepsilon'\).
The old field is in the new regularity class at radius
\(\alpha_0\).  The datum displacement vanishes at every older
row, and at the new rows it is
\[
 B_n(c)=\frac1i\log\{V_n(c)(\mathcal M^nU_0(c))^{-1}\},
 \qquad \sup|B_n|\le\frac\pi2\varepsilon'\le\alpha_0.
\]
The boundary derivative of the inner Wilson action is bounded
in the force norm by \(6\alpha_0\): at a full bond this is
the current estimate, and at a boundary bond it follows by
summing its at most six incident inner plaquettes.  Such a bond
is on the bare layer, so the three required powers of the local
spacing are exactly the bare powers.

%
Work on the original bare region of the new truncation.
Write \(q(A)\) for the straightened rows and \(b^\natural\)
for their prescribed values. The outward factor in
\eqref{mb:outward-node-identity} is determined by the
already prescribed interior rows of the next finer block.
Those displacements are zero in this ascent, so the factor
is one and \(\|b^\natural\|_\infty\le\alpha_0\).
Write \(K_A=q'(A)\) and \(F\) for the inner Wilson action
in the logarithmic chart based at \(U_0\). The equations are
\[
 F'(A)+DRD^\sharp A=K_A^\sharp\omega,\qquad q(A)=b^\natural .
\]
In the isometric row frames of the regional inverse, write
\(T^{-1}q(A)=Q^\natural A+C(A)\) and
\(\widehat b=T^{-1}b^\natural\), using dimensionless row
coordinates. Then
\[
 F'(A)+DRD^\sharp A=F'(0)+\Delta_R A+\mathcal V(A),
 \qquad \Delta_R=D^\sharp D+DRD^\sharp+\Delta'.
\]
Here \(D\) in the curl term is covariant at \(U_0\), and
\(\Delta'\) is the curvature correction to the covariant
curl Hessian of the inner Wilson action. Let \(H_0\) be
the energetic row lift for the regional inverse and
\(\widetilde G_0\) its constrained inverse. Thus
\(Q^\natural H_0=1\), \(Q^\natural\widetilde G_0=0\), and
\(\ker\widetilde G_0=\operatorname{ran}Q^{\natural\sharp}\).
With \(\mathcal G(A)=F'(A)+DRD^\sharp A\), the field equation is
\[
 \begin{split}
 w(A)&=(1+H_0^\sharp C'(A)^\sharp)^{-1}H_0^\sharp\mathcal G(A),\\
 \Psi(A)&=\widetilde G_0[-F'(0)-\mathcal V(A)+C'(A)^\sharp w(A)]
             +H_0[\widehat b-C(A)] .
 \end{split}
\]
Applying \(Q^\natural\) to \(A=\Psi(A)\) gives the exact row
equation. The identity \(\widetilde G_0\Delta_R=1-H_0Q^\natural\) puts the
residual of the linear field equation in
\(\ker\widetilde G_0\). Applying \(H_0^\sharp\) identifies
its multiplier as \(w(A)\), giving \(\omega=Tw(A)\).

The multiplier estimate uses only first differences of \(A\), by
the identity
\[
 H_0^\sharp\mathcal G(A)=
 H_0^\sharp[F'(0)+\Delta' A+\mathcal V(A)]
 +(DH_0)^\sharp DA+(D^\sharp H_0)^\sharp RD^\sharp A .
\]
The weighted column estimates for the lifts bound this by
\(K_8(\mathsf j+\|A\|_X+|\mathcal V(A)|_{(-3)})\), where
\(\mathsf j=|F'(0)|_{(-3)}\). In the difference of the
expressions at two fields, \(F'(0)\) cancels.
Let \(K_1,K_7\) be the bounds of \(H_0,\widetilde G_0\)
in the row-to-field and force-to-field norms, and set
\(K_8=K_1(19+12K_1^3)\), \(K_5=2(K_7+K_1)\) and
\(\Theta=4K_8\). The bound on the Neumann perturbation
is \(4K_1L^3\epsilon^*\le1/8\), where
\(\|C'(A)^\sharp\|_{\infty\to(-3)}\le4L^3\epsilon^*\).
Use \(\operatorname{Lip}(\mathcal V)\le4C_{4,\mathrm{svar}}L^4r\),
\(\operatorname{Lip}(C'{}^\sharp)\le2L^4 C_3^\natural\)
and \(\operatorname{Lip}(C)\le L\epsilon\).
Here \(\epsilon,\epsilon^*\) are the row and transpose error
bounds at \((\alpha_0,r_\circ)\), both \(O(\alpha_0+r_\circ)\).
The inverse difference identity gives
\(\|w(A)\|_\infty\le(17/7)K_8r\) and
\(\operatorname{Lip}(w)\le3K_8\) for
\(\mathsf j\le r\le r_\circ\).
In the difference of \(C'(A)^\sharp w(A)\) use the small
bound on \(C'(A)^\sharp\) with the difference of \(w\),
and the small bound on \(w\) with the difference of \(C'\).
The resulting contraction constant is at most
\[
 K_7[4C_{4,\mathrm{svar}}L^4r_\circ
        +6K_8L^4C_3^\natural r_\circ+12K_8L^3\epsilon^*]
      +K_1L\epsilon .
\]
Choose the fixed chart thresholds as in the regional
contraction estimates, so this is at most \(1/2\).
They also ensure that
\(\alpha_0+3L^2r_\circ<L^{-2}c_{2,\mathrm{sobj}}\),
the curvature ceiling for the physical stationarity argument.
At zero the inverse bounds give
\(\|\Psi(0)\|_X\le(K_5/2)(\mathsf j+\|b^\natural\|_\infty)\).
The threshold \(c_{10}\) keeps
\(r=K_5(\mathsf j+\|b^\natural\|_\infty)\le7K_5\alpha_0\)
within the fixed contraction radius \(r_\circ\).
Both closed balls are mapped into themselves, so their
unique fixed points coincide.

For a real infinitesimal gauge \(\lambda\), write
\(\delta_\lambda A\) for its bond tangent and
\(\varrho_A\lambda\) for the derivative of the normalized
pillar-top values. After the site row conjugation the
linear part of the site pair is
\((-R\Delta_{U_0}\lambda,Q'\lambda)\). Differentiating
the real gauge change and the node factors gives errors
bounded by \(C(\alpha_0+\|A\|_X)\|\lambda\|_Y\) in the
weighted site force and top-row norms. Compose with the
covariant site inverse \eqref{mb:site-pair-inverse} with
the first sign reversed. Its uniform site bound gives
identity plus an error of norm less than \(1/2\) on the
fixed chart. This proves that
\(\lambda\mapsto(RD^\sharp\delta_\lambda A,\varrho_A\lambda)\)
is onto. When \(\varrho_A\lambda=0\), the derivative \(K_A^\flat\) of
the normalized reduced rows vanishes on \(\delta_\lambda A\)
by covariance. The node relation changes each outward row
by a factor depending only on the already retained interior
rows of the next finer block, giving a triangular change
of row coordinates with invertible diagonal. Hence the
kernel of \(K_A^\flat\) is \(\ker K_A\). Gauge invariance
also gives \(\langle F'(A),\delta_\lambda A\rangle=0\).
Pairing the field equation with these tangents gives
\(\langle RD^\sharp A,RD^\sharp\delta_\lambda A\rangle=0\).
With \(\varrho_A\lambda=0\) the first entry of the site
pair can be any vector in \(\operatorname{ran}R\).
Taking it to be \(RD^\sharp A\) gives \(RD^\sharp A=0\).
Hence \(F'(A)=K_A^\sharp\omega\).
%
Differentiating the physical rows adds the infinitesimal
top-site rotations to the normalized derivative. The
derivative of the normalized top map on gauge tangents is
onto by the site pair above. Consequently a physical
constraint tangent can be made a normalized reduced
constraint tangent by adding a gauge tangent. Both parts
are annihilated by \(F'(A)\), giving stationarity on the
physical constraint fibre.
In the solved field equation the second difference is now
\(K_A^\sharp\omega-F'(0)-\Delta'A-\mathcal V(A)\).
The regional row and force estimates bound this by
\(4(1+L^3)\Theta r\) in the weighted force norm. This gives
\[
 \|A\|_{X,\mathbf B}\le7K_5\alpha_0,
 \qquad
 |D^{\mathsf T}DA|_{(-3)}
       \le28(1+L^3)\Theta K_5\alpha_0.
\]
Restore the node values by taking the gauge equal to
\(v_j[A]^{-1}\) at tops of level \(j\). On interior rows
\eqref{mb:node-diagram} gives the prescribed averages.
On an outward row the gauge at its outer end is
\(v_{j-1}[A]^{-1}\). The one-step node factor in
\eqref{mb:outward-node-identity} cancels against the same
factor of the straightened datum, giving the exact physical
datum for both orientations. Gauge covariance carries
constrained stationarity to the restored field.

%
Set \(\alpha_2=7L^2\max\{1,4(1+L^3)\Theta\}K_5\alpha_0\).
The threshold \(c_{10}\le c_8,c_9\) also ensures
\(\alpha_2\le\rho_{1,\mathrm{svar}}\),
\(C_{4,\mathrm{svar}}\alpha_2\le1\),
\(4\alpha_2+C_{4,\mathrm{svar}}(\alpha_0+\alpha_2)\le1/2\),
and \(2\alpha_0+3\alpha_2\le1/3\).
The curvature and current expansions with the displayed
first and second difference bounds give provisional
regularity radius \(2(2\alpha_0+3\alpha_2)=K_9\alpha_0\).
The same bounds hold after node restoration by covariance.
By the datum threshold \(K_9\alpha_0\le a_5(\rho)\). To
reduce this provisional radius use the improvement in the
class with full matrix divergence. On a critical field
with input radius \(e_0\le a_5(\rho)\) and datum bound
\(\varepsilon'\le e_0\), its output radius is
\[
 \max\{6\varepsilon',\,3K^\sharp\varepsilon'+3\vartheta e_0\},
 \qquad
 \vartheta=C_b\rho^2e^{-\delta_*\rho/24}+C_c\rho^4e_0 .
\]
The estimate at positive index uses the localized critical
equation. At index zero the plaquette is prescribed, and
the full-bond divergence is bounded by the six incident
plaquettes of index at most one. Choose the separation
floor and \(a_5(\rho)\) so the two terms of \(\vartheta\)
are at most \(1/24\) each. Since
\(B_{\mathrm I}\ge\max\{6,L^2,12K^\sharp\}\), the output
radius is at most \(\max\{B_{\mathrm I}\varepsilon',e_0/2\}\).
Iterating until the first term dominates gives radius
\(B_{\mathrm I}\varepsilon'\le B_3^{\mathrm s}\varepsilon'\).
This last return, including its current estimate, is essential;
the factor \(L^3\) cannot be iterated as a loss of regularity.

\begin{proposition}[The finite-depth induction]
\label{mb:finite-depth-orbit-induction}
Assume the one-level analytic input just stated, uniformly in
the admitted geometry, and the regularity improvement and
regular-orbit comparison estimates used in that input.  There
are radii \(a_0^{\mathrm s}>0\) and
\(a_1^{\mathrm s}(\rho)>0\), independent of the depth and
torus, with
\(B_3^{\mathrm s}a_1^{\mathrm s}(\rho)\le a_0^{\mathrm s}\),
such that a real datum whose test holonomies are all within
\(0<\varepsilon_1\le a_1^{\mathrm s}(\rho)\) one has a
critical configuration in
\(\mathfrak U'_k(B_3^{\mathrm s}\varepsilon_1)\).
Any two critical configurations for that datum in
\(\mathfrak U'_k(\varepsilon_0)\), with
\(\varepsilon_0\le a_0^{\mathrm s}\), belong to the same
residual gauge orbit.
\end{proposition}

\begin{proof}
Take \(\varepsilon'>C_\downarrow\varepsilon_1\) in the
uniform smallness range, with \(\varepsilon'\le2C_\downarrow\varepsilon_1\),
and first
descend the datum as in \eqref{mb:descent-old-rows}--
\eqref{mb:descent-defect}.  At depth zero the datum itself
specifies every inner bare bond.  Its curvature is bounded by
\(\varepsilon'\); its covariant plaquette divergence is bounded
by \(6\varepsilon'\).  Since \(B_3^{\mathrm s}\ge6\), this
starts the upward induction.

Suppose \(U^{(n-1)}\) has been constructed.  On every double
block read by a level-\(n\) row, all level-\((n-1)\) bonds
belong to the last member of the preceding truncation.  Exact
locality of the one-step average therefore gives
\[
 \mathcal M^nU^{(n-1)}(c)
    =\mathcal M(V^{(n-1)}_{n-1})(c).
\]
Together with \eqref{mb:descent-defect}, this verifies the new
datum error required by the analytic step.  That step produces
\(U^{(n)}\in\mathfrak U'_n(B_3^{\mathrm s}\varepsilon')\).
The older constraints are exact by
\eqref{mb:descent-old-rows}.  The same bound consequently holds
at every depth.  At depth \(k\) the radius is at most
\(B_3^{\mathrm s}\varepsilon'\). Choose \(a_1^{\mathrm s}(\rho)\)
small enough that \(2B_3^{\mathrm s}C_\downarrow a_1^{\mathrm s}(\rho)\)
is below the initial radius allowed by the improvement estimate
on the full sequence. Applying that estimate at the original
datum radius \(\varepsilon_1\) gives \(B_3^{\mathrm s}\varepsilon_1\).

For uniqueness, take one regular critical field as reference.
The regular-orbit comparison places any second one in its
Landau chart, with zero reduced datum and field norm at most
\(\mathsf B_1\varepsilon_0\).  Choose
\(a_0^{\mathrm s}\le r_\circ/\mathsf B_1\) and
\(6K_5a_0^{\mathrm s}\le r_\circ\), in addition to the
curvature ceilings of that chart.  Both the zero field and the
second field then solve the same coupled equation in its
uniqueness ball.  The field component of the second solution
is zero.  Its gauge fixes the tops, and the common prescribed
datum fixes the retained seam nodes; hence it is a residual
gauge.  This proves orbit uniqueness.
\end{proof}

For the actual determining sets \eqref{mb:determining-rows},
\cref{bgc:proper-background} verifies these analytic premises using
\cref{prop:psr-strong-inverse,prop:rtc-nonstationary-trial} and the
comparison and reset estimates
\cref{bgc:strong-comparison,bgc:global-comparison,bgc:reset}.
The proof retains the complete outward blocks and the inner-action
boundary force. Its uniqueness conclusion concerns all regular
critical configurations in the stated radius. It does not identify
an unrestricted global minimum or assert an infinite-volume or
continuum limit.

% Theorem 49.1, and Theorem 49.3.
\paragraph{Recovering the regularity radius from the datum.}
\hypertarget{hi:real-datum-profile}{}
In the separated geometry choose a power \(\sigma\) of \(L\)
satisfying \(\sigma\ge L^2\rho_{\mathrm c}\), with
\(\mathfrak M\ge\mathfrak M_{\mathrm c}\) and
\(\rho\ge L^2\sigma\), and with at least \(6\rho\) sites
in each direction at the top level. Let \(U\) be
critical and already regular at a radius \(a_{\mathrm c}\).
The fixed floors \(\mathfrak M_{\mathrm c},\rho_{\mathrm c}\)
come from the local flat inverse construction. At a bare
site \(x\) let \(r_U(x)\) be the maximum of
\(\eta^{-2}\mathsf L(p)^2|\partial U(p)-1|\) and
\(\mathsf L(b)^3|J(U)(b)|\) over the plaquettes and full
bonds based at \(x\). Here \(\mathsf L\) is the local layer
spacing and the current in this test is formed at full bonds.
For \(R=20L\sigma\) in the weighted distance, let \(r_R(x)\)
be the supremum of \(r_U\) over the ball and \(v_R(x)\) the
largest defect in the datum tests meeting it, including
the pinch tests if inward rows are kept. The local estimate at a regular critical
field has the form
\[
       r_U(x)\le K_{\mathrm{loc}}v_R(x)+\tfrac12 r_R(x).
\]
Here are the local estimates behind this inequality.
If \(x\) is in the bare layer away from the next layer,
the bonds in the relevant plaquettes are datum bonds.
The current there is bounded by the sum of at most six
such plaquette defects. Otherwise choose level \(k\)
to be the layer of \(x\), or the next layer when \(x\)
is within \(5\sigma\) in the units of its own layer
of that next member. In the units of level \(k\), the
enlarged cube reads only layers \(k-1,k\).
Write \(Q_s\) for the \(s\)-neighborhood of the
partition cube of side \(\sigma\) containing \(x\).
For the auxiliary sequence take the full torus at
level zero and \(Q_{s_i}\) for \(1\le i<k\), where
\[
 s_k=\sigma-\rho_{\mathrm c},\qquad
 s_i=s_k+\sum_{i\le l<k}
             (\rho_{\mathrm c}+\mathfrak M_{\mathrm c})L^{l-k}.
\]
The member at level \(k\) is \(Q_{s_k}\cap\Omega_k\).
Keep an inward row of this auxiliary sequence only
when the original determining set keeps it. Its
positive-level members are subordinate to the original
ones, and the largest is inside \(Q_{\sigma-2}\).
If \(v_R(x)\) is above the fixed ceiling of the
datum transport gauge, the local inequality follows
from \(r_U(x)\le a_{\mathrm c}\) by increasing
\(K_{\mathrm{loc}}\). Use that gauge below its ceiling.

For the real local gauge on \(Q_\sigma\), take the
enlargement of side \(3\sigma+4L\mathfrak r_{\mathrm p}\),
where \(\mathfrak r_{\mathrm p}\) is the fixed cube floor.
It lies in \(Q_{3\sigma}\) and lifts to a single cube.
The two local scales give curvature at most
\(L^2 r_R(x)\eta^2\), and the current comparison gives
covariant divergence at most \(2L^3r_R(x)\eta^2\).
The local regular-gauge estimate with these two inputs
bounds the potential, its first difference and its flat
second-order reads on \(Q_\sigma\) by \(C\sigma r_R(x)\)
in level-\(k\) units. Its smallness condition
\(C\sigma r_R(x)\le c\) is included in the radius choice. Extend this potential by zero on the other bonds
of the auxiliary torus. A constrained variation there
is fixed on the bare determining bonds outside the
largest positive-level member. Patch it to the original
field across \(Q_\sigma\). The patched curve fixes each
old row. A row whose two endpoints are below its
reading level in the auxiliary sequence is fixed by
induction in the averaging equation. A row with an
endpoint of exactly its reading level is an auxiliary
constraint, and its two-block read is inside \(Q_\sigma\).
The old inward rows meeting the varied set, when
present, are inherited auxiliary constraints at the
interface of levels \(k-1,k\). Every varied plaquette
lies within \(Q_\sigma\), so the two actions have the
same derivative along the curves. This proves
criticality for the extension at its own auxiliary
datum. The cutoff region retains the actual current of
the extension. Its local length there is the bare
spacing, so the zero extension costs \(C\sigma r_R(x)\)
in the weighted current norm. Passing to the flat
Landau slice with the prescribed node factors gives a
datum bound \(b_{\rm f}\le C\sigma^2r_R(x)\) and the
flat response bound \(r_{\rm f}\le C\sigma^2r_R(x)\).

In the inner cube the datum in that slice is smaller.
The transport gauge of the datum writes each nearby row
in terms of the tested plaquettes and, at a retained
seam, the tested holonomies. A row based at \(y\) has
size at most \(C(1+|y-x|)v_R(x)\) in the chosen scale.
At the outer endpoint of a coarse row the extra
straightening factor is one, since the prescribed
datum gauge makes the corresponding lower contour one.
At an edge row the two factors at the common seam node
cancel, leaving the tested seam transport.

The flat response uses the full flat linear package of
\cref{lem:mb-retained-edge-Schur} when inward edge rows are
retained. It bounds the weighted potential and its first
difference at \(|y-x|\le\sigma/8\) by
\[
 C(1+|y-x|)v_R(x)+Ce^{-c\sigma}b_{\rm f}+Cr_{\rm f}^2 .
\]
This is the near and far split of the datum sum. The
linear distance factor is summable against the kernel,
and the far rows are at distance at least \(\sigma/8-2\).
In the multiplier equation split the differentiated
potential sum at the same inner cube. For rows within
\(4L+2\) of \(x\) this gives
\[
 |\omega(r)|\le Cv_R(x)+Ce^{-c\sigma}(b_{\rm f}+r_{\rm f})
                         +Cr_{\rm f}^2 .
\]
The actual row differential in \(F'(A)=K_A^{\natural *}\omega\)
reads only these nearby multipliers at the full bond.
The plaquette estimate uses the potential bound on the
neighboring tops. The two readouts agree with those
of the original field up to gauge, and the old local
length is at most \(L\) times the auxiliary length at
these reads. The resulting coefficient of \(r_R(x)\)
is at most
\(C\sigma^2 e^{-c\sigma}+C\sigma^4a_{\mathrm c}\).
Choose \(\sigma\) first and then \(a_{\mathrm c}\) to
make this at most \(1/2\), with the ceilings of the
local gauge and flat inverse also imposed at that choice.

The estimate implies regularity improvement with no
accumulation in the number of layers. Let
\[
 P(x)=\sup_\tau e^{-\delta d_{\mathbf B}(x,\tau)}
                   |\partial V(\tau)-\mathbf 1|,
       \qquad e^{\delta R}=\sqrt 2 ,
\]
where \(\tau\) ranges over the datum tests. The triangle
inequality gives \(v_R(x)\le\sqrt2 P(x)\) and
\(P(y)\le\sqrt2 P(x)\) for \(y\) in the ball.
For \(a>0\), maximize \(r_U(x)/(P(x)+a)\) over the bare sites.
If \(S\) is the maximum, the local inequality at a maximizing
site gives \(S\le\sqrt2 K_{\mathrm{loc}}+S/\sqrt2\).
Letting \(a\downarrow0\) yields
\(r_U(x)\le5K_{\mathrm{loc}}P(x)\).
Moving from the base point of a bond or plaquette to a
corner costs at most two bare steps in the weighted distance,
so the same estimate holds there with an enlarged fixed
constant. In particular, when every datum test is bounded
by \(\varepsilon'\), the regularity radius is at most
\(C\varepsilon'\). The flat chart admission in this calculation
is verified in \cref{bgc:reset}; the provisional one-level output
lies in its permitted radius by \eqref{bgc:radius-order}.

\subsubsection{Norms and the remaining finite estimates}

The natural weights in the nonlinear
estimates can be stated explicitly.
One field norm and force norm are
\begin{equation}
 \|A\|_{X_1}
   =\max\{\|\mathsf L A\|_\infty,
                    \|\mathsf L^2\nabla_{U_0}A\|_\infty\},
 \qquad
 \|f\|_{\mathcal F}
                    =\|\mathsf L^3f\|_\infty .
 \label{mb:natural-field-force}
\end{equation}
The force is the physical density
in the fine pairing
\(\eta^4\sum_b(f_b,A_b)\).
The datum \(b\) is a raw
row logarithm, with norm
\(\|b\|_Y=\sup_c|b_c|\).
Use the counting pairing
\(\sum_c(\omega_c,b_c)\)
for its multiplier and
\(\|\omega\|_Z=\sup_c|\omega_c|\).
For differentiated conclusions
append the separate natural
second-order and current
readouts; do not replace
them by the supremum of
every second difference.

This convention explains a
normalization in the named
propagator formulas. The
per-length row field is
\(B_Q(c)=\xi_c^{-1}b_c\)
with cell pairing
\(\sum_c\xi_c^4(\omega_Q(c),B_Q(c))\).
Equality with the raw
pairing requires
\begin{equation}
 \omega_Q(c)=\xi_c^{-3}\omega_c,\qquad
       \|\omega\|_Z=\|\xi_c^3\omega_Q\|_\infty .
 \label{mb:multiplier-normalization}
\end{equation}
It is the bilinear weighted
transpose of \(b=\mathsf mB_Q\).
Similarly the per-length
penalty
\(\sum_c\xi_c^4a_0\xi_c^{-2}|B_Q(c)|^2\)
equals \(\sum_ca_0|b_c|^2\).
Thus its raw mass is
\(a_0\), while its
per-length mass is
\(a_0\xi_c^{-2}\).
The generic identity in \eqref{mb:soft-readout} uses one
consistent pairing;
\eqref{mb:named-gauge-operators}--\eqref{mb:next-step-form}
use the per-length
convention. Converting
their multipliers and
quadratic readouts
uses \eqref{mb:multiplier-normalization}.

The requirements for a general background construction can be
organized as the local inequalities below. For the proper
determining sets in force, the later construction supplies the
estimates needed for the physical problem on its stated domains.
The strong row saddle is proved in
\cref{prop:psr-strong-inverse}; the actual Wilson and constraint
remainders are bounded in
\eqref{pnc:Wilson-residual-bound} and
\eqref{pnc:constraint-residual-bounds}. The site inverse and
finite gauge construction are given by
\eqref{pnc:site-block-inverse}--\eqref{pnc:actual-mean-slice}.
These estimates use the physical amplitudes
\eqref{pnc:physical-amplitudes}, retain every outward factor,
and apply about real references with the required actual
curvature and current bounds. Their uniform holomorphic
continuation and weighted response are proved in
\cref{prop:pnc-critical-continuation} and
\eqref{pnc:weighted-first-response}.
\Cref{prop:rtc-nonstationary-trial,bgc:proper-background}
then supply the regular-trial correction and the real
background induction on those determining sets.

The canonical independent-current completion is fixed separately
in \cref{prop:acc-current-identification}, with its strong bounds
in \cref{prop:acc-strong-propagators}. These results apply to that
specified operator family. For any other primitive package,
the list and conditional assembly below retain its own domain,
norms and complete projected comparison
\eqref{mb:complete-normal-form}; no identification follows
merely from using the same operator names.

\emph{The site estimate.}
For the actual \(Q'\), including its outward endpoint
convention, the real site correction uses the weighted inverse
on the admitted separated sequence. On the primitive complex
ball one needs the corresponding weighted kernel bounds for
\eqref{mb:site-pair-inverse} and its first differences.
Prove the nonlinear gauge
and node remainders there
with fixed admission
margins. This is the
quantitative version
of the completely
proved projector
algebra in \eqref{mb:site-range}--\eqref{mb:site-pair-inverse}.

\emph{The constrained inverse.}
For the actual Lagrangian
and all actual row and
gauge terms, prove
\begin{equation}
 \|\mathcal S_0^{-1}\|_{\mathcal F\oplus Y
                         \,\longrightarrow\,X_1\oplus Z}
        \le C,\qquad
 \sup_y\sum_c e^{m d_{\mathbf B}(y,y_c)}
            \|(\mathcal S_0^{-1})_{yc}\|_{\rm norm}\le C,
 \label{mb:inverse-frontier}
\end{equation}
and the corresponding column
and readout estimates.
The normalized entry includes
the physical cell weights
of its source species.
This is required first on
the primitive real class,
then by a whole-ball
complex perturbation.
The complete projected
current and gauge
comparison in \eqref{mb:complete-normal-form},
with every term of
\eqref{mb:named-gauge-operators}--\eqref{mb:next-step-form},
is a separate
identification task.

\emph{The local coefficient estimate.}
For the actual fine Wilson
action, actual straightened
constraint and actual datum
factors, prove \eqref{mb:coefficient-list} in
these norms, with
\begin{equation}
 c_F(r)+c_q(r)+c_{q\omega}(r)
                  +c_{\rm mix}(r)\le C(r+\gamma_0+\epsilon_0)
 \label{mb:coefficient-frontier}
\end{equation}
in the chosen linearization,
or the corresponding
small triangular matrix
just treated. The finite
coordinate remainders
must be summed using
their actual supports.
For example bare
second differences
created by a group
product have to be
paired with their
actual powers of
\(\eta\), or summed
by parts before a
kernel is applied.
Counting all fine
variables of a
pillar as independent
losses is not a
proof of \eqref{mb:coefficient-frontier}.

\emph{The slice improvement.}
The current comparison and
\eqref{mb:slice-absorption} give the reset in the real slice
under its fixed potential ceiling. The descent and ascent
also require the preliminary contour
alignment at seams
with its single
defect cost.
The exact rigidity
and summability
\eqref{mb:rigidity-defect}--\eqref{mb:rigidity-sum}
already supply
the group-valued
part.

\emph{The critical patch estimate.}
The auxiliary sequence and the flat row readouts give
\eqref{mb:patch-forcing-list} and the pointwise profile on the
admitted real geometry. The real cube gauge supplies the
potential and second-order bounds used on that geometry. The
maximum argument transfers the local inequality to the profile.

\emph{The continuation and tower estimate.}
\Cref{htc:common-ball,htc:tower-bound} construct the common analytic
ball and bound the ratio $M/r_0^2$ in
\eqref{mb:complex-readout-disc}, including the actual representative
gauge and physical current. \Cref{htc:all-towers} supplies both
erosion rules and every truncation. The outer and inner substitutions,
with their full parameter domains and nested tower tests, are proved
in \cref{htc:outer-substitution,htc:inner-substitution}. The geometric
and radius comparisons below use these estimates with their stated
localization collars.

These are finite lists of
specific actual inequalities
and identities with their
uniformities. Their
proofs can share the
same matrix estimates,
and some can be
proved in a different
order. They are
not all extra
assumptions on
every possible
route.
For example a
direct smooth
fine section
can be replaced
by the real
coarse datum
section \eqref{mb:real-seam-section}
and the lower
background induction.
The latter avoids
requiring an
independent fine
trial construction.

\subsubsection{The six conditional background conclusions}

Here is a precise way to assemble
the reductions without promoting
an estimate that has not been
proved. Suppose the preceding
finite estimate list is
established on the declared
primitive domains, with its
actual normal-form matches
and fixed margins.
Suppose also that the
history regions satisfy
their finite separation
certificates, and that
the one-step section
uses the proved pivot
support geometry.
All constants below
may depend on
\(L,\rho,\mathfrak M\)
and the stated local
geometry, but not on
depth, volume, region
shape, datum position
or real holonomy.
Then the following
six conclusions hold.

\paragraph{Regular orbit.}
Every datum in the specified
small regular class has
one regular critical
residual orbit, with
\begin{equation}
 \sup_p\eta^{-2}\mathsf L(p)^2|\partial U(p)-1|
       +\sup_b\mathsf L(b)^3|J_A(U)(b)|
                        \le C\sup_\tau|\partial\mathbf V(\tau)-1|.
 \label{mb:regular-orbit-conclusion}
\end{equation}
It has the unique comb
representative outside
all nodes. Existence
is the datum descent,
the actual induced
action step and
\eqref{mb:criticality-transitivity}, with the
uniform radius reset.
The induction is
on the marked
two-radius data
class of \eqref{mb:test-dependency},
so each recursive
completion encounters
at most one loss
on a test path.
The single global
small base radius
is chosen after
that fixed
\(C_\downarrow\).
Uniqueness is the
slice argument
following \eqref{mb:induction-radius}.
The representative
is the exact
construction in \eqref{mb:comb-axial}--\eqref{mb:axial-uniqueness}.
At trivial complete
datum the orbit
is flat and
is an absolute
minimum of the
nonnegative Wilson
action.
At other data
the conclusion
is regular
criticality.

\paragraph{Local regularity.}
The field has the profile
\eqref{mb:weighted-maximum}--\eqref{mb:layer-profile}.
On a cube of side
\(m\xi_j\) with its
prescribed enlargement
inside \(\Omega_j\)
and shorter than
half the torus,
\(m\epsilon\le c\)
gives a real local
gauge with
\begin{equation}
 \xi_j|A|+\xi_j^2|\nabla A|
       +\xi_j^3|\Delta A|
       +\xi_j^3|\partial^\sharp\partial A|
                       \le C m\epsilon .
 \label{mb:local-regularity-conclusion}
\end{equation}
Its first differences
have the Hölder
bound in \eqref{mb:holder-reduction}.
Here \(\Delta=-\Delta_\eta\)
is the positive
componentwise Laplacian;
the heat generator
\(\Delta_\eta\) in
\eqref{mb:laplacian-comparison} has the
opposite sign.
The full-current radius is required on the enlarged cube.
The single real gauge construction supplies the ordinary
Laplacian in \eqref{mb:local-regularity-conclusion}, and the
covariant divergence comparison supplies its flat curl term.
The constant-layer site estimate then gives the Hölder bound.

\paragraph{Complex continuation.}
For each admitted real
datum center there is
one holomorphic
comb background
on its specified
physical datum
ball, agreeing
with the real
representative.
It has a Landau
pair in \eqref{mb:landau-pair}
with its actual
node normalization,
and a possibly
different comb
pair. The real
conjugation and
admitted gauge
covariance follow
from uniqueness.
The coupled
construction,
site correction
and real/complex
gluing prove
these assertions.
A global all-site
bound on the
comb gauge is
included only
when its actual
readout estimate
is supplied;
a top bound
for the Landau
gauge is not
silently enlarged.

\paragraph{Response.}
On the entire admitted
datum ball, a unit
source \(e\) at
row \(c\) has
the bounds
\begin{equation}
 \begin{split}
 &\xi_y|h|+\xi_y^2|\nabla h|
       +\xi_y^3|D^\sharp Dh|
       +\xi_y^3|\Delta_{U_0}h|
       +\xi_y^3|\delta J_A|\\
 &\hspace{35mm}\le
                  C e^{-m d_{\mathbf B}(y,y_c)}|e| .
 \end{split}
 \label{mb:response-conclusion}
\end{equation}
The multiplier and top
node readouts have
their corresponding
normalized bounds.
The Hölder member
follows from the
actual pointwise
gradient and
Laplacian estimates.
At a noncritical
independent-current
pair the second-order
members are the
combinations \eqref{mb:off-current-readouts};
the critical
slice identity
is not asserted
there.
\eqref{mb:datum-response-system},
the actual kernel
estimates and
the local
readout formulas
prove these
claims.

\paragraph{Hereditary complex domains.}
For sufficiently small radii $(e_0,e_1)$, \cref{htc:all-towers}
gives tower curvature and current bounded by $C(e_0+e_1^2)$
in their declared scales. The common complex ball and its
uniform crude readout bound are proved in \cref{htc:common-ball}.
At settled elements
of an ambient
sequence the
same readouts
use \eqref{mb:ambient-tower-weights}
and have the
ambient exponentially
weighted Lipschitz
bound.
Outer and inner
substitutions
preserve the
specified cores
under their
all-clause
numerical margins.
These follow
from \eqref{mb:complex-linear-curvature}--\eqref{mb:complex-heredity},
\cref{atc:theorem},
and \eqref{mb:all-clause-path},
with the true
radius factors
retained.

\paragraph{Construction, restriction and comparison.}
The certified label,
healed, prepared,
stage, tower and
joined sequences
are separated.
Restriction gives
\eqref{mb:restriction-identity}, and
the next-step
background nests
with its truncation.
For a collared
region the joined
problem and
exact lifted
datum define
its localized
background.
The common
admitted path
then gives
\eqref{mb:boundary-comparison}, with
its potential,
gradient and
second-order
readouts.
The comparison
gauge is prescribed
only at the
stated tops.
Separation closure,
the adaptive
factorization,
constrained criticality,
regular uniqueness
and the weighted
response prove
this conclusion.

\begin{samepage}
The numerical constants in this
theorem must have a solved
order of choice.
One admissible order
fixes the local
geometry and its
uniform coefficient
degrees, the inverse
and response
constants, then
the cutoff scale
\(\sigma\), and
finally the small
regularity and
analytic radii.
If an exponent
or reserve parameter
must be chosen
before later
geometry, the
needed degree
bound must be
uniform in
that geometry,
or the whole
finite feasibility
list must be
solved jointly.
Logarithmic
smallness at
already fixed
geometry does
not prove a
different
choice order.

The first-step chart,
classical flat inverse
and kernels, and
physical first-step
backgrounds provide
proved instances
of these reductions.
The site, layered,
full-current and
window estimates
listed above remain
their precise
multilevel
application tasks.
No claim about
an all-current
completion,
an interacting
Gaussian law
or a history
formation bound
follows from
the conditional
background theorem
without those
separate matches
and estimates.
\end{samepage}


% END MODULE chapters/multilevel_background_reductions.tex

% BEGIN MODULE chapters/04_propagators.tex
\subsection{Constrained propagators and local decay}
\label{candidate:propagators}
\subsubsection{Propagators in a background field: exponential decay produced by the averages}
\label{sec:propagators-propagators-in-a-background-field-exponential-decay-produced-by-the}

Every object produced by a renormalisation step is assembled from inverses of ``Hessian of the action + gauge term + mass made of block averages''. These inverses are taken around a background on a decreasing sequence of regions, in which each layer has a scale of its own. The propagator hypotheses are: they exist; they depend analytically on the background; in each layer they behave like massive Green's functions measured in the unit of that layer; they decay exponentially; and they are sums of terms that are local both in the field and in the regions. A single requirement governs the whole construction. Since the regions are supplied by the large-field analysis, none of the constants may vary with their shape or their number, nor with the level, the cutoff or the torus. Two numbers are exceptions, and neither of them enters any localised bound, rate, weight or sum of weights. (Ba{\l}aban's papers are referred to by the volume of Commun. Math. Phys. in bold. Where two papers must be told apart, the first page is added: \textbf{99}, 75 and \textbf{99}, 389; \textbf{102}, 277; \textbf{122}, 175. The volume and page references identify the published papers. Clauses 0--5 appear below; the following list defines the references to later requirements. The mechanisms in (iii) carry the letters (a)--(k).)

The numbered requirements concern six propagators, their covariances and walks, and parameter decoupling. The vector and current conclusions require the specified operator and domain estimates. \Cref{ch:layered-current-precision} proves its stated site, row, averaging-coefficient and physical-Hessian reductions; those reductions establish only the requirements stated there.

\paragraph{References to the propagator hypotheses}
References to statements 6--8 mean the requirements below.
The labels 1c.1--1c.10 in the strategy name stages of the
conditional argument, as follows.

\begin{description}
\item[Conditioned covariance (statement 6).]
This means the actual fluctuation precision \(\mathsf A\), its
conditioning \(\chi_Y\mathsf A\chi_Y\), and the inverse \(C^{(k)}(Y)\)
defined below, with the coercivity, complex-domain and boundary
requirements of Input A2 and the kernel/locality requirements of
Input A3 in the dependency catalogue. Their walk requirements are
included in statement 5 here. Physical identification uses the
complete current and gauge comparison below and the three Hessian
readouts of \cref{ch:layered-current-precision}, not a Ward
identity alone. The canonical proper covariance, including
its outer constraint curvature and relative contour form, is
constructed in \cref{pgc:proper-covariance,pgc:section}.
\item[Walks as parameter-dependent operator data (statement 7).]
For a fixed cube partition \(\mathsf P\), each supplied walk term
has its specified finite parameter carrier \(\mathsf P(\omega)\).
Write \(s^\omega=\prod_{\Delta\in\mathsf P(\omega)}s(\Delta)\)
and \(\mathsf O(s)=\sum_\omega s^\omega\mathsf O_\omega\).
The required data include the total-incidence bound in mechanism
(i), absolute locally uniform convergence on
\(|s(\Delta)|<e^{\bar\kappa}\), the stated weighted bounds,
reality and covariance, compatibility on common lifts, and
wall-component splitting when separating parameters vanish.
The exceptional convergence-only bound mentioned below is part
of these supplied data. These are the operator-data hypotheses
following \eqref{eq:rg-objects-R9} in \cref{ch:rg-objects-v2}; the substitutions,
their domain condition and their alternating brackets are
\eqref{eq:rg-objects-R10}--\eqref{eq:rg-objects-R12} there. None of this asserts \(QH(s)=1\).
\item[Decoupled response (statement 8 and system (I)).]
The operative system (I) is \eqref{eq:rg-objects-R9} in
\cref{ch:rg-objects-v2}. There \(B'\) and \(\mathsf D\)
are raw rows; each map written \(\mathsf H_{(1)}(s)\) means
the raw-datum version \(\mathsf H_{(1)}(s)\mathsf m^{-1}\)
of the per-length map used here, and \(C_k\) means the raw
nonlinear part \(\mathcal C\) defined below.
Use the following hypotheses there
on the reduced current, holomorphy, bounds, splitting, transport
and actual substitution domains. Their existence and uniqueness
remain response inputs of A3. At \(s=1\), identification with
the physical branch uses the fixed-reference stationarity and
constraint equations in \eqref{eq:nonlinear-response-NR1}, the actual Lagrangian Hessian \eqref{eq:nonlinear-response-NR4},
and the coupled inverse in \eqref{eq:nonlinear-response-NR5}--\eqref{eq:nonlinear-response-NR9} in
\cref{ch:nonlinear-response-reductions}.
\item[The coupled alternative (system in \eqref{eq:propagators-II}).]
This is the explicitly defined operator-data map below. Its
coefficient bounds, radius and current requirements are additional
conditional data; it is not identified with \eqref{eq:nonlinear-response-NR1} at arbitrary
\(s\) or at an independent current. Equality of the two response
constructions is used only on their common domain at \(s=1\)
over the stated critical physical pair, after full operator
matching and the required uniqueness have been established.
\end{description}

\textbf{(i) The setting and the objects.}

\emph{Geometry.} The integer $L$ is odd and large. The cube size of the walk expansions, $M_1$, and the gap number, $R_1$, are both powers of $L$, with $R_1\ge L$. We work in level-$k$ units, in which $\mathbb T^{(j)}$ has spacing $\xi_j:=L^{j-k}$; moreover $\eta=\eta_k=\xi_0$. A sequence $\Omega_0\supset\dots\supset\Omega_k$ is called \emph{separated} when every $\Omega_j$ is a union of $M_1$-cubes at level $j$ and

\begin{gather*}
\FitMath{\operatorname{dist}_\infty(\Omega_j^c,\Omega_{j+1})>R_1M_1\xi_j .}
\end{gather*}

No further condition is placed on the regions. In particular $\Omega_0$ is arbitrary, and Dirichlet conditions are imposed outside it. \emph{Layers} are the sets $\mathcal L_j=\overline{\Omega_j\setminus\Omega_{j+1}}$. The \emph{tops} $y$ are the sites of $\mathbb T^{(j)}$ lying in $\mathcal L_j$. Such a top has local scale $\xi_y=\xi_j$ and block $\Delta(y)$, which is the cube of side $\xi_y$ with centre $y$. The \emph{rows} $c$ are the bonds of $\mathbb T^{(j)}$ that have one end point in $\mathcal L_j$ and the other end point either in $\mathcal L_j$ (the \emph{interior} rows) or, when $j\ge1$, outside $\Omega_j$ (the \emph{outward} rows); for them $\xi_c=\xi_j$. The \emph{inward} bonds, those that enter $\Omega_{j+1}$, carry no row. The distance

\begin{gather*}
\FitMath{d_{\mathbf B}(y,y')=\inf_{\Gamma\ni y,y'}|\Gamma|_{\mathbf B},} \\[3pt]
\FitMath{|\Gamma|_{\mathbf B}=\sum_{j\ge1}\xi_j^{-1}\operatorname{length}(\Gamma\cap\mathcal L_j)+\xi_0^{-1}\operatorname{length}(\Gamma\setminus\Omega_1)} \\[3pt]
\FitMath{(\Gamma\text{ connected})}
\end{gather*}

measures each layer in the unit of that layer. Outside $\Omega_1$, and hence also outside $\Omega_0$, length is counted in the bare unit. Therefore $d_{\mathbf B}$ is independent of $\Omega_0$, and a path that leaves $\Omega_0$ finds no cheaper unit of length out there. This distance differs from the distance (2.46) of \textbf{96}, whose contours are made up, in each layer, of bonds of the lattice of that layer. A larger class of paths is allowed here, and so $d_{\mathbf B}$ is the smaller of the two.

\emph{Regular pairs} $(\mathbf U,\mathbf J)$ consist of $\mathbf U=e^{i\eta A'}U$, with $U$ real, together with an \textbf{independent} $\mathfrak g^c$-valued $1$-form $\mathbf J$ that represents the current. Such a pair is required to satisfy, on $\mathcal L_j$,

\begin{gather*}
\FitMath{|\partial U-\mathbf 1|,|\partial\mathbf U-\mathbf 1|<\alpha_0\eta^2\xi_j^{-2},} \\[3pt]
\FitMath{\xi_j|A'|,\,\xi_j^2|\nabla_UA'|<\alpha_1,} \\[3pt]
\FitMath{|\mathbf J|<\gamma_0\xi_j^{-3},}
\end{gather*}

and, on every cube of side $O(M_1\xi_j)$ contained in $\mathcal L_j\cup\mathcal L_{j+1}$, to admit a real gauge $U=e^{i\eta A}$ with $\xi_j|A|,\xi_j^2|\nabla A|\le\alpha_0$. In this condition the bound is $\alpha_0$ on the entire cube, and not $O(M_1)\alpha_0$. Pairs in which the current is a separate variable, with a radius $\gamma_0$ of its own, are due to Ba{\l}aban (\textbf{109}, p. 262, (1.10)--(1.14)). In the present setting they are formed on a sequence, with the local scales. Throughout, $J(\mathbf U)$ is the first variation of the action of the \emph{inner} plaquettes, namely those inside $\Omega_0$. At every bond whose six plaquettes are all inner, it coincides with the current $D^*_{\mathbf U}F_{\mathbf U}$; when $\Omega_0=T_\ell$, every bond is of this kind. \emph{On shell}: the pair satisfies $\mathbf J=J(\mathbf U)$ at the bonds that touch $\Omega_1$. \emph{Real}: $A'=0$, and $\mathbf J$ is real. \emph{Local class on $W$}: the same conditions, required inside $W$.

\phantomsection\label{notdef:propagators}
\emph{Operators.} Relative to $\mathbf U$, the \emph{straightened} $j$-fold average of $e^{i\eta A}\mathbf U$ takes, at a row $c$ of level $j$, the value $\exp i\mathcal Q(A)(c)$. Here $\mathcal Q=\mathsf m\,(Q+C^{(2)}+\dots)$, and $\mathsf m$ is multiplication by $\xi_c$. Thus $Q$, which maps $1$-forms to rows, is the linear part per unit length, and $C^{(2)}$ is the quadratic part per unit length. When the level has to be displayed, the straightened average is written $\mathcal Q^{\rm st}_j(\mathbf U;A)$. In the units of level $j$, its linear part is the operator $Q^\sharp_j(\mathbf U)$ of the Notation, \ensuremath{\S}3, which is built from the means over the contour families of $\mathcal M$. The index st stands for ``straightened''; the \emph{straight} linear average $Q^{\rm st}_j(\mathbf U)$ of the Notation is a different operator. Consequently, in this part the unadorned letters $\mathcal Q$ and $Q$ differ from those of \ensuremath{\S}3 of the Notation: they denote neither the logarithm $\mathcal Q_j(\mathbf U;A)$ of $\mathcal M^j(e^{i\eta A}\mathbf U)\,\mathcal M^j(\mathbf U)^{-1}$ nor its linear part $Q_j(\mathbf U)$. The straightened and the plain averages are related by \emph{site factors} $v_n$; these are the node factors $v_n=v_n[A]$ of the Notation, \ensuremath{\S}3, and in 1c.10 they are also called node factors. These are group-valued functions on the sites of $\mathbb T^{(n)}$, and they accumulate from one step to the next (statement 3: $\mathcal M^n(e^{i\eta A}\mathbf U)=\big(e^{i\mathcal Q^{\rm st}_n(A)}\mathcal M^n\mathbf U\big)^{v_n}$). The linear parts of the straightened and of the plain average already differ at $\mathbf 1$, by a gauge term on the coarse lattice. For the plain averages, two facts of statement 3 fail: the identity $QD=\bar DQ'$, and the bound $1+C(\dots)$ on the derivative. In that identity, at a row of level $j$, $\bar D$ is the covariant difference on $\mathbb T^{(j)}$ formed with $\mathcal M^j\mathbf U$, and $Q'$ takes, at each of the two end points of the row, the $j$-fold average of a site function; this applies also to the outer end point of an outward row, which is a top of level $j-1$. The operator $Q'$ maps site functions to tops. It is the linearisation of the average of a gauge transformation, taken with the same contour variables. At a row of level $j$, the notation $\bar DQ'$ uses the $j$-fold site average at both endpoints. At the outer endpoint of an outward row, which is a top of level $j-1$, this value is the additional transported $L^4$-block average of its own-level top values. Thus $\bar D$ means the extended endpoint operator of \cref{ch:layered-current-precision}, not an ordinary difference of the two own-top components. The symbol ${}^{\mathsf T}$ denotes the transpose with respect to bilinear pairings, so that everything remains holomorphic. The operator $a$ acts on rows as multiplication by $a_0\xi_c^{-2}$, where $a_0$ is a fixed number (equal to $1$ in the sub-parts). On tops it acts as multiplication by a comparable weight $\approx\xi_y^{-2}$. We write $D=D_{\mathbf U}$. Finally, $\Delta^{\rm W}_{\mathbf U}$ is the Hessian of the action in $\Omega_0$. In statements 3--8 it is the Hessian of the action of the inner plaquettes, which is an exactly invariant function. In statements 1, 2 it is obtained by taking the Hessian on the entire lattice and compressing it to the bonds of $\Omega_0$.

\begin{gather*}
\FitMath{G'=(\Delta_{\mathbf U}+Q'^{\mathsf T}aQ')^{-1},} \\[3pt]
\FitMath{R=1-G'Q'^{\mathsf T}(Q'G'^2Q'^{\mathsf T})^{-1}Q'G',} \\[3pt]
\FitMath{\mathsf S=1-DG'RD^{\mathsf T},}
\end{gather*}


\begin{gather*}
\FitMath{\mathsf G_0=(\Delta^{\rm W}_{\mathbf U}+DRD^{\mathsf T}+Q^{\mathsf T}aQ)^{-1},} \\[3pt]
\FitMath{\mathsf G=(\mathsf G_0^{-1}-\Delta'_\pi[\mathbf J])^{-1},} \\[3pt]
\FitMath{\mathsf G_1=(\mathsf G^{-1}-\mathsf S^{\mathsf T}\Delta^{(2)}[\mathsf H^{\mathsf T}\mathbf J]\mathsf S)^{-1},} \\[3pt]
\FitMath{\mathsf H_{(1)}=\mathsf G_{(1)}Q^{\mathsf T}(Q\mathsf G_{(1)}Q^{\mathsf T})^{-1}}
\end{gather*}

(where a bracketed index occurs, the formula holds both with that index and without it). The operators are defined in this order: $\mathsf G$, then $\mathsf H$ from $\mathsf G$, then $\mathsf G_1$, then $\mathsf H_1$ from $\mathsf G_1$. The symbol $\Delta'_\pi[\mathbf J]$ denotes a specified current-linear completion whose required physical Ward identity is stated below. Its complete form is part of the operator data: linearity and that Ward identity alone do not determine its horizontal or independent-current terms. The distinct mean-constraint correction $\Delta^{(2)}$ is retained separately. Moreover, $(A,\Delta^{(2)}[\mathsf j]A)=2(C^{(2)}(A),\mathsf j)$. The \emph{six propagators} are 

\begin{gather*}
\FitMath{\mathsf G,\mathsf G_1,\mathsf H,\mathsf H_1,(Q\mathsf GQ^{\mathsf T})^{-1},(Q\mathsf G_1Q^{\mathsf T})^{-1}}
\end{gather*}

. In addition we set $\mathfrak G=\mathsf G_1-\mathsf G_1DRD^{\mathsf T}\mathsf G_1-\mathsf H_1Q\mathsf G_1$. For step $k+1$ we work on the bonds of $\mathbb T^{(k)}$ that touch $\Omega_{k+1}$, with $V=\mathcal M^k\mathbf U$. Let $\mathsf C_V$ be the elimination of the pivot bonds, and let $\mathsf G^{(2)}_V$ be the quadratic part of the soft axial gauge fixing. Let $\mathsf E^{(2)}$ be the quadratic current form of the one-step constraint, which is imposed through the plain $\mathcal M$ at $V$ and whose linear part is the $Q_V$ of the Notation, \ensuremath{\S}3. Its notation includes the actual pivot multiplier. More precisely, write $\mathfrak q_V$ for the raw constraint logarithm and $h_V$ for the pivot right inverse of $D\mathfrak q_V(0)$. In raw counting coordinates, its bilinear form contracts $D^2\mathfrak q_V(0)[u,v]$ with $h_V^{\mathsf T}\widehat j$, where $\widehat j$ is the raw coarse-bond current. The unhatted form uses the corresponding raw/per-length dual conversion. This is the curved-fibre readout derived in \cref{ch:layered-current-precision}; the coarse-bond current is not directly identified with a new-row multiplier. We set

\begin{gather*}
\FitMath{\Delta_k=(Q\mathsf G_1Q^{\mathsf T})^{-1}-a-\mathsf E^{(2)}[\mathsf H_1^{\mathsf T}\mathbf J],} \\[3pt]
\FitMath{\mathsf A=\mathsf C_V^{\mathsf T}(\mathsf G^{(2)}_V+\Delta_k)\mathsf C_V,} \\[3pt]
\FitMath{C^{(k)}(Y)=(\chi_Y\mathsf A\chi_Y)^{-1}.}
\end{gather*}

At the physical critical pair $(U_{k+1},J(U_{k+1}))$ of the step, the required identification is $\mathsf A=\mathsf C^{\mathsf T}\Delta^{(k)}\mathsf C$. This identity must hold for the complete projected forms, including the pivot graph, all current and mean-constraint corrections, and the specified soft-gauge term. The absolute or relative contour-gauge convention must be checked in that comparison; a Ward identity alone does not establish the identification. Its action-value component is the quadratic term of the action of the background, regarded as a function of its datum. This is Ba{\l}aban's expansion \textbf{102}, 277, (174), together with \textbf{109}, (2.4)--(2.8), (2.11). There the derivation is carried out on the whole lattice, using Ba{\l}aban's own rows and his averages. For $\mathcal M$ and the straightened rows, on the sequence of a label, the expansion is obtained from the system with a multiplier of \cref{ch:multilevel-backgrounds}. That multiplier coincides with the gradient of the action of the background with respect to its datum, as one sees by differentiating the constraint and pairing the result with the Euler--Lagrange equation. A single differentiation of the system yields the derivative of the multiplier. Here $\omega$ is the physical row multiplier in $J(\mathbf U)=Q^{\mathsf T}\omega$, with the declared cell-weighted transpose; it is unique on the independent datum space when $Q$ is onto. For the critical pair, the required identified form of $\mathsf G_1^{-1}$ is the projected Lagrangian Hessian $\mathsf S^{\mathsf T}K_{\rm L}\mathsf S$, together with the gauge term and the row mass, where $K_{\rm L}=\Delta^{\rm W}-\Delta^{(2)}[\omega]$. On the declared Landau slice the projections act as the identity. For the canonical completion this identity follows from
\cref{prop:acc-current-identification}; its proper outer-fibre
identification, including the pivot second derivative, is proved
in \cref{pgc:section}. Hence this derivative is obtained by applying $(Q\mathsf G_1Q^{\mathsf T})^{-1}-a$ to the datum. When the datum moves along a curve contained in the fibre of the one-step average, it picks up a second-order term, and this term is determined by the quadratic part of the one-step constraint. Paired with the gradient, this term yields $-\mathsf E^{(2)}[\mathsf H_1^{\mathsf T}\mathbf J]$. Ba{\l}aban also gives a second expression for the first two terms of $\Delta_k$, namely $\mathsf H_1^{\mathsf T}[\mathsf G_1^{-1}-DRD^{\mathsf T}-Q^{\mathsf T}aQ]\mathsf H_1$. It differs from them by $-\mathsf H_1^{\mathsf T}DRD^{\mathsf T}\mathsf H_1$, a term that vanishes on shell. The bounds and the positivity in statement 6 concern the form displayed above.

\paragraph{Positivity of the compressed fluctuation form.}
The direct argument in \cref{pgc:top-Schur,pgc:outer-geometry,pgc:proper-covariance}
proves the proper lower bound, complex continuation and local inverse
families for the precise relative-contour form below. The following
local-system comparison describes an alternative positivity mechanism.
Consider an admitted real layered pair in statement 6
with the complex potential slot equal to zero and with
real current. In this argument the local pieces have
width \(\mathfrak M\) on the fluctuation scale, with
\(\mathfrak M\ge M^{\mathrm{pcs}}\), \(R_1\ge L\), and
\(R_1\mathfrak M\) at least the floor of the local
propagator construction. The torus side in level-\(k\)
units is at least \(2L\mathfrak M\). Take \(Y\) to be
the free-bond set for a subfamily of the retained top
blocks in the conditioning geometry. The form \(T_Y=\chi_Y\mathsf A\chi_Y\)
acts on the free bonds in \(Y\), after pivot elimination.
To compare it with a positive form, use a partition
\(\sum_\kappa h_\kappa^2=1\) on those bonds, with each cutoff
supported inside a larger piece on which the local system
is defined. In each piece, continue the field in a gauge
with a small potential and keep the three current slots,
including the two multiplier readouts
\(\mathsf H^{\mathsf T}\mathbf J\) and
\(\mathsf H_1^{\mathsf T}\mathbf J\). These readouts obey
the smallness bounds in the units of their rows by the
weighted operator estimates. Perturbation of the flat
system gives the inverse \(C_\kappa\) in the fluctuation
slot with \(cI\le C_\kappa\le pI\), where \(c,p>0\) are
independent of the piece. The comparison is made for the
completed local system, including the pivot graph and
gauge block.

Write \(P=\sum_\kappa h_\kappa C_\kappa h_\kappa\),
with restriction and extension by zero in each summand.
Then \(cI\le P\le pI\). To compute \(PT_Y\), take a free
field \(B\), set \(f=T_YB\), and complete \(B\) to the
solution \(\Phi\) of the coupled equations with right-hand
side \(f\) in the fluctuation slot and zero in the other
slots. The uniform auxiliary inverse and pivot bounds
give \(\|\Phi\|_\infty\le C\|B\|_\infty\) in local units.
The local and global systems agree on inputs supported
in the intermediate piece, with output in the larger
piece. Choose \(\widetilde h_\kappa\) supported in that
intermediate piece and equal to one on the cutoff
support and its interaction neighborhood. The local
system applied to \(h_\kappa\Phi\)
has right-hand side \(h_\kappa f\) in the fluctuation slot
plus the cutoff commutator applied to
\(\widetilde h_\kappa\Phi\). Solving the local equations and
summing their free-field slots yields \(PT_Y=I-E\).
The commutator costs \(C/\mathfrak M\) in the scale-weighted source
norm, and bounded overlap of the cutoffs gives
\(\|E\|_{\infty\to\infty}\le C/\mathfrak M\le1/8\).

This last bound controls the eigenvalues of \(PT_Y\).
Since \(PT_Y\) is similar to the real symmetric matrix
\(P^{1/2}T_YP^{1/2}\), that matrix is at least \(7I/8\).
Thus \(T_Y\ge 7I/(8p)\), and the inverse on the free bonds
has norm at most \(8p/7\). On the axial slice the same
calculation uses the free fields vanishing on the box
trees. The reconstructed fields are supported on the
free bonds together with their pivot bonds. The bounded
norm of the pivot lift converts the free-field lower bound
into the lower bound in the norm of the reconstructed
field.

The walk terms for the covariance of the compressed
form are written against counting measure on its free
bonds. Their local comparison uses the traces of the
regions and of \(Y\) on the carrier, and the regular pair
on the enlarged pieces used in the two local systems.
For the history induction the forms to be integrated must
be identified with this compressed form, or compared with
it with an error smaller than its lower bound. In
particular, selecting part of a walk sum does not
preserve its positivity by this argument. To extend the
lower bound to the complex integration cores one must
bound the change in the complete form on the specified
cores.

\emph{Walks.} A \emph{cell} is a union of $O(1)$ cubes of side $O(M_1\xi_j)$ lying in two consecutive layers. In a \emph{walk} $\omega=(X_0,\dots,X_m)$, any two consecutive cells overlap. In a \emph{tree walk}, each $X_i$ overlaps an earlier cell, called its parent, and the cells are enumerated depth first. With arbitrary parent maps, one sequence of length $m$ is counted as many as $m!$ times. The \emph{hull} $[\![\omega]\!]$ enlarges the walk by $O(M_1)$ local units and is not cut off at $\Omega_0$. The \emph{carrier} is $[\omega]=[\![\omega]\!]\cap\Omega_0$. Finally, $d_{\mathbf B}(\omega;y,y')$ is the smallest weighted length of a tree that passes through $y,y'$ and meets the cells $X_1,\dots,X_m$. The tree is not required to meet $X_0$.

Several letters have a different meaning in the sub-parts 1c.1, 1c.6 and 1c.10, as follows. The cube size of the walks is denoted $\mathfrak M$ in 1c.1. In 1c.10 the cubes $\square$ that carry parameters (statement 7), whose side is $\ell_{\mathsf P}$, are denoted $\Delta$. Also there, $\Gamma_j$ denotes the set of rows of level $j$. In 1c.10, $Q$ stands for $\mathsf mQ$, and slanted $H,H_1$ stand for $\mathsf H\mathsf m^{-1},\mathsf H_1\mathsf m^{-1}$. Within 1c.6, which treats averaging operators at complex configurations together with the gauge identities on which the propagators rest, the operator $Q^{\rm st}_k(\mathbf U)=\eta\,\mathbf Q(W_{k-1})\cdots\mathbf Q(W_0)$, where $\mathbf Q(W)$ denotes the linear part of a single step, stands everywhere for the linear part of the straightened average, namely the $Q^\sharp_k(\mathbf U)$ of \ensuremath{\S}3 of the Notation; it is not the straight average defined in that section of the Notation.

In comparison with the Notation:
\begin{itemize}
\item $R_1$ denotes a number (as in \textbf{116}, p. 3) and not a region $R_t$.
\item A cell means a cell of a walk, as opposed to a level-$k$ cell of \ensuremath{\S}9.
\item The hull of a walk differs from the hull in a coarser partition.
\item $\mathfrak G$ denotes an operator, not the gauge-fixing factor $\mathfrak G_k$.
\item $\mathsf A$ is the fluctuation operator, which is written $\mathsf K$ in \ensuremath{\S}5 of the Notation.
\item $\Delta(y)$ is the block $B^j(y)$, regarded as a cube.
\item $s(\square)$ is a parameter, not a depth.
\item $\bar\kappa$ is the exponent in the radius of the parameters, the same $\bar\kappa$ as in the order of the constants fixed in the Notation. \textbf{116} and 1c.10 write this exponent as $\kappa_1$; that $\kappa_1$ is not the $\kappa_1$ of the Notation, which is the rate at which the large-field factor decays.
\item $\gamma,\gamma^{\rm pos}$ are lower bounds, not ceilings on couplings.
\item $\gamma_0$ is the radius of the current, which is also written $\alpha_0$ in \ensuremath{\S}11.
\item $\lambda$ denotes a site function (an infinitesimal gauge transformation); in statement 8 and in mechanisms (j), (k), however, it denotes a multiplier, that is, a function on rows.
\item $X$ denotes a localisation domain in statement 0, $X_i$ is a cell, $X$ in (I) is a function on rows, and $X_\chi$ is a gauge direction.
\item $A'$ is the complex factor of the triple $(U,A',\mathbf J)$; the one exception is 1c.10, where, in the system of clause 1 of the Theorem (the decoupled response at complex parameters), it stands for the unknown $H_1(s)B'+\mathbf A$.
\item $S$ denotes, in statement 6, a set of coarse cubes and, in \eqref{eq:propagators-II}, an operator; $\mathsf S$ denotes the operator $1-DG'RD^{\mathsf T}$.
\end{itemize}

In 1c.10 the following further conventions hold:
\begin{itemize}
\item $\theta_0$ denotes the decay rate with respect to $d_{\mathbf B}$; it is not one of the members of the sequence $\theta$ of \ensuremath{\S}6.
\item $\mathsf r=v[\mathbf H]$ is a function on the nodes with values in $G$, and not an isometry.
\item $\beta$ and $\beta_\natural$ denote a small real datum, not the bare inverse coupling.
\item $\omega(\cdot)$ is the map that sends a field to its multiplier; it is not a walk.
\item $\mathsf E_1=G_1DRD^{\mathsf T}G_1$ denotes an operator rather than an event.
\item $P_1(s)$ is the operator built from the walk data of $D(1-R)D^{\mathsf T}$, and neither a region nor a momentum.
\end{itemize}

\textbf{(ii) The principal statements.} All constants depend only on $L$ (and on H{\"o}lder exponents). There are two exceptions, and neither enters any localised bound, rate, weight or sum of weights. The first is the bound of statement 7, which serves only for convergence. The second is the entrywise bound of mechanism (i), which makes the family normal. The numbers $\gamma$, $\gamma^{\rm pos}$, $\delta_0$, $\delta_*$, $\delta'$ are positive constants of this kind.

In statement 2 the side of the torus, measured in top units, is at least a small fixed multiple of $R_1M_1$; in the sentence of that statement which concerns a single lattice, it is at least $C$. From statement 4 onwards the side is at least $CM_1$. In 5, the single assignment of terms valid for all tori needs room for three cubes of side $O(M_1)$ in each direction. Statements 7, 8 need room for three cubes of side $\ell_{\mathsf P}$. These are lower bounds on the volume, and no constant depends on the volume except the bound of statement 7 just mentioned.

In what follows:
\begin{itemize}
\item $x\in\Delta(y)$, and the sources $f$ lie near $\Delta(y')$;
\item $\mathsf e=e^{-\delta_0d_{\mathbf B}(y,y')}$, up to a fixed fraction of the rate;
\item $\|\cdot\|_\varepsilon$ denotes a H{\"o}lder seminorm;
\item kernels are taken with respect to the counting measure;
\item on $1$-forms, the covariant Laplacian is $\Delta_{\mathbf U}=\sum_\nu\nabla_\nu^{\mathsf T}\nabla_\nu$, where $\nabla_\nu$ are the covariant differences of the components;
\item $\Delta=\Delta_{\mathbf 1}$ is the full vector Laplacian $\partial^*\partial+\partial\partial^*$, with $\partial=D_{\mathbf 1}$. It acts componentwise and is not the $\Delta=\Delta^{\rm W}_U$ of the Notation, \ensuremath{\S}5.
\end{itemize}

\textbf{0. The domains of analyticity are contained in the classes.} Consider a pair in the domain of an activity on $X\in\mathcal D_j$. Such a pair satisfies $|\partial U-\mathbf 1|,|\partial\mathbf U-\mathbf 1|<\alpha_0\eta^2$, $|A'|,|\nabla_UA'|<\alpha_1$ and $|\mathbf J|<\alpha_0$, and on cubes of side $O(LM)$ it admits real gauges with $|A|,|\nabla A|<O(LM)\alpha_0$. It is a regular pair for the tower of regions obtained from $X$ by erosion (1b.8), and the radii with which it is regular are

\begin{gather*}
\FitMath{\big(O(LM)\,\alpha_0,\,\alpha_1,\,\alpha_0\big).}
\end{gather*}

Let a label be produced by steps in which each region $\Omega_n$, $n\ge1$, is a union of $LM$-cubes of level $n$. By a \emph{layered domain} we mean a domain defined by the layered bounds of \cref{sec:large-field-the-tests-the-nested-regions-and-the-content-of-a}, \emph{Structures}, whose radii do not exceed the running radii of the level of each layer. Only the following of these bounds enter. On the layer $\mathcal L_n$ of the label, $n\ge1$, one has $|\partial U-\mathbf 1|,|\partial\mathbf U-\mathbf 1|<\alpha_{0,n}\eta^2\xi_n^{-2}$, $|\mathbf J|<\alpha_{0,n}\xi_n^{-3}$ and $\xi_n|A'|,\xi_n^2|\nabla_UA'|<\alpha_{1,n}$. The cubes of level $n$ of the domain are the closed cubes of side $CM\xi_n$ that are contained in $\mathcal L_n\cup\mathcal L_{n+1}$ and meet $\Omega_{n+1}^c$; on the part in $X$ of each of them there is a real gauge with $\xi_n|A|,\xi_n^2|\nabla A|<CM\alpha_{0,n}$. Then on $X\cap\Omega_1$ a pair of a layered domain belongs to the local class of the \emph{label sequence}, which is the sequence formed by the $\Omega_n$ of the label. The radii are $O(M)\max_n\alpha_{0,n}$ and $\max_n\alpha_{1,n}$, together with $\max_n\alpha_{0,n}$ for the current. In that class the gauge clause is derived for those cubes of level-$n$ size which meet $\mathcal L_n$. For $n=0$, on the cubes inside $\mathcal L_1$, it is derived from the gauge of a cube of level $1$, which can be used at the index $0$ since $\xi_0<\xi_1$. The same remark applies to a cube of level-$n$ size inside $\mathcal L_{n+1}$, because some cube of level $n+1$ of the domain contains it. This smaller cube is contained in $\Omega_{n+1}$, intersects $\Omega_{n+2}^c$, and its side is below $M\xi_{n+1}$. The derivation of the clause on the cubes meeting $\mathcal L_n$ uses no other property of a cube; hence repeating the same steps with the grids of level $n+1$ produces the larger cube. The factors $O(LM)$, $O(M)$ are the cost of the gauge clause, since on a cube of side $O(M)$ a gauge has $|A|\le O(M)\alpha_0$. The field strengths are fixed after $M$.

\textbf{1. The flat theorem.} Let $\mathbf U=\mathbf 1$, $\mathbf J=0$, and let $Q$, $Q'$ be the straight averages taken over the centred blocks. Then on every separated sequence

\begin{gather*}
\FitMath{|\mathsf G_0f|,|\nabla\mathsf G_0f|,|\mathsf G_0\nabla^*f|,|\Delta\mathsf G_0f|\,(x)\le C[\xi_y^2,\xi_y,\xi_y,1]\,\mathsf e\,|f|,} \\[3pt]
\FitMath{|\nabla\mathsf G_0\nabla^*f(x)|\le C_\varepsilon\,\mathsf e\,\big(\xi_{y'}^\varepsilon\|f\|_\varepsilon+|f|\big),}
\end{gather*}

together with matching bounds in H{\"o}lder and $L^2$ norms. The same bounds hold for $G'$, and

\begin{gather*}
\FitMath{|(Q'G'^2Q'^*)^{-1}(y,y')|\le C\xi_y^{-4}\mathsf e,} \\[3pt]
\FitMath{|(Q\mathsf G_0Q^*)^{-1}(c,c')|\le C\xi_c^{-2}\mathsf e,} \\[3pt]
\FitMath{|\mathsf H(b,c)|,|\nabla\mathsf H(b,c)|\le C[1,\xi_y^{-1}]\mathsf e,}
\end{gather*}


\begin{gather*}
\FitMath{Q\mathsf G_0Q^*\ge\gamma\,\xi_c^2,} \\[3pt]
\FitMath{\big\langle B,\big[(Q\mathsf G_0Q^*)^{-1}-a\big]B\big\rangle\,\ge\,c\,L^{-5}\|B\|^2}
\end{gather*}

for every $B$ with three properties: it lives on the bonds of $\mathbb T^{(k)}$ that touch a union of $L$-blocks lying well inside $\Omega_k$; its straight one-step average vanishes; and it vanishes on the axial trees of the blocks. Each block has one such tree. It consists of the contours joining the site $y\in\mathbb T^{(k+1)}$ of the block to each of its points $x$, along which the coordinates are changed successively, always in the same fixed order. The gauge so defined is the strict axial gauge, which is a hard gauge fixing. It is not the block axial gauge (Notation, \ensuremath{\S}3), where the means $V(y,x)$ over contour families are all equal to $\mathbf 1$. For the plain one-step constraint, in that gauge or on trees, the lower bound is proved again in 1c.8. If the gauge condition is dropped, the form has null directions.

\textbf{2. Real backgrounds with Ba{\l}aban's average.} Let $Q,Q'$ be formed with one contour per site and moved to centred blocks by a translation. Let the Hessian be the compressed one. Let $U$ be real, with local gauges and with small plaquettes; for $\mathsf G_0$ this includes the plaquettes that stick out of $\Omega_0$. Then $\Delta_U+Q'^*aQ'$, $Q'G'^2Q'^*$ and $\mathsf G_0^{-1}$ are positive. The operators $G',\nabla G',\mathsf G_0,\nabla\mathsf G_0,(Q'G'^2Q'^*)^{-1}$ satisfy the bounds of 1 for every $\Omega_0$, including a collar of $\Omega_1$ and a concentric cube. On a single lattice, and for $U$ with $|\partial U-\mathbf 1|,|D_U^*\partial U|<\alpha\eta^2$, where $\alpha$ is small, also $Q\mathsf G_0Q^*\ge\gamma>0$, and $\mathsf G_0Q^*(Q\mathsf G_0Q^*)^{-1}$ is a right inverse of the linearised constraint; this right inverse and its gradient decay like $e^{-\delta_*|y-y'|}$. The rows of $\mathcal M$ agree with these rows up to a transport and an error $O(\alpha_0)$ in the row and column norms, uniformly in the level. Statements 3--8 concern different operators, built from the rows of $\mathcal M$ and the Hessian of the inner plaquettes. At $\mathbf 1$ the two kinds of rows coincide.

\textbf{3. The averages and the gauge identities.} The operators $Q,Q'$ are local and holomorphic in $\mathbf U$. Their row sums are $\le1+C(\alpha_0+\alpha_1)$ and their entries are $\le C(\eta/\xi_c)^4$, independently of the number of levels. At each row the non-linear average is holomorphic in $(A',A)$ on a polydisc $\sup|A'|,\sup|A|<r_L$, where $r_L$ is determined by $L$ and measured in units of the row's level. Fix a bond of $\mathbb T^{(k)}$, and write $\delta/\delta A_b=\eta^{-4}\partial/\partial A(b)$. Let $\mathcal Q^{\rm st}_n$ be the $n$-fold straightened average, and let $v_n$ be the accumulated site factors belonging to it. Then on this polydisc

\begin{gather*}
\FitMath{\Big|\frac{\delta\mathcal Q(A)(c)}{\delta A_b}\Big|\le1+C(\alpha_0+\sup|A'|+\sup|A|),} \\[3pt]
\FitMath{\Big|\frac{\delta(\mathcal Q-\mathsf mQ)(A)(c)}{\delta A_b}\Big|\le C\sup|A|,} \\[3pt]
\FitMath{\mathcal M^n(e^{i\eta A}\mathbf U)=\big(e^{i\mathcal Q^{\rm st}_n(A)}\mathcal M^n\mathbf U\big)^{v_n}.}
\end{gather*}

Let $\bar D$ be the covariant difference of $\mathcal M^j\mathbf U$ on $\mathbb T^{(j)}$. At a row of level $j$, let $Q'\lambda$ denote the $j$-fold average taken at both ends of the row. For an outward row, the outer end is a top of level $j-1$, not a top of level $j$. Then

\begin{gather*}
\FitMath{QD\lambda=\bar DQ'\lambda,} \\[3pt]
\FitMath{\Delta^{\rm W}_{\mathbf U}-\Delta'_\pi[J(\mathbf U)]=\mathsf S^{\mathsf T}\Delta^{\rm W}_{\mathbf U}\mathsf S,} \\[3pt]
\FitMath{\mathsf G_{(1)}DR=DG'R+\mathsf G_{(1)}\Delta'_\pi[\mathbf J-J(\mathbf U)]DG'R.}
\end{gather*}

On shell, the defect is zero, $RD^{\mathsf T}\mathsf H_{(1)}=0$, and $\mathsf HB$ is a critical point of $\frac12(A,\Delta^{\rm W}_{\mathbf U}A)$ on $\{QA=B,\ RD^{\mathsf T}A=0\}$, with critical value $\frac12(B,[(Q\mathsf GQ^{\mathsf T})^{-1}-a]B)$.

\textbf{4. The six propagators at regular pairs.} The ceiling on the radii $\alpha_0,\alpha_1,\gamma_0$ is fixed before $M_1$ and is \textbf{independent of $M_1$}. At every regular pair all the inverses exist, they are holomorphic in $(\mathbf U,\mathbf J)$, and they are real at real pairs. At real pairs $\mathsf G_0,\mathsf G,\mathsf G_1>0$, both on shell and off shell. The six propagators and $(Q'G'^2Q'^{\mathsf T})^{-1}$ satisfy the upper bounds of 1 with covariant differences, and $\mathsf G$ and $\mathsf G_1$ satisfy those of $\mathsf G_0$. The exception is the bound on $\Delta_{\mathbf U}\mathsf G_{(1)}$, and the H{\"o}lder quotients are taken only over $|x-x'|\le\xi_y$. For $G'$ at complex pairs the bounds reduce to the supremum bounds on $G'f$ and $DG'f$. Moreover $|D(1-R)D^{\mathsf T}(b,b')|\le C(\eta/\xi_{y'})^4\xi_y^{-2}\mathsf e$.

\textbf{5. Walk expansions.} Let $\mathsf O$ be any of the following operators:
\begin{itemize}
\item one of the six propagators;
\item $1-R$ or one of its derivatives;
\item $\mathfrak G$;
\item Ba{\l}aban's $H_0,\tilde G$, namely $\mathsf H\mathsf m^{-1}$ and $\mathsf G-\mathsf HQ\mathsf G$ at $\mathbf J=0$;
\item $C^{(k)}(Y)$ or $\chi_Y\mathsf A\chi_Y$, these two taken at the pairs that are regular for the longer sequence $(\Omega_j)_{j\le k+1}$, where $M_1$ is larger and the ceiling on the radii is smaller than in 4.
\end{itemize}

For $\tilde G$ the indices are the walks together with the pairs of walks in which the second walk starts at a cell that a row joins to some cell of the first walk, its cells being placed after those of the first. In the display, $|\omega|$ is then the total length of the two walks, and the step that joins them contributes no small factor. Each such operator is then a sum $\sum_\omega\mathsf O_\omega$, where, in the norms in which $\mathsf O$ is bounded,

\begin{gather*}
\FitMath{\|\mathsf O_\omega\|\le(\text{bound of }\mathsf O)\cdot(CM_1^{-1/2})^{|\omega|}M_1^{-|\omega|/2},} \\[3pt]
\FitMath{d_{\mathbf B}(y,y')\,\text{replaced by}\,d_{\mathbf B}(\omega;y,y').}
\end{gather*}

The first factor makes the sum over the walks converge. The second factor is left over, and it absorbs the factor $|\omega|+1$ of mechanism (i). The bounds of $\mathsf O$ are as follows:
\begin{itemize}
\item For $\mathfrak G$ and $\tilde G$ there are exactly four: the bounds on $|\mathsf Of|$, $|\nabla\mathsf Of|$, $|\mathsf O\nabla^{\mathsf T}f|$ and on the H{\"o}lder seminorm of $\nabla\mathsf Of$ that 4 provides for $\mathsf G$, $\mathsf G_1$. For $\tilde G$ they hold with a smaller fraction of $\delta_0$.
\item For $1-R$ and for its derivatives $D(1-R)$, $(1-R)D^{\mathsf T}$, $D(1-R)D^{\mathsf T}$, the relevant bounds are the kernel bounds $C(\eta/\xi_{y'})^4[1,\xi_y^{-1},\xi_y^{-1},\xi_y^{-2}]\,\mathsf e$. The last of these is displayed in 4.
\item For $\chi_Y\mathsf A\chi_Y$, the bound is the kernel bound that 6 states for $C^{(k)}(Y)$.
\end{itemize}

Each $\mathsf O_\omega$ has the following properties:
\begin{itemize}
\item it is holomorphic on the local class on $[\omega]$;
\item its dependence on the regions is through their traces on $[\![\omega]\!]$ alone;
\item it is gauge covariant and real at real pairs;
\item it is \textbf{the same function} for any two sequences, tori or cutoffs that a \emph{shift}, namely a translation in $M_1\xi_k\mathbb Z^4$, relates on the hull of (i), \emph{Walks}, which is the set just named; in the Notation, \ensuremath{\S}5, this hull is called the surroundings of $\omega$. For $C^{(k)}(Y)$, $\chi_Y\mathsf A\chi_Y$ the translation is in $M_1\xi_{k+1}\mathbb Z^4$ and must also carry $\Omega_{k+1}$ and $Y$.
\end{itemize}

The operators are covariant under every isometry that preserves
$\mathbb T^{(k)}$ and the regions, using the induced action on each
of their input and output spaces. At a reversed bare bond $b$ this
action sends $A(b)$ to $-\operatorname{Ad}(\mathbf U(b)^{-1})A(b)$.
At a reversed row or coarse bond $c$ of level $n$, it instead sends
$B(c)$ to $-\operatorname{Ad}(V_n(c)^{-1})B(c)$, where
$V_n(c)=(\mathcal M^n\mathbf U)(c)$ is the actual level-$n$ average;
in particular, the covariance and fluctuation operators on level $k$
use $V_k$. Site slots are transported without an orientation factor.
These formulas act on the independent positively oriented tangent
coordinates. The skew extension $\nu(\bar c)=-\nu(c)$ within a fixed
choice of orientation remains a separate convention.
Thus an operator from one slot to another transforms as
$\mathsf r_{\rm out}\mathsf O\mathsf r_{\rm in}^{-1}$.
For a term evaluated at a pair given only on its carrier, every
nonzero kernel entry has the complete averaging reads needed to form
both orientation factors. The transformed kernel is zero at zero
entries; filling undefined averages outside those reads by arbitrary
elements of $G^c$ does not change it. The maps preserve the complex
bilinear pairings:
$\operatorname{Ad}(g)^{\mathsf T}=\operatorname{Ad}(g^{-1})$,
and the isometry preserves each row's level and hence its cell weight.
Accordingly the same transformation law applies to the declared
transposes, with either physical or counting row measure, without an
additional density factor. At complex pairs this is a bilinear identity,
not a claim that the action preserves the norm. For $C^{(k)}(Y)$, $\chi_Y\mathsf A\chi_Y$ the isometry must also preserve $\mathbb T^{(k+1)}$ and $\Omega_{k+1}$, and $Y$ is carried along. Pivot bonds go to pivot bonds, since each pivot bond is the middle bond of its path. For single terms the covariance holds only under those isometries which in addition preserve the $M_1$-cubes of the top level. For $C^{(k)}(Y)$, $\chi_Y\mathsf A\chi_Y$ that level is $k+1$, and in the strict axial gauge the isometry must also preserve the axial trees. This covariance holds at the pairs whose image again lies in the class.


\textbf{(iii) The strategy.}

\emph{The chain.} The parts have the following content.
\begin{itemize}
\item 1c.1 (the operators and the statements about them) fixes the operators, the classes and the form of the statements. It counts the walks and derives statement 0. It also shows that the identification of $\mathsf A$ at the step reduces to Ba{\l}aban's expansion of the background action.
\item 1c.2 (the flat case, I) states the flat theorem and compares it clause by clause with the published statements. It also contains the exact algebra at $\mathbf 1$: existence, identities with boundary terms, and the independence of $R$ and $\mathsf H$ from $\Omega_0$.
\item In 1c.3 (the second flat case) the one-scale operator is treated uniformly in the spacing; then come $G'$, $(Q'G'^2Q'^*)^{-1}$ and $R$.
\item In 1c.4 (the flat case, III) the local two-scale operators are treated, followed by $\mathsf G_0$, $(Q\mathsf G_0Q^*)^{-1}$ and $\mathsf H$. With this, statement 1 is complete.
\item 1c.5 (Ba{\l}aban's theorems in a background) yields statement 2. On the whole lattice it comes from the published theorems; on sequences it comes from the flat case, by a shortened form of mechanism (f).
\item 1c.6 (the averaging operators at complex pairs; the algebraic identities) yields statement 3.
\item 1c.7 (the six propagators) uses mechanisms (f), (g) to derive, from the flat bounds and 1c.6, statement 4 and the expansions of 5.
\item In 1c.8 (the covariance; positivity) these two mechanisms are applied to a larger system. This yields statement 6, except for the expansion of $\mathsf A$ over sets of coarse cubes.
\item In 1c.9 (walk expansions as data) an expansion is regarded as a datum enjoying the short list of properties of 7. Statement 7 is deduced from these properties together with the geometry. Moreover, 1c.9 \emph{makes} terms that are real, covariant and identical on all tori with room; these properties do not follow from the list.
\item 1c.10 (the decoupled response) takes the operators as data and solves the systems (I), \eqref{eq:propagators-II}. Next it inserts the local maps and operators of 1c.6--1c.9; for $\mathfrak G$ and $D(1-R)D^{\mathsf T}$ the inserted data are those named at the end of 7. Last, it treats the case $s\equiv1$.
\end{itemize}

\emph{The order of \cref{ch:multilevel-backgrounds} and \cref{sec:propagators-propagators-in-a-background-field-exponential-decay-produced-by-the}, which involves no circle.} None of statements 1--7 mentions a background of \cref{ch:multilevel-backgrounds}, since the classes are defined by inequalities. The results are obtained in the following order:
\begin{itemize}
\item the geometry and the real calculus of $\mathcal M$ (1b.1--1b.2);
\item the flat case;
\item the lemmas of 1b.3 that are geometric or algebraic (folding, cells and rows, the algebra of $Q',Q,R$, sums over tops, weights for a general metric), and none of them uses any bound on a propagator;
\item 1c.5;
\item the rest of 1b.3; then existence, regularity (which provides the backgrounds of \cref{ch:multilevel-backgrounds} with their local gauges) and continuation (1b.4--1b.6);
\item 1c.6--1c.9. Their inputs are the flat case, the alignment of the blocks in 1c.5, and 1b.1--1b.2, together with, besides these, geometric facts in which neither a propagator nor a background of \cref{ch:multilevel-backgrounds} enters (meets of sequences; the covariance of $\mathcal M$ with respect to the isometries of the lattice);
\item the decay of the response, and the membership in the classes of the backgrounds of \cref{ch:multilevel-backgrounds} (1b.7--1b.9);
\item finally, the identification at $s\equiv1$.
\end{itemize}

\emph{Mechanisms.}

(a) \textbf{The metric matches the mass.} The form satisfies $\langle\lambda,(\Delta+Q'^*aQ')\lambda\rangle\ge\gamma\sum_j\xi_j^{-2}\|\lambda\|^2_{\mathcal L_j}$, and on each layer $d_{\mathbf B}$ has slope $\xi_j^{-1}$. A single conjugation by $e^{\varsigma\varphi}$, with $\varphi$ Lipschitz with respect to $d_{\mathbf B}$ (Combes--Thomas), therefore serves for all layers at once. It costs $O(\varsigma L^2)$, and this cost stays below $\gamma$ under a ceiling on the rate, $\varsigma\le c\gamma L^{-2}$, where $L$ is chosen first. Every layer that has to be traversed in full costs the gap, that is, the separation of the boundaries of consecutive regions, which exceeds $R_1M_1$ when measured in the unit of that layer; thus $d_{\mathbf B}(y,y')\ge R_1M_1(|j-j'|-1)_+$. Let a rate $\tau>0$ and a power $q\ge0$ be given, and choose the product $R_1M_1$ after them so that $\tau R_1M_1\ge C(1+q)\log L$. Then

\begin{gather*}
\FitMath{\sum_{y'}e^{-\tau d_{\mathbf B}(y,y')}\le C(\tau,L),} \\[3pt]
\FitMath{(\xi_{y'}/\xi_y)^q\le L^qe^{\tau d_{\mathbf B}(y,y')}:}
\end{gather*}

the factor $e^{-c\tau R_1M_1}$ for each level outweighs the growth $L^4$ in the number of sites. Conjugation and walk expansions do different work. Conjugation yields decay of the inverse as a whole, and it requires no cube size to be large; its rate is a fraction of the spectral gap, which is the lower bound $\gamma$ of the form. A walk expansion gives, in addition, a sum of terms, each of which depends on the background and on the regions only near its own walk. This property is what statements 5--8, the comparison of tori and \cref{ch:local-formation-program} make use of. Both methods are used. Conjugation serves for the scalar, one-scale and local operators and for the inverses on tops and rows. Walks serve for the assembly, namely of $\mathsf G_0$ in (d) and of the propagators in a background in (f). In \textbf{95}, p. 36, Ba{\l}aban describes conjugation as ``probably the simplest proof'' of the decay; he builds a walk representation instead, since ``such representations will be our basic tool in investigation of more complicated operators''.

(b) \textbf{No Dirichlet condition on a lattice finer than the blocks.} Because of the gap that separates $\Omega_0^c$ from $\Omega_1$, $\partial\Omega_0$ lies in $\mathcal L_0$. On that layer $Q=Q'=1$ and the gauge term is zero, and the local operator $G_\square$, which inverts the form on a cube $\square$ (see (d)), obeys $a_0\le\eta^2G_\square^{-1}\le a_0+16$. Finite range together with a mass comparable with the norm gives decay for an arbitrary set of bonds. The following examples show that this placement of the boundary matters. Let $G$ be a one-scale operator (spacing $\eta$, unit blocks) whose order is $\ge1$, with Dirichlet conditions imposed outside a union of \emph{whole} blocks, so that the lattice is finer than the blocks and no block is cut. Two examples, a flat face and a re-entrant edge, rule out constants uniform in the order, by powers of $\eta^{-1}$. This applies to the $L^2$ bounds on second differences, to the bounds on $\Delta G$, to those on $\nabla G$ and to its H{\"o}lder norm. It applies equally to the pointwise bound and to the H{\"o}lder bound involving $\nabla G\nabla^*$, taken in the form of statement 1, the powers being those that occur for $\nabla G$ and for its H{\"o}lder norm. Indeed, on the region every bounded $f$ supported in a block can be written as $\nabla^*F$ with $F=\nabla\Delta^{-1}f$, where $\Delta^{-1}$ denotes the free inverse on $\eta\mathbb Z^4$ (on a torus, the inverse on functions of mean zero, after $f$ has been modified on a block lying outside the region); moreover $|F|,\|F\|_\varepsilon\le C_\varepsilon\sup|f|$ uniformly in $\eta$, and the sum over the blocks is controlled by the exponential factor in the bound. For the $L^2$ bound involving $\nabla G\nabla^*$, by contrast, the two examples rule out no uniform constant. On a separated sequence no operator of this kind arises.

(c) \textbf{One scale, two scales.} Uniformity in $\eta$ comes from two sources. The first is the representation $(\Delta+1)^{-1}=\int_0^\infty e^{-t}e^{-t\Delta}dt$, combined with heat-kernel bounds that hold both below and above $t=\eta^2$. The second is a factorisation of $\partial(1-R)\partial^*$ through the unit lattice. At a seam one makes an exact change of gauge $A=A_2+H_jB-\partial H'_j\vartheta$. Here $A_2\in\ker Q_j$, and $B$ is a $1$-form on the coarse lattice whose values on the axial trees of the $L$-blocks of the coarse region are zero. Further, $\vartheta$ is a site function on the coarse lattice, supported in the coarse region, whose mean over every one of those blocks is zero; for these two classes the representation is unique. Finally, $H_j,H'_j$ are the one-scale minimisers, with $H'_j$ the minimiser of $\|\Delta\lambda\|^2$ rather than of $\|\partial\lambda\|^2$. This change of gauge splits the form orthogonally into three parts: a form in $A_2$, a form in $B$ and a form in $\vartheta-\vartheta_1$, where $\vartheta_1$ is linear in $A_2$ and $B$. If the coarse region is all of the torus, the two-scale representation of the inverse of the form so obtained becomes a recursion $G_j\to G_{j+1}$, where $G_j$ denotes the one-scale propagator of order $j$. It is applied for a single step and is never iterated. In the scalar theory, recursions of this kind add up to a series over the scales, and they served as the method of proof of regularity (\textbf{89}, (2.34)). Here $G_j$ occurs multiplied by other operators. Iteration then leads to long products of operators whose gauge conditions are not the same at different scales (a remark of Ba{\l}aban, \textbf{96}, pp. 246--247). For gauge propagators, accordingly, no series over the scales is employed; instead the one-scale $G_j$ has to be known simultaneously for all $j$ (1c.3). Three spectral gaps on the unit lattice take the place of a small parameter. The sharp one is stated for $B$ on a unit lattice that vanishes on the axial trees. In it, $Q_1$ and $\partial_1$ are the straight one-step average and the exterior derivative; the $Q_1(\mathbf 1)$ of the Notation, \ensuremath{\S}3, also contains site terms. The sharp gap reads

\begin{gather*}
\FitMath{L^{2}\sum_c|(Q_1B)(c)|^2+\sum_p|(\partial_1B)(p)|^2\ge cL^{-5}\|B\|^2.}
\end{gather*}


(d) \textbf{Many scales at $\mathbf 1$.} No expansion of $G'$ is required. For $u=G'(f+\nabla^*F)$ the conjugation of (a) yields 

\begin{gather*}
\FitMath{\|\xi^{-1}e^{\varsigma\varphi}u\|+\|e^{\varsigma\varphi}\nabla u\|\le C\gamma^{-1}(\|\xi e^{\varsigma\varphi}f\|+\|e^{\varsigma\varphi}F\|)}
\end{gather*}

, where $\xi$ denotes multiplication by the local scale. From this one passes to local $L^2$ norms, and then to interior estimates at the scale $\xi_y$, with the mass term treated as a source. The walk expansions are obtained afterwards, as data. For $\mathsf G_0$ we take $\sum_\square h_\square^2=1$, with $\zeta_\square=1$ near $\operatorname{supp}h_\square$, and we let $G_\square$ be the inverse, on the periodised cube $\square$, of the form of (c), taken with one scale or with two. For the cubes of level $0$ that extend beyond $\Omega_0$, $G_\square$ is instead the operator of (b). For a cube of level $\ge1$ the periodic two-scale model is retained, irrespective of where $\Omega_0^c$ and the finer layers lie; this is legitimate since across the two seams to $\mathcal L_j\cup\mathcal L_{j+1}$ the function $h_\square$ is cut off. With $P=1-R$ and $\mathcal P=\partial P\partial^*$ one has

\begin{gather*}
\FitMath{\mathsf G_0^{-1}\sum_\square h_\square G_\square h_\square=1-\mathcal R,} \\[3pt]
\FitMath{\mathcal Ph-h\mathcal P_\square=\zeta[\mathcal P_\square,h]+\zeta(\mathcal P-\mathcal P_\square)h+(1-\zeta)\mathcal Ph:}
\end{gather*}

of all these operators, only the non-local $\mathcal P$ cannot be transplanted. Each commutator costs $M_1^{-1}$, because $|\nabla h|\le C(M_1\xi_j)^{-1}$ stands against $|\nabla G_\square|\le C\xi_j$. The change of domain costs $e^{-cM_1}$; this follows from resolvent identities on the two domains, recombined by means of $1-P=R$. The tail is likewise exponentially small, and it is turned into a step $\square\to\square'$ by means of $\sum h_{\square'}^2=1$. The factor $M_1^{-1}$ for each step outweighs the number $C(\tau,L)$ of neighbours.

Lower bounds independent of the number of levels underlie the inverses on tops and rows. The bound $Q'G'^2Q'^*\ge\gamma\xi_y^4$ follows by duality with test functions $\psi_y$ on the blocks, with $Q'\psi_y=\delta_y$, whose second differences are $\le C\xi_y^{-2}$, across the faces as well. The bound $Q\mathsf G_0Q^*\ge\gamma\xi_c^2$, with $\gamma$ independent of $L$, follows from a right inverse $\mathsf E$ for all rows of all levels at once. This right inverse has scale-invariant energy, $\|\nabla\mathsf EB\|^2\le C\|\mathsf m^{-1}B\|^2$. The construction below gives the finite energy lift for the cornered-comb rows. Applying its estimate to the actual averaged rows requires the row comparison and the indicated pairing change. In 2, where there is a single lattice, the test forms are placed on the face separating the two blocks of a bond; there the share $\theta_0>\frac12$ of its own chains outweighs $1-\theta_0$. Neither matrix is bounded below by a number independent of the scales $\xi$. They are therefore rescaled, to $\Sigma=\xi^{-2}Q'G'^2Q'^*\xi^{-2}$ or $\xi^{-1}Q\mathsf G_0Q^*\xi^{-1}$. A matrix $\Sigma\ge\gamma$ with $|\Sigma|\le Ce^{-\tau d_{\mathbf B}}$ then satisfies $|\Sigma^{-1}|\le2\gamma^{-1}e^{-\tau'd_{\mathbf B}}$, where $\tau'$ is small compared with the spectral gap $\gamma$ (Combes--Thomas, Schur). Thus both inverses decay without any expansion, and $\mathsf H$ decays column by column. A sequence is compared with the same sequence cut down around a cube. The cut-down sequence is again separated and coincides with the given one on the doubled cube, and its $Q\mathsf G_0Q^*$ provides the local operators that enter the walk expansion of $(Q\mathsf G_0Q^*)^{-1}$. The comparison uses Green's identity with a cut-off, by way of $(\partial^*\partial+Q^*aQ)\,\mathsf G_0Q^*=Q^*$, and the gauge term vanishes on the columns.

\paragraph{An energy-bounded lift of all determining rows.}
\label{mb:all-row-energy-lift}

The row estimate used to invert the coarse Schur operator has a
finite construction.  It is useful to retain the distinction between
the cornered-comb rows used in this construction and the differential
of the averaged nonlinear constraint.  Write
\(Q^+_{\mathrm c}=(Q^{\mathrm c}_j|_{\widehat\Gamma_j})_j\)
for the former, with all inward rows retained in
\(\widehat\Gamma_j\).  Its output is a raw row logarithm, with
the counting pairing \(\sum_c\langle\omega(c),b(c)\rangle\).
In particular, \(Q^{\mathrm c}_0 A=\eta A\).  Fine fields have
the pairing \(\eta^4\sum_b\langle A(b),A'(b)\rangle\).
The rows here are the linear parts of the iterated cornered-comb
average, with the endpoint factors kept in its nonlinear
factorization.  They are not the response operator of the critical
background.

Assume \(L\ge13\), the separated aligned geometry of
\eqref{mb:scales}--\eqref{mb:separation}, with gap at least
\(L^2\xi_j\), and a real background satisfying
\[
 |\partial U(p)-1|\le\alpha\eta^2\mathsf L(p)^{-2}.
\]
The required smallness threshold for \(\alpha\) depends on \(L\),
but not on the depth.  In the construction below all norms of row
matrices are the counting \(\ell^2\) norms.

For a positive row \(c\) of level \(j\ge1\), choose its endpoint
\(y_c\) in the level-\(j\) layer.  Put \(N=L^j\),
\(n=(N-1)/2\), and \(w=\lfloor N/8\rfloor\).  In integer
coordinates \(P_\nu\) measured from the lowest corner of
\(B^j(y_c)\), set
\begin{align*}
 h(P)&=(w-|P-n|)_+,\\
 \varphi_\pm(P)&=(w-|P-(n\pm w)|)_+,\\
 g^+(P)&=(n+w)\varphi_+(P)-(n-w)\varphi_-(P),\\
 g^-(P)&=(n+w+1)\varphi_-(P)-(n-w+1)\varphi_+(P).
\end{align*}
The sign is positive when \(y_c=c_-\), and negative otherwise.
For a bare bond \(b\) in the direction \(\mu\) of \(c\), define
\[
 \psi_c(b)=\frac{N^3}{2\eta w^9}
      g^{\pm}(P_\mu(b_-))\prod_{\nu\ne\mu}h(P_\nu(b_-)),
\]
and set it equal to zero on the other bonds.  Transport a row vector
to \(b_-\) along the cornered comb in this cell; for a negative
row first transport it along the straight segment to that corner.
The resulting test form is denoted by \(Z_c e\).  At level zero
put \(Z_c e=\eta^{-1}e\) on \(c\), and zero elsewhere.  Set
\(Z\omega=\sum_c Z_c\omega(c)\) and \(\Theta=Q^+_{\mathrm c}Z\).

The profiles are supported a positive fraction of the cell width
away from its faces.  Extending them by zero introduces no jump
across a cell face.  The identities
\begin{gather*}
 \sum_P h(P)=w^2,\\
 \sum_Pg^+(P)(P+1)=\sum_Pg^-(P)(N-1-P)=2Nw^3,\\
 \sum_Pg^+(P)(N-1-P)=\sum_Pg^-(P)(P+1)=0
\end{gather*}
follow by summing each triangular profile about its center.  These
are the two moments needed for the positive and negative rows at
one site.  The chain multiplicities in the cornered comb are
\(P+1\) and \(N-1-P\), respectively.  Thus at a flat field the
test forms are biorthogonal to the rows at the same level.

For a small-curvature field, comparing each chain transport with
the cornered comb changes this calculation by \(O_L(\alpha)\).
There are at most sixteen nonzero entries in each row and eight
in each column of the same-level error.  Decreasing the curvature
threshold gives
\begin{equation}
 \Theta=D_0+N_0,\qquad
 \|D_0-1\|\le\frac1{48},\qquad
 \|N_0\|\le\frac{288}{L}<23,\qquad N_0^2=0.
 \label{mb:row-lift-triangular}
\end{equation}
Here \(D_0\) preserves the level and \(N_0\) raises it by one.
To check the last two assertions, a level-\(j\) test form can be
read at a higher level only by an outward row of level \(j+1\).
Its entry has size at most \(36L^{-3}\); at most \(8L^4\)
entries occur in one row and at most eight in one column.  The
Schur test gives \(288/L\).  Two consecutive upward transitions
would place \(\Omega_{j+2}\) within
\((1+(L-1)/2+L)\xi_{j+1}<L^2\xi_{j+1}\) of
\(\Omega_{j+1}^c\), contrary to separation.  Downward
transitions vanish because the supports avoid the outer sheet of
each cell.  This proves the asserted form and nilpotence.

\begin{proposition}[A lift in the energy norm]
\label{mb:row-lift-energy}
For sufficiently small \(\alpha\), there is a right inverse
\(Z^\sharp\) of \(Q^+_{\mathrm c}\) satisfying
\[
 Q^+_{\mathrm c}Z^\sharp=1,\qquad
 \mathcal N(Z^\sharp\omega)
       \le48^2c_Z\|\omega\|_2^2,
 \qquad c_Z=21\cdot18^{12}.
\]
Here
\[
 \mathcal N(A)=\|DA\|_{2,\sim}^2+\|D^*A\|_2^2
       +\|\mathsf L^{-1}A\|_2^2+\|Q^+_{\mathrm c}A\|_2^2.
\]
The first term is the usual inner-plaquette energy together with
\(\eta^2\sum_p(\sum_{b\subset\partial p}|A(b)|)^2\) over
plaquettes which meet an inner bond and are not inner plaquettes.
All constants are uniform in the depth and admitted torus.
\end{proposition}

\begin{proof}
Since \(N_0^2=0\), the exact factorization is
\[
 \Theta=(1+N_0)\{1+(1-N_0)(D_0-1)\}.
\]
The perturbation in braces has norm at most \(24/48=1/2\).
Its Neumann inverse therefore has norm at most two, and
\(\|\Theta^{-1}\|\le48\).  Notice that the inverse is
\[
 \Theta^{-1}=
 \{1+(1-N_0)(D_0-1)\}^{-1}(1-N_0).
\]
The same-level error must be included at a curved field.

For completeness, the energy estimate is local.  The profile
bounds and the curvature comparison give
\[
 |Z_c e|\le\tfrac12 18^5\xi_j^{-1}|e|,
 \qquad |\nabla_U Z_c e|\le18^6\xi_j^{-2}|e|.
\]
A cell contains \(N^4\) sites and has physical volume
\(\eta^4N^4=\xi_j^4\).  At most two rows of each direction
are assigned to one cell.  The curl, divergence and weighted
field terms consequently contribute at most
\(20\cdot18^{12}+2(18^5/2)^2\) times the squared row norm.
The level-zero terms obey the smaller bounds \(24\), \(16\)
and \(1\).  Disjoint cell cores prevent sums over different
levels in this calculation.  Finally
\(\|Q^+_{\mathrm c}Z\omega\|_2^2
 =\|\Theta\omega\|_2^2\le625\|\omega\|_2^2\).
Adding the estimates gives \(\mathcal N(Z\omega)\le
c_Z\|\omega\|_2^2\).  Take \(Z^\sharp=Z\Theta^{-1}\).
\end{proof}

This lift also gives the lower bound for the coarse Gram operator.
Suppose the real positive operator \(G_0\) is defined by the
Wilson Hessian, its specified gauge term and the retained row
penalty, and suppose its form obeys
\[
 \langle A,G_0^{-1}A\rangle
       \le C_E\mathcal N(A).
\]
For the inner Wilson Hessian one can take
\(C_E=1+19C_{4,\mathrm{svar}}\), where
\(C_{4,\mathrm{svar}}\) is the constant in its local
curvature-remainder estimate.  With
\(\gamma_Q=(48^2 C_Ec_Z)^{-1}\), Cauchy--Schwarz gives
\begin{align*}
 \|\omega\|_2^2
 &=\langle G_0^{-1}Z^\sharp\omega,
                G_0(Q^+_{\mathrm c})^*\omega\rangle\\
 &\le\gamma_Q^{-1/2}\|\omega\|_2
       \langle\omega,Q^+_{\mathrm c}G_0
                        (Q^+_{\mathrm c})^*\omega\rangle^{1/2}.
\end{align*}
Thus \(Q^+_{\mathrm c}G_0(Q^+_{\mathrm c})^*\ge\gamma_Q\).
This conclusion uses positivity of the specified fine operator;
the right inverse by itself does not prove that positivity.

Removing inward rows is restriction in the counting row space.
The lift and the Gram estimate restrict with the same constants.
Their use for the averaged constraint additionally requires the
comparison of its rows and endpoint transports with the
cornered-comb rows.  Their use for a current-corrected inverse
requires the corresponding form comparison.  Neither comparison
is an identification of this lift with the derivative of the
nonlinear background map.

For the scale convention in this section, the row pairing
is weighted by \(\mathsf m^4\), where \(\mathsf m(c)=\xi_c\).
If \(\widehat Q\) is the raw operator and \(Q=\mathsf m^{-1}\widehat Q\),
then \(Q^*=\widehat Q^{\mathsf T}\mathsf m^3\), with
\(\widehat Q^{\mathsf T}\) taken in the counting row pairing.
Thus multiplication by \(\mathsf m^2\), which identifies the
weighted row Hilbert space with the counting row space,
conjugates \(QG_0Q^*\) to
\(\mathsf m(\widehat QG_0\widehat Q^{\mathsf T})\mathsf m\).
This is the source of the two factors of the local spacing
in the target Gram lower bound. The constants from the
energy comparison must also have the uniformity required
for the intended application.


\paragraph{From comb rows to the averaging constraint.}
\label{mb:averaging-row-transfer}

The row lift above applies first to a specified linear family.
We describe the change to the straightened mean constraint,
including its error at the reference field.  Throughout this
comparison the reference configuration \(U_0\) is real, the
regions satisfy the separated-layer hypotheses, and all local
boxes used by an average are admitted boxes.  Constants depend
on \(L\), and the smallness and separation thresholds are fixed
independently of the number of layers.

Write \(Q_{\mathrm c}^+\) for the raw cornered-comb rows on all
retained positively oriented bonds and \(Q_{\mathrm c}\) for
their restriction to the interior and outward bonds.  The
pairing of raw rows is counting measure.  The fine-field
pairing has weight \(\eta^4\), and the site pairing on tops
has weight \(\xi_y^4\).  In these conventions the comparison
at level \(j\) reads
\begin{equation}
 Q_{\mathrm m,j}=T_j(Q_{\mathrm c,j}+E_j),\qquad
 Q'_{\mathrm m,j}=T'_j(Q'_{\mathrm c,j}+E'_j).
 \label{mb:mean-comb-comparison}
\end{equation}
Here \(T_j\omega(c)=\operatorname{Ad}(t_j(c_-))\omega(c)\)
and \(T'_j\mu(y)=\operatorname{Ad}(t_j(y))\mu(y)\), with
transports of \(U_0\) in the appropriate local boxes.  They
preserve the indicated real pairings.  The field-row error
vanishes outside the double block under \(c\), and
\begin{align}
 \sum_b|E_j(c,b)|&\le\kappa_F\alpha_0\xi_j,
 &\sum_{c\in\Gamma_j^+}|E_j(c,b)|
       &\le\kappa_F\alpha_0\eta^4\xi_j^{-3},
 \label{mb:mean-comb-moments}\\
 |E'_j\lambda(y)|&\le\kappa_F\alpha_0
                         \max_{B^j(y)}|\lambda|.
 \notag
\end{align}
At level zero these errors vanish.  At higher levels, the
curvature bound on the double block is
\(L^2\alpha_0L^{-2j}\).  The local comparison of the two
averages gives the first bound and the corresponding bound
for each kernel entry.  A fine bond belongs to at most eight
such rows at a fixed level, which gives the second bound.
In level-\(j\) units the field is \(\xi_jA\), so the row
and column bounds become \(\kappa_F\alpha_0\) and
\(\kappa_F\alpha_0(\eta/\xi_j)^4\), respectively.

The positive fine operator used for the comb Gram estimate is
\begin{equation}
 G_{\mathrm c}^{-1}
  =\Delta^{\mathrm W}_{U_0}+B+D R_{\mathrm c}D^*
                              +Q_{\mathrm c}^*Q_{\mathrm c},
 \qquad
 R_{\mathrm c}=\operatorname{proj}_{\Delta_{U_0}
                                       \ker Q'_{\mathrm c}}.
 \label{mb:comb-fine-operator}
\end{equation}
The boundary operator \(B\) is positive and symmetric, its
kernel joins only bonds of a common non-inner plaquette, and
its row sums are at most \(12\eta^{-2}\).  The projection in
\eqref{mb:comb-fine-operator} is taken in the fine site pairing.
The scale-weighted comb propagator estimates give
\begin{equation}
 M_{\mathrm c}=Q_{\mathrm c}G_{\mathrm c}
                    (Q_{\mathrm c})^*\ge\gamma,
 \qquad
 H_{\mathrm c}=G_{\mathrm c}(Q_{\mathrm c})^*
                    M_{\mathrm c}^{-1},\qquad
 Q_{\mathrm c}H_{\mathrm c}=1,
 \label{mb:comb-natural-lift}
\end{equation}
and the scale-weighted bounds
\[
 \|H_{\mathrm c}\beta\|_X\le K_1\|\beta\|_\infty,
 \qquad
 \|H_{\mathrm c}^*f\|_\infty\le K_1|f|_{(-3)}.
\]
These are estimates for the operator in
\eqref{mb:comb-fine-operator}.
In the nonlinear estimates below,
\(|f|_{(-s)}=\sup \mathsf L^s|f|\) uses the scale assigned
to the fine bond or plaquette, and
\(\|A\|_X=\max\{|A|_{(-1)},|\nabla A|_{(-2)}\}\).
The version using the scale of the top of a cell is equivalent
with factors at most \(L^s\).  Enlarge \(K_1\) by these
fixed factors when using the linear estimates in these norms.

Let \(\mathcal Q^{\natural}(A)\) be the logarithmic
straightened mean on the interior and outward rows, with
\(K_A=(\mathcal Q^{\natural})'(A)\).
With the corresponding fixed transport \(T\), one has
\begin{equation}
 T^{-1}\mathcal Q^{\natural}(A)
       =Q_{\mathrm c}A+C(A),\qquad
 T^{-1}K_A=Q_{\mathrm c}+C'(A),\qquad
 C'(0)=E.
 \label{mb:nonlinear-comb-comparison}
\end{equation}
In particular, the remainder contains a linear term at a
curved reference field.

Here are the quantitative estimates in these row coordinates.  Let \(c_4,C_2\) be the local logarithm and mean
kernel constants, let \(\kappa_{F,\mathrm{slin}}\) be the
constant in the corresponding local comparison, and put
\[
 C_3^{\natural}=104/c_4,\qquad
 \epsilon=\kappa_{F,\mathrm{slin}}L^2\alpha_0
                         +4(C_2+1)Lr,\qquad
 \epsilon^*=\kappa_{F,\mathrm{slin}}L^2\alpha_0
                         +C_3^{\natural}Lr.
\]
For \(|A|_{(-1)},|A_i|_{(-1)}\le r\le c_4/(4L)\),
\begin{align}
 \|C(A)\|_\infty&\le L\epsilon|A|_{(-1)},&
 \|C'(A)a\|_\infty&\le L\epsilon|a|_{(-1)},
 \label{mb:nonlinear-row-bounds}\\
 |C'(A)^*\omega|_{(-3)}&\le4L^3\epsilon^*\|\omega\|_\infty,&
 |(C'(A_1)^*-C'(A_2)^*)\omega|_{(-3)}
 &\le2L^4C_3^{\natural}|A_1-A_2|_{(-1)}\|\omega\|_\infty.
 \notag
\end{align}
The first line also gives the Lipschitz bound for \(C\).
For the transpose estimates, a fine bond of scale \(\xi_n\)
occurs only in rows of levels \(n-1,n,n+1\); their weights
sum to at most \((L^3+2)\xi_n^{-3}\).  The derivatives of the local mean kernel bound the
column sum.  Applying Cauchy's estimate in the local field ball
to that column norm gives the last bound in
\eqref{mb:nonlinear-row-bounds}.  In particular these
estimates do not introduce a factor equal to the depth.

The right inverse now transfers on the small chart.  If
\(K_1L\epsilon\le1/2\), then
\[
 T^{-1}K_AH_{\mathrm c}=1+C'(A)H_{\mathrm c},\qquad
 \|C'(A)H_{\mathrm c}\|_{\infty\to\infty}\le1/2.
\]
Its Neumann inverse therefore gives
\begin{equation}
 H_A=H_{\mathrm c}(1+C'(A)H_{\mathrm c})^{-1}T^{-1},
 \qquad K_AH_A=1,
 \qquad \|H_A\beta\|_X\le2K_1\|\beta\|_\infty.
 \label{mb:mean-local-right-inverse}
\end{equation}
The dual calculation controls the multiplier.  If
\(q=4K_1L^3\epsilon^*\le1/8\), then
\[
 S_A=1+H_{\mathrm c}^*C'(A)^*,\qquad
 \|S_A^{-1}\|_{\infty\to\infty}\le8/7,
 \qquad H_{\mathrm c}^*(T^{-1}K_A)^*=S_A.
\]
Consequently, whenever
\(F'(A)+DR_{\mathrm c}D^*A=K_A^*\omega\), its multiplier
is uniquely determined by
\begin{equation}
 T^{-1}\omega
  =S_A^{-1}H_{\mathrm c}^*(F'(A)+DR_{\mathrm c}D^*A).
 \label{mb:mean-multiplier-readout}
\end{equation}
This argument does not assume a bound for the multiplier
before solving for it.  Here \(F(A)\) is the inner Wilson
action evaluated at \(e^{i\eta A}U_0\).

There is a further coordinate change when the datum is
written with its endpoint gauges.  The seam coordinate
change \(\Xi\) gives \(K_A^{\flat}=D\Xi\,K_A\).
For the actual reduced datum, write \(n(A)=x(A)^{v[A]}\).
Its derivative is
\[
 dn[A]a=[dx[A]a]^{v[A]}+\operatorname{rot}_{n(A)}(\mu),
 \qquad \mu=i^{-1}(dv[A]a)v[A]^{-1}.
\]
Gauge covariance gives
\(dn[A]\delta_\lambda A=\operatorname{rot}_{n(A)}
(\lambda|_{\Lambda_*})\), where \(\delta_\lambda A\)
is the infinitesimal bare gauge variation and
\(\operatorname{rot}\) denotes the infinitesimal action on
the datum.  To attain a prescribed tangent
datum, first choose \(a\) using the surjectivity of the
straightened rows and then choose \(\lambda=\mu\) on the
tops.  The vector \(a-\delta_\lambda A\) attains the datum.
This proves surjectivity of the finite-dimensional datum map.
A uniform bound for its right inverse additionally
uses the estimates for this gauge extension.

The first line of \eqref{mb:nonlinear-row-bounds} holds for
complex increments around the fixed real reference field.
The transpose bounds hold with an additional factor two,
using the complex bilinear transpose.  Extension of a real
linear estimate to complex vectors is distinct from allowing
the reference field and the independent current to vary on
the full complex domain of the propagator statements.

Thus the local averaging constraint has a bounded right
inverse on its specified chart.  To obtain the target Gram
estimate for the corrected rows and current-completed fine
operator, one must still match the row coordinates and
pairings, establish the energy bound for the resulting lift,
and compare that fine quadratic form with
\eqref{mb:comb-fine-operator}.  The construction in
\eqref{px:actual-layered-row-lift} addresses the actual row
energy; the row comparison above supplies its source
interpretation when the averaging conventions agree.
Neither surjectivity in the supremum norm nor
\eqref{mb:comb-natural-lift} alone performs the remaining
quadratic-form comparison.

(e) \textbf{Gauge algebra without Gaussian integrals.} On a pure gauge, take the derivative of the covariance of the straightened average; the site factors then convert point values into block means: $QD=\bar DQ'$. Consequently the rows annihilate $D\ker Q'$. Combine this with the invariance of the action to second order, which is exact for the inner plaquettes; this accounts for that Hessian. The result is $\mathsf G^{-1}D\lambda=D\Delta_{\mathbf U}\lambda-\Delta'_\pi[\mathbf J-J(\mathbf U)]D\lambda$ on $\ker Q'$. The propagators are taken at pairs whose current is an independent variable carrying a bound of its own, as in the spaces of \textbf{109}, p. 262. The reason is that bounds on the plaquettes, on the potential and on the first differences do not control $J(U)$ uniformly in the level: a staggered field satisfies all of them with $|J(U)|\ge\alpha\eta^{-1}$. In \textbf{99}, 389, by contrast, $J=J(U)$, and a condition on second differences is imposed in addition ((3.36), (3.38)). The price of a free current is the defect in 3. Uniformity in the level holds for the following reason. Row norms and column norms are both submultiplicative, and for a single step they lie near $L$ and $L^{-3}$. The relative errors $O(\alpha_0(L^n\eta)^2+\alpha_1L^n\eta)$ of the steps form geometric series that are dominated by the top level.

% Earlier Part I, Section 4.6, two substitutions and common-region
% cancellation (28276--28431, 30131--30168, 30719--30850, 31275--31322).
\paragraph{Matching the local factors of the background propagator.}
For the native propagator construction use the same candidate cubes,
trace-local bump rule and rules for modifying the cover in the two
nests. Let
\(G_0=\sum_\square h_\square G_\square h_\square\) be the
native vector parametrix. Divide each candidate bump by the
square root of the sum of all retained and inserted bump
squares at that point, so that \(\sum_\square h_\square^2=1\).
Use local inverses whose levels agree with the global levels
on the tested stencils. For a level-\(j\) cube this agreement
holds when its twice-enlarged cube lies in
\(\Omega_j\setminus\Omega_{j+2}\). Separation supplies
the inner margin. At the outer face one has to check the
same levels or use the corresponding boundary model.
Write \(D\) for the covariant site derivative and \(P\)
for the projection occurring in the gauge term. Write
\(\Delta_aG_0=I-R\). The four summands of \(R\) at a cube are
\[
\begin{aligned}
 &K(h_\square)G_\square h_\square,\qquad
 (1-\zeta_\square)DPD^*h_\square G_\square h_\square,\\
 &\zeta_\square(DPD^*-DP_\square D^*)h_\square G_\square h_\square,
 \qquad
 \zeta_\square[DP_\square D^*,h_\square]G_\square h_\square .
\end{aligned}
\]
Here \(K(h_\square)\) is the local differential and averaging
commutator. The last term uses the local projection and the
first difference of the bump. In the two occurrences of the
global projection substitute
\[
 P=G'Q'^*(Q'G'^2Q'^*)^{-1}Q'G'.
\]
For \(G'\), the scalar parametrix with the same square
partition gives the series of local inverse and commutator
factors. For the site Gram inverse, glue its local inverses.
The resulting error consists of a tail outside the cut,
a difference of scalar propagator squares and a local
commutator. Insert the scalar expansion with the cover
modified at the cut. The difference of squares is expanded
as \((A-B)A+B(A-B)\), with the corresponding local
scalar propagators \(A,B\). In the vector error with
\(\zeta_\square\), the term containing no small commutator
or separated supports cancels against the local projection
term. To get this cancellation, replace the inserted bumps
by \(\zeta_\square\), compare the modified local inverses
with the original ones, and commute the cutoffs to the
right. The errors have a cutoff difference, a localized
inverse difference or a commutator. The term left after
these operations uses \(\zeta_\square h_\square=h_\square\).
Consequently \(G_0\sum_{m\ge0}R^m\), with these substitutions,
has only local inverse, projection and commutator factors
\cite[(3.87)--(3.108)]{B85c}.

For comparison on a common region, fix one list of
possible words before imposing either nest. Take as the
nominal footprint the union of the constituent cubes
with five layers of padding at their own scales.
Admission of the word also reads the geometric
decisions that modify the cover. In top units the index
rule reads the region data less than \(2M\) beyond the
footprint, since it tests the neighboring cube in deciding
which cover to use. Enlargement by \(2M\) is also
sufficient for the coefficients, including the bumps and
the full field reads. If the outer averaging and
derivative factors read entirely inside the footprint,
the field reads require only the original footprint.
Use this enlarged region for comparison. Equality of the
two nests and their region data there gives the same
cover decisions, local sequences and cutoffs. The local
operators are therefore the same functions of the common
background reads. At the same background on those reads
their values are equal. Thus all the words whose enlarged
reads lie in the common region cancel, term by term.
Under the native regularity and separation assumptions
the absolute walk majorant permits this cancellation
inside the sum. It estimates the remaining terms by
the first escape from the common region, with the
geometric padding subtracted before measuring the escape.
An operator formed with an additional region restriction
uses its required margin to that region's boundary.
For the history comparison, one neighboring cell supplies
this geometric padding when \(\check R_s\ge2\).
The trace data must record the full nest inside each
component meeting the filed domain. The clean neighboring
cells have trivial trace when they belong to the
regular region of every nest in force. Where the padding
reaches a different component, the comparison uses its
nest and region trace on the cells met by the padding.
For an application to a switched or compressed inverse
one must also match its actual masks and the data read
by its additional vertices.

(f) \textbf{Unfolding and parametrix.} With auxiliary unknowns 

\begin{gather*}
\FitMath{\Psi=(A;v,w,\sigma;p,q,r;\nu)}
\end{gather*}

, the operators $R$, $\Delta'_\pi$, $\mathsf S^{\mathsf T}\Delta^{(2)}\mathsf S$ and the constraint are combined into a single system $\mathbb K$ of finite range. The inverse of this system has the propagators as blocks. Let $\varkappa$ be an $M_1$-cube; these cubes have side $M_1$ in the local unit and tile $\Omega_0$, whereas the $\square$ of (b), (d) are the cubes of the flat expansions. On a regular interior piece, in the gauge of the class, $\mathbb K_\varkappa$ equals the flat system plus $O(\alpha_0+\alpha_1+\gamma_0)$, and this perturbation need only be small compared with flat constants. A piece meeting the true bare boundary instead uses the direct massive inverses of the actual compressed Wilson Hessian and the covariant site Laplacian. In bare units the real operators $1+\eta^2\Delta_Y(U)$ and $1+\eta^2\Delta_U^\circ$ have uniform positive lower bounds, respectively $1/2$ and $1$, and uniformly bounded finite-range kernels. A small complex perturbation preserves their inverses and exponential decay; local elimination then supplies the remaining slots. This boundary construction uses the actual incident inner plaquettes and requires no gauge across the boundary. Let $\theta_\varkappa$ be a partition of unity, let the cut-offs $\tilde h_\varkappa=1$ near $\operatorname{supp}\theta_\varkappa$ be such that $\mathbb K=\mathbb K_\varkappa$ on $\operatorname{supp}\tilde h_\varkappa$, and let $\delta$ denote the rate of the local inverses. Then

\begin{gather*}
\FitMath{\mathbb K\sum_\varkappa\tilde h_\varkappa\mathbb K_\varkappa^{-1}\theta_\varkappa=1-\sum_\varkappa\mathbb R_\varkappa,} \\[3pt]
\FitMath{\mathbb R_\varkappa=-[\mathbb K_\varkappa,\tilde h_\varkappa]\mathbb K_\varkappa^{-1}\theta_\varkappa=O\big(M_1^{-1}e^{-c\delta M_1/L}\big),}
\end{gather*}

where the power comes from the commutator, and the exponential from how far $\operatorname{supp}\theta_\varkappa$ lies from $\operatorname{supp}\nabla\tilde h_\varkappa$. This bound outweighs the $O(1)$ neighbours. Three features keep the perturbation and the commutators from losing a derivative: scaling weights, one component of the source that admits a divergence, and the absence of any derivative on $p$.

(g) \textbf{Trees.} The system $\mathbb K$ depends affinely on the non-local coefficient $\mathsf H^{\mathsf T}\mathbf J$. The part of walk length zero stays in the local operators. The remainder, of size $O(\gamma_0M_1^{-1}e^{-cM_1/L})$, forms a vertex that carries a walk of $\mathsf H$. Thus vertices are small because of $M_1$ and not because of the field, and each term is indexed by a trunk with branches. When the cells are listed depth first, the number of such trees grows geometrically: for each first cell, there are $C^m$ tree walks of length $m$ starting from it. Both this count and the size of $M_1$ needed for the sum over trees to converge coincide with the ones in 1c.9. In 1c.7 and 1c.8 a branch is listed as soon as its cut is reached, and the trunk is resumed afterwards. For the data of 1c.9 the listing first continues at every cell along the trunk (within a branch, along the branch). So the whole trunk comes first, and then the branches come in the reverse order of their cuts, beginning with the last cut and ending with the first. Each of the two orders is a depth-first order. The order of listing changes neither the cells nor their parents, and so it changes neither the hull, nor the tree length, nor the bounds.

(h) \textbf{Multiplier and positivity.} Adding the equation ``$QA=$ fluctuation field'' brings in $(Q\mathsf G_1Q^{\mathsf T})^{-1}$ as a Lagrange multiplier. Consequently $C^{(k)}(Y)$ occurs as a block in the inverse of a system of finite range. In the flat case, positivity follows from $\sum_{p'}|\bar\partial QA(p')|^2\le\eta^4\sum_p|\partial A(p)|^2$, with $\bar\partial$ taken on $\mathbb T^{(k)}$, together with a definite form on double blocks. In a background, a homotopy to $\mathbf 1$ preserves $C_\varkappa>0$ for the local inverses. The eigenvalues $z$ of $\chi_Y\mathsf A\chi_Y\sum h_\varkappa C_\varkappa h_\varkappa$ satisfy $|1-z|\le CM_1^{-1}$.

(i) \textbf{Length against parameters.} A tree of length $t$ comes close to at most $2^4\max\{1,Ct/\ell_{\mathsf P}\}$ cubes $\square$ of the partition of side $\ell_{\mathsf P}$ in statement 7; these are not the cubes $\square$ of (b), (d). Hence $\bar\kappa\#\mathsf P(\omega)\le2^4\bar\kappa+(C\bar\kappa/\ell_{\mathsf P})|\Gamma|_{\mathbf B}$ for every tree $\Gamma$ admitted in $d_{\mathbf B}(\omega;y,y')$, provided that $X_0$, which $\Gamma$ need not meet, lies within $O(1)$ top units of $y$ (no other term enters a functional near $y$, since the left argument of the kernel of $\mathsf O_\omega$ touches $X_0$). Consequently, for $|s(\square)|\le e^{\bar\kappa}$ and $\delta\ge C\bar\kappa/\ell_{\mathsf P}$,

\begin{gather*}
\FitMath{|s^\omega|\,e^{-\delta d_{\mathbf B}(\omega;y,y')}\le e^{2^4\bar\kappa}\,e^{-(\delta-C\bar\kappa/\ell_{\mathsf P})d_{\mathbf B}(y,y')}.}
\end{gather*}

Each cube beyond the first $2^4$ adds to the tree a length $\ge c\ell_{\mathsf P}$. Since $\ell_{\mathsf P}$ is chosen last, $e^{-c\delta\ell_{\mathsf P}}$ outweighs $e^{\bar\kappa}$. The factor $e^{2^4\bar\kappa}$ is left over and cannot be avoided, as a cell at a corner of the grid shows. The displayed bound uses only the total-incidence count: $|s^\omega|\le e^{\bar\kappa\#\mathsf P(\omega)}$, followed by $d_{\mathbf B}(\omega;y,y')\ge d_{\mathbf B}(y,y')$. Its parameter loss is the stated reduction of the spatial decay rate. The estimate is Ba{\l}aban's (\textbf{116}, p. 5, after (1.11)); in his setting as well, the exponential by itself does the paying, and $e^{16\bar\kappa}$ remains. He keeps the factor $M_1^{-|\omega|/2}$ in reserve together with the exponential (p. 3, after (1.7)). This factor is not used up on the parameters, and here it is still at our disposal for the factor $|\omega|+1$ that occurs below. Since the uncut hull, which is his localisation domain, is connected, $\mathsf P(\omega)$ is connected through walls; this yields the splitting. The carrier is needed because the fields are given only on $\Omega_0$, and it need not be connected when $\Omega_0\ne T_\ell$.

The properties of single terms are produced by construction; they are not deduced. The sum is itself covariant and real. Therefore the Haar mean of the terms, made symmetric under the conjugation fixing the real pairs, is once more an expansion of the sum. By convexity it satisfies the same bounds, provided that the family of bounds is invariant under the real gauge transformations and under this conjugation; the family of localised bounds has this invariance. A single assignment for all tori is obtained in three steps: a mean over translations; a push-forward from covering tori; and a passage to the limit by Montel's theorem, along a subsequence fixed once, at the given top level and units. In the push-forward, a walk wrapping around the small torus gains at most $|\omega|+1$ blocks; the weight therefore loses a factor $|\omega|+1$, and nothing else is lost. The family is made normal by an entrywise bound, which is one and the same number for all tori with room. This number is allowed to depend on the level, since the entries of $(Q\mathsf GQ^{\mathsf T})^{-1}$ are $O(\eta^{-2})$, and none of the norms, rates, weights or sums of weights involves it. If a family of terms is already carried by all shifts and compatible with all coverings, it is not changed.

\paragraph{The coupled map \eqref{eq:propagators-II}.}
The datum in this map is the raw row logarithm \(B'\).
The fine pairing retains its physical cell weights, while the
raw row pairing is counting measure. At \(s=1\), define
\(\hat Q=\mathsf mQ\), \(H=\mathsf H\mathsf m^{-1}\),
\(S=(\hat Q\mathsf G\hat Q^{\mathsf T})^{-1}\),
\(H_0=\mathsf G_0\hat Q^{\mathsf T}
(\hat Q\mathsf G_0\hat Q^{\mathsf T})^{-1}\), and
\(\tilde G=\mathsf G_0-H_0\hat Q\mathsf G_0\).
Here \(S\) acts on raw rows and is distinct from the fine
projector \(\mathsf S\). The multiplication operator
\(\mathsf a\) is fixed by
\(\hat Q^{\mathsf T}\mathsf a\hat Q=Q^{\mathsf T}aQ\);
for the row volumes \(\xi_c^4\) and
\(a(c)=a_0\xi_c^{-2}\), it is \(a_0 I\).
All transposes in these formulas use the indicated two pairings.

Let \(F(T)=A_{\Omega_0}(e^{i\eta T}\mathbf U)\), using its
regular holomorphic continuation, and define
\(\mathcal C(T)=\mathcal Q(T)-\hat QT\) and
\(\mathcal N(T)=F'(T)-F'(0)-\Delta^{\rm W}_{\mathbf U}T\).
Both remainders vanish to second order at zero on their
actual admitted balls. The independent \(\mathbf J\) is not
set equal to \(F'(0)\). The transpose of
\(\mathcal C'_T\) maps a raw row multiplier \(\lambda\)
to a fine force. For \(s\ne1\), the four maps
\(H_0(s),\tilde G(s),H(s),S(s)\) are their specified
walk-data sums, not inverses recomputed using identities
that may fail after decoupling.

The map on a fine potential \(T\) and raw row multiplier
\(\lambda\) is
\refstepcounter{equation}
\begin{gather*}
\FitMath{\begin{aligned}
T={}&H_0(s)[B'-\mathcal C(T)]\\
 &-\tilde G(s)[\mathcal N(T)-\mathcal C'_T{}^{\mathsf T}\lambda],\\
\lambda={}&H(s)^{\mathsf T}\mathbf J
       +[S(s)-\mathsf a][B'-\mathcal C(T)]\\
 &+H(s)^{\mathsf T}[\mathcal N(T)-\mathcal C'_T{}^{\mathsf T}\lambda].
\end{aligned}}\tag{\theequation}\label{eq:propagators-II}
\end{gather*}
Thus \(H_0,\tilde G\) are the zero-current slots, whereas
\(H,S\) use the current-dependent propagator.
The centred multiplier variable is
\(\lambda-H(s)^{\mathsf T}\mathbf J\).
With these raw-row pairings, the physical multiplier in the
identified case is \(\mathsf m^{-3}\lambda\).
\Cref{dnr:coupled} proves convergence and uniqueness on the full
first-order regular-pair domain and enlarged parameter discs,
with the additional independent-current smallness required by
this coupled system. Its decorated maps are the actual carrier
sums of \cref{dnr:parameter-families}.

(j) \textbf{Operators as data.} The principle goes back to Ba{\l}aban, who observes that in the equations for the response ``we can replace the propagators by arbitrary operators having the same regularity properties and satisfying the same bounds'' (\textbf{116}, p. 3, (1.5)). In its published form, however, the second equation relies on the identities $RD^{\mathsf T}\mathsf H=0$ and $\mathsf G^{-1}\mathsf H=Q^{\mathsf T}(Q\mathsf GQ^{\mathsf T})^{-1}$, and therefore on their consequences $Q\mathsf H=1$ and $\Delta_\pi\mathsf H=Q^{\mathsf T}[(Q\mathsf GQ^{\mathsf T})^{-1}-a]$ ($\Delta_\pi=\Delta^{\rm W}_{\mathbf U}-\Delta'_\pi$). None of these is available for the sums $\sum_\omega s^\omega\mathsf O_\omega$: on the cubes around a row where $s$ vanishes, every term is absent, so that $Q\mathsf H(s)$ equals $0$ there and not $1$. The first identity, moreover, holds only for the pairs on shell, even at $s\equiv1$. Consequently no identity enters the argument. The identities are applied earlier, where they hold. With their help the gradient of the non-quadratic part of the action is put into a shape from which its second-order term $\Delta_\pi\mathsf H$ is absent, and this expression \emph{defines} $\mathsf W_s$ for every $s$ and every pair (it is a gradient only at $s\equiv1$ and on shell). From then on only bounds are used, in the norms

\begin{gather*}
\FitMath{\|A\|_\varpi=\sup_y\varpi(y)\max\{\xi_y|A|,\,\xi_y^2|\nabla_{\mathbf U}A|\},} \\[3pt]
\FitMath{|f|^\sharp_\varpi=\sup_y\varpi(y)\,\xi_y^3|f|,} \\[3pt]
\FitMath{\varpi(y)\le e^{\theta d_{\mathbf B}(y,y')}\varpi(y'):}
\end{gather*}

the first measures a potential together with its first derivative, the second a force. The operator $\mathfrak G(s)$ takes the second norm back to the first, with norm $\hat{\mathsf b}\sim e^{2^4\bar\kappa}$, which is the size of every one of the operators; this closes the fixed-point argument.

\paragraph{The reduced force on the complex ball.}
Use raw row coordinates in this calculation. Write
\(\widehat Q=\mathsf m Q\) and let \(H_s\) be the walk sum
for the raw lift \(\mathsf H\mathsf m^{-1}\). Let \(S_s\) be the
row inverse sum in these coordinates. At \(s=1\) it equals
\((\widehat Q\mathsf G\widehat Q^{\mathsf T})^{-1}\), where the
transpose uses the raw row pairing. If \(\lambda\) is a raw
multiplier, the multiplier in the row-volume pairing is
\(\mathsf m^{-3}\lambda\). Indeed the row amplitude is multiplied
by \(\mathsf m\) and the row volume is \(\mathsf m^4\), so the
corresponding per-length transpose is
\(\widehat Q^{\mathsf T}\mathsf m^3\).
Put \(P_s=S_s-\mathsf a\) with the row mass \(\mathsf a\) used
in \eqref{eq:propagators-II}.

Let \(\mathcal C_2\) be the quadratic part of the raw remainder
\(\mathcal C\). Given an auxiliary potential \(T\), solve
\[
 X(T)=\mathcal C(T-H_sX(T)),\qquad
 Y(T)=T-H_sX(T),\qquad
 E_T=(\mathrm{d}X(T))^{\mathsf T}.
\]
The bounded local row map and the quadratic estimate on
\(\mathcal C\) make this a contraction in \(X\) on a ball of
radius \(C r^2\) when \(\|T\|\le r\le c/\hat{\mathsf b}\).
Indeed the Lipschitz constant is \(C\hat{\mathsf b}r\), and
\(\|Y(T)\|\le 2r\) after decreasing \(c\).
Differentiating the equation for \(X\) gives
\[
 E_T=\mathcal C'_{Y(T)}{}^{\mathsf T}
       (1+H_s^{\mathsf T}\mathcal C'_{Y(T)}{}^{\mathsf T})^{-1}.
\]
In particular \(X\) starts with \(\mathcal C_2(T)\).
In the following formulas \(\mathcal V\) is the quadratic
remainder of the action gradient, restricted to the force
bonds of the local calculation as in \(\mathcal N\) above.
Write \(\phi=H_s^{\mathsf T}\mathbf J\). The formula used with
the walk data is
\begin{equation}\label{eq:elab-reduced-force}
\begin{split}
 \mathsf W_s(T)
  ={}&-(E_T-(\mathrm{d}\mathcal C_2(T))^{\mathsf T})\phi
      -\widehat Q^{\mathsf T}P_s X(T)\\
    &-E_T P_s\widehat Q T+E_T P_s X(T)
      +\mathcal V(Y(T))-E_T H_s^{\mathsf T}\mathcal V(Y(T)).
\end{split}
\end{equation}
For the reduction take the full carrier at \(s=1\),
with the on-shell identity \(RD^{\mathsf T}H_{s=1}=0\). Use
\(\Delta_\pi H_{s=1}=\widehat Q^{\mathsf T}P_{s=1}\),
where \(H_{s=1}\) is the raw lift from \(\mathsf G\). Thus \(H_{s=1}^{\mathsf T}\Delta_\pi
H_{s=1}=P_{s=1}\). The terms in the differentiated reduced action
come from
\[
 -(\phi,X-\mathcal C_2(T))
 -(\widehat Q T,P_{s=1}X)+\tfrac12(X,P_{s=1}X)+V_0(Y),
 \qquad V_0'=\mathcal V
\]
on the full physical domain. Symmetry of \(P_{s=1}\) gives the
fourth term in \eqref{eq:elab-reduced-force}, and the chain
rule for \(Y\) gives its last two terms. For general \(s\)
\eqref{eq:elab-reduced-force} is evaluated directly using
the designated operator sums.

For the estimate, the two small arguments \(T_i\) have
\(\|T_i\|\le r\). Set \(d=\|T_1-T_2\|_\varpi\). In the raw
row norm the fixed point and its derivative satisfy
\[
 |X_i|\le Cr^2,\quad |X_1-X_2|_\varpi\le Crd,\quad
 |E_{T_i}\psi|^\sharp_\varpi\le Cr|\psi|_\varpi,\quad
 |(E_{T_1}-E_{T_2})\psi|^\sharp_\varpi\le Cd|\psi|.
\]
Here \(|\cdot|\) on rows is the unweighted supremum and
\(|\cdot|_\varpi\) uses the weight at the assigned top.
The first two bounds follow from the row contraction.
The two derivative bounds use the local estimate for the
transpose of the average differential and the resolvent
identity for \(1+H_s^{\mathsf T}\mathcal C'_Y{}^{\mathsf T}\).
Each difference in the resolvent identity contains one
difference of \(\mathcal C'_Y\), bounded by \(Cd\) with its
weight on that factor. Both resolvents are bounded on
the same ball since \(C\hat{\mathsf b}r\) is small.

The second and third terms of \eqref{eq:elab-reduced-force}
have differences bounded by \(C\hat{\mathsf b}rd\).
In the fourth the additional small factor \(r\) comes
from \(E_T\) acting on the quadratic \(X\), or from \(X\)
when the difference falls on \(E_T\). The fifth term is
bounded by \(Crd\) by the local action remainder estimate
and \(\|Y(T_1)-Y(T_2)\|_\varpi\le 2d\). For the sixth,
the undifferenced action remainder is \(O(r^2)\), its
difference is \(O(rd)\), and \(H_s^{\mathsf T}\) costs
\(C\hat{\mathsf b}\). These contributions are therefore
bounded by \(C\hat{\mathsf b}r^2d\).

For the first term let \(r_0=c/\hat{\mathsf b}\) be the
radius of the full holomorphic ball and assume \(0<r<r_0\).
Put \(N(T)=E_T-(\mathrm{d}\mathcal C_2(T))^{\mathsf T}\).
The function \(g(z)=(N(zT_1)-N(zT_2))\phi\) is holomorphic
on \(|z|<r_0/r\), with \(g(0)=g'(0)=0\), because
\(\mathcal C_2\) is the quadratic term of \(X\).
The derivative difference bound gives
\(|g(z)|^\sharp_\varpi\le C|z|d\,|\phi|\).
Apply Schwarz's lemma to \(g(z)/z\), after pairing with a
dual functional of norm one, and take the supremum over
the dual unit ball. This gives
\[
 |g(1)|^\sharp_\varpi
   \le C(r/r_0)d\,|\phi|
   \le C\hat{\mathsf b}^{\,2}\gamma_0 rd.
\]
The resulting bound is \(C\hat{\mathsf b}
(1+\gamma_0\hat{\mathsf b})rd\). The coefficient and operator
bounds here are on the local class used for the walk sums,
with the force bonds, row supports, and the neighboring
plaquettes inside its enlarged carrier. In applying this
calculation to a history, membership of the actual restricted
fields in that common class must still be checked.


The Lipschitz constant of the fixed-point map is $\hat{\mathsf b}\cdot\hat{\mathsf b}^2\cdot\hat{\mathsf b}\sup|B'|$, with contributions from $\mathfrak G$, $\mathsf W_s$ and $\mathsf H_1$, and this is why the radius is $\hat{\mathsf b}^{-4}$. Decay amounts to choosing the weight $\varpi=e^{\theta d_{\mathbf B}(\cdot,y')}$. When $s\not\equiv1$, bounds of first order on the operators do not yield a bound on $\Delta_{\mathbf U}\mathbf H$ independent of the level; an oscillating kernel belonging to the class of data is an example. The reduced current takes its place. It is built from the multiplier, and in it the second-order part of the gradient occurs only in paired form, $(D^{\mathsf T}DA,\mathsf He_c)=(DA,D\mathsf He_c)$, where $e_c$ is an element of $\mathfrak g^c$ placed at the row $c$. As a result, only first-order norms occur.

% and Proposition 79.1(c).
\paragraph{Restriction of the response.}
A surviving walk has its full parameter support in one
wall component of \(\{s\ne0\}\). It gives the same summand
when all parameters outside that component are set to zero.
For the nonlinear part, use the two nested collars in the
definition of the local class. The potential and datum
slots of the walk operators are supported inside the first
collar boundary, at distance greater than the reach of
the local average and force from the other components.
Thus an iterate supported there has its averages, action
remainder and transposed differential computed from data
in the enlarged collar. The next iterate is again supported
in the same slots. Starting at zero, the iteration for the
union is the sum of the component iterations. Uniform
contraction proves the same identity for the limits and
allows comparison of settings by checking the walk terms
and the local coefficients on these collars. A match of the
decay rates alone would not identify the two responses.

For identification with the actual background, use all
parameters equal to one at a real critical reference with
its actual current. The reduced datum is obtained from
the physical datum by the seam and node factors of the
straightened average. Let \(C_{\mathsf b}\) bound that map, \(r_{\mathsf b}\)
be its chart radius, \(r_B\) the radius of the holomorphic
background chart, and \(\varepsilon_{\mathrm D}\) the
response datum radius at \(s=1\). Take the physical datum
in the ball of radius
\(\min\{r_B,r_{\mathsf b},\varepsilon_{\mathrm D}/C_{\mathsf b}\}\).
The reference pair must belong to the common regular
class at the admitted response radii. The holomorphic
background in the normalized node gauge must lie in the
response uniqueness ball and solve the constraint and
slice equations there.
The multiplier argument below and uniqueness identify it
with the response to the reduced datum. On the constant
sequence the two datum coordinates agree. The identification
on a history requires the same determining rows and the
same current correction in both constructions.

\paragraph{The coupled response in raw row coordinates.}
For \eqref{eq:propagators-II} the walk sums in the four
slots are \(H_0(s),\tilde G(s),H(s)^{\mathsf T}\), and
\(S(s)-\mathsf a\). The mass subtraction is a local
diagonal term and does not carry a decoupling parameter.
In the estimate of these slots, the potentials and the
row outputs of a datum at \(y'\) carry the factor
\(\xi_y/\xi_{y'}\). The potential from a force supported
over \(y'\) carries \((\xi_y/\xi_{y'})^3\). The two bounds
have the same exponential decay in \(d_{\mathbf B}(y,y')\)
and hold for partial sums. To see how the normalization
enters, the datum at \(y'\) is a raw row logarithm,
whereas the output potential and its first difference
are measured in the units at \(y\). Transpose the
\(H\)-kernel with its fine pairing and the counting
row pairing before estimating the multiplier slot.
To absorb a ratio of power \(p\le3\), use the distance
lower bound \(\rho(|j-j'|-1)_+\) for tops on levels
\(j,j'\), with \(\rho=R_1\mathfrak M\). Then the ratio
costs at most \(L^p e^{p(\log L)d_{\mathbf B}/\rho}\).
For the kernel rate \(3\dot\delta_0/16\), the separation
floor \(\dot\delta_0\rho\ge512\log L\) leaves decay at
rate \(11\dot\delta_0/64\). Summing with the admissible
weights at smaller rates gives the common operator
bound \(\hat{\mathsf b}\) of order \(e^{16\bar\kappa}\).
The carrier construction uses parameter cubes of side
at least \(88\,\mathsf r_{\mathrm w}\mathfrak M\), where
\(\mathsf r_{\mathrm w}\) bounds the number of cubes
in a walk cell, with at least three parameter cubes
in each direction. The side is also chosen so that
its product with \(\dot\delta_0\) is at least
\(4352\bar\kappa\). These are floors on the regular
pair and walk construction, which is used throughout
the contraction.

Put \(\lambda_0=H(s)^{\mathsf T}\mathbf J\) and
\(\nu=\lambda-\lambda_0\). For \(|B'|_\infty\le r\)
consider the ball of radius \(R=2\hat{\mathsf b}r\)
in \(\max\{\|T\|,|\nu|_\infty\}\).
On this ball let \(\Delta\) be the maximum of the
weighted differences of \(T\) and \(\nu\) at two points.
The differences of the two quadratic remainders
\(\mathcal C,\mathcal N\) are bounded by \(CR\Delta\).
For the remaining term take the difference of
\(\mathcal C'_T{}^{\mathsf T}(\lambda_0+\nu)\) by putting
the weight on the changed factor. A difference in \(T\)
costs \(C(|\lambda_0|_\infty+R)\Delta\), and a difference
in \(\nu\) costs \(CR\Delta\). Since
\(|\lambda_0|_\infty\le \hat{\mathsf b}\gamma_0\), the
Lipschitz constant of the map is bounded by
\[
            C\hat{\mathsf b}(R+\hat{\mathsf b}\gamma_0).
\]
Choose both \(r\) and \(\gamma_0\) at most a sufficiently
small constant times \(\hat{\mathsf b}^{-2}\). The map has
Lipschitz constant at most \(1/2\), and its value at zero
has norm at most \(\hat{\mathsf b}r\). It therefore
preserves the ball. Iteration gives the holomorphic
solution, with the same contraction estimates in the
weighted norms. Subtracting the equations at two data
then bounds their solution difference by
\(2\hat{\mathsf b}|B'_1-B'_2|_\varpi\). The choice of
weight centred at the changed row gives the decay.

These estimates concern the reduced system on the interior and
outward rows. The identification at parameters one is proved below
at a real critical reference with its own current.

(k) \textbf{From the slice to a multiplier.} Let $s\equiv1$, and fix a critical real $U_0$ such that $(U_0,J(U_0))$ is a regular pair whose radii lie below the ceiling. Then (I) is equivalent to the constraint, the Landau condition and $\mathsf z:=F'(\mathbf H)-K_{\mathbf H}^{\mathsf T}\lambda\in DR(\text{site functions})$, where $\lambda$ is the multiplier formed from $\mathbf H$. In other words, (I) amounts to stationarity on the gauge slice alone. If $B'$ is real and close to $0$, then $\mathbf H(B')$ is real and lies in a neighbourhood of $0$ whose size is not uniform. On that neighbourhood the invariance of the action and the factorisation in 3 provide gauge directions $X_\chi$ with $(F'(\mathbf H),X_\chi)=0$ and $K_{\mathbf H}X_\chi=0$, and the $RD^{\mathsf T}X_\chi$ fill out the range of $R$. Therefore $\mathsf z=0$ on this neighbourhood, and by the identity theorem also on the whole polydisc, so the neighbourhood enters no constant. The factorisation of the averages holds only up to site factors and, at rows that leave their layer, up to seam factors; this is the reason for the row datum.

\paragraph{The reduced response over a real critical reference.}
For the remainder of this argument put $s=\mathbf 1$ and
take the real regular pair $(U_0,J)$ with
$J=J_{\Omega_0}(U_0)=\hat Q^{\mathsf T}\lambda_0$.
Use the closed row set: the outer end of an outward row
has its full block of interior rows at the preceding level.
The transposes use the fine weighted pairing and the counting
row pairing. Put $\mathcal Q(T)=\hat QT+\mathcal C(T)$ and
$K_T=\hat Q+\mathcal C'_T$. If $F(T)$ is the action of the
inner plaquettes at $e^{i\xi T}U_0$, then
$F'(T)=J+\Delta T+\mathcal V(T)$, with the same Hessian and
current convention as in the propagators. Here
$\mathcal V$ is the nonlinear force remainder.

Define $\mathsf z=F'(T)-K_T^{\mathsf T}\lambda$ and
$\phi=RD^{\mathsf T}T$. Write $r_1,r_2$ for the left sides
minus the right sides of the two reduced equations,
evaluated at $B'=\mathcal Q(T)$. Their exact expressions are
\begin{equation}\label{mb:reduced-response-defects}
 r_1=\widetilde G(\mathsf z+D\phi),\qquad
 r_2=\mathsf L G'\phi-H^{\mathsf T}(\mathsf z+D\phi).
\end{equation}
Indeed $\widetilde G\hat Q^{\mathsf T}=0$ and
$T=H_0\hat QT+\widetilde G G_0^{-1}T$, where
$G_0^{-1}T=\Delta T+D\phi+\hat Q^{\mathsf T}\mathsf a\hat QT$.
This gives $r_1$. For $r_2$ use $\hat QH=1$ and
\[
 H^{\mathsf T}\Delta T=(\hat S-\mathsf a)\hat QT
                   +H^{\mathsf T}\Delta'_{\pi}T-H^{\mathsf T}D\phi.
\]
To identify the middle term, set $\Upsilon=G'RD^{\mathsf T}$.
The quadratic form of $\Delta'_{\pi}$ has terms bilinear in
$(A,\Upsilon A)$ or quadratic in $\Upsilon A$. In the
polarization at $(H\beta,T)$ only the term with
$(H\beta,\Upsilon T)$ remains, because $RD^{\mathsf T}H=0$
at the real current. If $\mathfrak b(v,A)$ denotes this
bilinear part with the arguments in that order, define
$\langle\beta,\mathsf L v\rangle=\mathfrak b(v,H\beta)$.
This gives $H^{\mathsf T}\Delta'_{\pi}T=\mathsf L G'\phi$.
Also $H^{\mathsf T}J=\lambda_0$.

For the reduced solution the first equation gives
$\hat QT=B'-\mathcal C(T)$. Since the kernel of $\widetilde G$
is $\operatorname{ran}\hat Q^{\mathsf T}$, the two vanishing
defects in \eqref{mb:reduced-response-defects} give
\begin{equation}\label{mb:reduced-response-gauge}
 F'(T)+D\phi=K_T^{\mathsf T}\lambda+\hat Q^{\mathsf T}\mathsf L G'\phi.
\end{equation}
For real $B'$ near zero, let $X_\chi(T)$ be the variation
of $T$ by the real gauge $e^{i\tau\chi}$. The accompanying
node variation is $\varrho(T;\chi)$, obtained by undoing
the averaging factors of the two fields. The straightened
constraint varies by $\ell^\natural_{B'}[\varrho(T;\chi)]$,
the infinitesimal row action with its outward endpoint
transport. At $T=0$ the maps $X_\chi$ and $\varrho$ are
$-D\chi$ and $Q'\chi$. The map
\[
 \chi\longmapsto(\varrho(T;\chi),RD^{\mathsf T}X_\chi(T))
\]
is invertible near zero. At zero this follows from the
onto site average and the inverse of $RD^{\mathsf T}D$ on
$\ker Q'$. The finite-dimensional inverse persists on
an open real neighborhood. Choose $\chi$ with first
component zero and second component $\phi$. Then
$K_TX_\chi=0$ and $(F'(T),X_\chi)_\xi=0$ by gauge invariance.
Pairing \eqref{mb:reduced-response-gauge} with $X_\chi$ gives
\[
 (\phi,\phi)_\xi=-\langle\mathsf L G'\phi,\mathcal C'_T X_\chi\rangle.
\]
At fixed finite volume the right side is bounded by
$C\|T\|(\phi,\phi)_\xi$, since $\mathcal C'_0=0$ and
$\|X_\chi\|\le C\|\phi\|$. Thus $\phi=0$ near the real
origin. Holomorphy extends this identity to the entire
reduced datum ball, so $F'(T)=K_T^{\mathsf T}\lambda$ there.
This argument uses no uniform size for the real neighborhood.
In the reverse direction a solution of the three equations
has both defects zero. A sufficiently small $T$ puts its
multiplier in the same uniqueness ball: the second equation
bounds $|\lambda-\lambda_0|$ by
$C\hat{\mathsf b}\|T\|(1+|\lambda_0|+|\lambda-\lambda_0|)$,
where the last term is absorbed. This gives the comparison
on a common smaller uniform radius.

To return to physical data, the factorization of $\mathcal Q$
gives the plain datum $(e^{iB'^\natural}V_0)^{v[T]}$, where
$V_0=M_{\mathbf B}(U_0)$. At an outward row the right endpoint
factor is $\theta_z(B')^{-1}$ if the outer endpoint is $z=c_+$,
and the left factor is $\theta_z(B')$ if $z=c_-$. This is
the mean formed from the interior rows in the block at $z$.
Use the inverse conversion of the physical logarithm to $B'$
on the common smaller chart. To check criticality on the
physical fibre, pair $\lambda$ with $\ell^\natural_{B'}[\sigma]$.
At small real $B'$ the map of gauge variations above supplies
every node value $\sigma$. The pairing is zero since it equals
$(F'(T),X_\chi)_\xi$. Holomorphy extends its vanishing to the
datum ball. On a differentiable curve of real fields with
fixed physical datum, the derivative of the reduced datum
is a node variation of this form, obtained by differentiating
the factors $v[T]$ with the physical datum held fixed. The
derivative of the action on the curve therefore vanishes.
After undoing $v[T]$ at the nodes this gives a critical field
at the specified physical datum. Where the response lies in
the regularity class, uniqueness of the real orbit identifies
it with the continued regular background in the normalized
slice. For the decoupled response at an independent current
we use the reduced system and its analytic estimates. The
identification with a physical critical orbit uses the real
critical reference and the closed rows as above.

\emph{Constants} are chosen in the following order, each large or small enough relative to those chosen before it. First comes $L$, and with it, as a single group, the numbers determined by $L$ alone: the flat constants and the flat rate, the lower thresholds that the flat case imposes on $M_1$ and on $R_1M_1$, the radius of the local perturbation theory, the rates $\delta_0$, $\delta_*$, $\delta'$, and the ceiling on $\alpha_0,\alpha_1,\gamma_0$. Next come $M_1$ and $R_1M_1$; the ceiling does not depend on them, and they do not depend on the ceiling. Then comes $\bar\kappa$. It does not depend on $M_1,R_1$, and they do not depend on it, so it may just as well be chosen before them. After that come $\ell_{\mathsf P}$, the radius in $B'$, and the second ceiling on $\gamma_0$, which concerns \eqref{eq:propagators-II}. The final choices are $M$ and the field strengths of the construction, taken so small that $O(LM)\alpha_0$ is below the ceiling. Two numbers are exceptions: among the bounds of statement 7, the one bound that is used for convergence alone; and, in mechanism (i), the entrywise bound by which the family is normal. Neither of them appears in a localised bound, a rate, a weight or a sum of weights. With the exception of these two, every constant, rate and ceiling, and the radius in $B'$, is independent of $k$, $K$, $\ell$, of the sequence (in particular of $\Omega_0$), of $Y$ and of the pair, and also of $M_1,R_1$ once these are large enough. The sizes $\ell_{\mathsf P}$ and $M$, and the side of the coarse cubes of 6, are each at least a multiple of $M_1$. The field strengths are small compared with $(LM)^{-1}$. Where statement 8 makes the identification with the response of \cref{ch:multilevel-backgrounds}, the datum of the background is small compared with $M_1^{-1}$.


\subsubsection{The additional reduced rows at a retained inward seam}
\label{mb:retained-edge-linear-package}

This paragraph supplies the linear inverse for the optional convention
that retains inward bonds. The determining sets in
\eqref{mb:determining-rows} still use only interior and outward rows.
At each additional seam node choose one retained inward bond as its
tying bond, oriented towards that node. The tying equation fixes its
node factor. The remaining independent equations are the quotients
of the other retained inward bonds by this tying bond. Call these
the edge rows. If only one inward bond is retained at a node, it
adds a node factor but no edge row.

Write \(K_{\rm p}\) for the differential of the proper reduced
rows and \(K_{\rm e}\) for that of these edge rows, all in raw
dimensionless logarithmic coordinates. Both row spaces have counting
measure. In this paragraph \(K^*\) denotes the transpose for that
counting measure and the physical fine pairing
\(\eta^4\sum_b\langle A(b),A'(b)\rangle\). In particular it is
not the physical top transpose with its row-length weight omitted.
Let \(a_q(A)\) be the iterated reduced logarithm at level \(q\)
in the reference chart, and extend it to reversed bonds by
\(\widehat a_q(\bar b)=-\operatorname{Ad}(W_q(b)^{-1})a_q(b)\),
where \(W_q=\mathcal M^qU_0\). For an edge
\(e=(b_i(x),b_1(x))\), put
\(V_0(e)=W_q(b_i(x))W_q(b_1(x))^{-1}\). Differentiation of its
quotient gives the exact formula
\begin{equation}\label{mb:edge-row-differential}
 K_{\rm e}A(e)=\widehat{\dot a_q[A]}(b_i(x))
       -\operatorname{Ad}(V_0(e))\widehat{\dot a_q[A]}(b_1(x)).
\end{equation}

Put \(\mathfrak M=M_1\). The auxiliary peeled sequence has
\(\Omega'_0=\Omega_0\);
for \(q\ge1\), \(\Omega'_q\) is the union of the aligned
\(\mathfrak M\xi_q\)-cubes in \(\Omega_q\) whose sup distance
from \(\Omega_q^c\) is at least \(\mathfrak M\xi_q\).
Its determining set contains only proper rows. A full torus is
unchanged, and empty members are allowed. Use the same large
separation floors for both sequences. The elementary consequences are
\[
 \Omega_{q+1}\subset\Omega'_q\subset\Omega_q,
 \qquad \xi'(z)\le\xi(z)\le L\xi'(z).
\]
Indeed a discarded cube touches the old boundary and has width
\(\mathfrak M\xi_q\). Separation puts the next member beyond
this discarded strip; erosion of that next member adds
\(\mathfrak M\xi_{q+1}\) to its distance from the old boundary,
so the original separation is preserved. A site loses at most one
level. Each peeled pillar is contained in its old pillar.

\begin{lemma}[Uniform edge Schur inverse]
\label{lem:mb-retained-edge-Schur}
Let \(U_0\) be real. Assume that \((U_0,J(U_0))\) is a real
regular pair on the original sequence and that \((U_0,0)\) is a
real regular pair on the peeled sequence, in the preceding definition
with \(A'=0\) and \(\alpha_0=\alpha_1=\gamma_0=\alpha\).
Require the local-gauge bounds on every complete box used in the
proper-row inverse estimates, with a common sufficiently small
radius \(\alpha\). These are two primitive admissions: the original
pair uses the actual physical current, and the peeled pair uses
independent current zero. The original physical-current premise
ensures positivity of the zero-independent-current completed Wilson
form. Both admissions hold at the flat reference.
The separation floor may be increased by a fixed number depending
only on the inverse constants and \(L\).

Let \(G_{\rm p}=T^{-1}\) be the original positive completed
Wilson inverse, in raw row normalization, and set
\[
 \begin{split}
 T&=\Delta^{\rm W}_{U_0}+DRD^\sharp+K_{\rm p}^*K_{\rm p},\\
 H_{\rm p}&=G_{\rm p}K_{\rm p}^*
               (K_{\rm p}G_{\rm p}K_{\rm p}^*)^{-1},\qquad
 C_{\rm p}=G_{\rm p}-H_{\rm p}K_{\rm p}G_{\rm p},\\
 M_{\rm e}&=K_{\rm e}C_{\rm p}K_{\rm e}^*.
 \end{split}
\]
There are \(c,C,s>0\), independent of depth, mesh, volume and
the subset of inward bonds retained, such that
\[
 cI\le M_{\rm e}\le CI,\qquad
              |M_{\rm e}^{-1}(e,e')|\le Ce^{-s d(y_e,y_{e'})}.
\]
Here \(y_e\) is the outer top of the nontying bond of \(e\),
and \(d\) is the original weighted distance. If there are no
edge rows, the edge space is zero and the following corrections
are zero. Otherwise define
\begin{equation}\label{mb:edge-Schur-package}
 \begin{split}
 H_{\rm all}(b_{\rm p},b_{\rm e})
   &=H_{\rm p}b_{\rm p}+C_{\rm p}K_{\rm e}^*M_{\rm e}^{-1}
                         (b_{\rm e}-K_{\rm e}H_{\rm p}b_{\rm p}),\\
 C_{\rm all}
   &=C_{\rm p}-C_{\rm p}K_{\rm e}^*M_{\rm e}^{-1}K_{\rm e}C_{\rm p}.
 \end{split}
\end{equation}
With \(K=(K_{\rm p},K_{\rm e})\) and
\(\mathcal X=\Delta^{\rm W}_{U_0}+DRD^\sharp\), these maps satisfy
\[
 KH_{\rm all}=1,\quad KC_{\rm all}=0,\quad
 C_{\rm all}\mathcal X=1-H_{\rm all}K,\quad
                    \ker C_{\rm all}=\operatorname{ran}K^*.
\]
They have the same raw-datum-to-field and force-to-field bounds
as the proper package, with a smaller positive exponential rate.
More precisely, use the cell field norm
\(\|A\|_y=\max\{\xi_y|A|_y,\xi_y^2|\nabla_{U_0}A|_y\}\)
and force norm \(|f|^\sharp_y=\xi_y^3|f|_y\). Here the maxima
are over bonds or difference sites assigned to the pillar of \(y\),
and \(|b|_y\) is the maximum of the Euclidean raw row norms at
that top, taken as zero if there is no row. Complete fixed
enlargements in local estimates have bounded overlap. Then
\[
 \begin{split}
 \|H_{\rm all}b\|_y&\le C\sum_z e^{-s d(y,z)}|b|_z,\\
 \|C_{\rm all}f\|_y&\le C\sum_z e^{-s d(y,z)}|f|^\sharp_z.
 \end{split}
\]
The localized gradient H\"older bounds hold for each fixed
\(0<\beta<1\), with constants additionally depending on \(\beta\).
This is a linear package for the completed Wilson Hessian; it is
not an identification with a current-corrected background response.
\end{lemma}

\begin{proof}
Let \(H'_{\rm p}\) be the proper raw-row lift on the peeled
sequence. Use the Euclidean Lie algebra norm induced by the real
trace pairing and put
\(\|A\|_{\mathcal E}^2=\sum_y\|A\|_y^2\).
The proper exponential estimates and the bounded weighted metric
volume imply
\[
 \|H_{\rm p}b\|_{\mathcal E}\le C\|b\|_2,
 \quad \|H'_{\rm p}b'\|_{\mathcal E'}\le C\|b'\|_2,
 \quad \|A\|_{\mathcal E}\le L^2\|A\|_{\mathcal E'}.
\]
The last inequality follows by replacing the maximum on an old
pillar by the sum of squares on its disjoint peeled subpillars.
All constants in this proof are fixed before the depth is chosen.

We construct a test for each edge row. Write
\(b_i(x)=(y_i,x)\), with \(x=y_i+\sigma\xi_qe_\mu\), and
let \(B(X)\) be the level-\((q+1)\) block containing \(x\).
The bonds
\[
 f_0=(y_i,x),\quad f_1=(x,x+\sigma\xi_qe_\mu),\quad
 f_2=(x+\sigma\xi_qe_\mu,x+2\sigma\xi_qe_\mu)
\]
are proper interior rows of the peeled sequence. In fact this
entire block is in the strip removed from \(\Omega_{q+1}\),
while it and \(y_i\) lie in \(\Omega'_q\). On their common
line transport a vector \(v\) from \(y_i\) to the initial
point of each bond, and assign the oriented values \((1,-2,1)v\)
in these transported frames. On a reversed bond use the covariant
minus sign. Denote this pattern by \(J_ev\).
Different edge rows have disjoint pattern supports: inward bonds
at a common seam node have distinct directions, and parallel
inward bonds on the same block line are separated by at least
\(L-1\) bonds. Since \(L\ge13\), the three-bond patterns cannot
meet. None uses a tying bond. Thus
\[
                  \|Jb\|_2^2=6\|b\|_2^2,
                \qquad Jb=\sum_eJ_eb(e).
\]

The cancellation at the next scale is essential. Along a line,
the \(L\)-step paths starting in the first endpoint block meet
the pattern in the sets \(\{0\},\{0,1\},\{0,1,2\}\).
Their sums are \(1,-1,0\), the first two occurring once.
For the next coarse bond the nonzero sums are \(-1,1\), again
once each. At a flat reference each pair cancels exactly in
the \(L^{-4}\)-normalized straight-path average. At a curved
reference the two transports differ by a loop in the same double
block. The loop comparison bounds that difference by \(C_L\alpha\).
Let \(\mathbf Q_q=D_aF_{W_q}(0)\), where \(F_W\) is the
straightened one-step logarithm of
\eqref{px:actual-one-step-straightening}. It is the derivative in
the chain rule for the reduced logarithms \(a_{q+1}\).
At a flat reference, differentiating the two endpoint factors in
that formula cancels the internal contour derivatives of the plain
mean. The remaining map is the covariant straight-path average.
For curvature \(O(\alpha)\), differentiate the same endpoint
factors, local logarithms and defining mean equation on the complete
double block. In its real comb gauge their coefficients and the
inverse derivative of the mean equation differ from their flat
values by \(O_L(\alpha)\), using the logarithm series on its fixed
smaller ball. Unitary covariance restores the original frames.
Thus \(\mathbf Q_q\) differs from the covariant straight-path map
by \(C_L\alpha\) in its fixed finite row and column sums.
The one-step row has at most \(8L^4\) input bonds, and each input
belongs to at most eight double blocks. These are one-step counts,
with no dependence on \(q\).

Set \(A_b=H'_{\rm p}Jb\). The iterated reduced logarithms for
the two sequences agree on every common complete read, by induction
in the averaging equation. Hence
\[
          K_{\rm e}A_b=b,\qquad
          \|K_{\rm p}A_b\|_2\le C_L\alpha\|b\|_2,
          \qquad \|A_b\|_{\mathcal E}\le C\|b\|_2.
\]
For the first identity, \eqref{mb:edge-row-differential} reads the
assigned value on \(b_i\) and zero on the tying bond. An old proper
row shared by the peeled sequence reads zero. Every other old
proper row has its full predecessor double block among the peeled
interior rows, so the preceding cancellation and the chain rule
give the second inequality. Bounded double-block overlap gives
both Schur sums needed for that \(\ell^2\) estimate.

Correct into the proper kernel by
\(A=A_b-H_{\rm p}K_{\rm p}A_b\). The actual local row bounds give
\[
 K_{\rm p}A=0,\qquad \|A\|_{\mathcal E}\le C_A\|b\|_2,
 \qquad \|K_{\rm e}A-b\|_2\le C_\mathrm e\alpha\|b\|_2.
\]
Choose \(C_\mathrm e\alpha\le1/2\). On this proper kernel,
\[
                (A,TA)_\eta\le C_T\|A\|_{\mathcal E}^2.
\]
To check this bound, the row penalty vanishes, \(R\) is an
orthogonal projection, the curl and divergence use finitely many
neighboring differences, and the Wilson-minus-curl remainder is
bounded by \(C_L\alpha\xi_y^{-2}|A|\). The physical number of
bonds in a pillar is \(O(\xi_y^4)\); this cancels the four powers
in the squared first-difference bound. Neighboring layers differ
by at most \(L\).

Since \(T\) is positive, its constrained inverse satisfies
\[
 (f,C_{\rm p}f)_\eta
   =\sup_{K_{\rm p}u=0}\{2(f,u)_\eta-(u,Tu)_\eta\}.
\]
Put \(f=K_{\rm e}^*b\), test with \(u=tA\), and use
\(\langle b,K_{\rm e}A\rangle\ge\|b\|_2^2/2\). Taking
\(t=(2C_TC_A^2)^{-1}\) proves
\[
          \langle b,M_{\rm e}b\rangle
                      \ge(4C_TC_A^2)^{-1}\|b\|_2^2.
\]
The upper bound and an exponentially decreasing raw kernel for
\(M_{\rm e}\) follow by composing the proper constrained inverse
with the two actual local edge maps. Their counting transposes
have the physical fine density; no mesh-counting factor is introduced.

For completeness, fix an edge \(e_0\) and conjugate this raw
matrix by multiplication by
\(\exp(s\min\{d(y_e,y_{e_0}),N\})\).
The distance is one-Lipschitz in itself. Let \(a>0\) be the
kernel rate already obtained. For \(0<s\le a/4\), the exponential
kernel bound and the metric sum at the fixed rate \(a/2\) bound
the row and column sums of the conjugation error by \(C s\),
uniformly in \(N\). The constant \(C\) is fixed once the
separation exceeds the floor for this rate. Choose
\(s\le\min\{a/4,c/(2C)\}\).
The conjugated inverse has norm at most \(2/c\). Apply it to a
unit coordinate at \(e_0\) and let \(N\) increase. This proves
the displayed decay of \(M_{\rm e}^{-1}\). Now increase the
separation floor also to bound the metric exponential sums at this
smaller rate, using the same constants already fixed above.
Only a fixed number of edge species occur at a top.

Direct multiplication in \eqref{mb:edge-Schur-package} proves
the first three identities. Positivity on \(\ker K\) identifies
\(C_{\rm all}\) as its constrained inverse, giving
\(\ker C_{\rm all}=\operatorname{ran}K^*\).
Composition of the exponential kernels proves the two cell bounds
after one further fixed decrease of the rate. For the localized
gradient H\"older estimate, write each correction as
\(C_{\rm p}f\) and apply the proper estimate to this force;
its bound follows from the same convolution. This uses no derivative
of a row datum. All maps are real at real coefficients. On a
common smaller admitted complex tube their bilinear matrix formulas
are holomorphic; the uniform weighted perturbation estimates for
the proper inverse and local rows preserve the edge gap by Neumann
series. At the flat reference the loop error is zero and both
primitive admissions are automatic. The resulting flat package
therefore includes any retained inward edge rows in the local
regularity argument.
\end{proof}

\paragraph{The derivative at a retained-inward critical background.}
Suppose additionally that \(U_0\) is a real critical reference
for these reduced rows. Write \(q(A)\) for their raw nonlinear
logarithm, so \(q'(0)=K\), and write \(F\) for the inner
Wilson action in this chart. In this paragraph put
\(H=H_{\rm all}\), \(C=C_{\rm all}\) and
\(Z=DRD^\sharp\). The nonlinear system to be solved is
\[
 F'(A)+ZA=q'(A)^*\omega,\qquad q(A)=b.
\]
The gauge term is retained until the physical-fibre argument.
Since \(KH=1\), the multiplier at the reference is unique and
\[
 F'(0)=K^*\omega_0,\qquad
 \omega_0=H^*F'(0),\qquad |\omega_0|_\infty\le C_0\alpha.
\]
Here and below the constants are uniform over the admitted data.
The last bound uses the physical current admission and the absolute
column estimate: in a cell the pairing has scale
\(\xi_y^{-1}\xi_y^{-3}\xi_y^4=1\).

Define \(B\) by
\[
 (v,Bu)_\eta=\langle\omega_0,q''(0)[u,v]\rangle.
\]
The local two-leg bounds for the actual logarithms, with both leg
sums and the fine density, give
\(|Bu|_{(-3)}\le C_1|\omega_0|_\infty\|u\|_1\).
Differentiating both factors of an edge quotient gives the same
bound; only a fixed number of inward bonds share a tying bond.
The differential of the system at \((0,\omega_0)\) is
\[
 \mathcal L(u,\nu)=((\mathcal X-B)u-K^*\nu,Ku).
\]
For a force \(f\) and raw datum \(b\), its inverse has field part
\[
 u=(1-CB)^{-1}(Cf+Hb),\qquad
 \nu=H^*((\mathcal X-B)u-f).
\]
Applying \(C\) proves the first formula. Conversely it gives
\(Ku=b\), and the first-row residual belongs to
\(\ker C=\operatorname{ran}K^*\); applying \(H^*\)
then gives the unique multiplier in the second formula.
Decrease \(\alpha\) so that \(\|CB\|_{1\to1}\le1/2\),
also in a fixed smaller exponentially weighted norm.

The weak multiplier estimate uses the unchanged site maps of the
proper package. The proper exponential bounds for \(P=1-R\) on
site functions and for \(DPD^\sharp\) remain valid.
Write \(\mathrm{d}\) for the covariant bond-to-plaquette curl, reserving
\(D\) for the site-to-bond derivative. The divergence maps
a field controlled in \(\|\cdot\|_1\)
to a site function with its two powers of local length. In
\(\langle e,H^*\mathrm{d}^\sharp\mathrm{d}u\rangle
=(\mathrm{d}He,\mathrm{d}u)_\eta\)
move one difference to each field; for the gauge term use
\(\langle e,H^*DRD^\sharp u\rangle
=(D^\sharp He,RD^\sharp u)_\eta\) and \(R=1-P\).
The cell counts cancel the four inverse powers of length.
The Wilson-minus-curl term uses its local force bound.
Together with the column bound for \(H^*f\), this proves
\[
 \max\{\|u\|_1,|\nu|_\infty\}
 \le C_I\max\{|f|_{(-3)},|b|_\infty\}.
\]
The same argument with the weight on the changed factor proves
the weighted version, with a fixed smaller rate.

Here is an explicit contraction using this inverse. Write
\(\omega=\omega_0+\nu\), \(X=(A,\nu)\),
\(\|X\|=\max\{\|A\|_1,|\nu|_\infty\}\), and put
\[
 \begin{split}
 N_1(X)&=F'(A)+ZA-q'(A)^*(\omega_0+\nu)
                 -[(\mathcal X-B)A-K^*\nu],\\
 N_2(X)&=q(A)-KA.
 \end{split}
\]
Both remainders and their first derivatives vanish at zero.
The local Taylor and transpose bounds give, on
\(\|X_i\|\le r\) in the common chart,
\[
 \max\{|N_1(X_2)-N_1(X_1)|_{(-3)},
             |N_2(X_2)-N_2(X_1)|_\infty\}
 \le C_N r\|X_2-X_1\|.
\]
The mixed term is \(-(q'(A)-K)^*\nu\); its difference has
one small factor and one difference. The part paired with
\(\omega_0\) uses the second difference of \(q'\).
Choose \(r\) below the common row and action chart radii and
so that \(C_IC_Nr\le1/4\). For
\(|b|_\infty<r/(2C_I)\), the map
\[
 X\longmapsto\mathcal L^{-1}(-N_1(X),b-N_2(X))
\]
maps the closed \(r\)-ball into itself and contracts by at
most \(1/4\). Iteration gives its unique solution there,
holomorphic on the open datum ball and real for real data.
Weighted differences have the same bound after absorption.
All constants and radii were fixed before the depth and the
retained subset were chosen.

Physical identification uses the same original and peeled
admissions, the actual node and seam datum conversion, and the
common smaller chart described above. The real gauge-variation
argument removes \(ZA\) near the origin and identifies
criticality on the physical fibre; its analytic identities extend
over the connected datum ball. The compared physical branch must
lie in this response uniqueness ball and in the admitted regularity
class. Uniqueness of its real critical orbit then gives the
normalized branch identification. Thus the optional continuation
uses the corrected derivative \(\mathcal L\) and the stated
admissions throughout.

% END MODULE chapters/04_propagators.tex

% BEGIN MODULE chapters/layered_current_precision.tex
\subsection{Layered precision and current terms}
\label{ch:layered-current-precision}

The precision of a constrained fluctuation is obtained only after
three different operations. One forms an operator on the fine
field, reads its induced action on the coarse datum, and then
restricts that coarse action to the next actual fibre. Gauge
and current terms enter at all three places. A positive
matrix at one of them does not identify the others.

We retain this order. The fine site and field operators,
their row lift and their constraint-curvature terms are
typed separately. Their complete physical match is an
identity of projected forms, with the actual soft gauge.
The alternative explicitly defined relative-current
first-step extension gives another coherent starting
point. It supplies its stated first-step estimates;
using that definition at later levels requires its
own consistency and closure.

The reductions below use the domain and background
sections. In particular, all transposes are bilinear
weighted transposes on complex domains, while real
lower bounds use Hermitian quadratic forms. The
primitive domains are specified before the unknown
critical field. No arbitrary independent current is
replaced by the physical action current.

\subsubsection{Cell norms retain the ultraviolet scale}

Work on a finite compatible torus, with a separated
aligned sequence and the spacings of
\cref{ch:multilevel-backgrounds}.
Let \(\Delta(y)\) be the pillar of a top \(y\),
of side \(\xi_y\). Outside the action region
one can add bare cells and extend the local
scale to \(\eta\); a compressed operator is
zero on that added space. The point scale
\(\xi(x)\) is the scale of its pillar.
Use the exact path metric
\begin{equation}
 d(x,z)=\inf_{\gamma:x\to z}
                   \int_\gamma\frac{ds_\infty}{\xi(s)} .
 \label{px:layer-metric}
\end{equation}
The half-open convention only assigns a
measure-zero cell boundary; paths and
their infima can be approximated on
either side. The corresponding lattice
metric is comparable after the fixed
collar constants have been inserted.

For each field species \(s\), declare a
positive amplitude \(a_s(y)\) and use
\(\|f\|_{s,y}=a_s(y)^{-1}
\|f\|_{\infty,\Delta(y)}\).
For example a physical potential has
amplitude \(\xi_y^{-1}\), a physical
force density \(\xi_y^{-3}\), and a
raw datum logarithm amplitude one.
The site gauge parameter also has
amplitude one. A raw multiplier uses
counting measure and amplitude one;
its per-length density has amplitude
\(\xi_y^{-3}\). The conversion is
the weighted dual conversion of
the domain section.

If \(P_y\) is restriction to a cell,
define the normalized cell entry of
an operator \(K:t\to s\) by
\begin{equation}
 k_{s t}(y,z)=
       \|P_yKP_z\|_{\,\ell^\infty(t,z)\to\ell^\infty(s,y)} .
 \label{px:cell-entry}
\end{equation}
All physical cell measures in the
action of \(K\) are included before
taking this norm. It is a block
operator norm, not a pointwise
bound on the continuum kernel.
A Green kernel with its fine
Coulomb singularity can have a
bounded cell norm. Replacing it
by a pointwise kernel bound
independent of \(\eta\) would
be a stronger and generally
false assertion.

For compatible species, the cell
partition and triangle inequality
give
\begin{equation}
 (k\ell)_{s u}(y,z)
       \le\sum_w k_{s t}(y,w)\ell_{t u}(w,z).
 \label{px:cell-composition}
\end{equation}
This follows first for finite
lattices by inserting
\(\sum_wP_w=1\). In an
infinite-volume application
the same formula holds when
the row sums below give
absolute convergence.
Thus a two-sided exponential
cell norm is an algebra norm,
without discarding the fine
singularity inside a block.

\subsubsection{The growth of a ball in the layer metric}

The layer metric is not assumed
to have polynomial volume growth.
A small-cell region can contain
many more pillars than a nearby
large-cell region. Separation
controls that growth quantitatively.

Suppose \(x\) has layer index \(j\)
and \(z\) has index \(i\).
Any path between them crosses
each complete intervening layer.
For a layer \(r\), a crossing
from \(\Omega_{r+1}\) to
\(\Omega_r^c\) has an arc
whose interior lies in
\(\Omega_r\setminus\Omega_{r+1}\).
The separation makes its
physical length greater
than \(\rho\xi_r\).
On that arc the metric
density is \(\xi_r^{-1}\),
so its length in \eqref{px:layer-metric}
is at least \(\rho\).
Arcs belonging to distinct
layers have disjoint
interiors. Consequently
\begin{equation}
 d(x,z)\ge\rho(|i-j|-1)_+ .
 \label{px:level-barrier}
\end{equation}
One layer is lost because an
endpoint can already lie next
to its neighboring layer.
The same proof applies to
an intermediate point of a
path, so a path of metric
length at most \(R+1\)
starting in layer \(j\)
never uses a scale larger
than
\begin{equation}
 \xi_{\max}\le
            L\xi_j L^{(R+1)/\rho}.
 \label{px:path-max-scale}
\end{equation}
If the top index would exceed
the depth, this is simply
an overestimate.

Its physical length is at
most \((R+1)\xi_{\max}\).
Every pillar center in
the metric \(R\)-ball
therefore has a lift
in the physical sup-norm
ball of radius
\begin{equation}
 H_R=L(R+1)\xi_j L^{(R+1)/\rho}.
 \label{px:physical-ball}
\end{equation}
For each fixed level \(i\)
the centers lie on its
grid of spacing \(\xi_i\).
The number of grid points
in that physical ball
is at most \(C(1+H_R/\xi_i)^4\).
\eqref{px:level-barrier} restricts
the possible indices to
\(|i-j|\le1+(R+1)/\rho\).
Using the smallest such
\(\xi_i\) in the grid
count and then summing
the possible levels
gives
\begin{equation}
 \#\{z:d(y,z)\le R\}
   \le C_L(1+R)^5
                  \exp\!\left(\frac{8\log L}{\rho}R\right).
 \label{px:metric-ball-count}
\end{equation}
Indeed \(H_R/\xi_i\le
C_L(R+1)L^{2(R+1)/\rho}\);
its fourth power gives
the exponential in
\eqref{px:metric-ball-count}, and the number
of levels costs one
more polynomial power.
Since \(\rho\ge1\),
the harmless factors
\(L^{8/\rho}\) and
the bound on the
number of levels
can be included
in \(C_L\) uniformly
in \(\rho\).

The argument works on a torus
without a nonwrapping-ball
assumption. For every
reached center choose
the lift reached by
its path. Distinct
periodic centers
give distinct lifted
grid points. Counting
all grid points in
the cover ball is
therefore an upper
bound, even if
that ball crosses
many fundamental
domains.
Restricting the
allowed regions
only decreases
the number counted.

An entry bound of the form
\begin{equation}
 k_{s t}(y,z)\le
   C\left(\frac{\xi_y}{\xi_z}\right)^p e^{-m d(y,z)}
 \label{px:entry-with-scale}
\end{equation}
now implies a weighted row and
column sum, provided the
rate pays for scale growth.
By \eqref{px:level-barrier} the scale ratio
on an annulus of metric
radius \(R\) is at most
\(L^{|p|}e^{|p|(\log L)(R+1)/\rho}\).
Break the space into
unit annuli and use
\eqref{px:metric-ball-count}. Thus
\begin{equation}
 \sup_y\sum_z e^{\delta d(y,z)}k_{s t}(y,z)
 +\sup_z\sum_y e^{\delta d(y,z)}k_{s t}(y,z)
       \le C'
 \label{px:entry-to-schur}
\end{equation}
whenever
\begin{equation}
 m-\delta>\frac{(8+|p|)\log L}{\rho}.
 \label{px:summability-rate}
\end{equation}
For the column sum use the
same ball argument centered
at \(z\); the absolute
scale exponent is
unchanged.
The convergent scalar
sum is bounded by
\(C\sum_{n\ge0}(n+2)^5
 e^{-[m-\delta-(8+|p|)\log L/\rho]n}\).
This proves the claim
with its precise
uniformity.

If the amplitude normalization
has already removed the
scale ratio, take \(p=0\).
For several species use
the largest of their
finitely many exponents.
An exponential bound
with an arbitrarily
small \(m\) does not
automatically imply
\eqref{px:entry-to-schur}. The geometry
and rate must satisfy
\eqref{px:summability-rate}, or a stronger
actual row estimate
must be proved
directly.
This is a numerical
condition in the
order of choice
of separation and
kernel rates.

\subsubsection{The exponential-chart Ward identity}

We derive a complete current-linear normal form.
This also states exactly what a separately
specified current completion would have to match.
For this calculation use raw bond logarithms
and counting pairings. Equivalently use physical
potentials and their dual fine cell weights,
with the conversions stated below.

Let \(F_U(a)=A_{\Omega_0}(e^{ia}U)\), with
all plaquettes of this action inner to its
gauge-invariant active graph, and put
\(j_U=DF_U(0)\), \(\mathscr D_U=D^2F_U(0)\).
For a site function \(\lambda\), define
the raw infinitesimal gauge generator by
\begin{equation}
 (D_U\lambda)(b)
       =\lambda(b_-)-\operatorname{Ad}(U_b)\lambda(b_+).
 \label{px:gauge-generator}
\end{equation}
This equation fixes the sign of \(D_U\)
in this section. A convention with
the opposite covariant difference
changes both occurrences consistently.
The site Laplacian \(D_U^\sharp D_U\)
is unchanged.

At a nonzero reference logarithm \(a\)
the gauge generator in that same
logarithmic coordinate is
\begin{equation}
 \mathfrak d_U(a)\lambda
   =\mathcal E_a^{-1}
          \{\lambda(b_-)
           -\operatorname{Ad}(e^{ia_b}U_b)\lambda(b_+)\},
 \qquad
 \mathcal E_a h=\int_0^1
                  \operatorname{Ad}(e^{ita})h\,dt .
 \label{px:nonlinear-gauge-generator}
\end{equation}
The equality follows by differentiating
\(e^{it\lambda_-}e^{ia_b}U_be^{-it\lambda_+}\)
and multiplying by its inverse on the
right. The factor \(1/i\) in the
right-logarithmic derivative is
included in \(\mathcal E_a\).
The integral formula follows by
differentiating the exponential
series, or by integrating its
constant-coefficient differential
equation. It is holomorphic on
the admitted logarithm chart.

The derivative at zero is
\begin{equation}
 D\mathfrak d_U(0)[v]\lambda
  =-\frac i2
        [v_b,\lambda(b_-)+\operatorname{Ad}(U_b)\lambda(b_+)] .
 \label{px:gauge-generator-derivative}
\end{equation}
Indeed \(D\mathcal E_0[v]h=(i/2)[v,h]\),
while differentiating the endpoint
transport contributes
\(-i[v,\operatorname{Ad}(U_b)\lambda_+]\).
Subtracting the derivative of
\(\mathcal E_a\) acting on
\(\lambda_- -\operatorname{Ad}(U_b)\lambda_+\)
gives exactly \eqref{px:gauge-generator-derivative}, with a plus
between the two endpoint site
values inside its commutator.

Gauge invariance implies
\(DF_U(a)[\mathfrak d_U(a)\lambda]=0\).
Differentiating once at zero proves
\begin{equation}
 D_U^\sharp j_U=0,\qquad
              \mathscr D_UD_U=\mathcal W_U[j_U],
 \label{px:physical-ward}
\end{equation}
where the explicit linear-current operator is
\begin{equation}
 (\mathcal W_U[J]\lambda)(b)
      =\frac i2[
           \lambda(b_-)+\operatorname{Ad}(U_b)\lambda(b_+),J_b].
 \label{px:current-ward-map}
\end{equation}
To verify its order and sign, cyclicity gives
\(\langle J,(i/2)[v,S]\rangle
 =\langle v,(i/2)[S,J]\rangle\).
Together with \eqref{px:gauge-generator-derivative} this is the
differentiated invariance identity.
The formula uses the actual action
current. No constrained criticality
has been assumed.

The operator \(\mathcal W_U[J]\) is
defined for an independent \(J\)
by \eqref{px:current-ward-map}. It is local and
linear in that current. It
does not differentiate \(J\).
For an arbitrary independent
current, \(D_U^\sharp J\)
need not vanish, and
\(D_U^\sharp\mathcal W_U[J]\)
need not be symmetric.
Those physical identities
cannot be inserted in
an off-current definition.

If the spacing is \(\eta\),
the physical generator is
\(D_U^{\rm phys}=\eta^{-1}D_U\).
The raw and physical currents satisfy
\(j_{\rm raw}=\eta^3J_{\rm phys}\),
and the physical Hessian operator
is \(\eta^{-2}\mathscr D_{\rm raw}\)
with respect to the \(\eta^4\)
pairing. Thus
\begin{equation}
 \mathscr D_{\rm phys}D_U^{\rm phys}
       =\mathcal W_U[J_{\rm phys}].
 \label{px:physical-ward-units}
\end{equation}
There is no extra inverse power of
\(\eta\) in \(\mathcal W_U[J_{\rm phys}]\).
Both sides have the dimension of
a force density applied to a
dimensionless site gauge.

For the operator formulas from this
point onward restore the physical
cell pairings and write \(D_U\)
for \(D_U^{\rm phys}\).
Write \(j_U\) for the physical
current density and
\(\mathscr D_U\) for the
physical Hessian operator.
Their conversion is exactly
the one in \eqref{px:physical-ward-units}.
The formula in \eqref{px:current-ward-map} is
unchanged. In bare units
\(\eta=1\) all these
conventions coincide.

\subsubsection{An explicit symmetric current correction}

Use the site operators of the
background section and set
\begin{equation}
 \Lambda=G'RD_U^\sharp,\qquad
                     \mathsf S=1-D_U\Lambda .
 \label{px:lambda-projection}
\end{equation}
On their admitted domain the
proved projector identities
give \(\mathsf SD_UG'R=0\).
Define the following specific
symmetric current-linear form:
\begin{equation}
 \begin{split}
 \mathcal B_U[J]={}&
       \mathcal W_U[J]\Lambda
           +\Lambda^\sharp\mathcal W_U[J]^\sharp\\
 &-\tfrac12\Lambda^\sharp
       \{D_U^\sharp\mathcal W_U[J]
                    +\mathcal W_U[J]^\sharp D_U\}\Lambda .
 \end{split}
 \label{px:symmetric-current-correction}
\end{equation}
This is a definition on the
whole independent-current
domain. Its symmetry is
visible term by term.
For the physical current,
\eqref{px:physical-ward} and symmetry of
\(\mathscr D_U\) imply
\(D_U^\sharp\mathcal W_U[j_U]
 =D_U^\sharp\mathscr D_UD_U\),
which is symmetric. Expanding
\((1-D_U\Lambda)^\sharp
 \mathscr D_U(1-D_U\Lambda)\)
therefore proves
\begin{equation}
 \mathscr D_U-\mathcal B_U[j_U]
                   =\mathsf S^\sharp\mathscr D_U\mathsf S .
 \label{px:physical-invariant-hessian}
\end{equation}
This is an exact identity at
every admitted real or complex
field with its physical
current. It requires no
inner minimization.

The apparent derivative of \(J\)
in the middle term of \eqref{px:symmetric-current-correction}
is never needed for a norm
estimate. Its bilinear
form can be written
\begin{equation}
 \begin{split}
 \langle u,\mathcal B_U[J]v\rangle={}&
   \langle u,\mathcal W_U[J]\Lambda v\rangle
     +\langle\mathcal W_U[J]\Lambda u,v\rangle\\
 &-\tfrac12\langle D_U\Lambda u,
                           \mathcal W_U[J]\Lambda v\rangle
   -\tfrac12\langle\mathcal W_U[J]\Lambda u,D_U\Lambda v\rangle .
 \end{split}
 \label{px:current-form-no-current-gradient}
\end{equation}
Derivatives are kept on \(\Lambda\).
Thus bounded current, a site
inverse bound and its
differentiated field
readouts are the actual
ingredients in a
current perturbation
estimate. Replacing
this form by a
pointwise bound on
\(D_U^\sharp\mathcal W_U[J]\)
would introduce an
unassumed derivative
of the current.

%
%
\paragraph{Comparison with the named gauge-restoration form.}
\hypertarget{named-gauge-restoration}{}
Use the same active graph, background, site inverse and projection
in the two definitions. The named form uses the covariant
difference \(D_{\rm n}=-D_U\). This leaves the site operators
unchanged and gives its gauge read
\(G'RD_{\rm n}^\sharp A=-\Lambda A\).
Write \(w=\Lambda A\) and
\(w_+=\operatorname{Ad}(U_b)w(b_+)\), \(w_-=w(b_-)\) on a bond.
With the spacing \(\xi\), \(D_Uw=\xi^{-1}(w_--w_+)\).
Substitution of the gauge read \(-w\) in the named quadratic
form gives the bond expression
\[
 i[A,w_-+w_+]-i\xi^{-1}[w_-,w_+] ,
\]
paired with \(J\) in the physical cell pairing.

The diagonal of \eqref{px:symmetric-current-correction} is
\(2\langle A,\mathcal W_U[J]w\rangle
-\langle D_Uw,\mathcal W_U[J]w\rangle\).
Cyclicity in the first pairing gives \(i[A,w_-+w_+]\).
In the second, use
\([w_--w_+,w_-+w_+]=2[w_-,w_+]\).
Thus its bond expression is the same. Polarization proves
\(\mathcal B_U[J]=\Delta'_\pi[U,J]\) for every independent
current on the common domain. In particular the equality uses
the named quadratic form itself. The physical Ward identity
is used separately when \(J=j_U\).

The site output \(w\) lies in \(\ker Q'\). At each site of
the bare layer outside \(\Omega_1\) the site row is the
identity, so \(w\) vanishes there. The displayed commutators
therefore use the current only on the inner bonds with an
endpoint in \(\Omega_1\). This is the read set of the
gauge-restoration term, before adding the separate
constraint-curvature subtraction.

For the explicit choice \eqref{px:symmetric-current-correction},
\(R=0\) implies
\(\Lambda=0\) and
\(\mathcal B_U[J]=0\)
for every \(J\).
For an unspecified
completion this
conclusion does
not follow merely
from zero-fold
averaging maps.
The physical
identity in \eqref{px:physical-invariant-hessian}
and the choice
of an off-current
extension are
different statements.

In particular \(\mathcal B_U[j_U]\)
is a gauge-Ward correction, not
the full constrained-action
curvature correction even on
the physical graph.
For horizontal vectors \(u,v\)
with \(\Lambda u=\Lambda v=0\),
its bilinear form is zero.
The Lagrangian Hessian still
contains
\(-\langle\omega,D^2q[u,v]\rangle\).
The distinct subtraction
in \eqref{px:second-current-variant} is precisely
where that mean-constraint
term belongs after its
actual multiplier has
been identified. No
gap for a constrained
fluctuation is transferred
from the Ward identity
alone.

\subsubsection{The off-current gauge defect is computable}

The preceding explicit correction has useful
identities even before inverse estimates.
Use the actual row intertwining
\(QD_U=\bar DQ'\), on the
specified site and row spaces.
This is an input identity to
be checked for the actual
straightened rows, including
all outward endpoint values.
Set
\begin{equation}
 \widehat G[J]^{-1}
    =\mathscr D_U-\mathcal B_U[J]
                           +D_URD_U^\sharp+Q^\sharp aQ ,
 \qquad \delta J=J-j_U .
 \label{px:explicit-current-green}
\end{equation}
Every assertion about an inverse is
restricted to its admitted
invertibility domain.
The identities
\(\mathsf SD_UG'R=0\),
\(Q'G'R=0\) and
\(\Delta G'R=R\), with
\(\Delta=D_U^\sharp D_U\),
give
\begin{equation}
 \widehat G[J]^{-1}D_UG'R
       =D_UR-\mathcal B_U[\delta J]D_UG'R .
 \label{px:current-defect-before-inversion}
\end{equation}
For the action part use
\eqref{px:physical-invariant-hessian}, whose right side
annihilates \(D_UG'R\).
The gauge part is
\(D_UR\Delta G'R=D_UR\).
The row mass is zero
because
\(QD_UG'R=\bar DQ'G'R=0\).
Linearity of \(\mathcal B\)
supplies the remaining
current difference.
Multiplying by the inverse
proves
\begin{equation}
 \widehat G[J]D_UR
      =D_UG'R+
           \widehat G[J]\mathcal B_U[\delta J]D_UG'R .
 \label{px:current-defect}
\end{equation}
Thus the failure of the
physical gauge identity
is exactly linear in
the independent-current
defect before inversion,
and has the displayed
resolvent dependence
after inversion.

If \(Q\widehat G[J]Q^\sharp\)
is invertible, define
\begin{equation}
 \widehat H[J]=\widehat G[J]Q^\sharp
                         (Q\widehat G[J]Q^\sharp)^{-1}.
 \label{px:explicit-current-lift}
\end{equation}
Transposing \eqref{px:current-defect} and using
\(RG'Q'^\sharp=0\)
proves
\begin{equation}
 RD_U^\sharp\widehat H[J]
        =RG'D_U^\sharp\mathcal B_U[\delta J]\widehat H[J].
 \label{px:off-current-slice-defect}
\end{equation}
The first term after
transposition vanishes
because
\(D_U^\sharp Q^\sharp
=Q'^\sharp\bar D^\sharp\).
At \(J=j_U\) the
right side is zero.
An arbitrary independent
current does not
obey that slice
condition.
Norm estimates for
its defect use the
actual site and
field kernels and
the form in \eqref{px:current-form-no-current-gradient}.

There is an equally explicit
second variant. Let
\(\mathcal C_U^{(2)}\)
be the homogeneous
quadratic part per
unit length of
the actual nonlinear
straightened average,
and normalize its
current form by
\begin{equation}
 \langle v,\Delta^{(2)}[\omega]v\rangle
                      =2\langle\omega,\mathcal C_U^{(2)}(v)\rangle .
 \label{px:constraint-current-form}
\end{equation}
Define it bilinearly by
polarization, and set
\begin{equation}
 \widehat G_1[J]^{-1}
   =\widehat G[J]^{-1}
      -\mathsf S^\sharp
             \Delta^{(2)}[\widehat H[J]^\sharp J]\mathsf S .
 \label{px:second-current-variant}
\end{equation}
The added form annihilates
\(D_UG'R\) on its
right, by \(\mathsf SD_UG'R=0\).
Therefore \eqref{px:current-defect-before-inversion}--\eqref{px:current-defect}
hold with \(\widehat G_1\)
in place of \(\widehat G\).
Its row lift has
the corresponding
slice-defect identity.
The \(\widehat H\)
inside \eqref{px:second-current-variant} is
formed from the
first variant;
replacing it by
the new row lift
would change the
definition.

These identities do not prove
invertibility or bounds.
They do prove the correct
physical cancellations
once those estimates
hold. They therefore apply to the
\hyperlink{named-gauge-restoration}{named current completion}
on the common primitive domain.
If a different extension
is chosen, its own
defect must be
calculated from
its definition.

\subsubsection{An inner action and a compression have different Ward terms}

For an inner-plaquette action
the gauge calculation in \eqref{px:physical-ward}
is exact on its active
graph. A different
operation first takes
the Hessian of a
larger action and
then compresses it.
Let \(\chi\) be
restriction of fine
fields to a set
of bonds, and
let \(\chi_0\)
restrict sites.
Writing \(D_{\rm loc}\)
for a proposed
local gauge generator,
the defect of
commuting compression
and gauge variation is
\begin{equation}
 E_D=\chi D_U\chi_0^\sharp-D_{\rm loc}.
 \label{px:gauge-compression-defect}
\end{equation}
Even if \(E_D=0\), the
full Hessian can couple
through removed bonds.
The compressed Ward
identity reads
\begin{equation}
 \begin{split}
 (\chi\mathscr D_U\chi^\sharp)D_{\rm loc}
   ={}&\chi\mathcal W_U[j_U]\chi_0^\sharp\\
     &-\chi\mathscr D_U(1-\chi^\sharp\chi)
                          D_U\chi_0^\sharp
       -\chi\mathscr D_U\chi^\sharp E_D .
 \end{split}
 \label{px:compressed-ward}
\end{equation}
This follows by substituting
\eqref{px:gauge-compression-defect}, inserting
\(1=\chi^\sharp\chi+
(1-\chi^\sharp\chi)\),
and applying \eqref{px:physical-ward}.
The extra terms are
supported near the
boundary when the
full operator is
local; with an
exponentially local
operator they have
the corresponding
boundary tails.
They are not part
of the physical
interior current
unless the action
has actually been
changed to the
inner action.

There is a parallel distinction
for lower bounds.
Compressing a positive
form to free coordinates
by zero extension
preserves its gap:
\begin{equation}
 A\ge cI\quad\Longrightarrow\quad
                         \chi A\chi^\sharp\ge cI .
 \label{px:compression-gap}
\end{equation}
Rebuilding its Laplacian,
row averages or block
means on the clipped
region is another
operator. For example
take a unit physical
block with fine spacing
\(\eta\), its discrete
Laplacian and its
full-block constant
projection. Their sum
has a fixed physical
gap after a fixed
normalization.
Keep one fine site,
drop all crossing
Laplacian edges as
in a Neumann cut,
and retain only
the compressed
full-block mean.
The remaining
eigenvalue is
the fine-volume
fraction \(\eta^4\),
which tends to
zero. In the
true Dirichlet
compression the
crossing-edge
energy remains;
this example
does not
contradict
\eqref{px:compression-gap}.

Thus every local inverse theorem
specifies the order:
form the actual
operator with its
declared rows and
gauge, then compress
to its free region.
A local construction
with modified rows
instead supplies
its complete
boundary comparison
and its own gap
proof. A collar
can make those
differences small
at the rows
that are used;
it does not
erase them
as identities.

\subsubsection{The site readout at an outward row}

The site average used in the
Landau projector has one value
at each top, taken at that
top's assigned level. The
endpoint readout of a row
is a different list when
the row is outward.
We now give their exact
relation.

Let \(\mathsf T_j(y,x)\)
be the product of the
admitted internal mean
transports from a level-\(j\)
center \(y\) down its
pillar to a bare site \(x\).
For a top of level \(j\)
define
\begin{equation}
 (Q'_{\rm pill}\lambda)(y)
   =L^{-4j}\sum_{x\in B^j(y)}
                 \operatorname{Ad}(\mathsf T_j(y,x))\lambda(x).
 \label{px:own-pillar-site-average}
\end{equation}
These pillars are disjoint.
There is therefore an explicit
right inverse on node fields:
\begin{equation}
 (E_{\rm pill}\psi)(x)
      =\operatorname{Ad}(\mathsf T_j(y,x))^{-1}\psi(y)
       \quad (x\in B^j(y)),\qquad
                 Q'_{\rm pill}E_{\rm pill}=1 .
 \label{px:site-right-inverse}
\end{equation}
This proves rank without
an elliptic estimate.
For real transports the
right inverse has
supremum norm one.
With fine site measure
\(\eta^4\) and top
measure \(\xi_y^4\),
it is also an
isometry in \(L^2\).
Indeed a pillar
contains \(L^{4j}\)
bare sites, so
\(\eta^4L^{4j}=\xi_y^4\).
The invariant bilinear
color pairing gives
\begin{equation}
 Q'_{\rm pill}=E_{\rm pill}^\sharp,\qquad
                  Q'_{\rm pill}(Q'_{\rm pill})^\sharp=1.
 \label{px:site-weighted-adjoint}
\end{equation}
The algebraic identities
hold for admitted complex
transports as well.
Their Hermitian norm
bounds then include
the actual condition
numbers of those
transports.
A counting kernel
must include the
ratio \(\xi_y^4/\eta^4\)
in the transpose;
an ordinary
unweighted array
transpose would
not be \eqref{px:site-weighted-adjoint}.

For any level-\(j\) center
\(z\), write \(Q'_j\lambda(z)\)
for its full \(j\)-fold
transported site average,
whether or not its assigned
top level is \(j\).
If \(z\) is the outer
endpoint of a determining
level-\(j\) row, its
assigned level is \(j-1\).
The entire block \(B^j(z)\)
lies in that outer
layer. Decompose it
into its \(L^4\)
level-\((j-1)\)
pillars. The product
definition of the
transport gives
\begin{equation}
 Q'_j\lambda(z)
     =L^{-4}\sum_{w\in B(z)\cap\mathbb T^{(j-1)}}
        \operatorname{Ad}(T_{z,w}(\mathcal M^{j-1}U))
                               (Q'_{\rm pill}\lambda)(w).
 \label{px:outward-site-extension}
\end{equation}
To verify the coefficient,
each smaller average has
weight \(L^{-4(j-1)}\);
multiplying it by
\(L^{-4}\) gives
\(L^{-4j}\).
The two transport
factors multiply
in their path order.
This proves \eqref{px:outward-site-extension}
exactly, without
an asymptotic
or small-field
approximation.

Define \(\mathcal E_{j,z}\psi\)
to be \(\psi(z)\) at an
assigned level-\(j\) top,
and to be the finite
average on the right
of \eqref{px:outward-site-extension} at an
outer endpoint.
Then define the extended
coarse difference by
\begin{equation}
 (\bar D_{\rm ext}\psi)(c)
   =\xi_c^{-1}\{
          \mathcal E_{j,c_-}\psi
               -\operatorname{Ad}((\mathcal M^jU)(c))
                          \mathcal E_{j,c_+}\psi\},
                         \qquad c\in\Gamma_j .
 \label{px:extended-coarse-difference}
\end{equation}
For an interior row it
is the ordinary covariant
difference of its two
top values. For an
outward row it
first averages the
outer finer tops.
There is no need
to duplicate that
site as a new
independent top
in the site Gram
matrix.

Here is a direct verification
of the field intertwining.
Under an infinitesimal
gauge \(e^{it\lambda}\)
the derivative of the
level-\(j\) node factor
is
\begin{equation}
 \left.\frac1i\frac d{dt}v_j[tD_U^{\rm phys}\lambda](y)\right|_{t=0}
                         =\lambda(y)-Q'_j\lambda(y).
 \label{px:node-factor-gauge-derivative}
\end{equation}
At one step the change in
each relative contour has
derivative
\(\lambda(y)-\operatorname{Ad}(T_{y,x})\lambda(x)\).
The implicit mean at the
identity has differential
the arithmetic mean of
its inputs, so this
gives the one-step
formula. Iterating
and using the product
rule for the
accumulated node
factors gives
\eqref{px:node-factor-gauge-derivative}, with the
hierarchical site
average in \eqref{px:own-pillar-site-average}.

Differentiate the actual
straightening identity
\(\mathcal M^j(e^{i\eta A}U)
 =v_j[A]\,e^{iq_j^{\rm st}(A)}
          \mathcal M^jU\,v_j[A]^{-1}\)
on a row.
The plain coarse
variation is
\(\lambda(c_-)
-\operatorname{Ad}(\mathcal M^jU(c))\lambda(c_+)\).
Subtract the two
node derivatives
\eqref{px:node-factor-gauge-derivative}, with their
endpoint transports.
The remaining variation
is
\begin{equation}
 Dq_j^{\rm st}(0)\,D_U^{\rm phys}\lambda
     =Q'_j\lambda(c_-)
             -\operatorname{Ad}(\mathcal M^jU(c))Q'_j\lambda(c_+).
 \label{px:raw-straightened-intertwining}
\end{equation}
Divide by \(\xi_c\)
to obtain the
per-length row \(Q\).
Using \eqref{px:outward-site-extension} proves
the fully typed identity
\begin{equation}
                 QD_U^{\rm phys}
                          =\bar D_{\rm ext}Q'_{\rm pill}.
 \label{px:typed-seam-intertwining}
\end{equation}
This supplies the
intertwining needed
in \eqref{px:explicit-current-green}--\eqref{px:off-current-slice-defect}.
There \(\bar D\)
means this extended
row difference and
\(Q'\) means the
own-pillar site map.
Using a single
unextended value
of the outer
top would give
the wrong formula.

Each endpoint extension reads
at most \(L^4\) tops in
one adjacent layer.
Its real row sum is
at most one, since
the transports are
unitary and the
weights add to one.
A finer top occurs
in one such outer
block and at most
eight incident
coarse rows.
Thus its two-sided
normalized cell
counts are bounded
at fixed \(L\),
without a factor
depending on the
depth.
For complex pairs
the same statement
uses the actual
bounded transport
condition numbers.
The latter are
atomic domain
estimates; the
algebraic
intertwining does
not prove them.

\subsubsection{Finite real existence of the site operators}

There is also a real finite positivity
statement independent of a quantitative
inverse bound. Assume the actual
mean transports are admitted and
unitary, and that every active fine
site lies in a connected pillar
with its top. Let \(a\) be
strictly positive on the
top space. Then
\begin{equation}
 L_{\rm site}=D_U^\sharp D_U+
                    (Q'_{\rm pill})^\sharp aQ'_{\rm pill}
                         \quad\text{is strictly positive}.
 \label{px:finite-site-positivity}
\end{equation}
To prove it, a vector of zero
Hermitian energy must obey
\(D_U\lambda=0\) and
\(Q'_{\rm pill}\lambda=0\).
Split into Hermitian real
and imaginary parts;
it suffices to treat
real \(\lambda\).
The first identity means that
\(e^{it\lambda(x)}\) leaves
every active bond unchanged:
\begin{equation}
 e^{it\lambda(b_-)}U_b
                        e^{-it\lambda(b_+)}=U_b.
 \label{px:parallel-gauge-stabilizer}
\end{equation}
Gauge covariance of each actual
mean transport therefore gives
\(\operatorname{Ad}(\mathsf T_j(y,x))\lambda(x)
 =\lambda(y)\).
Average over the pillar:
\(Q'_{\rm pill}\lambda(y)=\lambda(y)\).
Its zero value makes
\(\lambda(y)=0\).
Parallel transport along
the active graph from
the top to every
site makes
\(\lambda=0\).
This proves \eqref{px:finite-site-positivity}.

It follows that \(G'=L_{\rm site}^{-1}\)
exists on every finite real
problem of this type.
The site Gram matrix
\begin{equation}
 Q'_{\rm pill}G'^2(Q'_{\rm pill})^\sharp
 \label{px:finite-site-gram}
\end{equation}
is strictly positive as well.
For a nonzero node vector
\(\psi\), surjectivity makes
\((Q'_{\rm pill})^\sharp\psi\ne0\);
invertibility of \(G'\)
then makes its squared
norm positive.
Thus the formulas for
\(R\) and \(\mathsf S\)
are genuine finite
real matrices before
their uniform bounds
are established.

No volume-uniform or scale-weighted
lower bound has been deduced
from this nullspace argument.
That bound needs the
quantitative site Poincaré
and local-gauge estimates.
Nor does \eqref{px:finite-site-positivity} prove
positivity of the fine
Wilson Hessian or
of its independent-current
completion.
The distinctions between
finite invertibility,
uniform coercivity and
exponentially local
readouts remain
necessary.

\subsubsection{Transport control on one pillar}

The site gap has a quantitative proof
from the actual averages. Let a
pillar have side \(\xi=n\eta\),
where \(n=L^j\), and assume
\begin{equation}
 |\partial U(p)-1|\le
                 \epsilon\,\eta^2\xi^{-2}
 \label{px:pillar-curvature}
\end{equation}
on its determining neighborhood.
The fixed path neighborhood has
an embedded rectangular lift.
If an extra fine collar is
needed, the regularity bounds
in the adjacent layer give
the same inequality with a
fixed \(L\)-dependent factor.
We include that factor in
the constants below.

Take the fine coordinate comb
from the pillar center and
put \(g(x)=U(\Gamma_{y,x})\).
For adjacent \(x,x+\eta e_\mu\),
the loop made by the two
root combs and that bond
is filled by strips of
width one fine cell.
Commuting the extra
\(\mu\)-step through
the later coordinate
segments uses at most
\(Cn\) plaquettes,
not \(Cn^2\).
The unitary Stokes
product inequality
therefore gives
\begin{equation}
 |U^g(b)-1|\le C\epsilon/n
                  \quad\text{on internal fine bonds}.
 \label{px:fine-comb-strip}
\end{equation}
Here \(g(y)=1\).
Every curvature norm is unchanged.
The gauge need not be
close to one, and its
derivatives are not
estimated in this step.

Write \(\xi_r=L^r\eta\),
and let \(a_r\) be the
supremum of
\(|\mathcal M^r(U^g)(c)-1|\)
on the level-\(r\)
internal edges used
in the hierarchical
transport. Let
\(\varpi_r\) be the
corresponding coarse
plaquette bound.
The determining
neighborhoods remain
in the rectangular
pillar neighborhood:
each parent block
is a union of
its child blocks,
and each contour
stays in its
assigned block.
The real curvature
iteration proved
in the background
section gives
\begin{equation}
 \varpi_r\le C_L\epsilon(\xi_r/\xi)^2 .
 \label{px:pillar-coarse-curvature}
\end{equation}
The same finite comparison
shows that one actual
averaged edge differs
from its central
straight transport
by at most
\(C_L\varpi_r\).
Indeed every closed
logarithm in its
definition is a
comparison loop
with that central
transport. The
unitary strip
estimate bounds
those logarithms,
and then their
exponential average.
The central route
has \(L\) edges,
so
\begin{equation}
       a_{r+1}\le L a_r+C_L\varpi_r .
 \label{px:pillar-open-recurrence}
\end{equation}
The coefficient of \(a_r\)
is exactly \(L\), by
the unitary product
inequality.

Iterating \eqref{px:pillar-open-recurrence}, using
\(a_0\le C\epsilon\eta/\xi\),
gives
\begin{equation}
 \begin{split}
 a_r&\le L^ra_0+
       C_L\epsilon\sum_{s<r}
                 L^{r-1-s}(\xi_s/\xi)^2\\
     &\le C_L\epsilon\,\xi_r/\xi .
 \end{split}
 \label{px:pillar-open-bound}
\end{equation}
For the sum, insert
\(\xi_s=\eta L^s\);
its geometric factor is
bounded by a constant
times \((\xi_r/\xi)^2\),
which is no larger
than \(\xi_r/\xi\).
The constants do
not multiply with
the number of
levels.

The internal mean on
level \(r\) is a
finite smooth mean
of paths of \(O(L)\)
edges. On its
regular near-identity
chart it differs
from one by at
most \(C_La_r\).
The hierarchical
transport has one
such factor at
each level below
\(j\). All are
unitary, hence
\begin{equation}
 |\mathsf T_j(U^g;y,x)-1|
       \le C_L\sum_{r<j}a_r
       \le C_L\epsilon
                    \quad(x\in B^j(y)).
 \label{px:uniform-pillar-transport}
\end{equation}
The last sum uses
\(\sum_{r<j}\xi_r/\xi
  \le1/(L-1)\).
The fixed smallness
of \(\epsilon\)
also admits every
mean used in this
argument. Thus
\eqref{px:uniform-pillar-transport} is a
uniform-depth
property of
the actual
iterated transports,
not a hypothesis
of such a property.

\subsubsection{A real weighted site gap}

We first prove the flat block
inequality with its scale.
Let \(\bar\lambda\) be the
arithmetic mean on the
\(n^4\) sites of a block.
The exact variance identity is
\begin{equation}
 \sum_x|\lambda(x)-\bar\lambda|^2
       =\frac1{2n^4}\sum_{x,z}|\lambda(x)-\lambda(z)|^2 .
 \label{px:block-variance}
\end{equation}
Connect each pair by the
coordinate-order path.
Its length is at most
\(4n\). A fixed
oriented edge lies
in at most \(n^5\)
of these paths:
its two longitudinal
endpoint coordinates
give at most \(n^2\)
choices, and the
three other coordinates
give at most \(n^3\).
Cauchy on the path
and then \eqref{px:block-variance}
give a constant
times \(n^2\)
times the sum
of squared
edge differences.
Multiplying by
\(\eta^4\xi^{-2}\)
proves
\begin{equation}
 \xi^{-2}\|\lambda\|_{2,B}^2
  \le C\big\{
         \|\partial_{\rm int}\lambda\|_{2,B}^2
                              +\xi^2|\bar\lambda|^2\big\}.
 \label{px:flat-block-poincare}
\end{equation}
Here the fine norm uses
\(\eta^4\), and
\(\partial\) divides a
bond difference by
\(\eta\).
The mean part is
\(\eta^4n^4\xi^{-2}
|\bar\lambda|^2
=\xi^2|\bar\lambda|^2\).
For \(n=1\) the
variance vanishes
and the same
inequality holds.

Apply the real comb gauge
of \eqref{px:fine-comb-strip} in each pillar.
Only internal bonds are
used in the block
inequality, so the
gauges of neighboring
pillars may be
different. They
define a unitary
site gauge globally,
and all covariant
norms are unchanged.
In this frame
\begin{equation}
 \|(\partial_{\rm int}-D_{U,{\rm int}})\lambda\|_{2,B}
                       \le C_L\epsilon\,\xi^{-1}\|\lambda\|_{2,B}.
 \label{px:block-gradient-comparison}
\end{equation}
The estimate follows from
\(|\operatorname{Ad}(U_b)-1|
 \le C_L\epsilon\eta/\xi\)
and the at-most-eight
incidences at a site.

Let \(Q'_0\) be arithmetic
block average. \eqref{px:uniform-pillar-transport}, Cauchy for
the finite average
and \(\xi^4=n^4\eta^4\)
give
\begin{equation}
 \xi^4|(Q'_{\rm pill}-Q'_0)\lambda|^2
                        \le C_L\epsilon^2\|\lambda\|_{2,B}^2 .
 \label{px:block-average-comparison}
\end{equation}
Use \eqref{px:block-gradient-comparison}--\eqref{px:block-average-comparison} on
the right of \eqref{px:flat-block-poincare}.
The result is the
same two terms
with \(D_U\) and
\(Q'_{\rm pill}\),
plus
\(C_L\epsilon^2\xi^{-2}
\|\lambda\|_{2,B}^2\).
Choose the fixed
curvature ceiling
so that this
coefficient is
at most one half
of the left
coefficient.
Absorption proves
\begin{equation}
 \xi^{-2}\|\lambda\|_{2,B}^2
 \le C_L\big\{
       \|D_{U,{\rm int}}\lambda\|_{2,B}^2
                         +\xi^2|Q'_{\rm pill}\lambda|^2\big\}.
 \label{px:covariant-block-poincare}
\end{equation}

Sum over the disjoint pillars.
Every internal fine bond
is counted at most once;
all other gradient
terms in the full
energy are nonnegative.
Suppose
\(a(y)=c_y\xi_y^{-2}\)
with \(c_y\ge a_*>0\).
Then
\begin{equation}
 \begin{split}
 \sum_x\eta^4\xi(x)^{-2}|\lambda(x)|^2
 &\le C_{L,a_*}\big\{
        \|D_U\lambda\|_2^2
             +\langle Q'_{\rm pill}\lambda,
                              aQ'_{\rm pill}\lambda\rangle_{\rm tops}\big\}\\
 &=C_{L,a_*}\langle\lambda,L_{\rm site}\lambda\rangle .
 \end{split}
 \label{px:real-weighted-site-gap}
\end{equation}
This is the actual site
weighted gap, uniform
in depth, volume,
region shape and
real holonomy.
It uses the small
real curvature
and actual transport
construction; no
physical current
or criticality
is used.
Crossing-edge
energy is kept,
not replaced
by a clipped
block mean.

For a complex site vector
split into Hermitian
real and imaginary
parts. Since the
real unitary
connection is a
real matrix in
an orthonormal
Lie-algebra basis,
the same inequality
holds for its
Hermitian energy.
If the top scale
is at most one,
it in particular
gives an ordinary
finite \(L^2\)
gap independent
of the bare
mesh.
In other physical
units its smallest
unweighted gap
is scaled by
the inverse
square of the
largest local
scale.

\subsubsection{The complex site domain uses the actual iterated means}

The preceding transport argument has a
complex version with a fixed physical
radius, independent of the number
of fine sites. We give it to avoid
assuming a uniform radius for the
iterated site average.

Consider one step on a fixed
two-block neighborhood. Suppose
all its edge matrices satisfy
\(|W_b-1|\le a\), and let
\(\varpi=\sup|\partial W-1|\).
There are constants \(r_L,C_L\)
for which the actual means
and logarithms are admitted
if \(a,\varpi<r_L\), and
\begin{equation}
 \begin{split}
 a_{\rm next}&\le La+C_La^2+C_L\varpi,\\
 \varpi_{\rm next}&\le
       L^2(1+C_La)\varpi+C_L\varpi^2 .
 \end{split}
 \label{px:complex-one-step-recurrences}
\end{equation}
For the first line, a central
straight product differs
from one by at most
\(La+C_La^2\).
Plaquette sliding bounds
each closed comparison
logarithm by \(C_L\varpi\):
all prefix transports
have condition number
at most \(1+C_La\)
on this finite
neighborhood.
The defining exponential
average therefore changes
the central product
by at most that
amount.

For the second line, take the
local complex comb gauge.
Its values and inverses
are \(1+O(C_La)\).
It sets the tree
bonds to one; each
remaining bond is
then \(1+O(C_L\varpi)\)
by the finite
plaquette sliding
identity.
In that gauge the
Taylor/flux argument
for the actual
one-step map gives
the coefficient
\(L^2\) and a
quadratic error
\(C_L\varpi^2\).
Conjugating the
fine and coarse
plaquettes through
the small complex
gauge changes their
norms by
\(1+C_La\).
This proves
\eqref{px:complex-one-step-recurrences}.
Complex covariance
here follows directly
from the conjugation
identity for the
small logarithm,
or from uniqueness
of the admitted
finite mean.
No absolute-value
formula is
continued
holomorphically.

These constants can be bounded
before an iteration is run.
The finite mean construction
has a joint raw-edge
holomorphic ball of radius
\(c/L\), with universal
supremum bounds; two scalar
Cauchy circles bound its
second derivative by
\(CL^2\).
A comb on the one-step
patch sweeps at most
\(CL^2\) plaquettes,
so its non-tree
logarithms have size
\(CL^2\varpi\).
The quadratic Taylor
error in the coarse
plaquette is therefore
at most
\(CL^2(L^2\varpi)^2
=CL^6\varpi^2\).
The fine logarithmic
curl error times
its \(L^2\) flux
factor has the
same upper order.
Prefix conjugations
cost \(1+CLa\).
For example
\(r_L=cL^{-3}\)
and \(C_L=CL^8\),
with fixed universal
shrinkages and
enlargements, suffice
in \eqref{px:complex-one-step-recurrences}.
This count uses
directional Cauchy
on the proved
full supremum
ball, not one
coordinate estimate
summed over all
field coordinates.

On a pillar let \(t_r=\xi_r/\xi\),
and normalize
\(A_r=a_r/t_r\),
\(V_r=\varpi_r/t_r^2\).
\eqref{px:complex-one-step-recurrences} implies,
with an increased fixed constant,
\begin{equation}
 Z_{r+1}\le Z_r+C_Lt_r(Z_r+Z_r^2),
            \qquad Z_r=A_r+V_r,\qquad
                 \sum_{r<j}t_r\le\frac1{L-1}.
 \label{px:complex-transport-gronwall}
\end{equation}
For \(Z_r\le1\), the factor
is at most \(1+2C_Lt_r\).
Its product is at most
\(\exp(2C_L/(L-1))\).
Choose \(Z_0\) below
one half the minimum
one-step chart radius
times the inverse
of this number.
Induction then gives
\(Z_r\le C_LZ_0\)
and keeps every step
in its chart.
The constants depend
on \(L\), not \(j\).
With the preceding finite
choices a sufficient
common small input
radius has the explicit
form
\(cL^{-3}\exp(-CL^8)\).
The precise numerical
constants are not
needed below, but
this gives a
positive lower
radius rather
than an asserted
upper cap.
The product of the
internal means and
its inverse have
norm at most
\(\exp(C_L\sum_ra_r)\).
Hence
\begin{equation}
 |\mathsf T_j(W;y,x)-1|\le C_LZ_0,\qquad
 |\mathsf T_j(W;y,x)|+
                  |\mathsf T_j(W;y,x)^{-1}|\le C_L .
 \label{px:complex-pillar-transport}
\end{equation}
This proves the existence
and boundedness of
the whole finite
hierarchical formula
on that domain.

Now start with a real field
satisfying \eqref{px:pillar-curvature}, put it
in the real comb gauge
\eqref{px:fine-comb-strip}, and let
\(W=e^{i\eta E}U\).
Assume the actual physical
first-order input satisfies
\begin{equation}
 h=\max\{\|\xi E\|_\infty,
                    \|\xi^2\nabla_UE\|_\infty\}\le h_L .
 \label{px:site-complex-input}
\end{equation}
In the comb frame the ordinary
first difference of \(E\)
differs from its covariant
one by
\(O(\epsilon\xi^{-1}|E|)\).
The elementary plaquette
product then gives
\(a_0\le C(\epsilon+h)\eta/\xi\)
and
\(\varpi_0\le
C(\epsilon+h+h^2)(\eta/\xi)^2\).
Thus the starting \(Z_0\)
in \eqref{px:complex-transport-gronwall} is at most
\(C(\epsilon+h)\).

Choose a fixed \(h_L>0\)
so that the same argument
holds with \(E\) replaced
by \(sE\) for
\(|s|<h_L/h\), and
assume \(h\le h_L/4\).
Reduce the real curvature
ceiling at the same
time so that
\(\epsilon\le h_L\)
and the resulting
starting \(Z_0\)
satisfies \eqref{px:complex-transport-gronwall}.
The transport is
holomorphic in \(s\)
on this disc, with
a uniform bound.
Cauchy's formula
on a smaller
concentric disc,
or the geometric
sum of its Taylor
coefficients, yields
\begin{equation}
 |\mathsf T_j(e^{i\eta E}U;y,x)
                       -\mathsf T_j(U;y,x)|\le C_Lh .
 \label{px:site-transport-lipschitz}
\end{equation}
The bound in \eqref{px:complex-pillar-transport} is
uniform in \(j\) and
in the number of
coordinates. This
is a single
directional Cauchy
argument on an
actual admitted
physical disc.
It does not
assert an
entrywise absolute
Hessian sum for
an arbitrary
many-variable
functional.

For the complex site average
write \(Q'_E\) and for
the real reference write
\(Q'_U\). The transported
finite-average calculation
used in \eqref{px:block-average-comparison} gives
\begin{equation}
 \|a^{1/2}(Q'_E-Q'_U)\lambda\|_{2,{\rm tops}}
                   \le C_Lh\,\|\xi^{-1}\lambda\|_2 .
 \label{px:complex-site-average-difference}
\end{equation}
Here the real row weights
satisfy
\(a(y)\le a^*\xi_y^{-2}\).
The proof squares the
average difference in
each disjoint pillar,
uses \eqref{px:site-transport-lipschitz}, and
sums its fine volume.
There is no factor
equal to the
number of
fine sites.
The same argument
applies to the
outer endpoint
extension \eqref{px:outward-site-extension}:
its neighborhood
uses one adjacent
scale, so the
input constants
change by fixed
powers of \(L\).

\subsubsection{Relative forms give the complex site inverse}

It remains to compare the actual
site operators, not their
unscaled matrix coefficients.
Let \(L_U\) be the real
site operator and let
\(L_E=D_{e^{i\eta E}U}^\sharp
D_{e^{i\eta E}U}+Q_E'^\sharp aQ'_E\).
The transposes in \(L_E\)
are bilinear analytic
transposes. At a real
reference they are
the usual Hermitian
adjoints.

For a fine bond the difference
of the two site derivatives
is bounded by
\(C|E_b||\lambda(b_+)|\).
Adjacent pillar scales
differ by at most the
fixed factor \(L\);
the region gaps exclude
a bare bond jumping
over another layer.
The finite incidence
count therefore gives
\begin{equation}
 \|(D_{e^{i\eta E}U}-D_U)\lambda\|_2
                        \le C_Lh\|\xi^{-1}\lambda\|_2 .
 \label{px:complex-site-gradient-difference}
\end{equation}
Keeping derivatives on the
two arguments, expand
the kinetic difference.
Use \eqref{px:complex-site-gradient-difference}, then
\eqref{px:complex-site-average-difference} for the mass
difference, and finally
the real weighted
site gap in \eqref{px:real-weighted-site-gap}.
For arbitrary complex
vectors \(v,w\),
this gives
\begin{equation}
 |\langle v,(L_E-L_U)w\rangle_{\rm H}|
 \le C_Lh\,
      \langle v,L_Uv\rangle_{\rm H}^{1/2}
      \langle w,L_Uw\rangle_{\rm H}^{1/2}.
 \label{px:relative-complex-site-form}
\end{equation}
For example the quadratic mass
term is
\((Q'_E-Q'_U)^\sharp
a(Q'_E-Q'_U)\); the
two cross terms contain
only one difference.
On replacing the left
argument by its complex
conjugate, the same
absolute bounds apply
to the analytic
bilinear transpose.
No Hermitian
conjugation of
the complex
coefficients is
inserted into
their definition.

Let
\(B_E=L_U^{-1/2}(L_E-L_U)L_U^{-1/2}\).
\eqref{px:relative-complex-site-form} proves
\(\|B_E\|_{2\to2}\le C_Lh\).
For \(C_Lh\le1/2\)
one has
\begin{equation}
 G'_E=L_U^{-1/2}(1+B_E)^{-1}L_U^{-1/2},\qquad
 \operatorname{Re}\langle\lambda,L_E\lambda\rangle_{\rm H}
                  \ge\tfrac12\langle\lambda,L_U\lambda\rangle_{\rm H}.
 \label{px:complex-site-inverse}
\end{equation}
The inverse is holomorphic
on the whole admitted
physical ball.
It has the same
weighted coercivity,
up to a fixed
factor, uniformly
in depth and
volume.
This is a
relative-form
argument. The
unscaled operator
norm of a
coefficient
perturbation
need not be
bounded as
\(\eta\downarrow0\).

The proof establishes the site
inverse part of the
primitive package.
It has not yet
established the
site Gram inverse,
the fine vector
Green function,
or the layered
current precision.
Those require
their separate
row, readout
and gap
arguments.

\subsubsection{Exponential decay from the weighted site form}

Let \(V\) denote multiplication by
\(\xi(x)^{-2}\) on sites.
The real gap in \eqref{px:real-weighted-site-gap} is
\(L_U\ge cV\).
We prove a weighted inverse
bound before converting it
to pointwise regularity.
This avoids estimating a
bare kinetic operator
by its unscaled norm.

Let \(\psi\) be bounded and
one-Lipschitz for \eqref{px:layer-metric},
and let \(M_s=e^{s\psi}\)
on sites. Use the value
\(\psi(b_-)\) on an edge
and \(\psi(y)\) on a
top. Put
\begin{equation}
 D_s=M_s^{\rm edge}D_UM_s^{-1},\qquad
 Q'_s=M_s^{\rm top}Q'_UM_s^{-1}.
 \label{px:site-weighted-rows}
\end{equation}
On an edge the two endpoint
values of \(\psi\) differ by
at most \(C_L\eta/\xi_b\).
Inside a pillar their
difference from the
top value is at most
a fixed constant.
The difference from
the unweighted row
therefore satisfies
\begin{equation}
 \|(D_s-D_U)\lambda\|_2
  +\|a^{1/2}(Q'_s-Q'_U)\lambda\|_{2,{\rm tops}}
               \le C_L|s|\|V^{1/2}\lambda\|_2
 \label{px:site-weighted-row-difference}
\end{equation}
for \(|s|\le s_L\).
The edge estimate follows
from \(|e^{st}-1|\le
C|s||t|\) on the
stated bounded interval.
For the site row use
the same estimate
in its finite
transported average,
then Cauchy with
the actual pillar
volume. Its
constant is
independent of
\(\|\psi\|_\infty\).

The conjugated form has the
exact factorization
\begin{equation}
 M_sL_UM_s^{-1}
              =D_{-s}^\sharp D_s+(Q'_{-s})^\sharp aQ'_s .
 \label{px:conjugated-site-form}
\end{equation}
Expanding it and using
\eqref{px:site-weighted-row-difference}, \eqref{px:real-weighted-site-gap} proves
a relative form error
at most \(C_L|s|\).
With a complex field
perturbation in the
ball of \eqref{px:complex-site-inverse},
the same calculation
gives error
\(C_L(h+|s|)\)
relative to the
real \(L_U\).
Choose both small
enough for that
number to be
at most one half.
All these forms
are invertible.

For any such inverse \(\mathcal G_s\),
testing the equation
\(L_s\lambda=V^{1/2}f\)
against \(\lambda\)
gives
\begin{equation}
 \|V^{1/2}\mathcal G_sV^{1/2}\|_{2\to2}
                                        \le 2/c .
 \label{px:weighted-site-inverse}
\end{equation}
Indeed its real part is at
least \((c/2)\|V^{1/2}\lambda\|_2^2\),
whereas the right
pairing is at most
\(\|V^{1/2}\lambda\|_2\|f\|_2\).
This proof uses neither
a normal matrix nor
a pointwise positive
complex kernel.

Choose
\(\psi(x)=\min\{d(x,\Delta(z)),N\}\).
It is one-Lipschitz.
The diameter of a
pillar in \eqref{px:layer-metric}
is bounded, so
on \(\Delta(y)\)
its value is at
least \(d(y,z)-C\)
after \(N\) is
large enough.
\eqref{px:weighted-site-inverse}
then yields
\begin{equation}
 \|P_yG'_EP_z\|_{2\to2}
       \le C_L\xi_y\xi_z e^{-m d(y,z)}
 \label{px:site-cell-l2-decay}
\end{equation}
for a fixed \(m=m(L)>0\).
On a finite torus take
the periodic distance.
On the cover first
use bounded \(N\)
and then increase
it. The constants
are independent
of this truncation.

In particular the normalized
operator \(V^{1/2}G'_EV^{1/2}\)
has exponentially decaying
cell \(L^2\) entries.
Its two-sided cell
Schur norm follows
from \eqref{px:metric-ball-count}--\eqref{px:summability-rate}
when the separation
is chosen with
\begin{equation}
             \frac{8\log L}{\rho}<m/2 .
 \label{px:site-separation-choice}
\end{equation}
The rate \(m\) above depends
only on the fixed one-step
constants and the site
gap; it does not depend
on \(\rho\). Thus this
is a genuine order of
choice: increase the
separation after those
constants, or verify
that the already chosen
separation satisfies
this finite inequality.
It is not an appeal
to unspecified
polynomial volume
growth.

\subsubsection{A scalar subsolution estimate at the fine mesh}

We use a local estimate whose proof
also explains why a fine Coulomb
singularity causes no loss in
\eqref{px:site-cell-l2-decay}. On the infinite scalar
fine lattice in four dimensions,
let \(\mathcal L=-\Delta_\eta\).
The positive heat kernel gives
a positive Green kernel
\begin{equation}
 G_0(x)=\int_0^\infty p_t^\eta(x)\,dt,\qquad
 |G_0(x)|\le C(|x|+\eta)^{-2},\quad
 |\nabla G_0(x)|\le C(|x|+\eta)^{-3}.
 \label{px:scalar-massless-kernel}
\end{equation}
These inequalities follow by
integrating the mesh-uniform
heat kernel derivative
bounds. For example split
at \(t=|x|^2\), use
a polynomial off-diagonal
tail before the split,
and \(Ct^{-2}\), respectively
\(Ct^{-5/2}\), afterwards.
At \(|x|\le\eta\)
the bounds are the
corresponding lattice
scaling bounds.
Also \(\mathcal LG_0\)
is the normalized
lattice delta, by
integrating
\(-\partial_tp_t^\eta\).

Suppose \(w\ge0\) and
\(\mathcal Lw\le F_0\)
on a ball of radius
\(2r\), \(r\ge8\eta\).
Choose a nonnegative
cutoff \(\chi\),
equal to one on
the \(r\)-ball,
supported in the
\(2r\)-ball, with
its first two
differences bounded
by \(Cr^{-1}\),
\(Cr^{-2}\).
At the center
\(x\), summation
by parts gives
\begin{equation}
 w(x)=
       \langle\chi G_0(x-\cdot),\mathcal Lw\rangle
             -\langle[\mathcal L,\chi]G_0(x-\cdot),w\rangle .
 \label{px:subsolution-cutoff-identity}
\end{equation}
All sums use \(\eta^4\).
This identity follows by
moving \(\mathcal L\)
off \(w\) and using
\(\mathcal LG_0=\delta_x\).
There is no derivative
of \(w\) in its
last term.
Since the first
test is nonnegative,
its contribution
is at most
\(CF_0r^2\).
Indeed the lattice
sum of \(G_0\)
on that ball
is at most
\(Cr^2\), by
dyadic shells.

The commutator is supported
on the cutoff annulus.
Writing it as
\((\mathcal L\chi)G_0\)
plus products of
first differences
shows that its
absolute value is
at most \(Cr^{-4}\)
there, by \eqref{px:scalar-massless-kernel}.
Cauchy and the
annulus volume
therefore prove
\begin{equation}
 w(x)\le C\{r^{-2}\|w\|_{2,B_{2r}}+r^2F_0\}.
 \label{px:scalar-subsolution-estimate}
\end{equation}
If \(r/\eta\) is bounded,
the point-value estimate
\(|w(x)|\le\eta^{-2}\|w\|_2\)
gives the same inequality
with a changed constant.
Thus the local
estimate is
uniform at
coarse meshes
as well.

For a real unitary connection,
let \(\lambda\) be Lie-algebra
valued and put \(w=|\lambda|\).
For \(\lambda(x)\ne0\), pair
each term
\(\lambda(x)-\operatorname{Ad}(U_{x,x'})\lambda(x')\)
with \(\lambda(x)/|\lambda(x)|\).
Unitarity gives
\begin{equation}
 -\Delta_\eta|\lambda|(x)
       \le |D_U^\sharp D_U\lambda(x)|.
 \label{px:discrete-kato}
\end{equation}
At a zero value of \(\lambda\)
the left side is nonpositive,
so the same inequality
holds. This proves
the covariant scalar
subsolution inequality
without differentiating
a gauge.
It is used only
on a patch with
the full covariant
Laplacian stencil.
Bare boundary
cells will use
the point-value
fallback instead.

\subsubsection{Pointwise site Green bounds}

First take a real \(U\) and
\(\lambda=G'_Uf\).
At a site in pillar
\(\Delta(y)\), the
mass part of its
equation obeys
\begin{equation}
 |(Q_U'^\sharp aQ'_U\lambda)(x)|
                 \le C_L\xi_y^{-4}\|\lambda\|_{2,\Delta(y)} .
 \label{px:site-mass-pointwise}
\end{equation}
In fact it is the transported
vector \(a(y)Q'_U\lambda(y)\),
whose norm is at most
\(a^*\xi_y^{-2}|Q'_U\lambda(y)|\).
Cauchy in the average
gives the additional
\(\xi_y^{-2}\).

Take a local ball of radius
comparable to \(\xi_y\).
It meets a bounded
\(L\)-dependent number
of pillars, all with
scales in
\([\xi_y/L,L\xi_y]\).
The gaps exclude a
second intervening
layer in that
ball once
\(\rho>C_LL\).
Call the corresponding
whole-pillar neighborhood
\(\mathcal N(y)\).
Its cells are at
bounded \(d\)-distance
from \(y\).
Apply \eqref{px:scalar-subsolution-estimate}--\eqref{px:site-mass-pointwise}
on a lift of this
ball. This proves
\begin{equation}
 \|\lambda\|_{\infty,\Delta(y)}
   \le C_L\xi_y^{-2}
                 \sum_{w\in\mathcal N(y)}\|\lambda\|_{2,\Delta(w)}
         +C_L\xi_y^2\|f\|_{\infty,\mathcal N(y)} .
 \label{px:site-local-l2-to-infinity}
\end{equation}
Near a boundary of the action
region, every pillar lacking
this full patch is bare:
the nonbare regions have
the prescribed gap from
that boundary. The
single-cell fallback
in \eqref{px:scalar-subsolution-estimate} proves
\eqref{px:site-local-l2-to-infinity} there directly.
On a torus the local
ball is lifted to
the periodic cover.
Its fixed radius
can meet only a
fixed number of
copies of a cell,
since every period
contains at least
two top blocks.
This changes only
the constant.

If \(f\) is supported in
\(\Delta(z)\), use \eqref{px:site-cell-l2-decay}
and
\(\|f\|_{2,\Delta(z)}
 \le\xi_z^2\|f\|_\infty\)
on the right side.
The far contribution is
\(C\xi_z^3\xi_y^{-1}
e^{-m d(y,z)}\|f\|_\infty\).
The local forcing term
has \(d(y,z)\le C_L\).
\eqref{px:level-barrier} absorbs
the scale ratio:
\begin{equation}
 \|G'_Uf\|_{\infty,\Delta(y)}
       \le C_L\xi_y^2
                e^{-m_1d(y,z)}\|f\|_\infty,
 \qquad 0<m_1<m-\frac{3\log L}{\rho}.
 \label{px:site-pointwise-green}
\end{equation}
Choosing the separation large
enough leaves a positive
fixed \(m_1\).
This is a pointwise
estimate for a
bounded cell source.
It is not a
uniform bound
on the Green
kernel evaluated
at two individual
fine sites.

The actual equation also gives
\begin{equation}
 \|D_U^\sharp D_UG'_Uf\|_{\infty,\Delta(y)}
                  \le C_Le^{-m_1d(y,z)}\|f\|_\infty .
 \label{px:site-laplacian-green}
\end{equation}
Use \(D_U^\sharp D_U\lambda
=f-Q_U'^\sharp aQ'_U\lambda\),
and bound its mass term
by the supremum on its
same pillar. The
cell diameter in
the layer metric
is bounded.

\subsubsection{Site gradients and their complex continuation}

For differentiated pointwise bounds use
the actual local gauge clause of the
regular-pair class. On the required
patch around a cell it supplies
a unitary gauge with
\begin{equation}
 U=e^{i\eta A},\qquad
 \xi|A|+\xi^2|\nabla A|\le a_{\rm loc},
                         \qquad a_{\rm loc}\le a_L .
 \label{px:site-local-coefficient-class}
\end{equation}
The constant \(a_L\) is fixed small.
Near one interface the patch may
be described in the finer
adjacent scale; converting
to \(\xi\) changes these
bounds by fixed powers
of \(L\).
The region gaps and
the fixed patch collar
are chosen to include
no third layer.
This gauge clause is
an input defining
the regular-pair
class. Deriving it
for an actual
background from
weaker tests is
the separate
background
estimate.

In this local gauge the site
equation is a scalar
componentwise Laplace
equation with first-
and zeroth-order
errors. Expanding
the transported
differences gives
\begin{equation}
 |\Delta_\eta\lambda|
 \le |f-Q'^\sharp aQ'\lambda|
       +C_La_{\rm loc}
          \{\xi^{-1}|\nabla\lambda|
                                      +\xi^{-2}|\lambda|\}.
 \label{px:site-scalar-equation}
\end{equation}
Each difference of a
link exponential costs
\(\eta|\nabla A|\);
each undifferentiated
link error costs
\(\eta|A|\).
After division by
the two powers of
\(\eta\) in the
Laplacian the
terms have exactly
the displayed
dimensions.

Here is the interior absorption
that supplies the gradient.
On a slightly larger ball
of radius \(C\xi\), put
\(M=\|\lambda\|_\infty\)
and
\(F=\|f-Q'^\sharp aQ'\lambda\|_\infty\).
Let \(d_b\) be distance
of a fine edge to
the ball boundary
and take the finite
maximum of
\(d_b|\nabla\lambda(b)|\).
If \(d_b\) is only
a fixed multiple
of \(\eta\), the
point-value bound
on a difference
controls it by
\(CM\).
Otherwise apply the
scalar interior gradient
estimate to a ball
of radius \(d_b/4\)
around an edge
where this maximum
is approached.
Distances in that
smaller ball are
comparable to
\(d_b\). Its
gradient supremum
is therefore at
most the weighted
maximum divided
by \(c d_b\).
Using \eqref{px:site-scalar-equation}
in the scalar
estimate gives
\begin{equation}
 \sup_b d_b|\nabla\lambda(b)|
       \le C_L(M+\xi^2F)
                       +C_La_{\rm loc}
                           \sup_b d_b|\nabla\lambda(b)|.
 \label{px:site-gradient-absorption}
\end{equation}
Choose \(C_La_{\rm loc}\le1/2\).
Absorption, \eqref{px:site-pointwise-green}--\eqref{px:site-laplacian-green}
and the scalar interior
Hölder estimate then give,
for a bounded source
in \(\Delta(z)\),
\begin{equation}
 \begin{split}
 \|\nabla_UG'_Uf\|_{\infty,\Delta(y)}
       &\le C_L\xi_y e^{-\delta d(y,z)}\|f\|_\infty,\\
 [\nabla_UG'_Uf]_{\alpha,\Delta(y)}
       &\le C_{L,\alpha}\xi_y^{1-\alpha}
                           e^{-\delta d(y,z)}\|f\|_\infty,
                       \qquad 0<\alpha<1 .
 \end{split}
 \label{px:site-gradient-green}
\end{equation}
The Hölder comparison is made
in the assigned local frame,
or with its declared
transport to a common
base. The bounded
potential and pointwise
gradient give the
transported version.
The constant is not
asserted bounded as
\(\alpha\uparrow1\).
Cells with a bounded
number of fine sites
use the direct
difference bound;
near the action
boundary they
are bare, as
already noted.

These pointwise estimates continue
to the actual complex site
operator by a stronger
norm argument than a
pointwise matrix
coefficient estimate.
For a fixed source cell
\(z\), let
\begin{equation}
 \begin{split}
 \|\lambda\|_{\mathcal Y_z}
   &=\sup_y e^{\delta d(y,z)}
       \max\{\xi_y^{-2}\|\lambda\|_{\infty,\Delta(y)},
              \xi_y^{-1}\|\nabla_U\lambda\|_{\infty,\Delta(y)}\},\\
 \|f\|_{\mathcal F_z}
   &=\sup_y e^{\delta d(y,z)}\|f\|_{\infty,\Delta(y)} .
 \end{split}
 \label{px:site-profile-spaces}
\end{equation}
Use a rate \(\delta\) strictly
below the cell Green rate,
with its spare rate
larger than the
ball-growth rate
in \eqref{px:metric-ball-count}. Summing
the bounded-source
cell estimates
then proves
\(\|G'_U\|_{\mathcal F_z\to\mathcal Y_z}
\le C_L\), uniformly
in \(z\).

For the actual left perturbation
\(U_E=e^{i\eta E}U\), expand
the two covariant site
differences and the
two site averages.
The kinetic difference
has a coefficient
bounded by \(|E|\)
on a first difference
of \(\lambda\), and
a coefficient bounded
by
\(|\nabla_UE|+|E|^2\)
on \(\lambda\).
The mass difference
uses \eqref{px:site-transport-lipschitz}, has
size \(C_Lh\xi^{-2}\),
and reads only
the same pillar.
Hence
\begin{equation}
 \|(L_E-L_U)\lambda\|_{\mathcal F_z}
                         \le C_Lh\|\lambda\|_{\mathcal Y_z}.
 \label{px:strong-site-perturbation}
\end{equation}
No second difference of
\(\lambda\) occurs:
the difference of two
connections in a
site Laplacian is
lower order.
The first-order
input on \(E\)
is used here;
its supremum
alone would
not bound the
zeroth-order
coefficient.
The metric diameter
of a pillar is
bounded, so the
profile weight
passes through
the mass average
with a fixed
factor.

For \(h\) below the resulting
fixed threshold, the series
\begin{equation}
 G'_E=\sum_{n\ge0}
            [-G'_U(L_E-L_U)]^nG'_U
 \label{px:strong-site-neumann}
\end{equation}
converges from \(\mathcal F_z\)
to \(\mathcal Y_z\).
At finite volume it
coincides with the
inverse already
constructed in
\eqref{px:complex-site-inverse}, by
uniqueness.
It therefore proves
the complex versions
of the pointwise
potential and
gradient estimates.
The equation supplies
the complex covariant
Laplacian bound.
In the same real
local frame,
\(e^{i\eta E}e^{i\eta A}\)
has a physical
logarithm with
the first-order
bounds
\(C(a_{\rm loc}+h)\).
The scalar interior
estimate then
gives the complex
Hölder member
as well.

The argument establishes the
site estimates for bounded
cell sources. A mixed
gradient with a
divergence source
has a Calderón--Zygmund
singularity and
requires its
specified cancellation
or Hölder input.
No uniform bound
for that mixed
operator on all
bounded raw
sources has
been inferred.

\subsubsection{The site Gram inverse has its own lower-bound proof}

Invertibility of \(G'\) does not by
itself give a uniform inverse for
\(Q'G'^2Q'^\sharp\). We prove
that estimate using site
bump functions.
Let \(\Xi\) multiply a
top value by \(\xi_y\).

Choose a fixed smooth scalar
bump \(\chi_y\), supported
in the middle portion
of its pillar and equal
to one on a smaller
middle portion.
Its first two
differences are bounded
by \(C\xi_y^{-1}\)
and \(C\xi_y^{-2}\).
The arithmetic average
\(\beta_y=Q'_0\chi_y\)
has a positive
lower bound
independent of
the fine mesh:
count the grid
points in that
smaller cube.
The finitely
coarse cases
use the center
point.

Use the real local gauge
\eqref{px:site-local-coefficient-class} on a neighborhood
of the pillar and
its one-fine-cell
collar.
In this gauge
the initial open
bond bound is
\(C a_{\rm loc}\eta/\xi_y\).
The recurrence
\eqref{px:pillar-open-recurrence}, with
the same real
curvature bound,
therefore puts
the hierarchical
transports within
\(C_L(a_{\rm loc}+\epsilon)\)
of one.
Consequently the
three by three
matrix
\(m_y\psi=Q'_U(\chi_y\psi)\)
equals
\(\beta_y I+
O(C_L(a_{\rm loc}+\epsilon))\).
For the fixed
small class
ceiling it is
invertible with
uniform bound.

Define a site lift on that
pillar by
\(\chi_y m_y^{-1}\psi(y)\)
in the local gauge and
undo the gauge.
At the node use
its transformed
frame as well.
The resulting
operator \(E_{\rm sm}\)
satisfies
\begin{equation}
 Q'_UE_{\rm sm}=1,\qquad
 \|E_{\rm sm}\psi\|_2\le C_L\|\psi\|_{2,{\rm tops}},\qquad
 \|L_UE_{\rm sm}\psi\|_2
                         \le C_L\|\Xi^{-2}\psi\|_{2,{\rm tops}} .
 \label{px:site-smooth-lift}
\end{equation}
The first equality holds
pillar by pillar, since
the site averages have
disjoint supports.
For the last inequality,
the covariant Laplacian
of one bump is
bounded by
\(C_L\xi_y^{-2}|\psi(y)|\).
Expand it in the
local gauge: the
second cutoff
differences, the
connection times
first differences,
and the differentiated
connection all have
that bound by
\eqref{px:site-local-coefficient-class}.
Its mass term
is the transported
vector \(a(y)\psi(y)\)
on the same pillar.
The volume is
\(\xi_y^4\), giving
the last \(L^2\)
scale in \eqref{px:site-smooth-lift}.
One fine collar
of different bumps
has bounded overlap,
since each fine
site has a fixed
number of neighbors.
For a bare pillar,
take the single
site bump; its
Laplacian is
bounded by
\(C\eta^{-2}|\psi|\),
which is exactly
the same scale.
Thus no missing
local gauge at
the outer bare
boundary is used.

For a real top vector \(\psi\),
use \eqref{px:site-smooth-lift} with
\(\Xi^4\psi\) and
the equation
\(L_UG'_U=1\):
\begin{equation}
 \begin{split}
 \|\Xi^2\psi\|_{2,{\rm tops}}^2
   &=\langle(Q'_U)^\sharp\psi,
                         E_{\rm sm}\Xi^4\psi\rangle\\
   &=\langle G'_U(Q'_U)^\sharp\psi,
                           L_UE_{\rm sm}\Xi^4\psi\rangle\\
   &\le C_L\|G'_U(Q'_U)^\sharp\psi\|_2
                            \|\Xi^2\psi\|_{2,{\rm tops}} .
 \end{split}
 \label{px:site-gram-duality}
\end{equation}
Divide by the last node norm,
with the zero case immediate.
This proves the operator lower
bound
\begin{equation}
 Q'_UG_U'^2(Q'_U)^\sharp\ge c_L\Xi^4 .
 \label{px:site-gram-lower}
\end{equation}
The \(\Xi^4\) test in \eqref{px:site-gram-duality}
is essential. Testing only
with \(E_{\rm sm}\psi\)
would give a weaker
ratio of two norms
and would not prove
the stated weighted
lower bound at
different local
scales.

For the upper bound and
locality, set
\begin{equation}
 \mathcal F_U=G'_U(Q'_U)^\sharp\Xi^{-2},\qquad
 M_U=\Xi^{-2}Q'_UG_U'^2(Q'_U)^\sharp\Xi^{-2}
                                      =\mathcal F_U^\sharp\mathcal F_U .
 \label{px:normalized-site-gram}
\end{equation}
The site injection \((Q'_U)^\sharp\)
is an isometry between
node \(L^2\) and
its transported constant
pillar fields.
\eqref{px:site-cell-l2-decay}
therefore bounds the
normalized cell entry
of \(\mathcal F_U\)
by
\begin{equation}
          C_L(\xi_y/\xi_z)e^{-m d(y,z)} .
 \label{px:site-column-cell-bound}
\end{equation}
The same estimate holds
for the complex
\(\mathcal F_E\), with
the bounded transport
condition numbers
already proved.
The cell summability
calculation with
\(p=1\) gives a
uniform exponential
two-sided Schur
norm of these
operators, after
increasing the
fixed separation
if needed.
Their products
therefore have
the same type
of norm.
In particular
\(c_LI\le M_U\le C_LI\)
at a real pair.

The complex perturbation is
uniform as well.
The finite matrix
\(\mathcal F_E\) is
holomorphic on the
actual physical
first-order ball
constructed above.
Its operator norm
is uniformly bounded
there by \eqref{px:site-column-cell-bound}.
Apply one scalar
Cauchy circle to
the line through
\(E\). This gives
\(\|\mathcal F_E-\mathcal F_U\|
\le C_Lh\),
independent of
the number of
coordinates.
Hence
\begin{equation}
 \|M_E-M_U\|_{2\to2}\le C_Lh,\qquad
                         \|M_E^{-1}\|_{2\to2}\le 2/c_L
 \label{px:complex-site-gram}
\end{equation}
on a fixed smaller ball.
Here \(M_E=\mathcal F_E^\sharp\mathcal F_E\)
uses the analytic
bilinear transpose.
Its norm difference
is bounded by
\((\|\mathcal F_E\|+\|\mathcal F_U\|)
\|\mathcal F_E-\mathcal F_U\|\).
The transpose has
the same Hilbert
operator norm in
the real weighted
orthonormal bases.
No Hermitian
positivity of
the complex
Gram matrix is
assumed.

An exponentially weighted
Schur norm also controls
its conjugation. For
a bounded node function
\(\psi\) that is one-Lipschitz
in \(d\), its cell
entries change by
the factor
\(e^{s(\psi(y)-\psi(z))}-1\).
The inequality
\begin{equation}
 |e^{st}-1|\le |s|\,|t|e^{|s||t|}
 \label{px:nonlocal-conjugation}
\end{equation}
and a spare fraction of
the exponential rate
bound this perturbation
by \(C_L|s|\) in
operator norm.
Indeed
\(\sup_{r\ge0}r e^{-\delta r/2}<\infty\).
Choose \(s\) small
and use \eqref{px:complex-site-gram}.
The inverse of the
conjugated \(M_E\)
is uniformly bounded.
The distance-weight
argument used for
\eqref{px:site-cell-l2-decay} gives
exponential entries
of \(M_E^{-1}\).

In the physical top pairing
the weighted orthonormal
basis at \(y\) is
\(\xi_y^{-2}\) times
the color basis.
If \(B_{yz}\) are
the entries of
\(M_E^{-1}\) in
those orthonormal
coordinates, then
\(\|B_{yz}\|\le C_Le^{-m_2d(y,z)}\).
Undoing both
\(\Xi^{-2}\)
factors and
the orthonormal
bases gives the
counting kernel
\begin{equation}
 \left|
 (Q'_EG_E'^2(Q'_E)^\sharp)^{-1}(y,z)
 \right|
                   \le C_L\xi_y^{-4}e^{-m_2d(y,z)} .
 \label{px:site-gram-counting-kernel}
\end{equation}
Explicitly the unnormalized
entries of \(M_E^{-1}\)
are
\((\xi_z^2/\xi_y^2)B_{yz}\).
Multiplying on the
left and right
by \(\xi_y^{-2}\),
\(\xi_z^{-2}\)
leaves \(\xi_y^{-4}B_{yz}\).
This explains the
asymmetric fourth
power in a
counting-kernel
formula, despite
symmetry in the
physical weighted
pairing.

Finally the actual site projector
has the representation
\begin{equation}
 R_E=1-\mathcal F_EM_E^{-1}\mathcal F_E^\sharp .
 \label{px:site-projector-kernel}
\end{equation}
It agrees with its original
Gram formula by \eqref{px:normalized-site-gram}.
The smoothing part has
an exponential cell
\(L^2\) Schur norm,
and the matrix is
holomorphic on the
whole smaller
physical ball.
At a real field
it is the orthogonal
projector already
identified.
This completes the
site Gram and
projector part
of the package.
Fine vector
coercivity,
mixed derivative
source norms
and the current
constraint terms
are not consequences
of \eqref{px:site-projector-kernel}.

\subsubsection{The physical fine-to-coarse readout is a Hessian identity}

We next identify the
three quadratic forms
in the physical
calculation.
This is a finite-dimensional
identity on any
admitted chart;
it precedes a
uniform inverse
estimate.
Let \(U\) be
the reference
critical field,
\(j_U\) its action
current, \(Q\)
the derivative
of the actual
straightened
constraint, and
\(\omega\) its
multiplier.
Take a fixed
reference slice,
so that its
linear condition
is \(RD_U^\sharp a=0\).

Criticality on that slice
already gives the full
stationarity equation
\begin{equation}
                         j_U=Q^\sharp\omega .
 \label{px:physical-stationarity}
\end{equation}
To see this, first note
that \(Q\mathsf S=Q\),
\(\mathsf S a\) is
in the slice, and
\(a-\mathsf S a=D_UG'RD_U^\sharp a\).
For \(a\in\ker Q\),
criticality annihilates
\(\mathsf Sa\);
the Ward identity
annihilates its
remaining gauge
part.
Thus \(j_U\)
annihilates
\(\ker Q\).
The finite-dimensional
annihilator identity
gives \eqref{px:physical-stationarity}.
When the row
space is the
actual independent
datum space,
\(Q\) is onto
and \(\omega\)
is unique.
Equivalently it
is unique whenever
the row Gram
inverse used in
\eqref{px:explicit-current-lift} exists.

Because \(Q\widehat H=1\),
\eqref{px:physical-stationarity} implies
\(\widehat H^\sharp j_U=\omega\).
Define the actual
Lagrangian Hessian
\begin{equation}
             K_{\rm L}=\mathscr D_U-\Delta^{(2)}[\omega].
 \label{px:physical-lagrangian-hessian}
\end{equation}
The quadratic constraint
term here is the
second derivative
of the same
straightened
constraint as \(Q\),
with the same
per-length pairing.
It cannot be
replaced by
the quadratic
part of a
different row map.

For the explicitly chosen
Ward correction \eqref{px:symmetric-current-correction},
the physical second
variant is therefore
\begin{equation}
 \widehat G_1[j_U]^{-1}
    =\mathsf S^\sharp K_{\rm L}\mathsf S
                      +D_URD_U^\sharp+Q^\sharp aQ .
 \label{px:physical-full-form}
\end{equation}
Let \(W_1=Q\widehat G_1[j_U]Q^\sharp\)
and \(H_1=\widehat G_1[j_U]Q^\sharp W_1^{-1}\).
The physical version
of \eqref{px:off-current-slice-defect} gives
\(RD_U^\sharp H_1=0\).
Consequently
\(\mathsf SH_1=H_1\),
and direct multiplication
of \eqref{px:physical-full-form} gives
\begin{equation}
                   H_1^\sharp K_{\rm L}H_1=W_1^{-1}-a .
 \label{px:physical-row-readout}
\end{equation}
For example,
\(H_1^\sharp\widehat G_1^{-1}H_1=W_1^{-1}\),
the gauge contribution
is zero, and the
row mass contributes
\(a\), since \(QH_1=1\).
This proves all
signs in \eqref{px:physical-row-readout}.
It does not
prove the existence
of either inverse.

Let the actual local
critical branch be
parametrized by a
coarse variation \(B\),
with reduced datum \(b(B)\).
Its derivative \(h_B\)
at zero is taken
in the fixed
reference slice.
Differentiated
stationarity gives
the constrained
linear equations
for \(K_{\rm L}\),
with row value
\(Db(0)\dot B\).
They imply
\begin{equation}
        h_B=H_1Db(0),\qquad
 D^2\mathcal A_{\rm eff}(0)
      =Db(0)^\sharp(W_1^{-1}-a)Db(0)
                            +\langle\omega,D^2b(0)\rangle .
 \label{px:coarse-action-hessian}
\end{equation}
Here \(\mathcal A_{\rm eff}\)
is the action evaluated
on that branch.
For the first assertion,
the differentiated
equation tested against
slice vectors in
\(\ker Q\) says
\(\mathsf S^\sharp K_{\rm L}h_B\)
lies in the range
of \(Q^\sharp\).
Adding \(Q^\sharp aQh_B\)
and using \eqref{px:physical-full-form}
then solves the
row problem by
its defining inverse.
For the second,
differentiate the
value function
twice.
The second derivative
of the branch
contracts with
\(j_U=Q^\sharp\omega\);
differentiating
its constraint
replaces that
contraction by
\(\omega D^2b-
\omega D^2q[h_B,h_B]\).
The latter term
is exactly the
subtraction in
\(K_{\rm L}\).

\subsubsection{Which coarse coordinates remove the datum curvature}

The statement \(b(B)=B\)
has a specific coordinate
meaning. In the
present chart the
actual coarse group
field is
\begin{equation}
                             V(B)(c)=e^{i\xi_cB(c)}V_0(c).
 \label{px:left-coarse-chart}
\end{equation}
The fine straightened
row logarithm is
relative to the
same reference
\(V_0\).
On an interior
row the exact
node identity
then gives \(b(B)(c)=B(c)\).
It holds for
nontrivial \(V_0(c)\)
as well.
No assertion
about the absolute
logarithm of
\(V(B)(c)\)
is involved.

An outward row instead has
its prescribed factor
at the finer outer
endpoint.
In raw notation
its equation is
\begin{equation}
       e^{i\xi_cb(B)(c)}
           =e^{i\xi_cB(c)}\operatorname{Ad}_{V_0(c)}
                                      v_z(B_{\rm int}) .
 \label{px:outward-datum-factor}
\end{equation}
It can have both
linear and quadratic
datum corrections.
If only the
highest active
level is varied,
the lower-level
data \(B_{\rm int}\)
in this factor
are fixed to zero.
Then \(v_z(0)=1\)
and the factor
is constant.
More generally
the next-step
support mask
must be checked
against every
outward carrier.
Deep interior
rows give an
immediate instance;
a mask by itself
does not change
a logarithmic
coordinate.

For contrast, if an absolute
coordinate \(W_c=\log(V_0(c))/i\)
is used and the variable
is still a relative
increment \(D\), its
new absolute logarithm is
\(\log(e^{iD}e^{iW_c})/i\).
Its second derivative
in \(D\) generally
does not vanish.
\eqref{px:coarse-action-hessian}
then retains its
datum-curvature term.
This distinction is
independent of
the support mask.

Let \(\iota\) inject the actual
active next-step coarse
bond variables into
the row space.
When the literal
map in \eqref{px:left-coarse-chart}--\eqref{px:outward-datum-factor}
satisfies
\begin{equation}
 Db(0)\iota=\iota,\qquad
                     D^2b(0)[\iota\,\cdot,\iota\,\cdot]=0,
 \label{px:active-datum-flatness}
\end{equation}
its coarse action
precision on those
variables is
\(\iota^\sharp(W_1^{-1}-a)\iota\).
Its coarse current
is
\(\iota^\sharp H_1^\sharp j_U\).
The restriction
\(\iota^\sharp\)
and the physical
dual weights are
part of the
readout, even
when notation
suppresses them.

\subsubsection{The next curved fibre inserts a different multiplier}

At the next lattice,
write raw bond increments
as \(D\), so that
the group field is
\(e^{iD}V\).
The actual plain
constraint logarithm
is
\begin{equation}
 \mathfrak q_V(D)
    =\log\!\left(M(e^{iD}V)M(V)^{-1}\right)/i .
 \label{px:next-raw-constraint}
\end{equation}
Its local pivot graph
has derivative
\(C_V\), and
the pivot injection
\(h_V\) satisfies
\(D\mathfrak q_V(0)h_V=1\).
Twice differentiating
the graph equation
gives
\begin{equation}
 D^2\Phi_V(0)[u,v]
                =-h_VD^2\mathfrak q_V(0)[C_Vu,C_Vv].
 \label{px:next-graph-second-derivative}
\end{equation}
This is the actual
pivot graph, not
a free extension
chosen after the
Hessian.

Let \(\widehat K\) and
\(\widehat J\) be the
coarse effective action
Hessian and gradient
in these raw
counting coordinates.
The chain rule on
that graph is
\begin{equation}
 \begin{split}
 D^2(\mathcal A_{\rm eff}\circ\Phi_V)(0)[u,v]
   ={}&\langle C_Vu,\widehat K C_Vv\rangle\\
     &-\langle h_V^\sharp\widehat J,
                   D^2\mathfrak q_V(0)[C_Vu,C_Vv]\rangle .
 \end{split}
 \label{px:next-constrained-hessian}
\end{equation}
Thus the next constraint
multiplier is
\(h_V^\sharp\widehat J\).
It is not
\(\widehat J\)
on the new
row space unless
the notation
already includes
this \(h_V^\sharp\).
Define \(E_V^{(2)}[\widehat J]\)
by this last
contraction.
With the actual
soft gauge Hessian
\(\widehat K_{\rm GF}\),
the next free
precision is
\begin{equation}
             C_V^\sharp
                \bigl(\widehat K-E_V^{(2)}[\widehat J]
                                      +\widehat K_{\rm GF}\bigr)C_V .
 \label{px:complete-next-precision}
\end{equation}
A gauge term relative
to \(V\) and the
absolute contour
gauge have equal
values on an
axial physical
reference.
Off that locus
their Hessians
can differ, as
shown in the
domain section.
\eqref{px:complete-next-precision}
uses the actual
gauge definition
of the model
being considered.

For a uniform next-lattice
spacing \(\xi\),
the raw and physical
conversions are
\begin{equation}
       \widehat K=\xi^2 K_{\rm phys,count},\qquad
                              \widehat J=\xi^3J_{\rm phys}.
 \label{px:coarse-raw-duality}
\end{equation}
Indeed
\(\xi^4\sum B(KB)=
\xi^2\sum D(KD)\)
when \(D=\xi B\),
and
\(\xi^4\sum J_{\rm phys}B=
\sum \xi^3J_{\rm phys}D\).
At nonuniform row
lengths the corresponding
diagonal length
matrices must be
kept on both
sides; their
commutation is
not assumed.
The coupling
coefficients in
the original
action remain
inside \(K\)
and \(J\).

Combining \eqref{px:physical-row-readout}, \eqref{px:active-datum-flatness}
and \eqref{px:next-constrained-hessian} gives the
physical diagram:
fine Lagrangian
Hessian, coarse
value Hessian,
then next-fibre
Hessian.
Each arrow has
its own row map,
its own dual
current, and
its own gauge
term.
The physical
identity proves
none of those
three objects
interchangeable.

For completeness the chosen
off-current extension
has an exact
additional readout.
Put \(h=H_1[J]\),
\(r=RD_U^\sharp h\),
\(\delta J=J-j_U\),
and \(K[J]=\mathscr D_U-
\Delta^{(2)}[\widehat H[J]^\sharp J]\).
Using \eqref{px:physical-invariant-hessian} in
\eqref{px:second-current-variant} gives
\begin{equation}
 \begin{split}
 (Q\widehat G_1[J]Q^\sharp)^{-1}-a
   ={}&h^\sharp\mathsf S^\sharp K[J]\mathsf Sh
         -h^\sharp\mathcal B_U[\delta J]h+r^\sharp r .
 \end{split}
 \label{px:off-current-row-readout}
\end{equation}
All pairings are
bilinear on the
complex domain.
On a real pair
the last term
is nonnegative;
it vanishes on
the physical
graph.
The middle term
is a current
defect and does
not disappear
from a Ward
property alone.
For a separately
named normal form,
the complete
projected difference
between \eqref{px:complete-next-precision}
and that form
must be identified,
including its
horizontal multiplier
and gauge terms.

\subsubsection{A flat Feynman comparison on all the layers}

There is a useful vector
estimate which does not
involve the Landau
projection or an
independent current.
It uses the actual
flat straightened
averages.
We prove it to
specify exactly
what is available
before those two
additional operations.

First take the full periodic
fine lattice as the
ambient graph.
Let \(a_\mu(x)\)
be a fine vector
one-form, with
physical \(L^2\)
measure \(\eta^4\).
For a positive
level-\(j\) row
of length
\(\xi_j=n\eta\),
\(n=L^j\), its
flat per-length
average is
\begin{equation}
 (Q_j^0a)_\mu(y)
    =n^{-5}\sum_{u\in\mathcal B_n}
                      \sum_{t=0}^{n-1}a_\mu(y+\eta u+\eta t e_\mu).
 \label{px:flat-field-average}
\end{equation}
Here \(\mathcal B_n\)
is the centered
set of \(n^4\)
fine block sites.
The formula is
the flat straightened
map, not the
plain path map.
At \(j=0\)
it is the
identity.
The coefficient
of each fine
bond is nonnegative,
the coefficients
sum to one,
and each is at
most \(n^{-4}\).
The last fact
follows because
at most \(n\)
pairs \((u_\mu,t)\)
give a fixed
longitudinal
position, while
the three other
coordinates are
fixed.

Its support lies in the
union of the two
endpoint blocks.
Consequently
\begin{equation}
 |Q_j^0a(c)|\le
               C\xi_j^{-2}\|a_\mu\|_{2,\mathcal D_c} .
 \label{px:flat-row-ltwo}
\end{equation}
For example Cauchy
gives
\(\sum_b w_b^2\le n^{-4}\)
from \(\sum_bw_b=1\)
and \(w_b\le n^{-4}\);
converting counting
\(L^2\) to physical
\(L^2\) contributes
\(\eta^{-2}\).
Thus the resulting
coefficient is
\(\eta^{-2}n^{-2}=\xi_j^{-2}\).
This bound is
specific to this
flat row formula.

For a fine bond \(b\),
put \(\xi_b\) equal
to the larger of
its endpoint pillar
scales and define
\begin{equation}
 \|a\|_{\rm mass}^2
          =\sum_b\eta^4\xi_b^{-2}|a(b)|^2,\qquad
 \mathcal E_{\rm F}^0(a)
     =\|\partial a\|_2^2+
                  \sum_{c\in\Gamma}c_c\xi_c^2|Q^0a(c)|^2 ,
 \quad 0<c_*\le c_c\le c^* .
 \label{px:flat-feynman-energy}
\end{equation}
The component gradient
contains all fine
coordinate differences
on the ambient
torus. The second
term is exactly
the physical row
pairing with
\(a(c)=c_c\xi_c^{-2}\).
In particular it
has neither an
extra fine volume
factor nor an
unweighted counting
transpose.

For each nonbare top \(y\)
and each direction \(\mu\),
one of the adjacent
level-\(j\) rows
has \(y\) as
an endpoint and
is a determining
row. Indeed the
excluded higher
region is a
union of aligned
blocks of side
at least \(L\xi_j\).
Along that
coordinate line,
the excluded
block containing
an exterior
point has at
least \(L\ge3\)
level-\(j\) sites.
Both neighboring
sites cannot
belong to the
higher region.
Choose one that
does not.
If the other
endpoint is
outside \(\Omega_j\),
this is precisely
an outward row,
which is retained.
At the highest
level there is
no higher
region to
exclude.

Let \(m_{y,\mu}\) be the
arithmetic mean of
\(a_\mu\) on the
pillar enlarged by
one fine cell.
A rectangular
carrier containing
that enlarged
pillar and the
chosen row's
two endpoint
blocks has side
length at most
\(C\xi_j\).
For \(j\ge1\)
it lies inside
the ambient
fine graph by
the nesting
gap.
The ordinary
box Poincare
inequality,
proved by the
coordinate-path
count used for
\eqref{px:uniform-pillar-transport}, gives
\begin{equation}
 |Q_j^0a(c(y,\mu))-m_{y,\mu}|
           \le C\xi_j^{-1}\|\partial a_\mu\|_{2,\mathcal D_{y,\mu}} .
 \label{px:flat-row-mean-comparison}
\end{equation}
To verify this without
a pointwise trace
estimate, subtract
the carrier's
arithmetic mean.
Both functionals
annihilate that
constant in
their difference
and have \(L^2\)
norm at most
\(C\xi_j^{-2}\).
Poincare then
supplies one
factor \(\xi_j\).

A second use of Poincare
on the enlarged
pillar gives
\begin{equation}
 \xi_j^{-2}\|a\|_{2,\mathcal B_y^+}^2
       \le C\sum_\mu\|\partial a_\mu\|_{2,\mathcal D_{y,\mu}}^2
                    +C\xi_j^2\sum_\mu|Q_j^0a(c(y,\mu))|^2 .
 \label{px:flat-field-local-poincare}
\end{equation}
Every fine bond with
a nonbare endpoint
can be assigned
to a pillar of
its larger endpoint
scale and is
included in
that pillar's
enlargement.
If both endpoints
are bare, its
level-zero row
is the identity
and controls it
directly.

The overlap in the gradient
sum is uniform.
For a fixed level,
all carrier centers
are on its grid
and their diameters
are \(C\xi_j\),
so a fine point
belongs to a
fixed number
of them.
A carrier of
an outward
level-\(j\) row
is inside
\(\Omega_{j-1}\),
and its endpoint
blocks are
outside
\(\Omega_{j+1}\).
Its one-fine-cell
collar can enter
only an adjacent
level.
Thus a fixed
fine point
is counted at
only a fixed
number of levels,
not at every
level of the
sequence.
Each chosen row
has at most its
two endpoints
as selecting
tops.
Summing \eqref{px:flat-field-local-poincare}
therefore proves
\begin{equation}
                         \|a\|_{\rm mass}^2
                                    \le C_L\mathcal E_{\rm F}^0(a).
 \label{px:flat-feynman-gap}
\end{equation}
The constants are
independent of
depth, fine mesh,
volume and
the shapes of
the separated
regions.
Periodic carriers
which wrap are
lifted to their
rectangles; the
compatible torus
has at least
two top blocks,
so the multiplicity
of that lift
is bounded.

There is also the upper
row estimate
\begin{equation}
              \sum_{c\in\Gamma}\xi_c^2|Q^0a(c)|^2
                                      \le C_L\|a\|_{\rm mass}^2 .
 \label{px:flat-field-row-upper}
\end{equation}
Use \eqref{px:flat-row-ltwo} on
each row carrier.
All fine endpoint
scales there are
comparable to
\(\xi_c\), within
the fixed adjacent
scale factor.
The preceding
overlap count
then applies.
\eqref{px:flat-feynman-gap}
and \eqref{px:flat-field-row-upper}
are estimates
for the
flat component
Feynman form.

For a proper active region,
extend the sequence
by bare cells on
the ambient torus
and compress this
full positive
form by zero
extension of its
active bond
coordinates.
Its lower bound
survives exactly:
apply \eqref{px:flat-feynman-gap}
to the extended
vector.
This is the
Dirichlet compression
of the component
Laplacian.
An inner-plaquette
Maxwell Hessian
has a different
boundary form.
Nor has the
nonlocal subtraction
\(D_U(1-R)D_U^\sharp\)
been included.
Those comparisons
are separate
in the vector
precision problem.

\subsubsection{What small curvature does to the actual Wilson Hessian}

The Wilson part admits a
relative form estimate
without assuming a
matrix norm bound.
We derive it directly
from the plaquette
product.
Use
\(\langle X,Y\rangle_{\rm col}
=\tfrac12\operatorname{Tr}(XY)\)
on Hermitian traceless
matrices.
For a plaquette
with oriented bonds
\(b_1,\ldots,b_4\),
put \(P=U(b_1)\cdots U(b_4)\)
and \(T_{r-1}=U(b_1)\cdots U(b_{r-1})\).
A left variation
\(e^{i\eta v(b)}U(b)\)
inserts
\(X_r=\eta\operatorname{Ad}(T_{r-1})v(b_r)\).
On reversed bonds
the variation is
\(v(\bar b)=-\operatorname{Ad}(U(b)^{-1})v(b)\),
as follows by
inverting the
group field.

The varied plaquette is
\(\prod_{r=1}^4e^{iX_r}P\).
At \(P=1\), the
trace of the
mixed second
derivative is
the trace of
the product of
the two sums
of insertions.
Cyclicity of
the trace
cancels their
commutators.
For general real
\(P\), subtract
that calculation
and use the
factor \(P-1\)
in every
remaining term.
Thus the
plaquette action
\(1-\tfrac12\operatorname{Tr}P\)
satisfies
\begin{equation}
 \begin{split}
 D^2\mathcal A_p(U)[v,w]
     &=\left\langle\sum_rX_r,\sum_rY_r\right\rangle_{\rm col}
                                      +\mathcal R_p(v,w),\\
 |\mathcal R_p(v,w)|
     &\le C\|P-1\|
           \left(\sum_r|X_r|^2\right)^{1/2}
           \left(\sum_r|Y_r|^2\right)^{1/2}.
 \end{split}
 \label{px:plaquette-hessian-exact}
\end{equation}
There are only sixteen
mixed insertion terms.
Unitarity of the
prefixes makes
\(\sum|X_r|^2=\eta^2\sum|v(b_r)|^2\).
The same polynomial
trace expression,
with no real-part
operation, is its
holomorphic continuation.

Define the plaquette
linearized curl by
\begin{equation}
 (\mathfrak d_Uv)(p)
          =\eta^{-1}\sum_{r=1}^4
                              \operatorname{Ad}(T_{r-1})v(b_r).
 \label{px:plaquette-linearized-curl}
\end{equation}
Then \(\sum_rX_r=\eta^2\mathfrak d_Uv(p)\).
If
\(\|P-1\|\le\epsilon\eta^2\xi_p^{-2}\),
with \(\xi_p\)
the largest
corner scale,
summing \eqref{px:plaquette-hessian-exact}
over the actual
plaquettes gives
\begin{equation}
 \left|\langle v,\mathscr D_Uw\rangle
                       -\langle\mathfrak d_Uv,\mathfrak d_Uw\rangle_2\right|
              \le C_L\epsilon\|v\|_{\rm mass}\|w\|_{\rm mass}.
 \label{px:wilson-form-remainder}
\end{equation}
Every bond is in
at most six
unoriented
plaquettes.
Each error has
the factor
\(\eta^2\epsilon\xi_p^{-2}\)
from \(P-1\)
and another
\(\eta^2\)
from the two
insertions.
Their product
is the physical
\(\eta^4\) measure
times the
natural inverse
scale squared.
There is no
loss \(\epsilon\eta^{-2}\)
in this form
estimate.

This curl is close to the
ordinary forward
covariant curl
\(\mathrm{d}_Uv=\nabla_{U,\mu}v_\nu-
\nabla_{U,\nu}v_\mu\).
For example on a
\(\mu\nu\) plaquette
the difference
of their numerator
sums is
\begin{equation}
 (1-\operatorname{Ad}P)
           \bigl(\operatorname{Ad}U_\nu(x)v_\mu(x+\eta e_\nu)
                                                    +v_\nu(x)\bigr).
 \label{px:curl-convention-difference}
\end{equation}
Direct multiplication of
the four prefixes
proves this
identity.
It gives
\(\|\mathfrak d_Uv-\mathrm{d}_Uv\|_2
\le C_L\epsilon\|v\|_{\rm mass}\),
since \(\eta\le\xi_p\).
In particular
\begin{equation}
 \langle v,\mathscr D_Uv\rangle
     \ge(1-C_L\epsilon)\|\mathrm{d}_Uv\|_2^2
                                     -C_L\epsilon\|v\|_{\rm mass}^2
 \label{px:wilson-lower-form}
\end{equation}
for a real field,
after decreasing
the fixed
curvature ceiling.
Positive bounded
plaquette weights
give the same
statement with
their lower and
upper bounds.
Signed layer
profiles or
additional current
subtractions
must be estimated
separately.

On the full periodic graph,
the covariant Hodge
calculation is also
explicit:
\begin{equation}
 \left|
   \|\mathrm{d}_Uv\|_2^2+\|D_U^\sharp v\|_2^2
                              -\sum_\mu\|\nabla_{U,\mu}v\|_2^2
 \right|
                          \le C_L\epsilon\|v\|_{\rm mass}^2 .
 \label{px:covariant-hodge-error}
\end{equation}
Expand both squares.
The diagonal terms
agree because each
covariant shift is
unitary.
For different
directions the
uncancelled terms
are commutators
of the two shifts,
or of a shift
and an inverse
shift.
Their coefficients
are conjugates
of \(\operatorname{Ad}P-1\).
The two difference
denominators
contribute
\(\eta^{-2}\);
the plaquette
bound leaves
\(C\epsilon\xi^{-2}\).
There are a
fixed number
of neighboring
terms at each
bond.
Adjacent scales
are comparable,
so Cauchy proves
\eqref{px:covariant-hodge-error}.
On a proper
inner graph
the summation
by parts has
additional
boundary terms.

The full Feynman gauge
uses \(D_UD_U^\sharp\).
The actual fine
Landau form
instead uses
\(D_URD_U^\sharp\).
Its difference
is the nonnegative
real smoothing
form
\begin{equation}
          D_U(1-R)D_U^\sharp
             =(D_U\mathcal F_U)M_U^{-1}
                                      (D_U\mathcal F_U)^\sharp .
 \label{px:landau-feynman-difference}
\end{equation}
A lower bound for
the Feynman form
does not survive
this subtraction
by monotonicity.
This is one
reason the
vector coercivity
step remains
distinct from
the site gap
and \eqref{px:flat-feynman-gap}.

For the admitted complex
perturbation \(E\),
the actual plaquette
Taylor formula gives
\(\|P(U_E)-P(U)\|\le C_Lh\eta^2\xi_p^{-2}\).
Indeed the first
order boundary
sum of \(E\)
is its covariant
curl; the remaining
terms contain
either two
increments or
one increment
times a real
connection on
a local gauge
patch.
The first-order
core bounds
control each
by that scale.
Also
\(\|\mathrm{d}_{U_E}v-\mathrm{d}_Uv\|_2
\le C_Lh\|v\|_{\rm mass}\).
Subtracting \eqref{px:plaquette-hessian-exact}
therefore proves
\begin{equation}
 |\langle v,(\mathscr D_{U_E}-\mathscr D_U)w\rangle|
   \le C_Lh
       (\|\mathrm{d}_Uv\|_2+\|v\|_{\rm mass})
       (\|\mathrm{d}_Uw\|_2+\|w\|_{\rm mass}) .
 \label{px:complex-wilson-relative-form}
\end{equation}
This uses a physical
first-order
perturbation
radius, not
an unscaled
matrix norm
for the fine
precision.
It is a genuine
part of the
complex comparison.
The corresponding
actual field-row,
projected-current
and constraint
form comparisons
remain to be
joined to it.

\subsubsection{A mixed derivative needs cancellation}

One cannot replace those
comparisons by a
universal supremum
bound for every
order-zero
projection.
Already at a
flat field the
operator
\(\partial_1\partial_2(-\Delta+1)^{-1}\)
is not bounded
on \(L^\infty\).
Its kernel near
zero is
\(c\,x_1x_2|x|^{-6}\)
plus a locally
integrable term
in four
dimensions.
Let \(f_\eta\)
equal the sign
of this leading
kernel on the
annulus
\(C\eta<|x|<1/4\),
with a bounded
smooth cutoff
at its edges.
Then
\begin{equation}
 \|f_\eta\|_\infty\le1,\qquad
 \left|(\partial_1\partial_2(-\Delta+1)^{-1}f_\eta)(0)\right|
                                      \ge c\log(\eta^{-1})-C .
 \label{px:mixed-derivative-counterexample}
\end{equation}
The radial integral
is \(\int_{C\eta}^{1/4}dr/r\);
the angular integral
of \(|x_1x_2|/|x|^2\)
is positive.
The integrable
remainder is
uniformly bounded.
Discrete kernels
have the same
annular estimate
away from a
fixed number
of fine cells,
or the continuum
example follows
as a limit
of any proposed
uniform lattice
supremum estimate.

The required replacement has
a precise cancellation
form.
If a kernel \(K(x,z)\)
has size
\(C|x-z|^{-4}\)
on a cell of
scale \(\xi\),
and a uniformly
bounded constant-source
response, write
its action as
\begin{equation}
 \int K(x,z)(f(z)-f(x))\,dz
                  +f(x)\int K(x,z)\,dz .
 \label{px:holder-cancellation}
\end{equation}
For a transported
Hölder test,
the first term
is bounded by
\begin{equation}
        C[f]_\alpha
                 \int_0^{C\xi}r^{\alpha-1}\,dr
                  \le (C_\alpha\xi^\alpha)[f]_\alpha .
 \label{px:holder-singular-integral}
\end{equation}
The same shell
calculation applies
to fine sums,
with the first
shell at \(\eta\);
the bound remains
uniform as its lower
endpoint decreases.
The second term
is the stated
constant-source
response.
Thus an actual
mixed-kernel
estimate plus
its cancellation
clause gives
the required
Hölder-to-supremum
bound.
Size alone
does not.

For the fine gauge projection
\(\mathsf S=1-D_UG'RD_U^\sharp\),
this distinction is
material.
Its differentiated
site inverse
has order zero
at the fine
diagonal.
Energy estimates
use its \(L^2\)
form bounds;
point estimates
must either
keep it composed
with a smoothing
field Green
operator or
use actual
Hölder data
and \eqref{px:holder-cancellation}.
No unqualified
supremum Schur
norm of
\(\mathsf S\)
on arbitrary
bounded fine
fields is
being asserted.

\subsubsection{Compact dual bumps for the actual flat field rows}

We can also prove a uniform
right inverse for the
flat layered field
rows. This is a
different statement
from the site
lift \eqref{px:site-smooth-lift}.
The proof uses
two longitudinal
moments and a
strictly triangular
correction between
adjacent levels.

Fix a top \(y\) of
level \(j\ge1\),
put \(n=L^j=2m+1\),
and work on its
fine block
\(\eta\mathcal B_n+y\).
Choose a field
with just its
\(\mu\)-component
nonzero there:
\begin{equation}
 a_\mu(y+\eta u)
       =\psi\,\phi(u_\mu)
                         \prod_{\nu\ne\mu}\chi(u_\nu).
 \label{px:row-dual-bump}
\end{equation}
It is zero outside
the block and
on bonds whose
other endpoint
is outside.
For sufficiently
large \(n\),
take \(\chi(u)\)
from a fixed
even smooth
profile of
\(u/n\),
supported in
\((-1/4,1/4)\)
and positive
on a smaller
interval.
Write
\begin{equation}
 \beta=n^{-1}\sum_u\chi(u),\qquad
       \sigma=n^{-1}\sum_u(u/n)^2\chi(u).
 \label{px:row-profile-moments}
\end{equation}
Both have positive
lower and finite
upper bounds
independent of
\(n\).

For a fine longitudinal
position \(u\) in
this block, the
right coarse row
in \eqref{px:flat-field-average} counts
that bond
\(u+m+1\) times.
The left row
counts it
\(m-u\) times.
No other coarse
row sees it.
After the transverse
sums, their
two response
functionals are
\begin{equation}
 \ell_+(\phi)=\beta^3n^{-2}\sum_u(u+m+1)\phi(u),
 \qquad
 \ell_-(\phi)=\beta^3n^{-2}\sum_u(m-u)\phi(u).
 \label{px:two-row-moments}
\end{equation}
Define
\begin{equation}
 \begin{split}
 \phi_+(u)&=\beta^{-4}\chi(u)
           +\frac{m}{n}\frac{\beta^{-3}}{\sigma}
                                             \frac{u}{n}\chi(u),\\
 \phi_-(u)&=\beta^{-4}\chi(u)
           -\frac{m+1}{n}\frac{\beta^{-3}}{\sigma}
                                             \frac{u}{n}\chi(u).
 \end{split}
 \label{px:explicit-dual-profiles}
\end{equation}
Evenness gives
\(\sum u\chi(u)=0\).
Substitution into
\eqref{px:two-row-moments} now gives
\begin{equation}
 (\ell_+(\phi_+),\ell_-(\phi_+))=(1,0),\qquad
       (\ell_+(\phi_-),\ell_-(\phi_-))=(0,1).
 \label{px:dual-profile-exactness}
\end{equation}
For example the sum of
the two functionals
is \(\beta^3n^{-1}\sum\phi\).
For either profile
this is one;
the odd term
in \(\phi_+\)
sets the left
functional to
zero. The
right one is
then one.
The calculation
for \(\phi_-\)
is the same
with the other
moment.
The signed
profiles are
uniformly bounded
and their first
differences are
\(O(n^{-1})\).

The finitely many odd
values below this
smooth-profile range
need no limiting
argument.
Put the three
transverse profiles
at their center
point and use
longitudinal bonds
with positions
\(-1,0\).
Both endpoints
of these bonds
are in the
block for
\(n\ge3\).
Their two response
columns form
\begin{equation}
             n^{-5}
                    \begin{pmatrix}m&m+1\\m+1&m\end{pmatrix}.
 \label{px:finite-row-profile-matrix}
\end{equation}
Its determinant is
\(-n^{-9}\).
Its inverse
therefore gives
the required
two exact dual
profiles, with
a bound depending
only on the
fixed cutoff
between this
finite set
and the smooth
range.
That is a
universal constant,
not a depth
dependence.

Assign a determining row \(c\)
to one of its
level-\(j\) top
endpoints \(y\).
An outward row
has its interior
endpoint as
that top.
Use \(\phi_+\)
or \(\phi_-\)
according to
the chosen
endpoint.
Its seed \(E_c^0\psi\)
is supported on
fine bonds
strictly inside
that pillar and
satisfies, on
the full regular
level-\(j\) grid,
\begin{equation}
                       Q_j^0E_c^0\psi=\delta_c\psi .
 \label{px:one-row-seed}
\end{equation}
The transverse support
lies in a single
coarse transverse
block, and
\eqref{px:dual-profile-exactness} removes
the only other
longitudinal
row.
For a bare
row take its
single fine
bond value.
It has both
endpoints in
the bare layer.

Let \(E^0\psi\) be the
sum of these
seeds at every
level.
There are at
most eight
assigned seeds
at a top.
Their supports
belong to
distinct pillars,
and the collars
in their gradient
norm have bounded
overlap.
The amplitude
and first-difference
bounds just proved
give
\begin{equation}
 \|E^0\psi\|_{\rm mass}^2+
                \|\partial E^0\psi\|_2^2
                     \le C\sum_{c\in\Gamma}\xi_c^2|\psi(c)|^2 .
 \label{px:seed-energy}
\end{equation}
The factor on the
right is forced
by the dimensions:
the supporting
volume is
\(\xi_c^4\),
and a first
difference or
mass factor
is \(\xi_c^{-1}\).
For a bare
seed the same
calculation is
\(\eta^4\eta^{-2}|\psi|^2\).

\subsubsection{A contraction estimate between adjacent levels}

A seed of level \(j\)
has zero value
under every lower
determining row.
Indeed its fine
support is in
\(\Omega_j\), whereas
a lower row's
entire carrier
is outside the
next region of
that lower level.
For \(i<j\),
that next region
contains \(\Omega_j\).
It also has
zero value
under every
row of level
\(i\ge j+2\):
such a carrier
is inside
\(\Omega_{i-1}
\subset\Omega_{j+1}\),
whereas the
seed pillar
is outside
\(\Omega_{j+1}\).
The same argument
uses the two
bare endpoints
for \(j=0\).
Thus only the
next level
can receive
an unwanted
value.

The full flat averages compose
exactly:
\begin{equation}
                  Q_{j+1}^{0,{\rm full}}
                         =Q_L^0Q_j^{0,{\rm full}} .
 \label{px:flat-average-composition}
\end{equation}
In the defining sums,
the transverse
digits are
\(nr+u\), where
\(r\) is in
the centered
\(L\)-range
and \(u\)
in the centered
\(n\)-range.
They tile the
centered \(Ln\)-range.
The longitudinal
digits are
\(nt+s\), with
\(0\le t<L\),
\(0\le s<n\).
They tile
\(0,\ldots,Ln-1\).
The product
normalization is
\(L^{-5}n^{-5}\).
This proves
\eqref{px:flat-average-composition}, including
the per-length
normalization.

Let the row Hilbert norm be
\(\|\psi\|_{\rm row}^2
=\sum_c\xi_c^2|\psi(c)|^2\).
A regular \(L\)-step
average has row
coefficient sum
one and column
coefficient sum
\(L^{-4}\).
The column identity
can be checked
directly: the
two longitudinal
counts add to
\(L\), while
the three transverse
coordinates select
one block.
The ordinary
counting \(\ell^2\)
operator norm
is therefore
at most \(L^{-2}\)
by Schur.
The row norm
at the next
level has one
extra factor
\(L\) in its
square root.
Hence the
map between
these row
norms has
norm at
most \(L^{-1}\).

Writing \(T\) for the
next-level leakage,
\eqref{px:one-row-seed}--\eqref{px:flat-average-composition}
give
\begin{equation}
              Q^0E^0=1+T,\qquad
                  T_{ij}=0\ (i\ne j+1),\qquad
                         \|T\|_{{\rm row}\to{\rm row}}\le L^{-1}.
 \label{px:row-triangular-leakage}
\end{equation}
Restriction to actual
determining rows
and injection
from their
subset of the
regular grid
cannot increase
the norm.
Different level
blocks of \(T\)
have different
input and output
levels, so the
whole direct-sum
norm is the
supremum of
those block
norms.
No sum over
the number
of levels is
introduced.

We have proved the explicit
flat row lift
\begin{equation}
 E_{\rm row}^0=E^0(1+T)^{-1},\qquad
 Q^0E_{\rm row}^0=1,\qquad
 \|E_{\rm row}^0\psi\|_{\rm mass}^2+
       \|\partial E_{\rm row}^0\psi\|_2^2
       \le \frac{C}{(1-L^{-1})^2}\|\psi\|_{\rm row}^2 .
 \label{px:flat-layered-row-lift}
\end{equation}
At finite depth the
inverse is the
finite triangular
series; its norm
is bounded by
the displayed
geometric series.
The same series
converges in
the row Hilbert
space for an
infinite compatible
sequence.
The assertion
is a right
inverse with
an energy bound,
not an estimate
of a nonlinear
background map.

At a real flat connection
the same construction
uses a flat unitary
frame on each
pillar.
The response at
the other coarse
endpoint differs
by its fixed
unitary transport.
The two scalar
moments still
annihilate the
unwanted row;
the desired row
is normalized
in its own
color frame.
Schur uses
the same positive
scalar coefficient
majorants, so
its constants
are unchanged.
All profile
supports lie
inside their
pillars.
This handles
periodic flat
holonomies without
choosing a
single globally
small logarithm.
Under this gauge transfer the
gradient in the energy bound
is \(\nabla_U\), or equivalently
the ordinary component
gradient in its flat parallel
frame. An ordinary gradient
in an arbitrary original
gauge is not bounded by
that statement.
No analogous
bare \(L^2\)
field-row estimate
at a general
nonflat field
has been assumed.

\subsubsection{What a field-row Gram estimate actually needs}

The last proof gives the
flat instance of a
general useful
duality reduction.
Let \(K\) be a
real symmetric
positive fine
precision and
\(G=K^{-1}\).
Let \(\Xi\)
multiply each
field row by
its length.
Suppose the
actual row map
has an extension
\(E_{\rm row}\)
such that
\begin{equation}
 QE_{\rm row}=1,\qquad
 \langle E_{\rm row}\psi,K E_{\rm row}\psi\rangle
                            \le C\|\Xi^{-1}\psi\|_{2,{\rm rows}}^2 .
 \label{px:field-row-extension-premise}
\end{equation}
The node pairing
here has volume
\(\xi_c^4\);
its last norm
is just
\(\sum_c\xi_c^2|\psi(c)|^2\).
Testing with
\(E_{\rm row}\Xi^2\psi\)
gives
\begin{equation}
 \|\Xi\psi\|_{2,{\rm rows}}^2
   =\langle Q^\sharp\psi,E_{\rm row}\Xi^2\psi\rangle
   \le C^{1/2}\langle\psi,QGQ^\sharp\psi\rangle^{1/2}
                                       \|\Xi\psi\|_{2,{\rm rows}} .
 \label{px:field-row-gram-lower}
\end{equation}
Thus
\(QGQ^\sharp\ge C^{-1}\Xi^2\).
No pointwise
delta-function
test in the
fine field
was used.

An upper bound needs the
separate actual trace
estimate
\begin{equation}
               \|\Xi^{-1}Qa\|_{2,{\rm rows}}^2
                                      \le C\langle a,Ka\rangle .
 \label{px:field-row-trace-premise}
\end{equation}
Indeed apply it
to \(a=G^{1/2}f\).
It proves
\(\|\Xi^{-1}QG^{1/2}\|\le C^{1/2}\)
and therefore
\(QGQ^\sharp\le C\Xi^2\).
For the flat
Feynman form,
\eqref{px:flat-field-row-upper}, \eqref{px:flat-feynman-gap}
and \eqref{px:flat-layered-row-lift}
prove both
premises.
For the flat
Landau form,
\eqref{px:flat-layered-row-lift} proves
the upper energy
cost in \eqref{px:field-row-extension-premise},
since \(0\le R\le1\);
the needed fine
coercivity is
still a distinct
step.

Normalize the row covariance
as
\(N=\Xi^{-1}QGQ^\sharp\Xi^{-1}\).
If its exponential
cell Schur bound
has also been
proved, the
conjugation argument in \eqref{px:nonlocal-conjugation} applies
verbatim to \(N\).
It gives exponential
entries of \(N^{-1}\).
In physical weighted
orthonormal row
coordinates,
undoing normalization
gives the counting
kernel estimate
\begin{equation}
                   |(QGQ^\sharp)^{-1}(c,d)|
                \le C\,\xi_c^{-3}\xi_d\,e^{-m_3d(c,d)} .
 \label{px:field-row-inverse-counting}
\end{equation}
To check the powers,
an orthonormal
entry of \(N^{-1}\)
becomes
\((\xi_d^2/\xi_c^2)\)
times that entry
on unnormalized
values.
The two exterior
\(\Xi^{-1}\)
factors leave
\(\xi_c^{-3}\xi_d\).
This differs
from the
\(\xi_y^{-4}\)
in \eqref{px:site-gram-counting-kernel}:
the site Gram
contains \(G'^2\),
whereas this
row covariance
contains \(G\).

Finally a nonflat row
extension is reduced
to a definite
error estimate.
Starting from a
specified flat
or locally flat
extension \(E_0\),
set
\(D_Q=QE_0-1\).
If the actual
primitive comparison
proves
\begin{equation}
          \|D_Q\|_{{\rm row}\to{\rm row}}\le\vartheta<1,
 \label{px:nonflat-row-error}
\end{equation}
then \(E_{\rm row}=E_0(1+D_Q)^{-1}\)
is an actual
right inverse,
with its energy
cost multiplied
by
\((1-\vartheta)^{-2}\),
provided that
cost for \(E_0\)
has been proved
in the actual
precision.
The remaining
task is \eqref{px:nonflat-row-error}
for the full
actual row map
on the admitted
domain, together
with \eqref{px:field-row-trace-premise};
it is not
an assumed bound
for the desired
Green kernel.

\subsubsection{The actual straightened row is a finite analytic recursion}

We now obtain the needed
row coefficient bounds
from the actual
formula, rather
than from a
proposed inverse.
Use raw one-step
link logarithms
on a lattice
whose spacing
has been set
to one.
Let \(T_{yx}(W)\)
be the internal
mean transport
of the block
chart, with
its specified
coordinate-order
path multiset.
For a reference
\(W\) and a
raw increment \(a\),
put
\begin{equation}
 \begin{split}
 u_W(a)(y)
   &=\exp\!\left[L^{-4}\sum_{x\in B_L(y)}
       \log\!\left(T_{yx}(e^{ia}W)T_{yx}(W)^{-1}\right)\right],\\
 F_W(a)(c)
   &=\frac1i\log\!\left[
        u_W(a)(c_-)^{-1}M(e^{ia}W)(c)
                    u_W(a)(c_+)M(W)(c)^{-1}\right].
 \end{split}
 \label{px:actual-one-step-straightening}
\end{equation}
All logarithms are
the regular branch.
The first line
is an exponential
of averaged
relative logarithms;
it is not
another unspecified
mean.
The internal
transports themselves
use the actual
Federbush mean
defined in the
block chart.
This distinction
fixes both
levels of
averaging.

For \(\|\log W/i\|_\infty\)
and \(\|a\|_\infty\)
below \(c/L\),
these are
holomorphic
finite functions.
Every path has
at most \(CL\)
bonds.
Products of
their exponentials
remain in a
fixed logarithmic
neighborhood;
the mean contraction
proved in the
block chart
is uniform
in the number
of paths.
After a fixed
decrease of \(c\),
the expression
inside the
second logarithm
remains in that
same neighborhood.
Thus \(\|F_W(a)\|_\infty\le C\)
on this joint
raw ball, with
universal \(C\).

Differentiating at \(W=1,a=0\)
cancels the two
internal path
sums from the
plain average
against the
two factors \(u\).
The remaining
sum is exactly
the straight
segment average
\(S_L\), with
raw row sum
\(L\) and
raw column
sum \(L^{-3}\).
It is essential
to use both
factors in
\eqref{px:actual-one-step-straightening};
the plain
average has
a different
derivative.

Let \(A_W(a)=D_aF_W(a)\).
Two scalar Cauchy
circles, in
arbitrary sup-norm
directions of
the reference
logarithm and
of \(a\), give
\begin{equation}
 \begin{aligned}
 \|A_W(a)\|_{\rm row}
    &\le L(1+C_L\delta),&
 \|A_W(a)\|_{\rm col}
    &\le L^{-3}(1+C_L\delta),\\
 \sup_b\sum_{c,b'}\|D_bD_{b'}F_W(a)(c)\|
    &\le C_L,&
 \delta&=\|\log W/i\|_\infty+\|a\|_\infty .
 \end{aligned}
 \label{px:finite-one-step-schur}
\end{equation}
Here and below the
norm of a color
entry is its
finite matrix
norm.
We explain why
the derivative
estimate is
independent of
a field-coordinate
count.
The two circles
bound a second
directional
derivative by
\(C L^2\).
For a fixed
output color
row, its
operator norm
from bounded
input colors
controls the
sum of its
entry norms,
up to the
fixed color
dimension:
sum the dual
norms of
its three
scalar rows.
A fixed input
bond affects
only the
coarse rows
at the endpoints
of the one or
two blocks
containing it.
Their number
is fixed.
Apply the
same argument
with one input
bond fixed
to obtain the
last line.
Dividing the
row and column
errors by their
flat values
costs only
powers of \(L\).
For definiteness,
all one-step
constants in
\eqref{px:finite-one-step-schur} may
be enlarged
to \(C L^{12}\).
This deliberately
loose polynomial
bound is before
any number
of levels is
introduced.

The multilevel recursion is
literal. Write
\(W_r=M^r(W_0)\),
\(a_0=\eta A\), and
\begin{equation}
                 a_{r+1}=F_{W_r}(a_r),\qquad
                   v_{r+1}=v_r|_{\rm next}\,u_{W_r}(a_r).
 \label{px:actual-straightened-recursion}
\end{equation}
Equivariance of the
plain mean gives,
by induction,
\begin{equation}
                M^r(e^{i\eta A}W_0)
                                  =(e^{ia_r}W_r)^{v_r}.
 \label{px:actual-factorization}
\end{equation}
Indeed apply \(M\) to the
previous factorization,
restrict its gauge
to the next
nodes, and then
use \eqref{px:actual-one-step-straightening}.
The gauge factors
multiply in the
order in \eqref{px:actual-straightened-recursion}.
They do not
alter the
recursion for
\(a_r\).

\subsubsection{Summability of the raw radii gives a depth-uniform row bound}

For the row estimates from this point on,
\(\epsilon\) is a common ceiling for the
real curvature and the local real-gauge
bounds. Thus the earlier \(a_{\rm loc}\)
and the curvature radius are both at
most \(\epsilon\); when they are supplied
separately, take their maximum here.
Existence of those local gauges remains
the primitive regular-pair hypothesis.
No bound for \(a_{\rm loc}\) is inferred
from curvature alone.

Fix a row of length
\(\xi=n\eta\),
\(n=L^j\).
On its determining
carrier, the
regular-pair gauge
clause gives
a real local
potential and
its first
difference at
scale \(\xi\),
with at most
the fixed
adjacent-layer
factor at a
seam.
The open-bond
and curvature
recurrences
\eqref{px:pillar-open-recurrence}, \eqref{px:complex-one-step-recurrences}
therefore give
\begin{equation}
 \|\log W_r/i\|_\infty
                    \le C_L(\epsilon+h)t_r,\qquad
                    t_r=L^r/n,\qquad r<j .
 \label{px:reference-raw-radii}
\end{equation}
These estimates are
used only on
the dependence
cone of that
row.
They also follow
by extending its
local small
potential with
a cutoff on
a fixed larger
cube and then
using locality.
The cutoff
keeps its first
difference at
the same scale.
The torus room
required for
nonflat local
classes permits
this cube;
the flat
periodic formulas
need no such
extension.
No gauge bound
growing with
\(M_1\) is
used.

Allow an arbitrary fine
increment with
\(\xi\|A\|_\infty\le r\).
It need not
have small
fine derivatives.
On the recursion
ball, \eqref{px:finite-one-step-schur}
and \(F_W(0)=0\)
give
\begin{equation}
 \|a_{r_0}\|_\infty\le 2r\,t_{r_0}
                  \quad(0\le r_0\le j)
 \label{px:fluctuation-raw-radii}
\end{equation}
when \(\epsilon+h+r\)
is below a
fixed threshold
depending on
\(L\).
To avoid confusing
the index with
the radius,
write
\(p_s=\|a_s\|_\infty/t_s\).
As long as
\(p_s\le2r\),
Taylor's formula
and \eqref{px:finite-one-step-schur}
imply
\begin{equation}
 p_{s+1}\le p_s
       [1+C_Lt_s(\epsilon+h+p_s)] .
 \label{px:radius-discrete-gronwall}
\end{equation}
The sum of \(t_s\)
is below
\((L-1)^{-1}\).
Choose the
fixed ceiling
so that
\(C_L(\epsilon+h+2r)/(L-1)<\log2\).
Multiplying
the factors
proves \eqref{px:fluctuation-raw-radii}
by first exit.
The raw radii
also stay
strictly inside
the \(c/L\)
one-step balls.
Finally
\(\|u_{W_s}(a_s)-1\|\le C_L\|a_s\|_\infty\);
their condition
numbers have
product at
most
\(\exp(C_L\sum_s\|a_s\|_\infty)\).
A further fixed
decrease makes
this at most
two, supplying
the room for
the complex
covariance in
\eqref{px:actual-factorization}.

Put \(A_s=D_aF_{W_s}(a_s)\)
and \(P_{b,a}=A_{b-1}\cdots A_a\).
The same summable
radii in \eqref{px:finite-one-step-schur}
give
\begin{equation}
 \|P_{b,a}\|_{\rm row}\le C L^{b-a},\qquad
                   \|P_{b,a}\|_{\rm col}\le C L^{-3(b-a)} .
 \label{px:multilevel-row-column-products}
\end{equation}
The constant is bounded
by the product
of
\(1+C_L(\|\log W_s/i\|_\infty+\|a_s\|_\infty)\).
It is independent
of \(a,b,j\).
Multiplying only
uniform bounds
larger than one
at each level
would miss this
point.

The per-length row logarithm
is \(q_j(A)=a_j/\xi\).
Its derivative with
respect to the
physical fine
coordinates is
\((\eta/\xi)P_{j,0}=n^{-1}P_{j,0}\).
Consequently
\begin{equation}
 \|Dq_j(A)\|_{\rm row}\le C,\qquad
 \|Dq_j(A)\|_{\rm col}\le C n^{-4},\qquad
                   |D_bq_j(A)(c)|\le C(\eta/\xi)^4 .
 \label{px:actual-row-density}
\end{equation}
These are counting-coordinate
derivatives.
A functional derivative
uses the additional
\(\eta^{-4}\).
Telescoping the
products against
their flat factors
also gives row
error \(C_L(\epsilon+h+r)\)
and column error
\(C_L(\epsilon+h+r)n^{-4}\),
in the common
local unitary
frames.

An actual row carrier has
only its two
endpoint blocks.
A fixed fine
bond in a
layer can occur
in carriers
of only adjacent
levels: the
outer block
is inside the
previous region
and both
endpoint blocks
are outside
the next.
At each of
those levels
the number
of rows is
fixed.
Thus \eqref{px:actual-row-density}
implies for
the full
layered field
map
\begin{equation}
       \|\Xi^{-1}Qa\|_{2,{\rm rows}}
                                  \le C_L\|a\|_{\rm mass}.
 \label{px:actual-row-mass-trace}
\end{equation}
For a single row,
its coefficient
sum is bounded
and each entry
is at most
\(C(\eta/\xi_c)^4\).
Cauchy therefore
gives the
analogue of
\eqref{px:flat-row-ltwo};
the carrier
overlap argument
of \eqref{px:flat-field-row-upper}
then proves
\eqref{px:actual-row-mass-trace}.
This estimate
has been
obtained from
the actual
finite recursion,
not by treating
an infinity
norm bound
alone as an
\(L^2\) bound.

\subsubsection{The two-legged constraint coefficient has a uniform moment}

The same computation gives
the coefficient estimate
needed for the
quadratic current
term.
Let \(n_s=L^s\)
and
\(m_s=L^{j-s-1}\).
Twice differentiating
the recursion gives
\begin{equation}
 D^2_{a_0}a_j
      =\sum_{s<j}P_{j,s+1}
                      D^2F_{W_s}(a_s)[P_{s,0}\,\cdot,P_{s,0}\,\cdot].
 \label{px:actual-hessian-chain}
\end{equation}
Fix the first fine
bond \(b\) and
sum absolutely
over the second.
For the term
indexed by \(s\),
the largest
entry of the
outer product
is at most
\(C m_s^{-3}\),
by its column
sum.
The sum over
the first
inner leg is
at most
\(C n_s^{-3}\).
The second
inner row
sum is at
most \(C n_s\).
The remaining
sum over the
one-step Hessian,
with its first
input fixed,
is bounded
by \eqref{px:finite-one-step-schur}.
Thus
\begin{equation}
 \sum_{b'}\|D_{a_0(b)}D_{a_0(b')}a_j(c)\|
                 \le C_L\sum_{s<j}m_s^{-3}n_s^{-2}.
 \label{px:raw-hessian-leg-sum}
\end{equation}
This step uses a
column sum on
the fixed leg
and a row sum
on the summed
leg. Using
two row sums
would lose the
fine density.

Physical input derivatives
contribute \(\eta^2\);
the output per-length
normalization contributes
\(\xi^{-1}\).
Since
\(m_s^{-3}n_s^{-2}
=L^3n_s/n^3\)
and
\(\sum_{s<j}n_s\le n/(L-1)\),
we obtain
\begin{equation}
 \sup_b\sum_{b'}\|D_bD_{b'}q_j(A)(c)\|
                  \le C_L\,\xi\,(\eta/\xi)^4 .
 \label{px:constraint-hessian-leg-moment}
\end{equation}
There is no
number-of-levels
factor.
The geometric
sum is dominated
by its final
scales.
Symmetry of
the second
derivative gives
the corresponding
other-leg sum.

Let \(\omega(c)\) be
an arbitrary row
current with
\(|\omega(c)|\le\gamma_\omega\xi_c^{-3}\).
Pair the Hessian
at zero with
the physical
row volume.
\eqref{px:constraint-hessian-leg-moment},
Schur and the
same bounded
carrier overlap
give
\begin{equation}
 \left|\sum_c\xi_c^4
              \langle\omega(c),D^2q_c(0)[u,v]\rangle_{\rm col}\right|
                  \le C_L\gamma_\omega
                                      \|u\|_{\rm mass}\|v\|_{\rm mass}.
 \label{px:quadratic-current-form-bound}
\end{equation}
For clarity, the
row prefactor
\(\xi_c^4|\omega(c)|\)
is bounded by
\(\gamma_\omega\xi_c\).
Multiplying it
by the leg
moment
\(C_L\xi_c(\eta/\xi_c)^4\)
leaves
\(C_L\gamma_\omega\eta^4\xi_c^{-2}\)
on each fine
leg.
Adjacent scales
replace \(\xi_c\)
by the fine
bond scale
at only a
fixed cost.

The holomorphic ball was
used to extract
these coefficients.
The quadratic
form in \eqref{px:quadratic-current-form-bound}
then extends
to every finite
fine vector by
its coefficient
formula.
It does not
require a
Gaussian sample
to lie in
that holomorphic
ball.
For the second
variant one
must still
bound the
actual induced
current
\(\widehat H[J]^\sharp J\)
in this row
norm, and
bound \(\mathsf S\)
in the weighted
fine energy
spaces.
Those are
separate operator
readouts.

\subsubsection{The nonflat field-row lift follows from the row comparison}

The coefficient estimates allow
one to complete
\eqref{px:nonflat-row-error} at a
regular pair.
Keep the real
local gauges
of its center.
Construct the
profiles \eqref{px:row-dual-bump}
in each pillar
in such a gauge,
and undo the
gauge.
For a seed
assigned to the
far endpoint
of a row,
normalize its
color by the
fixed unitary
transport of
that row.
Denote this
operator again
by \(E^0\).
It is fixed
at the real
center while
the complex
increment varies.

Two local real gauges used
on the same
pillar differ
there by a
nearly constant
unitary matrix.
To prove this,
write their
connection equation
on each fine
bond and telescope
along a coordinate
path of length
\(C\xi_j\).
Both gauge
potentials have
size at most
\(C_L\epsilon/\xi_j\).
The transition
therefore differs
from its root
value by at
most \(C_L\epsilon\).
The first-order
class on a
fixed larger
cube supplies
both gauges.
Only a path
of size
\(C_L\xi_j\)
is used, not
the full
\(M_1\xi_j\)
cube.

It follows from \eqref{px:actual-row-density}
and this transition
bound that the
same-level seed
response is the
identity with
an error of
row and column
norm
\(C_L(\epsilon+h)\)
in the row
energy norm.
Only a fixed
number of
same-level rows
meet a seed
pillar.
The next-level
response has
a leading part
with the
positive scalar
majorant of
\(Q_L^0\), multiplied
by fixed unitary
color matrices,
and an error.
There are at
most \(C L^4\)
lower-level
seeds in one
such carrier;
a seed affects
at most a
fixed number
of next-level
rows.
An error entry
of size
\(C_L(\epsilon+h)\)
therefore has
row sum
\(C_LL^5(\epsilon+h)\)
and column sum
\(C_LL(\epsilon+h)\)
after the
row-energy
normalization.
Schur bounds
this error by
\(C_LL^3(\epsilon+h)\).
Absorb this
fixed power
into \(C_L\).

All other level blocks
are zero by
the carrier
argument preceding
\eqref{px:flat-average-composition}.
The leading
next-level part
still has
norm at most
\(L^{-1}\);
Schur uses
only the
scalar majorants,
not a compatible
global choice
of flat color
frames.
We have proved
\begin{equation}
 \|Q_{\mathbf U}E^0-1\|_{{\rm row}\to{\rm row}}
                 \le L^{-1}+C_L(\epsilon+h)=:\vartheta<1 .
 \label{px:actual-row-lift-error}
\end{equation}
The ceiling needed here
depends on \(L\)
and not on
depth, volume
or \(M_1\).
The actual
right inverse is
\begin{equation}
             E_{\rm row}=E^0(Q_{\mathbf U}E^0)^{-1}.
 \label{px:actual-layered-row-lift}
\end{equation}
Both its squared mass norm
and its covariant
gradient energy
are bounded by
\(C_L(1-\vartheta)^{-2}\|\psi\|_{\rm row}^2\).
For the gradient
of a seed,
expand in its
local gauge:
the profile
derivative is
of size
\(C\xi_j^{-1}|\psi|\)
and the
connection term
has the same
bound.
Supports and
their fine
collars have
bounded overlap.
The small
complex perturbation
adds at most
\(C_Lh\) times
the mass norm.
This proves the
stated energy
bound, including
bare seeds.

There is a corresponding
actual Feynman
Poincare inequality.
Repeat \eqref{px:flat-row-mean-comparison}
in a common
local real gauge
for a pillar
and its chosen
row carriers.
The row comparison
after \eqref{px:actual-row-density}
has both row
and column error
bounds.
Schur therefore
bounds its
local \(L^2\)
functional norm
by
\(C_L\epsilon\xi_j^{-2}\).
Ordinary gradients
and covariant
gradients differ
there by
\(C_L\epsilon\xi_j^{-1}\)
times the field.
Consequently
\eqref{px:flat-field-local-poincare}, summed
over the carriers,
gives
\begin{equation}
 \|a\|_{\rm mass}^2
       \le C_L\bigl(\|\nabla_Ua\|_2^2+
                         \|a^{1/2}Q_Ua\|_{2,{\rm rows}}^2\bigr)
                         +C_L\epsilon^2\|a\|_{\rm mass}^2 .
 \label{px:actual-feynman-poincare}
\end{equation}
The repeated letter \(a\)
in \(a^{1/2}\)
is the fixed
positive row
weight, whereas
the arguments
are fine fields.
Equivalently that
term is
\(\sum_cc_c\xi_c^2|Q_Ua(c)|^2\).
Absorb the last
term by the
fixed weak-field
ceiling.
Only carrier
differences in
the interior
are used;
the bare
boundary fields
are controlled
by their
identity rows.
Thus the same
argument applies
to an active
component-gradient
form containing
all those
interior
differences.

The full complex Feynman
form is stable
on a fixed
physical ball.
The actual row
comparison gives
\begin{equation}
 \|a^{1/2}(Q_{\mathbf U}-Q_U)v\|_2
                       +\|(\nabla_{\mathbf U}-\nabla_U)v\|_2
                                \le C_Lh\|v\|_{\rm mass}.
 \label{px:feynman-coefficient-comparison}
\end{equation}
For the first term
use the row
and column
differences, the
physical row
volumes and
the bounded
carrier overlap.
For the second
use
\(\operatorname{Ad}(e^{i\eta E})-1=O(\eta E)\).
Expanding the
two bilinear
squares and
using \eqref{px:actual-feynman-poincare}
gives a
relative form
perturbation
\(C_Lh\).
Its inverse
exists by
the normalized
Neumann series.
This proves a
component Feynman
inverse statement.
The actual
Landau subtraction
\eqref{px:landau-feynman-difference} is
still absent.


% END MODULE chapters/layered_current_precision.tex

% BEGIN MODULE chapters/nonlinear_response_reductions.tex
\subsection{Response equations and weighted inverse estimates}
\label{ch:nonlinear-response-reductions}

The response of a constrained background is a solution of a coupled
field--multiplier system.  Bounds on the inverse of one positive quadratic
form give only part of that response.  Constraints, gauges, the current
readout, the local length factors, and admissibility of the substituted data
must all be included.  This section derives the exact reductions needed for
these tasks. Its linear algebra and analytic estimates are proved.
On an admitted strong reference domain,
\cref{prop:ncw-named-families,prop:scf-strong-current-families,prop:ngc-joint-connected-family}
construct the primitive coefficients and the selected nonlinear
current germ. The construction of a reference from homogeneous
curvature-small data is \cref{prop:cta-curvature-critical-atlas}.
Applications to complementary histories retain their specified
geometry, source and admission hypotheses.

All calculations are first made in a finite compatible geometry.  The
constants in an infinite-volume or arbitrary-depth application must be
uniform in that geometry before a limit is taken.  In particular, a
finite-dimensional implicit-function theorem by itself supplies no uniform
radius.

\subsubsection*{The actual equations and their differentiated right-hand sides}

Use the real reference field $U_0$ and write the varied fine field as
$e^{i\eta A}U_0$.  For the multilevel determining set $\mathbf B$, the
straightened average is denoted by $\mathcal Q^{\rm st}(A)$ and its
derivative by $K_A$.  This is the nonlinear straightening of the group
average, rather than the straight linear average.  With the reduced datum
$\mathsf b(B)$, the physical critical pair satisfies
\refstepcounter{equation}
\begin{align*}
 F'(A)&=K_A^{\mathsf T}\omega,
       &\mathcal Q^{\rm st}(A)&=\mathsf b(B),\\
 RD^{\mathsf T}A&=0,
       &F(A)&=A_{\Omega_0}(e^{i\eta A}U_0).
 \tag{\theequation}\label{eq:nonlinear-response-NR1}
\end{align*}
Here $R$ and $D$ are formed at the reference field.  The first equation is
the full physical stationarity identity when that identity has been
proved; its restriction to the slice is sufficient for the algebra below.
When the reference is changed with the datum, differentiating the slice
also contributes $(\dot R D^{\mathsf T}+R\dot D^{\mathsf T})A$.
We keep the reference fixed in \eqref{eq:nonlinear-response-NR1}.  A calculation that moves the
reference and then omits this contribution does not differentiate \eqref{eq:nonlinear-response-NR1}.

The reduced datum has a local triangular correction at an outward bond.
For an outer endpoint $z=c_+$ it is defined by
\begin{equation}
 e^{i\mathsf b(B)(c)}
   =e^{iB(c)}\operatorname{Ad}(\mathbf V_0(c))\mathsf v_z(B).
 \label{eq:nonlinear-response-NR2}
\end{equation}
At an interior bond it equals $B(c)$.  The block factor $\mathsf v_z$
depends on the interior bonds of its block.  Consequently the derivative
of $\mathsf b$ is triangular in the same ordering.  On a smaller regular
logarithm chart, differentiation of the exponential and logarithm gives
\begin{equation}
 \dot b= D\mathsf b(B)\dot B,\qquad
 \dot b(c)=L_c(B)\dot B(c)+\sum_{c'\ {\rm interior}}
                     T_{cc'}(B)\dot B(c')
 \label{eq:nonlinear-response-NR3}
\end{equation}
on the outward rows.  Define the right logarithmic differential
$\mathcal E_X(V)=(1/i)
 [\frac{d}{dt}\exp(i(X+tV))]_{t=0}\exp(-iX)$.
Differentiating \eqref{eq:nonlinear-response-NR2} after right multiplication by $\exp(-i\mathsf b)$
gives $L_c=\mathcal E_{\mathsf b(B)(c)}^{-1}\mathcal E_{B(c)}$.
The same differential and the derivative of the block factor give
$T$.  Thus the outward diagonal is invertible near the regular chart;
it need not be the identity.  The bounds require a fixed distance to the logarithm
cut and the actual local block-factor estimate.  Triangularity establishes
invertibility once those diagonal logarithmic derivatives are invertible;
it does not bound an unrestricted product of triangular operators over
an arbitrary number of levels.

Let $\mathcal V$ be the fixed linear slice, after the residual node gauge
has been fixed.  Its fine pairing and the datum pairing are the actual
cell-weighted pairings.  Write $q=K_A|_{\mathcal V}$ and let $q^\sharp$
be its transpose for these pairings.  The constrained second variation is
\begin{equation}
 \mathcal K=F''(A)
       -\sum_a \mu_a\omega_a
                       (\mathcal Q^{\rm st}_a)''(A)
 \quad\hbox{on }\mathcal V.
 \label{eq:nonlinear-response-NR4}
\end{equation}
The weights $\mu_a$ in this coordinate expression are the datum cell
weights.  Equivalently the second term is the Hessian of the scalar
pairing $\langle\omega,\mathcal Q^{\rm st}(A)\rangle$.  This latter
description fixes the convention if a different coordinate basis is used.
The curvature term in \eqref{eq:nonlinear-response-NR4} cannot be replaced by the Hessian of $F$
alone.  It is the usual second fundamental form of the constraint fibre.

Differentiating \eqref{eq:nonlinear-response-NR1} at a fixed reference gives the exact system
\begin{equation}
 \mathcal K\dot A-q^\sharp\dot\omega=0,
       \qquad q\dot A=D\mathsf b(B)\dot B.
 \label{eq:nonlinear-response-NR5}
\end{equation}
More generally, for explicit parameters $p$ in the action, gauge or
constraints, put $\mathcal C(A,p)=0$ and
$\mathcal E(A,\omega,p)=F'_p(A)-\mathcal C_A^\sharp\omega$.
The right-hand sides are then
\begin{equation}
 f=-D_p\mathcal E\,\dot p,\qquad
 b=-D_p\mathcal C\,\dot p,\qquad
 \mathcal K\dot A-q^\sharp\dot\omega=f,\quad q\dot A=b.
 \label{eq:nonlinear-response-NR6}
\end{equation}
Every differentiation in $D_p$ is made before substituting the solution
$A(p),\omega(p)$.  This distinction prevents an implicit extra response
term on the right-hand side.

An independent-current extension is an additional specified operator.
The formula in \eqref{eq:nonlinear-response-NR4} identifies the physical constrained Hessian.  It does not
identify every separately defined $\Delta'_\pi[J]$ with that Hessian.
Likewise, an absolute contour gauge functional and its relative version
have the same quadratic gauge block at an axial field, but generally have
different blocks away from that field.  The complete gauge and current
normal form must be matched before an independently named propagator can
be inserted into \eqref{eq:nonlinear-response-NR5}.

\subsubsection*{A complete block inverse, including the multiplier}

The following result is finite-dimensional and exact.  It also applies to
bounded operators when the displayed inverses and compositions exist on
the stated spaces.  Every transpose in the complex calculation is
bilinear, with no conjugation.

\begin{proposition}[Response from a primitive constrained inverse]
Suppose $q:\mathcal V\to\mathcal D$ is onto, $E:\mathcal D\to\mathcal V$
satisfies $qE=I$, and $P:\mathcal X\to\mathcal V$ is an isomorphism onto
$\ker q$.  Assume
\begin{equation}
 H=P^\sharp\mathcal K P\quad\hbox{is invertible},\qquad
 L=P H^{-1}P^\sharp,\qquad T=E-L\mathcal K E.
 \label{eq:nonlinear-response-NR7}
\end{equation}
Then the unique solution of \eqref{eq:nonlinear-response-NR6} is
\refstepcounter{equation}
\begin{align*}
 \dot A&=Lf+Tb,\\
 \dot\omega&=E^\sharp(\mathcal K L-I)f
                          +E^\sharp\mathcal K T b.
 \tag{\theequation}\label{eq:nonlinear-response-NR8}
\end{align*}
If $\mathcal K^\sharp=\mathcal K$, this simplifies to
\begin{equation}
 \dot\omega=-T^\sharp f+M b,\qquad
 M=E^\sharp\mathcal K E
                     -E^\sharp\mathcal K L\mathcal K E,
       \qquad M^\sharp=M.
 \label{eq:nonlinear-response-NR9}
\end{equation}
\end{proposition}

\begin{proof}
Write $\dot A=Eb+Pv$.  Applying $P^\sharp$ to the first equation of
\eqref{eq:nonlinear-response-NR6}, and using $qP=0$, gives
$Hv=P^\sharp(f-\mathcal K Eb)$.  This determines $v$ and gives the first
formula of \eqref{eq:nonlinear-response-NR8}.  The residual $\mathcal K\dot A-f$ annihilates $\ker q$.
It therefore belongs to the range of $q^\sharp$.  Since
$E^\sharp q^\sharp=I$, its unique multiplier is the second formula of
\eqref{eq:nonlinear-response-NR8}.  Conversely these formulas satisfy both equations, so there is
no lost multiplier equation.  For a homogeneous solution, the same
argument gives $v=0$, then $\dot A=0$ and $\dot\omega=0$.

For the symmetric assertion, $H^\sharp=H$ and $L^\sharp=L$.
Thus $T^\sharp=E^\sharp-E^\sharp\mathcal K L$, which gives \eqref{eq:nonlinear-response-NR9}.
Finally $E^\sharp\mathcal K T$ equals the displayed expression for $M$;
both its terms are symmetric.  Positivity of the full $\mathcal K$ was
not used.  Invertibility was needed on the actual constrained slice.
\end{proof}

The inverse $L$ acts on fine forces; $T$ lifts the constraint datum after
the energetic correction; $M$ is the multiplier response.  Bounding $L$
alone leaves the products $L\mathcal K E$ and
$E^\sharp\mathcal K T$ to be controlled.  The size and locality of $E$
also matter.  A right inverse with a norm growing in the number of layers
would not prove a uniform response theorem.

At one step the actual raw pivot derivative has size $L^{-3}$, whereas
the derivative after division of the coarse logarithm by $L$ has size
$L^{-4}$.  Accordingly a change of the constraint normalization changes
both $E$ and the dual multiplier.  For example, replacing $q$ by
$\widetilde q=q/L$ while keeping the same pairings replaces
$\omega$ by $\widetilde\omega=L\omega$ in the scalar Lagrangian.
Using the new right inverse with the old multiplier would change
\eqref{eq:nonlinear-response-NR4}, not merely its notation.

\subsubsection*{Cell weights and a kernel algebra}

We describe the norm in which the preceding formulas can be estimated
without a bare-mesh counting loss.  Let each component $i$ have a positive
cell weight $\mu_i$ and an amplitude $s_i>0$.  A vector is measured by
$\|v\|_s=\sup_i |v_i|/s_i$.  A kernel is written as
\begin{equation}
 (Av)_i=\sum_j\mu_j A_{ij}v_j,
       \qquad \widehat A_{ij}=s_i^{-1}A_{ij}s_j.
 \label{eq:nonlinear-response-NR10}
\end{equation}
For different input and output species their respective amplitudes and
cell weights are used.  The identity kernel is $\delta_{ij}/\mu_j$.
Thus even the identity is not represented by $\delta_{ij}$ when it is
used with a weighted summation convention.

For a metric $d$ common to the component locations define
\begin{equation}\begin{aligned}
 \|A\|_m=\max\bigg\{&\sup_i\sum_j\mu_j
                 |\widehat A_{ij}|e^{m d(i,j)},\\
                 &\sup_j\sum_i\mu_i
                 |\widehat A_{ij}|e^{m d(i,j)}\bigg\}.
\end{aligned}\label{eq:nonlinear-response-NR11}\end{equation}
Matrix norms in a colour or local finite-dimensional fibre are operator
norms.  The two sums use the correct weights for their two species.
The triangle inequality for $d$ and absolute convergence give
\begin{equation}
 \|AB\|_m\le\|A\|_m\|B\|_m.
 \label{eq:nonlinear-response-NR12}
\end{equation}
Indeed, insert the intermediate sum with weight $\mu_k$, use
$d(i,j)\le d(i,k)+d(k,j)$ and cancel the intermediate amplitudes.
For the row bound take the supremum of the inner $B$ row sum; for the
column bound use the corresponding $A$ column sum.  These are distinct
calculations and neither counts the number of fine bonds.

The first row sum bounds $A$ on the normalized supremum norm.  The two
sums together bound it on the normalized, cell-weighted $\ell^2$ space
by the Schur estimate.  To see the latter directly, apply weighted
Cauchy--Schwarz to each row with weights
$\mu_j|\widehat A_{ij}|$, then sum in $i$ and use the column bound.
The result is $\|Av\|_{2,s}\le\|A\|_0\|v\|_{2,s}$.

For the natural response components above a top of scale $\xi_y$, the
amplitudes are
\begin{equation}
 s_H(y)=\xi_y^{-1},\qquad
 s_{\nabla H}(y)=\xi_y^{-2},\qquad
 s_J(y)=\xi_y^{-3}.
 \label{eq:nonlinear-response-NR13}
\end{equation}
Equivalently the output is
$(\xi_yH,\xi_y^2\nabla H,\xi_y^3J)$.  The cell weights are determined by
the fine, coarse and pillar pairings.  Each constraint length factor belongs in the
matrix before \eqref{eq:nonlinear-response-NR10} is formed.  The domain section gives the resulting
transposes.  These transpose conventions are needed when applying \eqref{eq:nonlinear-response-NR12} to a product
that changes species at a seam.

For a raw logarithm $a=g y$, a physical radius $\alpha$ becomes
$\alpha/g$ in the canonical Gaussian variable $y$.  A derivative with
respect to $y$ contributes $g$ times the derivative with respect to $a$.
Thus the second derivative of a scalar activity in Gaussian coordinates
is $g^2$ times its raw-logarithm Hessian.  This conversion occurs in
addition to the length and cell-weight conversions in \eqref{eq:nonlinear-response-NR10}--\eqref{eq:nonlinear-response-NR13}.

\begin{corollary}[Uniform response after the primitive bounds]
Suppose the actual normalized kernels in \eqref{eq:nonlinear-response-NR7} satisfy
\begin{equation}
 \|L\|_m\le a,\quad \|E\|_m\le e,\quad
 \|\mathcal K E\|_m\le k,
       \quad \|E^\sharp\mathcal K\|_m\le k'.
 \label{eq:nonlinear-response-NR14}
\end{equation}
Then
\begin{equation}
 \|T\|_m\le e+ak,\qquad
 \|M\|_m\le k'(e+ak).
 \label{eq:nonlinear-response-NR15}
\end{equation}
In the nonsymmetric case the force-to-multiplier block has norm at most
$k'a+\|E^\sharp\|_m$.  In the symmetric case it is $-T^\sharp$.
Its norm equals the norm of $T$ with reversed species if the dual
amplitudes are chosen reciprocally for the same cell-weighted pairings.
For other force and current amplitudes its weighted transpose bound
must be checked separately.
\end{corollary}

The last qualification is necessary.  A raw transpose and a transpose
after conjugation by amplitude weights are different matrices.  If
$\widehat T_{ij}=s_i^{-1}T_{ij}s_j$, its transposed normalized kernel
is obtained by giving the dual components amplitudes $s_i^{-1}$ and
$s_j^{-1}$.  A natural force weight of a different physical dimension
does not satisfy this rule automatically.  Its conversion belongs in
the kernel before a Schur norm is claimed.

The bounds in \eqref{eq:nonlinear-response-NR14} are the primitive targets.  They must hold for the
actual enlarged saddle system on the retained complex domain and at
every allowed seam.  The corollary does not assert them.  In particular,
an unweighted flat inverse does not supply $k,k'$ on a current-dependent
layered domain.  It does establish an exact-model anchor from which a
perturbative comparison can start.

\subsubsection*{Invertibility on the whole complex ball}

Let $\mathcal S$ denote the full field--constraint operator in \eqref{eq:nonlinear-response-NR6},
with all gauge rows included.  Fix a reference operator $\mathcal S_0$
on the same spaces.  If the primitive estimates give
\begin{equation}
 \|\mathcal S_0^{-1}\|_m\le C_0,
 \qquad
 \|\mathcal S(B)-\mathcal S_0\|_m\le\varepsilon,
       \qquad C_0\varepsilon<1,
 \label{eq:nonlinear-response-NR16}
\end{equation}
then the Neumann series converges in the kernel algebra and proves
\begin{equation}
 \|\mathcal S(B)^{-1}\|_m
          \le {C_0\over1-C_0\varepsilon}.
 \label{eq:nonlinear-response-NR17}
\end{equation}
It is essential that the second bound in \eqref{eq:nonlinear-response-NR16} holds on the whole
retained complex ball.  A bound only along the real solutions does not
give the analytic disc needed later for Cauchy estimates.

If $\mathcal S_0$ is unbounded in the chosen norm, the correct statement
uses the relative perturbation
$\mathcal S_0^{-1}(\mathcal S-\mathcal S_0)$ between its actual graph
spaces.  The same series applies there.  In particular an elliptic
second-order operator cannot be treated as a bounded supremum-norm
matrix with a mesh-independent norm merely because its inverse is
bounded.

If the coefficients of $\mathcal S(B)$ are holomorphic in the indicated
operator norm, holomorphy follows term by term on a strictly smaller
disc, or by the resolvent identity
\begin{equation}
 \mathcal S(B')^{-1}-\mathcal S(B)^{-1}
   =-\mathcal S(B')^{-1}
              [\mathcal S(B')-\mathcal S(B)]\mathcal S(B)^{-1}.
 \label{eq:nonlinear-response-NR18}
\end{equation}
The identity also proves continuity and the derivative formula wherever
the coefficient derivative exists in the indicated operator norm.
No complex positivity assertion is used.  Real coercivity may construct
$\mathcal S_0^{-1}$; the complex inverse then requires a genuine norm
comparison or another explicitly convergent construction.

For a complete identification of the proposed response, \eqref{eq:nonlinear-response-NR16} has to be
checked with the nonlinear $K_A$, the full constraint-curvature block,
the prescribed relative or absolute gauge block, and the seam rows.
Changing any of these while keeping the same notation for
$\mathcal S$ invalidates the comparison.

\subsubsection*{From local forcing to first response}

Suppose a change in one datum coordinate $B(c)$ gives a forcing column
$F_c$ in the enlarged system.  The column consists of both the field
force and the constraint force in \eqref{eq:nonlinear-response-NR6}.  If it is supported within
distance $r_0$ of the top $y_c$ and has normalized supremum norm at most
$b_0$, then \eqref{eq:nonlinear-response-NR17} gives
\begin{equation}
 |D_{B(c)}Z_y[e]|/s_y
   \le {C_0b_0 e^{mr_0}\over1-C_0\varepsilon}
                                  e^{-m d(y,y_c)}|e|.
 \label{eq:nonlinear-response-NR19}
\end{equation}
Here $Z$ contains the output species actually controlled by the inverse,
with amplitude $s_y$ from \eqref{eq:nonlinear-response-NR10}. Equivalently apply the bound to
the normalized readout $Z_y/s_y$; subsequent Cauchy estimates use
normalized readouts.
The proof is the row estimate: for each contributing component $j$,
$d(y,y_c)\le d(y,j)+r_0$, so multiply by $e^{m d(y,y_c)}$ and sum the
kernel against $F_c$.  A triangular datum correction from \eqref{eq:nonlinear-response-NR3} changes
$b_0,r_0$ by its actual block bounds.  It must be included in $F_c$.

For an exponentially local forcing column, a bound
$|F_c(j)|/s_j\le b_0e^{-a d(j,y_c)}$ gives \eqref{eq:nonlinear-response-NR19} with any
$0<m'<\min(m,a)$ and $m'$ in place of $m$.  Indeed the triangle
inequality leaves a factor $e^{-(a-m')d(j,y_c)}\le1$, and the remaining
sum is a weighted row sum at rate $m'$.  No uncontrolled number of
sites appears.

This proves response decay from three smaller tasks: construct the
actual primitive inverse, verify its weighted row norm, and compute the
actual datum forcing.  It does not assume response decay as a premise.
If the output inverse controls only $H$ and the multiplier, the estimate
has not yet been proved for $\nabla H$ or the current.  Those outputs
require their own readout maps.

For example the physical current is a specified nonlinear function
$J(U)$.  In a fixed logarithmic chart its derivative is
\begin{equation}
 \dot J=D_A[J(e^{i\eta A}U_0)]\dot A.
 \label{eq:nonlinear-response-NR20}
\end{equation}
An implementation may replace this expression by the multiplier and
lower-order terms through the differentiated field equation.  Such a
replacement must be derived from \eqref{eq:nonlinear-response-NR1} and the definition of $J$.
The equality is not implied by the fact that both quantities have
dimension three.  On an independent-current core the variable $\mathbf J$
and the computed current $J(\mathbf U)$ are distinct until the physical
pair is imposed.

\subsubsection*{Why a kernel bound does not prove a difference bound}

Uniform locality and a gap do not control variation at the bare scale.
On a one-dimensional mesh, the positive diagonal inverse
$R_{ii}=1+\tfrac12(-1)^i$ has norm at most $3/2$ and zero range.
Applied to a constant force it produces a difference of size
$\eta^{-1}$.  Its inverse is positive and uniformly bounded as well.
This example shows exactly why a separate regularity estimate is needed.

There is a useful exact identity for that estimate.  On a translation
patch let $T$ shift by one mesh step, $\nabla=(T-I)/\eta$ and
$\mathcal K_T=T\mathcal K T^{-1}$.  If $\mathcal K v=f$, then
\begin{equation}
 \mathcal K_T\nabla v
    =\nabla f-{\mathcal K_T-\mathcal K\over\eta}v.
 \label{eq:nonlinear-response-NR21}
\end{equation}
To prove it, write $T(\mathcal K v)=\mathcal K_TTv$, subtract
$\mathcal K v$ and divide by $\eta$.  At a boundary this calculation
must be made on the common interior patch; it does not shift an
inadmissible boundary row into the equation.

For a scalar divergence-form example
$\mathcal K=D^\sharp aD+V$, with constant mesh difference $D$,
the coefficient commutator in \eqref{eq:nonlinear-response-NR21} is
\begin{equation}
 {\mathcal K_T-\mathcal K\over\eta}
      =D^\sharp(\nabla a)D+\nabla V.
 \label{eq:nonlinear-response-NR22}
\end{equation}
If the energy norm is
$\|v\|_{\mathcal E}^2=\|Dv\|_2^2+\|v\|_2^2$, the bilinear pairing
and Cauchy--Schwarz imply
\begin{equation}
 \big\|[(\mathcal K_T-\mathcal K)/\eta]v\big\|_{\mathcal E^*}
 \le (\|\nabla a\|_\infty+\|\nabla V\|_\infty)
                                               \|v\|_{\mathcal E}.
 \label{eq:nonlinear-response-NR23}
\end{equation}
Thus a bounded inverse from $\mathcal E^*$ to $\mathcal E$, together
with coefficient regularity, bounds $\nabla v$ in energy norm.  This
does not assert a supremum bound; passing to a supremum or H\"older
bound requires the corresponding local elliptic estimate.

In the gauge system, shifting the covariant differences adds transport
commutators, the average rows add contour terms, and a seam adds its
own boundary comparison.  The required finite list is: the translated
interior equation, the gauge-row commutator, the constraint-row
commutator, and each seam or exterior boundary row.  Each term must
be bounded between its actual force and output norms.  The uniform
primitive regularity estimate is precisely this list of bounds and
the associated local elliptic estimate.  It is not supplied by
exponential decay of $H^{-1}$ alone.

\subsubsection*{Joint Cauchy estimates with joining distances}

Once \eqref{eq:nonlinear-response-NR19} has been established uniformly on a complex product ball,
several useful higher response bounds require no further model
estimate.  We give the argument with all radius losses visible.

Let $\mathcal B_r$ be a ball in the supremum of the dimensionless
bond logarithms, with a fixed norm on the three colour coordinates.
Suppose a holomorphic map $Z_y$ on $\mathcal B_r$ satisfies
\begin{equation}
 |\partial_c Z_y(B)[e]|
        \le C e^{-m d(y,y_c)}|e|,
                       \qquad B\in\mathcal B_r.
 \label{eq:nonlinear-response-NR24}
\end{equation}
Fix $0<\theta<1$.  There is a colour constant $c_*$ such that the
coordinate polydisc of radius $c_*(1-\theta)r$ about every
$B\in\mathcal B_{\theta r}$ lies in $\mathcal B_r$.  The number of
bonds does not enter this assertion, since the norm is a supremum over
bonds; simultaneous variation of three colour components changes only
$c_*$.  Cauchy's formula applied to the first derivative gives
\begin{equation}
 |\partial_{c_1}\cdots\partial_{c_q} Z_y(B)|
 \le C(q-1)!
       [c_*(1-\theta)r]^{-(q-1)}
                       e^{-m\max_j d(y,y_{c_j})}.
 \label{eq:nonlinear-response-NR25}
\end{equation}
Repeated coordinates are allowed.  The relevant multi-index factorial
is at most $(q-1)!$.  Holomorphic derivatives commute, so each one of
the $q$ coordinates can be selected as the first derivative in
\eqref{eq:nonlinear-response-NR24}; taking the minimum of the resulting exponential bounds proves
the maximum in \eqref{eq:nonlinear-response-NR25}.

Let $\mathcal T(y,c_1,\ldots,c_q)$ be the minimum spanning-tree length
in the underlying metric.  The star through $y$ is a competing tree,
and therefore
\begin{equation}
 \mathcal T\le\sum_{j=1}^q d(y,y_{c_j})
                \le q\max_jd(y,y_{c_j}).
 \label{eq:nonlinear-response-NR26}
\end{equation}
Consequently \eqref{eq:nonlinear-response-NR25} gives a joining-tree exponential with rate $m/q$.
For the second derivative it also gives the more explicit bound
\begin{equation}
 |\partial_c\partial_{c'}Z_y(B)|
 \le {C\over(1-\theta)r}
       e^{-\frac m2[d(y,y_c)+d(y,y_{c'})]}.
 \label{eq:nonlinear-response-NR27}
\end{equation}
This is the estimate used in the pointed Taylor remainder.  It costs
one datum radius.  The formula in \eqref{eq:nonlinear-response-NR25} is not a uniform-in-$q$ tree bound
at a fixed rate: its rate deteriorates with $q$.  If an arbitrary-order
polymer estimate needs a fixed rate, one must use the local nonlinear
equations and their tree expansion, with the resulting factorial and
entropy costs.  A fixed-order application may use \eqref{eq:nonlinear-response-NR25} directly.

The same Cauchy argument applies to an output H\"older seminorm if
the first response is already bounded in that seminorm on the larger
ball.  Cauchy's formula can be applied to every normalized difference
$[Z_y-Z_{y'}]/|y-y'|^\gamma$ and then a supremum taken.  It cannot
create spatial H\"older regularity from a pointwise bound.

\subsubsection*{Finite boxes and compatible boundary forcing}

Let $\mathcal S$ be the actual linearized saddle operator and
$\mathcal S_\Lambda$ its prescribed finite-box version.  To compare
solutions, their equations must be put on common variables and rows
in an interior collar.  Write the resulting exact difference equation
as
\begin{equation}
 \mathcal S_\Lambda\,\delta X=f_\Lambda.
 \label{eq:nonlinear-response-NR28}
\end{equation}
The force includes the difference of boundary rows, the datum
restriction, the gauge-restoration term and any cutoff commutator.
The identity in \eqref{eq:nonlinear-response-NR28} is checked before its right-hand side is bounded.
Replacing an actual box boundary by a convenient Dirichlet boundary
changes this force and requires a new comparison.

Suppose the inverse has weighted row bound $C$ at rate $m$ and the
force is supported in a collar $\mathcal C$.  If
$\|f_\Lambda\|_s\le A R^p$, then for a component at $y$
\begin{equation}
 |\delta X_y|/s_y
       \le CA R^p e^{-m\operatorname{dist}(y,\mathcal C)}.
 \label{eq:nonlinear-response-NR29}
\end{equation}
This follows by extracting the exponential distance from the weighted
row sum.  In particular no factor equal to the number of bare sites
in the collar is necessary.  Such a factor would generally destroy
uniformity as the mesh tends to zero.

The same conclusion, at any rate $m'<\min(m,a)$, holds if
\begin{equation}
 |f_\Lambda(j)|/s_j
             \le A R^p e^{-a\operatorname{dist}(j,\mathcal C)}.
 \label{eq:nonlinear-response-NR30}
\end{equation}
For each $j$, the distance from $y$ to the collar is at most
$d(y,j)+\operatorname{dist}(j,\mathcal C)$.  Multiplying the summand
by $e^{m'\operatorname{dist}(y,\mathcal C)}$ leaves a weighted row
sum and a remaining exponential at most one.  This proves the claim.

A cutoff must be made in physical coordinates.  For a smooth function
$\chi(x/R)$ on a mesh $\eta$, its first difference divided by $\eta$
is $O(R^{-1})$ and its second difference divided by $\eta^2$ is
$O(R^{-2})$, by the integral form of the mean-value formula.
The unnormalized difference itself has an extra factor $\eta$.
These are different quantities.  A cutoff of an affine potential of
size $O(R)$ can therefore have bounded first derivatives and
polynomially bounded forcing, rather than a bare-mesh singularity.

For a nonlinear equation a cutoff of a solution is not a solution even
where the cutoff is a constant strictly between zero and one.  All
such points belong to the transition collar.  On that collar the
force is computed from the actual nonlinear formula or its integral
Taylor formula.  The admissibility of the cutoff field, its averages,
and its tower backgrounds is a separate domain check.  A bounded
formal commutator cannot replace it.

For two nonlinear backgrounds, subtracting their equations gives an
averaged Jacobian along a common admissible coordinate segment.  The
inverse estimate for that averaged operator must be justified, for
example by a uniform comparison of every Jacobian with the same
reference as in \eqref{eq:nonlinear-response-NR16}.  An arbitrary average of invertible complex
matrices need not be invertible.  The first-step background section
uses just such a common comparison; a multilevel application needs
it on the actual multilevel ball.

\subsubsection*{The full substitution and the test for staying inside a core}

An old-origin source is evaluated on data of the form
\begin{equation}
 V_j(\nu)=\mathcal M^j
        \big(e^{i\eta(t+\nu\zeta)H}U\big)|_{\Gamma_j}.
 \label{eq:nonlinear-response-NR31}
\end{equation}
The parameters on the fine bonds in $\zeta$ may vary independently.
The inner analytic argument uses $|\nu|<c/s$, not merely the physical
segment between zero and one.  There are also outer and removed
substitutions.  Each must satisfy the same list of field, independent
current, node, gauge and tower conditions that defines its assigned
core.  A proof only of the curvature clause is incomplete.

Here is a precise useful criterion.  Express the core as inequalities
$\|\mathcal R_a(X)\|<r_a$ for its actual clauses, including all
hereditary tower clauses.  Let $X_0$ be the complete reference pair and
assume margins $\|\mathcal R_a(X_0)\|\le(1-\epsilon_a)r_a$.
If a common analytic curve satisfies
\begin{equation}
 \sup_{\nu\ {\rm on\ the\ retained\ disc}}
 \|\mathcal R_a(X_\nu)-\mathcal R_a(X_0)\|
                         <\epsilon_a r_a
                         \quad\hbox{for every }a,
 \label{eq:nonlinear-response-NR32}
\end{equation}
then it stays in the core.  The proof is the triangle inequality for
every defining clause.  Its substantive content in an application is
the uniform bound for the infinite hereditary list, not that triangle
inequality.

If the actual response and average derivatives imply
\begin{equation}
 {\|\mathcal R_a(X_\nu)-\mathcal R_a(X_0)\|\over r_a}
      \le C_a|\nu|\,\delta_k\sigma_j/\alpha_j,
 \label{eq:nonlinear-response-NR33}
\end{equation}
then the inner disc is admitted whenever
$c\sup_a C_a\delta_k\sigma_j/(s\alpha_j)<\inf_a\epsilon_a$.
This calculation shows the exact estimate required; it does not
establish \eqref{eq:nonlinear-response-NR33} for the actual maps.  The ratio $\sigma_j$ is the
old-to-current physical length ratio and must be inserted with the
units of the relevant average.  It cannot be guessed by identifying
raw logarithms with potentials.

One part of this check is explicit on the fixed real-base
tower class. At a positive bare bond \(b\) let \(l(b)\)
be its layer spacing in the ambient sequence. Let \(n_b\)
be the maximum of the potential and fixed real-base
covariant first-difference seminorms scaled by \(l(b)\)
and \(l(b)^2\), respectively. For a shift
\(\mathbf A\) of \(\mathbf U=e^{i\eta A'}U\), set
\[
 u(b)=\max\{l(b)|\mathbf A(b)|,\,
                 l(b)^2|\nabla_{\mathbf U}\mathbf A|(b)\},
 \qquad
 A'_{\mathrm{new}}(b)
 =\frac1{i\eta}\log\bigl(e^{i\eta\mathbf A(b)}e^{i\eta A'(b)}\bigr).
\]
Here the first differences are those inside the domain of
the tower. Write \(u^+(b)\) for the maximum of \(u\) at
\(b\) and its available translates by one positive bare
step. The chart estimate gives
\(n_b(A'_{\mathrm{new}}-A')\le2u^+(b)\) when
\(u(b)\le (256L^2)^{-1}\) throughout the chart and the base
potential belongs to the admitted tower class. The
neighboring entries occur because a first difference of the
product of the two exponentials uses both endpoints.
Thus if \(n_b(A')+2u^+(b)<a_{\mathrm T}\), the new potential
lies in the same fixed-base class with radius \(a_{\mathrm T}\).
The tower response estimate then gives the increment at
an element whose tower index agrees with the ambient
index capped at the tower height, in the form
\(2e^\delta C\sup_b e^{-\delta d^\flat(x,x_b)}u(b)\)
at its base point \(x\), using the weighted distance
\(d^\flat\) of the ambient sequence. The factor \(e^\delta\)
absorbs the one bare step in \(u^+\).
For the inner contour put \(\mathbf A=\nu\zeta H\) and
check these bounds on the full retained disc, using the
actual fine-bond differences of \(\zeta H\).
This verifies the fixed real-base potential clause when
there is the stated margin. The induced tower current
is part of that response. The clauses for independently
varied currents and outer and removed substitutions still
require their own bounds.

There are two logically different uses of a displacement estimate.
A bound for $X_\nu-X_0$ on a disc already known to be admitted proves
a Taylor estimate.  To prove admission, one needs the bound while
continuing up to a possible first exit.  On a finite geometry this is
done on the connected part of the analytic solution branch before
exit.  If \eqref{eq:nonlinear-response-NR32} is strict up to its closure and the primitive inverse
has a uniform continuation radius there, no first exit occurs.
The uniform continuation radius and the strict margin are both
needed.  Without them the first-exit argument is circular.

At young ages the factor $\sigma_j$ may be too large to make \eqref{eq:nonlinear-response-NR33}
small.  The permitted fallback is a direct estimate at the finitely
many such ages, using their actual local disc and radius.  It must
bound the whole reference value as well as the increment.  An
increment that vanishes at $\nu=0$ does not prove that $X_0$ itself
satisfies its core inequalities.

For removed substitutions the current and the background may change
together.  The path used in the comparison must then remain in the
joint core.  Separate coordinate Cauchy estimates at the two endpoints
do not imply that the segment connecting them is admitted, especially
when a hereditary condition is not convex.  One can instead use the
actual common analytic chart and the strict clause margins in \eqref{eq:nonlinear-response-NR32}.

\subsubsection*{The remaining quantitative tasks}

The preceding arguments reduce the multilevel response problem to
the following specific estimates on the primitive domain.  First,
the actual constraint rows and gauge slice must admit $E,P$ with
uniform weighted bounds, including every outward seam row.  Second,
the constrained inverse and its energetic compositions must satisfy
\eqref{eq:nonlinear-response-NR14}--\eqref{eq:nonlinear-response-NR17} on the whole retained complex ball.  Third, the datum
forcing and current readout must satisfy their local weighted bounds.
Fourth, the translated equations and boundary rows must give the
required first-difference and H\"older estimates.  Fifth, all inner,
outer and removed substitutions must satisfy \eqref{eq:nonlinear-response-NR32}, with a supremum
over the actual tower list and independent fine parameters.

These tasks have different mechanisms and should be checked
separately.  Proving the first-step versions does not prove their
constants uniform in depth.  Once these estimates hold, the field
and multiplier response follows from the block inverse, joining
distance estimates follow from Cauchy on the admitted ball, and
window-background comparison follows from the exact boundary-force
equation.  These last implications are the proved part of the
section and may be used without adding a new response assumption.


% END MODULE chapters/nonlinear_response_reductions.tex

% BEGIN MODULE chapters/discrete_holder_estimates.tex
\subsection{A mesh-uniform interior H\"older estimate}
\label{ch:discrete-holder-estimate}

The response system needs spatial regularity as well as decay.  We prove
the particular scalar estimate that converts a bounded potential and a
bounded componentwise Laplacian into a H\"older bound for its first
differences.  The estimate is uniform in the mesh.  It also applies in a
parallel frame for an unweighted flat unitary connection, including
arbitrary flat holonomies on a compatible torus.  The nonlinear gauge
application is then reduced to explicit coefficient estimates.  The
construction of those local gauges and the layered equations supplying
the Laplacian bound remain model-dependent tasks.

Fix $0<\alpha<1$.  On $\eta\mathbb Z^4$ define
\begin{equation}
 D_i^+u(x)={u(x+\eta e_i)-u(x)\over\eta},\quad
 D_i^-u(x)={u(x)-u(x-\eta e_i)\over\eta},\quad
 \Delta_\eta=\sum_{i=1}^4D_i^-D_i^+.
 \label{eq:discrete-holder-estimates-DH1}
\end{equation}
The Laplacian has nonpositive Fourier symbol.  The counting measure in
all convolution kernels is $\eta^4$ times counting measure.  We use the
Euclidean distance, allowing dimension-dependent constants when a
coordinate path is used.  A derivative of order $r$ means an arbitrary
product of $r$ of the forward or backward differences in \eqref{eq:discrete-holder-estimates-DH1}.

\subsubsection*{Heat kernels with enough derivative and tail control}

Start on the unit lattice.  For $T\ge0$ set
\begin{equation}
 q_T(n)={1\over(2\pi)^4}\int_{[-\pi,\pi]^4}
 e^{-T\omega(\theta)}e^{in\cdot\theta}\,d\theta,
 \qquad\omega(\theta)=4\sum_i\sin^2(\theta_i/2).
 \label{eq:discrete-holder-estimates-DH2}
\end{equation}
This is a nonnegative probability kernel.  Indeed, if $\mathsf A$ is
the sum of the eight nearest-neighbour translations, then
$\Delta_1=\mathsf A-8I$ and
\begin{equation}
 e^{T\Delta_1}=e^{-8T}\sum_{k=0}^{\infty}{T^k\mathsf A^k\over k!}.
 \label{eq:discrete-holder-estimates-DH3}
\end{equation}
Every coefficient is nonnegative, and applying the operator to the
constant function gives one.  Fourier transformation gives \eqref{eq:discrete-holder-estimates-DH2}, so
the two constructions agree.  The physical heat-kernel density is
\begin{equation}
 p_t^\eta(\eta n)=\eta^{-4}q_{t/\eta^2}(n).
 \label{eq:discrete-holder-estimates-DH4}
\end{equation}
It follows directly that convolution with $p_t^\eta$ is a contraction
on bounded functions.

The following polynomial tail estimates suffice.

\begin{lemma}[Uniform differentiated heat-kernel tails]
For integers $r,N\ge0$ there is $C_{r,N}$ independent of $\eta,t,a$
such that
\par\noindent
\begin{minipage}{\linewidth}
\refstepcounter{equation}
\begin{align*}
 \sum_{|x|\ge a}\eta^4|D^r p_t^\eta(x)|
 &\le C_{r,N}t^{-r/2}(1+a/\sqrt t)^{-N},
                                      &&t\ge\eta^2,\\
 \sum_{|x|\ge a}\eta^4|D^r p_t^\eta(x)|
 &\le C_{r,N}\eta^{-r}(1+a/\eta)^{-N},
                                      &&0\le t\le\eta^2.
 \tag{\theequation}\label{eq:discrete-holder-estimates-DH5}
\end{align*}
\popQED\par
\end{minipage}
\end{lemma}

\begin{proof}
A product of $r$ unit differences has Fourier multiplier $m_r(\theta)$,
a product of factors $e^{\pm i\theta_i}-1$.  For every multi-index $b$,
\begin{equation}
 |\partial^b m_r(\theta)|
           \le C_{r,b}|\theta|^{(r-|b|)_+}.
 \label{eq:discrete-holder-estimates-DH6}
\end{equation}
The bound follows by assigning the derivatives to factors.  A factor
that has not been differentiated vanishes to first order at zero;
all its positive-order derivatives are bounded.  If more than $r$
derivatives occur the asserted power is zero.

For $\theta\in[-\pi,\pi]^4$, $\omega(\theta)\ge c|\theta|^2$.
Its first derivatives are $O(|\theta|)$ and its higher derivatives
are bounded.  Differentiating $e^{-T\omega}$ for $T\ge1$ gives
\begin{equation}
 |\partial^b e^{-T\omega(\theta)}|
 \le C_b T^{|b|/2}
           (1+|\sqrt T\theta|)^{M_b}e^{-c_bT|\theta|^2}.
 \label{eq:discrete-holder-estimates-DH7}
\end{equation}
For completeness, each term of the repeated product rule consists of
$e^{-T\omega}$ times factors $T\partial^{b_j}\omega$, with
$b_j\ne0$ and $\sum b_j=b$.  A first derivative contributes
$\sqrt T\,O(|\sqrt T\theta|)$.  A factor of order at least two
is at most a constant times $T^{|b_j|/2}$ because $T\ge1$.
Multiplying proves \eqref{eq:discrete-holder-estimates-DH7}; the finite polynomial can be absorbed by
a smaller Gaussian exponent when integrating.

Leibniz' rule, \eqref{eq:discrete-holder-estimates-DH6}, and \eqref{eq:discrete-holder-estimates-DH7} now imply
\begin{equation}
 \int_{[-\pi,\pi]^4}
  |\partial^b(m_r e^{-T\omega})|^2\,d\theta
                    \le C_{r,b}T^{|b|-r-2},\qquad T\ge1.
 \label{eq:discrete-holder-estimates-DH8}
\end{equation}
When at most $r$ derivatives hit the multiplier, scale
$\theta=T^{-1/2}v$ in that term.  When more than $r$ hit it, its
bound is smaller than the displayed power since $T\ge1$.  The
four-dimensional volume factor is $T^{-2}$.  This explains every
power in \eqref{eq:discrete-holder-estimates-DH8}.

Let $\delta^r$ denote the corresponding unit differences.  Parseval's
identity and differentiation of the periodic Fourier multiplier give,
for every integer $M\ge0$,
\begin{equation}
 \sum_n (1+|n|^2/T)^M|\delta^r q_T(n)|^2
                    \le C_{r,M}T^{-r-2},\qquad T\ge1.
 \label{eq:discrete-holder-estimates-DH9}
\end{equation}
Expand the weight into coordinate monomials.  The term $|n^b|^2T^{-|b|}$
is bounded by \eqref{eq:discrete-holder-estimates-DH8} times $T^{-|b|}$.  All multiplier functions are
smooth and periodic, so there are no boundary terms at $\pm\pi$.

Choose $M>N+2$, increasing it by one if necessary.  Comparing the
lattice sum with the radial integral in four dimensions gives
\begin{equation}
 \sum_{|n|\ge A}(1+|n|^2/T)^{-M}
       \le C_{M,N}T^2(1+A/\sqrt T)^{-2N},\qquad T\ge1.
 \label{eq:discrete-holder-estimates-DH10}
\end{equation}
One way to check the comparison is to cover the lattice by unit cubes.
For $T\ge1$ the values of $1+|x|/\sqrt T$ at any two points of a cube
are comparable by a fixed dimensional factor.  The integral of
$(1+|x|^2/T)^{-M}$ outside a ball of radius comparable with $A$ is
$T^2$ times the corresponding integral in scaled coordinates.  That
integral decays at least as $(1+A/\sqrt T)^{-2N}$ for this choice of $M$.
Small values of $A$ are absorbed by the full integral bound.

Cauchy--Schwarz with \eqref{eq:discrete-holder-estimates-DH9} and \eqref{eq:discrete-holder-estimates-DH10} proves
$\sum_{|n|\ge A}|\delta^rq_T(n)|
\le C T^{-r/2}(1+A/\sqrt T)^{-N}$.
For $0\le T\le1$, all derivatives of
$m_re^{-T\omega}$ of any fixed order are uniformly bounded.
The same Parseval argument, with weight $(1+|n|^2)^M$, gives instead
$C(1+A)^{-N}$.  Finally \eqref{eq:discrete-holder-estimates-DH4} contributes $\eta^{-r}$ to the
weighted $\ell^1$ norm.  Substitution of $T=t/\eta^2$ proves \eqref{eq:discrete-holder-estimates-DH5}.
\end{proof}

These constants depend on the fixed derivative order and the desired
tail power.  They do not depend on the number of lattice sites.
The small-time part of \eqref{eq:discrete-holder-estimates-DH5} is important: replacing it by a
continuum heat-kernel formula all the way to $t=0$ would not justify
the uniform estimates for the discrete second and third differences.

\subsubsection*{A massive resolvent and its first-difference modulus}

For $\mu>0$ define
\begin{equation}
 G_\mu^\eta(x)=\int_0^\infty e^{-\mu^2t}p_t^\eta(x)\,dt,
       \qquad R_\mu f=G_\mu^\eta*f.
 \label{eq:discrete-holder-estimates-DH11}
\end{equation}
Its kernel is nonnegative and has total mass $\mu^{-2}$.  For fixed
$\eta$, $\Delta_\eta$ is a bounded operator on bounded functions.
Integrating the derivative of
$e^{-\mu^2t}e^{t\Delta_\eta}$ in operator norm shows
\begin{equation}
 (-\Delta_\eta+\mu^2)R_\mu f=f,
       \qquad \|R_\mu f\|_\infty\le\mu^{-2}\|f\|_\infty.
 \label{eq:discrete-holder-estimates-DH12}
\end{equation}
Conversely, any bounded solution of this equation equals $R_\mu f$.
Indeed integrate the same identity applied to the bounded solution;
the terminal term tends to zero because of $e^{-\mu^2t}$ and the
heat contraction.  This also proves uniqueness without a maximum
attainment argument on an infinite lattice.

Assume $\mu\eta\le1$.  Integrating \eqref{eq:discrete-holder-estimates-DH5} for $r=1$, first on
$[0,\eta^2]$ and then on $[\eta^2,\infty)$, gives
\begin{equation}
 \sum_x\eta^4|DG_\mu^\eta(x)|\le C\mu^{-1}.
 \label{eq:discrete-holder-estimates-DH13}
\end{equation}
The first interval contributes $C\eta$ and the second is bounded
by $C\int_0^\infty e^{-\mu^2t}t^{-1/2}\,dt=C'\mu^{-1}$.

For a lattice displacement $h$, telescoping along a coordinate path
and using \eqref{eq:discrete-holder-estimates-DH5} with $r=1,2$ gives, for $t\ge\eta^2$,
\begin{equation}
 \|\tau_h Dp_t^\eta-Dp_t^\eta\|_1
       \le C\min(t^{-1/2},|h|t^{-1}).
 \label{eq:discrete-holder-estimates-DH14}
\end{equation}
The second bound follows by summing the second difference over the
path of physical length at most $2|h|$ in four dimensions; the first
is twice the first-difference norm.  For $t\le\eta^2$ use
$C\eta^{-1}$ instead.  Every nonzero lattice displacement has
$|h|\ge\eta$.

Suppose first that $\eta\le|h|\le\mu^{-1}$.  Integrate \eqref{eq:discrete-holder-estimates-DH14} on
$[\eta^2,|h|^2]$, on $[|h|^2,\mu^{-2}]$ and on the remaining
tail.  Together with the small-time interval, this gives
\begin{equation}
 \|\tau_h DG_\mu^\eta-DG_\mu^\eta\|_1
       \le C|h|\big(1+\log(1/(\mu|h|))\big).
 \label{eq:discrete-holder-estimates-DH15}
\end{equation}
On the middle interval the integral is
$|h|\int_{|h|^2}^{\mu^{-2}}t^{-1}dt$.
The first interval contributes $O(|h|)$, and the exponentially
damped last interval does as well.  If $|h|>\mu^{-1}$ use twice
\eqref{eq:discrete-holder-estimates-DH13}.  Since for $0<x\le1$
\begin{equation}
 x^{1-\alpha}(1+\log(1/x))
        \le 1+{1\over e(1-\alpha)},
 \label{eq:discrete-holder-estimates-DH16}
\end{equation}
both ranges imply
\begin{equation}
 \|\tau_h DG_\mu^\eta-DG_\mu^\eta\|_1
          \le C_\alpha\mu^{\alpha-1}|h|^\alpha.
 \label{eq:discrete-holder-estimates-DH17}
\end{equation}
The value at $h=0$ is zero.  It follows for bounded $f$ that
\begin{equation}
 \|DR_\mu f\|_\infty\le C\mu^{-1}\|f\|_\infty,
 \qquad [DR_\mu f]_\alpha
                \le C_\alpha\mu^{\alpha-1}\|f\|_\infty.
 \label{eq:discrete-holder-estimates-DH18}
\end{equation}
The seminorm compares lattice points, using their physical distance.
The logarithmic modulus \eqref{eq:discrete-holder-estimates-DH15} explains why these estimates do not
give a mesh-uniform Lipschitz bound for $DR_\mu f$ from bounded $f$
alone.  The constants are permitted to diverge as $\alpha$ tends to one.

\subsubsection*{Off-diagonal resolvent derivatives}

Let $R\ge\eta$ and $\mu=R^{-1}$.  For every fixed $c>0$ and
$r=0,1,2,3$,
\begin{equation}
 \sum_{|x|\ge cR}\eta^4|D^rG_{1/R}^\eta(x)|
                           \le C_{r,c}R^{2-r}.
 \label{eq:discrete-holder-estimates-DH19}
\end{equation}
Choose $N>\max(0,r-2)$.  On the time interval $[\eta^2,R^2]$,
the large-time part of \eqref{eq:discrete-holder-estimates-DH5} is bounded by
$Ct^{-r/2}(\sqrt t/R)^N$.  Its integral is $O(R^{2-r})$.
On $[R^2,\infty)$ the exponential $e^{-t/R^2}$ gives the same
bound by scaling.  The interval $[0,\eta^2]$ contributes at most
$C\eta^{2-r}(\eta/R)^N$, which is also $O(R^{2-r})$.
This proves \eqref{eq:discrete-holder-estimates-DH19}, including the second and third derivatives for
which the uncut global norms would have a worse small-scale behaviour.

There is no assertion here that the second resolvent derivative has
a uniformly bounded global absolute sum.  Its off-diagonal bound is
what makes the interior cutoff argument work.

\subsubsection*{An interior estimate with no unknown cutoff gradient}

Let $Q_R$ be the cube of side $2R$ centred at a lattice point.  The
following formulation uses a fixed ratio of cubes; changing that ratio
changes only the constants.

\begin{proposition}[Interior first-difference H\"older bound]
Let $R\ge64\eta$.  Suppose $u$ is defined on $Q_{2R}$ and
$\Delta_\eta u=f$ at every point needed by the equation there.
Put $M=\|u\|_{\infty,Q_{2R}}$ and
$F=\|f\|_{\infty,Q_{2R}}$.  On $Q_{R/2}$, with all first differences
and comparisons taken inside its one-mesh enlargement,
\refstepcounter{equation}
\begin{align*}
 \|Du\|_\infty&\le C(R^{-1}M+RF),\\
 [Du]_\alpha&\le C_\alpha
                    (R^{-1-\alpha}M+R^{1-\alpha}F).
 \tag{\theequation}\label{eq:discrete-holder-estimates-DH20}
\end{align*}
The constants are independent of $R,\eta,u$.
\end{proposition}

\begin{proof}
Choose the restriction of a smooth cutoff $\chi(x/R)$, equal to one
on $Q_R$ and supported in $Q_{3R/2}$.  It satisfies
$|D\chi|\le C/R$ and $|\Delta_\eta\chi|\le C/R^2$.
Its difference supports lie in a collar at distance at least $cR$
from the evaluation cube and its one-mesh enlargement.  The condition
$R\ge64\eta$ makes this separation uniform after every difference
used below.  Extend $u$ arbitrarily outside $Q_{2R}$; the value of
$v=\chi u$ and all forces below are independent of that extension.

With $\mu=1/R$ the exact product identity is
\begin{equation}
 (-\Delta_\eta+\mu^2)v
      =a-[\Delta_\eta,\chi]u,
       \qquad a=\chi(-f+\mu^2u).
 \label{eq:discrete-holder-estimates-DH21}
\end{equation}
Both sides are extended by their actual compactly supported formulas.
The equation for $u$ is used only where $\chi$ or its differences are
nonzero, inside the given domain.  By \eqref{eq:discrete-holder-estimates-DH12},
$v=R_\mu a-R_\mu[\Delta_\eta,\chi]u$.
The first term has force norm at most $F+R^{-2}M$, so \eqref{eq:discrete-holder-estimates-DH18}
gives exactly its contributions to \eqref{eq:discrete-holder-estimates-DH20}.

To treat the commutator one must avoid estimating the unknown
$Du$ in the cutoff collar.  Move the commutator onto the kernel.
For scalar multiplication by $\chi$ and the constant cell-weighted
pairing,
\begin{equation}
 [\Delta_\eta,\chi]^{\mathsf T}
       =[\chi,\Delta_\eta]=-[\Delta_\eta,\chi].
 \label{eq:discrete-holder-estimates-DH22}
\end{equation}
For $j=1,2$, apply this identity to
$K_j(y)=D_x^jG_\mu^\eta(x-y)$ in the convolution.  The discrete
product rule shows that its absolute value is bounded by
\begin{equation}
 C M\left(R^{-2}\|1_{|\cdot|\ge cR}D^jG_\mu^\eta\|_1
       +R^{-1}\|1_{|\cdot|\ge cR}D^{j+1}G_\mu^\eta\|_1\right).
 \label{eq:discrete-holder-estimates-DH23}
\end{equation}
To check the rule explicitly,
$\Delta(\chi K)-\chi\Delta K-K\Delta\chi$
is the sum over $i$ of
$(D_i^+\chi)(D_i^+K)+(D_i^-\chi)(D_i^-K)$.
All factors containing a difference of $\chi$ are supported in the
separated collar.  Changes of the evaluation point by at most the
fixed number of mesh steps in a derivative merely change $c$.
The kernel is absolutely summable on that collar by \eqref{eq:discrete-holder-estimates-DH19}, which
also justifies the summation by parts.

For $j=1$, \eqref{eq:discrete-holder-estimates-DH19} bounds in \eqref{eq:discrete-holder-estimates-DH23} by $CM/R$.
For $j=2$ it bounds it by $CM/R^2$.  Let
$w=R_\mu[\Delta_\eta,\chi]u$.  A coordinate path between two points
of $Q_{R/2}$ stays in that cube; telescoping the first differences
of $w$ gives
$|Dw(x)-Dw(x')|\le CM R^{-2}|x-x'|$.
Its supremum bound is $CM/R$.  Since $|x-x'|\le C R$ in the cube,
this gives $[Dw]_\alpha\le C_\alpha M R^{-1-\alpha}$.
Adding the contribution of $R_\mu a$ proves both estimates.
On the evaluation cube $v=u$, including all the differences used,
so the estimates apply to $u$ itself.
\end{proof}

The proof has not used a previously known bound on $Du$ in the
outer collar.  Estimating the original cutoff commutator directly
by $R^{-1}\|Du\|$ would have introduced that extra premise.  Moving
the commutator onto the separated resolvent kernel is the step
that removes it.

The result holds componentwise for vectors in any fixed
finite-dimensional normed space.  Equivalent coordinate norms
contribute a constant depending on that fixed dimension.  No number
of lattice sites appears.

\subsubsection*{Tori and flat unitary holonomies}

On a torus with periods $\ell_i$, the scalar heat kernel and massive
resolvent are periodizations of their infinite-lattice kernels.
Their absolute derivative sums are bounded by the corresponding
infinite sums.  For a set at periodic distance at least $cR$ from
the origin, every lift has Euclidean distance at least $cR$.
Thus \eqref{eq:discrete-holder-estimates-DH19} holds for periodic distance as well.  If the cubes and
their collars in the interior theorem embed without wrapping, the
same cutoff proof gives \eqref{eq:discrete-holder-estimates-DH20} with the same constants.  For example
$\min_i\ell_i\ge 8R$ suffices for the stated cubes.

A flat unitary connection can be made trivial on the lifted lattice.
The deck transformations act by its commuting unitary holonomies.
The kernel on the torus is the image sum with the corresponding
unitary factors.  On a common parallel frame in the embedded cube,
differences act on the lifted kernel.  Taking the norm and using
unitarity bounds the image sum by the scalar sum already estimated.
Hence all the preceding local bounds remain uniform in the
holonomies.  Their distance from the identity is irrelevant.

The argument requires flatness.  Curved parallel transports have
additional commutators.  Nor does it give the same estimate for an
arbitrary nonlocal or layered positive operator: the proof used the
componentwise constant-coefficient Laplacian in \eqref{eq:discrete-holder-estimates-DH1}.

\subsubsection*{The local gauge application}

Here is how \eqref{eq:discrete-holder-estimates-DH20} enters the background and response estimates.
Suppose a local real gauge has supplied a potential on a cube
enlarged by a fixed multiple of its local scale $\xi$, with
\begin{equation}
 \|A\|_\infty\le C_0m\varepsilon\xi^{-1},\qquad
 \|\Delta_\eta A\|_\infty
                       \le C_1m\varepsilon\xi^{-3}.
 \label{eq:discrete-holder-estimates-DH24}
\end{equation}
Apply the theorem on cubes of radius a fixed fraction of $\xi$,
covering the interior.  It proves
\begin{equation}
 [DA]_\alpha\le C_\alpha m\varepsilon\xi^{-2-\alpha},
       \qquad \|DA\|_\infty\le C m\varepsilon\xi^{-2}.
 \label{eq:discrete-holder-estimates-DH25}
\end{equation}
For pairs farther apart than a fixed fraction of $\xi$, the
supremum bound gives the same H\"older estimate.  Thus local cubes
give the full indicated interior seminorm.  On the enlarged real cubes of the layered profile,
\hyperlink{mb:real-cube-gauge}{the single-gauge construction}
supplies both premises of \eqref{eq:discrete-holder-estimates-DH24}
from curvature and full-current bounds.

There is a finite coarse-mesh case in this application.  If the
chosen fraction of $\xi$ is less than $64\eta$, the ratio
$\xi/\eta$ is bounded by a fixed constant.  Directly,
$|DA|\le2\|A\|_\infty/\eta$ and
$[DA]_\alpha\le4\|A\|_\infty\eta^{-1-\alpha}$.
They give \eqref{eq:discrete-holder-estimates-DH25} with the same powers of $\xi$ and a larger
fixed constant.  The same fallback applies to the response
application below.  Thus no unmentioned lower bound on
$\xi/\eta$ is needed.

For a response, the available equation often controls the positive
covariant componentwise Laplacian
$\mathcal L_{U_0}=\sum_i\nabla_{U_0,i}^{\mathsf T}\nabla_{U_0,i}$.
Thus $\mathcal L_{\mathbf 1}=-\Delta_\eta$ with the heat-generator
sign in \eqref{eq:discrete-holder-estimates-DH1}.  Its conversion can also be made quantitative.
In a local gauge write $U_0(b)=e^{i\eta A_0(b)}$
and let $T_b=\operatorname{Ad}(U_0(b))$ in the chosen unitary colour
representation.  If $\eta|A_0|$ stays in a fixed small exponential
chart, then
\begin{equation}
 \|(T_b-I)/\eta\|\le C|A_0(b)|.
 \label{eq:discrete-holder-estimates-DH26}
\end{equation}
Differencing $T_b$ in a transverse direction uses the integral
derivative of the exponential.  It gives
$|D T_b|\le C\eta|D A_0|$, with the adjacent bond values interpreted
in the same local frame.  The inverse transport satisfies the same
bound.  Expanding the forward and backward covariant differences
therefore gives
\begin{equation}
 |(\mathcal L_{U_0}+\Delta_\eta)h|
 \le C\big(|A_0|\,|Dh|
                     +( |DA_0|+|A_0|^2)|h|\big).
 \label{eq:discrete-holder-estimates-DH27}
\end{equation}
One can verify this by first writing each covariant first difference
as the ordinary difference plus $(T_b-I)/\eta$ times a neighbouring
value.  In the second difference the apparent $\eta^{-1}$ terms
combine into a first difference of $h$, a difference of the
coefficient $(T_b-I)/\eta$, and a product of two such coefficients.
These are exactly the three terms on the right of \eqref{eq:discrete-holder-estimates-DH27}.
There is no estimate of an individual uncombined $\eta^{-1}A_0h$
term.

Assume on a local-scale cube that
\refstepcounter{equation}
\begin{align*}
 \xi|A_0|+\xi^2|DA_0|&\le a_0,\\
 \xi|h|+\xi^2|\nabla_{U_0}h|
                  +\xi^3|\mathcal L_{U_0}h|&\le B.
 \tag{\theequation}\label{eq:discrete-holder-estimates-DH28}
\end{align*}
The formula in \eqref{eq:discrete-holder-estimates-DH26} first gives
$\xi^2|Dh|\le C(1+a_0)B$.
Then \eqref{eq:discrete-holder-estimates-DH27} implies
\begin{equation}
 \xi^3|\Delta_\eta h|
                   \le C(1+a_0+a_0^2)B.
 \label{eq:discrete-holder-estimates-DH29}
\end{equation}
Applying \eqref{eq:discrete-holder-estimates-DH20} at scale $\xi$ yields
$[Dh]_\alpha\le C_{\alpha,a_0}B\xi^{-2-\alpha}$.
The difference between $Dh$ and $\nabla_{U_0}h$ is a bounded
coefficient times a neighbouring $h$.  Its H\"older seminorm is
controlled by the same local bounds, using the Lipschitz coefficient
bound from $DA_0$ and the first-difference bound for $h$.
Thus
\begin{equation}
 [\nabla_{U_0}h]_\alpha
                         \le C_{\alpha,a_0}B\xi^{-2-\alpha}.
 \label{eq:discrete-holder-estimates-DH30}
\end{equation}

This initially compares gradient components in the chosen local
gauge.  If the assigned H\"older seminorm transports one component
to the other along paths of length at most $C|x-x'|$, unitary
product telescoping and \eqref{eq:discrete-holder-estimates-DH28} give
$\|\operatorname{Ad}U_0(\Gamma)-I\|
\le C a_0|x-x'|/\xi$.
Multiplying by $|\nabla_{U_0}h|\le B/\xi^2$ adds at most
$C a_0B\xi^{-2-\alpha}|x-x'|^\alpha$ on the scale-$\xi$
cube.  Hence \eqref{eq:discrete-holder-estimates-DH30} also holds in that transported seminorm.
Its premise still includes the pointwise covariant-gradient
bound in \eqref{eq:discrete-holder-estimates-DH28}; the argument removes an additional H\"older
premise, rather than that gradient bound.

If the right-hand side of \eqref{eq:discrete-holder-estimates-DH28} has the response factor
$B=C e^{-m d(y,y_c)}|e|$, that factor persists on a fixed-size
local-scale cube up to the fixed exponential comparison of distances
within the cube.  This gives the corresponding H\"older member of
the response estimate.  The primitive model input is now explicit:
the actual local gauge and the three bounds in \eqref{eq:discrete-holder-estimates-DH28}, with the
proper current and boundary readouts.  No separate H\"older estimate
is needed once those inputs have been proved.  At a seam the cube,
frame and difference equation must be compatible on both sides;
the scalar theorem cannot supply that compatibility.


% END MODULE chapters/discrete_holder_estimates.tex

% BEGIN PROOF EXPANSION background/local_coulomb_gauge.tex
\subsection{Local gauges from curvature and the Wilson current}
\label{ch:bg-local-coulomb}

The background estimates require a local potential and its first
differences in physical units. A bound on plaquettes gives a small
potential in a comb gauge, but does not directly bound its first
differences. We construct a Coulomb gauge in a finite $L^p$ space and
then use the actual Wilson current to obtain the additional derivative.
The distinction between these two steps is essential: no boundedness
of an order-zero singular integral on $L^\infty$ is used.

The problem of choosing regular local representatives is part of
Ba{\l}aban's gauge-fixing construction \cite{B85b}. The proof here
specifies a different local gauge and normalizes the current directly
by differentiating the Wilson action. The representative need not
preserve block averages or a previously imposed contour gauge. Those
additional requirements are discussed at the end.

\subsubsection{The statement and its physical normalization}

Let $\eta>0$ be a lattice spacing. For $R>0$ with $R/\eta$ an integer,
write
\[
 Q_R=\{x\in\eta\mathbb Z^4:|x|_\infty\le R\}.
\]
An assertion on bonds or plaquettes of a cube concerns those having
all their vertices in the cube. A connection assigns $U_\mu(x)\in
\operatorname{SU}(2)$ to the positive bond from $x$ to $x+\eta e_\mu$;
the reverse bond has the inverse value. Gauge transformations act by
\begin{equation}
 U_\mu^g(x)=g(x)U_\mu(x)g(x+\eta e_\mu)^{-1}.
 \label{bg:gauge-action}
\end{equation}
Lie-algebra elements are traceless Hermitian matrices. We use
$\langle X,Y\rangle=\tfrac12\operatorname{tr}(XY)$ and its Euclidean
norm. All estimates below remain valid for the equivalent operator
norm, with a changed absolute constant.

For distinct $\mu,\nu$, put
\begin{equation}
 \begin{split}
 P_{\mu\nu}(x)&=U_\mu(x)U_\nu(x+\eta e_\mu)\\
 &\qquad{}\times U_\mu(x+\eta e_\nu)^{-1}U_\nu(x)^{-1},\\
 F_{\mu\nu}(x)&=\frac{1}{i\eta^2}\log P_{\mu\nu}(x),\\
 W_{\mu\nu}(x)&=\frac{P_{\mu\nu}(x)-P_{\mu\nu}(x)^*}{2i\eta^2}.
 \end{split}
 \label{bg:curvature-and-sine}
\end{equation}
The logarithm is the branch near the identity. On this branch
$F_{\nu\mu}=-F_{\mu\nu}$, and $W_{\mu\nu}$ is the matrix sine
$\eta^{-2}\sin(\eta^2F_{\mu\nu})$.

The Wilson current is the physical gradient of
$\sum_{p}(1-\tfrac12\operatorname{tr}U_p)$. At every bond whose full
plaquette stencil is present, it is normalized by
\begin{equation}
 \left.\frac{d}{dt}A(e^{i\eta t a}U)\right|_{t=0}
       =\eta^4\sum_{x,\mu}\langle J_\mu(x),a_\mu(x)\rangle.
 \label{bg:current-normalization}
\end{equation}
This is a local definition: compactly supported variations change
only finitely many plaquettes. No finite total action on an infinite
lattice is assumed.

For a matrix-valued function let
\begin{equation}
 D_\mu^+f(x)=\frac{f(x+\eta e_\mu)-f(x)}{\eta},\quad
 D_\mu^-f(x)=\frac{f(x)-f(x-\eta e_\mu)}{\eta},\quad
 L_\eta=-\sum_\mu D_\mu^-D_\mu^+.
 \label{bg:ordinary-differences}
\end{equation}
The Coulomb condition on a potential is
$\operatorname{div}_\eta A:=\sum_\mu D_\mu^-A_\mu=0$.
The notation $DA$ includes every forward first difference of every
component. For a function on a cube, $[f]_{\alpha,Q}$ is the ordinary
discrete H\"older seminorm, with physical Euclidean distance.

\begin{proposition}[A local gauge controlled by the physical current]
\label{prop:bg-local-gauge}
There are absolute constants $\varepsilon_*>0$ and $C<\infty$ with
the following property. Let $0<\eta\le r$, $r/\eta\in\mathbb N$, and
let a real connection be given on $Q_{8r}$. Assume that all its
plaquettes lie in the logarithm branch of
\eqref{bg:curvature-and-sine}, and that
\begin{equation}
 \|F\|_{\infty,Q_{8r}}\le\kappa,\qquad
                    r^2\kappa\le\varepsilon_*.
 \label{bg:local-small-curvature}
\end{equation}
Suppose the actual current in \eqref{bg:current-normalization}
satisfies $\|J\|_{\infty,Q_{7r}}\le j$.
Then there is a real site gauge $g$ such that $U^g=e^{i\eta A}$
on $Q_{6r}$, in the logarithm branch near the identity, and
$\operatorname{div}_\eta A=0$ wherever its stencil lies in
$Q_{6r}$. On $Q_r$,
\begin{equation}
 \begin{split}
 \|A\|_\infty&\le Cr\kappa,\\
 \|DA\|_\infty&\le C(\kappa+rj),\\
 \|L_\eta A\|_\infty&\le C(j+r\kappa^2).
 \end{split}
 \label{bg:local-gauge-bounds}
\end{equation}
For each fixed $0<\alpha<1$,
\begin{equation}
 [DA]_{\alpha,Q_r}\le
                   C_\alpha r^{-\alpha}(\kappa+rj).
 \label{bg:local-holder-bound}
\end{equation}
The constants are independent of $\eta$, $r$, the connection and its
holonomies outside the cube. No smallness condition on $j$ is needed.
The gauge and the bound on $A$ use only the curvature assumption;
the current enters the derivative estimates.
\end{proposition}

The condition can equivalently be imposed using
$\eta^{-2}\|P-1\|_\infty$: on a fixed small logarithm branch it is
comparable with $\|F\|_\infty$. In particular the hypothesis has the
physical curvature scaling of the background domains, rather than
a bound on an unscaled bond potential.

We prove the proposition after the analytic estimates below. Under
$x=r\widetilde x$, the scaled quantities are
\begin{equation}
 \widetilde\eta=\eta/r,\quad
 \widetilde A=rA,\quad
 \widetilde F=r^2F,\quad
 \widetilde J=r^3J.
 \label{bg:scaling}
\end{equation}
Thus it suffices to prove the statement at $r=1$.

\subsubsection{Uniform finite-$L^p$ estimates on a lattice torus}

The auxiliary torus in this proof has physical side $32$. Its spacing
is $h=1/N$, $N\in\mathbb N$, and integration means $h^4$ times
counting measure. All constants in this subsection may depend on
$p$ and on that fixed side, but not on $N$. Let $\Pi$ be the
arithmetic mean and let $G_h=L_h^{-1}(1-\Pi)$.

\begin{lemma}[Uniform Green and Riesz estimates]
\label{lem:bg-finite-p}
For $1<p<\infty$ and any choices of forward or backward differences,
\begin{equation}
 \|D_iD_jG_h f\|_p\le C_p\|f\|_p.
 \label{bg:riesz-lp}
\end{equation}
For $p>4$,
\begin{equation}
 \|DG_h f\|_\infty\le C_p\|f\|_p,
 \qquad
 \|u-\Pi u\|_\infty\le C_p\|Du\|_p.
 \label{bg:green-sobolev}
\end{equation}
For a one-form $a$,
\begin{equation}
 \|Da\|_p\le C_p
       \bigl(\|\mathrm{d}a\|_p+\|\operatorname{div}_h a\|_p\bigr),
 \quad
 (\mathrm{d}a)_{\mu\nu}=D_\mu^+a_\nu-D_\nu^+a_\mu.
 \label{bg:hodge-lp}
\end{equation}
These estimates hold componentwise for the fixed colour space.
\end{lemma}

\begin{proof}
We give the kernel argument, including the mesh uniformity at the
diagonal. The Fourier symbol of $L_h$ is
\begin{equation}
 \lambda_h(k)=4h^{-2}\sum_{\mu=1}^4\sin^2(hk_\mu/2),
 \qquad k\in(2\pi/32)\mathbb Z^4
                       \cap[-\pi/h,\pi/h)^4.
 \label{bg:laplacian-symbol}
\end{equation}
It is comparable with $|k|^2$ on the Brillouin cube. Fourier
orthogonality proves the $L^2$ estimate in
\eqref{bg:riesz-lp}; each first-difference symbol has modulus at most
$\lambda_h^{1/2}$.

Let $p_t^h$ be the physical heat kernel. Direct Fourier integration
gives, for a product of $s$ differences,
\begin{equation}
 \|D^sp_t^h\|_\infty\le
          C_s(t+h^2)^{-(4+s)/2},\qquad 0<t\le1.
 \label{bg:heat-linfty}
\end{equation}
For $t\le h^2$, use the volume of the Brillouin cube and
$|D|\le2/h$. For $t\ge h^2$, put $q=k\sqrt t$ and use
$\lambda_h\ge c|k|^2$. The corresponding physical $L^1$ tails are
\begin{equation}
 \sum_{|x|\ge a}h^4|D^sp_t^h(x)|\le
 C_{s,M}(t+h^2)^{-s/2}
        (1+a/\sqrt{t+h^2})^{-M}.
 \label{bg:heat-tails}
\end{equation}
Here and below the distance is periodic. One may obtain this estimate
from the differentiated heat-kernel tail calculation in
\cref{ch:discrete-holder-estimate}. For clarity, its Fourier argument
uses integration by parts in a coordinate with
$|x_\mu|\ge|x|/2$ when $t\ge h^2$. Each differentiation of the
Gaussian multiplier is bounded by a polynomial in $k\sqrt t$ times
$e^{-c t|k|^2}$, and supplies one factor $\sqrt t$.
At $t\le h^2$, the random-walk series
$e^{t\Delta_h}=e^{-8t/h^2}\sum_{n\ge0}(t/h^2)^n\mathcal A^n/n!$
has uniformly bounded moments of every order; a product of $s$
differences costs $C_sh^{-s}$. Periodizing the infinite-lattice
estimate proves the displayed torus version.

The semigroup law turns these two bounds into a pointwise bound.
In $D^sp_t=(D^sp_{t/2})*p_{t/2}$, either the first or the second
variable has distance at least $|x|/2$. Use the $L^1$ tail for that
factor and the $L^\infty$ bound for the other. Consequently
\begin{equation}
 |D^sp_t^h(x)|\le C_{s,M}(t+h^2)^{-(4+s)/2}
                 (1+|x|/\sqrt{t+h^2})^{-M}.
 \label{bg:heat-pointwise}
\end{equation}
For $t\ge1$, subtract the constant mode. The nonzero frequencies
have a fixed positive lower spectral bound, and the same Fourier
estimate gives $C_se^{-ct}$ for every fixed derivative order.

Integrate the heat kernel in time. With $\rho(x)$ the periodic
distance from zero, the kernels with physical summation convention
satisfy
\begin{equation}
 \begin{split}
 |DG_h(x)|&\le C(\rho(x)+h)^{-3}+C,\\
 |D_iD_jG_h(x)|&\le C(\rho(x)+h)^{-4}+C,\\
 |D_\ell D_iD_jG_h(x)|&\le
                      C(\rho(x)+h)^{-5}+C.
 \end{split}
 \label{bg:green-kernel-bounds}
\end{equation}
For example split the integral for a second derivative at
$t=(\rho+h)^2$, use $M>6$ below that point and the power $t^{-3}$
above it. The interval $0<t<h^2$ contributes at most $Ch^{-4}$ at
the diagonal. The first and third derivative estimates follow with
powers $t^{-5/2}$ and $t^{-7/2}$. Thus a diagonal matrix entry is
large, but its physical mass $h^4$ is retained.

We now prove the $L^p$ consequence for $T=D_iD_jG_h$.
Its physical kernel has sum zero, its $L^2$ norm is at most one,
and \eqref{bg:green-kernel-bounds} implies
\begin{equation}
 \sum_{\rho(x-y)>4|y-z|}h^4
       |K_h(x-y)-K_h(x-z)|\le C
 \quad(y\ne z).
 \label{bg:kernel-cancellation-bound}
\end{equation}
Indeed join $y$ to $z$ by a shortest coordinate path. The third
difference estimate, summed along that path, bounds the difference
by $C|y-z|\rho(x-y)^{-5}$ away from the displayed ball. Summing
annuli in four dimensions gives a uniform constant. The bounded
large-distance term is harmless on the fixed torus.

Use the Calder\'on--Zygmund decomposition \cite{CZ52}, which we spell
out for these finite lattices. Split the coordinate intervals of the
torus repeatedly into two integer intervals of lengths differing by
at most one. The resulting rectangles have uniformly bounded aspect
ratio, and each parent has volume at most $3^4$ times that of any
child. At the final generation they are single lattice sites.
For a level $\lambda$ above the mean of $|f|$, take the maximal
rectangles $Q$ with average greater than $\lambda$. They are disjoint,
and
\begin{equation}
 \sum_Q|Q|\le\lambda^{-1}\|f\|_1,\qquad
 f=g+\sum_Qb_Q,\quad
 |g|\le3^4\lambda,\quad
 \|g\|_2^2\le C\lambda\|f\|_1.
 \label{bg:dyadic-decomposition}
\end{equation}
Here $b_Q=(f-f_Q)1_Q$, so it has mean zero and
$\sum_Q\|b_Q\|_1\le2\|f\|_1$.
The enlarged rectangles have total measure at most
$C\lambda^{-1}\|f\|_1$. Outside them,
\eqref{bg:kernel-cancellation-bound}, with the cancellation of each
$b_Q$, gives
$\|\sum_QTb_Q\|_{L^1(\mathrm{outside})}\le C\|f\|_1$.
The $L^2$ estimate controls $Tg$. Chebyshev's inequality proves
\begin{equation}
       |\{|Tf|>\lambda\}|\le C\lambda^{-1}\|f\|_1.
 \label{bg:weak-one}
\end{equation}
If $\lambda$ is below the mean, the same bound is immediate from
the total volume of the torus. This includes the last atomic scale;
no continuum limit has been used.

For completeness the interpolation is uniform as well. Split $f$
at height proportional to $t$ when estimating $|Tf|>t$, use the
weak $L^1$ bound on the large part and the $L^2$ bound on the small
part, and integrate $p t^{p-1}$ in $t$. For $1<p<2$, the two
integrals are bounded by constants times
\[
 \int|f(x)|\int_0^{|f(x)|}t^{p-2}\,dt\,dx,
 \qquad
 \int|f(x)|^2\int_{|f(x)|}^\infty t^{p-3}\,dt\,dx.
\]
They equal $\|f\|_p^p/(p-1)$ and
$\|f\|_p^p/(2-p)$. Duality proves $2<p<\infty$, since the
adjoint has the same kernel estimates. This proves
\eqref{bg:riesz-lp} for all forward/backward choices.

For $p>4$, $p'=p/(p-1)<4/3$. The first kernel in
\eqref{bg:green-kernel-bounds} has uniformly bounded physical
$L^{p'}$ norm, because
$\sum_xh^4(\rho(x)+h)^{-3p'}\le C_p$. H\"older's inequality
therefore proves the first assertion in
\eqref{bg:green-sobolev}. The identity
$u-\Pi u=-\sum_iD_i^-G_hD_i^+u$ proves the second.

Finally ordinary differences commute and
\begin{equation}
 L_ha_\mu=\sum_\nu D_\nu^-(\mathrm{d}a)_{\mu\nu}
                     -D_\mu^+\operatorname{div}_h a.
 \label{bg:ordinary-hodge-identity}
\end{equation}
Apply $D_i^+G_h$ and \eqref{bg:riesz-lp}. Constant one-forms have
zero derivative, so their mean causes no additional term.
This proves \eqref{bg:hodge-lp}.
\end{proof}

The use of $p>4$ in \eqref{bg:green-sobolev} is strict. The
order-zero operator in \eqref{bg:riesz-lp} is never applied with
$p=\infty$. Its kernel has the borderline singularity $|x|^{-4}$,
whereas the kernel used to bound a first derivative from a Laplacian
has the integrable power $|x|^{-3}$ in the required dual space.

\subsubsection{A small-potential Coulomb gauge in a finite $L^p$ space}

\begin{lemma}[A uniform nonlinear gauge calculation]
\label{lem:bg-gauge-remainder}
There are absolute $a_*>0$ and $C<\infty$ such that the following
holds for Hermitian traceless matrices $\lambda,A,w$ whenever
$|\lambda|+h|A|+h|w|<a_*$. Set
\begin{equation}
 \mathcal B_h(\lambda,A,w)=\frac1{ih}
       \log\big(e^{i\lambda}e^{ihA}e^{-i(\lambda+hw)}\big)
       =A-w+N_h(\lambda,A,w).
 \label{bg:gauge-bch}
\end{equation}
Then
\begin{equation}
 |N_h|\le C\bigl(|\lambda|+h(|A|+|w|)\bigr)(|A|+|w|).
 \label{bg:gauge-remainder-pointwise}
\end{equation}
On a fixed smaller domain its first two derivatives satisfy the
corresponding product bounds obtained by differentiating the right
side. All constants are uniform in $0<h\le1$.
\end{lemma}

\begin{proof}
Write $X=hA$, $Y=hw$ and consider
$\mathcal L(\lambda,X,Y)=i^{-1}
\log(e^{i\lambda}e^{iX}e^{-i(\lambda+Y)})$.
This is analytic on a fixed matrix polydisc. It vanishes for
$X=Y=0$, for every admitted $\lambda$, and its differential in
$(X,Y)$ at $(0,0,0)$ is $X-Y$.
Taylor's integral formula in $(X,Y)$ therefore gives
\[
 |\mathcal L-(X-Y)|\le
             C|\lambda|(|X|+|Y|)+C(|X|+|Y|)^2.
\]
Divide by $h$. The same integral formula and the Cauchy estimate
on a fixed larger polydisc prove the derivative bounds. The argument
uses finite matrix analyticity, so its constants do not depend on
the number of lattice sites.
\end{proof}

\begin{proposition}[Coulomb gauge from a small global potential]
\label{prop:bg-small-potential-coulomb}
Fix $p>4$. On the physical side-$32$ torus there are constants
$b_p>0$ and $C_p<\infty$, uniform in $h=1/N$, with this property.
If $V_\mu=e^{ih a_\mu}$ is real and
$\|a\|_\infty\le b_p$, there is a mean-zero real site function
$\lambda$ such that $g=e^{i\lambda}$ and
\begin{equation}
 V^g=e^{ih B},\quad
 \operatorname{div}_hB=0,\quad
 \|D\lambda\|_p+\|\lambda\|_\infty+\|B\|_p
                                      \le C_p\|a\|_\infty.
 \label{bg:coulomb-lp}
\end{equation}
The solution is unique in the ball specified in the proof. This is
a uniqueness assertion for this small gauge parameter, not for all
gauge representatives of the connection.
\end{proposition}

\begin{proof}
At the bond $(x,\mu)$ use \eqref{bg:gauge-bch} with
$w=D_\mu^+\lambda(x)$ and $A=a_\mu(x)$. Thus
\begin{equation}
 B=a-D\lambda+N(\lambda),\qquad
 \lambda=-G_h\operatorname{div}_h(a+N(\lambda))
 \label{bg:coulomb-fixed-point}
\end{equation}
is exactly the Coulomb equation. The second map has mean zero.
Use the Banach norm $\|D\lambda\|_p$ on mean-zero site functions.
The Sobolev estimate in \eqref{bg:green-sobolev} controls
$\|\lambda\|_\infty$ by this norm. The elementary atomic bound is
\begin{equation}
 h\|D\lambda\|_\infty
       \le h^{1-4/p}\|D\lambda\|_p
       \le\|D\lambda\|_p.
 \label{bg:atomic-small-increment}
\end{equation}
It controls the group increment without asserting a uniform bound
on the unscaled $L^\infty$ derivative of $\lambda$.

On the ball $\|D\lambda\|_p\le R$ with
$\|a\|_\infty\le b$ and $b+R$ small,
\cref{lem:bg-gauge-remainder} gives
\begin{equation}
 \begin{split}
 \|N(\lambda)\|_p&\le C_p(b+R)(b+R),\\
 \|N(\lambda)-N(\widetilde\lambda)\|_p
      &\le C_p(b+R)\|D\lambda-D\widetilde\lambda\|_p.
 \end{split}
 \label{bg:gauge-remainder-lp}
\end{equation}
To check the quadratic increment term, use
$h\|w^2\|_p\le h\|w\|_\infty\|w\|_p$ and
\eqref{bg:atomic-small-increment}. For a difference of two such
terms use $w^2-\widetilde w^2$ as two ordered matrix products.
The terms containing $\lambda$ use its uniform Sobolev bound.
These arguments also prove the displayed Lipschitz estimate when
the matrices do not commute.

By \eqref{bg:riesz-lp}, the map on the right of
\eqref{bg:coulomb-fixed-point} is bounded in this norm by
$C_p\|a+N\|_p$. Let $K_p$ dominate this constant, the fixed torus
volume factor and the constants in
\eqref{bg:gauge-remainder-lp}. Take $R=2K_pb$, and then take $b_p$
so small that the domain of \cref{lem:bg-gauge-remainder} contains
the whole ball and
\begin{equation}
 K_pC_p(b_p+2K_pb_p)\le\tfrac14,
 \qquad K_pC_p(b_p+2K_pb_p)^2\le K_pb_p/2.
 \label{bg:coulomb-smallness}
\end{equation}
The map is a contraction and maps the closed radius-$R$ ball into
itself. Its real iterates converge to a real solution. The logarithms
remain on the common branch by the same inequalities. Taking the
divergence of the first identity in
\eqref{bg:coulomb-fixed-point} proves the Coulomb equation exactly.
The estimates for $\lambda$ and $B$ now follow from the ball bound,
Sobolev and \eqref{bg:gauge-remainder-lp}.
\end{proof}

The parameter ball is uniform in the mesh. Its proof used a fixed
finite $p$ throughout. The smallness of $hD\lambda$ in
\eqref{bg:atomic-small-increment} is enough to keep the link logarithm
on its branch; an $L^\infty$ estimate for the order-zero gauge
projection was not needed.

\subsubsection{Comb gauges and an auxiliary periodic extension}

\begin{lemma}[A comb potential and its curvature-preserving extension]
\label{lem:bg-comb-extension}
Suppose $U$ is real on $Q_8$, $h=1/N\le1$, and
$\|F\|_{\infty,Q_8}\le\kappa$ with $\kappa$ sufficiently small.
There is a real gauge $g_0$ on $Q_8$ for which
$U^{g_0}=e^{ih a}$ on $Q_7$ and $\|a\|_\infty\le C\kappa$.
There is also a side-$32$ periodic potential $a^{\rm ext}$ agreeing
with $a$ on $Q_6$, satisfying
\begin{equation}
 \|a^{\rm ext}\|_\infty\le C\kappa,\qquad
 \|F(e^{ih a^{\rm ext}})\|_\infty\le C\kappa.
 \label{bg:comb-extension-bounds}
\end{equation}
The extension does not assert a current bound outside $Q_6$.
\end{lemma}

\begin{proof}
Use the fixed coordinate-order comb from one corner of $Q_8$ to
each site. Choose $g_0$ so its tree transports become one.
The transformed value of a bond is the closed transport made from
the two combs and that bond. Moving the bond past the later coordinate
segments expresses this loop as a product of at most $C/h$
conjugated plaquettes. The estimate
$\|AB-1\|\le\|A-1\|+\|B-1\|$ for unitary matrices therefore gives
\begin{equation}
                   \|U^{g_0}(b)-1\|\le Ch\kappa.
 \label{bg:comb-bound}
\end{equation}
Take the logarithm and divide by $h$. No bound on $g_0-1$ is used.

We record a local analytic identity that is also used later. If
$V_\mu=e^{ih B_\mu}$ and all $h|B_\mu|$ are in a fixed small
ball, then
\begin{equation}
 F_{\mu\nu}(V)=(\mathrm{d}B)_{\mu\nu}+E_{\mu\nu}(B),\qquad
 |E_{\mu\nu}(B)(x)|\le C\max_{b\in\partial p}|B(b)|^2.
 \label{bg:curvature-expansion}
\end{equation}
Indeed apply the matrix logarithm to the four exponential factors
of the plaquette. Its first-order term is
$h[B_\mu(x)+B_\nu(x+he_\mu)-B_\mu(x+he_\nu)-B_\nu(x)]
=h^2(\mathrm{d}B)_{\mu\nu}$. All remaining terms have degree at least two
in these four variables, and division by $h^2$ gives the stated
uniform bound. Its first derivatives with respect to the four
$B$ variables are bounded by $C\max|B|$ on a smaller chart.

It follows that $\|\mathrm{d}a\|_\infty\le C\kappa$ on $Q_7$.
Choose a fixed smooth scalar cutoff $\chi$, equal to one on $Q_6$
and zero outside the interior of $Q_7$, and put
$a^{\rm ext}(b)=\chi(b_-)a(b)$ on the available bonds and zero
elsewhere. Embed its support in the side-$32$ torus.
The discrete product rule gives
$\|d a^{\rm ext}\|_\infty\le C\kappa$.
Apply \eqref{bg:curvature-expansion} once more to the extended
potential. This proves \eqref{bg:comb-extension-bounds}.
At the finest exceptional collar positions the same formula uses
only bonds in $Q_8$, since $h\le1$. Thus the construction has the
claimed domain even when $N=1$.
\end{proof}

The extension is an analytic tool on a fixed auxiliary torus. On
the inner cube it agrees with a gauge transform of the given field.
It imposes no condition on holonomies of any larger ambient torus.
The current of the extension is deliberately not used near its cutoff.

\subsubsection{The exact current identity and its elliptic consequence}

\begin{lemma}[The Wilson current as a covariant difference]
\label{lem:bg-exact-current}
With the normalization in \eqref{bg:current-normalization},
\begin{equation}
 J_\mu(x)=\frac1h\sum_{\nu\ne\mu}
 \left[W_{\mu\nu}(x)
 -\operatorname{Ad}(U_\nu(x-he_\nu)^{-1})
                          W_{\mu\nu}(x-he_\nu)\right].
 \label{bg:exact-wilson-current}
\end{equation}
This identity is exact for real $\operatorname{SU}(2)$ fields on the
logarithm branch. It is not a small-field replacement of the action.
If $U=e^{ihB}$, $\operatorname{div}_hB=0$ and
$M=\|F\|_\infty$, then, wherever the full local stencil is present,
\begin{equation}
 \begin{split}
 L_hB_\mu={}&J_\mu-\sum_{\nu\ne\mu}D_\nu^-E_{\mu\nu}(B)
          -\sum_{\nu\ne\mu}D_\nu^-S_h(F_{\mu\nu})
          -\mathcal R_\mu,\\
 S_h(X)&=h^{-2}\sin(h^2X)-X,\\
 |\mathcal R(x)|&\le C\|B\|_{\infty,\mathrm{stencil}}M,\qquad
 |\operatorname{div}S_h(F)(x)|\le Ch^3M^3.
 \end{split}
 \label{bg:current-laplacian-equation}
\end{equation}
In the same local norms,
\par\noindent
\begin{minipage}{\linewidth}
\begin{equation}
 |DE(B)|\le C\|B\|_\infty\max|DB|.
 \label{bg:curvature-remainder-derivative}
\end{equation}
\popQED\par
\end{minipage}
\end{lemma}

\begin{proof}
Differentiate a plaquette containing the positive bond $(x,\mu)$.
For the plaquette based at $x$ whose first edge is that bond, the
derivative is
$h\langle a_\mu(x),h^2W_{\mu\nu}(x)\rangle$.
For the plaquette based at $x-he_\nu$, the bond occurs reversed.
Move the insertion to the initial point by cyclicity of the trace.
Its contribution is
\[
 -h^3\left\langle a_\mu(x),
 \operatorname{Ad}(U_\nu(x-he_\nu)^{-1})
                    W_{\mu\nu}(x-he_\nu)\right\rangle.
\]
The apparent additional conjugation by the plaquette itself drops
out: $W(P)=(P-P^*)/(2ih^2)$ commutes with $P$.
The same calculation with the inverse orientation handles $\nu<\mu$,
using $W_{\nu\mu}=-W_{\mu\nu}$.
Divide the total derivative by the physical factor $h^4$.
This proves \eqref{bg:exact-wilson-current}, with its sign and
all spacing powers fixed.

Replace its covariant difference by the ordinary backward difference.
The error is bounded by $C\|B\|_\infty M$, since
$\|1-\operatorname{Ad}(e^{-ihB})\|\le Ch|B|$ and
$|W(F)|\le C|F|$ on the branch. Next insert
$F=\mathrm{d}B+E(B)$ and $W(F)=F+S_h(F)$, and use
\eqref{bg:ordinary-hodge-identity} with zero divergence.

The sine series gives $|S_h(X)|\le Ch^4|X|^3$.
An ordinary first difference of any bounded lattice function is at
most twice its supremum divided by $h$. Therefore
\begin{equation}
                  |D^-S_h(F)|\le Ch^3M^3.
 \label{bg:sine-difference}
\end{equation}
No derivative bound for $F$ is hidden in this estimate.
Finally $E(B)(x)$ is the same analytic function of the four
neighbouring $B$ values at every $x$. Subtract its values at two
neighbouring plaquettes and use the integral derivative of that
function, bounded as in \eqref{bg:curvature-expansion}.
This proves \eqref{bg:curvature-remainder-derivative}.
\end{proof}

The factor $h^3$ in \eqref{bg:sine-difference} is useful. A bound
on the physical curvature is much stronger than a bound on an
unscaled plaquette logarithm when the mesh is small. This factor
permits the Wilson current to replace the divergence of logarithmic
curvature without requiring a modulus of continuity for that
curvature.

\subsubsection{Proof of the local gauge bounds}

\begin{proof}[Proof of \cref{prop:bg-local-gauge}]
Scale to $r=1$ as in \eqref{bg:scaling}; write $h=\eta/r$.
Apply \cref{lem:bg-comb-extension}, then
\cref{prop:bg-small-potential-coulomb} with $p=16$ to the auxiliary
connection. Write its transformed potential as $B$.
Choose $\varepsilon_*$ small enough that $C\kappa$ in the extension
is below $b_{16}$ and all logarithms stay in the fixed branches.
On the torus,
\begin{equation}
 \operatorname{div}_h B=0,\qquad
 \|B\|_{16}\le C\kappa,\qquad
 \|F(B)\|_\infty\le C\kappa.
 \label{bg:first-bootstrap-data}
\end{equation}
The last bound follows from gauge covariance of the actual
plaquettes, not from a separate estimate of differentiated gauge
parameters.

The group chart also gives $h\|B\|_\infty\le C\kappa$, using
\eqref{bg:atomic-small-increment}. Hence
\eqref{bg:curvature-expansion} is valid at every torus plaquette.
By \eqref{bg:hodge-lp} with $p=8$,
\begin{equation}
 \|DB\|_8\le C\bigl(\|F\|_8+\|B\|_{16}^2\bigr)
                                             \le C(\kappa+\kappa^2).
 \label{bg:first-finite-p-bootstrap}
\end{equation}
Apply the Sobolev estimate in \eqref{bg:green-sobolev} to each
component of $B$, retaining its mean, which is bounded by
$C\|B\|_{16}$. We obtain
\begin{equation}
                          \|B\|_\infty\le C\kappa.
 \label{bg:potential-linfty}
\end{equation}
This is a finite-$p$ estimate. It does not estimate a Riesz transform
of $F$ in $L^\infty$.
Return to \eqref{bg:hodge-lp}, now with $p=16$ and the bound
$\|E(B)\|_{16}\le C\|B\|_\infty\|B\|_{16}$. Thus
\begin{equation}
                              \|DB\|_{16}\le C\kappa.
 \label{bg:gradient-finite-p}
\end{equation}

On $Q_6$ the auxiliary connection before the last gauge agrees with
$U^{g_0}$. Its current on $Q_5$ is therefore the gauge transform of
the given actual current. Norms of real currents are unchanged by
that unitary conjugation. Use
\cref{lem:bg-exact-current}, \eqref{bg:potential-linfty} and
\eqref{bg:gradient-finite-p}; the full current stencil lies in
the region of agreement. On a fixed cube between $Q_4$ and $Q_5$,
\begin{equation}
                         \|L_hB\|_{16}\le C(j+\kappa+\kappa^2).
 \label{bg:laplacian-finite-p}
\end{equation}
The terms are respectively the actual current, $BDB$, the transport
term $BF$, and the bounded sine correction $h^3F^3$.
Since $h\le1$ and $\kappa\le\varepsilon_*$, the last two are
bounded in the indicated fixed volume. The extension's current
near its cutoff does not enter this estimate.

Let $\psi$ be a fixed scalar cutoff supported in $Q_4$ and equal
to one on $Q_3$. Its sampled first and second differences are
uniformly bounded, including when $h=1$. The discrete product rule,
\eqref{bg:gradient-finite-p}, and
\eqref{bg:laplacian-finite-p} give
\[
                       \|L_h(\psi B)\|_{16}\le C(j+\kappa).
\]
All neighbours read by the product rule are inside $Q_5$.
The constant mode is killed by $D$, so
$D(\psi B)=DG_hL_h(\psi B)$.
The order-minus-one bound in \eqref{bg:green-sobolev} proves
\begin{equation}
                      \|DB\|_{\infty,Q_2}\le C(\kappa+j).
 \label{bg:gradient-linfty}
\end{equation}
The margin between $Q_2$ and $Q_3$ contains every difference stencil
even at $h=1$. This is the step that obtains a bounded first
derivative. It uses the current to control a Laplacian in $L^{16}$,
then an order-minus-one kernel. It is not an endpoint consequence
of $F\in L^\infty$ alone.

Use the exact equation \eqref{bg:current-laplacian-equation} again,
now pointwise on $Q_1$. If $h\le1/1024$, the same estimate holds on
$Q_{3/2}$: every first-difference base point read by
\eqref{bg:curvature-remainder-derivative} then lies in $Q_2$.
Equations \eqref{bg:potential-linfty},
\eqref{bg:gradient-linfty}, and \eqref{bg:sine-difference} give
\begin{equation}
 \|L_hB\|_\infty\le
 C\bigl(j+\kappa(\kappa+j)+h^3\kappa^3\bigr)
                                      \le C(j+\kappa^2).
 \label{bg:laplacian-linfty}
\end{equation}
Decrease $\varepsilon_*$ once to absorb the harmless factor
$1+C\kappa$ into the constant multiplying $j$.

When $h\le1/1024$, apply the mesh-uniform interior H\"older estimate of
\cref{ch:discrete-holder-estimate} to the ordinary componentwise
Laplacian, using \eqref{bg:potential-linfty} and
\eqref{bg:laplacian-linfty} on $Q_{3/2}$.
It yields $[DB]_{\alpha,Q_1}\le C_\alpha(\kappa+j)$.
If its fixed interior radius is less than $64h$, then $h$ has a
fixed positive lower bound. In that finite coarse-mesh range,
$|DB(x)-DB(y)|\le2\|DB\|_\infty$ and $|x-y|\ge h$ prove the
same H\"older bound directly from
\eqref{bg:gradient-linfty}.

The combined gauge is $g=e^{i\lambda}g_0$ on the original cube.
It is real and its transformed links agree with $e^{ihB}$ on the
inner region. Undo \eqref{bg:scaling}. This gives every power of
$r$ in \eqref{bg:local-gauge-bounds} and
\eqref{bg:local-holder-bound}.
\end{proof}

The smallness constant is chosen in a definite order: first fix the
matrix logarithm polydisc and the side-$32$ torus; then the finite
constants in \cref{lem:bg-finite-p} for $p=8,16$; then the contraction
ball in \eqref{bg:coulomb-smallness}; finally choose
$\varepsilon_*$ below the finitely many branch and absorption
ceilings in the proof. None of these choices depends on the mesh,
the value of $j$, or the connection outside the original cube.

\subsubsection{Comparison with the background regularity hypotheses}

\begin{corollary}[Local-scale regularity]
\label{cor:bg-physical-local-scale}
Suppose a cube of radius $8m\xi$ lies inside the action region,
including every plaquette used by the current, and
\begin{equation}
 \|F\|_\infty\le\alpha\xi^{-2},\qquad
 \|J_A(U)\|_\infty\le\alpha\xi^{-3},\qquad
                         m^2\alpha\le\varepsilon_*.
 \label{bg:model-local-input}
\end{equation}
Here $m\ge1$, $m\xi/\eta\in\mathbb N$, and $0<\eta\le\xi$.
There is a local gauge on the inner cube for which
\par\noindent
\begin{minipage}{\linewidth}
\begin{equation}
 \begin{split}
 \xi\|A\|_\infty&\le Cm\alpha,\\
 \xi^2\|DA\|_\infty&\le C(1+m)\alpha,\\
 \xi^3\|L_\eta A\|_\infty&\le C(\alpha+m\alpha^2),\\
 \xi^{2+\beta}[DA]_\beta&\le
             C_\beta m^{-\beta}(1+m)\alpha
                    \quad(0<\beta<1).
 \end{split}
 \label{bg:model-local-output}
\end{equation}
\popQED\par
\end{minipage}
\end{corollary}

\begin{proof}
Use \cref{prop:bg-local-gauge} with $r=m\xi$,
$\kappa=\alpha\xi^{-2}$ and $j=\alpha\xi^{-3}$, then substitute.
The first-order bounds are linear in $m$. The hypothesis
$m^2\alpha\le\varepsilon_*$ is a cost of this proof and is retained.
\end{proof}

The corollary supplies the potential, gradient and Laplacian premises
needed for the local H\"older deduction. It applies to the physical
Wilson current. A separately prescribed auxiliary current satisfying
a numerical bound cannot replace $J_A(U)$ in
\eqref{bg:model-local-input}. At a boundary where some plaquettes of
the bond stencil are absent, the exact equation has different terms;
the statement above does not remove that boundary problem.

The connection may already have a block-axial or a Landau representative
required by its global construction. The local gauge just constructed
is an additional frame in which to estimate that connection. Passing a
non-covariant gauge functional to this frame requires its own exact
transformation law. In particular, the absolute soft contour gauge is
only invariant under block-constant gauges. The present $g$ is not
claimed to be block constant.

\subsubsection{The coupled inverse and its analytic inputs}

For the actual constrained action, the field operator on its fixed
slice is the Lagrangian Hessian
\begin{equation}
 K_{\rm L}=F''(A)-D^2\langle\omega,q(A)\rangle,
 \qquad
 \mathcal S=\begin{pmatrix}K_{\rm L}&-q'(A)^\sharp\\q'(A)&0\end{pmatrix}.
 \label{bg:remaining-coupled-operator}
\end{equation}
Here $q$ is the actual nonlinear straightened block constraint, the
pairing contains its physical row weights, and $\omega$ is the
corresponding row multiplier. A uniform bound for
$\mathcal S^{-1}$ uses the following quantitative ingredients:
\begin{enumerate}
\item a quantitative inverse on the true constrained tangent space,
including the boundary and seam conditions;
\item a controlled right inverse of the actual nonflat field rows,
with their node factors and outward-row correction;
\item a relative bound for the constraint-curvature term paired with
the actual multiplier, and for every specified independent-current
completion;
\item the exact comparison of the chosen contour gauge with the
gauge block used in the precision;
\item control of the nonlocal subtraction that replaces the full
Feynman term $D_UD_U^\sharp$ by $D_URD_U^\sharp$.
\end{enumerate}
The local gauge result proves an input to these estimates. Its
derivative bounds do not by themselves settle the five requirements.
The following single-scale results prove a separate projected-Landau
gap and the corresponding real constrained saddle inverse in the
homogeneous periodic energy spaces; see
\cref{prop:loc-periodic-nonflat-gap,cor:slp-saddle-inverse}.
They retain the actual constraint curvature. The proper strong
response and its coefficient families on a given admitted reference
are constructed in
\cref{prop:scf-strong-current-families,prop:ngc-joint-connected-family}.
The forced trial residual is treated in
\cref{prop:rtc-nonstationary-trial}, and
\cref{prop:cta-curvature-critical-atlas}
produces the required homogeneous trial from coarse curvature.
The proper curvature-data construction is proved in
\cref{bgc:proper-background}, its tower admission in
\cref{htc:all-towers,htc:parameter-core}, and the response on the
full first-order regular-pair domain in \cref{dnr:response}.
The local gauge theorem itself remains an estimate for the
given connection.

% END PROOF EXPANSION background/local_coulomb_gauge.tex

% BEGIN PROOF EXPANSION background/single_scale_landau_precision.tex
\subsection{The Landau precision on one physical block scale}
\label{ch:bg-single-scale-landau}

The full Feynman form controls every component of the divergence.
The background precision retains only its projection by $R$.
Subtracting the omitted component is therefore a substantive step.
We first prove a gap for that projected flat form, uniform as the fine
mesh tends to zero. We then compare the actual nonflat coefficients
with the flat ones in a common small physical gauge. The result
includes the nonlinear constraint Hessian and an inverse on the
actual simultaneous constraint and Landau slice.

There is one physical block scale throughout this section. All coarse
sites have the same assigned level, all rows are present, and the
ambient lattice is periodic. Thus there are no outward rows, layer
seams or proper-region boundary terms. The assumptions below do not
replace those cases by a periodic problem.

\subsubsection{Operators and pairings}

Set the coarse spacing to one and the fine spacing to $h=L^{-k}$,
where $L\ge3$ is a fixed odd integer and $k\ge0$. The periodic side
length is an integer $M\ge2$. Fine site, bond and plaquette norms
use measure $h^4$; coarse site and bond norms use counting measure,
which is their physical volume measure at this scale.
The symbol $\sharp$ denotes the transpose for these pairings, and
is the Hermitian adjoint when taking norms of real coefficients on
their complexification.

At a flat connection choose a parallel frame on the universal cover.
The resulting boundary twists are unitary. Write $D$ for the fine
gradient, $d$ for its curl and $D^\sharp$ for its divergence adjoint.
Let $q$ be the arithmetic average on each centered fine block, and
let $Q$ be the physical straight bond derivative:
\begin{equation}\label{slp:flat-rows}
 (q\lambda)(y)=h^4\sum_{x\in B(y)}\lambda(x),\qquad
 (Qv)_\mu(y)=h^5\sum_{x\in B(y)}\sum_{t=0}^{h^{-1}-1}
                                      v_\mu(x+th e_\mu).
\end{equation}
These are the flat derivatives of the actual site and straightened
field maps. Telescoping proves the exact identity
\begin{equation}\label{slp:flat-intertwining}
                            QD=\mathrm{d}_cq,
\end{equation}
where $\mathrm{d}_c$ is the gradient with unit coarse step. In particular
the site map in this identity is an average, not evaluation at the
block center.

Put
\begin{equation}\label{slp:flat-projector}
 L=D^\sharp D,
 \quad G=(L+q^\sharp q)^{-1},
 \quad F=Gq^\sharp,
 \quad \mathcal M=F^\sharp F,
 \quad R=1-F\mathcal M^{-1}F^\sharp.
\end{equation}
The positive coefficient multiplying $q^\sharp q$ can be changed
without changing the range of $F$, hence without changing $R$.
Indeed, for any $z\in\ker q$, $GLz=z$. Consequently
\begin{equation}\label{slp:projector-range}
              \operatorname{ran}R=L(\ker q),
 \qquad qGR=0,\qquad LGR=R.
\end{equation}
For the first assertion, $F^\sharp Lz=qz=0$. Conversely, if
$F^\sharp f=0$, then $z=Gf$ satisfies $qz=0$ and $Lz=f$.
The other two identities follow from the same calculation.

The omitted projector is not $q^\sharp q$. The factor $G$ in
\eqref{slp:flat-projector} is important. It makes $R$ the projection
onto $L(\ker q)$, rather than onto $\ker q$.

\subsubsection{A gap for the flat projected form}

\begin{proposition}[A flat Landau gap uniform in the fine mesh]
\label{prop:slp-flat-gap}
There is a constant $c_0>0$, independent of $k$, $M$ and the real
flat holonomies, such that
\begin{equation}\label{slp:flat-gap}
 \|\mathrm{d}v\|_2^2+\|RD^\sharp v\|_2^2+\|Qv\|_2^2
              \ge c_0\bigl(\|Dv\|_2^2+\|v\|_2^2\bigr).
\end{equation}
Here $Dv$ denotes all componentwise fine first differences.
The assertion holds for the three real colour components and for
their complexification.
\end{proposition}

\begin{proof}
We include the fine-mesh limiting argument. It avoids substituting
a fixed-$L$ estimate for a bound uniform in $L^k$.

Translation by one coarse unit commutes with all the flat operators.
Bloch decomposition reduces the problem to one coarse block with
quasiperiodic phase $\theta\in[-\pi,\pi]^4$. It is enough to prove
one bound for every such phase. The allowed phases of any finite
periodic torus form a subset. A flat adjoint holonomy merely shifts
these phases after simultaneous diagonalization of its commuting
unitary translations. It therefore suffices to treat one complex
colour component.

On a Bloch block the coarse site space has dimension one.
The site inequality from the block variance proof gives
\begin{equation}\label{slp:bloch-site-gap}
 \|u\|_2^2\le C\bigl(\|Du\|_2^2+|qu|^2\bigr),
\end{equation}
uniformly in $h$ and $\theta$: it uses only internal edges of the
block, so no boundary phase enters its constant. Hence
$\|G\|\le C$, $\|F\|\le C$ and $\|DF\|\le C$.
Also $\|F\|\ge c>0$ uniformly. To verify the last assertion,
choose a fixed smooth function $\varphi$ supported strictly inside
the cell and of integral one. For all sufficiently fine meshes its
sample has $|q\varphi|\ge1/2$ and
$\|(L+q^\sharp q)\varphi\|_2\le C$. Then
\[
 \tfrac12\le |q\varphi|
   =|\langle F,(L+q^\sharp q)\varphi\rangle|
   \le C\|F\|_2.
\]
For the finitely many remaining values of $h=L^{-k}$,
$L+q^\sharp q$ has operator norm at most $16h^{-2}+1$ and
$\|q^\sharp1\|=1$, which gives the same conclusion with a fixed
smaller constant. The argument is uniform in $\theta$.

Write $P=1-R$. Since $P$ has range spanned by $F$, summation by
parts gives
\begin{equation}\label{slp:missing-divergence}
 \|PD^\sharp v\|_2
  =\frac{|\langle F,D^\sharp v\rangle|}{\|F\|_2}
  \le\frac{\|DF\|_2}{\|F\|_2}\|v\|_2
  \le C\|v\|_2.
\end{equation}
The exact flat Hodge identity thus yields the uniform G\aa rding
inequality
\begin{equation}\label{slp:garding}
 \|Dv\|_2^2\le
     \|\mathrm{d}v\|_2^2+\|RD^\sharp v\|_2^2+C\|v\|_2^2.
\end{equation}
There is no fine-spacing loss in
\eqref{slp:missing-divergence}; the derivative falls on the smoothed
column $F$, not on an arbitrary bounded site function.

Suppose there were no uniform mass gap. Choose Bloch data with
$\|v_n\|_2=1$ and with the left side of
\eqref{slp:flat-gap} tending to zero. Pass to a subsequence with
$\theta_n\to\theta$. If the corresponding integers $h_n^{-1}$
remain bounded, all the matrices vary continuously in a fixed
finite-dimensional space. A zero-energy limit would then contradict
the kernel calculation below, applied on that finite lattice.
We may therefore assume $h_n\to0$.

Inequality \eqref{slp:garding} bounds every first difference.
Here is a useful precise compactness realization. Interpolate a
Bloch grid function by its trigonometric polynomial with frequencies
$\theta_n+2\pi\ell$ in the Brillouin cube. Parseval gives the same
$L^2$ norm, and
$4h_n^{-2}\sin^2(h_n k_\mu/2)\ge c|k_\mu|^2$ gives a uniform
continuum $H^1$ bound. After multiplication by $e^{-i\theta_n x}$
these are periodic $H^1$ functions on a fixed cell. Rellich
compactness gives strong $L^2$ convergence of a subsequence.
The piecewise constant interpolation has the same limit. Indeed
discrete orthogonality gives, for $|z|\le Ch_n$,
\[
 h_n^4\sum_x|u^{\rm trig}(x+z)-u^{\rm trig}(x)|^2
        \le |z|^2\|\nabla u^{\rm trig}\|_2^2.
\]
Integrating $z$ over one fine cell shows that its $L^2$ difference
from the piecewise constant interpolation is at most
$Ch_n\|Du\|_2$. The limiting field $v$ consequently has norm one
and lies in the quasiperiodic $H^1$ space.

Ordinary differences converge weakly to the continuum differences,
as is seen first on a fixed finite set of Fourier modes and then by
their uniform $L^2$ bound. Thus $\mathrm{d}v=0$.
The site averages converge to integration over the cell. The straight
rows converge to
\begin{equation}\label{slp:continuum-row}
 (Q_\theta v)_\mu=
       \int_B\int_0^1v_\mu(x+t e_\mu)\,dt\,dx.
\end{equation}
To check this convergence, rewrite the discrete row as
$h_n\sum_t qT_{th_n}v_n$ on piecewise constant fields.
Replacing $th_n$ by $t$ changes the result by at most
$Ch_n\|Dv_n\|_2$, using the translation estimate just proved.
Strong $L^2$ convergence handles the remaining integral.
It follows that $Q_\theta v=0$.

The smoothed columns $F_n$ converge strongly as well. The equation
$(L_n+q_n^\sharp q_n)F_n=q_n^\sharp1$, its uniform energy bound,
and the same compactness give a strong $L^2$ limit $F_\theta$.
Testing against sampled smooth Bloch functions identifies its equation
\begin{equation}\label{slp:continuum-site-equation}
 \langle \nabla\phi,\nabla F_\theta\rangle
 +\overline{q_\theta\phi}\,q_\theta F_\theta
                         =\overline{q_\theta\phi}.
\end{equation}
The continuum block Poincar\'e inequality proves uniqueness.
The uniform lower bound on $\|F_n\|$ passes to the limit.
Therefore their rank-one orthogonal projections converge in operator
norm under these interpolations. Since $D^\sharp v_n$ is weakly
bounded and $R_nD^\sharp v_n\to0$, we obtain
$R_\theta\nabla^\sharp v=0$.
The continuum projector satisfies the same range identity as
\eqref{slp:projector-range}; its proof uses the continuum resolvent
and is unchanged.

We finish by identifying the zero kernel, both in the continuum and
on each fixed mesh. Fourier decomposition of a curl-free Bloch field
gives $v=D\lambda$ unless $\theta=0$, in which case
$v=D\lambda+c$ with $c$ a constant vector. In the continuum $D$
means $\nabla$ and $\lambda\in H^2$; on a fixed mesh this is the
same finite Fourier calculation. For each nonzero frequency the
curl condition says that its vector coefficient is a scalar multiple
of the gradient symbol. The only zero gradient symbol is the constant
mode at $\theta=0$.

At $\theta=0$ the row of a constant vector is that vector, while
$QD\lambda=\mathrm{d}_cq\lambda=0$. Hence $Qv=0$ forces $c=0$.
At $\theta\ne0$, some component of $\mathrm{d}_c$ is nonzero, so
$QD\lambda=0$ forces $q\lambda=0$. At $\theta=0$ we obtain the
same normalization by subtracting a scalar constant from $\lambda$.
Now $D^\sharp v=L\lambda$ belongs to $L(\ker q)$, and
\eqref{slp:projector-range} gives
\[
          0=RD^\sharp v=L\lambda,\qquad D\lambda=0.
\]
Thus $v=0$, contradicting its unit norm. This proves a uniform
mass gap. Combine that gap with \eqref{slp:garding} to obtain
the full estimate \eqref{slp:flat-gap}.
\end{proof}

The compactness proof determines a positive constant for the fixed
four-dimensional cell geometry. It gives no claimed numerical
optimal value. Its quantifiers include all $k$ and every Bloch
phase, hence every compatible periodic volume and flat holonomy.

\subsubsection{Actual nonflat coefficients in a common gauge}

For a real connection $U$, let $D_U$ be its covariant site gradient.
Let $q_U$ be the actual hierarchical site average on each assigned
pillar and let $Q_U$ be the derivative of the actual straightened
nonlinear field constraint $\mathcal Q_U$. These are the operators
whose coefficients were obtained from the finite recursion in
\cref{ch:layered-current-precision}. In particular
\begin{equation}\label{slp:actual-intertwining}
               Q_UD_U=d_{\mathcal M^kU}q_U.
\end{equation}
The derivative of the plain nonlinear average has a different site
readout and is not substituted here.
Define
\begin{equation}\label{slp:actual-projector}
 G_U=(D_U^\sharp D_U+q_U^\sharp q_U)^{-1},\quad
 F_U=G_Uq_U^\sharp,\quad
 R_U=1-F_U(F_U^\sharp F_U)^{-1}F_U^\sharp.
\end{equation}
These are the actual site resolvent and Landau projector, with the
physical pairings already declared.

\begin{lemma}[Comparison of the actual rows and projector]
\label{lem:slp-coefficients}
Let $V$ be real and flat, and suppose in a common parallel frame
\begin{equation}\label{slp:common-gauge-smallness}
                  U_b=e^{ih a_b}V_b,\qquad\|a\|_\infty\le\delta.
\end{equation}
For a fixed sufficiently small $\delta_L>0$, the actual regular
mean branches are admitted. Uniformly in $k$ and $M$,
\begin{equation}\label{slp:coefficient-comparison}
 \begin{split}
 \|D_U-D_V\|_{2\to2}&\le C\delta,\\
 \|q_U-q_V\|_{2\to2}+\|Q_U-Q_V\|_{2\to2}&\le C_L\delta,\\
 \|R_U-R_V\|_{2\to2}&\le C_L\delta,\\
 \|(R_UD_U^\sharp-R_VD_V^\sharp)v\|_2
                         &\le C_L\delta(\|D_Vv\|_2+\|v\|_2).
 \end{split}
\end{equation}
There is also a right inverse $E_U$ of $Q_U$ with
\par\noindent
\begin{minipage}{\linewidth}
\begin{equation}\label{slp:field-lift-bounds}
 \|E_Ub\|_2+\|\nabla_UE_Ub\|_2\le C_L\|b\|_2,
 \qquad
 \sup_c\sum_{b\ {\rm fine}}h^4|E_U(b,c)|\le C_L.
\end{equation}
\popQED\par
\end{minipage}
\end{lemma}

\begin{proof}
The first comparison follows from
$\|\operatorname{Ad}(e^{ih a})-1\|\le Ch\delta$.
The actual internal path transports over physical length at most
one differ from those at $V$ by $C_L\delta$ in this frame.
For hierarchical transports, the perturbation at intermediate
physical scale $\xi_s$ is at most $C_L\delta\xi_s$, and
$\sum_{s<k}\xi_s\le1/(L-1)$. The site-row comparison follows
by averaging their transported coefficients.

For the field rows use the actual finite analytic recursion and its
row and column estimates in
\cref{ch:layered-current-precision}. Its one-step relative error at
scale $s$ is $C_L\delta\xi_s$; their geometric sum is bounded by
$C_L\delta$. This bounds both Schur sums for the declared physical
measures by that amount; the fine column density retains its factor
$h^4$. Schur's test
gives the second line of \eqref{slp:coefficient-comparison}.
This is the row comparison for $Q_U$, including the differentiated
straightening factors. It is not an infinity-norm estimate for a
transported straight main term in place of $Q_U$.

The flat site gap and the first two comparisons make
$D_U^\sharp D_U+q_U^\sharp q_U$ coercive on the physical site
$H^1$ space, uniformly in volume. Its form differs from the flat
form by at most $C_L\delta$ times the two $H^1$ norms.
The resolvent identity in the $H^1$--$H^{-1}$ pairings gives
\[
 \|G_U-G_V\|_{L^2\to H^1}\le C_L\delta,
 \qquad \|F_U-F_V\|_{2\to H^1}\le C_L\delta.
\]
The flat Gram lower bound used in the Bloch proof, or equivalently
the normalized site Gram estimate already proved for the actual
pillars, implies $F_V^\sharp F_V\ge cI$. Hence its inverse and the
inverse at $U$ are uniformly bounded, and their difference is
$O_L(\delta)$. Substitute these facts in
\eqref{slp:actual-projector}. This proves the projector comparison.
Its last stated line now follows by inserting and subtracting
$R_UD_V^\sharp v$. It loses no factor $h^{-1}$ because the
remaining divergence is measured by the physical first-difference
norm of $v$.

For the right inverse, use the compact signed two-moment row seeds
of \cref{ch:layered-current-precision}. At one assigned scale their
sum $E^0$ satisfies $Q_VE^0=1$ exactly. Every column is supported
inside a physical block, has bounded amplitude and bounded first
differences, and has physical $L^1$ norm at most a fixed constant.
Their supports and gradient collars have bounded overlap.
The field-row comparison therefore gives
\[
 S_U=Q_UE^0=1+B_U,\qquad
      \|B_U\|_{\ell^1\to\ell^1}
       +\|B_U\|_{\ell^\infty\to\ell^\infty}\le C_L\delta.
\]
Reduce $\delta_L$ so each norm is less than $1/2$.
Then $E_U=E^0S_U^{-1}$ is an exact right inverse.
The Neumann series is bounded in both coarse Schur norms. The seed
energy and column $L^1$ bounds prove \eqref{slp:field-lift-bounds}.
The finitely many small block sizes, including a bare identity row,
have the explicit finite seed inverses given in that construction.
All constants are therefore independent of $k$.
\end{proof}

\subsubsection{The physical constraint curvature and the full form}

For a row multiplier $\omega$, pair the Hessian of the actual
per-length nonlinear constraint with the physical coarse measure:
\begin{equation}\label{slp:constraint-curvature}
 B_\omega(v,w)=\sum_c
             \langle\omega(c),D^2\mathcal Q_{U,c}(0)[v,w]\rangle.
\end{equation}
The two-leg coefficient moment of the actual recursion gives
\begin{equation}\label{slp:constraint-bound}
 |B_\omega(v,w)|\le C_L\|\omega\|_\infty\|v\|_2\|w\|_2.
\end{equation}
This is \eqref{px:quadratic-current-form-bound} at physical scale
one. No Gaussian sample is required to lie in a holomorphic chart:
the estimate is a bound on the resulting finite bilinear coefficient
matrix.

Let $a$ now denote a fixed positive row weight, with
$0<a_*\le a(c)\le a^*<\infty$. It is unrelated to the potential
$a_b$ in \eqref{slp:common-gauge-smallness}. Define the real form
\begin{equation}\label{slp:full-form}
 \begin{split}
 \mathcal K_{U,\omega}(v,w)={}&D^2A_W(U)[v,w]-B_\omega(v,w)\\
 &+\langle R_UD_U^\sharp v,R_UD_U^\sharp w\rangle
                       +\langle Q_Uv,aQ_Uw\rangle.
 \end{split}
\end{equation}
The action derivatives use physical left variations
$e^{ihtv}U$. Thus the action Hessian has its physical $L^2$
normalization; no unscaled bond Hessian is inserted in this formula.

\begin{proposition}[Uniform nonflat precision in a common small gauge]
\label{prop:slp-nonflat-gap}
Fix $L,a_*,a^*$. There are $\varepsilon_L>0$ and $c_L,C_L>0$
independent of $k,M$ and the flat holonomies such that the following
holds. Suppose \eqref{slp:common-gauge-smallness} holds relative to
a real flat connection $V$,
\begin{equation}\label{slp:nonflat-inputs}
 \|F(U)\|_\infty\le\kappa,
 \qquad\|\omega\|_\infty\le\gamma,
 \qquad\delta+\kappa+\gamma\le\varepsilon_L.
\end{equation}
Then, with
$\|v\|_X^2=\|\nabla_Uv\|_2^2+\|v\|_2^2$,
\begin{equation}\label{slp:nonflat-gap}
 c_L\|v\|_X^2\le\mathcal K_{U,\omega}(v,v)
                       \le C_L\|v\|_X^2.
\end{equation}
Consequently its full penalized inverse has
$L^2\to L^2$ norm at most $c_L^{-1}$ and
$X^*\to X$ norm at most $c_L^{-1}$.

If $U$ is a critical point of the actual constrained action and
its full physical stationarity equation is
$J_A(U)=Q_U^\sharp\omega$, then
\begin{equation}\label{slp:multiplier-from-current}
                      \|\omega\|_\infty\le C_L\|J_A(U)\|_\infty.
\end{equation}
Thus a sufficiently small bound on this actual current supplies
the multiplier hypothesis. Criticality and that full stationarity
identity are hypotheses, not conclusions of the coercivity estimate.
\end{proposition}

\begin{proof}
Apply \cref{prop:slp-flat-gap}, with the positive lower row weight,
to the flat reference. It controls both the flat component gradient
and the physical mass norm. The two gradient norms at $U$ and $V$
are equivalent for small $\delta$ by
\cref{lem:slp-coefficients}.

For the Wilson term use the exact plaquette Hessian and its curl
comparison from \cref{ch:layered-current-precision}. In the common
frame its linearized plaquette curl differs from the flat curl by
at most $C\delta\|v\|_2$. Its remainder is bounded by
$C\kappa\|v\|_2\|w\|_2$, since the two insertions contribute
$h^2$, the plaquette defect contributes $\kappa h^2$, and their
product is the physical measure $h^4$. Thus
\begin{equation}\label{slp:wilson-comparison}
 \left|D^2A_W(U)[v,w]-\langle \mathrm{d}_Vv,\mathrm{d}_Vw\rangle\right|
 \le C(\delta+\kappa)\|v\|_{X_V}\|w\|_{X_V}.
\end{equation}
The gauge and row comparisons in
\eqref{slp:coefficient-comparison}, followed by expansion of the
two squares, bound their form differences by
$C_L\delta\|v\|_{X_V}\|w\|_{X_V}$.
Equation \eqref{slp:constraint-bound} bounds the actual constraint
curvature by $C_L\gamma$ times these norms. Choose
$\varepsilon_L$ so the total relative error is below one half of
the flat lower constant. This proves both inequalities in
\eqref{slp:nonflat-gap}. The inverse estimates are the ordinary
finite-dimensional energy estimates in the stated physical norms.

Finally apply the exact right inverse in
\cref{lem:slp-coefficients} to the full stationarity equation:
$\omega=E_U^\sharp J_A(U)$. The last bound in
\eqref{slp:field-lift-bounds}, with fine measure $h^4$ and unit
coarse measure, proves \eqref{slp:multiplier-from-current}.
There is no factor $h^{-4}$ in that readout.
\end{proof}

This proves positivity of a form with the actual projected gauge.
It does not infer that positivity by monotonicity from the Feynman
form. The flat projected gap was proved separately, before the
nonflat perturbation was made.

\subsubsection{The actual constraint and Landau slice}

The positive extension in \eqref{slp:full-form} gives an inverse
for the actual constrained Hessian as well. Let
\begin{equation}\label{slp:actual-spaces}
 X_L=\{v:R_UD_U^\sharp v=0\},\qquad
 X_0=X_L\cap\ker Q_U,
 \qquad K_L=D^2A_W(U)-B_\omega.
\end{equation}
Give both subspaces the covariant $X$ norm. On $X_0$ the two
penalty terms vanish exactly, so
$K_L(v,v)\ge c_L\|v\|_X^2$.

\begin{corollary}[A uniform saddle inverse on the actual slice]
\label{cor:slp-saddle-inverse}
Under the assumptions of \cref{prop:slp-nonflat-gap}, the equations
\begin{equation}\label{slp:saddle}
 K_Lu-Q_U^\sharp\mu=f\quad\hbox{on }X_L,
 \qquad Q_Uu=b
\end{equation}
have a unique solution $u\in X_L$, $\mu$ in the actual coarse
field-row space, and
\begin{equation}\label{slp:saddle-bound}
                  \|u\|_X+\|\mu\|_2
                        \le C_L(\|f\|_{X_L^*}+\|b\|_2).
\end{equation}
The constant is independent of fine spacing and compatible volume.
\end{corollary}

\begin{proof}
First construct a bounded right inverse taking values in $X_L$.
The exact site range identities in
\eqref{slp:projector-range} also hold at $U$. Thus
\begin{equation}\label{slp:gauge-restoration}
 S_U=1-D_UG_UR_UD_U^\sharp,
 \quad S_UX\subset X_L,
 \quad Q_US_U=Q_U.
\end{equation}
The last identity uses
\eqref{slp:actual-intertwining} and $q_UG_UR_U=0$.

We verify the energy bound, rather than treating $D_U^\sharp$ as
a bounded operator on bare $L^2$. For $\lambda=G_UR_UD_U^\sharp v$,
the range identity gives
$q_U\lambda=0$ and $D_U^\sharp D_U\lambda=R_UD_U^\sharp v$.
The site gap gives
$\|\lambda\|_2+\|D_U\lambda\|_2\le C_L\|v\|_X$.
The covariant Hodge identity applied to $D_U\lambda$ yields
\[
 \|\nabla_UD_U\lambda\|_2
 \le C_L\bigl(\|D_U^\sharp D_U\lambda\|_2
               +\kappa\|\lambda\|_2
               +\sqrt\kappa\|D_U\lambda\|_2\bigr).
\]
Indeed the covariant curl of $D_U\lambda$ is a shift commutator
bounded by $C\kappa\|\lambda\|_2$, and the Hodge remainder is
$C\kappa\|D_U\lambda\|_2^2$. These are precisely the complete
periodic plaquette stencils, so there is no boundary term.
This proves $\|S_Uv\|_X\le C_L\|v\|_X$.
Therefore $E_L=S_UE_U$ is a bounded right inverse of $Q_U$ into
$X_L$.

Write $u=E_Lb+z$, $z\in X_0$. The coercive equation
\[
 K_L(z,w)=f(w)-K_L(E_Lb,w)\quad(w\in X_0)
\]
has a unique solution and bounds $\|z\|_X$ by
$C_L(\|f\|_{X_L^*}+\|b\|_2)$.
Define the row vector $\mu$ by
$\langle\mu,b'\rangle=(K_Lu-f)(E_Lb')$ for every row vector $b'$;
abbreviate this as $\mu=E_L^\sharp(K_Lu-f)$.
The functional $K_Lu-f$
annihilates $X_0$, and every $w\in X_L$ decomposes as
$w-E_LQ_Uw+E_LQ_Uw$. Hence it equals $Q_U^\sharp\mu$ on
$X_L$, proving \eqref{slp:saddle}. The boundedness of $K_L$ and
$E_L$ proves the multiplier estimate. The same kernel coercivity
and the surjectivity of $Q_U|_{X_L}$ give uniqueness.
\end{proof}

\subsubsection{Complex perturbations and localization}

On a physical first-order complex ball about the real reference,
the already established coefficient estimates for the actual means,
site resolvent and Gram projector give a form perturbation of size
$C_Ls$, where
\begin{equation}\label{slp:complex-radius}
 \|E\|_\infty+\|\nabla_UE\|_\infty+\|\zeta\|_\infty\le s,
 \qquad U_E=e^{ihE}U,\quad\omega_E=\omega+\zeta.
\end{equation}
The first derivative of $E$ controls the complex plaquette change;
its presence cannot be replaced by a fixed unscaled bond radius.
For $s$ smaller than a fixed fraction of the real lower constant,
the full penalized inverse exists holomorphically by the normalized
Neumann series. This statement uses analytic bilinear transposes
for the continued coefficients; it asserts no Hermitian positivity
at a complex field. A constrained inverse on a moving complex slice
additionally requires a uniformly bounded holomorphic identification
of that slice. Only the real-slice saddle estimate and the full
penalized complex inverse are asserted here.

The common-gauge assumption in
\eqref{slp:common-gauge-smallness} has global scope. The local gauge
construction supplies it on bounded patches. The following
localization argument combines those frames without requiring one
global small gauge; \cref{prop:loc-periodic-nonflat-gap} proves the
real curvature-small homogeneous periodic case.

To patch the estimates, one needs an operator comparison for the
nonlocal $R_U$. For a scalar partition function $\chi$, its gauge
term contains
\begin{equation}\label{slp:localization-commutator}
 R_UD_U^\sharp(\chi v)-\chi R_UD_U^\sharp v
 =R_U[D_U^\sharp,\chi]v+[R_U,\chi]D_U^\sharp v.
\end{equation}
An $L^2$ bound for $[R_U,\chi]$ alone leaves a fine derivative
on $v$. The required composed bound and the sharper quadratic IMS
estimate are proved in \cref{lem:loc-projector-ims} using the
smoothed projector's physical cell operators. The comparison
between the global site inverse and the periodic extension is
proved in \cref{lem:loc-projector-shielding}. Both arguments use
the exact actual site rows and Gram matrix, independently of the
fine vector inverse.

Together these results close a real nonflat periodic one-scale
energy inverse with the actual nonlinear constraint curvature.
Different assigned scales, proper-region boundary forms, a moving
complex slice and the independent-current completion remain
distinct tasks. The equality with a separately specified absolute
contour precision and the stronger pointwise response norms also
require their own estimates.

% END PROOF EXPANSION background/single_scale_landau_precision.tex

% BEGIN PROOF EXPANSION background/single_scale_localization.tex
\subsection{Localization of the one-scale Landau precision}
\label{ch:bg-single-scale-localization}

The common-gauge theorem leaves a nonlocal comparison: a local frame
does not change the global projector into a projector computed on
that patch. We prove the comparison using only the site inverse and
its Gram matrix. The proof uses physical $L^2$ operators between
unit cells. It requires neither pointwise mixed-derivative estimates
nor a fine vector inverse as an input.

We retain the periodic, single-assigned-scale setting of
\cref{ch:bg-single-scale-landau}. The coarse spacing is one, the
fine spacing is $h=L^{-k}$, and every site and field row is present.
All constants below can depend on $L$ and on the fixed upper and
lower row weights, but not on $h$ or the compatible periodic volume.
The real curvature ceiling is fixed below the one for the actual
site gap and normalized Gram gap already proved in
\cref{ch:layered-current-precision}.

\subsubsection{Derivative bounds for the smoothed projector}

Let $B_y$ be the fine cell belonging to coarse site $y$. A fine bond
is assigned to the cell of its initial site. Distances between cells
are physical coarse lattice distances, denoted by $d(y,z)$.
Write
\begin{equation}\label{loc:site-operators}
 A_s=D_U^\sharp D_U+q_U^\sharp q_U,\quad
 G=A_s^{-1},\quad F=Gq_U^\sharp,\quad
 \mathcal M=F^\sharp F,\quad
 P=F\mathcal M^{-1}F^\sharp=1-R.
\end{equation}
Here the $q_U$ and all transports are the actual ones. In this
homogeneous setting the site gap and Gram gap read
\begin{equation}\label{loc:site-input-gaps}
 \langle u,A_su\rangle\ge\beta\|u\|_2^2,
 \qquad m_0I\le\mathcal M\le m_1I,
\end{equation}
with fixed positive constants. The first bound is the current-free
site gap. We justify the second bound here without importing a
site lift that requires differentiated background regularity.

\begin{lemma}[A current-free energy lift and Gram lower bound]
\label{lem:loc-energy-lift-gram}
At a fixed sufficiently small real curvature ceiling there is a
right inverse $E$ of the actual site average with
\begin{equation}\label{loc:site-energy-lift}
 q_UE=1,\qquad
 \|Eb\|_{A_s}^2:=\|D_UEb\|_2^2+\|q_UEb\|_2^2
                                                   \le C_E^2\|b\|_2^2.
\end{equation}
It follows that the actual Gram matrix satisfies
\begin{equation}\label{loc:current-free-gram-gap}
                  C_E^{-4}I\le F^\sharp F\le\beta^{-2}I.
\end{equation}
No bound on the physical current or on a derivative of a local
potential is used.
\end{lemma}

\begin{proof}
The site rows use disjoint pillars. Since each coefficient is a
unitary adjoint transport with averaging weight $h^4$, their
physical transpose obeys $q_Uq_U^\sharp=1$. Thus
$\|q_U\|=\|q_U^\sharp\|=1$.

First work on a periodic volume large enough to embed the fixed
surrounding cubes used below. For each unit pillar choose a real
comb frame on such a cube. Its potential has supremum at most
$C_L\kappa$, using
curvature only. In this frame every actual hierarchical internal
transport differs from one by $C_L\kappa$: the admitted intermediate
open transports have size $C_L\kappa\xi_s$, and the physical
scales obey $\sum_{s<k}\xi_s\le1/(L-1)$. This is the current-free
pillar transport estimate. It does not differentiate that potential.

Choose one fixed nonnegative smooth scalar bump $\phi$ supported
strictly inside the centered unit cell, with $\phi(0)>0$.
Its sampled mean $m_h=h^4\sum_x\phi(x)$ has a uniform positive
lower bound for $h=L^{-k}$. For fine meshes this follows from its
positive integral; the finitely many larger meshes contain the
center site and have a positive mean as well. Its physical $L^2$
norm and the physical $L^2$ norm of its first differences are
uniformly bounded, by the smooth bounds and the fixed support.

On the pillar at $y$ put $u^{(y)}(x)=\phi(x-y)\psi_y$ in this
local frame and zero outside. The actual site readout there is
$N_y\psi_y$, where
$\|N_y-m_hI\|\le C_L\kappa m_h$. Decrease the curvature ceiling
so that $N_y$ is invertible with $\|N_y^{-1}\|\le2/m_h$.
Choose $\psi_y$ to give the prescribed $b_y$, including the
unitary coarse-frame rotation. No other site row sees this bump,
so the sum of these fields is an exact right inverse $E$.

The first-difference product rule in its local frame gives
\[
 \|D_Uu^{(y)}\|_2
 \le C\bigl(\|D\phi\|_2
                  +\kappa\|\phi\|_2\bigr)|\psi_y|.
\]
Only the potential's supremum enters. The first-difference collars
of the bumps have bounded overlap, including the finitely coarse
meshes. Sum over pillars and use $q_UEb=b$ to obtain
\eqref{loc:site-energy-lift}.

For a smaller torus perform this construction on a sufficiently
large finite periodic cover. Average the resulting right inverse
over the finite deck group, acting simultaneously on its fine
output and coarse input. Every conjugate remains a right inverse
and has the same energy bound, since the pulled-back coefficients
commute with those unitary deck translations. Their average has
the same bound and preserves the deck-invariant subspaces. It
therefore descends to an exact right inverse on the original
torus. Its squared input and output norms acquire the same number
of sheets, so the constant in \eqref{loc:site-energy-lift} is
unchanged. This proves the energy lift even when a surrounding
comb cube does not embed in the original small torus.

Finally $A_sFb=q_U^\sharp b$, and hence
\[
 \|b\|_2^2=\langle A_sFb,Eb\rangle
       \le C_E\|Fb\|_{A_s}\|b\|_2,
 \qquad
 \|Fb\|_{A_s}^2=\langle q_U^\sharp b,Fb\rangle
                                      \le\|b\|_2\|Fb\|_2.
\]
For $b\ne0$ these imply $\|Fb\|_2\ge C_E^{-2}\|b\|_2$.
The upper bound follows from $\|G\|\le\beta^{-1}$ and
$\|q_U^\sharp\|=1$. This proves
\eqref{loc:current-free-gram-gap}.
\end{proof}

\begin{lemma}[Exponential cell bounds with one derivative]
\label{lem:loc-smoothed-derivatives}
There are positive constants $s_0$ and $C$ with the following
property. The operators $F$, $D_UF$ and $\mathcal M^{-1}$ have
exponentially decreasing cell $L^2$ operator entries. In particular
\begin{equation}\label{loc:smoothing-kernels}
 \begin{split}
 T&=PD_U^\sharp=F\mathcal M^{-1}(D_UF)^\sharp,\\
 K_s&=D_UPD_U^\sharp=(D_UF)\mathcal M^{-1}(D_UF)^\sharp
 \end{split}
\end{equation}
are bounded on their physical $L^2$ spaces, with
\begin{equation}\label{loc:cell-exponential-moments}
 \|1_{B_y}T1_{B_z}\|+\|1_{B_y}K_s1_{B_z}\|
                                     \le Ce^{-s_0d(y,z)}.
\end{equation}
Every fixed polynomial cell moment of these bounds is finite
uniformly in the fine mesh and volume.
\end{lemma}

\begin{proof}
Let $\psi$ be bounded and physically one-Lipschitz on fine sites,
and give coarse sites the sampled values $\psi(y)$.
For $w=e^{s\psi}$, the fine-gradient commutator satisfies
\[
 \|wD_Uw^{-1}-D_U\|_{2\to2}\le C|s|\qquad(|s|\le1).
\]
The difference of weights across one fine bond is $O(|s|h)$;
division by $h$ therefore leaves a physical constant.
Similarly $w_fq_U^\sharp w_c^{-1}-q_U^\sharp$ has norm
$O_L(|s|)$, since the actual site-row carrier has physical
diameter bounded independently of $k$, and its physical Schur sums
are bounded. Conjugating the site form changes it by at most
$C_L|s|$ times the two site $H^1$ norms.

By \eqref{loc:site-input-gaps}, for fixed sufficiently small $|s|$
the conjugated inverse is bounded from $H^{-1}$ to $H^1$.
Consequently
\begin{equation}\label{loc:weighted-F-energy}
 \|w_fFw_c^{-1}\|_{2\to2}
       +\|w_fD_UFw_c^{-1}\|_{2\to2}\le C_L.
\end{equation}
For the derivative term, differentiate the weighted output and
subtract the first commutator above. Thus only one site energy
derivative is used.
The same conjugation argument and resolvent identity give
$\|w_c\mathcal Mw_c^{-1}-\mathcal M\|\le C_L|s|$.
Its inverse is uniformly bounded for smaller $|s|$ by the real
Gram lower bound. This is the weighted Gram argument, now keeping
the energy derivative in \eqref{loc:weighted-F-energy}.

Insert these weighted bounds in \eqref{loc:smoothing-kernels}.
The adjoint factors use the same bounds with $-\psi$.
Choose $\psi$ to be a truncated distance from the input cell;
its values differ from cell distance by at most a fixed constant.
This gives \eqref{loc:cell-exponential-moments}.
Homogeneous four-dimensional cells have at most $C(n+1)^3$
neighbors at distance between $n$ and $n+1$. Summing against the
exponential proves every fixed polynomial moment.

All weights can be chosen on the finite periodic lattice. Truncated
distance functions have the required physical Lipschitz bounds,
independently of its size.
\end{proof}

The derivative in $PD_U^\sharp$ has fallen on the smoothed column
$F$. Its boundedness here is not an order-zero $L^\infty$
projection estimate. All the bounds in the lemma are physical
$L^2$ cell bounds.

\subsubsection{Cutoffs and the nonlocal IMS bound}

Choose a smooth partition $\{\chi_i\}$ on the covering lattice,
with supports of radius $CR_0$, where $R_0\ge2$ is an integer.
It may be chosen so that
\begin{equation}\label{loc:partition-bounds}
 \begin{split}
 &\sum_i\chi_i(x)^2=1,\qquad
 \sum_i|D\chi_i(x)|^2\le C R_0^{-2},\\
 &\sum_i|\chi_i(x)-\chi_i(z)|^2
                         \le C\min\{1,|x-z|^2/R_0^2\}.
 \end{split}
\end{equation}
For instance translate one fixed smooth bump at every integer
coarse site, scale its radius by $R_0$, and divide by the square
root of the sum of its squared translates. This sum is comparable
with $R_0^4$. The construction is invariant under integer
translations and has the bounds displayed, including
\begin{equation}\label{loc:within-cell-cutoffs}
 \sum_i\sup_{x\in B_y}|\chi_i(x)-\chi_i(y)|^2\le C R_0^{-2}.
\end{equation}
The number of summands can grow with $R_0$; their squared derivative
sum has the stated bound. Cutoffs on fine bonds are sampled at their
initial sites.

\begin{lemma}[A composed commutator and an IMS estimate]
\label{lem:loc-projector-ims}
For these cutoffs,
\begin{equation}\label{loc:composed-projector-commutator}
 \sum_i\|[P,\chi_i]D_U^\sharp v\|_2^2
                                      \le C R_0^{-2}\|v\|_2^2.
\end{equation}
The projected gauge form satisfies the sharper quadratic identity
with error
\begin{equation}\label{loc:gauge-ims}
 \left|\sum_i\|RD_U^\sharp(\chi_iv)\|_2^2
                         -\|RD_U^\sharp v\|_2^2\right|
                                      \le C R_0^{-2}\|v\|_2^2.
\end{equation}
The sums use the corresponding finite periodic cutoff family.
\end{lemma}

\begin{proof}
Put $B_i=D_U^\sharp\chi_i-\chi_iD_U^\sharp$, with the site and
bond multiplications understood. The local product rule and
\eqref{loc:partition-bounds} give
$\sum_i\|B_iv\|_2^2\le C R_0^{-2}\|v\|_2^2$.
The exact algebra is
\begin{equation}\label{loc:commutator-algebra}
          [P,\chi_i]D_U^\sharp=[T,\chi_i]-PB_i,
                  \qquad T=PD_U^\sharp.
\end{equation}

To estimate $[T,\chi_i]$, first replace each cutoff on each cell
by its center value. The input and output remainders satisfy the
vector-valued bound $CR_0^{-1}$ by
\eqref{loc:within-cell-cutoffs} and the $L^2$ bound on $T$.
For the center values, the cell entry of the commutator is
$(\chi_i(z)-\chi_i(y))T_{yz}$. Let
$t_{yz}=\|T_{yz}\|$ and $\rho_{yz}=1+d(y,z)$.
Cauchy--Schwarz with weights $\rho_{yz}t_{yz}$ gives
\[
 \sum_i\left\|\sum_z(\chi_i(z)-\chi_i(y))T_{yz}v_z\right\|^2
 \le\left(\sum_z\rho_{yz}t_{yz}\right)
       \sum_z\frac{t_{yz}}{\rho_{yz}}
          \sum_i|\chi_i(z)-\chi_i(y)|^2\|v_z\|^2.
\]
Sum over $y$. The first row moment and first column moment are
finite by \cref{lem:loc-smoothed-derivatives};
\eqref{loc:partition-bounds} supplies $C\rho_{yz}^2/R_0^2$.
This proves \eqref{loc:composed-projector-commutator}.

For the quadratic estimate use
\begin{equation}\label{loc:gauge-smoothing-split}
 \|RD_U^\sharp v\|_2^2
      =\|D_U^\sharp v\|_2^2-\langle v,K_sv\rangle.
\end{equation}
The first term is the form of the local operator $D_UD_U^\sharp$.
Its double-cutoff identity has error at most
$CR_0^{-2}\|v\|_2^2$: its coefficients have range at most $2h$
and size at most $Ch^{-2}$, while the squared cutoff difference
is $O(h^2/R_0^2)$.

We give the double-cutoff estimate for the smoothing term in
\eqref{loc:gauge-smoothing-split}, since an operator norm of a
commutator alone would not give its quadratic improvement.
For a cell entry $K_{yz}$ of $K_s$, the contribution is
\[
           \sum_i\chi_i|_{B_y}\,K_{yz}\,\chi_i|_{B_z}-K_{yz}.
\]
Write $\chi_i|_{B_y}=c_i(y)+\delta_i^y$, where
$c_i(y)=\chi_i(y)$. Use
$\sum_i(\chi_i|_{B_y})^2=1$ on each cell to expand it as minus
one half of
\begin{equation}\label{loc:double-cutoff-entry}
 \begin{split}
 &\sum_i(c_i(y)-c_i(z))^2K_{yz}\\
 &+2\sum_i(c_i(y)-c_i(z))
                             (\delta_i^yK_{yz}-K_{yz}\delta_i^z)\\
 &+\sum_i\bigl((\delta_i^y)^2K_{yz}
            -2\delta_i^yK_{yz}\delta_i^z+K_{yz}(\delta_i^z)^2\bigr).
 \end{split}
\end{equation}
The first line is bounded by
$C d(y,z)^2R_0^{-2}\|K_{yz}\|$.
For the second line use Cauchy--Schwarz in $i$,
\eqref{loc:partition-bounds} and
\eqref{loc:within-cell-cutoffs}. Its bound is
$C(1+d(y,z))R_0^{-2}\|K_{yz}\|$.
For the last line use the squared within-cell bounds at both
ends. Its bound is $CR_0^{-2}\|K_{yz}\|$.
Thus the entry is at most
$CR_0^{-2}(1+d(y,z))^2\|K_{yz}\|$ in operator norm.
The exponential second moment and Schur's test prove
\eqref{loc:gauge-ims}.
\end{proof}

The same partition gives
\begin{equation}\label{loc:full-ims}
 \left|\sum_i\mathcal K_{U,\omega}(\chi_iv,\chi_iv)
                      -\mathcal K_{U,\omega}(v,v)\right|
                                  \le C R_0^{-2}\|v\|_2^2.
\end{equation}
Here $\mathcal K$ is exactly \eqref{slp:full-form}.
For the Wilson Hessian the coefficient matrix has range $O(h)$
and physical size $O(h^{-2})$, so the same local double-cutoff
calculation applies. The operator $Q_U^\sharp aQ_U$ has bounded
physical Schur sums and bounded physical carrier diameter, by the
actual row and column estimates. The actual constraint Hessian has
the same properties, with an extra factor $\|\omega\|_\infty$,
by its two-leg moment. Their double-cutoff errors are therefore
$O(R_0^{-2})$ as well. This proves
\eqref{loc:full-ims} without replacing either row by a straight
main term.

\subsubsection{Shielding the projector from the exterior}

\begin{lemma}[Local comparison of the actual projector]
\label{lem:loc-projector-shielding}
Consider two real admitted single-scale connections whose operators
have the uniform site, Gram and weighted energy bounds above.
They may live on different periodic volumes, with the common cubes
embedded injectively in both. Suppose, in common local frames, the
connections and their actual site-row carriers agree on $Q_{24R}$,
where $R$ exceeds a fixed carrier-diameter constant.
For fine sources supported in $Q_{2R}$,
\begin{equation}\label{loc:projector-shielding}
 \|1_{Q_{2R}}(R_1-R_2)1_{Q_{2R}}\|_{2\to2}
                                      \le Ce^{-cR}.
\end{equation}
The compression is identified by their common fine coordinates.
\end{lemma}

\begin{proof}
We separate the fine resolvent and coarse Gram inverse comparisons.
All cutoffs below are supported in the common region. Their
commutators use only coefficient stencils in that region.

First let $f,g$ be fine sources supported in $Q_{10R}$.
Choose $\theta=1$ on $Q_{12R}$ and supported in $Q_{16R}$.
The geometric resolvent identity, with the cutoff embedding the
two spaces, is
\begin{equation}\label{loc:geometric-resolvent}
 \theta G_1-G_2\theta
                  =G_2(A_{s,2}\theta-\theta A_{s,1})G_1.
\end{equation}
The middle operator is supported in an annulus separated from the
source region. Its $H^1\to H^{-1}$ norm is bounded uniformly:
the gradient commutator contains a physical derivative of $\theta$;
the site-row commutator has bounded range and physical Schur sums.
There is no subtraction of different connection coefficients
outside this annulus. The weighted site energy bound makes both
resolvent solutions exponentially small in the annulus.
Pairing the identity with $g$ proves a compressed resolvent
difference $Ce^{-cR}$. The same energy argument gives its localized
one-derivative version. The distances can be decreased by a fixed
row-carrier radius without changing this conclusion.

Now put $\Lambda=\{y:B_y\subset Q_{6R}\}$ for coarse nodes.
For a node vector supported in $\Lambda$, its source $q^\sharp b$
is supported in $Q_{7R}$, and the two actual $q^\sharp$ agree
there. The just-proved comparison controls the difference of the
two $F$ outputs in $Q_{10R}$. Their tails outside $Q_{10R}$ are
exponentially small by the weighted $F$ bound. Therefore
\begin{equation}\label{loc:compressed-gram-difference}
 \|1_\Lambda(\mathcal M_1-\mathcal M_2)1_\Lambda\|
                                      \le Ce^{-cR}.
\end{equation}
This is an operator bound for arbitrary vectors supported in
$\Lambda$, not a sum of separate entry errors over its volume.

Each compressed Gram matrix
$\mathcal M_{j,\Lambda}=1_\Lambda\mathcal M_j1_\Lambda$
retains the lower bound $m_0$. Its inverse is bounded by $m_0^{-1}$,
so the resolvent identity bounds the difference of their inverses
by $Ce^{-cR}$.
For a node source $b$ supported in $Q_{3R}$, compare that finite
compression inverse with the global inverse. The exact identity is
\begin{equation}\label{loc:compression-inverse-identity}
 \mathcal M_{j,\Lambda}^{-1}b-1_\Lambda\mathcal M_j^{-1}b
 =\mathcal M_{j,\Lambda}^{-1}1_\Lambda\mathcal M_j
                    1_{\Lambda^c}\mathcal M_j^{-1}b.
\end{equation}
The global inverse has an exponentially small tail on
$\Lambda^c$, by the weighted Gram estimate. The other two
operators on the right are uniformly bounded. It follows that
the compressions of $\mathcal M_1^{-1}$ and
$\mathcal M_2^{-1}$ to $Q_{3R}$ differ by $Ce^{-cR}$.

Finally write $P_j=F_j\mathcal M_j^{-1}F_j^\sharp$.
If $f$ is supported in $Q_{2R}$, $F_j^\sharp f$ has an exponentially
small tail outside $Q_{3R}$. Their retained parts differ by
$Ce^{-cR}$ from the fine resolvent comparison. Insert these facts
and the compressed Gram-inverse comparison into
$\langle f,P_jg\rangle$. All uncompressed factors have uniform
norms. We obtain \eqref{loc:projector-shielding}, since
$R_j=1-P_j$.
\end{proof}

Only the existing site and Gram estimates entered this comparison.
No fine vector Green function or lower bound for the full precision
has been used in its proof.

\subsubsection{The nonflat periodic gap from local curvature bounds}

\begin{proposition}[A one-scale curvature-small Landau inverse]
\label{prop:loc-periodic-nonflat-gap}
Fix $L$ and the row bounds $a_*,a^*$. There are positive constants
$\kappa_*$, $\gamma_*$, $c$ and $C$,
independent of $k$, periodic volume and real holonomy, with the
following property. Let $U$ be real on a compatible periodic fine
lattice, with every site and field row at the same assigned
physical scale one. Suppose
\begin{equation}\label{loc:global-smallness}
        \|F(U)\|_\infty\le\kappa_* \quad\hbox{and}\quad
                     \|\omega\|_\infty\le\gamma_*.
\end{equation}
Use the actual straightened nonlinear constraint and actual
projected Landau gauge in \eqref{slp:full-form}. Then
\begin{equation}\label{loc:global-gap}
 c\bigl(\|\nabla_Uv\|_2^2+\|v\|_2^2\bigr)
 \le\mathcal K_{U,\omega}(v,v)
 \le C\bigl(\|\nabla_Uv\|_2^2+\|v\|_2^2\bigr).
\end{equation}
The full penalized inverse and the real actual-slice saddle inverse
have the uniform physical bounds stated in
\cref{prop:slp-nonflat-gap,cor:slp-saddle-inverse}, with no common
global small-gauge assumption.
\end{proposition}

\begin{proof}
Choose an integer $R$ large enough for the carrier margins below.
It will be fixed before the curvature ceiling is chosen.
First assume that the periodic side is at least $512R$.
At every coarse center use the local comb gauge and extension
construction on a cube of radius $32R$. Apply the scaled version
of \cref{lem:bg-comb-extension} with its radius parameter $4R$.
The resulting auxiliary periodic potential, on a torus of side
$128R$, satisfies
\begin{equation}\label{loc:local-periodic-extension}
 \|A_i^{\rm ext}\|_\infty\le CR\kappa,
 \qquad \|F_i^{\rm ext}\|_\infty\le C\kappa,
\end{equation}
and agrees with the locally gauged original field on $Q_{24R}$,
provided $R^2\kappa$ is sufficiently small. The estimates use
curvature only. The current derivative bounds of the Coulomb gauge
theorem are not needed for this coercivity step.

All operators in the present precision transform covariantly under
this local frame change. For the nonlinear straightened constraint,
both its internal means and its endpoint means are bi-equivariant;
therefore its output logarithm is conjugated at the coarse initial
site. Its first and second derivatives transform accordingly.
Transform $\omega$ by that same unitary conjugation. The physical
pairings and its supremum bound are unchanged. This covariance
statement concerns the projected Landau form. It makes no assertion
about the separately specified absolute soft contour gauge.

Take the cutoffs in \eqref{loc:partition-bounds} supported in cubes
of radius $R$ by a fixed rescaling of their defining bump. For
$v_i=\chi_iv$, the full stencil of $D_U^\sharp v_i$ is in
$Q_{2R}$. Extend $\omega$ and the row weights to the local torus,
keeping their actual values on every row whose carrier meets that
support, and keeping $\|\omega\|_\infty\le\gamma$ and
$a_*\le a\le a^*$. The Wilson, row-mass and constraint-curvature
forms of $v_i$ then agree exactly with those in the local extension.
The only difference is the projector.

For a real orthogonal projector,
$\|Rf\|^2=\langle f,Rf\rangle$. Consequently
\cref{lem:loc-projector-shielding}, with
$f=D_U^\sharp v_i$, gives
\begin{equation}\label{loc:local-form-comparison}
 |\mathcal K_U(v_i,v_i)-\mathcal K_i^{\rm ext}(v_i,v_i)|
                             \le Ce^{-cR}\|v_i\|_X^2.
\end{equation}
The coefficient values outside the common region have not been
subtracted in a bare fine operator norm.

Choose $\gamma_*$ smaller than the multiplier ceiling of
\cref{prop:slp-nonflat-gap}. Once $R$ is fixed, choose
$\kappa_*$ so that $R^2\kappa_*$ admits
\eqref{loc:local-periodic-extension} and
$CR\kappa_*+C\kappa_*+\gamma_*$ is below the common-gauge
ceiling in that proposition. Each auxiliary periodic precision
then has the same lower constant $c_1>0$.
It follows that
\begin{equation}\label{loc:localized-lower-bound}
             \mathcal K_U(v_i,v_i)
                      \ge(c_1-Ce^{-cR})\|v_i\|_X^2.
\end{equation}

The usual covariant gradient double-cutoff calculation also gives
\begin{equation}\label{loc:gradient-ims}
 \left|\sum_i\|v_i\|_X^2-\|v\|_X^2\right|
                                      \le C R^{-2}\|v\|_2^2.
\end{equation}
The covariant shifts are unitary. Their fine denominator cancels
the squared cutoff difference just as in the local term of
\cref{lem:loc-projector-ims}.
Sum \eqref{loc:localized-lower-bound}, use
\eqref{loc:full-ims} and \eqref{loc:gradient-ims}, and first
choose $R$ so that $Ce^{-cR}\le c_1/4$ and the total
$CR^{-2}$ error is at most $c_1/4$. The constants used for
these choices are uniform throughout a fixed initial small
curvature and multiplier ceiling. The subsequent decrease of
$\kappa_*,\gamma_*$ does not change them. This proves the lower
bound in \eqref{loc:global-gap}. The upper bound follows from
the local Wilson Hessian bound, $\|RD_U^\sharp v\|\le
\|D_U^\sharp v\|$, and the actual row and two-leg Schur bounds.

For a smaller periodic volume, choose an integer $t$ such that its
$t^4$-sheeted periodic cover has side at least $512R$. Pull back
the connection, the actual rows and the row multiplier to that
finite cover. The local averaging formulas and covariant
differences commute with pullback, including every periodic path
multiplicity. The same is true of their physical transposes.
The site inverse preserves the deck-invariant subspace and equals
the pullback of the smaller site's inverse there, by uniqueness
of its gapped equation. The same argument applies successively
to $F$, its Gram matrix and its Gram inverse. Consequently the
actual $R$ and the full form descend exactly. For a lifted vector,
the full form and its squared $X$ norm both multiply by $t^4$.
Divide the already proved estimate on the cover by that factor.
This proves the original finite-volume estimate without selecting
a global nearly flat gauge or restricting its torus holonomies.

The actual one-scale row lift is constructed from the same compact
two-moment seeds in local unit-cell gauges. The actual row/column
comparison makes its coarse error $O_L(\kappa)$ in both Schur
norms; its Neumann correction gives the bounds in
\eqref{slp:field-lift-bounds}. There are no adjacent-level leakage
terms in this homogeneous setting. The gauge-restoration energy
argument in \cref{cor:slp-saddle-inverse} uses only the site gap,
curvature bound and full periodic covariant Hodge identity. It is
therefore valid here as well. Repeating that corollary's kernel
coercivity argument proves the actual-slice saddle assertion.
\end{proof}

At a physical critical background whose full stationarity equation
is $J_A(U)=Q_U^\sharp\omega$, the column $L^1$ bound for this
actual row lift again gives
$\|\omega\|_\infty\le C_L\|J_A(U)\|_\infty$.
Thus sufficiently small physical curvature and actual current
supply the two numerical ceilings. This does not construct the
critical background or infer the equation from criticality on an
unidentified slice.

Restoring a common physical coarse spacing $\xi$ replaces the
mass norm by $\xi^{-2}\|v\|_2^2$, the curvature ceiling by
$\xi^2\|F\|_\infty\le\kappa_*$, and the multiplier ceiling by
$\xi^3\|\omega\|_\infty\le\gamma_*$. The full inverse then has
$L^2\to L^2$ norm at most $C\xi^2$. These powers follow from the
declared physical pairings and do not depend on the number of
fine bonds in a coarse cell.

\subsubsection{The scope of the recovered inverse}

The new conclusion is a real, unweighted Wilson/Lagrangian
precision estimate at one assigned physical scale, with its actual
straightened field rows and actual projected Landau gauge. The
proof of the nonlocal localization used only the separately proved
site and Gram package. In particular it does not assume the vector
inverse in order to prove that inverse.

Several tasks remain in the general background problem. Different
assigned scales change the cell metric and carrier bookkeeping.
Proper-region inner-boundary forms have boundary terms in both
Hodge and shielding identities. Independent analytic-current
completions contain additional specified defects. An identification
with the absolute contour precision has its own gauge-restoration
terms. None of these is an equality supplied by the periodic
one-scale result. The nonlinear trial residual and construction of
the many-level critical orbit remain separate from the coercive
linear inverse established here.

% END PROOF EXPANSION background/single_scale_localization.tex

% BEGIN LAYERED EXPANSION inverse/weighted_site_restoration.tex
\subsection{Weighted site inverses on the actual determining region}
\label{ch:layered-site-restoration}

The projected gauge removes site directions with zero assigned
average. Its inverse must therefore control these directions in the
physical units of their own pillars. We prove that control on the
actual active graph, including the outer boundary. We also prove a
normalized Gram gap from curvature alone. A differentiated local
potential and a physical-current bound are not premises of this
Gram estimate.

The setting is the separated aligned sequence of
\cref{ch:multilevel-backgrounds}. The fine spacing is
$\eta=L^{-k}$ and the assigned spacings are $\xi_j=L^{j-k}$.
The pillars $B_y$ partition $\Omega_0$; write $\xi(x)=\xi_y$ on
$B_y$. Outside $\Omega_0$ put $\xi(x)=\eta$ for geometric purposes
only. All operators below act on the original graph of sites and
bonds inner to $\Omega_0$. The action still contains precisely its
inner plaquettes.

Fine scalar and vector norms use measure $\eta^4$. Top norms use
the physical measure $\xi_y^4$, and field-row norms use
$\xi_c^4$. The actual own-pillar site average is denoted by $q$.
Thus $q$ is the transported average in
\eqref{px:own-pillar-site-average}, not an unextended endpoint
readout. Its physical transpose satisfies $qq^\sharp=1$.
Let
\begin{equation}\label{lsi:site-operators}
 L_U=D_U^\sharp D_U,\qquad
 A_s=L_U+q^\sharp a_sq,\qquad
 a_s(y)=c_y\xi_y^{-2},\qquad
          0<c_*\le c_y\le c^*<\infty.
\end{equation}
The physical gradient is
$D_U\lambda(b)=\eta^{-1}(\lambda(b_-)-\operatorname{Ad}U_b\lambda(b_+))$.
Adjoints in real positivity statements are the Hilbert adjoints
for these physical measures. All the estimates extend to the
complexification of the real coefficient spaces.

Assume the actual means are on their regular branches and
\begin{equation}\label{lsi:curvature-ceiling}
 \left\|\frac1{i\eta^2}\log U(\partial p)\right\|
                          \le\kappa\,\mathsf L(p)^{-2}
                     \quad(p\subset\Omega_0),
\end{equation}
with a fixed sufficiently small $\kappa$. The branch admission
and the estimates on the actual pillar transports follow from
the current-free comb comparison already proved in
\cref{ch:layered-current-precision}. Constants below depend only
on $L,c_*,c^*$ and fixed geometric collar choices. We specify the
additional lower bound on the separation $\rho$ below. There is
no dependence on depth, compatible volume or number of components.
Enlarge the fixed $C_{\rm geom}$ once so it exceeds four times
the chosen comb-cube radius in units of a pillar side.

\subsubsection{A slowly varying length comparable to the assigned scale}

The assigned scale jumps by a factor $L$ at a layer boundary.
Differentiating that step function would introduce a fine-mesh
loss. We instead use a comparable Lipschitz length.

\begin{lemma}[A smooth comparison length]
\label{lem:lsi-comparison-length}
For $0<\delta<1$, define on the ambient torus
\begin{equation}\label{lsi:comparison-length}
 \ell(x)=\inf_z\{\xi(z)+\delta\operatorname{dist}_\infty(x,z)\}.
\end{equation}
If $\rho\delta\ge1$, then
\begin{equation}\label{lsi:comparison-length-bounds}
            L^{-1}\xi(x)\le\ell(x)\le\xi(x),
 \qquad |\ell(x)-\ell(z)|\le
                            \delta\operatorname{dist}_\infty(x,z).
\end{equation}
On the bare layer and outside the active region, $\ell=\eta$.
For $w=\ell^{-1}$, neighboring fine sites satisfy
\begin{equation}\label{lsi:weight-differences}
 |D\ell|\le\delta,\qquad |Dw|\le C\delta\ell^{-2}.
\end{equation}
The site and bond choices of $\ell$ are equivalent with absolute
constants. The actual bond scale $\mathsf L(b)$ is comparable to
them with a constant depending only on $L$.
\end{lemma}

\begin{proof}
The upper bound follows by choosing $z=x$. Infima of functions
with Lipschitz constant $\delta$ have that same Lipschitz bound.
Suppose $x$ belongs to layer $j$. If $z$ has assigned level at
least $j-1$, then $\xi(z)\ge\xi_j/L$. If its level is at most
$j-2$, it lies outside $\Omega_{j-1}$, whereas $x\in\Omega_j$.
Separation gives
$\operatorname{dist}_\infty(x,z)>\rho\xi_{j-1}$, so the distance
term is at least $\xi_j/L$. The same argument includes exterior
points, assigned the bare scale. For $j=0$, the global lower
bound $\xi(z)\ge\eta$ suffices. This proves the lower bound.
At a bare point the lower bound $\ell\ge\eta$ and the upper
bound $\ell\le\eta$ coincide.

The inverse-weight estimate follows by subtracting reciprocals.
Since $\ell\ge\eta$ and its change over one fine step is at
most $\delta\eta$, neighboring values have a fixed ratio.
Fine neighbors can have assigned levels differing by at most one:
otherwise they would cross the separated intervening layer in
physical distance at most $\eta$. Thus their assigned scales,
and the largest corner scale used on a bond or plaquette, differ
by at most a fixed power of $L$. This proves the last assertion.
\end{proof}

The comparison metric obtained by integrating $ds/\ell$ is
equivalent to the actual layer metric within a factor $L$.
We will choose $\delta$ small using fixed block constants, and
then choose $\rho\ge\max\{C_{\rm geom}L,\delta^{-1}\}$.
These choices precede the depth and volume.

\subsubsection{Current-free block inequalities and a weighted site lift}

The real block argument gives two useful versions of Poincar\'e:
\begin{equation}\label{lsi:block-poincare}
 \begin{split}
 \|\xi^{-1}u\|_2^2
   &\le C\bigl(\|D_Uu\|_2^2+
                                  \|\Xi^{-1}qu\|_{\rm top}^2\bigr),\\
 \|u\|_2^2
   &\le C\bigl(\|\ell D_Uu\|_2^2+\|qu\|_{\rm top}^2\bigr).
 \end{split}
\end{equation}
Here $\Xi$ is multiplication by $\xi_y$ on tops. The first is
the actual weighted site gap. For the second multiply each
pillar's block inequality by $\xi_y^2$ before summing, and use
\cref{lem:lsi-comparison-length}. Only internal pillar bonds are
used in either estimate, so crossing energies have a favorable
sign. Both inequalities use curvature and unitarity, not a current.
In particular $A_s$ is strictly positive on its finite site space.

\begin{lemma}[A weighted right inverse for the actual site average]
\label{lem:lsi-weighted-site-lift}
There is a real linear map $E$ with $qE=1$ such that, for each
fixed $r\in\{-1,0,1\}$,
\begin{equation}\label{lsi:weighted-lift-bounds}
 \|\ell^{-r}E\psi\|_2\le C_r\|\Xi^{-r}\psi\|_{\rm top},
 \qquad
 \|\ell^{-r}D_UE\psi\|_2
                       \le C_r\|\Xi^{-r-1}\psi\|_{\rm top}.
\end{equation}
The constants are uniform over all layers and proper determining
regions in this setting.
\end{lemma}

\begin{proof}
For a nonbare pillar use a comb frame on a fixed surrounding cube
of radius $C_L\xi_y$. Separation places the required cube inside
$\Omega_{j-1}\subset\Omega_0$ when the pillar has level $j\ge1$.
Its local potential has supremum at most $C_L\kappa\xi_y^{-1}$.
The actual hierarchical internal transports in the pillar differ
from one by $C_L\kappa$, by their geometric sum over intermediate
physical scales. This comparison differentiates no potential.

Choose a fixed nonnegative smooth scalar bump supported strictly
inside the unit centered cube, positive at its center, and sample
its rescaling by $\xi_y$. Its mean $m_y$ on the pillar has positive
lower and finite upper bounds independent of its level. Fine meshes
converge to the positive integral; the finitely coarse possibilities
contain the center site. The actual readout of this bump times a
constant colour is $N_y$ times that colour, with
$\|N_y-m_yI\|\le C_L\kappa m_y$. Normalize by $N_y^{-1}$ and
the unitary top-frame rotation. This gives the prescribed site
value exactly. No other site row reads this bump.

The bump has amplitude $O(|\psi_y|)$ and first differences
$O(\xi_y^{-1}|\psi_y|)$ in its local frame. The connection term
has the same first-difference bound. Its physical support volume
is $O(\xi_y^4)$. Multiplying by $\ell^{-r}$ and using
\eqref{lsi:comparison-length-bounds} gives exactly the two powers
of $\Xi$ in \eqref{lsi:weighted-lift-bounds}.
For a bare pillar take its one site value: $q_0$ is the identity.
The first-difference collar of this seed has the same estimates,
since $\xi_y=\eta$ and all transports are unitary.
The collars of different pillars have bounded overlap; adjacent
scales have ratio at most $L$. Summing proves the claims.

If a surrounding cube does not embed in the original torus, make
the construction on a sufficiently large finite periodic cover
and average $E$ over the deck group, acting on its fine output
and coarse input simultaneously. Pullback preserves the actual
rows, $\Xi$ and $\ell$. Deck translations are unitary in each
of the weighted norms above. The averaged map retains $qE=1$
and all the bounds, preserves deck-invariant vectors and descends.
The squared norms have the same sheet-count factor. This also
handles nontrivial flat holonomies without a global gauge choice.
\end{proof}

\subsubsection{The normalized layered Gram matrix from weighted energy}

Put
\begin{equation}\label{lsi:normalized-F-M}
 G=A_s^{-1},\qquad F=Gq^\sharp\Xi^{-2},\qquad
                         M=F^\sharp F.
\end{equation}
The normalization is the one in the actual site projector. It
retains the physical top measures; an unweighted array transpose
would give a different matrix.

\begin{proposition}[A current-free normalized Gram gap]
\label{prop:lsi-gram-gap}
For a sufficiently small fixed $\delta$ and the corresponding
separation above, there are $m_*,m^*>0$ such that
\begin{equation}\label{lsi:gram-gap}
                         m_*I\le M\le m^*I.
\end{equation}
They are independent of depth, region shape and compatible volume.
Only \eqref{lsi:curvature-ceiling} and the actual site geometry
are required; no physical current or differentiated local-gauge
bound is a premise.
\end{proposition}

\begin{proof}
Write $u=Fb$ and introduce
\begin{equation}\label{lsi:plus-energy}
             \mathcal E_+(u)=\|\ell D_Uu\|_2^2+\|qu\|_{\rm top}^2.
\end{equation}
By \eqref{lsi:block-poincare}, $\|u\|_2^2\le C\mathcal E_+(u)$.
Test $A_su=q^\sharp\Xi^{-2}b$ against $\ell^2u$.
The exact discrete product identity is
\begin{equation}\label{lsi:weighted-green-identity}
 \operatorname{Re}\langle D_Uu,D_U(\ell^2u)\rangle
 =\|D_U(\ell u)\|_2^2
 -\sum_{b}\eta^4|D\ell(b)|^2
          \operatorname{Re}\langle u(b_-),\operatorname{Ad}U_bu(b_+)\rangle.
\end{equation}
Using $|D\ell|\le\delta$ and one more product rule bounds
this from below by
$\tfrac12\|\ell D_Uu\|_2^2-C\delta^2\|u\|_2^2$.

Within one pillar the variation of $\ell$ is at most
$C\delta\xi_y$, and $\ell(y)\ge\xi_y/L$. Therefore
\begin{equation}\label{lsi:q-commutator-plus}
 \|\Xi^{-2}\{q(\ell^2u)-\ell(y)^2qu\}\|_{\rm top}
                                             \le C_L\delta\|u\|_2.
\end{equation}
This follows by Cauchy--Schwarz for each unitary transported
average and the disjointness of pillars. The site-mass term in
the tested equation is consequently bounded below by
$c\|qu\|_{\rm top}^2-C_L\delta^2\|u\|_2^2$, after a fixed
Young inequality. Moreover,
\begin{equation}\label{lsi:q-weight-contraction}
                    \|\Xi^{-2}q(\ell^2u)\|_{\rm top}\le\|u\|_2,
\end{equation}
because $\ell(x)\le\xi_y$ at every site in $B_y$ and the
absolute transported averaging operator is an $L^2$ contraction.
We obtain
\[
 c\mathcal E_+(u)\le\|b\|_{\rm top}\|u\|_2
                                      +C_L\delta^2\|u\|_2^2.
\]
Choose $\delta$ so the last term is absorbed using
\eqref{lsi:block-poincare}. Then
\begin{equation}\label{lsi:plus-energy-upper}
 \mathcal E_+(Fb)\le C\|b\|_{\rm top}\|Fb\|_2,
 \qquad \|Fb\|_2\le C\|b\|_{\rm top}.
\end{equation}
This proves the Gram upper bound.

For the lower bound test the same equation against
$\varphi=E\Xi^2b$, with the exact lift from
\cref{lem:lsi-weighted-site-lift}. Then $q\varphi=\Xi^2b$ and
$\|\ell^{-1}D_U\varphi\|_2\le C\|b\|_{\rm top}$.
Thus
\begin{equation}\label{lsi:gram-energy-duality}
 \begin{split}
 \|b\|_{\rm top}^2
 &=\langle A_sFb,E\Xi^2b\rangle\\
 &\le\bigl(\|\ell D_UFb\|_2
                         +C\|qFb\|_{\rm top}\bigr)C\|b\|_{\rm top}.
 \end{split}
\end{equation}
Here the site-mass pairing is
$\langle qFb,c_yb\rangle_{\rm top}$, with no extra length factor.
For $b\ne0$, the last display gives
$\mathcal E_+(Fb)\ge c\|b\|_{\rm top}^2$.
Combine this with \eqref{lsi:plus-energy-upper} to obtain
$\|Fb\|_2\ge c\|b\|_{\rm top}$. Squaring proves the lower
bound in \eqref{lsi:gram-gap}.
\end{proof}

This argument uses first differences of a compact site lift.
It does not estimate $L_UE$ by differentiating a local connection.
In particular the weighted Gram normalization has not introduced
a hidden current premise.

\subsubsection{Two weighted derivatives on the zero-average directions}

Let $\lambda\in\ker q$. At every bare top it vanishes exactly,
because the assigned zero-level average is the identity. Therefore
\begin{equation}\label{lsi:bare-collar}
          \lambda=0\quad\hbox{on }\Omega_0\setminus\Omega_1.
\end{equation}
The separation of $\Omega_1$ from $\Omega_0^c$ supplies more than
the fixed number of fine layers needed by the following stencils.
This fact concerns the actual assigned averages. It would fail
for an arbitrary clipped average on an arbitrary subset of a block.

\begin{proposition}[A weighted inverse on $\ker q$]
\label{prop:lsi-zero-average-inverse}
For every site vector $\lambda$ with $q\lambda=0$,
\begin{equation}\label{lsi:zero-average-estimate}
 \|\xi^{-2}\lambda\|_2
       +\|\mathsf L^{-1}D_U\lambda\|_2
       +\|\nabla_UD_U\lambda\|_2
                                      \le C\|L_U\lambda\|_2.
\end{equation}
The last norm is the full component-gradient norm after zero
extension across the action boundary. Every nonzero stencil in
that norm lies inside the action region. Thus it is an estimate
for the true inner site operator, not a replacement by a compressed
periodic Laplacian.
\end{proposition}

\begin{proof}
Put $f=L_U\lambda$ and $w=\ell^{-1}$. Apply the site gap in
\eqref{lsi:block-poincare} to $w\lambda$. Since $q\lambda=0$,
\begin{equation}\label{lsi:q-commutator-minus}
 \|\Xi^{-1}q(w\lambda)\|_{\rm top}
                    \le C_L\delta\|\ell^{-2}\lambda\|_2.
\end{equation}
Indeed subtract $w(y)q\lambda$ in each pillar and use
$|w(x)-w(y)|\le C_L\delta\xi_y^{-1}$.
Since $\ell\ge\xi/L$, the left mass term in the site gap
controls $\|\ell^{-2}\lambda\|_2$ up to $L$.
Choose $\delta$ smaller if necessary and absorb
\eqref{lsi:q-commutator-minus}. We get
\begin{equation}\label{lsi:minus-mass-by-gradient}
          \|\ell^{-2}\lambda\|_2^2
                                     \le C\|D_U(w\lambda)\|_2^2.
\end{equation}
Apply the weighted Green identity with weight $w$ to the equation
$L_U\lambda=f$. The weight difference in
\eqref{lsi:weight-differences} gives
\[
 \|D_U(w\lambda)\|_2^2
 \le\|f\|_2\|\ell^{-2}\lambda\|_2
                           +C_L\delta^2\|\ell^{-2}\lambda\|_2^2.
\]
Absorb its last term using
\eqref{lsi:minus-mass-by-gradient}. This proves the weighted
mass bound. The product rule then proves
$\|\ell^{-1}D_U\lambda\|_2\le C\|f\|_2$.
The comparison length converts them into the first two terms
of \eqref{lsi:zero-average-estimate}.

For the last term extend $\lambda$ by zero outside $\Omega_0$
and extend the connection there by arbitrary unitary bonds.
By \eqref{lsi:bare-collar}, all nonzero two-step stencils are
still inner to $\Omega_0$, where the extended Laplacian agrees
with $L_U\lambda$. The covariant curl of $D_U\lambda$ is a
commutator of two shifts. Its coefficient is a conjugate of
$(\operatorname{Ad}U_p-1)/\eta^2$, so
\begin{equation}\label{lsi:curl-of-gradient}
          \|\mathrm{d}_UD_U\lambda\|_2
                           \le C_L\kappa\|\xi^{-2}\lambda\|_2.
\end{equation}
The covariant Hodge calculation in
\eqref{px:covariant-hodge-error}, applied to this zero-extended
gradient, therefore gives
\begin{equation}\label{lsi:second-derivative-hodge}
 \|\nabla_UD_U\lambda\|_2^2
 \le C\bigl(\|f\|_2^2
       +\kappa^2\|\xi^{-2}\lambda\|_2^2
       +\kappa\|\mathsf L^{-1}D_U\lambda\|_2^2\bigr).
\end{equation}
The corner scales in each plaquette are comparable within a fixed
power of $L$. The unknown exterior curvature contributes nothing:
the relevant field and its shifted values vanish in its collar.
Substitute the first two estimates to finish the proof.
\end{proof}

No derivative of the curvature, of a gauge potential or of a current
has entered this estimate. The extra boundary control comes from
the actual zero-level site constraints.

\subsubsection{Gauge restoration and the actual outward rows}

Define the actual orthogonal projector and restoration map by
\begin{equation}\label{lsi:restoration-maps}
 R=1-FM^{-1}F^\sharp,\qquad
 \Lambda=GRD_U^\sharp,\qquad S=1-D_U\Lambda.
\end{equation}
The exact range identities are
\begin{equation}\label{lsi:range-identities}
         \operatorname{ran}R=L_U(\ker q),\qquad
                      qGR=0,\qquad L_UGR=R.
\end{equation}
For $qz=0$, $GL_Uz=z$ and $F^\sharp L_Uz=\Xi^{-2}qz=0$.
Conversely $F^\sharp f=0$ implies $qGf=0$ and $L_UGf=f$.
This proves the range statement and the two identities.

For a field $v$ define
\begin{equation}\label{lsi:field-X}
 \|v\|_X^2=\|\nabla_Uv^0\|_2^2+
                                    \|\mathsf L^{-1}v\|_2^2,
\end{equation}
where $v^0$ is zero on all bonds outside the active graph.
This comparison norm counts its boundary jumps; it does not add
those jumps to the action. Its value is independent of an arbitrary
unitary extension of the connection: when a difference has two
nonzero field values its connecting bond is active, and otherwise
unitarity fixes the norm of its single nonzero value.

\begin{corollary}[Uniform restoration on every assigned scale]
\label{cor:lsi-restoration-bound}
The actual maps satisfy
\begin{equation}\label{lsi:restoration-bound}
 \|\xi^{-2}\Lambda v\|_2+\|D_U\Lambda v\|_X
                          +\|Sv\|_X\le C\|v\|_X,
 \qquad RD_U^\sharp S=0,\qquad QS=Q.
\end{equation}
Here $Q$ is the derivative of the actual straightened field map
on the literal determining rows, including its outward rows.
\end{corollary}

\begin{proof}
The site vector $\lambda=\Lambda v$ obeys $q\lambda=0$ and
$L_U\lambda=RD_U^\sharp v$. Since $R$ is orthogonal and the
inner divergence is bounded by the component-gradient comparison
norm, \cref{prop:lsi-zero-average-inverse} proves the estimates.
The identity $RD_U^\sharp S=0$ follows from
$RL_UGR=R$.

The actual field intertwining is
\begin{equation}\label{lsi:actual-outward-intertwining}
                              QD_U=\bar D_{\rm ext}q.
\end{equation}
At an outward endpoint of a level-$j$ row, $\bar D_{\rm ext}$
averages the $L^4$ assigned level-$(j-1)$ top values in that
whole outer block, with their actual mean transports. That is
the exact extension in \eqref{px:outward-site-extension}; it is
not evaluation at one outer top. Hence
$QD_U\Lambda=\bar D_{\rm ext}qGRD_U^\sharp=0$.
This proves $QS=Q$ on the true row spaces.
\end{proof}

Let $a_f(c)=d_c\xi_c^{-2}$, with fixed positive lower and upper
bounds for $d_c$, and give data the norm
\begin{equation}\label{lsi:row-data-norm}
                  \|b\|_Y^2=\sum_c\xi_c^4a_f(c)|b(c)|^2.
\end{equation}
The actual field-row construction in
\eqref{px:actual-layered-row-lift} supplies a right inverse
$E_Q$ with $\|E_Qb\|_X\le C\|b\|_Y$. Its compact seeds and
adjacent-level Neumann correction are the actual ones; they
include the identity seeds in the bare boundary layer. The
zero-extension part of $X$ adds only a bounded bare mass term
to their stated inner energy estimate.
Thus $E_L=SE_Q$ is an exact right inverse into
$X_L=\ker(RD_U^\sharp)$ and
\begin{equation}\label{lsi:slice-infsup}
 Q E_L=1,\quad\|E_Lb\|_X\le C\|b\|_Y,\quad
 \inf_{\mu\ne0}\sup_{v\in X_L\setminus\{0\}}
 \frac{|\langle\mu,Qv\rangle_{\rm rows}|}
                   {\|\mu\|_{Y^*}\|v\|_X}\ge C^{-1}.
\end{equation}
The dual norm is $\|\mu\|_{Y^*}=\|a_f^{-1/2}\mu\|_{2,\rm rows}$.
For the last inequality choose a unit data vector attaining the
dual norm and apply $E_L$. This proves the uniform inf-sup part
of the actual saddle problem. Coercivity on its simultaneous
constraint kernel is a separate estimate.

\subsubsection{The explicit current correction has an energy bound}

For an independent current define $\mathcal W_U[J]$ and the
symmetric correction $\mathcal B_U[J]$ by their actual explicit
formulas \eqref{px:current-ward-map} and
\eqref{px:symmetric-current-correction}. Suppose
\begin{equation}\label{lsi:current-amplitude}
                  \gamma_J=\sup_b\mathsf L(b)^3|J(b)|<\infty.
\end{equation}
Then
\begin{equation}\label{lsi:current-form-bound}
                 |\langle u,\mathcal B_U[J]v\rangle|
                                   \le C\gamma_J\|u\|_X\|v\|_X.
\end{equation}
Indeed, the endpoint formula for $\mathcal W$ and comparable
corner scales give
\[
 |\langle u,\mathcal W_U[J]\lambda\rangle|
       \le C_L\gamma_J\|\mathsf L^{-1}u\|_2
                                      \|\xi^{-2}\lambda\|_2.
\]
Insert this inequality and \cref{cor:lsi-restoration-bound} in
the four undifferentiated-current pairings of
\eqref{px:current-form-no-current-gradient}. This proves
\eqref{lsi:current-form-bound} without a derivative of $J$.
Likewise the actual quadratic constraint form satisfies
\[
 |\langle Su,\Delta^{(2)}[\omega]Sv\rangle|
       \le C\gamma_\omega\|u\|_X\|v\|_X,
 \qquad \gamma_\omega=\sup_c\xi_c^3|\omega(c)|,
\]
by its already proved two-leg moment. The particular induced
readout is bounded in \cref{thm:scr-induced-current} and
\eqref{ncw:multiplier-family-bound}; the exact canonical current
identification is \cref{prop:acc-current-identification}.

The results in this section supply the current-free normalized
site Gram gap, two weighted derivatives on the actual zero-average
site directions, a uniformly bounded field restoration, and the
inf-sup part of the all-layer proper-region saddle problem. They
use the literal outward-row extension and the actual bare boundary
collar. The remaining mass estimate is a vector estimate, not
another site gap.

% END LAYERED EXPANSION inverse/weighted_site_restoration.tex

% BEGIN LAYERED EXPANSION inverse/flat_mixed_scale_decomposition.tex
\subsection{A flat regular decomposition on the determining regions}
\label{ch:flat-mixed-scale-decomposition}

The missing vector estimate is a div--curl estimate with the actual
changing-scale averages. The outward rows are useful here: their
outer endpoint average is taken over an entire block whose smaller
internal rows are also prescribed. Those rows control the
fluctuations hidden by the endpoint average. We first quantify this
fact and then glue local scalar primitives. The construction retains
holes, disconnected components and the true bare boundary layer.

Use the sites, active bonds and inner plaquettes of
\cref{ch:layered-site-restoration}. Let $V$ be a real flat
connection on this active cubical complex. Flatness means every
inner plaquette is one. There is no assumption that $V$ is a
global pure gauge. In a convex fine box it can be put in a parallel
frame, so all calculations there reduce to ordinary cochains.
All norms and transposes have their previously declared physical
measures. The field-row map $Q$ is the actual straightened
derivative; at a flat field it is the transported straight
average, not the derivative of the plain average.

\subsubsection{What the outward endpoint average controls}

Make a graph with the assigned tops as vertices. Two vertices are
joined when their pillars share a face containing a fine bond.
Their levels are equal or adjacent. Let $\sigma_{yz}$ be the
larger of their two scales and let $T_{yz}$ be the parallel
transport between them inside the local one- or two-scale
rectangle. The local rectangle is contractible in the active
region, so flatness makes this transport independent of its
chosen path there.

For a fixed real $r$ define
\begin{equation}\label{fms:graph-and-row-energy}
 \begin{split}
 J_r(\psi)&=\sum_{y\sim z}\sigma_{yz}^{2-2r}
             |\psi(y)-\operatorname{Ad}T_{yz}\psi(z)|^2,\\
 D_r(\psi)&=\sum_{c\in\mathbf B}\xi_c^{4-2r}
                                     |\bar D_{\rm ext}\psi(c)|^2.
 \end{split}
\end{equation}
Count each unoriented graph edge once. The coarse difference in
the second line has the exact averaged outer endpoint from
\eqref{px:extended-coarse-difference}.

\begin{lemma}[Completion of the outward rows]
\label{lem:fms-outward-completion}
For every fixed $r$,
\begin{equation}\label{fms:outward-equivalence}
                         C_{L,r}^{-1}J_r\le D_r\le C_{L,r}J_r.
\end{equation}
The constants have no depth or volume factor. The comparison is
local: a graph edge uses only its associated determining row and
the smaller internal rows in its outer endpoint block.
\end{lemma}

\begin{proof}
An equal-level graph edge is precisely an interior determining
row, so the two differences coincide, with their length factor.
For a face between a level-$j$ pillar at $y$ and smaller pillars,
put $s=\xi_j$. The determining outward bond from $y$ ends at the
level-$j$ grid point $z$ across that face. Its entire block
$B^j(z)$ is outside $\Omega_j$ and inside $\Omega_{j-1}$.
It consists of $L^4$ assigned level-$(j-1)$ pillars. Every
level-$(j-1)$ bond internal to this block is an actual prescribed
interior row. This is the no-inward-row geometry; it is not true
for an arbitrary clipped block.

Choose a parallel frame on the two-block rectangle. Write $a$ for
the value at $y$, $b_w$ for the smaller top values in $B^j(z)$,
and $\bar b=L^{-4}\sum_wb_w$. The outward difference is
$s^{-1}(a-\bar b)$. For every smaller pillar on the face,
\[
 |a-b_w|^2\le2|a-\bar b|^2+2|b_w-\bar b|^2,
 \qquad
 \sum_w|b_w-\bar b|^2\le C L^2
                    \sum_{\substack{w\sim w'\\ \mathrm{internal}}}|b_w-b_{w'}|^2.
\]
The last inequality is the finite $L$-grid path Poincar\'e
inequality. Multiply by $s^{2-2r}$ and sum over at most $L^3$
face neighbors. The smaller row weights differ by the fixed
factor $L^{2-2r}$. This bounds the corresponding graph edges by
the outward row and the smaller internal rows.

Conversely choose one smaller face neighbor $w_0$. Write
$a-\bar b=(a-b_{w_0})+L^{-4}\sum_w(b_{w_0}-b_w)$ and join
$w_0$ to $w$ by coordinate paths in the smaller block. Their
length is at most $4L$. Cauchy--Schwarz gives the reverse bound.
The same argument works after restoring the unitary parallel
transports.

A top belongs to only one such outer large block at the next
level. That block can be read by at most eight incident outward
rows. Internal smaller edges and all the constructed local paths
therefore have bounded overlap depending only on $L$. Summing
the local inequalities proves \eqref{fms:outward-equivalence}.
\end{proof}

In particular the kernel of $\bar D_{\rm ext}$ is the space of
parallel sections of this pillar adjacency graph. Such a section
lifts to a parallel fine site field by the actual own-pillar
transport, and conversely every parallel fine field gives such
a section. For a fine edge crossing a pillar face, this follows
in the same two-block rectangle used in the proof. Thus the
averaged outward endpoint does not introduce an extra coarse
gauge kernel. This statement retains the actual holonomy
conditions on each connected component.

\subsubsection{A local box estimate with its actual boundary adjoint}

\begin{lemma}[Local cochain Hodge decomposition]
\label{lem:fms-box-hodge}
Let $D$ be a convex rectangular fine box of diameter comparable to
$s\ge\eta$, with bounded aspect ratio. For any one-cochain $v$
on all its inner bonds there is a site potential $\lambda$ and
a remainder $r=v-D_V\lambda$ such that
\begin{equation}\label{fms:box-hodge-estimates}
 \|r\|_{2,D}\le Cs\|\mathrm{d}_Vv\|_{2,D},\qquad
 \|\nabla_Vr\|_{2,D'}\le C\|\mathrm{d}_Vv\|_{2,D},
\end{equation}
where $D'$ is any fixed smaller concentric box. The potential is
unique up to a parallel constant after its gradient is chosen.
It may be normalized to have zero transported average on a
specified centered pillar of scale comparable to $s$ in $D$.
The constants include only the fixed size ratios. They are
independent of the number of fine sites in the box.
\end{lemma}

\begin{proof}
Choose a parallel frame on $D$. Let $\mathrm{d}_0$ be the fine gradient
on the path-product graph and $\mathrm{d}_1$ its plaquette coboundary.
Project $v$ orthogonally onto $\operatorname{ran}\mathrm{d}_0$ and write
$r=v-\mathrm{d}_0\lambda$. The projection equation is
$\mathrm{d}_0^\sharp r=0$, with the full graph adjoint, including its
boundary vertices. Also $\mathrm{d}_1r=\mathrm{d}_1v$.

Here the required graph Hodge estimate is elementary and includes
the boundary condition. In one direction with $N+1$ vertices,
$\mathrm{d}_0^\sharp \mathrm{d}_0$ is the path Neumann matrix, whereas
$\mathrm{d}_0\mathrm{d}_0^\sharp$ on its $N$ edges is the tridiagonal matrix with
diagonal $2\eta^{-2}$ and off-diagonal $-\eta^{-2}$.
Their nonzero eigenvalues are
\[
              4\eta^{-2}\sin^2\frac{\pi m}{2(N+1)}
                         \quad(1\le m\le N).
\]
The scalar matrix also has its constant zero mode. Tensor products
give the one-cochain Hodge matrix
$\mathrm{d}_0\mathrm{d}_0^\sharp+\mathrm{d}_1^\sharp \mathrm{d}_1$: in component $\mu$ its
longitudinal factor is the edge matrix and its other factors are
the scalar path matrices. Every component therefore has a
nonzero longitudinal frequency and a gap at least $c s^{-2}$.
Its quadratic form is the sum of the interior component first
differences, with additional nonnegative longitudinal boundary
edge terms. This follows either from those tensor factors or by
expanding their squares.

It follows that
$\|r\|_2\le Cs\|\mathrm{d}_1r\|_2$, and that every first difference
whose stencil is inside $D'$ is bounded by $C\|\mathrm{d}_1r\|_2$.
This proves \eqref{fms:box-hodge-estimates}.
Changing the parallel constant in $\lambda$ does not change $r$.
The transported average of a parallel constant is that constant,
so it gives the stated normalization. The local Poincar\'e
inequality with a mean on the specified positive-volume pillar
has a constant depending only on the size ratio. It includes a
bare one-site pillar, whose physical volume is then $s^4$ up to
that same fixed ratio. Restore the parallel frame to finish.
\end{proof}

The divergence in this lemma is the full adjoint on the finite
box graph. An interior formula that discards its boundary vertex
equations would allow additional harmonic gradients and would
not prove the lemma.

\subsubsection{Adaptive patches and removal of the bare boundary collar}

Let $C_{\rm row}$ be a fixed physical carrier-radius constant for
the flat rows and their completed outer blocks. Fix an integer
$K\ge64L^2(C_{\rm row}+1)$ and require $\rho\ge C K L$.
A box of radius $K\xi_y$
about an assigned top can meet at most two adjacent assigned
levels. Indeed it cannot reach level $j-2$ from a level-$j$ top
without crossing distance $\rho\xi_j/L$, nor level $j+2$
without crossing distance $\rho L\xi_j$. It cannot meet both
levels $j-1$ and $j+1$, because those sets are separated by
more than $\rho\xi_j$. Increase the fixed constant in the
separation for the enlarged boxes used below.

Consequently these boxes have uniformly bounded overlap, and
there is a partition $\sum_y\chi_y=1$ on $\Omega_0$ with
\begin{equation}\label{fms:adaptive-partition}
 \operatorname{supp}\chi_y\subset Q_{2\xi_y}(y),\quad
 |D\chi_y|\le C_L\xi_y^{-1},\quad
 |D^2\chi_y|\le C_L\xi_y^{-2}.
\end{equation}
One construction uses a fixed smooth bump equal to one on a
centered pillar, supported in its doubled cube, and divides by
the sum of all such bumps. Pillars are disjoint; every overlapping
scale has ratio at most $L$. Hence that sum is bounded above and
below and its first two differences have the stated bounds.
The same counting proves bounded overlap for the larger
$K\xi_y$ boxes and the pair neighborhoods below. Whenever cutoff
overlaps are used, include their two-fine-step collars as well;
the same estimates apply since $\eta\le\xi_y$.

We first remove a genuine boundary difficulty. Choose $H>10K$
fixed and a scalar cutoff $\theta$, equal to one within
$H\eta$ of $\Omega_0^c$, zero beyond $2H\eta$, with
$|D\theta|\le C\eta^{-1}$. Choose $\rho>4H+C_L$.
Every bond read by $v_b=\theta v$ is a bare determining row,
and $Q_0v=v$ on these bonds. The elementary finite-difference
bound gives
\begin{equation}\label{fms:boundary-removal}
 \|v_b\|_X\le C\|Qv\|_Y,\quad
 \|\mathrm{d}_V(v-v_b)\|_2+\|Q(v-v_b)\|_Y
                        \le C(\|\mathrm{d}_Vv\|_2+\|Qv\|_Y).
\end{equation}
For the curl use the product rule; its cutoff term has size
$C\eta^{-1}|v|$ in the bare collar. For $Q$ use its bounded
physical row and column sums, or note that the higher row
carriers miss this collar. The norm $X$ includes zero-extension
boundary jumps, also bounded by $\eta^{-1}|v|$ there.

Put $w=v-v_b$. It vanishes in this boundary collar. Every
nontrivial box used to decompose $w$ is therefore inside
$\Omega_0$: nonbare boxes are inside by separation, and a
bare box meeting the exterior lies entirely in the zero collar.
For such a zero box choose potential and remainder zero.
This convention never requires a flat extension across a hole.

\subsubsection{Gluing the normalized local potentials}

\begin{proposition}[A regular decomposition with the actual averages]
\label{prop:fms-regular-decomposition}
For every field $v$ on the true active bonds there is a site
potential $\lambda_0$ with $q\lambda_0=0$ and a field $r_0$ such
that
\begin{equation}\label{fms:regular-decomposition}
 v=D_V\lambda_0+r_0,\qquad
 \|r_0\|_X\le C(\|\mathrm{d}_Vv\|_2+\|Qv\|_Y).
\end{equation}
The constant is independent of all assigned levels and the
topology of the determining region.
\end{proposition}

\begin{proof}
First suppose every ambient period exceeds $8K\max_y\xi_y$,
so all local boxes below embed. A finite-cover argument at the
end removes this temporary restriction.
It suffices by \eqref{fms:boundary-removal} to decompose $w$.
For every nonzero patch apply \cref{lem:fms-box-hodge} on
$D_y=Q_{K\xi_y}(y)$, choosing
\begin{equation}\label{fms:local-normalized-potentials}
 w=D_V\lambda_y+r_y,
 \quad q_y\lambda_y=0,
 \quad\|r_y\|_{2,D_y}\le C_L\xi_y\|\mathrm{d}_Vw\|_{2,D_y},
 \quad\|\nabla_Vr_y\|_{2,D_y/2}\le C_L\|\mathrm{d}_Vw\|_{2,D_y}.
\end{equation}
The local normalization uses the actual pillar average.

We spell out the comparison of constants on overlaps. If the
cutoff supports, including the stated fine collars, meet at $y,z$,
their scales have ratio at most $L$ and their centers have
distance at most $8\sigma_{yz}$. A coordinate route and every
completed outer block lie in a rectangle of radius
$C(C_{\rm row}+1)\sigma_{yz}$ about these centers, with the
dimensional constant absorbed by the choice of $K$ above.
Since the smaller radius of $D_y/2,D_z/2$ is at least
$K\sigma_{yz}/(2L)$, this rectangle lies in their intersection.
Choose a convex pair neighborhood there
containing the two central pillars and the following finite row
route. Follow a coordinate fine path from one pillar to the
other and record the pillars it crosses. There are at most
$C_L$ of them because their scales are comparable. Replace each
unequal-level face crossing by its determining outward row and
the internal smaller rows in its full outer block, as in
\cref{lem:fms-outward-completion}. All these carriers fit in the
pair neighborhood when $K$ is chosen large enough.

For $f_a=q_a\lambda_z$, the local intertwining gives
$\bar D_{\rm ext}f=Q(w-r_z)$ on that route. Since $f_z=0$,
the finite-route version of the completion lemma implies
\begin{equation}\label{fms:local-constant-comparison}
 |q_y\lambda_z|
 \le C_L\left\{\sum_{c\in\mathcal C_{yz}}\xi_c^2|Qw(c)|^2
                    +\|\mathrm{d}_Vw\|_{2,D_y\cup D_z}^2\right\}^{1/2}
                         =:C_L e_{yz}.
\end{equation}
For the remainder term used here, a physical flat row has
$L^2$ functional norm at most $C\xi_c^{-2}$: its coefficient
sum is bounded and its largest fine coefficient is at most
$C(\eta/\xi_c)^4$. Thus
$\xi_c|Qr_z(c)|\le C_L\|\mathrm{d}_Vw\|_{2,D_z}$, by
\eqref{fms:local-normalized-potentials}. Only $C_L$ rows enter.

On the pair neighborhood let $h_{yz}=\lambda_y-\lambda_z$.
Its derivative is $r_z-r_y$ and its average at $y$ is
$-q_y\lambda_z$. Poincar\'e with that central pillar mean gives
\begin{equation}\label{fms:potential-overlap-bounds}
 \|h_{yz}\|_2\le C_L\sigma_{yz}^2e_{yz},
 \qquad \|D_Vh_{yz}\|_2\le C_L\sigma_{yz}e_{yz}.
\end{equation}
In particular the constants of the primitives have been
estimated; a bound only on their gradient would not suffice.
The pair neighborhoods and their routes have bounded overlap, so
\begin{equation}\label{fms:overlap-energy-sum}
              \sum_{y,z}e_{yz}^2
                      \le C_L(\|\mathrm{d}_Vw\|_2^2+\|Qw\|_Y^2),
\end{equation}
where only intersecting cutoff supports are summed.

Glue $\lambda=\sum_y\chi_y\lambda_y$ and put $r=w-D_V\lambda$.
The exact product rule writes $r$ as the sum of the cutoff
remainders $\chi_y r_y$ and terms $D\chi_y$ times shifted
primitive differences. The latter are differences because
$\sum_yD\chi_y=0$. A second product rule for its gradient,
Cauchy--Schwarz and bounded overlap give
\begin{equation}\label{fms:gluing-product-bounds}
 \begin{split}
 \|r\|_X^2\le C_L\sum_y\bigl(
     \|\nabla_Vr_y\|_{2,D_y/2}^2+\xi_y^{-2}\|r_y\|_{2,D_y}^2\bigr)\\
 {}+C_L\sum_{y,z}\bigl(
     \sigma_{yz}^{-4}\|h_{yz}\|_2^2
                 +\sigma_{yz}^{-2}\|D_Vh_{yz}\|_2^2\bigr).
 \end{split}
\end{equation}
All shifts in this formula remain in the specified interior
boxes. Equations \eqref{fms:local-normalized-potentials} through
\eqref{fms:overlap-energy-sum} prove
$\|r\|_X\le C_L(\|\mathrm{d}_Vw\|_2+\|Qw\|_Y)$.

The glued potential need not have zero assigned average. On
$B_y$, subtract $\lambda_y$ in its average and use its zero
normalization. The averaging functional has norm $\xi_y^{-2}$,
so \eqref{fms:potential-overlap-bounds} gives
\begin{equation}\label{fms:glued-average-bound}
              \sum_y|q_y\lambda|^2
                       \le C_L(\|\mathrm{d}_Vw\|_2^2+\|Qw\|_Y^2).
\end{equation}
This is a counting norm of the site values, with no missing
physical volume factor: the factor $\xi_y^{-2}$ cancels the
$\xi_y^2$ in the potential overlap estimate.

At a flat field the exact bump lift of the site averages can be
chosen with the additional estimate
\begin{equation}\label{fms:flat-lift-second-difference}
             qE=1,\qquad
                  \|D_VE\psi\|_X^2\le C_L\sum_y|\psi_y|^2.
\end{equation}
To verify it, use the same normalized compact scalar profiles in
parallel frames. Their second differences are
$O(\xi_y^{-2}|\psi_y|)$ and their physical volume is
$O(\xi_y^4)$. The weighted first differences have that same
squared counting bound. For a bare seed the two differences cost
at most $C\eta^{-2}$, canceled by its physical volume
$\eta^4$ after squaring. The two-step collars have bounded
overlap. A large finite cover and deck averaging handle any
nonembedded surrounding box, preserving these energy estimates.

Set $\lambda_0=\lambda-Eq\lambda$. Its average is zero exactly.
Add $D_VEq\lambda$ to $r$ and then add back $v_b$ to the
remainder. Equations \eqref{fms:glued-average-bound},
\eqref{fms:flat-lift-second-difference} and
\eqref{fms:boundary-removal} prove
\eqref{fms:regular-decomposition} on this large volume.

For an arbitrary smaller compatible torus, pull back the entire
determining geometry, flat connection and input field to a finite
periodic cover whose periods exceed $8K\max_y\xi_y$.
Construct the potential $\lambda_0$ there by the preceding proof.
Average it over the deck group. The lifted input is invariant,
and the actual $D_V,q,Q$, all physical weights and the boundary
collar commute with these translations. The averaged potential
still has $q\lambda_0=0$. Its remainder is the average of the
translated remainders, so its $X$ norm obeys the same bound by
the triangle inequality. It descends to the original torus.
The squared remainder norm and both squared input norms multiply
by the number of sheets under pullback; these factors cancel.
This proves the full regular decomposition at every compatible
volume, not only the auxiliary site lift. The argument does not
require trivial holonomy on the original active complex.
\end{proof}

The argument uses local primitives on convex boxes and compares
their constants through the actual determining rows. It therefore
also controls closed fields with periods around holes: their
constants can be glued only if the corresponding row information
is retained. No assertion that the region has trivial first
cohomology was made.

\subsubsection{The flat layered Landau gap}

\begin{proposition}[Flat coercivity on the true determining region]
\label{thm:fms-flat-landau-gap}
For the separated aligned determining geometry, with separation
larger than the fixed constants above, there are $c,C>0$ such that
\begin{equation}\label{fms:flat-landau-gap}
 c\|v\|_X^2\le
       \|\mathrm{d}_Vv\|_{2,\Omega_0}^2+\|RD_V^\sharp v\|_2^2+\|Qv\|_Y^2
                                             \le C\|v\|_X^2.
\end{equation}
This is uniform in depth, compatible volume, region topology
and real flat holonomies. The curl includes exactly the inner
plaquettes, the divergence is the adjoint on the true active
graph, and $R,Q$ are the actual assigned-site projector and
straightened determining rows.
\end{proposition}

\begin{proof}
Take the regular decomposition of
\cref{prop:fms-regular-decomposition}. Since $q\lambda_0=0$,
the actual range identity gives
\[
 L_V\lambda_0=R L_V\lambda_0
                    =RD_V^\sharp v-RD_V^\sharp r_0.
\]
The divergence is bounded by the $X$ norm. Apply
\cref{prop:lsi-zero-average-inverse} to conclude
\[
 \|D_V\lambda_0\|_X
       \le C(\|RD_V^\sharp v\|_2+\|r_0\|_X).
\]
Add the remainder bound. This proves the lower estimate in
\eqref{fms:flat-landau-gap}. The upper estimate follows from
the local curl and divergence bounds and the physical row and
column bounds for $Q$. The boundary part of $X$ is a comparison
norm only; no boundary plaquettes or crossing action bonds have
been added in deriving its lower bound.
\end{proof}

At a flat field the physical Wilson Hessian is precisely the
inner curl form. Hence the theorem proves the full flat
projected-Landau precision gap on the actual multi-level region.
Together with the actual-slice inf-sup estimate in
\eqref{lsi:slice-infsup}, it gives the flat constrained saddle
inverse in the declared energy and dual-energy norms. The passage
to a nonflat regular field still requires coefficient comparison
and localization of this same geometry. It is not a consequence
of a site gap alone.

% END LAYERED EXPANSION inverse/flat_mixed_scale_decomposition.tex

% BEGIN LAYERED EXPANSION inverse/nonflat_proper_landau_precision.tex
\subsection{A nonflat Landau inverse on the determining regions}
\label{ch:nonflat-proper-landau}

We retain the same determining geometry in every comparison.
A local comparison field equals the identity outside its gauge
patch, but it has the original assigned averages, determining
rows and inner action region. No periodic inner operator is
substituted at a seam or an action boundary.

Use the hypotheses and physical spaces of
\cref{ch:layered-site-restoration,ch:flat-mixed-scale-decomposition}.
All row weights are positive and comparable with the assigned
inverse scale squared. The estimates concern the unweighted
inner Wilson action and actual straightened constraint.
Signed profiles and separately specified current completions
are not included by a change of notation.

\subsubsection{A common weighted gauge}

Let \(V\) be real and flat on the active graph and write
\begin{equation}\label{npl:common-gauge}
 U_b=e^{i\eta A_b}V_b,\qquad
                  \varepsilon=\sup_b\mathsf L(b)|A_b|.
\end{equation}
The reference may have nontrivial holonomy around holes or torus
cycles. Site and row spaces, \(\Xi\), and \(\ell\) are the same
for both backgrounds.

\begin{lemma}[Actual weighted coefficient comparison]
\label{lem:npl-coefficients}
For sufficiently small fixed \(\varepsilon\) and curvature ceiling,
\begin{equation}\label{npl:coefficient-comparison}
 \begin{split}
 \|\ell(D_U-D_V)u\|_2&\le C\varepsilon\|u\|_2,\\
 \|(q_U-q_V)u\|_{\rm top}
 +\|\Xi^{-2}(q_U-q_V)(\ell^2u)\|_{\rm top}
                                      &\le C\varepsilon\|u\|_2,\\
 \|a_f^{1/2}(Q_U-Q_V)v\|_{2,\rm rows}
                           &\le C\varepsilon\|\mathsf L^{-1}v\|_2,\\
 \|(F_U-F_V)b\|_2
 +\|\ell(D_UF_U-D_VF_V)b\|_2
                                      &\le C\varepsilon\|b\|_{\rm top},\\
 \|R_U-R_V\|_{2\to2}&\le C\varepsilon,\\
 \|(R_UD_U^\sharp-R_VD_V^\sharp)v\|_2
                                      &\le C\varepsilon\|v\|_{X_V}.
 \end{split}
\end{equation}
The \(Q\)'s are the actual straightened rows, including their
node derivatives and extended outward endpoint averages.
\end{lemma}

\begin{proof}
The first bound follows from
\(\|\operatorname{Ad}(e^{i\eta A})-1\|\le C\eta|A|\).
Each actual row carrier meets at most two adjacent assigned
scales. In a local parallel frame for \(V\), intermediate
transports at physical scale \(\xi_s\) change by at most
\(C_L\varepsilon\xi_s/\xi_c\). Their geometric sum is bounded
uniformly in the number of averaging levels. Differentiating
the finite actual recursion gives both physical Schur sums for
the field-row difference. This is the row comparison proved in
\cref{ch:layered-current-precision}, including the straightening
factors. The site-row comparison is the same transported-average
calculation. Inserting \(\ell^2/\xi_y^2\), or the field mass
weights on a two-scale carrier, preserves its \(O_L(\varepsilon)\)
bound. Thus the first three lines have no fine-density loss.

For completeness, the normalized resolvent comparison uses the
weighted energy, rather than an unscaled fine matrix norm.
Put \(u=(F_U-F_V)b\). Then
\begin{equation}\label{npl:column-difference-equation}
 A_{s,U}u=(q_U^\sharp-q_V^\sharp)\Xi^{-2}b
                         -(A_{s,U}-A_{s,V})F_Vb.
\end{equation}
Test against \(\ell^2u\). The real part on the left controls
\(\|\ell D_Uu\|_2^2+\|q_Uu\|_{\rm top}^2\), and this energy
controls \(\|u\|_2^2\), by the proof of
\cref{prop:lsi-gram-gap}. The first right-hand term is bounded
by \(C\varepsilon\|b\|_{\rm top}\|u\|_2\).
In the gradient-form difference, the factor \(D_U-D_V\)
has size \(C\varepsilon\ell^{-1}\), while
\(D_U(\ell^2u)\) is bounded by
\(C\ell(|\ell D_Uu|+\delta|u|)\).
The alternative placement of that difference is bounded by
\(C\varepsilon\|\ell D_VF_Vb\|_2\|u\|_2\).
The two site-mass differences are bounded by
\(C\varepsilon(\|F_Vb\|_2+\|q_VF_Vb\|_{\rm top})\|u\|_2\).
The weighted column energy bounds therefore give
\(\|u\|_2+\|\ell D_Uu\|_2\le C\varepsilon\|b\|_{\rm top}\).
This proves the fourth line after adding the gradient coefficient
difference.

Consequently \(M_U-M_V=O(\varepsilon)\) in normalized top norm.
The Gram gaps give the same bound for the inverse difference.
Insert these in \(R=1-FM^{-1}F^\sharp\).
For the last line insert and subtract \(R_UD_V^\sharp v\).
The remaining adjoint coefficient difference is bounded by
\(C\varepsilon\|\mathsf L^{-1}v\|_2\).
No derivative of the potential is used.
\end{proof}

Pair the Hessian of the actual per-length nonlinear constraint by
\begin{equation}\label{npl:constraint-form}
 B_\omega(v,w)=\sum_{c\in\mathbf B}\xi_c^4
 \langle\omega(c),D^2\mathcal Q_{U,c}(0)[v,w]\rangle,\qquad
                    \gamma_\omega=\sup_c\xi_c^3|\omega(c)|.
\end{equation}
Its actual two-leg moment bounds it by
\(C\gamma_\omega\|\mathsf L^{-1}v\|_2\|\mathsf L^{-1}w\|_2\).
Define the physical form
\begin{equation}\label{npl:physical-form}
 \begin{split}
 K_{U,\omega}(v,w)={}&D^2A_{\Omega_0}(U)[v,w]-B_\omega(v,w)\\
 &+\langle R_UD_U^\sharp v,R_UD_U^\sharp w\rangle
                       +\langle Q_Uv,a_fQ_Uw\rangle_{\rm rows}.
 \end{split}
\end{equation}
Action derivatives use physical left variations \(e^{i\eta tv}U\).

\begin{corollary}[Coercivity in the common gauge]
\label{cor:npl-common-gap}
There are fixed \(c_1,\epsilon_1>0\) such that
\par\noindent
\begin{minipage}{\linewidth}
\begin{equation}\label{npl:common-gap}
 K_{U,\omega}(v,v)\ge c_1\|v\|_{X_U}^2
 \quad\hbox{if}\quad
                    \varepsilon+\kappa+\gamma_\omega\le\epsilon_1.
\end{equation}
\popQED\par
\end{minipage}
\end{corollary}

\begin{proof}
The exact Wilson plaquette Hessian differs from its transported
linearized curl form by \(C\kappa\) times the mass norms. The
curl differs from its flat reference by \(C\varepsilon\) times
the mass norm: each prefix transport changes by
\(O(\eta\varepsilon/\mathsf L)\) before division by \(\eta\).
This is a local inner-plaquette calculation, including the action
boundary. Thus
\[
 |D^2A_{\Omega_0}(U)[v,w]-\langle \mathrm{d}_Vv,\mathrm{d}_Vw\rangle|
       \le C(\varepsilon+\kappa)\|v\|_{X_V}\|w\|_{X_V}.
\]
Use \cref{lem:npl-coefficients} for the gauge and row squares
and the two-leg estimate for the actual constraint form. Absorb
the total error in \cref{thm:fms-flat-landau-gap}.
The two component-gradient norms are equivalent by their
coefficient comparison. This proves \eqref{npl:common-gap}.
\end{proof}

\subsubsection{Weighted cell kernels and IMS}

Let \(\psi\) be bounded and one-Lipschitz for
\(d_\ell=\inf\int ds/\ell\), with its sampled top values.
Conjugate site and top vectors by \(e^{s\psi}\). A derivative
commutator has size \(C_L|s|\ell^{-1}\), and a normalized site
row changes by \(C_L|s|\). Testing the conjugated equation
against \(\ell^2u\) perturbs the coercive weighted-energy form
by \(C_L|s|\). The preceding weighted proof consequently gives,
for a fixed \(|s|\le s_0\),
\begin{equation}\label{npl:weighted-columns}
 \|e^{s\psi_f}Fe^{-s\psi_t}\|_{2\to2}
       +\|e^{s\psi_f}\ell D_UFe^{-s\psi_t}\|_{2\to2}\le C.
\end{equation}
The difference from the original column is \(O(|s|)\).
Therefore the conjugated normalized Gram matrix changes by
\(O(|s|)\), and its inverse stays bounded by the Gram gap.
It follows that the normalized smoothing operator
\begin{equation}\label{npl:normalized-smoothing}
 \widetilde K_s=\ell D_UPD_U^\sharp\ell
        =(\ell D_UF)M^{-1}(\ell D_UF)^\sharp,\qquad P=1-R,
\end{equation}
has exponential physical cell-\(L^2\) entries. This is not a
pointwise fine-kernel assertion.

The actual layer metric has cell growth bounded by
\(C_L(1+t)^5e^{8t\log L/\rho}\).
Since \(d_\ell\ge d\), choose the separation so that
\begin{equation}\label{npl:rate-budget}
                               8\log L/\rho<s_0/2.
\end{equation}
All fixed polynomial row and column moments of the smoothing
entries are then bounded. This explicitly pays for the growth
of the changing-scale metric.

Choose a large fixed \(R\). A useful partition can be constructed
without assumptions on boundary density. Select tops greedily
in decreasing scale so that balls of radius \(R\xi_y/(20L)\)
are disjoint and the selection is maximal. If
\(\rho\ge C_LR\), intersecting enlarged balls have only
adjacent scales. The larger balls of radius \(R\xi_y\)
cover the region, allowing a fixed enlargement for fine cells.
Packing the smaller disjoint balls gives overlap at most \(C_L\).
Take bumps equal to one on the covering balls and supported in
their doubles, and normalize their squared sum. They satisfy
\begin{equation}\label{npl:square-partition}
 \begin{split}
 &\sum_i\chi_i(x)^2=1,\qquad
   \sum_i|D\chi_i(x)|^2\le C_LR^{-2}\xi(x)^{-2},\\
 &\sum_i|\chi_i(x)-\chi_i(z)|^2
                  \le C_L\min\{1,d_\ell(x,z)^2/R^2\},\\
 &\sum_i\sup_{x\in B_y}|\chi_i(x)-\chi_i(y)|^2\le C_LR^{-2}.
 \end{split}
\end{equation}
The pair difference follows by integrating the vector of cutoff
gradients along a metric path. The supremum bound follows from
bounded overlap and the bounded metric diameter of a pillar.

\begin{lemma}[IMS for the actual layered form]
\label{lem:npl-ims}
For the true inner form,
\begin{equation}\label{npl:ims}
 \left|\sum_iK_{U,\omega}(\chi_iv,\chi_iv)-K_{U,\omega}(v,v)\right|
                         \le CR^{-2}\|\mathsf L^{-1}v\|_2^2.
\end{equation}
The same error bounds the difference between
\(\sum_i\|\chi_iv\|_X^2\) and \(\|v\|_X^2\).
\end{lemma}

\begin{proof}
Split the gauge form as
\(\|D_U^\sharp v\|_2^2-\langle v,D_UPD_U^\sharp v\rangle\).
For the smoothing term put \(u=\ell^{-1}v\) and use
\eqref{npl:normalized-smoothing}. The exact double-cutoff
expansion \eqref{loc:double-cutoff-entry} bounds its cell error
by \(CR^{-2}(1+d_\ell(y,z))^2\) times the entry norm, using
the within-cell and pair estimates in \eqref{npl:square-partition}.
The exponential moments give \(CR^{-2}\|u\|_2^2\) by Schur.

The remaining divergence matrix has fine range \(2\eta\) and
physical coefficients \(C\eta^{-2}\). Its squared cutoff
difference is \(C\eta^2/(R^2\mathsf L^2)\), cancelling the fine
denominator. The Wilson Hessian is a sum of the same kind of
finite plaquette matrices, including at its true boundary.
The row-mass and actual constraint-curvature matrices have
bounded mass-normalized physical Schur sums and bounded carrier
diameter in the layer metric. Their double-cutoff errors have
the same mass bound. This proves \eqref{npl:ims}.
The local unitary-shift calculation gives the \(X\) identity.
A zero-extension boundary jump has only one nonzero field
value, so the identity \(\sum_i\chi_i^2=1\) makes its contribution
exact. No exterior cutoff extension or action bond is required.
\end{proof}

\subsubsection{Shielding on the same determining geometry}

\begin{lemma}[Localized comparison of the actual projector]
\label{lem:npl-shielding}
Let two admitted real backgrounds have the same determining
geometry and agree in common local frames on
\(Q_{24R\xi_0}\). Suppose all assigned scales there are comparable
to \(\xi_0\) within a fixed factor depending on \(L\).
For its core \(Q_{2R\xi_0}\),
\begin{equation}\label{npl:shielding}
       \|1_{\rm core}(R_1-R_2)1_{\rm core}\|_{2\to2}\le Ce^{-cR}.
\end{equation}
Only the current-free site and normalized Gram package is used.
\end{lemma}

\begin{proof}
Use a cutoff supported strictly inside the agreement cube in
\begin{equation}\label{npl:geometric-resolvent}
 \theta G_1-G_2\theta
                 =G_2(A_{s,2}\theta-\theta A_{s,1})G_1.
\end{equation}
The middle operator is a gradient or actual site-row commutator
in a buffer annulus. Its energy-to-dual-energy norm has the
physical scale. Weighted site estimates make the two solutions
exponentially small there. The compressed fine inverse difference
is consequently \(C\xi_0^2e^{-cR}\), with the corresponding
localized energy-derivative bound. No exterior connection
coefficients have been subtracted.

Apply this comparison to \(F=Gq^\sharp\Xi^{-2}\).
All local scale factors are comparable to \(\xi_0\), so its
inverse-length factors cancel the site-inverse scale.
The normalized cell-\(L^2\) comparison is \(Ce^{-cR}\).
Compress the Gram matrices to an intermediate top box. Their
difference is \(Ce^{-cR}\) in operator norm, by splitting
the \(F\) outputs into the common interior and their tails.
Both compressions retain the uniform positive gap.

Their inverse difference is therefore \(Ce^{-cR}\).
The exact compression-inverse identity
\eqref{loc:compression-inverse-identity}, applied to these
normalized matrices with the physical top measures, compares
each compressed inverse with its global inverse. The boundary
term is exponentially small by the weighted Gram estimate.
Insert this and the \(F^\sharp\) source tails into
\(P=FM^{-1}F^\sharp\). This proves \eqref{npl:shielding}.
These are operator estimates on arbitrary supported vectors,
not sums of separate errors over a patch volume.
\end{proof}

\subsubsection{Local fields and the bare action boundary}

For an interior patch of scale \(\xi_0\), choose a comb frame
on \(Q_{32R\xi_0}\). At most two adjacent scales occur there
when \(\rho\ge C_LR\). The local potential is bounded by
\(C_LR\kappa\xi_0^{-1}\). Multiply it by a cutoff equal to
one on \(Q_{24R\xi_0}\), vanishing farther out while still
inside \(\Omega_0\). On the original active graph put
\(U_i^{\rm ext}=e^{i\eta A_i^{\rm ext}}\), and put the identity
connection elsewhere. The local logarithm and product rules give
\begin{equation}\label{npl:cutoff-field-bounds}
 \sup_b\mathsf L(b)|A_i^{\rm ext}(b)|\le C_LR\kappa,\qquad
 \sup_p\mathsf L(p)^2|F_i^{\rm ext}(p)|
                              \le C_L(\kappa+R^2\kappa^2).
\end{equation}
The cutoff derivative contributes
\((R\xi_0)^{-1}(R\kappa\xi_0^{-1})\), with the required
curvature units. All its averages use the original geometry.

A nonbare center has the surrounding cube inside \(\Omega_0\)
by separation. For a bare center either it is also interior,
or its whole field support is in a bare outer collar.
In that case \(Q_0v_i=v_i\), no higher row carrier reaches its
support, and its nonlinear constraint Hessian vanishes.
The true inner Wilson lower bound and positive gauge give
\begin{equation}\label{npl:boundary-gap}
 K_{U,\omega}(v_i,v_i)\ge(d_*-C\kappa)\eta^{-2}\|v_i\|_2^2
                                                \ge c_b\|v_i\|_X^2.
\end{equation}
Here \(d_*=\inf_c\xi_c^2a_f(c)>0\) is the lower field-row weight.
The last inequality uses only the fine finite-difference bound
\(\|\nabla v_i^0\|_2\le C\eta^{-1}\|v_i\|_2\).
This is a boundary estimate for the actual inner action.

\subsubsection{The real all-layer precision and saddle inverse}

\begin{proposition}[Uniform nonflat precision on the determining region]
\label{thm:npl-proper-gap}
There are fixed positive constants \(\rho_*\), \(\kappa_*\),
\(\gamma_*\), \(c\) and \(C\) such that,
for every depth, compatible volume and separated aligned
determining geometry with \(\rho\ge\rho_*\), a real field with
\begin{equation}\label{npl:final-assumptions}
 \sup_{p\subset\Omega_0}\mathsf L(p)^2|F_U(p)|\le\kappa_*,
 \qquad \sup_{c\in\mathbf B}\xi_c^3|\omega(c)|\le\gamma_*
\end{equation}
satisfies
\begin{equation}\label{npl:proper-gap}
                     c\|v\|_X^2\le K_{U,\omega}(v,v)\le C\|v\|_X^2.
\end{equation}
The constants depend only on \(L\) and the positive weight
bounds. The form uses the actual inner Wilson Hessian,
projected gauge and nonlinear straightened-constraint curvature.
\end{proposition}

\begin{proof}
Initially pass to a finite periodic cover large enough for the
patches. Fix the initial curvature ceiling for the site estimates.
It gives fixed normalized column and Gram constants and a rate
\(s_0\). The flat gap gives a fixed lower constant.
Choose \(R\) large enough that \(CR^{-2}\) and \(Ce^{-cR}\)
are smaller than a quarter of the common-gauge lower constant.
Choose \(\rho\) to satisfy the flat patch conditions,
\eqref{npl:rate-budget}, and every fixed \(C_LR\) collar
inequality. These prior constants are uniform above their
initial separation thresholds. Finally reduce
\(\kappa_*,\gamma_*\) so \eqref{npl:cutoff-field-bounds}
admits every local mean and meets \eqref{npl:common-gap}.
This order does not use the depth.

Use \eqref{npl:square-partition}. For an interior \(v_i=\chi_iv\),
change to its local comb frame and retain \(\omega,a_f\) on
every carrier meeting its support. Extend those coefficients
elsewhere with the same bounds. Fixed gauge changes conjugate
the actual means, their first and second derivatives and the
site projector by the matching unitary fine and top rotations.
The forms are unchanged. All local terms of the comparison field
therefore agree exactly on \(v_i\). Only \(R\) differs.

For a real orthogonal projector, \(\|Rf\|^2=\langle f,Rf\rangle\).
Apply \cref{lem:npl-shielding} to \(f=D_U^\sharp v_i\);
its support is inside the core. Thus
\[
 |K_{U,\omega}(v_i,v_i)-K_{U_i^{\rm ext},\omega_i}(v_i,v_i)|
                                      \le Ce^{-cR}\|v_i\|_X^2.
\]
The extension has a global small weighted potential relative
to the identity on the same active graph, so
\cref{cor:npl-common-gap} applies. Boundary patches have the
direct estimate \eqref{npl:boundary-gap}. Sum these local lower
bounds and use \cref{lem:npl-ims}; the chosen errors are absorbed.
The upper bound follows from the local Wilson coefficient
estimate, orthogonality of \(R\), physical row bounds and
the two-leg constraint moment.

On a smaller original torus the actual local rows, physical
transposes and \(\Xi\) commute with pullback to the finite cover.
Uniqueness of the gapped site equation and then the normalized
Gram equation makes their inverses preserve deck-invariant
subspaces. Thus the actual projector descends exactly.
The full form and the squared \(X\) norm have the same
sheet-count factor. Divide the cover estimate by that factor.
No global almost-flat gauge or topology restriction is used.
\end{proof}

Let \(X_L=\ker(R_UD_U^\sharp)\) and
\(K_L=D^2A_{\Omega_0}(U)-B_\omega\).
On \(X_L\cap\ker Q_U\), both positive penalties vanish exactly.
Coercivity there and the true-row inf-sup
\eqref{lsi:slice-infsup} prove the actual saddle estimate
\begin{equation}\label{npl:saddle-estimate}
 K_Lu-Q_U^\sharp\mu=f\ \hbox{on }X_L,\qquad Q_Uu=b,\qquad
 \|u\|_X+\|\mu\|_{Y^*}\le C(\|f\|_{X_L^*}+\|b\|_Y).
\end{equation}
Explicitly write \(u=E_Lb+z\), solve the coercive kernel
equation for \(z\), and define
\(\langle\mu,b'\rangle=(K_Lu-f)(E_Lb')\).
Every \(w\in X_L\) splits as
\(w-E_LQ_Uw+E_LQ_Uw\), proving the full equation.
Kernel coercivity and surjectivity give uniqueness.
This bounds a linearized saddle with its actual small multiplier;
it does not construct the critical background.

\subsubsection{Exponential energy kernels in the physical layer metric}

The energy inverse also has a quantitative spatial consequence.
Let \(G_f=K_{U,\omega}^{-1}\), and let \(1_y\) restrict fields
to an assigned cell. With \(d\) the original physical layer metric,
\begin{equation}\label{npl:field-cell-kernels}
 \begin{split}
 \|1_yG_f1_z\|_{2\to2}
                         &\le C\xi_y\xi_z e^{-m d(y,z)},\\
 \|1_y\nabla_UG_f1_z\|_{2\to2}
                         &\le C\xi_z e^{-m d(y,z)},\\
 \|1_yG_f\nabla_U^\sharp1_z\|_{2\to2}
                         &\le C\xi_y e^{-m d(y,z)},\\
 \|1_y\nabla_UG_f\nabla_U^\sharp1_z\|_{2\to2}
                         &\le C e^{-m d(y,z)}.
 \end{split}
\end{equation}
Here \(\nabla_U\) is the comparison component gradient defining
\(X\); its adjoint in the last two lines is the corresponding
typed force source. The norms include all physical cell measures.
These are cell operator estimates, not pointwise bounds at the
fine Green singularity.

To prove them, conjugate by \(w=e^{s\psi}\) for bounded
one-Lipschitz \(\psi\) in \(d_\ell\). A local derivative
commutator has size \(C|s|\ell^{-1}\), so the local Wilson and
divergence forms change by \(C|s|\) in the two \(X\) norms.
The normalized nonlocal term in
\eqref{npl:normalized-smoothing} has conjugated difference
\(C|s|\) in physical \(L^2\) norm, by the weighted column and
Gram resolvent arguments. The row and constraint terms have
bounded normalized Schur sums and fixed metric carrier radius;
their differences have that same \(C|s|\) bound. Thus
\begin{equation}\label{npl:weighted-field-form}
 \|wK_{U,\omega}w^{-1}-K_{U,\omega}\|_{X\to X^*}
                                                        \le C|s|.
\end{equation}
The real lower bound and a normalized Neumann series imply
\(\|wG_fw^{-1}\|_{X^*\to X}\le C\) for a fixed small \(s\).
Converting a force supported in \(B_z\) to the dual energy norm
costs at most \(C\xi_z\); converting output energy to \(L^2(B_y)\)
costs at most \(C\xi_y\). Choose \(\psi\) to be a truncated
distance from the source cell. Its within-cell oscillation is
bounded. This proves the first two lines of
\eqref{npl:field-cell-kernels}. A gradient-adjoint source has
dual norm bounded by its \(L^2\) norm, with no scale factor.
The same argument proves the last two lines, including the
single-step support enlargement of that source.

After increasing the fixed separation once more so that
\(8\log L/\rho<m/2\), the metric ball count makes their normalized
row and column sums finite. In particular, if
\(\sup_z\xi_z\|f\|_{2,B_z}\le B\), then
\begin{equation}\label{npl:cell-sup-response}
 \sup_y\left\{\xi_y^{-1}\|G_ff\|_{2,B_y}
                      +\|\nabla_UG_ff\|_{2,B_y}\right\}\le CB.
\end{equation}
Indeed the scale factors in the first two kernel lines cancel
these input and output amplitudes before the exponential cell
sum is taken. This is an energy response estimate uniform in
the number of fine bonds in each pillar.

\subsubsection{The scope of the current extension}

The explicit current bound \eqref{lsi:current-form-bound}
also makes
\begin{equation}\label{npl:explicit-current-form}
 \widetilde K[J,\omega]=D^2A_{\Omega_0}(U)-\mathcal B_U[J]
       +D_UR_UD_U^\sharp+Q_U^\sharp a_fQ_U
                             -S^\sharp\Delta^{(2)}[\omega]S
\end{equation}
coercive for sufficiently small
\(\sup_b\mathsf L(b)^3|J(b)|\) and
\(\sup_c\xi_c^3|\omega(c)|\).
First apply \eqref{npl:proper-gap} at zero multiplier and
then absorb the two explicit relative form estimates.
On \(X_L\), \(S\) is the identity and every term in
\(\mathcal B_U[J]\) has a zero horizontal gauge factor.
Its restriction to the constraint kernel is therefore the same
Lagrangian Hessian when \(\omega\) is the actual multiplier.

\Cref{prop:acc-current-identification} identifies this completion
with the specified canonical current correction.
\Cref{thm:scr-induced-current} supplies the induced readout, and
\cref{prop:acc-strong-propagators,dnr:linear-families,dnr:response}
supply the strong outputs, carrier families and nonlinear response.
The absolute soft contour precision retains its separate
gauge-restoration comparison.

% END LAYERED EXPANSION inverse/nonflat_proper_landau_precision.tex

% BEGIN LAYERED EXPANSION inverse/coarse_schur_current_readout.tex
\subsection{Coarse Schur kernels and the induced current}
\label{ch:layered-coarse-current-readout}

The current in the second explicit precision is obtained by
transposing a minimizing field lift. An energy bound for that
lift alone would charge the whole volume of a bounded current.
We prove exponential energy columns and sum them in the actual
layer metric. This gives the required weighted supremum bound
for the induced row multiplier without a fine pointwise Green
estimate.

Keep the real background and determining geometry of
\cref{thm:npl-proper-gap}. Currents below may be complex, but
the background is real. Every transpose in an analytic formula
is the physical bilinear transpose. Hilbert norms use the
complexification of the real physical spaces. Let
\begin{equation}\label{scr:first-precision}
 \begin{split}
 \mathcal K_0&=D^2A_{\Omega_0}(U)+D_URD_U^\sharp
                                             +Q^\sharp a_fQ,\\
 \mathcal K_J&=\mathcal K_0-\mathcal B_U[J],\qquad
                              G_J=\mathcal K_J^{-1},\\
 \gamma_J&=\sup_b\mathsf L(b)^3|J(b)|.
 \end{split}
\end{equation}
This is the displayed symmetric current extension, with
\(\mathcal B_U[J]\) defined by its endpoint formula. It is not
identified here with another named completion.

The field-row Hilbert space \(\mathcal H_\Gamma\) has measure
\(\xi_c^4\). Write \(\Xi_\Gamma b(c)=\xi_cb(c)\), so that
\(\mathcal U_\Gamma=\Xi_\Gamma^2\) is unitary from
\(\mathcal H_\Gamma\) to the counting \(\ell^2\) space.
This is a field-row scale. At an outward row it is the row's
assigned level, even when one endpoint is a top of the smaller
level. It is not the site-scale operator \(\Xi\).

\subsubsection{Weighted restoration for the current correction}

\begin{lemma}[Restoration with a spatial exponential weight]
\label{lem:scr-weighted-restoration}
For a fixed small \(s_*>0\), a bounded one-Lipschitz function
\(\psi\) for \(d_\ell\), and \(w=e^{s\psi}\), \(|s|\le s_*\),
the actual maps satisfy
\begin{equation}\label{scr:weighted-restoration}
 \|\ell^{-2}w\Lambda w^{-1}v\|_2
       +\|wD_U\Lambda w^{-1}v\|_X
       +\|wSw^{-1}v\|_X\le C\|v\|_X.
\end{equation}
Consequently
\begin{equation}\label{scr:weighted-current-bound}
 \|w\mathcal B_U[J]w^{-1}\|_{X\to X^*}\le C\gamma_J.
\end{equation}
No derivative of \(J\) or second derivative of \(\psi\) is a
premise.
\end{lemma}

\begin{proof}
For \(\lambda=\Lambda w^{-1}v\), put
\(f=L_U\lambda=RD_U^\sharp w^{-1}v\).
The weighted site projector bound and the first derivative
commutator give \(\|wf\|_2\le C\|v\|_X\).
Also \(q\lambda=0\), including the bare collar.

Repeat the zero-average proof with the combined scalar weight
\(h_w=\ell^{-1}e^{s\psi}\).
Its relative first difference is bounded by
\(C(\delta+|s|)/\ell\). Subtract \(h_w(y)q\lambda=0\) in each
pillar. The site-gap average error is now at most
\[
 C_L(\delta+|s|)\|\ell^{-2}w\lambda\|_2.
\]
Testing \(L_U\lambda=f\) against \(h_w^2\lambda\) and using the
exact weighted Green identity proves
\[
 \|\ell^{-2}w\lambda\|_2+
               \|\ell^{-1}wD_U\lambda\|_2\le C\|wf\|_2,
\]
after reducing \(s_*\). The weighted Hodge calculation for
\(wD_U\lambda\) contains only first weight differences.
Its divergence is \(wf\) plus \(O(|s|\ell^{-1}wD_U\lambda)\);
its curl is weighted curvature times \(\lambda\), plus that
same first-weight error. Thus its component-gradient energy
is bounded by \(C\|wf\|_2^2\). The zero-average bare collar
keeps every nonzero stencil inside the action region, exactly
as before. This proves \eqref{scr:weighted-restoration}.

Use the four undifferentiated-current pairings in
\eqref{px:current-form-no-current-gradient}. Apply the
restoration bounds with both signs of \(s\). The endpoint
ratios of \(w\) in \(\mathcal W_U[J]\) are bounded uniformly,
since an endpoint step has bounded metric length.
Each pairing is therefore bounded by
\(C\gamma_J\|u\|_X\|v\|_X\).
This proves \eqref{scr:weighted-current-bound}.
\end{proof}

The weighted inverse of \(\mathcal K_0\) is bounded from \(X^*\)
to \(X\) by the already proved field Combes argument.
For sufficiently small \(\gamma_J\), the normalized Neumann
series using \eqref{scr:weighted-current-bound} gives
\begin{equation}\label{scr:weighted-GJ}
                  \|wG_Jw^{-1}\|_{X^*\to X}\le C.
\end{equation}
This proof applies to complex currents. It invokes positivity
only at the real zero-current reference; it asserts no
Hermitian positivity at a complex current.

\subsubsection{The normalized coarse Schur matrix}

Define
\begin{equation}\label{scr:coarse-definitions}
 S_J=QG_JQ^\sharp,\quad
 T_J=a_f^{1/2}S_Ja_f^{1/2},\quad
 T_J^{\rm n}=\Xi_\Gamma^2T_J\Xi_\Gamma^{-2}.
\end{equation}
The last matrix acts in counting orthonormal coordinates.
Give a row \(c\) the metric location of one incident top at its
row level. Its carrier is within bounded layer-metric distance
of that top, and at most eight rows are assigned to a top.
The row ball count has the same growth rate as the top count.

\begin{proposition}[Uniform normalized Schur inverse]
\label{prop:scr-schur}
For sufficiently small \(\gamma_J\), \(S_J\) is invertible and
\begin{equation}\label{scr:schur-kernels}
 |(T_J^{\rm n})^{-1}(c,d)|\le Ce^{-m d(c,d)}.
\end{equation}
The constants are independent of depth, compatible volume and
the current. After increasing the fixed separation so the
metric growth rate is less than \(m/2\), this inverse is bounded
on counting \(\ell^\infty\). Equivalently \(T_J^{-1}\) is bounded
in the coarse amplitude norm \(\sup_c\xi_c^2|\beta(c)|\).
\end{proposition}

\begin{proof}
At \(J=0\), \(T_0\) has a positive lower and upper bound.
For the lower bound take the exact field-row right inverse
\(E_Q\), or \(E_L=SE_Q\), and put
\(v=E_Qa_f^{-1/2}\beta\). Its \(\mathcal K_0\) energy is at
most \(C\|\beta\|_{\mathcal H_\Gamma}^2\). Since
\(Qv=a_f^{-1/2}\beta\),
\begin{equation}\label{scr:schur-energy-lower}
 \begin{split}
 \|\beta\|_{\mathcal H_\Gamma}^2
 &=\langle Q^\sharp a_f^{1/2}\beta,
                                    E_Qa_f^{-1/2}\beta\rangle\\
 &\le C\langle\beta,T_0\beta\rangle^{1/2}
                                      \|\beta\|_{\mathcal H_\Gamma}.
 \end{split}
\end{equation}
The upper bound follows from the energy inverse and the bounded
trace \(a_f^{1/2}Q:X\to\mathcal H_\Gamma\).
No positivity of a different effective coarse action has been
used.

Let \(w_\Gamma\) sample \(e^{s\psi}\) at the row locations.
The finite metric carriers of the actual rows imply
\[
 \|w_\Gamma a_f^{1/2}Qw^{-1}-a_f^{1/2}Q\|_{X\to\mathcal H_\Gamma}
                                                        \le C|s|.
\]
The weighted base inverse differs from its original value by
\(C|s|\). Its current perturbation differs from the weighted
base inverse by \(C\gamma_J\), by the same normalized Neumann
series as in \eqref{scr:weighted-GJ}. Therefore
\begin{equation}\label{scr:schur-weight-perturbation}
 \|w_\Gamma T_Jw_\Gamma^{-1}-T_0\|
                                             \le C(|s|+\gamma_J).
\end{equation}
Choose \(s\) and the current ceiling so the right side is less
than half the real lower constant of \(T_0\).
The inverse of the conjugated \(T_J\) is uniformly bounded.
Truncated distance weights give exponential entries between
physical orthonormal row cells, which are precisely
\((T_J^{\rm n})^{-1}(c,d)\). The cell count and a strict rate
reserve give the stated \(\ell^\infty\) bound.
\end{proof}

For reference, these normalizations imply the counting-array
and physical-kernel bounds
\begin{equation}\label{scr:unnormalized-Schur-kernels}
 |S_J^{-1}(c,d)|_{\rm array}
       \le C\xi_c^{-3}\xi_d e^{-m d(c,d)},\qquad
 |S_J^{-1}(c,d)|_{\rm physical}
       \le C\xi_c^{-3}\xi_d^{-3}e^{-m d(c,d)}.
\end{equation}
Indeed \(S_J^{-1}=a_f^{1/2}T_J^{-1}a_f^{1/2}\), and
converting its array to a physical kernel divides by the input
measure \(\xi_d^4\). These are different conventions for the
same operator.

\subsubsection{The actual minimizing lift and the induced multiplier}

Define the actual first-variant lift by
\begin{equation}\label{scr:actual-lift}
                   \widehat H_J=G_JQ^\sharp S_J^{-1},
                         \qquad Q\widehat H_J=1.
\end{equation}
It has the useful normalized representation
\begin{equation}\label{scr:normalized-lift}
       \widehat H_Ja_f^{-1/2}
                       =G_JQ^\sharp a_f^{1/2}T_J^{-1}.
\end{equation}
The weighted inverse, trace and normalized Schur bounds make
this a bounded weighted operator from \(\mathcal H_\Gamma\)
to \(X\). In particular, with \(X_y\) the cell component energy
and a fixed one-fine-step collar for its derivatives,
\begin{equation}\label{scr:lift-energy-columns}
 \|1_y\widehat H_J1_c b\|_{X_y}
                         \le Ce^{-m d(y,c)}\|1_cb\|_Y.
\end{equation}
This follows directly by a distance weight centered at \(c\);
it does not require summing an unnormalized fine kernel.

\begin{proposition}[The actual induced row current has its physical amplitude]
\label{thm:scr-induced-current}
For the displayed first precision and its actual lift, put
\begin{equation}\label{scr:induced-current}
                         \omega_J=\widehat H_J^\sharp J.
\end{equation}
For sufficiently small \(\gamma_J\),
\begin{equation}\label{scr:induced-current-bound}
                         \sup_c\xi_c^3|\omega_J(c)|\le C\gamma_J.
\end{equation}
This holds uniformly over all assigned scales, determining
regions and compatible volumes, including bare and outward
rows. It is valid for complex independent currents on the
stated small ball.
\end{proposition}

\begin{proof}
Fix a row \(c\) and a unit colour \(e\). Take the datum
\(b^{(c)}(c)=\xi_c^{-1}e\), zero elsewhere.
Its squared \(Y\) norm is \(\xi_c^2a_f(c)\), bounded above
and below by fixed constants. Physical transposition gives
the exact dimensional identity
\begin{equation}\label{scr:column-current-duality}
 \xi_c^3\langle e,\omega_J(c)\rangle
                         =\langle\widehat H_Jb^{(c)},J\rangle_{\rm fine}.
\end{equation}
On a cell of side \(\xi_y\), the current assumption implies
\begin{equation}\label{scr:current-cell-energy}
                           \|\xi J\|_{2,B_y}\le C_L\gamma_J.
\end{equation}
Indeed its physical volume is \(\xi_y^4\), while the current
amplitude is at most \(C_L\gamma_J\xi_y^{-3}\).
Pair cell by cell in \eqref{scr:column-current-duality}:
\[
 |\xi_c^3\langle e,\omega_J(c)\rangle|
 \le\sum_y\|\xi J\|_{2,B_y}
                  \|\xi^{-1}\widehat H_Jb^{(c)}\|_{2,B_y}
 \le C\gamma_J\sum_y e^{-m d(y,c)}.
\]
Increase the fixed separation, if necessary, so that its growth
rate is below \(m/2\). The last sum is uniform. Taking the three
colour components proves \eqref{scr:induced-current-bound}.
For complex currents the same bilinear identity and absolute
Cauchy--Schwarz bound apply; no current was replaced by the
Wilson gradient.
\end{proof}

At a bare row this proof is still dimensionally exact.
The cell force norm is \(\eta\|J\|_2=O(\gamma_J)\), the
unit energy datum has size \(\eta^{-1}\), and its physical
row dual has size \(\eta^3\omega\). No ultraviolet factor was
discarded.

\subsubsection{Exponential coupled kernels}

The full penalized saddle equations have an explicit inverse:
\begin{equation}\label{scr:penalized-saddle}
 \begin{split}
 \mathcal K_Ju-Q^\sharp\mu&=f,\qquad Qu=b,\\
 u&=(G_J-\widehat H_JQG_J)f+\widehat H_Jb,\\
 \mu&=S_J^{-1}(b-QG_Jf).
 \end{split}
\end{equation}
All normalized factors on the right have uniformly bounded
weighted energy norms. The same holds for the resulting block
inverse. It consequently has exponential cell kernels with the
physical force and datum normalizations: a force cell costs
\(\xi_z\|f\|_{2,B_z}\), a datum row costs \(\|b\|_{Y_c}\),
a field output is measured in \(X_y\), and a multiplier output
is measured by \(\xi_c^3|\mu(c)|\), up to its fixed weight
constant. In particular
\begin{equation}\label{scr:saddle-energy-bound}
                    \|u\|_X+\|\mu\|_{Y^*}
                         \le C(\|f\|_{X^*}+\|b\|_Y).
\end{equation}
Removing the artificial row-mass term from the first equation
changes its multiplier to \(\mu-a_fb\), a local bounded map
in \(Y^*\). It preserves these bounds and spatial kernels.

For an independent current the minimizing lift is not asserted
to lie in the Landau slice. Its exact off-current slice defect
is the one in \eqref{px:off-current-slice-defect}.
The coupled formula above keeps the displayed gauge penalty.

\subsubsection{The second explicit variant is admitted}

Now define, in the prescribed order,
\begin{equation}\label{scr:second-precision}
 \mathcal K_{1,J}=\mathcal K_J
                        -S^\sharp\Delta^{(2)}[\omega_J]S,\qquad
                \omega_J=\widehat H_J^\sharp J.
\end{equation}
The lift inside \(\omega_J\) is the first-variant lift in
\eqref{scr:actual-lift}. Replacing it by the lift formed from
\(\mathcal K_{1,J}\) would be a different definition.

The new term has energy norm at most \(C\gamma_J\), by
\eqref{scr:induced-current-bound} and the actual two-leg
constraint estimate. Its conjugated norm has that same bound:
the row carrier has fixed metric radius and the two restoration
factors have \eqref{scr:weighted-restoration}.
Thus a further fixed decrease of the current ceiling makes
\(\mathcal K_{1,J}\) invertible with the same weighted energy
bounds. Its normalized coarse Schur matrix, lift and coupled
inverse satisfy the same conclusions. This proves admission
of the second displayed variant, rather than merely bounding a
constraint form at a freely chosen small multiplier.

The maps are holomorphic in the independent current ball.
The first inverse, normalized Schur inverse and lift are
convergent analytic inverses there; the readout
\(\widehat H_J^\sharp J\) is consequently holomorphic, and
so is the second precision and its inverse. The transpose is
bilinear throughout this continuation.

\subsubsection{The physical slice}

If \(J=j_U\) is the actual Wilson current, the inner-action
Ward identity and the displayed current correction give
\begin{equation}\label{scr:physical-gauge-inverse}
       G_JD_UR=D_UGR.
\end{equation}
The second subtraction annihilates \(D_UGR\), since
\(SD_UGR=0\), so this identity also holds for
\(G_{1,J}=\mathcal K_{1,J}^{-1}\).
Transposition and \(QD_UGR=\bar D_{\rm ext}qGR=0\) imply
that its row lift is horizontal. With that second lift
\(\widehat H_{1,J}\), set
\(C_1=G_{1,J}-\widehat H_{1,J}QG_{1,J}\).
Then \(C_1S^\sharp f\) is both horizontal and in \(\ker Q\):
the gauge cancellation is
\((SD_UGR)^\sharp=0\).
These identities use the actual extended outward difference.

At an actual constrained critical point, full stationarity is
\(j_U=Q^\sharp\omega_{\rm phys}\).
Since \(Q\widehat H_J=1\), it follows exactly that
\(\omega_J=\omega_{\rm phys}\).
In particular the induced-current theorem bounds the actual
multiplier at every layer if the actual current has its
stated small amplitude. The full stationarity equation can
be deduced from criticality on the actual horizontal constraint
kernel: \(S\) preserves \(Q\), and the Ward identity annihilates
the gauge part \(1-S\). It is not inferred from criticality
on an unidentified slice.

On horizontal vectors the Ward correction vanishes and
\(S=1\). Therefore the second form restricts to the actual
Lagrangian Hessian plus the row-mass term.
The formula
\[
 u=C_1S^\sharp f+\widehat H_{1,J}b
\]
is the physical horizontal saddle solution, with its multiplier
adjusted by \(-a_fb\) when the row mass is removed.
Its exponential energy kernels follow from the weighted
factors already proved. No fine pointwise inverse has been
used to obtain this physical slice statement.

The conclusions concern the two displayed current extensions.
They do not identify another independently named
\(\Delta'_\pi[J]\), nor a signed-profile precision or the
absolute contour precision. Fine pointwise and H\"older
responses, moving complex background slices and the nonlinear
trial-field construction remain further tasks.

% END LAYERED EXPANSION inverse/coarse_schur_current_readout.tex

% BEGIN LAYERED EXPANSION inverse/complex_moving_landau_inverse.tex
\subsection{A complex inverse on the moving Landau slice}
\label{ch:complex-moving-landau}

The real proper-region estimate has a uniform complex
neighbourhood. The relevant slice moves because its site
projection and covariant divergence depend on the background.
We first continue the site operators by a balanced physical
form. We then solve a saddle problem with a fixed target space
for both constraints. Its zero gauge datum is exactly the
moving Landau condition.

The geometry, scales, bare collars and outward rows remain
those of \cref{thm:npl-proper-gap}. Fix a real background \(U_0\)
with the curvature ceiling there reduced by a factor of two.
All comparison norms and \(\ell\) in this section refer to
this fixed geometry and real frame. Write
\begin{equation}\label{cml:tube}
 U_A(b)=e^{i\eta A(b)}U_0(b),\qquad
 \|A\|_{\mathcal T}
   =\sup_b\ell(b)|A(b)|
             +\sup\ell^2|\nabla_{U_0}A|<\varepsilon_* .
\end{equation}
The comparison derivatives use the same zero-extension
convention as \(X\). At the bare boundary \(\ell=\eta\), so
the boundary step also has this scale. The field \(A\) may
be complex. Transposes of analytic operators are bilinear
physical transposes, not conjugate transposes.

\subsubsection{Uniform continuation of the actual site projection}

For this section use a subscript \(A\) for actual objects at
\(U_A\), and a subscript \(0\) for the real reference. In
particular
\begin{equation}\label{cml:site-definitions}
 \begin{split}
 A_{s,A}&=D_A^\sharp D_A+q_A^\sharp a_s q_A,\qquad
 G_{s,A}=A_{s,A}^{-1},\\
 F_A&=G_{s,A}q_A^\sharp\Xi^{-2},\qquad
 M_A=F_A^\sharp F_A,\qquad
 R_A=1-F_AM_A^{-1}F_A^\sharp .
 \end{split}
\end{equation}
Here \(\Xi\) is the site-top scale, not the field-row scale.

\begin{lemma}[Balanced complex site form]
\label{lem:cml-site}
For a fixed small \(\varepsilon_*\), all operators in
\eqref{cml:site-definitions} are defined and holomorphic
throughout \eqref{cml:tube}. With physical Hilbert operator
norms,
\begin{equation}\label{cml:site-comparison}
 \|F_A-F_0\|+\|\ell D_0(F_A-F_0)\|
           +\|M_A-M_0\|+\|R_A-R_0\|\le C\|A\|_{\mathcal T}.
\end{equation}
Moreover \(R_A^\sharp=R_A\), \(R_A^2=R_A\), and
\begin{equation}\label{cml:site-algebra}
 q_AG_{s,A}R_A=0,\qquad
         D_A^\sharp D_AG_{s,A}R_A=R_A .
\end{equation}
No complex orthogonality is claimed.
\end{lemma}

\begin{proof}
The actual transport paths have length at most a fixed
multiple of their row scale. On their carriers \(\ell\)
is comparable to that scale. Telescoping their transports
therefore gives, for each of the finitely many powers
\(r=-2,-1,0,1,2\),
\begin{equation}\label{cml:actual-row-comparison}
 \begin{split}
 \|\Xi^{-r}(q_A-q_0)\ell^r\|&\le C\|A\|_{\mathcal T},\\
 \|\ell(D_A-D_0)u\|_2
                         &\le C\|A\|_{\mathcal T}\|u\|_2,\\
 \|a_f^{1/2}(Q_A-Q_0)v\|_{\mathcal H_\Gamma}
                         &\le C\|A\|_{\mathcal T}\|\ell^{-1}v\|_2 .
 \end{split}
\end{equation}
The same estimates follow for the transposes. Physical
carrier Schur sums, rather than the number of fine sites
in a carrier, are used. The local logarithms remain in
their common analytic charts after a fixed decrease of
\(\varepsilon_*\). The geometric sum of path lengths
through successive homogeneous subscales is bounded by
the final path scale. Thus these comparisons are
independent of depth. They include completed outward rows;
the identity \(Q_AD_A=\bar D_{{\rm ext},A}q_A\) is kept
with its actual extension.

Put
\[
 \|u\|_{Z_+}^2=\|\ell D_0u\|_2^2+\|q_0u\|_{\rm top}^2 .
\]
The real block inequality gives \(\|u\|_2\le C\|u\|_{Z_+}\).
Consider the sesquilinear testing form
\(\mathfrak a_A(u,v)=
 \langle A_{s,A}u,\ell^2\overline v\rangle_{\rm bil}\).
At the real reference its weighted Green identity, and
the \(O(\delta)\) average commutator, give
\begin{equation}\label{cml:balanced-coercivity}
 \operatorname{Re}\mathfrak a_0(u,u)\ge c\|u\|_{Z_+}^2,\qquad
 |\mathfrak a_A(u,v)-\mathfrak a_0(u,v)|
       \le C\|A\|_{\mathcal T}\|u\|_{Z_+}\|v\|_{Z_+}.
\end{equation}
Indeed a derivative error is \(O(\varepsilon_*\ell^{-1}u)\)
before multiplication by \(\ell\); the row errors in the
test \(q_A(\ell^2v)\) are cancelled by the \(\Xi^{-2}\)
site mass. The first differences of \(\ell\) cost \(O(\delta)\).
Choose \(\delta\) as in the real proof and then
\(\varepsilon_*\) small. This is a coercive perturbation
on \(Z_+\).

The source for \(F_A b\) obeys
\[
 |\langle q_A^\sharp\Xi^{-2}b,\ell^2\overline v\rangle_{\rm bil}|
                           \le C\|b\|_{\rm top}\|v\|_{Z_+}.
\]
Its difference from the real source is \(O(\|A\|_{\mathcal T})\)
in that dual norm. Solving the balanced form and subtracting
the two equations proves the first two bounds in
\eqref{cml:site-comparison}. This also proves holomorphy by
a form Neumann series on fixed spaces; the testing weight
does not depend on \(A\).

The real normalized Gram lower bound implies
\(M_0^{-1}\) is uniformly bounded. The \(L^2\) bound for
\(F_A-F_0\) gives \(M_A-M_0=O(\|A\|_{\mathcal T})\).
Its inverse and the last bound in
\eqref{cml:site-comparison} follow by another Neumann
series. The projection algebra in
\eqref{cml:site-definitions} is bilinear algebra, so it
continues exactly. Since \(F_A^\sharp R_A=0\) and
\(G_{s,A}^\sharp=G_{s,A}\), the first identity in
\eqref{cml:site-algebra} follows. Apply \(A_{s,A}G_{s,A}=1\)
to obtain the second.
\end{proof}

\subsubsection{Two physical derivatives in the complex zero-average space}

\begin{lemma}[Complex gauge restoration]
\label{lem:cml-restoration}
If \(q_A\lambda=0\) and \(D_A^\sharp D_A\lambda=f\), then
\begin{equation}\label{cml:zero-average-H2}
 \|\ell^{-2}\lambda\|_2+\|\ell^{-1}D_0\lambda\|_2
                          +\|\nabla_0D_A\lambda\|_2\le C\|f\|_2 .
\end{equation}
Consequently
\begin{equation}\label{cml:restoration}
 \Lambda_A=G_{s,A}R_AD_A^\sharp,\qquad S_A=1-D_A\Lambda_A
\end{equation}
are bounded holomorphic maps in their physical norms, with
\(\|D_A\Lambda_Av\|_X+\|S_Av\|_X\le C\|v\|_X\).
Exactly
\[
 R_AD_A^\sharp S_A=0,\qquad Q_AS_A=Q_A .\qedhere
\]
\end{lemma}

\begin{proof}
Apply the real weighted site gap to \(\ell^{-1}\lambda\).
In its average term insert \(q_A\lambda=0\).
The weight commutator costs \(C\delta\|\ell^{-2}\lambda\|_2\),
and \eqref{cml:actual-row-comparison} costs
\(C\varepsilon_*\|\ell^{-2}\lambda\|_2\).
After absorption,
\[
 \|\ell^{-2}\lambda\|_2
                         \le C\|D_0(\ell^{-1}\lambda)\|_2 .
\]
Test the complex equation against
\(\ell^{-2}\overline\lambda\), and compare its bilinear form
with the real reference form. The exact real weighted Green
identity has an \(O(\delta^2)\|\ell^{-2}\lambda\|_2^2\)
error. Each complex derivative perturbation contributes at
most
\[
 C\varepsilon_*
 \bigl(\|\ell^{-1}D_0\lambda\|_2+
       \|\ell^{-2}\lambda\|_2\bigr)^2 .
\]
Absorb these errors using the preceding gap. The forcing
pairing is bounded by
\(\|f\|_2\|\ell^{-2}\lambda\|_2\).
This proves the first two terms of
\eqref{cml:zero-average-H2}, without a derivative of \(A\).

For the last term the \(C^1\) part of the tube is used.
Expanding the actual covariant stencils gives
\begin{equation}\label{cml:laplace-difference}
 \|(D_0^\sharp D_0-D_A^\sharp D_A)\lambda\|_2
       +\|\nabla_0(D_A-D_0)\lambda\|_2
 \le C\varepsilon_*
       \bigl(\|\ell^{-1}D_0\lambda\|_2+
                         \|\ell^{-2}\lambda\|_2\bigr).
\end{equation}
The differentiated transport difference contains
\(\nabla_0A\), \(A D_0\lambda\), and \(A^2\lambda\).
Its two explicit tube powers are respectively
\(\ell^{-2}\), \(\ell^{-1}\), and \(\ell^{-2}\);
there is no second derivative of \(A\).
The bare row is still the identity, so \(\lambda=0\) on
the same bare collar. Extend it by zero and apply the
real covariant Hodge estimate \eqref{px:covariant-hodge-error}
to \(D_0\lambda\). Every nonzero second-derivative stencil
stays inside the inner-plaquette region. Equation
\eqref{cml:laplace-difference} and the first two bounds
give the final term in \eqref{cml:zero-average-H2}.

Take \(\lambda=G_{s,A}R_AD_A^\sharp v\).
Its source has norm at most \(C\|v\|_X\), and
\eqref{cml:site-algebra} supplies both equations just used.
The mass term of \(D_A\lambda\) follows from the same
comparison. This proves the norm assertions.
All maps are holomorphic in the finite coordinate system,
and their uniform estimates are in the displayed spaces.
Finally \eqref{cml:site-algebra} and the actual intertwining
give the two exact identities.
\end{proof}

\subsubsection{The combined constraint has a fixed target}

Let \(Z_0=\operatorname{ran}R_0\), with its physical \(L^2\)
norm. Define the combined constraint on the fixed field
space \(X\) by
\begin{equation}\label{cml:combined-constraint}
 \mathcal C_Au=(Q_Au,\ R_0R_AD_A^\sharp u)
                         \in Y\oplus Z_0 .
\end{equation}
This use of \(R_0\) fixes the target only. It does not
replace the moving projection \(R_A\).

\begin{lemma}\label{lem:cml-combined-infsup}
The reference constraint has a bounded exact right inverse
\begin{equation}\label{cml:combined-right-inverse}
 \mathcal E_0(b,\rho)=E_{L,0}b+D_0G_{s,0}R_0\rho .
\end{equation}
Also
\begin{equation}\label{cml:combined-comparison}
          \|\mathcal C_A-\mathcal C_0\|_{X\to Y\oplus Z_0}
                                          \le C\|A\|_{\mathcal T}.
\end{equation}
Thus \(\mathcal C_A\) is onto with a uniform right inverse.
Its second component is zero if and only if
\(R_AD_A^\sharp u=0\).
\end{lemma}

\begin{proof}
The first summand of \eqref{cml:combined-right-inverse}
has row value \(b\) and zero reference divergence.
The second has zero row value by
\(Q_0D_0G_{s,0}R_0=\bar D_{{\rm ext},0}q_0G_{s,0}R_0=0\).
Its reference projected divergence is \(\rho\), by
\eqref{cml:site-algebra}. Its energy bound is the real
zero-average restoration estimate.
This proves the right-inverse assertion.

For the comparison use
\eqref{cml:actual-row-comparison},
\eqref{cml:site-comparison}, and
\(\|D_0^\sharp u\|_2\le C\|u\|_X\).
Then \(\mathcal C_A\mathcal E_0=1+O(\varepsilon_*)\);
inverting it gives a right inverse for \(\mathcal C_A\).

If \(z\in\operatorname{ran}R_A\) and \(R_0z=0\), then
\(z=(R_A-R_0)z\), so \(z=0\).
Also \(R_0R_A:Z_0\to Z_0\) is \(1+O(\varepsilon_*)\).
It follows that \(R_0:\operatorname{ran}R_A\to Z_0\)
is an isomorphism. Apply its injectivity to
\(z=R_AD_A^\sharp u\) to prove the final equivalence.
\end{proof}

\subsubsection{The full complex saddle and its actual slice}

Let \(\omega\) be a complex field-row multiplier with
\(\gamma_\omega=\sup_c\xi_c^3|\omega(c)|\), and define the
displayed Lagrangian Hessian
\begin{equation}\label{cml:lagrangian}
           H_{A,\omega}=D^2A_{\Omega_0}(U_A)
                                  -\Delta_A^{(2)}[\omega].
\end{equation}
No current extension is part of this definition.

\begin{proposition}[Complex saddle inverse]
\label{thm:cml-moving-slice}
There are depth- and volume-independent
constants \(\varepsilon_*>0\) and \(\gamma_*>0\).
If \eqref{cml:tube} holds and
\(\gamma_\omega<\gamma_*\), the system
\begin{equation}\label{cml:full-saddle}
 \begin{split}
 H_{A,\omega}u-Q_A^\sharp\mu-D_AR_AR_0\nu&=f,\\
 Q_Au=b,\qquad R_0R_AD_A^\sharp u&=\rho
 \end{split}
\end{equation}
has a unique solution
\((u,\mu,\nu)\in X\oplus Y^*\oplus Z_0\) and
\begin{equation}\label{cml:full-saddle-bound}
 \|u\|_X+\|\mu\|_{Y^*}+\|\nu\|_2
                \le C(\|f\|_{X^*}+\|b\|_Y+\|\rho\|_2).
\end{equation}
Its inverse is holomorphic in \((A,\omega)\) on this fixed
product of spaces. At \(\rho=0\) it gives exactly the
actual moving-slice saddle
\begin{equation}\label{cml:actual-slice-saddle}
 \begin{gathered}
 X_L(A)=\ker(R_AD_A^\sharp),\qquad Q_Au=b,\\
 H_{A,\omega}u-Q_A^\sharp\mu=f
                                      \quad\hbox{on }X_L(A).
 \end{gathered}
\end{equation}
This saddle has a unique \((u,\mu)\) with the bound
\par\noindent
\begin{minipage}{\linewidth}
\begin{equation}\label{cml:actual-slice-bound}
       \|u\|_X+\|\mu\|_{Y^*}
                   \le C(\|f\|_{X_L(A)^*}+\|b\|_Y).
\end{equation}
\popQED\par
\end{minipage}
\end{proposition}

\begin{proof}
At \(A=0,\omega=0\), the real proper-region gap gives
coercivity of \(H_{0,0}\) on
\(\ker\mathcal C_0=X_L(0)\cap\ker Q_0\).
Its full form is bounded on \(X\).
The right inverse \eqref{cml:combined-right-inverse}
then proves invertibility of the full reference saddle.
For completeness, lift \((b,\rho)\), solve the coercive
equation on \(\ker\mathcal C_0\), and recover the multiplier
by evaluating the remaining functional on the same right
inverse. This gives the full estimate, with no positivity
assertion off the constraint kernel.

The actual plaquette expansion implies
\begin{equation}\label{cml:hessian-comparison}
             \|H_{A,\omega}-H_{0,0}\|_{X\to X^*}
                                  \le C(\|A\|_{\mathcal T}+\gamma_\omega).
\end{equation}
To check the scale, the curl-row difference is bounded
by \(C\|A\|_{\mathcal T}\|\ell^{-1}v\|_2\).
The physical plaquette-curvature difference is bounded by
\(C\|A\|_{\mathcal T}\ell^{-2}\), from the first covariant
differences of \(A\) and the quadratic \(A\) term.
The exact Wilson Hessian remainder is consequently a
bounded mass-form perturbation. This is a local
inner-plaquette calculation and has no boundary
integration-by-parts term. The actual constraint
two-leg estimate bounds the last term of
\eqref{cml:lagrangian} by \(C\gamma_\omega\) in the two
energy norms. Its carrier estimates continue on the same
complex charts with a uniform constant.

The transpose of the second component of
\(\mathcal C_A\) is \(D_AR_AR_0\), because \(R_A\) is
bilinearly symmetric. The whole operator in
\eqref{cml:full-saddle} therefore differs from the
reference saddle by
\(C(\|A\|_{\mathcal T}+\gamma_\omega)\) between fixed
energy and dual-energy product spaces.
The reference inverse and a Neumann series prove the
first assertion and \eqref{cml:full-saddle-bound}.
Every entry is holomorphic, so its inverse is
holomorphic there as well. Complex orthogonality and a
parametrization of the moving field slice were not used.

Set \(\rho=0\). The previous lemma identifies its kernel
with the actual \(X_L(A)\). Testing the first equation
there eliminates \(\nu\), proving
\eqref{cml:actual-slice-saddle}.
Conversely, any solution of that restricted equation
extends to \eqref{cml:full-saddle} by a unique gauge
multiplier. Indeed \(R_0R_AD_A^\sharp:X\to Z_0\) is
onto, and the annihilator of its kernel is its
transpose range. The transpose is injective by the
same right-inverse bound. Full saddle uniqueness
therefore proves restricted uniqueness.

For a functional given only on \(X_L(A)\), take a
norm-preserving Hahn--Banach extension to \(X\).
The full bound gives \eqref{cml:actual-slice-bound}.
Another extension changes only the eliminated gauge
multiplier, so the field and physical row multiplier
are independent of that choice. Holomorphic dependence
is asserted for fixed-space source extensions, and
requires no holomorphic choice of Hahn--Banach
extensions.
\end{proof}

On any smaller fixed tube, Cauchy's inequality gives uniform
operator jets in the displayed physical energy norms, with factorial
bounds in the tube norm and in \(\gamma_\omega\).
The result permits complex backgrounds and the actual moving
Landau constraint at all assigned scales on proper determining
regions. Strong outputs and the regular-pair nonlinear response
are proved in \cref{prop:acc-strong-propagators,dnr:linear-families,dnr:response}.
The $C^1$ background tube itself does not bound the physical Wilson
current; the current is independent in these constructions.

% END LAYERED EXPANSION inverse/complex_moving_landau_inverse.tex

% BEGIN LAYERED EXPANSION discrete_calderon_zygmund.tex
\subsection{Uniform discrete singular-integral estimates}
\label{dcz:section}

The strong proper-region estimates require finite-\(p\) bounds for
second derivatives of a scalar inverse. Supremum-norm bounds for
these order-zero operators would be false. We extend
\cref{lem:bg-finite-p} to arbitrary masses and periods, and
derive the local heat embeddings needed for the response
bootstrap. The singular-integral proof itself was given in
that lemma, including the physical measure at the last
atomic scale; only its new uniform parameters are checked here.

Let \(\Gamma_h=h\mathbb Z^4\), \(0<h\le1\), with measure
\(h^4\) times counting measure. Put
\[
 D_i^+u(x)=h^{-1}(u(x+he_i)-u(x)),\qquad
 D_i^-u(x)=h^{-1}(u(x)-u(x-he_i)),\qquad
                         L_h=-\sum_iD_i^-D_i^+.
\]
Either sign of a difference will be allowed below. The notation
\(\|D^2u\|_p\) includes all ordered component second
differences. Replacing a forward by a backward difference
only shifts its location and changes fixed constants.
All convolutions use physical measure:
\[
                  (k*u)(x)=h^4\sum_y k(x-y)u(y).
\]

\subsubsection{The scalar kernels}

Let \(p_t^h\) be the physical heat density of \(e^{-tL_h}\).
The differentiated heat estimates already proved in
\eqref{bg:heat-pointwise} imply, for every derivative
order \(a\ge0\) and every integer \(N\),
\begin{equation}\label{dcz:heat-pointwise}
 |D^ap_t^h(x)|\le C_{a,N}(t+h^2)^{-(4+a)/2}
             \left(1+\frac{|x|}{\sqrt{t+h^2}}\right)^{-N}.
\end{equation}
The quoted local-time estimate extends to \(t\ge1\)
on the infinite lattice by the exact scaling
\[
       p_t^h(x)=t^{-2}p_1^{\,h/\sqrt t}(x/\sqrt t).
\]
Each derivative contributes a further \(t^{-1/2}\).
The rescaled mesh is still at most one, so the
same pointwise estimate holds at every time.
They also give
\begin{equation}\label{dcz:heat-Lp}
 \|D^ap_t^h\|_{r,h}
       \le C_{a,r}(t+h^2)^{-a/2-2(1-1/r)},
                  \qquad1\le r\le\infty.
\end{equation}
For finite \(r\), sum \eqref{dcz:heat-pointwise}
with \(Nr>4\); the lattice sum of the spatial factor is
at most \(C(t+h^2)^2\). The case \(r=\infty\) is the
pointwise bound. Thus these estimates include \(t<h^2\).

For \(\mu\ge0\), define the second-resolvent kernel
\begin{equation}\label{dcz:kernel-definition}
 k_{ij}^{h,\mu}(x)
   =\int_0^\infty e^{-\mu^2t}D_i^\sigma D_j^\tau p_t^h(x)\,dt,
                    \qquad \sigma,\tau\in\{+,-\}.
\end{equation}
This integral converges at every lattice point, including
zero: the mesh is fixed at the lower limit, and the
large-time derivative bound is \(Ct^{-3}\).
For compactly supported input, its convolution is the
Fourier multiplier
\[
 \frac{\mathrm{d}_i^\sigma(\theta)\mathrm{d}_j^\tau(\theta)}
      {\lambda_h(\theta)+\mu^2},\qquad
 \lambda_h(\theta)=4h^{-2}\sum_i\sin^2(h\theta_i/2),
       \quad \theta\in[-\pi/h,\pi/h]^4.
\]
At \(\mu=0,\theta=0\) the value of the multiplier can be
assigned arbitrarily on the infinite lattice.
Since \(|\mathrm{d}_i^\sigma \mathrm{d}_j^\tau|\le\lambda_h\),
its \(L^2\) operator norm is at most one.

\begin{lemma}\label{dcz:kernel-regularity}
Uniformly in \(h,\mu\), these kernels satisfy
\begin{equation}\label{dcz:kernel-bounds}
 |k_{ij}^{h,\mu}(x)|\le C(|x|+h)^{-4},\qquad
 |D_\ell^\pm k_{ij}^{h,\mu}(x)|
                                      \le C(|x|+h)^{-5}.
\end{equation}
In particular, for a lattice displacement \(y\),
\begin{equation}\label{dcz:Hormander}
 h^4\sum_{|x|\ge C_0(|y|+h)}
      |k_{ij}^{h,\mu}(x-y)-k_{ij}^{h,\mu}(x)|\le C
\end{equation}
with a fixed sufficiently large \(C_0\).
\end{lemma}

\begin{proof}
The mass factor can be omitted for these upper bounds.
For \(|x|=r\ge h\), integrate
\eqref{dcz:heat-pointwise} with \(a=2\), using
\(v=t+h^2\). For \(v\le r^2\), its bound is
\(C r^{-N}v^{N/2-3}\), which has integral
\(Cr^{-4}\) if \(N>4\). For \(v\ge r^2\), the
integral of \(Cv^{-3}\) has the same bound.
At \(r\le h\) the full integral is \(Ch^{-4}\).
Using \(a=3,N>5\) proves the derivative bound.

Join \(0\) to \(y\) by coordinate steps. On the region
in \eqref{dcz:Hormander}, all intermediate arguments
are comparable in distance to \(x\).
Telescoping the derivative bound gives
\[
        |k(x-y)-k(x)|\le C|y|\,|x|^{-5}.
\]
On the dyadic annulus with radius \(2^q(|y|+h)\),
the physical number of lattice points is at most
\(C[2^q(|y|+h)]^4\). Summing the resulting
\(C|y|/[2^q(|y|+h)]\) proves the assertion.
For \(y=0\) the difference is zero.
\end{proof}

\subsubsection{Weak type and finite-\(p\) bounds}

\begin{proposition}[Uniform finite-\(p\) scalar estimate]
\label{dcz:uniform-Lp}
For every \(1<p<\infty\),
\begin{equation}\label{dcz:uniform-Lp-bound}
 \|D_i^\sigma D_j^\tau(L_h+\mu^2)^{-1}f\|_{p,h}
                                     \le C_p\|f\|_{p,h},
                           \qquad \mu\ge0.
\end{equation}
At zero mass on a finite periodic lattice the inverse is
the inverse on constants' orthogonal complement; the
second derivative annihilates its constant mode.
The same constant is valid on every compatible periodic
torus, independently of its size.
\end{proposition}

\begin{proof}
Write \(T\) for the infinite-lattice multiplier above.
The proof of \eqref{bg:weak-one} uses only the
\(L^2\) bound and its displayed kernel cancellation
condition. Both have just been established for
\(T\), with constants independent of \(h,\mu\).
On the infinite lattice take dyadic cubes of
sides \(2^nh\), \(n\ge0\). For compactly supported
input their averages eventually tend to zero,
so the same maximal-cube decomposition exists
at every positive threshold. A parent has
sixteen children of equal physical volume,
which improves its average bound to \(16\lambda\).
The bad part of a single-site cube is zero.
Its good-part estimate and the cancellation
sum outside enlarged cubes are therefore
exactly the two estimates in that proof.
They give weak-\((1,1)\) with the same uniform
parameters. The distribution-function
interpolation and adjoint duality proved
there apply unchanged, giving every
\(1<p<\infty\). Density defines these operators
on \(L^p\). This step reuses the proof rather
than assuming a continuous singular-integral
theorem for the lattice.

For \(\mu>0\), the kernel is absolutely summable
at every fixed mesh and mass: by
\eqref{dcz:heat-Lp},
\[
 \|k^{h,\mu}\|_{1,h}
       \le C\int_0^\infty
                       e^{-\mu^2t}(t+h^2)^{-1}\,dt<\infty.
\]
Periodize that kernel. To transfer its \(L^p\)
bound, truncate a periodic input to an increasing
union of fundamental cells. Absolute summability
makes the truncated convolution converge to the
periodic one on cells a large fixed distance
from the truncation boundary. Apply the
infinite-lattice estimate, divide its \(p\)-th
power by the number of cells, and let the
truncation grow. The proportion of boundary
cells tends to zero; a further limit of the
fixed distance removes the summable tail.
This proves the same periodic bound. Finally,
on any fixed finite torus let \(\mu\downarrow0\).
Every nonconstant Fourier mode converges to
the massless second-resolvent multiplier and
the constant second-derivative mode is zero.
The uniform bound passes to this finite
matrix limit.
\end{proof}

No assertion at \(p=\infty\) follows from this
proposition. The proof uses cancellation of the
bad parts, precisely where an absolute
supremum-norm kernel estimate would fail.

\subsubsection{Interior second derivatives and the flat Hodge estimate}

Let \(Q_r\subset Q_R\) be concentric physical fine
boxes, and assume that their distance is large
enough to contain the fixed two-step stencils.
Choose a sampled smooth cutoff \(\chi=1\) on
\(Q_r\), supported strictly inside \(Q_R\), with
\(|D\chi|\le C(R-r)^{-1}\) and
\(|D^2\chi|\le C(R-r)^{-2}\).
Extend \(\chi u\) by zero to the full lattice.
The product rule and \cref{dcz:uniform-Lp} give
the following precise local estimate.

\begin{corollary}\label{dcz:interior-Lp}
For \(1<p<\infty\),
\begin{equation}\label{dcz:interior-second-bound}
 \begin{split}
 \|D^2u\|_{p,Q_r}\le C_p\{&
       \|L_hu\|_{p,Q_R}
       +(R-r)^{-1}\|Du\|_{p,Q_R}\\
       &+(R-r)^{-2}\|u\|_{p,Q_R}\}.
 \end{split}
\end{equation}
The constants are independent of the fine mesh.
\end{corollary}

\begin{proof}
For a compactly supported \(v\),
\(D_i^\sigma D_j^\tau v
  =D_i^\sigma D_j^\tau L_h^{-1}(L_hv)\)
as Fourier multipliers. Thus
\(\|D^2v\|_p\le C_p\|L_hv\|_p\).
Apply this to \(v=\chi u\).
The exact product rule for \(L_h(\chi u)\)
has the term \(\chi L_hu\), terms containing
one difference of \(\chi\) and one of \(u\),
and terms containing two differences of
\(\chi\) and a shifted value of \(u\).
Their indicated supports lie in \(Q_R\).
This proves the estimate.
\end{proof}

For a compactly supported one-cochain \(v\), the
flat operator identity
\[
              \mathrm{d}^\sharp \mathrm{d}+DD^\sharp=L_hI
\]
and \cref{dcz:uniform-Lp} imply
\begin{equation}\label{dcz:whole-Hodge}
                   \|Dv\|_p
                     \le C_p(\|\mathrm{d}v\|_p+\|D^\sharp v\|_p).
\end{equation}
To see the operator types, write
\(v=L_h^{-1}(\mathrm{d}^\sharp \mathrm{d}v+DD^\sharp v)\)
and differentiate. Each resulting coefficient
is a second-resolvent multiplier of
\cref{dcz:uniform-Lp}; there are only finitely
many component pairs. Applying this to
\(\chi v\) proves
\begin{equation}\label{dcz:interior-Hodge}
 \|Dv\|_{p,Q_r}\le C_p\{
     \|\mathrm{d}v\|_{p,Q_R}+\|D^\sharp v\|_{p,Q_R}
                         +(R-r)^{-1}\|v\|_{p,Q_R}\}.
\end{equation}
These are full-lattice cutoff identities.
They impose no boundary condition on an
arbitrary proper determining region.

\subsubsection{Finite-time heat Sobolev and Morrey estimates}

We state the strict-exponent estimates needed
for bootstrap, so no critical endpoint is
used implicitly. For compactly supported \(u\)
put \(g=L_hu\). The exact finite-time identity is
\[
                         u=P_1u+\int_0^1P_tg\,dt,
                         \qquad P_t=e^{-tL_h}.
\]
Young's inequality and \eqref{dcz:heat-Lp}
give the following consequences.

\begin{proposition}\label{dcz:heat-Sobolev}
Let \(1<p\le q\le\infty\). If
\(2(1/p-1/q)<1\), then
\begin{equation}\label{dcz:heat-function-embedding}
                  \|u\|_q\le C_{p,q}(\|u\|_p+\|L_hu\|_p).
\end{equation}
If \(1/2+2(1/p-1/q)<1\), then
\begin{equation}\label{dcz:heat-gradient-embedding}
                 \|Du\|_q\le C_{p,q}(\|u\|_p+\|L_hu\|_p).
\end{equation}
The constants are uniform in \(h\). The corresponding
interior estimates follow by a cutoff, with the
commutator terms of \eqref{dcz:interior-second-bound}.
\end{proposition}

\begin{proof}
The fixed-time operators \(P_1,D P_1\) map \(L^p\)
to \(L^q\) with a uniform bound. The integrals
of their variable-time bounds are respectively
\[
 \int_0^1(t+h^2)^{-2(1/p-1/q)}\,dt,\qquad
 \int_0^1(t+h^2)^{-1/2-2(1/p-1/q)}\,dt.
\]
The stated strict inequalities make both
integrals bounded independently of \(h\).
The cutoff product rule proves the local version.
\end{proof}

\begin{proposition}[Uniform gradient Morrey estimate]
\label{dcz:gradient-Morrey}
If \(4<p<\infty\), then
\begin{equation}\label{dcz:gradient-Holder}
 \|u\|_\infty+\|Du\|_\infty
                +[Du]_{1-4/p,\rm loc}
                      \le C_p(\|u\|_p+\|L_hu\|_p).
\end{equation}
The seminorm uses distinct fine lattice points
at physical distance at most one. This estimate
also has the cutoff interior version.
\end{proposition}

\begin{proof}
The two supremum bounds follow from
\cref{dcz:heat-Sobolev}. For a fine lattice
translation \(z\), the derivative kernel obeys
\[
 \|\tau_zDp_t^h-Dp_t^h\|_{p'}
 \le C_p\min\{(t+h^2)^{-1/2-2/p},
                      |z|(t+h^2)^{-1-2/p}\}.
\]
The first bound follows from the triangle
inequality and \eqref{dcz:heat-Lp}; the second
telescopes one further difference along a
coordinate path. Distinct grid points have
\(|z|\ge h\). Split its time integral at
\(t=|z|^2\). The small-time contribution is
at most \(C_p|z|^{1-4/p}\), since
\(1/2+2/p<1\). The large-time contribution
has the same bound by direct integration.
The fixed-time term \(D P_1u\) has increment
at most \(C_p|z|\|u\|_p\), which is bounded
by \(C_p|z|^{1-4/p}\|u\|_p\) for
\(|z|\le1\). This proves the seminorm.
Cutoffs again use only the displayed first
and second difference commutators.
\end{proof}

For example, a bounded physical \(H^1\)
solution with a locally \(L^2\) Laplacian
first has gradient in \(L^3\) and itself
in every fixed finite \(L^q\). A subsequent
\(L^3\) forcing bound gives a bounded
solution and gradient in \(L^8\).
An \(L^8\) forcing bound then gives bounded
gradient. Once the equation supplies an
\(L^p\) Laplacian bound at any chosen finite
\(p>4/(1-\alpha)\), the same estimate gives a gradient
H\"older exponent at least \(\alpha\).
All these implications use strict inequalities
in \cref{dcz:heat-Sobolev}; none uses an
endpoint singular-integral estimate.

\subsubsection{All finite-\(p\) derivatives in an admitted Coulomb gauge}

We finish with the coefficient estimate needed
in the vector bootstrap. Suppose the real
connection on a fine box has an admitted
logarithmic Coulomb representative \(U=e^{ihA}\)
with
\[
 \|A\|_\infty\le C\kappa/R,\qquad
 D^\sharp A=0,\qquad
              \|F_U\|_\infty\le\kappa/R^2,
\]
on a box of size comparable to \(R\), with
the fixed larger collars of the current-free
local gauge construction. The existence of
that representative was proved in
\cref{ch:bg-local-coulomb}; its initial
construction does not use a current bound.

\begin{corollary}\label{dcz:Coulomb-Lp}
For every finite \(1<p<\infty\), on a fixed
smaller box,
\begin{equation}\label{dcz:Coulomb-Lp-bound}
                   \|DA\|_{p}\le C_p\kappa R^{4/p-2}.
\end{equation}
There is no derivative of curvature or
physical current in this assertion.
\end{corollary}

\begin{proof}
The exact four-link logarithm gives
\[
           F_U=\mathrm{d}A+E_h(A),\qquad
                         \|E_h(A)\|_\infty\le C\|A\|_\infty^2.
\]
Every nonlinear term has at least two
factors \(hA\) before the factor \(h^{-2}\)
is removed. This bound therefore requires
no derivative of \(A\). Insert this curl
identity and \(D^\sharp A=0\) into
\eqref{dcz:interior-Hodge}, using an
intermediate cutoff box. Its three
terms are bounded by
\[
 C_p\left\{\frac{\kappa}{R^2}R^{4/p}
       +\frac{\kappa^2}{R^2}R^{4/p}
       +R^{-1}\frac{\kappa}{R}R^{4/p}\right\}.
\]
The fixed small curvature ceiling bounds
\(\kappa^2\) by a constant times \(\kappa\).
This proves the claim. In a complex
\(C^1\) tube about this real representative,
the added potential and its first
derivative have the same finite-\(p\)
scale bounds by their stated supremum
norms. This last observation changes
coefficients only; it does not assert
unitarity of the complex connection.
\end{proof}

% END LAYERED EXPANSION discrete_calderon_zygmund.tex

% BEGIN LAYERED EXPANSION inverse/proper_strong_response.tex
\subsection{Strong responses on actual determining regions}
\label{ch:proper-strong-response}

The proper-region energy inverse controls a local physical
potential and its gradient. We upgrade it to pointwise and
H\"older control for bounded forces. The proof uses finite
\(L^p\) estimates in current-free Coulomb frames. Curvature
is never differentiated. The actual Wilson stencil, the
smoothed site projection and the nonlinear constraint
Hessian enter the argument in different ways.

Keep the actual determining regions, outward rows and
physical pairings of \cref{thm:npl-proper-gap}. Increase
their fixed separation and interior collar constants
once as specified below. All constants may depend on
the fixed \(L\), and H\"older constants may depend on
their exponent. They are independent of depth, mesh,
compatible volume and the shapes of the separated regions.
The frozen geometry and positive row weights are used
throughout.

\subsubsection{The forces and the local strong norm}

Let \(B_y\) be an assigned pillar of side \(\xi_y\).
At a nonbare pillar choose concentric interior boxes
of radii fixed multiples of \(\xi_y\). The largest
multiple \(R_*\) is fixed after the actual carrier
constant \(C_L\), with enough room for four successive
cutoffs and four complete carrier enlargements.
Take \(\rho\) larger than their fixed geometric margin.
These boxes meet at most two adjacent assigned levels;
all their scales are comparable within a fixed \(L\)
factor. Their fixed enlargements meet only a bounded
number of assigned pillars. A nonbare such box lies
inside the inner-plaquette region. The bare case has
\(\xi_y=\eta\) and will be estimated directly.

For finite \(4<p<\infty\) define
\begin{equation}\label{psr:force-norm}
 \mathcal F_p(f)=
    \sup_y\xi_y^{\,3-4/p}\|f\|_{p,B_y},\qquad
 \mathcal F_\infty(f)=\sup_b\mathsf L(b)^3|f(b)|.
\end{equation}
Replacing \(B_y\) by any of the fixed boxes changes
this norm by a constant. In particular
\(\mathcal F_p(f)\le C_p\mathcal F_\infty(f)\).
The cell energy amplitude is
\begin{equation}\label{psr:energy-amplitude}
 \mathcal E_{\rm cell}(v)
   =\sup_y\{\xi_y^{-1}\|v\|_{2,B_y}
                           +\|\nabla_Uv\|_{2,B_y^+}\}.
\end{equation}
One fine-step collars are included in \(B_y^+\).

The local gradient H\"older seminorm below is
covariant: two values based at \(x,x'\) are compared
after adjoint transport along a fixed coordinate-order
path inside the same interior box. Denote it by
\([\,\nabla_Uv\,]_{\alpha,y}^{U}\).
At bare cells it is the corresponding finite
one-step collar seminorm on present stencils, with
the zero-extension convention of \(X\).
No regularity across a rough exterior boundary or
between disconnected components is asserted.
Set
\begin{equation}\label{psr:strong-norm}
 \|v\|_{{\rm str},\alpha}
 =\sup_y\{\xi_y\|v\|_{\infty,B_y}
       +\xi_y^2\|\nabla_Uv\|_{\infty,B_y}
              +\xi_y^{2+\alpha}
                           [\,\nabla_Uv\,]_{\alpha,y}^{U}\}.
\end{equation}
This definition uses actual short transports and is
unchanged by real site gauges. For complex connections
in the tube considered below, the short transports
and their inverses have uniformly bounded norms.

\subsubsection{Scalar finite-\(p\) regularity without a current bound}

On a nonbare box use the current-free local Coulomb
construction. In its frame, with \(r=\xi_y\),
\begin{equation}\label{psr:Coulomb-data}
 U_\mu=e^{i\eta a_\mu},\qquad
 \|a\|_\infty\le C\kappa r^{-1},\qquad
       \|Da\|_q\le C_q\kappa r^{4/q-2}
                  \quad(1<q<\infty).
\end{equation}
These are the curvature-only conclusions of
\cref{dcz:Coulomb-Lp}, not the additional
current-dependent \(L^\infty\) derivative estimate.
The largest box still has small scaled curvature
after one fixed reduction of the curvature ceiling.

For site fields the exact covariant Laplacian is
\begin{equation}\label{psr:site-Laplacian}
 (\Delta_Uu)(x)=\eta^{-2}\sum_\mu
 \{2u(x)-\operatorname{Ad}U_\mu(x)u(x+\eta e_\mu)
     -\operatorname{Ad}U_\mu(x-\eta e_\mu)^{-1}u(x-\eta e_\mu)\}.
\end{equation}
Rescale the box by \(r\), putting \(h=\eta/r\).
In this frame it can be written exactly as
\begin{equation}\label{psr:scalar-decomposition}
             \Delta_U=L_h+\sum_\mu a_\mu^{(1)}D_\mu+b^{(0)},
 \qquad
 \|a^{(1)}\|_\infty\le C\kappa,\qquad
                           \|b^{(0)}\|_q\le C_q(\kappa+\kappa^2).
\end{equation}
Forward and backward shifted terms may both occur
in this notation. To verify it, put
\(E_+(x)=(\operatorname{Ad}e^{ih a(x)}-1)/h\)
and define \(E_-\) with the inverse exponential.
The difference from \(L_h\) is a first difference
of \(u\) with coefficients \(E_\pm=O(a)\), and
a zeroth coefficient
\(-[E_+(x)+E_-(x-he_\mu)]/h\).
Since \(E_-=-E_++O(h a^2)\), that coefficient is
bounded by first differences of \(a\) and \(Ca^2\).
This proves \eqref{psr:scalar-decomposition} with
no second derivative of the connection.

\begin{lemma}[Two local bootstraps]
\label{lem:psr-local-bootstraps}
In the rescaled fixed boxes the following estimates
are uniform in \(h\).

If \(\Delta_Uu=g\), \(g\in L^p\), \(4<p<\infty\), and
the outer \(H^1\) norm of \(u\) is at most \(B\),
then, for \(0<\alpha<1-4/p\),
\begin{equation}\label{psr:scalar-point-bootstrap}
 \|u\|_\infty+\|Du\|_\infty+[Du]_\alpha
                                      \le C_{p,\alpha}(B+\|g\|_p).
\end{equation}
The same conclusion holds if additional zeroth-order
operators with complete carrier radius \(C_L\)
have bounded local \(L^q\)-operator norms for every
finite \(q\), using the successive carrier enlargements.

If \(\Delta_Uu=D_U^\sharp F+g_0\), the same
outer energy is at most \(B\), \(F\in L^p\), and
\(g_0\) is bounded by \(B\) in all the finite
norms used, then
\begin{equation}\label{psr:divergence-bootstrap}
                 \|u\|_p+\|Du\|_p\le C_p(B+\|F\|_p).
\end{equation}
This is a finite-\(p\) divergence-source estimate.
It includes no \(p=\infty\) derivative assertion.
\end{lemma}

\begin{proof}
All cutoffs have fixed physical derivative bounds.
The energy heat embedding first gives \(u\in L^3\).
For example insert
\(u=P_1u+\int_0^1P_tD^\sharp Du\,dt\);
the derivative heat bound is integrable from
\(L^2\) to \(L^3\). Thus
\(b^{(0)}u\in L^2\), using \(b^{(0)}\in L^6\).
The other lower term is in \(L^2\).
The interior estimate \cref{dcz:interior-Lp}
and the strict heat embeddings
\cref{dcz:heat-Sobolev} give
\(D^2u\in L^2\), \(Du\in L^3\), and
\(u\in L^{12}\) on a smaller box.
Now \(b^{(0)}u\in L^3\), using its \(L^4\)
bound. A second cutoff gives \(D^2u\in L^3\),
then \(u\in L^\infty\) and
\(Du\in L^q\) for every \(q<12\).
If \(4<p\le8\), take \(q=p\).
If \(p>8\), first take \(q=8\) and use another
cutoff to obtain \(D^2u\in L^8\) and
\(Du\in L^\infty\). The final \(L^p\)
equation and \cref{dcz:gradient-Morrey}
prove \eqref{psr:scalar-point-bootstrap}.
The extra zeroth-order operator reads the
already controlled outer \(L^q\) field at each
step. The reserved complete carrier collars
give exactly these larger input boxes.
All Sobolev gains just used are strict.

For the divergence source, apply \(D(L_h+1)^{-1}\)
to the cutoff equation. Its divergence term
is \(D(L_h+1)^{-1}D^\sharp(\chi F)\),
bounded on each finite \(L^q\) by
\cref{dcz:uniform-Lp}. Its other terms have
the order-minus-one heat bound. Starting with
the outer energy, these terms are in \(L^2\);
the derivative output is in \(L^3\), since
\(1/2+2(1/2-1/3)<1\).
The field output is also in \(L^{12}\).
The next lower terms are in \(L^3\), giving
derivatives in every \(q<12\), in particular
\(q=\min\{8,p\}\). The field becomes bounded.
If \(p>8\), the order-minus-one terms from
\(L^8\) are bounded, while the divergence
term is estimated in the requested finite
\(L^p\). This proves \eqref{psr:divergence-bootstrap}.
Covariant \(D_U^\sharp F\) differs from the
flat divergence by the bounded product \(aF\).
Cutoff commutators add only \(F,Du,u\) on
the outer boxes. No difference of \(F\) is
estimated in an \(L^p\) norm.
\end{proof}

First-order Morrey, obtained from the same
order-minus-one translation kernels, gives
\(\|a\|_{C^\beta}\le C_q(\|a\|_q+\|Da\|_q)\)
for \(q>4\), \(\beta=1-4/q\).
Thus the ordinary derivative estimate in
\eqref{psr:scalar-point-bootstrap} also
controls the covariant derivative and its
transported H\"older seminorm. The short
transport differs from the identity by
\(C\kappa|x-x'|\) in these rescaled boxes.

\subsubsection{Strong columns of the actual site smoothing}

Let \(F=G_sq^\sharp\Xi^{-2}\), \(M=F^\sharp F\),
\(P=1-R=FM^{-1}F^\sharp\), with all the actual
site definitions. A column input supported at
top \(z\) is measured in its physical top
Hilbert norm. Its unit vector has value
\(\xi_z^{-2}e\), not value \(e\).

\begin{proposition}[Strong normalized site columns]
\label{prop:psr-site-columns}
For every \(0<\alpha<1\), a fixed \(m>0\)
can be chosen such that
\begin{equation}\label{psr:F-strong-columns}
 \begin{split}
 &\xi_y^2\|F1_zb\|_{\infty,B_y}
      +\xi_y^3\|D_UF1_zb\|_{\infty,B_y}\\
 &\qquad+\xi_y^{3+\alpha}
             [D_UF1_zb]_{\alpha,y}^{U}
               \le C_\alpha e^{-m d(y,z)}\|1_zb\|_{\rm top}.
 \end{split}
\end{equation}
No current premise or differentiated curvature is used.
\end{proposition}

\begin{proof}
The weighted balanced site form and normalized Gram
estimates give the column energy bounds
\[
 \|u\|_{2,\mathcal D_y}+
        \xi_y\|D_Uu\|_{2,\mathcal D_y}
                    \le Ce^{-m_0d(y,z)}\|1_zb\|_{\rm top},
 \qquad u=F1_zb,
\]
on each fixed enlarged box. Its equation is
\[
              \Delta_Uu=q^\sharp\Xi^{-2}1_zb-q^\sharp a_squ.
\]
The actual adjoint \(q^\sharp\) extends a
top colour as its transported density on its
own pillar. Its pointwise coefficient has norm
one for real fields; there is no fine factor
\(\eta^{-4}\). Also
\(|qu(t)|\le C\xi_t^{-2}\|u\|_{2,B_t}\).
Consequently the right side on \(\mathcal D_y\)
is bounded by
\(C\xi_y^{-4}e^{-m_0d(y,z)}\|1_zb\|_{\rm top}\).
The direct source vanishes unless the box meets
the source pillar, where its normalized amplitude
has this same bound. Scale \(u\) as
\(u'=\xi_y^2u\). Its outer \(H^1\) norm and
bounded forcing have the same exponential bound.
Apply \cref{lem:psr-local-bootstraps} with
a finite \(p>4/(1-\alpha)\) and scale back.

At a bare cell the physical atomic inequality
\(|u(x)|\le\eta^{-2}\|u\|_{2,B_x}\)
and its one-step neighbours prove the result
directly, including the finite collar seminorm.
This gives \eqref{psr:F-strong-columns}, after
one fixed decrease of the exponential rate.
\end{proof}

\begin{lemma}[The smoothed gauge term is a pointwise force]
\label{lem:psr-smoothed-force}
If \(\mathcal E_{\rm cell}(v)\le B\), then
\begin{equation}\label{psr:DPD-pointwise}
 \|D_UPD_U^\sharp v\|_{\infty,B_y}
       +\xi_y^\alpha[D_UPD_U^\sharp v]_{\alpha,y}^{U}
                                      \le C_\alpha\xi_y^{-3}B.
\end{equation}
The actual field row mass also satisfies
\par\noindent
\begin{minipage}{\linewidth}
\begin{equation}\label{psr:row-pointwise}
 \|Q^\sharp a_fQv\|_{\infty,B_y}\le C\xi_y^{-3}B,\qquad
 \|Q^\sharp\mu\|_{\infty,B_y}
                         \le C\xi_y^{-3}\sup_c\xi_c^3|\mu(c)|.
\end{equation}
\popQED\par
\end{minipage}
\end{lemma}

\begin{proof}
Use physical orthonormal top columns \(F_z\).
Their derivative \(L^2\) norm on \(B_y\) is
at most \(C\xi_y^{-1}e^{-md(y,z)}\).
Thus
\[
 |(D_UF)^\sharp v(z)|_{\rm ON}
 \le\sum_y C\xi_y^{-1}e^{-md(y,z)}
                                      \|v\|_{2,B_y}\le CB.
\]
The normalized Gram inverse has exponential
orthonormal entries and bounded row sums.
Its output coefficients are therefore \(O(B)\).
Sum the pointwise and H\"older derivative
columns \eqref{psr:F-strong-columns} to obtain
\eqref{psr:DPD-pointwise}.
The metric growth rate is chosen below half
the retained exponential rate.

For an actual field row,
\(|Qv(c)|\le C\xi_c^{-2}\|v\|_{2,\operatorname{car}c}
\le C\xi_c^{-1}B\).
Its coefficient at a fixed fine bond is at most
\(C(\eta/\xi_c)^4\). The physical adjoint
multiplies this by \(\eta^{-4}\xi_c^4\);
these factors cancel. Only a bounded number of
rows with comparable scales meet that bond.
Since \(a_f(c)=O(\xi_c^{-2})\), both assertions
in \eqref{psr:row-pointwise} follow. This
argument uses the full outward carrier, not
an individual endpoint value.
\end{proof}

\subsubsection{The exact Wilson stencil in Feynman form}

\begin{lemma}[A bounded curvature kernel]
\label{lem:psr-Wilson-stencil}
On every interior full stencil,
\begin{equation}\label{psr:Wilson-stencil}
        D^2A_{\Omega_0}(U)+D_UD_U^\sharp
                            =\Delta_U^{\rm comp}+\mathcal Z_U,
\end{equation}
where \(\Delta_U^{\rm comp}\) is the site
covariant Laplacian acting on each vector
component. The operator \(\mathcal Z_U\)
reads at most a fixed two-fine-step collar
and has row and column coefficient sums
at most \(C\kappa\xi_y^{-2}\) on a
scale-\(\xi_y\) box. In particular it is
a local bounded \(L^p\) zeroth-order
operator, uniformly for \(1\le p\le\infty\).
\end{lemma}

\begin{proof}
Use the exact plaquette insertion
\eqref{px:plaquette-hessian-exact}.
After division by the physical fine
measure, its remainder is a matrix on
four plaquette edges with each coefficient
bounded by \(C|P_p-1|\eta^{-2}\).
Hence this part is a bounded finite-shift
kernel of size \(C|F_p|\).

Let \(\mathfrak d_U=\mathrm{d}_U+E\), with the
actual curl convention in
\eqref{px:curl-convention-difference}.
Each coefficient of \(E\) is bounded by
\(C|P_p-1|/\eta\le C\eta|F_p|\);
each coefficient of \(\mathrm{d}_U\) is \(O(\eta^{-1})\).
Expand
\[
 \mathfrak d_U^\sharp\mathfrak d_U-\mathrm{d}_U^\sharp \mathrm{d}_U
                         =\mathrm{d}_U^\sharp E+E^\sharp \mathrm{d}_U+E^\sharp E.
\]
Multiply these stencil coefficients directly.
The first two products have coefficients
\(O(|F|)\) at a bounded number of shifted
sites; the last is \(O(\eta^2|F|^2)\).
Taking a finite difference of \(F\) is
unnecessary. This also proves the absolute
column bounds, since each product has a
bounded number of incoming as well as
outgoing stencil entries.

Finally set \(T_\mu=\operatorname{Ad}U_\mu\,\tau_\mu\).
The off-diagonal entries of the exact
covariant curl-divergence sum are
\[
 \eta^{-2}(T_\mu T_\nu^{-1}-T_\nu^{-1}T_\mu).
\]
The two paths differ by an elementary
plaquette. Their difference is a
two-step shift times a conjugate of
\(\operatorname{Ad}P-1\), with bounded
transport coefficients. It has size
\(C|F|\). The diagonal entries equal
\(\sum_\nu\eta^{-2}(2-T_\nu-T_\nu^{-1})\).
Combining these three finite-stencil
calculations proves \eqref{psr:Wilson-stencil}.
All identities hold bilinearly on the
complex continuation as well.
\end{proof}

The actual constraint Hessian
\(B_\omega=\Delta^{(2)}[\omega]\) has a
different carrier, of radius \(C_L\xi_y\).
The two-leg moment
\eqref{px:constraint-hessian-leg-moment}
and its symmetric counterpart give
\begin{equation}\label{psr:constraint-Lp}
 \|B_\omega v\|_{p,\mathcal D_y}
 \le C_L\gamma_\omega\xi_y^{-2}
                   \|v\|_{p,\mathcal D_y^{+C_L}},
 \qquad 1\le p\le\infty .
\end{equation}
Indeed its physical fine operator is
\(\eta^{-4}\sum_c\xi_c^4\omega(c)D^2q_c\).
For a fixed fine leg the absolute other-leg
sum is at most \(C\gamma_\omega\xi_y^{-2}\);
symmetry gives the column bound. Weighted
Schur or Jensen gives \eqref{psr:constraint-Lp}.
Scales on a complete carrier are comparable.
This proves an \(L^p\)-to-\(L^p\) estimate,
not an energy-to-pointwise assertion.

\subsubsection{The penalized and full row-constrained strong inverses}

Write the actual displayed precision as
\begin{equation}\label{psr:actual-precision}
 K_{U,\omega}=D^2A_{\Omega_0}(U)-B_\omega
                           +D_URD_U^\sharp+Q^\sharp a_fQ .
\end{equation}
The multiplier is small in
\(\gamma_\omega=\sup_c\xi_c^3|\omega(c)|\).
This is the local constraint Hessian, with
no independently named current completion.

\begin{proposition}[The proper strong inverse]
\label{prop:psr-strong-inverse}
For the real backgrounds of
\cref{thm:npl-proper-gap}, and finite \(4<p<\infty\),
\(0<\alpha<1-4/p\),
\begin{equation}\label{psr:strong-inverse-bound}
                  \|K_{U,\omega}^{-1}f\|_{{\rm str},\alpha}
                                         \le C_{p,\alpha}\mathcal F_p(f).
\end{equation}
In particular bounded physical forces
give this estimate for every \(\alpha<1\).
For the full row saddle
\[
             K_{U,\omega}v-Q^\sharp\mu=f,\qquad Qv=b,
\]
one has
\begin{equation}\label{psr:strong-row-saddle}
 \|v\|_{{\rm str},\alpha}+\sup_c\xi_c^3|\mu(c)|
            \le C_{p,\alpha}
                  \{\mathcal F_p(f)+\sup_c\xi_c|b(c)|\}.
\end{equation}
Removing the row mass from the saddle
adjusts its multiplier by \(-a_fb\)
and preserves this bound.
\end{proposition}

\begin{proof}
The exponential energy inverse gives
\(\mathcal E_{\rm cell}(v)\le C\mathcal F_p(f)\);
on a pillar the dual forcing norm is
\(\xi_y\|f\|_2\le C\mathcal F_p(f)\).
Use \(R=1-P\) and \eqref{psr:Wilson-stencil}
to write its interior equation as
\begin{equation}\label{psr:vector-local-equation}
 \Delta_U^{\rm comp}v
 =f+D_UPD_U^\sharp v-Q^\sharp a_fQv
                                      +B_\omega v-\mathcal Z_Uv.
\end{equation}
The first two nonlocal smoothing and row
terms have bounded pointwise amplitude
\(C\xi_y^{-3}\mathcal E_{\rm cell}(v)\).
The remaining two operators have the
local \(L^q\) bounds just proved.
Set \(v'=\xi_yv\), \(f'=\xi_y^3f\).
Its outer \(H^1\) bound is the cell
energy amplitude and its scalar component
equation is precisely the one in
\cref{lem:psr-local-bootstraps}.
The four nested boxes retain every
constraint carrier at each step. That
lemma proves the required local bound.
Ordinary and covariant derivative
seminorms are equivalent by
\eqref{psr:Coulomb-data} and first-order
Morrey. Bare cells follow directly from
their physical atomic energy bounds;
only a fixed number of neighbouring
fine values is read there.

For the row saddle, the normalized
Schur construction of
\cref{prop:scr-schur} applies also to
this precision: its proof uses its
energy inverse, exact row lift and
weighted coefficients, which
\eqref{psr:actual-precision} has.
The array kernel of its unnormalized
inverse is bounded by
\(C\xi_c^{-3}\xi_d e^{-md(c,d)}\).
Thus its action on \(b\) has
\(\xi^3\) amplitude at most
\(C\sup_d\xi_d|b(d)|\).
The other multiplier contribution is
the transpose minimizing lift applied
to \(f\). Its energy column proof
\eqref{scr:column-current-duality}
requires only
\(\sup_y\xi_y\|f\|_{2,B_y}\le C\mathcal F_p(f)\);
it gives the same pointwise multiplier
bound. The row-lift columns and energy
inverse give the corresponding cell
energy estimate for \(v\). Finally
\eqref{psr:row-pointwise} bounds
\(Q^\sharp\mu\) as an additional
point force. Apply the same local
bootstrap to obtain
\eqref{psr:strong-row-saddle}.
\end{proof}

\subsubsection{Finite-\(p\) restoration of a force}

The adjoint restoration is not bounded
on fine \(L^\infty\) in general.
The finite-\(p\) substitute is sufficient.

\begin{proposition}\label{prop:psr-force-restoration}
For every finite \(4<p<\infty\),
\begin{equation}\label{psr:Ssharp-finite-p}
                    \mathcal F_p(S^\sharp f)\le C_p\mathcal F_p(f).
\end{equation}
In particular a bounded physical force
has a restored force with the correct
local amplitude in every finite \(L^p\).
\end{proposition}

\begin{proof}
Put \(u_s=G_sD_U^\sharp f\). The weighted
site energy form gives the four cell
bounds for \(G_s\), \(D_UG_s\),
\(G_sD_U^\sharp\) and \(D_UG_sD_U^\sharp\), with scale factors
\(\xi_y\xi_z,\xi_z,\xi_y,1\),
respectively, times \(Ce^{-md(y,z)}\).
This is the same weighted form and
distance argument as for the field
energy inverse. The site gap supplies
its mass term.
Choose the fixed separation so the
metric growth and one scale-ratio cost
\(\xi_y/\xi_z\) together are less than
half the retained decay rate.
Since \(\|f\|_{2,B_z}\le C\xi_z^{-1}
\mathcal F_p(f)\), summing these columns gives
\[
 \|u_s\|_{2,\mathcal D_y}
          +\xi_y\|D_Uu_s\|_{2,\mathcal D_y}
                                  \le C\mathcal F_p(f).
\]
Its local equation is
\(\Delta_Uu_s=D_U^\sharp f-q^\sharp a_squ_s\).
The average term is a bounded point
forcing of amplitude
\(C\xi_y^{-4}\mathcal F_p(f)\).
Rescale \(u_s'=\xi_y^2u_s\),
\(f'=\xi_y^3f\), and apply the
divergence-source estimate
\eqref{psr:divergence-bootstrap}.
It follows that
\[
 \xi_y^{3-4/p}\|D_Uu_s\|_{p,\mathcal D_y}
                                  \le C_p\mathcal F_p(f).
\]

Exactly \(S^\sharp f=f-D_URu_s\).
Write \(Ru_s=u_s-FM^{-1}F^\sharp u_s\).
The just proved \(L^2\) cell amplitude
of \(u_s\), the exponential \(F\)
columns and the normalized Gram
inverse bound its last derivative
pointwise by
\(C\xi_y^{-3}\mathcal F_p(f)\).
The displayed finite-\(p\) derivative
bound controls the other term.
This proves \eqref{psr:Ssharp-finite-p}.
The proof used the divergence-source
Riesz estimate only at finite exponents.
\end{proof}

\subsubsection{Physical critical pairs and the actual Landau saddle}

Let \(H_L=D^2A_{\Omega_0}(U)-B_\omega\).
At a physical critical pair full stationarity is
\(j_U=Q^\sharp\omega\). Let \(\mathfrak q_U(a)\)
be the actual nonlinear field-mean coordinate,
with \(D\mathfrak q_U(0)=Q\).
Its actual gauge cocycle gives
\[
             D\mathfrak q_U(a)\Gamma_a\lambda
                           =\mathrm{d}(a)\tau_a\lambda,\qquad
          \tau_0=q,\quad \mathrm{d}(0)=\bar D_{\rm ext}.
\]
The extended outward endpoint values are part of this
identity. Fine gauge invariance and onto \(q\) give
\(\bar D_{\rm ext}^\sharp\omega=0\).
For \(q\lambda=0\), differentiate the displayed
identity in an arbitrary field direction.
Its \(\mathrm{d}'(a)\tau_0\lambda\) term vanishes, and its
\(\mathrm{d}(0)\tau'_0\lambda\) term pairs to zero with
\(\omega\). Differentiating the action Ward identity
produces the same gauge-generator acceleration;
full stationarity cancels it against the constraint
term. Therefore
\begin{equation}\label{psr:physical-Lagrangian-Ward}
                          H_LD_U\lambda=0\qquad(q\lambda=0).
\end{equation}
This proves the full bilinear null identity, not
only its restriction to the constraint kernel.

Since \(qG_sR=0\), it follows that
\(K_{U,\omega}D_UG_sR=D_UR\).
Its inverse and row lift have the physical
horizontality identities used in
\cref{ch:layered-coarse-current-readout}.
For a force represented by \(f\) on the actual
Landau slice, use the extension \(S^\sharp f\);
it agrees there with \(f\).
The full row saddle with this source is horizontal,
by the Ward identity and \(SD_UG_sR=0\).
It is consequently the actual physical Landau
saddle after removal of the row mass.
\Cref{prop:psr-force-restoration} and
\eqref{psr:strong-row-saddle} give
\begin{equation}\label{psr:physical-Landau-strong}
 \|v\|_{{\rm str},\alpha}+\sup_c\xi_c^3|\mu(c)|
       \le C_{p,\alpha}
                  \{\mathcal F_p(f)+\sup_c\xi_c|b(c)|\}.
\end{equation}
There is no fine \(L^\infty\) assertion for
\(S^\sharp f\) in this argument.

\subsubsection{Complex tubes, small volumes and scope}

The same full penalized and full row-saddle
strong estimates hold in a sufficiently small
weighted \(C^1\) tube \eqref{cml:tube} about a
real reference. Here is the required weighted
energy justification. Repeat the balanced site
form with a spatial exponential weight on both
the source and the output. Its normalized
perturbation is \(O(\varepsilon_*+|s|)\).
The normalized \(F\), Gram and \(R\) therefore
have the same weighted bounds by Neumann
series from the real reference. The full
penalized field form changes by
\(C(\varepsilon_*+\gamma_\omega+|s|)\).
Its real reference gap gives the complex
weighted inverse and all cell energy
estimates used above. This argument concerns
the full penalized form; it does not assert
a new spatial exponential theorem for an
arbitrary moving-slice saddle.

Use the real reference's Coulomb frame on
each interior box. The complex perturbation
has bounded potential and first derivative
in that frame by the \(C^1\) tube.
The logarithm of its product with the
real reference links therefore has
\(\|a\|_\infty\le C(\kappa+\varepsilon_*)/\xi\)
and
\(\|Da\|_q\le C_q(\kappa+\varepsilon_*)\xi^{4/q-2}\).
All scalar and stencil estimates continue
with these bounds. Their proofs use
bilinear algebra and absolute estimates,
not positivity at a complex background.
The strong conclusions follow unchanged.
For actual complex physical critical
pairs the analytic Ward identity
\eqref{psr:physical-Lagrangian-Ward}
also proves \eqref{psr:physical-Landau-strong}.
At an arbitrary noncritical complex
pair only the full penalized and full
row-constrained strong conclusions have
been established here.

If a local box does not embed in a small
torus, lift all actual rows, data and
backgrounds to a finite periodic cover
with periods exceeding \(8R_*\max\xi\).
Site and Gram inverses preserve the
deck-invariant subspace by uniqueness;
so do the full field and saddle
inverses. The cover has a bounded
number of sheets because the original
periods contain at least two largest
blocks. Apply the local estimates on
the cover and descend. For a single
column sum its finitely many deck
images; each distance is at least the
original torus distance, and one fixed
decrease of the rate absorbs the
bounded multiplicity.

We have proved bounded-force pointwise
and local gradient H\"older responses
on the actual multiscale determining
regions. The estimates include the
physical critical Landau saddle and
small complex \(C^1\) continuations in
the ranges explicitly stated. They do
not infer a pointwise singular Green
kernel from a bounded-force bound.
No nonlinear proper-region critical
background or independently named
current completion has been
constructed by these linear estimates.

% END LAYERED EXPANSION inverse/proper_strong_response.tex

% BEGIN LAYERED EXPANSION inverse/actual_current_completion.tex
\subsection{An explicit current completion of the projected Hessian}
\label{ch:actual-current-completion}

The invariant physical Hessian is determined by its
Landau projection. Its continuation to an independent
current requires an additional convention. We fix that
convention by the second jet of an explicit finite
gauge chain. This identifies the current correction
with the symmetric form already estimated, fixes the
order of its two propagator variants, and gives their
bounded-force regularity on the actual determining
geometry.

Throughout this section the action is the sum of inner
Wilson plaquettes in \(\Omega_0\). Its physical current
is \(j_U\), its ambient Hessian is \(\mathscr D_U\),
and transposes use the physical bilinear pairings.
All means are re-centered at \(U\). The field row
\(Q\) is the derivative per unit row length of the
actual straightened mean; it is not the derivative
of the plain mean. The site map \(q\) is its actual
own-pillar gauge average, with the full extension at
outward endpoints. Increase the fixed complete-carrier
collars and separation once for the finitely many
scalar stages below, independently of depth.

\subsubsection{The physical projection and its normalization}

Write
\begin{equation}\label{acc:site-operators}
 \begin{split}
 G_s&=(D^\sharp D+q^\sharp a_sq)^{-1},\quad
 F=G_sq^\sharp\Xi^{-2},\quad M=F^\sharp F,\\
 R&=1-FM^{-1}F^\sharp,\qquad
 \Lambda=G_sRD^\sharp,\qquad S=1-D\Lambda .
 \end{split}
\end{equation}
The site-top scale \(\Xi\) differs from the field-row
length multiplication \(\mathsf m\).
The intrinsic projection in this formula is precisely
\[
 R=1-G_sq^\sharp(qG_s^2q^\sharp)^{-1}qG_s .
\]
Indeed the two factors \(\Xi^{-2}\) in \(M\) cancel
with its inverse in their indicated order. This is
an algebraic normalization, not a different projection.

The exact identities are
\begin{equation}\label{acc:projection-algebra}
 \begin{gathered}
 qG_sR=0,\qquad D^\sharp DG_sR=R,\qquad
 RD^\sharp S=0,\qquad S^2=S,\\
 QD=\bar D_{\rm ext}q,\qquad QS=Q,\qquad
                               SDG_sR=0 .
 \end{gathered}
\end{equation}
Thus \(\lambda=\Lambda A\) is the unique
zero-average parameter for which
\(RD^\sharp(A-D\lambda)=0\).
For uniqueness, \(D^\sharp D\) maps \(\ker q\)
onto \(\operatorname{ran}R\), with inverse
\(G_sR\). The site gap makes its kernel zero.
These statements hold on the actual proper graph,
including all outward rows and bare collars.

The physically projected Hessian is
\begin{equation}\label{acc:physical-projection}
                         \Delta_{\pi,U}^{\rm phys}
                                           =S^\sharp\mathscr D_US .
\end{equation}
This determines the correction at \(J=j_U\).
It alone does not determine a linear map at every
independent \(J\): the physical current graph is
not an open subset of the field-current product
domain. We use the following finite-gauge
convention to make that continuation explicit.

\subsubsection{A finite gauge chain fixes the independent-current map}

For a fine vector \(A\) put \(\lambda_A=\Lambda A\).
At a bond \(b\), let
\(\ell_A=\lambda_A(b_-)\) and
\(r_A=\operatorname{Ad}U_b\,\lambda_A(b_+)\).
Define the germ
\begin{equation}\label{acc:finite-gauge-chain}
 \Pi_U(A)(b)=\frac1{i\eta}
       \log\left(e^{-i\ell_A}e^{i\eta A(b)}e^{ir_A}\right).
\end{equation}
It is the relative potential of the actual finite
gauge transform by \(e^{-i\lambda_A}\).
Only its second jet at zero is used.
Its first derivative is \(S\). No nonlinear
Landau-slice assertion about \(\Pi_U(A)\)
is needed or made.

Set \(T_U[u,v]=D^2\Pi_U(0)[u,v]\), and define
the independent-current completion by
\begin{equation}\label{acc:canonical-current-definition}
          \langle u,\Delta'_{\pi,U}[J]v\rangle
                                  :=\langle J,T_U[u,v]\rangle .
\end{equation}
This is our explicit convention for the named
independent-current operator. It retains the
physical projection \eqref{acc:physical-projection}.
It does not assert uniqueness among all possible
current-linear extensions with that physical value.

\begin{proposition}[The explicit gauge-chain jet]
\label{prop:acc-current-identification}
The completion \eqref{acc:canonical-current-definition}
equals exactly \(\mathcal B_U[J]\) of
\eqref{px:symmetric-current-correction}, for every
independent real or complex current. More explicitly,
\begin{equation}\label{acc:explicit-current-map}
 \begin{split}
 W_U[J]\lambda(b)&=\frac i2
      [\lambda(b_-)+\operatorname{Ad}U_b\lambda(b_+),J(b)],\\
 \Delta'_{\pi,U}[J]
   &=W_U[J]\Lambda+\Lambda^\sharp W_U[J]^\sharp\\
   &\quad-\tfrac12\Lambda^\sharp
             (D^\sharp W_U[J]+W_U[J]^\sharp D)\Lambda .
 \end{split}
\end{equation}
In particular it is symmetric and linear in \(J\).
The identity does not assume \(D^\sharp J=0\).
\end{proposition}

\begin{proof}
The second-order logarithm of the three
exponentials in \eqref{acc:finite-gauge-chain}
has the quadratic part
\[
 \frac i2[A,\ell_A+r_A]
                          -\frac{i}{2\eta}[\ell_A,r_A].
\]
Polarizing gives
\begin{equation}\label{acc:second-jet}
 \begin{split}
 T_U[u,v](b)={}&\frac i2
       \{[u(b),\ell_v+r_v]+[v(b),\ell_u+r_u]\}\\
       &-\frac{i}{2\eta}
                       \{[\ell_u,r_v]+[\ell_v,r_u]\}.
 \end{split}
\end{equation}
The apparent inverse spacing in the second term
does not produce an inverse-spacing norm loss.
Since \(D\lambda=(\ell-r)/\eta\), cyclicity
of the colour trace rewrites its pairing as
\begin{equation}\label{acc:undifferentiated-current-form}
 \begin{split}
 \langle J,T_U[u,v]\rangle={}&
  \langle u,W_U[J]\Lambda v\rangle+
  \langle W_U[J]\Lambda u,v\rangle\\
  &-\tfrac12\langle D\Lambda u,W_U[J]\Lambda v\rangle
   -\tfrac12\langle W_U[J]\Lambda u,D\Lambda v\rangle .
 \end{split}
\end{equation}
For example
\([\ell_u-r_u,\ell_v+r_v]+
  [\ell_v-r_v,\ell_u+r_u]
 =2([\ell_u,r_v]+[\ell_v,r_u])\).
This verifies both the factor one half and the
physical power of \(\eta\).
Taking bilinear transposes in
\eqref{acc:undifferentiated-current-form}
proves \eqref{acc:explicit-current-map}.
No divergence identity for the independent
current was inserted.
\end{proof}

Gauge invariance gives
\(A_{\Omega_0}(e^{i\eta A}U)
 =A_{\Omega_0}(e^{i\eta\Pi_U(A)}U)\).
The second chain rule at zero therefore proves
\begin{equation}\label{acc:on-current-identity}
             \mathscr D_U-\Delta'_{\pi,U}[j_U]
                                            =S^\sharp\mathscr D_US .
\end{equation}
This is an exact physical identity even when
\(U\) is not constrained critical.
It also follows from
\(\mathscr D_UD=W_U[j_U]\).
The bilinear chain calculation extends
holomorphically in the complex field.
After taking its jet, the coefficient formula
defines the form on all finite fine vectors;
an arbitrary Gaussian vector need not lie in
a nonlinear gauge chart.

The current dependence is now fully specified:
\(D_J\Delta'_{\pi,U}[J]\dot J
=\Delta'_{\pi,U}[\dot J]\), and higher current
derivatives vanish. A definition using a
nonlinear gauge parameter with the same first
jet could add a term pairing \(D^\sharp J\)
with its second gauge jet. Such a term is
invisible on physical currents. Formula
\eqref{acc:canonical-current-definition}
fixes this freedom before using the physical
current identity.

\subsubsection{The two named variants in their actual order}

Let \(C_U^{(2)}\) be the quadratic part per
unit length of the re-centered straightened
field mean. Thus its group-angle logarithm
has expansion
\(\mathsf m(QA+C_U^{(2)}(A)+\cdots)\).
Define
\begin{equation}\label{acc:constraint-current-form}
       \langle A,\Delta_U^{(2)}[\omega]A\rangle
                          =2\langle\omega,C_U^{(2)}(A)\rangle .
\end{equation}
In particular this is the quadratic part of the
straightened mean. With the row pairing held fixed,
multiplication by \(\mathsf m\) converts a multiplier
for the group angle to a per-length multiplier.
The conversion to the counting pairing is given below.
All transposes below include the physical
fine and row volumes.

The precise order is
\begin{equation}\label{acc:named-variants}
 \begin{split}
 G_0^{-1}&=\mathscr D_U+DRD^\sharp+Q^\sharp a_fQ,\\
 G^{-1}&=G_0^{-1}-\Delta'_{\pi,U}[J],\\
 H&=GQ^\sharp(QGQ^\sharp)^{-1},\qquad \omega_*=H^\sharp J,\\
 G_1^{-1}&=G^{-1}-S^\sharp\Delta_U^{(2)}[\omega_*]S,\\
 H_1&=G_1Q^\sharp(QG_1Q^\sharp)^{-1}.
 \end{split}
\end{equation}
\paragraph{The same completion in the propagator convention.}
There is a sign change between the covariant forward
difference used for the propagator construction and the
infinitesimal gauge generator of \eqref{px:gauge-generator}.
In per-length units the forward difference at a bond is
\((r-\ell)/\eta\), while the gauge generator is
\((\ell-r)/\eta\). The site inverse and \(R\) are unchanged
by this change of sign when the site average and mass are
the same. Hence the site parameter in the forward-difference
formula for the projection is \(v=-\lambda_A\).
In that convention the current term is the pairing of
\(J\) with
\[
        i[v_-,A]-i[A,v_+]-i[v_-,D_{\rm forw}v].
\]
Here \(v_-\) and \(v_+\) are its values at the initial
and transported terminal endpoints and
\(D_{\rm forw}v=(v_+-v_-)/\eta\). Substitution gives
\[
             i[A,\ell_A+r_A]-\frac{i}{\eta}[\ell_A,r_A],
\]
which is \(T_U[A,A]\) from \eqref{acc:second-jet}.
The fine pairing in both formulas is the trace pairing
with the bare cell volume. By polarization the
two current completions agree also for independent complex
\(J\) on the common graph with the same site average.
For the separate correction from the second derivative
of the mean, let \(C^{(2)}\) be the quadratic term in the
per-length row coordinates of the same mean. The raw
quadratic term is \(\mathsf m C^{(2)}\). The raw lift
is \(\mathsf H\mathsf m^{-1}\), with counting-row
transpose \(\mathsf m^3\mathsf H^*\), where the star
uses row volumes \(\mathsf m^4\). Therefore its contraction
with the current is
\[
  (\mathsf m C^{(2)}(A),\mathsf m^3\mathsf H^*J)_{\rm count}
        =(C^{(2)}(A),\mathsf H^*J)_{\rm vol}.
\]
The two corrections thus have the same readouts when
the site maps, row maps, and local domains have been
identified. The calculation is unchanged on the
projected field with \(A\) replaced by \(SA\).

The additional quadratic form is exactly
\[
 -2\langle J,H C_U^{(2)}(SA)\rangle
                  =-\langle A,S^\sharp
                         \Delta_U^{(2)}[H^\sharp J]SA\rangle .
\]
Thus it is precisely the second explicit variant.
The \(H\) inside the second subtraction is formed
from the first \(G\), and is not \(H_1\).
Nor does \(C_U^{(2)}(SA)\) mean the second
derivative of the nonlinear composite
mean with \(\Pi_U(A)\).

With the convention
\eqref{acc:canonical-current-definition}, all
operators in \eqref{acc:named-variants} are
the already constructed explicit operators.
In particular
\cref{thm:scr-induced-current} proves
\begin{equation}\label{acc:induced-multiplier}
          \sup_c\xi_c^3|\omega_*(c)|\le C\gamma_J,\qquad
                      \gamma_J=\sup_b\mathsf L(b)^3|J(b)|.
\end{equation}
At a small full-current radius this
already admits the second subtraction.
The support refinement below requires
only a small current on bonds touching
\(\Omega_1\) for the propagators,
while retaining the full \(\gamma_J\)
in this all-row amplitude estimate.

\subsubsection{Gauge defects and the meaning of on shell}

Put \(\delta J=J-j_U\). Equations
\eqref{acc:projection-algebra} and
\eqref{acc:on-current-identity} give
\begin{equation}\label{acc:off-current-defect}
 \begin{split}
 G^{-1}DG_sR&=DR-\Delta'_{\pi,U}[\delta J]DG_sR,\\
 G_{(1)}DR&=DG_sR+
              G_{(1)}\Delta'_{\pi,U}[\delta J]DG_sR .
 \end{split}
\end{equation}
The parenthesized index permits either
variant. The second subtraction vanishes
on \(DG_sR\) because \(SDG_sR=0\).
Transposition gives its corresponding
row-lift slice defect. On shell that
defect is zero; no independent current
is declared divergence free.

Only bonds touching \(\Omega_1\) matter
for the first current correction.
Write \(J^{\rm t}\) for its restriction
to those bonds and \(J^{\rm b}=J-J^{\rm t}\).
Set
\begin{equation}\label{acc:touching-current-radius}
 \gamma_{\rm t}
 =\sup_{b:\,b\text{ touches }\Omega_1}
                               \mathsf L(b)^3|J(b)|.
\end{equation}
The supremum is zero if that set is empty.
Indeed \(\Lambda A\in\ker q\), so it
vanishes at all bare sites. Both
\(W_U[\delta J]\Lambda A\) and
\(D\Lambda A\) therefore vanish on
bonds with neither endpoint in
\(\Omega_1\).
Consequently equality \(J=j_U\) on
bonds touching \(\Omega_1\) already
implies
\(\Delta'_{\pi,U}[J]=\Delta'_{\pi,U}[j_U]\).
In particular
\(\Delta'_{\pi,U}[J]=\Delta'_{\pi,U}[J^{\rm t}]\).

The same support statement suffices
for the second variant. Let \(P_0\)
restrict row data to the bare rows.
The exact identity is \(Q_0H=P_0\).
Thus a datum supported on a nonbare
row lifts to a vector which must
vanish on bare bonds. Under the
touching-bond on-shell condition,
\(H^\sharp(J-j_U)\) has only bare
row components. More precisely, because
the physical fine and bare-row measures
are both \(\eta^4\), that same restriction
identity gives
\begin{equation}\label{acc:bare-current-readout}
           \omega_*[J]=\omega_*[J^{\rm t}]+\iota_0J^{\rm b},
\end{equation}
where \(\iota_0\) identifies a bare bond
with its level-zero row.
The actual bare
mean is linear, so its quadratic
part is zero. These components do
not contribute to
\(\Delta_U^{(2)}[\omega_*]\).
This retains the on-shell condition
on bonds touching \(\Omega_1\);
it does not strengthen it to a
condition on every bare bond.
Applying \eqref{acc:induced-multiplier}
to \(J^{\rm t}\) gives
\begin{equation}\label{acc:nonbare-induced-radius}
           \sup_{c\text{ nonbare}}\xi_c^3|\omega_*(c)|
                                             \le C\gamma_{\rm t}.
\end{equation}
Thus \(G,H,G_1,H_1\) are independent
of \(J^{\rm b}\). Their admission and
strong bounds use only a small
\(\gamma_{\rm t}\). Their all-row
induced current still has the full
\(\gamma_J\) bound
\eqref{acc:induced-multiplier}; arbitrary
bare currents are not thereby bounded.

If \(U\) is an actual physical critical
point with
\(j_U=Q^\sharp\omega_{\rm phys}\),
then, for the full physical current,
\(H^\sharp j_U=\omega_{\rm phys}\).
The physical Lagrangian Ward identity
\eqref{psr:physical-Lagrangian-Ward}
gives
\[
 S^\sharp(\mathscr D_U-
              \Delta_U^{(2)}[\omega_{\rm phys}])S
                  =\mathscr D_U-\Delta_U^{(2)}[\omega_{\rm phys}].
\]
Therefore the named second precision
is exactly
\begin{equation}\label{acc:physical-second-precision}
 G_1^{-1}=\mathscr D_U-
       \Delta_U^{(2)}[\omega_{\rm phys}]
                              +DRD^\sharp+Q^\sharp a_fQ .
\end{equation}
The same conclusion holds under the
on-shell condition just described,
since the remaining induced-current
change is bare-only. This identifies
the actual physical second variant
with the local Lagrangian precision,
not merely their restrictions to
horizontal vectors.

\subsubsection{Regularity at independent currents}

The current form
\eqref{acc:undifferentiated-current-form}
and weighted restoration give
\[
 |\langle u,\Delta'_{\pi,U}[J]v\rangle|
                       \le C\gamma_{\rm t}\|u\|_X\|v\|_X .
\]
Its conjugated version has that same
bound. The first and second variants
therefore have the energy and
exponential cell estimates already
proved, for a fixed ceiling on
\(\gamma_{\rm t}\). Only the nonbare
part of \(\omega_*\) enters the
second subtraction.
For real pairs their forms are
positive. On a small complex \(C^1\)
tube the inverses are bounded analytic
perturbations of the real references.
No physical-current identification is
used for this admission.

We next justify the pointwise and
gradient H\"older bounds for these
independent-current variants. Keeping
the apparent derivative of \(J\)
out of this proof is essential.
Use the normalized \(\mathcal F_p\)
and cell energy amplitudes of
\cref{ch:proper-strong-response}.
All exponents \(p\) in the following
argument are finite.

\begin{lemma}[Current terms from cell energy]
\label{lem:acc-current-bootstrap}
Suppose \(\mathcal E_{\rm cell}(v)\le B\).
Then
\begin{equation}\label{acc:first-current-F8}
                \mathcal F_8(\Delta'_{\pi,U}[J]v)
                                             \le C\gamma_{\rm t}B .
\end{equation}
For a row current with small
\(\gamma_{\omega,\rm nb}
 =\sup_{c\text{ nonbare}}\xi_c^3|\omega(c)|\), the
second form initially has the bound
\begin{equation}\label{acc:second-current-F3}
         \mathcal F_3(S^\sharp\Delta_U^{(2)}[\omega]Sv)
                                             \le C\gamma_{\omega,\rm nb} B .
\end{equation}
Here \(\mathcal F_3\) means the same
local physical formula as
\eqref{psr:force-norm} at exponent \(3\).
If in addition the normalized local
\(W^{2,3}\) norm of \(v\) is bounded
by \(B_1\), both current terms belong
to every finite \(\mathcal F_p\),
\(p>4\), with bound
\(C_p(\gamma_{\rm t}+\gamma_{\omega,\rm nb})(B+B_1)\).
\end{lemma}

\begin{proof}
By the exact support identities, replace
\(J\) by \(J^{\rm t}\) and discard
bare row components of \(\omega\).
Set \(\lambda=\Lambda v\).
The local divergence has site \(L^2\)
amplitude of grade two bounded by \(B\).
Its \(P\) part has the pointwise
grade-two bound: use
\(F^\sharp D^\sharp v=(DF)^\sharp v\),
whose normalized top coefficients
are \(O(B)\), and sum the strong
site columns.
Thus \(RD^\sharp v\) has the same
local \(L^2\) amplitude.
Site energy columns, with one
scale-ratio reserve, give
\[
 \xi_y^{-2}\|\lambda\|_{2,\mathcal D_y}
        +\xi_y^{-1}\|D\lambda\|_{2,\mathcal D_y}\le CB.
\]
Since \(\Delta_U\lambda=RD^\sharp v\),
the first scalar \(L^2\) cutoff step
gives normalized \(W^{2,2}\) bounds.
The strict heat embeddings imply
that \(\lambda\) lies in every finite
local \(L^q\) and \(D\lambda\) lies
in local \(L^3\), with the physical
dimensions \(1\) and \(\xi^{-1}\).
The same \(L^3\) bound holds for \(v\).

Group the exact correction as
\begin{equation}\label{acc:grouped-current-operator}
 \Delta'_{\pi,U}[J]v
  =W_U[J]\lambda+
       D R G_s\{W_U[J]^\sharp(v-\tfrac12D\lambda)
                           -\tfrac12D^\sharp W_U[J]\lambda\}.
\end{equation}
The site force
\(a=W_U[J]^\sharp(v-D\lambda/2)\)
has local \(L^3\) grade-four
amplitude at most \(C\gamma_{\rm t}B\).
The site energy estimate for \(G_sa\)
and the scalar cutoff bootstrap
through \(L^3\) give a derivative
in \(L^8\) of grade three.
Subtracting its \(P\) part costs a
bounded pointwise derivative, by
the site columns and the Gram
inverse. Hence \(DRG_sa\in\mathcal F_8\).
The other force
\(h=W_U[J]\lambda\) is in every
finite \(\mathcal F_q\), because
\(\lambda\) is in every finite
local \(L^q\). Exactly
\[
                 DRG_sD^\sharp h=h-S^\sharp h .
\]
The finite-\(q\) restoration bound
\eqref{psr:Ssharp-finite-p}, at \(q=8\),
controls this term and the direct
\(h\) term. This proves
\eqref{acc:first-current-F8}.

For the second form, \(Sv=v-D\lambda\)
has its normalized \(L^3\) bound.
The actual two-leg Schur estimate
\eqref{psr:constraint-Lp} gives
\(\Delta_U^{(2)}[\omega]Sv\in\mathcal F_3\).
The force-restoration proof at this
one exponent uses just its first
divergence-source step: the local
energy of \(G_sD^\sharp h\) and
\(h\in L^3\) give \(DG_sD^\sharp h\)
in \(L^3\) by finite-\(3\) Riesz,
while the remaining point source
in \(L^2\) has derivative in \(L^3\).
Its \(P\) part is bounded pointwise.
This proves the finite-\(3\)
version of restoration and
\eqref{acc:second-current-F3}.
No endpoint statement is needed.

Finally a local \(W^{2,3}\) bound
gives \(v\in L^\infty\) and
\(Dv\in L^8\). Then
\(RD^\sharp v\in L^8\) at grade
two, since its \(P\) part remains
pointwise bounded. The scalar
bootstrap for \(\lambda\) gives
both \(\lambda\) and \(D\lambda\)
bounded at their physical grades.
Thus \(h\) and \(a\) in
\eqref{acc:grouped-current-operator}
are bounded. A bounded site force
of grade four has
\(DRG_sa\) pointwise bounded at
grade three, by the scalar point
bootstrap and its smoothed \(P\)
part. The divergence-source term
is controlled in every finite
\(\mathcal F_p\) by force restoration.
Also \(Sv\) is bounded, so the
constraint term is bounded before
the final finite-\(p\) restoration.
This proves the last assertion.
Every derivative has fallen on
a resolvent or the unknown field,
not on the independent current.
\end{proof}

\begin{proposition}[Strong propagators under the explicit convention]
\label{prop:acc-strong-propagators}
For \(G_\bullet\in\{G,G_1\}\), finite
\(4<p<\infty\) and
\(0<\alpha<1-4/p\),
\begin{equation}\label{acc:strong-named-bound}
 \|G_\bullet f\|_{{\rm str},\alpha}
                         \le C_{p,\alpha}\mathcal F_p(f).
\end{equation}
Their full row saddles have the
same strong bound with
\(\sup_c\xi_c|b(c)|\) added, and
their multiplier has bounded
\(\xi^3\) amplitude. The ceiling on
\(\gamma_{\rm t}\) can be fixed independently
of \(\alpha\). The result holds
on real regular pairs and their
small complex \(C^1\) tubes.
\end{proposition}

\begin{proof}
The energy estimate supplies
\(\mathcal E_{\rm cell}(v)\le
C\mathcal F_p(f)\) for either
variant. Relative to \(G_0^{-1}\),
their equations have the first
correction in \(\mathcal F_8\)
and, for the second variant,
the second correction in
\(\mathcal F_3\), with
\(\gamma_{\omega_*,\rm nb}\le C\gamma_{\rm t}\)
by \eqref{acc:nonbare-induced-radius}.
The local scalar vector proof of
\cref{prop:psr-strong-inverse},
stopped after its \(L^3\) equation,
therefore gives normalized
\(W^{2,3}\) control. All carrier
enlargements are retained.
The last part of
\cref{lem:acc-current-bootstrap}
now puts both corrections in the
requested finite \(\mathcal F_p\).
Applying the full \(G_0\) strong
estimate proves
\eqref{acc:strong-named-bound}.
This is a bootstrap, not an
\(\alpha\)-dependent absorption of
a singular-integral constant into
the current radius.

For the full row saddle the
already proved coarse Schur
kernels bound its multiplier
by the normalized force and
row datum amplitudes. Its
\(Q^\sharp\mu\) is a bounded
point force by
\eqref{psr:row-pointwise}.
The same two stages prove its
strong field estimate.
Complex coefficients have the
same absolute bounds in the
real reference's local Coulomb
frames, and the energy inverse
there was obtained by Neumann
series. The proof does not use
Hermitian positivity at a
complex pair.
\end{proof}

The argument can retain a small spatial
exponential weight at every step.
Its local oscillation on complete
carriers is bounded, and the site,
Gram and row sums retain a strict
metric growth reserve. This gives
the corresponding cell-to-cell
strong response estimates. For a
bounded source supported in \(B_z\),
the normalized output in
\eqref{psr:strong-norm} on \(B_y\)
is at most
\(C_\alpha e^{-md(y,z)}\xi_z^3
\|f\|_{\infty,B_z}\).
If desired the factor
\((\xi_z/\xi_y)^3\) is absorbed
by one further fixed scale-ratio
reserve to express the potential
bound as
\(C\xi_y^2e^{-m'd(y,z)}\|f\|_\infty\).

These bounded-force conclusions do
not assert an \(L^\infty\) bound
for \(\Delta_UG\) or
\(\Delta_UG_1\). The divergence
term containing \(W_U[J]\lambda\)
was controlled only in finite
\(L^p\), as required when no
derivative of \(J\) is prescribed.
Nor does a weighted cell estimate
alone construct a background-local
walk expansion.
The canonical continuation
\eqref{acc:canonical-current-definition}
has been made explicit; identifying
another independently prescribed
off-current convention would still
require its own formula comparison.

% END LAYERED EXPANSION inverse/actual_current_completion.tex

% BEGIN LAYERED EXPANSION inverse/local_current_unfolding.tex
\subsection{Finite-carrier equations for the actual current inverse}
\label{ch:lcu-local-current-unfolding}

The projected gauge and current terms conceal global site inverses.
We display an equivalent system whose coefficients are actual local
rows and covariant differences. We prove its inverse on every source
component, including divergence sources. This also identifies the
cutoff estimates needed for a subsequent local parametrix. No
compression of an indefinite system is presumed invertible.

Use the actual proper geometry, physical bilinear transposes and
comparison length \(\ell\) of
\cref{ch:layered-site-restoration,ch:actual-current-completion}.
The field is real and has the curvature ceiling of
\cref{thm:npl-proper-gap}. Put
\begin{equation}\label{lcu:site-objects}
 A_s=D^\sharp D+q^\sharp a_s q,\quad G_s=A_s^{-1},\quad
 F=G_sq^\sharp\Xi^{-2},\quad M=F^\sharp F,\quad
 P=FM^{-1}F^\sharp,\quad R=1-P.
\end{equation}
All site rows, field rows and constraint derivatives are the actual
re-centered ones, including complete outward carriers. As permitted
by the exact support identities in
\cref{ch:actual-current-completion}, replace \(J\) by its restriction
to bonds touching \(\Omega_1\), and discard bare components of the
constraint multiplier. These changes leave both named field
operators unchanged. They avoid putting irrelevant, possibly
unbounded bare currents into auxiliary coefficients.

Write \(W=W_U[J]\) for the actual local current map, and
\(C_\omega=\Delta_U^{(2)}[\omega]\) for the unprojected
constraint Hessian. In the second named variant \(\omega=H^\sharp J\)
uses the first lift. The following arguments also allow an
independent nonbare \(\omega\) within the same small ceiling.
Define
\begin{equation}\label{lcu:field-operators}
 K_F=H_U^{\rm W}+DD^\sharp+Q^\sharp a_fQ,\qquad
 \Lambda=G_sRD^\sharp,\qquad S=1-D\Lambda.
\end{equation}
Here \(H_U^{\rm W}\) is the actual inner Wilson Hessian.
The two field operators are
\begin{equation}\label{lcu:field-variants}
 K^{(0)}=K_F-DPD^\sharp-B_U[J],\qquad
 K^{(1)}=K^{(0)}-S^\sharp C_\omega S.
\end{equation}
They are respectively the first and second canonical named
precisions. Their energy inverses follow from
\cref{ch:actual-current-completion}; the retained current ceilings
are sufficiently small fixed constants.

\subsubsection{Site source spaces with their physical weights}

Let \(X\) be the actual field energy space. Introduce the site norms
\begin{align}
 \|u\|_Y^2&=\|u\|_2^2+\|\ell Du\|_2^2,\notag\\
 \|\lambda\|_Z&=\|\ell^{-2}\lambda\|_2+
       \|\ell^{-1}D\lambda\|_2+\|\nabla_U D\lambda\|_2,\notag\\
 \|h\|_{Y^\flat}&=\sup_{u\ne0}
       \frac{|\langle h,\ell^2u\rangle|}{\|u\|_Y},\qquad
 \|c\|_{T_{-2}}=\|\Xi^{-2}c\|_{\rm top}.
 \label{lcu:site-norms}
\end{align}
The second differences and the bare collar have the convention of
\eqref{lsi:zero-average-estimate}. In particular this is a norm
on the actual finite site space, with the atomic convention at
bare sites. \(Y^\flat\) uses the displayed weighted pairing;
it is not the unweighted dual norm of \(Y\).
If \(h=h_0+D^\sharp h_1\), the product rule gives
\begin{equation}\label{lcu:divergence-source}
 \|h\|_{Y^\flat}\le
 C\{\|\ell^2h_0\|_2+\|\ell h_1\|_2\}.
\end{equation}
Thus a divergence source need not have a pointwise derivative.
The constants below are uniform in the depth and compatible
volume, and in all separated proper shapes under consideration.

\begin{lemma}[A second-difference site lift]
\label{lem:lcu-H2-site-lift}
There is a real linear right inverse \(E\) of \(q\) such that
\begin{equation}\label{lcu:H2-site-lift}
                       qE=1,\qquad \|Ec\|_Z\le C\|c\|_{T_{-2}}.
\end{equation}
Only the small curvature and the actual pillar geometry are
required; no physical-current bound is used.
\end{lemma}
\begin{proof}
Use the compact interior bump construction of
\cref{lem:lsi-weighted-site-lift}, with its exact transported
normalization on each pillar. On the fixed surrounding cube
of length \(C_L\xi_y\), a current-free Coulomb frame has
\[
 \|a\|_\infty\le C_L\kappa\xi_y^{-1},\qquad
                  \|Da\|_2\le C_L\kappa .
\]
The second bound is the finite exponent \(p=2\) case of
\cref{dcz:Coulomb-Lp}; using a smaller cube keeps its collar
inside the nonbare determining region.
For one profile of amplitude \(c_y\), two covariant
differences contain \(D^2\phi\), \(aD\phi\), \(Da\,\phi\)
and \(a^2\phi\). Their \(L^2\) sum is at most \(C_L|c_y|\).
The mass and first-gradient terms with weights
\(\ell^{-2},\ell^{-1}\) have that same bound. The normalization
matrix is constant within that pillar, so differentiating the
profile does not differentiate this matrix.

For the finitely coarse meshes use its center seed; its weighted
second differences cost a fixed constant depending on \(L\).
At a bare top this is the one-site seed, and every difference
cost is paid by its scale \(\eta\). The complete derivative
collars have bounded overlap. Summing the squared estimates
gives \(\sum_y|c_y|^2=\|\Xi^{-2}c\|_{\rm top}^2\).
If the surrounding frame cube does not embed, construct on a
finite periodic cover and average the operator over simultaneous
fine and top deck shifts. The right-inverse identity and these
Hilbert norms descend, with the same sheet factor on both sides.
\end{proof}

\begin{lemma}[Two actual site blocks]
\label{lem:lcu-site-blocks}
The following systems have unique solutions:
\begin{align}
 A_s\lambda+v=f,\qquad
 A_sv-q^\sharp\Xi^{-2}\beta=h,\qquad q\lambda=c,
 \label{lcu:primal-site-block}\\
 A_st+q^\sharp\Xi^{-2}\beta'=h',\qquad
 A_s\kappa-t=r,\qquad q\kappa=d.
 \label{lcu:transpose-site-block}
\end{align}
Their respective estimates are
\par\noindent
\begin{minipage}{\linewidth}
\begin{align}
 \|\lambda\|_Z+\|v\|_Y+\|\beta\|_{\rm top}
 &\le C\{\|f\|_2+\|h\|_{Y^\flat}+\|c\|_{T_{-2}}\},\notag\\
 \|\kappa\|_Z+\|t\|_Y+\|\beta'\|_{\rm top}
 &\le C\{\|r\|_2+\|h'\|_{Y^\flat}+\|d\|_{T_{-2}}\}.
 \label{lcu:site-block-bounds}
\end{align}
\popQED\par
\end{minipage}
\end{lemma}
\begin{proof}
The balanced test in \eqref{lsi:weighted-green-identity},
with the same small Lipschitz constant for \(\ell\), proves
\begin{equation}\label{lcu:balanced-site-inverse}
                         \|G_sh\|_Y\le C\|h\|_{Y^\flat}.
\end{equation}
Indeed its real left side controls
\(\|\ell Du\|_2^2+\|qu\|_{\rm top}^2\), and the actual
block Poincare estimate controls \(\|u\|_2^2\).
This argument applies to every source \(h\), not only to a
site-row column.

For the first block put \(w=G_sh\). Its exact solution is
\begin{equation}\label{lcu:primal-elimination}
 \begin{split}
 \beta&=M^{-1}\{F^\sharp(f-w)-\Xi^{-2}c\},\\
 v&=w+F\beta,\qquad \lambda=G_s(f-v).
 \end{split}
\end{equation}
The identity \(qG_sF=\Xi^2M\) proves \(q\lambda=c\).
The normalized Gram gap and the balanced column estimate for
\(F\) bound \(\beta,v\). For the \(Z\) bound set
\(\lambda_0=\lambda-Ec\). It has zero \(q\) average and
\[
 \Delta_U\lambda_0=f-v-q^\sharp a_sc-\Delta_UEc .
 \]
The right side has \(L^2\) norm at most the first displayed
source bound: \(q^\sharp\) is the actual coisometric
transported extension, \(a_s\asymp\Xi^{-2}\), and
\cref{lem:lcu-H2-site-lift} controls \(\Delta_UEc\).
Apply \eqref{lsi:zero-average-estimate} and then add \(Ec\).

For the second block put \(w'=G_sh'\). The formulas are
\begin{equation}\label{lcu:transpose-elimination}
 \begin{split}
 \beta'&=M^{-1}\{F^\sharp(w'+r)-\Xi^{-2}d\},\\
 t&=w'-F\beta',\qquad \kappa=G_s(t+r).
 \end{split}
\end{equation}
The same argument, with \(\kappa-Ed\), proves its bound.
These formulas prove existence and uniqueness as well.
All transposes and weights are the physical ones; replacing
the top transpose by a counting-array transpose would change
the displayed normalizations.
\end{proof}

\subsubsection{The complete finite-carrier system}

Use \(\epsilon=0\) or \(1\) to select a variant, and set
\begin{equation}\label{lcu:delta-h}
 \delta=\epsilon C_\omega(A-D\lambda),\qquad
 h_J(A,\lambda)=W^\sharp A-\tfrac12
       \{D^\sharp W\lambda+W^\sharp D\lambda\}-D^\sharp\delta .
\end{equation}
The unknowns are \(A,\lambda,v,\beta,t,\kappa,\beta'\).
With arbitrary sources they satisfy
\begin{equation}\label{lcu:full-system}
 \begin{split}
 K_FA-Dv-W\lambda-\delta-Dt&=f_0,\\
 A_s\lambda+v-D^\sharp A&=g,\\
 A_sv-q^\sharp\Xi^{-2}\beta&=h,\qquad q\lambda=c,\\
 A_st+q^\sharp\Xi^{-2}\beta'-h_J(A,\lambda)&=h',\\
 A_s\kappa-t&=r,\qquad q\kappa=d .
 \end{split}
\end{equation}
In particular every coefficient is a covariant difference,
an actual site or field row, a local current multiplication
or an actual two-leg constraint coefficient. No inverse,
projector or induced multiplier is hidden in this matrix.
For the second variant its prescribed coefficient \(\omega\)
must still be constructed from the first lift.

\begin{proposition}[Unfolding with all auxiliary sources]
\label{prop:lcu-unfolding}
For the actual real curvature-small pairs and current ceilings
above, \eqref{lcu:full-system} is invertible and
\begin{equation}\label{lcu:full-inverse-bound}
 \begin{split}
 &\|A\|_X+\|\lambda\|_Z+\|\kappa\|_Z+
      \|v\|_Y+\|t\|_Y+\|\beta\|_{\rm top}+\|\beta'\|_{\rm top}\\
 &\quad\le C\{\|f_0\|_{X^*}+\|g\|_2+\|r\|_2+
          \|h\|_{Y^\flat}+\|h'\|_{Y^\flat}
                       +\|c\|_{T_{-2}}+\|d\|_{T_{-2}}\}.
 \end{split}
\end{equation}
At zero auxiliary sources its field block is precisely
\((K^{(\epsilon)})^{-1}\). Appending \(-Q^\sharp\mu\) to
the field equation and \(QA=b\) gives the full actual
row saddle, with the additional bound
\(\|a_f^{-1/2}\mu\|_{\rm rows}\) on the left and
\(\|a_f^{1/2}b\|_{\rm rows}\) on the right.
\end{proposition}
\begin{proof}
At zero auxiliary sources, the primal block gives
\[
             v=PD^\sharp A,\qquad \lambda=\Lambda A.
\]
The transpose block gives \(t=RG_sh_J(A,\lambda)\).
Since \(\Lambda^\sharp=DRG_s\), the exact symmetric
current formula in \eqref{acc:grouped-current-operator}
shows
\[
 W\lambda+DRG_s\{W^\sharp A-
          (D^\sharp W\lambda+W^\sharp D\lambda)/2\}=B_U[J]A .
\]
The remaining two terms are
\(-\delta+DRG_sD^\sharp\delta=-S^\sharp\delta\).
Thus the field equation is exactly
\eqref{lcu:field-variants}, with the stated sign and order.

Here are the bounds needed for nonzero auxiliary sources.
The actual undifferentiated current and two-leg coefficient
estimates give
\begin{equation}\label{lcu:local-coefficient-bounds}
 \begin{split}
 \|\ell W\lambda\|_2&\le C\gamma_J\|\lambda\|_Z,\\
 \|W^\sharp A\|_{Y^\flat}
       +\|W^\sharp D\lambda\|_{Y^\flat}
       +\|D^\sharp W\lambda\|_{Y^\flat}
                         &\le C\gamma_J(\|A\|_X+\|\lambda\|_Z),\\
 \|\ell C_\omega u\|_2&\le
                           C\gamma_\omega\|\ell^{-1}u\|_2.
 \end{split}
\end{equation}
The last line is a two-leg Schur estimate, not an
\(L^2\)-to-supremum estimate. A fixed-leg row has physical
density \((\eta/\xi)^4\); the transpose cancels it with
\(\eta^{-4}\xi^4\). Both absolute sums are bounded and
all lengths on the carrier are comparable. This proves that
line also for an arbitrary field \(u\).
The middle line uses \eqref{lcu:divergence-source} for
the term \(D^\sharp W\lambda\), and differentiates no current.
Also \(Dv\) and \(Dt\) have \(X^*\) norms at most
\(C\|v\|_2,C\|t\|_2\), by the full physical adjoint.

Solve the primal block for a given \(A\). Subtract its
zero-source solution to obtain remainders \(\lambda_e,v_e\)
bounded by \(\|g\|_2+\|h\|_{Y^\flat}+\|c\|_{T_{-2}}\).
Solve the transpose block next. Its part depending on \(A\)
is exactly the one just eliminated. Its remaining source
is bounded by the auxiliary source norms, using
\eqref{lcu:local-coefficient-bounds}. The field equation
therefore reduces to \(K^{(\epsilon)}A=f_0+f_e\),
where \(\|f_e\|_{X^*}\) has that same bound.
The actual canonical energy inverse bounds \(A\).
The site-block estimates then bound every remaining variable.
For the row saddle use its actual energy inverse instead;
the right inverse of \(Q\) and its dual row norm give
the displayed multiplier estimate. No stationarity or
physical-current identity is used in this argument.
\end{proof}

\begin{corollary}[Bilinear complex continuation]
\label{cor:lcu-complex-inverse}
The full augmented inverse and the appended row saddle continue
holomorphically, with the same source and solution spaces, to a
common smaller weighted \(C^1\) tube about a real reference.
The independent touching current and nonbare multiplier stay
within their indicated complex ceilings.
\end{corollary}
\begin{proof}
Use the fixed real physical norms. In the local coefficients,
the weighted \(C^1\) tube controls \(A_s:Z\to L^2\) and
\(A_s:Y\to Y^\flat\); it controls \(K_F:X\to X^*\)
by its weak plaquette form. Actual transported row differences
have both normalized sums bounded by \(C\varepsilon\).
The same holds for the actual two-leg coefficient by its
analytic row bound. The current terms obey
\eqref{lcu:local-coefficient-bounds} without a current
derivative. Thus the augmented operator is a small
holomorphic perturbation, between exactly the displayed
spaces, of its real reference operator. A Neumann inverse
proves the assertion. Complex transposes are bilinear
throughout; no orthogonal positivity at a complex pair
is used.
\end{proof}

\subsubsection{Slow cutoffs and the bare boundary}

Let \(\chi\) be a real scalar cutoff on the actual sites.
Use its endpoint average on fine bonds and its center or
midpoint value on tops and rows. Suppose
\[
 |\chi|\le1,\qquad |\ell D\chi|\le C/m,\qquad
             |\ell^2\Delta_{\rm graph}\chi|\le C/m .
\]
The last operator is the scalar Laplacian on the actual
site graph. At an action boundary its first-difference
boundary term satisfies this bound because \(\ell\asymp\eta\)
there. Inside a smooth patch a second-difference cutoff
bound \(C/(m^2\ell^2)\) is sufficient.
Every complete row carrier must have \(\chi\) oscillation
at most \(C_L/m\). Cutoffs on patches of radius \(m\ell\)
have this property when \(m\,\operatorname{Lip}(\ell)\)
is sufficiently small.

\begin{proposition}[The actual unfolded commutator]
\label{prop:lcu-commutator}
Let \(\mathbb K^{(\epsilon)}\) be
\eqref{lcu:full-system}, and multiply each component
by the corresponding scalar cutoff. In the source
and solution norms of \eqref{lcu:full-inverse-bound},
\begin{equation}\label{lcu:commutator}
        \|[\mathbb K^{(\epsilon)},\chi]\Psi\|_{\rm source}
                              \le C m^{-1}\|\Psi\|_{\rm solution}.
\end{equation}
The same assertion holds with the actual row saddle appended.
\end{proposition}
\begin{proof}
The product rules and the complete row oscillation give
\[
 \begin{split}
 \|[A_s,\chi]\lambda\|_2&\le C m^{-1}\|\lambda\|_Z,\\
 \|[A_s,\chi]v\|_{Y^\flat}&\le C m^{-1}\|v\|_Y,\\
 \|[q,\chi]\lambda\|_{T_{-2}}&\le C m^{-1}\|\lambda\|_Z,\\
 \|[q^\sharp\Xi^{-2},\chi]\beta\|_{Y^\flat}
                                      &\le C m^{-1}\|\beta\|_{\rm top}.
 \end{split}
\]
For the first line the Laplacian product contains
\(D\chi\,D\lambda\) and
\(\Delta_{\rm graph}\chi\,\lambda\).
For the second line use the weak gradient form and test
against \(\ell^2u\); only \(D\chi\) is required.
For the two row lines use the actual positive absolute
averaging weights and the displayed carrier oscillation.
The site mass commutator satisfies the first two bounds
by these same two row estimates.

The field form \(K_F\) has
\(\|[K_F,\chi]A\|_{X^*}\le C m^{-1}\|A\|_X\).
This follows in its actual weak plaquette and full-divergence
forms, together with the two field-row sums. It uses no
derivative of curvature. The commutators of \(D,D^\sharp\)
cost \(|D\chi|\), with exactly the powers of \(\ell\)
in the displayed spaces. In the current terms insert
the same carrier oscillation in
\eqref{lcu:local-coefficient-bounds}.
For the constraint term,
\(\|\ell[C_\omega,\chi]u\|_2
\le C\gamma_\omega m^{-1}\|\ell^{-1}u\|_2\).
The extra expression from \(A-D\lambda\) is
\(-[D,\chi]\lambda\), whose weighted mass norm is
at most \(C m^{-1}\|\lambda\|_Z\).
Finally treat every \(D^\sharp\delta\) as a divergence
source using \eqref{lcu:divergence-source}.
These bounds cover each row of the system.

At a rough action boundary use the actual full adjoints
and retain every incident inner plaquette. The preceding
weak product identities still hold. The atomic scale
\(\ell\asymp\eta\) pays every one-edge boundary difference.
No continuation of a Coulomb frame across a missing
plaquette is part of this estimate.
\end{proof}

\subsubsection{A positive factorization for the projected gauge}

The following alternative separates the actual nonlocal
Landau gauge into local positive operators and a positive
coarse inverse.

\begin{proposition}[The actual gauge Schur matrix]
\label{prop:lcu-positive-factorization}
At a real curvature-small background put \(G_F=K_F^{-1}\) and
\begin{equation}\label{lcu:gauge-Schur}
              \Theta=M-(DF)^\sharp G_F(DF).
\end{equation}
There is a uniform \(0<\delta<1\) such that
\begin{equation}\label{lcu:gauge-Schur-gap}
                  \delta M\le\Theta\le M.
\end{equation}
For the zero-current Landau inverse one has exactly
\begin{equation}\label{lcu:Woodbury}
 G_0=G_F+G_FDF\,\Theta^{-1}(DF)^\sharp G_F .
\end{equation}
All quantities admit bounded holomorphic continuation to
a common smaller weighted \(C^1\) tube.
\end{proposition}
\begin{proof}
The actual zero-current precision is
\(K_0=K_F-DFM^{-1}(DF)^\sharp\).
It controls \(c\|A\|_X^2\) by
\cref{thm:npl-proper-gap}, while the local form \(K_F\)
is at most \(C\|A\|_X^2\). Thus \(K_0\ge\delta K_F\).
For
\(T=G_F^{1/2}DFM^{-1/2}\) this says
\(1-TT^\sharp\ge\delta\), hence
\(1-T^\sharp T\ge\delta\).
Conjugating by \(M^{1/2}\) proves
\eqref{lcu:gauge-Schur-gap}. The normalized Gram gap
makes it a genuine depth-uniform coarse gap.
The resolvent multiplication identity for
\(K_F-DFM^{-1}(DF)^\sharp\) now gives
\eqref{lcu:Woodbury}. Every transpose here uses the
physical fine or top measure.

In a weighted \(C^1\) tube the operators between the
displayed fixed physical spaces are small holomorphic
perturbations of their real reference values. This is
the site/Gram continuation of
\cref{ch:complex-moving-landau} and the actual
current continuation. In particular the gaps supply
bounded inverse maps there, and the algebraic identity
continues by multiplication. Positivity is used only
at the real reference.
\end{proof}

For a set of sites or bonds, the Dirichlet compression
of \(A_s\) or \(K_F\) is obtained by zero extension
into the actual form. Retain the full coefficient of
every crossing row. Each positive compression inherits
the corresponding global coercivity in its zero-extension
energy norm. Its entries depend only on the complete
carriers of rows meeting that set, together with the
incident inner plaquettes. Thus its geometry-decision
collar is finite. In contrast, deleting the exterior
part of a row and asking for a new row right inverse
can destroy its inf-sup constant; and an indefinite
compression of \eqref{lcu:full-system} is not asserted
to satisfy \eqref{lcu:full-inverse-bound}.

The positive local walk construction in
\cref{lpw:positive-local-walks} supplies the local parametrix
step with these crossing coefficients and a deterministic
finite-decision cover. Coarse inverse carriers for
\(M,\Theta\), the small current corrections and the
first-lift readout must then be assembled in their
actual physical normalizations.
The present equations, estimates and factorization
do not by themselves produce the enlarged independent
cube-parameter family or identify a clean conditional
history law.

% END LAYERED EXPANSION inverse/local_current_unfolding.tex

% BEGIN LAYERED EXPANSION local_positive_inverse_walks.tex
\subsection{Local walks for the positive finite-carrier inverses}
\label{lpw:section}

An exponentially decreasing inverse need not have
coefficients determined by local background data.
We now construct such coefficients for the two
positive finite-carrier operators that enter the
projected precision. The local inverse in the
construction is a genuine Dirichlet compression.
No global inverse with a locally changed
background is declared to be a local coefficient.

The geometry, physical measures and true outer
boundary are those of
\cref{ch:layered-site-restoration}.
All coarse rows retain their complete carriers.
For a real regular connection set
\begin{equation}\label{lpw:positive-operators}
 \mathcal A_s=D_U^\sharp D_U+q^\sharp a_sq,\qquad
 \mathcal A_F=\mathscr D_U+D_UD_U^\sharp+Q^\sharp a_fQ.
\end{equation}
Here \(\mathscr D_U\) is the actual inner
Wilson Hessian. The second operator in
\eqref{lpw:positive-operators} is the full
Feynman operator, before a site-projection
subtraction or an independent-current
correction. The result of this subsection
does not assert locality of those
nonlocal subtractions.

\subsubsection{The positive forms and the cell norms}

For site or vector fields write
\[
 \|u\|_{\mathcal E}^2
       =\|D_Uu\|_2^2+\|\xi^{-1}u\|_2^2,
\]
using the componentwise covariant gradient
in the vector case. Its local version
\(\mathcal E_y\) uses \(B_y\) and one
fine-step collar. For an \(L^2\) source
supported in \(B_z\), put
\[
                   \|f\|_{\mathcal F_z}
                            =\xi_z\|f\|_{2,B_z}.
\]
Thus \(|\langle f,u\rangle|
\le\|f\|_{\mathcal F_z}\|u\|_{\mathcal E_z}\).
Cutoffs and collars have bounded overlap,
so the squared global energy is equivalent
to the sum of the squared cell energies.
No pointwise Green-kernel norm is implicit
in this notation.

There are uniform constants \(c,C>0\) with
\begin{equation}\label{lpw:positive-gap}
 c\|u\|_{\mathcal E}^2
          \le\langle u,\mathcal A_\bullet u\rangle
          \le C\|u\|_{\mathcal E}^2,
                             \qquad \bullet=s,F.
\end{equation}
For the site operator this is the actual
block inequality
\eqref{lsi:block-poincare} and the
transported-average upper bound.
For the vector operator,
\[
 \mathcal A_F=K_{U,0}+D_UPD_U^\sharp,\qquad P=1-R.
\]
The real projection satisfies \(0\le P\le1\).
The lower bound follows from
\eqref{npl:proper-gap}, and the additional
term is at most \(\|D_U^\sharp u\|_2^2\),
bounded by the component gradient and
bare mass energy. The latter pays any
one-edge term at a rough true boundary,
where \(\xi=\eta\).
This proves the upper bound too.
In particular this argument does not
require the connection to be stationary.

Let \(E_D\) inject fields supported on
any chosen fine set \(D\), extending
them by zero in the original field
space. Define the compressed operator
\begin{equation}\label{lpw:compression}
                  \mathcal A_{\bullet,D}
                              =E_D^\sharp\mathcal A_\bullet E_D.
\end{equation}
Its norm is the energy of that zero
extension, including the artificial
boundary gradients. Inequality
\eqref{lpw:positive-gap} holds on this
space with exactly the same constants.
For an averaging row crossing the
boundary of \(D\), its contribution
is its full row evaluated on the
zero extension. It is not replaced
by a renormalized partial average.
There is consequently no new
partial-row inf-sup assertion.

The matrix \eqref{lpw:compression}
depends only on the fine coefficients
on \(D\), their finite difference
collar and the complete rows meeting
that collar. It also retains which
original inner plaquettes are present.
These observations apply at the
actual outer boundary as well as
at an artificial interior boundary.

\subsubsection{A locally determined slow partition}

We give an explicit weighted cover so
its choice does not depend on a
globally ordered greedy selection.
Fix a large integer \(m\), and then
choose the comparison-length slope
\(\delta>0\) so \(m\delta\) is
sufficiently small. Increase
the separation so
\(\rho\delta\ge C_L\).
The comparison length \(\ell\) is
\eqref{lsi:comparison-length}.
Put
\[
                           r_x=m\ell(x)
\]
at every ambient fine site \(x\),
including the geometric exterior.
Initially suppose the periods are
large enough that the balls below
embed. The periodic case is handled
after the construction.

Choose a fixed nonnegative smooth
bump \(\varphi\), equal to one on
the ball of radius \(1/4\) and
supported in the ball of radius
\(1/2\). Define
\begin{equation}\label{lpw:weighted-partition}
 \begin{split}
 \omega_x&=\eta^4r_x^{-4},\\
 N(y)&=\sum_x\omega_x
                  \varphi((y-x)/r_x)^2,\\
 \chi_x(y)&=N(y)^{-1/2}\varphi((y-x)/r_x).
 \end{split}
\end{equation}
The sums are ambient fine-site sums.
The radius and weight at a centre
are fixed while \(y\) varies.

\begin{lemma}[Weighted slow partition]
\label{lpw:slow-partition}
Uniformly in mesh and depth,
\begin{equation}\label{lpw:partition-bounds}
 \sum_x\omega_x\chi_x(y)^2=1,\qquad
 c\le N(y)\le C,\qquad
 |D^a\chi_x(y)|\le C_ar_x^{-a}\quad(a=0,1,2).
\end{equation}
Only centres with \(|x-y|\le Cr_y\)
contribute at \(y\), and
\begin{equation}\label{lpw:weighted-overlap}
                 \sum_{x:\,y\in B(x,2r_x)}\omega_x\le C.
\end{equation}
The same holds with any fixed enlargement
of those balls, after reducing \(m\delta\).
All assigned scales in an enlarged ball
are comparable within a fixed \(L\) factor.
\end{lemma}

\begin{proof}
If \(|x-y|\le ar_x\) for a fixed
enlargement factor \(a\), then
\[
                 |r_x-r_y|\le m\delta|x-y|\le am\delta r_x.
\]
Thus \(r_x\) and \(r_y\) are comparable,
and all such centres lie in a fixed
multiple of the radius-\(r_y\) ball.
Conversely, in a sufficiently small
fixed multiple of that ball,
the bump is bounded below.
Since \(\ell\ge\eta\), \(r_y\ge m\eta\).
Fine lattice counting in these balls
therefore gives
\[
 \sum_{|x-y|\le cr_y}\eta^4r_x^{-4}\ge c,\qquad
 \sum_{|x-y|\le Cr_y}\eta^4r_x^{-4}\le C.
\]
This proves the bounds for \(N\) and
the weighted overlap. In particular
the large number of fine centres
has been canceled by their physical
weights, rather than estimated by
an unweighted count.

Differentiate \(N\) in \(y\).
Only the bump is differentiated;
each derivative costs \(r_x^{-1}\).
The preceding comparisons give
\(|D^aN(y)|\le C_ar_y^{-a}\)
for \(a\le2\). The product rule
and its positive lower bound prove
the cutoff estimates.
The partition identity is its
definition. Finally
\(\ell\le\xi\le L\ell\);
the same Lipschitz comparison makes
all assigned scales in these balls
comparable. The separation excludes
nonadjacent layers on the fixed
enlargements used here.
\end{proof}

The partition is local in the geometry.
Here is the point needed for exact
carrier dependence. If \(x\) is assigned
level \(j\), a candidate in the infimum
defining \(\ell(x)\) at level at most
\(j-2\) has distance at least
\(\rho\xi_{j-1}\). For
\(\rho\delta>2L\), its distance term
alone exceeds \(2\xi_j\), so it
cannot attain an infimum at most
\(\xi_j\). A candidate at level
at least \(j\) cannot improve the
zero-distance candidate \(x\).
Only the adjacent lower level need
be examined, at distance at most
\(\xi_j/\delta\).
Consequently all radii and cutoff
normalizations used near \(x\) are
decided by the layer labels in
\begin{equation}\label{lpw:geometry-decision-collar}
                         B(x,C_L\ell(x)/\delta).
\end{equation}
Increase \(\rho\delta\) so this
collar meets only the fixed adjacent
levels. This may be a larger collar
than the support of \(\chi_x\);
it is included in the carrier.
It cannot be omitted on the ground
that \(\ell\) is merely comparable
to the assigned scale.

\subsubsection{Local inverse columns and the commutator}

Let \(D_x\) contain the fine variables
in \(B(x,2r_x)\), intersected with
the original field space. Its
finite carrier collar is fixed
large enough for the complete rows.
Write
\[
 I_x=E_{D_x}\mathcal A_{\bullet,D_x}^{-1}E_{D_x}^\sharp.
\]
The two outside injections specify
zero extension; \(I_x\) is not an
inverse on a newly defined determining
region. In particular the original
outward rows and bare boundary
are still used in its compression.

\begin{lemma}[Uniform local columns]
\label{lpw:local-columns}
For some \(s_0>0\),
\begin{equation}\label{lpw:local-column-bound}
       \|1_y I_x1_z\|_{\mathcal F_z\to\mathcal E_y}
                  \le Ce^{-s_0d(y,z)}.
\end{equation}
The constants are independent of
the patch radius, mesh and depth.
The same conclusion holds for the
small complex \(C^1\) tubes of the
actual regular real coefficients.
\end{lemma}

\begin{proof}
The unweighted energy estimate is
\eqref{lpw:positive-gap} restricted
to the zero-extension subspace.
For a spatial exponential weight
whose logarithm changes by at most
\(s\eta/\xi\) on a fine step,
the covariant product rule changes
the principal form by at most
\(Cs\|u\|_{\mathcal E}^2\).
The oscillation of that weight on
one complete \(q\) or \(Q\) carrier
is \(O(s)\), so the same estimate
holds for their positive row forms.
The Wilson plaquette form has this
bound by its exact finite insertion
formula. Taking \(s\) small compared
with \(c\) preserves coercivity.
Multiplication by the weight
preserves zero extension, so this
argument is unchanged by the
artificial Dirichlet boundary.
Choose the logarithm to be the
truncated distance from \(B_z\)
in the actual comparison metric.
The weighted energy bound gives
\eqref{lpw:local-column-bound}.

For a complex tube, the local
operator coefficients are analytic.
The relative form estimates used
in the complex inverse construction
apply to the two local parts of
\eqref{lpw:positive-operators}.
Restricting their inputs by zero
extension cannot enlarge those
relative bounds. A fixed sufficiently
small tube is therefore a Neumann
perturbation of every weighted real
compression, with the same constants.
No Hermitian positivity at the
complex coefficient is asserted.
\end{proof}

\begin{lemma}[Slow commutator]
\label{lpw:slow-commutator}
The local columns also satisfy
\begin{equation}\label{lpw:commutator-column-bound}
 \|1_y[\mathcal A_\bullet,\chi_x]I_x\chi_x1_z
                  \|_{\mathcal F_z\to\mathcal F_y}
                       \le \frac C m e^{-s_1d(y,z)}
\end{equation}
for a fixed \(s_1>0\), whenever the
indicated supports meet. Its output
and input lie in fixed enlargements
of the patch. Its coefficients
depend only on the data there and
in \eqref{lpw:geometry-decision-collar}.
\end{lemma}

\begin{proof}
In a covariant second difference,
the exact scalar product rule
expresses the commutator as
first differences of \(\chi_x\)
times first covariant differences
of the unknown, and second
differences of \(\chi_x\) times
the unknown. There is no
derivative of a connection
coefficient in this identity.
The same rule for the
curl-divergence principal form
of the vector operator has
these two types of terms.
The remaining plaquette kernel
has bounded finite-shift
coefficients by
\cref{lem:psr-Wilson-stencil};
its commutator is smaller.
After multiplying the output
\(L^2\) norm by \(\xi_y\),
these bounds are
\[
 \frac C m\{\|D_Uu\|_{2,\mathcal D_y}
                         +\xi_y^{-1}\|u\|_{2,\mathcal D_y}\}.
\]
Here \(\mathcal D_y\) is a fixed
cell enlargement, with comparable
assigned scales.

For the averaging terms introduce
the scalar cutoff value at the
row's top. The row commutators are
\(q\chi_x-\chi_{x,\rm top}q\)
and \(Q\chi_x-\chi_{x,\rm row}Q\).
Every difference of cutoff values
on their complete carrier is
at most \(C/m\). Their normalized
row and column bounds therefore
give that same factor in the
physical \(L^2\) mass norm.
The physical transpose is essential
here; no fine row count remains.
Use \cref{lpw:local-columns} and
sum the finitely many neighbouring
cells. This proves the displayed
exponential bound with a slightly
smaller rate. The local matrices,
cutoffs and complete row carriers
give its stated exact dependence.

At a true boundary, the full
interior curl-divergence identity
need not hold. Every such stencil
is in the bare layer. There
\(\xi=\eta\); each present matrix
entry is \(O(\eta^{-2})\), its
cutoff difference is \(O(1/m)\),
and it reads only finitely many
neighbours. Thus its force norm
is bounded by \(C/m\) times the
bare mass energy directly.
This also covers a missing
neighbour in the scalar graph
Laplacian. No second-difference
cancellation across an absent
stencil is required.
\end{proof}

This proof works next to a rough
outer boundary. The identity is
the commutator of the original
finite matrices and their diagonal
cutoffs. Missing exterior plaquettes
are not filled in, and no local
Coulomb gauge across that boundary
has been assumed.

\subsubsection{An exact local parametrix}

Define
\begin{equation}\label{lpw:parametrix}
                 \mathcal M=\sum_x\omega_x\chi_x I_x\chi_x.
\end{equation}
All sums are finite on a finite
torus. The preceding bounds make
them uniformly bounded between
cell source and energy norms.
Because \(\chi_x\) has support
strictly inside \(D_x\),
compression and extension give
the exact identity
\begin{equation}\label{lpw:parametrix-identity}
 \mathcal A_\bullet\mathcal M=1+\mathcal T,\qquad
 \mathcal T=\sum_x\omega_x
                  [\mathcal A_\bullet,\chi_x]I_x\chi_x.
\end{equation}
To verify its main term, multiply
by \(\chi_x\) on the output of
the compressed equation:
\[
 \chi_x\mathcal A_\bullet I_x\chi_x
                               =\chi_x^2.
\]
The weighted sum of these squares
is one by \eqref{lpw:partition-bounds}.
Thus there is no residual
background-replacement term in
\eqref{lpw:parametrix-identity}.
The operator is exactly the
same local operator inside each
compression.

For a block operator \(\mathcal B\)
put
\[
 \|\mathcal B\|_{\rm Sch,s}
   =\max\left\{\sup_y\sum_z e^{sd(y,z)}\|\mathcal B_{yz}\|,
               \sup_z\sum_y e^{sd(y,z)}\|\mathcal B_{yz}\|\right\},
\]
using \(\mathcal F\to\mathcal F\)
for \(\mathcal T\) and
\(\mathcal F\to\mathcal E\) for
\(\mathcal M\).
The metric growth reserve is
chosen below the column decay.
Equations \eqref{lpw:weighted-overlap}
and \eqref{lpw:commutator-column-bound}
give
\begin{equation}\label{lpw:parametrix-Schur}
                 \|\mathcal T\|_{\rm Sch,s}\le C/m,\qquad
                 \|\mathcal M\|_{\rm Sch,s}\le C.
\end{equation}
Indeed for fixed output cell,
all contributing centres have
radius comparable to \(m\xi_y\);
their physical weights have
bounded sum. Sum the remaining
exponential input-cell kernel.
The column calculation is the
same with the two cells reversed.
The local column estimate,
rather than the total energy
norm of a whole patch, prevents
a spurious factor \(m^2\).

\subsubsection{Carriers and absolute walk sums}

Choose a tagged mixed-scale carrier
for each centre as follows. Include
its inverse patch, the complete
row collars and the geometric
decision collar
\eqref{lpw:geometry-decision-collar}.
Cover that set by the ambient
\(M_{\rm c}\xi_j\) cubes at its adjacent
assigned levels. Choose the fixed
carrier cube integer \(M_{\rm c}\) large
enough compared with \(C_L/\delta\),
and then enlarge \(\rho\) once
more compared with \(LM_{\rm c}\).
Only a bounded number \(B_L\)
of tagged cubes is needed.
Connect intersecting or neighbouring
cubes at the same or adjacent
level. Their connected carrier
has length at most \(B_L\).
This records geometry dependence
as well as matrix endpoints.
Include also the canonical tagged
\(M_{\rm c}\)-cube of each input and output
assigned cell. Any short joining
paths cost a fixed number of tags,
absorbed into \(B_L\).
For multiplication use the shifted
length
\(\sigma_c(X)=b_L+d_c(X)\),
where \(b_L\ge1\) exceeds the
fixed cost of joining two carriers
at a common intermediate cell.
Then
\[
             \sigma_c(X\vee Y)\le\sigma_c(X)+\sigma_c(Y).
\]
Here \(X\vee Y\) includes that
short join if the original tagged
vertex sets were disjoint despite
meeting physically. This convention
is equivalent, up to one fixed
factor in the norm, to the usual
connected tree length.
Enlarge \(B_L\) once to bound
\(\sigma_c\) for every local carrier.

Fix any desired finite carrier
rate \(\kappa_c>0\). The parameters
are chosen in this order: first
\(\kappa_c\) and the operator gap
constants, then \(m\) sufficiently
large, then \(\delta\), \(M_{\rm c}\),
and \(\rho\) as above. The bound
on the number of carrier cubes
is independent of their physical
radius after this choice of \(M_{\rm c}\).
Thus \(m\) can be chosen so
\begin{equation}\label{lpw:carrier-contraction}
                         (C/m)e^{\kappa_cB_L}<1/2.
\end{equation}
The global curvature and complex
tube ceilings are reduced only
after these finitely many fixed
parameters. None depends on depth
or volume.
The carrier cube size \(M_{\rm c}\)
is distinct from a prescribed history
cube size. A statement giving only
one history cube of separation must
verify that its physical width meets
the present \(\rho\) bound. That
history-admission comparison is not
part of the local inverse theorem.

Group local terms with the same
carrier. For an operator-valued
family \(\mathbf B=(B_X)\), use
the weighted Schur norm with the
additional sum over \(X\) and
factor \(e^{\kappa_c\sigma_c(X)}\).
The two estimates
\eqref{lpw:parametrix-Schur}
then become
\[
             \|\mathbf T\|_{\rm Sch,s,\kappa_c}<1/2,\qquad
             \|\mathbf M\|_{\rm Sch,s,\kappa_c}\le C_{\kappa_c}.
\]
These are sums of norms of actual
coefficients. No cancellation
between distinct carriers was
used to make the first estimate
small.

\begin{proposition}[Positive inverse walks on the determining geometry]
\label{lpw:positive-local-walks}
For \(\bullet=s,F\), the actual inverse
has the exact absolutely convergent
walk representation
\begin{equation}\label{lpw:walk-series}
                \mathcal A_\bullet^{-1}
                   =\mathcal M\sum_{n\ge0}(-\mathcal T)^n.
\end{equation}
Its carrier coefficients are
holomorphic functions of the
local background data and the
declared local determining geometry.
They have a uniform weighted
cell-to-cell energy norm with
the carrier rate just fixed.
This holds on the real regular
domain and its specified small
complex \(C^1\) tubes.
\end{proposition}

\begin{proof}
The Schur family norm is
submultiplicative. In a product,
the intermediate cell joins the
two contributing carriers. Use
their recorded short join and
the subadditive shifted length.
Its carrier weight is at most
the product of the two weights,
and the output-input distance
obeys the triangle inequality.
Thus the absolute sum of
walks of length \(n\) is at most
\(C_{\kappa_c}2^{-n}\).
This proves convergence in the
stated family norm.
Multiply \eqref{lpw:walk-series}
by \(\mathcal A_\bullet\) and
use \eqref{lpw:parametrix-identity}.
The product is the identity.
The positive operator is
invertible by
\eqref{lpw:positive-gap},
so the series is its inverse.

Each local factor is a finite
matrix inverse of
\eqref{lpw:compression}, multiplied
by local cutoff and stencil
coefficients. Its matrix entries,
original row list and original
plaquette list are determined
inside its recorded carrier.
Agreement of two determining
geometries and backgrounds there
therefore gives exactly the same
coefficient. This is stronger
than a shielding estimate.
The local inverse is analytic
on its uniformly admitted
complex tube. Products are
analytic on the union carrier,
and the normally convergent
walk sum is analytic there.
\end{proof}

\subsubsection{Small periods and the scope of the locality statement}

If the physical balls wrap around
a small compatible torus, first
lift the local coefficients and
geometry to a finite periodic
cover whose periods exceed
the fixed patch and decision
diameters. A bounded number of
copies suffices, since the original
periods contain the largest
assigned blocks. The lifted
positive form has the same
uniform gap.
Its inverse preserves the
deck-invariant subspace by
uniqueness. On that subspace
it is exactly the original
inverse.

There is no multiplicity loss
in the Schur contraction.
Before projecting a filed family,
average it over the finite deck
group, translating its carrier tags
as well. This step is needed if
the auxiliary carrier grid itself
is not deck invariant. The total
inverse is unchanged by uniqueness,
and the absolute family norm does
not increase. The averaged family
is now deck covariant.
For a base source cell, its lift
is the sum of its deck images.
The projected kernel is the
sum of the corresponding cover
kernel columns. A base row sum
is consequently bounded by the
cover row sum; the column statement
follows by deck covariance.
Projection of a connected
carrier cannot increase its
length. The same absolute
walk estimates descend with
the same contraction constant.
This argument concerns classical
linear inverses; it does not
identify a quantum partition
function with the partition
function on a larger cover.

The result constructs local walks for $\mathcal A_s^{-1}$ and
the full Feynman inverse. The exact compositions with the site
projection, coarse inverses and independent current are supplied
by \cref{prop:ncw-named-families}. Their strong output and
cube-parameter versions are \cref{dnr:linear-families,dnr:parameter-families}.
The carriers retain their tagged mixed-scale determining geometry.
Equality with a separately prescribed signed or absolute-contour
precision requires that precision's gauge comparison.

% END LAYERED EXPANSION local_positive_inverse_walks.tex

% BEGIN LAYERED EXPANSION background/coarse_carrier_inverse.tex
\subsection{Local carrier expansions for the positive coarse inverses}
\label{ch:cci-coarse-carrier-inverse}

We invert the actual positive coarse matrices while retaining
the full background carriers of their coefficients. The short
part is cut by its recorded carrier, and its matrix powers are
combined before taking absolute values. This permits the real
Hilbert-space gap to control the absolute walk expansion.
The local positive inverses used as inputs are those of
\cref{lpw:positive-local-walks}; no compression of an indefinite
projected operator is inverted here.

The initial carrier tile size is $M_0$. The final size $M_c$
will be chosen larger after a finite short-power construction.
Both are aligned powers of $L$. Neither is silently identified
with a prescribed history size $M_h$. The result requires the
separation specified below. Admission of a separately fixed
history geometry remains a comparison to be checked.

\subsubsection{The actual matrices and their physical normalization}

Retain the physical site and field pairings. Put
\begin{equation}
 \begin{split}
 G_s&=\mathcal A_s^{-1},\qquad
 F=G_sq^\sharp\Xi^{-2},\qquad M=F^\sharp F,\\
 G_F&=\mathcal A_F^{-1},\qquad
 \Theta=M-(D_UF)^\sharp G_F(D_UF).
 \end{split}
 \label{cci:actual-coarse-matrices}
\end{equation}
Here $\mathcal A_s,\mathcal A_F$ are the positive local
operators in \eqref{lpw:positive-operators}. The coarse
Hilbert norm is $\sum_y\xi_y^4|b(y)|^2$. Thus
\begin{equation}
 (Jb)(y)=\xi_y^2b(y),\qquad
             \widehat M=JMJ^{-1},\qquad
             \widehat\Theta=J\Theta J^{-1}
 \label{cci:physical-unitary}
\end{equation}
is the unitary passage to counting $\ell^2$. All matrix
families and ordinary transposes below refer to this
normalization. Physical transposes have not been replaced
by unweighted transposes in \eqref{cci:actual-coarse-matrices}.

Both normalized real matrices have fixed positive gaps.
The Gram bound is \eqref{lsi:gram-gap}. To check the
second one define
\[
 K_L=\mathscr D_U+D_URD_U^\sharp+Q^\sharp a_fQ.
\]
The real lower bound in \eqref{npl:proper-gap} and the upper energy bound for
$\mathcal A_F$ give $K_L\ge\delta\mathcal A_F$ for one
$\delta>0$. Since
\[
 K_L=\mathcal A_F-(D_UF)M^{-1}(D_UF)^\sharp,
\]
the operator $C=G_F^{1/2}(D_UF)M^{-1/2}$ obeys
$I-CC^\sharp\ge\delta I$. Hence $\|C\|^2\le1-\delta$
and $I-C^\sharp C\ge\delta I$. Conjugating gives
\begin{equation}
                \delta M\le\Theta\le M.
 \label{cci:Theta-gap}
\end{equation}
Consequently, for $A\in\{\widehat M,\widehat\Theta\}$,
\begin{equation}
                         aI\le A\le bI
 \label{cci:common-real-gap}
\end{equation}
with $a,b>0$ independent of depth, compatible volume and
the sufficiently separated geometry. This is a real
positivity argument; complex continuation will use a
perturbation of these real matrices.

\subsubsection{Actual coefficient families before coarse inversion}

A tag is a complete ambient $M_0\xi_j$ cube with its level,
row weights and determining-geometry data retained. A carrier $X$
records a finite connected collection of tags and its
fixed short joins. Its endpoints include the canonical
assigned-level tags of the matrix input and output cells.
Use the guarded length
\[
                       \sigma_0(X)=b_L+d_{c,0}(X)
\]
from the positive-walk construction, with $b_L$ paying
the fixed cost of joining two carriers at a common cell.
In particular
$\sigma_0(X\vee Y)\le\sigma_0(X)+\sigma_0(Y)$.
The length is that of the recorded full carrier, not a
shorter connection through unrecorded geometry.
Each tag covers assigned cells of its declared level,
or is one of the fixed short joining tags. This is the
covering used for the positive-walk coefficients.

For a counting-normalized matrix family $\mathbf A$
write $A_{yz}=\sum_X A_{yz}[X]$ and set
\begin{equation}
 \|\mathbf A\|_{s,\kappa,0}
 =\max\left\{
 \sup_y\sum_{z,X}e^{s d(y,z)+\kappa\sigma_0(X)}\|A_{yz}[X]\|,
 \sup_z\sum_{y,X}e^{s d(y,z)+\kappa\sigma_0(X)}\|A_{yz}[X]\|
                                                        \right\}.
 \label{cci:carrier-family-norm}
\end{equation}
Joining full carriers in a matrix product and using the
triangle inequality for $d$ makes this an algebra norm.
The identity has its canonical one-cell carrier and a
bounded norm. Transpose preserves the norm. Replacing a
family by half its sum with its transpose makes each
coefficient symmetric under endpoint reversal without
changing a symmetric total matrix or increasing the norm.

For every fixed finite $\kappa_0$, the actual matrices in
\eqref{cci:physical-unitary} have families with
\begin{equation}
       \|\mathbf A\|_{s_0,\kappa_0,0}\le K_{\kappa_0},
                      \qquad s_0>0.
 \label{cci:actual-input-family}
\end{equation}
Here $M_0$ and the primitive parameters are first chosen
as in the positive-walk theorem. The constants are uniform
above their separation threshold. We verify the physical
scaling in this assertion, because an unweighted fine
entry sum would not prove it.

A unit counting top datum is $b_z=\xi_z^{-2}e_z$ in
physical coordinates. The source
$q^\sharp\Xi^{-2}b_z$ has cell $L^2$ norm at most
$C\xi_z^{-2}$, and thus source norm $C\xi_z^{-1}$ in
the notation of \cref{lpw:local-columns}. The actual
site walk gives
\begin{equation}
 \|1_yFb_z\|_2\le C(\xi_y/\xi_z)e^{-s d(y,z)},\qquad
                    \|1_yD_UFb_z\|_2\le C\xi_z^{-1}e^{-s d(y,z)}.
 \label{cci:normalized-F-columns}
\end{equation}
These inequalities also hold as sums over the recorded
walk carriers, at the prescribed rate. Applying a fine
derivative uses the energy output of that walk and its
already included one-step collar.
Similarly the field walk gives
\begin{equation}
                 \|1_wG_F1_v\|_{2\to2}
                              \le C\xi_w\xi_v e^{-s d(w,v)}.
 \label{cci:normalized-GF-columns}
\end{equation}
The local $q$ injection and all complete row carriers
are retained in these coefficients.

For $\widehat M_{yz}$, Cauchy--Schwarz inside each physical
fine cell leaves the factor $\xi_w^2/(\xi_y\xi_z)$.
For the subtracted term in $\widehat\Theta_{yz}$ it leaves
$\xi_w\xi_v/(\xi_y\xi_z)$ on the chain $y,w,v,z$.
The level barrier \eqref{px:level-barrier} bounds each
scale ratio by a fixed $L$ factor times
$\exp\{(\log L)d/\rho\}$. Reserve that small rate before
summing the exponentially decreasing cell chains with
\eqref{px:metric-ball-count}. This proves
\eqref{cci:actual-input-family}, including its absolute
row and column sums. No fine lattice count or power of
$\eta^{-1}$ is inserted. Carrier unions are the actual
unions of the site, field and injection carriers; the
shared fine cell and the guarded join preserve locality.

The positive-walk coefficients are uniformly holomorphic
on their common small complex $C^1$ tubes. Their normally
convergent family sums and the preceding bounded
compositions give the same assertion for
\eqref{cci:actual-input-family}. On a smaller tube,
Cauchy's estimate gives a Lipschitz bound in that family
norm. This supplies the complex perturbation used below;
it supplies no independent enlarged physical cube domains.

\subsubsection{Short connected carriers and the actual local counting}

Fix finite integers $m,R$ for the moment. On the final
separation threshold, every old connected carrier of
length at most $R$ has metric diameter at most
$C_LM_0(R+1)$. A chain of $m$ of them is contained in a
ball of radius
\begin{equation}
                         B=C_LM_0m(R+1).
 \label{cci:short-ball-radius}
\end{equation}
Choose the final separation much larger than $LB$.
A metric ball of radius $4B+C_L$ then meets at most
two adjacent levels. Otherwise a path between its
extreme levels would cross a complete intervening
layer at cost at least $\rho$, contradicting the
bound on the ball diameter. Its reached scales are
comparable within $L$, and its physical extent is
$C_L(B+1)\xi_y$. Counting their two grids proves
\begin{equation}
                 \#\{z:d(y,z)\le 2B+C_L\}
                                  \le C_L(1+B)^4.
 \label{cci:local-polynomial-count}
\end{equation}
This is a local consequence of separation. The global
growth bound remains \eqref{px:metric-ball-count}.

The same observation proves the diameter assertion for
the short carriers: neighboring complete old tiles
have metric span $C_LM_0$, and fewer than $m(R+1)$
such steps cannot cross two separated interfaces.
All tag eligibility, row lists and geometry decisions
are read in the resulting ball, since those data were
included in each old carrier. Cutoff by endpoint
distance alone would not have this property.

Now choose a final aligned carrier tile size $M_c$,
an $L$-power multiple of $M_0$, with
\begin{equation}
                   M_c\ge C_L(B+1),\qquad \rho\ge C_LLM_c.
 \label{cci:final-carrier-tiles}
\end{equation}
The union of two radius-$B$ balls with centers at
distance at most $B$ has a connected covering with
at most $B_L$ final tags, including its endpoint tags
and short joins. The constant $B_L$ depends on $L$
and the fixed collar conventions, not on $m,R$.
Write $\sigma_c$ for the same guarded length on these
final tags. Reblocking an old carrier into its aligned
parents preserves its connections and satisfies
$\sigma_c(\pi X)\le\sigma_0(X)$.
Thus an old strong carrier weight pays the final one.

\subsubsection{Combine a short power before taking absolute values}

Fix the desired finite final carrier rate $\kappa>0$ and
use \eqref{cci:actual-input-family} at a stronger rate
$\kappa_0>\kappa+1$. Choose $\beta\ge2b$ and put
\begin{equation}
 H=I-A/\beta,\qquad S=H_{\le R},\qquad E=H_{>R},
 \qquad H_{\le R}=
                  \sum_{\sigma_0(X)\le R}H[X].
 \label{cci:full-carrier-cutoff}
\end{equation}
The families are first transpose symmetrized. This
cutoff consequently preserves the real symmetry of $S$.
Let $K_H\ge2$ bound the old family norms of $H,S$,
including the identity carrier. The tail obeys
\begin{equation}
 \|\mathbf E\|_{s,\kappa,c}
             \le K_H e^{-(\kappa_0-\kappa)R}
                       \quad(0\le s\le s_0).
 \label{cci:actual-carrier-tail}
\end{equation}
Here and below the old coefficients are reblocked into
final tags when the subscript $c$ is used.
This estimate cuts by the full old carrier; a long
dependency with nearby matrix endpoints remains in
the tail.

For $R$ above a fixed threshold the unweighted tail
operator norm is at most $a/(2\beta)$. Since
$I/2\le H\le(1-a/\beta)I$ on the real domain,
\begin{equation}
                0\le S\le qI,\qquad
                              q=1-a/(2\beta)<1.
 \label{cci:real-short-contraction}
\end{equation}
The real spectral estimate is used only at this point.
The short matrix has propagation at most $B/m$.
Therefore its $m$th power has range $B$, and
\eqref{cci:local-polynomial-count} and the Hilbert row
estimate give
\begin{equation}
                  \|S^m\|_{\mathrm{Sch}}
                               \le C_L(1+B)^2q^m.
 \label{cci:combined-real-Schur}
\end{equation}
Indeed the squared Hilbert--Schmidt norms of the blocks
in one row sum to at most $r_0\|S^m\|_{2\to2}^2$, where
$r_0$ bounds the colour block dimension. Cauchy--Schwarz
costs the square root of the number of reached cells.
The same holds for columns. This estimate concerns the
already combined matrix entries of $S^m$, not the
absolute values of its individual paths.

Assign that combined entry $(S^m)_{yz}$ the symmetric
full carrier covering the two radius-$B$ balls of
$y,z$. It includes every short path and its full
coefficient dependencies. This is a local function:
agreement of backgrounds and recorded geometry on
that carrier leaves all contributing short factors
and their eligibility unchanged. Its final length is
at most $B_L$. Denote the resulting family by
$\mathbf N_m$. At any $0<s\le s_0$ it has bound
\begin{equation}
        \|\mathbf N_m\|_{s,\kappa,c}
              \le C_Le^{sB+\kappa B_L}(1+B)^2q^m.
 \label{cci:combined-local-family}
\end{equation}
In particular the cancellation used in
\eqref{cci:combined-real-Schur} takes place inside a
specified local coefficient before filing it.

Keep the remaining part as an actual word family:
\begin{equation}
 \begin{split}
 E_j&=H^j-S^j
       =\sum_{i=0}^{j-1}S^i E H^{j-1-i},\\
       \|\mathbf E_j\|_{s,\kappa,c}
                    &\le jK_H^j e^{-(\kappa_0-\kappa)R}.
 \end{split}
 \label{cci:exact-word-tail}
\end{equation}
The family products file into their full joined
carriers. Transpose symmetrization preserves the
bound and the total operator $E_j$, since both
$H^j$ and $S^j$ are symmetric. Define
\begin{equation}
                     \mathbf T=\mathbf N_m+\mathbf E_m.
 \label{cci:local-contracted-step}
\end{equation}
Its total operator is exactly $H^m$.

Here is a noncircular choice of all parameters.
First fix the real gaps and $\kappa_0>\kappa+1$,
using the earlier fixed energy estimates. The comparison
length slope and the carrier grids are auxiliary; changing
them for the walks does not change the actual coarse matrices
or their already proved gaps.
Next construct the positive-walk inputs and their
initial $M_0,K_H,s_0$. Choose a positive integer $D$ so that
$(\kappa_0-\kappa)D>2\log K_H+2$ and $R=Dm$
already meets the spectral-tail threshold for $m\ge1$.
Then choose $m$ large. The quantity
\begin{equation}
 \Gamma_m=2C_Le^{\kappa B_L}(1+C_LM_0m(Dm+1))^2q^m
                 +mK_H^m e^{-(\kappa_0-\kappa)Dm}
                                  \longrightarrow0.
 \label{cci:finite-parameter-choice}
\end{equation}
Choose $m$ with $\Gamma_m<1/4$, put $R=Dm$, and
choose $M_c$ as in \eqref{cci:final-carrier-tiles}.
Set
\begin{equation}
                   s_*\le\min\{s_0/2,\log2/(B+1)\}.
 \label{cci:output-spatial-rate}
\end{equation}
Only then require
\begin{equation}
 \rho\ge\max\{\rho_0,C_LLM_c,64(\log L)/s_*\}.
 \label{cci:actual-separation-admission}
\end{equation}
The last reserve allows later scale conversions and
global cell sums at a smaller spatial rate. All gap
and input-family bounds used to choose $m$ are uniform
above their initial separation threshold $\rho_0$.
Thus no constant used earlier is being determined
by the final geometry. The resulting family satisfies
\begin{equation}
                         \|\mathbf T\|_{s_*,\kappa,c}<1/4
                         \quad\hbox{at real inputs}.
 \label{cci:real-family-contraction}
\end{equation}
This retains the prescribed finite carrier rate
$\kappa$ through the enlarged carrier tile $M_c$.
It does not retain a prescribed spatial rate: the
positive $s_*$ has been decreased explicitly.

\subsubsection{The exact coarse inverse and complex continuation}

\begin{proposition}[Carrier inverses of the actual positive coarse matrices]
\label{prop:cci-positive-coarse-inverse}
For every prescribed finite carrier rate $\kappa$,
choose the parameters and require the separation in
\eqref{cci:actual-separation-admission}. Then both
$M^{-1}$ and $\Theta^{-1}$ have exact local carrier
families. In counting normalization these satisfy
\begin{equation}
     \|\mathbf{\widehat M^{-1}}\|_{s_*,\kappa,c}
       +\|\mathbf{\widehat\Theta^{-1}}\|_{s_*,\kappa,c}
                                                   \le C_\kappa.
 \label{cci:inverse-family-bound}
\end{equation}
The constants are uniform in depth, compatible volume
and admitted geometry. The families are holomorphic
on a common sufficiently small complex $C^1$ tube
about the real regular domain and retain exact
dependence on their recorded background and geometry.
The physical inverses are obtained by $J^{-1}(\cdot)J$.
\end{proposition}
\begin{proof}
For either normalized $A$, let $\star$ denote the
joined-carrier product. The explicit inverse family is
\begin{equation}
 \mathbf B=\beta^{-1}
        \left(\sum_{j=0}^{m-1}\mathbf H^{\star j}\right)
                       \star\sum_{v\ge0}\mathbf T^{\star v}.
 \label{cci:actual-inverse-walk}
\end{equation}
Its norm is at most
$\beta^{-1}\sum_{j<m}K_H^j/(1-\|\mathbf T\|)$.
Every factor is an actual local coefficient and
every product carries its full joined union.
The sum is absolutely convergent in the family norm.
At complete coefficients its operator is
\[
       \beta^{-1}\left(\sum_{j<m}H^j\right)(I-H^m)^{-1}=A^{-1}.
\]
This follows from
$(I-H)\sum_{j<m}H^j=I-H^m$; all factors are
polynomials or convergent functions of $H$.
The real positive matrix is invertible by
\eqref{cci:common-real-gap}. Final transpose
symmetrization of the family preserves this inverse,
its locality and its two-sided norm.

For complex continuation keep the same geometry,
cutoff, powers and filing balls. Suppose
\begin{equation}
          \|\mathbf A(z)-\mathbf A(z_0)\|_{s_0,\kappa_0,0}
                                      <\varepsilon,
 \label{cci:carrier-perturbation}
\end{equation}
where $z_0$ is real. Tail projection is bounded in
the input family norm. Each finite product is
Lipschitz with bound $j(K_H+1)^{j-1}$ times the
perturbation. For the combined near entry, the
same finite matrix identity bounds
$S(z)^m-S(z_0)^m$ in unweighted Schur norm; its
single local ball carrier costs at most
$e^{s_*B+\kappa B_L}$. Thus
$\|\mathbf T(z)-\mathbf T(z_0)\|_{s_*,\kappa,c}
\le C_{m,R,\kappa}\varepsilon$.
Choose $\varepsilon$ so this is below $1/4$.
The actual input-family analyticity proved above
then provides a common physical $C^1$ tube on
which $\|\mathbf T(z)\|<1/2$.

Series \eqref{cci:actual-inverse-walk} is normally
convergent there. Its polynomial identity still
makes it the inverse of the actual complex matrix.
No Hermitian inequality is asserted at a complex
input. Each filed coefficient is analytic in the
data on its union carrier; the absolutely convergent
sum over words with that recorded carrier preserves
this local dependence. Banach Cauchy estimates give
local derivative families on a smaller tube.
Equivalently the actual first derivative has the
local representation $-A^{-1}(DA)A^{-1}$, and higher
derivatives are the finite ordered resolvent
products. Their joined carriers retain every
marked input. This proves the proposition.
\end{proof}

\begin{corollary}[Positive row Schur inverses]
\label{cor:cci-row-inverse}
The same construction applies to an actual coarse
matrix with a bounded number of row types per
assigned cell, after its physical counting-unitary
normalization, whenever its real gap and its actual
input family satisfy \eqref{cci:common-real-gap} and
\eqref{cci:actual-input-family}. Its inverse then has
the bound in \eqref{cci:inverse-family-bound}, with
the corresponding finite constants and carrier tiles.
In particular it applies to
\begin{equation}
 \begin{split}
 T_0^{\mathrm n}&=\Xi_\Gamma^2a_f^{1/2}
                     QG_0Q^\sharp a_f^{1/2}\Xi_\Gamma^{-2},
 \\
 G_0&=G_F+G_F(D_UF)\Theta^{-1}(D_UF)^\sharp G_F.
 \end{split}
 \label{cci:actual-row-Schur}
\end{equation}
There are at most eight rows per assigned top, as in
\eqref{scr:coarse-definitions}. Small symmetric actual family-norm
perturbations of this matrix have the same inverse
construction through \eqref{cci:carrier-perturbation}.
\end{corollary}
\begin{proof}
Bundle the actual row species at each top; their
colour block dimension is at most twenty-four.
The counting proof and Hilbert row estimate use
only this uniform bound. The row metric locations
and their complete carriers differ from their
incident tops by a fixed distance, so the same
local and global geometry estimates apply.

For \eqref{cci:actual-row-Schur}, multiplication by
$K_L=\mathcal A_F-(D_UF)M^{-1}(D_UF)^\sharp$
verifies the displayed Woodbury expression for
$G_0=K_L^{-1}$, using
$\Theta=M-(D_UF)^\sharp G_F(D_UF)$.
It is an exact composition of the local families
already constructed. More explicitly,
\eqref{cci:normalized-F-columns} and
\eqref{cci:normalized-GF-columns} give
$\|1_yG_FD_UFJ^{-1}e_z\|_{\mathcal E_y}
\le Ce^{-s'd(y,z)}$. The only intermediate scale
ratio is $\xi_w/\xi_z$, absorbed by the retained
separation reserve. Thus $G_0$ has the same
source-to-energy family grade as $G_F$.
The normalized trace
$\mathcal Q_n=\Xi_\Gamma^2a_f^{1/2}Q$ is bounded from
field energy to counting row space and has the actual
finite complete carriers. Its counting-to-physical
transpose is $\mathcal Q_n^\sharp
=Q^\sharp a_f^{1/2}\Xi_\Gamma^{-2}$.
The actual row density gives the stronger source bound
\begin{equation}
 |(\mathcal Q_n^\sharp e_c)(b)|\le C\xi_c^{-3}
                         \mathbf1_{\mathrm{carrier}(c)}(b),
 \qquad
 \xi_y\|1_y\mathcal Q_n^\sharp e_c\|_2\le C.
 \label{cci:row-force-cell-injection}
\end{equation}
On its complete carrier $\xi_y/\xi_c$ is bounded by a
fixed $L$ factor, and physical cell volume is $\xi_y^4$.
Thus the injection belongs to the force-cell space
already admitted by the positive walks. Composing these
local injections with $G_0$ gives the actual input
family of $T_0^{\mathrm n}$ in the proved norms.
The real lower bound is the smooth-row-seed
argument in \eqref{scr:schur-energy-lower}; the
upper bound follows from the same trace and
energy inverse bounds.

Choose the Gram and $\Theta$ inverse families
at a stronger intermediate carrier rate first.
Their finite compositions then meet the stronger
input rate needed for this second application.
This adds only finitely many parameter choices:
the final carrier tile and separation are
increased after both constructions. The proof
does not identify that tile with $M_h$.
Finally the perturbation assertion is precisely
the common-tube argument above, once the stated
actual family difference has been bounded.
\end{proof}

\subsubsection{Periods and the admission retained by the result}

For small periods choose one finite periodic cover large
enough for all primitive and coarse patches and both aligned
tile grids. Construct the primitive positive-walk families,
their coarse compositions and the short-power inversion on
that same cover, before any family projection. Thus the
aligned $M_0$ to $M_c$ parent relation is used before
projected translated tag phases appear. The positive
operators and actual coarse matrices descend on the
deck-invariant subspace. If the auxiliary tile grid is
not itself invariant under deck shifts, first average
the final filed inverse family over the deck group,
translating its carrier tags as well. Its total inverse
is unchanged by uniqueness, and its family norm does
not increase. The average is now covariant as a family.

Project its carriers with their level and periodic
identification data retained. A base row sum is bounded
by the cover row sum over all source lifts. By deck
covariance the base column sum is bounded by the cover
column sum over all output lifts. Connected projection
cannot increase carrier length, and base distances
cannot increase. Thus there is no sheet-count loss.
Projected translated-cover tags are part of the
recorded geometry; agreement with an independently
prescribed history tiling is not assumed.

The result supplies actual local coefficients for the
positive Gram and coarse Schur inverses on the admitted
proper row graph. All complex bounds concern the common
physical tube just proved. Independent enlarged cube
input domains, separately prescribed indicator fillings,
and hereditary tower comparisons require their own
admission arguments. In particular a source history
with only its stated $M_h$ separation is not declared
to satisfy \eqref{cci:actual-separation-admission} by
renaming $M_h$ as the carrier size $M_c$.

% END LAYERED EXPANSION background/coarse_carrier_inverse.tex

% BEGIN LAYERED EXPANSION inverse/named_current_local_walks.tex
\subsection{Local carriers for the named current propagators}
\label{ch:ncw-named-current-local-walks}

We assemble the actual site projection, gauge restoration and
two current inverses in their prescribed order. The local
positive walks are those of \cref{lpw:positive-local-walks}.
All coefficients retain the actual complete rows, outward
endpoints and inner plaquette set. The carrier cube size is
\(M_{\rm c}\); it is independent of a prescribed history
cube size \(M_{\rm h}\).
Start with the actual real curvature-small reference and
determining geometry of that positive-walk theorem.
The final separation requirements of the finitely many
coarse-inversion stages are retained. Complex coefficients
range over their common weighted \(C^1\) tube about this
real reference. The independent current is specified below.
The fixed positive site and field row weights are included
in the declared local coefficient data.

Use the shifted connected carrier length \(\sigma_c\), canonical
endpoint tags and short joins of that construction. For a
family \(\mathbf T=(T_C)\) between normalized cell spaces,
write
\begin{equation}\label{ncw:family-norm}
 \begin{split}
 \|\mathbf T\|_{s,\kappa}
 =\max\bigg\{&
 \sup_y\sum_{z,C}e^{s d(y,z)+\kappa\sigma_c(C)}
                                     \|1_yT_C1_z\|,\\
 &\sup_z\sum_{y,C}e^{s d(y,z)+\kappa\sigma_c(C)}
                                     \|1_yT_C1_z\|\bigg\}.
 \end{split}
\end{equation}
Only families whose carriers contain their matrix endpoints
are used. Products retain the indicated short join at each
intermediate cell. The carrier inequality and triangle
inequality make this norm submultiplicative whenever the
intermediate normalized spaces agree.

Site \(L^2\) uses the physical fine measure. The field energy
is \(X\), its dual is \(X^*\), and \(Z\) is the actual site
second-difference space in \eqref{lcu:site-norms}.
Their local norms use the assigned cells and fixed
derivative collars. Energy spaces and their duals are localized
by a smooth square partition on fixed enlargements of the
disjoint pillars. Each profile has first difference at most
\(C_L/\xi_y\), and their overlaps are bounded. Such a
partition is obtained by normalizing squared smooth bumps
on the pillar cubes; overlapping scales are comparable
within \(L\). Multiplication is bounded in the physical
energy and dual norms by the product rule and the mass term.
Inserting this square partition in a product proves the
Schur multiplication estimate, with one fixed localization
constant. Absorb that constant in the algebra norm used
in \eqref{ncw:family-norm}. This avoids a sharp
characteristic-function cutoff in a Sobolev space.
Pure \(L^2\) blocks may use disjoint cell ownership.
All localization collars are retained in the carriers.
Coarse spaces will always specify their
physical-to-counting unitary. No unweighted array transpose
is substituted for a physical transpose.

Fix a finite carrier rate \(\kappa>0\), and leave a strict
reserve in the positive-walk spatial rate. Separation can
be increased so that, for some sufficiently small \(\varepsilon_g\),
\begin{equation}\label{ncw:scale-reserve}
 \max\{\xi_y/\xi_z,\xi_z/\xi_y\}
             \le C_L e^{\varepsilon_g d(y,z)}.
\end{equation}
Indeed crossing one nonadjacent layer costs a fixed multiple
of \(\rho\) in the layer distance, while its scale changes
by a factor \(L\); adjacent-cell ratios cost only \(C_L\).
The same choice leaves the cell-count reserve used below.
All geometric constants, tile choices and rate reserves
are fixed before the depth or volume.

\subsubsection{The output site factor has two differences}

The positive site parametrix uses the local output factor
\[
 L_x=\chi_x E_{D_x}A_{s,D_x}^{-1}E_{D_x}^{\sharp}\chi_x .
\]
Its artificial Dirichlet boundary lies outside the support
of \(\chi_x\), with room for every complete row and
derivative collar. Let \(Z_y\) denote the local \(Z\) norm.

\begin{lemma}[A current-free second-difference output]
\label{lem:ncw-site-output}
For a source in the assigned cell \(B_z\),
\begin{equation}\label{ncw:site-output-H2}
 \|1_yL_x1_zf\|_{Z_y}
       \le C\frac{\xi_z}{\xi_y}
                    e^{-s_1d(y,z)}\|f\|_{2,B_z}
\end{equation}
whenever this local factor is present.
The constants are uniform in the patch radius, mesh and depth.
The coefficient retains the positive-walk carrier, enlarged
only by its fixed derivative collar. No physical-current
bound is a premise.
\end{lemma}
\begin{proof}
Put \(v=E_{D_x}A_{s,D_x}^{-1}E_{D_x}^{\sharp}\chi_x1_zf\)
and \(u=\chi_xv\). The local positive column bound gives
\[
 \|v\|_{\mathcal E_y}
          \le C\xi_z e^{-s_0d(y,z)}\|f\|_{2,B_z}.
\]
Multiplication by \(\chi_x\) preserves that bound on fixed
cell enlargements. Consequently the mass and first-gradient
parts of \(Z_y\) have the stated additional factor
\(\xi_y^{-1}\).

For the second differences the exact compressed equation,
multiplied by the interior cutoff, gives in the original
site space
\begin{equation}\label{ncw:output-equation}
 \Delta_Uu=\chi_x^2 1_zf+
             [A_s,\chi_x]v-q^\sharp a_squ .
\end{equation}
There is no artificial boundary source on the support of
this equation: its cutoff is separated from the compression
boundary by more than the row range.
The slow-commutator estimate bounds its middle term locally
by \(C\xi_y^{-1}\|v\|_{\mathcal E_{\mathcal D_y}}\).
For the last term the actual transported row bound gives
\[
 \|q^\sharp a_squ\|_{2,B_y}
                \le C\xi_y^{-2}\|u\|_{2,\mathcal D_y}
                \le C\xi_y^{-1}\|u\|_{\mathcal E_{\mathcal D_y}} .
\]
The enlargement \(\mathcal D_y\) has bounded metric radius
and bounded cell count, with comparable scales. This uses
the complete average; it is not a point evaluation of \(u\).

At a nonbare cell take a cutoff on a fixed surrounding
cube, equal to one on the derivative collar of \(B_y\).
That cube is inside the actual action region by separation.
The unitary discrete Hodge calculation applied to its
compactly supported scalar field gives
\begin{equation}\label{ncw:local-scalar-H2}
 \|\nabla_UD_Uu\|_{2,B_y}
 \le C\{\|\Delta_Uu\|_{2,\mathcal D_y}
        +\xi_y^{-1}\|D_Uu\|_{2,\mathcal D_y}
        +\xi_y^{-2}\|u\|_{2,\mathcal D_y}\}.
\end{equation}
In that calculation the curl of \(D_Uu\) is a finite
curvature multiplication and shift of \(u\), and the
covariant Hodge commutator is curvature times \(D_Uu\).
Thus it uses no derivative of curvature or of the current.
It is the local compact-support version of the calculation
in \eqref{lsi:second-derivative-hodge}.

Insert \eqref{ncw:output-equation} and the column bound.
The direct source term is present only on a bounded
enlargement of \(B_z\), where its scale ratio and distance
are bounded. All other terms are at most
\(C(\xi_z/\xi_y)e^{-s_1d(y,z)}\|f\|_2\), with a fixed
smaller rate.

At a true outer boundary the assigned scale is \(\eta\).
Each first or second difference has finitely many
neighbours and costs \(\eta^{-1}\) or \(\eta^{-2}\).
The atomic mass and gradient parts of the local energy
therefore bound the same \(Z_y\) norm by
\(C\eta^{-1}\|v\|_{\mathcal E_{\mathcal D_y}}\).
This proves the result without applying an interior
Hodge identity across a missing plaquette.
\end{proof}

\begin{proposition}[A local site inverse in the actual source norm]
\label{prop:ncw-site-H2-family}
The actual site inverse has an exact carrier family with
\begin{equation}\label{ncw:Gs-H2-family}
             \|\mathbf G_s\|_{L^2\to Z;s,\kappa}\le C_\kappa,
 \qquad
             \|D\mathbf G_s\|_{L^2\to X;s,\kappa}\le C_\kappa
\end{equation}
for a fixed \(s>0\).
It is holomorphic on a common smaller complex \(C^1\)
tube about each actual real reference, in the fixed real
physical norms.
\end{proposition}
\begin{proof}
In the positive-walk series replace only its leftmost
output factor by \eqref{ncw:site-output-H2}.
An error factor between source cells, expressed in
unweighted \(L^2\) norms, carries the ratio
\(\xi_{\rm input}/\xi_{\rm output}\), since the positive
source norm is \(\xi\|f\|_2\). The output upgrade has
exactly that same endpoint ratio. In every product these
ratios telescope. The whole walk therefore pays
\(\xi_z/\xi_y\) once, not once at every step.
Use \eqref{ncw:scale-reserve} to absorb it by the fixed
spatial-rate reserve. The positive error-family
contraction and the shifted carrier weight are unchanged.
Its absolute geometric sum proves the first bound.
The local map \(D:Z\to X\) proves the second.

For complex coefficients, use the real reference
Hodge estimate. The weighted \(C^1\) tube gives locally
\[
 \|(\Delta_U-\Delta_{U_0})u\|_2
 \le C\varepsilon\{\xi_y^{-1}\|D_{U_0}u\|_2+
                                  \xi_y^{-2}\|u\|_2\}.
\]
This is the actual first-difference coefficient expansion
used in the complex site comparison; it needs only the
tube's first covariant difference. The local inverse
column bound is already uniform on that tube. The
preceding proof and the equivalence of its fixed norms
therefore apply unchanged. Each factor is a local finite
matrix inverse and a local differentiated cutoff, so it
is holomorphic; normal convergence gives the claimed
carrier-family holomorphy.
\end{proof}

\begin{corollary}[A divergence source without its derivative]
\label{cor:ncw-divergence-source}
The same family gives
\begin{equation}\label{ncw:divergence-family}
                 \|\,\mathbf G_sD^\sharp\,\|_{X^*\to L^2;s,\kappa}
                                                       \le C_\kappa .
\end{equation}
This is the physical transpose of \(D\mathbf G_s:L^2\to X\).
The source \(D^\sharp f\) is interpreted by its variational
pairing. The assertion requires no fine pointwise value
of that divergence.
\end{corollary}

\subsubsection{Variational sources and balanced site fields}

Let \(\mathfrak F\) be the fine force space with local norm
\(\|f\|_{\mathfrak F_y}=\xi_y\|f\|_{2,B_y}\).
It embeds in \(X^*\), with a uniform constant.
Use the site spaces \(Y,Y^\flat\) of \eqref{lcu:site-norms}.

\begin{lemma}[Local positive factors accept variational sources]
\label{lem:ncw-variational-families}
The positive inverse families also have the bounds
\begin{equation}\label{ncw:variational-families}
 \|\mathbf G_F\|_{X^*\to X;s,\kappa}
       +\|\mathbf G_s\|_{Y^\flat\to Y;s,\kappa}\le C_\kappa .
\end{equation}
The site divergence \(D^\sharp f\) can be used as a
\(Y^\flat\) source whenever \(\|\ell f\|_2\) is finite.
All these statements use physical variational pairings.
\end{lemma}
\begin{proof}
The local weighted energy proof in
\cref{lpw:local-columns} applies to a functional
supported on a cell, with its local energy-dual norm:
testing the compressed equation and its exponential
conjugate uses only that norm. Multiplying the functional
by the smooth input cutoff is bounded by the same product
rule. Thus the local inverse has energy-dual to energy
columns with the same decay.

The commutator in \cref{lpw:slow-commutator} maps that
energy output into the \(L^2\) force space \(\mathfrak F\).
Consequently the rightmost error factor in a nonempty
positive walk maps an energy-dual source to \(\mathfrak F\);
all subsequent error factors act in the already controlled
\(\mathfrak F\) space. The zero-error term is the local
variational inverse itself. The same absolute geometric
sum therefore proves the energy-dual inverse family.
For the vector operator its energy space is \(X\).

For the scalar operator the local unbalanced energy is
\(\mathcal E_y\simeq\xi_y^{-1}Y_y\). Its local dual,
when used as a source in the unweighted equation,
has norm \(\xi_y^{-1}Y_y^\flat\), because
\(Y^\flat\) uses the pairing with \(\ell^2u\).
Thus the scalar column picks up the ratio
\(\xi_y/\xi_z\). Absorb this once by
\eqref{ncw:scale-reserve}, as in the preceding proof.
This gives the second bound in
\eqref{ncw:variational-families}.
Finally
\[
 |\langle D^\sharp f,\ell^2u\rangle|
       =|\langle f,D(\ell^2u)\rangle|
       \le C\|\ell f\|_2\|u\|_Y.
 \]
No derivative of \(f\), and hence no derivative of an
independent current, is taken.
\end{proof}

\subsubsection{Physical top units, projection and restoration}

Write \(J_s=\Xi^2\) for the unitary from the physical top
space to counting \(\ell^2\), and define
\begin{equation}\label{ncw:normalized-site-columns}
 \widehat F=FJ_s^{-1}=G_sq^\sharp\Xi^{-4},\qquad
 \widehat M=J_sMJ_s^{-1}=\widehat F^\sharp\widehat F .
\end{equation}
The transpose in the last equality is between counting
top space and the physical fine site space. In particular
\(\widehat F^\sharp=J_sF^\sharp\); this formula keeps
the physical top measure.

The actual injection \(q^\sharp\Xi^{-4}\) maps a unit
counting top datum to density \(C\xi_z^{-4}\) on its
pillar. Its local \(Y^\flat\) norm is bounded, since
that norm supplies \(\xi_z^2\) and the physical
\(L^2\) norm supplies another \(\xi_z^2\).
Consequently \eqref{ncw:variational-families} proves
\begin{equation}\label{ncw:F-carrier-grades}
 \|\mathbf{\widehat F}\|_{\ell^2_{\rm top}\to Y;s,\kappa}
 +\|\mathbf{\widehat F}\|_{\ell^2_{\rm top}\to L^2;s,\kappa}
 +\|D\mathbf{\widehat F}\|_{\ell^2_{\rm top}\to\mathfrak F;s,\kappa}
                                                     \le C_\kappa .
\end{equation}
Equivalently its unabsorbed local \(L^2\) column has
factor \(\xi_y/\xi_z\), as in
\eqref{cci:normalized-F-columns}. No fine-coordinate
dimension appears in this normalization.

Choose the positive coarse inverse families of
\cref{prop:cci-positive-coarse-inverse} at a stronger
intermediate carrier rate \(\kappa_1>\kappa+1\).
Their input positive families are first obtained at
another fixed rate \(\kappa_2>\kappa_1+1\).
All ensuing finite reblockings use aligned parents and
the recorded joins. Work on a common final periodic
cover when needed, so the canonical tags for all
products agree before descent.
Until the row-Schur inversion, the following compositions
are made at \(\kappa_1\); their displayed weaker bounds at
\(\kappa\) follow by decreasing the carrier weight.

\begin{proposition}[Local projection and gauge restoration]
\label{prop:ncw-projection-restoration}
The actual maps
\begin{equation}\label{ncw:projection-restoration}
 P=\widehat F\widehat M^{-1}\widehat F^\sharp,\quad
 R=1-P,\quad \Lambda=G_sRD^\sharp,\quad S=1-D\Lambda
\end{equation}
have exact local carrier families, with bounds
\par\noindent
\begin{minipage}{\linewidth}
\begin{align}
 \|\mathbf P\|_{L^2\to Y;s,\kappa}
 +\|\mathbf R\|_{L^2\to L^2;s,\kappa}
 +\|\mathbf R\|_{Y\to Y;s,\kappa}&\le C_\kappa,\notag\\
 \|\mathbf\Lambda\|_{X\to Z;s,\kappa}
 +\|D\mathbf\Lambda\|_{X\to X;s,\kappa}
 +\|\mathbf S\|_{X\to X;s,\kappa}
 +\|\mathbf S^\sharp\|_{\mathfrak F\to\mathfrak F;s,\kappa}
                                                      &\le C_\kappa .
 \label{ncw:projection-family-bounds}
\end{align}
\popQED\par
\end{minipage}
\end{proposition}
\begin{proof}
Compose \eqref{ncw:F-carrier-grades} with the actual
normalized Gram inverse. Its transpose has
\(\widehat F^\sharp:L^2\to\ell^2_{\rm top}\).
This proves \(P:L^2\to Y\), hence the two bounds for \(R\).
The full physical divergence maps \(X\) locally into
site \(L^2\). Near a true boundary its one-edge term
is paid by the bare mass. Composing it with \(R\)
and \eqref{ncw:Gs-H2-family} proves the bounds for
\(\Lambda,D\Lambda,S\).

For the last bound use the exact transpose
\[
                    S^\sharp=1-DRG_sD^\sharp .
 \]
Here \(D^\sharp:\mathfrak F\to Y^\flat\) is a
variational divergence, \(G_s:Y^\flat\to Y\),
\(R:Y\to Y\), and \(D:Y\to\mathfrak F\).
Every factor has the claimed local family. This is
a weighted \(L^2\) force assertion; it is not an
\(L^\infty\) Riesz assertion.
\end{proof}

\subsubsection{The zero-current Landau inverse}

Put \(\widehat B=D\widehat F\) and
\(\widehat\Theta=J_s\Theta J_s^{-1}\).
The exact positive factorization is
\begin{equation}\label{ncw:zero-current-Woodbury}
 G_0=G_F+
          G_F\widehat B\widehat\Theta^{-1}\widehat B^\sharp G_F .
\end{equation}
It is the physical identity \eqref{lcu:Woodbury} in
counting top units. The injection
\(\widehat B:\ell^2_{\rm top}\to\mathfrak F\)
is \eqref{ncw:F-carrier-grades}; its transpose acts
from field energy into counting top space.
Thus the positive coarse families and the positive
field families prove
\begin{equation}\label{ncw:G0-family}
 \|\mathbf G_0\|_{X^*\to X;s,\kappa_1}
       +\|\mathbf G_0\|_{\mathfrak F\to X;s,\kappa_1}
                                                       \le C_{\kappa_1}.
\end{equation}
There is no replacement of the actual projected
precision by a periodic model in this argument.
The inverse \(\widehat\Theta^{-1}\) contains its
complete geometry carrier.

\subsubsection{The first current precision and its lift}

Restrict \(J\) to bonds touching \(\Omega_1\), as allowed
by the exact identities in
\cref{ch:actual-current-completion}, and put
\[
             \gamma_{\rm t}=\sup_{b\ {\rm touches}\ \Omega_1}
                                        \mathsf L(b)^3|J(b)|.
\]
We make this restriction before forming the coefficient
families. It is a global exact identity for the field
operator; an individual inverse coefficient need not
vanish on bare bonds.

\begin{lemma}[The actual current family is small]
\label{lem:ncw-current-family}
The canonical symmetric correction has an exact carrier
family with
\begin{equation}\label{ncw:current-family-bound}
 \|\mathbf B_U[J]\|_{X\to\mathfrak F;s,\kappa_1}
       +\|\mathbf B_U[J]\|_{X\to X^*;s,\kappa_1}
                                      \le C_{\kappa_1}\gamma_{\rm t}.
\end{equation}
No derivative of \(J\) is used.
\end{lemma}
\begin{proof}
For \(\lambda=\Lambda v\), the exact grouped formula is
\begin{equation}\label{ncw:grouped-B}
 B_U[J]v=W[J]\lambda+DRG_s h_J,\qquad
 h_J=W[J]^\sharp v-\tfrac12
       \{D^\sharp W[J]\lambda+W[J]^\sharp D\lambda\}.
\end{equation}
The actual endpoint current bounds in
\eqref{lcu:local-coefficient-bounds} show that
\(W[J]\lambda\) is in \(\mathfrak F\), with norm
\(C\gamma_{\rm t}\|\lambda\|_Z\), while \(h_J\)
is in \(Y^\flat\), with norm
\(C\gamma_{\rm t}(\|v\|_X+\|\lambda\|_Z)\).
These bounds hold as carrier bounds: \(W[J]\) reads
one bond and its two endpoint values, and
\(D^\sharp W[J]\lambda\) is kept as a divergence.
Compose with the families
\(G_s:Y^\flat\to Y\), \(R:Y\to Y\) and
\(D:Y\to\mathfrak F\).
Every union keeps the actual input current bond and
both matrix endpoints. The two Schur sums give
\eqref{ncw:current-family-bound}.
\end{proof}

The first named inverse is now defined without an
implicit change of completion:
\begin{equation}\label{ncw:first-inverse-series}
 \begin{split}
 \mathbf G&=\mathbf G_0\star
                   \sum_{n\ge0}(\mathbf B_U[J]\star\mathbf G_0)^{\star n},\\
               \sum_CG_C&=G=(G_0^{-1}-B_U[J])^{-1}.
 \end{split}
\end{equation}
The right side denotes its joined-carrier family;
at complete coefficients it is the displayed matrix
inverse. Choose the touching-current ceiling so
the \(\mathfrak F\to\mathfrak F\) error norm is less
than \(1/4\). For a general \(X^*\) source only the
rightmost current factor uses
\(B_U[J]G_0:X^*\to\mathfrak F\); all subsequent
loops use \(\mathfrak F\to\mathfrak F\).
This proves exact normally convergent families for
\[
                         G:X^*\to X,\qquad G:\mathfrak F\to X,
\]
and the same family norm for \(G-G_0\) is
\(C_{\kappa_1}\gamma_{\rm t}\).

Let the physical field-row measure be \(\xi_c^4\) and put
\begin{equation}\label{ncw:normalized-field-row}
 J_\Gamma=\Xi_\Gamma^2,\qquad
 \mathcal Q_n=J_\Gamma a_f^{1/2}Q,\qquad
       \mathcal Q_n^\sharp=Q^\sharp a_f^{1/2}J_\Gamma^{-1}.
\end{equation}
This trace maps \(X\) to counting row space; its
transpose maps counting rows into \(\mathfrak F\).
The latter is the actual density computation in
\eqref{cci:row-force-cell-injection}. Define
\begin{equation}\label{ncw:first-row-Schur}
 \mathcal T_0=\mathcal Q_nG_0\mathcal Q_n^\sharp,\qquad
 \mathcal T_J=\mathcal Q_nG\mathcal Q_n^\sharp .
\end{equation}
The real \(\mathcal T_0\) has the actual SCR lower
and upper gap. Apply \cref{cor:cci-row-inverse} to
its actual family, using the intermediate rate
\(\kappa_1\). It gives a family for
\(\mathcal T_0^{-1}\) at the final rate \(\kappa\),
with a final aligned tile size \(M_{\rm c}\) and
a positive final spatial rate. Reblock all preceding
families into those parents and reduce their spatial
rates to a common positive value.
The actual difference satisfies
\[
       \|\mathbf{\mathcal T_J}-\mathbf{\mathcal T_0}\|_{s,\kappa}
                                  \le C_\kappa\gamma_{\rm t}.
\]
A second fixed decrease of the current ceiling therefore
gives the exact inverse family
\begin{equation}\label{ncw:first-Schur-inverse}
 \mathbf{\mathcal T_J^{-1}}
       =\mathbf{\mathcal T_0^{-1}}\star
           \sum_{n\ge0}
           \bigl[-(\mathbf{\mathcal T_J}-\mathbf{\mathcal T_0})
                    \star\mathbf{\mathcal T_0^{-1}}\bigr]^{\star n}.
\end{equation}
At a complex background the same argument is centered
at the real reference matrix and includes its small
actual physical-tube family perturbation.

Write \(\mathcal S_J=QGQ^\sharp\).
The normalized first lift and the physical lift are
\begin{equation}\label{ncw:first-lift}
 \mathcal H_J=G\mathcal Q_n^\sharp\mathcal T_J^{-1},\qquad
 H=\mathcal H_JJ_\Gamma a_f^{1/2}
                         =GQ^\sharp\mathcal S_J^{-1}.
\end{equation}
Thus \(\mathcal H_J\) has a local carrier family from
counting rows into \(X\), and \(H\) has one from the
physical datum norm \(\|a_f^{1/2}b\|_{\rm rows}\).
In particular \(QH=1\) as an exact identity of the sums.
No termwise row or gauge identity has been assumed.

\subsubsection{The first-lift induced multiplier retains its root}

Let \(J_{\rm full}\) be the original current and
\(J=J_{\rm full}^{\rm t}\) the restriction used above.
Define
\[
                       \omega^{\rm t}=H^\sharp J .
\]
For a carrier coefficient \(H_C\), retain the coefficient
\(\omega^{\rm t}_{c,C}=((H_C)^\sharp J)(c)\).
The carrier contains the row \(c\), every paired fine
current bond, and the background data read by \(H_C\).

\begin{proposition}[An absolutely rooted multiplier family]
\label{prop:ncw-rooted-multiplier}
The first-lift readout satisfies
\begin{equation}\label{ncw:multiplier-family-bound}
        \sup_c\sum_C e^{\kappa\sigma_c(C)}
                       \xi_c^3|\omega^{\rm t}_{c,C}|
                                           \le C_\kappa\gamma_{\rm t}.
\end{equation}
The actual multiplier for the unrestricted bare current is
\begin{equation}\label{ncw:bare-multiplier-split}
 \omega_* =H^\sharp J_{\rm full}
             =\omega^{\rm t}+\iota_0J_{\rm full}^{\rm b}.
\end{equation}
In particular its nonbare family obeys
\eqref{ncw:multiplier-family-bound}. An all-row bound
retains the full current radius
\(\gamma_{\rm full}=\sup_b\mathsf L(b)^3|J_{\rm full}(b)|\).
\end{proposition}
\begin{proof}
Fix a unit colour \(e\) at row \(c\) and use
\(b^{(c)}=\xi_c^{-1}e_c\).
Its physical datum energy is bounded above and below
by constants. The exact physical transpose gives
\begin{equation}\label{ncw:rooted-duality}
 \xi_c^3\langle\omega^{\rm t}_{c,C},e\rangle
                          =\langle H_Cb^{(c)},J\rangle_{\rm fine}.
\end{equation}
Split the fine pairing by disjoint bond ownership cells.
On such a cell
\(\|\xi J\|_2\le C_L\gamma_{\rm t}\), while the mass
part of \(X_y\) bounds
\(\|\xi^{-1}H_Cb^{(c)}\|_2\).
Therefore the left side of
\eqref{ncw:multiplier-family-bound} is at most
\[
 C\gamma_{\rm t}
       \sum_{y,C}e^{\kappa\sigma_c(C)}
                   \|1_yH_Cb^{(c)}\|_{X_y}\le C_\kappa\gamma_{\rm t}.
\]
This is the actual family column sum; no whole-volume
energy norm is charged to \(J\).

For the bare split use the exact identities
\(Q_0H=P_0\) and equality of the physical fine and
level-zero row measures, both \(\eta^4\).
They give \(H^\sharp J_{\rm full}^{\rm b}
=\iota_0J_{\rm full}^{\rm b}\), as in
\eqref{acc:bare-current-readout}.
This is an identity of the complete sums. It does
not assert that a nonbare column of each \(H_C\)
vanishes at every bare bond. We used the touching
restriction before forming these families for
precisely this reason. The extra bare family is
local and has amplitude bounded by the full
\(\gamma_{\rm full}\), rather than by
\(\gamma_{\rm t}\).
\end{proof}

\subsubsection{Only then form the second precision}

The actual bare re-centered mean is linear, so its
second derivative is zero. Thus
\(\Delta_U^{(2)}[\omega_*]
=\Delta_U^{(2)}[\omega^{\rm t}]\).
Form the constraint family by joining each rooted
coefficient \(\omega^{\rm t}_{c,C}\) to the complete
two-leg carrier of its actual row \(c\).
The physical two-leg bound then gives
\begin{equation}\label{ncw:constraint-family}
 \|\mathbf C_{\omega_*}\|_{X\to\mathfrak F;s,\kappa}
              \le C_\kappa\gamma_{\rm t},\qquad
 \mathbf C_{\omega_*}=\mathbf{\Delta_U^{(2)}[\omega_*]}.
\end{equation}
To check this normalization, the row coefficient has
fixed-leg density \(C_L\xi_c(\eta/\xi_c)^4\).
Its physical transpose has factor \(\eta^{-4}\xi_c^4\).
Inserting \(|\omega(c)|\le C\gamma_{\rm t}\xi_c^{-3}\)
and the field mass norm gives
\(\|\ell C_{\omega_*}v\|_2\le
C\gamma_{\rm t}\|\ell^{-1}v\|_2\).
Both leg sums are retained. There are only a bounded
number of incident rows and neighbouring assigned cells.
For a coefficient family, use the absolute sum
\eqref{ncw:multiplier-family-bound} in place of the
single multiplier bound. The row star changes the
carrier length by a fixed constant and its input-output
metric radius is bounded. This proves
\eqref{ncw:constraint-family} as a family estimate.

By \eqref{ncw:projection-family-bounds},
\begin{equation}\label{ncw:second-small-family}
 \mathbf V_*=\mathbf S^\sharp\star
                         \mathbf C_{\omega_*}\star\mathbf S,
 \qquad
              \|\mathbf V_*\|_{X\to\mathfrak F;s,\kappa}
                                      \le C_\kappa\gamma_{\rm t}.
\end{equation}
The exact second inverse is consequently
\begin{equation}\label{ncw:second-inverse-series}
 \begin{split}
 \mathbf G_1&=\mathbf G\star\sum_{n\ge0}
                                  (\mathbf V_*\star\mathbf G)^{\star n},\\
 \sum_C(G_1)_C&=G_1
        =(G^{-1}-S^\sharp\Delta_U^{(2)}[\omega_*]S)^{-1}.
 \end{split}
\end{equation}
Decrease the same fixed touching-current ceiling so
the force-space error has family norm less than \(1/4\).
As for the first inverse, a general \(X^*\) source
is used only at the rightmost factor; subsequent
loops take place in \(\mathfrak F\).
This proves carrier bounds for \(G_1:X^*\to X\)
and \(G_1:\mathfrak F\to X\), and for
\(G_1-G_0\) with the additional factor
\(\gamma_{\rm t}\).
The lift inside \(\omega_*\) throughout this
construction is \(H\), not the following \(H_1\).

Put
\begin{equation}\label{ncw:second-row-lift}
 \mathcal T_1=\mathcal Q_nG_1\mathcal Q_n^\sharp,\qquad
 \mathcal H_1=G_1\mathcal Q_n^\sharp\mathcal T_1^{-1},
 \qquad H_1=\mathcal H_1J_\Gamma a_f^{1/2}.
\end{equation}
The difference \(\mathcal T_1-\mathcal T_0\) has
family norm \(C_\kappa\gamma_{\rm t}\).
Use the already constructed \(\mathcal T_0^{-1}\)
in the same Neumann formula as
\eqref{ncw:first-Schur-inverse}.
This constructs the exact \(\mathcal T_1^{-1}\)
and \(H_1\) families without a new positivity
assumption at an independent current.
At complete coefficients,
\[
 H_1=G_1Q^\sharp(QG_1Q^\sharp)^{-1},
                         \qquad QH_1=1 .
\]

\subsubsection{The assembled range and its limits}

\begin{proposition}[Named current carriers in physical norms]
\label{prop:ncw-named-families}
Fix a finite carrier rate \(\kappa\).
Choose the two positive coarse-inversion stages,
their carrier tiles, separation reserves and common
complex \(C^1\) tube as above. There is a fixed
touching-current ceiling \(\gamma_*>0\) such that
for \(\gamma_{\rm t}<\gamma_*\) the actual canonical
operators have exact, normally convergent local
carrier families in the following spaces:
\begin{align}
 &G_s:L^2\to Z,\quad R:L^2\to L^2,\quad
                     \Lambda:X\to Z,\quad S:X\to X,\notag\\
 &G_0,G,G_1:X^*\to X,\qquad
                   G_0,G,G_1:\mathfrak F\to X,\notag\\
 &\mathcal H_J,\mathcal H_1:
                         \ell^2_{\rm rows}\to X,\qquad
 \mathcal T_J^{-1},\mathcal T_1^{-1}:
                      \ell^2_{\rm rows}\to\ell^2_{\rm rows}.
 \label{ncw:assembled-families}
\end{align}
Their family norms are bounded by fixed \(C_\kappa\).
The nonbare first-lift readout has the rooted bound
\eqref{ncw:multiplier-family-bound}. The physical
Schur inverses and lifts are obtained by the displayed
row unitaries and weights, not by changing transposes.

Every coefficient is holomorphic in the actual local
background and touching-current data on the stated
tubes, with its complete recorded determining geometry.
Agreement of these data on its carrier gives the same
coefficient. The constants are independent of depth,
compatible volume and admitted region shape.
The current ceiling depends on the fixed carrier
constants and has no H\"older-exponent dependence.
\end{proposition}
\begin{proof}
The successive constructions prove each assertion in
the order \(R,\Lambda,G_0,G,H,\omega_*,G_1,H_1\).
All operations are finite local coefficients, positive
coarse inverse families or normally convergent
small-current series in the same joined-carrier
algebra. Each product retains its endpoint tags and
fixed short joins. The cutoff-source and scale-ratio
reserves were paid before these compositions.

Only finitely many positive inversion and reblocking
stages are used. Choose the intermediate carrier
rates first, then their finite short-power parameters,
the final aligned \(M_{\rm c}\), and the common
separation threshold. Finally choose the common
physical tube and the current ceiling smaller than
the finitely many displayed Neumann bounds. None
of these steps uses a H\"older exponent or a current
derivative. The \(H^2\) output argument uses only
the \(L^2\) unitary Hodge estimate.

For small periods carry out this entire finite
construction on one sufficiently large compatible
periodic cover with common tag grids. Average the
final families over simultaneous deck shifts and
project their recorded tags. The row and column
absolute sums descend as in the positive and coarse
walk proofs. Their shifted joins and physical
unitaries are retained. Thus no sheet-count
factor enters the bound or the current ceiling.
\end{proof}

For clarity, write \(G_J=G\). The physical row Schur inverse is
\[
 (QG_iQ^\sharp)^{-1}
    =a_f^{1/2}J_\Gamma^{-1}\mathcal T_i^{-1}
                                J_\Gamma a_f^{1/2},
                     \qquad i=J,1 .
\]
In the physical angle and dual-current units
\(\beta=\Xi_\Gamma b\), \(\nu=\Xi_\Gamma^3\mu\),
put \(c_a=\Xi_\Gamma a_f^{1/2}\). The mass-removed
coarse operator therefore has the exact normalized family
\begin{equation}\label{ncw:mass-removed-coarse-family}
 \Xi_\Gamma^3\{(QG_iQ^\sharp)^{-1}-a_f\}\Xi_\Gamma^{-1}
                      =c_a(\mathcal T_i^{-1}-1)c_a .
\end{equation}
This identity retains the artificial mass subtraction.
The separate lower bound on the actual outer plain-constraint
fibre is proved in \cref{pgc:proper-covariance}, with the full
outer constraint curvature and the specified relative contour
quadratic form. Its physical absolute-contour identification
holds in an averaged-contour axial representative.

The result gives the canonical completion selected in
\cref{ch:actual-current-completion}. The first-order output
families on the same $C^1$ domain and their enlarged cube-parameter
discs are constructed in \cref{dnr:linear-families,dnr:parameter-families}.
All prescribed histories must use the common separation and
filing constants of these constructions. The positive clean
conditional measure and its normalizer require their own
measure and gauge identifications.

% END LAYERED EXPANSION inverse/named_current_local_walks.tex

% BEGIN LAYERED EXPANSION discrete_schauder_estimates.tex
\subsection{Second derivatives in a discrete H\"older norm}
\label{dsc:section}

The bounded-force estimates give a first derivative
in every fixed H\"older class below exponent one.
For the nonlinear gauge slice, a site gauge
requires two derivatives in the same H\"older
class. This needs cancellation of the
second-resolvent kernel, in addition to its
finite-\(p\) bounds. We prove that estimate and
then apply it to the actual site inverse.

Fix \(0<\alpha<1\). On \(h\mathbb Z^4\) write
\[
 \|f\|_{C^\alpha_h}
   =\|f\|_\infty+
       \sup_{0<|x-x'|\le1}
          \frac{|f(x)-f(x')|}{|x-x'|^\alpha}.
\]
Only distinct lattice points occur in the
seminorm. First and second difference norms
use all forward or backward component
differences. On a fixed interior box the
same notation restricts to present stencils.
All statements concern the physical lattice
measure of \cref{dcz:section}.

\subsubsection{Cancellation of a truncated second kernel}

Let \(k_h=D_i^\sigma D_j^\tau(L_h+1)^{-1}\)
denote its physical convolution kernel.
In addition to \cref{dcz:kernel-regularity},
it obeys
\begin{equation}\label{dsc:far-kernel}
 |k_h(x)|+|Dk_h(x)|\le C_N|x|^{-N},
                           \qquad |x|\ge1,
\end{equation}
for any fixed sufficiently large \(N\).
Indeed insert the heat representation and
use \eqref{dcz:heat-pointwise} with spatial
power \(N\), retaining the factor \(e^{-t}\).
For a derivative of order \(a=2,3\), its
integral is bounded by
\[
 C_N|x|^{-N}\int_0^\infty
         e^{-t}(t+h^2)^{(N-4-a)/2}\,dt.
\]
Choose \(N>4+a\). The integral is uniformly
bounded for \(h\le1\). No exponential
off-diagonal assertion is needed here.

\begin{lemma}[Uniform truncated cancellation]
\label{dsc:truncated-cancellation}
The kernel has total physical sum zero and
\begin{equation}\label{dsc:truncated-cancellation-bound}
       \left|h^4\sum_{|x|\le r}k_h(x)\right|\le C,
                                           \qquad r\ge h.
\end{equation}
The constant is independent of the mesh and
of the truncation radius.
\end{lemma}

\begin{proof}
For fixed \(h\) the massive kernel is absolutely
summable by the heat \(L^1\) estimate in
\cref{dcz:uniform-Lp}. Each differentiated heat
kernel has sum zero, so its time integral has
sum zero as well. For the truncated sum split
the time integral at \(r^2\).

For \(0<t<r^2\), replace the sum over the ball
by the negative sum over its complement.
The differentiated heat tail bound
\eqref{bg:heat-tails}, in its infinite-lattice
version used in \cref{dcz:section}, gives
\[
 \left|h^4\sum_{|x|\le r}D_i^\sigma D_j^\tau p_t^h(x)\right|
 \le C_N(t+h^2)^{-1}
           (1+r/\sqrt{t+h^2})^{-N}.
\]
Its integral is bounded uniformly: put
\(v=t+h^2\), use \(h\le r\), and bound it by
\(C_Nr^{-N}\int_0^{2r^2}v^{N/2-1}\,dv\).
For \(t\ge r^2\), use the physical volume
\(Cr^4\) of the lattice ball and the
pointwise derivative bound \(Ct^{-3}\).
The integral is at most
\(Cr^4\int_{r^2}^\infty t^{-3}\,dt=C\).
The mass factor only improves both bounds.
\end{proof}

The same bound holds for a ball whose centre
is displaced from the kernel centre by at
most one quarter of its radius. Compare
concentric inner and outer balls of comparable
radii; their difference is an annulus on
which the absolute kernel bound has total
mass at most \(C\). This observation will
control the constant part of a nearby
H\"older increment.

\subsubsection{The scalar H\"older multiplier}

\begin{proposition}[Uniform discrete Schauder multiplier]
\label{dsc:holder-multiplier}
For every \(0<\alpha<1\),
\begin{equation}\label{dsc:holder-multiplier-bound}
 \|D_i^\sigma D_j^\tau(L_h+1)^{-1}f\|_{C^\alpha_h}
                                   \le C_\alpha\|f\|_{C^\alpha_h}.
\end{equation}
The same constant holds on every compatible
periodic torus. At fixed mesh the convolution
is the ordinary absolutely summable massive
kernel; no principal value is left undefined.
\end{proposition}

\begin{proof}
Write \(Tf=k_h*f\). Cancellation gives
\[
 Tf(x)=h^4\sum_w k_h(x-w)[f(w)-f(x)].
\]
For \(|w-x|\le1\), use the H\"older
increment and
\(|k_h(z)|\le C(|z|+h)^{-4}\).
The sum over dyadic annuli is at most
\(C_\alpha[f]_\alpha\), including its
last atomic annulus. For larger distances,
\eqref{dsc:far-kernel} gives
\(C\|f\|_\infty\). This proves the
supremum bound.

For the seminorm let
\(\delta=|x-x'|\), \(h\le\delta\le1/16\).
Since each translate of the kernel has
sum zero,
\[
 Tf(x)-Tf(x')
   =h^4\sum_w[k_h(x-w)-k_h(x'-w)]
                                      [f(w)-f(x)].
\]
Split the sum into \(|w-x|\le4\delta\)
and its complement. In the near part,
the first kernel is multiplied by a
H\"older increment at its own centre,
so its absolute sum is
\(C_\alpha[f]_\alpha\delta^\alpha\).
For the second kernel write
\[
 f(w)-f(x)=[f(w)-f(x')]+[f(x')-f(x)].
\]
The first term has the same bound. The
constant second term is multiplied by
the kernel sum over a ball of radius
\(4\delta\) centred a distance \(\delta\)
from its kernel centre.
\Cref{dsc:truncated-cancellation} and
the following annulus comparison bound
that sum by \(C\). Its contribution is
therefore also at most
\(C[f]_\alpha\delta^\alpha\).

In the far part, telescope a coordinate
path from \(x\) to \(x'\). The derivative
bound of \cref{dcz:kernel-regularity}
gives
\[
 |k_h(x-w)-k_h(x'-w)|
       \le C\delta|x-w|^{-5}
                         \quad(4\delta<|x-w|\le1).
\]
Multiplication by
\([f]_\alpha|x-w|^\alpha\) and
dyadic summation gives
\[
 C_\alpha[f]_\alpha\delta
             \sum_{\delta\le r\le1,\ {\rm dyadic}}r^{\alpha-1}
                \le C_\alpha[f]_\alpha\delta^\alpha.
\]
At distance at least one, the differentiated
far bound \eqref{dsc:far-kernel} gives
\(C\delta\|f\|_\infty\), which is bounded
by \(C\delta^\alpha\|f\|_\infty\).
For \(1/16<\delta\le1\), the already
proved supremum bound suffices.
If the mesh exceeds \(1/16\), this last
case covers every distinct pair.

A periodic input extends to the infinite
lattice with the same supremum norm and
no larger local H\"older seminorm: the
periodic distance is at most the distance
of any chosen lifts. Convolution by the
absolutely summable kernel is its periodic
massive multiplier. Applying the estimate
to that extension proves the periodic
claim, independently of the periods.
\end{proof}

\subsubsection{Interior flat and covariant second derivatives}

\begin{corollary}[Interior second-derivative estimate]
\label{dsc:interior-Schauder}
Let \(Q_r\subset Q_R\) be fixed nested physical
boxes with room for the difference stencils.
For their sampled smooth cutoffs,
\begin{equation}\label{dsc:interior-Schauder-bound}
 \|D^2u\|_{C^\alpha(Q_r)}
       \le C_{\alpha,r,R}
          \{\|L_hu\|_{C^\alpha(Q_R)}
                             +\|u\|_{C^{1,\alpha}(Q_R)}\}.
\end{equation}
The constant is uniform in the fine mesh.
Here \(\|u\|_{C^{1,\alpha}}\) means
\(\|u\|_\infty+\|Du\|_\infty+[Du]_\alpha\).
\end{corollary}

\begin{proof}
Take a cutoff \(\chi=1\) on \(Q_r\), with
support a fixed distance inside \(Q_R\).
Extend \(v=\chi u\) by zero. Then
\[
 D^2v=D^2(L_h+1)^{-1}(L_h+1)v.
\]
The discrete product rule writes its source
as \(\chi L_hu\), \(\chi u\), and finitely
many shifted products of \(D\chi,Du\) or
\(D^2\chi,u\). Their global \(C^\alpha\)
norms have the displayed bound. All
neighbours lie in the indicated outer
box, including at the support boundary.
Apply \cref{dsc:holder-multiplier}.
\end{proof}

Suppose now that \(U_0=e^{iha_0}\) in a
local real Coulomb frame and that
\begin{equation}\label{dsc:regular-reference}
                  \|a_0\|_{C^{1,\alpha}}\le C_0.
\end{equation}
This is available at an actual small-curvature,
small-current reference from
\eqref{bg:local-holder-bound}; the current is
used here, in contrast with the
curvature-only estimates in
\cref{ch:proper-strong-response}.
The exact scalar decomposition
\eqref{psr:scalar-decomposition} has
first and zeroth coefficients bounded
in \(C^\alpha\) by a constant depending
only on \(C_0\).
For the zeroth coefficient this uses
one difference of \(a_0\), not two:
the divided difference of
\((\operatorname{Ad}e^{iha_0}-1)/h\)
is an analytic bounded coefficient
times a first difference of \(a_0\).

Consequently a site solution
\(\Delta_{U_0}u=g\) satisfies
\begin{equation}\label{dsc:covariant-interior}
 \|D_{U_0}^2u\|_{C^\alpha(Q_r)}
       \le C_{\alpha,C_0,r,R}
                 \{\|g\|_{C^\alpha(Q_R)}
                             +\|u\|_{C^{1,\alpha}(Q_R)}\}.
\end{equation}
Indeed its flat Laplacian is \(g\)
minus the lower-order terms in that
decomposition. Products preserve
\(C^\alpha\), so
\cref{dsc:interior-Schauder} applies.
Covariant second differences differ
from ordinary second differences by
terms containing \(a_0Du\), \(Da_0u\)
and bounded quadratic coefficients
in \(a_0\). Their \(C^\alpha\) norms
have the same bound. All comparisons
are in this fixed local frame;
transported seminorms give the
gauge-covariant statement.

\subsubsection{The proper site gauge inverse}

Fix the real proper reference and its actual
\(q_0,G_s,F,M,R_0\), with their physical
site and top pairings. Assume the actual
local curvature and current bounds needed
for \eqref{dsc:regular-reference} on the
same complete interior collars.
Let
\begin{equation}\label{dsc:site-source-norm}
 \|\rho\|_{\alpha,2}
   =\sup_y\{\xi_y^2\|\rho\|_{\infty,B_y}
                   +\xi_y^{2+\alpha}[\rho]_{\alpha,y}^{U_0}\}.
\end{equation}
At bare cells the finite collar convention
of \cref{ch:proper-strong-response} is used.
The seminorm compares pairs across adjoining
pillar faces on the same complete interior
boxes used there. Taking it separately
inside each pillar would admit discontinuous
piecewise-constant sources and is not the
norm in this statement. This same complete-box
convention is used for the derivatives below.
For a site gauge parameter define
\begin{equation}\label{dsc:site-gauge-norm}
 \begin{split}
 \|\lambda\|_{Y^{2,\alpha}}=\sup_y\{&
       \|\lambda\|_{\infty,B_y}
       +\xi_y\|D_{U_0}\lambda\|_{\infty,B_y}\\
       &+\xi_y^2\|D_{U_0}^2\lambda\|_{\infty,B_y}
       +\xi_y^{2+\alpha}
                         [D_{U_0}^2\lambda]_{\alpha,y}^{U_0}\}.
 \end{split}
\end{equation}
The short coordinate transports are those
already fixed in the strong norm. A fixed
increase of the separation absorbs the
scale-ratio factors in the following sums.

\begin{proposition}[Strong inverse of the proper site gauge block]
\label{dsc:proper-site-gauge-inverse}
The actual projection and inverse satisfy
\begin{equation}\label{dsc:R-holder}
                   \|R_0\rho\|_{\alpha,2}\le C_\alpha\|\rho\|_{\alpha,2},
\end{equation}
and
\begin{equation}\label{dsc:proper-site-gauge-bound}
             \|G_sR_0\rho\|_{Y^{2,\alpha}}
                                      \le C_\alpha\|\rho\|_{\alpha,2}.
\end{equation}
In particular, for \(\rho\in\operatorname{ran}R_0\),
\(\lambda=G_s\rho\) is the unique site solution of
\begin{equation}\label{dsc:proper-site-gauge-identities}
                  q_0\lambda=0,\qquad R_0\Delta_{U_0}\lambda=\rho.
\end{equation}
The true bare-site zero condition is retained.
\end{proposition}

\begin{proof}
Use physical orthonormal top columns \(F_z\).
\Cref{prop:psr-site-columns} gives their
pointwise size \(C\xi_y^{-2}e^{-md(y,z)}\)
and their local H\"older size
\(C_\alpha\xi_y^{-2-\alpha}e^{-md(y,z)}\).
A source with unit \(\|\cdot\|_{\alpha,2}\)
norm has \(L^2\) norm at most \(C\) on
each pillar. Thus its pairing with \(F_z\)
is bounded by the uniformly summable
exponential cell sum. The normalized
Gram inverse preserves that coefficient
bound. Summing the output columns proves
\eqref{dsc:R-holder} for
\(R_0=1-FM^{-1}F^\sharp\).
For the H\"older bound use the column
increment before summing; no derivative
of the input is taken.

Set \(u=G_sR_0\rho\). The exact site
identity \(q_0G_sR_0=0\) gives
\[
                            q_0u=0,\qquad
                            \Delta_{U_0}u=R_0\rho.
\]
The second identity follows from
\(G_s^{-1}=\Delta_{U_0}+q_0^\sharp a_sq_0\).
Its exponential energy columns give
\[
 \xi_y^{-2}\|u\|_{2,B_y}
               +\xi_y^{-1}\|D_{U_0}u\|_{2,B_y^+}
                            \le C\|\rho\|_{\alpha,2}.
\]
For example the \(G_s\) block carries
\(\xi_y\xi_z e^{-md(y,z)}\), while the
source \(L^2\) cell norm is bounded.
The ratio \(\xi_z/\xi_y\) is absorbed
by the retained metric rate. The
derivative block is treated in the
same way. The projection \(R_0\)
has already been included in the
source bound.

Rescale a nonbare box by \(\xi_y\).
The bounded-force site bootstrap
\cref{lem:psr-local-bootstraps} gives
its \(C^{1,\alpha}\) norm, by choosing
a finite \(p>4/(1-\alpha)\).
Its source has the rescaled
\(C^\alpha\) bound from
\eqref{dsc:R-holder}. The local
reference potential has the fixed
\(C^{1,\alpha}\) bound
\eqref{dsc:regular-reference}.
Apply \eqref{dsc:covariant-interior}
on one further nested box. Rescaling
back gives every power in
\eqref{dsc:proper-site-gauge-bound}.
At a bare cell the physical atomic
energy inequality and a fixed
two-step collar prove those same
weighted bounds directly. In
particular \(q_0u=0\) forces the
bare site values to be zero.

For uniqueness in
\eqref{dsc:proper-site-gauge-identities},
suppose \(q_0\lambda=0\) and
\(R_0\Delta_{U_0}\lambda=0\).
Then \(\Delta_{U_0}\lambda\) belongs
to \(\operatorname{ran}F\), say
\(\Delta_{U_0}\lambda=F a\).
Since \(q_0\lambda=0\),
\(\lambda=G_sFa\). Applying \(q_0\)
gives \(q_0G_sFa=0\). But
\[
                   q_0G_sF=\Xi^2M
\]
by \(F=G_sq_0^\sharp\Xi^{-2}\)
and \(M=F^\sharp F\).
The normalized Gram is invertible,
so \(a=0\), then \(\lambda=0\).
This proves uniqueness with the
same actual boundary convention.
\end{proof}

The second-derivative site estimate uses
the physical current of the reference
to obtain its \(C^{1,\alpha}\) local
potential. It does not upgrade an
arbitrary curvature-only connection
to that class. Its purpose is the
gauge restoration near an actual
regular critical reference, while
the curvature-only bounded-force
field estimates remain as stated in
\cref{ch:proper-strong-response}.

% END LAYERED EXPANSION discrete_schauder_estimates.tex

% BEGIN LAYERED EXPANSION background/proper_nonlinear_critical_continuation.tex
\subsection{Nonlinear continuation of an actual proper critical reference}
\label{ch:pnc-proper-critical-continuation}

We construct a nearby constrained-critical orbit on the actual
determining region, starting from a given real physical critical
reference. The region, its outward rows and its bare boundary are
kept. The reference is not inferred from small curvature of a datum.
The input is an actual solution, and the conclusions are a uniform
analytic continuation, full stationarity, literal constraint and
local orbit uniqueness.

Use the proper geometry and norms of
\cref{ch:proper-strong-response}. Write $U_0$ for the real reference,
$\eta$ for the fine spacing, $\xi_c$ for a field-row length and
$\xi_y$ for a site-top length. Let $q_0$ be the actual own-pillar
site average, $Q_0$ the actual straightened field derivative and
$R_0$ the real site projection. The site Laplacian and inverse are
$\Delta_0=D_0^\sharp D_0$ and
$G_s=(\Delta_0+q_0^\sharp a_sq_0)^{-1}$.
Assume the actual bounds
\begin{equation}
 \sup_p\mathsf L(p)^2|F(U_0)(p)|\le\kappa_0,\qquad
 \sup_b\mathsf L(b)^3|j_0(b)|\le\gamma_0,\qquad
                         j_0=Q_0^\sharp\omega_0.
 \label{pnc:actual-reference}
\end{equation}
Here $j_0$ is the actual Wilson current. The last equality is full
physical stationarity in the declared straightened coordinates.
The induced-current theorem gives
$\sup_c\xi_c^3|\omega_0(c)|\le C\gamma_0$.
We take $\kappa_0,\gamma_0$ below fixed sufficiently small ceilings,
with the strict collar room required by the strong and Schauder
estimates. The separation is fixed large enough for their metric
growth and scale-ratio reserves.

For a fine increment use $\|A\|_{\mathrm{str},\alpha}$ of
\eqref{psr:strong-norm}, with its derivatives and short transports
measured in the fixed reference $U_0$. Set
\begin{equation}
 \|f\|_{3,\infty}=\sup_b\mathsf L(b)^3|f(b)|,\qquad
 \|b\|_{1,\infty}=\sup_c\xi_c|b(c)|,\qquad
 \|\nu\|_{3,\infty}=\sup_c\xi_c^3|\nu(c)|.
 \label{pnc:physical-amplitudes}
\end{equation}
The same symbol on fine forces and row multipliers denotes their
respective displayed index sets. The site source and gauge norms
are \eqref{dsc:site-source-norm} and \eqref{dsc:site-gauge-norm}.

Put $U_A=e^{i\eta A}U_0$ and
$F(A)=A_{\Omega_0}(U_A)$. Let $q(A)$ be the actual re-centered
straightened per-length constraint, so $q(0)=0$, $q'(0)=Q_0$.
We use the sign convention
$D_0\lambda=\eta^{-1}(\lambda_- -\operatorname{Ad}U_0\lambda_+)$;
its gauge curve is $g_t=e^{it\lambda}$. All complex transposes below
are physical bilinear transposes.

\begin{proposition}[Uniform proper critical continuation]
\label{prop:pnc-critical-continuation}
There are $r,\rho,C>0$, independent of depth, compatible volume
and the shapes of the separated regions, such that the following
holds for the actual reference \eqref{pnc:actual-reference}.
For every literal mixed datum
$V_B(c)=e^{i\xi_cB(c)}V_0(c)$ with
$\|B\|_{1,\infty}<\rho$, there is a holomorphic pair
$(A(B),\nu(B))$ satisfying
\begin{equation}
 \begin{split}
 q(A(B))&=b(B),\qquad R_0D_0^\sharp A(B)=0,\\
 F'(A(B))&=q'(A(B))^\sharp(\omega_0+\nu(B)),\\
 \|A(B)\|_{\mathrm{str},\alpha}
                  +\|\nu(B)\|_{3,\infty}&\le C\|B\|_{1,\infty}.
 \end{split}
 \label{pnc:continuation-equations}
\end{equation}
The reduced datum $b(B)$, defined next, includes every outward
factor. A specified admitted fine gauge gives a literal field
with determining mean exactly $V_B$. It is holomorphically
constrained critical, and real for real $B$.
It is the unique nearby critical orbit modulo fine gauges equal
to one at all assigned tops. Its actual curvature and current,
and its physical multipliers, obey the uniform continuation
bounds proved below. Localized analytic response holds in the
fixed determining geometry.
\end{proposition}

\subsubsection{The reduced datum at an outward endpoint}

Write $a_j(A)$ for the raw level-$j$ straightened coordinate and
$v_j[A]$ for its actual node factor. For each assigned top $y$
put $v[A](y)=v_{j_y}[A](y)$. At an outward end $z$ of a level-$j$
row, $z$ is assigned level $j-1$, and
\begin{equation}
 v_j[A](z)=v_{j-1}[A](z)\,u_{W_{j-1}}(a_{j-1}(A))(z),
                         \qquad W_{j-1}=M^{j-1}U_0.
 \label{pnc:outward-node-factor}
\end{equation}
The last factor is the actual one-step factor of
\eqref{hb:one-step-actual-map}. Its entire input block consists
of the prescribed level-$(j-1)$ interior rows. There is no
missing inward datum in this formula.

For a row $c$, let $u_-(q)$ or $u_+(q)$ be this factor at a lower
assigned outward endpoint, and one at a same-level endpoint.
Only one endpoint is lower. Removing the assigned node factor
from the actual plain-mean diagram gives exactly
\begin{equation}
 \mathcal T(q)_c
   =u_-(q)\{e^{i\xi_cq_c}V_0(c)\}u_+(q)^{-1},\qquad
                  M_{\mathbf B}(U_A)=\mathcal T(q(A))^{v[A]}.
 \label{pnc:mixed-datum-diagram}
\end{equation}
This is \eqref{mb:outward-node-identity} with its endpoint order
retained. Interior rows are unchanged. Each outward factor reads
only finer interior rows, which are themselves unchanged by
$\mathcal T$. Its inverse is therefore explicit:
\begin{equation}
 e^{i\xi_cb_c(B)}V_0(c)
   =u_-(B)^{-1}\{e^{i\xi_cB_c}V_0(c)\}u_+(B).
 \label{pnc:actual-reduced-datum}
\end{equation}
For an interior-to-outer orientation this says
$e^{i\xi_cb_c}=e^{i\xi_cB_c}\operatorname{Ad}(V_0(c))u_+(B)$,
as in the original outward diagram. Reversing the stored row
uses the inverse group equation and requires no new convention.

The one-step block contains $C_L$ coarse coordinates, has a
uniform relative analytic branch at the real reference and
reads adjacent scales with ratio $L$. Consequently
\begin{equation}
 b(0)=0,\qquad\|b(B)\|_{1,\infty}\le C_L\|B\|_{1,\infty},
 \qquad \|D^mb(B)\|\le C_{L,m}
 \label{pnc:reduced-datum-bounds}
\end{equation}
on a fixed smaller raw-angle ball. These derivative bounds use
the displayed amplitude norms and bounded complete carriers.
The same holds for the inverse datum chart. There is no product
over levels: the factors depend on interior rows, never on an
already corrected outward row. In particular solving $q(A)=B$
on outward rows without this reduction would solve a different
literal problem.

\subsubsection{A quadratic residual for the actual Wilson equation}

The actual proper reference has local Coulomb potentials
$a_0$ with uniformly small $C^{1,\alpha}$ norm after rescaling
each nonbare box by its assigned length. This uses both the
physical curvature and the actual current in
\eqref{pnc:actual-reference}; it is not a curvature-only
$C^{1,\alpha}$ assertion.

\begin{lemma}[The complete re-centered Wilson residual]
\label{lem:pnc-Wilson-residual}
Let $J(A)=F'(A)$ denote the physical normal-coordinate gradient.
Then
\begin{equation}
 N_F(A)=J(A)-J(0)-F''(0)A,\qquad
 \|N_F(A)\|_{3,\infty}\le C\|A\|_{\mathrm{str},\alpha}^2.
 \label{pnc:Wilson-residual-bound}
\end{equation}
It is uniformly holomorphic on a fixed strong ball and has
Lipschitz constant $Cr$ there from the strong norm to the
weighted force norm.
\end{lemma}
\begin{proof}
On a nonbare box scale its length to one and set
$h=\eta/\xi$. In its Coulomb frame the actual product has
potential
\[
        P_h(A)=\frac1{ih}\log(e^{ihA}e^{iha_0}).
\]
The actual gradient in this frame is
$D_AP_h(A)^\sharp J_h(P_h(A))$, with $J_h$ the flat-coordinate
Wilson gradient. The exact current calculation of
\cref{lem:hb-Wilson-residual} gives
$J_h(a)=\mathrm{d}^\sharp \mathrm{d}a+N_h(a)$, with uniformly analytic
$N_h:C^{1,\alpha}\to L^\infty$.
Use cutoffs on one larger box if necessary; the current on the
inner full stencil is unchanged.

Here are the divided-difference bounds needed for the coordinate
composition. Put $P_0=D_AP_h(0)$ and
$E_h(A)=P_h(A)-a_0-P_0A$. Matrix multiplication and its logarithm
are analytic functions of $ha_0,hA$, and division by $h$ removes
their zero of first order. Taylor's formula gives
\begin{equation}
 \begin{split}
 |E_h(A)|&\le Ch|A|^2,\qquad \|DE_h(A)\|_\infty\le Chr^2,\\
 \|D^2E_h(A)\|_\infty&\le Cr^2,\qquad
 |D_AP_h(A)-P_0|\le Chr,\\
 |D_AP_h(A)-P_0-D_A^2P_h(0)[A]|&\le Ch^2r^2.
 \end{split}
 \label{pnc:BCH-residual-scales}
\end{equation}
Here $D$ denotes ordinary spatial differences and the parameter
derivatives are written $D_A$; $r$ bounds the rescaled strong
norm of $A$. To check the only delicate second-difference line,
each leading term is $hA D^2A$ or $h(DA)^2$.
Differentiating a coefficient with respect to $a_0$ supplies
another $h$, so the remaining terms contain
$h^2A^2D^2a_0$, $h^2A DA Da_0$ or higher such powers.
The elementary bounds $\|D^2A\|_\infty\le Ch^{-1}\|DA\|_\infty$
and $\|D^2a_0\|_\infty\le Ch^{-1}\|Da_0\|_\infty$
cancel every apparent mesh loss. The first differences of
$E_h$ and its bounded second differences also bound its
$C^{1,\alpha}$ norm by $Cr^2$.

Now Taylor-expand
$D_AP_h(A)^\sharp J_h(P_h(A))$ at $A=0$.
The constant and linear terms are exactly $J(0)$ and $F''(0)A$.
The flat second-order term acting on $E_h$ is bounded by
\eqref{pnc:BCH-residual-scales}. The term
$(D_AP_h(A)-P_0)^\sharp \mathrm{d}^\sharp \mathrm{d}(P_0A)$ is also bounded:
its first factor is $O(hr)$, while the second is at most
$Ch^{-1}r$. Spatial differences of $P_0$ have an additional
$h$ for each derivative of its dependence on $a_0$; the same
rough second-difference estimate covers them. The quadratic
remainder of $D_AP_h$ times $J_h(a_0)$ has bound $Ch^2r^2$
times $Ch^{-1}\|Da_0\|_\infty+C$, and is bounded by $Cr^2$.
All terms containing $N_h$ are bounded by its uniform Banach
derivative estimates and the $C^{1,\alpha}$ coordinate bounds.
This proves the complete gradient residual, rather than merely
an action remainder.

At a bare cell $\xi=\eta$, its rescaled mesh is one. There are
only a fixed number of incident inner plaquettes. Their exact
finite matrix gradient has uniformly bounded analytic derivatives
at a unitary reference. Direct Taylor subtraction therefore gives
the same quadratic estimate, including a rough action boundary;
no exterior Coulomb frame is used.

Rescaling back multiplies the gradient by $\xi^{-3}$.
The complete carriers have comparable assigned lengths, proving
\eqref{pnc:Wilson-residual-bound}. The same estimates hold for
complex $A$ on a fixed larger ball. Banach-valued Cauchy estimates
and $N_F(0)=DN_F(0)=0$ give its $Cr$ Lipschitz bound.
No second derivative of $A$ or of $a_0$ was assumed.
\end{proof}

\subsubsection{Actual constraint coefficients in the physical norms}

Put $C_q(A)=q(A)-Q_0A$. On each determining row the actual
two-leg coefficient bound is
\begin{equation}
 \sup_e\sum_{e'}\|\partial_e\partial_{e'}q_c(A)\|
                           \le C_L\xi_c(\eta/\xi_c)^4.
 \label{pnc:actual-two-leg-bound}
\end{equation}
It is \eqref{px:constraint-hessian-leg-moment} on the admitted
complex increment ball about this actual regular reference.
All fine input scales on the complete carrier are comparable
to $\xi_c$ within a fixed $L$ factor. The physical transpose
uses $\eta^{-4}\xi_c^4$. Thus, for joint amplitude at most $r$,
\begin{equation}
 \begin{split}
 \|C_q(A)\|_{1,\infty}&\le Cr^2,\\
 \|(q'(A)-Q_0)^\sharp\nu\|_{3,\infty}&\le Cr\|\nu\|_{3,\infty},\\
 \|(q'(A)-Q_0-q''(0)[A])^\sharp\omega_0\|_{3,\infty}
                                      &\le C\|\omega_0\|_{3,\infty}r^2.
 \end{split}
 \label{pnc:constraint-residual-bounds}
\end{equation}
For the first line sum the fixed-leg bound over the
$O((\xi_c/\eta)^4)$ legs and insert $|A|\le Cr/\xi_c$.
For the second line its column density is
$Cr(\eta/\xi_c)^4$, which cancels the physical transpose factors.
For the last line apply Cauchy's estimate to the actual Hessian
in its fixed-leg sum norm on the larger weighted increment ball,
then integrate twice. Its column remainder is
$Cr^2(\eta/\xi_c)^4$. This is not an inference of an absolute
coefficient sum from an arbitrary second directional bound.
There are only a bounded number of rows reading a fixed fine leg.
All three maps have the corresponding quadratic Lipschitz bounds.

\subsubsection{The joint proper fixed point}

Let
$H_{L,0}=F''(0)-\langle\omega_0,q''(0)\rangle$.
The actual physical reference and the strong row-saddle theorem
give an inverse $\mathcal L_0^{-1}$ of
\begin{equation}
 (H_{L,0}+D_0R_0D_0^\sharp)A-Q_0^\sharp\nu=f,\qquad Q_0A=b,
 \qquad
 \|A\|_{\mathrm{str},\alpha}+\|\nu\|_{3,\infty}
                  \le C(\|f\|_{3,\infty}+\|b\|_{1,\infty}).
 \label{pnc:base-joint-inverse}
\end{equation}
Choose a finite $p>4/(1-\alpha)$ in
\eqref{psr:strong-row-saddle}; its force norm is bounded by the
displayed supremum norm. Removing the artificial row mass changes
its multiplier by $-a_fb$, which has exactly the displayed
amplitude. No moving-slice inverse is being presumed here.

The actual penalized equations are
\begin{equation}
 F'(A)+D_0R_0D_0^\sharp A-q'(A)^\sharp(\omega_0+\nu)=0,
                                  \qquad q(A)=b(B).
 \label{pnc:actual-penalized-system}
\end{equation}
After subtracting the base equation they are precisely
\begin{equation}
 \begin{split}
 (A,\nu)=\mathcal L_0^{-1}\bigl(&-N_F(A)
 +(q'(A)-Q_0-q''(0)[A])^\sharp\omega_0\\
 &+(q'(A)-Q_0)^\sharp\nu,\quad b(B)-C_q(A)\bigr).
 \end{split}
 \label{pnc:actual-fixed-point}
\end{equation}
The residual estimates give quadratic size $Cr^2$ and Lipschitz
constant $Cr$ in the joint strong/multiplier ball. Choose $r$
below all primitive ceilings with $Cr<1/4$, and then choose
$\rho$ so $C\|b(B)\|_{1,\infty}<r/4$. The radius-$r/2$ ball
is invariant and contractive. Iteration constructs a unique
holomorphic pair, real for real data, with
\begin{equation}
 \|A(B)\|_{\mathrm{str},\alpha}+\|\nu(B)\|_{3,\infty}
                                    \le C\|B\|_{1,\infty}.
 \label{pnc:fixed-point-size}
\end{equation}
It remains to remove the penalty and identify the literal fibre.

\subsubsection{Strong site extensions in the actual reference}

\begin{lemma}[Average and node extensions]
\label{lem:pnc-site-extensions}
There are real linear maps $E_q,E_n$ from assigned top values to
site functions such that
\begin{equation}
 q_0E_q=1,\qquad \operatorname{ev}_{\mathrm{tops}}E_n=1,\qquad
 \|E_qc\|_{Y^{2,\alpha}}+\|E_nc\|_{Y^{2,\alpha}}
                                                   \le C\|c\|_\infty.
 \label{pnc:site-extension-bounds}
\end{equation}
Top data zero on the bare tops give site functions zero on the
bare layer. Their supports and fine derivative collars have bounded
overlap and bounded layer-metric radius.
\end{lemma}
\begin{proof}
In a nonbare pillar use a local Coulomb frame of $U_0$ on its
fixed complete surrounding box. The rescaled reference potential
has its $C^{1,\alpha}$ bound from the actual current. Choose a
smooth scalar bump supported strictly inside the pillar and
positive at its center. Its sampled scalar average $m_y$ is
bounded above and below: fine meshes tend to its positive integral,
and the finitely many coarse possibilities contain the center.
The actual transported average of this bump times a constant
colour is a three-by-three matrix $N_y$ with
$\|N_y-m_yI\|\le C\kappa_0m_y$. The actual hierarchical
transports have this bound by their geometric sum of physical
lengths in the small local potential. Normalize by $N_y^{-1}$,
including the unitary top-frame rotation, to obtain $E_q$.
Normalize the bump's center value instead to obtain $E_n$.

Each covariant first derivative is a profile derivative plus the
local reference potential times the profile. Two derivatives
contain only $D^2\phi$, $a_0D\phi$, $Da_0\phi$ and $a_0^2\phi$,
with the indicated rescaled bounds. Their $\alpha$ differences
therefore use $Da_0\in C^\alpha$, not $D^2a_0$.
This proves \eqref{pnc:site-extension-bounds} on the fine meshes.
For the finitely many smaller ratios $\xi_y/\eta$, direct center
seeds have bounded weighted differences of every fixed order.
At a bare top use its one site. Adjacent assigned scales differ
by at most $L$, so its fine collars cost only a fixed $L$ factor.
The supports are disjoint inside pillars; their derivative
collars have bounded overlap. These estimates are covariant
under the real local frame changes and hence give the global
norms stated above. They also prove the assertion about bare
zeros and carriers.

If a required surrounding box does not embed in a small periodic
volume, pass first to a finite periodic cover where it does.
Construct the two extension operators there, and average each
operator over simultaneous fine and top deck shifts. The actual
reference, its averages and node evaluation commute with these
shifts, so both right-inverse identities are retained. The local
strong norms are convex and invariant under the shifts; the
averaged operator has the same bound. It descends on invariant
data with the same supremum and H\"older bounds and with the
same carriers modulo the torus. No sheet count enters these
norms. This treats small volumes without a nonembedded frame
or a limit in the volume.
\end{proof}

The site block needed for a finite gauge is now an actual bounded
inverse:
\begin{equation}
 \lambda\longmapsto(q_0\lambda,R_0\Delta_0\lambda),\qquad
 \lambda=E_qc+G_s\{\rho-R_0\Delta_0E_qc\}.
 \label{pnc:site-block-inverse}
\end{equation}
The second formula uses \eqref{dsc:proper-site-gauge-bound}; its
braces are in $\operatorname{ran}R_0$ and have bounded
$\|\cdot\|_{\alpha,2}$ norm. The exact identities
$q_0G_sR_0=0$, $\Delta_0G_sR_0=R_0$ prove both equations.
Uniqueness is \eqref{dsc:proper-site-gauge-identities}.
If $c$ is zero at bare tops, this inverse retains bare site zeros.

\subsubsection{The actual assigned effective gauge and its outward extension}

For a fine gauge $g$ define at assigned tops
\begin{equation}
 h_A(g)=v[A^g]^{-1}g|_{\mathrm{tops}}v[A],\qquad
 h_A(e^{it\lambda})=1+it\tau_A\lambda+O(t^2).
 \label{pnc:assigned-effective-gauge}
\end{equation}
Apply plain gauge covariance to \eqref{pnc:mixed-datum-diagram}.
It gives the exact nonlinear identity
\begin{equation}
       \mathcal T(q(A^g))=\mathcal T(q(A))^{h_A(g)}.
 \label{pnc:actual-proper-gauge-cocycle}
\end{equation}
This identity fixes the coarse action being differentiated; it
does not replace an outward endpoint by a chosen top value.

For completeness the one-step derivative can be computed directly.
Put $T_{\mathrm{cur}}T_{\mathrm{ref}}^{-1}=e^{iz_x}$,
$m=L^{-4}\sum_xz_x$ and $u=e^{im}$, and let
$\mathcal E_z w=\int_0^1\operatorname{Ad}(e^{itz})w\,dt$.
Then
\begin{equation}
 \tau^{(1)}_{W,a}\lambda(y)
 =\operatorname{Ad}(u^{-1})\left[\lambda(y)
 +\mathcal E_mL^{-4}\sum_x\mathcal E_{z_x}^{-1}
       \{-\lambda(y)+\operatorname{Ad}(T_{\mathrm{cur}})\lambda(x)\}
                                                     \right].
 \label{pnc:one-step-effective-site-map}
\end{equation}
The left derivative of the relative transport is
$i(\lambda(y)-\operatorname{Ad}(T_{\mathrm{cur}})\lambda(x))$;
differentiating $u(A^g)^{-1}g(y)u(A)$ gives the displayed sign.
At $a=0$ all $z_x,m$ are zero and this is exactly the real
transported site average. One-step analytic differentiation gives
an error at most $C_L\|a\|_\infty$ in the site supremum norm.
The actual raw radii on a top's hierarchy are bounded by
$C_Lr\xi_s/\xi_y$, with summable ratios. Telescoping the
products, whose zero-increment factors are unitary averages,
therefore proves
\begin{equation}
             \|\tau_A-q_0\|_{\infty\to\infty}\le Cr,
                        \qquad\tau_0=q_0.
 \label{pnc:effective-site-comparison}
\end{equation}
At a bare top $\tau_A\lambda=\lambda$ exactly.

At an outward endpoint, the higher-level effective parameter is
\eqref{pnc:one-step-effective-site-map} applied to all assigned
finer parameters on its complete outer block. At zero this is
precisely
$L^{-4}\sum_w\operatorname{Ad}(T_{z,w}(W_{j-1}))\lambda(w)$,
the extension in \eqref{px:outward-site-extension}.
At nonzero $A$ the evaluation terms and the logarithmic
differentials in \eqref{pnc:one-step-effective-site-map} must
remain. In particular it is not an independently chosen site
average at the new field.

Let $\Gamma_A$ be the normal-coordinate fine gauge generator.
Differentiating \eqref{pnc:actual-proper-gauge-cocycle} gives
\begin{equation}
 q'(A)\Gamma_A\lambda=\mathcal D_A\tau_A\lambda,
 \qquad
\mathcal D_A=\bigl[D\mathcal T(q(A))\bigr]^{-1}
                         d_{\mathcal T(q(A))}^{\mathrm{lit}}.
 \label{pnc:nonlinear-proper-intertwining}
\end{equation}
The derivative is left trivialized in the declared per-length
group coordinates, as is the literal gauge gradient. At zero,
$\mathcal D_0=\bar D_{\mathrm{ext}}$ by the preceding explicit
outer-block derivative and \eqref{px:typed-seam-intertwining}.
Most importantly, $\tau_A\lambda=0$ makes the left side zero
on every interior and outward row.

\subsubsection{Surjective gauge tests and removal of the penalty}

The exact fine generator is
\begin{equation}
                 \Gamma_A\lambda
                    =\mathcal E_{\eta A}^{-1}D_{U_A}\lambda.
 \label{pnc:fine-gauge-generator}
\end{equation}
Covariant product differences give
\begin{equation}
 \|(\Gamma_A-D_0)\lambda\|_{\mathrm{str},\alpha}
                    \le Cr\|\lambda\|_{Y^{2,\alpha}}.
 \label{pnc:generator-comparison}
\end{equation}
Terms have the forms $A\lambda$, $\eta A D_0\lambda$ and their
first covariant differences. The second site derivative and its
$\alpha$ seminorm in $Y^{2,\alpha}$ control exactly those
differences. On a carrier $\eta\le\xi$ and the physical powers
cancel. This argument uses the fixed reference transports, not
ordinary differences of an uncontrolled frame change.

For $\rho\in\operatorname{ran}R_0$ with finite
$\|\rho\|_{\alpha,2}$ put $\lambda_0=G_s\rho$ and
\begin{equation}
 \lambda_A(\rho)=\lambda_0
                   -E_q(\tau_AE_q)^{-1}\tau_A\lambda_0.
 \label{pnc:kernel-corrected-site-lift}
\end{equation}
The inverse exists since $\tau_AE_q=1+O(r)$.
Then $\tau_A\lambda_A=0$ and
$\|\lambda_A-\lambda_0\|_{Y^{2,\alpha}}
 \le Cr\|\rho\|_{\alpha,2}$.
Bare zeros are exact: $\lambda_0$ is zero at bare sites;
$\tau_A$ there is evaluation; and the bare-output block of
$\tau_AE_q$ is the identity with no nonbare input. Thus the
correction's bare coefficients and site values are zero.

The boundedness of $R_0$ on the site source space, its range
identities, and \eqref{pnc:generator-comparison} give
\begin{equation}
 \mathscr R_A\rho=R_0D_0^\sharp\Gamma_A\lambda_A(\rho),\qquad
                \|\mathscr R_A-1\|_{\alpha,2\to\alpha,2}\le Cr
              \quad\hbox{on }\operatorname{ran}R_0.
 \label{pnc:gauge-test-surjectivity}
\end{equation}
A smaller $r$ makes this invertible. For every site source
$\sigma\in\operatorname{ran}R_0$ there is consequently a gauge
parameter with $\tau_A\lambda=0$ and
$R_0D_0^\sharp\Gamma_A\lambda=\sigma$.

For the fixed point, $\rho=R_0D_0^\sharp A$ belongs to this
source space, with norm at most $C\|A\|_{\mathrm{str},\alpha}$.
Test \eqref{pnc:actual-penalized-system} against these actual gauge
vectors. Gauge invariance of the inner-plaquette action kills
$F'(A)$; \eqref{pnc:nonlinear-proper-intertwining} kills the
constraint term. Thus
\begin{equation}
                   \langle\rho,\sigma\rangle_{\mathrm{bil},\eta}=0
                      \quad(\sigma\in\operatorname{ran}R_0).
 \label{pnc:penalty-annihilator}
\end{equation}
This forces $\rho=0$ also at complex data. Indeed $R_0$ is the
real reference projection in Pauli coordinates, so the
coefficientwise conjugate $\overline\rho$ remains in its range
and source space. Taking $\sigma=\overline\rho$ gives its
positive physical squared norm. We have not used the false
implication that a complex vector with zero bilinear self-pairing
must vanish. The full equations
\eqref{pnc:continuation-equations} are now proved.

\subsubsection{An admitted literal representative and the local mean slice}

For assigned node values $w=e^{i\zeta}$ define the actual fine
extension $\mathfrak E(w)=e^{iE_n\zeta}$. Its node values are
exactly $w$ and its logarithm has the stated $Y^{2,\alpha}$
bound. Set
\begin{equation}
                    \widehat U(B)=U_{A(B)}^{\mathfrak E(v[A(B)]^{-1})}.
 \label{pnc:literal-critical-representative}
\end{equation}
The datum diagram and \eqref{pnc:actual-reduced-datum} give
$M_{\mathbf B}\widehat U(B)=V_B$ exactly.

The finite gauge map has the uniform estimate
\begin{equation}
 \|A^{e^{i\lambda}}-A-D_0\lambda\|_{\mathrm{str},\alpha}
       \le C(\|A\|_{\mathrm{str},\alpha}
                         +\|\lambda\|_{Y^{2,\alpha}})\|\lambda\|_{Y^{2,\alpha}}.
 \label{pnc:finite-gauge-bound}
\end{equation}
To see the mesh cancellation, write
$\operatorname{Ad}(U_0)\lambda_+=\lambda_- -\eta D_0\lambda$
inside the exact three exponential factors defining $A^g$.
At $\eta=0$ their product is one. Dividing its logarithm by
$\eta$ is analytic in $\lambda,A,D_0\lambda$, uniformly down
to zero mesh. Covariant product rules use only the site
$Y^{2,\alpha}$ norm for its first field derivative and its
H\"older increment. This proves the bound and its local
Lipschitz version. Thus the specified literal gauge leaves
the increment in a fixed admitted strong ball after one
decrease of the continuation radius. No general complex
gauge saturation of a domain is assumed.

For an arbitrary nearby literal field define the assigned
mean-slice map
\begin{equation}
 \lambda\longmapsto
  \bigl(i^{-1}\log h_A(e^{i\lambda}),\ R_0D_0^\sharp A^{e^{i\lambda}}\bigr).
 \label{pnc:actual-mean-slice-map}
\end{equation}
Its derivative at zero reference and zero parameter is
$(q_0\lambda,R_0\Delta_0\lambda)$, whose actual inverse is
\eqref{pnc:site-block-inverse}. The effective site comparison
and finite gauge bound make its derivative a small perturbation
on the same fixed spaces. The explicit contraction with that
linear inverse solves
\begin{equation}
 h_A(e^{i\lambda})=v[A],\qquad R_0D_0^\sharp A^{e^{i\lambda}}=0,
                 \qquad\|\lambda\|_{Y^{2,\alpha}}
                                     \le C\|A\|_{\mathrm{str},\alpha}.
 \label{pnc:actual-mean-slice}
\end{equation}
At bare tops its target is one and its first coordinate is
exactly $\lambda$, so bare zeros are retained during this
contraction. The resulting potential $\overline A$ has
$q(\overline A)=b(B)$ if the original literal datum is $V_B$,
by \eqref{pnc:actual-proper-gauge-cocycle}. Its node factor
equals the node values of the chosen gauge. The subsequent
literal extension therefore differs from the original field
by a gauge exactly one at every assigned top.

No tangent of the full $q$ fibre is lost by imposing Landau
gauge: \eqref{pnc:gauge-test-surjectivity} corrects its Landau
component by a gauge vector in $\ker\tau_A$, and hence in
$\ker q'(A)$. The action is constant along that vector.
These two analytic maps consequently identify all tangent directions of
the local $q=b(B)$ and literal fibres modulo that residual
gauge group, and preserve the action. The inverse of the
specified literal extension solves the mean-slice equations
at a mean-slice point; uniqueness of their contraction proves
this assertion, including the full group values. Differentiating
the action identities proves physical constrained criticality
of \eqref{pnc:literal-critical-representative} at complex as
well as real data.

\subsubsection{Competing multipliers and uniqueness of the critical orbit}

Smallness of a competing field does not by itself assert a
small multiplier. Here the necessary readout is available from
the actual row lift. Let $H_0$ be the zero-force physical row
lift at $U_0$, so $Q_0H_0=1$. The normalized energy columns in
\eqref{scr:lift-energy-columns} and their proper strong upgrade
in \eqref{psr:strong-row-saddle} give a uniformly bounded
right inverse from the row amplitude norm to the strong norm.
The actual two-leg density estimate, paired with those columns,
also gives, for some $s>0$,
\begin{equation}
 T_A=q'(A)H_0=1+E_A^{\mathrm{row}},\qquad
 \|E_A^{\mathrm{row}}\|_{\mathrm{Schur},s}\le Cr.
 \label{pnc:perturbed-row-readout}
\end{equation}
The matrix here is normalized by $b_c=\xi_c^{-1}\beta_c$;
the Schur norm is the maximum of its absolute weighted row
and column sums in $d(c,c')$, the layer distance between
completed carriers. To verify the estimate, fix a fine leg
in the variation of $q'$ and use
\eqref{pnc:actual-two-leg-bound}. The factor
$(\eta/\xi_c)^4$ sums its fine carrier and the remaining
increment gives $r$. The lift column has exponential cell
energy decay. Carrier enlargement and its finitely many
neighboring scales change this exponential by a fixed factor.
The geometric growth and scale-ratio reserves leave a summable
exponential after taking both normalized sums. Thus this is
an absolute matrix estimate, not merely an unweighted operator
bound. A Neumann series in this Schur algebra proves that
\begin{equation}
                 \mathscr H_A=H_0T_A^{-1},\qquad
                         q'(A)\mathscr H_A=1
 \label{pnc:actual-row-right-inverse}
\end{equation}
has the same exponentially decaying energy columns.

Suppose now that a competing mean-slice point $A$ is critical
on $q=b(B)$. Surjectivity just proved supplies its unique
multiplier $\omega$. Write $\nu=\omega-\omega_0$.
For the unit energy datum $b_{c,e}=\xi_c^{-1}e$, the physical
coarse pairing gives the exact identity
\begin{equation}
 \xi_c^3\langle\nu(c),e\rangle
 =\left\langle F'(A)-F'(0)
             -(q'(A)-Q_0)^\sharp\omega_0,
                                      \mathscr H_A b_{c,e}\right\rangle.
 \label{pnc:competing-multiplier-readout}
\end{equation}
This identity does not estimate $F'(A)$ in the fine supremum
norm. Instead integrate the actual Wilson Hessian along $tA$
and use its weak plaquette form. On a rescaled complete cell
of length $\xi$, its absolute pairing is bounded by
\begin{equation}
 C\left\{\|D_0A\|_\infty\|D_0w\|_1
 +\xi^{-1}\bigl(\|A\|_\infty\|D_0w\|_1
                   +\|D_0A\|_\infty\|w\|_1\bigr)
 +\xi^{-2}\|A\|_\infty\|w\|_1\right\}.
 \label{pnc:Wilson-weak-readout}
\end{equation}
Indeed differentiating a plaquette twice gives the two
first covariant differences, and terms without such a
difference carry its small plaquette curvature. The normal
coordinate correction contains $\eta[A,w]$ in the first
variation. Sum its difference by parts onto this product;
$\eta\le\xi$ bounds the resulting terms by the same display.
This calculation uses only first differences of $A,w$ and
the actual curvature along $tA$, never $D^2A$ or an estimate
of its current. It is also valid on a bare cell, where the
rescaled plaquette is a fixed finite matrix product.

For $w=\mathscr H_Ab_{c,e}$, the column energy bound and
Cauchy--Schwarz on a physical cell give
$\|w\|_1\le C\xi^3e^{-s d}$ and
$\|D_0w\|_1\le C\xi^2e^{-s d}$.
Every term in \eqref{pnc:Wilson-weak-readout} is therefore
$C\|A\|_{\mathrm{str},\alpha}e^{-s d}$.
The other term in \eqref{pnc:competing-multiplier-readout}
has force bound $C\gamma_0\|A\|_{\mathrm{str},\alpha}\xi^{-3}$
and the same summable column estimate. Summing proves
\begin{equation}
               \|\omega-\omega_0\|_{3,\infty}
                              \le C\|A\|_{\mathrm{str},\alpha}.
 \label{pnc:competing-multiplier-bound}
\end{equation}
This also holds for complex competing points, since all
pairings used to obtain the identity are bilinear and the
estimates use their absolute values.

Choose the fixed orbit neighborhood small enough that its
mean-slice field and \eqref{pnc:competing-multiplier-bound}
lie in the contraction ball. Both solve the complete
penalized system with zero penalty, so the fixed-point
uniqueness identifies them. The literal-to-mean-slice
construction then identifies the two literal representatives
by a gauge exactly one at all assigned tops. This proves
the claimed local orbit uniqueness, without assuming a
multiplier bound for a competing critical point.

\subsubsection{Physical continuation bounds and localized response}

The normal and left gradients are related by
$F'(A)=\mathcal E_{\eta A}^{\sharp}j(U_A)$.
The exact stationarity equation and the fixed-leg density of
$q'$ give
\begin{equation}
 \|j(U_{A(B)})-j_0\|_{3,\infty}
            +\|F(U_{A(B)})-F(U_0)\|_{2,\infty}
                                     \le C\|B\|_{1,\infty}.
 \label{pnc:physical-fine-bounds}
\end{equation}
Here the curvature norm is the assigned plaquette norm with
weight $\mathsf L^2$. For the current, subtract
$Q_0^\sharp\omega_0$ from
$q'(A)^\sharp(\omega_0+\nu)$, use
\eqref{pnc:constraint-residual-bounds}, and multiply by
$\mathcal E_{\eta A}^{-\sharp}=1+O(\eta A)$.
For curvature, differentiate the four exact link factors;
their first differences give $D_0A$ and their commutators
give $A^2$, with weights $\mathsf L^{-2}$. The bare
calculation is the same finite product with $\mathsf L=\eta$.
The specified literal gauge has supremum distance $C\|B\|$
from one. Its adjoint action changes the reference curvature
and current by at most $C\|B\|\kappa_0$ and
$C\|B\|\gamma_0$. Thus \eqref{pnc:physical-fine-bounds}
holds for $\widehat U(B)$ as well. In particular its actual
curvature and current ceilings increase by at most
$C\|B\|_{1,\infty}$.

There are two different coarse multipliers. The one in
\eqref{pnc:continuation-equations} belongs to reduced
straightened coordinates. The derivative of the literal
critical action $S(B)=F(A(B))$ is instead
\begin{equation}
 D_BS=Db(B)^\sharp(\omega_0+\nu(B))=:\mu_B,
             \qquad \mu_{\mathrm{left}}(B)
                          =\mathcal E_{\xi B}^{-\sharp}\mu_B.
 \label{pnc:literal-physical-multiplier}
\end{equation}
The inverse datum chart has bounded normalized row and
column sums on its adjacent complete carriers. Physical
transpose therefore maps amplitude three to amplitude
three: normalized dual variables are $\xi b$ and
$\xi^3\omega$. Consequently
\begin{equation}
 \|\mu_B-Db(0)^\sharp\omega_0\|_{3,\infty}
 +\|\mu_{\mathrm{left}}(B)-Db(0)^\sharp\omega_0\|_{3,\infty}
                                         \le C\|B\|_{1,\infty}.
 \label{pnc:literal-multiplier-bound}
\end{equation}
This formula retains the outward datum derivative at zero;
it need not be the identity. The local fibre coordinates
and stationarity proved above show that this is the actual
literal physical multiplier, not only a formal action
derivative.

For clarity, the localization claim can be stated without
assuming an all-order absolute tree sum. Give each completed
cell a weight $w=e^{s\psi}$, where $\psi$ is one-Lipschitz
in the proper layer metric. Use $w$ times each local norm,
and the same weight at the associated row. The base
exponential saddle columns, followed by the local finite-$p$
bootstrap, give \eqref{pnc:base-joint-inverse} in these
weighted spaces for a fixed small $s>0$. Enlarging a carrier
changes $w$ by at most a fixed factor. The residual's
derivative consequently has weighted operator norm $Cr$.
No derivatives of the artificial weight are needed: this
is a cellwise norm estimate with bounded carrier comparisons.
The differentiated fixed point therefore gives
\begin{equation}
 \|DA(B)[h]\|_{\mathrm{str},\alpha;w}
       +\|D\nu(B)[h]\|_{3,\infty;w}
                            \le C\|h\|_{1,\infty;w}.
 \label{pnc:weighted-first-response}
\end{equation}
The same bound holds at every point of a smaller fixed
complex datum ball. Banach Cauchy estimates there give
one weighted input slot and any fixed number $m-1$ of
unweighted slots, with constant $C_m$ independent of
depth and volume. The explicit datum reduction and literal
extension have bounded carriers and the same estimates.
Choosing $\psi$ to be distance from any selected source
carrier proves, for point data $c_1,\ldots,c_m$, normalized
response at an output cell $y$ bounded by
\begin{equation}
                         C_m e^{-s\max_i d(y,c_i)}.
 \label{pnc:localized-analytic-jets}
\end{equation}
This establishes localized holomorphic response in the
fixed proper geometry. It asserts neither an all-arity
summed connected kernel bound nor existence of a proper
reference from arbitrary curvature-only coarse data.
All inverse constants, branch choices, and admitted
complex increments used above are uniform around the
given real physical reference. This completes the proof
of \cref{prop:pnc-critical-continuation}.

% END LAYERED EXPANSION background/proper_nonlinear_critical_continuation.tex

% BEGIN LAYERED EXPANSION background/one_step_nonflat_background.tex
\subsection{Critical correction of a regular trial on the actual proper region}
\label{ch:rtc-regular-trial-correction}

The reference in \cref{prop:pnc-critical-continuation} can be
a regular trial rather than a critical point. Its stationarity
defect must then remain in the equation. This gives actual
critical fields on the proper image of sufficiently regular
fine trials. It does not supply a fine trial for every proposed
curvature-only mixed datum. The fixed one-step face lift is an
explicit example; its existence theorem was already proved in
\cref{ch:first-step-background-expanded}.

Retain the determining geometry, actual inner action, complete
carriers, true bare boundary, physical pairings and strong norms
of \cref{ch:pnc-proper-critical-continuation}. The reference
$U_0$ is real, with its actual curvature and Wilson current
satisfying
\begin{equation}
 \|F(U_0)\|_{2,\infty}\le\kappa_*,\qquad
                 \|j_0\|_{3,\infty}\le\gamma_*.
 \label{rtc:regular-trial}
\end{equation}
The fixed ceilings have strict room inside the bounds for the
current-dependent local $C^{1,\alpha}$ regularity and the site
Schauder estimates. An independent real row multiplier
$\omega_0$ may be chosen with
$\|\omega_0\|_{3,\infty}\le\gamma_{\omega,*}$, the precision
ceiling. Define the actual defect
\begin{equation}
 e_0=j_0-Q_0^\sharp\omega_0,\qquad
                            \delta=\|e_0\|_{3,\infty}.
 \label{rtc:actual-defect}
\end{equation}
There is no stationarity assumption. In particular
$\omega_0=0$ is allowed and gives $\delta=\|j_0\|_{3,\infty}$.

\begin{proposition}[Correction of an actual regular trial]
\label{prop:rtc-nonstationary-trial}
There are positive $\delta_*,\rho,r,C$, independent of depth,
compatible volume and the separated geometry, such that a
trial satisfying \eqref{rtc:regular-trial} with $\delta<\delta_*$ has the
following property. Write $V_0=M_{\mathbf B}(U_0)$ and let
$V_B(c)=e^{i\xi_cB(c)}V_0(c)$, with
$\|B\|_{1,\infty}<\rho$. There is a holomorphic pair
$(A(B),\nu(B))$, real at real $B$, satisfying
\begin{equation}
 \begin{split}
 q(A(B))&=b(B),\qquad R_0D_0^\sharp A(B)=0,\\
 F'(A(B))&=q'(A(B))^\sharp(\omega_0+\nu(B)),\\
 \|A(B)\|_{\mathrm{str},\alpha}
                  +\|\nu(B)\|_{3,\infty}
                              &\le C(\delta+\|B\|_{1,\infty}).
 \end{split}
 \label{rtc:corrected-critical-system}
\end{equation}
Here $F(A)=A_{\Omega_0}(e^{i\eta A}U_0)$ and $b(B)$ is the
actual outward-corrected datum in
\eqref{pnc:actual-reduced-datum}. An admitted specified fine
gauge gives a literal constrained-critical representative
$\widehat U(B)$ with $M_{\mathbf B}\widehat U(B)=V_B$.
It is the unique critical orbit in a fixed neighborhood of
the trial, modulo gauges one at all assigned tops.
The curvature and current differences from the trial have
weighted bounds $C(\delta+\|B\|_{1,\infty})$. The literal
multiplier and fixed-order localized analytic responses have
the corresponding bounds of the critical-reference theorem.
\end{proposition}

\begin{proof}
We give the changes to the model proof explicitly. None of the
following uses an inverse at a desired, as yet unknown orbit.
Put $H_{L,0}=F''(0)-\langle\omega_0,q''(0)\rangle$.
The real gap in \eqref{npl:proper-gap} and the full strong saddle
in \eqref{psr:strong-row-saddle} apply to
$H_{L,0}+D_0R_0D_0^\sharp+Q_0^\sharp a_fQ_0$ without
stationarity. Removing its row mass by the stated multiplier
shift gives the bounded full inverse $\mathcal L_0^{-1}$ in
\eqref{pnc:base-joint-inverse}. In particular no physical
Ward identity is used to turn a slice equation into a full
equation at the trial.

Use the complete actual gradient residual
$N_F(A)=F'(A)-j_0-F''(0)A$. The local Coulomb regularity
needed in \cref{lem:pnc-Wilson-residual} comes from
\eqref{rtc:regular-trial}, not from criticality. Its proof
therefore gives exactly the same quadratic bound and
$Cr$ Lipschitz constant. The actual constraint densities
in \eqref{pnc:constraint-residual-bounds} also hold about
this real regular reference, including the $\omega_0q'''$
remainder. Subtracting the trial equation now gives the
following complete fixed point:
\begin{equation}
 \begin{split}
 (A,\nu)=\mathcal L_0^{-1}\bigl(&-e_0-N_F(A)
        +(q'(A)-Q_0-q''(0)[A])^\sharp\omega_0\\
       &+(q'(A)-Q_0)^\sharp\nu,\quad b(B)-C_q(A)\bigr).
 \end{split}
 \label{rtc:actual-forced-fixed-point}
\end{equation}
The new term is the actual force $-e_0$, not an omitted
linear remainder. The nonlinear part has size $Cr^2$ and
Lipschitz constant $Cr$ in the same joint strong/multiplier
ball. Choose $r$ below the primitive ceilings with $Cr<1/4$,
then choose $\delta_*,\rho$ so that
$C(\delta_*+\rho)<r/8$. The same explicit iteration is
invariant and contractive and gives the last bound in
\eqref{rtc:corrected-critical-system}. Its affine source
does not affect differentiated contraction estimates.

The actual node diagram, outward datum map, and effective
gauge cocycle are reference identities. Thus
\eqref{pnc:actual-proper-gauge-cocycle} and
\eqref{pnc:nonlinear-proper-intertwining} remain exact.
The site extensions and the bound in \eqref{dsc:proper-site-gauge-bound}
use the actual trial current already assumed in
\eqref{rtc:regular-trial}. They give the same corrected
gauge tests in $\ker\tau_A$ and the surjectivity in
\eqref{pnc:gauge-test-surjectivity}. Testing the full
penalized equation at the fixed point therefore gives
\eqref{pnc:penalty-annihilator}, with arbitrary test
$\sigma\in\operatorname{ran}R_0$. Its conjugate test
removes the penalty also at complex data. There is no
extra $e_0$ in this gauge identity: it is the exact
Wilson equation being tested, not only its Taylor part.

The admitted literal extension and the mean-slice
construction in \eqref{pnc:literal-critical-representative}
and \eqref{pnc:actual-mean-slice} likewise use no trial
stationarity. They identify full fibres modulo precisely
the top-fixed gauge group, including bare zeros and
outward factors.

For uniqueness use the FULL penalized row lift $H_0$
at the trial. It has $Q_0H_0=1$, the strong bound and
the normalized exponential energy columns. It need not
be a horizontal physical lift. Those are all the
properties used to construct the right inverse
$\mathscr H_A=H_0(q'(A)H_0)^{-1}$. For a competing
critical mean-slice point with multiplier
$\omega_0+\nu$, its correct readout is
\begin{equation}
 \begin{split}
 \xi_c^3\langle\nu(c),e\rangle
  =\big\langle e_0+F'(A)-j_0
       -(q'(A)-Q_0)^\sharp\omega_0,
                         \mathscr H_A b_{c,e}\big\rangle,\qquad
                         b_{c,e}=\xi_c^{-1}e.
 \end{split}
 \label{rtc:competing-multiplier-readout}
\end{equation}
The new first pairing is bounded by $C\delta$: its
fine force is at most $\delta\xi^{-3}$ on a cell and
the column's physical $L^1$ norm is
$C\xi^3e^{-s d}$. The weak Wilson readout in
\eqref{pnc:Wilson-weak-readout} bounds the other terms.
Thus
\begin{equation}
 \|\omega-\omega_0\|_{3,\infty}
                      \le C(\delta+\|A\|_{\mathrm{str},\alpha}).
 \label{rtc:competing-multiplier-bound}
\end{equation}
A smaller fixed orbit neighborhood and $\delta_*$ put
this pair in the contraction ball. The mean-slice and
fixed-point uniqueness prove the claimed orbit uniqueness.

Finally subtract $j_0=Q_0^\sharp\omega_0+e_0$ from the
new stationarity equation. The transpose density,
normal-to-left gradient differential and fine literal
gauge give
\begin{equation}
 \|j(\widehat U(B))-j_0\|_{3,\infty}
 +\|F(\widehat U(B))-F(U_0)\|_{2,\infty}
                                \le C(\delta+\|B\|_{1,\infty}).
 \label{rtc:physical-continuation-bounds}
\end{equation}
The literal parameter multiplier is exactly
$Db(B)^\sharp(\omega_0+\nu(B))$, with the left-gradient
conversion in \eqref{pnc:literal-physical-multiplier}.
Its difference from $Db(0)^\sharp\omega_0$ has the same
defect-plus-datum bound. The weighted differentiated
contraction and fixed-order Cauchy argument in
\eqref{pnc:weighted-first-response} are unchanged.
This proves the proposition.
\end{proof}

\subsubsection{What is obtained on the real proper image}

For the fixed determining geometry let $\mathscr T_{\rm reg}$
be the real fine trials obeying \eqref{rtc:regular-trial},
with $\|j(U)\|_{3,\infty}<\delta_*/4$. Define the actual image
\begin{equation}
        \mathscr I_{\rm reg}
                   =\{M_{\mathbf B}(U):U\in\mathscr T_{\rm reg}\}.
 \label{rtc:actual-proper-image}
\end{equation}
This is a definition by the given fine fields, not a
characterization using coarse plaquettes alone.
For every $V\in\mathscr I_{\rm reg}$, choose an actual preimage
$U_0$. Apply \cref{prop:rtc-nonstationary-trial} with
$\omega_0=0$ and $B=0$. It produces a critical fine field
with literal mean $V$. Thus the constrained-critical fibre
is nonempty at every point of this real image. The same
preimage gives a holomorphic critical branch on a uniform
complex raw-angle ball about $V$, without presupposing a
critical reference there.

The image is nonempty and contains genuinely nonflat
examples. For instance use the identity reference and
a smooth one-colour potential supported in a complete
homogeneous layer patch of scale $\xi$. A rescaled fixed
profile with amplitude $t/\xi$ has actual fine curvature
at most $C|t|\xi^{-2}$ and Wilson current at most
$C|t|\xi^{-3}$. These estimates follow directly from its
two bounded physical differences and the exact plaquette
current; no unknown inverse is used. A cutoff equal to
one on the complete carriers of a coarse plaquette may
contain $a_2(x)=t(x_1-y_1)H/\xi^2$, $a_1=0$, for one
fixed colour $H$. There all factors commute. The actual
means reduce to line averages plus contour gradients;
the latter cancel around the coarse plaquette. Its
mean plaquette is therefore $e^{itH}$, whereas the
cutoff makes the trial the identity outside the patch.
For sufficiently small nonzero $t$ this is a nonflat
point of the image in \eqref{rtc:actual-proper-image}, with bounds
independent of the layer depth. This example requires
the indicated complete patch; flat trials require none.

There is also a precise local independence of the chosen
trial. Suppose two real trials are related by
\begin{equation}
             U_1^g=e^{i\eta C}U_0,
             \qquad \|C\|_{\mathrm{str},\alpha;U_0}<\varepsilon_r,
 \label{rtc:controlled-trial-overlap}
\end{equation}
where $g$ is unitary and one at all assigned tops.
Reduce the trial-current ceiling and working datum balls
once. The branches built about these trials then agree
as physical orbits on their common datum domain, and
their total Wilson actions agree.

Indeed covariance of the reference transports identifies
the strong norms before the real gauge $g$. The relative
connection after it is the divided logarithm of
$e^{i\eta A_1}e^{i\eta C}$. Its Banach product estimate
gives strong norm at most $C(\|A_1\|_{\mathrm{str},\alpha}
+\|C\|_{\mathrm{str},\alpha})$. This is the same
mesh-uniform relative BCH calculation used in
\eqref{pnc:finite-gauge-bound}. Choose $\varepsilon_r$
and the working radii so this and the specified literal
extension fit the $U_0$ uniqueness neighborhood. Both
fields are then critical with the same literal datum.
The mean-slice map and the competing-multiplier bound in
\eqref{rtc:competing-multiplier-bound} identify them
modulo a top-fixed fine gauge. The slice and extension
are holomorphic in the datum, so the identification is
holomorphic on the overlap. Gauge invariance identifies
the actions. These transitions satisfy the cocycle
where the controlled overlaps meet: their quotient
fixes an invertible connection and all assigned tops,
and propagation along each connected fine component
forces it to be the identity.

The construction is also covariant under simultaneous
arbitrary real unitary fine gauge changes of a trial and
its datum. They conjugate the actual full saddle and
all relative coordinates isometrically. Its unique
small fixed point consequently transforms covariantly;
different admitted literal extensions differ by a
top-fixed gauge. The action is invariant. No smallness
of the gauge function itself is required when the
reference is transformed at the same time.

These facts give existence on the real image and
holomorphic descent on controlled trial overlaps.
\Cref{bgc:global-comparison,bgc:proper-background} identify
arbitrary real regular critical preimages of the same mixed
datum, including those not initially given in a common trial
chart. Their action is consequently canonical on the stated
regular class. The existence proof there constructs regular
backgrounds for every sufficiently small curvature datum by
datum descent and critical ascent; it does not require a
separately smoothed fine trial for that datum.

\subsubsection{The fixed one-step face lift as a concrete corollary}

For comparison with the already proved first-step result,
fix $\eta=1/L$ and the actual centered blocks. Given any
real coarse $V$, set $U_0(V)=I$ on internal fine links
and $U_0(V)(x,\mu)=V_\mu(y)$ on every positive edge
crossing from $B(y)$ to $B(y+e_\mu)$; use the inverse
on reversed edges. All coordinate-order internal paths
stay in their block. Each length-one coarse line used
in the actual mean crosses exactly one face. Thus its
internal contour transport is $I$, its line transport
is $V_\mu(y)$, and its central reference line is the
same matrix. In the notation of the actual mean,
\begin{equation}
                  T_{y,x}(U_0(V))=I,\qquad Z_x(U_0(V))=I,
                        \qquad\mathcal M U_0(V)=V.
 \label{rtc:exact-face-trial}
\end{equation}
The logarithm taken here is $\log Z_x=0$, never a
logarithm or a group root of $V_\mu(y)$. Large and
central holonomies present no obstruction.

Every fine plaquette is $I$ unless both of its
directions cross their block faces. In that case
its four ordered factors are exactly the coarse
plaquette, up to the frame conjugation or orientation.
With $\epsilon=\sup_P\|V(\partial P)-I\|$ this gives
\begin{equation}
       \|F(U_0(V))\|_\infty\le CL^2\epsilon,
                    \qquad\|j(U_0(V))\|_\infty\le CL^3\epsilon.
 \label{rtc:fixed-L-trial-bounds}
\end{equation}
The current estimate uses the bounded number of
incident fine plaquettes and their exact sine terms
with prefactor $\eta^{-3}$. For the homogeneous one-step
geometry the assigned scale is one. Choose
$\epsilon<\epsilon_L$ so these two bounds meet the
trial ceilings. The proposition, at $B=0$ and
$\omega_0=0$, then constructs the actual one-step
critical orbit of mean $V$.

This recovers the existence in
\cref{ch:first-step-background-expanded} through the
new full proper-region inverse. The bounded-area
block-contour comparison there puts every competing
fixed-$L$ regular field of that same mean into this
trial neighborhood after a node-fixed axial gauge;
all fixed strong differences cost constants depending
only on $L$. Local uniqueness therefore identifies
this branch with that previously constructed background,
including its compatible complex continuations.
Here coarse chart closeness really does give the required
fine comparison: if $V_1=e^{iB}V_0$, the relative fine
potential of their face lifts is $\eta^{-1}B$ on crossing
edges and zero internally. Since $\eta=1/L$ is fixed,
its covariant strong norm is at most $C_L\|B\|_\infty$;
ordinary finite differences cost at most fixed powers
of $L$, with unitary reference transports. Thus shrinking
the coarse working balls makes their overlaps controlled
in the sense of \eqref{rtc:controlled-trial-overlap}.
The face-lift corollary concerns one bare step at fixed
$L$. The multilevel construction in \cref{bgc:proper-background}
uses the regular solution of the preceding truncation as
its next real trial, with the explicit $L^3$ current loss
and the subsequent uniform reset.

Indeed an iterated face lift concentrates rather than spreads
curvature. Take one positive coarse link $e^{i\theta H}$,
all other links $I$, on a coarse torus of side at least two,
with $H^2=I$ and small nonzero $\theta$. After $j$ face lifts,
$\eta=L^{-j}$ and a fine corner plaquette still equals
$e^{i\theta H}$. A face edge adjacent to exactly that
nontrivial plaquette has current $\pm\eta^{-3}\sin\theta H$;
the other incident plaquettes are $I$. Its action variation
is $\eta\sin\theta$, and the physical pairing is $\eta^4$.
Thus at assigned scale one,
\begin{equation}
 \|F(U_{\mathrm{face},j})\|_\infty\ge L^{2j}|\theta|,
 \qquad
 \|j(U_{\mathrm{face},j})\|_\infty\ge L^{3j}|\sin\theta|.
 \label{rtc:thin-depth-growth}
\end{equation}
Fixed small coarse curvature consequently does not make
these iterated trials uniformly regular.
\Cref{bgc:proper-background} instead descends the coarse datum
and solves a critical equation at each ascent, applying
\cref{bgc:reset} before the next level. This controls the
current as well as curvature uniformly in the depth.

% END LAYERED EXPANSION background/one_step_nonflat_background.tex

\subsection{Regular backgrounds on separated regions}
\label{bgc:section}

The preceding proper-region inverse and trial correction complete
the ascent through the determining sets of
\eqref{mb:determining-rows}. We give the matching argument, including
the current bound at the action boundary and uniqueness among all
regular critical configurations. These determining sets have no inward
rows. Every outward row and its entire lower endpoint block remain
part of the construction.

Fix $0<\alpha<1$ and the geometric floors needed for
\cref{prop:psr-strong-inverse,prop:rtc-nonstationary-trial}.
Use $L\ge13$, and the aligned separated geometry of
\cref{mb:datum-descent-critical-ascent}; increase its floors below
when choosing a fixed local regularity patch. Constants may depend
on these choices and $L,\alpha$, but not on depth, the number of
pillars, the volume, or the shapes and number of components of the
regions. They apply to every truncation with these floors.

For a real field define
\begin{equation}
 \mathcal R(U)=\max\left\{
   \sup_p\eta^{-2}\mathsf L(p)^2|\partial U(p)-1|,
   \sup_{b\ \mathrm{full}}\mathsf L(b)^3|J(U)(b)|\right\}.
 \label{bgc:regularity}
\end{equation}
Here the first supremum is over inner plaquettes, and a full bond
has all six incident plaquettes in the action region. Write
$\mathcal R_{\rm mat}$ for the same expression with
$\eta^{-2}|D_U^{\mathsf T}\partial U|$ in place of $|J(U)|$.
The scalar and traceless parts of a plaquette give
\begin{equation}
 \mathcal R(U)\le\mathcal R_{\rm mat}(U),\qquad
 \mathcal R_{\rm mat}(U)\le r+\tfrac32r^2
                  \quad\hbox{if }\mathcal R(U)\le r.
 \label{bgc:matrix-current}
\end{equation}
Indeed write $\partial U(p)=c_p1+i\eta^2F_p$.
The scalar divergence at a full bond is a sum of three opposite
differences of $c_p$, with $1-c_p=|\partial U(p)-1|^2/2$.
Use $\eta\le\mathsf L(b)\le\mathsf L(p)$ on each incident
plaquette. This proves the second inequality; the traceless
part proves the first. Thus either regularity convention gives
the other after a fixed reduction of its ceiling.

\subsubsection{The strong norm of a regular comparison}

The following observation permits use of the strong nonlinear
chart for competitors initially specified only by curvature and
current. It is an elliptic consequence of those currents, not an
embedding of a discrete $C^1$ ball into a uniform $C^{1,\alpha}$ ball.

For a comparison based at $U_0$, write
\begin{equation}
 \mathcal N_1(A)=\max\left\{
       \|\mathsf L A\|_\infty,
       \|\mathsf L^2\nabla_{U_0}A\|_\infty\right\}.
 \label{bgc:local-first-norm}
\end{equation}
This is the weighted first-order supremum norm used in the slice
absorption estimate. The energy space $X$ retains its physical
$L^2$ definition in the inverse estimates.

\begin{lemma}[Improvement of a comparison in the slice]
\label{bgc:strong-comparison}
Let $U_0$ and $U_1=e^{i\eta A}U_0$ be real, with
$\mathcal R(U_i)\le\varepsilon$. Suppose that
$R_0D_0^\sharp A=0$ and $a=\|\mathsf LA\|_\infty$ is below
a fixed small ceiling. Let $b=q(A)$ be the actual straightened
per-length row coordinate. Put $\beta=\sup_c\xi_c|b_c|$.
Then
\begin{equation}
 \mathcal N_1(A)+\|\mathsf L^3D_0^\sharp D_0A\|_\infty
       +\|A\|_{{\rm str},\alpha}
                         \le C(\varepsilon+\beta).
 \label{bgc:comparison-bound}
\end{equation}
The ceiling on $a$ and the constant are uniform in the geometry.
No critical equation for $U_0$ or $U_1$ is required.
\end{lemma}

\begin{proof}
The covariant $D_0$ in $D_0^\sharp D_0A$ is the curl, whereas
the $D_0$ in the Landau condition is the site derivative.
We first use the actual first-order slice calculation
\eqref{mb:slice-absorption}. Its boundary and normalization
checks are useful here. The exact differentiated plaquette product
gives, at a full bond,
\[
 \mathsf L(b)^3|D_0^\sharp D_0A(b)|
                  \le C\varepsilon+C_La\mathcal N_1(A).
\]
One obtains this by subtracting the two currents in their common
normal coordinate. The normal-to-left logarithmic differential
is $1+O(\eta A)$. The linear curvature correction has size
$C\varepsilon a$, and the nonlinear remainder has size
$C_La\mathcal N_1(A)$. These estimates use first covariant differences
of $A$. In the paired plaquette terms a difference of a product
places one difference on one factor; terms with two undifferenced
potentials carry the remaining mesh factor. The elementary
inequality $\eta|\nabla A|\le C\sup|A|$ on the same stencil
pays that factor. It gives no inverse power of the fine mesh.

At a nonfull inner bond the bond and every bond in its action
plaquettes belong to the bare determining layer. There
$\eta A$ is the bare row coordinate. The at most six plaquette
stencils therefore give the bound $24\beta$ for
$\mathsf L^3|D_0^\sharp D_0A|$. On an outward row the
additional factor reads only the prescribed interior rows of
the preceding level, so its group angle is at most $C_L\beta$.
The same conclusion holds if $\beta$ is initially measured
after the triangular datum change
\eqref{pnc:actual-reduced-datum}.

For clarity, use the actual derivative $Q_0=q'(0)$ in the
completed operator, rather than an auxiliary comb derivative.
The actual two-leg bound gives
\[
 \sup_c\xi_c|q(A)_c-Q_0A(c)|\le Ca^2,
 \qquad
 \|\mathsf L^3Q_0^\sharp a_fQ_0A\|_\infty
                                      \le C(\beta+a^2).
\]
In the second estimate $a_f(c)$ has size $\xi_c^{-2}$.
The physical transpose density cancels the fine counting
measure on the complete row carrier. This is precisely the
fixed-leg estimate used in \eqref{psr:row-pointwise}; it includes
the outer block of an outward row.

Let $K_0=K_{U_0,0}$ be \eqref{psr:actual-precision} with zero
independent multiplier. The exact identity in the slice is
\begin{equation}
 K_0A=D_0^\sharp D_0A+\Delta'_{U_0}A
                                  +Q_0^\sharp a_fQ_0A,
 \qquad
 \|\mathsf L^3\Delta'_{U_0}A\|_\infty
                                                   \le C_L\varepsilon a.
 \label{bgc:comparison-equation}
\end{equation}
To verify the last bound directly, the difference between the
quadratic Wilson plaquette and the covariant curl square has
coefficients bounded by $C\eta^{-2}|\partial U_0-1|$
on its finite four-link stencil. It therefore contributes at
most $C\varepsilon\mathsf L^{-2}$ times the local maximum
of $|A|$. Adjacent scales cost at most $L$.

The first-order consequence of
\cref{prop:psr-strong-inverse} applied to this equation gives
\[
 \mathcal N_1(A)\le C_1(\varepsilon+\beta)
                     +C_2(\varepsilon+a)\mathcal N_1(A).
\]
Every norm is finite on the finite lattice, so absorption is
legitimate before any uniform estimate has been obtained.
Choose $\varepsilon,a$ so that the last coefficient is at most
$1/2$. This proves the first-order bound and then the displayed
bound on $D_0^\sharp D_0A$. Now the entire right side of
\eqref{bgc:comparison-equation} has weighted supremum at most
$C(\varepsilon+\beta)$. Its $\mathcal F_p$ norm has the same
bound for each fixed finite $p>4/(1-\alpha)$. Apply
\eqref{psr:strong-inverse-bound} again to obtain the strong
bound in \eqref{bgc:comparison-bound}.
\end{proof}

\begin{lemma}[Comparison of arbitrary regular fibres]
\label{bgc:global-comparison}
Suppose
\[
 \mathcal R(U_i)\le\varepsilon,\qquad
 \sup_c|M_{\mathbf B}U_1(c)(M_{\mathbf B}U_0(c))^{-1}-1|
                                                        \le\delta.
\]
For $\varepsilon+\delta$ below a fixed ceiling,
there are a real $A$ and a unitary gauge $g$ such that
\begin{equation}
 U_1^g=e^{i\eta A}U_0,\qquad
 R_0D_0^\sharp A=0,\qquad g(y)=v[A](y)\quad(y\hbox{ a top}),
 \qquad \|A\|_{{\rm str},\alpha}\le C(\varepsilon+\delta).
 \label{bgc:global-slice}
\end{equation}
The reduced datum is the literal datum difference under
\eqref{pnc:actual-reduced-datum}. In particular, equal literal
data give $q(A)=0$ on every row. If also $A=0$, $g$ belongs
to $\mathcal G_{\rm res}$.
\end{lemma}

\begin{proof}
First use the preliminary gauge constructed in
\eqref{mb:rigidity-defect}--\eqref{mb:rigidity-sum}.
Here are the uniformity and seam checks. With
$W_n=\mathcal M^nU_1$ and $W_{0,n}=\mathcal M^nU_0$,
descend from $g_k=1$ by
\[
 g_n=1\text{ on the level-$n$ tops},\qquad
 g_n(x)=W_{0,n}(y,x)^{-1}g_{n+1}(y)W_n(y,x)
                                      \quad(x\in B(y)).
\]
The blocks of the next level and the current tops partition
the current member, and the zero-length contour equals one.
Thus these formulas extend one another and define one fine
gauge $\widehat g$. They align all the averaged internal
contours and fix every final top.

In a double block the reconstructed link differs from its
coarse mean by $C_L$ times the local plaquette defect, by
\eqref{mb:bond-reconstruction}. The coefficient of a difference
of the coarse means is one, since all other factors are unitary.
At a seam the inward comparison link belongs to the corridor
of the outward row one level above. The latter's outer block
lies in the lower layer; its internal paths use at most
$2(L-1)$ prescribed lower rows. Telescoping gives an endpoint
defect at most $2(L-1)\delta+C_L\varepsilon$.
The other endpoint is already aligned. Along a deeper pillar
the remaining errors sum as
$C_L\varepsilon\sum_{n<j}L^{2(n-j)}$.
Consequently every carried link, at every intermediate level,
has relative group error at most $C_L(\varepsilon+\delta)$.
This argument also controls the omitted inward comparison
links; no inward constraint is added.

Ascend through the truncations, keeping these same two fine
configurations; the fields are not replaced by their averages.
If $\mathsf L^{(j)}$ is a truncated assigned scale in physical
units, then $\mathsf L^{(j)}\le\mathsf L^{(k)}$ on each
bond and plaquette. Both fields consequently have truncated
regularity at most $\varepsilon$ at every stage, directly
from their full-depth bounds. The real gauge corrections
preserve these bounds. At an intermediate truncation the
carried-link estimate just proved bounds its literal datum
difference by $C_L(\varepsilon+\delta)$; the triangular
outward change therefore gives
$\beta_j=\sup\xi|q^{(j)}(A)|\le C_L(\varepsilon+\delta)$.
At the last truncation its bound is instead $C_L\delta$.
At depth zero the relative
logarithm has $\mathcal N_1$ norm at most $C_L(\varepsilon+\delta)$,
because the mesh in its units is one. At the next level replace
$A^-$ by $A_0=LA^-$; its $\mathcal N_1$ norm increases by at most
$L^2$, and $\eta A_0=\eta^-A^-$ leaves all old node factors
unchanged. Apply the real site correction following
\eqref{mb:site-pair-inverse}. Prescribe the correction factor
one at old tops and $v_{j-1}[A_0]^{-1}v_j[A_0]$ at promoted
tops, with the multiplication order of the node diagram.
The latter defect is at most $C_L\|\mathsf LA_0\|_\infty$.
The site equation has the actual weighted inverse and fixes
the Landau condition; its nonlinear remainders have the
quadratic bounds displayed there. It therefore gives
$\|\mathsf LA\|_\infty\le C_L\mathcal N_1(A_0)$ and the required
normalized node values.

The first-order absorption in
\cref{bgc:strong-comparison} now resets $\mathcal N_1(A)$ to
$C_*(\varepsilon+\delta)$ at each level. Choose $C_*$ first,
then $\varepsilon+\delta$ so that
$C_LL^2C_*(\varepsilon+\delta)$ lies below the fixed site and
slice ceilings. This is one choice for all levels, rather than
an iteration of the factor $L^2$. On the last truncation the
normalized datum is the actual literal difference, because
$\widehat g$ fixes both endpoints of each determining row.
The explicit triangular outward formula gives
$\sup\xi|q(A)|\le C_L\delta$, including a reversed outward
orientation. The strong conclusion follows from
\cref{bgc:strong-comparison}. If $\delta=0$, all interior
coordinates vanish, hence every outward node factor and
outward coordinate also vanish. Finally $v[0]=1$, and every
node of the present determining set is a top, so $A=0$ makes
$g$ a residual gauge.
\end{proof}

\subsubsection{Resetting the current as well as curvature}

\begin{lemma}[Uniform regularity reset]
\label{bgc:reset}
There are a fixed geometric floor, a ceiling $a_c>0$, and
$C_R$ such that every real constrained-critical field with
$\mathcal R(U)\le a_c$ satisfies
\begin{equation}
 \mathcal R(U)\le C_R\varepsilon
 \label{bgc:reset-bound}
\end{equation}
when all its ordinary and mixed datum test holonomies are
within $\varepsilon$ of one. The same conclusion, with twice
the constant, holds for $\mathcal R_{\rm mat}$ after decreasing
the ceiling. The constants are uniform over all truncations.
\end{lemma}

\begin{proof}
We apply the local calculation underlying
\eqref{mb:one-patch-improvement}, with the actual flat inverse
and nonlinear equation now available. We record the admission
checks needed for that application. First fix the auxiliary
flat construction and its collar radius $\rho_c$. Next choose
a fixed patch parameter $\sigma\ge L^2\rho_c$ that is a power of $L$, and only then
increase the ambient separation floor to $\rho\ge L^2\sigma$.
Put $R=20L\sigma$ in layer distance.
The auxiliary nested cubes described in
\hyperlink{hi:real-datum-profile}{the local datum-profile calculation}
have positive-level union inside $Q_{\sigma-2}$, a complete
two-level inner patch, and a flat bare exterior. The local
real gauge on $Q_\sigma$ uses curvature and the actual full
current. It gives potential, first difference, and flat
second-order reads at most $C\sigma r_R(x)$ in local units.
Its domain condition is $C\sigma a_c<c$.

Extend that potential by zero on the auxiliary torus. All
cutoff plaquettes belong to the bare layer, so their physical
current, including the terms created by the cutoff, has
weighted bound $C\sigma r_R(x)$. In the positive-level
members there was no cutoff. A constrained variation on
the auxiliary geometry patches to a constrained variation
of the original field: a row of lower assigned endpoints
is fixed by induction in the actual averaging recursion;
a row with an endpoint at its reading level is an auxiliary
row and has its entire double block inside $Q_\sigma$.
The outward endpoint block is included in this assertion.
Every changed action plaquette is also inside $Q_\sigma$.
The extended field is therefore critical at its own
auxiliary datum.

Apply the local datum transport gauge and the real site
correction relative to the identity. They give a preliminary
slice potential of small weighted supremum and a datum of
size $b_f\le C\sigma^2r_R(x)$. The zero extension's current
and \cref{bgc:strong-comparison} put this slice potential in
the strong ball, with size $r_f\le C\sigma^2r_R(x)$.
Thus the actual flat case of
\cref{prop:pnc-critical-continuation} applies, and its local
uniqueness identifies the auxiliary critical field with that
solution. The required conditions are
$C\sigma^2a_c<\min(r,\rho)$, with the fixed flat chart
radii $r,\rho$. They can be imposed after $\sigma$ is fixed.

Split the flat saddle's exponentially decaying datum sum
at distance $\sigma/8$. In the inner part a row at $y$ has
size $C(1+|y-x|)v_R(x)$ in the datum transport gauge; its
outward factor is one because its entire lower internal
contour has been prescribed. The linear distance is
summable against the kernel. The outer part contributes
$Ce^{-c\sigma}b_f$. The exact nonlinear force and row
remainders contribute $Cr_f^2$. The multiplier readout
uses the same split and contributes
$Ce^{-c\sigma}(b_f+r_f)$ outside and $Cr_f^2$ nonlinearly.
At the full bond the exact equation
$F'(A)=q'(A)^\sharp\omega$ reads only the complete carriers
of nearby rows. Its physical transpose density bounds the
current by that multiplier estimate. The first-order field
estimate bounds curvature. Thus both readouts obey
\[
 r_U(x)\le K v_R(x)
       +\{C\sigma^2e^{-c\sigma}+C\sigma^4a_c\}r_R(x).
\]
The original and auxiliary local scales differ by at most
$L$ at the readouts; this fixed factor is included in $C,K$.
Bare points away from the next member use prescribed
plaquettes directly and the sum of their six current terms.

Choose $\sigma$ so the exponential term is at most $1/4$,
and then $a_c$ so the quadratic coefficient is at most
$1/4$ and all preceding local and strong chart conditions
hold. The result is $r_U(x)\le Kv_R(x)+r_R(x)/2$.
Let
$P(x)=\sup_\tau e^{-\delta d_{\mathbf B}(x,\tau)}
|\partial V(\tau)-1|$, where $e^{\delta R}=\sqrt2$.
For $t>0$ maximize $r_U(x)/(P(x)+t)$.
At a maximizing point its value $S$ satisfies
$S\le\sqrt2K+S/\sqrt2$; this uses
$P(y)+t\le\sqrt2(P(x)+t)$ in the patch.
Letting $t$ decrease to zero gives $r_U\le5KP$.
The bounded displacement from a corner to the base of a
readout changes this constant by a fixed factor.
Since $P\le\varepsilon$, this proves
\eqref{bgc:reset-bound}. Finally,
\eqref{bgc:matrix-current} gives the current assertion.
\end{proof}

\subsubsection{The mixed datum and the critical ascent}

\begin{proposition}[The regular proper background]
\label{bgc:proper-background}
For the geometry fixed above there are $a_0,a_1,C_B>0$,
independent of depth and volume, with $C_Ba_1<a_0$, such
that every real mixed datum whose ordinary and mixed test
holonomies are within $\varepsilon\le a_1$ of one has a
constrained-critical configuration $U(V)$ with
\begin{equation}
              M_{\mathbf B}U(V)=V,\qquad
              \mathcal R_{\rm mat}(U(V))\le C_B\varepsilon.
 \label{bgc:background-bound}
\end{equation}
Any two critical configurations for the same datum with
$\mathcal R\le a_0$ lie in the same residual gauge orbit.
The orbit is covariant under real gauges of the datum.
\end{proposition}

\begin{proof}
The constants are chosen in the following order. Fix the
inverse and site bounds, the regularity patch and its
ceiling $a_c$, and a reset constant $B\ge\max(12,2C_R)$.
Fix the one-level trial-correction constant $C_A$ below.
Finally choose $a_0$ below the comparison, multiplier and
PNC uniqueness ceilings, and choose $a_1$ so that
\begin{equation}
 C_\downarrow a_1<\varepsilon_\downarrow,\qquad
 C_AL^3BC_\downarrow a_1<\min(a_c,a_0),\qquad
 6L^3BC_\downarrow a_1<\min(\kappa_*,\gamma_*,\delta_*),
 \label{bgc:radius-order}
\end{equation}
as well as the fixed RTC datum radius. Here
$\kappa_*,\gamma_*,\delta_*$ are its real-trial ceilings.
Enlarging a fixed constant or shrinking $a_1$ makes the
inequalities strict without changing their uniformity.

Descend the datum by
\eqref{mb:descent-old-rows}--\eqref{mb:descent-defect}.
The two-radius argument there gives $V^{(n)}$ at every
truncation, with all tests bounded by
$\varepsilon'=C_\downarrow\varepsilon$ and with the
new one-step mismatch at most $\varepsilon'$.
There is no accumulated factor: a deep new plaquette is
one or a single old coarse plaquette; every seam comparison
uses the smaller-radius tests within six current-level
spacings of the seam. Separation keeps that neighborhood
outside the part changed at the previous descent.

At depth zero every inner bond is prescribed. The field
$V^{(0)}_0$ is therefore critical in its singleton fibre.
Its curvature is at most $\varepsilon'$ and its full
current is at most $6\varepsilon'$, in bare units. This
starts the bound $\mathcal R_{\rm mat}\le B\varepsilon'$.

Suppose $U^{(n-1)}$ has that bound. Regard the same physical
field on the determining set of depth $n$. An assigned
scale increases by at most $L$. The curvature radius
increases by at most $L^2$ and the full-current radius by
at most $L^3$. Hence its new radius is at most
$\alpha_0=L^3B\varepsilon'$.
The inner action has not changed. A nonfull inner bond
has a missing incident plaquette. Its endpoints and every
bond in its remaining inner plaquettes are within $3\eta$
of $\Omega_0^c$. The separation
$\operatorname{dist}_\infty(\Omega_0^c,\Omega_1)
>R_1\mathfrak M\eta>L^2\eta$ puts this entire stencil
outside $\Omega_1$. Every local scale on it is therefore
$\eta$, and all its inner edges are bare determining rows.
The exact current is a sum
of at most six traceless plaquette terms with prefactor
$\eta^{-3}$, so
\begin{equation}
 \|J_{\Omega_0}(U^{(n-1)})\|_{3,\infty}
                                                   \le6\alpha_0.
 \label{bgc:ascent-boundary-current}
\end{equation}
This checks the actual force required by RTC, including
the action boundary. A bound only on full bonds would
not suffice for that application.

On a double block of a new row, every level-$(n-1)$ bond
read by its one-step average lies in the last determining
member of the old truncation. Exact local recursion gives
\[
 \mathcal M^nU^{(n-1)}(c)
                         =\mathcal M(V^{(n-1)}_{n-1})(c).
\]
All older row values remain the target values. Define the
literal per-length displacement by
$\xi_cB_c=i^{-1}\log(V^{(n)}_n(c)
(\mathcal M^nU^{(n-1)}(c))^{-1})$ on new rows and by
zero on older rows. Then
$\|B\|_{1,\infty}\le(\pi/2)\varepsilon'$.
At a new outward endpoint the reduced datum factor reads
only older interior displacements, which are zero.
Thus \eqref{pnc:actual-reduced-datum} gives the same raw
new displacement, with the exact orientation order.

Apply \cref{prop:rtc-nonstationary-trial} with this actual
real trial and independent row multiplier zero. Its defect
is the force in \eqref{bgc:ascent-boundary-current}; the
forced equation retains this entire defect. The physical
continuation bounds give a real critical field at the
target literal datum, of provisional regularity at most
$C_A\alpha_0$. The literal node restoration is unitary,
so it preserves curvature and the full current; the
outward factors cancel exactly. The inequalities
\eqref{bgc:radius-order} put this output inside $a_c$.
Apply \cref{bgc:reset} at its new datum to return to
$\mathcal R_{\rm mat}\le B\varepsilon'$.
This completes the induction with one radius for all $n$.
At the final level apply the reset once more using the
original datum radius $\varepsilon$, obtaining
\eqref{bgc:background-bound} with a fixed $C_B$.

For uniqueness take any two real regular critical fields
of the same datum. The action boundary obeys the same
bound as \eqref{bgc:ascent-boundary-current}, so either
field is an admitted physical critical reference after
choosing $a_0$ sufficiently small. Apply
\cref{bgc:global-comparison} to the pair with $\delta=0$.
It gives a strong mean-slice point of size at most $Ca_0$
and zero reduced datum. This is a coordinate on the actual
constraint fibre by the gauge parametrization in
\eqref{pnc:actual-mean-slice}; constrained criticality
therefore gives its full multiplier equation.
The competing-multiplier readout
\eqref{pnc:competing-multiplier-readout} bounds its
multiplier displacement by $Ca_0$. Choose $a_0$ so both
components lie in the fixed PNC contraction ball.
They and the zero solution have the same datum; uniqueness
gives $A=0$. The last assertion of
\cref{bgc:global-comparison} makes the gauge residual.
Covariance follows by transforming the entire actual
equation and then applying this uniqueness.
\end{proof}

\begin{corollary}[A fixed physical continuation ball]
\label{bgc:uniform-physical-chart}
For every orbit in \cref{bgc:proper-background}, choose a
real representative $U_0$. There is $r_G>0$, uniform in
the datum, depth and admitted geometry, such that for
$V_B(c)=e^{i\xi_cB_c}V(c)$ and
$\|B\|_{1,\infty}<r_G$ there is a holomorphic literal
critical representative $\widehat U(B)$ with
$M_{\mathbf B}\widehat U(B)=V_B$. In its prescribed
Landau coordinate,
\begin{equation}
 \|A(B)\|_{{\rm str},\alpha}
       +\|\omega(B)-\omega(0)\|_{3,\infty}
                         \le C\|B\|_{1,\infty},
 \qquad
 \|F(\widehat U(B))\|_{2,\infty}
       +\|J_{\Omega_0}(\widehat U(B))\|_{3,\infty}
                         \le C(\varepsilon+\|B\|_{1,\infty}).
 \label{bgc:physical-chart-bounds}
\end{equation}
The physical bounds hold on complete fine plaquette and
current stencils. Their fixed-order localized responses
are those of \eqref{pnc:localized-analytic-jets}.
The construction requires no hereditary complex-domain
hypothesis.
\end{corollary}

\begin{proof}
The regular real reference just constructed has bounded
inner current by the bare-boundary calculation above and
bounded induced multiplier by the actual row lift.
All hypotheses of \cref{prop:pnc-critical-continuation}
therefore hold, uniformly. Its literal node extension and
\eqref{pnc:physical-fine-bounds} prove the assertion after
one fixed decrease of its datum radius. If desired choose
$r_G$ so the right side of \eqref{bgc:physical-chart-bounds}
is below any fixed small regularity ceiling on the full
complex ball. This decreases a fixed radius and introduces
no dependence on depth or $\varepsilon$.

For real $B$ in this ball the output is a real regular
critical orbit. Wherever it also lies in the real datum
domain of \cref{bgc:proper-background}, that theorem's
uniqueness identifies it with the real background there.
Thus these are continuations of the canonical real orbit.
Changing a real representative conjugates the equations
isometrically. On a common real datum neighborhood the
physical orbits agree by uniqueness; their holomorphic
germs, and their actions, agree by the identity theorem.
No choice of a distant complex gauge representative is
needed for these local continuations.
\end{proof}


\subsection{Uniform continuation on every tower}
\label{htc:section}

The tower tests require a gain that is quadratic in the complex
potential radius. A first derivative estimate alone gives a linear
loss. The argument below obtains the quadratic gain from three real
configurations, after constructing a common complex neighborhood on
which the normalized curvature and current are bounded. The real
background theorem supplies the three real estimates; the constrained
Euler equation supplies the crude complex current estimate. The
independent current of a regular pair does not enter this construction.
This is the tower argument underlying Ba{\l}aban's construction
\cite{B87}, with the determining rows, complex gauges and
parameter neighborhoods specified here.

\subsubsection{The two erosion rules}

We use the two cube rules fixed in the
\hyperref[notdef:localization-cubes]{localization geometry}.
Write $M_1=M_1'$ in this subsection.
All continuation and restriction statements below use determining
sets with all inward rows removed. The box-datum lemma also states
the variant in which those rows are retained.
All lengths in the next lemma are in bare units. Thus level $q$ has
spacing $L^q$. The constants satisfy
\[
 L\ge13,\qquad R_1\ge L,\qquad M_1\ge R_1\mathfrak M,
                    \qquad M\ge L^4M_1.
\]
The background floors are imposed before these inequalities. The
localization size $M$ is then increased, if necessary; no constant
chosen below depends on the cutoff, the tower depth or a coupling.

\begin{lemma}[Protected stencils]\label{htc:geometry}
Let $Z$ be a union of level-$(p-1)$ localization cubes, and construct
either of its two towers $Z=\Omega_0\supset\cdots\supset\Omega_p$.
For every $n\le p$, the truncation at $n$ is an admitted separated
sequence with no inward determining rows. Its two-layer inset
$Z^{\ominus2}$ has distance at least $ML^{p-1}$ from $\Omega_n^c$.
Every bare site within distance $4L^n$ of that inset lies in a free
level-$n$ pillar. In particular every plaquette and every full
six-plaquette current stencil based in the inset lies in these pillars.
\end{lemma}
\begin{proof}
Let $w_q$ be an upper bound on the distance from a discarded point to
$Z^c$. The discarded cube at the next step either meets the old
complement or fails its distance test. Its side plus the rejected gap
therefore gives
\[
 w_0=0,\qquad
 w_{q+1}-w_q=
 \begin{cases}
 R_1\mathfrak M L^q+\mathfrak M L^{q+1},&\text{first rule},\\
 2M_1L^{q+1},&\text{second rule}.
 \end{cases}
\]
Nested grids ensure that this argument also applies to a cube not
contained in the previous member. Summing the geometric series gives
$w_q\le ML^{p-3}$ for either rule and all $q\le p$.
Consequently
\begin{equation}\label{htc:erosion-loss}
 \operatorname{dist}_\infty(x,\Omega_q^c)
 \ge\operatorname{dist}_\infty(x,Z^c)-ML^{p-3}.
\end{equation}
The two-layer inset has distance at least $2ML^{p-1}$ from $Z^c$,
which proves the asserted distance bound.

Both rules give unions of the required aligned cubes and gaps greater
than $R_1\mathfrak M L^{q-1}$ at level $q$. At a nonempty localization
the level-$n$ torus side has at least $M/L\ge L^3R_1\mathfrak M$
sites, so the torus floor for the background theorem is satisfied.
The statements remain valid when a member is empty; for the full
torus every member is the full torus.

If $x$ is within $4L^n$ of the inset, its level-$n$ block has distance
at least $(M/L-5)L^n>L^{n-1}$ from $\Omega_n^c$. Its only determining
node is its top: a lower row has an endpoint outside $\Omega_n$ and
length at most $L^{n-1}$, whereas a level-$n$ node is a block center.
This proves that the pillar is free. A current stencil reaches at
most two bare steps, which are contained in the stated halo.
\end{proof}

For a separated sequence of depth $n$, put $\eta=L^{-n}$. Let
$\ell_-(b)$ be the smaller assigned physical pillar length at the
endpoints of $b$, and take the minimum over the corners for a
plaquette. A level-$q$ determining row and its full double-block read
have $\ell_-\ge L^{q-1-n}$. This follows because the read has radius
at most $3L^{q-n}/2$, while the preceding layer gap is greater than
$R_1\mathfrak M L^{q-1-n}$. Define
\begin{equation}\label{htc:input-norm}
                  \|w\|_-:=\sup_b\ell_-(b)|w(b)|.
\end{equation}
The real and imaginary Hermitian parts of $w$ have no larger norm.

\subsubsection{A common complex ball}

The next lemma uses only the real mixed background and its physical
datum chart, proved in
\cref{bgc:proper-background,bgc:uniform-physical-chart}. In particular,
neither result is applied with a hereditary core as an input.
We write $b$ for a raw row logarithm: the datum is $e^{ib}V_0$.
Thus $\|b\|_\infty$ is the norm called $\|B\|_{1,\infty}$ when
$b(c)=\xi_c B(c)$.

Throughout this subsection, $\mathcal H_{\mathbf B,V_0}(b)$ means
the Landau component of the physical chart with \emph{literal}
datum $e^{ib}V_0$. Its evaluation includes the triangular conversion
\eqref{pnc:actual-reduced-datum}: if $\mathfrak n_{\mathbf B,V_0}(b)$
denotes that converted per-length datum, the solved equation is
$q(H)=\mathfrak n_{\mathbf B,V_0}(b)$.
The completed representative is
\[
 \Phi_{\mathbf B,V_0}(b)=(e^{i\eta H}U_0)^{u_b},\qquad
 M_{\mathbf B}(e^{i\eta H}U_0)
                         =(e^{ib}V_0)^{u_b^{-1}|_{\rm nodes}}.
\]
Thus $H$ alone need not have the literal datum. Every estimate in
the literal variable $b$ includes this conversion and the node
factors. Their finite carriers and uniform derivatives are part
of \eqref{pnc:reduced-datum-bounds} and
\eqref{pnc:weighted-first-response}.

\begin{lemma}[Uniform analytic readouts]\label{htc:common-ball}
There are $a_*,r_*,C_*>0$, depending only on the fixed background
geometry, with the following property. Let $U$ be real and satisfy
\begin{equation}\label{htc:real-input}
 |\partial U(P)-1|<e_0\eta^2\ell_-(P)^{-2},\qquad e_0<a_*.
\end{equation}
For every $\|w\|_-<2r_*$ the averaged datum
$M_{\mathbf B}(e^{i\eta w}U)$ and its continued critical configuration
$\Phi_U(w)$ exist and depend holomorphically on $w$. For a plaquette
or full current stencil contained in free pillars of one level $q$,
put $\xi=L^{q-n}$. Then
\begin{equation}\label{htc:crude-readouts}
 \left|\eta^{-2}\xi^2(\partial\Phi_U(w)-1)\right|\le1,
                   \qquad |\xi^3J_\eta(\Phi_U(w))|\le1.
\end{equation}
The same $r_*$ works for every depth, compatible volume, real base
$U$ satisfying \eqref{htc:real-input}, and either tower rule.
\end{lemma}
\begin{proof}
The actual averaging estimates give constants $a_{\rm av}>0$ and
$C_Q\ge1$, independent of the number of iterations, such that
\begin{equation}\label{htc:mean-ball}
 b(w)=\frac1i\log\{M_{\mathbf B}(e^{i\eta w}U)
                                  M_{\mathbf B}(U)^{-1}\},
 \qquad \|b(w)\|_\infty\le C_QL\|w\|_-
\end{equation}
is holomorphic when $L^2e_0<a_{\rm av}$ and
$L\|w\|_-<a_{\rm av}$. Indeed the order-$q$ read has
$L^q\eta|w|\le L\|w\|_-$ and curvature at most
$L^2e_0L^{-2q}$ on its whole double block. Apply
\cref{prop:averaging-1a-1-6} on that block. At level zero the formula
is $b(w)=\eta w$. The same averaging estimate shows that the real
datum $V_0=M_{\mathbf B}(U)$ has every ordinary and mixed test
bounded by $C_Ve_0$. On a mixed plaquette the coarse bond reconstructed
from the lower rows equals the actual next average, by the semigroup
identity. No new seam defect occurs.

Choose $a_*$ so $C_Va_*$ lies below the real datum ceiling. The real
background theorem supplies its normalized critical representative
$U_0$ and the bound
\begin{equation}\label{htc:real-background-bound}
 \eta^{-2}\ell(P)^2|\partial U_0(P)-1|
       +\ell(b)^3|J_\eta(U_0)(b)|\le C_Be_0.
\end{equation}
Here and below a current is evaluated only at a full stencil; boundary
forcing is retained in the regional equations used to construct $U_0$.
The physical chart gives $W(b)=e^{i\eta H(b)}U_0$, holomorphic for
$\|b\|_\infty<r_B$, with weighted potential and first differences at
most $C_H\|b\|_\infty$. It also bounds the change of multiplier by
$C_H\|b\|_\infty$.

We record why its current is bounded on that entire ball. If $T_H$
is the differential of the exponential coordinate in left-trivialized
physical variables, the equation after removal of the gauge penalty is
\begin{equation}\label{htc:Euler-current}
 T_H^\sharp J_\eta(W(b))
                  =q'(H)^\sharp(\omega_0+\nu(b)).
\end{equation}
The adjoints in this identity include the physical row and bond
densities. The actual row-transpose estimate and the multiplier
readout of \eqref{pnc:literal-multiplier-bound} give
\[
 \|\ell^3J_\eta(W(b))\|_\infty
       \le C(e_0+\|b\|_\infty).
\]
At $H=0$, the multiplier also has a uniform bound. If $R_0$ is the
actual right inverse of the retained row map, apply $R_0^\sharp$ to
$J_\eta(U_0)=q'(0)^\sharp\omega_0$. The weighted row-lift estimate
in \eqref{pnc:actual-row-right-inverse} gives
$\|\omega_0\|_{3,\infty}\le C\|\ell^3J_\eta(U_0)\|_\infty
\le CC_Be_0$. Thus both the base multiplier and its increment have
the claimed bounds. The bondwise factors $T_H^{\pm1}$ are bounded because
$\eta|H|\le\ell|H|$ is small. Thus this estimate comes from the
Euler equation and its multiplier bound. It does not require a
second-difference estimate inferred from the first-difference norm.
The plaquette product similarly gives
$\eta^{-2}\ell^2|\partial W(b)-1|\le C(e_0+\|b\|_\infty)$.

We must also control the prescribed normalized representative.
Fix a free pillar with top $y$ at level $t$ and let
$h_y=\sup_{B^t(y)}\xi_y|H|$. For a level-$q$ segment of the ascent
to $y$, retain its real reference contour as a unitary factor.
Its condition number is $1$; it need not be the identity.
The actual relative contour logarithm is bounded by
$C_\Gamma C_Q L^{q-t+1}h_y$. To see the scale, apply the actual mean
estimate to the order-$q$ double block: its relative row logarithm
is at most $C_QL^{q-t}h_y$. A one-step comb has at most $4(L-1)$
bonds; the finite product logarithm on the fixed chart costs at most
twice their sum. Indeed, factoring each bond into its relative
factor and its real reference factor rewrites the relative contour
as a product of unitary conjugates of the relative factors.
The same bound for the mean contour follows from
the contraction defining that contour. The node factor at the top
is a product of exponentials with logarithms bounded by
$C_NL^{q-t}h_y$. These are the one-step site corrections in
\cref{prop:averaging-1a-1-6}: at level $q$, their normalized input
is $L^{q-t}h_y$, and the inverse derivative of the site mean is
bounded on its fixed chart. The constant $C_N$ depends on $L$,
including the finite comb length, and not on $q$ or $t$.
The real reference node factors are unitary as well. Telescope
these factors through the ascent; their conjugations preserve all
norms in the estimates just obtained. Each ascent meets a level at
most once. Therefore the gauge $u_b$
which takes $W(b)$ to the prescribed representative satisfies
\begin{equation}\label{htc:pillar-gauge}
 |u_b(x)^{\pm1}|\le
 \exp\!\left(C_Gh_y\sum_{q=0}^{t-1}L^{q-t}\right)
 \le e^{C_Gh_y/(L-1)},\qquad
                    \Phi_U(w)=W(b(w))^{u_{b(w)}}.
\end{equation}
The identity holds at real data by the defining axial normalization
and at complex data by holomorphy. All factors use the actual
averaging contours and node correction. Choose $r_G\le r_B/2$ so
$\|b\|_\infty<r_G$ makes the condition number of this gauge at most
$2$. Conjugation then costs at most $2$ in both readouts. The entire
stencil lies in the indicated free pillars, so it has no additional
boundary term.

Finally choose, in this order,
\begin{equation}\label{htc:fixed-radius}
 r_*\le\frac1{4L}\min\{a_{\rm av},r_G/C_Q\},\qquad
 C(a_*+2C_QLr_*)\le\frac12.
\end{equation}
The constants in the second inequality are the fixed plaquette,
Euler-current and gauge constants just obtained. Decrease $a_*$
and $r_*$ once to satisfy it. Equations
\eqref{htc:mean-ball}, \eqref{htc:Euler-current} and
\eqref{htc:pillar-gauge} now prove
\eqref{htc:crude-readouts} throughout the ball $\|w\|_-<2r_*$.
\end{proof}

\subsubsection{The quadratic gain}

\begin{lemma}[Three real configurations]\label{htc:three-real}
Let $F$ be a matrix-valued holomorphic map on $\|w\|<2r$, bounded
there by $1$, and write $w=w_{\rm r}+iw_{\rm i}$ with
$\|w_{\rm r}\|,\|w_{\rm i}\|\le\|w\|\le r$. Then
\begin{equation}\label{htc:three-real-estimate}
 |F(w)|\le |F(w_{\rm r})|+|F(w_{\rm i})|
                  +|F(0)|+\frac{3\|w\|^2}{2r^2}.
\end{equation}
\end{lemma}
\begin{proof}
For each of the three vectors, the one-variable Taylor series on
$|z|<2r/\|w\|$ bounds the remainder after degree one by
$\|w\|^2/(2r^2)$. Apply the triangle inequality to
$F(w)-F(w_{\rm r})-iF(w_{\rm i})+iF(0)$.
Its constant and linear terms vanish. The estimate follows by
duality with norm-one linear functionals on the matrix space.
\end{proof}

\begin{proposition}[Uniform tower bound]\label{htc:tower-bound}
There are $a_T>0$ and $K_T\ge1$ such that, on every separated sequence
of depth $n$, the assumptions
\begin{equation}\label{htc:quadratic-input}
 \begin{gathered}
 |\partial U(P)-1|,\quad
 |\partial(e^{i\eta f}U)(P)-1|
       <e_0\eta^2\ell_-(P)^{-2},\\
 \|f\|_-<e_1,\qquad e_0+e_1<a_T,
 \end{gathered}
\end{equation}
with $U$ real imply, at every free level-$q$ full read,
\begin{equation}\label{htc:quadratic-output}
 |\partial\Phi_U(f)-1|
       <K_T(e_0+e_1^2)\eta^2\xi^{-2},\qquad
 |J_\eta(\Phi_U(f))|
       <K_T(e_0+e_1^2)\xi^{-3},\quad \xi=L^{q-n}.
\end{equation}
The bounds also hold for $f=0$. The configuration is holomorphic
in $f$ on the fixed ball of \cref{htc:common-ball}, regardless of
whether the second curvature bound in \eqref{htc:quadratic-input}
holds at other points of that ball.
\end{proposition}
\begin{proof}
Transport the four increments of a plaquette to its initial corner
using $U$. Their sizes are at most $\eta e_1/\ell_-(P)$.
The remainder after the linear term in their exponential product is
bounded by $C\eta^2e_1^2\ell_-(P)^{-2}$.
The closing transport changes the linear curl by at most
$C\eta|f|\,|\partial U-1|$. Hence the two curvature bounds imply
\[
 \ell_-(P)^2|\mathrm d_U f(P)|\le C(e_0+e_1^2).
\]
Real unitary transports commute with taking Hermitian real and
imaginary parts. Applying the same product estimate to $f_{\rm r}$
and $f_{\rm i}$ proves
\[
 |\partial(e^{i\eta f_\sigma}U)-1|
       \le C_P(e_0+e_1^2)\eta^2\ell_-^{-2},
                         \qquad \sigma\in\{\mathrm r,\mathrm i\}.
\]
Choose $a_T\le\min\{r_*,a_*/(2C_P),1/2\}$ and reduce it, if
necessary, to meet the real datum ceiling for these two configurations.
Their averaged data are regular, and the continued real fields are
their unique regular critical orbits by
\cref{bgc:proper-background,bgc:uniform-physical-chart}.
Thus each normalized readout at $0,f_{\rm r},f_{\rm i}$ is bounded by
$C_RC_P(e_0+e_1^2)$. Here the prescribed gauges at real points are
unitary, so they preserve the norm.

Apply \cref{htc:three-real} to the two readouts in
\eqref{htc:crude-readouts}. It gives \eqref{htc:quadratic-output}
with, for example,
$K_T=1+4C_RC_P+2r_*^{-2}$. This number has no depth or real-base
dependence. In the notation of \eqref{mb:complex-readout-disc},
one may take $M=1$ and $r_0=r_*$ after the harmless conversion from
a ball to the displayed two-variable polydisc. In particular
$M/r_0^2=r_*^{-2}$ is a fixed constant. No running normalization is
part of this radius.
\end{proof}

Write $\mathcal Q_j(X;e_0,e_1,c_0)$ for the simple triple class of
\eqref{dc:simple-core}, with independent current bound $c_0$.
We display this third radius in $\mathcal C_j$ when it differs
from the curvature radius.

\begin{corollary}[Both towers at every truncation]\label{htc:all-towers}
Let $(U,A',\mathbf J)$ satisfy the simple conditions on a level-$j$
domain $X$, with radii $(e_0,e_1,c_0)$ and $e_0+e_1<a_T$.
For every tower address $(p,Z)$, both erosion rules and $1\le n\le p$,
the tower at the actual averaged datum exists. In level-$n$ units,
on $Z^{\ominus2}$ it satisfies
\begin{equation}\label{htc:address-bound}
 \begin{aligned}
 |\partial U_n(Z;\mathbf U)-1|
   &<K_T(e_0+e_1^2)L^{-(j-p)}L^{-2p},\\
 |J^{(n)}(U_n(Z;\mathbf U))|
   &<K_T(e_0+e_1^2)L^{-(j-p)}L^{2(n-p)}.
 \end{aligned}
\end{equation}
The same is true for the real factor. Consequently
\begin{equation}\label{htc:core-inclusion}
 e_0\le b_0,\quad e_1\le b_1,\quad
 K_T(e_0+e_1^2)\le b_0
 \quad\Longrightarrow\quad
 \mathcal Q_j(X;e_0,e_1,c_0)
             \subset\mathcal C_j(X;b_0,b_1,c_0).
\end{equation}
The third radius denotes the independent current bound; when omitted
it equals the first radius.
\end{corollary}
\begin{proof}
First take $j=p$ and put $\lambda=L^{n-p}$. On changing from level
$p$ to level $n$, the potential becomes $\lambda A'$, the physical
curvature gains $\lambda^2$, and the real gauge clauses remain valid.
The simple input therefore satisfies \eqref{htc:quadratic-input} with
$\epsilon_0=\lambda^2e_0$ and $\epsilon_1=\lambda e_1$.
The minimum length is at most $1$, so these uniform bounds imply its
weighted bounds. By \cref{htc:geometry}, every required output
stencil lies in free top pillars, where $\xi=1$. The output curvature
has factor $\lambda^2L^{-2n}=L^{-2p}$, and its level-$n$ current has
factor $\lambda^2=L^{2(n-p)}$.

For a general address of a level-$j$ domain, restriction to level $p$
has simple radii $\mu e_0,\mu e_1$ with $\mu=L^{-(j-p)}$.
The curvature and first difference actually gain $\mu^2$, which is
stronger than the radius asserted here. Apply the preceding result
and use $\mu e_0+\mu^2e_1^2\le\mu(e_0+e_1^2)$. The independent
current restricts by $\mu^3$ and has no role in a tower's physical
current. This proves every clause, without summing over addresses.
For a bounded localization domain in $\mathbb R^4$, choose a torus
large enough to contain it and every determining read injectively.
The regional action, row constraints and normalization paths are
then exactly the same finite objects. Uniqueness makes the tower
independent of the chosen torus, so all assertions hold on that
domain with the same constants.
\end{proof}

\subsubsection{Membership throughout parameter discs}

\begin{corollary}[Uniform parameter domains]\label{htc:parameter-core}
Let $U$ be fixed and let $A'(z),\mathbf J(z)$ be holomorphic on an
open polydisc $D$, with their values satisfying the simple bounds
at fixed radii $(e_0,e_1,c_0)$ throughout $D$. If the inequalities
of \eqref{htc:core-inclusion} hold, the whole image of $D$ lies in the
target core. All tower maps are holomorphic on $D$, with the common
bounds displayed in \eqref{htc:address-bound}. A strict inequality
$K_T(e_0+e_1^2)<b_0$ supplies a common strict margin, also on every
smaller closed polydisc. No radius loss
depending on the number of tower addresses is required.
\end{corollary}
\begin{proof}
The mean and background compositions in \cref{htc:common-ball} are
defined before testing core membership. The simple bounds put each
composition in its fixed ball at every point of $D$.
\Cref{htc:all-towers} therefore applies pointwise with the same
constants, and composition proves holomorphy. Taking the supremum
over all addresses preserves this common bound; no uniform strict
margin is inferred merely from individual strict inequalities.
The real factor and its
real gauges are fixed; the independent current uses its own assumed
simple inequality.
\end{proof}

\begin{proposition}[Outer substitutions]\label{htc:outer-substitution}
Fix $c_\natural\ge1$. There is $c_O>0$, depending only on the fixed
geometry and $c_\natural$, such that the following holds. Suppose
$0<b_0\le b_1$, $b_1^2\le c_\natural b_0$, and $b_0+b_1$ is below
a fixed ceiling. Start from a triple in the core with radii
$(b_0/2,b_1/2,c_0)$ and put $\mathbf U=e^{i\eta_jA'}U$.
If
\[
 \max\{|A|,|\nabla_{\mathbf U}A|\}<c_Ob_0,
                       \qquad |\mathbf J'|<b_0,
\]
then $(e^{i\eta_j zA}\mathbf U,\mathbf J')$ belongs to the core
at radii $(b_0,b_1,b_0)$ for every $|z|\le1$, with the real factor
$U$ and the principal Baker--Campbell--Hausdorff potential. The same
conclusion is uniform in any other parameters for which the stated
bounds hold on their whole polydisc.
\end{proposition}
\begin{proof}
Choose $R=64K_T(1+c_\natural)$, and then choose
$c_O\le(16C_OR)^{-1}$, where $C_O$ bounds the first-order product
and plaquette estimates on the fixed logarithm chart. For $|z|<R$,
\[
 A'_z=(i\eta_j)^{-1}
          \log(e^{i\eta_jzA}e^{i\eta_jA'})
\]
is holomorphic. The product rule, including the change from
$\nabla_{\mathbf U}$ to $\nabla_U$, gives
$|A'_z|,|\nabla_UA'_z|<b_1$ and
$|\partial(e^{i\eta_jA'_z}U)-1|<b_0\eta_j^2$.
Indeed the additional terms are bounded by
$C_O(Rc_Ob_0+(Rc_Ob_0)^2+Rc_Ob_0b_1)$; reducing the fixed ceiling
makes their sum smaller than $b_0/4$. The real gauges remain those
of $U$. Thus the simple conditions hold on this entire auxiliary
disc, although membership in the target hereditary core has not
yet been used there.

At an address $(p,Z)$ divide the two tower readouts by the factors
$L^{-(j-p)}L^{-2p}$ and $L^{-(j-p)}L^{2(n-p)}$, respectively.
By \cref{htc:all-towers} each resulting holomorphic function $F(z)$
is bounded on $|z|<R$ by
$K_T(1+c_\natural)b_0$. Schwarz applied to $F(z)-F(0)$ gives
\[
 |F(z)-F(0)|\le\frac{2K_T(1+c_\natural)b_0}{R}
                    \le b_0/32\qquad(|z|\le1).
\]
The old core gives $|F(0)|<b_0/2$. This proves every complex tower
test with room. Its real counterpart is unchanged. The simple
conditions and $|\mathbf J'|<b_0$ complete the inclusion, uniformly
over the auxiliary parameters.
\end{proof}

\subsubsection{Inner window substitutions}

We now prove the core inclusion for the actual inner window map.
This application uses the upper triple's tower clause. That clause
has either been retained from the upper core or has been established
by \cref{htc:parameter-core}; it is not an assumption used to prove
\cref{htc:tower-bound}.

Let $\square$ be the union of the $2^4$ level-$k$ localization cubes
meeting a vertex, and let $\square_0=\square^{\sim5}$ be a box with
an injective torus lift. Put $N=k+1$, $\nu=N-j\ge1$ and
$\sigma=L^{-\nu}$. Let $\mathbf B^+$ be the first-rule tower of
$\square_0$ to depth $N$, and $\mathbf B_j$ its truncation at $j$.
For a dimensionless fine logarithm $a$ write
$\mathsf q_q(a)=i^{-1}\log\mathcal M^q(e^{ia})$ on the fixed
near-trivial chart. This is the plain nonlinear mean logarithm;
it is not the linear row map or the straightened constraint.
The upper triple has radii
$\alpha_\iota=(1+2\beta_s)A_\iota$, $0<\beta_s\le1/4$.
We require its simple conditions on $\square_0$ and its tower
clause at $(N,\square_0)$. Thus its upper tower has curvature
less than $\alpha_0L^{-2N}$ and level-$N$ current less than
$\alpha_0$ on $\square^{\sim3}$.

\begin{lemma}[The box datum]\label{htc:box-datum}
For $M(\alpha_0+\alpha_1)$ below a fixed ceiling, there are
holomorphic maps at fixed real $U$, denoted $V$ and $v$, with
\begin{equation}\label{htc:box-data}
 M_{\mathbf B^+}(\mathbf U)=V^v,\qquad
 \|\log V\|_\infty\le C_\square M\alpha_0,
          \qquad |v(x)|\,|v(x)^{-1}|\le e^{C_\square M\alpha_1}.
\end{equation}
They use only the bare field on $\square_0$, agree under restriction
from a larger domain, and are real at real potentials. The node
gauge is normalized at the prescribed base point.
The construction and these bounds also hold when all inward rows
of the determining set are retained. At fixed real base and fixed
radii, the domain defined by the simple conditions and the upper
tower clauses is open. Both maps commute with translations of the
aligned box and with its injective lifts between tori. Under a real
gauge $h$, their values become
$h(y_0)Vh(y_0)^{-1}$ and $h(x)v(x)h(y_0)^{-1}$, respectively.
\end{lemma}
\begin{proof}
Choose the prescribed comb from the base point to each site and
write $g(x)$ for its transport in $\mathbf U$. The comb has length
at most $C M$ in level-$N$ units. Along it the real transports have
norm one, while a complex increment has norm and inverse norm at
most $e^{L^{-N}|A'|}$. Consequently
$\operatorname{cn}(g)\le e^{CM\alpha_1}$.
In the comb gauge, a bare link is the product of the plaquettes in
the strip between the two combs. There are at most $CM L^N$ such
plaquettes. Their transports have the preceding condition-number
bound, and their defects are at most $\alpha_0L^{-2N}$.
The product bound therefore gives
$|\mathbf U^g(b)-1|\le CM\alpha_0 L^{-N}$, after reducing the
fixed ceiling. This argument uses the complex curvature clause as
well as the real base and the potential bound.

Write $\mathbf U^g=e^{ia}$ on this chart. Every order-$q$ determining
read has $L^q\sup|a|\le CM\alpha_0$, since $q\le N$.
The actual near-trivial averaging estimate therefore gives
$V=M_{\mathbf B^+}(\mathbf U^g)$ and
$\|\log V\|_\infty\le C_\square M\alpha_0$.
Complex covariance of each mean is valid on this same chart: the
condition number above times its input diameter remains below the
mean radius. Set $v=g^{-1}$ at the nodes. This proves the exact
identity in \eqref{htc:box-data}. Paths, products, logarithms and
means are holomorphic on the stated domain. Their definitions prove
the restriction, reality and base-point assertions.
Every determining row, including an inward row if it is retained,
has its full read inside the box. The same construction therefore
proves the asserted variant. All simple inequalities are strict
inequalities between continuous functions on finitely many cells.
For a fixed real base, each admitted tower clause uses a fixed
holomorphic datum chart; its defining raw-log inequalities and its
curvature and current inequalities are also strict and continuous.
There are only finitely many truncations and cells on this box.
Their intersection is open, including when it is empty because a
real-base clause fails. This is a finite-dimensional openness
argument; it makes no assertion of a uniform margin over increasing
cutoffs. The common bounds used for parameter discs were established
separately in \cref{htc:parameter-core}.
Finally the prescribed comb transports obey
$g^h(x)=h(y_0)g(x)h(x)^{-1}$. This gives the displayed gauge laws;
their paths and all averaging reads are preserved by the indicated
translations and injective lifts, which proves the remaining claims.
\end{proof}

\begin{lemma}[Restriction of a constrained critical field]
\label{htc:critical-restriction}
Let $\mathbf B$ be the determining set of a separated sequence
$\Omega_0\supset\cdots\supset\Omega_m$, with all inward rows
removed. Let $Z\subset\Omega_m$, and let $\mathbf B'$ be the
determining set of a separated sequence
$Z=\Omega'_0\supset\cdots\supset\Omega'_{m'}$, $m'\le m$,
also with all inward rows removed. A critical field for
$\mathbf B$ restricts to a critical field for $\mathbf B'$ at
its induced datum. The same assertion holds when $\mathbf B'$
is the determining set of a truncation of the original sequence.
\end{lemma}
\begin{proof}
Write $W$ for the old field and extend a competing field $W'$ on
$Z$ by $W$ outside the inner bare bonds of $Z$, obtaining
$\widetilde W'$. Suppose $M_{\mathbf B'}W'=M_{\mathbf B'}W$.
We first show $M_{\mathbf B}\widetilde W'=M_{\mathbf B}W$.
Put $\Omega'_{m'+1}=\varnothing$. For $0\le q\le m'$, let
$P(q)$ assert equality of the two order-$q$ averages at every
row whose full blocks lie in $\Omega_0$ and whose endpoints both
avoid $\Omega'_{q+1}$.

For $q=0$, a changed bare bond satisfying this condition belongs
to the frozen set $\Gamma'_0$, so $P(0)$ holds. Suppose $P(q-1)$
holds. If a row in $P(q)$ has an endpoint in $\Omega'_q$, this
endpoint is in $\Omega'_q\setminus\Omega'_{q+1}$. Separation
puts its other endpoint and its complete read inside
$\Omega'_{q-1}$, and the row belongs to $\Gamma'_q$, including
when it points outward. Its average is thus fixed by the induced
datum. Otherwise both endpoints are outside $\Omega'_q$.
Since $\Omega'_q$ is a union of aligned level-$q$ cubes, both
predecessor blocks avoid $\Omega'_q$. The one-step mean reads
only these blocks, where $P(q-1)$ applies. This proves $P(q)$.
At $q=m'$ all order-$m'$ averages agree. The averaging semigroup
therefore gives equality at every higher-order old row as well.
An old row of order $q<m'$ has an endpoint outside
$\Omega_{q+1}\supset Z$. Its length is $L^q$, whereas
$\Omega'_{q+1}$ has distance greater than $L^q$ from $Z^c$.
Neither endpoint can lie in $\Omega'_{q+1}$, so $P(q)$ also
covers every such old row.

Every bare plaquette meeting $Z^c$ has each corner within one
bare step of $Z^c$. Separation therefore puts any of its inner
bonds in $\Gamma'_0$, where $W'=W$. All action terms crossing
the boundary are constant along a constrained variation of $W'$.
Extending that variation consequently preserves the old datum,
and its old action derivative equals its regional action derivative.
The first is zero by criticality, which proves the assertion.

For a truncation at depth $t$, the retained rows below $t$ are
unchanged. Its top rows include every order-$t$ row needed in the
read of an old row of order greater than $t$; this follows from
the separation collar and the nested block identity. Equality
of the top averages and the semigroup preserve these old rows.
The underlying region is unchanged, so there is no crossing
action term. The same variation argument applies.
\end{proof}

\begin{proposition}[Inner substitution on its full domain]
\label{htc:inner-substitution}
There are fixed $M_0,C_I,c_2,c_3>0$ and, for $M\ge M_0$, a ceiling
$a_I(M)>0$ with the following property. Let
$A_0,A_1,a_0,a_1>0$, and suppose
\begin{equation}\label{htc:inner-radii}
 A_0+A_1<a_I(M),\quad a_0+a_1<a_I(M),\quad
 C_IMa_0<a_1,\quad A_0\le2\sqrt{1+\nu}\,a_0.
\end{equation}
Put $C_V=2C_\square$. First allow an independent complex raw
datum $b_V$ on $\mathbf B^+$ in the ball
\begin{equation}\label{htc:upper-datum-ball}
                         \|b_V\|_\infty<C_VMA_0,
\end{equation}
and set $H_+=\mathcal H_{\mathbf B^+,1}(b_V)$.
For the physical application, take the upper input and the data
of \cref{htc:box-datum} and put $b_V=i^{-1}\log V$.
Equation~\eqref{htc:box-data} places this datum in
\eqref{htc:upper-datum-ball}, since
$\alpha_0\le3A_0/2$.
For a complex field $A$ in level-$k$ units and a raw row correction
$b'$ on $\mathbf B_j$, suppose
$\|A\|_\infty<c_2A_0$ and $\|b'\|_\infty<c_3a_0$.
Define the actual raw fine potential and row datum by
\begin{equation}\label{htc:inner-map}
 a=(i)^{-1}\log(e^{iL^{-k}A}e^{iL^{-N}H_+}),\qquad
 b(c)=\mathsf q_q(a)(c)\quad(c\in\Gamma_q),\qquad
 H=\mathcal H_{\mathbf B_j,1}(\tau b+b'),\quad0\le\tau\le1.
\end{equation}
The application of $\mathcal H$ in this formula includes
$\mathfrak n_{\mathbf B_j,1}(\tau b+b')$ before the reduced saddle
equations are solved. All maps are defined and holomorphic on the full indicated open
balls in $(b_V,A,b')$. Whenever $b_V$ comes from an upper input
as above, their output on $\square^{\sim3}$ satisfies
\begin{equation}\label{htc:inner-readouts}
 |H|,|\nabla H|<a_1/2,\qquad
 |\partial e^{iL^{-j}H}-1|<a_0L^{-2j}/8,\qquad
 |J^{(j)}(e^{iL^{-j}H})|<a_0/16.
\end{equation}
For every $X\subset\square^{\sim2}$ of level $j-1$, the triple
$(1,H,J^{(j)}(e^{iL^{-j}H}))|_X$ lies in the core with radii
$(a_0/4,a_1/2,a_0/4)$. These estimates and inclusions hold throughout
every auxiliary parameter polydisc on which the upper input and
the two increment bounds hold.
The same inclusion holds on a larger level-$(j-1)$ domain whose
complete simple-condition stencils lie in $\square^{\sim3}$ and
which is contained in the top member of the truncated window tower.
\end{proposition}
\begin{proof}
We give the choices and the estimates in their dependency order.
The physical datum chart and \eqref{pnc:weighted-first-response}
give constants $C,r,s>0$ for the flat reference on every tower.
Its potential, first difference, normalized curvature and physical
current have Lipschitz constant $C$ on $\|d\|_\infty<r$.
If a change of $d$ is supported at distance at least $D$ from an
output cell, its contribution is at most $Ce^{-sD}$ times its
supremum. For the current this follows by differentiating
\eqref{htc:Euler-current}; the weighted row-transpose density sums
the differentiated multiplier columns. It therefore retains the
same exponential bound. The output cell and the full current
stencil are included in the distance convention.

We first admit the full product of input balls, without requiring
$b_V$ to arise from an upper field. On a complete order-$q$ read of
the truncated tower, the upper assigned length is at least
$L^{q-1-N}$, by the geometry used in \eqref{htc:mean-ball}.
The weighted flat response therefore gives
\[
 L^q\sup_{\mathsf P^q(c)}|L^{-N}H_+|
       \le C\|b_V\|_\infty,\qquad c\in\Gamma_q,\quad q\le j.
\]
The product logarithm and the actual mean on this read consequently
give, uniformly for $0\le\tau\le1$,
\[
 \|\tau b+b'\|_\infty
 \le C\bigl(\|b_V\|_\infty+L^{j-k}\|A\|_\infty\bigr)
                                      +\|b'\|_\infty.
\]
Since $j\le k$, decreasing $a_I(M)$ makes the right side smaller
than $r/2$ on the entire product of the three stated balls,
and keeps all preceding raw logarithms and means in their fixed
charts. This proves the full-domain holomorphy. Only the sharper
readout estimates below use the actual upper input.

For any one of these readouts $R_x(d)$, with $R_x(0)=0$, Cauchy on
the fixed datum ball gives
\begin{equation}\label{htc:local-Taylor}
 |R_x(d)-DR_x(0)d|\le C\|d\|_\infty^2/r.
\end{equation}
For a split $d=d_{\rm in}+d_{\rm out}$, where the outer part is
at distance $D$, apply this bound only to $d_{\rm in}$ and use
the preceding Lipschitz estimate for $d_{\rm out}$. This gives
\begin{equation}\label{htc:near-far-Taylor}
 |R_x(\tau d)-\tau R_x(d)|
 \le C\bigl(e^{-sD}\|d_{\rm out}\|_\infty
                         +\|d_{\rm in}\|_\infty^2/r\bigr),
                         \qquad 0\le\tau\le1.
\end{equation}
These are estimates of the actual current map, so no bound on
second differences of an arbitrary complex input is being assumed.

Let $q_0=\log M_{\mathbf B_j}(e^{iL^{-N}H_+})/i$.
The uniform row bounds give
\[
 \|q_0\|_\infty\le C M A_0,
 \qquad \|q_{\rm in}\|_\infty\le C M A_0\sigma,
 \qquad \sup_{\square^{\sim4}}|H_+|\le C M A_0.
\]
Here an inner row is one whose complete double block is in
$\square^{\sim4}$. Such a row is a level-$j$ row, by
\cref{htc:geometry}; its raw input has the extra factor
$L^{j-N}=\sigma$. Every other row is at distance
$D\ge c M L^{\nu-1}$ from $\square^{\sim3}$ in the level-$j$
metric. This follows by stopping a path at the edge of
$\square^{\sim4}$; before that edge all pillar lengths are one
in level-$j$ units.

Choose $M_0$ so, for every $M\ge M_0$ and $\nu\ge1$,
\begin{equation}\label{htc:inner-small-choices}
 C M e^{-scM L^{\nu-1}}\le10^{-3}L^{-2\nu}.
\end{equation}
Such a choice exists: with $t=L^{\nu-1}\ge1$, the quantity
$t^2e^{-scMt}$ decreases for $scM\ge2$.
Next choose $a_I(M)$ so $CMa_I<r/8$ and
$CM^2a_I/r<10^{-3}$, as well as all preceding logarithm and
gauge ceilings. Combining \eqref{htc:local-Taylor} and
\eqref{htc:near-far-Taylor}
then gives, for curvature and current,
\begin{equation}\label{htc:inner-ray-error}
 |R_x(\tau q_0)-\tau R_x(q_0)|
                         \le A_0\sigma^2/100.
\end{equation}
For the potential and first difference they give the simpler bound
$CMA_0\sigma$.

We identify the endpoint $\tau=1$. Restriction of the upper
critical field to the truncated sequence is critical at its full
induced datum by \cref{htc:critical-restriction}. Real regular uniqueness
identifies the restricted orbit. The equality extends to the
common complex datum ball by holomorphy. Restoring the upper node
gauge and both axial normalizations gives
\[
 e^{iL^{-j}\mathcal H_{\mathbf B_j,1}(q_0)}
                         =(U_N(\square_0;\mathbf U))^G.
\]
The local products in \eqref{htc:pillar-gauge}, the corresponding
product for the flat nested tower, and \eqref{htc:box-data} give
$\operatorname{cn}(G)\le\exp(CM(A_0+A_1))\le4/3$ after decreasing
$a_I(M)$. We check the real-base change in this identification.
Write $V_0=M_{\mathbf B^+}(U)$ and
$e^{ib^U}V_0=M_{\mathbf B^+}(\mathbf U)$; the mean estimate gives
$\|b^U\|_\infty\le C\alpha_1$. At the point in question use the
polar factorization $v=v_r e^{i\mu}$, with $v_r$ unitary and
$\|\mu\|_\infty\le CM\alpha_1$. Conjugating the exact datum
identity by $v_r^{-1}$ gives
\[
 e^{i\widetilde b^U(c)}V'_0(c)
       =e^{i\mu(c_-)}e^{ib_V(c)}e^{-i\mu(c_+)},
 \quad V'_0=V_0^{v_r^{-1}},\quad
 \widetilde b^U=\operatorname{Ad}(v_r^{-1})b^U.
\]
Thus $V'_0=e^{ib_0}$ with $b_0$ real and
$\|b_0\|_\infty\le CM(A_0+A_1)$.
The four logarithms $b_0,\widetilde b^U,\mu,b_V$ lie in a fixed
smaller datum ball. At real values, global regular uniqueness gives
both change of base and covariance under the node gauge. The two
sides are holomorphic functions of the small logarithms on that
ball, so the identity theorem gives these identities for their
complex values. The polar factorization is used only to locate
the given point in the common charts; no holomorphic dependence
of its unitary factor is claimed. This proves the required base
change and the endpoint identification on the actual domain.
The upper tower clause and $\alpha_0\le3A_0/2$ now imply
\[
 L^{2j}|\partial e^{iL^{-j}\mathcal H_j(q_0)}-1|
                  \le2A_0\sigma^2,\qquad
 |J^{(j)}(e^{iL^{-j}\mathcal H_j(q_0)})|
                  \le2A_0\sigma^3.
\]

On their fixed common logarithm chart the actual means obey
\begin{equation}\label{htc:inner-row-difference}
 |\mathsf q_q(a)(c)-\mathsf q_q(a_{\rm ref})(c)|
       \le C L^q\sup_{\mathsf P^q(c)}|a-a_{\rm ref}|.
\end{equation}
Here $a_{\rm ref}=L^{-N}H_+$. Apply Cauchy to the segment
$a_{\rm ref}+z(a-a_{\rm ref})$ on the radius
$c_{\rm av}/(2L^q\sup|a-a_{\rm ref}|)$, where the mean is bounded by
$C_Qc_{\rm av}$; Schwarz applied after subtracting its value at
zero gives the displayed inequality. The product logarithm in
\eqref{htc:inner-map} satisfies
$|a-a_{\rm ref}|\le C L^{-k}|A|$.
Consequently $\|b-q_0\|_\infty\le C L^{j-k}\|A\|_\infty$.
Choose $c_2,c_3$ so
\[
 C\|\tau(b-q_0)+b'\|_\infty<a_0/128
\]
in every readout norm above. This is a fixed choice, because
$2\sqrt{1+\nu}L^{-(\nu-1)}\le2\sqrt2$.
It also puts the entire datum in $\|d\|_\infty<r/2$, and so proves
the asserted holomorphy on the full product of increment balls.
The resulting curvature and current bounds are, respectively,
\[
 (2+1/100)A_0\sigma^2+a_0/128,
 \qquad 2A_0\sigma^3+A_0\sigma^2/100+a_0/128.
\]
Since $2\sqrt{1+\nu}L^{-2\nu}<0.0168$ and
$2\sqrt{1+\nu}L^{-3\nu}<0.0013$, these are smaller than
$a_0/8$ and $a_0/16$. For the potential and first difference the
bound is $CMA_0\sigma+a_0/128$. Choose $C_I$ large enough that
\eqref{htc:inner-radii}, together with
$2\sqrt{1+\nu}L^{-\nu}<0.218$, makes it smaller than $a_1/2$.
This proves \eqref{htc:inner-readouts}.

It remains to check every nested tower, without paying the large
constant $K_T$ against this fixed fractional radius. Let $(p,Z)$ be
any address of $X$, use either rule, and truncate at $n\le p$.
The whole $Z$ is inside the top member of the level-$j$ window
tower. The constraint and boundary-action checks in
\cref{htc:critical-restriction} identify its normalized
critical continuation with $(e^{iL^{-j}H}|_Z)^w$.
All its mean logarithms have size at most $C_Q\sup_Z|H|$ and are
inside the common datum ball by the choice of $a_I$. Its normalization
gauge is the product of its free ascent contours. At level $q$ their
logarithms are bounded by $CL^{q-j}\sup_Z|H|$; summing $q<n$ and
including the top node factor yields
$\operatorname{cn}(w)\le e^{C a_1}<6/5$.
Thus the identity with its bound is
\begin{equation}\label{htc:nested-identity}
 U_n(Z;e^{iL^{-j}H})=(e^{iL^{-j}H}|_Z)^w,
                         \qquad\operatorname{cn}(w)<6/5.
\end{equation}
For real window data this is the residual-gauge identity from
regular uniqueness. Holomorphy on the common product ball proves
it for all the present complex parameters.

By \cref{htc:geometry}, the full stencils of every inset lie in $Z$.
Conjugation and the change from level $j$ to level $n$ therefore give
\[
 |\partial U_n-1|<\frac65\frac{a_0}{8}L^{-2j}
             <\frac{a_0}{4}L^{-(j-p)}L^{-2p},
\]
\[
 |J^{(n)}(U_n)|<\frac65\frac{a_0}{16}L^{-3(j-n)}
             <\frac{a_0}{4}L^{-(j-p)}L^{2(n-p)}.
\]
The second comparison uses $2j-p-n\ge0$; the first uses $p\le j$.
The real factor is $1$, whose towers are identically $1$. Its gauge
clause holds with zero potential. The independent current chosen in
the statement is the physical current already bounded in
\eqref{htc:inner-readouts}. Thus every core clause holds on the
whole parameter domain, including every nested tower and both rules.
The proof used $X\subset\square^{\sim2}$ only to put its complete
simple stencils in $\square^{\sim3}$ and its nested regions in
the top member. This proves the last assertion as well.
\end{proof}



\subsection{Tower response in the ambient metric}
\label{atc:section}

We now prove the weighted tower comparison used at settled elements
of a deeper sequence. The determining tower is fixed during this
comparison. The ambient sequence specifies the input norm and the
output weights; it does not change the constrained problem.

Let $1\le p\le i$. Take separated sequences
$(\Omega'_q)_{q\le p}$ and $(\Omega^\flat_q)_{q\le i}$, with
$\Omega'_q\subset\Omega^\flat_q$ for $q\le p$, and put
$Z=\Omega'_0$. The determining set $\mathbf B'$ has all inward
rows removed. All rows, outward endpoint blocks and physical
pairings retain their actual conventions. Work in level-$p$ units,
so $\eta=L^{-p}$ and $\xi_q=L^{q-p}$.
Assume each torus period contains at least $6\rho$ level-$p$
sites, where $\rho=R_1\mathfrak M$. All volume-uniform assertions
below refer to this admitted class. Sup distances between torus
points mean their shortest periodic sup distance.
At a bare site $x$, let $\lambda'(x)$ and $\lambda^\flat(x)$ be
the largest indices of the two sequences containing $x$, with
value zero outside their first members, and set
\[
 \xi'(x)=\xi_{\lambda'(x)},\qquad
 l(x)=\xi_{\lambda^\flat(x)},\qquad
 s(x)=l(x)/\xi'(x)\ge1.
\]
Write $d'$ and $d^\flat$ for the corresponding weighted path
metrics, with density the reciprocal of the assigned scale.
For a real field $U$ and a complex potential $v$ on the inner
bare bonds of $Z$, define
\begin{equation}\label{atc:input}
 \epsilon_U(P)=\eta^{-2}l(x_P)^2|\partial U(P)-1|,\qquad
 n_b(v)=\max\{l(x_b)|v(b)|,\ l(x_b)^2|\nabla_Uv|_Z(b)\}.
\end{equation}
Only differences whose entire three-bond read lies in $Z$ enter
the last term. Put $n(v)=\max_b n_b(v)$ and
\[
 E_U^\delta(x)=\max_P e^{-\delta d^\flat(x,x_P)}\epsilon_U(P),
 \qquad E_v^\delta(x)=\max_b e^{-\delta d^\flat(x,x_b)}n_b(v).
\]
The maxima use positively oriented inner plaquettes and bonds.
An output element is an inner plaquette or a full bond, the latter
meaning that all six incident plaquettes are inner. Its base $x$
is settled if $\lambda'(x)=\min\{\lambda^\flat(x),p\}$.
For the normalized critical continuation
$\Phi_U(v)=\Phi_{\mathbf B',M_{\mathbf B'}U}(b(v))$, with the
literal raw datum $b(v)$ of \eqref{htc:mean-ball}, define
\begin{equation}\label{atc:output}
 \mathcal G_U(v;P)=\eta^{-2}l(x_P)^2(\partial\Phi_U(v)(P)-1),
 \qquad
 \mathcal G_U(v;b)=l(x_b)^2\xi'(x_b)J_\eta(\Phi_U(v))(b).
\end{equation}
In particular, the current weight is $l^2\xi'$. These are the
weights of \eqref{mb:ambient-tower-weights}.

\begin{lemma}[The two metrics and local scales]\label{atc:geometry}
Let $\rho=R_1\mathfrak M$ be the separation in layer units.
For bare sites $x,y$ and for a row $c$ of level $q$, with a bond
$b$ in its complete read, one has
\begin{equation}\label{atc:metric-bounds}
 \begin{gathered}
 d'(x,y)\ge d^\flat(x,y),\qquad
 d'(x,y)\ge |x-y|_\infty/\xi_p,\\
 s(x)/s(y)\le L^{2+2d'(x,y)/\rho},\qquad
 d'(y_c,x_b)\le C_L,\qquad \xi_q\le L\xi'(x_b).
 \end{gathered}
\end{equation}
If $|x-y|_\infty\le3l(x)$, then
$L^{-1}l(x)\le l(y)\le Ll(x)$ and
$d^\flat(x,y)\le C_L$. If $x$ is settled and
$|x-y|_\infty>l(x)$, then $s(x)\le1+d'(x,y)$.
\end{lemma}
\begin{proof}
Containment of the sequences gives $\xi'\le l$ at every
point of each bare path. Also $\xi'\le\xi_p$. Integrating
these two pointwise inequalities and taking infima proves
the first two statements.
To pass between layers with indices $a<a'$, a connected path
must cross each intermediate collar. The collar at level $q$
has width greater than $\rho\xi_q$ and path density
$\xi_q^{-1}$, so each intermediate layer costs at least
$\rho$. Thus the change of either index is at most one plus
the corresponding weighted distance divided by $\rho$.
Apply this to both indices and use $d^\flat\le d'$ to bound
their difference, giving the ratio of $s$ in the display.

The complete order-$q$ read lies in the preceding member;
its endpoint blocks have diameter $O(\xi_q)$ and every site
there has scale at least $\xi_{q-1}$. A path through these
blocks has weighted length $C_L$. This includes the lower
block at an outward endpoint and proves both row-read bounds.
A ball of radius $3l(x)$ cannot cross two ambient layers,
since the intervening collar is wider than that ball once
$\rho\ge6L$. Its scales differ by at most $L$, and coordinate
paths in that ball give the local distance estimate.
Finally, at a settled site either $s(x)=1$ or
$\lambda'(x)=p$. In the latter case, a far point satisfies
$d'(x,y)\ge|x-y|_\infty/\xi_p>l(x)/\xi_p=s(x)$.
\end{proof}

\begin{lemma}[Physical readouts on the fixed datum ball]
\label{atc:datum-response}
There are fixed $r,C,m>0$ such that, at every sufficiently
regular real datum $V_0$, the following holds on $\|b\|_\infty<r$.
At any output element with base $x$, define $F_x(b)$ as the
curvature or current readout normalized by $\xi'(x)$ rather
than $l(x)$ in \eqref{atc:output}. Then $F_x$ is holomorphic,
$|F_x(b)|\le C$, and
\begin{equation}\label{atc:datum-Lipschitz}
 |F_x(b_1)-F_x(b_0)|
 \le C\max_c e^{-m d'(x,y_c)}|b_1(c)-b_0(c)|.
\end{equation}
The constants cover full current stencils across layer seams.
\end{lemma}
\begin{proof}
Use \cref{bgc:uniform-physical-chart} on a smaller fixed raw
datum ball. The naturalizer and literal node extension are
included, as specified before \cref{htc:common-ball}.
Apply \eqref{pnc:weighted-first-response} with the cell weight
$e^{-m d'(x,\cdot)}$. It is an admitted Lipschitz weight.
The potential, first difference and multiplier derivatives
therefore have the right-hand side of
\eqref{atc:datum-Lipschitz} in their physical local norms.

For the current differentiate the literal constrained Euler
identity \eqref{htc:Euler-current}. The terms differentiate
the local bond factor, the actual constraint transpose and
the multiplier. The transpose retains its physical row density
and the full outward endpoint read. Its weighted row sum,
and the corresponding sum of its differentiated coefficients,
are bounded by the actual carrier estimates. Thus the current
derivative has the same bound in grade three. This step uses
the equation of the continued critical field; it does not
bound second differences of an arbitrary $C^1$ potential.
Plaquette products give the curvature derivative in grade two.

The representative's gauge must also be differentiated.
On each free ascent its factors have logarithmic differences
bounded by $CL^{q-t}$ times the weighted local difference of
the Landau fields, as in \eqref{htc:pillar-gauge}.
Each factor's complete read lies within distance $C_L$ of
the pillar top, and the site being normalized has distance
at most $2$ from that top. The geometric sum over $q<t$ is
bounded. The top node factor has the same local bound from
the site-mean inverse. Real reference factors are unitary;
telescoping relative factors preserves the bounds. Consequently
the gauge derivative and its inverse have the same weighted
bound. Differentiating conjugation of the physical readout
adds its already bounded value times this gauge derivative.

Every corner in a full stencil is at most two bare steps
from $x$. Its assigned scale differs from $\xi'(x)$ by at
most $L$. Absorb these fixed grade-two and grade-three factors
and the bounded metric displacement in $C$. This treats seams
without omitting any plaquette. At a nonfull boundary bond no
full-current claim is made; the regional Euler equation still
retains its boundary terms. Integrating the derivative along
the segment from $b_0$ to $b_1$ proves the display.
\end{proof}

\begin{proposition}[Ambient tower bound and response]\label{atc:theorem}
There are fixed $a,K_a,K_T,\delta>0$ such that the following
holds, uniformly in both sequences, their depths and volume.
Suppose $\max_P\epsilon_U(P)<a$ and $n(w)<a$.
The actual tower configuration and its full current exist.
For fixed $U$ they are holomorphic on this open convex set
of potentials. At each settled output element with base $x$,
\begin{equation}\label{atc:apriori}
 |\mathcal G_U(w)|\le K_a(E_U^\delta(x)+E_w^\delta(x))\le1.
\end{equation}
For $n(w_0),n(w_1)<a$,
\begin{equation}\label{atc:Lipschitz}
 |\mathcal G_U(w_1)-\mathcal G_U(w_0)|
                 \le K_T E_{w_1-w_0}^\delta(x).
\end{equation}
These statements include both erosion rules and their truncations
whenever they are contained in the indicated ambient sequence.
\end{proposition}
\begin{proof}
Fix a common smaller rate $m$ for \cref{atc:datum-response}
and the real local profile proved in \cref{bgc:reset}.
These primitive estimates are uniform above their fixed
geometric floors. Increase the separation floor so that
$4\log L/\rho\le m/16$, and put $\delta=m/8$.
Choose all amplitude ceilings after these choices.

The input norm $\|w\|_-$ of \eqref{htc:input-norm} is at
most $n(w)$. The averaged real datum is regular by the same
full-read argument as \eqref{htc:mean-ball}. That equation
and \cref{bgc:proper-background,bgc:uniform-physical-chart}
therefore define the tower holomorphically on a fixed ball
$\|w\|_-<r_w$. Its current is a polynomial in its bond
matrices and inverses on a full stencil. The conditions
$n(w)<a$ are finitely many continuous linear seminorm bounds,
so their intersection is open and convex.

We first need a spatial real estimate. If $W$ is real with
$\sup\epsilon_W$ small, the actual mean on a tested mixed
plaquette $\tau$ of level $q$ has defect bounded by
\begin{equation}\label{atc:real-datum-profile}
 |\partial M_{\mathbf B'}W(\tau)-1|
 \le C\max_{P\text{ in its complete read}}
                         s(x_P)^{-2}\epsilon_W(P).
\end{equation}
Indeed the averaging estimate costs $\xi_q^2$ times the
physical fine curvature, and $\xi_q\le L\xi'(x_P)$ on
that read. At a mixed seam the reconstructed lower row is
the actual next average by the semigroup identity, so the
same calculation includes the outward endpoint block.
At level zero this is the prescribed plaquette itself.
The read has $d'$-diameter $C_L$.
In detail, the local real profile obtained from the proof
of \cref{bgc:reset} is
\begin{equation}\label{atc:critical-profile}
 r_{U(V)}(x)\le C\max_\tau
                 e^{-m d'(x,\tau)}|\partial V(\tau)-1|.
\end{equation}
Here $r$ is the maximum of the grade-two curvature and
grade-three full-current norms at $x$ in the tower scales,
and $\tau$ ranges over the actual ordinary and mixed datum
tests. To recall the implication, that proof gives
$r(x)\le K v_R(x)+\frac12\sup_{d'(x,y)\le R}r(y)$.
For $P(x)=\max_\tau e^{-m d'(x,\tau)}|\partial V(\tau)-1|$
choose $m$ small enough that $e^{mR}\le\sqrt2$.
Maximizing $r(x)/(P(x)+t)$, $t>0$, gives
$S\le\sqrt2K+S/\sqrt2$, hence $r\le5KP$ as $t\downarrow0$.
This proves the displayed profile with the stated metric
and tested datum, including its local full-current readout.
Multiply that profile
by $s(x)^2$ and use \cref{atc:geometry} in
\eqref{atc:real-datum-profile}. The factor
$s(x)^2/s(x_P)^2$ costs at most
$L^4e^{4(\log L)d'(x,x_P)/\rho}$, which the reserved
decay absorbs. Hence the real critical field at $M_{\mathbf B'}W$
has ambient readout at $x$ bounded by $C E_W^\delta(x)$.

Write $F(w)=\mathcal G_U(w;x)$ and $s=s(x)$.
For a real $g$ the plaquette product expansion gives
\[
 \epsilon_{e^{i\eta g}U}(P)
       \le2\epsilon_U(P)+C\max_{b\subset P}n_b(g),
\]
when $n(g)$ is small. The quadratic remainder is bounded by
$Cn(g)^2$ and is absorbed into the displayed linear term.
Real global uniqueness matches the critical field just
estimated to the real point of the fixed chart. Consequently
\begin{equation}\label{atc:real-readout}
 |F(g)|\le C(E_U^\delta(x)+n(g)),\qquad
                         |F(0)|\le C E_U^\delta(x).
\end{equation}
On the whole complex ball $\|w\|_-<r_w$, the crude bound
from \cref{atc:datum-response} is $|F(w)|\le Cs^2$.

Choose a scalar cutoff on positive bonds, equal to one for
$|x_b-x|_\infty\le l(x)$, equal to zero beyond $2l(x)$,
and linear in between. Denote it by $\chi(b)$ and extend
the resulting form by its reversal convention. Call a bond
near if $|x_b-x|_\infty\le3l(x)$, and write
$n_{\rm near}(v)=\max_{b\text{ near}}n_b(v)$.
The discrete product rule and \cref{atc:geometry} give
\begin{equation}\label{atc:cutoff}
 n(\chi v)\le C_\chi n_{\rm near}(v),\qquad
 \|\chi v\|_-\le C_Ls^{-1}n(\chi v),\qquad
 n_{\rm near}(v)\le e^{C_L\delta}E_v^\delta(x).
\end{equation}
For the first inequality the cutoff difference is at most
$\eta/l(x)$; adjacent ambient scales differ by at most $L$.
For the second, if $s>1$ the output is in the tower's top
layer and every tower scale is at most $\xi_p=l(x)/s$;
if $s=1$, use $\ell_-\le l$. These arguments also hold
when a difference stencil meets the edge of the support.

Apply \cref{htc:three-real} to $\chi w$ on the fixed
$\|\cdot\|_-$ ball. Its remainder is at most
$Cs^2\|\chi w\|_-^2/r_w^2\le Cn(\chi w)^2/r_w^2$.
Together with \eqref{atc:real-readout}, this proves
\begin{equation}\label{atc:near}
 |F(\chi w)|\le C(E_U^\delta(x)+n_{\rm near}(w))
\end{equation}
on a smaller fixed amplitude ball. This is the two-power
gain for near forcing; it uses no positive travel distance.

For far forcing let $u=w-\widetilde w$ vanish on
$|x_b-x|_\infty\le l(x)$. Cauchy on the actual averaging
chart gives, row by row,
\[
 |b(w)(c)-b(\widetilde w)(c)|
       \le C\xi_q\sup_{b\in\mathsf P^q(c)}|u(b)|
       \le C\max_{b\in\mathsf P^q(c)}s(x_b)^{-1}l(x_b)|u(b)|.
\]
Put $D=d'(x,x_b)$. A contributing bond is far, so
$s\le1+D$. The scale ratio and row-read distance give
\begin{equation}\label{atc:far-cancellation}
 s^2e^{-m d'(x,y_c)}s(x_b)^{-1}
 \le C(1+D)e^{-(m-2\log L/\rho)D}
 \le Ce^{-2\delta D}.
\end{equation}
Apply \eqref{atc:datum-Lipschitz}, remembering that
$F=s^2F_x$, and then use $d'\ge d^\flat$. We obtain
\begin{equation}\label{atc:far}
 |F(w)-F(\widetilde w)|\le C E_u^\delta(x).
\end{equation}
The bound is uniform on a common smaller fixed ball.
Combining it with \eqref{atc:near}, with
$\widetilde w=\chi w$, proves \eqref{atc:apriori} on
$n(w)<\bar a$ after fixing $K_a$ and reducing $\bar a$.

For the Lipschitz assertion put $v=w_1-w_0$ and split
the difference through $w_0+\chi v$. The far bracket is
bounded by \eqref{atc:far}. Put
$n_v=C_\chi n_{\rm near}(v)$ and
$E=E_U^\delta(x)+E_{w_0}^\delta(x)$.
If $n_v=0$, the near bracket vanishes. Otherwise set
$R=1+E/n_v$. For $|z|\le R$,
\[
 n(w_0+z\chi v)\le a+n_v+E\le(3+2C_\chi)a.
\]
Choose $a<\bar a/(4+2C_\chi)$. The a priori estimate
applies on a neighborhood of this disc and bounds
$F(w_0+z\chi v)$ by $K_a(2E+n_v)$. Schwarz's lemma,
after subtracting the value at zero and testing by norm-one
linear functionals, bounds the near bracket by
\[
 \frac{2K_a(2E+n_v)}{1+E/n_v}\le4K_a n_v
                                      \le C E_v^\delta(x).
\]
Together with the far bracket this proves
\eqref{atc:Lipschitz}. Finally choose $2K_a a\le1$.
No constant counts cells, addresses or truncations.
\end{proof}

\begin{corollary}[Covariance and reality]\label{atc:covariance}
For a real gauge $u$ on $Z$, let $\widetilde u$ equal $u(y)$
throughout each pillar with top $y$. Then the domain and bounds
above are unchanged by $(U,w)\mapsto(U^u,\operatorname{Ad}(u)w)$,
and
\[
 \Phi_{U^u}(\operatorname{Ad}(u)w)=\Phi_U(w)^{\widetilde u}.
\]
The output readout is conjugated by $\widetilde u(x)$.
For $w^*$ obtained by taking the adjoint at positive bonds,
\[
 \Phi_U(w^*)=\Phi_U(w)^{\mathfrak c},\qquad
 \mathfrak c(X)=(X^{-1})^*,\qquad
 |\mathcal G_U(w^*)|=|\mathcal G_U(w)|.
\]
For Hermitian $w$ the tower is its real normalized critical
representative. Also $\Phi_1(0)=1$ and $\mathcal G_1(0)=0$.
\end{corollary}
\begin{proof}
Unitary covariance preserves the input plaquette and derivative
norms. The actual means transform by $u$ at their nodes.
The prescribed axial normalization extends this node gauge
constantly on each pillar, so real uniqueness gives the first
identity at real potentials; holomorphy gives it on the complex
domain. Plaquettes and the complete current transform by
conjugation at the stated base.
The anti-holomorphic involution $\mathfrak c$ fixes the real
group. The map $w\mapsto\mathfrak c(\Phi_U(w^*))$ is
holomorphic and agrees with $\Phi_U(w)$ on the real slice,
which proves the second identity. The current is changed to
its adjoint. For $X\in\operatorname{SL}(2,\mathbb C)$,
$X^{-1}-1=\operatorname{adj}(X-1)$; the two singular values
of the adjugate are interchanged. Thus
$|X^{\mathfrak c}-1|=|X-1|$ in the operator norm, proving
the curvature norm assertion as well. Finally the identity
configuration has identity averages, zero action derivative
and the prescribed normalization; real uniqueness identifies
it with $\Phi_1(0)$. Its current and curvature vanish.
\end{proof}

\begin{corollary}[Localizations and deeper top metrics]
\label{atc:localization}
Let the ambient sequence extend from depth $i$ to depth $t\ge i$,
with the same members through $i$. Write $d_t$ for its metric
and $y_t(x)$ for the top of the pillar containing $x$. The bounds
of \cref{atc:theorem} remain valid with
$d_t(y_t(x),y_t(x_b))$ in the input envelopes, after multiplying
their constants by $e^{4\delta}$.

For a region $X$ made of localization cubes at level zero,
let $X_{[q]}$ be its second-rule erosion and put
$\Omega'_q=\Omega^\flat_q\sqcap X_{[q]}$, $q\le p$.
Assume $M_1\ge\rho$ and $M\ge L^3M_1$.
Retain the preceding period floor at depth $p$; a period floor
at the deeper level $t$ implies it.
Define $X^{\wr p}$ as the union of the level-$(p-1)$
localization cubes whose two-layer enlargements lie in $X$.
The new sequence is separated. Every output element whose base
is within $2\xi_p$ of $X^{\wr p}$ is settled, provided its
full stencil belongs to $\Omega'_0$. Thus the same weighted
response holds for these actual localized towers.
\end{corollary}
\begin{proof}
Extending the ambient sequence only increases assigned scales,
so $d_t(x,z)\le d^\flat(x,z)$. Each site is within weighted
distance $2$ of its own top. The triangle inequality gives
\[
 d_t(y_t(x),y_t(z))\le d^\flat(x,z)+4,
\]
which proves the first assertion for both envelopes.

All intersections are taken cube-wise on nested aligned grids.
Their next-member gaps are bounded below by the minimum of the
two original gaps, so the localized sequence is separated.
For the second-rule erosion, every discarded point at stage $q$
is within distance
$2M_1\sum_{r=1}^q\xi_r<3M_1\xi_p$ of $X^c$.
This is the rejected-cube induction in \cref{htc:geometry},
which only requires that the initial region be a union of
the aligned smaller cubes. On the other hand,
$\operatorname{dist}_\infty(X^{\wr p},X^c)\ge2M\xi_{p-1}$.
A point within $2\xi_p$ of that inset therefore has distance
at least $(2M/L-2)\xi_p>3M_1\xi_p$ from $X^c$.
It lies in every $X_{[q]}$ through $p$. At such a bare site,
membership in $\Omega'_q$ is exactly membership in
$\Omega^\flat_q$, so the tower index is the ambient index
truncated at $p$. This is precisely the settled condition.
If the coarse partition is absent the inset is empty and
the assertion about its neighborhood is vacuous.
The assumption on the full stencil retains all physical
current terms at the localized boundary.
\end{proof}

\begin{corollary}[Multiplicative shifts in the ambient norm]
\label{atc:shift}
Let $(U,w)$ satisfy the hypotheses of \cref{atc:theorem}, and
put $\mathbf U=e^{i\eta w}U$. For a complex shift $A$, set
\[
 u_b=\max\{l(x_b)|A(b)|,\ l(x_b)^2|\nabla_{\mathbf U}A|_Z(b)\},
 \qquad w\boxplus A=(i\eta)^{-1}\log(e^{i\eta A}e^{i\eta w}).
\]
Let $u_b^+$ be the maximum over $b$ and its forward translates
that occur in a difference stencil in $Z$. After decreasing
the fixed ceiling $a$, the conditions
$u_b\le(256L^2)^{-1}$ and $n_b(w)+2u_b^+<a$ imply
\begin{equation}\label{atc:shift-bound}
 |\mathcal G_U(w\boxplus A)-\mathcal G_U(w)|
 \le2e^\delta K_T
             \max_b e^{-\delta d^\flat(x,x_b)}u_b
\end{equation}
at every settled output. The shifted tower exists on this
whole domain and has the literal product $e^{i\eta A}\mathbf U$
as the input to its actual average.
\end{corollary}
\begin{proof}
Take $a\le a_*=(64L^2)^{-1}$. The product logarithm lies in
the fixed chart since $\eta|w|\le a_*$ and
$\eta|A|\le a_*/4$. Write
$\log(e^Xe^Z)=X+Z+\ell(X,Z)$.
The absolutely convergent logarithm and exponential series
on $|X|,|Z|\le1/16$ give the derivative bound
$8(|X|+|Z|)$ for the remainder. As $\ell(0,Z)=0$,
Schwarz in $X$ also gives
$|\ell(X,Z)-\ell(X,Z')|\le16|X||Z-Z'|$.
Apply these estimates to neighboring bonds after transporting
them to the same base by $U$. The product rule gives
\[
 l(b)|(w\boxplus A-w)(b)|\le(1+10a_*)u_b,
\]
and the first-difference bound has coefficient at most
\[
 (1+10a_*)(1+2.1a_*L)+16a_*<2
\]
times $\max(u_b,u_{b^+})$. Here the term $2.1a_*L$
comes from replacing $\nabla_U A$ by $\nabla_{\mathbf U}A$:
$|1-\operatorname{Ad}(e^{i\eta w})|\le2.1\eta|w|$,
and neighboring scales differ by at most $L$.
Thus $n_b(w\boxplus A-w)\le2u_b^+$ and the stated margin
puts $w\boxplus A$ in the same domain as $w$.
The product identity holds at positive bonds by exponentiating
the logarithm and at reverse bonds by inversion.
Apply \eqref{atc:Lipschitz}. Adjacent bare sites have
ambient distance at most one, so replacing $u_b^+$ by $u_b$
in the weighted maximum costs at most $e^\delta$.
This proves the bound and the actual-input assertion.
\end{proof}


\subsection{Changing the localization of a tower}
\label{rtcg:section}

The response in \cref{atc:theorem} keeps the determining sequence
fixed. We next compare the two constrained problems obtained when
a localization is cut and its tower is rebuilt. All determining
sets in this subsection have their inward rows removed. A common
row means a common geometric row with its complete averaging read;
a common free pillar includes its prescribed normalization paths.

Use the notation of \cref{atc:section}. Let $\widehat{\mathbf B}$
be a finer localization of $\mathbf B'$, of the same depth $p$,
inside the same ambient sequence of depth $i$. Its first member is
$\widehat Z\subset Z$. Finer means that its members are contained
in the old members, so that the exact restriction identity
\eqref{mb:restriction-identity} applies. Denote their weighted
metrics by $d'$ and $\widehat d$. Both dominate $d^\flat$.
An old interior or outward row cannot become inward in
the finer sequence: an old layer endpoint stays outside
the next old member and hence outside the next finer
member. Thus removing all finer inward rows does not
remove an old retained row needed by the restriction identity.
Write $\mathcal C$ for the new determining rows together with
rows whose complete reads have changed. Enlarge $\mathcal C$
once to include the complete preceding blocks of its outward
rows and the old pillars containing its endpoints; call the
resulting carrier $\mathcal C^+$. It is a set of reads, not a new
set of constraints.

For $W=e^{i\eta w}U$, where $U$ is real, set
\begin{equation}\label{rtcg:profile}
 A(z)=\max\{\epsilon_U(P),n_b(w):x_P=z\ \hbox{or}\ x_b=z\},
 \qquad
 A_\delta(y)=\max_z e^{-\delta d^\flat(y,z)}A(z).
\end{equation}
Empty maxima are zero. Complete stencils can instead be based
at any of their corners, at a fixed change in the constant.
All amplitudes below lie below one fixed smaller ceiling than
the ceilings of \cref{bgc:uniform-physical-chart,atc:theorem}.

\begin{lemma}[Alignment of the induced and literal data]
\label{rtcg:alignment}
There are fixed $C,r,\delta>0$ with the following property.
The literal finer datum of $W$ and the datum induced by the
old critical continuation can be put in one physical datum
chart on $\widehat{\mathbf B}$, holomorphically in $w$.
In this chart their logarithms $b_0(w),b_1(w)$ satisfy
$\|b_\nu(w)\|_\infty<r$, and
\begin{equation}\label{rtcg:collar-datum}
 b_1(c)-b_0(c)=0\quad(c\notin\mathcal C),\qquad
 |b_1(c)-b_0(c)|\le C A_\delta(y_c)\quad(c\in\mathcal C).
\end{equation}
The segment joining these logarithms stays in the chart.
At an unchanged free pillar, the normalized finer continuations
at its two endpoints are respectively the new and the old tower
in exactly the same node frame.
\end{lemma}

\begin{proof}
First fix the gauges and the real centre of the chart. On every
old pillar of height $j$ put the ascent gauge equal to one at
its top. If $W_q=\mathcal M^qW$ and $T_q(z,x;W_q)$ is the
actual level-$q$ mean contour from the parent centre $z$, its
descent is
\begin{equation}\label{rtcg:ascent}
 g_j=1,\qquad g_q(x)=g_{q+1}(z)T_q(z,x;W_q),\quad q<j.
\end{equation}
This is the contour convention of
\eqref{mb:transported-bond}; reverse orientations use inverses.
The contour from a centre to itself is one, so the descent
extends the gauge already assigned at coarser sites. Covariance
of the actual mean makes the internal mean contours one.
At layer tops set $g=1$. The pillar partition makes these
assignments compatible. Thus every retained old datum is
unchanged, including an outward row and its lower endpoint
block. The old critical field is already in the canonical
old normalization, so its own descent is the identity on
its old pillars. We do not apply the raw input's gauge to
that critical field. The two separately normalized fields
still have the same old datum.

For completeness this construction is uniform on the complex
ball. Retain each real contour as a unitary factor. The relative
level-$q$ factor on a pillar of height $j$ has logarithm bounded
by $CaL^{q-j}$ when $n(w)<a$. The local mean estimate and its
site correction give this bound exactly as in
\eqref{htc:pillar-gauge}. Moving the real factors to the right
and telescoping the relative factors gives
\[
 \|g(w)g(0)^{-1}-1\|\le
 C a\sum_{q<j}L^{q-j}\exp\left(Ca\sum_{q<j}L^{q-j}\right),
 \qquad \|g(w)^{\pm1}\|\le e^{Ca/(L-1)}.
\]
The real factors have norm one. All means and inverses are
holomorphic on their fixed charts, so this also proves the
holomorphy and the uniform condition number of the descent.
There is no factor raised to the number of levels.

We spell out the comparison of new rows. At a real input the
two fields just normalized have the same old datum.
In an old pillar of height $j$, averaging a curvature bound
at level $q<j$ gives $C\epsilon L^{2(q-j)}$.
The finite contour reconstruction
\eqref{mb:transported-bond}--\eqref{mb:rigidity-sum}
therefore compares both fields to the same face lift of their
old datum. For the complex version use relative products,
rather than unitary invariance. If $E_q$ is the normalized
bond discrepancy at the current level, its exact factorization
has the form
\[
 E_q=(\widehat W_q\widehat W_{q+1}^{-1})
                (\widehat W_{q+1}L_{q+1}^{-1}).
\]
Here $L_{q+1}$ denotes the face lift in the common old frame.
The two level-$(q+1)$ quantities in this identity mean
their actual face lifts evaluated at the same level-$q$
bond as $\widehat W_q$. Thus all three factors belong to
one matrix fibre. The one-step lift consistency proved
with \eqref{mb:collar-face-lift} gives this exact product.
For matrices near one,
$|AB-1|\le |A-1|+(1+|A-1|)|B-1|$. Consequently
\begin{equation}\label{rtcg:product-descent}
 e_q\le C\epsilon_q+(1+C\epsilon_q)e_{q+1},\quad
 \epsilon_q\le C\epsilon L^{2(q-j)},\qquad
 \prod_{q<j}(1+C\epsilon_q)\le
             \exp\left(\frac{C\epsilon}{L^2-1}\right).
\end{equation}
The endpoint discrepancy at an old retained row is zero.
At an inward link, use the outward link one level higher and
its complete preceding block, as in the inward calculation
following \eqref{mb:rigidity-sum}. This adds a fixed number
of lower internal links, not another constraint. Summing
\eqref{rtcg:product-descent} proves a uniform raw discrepancy.

Here the local $\epsilon$ can be replaced by $C A_\delta(y_c)$.
For the real critical field this follows from
\eqref{atc:critical-profile} and the complete-read datum estimate
\eqref{atc:real-datum-profile}. For its complex variation,
differentiate the physical datum chart from the real input to
$w$. The raw variation of a level-$q$ input row is bounded
by $C\xi_q|w|$ on its complete read, hence by $Cn(w)$.
The weighted potential and first-difference response in
\eqref{pnc:weighted-first-response}, and the differentiated
literal pillar factors from the proof of
\cref{atc:datum-response}, give the same weighted estimate
for the continued field. A plaquette in an old pillar of
height $j$ then has defect bounded by
$C A_\delta(y_c)L^{2(q-j)}$ at level $q$.
Indeed its physical curvature is bounded in scale $\xi_j$;
averaging to level $q$ multiplies it by $\xi_q^2$.
The complete reconstruction reads remain within bounded
$d'$-distance of that pillar top. The ambient metric is
smaller, and the convolution of the two weighted suprema
uses a smaller fixed $\delta$. This proves the local form
of the estimate in \eqref{rtcg:product-descent}.
The input field obeys the same curvature bound by the
plaquette product formula and the $C^1$ norm in
\eqref{atc:input}. No second difference of the input is used.

Take as finer real reference the regular critical orbit at
the finer datum of $U^{g(0)}$. Its existence follows from
\cref{bgc:proper-background}. At $w=0$ the preceding descent
places the datum induced by the old real critical field within
$C\sup A$ of this reference datum. The variation of either
endpoint datum on the whole complex ball is $Ca$: this follows
for the input from the mean derivative and the bound on $g$,
and for the old field from its physical chart and the fact
that a new-row scale does not exceed the old scale on its
complete read, up to a single factor $L$.
Choose $C(\sup\epsilon_U+a)<r/4$. Both endpoints then lie
in the concentric half-ball of one finer physical chart.
Their logarithms and the segment joining them lie in that
chart. The logarithm has bounded derivative there, so the
raw discrepancy just proved gives \eqref{rtcg:collar-datum}.

Every common old row is fixed exactly by the gauge at its
endpoints, and the two fields have identical old data.
Thus their finer data, and their raw logarithms relative
to the same reference, agree on common rows. Passing to the
reduced datum does not produce a chain of new supports:
\eqref{mb:seam-datum-nilpotence} says that its off-diagonal
part reads interior columns into outward rows and that
$TD^{-1}T=0$. Hence differentiation or subtraction enlarges
the changed support by at most one complete preceding block.

It remains to identify the endpoints as fields, rather
than merely as data. For real $w$, restriction of the old
critical field is critical at its induced finer datum by
\eqref{mb:restriction-identity}. Its regularity improves in
the finer scales. Global regular-fibre uniqueness from
\cref{bgc:global-comparison} identifies it with the finer
physical chart at that datum. These identities extend to
the whole complex ball by the identity theorem: both sides,
their constraints, and all normalization factors are
holomorphic there, and the real ball is a uniqueness set.
Only the transformed raw data are admitted to this common
physical chart. No invariance of a hereditary core under
complex gauge transformations is asserted or required.

Finally the gauge \eqref{rtcg:ascent} fixes every old top.
Gauge covariance of a normalized finer representative uses
the pillar-constant extension of its values at the new tops.
On a common free pillar its top is an old top, so this
extension is exactly one. Normalizing the old restriction
on the finer sequence also acts identically on such a
pillar: its mean contours there already have the prescribed
old normalization. These statements hold as matrix identities
also at complex data. In particular the readout frames agree;
a bound on the condition number alone is not used to identify
their magnitudes.
\end{proof}

\begin{proposition}[Comparison after a geometric cut]
\label{rtcg:cut}
Suppose the complete read at $x$ lies in unchanged free
pillars and is settled for both towers. Work in the units
of the earlier ambient level $i$, so that $l\le1$.
Suppose every changed complete datum carrier is at physical
distance at least $D\ge2$ from $x$. Let $\mathcal R(z)>0$
be an allowance profile satisfying
\begin{equation}\label{rtcg:allowance}
 A(z)\le C_0\mathcal R(z),\qquad
 \mathcal R(z)\le C_1\mathcal R(x)
                    e^{\kappa d^\flat(x,z)},\quad
                 \kappa<\delta/4.
\end{equation}
Then the difference of the two matrix readouts in
\eqref{atc:output}, in their prescribed common frame, is at most
\begin{equation}\label{rtcg:cut-bound}
 C C_0C_1 e^{-c\delta D}\mathcal R(x).
\end{equation}
The assertion includes full current stencils across a common
layer seam. All constants are independent of depths and volume.
\end{proposition}

\begin{proof}
Apply \cref{atc:datum-response} on the finer chart of
\cref{rtcg:alignment}. Its proof includes the reduced datum,
the literal node correction and the full physical Euler
current. Thus it gives the native estimate
\begin{equation}\label{rtcg:screened-profile}
 |\widehat F_x-F_x|\le
 C\max_{c\in\mathcal C^+}
             e^{-m\widehat d(x,y_c)}A_\delta(y_c).
\end{equation}
All source measures are included in its weighted row sum.
There is no cardinality factor for the collar.
The support enlargement described above is why the distance
hypothesis uses the complete carrier.

The ambient readout is the native one times $s(x)^2$ for
both curvature and current. At a settled read either $s=1$
or $\xi'(x)=\xi_p$. In the latter case
$\widehat d(x,y_c)\ge |x-y_c|_\infty/\xi_p\ge D/\xi_p$
and $s=l/\xi_p\le1/\xi_p$, so
$s\le1+\widehat d(x,y_c)$.
A fixed part of the exponential in
\eqref{rtcg:screened-profile} pays $s^2$.
This is the extra scale argument needed for a changed
geometry; the current retains the weight $l^2\xi'$.

For any $c,z$, the triangle inequality and
$\widehat d\ge d^\flat$ give
\[
 d^\flat(x,z)\le
         \widehat d(x,y_c)+d^\flat(y_c,z).
\]
Insert \eqref{rtcg:allowance} in the definition of $A_\delta$.
Reserve one part of each exponential for this triangle
inequality and another part for the factor $s^2$.
The remaining decay is at most
$\exp(-c\delta\widehat d(x,y_c))\le\exp(-c\delta D)$,
since $\xi_p\le1$ in the present units.
This proves \eqref{rtcg:cut-bound}.
The frame identities in \cref{rtcg:alignment} make it a
matrix difference estimate. The reverse triangle inequality
therefore gives the same estimate for the change of magnitude.
The full-stencil convention and the physical Euler derivative
in \cref{atc:datum-response} cover a common seam without
discarding any incident plaquette.
\end{proof}

The distance hypothesis in \cref{rtcg:cut} holds for the
geometric restriction of an earlier address. Here
$p\le r\le i$, its nominal layer is at most $i$, and the
read retains its $r$-window. The accumulated erosion reach
through level $q$ is at most $C_LM_1\xi_q$, by the geometric
sum in the proof of \cref{htc:geometry}. Therefore a new
row at level $q\le p$ is within $C_LM_1\xi_p$ of the deleted
set. Its complete backward read adds $C_L\xi_p$.
An old parent pillar adds at most $C_L\xi_i=C_L$ in physical
distance, even when its new endpoint has smaller assigned
scale. Thus increasing the fixed cut halo from $3M_1$ to
$C_{\rm cut}M_1$ puts every changed complete carrier at
distance at least $D$ when the read starts at distance
$D+C_{\rm cut}M_1$ from the deleted set.
There is no additional spacing assumption on discarded
components. Indeed every old erosion test and complete read
has a connected backward carrier contained in $Z$ when it
is retained. Follow such a carrier from the retained
component toward a point removed by the component selection.
If it stays in $Z\cap R$, it connects the two components,
contrary to their definition. Its first exit must therefore
meet $Z\setminus R$. A contact with $Z^c$ was already an old
boundary and cannot cause a new membership change.
The same geometric bound on the backward carrier therefore
locates the changed row near the actual cut, also in this case.
The constant $C_{\rm cut}$ is chosen before the localization
floor $M$ is increased. Equivalently retain the original
hypothesis $D+3M_1$ and increase $K_{\rm cl}$ in
$D\ge K_{\rm cl}M$ so that
$(C_{\rm cut}-3)_+M_1\le D/2$. Then every changed complete
carrier has distance at least $D/2$, which changes only
the fixed decay rate in \eqref{rtcg:cut-bound}.
This argument uses physical distance
and $l\le1$. It does not assert that an arbitrary Euclidean
collar has bounded width in the ambient weighted metric.




% BEGIN LAYERED EXPANSION inverse/canonical_nonlinear_germ.tex
\subsection{A selected nonlinear extension at an independent current}
\label{ch:cng-canonical-nonlinear-germ}

We specify a nonlinear auxiliary action whose Hessian is the
canonical second precision. Its small constrained solution
agrees with the physical critical germ at physical inputs.
Away from those inputs it is a selected independent-current
extension, not a claim of Wilson stationarity.

The distinction between two output fields is essential.
The auxiliary potential has first derivative \(H_1\) in
the reduced datum. Its finite gauge completion has first
derivative \(SH_1\). These agree at a physical pair, where
the second lift is horizontal, and can differ off current.
Literal outward data also retain their actual datum
differential.

\subsubsection{The fixed chart and the strong spaces}

Fix \(0<\alpha<1\), and a given real regular reference
\(U_*\) on the actual determining geometry, with strict
room in its curvature and physical-current ceilings.
The reference need not be critical. Its local real
Coulomb potentials have the scaled \(C^{1,\alpha}\)
bounds supplied by the current-regular construction.
All choices of local frames and node extensions in
this subsection are made in this fixed reference chart.

Let
\[
              U=e^{i\eta E}U_*,\qquad
                    \|E\|_{{\rm str},\alpha;U_*}<\varepsilon .
\]
This is a weighted \(C^{1,\alpha}\) tube, stronger than
the \(C^1\) tube used for energy inverses. Let
\(\mathcal A_\alpha\) be the fine strong space in
\eqref{psr:strong-norm}, measured in \(U_*\), and let
\(\mathcal Y_{2,\alpha}\) be the site gauge space in
\eqref{dsc:site-gauge-norm}. Its source space
\(\mathcal C_{\alpha,2}\) is
\eqref{dsc:site-source-norm}, including comparisons
across neighbouring pillar faces on complete boxes.
Use the amplitude norms
\begin{equation}\label{cng:amplitude-norms}
 \|f\|_{\infty,3}=\sup_b\mathsf L(b)^3|f(b)|,\quad
 \|h\|_{\infty,4}=\sup_x\xi_x^4|h(x)|,\quad
 \|b\|_{\mathcal B}=\sup_c\xi_c|b(c)|,\quad
 \|\nu\|_{\mathcal W}=\sup_c\xi_c^3|\nu(c)|.
\end{equation}
The site scale and the field-row scale retain their
different assigned levels at outward endpoints.
The finite-\(p\) force norm is \(\mathcal F_p\) of
\cref{ch:proper-strong-response}.

The independent fine current has a full amplitude
\(\|J\|_{\infty,3}<\gamma_*\).
This is convenient for the updated-current assertion
below. The actual Wilson current \(j_U\) at a complex
\(U\) is not assumed to have this bound. Its possibly
large value will be subtracted only as an exact linear
centering. All transposes and scalar pairings in an
analytic formula are physical bilinear ones.

For this fixed \(U\), retain its actual \(D,q,Q,G_s,R\)
and put
\begin{equation}\label{cng:restoration-objects}
              \Lambda=G_sRD^\sharp,\qquad S=1-D\Lambda .
\end{equation}
Define \(G,H,\omega,G_1,H_1\) in the order of
\cref{ch:actual-current-completion}, with
\(\omega=H^\sharp J\). In particular the lift in
\(\omega\) is the first lift. Their completed operators
are uniformly admitted on a smaller tube by the
energy and strong inverse results already proved.

\subsubsection{Completed restoration and the two adjoint bounds}

\begin{lemma}[Strong restoration on the stronger tube]
\label{lem:cng-strong-restoration}
After shrinking the fixed \(C^{1,\alpha}\) tube,
\begin{equation}\label{cng:strong-restoration}
 \|\Lambda A\|_{\mathcal Y_{2,\alpha}}
                         \le C_\alpha\|A\|_{\mathcal A_\alpha}.
\end{equation}
This map is holomorphic in the background on that
tube. It is an estimate for the completed operator.
It does not assert a \(C^{2,\alpha}\) bound for
every individual inverse-walk coefficient.
\end{lemma}
\begin{proof}
The first physical divergence maps
\(\mathcal A_\alpha\) into \(\mathcal C_{\alpha,2}\).
The site projection \(R\) preserves that source space.
Indeed its nonidentity part is
\(\widehat F\widehat M^{-1}\widehat F^\sharp\).
The completed site-column point and first-difference
bounds, the normalized coarse inverse and their
exponential sums control its scaled \(C^\alpha\)
norm. All normalizations are those of
\eqref{ncw:normalized-site-columns}.

For \(\lambda=G_sRD^\sharp A\), the exact identities
\(qG_sR=0\) and
\(\Delta_UG_sR=R\) give
\[
               q\lambda=0,\qquad
                         \Delta_U\lambda=RD^\sharp A .
 \]
The zero-average energy estimate supplies its weighted
mass and first differences. The local scalar bootstrap
supplies the \(C^{1,\alpha}\) bound needed on the
right of the interior Schauder estimate.
On a fixed complete nonbare box the coefficients of
\(\Delta_U\) have their \(C^\alpha\) bounds because
the potential is \(C^{1,\alpha}\). At a true boundary
the bare values of \(\lambda\) are exactly zero and
the remaining atomic stencils have their scaled bounds.
The proof of \eqref{dsc:proper-site-gauge-bound}
therefore gives \eqref{cng:strong-restoration}.

For completeness, at complex \(U\) compare with the
real reference Laplacian. Its difference is a
first-order term with \(C^\alpha\) coefficients and
a zeroth-order term with \(C^\alpha\) coefficients,
bounded by the \(C^{1,\alpha}\) tube. Apply the real
Schauder estimate to this equation, using the already
controlled \(C^{1,\alpha}\) solution norm. No
\(C^{2,\alpha}\) conclusion from merely \(C^1\)
coefficients is used. The energy inverse is
holomorphic, and the uniform stronger bound on
the smaller tube gives its holomorphy in these
displayed Banach norms.
\end{proof}

\begin{lemma}[Adjoints in their actual force classes]
\label{lem:cng-adjoint-restoration}
For site and field forces respectively,
\begin{align}
 \|\Lambda^\sharp h\|_{\infty,3}
                            &\le C\|h\|_{\infty,4},\notag\\
 \mathcal F_p(\Lambda^\sharp D^\sharp f)
                            &\le C_p\|f\|_{\infty,3},
                    \qquad 4<p<\infty .
 \label{cng:adjoint-force-bounds}
\end{align}
The completed operators in these assertions are
holomorphic on a common smaller physical tube.
There is no endpoint \(p=\infty\) in the second line.
\end{lemma}
\begin{proof}
The balanced site family
\(G_s:Y^\flat\to Y\) gives bounded local \(Y\)
norm for \(u=G_sh\): a site force of amplitude
four has uniformly bounded local \(Y^\flat\)
norm. Its equation is
\[
                       \Delta_Uu=h-q^\sharp a_squ .
 \]
The average \(qu\) is bounded by its cell \(L^2\)
norm, and \(q^\sharp\) extends the transported
density. Thus the rescaled forcing is bounded.
The scalar \(L^2\)-to-\(L^3\)-to-\(L^8\) bootstrap
gives
\(\xi_y^2|u|+\xi_y^3|Du|\le C\|h\|_{\infty,4}\).
It uses no derivative of \(h\).
Subtract \(Pu=\widehat F\widehat M^{-1}\widehat F^\sharp u\).
Its normalized coarse coefficients are bounded by
the column sum, then by the coarse \(\ell^\infty\)
inverse. The completed \(F\) columns and their first
differences give that same point bound for \(Pu,DPu\).
Since \(\Lambda^\sharp=DRG_s\), this proves the first line.

For the second line use the exact identity
\[
                \Lambda^\sharp D^\sharp f=f-S^\sharp f .
 \]
The finite-\(p\) force-restoration estimate of
\cref{ch:proper-strong-response} applies to this
variational divergence source. It gives the stated
bound without a pointwise differentiation of \(f\).
Its proof on the complex physical tube uses the
same energy inverse and local finite-\(p\) bootstrap.
The absolute estimates are uniform on a smaller
tube; the analytic inverse formulas give the
claimed holomorphy.
\end{proof}

\subsubsection{The finite gauge chain as a strong analytic map}

For \(A\in\mathcal A_\alpha\), set \(\lambda=\Lambda A\)
and define
\begin{equation}\label{cng:finite-gauge-map}
 \Pi_U(A)_b=\frac1{i\eta}\log
    \{e^{-i\lambda(b_-)}e^{i\eta A_b}
                     e^{i\operatorname{Ad}(U_b)\lambda(b_+)}\},
 \qquad \Xi_U(A)=\Pi_U(A)-SA .
\end{equation}
This is the finite chain already selected for the
canonical current convention. Its associated fine
gauge is \(e^{-i\lambda}\).
The exact equality
\[
 \operatorname{Ad}(U_b)\lambda(b_+)
                         =\lambda(b_-)-\eta D\lambda(b)
\]
makes division by \(\eta\) uniform: at \(\eta=0\)
the three factors multiply to one.
Thus \(\Pi_U\) is an analytic local function of
\(A,\lambda,D\lambda\), with
\begin{equation}\label{cng:Pi-strong-bounds}
 \Pi_U'(0)=S,\quad
 \|\Pi_U(A)\|_{\mathcal A_\alpha}\le C_\alpha\|A\|_{\mathcal A_\alpha},
 \quad
 \|\Xi_U(A)\|_{\mathcal A_\alpha}\le C_\alpha\|A\|_{\mathcal A_\alpha}^2 .
\end{equation}
Covariant product rules use at most two site
differences of \(\lambda\). They prove these bounds
also for the gradient H\"older seminorm.
If \(RD^\sharp A=0\), then \(\lambda=0\) and
\(\Pi_U(A)=A\) exactly. No claim that \(\Pi_U(A)\)
lies in that slice for a general \(A\) is made.

The physical adjoint of the derivative requires
more than an abstract strong-space norm. Pair a
local function of \(A,\lambda,D\lambda\) with a
bounded force \(f\). Differentiating its three
arguments gives the exact form
\begin{equation}\label{cng:adjoint-factory}
        f_0+\Lambda^\sharp f_\lambda
                         +\Lambda^\sharp D^\sharp f_d .
\end{equation}
Here \(f_0,f_d\) are field forces of amplitude
three and \(f_\lambda\) is a site force of
amplitude four. Endpoint injections cost only
the fixed number of incident bonds. The two
bounds in \eqref{cng:adjoint-force-bounds}
therefore control this expression in finite
\(\mathcal F_p\), without differentiating \(f\).

In particular \(\langle J,\Xi_U''(0)[u,v]\rangle
=\langle u,B_U[J]v\rangle\) by the exact finite-chain
calculation. Subtracting that quadratic jet in
\eqref{cng:adjoint-factory} gives
\begin{equation}\label{cng:current-nonlinear-residual}
 \mathcal F_p\{\nabla\langle J,\Xi_U(A)\rangle-B_U[J]A\}
                   \le C_{p,\alpha}\|J\|_{\infty,3}
                                             \|A\|_{\mathcal A_\alpha}^2 .
\end{equation}
Indeed the remaining three local coefficients are
quadratic in the scaled variables
\(\xi A,\lambda,\xi D\lambda\).
Their amplitude grades are exactly those just
listed. This proves the residual in a force norm;
an arbitrary Banach adjoint is not presumed to
preserve pointwise bounds.

\subsubsection{The centered Wilson residual at a complex base}

Put \(F_U(A)=A_{\Omega_0}(e^{i\eta A}U)\),
\(j_U=F_U'(0)\), and \(H_U=F_U''(0)\).
The centering in the next lemma is essential even
when \(U\) belongs to the complex strong tube.

\begin{lemma}[Uniform residual after linear subtraction]
\label{lem:cng-Wilson-residual}
On a fixed strong ball in the displayed tube,
\begin{equation}\label{cng:Wilson-residual}
 \|F_U'(A)-j_U-H_UA\|_{\infty,3}
                             \le C_\alpha\|A\|_{\mathcal A_\alpha}^2 .
\end{equation}
The difference is uniformly holomorphic in \(U,A\).
Its Lipschitz constant on a radius-\(r\) strong
ball is at most \(C_\alpha r\).
No amplitude bound for \(j_U\) at that complex base
is used.
\end{lemma}
\begin{proof}
On a nonbare box rescale its physical length to
one, and write \(h=\eta/\xi\). In a fixed real
reference frame the complex base has potential
\(a_0\) with bounded \(C^{1,\alpha}\) norm.
The relative logarithm is
\[
 P_h(A)=\frac1{ih}\log(e^{ihA}e^{iha_0}).
 \]
Let \(P_0=D_AP_h(0)\) and
\(E_h(A)=P_h(A)-a_0-P_0A\).
The divided finite matrix expansion gives
\[
 \|E_h(A)\|_\infty+\|DE_h(A)\|_\infty\le Chr^2,
 \qquad \|D^2E_h(A)\|_\infty\le Cr^2,
\]
and
\[
 \|D_AP_h(A)-P_0\|\le Chr,\qquad
 \|D_AP_h(A)-P_0-D_A^2P_h(0)[A]\|\le Ch^2r^2.
 \]
The second spatial difference contains \(hA D^2A\)
and \(h(DA)^2\); when a coefficient is differentiated
in \(a_0\), it supplies an additional \(h\).
The elementary bounds
\(\|D^2A\|_\infty\le Ch^{-1}\|DA\|_\infty\)
and
\(\|D^2a_0\|_\infty\le Ch^{-1}\|Da_0\|_\infty\)
therefore remove every apparent mesh loss.

In flat logarithmic coordinates the exact Wilson
gradient is \(\mathrm{d}^\sharp \mathrm{d}\,a+N_h(a)\), with the
uniform local analytic remainder proved in the
actual Wilson calculation. The relative gradient
is \(D_AP_h(A)^\sharp\{\mathrm{d}^\sharp \mathrm{d}P_h(A)+N_h(P_h(A))\}\).
Subtract its constant and linear terms before
estimating. The term with \(\mathrm{d}^\sharp \mathrm{d}E_h\) is
bounded by \(Cr^2\).
The first derivative difference \(O(hr)\) times
\(\mathrm{d}^\sharp \mathrm{d}(P_0A)\) is bounded by
\(Chr\,h^{-1}r\).
The remaining quadratic derivative coefficient
\(O(h^2r^2)\) multiplying the rough base
\(\mathrm{d}^\sharp \mathrm{d}a_0\) is bounded by
\(Ch^2r^2h^{-1}\|Da_0\|_\infty\).
All terms containing \(N_h\) obey its uniform
Banach derivative bounds. This is the same
relative-BCH calculation as in
\cref{lem:pnc-Wilson-residual}; it uses the complex
base's first differences, rather than a bound for
its second-difference current.

At a bare cell the rescaled mesh is one and the
finite incident plaquette calculation gives the
same analytic Taylor estimate. Rescaling the
gradient back multiplies it by \(\xi^{-3}\).
All estimates hold on a larger complex ball.
Their uniform holomorphy and the double zero
give the stated Lipschitz bound by Cauchy.
\end{proof}

\subsubsection{An explicit auxiliary action with the second precision}

Write \(\mathfrak q_U(A)\) for the actual re-centered
straightened per-length field constraint of PNC;
this notation distinguishes it from the site
average \(q\). Put
\begin{equation}\label{cng:nonlinear-constraint}
 N_q(x)=\mathfrak q_U(x)-Qx,\qquad
                    r_U(A)=\mathfrak q_U(\Pi_U(A)).
\end{equation}
Then \(r_U(0)=0\) and \(r_U'(0)=Q\), because
\(QS=Q\). The latter equality uses the actual
extended outward endpoint intertwining.

Define the following scalar function on the fine
strong ball:
\begin{equation}\label{cng:auxiliary-action}
 \begin{split}
 \mathcal L_{U,J}(A)={}&F_U(A)-F_U(0)-\langle j_U,A\rangle\\
 &-\langle J,\Xi_U(A)\rangle
            -\langle\omega,N_q(\Pi_U(A))\rangle\\
 &+\tfrac12\langle RD^\sharp A,RD^\sharp A\rangle_{\rm site}
            +\tfrac12\langle r_U(A),a_fr_U(A)\rangle_{\rm rows}.
 \end{split}
\end{equation}
Every pairing is bilinear, including the gauge
penalty. Thus this is a holomorphic scalar
function at complex inputs. The multiplier
\(\omega=H^\sharp J\) is fixed when differentiating
in \(A\).

\begin{lemma}[The exact auxiliary Hessian]
\label{lem:cng-exact-Hessian}
The action satisfies
\par\noindent
\begin{minipage}{\linewidth}
\begin{equation}\label{cng:exact-Hessian}
             \mathcal L_{U,J}'(0)=0,\qquad
                         \mathcal L_{U,J}''(0)=G_1^{-1}.
\end{equation}
\popQED\par
\end{minipage}
\end{lemma}
\begin{proof}
The first derivative of the nonlinear current
term is zero. Its second derivative is
\(-B_U[J]\), by the finite-chain identity.
Since \(N_q(0)=DN_q(0)=0\), the constraint
term contributes
\(-S^\sharp\Delta_U^{(2)}[\omega]S\).
The projection identities \(R^\sharp=R\),
\(R^2=R\) make the penalty contribution
\(DRD^\sharp\), also on the complex bilinear
continuation. Finally \(r_U'(0)=Q\) makes
the row-mass contribution \(Q^\sharp a_fQ\).
Together with \(H_U\), these are exactly the
terms in the canonical \(G_1^{-1}\).
\end{proof}

This formula selects a nonlinear off-current
convention. Equality of its second jet with the
canonical precision does not identify it with
a separately prescribed nonlinear core map.

\begin{lemma}[Constraint and composed-adjoint residuals]
\label{lem:cng-constraint-residuals}
For \(\|A\|_{\mathcal A_\alpha}\le r\), the actual
constraint obeys
\begin{align}
 \|r_U(A)-QA\|_{\mathcal B}&\le Cr^2,\notag\\
 \mathcal F_p\{(r_U'(A)-Q)^\sharp\nu\}
                                      &\le C_{p,\alpha}r\|\nu\|_{\mathcal W},
 \label{cng:constraint-residuals}\\
 \mathcal F_p\{r_U'(A)^\sharp a_fr_U(A)-Q^\sharp a_fQA\}
                                      &\le C_{p,\alpha}r^2,\notag\\
 \mathcal F_p\{\nabla\langle\omega,N_q(\Pi_U(A))\rangle
                       -S^\sharp\Delta_U^{(2)}[\omega]SA\}
                                      &\le C_{p,\alpha}\|\omega\|_{\mathcal W}r^2 .
 \notag
\end{align}
These are estimates for finite \(4<p<\infty\).
\end{lemma}
\begin{proof}
The actual fixed-leg two-leg estimate and its
analytic remainder, in
\eqref{pnc:actual-two-leg-bound} and
\eqref{pnc:constraint-residual-bounds}, give
\[
 \|N_q(x)\|_{\mathcal B}\le C\|x\|_{\mathcal A_\alpha}^2,\qquad
 \|(D\mathfrak q_U(x)-Q)^\sharp\nu\|_{\infty,3}
             \le C\|x\|_{\mathcal A_\alpha}\|\nu\|_{\mathcal W}.
 \]
After subtracting its linear term, \(DN_q(x)^\sharp\omega\)
has a force remainder bounded by
\(C\|\omega\|_{\mathcal W}\|x\|_{\mathcal A_\alpha}^2\).
These are actual column-density estimates on the
common admitted row ball, also when the base is in
the smaller complex tube. Physical row transposition
cancels the fine density \((\eta/\xi_c)^4\);
the recentering and complete outward carriers are retained.

The value estimate follows by
\(\Pi_U(A)=SA+O(A^2)\) and \(QS=Q\).
For the adjoints use
\[
 r_U'(A)^\sharp
       =\Pi_U'(A)^\sharp D\mathfrak q_U(\Pi_U(A))^\sharp .
 \]
The operator \(S^\sharp\) is bounded only in the
finite force class needed here.
The difference \(\Pi_U'(A)^\sharp-S^\sharp\)
has bound \(C_{p,\alpha}r\) from a bounded
amplitude-three force into \(\mathcal F_p\).
This is \eqref{cng:adjoint-factory}, before
subtracting its quadratic jet, together with
\cref{lem:cng-adjoint-restoration}.
In particular this argument never differentiates
the fine force supplied to that adjoint.

The identity \(S^\sharp Q^\sharp=Q^\sharp\)
cancels the constant part in the multiplier
residual. The remaining two factors give its
second line. For the row mass, use
\(a_f\asymp\Xi_\Gamma^{-2}\), so it converts
datum amplitude one into multiplier amplitude
three, and insert the first two lines.
For the last line put \(x=\Pi_U(A)\) in
the subtracted \(DN_q\) estimate and apply
the same adjoint formula. Terms with
\(\Xi_U(A)\) and
\(\Pi_U'(A)-S\) are quadratic after their
displayed leading term is removed.
\end{proof}

Combining the preceding lemmas, set
\[
                 N_{\mathcal L}(A)
                         =\mathcal L_{U,J}'(A)-G_1^{-1}A .
\]
For fixed small current ceiling they give
\begin{equation}\label{cng:full-gradient-residual}
              \mathcal F_p(N_{\mathcal L}(A))
                              \le C_{p,\alpha}\|A\|_{\mathcal A_\alpha}^2 .
\end{equation}
All maps in this subtracted expression are
uniformly holomorphic on a larger strong ball.
Their double zeros consequently give Lipschitz
constant \(C_{p,\alpha}r\) there.

\subsubsection{The quantitative auxiliary germ}

Let \(b_U(B)\) be the explicit actual outward
datum in \eqref{pnc:actual-reduced-datum}, relative
to the present base \(U\). Its finite complete
carriers and analytic bounds hold on a smaller
common tube. In particular \(b_U(0)=0\) and
\(\|b_U(B)\|_{\mathcal B}\le C_L\|B\|_{\mathcal B}\).
Choose a finite \(p>4/(1-\alpha)\).

\begin{proposition}[A selected independent-current critical germ]
\label{prop:cng-auxiliary-germ}
There are fixed positive radii \(\varepsilon,\gamma_*,r,\rho\)
such that the following holds on this fixed determining
geometry. For \(U\) in the smaller displayed
\(C^{1,\alpha}\) tube, \(\|J\|_{\infty,3}<\gamma_*\),
and \(\|b\|_{\mathcal B}<\rho\), the system
\begin{equation}\label{cng:auxiliary-system}
          \mathcal L_{U,J}'(A)-r_U'(A)^\sharp\nu=0,
                         \qquad r_U(A)=b
\end{equation}
has a unique solution in a fixed joint strong ball.
It is holomorphic in \(U,J,b\), real at real data, and
\begin{equation}\label{cng:germ-size}
       \|A\|_{\mathcal A_\alpha}+\|\nu\|_{\mathcal W}
                                      \le C_{p,\alpha}\|b\|_{\mathcal B}.
\end{equation}
The radii and constants are independent of depth,
compatible volume and admitted proper-region shape.
At \(b=0\), \(A=\nu=0\) for every admitted pair,
and the exact first derivatives are
\begin{equation}\label{cng:reduced-first-jets}
          D_bA(0)=H_1,\qquad
 D_b\nu(0)=(QG_1Q^\sharp)^{-1},\qquad
                     D_b(\Pi_U(A))(0)=SH_1 .
\end{equation}
For literal data set \(b=b_U(B)\); the first
and third derivatives then acquire the right factor
\(Db_U(0)\).
\end{proposition}
\begin{proof}
The actual second-variant row saddle has the strong
inverse
\[
 G_1^{-1}A-Q^\sharp\nu=f,\qquad QA=d,\qquad
 \|A\|_{\mathcal A_\alpha}+\|\nu\|_{\mathcal W}
          \le C_{p,\alpha}\{\mathcal F_p(f)+\|d\|_{\mathcal B}\}.
 \]
This is the completed canonical inverse, with its
row mass retained. It is valid for an independent
current and the common complex background tube.
Write this inverse as \(\mathfrak I_{U,J}\).
Then \eqref{cng:auxiliary-system} is exactly the
preconditioned fixed point
\begin{equation}\label{cng:preconditioned-fixed-point}
 (A,\nu)=\mathfrak I_{U,J}
   \bigl(-N_{\mathcal L}(A)
                   +(r_U'(A)-Q)^\sharp\nu,\quad
                                    b-(r_U(A)-QA)\bigr).
\end{equation}
The displayed residual estimates give size
\(C_{p,\alpha}r^2\) and Lipschitz constant
\(C_{p,\alpha}r\) in a radius-\(r\) joint ball.
Choose \(r\) so this constant is less than \(1/4\),
then \(\rho\) so its affine input is less than
\(r/4\). Iteration is invariant and contractive.
Its bound, realness and holomorphy prove the
first assertions. All choices are made after
fixing \(\alpha,p\), but use only the fixed
uniform constants. The current ceiling can be
kept inside the already admitted energy ceiling;
shrinking the strong and datum radii pays the
finite-\(p\) nonlinear constants.

At zero data the zero solution is exact. Every
nonlinear residual on the right has zero first
derivative there. Differentiation therefore gives
the full canonical row-saddle response,
\(A'=G_1Q^\sharp(QG_1Q^\sharp)^{-1}b'\)
and \(\nu'=(QG_1Q^\sharp)^{-1}b'\).
Applying \(\Pi_U'(0)=S\) proves the third
identity. Composition with the actual outward
datum proves its last assertion. The argument
uses the preconditioned residual map; no uniform
boundedness of an un-subtracted second-order
gradient on a \(C^{1,\alpha}\) ball is presumed.
\end{proof}

Denote \(C(B)=\Pi_U(A(b_U(B)))\). It satisfies
\(\mathfrak q_U(C(B))=b_U(B)\) and lies on the
same fine gauge orbit as the auxiliary field.
Off current it is not asserted to satisfy the
fixed linear Landau condition. In particular the
auxiliary identity \(D_bA(0)=H_1\) is not an
identity for this completed field. A specified
literal node extension adds its further fine
gauge vector to the first jet.

\subsubsection{Exact matching at physical pairs}

\begin{proposition}[The auxiliary germ equals the physical germ]
\label{prop:cng-physical-matching}
Suppose the base is an admitted real physical pair,
\begin{equation}\label{cng:physical-base}
             J=j_U=Q^\sharp\omega_{\rm phys}.
\end{equation}
Then \(\omega=H^\sharp J=\omega_{\rm phys}\).
For reduced datum \(b=b_U(B)\), let
\((A_{\rm phys},\nu_{\rm phys})\) be the physical
PNC germ in the Landau slice fixed at the present
base \(U\), namely \(R_UD_U^\sharp A=0\).
The reference \(U_*\) fixes comparison frames,
norms and the node extension; this slice and
\(\Lambda_U\) vary with the base \(U\).
On a common smaller datum ball,
\begin{equation}\label{cng:physical-matching}
 A=A_{\rm phys},\qquad C=A_{\rm phys},\qquad
                    \nu=\nu_{\rm phys}+a_fb .
\end{equation}
The same identity holds on the already proved
holomorphic physical critical graph wherever the
displayed strong and current admission holds.
\end{proposition}
\begin{proof}
The row identity \(QH=1\) identifies the first-lift
readout with the physical multiplier.
Fine gauge invariance gives
\(F_U(A)=F_U(\Pi_U(A))\) for the exact finite chain.
Its first Ward identity gives
\(\langle j_U,D\Lambda A\rangle=0\).
Substitute \(J=j_U=Q^\sharp\omega\) in
\eqref{cng:auxiliary-action}. The three linear
and current terms combine exactly as
\[
 -\langle j_U,A\rangle
 -\langle j_U,\Pi_U(A)-SA\rangle
 -\langle\omega,\mathfrak q_U(\Pi_U(A))-Q\Pi_U(A)\rangle
                       =-\langle\omega,r_U(A)\rangle .
 \]
Thus the scalar action is
\begin{equation}\label{cng:physical-scalar-identity}
 \begin{split}
 \mathcal L_{U,j_U}(A)={}&
       F_U(\Pi_U(A))-F_U(0)-\langle\omega,r_U(A)\rangle\\
 &+\tfrac12\langle RD^\sharp A,RD^\sharp A\rangle
                   +\tfrac12\langle r_U(A),a_fr_U(A)\rangle .
 \end{split}
\end{equation}
This is an equality of scalar functions on the
admitted ball, not merely a second-jet comparison.

The physical solution has \(RD^\sharp A_{\rm phys}=0\),
hence \(\Pi_U(A_{\rm phys})=A_{\rm phys}\).
Its complete stationarity equation is
\[
 F_U'(A_{\rm phys})
       =D\mathfrak q_U(A_{\rm phys})^\sharp
                                      (\omega+\nu_{\rm phys}).
 \]
Differentiate \eqref{cng:physical-scalar-identity}
at this field. The pullback of the first two
nonconstant terms is
\(r_U'(A_{\rm phys})^\sharp\nu_{\rm phys}\);
the gauge penalty is zero and the row mass
adds \(r_U'(A_{\rm phys})^\sharp a_fb\).
Therefore this field, with
\(\nu=\nu_{\rm phys}+a_fb\), solves
\eqref{cng:auxiliary-system}. The fixed-point
uniqueness proves \eqref{cng:physical-matching}.

For continuation along an already admitted
holomorphic physical graph, consider the
fixed-slice equation, constraint and full physical
stationarity just proved as fine-coordinate
functions of its graph parameters and \(B\).
They are holomorphic and vanish on the real
parameter neighbourhood. The identity theorem
therefore makes them vanish on the common
complex neighbourhood. This proves the stated
physical continuation without presuming a new
reference or stationarity for an arbitrary
off-current input.
\end{proof}

\subsubsection{A bounded updated current and a literal field}

Before the final literal gauge correction put
\begin{equation}\label{cng:updated-pair}
 \begin{split}
 U_{\rm new}&=e^{i\eta C(B)}U,\qquad
                \omega_{\rm new}=\omega+\nu-a_fb_U(B),\\
 J_{\rm new}&=\mathcal E_{\eta C(B)}^{-\sharp}
       \left\{D\mathfrak q_U(C(B))^\sharp\omega_{\rm new}
                                      +(J-Q^\sharp\omega)\right\}.
 \end{split}
\end{equation}
Here \(\mathcal E_zv=\int_0^1\operatorname{Ad}(e^{itz})v\,dt\).
The last expression is an explicit independent-current
selection at off-current inputs. Its coefficient
\(\omega_{\rm new}\) is in the displayed original
re-centered row coordinates; it is not declared to
be a new first-lift readout at the output pair.

\begin{proposition}[The updated independent current]
\label{prop:cng-current-update}
The update is holomorphic and satisfies
\begin{equation}\label{cng:current-update-bounds}
 J_{\rm new}(0)=J,\qquad
 \|J_{\rm new}-J\|_{\infty,3}\le C_{p,\alpha}\|B\|_{\mathcal B},
 \qquad
 \|J_{\rm new}\|_{\infty,3}
                  \le\|J\|_{\infty,3}+C_{p,\alpha}\|B\|_{\mathcal B}.
\end{equation}
Its field and curvature increments have the
corresponding strong and weighted bounds.
At a physical input it is exactly the left
Wilson current of \(U_{\rm new}\).
\end{proposition}
\begin{proof}
The actual transpose column bounds send
row amplitude three into fine force amplitude
three. The difference
\((D\mathfrak q_U(C)-Q)^\sharp\omega\)
has norm at most \(C\|\omega\|_{\mathcal W}\|C\|_{\mathcal A_\alpha}\).
Also
\(\|\nu-a_fb_U(B)\|_{\mathcal W}\le C\|B\|_{\mathcal B}\)
and
\(\mathcal E_{\eta C}^{-\sharp}-1\)
has point norm at most \(C\eta|C|\).
Subtract the value at \(C=0,\nu=0,B=0\) in
\eqref{cng:updated-pair}. These estimates and
\eqref{cng:germ-size} prove
\eqref{cng:current-update-bounds}. The curvature
increment follows from the four exact link
factors and the first difference of \(C\).
No actual Wilson derivative of \(U_{\rm new}\)
has been used in this off-current estimate.

At a physical input
\eqref{cng:physical-matching} gives
\(\omega_{\rm new}=\omega+\nu_{\rm phys}\),
and the residual \(J-Q^\sharp\omega\) vanishes.
The physical normal and left gradients satisfy
\[
 F_U'(C)=\mathcal E_{\eta C}^\sharp j(U_{\rm new}).
 \]
Insert physical stationarity from the preceding
proof. Equation \eqref{cng:updated-pair} is
then exactly \(J_{\rm new}=j(U_{\rm new})\).
\end{proof}

Finally choose the fixed real-reference node
extension \(E_n\) of
\cref{lem:pnc-site-extensions}. Apply it to the
small logarithm of the assigned factor
\(v_U[C(B)]^{-1}\), and exponentiate to obtain
the admitted fine gauge \(g(B)\). Its node
values are exactly that factor; its bare
values are one. The literal output pair is
\begin{equation}\label{cng:literal-update}
        \widehat U(B)=U_{\rm new}^{\,g(B)},\qquad
 \widehat J(B)_b=\operatorname{Ad}(g(B)_{b_-})J_{\rm new}(B)_b .
\end{equation}
The actual outward diagram gives
\(M_{\mathbf B}\widehat U(B)=V_B\).
The finite gauge estimate in
\eqref{pnc:finite-gauge-bound} preserves the
strong field bound, and its bounded point
adjoint preserves the current bounds.
At physical inputs this is the physical
critical orbit and its actual current in
the specified literal gauge.

All choices here belong to the fixed reference
chart. A different admissible node extension
differs by a fine gauge one at assigned tops,
but no global off-current chart-descent or
externally specified core convention has been
identified by this construction.

\subsubsection{Fixed-order spatial response}

The global strong estimates above can retain a
spatial weight. Let \(w_y=e^{s\psi(y)}\), where
\(\psi\) is one-Lipschitz in the proper layer
metric and \(s>0\) is fixed sufficiently small.
Weighted strong and source norms multiply each
local norm by \(w_y\). H\"older seminorms are
still taken on the same complete local boxes;
the artificial weight is not differentiated.

\begin{proposition}[Localized jets of the completed germ]
\label{prop:cng-localized-jets}
The first datum response satisfies
\begin{equation}\label{cng:weighted-response}
 \|D_BA[h]\|_{\mathcal A_\alpha;w}
 +\|D_B\nu[h]\|_{\mathcal W;w}
                                  \le C_{p,\alpha}\|h\|_{\mathcal B;w}.
\end{equation}
The same conclusion holds for the completed field
and the literal output pair in their displayed
strong and current norms. For each fixed \(m\),
the normalized response at an output cell \(y\)
to point data at \(c_1,\ldots,c_m\) is bounded by
\begin{equation}\label{cng:localized-jets}
                      C_m e^{-s\max_i d(y,c_i)} .
\end{equation}
These constants are uniform in depth and volume.
No all-arity absolute connected-kernel bound
is asserted.
\end{proposition}
\begin{proof}
The actual completed energy inverses, normalized
row kernels and readouts have exponential cell
bounds. Their local scalar and finite-\(p\)
upgrades retain the cell weight, since the
ratios \(w_z/w_y\) are bounded on each fixed
complete neighbouring box. This proves the
weighted canonical strong row-saddle inverse.
For \(\Lambda:\mathcal A_\alpha\to\mathcal Y_{2,\alpha}\),
use its weighted zero-average energy estimate
and the completed source equation
\(\Delta_U\lambda=RD^\sharp A\).
The projection \(R\) preserves the weighted
\(C^\alpha\) source norm by its completed
exponential column sums. The local Schauder
estimate then gives the weighted site bound.
The same argument gives the weighted first
adjoint bound in \eqref{cng:adjoint-force-bounds};
the finite-\(p\) force-restoration proof gives
the weighted second bound.

In the derivative of every local residual
keep one perturbation slot weighted and all
background slots in the unweighted small ball.
All point products, complete row moments and
the adjoint factory have the same bounds with
that one weight. The derivative of the
preconditioned fixed-point map therefore has
weighted norm at most \(C_{p,\alpha}r<1/2\).
Its derivative equation gives
\eqref{cng:weighted-response}. The explicit
outward datum, finite gauge chain, node
extension and current update have the same
one-slot weighted bounds.

These estimates hold on a fixed larger
complex datum ball. Cauchy in its other
slots gives the corresponding mixed derivative
with one weighted slot and \(m-1\) unweighted
slots. Choose \(\psi\) to be distance from the
carrier of any selected point datum. Symmetry
of the \(m\) derivative slots allows the
farthest such datum to be selected. This
proves \eqref{cng:localized-jets}.
The use of the completed source equation in
the Schauder step is essential; no individual
walk term was assigned a \(C^{2,\alpha}\)
site bound.
\end{proof}

This construction supplies a strong analytic
auxiliary germ, a literal mean-preserving
field/current update, and exact physical
matching on the admitted critical graph.
It keeps the \(H_1\), \(SH_1\) and outward
datum first jets distinct.
It does not construct a reference from
arbitrary curvature-only coarse data,
identify an externally prescribed off-current
core convention, or establish full proper
history domains and deletion laws.
Independent enlarged cube inputs and a
complete nonlinear connected-carrier norm
remain separate requirements.

% END LAYERED EXPANSION inverse/canonical_nonlinear_germ.tex

% BEGIN LAYERED EXPANSION inverse/biharmonic_restoration_carriers.tex
\subsection{Local biharmonic formulas for strong gauge restoration}
\label{ch:bhr-biharmonic-restoration}

The completed restoration \(G_sR\) gains two H\"older
derivatives on a H\"older source. A decomposition through
second-order site inverses need not have that property term
by term: an individual averaging-mass source can jump at a
pillar face. We obtain a different local decomposition from
a fourth-order positive operator. Its auxiliary Laplacian
absorbs that jump before the final Schauder estimate.

The construction uses the actual site graph, own-pillar
averages and physical measures. It neither clips an average
at an artificial boundary nor replaces the determining
region by a periodic model. The final carrier scale is
chosen separately from a prescribed history scale.

\subsubsection{The operator and its physical normalization}

Retain the geometry, comparison length \(\ell\), assigned
site scales \(\xi\), and actual covariant site Laplacian
\(\Delta=D_U^\sharp D_U\) of
\cref{ch:layered-site-restoration}. Let \(q\) take values
in raw top coordinates equipped, in this section, with
counting measure. Its adjoint for that measure is
\begin{equation}\label{bhr:counting-adjoint}
                 q^*=q^\sharp\Xi^{-4}.
\end{equation}
The sharp on the right is the physical top transpose.
Consequently the local fourth-order operator is
\begin{equation}\label{bhr:fourth-order-operator}
 K_4=\Delta^2+q^*q
      =\Delta^2+q^\sharp\Xi^{-4}q,
                  \qquad G_4=K_4^{-1}.
\end{equation}
At real unitary coefficients its form is
\[
 a_4(u,v)=\langle\Delta u,\Delta v\rangle_{
 \mathrm{fine}}+\langle qu,qv\rangle_{\mathrm{count}}.
\]
All analytic continuations use the bilinear form.
In particular \(q^*\) in an analytic formula is not a
conjugate transpose of a complex coefficient.

Use the graph norm
\begin{equation}\label{bhr:H2-energy}
 \|u\|_{V_2}^2=\|\Delta u\|_2^2
        +\|\ell^{-1}D_Uu\|_2^2+\|\ell^{-2}u\|_2^2.
\end{equation}
It controls the full covariant second differences on
complete interior stencils. At the true boundary we use
the original graph and the bare atomic norm; the latter
pays any omitted-neighbour stencil. Local norms \(V_{2,y}\)
use a fixed complete enlargement of the assigned cell.

\begin{lemma}[A site lift with two energy derivatives]
\label{lem:bhr-H2-site-lift}
For real curvature-small admitted data there is a local
right inverse \(E\) of the actual \(q\) such that
\begin{equation}\label{bhr:H2-site-lift}
                qE=1,\qquad
             \|Eb\|_{V_2}\le C\|b\|_{\ell^2(\mathrm{count})}.
\end{equation}
The support and collars of its seeds have bounded overlap.
No physical-current bound is needed for this energy estimate.
\end{lemma}
\begin{proof}
Use the normalized compact bump seed from
\cref{lem:lsi-weighted-site-lift} in a local real
Coulomb frame on a complete nonbare cube of side
\(C_L\xi_y\). Its normalization matrix is the actual
transported bump average and is uniformly invertible.
The supremum of the connection is \(C_L\kappa/\xi_y\).
The finite-\(p\) Coulomb estimate, used at \(p=2\), gives
\(\|Da\|_{2}\le C_L\kappa\) on the rescaled cube.
These facts follow from curvature alone by
\cref{dcz:Coulomb-Lp}.

For a seed of value \(b_y\), the ordinary second
differences of the scalar bump have \(L^2\) norm
\(C|b_y|\). Expanding the covariant second differences
adds \(aD\chi\), \((Da)\chi\) and \(a^2\chi\), with
the same bound. No pointwise bound for \(Da\) is used.
The first and zeroth weighted norms in
\eqref{bhr:H2-energy} have that bound by the existing
site-lift estimates. Its Laplacian is a sum of these
second differences. The physical support volume is
\(O(\xi_y^4)\), so all four powers of length cancel
against the two derivatives. This is precisely a
counting-top bound, rather than a physical-top bound.

For the finitely coarse meshes, including a bare
pillar, use the original transported atomic seed.
Here \(\eta/\xi_y\) has a fixed positive lower bound;
each second difference is bounded by \(C\eta^{-2}\)
times its value and the physical volume is
\(O(\xi_y^4)\). Thus the same bound holds. For fine
meshes the compact bump and its two-step collar lie
inside its pillar. The bounded overlap and adjacent
scale comparability sum these estimates.

If a frame cube does not embed, perform this entire
seed construction on a sufficiently large periodic
cover and average over simultaneous deck shifts.
The exact equation \(qE=1\), the quadratic energy
bound and its support description descend. This is
an averaging of linear maps, not an assertion about
partition functions on different volumes.
\end{proof}

\begin{proposition}[The full fourth-order gap]
\label{prop:bhr-full-H2-gap}
For the same real data,
\begin{equation}\label{bhr:H2-gap}
       c\|u\|_{V_2}^2\le a_4(u,u)\le C\|u\|_{V_2}^2.
\end{equation}
The constants are uniform in depth, compatible volume
and admitted proper-region shape. The corresponding
zero-extension Dirichlet compressions have the same
gap in the energy of their zero extension.
\end{proposition}
\begin{proof}
Write \(u=u_0+Equ\), where \(qu_0=0\). The exact
zero-average estimate \eqref{lsi:zero-average-estimate}
gives
\[
 \|u_0\|_{V_2}\le C\|\Delta u_0\|_2
       \le C\{\|\Delta u\|_2+\|qu\|_{\mathrm{count}}\}.
\]
Use \cref{lem:bhr-H2-site-lift} for the other term.
This proves the lower bound. Jensen's inequality
for the transported average gives
\[
 \sum_y|qu(y)|^2
       \le C\sum_y\xi_y^{-4}\|u\|_{2,B_y}^2
       \le C\|\ell^{-2}u\|_2^2,
\]
which proves the upper bound. For the component
second derivatives apply the zero-average Hodge
estimate \eqref{lsi:second-derivative-hodge} to
\(u_0\), and the seed bound to \(Equ\). The bare
mass controls its two-step outer collar.

If \(E_D\) injects site fields supported in a fine
set \(D\), the compression is
\(E_D^\sharp K_4E_D\). Its form is
\(\|\Delta E_Du\|_2^2+\|qE_Du\|_{\mathrm{count}}^2\).
It retains the full Laplacian and every full crossing
row evaluated on the zero extension. Restricting the
already proved inequality proves the assertion;
there is no partial-row inf-sup assumption.
\end{proof}

For the strong estimates fix \(0<\alpha<1\) and a
class \(\mathscr U_{\alpha,C_0}\) of these real
connections whose local Coulomb potentials on the
complete boxes have scaled \(C^{1,\alpha}\) norm
at most \(C_0\). The current-regular references used
for the nonlinear germ of
\cref{ch:cng-canonical-nonlinear-germ} belong to this
class. Write \(\mathcal A_\alpha\) for the fine strong
space of \eqref{psr:strong-norm}. We work on a fixed
smaller weighted \(C^{1,\alpha}\) tube about a
chosen member. All constants below may depend on
\(\alpha,C_0,L\). Curvature alone is the premise
of \eqref{bhr:H2-gap}; it is not being asserted to
give these stronger coefficient bounds.
On the complex tube, the Banach norms are measured
in the fixed real-reference comparison frames.
Covariant norms at the varying base are uniformly
equivalent on that tube.

\subsubsection{A source class suited to a local fourth-order inverse}

The source must allow two derivatives of a H\"older
function without requiring their pointwise bound.
Let \(\mathcal S_\alpha\) consist of finite site
arrays with a representation
\begin{equation}\label{bhr:source-class}
                 \rho=\Delta f+D_U^\sharp g+h,
\end{equation}
where
\[
 \|f\|_{\alpha,2}
 =\sup_y\{\xi_y^2\|f\|_{\infty,B_y^+}
              +\xi_y^{2+\alpha}[f]_{\alpha,B_y^+}^{U}\},
 \quad
 \|g\|_{\infty,3}=\sup_b\mathsf L(b)^3|g(b)|,
 \quad
 \|h\|_{\infty,4}=\sup_x\xi_x^4|h(x)|.
\]
The H\"older comparisons use complete boxes across
pillar faces. The norm \(\|\rho\|_{\mathcal S_\alpha}\)
is the infimum of the sum of these three norms over
all representations. Local source spaces
\(\mathcal S_{\alpha,z}\) use representations
supported on a fixed complete enlargement of \(B_z\).
The output space is \(Y^{2,\alpha}\) of
\eqref{dsc:site-gauge-norm}.
The source norms at nearby bases are uniformly
equivalent to the fixed-reference norm. Indeed
\((\Delta_U-\Delta_*)f\) can be put in
\(D_*^\sharp g_1+h_1\) using bounded values of
\(f\) and the connection difference and its first
difference. Likewise \((D_U^\sharp-D_*^\sharp)g\)
is a grade-four point source. Thus this comparison
does not differentiate the H\"older function \(f\)
or the bounded field \(g\).

Smooth localization is bounded on this source class.
Indeed \(\chi\Delta f=\Delta(\chi f)-[\Delta,\chi]f\).
The last commutator can be put in the form
\(D^\sharp g_\chi+h_\chi\), with coefficients
\(D\chi\) and \(D^2\chi\) multiplying shifted values
of \(f\). It differentiates no \(f\). The same
product rule applies to \(\chi D^\sharp g\).
Consequently a fixed smooth cell partition localizes
\(\mathcal S_\alpha\) with bounded overlap and bounded
norm. We use this localization for block products;
sharp characteristic functions are not multipliers
on the H\"older part of the source space.
Use its bounded reconstruction and restriction
maps in the cell-block norm. Multiplying that norm
by their fixed bound makes the product norm
submultiplicative; this constant is absorbed in
the local commutator constant below.

Every local source is a bounded functional on \(V_2\).
For example, for a unit source at scale \(\xi_z\),
\[
 \|f\|_2\le C,\qquad
 \|g\|_2\le C\xi_z^{-1},\qquad
 \|h\|_2\le C\xi_z^{-2}.
\]
Pair these with \(\Delta u,D_Uu,u\), respectively.
Thus its dual energy norm is uniformly bounded.

\begin{lemma}[Strong local fourth-order columns]
\label{lem:bhr-strong-local-columns}
Let \(I_D=E_D(E_D^\sharp K_4E_D)^{-1}E_D^\sharp\).
On a region strictly inside its artificial boundary,
\begin{equation}\label{bhr:strong-column}
 \|I_D\rho_z\|_{Y^{2,\alpha}(B_y^+)}
           \le C_\alpha e^{-s_0d(y,z)}
                             \|\rho_z\|_{\mathcal S_{\alpha,z}}.
\end{equation}
The interior collar contains all stencils and all
complete averaging carriers used in the estimate.
The same statement holds at a true boundary with
the bare atomic convention. It holds uniformly on
the stated smaller complex \(C^{1,\alpha}\) tube.
\end{lemma}
\begin{proof}
First obtain a weighted energy estimate. A raw metric
distance is only Lipschitz and is insufficient for
this fourth-order step. Smooth the truncated distance
from \(B_z\) at the local \(\ell\) scale using the
weighted smooth partition of
\cref{lpw:slow-partition}, with a fixed radius
multiple in place of \(m\). If its squared normalized
weights are \(p_x(y)\), use
\(\psi(y)=\sum_xp_x(y)d(x,B_z)\).
On an overlapping support, centre distances differ
by at most \(C_L\). Subtract one such constant before
differentiating \(\sum_xp_x=1\). The cutoff derivative
estimates, extended through order three by the same
bump calculation, give
\[
 |\psi-d(\cdot,B_z)|\le C_L,\quad
 |D\psi|\le C_L\ell^{-1},\quad
 |D^2\psi|\le C_L\ell^{-2}.
\]
For \(w=e^{s\psi}\), the exact covariant product rule
therefore gives
\[
 \|w^{-1}\Delta(wu)-\Delta u\|_2
       \le C_L|s|\{\|\ell^{-1}D_Uu\|_2
                                  +\|\ell^{-2}u\|_2\}.
\]
The oscillation of \(\log w\) on a full averaging
carrier is \(O_L(s)\). The conjugated bilinear form
\(a_4(w^{-1}u,wv)\) differs from \(a_4(u,v)\)
by at most \(C_L|s|\|u\|_{V_2}\|v\|_{V_2}\).
Choose \(s\) below the gap constant. Multiplication
by \(w\) preserves the zero-extension space, so the
weighted estimate holds for every compression.
It gives
\begin{equation}\label{bhr:energy-columns}
       \|I_D\rho_z\|_{V_{2,y}}
           \le Ce^{-s_1d(y,z)}
                                  \|\rho_z\|_{\mathcal S_{\alpha,z}}.
\end{equation}
There is no endpoint scale ratio: both the source
dual norm and the site \(H^2\) norm are dimensionless
in four dimensions.

Now put \(u=I_D\rho_z\) and \(v=\Delta u-f\).
On a complete interior box its exact equation is
\begin{equation}\label{bhr:auxiliary-second-order-equation}
       \Delta v=D_U^\sharp g+h-q^*qu,
                           \qquad \Delta u=f+v.
\end{equation}
The local \(L^2\) norm of \(v\) is controlled by
\eqref{bhr:energy-columns} and the local norm of
\(f\). The actual average satisfies
\(|qu(y)|\le C\xi_y^{-2}\|u\|_{2,B_y}\).
The physical density of \(q^*\) then bounds
\(q^*qu\) by \(C\xi_y^{-4}\) times the nearby
energy. Complete carriers change only fixed local
constants.

Rescale this box to length one. The finite-\(p\)
divergence-source bootstrap used in
\cref{prop:psr-force-restoration} bounds
\(v\) in \(W^{1,p}\) for every finite \(p\), using
its local \(L^2\) norm, \(g\) in \(L^\infty\), and
the bounded scalar forcing. First testing its
equation against a squared-cutoff multiple of
\(v\) supplies the local \(H^1\) bound from those
same quantities. On the complex tube use the
real-reference comparison in this energy test.
Lower-order covariant
coefficients have their admitted bounds. Choose
\(p>4/(1-\alpha)\). Morrey gives a \(C^\alpha\)
bound for \(v\); no derivative of \(g\) is taken.
Thus \(f+v\) is a \(C^\alpha\) source for the second
equation. The scalar finite-\(p\) bootstrap first
gives the lower \(C^{1,\alpha}\) norm of \(u\),
and the interior Schauder estimate
\eqref{dsc:covariant-interior} gives its full
\(C^{2,\alpha}\) norm. Its coefficient premise is
exactly the displayed \(C^{1,\alpha}\) background
class. Applying this on fixed nested boxes preserves
the exponential factor in \eqref{bhr:energy-columns}.

At the true boundary \(\xi=\eta\). Each required
cell has a fixed number of sites; its rescaled
values, differences and H\"older seminorm are
bounded directly by its local \(H^2\) energy.
No boundary Schauder theorem for a rough region is
being inferred. Artificial boundaries are outside
all boxes on which the strong assertion is used.

Finally, a small complex \(C^1\) perturbation changes
\(\Delta:V_2\to L^2\) and
\(q:V_2\to\ell^2(\mathrm{count})\) by a uniformly
small operator. The weighted forms and every
compression remain invertible by Neumann series.
The stronger \(C^{1,\alpha}\) tube bounds the
coefficients of the two scalar equations above in
their required H\"older norms. The same estimates
then prove \eqref{bhr:strong-column} there. This
uses bilinear perturbation of real positivity,
rather than complex Hermitian positivity.
\end{proof}

The gain in this lemma explains the fourth-order
choice. The discontinuous density \(q^*qu\) is a
source for \(v\), so \(v\) is H\"older. It is not a
discontinuous source for the final Schauder equation
\(\Delta u=f+v\).

\subsubsection{A strong local walk for \(G_4\)}

Use the actual weighted cover
\eqref{lpw:weighted-partition}, with
\(r_x=m\ell(x)\) and \(m\delta\) sufficiently small.
Its smooth bump proof gives, for every fixed order,
\(|D^a\chi_x|\le C_ar_x^{-a}\).
Take a compression patch \(D_x\) of radius \(4r_x\)
and retain all full crossing stencils, rows and
geometry-decision collars. The support of \(\chi_x\)
and its fixed complete cell enlargements lie
strictly inside that artificial boundary. Put
\(I_x=I_{D_x}\).

\begin{lemma}[The fourth-order cutoff error]
\label{lem:bhr-slow-commutator}
The error columns obey
\begin{equation}\label{bhr:strong-commutator}
 \|[K_4,\chi_x]I_x\chi_x\|_{
       \mathcal S_{\alpha,z}\to\mathcal S_{\alpha,y}}
                \le \frac{C_\alpha}{m}e^{-s_2d(y,z)}.
\end{equation}
Their representations and coefficients read only
the declared local compression data and collars.
\end{lemma}
\begin{proof}
For a local strong output \(u\), use the exact identity
\[
 [\Delta^2,\chi]u
       =\Delta[\Delta,\chi]u+[\Delta,\chi]\Delta u.
\]
The first term has the source form \(\Delta f_\chi\),
where \(f_\chi=[\Delta,\chi]u\). The covariant product
rule and \(u\in Y^{2,\alpha}\) show
\[
             \|f_\chi\|_{\alpha,2}
                              \le C_\alpha m^{-1}\|u\|_{Y^{2,\alpha}}.
\]
Only one cutoff derivative is needed for the leading
term; higher cutoff derivatives are smaller.
For the second term apply summation by parts to the
commutator itself, rather than differentiating
\(\Delta u\). For any site array \(w\), the exact
two-shift product rule writes
\[
 [\Delta,\chi]w=D_U^\sharp g_\chi(w)+h_\chi(w),
\]
where the coefficients are shifted values of
\(D\chi\) and \(D^2\chi\) times transported values
of \(w\). Explicitly, for the orientation
\(D_Uw(b)=\eta^{-1}(w_- -T_bw_+)\), take
\[
 g_\chi(w)(b)=(D\chi)(b)(w_-+T_bw_+),\qquad
 h_\chi(w)(x)=-(\Delta_{\mathrm{scalar}}\chi)(x)w(x).
\]
The scalar Laplacian uses exactly the edges of
the original graph. With \(w=\Delta u\), its weighted
\(L^\infty\) size has grade two. Therefore
\[
 \|g_\chi(\Delta u)\|_{\infty,3}
       +\|h_\chi(\Delta u)\|_{\infty,4}
                         \le C_\alpha m^{-1}\|u\|_{Y^{2,\alpha}}.
\]
This step contains no third difference of \(u\).
The summation-by-parts formula holds on the original
graph, so missing outer edges do not create an
unrecorded boundary term.

For \([q^*q,\chi]u\), insert the cutoff value at
the row top. Each difference on the complete row
carrier is at most \(C_L/m\). The actual physical
row and column density bounds give a point source
of grade four with norm at most
\(C_Lm^{-1}\|u\|_{Y^{2,\alpha}}\). Bare identity
rows commute exactly. At a rough true boundary all
remaining fourth-order entries have size
\(O(\eta^{-4})\); the atomic output and a cutoff
difference give the same source bound directly.

Apply \cref{lem:bhr-strong-local-columns} on the
cutoff support and its collars. Input localization
is bounded on \(\mathcal S_\alpha\), and all local
weights are comparable. This proves
\eqref{bhr:strong-commutator} with a smaller decay
rate. Every formula uses the original compressed
matrix and its declared finite carriers.
\end{proof}

Define
\begin{equation}\label{bhr:strong-parametrix}
 \mathcal M_4=\sum_x\omega_x\chi_xI_x\chi_x,
 \qquad T_4=\sum_x\omega_x[K_4,\chi_x]I_x\chi_x.
\end{equation}
Compression and the strict support collar give
\(\chi_xK_4I_x\chi_x=\chi_x^2\) as exact finite
operators. Thus
\begin{equation}\label{bhr:exact-parametrix}
                         K_4\mathcal M_4=1+T_4.
\end{equation}
In smooth cell blocks, use the maximum of the row
and column sums with weights \(e^{sd(y,z)}\).
The weighted overlap \eqref{lpw:weighted-overlap}
and the local column estimates give
\[
 \|\mathcal M_4\|_{
       \mathcal S_\alpha\to Y^{2,\alpha};\mathrm{Sch},s}\le C_\alpha,
 \qquad
 \|T_4\|_{
       \mathcal S_\alpha\to\mathcal S_\alpha;\mathrm{Sch},s}
                                               \le C_\alpha/m.
\]
For fixed output or input cell the sum of contributing
centre weights is bounded. The remaining cell sum
is exponential. A whole-patch energy norm, which
would introduce a power of \(m\), has not replaced
these column estimates.

File each local matrix on its full carrier as in
the positive-walk construction. Include its patch,
full crossing-row collar, local geometry-decision
collar, canonical input/output tags and fixed short
joins. For sufficiently large carrier cube integer
\(M_0\), each local factor uses at most \(B_L\)
tags, independently of \(m\). With
\(\sigma_0(X)=b_L+d_{c,0}(X)\), include the short
join in \(X\vee Y\), so
\(\sigma_0(X\vee Y)\le\sigma_0(X)+\sigma_0(Y)\).
For any prescribed finite rate \(\kappa_1\), choose
parameters in the order
\[
 \alpha,C_0,\kappa_1;\quad m;\quad
       \delta;\quad M_0;\quad\rho,
 \qquad (C_\alpha/m)e^{\kappa_1B_L}<1/2.
\]
Here \(m\delta\) is small, \(M_0\ge C_L/\delta\),
and \(\rho\) pays the full collars and metric
growth reserve. Only after these choices shrink
the common complex tube. There is no dependence
on depth or volume.

\begin{proposition}[Strong fourth-order inverse carriers]
\label{prop:bhr-fourth-order-carriers}
The exact series
\begin{equation}\label{bhr:G4-walk}
                 G_4=\mathcal M_4\sum_{n\ge0}(-T_4)^n
\end{equation}
is absolutely convergent in the strong carrier
Schur norm from \(\mathcal S_\alpha\) to
\(Y^{2,\alpha}\), at the chosen carrier rate.
Every coefficient is holomorphic on its complete
local background \(C^{1,\alpha}\) tube and depends
only on its declared local determining data.
\end{proposition}
\begin{proof}
The source-to-source family norm of \(T_4\) is
less than \(1/2\). Joined carrier products are
submultiplicative, so the series converges with
the asserted output norm. Multiplication by
\(K_4\) and \eqref{bhr:exact-parametrix} identify
its sum with the positive energy inverse.
Every local inverse is a finite Dirichlet matrix
with its full coefficients. Its uniform complex
invertibility and the strong column estimate give
holomorphy into the displayed local Banach spaces.
Every output factor includes its smooth cutoff,
so zero extension across its artificial boundary
preserves \(Y^{2,\alpha}\).

The coefficient difference under a small
\(C^{1,\alpha}\) background change is bounded in
the same family norm. Choose the final tube inside
a slightly larger common complex tube on which the
weighted strong local columns and cutoff commutators
have the bounds just proved. Their finite matrices
are holomorphic, and the source-space equivalence
identifies their input spaces with fixed Banach
spaces. Banach-valued Cauchy on the larger tube
therefore bounds the change of each cutoff local
column and error column by \(C\varepsilon\), with
the same exponential weight and full local carrier.
The weighted centre sum and the convergent walk
preserve this bound in the absolute family norm.
This argument uses strong regularity only on the
output and commutator collars. It does not assign
a global \(Y^{2,\alpha}\) norm to an uncut Dirichlet
solution across its artificial boundary.
\end{proof}

\subsubsection{A positive coarse inverse and the exact restoration}

Define on counting top space
\begin{equation}\label{bhr:coarse-M4}
                  M_4=qG_4q^*
                         =qG_4q^\sharp\Xi^{-4}.
\end{equation}
This is already in counting normalization.
It is not obtained by treating a physical
transpose as an ordinary transpose.

\begin{lemma}[The fourth-order coarse gap]
\label{lem:bhr-coarse-gap}
At real coefficients,
\begin{equation}\label{bhr:M4-gap}
                              cI\le M_4\le I.
\end{equation}
It has an actual local carrier family in counting
coordinates at any required fixed intermediate rate.
On a smaller common complex tube its inverse has
the local carrier family supplied by
\cref{prop:cci-positive-coarse-inverse}.
\end{lemma}
\begin{proof}
Set \(u=G_4q^*b\). Then
\(a_4(u,u)=\langle b,M_4b\rangle\).
The exact lift \(E\) gives
\[
 \|b\|^2=\langle q^*b,Eb\rangle
       =a_4(u,Eb)\le C a_4(u,u)^{1/2}\|b\|.
\]
This proves the lower bound. For the upper bound,
\(K_4\ge q^*q\); equivalently
\(\|qK_4^{-1/2}\|\le1\), giving \(M_4\le I\).

A unit counting top source \(q^*e_z\) has actual
point density at most \(C\xi_z^{-4}\) on its
complete pillar, so it is a unit grade-four
\(h\) source. The strong fourth-order family
sends it to a dimensionless site field; the
actual output average \(q\) is bounded on that
field. Its input/output columns therefore have
no missing physical length factors. Composition
gives the actual family of \(M_4\), with
exponentially summable cell blocks and its
declared full carriers. The number of top
species is fixed.

Apply the positive coarse inversion algorithm
of \cref{prop:cci-positive-coarse-inverse} to
this matrix. Request the primitive family at a
stronger intermediate rate first. Full-carrier
truncation, combined short powers and reblocking
then give the inverse at the requested final
rate, with its finite parameter and separation
choices. The family perturbation bound proved
above supplies its common complex tube. No
indefinite local compression is used.
\end{proof}

For a site source \(f\), put
\begin{equation}\label{bhr:exact-T4}
 \mathcal T_4 f
   =G_4\Delta f-G_4q^*M_4^{-1}qG_4\Delta f.
\end{equation}
The notation \(\mathcal T_4\) is distinct from
the parametrix error \(T_4\).

\begin{proposition}[Exact constrained least squares]
\label{prop:bhr-exact-restoration}
For every site source on the finite graph,
\begin{equation}\label{bhr:GsR-identity}
                          \mathcal T_4 f=G_sRf.
\end{equation}
Consequently
\begin{equation}\label{bhr:Lambda-identity}
 \Lambda=\{G_4-G_4q^*M_4^{-1}qG_4\}\Delta D_U^\sharp.
\end{equation}
These equalities hold also on the admitted
bilinear complex continuation.
\end{proposition}
\begin{proof}
At real data the unique minimum of
\(\|\Delta\lambda-f\|_2^2\) subject to
\(q\lambda=0\) exists by the zero-average gap.
Adding \(\|q\lambda\|^2\) to the normal operator
does not change its constrained equation.
Its saddle system is
\[
 K_4\lambda+q^*\mu=\Delta f,
                     \qquad q\lambda=0.
\]
Its Schur elimination is exactly
\eqref{bhr:exact-T4}. The residual
\(\Delta\lambda-f\) is orthogonal to
\(\Delta\ker q\). The actual identities
\(qG_sR=0\) and \(\Delta G_sR=R\) show
\(\operatorname{ran}R\subset\Delta\ker q\).
The zero-average gap makes \(\Delta\) injective
on \(\ker q\), and both spaces have dimension
equal to the site dimension minus the top dimension.
Hence \(\Delta\ker q=\operatorname{ran}R\).
Since \(\Delta\lambda\) belongs to that range,
the normal equation gives \(\Delta\lambda=Rf\).
Its zero-average solution is \(G_sRf\).

All matrices are holomorphic on the common
complex tube and their admitted inverses
remain defined. Alternatively replace real
orthogonality by the identities of the
bilinear projection \(R\). Either argument
extends the exact formula. Finally set
\(f=D_U^\sharp A\).
\end{proof}

\begin{proposition}[Local strong restoration coefficients]
\label{prop:bhr-strong-restoration-family}
Fix \(\alpha,C_0\) and a finite requested carrier
rate \(\kappa\). After the fixed parameter,
separation and tube choices above, \(G_sR\)
has an exact local family from
\(\mathcal C_{\alpha,2}\) to \(Y^{2,\alpha}\),
and \(\Lambda\) has an exact local family from
\(\mathcal A_\alpha\) to \(Y^{2,\alpha}\).
In smooth complete cell blocks these families
satisfy
\begin{equation}\label{bhr:strong-family-bound}
 \max\left\{
 \sup_y\sum_{z,X} e^{sd(y,z)+\kappa\sigma_c(X)}
                   \|\mathcal T_{4,yz}[X]\|_{\alpha,2\to Y^{2,\alpha}},
 \sup_z\sum_{y,X} e^{sd(y,z)+\kappa\sigma_c(X)}
                   \|\mathcal T_{4,yz}[X]\|_{\alpha,2\to Y^{2,\alpha}}
                         \right\}\le C_{\alpha,\kappa}.
\end{equation}
The analogous bound holds for \(\Lambda\) with
the strong field norm in its input. Each
coefficient, including its extended output,
has the stated \(C^{2,\alpha}\) site norm.
Its holomorphic dependence reads only its full
local background and determining-data carrier.
The constants are uniform in depth and volume.
\end{proposition}
\begin{proof}
The injection \(f\mapsto\Delta f\) is bounded
from \(\mathcal C_{\alpha,2}\) to
\(\mathcal S_\alpha\) by definition. Compose
the strong fourth-order family, the local
\(q,q^*\) maps and the counting inverse
\(M_4^{-1}\) in \eqref{bhr:exact-T4}.
Keep full endpoint tags and short joins in
each product. Their carrier norm is
submultiplicative and the stronger intermediate
rate pays the one coarse reblocking. The
final carrier cube integer \(M_c\) and its
separation are chosen after these finitely
many operations. Smooth cell localization
preserves the displayed source and output
spaces. This proves the first family bound.

The first divergence maps the strong field
space to the complete-box source space
\(\mathcal C_{\alpha,2}\). Its finite stencils
give a bounded local input family. Applying
\eqref{bhr:Lambda-identity} proves the second
assertion, without taking a second difference
of \(A\). The exact identities identify the
sums with the original restoration. A
coefficient of this new family need not obey
\(q\lambda=0\) by itself. Its strong bound
comes from the fourth-order local equation,
not cancellation of its averaging-mass source.
\end{proof}

All primitive patches, frame cubes and the
aligned coarse inversion grids are first
constructed on one sufficiently large periodic
cover. Perform the entire matrix construction
there before any carrier projection. The exact
operators commute with deck shifts; average
their families over those shifts and descend.
The \(C^{2,\alpha}\) local norms are convex,
and lifted endpoint sums cover the original
Schur sums, so no sheet factor is introduced.
Retain every lifted carrier before projection,
including winding paths and the full data
collar. Bilinear matrix identities and inverses
restrict to the deck-invariant sector by
uniqueness. This proves the small-period case
for both the primitive and coarse stages.

The new family supplies coefficientwise strong
restoration in the fixed \(C^{1,\alpha}\) chart.
It does not upgrade the old second-order walk
coefficients, or remove the need to construct
connected coefficient families for the nonlinear
germ of \cref{ch:cng-canonical-nonlinear-germ}.
Independent background cube domains, complete
history admission and an externally prescribed
off-current core convention remain separate.

% END LAYERED EXPANSION inverse/biharmonic_restoration_carriers.tex

% BEGIN LAYERED EXPANSION strong_current_families.tex
\subsection{Strong local families for the current row saddle}
\label{scf:section}

The energy families alone do not give
a strong Banach space in which to
expand a nonlinear equation.
We construct the required strong
families in the same prescribed order
as the current precisions. The
fourth-order restoration in
\cref{prop:bhr-strong-restoration-family}
is used for its new local
\(C^{2,\alpha}\) site coefficients.
No such conclusion is assigned to
the older second-order site pieces.

Fix \(0<\alpha<1\), a finite
\(p>4/(1-\alpha)\), and a finite
requested carrier rate \(\kappa\).
Let \(U_*\) be a real current-regular
reference on an admitted determining
geometry, with strict curvature
and current margins. Use a fixed
smaller weighted \(C^{1,\alpha}\)
tube \(U=e^{i\eta E}U_*\).
The independent current \(J\) is
bounded in its full amplitude norm.
The actual \(j_U\) need not have
that bound at a complex input.
All spaces and their comparison
frames are fixed at \(U_*\).

Write \(\mathcal A_\alpha\) for
the strong field space of CNG,
\(\mathcal Y_{2,\alpha}\) for its
site gauge space, and
\(\mathcal F_p\) for the force
space of
\cref{ch:proper-strong-response}.
More explicitly the force
cell norm is
\begin{equation}
             \|f\|_{\mathcal F_{p,y}}
                    =\xi_y^{\,3-4/p}\|f\|_{L^p(B_y^+)} .
 \label{scf:force-cell-norm}
\end{equation}
The cell enlargements include all
stencils and row collars.
A family norm is the maximum
of its weighted row and column
Schur sums between these normalized
cell Banach spaces, with the
joined full carriers of
\eqref{ncw:family-norm}.
Use a smooth local partition
to reconstruct strong fields
from their cell slots.
Its profiles, differences and
H\"older differences have their
scaled bounds on complete collars.
The product rule bounds this
reconstruction independently of
the number of fine coordinates.
The fixed reconstruction constant
is included in the algebra norm.

\subsubsection{The positive field primitive on finite-\(p\) forces}

Let
\[
                  K_F=H_U+D_UD_U^\sharp+Q_U^\sharp a_fQ_U
\]
be the actual positive field
operator of the local-walk
construction, before the site
projection and current terms.
Its local compressions retain
all crossing coefficients.
For a compression domain \(D\),
let \(I_D\) be its inverse
and let \(\chi\) be an output
cutoff separated from its
artificial boundary.

\begin{lemma}\label{scf:positive-field-columns}
For a force supported in an
assigned cell \(z\), the local
output factors satisfy
\begin{equation}
 \|\chi I_D\chi\,1_zf\|_{\mathcal A_{\alpha,y}}
       \le C_{p,\alpha}e^{-s d(y,z)}
                                      \|f\|_{\mathcal F_{p,z}} .
 \label{scf:local-field-strong-column}
\end{equation}
The constants are independent
of compression size, fine mesh,
depth and region shape.
The same estimate holds on
a common smaller complex tube.
\end{lemma}

\begin{proof}
For \(p>2\), the normalized
force cell embeds in the
\(L^2\) force cell, by the
physical volume \(O(\xi_z^4)\).
The positive local inverse
therefore has its proved
exponential energy column
with this source norm.
On a complete nonbare box,
the exact Wilson curl--divergence
stencil rewrites the equation
as a component covariant
Laplacian with a bounded
curvature multiplication and
fixed neighboring shifts.
The row-mass term has locally
bounded point forcing:
the complete average of the
solution is bounded by its
energy cell \(L^2\) norm,
and the actual transpose
extends its bounded density.
Its bound uses the physical
column density, not a
point evaluation of that
energy function.

In a real local Coulomb
frame the potential and
its first derivative are
bounded by the stated
current regularity. In the
complex tube they have the
same bounds with fixed
larger constants.
The component covariant
Laplacian is the ordinary
discrete Laplacian plus
first-order and zeroth-order
terms with these bounded
coefficients. This follows
by pairing the positive
and negative covariant
shifts before dividing by
the fine spacing.
There is no derivative of
curvature in that identity.

Apply the local finite-\(p\)
estimates on nested complete
boxes. Energy gives \(H^1\);
the first scalar estimate
gives \(W^{2,2}\), then
strict embeddings give
gradient \(L^3\).
The next estimate gives
\(W^{2,3}\), hence gradient
\(L^8\) and field \(L^\infty\).
If \(p\le8\), apply the
final \(p\) estimate there.
For \(p>8\), use the
\(p=8\) estimate first,
then the desired finite
\(p\). The bounded row
forcing and curvature terms
remain controlled at every
stage. The original force
belongs to each lower
exponent used on these
fixed boxes.
Morrey then gives the
gradient H\"older exponent
\(\alpha<1-4/p\).

All outer energy terms
retain the exponential
column weight, and neighboring
box weights differ by a
fixed factor. The source
is present only near \(z\).
This proves the displayed
column estimate with a
fixed smaller rate.
The output cutoff avoids
the artificial boundary.
At a true bare boundary
the scale is \(\eta\);
the finite atomic stencils
and the local mass bound
give the same scaled
estimate directly.
Multiplication by the
smooth output cutoff
preserves the strong norm.
\end{proof}

The slow commutator in the
positive parametrix maps
this strong output into
\(\mathcal F_p\), with norm
\(C_{p,\alpha}/m\).
Its derivative terms are
\(D\chi\,Dv\) and
\(\Delta\chi\,v\);
the row commutator has
the complete row difference
\(\chi(x)-\chi(\text{row})\).
The slow-cutoff bounds and
the actual row moments
give the same \(m^{-1}\)
after rescaling.
The finite-shift curvature
terms also have that bound:
their commutators contain
one cutoff difference and
a bounded curvature
coefficient.
No derivative of the
input force is used.
Choose \(m\) large after
fixing \(p,\alpha,\kappa\).
The weighted-overlap sum
in the local-walk proof
then makes this an
absolute contraction on
\(\mathcal F_p\).
Exactly as in
\eqref{lpw:walk-series},
the resulting identity is
\begin{equation}
         \mathbf G_F=\mathbf M
                    \sum_{n\ge0}(-\mathbf T)^{\star n},
       \qquad
       \|\mathbf G_F\|_{\mathcal F_p\to\mathcal A_\alpha}
                                              \le C_{p,\alpha,\kappa}.
 \label{scf:strong-positive-field-family}
\end{equation}
The determining geometry and
background data of every
coefficient are retained.
This is a new strong
parametrix; it need not
use the same local pieces
as a previously fixed
energy representation.

\subsubsection{The two site adjoints without a current derivative}

For a site field \(u\), use
the normalized first-order
cell norm
\[
 \|u\|_{\mathcal V_{p,y}}=
  \xi_y^{\,2-4/p}\|u\|_{L^p(B_y^+)}
       +\xi_y^{\,3-4/p}\|D_Uu\|_{L^p(B_y^+)} .
\]
For a site source write
\(\rho=D_U^\sharp f+h\),
where \(f\) is a field
force and \(h\) a site
force. Give this source
the infimum norm of
\[
  \|f\|_{\mathcal F_{p,y}}
            +\xi_y^{\,4-4/p}\|h\|_{L^p(B_y^+)}
\]
over such representations
on the complete source
collar.
The divergence is variational.
Its values need not be
bounded fine functions.

\begin{lemma}\label{scf:site-force-families}
The actual completed site
inverse has local families
from this source class
into \(\mathcal V_p\).
It also has a local family
with
\begin{equation}
 \sup_y\{\xi_y^2|G_sh|_y+
                    \xi_y^3|D_UG_sh|_y\}
               \le C\|h\|_{\infty,4}
 \label{scf:site-bounded-force-family}
\end{equation}
in the corresponding
weighted row/column norms.
Both assertions hold on
the common complex tube.
\end{lemma}

\begin{proof}
For the local positive
site compression, first
use its variational energy
column with the displayed
source representation.
In the interior it solves
\[
       \Delta_Uu=D_U^\sharp f+h-q^\sharp a_squ .
\]
The last term has bounded
point amplitude of grade
four, from the local
energy norm and the
actual complete row.
For the divergence term,
the uniform finite-\(p\)
double-gradient resolvent
in \cref{dcz:uniform-Lp}
gives the local first-order
estimate. To see its
lower-order bootstrap
explicitly, put each
potential-times-gradient
term into a divergence
of potential times \(u\),
plus a bounded first
potential derivative
times \(u\).
Initially \(u\in H^1\),
hence \(u\in L^3\).
The first-order estimate
at \(3\) gives
\(u\in W^{1,3}\), hence
\(u\in L^8\).
If \(p\le8\), finish
at \(p\) at this stage:
the outer field belongs
to \(L^p\), as do the
prescribed sources.
If \(p>8\), its next
application at \(8\)
gives \(u\in L^\infty\).
The final finite-\(p\)
estimate then uses \(f,h\)
at their prescribed
exponent and the now
bounded lower-order
coefficients.
All cutoffs are nested
inside the artificial
compression boundary.

The scalar commutator
maps \(\mathcal V_p\)
to a site \(L^p\)
force of grade four
with norm \(C_p/m\).
Thus the same positive
parametrix, now in this
source class, converges
after fixing \(m\)
large enough.
Multiplication of a
variational source by
a cutoff is legitimate:
\(\chi D^\sharp f\)
is \(D^\sharp(\chi f)\)
plus a bounded one-step
cutoff difference times
\(f\). This stays in
the declared source
class, without \(Df\).
This proves its exact
local family assertion.

For the second assertion
take \(f=0\) and bounded
\(h\). The scalar
second-order bootstrap,
with the same bounded
row forcing, gives
\(W^{2,p_0}\) for any
fixed \(p_0>4\) on the
output support.
Morrey gives \(u\)
and its gradient bounded
with the displayed
scaling. The slow
commutator takes that
output into a bounded
site force, so its
positive local series
converges in this class
as well. True bare
cells use their atomic
stencils throughout.
\end{proof}

Let \(\widehat F=G_sq^\sharp\Xi^{-4}\)
and let \(\widehat M\)
be the actual counting
site Gram matrix of
\eqref{ncw:normalized-site-columns}.
Its inverse has the CCI
local family.
The injection
\(q^\sharp\Xi^{-4}\)
has point density of
grade four.
Thus
\eqref{scf:site-bounded-force-family}
gives local families for
\(\widehat F,D\widehat F\)
with site amplitude
grades two and three.
Pairing a grade-two
site field with a
grade-two column costs
no cell volume:
\(\xi^{-2}\xi^{-2}\xi^4=1\).
H\"older's inequality gives
the same statement in
\(\mathcal V_p\).
Consequently the exact
projection
\[
          P=\widehat F\widehat M^{-1}\widehat F^\sharp,
                       \qquad R=1-P
\]
has a bounded family on
\(\mathcal V_p\) and on
the displayed bounded
site-amplitude/gradient
space.
All scale ratios are
absorbed once in the
fixed spatial reserve,
as in NCW.

The exact adjoint formulas
now give the coefficient
families
\begin{equation}
 \begin{split}
 \Lambda^\sharp&=DRG_s:
                   \{\text{site amplitude four}\}
                         \longrightarrow\{\text{field amplitude three}\},\\
 S^\sharp&=1-DRG_sD^\sharp:
                                      \mathcal F_p\longrightarrow\mathcal F_p .
 \end{split}
 \label{scf:strong-adjoint-families}
\end{equation}
In particular a bounded
field force may be used
in the second line by
its cell \(L^p\) norm.
The exact identity
\(\Lambda^\sharp D^\sharp=1-S^\sharp\)
also gives the explicit
family bound
\begin{equation}
 \|\mathbf{\Lambda^\sharp D^\sharp}f\|_{\mathcal F_p}
                          \le C_{p,\kappa}\|f\|_{\infty,3}.
 \label{scf:bounded-divergence-adjoint}
\end{equation}
There is no assertion
that \(S^\sharp\) maps
arbitrary bounded forces
to bounded forces.
This is the finite-\(p\)
distinction required by
the nonlinear adjoint
factory.

\subsubsection{Woodbury before either current correction}

Use the new strong
family for \(G_F\)
and the actual positive
coarse family for
\(\widehat\Theta^{-1}\).
The Woodbury identity is
\[
 G_0=G_F+
        G_F(D\widehat F)\widehat\Theta^{-1}
                                (D\widehat F)^\sharp G_F .
\]
The map \(D\widehat F\)
has field-force amplitude
three by the preceding
site lemma. Its transpose
maps a strong field to
counting top amplitudes:
the physical pairing on
a cell costs
\(\xi^{-3}\xi^{-1}\xi^4=1\).
The two exponential
Schur sums make these
maps bounded between
their cell spaces.
It follows that
\begin{equation}
               \|\mathbf G_0\|_{\mathcal F_p\to\mathcal A_\alpha}
                                   \le C_{p,\alpha,\kappa}.
 \label{scf:strong-G0-family}
\end{equation}

For \(\lambda=\Lambda v\),
the new fourth-order
restoration gives a
family in
\(\mathcal Y_{2,\alpha}\).
Hence \(S=1-D\Lambda\)
has a family on
\(\mathcal A_\alpha\).
The exact grouped current
formula
\eqref{ncw:grouped-B}
can be written as
\[
 B_U[J]v=W[J]\lambda+
   \Lambda^\sharp\!
       \left(W[J]^\sharp v
          -\tfrac12W[J]^\sharp D\lambda\right)
       -\tfrac12\Lambda^\sharp D^\sharp W[J]\lambda .
\]
The first term is a
bounded field force of
grade three. The two
site forces in parentheses
have amplitude four.
The last term uses a
bounded grade-three
field force inside its
variational divergence.
Equation
\eqref{scf:strong-adjoint-families}
therefore proves
\begin{equation}
          \|\mathbf B_U[J]\|_{\mathcal A_\alpha\to\mathcal F_p}
                           \le C_{p,\alpha,\kappa}\gamma,
               \qquad \gamma=\|J\|_{\infty,3}.
 \label{scf:strong-current-smallness}
\end{equation}
Every product keeps the
actual current bond,
both matrix endpoints,
and all background and
geometry data read by
its factors. No
derivative of \(J\)
is present.

Choose \(\gamma\) below
a fixed smaller positive
ceiling, after \(p,\alpha\)
and the carrier constants
have been chosen.
Then
\[
          \mathbf G=\mathbf G_0
                   \sum_{n\ge0}
                        (\mathbf B_U[J]\mathbf G_0)^{\star n}
\]
converges normally as a
family from \(\mathcal F_p\)
to \(\mathcal A_\alpha\).
Its complete sum is the
first canonical inverse.
The difference \(G-G_0\)
has the additional
factor \(C\gamma\).

\subsubsection{The two row inversions and the multiplier root}

The normalized field-row
trace and its transpose
are those in
\eqref{ncw:normalized-field-row}.
The trace maps a strong
field into counting row
amplitudes, by its
complete average and
outward row completion.
The transpose of a
unit counting row has
fine-force density of
grade three, as in
\eqref{cci:row-force-cell-injection}.
The positive family
for \(\mathcal T_0^{-1}\)
is already available.
Since \(G-G_0\) has
strong family norm
\(C\gamma\), its
normalized row-Schur
difference has that
same small bound.
The exact Neumann
formula
\eqref{ncw:first-Schur-inverse}
therefore gives the
first row inverse
and the first lift
as strong families.

For the induced
\(\omega_*=H^\sharp J\),
retain the original
fine-current bond in
each paired coefficient.
The exact physical
duality calculation
\eqref{ncw:rooted-duality}
gives its absolute
rooted bound
\[
          \sup_c\sum_C
             e^{\kappa\sigma_c(C)}
                           \xi_c^3|\omega_{c,C}|\le C\gamma .
\]
The bare readout is
the separately retained
local term described
in NCW. Its bare
second constraint jet
is identically zero.
Thus the actual
\(\Delta_U^{(2)}[\omega_*]\)
has a family from
\(\mathcal A_\alpha\)
into bounded forces
of amplitude three.
This uses both absolute
leg sums of the
actual row coefficient,
its fixed-leg fine
density and the
physical transpose.
No individual invariant
Taylor piece is used.

Compose it with \(S\)
and \(S^\sharp\).
The second correction
then has family norm
\[
 \|\mathbf S^\sharp
       \mathbf{\Delta_U^{(2)}[\omega_*]}\mathbf S
                  \|_{\mathcal A_\alpha\to\mathcal F_p}
                                        \le C\gamma .
\]
A further fixed decrease
of the current ceiling
admits the second
Neumann series for
\(G_1\). Its row
matrix differs from
the same positive
\(\mathcal T_0\)
by \(C\gamma\).
The second row inverse
and \(H_1\) follow by
the same exact formula.
The readout in this
second correction has
throughout used the
first lift \(H\).

\begin{proposition}\label{prop:scf-strong-current-families}
Fix the stated
\(\alpha,p,\kappa\).
On a smaller common
\(C^{1,\alpha}\)
background tube and
for a sufficiently
small full independent
current radius, the
canonical second-variant
row saddle has an
exact local family
\begin{equation}
 \mathfrak I_{U,J}:
      \mathcal F_p\oplus\mathcal B
                \longrightarrow
          \mathcal A_\alpha\oplus\mathcal W .
 \label{scf:strong-row-saddle-family}
\end{equation}
Its row/column carrier
norm is bounded by a
fixed constant.
The same construction
gives strong families
for \(G_F,G_0,G,G_1\),
\(H,H_1,S,\Lambda\)
and the two adjoint
maps in
\eqref{scf:strong-adjoint-families},
in their stated spaces.
Each coefficient is
holomorphic in its
actual local background
and independent current
data, with its complete
determining geometry.
The constants and
radii are independent
of depth, compatible
volume and admitted
region shape.
\end{proposition}

\begin{proof}
The preceding constructions
give the displayed maps
in their prescribed order.
The full row saddle
is obtained from
the exact formulas
\[
 A=G_1f+H_1(d-QG_1f),\qquad
 \nu=(QG_1Q^\sharp)^{-1}(d-QG_1f),
\]
with the physical
normalizations of
\eqref{ncw:mass-removed-coarse-family}.
They give precisely
the datum amplitude
one and multiplier
amplitude three in
\(\mathcal B,\mathcal W\).
Products use the
smooth cell reconstruction
and the joined-carrier
algebra. Local
primitive matrices,
the new restoration
families and all
coarse inverses are
holomorphic on their
common smaller tube.
Normal convergence
preserves this property.
Only finitely many
strong primitive,
coarse inversion and
current Neumann stages
are used, so the
final tube and ceiling
are positive.

When a cover is
needed, use one
common sufficiently
large cover for
all stages, keep
full lifted carriers,
and deck-average
before projection.
The strong local
cell norms and
physical measures
descend with the
same row/column
bounds, as in NCW.
No quantum-volume
identity is involved.
\end{proof}

The current ceiling
in this strong
family construction
may depend on the
fixed \(p,\alpha\)
and carrier rate.
The larger energy
ceiling from the
completed inverse
theorems is unchanged.
The proposition supplies
the stated stronger
local representation
on its smaller domain;
it makes no endpoint
Riesz assertion and
no history-deletion
identification.

% END LAYERED EXPANSION strong_current_families.tex

% BEGIN LAYERED EXPANSION inverse/nonlinear_germ_connected_families.tex
\subsection{Connected families for the selected nonlinear germ}
\label{ch:ngc-nonlinear-connected-families}

The strong inverse families let us expand the selected
nonlinear equation itself. Each coefficient below retains
the background carrier of every inverse factor, the
actual row carrier of every nonlinear vertex, and all
datum inputs. Their absolute connected norm is summable
over every Taylor degree. This is stronger than a separate
bound for each fixed derivative.

We use the auxiliary action and current convention of
\cref{ch:cng-canonical-nonlinear-germ}. The resulting sum
is that selected germ, rather than a new equation with
generic local coefficients. Its physical identification
is inherited from the exact scalar identity proved there.

\subsubsection{Inputs, cell slots and carrier norms}

Fix \(0<\alpha<1\), a finite \(p>4/(1-\alpha)\),
and a finite requested carrier rate \(\kappa\).
Choose the primitive families at a stronger fixed rate
\(\kappa_+\). The final rate reserve will pay fixed
local stars, supported-slot collars and the base-variable
Taylor expansion. All choices precede depth and volume.

Let \(\mathfrak D\) be a smaller common domain of pairs
\(P=(U,J)\) supplied by
\cref{prop:scf-strong-current-families,prop:bhr-strong-restoration-family}.
The background is in a weighted \(C^{1,\alpha}\) tube
about a fixed given real current-regular \(U_*\).
The full independent-current amplitude is smaller than
the strong-family ceiling. That ceiling may depend on
\(\alpha,p,\kappa_+\); the larger energy ceiling is
not asserted to suffice for this representation.
Comparison frames, norms and the node extension are
fixed at \(U_*\). At each present base \(U\), its
restoration, constraint and fixed-for-this-germ slice
are the actual \(\Lambda_U,\mathfrak q_U\), and
\(R_UD_U^\sharp A=0\).

Use the spaces
\begin{equation}\label{ngc:joint-spaces}
 \mathcal X=\mathcal A_\alpha\oplus\mathcal W,
 \qquad \mathcal Z=\mathcal F_p\oplus\mathcal B,
 \qquad x=(A,\nu).
\end{equation}
Here the fine strong norm, row datum amplitude one
and multiplier amplitude three are those of
\cref{ch:cng-canonical-nonlinear-germ}.
The local norm is the maximum of its field and
row components at the assigned cell. There are only
a fixed number of row species at a top.

Local field slots are supported strong fields from
a fixed smooth cell partition with complete collars.
Their reconstruction has bounded overlap and is
bounded from the maximum slot norm into the global
strong norm. Use the corresponding bounded smooth
localization for force slots. Row data and currents,
which are not differentiated, can use bounded
ownership slots. Multiply block norms by the fixed
reconstruction bound so their carrier product norm
is submultiplicative. None of these slots counts
its individual fine coordinates.

A carrier \(X\) records full tags, short joins,
local determining geometry and complete coefficient
read sets. Its length is
\(\sigma_c(X)=b_L+d_c(X)\), with the guarded joining
rule used in the linear families. Every multilinear
coefficient has its output tag, all input tags and
its primitive coefficient carriers in \(X\).
For an \(m\)-linear family \(\mathbf V_m\), use
\begin{equation}\label{ngc:multilinear-family-norm}
 \begin{split}
 \|\mathbf V_m\|_\kappa=\max\bigg\{&
 \sup_y\sum_{z_1,\ldots,z_m,X}
              e^{\kappa\sigma_c(X)}\|V_{m;y,\mathbf z}[X]\|,\\
 &\max_i\sup_{z_i}\sum_{y,\mathbf z_{\ne i},X}
              e^{\kappa\sigma_c(X)}\|V_{m;y,\mathbf z}[X]\|
                                                    \bigg\}.
 \end{split}
\end{equation}
The operator norm inside the sum uses the displayed
local Banach slots and is the supremum on a fixed
larger common local base domain \(\mathfrak D_X^+\).
Thus the supremum is taken separately for each
coefficient before the carrier sum. The final
base domain has a strict margin inside these
larger domains. It is a norm of divided Taylor
coefficients \(D^m/m!\). Ordered products and joins
of such families obey the usual substitution bound:
sum each internal cell index using a row sum, or
the distinguished column sum along its path to a
selected input. The remaining factors use row sums.
The resulting norm is at most the product of the
factor norms. Fixed local stars cost \(C_\kappa^m\)
at an \(m\)-linear vertex.

\begin{lemma}[Uniform local-data majorants]
\label{lem:ngc-local-data-majorants}
The required strong linear families can be chosen
with bounded coefficientwise local-data supremum
norms as just defined. Their constructive bounds
are positive majorants uniform on the larger
common base domain. The same convention applies
to every tree and nonlinear coefficient below.
\end{lemma}
\begin{proof}
The local Dirichlet columns and cutoff errors
have bounds \(Ce^{-sd(y,z)}\) and
\((C/m)e^{-sd(y,z)}\), respectively, uniformly
on their entire larger local coefficient tube.
Taking each coefficient's local-data supremum
therefore preserves those bounds. The weighted
centre count and the positive word sums for
the primitive walks give the sum of these
suprema, rather than only a bound at one
common global base point.

There is one extra count in a positive coarse
inversion. A combined near power \(H_<^m\)
has \(2\)-operator norm at most \((q')^m\)
uniformly on the larger complex tube. Its
entrywise local-data supremum is consequently
at most \(C(q')^m\). The declared endpoint
balls have radius \(B_m\), so the number of
possible endpoints in a row or column is at
most \(C(1+B_m)^4\). Hence its stronger
supremum family bound is
\[
                         C(1+B_m)^4(q')^m.
\]
The smaller \((1+B_m)^2(q')^m\) Hilbert-row
estimate at one base point would not by
itself justify moving the supremum through
the endpoint sum. The larger polynomial
above suffices: \(B_m\) grows polynomially
in \(m\), and the contraction is exponential.
Use it in every positive coarse short-power
stage, choose the long cutoff \(R=Dm\) as
before, and increase \(m\) until both this
term and the positive long-word bound are
small in the carrier norm. Reblocking still
uses a bounded number of final tags for the
combined entry. This only enlarges fixed
carrier and separation choices.

The remaining coarse words, inverse Neumann
series, rooted multiplier pairings and
strong-current perturbations are bounded
by products of these positive coefficient
majorants. Choose the strong current radius
after their constants. This proves the
claim through the finitely many primitive,
Gram, gauge-Schur and row-Schur stages.
No exchange of a pointwise base supremum
with an infinite carrier sum is presumed.
All later analytic local vertex bounds
hold uniformly on a larger local strong
ball, so their Cauchy majorants and ordered
tree bounds have the same property.
\end{proof}

The strong linear inputs are the exact families
\begin{equation}\label{ngc:linear-inputs}
 \mathfrak I_P:\mathcal Z\longrightarrow\mathcal X,
 \quad \Lambda:\mathcal A_\alpha\longrightarrow Y^{2,\alpha},
 \quad S:\mathcal A_\alpha\longrightarrow\mathcal A_\alpha,
\end{equation}
and the adjoint families
\begin{equation}\label{ngc:adjoint-inputs}
 \Lambda^\sharp:\{h\}_{\infty,4}\longrightarrow\{f\}_{\infty,3},
 \quad S^\sharp:\mathcal F_p\longrightarrow\mathcal F_p,
 \quad \Lambda^\sharp D^\sharp:
                 \{f\}_{\infty,3}\longrightarrow\mathcal F_p.
\end{equation}
They are precisely
\eqref{scf:strong-row-saddle-family},
\eqref{scf:strong-adjoint-families} and
\eqref{scf:bounded-divergence-adjoint}, with the new
strong site restoration. The first-lift readout
\(\omega=H^\sharp J\) has its rooted local family
with multiplier amplitude \(C\|J\|_{\infty,3}\).
The original fine-current input stays in each such
carrier. It is not replaced by a second-lift readout.

\subsubsection{The actual nonlinear vertices}

For this fixed \(P\) put
\[
 \Pi(A)=\Pi_U(A),\quad \Xi(A)=\Pi(A)-SA,\quad
 N_q(a)=\mathfrak q_U(a)-Qa,\quad r(A)=\mathfrak q_U(\Pi(A)).
\]
These are the finite gauge chain and actual
re-centered straightened constraint already defined.
The local formula for \(\Pi\) is analytic in
\(A,\Lambda A,D\Lambda A\). Its divided fine-mesh
factor is removed before taking norms.
\Cref{prop:bhr-strong-restoration-family} controls
the site argument through two H\"older derivatives,
coefficient by coefficient.
Use \(S=1-D\Lambda\) from this same restoration
family when expanding the finite chain. Its local
linear jet \(A-D\Lambda A\) is then removed before
filing the quadratic remainder \(\Xi\).

\begin{lemma}[Local analytic vertices without fine counts]
\label{lem:ngc-local-vertices}
The local Wilson, constraint, finite-gauge and
adjoint remainders used below have exact multilinear
carrier families with coefficients bounded by
\(C_0C_1^m\) at degree \(m\). Their coefficients are
holomorphic in their actual local \((U,J)\) data.
The constants are uniform in depth and fine mesh.
\end{lemma}
\begin{proof}
The centered Wilson remainder
\(F_U'(A)-j_U-F_U''(0)A\) is a local map from
the strong ball into bounded grade-three forces.
Its uniform holomorphic bound is
\cref{lem:cng-Wilson-residual}. Its constant and
linear jets are subtracted before expansion.
This is essential when the actual \(j_U\) at a
complex base is large. Each plaquette expression
reads a fixed fine stencil, hence a fixed number
of complete strong cell slots. Banach-valued
Cauchy and polarization give \(C^m\), independently
of the dimensions of those slots.

The actual mean and its adjoint satisfy the
holomorphic bounds in
\cref{lem:cng-constraint-residuals}. An assigned
row and its complete carrier read at most a fixed
number of adjacent-scale cells. Proper separation
excludes a second intervening level. The fixed-leg
fine density in \eqref{pnc:actual-two-leg-bound}
cancels the inverse fine volume in the physical
transpose. Thus the bound for
\(DN_q(a)^\sharp\nu\) is a bounded point-force
bound before any adjoint restoration is applied.
Cauchy on its uniform larger strong ball gives
all higher Taylor degrees in the same point-force
class. It does not count fine coordinates or
postulate a new higher-leg density estimate.

The finite gauge function uses local point
products of the scaled variables
\(\xi A,\Lambda A,\xi D\Lambda A\).
Its local power series has \(C^m\) coefficients.
Substitute the strong linear families in
\eqref{ngc:linear-inputs}; each occurrence brings
its entire data carrier. A point derivative
paired with a bounded fine force has exactly
the factory
\[
                f_0+\Lambda^\sharp h_\lambda
                             +\Lambda^\sharp D^\sharp f_d.
\]
Its local fields \(f_0,f_d\) have grade three
and \(h_\lambda\) has site grade four. Endpoint
injection uses only the fixed graph degree.
The families in \eqref{ngc:adjoint-inputs}
give the required finite-\(p\) output bounds.
No derivative is taken of the paired force or
of \(J\).

These local analytic maps involve a fixed number
of strong cell slots. Substitution of an
\(m\)-linear coefficient attaches at most a
fixed multiple of \(m+1\) linear factors or
fixed local stars. Their family norms and the
guarded joins change \(C^m\) by another fixed
exponential in \(m\). The rooted family of
\(\omega(P)\), where it occurs, attaches its
full base-current carrier and has bounded norm.
This proves the assertion in
\eqref{ngc:multilinear-family-norm}.
\end{proof}

We next make explicit which force classes enter
these adjoints. The identity \(QS=Q\) gives
\begin{equation}\label{ngc:quadratic-constraint-value}
                       r(A)-QA=Q\Xi(A)+N_q(\Pi(A)).
\end{equation}
Use this equality before filing coefficients.
It avoids a linear cancellation between different
restoration carriers. Similarly,
\begin{equation}\label{ngc:multiplier-force-rewrite}
 \begin{split}
 (r'(A)-Q)^\sharp\nu={}&\Xi'(A)^\sharp Q^\sharp\nu
                +S^\sharp DN_q(\Pi(A))^\sharp\nu\\
             &+\Xi'(A)^\sharp DN_q(\Pi(A))^\sharp\nu.
 \end{split}
\end{equation}
Every input to \(S^\sharp\) or the nonlinear
adjoint factory on the right is a bounded
grade-three force, supplied by the actual row
columns. A generic \(\mathcal F_p\) force is
not passed to a factory that requires a point
force.

For the constraint-current term put
\(C_\omega=\Delta_U^{(2)}[\omega]\), the actual
second field-mean jet. Its residual has the
exact expression
\begin{equation}\label{ngc:constraint-current-rewrite}
 \begin{split}
 &\nabla\langle\omega,N_q(\Pi(A))\rangle
                         -S^\sharp C_\omega SA\\
 &\quad=S^\sharp\{C_\omega\Xi(A)
           +DN_q(\Pi(A))^\sharp\omega-C_\omega\Pi(A)\}
       +\Xi'(A)^\sharp DN_q(\Pi(A))^\sharp\omega.
 \end{split}
\end{equation}
The braces are bounded point forces and are
quadratic in \(A\). In the row mass use
\begin{equation}\label{ngc:row-mass-rewrite}
 r'(A)^\sharp a_fr(A)-Q^\sharp a_fQA
  =(r'(A)-Q)^\sharp a_fr(A)
                         +Q^\sharp a_f(r(A)-QA).
\end{equation}
The weight \(a_f\asymp\xi^{-2}\) converts datum
amplitude one into multiplier amplitude three.
The current term is the exact finite-chain
adjoint remainder after its quadratic scalar
jet \(\langle A,B_U[J]A\rangle/2\) is removed.
Its factory inputs are again bounded forces.
These formulas show both the quadratic zeros
and the valid source classes coefficient by
coefficient. There is no \(L^\infty\) Riesz
claim.

\subsubsection{An exact tree solution of the selected equation}

Let \(N_{\mathcal L}(A)=\mathcal L_P'(A)-G_1^{-1}A\)
be the subtracted auxiliary gradient of
\eqref{cng:full-gradient-residual}.
Define
\begin{equation}\label{ngc:preconditioned-vertices}
 L_Pb=\mathfrak I_P(0,b),\qquad
 \mathfrak N_P(A,\nu)=\mathfrak I_P
 \binom{-N_{\mathcal L}(A)+(r'(A)-Q)^\sharp\nu}
                  {-(r(A)-QA)}.
\end{equation}
The actual preconditioned equation is
\begin{equation}\label{ngc:exact-fixed-point}
                         x=L_Pb+\mathfrak N_P(x).
\end{equation}
It is exactly \eqref{cng:preconditioned-fixed-point}.
In particular its multiplier and curvature terms
have not been replaced by model vertices.

\begin{lemma}[The nonlinear vertex family]
\label{lem:ngc-preconditioned-family}
On a fixed smaller joint strong ball,
\begin{equation}\label{ngc:vertex-series}
         \mathfrak N_P(x)=\sum_{m\ge2}V_{m,P}[x^m],
               \qquad \|\mathbf V_{m,P}\|_{\kappa_+}\le C_0C_1^m.
\end{equation}
The series and carrier sums converge absolutely
on a uniform ball. The linear family \(L_P\)
has bounded row and column carrier norm.
\end{lemma}
\begin{proof}
Use the explicit remainder formulas above and
\cref{lem:ngc-local-vertices}. The source tuple
is in \(\mathcal Z\): its first component is
in finite \(\mathcal F_p\), and its second
has datum amplitude one. Apply the exact strong
row-saddle family \(\mathfrak I_P\).
All constant and linear jets were removed before
filing. For example \eqref{ngc:quadratic-constraint-value}
has two manifestly quadratic terms; the second
line of \eqref{ngc:multiplier-force-rewrite} has
at least one field factor and one multiplier
factor. The removed finite-chain scalar
quadratic jet has exactly the canonical
\(B_U[J]\) coefficient. Consequently no
degree-zero or degree-one representative is
left for cancellation between carriers.

Linear primitive family products and the local
\(C^m\) coefficient bounds give the displayed
bound, after enlarging \(C_1\). Full base
coefficient carriers are retained in every
product. The linear assertion is the same
row-saddle theorem with zero force source.
\end{proof}

\begin{proposition}[All-degree connected datum families]
\label{prop:ngc-datum-connected-family}
For \(P\in\mathfrak D\) and sufficiently small
\(\|b\|_{\mathcal B}<r_b\), the selected auxiliary
solution has an exact expansion
\begin{equation}\label{ngc:datum-family}
 x(P,b)=\sum_{n\ge1}\sum_{\mathbf z,X}
                    x_{n;\mathbf z}[X](P_X)[b_{z_1},\ldots,b_{z_n}],
 \qquad
 \sum_{n\ge1}r^n\|\mathbf x_n(P)\|_{\kappa_+}\le Cr,
                      \quad 0<r\le r_b.
\end{equation}
Every \(X\) is connected and contains its output,
all datum inputs and all background/current read
sets. Each coefficient is holomorphic in those
local base data. Constants and radii are uniform
in depth, compatible volume and admitted shape.
\end{proposition}
\begin{proof}
Expand \eqref{ngc:exact-fixed-point} in ordered
rooted trees. A leaf carries the exact linear
family \(L_Pb\); an internal vertex of arity
\(m\ge2\) carries \(V_{m,P}\). Join every edge
carrier at its shared cell, including the
fixed short join when physical overlap does
not mean an identical tagged vertex. The
result contains every primitive base read set,
not merely the geometric hull of the datum
leaves. An \(n\)-leaf tree has at most
\(n-1\) internal vertices. All fixed stars and
linear factors can therefore be paid by a
fixed exponential in its degree.

To prove convergence, replace every coefficient
by its family norm. The positive scalar majorant is
\begin{equation}\label{ngc:tree-majorant}
              t=C_{\rm lin}r+\sum_{m\ge2}C_0C_1^m t^m.
\end{equation}
Choose \(r_b\) so this map preserves
\([0,2C_{\rm lin}r_b]\) and its derivative there is
less than \(1/2\). Iteration from zero gives
a convergent positive series with
\(t(r)\le2C_{\rm lin}r\). Its coefficients are the
sums of all ordered tree weights. The
substitution estimate for
\eqref{ngc:multilinear-family-norm} bounds the
actual coefficient families by this series.
There is no extra factorial: the vertices
were divided Taylor coefficients, and the
ordered tree sum is precisely the expansion
of the powers in \eqref{ngc:tree-majorant}.

Normal convergence permits substitution in
the actual equation. Its sum lies in the
strong uniqueness ball of
\cref{prop:cng-auxiliary-germ}, so it equals
that already constructed selected solution.
This establishes the claimed identity, not
only a formal series solution. The local
holomorphy of every primitive and normal
convergence give the final assertion.
\end{proof}

For functions rather than multilinear kernels,
group all terms with the same carrier. If
\(x_y(P,b)=\sum_X x_{y,X}(P_X,b_X)\), then
\begin{equation}\label{ngc:function-family-bound}
 \sup_y\sum_X e^{\kappa_+\sigma_c(X)}
      \sup_{\|b_X\|<r}\|x_{y,X}(P_X,b_X)\|_{\mathcal X_y}
                                                        \le Cr.
\end{equation}
This is an absolute connected family bound
for the whole analytic function. It does
not follow merely by summing separate
fixed-order decay estimates.

\subsubsection{The literal datum and the field/current outputs}

The explicit inverse outward chart \(b_U(B)\)
in \eqref{pnc:actual-reduced-datum} is local
and analytic with uniform \(C_L^m\) Taylor
bounds. Its outward factor reads only finer
interior rows, which are unchanged by that
chart. There is no recursive product over
already corrected outward data. A one-step
complete carrier therefore supplies its
connected family, including its actual
linear derivative \(Db_U(0)\). Substitute
that family in \eqref{ngc:datum-family}, and
reduce the literal radius \(r_B\) so its
absolute image lies in the reduced datum ball.
The same tree majorant then gives
\eqref{ngc:function-family-bound} with \(B\)
and a changed fixed constant.

Put \(C=\Pi_U(A)\) and
\(\delta\omega=\nu-a_fb_U(B)\).
The point-force identity for the new current
can be written without cancellation of its
degree-zero terms:
\begin{equation}\label{ngc:current-increment}
 \begin{split}
 J_{\rm new}-J={}&(\mathcal E_{\eta C}^{-\sharp}-1)J\\
 &+\mathcal E_{\eta C}^{-\sharp}
       \{DN_q(C)^\sharp\omega
                       +D\mathfrak q_U(C)^\sharp\delta\omega\}.
 \end{split}
\end{equation}
Here \(\omega=H^\sharp J\) still uses the first
lift. Every term on the right has positive
datum degree. Its row transpose inputs have
the actual fine fixed-leg density, so its
output is a point force of amplitude three.
No adjoint-restoration map is needed after
this point-force expression. The divided
exponential differential is a bounded local
analytic point function.

\begin{proposition}[Connected families for the literal update]
\label{prop:ngc-literal-connected-family}
The following outputs have exact holomorphic
connected families satisfying a bound \(Cr\)
of the form \eqref{ngc:function-family-bound}
on \(\|B\|_{\mathcal B}<r\le r_B\):
\[
 A,\quad C,\quad \nu,\quad\delta\omega,
       \quad J_{\rm new}-J,\quad
                \widehat A,\quad\widehat J-J.
\]
Field outputs have the strong norm, multiplier
outputs amplitude three, and current outputs
point-force amplitude three. The literal
potential is defined by
\(e^{i\eta\widehat A}U=\widehat U\), with
the fixed node extension \(E_n\). Every
coefficient retains all actual datum and
base read sets.
\end{proposition}
\begin{proof}
For \(C=SA+\Xi(A)\), substitute the connected
solution family in the local finite-gauge
series and its strong linear factors. Its
degree-one term is \(SH_1b\); higher terms
are controlled by the same analytic majorant.
The multiplier increment is an exact linear
combination of known families. Equation
\eqref{ngc:current-increment} and the point-force
row bounds prove the current assertion by
local analytic substitution. Multiplication
of two such normalized families has its
usual joined-carrier bound.

For literalization let
\(\zeta=E_n(i^{-1}\log v_U[C]^{-1})\),
and \(g=e^{i\zeta}\). The actual assigned
node factor is a local analytic function
of the fine strong field on its complete
physical star. At an assigned scale \(\xi_y\),
its relative one-step factors at intermediate
scale \(\xi_t\) have variations bounded by
\(C_Lr\xi_t/\xi_y\) on a uniform complex
strong ball. The actual path-row bounds for
the group mean give this estimate before
any fine coordinates are counted. Since
\(\sum_{t\le j_y}\xi_t/\xi_y\le L/(L-1)\),
their product is uniformly analytic there.
The union of its nested read sets lies in
a star of radius \(C_L\xi_y\), by the same
geometric sum. Cauchy gives its \(C^m\)
strong-cell coefficients. These are the
actual relative node products, including
the outward one-step factor of
\eqref{pnc:outward-node-factor}. The fixed-reference
\(E_n\) has local \(Y^{2,\alpha}\) seeds.
Thus \(\zeta\) has a connected strong site
family with norm \(Cr\).
The exact relative potential is
\[
 \widehat A_b=(i\eta)^{-1}
       \log\{g_-e^{i\eta C_b}U_bg_+^{-1}U_b^{-1}\}.
\]
It is an analytic local function of
\(C,\zeta,D\zeta\); its apparent fine-mesh
division cancels before expansion, as in
\eqref{cng:Pi-strong-bounds}.
The same local substitution proves the
strong field assertion. Finally use
\[
 \widehat J-J
   =\operatorname{Ad}(g_-)(J_{\rm new}-J)
                          +(\operatorname{Ad}(g_-)-1)J.
\]
This is a local point product of families
already controlled and a bounded base
current. Its constant datum jet is zero
before filing. The bare node values are
one and no boundary gauge term is added.
\end{proof}

\subsubsection{Joint base and datum Taylor coefficients}

The preceding coefficients are holomorphic in
their local base data. We now record an
absolute all-degree statement in those base
variables too. The variables are supported
Banach slots, not independent raw fine-bond
discs.

Fix an interior pair \(P_0=(e^{i\eta E_0}U_*,J_0)\)
of the smaller common domain. Let \(e_y\) be
supported strong field slots and \(j_y\) bounded
current ownership slots. Write
\begin{equation}\label{ngc:base-reconstruction}
 E=E_0+\sum_y\mathcal R_y e_y,\qquad
 J=J_0+\sum_y\mathcal O_y j_y.
\end{equation}
The smooth reconstruction \(\mathcal R_y\)
has complete collars and bounded overlap.
The ownership \(\mathcal O_y\) is bounded
in the point-current norm and need not be
smooth, since no current is differentiated.
Normalize the slot radii by fixed positive
fractions of their admitted tube and current
margins.

Every tag and its read collar meet only
their own and adjacent assigned levels when
\(\rho\gg LM_c\). A tag then meets at most
\(C_LM_c^4\) assigned strong cells. In particular,
after one fixed collar enlargement, the number
\(n_X\) of base slots read by a full carrier
satisfies
\begin{equation}\label{ngc:base-slot-count}
                       n_X\le C_{L,M_c}\sigma_c(X).
\end{equation}
Each slot may contain arbitrarily many fine
coordinates. This count does not count them.
The enlargement and its canonical input tags
are charged in the initial rate choice.
The tag graph has bounded degree \(D_L\).
At most \(C_{L,M_c}\) assigned output cells belong
to one tag. Thus a radius-\(j\) tag ball has at
most \(C_{L,M_c}D_L^{j+1}\) possible output cells.
Reserve also a rate \(a>\log D_L+1\) for columns
anchored at a base input, in addition to the
supported-collar and Taylor rate costs.

\begin{proposition}[Joint all-degree connected expansion]
\label{prop:ngc-joint-connected-family}
On fixed sufficiently small reconstructed base
slot and literal datum balls, all outputs in
\cref{prop:ngc-literal-connected-family} have
an exact joint Taylor expansion with absolutely
summable connected coefficients at the requested
final rate \(\kappa\). If \(n\) is the datum
degree and \(\ell\) the combined base-slot
degree, then
\begin{equation}\label{ngc:joint-degree-bound}
       \sum_{n\ge1,\ell\ge0}r^n\delta^\ell
                   \|\mathbf F_{n,\ell}\|_\kappa\le Cr.
\end{equation}
Every coefficient contains its output, all datum
input cells, all selected base input cells and
its full determining-geometry carrier. All
coefficients in this display are divided Taylor
coefficients in their indicated Banach slots.
\end{proposition}
\begin{proof}
Begin with the datum family at a reserved rate
larger than the final rate, including the
one-time supported-slot collar enlargement.
For a fixed carrier its coefficient function
is holomorphic on a product of \(n_X\)
normalized Banach slot balls of a common
positive radius. Bounded-overlap reconstruction
maps that product into the admitted global
strong/current tube. Locality means slots
outside its complete read set do not enter.

Multivariate Cauchy and dimension-free
polarization bound the coefficient of total
base degree \(\ell\), for each distribution
of degrees among slots, by \(C^\ell\)
times its larger-ball bound. The polarization
factor at degree \(j\) is at most
\(j^j/j!\le e^j\), so it is included in
\(C^\ell\). Summing those distributions on
a radius-\(\delta_1\) ball costs
\[
                  (1-C\delta_1)^{-n_X}
       \le\exp(C'\delta_1 C_{L,M_c}\sigma_c(X)),
                    \qquad C\delta_1<1/2.
\]
Choose \(\delta_1\) so this exponent fits the
reserved Taylor rate, while leaving the additional
rate \(a\) unused. Take the actual base radius
\(\delta\le\delta_1/2\). The coefficient of
degree \(\ell\) then gains
\((\delta/\delta_1)^\ell\), so summing the
degreewise supremum norms is justified by a
geometric series. The original datum row and
datum-column norms retain their bounds.

For a column anchored at a base slot \(w\),
every contributing carrier contains its tag.
The unused rate gives
\(e^{-a\operatorname{dist}_{\rm tag}(y,w)}\)
for its possible output cell \(y\). The
tag-ball count above makes the sum over those
outputs finite, uniformly in volume and depth.
Use the degreewise datum row bound for the
remaining carrier and input sums. The same
geometric factor sums the base degrees.
This proves the base-column bound as well,
and hence \eqref{ngc:joint-degree-bound} in
the row/column norm, not only its row half.
All selected slots meet the declared carrier
collar, which already contains their input
tags and short joins. Thus the resulting
joint Taylor carrier is connected. There
is no fine-coordinate Cauchy count.

Every output increment under consideration
vanishes at \(B=0\) for every admitted base
pair. Its datum expansion was filed with
that zero removed before the base expansion.
Consequently every resulting joint coefficient
has \(n\ge1\), as asserted. Normal convergence
identifies the joint series with the same
selected germ and literal update.
\end{proof}

\subsubsection{Physical identification and the scope retained}

At a physical pair \(J=j_U=Q^\sharp\omega_{\rm phys}\),
the sums are the physical PNC solution and
its actual left current after the specified
literal gauge. This follows from uniqueness
and \cref{prop:cng-physical-matching,prop:cng-current-update};
the tree construction does not impose a new
physical equation. The first reduced-datum
derivatives remain
\[
              D_bA(0)=H_1,\qquad D_bC(0)=SH_1.
\]
Literal data add the actual right factor
\(Db_U(0)\), and the literal node gauge adds
its fine gauge vector to the field derivative.
Off current these jets are not equated.

For small periods construct every primitive,
local vertex and tree on one common cover
before projecting its carrier. Retain each
lifted input and every winding join. The
absolute tree norm is bounded by the same
cover sums. Projection solves the actual
finite periodic fixed-point equation, and
strong contraction uniqueness identifies
its sum. Deck averaging preserves the
supported strong bounds and the original
physical units. No assertion about a
quantum partition function on different
volumes is used.

The result supplies all-degree connected
families for the selected off-current germ
on its fixed strong reference chart.
It does not identify a different prescribed
core convention, manufacture a reference
from arbitrary curvature-only coarse data,
or prove history/tower admission or deletion
laws. Its local background domains are
the stated strong Banach domains. The
auxiliary cube-Gaussian discs of a separate
formation construction are not being
declared independent fine-background discs.

% END LAYERED EXPANSION inverse/nonlinear_germ_connected_families.tex

\subsection{The decoupled response on the regular-pair domain}
\label{dnr:section}

The equations in \eqref{eq:rg-objects-R9} use a field and its first
covariant differences.  Their background domain is consequently the
regular-pair domain of \cref{notdef:propagators}.  We construct the
response on that domain.  The H\"older estimates below are estimates
on inverse outputs; no H\"older seminorm of the complex background
perturbation is assumed.  The construction follows the constraint
straightening and replacement of propagators in Ba{\l}aban's background
and cluster arguments \cite{B85e,B88a}.  We give the first-order estimates,
the independent-current completion, and the complete carrier convention
needed for the present averages.

Fix the blocking factor and the separation constants of the positive
local inverse construction.  All its complete gauge cubes, row reads,
outward completions, and local compression collars are retained.
A pair is written \(\mathbf U=e^{i\eta E}U\), \(\mathbf J\), where \(U\)
is real, its local gauge potentials and first differences have the
regular-pair bounds, and
\[
 \sup_y\{\xi_y|E|_y+\xi_y^2|\nabla_U E|_y\}<a_*,\qquad
 \sup_y\xi_y^3|\mathbf J|_y<\gamma_* .
\]
The curvature bounds in that domain are imposed as well.  The positive
constants \(a_*,\gamma_*\) are fixed after the finitely many proof
carrier rates and minimal tile floors.  After the auxiliary inverse construction and its proof rates have
been fixed, the local column estimates are uniform in every larger
admissible final tile size and separation. Enlarging those final
parameters does not decrease these numerical radii: the already
constructed families remain available, and reblocking into larger
aligned tiles does not increase their carrier norms. At each such
choice the regular-pair hypothesis is imposed on that choice's
complete gauge cubes and collars. A pair admitted for smaller cubes
is not thereby admitted for larger cubes. Obtaining the required
local potentials from curvature alone can require a smaller curvature
ceiling after the cube size has been fixed. The radii above are
independent of depth, fine mesh, compatible volume, and admitted
region shape on the resulting regular-pair domain. The current is
an independent variable.

Use the physical fine pairing and counting measure on raw rows.  A raw
row is a logarithmic angle, without division by its length.  If \(Q\)
denotes the per-length row of the propagator construction, set
\(\widehat Q=\mathsf m Q\).  In this subsection the letters
\(H,H_1,\Sigma\) denote the raw-coordinate maps
\[
 H=\mathsf H\mathsf m^{-1},\qquad
 H_1=\mathsf H_1\mathsf m^{-1},\qquad
 \Sigma=(\widehat QG\widehat Q^\sharp)^{-1}.
\]
Thus \(\widehat Q^\sharp\mathsf a\widehat Q=Q^\sharp a_fQ\)
defines the raw row mass \(\mathsf a\).  A sharp denotes the transpose
in these pairings, continued complex linearly.  We reserve \(S\)
for the fine projection \(1-D\Lambda\).
Every field derivative, row, site inverse and current completion in
this subsection is evaluated at the complex configuration
\(\mathbf U\).  Thus an unsubscripted \(D\) means
\(D_{\mathbf U}\), and \(\mathrm d\) means its covariant curl.
The real factor \(U\) is used only to describe the input chart.

For a positive weight satisfying
\(w(y)/w(z)\le e^{\theta d(y,z)}\), put
\begin{equation}\label{dnr:norms}
 \begin{split}
 \|A\|_{1,w}&=\sup_yw(y)
       \max\{\xi_y|A|_y,\xi_y^2|\nabla_{\mathbf U}A|_y\},\\
 |f|_{3,w}&=\sup_yw(y)\xi_y^3|f|_y,\qquad
 |b|_w=\sup_cw(y_c)|b(c)|.
 \end{split}
\end{equation}
All differences and suprema use complete assigned stencils.
The fixed neighborhood radius \(r_0\) below includes these stencils
and the complete nonlinear row read.  Local scale ratios on this
neighborhood are bounded by constants depending only on the blocking
factor.  The positive inverse decay leaves a fixed admissible
\(\theta>0\) after these ratios have been absorbed.

\subsubsection{Strong outputs without a stronger background domain}

\begin{lemma}\label{dnr:divergence-holder}
Fix \(p=8\) and \(\alpha=1/4\).  The positive field inverse \(G_F\)
has complete local carrier families from cell \(L^p\) forces to
\(\mathcal A_\alpha\), and from covariant divergences of cell
\(C^\alpha\) tensors to \(\mathcal A_\alpha\).  The assertions hold
on the complex \(C^1\) tube just specified.  The second assertion
uses the infimum norm on representations
\(f=D^\sharp F+g\), with \(F\) of grade two in \(C^\alpha\)
and \(g\) of grade three in cell \(L^p\).
\end{lemma}

\begin{proof}
The positive field primitive and its local compressions are those of
\cref{scf:positive-field-columns}.  Its proof uses the background
potential and its first difference only.  Indeed, after pairing the
oppositely oriented covariant shifts, its equation in a local gauge
is the ordinary component lattice Laplacian plus first-order and
zeroth-order terms with bounded coefficients.  The Wilson correction
is a bounded curvature term with fixed neighboring shifts.  The
complete row mass is bounded point forcing from the local energy
norm.  None of these statements differentiates the curvature.
The finite-\(p\) bootstrap in
\eqref{psr:scalar-point-bootstrap} therefore proves the first assertion
on this \(C^1\) tube, coefficient by coefficient.

Here is the additional divergence estimate.  In unit cell coordinates
let \(P_t\) be the flat lattice heat operator.  The kernel of
\(\nabla(L_h+1)^{-1}\nabla^\sharp\) is
\(\int_0^\infty e^{-t}\nabla\nabla^\sharp P_t\,dt\).
The sum of each differentiated heat kernel is zero.  In its action
on \(F\), subtract \(F(x)\) before taking absolute values.  The
heat derivative bounds already proved in the discrete estimates give
\[
 \sum_z|\nabla^2P_t(x,z)|\,|F(z)-F(x)|
       \le C t^{-1+\alpha/2}[F]_\alpha
       \quad(t\ge h^2).
\]
For \(t<h^2\), the same estimate with \(t\) replaced by \(h^2\)
and the length of that interval gives a bounded contribution.
For the difference between outputs at points a distance \(r\) apart,
split the integral at \(r^2\).  The near part is at most
\(C r^\alpha[F]_\alpha\); the far part uses one further heat
derivative and is bounded by
\[
 C r[F]_\alpha\int_{r^2}^1t^{(\alpha-3)/2}\,dt
       +Cr\|F\|_\infty\le C r^\alpha\|F\|_{C^\alpha}.
\]
The mass controls the integral beyond one.  These estimates are
uniform down to \(r=h\).

Apply the finite-\(p\) divergence bootstrap
\eqref{psr:divergence-bootstrap} first.  It bounds \(A,\nabla A\)
in \(L^8\), and \(A\) in the supremum norm on a smaller box.
Every lower-order term is now an \(L^8\) force.  Its flat inverse
has gradient in \(C^\beta\) for every \(\beta<1/2\).
The preceding heat calculation handles the prescribed divergence.
The difference between a covariant and an ordinary divergence is a
bounded potential times \(F\).  Cutoff commutators involve only
\(F,A,\nabla A\) on the next reserved collar.  This proves the
interior \(C^{1,\alpha}\) estimate without a derivative of \(F\)
or a H\"older bound on a derivative of the background.

For completeness the estimate also gives a family, not only a bound
on its complete sum.  Use the same local parametrix \(M\) and
commutator \(T=K_FM-1\) as in the positive field construction.
For an initial divergence source, the just proved estimate bounds
\(Mf\) strongly.  The slow cutoff commutator maps that output to
a cell \(L^8\) force with norm \(C/m\).  Hence
\(G_Ff=Mf-G_FTf\), where the second \(G_F\) is its already
convergent \(L^8\)-force family.  This is an exact normally convergent
carrier formula.  All collars used in the estimate are recorded in
its carrier.  At a true bare cell the spacing equals the cell scale
and the same estimates are finite atomic estimates.  The cover and
deck averaging prescription of \cref{prop:ncw-named-families}
preserves these bounds.
\end{proof}

\begin{proposition}\label{dnr:linear-families}
On the preceding regular-pair domain, the raw maps
\(H,H_1,\Sigma\), and
\[
 \mathcal G=G_1-G_1DRD^\sharp G_1-H_1\widehat QG_1
\]
have exact complete local carrier families with exponentially
summable norms from raw row amplitudes, or bounded grade-three
forces, to the first-order norm in \eqref{dnr:norms}.
The same assertion holds for \(H_0\),
\(\widetilde G=G_0-H_0\widehat QG_0\), and
\(P_1=D(1-R)D^\sharp\), the last as a map from first-order fields
to grade-three forces.  The physical transpose \(H^\sharp\)
maps bounded grade-three forces to raw row amplitudes with the
corresponding column bound.  Every coefficient is holomorphic in
the background and independent current on its complete local domain.
\end{proposition}

\begin{proof}
We specify the changes needed in the strong-family argument;
the positive coarse inversions and complete current coefficients
remain those of \cref{prop:ncw-named-families}.
The site inverse, its normalized columns, and the two adjoint
families in \eqref{scf:strong-adjoint-families} use only bounded
background potentials and first differences.  Their proofs through
\eqref{scf:strong-G0-family} therefore apply on the present domain.
The positive field factor is supplied by
\cref{dnr:divergence-holder}.  Woodbury's formula gives
\(G_0:\mathcal F_8\to\mathcal A_\alpha\).

For \(v\in\mathcal A_\alpha\), set
\(\lambda=\Lambda v=G_sRD^\sharp v\).
The source \(D^\sharp v\) is bounded of grade two.
The projection \(R=1-\widehat F\widehat M^{-1}\widehat F^\sharp\)
preserves this bounded source class, since its column factor is
bounded and its two absolute cell sums converge.  The bounded-source
site bootstrap gives bounded \(\lambda,D\lambda\), of grades zero
and one.  This is the only regularity of \(\Lambda v\) needed here.
In particular \(Sv=v-D\lambda\) is bounded of grade one.

The exact grouped first current correction is
\begin{equation}\label{dnr:current-correction}
 B_{\mathbf U}[J]v=W[J]\lambda+
  \Lambda^\sharp\{W[J]^\sharp v-\tfrac12W[J]^\sharp D\lambda\}
  -\tfrac12\Lambda^\sharp D^\sharp W[J]\lambda .
\end{equation}
Here \(W[J]\lambda(b)=\tfrac i2
[\lambda(b_-)+\operatorname{Ad}(\mathbf U(b))\lambda(b_+),J(b)]\).
The first term is a bounded grade-three force.  The source of
\(\Lambda^\sharp\) in braces is bounded of grade four; the last
term uses the finite-\(p\) bound
\eqref{scf:bounded-divergence-adjoint}.  Thus
\(B_{\mathbf U}[J]:\mathcal A_\alpha\to\mathcal F_8\) has family norm
\(C\gamma_*\), with no difference of \(J\).
The Neumann series for \(G\) converges in these spaces after a
fixed decrease of \(\gamma_*\).

The first row inverse and \(H\) now follow from
\eqref{ncw:first-Schur-inverse}, using the actual bounded row
injection.  Its multiplier \(H^\sharp J\) keeps the fine-current
root and has \eqref{ncw:multiplier-family-bound}.
The second row jet has both absolute leg sums and the physical
fine density of \eqref{ncw:constraint-family}; it consequently
maps a bounded grade-one field to a bounded grade-three force.
Since \(Sv\) is bounded, and \(S^\sharp\) is bounded on
\(\mathcal F_8\),
\[
 \|\mathbf S^\sharp\mathbf{\Delta^{(2)}[H^\sharp J]}
                  \mathbf S\|_{\mathcal A_\alpha\to\mathcal F_8}
          \le C\gamma_* .
\]
The second Neumann series and second row inversion give
\(G_1,H_1\).  This argument does not assert that \(S\) preserves
\(\mathcal A_\alpha\) on the larger background domain.

It remains to check \(G_1DRD^\sharp G_1\).
For a bounded grade-three source \(f\), the first \(G_1\)
has \(C^{1,\alpha}\) output.  Therefore \(D^\sharp G_1f\)
is a grade-two \(C^\alpha\) site source.  The potential in
\(D^\sharp\) is Lipschitz, which suffices for this conclusion.
The same column formula for \(R\), whose site columns have bounded
first derivatives, preserves this \(C^\alpha\) class.
The final \(G_1D\) is bounded on it by
\cref{dnr:divergence-holder}, followed by Woodbury and the two
small-current series.  These are coefficientwise estimates in the
joined-carrier algebra.  They prove the asserted family for
\(\mathcal G\).  Likewise
\(P_1=D\widehat F\widehat M^{-1}\widehat F^\sharp D^\sharp\)
uses only bounded first derivatives of the site columns and of
the input field.

Finally each cell pairing for \(H^\sharp f\) costs
\(\xi^{-1}\xi^{-3}\xi^4=1\).  The absolute column sum proves
the raw row bound.  The raw-coordinate changes are precisely
the physical weights in \eqref{ncw:mass-removed-coarse-family}.
All series above converge normally on the same smaller \(C^1\)
tube.  They preserve their input and output collars, the read of
the independent current, and the specified coarse Gram inverses.
This proves both holomorphy and locality.
\end{proof}

\subsubsection{Cube parameters and local nonlinear coefficients}

Choose the partition \(\mathsf P\) after the inverse tiles.
Each coefficient is assigned its complete connected hull before
intersection with \(\Omega_0\).  Use the following cell convention.
For a tile \(q\) at scale \(\xi_q\), its \(j\)-enlargement has
width \(j\mathfrak M\xi_q\), \(0\le j\le5\), and its uncut hull
has width \(5\mathfrak M\xi_q\).  A cell hull is the union of these
tile hulls.  Choose the fixed lower floor for \(\mathfrak M\) so
that output and input supports fit the one-enlargements and
each successive nonlinear read fits the next enlargement.
The local inverse compressions and their bootstrap collars are
recorded within the five-enlargements; fixed reblocking, before
choosing the partition, achieves this for the finitely many proof
primitives.  This is a choice of the common geometric hierarchy,
not a deletion of rows from a previously specified sequence.
All factors in a product, including output
cutoffs, endpoint columns, current roots, and short joins, are
retained.  Let \(\mathsf P(C)\) be the cubes meeting this hull.
Use face connectivity, adding the finitely many intervening cubes
at a face or corner when necessary.  This convention is fixed
before introducing the parameters.

\begin{lemma}\label{dnr:parameter-families}
For \(\bar\kappa\ge0\), the partition side can be chosen so that
each family in \cref{dnr:linear-families} defines
\begin{equation}\label{dnr:decoupling}
 O_Y(s)=\sum_{\mathsf P(C)\subset Y}
           \Bigl(\prod_{\Delta\in\mathsf P(C)}s(\Delta)\Bigr)O_C,
       \qquad |s(\Delta)|\le e^{\bar\kappa},
\end{equation}
with normally convergent row and column sums.  For one constant
\(b=C e^{C\bar\kappa}\ge2\), all its weighted operator norms
needed in \eqref{dnr:norms} are at most \(b\), including
\(H_s^\sharp\) and \(P_{1,s}\).  A fixed positive exponential
decay rate remains.  No algebraic inverse identity is claimed for
the decorated maps at \(s\ne1\).
\end{lemma}

\begin{proof}
Choose a target side \(\ell_*\), a compatible power of \(L\),
at least \(C(1+\bar\kappa)\) times the largest fixed tile size
and the inverse of a reserved carrier rate.  Write \(p\) for
the side of the present cubic torus in the same units.
First suppose \(p\ge3\ell_*\), and take
\(\ell_{\mathsf P}=\ell_*\).
Cover the connected hull by balls whose radius is a fixed fraction
of \(\ell_{\mathsf P}\).  Starting at an endpoint, charge a new
ball to the first tree segment leaving the previous ball.
Disjoint charged segments have total length at most a fixed
multiple of the hull tree length \(\ell(C)\).  Each ball meets
at most a dimension-dependent number of partition cubes, and
the fixed carrier thickness is smaller than that radius.  Hence
\[
 \#\mathsf P(C)\le C_0+C_1\ell(C)/\ell_{\mathsf P}.
\]
The factor \(e^{\bar\kappa\#\mathsf P(C)}\) consumes at most
half the reserved rate and leaves \(e^{C_0\bar\kappa}\).

If \(p<3\ell_*\), use instead the coarsest compatible partition,
of side \(p/2\).  The declared tori have equal sides that are
twice a power of \(L\), so this partition has exactly \(2^4\)
cubes.  For every carrier,
\[
 \#\mathsf P(C)\le16,
 \qquad
 \left|\prod_{\Delta\in\mathsf P(C)}s(\Delta)\right|
       \le e^{16\bar\kappa}.
\]
No spatial decay is spent in this case.
When primitive patches require a cover, first construct and
deck-average the undecorated families on one common cover,
as in \cref{prop:ncw-named-families}.  Retain the full lifted
carrier and its collars when projecting it.  Only then assign
the product of parameters, once for each distinct base cube
meeting the projected complete hull.  Several lifted cubes or
carrier visits can therefore use the same base parameter only
once in this coefficient.  This prescription is squarefree in
the independent base parameters; it is not substitution of
tied variables into an already decorated cover monomial.
The parameter factor is unchanged by a deck translation, and
its absolute bound multiplies the descended row and column
majorants without a sheet-count factor.

The projected complete hull still contains every projected
input, output and nonlinear-read collar.  It is connected,
and its cube set is face connected with the fixed joins.
Thus a coefficient surviving an active mask lies in one wall
component.  The five-enlargement stocks are formed from these
complete projected hulls; a neighborhood path lifts starting
at its chosen carrier lift, so their successive-read inclusions
and component separation persist.  This changes only the
auxiliary parameter partition and supplies no identification
with a separately prescribed history partition.

In either case the remaining row and column sums are the
convergent carrier sums of \cref{dnr:linear-families}.
Increasing the fixed constant in \(b=C e^{C\bar\kappa}\)
to include \(16\) proves the asserted weighted bounds.
Normal convergence on smaller polydiscs proves holomorphy.
The coefficient prescription on complete common lifts proves
transport compatibility whenever the parameter partitions and
recorded geometry are carried as well.  Gauge averaging and
averaging with complex conjugation preserve the invariant
coefficient norms and the complete sums.
\end{proof}

For a union \(Z\) of active cubes, define \(\mathbb Y^j(Z)\)
to be the assigned sites in the \(j\)-enlargements of tiles
belonging to cells \(X\) whose complete uncut hull has
\(\mathsf P(X)\subset Z\).  This is the five-enlargement
convention for \(j=0,\ldots,5\), rather than erosion of
\(\Omega_0\cap\bigcup Z\).  The chosen floor gives
\(N_{r_0}(\mathbb Y^j(Z))\subset\mathbb Y^{j+1}(Z)\) for
\(j\le4\), and puts all these reads in \(\Omega_0\cap\bigcup Z\).
This preserves the actual determining sequence and the indicated
\(\mathbb Y^1,\mathbb Y^2,\mathbb Y^3\) stocks.
Retain nonlinear rows and the plaquette pairing at
\(\mathbb Y^2(Y)\), force values at \(\mathbb Y^1(Y)\), and
first-difference reads at \(\mathbb Y^3(Y)\).
A retained nonbare force has
all six incident plaquettes.  At a true bare boundary its cell
scale is the fine spacing.  Zero extension outside these stocks
defines the local maps.

Write \(\mathcal Q(A)=\widehat QA+C(A)\) for the actual iterated
straightened logarithmic average and
\[
 F(A)=\sum_{p\subset\Omega_0\cap\bigcup Y}a_p(e^{i\eta A}\mathbf U),
 \qquad V(A)=F'(A)-F'(0)-F''(0)A,
 \qquad \Delta'=F''(0)-\mathrm d^\sharp\mathrm d.
\]
Only the retained force values enter \(V,\Delta'\).
The associated bilinear form is
\(\mathcal E(A,B)=(\mathrm dA,\mathrm dB)_\eta\) on the
retained plaquette stock, with the same complete read convention.

\begin{lemma}\label{dnr:local-coefficients}
There are fixed \(c,C>0\) such that, for \(\|A_i\|_1\le r<c\),
the local maps just defined obey
\begin{align}
 |C(A)|_y&\le C\max_{N_{r_0}(y)}|A|_{\rm sc}^{,2},&
 |C(A_1)-C(A_2)|_w&\le Cr\|A_1-A_2\|_{1,w},\notag\\
 |C'_A{}^\sharp z|_{3,w}&\le Cr|z|_w,&
 |(C'_{A_1}-C'_{A_2})^\sharp z|_{3,w}
     &\le C\|A_1-A_2\|_{1,w}|z|,\label{dnr:local-bounds}\\
 |V(A)|_{3,w}&\le C\|A\|_1\|A\|_{1,w},&
 |V(A_1)-V(A_2)|_{3,w}&\le Cr\|A_1-A_2\|_{1,w}.\notag
\end{align}
The same bounds hold with the weight on the multiplier in the
second line.  The linear maps \(\widehat Q,\widehat Q^\sharp,
\Delta'\) have the corresponding fixed local bounds.  All maps
are holomorphic in the actual local data, covariant, real on real
data, and determined by their complete reads.  In particular
\(C(0)=C'_0=V(0)=V'(0)=0\).
\end{lemma}

\begin{proof}
For the row maps use the complete iterated-average expansion in
\cref{prop:averaging-1a-1-6}, with its absolute row and column sums
and physical transpose.  Iteration is in raw angles; the factors
of successive row lengths form the geometric sums in that proof.
The first nonlinear coefficient is quadratic.  Cauchy's formula
on a fixed larger amplitude ball gives the two derivative bounds;
one uses the absolute fixed-leg column sum for the transpose.
Local weights on a complete read differ by at most
\(e^{\theta r_0}\), which is incorporated in \(C\).

We record the cancellation needed for the action bound.  Rescale
one assigned cell to unit size and denote its fine mesh by \(h\).
For a plaquette transport the four field values to its base, writing
them as \(\alpha_1,\ldots,\alpha_4\), with the last two reversed.
Set
\[
 ih^2\Theta_p=\log(e^{-ih\alpha_3}e^{-ih\alpha_4}
                         e^{ih\alpha_1}e^{ih\alpha_2}),
 \qquad \Theta_p=(\mathrm dA)(p)+N_p.
\]
The exponential and logarithm series on a fixed small ball give
\(|N_p|\le Cs^2\), \(|N'_p|\le Cs\), where
\(s=\max|A|\).  Also \(|\Theta_p|\le C(g+s^2)\), where
\(g=\max|\nabla_{\mathbf U}A|\), and \(h|\Theta_p|\le Cs\).
Write the background plaquette as
\(c_p1+ih^2F_p\).  The two by two determinant identity gives
\(|c_p-1|\le C a_*^2h^4\).  Expansion of its trace against
\(e^{ih^2\Theta_p}\) shows that the plaquette derivative is
\[
 h^4\langle F_p+c_p\Theta_p+M_p,\delta\Theta_p\rangle,
 \qquad |M_p|\le Ch^4(|\Theta_p|^3+a_*|\Theta_p|^2).
\]
At a nonbare force bond pair the opposite incident plaquettes before
dividing by \(h\).  Transported differences of their \(N_p\)'s
are bounded by \(Chsg\), by the bound on \(N'_p\).
The terms \((c_p-1)\mathrm dA\) are bounded after the two
differences by \(Ca_*s\), using their \(h^4\) factor.
The terms \(M_p\) are bounded after one difference by
\(Cs^2(s+a_*)\), using \(h|\Theta_p|\le Cs\).
The remaining differentiated \(N_p\) factor costs
\(Cs(a_*+g+s^2)\).  Therefore
\[
 |F'(A)-F'(0)-\mathrm d^\sharp\mathrm dA|
       \le Cs(a_*+\max\{s,g\}).
\]
Put \(n=\max\{s,g\}\) and
\(w(z)=F'(zA)-F'(0)-z\mathrm d^\sharp\mathrm dA\).
The preceding estimate gives
\(|w(z)|\le C|z|s(a_*+|z|n)\) whenever \(|z|s\le c_0\),
for a fixed \(c_0>0\).  The case \(s=0\) is immediate.
For \(s>0\), Cauchy's formula on \(|z|=R\) bounds the
linear coefficient by \(Cs(a_*+Rn)\).  Letting \(R\downarrow0\)
gives \(|\Delta'A|\le Ca_*s\), without a first-difference term.
To bound the higher coefficients, take \(T=c_0/s\ge2\),
decreasing the field radius if necessary.  Cauchy's formula gives
\(|[z^q]w|\le CsT^{1-q}(a_*+Tn)\), so
\[
 |V(A)|\le2Cs(a_*/T+n)\le Csn .
\]
The last inequality uses the fixed bound for \(a_*\) and \(s\le n\).
This proves the bounds on \(\Delta'\) and \(V\).  At a bare cell \(h=1\); direct absolute estimates
on the at most six plaquettes give the same bounds.  The small
amplitude ball is fixed before the mesh is chosen.
Polarization, or Cauchy's formula along the segment between two
fields on a ball twice as large, gives the Lipschitz estimate.
Rescaling supplies the factors \(\xi_y,\xi_y^2,\xi_y^3\).
All identities used are local products, convergent logarithms,
and physical transposes.  They prove the final assertions as well.
\end{proof}

\subsubsection{The two fixed points and the current}

Write \(C''[a]=(\mathrm d C_2)_a\), where \(C_2\) is the
quadratic part of \(C\).  For \(\|a\|_1<c/b\) define
\begin{equation}\label{dnr:straightening}
 D_s(a)=C(a-H_sD_s(a)),\qquad T_s(a)=a-H_sD_s(a),\qquad
 E_s(a)=C'_{T_s(a)}{}^\sharp
             (1+H_s^\sharp C'_{T_s(a)}{}^\sharp)^{-1}.
\end{equation}
The first equation is solved by iteration from zero.  Define, for
every parameter value, the following particular nonlinear force:
\begin{align}
 W_s(a)={}&-[E_s(a)-C''[a]^\sharp]H_s^\sharp\mathbf J
       -\widehat Q^\sharp(\Sigma_s-\mathsf a)D_s(a)\notag\\
 &-E_s(a)(\Sigma_s-\mathsf a)\widehat Qa
       +E_s(a)(\Sigma_s-\mathsf a)D_s(a)\notag\\
 &+V(T_s(a))-E_s(a)H_s^\sharp V(T_s(a)).
 \label{dnr:force}
\end{align}
This is the force specified in \eqref{eq:elab-reduced-force}.
Its definition uses bounds on the decorated maps, and does not use
\(\widehat QH_s=1\) or a slice identity at general \(s\).

\begin{proposition}\label{dnr:response}
There are fixed \(c,C>0\) such that the following holds with
\(r=c b^{-8}\), where \(b\) is the constant of
\cref{dnr:parameter-families}.  For \(|B'|<r\), the equations
\begin{equation}\label{dnr:response-equations}
 a=H_{1,s}B'+A_1,\qquad A_1=-\mathcal G_sW_s(a),\qquad
 \mathbf H=T_s(a)
\end{equation}
and \eqref{dnr:straightening} have a unique solution in
\(\|A_1\|_1\le br\), \(|D_s|\le Cb^2r^2\).
They are exactly \eqref{eq:rg-objects-R9} in raw coordinates.
The solution is holomorphic jointly in \(s,B',\mathbf U,\mathbf J\)
on the indicated open domains and satisfies
\begin{equation}\label{dnr:response-bounds}
 \|\mathbf H\|_{1,w}\le Cb|B'|_w,\qquad
 \|A_1\|_{1,w}\le Cb^5|B'||B'|_w,
 \qquad \|\delta\mathbf H\|_{1,w}\le Cb|\delta B'|_w.
\end{equation}
Here the difference is between two data in the same ball.
The field and its datum derivative also satisfy the corresponding
\(\mathcal A_{1/4}\) output estimates, with a fixed larger power
of \(b\).

For \(\|A\|_1\le4br\), put
\begin{align}
 \langle H_s^\circledast(A)(c),e\rangle
   & =\mathcal E(A,H_se_c)
       +(\Delta'A+V(A),H_se_c)_\eta,\notag\\
 \Psi_{s,A}&=(1+H_s^\sharp C'_A{}^\sharp)^{-1},\qquad
 \omega_s(A)=\Psi_{s,A}[H_s^\sharp\mathbf J+H_s^\circledast(A)],
 \label{dnr:multiplier}\\
 j_s(A)&=f_A^{-1}(\widehat Q+C'_A)^\sharp\omega_s(A)
                 -\widehat Q^\sharp H_s^\sharp\mathbf J,\notag
 \end{align}
where \(f_A(b)=\sum_{n\ge0}(-i\eta\operatorname{ad}_{A(b)})^n/(n+1)!\).
Then the reduced response is
\(\mathfrak h_Y=(\mathbf H,j_s(\mathbf H))\).  It vanishes at
\(B'=0\), and
\begin{equation}\label{dnr:current-bound}
 |j_s(\mathbf H)|_{3,w}
       +|P_{1,s}\mathbf H|_{3,w}\le Cb^4|B'|_w,
 \qquad |\delta j_s(\mathbf H)|_{3,w}\le Cb^4|\delta B'|_w.
\end{equation}
The shifted independent current is \(\mathbf J+j_s(\mathbf H)\).
\end{proposition}

\begin{proof}
Let \(\|a\|_1\le\varepsilon<c/b\).
By \cref{dnr:local-coefficients}, the first equation of
\eqref{dnr:straightening} maps the row ball \(|D|\le C\varepsilon^2\)
into itself and has Lipschitz constant \(Cb\varepsilon<1/4\).
Its solution satisfies
\[
 |D_s(a)|_w\le C\varepsilon\|a\|_{1,w},\quad
 |\delta D_s|_w\le C\varepsilon\|\delta a\|_{1,w},\quad
 \|T_s(a)\|_1\le2\varepsilon,
 \quad\|\delta T_s\|_{1,w}\le2\|\delta a\|_{1,w}.
\]
Differentiation of that fixed-point equation gives the last
expression in \eqref{dnr:straightening}.  In particular
\(E_s(a)=C''[a]^\sharp+O(b\varepsilon^2)\) as a map from
raw multipliers to bounded forces.  To check the quadratic error,
write
\(C'_{T_s(a)}-C''[a]=(C'_{T_s(a)}-C''[T_s(a)])
+C''[T_s(a)-a]\).
The first term is quadratic by Cauchy's formula and the second
is bounded by \(Cb\varepsilon^2\); the inverse in
\eqref{dnr:straightening} differs from one by \(Cb\varepsilon\).
The same decomposition gives its difference bound
\(Cb\varepsilon\|\delta a\|_{1,w}\).

Insert these bounds into each of the six terms of
\eqref{dnr:force}.  The first term uses
\(|H_s^\sharp\mathbf J|\le b\gamma_*\).
The result, with a fixed larger \(C\), is
\begin{equation}\label{dnr:force-estimate}
 |W_s(a)|_{3,w}\le Cb^2\varepsilon\|a\|_{1,w},\qquad
 |\delta W_s|_{3,w}\le Cb^2\varepsilon\|\delta a\|_{1,w}.
\end{equation}
For \(\|A_1\|_1\le br\), \(|B'|\le r\), take
\(\varepsilon=2br\).  The outer map in
\eqref{dnr:response-equations} has norm at most \(Cb^5r^2\)
and Lipschitz constant at most \(Cb^4r<1/4\).
It maps the stated ball into itself.  Banach's theorem proves
existence and uniqueness.  Subtracting the two fixed-point
equations in a weighted norm and absorbing the same constant
proves \eqref{dnr:response-bounds}.  The iterates are holomorphic
and converge uniformly on smaller polydiscs, proving joint
holomorphy.  The construction remains valid on each local chart
of the regular-pair domain; agreement on overlapping charts
follows from the identical equations and contraction uniqueness.
For the stronger output assertion, let
\(\|\cdot\|_{\mathcal A_{1/4};w}\) be the norm in
\eqref{psr:strong-norm}, with each complete local-box norm
multiplied by \(w(y)\).  The weight is not differentiated.
The strong families in \cref{dnr:linear-families}, with the
parameter summation of \cref{dnr:parameter-families}, give
\[
 \|H_s z\|_{\mathcal A_{1/4};w},\quad
 \|H_{1,s}z\|_{\mathcal A_{1/4};w}
       \le b_{\rm str}|z|_w,\qquad
 \|\mathcal G_s f\|_{\mathcal A_{1/4};w}
       \le b_{\rm str}|f|_{3,w},
 \qquad b_{\rm str}\le Cb^q
\]
for a fixed exponent \(q\).  These are first-gradient
H\"older output norms on the same background domain.
For two data in the indicated ball, subtract the exact equations:
\[
 \delta\mathbf H=H_{1,s}\delta B'
                 -\mathcal G_s\delta W_s-H_s\delta D_s.
\]
The weighted fixed-point estimates already give
\(\|\delta a\|_{1,w}\le Cb|\delta B'|_w\).
With \(\varepsilon=2br\), the force and straightening bounds yield
\[
 |\delta W_s|_{3,w}\le Cb^4r|\delta B'|_w,
 \qquad |\delta D_s|_w\le Cb^2r|\delta B'|_w.
\]
Consequently
\(\|\delta\mathbf H\|_{\mathcal A_{1/4};w}
\le Cb^q|\delta B'|_w\), since \(r=cb^{-8}\).
Taking one datum to be zero bounds the field itself.
At every datum in the open ball, divide by the size of a
complex datum increment and pass to the derivative.  Holomorphy
and continuity of the finite-lattice norm give the same bound
for the datum derivative on that whole ball.

For the current, the physical gradient pairing gives
\[
 |H_s^\circledast(A)|_w\le Cb(\|A\|_{1,w}
                            +\|A\|_1\|A\|_{1,w}).
\]
In the curl term the two first differences have grades two and
two, so the cell factor is \(\xi^{-2}\xi^{-2}\xi^4=1\).
The other terms use the field/force pairing and its column sum.
The same estimate holds for differences.  Since
\(\|H_s^\sharp C'_A{}^\sharp\|\le Cb\|A\|_1<1/4\),
\(\Psi_{s,A}\) is uniformly bounded.  It follows that
\[
 |\omega_s(A)-H_s^\sharp\mathbf J|_w
       \le Cb(1+b\gamma_*)\|A\|_{1,w}.
\]
The exponential series gives \(\|f_A^{-1}\|\le2\) and
\(\|f_A^{-1}-1\|\le C\eta|A|\), with the corresponding
difference bounds.  Expand the last line of
\eqref{dnr:multiplier} around \(A=0\); its constant term cancels.
The local transpose estimate and the preceding multiplier bound
give \(|j_s(A)|_{3,w}\le Cb^3\|A\|_{1,w}\), and the same
Lipschitz estimate.  Combine this with
\eqref{dnr:response-bounds} and the bound on \(P_{1,s}\).
\end{proof}

\begin{proposition}\label{dnr:walls}
The response of \cref{dnr:response} is real on real data and real
parameters, gauge covariant, and compatible with common-lift
transport.  If \(Z\) is a union of wall components of
\(\{\Delta:s(\Delta)\ne0\}\), its restriction is the response
formed using precisely those ambient carrier coefficients whose
complete cube sets lie in \(Z\).  The field is supported over
\(\mathbb Y^1(Z)\); its reduced current and projected field are
supported over \(\mathbb Y^2(Z)\).  All their complete native
reads lie over \(\mathbb Y^3(Z)\).  A component with no retained
datum has zero response.

If every cube meeting any complete cell hull has parameter one,
all three fields equal their all-one fields on that same sequence.
At the all-one point their complete sums are independent of the
auxiliary carrier tiling and parameter partition.  They commute
with an identification of two regional sequences that preserves
all local field, row, gauge, and geometry data in common lifts.
For a general proper mask the derivative of the first response at
zero is \(H_{1,s}\); its linearized constraint is
\(\widehat QH_{1,s}B'\).
\end{proposition}

\begin{proof}
The cube set of a surviving coefficient is face connected and
contains its complete input and output read collars.  It therefore
lies in one active wall component.  Formula
\eqref{dnr:decoupling} splits into a direct sum over these components.
The local maps in \cref{dnr:local-coefficients} cannot cross their
reserved collars.  Starting both iterations at zero proves the
same direct-sum assertion for every iterate and hence for its limit.
The transpose, curl pairing, and local row derivative in
\eqref{dnr:multiplier} use at most the next two reads; this proves
the asserted support of the current.  The field, current and
projected field vanish on a component with zero datum because
the zero iterate remains zero there.  An indicator of a component
consequently creates no commutator on these complete native reads.

Every coefficient and local operation is real, covariant and
transported with its complete determining data.  The same iteration
argument proves these properties.  At the specified one-valued
cubes every coefficient weight is one and the local structure is
the whole-sequence structure.  Absolute summation gives the actual
operators, independently of their expansion.  Their physical
formulas and local structure commute with the stated identification;
uniqueness identifies the responses.  Finally
\(D_s(0)=D'_s(0)=W_s(0)=W'_s(0)=0\), which gives the last
assertion by differentiation.  Thus restriction to an active
component identifies its restricted ambient solve.  Identification
with a different sequence's solve requires the indicated equality
of determining data and operators.
\end{proof}

\subsubsection{Physical identification and the coupled response}

\begin{proposition}\label{dnr:physical}
Let \(Y\) contain every complete cell hull of the determining
sequence, let \(s=1\) there, and let \(U_0\) be a real critical reference in the
regular class, with its physical current
\(J(U_0)=\widehat Q^\sharp\lambda_0\).
Use the closed noninward determining stock, including the complete
preceding block at every outward row.
For the response of \cref{dnr:response},
\begin{equation}\label{dnr:physical-equations}
 \mathcal Q(\mathbf H)=B',\qquad RD^\sharp\mathbf H=0,\qquad
 F'(\mathbf H)=(\widehat Q+C'_{\mathbf H})^\sharp
                         \omega_1(\mathbf H).
\end{equation}
Consequently \(J(U_0)+j_1(\mathbf H)=J(e^{i\eta\mathbf H}U_0)\)
and \(P_{1,1}\mathbf H=DD^\sharp\mathbf H\).
Every sufficiently small critical field on this slice is this
response at its own reduced datum.  The node factors and the
change between ordinary and reduced outward data are the actual
ones of \eqref{mb:reduced-datum}--\eqref{mb:reduced-bounds}.
They depend only on the complete corresponding row or top pillar,
are holomorphic on a fixed larger small ball, and commute with
restriction and common-lift transport on those reads.
\end{proposition}

\begin{proof}
The current identification in
\cref{prop:acc-current-identification} identifies the second
precision on the Landau slice with the Hessian of the constrained
action.  In particular \(\mathcal G\) is its inverse on
\(\mathcal K=\ker\widehat Q\cap\ker RD^\sharp\), and
\(H,H_1\) lie in the Landau slice, with
\(\widehat QH=\widehat QH_1=1\).
The first equation of \eqref{dnr:straightening} therefore gives
the first two equations of \eqref{dnr:physical-equations}.
Differentiate the straightening map.  On \(\mathcal K\) its
derivative maps bijectively onto
\(\ker(\widehat Q+C'_{\mathbf H})\cap\ker RD^\sharp\), since
it differs from the identity by \(O(b\|a\|_1)\).
Here is the exact action calculation.  Write
\(\Delta_\pi=F''(0)-B_{U_0}[J]\),
\(\Delta_1=\Delta_\pi-S^\sharp\Delta^{(2)}[H^\sharp J]S\), and
\(V_0(T)=F(T)-F(0)-(J,T)-\tfrac12(T,F''(0)T)\).
For \(a\) on the Landau slice, put \(D=D_1(a)\), \(T=a-HD\),
and \(P=\Sigma-\mathsf a\).  The exact identities
\(\Delta_\pi H=\widehat Q^\sharp P\),
\(H^\sharp\Delta_\pi H=P\), and
\(\tfrac12(a,\Delta^{(2)}[H^\sharp J]a)
=(C_2(a),H^\sharp J)\) give
\begin{align}
 F(T)&=F(0)+(J,a)+\tfrac12(a,\Delta_1a)+\mathcal R(a),\notag\\
 \mathcal R(a)&=-(D-C_2(a),H^\sharp J)
       -(\widehat Qa,PD)+\tfrac12(D,PD)+V_0(T).
 \label{dnr:pullback-action}
\end{align}
The pairings in each term are its physical fine or raw row pairing.
The first two identities follow from
\(G^{-1}H=\widehat Q^\sharp\Sigma\) and \(RD^\sharp H=0\).
The current completion vanishes bilinearly on two Landau fields,
so \(F''(0)\) may be replaced by \(\Delta_\pi\) in their pairing.
Differentiating \eqref{dnr:pullback-action}, using
\((\mathrm dD)_a^\sharp=E_1(a)\), yields precisely
\[
 (F'(T),(\mathrm dT)_a\delta)
       =(J+\Delta_1a+W_1(a),\delta)
\]
for every slice vector \(\delta\).  The six derivatives of
\(\mathcal R\) are the six terms of \eqref{dnr:force}, with
their displayed signs.  On \(\mathcal K\),
\(\Delta_1\) is the precision inverted by \(\mathcal G\),
while \(\Delta_1H_1B'\) and \(J\) annihilate \(\mathcal K\).
Thus the outer fixed-point equation makes \(F'(\mathbf H)\)
annihilate the constrained tangent space.

Finite-dimensional annihilator duality now gives
\(F'(\mathbf H)=(\widehat Q+C'_{\mathbf H})^\sharp\lambda
+DR\phi\).  Pairing with \(H\) gives
\(\lambda=\omega_1(\mathbf H)\), because \(H^\sharp DR=0\)
and \(H^\sharp F'(\mathbf H)=H^\sharp J+H^\circledast(\mathbf H)\).
The remaining gauge multiplier vanishes.  Here no uniform
second-derivative gauge estimate is needed: on each fixed finite
lattice, the infinitesimal gauge directions with trivial reduced
site averages map onto \(\operatorname{ran}R\) under
\(RD^\sharp\) at zero, by the actual site inverse.  Surjectivity
persists in some neighborhood of zero.  The action is invariant
on those directions and their reduced constraint derivative is
zero.  For real data in this neighborhood, pair with a gauge
direction whose \(RD^\sharp\) value is \(R\phi\).  The result is
\(\|R\phi\|_2^2=0\).  The residual is holomorphic throughout
the uniform datum ball and vanishes on its real neighborhood of
zero; the identity theorem extends its vanishing to that ball.
This proves the last equation in \eqref{dnr:physical-equations}.

Differentiating \(e^{i\eta A}\) shows directly that
\(F'(A)=f_AJ(e^{i\eta A}U_0)\).  Inserting this in
\eqref{dnr:multiplier} gives the physical current identity.
The projection identity follows from \(RD^\sharp\mathbf H=0\).
Conversely a sufficiently small critical field on this slice has
a unique multiplier, since
\((\widehat Q+C'_A)H=1+C'_AH\) is invertible.  Reversing the
straightening calculation makes it a solution of the same two
contractive equations; uniqueness identifies it.

The node statements follow from the exact iterated averaging
identity: each elementary node factor is the small logarithmic
mean of finitely many transports in its block.  Its convergent
exponential and logarithm formulas read only that block.
Successive factor errors have the length factors of the averaging
iteration; their geometric sum gives the fixed ball and bound
in \eqref{mb:reduced-bounds}.  At an outward row the only additional
factor is the node factor of its preceding lower block.  Solving
the triangular near-identity change of raw data gives
\eqref{mb:reduced-inverse}, with the same complete read.  These
identities also prove the covariance and transport assertions.
\end{proof}

\begin{proposition}\label{dnr:coupled}
The actual coupled system \eqref{eq:propagators-II}, using the
families of \cref{dnr:parameter-families}, has a unique small
solution when \(\gamma+|B'|\le c b^{-3}\), where
\(\gamma=|\mathbf J|_3\).  In the centered multiplier variable
\(\mu=\lambda-H_s^\sharp\mathbf J\), it satisfies
\[
 \|T\|_{1,w}+|\mu|_w\le Cb|B'|_w.
\]
It has the same holomorphy, weighted locality, covariance and
transport properties on its own complete reads.  The component
restriction and support assertions apply to \(T\) and
\(\mu+\mathsf aB'\).  The centered multiplier \(\mu\) itself
retains the local term \(-\mathsf aB'\) outside the active components.
At \(B'=0\), \(T=\mu=0\).  At the whole-sequence all-one
point specified in \cref{dnr:physical}, over a physical critical
reference and on the same closed determining stock, it is the same
sufficiently small constrained critical
field as \cref{dnr:physical}, after the prescribed change of raw
data.  Equality of the two decoupled systems away from \(s=1\)
is not asserted.
\end{proposition}

\begin{proof}
Substitute \(\lambda=H_s^\sharp\mathbf J+\mu\) into
\eqref{eq:propagators-II}.  Its constant term is
\((H_{0,s}B',(\Sigma_s-\mathsf a)B')\).
On the ball \(\|T\|_1+|\mu|\le Cb r\), the remaining terms
have Lipschitz constant at most \(Cb^2(\gamma+r)\), by
\cref{dnr:local-coefficients}.  Decrease \(c\) so this is below
\(1/4\); the ball is invariant.  The weighted subtraction proof
and normally convergent iterates give all the quantitative and
analytic assertions.  Carrier supports give transport as in \cref{dnr:walls}.
For the component assertion retain the local multiplier term.
Writing \(N_T=V(T)-C'_T{}^\sharp\lambda\), the second equation is
\[
 \mu+\mathsf aB'
   =\Sigma_s[B'-C(T)]+H_s^\sharp N_T+\mathsf aC(T).
\]
The first two terms have their indicated carrier supports, and
the last reads only the next complete row neighborhood of the
support of \(T\).  The reserved collars separate these reads
between active wall components.  For a component carrier
\(\mathbb Y_i=\mathbb Y^1(c_i)\), the field \(T_i\) is supported
over \(\mathbb Y_i\), and \(\mu_i+\mathsf aB'\) is supported on
its first complete row neighborhood.  The datum is read on that
neighborhood: the local term \(-\mathsf aB'\) also enters
\(C'_{T_i}{}^\sharp\lambda_i\).  These reads lie in
\(\mathbb Y^2(c_i)\).  If \(T_i,\mu_i\) are obtained
using the maps of component \(i\), with the same full datum,
the fixed-point equations and uniqueness give
\[
 T=\sum_iT_i,\qquad
 \mu=-\mathsf aB'+\sum_i(\mu_i+\mathsf aB').
\]
In verifying these equations, \(\lambda\) and its component
counterpart agree on every row read by \(C'_T{}^\sharp\) in
that component; their common remaining term is
\(-\mathsf aB'\).  Thus the assertion also holds when no
carrier survives: then \(T=0\) and \(\mu=-\mathsf aB'\).
For the physical identification start with the already constructed
response of \cref{dnr:physical} on the common smaller datum ball.
Put \(K_T=\widehat Q+C'_T\) and
\(R_T=V(T)-C'_T{}^\sharp\lambda\).  Its Euler equation and slice
condition give
\[
 G_0^{-1}T=\widehat Q^\sharp
          [\lambda-\lambda_0+\mathsf a\widehat QT]-R_T.
\]
The exact identities
\(\widetilde G\widehat Q^\sharp=0\) and
\(\widetilde GG_0^{-1}=1-H_0\widehat Q\) give the first coupled
equation.  For the second, pair the Euler equation with \(H\).
Since \(H,T\) are Landau fields, their current-completion pairing
vanishes, and
\[
 H^\sharp F''(0)T=H^\sharp\Delta_\pi T
       =(\Sigma-\mathsf a)\widehat QT,\qquad
 H^\sharp J=\lambda_0.
\]
It follows that
\(\lambda=H^\sharp J+(\Sigma-\mathsf a)\widehat QT
+H^\sharp R_T\), the second coupled equation.
Contraction uniqueness identifies the two fields on this common
ball.  The coupled field's Euler, constraint, and Landau residuals
are holomorphic on its entire connected datum ball and vanish
on that smaller ball.  The identity theorem proves their
vanishing throughout the coupled ball.  The smaller
independent-current radius is part of this coupled-system
statement and is not imposed on \cref{dnr:response}.
\end{proof}

\begin{corollary}\label{dnr:flat-identification}
At the trivial reference, the Landau representative of a sufficiently
small regular box minimizer is
\(\mathbf H(1;b(B))\), with its actual ordinary row logarithm \(B\),
the reduced datum \(b(B)\) of \eqref{mb:reduced-datum}, and its
actual node factors.  Suppose a common local gauge gives
\(|B(c)|\le C\xi_c a_0\) on the full determining stock.
Then the potential has size \(Ca_0\), and its first differences
have size \(Ca_0/\xi_y\) at cell \(y\).  On complete stencils
its normalized first differences have a fixed H\"older bound.
The datum derivative has exponentially screened first-order and
H\"older output bounds.  Node factors on a complete common top
pillar have a bounded holomorphic chart and the corresponding
Cauchy difference bound.
\end{corollary}
\begin{proof}
The trivial reference is critical and has zero current.  The
characterization in \cref{dnr:physical} applies to the small Landau
representative.  The outward datum change reads one preceding
block only, with neighboring row lengths in a fixed ratio; it
preserves the profile \(\xi_c a_0\).  Integrate the weighted
datum derivative in \eqref{dnr:response-bounds} from zero.
The ratio \(\xi_c/\xi_y\) is absorbed in the fixed reserved
exponential rate of the assigned-cell kernels.  This gives the
two displayed sizes.  The strong output conclusion following
\eqref{dnr:response-bounds} gives the H\"older estimate.
For the node factors use their complete-read holomorphic formula
and fixed larger ball in \cref{dnr:physical}; Cauchy's formula
gives the difference bound on a smaller ball.  These statements
apply to the actual ordinary datum, and include its outward
change of variables.
\end{proof}

\begin{corollary}\label{dnr:interpolation}
The response hypotheses in \eqref{eq:rg-objects-R9} are supplied
by \cref{dnr:response,dnr:walls,dnr:physical} on the original
first-order regular-pair domain, with the fixed small radii above.
For any function \(F\) whose complete core contains
\((\rho_XP)\dotplus\nu\mathfrak h_Y\) on its stated \(\nu\)-disc,
\eqref{eq:rg-objects-R10} is a holomorphic interpolation with the
wall splitting and rooted alternating-bracket properties stated
there.  A change of units multiplies its potential by \(L^{-m}\)
and its current by \(L^{-3m}\), as in
\eqref{eq:rg-objects-R7}.  Core inclusion is checked for the
particular function and its full auxiliary disc; it is separate
from solving the response equations.
\end{corollary}

\begin{proof}
The first assertion collects the preceding results in the specified
raw-coordinate convention.  The shift is the small-branch product
in \eqref{eq:rg-objects-R8}; its logarithm is holomorphic on the
indicated core.  Composition gives the interpolation.  The support
and splitting proof in \cref{dnr:walls} gives the vanishing required
for the alternating brackets, and the Cauchy calculation following
\eqref{eq:rg-objects-R12} applies on the complete disc.  The last
assertion is the physical scaling of a potential and a current.
\end{proof}

% BEGIN LAYERED EXPANSION background/homogeneous_strong_inverse.tex
\subsection{A strong field--multiplier inverse at every homogeneous depth}
\label{ch:hb-strong-inverse}

The energy gap of the homogeneous Landau precision can be upgraded to
the strong physical norm needed for a nonlinear background equation.
We prove that upgrade for the actual flat rows, uniformly in the number
of fine bonds in a block.  The argument uses the already proved energy
gap, but assumes no supremum-norm vector inverse.  In particular an
order-zero singular integral is not applied at the $L^\infty$ endpoint.

Fix the odd blocking factor $L\ge3$ and an integer depth $k\ge0$.
The coarse spacing is one, the fine spacing is $h=L^{-k}$, and the
periodic side is an integer $M\ge2$.  A real flat reference may have
arbitrary commuting unitary holonomies.  Work in its parallel frame on
the universal cover.  Fields are then quasiperiodic; all local norms and
matrix products below are understood in that frame.  Fine pairings use
$h^4$ times counting measure and coarse pairings use counting measure.
The transpose $\sharp$ always uses those pairings.  Complexification is
linear, with the analytic bilinear transpose, rather than conjugating a
complex parameter.

Fix $0<\alpha<1$.  For a fine field or site function define
\begin{equation}
 \begin{split}
 \|u\|_{X_\alpha}&=\|u\|_\infty+\|Du\|_\infty+[Du]_{\alpha,\rm loc},\\
 \|u\|_{Y_\alpha}&=\|u\|_\infty+\|Du\|_\infty
                         +\|D^2u\|_\infty+[D^2u]_{\alpha,\rm loc},\\
 \|u\|_{Z_\alpha}&=\|u\|_\infty+[u]_{\alpha,\rm loc}.
 \end{split}
 \label{hb:norms}
\end{equation}
All products of forward differences of the indicated order are included.
The local seminorm compares pairs at physical distance at most one,
using parallel transport of the flat reference.  Differences at larger
distances are controlled by the supremum term.  This definition is
independent of a choice of fundamental cube, up to fixed constants.

Let $q$ be the actual flat site average, $Q$ the actual flat
straightened field row, and $D,d$ the flat gradient and curl.  Thus
\begin{equation}
 (q\lambda)(y)=h^4\sum_{x\in B(y)}\lambda(x),\qquad
 (Qv)_\mu(y)=h^5\sum_{x\in B(y)}\sum_{t=0}^{h^{-1}-1}
                                      v_\mu(x+th e_\mu),
 \quad QD=\mathrm{d}_cq.
 \label{hb:actual-flat-rows}
\end{equation}
The formula is in the parallel frame; the corresponding transports are
restored on the quotient torus.  The plain nonlinear group average has
a different derivative before straightening.

\subsubsection{Smooth exact right inverses of the actual rows}

\begin{lemma}[Uniform physical row seeds]
\label{lem:hb-row-seeds}
There are exact right inverses $E_q,E_Q$ with
\begin{equation}
 qE_q=1,\qquad QE_Q=1,\qquad
 \|E_qb\|_{Y_\alpha}+\|E_Qc\|_{Y_\alpha}
                         \le C_\alpha(\|b\|_\infty+\|c\|_\infty).
 \label{hb:seed-strong-bounds}
\end{equation}
Their columns have fixed physical support and bounded overlap.  They
also have uniformly bounded physical $H^1$ energy from coarse $\ell^2$.
These constants are independent of $h,M$ and flat holonomy.
\end{lemma}
\begin{proof}
Write $h=1/n$, with $n$ odd.  Choose an even nonnegative smooth
function $\varphi$ supported in $(-1/4,1/4)$ and positive on
$[-1/8,1/8]$.  Set
\begin{equation}
 m_0=h\sum_{j}\varphi(jh),\qquad
 m_2=h\sum_j(jh)^2\varphi(jh).
 \label{hb:seed-moments}
\end{equation}
For $n\ge9$ both have uniform positive lower bounds.  Their limits
are the positive integrals; the finitely many earlier values are
positive, since $\pm h$ lie in the positive core.  They have uniform
upper bounds by the fixed support.  The scalar site bump
$m_0^{-4}\prod_\nu\varphi(x_\nu-y_\nu)$ is an exact column of
$E_q$.  Different pillars have disjoint supports.

For a positive coarse bond $c=(y,y+e_\mu)$ put
\begin{equation}
 (E_Qc)_\mu(x)
 =\left(a_h+b_h(x_\mu-y_\mu)\right)
                    \prod_\nu\varphi(x_\nu-y_\nu),\qquad
 a_h=m_0^{-4},\quad b_h=\frac{1-h}{2m_0^3m_2},
 \label{hb:signed-field-seed}
\end{equation}
with all other components zero, and multiply by the desired colour.
Only the two rows $c$ and $(y-e_\mu,y)$ can see this seed.  At
longitudinal coordinate $s=x_\mu-y_\mu$, their exact weights in
\eqref{hb:actual-flat-rows} are respectively
$s+(1+h)/2$ and $(1-h)/2-s$, with the common physical fine
weight $h^4$.  This follows by counting the possible starting points
of a straight segment.  In one dimension, at site index
$j\in[-(n-1)/2,(n-1)/2]$, the two counts are
$j+(n+1)/2$ and $(n-1)/2-j$.

The transverse mean of the seed is $m_0^3$.  Its longitudinal
zeroth and first moments are $a_hm_0$ and $b_hm_2$.
The two row values are therefore one and zero, exactly.  The
coefficients $a_h,b_h$ are uniformly bounded, so these signed
smooth seeds and all their fixed-order physical differences are
uniformly bounded.

For the finitely many $n<9$, use direct seeds.  The site row has
the right inverse $q^\sharp$.  For a field row use its middle
positive face bond, at longitudinal index $(n-1)/2$ and transverse
block centres.  Its coefficient in that row is $h^4$, and its
coefficient in the preceding row is zero.  Put $h^{-4}$ times
the prescribed colour there.  This includes $n=1$, when $Q=1$.
All finite-difference norms are bounded by a fixed constant in this
finite range of $n$.  On a two-cell periodic direction the two
positive coarse bonds remain distinct, so the same count applies.

The columns have bounded support and overlap.  Squaring their
pointwise first-difference bounds and summing with physical measure
proves the $\ell^2\to H^1$ energy bounds.  Their continuous smooth
profiles, or the fixed coarse-mesh difference bounds, prove
\eqref{hb:seed-strong-bounds}.  Unit parallel transports preserve
all these estimates and give the construction for arbitrary flat
holonomy.
\end{proof}

\subsubsection{The heat regularization used in the strong norm}

For a field $u$ write
$\|u\|_{2,\rm cell,\infty}=\sup_y\|u\|_{L^2(B_y)}$.
Let $L=-\sum_\mu D_\mu^-D_\mu^+$ and $P_t=e^{-tL}$.
The differentiated heat kernels in the discrete H\"older estimate
give, for $0<t\le1$,
\begin{equation}
 \|D^jp_t\|_1\le C_j(t+h^2)^{-j/2}.
 \label{hb:heat-L1}
\end{equation}
At $t=1$ their physical $L^2$ norms on successive unit cells have
summable tails. The differentiated fixed-time Fourier bound gives a
uniform pointwise bound, and the proved heat-tail estimate gives
arbitrarily fast decreasing cell $L^1$ tails. Apply
$\|k\|_{L^2(C)}^2\le\|k\|_\infty\|k\|_{L^1(C)}$ with a sufficiently
high tail power. Equivalently use the pointwise estimate
\eqref{bg:heat-pointwise} directly.
Consequently $D^jP_1$ maps the cell-$L^2$ supremum norm to
$L^\infty$, uniformly in the mesh.  Twisted torus kernels are
periodizations with unitary phases; their absolute bounds are
dominated by the same infinite-lattice sums.

\begin{lemma}[Strong estimates from a bounded scalar Laplacian]
\label{lem:hb-heat-strong}
If $Lu=g$, then
\begin{equation}
 \|u\|_{X_\alpha}
       \le C_\alpha(\|u\|_{2,\rm cell,\infty}+\|g\|_\infty).
 \label{hb:heat-X-bound}
\end{equation}
If $g\in Z_\alpha$, then
\begin{equation}
 \|u\|_{Y_\alpha}
       \le C_\alpha(\|u\|_{2,\rm cell,\infty}+\|g\|_{Z_\alpha}).
 \label{hb:heat-Y-bound}
\end{equation}
The same estimates hold with a fixed small exponential spatial weight,
using local cell seminorms for the weighted norms.
\end{lemma}
\begin{proof}
The exact finite-time identity is
\begin{equation}
                 u=P_1u+\int_0^1P_tg\,dt.
 \label{hb:finite-time-identity}
\end{equation}
The first term and all its fixed derivatives are controlled by the
cell-$L^2$ norm.  The zeroth and first derivatives of the integral
are controlled by \eqref{hb:heat-L1}, since the power $1/2$ is
integrable.  For a grid translation of length $a\le1$, the
first-derivative kernel difference is bounded in $L^1$ by
\begin{equation}
 C\min\{(t+h^2)^{-1/2},\ a(t+h^2)^{-1}\}.
 \label{hb:first-kernel-increment}
\end{equation}
Use the first estimate directly and the second by telescoping one
more difference.  Splitting the time integral at $t=a^2$ gives
$Ca(1+\log(1/a))\le C_\alpha a^\alpha$.
This proves \eqref{hb:heat-X-bound}.

For two derivatives use the cancellation of $D^2p_t$ against a
constant.  Subtract $g(x)$ inside its convolution.  The spatial
$\alpha$ moment of the differentiated heat bound gives
\begin{equation}
 \|D^2P_tg\|_\infty
       \le C_\alpha(t+h^2)^{\alpha/2-1}\|g\|_{Z_\alpha}.
 \label{hb:second-derivative-cancellation}
\end{equation}
This is integrable because $\alpha>0$. A translated difference is
bounded by twice that expression. When $a\le\sqrt{t+h^2}$ one
can instead use
$Ca(t+h^2)^{(\alpha-3)/2}\|g\|_{Z_\alpha}$: telescope an extra
kernel difference and use its cancellation, subtracting the forcing
at a fixed endpoint. The translation is within the diffusion length,
so its moment has the same scale. Thus the minimum of these two
bounds is valid. Splitting again at $a^2$
gives $C_\alpha a^\alpha$.  This proves the last seminorm in
\eqref{hb:heat-Y-bound}.  The other terms were already bounded.

We specify why a small exponential weight causes no mesh loss.
For a constant vector $s$, conjugation by $e^{s\cdot x}$ replaces
the symbol of $L$ by its value at $\xi+is$.  Its real part is
bounded below by $c|\xi|^2-C|s|^2$ on the Brillouin cube.
The same differentiated Fourier calculation as for the heat tails,
now on the shifted contour, gives
\begin{equation}
 \|e^{s\cdot x}D^jp_t(x)\|_1
       \le C_j(t+h^2)^{-j/2}\quad(0<t\le1,\ |s|\le s_0),
 \label{hb:tilted-heat-bound}
\end{equation}
with the analogous fixed-time cell-$L^2$ and increment bounds.
One can see the uniformity directly from
$\operatorname{Re}\lambda_h(\xi+is)
=2h^{-2}\sum_\mu\{1-\cosh(hs_\mu)\cos(h\xi_\mu)\}$:
the negative part is at most $C|s|^2$, and every derivative of the
remaining multiplier has the same scaled Gaussian bound.
The estimate for $e^{s|x|_1}$ follows by summing the sixteen sign
choices.  A one-Lipschitz spatial weight has ratio at two points
at most $e^{s|x-y|_1}$, so the same convolution estimates apply.
Local H\"older differences read weights at neighboring cells, whose
ratios are bounded.  This proves the weighted assertion.
\end{proof}

\subsubsection{Strong smoothing by the actual site projector}

Put
\begin{equation}
 G_s=(L+q^\sharp q)^{-1},\quad F=G_sq^\sharp,
 \quad M_s=F^\sharp F,\quad P=FM_s^{-1}F^\sharp,
 \quad R=1-P.
 \label{hb:site-projector}
\end{equation}
The proved site/Gram energy package gives exponential physical
cell-$L^2$ entries of $G_s,F,DF,M_s^{-1}$.  These statements
are applicable here before any vector inverse is constructed.

For coarse $b\in\ell^\infty$, weighted cell summation gives
$\|Fb\|_{2,\rm cell,\infty}+\|qFb\|_\infty\le C\|b\|_\infty$.
The actual scalar equation is
\begin{equation}
                         L(Fb)=q^\sharp(b-qFb).
 \label{hb:F-scalar-equation}
\end{equation}
Here $q^\sharp$ extends a coarse colour as a constant in its
cell in the parallel frame, so its supremum norm is unchanged.
Thus \cref{lem:hb-heat-strong} proves
\begin{equation}
 \|Fb\|_{X_\alpha}\le C_\alpha\|b\|_\infty,
 \qquad \|DFb\|_{Z_\alpha}\le C_\alpha\|b\|_\infty.
 \label{hb:F-strong}
\end{equation}
The weighted forms follow from the same scalar equation and
\eqref{hb:tilted-heat-bound}.

The adjoints $F^\sharp,(DF)^\sharp$ map fine cell-$L^2$
supremum norms to coarse $\ell^\infty$: sum their exponentially
decreasing cell operator entries.  Composing these estimates gives
\begin{equation}
 \begin{split}
 \|Pu\|_{X_\alpha}&\le C_\alpha\|u\|_{2,\rm cell,\infty},\\
 \|D P D^\sharp u\|_\infty
             &\le C\|u\|_{2,\rm cell,\infty}.
 \end{split}
 \label{hb:projector-strong}
\end{equation}
In the second line use the exact factorization
$DPD^\sharp=(DF)M_s^{-1}(DF)^\sharp$.
No absolute pointwise kernel bound was inferred from a cell-$L^2$
bound.  Strong estimates were obtained only after the scalar heat
regularization.  In particular $R$ is bounded on $Z_\alpha$.

If $\rho\in\operatorname{ran}R\cap Z_\alpha$, the exact range
identities give $qG_s\rho=0$ and $LG_s\rho=\rho$.
Using the site energy bound and \eqref{hb:heat-Y-bound},
\begin{equation}
                   \|G_s\rho\|_{Y_\alpha}\le C_\alpha\|\rho\|_{Z_\alpha}.
 \label{hb:site-restoration-strong}
\end{equation}
This estimate will be used for gauge restoration, rather than an
unbounded point-evaluation map on an energy space.

\subsubsection{The full vector inverse and its multiplier}

Define the actual flat penalized operator
\begin{equation}
 H_0=\mathrm{d}^\sharp \mathrm{d}+DRD^\sharp,\qquad
                         K_0=H_0+Q^\sharp Q.
 \label{hb:flat-vector-operator}
\end{equation}
The uniform flat Landau gap proves
$\langle v,K_0v\rangle\ge c(\|Dv\|_2^2+\|v\|_2^2)$.
Conjugating by a small physical exponential weight changes this
form by at most $C|s|$ times its two $H^1$ norms.  For the
nonlocal part use the weighted factorization in
\eqref{hb:projector-strong}; for $Q$ use its fixed carrier and
physical row/column bounds.  The normalized Neumann series in
$H^1,H^{-1}$ therefore gives exponential cell-$L^2$ bounds for
$K_0^{-1}$, uniformly in $h$ and the torus.
In particular, for $u=K_0^{-1}f$,
\begin{equation}
                 \|u\|_{2,\rm cell,\infty}\le C\|f\|_\infty.
 \label{hb:vector-cell-bound}
\end{equation}

Flat Hodge theory is the exact operator identity
$\mathrm{d}^\sharp \mathrm{d}+DD^\sharp=L I$.  Hence
\begin{equation}
                     Lu=f+DPD^\sharp u-Q^\sharp Qu.
 \label{hb:vector-scalar-equation}
\end{equation}
The second term is bounded by \eqref{hb:projector-strong}.
The row $Q$ maps the cell-$L^2$ supremum norm to coarse
$\ell^\infty$.  Its transpose maps coarse $\ell^\infty$ to
fine $\ell^\infty$: in \eqref{hb:actual-flat-rows} each fine
coefficient is at most $h^4$ and each fine bond occurs in a fixed
number of coarse rows.  The physical transpose divides by $h^4$,
giving a bounded density, not $h^{-4}$ growth.
Therefore the right side of \eqref{hb:vector-scalar-equation}
is bounded by $C\|f\|_\infty$.
Apply \cref{lem:hb-heat-strong} to obtain the actual strong inverse
\begin{equation}
                       \|K_0^{-1}f\|_{X_\alpha}\le C_\alpha\|f\|_\infty.
 \label{hb:full-vector-strong-inverse}
\end{equation}
The argument with spatial weights gives its exponentially weighted
version as well.

\begin{proposition}[Uniform strong joint inverse]
\label{prop:hb-joint-strong-inverse}
For every bounded fine $f$ and coarse $b$, the system
\begin{equation}
                    H_0A-Q^\sharp\omega=f,\qquad QA=b
 \label{hb:joint-linear-system}
\end{equation}
has a unique solution, and
\begin{equation}
             \|A\|_{X_\alpha}+\|\omega\|_\infty
                       \le C_\alpha(\|f\|_\infty+\|b\|_\infty).
 \label{hb:joint-strong-bound}
\end{equation}
It satisfies the corresponding small exponentially weighted estimate.
The constants are uniform in homogeneous depth, volume and flat
holonomy.  This is an inverse for the penalized equations; removing
their gauge penalty in a nonlinear solution requires a separate Ward
argument.
\end{proposition}
\begin{proof}
Put $S=QK_0^{-1}Q^\sharp$.  It has a uniform positive lower
bound.  Indeed the actual lift $E_Q$ from
\cref{lem:hb-row-seeds} has bounded $K_0$ energy, and
\begin{equation}
 \|b\|_2^2
 =\langle K_0K_0^{-1}Q^\sharp b,E_Qb\rangle
 \le C\langle b,Sb\rangle^{1/2}\|b\|_2.
 \label{hb:Schur-lower-bound}
\end{equation}
The upper bound follows from the full gap and the bounded row norm.
The weighted vector inverse and the fixed row carrier show that
the conjugate of $S$ differs from $S$ by $O(|s|)$ in operator
norm.  Its inverse is therefore exponentially decreasing on the
coarse lattice, and $S^{-1}$ is bounded on coarse $\ell^\infty$
and its small weighted variants.

Set
\begin{equation}
 \sigma=S^{-1}(b-QK_0^{-1}f),\qquad
 A=K_0^{-1}(f+Q^\sharp\sigma),\qquad \omega=\sigma-b.
 \label{hb:joint-inverse-formula}
\end{equation}
Direct substitution gives \eqref{hb:joint-linear-system}.
The coarse bound on $\sigma$ and
\eqref{hb:full-vector-strong-inverse} give
\eqref{hb:joint-strong-bound}.  Uniqueness follows from positivity
of $S$ and $K_0$.  Every step retains the full multiplier; no
supremum bound for a field-only inverse is used in its place.
\end{proof}

There is one more strong inverse needed for the nonlinear gauge
argument.  On site functions the map
\begin{equation}
 \lambda\longmapsto(q\lambda,RL\lambda)
 :Y_\alpha\longrightarrow
     \ell^\infty(\text{coarse sites})\oplus
                         (\operatorname{ran}R\cap Z_\alpha)
 \label{hb:site-gauge-map}
\end{equation}
has the explicit inverse
\begin{equation}
 \lambda=E_qc+G_s\{\rho-RLE_qc\}.
 \label{hb:site-gauge-inverse}
\end{equation}
The braces belong to $\operatorname{ran}R$; $qG_sR=0$ and
$LG_sR=R$ prove both equations.  The bounds for $E_q$, $R$ and
\eqref{hb:site-restoration-strong} prove a uniform $Y_\alpha$
bound.  If both outputs vanish, $\lambda\in\ker q$ implies
$L\lambda\in\operatorname{ran}R$, so $L\lambda=0$ and the site
gap gives $\lambda=0$.  This establishes uniqueness in the
stated strong space, including nontrivial flat holonomy.

% END LAYERED EXPANSION background/homogeneous_strong_inverse.tex

% BEGIN LAYERED EXPANSION background/homogeneous_critical_orbit.tex
\subsection{An actual homogeneous Wilson background on a fixed coarse ball}
\label{ch:hb-critical-orbit}

The strong inverse just proved permits a construction of the nonlinear
background itself. The data are arbitrary small coarse bond logarithms
relative to a flat connection. Their smallness is measured in the
physical coarse supremum norm; neither volume nor fine depth reduces
the radius. This class includes spatially varying noncommuting coarse
connections. It is smaller than the class of all small-curvature coarse
fields on arbitrarily large tori. Proper regions and hereditary
decoupling domains are separate questions.

Keep the setting of \cref{ch:hb-strong-inverse}. Let $V_0$ be a real
flat coarse connection and choose a flat fine extension $W_0$ with
$M^kW_0=V_0$. To obtain one, choose commuting logarithms of its four
holonomies, use their constant fine representatives, and extend the
coarse node gauge to the fine nodes. Gauge covariance of the actual
mean, and its value on identical transported paths, give the equality.
Use the parallel frame of $W_0$ and put
\begin{equation}
 U_A=e^{ihA}W_0,\qquad V_B=e^{iB}V_0,\qquad
 \mathscr A_h(A)=\sum_{p>0}\left(1-\tfrac12\operatorname{tr}U_A(p)\right).
 \label{hb:actual-action-data}
\end{equation}
The sum uses each positive unoriented plaquette once. The colour
pairing is $\langle X,Y\rangle=\tfrac12\operatorname{tr}(XY)$.
Thus the flat quadratic action is $\tfrac12\|\mathrm{d}A\|_2^2$ with fine
measure $h^4$. Complex parameters mean holomorphic continuation to
$\operatorname{SL}(2,\mathbb C)$, without inserting a real-part operation.

\begin{proposition}[A depth-uniform homogeneous critical orbit]
\label{prop:hb-actual-critical-orbit}
Fix $L$ and $0<\alpha<1$. There are $r,\rho,\delta,C>0$, independent
of depth, volume and flat holonomies, such that for
$\|B\|_\infty<\rho$ there is a holomorphic pair $(A(B),\omega(B))$
with
\begin{equation}
 \begin{split}
 q_h(A(B))&=B,\qquad RD^\sharp A(B)=0,\\
 \mathscr A_h'(A(B))&=q_h'(A(B))^\sharp\omega(B),\\
 \|A(B)\|_{X_\alpha}+\|\omega(B)\|_\infty&\le C\|B\|_\infty.
 \end{split}
 \label{hb:critical-pair-conclusion}
\end{equation}
The actual straightened nonlinear average $q_h$ is defined below;
$R,D$ are the fixed flat site projector and gradient. The pair is
real for real $B$. A specified holomorphic fine gauge converts $U_A$
into $\widehat U(B)$ satisfying the literal constraint
$M^k\widehat U(B)=V_B$. For real data this is a critical point of the
actual Wilson action on that constraint manifold.

Among node-fixed gauge orbits possessing a representative
$U=e^{ihA}W_0$ with $\|A\|_{X_\alpha}<\delta$, its orbit is the
only constrained-critical orbit. The representative with $q_h(A)=B$
and $RD^\sharp A=0$ has a strict local action minimum on that
simultaneous slice. The literal representative is a local minimum
modulo node-fixed gauges. All maps have fixed absolute complex radii
inside the actual logarithm charts and exponentially localized
response bounds.
\end{proposition}

The fixed reference gauge distinguishes the orbit during construction.
It is not identified with a background-dependent Landau gauge. Its
penalty will be removed by the exact nonlinear Ward identity.

\subsubsection{The actual nonlinear row and its physical transpose}

On a one-step lattice with spacing set to one, let $T_{yx}(W)$ be
the specified internal mean transport and put
\begin{equation}
 \begin{split}
 u_W(a)(y)&=\exp\left[L^{-4}\sum_{x\in B_L(y)}
       \log\{T_{yx}(e^{ia}W)T_{yx}(W)^{-1}\}\right],\\
 F_W(a)(c)&=\frac1i\log\left[u_W(a)(c_-)^{-1}
       M(e^{ia}W)(c)u_W(a)(c_+)M(W)(c)^{-1}\right].
 \end{split}
 \label{hb:one-step-actual-map}
\end{equation}
The transports use the actual Federbush mean. The auxiliary $u_W$
is the explicit exponential of averaged relative logarithms. Set
$a_0=hA$, $v_0=1$, $W_s=M^sW_0$, and define
\begin{equation}
 a_{s+1}=F_{W_s}(a_s),\qquad
 v_{s+1}=v_s|_{\mathrm{next}}u_{W_s}(a_s),\qquad
 q_h(A)=a_k,\qquad v_h(A)=v_k.
 \label{hb:actual-recursion}
\end{equation}
At depth zero $q_1(A)=A$ and $v_1=1$. Solving the second line of
\eqref{hb:one-step-actual-map} for the plain mean, and applying its
gauge covariance inductively, proves
\begin{equation}
 M^k(U_A)=\left(e^{iq_h(A)}V_0\right)^{v_h(A)},\qquad
 U^g(c)=g(c_-)U(c)g(c_+)^{-1}.
 \label{hb:actual-plain-factorization}
\end{equation}
In particular the order of the two factors $v_su_s$ is fixed.

All $W_s$ are one in the parallel frame. The finite analytic
coefficient estimates proved for this actual recursion give
\begin{equation}
 \|D_aF_1(a)\|_{\mathrm{row}}\le L(1+C_L\|a\|_\infty),\qquad
 \|D_aF_1(a)\|_{\mathrm{col}}\le L^{-3}(1+C_L\|a\|_\infty).
 \label{hb:one-step-bounds}
\end{equation}
These are counting-coordinate sums of colour matrix norms. At zero
they are the exact straight-row values. The analytic chart is
$\|a\|_\infty<c_L/L$. Put $\xi_s=hL^s$. A first-exit argument gives
\begin{equation}
 \|a_s\|_\infty\le2r\xi_s\quad(0\le s\le k),\qquad
                     \sum_{s<k}\xi_s\le(L-1)^{-1}.
 \label{hb:uniform-raw-domain}
\end{equation}
Indeed $p_s=\|a_s\|_\infty/\xi_s\le2r$ implies
$p_{s+1}\le p_s(1+C_L\xi_sp_s)$. Choose
$2C_Lr/(L-1)<\log2$ and multiply. The first exit contradicts this
product. Decrease $r$ to keep every raw input strictly inside its
chart. Since $\|u_{W_s}(a_s)-1\|\le C_L\|a_s\|_\infty$, also
\begin{equation}
 \|v_h(A)-1\|_\infty\le C_Lr,\qquad
 \sup_y\|v_h(A)(y)\|\|v_h(A)(y)^{-1}\|\le e^{C_Lr}.
 \label{hb:v-domain}
\end{equation}
These are complex-ball bounds; unitary reference twists preserve them.

The actual row calculation and its fixed-input-leg Hessian bound
give on this entire fine supremum ball
\begin{equation}
 \begin{split}
 q_h(0)&=0,\qquad q_h'(0)=Q,\\
 \|q_h'(A)-Q\|_{\mathrm{row}}&\le C_L\|A\|_\infty,\qquad
 \|q_h'(A)-Q\|_{\mathrm{col}}\le C_Lh^4\|A\|_\infty,\\
 \sup_b\sum_{b'}\|\partial_b\partial_{b'}q_{h,c}(A)\|
                                                   &\le C_Lh^4.
 \end{split}
 \label{hb:actual-constraint-moments}
\end{equation}
Every fine input affects a fixed number of coarse row carriers.
The last line is \eqref{px:constraint-hessian-leg-moment} at output
scale one. Its fine density follows from one inner column sum,
one inner row sum and the geometric sum over the differentiated
scale. Integrate that bound from zero to $A$ for the column error;
sum over the $O(h^{-4})$ row inputs for the row error. A bound on
second directional derivatives alone would not justify this step.

Writing $C_h(A)=q_h(A)-QA$, one obtains
\begin{equation}
 \|C_h(A)\|_\infty\le C_L\|A\|_\infty^2,\qquad
 \|(q_h'(A)-Q)^\sharp\omega\|_\infty
                         \le C_L\|A\|_\infty\|\omega\|_\infty.
 \label{hb:nonlinear-row-bounds}
\end{equation}
The fine component of the transpose is
$h^{-4}\sum_c(q_h'(A)-Q)_{c,b}^{\mathsf T}\omega(c)$.
The column density in \eqref{hb:actual-constraint-moments} cancels
this $h^{-4}$. The same moment gives, on the joint radius-$r$ ball,
\begin{equation}
 \begin{split}
 \|C_h(A)-C_h(\widetilde A)\|_\infty
     &\le C_Lr\|A-\widetilde A\|_\infty,\\
 \|(q_h'(A)-Q)^\sharp\omega
 -(q_h'(\widetilde A)-Q)^\sharp\widetilde\omega\|_\infty
     &\le C_Lr(\|A-\widetilde A\|_\infty
                              +\|\omega-\widetilde\omega\|_\infty).
 \end{split}
 \label{hb:nonlinear-row-Lipschitz}
\end{equation}
These estimates use the full finite-lattice constraint.

\subsubsection{The complete Wilson residual}

Let $J_h(A)$ be the normal-coordinate physical gradient:
$\mathscr A_h'(A)[a]=h^4\sum_b\langle J_h(A)(b),a(b)\rangle$.
The field equation needs this entire gradient, including its
coordinate differential, rather than an action remainder alone.

\begin{lemma}[Uniform quadratic Wilson residual]
\label{lem:hb-Wilson-residual}
The map $N_h(A)=J_h(A)-\mathrm{d}^\sharp \mathrm{d}A$ is holomorphic from a fixed
$X_\alpha$ ball to fine $\ell^\infty$. On a smaller ball,
\begin{equation}
 \|N_h(A)\|_\infty\le C\|A\|_{X_\alpha}^2,\qquad
 \|N_h(A)-N_h(\widetilde A)\|_\infty
                       \le Cr\|A-\widetilde A\|_{X_\alpha}
 \label{hb:Wilson-residual-bound}
\end{equation}
when both arguments have norm at most $r$. The constants are uniform
in mesh and volume.
\end{lemma}
\begin{proof}
For $U_A(p)=e^{ih^2F_p(A)}$, the exact four-factor logarithm gives
\begin{equation}
 F(A)=\mathrm{d}A+E_h(A),\qquad
 \|E_h(A)\|_\infty\le C\|A\|_\infty^2,\qquad
 \|DE_h(A)\|_\infty\le C\|A\|_\infty\|DA\|_\infty.
 \label{hb:plaquette-quadratic}
\end{equation}
The linear logarithm is $h^2\mathrm{d}A$. Each remaining monomial contains
at least two factors $hA$, so division by $h^2$ leaves a uniformly
analytic function of the four neighbouring $A$ values. Its derivative
is bounded by $C\max|A|$. Subtract the same function at neighbouring
plaquettes to prove the last bound. No $D^2A$ is used.

The exact physical left current is
\eqref{bg:exact-wilson-current}. Its covariant difference differs
from the ordinary difference by at most
$C\|A\|_\infty\|F\|_\infty$. Its sine factor obeys
\begin{equation}
 W_h(F)=h^{-2}\sin(h^2F)=F+S_h(F),\quad
 |S_h(F)|\le Ch^4|F|^3,\quad
                 |D^-S_h(F)|\le Ch^3\|F\|_\infty^3.
 \label{hb:sine-remainder}
\end{equation}
The last estimate uses two bounded values, not a derivative of $F$.
It follows that
\begin{equation}
 J_h^{\mathrm{left}}(A)=\mathrm{d}^\sharp \mathrm{d}A+R_h(A),\qquad
 \|R_h(A)\|_\infty\le C\bigl(\|A\|_\infty\|DA\|_\infty
               +\|A\|_\infty^3+h^3\|F\|_\infty^3\bigr)
                       \le C\|A\|_{X_\alpha}^2.
 \label{hb:left-current-residual}
\end{equation}
No zero-divergence assumption is required: the ordinary divergence
of the linear plaquette curl is $\mathrm{d}^\sharp \mathrm{d}A$ itself.

The exact change from normal to left coordinates is
\begin{equation}
 \mathcal E_{hA}a=\int_0^1\operatorname{Ad}(e^{ithA})a\,dt,\quad
 J_h(A)=\mathcal E_{hA}^\sharp J_h^{\mathrm{left}}(A),\quad
                  \|\mathcal E_{hA}-1\|\le Ch\|A\|_\infty.
 \label{hb:normal-current-differential}
\end{equation}
Although $\mathrm{d}^\sharp \mathrm{d}$ is not uniformly bounded from $X_\alpha$
to $\ell^\infty$, the elementary fine difference bound gives
$\|\mathrm{d}^\sharp \mathrm{d}A\|_\infty\le Ch^{-1}\|DA\|_\infty$. The explicit
$h$ in the coordinate differential cancels that loss. This proves
the first bound in \eqref{hb:Wilson-residual-bound} for the complete
gradient. The same holomorphic estimates hold on a larger fixed
complex ball. Banach-valued Cauchy estimates bound its second
derivative uniformly on a smaller ball. Integrating that derivative,
using $N_h(0)=DN_h(0)=0$, proves the stated $Cr$ Lipschitz estimate.
No coordinate-count factor enters this operator-norm argument.
\end{proof}

\subsubsection{The actual joint field--multiplier contraction}

Let $\mathcal L_0^{-1}(f,b)$ be the inverse in
\cref{prop:hb-joint-strong-inverse}, with norm at most $C_0$.
The actual penalized system
\begin{equation}
 J_h(A)+DRD^\sharp A-q_h'(A)^\sharp\omega=0,\qquad q_h(A)=B
 \label{hb:actual-penalized-equations}
\end{equation}
is exactly equivalent to
\begin{equation}
 (A,\omega)=\mathcal T_B(A,\omega)
 :=\mathcal L_0^{-1}\bigl(-N_h(A)+(q_h'(A)-Q)^\sharp\omega,
                                             B-C_h(A)\bigr).
 \label{hb:actual-fixed-point}
\end{equation}
For $\|(A,\omega)\|=\|A\|_{X_\alpha}+\|\omega\|_\infty\le r$,
the preceding estimates imply
\begin{equation}
 \|\mathcal T_B(x)-\mathcal T_B(\widetilde x)\|
       \le C_1r\|x-\widetilde x\|,\qquad
                       \|\mathcal T_B(0)\|\le C_0\|B\|_\infty.
 \label{hb:contraction-constants}
\end{equation}
Choose $r$ below all primitive ceilings and $C_1r\le1/4$, and
$\rho\le r/(4C_0)$. The radius-$r/2$ ball is invariant and
contractive. Iteration from zero gives a unique pair there. The
same Lipschitz inequality identifies it with any solution in the
radius-$r$ ball, and
\begin{equation}
 \|A(B)\|_{X_\alpha}+\|\omega(B)\|_\infty
                                     \le2C_0\|B\|_\infty.
 \label{hb:actual-fixed-point-bound}
\end{equation}
The holomorphic iterations converge locally uniformly on the fixed
complex polydisc. Their limit is holomorphic; reality and uniqueness
give reality for real $B$. An unknown orbit or inverse is not a
hypothesis of this construction.

\subsubsection{The effective coarse gauge and removal of the penalty}

The nonlinear straightened average has its own effective coarse gauge.
One cannot obtain its Ward identity by replacing the flat site map
by a convenient average. For $g=e^{-it\lambda}$ define
\begin{equation}
 h_A(g)=v_h(A^g)^{-1}g|_{\mathrm{coarse}}v_h(A),\qquad
 h_A(e^{-it\lambda})=1-it\tau_A\lambda+O(t^2),
 \label{hb:effective-coarse-gauge}
\end{equation}
where $U_{A^g}=U_A^g$. The exact factorization proves
\begin{equation}
 e^{iq_h(A^g)}V_0=(e^{iq_h(A)}V_0)^{h_A(g)},\qquad
 q_h'(A)\Gamma_A\lambda=\Gamma^c_{q_h(A)}\tau_A\lambda.
 \label{hb:exact-nonlinear-intertwining}
\end{equation}
The symbols $\Gamma_A,\Gamma_B^c$ are normal-coordinate infinitesimal
fine and coarse gauge generators, with the same sign convention.

We derive estimates for these actual maps. Write
$\mathcal E_zv=\int_0^1\operatorname{Ad}(e^{itz})v\,dt$.
At one step in the flat frame put $T_{yx}(e^{ia})=e^{iz_x}$,
$m=L^{-4}\sum_xz_x$ and $u=e^{im}$. Direct differentiation gives
\begin{equation}
 \begin{split}
 \tau_a^{(1)}\lambda(y)
 =\operatorname{Ad}(u(y)^{-1})\biggl[\lambda(y)
 +\mathcal E_m L^{-4}\sum_{x\in B_L(y)}\mathcal E_{z_x}^{-1}
 \{-\lambda(y)+\operatorname{Ad}(T_{yx}(e^{ia}))\lambda(x)\}\biggr].
 \end{split}
 \label{hb:one-step-effective-gauge}
\end{equation}
Indeed the left logarithmic derivative of
$T_{yx}(U^g)=g(y)T_{yx}(U)g(x)^{-1}$ is
$i[-\lambda(y)+\operatorname{Ad}(T_{yx}(U))\lambda(x)]$.
Convert it by $\mathcal E_{z_x}^{-1}$, average, and differentiate
$u(A^g)^{-1}g(y)u(A)$. This gives the formula, including its sign.
At zero it is exactly the one-step site average $q^{(1)}\lambda$.
The finite one-step logarithm domain therefore gives
\begin{equation}
 \|\tau_a^{(1)}-q^{(1)}\|_{\infty\to\infty}\le C_L\|a\|_\infty,
 \qquad \|\tau_a^{(1)}\|_{\infty\to\infty}
                                     \le1+C_L\|a\|_\infty.
 \label{hb:one-step-tau-bound}
\end{equation}
The effective gauge parameters compose through successive scales.
Telescoping that product and using \eqref{hb:uniform-raw-domain} yields
\begin{equation}
 \|\tau_A-q\|_{\infty\to\infty}\le C_L\|A\|_\infty.
 \label{hb:full-tau-bound}
\end{equation}
This retains the node evaluation terms in
\eqref{hb:one-step-effective-gauge}; no fine $L^2$ point-evaluation
bound is inferred from it.

The fine generator has the exact expression
\begin{equation}
 \begin{split}
 \Gamma_A\lambda&=\mathcal E_{hA}^{-1}D_{U_A}\lambda,\\
 (D_{U_A}\lambda)_\mu(x)
   &=h^{-1}\{\operatorname{Ad}(U_A(x,\mu))\lambda(x+he_\mu)-\lambda(x)\}.
 \end{split}
 \label{hb:actual-fine-generator}
\end{equation}
It implies
\begin{equation}
 \|(\Gamma_A-D)\lambda\|_{X_\alpha}
                     \le C\|A\|_{X_\alpha}\|\lambda\|_{Y_\alpha}.
 \label{hb:fine-generator-bound}
\end{equation}
Its zeroth-order difference consists of $A\lambda$ and $hA D\lambda$
terms. One fine difference gives terms of the forms $DA\lambda$
and $A D\lambda$, together with $hDA D\lambda$ and $hA D^2\lambda$.
Their analytic coefficients are uniformly bounded.
The discrete product rule for the local H\"older
seminorm proves the stated last seminorm as well. The inverse of
$\mathcal E_{hA}$ is uniformly analytic on the chosen chart.

For $\rho\in\operatorname{ran}R\cap Z_\alpha$ put
$\lambda_0=G_s\rho$. The strong site inverse gives
$q\lambda_0=0$, $RL\lambda_0=\rho$ and
$\|\lambda_0\|_{Y_\alpha}\le C\|\rho\|_{Z_\alpha}$.
Since $\tau_AE_q=1+O(r)$, define
\begin{equation}
 \lambda_A(\rho)=\lambda_0
                   -E_q(\tau_AE_q)^{-1}\tau_A\lambda_0.
 \label{hb:effective-kernel-lift}
\end{equation}
Then $\tau_A\lambda_A=0$ and
$\|\lambda_A-\lambda_0\|_{Y_\alpha}\le Cr\|\rho\|_{Z_\alpha}$.
Boundedness of $R$ on $Z_\alpha$ and
\eqref{hb:fine-generator-bound} show that
\begin{equation}
 \mathscr R_A\rho=RD^\sharp\Gamma_A\lambda_A(\rho),\qquad
 \|\mathscr R_A-1\|_{\operatorname{ran}R\cap Z_\alpha}\le Cr.
 \label{hb:actual-gauge-restoration-map}
\end{equation}
After decreasing $r$ this is invertible by a Neumann series.
Consequently for every such $\rho$ there is a site function with
\begin{equation}
             \tau_A\lambda=0,\qquad RD^\sharp\Gamma_A\lambda=\rho.
 \label{hb:exact-gauge-test}
\end{equation}

For the real fixed point choose $\rho=RD^\sharp A$. Test
\eqref{hb:actual-penalized-equations} against the actual gauge vector
$\Gamma_A\lambda$ from \eqref{hb:exact-gauge-test}.
Wilson gauge invariance annihilates its current term. The constraint
term vanishes by \eqref{hb:exact-nonlinear-intertwining} and
$\tau_A\lambda=0$. The remaining equality is
\begin{equation}
                         0=\langle\rho,\rho\rangle_{2,h}.
 \label{hb:penalty-vanishing}
\end{equation}
Thus $RD^\sharp A=0$ and the full unpenalized stationarity equation
holds. The identity extends to the complex data ball by holomorphy
from its real open part. This proves the equations in
\eqref{hb:critical-pair-conclusion}.

The same argument with arbitrary $\lambda$ gives
\begin{equation}
                     (\Gamma_B^c)^\sharp\omega=0.
 \label{hb:coarse-multiplier-Ward}
\end{equation}
Indeed $\tau_A$ is onto with right inverse
$E_q(\tau_AE_q)^{-1}$. The coarse frame term arising on
differentiation of \eqref{hb:actual-plain-factorization} therefore
contributes no extra physical force.

\subsubsection{A literal trial lift and admission of its fine gauge}

There is an explicit bounded extension of coarse node values to
fine site functions. Choose fixed scalar smooth bumps supported
inside the pillars and equal to one at their centers, transporting
colours in the flat frame. Their supports are disjoint. Sample them
on the fine lattice. At the finitely many coarse meshes use center
Kronecker bumps, whose physical derivatives have a fixed bound.
This gives
\begin{equation}
 \operatorname{ev}_{\mathrm{coarse}}E_n=1,\qquad
                        \|E_n\nu\|_{Y_\alpha}\le C\|\nu\|_\infty.
 \label{hb:node-extension}
\end{equation}
For $w=e^{i\nu}$ close to one define $\mathfrak E(w)=e^{iE_n\nu}$.
It is unitary for real $\nu$, holomorphic for complex $\nu$, and has
exactly the prescribed node values.

The finite fine gauge transformation satisfies
\begin{equation}
 \begin{split}
 A^{e^{-i\lambda}}&=A+D\lambda+\mathcal B_h(A,\lambda),\\
 \|\mathcal B_h(A,\lambda)\|_{X_\alpha}
   &\le C(\|A\|_{X_\alpha}+\|\lambda\|_{Y_\alpha})\|\lambda\|_{Y_\alpha}.
 \end{split}
 \label{hb:finite-gauge-BCH}
\end{equation}
The dependence is uniformly analytic. To check the apparent $h^{-1}$,
write $\lambda(x+he_\mu)=\lambda(x)+hD_\mu\lambda(x)$ before
taking the product and logarithm. At $h=0$ the numerator is exactly
one. Its logarithm divided by $h$ therefore extends analytically
to $h=0$. Differences of this finite formula use only
$A,DA,\lambda,D\lambda,D^2\lambda$ and their stated local seminorms.
Its first-order terms are $A+D\lambda$. Taylor's formula proves
\eqref{hb:finite-gauge-BCH} and the corresponding local Lipschitz
bound, uniformly in the mesh.

For any small $A$ put
\begin{equation}
                     \widehat U_A=U_A^{\mathfrak E(v_h(A)^{-1})}.
 \label{hb:literal-representative}
\end{equation}
Its literal mean is $e^{iq_h(A)}V_0$. Its normal coordinate has
$X_\alpha$ norm at most $C\|A\|_{X_\alpha}$ by
\eqref{hb:v-domain}, \eqref{hb:node-extension} and
\eqref{hb:finite-gauge-BCH}. Choose the fixed-point radius below
the primitive admission ceiling divided by this fixed constant.
Then the original field, its literal representative, the specified
fine gauges and all intermediate raw logarithms stay inside the
actual complex charts. This proves admission for this gauge without
assuming arbitrary complex gauge saturation of the domain.

A literal section can also be constructed independently of the
action. On the range of the actual smooth seed $E_Q$, solve
\begin{equation}
                       A_{\mathrm{tr}}=E_Q\{B-C_h(A_{\mathrm{tr}})\}.
 \label{hb:actual-trial-section}
\end{equation}
The seed bound and \eqref{hb:nonlinear-row-Lipschitz} make this a
supremum-norm contraction. Its solution has $q_h(A_{\mathrm{tr}})=B$.
Since it is $E_Q$ applied to a bounded coarse vector, its whole
$Y_\alpha$ norm is at most $C\|B\|_\infty$. Equation
\eqref{hb:literal-representative} gives an actual literal trial
lift. Its physical curvature is at most $C\|B\|_\infty$.
Before the last gauge its physical current has the same bound by
the exact current formula and the $Y_\alpha$ estimate. Gauge
covariance preserves this bound for real data; the bounded condition
number of the specified complex gauge gives it in the complex case.
This supplies a genuine deep trial field with physical derivative
bounds, rather than assuming the required trial lift.

\subsubsection{The mean-gauge slice and physical criticality}

For $\lambda$ small in $Y_\alpha$ define
\begin{equation}
 \mathscr S_A(\lambda)=
    \bigl(i\log h_A(e^{-i\lambda}),\ RD^\sharp A^{e^{-i\lambda}}\bigr).
 \label{hb:nonlinear-mean-slice-map}
\end{equation}
At $(A,\lambda)=(0,0)$ its derivative is $(q\lambda,RL\lambda)$,
with explicit bounded inverse \eqref{hb:site-gauge-inverse}.
On a small $X_\alpha\times Y_\alpha$ ball that derivative differs
from this fixed map by at most
$C(\|A\|_{X_\alpha}+\|\lambda\|_{Y_\alpha})$.
For the first component differentiate the exact $h_A$ at the current
point: its derivative is $\tau_{A^g}$ composed with the bounded
near-identity exponential gauge differential and a near-identity
conjugation. Use \eqref{hb:full-tau-bound}. For the second use
$RD^\sharp:X_\alpha\to Z_\alpha$ and
\eqref{hb:finite-gauge-BCH}. These bounds give the claimed operator
norm between the precise strong spaces.

The contraction with linear part \eqref{hb:site-gauge-inverse}
therefore solves
\begin{equation}
 h_A(e^{-i\lambda})=v_h(A),\qquad RD^\sharp A^{e^{-i\lambda}}=0,
 \qquad\|\lambda\|_{Y_\alpha}\le C\|A\|_{X_\alpha}.
 \label{hb:mean-slice-solution}
\end{equation}
Explicitly subtract $\mathscr S_A(0)$ and move the quadratic
remainder to the right after applying that inverse. The Lipschitz
constant is $C(\|A\|_{X_\alpha}+\|\lambda\|_{Y_\alpha})$.
A ball of radius $C'\|A\|_{X_\alpha}$ is invariant and contractive
for a fixed smaller input radius. This proves existence, analytic
dependence and uniqueness in this gauge chart.

If $M^kU_A=V_B$, write $g=e^{-i\lambda}$ for this solution and
$\overline A=A^g$. Exact factorization gives
\begin{equation}
 q_h(\overline A)=B,\qquad RD^\sharp\overline A=0,\qquad
                           v_h(\overline A)=g|_{\mathrm{coarse}}.
 \label{hb:literal-to-mean-slice}
\end{equation}
The last equality follows from
$v_h(\overline A)^{-1}g|_{\mathrm{coarse}}v_h(A)=v_h(A)$.
No coarse stabilizer is divided out. Consequently
$\widehat U_{\overline A}$ differs from $U_A$ by a gauge exactly
equal to one at every coarse node.

These maps establish the critical-point correspondence with all
tangent directions included. The map $A\mapsto\widehat U_A$ sends
the entire local $q_h=B$ manifold into the literal fibre and preserves
the action. Conversely \eqref{hb:mean-slice-solution} sends the local
literal fibre into $q_h=B$ and preserves the action. At a mean-slice
point the two maps compose to the identity up to a node-fixed gauge:
the inverse of its specified extension gauge solves
\eqref{hb:mean-slice-solution}, and uniqueness in the small gauge
chart identifies it. Differentiating these action identities proves
both directions of the critical-point correspondence. In particular
the constructed fixed point is a literal constrained-critical orbit.

\subsubsection{Competing multipliers and orbit uniqueness}

Let $A$ be any small mean-slice critical point. Its actual derivative
is onto, with the exact right inverse
\begin{equation}
 T_A=q_h'(A)E_Q=1+O(r),\qquad E_A=E_QT_A^{-1},\qquad q_h'(A)E_A=1.
 \label{hb:actual-nonlinear-row-lift}
\end{equation}
The error has coarse row and column sums at most $Cr$, by
\eqref{hb:actual-constraint-moments} and the seed carriers. Thus its
Neumann inverse is bounded in both Schur norms. Columns of $E_A$
have uniformly bounded physical $W^{1,1}$ norms and $H^1$ energy:
sum the column $\ell^1$ bound of $T_A^{-1}$ against the corresponding
fixed seed-column bounds.

Since $q_h'(A)$ is onto, constrained criticality gives a unique
$\omega$ with $J_h(A)=q_h'(A)^\sharp\omega$. To bound it without
an assumed $D^2A$ estimate, evaluate the first variation on $E_Ae_c$.
The exact current identity, summed by parts once, gives
\begin{equation}
 |\mathscr A_h'(A)[a]|
 \le C\|F(A)\|_\infty\bigl(
       \|D(\mathcal E_{hA}a)\|_{1,h}
             +\|A\|_\infty\|\mathcal E_{hA}a\|_{1,h}\bigr).
 \label{hb:first-variation-readout}
\end{equation}
The measure is $h^4$. This is summation by parts of the covariant
difference of $W_h(F)$, whose supremum is at most $C\|F\|_\infty$.
Its transport differs from the flat one by $Ch\|A\|_\infty$.
Moreover $\|\mathcal E_{hA}\|\le C$ and
$\|D\mathcal E_{hA}\|\le Ch\|DA\|_\infty$. The seed column bounds
and \eqref{hb:plaquette-quadratic} therefore imply
\begin{equation}
                               \|\omega\|_\infty\le C\|A\|_{X_\alpha}.
 \label{hb:competing-multiplier-bound}
\end{equation}
There is no fine point-mass readout or hidden $h^{-1}$ here.

A literal critical connection in the radius-$\delta$ chart maps
to a mean-slice critical point of norm at most $C\delta$.
Choose $\delta$ so that its field and the multiplier in
\eqref{hb:competing-multiplier-bound} have joint norm less than $r$.
Decrease $\rho$ further so that the literal constructed representative
has strong norm below $\delta/2$.
They solve \eqref{hb:actual-fixed-point}, since the fixed gauge
penalty vanishes. Contraction uniqueness identifies them with
$(A(B),\omega(B))$. Equation \eqref{hb:literal-to-mean-slice}
then identifies the original connection with $\widehat U(B)$ modulo
a node-fixed gauge. This proves the stated orbit uniqueness.

\subsubsection{A real minimum and the effective coarse action}

Keep $\omega=\omega(B)$ fixed and set
$\mathscr L(A)=\mathscr A_h(A)-\langle\omega,q_h(A)-B\rangle$.
On a fixed small strong ball,
\begin{equation}
 \mathscr L''(A)[a,a]+\|RD^\sharp a\|_2^2+\|q_h'(A)a\|_2^2
                          \ge c(\|Da\|_2^2+\|a\|_2^2).
 \label{hb:fixed-gauge-extended-coercivity}
\end{equation}
Here is the comparison with the proved flat gap. The exact
twice-varied plaquette estimate \eqref{slp:wilson-comparison} uses
physical left variations. Normal variations are $\mathcal E_{hA}a$;
their difference from $a$ has coefficient $O(h\|A\|_\infty)$ and
first-difference coefficients $O(h\|DA\|_\infty)$ and
$O(h\|A\|_\infty)D$. The change in the curl form is at most
$Cr\|a\|_{H^1}^2$. The second coordinate differential contributes
a local coefficient $O(h)J_h^{\mathrm{left}}(A)$, bounded by $Cr$:
the exact current formula gives
$h\|J_h^{\mathrm{left}}(A)\|_\infty
 \le C\|DA\|_\infty+Ch\|A\|_{X_\alpha}^2$.
The two-leg moment bounds the constraint curvature by
$C\|\omega\|_\infty\|a\|_2^2$. Finally physical Schur applied to
\eqref{hb:actual-constraint-moments} gives
$\|q_h'(A)-Q\|_{2\to2}\le Cr$. Thus the entire left side differs
from $\langle a,K_0a\rangle$ by at most
$C(r+\|\omega\|_\infty)\|a\|_{H^1}^2$, proving
\eqref{hb:fixed-gauge-extended-coercivity} after a fixed decrease.

Let $A_*=A(B)$ and let $A_*+a$ be another real mean-slice point
with the same $q_h=B$. Then $RD^\sharp a=0$. The equality of the
two constraints and the two-leg moment imply, for $A_t=A_*+ta$,
\begin{equation}
                       \|q_h'(A_t)a\|_2\le C\|a\|_\infty\|a\|_2.
 \label{hb:curved-fibre-chord}
\end{equation}
Indeed $D^2q_h(A)[a,\cdot]$ has physical $L^2$ operator norm at
most $C\|a\|_\infty$: the fixed-leg moment supplies its column
density $Ch^4\|a\|_\infty$, and summing over a row's inputs
supplies row sum $C\|a\|_\infty$. Apply physical Schur, then
integrate in $t$. Since $\mathscr L'(A_*)=0$, Taylor's formula gives
\begin{equation}
 \mathscr A_h(A_*+a)-\mathscr A_h(A_*)
                  \ge c'(\|Da\|_2^2+\|a\|_2^2)
 \label{hb:actual-local-minimum}
\end{equation}
on a smaller fixed ball. The mean-gauge slice preserves action and
is a local node-fixed orbit section of the literal fibre. This proves
the minimum assertion.

Define $S_c(V_B)=\mathscr A_h(A(B))$ and continue holomorphically.
It is locally coarse-gauge invariant. A sufficiently small coarse
gauge has the fine extension \eqref{hb:node-extension}. The transformed
critical field remains in the strong chart and is critical for the
transformed literal coarse data. Orbit uniqueness identifies it with
the new critical orbit, and fine gauge invariance identifies their
actions. Thus the coarse invariance required when subtracting
$S_c\circ M^k$ is a conclusion of this construction.

\subsubsection{Localized analytic response and the exact domain supplied}

Every carrier in the depth-$k$ recursion has bounded physical diameter:
the radii grow by a fixed multiple of $\xi_s$, and
$\sum_{s<k}\xi_s\le(L-1)^{-1}$. The nonlinear Wilson term is local
on the fine stencil. Therefore all the residual and transpose
Lipschitz bounds above remain valid with a small fixed exponential
weight, with constants multiplied by at most $e^{C_Ls}$. For the
Wilson residual use its local stencil and the same product estimates;
for $q_h$ and its transpose use their actual finite carriers and
\eqref{hb:actual-constraint-moments}. Write $w_z(x)=e^{s d(x,z)}$,
where $d$ is physical periodic distance and $s>0$ is small enough
for the weighted inverse of \cref{prop:hb-joint-strong-inverse}.
Here $\|w_zu\|_{X_\alpha}$ denotes the cellwise weighted norm of
$u$ and its indicated local derivatives and seminorms; it does not
mean differentiating the product $w_zu$.

Decrease $r$ so that the derivative of the nonlinear fixed-point
map has norm below $1/2$ in these weighted spaces as well.
Differentiating \eqref{hb:actual-fixed-point} gives the explicit
Neumann inverse of $1-D\mathcal T_B$. Thus
\begin{equation}
 \|w_zDA(B)b\|_{X_\alpha}
       +\|w_zD\omega(B)b\|_\infty\le C\|w_zb\|_\infty.
 \label{hb:localized-background-response}
\end{equation}
The constant is uniform on a smaller unweighted complex data ball,
and in $z,k,M$. The same proof shows analytic dependence on the
unweighted parameter $B$ with values in the weighted operator space.
Banach-valued Cauchy differentiation in the remaining $m-1$
directions consequently gives
\begin{equation}
 \begin{split}
 &\|w_zD^mA(B)[b_1,\ldots,b_m]\|_{X_\alpha}
   +\|w_zD^m\omega(B)[b_1,\ldots,b_m]\|_\infty\\
 &\hspace{20mm}\le C_m\|w_zb_1\|_\infty
                                  \prod_{j=2}^m\|b_j\|_\infty.
 \end{split}
 \label{hb:localized-background-jets}
\end{equation}
Any input can occupy the first slot by symmetry. For unit inputs at
coarse bonds $y_1,\ldots,y_m$, choosing $z=y_j$ for each $j$
gives decay $C_m e^{-s\max_jd(x,y_j)}$. Since
$\max_jd(x,y_j)\ge m^{-1}\sum_jd(x,y_j)$, every fixed jet has
all spatial moments, with constants uniform in depth and volume.
The literal representative has the same property: $v_h$ has uniformly
bounded analytic derivatives and bounded carriers by its explicit
formula; $E_n$ is local; and the gauge change is uniformly analytic
in the stated strong norms.

These estimates provide an actual fixed-radius homogeneous background,
its multiplier, a literal trial lift, a local minimum, local coarse
gauge invariance and localized analytic response. They do not provide
independent fine decoupling domains, a hereditary tower, mixed-scale
proper-region inversion, or extension to every curvature-only regular
coarse field. Those are separate model estimates, not conclusions
hidden in the domain of this proposition.

% END LAYERED EXPANSION background/homogeneous_critical_orbit.tex

% BEGIN LAYERED EXPANSION formation/hierarchical_mean_second_jet.tex
\subsection{The second jet and fixed-leg density of the actual hierarchy}
\label{layered:hierarchical-mean-second-jet}

The straightened nonlinear row in
\eqref{hb:actual-recursion} is obtained by composing the
actual one-step rows. This fact is essential for its
fixed-fine-leg moment. We compute the one-step second
jet before estimating the composition. A single mean
over a block of side \(L^k\) would be a different
nonlinear map and does not replace this computation.

For Hermitian logarithms write
\([X,Y]_{\rm H}=i[X,Y]\); thus the quadratic term in
\(i^{-1}\log(e^{iX}e^{iY})\) is
\(\tfrac12[X,Y]_{\rm H}\).
At one output row let \(t_x,t'_x,\ell_x\) denote the
degree-one logarithms of its source internal transport,
target internal transport, and straight length-\(L\)
line. A bar denotes the normalized average over the
\(L^4\) starting sites. Subscripts two denote
homogeneous degree-two terms. Define
\[
 {\rm Cov}_{\rm H}(a,b)
       =\overline{[a_x,b_x]_{\rm H}}
                         -[\bar a,\bar b]_{\rm H}.
\]

\begin{lemma}[The actual one-step quadratic coefficient]
\label{lem:hfm-actual-second-jet}
The quadratic coefficient of the literal plain mean is
\begin{equation}
 \begin{split}
 (M_{\log})_2={}&\bar t_2+\bar\ell_2-\bar t'_2\\
 &+\frac12\overline{
       [t,\ell]_{\rm H}-[t,t']_{\rm H}-[\ell,t']_{\rm H}}.
 \end{split}
 \label{hfm:literal-one-step-second-jet}
\end{equation}
For the actual straightened map \(F_1\) in
\eqref{hb:one-step-actual-map}, its linear coefficient
is the straight row \(Q_1=\bar\ell\), and
\begin{equation}
 (F_1)_2=\bar\ell_2+\frac12\{
  {\rm Cov}_{\rm H}(t,\ell)
  -{\rm Cov}_{\rm H}(t,t')
  -{\rm Cov}_{\rm H}(\ell,t')\}.
 \label{hfm:straightened-one-step-second-jet}
\end{equation}
In particular all second jets of the internal
transports cancel from the straightened row.
\end{lemma}

\begin{proof}
For a normalized group mean, the logarithm of its
output through degree two is the average of the
input logarithms. Indeed its defining equation is
the average of \(\log(e^{iz_x}e^{-im})=0\);
the degree-two commutator containing \(m_1=\bar z_1\)
averages to \([\bar z_1,\bar z_1]=0\).
This proves the assertion for the internal
implicit mean. The outer operation is explicitly
the exponential of the average of closed
logarithms, so its logarithm is that arithmetic
average exactly. Its second jet has the same
property without replacing it by an implicit mean.

Let \(c\) be the logarithm of the central straight
reference path. The actual closed inputs of the outer
mean are \(T_x e^{i\ell_x}(T'_x)^{-1}e^{-ic}\).
Expand their logarithms and then multiply the output
by \(e^{ic}\), as in the actual definition of \(M\).
The terms containing \(c_1\) cancel: the averaged
commutator with \(-c_1\) is cancelled by the final
product with \(e^{ic_1}\). The linear result is
\(\bar t+\bar\ell-\bar t'\), and its quadratic
coefficient is \eqref{hfm:literal-one-step-second-jet}.
This also checks that the central reference path
has not been silently discarded before taking a jet.

The straightening is
\(e^{-i\bar t_{\log}}M(e^{ia})e^{i\bar t'_{\log}}\),
where the barred logarithms include their second
coefficients. Its first coefficient is \(\bar\ell\).
The two endpoint exponentials cancel
\(\bar t_2-\bar t'_2\) exactly. Their additional
quadratic commutators are
\[
 \frac12\{-[\bar t,\bar\ell]_{\rm H}
              +[\bar t,\bar t']_{\rm H}
              +[\bar\ell,\bar t']_{\rm H}\}.
\]
Combining them with the averaged commutators gives
\eqref{hfm:straightened-one-step-second-jet}.
\end{proof}

\begin{lemma}[One-step rows and one fixed input leg]
\label{lem:hfm-one-step-leg}
On a raw fine-angle ball \(\|a\|_\infty<c/L\),
with \(c>0\) universal, the actual one-step map has
\begin{align}
 \|DF_1(a)\|_{\rm row}
                  &\le L+C L^2\|a\|_\infty,\notag\\
 \|DF_1(a)\|_{\rm col}
                  &\le L^{-3}+C L\|a\|_\infty,\notag\\
 \sup_b\sum_{c,b'}
      \|D^2F_1(a)_{c;b,b'}\|&\le C L .
 \label{hfm:actual-one-step-absolute-moments}
\end{align}
The constants are independent of \(L\ge3\), compatible
volume and the number of path-mean inputs.
\end{lemma}

\begin{proof}
Every coordinate-order path has length at most \(C L\)
and contains a fixed fine bond at most once. Its first
derivative therefore has absolute row sum \(CL\);
with one input bond fixed its norm is bounded by a
constant. Its second derivative with one fixed leg
and the other leg summed is at most \(CL\).
These statements remain valid on the indicated
complex ball, since products of complex transports
have norm at most \(e^{CL\|a\|_\infty}\).

All averaging weights are normalized. The inverse
in the implicit group-mean equation is within a
fixed neighborhood of the identity on this ball.
Implicit differentiation once and twice preserves
the preceding bounds: terms have either one
second path derivative or two first derivatives,
with the fixed leg in one of them. Products of
separate normalized averages obey the same bounds.
Only a fixed number of coarse rows can use a
given internal bond or a given crossing bond.
This proves the last line of
\eqref{hfm:actual-one-step-absolute-moments}.
Without a fixed leg, the same calculation gives
the absolute row Hessian bound \(CL^2\).

At zero the exact derivative is \(Q_1\), with row
norm \(L\) and column norm \(L^{-3}\). Integrate the
two Hessian bounds on the segment from zero to
\(a\). This gives the first two lines. The column
estimate can be written
\[
       L^{-3}(1+C L^4\|a\|_\infty).
\]
Thus uniform path-row bounds do not assert a
uniform nonlinear column constant after factoring
out \(L^{-3}\).
\end{proof}

\begin{proposition}[Hierarchical fixed-fine-leg density]
\label{prop:hfm-fixed-leg-density}
Let \(h=L^{-k}\), and let \(q_h\) be the actual
straightened hierarchy of \eqref{hb:actual-recursion},
with physical fine potential as input. There are
universal \(c,C>0\) such that, for
\(\|A\|_\infty<cL^{-3}\),
\begin{equation}
 \sup_b\sum_{c,b'}
       \|D^2q_h(A)_{c;b,b'}\|\le C L^3h^4 .
 \label{hfm:hierarchical-fixed-leg-density}
\end{equation}
Consequently
\begin{equation}
 \|q_h'(A)-Q_h\|_{\rm col}
                      \le C L^3h^4\|A\|_\infty .
 \label{hfm:hierarchical-column-error}
\end{equation}
For each fixed \(L\) these bounds are uniform in
depth. They require no projection of the multiplier
onto a divergence-free subspace.
\end{proposition}

\begin{proof}
At zero the chain rule is the exact identity
\begin{equation}
 D^2q_h(0)[u,v]
  =h^2\sum_{t=0}^{k-1}
       Q_1^{\,k-1-t}
       D^2F_1(0)[Q_1^tu,Q_1^tv].
 \label{hfm:exact-hierarchical-second-jet}
\end{equation}
The powers denote the actual linear row at its
successive grids. The factor \(h^2\) comes from
the two physical inputs \(a_0=hA\).
Write
\(B_L=\sup_b\sum_{c,b'}\|D^2F_1(0)_{c;b,b'}\|\).
Propagate the fixed inner leg by the column bound
\(L^{-3t}\), the other inner leg by the row bound
\(L^t\), and the output by the remaining column
bound \(L^{-3(k-1-t)}\). The resulting bound is
\begin{equation}
 \begin{split}
 h^2B_L\sum_{t=0}^{k-1}
       L^{-3(k-1-t)}L^{-3t}L^t
 &=h^4\,\frac{B_LL^3}{L-1}(1-h)\\
 &\le C L^3h^4 .
 \end{split}
 \label{hfm:exact-density-sum}
\end{equation}
For \(k=0\) the Hessian is zero.
The central-edge contribution in the one-step
row is included in \(B_L\); it is not estimated
as a density of order \(L^{-4}\).

At a nonzero field, the second-order chain rule
has the same sum with the actual Jacobian products
on its three branches and \(D^2F_1(a_t)\) at its
vertex. The row majorant gives
\(\|a_t\|_\infty\le2\xi_t\|A\|_\infty\), where
\(\xi_t=hL^t\). The row products, divided by
their powers of \(L\), are bounded by
\(\exp(C\|A\|_\infty\sum_tL\xi_t)\le e^{C\|A\|_\infty}\).
The analogous column products are bounded by
\[
 \exp\!\left(CL^4\|A\|_\infty\sum_{t<k}\xi_t\right)
                         \le e^{C L^3\|A\|_\infty}.
\]
This is a fixed constant on the stated ball.
The same estimate applies to every subproduct,
and the one-step fixed-leg Hessian remains \(CL\).
Thus \eqref{hfm:exact-density-sum}, multiplied by
fixed constants, proves
\eqref{hfm:hierarchical-fixed-leg-density}.
Integration in \(A\) proves
\eqref{hfm:hierarchical-column-error}.
\end{proof}

The factor \(L^3\) in this absolute estimate can be
necessary. For the symmetric coordinate-order
contours, choose different axes \(\mu,\nu\) and a
background, on the two relevant blocks, of the form
\[
          A_\mu(X)=rX_\nu\sigma_2,\qquad A_\lambda=0
                  \quad(\lambda\ne\mu).
\]
Extend it periodically with supremum at most \(Cr\);
the extension can be smooth with uniformly bounded
physical derivatives. Vary only the positive fine
\(\nu\)-bond from the source block center, with
physical variation \(\delta A=\sigma_1\).
Put \(x\in\{-(L-1)/2,\ldots,(L-1)/2\}^4\) and
\(u=x/L\). The first internal-transport variation
at that bond has weight
\[
 w_x=\frac{\mathbf1_{\{x_\nu>0\}}}{d(x)},\qquad
                 d(x)=\#\{\lambda:x_\lambda\ne0\}.
\]
Indeed the \(\nu\)-segment must be first among the
nonzero coordinate segments. The target transport
does not contain this fine bond.
The background has
\(\ell_x=r u_\nu\sigma_2\) and
\(t'_x=(r/2)u_\mu u_\nu\sigma_2\).
Their averages vanish. The covariance formula
\eqref{hfm:straightened-one-step-second-jet} gives
\begin{equation}
 D^2q_{1/L}(0)[\delta A,A]_\mu(0)
              =-\frac rL a_L\sigma_3,\qquad
 a_L=L^{-4}\sum_xw_x\frac{x_\nu}{L}
              =\frac{(L+1)(L^4-1)}{32L^5}.
 \label{hfm:central-edge-lower-bound}
\end{equation}
The \(u_\mu u_\nu\) term cancels by reflection in
\(x_\mu\). In the remaining sum there are
\(\binom3j(L-1)^j\) choices with exactly \(j\)
nonzero other coordinates. Use
\[
 \sum_{j=0}^3\frac{\binom3j(L-1)^j}{j+1}
                         =\frac{L^4-1}{4(L-1)}
\]
and sum \(x_\nu=1,\ldots,(L-1)/2\) to obtain
\eqref{hfm:central-edge-lower-bound}.
Since \(a_L\to1/32\), a bound written as
\(C_Lh^4\|A\|_\infty\) cannot in general have
\(C_L\) bounded as \(L\) grows. This example does
not contradict its uniformity in \(k\) for fixed \(L\).

\subsubsection{Uniform row amplitudes and a radius that may depend on \(L\)}

For completeness the row majorant has a different
dependence on \(L\). Normalized path and mean
expansions give higher absolute row coefficients
\((CL)^m\) at one step. Their proof uses products
of at most \(CL\) edge factors and normalized
averages, and has no factor equal to the number
of sites in the block. The combined linear row
is exactly \(L\). Thus the positive majorant is
\[
 b_{t+1}\le Lb_t+
                  \frac{C(Lb_t)^2}{1-CLb_t}.
\]
For a partial output spacing \(\xi=hL^s\),
\(\sum_{t<s}L\xi_t\le L\xi/(L-1)\le3\xi/2\).
The first-exit product, the node-frame product
and the final plain logarithm therefore give
the following bound for the divided degree-\(m\)
Taylor row coefficients \(T^{(m)}=D^m/m!\):
\begin{equation}
 \sum_{b_1,\ldots,b_m}\|T_{h,s;b;b_1\ldots b_m}^{(m)}\|
                      \le C^m\xi^m,
 \label{hfm:uniform-L-plain-rows}
\end{equation}
with \(C\) independent of \(L\ge3\).
The physical row support has radius \(C\xi\),
by the same sum. This statement applies to the
literal mean as well as to the straightened row
and its node correction.

The free inverse and its smooth seeds in
\cref{ch:hb-strong-inverse} use the direct rows
\eqref{hb:actual-flat-rows} at \(h=1/N\).
Their constants are uniform over odd \(N\);
they do not depend on a factorization \(N=L^k\).
The Wilson residual and the smooth node extension
also have such mesh-uniform constants.
The nonlinear transpose constant in
\eqref{hfm:hierarchical-column-error} may instead
grow as \(L^3\). Decrease the nonlinear ball
after fixing \(L\), so that its contribution to
the fixed-point Lipschitz bound is at most one half.
The solution estimate is still
\[
       \|A(B)\|_{X_\alpha}+\|\omega(B)\|_\infty
                    \le2\|\mathcal L_0^{-1}\|\,\|B\|_\infty .
\]
Its amplitude constant is uniform in \(L\);
its allowed radius need not be.
The actual node product, its local extension
and finite-gauge formula preserve this uniform
amplitude estimate for the literal field.
Equation \eqref{hfm:uniform-L-plain-rows}, stopped
one level before the output, consequently gives
\begin{equation}
                 \|\log p_j^{\rm lit}(e^{iB})\|_\infty
                                \le \frac C L\|B\|_\infty .
 \label{hfm:uniform-amplitude-coarse-contraction}
\end{equation}
Here \(p_j^{\rm lit}=M^{j-1}U_j\) is the literal
partial mean, before the one-step axial gauge,
and \(B\) lies in the smaller ball just specified.
The axial representative need not have this raw
angle bound: a flat constant Abelian input can
become a field with identity internal links and
full output angle on the face-crossing links.
For a complete gauge-invariant physical function,
the two representatives give the same value
on their admitted overlap. Apply that complete
identity before composing any noninvariant local
representatives with \(p_j^{\rm lit}\).
A fixed analytic-radius shrink in these literal
coordinates can be paid by choosing a sufficiently
large blocking integer once. A full iteration
still requires the representative identity and
the local analytic-family and source-domain
bounds at its buffered radii.

% END LAYERED EXPANSION formation/hierarchical_mean_second_jet.tex

% BEGIN LAYERED EXPANSION background/flat_reference_descent.tex
\subsection{Descent between real flat reference charts}
\label{ch:frd-flat-reference-descent}

The homogeneous background was constructed on a fixed complex bond
ball about each real flat connection. We now show that these physical
background orbits agree wherever the balls overlap. A comparison of
coarse flat connections must retain their holonomies. A globally
periodic pure gauge is generally unavailable; the construction below
uses a smooth gauge on the universal cover with its exact two-sided
deck law. It selects no logarithm or root of a global holonomy.

The physical coarse torus has integer side $M\ge2$, and the fine
spacing is $h=L^{-k}$, $k\ge0$. Coarse links have unit length. All
connections in the first geometric construction are real
$\operatorname{SU}(2)$ connections. Complexification will be used only
after a uniformly close real flat fine extension has been produced.
Constants depend on $L$, $0<\alpha<1$ and a fixed scalar bump, but
not on depth, volume or holonomy.
Bond-ball distances use a unitary-invariant matrix operator norm;
it is equivalent to the Pauli norm in the strong background estimates.

\subsubsection{Parallel frames with their based holonomies}

For a flat connection on a periodic lattice, choose its parallel
frame on the universal cover, based at the coarse node zero. It is
convenient to write the links as
\begin{equation}
                         W(x,x')=p(x)p(x')^{-1},\qquad p(0)=1.
 \label{frd:parallel-potential}
\end{equation}
Thus $p(x)$ is the inverse of the path transport from zero to $x$.
Let $H_\nu$ denote the actual positively oriented holonomy from
zero to $Me_\nu$. Periodicity and flatness give
\begin{equation}
          p(x+Me_\nu)=p(x)H_\nu^{-1}.
 \label{frd:parallel-deck}
\end{equation}
Indeed the two parallel potentials for the translated connection
differ by a constant right factor, determined at zero. The four
$H_\nu$ commute, since the fundamental cycles commute and the
connection is flat. They need not be close to the identity.

Let $V_0,V_1$ be real flat coarse connections, and let $W_0$ be
any real flat fine extension of $V_0$. On coarse nodes the potential
$p_0$ of $W_0$ also represents $V_0$: every path in a flat averaging
row has the same endpoint transport, so $M^kW_0=V_0$ is that
transport. Choose the coarse parallel potential $p_1$ for $V_1$.
In the parallel frame of $W_0$, put
\begin{equation}
 g_z=p_0(z)^{-1}p_1(z),\qquad
 \widetilde V_1(z,z+e_\mu)=g_zg_{z+e_\mu}^{-1}.
 \label{frd:relative-coarse-gauge}
\end{equation}
Here a tilde denotes the frame change $p_0^{-1}$, which makes
$W_0$ the identity on the cover. With the actual based holonomies
$H_{i,\nu}$, equations \eqref{frd:parallel-deck} give exactly
\begin{equation}
             g_{z+Me_\nu}=H_{0,\nu}g_zH_{1,\nu}^{-1}.
 \label{frd:mixed-deck-law}
\end{equation}
The order of these factors matters. In particular $g_z$ need not
be periodic or globally close to one.

Suppose
\begin{equation}
              \sup_c\|V_1(c)-V_0(c)\|\le\varepsilon.
 \label{frd:coarse-flat-distance}
\end{equation}
The two parallel-frame factors are unitary, so
$\|\widetilde V_1(c)-1\|\le\varepsilon$. Consequently
$\|g_z-g_{z+e_\mu}\|\le\varepsilon$. Any fixed-radius unit
star of nodes has
\begin{equation}
                     \|g_z-g_{z'}\|\le C_*\varepsilon
 \label{frd:local-gauge-cluster}
\end{equation}
by a coordinate path with a fixed number of edges. This estimate
uses no path going around the torus. In particular it does not
bound $H_{1,\nu}-H_{0,\nu}$ by a volume-independent multiple of
$\varepsilon$.

\subsubsection{A smooth quaternion interpolation respecting both deck factors}

Choose a nonnegative smooth product bump $\rho$, supported in
$(-1,1)^4$ and positive on $[-3/4,3/4]^4$. On $\mathbb R^4$ set
\begin{equation}
 \eta_z(X)=\frac{\rho(X-z)}{\sum_{w\in\mathbb Z^4}\rho(X-w)},\qquad
                         M_g(X)=\sum_z\eta_z(X)g_z.
 \label{frd:cardinal-partition}
\end{equation}
Only a fixed number of terms are nonzero. The denominator has
uniform positive lower and finite upper bounds. Thus
$\sum_z\eta_z=1$, $|\partial^r\eta_z|\le C_r$ and
$\eta_z(y)=\delta_{zy}$ at the integer nodes. The last assertion
uses the open support: every other bump and all its derivatives
vanish at that boundary point. Also
$\eta_{z+Mn}(X+Mn)=\eta_z(X)$.

\begin{lemma}[Normalized quaternion interpolation]
\label{lem:frd-quaternion-interpolation}
For a fixed small $\varepsilon$, the formula
\begin{equation}
                    g(X)=\frac{M_g(X)}{\sqrt{\det M_g(X)}}
 \label{frd:normalized-quaternion-interpolation}
\end{equation}
defines a smooth $\operatorname{SU}(2)$ function, using the positive
real scalar square root. It has $g(z)=g_z$ and the exact deck law
\begin{equation}
                  g(X+Me_\nu)=H_{0,\nu}g(X)H_{1,\nu}^{-1}.
 \label{frd:smooth-deck-law}
\end{equation}
For $1\le r\le3$,
\begin{equation}
                            \|\partial^r g\|_\infty\le C_r\varepsilon.
 \label{frd:smooth-gauge-bounds}
\end{equation}
These bounds also hold for the right-translated derivatives.
\end{lemma}
\begin{proof}
Every $g_z$ has real quaternion form $a_zI+i b_z\cdot\sigma$,
with $a_z^2+|b_z|^2=1$. Real scalar averaging preserves that form.
For a matrix $M=aI+i b\cdot\sigma$ with real coefficients,
$M^*M=(a^2+|b|^2)I$ and $\det M=a^2+|b|^2$.
At a given $X$ fix one active node $z_0$. All other active
nodes belong to a fixed star, so
$\|M_g(X)g_{z_0}^{-1}-1\|\le C\varepsilon$ by
\eqref{frd:local-gauge-cluster}. Hence
$\det M_g\ge(1-C\varepsilon)^2>0$. Formula
\eqref{frd:normalized-quaternion-interpolation} is therefore
globally smooth and lies in $\operatorname{SU}(2)$.

For $r\ge1$, subtract $g_{z_0}$ in the differentiated partition
sum, since $\sum_z\partial^r\eta_z=0$. It follows that
$\|\partial^r M_g\|\le C_r\varepsilon$ for $r\le3$.
Differentiating the determinant and its positive scalar square
root gives the same bounds for their derivatives. Their uniform
positive lower bound proves \eqref{frd:smooth-gauge-bounds}.
No derivative of a globally selected root node is taken: on each
star $z_0$ is fixed, and the resulting uniform estimates cover
the whole space.

Normalization by this determinant is equivariant under constant
left and right $\operatorname{SU}(2)$ multiplication. Both determinant
factors are one, and the positive square root is unchanged.
Reindexing \eqref{frd:cardinal-partition} by a deck shift and using
\eqref{frd:mixed-deck-law} proves \eqref{frd:smooth-deck-law}.
Cardinality gives $g(z)=g_z$. Finally right translation by the
unitary $g(X)^{-1}$ preserves the derivative norms.
\end{proof}

Only a positive real number is square-rooted in this argument.
There is no choice of a continuous matrix logarithm of a holonomy,
including at a central holonomy, and no choice of an
$h$-dependent group root.

\begin{proposition}[A uniformly close flat fine extension]
\label{prop:frd-close-flat-extension}
There is a real flat fine extension $W_1'$ of $V_1$ such that
\begin{equation}
 W_1'=e^{ihC_h}W_0,\qquad
                         \|C_h\|_{X_\alpha(W_0)}\le C_\alpha\varepsilon.
 \label{frd:close-flat-extension-bound}
\end{equation}
It is periodic in the original frame, has the actual holonomies
of $V_1$, and has $M^kW_1'=V_1$ exactly.
\end{proposition}
\begin{proof}
In the $W_0$ parallel frame define
$\widetilde W_1'(x,x+he_\mu)=g(x)g(x+he_\mu)^{-1}$.
Its plaquettes are one by telescoping. The deck law is conjugation
by $H_{0,\nu}$, so undoing the $p_0^{-1}$ frame gives an ordinary
periodic fine connection.

Every flat path from a coarse node $y$ to $y+e_\mu$ has transport
$g_y g_{y+e_\mu}^{-1}$ in this frame. Its actual contour means
therefore give that same matrix. Identical transported inputs have
their exact common mean even when that common value is far from
one: the defining relative products are one, and unitary covariance
translates their regular branch. Thus the full average is exactly
$V_1$. Its coarse cycle products are the fine cycle products, so
the based fine holonomies are $H_{1,\nu}$.

Set
$F_{h,\mu}(X)=g(X)g(X+he_\mu)^{-1}$.
The integral difference formula and
\eqref{frd:smooth-gauge-bounds} imply
$\|\partial^r(F_{h,\mu}-1)\|\le Ch\varepsilon$ for $r=0,1,2$.
For $r=2$ this uses the third derivative of $g$; products of
lower derivatives are also bounded by $C\varepsilon$.
The logarithm near one consequently satisfies
\[
 \left\|\partial^r\frac1{ih}\log F_{h,\mu}\right\|_\infty
                                  \le C\varepsilon,\qquad r=0,1,2.
\]
Sampling this continuous expression on fine sites bounds its
physical first differences and their local $\alpha$ seminorm,
for every $\alpha<1$. The frame change is unitary and converts
ordinary differences into the flat covariant differences of
$W_0$. This proves \eqref{frd:close-flat-extension-bound}.
The logarithm is traceless anti-Hermitian, so $C_h$ is real.
\end{proof}

\subsubsection{Independence of the chosen real flat fine extension}

\begin{lemma}[Node-fixed equivalence of flat extensions]
\label{lem:frd-flat-extension-equivalence}
If $W_1$ is any other real flat fine extension of the same literal
coarse connection $V_1$, then there is a periodic unitary fine
gauge $f$, equal to one at every coarse node, with $W_1^f=W_1'$.
Applying $f$ simultaneously to a field and its reference preserves
its reference-covariant strong norm.
\end{lemma}
\begin{proof}
Let $T_1(0,x)$ and $T'(0,x)$ be their path transports on the
cover and set $f(x)=T'(0,x)^{-1}T_1(0,x)$.
Flatness shows directly on every edge that $W_1^f=W_1'$.
On a coarse path both transports are the same product of the
links of $V_1$. Hence $f$ is one at all coarse nodes. Equality
of the based holonomies makes $f$ periodic: their left deck
factors cancel in its displayed formula. The argument uses
equality of every coarse link. Four equal holonomies alone
would not force the gauge to be one at all coarse nodes.

For $U=e^{iha}W_1$, one has
$U^f=e^{ih\operatorname{Ad}(f)a}W_1^f$. The unitary frame change
intertwines every reference-covariant difference and parallel
comparison. Thus its supremum and local H\"older norms are
preserved. No ordinary fine derivative of this possibly large
gauge is estimated or required.
Unitary covariance also preserves the admitted complex mean branches.
\end{proof}

\subsubsection{Complex uniqueness inside a flat-reference strong chart}

We record the complex uniqueness statement needed on overlaps.
It is stronger than equality obtained by analytic continuation
from a possibly empty real overlap.

\begin{lemma}[Complex critical-orbit uniqueness]
\label{lem:frd-complex-uniqueness}
There are fixed $\delta_c,\rho_c>0$, uniform in depth, volume and
real flat reference, with the following property. Let
$V=e^{iB}V_0$, $\|B\|_\infty<\rho_c$. A complex literal
constrained-critical connection with a representative
$U=e^{ihA}W_0$, $\|A\|_{X_\alpha(W_0)}<\delta_c$, is in the
node-fixed complex gauge orbit of the constructed background.
The comparison gauge is holomorphic in any holomorphic family
of such representatives.
\end{lemma}
\begin{proof}
The canonical local branch is holomorphically critical: its real
stationarity identities, written using the analytic right inverse,
extend to its complex parameter ball. This is continuation inside
one chart about a real center, not a claim about a real overlap.

The actual mean-slice contraction
\eqref{hb:mean-slice-solution} is an analytic contraction between
the displayed complex Banach spaces. It sends $U$ to
$\overline A$ with $q_h(\overline A)=B$ and
$RD^\sharp\overline A=0$. Its group equation fixes the complete
effective coarse gauge, so the subsequent literal correction
differs from $U$ by an exactly node-fixed gauge, including at
reducible coarse connections.

The two analytic fibre maps in the proof of
\eqref{hb:literal-to-mean-slice} identify all tangent directions
and preserve the action. Thus holomorphic constrained criticality
implies $J_h(\overline A)=q_h'(\overline A)^\sharp\omega$.
The actual derivative is onto with right inverse
\eqref{hb:actual-nonlinear-row-lift}. The multiplier readout
\eqref{hb:first-variation-readout} uses absolute values and the
physical $W^{1,1}$ seed-column bound, so it remains valid for
complex fields and gives
$\|\omega\|_\infty\le C\|\overline A\|_{X_\alpha}$.
This estimate uses no positivity.

Choose $\delta_c$ so the resulting pair lies in the joint
contraction ball of \eqref{hb:actual-fixed-point}. Its reference
Landau penalty is zero by the imposed slice. It therefore solves
that actual fixed-point equation and equals its unique complex
solution. The explicit mean-slice gauge and literal node
extension give the claimed comparison. They depend
holomorphically on the representative. Real positivity was used
to prove a minimum, but is not used in this uniqueness argument.
\end{proof}

\subsubsection{Changing the reference in a strong complex neighbourhood}

In the smooth relative frame of
\cref{prop:frd-close-flat-extension}, an increment $a$ measured
relative to $W_1'$ has the same strong norm as
$\operatorname{Ad}(g^{-1})a$ in the trivial frame. Product
differences and \eqref{frd:smooth-gauge-bounds} give
\begin{equation}
 \|a\|_{X_\alpha(W_0)}\le C_\alpha\|a\|_{X_\alpha(W_1')}.
 \label{frd:reference-norm-comparison}
\end{equation}
Only the first two bounded derivatives of $g$ enter. Its values
may range throughout the unitary group. Locally its variation
is $O(\varepsilon)$, and unitary conjugation does not change
the zeroth-order size.

The pointwise matrix map
\begin{equation}
                Z_h(a,C)=\frac1{ih}\log(e^{iha}e^{ihC})
 \label{frd:relative-BCH}
\end{equation}
extends analytically to $h=0$ with value $a+C$. Its nonlinear
part is $h$ times a uniformly analytic quadratic expression.
The discrete product rule and the Banach algebra estimate for
$X_\alpha$ give, on a fixed small ball,
\begin{equation}
 \|Z_h(a,C_h)\|_{X_\alpha(W_0)}
        \le C_\alpha(\|a\|_{X_\alpha(W_1')}+\varepsilon).
 \label{frd:strong-BCH-comparison}
\end{equation}
There is no $h^{-1}$ loss: the logarithm in
\eqref{frd:relative-BCH} vanishes to first order in $h$.

Let $\mathcal B_\rho(V_f)$ be the coarse complex ball
$V=e^{iB}V_f$, $\|B\|_\infty<\rho$, about a real flat $V_f$.
If two such balls intersect, their real centers satisfy
\begin{equation}
                 \|V_1-V_0\|_\infty\le2e^\rho\rho\le3\rho
 \label{frd:overlap-centers-close}
\end{equation}
for a fixed small $\rho$. Choose $\rho$ below the homogeneous
existence radius, the real interpolation threshold, and the
complex uniqueness radius divided by the constants in
\eqref{frd:strong-BCH-comparison}. All these choices are uniform.

Take the $V_1$-based background and use the real unitary
node-fixed gauge of \cref{lem:frd-flat-extension-equivalence}
to put its flat fine reference at $W_1'$. Its increment has
$X_\alpha(W_1')$ norm at most $C\rho$. Equations
\eqref{frd:strong-BCH-comparison} and
\eqref{frd:overlap-centers-close} put that actual connection
inside the $W_0$ complex uniqueness chart, with strong norm
at most $C'\rho<\delta_c$. Its coarse mean is still $V$,
and it remains holomorphically constrained critical.
\Cref{lem:frd-complex-uniqueness} identifies it with the
$V_0$-based background orbit.

\subsubsection{Holomorphic orbit and action descent}

\begin{proposition}[Uniform descent over real flat centers]
\label{prop:frd-flat-atlas-descent}
For a fixed $\rho>0$ independent of $M,k$, the physical
homogeneous background orbits and their total actions descend
to the open set
\begin{equation}
 \mathcal D_\rho^{\mathbb C}
       =\bigcup_{V_f\ \mathrm{real\ flat}}\mathcal B_\rho(V_f).
 \label{frd:flat-atlas-domain}
\end{equation}
The total action is a single holomorphic function there. The
local background representatives have holomorphic node-fixed
transition gauges, with the cocycle identity on triple
overlaps. On real data these gauges are unitary. The descended
orbit is covariant under all real coarse gauge transformations,
and its total action is invariant under them.
\end{proposition}
\begin{proof}
The preceding argument gives equality of the physical orbits
on each complex overlap directly, without a real-overlap
identity theorem. The transition gauge is explicit: apply the
fixed real comparison gauge $f$, then the actual holomorphic
mean-slice gauge for the resulting $W_0$ representative, then
the canonical literal extension at its mean-slice solution.
Their product has value one at every coarse node. The exact
group equations in the mean-slice construction prove this,
not merely a first-order calculation.
After factoring out the fixed real comparison gauge, the remaining
site gauge has a logarithm of $Y_\alpha(W_0)$ norm at most $C\rho$.
The comparison formula is analytic on a fixed larger neighborhood
of the inner overlap: choose an intermediate coarse radius so its
background and BCH image still lie inside $\delta_c$, and then
choose the atlas radius much smaller. Cauchy's formula on that
intermediate ball gives uniform parameter jets of the transition
in these relative norms. The possible thinness of the final overlap
therefore introduces no volume- or depth-dependent constant.

The node-fixed gauge action is free on a connected fine graph.
Indeed a stabilizing gauge is transported from its value at
one coarse node along every fine edge; that value is one.
Consequently the transition gauge between two given literal
representatives is unique. If $U_i=U_j^{t_{ij}}$, this implies
$t_{02}=t_{01}t_{12}$ on triple overlaps. This argument does
not require irreducibility of the coarse holonomies.
Local representatives and these holomorphic transitions give
the asserted holomorphic orbit descent. Fine Wilson trace
invariance gives equality of their holomorphic total actions,
so ordinary function gluing gives a single holomorphic action.
Only its local strong representatives are asserted; a globally
chosen gauge with the same strong bound is not needed.

For a real coarse gauge $u$, the center $V_f^u$ is again real
flat, and
$V^u(V_f^u)^{-1}=u(VV_f^{-1})u^{-1}$ at each initial node.
Thus the ball is carried isometrically to the ball about
$V_f^u$. Extend $u$ by any unitary fine site values with the
prescribed coarse values, and transform a field and its flat
reference together. Their reference-covariant strong norms
are preserved exactly. Their literal means transform by $u$
and their actions are unchanged. Complex critical-orbit
uniqueness at the transformed reference identifies the
transformed branch with the descended branch. Two choices
of the fine extension differ by a node-fixed gauge, so the
orbit statement is independent of that choice. This proves
the last assertions, including arbitrary real coarse gauges.
\end{proof}

The domain \eqref{frd:flat-atlas-domain} is a union of small
complex balls about real flat connections. It is not asserted
to contain all curvature-only regular data or to be saturated
under arbitrary complex gauges. On an overlap the equality is
the nonlinear total identity $S_0(B_0(V))=S_1(B_1(V))$.
Taylor polynomials about the two centers need not agree
term by term. This descent therefore does not identify local
fixed-reference Taylor pieces, their filing, their contour
cutoffs, or older fine-core activity transports. Those require
their own coordinate and domain comparisons.

% END LAYERED EXPANSION background/flat_reference_descent.tex

% BEGIN LAYERED EXPANSION background/curvature_trial_atlas.tex
\subsection{A homogeneous critical atlas from coarse curvature alone}
\label{ch:cta-curvature-trial-atlas}

We construct the regular trial which was an input to
\cref{prop:rtc-nonstationary-trial}. The coarse connection need
not be close to a flat connection in a global bond gauge. Its
plaquettes are small, but its holonomies and topological sector
are unrestricted. A smooth quaternionic line bundle interpolates
the coarse connection. Sampling its connection gives fine trials
with bounded physical curvature and current at every depth.

The setting is homogeneous. Coarse spacing is one, the fixed
blocking integer is the odd number $L$, and $h=L^{-k}$ is the
fine spacing. The coarse torus has compatible integer periods
at least two. Write $V_{ij}$ for transport from the head $j$
to the tail $i$, and $V_{ji}=V_{ij}^{-1}$. Assume
\begin{equation}\label{cta:coarse-curvature}
       \|V(\partial p)-1\|\le\kappa
                         \quad\hbox{for every coarse plaquette }p.
\end{equation}
For small angles this norm and the admitted plaquette logarithm
are equivalent. All smallness choices below are independent of
$h$, the periods and the number of coarse cells. Constants may
depend on $L$ and on the fixed regularity exponents and collars.

\subsubsection{Local transports without logarithms of coarse links}

We work first on the universal covering lattice $\mathbb Z^4$.
Lift $V$ periodically. Distinct lifted vertices remain distinct
indices, even if they have the same image on a small torus.
This convention is important on a two-cell period: its two
oppositely placed crossing edges can have different holonomy.

A bounded closed edge word in a box of fixed diameter can be
reduced by interchanging adjacent coordinate steps and canceling
inverse pairs. An interchange inserts a conjugate of one
plaquette. A word of bounded length requires a bounded number of
such interchanges. Since all transports are unitary,
\begin{equation}\label{cta:short-loop-bound}
 \|V(\gamma)-1\|\le C_R\kappa
\end{equation}
for every contractible edge word of length and diameter at most
$R$. This is a finite word estimate on the covering lattice.
It is not a bound for a noncontractible torus loop.

One may fix short paths with reverse paths paired. We specify
a symmetric choice to retain lattice covariance. For two vertices
in a fixed bounded box, take the at most $4!$ coordinate-order
paths, traversing each nonzero coordinate displacement in one
block. Let their backward transports be $W_\gamma$ and set
\begin{equation}\label{cta:short-transport}
 M_{ij}=\frac1{n_{ij}}\sum_\gamma W_\gamma,\qquad
 U_{ij}=M_{ij}(\det M_{ij})^{-1/2},\qquad U_{ii}=1.
\end{equation}
An average of real quaternion matrices is quaternionic and
$M_{ij}^\sharp M_{ij}=(\det M_{ij})I_2$. Here $\sharp$ on a
$2$ by $2$ matrix means its quaternionic adjugate; on real
quaternion matrices it is the Hermitian transpose. The paths
differ by bounded contractible loops. Thus $M_{ij}$ is within
$C_R\kappa$ of any one $W_\gamma$, and its determinant is
positive and close to one. The square root in
\eqref{cta:short-transport} is the positive scalar root.
It is not a root or logarithm of a coarse group element.
Only overlapping partition indices are used in the projection
below. Auxiliary paths to an anchor are used only on the fixed
read boxes. Fix their maximum radius after the required mean and
regularity collars, and choose $\kappa$ afterward. No estimate
for unbounded-length paths is asserted.

This construction is equivariant under left and right unitary
multiplication, and under reversal. Indeed the reversed path
average is $M_{ij}^\sharp$, its determinant is unchanged, and
$U_{ji}=U_{ij}^\sharp=U_{ij}^{-1}$. For adjacent vertices there
is one path, so $U_{ij}$ is exactly the given coarse link.
Coordinate isometries permute the path multiset. In particular
the choice is translation and hypercubic covariant. The short
loop bound also gives, for three vertices in a fixed box,
\begin{equation}\label{cta:triangle-defect}
                  \|U_{ij}-U_{iy}U_{yj}\|\le C_R\kappa.
\end{equation}
To see this for the normalized averages, first use
\eqref{cta:short-loop-bound} for one chosen path in each factor,
then use their $C_R\kappa$ difference from the averages.

\subsubsection{A smooth rank-one projection}

Choose an even nonnegative smooth function $\varphi$, supported
in $(-1,1)$, positive on $[-1/2,1/2]$, with $\varphi(0)=1$.
For $i\in\mathbb Z^4$ put
\begin{equation}\label{cta:cardinal-partition}
 f_i(x)=\prod_{\mu=1}^4\varphi(x_\mu-i_\mu),\qquad
 s_i(x)=\frac{f_i(x)}{(\sum_j f_j(x)^2)^{1/2}},\qquad
                       \psi_i=s_i^2.
\end{equation}
The denominator is uniformly bounded below. The $s_i$ are
smooth, including at their support boundaries; all their fixed
derivatives are uniformly bounded. At most sixteen indices are
active at any point, $\sum_i s_i^2=1$, and $s_i(j)=\delta_{ij}$.
Along a coarse edge only its two endpoints can be active.

Use these indices as quaternionic coordinates and form
\begin{equation}\label{cta:almost-projection}
                         P_{ij}(x)=s_i(x)s_j(x)U_{ij}.
\end{equation}
This is a bounded self-adjoint quaternionic operator on
$\ell^2(\mathbb Z^4;\mathbb H)$, but every calculation below
is finite: its nonzero block at a point has at most sixteen
quaternionic rows. Derivatives on a fixed small neighborhood
use a fixed finite union of those indices. No infinite
determinant or infinite-dimensional Gaussian is introduced.

Fix an anchor $y$ and a physical box of fixed radius $R$ about
it. Define, throughout that box,
\begin{equation}\label{cta:comparison-frame}
              u_y(x)_i=s_i(x)U_{iy}.
\end{equation}
The anchor need not be active. Its transports have bounded
length in a fixed larger box. The exact identities are
$u_y^\sharp u_y=I_2$ and $u_y^\sharp \mathrm{d}u_y=0$; the latter follows
from $\sum_i s_i\,\mathrm{d}s_i=0$, since each $U_{iy}$ is independent
of $x$. By \eqref{cta:triangle-defect},
\begin{equation}\label{cta:projection-comparison}
           \|P-u_yu_y^\sharp\|_{C^m}\le C_{m,R}\kappa
\end{equation}
for every fixed $m$. The matrix norm has no volume factor,
because the active block is bounded in size.

Choose $\kappa$ below the finitely many margins used below.
At each point the spectrum of $P$ lies in disjoint small
neighborhoods of zero and one. Define
\begin{equation}\label{cta:spectral-projection}
 \Pi(x)=\frac1{2\pi i}\int_{|z-1|=1/4}(z-P(x))^{-1}\,dz.
\end{equation}
The contour resolvent is computed on the finite active block;
the integral on its zero complement is zero. The resolvent
identity, or its Neumann expansion about $u_yu_y^\sharp$, proves
\begin{equation}\label{cta:Pi-local-bound}
 \Pi^2=\Pi=\Pi^\sharp,\quad
 \operatorname{rank}_{\mathbb H}\Pi=1,\quad
       \|\Pi-u_yu_y^\sharp\|_{C^m}\le C_{m,R}\kappa.
\end{equation}
For example every differentiated resolvent is a finite sum of
products of bounded resolvents and derivatives of $P$; a
difference from the comparison projector contains at least
one $O(\kappa)$ difference. This proves the derivative estimate
without a spectral-eigenvector choice. The quaternionic rank
is one, equivalently its complex rank is two.

The ranges $E_x=\operatorname{ran}\Pi(x)$ form a smooth
quaternionic line bundle on the cover. Put on it the connection
$\nabla=\Pi \mathrm{d}$. A normalized local frame on the whole fixed
box is
\begin{equation}\label{cta:local-frame}
 v_y=\Pi u_y\,(u_y^\sharp\Pi u_y)^{-1/2}.
\end{equation}
The middle $2$ by $2$ matrix is a positive scalar multiple of
$I_2$: it is a self-adjoint quaternion of size one. It differs
from $I_2$ by $O(\kappa)$. Thus
$v_y^\sharp v_y=I_2$, $v_yv_y^\sharp=\Pi$, and
$v_y-u_y=O(\kappa)$ in every fixed local $C^m$ norm. Its connection
matrix and its curvature satisfy
\begin{equation}\label{cta:smooth-connection}
 A_y=v_y^\sharp \mathrm{d}v_y,\qquad
 \|A_y\|_{C^m}+\|\mathrm{d}A_y+A_y\wedge A_y\|_{C^{m-1}}
                                             \le C_{m,R}\kappa.
\end{equation}
Indeed $u_y^\sharp \mathrm{d}u_y=0$ cancels the order-one terms from
the partition. The remaining terms have a factor $v_y-u_y$
or its derivative. The matrices $A_y$ are traceless
anti-Hermitian. Local frames have exact unitary transition
functions $v_z^\sharp v_y$. We have not chosen a global smooth
frame, which need not exist.

\subsubsection{Exact edge interpolation and periodic descent}

On a coarse edge $(y,z)$ the two-index block of $P$ is exactly
a projection. Its unit frame is
\[
             u=(s_y I_2,\ s_zU_{zy})^{\mathsf T}.
\]
Its connection $u^\sharp \mathrm{d}u$ is zero along that edge. At the
endpoints this frame is $e_y$ and $e_zU_{zy}$, respectively.
Backward connection transport therefore sends
$e_z\zeta$ to $e_yU_{yz}\zeta$. Consequently
\begin{equation}\label{cta:exact-edge}
    \text{backward transport in }(E,\nabla)\text{ along }(y,z)
                         \text{ in canonical endpoint frames}=V_{yz}.
\end{equation}
This checks the orientation against the lattice convention.
It also explains why a potentially large $V_{yz}$ requires
neither a logarithm nor a fractional power.

For each coarse-period translation $t$, let $S_t$ shift the
ambient vertex indices. The construction is exactly equivariant:
\begin{equation}\label{cta:deck-law}
   P(x+t)=S_tP(x)S_t^{-1},\qquad
   \Pi(x+t)=S_t\Pi(x)S_t^{-1}.
\end{equation}
Thus the bundle and connection descend under
$(x,w)\sim(x+t,S_tw)$. At every coarse vertex their canonical
frame $e_y$ descends as well. This gives an ordinary rank-one
quaternionic bundle over the coarse torus. The use of a
covering-coordinate Hilbert space is only a uniform way to
retain distinct local indices; all its bundle charts and
connection computations are finite-dimensional.

Equivalently, any fixed finite collection of read stars can
first be placed in a sufficiently large compatible finite
cover. The finite active-block formulas are the same there.
One descends the equivariant bundle and connection. One does
not average projections or compare partition functions on
different volumes. In particular no short torus loop is
silently declared contractible: the loop filling in
\eqref{cta:short-loop-bound} always occurs between distinct
lifted indices.

\subsubsection{Fine trials and the actual Wilson-current cancellation}

At the fine vertices choose arbitrary unit quaternionic frames
of $E$, with frame $e_y$ at the coarse vertices. A choice is
possible on this finite set even when the bundle is nontrivial.
Define $W_h$ on a fine edge by backward parallel transport of
$\nabla$ between those frames. This is a real $\operatorname{SU}(2)$ lattice
connection. Changing the fine frames changes it by a fine gauge
equal to one at coarse vertices.

On every fixed physical read box use the smooth frame $v_y$.
In this frame $A_y=i a_y$, where $a_y$ is Hermitian and has
all fixed $C^m$ norms bounded by $C_m\kappa$. The parallel
transport equation and its integral iteration give fine links
$W_h(x,\mu)=e^{ih a_{h,\mu}(x)}$ with
\begin{equation}\label{cta:sampled-potential}
                         \|a_h\|_{C^m}\le C_m\kappa.
\end{equation}
Here differences or derivatives of the sampled function have
physical units. To justify division by $h$, the transport
matrix is $1+\int_0^h A_y(x+te_\mu)\,dt+O(h^2\kappa^2)$.
Its admitted logarithm vanishes at $h=0$, so its logarithm
divided by $h$ is smooth with these bounds. Differentiating
the integral equation proves every fixed spatial derivative
estimate. This is a local frame statement, not a global small
potential for the lattice connection.

The same equation applied to a square gives, uniformly in
the local box,
\begin{equation}\label{cta:small-square}
 \|W_h(\partial p)-1\|\le C\kappa h^2,\qquad
 \|\partial_x^m(W_h(\partial p)-1)\|
                                      \le C_m\kappa h^2.
\end{equation}
The missing linear term follows either by multiplying the
four integral expansions or by cancellation of the two
oppositely oriented segments in each direction. Their
second-order coefficient is the connection curvature.
The integral remainder and all its fixed derivatives are
bounded by $C_m\kappa h^3$. These estimates follow from
\eqref{cta:smooth-connection}; no field equation is assumed.

It is essential to pair the plaquettes in the current. Put
$\mathcal S(P)=(P-P^{-1})/(2i)$. For the positive bond
$b=(x,x+he_\mu)$, the two plaquettes in the $\mu\nu$ plane
contribute, in the tail frame,
\begin{equation}\label{cta:paired-current}
 h^{-3}\left\{\mathcal S(P_{\mu\nu}(x))
   -\operatorname{Ad}(W_h(x-he_\nu,\nu)^{-1})
                    \mathcal S(P_{\mu\nu}(x-he_\nu))\right\}.
\end{equation}
This follows by varying the first link in the upper
plaquette and the inverse third link in the lower one,
then using cyclicity of the trace. It is the actual physical
left Wilson current normalization. The translated difference
of the two $\mathcal S$ terms is $O(\kappa h^3)$ by
\eqref{cta:small-square}. The transporting link in this
smooth frame is $1+O(\kappa h)$; its additional commutator
is $O(\kappa^2h^3)$. Summing the three planes proves
\begin{equation}\label{cta:trial-curvature-current}
       \|F(W_h)\|_\infty+\|j(W_h)\|_\infty\le C\kappa,
\end{equation}
with physical fine spacing $h$. The statement is gauge
covariant, so arbitrary vertex frames cause no loss.
Bounding each summand before this pairing would instead
give the false bound $C\kappa/h$ for this purpose.

\subsubsection{The complete hierarchical mean has a uniform error}

The actual $k$-fold mean of a coarse edge reads a fixed
physical star of radius $R_L$, independent of $k$:
the successive stencil radii sum as a geometric series.
Choose the preceding local boxes larger than that star,
including every internal contour and a fixed collar.
Use its tail vertex $y$ as anchor. The local smooth
frame at its head is $e_zU_{zy}$, so the central coarse
edge is $1$ in this frame, by \eqref{cta:exact-edge}.

Extend $a_h$ from the read star by a fixed smooth cutoff
on a larger box. Its supremum and fixed physical
derivatives remain $C_L\kappa$. On a small period one
can do this on a large compatible cover; the mean's
read set is kept before descent. The actual mean is
local and gauge covariant, so this extension does not
alter the edge under consideration.

Apply the absolute plain-mean row majorant
\eqref{hit:plain-mean-absolute-Taylor}, with physical
output length $\xi_s=hL^s$. It gives
\begin{equation}\label{cta:partial-mean-chart}
       \|\log M^s(e^{ih a_h})\|_\infty
                           \le C_L\kappa\xi_s,\qquad 0\le s\le k
\end{equation}
on every intermediate read edge needed for this final
edge. The majorant also controls the actual internal
node factors and closed logarithms. Choose $\kappa$
below its fixed raw input ceilings; thus every actual
mean operation is on its intended analytic branch.
At $s=k$ the local output is $1+O(\kappa)$, while its
central original edge was exactly $1$. Transforming
the endpoints back gives
\begin{equation}\label{cta:mean-error}
 \sup_c\left\|\frac1i\log\{M^k(W_h)(c)V(c)^{-1}\}\right\|
                                                    \le C_L\kappa.
\end{equation}
No factor $C_L^k$ has been iterated. This estimate
uses the full actual plain mean, not a straight-line
linear replacement. Its gauge covariance also shows
that the answer is independent of the arbitrary fine
vertex frames with canonical coarse endpoints.

\begin{proposition}[Actual smooth homogeneous trials]
\label{prop:cta-uniform-trial}
For sufficiently small coarse curvature in
\eqref{cta:coarse-curvature}, every $h=L^{-k}$ has a
real fine trial $W_h(V)$, defined up to a gauge one at
coarse nodes, with \eqref{cta:trial-curvature-current}
and \eqref{cta:mean-error}. Its local smooth-frame
potentials have all fixed $C^m$ bounds $C_m\kappa$.
Its connection has exactly the original coarse-edge
transports before hierarchical averaging. All actual
hierarchical mean charts are admitted. The constants
are uniform in depth and compatible volume.
\end{proposition}

\subsubsection{A holomorphic comparison which has no fine-frame jump}

Fix a real coarse centre $V_*$. Write nearby coarse
data as $V_B=e^{iB}V_*$, with the raw coarse bond
supremum norm of $B$ small. Products along every
fixed short path are holomorphic in $B$. In complex
formulas $\sharp$ is the analytic quaternionic
adjugate, followed by block transpose; it does not
conjugate a parameter. The normalized averages in
\eqref{cta:short-transport} use the scalar branch
near determinant one. Reversal and $P^\sharp=P$
remain algebraic identities. The fixed contour in
\eqref{cta:spectral-projection} therefore defines
a holomorphic idempotent $\Pi_B$, with
\begin{equation}\label{cta:projection-parameter-bound}
       \|\Pi_B-\Pi_*\|_{C^m}\le C_m\|B\|_\infty.
\end{equation}
All these norms are local finite-block norms and
their supremum over physical boxes. The constants
do not count the coarse volume. In particular
$\Pi_B(y)=\Pi_*(y)=e_ye_y^\sharp$ at coarse vertices.

For two close idempotents put
\begin{equation}\label{cta:Kato-map}
 \begin{split}
 T_B={}&\{\Pi_B\Pi_*+(1-\Pi_B)(1-\Pi_*)\}
                  \{1-(\Pi_B-\Pi_*)^2\}^{-1/2}.
 \end{split}
\end{equation}
The inverse square root is its norm-convergent
binomial series. Direct multiplication of the two
idempotents shows
\begin{equation}\label{cta:Kato-identities}
 T_B^\sharp T_B=1,\qquad
 T_B\Pi_*=\Pi_BT_B,\qquad
 T_B(y)=1,\qquad
                   \|T_B-1\|_{C^m}\le C_m\|B\|_\infty.
\end{equation}
For the first equality, if
$S=\Pi_B\Pi_*+(1-\Pi_B)(1-\Pi_*)$, then
$S^\sharp S=1-(\Pi_B-\Pi_*)^2$; the latter commutes
with both idempotents and with $S$. Real data make
$T_B$ unitary. These formulas remain valid on the
smaller complex ball by the same algebra and series.

Choose the fine frames of the deformed trial to be
$T_B(x)e_x^*$, where $e_x^*$ are the fixed frames
of $W_h(V_*)$. Pulling its connection back to $E_*$
gives
\begin{equation}\label{cta:aligned-connection}
     \nabla_B^{\rm al}=\nabla_*+C_B,\qquad
       C_B=\Pi_*T_B^\sharp(\mathrm{d}T_B)\Pi_*\big|_{E_*},\qquad
                           \|C_B\|_{C^m}\le C_m\|B\|_\infty.
\end{equation}
This follows by applying $T_B^\sharp\Pi_B\mathrm{d}$ to
$T_Bs$ for a section $s$ of $E_*$. It compares
connections, not discontinuously chosen vertex frames.

The relative edge transport of these two smooth
connections is $1+O(h\|B\|)$ in the tail reference
frame. The integral transport equation and its
differentiated equations give an admitted logarithm
\begin{equation}\label{cta:aligned-trial-bound}
 W_h(V_B)(b)W_h(V_*)(b)^{-1}=e^{ih E_h(B)(b)},\qquad
       \|E_h(B)\|_{\mathrm{str},\alpha;W_h(V_*)}
                                         \le C_\alpha\|B\|_\infty.
\end{equation}
The local reference connection has its uniform
$C^m$ bounds, so ordinary derivatives of the
relative transport give the stated covariant
differences and H\"older seminorm. Division by $h$
is removable in this relative equation, just as in
\eqref{cta:sampled-potential}. Both the map and its
fixed-order derivatives are holomorphic uniformly
in $h$. A new arbitrary anchor choice at each
$B$ would not justify this estimate: it could
create a jump of order $B/h$ in relative coordinates.
The prescribed aligned frames avoid that jump.

The construction is local in the coarse data on
each fixed physical box: $P$ reads finitely many
short paths, and the functional calculus and $T_B$
act pointwise on its finite active block. Thus its
parameter estimates do not hide a volume-dependent
global gauge choice.

\subsubsection{Critical correction at every curvature-small datum}

Fix the real trial $W_h(V_*)$, with approximate
multiplier zero. By \eqref{cta:trial-curvature-current}
it satisfies the actual regularity and defect
hypotheses of \cref{prop:rtc-nonstationary-trial},
with defect at most $C\kappa$. The homogeneous
determining region has no artificial inner boundary;
the actual site and field rows are admitted by
\eqref{cta:partial-mean-chart}. All fixed inverse
and regularity collars can be chosen before
decreasing $\kappa$.

Let $V_*^{\rm tr}=M^kW_h(V_*)$. For target
$V_B=e^{iB}V_*$, its literal datum relative to
the trial is
\begin{equation}\label{cta:RTC-datum}
       B^{\rm tr}(B)=\frac1i\log(V_B(V_*^{\rm tr})^{-1}),
       \qquad \|B^{\rm tr}(B)\|_\infty
                              \le C(\kappa+\|B\|_\infty).
\end{equation}
This logarithm is close to the identity even
when the individual coarse links are not.
Apply the actual forced fixed point of RTC
to this datum. It retains the source
$-j(W_h(V_*))$, removes its gauge penalty
by the exact gauge tests, and uses the actual
mean slice and literal correction. It gives
a holomorphic critical orbit with mean
exactly $V_B$. In its trial chart its field
and multiplier are bounded by
$C(\kappa+\|B\|_\infty)$.
In particular a critical orbit at $V_*$
has actual physical curvature and current
at most $C\kappa$. No critical reference
or globally small coarse potential was
assumed in obtaining it.

\subsubsection{Identification on overlaps and the resulting atlas}

Consider two real centres $V_*,\widetilde V_*$
whose sufficiently small complex datum balls
overlap. Their raw relative coarse distance
is at most $C\rho$, where $\rho$ is the
common radius. The real unitary map in
\eqref{cta:Kato-map} aligns their trial bundles.
It is the identity at coarse nodes. Fine
frames chosen differently from this alignment
are related by a real fine gauge one at those
nodes. After that fixed gauge, their real
trials differ by at most $C\rho$ in the
covariant strong norm of
\eqref{cta:aligned-trial-bound}.

Both complex critical germs on the overlap
differ from their respective trials by at
most $C(\kappa+\rho)$. Reference strong norms
are equivalent under the above small real
connection change. Choose $\rho$ and then
$\kappa$ so that their combined bound lies
inside the larger RTC uniqueness chart of
the first trial. They have the same literal
coarse datum. The actual complex mean-slice
and competing-multiplier uniqueness of RTC
therefore identify their orbits modulo
fine gauges one at coarse nodes. No positivity
of a complex square is needed in this step.

These orbit identifications are holomorphic.
Indeed the fixed-reference mean-slice maps
and the specified literal correction used
in RTC are holomorphic. They send the two
germs to the same unique slice solution;
composing their admitted gauges gives the
transition gauge, with identity coarse-node
values. On a triple overlap this describes
the same orbit, irrespective of which
sequence of Kato maps was used. Kato maps
themselves need not obey a multiplicative
cocycle. Alternatively both local trial
frames can be placed in the common bundle
$E(V)$ by $T(\Pi(V),\Pi(V_*))$; ordinary
frame changes there obey the exact cocycle.
The node-fixed fine gauge action is free:
an element fixing the fine connection is
propagated from its identity value at a
coarse node along every fine edge. Thus
the transition gauges for specified local
representatives also obey that cocycle.
The total Wilson action is unchanged by
all these gauges.

Changing the bounded short-path policy does
not change this orbit. Its used transports
differ by $O(\kappa)$ by
\eqref{cta:short-loop-bound}; hence the two
projections differ by $O(\kappa)$ in every
fixed local derivative norm. Their Kato
identification is the identity at coarse
nodes and their trial connections differ
by $O(\kappa)$ in the strong norm. The same
RTC uniqueness argument identifies the two
critical corrections. The partition and
all finite read collars remain fixed in
this comparison.

The interpolating construction is equivariant
under real coarse gauge transformations.
For $g=(g_i)$, the ambient constant diagonal
map $G=\operatorname{diag}(g_i)$ gives
$P(V^g)=GP(V)G^{-1}$, and similarly for
$\Pi$ and its connection. At coarse nodes
$G e_i=e_i g_i$, exactly the required
coarse gauge. Coordinate isometries likewise
permute the chosen path averages and
partition. Arbitrary fine frame choices
only introduce the gauges already factored
out. RTC uniqueness thus gives real coarse
gauge and lattice covariance of the total
critical orbit and action.

\begin{proposition}[A physical homogeneous curvature atlas]
\label{prop:cta-curvature-critical-atlas}
For sufficiently small $\kappa$, every
real homogeneous coarse field satisfying
\eqref{cta:coarse-curvature} has an actual
constrained-critical fine orbit at every
$h=L^{-k}$. Its literal mean is the given
coarse field, and its physical fine
curvature and current are bounded by
$C\kappa$, uniformly in depth and volume.
Large holonomies and nontrivial bundle
topology are allowed.

There is a fixed positive datum radius
about every such real field on which the
orbit has a holomorphic germ. These germs
identify on overlaps as above. The total
action therefore defines a holomorphic
function on their union, with real coarse
gauge and lattice covariance. On real
curvature-small data,
\begin{equation}\label{cta:action-bound}
                         |S_h(V)|\le C\kappa^2\mathcal V.
\end{equation}
Local fixed-order analytic responses have
the uniform localized bounds of RTC in
each chosen real trial chart.
\end{proposition}

\begin{proof}
The trial proposition supplies all the
real reference hypotheses of RTC and
the small literal datum error. Its
actual forced fixed point gives the
critical orbit and its physical bounds.
The overlap argument identifies the
germs, so at any real curvature-small
datum the estimate obtained by using
that datum itself as centre is the
estimate for the glued orbit. Its
fine plaquette action is at most
$Ch^4\kappa^2$; summing
$O(h^{-4}\mathcal V)$ plaquettes proves
\eqref{cta:action-bound}. Localized
responses follow from the RTC fixed
point on the same reference, composed
with the pointwise analytic datum
\eqref{cta:RTC-datum}. The deck-equivariant
geometric construction gives the same
statements at small periods without
any quantum-cover comparison.
\end{proof}

The uniqueness asserted here is the
controlled trial-atlas uniqueness.
It does not classify arbitrary
disconnected fine critical sectors.
The result removes the need to assume
a regular trial for homogeneous
curvature-small data; it does not
construct a mixed-scale proper trial,
an independent-current history core,
or a positive conditional measure.

% END LAYERED EXPANSION background/curvature_trial_atlas.tex

% BEGIN LAYERED EXPANSION background/weighted_curvature_atlas.tex
\subsection{Constant diagonal weights on the curvature atlas}
\label{ch:wca-weighted-curvature-atlas}

The curvature atlas supplies smooth fine trials without a
global small-potential assumption. We keep those trials and
change the six plaquette coefficients. The change of principal
part is bounded from \(W^{2,p}\) to \(L^p\). We first retain
that stronger field norm in the joint inverse, and then solve
the weighted equation about the given trial.

Coarse spacing is one, \(h=L^{-k}\), and the homogeneous
periods are compatible with the fixed odd \(L\).
Let
\begin{equation}
 \mathcal A_h^a(U)=
       \sum_{\mu<\nu}a_\mu a_\nu\sum_x
                  \left(1-\tfrac12\operatorname{tr}U_{\mu\nu}(x)\right),
 \qquad a=(a_1,\ldots,a_4).
 \label{wca:weighted-action}
\end{equation}
The parameter \(a\) is constant in space.
Its complex domain is a polydisc about \(1\);
the real weights in that disc are positive.
The group mean, its straightening, the site restoration
and the literal node correction do not change.

\subsubsection{A joint inverse with two \(L^p\) derivatives}

Fix \(0<\alpha<1\) and a finite
\(p>4/(1-\alpha)\).
Use complete enlarged unit cells \(C_y^+\), with all
stencils and averaging collars.
In the fixed comparison frames of a real regular reference
put
\begin{equation}
 \begin{split}
 \|f\|_{\mathcal F_p}
     &=\sup_y\|f\|_{L^p(C_y^+)},\\
 \|A\|_{\mathcal W_p}
     &=\sup_y\sum_{r=0}^2
                              \|D^rA\|_{L^p(C_y^+)} .
 \end{split}
 \label{wca:Wp-norms}
\end{equation}
The fine measure is \(h^4\), and all mixed covariant
second differences are included.
The fixed enlargement can accommodate the finitely many
nested cutoffs used below.
Discrete Morrey gives
\begin{equation}
                 \|A\|_{\mathcal A_\alpha}
                              \le C_{p,\alpha}\|A\|_{\mathcal W_p}.
 \label{wca:Morrey}
\end{equation}
These are cell Banach norms, not sums over fine coordinates.
On a sufficiently small \(C^{1,\alpha}\) tube the
second-difference norms at the present base and at the
fixed real reference are equivalent: their difference
contains only the connection difference, its first
difference, and at most one difference of the field.

\begin{lemma}[The strong row saddle in \(W^{2,p}\)]
\label{lem:wca-Wp-joint-inverse}
On the homogeneous restriction of the domain of
\cref{prop:scf-strong-current-families}, the field output
of the full row saddle satisfies
\begin{equation}
 \|A\|_{\mathcal W_p}+\|\mu\|_\infty
          \le C_{p,\alpha}
                   \bigl(\|f\|_{\mathcal F_p}+\|b\|_\infty\bigr).
 \label{wca:Wp-saddle-bound}
\end{equation}
It retains a sufficiently small exponential weight
on the cell norms. The constants are uniform in depth
and compatible volume. The present base may have
a rough unweighted Wilson current; it must be in
the stated strong tube, and the independent current
has the smaller strong-family bound.
\end{lemma}

\begin{proof}
We give the second-derivative step in the strong
column proof, rather than inferring it from
\(C^{1,\alpha}\) regularity.
For the positive field primitive, its local energy
columns and coarse row bounds first control
the local \(H^1\) norm and the multiplier.
A localized \(\mathcal F_p\) force is in the
physical \(L^2\) force class because \(p>2\)
and these cells have bounded physical volume.

After adding the full divergence term, the exact
Wilson curl--divergence stencil is a component
covariant Laplacian plus bounded curvature
multiplications and fixed neighboring shifts,
as in \cref{scf:positive-field-columns}.
The row-mass forcing is bounded pointwise by
the energy cell norm of the complete average.
The physical transpose has its actual fine
density, so no energy field is evaluated at
a point.
In the local real comparison frame the scalar
equation has bounded first-order and zeroth-order
coefficients. On the complex tube the same
coefficient bounds hold.
Pairing opposite shifts before dividing by \(h\)
gives these bounds without differentiating curvature.
Possible normal-coordinate current coefficients
carry \(h\); the plaquette bound controls them.
No pointwise estimate for the present unweighted
current is used.

Apply the local scalar estimates on nested complete
boxes. Energy gives \(H^1\). The \(2\) estimate gives
\(W^{2,2}\), hence gradient \(L^3\) by a strict
four-dimensional embedding. The \(3\) estimate then
gives \(W^{2,3}\), hence gradient \(L^8\) and bounded
field. If \(p\le8\), apply the final \(p\) estimate.
If \(p>8\), first apply the \(8\) estimate, obtaining
bounded gradient, and then the desired finite
\(p\) estimate. Retain the \(D^2A\) term in that
last estimate, before the Morrey step in the
strong-field proof. Each prescribed force belongs
to the lower exponents used in this finite
bootstrap. The cutoff commutators contain only
the controlled outer field and first differences.

The slow positive parametrix has the same
\(C_p/m\) force error when its output slot
is supplied with this additional \(W^{2,p}\)
norm. Its local output factors have just been
bounded in that norm. Thus its positive series
still converges after increasing the fixed \(m\).
In the Woodbury correction the field output
again has a positive field inverse on the left.
The current corrections map strong fields to
\(\mathcal F_p\), by
\eqref{scf:strong-current-smallness}
and the subsequent second-constraint estimate.
Their physical adjoints use the finite-\(p\)
force classes of
\eqref{scf:strong-adjoint-families}.
Consequently the current Neumann series retains
the \(W^{2,p}\) output: apply this improved
leftmost field factor, and use the already
proved strong bounds for its remaining factors.
There is no assertion that the restoration
projection by itself preserves three derivatives.

The row lift has the same leftmost field factor.
The normalized coarse row inverse bounds its
multiplier in \(\ell^\infty\), while
\(Q^\sharp\) sends a bounded row multiplier
to the bounded fine force supplied by its
physical column density.
The full row-saddle formula in
\eqref{scf:strong-row-saddle-family} therefore
gives \eqref{wca:Wp-saddle-bound}.

For a spatial weight use the energy columns first
and apply the finite sequence of scalar estimates
on each output cell. Weights on neighboring
enlarged cells have bounded ratio.
No artificial weight is differentiated.
The other primitive and coarse factors have
their two exponential Schur sums.
This proves the weighted assertion.
\end{proof}

We shall use this lemma at a smooth trial \(U_0\),
with independent current zero.
Write
\[
 H_0=F_0''(0),\qquad
       F_0(A)=\mathcal A_h^1(e^{ihA}U_0).
 \]
Remove the artificial row mass from the saddle
by its exact multiplier shift.
If the mass-added saddle returns \((A,\mu)\)
with \(QA=b\), put \(\nu=\mu-a_fb\).
Then the inverse of
\begin{equation}
 \mathcal L_0\binom A\nu
      =\binom{(H_0+DRD^\sharp)A-Q^\sharp\nu}{QA}
 \label{wca:trial-joint-operator}
\end{equation}
has the same bound
\eqref{wca:Wp-saddle-bound}.
Here \(a_f\) is bounded above and below,
since coarse spacing is one.
This is the full saddle; no horizontality
of its trial lift is assumed.

\subsubsection{The principal perturbation and the weighted trial force}

Let \(V_*\) be real with
\(\|V_*(\partial p)-1\|\le\kappa\).
Choose the smooth trial
\begin{equation}
                         U_0=W_h(V_*)
 \label{wca:CTA-trial}
\end{equation}
of \cref{prop:cta-uniform-trial}.
Its local potentials have all fixed
\(C^m\) bounds \(C_m\kappa\), and
its ordinary Wilson current and curvature
are bounded by \(C\kappa\).
All finite read collars are chosen
before \(\kappa\) is decreased.
Holonomies and bundle topology are not restricted.

Put
\[
 F_0^a(A)=\mathcal A_h^a(e^{ihA}U_0),\qquad
       j_0^a=(F_0^a)'(0),\qquad H_{0,a}=(F_0^a)''(0).
 \]

\begin{lemma}[Bounds at the unchanged trial]
\label{lem:wca-trial-weight-bounds}
On a fixed complex weight polydisc,
\begin{align}
 \|j_0^a\|_\infty&\le C\kappa,\notag\\
 \|(H_{0,a}-H_0)A\|_{\mathcal F_p}
                       &\le C|a-1|\|A\|_{\mathcal W_p},\notag\\
 \|(F_0^a)'(A)-j_0^a-H_{0,a}A\|_{\mathcal F_p}
                       &\le C\|A\|_{\mathcal W_p}^2 .
 \label{wca:trial-weight-bounds}
\end{align}
The last remainder has Lipschitz constant \(Cr\)
on a radius-\(r\) ball.
The maps are holomorphic on fixed larger
product domains, uniformly in \(h\) and volume.
\end{lemma}

\begin{proof}
For each plane multiply the exact paired-current
formula \eqref{cta:paired-current} by
\(a_\mu a_\nu\).
The two plaquettes retain the same constant
coefficient. Their transported difference
is \(O(\kappa h^3)\), including the
commutator from the connecting link.
Division by \(h^3\) and summation over the
three planes incident to a bond give the
first bound. A bound on separate
unpaired plaquettes would lose a factor \(h^{-1}\).

Each plane Hessian at this smooth trial is
a second-difference operator with bounded
principal coefficients and bounded lower-order
coefficients in the fixed local frames.
This also follows directly by differentiating
its paired plaquette expression: covariant
second differences act on the variation,
and differences of trial coefficients have
the bounds \eqref{cta:sampled-potential}.
Thus its \(L^p\) output is bounded by
the three terms in \eqref{wca:Wp-norms}.
The difference of the weights is
\(a_\mu a_\nu-1\).
There are six planes, which proves
the second bound.
It is not a bounded perturbation
on \(C^{1,\alpha}\) alone.

For the final bound apply the centered
Wilson residual calculation
\cref{lem:cng-Wilson-residual}
separately to those six planes.
Constant coefficients \(a_\mu a_\nu\)
change only its constants.
Subtract both the constant and the full
linear jet before estimating.
The remainder is bounded pointwise by
\(C\|A\|_{\mathcal A_\alpha}^2\).
Use \eqref{wca:Morrey} and the cell volume
to obtain its \(\mathcal F_p\) bound.
Cauchy on the larger strong ball gives
the Lipschitz estimate.
\end{proof}

Define \(\mathcal L_{0,a}\) by
\eqref{wca:trial-joint-operator} with
\(H_{0,a}\) in place of \(H_0\).
The perturbation affects only its first
component. By the second bound in
\eqref{wca:trial-weight-bounds},
\begin{equation}
 \|\mathcal L_0^{-1}
                  (\mathcal L_{0,a}-\mathcal L_0)\|
          _{\mathcal W_p\oplus\ell^\infty\to
                            \mathcal W_p\oplus\ell^\infty}
                                      \le C|a-1|.
 \label{wca:principal-Neumann-bound}
\end{equation}
Choose the weight disc so that this
quantity is at most \(1/4\).
Its Neumann series gives the actual
\(\mathcal L_{0,a}^{-1}\), with uniform
bound and a sufficiently small exponential
cell weight. This choice is independent
of depth and volume.

\subsubsection{The forced weighted equation}

Let \(V_*^{\rm tr}=M^kU_0\).
For target \(V_B=e^{iB}V_*\), use the actual
literal offset and its reduced datum
\begin{equation}
 \beta_*(B)=i^{-1}\log\{V_B(V_*^{\rm tr})^{-1}\},
 \qquad \mathsf b_*(B)=b_{U_0}(\beta_*(B)),
 \qquad
 \|\mathsf b_*(B)\|_\infty\le C_L(\kappa+\|B\|_\infty).
 \label{wca:actual-datum}
\end{equation}
Here \(\beta_*\) is \eqref{cta:RTC-datum},
and \(b_{U_0}\) is the actual datum map
\eqref{pnc:actual-reduced-datum}.
The relative product inside the logarithm is
near one. Neither individual coarse transport
needs a small logarithm.

Write \(\mathfrak q_0(A)\) for the actual
re-centered straightened constraint at \(U_0\),
and put
\[
 C_q(A)=\mathfrak q_0(A)-QA,\qquad
       N_a(A)=(F_0^a)'(A)-j_0^a-H_{0,a}A .
 \]
The weighted counterpart of the forced RTC
equation, with approximate multiplier zero, is
\begin{equation}
 \binom A\nu=\mathcal L_{0,a}^{-1}
  \binom{-j_0^a-N_a(A)
                    +(\mathfrak q_0'(A)-Q)^\sharp\nu}
        {\mathsf b_*(B)-C_q(A)}.
 \label{wca:forced-fixed-point}
\end{equation}
Its affine source is the weighted trial current.
It has not been removed by a centering convention.

\begin{lemma}[A weighted correction of the smooth trial]
\label{lem:wca-weighted-trial-correction}
After fixed reductions of the weight disc,
the curvature ceiling and a datum radius \(\rho\),
\eqref{wca:forced-fixed-point} has a unique solution
in a fixed \(\mathcal W_p\oplus\ell^\infty\) ball.
It is holomorphic in \(a,B\), real at real data,
and
\begin{equation}
 \|A(a,B)\|_{\mathcal W_p}+\|\nu(a,B)\|_\infty
                         \le C(\kappa+\|B\|_\infty).
 \label{wca:trial-correction-bound}
\end{equation}
It solves the full weighted critical equation
after the gauge penalty is removed.
The specified literal correction gives mean
exactly \(V_B\).
\end{lemma}

\begin{proof}
The actual constraint estimates give
\[
 \|C_q(A)\|_\infty\le C\|A\|_{\mathcal A_\alpha}^2,\qquad
 \|(\mathfrak q_0'(A)-Q)^\sharp\nu\|_\infty
              \le C\|A\|_{\mathcal A_\alpha}\|\nu\|_\infty .
 \]
These are the physical column-density bounds
in \cref{lem:cng-constraint-residuals},
before any restoration adjoint.
The fixed-\(L\) hierarchical density, not a
single direct large-block contour, supplies
their uniformity in depth.
Use \eqref{wca:Morrey} for the field
arguments and the cell volume for the force.
Together with \eqref{wca:trial-weight-bounds},
the nonlinear tuple in
\eqref{wca:forced-fixed-point} has size \(Cr^2\)
and Lipschitz constant \(Cr\) on a joint
radius-\(r\) ball.

Choose \(r\) inside the strong mean, site
and gauge-chart margins so that the
inverse bound times this Lipschitz constant
is below \(1/4\).
Include the SCF background tube among
these margins. The trial qualifies by
its unweighted bounds from CTA.
Choose \(\rho,\kappa\) afterward so the
inverse times the two affine sources
has norm below \(r/8\).
Require also \(C(\kappa+\rho)\) below
the smaller independent-current ceiling
of that strong-family tube.
Iteration from zero is invariant and
contractive. Normal convergence on
the larger product domain proves
holomorphy and \eqref{wca:trial-correction-bound}.
The construction used the inverse at the
given trial; no inverse at the desired
critical field was assumed.

The equation before penalty removal is exactly
\[
 (F_0^a)'(A)+DRD^\sharp A
                  =\mathfrak q_0'(A)^\sharp\nu,\qquad
                  \mathfrak q_0(A)=\mathsf b_*(B).
 \]
The actual gauge cocycle, site extensions
and gauge-test surjectivity in RTC depend
on the trial geometry and its strong
regularity. They do not depend on the
plaquette coefficients.
The trial already has the needed
unweighted regularity.
Test this equation on the actual fine
gauge directions in the kernel of its
effective coarse gauge. The first term
vanishes by gauge invariance of
\(\mathcal A_h^a\), and the constraint term
vanishes by the exact intertwining.
The site surjectivity used in
\eqref{pnc:gauge-test-surjectivity}
therefore annihilates the penalty.
For complex parameters use the same
range and bilinear nondegeneracy
argument as RTC, not positivity
of a complex square.
Thus \(RD^\sharp A=0\), and
the full equation is
\[
                  (F_0^a)'(A)=\mathfrak q_0'(A)^\sharp\nu .
 \]

The actual mean slice and literal node
extension are unchanged.
They convert the reduced datum in
\eqref{wca:actual-datum} to the prescribed
literal mean \(V_B\).
The extension has its fixed strong
site bound and gives a strong literal
field of size \(C(\kappa+\|B\|)\).
No \(W^{2,p}\) claim for an arbitrary
choice of literal extension is needed.
\end{proof}

At the literal output, denote by \(j_a\)
the weighted left current and by
\(\omega_a\) its physical row multiplier
in the declared coarse coordinates.
The row-coordinate conversion is the
actual derivative of the datum map and
the normal-to-left conversion of
\eqref{pnc:literal-physical-multiplier};
the trial multiplier \(\nu\) is not
silently identified with this covector.
The actual transpose density and
bounded divided exponential differential
give
\begin{equation}
 j_a(U_h^a(V_B))=Q^\sharp\omega_a,\qquad
 \|\omega_a\|_\infty+\|j_a(U_h^a(V_B))\|_\infty
                  +\|F(U_h^a(V_B))\|_\infty
                         \le C(\kappa+\|B\|_\infty).
 \label{wca:physical-weighted-current}
\end{equation}
Here \(Q\) is re-centered at the output
and the physical fine units are retained.
The first equality follows from actual
weighted stationarity and the admitted
coordinate conversions.
The current estimate uses that equality,
not an estimate of a second difference
of the corrected potential.
Its unweighted Wilson current is not
asserted pointwise bounded.

\begin{lemma}[Uniqueness on a strong trial neighborhood]
\label{lem:wca-strong-orbit-uniqueness}
For the same fixed smaller domains, the weighted
critical orbit is unique among orbits with
a sufficiently small \(C^{1,\alpha}\) representative
relative to the trial and with the same literal datum.
No \(W^{2,p}\) assumption on a competing
representative or point bound on its unweighted
current is required.
\end{lemma}

\begin{proof}
Put a competing orbit into the actual trial
mean slice by the unchanged site construction.
Its relative potential has strong norm at
most \(C\varepsilon\).
Let \(\mathcal H_{0,a}\) be the full row lift
of \(\mathcal L_{0,a}\).
It satisfies \(Q\mathcal H_{0,a}=1\).
The inverse
\[
 \mathscr H_A=\mathcal H_{0,a}
                 (\mathfrak q_0'(A)\mathcal H_{0,a})^{-1}
 \]
exists on a small strong ball, by the
actual row estimates. Its normalized
columns have the exponential energy
bounds of the full trial lift.

Pair stationarity with \(\mathscr H_A b_{c,e}\).
The trial term \(j_0^a\) gives \(C\kappa\).
For the remaining term use the weak
Wilson readout in
\eqref{pnc:Wilson-weak-readout},
separately in each weighted plane.
Its integration by parts pairs the
first field difference with the energy
column. Thus it is bounded by
\(C\|A\|_{\mathcal A_\alpha}\);
no derivative of the actual current
is evaluated pointwise. Consequently
\[
              \|\nu\|_\infty
                        \le C(\kappa+\|A\|_{\mathcal A_\alpha}).
 \]
This is the weighted version of
\eqref{rtc:competing-multiplier-readout}.

Every finite-array competing field has a
finite \(W^{2,p}\) norm, though it has not
yet been uniformly bounded.
Apply the exact inverse
\(\mathcal L_{0,a}^{-1}\) to its equation.
The nonlinear forces are bounded in
\(\mathcal F_p\) by its strong norm and
the multiplier bound just obtained.
It follows that
\[
 \|A\|_{\mathcal W_p}+\|\nu\|_\infty
            \le C\{\kappa+\|B\|_\infty
                                   +\varepsilon^2+\kappa\varepsilon\}.
 \]
Choose the strong uniqueness neighborhood
and then the working \(\rho,\kappa\)
so this lies inside the larger contraction
ball. The competing field is therefore
the same slice solution.
The slice and literal extension identify
precisely the top-fixed gauge orbit.
All these norm estimates and inverses
hold on the smaller complex domains,
so the same uniqueness applies there.
\end{proof}

\subsubsection{The weighted physical atlas}

\begin{proposition}[A diagonal deformation of the curvature atlas]
\label{prop:wca-weighted-curvature-atlas}
There are a curvature ceiling \(\kappa_*>0\),
a weight polydisc \(\mathcal A_+\) about \(1\),
and a datum radius \(\rho>0\), independent of
depth and compatible volume, with the following properties.
Every real coarse \(V\) with
\(\|V(\partial p)-1\|\le\kappa\le\kappa_*\)
has an actual weighted constrained-critical
fine orbit \(U_h^a(V)\).
For real positive \(a\) in a fixed smaller
cube, its literal mean is \(V\), and
\begin{equation}
 \|F(U_h^a(V))\|_\infty+
 \|j_a(U_h^a(V))\|_\infty+
 \|\omega_a(V)\|_\infty\le C\kappa .
 \label{wca:atlas-physical-bounds}
\end{equation}
It has a jointly holomorphic germ in \(a\)
and \(\|B\|_\infty<\rho\) about each real
coarse centre \(V_B=e^{iB}V\).
The germs identify as physical orbits on
controlled overlaps and define a holomorphic
total action
\[
                         S_h^a(V)=\mathcal A_h^a(U_h^a(V)).
 \]
For real data,
\begin{equation}
                        0\le S_h^a(V)\le C\kappa^2\mathcal V .
 \label{wca:weighted-action-bound}
\end{equation}
Large holonomies and nontrivial bundle topology
are allowed. At fixed \(a\) there is real coarse
gauge, translation and reflection covariance;
axis permutations also permute \(a\).
\end{proposition}

\begin{proof}
The smooth trial of CTA is independent of \(a\).
\cref{lem:wca-trial-weight-bounds} supplies its
weighted force, and
\cref{lem:wca-weighted-trial-correction}
corrects the actual datum error.
At the real centre \(B=0\), both errors
are \(O(\kappa)\).
Equation \eqref{wca:physical-weighted-current}
gives \eqref{wca:atlas-physical-bounds}.

For overlaps use the same Kato alignment of
the trial bundles as CTA.
The aligned trial connections differ by
\(C\rho\) in every fixed local derivative
norm, by \eqref{cta:aligned-connection}
and \eqref{cta:aligned-trial-bound}.
The two weighted corrections have strong
size \(C(\kappa+\rho)\).
Choose these constants inside the common
strong uniqueness neighborhood of
\cref{lem:wca-strong-orbit-uniqueness}.
The two solutions have the same literal
datum and weighted equation, so their
orbits coincide.
The actual mean slices and literal gauges
are holomorphic, giving holomorphic transitions.
Their top-fixed cocycle follows from the
free top-fixed gauge action, exactly as in CTA.
No multiplicative cocycle for the Kato map
itself is being assumed.

Real fine-frame changes and coarse gauge
changes conjugate the full equations.
The CTA short-path and partition policy
is equivariant under coordinate isometries.
The weighted action has the stated joint
permutation covariance, and the other
constructions are unchanged. Uniqueness
therefore proves all the covariance claims.

Finally a fine plaquette has
\(1-\tfrac12\operatorname{tr}U_p
                \le C h^4\kappa^2\).
The weights are bounded and positive,
and there are \(O(h^{-4}\mathcal V)\)
plaquettes. This proves
\eqref{wca:weighted-action-bound}.
Use the datum itself as real centre for
this estimate. No global bond gauge close
to a flat connection enters the proof.
\end{proof}

The corrected orbit remains in a fixed smaller
\(C^{1,\alpha}\) tube of its smooth trial.
Thus the existing off-current families may be
evaluated there with a sufficiently small
independent current, in particular with \(j_a\).
They are the operators of their stated convention.
This observation does not identify an unweighted
canonical precision with the weighted physical
Hessian. The latter is treated next.

\subsubsection{The physical weighted Hessian and its row mass}

Fix a real point of the atlas, put \(U=U_h^a(V)\),
and re-center the actual rows and derivatives there.
Let \(\mathscr D_a\) be the full normal Hessian of
\(\mathcal A_h^a\), and write
\(\Delta^{(2)}[\omega_a]\) for the actual
second constraint jet paired with its physical
multiplier. Define
\begin{equation}
 \mathscr H_{L,a}=\mathscr D_a-\Delta^{(2)}[\omega_a].
 \label{wca:physical-Lagrangian-Hessian}
\end{equation}
All these derivatives are in the declared
straightened normal coordinates. Gauge invariance
identifies the scalar action with the literal
coarse action, as in NEG.

We record the gauge-current term as well.
Let \(B_U[j_a]\) be the finite-gauge
curvature operator of the canonical convention.
The first inverse and its first row lift are
\begin{equation}
 G_a=\{\mathscr D_a-B_U[j_a]
                      +DRD^\sharp+Q^\sharp a_fQ\}^{-1},
 \qquad
 H_a=G_aQ^\sharp(QG_aQ^\sharp)^{-1}.
 \label{wca:weighted-first-variant}
\end{equation}
The second inverse is formed afterward:
\begin{equation}
 \mathscr G_a^{-1}
   =\mathscr D_a-B_U[j_a]+DRD^\sharp+Q^\sharp a_fQ
                 -S^\sharp\Delta^{(2)}[H_a^\sharp j_a]S .
 \label{wca:weighted-second-variant}
\end{equation}
The lift in the multiplier readout is the first lift.
The row covariance in the effective Hessian below
will use the second inverse.

\begin{lemma}[The physical identity and energy gap]
\label{lem:wca-physical-energy-gap}
For the real atlas points and the fixed real weight cube,
the two inverses in
\eqref{wca:weighted-first-variant}--\eqref{wca:weighted-second-variant}
exist with uniform energy bounds.
Moreover
\begin{equation}
 \mathscr G_a^{-1}
       =\mathscr H_{L,a}+DRD^\sharp+Q^\sharp a_fQ,\qquad
 c\|v\|_{\mathcal E}^2
       \le\langle v,\mathscr G_a^{-1}v\rangle
       \le C\|v\|_{\mathcal E}^2 .
 \label{wca:physical-second-precision}
\end{equation}
Here the physical homogeneous unit-scale norm is
\[
       \|v\|_{\mathcal E}^2
                 =\|D_Uv\|_{2,h}^2+\|v\|_{2,h}^2,
\]
equivalent to the energy norm in the nonflat
row-saddle theorem.
The fine Lagrangian form also satisfies
\par\noindent
\begin{minipage}{\linewidth}
\begin{equation}
 \langle v,\mathscr H_{L,a}v\rangle
      \ge c\|\mathrm{d}_Uv\|_{2,h}^2-C\kappa\|v\|_{\mathcal E}^2,
 \qquad
 |\langle v,\mathscr H_{L,a}w\rangle|
                        \le C\|v\|_{\mathcal E}\|w\|_{\mathcal E}.
 \label{wca:Lagrangian-energy-form}
\end{equation}
\popQED\par
\end{minipage}
\end{lemma}

\begin{proof}
The Wilson form calculation is applied plane by plane.
The real coefficients \(a_\mu a_\nu\) lie between
fixed positive constants. Its leading curl form
is therefore comparable to \(\|\mathrm{d}_Uv\|_2^2\).
The actual curvature terms have bound
\(C\|F(U)\|_\infty\|v\|_{\mathcal E}^2\).
They are retained in \(\mathscr D_a\).
The finite-gauge current operator has bound
\(C\|j_a\|_\infty\|v\|_{\mathcal E}^2\).
The constraint-curvature form has bound
\(C\|\omega_a\|_\infty\|v\|_{\mathcal E}^2\),
by the actual two-leg estimate.
All three small quantities are \(O(\kappa)\)
by \eqref{wca:atlas-physical-bounds}.

The geometric curl, restoration and row estimate
underlying \eqref{npl:proper-gap} is unchanged.
Combining it with these plane estimates proves
the gap for the first variant.
The actual coarse row lift and its physical
energy seeds give an invertible row Schur matrix.
Since \(QH_a=1\) and \(j_a=Q^\sharp\omega_a\),
its readout is exactly
\[
                         H_a^\sharp j_a=\omega_a .
 \]
The small second correction therefore preserves
the energy gap.

For its exact form use weighted gauge invariance
before simplifying any term. If \(\Pi\) is the
finite gauge chain, then
\(\mathcal A_h^a(e^{ih\Pi(A)}U)
              =\mathcal A_h^a(e^{ihA}U)\).
Twice differentiating gives the same
finite-chain identity as in ACC, with the
weighted physical current:
\[
       \mathscr D_a-B_U[j_a]
                              =S^\sharp\mathscr D_aS .
 \]
The physical stationarity and the exact mean
gauge cocycle give the Lagrangian Ward identity
used in \eqref{acc:physical-second-precision}.
Its algebra is unchanged by constant plaquette
weights: its first gradient is \(Q^\sharp\omega_a\),
and the coarse Ward covector annihilates
the induced coarse gauge. Consequently
\[
 S^\sharp\{\mathscr D_a-\Delta^{(2)}[\omega_a]\}S
                              =\mathscr H_{L,a}.
 \]
Substitution in
\eqref{wca:weighted-second-variant} proves
the equality in
\eqref{wca:physical-second-precision}.
Thus the gauge-current term was used in the
identity; it was not discarded on the
ground that the background is nearly flat.

The plane form estimate and the actual multiplier
bound now give
\eqref{wca:Lagrangian-energy-form}.
The same geometric curl--restoration--row estimate,
with \(\kappa\) decreased, proves the two energy
inequalities in
\eqref{wca:physical-second-precision}.
They concern the actual full constrained
precision. No positive fine Dirichlet
compression has been identified with it.
\end{proof}

Let
\begin{equation}
 \mathcal H_a=\mathscr G_aQ^\sharp
                         (Q\mathscr G_aQ^\sharp)^{-1}.
 \label{wca:physical-row-lift}
\end{equation}
The actual row seeds, their energy bound
and \eqref{wca:physical-second-precision} imply
\[
 Q\mathcal H_a=1,\qquad
 RD^\sharp\mathcal H_a=0,\qquad
             \|\mathcal H_a z\|_{\mathcal E}\le C\|z\|_{2,c}.
 \]
For completeness, the row-seed identity \(QE_Q=1\)
gives an inf-sup lower bound for \(Q^\sharp\)
from coarse \(L^2\) to \(\mathcal E^*\).
The bounded row map gives its upper bound.
The energy inverse then makes
\(Q\mathscr G_aQ^\sharp\) comparable to the
coarse identity, with no fine rank factor.
The middle identity follows from the
physical gauge test, as in NEG.

Twice differentiate the actual constrained graph.
Its first derivative in the fixed horizontal
slice is \(\mathcal H_a\), and its second
derivative satisfies the differentiated constraint.
The first fine action derivative is
\(Q^\sharp\omega_a\). The chain rule therefore gives
\begin{equation}
 \Gamma_{h,a}:=D^2S_h^a(V)
       =\mathcal H_a^\sharp\mathscr H_{L,a}\mathcal H_a
       =(Q\mathscr G_aQ^\sharp)^{-1}-a_f .
 \label{wca:mass-removed-effective-Hessian}
\end{equation}
The second equality follows by adding the
artificial row mass to the differentiated
stationarity equation and using
\(Q\mathcal H_a=1\) and horizontality.
Subtracting \(a_f\) is essential.
Positivity of the mass-added Schur matrix
alone would not give the next estimate.

\subsubsection{The effective curl and the outer Gaussian}

The covariant flux comparison
\cref{neg:covariant-flux} is independent of
the action coefficients:
\begin{equation}
 \|\mathrm{d}_VQv\|_{2,c}
       \le C_L\{\|\mathrm{d}_Uv\|_{2,h}+\kappa\|v\|_{2,h}\}.
 \label{wca:covariant-flux}
\end{equation}
The local gauge proof there needs curvature,
the actual normalized row moments and complete
carriers, not an unweighted current bound.
Alternatively the constructed trial frames
and the strong correction provide the same
small local potential for this atlas.
The comparison retains the flat positive
flux-row norm before its \(O(\kappa)\) error.
It does not count fine plaquettes.

Insert \(v=\mathcal H_a z\) in
\eqref{wca:Lagrangian-energy-form} and use
\eqref{wca:covariant-flux}.
The lift energy bound absorbs the error.
Together with the upper form bound this gives
\begin{equation}
 \Gamma_{h,a}[z,z]\ge
       c\|\mathrm{d}_Vz\|_{2,c}^2-C\kappa\|z\|_{2,c}^2,\qquad
 |\Gamma_{h,a}[z,w]|\le C\|z\|_{2,c}\|w\|_{2,c}.
 \label{wca:effective-curl-bound}
\end{equation}
This is a bound on the mass-removed
physical Hessian.

Use the actual one-step outer graph
\[
          V(u)=e^{i\Phi_V(u)}V,\qquad
          \mathsf C_V=D\Phi_V(0),\qquad
          \Phi_V(u)|_{\rm free}=u.
 \]
Let \(K_G(V)\) be its relative contour-gauge
Hessian. Define the complete composed action
and outer precision by
\begin{equation}
 K_{A,h,a}(V)=
       \mathsf C_V^\sharp\Gamma_{h,a}\mathsf C_V
                  +\langle\omega_a,D^2\Phi_V(0)\rangle,\qquad
 K_{h,a,\tau}(V)=\tau K_{A,h,a}(V)+K_G(V).
 \label{wca:actual-outer-precision}
\end{equation}
The second term is the outer nonlinear
constraint contribution. It is separate
from the inner
\(\Delta^{(2)}[\omega_a]\) in
\eqref{wca:physical-Lagrangian-Hessian}.
Centering the outer action by its linear
free-coordinate jet changes neither term.

\begin{proposition}[A gap throughout the weighted real atlas]
\label{prop:wca-outer-Gaussian-gap}
On a fixed smaller real weight cube and
for sufficiently small \(\kappa\), the actual
outer precision obeys
\begin{equation}
                 c_L I\le K_{h,a,\tau}(V)\le C_L I
 \label{wca:real-outer-gap}
\end{equation}
for real \(\tau\) in a fixed interval about one.
Constants are independent of the inner depth
and compatible volume.
\end{proposition}

\begin{proof}
The finite outer chart has
\(\|\mathsf C_V\|_{2\to2}\le C_L\),
\(\|\mathsf C_Vu\|_2\ge\|u\|_2\),
and its second derivative paired with a
bounded covector has norm at most
\(C_L\|\omega_a\|_\infty\), by
\eqref{neg:outer-chart-bounds}.
Thus \eqref{wca:effective-curl-bound} gives
\[
 K_{A,h,a}(V)[u,u]\ge
       c\|\mathrm{d}_V\mathsf C_Vu\|_2^2-C_L\kappa\|u\|_2^2,
 \qquad \|K_{A,h,a}(V)\|\le C_L .
 \]
The geometric outer curl--gauge inequality
\eqref{neg:outer-geometric-gap-bound}
has no action coefficient in it.
Combine it with this estimate.
The positive real coefficients and
\(\tau\) in a compact positive interval
retain a fixed multiple of that
geometric gap. Decrease \(\kappa\)
to absorb its negative mass error.
This proves the lower bound.
The finite chart, gauge and effective
Hessian bounds give the upper bound.

For a nonaxial real \(V\), the exact
node-fixed block gauge is a real
orthogonal change of free coordinates,
and it carries the relative contour
gauge to the axial one.
Thus the estimate holds in the
literal frame as well.
No global trivialization of the
coarse bundle is used.
\end{proof}

\subsubsection{Complex sectors and the response retained}

\begin{proposition}[A uniform sector near every real centre]
\label{prop:wca-complex-outer-sector}
Fix a real atlas centre \((a_0,V_0,\tau_0)\)
in smaller weight and \(\tau\) intervals.
There are positive radii, uniform in that
centre, depth and compatible volume, such that
the physical analytic continuation satisfies
\begin{equation}
 K_{h,a,\tau}(e^{iB}V_0)
           =K_0^{1/2}(I+E)K_0^{1/2},\qquad
 K_0=K_{h,a_0,\tau_0}(V_0),\qquad
                         \|E\|\le\tfrac12 ,
 \label{wca:complex-sector}
\end{equation}
when \(|a-a_0|,\|B\|_\infty,|\tau-\tau_0|\)
lie in those radii. The matrices are
complex symmetric; the inverse and the
determinant branch are holomorphic.
\end{proposition}

\begin{proof}
Use the common trial chart at \(V_0\).
The weighted fixed point and multiplier
are holomorphic on larger uniform
product domains.
Their Cauchy estimates give uniform
first variations in the strong field
and row norms.
The actual row column moments and
restoration comparisons then bound
their fine energy-operator variations.
In particular, after identifying the
fixed real comparison spaces, the
change of
\(\mathscr H_{L,a}+DRD^\sharp+Q^\sharp a_fQ\)
as an operator \(\mathcal E\to\mathcal E^*\)
is at most
\[
                  C\{|a-a_0|+\|B\|_\infty\}.
 \]
The Wilson estimate here is a derivative
form estimate with bounded coefficient
variations; it does not estimate a
rough unweighted current.
The multiplier term uses its actual
point-small analytic \(\omega_a\).
The unchanged row and gauge families
have their uniform complete-column
bounds on the same strong tube.

The real gap in
\eqref{wca:physical-second-precision}
therefore gives a complex energy
inverse by a Neumann series.
Its row Schur inverse has the
uniform inf-sup margin proved above.
The exact resolvent identity and
\eqref{wca:mass-removed-effective-Hessian}
give
\[
 \|\Gamma_{h,a}(e^{iB}V_0)-\Gamma_{h,a_0}(V_0)\|_{2,c\to2,c}
                      \le C\{|a-a_0|+\|B\|_\infty\}.
 \]
Subtracting the fixed artificial mass
does not alter this variation bound.
The finite outer chart, its second
derivative, the physical covector and
the relative gauge have uniform
analytic bounds as well.
Consequently the difference of the
two matrices in
\eqref{wca:complex-sector} is bounded by
\(C_L(|a-a_0|+\|B\|+|\tau-\tau_0|)\).
Choose the radii below half the
real lower gap.
This proves \eqref{wca:complex-sector}.

It follows, for example, that the
Hermitian real part of this matrix
is at least \(K_0/2\).
This is a sector estimate, not
positivity of an analytic bilinear
square. The finite-dimensional log
branch of \(I+E\) is its convergent
power series, fixing the determinant
and normalized Gaussian branches.
The radii came from operator norms,
so they do not shrink with volume.
\end{proof}

The construction also retains the following
fixed-order response. For a small exponential
cell weight \(w=e^{s\psi}\), with \(\psi\)
one-Lipschitz, its trial-chart first variation
satisfies
\begin{equation}
 \|D_BA[b]\|_{\mathcal W_p;w}
            +\|D_B\nu[b]\|_{\infty;w}
                       \le C\|b\|_{\infty;w}.
 \label{wca:weighted-first-response}
\end{equation}
Differentiate \eqref{wca:forced-fixed-point}.
Its weighted linear inverse is bounded,
and its weighted nonlinear derivative is
a contraction after a fixed further
reduction of the working ball.
The actual mean vertices have fixed
physical carriers, so the weight ratios
within them are bounded.
The same fixed-order Cauchy argument as
RTC bounds mixed derivatives with one
weighted datum slot and the other slots
unweighted, by constants \(C_n\).
The literal node extension and physical
weighted current inherit their strong
and point-force versions.
The geometric sum of the intermediate
physical path lengths keeps their
read stars uniform in depth.

The parameter \(a\) is global. Cauchy on
its larger fixed polydisc controls its
derivatives and their mixed datum
responses; no spatial location is
assigned to a weight derivative.
In particular at a real curvature
centre the first weight variation
of the trial-chart field and the
physical multiplier is \(O(\kappa)\).
These are fixed-order localized
responses. An all-degree connected
carrier expansion for this weighted
nonflat atlas is not inferred from
them.

The results concern the homogeneous curvature
atlas and its actual physical Gaussian.
They construct neither a mixed-scale trial
from arbitrary proper data nor independent
background cube discs, a history tower,
or a quantum identity between covers.

% END LAYERED EXPANSION background/weighted_curvature_atlas.tex

% BEGIN LAYERED EXPANSION background/weighted_connected_germ.tex
\subsection{Constant diagonal weights in the selected connected germ}
\label{ch:wcg-weighted-connected-germ}

We deform the selected off-current equation, including its
actual row inversions, by constant diagonal plaquette weights.
The principal change requires two finite-\(p\) derivatives
of the auxiliary field. It does not require two derivatives
of the base potential or of its input perturbations. The
finite gauge chain and the literal output retain their
\(C^{1,\alpha}\) regularity. These distinctions permit the
physical weighted curvature atlas to enter the construction.

\subsubsection{Fixed strong base data and the stronger output space}

Fix \(0<\alpha<1\), a finite \(p>4/(1-\alpha)\),
and a finite requested carrier rate \(\kappa\). Work on an
admitted SCF and NGC determining geometry with its complete
carriers, separation reserves and true bare boundary.
The reference \(U_*\) is a given real current-regular
field with strict margins. The base has the form
\[
                  U=e^{i\eta E}U_* ,
\]
in a smaller common \(C^{1,\alpha}\) tube, and \(J\)
is a bounded independent current. The full amplitude of
\(J\), rather than only its touching restriction, is used
here. No point bound for the actual unweighted or weighted
Wilson current at the present complex base is assumed.
All comparison frames, smooth reconstructions and the
literal node extension are fixed at \(U_*\).

The constant parameter is \(a=(a_1,\ldots,a_4)\), near
\(1\). Replace each \(\mu\nu\) Wilson plaquette
coefficient by \(a_\mu a_\nu\). Write
\begin{equation}\label{wcg:weighted-normal-action}
 F_U^a(A)=\mathcal A_{\Omega_0}^a(e^{i\eta A}U),\qquad
 j_U^a=(F_U^a)'(0),\qquad H_U^a=(F_U^a)''(0).
\end{equation}
The normalization is that of
\eqref{wca:weighted-action} and the actual normal
coordinates of CNG. The averaging map, site restoration,
row mass \(a_f\), and finite gauge chain are unchanged.
In particular \(a_f\) is not a component of \(a\).
Every adjoint below is the physical bilinear transpose;
analytic parameters are not conjugated.

Use the force space \(\mathcal F_p\), field strong
space \(\mathcal A_\alpha\), row datum space
\(\mathcal B\), and multiplier space \(\mathcal W\)
of SCF. On a complete enlarged scale-\(\xi_y\) cell set
\begin{equation}\label{wcg:Wp-cell-norm}
 \|A\|_{\mathcal W_{p,y}}
    =\sum_{r=0}^2\xi_y^{\,1+r-4/p}
                        \|D_*^rA\|_{L^p(B_y^+)} .
\end{equation}
All mixed second differences are included. The measure is
the physical fine measure. Fixed further cell enlargements
accommodate the finite nested estimates. On a bare cell
the rescaled mesh is one and this is a finite atomic norm.
Uniform discrete Morrey and fixed smooth reconstruction give
\begin{equation}\label{wcg:Wp-embedding}
       \mathcal W_p\longrightarrow\mathcal A_\alpha,
                \qquad \|A\|_{\mathcal A_\alpha}
                              \le C\|A\|_{\mathcal W_p}.
\end{equation}
Changing \(D_*\) to the covariant difference at a base
in the stated tube gives an equivalent norm: the difference
of the second differences contains only the connection
difference, its first difference, and at most one difference
of \(A\). It uses no second difference of \(E\).

Put \(\mathcal X_p=\mathcal W_p\oplus\mathcal W\)
and \(\mathcal Z=\mathcal F_p\oplus\mathcal B\).
Family norms are the row and each-input column sums of
\eqref{ngc:multilinear-family-norm}, using these cell
Banach spaces. They include the supremum of each coefficient
on its larger common local base domain before the carrier
sum. Full tags, short joins and read collars are charged
with \(\sigma_c(X)=b_L+d_c(X)\). Choose the primitive
families at a stronger rate \(\kappa_+\), reserving
the costs of supported base slots and their output counts.
The constant \(a\) is a global analytic parameter; it
is not a spatial input tag.

\subsubsection{The plane Hessian on a complex strong base}

\begin{lemma}[A local principal perturbation]
\label{lem:wcg-plane-principal-family}
Let \(V_a=H_U^a-H_U^1\). It has a local holomorphic
coefficient family satisfying
\begin{equation}\label{wcg:principal-family-bound}
 \|\mathbf V_a\|_{\mathcal W_p\to\mathcal F_p;\kappa_+}
                        \le C\max_\mu|a_\mu-1|.
\end{equation}
This is a coefficientwise local-data supremum bound on the
complex strong tube. The base variable is \(E\) in
\(C^{1,\alpha}\), not in \(W^{2,p}\). No current
derivative or point bound for \(j_U^a\) is required.
\end{lemma}

\begin{proof}
Rescale a nonbare cell to length one, with mesh
\(h=\eta/\xi_y\). Apply the exact four-edge insertion
formula \eqref{px:plaquette-hessian-exact} to one
orientation plane. In the same normalization as
\eqref{psr:Wilson-stencil}, its assembled Hessian is
\begin{equation}\label{wcg:plane-stencil}
 H_{U,\mu\nu}=\mathrm{d}_{U,\mu\nu}^{\sharp}\mathrm{d}_{U,\mu\nu}
                                      +Z_{U,\mu\nu}.
\end{equation}
Here \(\mathrm{d}_{U,\mu\nu}\) is the corresponding component
of the covariant curl; \(Z\) is a finite-shift kernel
on a complete two-step collar, with both coefficient sums
at most \(C\sup_p|P_p-I|/h^2\).

Indeed the exact plaquette insertion differs from the
covariant curl by a matrix \(E_p\) of size
\(C|P_p-I|/h\). Expand
\[
 (\mathrm{d}+E_p)^\sharp(\mathrm{d}+E_p)-\mathrm{d}^\sharp \mathrm{d}
                  =\mathrm{d}^\sharp E_p+E_p^\sharp \mathrm{d}+E_p^\sharp E_p.
\]
Multiply their finite stencil entries, whose curl entries
have size \(C/h\). The first two products have size
\(C|P_p-I|/h^2\); the last is smaller on the stated
plaquette domain. The remaining insertion term has that
same bound. This is multiplication of shifted matrices,
not differentiation of the curvature. The incoming and
outgoing stencil counts are bounded. The argument works
for each plane before any planes are added.

In the complex strong tube the rescaled plaquette satisfies
\(|P_p-I|\le Ch^2\). Each term of
\(\mathrm{d}_{U,\mu\nu}^{\sharp}\mathrm{d}_{U,\mu\nu}A\) is a
transported mixed second covariant difference, up to bounded
lower-order shifts. It is bounded in cell \(L^p\) by
the displayed \(W^{2,p}\) norm. Thus
\(H_{U,\mu\nu}:\mathcal W_p\to\mathcal F_p\)
has a uniform local bound. Restoring the scale gives
\(\xi_y^{3-4/p}\) on its force, exactly the second
derivative term in \eqref{wcg:Wp-cell-norm}.

This reasoning also explains the harmless normal-coordinate
current coefficients. The single-link current is a fixed
sum of plaquette insertions divided by \(h^3\), so
\(h|j_U|\le C\sup_p|P_p-I|/h^2\), even if
\(|j_U|\) is not bounded. Terms with that coefficient
are bounded finite-shift terms; no derivative of it is
taken. On true bare stencils the rescaled mesh is one
and the exact finite matrices give the same conclusion.

Relative to the fixed real frame, the first covariant
difference changes by a bounded local multiplication and
shift depending analytically on \(E\). Its first
difference depends only on \(E,DE\). Consequently the
second-difference operator in this proof is holomorphic
from the strong base ball into
\(\mathcal W_p\to\mathcal F_p\), with a uniform
bound on its larger local domain. Fixed smooth slot
reconstruction uses derivatives only through order two
on the field slot and through order one on the base slot.

Finally
\[
                 V_a=\sum_{\mu<\nu}(a_\mu a_\nu-1)H_{U,\mu\nu}.
\]
There are six planes. Assign the bounded local stars,
endpoint tags and actual background collar before taking
the coefficient norm. Their local-data suprema obey the
bounds just proved, so their finite row and column sums
give \eqref{wcg:principal-family-bound}.
\end{proof}

\subsubsection{Two derivatives in the unweighted primitive families}

\begin{lemma}[A coefficientwise stronger output]
\label{lem:wcg-Wp-linear-family}
The unweighted families needed for \(G_0\), the positive
normalized row inverse \(\mathcal T_0^{-1}\), and the
strong kinematic maps may be chosen so that
\begin{equation}\label{wcg:unweighted-Wp-family}
               \|\mathbf G_0\|_{\mathcal F_p\to\mathcal W_p;\kappa_+}
                                                          \le C.
\end{equation}
The same local-data supremum convention holds. Kinematic
maps have the spaces stated in SCF; in particular
\(\Lambda:\mathcal A_\alpha\to\mathcal Y_{2,\alpha}\)
and \(S:\mathcal A_\alpha\to\mathcal A_\alpha\).
No \(W^{2,p}\) assertion about \(S\) is made.
\end{lemma}

\begin{proof}
Retain the second differences at the output of the local
positive field inverse in
\cref{scf:positive-field-columns}. Its energy column and
complete row average bound first give local \(H^1\).
The component Laplacian equation then gives \(W^{2,2}\),
gradient \(L^3\), next \(W^{2,3}\), and gradient
\(L^8\) with bounded field. If \(p\le8\), finish
at \(p\); for \(p>8\), use \(8\) first and then
\(p\). All cutoffs lie in a fixed larger complete
output collar. Coefficient bounds in the comparison
frame use only its \(C^{1,\alpha}\) potential. The
complex tube has the same bounds. Before Morrey, this proves
\[
 \|\chi I_D\chi\,1_z f\|_{\mathcal W_{p,y}}
                  \le Ce^{-s d(y,z)}\|f\|_{\mathcal F_{p,z}}.
\]
This does not assert global \(W^{2,p}\) regularity for
an uncut Dirichlet solution across its artificial boundary.

The slow commutator still maps this output into
\(\mathcal F_p\) with norm \(C_p/m\): its terms
contain \(D\chi\,Dv\), \(D^2\chi\,v\), complete
row differences, and finite-shift curvature terms with
one cutoff difference. None uses a third output difference.
The weighted centre sums in the positive local walk
therefore give a normally convergent family
\(\mathbf G_F:\mathcal F_p\to\mathcal W_p\).
The column and error bounds are uniform on the larger
local coefficient domains, so taking their suprema
before summing leaves the same positive majorants.

The actual Woodbury identity is
\[
 G_0=G_F+G_F(D\widehat F)\widehat\Theta^{-1}
                                      (D\widehat F)^\sharp G_F.
\]
Keep the improved output at the leftmost \(G_F\).
The site column \(D\widehat F\) is a bounded force
of amplitude three. Its transpose on a strong field
has a bounded counting-row pairing, since
\(\xi^{-3}\xi^{-1}\xi^4=1\). All remaining factors
use their existing strong or energy bounds. This proves
\eqref{wcg:unweighted-Wp-family} without differentiating
a projection kernel. Unit counting field-row injection
has force density of amplitude three by
\eqref{cci:row-force-cell-injection}, so the improved
output also applies at each required row source.

For completeness the positive coarse inversions retain
the stronger local-data convention of
\cref{lem:ngc-local-data-majorants}. A combined short
power of length \(m\) has each entry at most
\(C(q')^m\) uniformly on the larger coefficient tube.
Its declared endpoint ball has polynomially many possible
endpoints. Summing their separate suprema costs
\(C(1+B_m)^4(q')^m\), not the smaller Hilbert-row
bound available at one common base point. Choose the
finite power and long cutoff using this polynomial;
exponential decay still makes the norm small. Reblocking
and the prescribed rate reserves proceed before any
weight or current radius is chosen. The exact positive
coarse matrices and their gaps have not changed.
This proves the required coefficientwise bound through
\(\widehat\Theta^{-1}\) and \(\mathcal T_0^{-1}\).
\end{proof}

\subsubsection{Weighted inverses in the prescribed order}

Write \(B_U[J]\) for the unchanged finite-gauge current
operator, and let \(S=1-D\Lambda\). Define complete
operators in the following order:
\begin{align}
 (G_0^a)^{-1}&=H_U^a+DRD^\sharp+Q^\sharp a_fQ,
                                                   \label{wcg:G0-definition}\\
 (G^a)^{-1}&=(G_0^a)^{-1}-B_U[J],
 \qquad H^a=G^aQ^\sharp(QG^aQ^\sharp)^{-1},
                                                   \label{wcg:first-definition}\\
 \omega^a&=(H^a)^\sharp J,\qquad
 C_a=S^\sharp\Delta_U^{(2)}[\omega^a]S,
                                                   \label{wcg:omega-definition}\\
 (G_1^a)^{-1}&=(G^a)^{-1}-C_a,\qquad
 H_1^a=G_1^aQ^\sharp(QG_1^aQ^\sharp)^{-1}.
                                                   \label{wcg:second-definition}
\end{align}
The first lift inside \(\omega^a\) is fixed before
the second inverse is formed.

\begin{proposition}[Weighted full row-saddle families]
\label{prop:wcg-weighted-linear-family}
On a fixed smaller base tube, weight polydisc and full
independent-current ball, the operators above have exact
holomorphic local families. The field inverses map
\(\mathcal F_p\) to \(\mathcal W_p\), and the
row lifts map \(\mathcal B\) to \(\mathcal W_p\).
The second full saddle has a bounded family
\begin{equation}\label{wcg:weighted-row-saddle}
 \mathfrak I^a_{U,J}:\mathcal Z\longrightarrow\mathcal X_p,
 \qquad
 (G_1^a)^{-1}A-Q^\sharp\nu=f,\quad QA=d.
\end{equation}
Its coefficientwise local-data supremum norm at the
reserved rate is bounded uniformly in depth, compatible
volume and admitted shape. The rooted multiplier family
has amplitude at most \(C\|J\|_{\infty,3}\).
\end{proposition}

\begin{proof}
First use the genuinely principal family from
\cref{lem:wcg-plane-principal-family}:
\begin{equation}\label{wcg:principal-Neumann-family}
 \mathbf G_0^a=\mathbf G_0\star
           \sum_{n\ge0}(-\mathbf V_a\star\mathbf G_0)^{\star n}.
\end{equation}
Its loop is on \(\mathcal F_p\); the field output
of each \(G_0\) is \(\mathcal W_p\). Choose the
larger weight polydisc so its loop norm is below \(1/4\).
The series is normally convergent and its complete sum is
the exact inverse in \eqref{wcg:G0-definition}.
Moreover \(\mathbf G_0^a-\mathbf G_0\) has norm
\(C|a-1|\). This step does not assume invertibility
of a new local weighted compression.

SCF gives
\(\mathbf B_U[J]:\mathcal A_\alpha\to\mathcal F_p\)
of norm \(C\gamma\), \(\gamma=\|J\|_{\infty,3}\),
with no derivative of \(J\). Use
\eqref{wcg:Wp-embedding} at its input. Then
\begin{equation}\label{wcg:first-Neumann-family}
 \mathbf G^a=\mathbf G_0^a\star
          \sum_{n\ge0}(\mathbf B_U[J]\star\mathbf G_0^a)^{\star n}
\end{equation}
converges after decreasing \(\gamma\). It retains
the stronger output at its leftmost factor.

The physical row normalization is kept explicit:
\begin{equation}\label{wcg:row-normalization}
 \mathcal Q_n=\Xi_\Gamma^2a_f^{1/2}Q,\qquad
 \mathcal Q_n^\sharp=Q^\sharp a_f^{1/2}\Xi_\Gamma^{-2}.
\end{equation}
The first map has counting-row output, and the second
is its counting-to-physical transpose. Put
\(\mathcal T_0=\mathcal Q_nG_0\mathcal Q_n^\sharp\).
Its inverse is already constructed. For
\(\mathcal T^a=\mathcal Q_nG^a\mathcal Q_n^\sharp\)
the complete local families satisfy
\[
      \|\mathbf{\mathcal T^a}-\mathbf{\mathcal T_0}\|
                                           \le C(|a-1|+\gamma).
\]
Decrease the weight disc and then the current radius
so the following exact family converges:
\begin{equation}\label{wcg:row-inverse-family}
 \mathbf{(\mathcal T^a)^{-1}}
 =\mathbf{\mathcal T_0^{-1}}\star
   \sum_{n\ge0}
    \bigl[-(\mathbf{\mathcal T^a}-\mathbf{\mathcal T_0})
                         \star\mathbf{\mathcal T_0^{-1}}\bigr]^{\star n}.
\end{equation}
Thus \(H^a\) is obtained by actual row inversion,
not by inferring a new gap from a completed operator bound.
The physical lift is
\(G^a\mathcal Q_n^\sharp(\mathcal T^a)^{-1}
                    \Xi_\Gamma^2a_f^{1/2}\).

Pair its current coefficients in that order. A datum
\(b_c=\xi_c^{-1}e\) at one row obeys
\[
          \langle H^ab_c,J\rangle=\xi_c^3\omega_c^a\cdot e.
\]
The same absolute row and column sums as in
\eqref{ncw:rooted-duality} prove
\[
 \sup_c\sum_X e^{\kappa_+\sigma_c(X)}
                              \xi_c^3|\omega_{c,X}^a|\le C\gamma.
\]
Keep the original fine-current bond in each carrier.
The complete bare contribution may be retained here since
the full current is small; its second constraint jet is
zero. This is the first-lift readout throughout.

The actual two-leg row estimate gives
\(\Delta_U^{(2)}[\omega^a]\) as a bounded point-force
family on strong fields. Composing with the SCF
\(S,S^\sharp\) families gives
\(\mathbf C_a:\mathcal A_\alpha\to\mathcal F_p\)
of norm \(C\gamma\). Consequently
\begin{equation}\label{wcg:second-Neumann-family}
 \mathbf G_1^a=\mathbf G^a\star
                   \sum_{n\ge0}(\mathbf C_a\star\mathbf G^a)^{\star n}
\end{equation}
has the required \(W^{2,p}\) output. Its normalized
row matrix differs from \(\mathcal T_0\) by
\(C(|a-1|+\gamma)\). Apply
\eqref{wcg:row-inverse-family} with that difference
to construct its exact inverse and then \(H_1^a\).

In physical row spaces write \(\mathcal S_1^a=QG_1^aQ^\sharp\).
The full saddle is exactly
\begin{equation}\label{wcg:full-saddle-formula}
 \begin{split}
 A&=G_1^af+H_1^a(d-QG_1^af),\\
 \nu&=(\mathcal S_1^a)^{-1}(d-QG_1^af).
 \end{split}
\end{equation}
The row injection is a force of amplitude three;
its normalized density is the one used in the preceding
step. These identities prove the spaces and bounds.
Each Neumann series is dominated by products of positive
coefficientwise suprema on the larger local domains.
There is no exchange of a global supremum with a
conditionally convergent carrier sum.
\end{proof}

The order of choices is important. Fix the rates, \(p\),
\(\alpha\), reference margins and all carrier geometry
first. Establish the stronger primitive constants. Choose
the weight polydisc from the principal and row loop bounds,
then the independent-current radius from both current
loops and both row inversions. Choose the nonlinear datum
and supported base radii last. None depends on fine mesh,
depth or volume. These choices do not identify the carrier
tile with a prescribed history cube or assert that a
one-history-cube separation supplies these larger collars.

\subsubsection{The centered weighted auxiliary equation}

Keep the actual finite gauge map \(\Pi_U\) of
\eqref{cng:finite-gauge-map}, and put
\begin{equation}\label{wcg:nonlinear-kinematics}
 \Xi(A)=\Pi_U(A)-SA,\quad
 N_q(c)=\mathfrak q_U(c)-Qc,\quad r(A)=\mathfrak q_U(\Pi_U(A)).
\end{equation}
The inputs to \(\Pi\) are \(A,\Lambda A,D\Lambda A\).
The identity
\(\operatorname{Ad}(U_b)(\Lambda A)_+
             =(\Lambda A)_- -\eta D\Lambda A\)
cancels its apparent division by \(\eta\) before
expansion. BHR and SCF therefore give
\(\Pi:\mathcal W_p\to\mathcal A_\alpha\), with
\(\Pi'(0)=S\) and \(\Xi(A)=O(\|A\|_{\mathcal W_p}^2)\)
in the strong norm. Two differences of \(\Pi\)
would require a third site difference of \(\Lambda A\).
No such estimate is part of this construction.

With \(\omega^a\) fixed in the \(A\) derivative, define
\begin{equation}\label{wcg:auxiliary-action}
 \begin{split}
 \mathcal L^a_{U,J}(A)={}&F_U^a(A)-F_U^a(0)-\langle j_U^a,A\rangle
                                  -\langle J,\Xi(A)\rangle\\
 &-\langle\omega^a,N_q(\Pi_U(A))\rangle
               +\tfrac12\langle RD^\sharp A,RD^\sharp A\rangle\\
 &+\tfrac12\langle r(A),a_fr(A)\rangle .
 \end{split}
\end{equation}
It is a finite-dimensional holomorphic scalar at each
finite lattice. Its uniformly bounded nonlinear force
families are obtained after centering, as proved next.
The actual \(j_U^a\) is subtracted even if it is rough.

\begin{lemma}[The Hessian and its subtracted force family]
\label{lem:wcg-centered-force-family}
The exact first two auxiliary jets are
\begin{equation}\label{wcg:exact-Hessian}
               (\mathcal L^a_{U,J})'(0)=0,\qquad
               (\mathcal L^a_{U,J})''(0)=(G_1^a)^{-1}.
\end{equation}
Let \(N_{\mathcal L}^a(A)=(\mathcal L^a_{U,J})'(A)
                                      -(G_1^a)^{-1}A\).
The source tuple
\begin{equation}\label{wcg:nonlinear-source}
 \mathcal R^a(A,\nu)=
 \binom{-N_{\mathcal L}^a(A)+(r'(A)-Q)^\sharp\nu}
                   {-(r(A)-QA)}
\end{equation}
has an exact analytic local family into \(\mathcal Z\)
with divided coefficient bounds \(C_0C_1^m\) for every
degree \(m\ge2\). Constant and linear jets are absent
in each filed representative. The coefficientwise larger
local-data norms are uniform on the weight polydisc.
\end{lemma}

\begin{proof}
The current quadratic jet is precisely \(B_U[J]\),
because \(\Pi\) is the selected finite chain. The
constraint-current jet is
\(S^\sharp\Delta_U^{(2)}[\omega^a]S\), since
\(N_q(0)=DN_q(0)=0\). The penalty contributes
\(DRD^\sharp\), and \(r'(0)=QS=Q\) gives
\(Q^\sharp a_fQ\) from the row mass. The remaining
Hessian is \(H_U^a\). These are exactly the terms
in \eqref{wcg:second-definition}, with their stated signs.

For the Wilson force perform the subtraction separately
in each orientation plane, pairing the two incident
plaquettes on an output bond before estimating. The
relative logarithmic calculation of
\cref{lem:cng-Wilson-residual} applies plane by plane.
It gives the local point-force bound
\begin{equation}\label{wcg:weighted-Wilson-residual}
 \|(F_U^a)'(A)-j_U^a-H_U^aA\|_{\infty,3}
                               \le C\|A\|_{\mathcal A_\alpha}^2.
\end{equation}
Here the constant and linear jets were removed before
filing the paired local force, not canceled between
separately estimated carriers. To recall the rough-base
step, in a unit cell write
\(P_h(A)=(ih)^{-1}\log(e^{ihA}e^{iha_0})\), with
\(a_0\) bounded in \(C^{1,\alpha}\). Its remainder
after the linear \(A\) jet has first spatial difference
\(O(hr^2)\) and second difference \(O(r^2)\).
The apparent terms \(hA D^2A\) and
\(h^2 A^2D^2a_0\) use
\(\|D^2A\|_\infty\le Ch^{-1}\|DA\|_\infty\)
and the corresponding estimate for \(a_0\).
The linearization difference multiplying \(D^2A\)
has its factor \(h\). Thus no second base derivative
or base-current amplitude survives. Multiplication by
the six bounded constants \(a_\mu a_\nu\) preserves
the bound and holomorphy. A bounded point force enters
\(\mathcal F_p\) with its physical cell normalization.

For completeness, the remaining force identities must
keep their valid source types. The constraint value is
\begin{equation}\label{wcg:quadratic-constraint}
                r(A)-QA=Q\Xi(A)+N_q(\Pi(A)).
\end{equation}
For the multiplier force one has
\begin{equation}\label{wcg:multiplier-rewrite}
 \begin{split}
 (r'(A)-Q)^\sharp\nu={}&\Xi'(A)^\sharp Q^\sharp\nu
                            +S^\sharp DN_q(\Pi(A))^\sharp\nu\\
                         &+\Xi'(A)^\sharp DN_q(\Pi(A))^\sharp\nu .
 \end{split}
\end{equation}
The row transposes in this expression are bounded point
forces by their actual fixed-leg density. A derivative
of \(\Xi\), paired with such a force, has the exact
adjoint-factory form
\[
                f_0+\Lambda^\sharp h+
                                      \Lambda^\sharp D^\sharp f_d.
\]
Its inputs have grades three, four and three, respectively.
SCF bounds the first two pointwise and the last in finite
\(\mathcal F_p\). It does not put a generic finite-\(p\)
force into a factory that requires a point force.

Let \(C_\omega=\Delta_U^{(2)}[\omega^a]\). The
subtracted constraint-current gradient is exactly
\begin{equation}\label{wcg:constraint-current-rewrite}
 \begin{split}
 &\nabla\langle\omega^a,N_q(\Pi(A))\rangle-S^\sharp C_\omega SA\\
 &\quad=S^\sharp\{C_\omega\Xi(A)
               +DN_q(\Pi(A))^\sharp\omega^a-C_\omega\Pi(A)\}
          +\Xi'(A)^\sharp DN_q(\Pi(A))^\sharp\omega^a .
 \end{split}
\end{equation}
The braces are point forces quadratic in \(A\).
The outside \(S^\sharp\) uses its finite-\(p\)
bound; the last term uses the same point-force factory.
The row-mass remainder is
\begin{equation}\label{wcg:row-mass-rewrite}
 r'(A)^\sharp a_fr(A)-Q^\sharp a_fQA
    =(r'(A)-Q)^\sharp a_fr(A)+Q^\sharp a_f(r(A)-QA).
\end{equation}
Since \(a_f\asymp\xi^{-2}\), its datum input has
the required multiplier amplitude three. Finally remove
the exact quadratic scalar jet
\(\langle A,B_U[J]A\rangle/2\) from the nonlinear
current term before its factory expansion. Its residual
is quadratic in \(A\), with size \(C\gamma\|A\|^2\).
The gauge penalty is exactly linear in the gradient and
has no remaining nonlinear vertex.

All these expressions have at least two joint factors
of \((A,\nu)\) after the displayed subtractions.
Their complete read sets and local strong bounds are
the NGC local vertices. The only new local operation,
the weighted Wilson sum, has been verified above.
Uniform Banach Cauchy bounds on a larger strong ball,
followed by the embedding from \(\mathcal W_p\),
give divided coefficients \(C_0C_1^m\) into
\(\mathcal Z\). Substitution and the SCF factory
bounds preserve their coefficientwise local-data sums.
This proves the claim with the arity zeros already
removed in each representative.
\end{proof}

\subsubsection{The actual selected germ and its connected trees}

The weighted auxiliary equation is
\begin{equation}\label{wcg:auxiliary-system}
           (\mathcal L^a_{U,J})'(A)-r'(A)^\sharp\nu=0,
                       \qquad r(A)=b.
\end{equation}
Using \eqref{wcg:full-saddle-formula}, it is exactly
\begin{equation}\label{wcg:fixed-point}
 x=L^a_Pb+\mathfrak N^a_P(x),\qquad
 L^a_Pb=\mathfrak I^a_{U,J}(0,b),\qquad
 \mathfrak N^a_P=\mathfrak I^a_{U,J}\mathcal R^a,
 \quad x=(A,\nu).
\end{equation}
The nonlinear family is a series
\(\sum_{m\ge2}V^a_{m,P}[x^m]\) with
\(\|\mathbf V^a_{m,P}\|_{\kappa_+}\le C_0C_1^m\).
Its linear output space is \(\mathcal X_p\), while
all finite-gauge vertices use the included strong field.

\begin{proposition}[All-degree weighted selected families]
\label{prop:wcg-connected-germ}
On smaller uniform base, weight and independent-current
domains, there is a unique small solution of
\eqref{wcg:auxiliary-system} with auxiliary field in
\(\mathcal W_p\). It is holomorphic and has exact
connected datum families satisfying
\begin{equation}\label{wcg:datum-family-bound}
 x(P,a,b)=\sum_{n\ge1}\mathbf x_n^a(P)[b^n],\qquad
       \sum_{n\ge1}r^n\|\mathbf x_n^a\|_{\kappa_+}\le Cr
                      \quad(0<r\le r_b).
\end{equation}
Every coefficient retains all output, input and base
read sets, including those of the inverse factors.
The supremum on the larger local base and weight domains
is included inside this family norm.
\end{proposition}

\begin{proof}
The exact fixed point is a contraction on a sufficiently
small \(\mathcal X_p\) ball. For the connected
representation expand it in ordered rooted trees.
A leaf is the exact family \(L^a_Pb\), and an internal
vertex of arity \(m\ge2\) is \(V^a_{m,P}\).
An \(n\)-leaf tree has at most \(n-1\) internal
vertices. Join its carriers at each shared internal
cell, retaining the canonical endpoint tags and fixed
short joins. This includes every background read set,
not just the tree joining datum positions.

The positive scalar majorant is
\begin{equation}\label{wcg:tree-majorant}
          t=C_{\rm lin}r+\sum_{m\ge2}C_0C_1^m t^m .
\end{equation}
Choose \(r_b\) so its map preserves
\([0,2C_{\rm lin}r_b]\) and its derivative there
is below \(1/2\). Iteration from zero gives the
convergent sum of all ordered tree weights, bounded
by \(2C_{\rm lin}r\). Row sums bound every internal
index except along a distinguished input path, where
the corresponding column sums apply. The guarded
carrier union is submultiplicative. It follows that
the actual divided coefficient norms are bounded by
this positive series. No additional factorial enters.

The coefficientwise local-data bounds of the linear
families and vertices were established before these
products. Thus the same positive majorant bounds the
sum of their separate suprema. Normal convergence
permits substitution into the actual equation. Its
sum is in the contraction ball and hence is the
selected solution, not merely a formal expansion.
\end{proof}

\subsubsection{Literal inputs and the bounded current update}

For literal coarse data use the actual map \(b=b_U(B)\)
of \eqref{pnc:actual-reduced-datum}, including
\(Db_U(0)\). Its explicit outward inverse reads the
prescribed finer interior rows without a depth product.
The NGC local analytic family for this map is unchanged
by \(a\). Substitute it into
\eqref{wcg:datum-family-bound}, decreasing the literal
radius so its absolute image lies in the reduced ball.
Put
\[
              C=\Pi_U(A),\qquad \delta\omega=\nu-a_fb_U(B).
\]
The selected point-current update is
\begin{equation}\label{wcg:current-update}
 J_{\rm new}^a=\mathcal E_{\eta C}^{-\sharp}
       \{D\mathfrak q_U(C)^\sharp(\omega^a+\delta\omega)
                                          +J-Q^\sharp\omega^a\}.
\end{equation}
Here \(\mathcal E\) is the same divided exponential
differential as in CNG. Its zero datum jet is \(J\).
For filing the increment use the exact expression
\begin{equation}\label{wcg:current-increment}
 \begin{split}
 J_{\rm new}^a-J={}&(\mathcal E_{\eta C}^{-\sharp}-1)J\\
 &+\mathcal E_{\eta C}^{-\sharp}
       \{DN_q(C)^\sharp\omega^a+
                          D\mathfrak q_U(C)^\sharp\delta\omega\}.
 \end{split}
\end{equation}
Every displayed summand has positive datum degree and
point-force amplitude three. In particular there is no
termwise cancellation of large bare currents in these
estimates; the present independent current is bounded
in its full amplitude norm.

Let \(g\) be the actual node correction with the fixed
extension \(E_n\), and define the literal output by
\begin{equation}\label{wcg:literal-output}
 \widehat U=(e^{i\eta C}U)^g=e^{i\eta\widehat A}U,
            \qquad \widehat J=\operatorname{Ad}(g_-)J_{\rm new}^a.
\end{equation}
The node argument has a connected
\(\mathcal Y_{2,\alpha}\) family. At assigned scale
\(\xi_y\), the actual intermediate node factors
have size \(C_Lr\xi_t/\xi_y\); their absolute sum
is bounded by \(C_Lr\), since
\(\sum_{t\le j_y}\xi_t/\xi_y\le L/(L-1)\).
Their read stars obey the same geometric bound.
Thus literalization has no \(C^k\) accumulation.
Its divided logarithm is a local analytic function of
\(C,\zeta,D\zeta\), where \(g=e^{i\zeta}\).
One field difference uses at most two site differences.

Consequently \(C,\widehat A\) have strong connected
families; \(\nu,\delta\omega\) have multiplier
families, and \(J_{\rm new}^a-J,\widehat J-J\) have
point-force families. All have a function-family bound
\(Cr\) with the same carrier convention as
\eqref{wcg:datum-family-bound}. The auxiliary field
\(A\) retains \(\mathcal W_p\). No \(W^{2,p}\)
bound for the general \(C\) or literal \(\widehat A\)
is inferred from this construction.

\subsubsection{Joint supported-base and global-weight coefficients}

Fix an interior base pair \(P_0=(e^{i\eta E_0}U_*,J_0)\)
and an interior weight \(a_0\). Reconstruct base changes
exactly as in NGC:
\[
 E=E_0+\sum_y\mathcal R_y e_y,\qquad
 J=J_0+\sum_y\mathcal O_y j_y,
\]
using supported strong \(C^{1,\alpha}\) field slots
and bounded current ownership slots. No \(W^{2,p}\)
condition is imposed on \(e_y\). Write
\(a=a_0+\epsilon\), with four global complex components.

\begin{proposition}[Joint connected Taylor bounds]
\label{prop:wcg-joint-connected-family}
On fixed smaller slot, literal datum and weight balls,
every output increment just listed has an exact joint
divided Taylor expansion with
\begin{equation}\label{wcg:joint-family-bound}
  \sum_{n\ge1,\ell,t\ge0}r^n\delta^\ell\zeta^t
              \|\mathbf F_{n,\ell,t}\|_\kappa\le Cr.
\end{equation}
Here \(n\) is datum degree, \(\ell\) is supported
base-slot degree and \(t\) is total degree in the four
global weight variables. The norm has output and all
spatial input row/column bounds; there is no spatial
column assigned to a global weight derivative.
\end{proposition}

\begin{proof}
For a full carrier \(X\), bounded-overlap reconstruction
of its base slots lies in the larger local strong/current
domain. The number of slots it reads is
\(n_X\le C_{L,M_c}\sigma_c(X)\), by the complete
collar counting of \eqref{ngc:base-slot-count}.
Their Banach-valued Cauchy and polarization bounds cost
at most \(C^\ell\) for each distribution of degree.
Summing these distributions at a smaller slot radius
costs \((1-C\delta_1)^{-n_X}\), bounded by
\(\exp(C'\delta_1C_{L,M_c}\sigma_c(X))\).
Choose \(\delta_1\) inside the reserved carrier-rate
budget and take the actual \(\delta\le\delta_1/2\).
The resulting factor \((\delta/\delta_1)^\ell\)
justifies summing degreewise supremum norms.

For a column anchored at a supported base input,
reserve an additional rate larger than
\(\log D_L+1\), where \(D_L\) is the bounded
degree of the tag graph. A possible output at tag
distance \(j\) then has its factor \(e^{-\upsilon j}\),
where \(\upsilon>\log D_L+1\),
which sums against the bounded-degree output count.
This is the same base-column estimate as in NGC;
it does not count fine coordinates within a Banach slot.

The coefficient functions are jointly holomorphic in
the four global weight variables on a larger fixed
polydisc. Weight Cauchy costs
\[
       \sum_{\beta\in\mathbb N^4}
                 (\zeta/\rho_a)^{|\beta|}
                              =(1-\zeta/\rho_a)^{-4}
\]
for a sufficiently small \(\zeta\). This is a fixed
factor, independent of the carrier and its size.
It requires no extra spatial tag. Combining these
bounds with the datum tree majorant proves
\eqref{wcg:joint-family-bound}. Every output increment
vanishes at \(B=0\) for every base and weight, with
that zero removed before filing; hence \(n\ge1\).
Normal convergence identifies the joint series with
the same selected germ and literal update.
\end{proof}

\subsubsection{Physical restriction through the curvature atlas}

Physical matching uses a supplied weighted physical
critical pair and its controlled local physical branch.
It is not a producer for arbitrary proper coarse data.
Suppose first that
\begin{equation}\label{wcg:physical-pair}
            J=j_U^a=Q^\sharp\omega_{\rm phys}
\end{equation}
is within the independent-current radius. Then
\(QH^a=1\) gives \(\omega^a=\omega_{\rm phys}\)
from the first lift. Weighted fine gauge invariance
implies the exact scalar identity
\begin{equation}\label{wcg:physical-scalar-identity}
 \begin{split}
 \mathcal L^a_{U,j_U^a}(A)={}&
       F_U^a(\Pi_U(A))-F_U^a(0)
                             -\langle\omega_{\rm phys},r(A)\rangle\\
 &+\tfrac12\langle RD^\sharp A,RD^\sharp A\rangle
                         +\tfrac12\langle r(A),a_fr(A)\rangle .
 \end{split}
\end{equation}
Indeed \(F_U^a(\Pi_U(A))=F_U^a(A)\), \(QS=Q\),
and the two centered linear/current terms combine to
\(-\langle\omega_{\rm phys},Q\Pi_U(A)\rangle\).
The nonlinear constraint term supplies its remainder.
This is an exact weighted identity, not a comparison
of Hessians alone.

\begin{proposition}[Admission and matching of WCA physical pairs]
\label{prop:wcg-WCA-physical-restriction}
Restrict the WCA weight polydisc, curvature ceiling
and datum radius by the fixed constants of this subsection.
Every real physical WCA pair, and its smaller complex
physical continuation, lies in the present strong-base
and independent-current domain through its smooth real
CTA trial. At that pair the selected weighted germ
and its literal update coincide with the continued
physical WCA orbit and actual weighted left current.
The all-degree connected and joint bounds above apply
to these local perturbation germs.
\end{proposition}

\begin{proof}
At a real coarse centre the CTA trial \(U_*\) has
uniform local all-\(C^m\) potentials and unweighted
curvature/current \(O(\kappa)\). It is a valid real
current-regular reference for SCF, BHR and NGC, even
with large holonomies and nontrivial bundle topology.
WCA puts its literal weighted orbit in a strong tube
of this trial of size \(C\kappa\), and proves
\(j_a(U)=Q_U^\sharp\omega_a\) with both point
amplitudes \(C\kappa\). On a smaller coarse complex
ball the bounds are \(C(\kappa+\rho)\). Choose these
below the present strong and current margins and use
the intersection of the two weight polydiscs. The
unweighted current of the weighted orbit need not
be point-small. The hypotheses concern its strong
distance from the smooth real trial, which WCA supplies.

Consider the supplied physical WCA branch about this
pair. The unchanged actual mean-slice chart puts it
in \(RD^\sharp A_{\rm phys}=0\), with
\(\mathfrak q_U(A_{\rm phys})=b_U(B)\), by its
actual assigned node correction. This is the kinematic
chart of PNC and NGC on the strong tube. Its site
estimates use the bounded local \(C^{1,\alpha}\)
coefficients and completed site families available
there. The finite gauge is the actual chart gauge;
the equation \(q\lambda=0\) alone is not being
treated as an exact top-fixed finite gauge. The given
branch has \(\|A_{\rm phys}\|_{\mathcal A_\alpha}
\le C\|B\|_{\mathcal B}\), by WCA's localized
strong response and this gauge chart.

Its multiplier difference is small in amplitude as
well. Pair its weighted equation with
\(H_1^a(\mathfrak q_U'(A_{\rm phys})H_1^a)^{-1}\).
This is an actual nearby row right inverse. Compare
with the base equation, integrating by parts plane
by plane in the Wilson difference. Only first field
differences and the exponentially summable energy
columns of this lift occur. The variation of the
right inverse has size \(C\|A_{\rm phys}\|_{\mathcal A_\alpha}\),
and the base weighted current is bounded. Thus the
multiplier difference has bound
\(C\|A_{\rm phys}\|_{\mathcal A_\alpha}\).
This is the weak readout used in RTC and WCA,
applied to the difference from the critical base;
there is no nonstationary defect here.

There is no assumed \(W^{2,p}\) estimate for this
competing branch. Each finite array has a finite
such norm. Its physical linearized full precision
is the one in \eqref{wcg:second-definition}: the
weighted finite-chain and Lagrangian Ward identities
are exactly those of \cref{lem:wca-physical-energy-gap}.
Apply \(\mathfrak I^a_{U,j_U^a}\) to its equation
after subtracting this linearization. The weighted
centered Wilson residual and the actual second
constraint remainder give force bound
\(C\|A_{\rm phys}\|_{\mathcal A_\alpha}^2\);
the multiplier product has the same bound by the
weak readout. The datum remainder is quadratic.
More explicitly, put
\(N_F(A)=(F_U^a)'(A)-j_U^a-H_U^aA\) and
\(C_\omega=\Delta_U^{(2)}[\omega_{\rm phys}]\).
In the physical slice the source and row unknown are
\begin{equation}\label{wcg:physical-upgrade-source}
 \begin{split}
 f={}&-N_F(A)+DN_q(A)^\sharp\omega_{\rm phys}-C_\omega A
       +DN_q(A)^\sharp\delta\omega_{\rm phys}
                                     -Q^\sharp a_fN_q(A),\\
 d={}&b_U(B)-N_q(A),\qquad
 \nu=\delta\omega_{\rm phys}+a_fb_U(B).
 \end{split}
\end{equation}
Every force summand is quadratic with the just established
weak multiplier bound. The last force is the retained
artificial row-mass correction; it has not been omitted.
Consequently
\[
 \|A_{\rm phys}\|_{\mathcal W_p}
       +\|\delta\omega_{\rm phys}\|_{\mathcal W}
       \le C\{\|b_U(B)\|_{\mathcal B}
                      +\|A_{\rm phys}\|_{\mathcal A_\alpha}^2\}.
\]
For sufficiently small \(B\) this lies inside the
auxiliary contraction ball. On its exact slice
\(\Pi_U(A_{\rm phys})=A_{\rm phys}\).
Differentiate \eqref{wcg:physical-scalar-identity}:
the physical pair solves the auxiliary system with
\(\nu=\delta\omega_{\rm phys}+a_fb_U(B)\).
The penalty derivative vanishes there. Contraction
uniqueness identifies the auxiliary and physical
slice solutions. The actual node gauge then identifies
their literal representatives.

At the physical pair, \eqref{wcg:current-update}
reduces to the actual normal-to-left differential
conversion of the weighted stationarity equation.
It is therefore the weighted left current of the
updated field. Its gauge transform in
\eqref{wcg:literal-output} is the actual literal
weighted current. All arguments use algebraic
identities and uniform norm estimates on the smaller
complex domains, so holomorphic continuation retains
the matching. No unrelated weighted PNC producer
outside the supplied WCA branch was imported.
\end{proof}

Off current the first reduced-datum jets are
\[
                   D_bA(0)=H_1^a,\qquad D_bC(0)=SH_1^a.
\]
They are not equated. Literal data add \(Db_U(0)\)
and the actual fine node-gauge derivative. On a
physical WCA overlap, both sums represent the same
weighted physical orbit by the exact matching, so
their total action and actual current agree under
the declared gauge transition. Individual filed
coefficients are not asserted gauge invariant or
reference independent.

For small periods perform the primitive construction,
coarse inverse, nonlinear vertices and tree products
on one sufficiently large common cover before any
carrier projection. Retain lifted data and winding
joins, then deck-average the classical coefficient
families and descend their sums. The norms have the
same positive majorants, and uniqueness identifies
the descended sum with the finite periodic selected
equation. This is descent of classical operators and
critical germs, not an equality of quantum partition
functions on different covers.

This result supplies the weighted all-degree selected
germ and its physical WCA restriction. It constructs
neither a trial for arbitrary proper coarse data nor
independent background cube discs. It proves no
weighted nonflat formation/tower/deletion statement,
and it does not identify a separately prescribed
off-current core convention merely from agreement
of its quadratic precision.

% END LAYERED EXPANSION background/weighted_connected_germ.tex



\section{Effective interactions and large fields}\label{sec:interactions}
We now organize one integration step and its iteration. The local estimates separate the relevant terms from the remainder, while the history sum controls regions in which the small-field expansion is unavailable. The source-dependent construction must preserve both normalization and the stated analytic domains. The estimates below specify where each of these requirements enters the induction.

% BEGIN MODULE chapters/rg_audit_chapter.tex
\subsection{Separated correlations from local source bounds}
\label{ch:conditional-source-bridge}
\subsubsection{From local source estimates to separated correlation limits}
\label{sec:conditional-source-bridge}

Local analytic energy estimates and a bound on the residual finite-volume integral imply bounds on separated correlations. This section proves that implication. Its hypotheses remain to be established for four-dimensional gauge theory: exact transport of a measure alone gives neither a partition-function bound nor decay of the dependence on older scales.

\paragraph{Finite-dimensional starting point}

Let $\mathcal P$ be a finite set of plaquettes, let $G=\operatorname{SU}(2)$, and assign one variable $U_b\in G$ to each positively oriented bond. Opposite orientations carry $U_{-b}=U_b^{-1}$. The product of normalized Haar measures is denoted by $dU$. Write

\begin{gather*}
\FitMath{a_p(U)=1-\tfrac12\operatorname{Re}\operatorname{tr}U_p,} \\[3pt]
\FitMath{A(U)=\sum_{p\in\mathcal P}a_p(U).}
\end{gather*}

Since $0\le a_p\le2$, the probability measure

\begin{gather*}
\FitMath{d\mu_\beta(U)=Z_\beta^{-1}e^{-\beta A(U)}dU,} \\[3pt]
\FitMath{Z_\beta=\int e^{-\beta A(U)}dU}
\end{gather*}

is well defined for every real $\beta$. Given bounded real functions $f_i$ on plaquette barycentres, set $X_i=\sum_pf_i(x_p)a_p(U)$ and

\begin{gather*}
\FitMath{\Gamma(z)=\int e^{\sum_i z_iX_i}d\mu_\beta.}
\end{gather*}

The bound $|X_i|\le2\sum_p|f_i(x_p)|$ implies uniform domination on every compact subset of $\mathbb C^n$. Differentiation under the integral therefore proves that $\Gamma$ is entire. Since $\Gamma(0)=1$, it has a holomorphic logarithm in a neighbourhood of the origin, and the joint cumulant is defined by

\begin{gather*}
\FitMath{\langle X_1;\ldots;X_n\rangle
 =\partial_{z_1}\cdots\partial_{z_n}\log\Gamma(0).}
\end{gather*}

The neighbourhood obtained this way can depend on the number of plaquettes.

If $T$ is a fixed linear operation on finite measures which preserves their total mass, then $T(e^{\sum z_iX_i}\mu_\beta)$ has total mass $\Gamma(z)$. For a finite sequence of such operations the same conclusion follows by induction. If $T$ is specified only on gauge-invariant densities, all insertions in this assertion must be gauge invariant. Neither linearity nor mass preservation supplies a uniform analytic bound on the resulting density.

\paragraph{A zero-free disc from a normalized bound}

\begin{lemma}\label{lem:normalized-zero-free-disc}
Let $c_i>0$ and put $\|z\|_c=\sum_i c_i|z_i|$. Suppose $h$ is holomorphic on $\|z\|_c<1$, $h(0)=1$, and $|h|\le Q$, where $Q\ge1$. Then

\begin{gather*}
\FitMath{|h(z)-1|\le(Q+1)\|z\|_c.}
\end{gather*}

In particular $h$ is nonzero on $\|z\|_c<(Q+1)^{-1}$. On the smaller ball $\|z\|_c<[2(Q+1)]^{-1}$ its principal logarithm satisfies $|\operatorname{Log}h|\le\log2$.
\end{lemma}

\begin{proof}
For $z\ne0$ write $z=tv$, where $t=\|z\|_c$ and $\|v\|_c=1$. The function $(h(\zeta v)-1)/(Q+1)$ maps the unit disc into its closed unit disc and vanishes at zero. The Schwarz lemma gives the first estimate. Strictness on the stated open balls follows by using slightly smaller compact discs and taking limits. For $|u|\le1/2$, the logarithmic series gives

\begin{gather*}
\FitMath{|\operatorname{Log}(1+u)|\le\sum_{k\ge1}\frac{|u|^k}{k}
 =-\log(1-|u|)\le\log2.}
\end{gather*}

The same series defines a single holomorphic branch on that ball.
\end{proof}

\begin{corollary}\label{cor:normalized-cumulant-bound}
Suppose $\Gamma=e^{\mathcal E}h$ near zero, where $\mathcal E(0)=0$ and the lemma applies to $h$. Then

\begin{gather*}
\FitMath{\left|\partial_{z_1}\cdots\partial_{z_n}
       (\log\Gamma-\mathcal E)(0)\right|
 \le\log2\,[2(Q+1)n]^n\prod_i c_i.}\qedhere
\end{gather*}


\end{corollary}

\begin{proof}
Use Cauchy's integral formula on a polydisc with radii strictly smaller than $[2(Q+1)nc_i]^{-1}$. The sum of the weighted radii is less than $[2(Q+1)]^{-1}$. Each variable is differentiated once, so no factorial occurs in this application of the formula. Let the radii increase to their stated values.
\end{proof}

The dependence on $Q$ is material. For example, $\cosh(Vz)$ is the moment-generating function of a random variable taking the values $V$ and $-V$ with equal probabilities. It is at most $e^V$ on the unit disc, but has a zero at $i\pi/(2V)$. Thus an exponential-in-volume estimate alone cannot give a volume-independent zero-free disc.

\paragraph{A precise local-energy hypothesis}

Fix $L>1$, a finite physical volume $V$, and a scale index $s\ge0$, with physical scale $a_s=L^{-s}$. A local domain $D$ has an integer size $m(D)\ge0$ and a dependence region $D^+$. Assume there are constants $C_0,C_1,\kappa_{\mathrm{ent}}$ such that

\begin{gather*}
\FitMath{\operatorname{diam}D^+\le C_0a_s(m(D)+1),}
\end{gather*}

and, for every fixed compact test support $B$, the number of domains of size $m$ whose dependence region meets $B$ is at most

\begin{gather*}
\FitMath{C_1V a_s^{-4}e^{\kappa_{\mathrm{ent}}m}.}
\end{gather*}

The factor $V$ is a finite-volume allowance; a more local counting estimate can replace it when available.

For cutoff $K$ let $e_{K,s,D}(w)$, $0\le s<K$, be holomorphic on the source ball

\begin{gather*}
\FitMath{\sum_i|z_i|\|f_i\|_*/r<1,} \\[3pt]
\FitMath{w=\sum_i z_if_i,}
\end{gather*}

where $r>0$ is independent of $K$. Suppose

\begin{gather*}
\FitMath{e_{K,s,D}(0)=0,} \\[3pt]
\FitMath{|e_{K,s,D}(w)|\le A e^{-\kappa m(D)},} \\[3pt]
\FitMath{\kappa>\kappa_{\mathrm{ent}}.}
\end{gather*}

Finally suppose $e_{K,s,D}$ depends on $w$ only through its values in $D^+$. These statements are the local-energy hypothesis. They include analyticity, an entropy gap, and geometric control of the dependence region.

\begin{proposition}\label{prop:separated-local-energy-bound}
Let $n\ge2$ and suppose the supports of $f_1,\ldots,f_n$ have pairwise distance at least $d>0$. Under the local-energy hypothesis,

\begin{gather*}
\FitMath{\left|\partial_{z_1}\cdots\partial_{z_n}
 \sum_{s<K}\sum_D e_{K,s,D}(w)\bigg|_{z=0}\right|
 \le CAV\max\{1,d^{-4}\}
       (n/r)^n\prod_i\|f_i\|_*.}
\end{gather*}

The constant is independent of $K$; it can depend on $L,C_0,C_1$ and the entropy gap.
\end{proposition}

\begin{proof}
If $D^+$ misses one support, the energy is independent of the corresponding variable, and the mixed derivative vanishes. Every contributing domain therefore meets both the first and second supports. Consequently

\begin{gather*}
\FitMath{m(D)\ge d/(C_0a_s)-1.}
\end{gather*}

Cauchy's formula with radii approaching $r/(n\|f_i\|_*)$ bounds each derivative by

\begin{gather*}
\FitMath{A(n/r)^n\prod_i\|f_i\|_*\,e^{-\kappa m(D)}.}
\end{gather*}

A zero test function makes the assertion immediate, so these radii may be assumed finite. The entropy estimate and a geometric series in $m$ bound the sum at scale $s$ by a constant times

\begin{gather*}
\FitMath{AVa_s^{-4}\exp(-c d/a_s),}
\end{gather*}

with a larger harmless constant when $d/a_s$ is bounded. It remains to estimate $\sum_{s\ge0}L^{4s}e^{-cdL^s}$.

For $0<d\le1$, let $s_*$ be the largest integer with $dL^{s_*}\le1$. The sum over $s\le s_*$ is at most $CL^{4s_*}\le Cd^{-4}$. Writing $s=s_*+j$ in the remaining sum bounds it by

\begin{gather*}
\FitMath{Cd^{-4}\sum_{j\ge1}L^{4j}e^{-c'L^j},}
\end{gather*}

which converges by the ratio test. For $d>1$, replace $d$ by $1$ in the exponential. This proves the stated uniform bound. At each fixed cutoff the domain sums converge normally on smaller source balls, permitting termwise differentiation there. Only the separated mixed-derivative estimate is asserted uniformly as the number of scales grows; the undifferentiated energy sums need not converge in that limit.
\end{proof}

If in addition $\Gamma_K=e^{\mathcal E_K}h_K$, where $\mathcal E_K=\sum_{s,D}e_{K,s,D}$ and $|h_K|\le Q(V)$ with $h_K(0)=1$, the corollary above gives a cutoff-independent bound on every separated cumulant. This conclusion remains a finite-volume conclusion when $Q$ depends on $V$.

\paragraph{What a comparison of cutoffs must prove}

Uniform boundedness gives compactness only after equicontinuity and a suitable countable test class have been supplied. Convergence along a prescribed whole sequence requires more. One sufficient condition is a summable comparison: for every fixed separated tuple,

\begin{gather*}
\FitMath{|S_{K+1}(f_1,\ldots,f_n)-S_K(f_1,\ldots,f_n)|
 \le C(f_1,\ldots,f_n)\epsilon_K,} \\[3pt]
\FitMath{\sum_K\epsilon_K<\infty.}
\end{gather*}

The triangle inequality then makes the sequence Cauchy. Its tail is bounded by $C\sum_{k\ge K}\epsilon_k$.

A local proof of this comparison can split the scales at $K/2$. At $s\le K/2$, it is enough to compare paired local energies with a bound

\begin{gather*}
\FitMath{|e_{K+1,s,D}-e_{K,s,D}|
 \le C(1+s)\epsilon_K e^{-\kappa m(D)}.}
\end{gather*}

At $s>K/2$, the individual energy bounds already give a tail controlled by $\sum_{s>K/2}L^{4s}e^{-cdL^s}$. This tail is summable in $K$. The factor $1+s$ in the paired bound is harmless because $\sum_s(1+s)L^{4s}e^{-cdL^s}<\infty$ for each $d>0$. The residual integrals must also be compared on a common zero-free source ball. If both residuals are within $1/2$ of $1$, then the logarithmic difference is bounded by twice the difference of the residuals, using the line segment between their values and the derivative $1/z$ of the logarithm. Cauchy's formula converts this bound into a cumulant comparison.

These are consequences of an actual paired estimate. A bound for two covariant constant-source lists does not automatically prove the estimate for position-dependent arrays. Similarly, ordinary holomorphy in an unweighted norm does not create geometric factors for the age of an inherited activity. Those factors must be present in a separately established weighted estimate or in the domain and bound of a weighted holomorphic extension.

\paragraph{A conditional tuning lemma}

Let $F_K$ be a real differentiable renormalized-coupling coordinate on an interval $I_K$, and suppose

\begin{gather*}
\FitMath{2/3\le F_K'(\beta)\le4/3} \\[3pt]
\FitMath{(\beta\in I_K).}
\end{gather*}

If $F_K$ takes values on both sides of $x_*$, the intermediate value theorem gives a unique $\beta_K\in I_K$ with $F_K(\beta_K)=x_*$. Moreover,

\begin{gather*}
\FitMath{\tfrac23|\beta-\beta'|
 \le|F_K(\beta)-F_K(\beta')|
 \le\tfrac43|\beta-\beta'|.}
\end{gather*}

Both inequalities follow from the mean value theorem on the interval between the two points. The inverse coordinate is therefore Lipschitz with constant $3/2$ on its image. If $F_K(\beta)=\beta-\sum_{j\le K}b_j(\beta)$, the derivative hypothesis follows from $\sum_j\sup_{I_K}|b_j'|\le1/3$. Existence of the interval, the crossing of the target, and the derivative estimate remain hypotheses to be established by the construction.

The conclusions of this section concern separated correlation functionals. They do not, without additional estimates, provide extensions across coincident points, reflection-positive distributions on the full test-function algebra, rotation covariance, or a characterization of the resulting scalar observable sector as a Yang--Mills theory. Those are additional mathematical steps with their own hypotheses.


% END MODULE chapters/rg_audit_chapter.tex

% BEGIN MODULE chapters/pointed_hessian_cancellation_v6.tex
\subsection{A pointed Hessian and its action counterterm}
\label{ch:pointed-hessian-cancellation}

For one local interaction, we identify the marginal quadratic term of
its pointed total composed with the free response defined below. This
term determines the correction to the Wilson coupling; the
renormalisation step must subtract it with that interaction's own
coefficient and action normalization before estimating the nonlinear
remainder. The calculation treats each base point separately. The bulk
sum is a convolution, although an individual pointed total need not be.
The nonlinear window remainder and its field-domain estimates are not
consequences of this quadratic calculation.

\subsubsection*{The pointed object and the analytic hypotheses}

Work in coarse lattice units on $\mathbb Z^4$. An independently
oriented bond $(x,\mu)$ starts at $x$ and ends at $x+e_\mu$.
A real field $B_\mu(x)$ is a traceless Hermitian matrix. Choose a
color basis $t_a$, $1\le a\le3$, orthonormal for
$\langle A,C\rangle=\frac12\operatorname{Tr}(AC)$. Thus
$B=\sum_a B^a t_a$. Complexification is bilinear, not sesquilinear.
The normalization gives the quadratic Wilson action
\begin{equation}
 A(e^{itB})=\frac{t^2}{2}\sum_{p}|\mathrm{d}B(p)|^2+O(t^3)
 \label{ph:action-normalization}
\end{equation}
on every finite set of plaquettes. Here $d$ is the oriented lattice
curl and each coordinate-plane plaquette is counted once.

An individual member $G$ has the pointed pair total, in its
exponential chart,
\begin{equation}
 \mathcal G_z(e^{i\eta A},J)=\sum_{X\ni z}
                  G(X;(e^{i\eta A},J),z),\qquad z\in\mathbb Z^4.
 \label{ph:pointed-total}
\end{equation}
The bare spacing $\eta$ can differ from the coarse unit spacing.
There are no stored spectators in a regular member. Assume the
following analytic and geometric properties of this actual family.

The sum in \eqref{ph:pointed-total} and its first two derivatives
converge normally in the neighborhoods where they are used. Every
summand is local in its pair variables and invariant under all local
gauge transformations for which the arguments remain in its declared
domain. In particular it is invariant under constant adjoint rotations.
Let $\mathcal P(B)$ be the actual background pair
formed from a sufficiently small finitely supported coarse field,
with $\mathcal P(0)=(1,0)$. It is twice differentiable, is gauge
covariant, and has first response, expressed in this potential chart,
\begin{equation}
 hB=D\mathcal P(0)B=(HB,D_1^*D_1HB).
 \label{ph:response}
\end{equation}
The precise linear map $H$ is part of the background construction;
it is not replaced by a convolution chosen for this proof. Put
$F_z(B)=\mathcal G_z(\mathcal P(B))$. Assume the pointed totals
are covariant under unit translations and the full hypercubic group,
including reflections. These are properties of one individual member
at all its base points. They do not require two different members
chosen at different sites to be related.

Finally assume an absolute exponential Hessian moment bound. With
$\delta_{x,\mu}^a$ the field of color $a$ at one coarse bond, let
\begin{equation}
 \mathcal K_z^{\mu\nu;ab}(x,y)
  =\sum_{X\ni z}D^2G(X,(1,0),z)
                      [h\delta_{x,\mu}^a,h\delta_{y,\nu}^b].
 \label{ph:actual-kernel}
\end{equation}
All derivatives of a pair function are taken in the $(A,J)$ chart
of \eqref{ph:pointed-total}, including its factor $\eta$.
For some $m>0$ assume
\begin{equation}
 \sum_{\mu,\nu,a,b}\sum_{x,y\in\mathbb Z^4}
 e^{m(|x-z|+|y-z|)}|\mathcal K_z^{\mu\nu;ab}(x,y)|\le C_m.
 \label{ph:exponential-moment}
\end{equation}
The corresponding derivative representation on finitely supported
fields is assumed, as is usual when the normally convergent
localized Hessians are composed with the free response. This bound
and the actual background/core hypotheses are model inputs here.
In particular, this section does not infer a uniformly bounded
coefficient extractor from a potential norm on a shrinking field ball.

We first justify that \eqref{ph:actual-kernel} really is the Hessian
of $F_z$. Constant adjoint invariance makes the first derivative of
each summand of \eqref{ph:pointed-total} zero at $(1,0)$: each
coefficient of a potential or current coordinate would be an
invariant linear functional on the adjoint color space, which has
no fixed vector. This applies to its finitely many local coordinates
and hence to every direction supplied by $h$. The normal convergence
permits summation. In the second-order chain rule the term containing
$D\mathcal G_z(1,0)D^2\mathcal P(0)$ is therefore zero.
Consequently \eqref{ph:actual-kernel}, with the displayed actual
$h$, is exactly $D^2F_z(0)$.

The exponential moment hypothesis has a concrete sufficient
localized form. Suppose the analytic core contains a usable pair
norm ball of radius $a_{\rm core}$, and on it
$|G(X,\cdot,z)|\le g e^{-\kappa d(X)}$. Suppose also
\[
 \|h\delta_{x,\mu}^a\|_{1;X}
       \le C e^{-\delta\operatorname{dist}(x,X)},\qquad
 \#\{X\ni z:q\le d(X)<q+1\}\le C e^{\kappa_D q},
 \qquad \kappa>\kappa_D,
\]
and $\operatorname{diam}(X)\le C_M(d(X)+1)$, with fixed
localization cube size. These are the core, response and entropy
conditions used in quadratic extraction. Two-variable Cauchy gives
\[
 |D^2G(X,0,z)[h\delta_x,h\delta_y]|
 \le Cg a_{\rm core}^{-2} e^{-\kappa d(X)}
 e^{-\delta(\operatorname{dist}(x,X)+
                 \operatorname{dist}(y,X))}.
\]
Choose $m<\delta/2$ so small that
$2mC_M<(\kappa-\kappa_D)/2$. Since $z\in X$, inserting the
weight in \eqref{ph:exponential-moment} costs at most
$e^{2mC_M(d(X)+1)}$ after retaining the two response distances.
The remaining sums over $x$ and $y$ cost at most
$C(d(X)+1)^8$: enclose $X$ in a box of diameter
$C_M(d(X)+1)$ and sum the exponentially decreasing distances
outside it. The rooted domain count now gives
\begin{equation}
 C_m\le Cg a_{\rm core}^{-2}
 \sum_{q\ge0}(q+1)^8
 e^{-(\kappa-\kappa_D-2mC_M)q}
 \le C' g a_{\rm core}^{-2}.
 \label{ph:moment-from-core}
\end{equation}
Thus \eqref{ph:exponential-moment} need not be introduced as a
new physical estimate when these localized hypotheses have already
been supplied. The factor $a_{\rm core}^{-2}$ is retained. The
coordinate normalization and near-zero core inclusions needed to
use this Cauchy estimate are treated in
\cref{ch:extraction-radius-normalization}.

\subsubsection*{Color, symmetry, and the local gauge Ward identity}

Constant adjoint invariance implies
\begin{equation}
 \mathcal K_z^{\mu\nu;ab}(x,y)=\delta_{ab}K_z^{\mu\nu}(x,y),
 \qquad
 K_z^{\mu\nu}(x,y)=\frac13\sum_{a=1}^3
                         \mathcal K_z^{\mu\nu;aa}(x,y).
 \label{ph:color}
\end{equation}
For a direct proof, fix the spatial indices and view the color
coefficients as a $3\times3$ complex matrix. Adjoint rotations
through $\pi$ about the color axes force its off-diagonal entries
to zero, and rotations interchanging the axes make its diagonal
entries equal. The adjoint action of $\operatorname{SU}(2)$ contains
these rotations. This also proves the assertion for complex
members. Equality of mixed derivatives gives
\begin{equation}
 K_z^{\mu\nu}(x,y)=K_z^{\nu\mu}(y,x).
 \label{ph:Hessian-symmetry}
\end{equation}
For a real member all these coefficients are real. Thus its bulk
kernel is self-adjoint in addition to being bilinearly symmetric.
No positivity of a Hessian is assumed.

For a finitely supported Lie-algebra-valued site function $\chi$,
write $(\mathrm{d}\chi)_\mu(x)=\chi(x+e_\mu)-\chi(x)$. Gauge covariance
of $\mathcal P$ and gauge invariance of the member give
$F_z(B^g)=F_z(B)$, where $B^g$ is the logarithmic coarse bond
field after a small gauge transformation. Choose the gauge
convention $U(x,\mu)^g=g(x)^{-1}U(x,\mu)g(x+e_\mu)$, so
the derivative at the identity of $g=e^{it\chi}$ is $\mathrm{d}\chi$.
Differentiate this identity once in $t$ and once in a finitely
supported field direction $C$. The mixed derivative of the gauge
coordinate itself is multiplied by $DF_z(0)=0$. Therefore
\begin{equation}
 D^2F_z(0)[\mathrm{d}\chi,C]=0.
 \label{ph:compact-Ward}
\end{equation}
The other sign convention for gauge transformations gives the same
zero identity. The argument initially uses real directions and
extends complex bilinearly. It uses neither a noncompact gauge
transformation nor a polynomial field in a nonlinear field domain.

\subsubsection*{Polynomial extensions justified by absolute moments}

For fields of polynomial growth define the extended Hessian by
\begin{equation}
 \mathcal H_z(B,C)=
 \sum_{\mu,\nu,a}\sum_{x,y}
          K_z^{\mu\nu}(x,y)B_\mu^a(x)C_\nu^a(y).
 \label{ph:polynomial-Hessian}
\end{equation}
Every such sum converges absolutely by
\eqref{ph:exponential-moment}. Compact cutoffs converge to it by
dominated convergence, independently of the cutoffs. This defines
an extension of the quadratic jet, not a value of the nonlinear
member at an unbounded polynomial field.

The Ward identity extends to polynomial gauge parameters and
polynomial second arguments. To prove it, let $\chi$ be a
polynomial of degree $q$ on sites and let
$\chi_R(x)=\chi(x)\vartheta(x/R)$ with a fixed smooth compact
cutoff equal to one near zero. Its forward differences converge
pointwise to $\mathrm{d}\chi$, and satisfy a bound
\[
 |\mathrm{d}\chi_R(x)|\le C_\chi(1+|x|)^{\max(q-1,0)}
                                      \quad(R\ge1).
\]
Indeed away from the cutoff annulus the usual polynomial difference
has this bound; on the annulus the extra difference of the cutoff
costs $R^{-1}$ and $|x|$ is comparable to $R$. Apply
\eqref{ph:compact-Ward} with a single-bond second argument and
pass to the limit using \eqref{ph:exponential-moment}. Summing
against any polynomial second argument is then legitimate by
absolute convergence. The radius of the nonlinear gauge chart
may depend on $R$: only its derivative at zero was used for each
fixed $R$.

In particular the linear polynomial $\chi(x)=x_\mu t_a$ gives
the separate row and column zeroth moments
\begin{equation}
 \sum_x K_z^{\mu\nu}(x,y)=0,\qquad
 \sum_y K_z^{\mu\nu}(x,y)=0.
 \label{ph:row-column-zero}
\end{equation}
The second identity follows also from
\eqref{ph:Hessian-symmetry}. Both hold for each fixed remaining
edge and pair of directions, not only after a full trace.

For the moment at $z=0$ set, initially in lower bond-start
coordinates,
\begin{equation}
 T_{\mu\nu,\kappa\lambda}
       =\sum_{x,y}K_0^{\mu\nu}(x,y)x_\kappa y_\lambda.
 \label{ph:T}
\end{equation}
For a symmetric matrix $S$, the quadratic gauge polynomial
$\chi(x)=\frac12\sum_{\mu,\kappa}S_{\mu\kappa}
x_\mu x_\kappa t_a$ has forward difference
\[
 (\mathrm{d}\chi)_\mu(x)=
       \sum_\kappa S_{\mu\kappa}x_\kappa+\tfrac12S_{\mu\mu}.
\]
Use the polynomial Ward identity with second argument
$C_\nu(y)=y_\lambda t_a$ at one direction $\nu$. The constant
term vanishes by \eqref{ph:row-column-zero}. Since $S$ is
arbitrary symmetric, one obtains
\begin{equation}
 T_{\mu\nu,\kappa\lambda}
       =-T_{\kappa\nu,\mu\lambda},\qquad
 T_{\mu\nu,\kappa\lambda}
       =-T_{\mu\lambda,\kappa\nu}.
 \label{ph:T-Ward}
\end{equation}
The second identity follows by interchanging the two Hessian
arguments and repeating the first. These are exactly the two
antisymmetries which remove symmetric gradients.

\subsubsection*{The bulk sum and its second moment}

Unit translation covariance of this individual member gives
$K_z^{\mu\nu}(x,y)=K_0^{\mu\nu}(x-z,y-z)$. The bulk quadratic
kernel is
\begin{equation}
 \Pi^{\mu\nu}(x,y)=\sum_zK_z^{\mu\nu}(x,y)
                         =\pi_{\mu\nu}(y-x),\qquad
 \pi_{\mu\nu}(u)=\sum_xK_0^{\mu\nu}(x,x+u).
 \label{ph:bulk}
\end{equation}
All sums converge absolutely. In fact, for $0<m'\le m$,
\begin{equation}
 \sum_u e^{m'|u|}|\pi_{\mu\nu}(u)|
 \le\sum_{x,y}e^{m'(|x|+|y|)}|K_0^{\mu\nu}(x,y)|<\infty.
 \label{ph:bulk-moment-bound}
\end{equation}
For fixed $x,y$, translation of $z$ in
\eqref{ph:bulk} is a sub-sum of the same absolute double sum.
This proves both the bulk formula and every rearrangement below.
Only this bulk object is a convolution. The original $K_0(x,y)$
continues to have two independent spatial arguments.

The zeroth and first bulk moments vanish by
\eqref{ph:row-column-zero}. Its second moment is
\begin{equation}
 \begin{split}
 P_{\mu\nu,\kappa\lambda}
  &:=\sum_u\pi_{\mu\nu}(u)u_\kappa u_\lambda\\
  &=\sum_{x,y}K_0^{\mu\nu}(x,y)
                    (y_\kappa-x_\kappa)(y_\lambda-x_\lambda)\\
  &=-T_{\mu\nu,\kappa\lambda}
                         -T_{\mu\nu,\lambda\kappa}.
 \end{split}
 \label{ph:P-versus-T}
\end{equation}
The pure $y_\kappa y_\lambda$ term is zero after summing $x$
by its row identity, and the pure $x_\kappa x_\lambda$ term is
zero after summing $y$. This is the reason for both minus signs;
no pointwise convolution identity for $K_0$ has been used.

The Fourier series
$\widehat\pi_{\mu\nu}(k)=\sum_u\pi_{\mu\nu}(u)e^{ik\cdot u}$
is holomorphic for sufficiently small $|\operatorname{Im}k|$
by \eqref{ph:bulk-moment-bound}. In particular
$\partial_\kappa\partial_\lambda\widehat\pi_{\mu\nu}(0)
=-P_{\mu\nu,\kappa\lambda}$. This optional Fourier description
concerns the bulk sum alone.

\subsubsection*{Bond midpoints and the hypercubic tensor}

Under reflections, lower bond starts undergo a direction-dependent
integer displacement. To apply tensor symmetry without dropping
that displacement, move first to bond-midpoint coordinates
$\widetilde x=x+e_\mu/2$ and
$\widetilde y=y+e_\nu/2$. The value of $T$ does not change:
expanding $\widetilde x_\kappa\widetilde y_\lambda$ adds only
terms of the forms $\sum_{x,y}K_0y_\lambda$,
$\sum_{x,y}K_0x_\kappa$ and $\sum_{x,y}K_0$, all zero by
\eqref{ph:row-column-zero}. The displacement $u$ changes by the
constant $(e_\nu-e_\mu)/2$. Because the bulk zeroth and first
moments are zero, its second moment $P$ is also unchanged.
Thus the lower-start extractor and the midpoint extractor agree
exactly.

The midpoint bond sets are
$\mathbb Z^4+e_\mu/2$. A signed permutation $w$ with
$we_\mu=s_\mu e_{\sigma(\mu)}$ maps this set to the midpoint
set of direction $\sigma(\mu)$; a negative orientation contributes
the sign $s_\mu$ to a one-form. Pointed covariance at the fixed
base point zero and the ordinary transformation of the midpoint
coordinates therefore give
\begin{equation}
 T_{\sigma(\mu)\sigma(\nu),\sigma(\kappa)\sigma(\lambda)}
       =s_\mu s_\nu s_\kappa s_\lambda
                                  T_{\mu\nu,\kappa\lambda}.
 \label{ph:W4}
\end{equation}
There is no affine correction left in this formula; the preceding
moment calculation proved that it vanished.

We now classify this tensor directly. By
\eqref{ph:T-Ward} it is zero if $\mu=\kappa$ or
$\nu=\lambda$. Reflecting one coordinate in
\eqref{ph:W4} shows that each coordinate label must occur an
even number of times among the four indices of a nonzero entry.
The two antisymmetric pairs must consequently be the same
unordered pair. The only possibilities are
$T_{\mu\mu,\kappa\kappa}$ and
$T_{\mu\kappa,\kappa\mu}$ for $\mu\ne\kappa$.
Permutations make all the first values equal to a number $b$;
\eqref{ph:T-Ward} makes the second values $-b$. Hence
\begin{equation}
 T_{\mu\nu,\kappa\lambda}
      =b\bigl(\delta_{\mu\nu}\delta_{\kappa\lambda}
                     -\delta_{\mu\lambda}\delta_{\kappa\nu}\bigr).
 \label{ph:Maxwell-T}
\end{equation}
Full reflections are used here: four distinct indices, including a
possible alternating tensor, are excluded by a single reflection.
This is not merely an argument using proper rotations.

Combining \eqref{ph:P-versus-T} with
\eqref{ph:Maxwell-T} gives
\begin{equation}
 P_{\mu\nu,\kappa\lambda}
    =b\bigl(\delta_{\mu\lambda}\delta_{\kappa\nu}
            +\delta_{\mu\kappa}\delta_{\lambda\nu}
            -2\delta_{\mu\nu}\delta_{\kappa\lambda}\bigr),
 \qquad b=P_{\mu\nu,\mu\nu}\quad(\mu\ne\nu).
 \label{ph:extractor}
\end{equation}
Thus the extractor is explicitly
$b[G]=\sum_u\pi_{\mu\nu}(u)u_\mu u_\nu$ for any distinct
$\mu,\nu$. It has no extra color factor, because
\eqref{ph:color} used the average of the three diagonal color
Hessians. The quadratic term of the bulk symbol is
$b(|k|^2\delta_{\mu\nu}-k_\mu k_\nu)$; higher lattice
derivative terms have not been set to zero.

For fixed actual $h$, the extractor is complex linear in the member
and ignores additive constants. It is continuous in the exponential
Hessian moment norm: the absolute second moment is bounded by that
norm times $\sup_{r\ge0}r^2e^{-mr}$, with a harmless choice of a
smaller exponential weight. Normal convergence in that norm also
permits termwise extraction and preserves holomorphy in a member
parameter. These are Hessian-norm statements, and make no assertion
that the Hessian norm is uniformly bounded by a potential norm on
coupling-dependent shrinking cores.

\subsubsection*{The affine quadratic jet and the tent action}

Let $B_\mu(x)=a_\mu+\sum_\kappa C_{\mu\kappa}(x-z)_\kappa$,
using midpoint coordinates or lower starts. The difference is a
constant one-form, which the Ward identity annihilates. \eqref{ph:row-column-zero} removes every occurrence of $a$, and
\eqref{ph:Maxwell-T} proves
\begin{equation}
 \frac12\mathcal H_z(B,B)
       =\frac b2\sum_{\mu<\nu}
                                  |C_{\nu\mu}-C_{\mu\nu}|^2.
 \label{ph:affine-jet}
\end{equation}
For complex fields the squared norm means the bilinear color
pairing; for real fields it is the ordinary squared norm. The
identity follows by expanding
$\sum_{\mu,\kappa}|C_{\mu\kappa}|^2
-\sum_{\mu,\nu}\langle C_{\mu\nu},C_{\nu\mu}\rangle$
into unordered pairs. In particular the axial affine choice
$\ell_\mu(x)=\sum_{\kappa<\mu}(x-z)_\kappa F_{\kappa\mu}$
has curl $F$ and quadratic jet $\frac b2\sum_{\mu<\nu}|F_{\mu\nu}|^2$.

To identify this with the member's own action counterterm, one more
normalization of the actual free response is needed. Assume that on
an affine coarse field $\ell$ of constant curl it reproduces the
same affine bare potential modulo a lattice gradient, and its
current component is zero. This is an affine-response statement
about the actual averaging/background problem, not a consequence
of the tensor classification. For this normalization argument
assume also that, at each bare bond in a bounded coarse region,
the matrix elements of $H$ in the coarse input bond decay
exponentially in its distance from that region. Their constants
must be uniform in its finitely many bare bonds at fixed $\eta$.
This is a precise sufficient free-response input. It makes
$H(\ell\vartheta(\cdot/R))$ converge at every such bond to an
absolutely convergent polynomially weighted matrix sum, by
dominated convergence. Thus its affine reproduction statement is
an equality of well-defined extended first responses, with no
nonlinear affine field inserted outside its analytic domain.

This input also gives the moment bound needed to apply the tensor
theorem to the pointed action itself. Its Hessian is a finite sum
over the bare plaquettes in a unit tent of products of first
differences of $H$. The exponential bounds at their finitely many
bonds give a product bound
$C_\eta e^{-m_H(|x-z|+|y-z|)}$. Summing in the two input bonds
with any smaller exponential weight proves
\eqref{ph:exponential-moment} for this action. This absolute
convergence statement does not assert that $C_\eta$ has a uniform
cutoff bound; any such bound is a further quantitative input.

There is a useful stronger scope. Write
$\mathrm{d}_\eta H(p)=\eta^{-1}\sum_{b\in\partial p}H(b)$ with oriented
signs. If the actual response satisfies the physical first-difference
bound
\[
 |\mathrm{d}_\eta H\delta_{x,\mu}^a(p)|
      \le C e^{-m_H|x-z|}\qquad (h_z(p)>0),
\]
uniformly in $\eta$, then the moment constant for the pointed
action is uniform in $\eta$ as well. Indeed, its kernel contains
\[
 \sum_p \eta^4 h_z(p)
 (\mathrm{d}_\eta H\delta_x)(p)\mathbin{\cdot}
 (\mathrm{d}_\eta H\delta_y)(p).
\]
The exact quadrature proved below makes the weighted plaquette
sum one in each orientation. Consequently this kernel is bounded
by $C e^{-m_H(|x-z|+|y-z|)}$, and any smaller exponential moment
is summable uniformly. This explains precisely why physical first
differences, rather than identity-only bounds on bare values, remove
the possible $C_\eta$ loss.

Let $\eta=1/N_0$ for an integer $N_0\ge1$, and use the midpoint
tents
\begin{equation}
 h_z(p)=\prod_{r=1}^4(1-|p_r-z_r|)_+,
 \qquad A_z(U)=\sum_{\mu<\nu}\sum_{p_{\mu\nu}}h_z(p)A_p(U).
 \label{ph:tent-action}
\end{equation}
For each fixed plaquette orientation,
\begin{equation}
                 \eta^4\sum_{p_{\mu\nu}}h_z(p)=1.
 \label{ph:quadrature}
\end{equation}
Here its barycenters have half-mesh shifts in directions $\mu,\nu$
and zero shifts in the other directions. The one-dimensional
unshifted identity is
$N_0^{-1}[1+2\sum_{r=1}^{N_0-1}(1-r/N_0)]=1$.
For the half-mesh shift it is
$2N_0^{-1}\sum_{r=0}^{N_0-1}(1-(r+1/2)/N_0)=1$.
Multiplication of these four identities proves
\eqref{ph:quadrature}, including the coarse case $N_0=1$.
Also $\sum_z h_z(p)=1$, since the two neighboring one-dimensional
tents sum to one in each coordinate.

For a bare affine potential of curl $F$, the first variation of a
plaquette holonomy in the chart $U_b=e^{it\eta A_b}$ is
$it\eta^2F_{\mu\nu}$. \eqref{ph:action-normalization}, gradient invariance and
\eqref{ph:quadrature} therefore give
\begin{equation}
       \frac12D^2(A_z\circ\mathcal P)(0)[\ell,\ell]
                      =\frac12\sum_{\mu<\nu}|F_{\mu\nu}|^2.
 \label{ph:action-b-one}
\end{equation}
The affine-response hypothesis was used exactly in this line.
The same pointed tensor theorem applies to this pointed action,
so \eqref{ph:affine-jet} identifies its extractor as $b[A_z]=1$.
Consequently
\begin{equation}
 D^2\bigl(F_z-b[G](A_z\circ\mathcal P)\bigr)(0)
                                    [\ell,\ell]=0
 \label{ph:own-counterterm}
\end{equation}
for every affine constant-curvature field. Polarization gives the
corresponding mixed affine identity. This proves the exact
extractor-to-counterterm match with the normalization
$A(e^{iB})=\frac12|\mathrm{d}B|^2+O(B^3)$.

\subsubsection*{Selecting different members at different sites}

The proof can be applied separately to each member $G^{[z]}$
chosen at a site. Within that member the family at all hypothetical
base points is translation covariant, so its own bulk extractor
$b[G^{[z]}]$ is defined by \eqref{ph:bulk} and
\eqref{ph:extractor}. Its pointed total at the actual site $z$
is canceled by that number times $A_z$, by
\eqref{ph:own-counterterm}. No translation covariance of the
array $z\mapsto G^{[z]}$ has been used.

This distinction is necessary for an inhomogeneous scalar source:
one may freeze it to a different constant-source member at each
site, but must use the coefficient of that same member at that
site. Equality of two unpointed totals alone does not identify
their separate site weights. Likewise the classification above
does not prove a nonlinear third-order or H\"older remainder.
That step still requires the actual localized response, its
uniform field-domain inclusion and its stated analytic derivative
bounds on the window.


% END MODULE chapters/pointed_hessian_cancellation_v6.tex

% BEGIN MODULE chapters/pointed_window_remainder_v6.tex
\subsection{The nonlinear remainder in a pointed window}
\label{ch:pointed-window-remainder}

The algebraic cancellation of a pointed Hessian concerns the second derivative
at the identity.  Its use at a nonlinear background requires a local Taylor
estimate.  We prove that estimate from the background hypotheses, retaining
the domain assertions that make the compositions meaningful.  In particular,
the datum is large compared with its value close to the base point.  A
Taylor bound using only its global supremum would lose the required power
of the scale ratio.

\subsubsection*{The background inputs and their coordinate radii}

Work first in units in which the coarse spacing of a single-scale window is
one.  The background package A1--A3 is used in the following explicit form.
On a box with collars there is a holomorphic potential/current representative
\(\Phi_\Lambda(B)\), with \(\Phi_\Lambda(0)=0\), for the actual constrained
critical orbit.  Its datum domain contains
\begin{equation}
 \mathcal B_{r_B}
   =\{B:\hbox{determining bonds}\longrightarrow\mathfrak{su}(2)^\mathbb C:
                        \sup_c|B(c)|<r_B\}.
 \label{eq:pointed-window-remainder-PW1}
\end{equation}
Here \(B(c)\) is a dimensionless logarithm of a group-valued coarse bond.
This is a product of balls at bonds.  It has no coupled derivative
restriction.  The orbit representative, its current, and the response
estimates hold throughout this complex domain.  Core and tower memberships
under restriction and substitution are also part of A1; they are not inferred
from the formal existence of a logarithm.  All differentiations and coordinate
sums below are first carried out on finite boxes.  Their uniform bounds permit
the stated coherent infinite-volume limits; no representation of a derivative
on an infinite \(\ell^\infty\) space by coordinate derivatives is assumed.

For clarity, the output norm over a top of local spacing \(\xi_y\) is
\begin{equation}
 Z_y=(\xi_yH,\xi_y^2\nabla H,\xi_y^3J),\qquad
 |Z_y|=\max\{\xi_y|H|,\ \xi_y^2|\nabla H|,\ \xi_y^3|J|\}.
 \label{eq:pointed-window-remainder-PW2}
\end{equation}
The first-response bound supplied by A3 is
\begin{equation}
 |\partial_{B(c)}Z_y(B)[e]|
       \le C e^{-m d(y,y_c)}|e|,
       \qquad B\in\mathcal B_{r_B},
 \label{eq:pointed-window-remainder-PW3}
\end{equation}
together with the stated spatial H\"older bound for \(\nabla H\).
The distance is the local-scale path distance. On the central
homogeneous part of the window, \eqref{eq:pointed-window-remainder-PW3}
uses level-\(j\) distances and the output weights are one.  Constants may depend on the
fixed window size, the decay rate and the H\"older exponent, but not on the
bare mesh, the torus or the number of preceding steps.

In the central part of the old box the top spacing is one in
the old units. The finer tops in the tower collar are handled
by the weighted row sum for the actual response. We compare
affine data on this same determining set, with its prescribed
seam rows.

Choose a fixed basis of the three-dimensional complex colour space.  If
\(\sup|B|\le r_B/2\), varying one basis coordinate by
\(|s|<c_{\rm col}r_B/2\) stays in \eqref{eq:pointed-window-remainder-PW1}, where \(c_{\rm col}>0\) depends
only on this choice of norm and basis.  Thus there is an admissible
coordinate Cauchy radius independent of the number of bonds.  The factors
\(\xi_y\) in \eqref{eq:pointed-window-remainder-PW2} are output weights and commute with this Cauchy estimate.
If instead a physical potential datum is written \(a(c)=B(c)/\xi_c\),
its coordinate radius is \(c_{\rm col}r_B/(2\xi_c)\), and
\(\partial_{a(c)}=\xi_c\partial_{B(c)}\).  Contracting a second derivative
against \(a(c)a(c')\) cancels both input factors \(\xi_c,\xi_{c'}\).
No inverse local spacing is lost by using dimensionless bond coordinates.

\begin{lemma}[Product decay for the second response]
On \(\mathcal B_{r_B/2}\), the first-response hypothesis implies
\begin{equation}
 |\partial_{B(c)}\partial_{B(c')}Z_y(B)[e,e']|
 \le \frac{C}{r_B}
       e^{-\frac m2(d(y,y_c)+d(y,y_{c'}))}|e|\,|e'|.
 \label{eq:pointed-window-remainder-PW4}
\end{equation}
The same conclusion holds for the local first-order pair norm on a domain
\(X\), with \(d(y,y_c)\) replaced by \(\operatorname{dist}(c,X)\).
\end{lemma}

\begin{proof}
Apply the coordinate Cauchy estimate to
\(\partial_{B(c)}Z_y(B)[e]\), using a slightly smaller fixed fraction of
the radius just specified.  This gives
\(Cr_B^{-1}e^{-m d(y,y_c)}|e|\,|e'|\).
The finite-dimensional holomorphic map has symmetric second derivatives,
so the same argument interchanging \(c,c'\) gives the corresponding bound
with \(c'\).  The smaller of two positive bounds is at most their geometric
mean, which proves \eqref{eq:pointed-window-remainder-PW4}.  Taking the supremum over the output bonds,
first differences and current components gives the stated local norm
version.  Colour indices cost only a fixed factor.
\end{proof}

This argument uses response decay at every complex datum in a ball.
Decay of the linear response at \(B=0\) alone is insufficient.

\subsubsection*{Weighted polynomial data and the local expansion}

Fix a base point \(z\), a scale ratio \(0<\sigma\le1\), and a window whose
distance from \(z\) to its boundary is comparable with \(\sigma^{-1}\).
Write \(a\) for a radius of a regular member, \(a_{\rm out}\) for the
output field radius and
\begin{equation}
 \rho=a_{\rm out}/a\ge1,\qquad
     a\le C_0r_B,\qquad a_{\rm out}\le c_0,\qquad
     0<\alpha\le1 .
 \label{eq:pointed-window-remainder-PW5}
\end{equation}
The constant \(c_0\) is chosen small enough for the matrix charts.
A fixed constant \(C_{\rm ax}\) in the axial datum estimate below is understood.
The local Taylor argument requires the full axial datum in
\(\mathcal B_{r_B/2}\). The central rows have bound
\(C_{\rm ax}a_{\rm out}\sigma\); the remote rows require their
own absolute bound in the same chart. The gain \(\sigma\) on
the central rows does not bound the remote ones. Once both
bounds are below \(r_B/2\), the entire segment from zero lies
in this polydisc.

For the flat-reference windows, this absolute domain condition
has a choice uniform in the depth. Use a fixed continuation
radius \(r_{\rm abs}>0\), smaller if necessary than the common
row and node-gauge radii, as \(r_B\). Fix the window geometry
first. The full induced datum of the upper background has
bound \(C_M A_0\), whereas its central rows have the additional
factor \(\sigma\), as in the proof of
\cref{htc:inner-substitution}. A fluctuation of amplitude
\(A_{\rm fl}\) changes a level-\(q\) raw row by at most
\(C L^{q-k}A_{\rm fl}\), with \(q\le j\le k\).
First omit the independent raw correction and call the physical
datum \(D\). If \(A_0,A_{\rm fl}\le a_{\rm abs}\), its full
norm is at most \(C_Ma_{\rm abs}\).
Construct the old axial change from this uncorrected datum.
Its comb in the settled layer has \(O(\sigma^{-1})\) links
of logarithmic size \(O(C_Ma_{\rm abs}\sigma)\), so the node
gauge has logarithm \(O(C_Ma_{\rm abs})\). Extend that
logarithm with a cutoff supported inside the settled layer.
On every row the transformed datum is the product of its
original value and its two endpoint gauge factors.

For an independent raw correction \(E\), hold this gauge
\(g\) fixed with respect to \(E\). It may still depend
holomorphically on the physical input. In additive datum
coordinates its action is the rowwise map
\[
 \mathcal G_g(D)(c)=\frac1i\log\!\left[
       g(c_-)e^{iD(c)}g(c_+)^{-1}\right].
\]
The endpoint factors and their inverses have bounded norms,
and these products lie in a fixed smaller logarithm chart.
The bounded derivatives of the exponential and logarithm
therefore give
\[
 \|\mathcal G_g(D+E)-\mathcal G_g(D)\|_\infty
       \le C\|E\|_\infty.
\]
There is no comb-length factor in this rowwise estimate.
With the prescribed correction norm below \(a_{\rm abs}\),
the full axial datum consequently satisfies
\begin{equation}\label{pw:absolute-datum-ceiling}
 \|\widehat B\|_\infty\le C_{\rm win}(M)a_{\rm abs}
       \le r_{\rm abs}/4 .
\end{equation}
This is absolute admission of the correction disc. No central
factor \(\sigma\) is assigned to an arbitrary correction,
and the axial gauge is not reconstructed from that correction.
Choose \(a_{\rm abs}\) after \(M\) and \(r_{\rm abs}\),
also below the ceiling \(a_I(M)\) of
\cref{htc:inner-substitution}. All contour logarithms and
the complex gauge covariance then stay in their fixed charts.
The constants do not depend on the depth: only the comb
length contains \(\sigma^{-1}\), and its link size supplies
the compensating \(\sigma\). This is an absolute-domain
check for the actual admitted window maps. It does not enlarge
a regular member's prescribed core: its membership and the
local Taylor ball of radius \(a\) remain separate requirements.
It also does not imply the proportional remote-datum bound
\(Ca_{\rm out}\) used below. That bound is needed to retain
\(a_{\rm out}\) in the exponentially small response from the
remote rows. If only an absolute remote bound is available,
its response must instead be retained as a separate tail.
For another family with radii \(r_B\) that vary with age,
the full-datum inclusion remains a separate hypothesis;
the central age estimate alone does not establish it.
A harmless enlargement of the radius ratio to at least one is allowed.
For a domain \(X\) containing \(z\), let \(D=2+\operatorname{diam}X\), in
the present units.

For the combined background and fluctuation we need, on the
central rows before choosing the old axial gauge, a logarithm
\(B\) with
\begin{equation}
 |B|\le C a_{\rm out}\sigma,\quad
 |\nabla B|\le C a_{\rm out}\sigma^2,\quad
 |\nabla B(x)-\nabla B(x')|
       \le C a_{\rm out}\sigma^{2+\alpha}|x-x'|^\alpha.
 \label{eq:pointed-window-remainder-PW6}
\end{equation}
Here is the passage from the common-frame field estimate to
these row bounds. Write \(\mathsf a\) for the bare bond logarithm
of the product field, with \(\eta_k^{-1}\mathsf a\) bounded by
\(Ca_{\rm out}\) in \(C^{1,\alpha}\) in level-\(k\) units. Translate
the central double-block carriers, including the one-row
translates used for the differences, to a common carrier.
The actual order-\(j\) row map on a smaller absolute sup ball
has first and second derivative norms bounded by \(CL^j\)
and \(CL^{2j}\), by directional Cauchy estimates. An old
step has length \(\sigma\) in upper units. The sup and
first-difference estimates therefore give the first two
bounds in \eqref{eq:pointed-window-remainder-PW6}. For two
rows \(r\) old spacings apart, applying the first derivative
to the difference of the increments costs
\(Ca_{\rm out}\sigma^{2+\alpha}|r|^\alpha\). The change of the
point at which that derivative is applied costs
\(Ca_{\rm out}^2\sigma^4|r|\) by the second derivative bound.
The latter fits the former on \(|r|\le C/\sigma\).

The decoupled field estimate is covariant for the input
connection. To use it here we must verify the spatial
regularity of that connection in the same frame as the
box field on the complete central carriers. A bound on
the gauge condition number controls the shifted link
sizes. The first differences also use the change of that
frame along the links. The localization uses the
\hyperlink{wf:relative-frame-restoration}{relative frame and its restoring identity}.

All statements in this section require the actual pair and its restrictions
to remain in the member's prescribed core.  The small chart ball used for
Taylor's formula must be contained in that core.  The radius \(a\) is a
fixed fraction of its curvature radius.  Throughout the admitted core
\(|\partial\mathbf U(p)-1|\le C a\eta_j^2\), so its own tent action obeys
\(|A(h_z,\mathbf U)|\le C a^2\), by the determinant identity and exact
tent plaquette count.  This is a bound on the whole core, not just its
contained chart ball.

In the axial gauge about \(z\), the datum has a local form
\begin{equation}
 \widehat B(c)=\ell_F(c)+r(c),\qquad
 \ell_{F,\mu}(x)=\sum_{\kappa<\mu}(x-z)_\kappa F_{\kappa\mu},
 \label{eq:pointed-window-remainder-PW7}
\end{equation}
where \(|F|\le C a_{\rm out}\sigma^2\), and
\begin{equation}
 |\widehat B(c)|\le C a_{\rm out}\sigma^2(1+|c-z|),\qquad
 |r(c)|\le C a_{\rm out}\sigma^{2+\alpha}(1+|c-z|)^2.
 \label{eq:pointed-window-remainder-PW8}
\end{equation}
These estimates are for the central rows. On the remote rows
the full axial datum is bounded by \(Ca_{\rm out}\) on the
admitted box chart. The distance of those rows from an inner
domain is at least \(c/\sigma\) in the local-scale metric.

Here are details of \eqref{eq:pointed-window-remainder-PW7}--\eqref{eq:pointed-window-remainder-PW8}.  Fixing the transports along a coordinate
comb to the identity expresses a remaining bond as a product of at most
\(C(1+|c-z|)\) transported plaquette holonomies.  With constant plaquette
strengths, the linear part of this product is exactly \(\ell_F\).
\eqref{eq:pointed-window-remainder-PW6} shows that a plaquette logarithm differs from its value at
\(z\) by at most
\[
 C\big[a_{\rm out}\sigma^{2+\alpha}(1+|c-z|)^\alpha
              +a_{\rm out}^2\sigma^3(1+|c-z|)\big].
\]
The second term includes differentiation of the non-Abelian commutator.
Conjugating a plaquette logarithm by a comb transport adds at most
\(C a_{\rm out}^2\sigma^3(1+|c-z|)\).
Since \(a_{\rm out}\) is bounded and \(\sigma^{1-\alpha}\le1\), both
errors are bounded by
\(C a_{\rm out}\sigma^{2+\alpha}(1+|c-z|)\).
Summing along the strip proves the second bound in \eqref{eq:pointed-window-remainder-PW8}.

For completeness, the finite-product estimate used here follows by
differentiating
\(P(t)=\prod_i e^{tX_i}\) twice: if \(\sum_i|X_i|\le c\), then
\(\sup_{0\le t\le1}|P''(t)|\le C(\sum_i|X_i|)^2\).
The analytic logarithm in this fixed matrix ball has bounded first two
derivatives, so
\[
 \left|\log\prod_i e^{X_i}-\sum_iX_i\right|
                         \le C\left(\sum_i|X_i|\right)^2.
\]
Its error here is at most
\(C a_{\rm out}^2\sigma^4(1+|c-z|)^2\), which is absorbed in \eqref{eq:pointed-window-remainder-PW8}.
The unexpanded product bound gives the first bound of \eqref{eq:pointed-window-remainder-PW8}.
All transports stay in a fixed matrix ball, because the original comb has
length \(O(\sigma^{-1})\) and its original bond logarithms are
\(O(a_{\rm out}\sigma)\).

We shall use the elementary weighted sum
\begin{equation}
 \sum_{c\in\mathbb Z^4}
       e^{-\mu\operatorname{dist}(c,X)}(1+|c-z|)^k
          \le C_{\mu,k}D^{k+4},\qquad k\ge0.
 \label{eq:pointed-window-remainder-PW9}
\end{equation}
Indeed \(|c-z|\le D+\operatorname{dist}(c,X)+C\).
The shell at distance between \(q\) and \(q+1\) is contained in a box
with at most \(C(D+q+1)^4\) bonds.  Bounding each shell by this larger
number and summing \(e^{-\mu q}(D+q+1)^{k+4}\) proves \eqref{eq:pointed-window-remainder-PW9}.
Only a polynomial in \(D\) is spent.  Every power needed below is fixed.

Let \(\widehat B^\circ\) agree with \(\widehat B\) on the
central rows and vanish on the remote ones. The weighted
first-response bound along the segment between the two data
gives
\[
 \|\Phi_\Lambda(\widehat B)-\Phi_\Lambda(\widehat B^\circ)\|_X
         \le C a_{\rm out}e^{-c/\sigma}.
\]
The support is remote from \(X\) in the tower metric. We use
the homogeneous sums on the central data.

Integrate the first response on the segment \(t\widehat B^\circ\), and then use
\eqref{eq:pointed-window-remainder-PW3}, \eqref{eq:pointed-window-remainder-PW8} and \eqref{eq:pointed-window-remainder-PW9}.  With a fixed integer \(s\), increased below when
necessary, this gives
\begin{equation}
 \|\Phi_\Lambda(\widehat B^\circ)\|_X\le C D^s a_{\rm out}\sigma^2,\qquad
 \|D\Phi_\Lambda(0)r\|_X
              \le C D^s a_{\rm out}\sigma^{2+\alpha}.
 \label{eq:pointed-window-remainder-PW10}
\end{equation}
The remote first-response bound gives the same first bound
for \(\Phi_\Lambda(\widehat B)\). Taylor's formula with integral remainder for \(\Phi_\Lambda\), followed by
\eqref{eq:pointed-window-remainder-PW4} and the product of the two weighted sums, gives
\begin{equation}
 \|\Phi_\Lambda(\widehat B^\circ)-D\Phi_\Lambda(0)\widehat B^\circ\|_X
       \le \frac{C D^s}{r_B}a_{\rm out}^2\sigma^4.
 \label{eq:pointed-window-remainder-PW11}
\end{equation}
The use of \eqref{eq:pointed-window-remainder-PW9} in \eqref{eq:pointed-window-remainder-PW11} is essential.  Using only
\(\sup|\widehat B|=O(a_{\rm out}\sigma)\) would produce the wrong power
\(\sigma^2\) in its right-hand side.

\subsubsection*{The finite-box linear response}

Take the depth-\(j\) truncation of the upper window tower.
Fix a constant curvature tensor \(F\) in one real colour
direction. Use the affine potential in \eqref{eq:pointed-window-remainder-PW7}
on the bare lattice, and on every inner bare bond take
\(U_t=\exp(i\eta_jt a_F)\). Prescribe the actual averages
of this field on every determining row, including the
finer rows of the tower. Every inner bond of a plaquette
crossing the exterior is a fixed level-zero row, since the
first inward layer is separated from that exterior.
Variations therefore act only at bonds where the full
current is defined. There the opposite abelian plaquette
contributions to the Wilson current cancel. Thus \(U_t\)
is critical for its own datum. On the fixed finite box it
is regular for small real \(t\), and regular uniqueness
identifies the box orbit at that datum.

On a level-\(q\) row \(c\), the first datum is
\(L^{q-j}a_F(c_-)\) in level-\(j\) units. The centered
block average reproduces the affine
transverse coordinates and the longitudinal line average
leaves \(a_{F,\mu}\) unchanged. The actual composition gives
the same result at every scale, as in \cref{abc:actual-affine-average}.
The literal first average gives the same affine datum.
In an endpoint block, a contour integral from the centre
has terms linear in a coordinate displacement or bilinear
in two distinct displacements. Each has zero centered
mean. Thus the endpoint contour terms cancel in the
linearized literal average and only the straight corridor
term remains. Differentiating the orbit identity gives the
linear pair \(L_F\) of this box, equal to the affine field
modulo a fine gradient, with zero linear current. Linearity
gives the same conclusion for complex colours. The pointed
Hessian kills that gradient, since the activity is gauge
invariant and its first derivative at the identity is zero.

On the central rows of level \(j\), the affine datum is
\(\ell_F\) of \eqref{eq:pointed-window-remainder-PW7} in
the same linear frame. Extend this central datum by zero
on the other rows in the expression below. If \(X\) is
at distance at least \(cR\) from those other rows, with
\(R\asymp\sigma^{-1}\), the difference from the full affine
datum is bounded by \(C|F|(1+R)^s\) and supported at that
distance. The first-response estimate on the same tower,
using the uniform weighted sum over tops, gives
\begin{equation}
 \|D\Phi_\Lambda(0)\ell_F-L_F\|_X
       \le C |F|(R+1)^s e^{-cR}.
 \label{eq:pointed-window-remainder-PW12}
\end{equation}
The central sum with its linear growth and this outer sum
also give \(\|L_F\|_X\le CD^s|F|\). The comparison uses
the literal affine datum on the prescribed finite tower.

For \(R\asymp\sigma^{-1}\),
\((R+1)^s e^{-cR}\le C_\alpha\sigma^\alpha\).  Combining
\eqref{eq:pointed-window-remainder-PW7}, \eqref{eq:pointed-window-remainder-PW10}--\eqref{eq:pointed-window-remainder-PW12} and the first-response bound on the remote datum, we obtain
\begin{equation}
 \begin{gathered}
 \Phi_\Lambda(\widehat B)=L_F+E_X,\qquad
 \|L_F\|_X\le C D^s a_{\rm out}\sigma^2,\\
 \|E_X\|_X
 \le C D^s\left[
 a_{\rm out}\sigma^{2+\alpha}
       +r_B^{-1}a_{\rm out}^2\sigma^4\right].
 \end{gathered}
 \label{eq:pointed-window-remainder-PW13}
\end{equation}
Domains too large for the interior condition in \eqref{eq:pointed-window-remainder-PW12} are put into the
exponentially long tail.  In applications their tree length is at least
\(c/\sigma-C\).

\subsubsection*{Taylor's formula for the actual member}

Let \(G_X\) denote a regular pointed member on \(X\), in its exponential
potential/current chart.  Assume its supremum on that core is
\(g_X=\mathsf g e^{-\kappa d_j(X)}\), and that its chart contains the
ball of radius \(a\).  A fixed change of constants allows the use of a
ball of radius \(2a\) and estimates on its half ball; this is the specified
room in the field domains.  Constant gauge transformations act by the
adjoint representation on both chart components.  There is no invariant
linear functional on \(\mathfrak{su}(2)\), so
\(DG_X(0)=0\).  In particular the second derivative of the background
does not enter the Hessian of this composition:
\begin{equation}
 D^2(G_X\circ\Phi_\Lambda)(0)[B,B']
       =D^2G_X(0)[D\Phi_\Lambda(0)B,D\Phi_\Lambda(0)B'].
 \label{eq:pointed-window-remainder-PW14}
\end{equation}
This assertion concerns regular members without stored colour spectators.

\begin{proposition}[Pointed nonlinear window remainder]
\label{thm:pointed-window-remainder}
Under the background and domain hypotheses above, for the near domains
there is a fixed integer \(s_1\) such that
\begin{equation}
 \begin{split}
 G_X(\Phi_\Lambda(\widehat B))-G_X(0)
    &=\tfrac12D^2G_X(0)[L_F,L_F]+\mathcal R_X,\\
 |\mathcal R_X|
    &\le C\mathsf g e^{-\kappa d_j(X)}D^{s_1}
       \big[\rho^2\sigma^{4+\alpha}+\rho^3\sigma^6\big].
 \end{split}
 \label{eq:pointed-window-remainder-PW15}
\end{equation}
The same estimate holds after absorbing the member's own tent-action
counterterm, with its own Hessian subtracted.
\end{proposition}

\begin{proof}
Put \(t=D^s\rho\sigma^2\), increasing \(s\) to cover every polynomial
factor in \eqref{eq:pointed-window-remainder-PW10}--\eqref{eq:pointed-window-remainder-PW13}.  First suppose \(t\) is less than a sufficiently
small fixed constant.  \eqref{eq:pointed-window-remainder-PW10} then places the composed pair in
the half chart ball of \(G_X\).  Cauchy's formula and Taylor's theorem give
\begin{equation}
 \left|G_X(P)-G_X(0)-\tfrac12D^2G_X(0)[P,P]\right|
                  \le Cg_X(\|P\|_X/a)^3,\qquad
 \|D^2G_X(0)\|\le Cg_X/a^2.
 \label{eq:pointed-window-remainder-PW16}
\end{equation}
In \eqref{eq:pointed-window-remainder-PW13}, divide by \(a\), use \(a/r_B\le C_0\), and expand the
difference of the two quadratic forms.  The product of \(L_F/a\) and
the first error term is bounded by
\(CD^{2s}\rho^2\sigma^{4+\alpha}\).
The product with the second error term is bounded by
\(CD^{2s}\rho^3\sigma^6\).  The square of the first error is smaller
than the first of these bounds.  The square of the second is bounded
by \(C t^4\), which is at most a constant times \(t^3\) in this case.
The mixed product is bounded by the sum of the squares.  Finally the
Taylor remainder in \eqref{eq:pointed-window-remainder-PW16} is at most \(C g_Xt^3\).
Increasing \(s_1\) proves \eqref{eq:pointed-window-remainder-PW15}.

If \(t\) exceeds that fixed small constant, use the actual core bound
\(|G_X(P)-G_X(0)|\le2g_X\), rather than Taylor's formula at \(P\).
The Hessian estimate still bounds its polynomial value at \(L_F\) by
\(C g_Xt^2\).  For \(t\ge c>0\),
\(1+t^2\le C_c t^3\), and this proves \eqref{eq:pointed-window-remainder-PW15} after increasing \(s_1\)
again.  Thus large radius ratios do not impose an upper bound on the age.

The coefficient of a member obeys
\(|b(G)|\le C\mathsf g/a^2\) by its Hessian moments.  The tent action
on its fixed local support has norm at most \(Ca^2\).
Consequently its counterterm is another holomorphic family of size
\(C\mathsf g\), with bounded support.  Apply the same proof to this
family and subtract.  The affine field and its curvature in the two
quadratic terms are the same.
\end{proof}

The exact own-counterterm identity in \eqref{ph:own-counterterm} of the
pointed Hessian cancellation theorem now applies to this last,
normalised family: its complete sum of pointed curvature tensors is zero.
A sum cut to the window is the complete sum minus the omitted domains.
The omitted domains have \(d_j(X)\ge c/\sigma-C\).
If the number through a cube of length \(q\) is at most \(Ce^{\kappa_Dq}\)
with \(\kappa>\kappa_D\), polynomial factors in \(D\le C M(d_j(X)+2)\)
are summable.  Thus \eqref{eq:pointed-window-remainder-PW15}, summed over near domains, gives
\begin{equation}
 C\mathsf g\big[\rho^2\sigma^{4+\alpha}+\rho^3\sigma^6\big],
 \label{eq:pointed-window-remainder-PW17}
\end{equation}
and completing the tensor sum costs at most
\(C\mathsf g\rho^2\sigma^4 e^{-c/\sigma}\).
If a far domain is not contained in any non-wrapping chart, its value has
only the norm bound and contributes an additional
\(C\mathsf g e^{-c/\sigma}\).  Such a term has no justified factor
proportional to the actual curvature: a flat field around a hole or a torus
cycle need not be a pure gauge.  It must be retained when a field-sensitive
tail estimate is claimed.

For the radius estimate used in a window, \(\rho\ge1\) permits both
exponentially small tails to be absorbed in \eqref{eq:pointed-window-remainder-PW17}.  Moreover
\begin{equation}
 \rho^2\sigma^{4+\alpha}+\rho^3\sigma^6
       \le2\rho^3\sigma^{4+\alpha},
       \qquad \rho\ge1,\quad \sigma\le1,\quad \alpha\le1.
 \label{eq:pointed-window-remainder-PW18}
\end{equation}
There are \(O(\sigma^{-4})\) level-\(j\) base points in one window of the
later scale.  Its remaining factor is therefore
\(C\mathsf g\rho^3\sigma^\alpha\).  With
\(\rho^2\le C(1+n-j)\) and \(\sigma=L^{-(n-j)}\), this is the summable
ordinary weight
\(C\mathsf g(1+n-j)^{3/2}L^{-\alpha(n-j)}\).
The central part of the input-domain condition holds at every
sufficiently old age. Indeed
\[
 \frac{C_{\rm ax}a_{\rm out}\sigma}{r_B}
       \le C_{\rm ax}C_0\rho\sigma
       \le C(1+n-j)^{1/2}L^{-(n-j)}.
\]
Choose a fixed age \(a_\star\) beyond which this is at most \(1/2\).
The remote datum must independently lie in the same half ball.
For the flat-reference windows this follows, at every age,
from \eqref{pw:absolute-datum-ceiling}; it is not a consequence
of the displayed age gain.
At each of the finitely many younger ages, the given core membership bounds
the member and its own counterterm by
\(C\mathsf g(1+\rho^2\sigma^4)\).
Here \(\rho\) is uniformly bounded and
\(\sigma\ge L^{-a_\star}\), so this is bounded by a fixed constant times
\(\mathsf g\rho^3\sigma^{4+\alpha}\).
This last argument uses the ordinary age bound; it is not valid for an
arbitrary unrelated pair \((\rho,\sigma)\).
Each site may carry a different Euclidean member: cancellation and the
estimate are performed at that site with that member's own coefficient.

\begin{proposition}[Vacuum-only window bound]
\label{cor:vacuum-only-window}
For an analytic gauge-invariant member without stored colour spectators,
no lattice covariance is needed to obtain
\begin{equation}
 |G_X(\Phi_\Lambda(\widehat B))-G_X(0)|
   \le C\mathsf g e^{-\kappa d_j(X)}D^{s_2}
                                      \rho^2\sigma^4.
 \label{eq:pointed-window-remainder-PW19}
\end{equation}
After the domains in a window have been summed, the bound is
\(C\mathsf g\rho^2\sigma^4\), with the same qualifications on remote,
non-wrapping charts and young ages as above.
\end{proposition}

\begin{proof}
Gauge invariance under constant colour rotations again gives \(DG_X(0)=0\),
independently of any spatial symmetry.  By \eqref{eq:pointed-window-remainder-PW10}, with
\(t=D^s\rho\sigma^2\), the local pair lies in the half chart ball whenever
\(t<c\).  Cauchy's estimate for the second-order Taylor remainder then gives
\(C g_Xt^2\).  For \(t\ge c\), use the core bound \(2g_X\le C_cg_Xt^2\).
This proves \eqref{eq:pointed-window-remainder-PW19}, and the domain entropy absorbs its fixed polynomial
in \(D\).  A remote domain not admitting a chart contributes only
\(C\mathsf g e^{-c/\sigma}\); it is bounded by
\(C\mathsf g\rho^2\sigma^4\) for \(\rho\ge1\).
If the central input-disc condition is unavailable at finitely
many young ages, while the remote condition is supplied as above,
the core bound is absorbed since
\(\rho\ge1\) and \(\sigma\ge L^{-a_\star}\).
Finally the \(O(\sigma^{-4})\) sites in the window cancel \(\sigma^4\),
leaving the ordinary weight \(\rho^2\) for a defect or a remainder.
\end{proof}


% END MODULE chapters/pointed_window_remainder_v6.tex

% BEGIN MODULE chapters/localized_gaussian_comparison_v7.tex
\subsection{Localized complex Gaussian comparison}
\label{ch:localized-gaussian-comparison}

The finite regular expansion uses a normalized Gaussian functional on a
local set of fluctuation coordinates. Its auxiliary Gaussian can have
as many coordinates as the surrounding volume. We prove its comparison
bound without paying for that auxiliary dimension. The argument permits
sharp, measurable cutoffs; it never shifts their arguments into the
complex domain.

\subsubsection*{The local matrices and the normalized functional}

All matrix norms below are operator norms for the Euclidean complex
space. Transposition is bilinear transposition, not conjugate
transposition. Let \(C=C^{\mathsf T}\) be an invertible complex
\(m\)-by-\(m\) matrix and let \(\Gamma\) be a complex
\(m\)-by-\(n\) matrix, where \(m\ge1\) and \(n\ge0\).
Assume real reference matrices \(C_0=C_0^{\mathsf T}>0\) and
\(\Gamma_0\), and one constant \(K\ge1\), such that
\begin{equation}
 \|C_0\|+\|C_0^{-1}\|+\|\Gamma_0\|\le K,
 \qquad
 \|C-C_0\|+\|\Gamma-\Gamma_0\|\le\varepsilon.
 \label{eq:localized-gaussian-comparison-FG1}
\end{equation}
The number \(\varepsilon\) will be smaller than a constant depending
only on \(K\). No dimension enters that choice. Define
\begin{equation}
 D=C+C\Gamma\Gamma^{\mathsf T}C,
 \qquad D_0=C_0+C_0\Gamma_0\Gamma_0^{\mathsf T}C_0.
 \label{eq:localized-gaussian-comparison-FG2}
\end{equation}
The real matrix \(D_0\) is positive definite. In the small neighborhood
in \eqref{eq:localized-gaussian-comparison-FG1}, normalized Gaussian densities are defined by
\begin{equation}
 d\mu_D(b)=(2\pi)^{-m/2}(\det D)^{-1/2}
             e^{-b^{\mathsf T}D^{-1}b/2}\,db,
 \qquad b\in\mathbb R^m.
 \label{eq:localized-gaussian-comparison-FG3}
\end{equation}
The square root is the holomorphic continuation of the positive root
at the reference. We justify both convergence and normalization below.

The auxiliary representation of the localized functional is
\begin{equation}
 \begin{split}
 {\cal J}_{C,\Gamma}[F]
  =\int d\mu_{I_n}(\xi)\,
       e^{-\xi^{\mathsf T}\Gamma^{\mathsf T}C\Gamma\xi/2}
       \int d\mu_C(b)\,e^{b^{\mathsf T}\Gamma\xi}F(b).
 \end{split}
 \label{eq:localized-gaussian-comparison-FG4}
\end{equation}
For \(n=0\) the outer integral is absent. An integrand \(F\) need
only be Borel measurable on \(\mathbb R^m\). Suppose
\begin{equation}
 |F(b)|\le F_0 e^{\rho|b|^2},\qquad F_0\ge0,\quad\rho\ge0.
 \label{eq:localized-gaussian-comparison-FG5}
\end{equation}

\begin{proposition}[Local Gaussian comparison]
\label{thm:localized-complex-Gaussian}
There are positive constants \(\varepsilon_0(K),\rho_0(K)\) and
\(A_K\), independent of \(m,n\), such that \eqref{eq:localized-gaussian-comparison-FG1},
\(\varepsilon\le\varepsilon_0\) and \(\rho\le\rho_0\) imply
absolute convergence of \eqref{eq:localized-gaussian-comparison-FG4},
\begin{equation}
 {\cal J}_{C,\Gamma}[F]=\int F\,d\mu_D,
 \qquad {\cal J}_{C,\Gamma}[1]=1,
 \qquad
 |{\cal J}_{C,\Gamma}[F]|
       \le F_0e^{A_K(\varepsilon+\rho)m}.
 \label{eq:localized-gaussian-comparison-FG6}
\end{equation}
The last estimate also bounds \(\int |F|\,|d\mu_D|\).
In particular its exponential cost uses the local dimension \(m\),
even when the auxiliary dimension \(n\) grows with the full volume.
\end{proposition}

\subsubsection*{Accretivity, integration and the determinant branch}

\begin{proof}
The joint precision in \eqref{eq:localized-gaussian-comparison-FG4} is the symmetric block matrix
\begin{equation}
 Q=\begin{pmatrix}
 C^{-1}&-\Gamma\\
 -\Gamma^{\mathsf T}&I_n+\Gamma^{\mathsf T}C\Gamma
 \end{pmatrix}.
 \label{eq:localized-gaussian-comparison-FG7}
\end{equation}
At the real reference it factors as
\begin{equation}
 Q_0=L_0^{\mathsf T}
       \begin{pmatrix}C_0^{-1}&0\\0&I_n\end{pmatrix}L_0,
 \qquad
 L_0=\begin{pmatrix}I_m&-C_0\Gamma_0\\0&I_n\end{pmatrix}.
 \label{eq:localized-gaussian-comparison-FG8}
\end{equation}
Both \(\|L_0\|\) and \(\|L_0^{-1}\|\) are at most
\(1+K^2\). Thus \(Q_0\ge c_K I\) for a positive \(c_K\).
The resolvent identity for \(C^{-1}-C_0^{-1}\), followed by
expanding the products in \eqref{eq:localized-gaussian-comparison-FG7}, gives
\(\|Q-Q_0\|\le A_K\varepsilon\) for small
\(\varepsilon\); the constants depend only on \(K\).
For real \((b,\xi)\), consequently,
\[
 \operatorname{Re}(b,\xi)^{\mathsf T}Q(b,\xi)
                   \ge(c_K/2)(|b|^2+|\xi|^2).
\]
Choose \(\rho_0<c_K/8\). This proves absolute convergence of the
joint integral with \eqref{eq:localized-gaussian-comparison-FG5}, hence Fubini's theorem applies. It does
not yet give the desired dimension bound; taking its absolute value
before integrating \(\xi\) would pay unnecessarily for \(n\).

Set \(A=I_n+\Gamma^{\mathsf T}C\Gamma\). It is a small perturbation
of the real positive matrix \(I_n+\Gamma_0^{\mathsf T}C_0\Gamma_0\).
Integrate \(\xi\) first. The normalized Gaussian exponential formula
gives the factor
\[
 (\det A)^{-1/2}
       \exp\{b^{\mathsf T}\Gamma A^{-1}\Gamma^{\mathsf T}b/2\}.
\]
This formula for a complex linear force does not require a contour
change. One proof normalizes the complex Gaussian first and
differentiates its exponential transform in the force. Integration
by parts, justified by accretivity, gives
\(\nabla_\ell G=A^{-1}\ell G\), with \(G(0)=1\), and hence
\(G(\ell)=e^{\ell^{\mathsf T}A^{-1}\ell/2}\).

For completeness, the normalization used in that proof follows
without a complex change of variables. On an accretive symmetric
precision neighborhood the real-space integral is holomorphic by a
fixed Gaussian majorant. It equals \((\det A)^{-1/2}\) at every
real positive symmetric matrix there by real orthogonal
diagonalization and the one-dimensional real Gaussian integral.
Successive one-variable identity theorems in the independent
symmetric matrix entries extend the equality to the complex
neighborhood, with the root continued from its positive value.
The real symmetric matrices have a full set of real coordinates
for this argument; no identity is being extended from a
lower-dimensional background graph.

The elementary inverse and determinant identities are
\begin{equation}
 \begin{split}
 D^{-1}&=C^{-1}-\Gamma
             (I_n+\Gamma^{\mathsf T}C\Gamma)^{-1}\Gamma^{\mathsf T},\\
 \det D&=\det C\,\det(I_n+\Gamma^{\mathsf T}C\Gamma).
 \end{split}
 \label{eq:localized-gaussian-comparison-FG9}
\end{equation}
The inverse identity follows by multiplication with \eqref{eq:localized-gaussian-comparison-FG2}.
For the determinant identity, eliminate either off-diagonal block in
\(\bigl(\begin{smallmatrix}I_m&-C\Gamma\\
\Gamma^{\mathsf T}&I_n\end{smallmatrix}\bigr)\); this proves
\(\det(I_m+C\Gamma\Gamma^{\mathsf T})
=\det(I_n+\Gamma^{\mathsf T}C\Gamma)\).
The roots in \eqref{eq:localized-gaussian-comparison-FG9} agree at the real reference, and agree throughout
the connected small perturbation neighborhood by continuity and
nonvanishing. The remaining \(b\) integral is precisely \eqref{eq:localized-gaussian-comparison-FG3}.
This proves the first identity in \eqref{eq:localized-gaussian-comparison-FG6}, including for Borel \(F\),
since no complex argument of \(F\) was introduced. The same
normalization argument gives \(\int d\mu_D=1\).

\subsubsection*{The local-dimension estimate}

The reference bounds imply
\(D_0\ge C_0\), \(\|D_0\|\le A_K\) and
\(\|D_0^{-1}\|\le K\). Expanding \eqref{eq:localized-gaussian-comparison-FG2} yields
\begin{equation}
 T=D_0^{-1/2}(D-D_0)D_0^{-1/2},\qquad
 \|T\|\le A_K\varepsilon=:\eta.
 \label{eq:localized-gaussian-comparison-FG10}
\end{equation}
Shrink \(\varepsilon_0\) so that \(\eta\le1/4\).
Then
\(D^{-1}=D_0^{-1/2}(I+T)^{-1}D_0^{-1/2}\) and
\(\|(I+T)^{-1}-I\|\le\eta/(1-\eta)\).
For the real symmetric matrix
\[
 R=D_0^{1/2}\bigl(\operatorname{Re}D^{-1}-2\rho I_m\bigr)D_0^{1/2}
\]
we obtain \(\|R-I_m\|\le A_K(\varepsilon+\rho)=:u\).
Shrink the constants further so that \(u\le1/2\).
Taking the absolute value of the \emph{marginal} density, and then
integrating the real positive Gaussian, gives
\begin{equation}
 \int e^{\rho|b|^2}|d\mu_D(b)|
   =|\det(I+T)|^{-1/2}(\det R)^{-1/2}
   \le(1-\eta)^{-m/2}(1-u)^{-m/2}.
 \label{eq:localized-gaussian-comparison-FG11}
\end{equation}
Every singular value of \(I+T\) is at least \(1-\eta\), and
every eigenvalue of \(R\) is at least \(1-u\). Finally
\(-\log(1-v)\le2v\) for \(0\le v\le1/2\) proves \eqref{eq:localized-gaussian-comparison-FG6}.
\end{proof}

\subsubsection*{Sharp failures, holomorphic parameters, and scope}

Partition the real coordinates into disjoint bond blocks. If a set
\(P_0\) of distinct blocks fails the tests \(|b_\beta|<\mathfrak b\),
then \(\sum_{\beta\in P_0}|b_\beta|^2
\ge|P_0|\mathfrak b^2\). Therefore, for \(\gamma\ge0\) with
\(\rho+\gamma/2\le\rho_0\), \eqref{eq:localized-gaussian-comparison-FG11} gives
\begin{equation}
 \int \prod_{\beta\in P_0}\mathbf1_{\{|b_\beta|\ge\mathfrak b\}}
       e^{\rho|b|^2}|d\mu_D|
 \le \exp\{A_K(\varepsilon+\rho+\gamma/2)m
                     -\gamma\mathfrak b^2|P_0|/2\}.
 \label{eq:localized-gaussian-comparison-FG12}
\end{equation}
This is a product-of-failures estimate, with one Gaussian tail per
distinct tested bond. Multiplication by additional indicators only
improves the bound.

If the matrices and the potential part of an integrand depend
holomorphically on external parameters, and the sharp cutoff sets
remain fixed, the same estimates give a parameter-independent
Gaussian majorant on each smaller parameter neighborhood.
Fubini and Morera then prove holomorphy of the integrals. No
holomorphy of a characteristic function in its real field argument
is asserted or required.

The rescaled structure \((\tau^{-1}C,\tau^{1/2}\Gamma)\) has
effective covariance \(\tau^{-1}D\). For real \(\tau\) in a fixed
compact subinterval of \((0,\infty)\), rescale the real reference
in the same way; all constants remain uniform. A small complex
displacement of \(\tau\) is included in the matrix perturbation.

In the regular step, let \(N_c(Z)\) be the number of coarse
localization cubes. The local free-bond color-coordinate count obeys
\(m\le c(LM)^4N_c(Z)\). Consequently \eqref{eq:localized-gaussian-comparison-FG6} supplies the factor
\(\exp\{A_K(LM)^4(\varepsilon+\rho)N_c(Z)\}\), and \eqref{eq:localized-gaussian-comparison-FG12}
supplies the failed-cutoff tail. If a Gaussian comparison instead uses
the symbol \(N(Z)\) for the local coordinate count itself, its bound is
\(\exp\{A_K(\varepsilon+\rho)N(Z)\}\), with no second factor
\((LM)^4\). These are two conventions for the same local cost. The matrix gap and kernel closeness
are the actual A2--A3 inputs; this section proves the resulting
finite-dimensional integration estimate.

Completeness is separate. Equality of this local functional with a
marginal of one whole-torus Gaussian requires the actual completeness
identity for its effective covariance. In the source construction
that identity is supplied on complete physical pairs and their
physical-data continuation. Arbitrary independent-current cores
need only their own local holomorphic functionals and the bounds
proved here. This theorem does not manufacture a global covariance
on those cores, verify nonlinear background/core estimates, or
establish a Gaussian law for actual averaging disintegrations.


% END MODULE chapters/localized_gaussian_comparison_v7.tex

% BEGIN MODULE chapters/local_formation_program.tex
\subsection{Local interactions and their source response}
\label{ch:local-formation-program}
\label{ch:unmarked-polymer-majorant}
\label{ch:finite-source-insertion}
\label{ch:renewed-source-pointing}
\label{ch:window-potential-formation}
\label{ch:linear-source-bootstrap}
\label{ch:old-activity-hessian}


The local part of a renormalization step has several distinct operations.
A function of a fluctuation is first localized, its Gaussian integral is
expanded, the connected logarithm is pointed at bonds, and an old
gauge-invariant function is rewritten on a new constrained background.
These operations have different domains.  Their order is part of the
definition of the step.

We work at a fixed finite cutoff and volume until uniform estimates have
been established.  A local pair \(P=(U,J)\) consists of the admitted
potential data and current data on the specified local domain.  Those
coordinates may be independent.  A complete physical pair is denoted
\(p_W\); its current is the one induced by the actual constrained
background belonging to the datum \(W\).  At such pairs the complete
Gaussian structure reproduces the physical fluctuation integral.
An identity on the set of pairs \(p_W\) is not an identity on an open
set of independent pairs \(P\).

We use three fluctuation coordinates.  The free Gaussian coordinates are
\(B\); the constrained physical fluctuation is \(\phi(B)\); and the
background response supplies the actual field \(U_P(B)\).  Their typical
sizes have the form
\begin{equation}
 g=\theta^{-1/2},\qquad
 \alpha=g\,A(\theta),\qquad
 \delta=s\alpha,\qquad 0<s\le s_\star.
 \label{fm:three-radii}
\end{equation}
The analytic size \(A(\theta)\), the cutoff in the coordinates \(B\), and
the physical radius \(\delta\) are separate quantities.  In particular a
bound on a cutoff box in \(B\) does not furnish a holomorphic continuation
to a ball of radius \(\alpha/g\).  The larger ball used below is obtained
from the actual unscaled constraint chart and the admitted datum
polydiscs of the background response.

At a fixed real reference \(\theta\), two scalar-source presentations are
useful.  In the variable-profile presentation the real action weight
\(x\), the calibration, Gaussian, determinant branches and cutoffs are
fixed; the profile is carried as an older potential member.  Its new
source energy is the difference of the localized fluctuation functions.
In a constant-shift presentation one may instead change the action
weight and compensate the free action, numerical terms and scalar
counterterm.  Differentiating that presentation changes more than the
older potential.  We use the first presentation for the insertion and
source-gradient estimates in this section.  Numerical or constant-shift
identities require their own reference-change rule.

For clarity, the primitive data needed for the local argument are the
following.
\begin{enumerate}
\item The actual local Gaussian matrices have a uniform real reference
gap and a small complex perturbation on every admitted local core and
every enlarged decoupling disc.  Their local scalar-coordinate dimension
is at most \(c(LM)^4\) times the coarse-cube count.  The auxiliary
dimension may be the volume of the surrounding torus.
\item The background and window maps are jointly holomorphic on the
specified datum and current domains, including all actual local
substitutions, fine brackets and fluctuation amplifications used below.
Their first-response, compatible-boundary and spatial regularity bounds
are uniform in cutoff and volume.
\item Fine restrictions are hereditary; at all interpolation parameters
equal to one they reproduce the original local function.  A component
separated from the declared root and containing a differenced cube gives
zero.  The actual coarse collar is a fixed-width union of neighboring
cubes, rather than a rectangular bounding box of a thin domain.
\item The actual pointing is formed from its free-potential, cutoff and
interaction legs.  The interaction leg uses its actual bond densities.
We prove a sufficient explicit convention below, and distinguish it
from an identification with a separately prescribed pointing.
\end{enumerate}
These are matrix, map and support assertions.  They do not assume the
activity bound, its source derivative, the many-scale recurrence, or the
old-origin Hessian concluded below.

The roles of the six constructions are successive.  The unmarked
expansion first provides a spare potential slot.  Cauchy's formula in
that slot gives derivatives in the norm of the inserted potential.
The pointing argument then differentiates the actual bond density.
Window-first normalization supplies the old potential direction with
its small radius factor before any fine bracket is applied.  The
source-gradient recurrence and the larger fluctuation-domain argument
finally propagate the first jet and bound the old Hessian.

There is one common normalization issue throughout.  If \(H\) is an old
action member, the new fluctuation direction is
\begin{equation}
 v(P,B)=H(U_P(B))-H(U_P(0)).
 \label{fm:carried-action-split}
\end{equation}
The omitted value \(H(U_P(0))\) is carried by the old background action.
It is not subtracted a second time from the logarithmic fluctuation
integral.  Equally, when a new scalar coefficient is extracted from a
pointed member, it is that member's own tent action that is subtracted.
A coefficient obtained from a different pointing or a different
background readout cannot be substituted by name.

All norm bounds in this section are on fixed admitted local cores.
When a complex parameter is varied, a cutoff remains fixed at its real
reference value and is a function of real integration coordinates.
The reference measure, geometric indices and filing rule are also
fixed.  A moving cutoff or changing history fibre must first be
expressed over a common reference operation before a norm-holomorphy
argument applies.  The last section records precisely how the local
estimates interact with a finite history and its retained inner Wilson
weight.

\subsubsection{The actual unmarked expansion}


We derive the activity and local logarithm bounds from the finite
Mayer, cutoff and coarse-decoupling formulas.  The Gaussian estimate
is used at a potential budget larger than the physical one.  No
logarithm or nonvanishing assertion is needed at that larger budget.
Only after the activities have been estimated do we apply the
proved hard-core polymer theorem.

All bounds below are local on the specified complex potential/current
cores.  At complete physical pairs the completeness identity also
identifies their gas with the regular fluctuation integral.  No
global Gaussian on an arbitrary independent-current pair is assumed.

\paragraph*{The actual formulas and their analytic inputs}

Fix a finite four-dimensional periodic cube lattice, with fine
localization cubes of side \(M\) and coarse cubes of side \(LM\).
Both grids tile the torus.  Lengths \(d_f\) and \(d_c\) are tree
lengths in fine and coarse cube units.  There are a fixed finite
number of potential types: the free, intrinsic, already formed
older, and optional inserted potential groups may each be
combined within their type before the Mayer expansion.
Thus this type count is independent of the number of older
levels or source profiles.
A tagged fine domain \(Y\) includes
its type in its name; geometry ignores this tag.
For the actual local potentials assume
\begin{equation}
 \begin{aligned}
 |V(Y;P,B)|&\le a_Ye^{-\lambda_Vd_f(Y)}
       \left(1+M^{-4}\sum_{b\subset Y}|B(b)|^2\right),
 \qquad a_Y\le a,\\
 a&=\sum_{\text{types}}\|\mathbf V^{\text{type}}\|_{\lambda_V}.
 \end{aligned}
 \label{eq:local-formation-UM1}
\end{equation}
The potentials are holomorphic on their declared complex fluctuation
boxes and external cores.  If a potential is defined only on a box,
the corresponding sharp cutoff is kept wherever that potential occurs.
The free potential, when present, is defined on all real fluctuation
coordinates.

Let \(\mathbf D\) be a selected set of tagged domains.
Write \(Y_0\) for the union of the free-bond sets requiring a field
box among those selections, and let \(P_0\) be a set of distinct
failed-cutoff bonds outside \(Y_0\).
The exact integrand is
\begin{equation}
 F_{\mathbf D,P_0}
   =(-1)^{|P_0|}\chi_{Y_0}
       \prod_{b\in P_0}(1-\chi_b)
       \prod_{Y\in\mathbf D}(e^{V(Y)}-1),
 \qquad \chi_b=\mathbf1_{\{|B(b)|<\mathfrak b\}}.
 \label{eq:local-formation-UM2}
\end{equation}
The cutoff product on \(Y_0\) makes every selected factor
well defined.  Outside that product its integrand is set to zero.
This convention introduces no complex argument into a characteristic
function.

For any set \(S\) of fine cubes, let \(\mathcal N_f(S)\)
consist of \(S\) and its fixed first layer of neighbors in the
given cubical adjacency.  Define \(S^\sharp\) to be the union
of the coarse cubes containing at least one fine cube of
\(\mathcal N_f(S)\).  This is a local cube union, not a
rectangular bounding box.  A bond uses its incident fine cubes
in this definition.  The collar preserves unions.  A fixed
number of additional neighboring cube layers is permitted only
if it is part of the assigned collar in every structural
identity; its fixed width is then included in the constants.

For the index \(\iota=(\mathbf D,P_0)\), let \(Z_0(\iota)\)
be its actual fine support, including the globally defined selected
potentials and the failed bonds.
Its coarse collar is \(Z_0^\sharp\).
For \(Z\supset Z_0^\sharp\), the actual terms and activities are
\begin{equation}
 \begin{split}
 T(Z,\iota)&=
 \prod_{\Delta\in Z\setminus Z_0^\sharp}
       \int_0^1ds_\Delta\,\partial_{s_\Delta}\,
       {\cal J}(Z_0;s_Z)[F_{\mathbf D,P_0}],\\
 H(Z)&=\sum_{\iota}T(Z,\iota).
 \end{split}\label{eq:local-formation-UM3}
\end{equation}
Here \(s_Z=1\) on \(Z_0^\sharp\), has the displayed parameters
on the remaining cubes of \(Z\), and is zero elsewhere.
Polymers \(Z\) are nonempty wall-connected coarse domains.
The empty selection contributes only the gas's leading one.
It is omitted from \(H\): on a nonempty \(Z\), its term has a
derivative of the normalized constant functional
\({\cal J}[1]=1\), and is zero.
Disconnected regions factor into their wall-components rather than
being counted again as single polymers.

The local functional \({\cal J}\) has the actual matrix
representation and reference bounds of the local Gaussian
comparison theorem.  Uniformly on independent coarse parameter
discs
\begin{equation}
 |s_\Delta|<R_s,\qquad R_s>2,
 \label{eq:local-formation-UM4}
\end{equation}
its matrix perturbation is at most \(\varepsilon\), its gap
constant is \(K\), and its local real scalar-coordinate dimension
obeys
\begin{equation}
 m(Z_0;s_Z)\le c_m(LM)^4N,\qquad N=\#Z.
 \label{eq:local-formation-UM5}
\end{equation}
The auxiliary Gaussian dimension is unrestricted.
For an integrand bounded by \(F_*e^{\rho|B|^2}\), including
distinct failed tests \(P_0\), the proved comparison gives
\begin{equation}
 |{\cal J}[F]|
 \le F_*
 \exp\!\left\{A_K(\varepsilon+\rho+\gamma/2)m
                    -\gamma\mathfrak b^2|P_0|/2\right\}
 \label{eq:local-formation-UM6}
\end{equation}
when \(\gamma>0\), \(\varepsilon\le\varepsilon_0(K)\), and
\(\rho+\gamma/2\le\rho_0(K)\).
This is the marginal local-coordinate estimate; no factor using
the exterior auxiliary dimension is allowed.
The potentials, matrices and integrals are jointly holomorphic
in the scalar potential couplings and \(s\), with the fixed
cutoff sets just described.

The remaining structural hypotheses are the actual heredity,
normalized-functional and wall-factorization identities.
They identify \eqref{eq:local-formation-UM3} with the finite expansion on complete
physical pairs and define coherent local terms on every local
complex core.  These hypotheses concern the assigned operators
and supports.  They are not consequences of a numerical
Gaussian bound.

\paragraph*{The two geometric counts}

Let \(Y^\sharp\) be the locally defined coarse collar above.
All neighborhood widths are fixed independently of \(L\)
and of the diameter or shape of \(Y\).
There is a fixed geometric constant \(C_h\) such that
\begin{equation}
 \#Y^\sharp\le C_h\left(1+\frac{d_f(Y)}L\right).
 \label{eq:local-formation-UM7}
\end{equation}
To prove this, take a fine covering tree for \(Y\), with length
within one of its infimum, and traverse every edge twice.
After adding a bounded initial segment, this gives a closed
walk of length at most \(C(d_f(Y)+1)\) passing within a fixed
fine-cube distance of every cube of \(Y\).
Split the walk into segments of length at most \(L\).
Each segment and its fixed fine collar meet a bounded number
of coarse cubes, independent of \(L\ge2\).
The number of segments is at most \(C(1+d_f(Y)/L)\).
This proves \eqref{eq:local-formation-UM7}, also on the periodic graph.
A failed bond has a coarse collar with at most a fixed number
of cubes.  Therefore, writing
\[
 r=\#\mathbf D,\qquad v=\#P_0,\qquad
 h=\#(Z\setminus Z_0^\sharp),
\]
and increasing \(C_h\) once,
\begin{equation}
 N\le C_h\left[
       r+\frac1L\sum_{Y\in\mathbf D}d_f(Y)+v+h\right].
 \label{eq:local-formation-UM8}
\end{equation}
There is no term proportional to the surrounding volume.

We also use the rooted fine-domain entropy bound.  For every
\(\xi>\kappa_f\),
\begin{equation}
 S_\xi:=
  \sup_{\text{fine cube }q}
       \sum_{\text{tagged }Y\ni q}e^{-\xi d_f(Y)}<\infty.
 \label{eq:local-formation-UM9}
\end{equation}
One direct counting proof rounds a covering tree to the
bounded-degree graph of fine cubes, traverses its edges twice
starting at \(q\), and records which neighboring cubes are
selected in \(Y\).
The walk has length at most \(C(d_f(Y)+1)\).
Its choices and the selection flags are bounded by a fixed
constant to that power.
This proves exponential rooted entropy, uniformly on the torus.
The fixed number of type tags changes that constant only.
In particular, if \({\cal Y}(Z)=\{Y:Y^\sharp\subset Z\}\),
\begin{equation}
 \sum_{Y\in{\cal Y}(Z)}e^{-\xi d_f(Y)}
        \le c_fL^4S_\xi N .
 \label{eq:local-formation-UM10}
\end{equation}
This follows by charging a domain to any one of its fine
cubes and summing over all fine cubes in \(Z\); overcounting
only enlarges the bound.
The number of free bonds whose collars fit in \(Z\) is at
most \(n_bN\), where
\begin{equation}
 n_b\le c_b(LM)^4 .
 \label{eq:local-formation-UM11}
\end{equation}
The bond count in \eqref{eq:local-formation-UM11} and the scalar-coordinate dimension
in \eqref{eq:local-formation-UM5} are distinct counts.  Neither is counted twice.

\paragraph*{A Gaussian-only budget and repeated Schwarz}

Choose a positive number \(W_0\), independent of cutoff and volume,
small enough for a Gaussian tilt, and choose
\(\lambda_0>\kappa_f\).
This is not yet the physical potential budget.
For every selected tagged domain introduce its own coupling
\(\zeta_Y\).  For \(a>0\) use the radii
\begin{equation}
 {\cal R}_Y=\frac{W_0}{a}
             e^{(\lambda_V-\lambda_0)d_f(Y)} .
 \label{eq:local-formation-UM12}
\end{equation}
We will require \(a\le W_0/2\), so every radius exceeds one.
On the product of these discs, the local density bound in \eqref{eq:local-formation-UM9}
gives
\begin{equation}
 \begin{split}
 \sum_{Y\in\mathbf D}|\zeta_YV(Y;P,B)|
 &\le W_0\sum_{Y\in\mathbf D}e^{-\lambda_0d_f(Y)}
             \left(1+M^{-4}\sum_{b\subset Y}|B(b)|^2\right)\\
 &\le C_pW_0L^4N+C_pW_0M^{-4}|B|^2 .
 \end{split}\label{eq:local-formation-UM13}
\end{equation}
The constant \(C_p\) includes the finite tags and their
rooted sums, but no number of selected domains or global
volume.
The quadratic term uses the rooted sum at each bond.
The constant term uses the fine-cube count in \eqref{eq:local-formation-UM10}.

Choose \(W_0,\gamma\) so that
\begin{equation}
 C_pW_0M^{-4}+\gamma/2\le\rho_0(K).
 \label{eq:local-formation-UM14}
\end{equation}
For complex \(\zeta_Y\) the Mayer product is bounded by
\(2^r\exp(\sum_Y|\zeta_YV(Y)|)\).
Apply \eqref{eq:local-formation-UM6} with the tilt in \eqref{eq:local-formation-UM13}.
Uniformly on both parameter polydiscs its integral is at most
\begin{equation}
 2^r e^{C_0N}\,p^v,\qquad
 p=e^{-\gamma\mathfrak b^2/2},
 \quad
 C_0=C_pW_0L^4+
 A_Kc_m(LM)^4
       \big(\varepsilon+C_pW_0M^{-4}+\gamma/2\big).
 \label{eq:local-formation-UM15}
\end{equation}
The failed tests are distinct bond blocks, so the tail is \(p^v\).
The Gaussian theorem was applied only to a fixed real integration
space with measurable cutoffs.
There is no zero-freeness assertion on this larger coupling
polydisc.

The integral vanishes whenever any selected \(\zeta_Y=0\).
Repeated Schwarz on the polydisc consequently divides the
bound by \(\prod_Y{\cal R}_Y\) at \(\zeta_Y=1\).
For clarity, division by \(\prod_Y\zeta_Y\) has removable
singularities; apply the maximum principle on smaller
polydiscs and let their radii increase to \eqref{eq:local-formation-UM12}.
This proves the product factor without a factorial.
Mixed coarse Cauchy on circles centered at
\(s_\Delta\in[0,1]\), of radii tending to \(R_s-1\),
then gives the literal term bound
\begin{equation}
 |T(Z,\mathbf D,P_0)|
 \le e^{C_0N}
       \left(\frac{2a}{W_0}\right)^r
       e^{-(\lambda_V-\lambda_0)
                           \sum_{Y\in\mathbf D}d_f(Y)}
       p^v (R_s-1)^{-h}.
 \label{eq:local-formation-UM16}
\end{equation}
The coupling derivatives and the coarse derivatives commute
by joint holomorphy and the same Gaussian domination.
If \(a=0\), all nonempty Mayer selections are zero.
Apply the same estimate to the cutoff-only terms, with
\(r=0\); this is also the limiting interpretation of \eqref{eq:local-formation-UM16}.

\paragraph*{The activity bound, including its linear amplitude}

Fix a desired activity rate \(\Lambda>0\).
For wall-connected coarse polymers there is a fixed \(c_\ell\)
such that
\begin{equation}
 \ell(Z):=d_c(Z)+5\le c_\ell N .
 \label{eq:local-formation-UM17}
\end{equation}
A cube-adjacency spanning tree proves \(d_c(Z)\le CN\);
the added five is absorbed because \(N\ge1\).
Choose
\begin{equation}
 K_*=2+c_\ell\Lambda,\qquad
 \Theta=C_h(C_0+K_*).
 \label{eq:local-formation-UM18}
\end{equation}
Choose one \(\sigma_0>\kappa_f\), and impose the explicit
rate conditions
\begin{equation}
 \lambda_V-\lambda_0\ge \Theta/L+\sigma_0,
 \qquad \log(R_s-1)\ge\Theta .
 \label{eq:local-formation-UM19}
\end{equation}
Using \eqref{eq:local-formation-UM8} to pay for
\(e^{(C_0+K_*)N}\), \eqref{eq:local-formation-UM16} becomes
\begin{equation}
 |T(Z,\mathbf D,P_0)|
 \le e^{-K_*N}
       \prod_{Y\in\mathbf D}
          \left(\frac{2a}{W_0}e^\Theta
                         e^{-\sigma_0d_f(Y)}\right)
       \prod_{b\in P_0}(pe^\Theta).
 \label{eq:local-formation-UM20}
\end{equation}
The unused factor
\(((R_s-1)^{-1}e^\Theta)^h\) is at most one.
Every selection has thereby paid its own part of the hull
cost.  This is the place where the coarse hull estimate
retains the favorable \(d_f(Y)/L\).

At fixed \(Z\), dropping the restriction
\(P_0\cap Y_0=\emptyset\) only enlarges the nonnegative sum.
All selected domains and failed bonds still have collars
inside \(Z\), so \eqref{eq:local-formation-UM10}--\eqref{eq:local-formation-UM11} give
\begin{equation}
 \begin{split}
 \sum_{\substack{\mathbf D,P_0\\
                     (\mathbf D,P_0)\ne(\emptyset,\emptyset)}}
   \prod_{Y\in\mathbf D}u_Y
   \prod_{b\in P_0}q
 &\le
  \prod_{Y\in{\cal Y}(Z)}(1+u_Y)
       \prod_{b:b^\sharp\subset Z}(1+q)-1\\
 &\le e^{\tau N}-1 ,
 \end{split}\label{eq:local-formation-UM21}
\end{equation}
where
\begin{equation}
 \begin{split}
 u_Y&=(2a/W_0)e^\Theta e^{-\sigma_0d_f(Y)},\qquad
 q=pe^\Theta,\\
 p_\star&=n_bp,\qquad
 \tau=e^\Theta\left[
        \frac{2c_fL^4S_{\sigma_0}}{W_0}a+p_\star\right].
 \end{split}\label{eq:local-formation-UM22}
\end{equation}
The products are finite in finite volume.  In particular no
factor from bonds outside \(Z\) has been inserted.
The subtraction of one is essential: it retains a selected
potential or failed cutoff, rather than creating a unit
amplitude for every \(Z\).

Take the physical smallness conditions
\begin{equation}
 a\le W_0/2,\qquad \tau\le1 .
 \label{eq:local-formation-UM23}
\end{equation}
For \(N\ge1\),
\(e^{\tau N}-1\le\tau N e^{\tau N}\) and \(N\le e^N\).
\eqref{eq:local-formation-UM20}--\eqref{eq:local-formation-UM23} prove
\begin{equation}
 \begin{aligned}
 |H(Z)|&\le e^{-K_*N}(e^{\tau N}-1)
   \le\tau e^{-(K_*-2)N}\\
   &\le C_\tau(a+p_\star)e^{-\Lambda\ell(Z)} ,\\
 C_\tau&=e^\Theta
       \max\!\left\{\frac{2c_fL^4S_{\sigma_0}}{W_0},1\right\}.
 \end{aligned}
 \label{eq:local-formation-UM24}
\end{equation}
This is an estimate for the actual activities \eqref{eq:local-formation-UM3}.
Their smallness has not been assumed.
The larger Gaussian budget \(W_0\) was used only to bound
the analytic integrals; the physical budget in \eqref{eq:local-formation-UM23} is
smaller and is what makes the activity amplitude at most one.

\paragraph*{The actual union filing and the local logarithm}

Choose a desired logarithmic output rate \(\mu>0\), and set
\(g_*=\Lambda-\mu\).
Let \(N\le C_v(d_c(Z)+1)\), and let the number of coarse
polymers containing a fixed cube with
\(m\le d_c(Z)<m+1\) be at most \(C_De^{\kappa_Dm}\).
These geometric bounds follow from bounded-degree cube
covering walks as above.
There is a fixed \(c_{\rm adj}\) such that polymers
incompatible with \(Z\) contain at least one cube in a
set of at most \(c_{\rm adj}N\) cubes: use \(Z\) and
its wall neighbors.

For these counting estimates use the positive norm majorants
\[
 h_Z=\sup_{P_Z}|H(Z;P_Z)|,\qquad
 \widehat w_Z=h_Ze^{\mu\ell(Z)},
\]
where the supremum is over the declared local core of \(Z\),
not over restrictions of one global field.
At a fixed core pair \(P_X\) on \(X\), actual activities are
used only for polymers \(Z\subset X\), with their values
obtained by the assigned hereditary restrictions.
Their weighted absolute values are bounded by
\(\widehat w_Z\).
The bound with amplitude \(\tau\le1\) in \eqref{eq:local-formation-UM24} gives
\begin{equation}
 \sup_q\sum_{Z\ni q}\widehat w_Ze^{\#Z}
 \le
 \tau\,C_De^{2C_v}e^{-5g_*}
       \sum_{m\ge0}
               e^{-(g_*-\kappa_D-C_v)m}
 =:\tau B_* .
 \label{eq:local-formation-UM25}
\end{equation}
Choose the rates so that
\begin{equation}
 g_*>\kappa_D+C_v,\qquad c_{\rm adj}B_*\le1 .
 \label{eq:local-formation-UM26}
\end{equation}
The explicit factor \(e^{-5g_*}\) allows the last
inequality by increasing the rate gap.
Only the upper bound \(\tau\le1\) is needed here.
It follows that
\begin{equation}
 \sum_{Z':Z'\not\sim Z}\widehat w_{Z'}e^{\#Z'}
       \le c_{\rm adj}N\tau B_*\le N .
 \label{eq:local-formation-UM27}
\end{equation}
This verifies, for the actual activities, the criterion
of the already proved hard-core tree theorem, with
the root weight \(\#Z\).
In particular it verifies that criterion on every hereditary
subsystem \(Z\subset X\), without assigning actual fields
outside \(X\).

If two polymers are incompatible, their covering trees
can be joined at a common cube or wall with length at
most five in four dimensions.
For a connected incompatibility list this gives
\begin{equation}
 \ell\!\left(\bigcup_{i=1}^nZ_i\right)
                         \le\sum_{i=1}^n\ell(Z_i).
 \label{eq:local-formation-UM28}
\end{equation}
Indeed a spanning tree of the incompatibility graph
has \(n-1\) joins, and
\(d_c(\cup_iZ_i)\le\sum_i d_c(Z_i)+5(n-1)\).
Thus the union filing has the same exponential weight.

Define the actual local logarithmic term by
\begin{equation}
 E(X)=\sum_{n\ge1}\frac1{n!}
       \sum_{\cup_iZ_i=X}
          \phi(Z_1,\ldots,Z_n)\prod_iH(Z_i),
 \label{eq:local-formation-UM29}
\end{equation}
where \(\phi\) is the connected hard-core graph coefficient.
For a fixed cube \(q\in X\), mark a vertex whose polymer
contains \(q\).  This only enlarges the absolute series.
Symmetry of the ordered list changes its factor
\(1/n!\) into \(1/(n-1)!\).
When dropping the union restriction, retain \(Z_i\subset X\)
for the actual terms at \(P_X\).
Use \eqref{eq:local-formation-UM28}, and then the rooted tree theorem for that
subsystem, dominated by the numerical majorants
\(\widehat w\).
Only the subsequent nonnegative majorant sum may be enlarged
to all geometric polymers containing \(q\).
The result is
\begin{equation}
 \begin{split}
 e^{\mu\ell(X)}|E(X)|
 &\le\sum_{\substack{Z\subset X\\Z\ni q}}
                       \widehat w_Ze^{\#Z}
 \le\tau B_*\\
 &\le C(a+p_\star).
 \end{split}\label{eq:local-formation-UM30}
\end{equation}
There is no factor \(\#X\) or ambient volume in this
local estimate.  It proves the unmarked potential-ball
majorant used in the finite insertion theorem.

Uniform convergence on smaller external and coupling
domains follows from these same positive majorants.
The actual local functions are holomorphic, so their
union-filed sum is holomorphic on each declared core.
At complete physical pairs, the complete-Gaussian and
wall-factorization identities identify the gas with the
finite regular integral; the hard-core theorem then gives
\(\exp(\sum_XE(X))=\Xi\ne0\).
On general local cores the activity and logarithm bounds
do not require that global Gaussian interpretation.

\paragraph*{Parameter choices and the available insertion slot}

All constants above concern fixed geometry and the actual
matrix/core package.  One must verify \eqref{eq:local-formation-UM4}--\eqref{eq:local-formation-UM6} and
\eqref{eq:local-formation-UM14}, \eqref{eq:local-formation-UM19}, \eqref{eq:local-formation-UM23}, \eqref{eq:local-formation-UM26} together.
For example, it suffices that the matrix package control
\(A_Kc_m(LM)^4\varepsilon\) by a fixed constant on the
chosen decoupling discs.
Choose \(W_0\) so its two contributions to \(C_0\) are
bounded, and choose positive \(\gamma\) of order at most
\((LM)^{-4}\) to bound its contribution.
Then \(C_0\) is bounded independently of exterior volume.
One can choose the target rate gap, the coarse parameter
radius and \(\lambda_V\) by \eqref{eq:local-formation-UM18}--\eqref{eq:local-formation-UM19}, and finally
make \(a\) small and \(\mathfrak b\) large to obtain \eqref{eq:local-formation-UM23}.
This lists explicit sufficient inequalities; it does not
assert them for a matrix family whose constants have not
been controlled on the required enlarged discs.

For an explicit insertion budget, put
\[
 A_a=\frac{2c_fL^4S_{\sigma_0}}{W_0},\qquad
 a_*=\min\!\left\{\frac{W_0}{2},
                           \frac1{4e^\Theta A_a}\right\},
 \qquad p_\star\le\frac1{4e^\Theta}.
\]
Then every total norm \(a\le a_*\) has \(\tau\le1/2\)
and satisfies \eqref{eq:local-formation-UM23}, uniformly above this cutoff threshold.
If the base total potential norm is at most \(a_*/2\), reserve
one additional potential slot of norm \(w_\star\le a_*/2\).
The entire slot stays within \eqref{eq:local-formation-UM23} and hence within
\eqref{eq:local-formation-UM24}--\eqref{eq:local-formation-UM30}.
This proves the local uninserted majorant uniformly on
that slot without assuming a zero-free activity domain.  The constants may depend on the fixed
localization parameters and rate choices, but not on
cutoff, torus size, old-level count or the auxiliary
Gaussian dimension.

\subsubsection{A small potential direction and its first jet}


We keep the calibration, the real action weight, the Gaussian structure,
the cutoffs and the background maps fixed.  The result below concerns the
finite regular fluctuation integral.  It does not enlarge the source domain
of a large-field history operation.  This distinction permits a small
potential direction to be used at its actual norm.
Localized terms and their insertion estimates are defined on the
specified complex local cores, with potential and current independent
there.  Their sum is identified with a complete Gaussian integral
only at physical complete background pairs \(P=p_W\).
An identity on that graph is not extended to arbitrary independent
currents.

\paragraph*{The fixed-reference integral}

Let \(\theta\ge\theta_{\min}\), \(g=\theta^{-1/2}\), and
\(\tau=x/\theta\), with \(|\tau-1|<\sigma_\star\).
The fixed small number \(\sigma_\star\), the allowed external pairs
and their Gaussian gap are restricted by the ordinary Gaussian
comparison and potential allowance below.  In particular the
decoupled precisions stay in their coercive regime and the free
potential leaves the declared allowance; no assertion is made
for arbitrary positive \(x/\theta\).  The sourced run considered
here keeps \(x>0\) fixed.
At a physical complete background pair \(P=p_W\), in the
fluctuation coordinates \(B\) on the non-pivot bonds, the reference
measure is the normalized Gaussian with covariance \(\tau^{-1}C(P)\).
The same identity holds on a holomorphic continuation in physical
background data wherever that continuation and the ordinary Gaussian
bounds have been supplied.  It is not an identity on every complex
core pair.  Put
\begin{equation}
 \Xi(P,z)=\int d\mu_{\tau^{-1}C(P)}(B)\,
       \prod_b\chi_b(B)\,
       \exp\left\{\sum_Y V(Y;P,B)+z\sum_Yv(Y;P,B)\right\}.
 \label{eq:local-formation-SI1}
\end{equation}
The parameter \(z\) enters only the potentials.  In particular the covariance,
the cutoff thresholds and the Gaussian determinant do not depend on \(z\).
Here the cutoffs are sharp indicators
\(\chi_b(B)=\mathbf1\{|B(b)|<\mathfrak b\}\) on free bonds.
For a bond set \(A\), write
\(\chi_A=\prod_{b\in A}\chi_b\) and
\(\chi_A^c=\prod_{b\in A}(1-\chi_b)\).
Thus the latter denotes a product of complements.
After the old action at the background has been carried separately, the
new fluctuation increment is
\begin{equation}
 {\cal E}(P,z)={\cal N}[\theta;x](P)+\log\Xi(P,z).
 \label{eq:local-formation-SI2}
\end{equation}
The field-independent Gaussian number includes
\(-\frac32|\mathfrak X|\log\tau\), as well as the Haar-coordinate,
scaling and gauge-fixing counts.  The localised normalizer
\(\mathbf N^\circ\) and the free-potential family \(\mathbf G_\tau\)
also use the fixed \((\theta,x)\).  All these quantities have zero
\(z\)-derivative.  Moving \(z\) into \(x\), and differentiating the resulting
Gaussian determinant, would describe a different presentation.

For an inserted exponent \(H\), the potential actually inserted is its
fluctuation difference.  Thus, when the fibre field is \(U_P(B)\) and
\(U_P(0)\) is its background, at a physical complete pair,
\begin{equation}
 \sum_Yv(Y;P,B)=H(U_P(B))-H(U_P(0)).
 \label{eq:local-formation-SI3}
\end{equation}
Consequently the background subtraction in the step functional is already
present in \eqref{eq:local-formation-SI1}; it must not be subtracted a second time.
Equivalently, the raw fibre integral with the additional exponent
\(zH(U_P(B))\) has the extra factor \(e^{zH(U_P(0))}\) outside \eqref{eq:local-formation-SI1}.
That factor belongs to the separately carried old action, rather than to
the newly generated cluster activities.

\paragraph*{The ordinary Gaussian expansion used here}

We state the analytic inputs at the level at which they are applied.
On an arbitrary complex local core pair \(P\), the actual localized
Gaussian functionals \({\cal J}(Z_0;s;P)\), potentials and
union-filed terms below are defined by their local rules.
We do not assume that the complete values for different \(Z_0\)
all arise as marginals of one global Gaussian at that \(P\).
A potential has norm
\begin{equation}
 |V(Y;P,B)|\le\|\mathbf V\|_\lambda e^{-\lambda d(Y)}
       \left(1+M^{-4}\sum_{b\subset Y}|B(b)|^2\right)
 \label{eq:local-formation-SI4}
\end{equation}
on its specified complex field box.  The potential may be defined only
there: the cutoff product on its domain is kept in every term in which
it occurs.  The ordinary Gaussian comparison estimate bounds a localized
decoupled functional by \(F_0e^{CuN(Z)}\) when its integrand is bounded by
\(F_0e^{\rho|B|^2}\), with \(u=\widehat\varepsilon+\rho\).
Here \(N(Z)\) counts the affected real scalar Gaussian coordinates,
of order \((LM)^4N_c(Z)\), where \(N_c(Z)\) is the number of
coarse cubes.  In the cube-count convention the same bound is
\(F_0e^{C(LM)^4uN_c(Z)}\); the dimension factor is counted once.
The Gaussian reference gap, radii and tilt are chosen so these factors are
bounded per coarse cube.  The decoupling parameters have a fixed enlarged
complex disc.  These are the ordinary background/Gaussian inputs; no
source derivative is contained in them.

Choose a potential allowance \(w_\star>0\), independent of the cutoff and
volume, small enough for the following ordinary expansion.  The base
potentials use at most a fixed number of slots, and leave one slot available
for an additional potential of norm at most \(w_\star\).  The same
Gaussian comparison, rooted domain count and decoupling majorant are valid
throughout this allowance.  The convergent hard-core expansion gives
local logarithmic terms \(E(X)\) and bounds
\begin{equation}
 |E(X)|\le C_\star e^{-\lambda_{\rm out}d(X)}.
 \label{eq:local-formation-SI5}
\end{equation}
The constant is uniform in all allowed external pairs, decoupling
parameters and potentials.  This is the ordinary uninserted polymer
majorant, not an assumed differentiability bound for a renormalization map.

Here is the explicit finite expansion to which that majorant applies.
Let \(\mathbf D\) be a set of selected interaction domains.  Write
\(Y_0\) for the union of their free-bond sets on which a field box is
required; a globally defined free potential does not contribute to this
union.  Let \(P_0\subset\mathfrak X\setminus Y_0\) be a set of failed
fluctuation cutoffs.  The indicator identities \(\chi_b^2=\chi_b\)
and \(\chi_b(1-\chi_b)=0\) give exactly
\(\prod_b\chi_b=\chi_{Y_0}
\sum_{P_0\subset\mathfrak X\setminus Y_0}
(-1)^{|P_0|}\chi_{P_0}^c\).
Before
decoupling, its integrand is
\begin{equation}
 F_{\mathbf D,P_0}(z)
    =(-1)^{|P_0|}\chi_{Y_0}\chi^c_{P_0}
             \prod_{Y\in\mathbf D}(e^{V(Y)+zv(Y)}-1).
 \label{eq:local-formation-SI6}
\end{equation}
The empty choice \((\mathbf D,P_0)=(\emptyset,\emptyset)\)
supplies the leading one in the gas and is omitted from its
activities.  All polymer domains \(Z\) are nonempty and
wall-connected, so no \(H(\emptyset)\) is defined.
Disconnected regions are products of their wall-component
activities, rather than additional single polymers.
Equivalently, for a nonempty \(Z\) the empty choice has at least
one decoupling derivative of the normalized constant functional
\({\cal J}[1]=1\) and contributes zero.
The index \(Y\) carries one of the fixed potential-type tags used in
the unmarked expansion.  Its cutoff requirement is fixed by its
support and type: a boxed type uses the common indicators on all
free bonds in that support, as in the definition of \(Y_0\), while
the globally defined free type requires none.  Older levels and
source origins do not add cutoff types.  Before forming the product,
sum all origins having the same support and type into that type's
potential.  This grouping is performed
before taking absolute majorants, and the norm in
\eqref{eq:local-formation-SI4} is the norm of the grouped potential.
Indeed, for finitely many origins on that same support,
\[
 e^{\sum_o V_o}-1
   =\sum_{\emptyset\ne I\subset\{o\}}
       \prod_{o\in I}(e^{V_o}-1).
\]
Linearity of the local Gaussian functional and of the decoupling
operations therefore preserves the complete integral and its union
filing.  The estimate is applied to the grouped left-hand side;
it does not require an absolute majorant for the separate terms on
the right.  An individual origin can still be selected as the marked
direction \(v\), without adding a Mayer type.  The geometric support
and its tree length ignore the fixed type tag.
For each support index \(\iota=(\mathbf D,P_0)\) let \(Z_0(\iota)\)
be its support.  The decoupled activities are formed by
\begin{equation}
 \begin{aligned}
 T(Z,\iota;z)&=
 \prod_{\Delta\subset Z\setminus Z_0(\iota)^\sharp}
          \int_0^1ds_\Delta\,\partial_{s_\Delta}\,
       {\cal J}(Z_0(\iota);s_Z)[F_{\mathbf D,P_0}(z)],\\
 H(Z,z)&=\sum_\iota T(Z,\iota;z).
 \end{aligned}
 \label{eq:local-formation-SI7}
\end{equation}
In this formula \(s_Z\) is one on \(Z_0(\iota)^\sharp\), is the
displayed free parameter on \(Z\setminus Z_0(\iota)^\sharp\), and
is zero outside \(Z\).  The normalized Gaussian functional
\({\cal J}\) is independent of \(z\).  At a physical complete
pair \(P=p_W\), its complete-structure identity is
\[
 {\cal J}(Z_0;1)[F]
       =\int d\mu_{\tau^{-1}C(P)}(B)\,F(B).
\]
This physical identity, together with interpolation of every remaining
coarse-cube parameter and the ordinary decoupling property, is what
identifies the gas with \eqref{eq:local-formation-SI1}.
Two polymers are compatible when they have neither a cube nor a wall
in common.  Multiplicativity over their wall-components yields
\(\Xi=\sum_{\Gamma\ {\rm compatible}}\prod_{Z\in\Gamma}H(Z)\),
with empty product one, at a physical complete pair.
On general local cores, the same local formulas define activities and
their union-filed terms without asserting one global Gaussian
representation.  Those logarithmic terms are
\begin{equation}
 E(X,z)=\sum_{n\ge1}\frac1{n!}
       \sum_{\cup_{i=1}^n Z_i=X}
            \rho^T(Z_1,\ldots,Z_n)\prod_{i=1}^n H(Z_i,z).
 \label{eq:local-formation-SI8}
\end{equation}
Here \(\rho^T(Z_1,\ldots,Z_n)\) is the sum over connected graphs on
\(\{1,\ldots,n\}\) of the products of their edge weights; an edge has
weight \(-1\) for an incompatible pair and zero otherwise.
For \(n=1\), \(\rho^T=1\).
Every sum is first formed in finite volume; the logarithmic series and
its differentiated series are justified by the same absolute majorants.

\begin{proposition}[The inserted-potential estimate]
\label{prop:finite-inserted-potential}
Let \(\varepsilon=\|\mathbf v\|_\lambda\).  Under the ordinary
local Gaussian/polymer hypotheses above, on every specified
complex local core pair,
\begin{equation}
 \left|\left.\partial_zE(X,z)\right|_{z=0}\right|
       \le C\varepsilon e^{-\lambda_{\rm out}d(X)}.
 \label{eq:local-formation-SI9}
\end{equation}
At a complete physical pair \(P=p_W\), its sum over \(X\)
is the normalized expectation
\begin{equation}
 \left.\partial_z\log\Xi(P,z)\right|_{z=0}
   =\frac{ \int d\mu\,\prod_b\chi_b\,e^{\sum_YV(Y)}
                            \sum_Yv(Y)}{
        \int d\mu\,\prod_b\chi_b\,e^{\sum_YV(Y)}}.
 \label{eq:local-formation-SI10}
\end{equation}
There is no volume factor in the local estimate in \eqref{eq:local-formation-SI9}.
\end{proposition}

\begin{proof}
If \(\varepsilon=0\), every inserted potential vanishes on its domain
and the assertion is immediate.  Otherwise take the actual auxiliary
potential \(z\mathbf v\) on
\(|z|<w_\star/\varepsilon\).  This is a legitimate potential in the
actual local Mayer and decoupling expressions,
regardless of whether it would be an admissible physical source for a
history construction.  The ordinary Gaussian comparison and the polymer
criterion hold on this disc.  Every integrand and every activity is
holomorphic there; Gaussian domination and the parameter-independent
majorant give holomorphy of the union-filed series \eqref{eq:local-formation-SI8}.
This local argument does not require one global Gaussian at a
general independent-current core pair.
Apply Cauchy's estimate on
\(|z|=w_\star/(2\varepsilon)\) to \eqref{eq:local-formation-SI5}.  It gives \eqref{eq:local-formation-SI9} with
constant \(2C_\star/w_\star\), which depends only on fixed ordinary
allowances.

The same argument gives a norm derivative, which is useful when the
actual source family of potentials is nonlinear.  On the potential
space with its fixed declared domains, joint holomorphy on every
two-dimensional affine plane makes the directional derivative
complex linear: the degree-one term of
\(E_X(\mathbf V+z_1\mathbf h_1+z_2\mathbf h_2)\) is the sum of
its two degree-one terms.  Cauchy's estimate makes this linear map
bounded.  For \(\|\mathbf h\|_\lambda<w_\star/2\), apply
one-variable Cauchy to \(t\mapsto E_X(\mathbf V+t\mathbf h)\)
on the disc \(|t|<w_\star/\|\mathbf h\|_\lambda\).
Summing its Taylor coefficients of degree at least two gives
\[
 \left|E_X(\mathbf V+\mathbf h)-E_X(\mathbf V)
                         -DE_X(\mathbf V)[\mathbf h]\right|
 \le C C_\star\frac{\|\mathbf h\|_\lambda^2}{w_\star^2}
                         e^{-\lambda_{\rm out}d(X)}.
\]
All constants are uniform in the external core pair, so this is
also the remainder bound in the unpointed supremum norm
\(\|F\|_{\rm out}=\sup_{X,P}
e^{\lambda_{\rm out}d(X)}|F_X(P)|\).
It does not assert a pointed or stronger rooted norm bound.
Uniform local Lipschitz bounds on smaller potential balls give the
unique continuous extension to the potential-norm completion.
The displayed remainder passes to that extension and proves
Fr\'echet differentiability there as well.
If an actual ordinary source family satisfies
\(\mathbf V(w)=\mathbf V(0)+w\dot{\mathbf V}+o(|w|)\)
in that norm, the remainder bound proves its chain rule
\(\partial_wE_X(\mathbf V(w))|_0=
DE_X(\mathbf V(0))[\dot{\mathbf V}]\).
Pointwise holomorphy of \(\mathbf V(w)\) on its ordinary source
disc, with a uniform potential-norm bound there, gives its
norm Taylor expansion by the same uniform Cauchy estimates.
Thus the auxiliary linear direction computes the actual first
source jet; the nonlinear physical source family need not exist
on the large auxiliary-potential disc.

This is also the marked-activity derivative of the explicit expansion.
The derivative of \eqref{eq:local-formation-SI6} replaces precisely one selected factor by
\begin{equation}
 v(Y)e^{V(Y)}
       \prod_{\substack{Y'\in\mathbf D\\Y'\ne Y}}
                                     (e^{V(Y')}-1).
 \label{eq:local-formation-SI11}
\end{equation}
Domains with \(V(Y)=0\) are included: their unmarked term is zero,
but their derivative need not be zero.  In a term with a marked potential
the cutoff on its domain is retained.  Derivatives commute with \eqref{eq:local-formation-SI7}
by the same Gaussian comparison and with \eqref{eq:local-formation-SI8} by absolute convergence.
Thus this proof differentiates the actual finite interaction formula,
not a hypothetical normed map without a specified integral.
The pinned cluster estimate is what keeps the local Cauchy bound in \eqref{eq:local-formation-SI5}
independent of volume.

At a complete physical pair, the sum rule for the ordinary
cluster expansion gives
\(\exp(\sum_XE(X,z))=\Xi(P,z)\ne0\).
Differentiating it and the finite integral proves \eqref{eq:local-formation-SI10}.
The denominator is this one regular Gaussian integral; no division by
the full sourced partition function or by a history mass is used.
\end{proof}

\paragraph*{The actual scalar action direction}

Let \(H_u(U)=A(u,U)\), where \(u\) is a bounded function on bare
plaquettes.  It is filed by assigning a plaquette to the union of the
fixed cubes containing its corners.  In the background field domain,
\(|\partial U(p)-1|\le C\alpha\eta^2\).
The identity
\(a_p(U)=\frac12\det(1-\partial U(p))\) gives a raw local action
member of norm at most \(CM^4\alpha^2\sup|u|\).
It has no coupling counterterm at this stage.

More precisely, at a complete physical pair the potential in \eqref{eq:local-formation-SI3}
is a first difference of this action.
Suppose the actual background response and its decoupled
substitutions stay in their stated field domains and shift a plaquette
by at most \(C\delta\eta^2\), where \(0<\delta\ll\alpha\).
Along this shift, the plaquette remains at distance at most
\(C\alpha\eta^2\) from the identity.  The derivative of the determinant
of a \(2\times2\) matrix is bounded by a constant times its size, so
\begin{equation}
 |a_p(U_P(B))-a_p(U_P(0))|
                          \le C\alpha\delta\eta^4.
 \label{eq:local-formation-SI12}
\end{equation}
This proof works at complex pairs.  It uses the actual curvature and
response bounds, not positivity.  On a local core the same determinant
bound applies to the two actual substituted plaquette matrices whenever
the declared response bounds give these curvature and difference
estimates; it does not require a global fibre map there.

We give the quantitative fine-response localization separately from
the coarse Gaussian decoupling in \eqref{eq:local-formation-SI7}.
For a plaquette \(p\), let \(X_0(p)\) be the union of its fixed
corner cubes.  Its number of cubes is at most a dimensional constant
\(C_{\rm base}\), and its tree length is zero.
Write \(\mathcal U(P')\) for the bare configuration used by the
actual background evaluation of a pair \(P'\), and set
\(F_p(P')=u(p)a_p(\mathcal U(P'))\).
Thus this local function is explicitly a function of a pair.
The ordinary response-bracket hypotheses used here are the following.
On each \(Y\supset X_0(p)\) the same response rule gives
\[
 \Phi_{p,Y}(s,\nu)=
 F_p\big(P|_{X_0(p)}\dotplus\nu\,\mathfrak h_Y(s;B)\big),
 \qquad
 |s_\Delta|<R=e^{\bar\kappa},\quad |\nu|<\rho,
\]
holomorphic jointly on these independent discs, uniformly in the
allowed complex external fields and fluctuation box, with
\[
 \rho\ge c_R\alpha/\delta>1,\qquad
 \sup|\Phi_{p,Y}|\le C\alpha^2\eta^4|u(p)|.
\]
Here \(c_R>0\) can contain the fixed loss
\(e^{-C\bar\kappa}\).
In particular all substituted pairs must stay in their actual cores
on the whole joint polydisc, not just for real interpolation parameters.
The operation \(\dotplus\) shifts the potential and current by the
actual response.  These are quantitative background and membership
hypotheses; the following Cauchy argument does not prove them.

For disjoint fine-cube sets \(K,N\), define the mixed bracket by
\[
 \Psi[K|N]=\int_{[0,1]^K}
       \prod_{\Delta\in K}\partial_{s_\Delta}
                        \Psi(s_K+1_N)\,ds_K ,
\]
with all other parameters zero.  Its scalar-action potential is
the bracket of
\(\Psi(s)=\Phi_{p,Y}(s,1)-\Phi_{p,Y}(s,0)\)
with \(K=Y\setminus X_0(p)\), \(N=X_0(p)\).
At \(\nu=0\) the field is the unshifted restricted background.
Schwarz's lemma in \(\nu\) first gives
\(\sup_s|\Psi(s)|\le
2C\rho^{-1}\alpha^2\eta^4|u(p)|\).
For every \(s_\Delta\in[0,1]\), a coordinate Cauchy circle has
radius arbitrarily close to \(R-1\).
Iterated Cauchy therefore bounds the bracket by
\[
 2C\rho^{-1}\alpha^2\eta^4|u(p)|
             (R-1)^{-|K|}
 \le C_R\alpha\delta\eta^4|u(p)|
              e^{-(\bar\kappa-1)|K|},
\]
where \(\bar\kappa\ge1\).
There is no factorial in this estimate, because each independent
cube parameter is differentiated once.
The ordinary rooting property says that a component containing a
differenced cube and missing \(X_0(p)\) makes the bracket zero.
For the present choice \(N=X_0(p)\), every component missing
the base has such a differenced cube.  No assertion about an
arbitrary remote component consisting only of undifferenced
parameters is being made.
The local response rule agrees with the response built on the
active wall-component.  Between local ambient domains on which
the substitutions are defined it has hereditary consistency,
\(\Phi_{p,Y}(s,\nu)=
\Phi_{p,\Lambda}(s\,\mathbf1_Y,\nu)\), with the base parameters
set to one.
At a complete physical pair on the full lattice its endpoints are
\[
 \Phi_{p,\Lambda}(1,1)=u(p)a_p(U_P(B)),\qquad
 \Phi_{p,\Lambda}(1,0)=u(p)a_p(U_P(0)).
\]
Apply the finite fundamental theorem of calculus in every fine-cube
parameter outside \(X_0(p)\), expanding evaluation at one as evaluation
at zero plus the integral of the derivative.
The resulting sum of brackets equals
\(\Phi_{p,\Lambda}(1,1)-\Phi_{p,\Lambda}(1,0)\).
At a complete physical pair this is exactly
\(u(p)[a_p(U_P(B))-a_p(U_P(0))]\).
Wall-component consistency identifies each bracket with the
corresponding \(\Phi_{p,Y}\); rooting removes the disconnected
terms.  Thus the construction gives a uniformly bounded local
extension on every declared core and localizes the actual
difference in \eqref{eq:local-formation-SI3} at complete physical pairs.

For a surviving bracket, the ordinary domain geometry gives
\(d(Y)\le C_{\rm geo}(|K|+1)\).
At fixed \(Y\) there are at most
\(CM^4\eta^{-4}(|K|+C_{\rm base})\) possible plaquette anchors.
Choose the fixed surplus \(\bar\kappa\) so that
\(\bar\kappa-1>C_{\rm geo}\lambda+2\), increasing it further
if a subsequent rooted-domain entropy sum is included.
Sum the displayed bracket bound over its plaquette anchors and
use that a linear polynomial times \(e^{-2|K|}\) is bounded.
This proves
\begin{equation}
 \|\mathbf v_u\|_\lambda
       \le C M^4\alpha\delta\,\sup|u|
       =C M^4\alpha^2 s\,\sup|u|,\qquad s=\delta/\alpha.
 \label{eq:local-formation-SI13}
\end{equation}
The factors \(\eta^4\) and \(\eta^{-4}\) cancel exactly.
The required membership and response estimates on all decoupling
substitutions are background inputs.  The bracket argument neither
creates those inclusions nor replaces them by bounds at real fields.

Combining \eqref{eq:local-formation-SI9} with \eqref{eq:local-formation-SI13} proves a one-step unpointed source response
of size \(C\alpha^2s\sup|u|\) for the actual local regular family
on every declared complex core.  Its physical sum is the source
response of the fixed-reference regular integral.
The Gaussian normalizer contributes zero derivative.
Its interaction vacuum value contributes
\(\partial_zE(X,(1,0),0)\), which is the source vacuum energy and obeys the
same bound.  Subtracting it leaves a factor at most two in the estimate.

\paragraph*{Own counterterms and the remaining iteration}

For a constant source, once the actual bond-density pointing is applied,
the derivative member is Euclidean.  Its own coefficient is
\(\partial_z b[E_z]|_0=b[\partial_zE_z|_0]\), because the free response
and the coefficient extractor are fixed and complex linear.
The counterterm to remove from the derivative action at site \(y\) is
therefore this number times \(A(h_y,U)\).  A coefficient belonging to
another member, or a derivative of a Gaussian presentation with a
different weight, would not supply that identity.

There are two further tasks before the one-step estimate is a complete
multilevel source-response bound.  First, renewing the actual pointing
uses three bond-indexed interpolation legs.  Differentiating a normalized
leg expectation produces both the derivative of its leg insertion and a
connected expectation of that insertion with \(\sum_Yv(Y)\).
A root-plus-insertion cluster bound and the path's admissible field room
must control those terms before the same estimate is asserted for each
pointed term.

Second, the old regular derivative members and the nonconstant source
defects have to be inserted again with their own coefficients.  A
constant-source derivative of size \(\alpha_j^2s_j\) can be normalized
by the pointed Hessian theorem.  Write \(\alpha_H>0\) for the fixed
spatial H\"older exponent, as distinct from an analytic field radius.
On a non-wrapping local chart, retaining
the \emph{true} radius ratio
\(\rho_j=\delta_{\rm out}/\alpha_j\) in the near-domain Taylor proof,
before any enlargement to \(\max\{1,\rho_j\}\), would give after the
fine-site count the contribution
\begin{equation}
 C\alpha_j^2s_j
   \left[\left(\frac{\delta_{\rm out}}{\alpha_j}\right)^2
                       L^{-\alpha_H(n-j)}
       +\left(\frac{\delta_{\rm out}}{\alpha_j}\right)^3
                       L^{-2(n-j)}\right].
 \label{eq:local-formation-SI14}
\end{equation}
Its quadratic part sums to \(C\delta_{\rm out}^2s_{\rm out}\) when
the ordinary flow and age weights apply.  For
\(\alpha_{\rm out}\asymp x_{\rm out}^{-1/2}(\log x_{\rm out})^q\)
and
\(\delta_{\rm out}\asymp x_{\rm out}^{-1/2}
(\log x_{\rm out})^{p_0}\), this is
\begin{equation}
 Cx_{\rm out}^{-1}(\log x_{\rm out})^{3p_0-q}.
 \label{eq:local-formation-SI15}
\end{equation}
This arithmetic is conditional on that homogeneous local refinement
and does not establish the derivative recurrence.  In
particular the defects must retain the factor \(\alpha_j^2s_j\) as
well as the scaled Lipschitz gain.  A fixed-disc source bound alone
does not provide it.  Marked and remote terms need their own uniform
fixed-disc bounds.  Field-independent tails around holes or torus cycles
cannot be assigned an extra \(\rho_j^2\) when \(\rho_j<1\); they require
a non-wrapping chart or a separately small class estimate.


\paragraph*{Identification with the fixed-reference source family}

In the variable-profile presentation the actual regular family has
the decomposition
\begin{equation}
       {\cal E}^{\,w}=N^\circ+G_\tau+F^{\,w},
 \qquad
       {\cal E}^{\,w}-{\cal E}^{\,0}=F^{\,w}-F^{\,0}.
 \label{fm:fixed-reference-source-family}
\end{equation}
Here \(x,\theta\), the Gaussian matrices, cutoffs and numerical
normalizers are fixed; the variable profile is an older member.
Thus \(N^\circ\) and \(G_\tau\) have no source derivative.
If its actual potential family is differentiable in the accepted
potential norm, the uniform norm-Taylor proof of \eqref{eq:local-formation-SI9} applies
to each unpointed local logarithm:
\begin{equation}
 \begin{aligned}
 \dot{\cal E}&=\dot F,\\
 \left.\partial_w E_X(V(w))\right|_0
       &=DE_X(V(0))[\,V'(0)\,].
 \end{aligned}
 \label{fm:actual-first-jet-chain}
\end{equation}
The derivative of the pointed density is computed in the next
section; it includes both terms of \eqref{sp:differentiated-leg}.  These are identities
for the fixed presentation, rather than derivative rules for
one that also changes its reference data.

For two auxiliary directions of norms \(A\) and
\(\varepsilon\), one spare potential family suffices:
\begin{equation}
        \|\lambda a+zv\|_{\lambda_V}
             \le|\lambda|A+|z|\varepsilon<w_\star.
 \label{fm:one-spare-family-load}
\end{equation}
The two directions are combined in this one slot; they
do not increase the allowed number of potential types.
The resulting auxiliary Cauchy polydisc is a domain
of the local potential expansion.  It need not be a
domain of physical Wilson sources or of history
operations.


\subsubsection{Renewing the actual bond pointing}


An estimate for an unpointed localized logarithm is not by itself an
estimate for a particular pointing of that logarithm. We retain the
actual three interpolation legs and differentiate their bond densities.
For a source filed as an older weight member, the Gaussian structure,
free potential, cutoffs and numerical normalizers are fixed. Thus the
first, free-potential leg and the second, cutoff leg are independent of
this source. Only the interaction leg contributes below. This choice
of presentation is essential to the statement.

\paragraph*{The finite local logarithm and its actual locality}

Fix a finite regular fluctuation integral, its Gaussian reference, its
cutoffs and all geometric and external data. Fix nested full complex
fluctuation boxes
\[
 \mathcal B_0=\{B:\sup_b|B(b)|<\mathfrak b_0\},\qquad
 \mathcal B_1=\{B:\sup_b|B(b)|<(1+\nu)\mathfrak b_0\},
 \qquad \mathfrak b_0\ge1,\quad \nu>0.
\]
The ordinary spare-potential estimate below accepts potentials
holomorphic and bounded in its growth norm on $\mathcal B_0$;
all its actual Gaussian cutoff supports lie there. The larger box
$\mathcal B_1$ is reserved for the one coordinate Cauchy estimate
on an inserted potential. Both boxes and the external-pair and
decoupling domains are fixed once for this step. Ordinary inputs
whose cutoffs fill half their analytic box leave this room by
choosing an intermediate $\mathcal B_0$. This uses that existing
room once, and is not a new shrinking at each source derivative.
A local potential family
has the Gaussian growth norm
\begin{equation}
 |V(Y;P,B)|\le \|\mathbf V\|_\lambda e^{-\lambda d(Y)}Q_Y(B),
 \qquad Q_Y(B)=1+M^{-4}\sum_{b\subset Y}|B(b)|^2.
 \label{sp:potential-norm}
\end{equation}
There are finitely many bond variables but arbitrarily many of them
as the volume grows. The norm and all constants used below are those
of the ordinary regular expansion, uniform in that number. A fixed
finite number of potential types can be combined by addition of their
norms.

Let $E_X(\mathbf V)$ denote the actual union-filed logarithmic term of
the local Gaussian functionals, Mayer expansion and connected-cluster
expansion. These local objects are defined on their independent-current
pair cores. Away from complete physical data, their wall-component
functionals need not be margins of one global Gaussian measure.
Normalized global Gaussian expectations will be used only at a complete
physical pair $P=p_W$, or by holomorphic continuation in the physical
datum $W$ where the same complete construction is defined. We use
the following consequence of its uninserted estimate, as established
for the fixed-reference integral in
\cref{ch:finite-source-insertion}. At every base system used
below there is an available potential allowance $w_\star>0$ such that
\begin{equation}
 |E_X(\mathbf V+\mathbf U)|
       \le C_\star e^{-\lambda_{\rm out}d(X)},
       \qquad \|\mathbf U\|_\lambda<w_\star.
 \label{sp:unmarked-allowance}
\end{equation}
The actual finite local Gaussian functionals and the absolutely
convergent logarithmic series are holomorphic in every finite collection of potential
parameters on this allowance. A single additional family is enough:
several auxiliary directions are added together in $\mathbf U$, with
the sum of their norm loads less than $w_\star$.

The locality needed here is part of this actual expansion, not a
choice made after taking its logarithm. A contribution filed at $X$
uses only selected interaction domains whose enlarged union is in
$X$. In particular, if a family $\mathbf a_b$ vanishes except on
domains containing $b^\sharp$, then
\begin{equation}
 DE_X(\mathbf V)[\mathbf a_b]=0,
 \qquad
 D^2E_X(\mathbf V)[\mathbf a_b,\mathbf v]=0
                  \quad\hbox{if }b^\sharp\not\subset X.
 \label{sp:root-locality}
\end{equation}
The assertions include domains with a zero unmarked potential. In a
Mayer derivative their zero factor can be replaced by the nonzero
inserted factor, and the marked domain must then be included in the
union. Omitting such terms would invalidate both locality and the
insertion identity.

\begin{lemma}[Two derivatives of the actual localized term]
\label{lem:root-source-local-Cauchy}
For potential families $\mathbf a,\mathbf v,\mathbf c$, the allowance in \eqref{sp:unmarked-allowance} gives
\begin{equation}
 \begin{split}
 |D^2E_X(\mathbf V)[\mathbf a,\mathbf v]|
 &\le \frac{16C_\star}{w_\star^2}
       \|\mathbf a\|_\lambda\|\mathbf v\|_\lambda
                      e^{-\lambda_{\rm out}d(X)},\\
 |DE_X(\mathbf V)[\mathbf c]|
 &\le \frac{2C_\star}{w_\star}
       \|\mathbf c\|_\lambda e^{-\lambda_{\rm out}d(X)}.
 \end{split}
 \label{sp:mixed-Cauchy}
\end{equation}
These derivatives preserve \eqref{sp:root-locality}.
\end{lemma}

\begin{proof}
If either direction in the first estimate has zero norm, the claim
is immediate. Otherwise put the single auxiliary family
$\zeta\mathbf a+z\mathbf v$ in the available slot. On the bidisc
with radii $w_\star/(4\|\mathbf a\|_\lambda)$ and
$w_\star/(4\|\mathbf v\|_\lambda)$ its total norm is at most
$w_\star/2$. Apply the two-variable Cauchy formula to the actual
holomorphic term in \eqref{sp:unmarked-allowance}. For the second
estimate use radius $w_\star/(2\|\mathbf c\|_\lambda)$.
The support assertion follows from the union filing before
differentiation. No radius for a full history construction has been
enlarged in this argument.
\end{proof}

Thus a separate second-mark estimate is not needed if the actual
ordinary bound is already uniform on the full potential allowance in \eqref{sp:unmarked-allowance}. If only a one-direction derivative
bound were known, without this allowance and locality, the lemma
would not follow from it.

\paragraph*{The interpolation paths and the differentiated density}

Let $\mathbf V_t$, $0\le t\le1$, be the actual interaction leg,
with local bond directions $\mathbf a_{t,b}$ satisfying
\begin{equation}
 \partial_t\mathbf V_t=\sum_b\mathbf a_{t,b}.
 \label{sp:actual-leg}
\end{equation}
Every summand is supported on domains containing $b^\sharp$.
At a finite volume the sum is finite. We assume that the ordinary
interpolation construction supplies the uniform bound
\begin{equation}
 \sup_b\int_0^1\|\mathbf a_{t,b}\|_\lambda\,dt\le A.
 \label{sp:ordinary-leg-norm}
\end{equation}
This is a property of the actual unmarked legs and their admissible
field domains. It is not inferred merely from an endpoint norm.
Uniform ordinary bounds and field room are required throughout the
path, with the same allowance in \eqref{sp:unmarked-allowance}.
The third-order vanishing of the intrinsic $V^{G,A}$ makes its
$t^{-2}$ radial path regular at zero. It does not by itself identify
the actual bond split. We do not replace that split by an Euler
decomposition or any other pointing to infer
\eqref{sp:ordinary-leg-norm}.

For the inserted source let
\begin{equation}
 \mathbf V_t(s)=\mathbf V_t+s\mathbf v_t,
 \qquad
 \mathbf a_{t,b}(s)=\mathbf a_{t,b}+s\mathbf c_{t,b},
 \qquad
 \partial_t\mathbf v_t=\sum_b\mathbf c_{t,b},
 \label{sp:source-leg}
\end{equation}
with the same bond support and with $\mathbf v_0=0$. It is enough
that these are the first source derivatives of holomorphic paths;
the affine notation in $s$ makes the formula easier to read. Suppose
\begin{equation}
 \sup_{0\le t\le1}\|\mathbf v_t\|_\lambda\le C_v\varepsilon,
 \qquad
 \sup_b\int_0^1\|\mathbf c_{t,b}\|_\lambda\,dt
                                      \le C_c\varepsilon.
 \label{sp:source-leg-norm}
\end{equation}
Choose $|s|C_v\varepsilon<w_\star/4$. The remaining allowance
in \eqref{sp:unmarked-allowance} bounds the first derivative in
direction $\mathbf a_{t,b}+s\mathbf c_{t,b}$ by a constant times
$\|\mathbf a_{t,b}\|_\lambda+|s|\|\mathbf c_{t,b}\|_\lambda$.
This is integrable by \eqref{sp:ordinary-leg-norm} and
\eqref{sp:source-leg-norm}. On a smaller source disc the mixed
Cauchy estimate similarly dominates its source derivative.
Consequently the actual path integral is holomorphic there and
may be differentiated under the integral.

The actual interaction-leg term carrying bond $b$ is
\begin{equation}
 E^\bullet_s(X,b)=\int_0^1
 DE_X(\mathbf V_t(s))[\mathbf a_{t,b}(s)]\,dt.
 \label{sp:actual-pointed-leg}
\end{equation}
Its source derivative contains two terms:
\begin{equation}
 \dot E^\bullet(X,b)=\int_0^1
 \left\{D^2E_X(\mathbf V_t)[\mathbf a_{t,b},\mathbf v_t]
             +DE_X(\mathbf V_t)[\mathbf c_{t,b}]\right\}dt.
 \label{sp:differentiated-leg}
\end{equation}
The first is the connected root-plus-source term. The second is
the direct change of the leg insertion. Neither can be dropped.
\eqref{sp:mixed-Cauchy}, \eqref{sp:ordinary-leg-norm}
and \eqref{sp:source-leg-norm} prove
\begin{equation}
 |\dot E^\bullet(X,b)|
 \le C_\star\left(\frac{16AC_v}{w_\star^2}
                     +\frac{2C_c}{w_\star}\right)
                \varepsilon e^{-\lambda_{\rm out}d(X)}.
 \label{sp:bond-bound}
\end{equation}
It is zero unless $b^\sharp\subset X$.

Here is the exact normalized-expectation interpretation on a complete
physical pair $P=p_W$. Write
$\langle\cdot\rangle_t$ for the normalized, possibly complex,
fixed-reference integral with interaction $\sum_YV_t(Y)$.
The denominator is nonzero by the same ordinary expansion.
With $a_{t,b}=\sum_Ya_{t,b}(Y)$,
$v_t=\sum_Yv_t(Y)$ and $c_{t,b}=\sum_Yc_{t,b}(Y)$,
\begin{equation}
 \partial_s\langle a_{t,b}(s)\rangle_{t,s}\big|_{s=0}
   =\langle c_{t,b}\rangle_t
           +\langle a_{t,b}v_t\rangle_t
           -\langle a_{t,b}\rangle_t\langle v_t\rangle_t.
 \label{sp:Duhamel}
\end{equation}
This follows by differentiating its numerator and denominator.
It requires no positivity. On these complete data, summing
\eqref{sp:differentiated-leg} over $X$ gives \eqref{sp:Duhamel}
by the physical sum rule for the actual connected expansion.
The disconnected product cancels; an unnormalized insertion would
give the wrong formula and could carry a volume factor. On an
arbitrary independent-current core,
\eqref{sp:differentiated-leg} remains the local derivative identity
and bound. It is not identified there with a covariance under an
unavailable global Gaussian measure.

\paragraph*{The older-weight path and its field room}

An older weight enters through its fluctuation difference
$v(Y;P,B)$, so $v(Y;P,0)=0$. Its interaction path is
$v_t(Y;P,B)=v(Y;P,tB)$, and
\begin{equation}
 c_{t,b}(Y;P,B)=\sum_{a=1}^3
                   B_b^a\partial_{B_b^a}v(Y;P,tB).
 \label{sp:older-weight-root}
\end{equation}
The intrinsic path $t^{-2}V^{G,A}(tB)$ belongs to the terms with
the action or gauge-fixing inverse-coupling coefficient. It must
not be applied to this older weight. In particular, no artificial
vanishing of its linear and quadratic fluctuation terms is assumed.

We spell out sufficient field room for \eqref{sp:source-leg-norm}.
Suppose $v$ is holomorphic and obeys
\eqref{sp:potential-norm} with norm $\varepsilon$ on
$\mathcal B_1$, for the actual external pairs and decoupling
parameters. We prove the bound on the entire $\mathcal B_0$
required by \eqref{sp:unmarked-allowance}, not merely on real
integration points. For real
$t\in[0,1]$, vary one coordinate at $tB$ by a complex number of
radius
\[
 r_b=c_\nu(1+|B_b^a|),
\]
where $c_\nu>0$ is a sufficiently small fixed fraction of the
available room and the fixed color norm. This disc lies in
$\mathcal B_1$, uniformly for $B\in\mathcal B_0$ and in
$t,b$ and the number of variables. Moreover its growth factor is
bounded by $C_\nu Q_Y(B)$: the one changed coordinate adds at
most $C_\nu(1+|B_b^a|^2)$ to the sum defining $Q_Y$, with $M$
fixed. Cauchy's formula therefore gives
\begin{equation}
 |B_b^a\partial_{B_b^a}v(Y;P,tB)|
       \le C_\nu\varepsilon e^{-\lambda d(Y)}Q_Y(B).
 \label{sp:coordinate-source-root}
\end{equation}
Also $Q_Y(tB)\le Q_Y(B)$, so
$\|\mathbf v_t\|_\lambda\le\varepsilon$ on the actual cutoff
box. Summing the three color coordinates in
\eqref{sp:older-weight-root} proves \eqref{sp:source-leg-norm}.
The original $B$ is used in this majorant, so no $1/t$ or
$\mathfrak b_0^2$ loss occurs as $t$ tends to zero. Thus the
differentiated root is a potential on exactly $\mathcal B_0$,
while the undifferentiated source used its one fixed larger box.

A complex disc for a common radial parameter alone would not
justify \eqref{sp:coordinate-source-root}. The individual bond
legs require the independent-coordinate field room just used.
Nor does an endpoint inserted-potential norm bound the variation
of an arbitrary interpolation path. The prescribed path and its
ordinary bounds in \eqref{sp:ordinary-leg-norm} remain necessary.

\paragraph*{Midpoint tents, entropy, and the exact sum rule}

Use the actual geometric tents at bond midpoints. The collars in
the union filing are assumed large enough that
$h_z(b)\ne0$ implies $z\in b^\sharp\subset X$ for a nonzero
bond-root term. The fixed block ratio $L$ gives
\begin{equation}
 \sum_{z\in X}h_z(b)=1,
 \qquad \sum_bh_z(b)\le4L^4.
 \label{sp:tents}
\end{equation}
These geometric identities are source independent. The renewed
pointed source family is therefore
\begin{equation}
 \dot F(X,z)=\sum_bh_z(b)\dot E^\bullet(X,b),
 \qquad
 \sup_{X,z}e^{\lambda_{\rm out}d(X)}|\dot F(X,z)|
                                      \le C\varepsilon.
 \label{sp:pointed-source-bound}
\end{equation}
The constant depends on $L$, the fixed ordinary allowances, $A$
and the field room, and is independent of volume and field-variable
count. No additional source-position factor is spent in this
supremum norm.

For completeness, the geometric root cost also remains harmless
when an application sums domains through $z$. If that count at
length $m$ is at most $C_D e^{\kappa_Dm}$ and
$\lambda_{\rm out}>\kappa_D$, then a smaller output rate costs
only
\[
 4L^4 C\varepsilon\sum_{m\ge0}(m+1)^q
                e^{-(\lambda_{\rm out}-\kappa_D)m}<\infty
\]
for every fixed polynomial power $q$. The two-variable Cauchy
estimate has already summed all cluster multiplicities and choices
of source domains through the ordinary rooted majorant. No second
unsummed count of them is justified or needed. Additional geometric
position sums cost a polynomial only when the actual domain
geometry supplies that bound. This does not permit unrecorded extra
domain types or an exponent loss larger than the available gap.

Finally sum \eqref{sp:actual-pointed-leg} over $b$ at a fixed
localization domain $X$. The chain rule
and \eqref{sp:actual-leg} give
$E_X(\mathbf V_1(s))-E_X(\mathbf V_0(s))$.
Because $\mathbf v_0=0$, its source derivative is
$DE_X(\mathbf V_1)[\mathbf v_1]$. On every declared local core,
\eqref{sp:tents} therefore proves the algebraic identity
\[
 \sum_{z\in X}\dot F(X,z)
                         =DE_X(\mathbf V_1)[\mathbf v_1].
\]
At a complete physical pair $p_W$, the physical Gaussian sum rule
then gives
\begin{equation}
 \sum_{X,z}\dot F(X,z)
       =\sum_XDE_X(\mathbf V_1)[\mathbf v_1]
       =\partial_s\log\Xi(\mathbf V_1+s\mathbf v_1)\big|_0.
 \label{sp:pointed-source-sum}
\end{equation}
The same formula is valid on the permitted analytic continuation
of the physical datum $W$. This is not continuation from the
physical graph to an arbitrary independent current. All sums are
finite-volume sums followed by absolutely convergent cluster sums.
Coherent infinite-volume extensions of the local families use the
same uniform local bounds and companion limits, when supplied.

If the actual reference, union filing and interpolation paths are
covariant under lattice isometries and a constant source has that
same covariance, these identities also hold after transporting
$X,b,P$. The midpoint tent then gives
$\dot F(\mathsf rX,\mathsf rz;\mathsf rP)
 =\dot F(X,z;P)$. The same reasoning preserves gauge invariance
and reality. These properties are inherited from the actual
construction, not inferred from a scalar norm estimate. A
nonconstant source does not acquire Euclidean covariance merely
from \eqref{sp:pointed-source-bound}.

If the source entered a Gaussian structure or the cutoff leg,
additional derivatives would be required. With a fixed reference,
a cutoff root is in general a boundary functional $N_{t,b}$, not the
bounded function $\partial_t\chi/\chi$. Its differentiated
normalized density is
$N_{t,b}(v e^V)/\Xi-(N_{t,b}(e^V)/\Xi)\langle v\rangle$,
plus any direct variation of $N_{t,b}$. If the reference measure
also varies, its score contributes to the denominator derivative
as well. The retained Gaussian tail would have to be proved under
insertion. The fixed-reference
older-weight presentation avoids those additional terms here.

The conclusion is one finite-step pointed insertion estimate.
It does not differentiate a full history class, prove closure of
the response under iteration, supply the spatial defect gain of
a nonconstant profile, or establish the clean-history Gaussian
disintegration needed for an interacting two-point law.


\paragraph*{The three legs and the normalized intrinsic path}

We make the finite structural issue in \eqref{sp:ordinary-leg-norm} explicit.  The free
potential is a compensation of the fixed gauge quadratic form.  If
\(Q^G=\sum_Y Q_Y^G\) and the forms are symmetric in the free real
coordinates, put
\begin{equation}
 q_b(Y;B)=\tfrac12 B(b)\mathbin{\cdot}(Q_Y^G B)(b),
 \qquad \sum_bq_b(Y;B)=Q_Y^G(B).
 \label{fm:gauge-quadratic-root}
\end{equation}
The dot product here is bilinear, so the formula has its holomorphic
meaning on a complex field box.  When the actual row bounds give
\(\|\mathbf q_b\|_{\lambda_V}\le A_G\), the path
\(-u(1-\tau)Q^G\), \(0\le u\le1\), has root
\(-(1-\tau)q_b\).  The inserted-potential estimate gives its pointed
bound as soon as this globally defined free family and its path fit in
the ordinary potential budget.  In particular one needs the actual
smallness of \(|1-\tau|A_G\); this is not a statement for arbitrary
positive \(\tau\).  The gauge form has not been multiplied by the
variable action weight.

The cutoff leg is next, before the intrinsic and older potentials have
been turned on.  Its only interaction is this globally defined
quadratic free potential.  Each sharp test is brought from infinity
to its final position.  Its boundary estimate is proved below.

The third leg turns on the Jacobian and older functions by
\(V_Y(tB)\).  The gauge and action remainders are normalized
differently:
\begin{equation}
                  V^{\rm int}_{t,Y}(B)=t^{-2}V_Y(tB),
       \qquad V_Y(0)=DV_Y(0)=D^2V_Y(0)=0.
 \label{fm:normalized-intrinsic-path}
\end{equation}
The three vanishings concern the intrinsic remainder after its
quadratic part has been put in the free Gaussian.  Its actual
bond derivative includes the derivative of the factor \(t^{-2}\).
The diagonal path in \eqref{fm:normalized-intrinsic-path} alone does
not specify how that derivative is distributed among bonds.

Here is a fully specified sufficient convention.  It is a construction
of an admissible local pointing; its identification with any separately
prescribed pointing remains an exact finite structural question.
For every intrinsic domain \(Y\), let \(\mathfrak B_Y\) be its actual
free bonds and choose fixed, field-independent numbers
\begin{equation}
 \vartheta_{Y,b}\ge0,\quad
 \vartheta_{Y,b}=0\ (b\notin\mathfrak B_Y),\quad
             \sum_{b\in\mathfrak B_Y}\vartheta_{Y,b}=1 .
 \label{fm:normalization-shares}
\end{equation}
The choices are covariant under the geometric symmetries of the
assigned domains.  Equal shares are possible on any nonempty such
domain.  If there are no free variables, the intrinsic member with
the vanishings in \eqref{fm:normalized-intrinsic-path} is zero and is
omitted.  Define
\begin{equation}
 \begin{split}
 T_Y(\mathbf t)&=\sum_b\vartheta_{Y,b}t_b,\\
 \widetilde V_Y(\mathbf t;B)
       &=T_Y(\mathbf t)^{-2}V_Y((t_b B(b))_b),\\
 a_{t,b}(Y;B)
       &=t^{-2}B(b)\mathbin{\cdot}D_bV_Y(tB)
                     -2\vartheta_{Y,b}t^{-3}V_Y(tB).
 \end{split}
 \label{fm:explicit-intrinsic-pointing}
\end{equation}
At \(\mathbf t=t\mathbf1\), the second line is exactly
\eqref{fm:normalized-intrinsic-path}.  Differentiating there gives
the third line, including its normalization term.  Summing over
bonds gives
\(\sum_ba_{t,b}(Y;B)=\partial_t[t^{-2}V_Y(tB)]\).
This is an independent-bond interpolation formula; no replacement by
an Euler pointing has been made.

We prove its uniform norm and domain assertions.  Use the nested
boxes in \eqref{sp:potential-norm}, and suppose that \(V_Y\) is holomorphic on the larger
one with potential norm at most \(\varepsilon\).  Fix \(B\) in the
smaller complex box.  A fixed radial disc \(|z|<1+c_\nu\) maps
\(zB\) into the larger box.  On that disc
\(Q_Y(zB)\le C_\nu Q_Y(B)\).  The function \(z\mapsto V_Y(zB)\)
has a zero of order at least three.  The maximum principle applied
to \(V_Y(zB)/z^3\) therefore gives
\begin{equation}
 |V_Y(tB)|\le C_\nu\varepsilon
             e^{-\lambda_Vd(Y)}t^3 Q_Y(B),
                      \qquad 0\le t\le1 .
 \label{fm:intrinsic-cubic-schwarz}
\end{equation}
To bound a coordinate derivative, leave a fixed further margin in
the radial disc and vary one color coordinate by a complex \(w\)
with \(|w|<c_\nu(1+|B_b^a|)\).  The resulting point is still in
the larger box, and
\(Q_Y(zB+we_{b,a})\le C_\nu Q_Y(B)\).
Coordinate Cauchy gives a bound
\(C_\nu\varepsilon e^{-\lambda_Vd(Y)}Q_Y(B)\) on
\(B_b^aD_{b,a}V_Y(zB)\), uniformly on a fixed disc of radius
greater than one.  This last function of \(z\) vanishes to order
at least two.  Applying the same maximum-principle argument and
summing the three colors proves
\begin{equation}
 |B(b)\mathbin{\cdot}D_bV_Y(tB)|
       \le C_\nu\varepsilon e^{-\lambda_Vd(Y)}
                           t^2Q_Y(B).
 \label{fm:intrinsic-root-schwarz}
\end{equation}
Consequently every summand in the last line of
\eqref{fm:explicit-intrinsic-pointing} is bounded in the original
potential norm, uniformly down to \(t=0\), and
\begin{equation}
 \sup_b\int_0^1\|\mathbf a_{t,b}\|_{\lambda_V}\,dt
                 \le C_\nu\varepsilon .
 \label{fm:constructed-pointing-integrability}
\end{equation}
The apparent negative powers of \(t\) have been canceled by the
actual orders of vanishing.

The independent parameters have a common admitted neighborhood
at every \(t>0\).  Choose \(c_\nu>0\) smaller if necessary and take
\(\max_b|t_b-t|<c_\nu t\).  Nonnegative shares imply
\(|T_Y(\mathbf t)-t|<c_\nu t\), so
\(|T_Y(\mathbf t)|\ge t/2\).  The scaled coordinates remain in the
larger fluctuation box.  Applying the radial argument to
\((t_b/t)B(b)\) gives
\(\|\widetilde{\mathbf V}(\mathbf t)\|_{\lambda_V}
\le C_\nu t\varepsilon\).  Thus the potential path also stays
inside its accepted slot if its initial allowance is chosen
small enough.  There is no singular parameter domain being
asserted at \(t=0\); the integrable diagonal limit is the one
just proved.

Each density is zero unless \(b\) belongs to its actual domain.
Its root therefore survives the actual union filing.  Fixed
covariant shares introduce neither field dependence nor a new
read outside that domain.  For Jacobian and older families the
unnormalized path \(V_Y(tB)\) uses only the coordinate Cauchy
bound in \eqref{sp:coordinate-source-root}; no third-order assumption is needed for its bond
derivative.

For a separately prescribed independent-bond convention, write
its normalization density as \(R_b(Y;B)\).  The same proof
applies if its exact derivative is
\(t^{-2}B(b)\cdot D_bV_Y(tB)-2t^{-3}R_b(Y;tB)\),
with \(\sum_bR_b=V_Y\), the same support and norm bound, and
third-order vanishing of each \(R_b\).  Those are concrete finite
identities and estimates to verify.  They do not follow from
the single equality \(\sum_ba_{t,b}=\partial_tV_t\): arbitrarily
large terms with zero bond sum would leave that equality
unchanged.  In particular, the construction in \eqref{fm:explicit-intrinsic-pointing} does not silently
identify inherited coefficients, own scalar readouts or
operation identities with a different convention.



\paragraph*{The moving sharp cutoff has an integrable rooted tail}

We next prove the boundary estimate for the actual cutoff leg.
The cutoffs act on three real coordinates of each tested free bond,
with the Euclidean Lie-algebra norm.  A cutoff on a further linear
image would require its corresponding uniformly nondegenerate
marginal bound; it is not identified with a free coordinate by name.
Put
\begin{equation}
 \chi_t(B(b))=\mathbf1_{\{t|B(b)|<\mathfrak b\}},
 \qquad 0\le t\le1,\quad \mathfrak b\ge1.
 \label{fm:moving-sharp-cutoff}
\end{equation}
At zero the cutoff is one.  As a measure derivative for \(t>0\),
\(-\partial_t\chi_t=|B(b)|\delta(\mathfrak b-t|B(b)|)\).
For every fixed \(B(b)\) away from the final boundary its integral
over \(0<t<1\) is exactly
\(\mathbf1_{\{|B(b)|>\mathfrak b\}}\).  This is the one-dimensional
change of variables \(u=t|B(b)|\).  It proves the total-variation
identity without dividing by a characteristic function.

For use with complex Gaussian matrices we require a uniform
majorant before integrating in \(t\).  Let \(D\) be the local
covariance after the auxiliary coordinates have been integrated
out as in \eqref{eq:localized-gaussian-comparison-FG6}.  Its real reference and perturbation obey the
same uniform gap bounds as there.  If
\(|F(B)|\le F_0e^{\rho|B|^2}\), define
\begin{equation}
 S=\operatorname{Re}(D^{-1})-2\rho I,\qquad
 Z_{\rm abs,\rho}=|\det D|^{-1/2}(\det S)^{-1/2}.
 \label{fm:absolute-tilted-Gaussian}
\end{equation}
For the permitted \(\rho\) and perturbation size,
\(S>0\), both \(S\) and \(S^{-1}\) have uniform operator
bounds, and
\begin{equation}
 |d\mu_D(B)|\,e^{\rho|B|^2}
        =Z_{\rm abs,\rho}\,d\mu_{S^{-1}}(B),\qquad
 Z_{\rm abs,\rho}\le e^{A_K(\widehat\varepsilon+\rho)m}.
 \label{fm:absolute-Gaussian-normalization}
\end{equation}
This follows by writing the density on the real contour.
The bound is the normalized local comparison bound, with the
local scalar-coordinate dimension \(m\).  No factor depending
on the auxiliary dimension has appeared.

Let \(P\) be a set of \(q\) other distinct failed free-bond tests.
For a sufficiently small fixed \(\gamma>0\),
\begin{equation}
 \prod_{a\in P}\mathbf1_{\{|B(a)|\ge\mathfrak b\}}
 \le
 e^{-\gamma q\mathfrak b^2/2}
     \exp\{\tfrac{\gamma}{2}\sum_{a\in P}|B(a)|^2\}.
 \label{fm:other-failures-tilt}
\end{equation}
Writing \(\Pi_P\) for the orthogonal projection onto their
\(3q\) real coordinates, the new real precision is
\(S-\gamma\Pi_P\).  Choose
\(\gamma\|S^{-1}\|\le1/2\).  It is positive, its covariance
has uniform upper and lower bounds, and its normalization
ratio is at most \(e^{C_K\gamma q}\).  To verify the last
claim, diagonalize
\(\gamma S^{-1/2}\Pi_PS^{-1/2}\).  There are at most \(3q\)
nonzero eigenvalues, each at most \(1/2\), and
\(-\log(1-u)\le2u\) on that interval.  Thus the logarithm of
the ratio is at most \(3K'\gamma q\).  This cost counts
failed coordinates, rather than all coordinates outside
the local Gaussian.

The marginal of the root block \(B(b)\) in this tilted
normalized real Gaussian is three dimensional, even if
the full local dimension is large.  Its radial density
is bounded by \(C_K u^2e^{-c_Ku^2}\).
Indeed its covariance is a principal block of the full
covariance, with the same uniform eigenvalue bounds.
The usual three-dimensional Gaussian density is therefore
at most \(C_Ke^{-c_K|x|^2}\); integration over the sphere
of radius \(u\) gives the assertion.  Consequently,
uniformly in all admitted complex matrix parameters,
\begin{equation}
 \begin{split}
 &\left|{\cal J}\left[
    |B(b)|\delta(\mathfrak b-t|B(b)|)\,
    F(B)\prod_{a\in P}\mathbf1_{\{|B(a)|\ge\mathfrak b\}}
                           \right]\right|\\
 &\quad\le
 C_KF_0e^{A_K(\widehat\varepsilon+\rho)m+C_K\gamma q}
       e^{-\gamma q\mathfrak b^2/2}
       \mathfrak b^3t^{-4}
                     e^{-c_K\mathfrak b^2/t^2}.
 \end{split}
 \label{fm:uniform-cutoff-boundary}
\end{equation}
This display uses the auxiliary-integrated local functional.
For a boundary distribution it follows either by polar
coordinates in \(B(b)\), or by first using a thin positive
shell in that coordinate and taking its limit.  The
Gaussian density and the free-potential factors are smooth;
the remaining cutoffs are bounded Borel functions.  Fubini
and the displayed integrable majorant justify that limit.

The singular-looking factor is integrable.  The substitution
\(u=\mathfrak b/t\) gives
\begin{equation}
 \int_0^1\mathfrak b^3t^{-4}
       e^{-c_K\mathfrak b^2/t^2}\,dt
       =\int_{\mathfrak b}^{\infty}u^2e^{-c_Ku^2}\,du
       \le C_K e^{-\gamma\mathfrak b^2/2},
 \label{fm:integrated-cutoff-tail}
\end{equation}
after reducing \(\gamma\) if necessary.  For example,
bound \(u^2e^{-c_Ku^2/2}\) by its finite maximum and
integrate the remaining Gaussian tail; with
\(\mathfrak b\ge1\) the extra polynomial is absorbed by
a still smaller fixed exponent.  The bound in
\eqref{fm:uniform-cutoff-boundary} is uniform before
the \(t\) integration.  It therefore also bounds an
integral whose integrand is a supremum over the admitted
matrix and external parameters.  This constructs the boundary envelope with the supremum inside
the integral.

We spell out its application to the literal finite
expansion.  During this leg only the globally defined
quadratic free family is present, so its selected
interaction domains require no small-field box
\(Y_0\).  Expand \(\prod_b\chi_{t_b}\) into failed
tests.  Differentiation in a tested bond \(b\) selects
terms whose failed set contains \(b\), replaces its
factor \(1-\chi_{t_b}\) by the boundary measure, and
leaves every other failed factor at threshold at least
\(\mathfrak b\).  Use
\eqref{fm:uniform-cutoff-boundary} for that root.
The independent coupling Schwarz estimates for the
selected free-potential Mayer factors, the enlarged
coarse-decoupling Cauchy estimates and the coarse
collar count are exactly those used in \eqref{eq:local-formation-UM13}--\eqref{eq:local-formation-UM24}.
The root bond is part of the actual support
\(Z_0(\mathbf D,P_0)\).  Thus no remote root is created.

For precision, the Schwarz argument is a holomorphic
one and never squares the boundary measure.  At fixed
\(t>0\), insert independent couplings \(\lambda_Y\)
in the selected free-potential Mayer factors.
The resulting boundary functional is holomorphic
on the discs in \eqref{eq:local-formation-UM12}, with the unsquared boundary
factor kept inside that linear functional.
Indeed the larger Gaussian-only budget bounds its
integrand by \(2^{\#\mathbf D}e^{CW_0L^4N}
e^{\rho|B|^2}\); \eqref{fm:uniform-cutoff-boundary} gives a common boundary
majorant on those discs.  Differentiation or a
contour integral in the couplings is therefore
justified by this same majorant.  The functional
vanishes when any selected \(\lambda_Y=0\).
Apply the scalar Schwarz lemma successively in
these variables, and then Cauchy's formula in
the coarse parameters.  Writing
\(r=\#\mathbf D\), \(v=|P_0|\ge1\),
\(h=|Z\setminus Z_0^\sharp|\),
\(p=e^{-\gamma\mathfrak b^2/2}\), and
\(g_{\rm bd}(t)=\mathfrak b^3t^{-4}
e^{-c_K\mathfrak b^2/t^2}\), this gives
\begin{equation}
 \begin{split}
 |\partial_{t_b}T_{Z,\mathbf D,P_0}|
 &\le C_K g_{\rm bd}(t)
       e^{C_0N+C_K\gamma(v-1)}
       \left({2a\over W_0}\right)^r\\
 &\quad{}\times
 e^{-(\lambda_V-\lambda_0)\sum_{Y\in\mathbf D}d_f(Y)}
          p^{v-1}(R_s-1)^{-h}.
 \end{split}
 \label{fm:analytic-boundary-Schwarz}
\end{equation}
The zero tests and the boundary factor have no
coupling derivative and are not subjected to an
\(L^2\) Cauchy--Schwarz inequality.  The coarse
parameters change the Gaussian matrices, not the
real boundary coordinate.  The formula in \eqref{fm:integrated-cutoff-tail}
integrates \(g_{\rm bd}\) to at most \(C_Kp\),
restoring \(p^v\) in the unmarked estimate in \eqref{eq:local-formation-UM16}.
The factor \(e^{C_K\gamma(v-1)}\) costs at most
\(e^{CN}\), since \(\gamma(LM)^4\) is bounded.
Enlarge \(C_0\) by that fixed constant before
choosing \(K_*,\Theta\) and the strict rates in
\eqref{eq:local-formation-UM18}--\eqref{eq:local-formation-UM19}.  Thus the same literal hull and
product sums apply with an integrable root
envelope throughout.

Allocate the same strict surplus of spatial and
decoupling rates as in the unmarked proof.  The
normalization cost in
\eqref{fm:absolute-Gaussian-normalization} is bounded
per coarse cube when
\((LM)^4(\widehat\varepsilon+\rho+\gamma)\) is bounded.
The rank-tilt cost \(e^{C_K\gamma q}\) is included in
this allocation, since \(q\le C(LM)^4\#Z\).
The failed-test position sum has at most
\(C(LM)^4\) choices per coarse cube.  Every selected
cutoff term still has its root failure.  In place of
an unmarked linear amplitude this gives
\begin{equation}
 \begin{aligned}
 p_{\rm cut}&=C_K(LM)^4e^\Theta
                         e^{-\gamma\mathfrak b^2/2},\\
 \int_0^1\sup_P
       |\partial_{t_b}H_Z(\mathbf t;P)|_{\mathbf t=t\mathbf1}
                \,dt
       &\le C p_{\rm cut}e^{-\lambda_2(d_c(Z)+5)} .
 \end{aligned}
 \label{fm:rooted-cutoff-activity}
\end{equation}
Here \(\Theta\) is the rate expenditure of \eqref{eq:local-formation-UM20}--\eqref{eq:local-formation-UM22};
all physical smallness conditions include this factor.
The supremum includes the local decoupling
parameters needed in forming this activity.  The
estimate follows by the same nonnegative sum over
selected domains and failed sites as before, with
one site pinned to \(b\).  Pinning removes a position
sum; any polynomial number of possible roots inside
a domain is absorbed by the already retained strict
rate gap.  An empty choice cannot contribute a root.

Finally differentiate the absolutely convergent
cluster series once.  In each connected cluster mark
the polymer that is differentiated.  The rooted
tree bound sums all its remaining polymers to at
most \(e^{a(Z)}\), where \(a\) is the size weight of
the hard-core criterion.  Sum its marked activity
with \eqref{fm:rooted-cutoff-activity}, spending the
strict difference between \(\lambda_2\) and the
retained output rate.  Positivity of these
majorants permits Tonelli in \(t\) and in the
cluster sum.  We obtain
\begin{equation}
 \int_0^1\sup_P
    |\partial_{t_b}E_X(\mathbf t;P)|_{\mathbf t=t\mathbf1}\,dt
        \le C p_{\rm cut}e^{-\lambda_{\rm out}d(X)} ,
 \quad
 \partial_{t_b}E_X=0\quad(b^\sharp\not\subset X).
 \label{fm:rooted-cutoff-logarithm}
\end{equation}
At a fixed pair on \(X\), only the hereditary subsystem
of polymers contained in \(X\) is used.  When the union
condition is dropped in the positive bound, only the
uniform local norm majorants are extended, not actual
activity functions outside their declared cores.
Its integral is the cutoff-leg density
\(E_\chi(X,b)\).  Summing in \(b\) recovers the
difference between the all-one and all-zero cutoff
endpoints.  Midpoint tents preserve its root as in
\eqref{sp:tents}--\eqref{sp:pointed-source-sum}.

All these estimates are local on their admitted
independent-pair cores.  The complete sum is
identified with the physical cutoff integral only
at \(p_W\) and its admitted continuation in \(W\).
No global Gaussian on arbitrary independent
currents, and no estimate for a history's outer
stored-variable integral, has been used.


\subsubsection{Window-first ordinary formation}


The estimate in this section is an ordinary formation estimate.
It concerns the actual functions obtained by rewriting old terms on
box backgrounds, before their fine-cube bracket localization.
The outer potential/current pair may have independent components.
This causes no substitution of an arbitrary pair into a theorem
about critical backgrounds: the inner window function first
constructs its own constrained box background and its own current.
The hypotheses below retain the domain and response assertions needed
to perform that construction on the whole complex parameter set.
We fix the window-first rule below as part of the localization:
at one old base point, its near pointed sum is one window
member before any fine bracket is applied.

\paragraph*{The maps and geometric inputs}

Fix a step from level \(k\) to level \(k+1\).
For an old level \(j\le k\), put
\begin{equation}
 a=k-j,\qquad \sigma_j=L^{-a},\qquad
 \rho_j=\max\{1,\alpha_k/\alpha_j\},\qquad
 s_k=\delta_k/\alpha_k .
 \label{eq:local-formation-WF1}
\end{equation}
The radii and the ordinary flow satisfy
\(\rho_j\le C(1+a)^{1/2}\).
In particular
\begin{equation}
 C_{\rm age}:=\sup_{a\ge0}C(1+a)^{1/2}L^{-a}<\infty,
 \qquad \rho_j\sigma_j\le C_{\rm age}.
 \label{eq:local-formation-WF2}
\end{equation}
Here \(\delta_k\) is the tested fluctuation amplitude, whereas
\(\alpha_k\) is the analytic field radius.  The ratio in \eqref{eq:local-formation-WF1}
is an ordinary full-radius ratio.  No small-field gain is being
deduced by replacing a smaller ratio with one.

We spell out the ordinary background and geometry inputs.  They
must concern the actual assigned maps, not arbitrary functions
with similar notation.

First, let \(P\) range over the local complex core used to form the
new potentials.  Let \(\mathbf H_k(P;s,B')\) be the actual response
to the tested free fluctuation coordinates, with the indicated
fine-cube parameters.  The plain nonlinear iterated block average
is denoted \(\mathcal M^j\).  On a window \(\square\), with a fixed
partition function \(\zeta_\square\), the datum used for an inner
term is the restriction of
\begin{equation}
 {\cal V}_{j,\square}(P;s,t,\nu)=
 M_{\mathbf B_j}\!\left(
 e^{\,i\eta_k(t+\nu\zeta_\square)
                  \mathbf H_k^{\rm win}(P;s,B')}\,U_+(P)\right),
 \qquad 0\le t\le1 .
 \label{eq:local-formation-WF3}
\end{equation}
The notation $U_+(P)$ denotes the exponential representative
of the upper box. The response $\mathbf H_k^{\rm win}$ is transported
to that same frame. The original node frame is recovered at the two
endpoints of each row. The factor $\eta_k$ converts a potential
in level-$k$ units to a bare bond logarithm. The map $M_{\mathbf B_j}$
reads the ordinary iterated average at each determining row, at
the level of that row.  The formula in \eqref{eq:local-formation-WF3} uses the
nonlinear average.  Replacing it by its differential would not
give the window identity.
The box-background map applied to this datum returns a pair
\begin{equation}
 P^{\rm box}_{j,\square}(P;s,t,\nu)
   =\big(U^{\rm box}_{j,\square}({\cal V}_{j,\square}),
                    J^{\rm box}_{j,\square}({\cal V}_{j,\square})\big)
 \label{eq:local-formation-WF4}
\end{equation}
in which \(J^{\rm box}\) is the current of that constrained critical
configuration.  The ordinary maps include their prescribed node
and gauge representatives.

Second, all these maps are holomorphic on the full independent
fine-parameter discs \(|s_\Delta|<R=e^{\bar\kappa}\).
The unamplified substitutions have a fixed positive amount of
datum, chart and core room.  If \(r_{B,j}\) is the dimensionless
datum radius, the background package supplies
\(\alpha_j/r_{B,j}\le C\).
Its response and nonlinear-average estimates bound the change
caused by the \(\nu\) term, in every relevant normalized datum,
chart and core norm, by
\begin{equation}
 C_R|\nu|\,\delta_k\sigma_j/\alpha_j
            \le C_R|\nu|s_k\rho_j\sigma_j .
 \label{eq:local-formation-WF5}
\end{equation}
The constant \(C_R\) includes the fixed loss from the enlarged
fine-parameter discs.  These are estimates for the actual
nonlinear average and response on their assigned analytic
domains.  For example, in a datum norm the estimate follows by
the uniform first-response bound, composed with the derivative
of the actual average, and integration along the perturbation.
For the inner terms it must also hold for the other core
inequalities and the tower substitutions. The outer and removed
terms use their own directional estimates below.
Existence and holomorphy on the corresponding datum balls are
part of the background package.

\paragraph{Restriction and the two substitutions.}
Use the two fixed tower rules on the admitted localization domains,
with inward rows removed. Here $L$ is odd, $L\ge13$, the geometric
sizes are powers of $L$, $R_1\ge L$, $M_1\ge R_1\mathfrak M$, and
$M\ge L^4M_1$. The localized tower estimates use the separated-region
floors and the fixed continuation and averaging radii of the
background construction. In particular they use
$\rho=L^{-2}R_1\mathfrak M\ge\mathfrak r^{\mathrm s}$, the
last-level torus width of at least $6\rho$ sites, and the floors
$\mathfrak M\ge\mathfrak m_{\mathcal M}$ and
$R_1\mathfrak M\ge\max\{\mathfrak r_{\mathcal M},\mathfrak r^{\mathrm{prop}}\}$.
The corresponding decay rates satisfy
\[
 \delta_{0,\mathrm{red}}R_1\mathfrak M\ge8(4\log L+\log 2),\qquad
 \mathring\delta_{0,\mathrm b}R_1\mathfrak M\ge512\log L,\qquad
 \delta_*^\natural R_1\mathfrak M\ge18(4\log L+\log 2).
\]

For nested localization domains $X\subset Y$ of levels $j-1$ and
$j'-1$, set $\lambda=L^{-(j'-j)}$. Restriction from $Y$ to $X$ maps
the core at radii $(b_0,b_1,c_0)$ into the old core at radii
$(\lambda b_0,\lambda b_1,\lambda^3c_0)$. In the simple clauses
of \eqref{eq:rg-objects-R4}, the curvature and the first difference
of the potential gain $\lambda^2$, the potential gains $\lambda$,
and the independent current gains $\lambda^3$. The tower of any
admitted $Z\subset X$ is computed from the same bare product in both
presentations. Its curvature and current, read at their own levels
as in \eqref{eq:rg-objects-R6}, inherit the ratio $\lambda$ in the
hereditary ceiling. If $j'=k+1$, $m=k+1-j\ge1$ and
$A_i\le2\sqrt{1+m}\,a_i$ for $i=0,1$, this restriction sends the
radii $((1+2\beta_{\mathrm s})A_0,(1+2\beta_{\mathrm s})A_1,A_0)$
into $(a_0/3,a_1/3,a_0/3)$ for $0<\beta_{\mathrm s}\le1/4$.
Indeed $(1+2\beta_{\mathrm s})2\sqrt{1+m}\,L^{-m}\le1/3$.

For the outer substitution in level-$j$ units, take the old triple
in $\mathcal U_j^\circ(X;b_0/2,b_1/2,c_0)$, with
$c_0>0$, $0<b_0\le b_1$, $b_1^2\le c_\natural b_0$, $c_\natural\ge1$ and
$b_0+b_1\le a^\sharp$. Choose $a^\sharp<a_T$ and below the
fixed ceiling of \cref{htc:outer-substitution}.
Take $K_{\mathrm{red}}^\sharp\ge K_T$ from \cref{htc:all-towers}, enlarged to cover the fixed substitution constants.
If the potential shift satisfies
\[
  \max\{|\mathbf A|,|\nabla_{\mathbf U}\mathbf A|\}
       <\frac{b_0}{128K_{\mathrm{red}}^\sharp(1+c_\natural)}
\]
on the internal bonds and differences, and $|\mathbf J'|<b_0$,
the triple of $(e^{i\eta_j\mathbf A}\mathbf U,\mathbf J')$ obtained
by \eqref{eq:rg-objects-R8} lies in
$\mathcal U_j^\circ(X;b_0,b_1,b_0)$. In fact, multiply the shift
by $z$ with $|z|\le R=4K_{\mathrm{red}}^\sharp(1+c_\natural)$.
Write $\beta_0=b_0/32$. The shift and its first difference
with transport by $\mathbf U$ have size at most $\beta_0$
on this disc. Transport by the real base $U$ changes the
difference by at most $\beta_0/4$. The local logarithm
estimate gives bounds below $b_1$ for the potential and its
first differences. In the plaquette product estimate, the
change relative to the old plaquette is bounded by
$\eta_j^2(2\beta_0+14\beta_0^2+4\eta_j\beta_0 b_0)$.
Multiplying back by the old plaquette gives a change at most
$3\beta_0\eta_j^2$, within the curvature margin.
For a tower of level $p\le j$ put
$\lambda_p=L^{-(j-p)}$. Each normalized curvature or current
readout of that tower is holomorphic in $z$ and bounded on the
disc by $K_{\mathrm{red}}^\sharp(1+c_\natural)\lambda_pb_0$.
Its value at zero has size less than $\lambda_pb_0/2$. Schwarz's
lemma applied to the difference from its value at zero gives a
change at $z=1$ of at most $\lambda_pb_0/2$. This checks the
hereditary clauses for both rules and every truncation. The real
base is unchanged and the independent current is tested separately.

The inner window estimate starts with the flat-reference
continuation $e^{i\mathsf a_0}$ of the upper window datum,
in the node frame of the window. For a fluctuation $\mathbf A$
in level-$k$ units put
\[
 \mathsf a=\frac1i\log\bigl(e^{iL^{-k}\mathbf A}e^{i\mathsf a_0}\bigr),
 \qquad B(c)=\frac1i\log\mathcal M^q(e^{i\mathsf a})(c),\quad c\in\Gamma_q.
\]
Here $\Gamma_q$ is the level-$q$ row set of the old truncated
window tower. The actual nonlinear average is used. If $B_0$
is the same datum at $\mathbf A=0$, the logarithm bound and
holomorphy of the average on its absolute row disc give
$|B(c)-B_0(c)|\le (9/2)C_Q L^{q-k}\sup|\mathbf A|$.
Indeed the bond logarithm changes by at most
$(9/8)L^{-k}|\mathbf A|$, and the Schwarz estimate on the
line joining the two bond logarithms costs $4C_QL^q$.
The bound holds for complex fluctuations in the admitted ball.
This also controls the whole datum with its actual amplitude.
Write $d=A_0+\sup|\mathbf A|$, where $A_0$ is the upper
curvature radius on the complete box. The upper comb estimate
gives $\|B_V\|_\infty\le C_MA_0$. The restriction identity and
the weighted response of the flat upper background then give
$\|B_0\|_\infty\le C_MA_0$ on every row of the old tower.
Since $q\le j\le k$, the preceding increment estimate yields
\begin{equation}\label{wf:whole-window-amplitude}
 \|B\|_\infty\le C_Md,\qquad
 \|B\|_{\mathrm{central}}\le C_Md L^{j-k}.
\end{equation}
The second bound uses that every central row has level $j$;
its upper background contribution has the stronger factor
$L^{j-(k+1)}$. These constants depend on the fixed geometry,
not on either depth. The same first bound holds after the
cutoff axial change: its comb has $O_M(L^{k-j})$ links of
size $O_M(dL^{j-k})$, hence its logarithm has size $O_M(d)$.
Its cutoff extension has that same supremum, and the endpoint
product on every row stays within $C_Md$ of the identity.
Lower rows outside its support retain their original values.
For the central axial bound require, in addition, the
complete upper field's first-difference bound of size $C_Md$
on the settled layer, as in the window theorem.
Its level-$j$ plaquette defects are consequently at most
$C_MdL^{2(j-k)}$. An axial link is a product of at most
$C_ML^{k-j}$ such plaquette factors, with bounded transports.
It is therefore within $C_MdL^{j-k}$ of the identity on
the central carrier, where the cutoff equals one and the
complete comb rectangles remain in the settled layer.
This proves the central bound after the axial change; it
uses that first-difference estimate as well as the datum bound.
On the outer cutoff transition only the global $C_Md$
bound is asserted here.
Thus the remote response has the proportional factor $d$;
the fixed continuation ceiling supplies only its domain.
An independent correction $B_{\mathrm{err}}$ contributes
$C\|B_{\mathrm{err}}\|_\infty$ separately. In particular,
the auxiliary correction ball used for Cauchy estimates does
not change the amplitude of the actual evaluation, where
$B_{\mathrm{err}}=0$.

The window continuation at the datum $\tau B+B_{\mathrm{err}}$,
$0\le\tau\le1$, gives the old curvature, potential and current
margins $(a_0/4,a_1/2,a_0/4)$ in its normalized core.
Use the actual upper input and literal-datum convention of
\cref{htc:inner-substitution}. Choose $M_\flat\ge M_0$,
$a^\flat<a_I(M)$ and $2B_3^2\mathsf K\ge C_I$ from that
proposition, and take its increment constants $c_2,c_3$ or
smaller ones. The required bounds are $M\ge M_\flat$,
$A_0+A_1\le a^\flat$, $a_0+a_1\le a^\flat$,
$2B_3^2\mathsf KMa_0<a_1$, and $A_0\le2\sqrt{1+m}\,a_0$.
Its fluctuation has size below $c_2A_0$ and the complex correction
has raw norm below $c_3a_0$. These are local input bounds on the
entire window. In \eqref{eq:local-formation-WF3}, the fluctuation
includes $(t+\nu\zeta_\square)\mathbf H_k$ in the same node frame.
To use the window estimate, identify its upper continuation and
determining rows with the ones in that formula. Check the bounds
on the transported fluctuation and any correction on the full
retained $\nu$-disc and all independent fine-parameter discs.
In particular, if the transported response has
\(\|\mathbf H_k\|_\infty\le C_R\delta_k\) on the fine-parameter
discs, then the window input on \(|\nu|<c_\nu/s_k\) is at most
\(C_R(s_k+c_\nu\|\zeta_\square\|_\infty)\alpha_k\).
When \(\alpha_k/A_0\) is bounded, small \(s_k\) and \(c_\nu\)
make this less than \(c_2A_0\). The same row derivative bound
on the admitted window gives a datum increment of size
\(C|\nu|L^{q-k}\delta_k\) relative to \(\nu=0\).
This is the datum estimate in \eqref{eq:local-formation-WF5}
after dividing by the old radius. Its use for the actual maps
requires the comparison just described. The other core readouts
must also obey their bounds throughout the disc.
The outer estimate starts with a shift already bounded in its
required norm. The removed substitution likewise requires its
own joint background and current bounds.

% Raw Part 3 Proposition 84.2 and its proof, 54045--54122;
% Theorem 83.1(g3), 52357--52403; Definition 85.14, 57689--57724;
% Theorem 85.2, 58644ff. The Cauchy argument uses the correction
% radius in the very same window theorem.
\paragraph{The window at the radii of the run.}
Let $x_i$ be the inverse couplings and use the running radii
$\alpha_\iota(x)=C_\iota x^{-1/2}(\log x)^q$, $\iota=0,1$.
For $m=k+1-j$ set $A_\iota=(1+\beta_{\mathrm s})\alpha_\iota(x_k)$
and $a_\iota=\min\{\alpha_\iota(x_{j-1}),\alpha_\iota(x_j)\}$.
The drift bounds are $-c_5\le x_i-x_k\le\bar B(k-i)$.
If $x_i\ge x_k$, the logarithmic factor in
$\alpha_\iota(x_k)/\alpha_\iota(x_i)$ is at most one, and the
remaining factor is at most $\sqrt{1+k-i}$ by
$\bar B\gamma^2\le1$ and $x_k\ge\gamma^{-2}$. If $x_i<x_k$,
the logarithmic ratio is at most $1+c_5\gamma^2$, using
$\log x_i\ge1$. The choice $(1+c_5\gamma^2)^q\le1+\beta_{\mathrm s}$
therefore gives $A_\iota\le 2\sqrt{1+m}\,a_\iota$. Both minima
are attained at the same index, so $a_1/a_0=C_1/C_0$.
The running ceiling gives $a_1\le1$ and both radius sums at most
$a^\flat\le a^\sharp$. Thus $a_1^2\le (C_1/C_0)a_0$, and the
choice of $C_1/C_0$ gives $2B_3^2\mathsf KMa_0<a_1$.

For the inner substitution take the window of $2^4$ level-$k$
cubes with a common vertex and put $\square_0=\square^{\sim5}$.
Use the rule-$\mathfrak r_1$ tower of $\square_0$ to depth $k+1$
and its truncation to depth $j$, with the same convention on
inward rows at both depths. The upper determining data
$M_{\mathbf B^+}(\mathbf U)$ on the first tower have a comb
representative $V$ and a node gauge $v$, with
\[
 M_{\mathbf B^+}(\mathbf U)=V^v,\qquad
 |V-1|< K_\square M\alpha_0^+,\qquad
 |v|\,|v^{-1}|< e^{K_\square M\alpha_1^+},
 \qquad \alpha_\iota^+=(1+2\beta_{\mathrm s})A_\iota .
\]
The comb construction uses the upper bare plaquette and potential
bounds and the $\mathfrak r_1$ tower clause for this box. It also
uses the fixed smallness ceiling on $(LM)^4(\alpha_0^++\alpha_1^+)$.
The representative is holomorphic at fixed real base. Its node
gauge extends by the tower's prescribed site rule. The tower at
the actual upper datum is the extended gauge transform of the
flat-reference continuation at $B_V=i^{-1}\log V$.

%
%
%
\paragraph{The ordinary response and its current.}
At the ordinary step on the constant sequence, the raw source
of the response is the corrected fluctuation logarithm at level
$k$. In the free variables of \eqref{eq:local-formation-OA9} it
is $B'=z-h_Wd$. The constant-sequence datum conversion is the
identity. On the tested fluctuation box the contraction gives
$\|B'\|_\infty\le C\delta_k$. On each admitted local restriction
use the same source on the corresponding rows, extended by zero
outside its support. For the assigned response settings take
$C\delta_k<\varepsilon_{\mathrm D}e^{-64\bar\kappa}$.

Write $\mathbf H_k$ for the potential of the decoupled response
and $\mathsf j_k$ for its increment of current, both in level-$k$
units. The bounds for the solution on the independent parameter
discs give
\[
 \|\mathbf H_k\|_\infty+
 \|\nabla_{\mathbf U}\mathbf H_k\|_\infty+
 \|\mathsf j_k\|_\infty\le C_R\|B'\|_\infty .
\]
For the current use the difference estimate between $B'$ and
zero. At a local layer these are bounds in that layer's units;
the assertion here uses the constant sequence. The potential and
current maps are holomorphic at fixed real base and have the
same restriction law on the interior of active wall components,
with their own support collars. Outside the components retain
the explicit current of the reduced system. The current estimate is for the current from
the decoupled equations. A second difference of a decoupled
potential cannot be estimated from the first-difference bound.

%
%
The weighted estimate for the current follows in the
same fixed pair and tower, with the switches held fixed.
For two sources \(B'_1,B'_2\), let \(\mathbf H_i\) be their
potentials and use the multiplier of \eqref{dnr:multiplier}.
Put
\[
 Z_s=\widehat Q^\sharp H_s^\sharp\mathbf J,\qquad
 K_i=(\widehat Q+C'_{\mathbf H_i})^\sharp
                         \omega_s(\mathbf H_i),\qquad
 R_i=\mathsf f(-i\xi\operatorname{ad}_{\mathbf H_i})^{-1}.
\]
The shifted current in \cref{dnr:response} is
\(\mathbf J_i^\flat=\mathbf J-Z_s+R_iK_i\).
Since \(Z_s\) is independent of the source, its difference is
\[
 \mathbf J_2^\flat-\mathbf J_1^\flat
   =R_2(K_2-K_1)+(R_2-R_1)K_1.
\]
The local transpose bound and the multiplier estimate in
\eqref{dnr:multiplier} give a bounded unweighted current
norm for \(K_i\) and the weighted Lipschitz bound for
\(K_2-K_1\). Indeed, subtract the two multiplier formulas,
use the resolvent identity for \(\Psi_{s,\mathbf H_i}\),
and then the first-difference bounds for
\(H_s^\circledast\) and \(C'_\mathbf H\), as in the proof
of \eqref{dnr:current-bound}. The inverse factors \(R_i\)
are bounded on the common bond disc, and their difference
costs \(C\xi|\mathbf H_2-\mathbf H_1|\). Thus the second
term, in the current units at a site \(y\), is at most
\(C\xi_y|\mathbf H_2-\mathbf H_1|\), since
\(\xi\le\xi_y\). This is the weighted potential
difference bound with the same spatial weight.
Thus a source difference supported away from the
full read has the same exponential gain in the
current norm as in the potential norm, on the
admitted complex switch disc. This applies to
the cutoff change within one old box after
including the stencils of its converted rows.
The weighted H\"older difference of the potential
on that same tower follows by using the local
first-difference H\"older norm in the propagator
estimate of the response theorem. It retains
the same source difference and spatial weight.
The identification with the physical current is
at a real critical point at full switches. Its
stencil is the native action's stencil before
restriction to the read.

Consider the outer evaluation obtained by shifting the restricted
pair by the ordinary response with multiplier $\lambda$. In
level-$j$ units its potential shift and current are
\[
 \mathbf A_{j,\lambda}=\lambda\sigma_j\mathbf H_k|_X,
 \qquad
 \mathbf J_{j,\lambda}=L^{-3(k+1-j)}\mathbf J_+|_X
                +\lambda\sigma_j^3\mathsf j_k|_X,
\]
where $\mathbf J_+$ is the input current in level-$(k+1)$ units.
The shift in the bare exponential is unchanged because
$\eta_j\sigma_j=\eta_k$. The covariant difference of the shift
gains $\sigma_j^2$ since the transport uses the same bare bond
field. If the upper triple is in the upper core used in the
two-level restriction, the first current on the right is within
$a_0/3$. The potential, its covariant first difference, and the
current increment are all bounded by
$C_R|\lambda|\sigma_j\delta_k$. On
$|\lambda|<c\,a_0/(C_R\sigma_j\delta_k)$, with $c$ chosen below
the explicit outer shift ceiling and the remaining current
margin, the outer substitution therefore belongs to the old
core. The running radius comparison makes this disc contain
$|\lambda|<c_\nu/s_k$ for a fixed small $c_\nu$. This checks the
outer shift that enters the ordinary bracket of the full
response. The extra distance gain in a window derivative, and
the two differences of responses in a removed term, require
estimates for those distinct shifts on their full carriers.

For real physical data with all fine parameters one, the
same response satisfies the exact constraint and slice equations.
Its potential produces the critical configuration at the
prescribed datum up to the node factors, and its completed
current is the actual current of that configuration. Thus the
two endpoints of the ordinary interpolation are the physical
backgrounds, after the stated node normalization. At intermediate
multipliers the pair can be off the critical graph. The current
slot in the displayed formula is retained there.

\paragraph*{The upper frame on the full window box.}
Let \(\mathbf B^+\) be the upper tower of \(\square_0\) and
\(\mathbf B_j\) its depth-\(j\) truncation, with the same
convention on inward bonds. The bracket domain \(Y\) contains
\(\square_0\). The upper input is in the ordinary core at
\(((1+2\beta_{\mathrm s})A_0,(1+2\beta_{\mathrm s})A_1,A_0)\);
its hereditary clause supplies the tower of \(\square_0\)
and its truncations used in the comb construction.
The window uses the upper tower
configuration at its comb datum \(V^v\). Put
\(B_V=i^{-1}\log V\). The comb estimate gives
\(\|B_V\|_\infty\le CMA_0\) and
\(\mathsf k(v)\le e^{CMA_1}\) at every node, where
\(\mathsf k(g)=\|g\|\|g^{-1}\|\). Its prescribed extension
\(\tilde v\) has the same bound at every site.

Apply the flat-reference Landau pair to the continuation
\(\Phi_{\mathbf B^+,\mathbf1}(B_V)\) on a fixed datum disc
inside the continuation and response radii. The identity
holds on every inner bare bond of \(\square_0\):
\[
 \Phi_{\mathbf B^+,\mathbf1}(B_V)
     =\big(e^{i\eta_{k+1}\mathbf H_{k+1}(B_V)}\big)^{u_{B_V}} .
\]
The derivative bound in this theorem has one inverse local
length \(\xi_y^{-1}\) for the potential on the pillar at \(y\).
Integrate it on the line from \(0\) to \(B_V\). Each top is an
endpoint of a bounded number of determining rows, and the
exponential sum over tops is uniform. Consequently
\[
 \sup_y \xi_y
        \sup_{b\ {\rm over}\ y}|\mathbf H_{k+1}(B_V)(b)|
       \le C\|B_V\|_\infty .
\]
In the ascent-gauge lemma the form \(\mathbf H_{k+1}\) therefore
lies in its weighted small ball, by weak coupling. The
gauge factor of the Landau pair equals
\(h[\mathbf H_{k+1}]\) of that lemma. At real data the two factors
have the same prescribed values at all tops and both satisfy the
ascent identities, which determine their remaining values.
The equality continues holomorphically on the datum disc.
In particular \(\mathsf k(u_{B_V})<3/2\) on the whole box.
Thus the upper tower configuration equals
\(E^{g_\square}\), where
\[
 E=e^{i\eta_{k+1}\mathbf H_{k+1}(B_V)},\qquad
 g_\square=\tilde v\,u_{B_V},\qquad
 \sup_{\square_0}\mathsf k(g_\square)\le C .
\]
Only the local-length bound for the background potential is
used here. A uniform upper-level bound for that potential on
the finest layers is unnecessary.

Use the assigned ordinary decoupled response on its own
admitted active region, then restrict it to \(\square_0\)
using the full response reads. The fine parameters vanish
off the active cells. The pivot-completed datum \(B'\) is
supported at level \(k\), with \(\|B'\|_\infty\le C\delta_k\).
The constant-sequence bound above uses local lengths one.
On an admitted layered response chart the support estimate
on the fine-parameter discs gives
\[
 \xi_y\sup_{b\ {\rm over}\ y}|\mathbf H_k(P;s,B')(b)|
       \le C\delta_k e^{-\theta d_{\mathbf B}(y,\Sigma_k)},
 \qquad \theta=\mathring\delta_0/8 .
\]
Here \(\Sigma_k\) consists of the tops at which the datum is
supported. Their local lengths in level-\(k\) units are at
least \(L^{-1}\), by the geometry of the level-\(k\) rows.
If there is no such top the response is zero. The scale
comparison in the same separated sequence gives
\(\xi_{y'} / \xi_y\le L
e^{(\mathring\delta_0/16)d_{\mathbf B}(y,y')}\).
Take \(y'\) in \(\Sigma_k\) at minimum distance. The loss from
dividing by \(\xi_y\) is therefore paid by the support
exponential. This proves
\(\sup|\mathbf H_k(P;s,B')|\le C\delta_k\) in level-\(k\) units
also on the fine layers of the box. The proof uses the support
estimate for the whole datum, without summing the responses
over its nonzero sites.

For an assigned shift in the upper tower frame, transport
it by \(\operatorname{Ad}(g_\square^{-1})\) at the initial
site of each bond. Its bound remains \(C\delta_k\) in
the exponential frame on the whole box. A shift supplied
directly in the exponential frame uses the same bound. The
window condition \(\|\mathbf A\|_\infty<c_2A_0\) then has
fixed room, since \(\delta_k/A_0\to0\). This comparison
concerns the two representatives of the upper tower. The
fluctuation originating at the input pair uses the relative
frame below.

\paragraph*{The rows in this frame.}
For \(c\in\Gamma_q\) in the truncated tower, every bare bond
of its full two-block carrier lies in \(\Omega_{q-1}\) of
the upper tower for \(q\ge1\). The local lengths there are
at least \(L^{q-1-(k+1)}\) in upper-level units. This is the
carrier assertion of the tower lemma, applied to the
truncated sequence, whose \(\Omega_{q-1}\) is the same as
in the upper tower. It yields
\[
 L^q \sup_{b\ {\rm in\ carrier}(c)}
       |\eta_{k+1}\mathbf H_{k+1}(b)|\le CLMA_0 .
\]
For \(q=0\) the corresponding estimate follows directly
from the local-length bound at a bare bond.
If \(\mathbf A=(t+\nu\zeta_\square)\mathbf H_k^{\rm win}\),
then \(L^q\sup|\eta_k\mathbf A|\le
C L^{q-k}(1+|\nu|\|\zeta_\square\|_\infty)\delta_k\).
In particular all the row discs stay within the absolute
average radius on \(|\nu|<c_\nu/s_k\) for small fixed \(c_\nu\).
On this disc the transported shift also satisfies
\(\|\mathbf A\|_\infty<c_2A_0\) on the whole box.

These checks identify the row of the actual shifted upper
tower with the row of \(e^{i\eta_k\mathbf A}E\), transformed at
its two endpoints by \(g_\square\). To check the regular branch,
all intermediate row carriers are inside the same two
centered blocks. The row estimate keeps their values close
to the identity in the exponential frame. Hence every
relative contour matrix and every argument of its local
mean lies in the prescribed logarithm chart. Conjugating the
relative matrices multiplies their distances from the
identity by at most the fixed condition number of \(g_\square\).
The same charts remain valid in the actual frame. In fact
for a relative contour diameter \(d\), smallness ensures
\(\mathsf k(g_\square)d<\omega_{\rm c}\) and
\(4\mathsf k(g_\square)d<\omega'_{\rm c}\). The transported
mean therefore lies in the ball that characterizes the
unique normalized mean. The relative logarithm is transported
by the same conjugation.
Applying this at each intermediate scale gives the claimed
endpoint rule on the full carrier, including the outer
determining rows. Thus no outer-row datum correction is needed
for the unamplified path of this upper box construction.

On the exponential chart write \(B(t,\nu)\) for the family
of logarithms of these transported row values. The old
continuation uses \(\tau=1\) in its datum \(\tau B+B_{\rm err}\).
Use \(B(t,0)\) and add \(B(t,\nu)-B(t,0)\) in the correction
slot, which is zero at \(\nu=0\). The bare bond logarithm
changes by at most \(C|\nu|\eta_k\delta_k\),
so the order-\(q\) row estimate gives a datum change at most
\(C|\nu|\delta_k L^{q-k}\).
Taking the supremum over all \(q\le j\) bounds the correction
by \(C|\nu|\delta_k L^{j-k}\). Decrease \(c_\nu\) so that
this is less than half the old correction radius \(c_3a_0\)
on the common disc, using
\(\alpha_k L^{j-k}/a_0\le C\) uniformly in \(j\le k\).
On an addressed old domain in \(\square^{\sim2}\) the
window theorem's normalized core readouts and the two
hereditary tower readouts are holomorphic and bounded on
that correction disc. Their Cauchy estimate gives
\eqref{eq:local-formation-WF5}.
The independent input current is used by the decoupled
response; the output current of the inner box is its
constrained current.

At a physical real critical endpoint use the nested restriction
on the constant-layer box, or on its truncated sequence. The
box rows determine every ambient row affected by extension
by the unchanged exterior. Any bond inside the box in a
crossing plaquette is a fixed row of level zero. Hence such
an extension preserves the ambient constraints and changes
only the action of the inner plaquettes. The restricted
field is critical, and regular uniqueness identifies its
orbit with the box orbit over its own datum. This applies
at the upper and old critical endpoints with their respective
towers. At other points the box reading is the prescribed
continuation from the row family. The two removed terms use
their assigned endpoints, response differences and core estimates.

\paragraph*{The lower rows at the edge of the window box.}
Write \(D_q(c)\) for the bare read set of the order-\(q\) average
at a row \(c\in\Gamma_q\). It is contained in the two centered
blocks of side \(L^q\) with centers at the endpoints of \(c\).
Indeed the one-step contour averages, pivot and two block
transports use only the double block. At the next finer level
each site of those blocks is replaced by its centered block of
children. Iterating gives the stated read set. In particular
the mask is determined by the actual row, including those
on the lower layers of the truncated tower. The second-kind
holonomies are read from these rows by the tower rule.

Fix the background input, the node frame, the fine parameters
and \(t\) as \(\nu\) varies. A row
whose \(D_q(c)\) misses the support of the bond direction
\(\zeta_\square\mathbf H_k^{\rm win}\) is constant in \(\nu\).

\paragraph*{The frame used for the fluctuation.}\hypertarget{wf:relative-frame-restoration}{}
Let \(E\) be the chosen representative of the upper box
background and let \(\mathbf U\) be the configuration of
the input pair. The relative smoothing construction
relates these two configurations. The single-lattice
window construction in \cite[(3.24)--(3.29)]{B87}
instead evaluates the response at the upper box field
extended by \(\mathbf U\) outside and carries the
independent current. In that extension use the upper
tower in the representative of its prescribed datum.
The separation of its first two members leaves a
bare rim inside the box. Each bare bond there is a
determining row at level zero and hence already has
the value of \(\mathbf U\). Therefore a plaquette
crossing the box face uses exactly the original
links, including its links on the inner side. The
same check applies to a stencil of fixed bare reach
inside this rim. If one uses the real tower of the
input real factor for the extension's real factor,
the two also agree on the rim and the relative bond
logarithms there are the original ones.

%
%
The same splice preserves every order-\(q\) row whose
full double block crosses the box face. More generally,
write \(W_i\) for the members of the inserted tower,
with \(W_0\) the box and \(W_i=\varnothing\) beyond
its depth. Consider rows with their full double
blocks in the ambient first member. The average of
the splice at a row \(c\)
is the original one if \(c\) is retained in the inserted
determining set, or if neither endpoint of \(c\) lies
in \(W_{q+1}\). For a retained row this follows from
the prescribed datum and its full carrier in the box.
For the second case, argue by induction over \(q\).
At level zero a bond that is not inner to the box
is unchanged, and an inner bond of this kind is
fixed by the bare datum. At positive level, a row with an endpoint in
\(W_q\) is an inner retained row, since its endpoints
avoid \(W_{q+1}\) and the layer gap fits its double
block inside \(W_{q-1}\). If both endpoints are outside
\(W_q\), their order-\(q\) blocks are outside \(W_q\)
by alignment. Each input row in the one-step average
then has no endpoint in \(W_q\), so its value at the
preceding level is unchanged. This proves equality on
the prescribed nonlinear branch. The same induction
works on the admitted complex row charts, because each
unchanged read has the same full input. If a double
block crosses the box face, an endpoint in \(W_{q+1}\)
would put that block inside the box by the layer gap.
This proves the asserted crossing-row equality.

For an auxiliary hereditary tower use the criterion on
every row of its actual truncated determining set,
including the rows at its top. If it holds, the full
real and complex input data for that tower agree with
the original data, so its chosen continuation and its
current read are unchanged. This covers, for example,
a containing tower of the same depth and the same
monotone rule. For an arbitrary straddling tower a
lower determining row can be wholly inside an active
member of the inserted tower and use a changed average.
Such hereditary reads still require admission and
estimates. For the relative-frame rule write
\[
 w_0(x)=E(\Gamma_{y_0,x})^{-1}a_\square
                         \mathbf U(\Gamma_{y_0,x}),\qquad
 w=(\mathcal K_{E,\mathbf U}+1)^{-1}
                  [(\theta_\square+\Delta\theta_\square)w_0].
\]
Here the anchor \(a_\square\) fixes the frame at the base
point, and \(\mathcal K_{E,\mathbf U}\) is the mixed
covariant operator on the frame. The inner fluctuation
uses \(\operatorname{Ad}(w)\) on the support where the
frame is invertible. The regularity estimate for this
frame concerns the input connection expressed in the
box coordinates. It is this estimate, together with the
covariant response bounds, that is needed for the ordinary
first-difference H\"older estimate. The comparison of the
two representatives of the upper tower above uses the
bound on their condition number.

At a real critical reference in the nested geometry,
restriction to the box is
critical at its own full tower datum. By regular uniqueness,
\(E=\mathbf U^g\) there for a real gauge \(g\), with anchor
\(a_\square=g(y_0)\). Transport along any comb path then
gives \(w_0(x)=g(x)\). For the mixed connection
\(E(x,y)g(y)\mathbf U(x,y)^{-1}=g(x)\).
It follows that
\((\mathcal K_{E,\mathbf U}+1)(\theta_\square g)
 =(\Delta\theta_\square+\theta_\square)g\)
with the prescribed cutoff boundary values.
Uniqueness of the smoothing equation gives
\(w=\theta_\square g\). Hence its adjoint action is the
one of the physical gauge wherever \(\theta_\square\ne0\).
This comparison applies to the whole real path of shifts
from that critical reference on those carriers, since the
frame is fixed before varying the fluctuation direction.

If a determining row \(c\) of level \(q\) has its complete
backward carrier in that chart, its changed datum is
\[
 \mathcal M^q(\mathbf U^g)(c)
        =g(c_-)\mathcal M^q(\mathbf U)(c)g(c_+)^{-1}.
\]
One must use both node factors at every such row,
including the lower rows. The old constrained
configuration at the transformed full datum is read in its
corresponding orbit. Holomorphic continuation of this
identity uses normalized frame representatives and the
regular logarithm charts on the complete carrier. In
particular the real quaternionic normalization of \(w\)
may be replaced in a complex chart by a local analytic
determinant normalization. The adjoint action itself
only requires invertibility. The endpoint identity does
not supply invertibility where the cutoff is zero.

%
For complex inputs, keep the operator on the same
box and with the same boundary condition as at the
real reference. Put \(\mathcal K_0=\mathcal K_{E_0,\mathbf U_0}\),
\(G_0=(\mathcal K_0+1)^{-1}\),
\(\Delta\mathcal K=\mathcal K_{E,\mathbf U}-\mathcal K_0\),
and \(F=(\theta_\square+\Delta\theta_\square)w_0\).
Subtract the equation for \(\theta_\square g\) at
the real reference. The identity is
\[
 w-\theta_\square g
   =(I+G_0\Delta\mathcal K)^{-1}G_0
       \bigl[F-F_0-\Delta\mathcal K(\theta_\square g)\bigr].
\]
To obtain a uniform complex frame from this identity,
one can bound \(G_0\Delta\mathcal K\) by \(1/2\)
in the sup operator norm on the box. On a carrier
with \(\theta_\square\ge\theta_*>0\), a sup bound
\(\theta_*/(4\sup\|g^{-1}\|)\) over the whole box for
\(G_0[F-F_0-\Delta\mathcal K(\theta_\square g)]\)
gives \(\|(\theta_\square g)^{-1}(w-\theta_\square g)\|
\le1/2\) there. The inverses in the formula and on
the carrier then follow by convergent Neumann series,
holomorphically and with a uniform condition number
on that carrier. We must prove the operator and
residual estimates on the admitted complex input
sections, with a matched reference in the perturbation
range and bounds uniform in the lattice spacing.
The first-difference H\"older estimate also uses
the relative-frame regularity in the box connection.
Near a vanishing cutoff the assigned outer reads
are part of the localization rule.

%
%
\paragraph*{The cutoff on the baseline field.}
In the single-lattice formula both the baseline and
the extra window direction are cut off. To compare
the baseline with the one in
\eqref{eq:local-formation-WF3}, let \(\chi_\square\)
be one on \(\square^{\sim3}\) and supported inside
the box. With the same fluctuation field in the
same box frame, join the two baseline choices by
\[
 \mathbf A_v=
 \bigl(t[\chi_\square+v(1-\chi_\square)]
                       +\nu\zeta_\square\bigr)
                 \mathbf H_k^{\rm win},
 \qquad 0\le v\le1 .
\]
With fixed room in the sup bound for the window
field, this path stays in the inner substitution
disc. On a level-\(q\) row the derivative of the
plain logarithmic datum in \(v\) has size at most
\(C_R\delta_k L^{q-k}\). It vanishes whenever the
complete double block lies in the plateau of
\(\chi_\square\), since there the bare product is
independent of \(v\). For an inner read contained in
\(\square^{\sim2}\), all remaining source rows are
separated from the read by at least \(cM L^{k-j}\)
in the old weighted distance, after reserving the
fixed stencil reach. Indeed the distance through
the physical margin is at least \(cML^k\) in bare
units and the old tower uses spacings at most \(L^j\).
Integrating the weighted source estimate along
the path gives the old potential and its first
differences an error bounded by
\[
 C_R\delta_k L^{j-k}\exp(-cM L^{k-j})
\]
in old units on the inner read. For the constrained
current in this same exponential representative,
first pass the plain row difference through the
local naturalizer. Its stencil enlarges the source
support by a fixed reach and preserves the distance
bound with a smaller constant. At the flat reference
the current of the reconstructed field is
\(\mathsf f(-i\xi\operatorname{ad}_{\mathbf H})^{-1}
\mathbf F^\flat\) at full old switches. The reduced
first variation \(\mathbf F^\flat\) has the same
weighted source estimate as \(\mathbf H\). Taking
the difference of this product, the inverse acts
bond by bond and its difference is bounded by the
local difference of \(\mathbf H\). This gives the
same bound for the normalized constrained current
on the inner read. This argument treats the
baseline-cutoff change inside the fixed old box.
The window identity retains the removed difference
and uses the assigned domain and box pair on each
of its other segments.

\paragraph*{The change of localization and its sum rule.}
For an old member \(F\), denote its original domain response
along the source interpolation by \(P^{\rm dom}(t,u)\).
For a near window the source's two removed segments go
from \(P^{\rm dom}\) to its prescribed intermediate pair
and from that pair to the frame-dependent inner pair
\(P^{\rm in}\). Calling the intermediate pair \(P^{\rm mid}\),
at fixed path parameter the restoring identity is
\[
 \partial_u F(P^{\rm dom})=\partial_u F(P^{\rm in})
  +\partial_u\bigl[F(P^{\rm dom})-F(P^{\rm mid})\bigr]
  +\partial_u\bigl[F(P^{\rm mid})-F(P^{\rm in})\bigr].
\]
Use the fundamental theorem of calculus on each removed
segment before differentiating at zero. The assigned current
is carried as a component of the pair throughout the
interpolation. Each removed direction uses its own full
carrier and complex disc.

For the common pointed sum, apply this identity to each
member and to its own tent, then add before the fine
bracket. For the resolver on the prescribed box, include
the complete box reads of the frame and the old inner
carriers in the base. This uses a fixed number of current
cubes. The locality estimate uses this base for the
entire relative-frame construction.
For the source's local fields, restriction to the wall
component of the base gives the same field on that
component. Once the whole frame read set is in the
base, the hereditary restriction in the source rule
therefore gives the same inner function on the smaller
ambient domain. The finite resummation of the brackets
recovers the physical sum, and the two differences
restore its original domain response. This argument
permits a common-base grouping of the complete localization
and preserves its summed physical action difference.

For \eqref{eq:local-formation-WF3} we must use the smoothed
frame on the full carriers where it is admitted and the
assigned outer-layer readings elsewhere. The cutoff can
vanish on outer carriers. The frame, old endpoint data and
both removed directions must therefore come from one
localization rule on the bracket's actual complex domains.
The same choice is required at both depths in the comparison.

It follows from \eqref{eq:local-formation-WF2}--\eqref{eq:local-formation-WF5} and the fixed room that there is a
constant \(c_\nu>0\), independent of \(j,k\), such that the
admitted inner substitutions have the common disc
\begin{equation}
 |\nu|<{\cal R}_\nu=c_\nu/s_k .
 \label{eq:local-formation-WF6}
\end{equation}
Indeed choose \(c_\nu\) so that
\(C_Rc_\nu C_{\rm age}\) is smaller than half of every
fixed room.  \eqref{eq:local-formation-WF5} then keeps each substitution inside
its assigned domain.  This argument verifies an inclusion
using the actual input-domain estimates; it does not extend
functions to unions of complex gauge orbits.
We take \(s_k<c_\nu/2\), so that \({\cal R}_\nu>2\).

Third, the critical box-background maps in \eqref{eq:local-formation-WF4} obey the
weighted first-response, H\"older and compatible boundary-comparison
estimates used in the nonlinear pointed window theorem.
The complete fields on \eqref{eq:local-formation-WF6}, including the unperturbed field,
have output size at most \(C\alpha_k\).
The central axial datum is bounded by \(C\alpha_k\sigma_j\).
Its half-datum-ball condition holds when \(C\rho_j\sigma_j\)
is sufficiently small, which fails at most at a fixed finite
set of young ages by \eqref{eq:local-formation-WF2}.
The remote datum has only an absolute \(C_M\alpha_k\) bound.
Use the fixed continuation radius and the once-for-all
ceiling \eqref{pw:absolute-datum-ceiling} to put the whole
axial datum in its half ball. This ceiling is imposed on the
upper input, the transported fluctuation and any raw correction
throughout the discs of \eqref{eq:local-formation-WF6}, with
the room already required in \eqref{eq:local-formation-WF5}.
The window stock satisfies it by reducing \(a^\flat<a_I(M)\)
after fixing \(M\); no bound on the number of old levels is used.
At the exceptional young ages the actual core memberships and
the whole-core bound for a member's own tent action remain available.
Finite boxes are treated first; their linearized boundary
comparison is uniform in the bare mesh.  Periodic or distant
tails are retained with their stated exponential lengths.

Finally, assume the actual window identities and elementary
localization geometry.  The partition \(\sum_\square
\zeta_\square=1\) is inserted along the direction of
differentiation.  For complete physical pairs the old action
difference is the sum over windows of
\(\int_0^1\partial_\nu|_0\) of its window evaluations \eqref{eq:local-formation-WF3}--\eqref{eq:local-formation-WF4},
plus the prescribed outer and removed terms.
The same local formulas define their extensions at every local
complex core.  Equality with a global physical fibre integral
is asserted only at complete physical pairs.

An inner window contains at most \(C_M\sigma_j^{-4}\) old base
points.  At a fixed base point the number of old domains of
tree length between \(m\) and \(m+1\) is at most \(C_De^{\kappa_Dm}\).
Diameters are bounded by \(C_M(1+m)\), and the input decay rate
has a fixed surplus over \(\kappa_D\).
Thus every fixed polynomial diameter factor is summable.
An outer field has the additional decay from its distance to
the window, while a removed field has the factor
\(e^{-cM/\sigma_j}\).
On the assigned single-lattice charts the bounds in the
member's local pair norm on the complete \(X^\sim\) are
\begin{equation}
 C_R\delta_k\sigma_j
 \begin{cases}
  e^{-\delta_F\mathsf d(X,\square)},&\text{outer cutoff direction},\\
  (1+|w|)e^{-\delta_F M/\sigma_j},&\text{each removed segment direction}.
 \end{cases}
 \label{eq:local-formation-WF7}
\end{equation}
Here \(w\) is the removed term's direction parameter with
\(|w|\delta_k\) small. For the directional estimates we use
admission in the old core on the respective parameter discs,
with a fixed fraction of its radius available along each
segment.
Long input domains retain their original tree decay: one
fraction supplies \(e^{-c_\kappa/\sigma_j}\), and the remaining
fraction supplies the declared output spatial rate.
Here \(c_\kappa>0\) comes from tree length in \(M\)-cubes and
the allocated input rate.  When these long-domain tails are
written together with removed-response tails as
\(e^{-cM/\sigma_j}\), the fixed \(c\) is allowed to depend on
\(M\) and the input rate; it is not asserted independent of them.
These are geometric and response hypotheses, not bounds on
the value of an old member.

\paragraph*{Renormalize the critical inner function first}

Let an old regular member have norm \(g_j\), be gauge invariant
and Euclidean as a pointed total, and have its own coefficient
\(b_j\).  At every base point remove its vacuum value and
\(b_jA(h_z,\cdot)\) before bracket interpolation.
The pointed Hessian theorem fixes this coefficient by the
unit normalization of the same tent action.
Site-frozen members may differ at different base points: the
argument uses each member's own hypothetical Euclidean family
and its own coefficient.  No translation invariance of that
array is required.

For \(j<k\), here is the quantitative window estimate on \eqref{eq:local-formation-WF4}.
The axial field near a base point has linear constant-curvature
part of size \(C\alpha_k\sigma_j^2\).
The H\"older and nonlinear response errors have sizes
\(C\alpha_k\sigma_j^{2+\alpha_H}\) and
\(C r_{B,j}^{-1}\alpha_k^2\sigma_j^4\), respectively.
Cauchy's second and third derivatives of a member cost
\(Cg_j/\alpha_j^2\) and \(Cg_j/\alpha_j^3\).
The linear derivative vanishes by constant-color invariance.
The quadratic value on the affine part is canceled on the
whole pointed total at that base point by the member's own
tent action.  No coefficient is assigned to an individual
domain \(X\).  The individual Taylor remainders have a fixed
polynomial in \(D=2+\operatorname{diam}_jX\) times their
original spatial decay; the domain count above sums these
polynomials.  Multiplying the remaining errors therefore
gives, at one base point after this domain sum,
\begin{equation}
 Cg_j\left[
       \rho_j^2\sigma_j^{4+\alpha_H}
                     +\rho_j^3\sigma_j^6\right].
 \label{eq:local-formation-WF8}
\end{equation}
This is the local Taylor argument on the constrained critical
box pair, with its uniform weighted response sums.
When its local Cauchy-size parameter is large, the whole-core
bound and the own-action bound give the same estimate after
increasing that summable polynomial, since \(1+T^2\le CT^3\)
for \(T\ge1\).
Any fixed finite set of young ages failing the central
half-datum-ball test is handled by the same whole-core bounds;
their scale ratios have a fixed positive lower bound. The
remote absolute-domain condition is retained separately.
The resulting constant is independent of the number of old
levels.

Now sum the \(C_M\sigma_j^{-4}\) base points.
Because \(\rho_j\ge1\), \(0<\alpha_H\le1\), and
\(\sigma_j\le1\), the renormalized inner window function has
bound
\begin{equation}
 Cg_j\rho_j^3\sigma_j^{\alpha_H}.
 \label{eq:local-formation-WF9}
\end{equation}
Truncating the pointed total at the enlarged window contributes
only the long-domain tail already set aside in the geometry.
It does not replace the pointed total by an unnormalized
individual term.

A defect or an aged remainder has no coupling counterterm.
After its value at the identity is removed, constant-color
invariance still eliminates the linear derivative, provided
there are no fixed color spectators.  The vacuum-only window
argument gives at one base point, after the same polynomial
domain sum, \(Cd_j\rho_j^2\sigma_j^4\), where \(d_j\) is that origin's norm.
The large local-size case uses \(1\le CT^2\) for \(T\ge1\).
After the base-point sum its inner bound is
\begin{equation}
 Cd_j\rho_j^2 .
 \label{eq:local-formation-WF10}
\end{equation}
Stored-field boundary origins are not included in this
statement.  Their spectators and their hull geometry need the
separate boundary formation estimate.

At \(j=k\) no age improvement is needed.
The vacuum-normalized old member together with its own action
has whole-core norm at most \(Cg_k\), by the coefficient bound
\(|b_k|\le Cg_k/\alpha_k^2\) and the core bound on its tent action.
The actual fluctuation change has size \(C\delta_k\) within
a core with fixed analytic room \(c\alpha_k\).
Cauchy's inequality therefore gives \(Cs_kg_k\), with its
original spatial decay.  There is no fine-volume age factor
at this level.  This is the term with weight \(\omega_1\);
no cancellation between different current-level domains is
used for it.

\paragraph*{The one small factor and the outer terms}

Fix \(t\in[0,1]\) and the fine parameters.
Subtract the window function's value at \(\nu=0\).
It vanishes there and is holomorphic on \eqref{eq:local-formation-WF6}.
The bounds in \eqref{eq:local-formation-WF9}--\eqref{eq:local-formation-WF10} hold on that whole disc.
Schwarz's lemma, or Cauchy's formula for its derivative at zero,
therefore supplies precisely one factor \(1/{\cal R}_\nu\le Cs_k\).
The actual window formula uses that derivative and its integral
over \(t\); the interval has length one.
Thus an old regular inner term contributes
\(Cs_kg_j\rho_j^3\sigma_j^{\alpha_H}\), and a vacuum-only
origin contributes \(Cs_kd_j\rho_j^2\).
No second amplification factor is extracted later.

\paragraph*{The outer direction and the two removed segments.}
Use the potential and the reduced current together in the
norm for the old member. The unamplified domain field has
size \(C\sigma_j\delta_k\), and the directions obey
\eqref{eq:local-formation-WF7}. These core estimates can omit
the H\"older quotient. Use the current specified by the
interpolation along each segment.

For the outer direction, at an admitted unperturbed point
with fixed old core margin, the normalized shift stays in
that margin on an affine multiplier disc of radius
\(ca_0(\sigma_j\delta_k)^{-1}
                 e^{\delta_F\mathsf d(X,\square)}\).
Cauchy's estimate therefore gives
the outer smallness and the distance factor together. If the
window is near \(X\) but \(X\) leaves \(\square^{\sim2}\), use
the spare old-domain length factor instead. For a removed
segment use \(|w|\le1\). On the admitted core the difference
between its endpoint values of a member of size \(g_j\) is
bounded by
\[
 C g_j e^{-\kappa d_j(X)}
       \frac{\sigma_j\delta_k}{a_0}
                         e^{-\delta_F M/\sigma_j}.
\]
This follows by integrating the old-field derivative along
that segment, with the fixed fraction of the old radius
available for Cauchy's inequality. A Cauchy estimate in \(w\)
on the circle \(|w|=1/2\) gives the same bound for the
directional derivative at zero.
The two endpoint differences telescope at real physical
arguments, with the exact domain and box endpoints of the
window identity. The fine-cube brackets use these bounds on
the entire parameter discs. This argument needs the bounds on
the assigned segment directions and endpoints. The
unamplified determining rows use the preceding comparison.
After the fine-site count the bound is
\begin{equation}
 Cs_kg_j\rho_j\sigma_j^{-3}e^{-cM/\sigma_j}.
 \label{eq:local-formation-WF11}
\end{equation}
The same reasoning applies to a vacuum-only origin with \(g_j\)
replaced by \(d_j\).
For \(0<\sigma\le1\) and \(\rho\ge1\),
\begin{equation}
 \rho\sigma^{-3}e^{-cM/\sigma}
       \le C\rho^3\sigma^{\alpha_H},
 \qquad
 \rho\sigma^{-3}e^{-cM/\sigma}\le C\rho^2.
 \label{eq:local-formation-WF12}
\end{equation}
Indeed the suprema of
\(\sigma^{-3-\alpha_H}e^{-cM/\sigma}\) and
\(\sigma^{-3}e^{-cM/\sigma}\) are finite.
Any fixed additional polynomial from spatial summation is
absorbed in the same way, by decreasing the exponential rate
once.  Long-domain tails are treated with their retained tree
decay and the same elementary inequality.
This proves the ordinary outer and removed bounds without
applying the pointed cancellation to an arbitrary pair
\((U,J)\).

% Earlier Part II, Section 5.5, 43719--43828, 44552--44574,
% 45281ff., 45440--45554 and 45629--45642.
\paragraph*{The sum of the removals on a capped window.}
For the native near unit whose own padded box fits the
box-read class, one can estimate the two removed terms
together. This class has \(\nu\ge3\) and its axial-datum
majorant is strictly smaller than the old box radius
\(\alpha_{3,j}\). Use the same capped determining
set at both endpoints. At the inner endpoint use the
averages of the cutoff layered field. At the other
endpoint splice the inner prescribed row datum with
the rows read from the domain response. In the splice
the exact nonlinear relative average \(Q_\star\) uses
both endpoint factors. The cutoff mismatch vanishes
on each full double block in the plateau. For the
splice itself, the capped and old determining sets
agree on the complete row reads wherever the splice
weight is nonzero. The old response has its prescribed
row datum there, so the two relative logarithms
being spliced are equal in the same node frame.
Outside the splice the datum is read from the response
itself. Thus on
the complete capped set the second endpoint is the
box restriction of the domain response over its own
datum. This identification uses all rows of the
splice, with the old current read in the same chart.

For a constituent \(T\) of the old member the two
endpoint differences telescope to
\(G=T(P^1)-T(P^0)\), where \(P^0\) is the
inner endpoint and \(P^1\) the spliced endpoint.
Take the comb logarithms \(B^0,B^1\) in the common
old box. The cap estimate gives
\(\|B^0\|_\infty<1.42\alpha_{3,j}/L\) and
\[
 \|B^1-B^0\|_\infty
   \le \tfrac12\alpha_{3,j}e^{-\delta D_v/8}.
\]
Here \(D_v\) is the old distance of the box read to
the cutoff change, and the native near-window geometry
gives \(D_v\ge ML^{n_0-j}/5\), with \(n_0\) the window
level. This is a multiple of \(M/\sigma_j\) when the
window levels agree. On the admitted old
section use \(B^0+\lambda(B^1-B^0)\) for
\(|\lambda|<e^{\delta D_v/8}\). If \(T\) has
majorant \(n_v\) there, its difference from the
value at \(\lambda=0\) is bounded by \(2n_v\)
and vanishes there. Schwarz's estimate at one
gives \(|G|\le2n_v e^{-\delta D_v/8}\).
On the native removed-parameter disc of radius
\(r_\square^{\rm rem}\) this holds for the whole
holomorphic difference. Cauchy's inequality gives
\[
 |\partial_u G(0)|
 \le 2n_v (r_\square^{\rm rem})^{-1} e^{-\delta D_v/8}.
\]
The same bound applies after integration of the real
interpolation parameter. This is a single-run estimate
of the combined jet. To use it for the actual brackets,
match their two endpoint jets with these capped reads
on one admitted old-core section and use their actual
removed-parameter radius.

\paragraph*{Fine brackets and the resulting formation estimate}

We now make the base of an inner bracket explicit.
For \(j<k\) and one old base point \(z\), let
\(\square=\square_k(z)\) be its assigned current window and
\(R_z=X'_k(z)\) its fixed enlarged box.
There are only a fixed number of direction-partition windows
near this base window.  Group their directions by linearity
before applying the box-background replacement, and choose
\(R_z\) to contain their fixed enlargements.  In \eqref{eq:local-formation-WF3} their
sum is the inner cutoff \(\zeta_\square\); the remaining
directions are outer terms.  This is part of the fixed
window-first rule.  The derivative at zero is linear in
these directions, so this grouping preserves the partition
identity without changing any incoming pointed member.
Let \(\mathsf A_{j,z}(X,\cdot)\) be the fixed action-family
filing with total \(A(h_z^{(j)},\cdot)\), and put
\(\widetilde G_{j,z}=G_{j,z}-b_{j,z}\mathsf A_{j,z}\).
The one inner window member is the function
\begin{equation}
 \begin{split}
 {\cal W}_{j,z}(P')=
 \sum_{\substack{X\ni z\\X\subset\square^{\sim2}}}
 \big[&
 \widetilde G_{j,z}\big(X,P^{\rm box}_{j,z}(P'),z\big)\\
 &-\widetilde G_{j,z}\big(X,(1,0),z\big)\big].
 \end{split}\label{eq:local-formation-WF13a}
\end{equation}
Here \(P'\) is a local potential/current pair of the current
resolved level \(k\), on \(R_z\).
The map \(P^{\rm box}_{j,z}\) first forms its actual nonlinear
averaged datum and then takes the old-level constrained
critical box pair of that datum.
The same box-background map is used for every \(X\) in this
sum.  The own tent action is included in the near sum; its
fixed support lies inside the assigned enlarged window.
The declared domain of this entire member is \(R_z\), and its
fine-cube base is
\(B_z=\operatorname{Hull}_{\pi_k}R_z\).
Its number of current fine cubes is bounded by a fixed
constant \(C_{\rm win}\), independent of \(j,k,z\).
This is a presentation of an inner member before localization;
we do not apply different old-\(X\) bracket operators to its
summands.

Let \({\cal W}_{j,z,Y}(s,t,\nu)\) mean this one function
evaluated by the assigned local response and the datum in \eqref{eq:local-formation-WF3} on \(Y\supset B_z\).
The inner potential contributed by this member is
\begin{equation}
 V^{\rm in}_{j,z}(Y)=
 \int_0^1
 \left[
   \left.\partial_\nu
       {\cal W}_{j,z,Y}(\,\cdot\,,t,\nu)\right|_{\nu=0}
 \right][\,Y\setminus B_z\mid B_z\,]\,dt .
 \label{eq:local-formation-WF13b}
\end{equation}
The brackets in this formula are linear in the whole function in \eqref{eq:local-formation-WF13a}.  Thus its cancellation and bound in \eqref{eq:local-formation-WF8} are established
before this single bracket is estimated.
Different \(z\)'s need not have a common base: there is no
cancellation between them, and their bounds can be added.
Hereditary restriction identifies the bracket on \(Y\) with
the bracket of the same ambient interpolation function whose
parameters outside \(Y\) are zero.  Finite fundamental theorem
of calculus therefore gives, at a complete physical pair,
\[
 \sum_{Y\supset B_z}V^{\rm in}_{j,z}(Y;p_W)
   =\int_0^1
       \left.\partial_\nu{\cal W}_{j,z,\Lambda}
                             (1,t,\nu;p_W)\right|_{\nu=0}\,dt .
\]
Every local term uses only the pair on \(Y\) and the input
members on its declared domain \(B_z\subset Y\).
This proves chart-coherent locality and the physical sum
identity for this fixed window-first rule.
Summing the window identities and the retained outer and
removed origins recovers the old action difference.
No equality of individual output domains with a different
filing is required.
The current-level origins \(j=k\), and the far and removed
origins, retain their own declared bases and the bounds
already proved.  For vacuum-only near origins the same
window-first grouping may be used without a marginal
subtraction.

The assigned maps are hereditary on active wall-components,
and complete at all parameters equal to one.
For an interpolation function \(\Psi\), its bracket is
\begin{equation}
 \Psi[K|N]=\int_{[0,1]^K}
        \prod_{\Delta\in K}\partial_{s_\Delta}
                            \Psi(s_K+1_N)\,ds_K .
 \label{eq:local-formation-WF13}
\end{equation}
For \eqref{eq:local-formation-WF13b}, \(K=Y\setminus B_z\), \(N=B_z\).
For an origin retaining its own base \(X_0\), use
\(K=Y\setminus X_0\), \(N=X_0\).
A component carrying a differenced cube and missing the base
makes this bracket zero.
Each coordinate has the disc \(R=e^{\bar\kappa}\), so mixed
Cauchy costs \((R-1)^{-|K|}
\le e^{-(\bar\kappa-1)|K|}\), with no factorial.
The common \(\nu\)-disc and its one derivative commute with
these brackets by joint holomorphy.
The preceding amplitude bounds therefore retain this additional
spatial factor.

Here is a quantitative rate allocation.
For a retained old domain \(X\) and its current fine hull
\(X_0\), put \(m=\#(Y\setminus X_0)\).
The actual cube geometry and rooting give a fixed constant
\(C_h\) such that
\begin{equation}
 d_k(Y)\le C_h\big[d_j(X)+m+1\big].
 \label{eq:local-formation-WF13c}
\end{equation}
To see this, view a tree covering \(X\) on the old lattice
in current units and cover it by current fine cubes.
Its hull has a covering tree of length at most
\(C(d_j(X)+1)\).  Each component of the additional cubes
meets that hull, or its bracket vanishes.
A rooted spanning forest adds at most a fixed number of
edges per additional cube.  This proves \eqref{eq:local-formation-WF13c}.
For an inner window member the same argument and
\(\#B_z\le C_{\rm win}\) give
\(d_k(Y)\le C_h(m+1)\), increasing \(C_h\) by a fixed
factor if necessary.

Split the old input rate into positive parts
\(\kappa_{\rm sum}+\kappa_{\rm far}+\kappa_{\rm ret}\le\kappa\).
Choose \(\kappa_{\rm sum}>\kappa_D\) with enough strict
gap to sum the fixed polynomial diameter factors, and
use \(\kappa_{\rm far}\) for the long-domain tail in
\eqref{eq:local-formation-WF7}--\eqref{eq:local-formation-WF12}.
The retained old-domain factor is
\(e^{-\kappa_{\rm ret}d_j(X)}\).
Put \(\beta=\bar\kappa-1\).
Choose an output rate \(\lambda\) and a number \(\epsilon>0\)
with
\begin{equation}
 C_h(\lambda+\epsilon)
                <\min\{\kappa_{\rm ret},\beta\}.
 \label{eq:local-formation-WF13d}
\end{equation}
All rates are chosen once, sufficiently large for the
subsequent ordinary polymer criterion.
Combining the retained old factor with \(e^{-\beta m}\)
and \eqref{eq:local-formation-WF13c} gives
\(C e^{-(\lambda+\epsilon)d_k(Y)}\).
For inner members the bounded base needs only the
\(\beta\) part of this inequality.
The number of possible current-window anchors at fixed
\(Y\) is polynomial in \(\#Y\); the old bases within one
such window have already been counted by
\(C_M\sigma_j^{-4}\).
The bound \(\#Y\le C(d_k(Y)+1)\) and the remaining
\(\epsilon\) absorb this polynomial.
This is a pointwise output-domain estimate; any later
sum over output domains spends its separately declared
entropy reserve.

The function
\(\rho_j^3\sigma_j^{\alpha_H}\) is bounded by
\begin{equation}
 \omega_{k+1-j}
       :=C(1+k-j)^{3/2}L^{-\alpha_H(k-j)},
 \qquad \sum_{n\ge1}\omega_n<\infty .
 \label{eq:local-formation-WF14}
\end{equation}
All inner age functions use the same enlarged fine-parameter
domains and the same \(\nu\)-disc. The outer and removed functions
use the enlarged fine-parameter domains with their respective
directional bounds.
At fixed \(j,z\), \eqref{eq:local-formation-WF13a}--\eqref{eq:local-formation-WF13b} provide the common-root
function before any bracket.  The resulting bounds for distinct
base points and levels are then added absolutely.
No bound for a sum of differently based bracket operators has
been inferred from a sup norm of their unbracketed sum.
The uniform bounds justify passage through the bracket
integrals and to limits of the old-level sum.

We have proved the following formation estimate.  Let \(g_j\)
be regular-member norms, \(d_j\) defect norms and \(r_j\)
aged-remainder norms.  Under the explicit background, window
and geometry inputs above, the actual old-origin potential
satisfies
\begin{equation}
 \|\mathbf V_k^{\rm old}\|_\lambda
 \le Cs_k\left[
       \sum_{j\le k}\omega_{k+1-j}g_j
       +\sum_{j\le k}\rho_j^2d_j
       +\sum_{j\in{\cal J}_k}\rho_j^2r_j\right].
 \label{wf:formation-bound}
\end{equation}
These are local bounds on every specified complex core.
At complete physical pairs, finite fundamental theorem of
calculus, the window partition, and hereditary resummation of
\eqref{eq:local-formation-WF13} identify their sum with the actual old action difference.
There is no assertion of a single global off-graph Gaussian
or a global torus background for an arbitrary local pair.

Adding a scalar Wilson-weight direction contributes at most
\(CM^4s_k\alpha_k^2\|u\|_\infty\), by the separate
determinant/response-bracket argument.
The formation rules are linear in their origins, so \eqref{wf:formation-bound}
applies to their first source derivatives.
For a frozen regular member that direction is
\(f(p_z){\cal L}_j[1]\) with its own \(b_j'(0)\);
the remaining direction is the source defect.
The finite potential-ball Taylor remainder justifies the
chain rule even when the actual ordinary source families are
nonlinear.  This gives the formation inequality used in the
linear source bootstrap, with no additional unexplained
source-response norm assumption.

What remains as input is the actual joint complex
window/core room, nonlinear averaged-datum estimates,
H\"older response and compatible boundary comparison,
and the outer/removed field locality for the assigned maps.
None follows from a scalar stability bound.
The result does not establish those background estimates,
replace them by a formal graph identity, or address a
large-field history operation.


%
%
%
%
\paragraph*{Direct windows for the ordinary regular terms.}
\hypertarget{wf:direct-window}{}
In the ordinary constant-sequence part of the step
one can localize the regular family by composing its
members with the fixed box read before taking brackets.
The complete box is then part of the term's carrier.
The carrier and the pointing determine where the
finished term is filed, as described below.

Write \((\mathbf A_s,\mathbf K_s)\) for the actual
decoupled potential and current increment, in
level-\(k\) units. On the ordinary constant-sequence
core the response estimate gives
\[
 \|\mathbf A_s\|+\|\nabla_{\mathbf U}\mathbf A_s\|
        +\|\mathbf K_s\|\le C_R\delta_k
 \qquad (|s_\Delta|<R).
\]
Fix the completed row source before varying the
switches. The base of the response uses the independent
current.
For a fixed upper real factor \(U\) define the shifted
pair by
\[
 P_{s,\lambda}=
  \bigl(e^{i\lambda\eta_k\mathbf A_s}\mathbf U,\,
                      \mathbf J+\lambda L^3\mathbf K_s\bigr)
\]
in level-\((k+1)\) units. The potential and its first
difference cost \(L|\lambda|\|\mathbf A_s\|\) and
\(L^2|\lambda|\|\nabla_{\mathbf U}\mathbf A_s\|\)
in these units. The current increment costs the stated
\(L^3\). Use upper target radii a fixed enlargement
of the incoming radii, below the fixed window ceiling.
On \(|\lambda|<c/s_k\) the three changes are within
the resulting reserve, for small fixed \(c\).
For the simple clauses this follows
from the logarithm product and plaquette estimates.
For the hereditary clauses use the proof of the outer
substitution with the actual fractional room: put the
same shift on a larger auxiliary amplitude disc,
on which the simple bounds at most double, and apply
the tower estimates there. Their holomorphic curvature
and current readouts have the depth factors of
\eqref{eq:rg-objects-R6} throughout that disc.
Schwarz then bounds their change by the reserved
fraction of the same factors. The real factor is fixed,
and the independent current is checked by its own
affine formula. This admits the shifted upper pair
on the full prescribed box, uniformly in the fine
switches.

For inherited levels \(1\le j\le k\), far terms and
a fixed number of the youngest ages
restrict the original pair to the actual old domain and
use the direct old substitution. Put \(a=k-j\) and
\(\sigma_j=L^{-a}\). Its potential, first difference and
current increment are \(\lambda\sigma_j\mathbf A_s\),
\(\lambda\sigma_j^2\nabla_{\mathbf U}\mathbf A_s\) and
\(\lambda\sigma_j^3\mathbf K_s\) in old units. The base
current is \(L^{-3(a+1)}\mathbf J\). The two-level
restriction leaves the old half-core room and the
outer substitution admits a disc of radius
\(c a_{0,j}/(\sigma_j\delta_k)\). Far domains have the
length loss \(d_j(X)\ge c/\sigma_j-C\), which pays for
the old base-point count. For the finitely many young
ages this count is bounded directly.

For the remaining near terms use the same fixed
window box, with its tower truncated at depth \(k\).
The comb map at level \(k\) applies to this determining
set in the window box. Use the full averaged datum at
\(P_{s,\lambda}\), including the outer rows, and
reconstruct its depth-\(k\) critical configuration
in the flat-reference Landau chart. Its induced
depth-\(j\) datum is formed on the truncation of
that same tower. The nested continuation identity
identifies the depth-\(j\) reconstruction with
the depth-\(k\) field in its corresponding orbit.
Read the field and its actual current in the
admitted exponential representative. Thus the
input current is used in making the shift and
the window reads the current of its own reconstruction.

Take the window radii
\(a^\flat_{\iota,j}=\min\{a_{\iota,j},c_\flat
a_{\iota,k+1}\}\) with fixed \(c_\flat>0\).
If \(b_\iota\) are the admitted radii of the shifted
upper pair, the running bounds give
\[
 \max_\iota\frac{b_\iota}{a^\flat_{\iota,j}}
    \le C\bigl(\sqrt{1+a}+c_\flat^{-1}\bigr).
\]
Consequently the requirements
\(CM(b_0+b_1)\sigma_j\le c'a^\flat_{0,j}\)
and \(CMb_0\le c'\) of the window map hold from a
fixed age onward, at weak coupling. In level-\(j\)
units the simple bounds of that representative
on the settled interior are at most \(CMb_0\sigma_j\).
Use the same bound for the curvature and current,
which have additional factors of \(\sigma_j\)
under rescaling. The tower estimate for simple
data then admits both hereditary rules at the
old window radii, since the age can absorb its
fixed constant as well.
On each of the two physical constant-sequence endpoints
the depth-\(k\) box sequence is finer than the ambient
one, and the depth-\(j\) box sequence is finer again.
The no-inward restriction argument applies at each
restriction. Criticality at the induced datum and
regular uniqueness then give the original activity
value by real gauge invariance. On the inner read the
six incident plaquettes of each current stencil are
inside the box, so the constrained current of the
restriction equals the ambient one in the same frame.
Use the full-switch current identity and its node
factors for the shifted physical endpoint.
This gives the window identity at both endpoints
of the bracket.

For the regular gain use the spatial estimate of the
reconstructed upper field. The flat-reference Landau
map in this box has the weighted first-datum-derivative
bound of the walk estimate, including the H\"older
seminorm of its first difference. Integrate this bound
on the line from zero to the complete comb datum.
The weighted sum over row sites is uniform. On the
settled upper interior the local lengths equal one in
upper units, and the reference connection is trivial.
For the Hölder seminorm take cutoffs equal to one
on the stencils of a fixed smaller neighborhood of
each upper cell. Close points use a common cutoff,
and the first-difference sup norm covers the remaining
separations. The result is the ordinary
\(C^{1,\alpha_H}\) bound there, of size \(CMb_0\).
The central old double blocks
are contained in this interior with their translated
stencils. Applying the nonlinear row derivative estimates
to this field gives \eqref{eq:pointed-window-remainder-PW6}
with \(\sigma=\sigma_j\). Take the root-\(z\) axial
gauge on the settled old node lattice and extend it by
the exponential of its logarithm times a cutoff
supported inside that layer. The cutoff is one on the
inner carriers and changes by \(O(\sigma_j)\) per
old link. The comb uses \(O(\sigma_j^{-1})\) links
of logarithmic size \(O(CMb_0\sigma_j)\), so its
logarithm is \(O(CMb_0)\). The axial link identity
also bounds its variation across one link by
\(O(CMb_0\sigma_j)\). The transition links therefore
stay inside the regular logarithm chart, and the
lower rows outside its support retain their complete
original data. Choose the fixed ceiling on \(b_0+b_1\)
to satisfy \eqref{pw:absolute-datum-ceiling}. This admits
the entire axial datum, including the remote rows, in the
fixed continuation half ball uniformly in \(j,k\).
Define the window function by the old Landau continuation
at this full axial datum.
The central top rows have size \(O(CMb_0\sigma_j)\);
all remote rows have size \(O(CMb_0)\), and are at
old weighted distance at least \(c/\sigma_j\) from
the inner read. The weighted Landau estimates for
the potential and its first and second-order reads
therefore give simple bounds
\(CMb_0(\sigma_j+e^{-c/\sigma_j})\) on the complete
old stencils. The same bound holds for the physical
current, by the full-switch first-variation formula.
The tower estimate for these simple bounds admits
both hereditary rules on the old window core,
after increasing the fixed age threshold. This is
the representative used in the Taylor calculation.
The axial change is a real gauge at the physical
endpoints, where invariance gives the same values.
Use the resulting affine curvature
for every member at the same base point and for its
own plaquette tent. The old complete
pointed subtraction and the retained domain tails
then give the bound
\[
 Cg_j\,(\rho_j^\flat)^3\sigma_j^{4+\alpha_H},
 \qquad
 \rho_j^\flat=\max\{1,\max_\iota b_\iota/a^\flat_{\iota,j}\},
\]
before the bracket and the base-point count. A
vacuum-only family uses its quadratic bound
\(C(\rho_j^\flat)^2\sigma_j^4\) times its own norm.

At a base point sum the near window functions before
forming the bracket. Pin every switch cell meeting the
full box, so a cut leaves all rows of this read in the
active component of that base. The switch locality of the
ordinary response then gives locality of the composed
function on the same cut. At \(\lambda=0\) the function is
independent of the switches. The alternating difference
therefore vanishes at zero and Schwarz on the common
amplitude disc gives the factor \(Cs_k\). The fine-cell
Cauchy factors pay for the added connected cells, and
the fixed enlargement of the base costs a fixed constant.
Summing the alternating differences gives the full
shifted value minus the unshifted value of the same
window function. The window identity identifies both
values on the physical pair. Finally \(L^{4a}\) old
base points cancel \(\sigma_j^4\), leaving the
summable H\"older gain for the regular family and the
vacuum-only loads of \eqref{wf:formation-bound}.
This supplies the ordinary regular bound on a source
data line when the same direct window and pointing are
used for all its systems.

%
%
The boundary assignment is made after formation.
A cluster carrier contains the complete carriers of
its vertices. If a window vertex uses a box crossing
\(\Lambda^c\), every cluster that contains it has a
carrier outside \(\Lambda\), even when its pointing
lies in \(\Lambda^0\). Its original endpoint
contribution is then in the boundary cluster family.
The domain bound for that family must use the complete
window read and its boundary pair. For the ordinary
step, put \(t=k+1\). The gap from \(\Lambda_t\) to
\(\Omega_t^c\) is at least \(2\mathsf a_t
=2LM\check R_t\) in level-\(k\) units. A window
box meeting \(\Lambda_t\) has diameter \(12M\) in
these units. The gap puts the whole box and any
fixed enlargement for the response reads inside
\(\Omega_t\), at weak coupling. This uses the
ordinary sequence before preparations. A prepared
stage can have \(\Lambda_t=\Omega_t\).

On the ordinary crossing box the post-step
pointwise and gauge bounds are those of the
top layer \(t\), in the upper units. They supply
the simple core bounds on its full reads. Apply
the tower estimate for simple data at fixed larger
radii to obtain the hereditary core there. Complete
the free source at the pivots on the actual level-\(k\)
double blocks and set it to zero on the remaining
bonds. Separation from \(\Omega_k^c\) puts these
double blocks in \(\Omega_k\). Use the response of
this layered label on \(\pi_k\), including when its
active component extends past the window. The
shift from this response uses the enlarged
fluctuation bound \(C_{\mathcal A}\delta_k\),
including its collar coordinates. In the preceding
upper substitution and the old window calculation
the common amplitude disc can therefore be
reduced to \(c/(C_{\mathcal A}s_k)\). The box
norms and old-age estimates remain the same, while
the bracket now gains \(C C_{\mathcal A}s_k\).
The bound \(C_{\mathcal A}\le C\varrho^{\rm st}(\theta_k)\)
on the enlarged box is sufficient, since
\(\varrho^{\rm st}(\theta_k)s_k\) tends to zero.

%
%
%
\paragraph{The ordinary conditional Gaussian.}
Put \(\mathsf A=\mathsf C^{\mathsf T}\Delta^{(k)}\mathsf C\)
on the free fluctuation coordinates. Write \({\rm I}\)
for the integration bonds, which have an endpoint
in \(\Lambda_t\), and \({\rm E}\) for the remaining
free bonds with an endpoint in \(\Omega_t\).
Let \(\mathsf a\) denote the retained scaled exterior
field. Completing the square at fixed \(\mathsf a\)
gives covariance \(\tau_k^{-1}\mathsf A_{\rm II}^{-1}\)
and mean \(-\mathsf A_{\rm II}^{-1}\mathsf A_{\rm IE}\mathsf a\).
The common ratio multiplies both blocks of the action
and cancels from the mean. For a local integration
set \(W\), let \(P_W\) be the projection onto its
integration coordinates. With the same guarded
switches \(\hat s\) in every factor, put
\[
 \begin{aligned}
 C&=(P_W\mathsf A_{\rm II}(\hat s)P_W)^{-1},
 & \Gamma&=P_W\mathsf A_{\rm II}(\hat s)(1-P_W)R,\\
 m&=\mathsf M\mathsf a,
 & \mathsf M&=C(P_W-\Gamma R^{\mathsf T})
                                   \mathsf A_{\rm IE}(\hat s).
 \end{aligned}
\]
Here \(R\) is the inverse root in the switched walk
expansion. The signed convention gives mean \(-m\).
At the full switches the block inverse identity gives
\(C(P_W-\Gamma R^{\mathsf T})=P_W\mathsf A_{\rm II}^{-1}\),
so this is the restriction of the same conditional
mean. The local covariance and coupling to the
auxiliary Gaussian are rescaled as
\((\tau_k^{-1}C,\tau_k^{1/2}\Gamma)\).

The weighted row bounds for \(\mathsf M\) and its
difference from the real reference \(\mathsf M_{\rm r}\)
have a common rate. Denote the small prefactor in
the difference bound by \(\varepsilon_*\). At a
crossing bond the threshold is
\(\mathfrak b_b=\theta_{\rm c}\mathfrak b\); at other
integration bonds it is \(\mathfrak b\). The exterior
datum is the sum of a real field obeying the three
tiers and a complex correction. The collar bound is
\(\mathfrak b/C_{\rm col}\), the rim bound is
\(\theta_{\rm c}\mathfrak b\), and the bound beyond
the rim is \(C\mathfrak b\). The correction, of size
at most \(\varepsilon_{\mathsf a}\mathfrak b\), starts
at distance \(W_t-O(1)\) from the integration bonds,
where \(W_t=LM\check R_t\). Splitting the row sum
at the collar width \(\varpi_T\) and at \(W_t\)
gives
\[
 \frac{|m(b)|}{\mathfrak b_b}
 \le C\left(\frac{1}{C_{\rm col}\theta_{\rm c}}
      +e^{-c\varpi_T}
      +\frac{1+\varepsilon_{\mathsf a}}{\theta_{\rm c}}
                                     e^{-cW_t}\right).
\]
The same estimate for the difference from
\(\mathsf M_{\rm r}\) on the real exterior field
has the factor \(\varepsilon_*\). Since that
reference mean is real,
\[
 \frac{|\operatorname{Im}m(b)|}{\mathfrak b_b}
 \le C\varepsilon_*+
        C\theta_{\rm c}^{-1}\varepsilon_{\mathsf a}e^{-cW_t}
 =:\varepsilon_\mu .
\]
Choose the collar threshold and width as in the
real step, then use the width of the rim to absorb
the far-field enlargement. This gives
\(|\operatorname{Re}m(b)|\le\theta_\mu\mathfrak b_b\)
with \(\theta_\mu<1/2\). The guard and small radii
of the switched structure give
\(\varepsilon_\mu\le c(LM)^{-2}\).

To separate the real cutoff factors, fix the selected
Mayer vertices and let \(Y_0\) contain the free bonds on
which they require a field box. In the signed convention
above \(B=\xi-m\) with \(\xi\) centered. On a bond of \(Y_0\)
split the displaced indicator into
\[
 \mathbf 1_{\{|\xi_b-m_b|<\mathfrak b_b\}}
 =h_b(|\xi_b|)
  +\big[\mathbf 1_{\{|\xi_b-m_b|<\mathfrak b_b\}}
        -h_b(|\xi_b|)\big],
\]
where \(h_b=1\) up to \(\mathfrak b_b/2-1\), decreases
linearly to zero at \(\mathfrak b_b/2\), and is zero
beyond. The taper is fixed as the mean varies. At real
data the bracket lies between zero and one, since the
support of the taper is inside the displaced ball. Its
support satisfies \(|\xi_b|\ge\mathfrak b_b/2-1\).
Outside \(Y_0\) use one minus the failed indicator, whose
failure implies
\(|\xi_b|\ge(1-\theta_\mu)\mathfrak b_b\).
Thus each selected remainder of a cutoff and each selected
failure has a Gaussian tail cost. Summing these choices
restores the physical product of cutoffs before taking the
polymer logarithm.

For complex parameters continue the integral over the
displaced ball by
\[
 f\longmapsto
 \int_{|A_b|<\mathfrak b_b} f(A_b+m_b)\,dA_b .
\]
Here \(f\) is the common integrand in that centered coordinate.
The taper integral keeps its real support fixed. The failure
integral is the full integral minus the ball integral, and
only coordinates outside the selected boxes use that full
integral. To estimate a continued difference, freeze the
integrand and subtract the ball integral with \(m_b\)
replaced by \(\operatorname{Re}m_b\). The difference
integrates over the side swept by the real sphere as the
imaginary part of \(m_b\) is turned on. Every point of
that side has real radius at least
\((1-\theta_\mu)\mathfrak b_b\). The remaining difference
is the real tapered shell or real complement above.
This permits the Gaussian tail estimate on the difference
itself. Joint holomorphy follows from the fixed-domain
formulas by domination on the admitted boxes.

For the sharp radial tests, integrate the Gaussian
auxiliaries and let \(Q\) be the marginal precision.
It is complex symmetric, with a uniform real floor,
bounded norm and small imaginary part. Write
\(m=u+iv\) and choose a smooth radial function
\(\psi\) with values in \([0,1]\), equal to one up to \(3/4\) and zero from
\(7/8\) onwards. In each integration coordinate
deform the real contour by
\[
 z_b=x_b-i v_b\psi(|x_b|/\mathfrak b_b).
\]
Apply the deformation to the physical ball and
complement integrals. The fixed taper is integrated
on its real support. The deformation fixes the test sphere, is the
identity on its complement, and stays inside the
complex ball of each boxed holomorphic factor.
Stokes' theorem therefore applies on each product
piece of the radial tests. The Gaussian majorant
permits passage to infinity on unrestricted
coordinates, where the free quadratic potential
is entire. The exterior test factor is constant
in these integration coordinates.

Put \(y=x+u\) and
\(w_b=v_b[1-\psi(|x_b|/\mathfrak b_b)]\).
Where \(w_b\ne0\) one has
\(|y_b|\ge |x_b|-|u_b|\ge\mathfrak b_b/4\),
so \(|w_b|\le4\varepsilon_\mu|y_b|\). It follows that
\[
 \operatorname{Re}\big((y+iw)^{\mathsf T}Q(y+iw)\big)
 \ge y^{\mathsf T}(\operatorname{Re}Q)y
   -C\big(\varepsilon_\mu^2+
      \varepsilon_\mu\|\operatorname{Im}Q\|\big)|y|^2 .
\]
The complex Jacobian on one bond is
\(1-i v_b\cdot\nabla[\psi(|x_b|/\mathfrak b_b)]\),
with modulus at most \(e^{C\varepsilon_\mu^2}\).
For the normalized Gaussian comparison use the real
precision obtained by subtracting this quadratic loss
from \(\operatorname{Re}Q\). The trace-log estimate
between the resulting normalization and that of \(Q\)
costs at most the dimension times
\(C(\|\operatorname{Im}Q\|+\varepsilon_\mu^2)\).
With the Jacobian this is \(O((LM)^{-4})\) per
variable, since the guarded comparison gives
\(\|\operatorname{Im}Q\|=O((LM)^{-4})\).
Each selected tail is estimated with a small centered
quadratic tilt. The tapered shell starts at
\(\mathfrak b_b/2-1\), and the failed cutoff and
the displaced side start at
\((1-\theta_\mu)\mathfrak b_b\) in the real part of the
centered coordinate. The shifted-side bound pays a factor
\(C(1+|v_b|)\exp(C|v_b|^2)\) on that bond in addition to
the real Gaussian tail. The bound on \(\varepsilon_\mu\)
and the weak-coupling floor on the thresholds absorb this
factor into the tail, leaving \(e^{-c\mathfrak b_b^2}\)
per selected bond. The free quadratic vertex is
estimated using its own local support. On a
fine cube \(Y\) its bound after the translation
is \(C|\tau_k-1|(M^4\mathfrak b^2+\sum_{b\in Y}|y_b|^2)\).
Its Mayer estimate uses weight \(C_M|\tau_k-1|\mathfrak b^2\) in a
boundary formation. Choosing the constant in
that weight large enough makes the centered
quadratic tilt \(O((LM)^{-4})\) per variable
after summing the local free cubes.
The boundary ratio disc keeps this weight small.
In formations contained in \(\Lambda_t\) the
mean is zero, and the centered estimate gives
the sharper interior weight. The fixed real
contour gives holomorphy by local domination.

%
%
%
%
%
\paragraph{Differences of displaced tests.}
An exterior selector of the ordinary integral stays
in the real retained kernel. In the unchanged collar
its fluctuation coordinate is the original one.
On the cutoff zone a sharp bond test at threshold
at most \(\delta_k\) is unchanged by the real
substitution. The test against an approximate
background is instead read using the localized
backgrounds on its retained cubes. These identities
apply before assigning switches to the analytic
terms. In particular a sharp selector is not
evaluated on a complex field.
For the ordinary step on the standing part
with its frames, sum out the passengers before
forming the conditional integral. The new
tests force the carried gauge indicators and
the old plaquette tests on the integration
region to equal one. The block-field tests
read the fixed upper datum. The adjacent
approximate bond factors that read the
integration variables also equal one there:
the ball test bounds the deviation from the
global background by \(\delta_k\), and the
interior comparison to the test background
leaves room below the threshold \(2\delta_k\).
The remaining factors of those tests are
carried by the real exterior kernel. This
removal is an identity on the support of the
whole term. It uses both the internal ball
tests and the retained fluctuation factors
in the unchanged part of the strip. At pivot
bonds use the local constraint reconstruction
on the full double block. This gives back the
value in the retained test. With the zone
equations the old whole-term description
puts the characteristic functions and
constraint distributions in the real kernel
as the analytic pair varies. The conditional
integral thus has precisely the product over
\(\mathfrak X_{\rm I}\) of its ordinary and
crossing balls for the conditional integral.

For an integrated selector whose domain does move,
write \(D_b\) for its nominal real ball and let
\(\phi_{b,\zeta_b}\) be a transport of the ball,
with \(\phi_{b,0}\) the identity. At a real
displacement it gives the tested domain
\(D_b(\zeta_b)=\phi_{b,\zeta_b}(D_b)\).
Suppose the transport and its spatial derivative
continue on the displacement disc, with their
images in the common holomorphic read. For the
calculation hold the Gaussian and the selected
potentials fixed as a common integrand. On the
coordinates in \(S\) take the difference between
the transported and nominal domains. Expanding
the product gives the alternating sum over
\(T\subset S\), in which exactly the balls of
\(T\) are transported. At the physical real
arguments the result is the integral of the
common integrand times
\[
 \prod_{b\in S}
    \bigl(\mathbf 1_{D_b(\zeta_b)}
                       -\mathbf 1_{D_b}\bigr)
\]
and the unchanged selectors. The product
vanishes unless each bond of \(S\) is in its
own difference of test domains. A change of
the integrand between endpoints belongs to
a separate analytic difference.

To continue this identity, use the pullback
by \(\phi_{b,\zeta_b}\) of the holomorphic
integrand and volume form to the fixed ball.
Equivalently, for \(b\in S\), integrate over
the chain swept by \(\partial D_b\) under
\(\phi_{b,r\zeta_b}\), \(0\le r\le1\), with
boundary orientation giving the difference
of the two balls. This follows by Stokes'
theorem in that coordinate, and iteration
gives one chain per moved test. A difference
of complementary selectors has the opposite
sign and the same chain. In the bounded
logarithm chart the real strip shift and
its first derivative have sizes
\(O(|H|)\) and \(O(|H|/\delta_k)\),
respectively, before scaling the field.
In the plain zone the group multiplier
itself continues holomorphically. In a
cutoff zone any occurrence of the smooth
cutoff in a transport formula must stay
at the real parameter of the chain.

For transports whose whole disc gives a
displacement \(o(\mathfrak b_b)\) and bounded
first derivatives in the scaled coordinates,
the chain of a nominal sphere has real
coordinate within \(o(\mathfrak b_b)\) of that
sphere, and its imaginary coordinate is
\(o(\mathfrak b_b)\). Thus the centered
real coordinate on the chain has norm
at least \(\mathfrak b_b/3\), by the tier
bound. In the other coordinates use the
ordinary Gaussian contour above. The
additional imaginary part on the chains
costs an arbitrarily small multiple of
\(\sum_{b\in S}\mathfrak b_b^2\) in the
quadratic estimate, by the uniform operator
bound and the real floor of \(Q\).
To see this, apply the quadratic inequality
to the cross term between the additional
imaginary part and the ordinary imaginary
part, using the bound \(4\varepsilon_\mu|y|\)
for the latter. The cross term with the
real vector has coefficient
\(\|\operatorname{Im}Q\|\). Together they
cost \(O((LM)^{-4})|y|^2\) and a fixed
multiple of the squared norm of the
additional imaginary part.
The chain volume is at most a fixed
polynomial in \(\mathfrak b_b\) per bond.
These costs leave \(e^{-c\mathfrak b_b^2}\)
for each moved test. The carrier includes
the test bond and the reads of its
transport. Summing the mixed differences
over \(S\) at their physical endpoints
recovers the full transported product
by inclusion and exclusion.

%
%
%
For the product balls of the ordinary step
one can bound the large switch discs by
recentering before forming these differences.
Use the prescribed radius
\(\rho=(\log\theta_k)^{6r+1}\) for the new switches.
Choose the constant in the enlarged fluctuation
box \(C_{\rm form}\rho\mathfrak b_b\) after the
absolute row bound for the mean. This fixed enlargement
fits in the reader estimates because
\(C_{\rm form}\rho\delta_k/\alpha_{\iota,k}\to0\).
It changes the bound on the cubic potential only by
a fixed factor, since
\((C_{\rm form}\rho)^3\theta_k^{-1/2}\mathfrak b^3\to0\).
The later stage therefore has the prescribed switch
radius available for its Cauchy estimate. Inherited
spectator switches keep their original discs.
In a local Gaussian activity let \(m_\mu\)
be the signed mean with each new near field
entering as \(\mu_{b'}\mathsf a(b')\). The
boxed factors in the centered coordinate \(y\)
read \(B=y-m_\mu\) and keep the same analytic
arguments throughout the following calculation.
Write their product, including the free
factors, as \(H_\mu(y)\). At a physical real
point the success domain in \(y_b\) is
\(D_b+m_b\), and the failure domain is its
complement. Replace each success by the
centered success plus its translated-domain
difference. Replace each failure by the
centered failure minus that same difference.
The integrand \(H_\mu\) stays fixed in each
of these operations.

For a row \(b\) of the mean split off the
exterior bonds within weighted distance
\(D=C_0(1+\log\rho)\) of \(b\), using the
distance of the kernel bound. Use only
columns of the local mean.
The weighted row estimate makes the contribution
from the other bonds, after amplification by
\(\rho\), smaller than \(c(LM)^{-2}\mathfrak b_b\)
if \(C_0\) is sufficiently large. For the close
near bonds \(N_b\) keep their real values as
spectators. In the translated-domain difference
for \(b\), use the shifts
\[
 \zeta_{b,T}=m_b\bigl(
   \mathsf a_{N_b}\mathbf1_T,\,
   (\mu\mathsf a)_{N_b^c}\bigr),\qquad T\subset N_b .
\]
Here only switched coordinates in \(N_b^c\)
carry multipliers, and far analytic fields use
their given arguments. Take the sum of the
resulting swept chains with coefficients
\[
 \prod_{b'\in T}\mu_{b'}\,
                  \prod_{b'\in N_b\setminus T}(1-\mu_{b'}).
\]
The tier estimate for the real close fields
still holds when some are set to zero. The
far tail estimate and the comparison to the
real kernel give
\(|\operatorname{Re}\zeta_{b,T}|\le\theta'\mathfrak b_b\)
with \(\theta'<1\), and
\(|\operatorname{Im}\zeta_{b,T}|
 \le C(LM)^{-2}\mathfrak b_b\).
Indeed the unamplified far data outside the
real tiers begin at distance \(W_t-O(1)\)
and give the same exponentially small error
as in the mean estimate. The translated
boundary thus remains at distance at least
\((1-\theta')\mathfrak b_b\) from zero in
its real coordinate. The chain estimate for
the centered Gaussian applies on the whole
switch disc.

In the ordinary step the distance in the
kernel bound controls the distance in
level-\(k\) units between these top-layer
bonds. Thus the number of close bonds is
at most \(C(D+1)^4\). The sum of the
absolute interpolation coefficients costs
at most \((1+2\rho)^{C(D+1)^4}\) for
one translated test. Its logarithm is
\(O((1+\log\rho)^5)=o(\mathfrak b^2)\).
This spends only part of the tail charge
for that test, even when many exterior
switches enter the same row. An original
failed test thus retains a tail factor,
whether read as a centered failure or as
a swept boundary. Differences of the
mandatory success tests of a boxed vertex
are on the union of the boxed reads. Count
each bond in that union once, even when
several vertices read it. Summing the
choices over a cube adds at most
\((1+e^{-c\mathfrak b^2})^{C(LM)^4}\)
to the local bound.

The recentered integrand is admitted on
these chains. At a mandatory boxed coordinate
\(|y_b|\le C\mathfrak b_b\), while
\(|m_\mu(b)|\le C\rho\mathfrak b_b\) by
the absolute row bound. The formation-box
constant accommodates the sum in the
enlarged fluctuation box. The free
quadratic factors on unrestricted
coordinates have the centered tilt and a
local cost bounded by
\(C_M|1-\tau_k|\rho^2\mathfrak b^2\).
This tends to zero in the boundary tube,
where \(|1-\tau_k|\le
C(\log\theta_k)^{-2q_{\rm t}}\) and
\(q_{\rm t}\ge q>\mathfrak p_0+\deg\rho\);
here \(\deg\rho\) is the power of
\(\log\theta_k\) defining the switch radius.
The same bound applies to the boxed
quadratic correction. The other vertices
use their enlarged-box estimates. The
estimates are uniform on the Gaussian
decoupling discs, so the usual Cauchy
bounds for those switches still apply.

At \(\mu\equiv1\) each chain sum selects
the original shift. The affine substitution
\(y=B+m\) puts each original cutoff ball
on the translated chain \(D_b+m_b\).
In a complementary test use full space
minus the ball. Where the domain is full
space the remaining factors are entire
in that coordinate, and the Gaussian
floor permits the affine contour to
return to the real space. This recovers
the original local activity before
resummation, also at the admitted complex
arguments. At \(\mu_{b'}=0\) the
analytic arguments and all chain sums
forget \(\mathsf a(b')\). At
\(\mathsf a(b')=0\) the two corresponding
endpoints agree and their weights add
to one, so the result forgets \(\mu_{b'}\).
The continuations are integrals over fixed
real parametrizations and are jointly
holomorphic with the real spectators fixed.
They factor over wall components because
the cut mean has zero columns off its
component and all the selectors are
coordinate factors. On the cutoff zone
the smooth function in the real strip
map and its Jacobian still use the retained
real variable. For a mixed-history
operation before passengers have been
removed, or with additional integrated
selectors, we must first reduce those
selectors to the same product or transport
their actual domains.

The window is inside \(\Omega_t\), while a cluster
containing it can also have other carriers outside
\(\Omega_t\). Use the new cluster estimate at
formation, since the long-carrier estimate for an
inherited boundary term requires support on both
sides of \(\Omega_t\). When a window is
unavailable, the corresponding origin still needs its
assigned boundary or direct read with the same physical
endpoint identity.

%
%
%
For the source comparison fix the label, Gaussian
centre and retained exterior data. For the regular
origins let \(D_E(X)\) be the age-weighted sum of
local differences of the inherited pointed families
on the full reads of \(X\). Vary each member with
its companion and its own tent coefficient. The
defects, aged remainders and local Wilson weights
have the corresponding loads in
\eqref{wf:formation-bound}, with the actual
layer radii for Wilson terms off the regular
region. Write \(D(X)\) for the sum of these loads.
The crossing vertex bound just proved is linear
in the old origin, so a changed vertex has bound
\[
 C\varrho^{\rm st}(\theta_k)s_k D(X)\,
                              e^{-r_Vd_k(Y)}
\]
on its actual carrier \(Y\). In a selected rooted
cluster interpolate its input arrays on the common
data line. Differentiation selects one changed
vertex; the derivative of its Mayer factor is
\(e^V\,\dot V\), with the same uniform sup bound on
\(V\) as in the unmarked expansion. The absolute
rooted summation at that vertex, with the unused
decay for the other branches, gives
\[
 |\Delta B_{\rm cross}(X)|
 \le C\varrho^{\rm st}(\theta_k)s_k D(X)\,
                    e^{-\kappa_{\mathbf T}d_{k+1}^\natural(X)} .
\]
The same proof bounds a fixed partial sum of
complete clusters, by the same positive majorants.
In particular one can select the clusters with a
crossing origin by its record in the formation list.
It applies to the source line over the changed
slots, with the full pointed member retained at
each slot. The small factor is uniform over the
enlarged box. The endpoint identity uses all filed
pieces at their actual physical switches.

%
%
The changed crossing clusters are used a second time in
the ordinary strip transfer. Write $E^\sharp$ for their
complete formed crossing sum and let $\mathbf S$ be the
sum of the nonempty strip-displacement brackets. It is
linear in the input potential, with each bracket formed
on the actual enlarged strip carrier. The strip input is
$E^\sharp+\mathsf K^{\rm kp}+\mathsf B^\times+\mathsf B^\tau$
minus the local Haar logarithms. The strip output is
$\mathbf S$ of this input plus its strip Jacobian differences.
At fixed geometry, fields, $\theta_k$ and scalar center $x_k$,
a change of the regular input arrays or of the spatial
Wilson weight enters through the formed vertices. Indeed
the nonconstant part $x_k(\cdot)-x_k-\zeta_k(\cdot)$ is a
Wilson potential, while the scalar center enters the Gaussian
ratio. Thus the direct crossing output changes by
$\Delta E^\sharp$, and the strip output changes by
$\mathbf S[\Delta E^\sharp]$. The separate form and normalizer
outputs and the strip Jacobians agree on this comparison.
A change of the scalar center requires their own ratio and
form estimates. Changes of old boundary or standing factors
in the same formula use their own loads and majorants.

In each nonempty strip bracket, subtract the evaluations
before summing the bracket. A displacement direction has
size controlled by the localized response from the strip
to its localization boundary. Cauchy's estimate on the
transfer disc bounds that difference by the input norm
times the displacement divided by its available radius.
The exponential in the strip-to-boundary distance pays
the enlargement factor of the strip disc and the ratio
of the field radius to the fluctuation radius. The total
strip term uses its own absolute majorant on the enlarged
data line. This yields the strip source slope in its own
decay norm. The direct crossing slope and the slope of
the displaced output are assigned on their respective
carriers.

For a carrier \(X\subset\Lambda\), the ordinary
formula is used with its actual vertices. It is filed
in the regular family at \(z\in\Lambda^0\), and in
the edge sum \(E^\partial(X)\) at \(z\notin\Lambda^0\).
At most \(CM^4(1+d_{k+1}(X))\) such sites can point
into \(X\), by the connected cube count. The spare
unit in the pointed rate pays this factor. Thus the
edge sum has the ordinary bound with one unit of
decay spent, and its source difference has the
corresponding data-line slope: the common
\(\mathbf N^\circ\) and \(\mathbf G_\tau\) cancel
before the sum. The Gaussian normalizer outside
\(\Lambda\) is a separate term. This assignment
retains the terms of the exponent with each complete
carrier counted once.

\paragraph*{Reconstructing the field of an old near term.}
On a box of the size of a level-\((\ell-1)\) cube about its base
point, a local level-\(\ell\) triple has curvature, potential and
current radii \(b_0,b_1,b_2\). Its box reading uses the full box
in the triple's domain and requires that the radii have small sum.
Take the averaged datum of the product field in the axial gauge of
that box and solve the old constrained critical problem for that
datum \cite[(3.26)--(3.27)]{B87}. Evaluate the old
term at the reconstructed field and at its current. This is a
holomorphic function of the input field on the box. It ignores
the independent input current. Its reconstructed potential has
positive H\"older exponent. The absolute smallness conditions
are \(CMb_0\le c\) and \(CMb_0 L^{-a}\le ca_{0,j}\),
where \(a=\ell-j\) is the age and \(a_{0,j}\) is the old curvature
radius. Independence of the input splitting uses the stronger
condition \(CM(b_0+b_1)L^{-a}\le ca_{0,j}\).
At a matched real critical endpoint the box reading agrees
modulo gauge by the restriction and uniqueness argument above.

On upper triples with \(b_0+b_1\le C\alpha_k\) the old
half-core condition holds at every sufficiently large age,
since \(\alpha_k/a_{0,j}\le C\sqrt{1+a}\). For the finite
set of younger ages, use the old core interpolation when the
whole old carrier is in \(\Omega_j\). For old term radii
\(a_{\iota,j}\ge\alpha_{\iota,j}\), the step uses at most
\(\phi_j^+\alpha_{\iota,j}\), with \(1-\phi_j^+\ge c_0/2\).
The allowed disc in the dimensionless amplitude of the old
response has radius
\[
 \left[C(1-\phi_j^+)^{-1}\frac{\delta_k}{\alpha_*}
 e^{C\varkappa}e^{-\delta_*\mathsf d_k(X)}\right]^{-1},
 \qquad \alpha_* = \min_\iota \alpha_{\iota,k}.
\]
Here \(\mathsf d_k(X)\) is the response distance of the old
carrier. At these ages the number of old base points per
window is bounded by a fixed constant. The ordinary core
estimate then supplies the corresponding young-age bound
without a H\"older improvement. The dimensionless amplitude
in this estimate is the old term's own interpolation parameter.

This use of the reconstruction requires an admitted triple
on the complete box in \eqref{eq:local-formation-WF3}.
For the ordinary bracket the full-box frame gives the row
comparison for the unamplified path and its cutoff
perturbation. To use the positive H\"older bound we also
need the prescribed reconstruction and the compatible
boundary comparison on the entire addressed old domain.

\paragraph*{Own readouts and the order of localization}

The scalar readout is taken before the old member is evaluated on
the new box background.  We record the resulting algebra because
it fixes both the normalization and the first jet.

At old level \(j\) and base point \(z\), write \(A_{j,z}\) for the
member's assigned tent action.  Let \(b_j\) be its pointed scalar
readout with the fixed free-response kernel and fixed chart.
Its definition is a linear combination of second derivatives at
the identity with absolutely convergent moment weights.  Hence
it is complex linear and continuous in that Hessian-moment norm.
Its action normalization is \(b_j(A_{j,z})=1\), as in the
pointed quadratic classification.  Set
\begin{equation}
 {\cal R}_{j,z}G
       =G-G(0)-b_j(G)A_{j,z}.
 \label{fm:own-old-renormalization}
\end{equation}
The zero denotes the identity point of the potential chart.
The own coefficient and own action are held together in this
definition.  The quadratic classification proves the
cancellation for this pair; it is not a cancellation between a
coefficient and a tent action taken from unrelated families.

Fix \(j,z\) and its current window \(R_{k,z}\).  Let
\(\mathcal N_{j,z}\) be the near domains retained in \eqref{eq:local-formation-WF13a}, and
let \(b_j(G)\) be the coefficient of the complete incoming pointed
member, extracted before this truncation.  Write \(\Phi_{G_X}\),
\(\Phi_{G_X(0)}\) and \(\Phi_{A_{j,z}}\) for evaluations through
its one common constrained box map.  The bracket
\(\mathcal L_{Y;R_{k,z}}\) has the fixed declared root
\(R_{k,z}\); it differentiates each selected fine-cube parameter
once and takes the prescribed fluctuation difference.  In finite
volume linearity gives exactly
\begin{equation}
 \begin{aligned}[b]
 &\mathcal L_{Y;R_{k,z}}\!\left[
   \sum_{X\in\mathcal N_{j,z}}(\Phi_{G_X}-\Phi_{G_X(0)})
                         -b_j(G)\Phi_{A_{j,z}}\right]\\
 &=\sum_{X\in\mathcal N_{j,z}}\mathcal L_{Y;R_{k,z}}
                  [\Phi_{G_X}-\Phi_{G_X(0)}]\\
 &\quad-b_j(G)\mathcal L_{Y;R_{k,z}}[\Phi_{A_{j,z}}].
 \end{aligned}
 \label{fm:common-window-linear-bracket}
\end{equation}
For an infinite absolutely convergent old-domain sum, the same
identity follows by uniform convergence on a slightly larger
interpolation polydisc and Cauchy's formula.  The own counterterm
occurs once with the complete member's coefficient.  The domains
omitted by the near truncation remain in the separately estimated
tail; no scalar coefficient is assigned to an individual domain.

This is the precise use of window-first normalization.  It
allows the pointed cancellation to be bounded before applying
one bracket at the fixed \(j,z\).  Different \(z\)'s may have
different roots; their bounds are summed afterwards with the
site and age counts of \eqref{eq:local-formation-WF13}.  No bound on a sum of functions
is applied to a sum of different bracket functionals.
Nor does \eqref{fm:common-window-linear-bracket} assert that
the newly formed bracket member has scalar readout \(b_j(G)\).
That readout belongs to the old member before localization.

If \(G(w)\) is a differentiable source family in this fixed
presentation, continuity and linearity give
\begin{equation}
 \left.\partial_w{\cal R}_{j,z}G(w)\right|_0
       =\dot G-\dot G(0)-b_j(\dot G)A_{j,z}.
 \label{fm:own-counterterm-first-jet}
\end{equation}
Thus a derivative of the formed member loses no normalization
merely because its scalar coefficient has already been
extracted.  The derivative has its own coefficient, computed
with the same free response.  This is the rule used in the
linear bootstrap.

If the presentation instead moves its readout or own action,
the complete elementary product rule is
\begin{equation}
 \left.\partial_w[b_j^w(G^w)\,A_{j,z}^w]\right|_0
  =\{\,\dot b_j(G^0)+b_j^0(\dot G)\,\}A_{j,z}^0
                  +b_j^0(G^0)\dot A_{j,z}.
 \label{fm:moving-readout-product-rule}
\end{equation}
In this display the superscript zero denotes source parameter
zero; \(G^0\) is a member, rather than its value at the identity
chart point.
The distinction matters in a reference-change argument.
\eqref{fm:own-counterterm-first-jet} uses the
fixed-presentation case; it cannot remove either extra term
of \eqref{fm:moving-readout-product-rule}.


\subsubsection{The many-scale first-source recurrence}


This section isolates a linear consequence of the ordinary local formation
bounds.  It uses the inserted-potential and renewed-pointing estimates at
one step, rather than a larger domain of physical sources.  In particular,
it distinguishes the derivative of a newly generated activity from the
old background action that is carried to the next scale.
The formation and domain hypotheses below concern the actual local
maps.  Their statement is not an assertion that a nonlinear history
construction satisfies them.

\paragraph*{Meshes, norms, and the fixed ordinary run}

Work first on a finite periodic lattice.  Let \(h_j=L^jh_0\) be the
physical mesh of level \(j\), with \(L>1\).  A source profile is a
bounded Lipschitz function on the physical torus; complex profiles are
allowed by complexification.  Put
\begin{equation}
 \|f\|_{\sharp,j}=\max\{\|f\|_\infty,C_Mh_j\operatorname{Lip}f\}.
 \label{eq:local-formation-LB1}
\end{equation}
At level \(j\), a pointed localization domain \(X\) contains its base
point \(z\).  Its enlarged source region \(X^\sim\) and the fixed
plaquette point \(p_z\) satisfy
\begin{equation}
 \sup_{p\in X^\sim}|p-p_z|_{\rm tor}
       \le C_Mh_j(1+d_j(X)).
 \label{eq:local-formation-LB2}
\end{equation}
Distances and diameters in this section use the torus metric.
The family norm with rate \(\gamma\) is the supremum of
\(e^{\gamma d_j(X)}|F(X,P,z)|\) over the actual specified complex
field domain and all \(X,z\).  Companion families, when present,
are included in the same norm.  A bound with rate \(\kappa+1\)
therefore leaves one unit of rate for a diameter.

The ordinary run fixes all Gaussian structures, background maps,
thresholds, label spaces and reference operations.  Sources change
local exponent members only.  Let
\begin{equation}
 \alpha_j=\theta_j^{-1/2}(\log\theta_j)^q,\qquad
 \delta_j=\theta_j^{-1/2}(\log\theta_j)^{p_0},\qquad
 s_j=\delta_j/\alpha_j,\qquad q>p_0\ge0.
 \label{eq:local-formation-LB3}
\end{equation}
Fixed multiplicative constants in these definitions can be included.
Assume the ordinary flow bounds, with constants independent of the
number of levels,
\begin{equation}
 c(k-j)-C_0\le\theta_j-\theta_k
                       \le C(k-j)+C_0\qquad(j\le k).
 \label{eq:local-formation-LB4}
\end{equation}
All \(\theta_j\) are sufficiently large.  These bounds imply, with a
uniform constant \(C_F\), that for \(1\le j\le k\),
\begin{equation}
 \begin{aligned}
 \alpha_{j-1}^2s_{j-1}&\le C_F\alpha_k^2s_k,\\
 \varrho_{k,j}^2\alpha_{j-1}^2s_{j-1}
                         &\le C_F\alpha_k^2s_k,\qquad
 \varrho_{k,j}=\alpha_k/\alpha_j,\\
 \varrho_{k,j}^2&\le C_F(1+k-j).
 \end{aligned}\label{eq:local-formation-LB5}
\end{equation}
Indeed the functions \(t^{-1}(\log t)^{2q}\) and
\((\log t)^{p_0-q}\) are decreasing for sufficiently large \(t\).
The bounded displacement \(C_0\) in \eqref{eq:local-formation-LB4} changes their ratios by
at most a constant.  The second line also uses the bounded
consecutive ratio \(\alpha_{j-1}/\alpha_j\).  The upper flow bound
and \(\theta_k\ge1\) prove the last line.  These observations also
give \(s_j\le C_Fs_k\) when \(j\le k\).

\paragraph*{Locality supplies the source-gradient gain}

Let \(E_j^w\) be an ordinary sourced regular family, holomorphic in a
neighborhood of \(w=0\).  Define its first derivative by
\begin{equation}
 {\cal L}_j[f]=\left.\frac{d}{dt}E_j^{tf}\right|_{t=0},
 \qquad
 A_j=\sup_{f\ne0}
        \frac{\|{\cal L}_j[f]\|_{\kappa+1}}{\|f\|_{\sharp,j}}.
 \label{eq:local-formation-LB6}
\end{equation}
The supremum is taken over the permitted profiles and the tori in
the run.  At a fixed finite lattice all differentiations are in
finite-dimensional source coordinates.  Locality means that the
value indexed by \(X\) depends only on the values of \(f\) on
\(X^\sim\).  To apply the following argument it suffices that the
source class admit bounded Lipschitz truncations of a profile and
constants, as does the class defined by \eqref{eq:local-formation-LB1}.  No density claim
about an infinite-dimensional algebraic dual is used.

The derivative defect at a base point is exactly
\begin{equation}
 {\cal D}_j[f](X,P,z)
       ={\cal L}_j[f](X,P,z)
                   -f(p_z){\cal L}_j[1](X,P,z).
 \label{eq:local-formation-LB7}
\end{equation}
We claim
\begin{equation}
 \|{\cal D}_j[f]\|_\kappa
                  \le C_M A_jh_j\operatorname{Lip}f.
 \label{eq:local-formation-LB8}
\end{equation}
To prove it, fix \(X,z\), and subtract \(f(p_z)\) from \(f\).
On \(X^\sim\) the resulting profile has modulus at most
\(C_Mh_j(1+d_j(X))\operatorname{Lip}f\).
For real \(f\), clip its values outside an interval twice this
large, using a Lipschitz scalar map that is the identity on that
interval and has Lipschitz constant at most a fixed constant.
For a complex profile clip the real and imaginary parts separately.
The resulting global profile \(g\) agrees with \(f-f(p_z)\) on
\(X^\sim\) and satisfies
\(\|g\|_{\sharp,j}\le
C_Mh_j(1+d_j(X))\operatorname{Lip}f\).
Locality and linearity give
\({\cal D}_j[f](X,P,z)={\cal L}_j[g](X,P,z)\).
Apply \eqref{eq:local-formation-LB6} and
\((1+d)e^{-d}\le1\) for \(d\ge0\) to obtain \eqref{eq:local-formation-LB8}.
Smooth profiles may instead use a smooth clipping map with the same
fixed bounds.  This proof applies uniformly at every external
complex field point in the family norm.
For a source measured at level \(m\ge j\), it gives in particular
\begin{equation}
 \|{\cal D}_j[f]\|_\kappa
       \le C A_jL^{-(m-j)}\|f\|_{\sharp,m}.
 \label{eq:local-formation-LB9}
\end{equation}
Thus the derivative defect retains the response amplitude \(A_j\)
and the scale-gradient gain together.

\paragraph*{The precise one-step inputs}

For clarity we list the ordinary inputs needed for the bootstrap.
The first two are estimates for local formation and the finite
regular gas, at fixed background maps; they are not estimates for
complex history masses.

The formation of potentials is linear in each origin.  When a
source is differentiated at zero, the regular-origin direction
at a site is \(f(p_z){\cal L}_j[1]\), accompanied by its own
coefficient, and the remaining direction is \eqref{eq:local-formation-LB7}.
We require the following ordinary formation estimate on the
actual localized complex pairs:
\begin{equation}
 \begin{split}
 \|\dot{\mathbf V}_k[f]\|_\lambda
 \le C s_k\bigg[
   &\alpha_k^2\|\dot u_k[f]\|_\infty
       +\sum_{j\le k}\omega_{k+1-j}
                          A_j\|f\|_\infty\\
   &+\sum_{j\le k}\varrho_{k,j}^2
                            \|{\cal D}_j[f]\|_\kappa
       +\sum_{j\in{\cal J}_k}\varrho_{k,j}^2
                            \|\dot R_j[f]\|_\kappa\bigg].
 \end{split}\label{eq:local-formation-LB10}
\end{equation}
Here \(\omega_n\le C(1+n)^CL^{-\alpha_H n}\) with
\(\alpha_H>0\), and \({\cal J}_k\) contains the aged remainder
origins that actually enter the regular part.  Younger remainders
remain in the remainder class.  The factor \(s_k\) is the
fluctuation-amplitude/field-radius ratio.  In a formulation with
additional positive powers of \(\theta_k^{-1}\) in this factor,
they are bounded by a constant times \(s_k\) at weak coupling.
The current scalar-origin term in \eqref{eq:local-formation-LB10} follows more sharply
from its fluctuation-difference estimate
\(C\alpha_k\delta_k\|\dot u_k[f]\|_\infty\).
All substitutions in \eqref{eq:local-formation-LB10} must lie in their actual specified
cores; a norm inequality does not create those inclusions.
In particular, a proof of the pointed Hessian cancellation and
nonlinear window estimate on constrained critical backgrounds
alone does not prove \eqref{eq:local-formation-LB10} for independent potential/current pairs
used in localization.  Passing from those background estimates
to all the actual localized pairs remains part of this ordinary
formation input.

The inserted-potential estimate for the actual finite
Mayer/decoupling series, followed by its actual renewed bond
pointing, gives
\begin{equation}
 \|{\cal L}_{k+1}[f]\|_{\kappa+1}
                  \le C_P\|\dot{\mathbf V}_k[f]\|_\lambda.
 \label{eq:local-formation-LB11}
\end{equation}
Its proof uses the ordinary spare-potential slot.  The auxiliary
potential parameter is not a physical source in a history box.
The normalized Gaussian, all reference measures, and
\(\mathbf N^\circ,\mathbf G_\tau\) have zero source derivative.
These facts are essential to \eqref{eq:local-formation-LB10}--\eqref{eq:local-formation-LB11}.

The ordinary constant-source construction supplies numbers \(b_j'(0)\)
with \(\sum_j|b_j'(0)|\le B_0\), uniformly in the cutoff.
Its fixed tents and stage factors satisfy
\begin{equation}
 \dot u_k[f](p)=f(p)+\sum_{j\le k}b_j'(0)\varphi_j(p)
                    \sum_yh_y^{(j)}(p)f(p_y),\qquad
 \|\dot u_k[f]\|_\infty\le(1+B_0)\|f\|_\infty.
 \label{eq:local-formation-LB12}
\end{equation}
The last bound uses \(0\le\varphi_j\le1\), \(h_y^{(j)}\ge0\),
and \(\sum_yh_y^{(j)}=1\).
The coefficient of \({\cal L}_j[1]\) is exactly \(b_j'(0)\):
the free response and coefficient extractor are fixed and
complex linear.  Thus no coefficient of another member enters
\eqref{eq:local-formation-LB10}, even when \(f(p_z)\) varies with the site.

Finally suppose the ordinary remainder families are bounded by
\(C_R\theta_j^{-p}e^{-\kappa d_j(X)}\), with \(p>3\), on the fixed
physical source disc \(|t|<r/\|f\|_{\sharp,j}\) around zero.
Cauchy's inequality at half that radius proves
\begin{equation}
 \|\dot R_j[f]\|_\kappa
       \le (2C_R/r)\theta_j^{-p}\|f\|_{\sharp,j}.
 \label{eq:local-formation-LB13}
\end{equation}
This is the only source-disc estimate needed for the small
remainder class.  It is not used to produce an \(\alpha_j^2\)
factor for regular activities.

\paragraph*{The regular family in a history.}
Let \(\Lambda\) be the small-field region of the intermediate
label and \(\Lambda^0\) its set of admitted base points.
A pointed term at \((X,z)\) is filed in the regular
family when \(X\subset\Lambda\), \(z\in\Lambda^0\) and all
its vertices are regular \cite[\S3]{B88b}. The fluctuation expansion places a
fine vertex with domain \(Y\) in this term only when the hull
of \(Y^\sim\) in the level-\((k+1)\) grid lies in \(X\).
Thus its complete determining set is interior. The boundary,
standing-operation and exterior-quadratic vertices vanish there,
and the Gaussian has no exterior mean. The local Gaussian and
threshold structure is that of the whole-lattice rule. At a
fixed local field the normaliser, free potential and Gaussian
ratio use the unsourced parameters \((\theta_k,x_k)\).
The physical source enters the local potentials.

The quadratic gauge correction also has a regular part and
a rim part. In the free fluctuation coordinates let \(K_A\)
and \(K_G\) be the Wilson and gauge-fixing forms at the
reference weight. With \(\tau_k=x_k/\theta_k\), the actual
precision is \(\tau_kK_A+K_G\). The reference Gaussian of
precision \(\tau_k(K_A+K_G)\) therefore requires the
potential \((\tau_k-1)B^{\mathsf T}K_GB/2\), as in
\eqref{agm:complete-tau-quadratic}. The gauge coordinates
\eqref{agm:contour-gauge-coordinate} use only bonds inside
their blocks. Thus this form splits into fine-cube terms
\((\tau_k-1)q_Y(B)\), where \(q_Y\) is the sum of the
local gauge squares assigned to \(Y\). The field \(B\)
includes the stored exterior coordinates. Completing the
square for the reference precision puts its mixed block in
the conditional mean. Substitute the shifted integration
coordinate into the entire \(q_Y(B)\), including its mixed
terms. A physical source in this presentation enters the
Wilson potential at fixed \(\tau_k\).

When the coarse hull of \(Y\) with its neighbors lies in
\(\Lambda\), use the relative quadratic bound of the
ordinary Gaussian leg. In the remaining cubes use the
bounded rim vertex. On the real cut the local row bound
and the number of bonds in a fine cube give
\[
 |(\tau_k-1)q_Y(B)|\le
 C|1-\tau_k|M^4\mathfrak b_k^2 .
\]
The centered relative quadratic estimate allows a fixed
disc about \(\tau=1\). At the rim use the disc of radius
\(2(\mathsf t_k+\bar c)/\theta_k\), where \(\mathsf t_k\)
is the real weight tolerance and \(\bar c\) the bounded
weight displacement allowed by the tube. The choices of
the tube pay for \(\mathfrak b_k^2\), also on the enlarged
field box used in the rim estimates. The two
assignments sum to the original quadratic correction,
and the choice between them reads the collared hull.
Every quadratic vertex in a regular output satisfies the
interior hull test, so it uses exactly the ordinary
relative quadratic term. The same test applies in the
ordinary recent-remainder contribution. Rim vertices enter
the boundary expansion with their own bounds. The global
cancellation of the marginal coefficient is the uncut
Gaussian identity \eqref{agm:compensating-Gaussian-integral}.

An old remainder of level \(j\) is aged here when
\(M\check R_j\le L^{k+1-j}\). It can enter a regular
term and then contributes to the marked part of that term.
Its origin carries the squared-radius weight
\(\varrho_{k,j}^2\). The weight \(\omega_{k+1-j}\) applies
to regular origins with their own pointed subtraction.
Terms containing a recent remainder go to the new remainder
on the one-layer inset of \(\Lambda\) and to the boundary
family elsewhere. Terms with all vertices regular that are
based at \(z\notin\Lambda^0\) use the edge rule.

On the complete reads of a regular term, the sourced array
uses the regular member for the constant \(w(p_y)\) at each
old base point \(y\), its own Wilson coefficient, and the
defect at \(y\). The Wilson shift is the actual plaquette
array with the stage factors and tents. The summable slopes of
the scalar increments and the local defect and remainder bounds
put this array in the ordinary source box about the same real
centre. All its interior potentials, including those for a
spatially varying shift, use the same local rule on every label.
Therefore the term filed at \((X,z)\) is the same function of
the same local array as in the whole-lattice rule. Its field
argument can vary with the history. For constant sources the
defects vanish and the same pointing gives the Euclidean family
with its source-dependent Wilson coefficient. A constant shift
\(c\) is included as a Wilson-weight vertex of the rule at
\((\theta_k,x_k)\).

To compare two arrays at this centre, join them inside the
ordinary box. The difference of the local vertex potentials
is linear in the input differences, with the four weights in
\eqref{eq:local-formation-LB10}, before differentiation. In
the convergent fluctuation expansion, differentiate along the
line by inserting that difference at one vertex. The Gaussian
terms are fixed, and the inserted-potential bound and the renewed
bond pointing give the output rate \(\kappa_{\mathbf T}\),
where \(\kappa_{\mathbf T}=(1+4\beta_{\rm s})\kappa\ge\kappa+1\).
Equivalently the finite difference is bounded by
\(C\vartheta_k^\sharp e^{-\kappa_{\mathbf T}d_{k+1}(X)}\)
times the weighted distance of the two arrays, where
\[
 \vartheta_k^\sharp=s_k+CM^4\theta_k^{-1/2}\mathfrak b_k^3
                          +\theta_k^{-1/4}\le C s_k
\]
at weak coupling. Apply this comparison to \(w=tf\) and to
zero. The difference quotients of the four inputs give the
actual shift derivative, the frozen regular derivative with
its own coefficient, the defect derivative, and the aged
remainder derivative. This gives the input to
\eqref{eq:local-formation-LB15} on the common regular family
once the vertex estimates on the assigned complex domains hold.

%
%
%
\paragraph*{The absolute bound needed on a source data line.}
\hypertarget{source:ordinary-absolute-line}{}
In the ordinary regular region \(X\subset\Lambda\), retain
the Gaussian compensation, cutoff and interaction legs
in that order. Use the same pointing for the step and
its constant-source companion. The comparison of two
arrays at fixed Gaussian proves a difference bound for
their interaction legs. To recover the absolute bound
first sum the interaction vertices on an admitted line.
For the regular, defect and aged-remainder norms in
\eqref{wf:formation-bound}, the enlarged box gives,
respectively, a fixed bound, \(CL^{-(k+1-j)}\), and
\(C\theta_j^{-p}\). The three sums are therefore bounded
by constants times
\[
 \sum_{a\ge0}\omega_{a+1},\qquad
 \sum_{a\ge0}(1+a)L^{-(a+1)},\qquad
 \sum_{a\ge0}(1+a)(\theta_k+ca)^{-p}.
\]
The first uses the positive H\"older gain and the third
is \(O(\theta_k^{2-p})\). In the Wilson direction the
weight norm is at most \(3r\); its vertex costs at most
\(CM^4s_k\alpha_k^2\,3r\).
Add the intrinsic and coarea vertex bounds on the
enlarged field box. They cost at most
\(CM^4\theta_k^{-1/2}\mathfrak b_k^3\), using
\(\mathfrak b_k\ge1\).

The coordinate Cauchy bounds on the old and coarea
legs and the normalized intrinsic pointing in
\eqref{fm:constructed-pointing-integrability}
give these bounds for the integrated bond roots
in \eqref{sp:ordinary-leg-norm}. Insert each root
as in \eqref{sp:actual-pointed-leg}. The rooted
cutoff term in \eqref{fm:rooted-cutoff-logarithm}
contributes at most \(C e^{-c\mathfrak b_k^2}\),
on its own leg. Midpoint tents cost a bounded number
of bond roots at a pointing, so retain an output
rate \(\kappa'\ge\kappa+1\). The free quadratic
root \eqref{fm:gauge-quadratic-root} separately gives
\(C_{\mathbf M}|1-\tau_k|e^{-\kappa'd_{k+1}(X)}\).
Consequently the three legs, after subtracting their
common normalizer \(\mathbf N^\circ\), give
\[
 |E(X,\cdot,z)-\mathbf N^\circ(X,\cdot,z)|
 \le\left(C\vartheta_k^\sharp+
             C_{\mathbf M}\mathsf t_k/\theta_k\right)
          e^{-\kappa'd_{k+1}(X)} .
\]
Take the tube and weak-coupling ceiling so that the
parenthesis is at most \(1/2\). For the ordinary
regular origins use the
\hyperlink{wf:direct-window}{direct window construction}
on their complete bases throughout the enlarged
data line. The other formation rules need their
own local domains and bond-root estimates. Use the
identical local pointing for the regular member in
the sourced and unsourced arrays.

\paragraph*{The induction and its constants}

\begin{proposition}[Linear regular-response bootstrap]
\label{prop:linear-source-bootstrap}
Under \eqref{eq:local-formation-LB1}--\eqref{eq:local-formation-LB13}, suppose the initial regular list is empty.
There are \(M_0<\infty\) and \(s_\star>0\), independent of the
number of levels, cutoff and volume, such that \(s_j\le s_\star\)
for all \(j\) implies
\begin{equation}
 A_j\le M_0\alpha_{j-1}^2s_{j-1}\qquad(j\ge1).
 \label{eq:local-formation-LB14}
\end{equation}
\end{proposition}

\begin{proof}
Divide \eqref{eq:local-formation-LB10} by \(\|f\|_{\sharp,k+1}\).
Use \eqref{eq:local-formation-LB9} with \(m=k+1\), \eqref{eq:local-formation-LB11}--\eqref{eq:local-formation-LB13}, and the monotonicity
\(\|f\|_{\sharp,j}\le\|f\|_{\sharp,k+1}\).  This gives
\begin{equation}
 A_{k+1}\le C_1s_k\left[
 \begin{aligned}
 &(1+B_0)\alpha_k^2+
 \sum_{j\le k}\omega_{k+1-j}A_j\\
 &{}+C\sum_{j\le k}\varrho_{k,j}^2A_jL^{-(k+1-j)}
 +C_r\sum_{j\le k}\varrho_{k,j}^2\theta_j^{-p}
 \end{aligned}\right].
 \label{eq:local-formation-LB15}
\end{equation}
There is no sum over the number of source labels in this inequality.
The norms of linear directions have already combined them into \(f\).
The flow bound and \eqref{eq:local-formation-LB5} imply
\begin{equation}
 \sum_{j\le k}\varrho_{k,j}^2\theta_j^{-p}
       \le C\sum_{a\ge0}(1+a)(\theta_k+ca)^{-p}
       \le C_p\theta_k^{2-p}\le C_p\alpha_k^2 .
 \label{eq:local-formation-LB16}
\end{equation}
To check the middle bound, compare the sum with its first term and
the integrals of \((\theta+ct)^{-p}\) and
\(t(\theta+ct)^{-p}\) on \([0,\infty)\).
They are fixed constants times
\(\theta^{1-p}\) and \(\theta^{2-p}\), respectively; the latter
integral is finite exactly when \(p>2\).
The final bound uses \(p>3\), \(q\ge0\), and large \(\theta_k\).

Assume \eqref{eq:local-formation-LB14} through level \(k\).  Both sums involving \(A_j\)
in \eqref{eq:local-formation-LB15} are at most
\[
 C M_0\alpha_k^2s_k
       \left(\sum_{n\ge1}\omega_n+\sum_{n\ge1}L^{-n}\right)
\]
by \eqref{eq:local-formation-LB5}.  Consequently, for constants \(C_2,C_3\) independent
of \(k\),
\[
 A_{k+1}\le\alpha_k^2s_k(C_2+C_3M_0s_k).
\]
Choose \(M_0\ge2C_2\) and
\(s_\star\le(2C_3)^{-1}\).  The right-hand side is at most
\(M_0\alpha_k^2s_k\), as required.  At the first step there are
no old regular members, so the same argument starts the induction.
The direct bottom-level scalar potential, if it occupies a separate
initial slot, must satisfy the same fixed-reference fluctuation
estimate used in \eqref{eq:local-formation-LB10}.
\end{proof}

\paragraph*{A consequence on the near, non-wrapping window class}

Let the final fluctuation size be \(\delta_n=\alpha_ns_n\), and
consider only terms for which the actual local window chart, the
pointed affine comparison and its complex domain inclusions hold
without a field-independent wrapping or holonomy tail.
In addition, the axial and H\"older bounds of the nonlinear window
proof must hold for the \emph{complete} near background-plus-fluctuation
pair with output size \(a_{\rm out}\le C\delta_n\).
A bound on its fluctuation increment alone is insufficient:
an unrestricted background of size \(\alpha_n\) would require that
larger size in the radius ratio.
The intended near class is the one on which the actual small-field
tests impose a fixed multiple of \(\delta_n\) on the complete field.
Here the complete box includes the upper comb, every lower-row
double block and the cutoff frame. The retained real profile
in this class is bounded by $C\delta_n$ on that box and on
the enlarged neighborhoods used by its gauge clauses.
Then \eqref{mb:real-layer-profile}, with its local gauge
estimate and tower restrictions, allows the real upper
curvature radius $A_0=C_M\delta_n$. The real reference has
zero potential slot; for its perturbations use
$A_1=C_M\delta_n$ and the already required complex-domain
inclusions. Those inclusions, rather than the real profile
alone, admit the varied potential and independent current.
The unamplified
transported fluctuation is at most $C_R\delta_{n-1}\le
C\delta_n$, together with its first difference on the
settled box; the last inequality follows from the adjacent
flow bound. Apply \eqref{wf:whole-window-amplitude} and its
central axial estimate with
$k=n-1$. It gives $C_M\delta_n$ on the whole raw and axial
datum, including the remote rows, and $C_M\delta_nL^{j-n}$
on the central rows, absorbing one factor $L$ in the constant.
At depth $j=n$ use the upper datum itself. The far response
therefore has size $C_M\delta_n e^{-cL^{n-j}}$, with the same
homogeneous amplitude as the central estimate. This verifies
the remote-amplitude requirement for this complete near class.
It does not apply to a box whose tests control only an inner
read or to an independent physical correction of larger size.
For ordinary formation on the enlarged disc
$|\nu|<c_\nu/s_k$, the fluctuation bound in the inner
substitution is instead $C\alpha_k$. The same calculation
therefore supplies the required whole-datum and remote bound
$C_M\alpha_k$ on that disc. No $\delta_k$ bound on the
amplified ordinary disc is used.
Keep the true ratio \(\rho_j=\delta_n/\alpha_j\); it need not
exceed one. The actual window-domain input includes
\(\alpha_j/r_{B,j}\le C\), where \(r_{B,j}\) is the
independent-coordinate datum radius, and the absolute full-datum
inclusion of \eqref{pw:absolute-datum-ceiling}. The central
half-datum-disc condition holds uniformly at all ages: with
\(\sigma=L^{-(n-j)}\), the flow gives
\[
 \frac{C_{\rm ax}\delta_n\sigma}{r_{B,j}}
 \le C\rho_j\sigma
 \le C s_n(1+n-j)^{1/2}L^{-(n-j)}.
\]
Choose \(s_\star\) smaller by a fixed factor so that the right-hand
side stays below the required half-disc constant. The remote rows
are admitted by the separate fixed absolute ceiling, with no
age gain asserted for them. On the complete near class just
specified, the preceding whole-window calculation supplies
their proportional amplitude \(C\delta_n\). Their response is
therefore \(C\delta_n e^{-c/\sigma}\) in the pointed-window
estimate. This amplitude comes from the complete input and
\eqref{wf:whole-window-amplitude}; the fixed ceiling supplies
its domain. For other classes without that proportional estimate,
the remote response remains a separate tail outside the
homogeneous near contribution.
This uses the near-domain proof for \(\rho_j>0\) itself and avoids
the young-age core-only fallback, which would lose a \(\rho_j^2\)
gain when \(\rho_j<1\).
After the four-dimensional fine-site count, the
renormalized constant-member contribution at age \(a=n-j\)
is bounded by
\begin{equation}
 C A_j\|f\|_\infty
       \left[\left(\frac{\delta_n}{\alpha_j}\right)^2
                              L^{-\alpha_Ha}
            +\left(\frac{\delta_n}{\alpha_j}\right)^3
                              L^{-2a}\right].
 \label{eq:local-formation-LB17}
\end{equation}
The value at the identity and this member's own
\(b_j'(0)A(h_z,\cdot)\) are removed here.
By \eqref{eq:local-formation-LB14}--\eqref{eq:local-formation-LB5},
\(A_j/\alpha_j^2\le Cs_n\) and
\(\delta_n/\alpha_j\le Cs_n(1+a)^{1/2}\).
The sum of \eqref{eq:local-formation-LB17} is therefore at most
\(C\delta_n^2s_n\|f\|_\infty\).

For the defect, the vacuum-only window estimate instead gives
\begin{equation}
 C\|{\cal D}_j[f]\|_\kappa
                 \left(\frac{\delta_n}{\alpha_j}\right)^2
 \le C\delta_n^2s_nh_j\operatorname{Lip}f.
 \label{eq:local-formation-LB18}
\end{equation}
Since \(h_j=h_nL^{-(n-j)}\), its sum is at most
\(C\delta_n^2s_n\|f\|_{\sharp,n}\).
Thus these near regular and source-gradient terms together have
size \(C\delta_n^2s_n\|f\|_{\sharp,n}\) per final window unit.
With the radii in \eqref{eq:local-formation-LB3}, the prefactor is
\(\theta_n^{-1}(\log\theta_n)^{3p_0-q}\).
The gain is obtained because the newly generated derivative member
already has amplitude \(\alpha_j^2s_{j-1}\).

No field-independent far tail is included in \eqref{eq:local-formation-LB17}--\eqref{eq:local-formation-LB18}.
Such a tail cannot acquire a factor \(\rho_j^2\) merely from a
source derivative.  Recent large-field terms, marked remainder
windows and torus or hole windings still require their own
ordinary source-disc bounds, domain comparisons or separately
small-class estimates.  Nor do the local implications in this
section verify the nonlinear formation hypotheses for every
history.  They show exactly how the \(\alpha^2\) normalization and
the source-gradient gain propagate once those ordinary local
hypotheses and the finite pointed estimate hold.


\paragraph*{The formation interface and the marked alternatives}

For the fixed window-first rule, the regular terms of \eqref{eq:local-formation-LB10}
are the linear application of \eqref{eq:local-formation-WF14}--\eqref{wf:formation-bound}.  The regular
direction is the source derivative of its near member with
its own coefficient, as in
\eqref{fm:own-counterterm-first-jet}; the defect is \eqref{eq:local-formation-LB7}.
The common-background and common-base rule in \eqref{eq:local-formation-WF13a}--\eqref{eq:local-formation-WF13b}
is kept while taking that derivative.  The bare scalar
direction is \eqref{eq:local-formation-SI13}, whose \(\alpha_k\delta_k=\alpha_k^2s_k\)
factor was obtained before the finite gas expansion.

For the regular near class the small physical radius of
the complete field is essential.  At age \(a=k-j\), use
the true ratio \(\delta_k/\alpha_j\) until the homogeneous
window estimate has been applied.  The complete
background-plus-fluctuation datum, rather than just the
increment, obeys that small-radius bound.  The all-age
half-disc admission in \eqref{eq:local-formation-LB17} then permits the
quadratic and cubic gains for ratios less than one.
Replacing that ratio by one before estimating the
window would lose the desired source normalization.
The full-radius amplification and its subsequent
Schwarz factor are the ones proved in the window
formation argument.

The vacuum-only defect terms use the quadratic
window estimate without a scalar-action cancellation.
Marked remainder and remote classes require a
different bound on their actual source domain.
One sufficient finite condition is joint holomorphy
and the ordinary estimate
\begin{equation}
 \sup_{\substack{|w|\|f\|_{\sharp,j}<r_0\\P\ {\rm admitted}}}
       |R_j^{wf}(X,P,z)|
       \le B_j e^{-\kappa d_j(X)}
 \quad\Longrightarrow\quad
 \|\dot R_j[f]\|_\kappa
       \le {2B_j\over r_0}\|f\|_{\sharp,j}.
 \label{fm:marked-fixed-disc-Cauchy}
\end{equation}
Here \(r_0\) is an actually admitted source radius
independent of \(X,P,z\); it is not inferred from a
real-source estimate.  The norm here is the pointwise
exponentially weighted member norm
\(\sup_{X,P,z}e^{\kappa d_j(X)}
|\dot R_j[f](X,P,z)|\), not a rooted domain sum
at that same exponent. Cauchy's formula on a circle
of radius \(r_0/(2\|f\|_{\sharp,j})\) proves the
implication.  The reference frame, source rule,
cutoff and read domain remain fixed on that circle.
If the remainder includes a history operation,
the corresponding bound is needed in its complete
kernel norm, with its inner Wilson weight; a bound
at one read datum is insufficient.

There is no extra \(\alpha_j^2\) on the right of
\eqref{fm:marked-fixed-disc-Cauchy}.  Any power or
logarithmic loss in using that bound must be
carried into the marked age sum.  For example,
if a later covariance calculation has a geometric
volume cost \(L^{8a}\), the elementary sufficient
rarity bound is
\begin{equation}
 q_a\le C\varepsilon_{\rm bad}\,q^a,\quad
        L^8q<1
 \quad\Longrightarrow\quad
       \sum_{a\ge0}L^{8a}q_a
          \le {C\varepsilon_{\rm bad}\over1-L^8q}.
 \label{fm:marked-age-volume-sum}
\end{equation}
The geometric sum is immediate; the point is its
actual weight and quantifiers.  An unweighted
small bad-class probability does not supply this
age moment.  A differentiated one-zone mass test
does not supply it either.  It is a separate
undifferentiated estimate on the allowed classes.

Remote and wrapping terms whose bound is independent
of the local field cannot be assigned an invented
factor \((\delta_k/\alpha_j)^2\).  They must either
be excluded by an actual non-wrapping near chart
or retained with their own fixed-disc and age
estimate.  Thus the source recurrence proved above
has a concrete formation route under its map
hypotheses, while its marked and remote alternatives
remain separate finite estimates to verify.


\subsubsection{The old-origin Hessian on the larger run window}


This section proves a conditional ordinary estimate.  Its inputs are
the complex domains and first-response estimates of the actual background
maps, the pointed quadratic cancellation, and the window formation rule.
The desired Hessian estimate is not one of the inputs.  In particular, an
estimate on the Gaussian cutoff box alone is not used as an analytic
continuation theorem.  We construct the larger fluctuation domain from
the unscaled constraint chart and the datum domains of the background
maps.

The conclusion concerns the regular old-origin part of a step exponent.
The conditional exterior-force estimate, boundary origins with stored
fields, and the identification of a clean-class conditional measure are
separate issues.

\paragraph*{Physical and Gaussian coordinates}

Fix a step from level \(k\) to level \(k+1\).  Write
\begin{equation}
 g_k=\theta_k^{-1/2},\qquad
 \alpha_i\asymp\theta_i^{-1/2}(\log\theta_i)^q,
 \qquad \mathfrak b_k\asymp(\log\theta_k)^p,
 \qquad \delta_k=g_k\mathfrak b_k .
 \label{eq:local-formation-OA1}
\end{equation}
Here \(q>p\), and \(\alpha_i\) is the smaller of the usable physical
radii in the datum, potential, curvature and current conditions.
Their fixed ratios are included in the constants.  A finite alteration
of an index in \(\alpha_i\), as needed for adjacent input layers, has the
same meaning.  We assume the ordinary flow gives
\begin{equation}
 \rho_i:=\max\{1,C\alpha_k/\alpha_i\}
       \le C_\rho(1+k-i)^{1/2},\qquad
 \sigma_i=L^{-(k-i)},\qquad i\le k .
 \label{eq:local-formation-OA2}
\end{equation}
Constants depend on the fixed block and collar choices, but not on the
number of old levels, the cutoff or the volume.  In particular,
\(C_{\rm a}:=\sup_{i\le k}\rho_i\sigma_i<\infty\).

The free Gaussian variable \(y_b\) has three real color coordinates at
each nonpivot bond of the level-\(k\) lattice.  The physical bond
logarithm is of size \(g_k y_b\).  Thus a physical analytic radius
\(\alpha_k\) corresponds to a Gaussian radius of order
\(\alpha_k/g_k\), rather than to \(\alpha_k\).
For complex variables we use the Euclidean norm of the three complex
components at a bond and its supremum over bonds.

Let \(I\) be the finite box of free bonds under consideration.  Its
outside free coordinates and other exterior data are fixed.  Every
real free coordinate appearing in the regular fields is bounded by
\(C_{\rm w}\mathfrak b_k\).  The number of coordinates satisfies
\begin{equation}
 m_I=3\#I\le C(rL^{-a_*}+D)^4\le C(r+D)^4.
 \label{eq:local-formation-OA3}
\end{equation}
This count will not enter the localized estimate below.  It is relevant
to the weaker estimate available without spatial locality.

\paragraph*{The primitive analytic inputs}

We specify the domain inputs used in the proof.  These are statements
about assigned maps and their first responses, not bounds on the
Hessian of an old activity.

For the actual one-step average, at its actual background
\(V^{\rm bg}(W)\), put
\begin{equation}
 q_W(a)=\frac1i\log\!\left[
       \mathcal M(e^{ia}V^{\rm bg}(W))W^{-1}\right]
        =LQ_Wa+\mathcal R_W(a).
 \label{eq:local-formation-OA4}
\end{equation}
The logarithm is taken on its assigned local branch.  The fixed-pivot
construction supplies linear maps \(\mathsf C_W\) and \(h_W\) with
\begin{equation}
 LQ_W\mathsf C_W=0,\qquad LQ_Wh_W=1,
 \qquad\|\mathsf C_W\|_{\infty\to\infty}
              +\|h_W\|_{\infty\to\infty}\le C_L .
 \label{eq:local-formation-OA5}
\end{equation}
The map \(\mathsf C_W\) is the pivot elimination map, not a Gaussian
covariance.  The maps in \eqref{eq:local-formation-OA4}--\eqref{eq:local-formation-OA5} are the ones used in the run.
Uniformly on an unscaled physical ball \(\|a\|_\infty<r_0\), the
remainder is holomorphic and obeys
\begin{equation}
 \|\mathcal R_W(a)\|_\infty\le C_0\|a\|_\infty^2,
 \qquad
 \|D\mathcal R_W(a)\|_{\infty\to\infty}
                                      \le C_0\|a\|_\infty .
 \label{eq:local-formation-OA6}
\end{equation}
Here \(r_0,C_0,C_L\) are independent of the box size.  The same bounds,
with fixed enlarged constants, apply to the assigned local charts on
the independent fine-parameter discs \(|s_\Delta|<e^{\bar\kappa}\).
At physical parameters the local charts agree with the run chart by
their restriction and completeness identities.

The ordinary background package supplies holomorphic maps on their
declared datum balls and their declared core slices.  The datum radii
satisfy \(r_{B,i}\ge c\alpha_i\).  At zero fluctuation all substituted
data and pairs have a fixed positive fraction \(\epsilon_0\) of room
in these domains.  This room includes the tower conditions, node
representatives, current bounds, and the inner, outer and removed
substitutions.  These are local complex assertions; they do not assert
the existence of a global Gaussian measure for independent currents.

We use the following precise first-response form of that package.
Measure each defining datum or core displacement at old level \(i\)
in units of its allowed radius.  Along an admitted physical perturbation
of supremum size \(u\), its derivative in any of these normalized
quantities is bounded by
\begin{equation}
 C_R\sigma_i\alpha_i^{-1}u.
 \label{eq:local-formation-OA7}
\end{equation}
The bound holds throughout the declared complex domains, for all the
fine parameters above.  It is the integrated first-response and
nonlinear-average estimate: the actual average, not its differential,
is used to construct the datum.  The constant \(C_R\) includes the
fixed enlarged parameter discs and the chart derivative.  For outer
and removed derivatives there are additionally their assigned
exponential distance factors.  As usual, a compatible boundary
comparison and the weighted H\"older response estimates are assumed
for the actual box backgrounds.

Finally, the field maps are hereditary on active wall components.
After the parameters outside a domain \(Y\) are set to zero, a local
potential uses only the free coordinates in \(Y\).  A fixed collar can
instead be included in \(Y\) throughout; it changes the constants
below.  The number of domains of tree length in \([m,m+1)\) containing
a fixed cube is at most \(C_De^{\kappa_Dm}\), and
\begin{equation}
 \operatorname{diam}Y\le C_M(1+d_k(Y)).
 \label{eq:local-formation-OA8}
\end{equation}
The spatial rates have the strict reserves needed to sum these
domains and the fixed polynomial diameter factors.

\paragraph*{The actual chart on the larger ball}

For clarity we give the contraction argument behind the radius claim.
Put \(z=g_k\mathsf C_Wy\).  The nonlinear correction is the solution
of
\begin{equation}
 d=\mathcal R_W(z-h_Wd),\qquad
 a_W(y)=z-h_Wd,
 \qquad V_W(y)=e^{ia_W(y)}V^{\rm bg}(W).
 \label{eq:local-formation-OA9}
\end{equation}
Choose \(r_1>0\) so small that
\(2r_1<r_0\) and \(4C_0C_Lr_1<1/2\).
For \(\|z\|_\infty\le r_1\), the right side maps the ball
\(\|d\|_\infty\le4C_0\|z\|_\infty^2\) to itself.
Indeed \(\|z-h_Wd\|_\infty\le2\|z\|_\infty\), and \eqref{eq:local-formation-OA6}
then gives that bound.  Its Lipschitz constant in \(d\) is at most
\(2C_0C_Lr_1<1/2\).  Iteration from zero therefore converges uniformly
to the unique solution, including locally uniformly in all complex
parameters.  Uniform limits of the holomorphic iterates are
holomorphic.  Differentiating the fixed-point equation, or its
uniformly convergent derivative iteration on a smaller ball, gives
\begin{equation}
 \|a_W(y)\|_\infty\le Cg_k\|y\|_\infty,
 \qquad \|D_ya_W(y)\|_{\infty\to\infty}\le Cg_k .
 \label{eq:local-formation-OA10}
\end{equation}
Moreover, \(q_W(a_W(y))=0\), because \(LQ_Wz=0\) and
\(LQ_Wh_Wd=d\).  This is the actual constraint, with its nonlinear
average.  Uniqueness shows that this chart agrees on the operational
cutoff box with the chart already used in the run.

There is consequently a fixed \(c_*>0\) for which all the old
substitutions are defined and holomorphic on
\begin{equation}
 \|y\|_\infty<R_k,
 \qquad R_k=c_*\alpha_k/g_k .
 \label{eq:local-formation-OA11}
\end{equation}
Here and below fixed exterior coordinates may be retained with norm
at most \(C_{\rm w}\mathfrak b_k\).  To prove the assertion, first make
\(\alpha_k\) small enough that \eqref{eq:local-formation-OA10} stays in its physical chart.
For any fixed complex \(y\) in \eqref{eq:local-formation-OA11}, follow the segment \(ty\),
\(0\le t\le1\).  On the portion where a composition of assigned
background maps is admitted, \eqref{eq:local-formation-OA7} and \eqref{eq:local-formation-OA10} bound each normalized
displacement by
\begin{equation}
 C_R g_k\|y\|_\infty\sigma_i/\alpha_i
       \le C_R c_*\rho_i\sigma_i
       \le C_R c_*C_{\rm a}.
 \label{eq:local-formation-OA12}
\end{equation}
Choose \(c_*\) so that the last expression is less than
\(\epsilon_0/4\).  The segment cannot reach a domain boundary:
every defining inequality still has at least half of its initial
room.  Atomic background maps are already given on those domains, so
this is an inclusion argument for a composition of assigned
holomorphic maps, not an assertion that a branch continues through
an unknown singularity.  The same argument gives the required joint
holomorphy with the fine parameters.  It also applies to local
independent-current formulas without postulating an off-graph global
Gaussian integral.

The window multiplier used to differentiate an inner direction can
range on a fixed disc \(|\nu|<c_1/c_*\), where \(c_1\) is chosen from
the same room estimates and \(c_*<c_1/4\).  In fact its normalized
increment is bounded by \(C_R|\nu|c_*\rho_i\sigma_i\).
The unamplified part costs only \(C_Rc_*C_{\rm a}\); choose \(c_1\)
so that their sum still spends less than half the room.  Thus this
disc contains \(\nu=1\) and permits the window derivative at zero.

Since \(q>p\), increasing the coupling floor gives
\(C_{\rm w}\mathfrak b_k<R_k/4\).  Two coordinate Cauchy circles
with radii \(cR_k\), centered anywhere on the real product window,
therefore fit in \eqref{eq:local-formation-OA11}.  There is no factor depending on \(m_I\)
in this inclusion.  No indicator function or cutoff is analytically
continued.  Only the activity functions are evaluated on the larger
ball; the real measure, Gaussian scaling and all numerical choices
of the run remain fixed.

\paragraph*{An amplitude bound before the Gaussian-growth norm}

Let the regular member at old level \(i\) have norm \(e_i\), with
\(e_i\le E_0\).  The norms include its infinite companion if one is
required.  Each pointed total has local gauge invariance, constant
color invariance and the full site symmetry needed for the pointed
Hessian theorem.  At every old base point subtract its vacuum and
its own extracted coefficient times its own tent action.  The tent
has unit quadrature; the action convention is
\(A(e^{iB})=\tfrac12|\mathrm{d}B|^2+O(B^3)\).
This is the regular part of the decomposition: its
counterterms have been transferred to the running Wilson coefficient.
A profile deficit is retained as a lag or boundary origin, rather
than silently subtracted from the measure.  The clean interior has
no such deficit in the canonical decomposition with redrawn profiles.
Thus the affine quadratic term cancels with coefficient one in that
normalization.  We do not replace the pointed object by a bulk
convolution or assign an action coefficient to an individual domain.

We recall explicitly the value estimate supplied by this cancellation.
Write \(D_X=2+\operatorname{diam}_iX\).  On an inner box background
the affine part and its error have sizes bounded, with fixed powers
of \(D_X\), by
\begin{equation}
 \alpha_k\sigma_i^2,
 \qquad
 \alpha_k\sigma_i^{2+\alpha_H}
                +r_{B,i}^{-1}\alpha_k^2\sigma_i^4,
 \qquad 0<\alpha_H\le1.
 \label{eq:local-formation-OA13}
\end{equation}
These estimates concern the complete substituted field, not just its
fluctuation increment.  They hold on the whole disc just constructed.
Second and third derivatives of a member cost
\(Ce_i\alpha_i^{-2}\) and \(Ce_i\alpha_i^{-3}\), with its
original tree decay.  Its linear derivative at the trivial pair is
zero by constant color invariance.

Set \(T=C D_X^s\rho_i\sigma_i^2\).  When \(T\) is small, Taylor's
formula, the cancellation of the affine quadratic value on the
pointed total, and \(r_{B,i}\ge c\alpha_i\) bound the remaining
terms by
\(Ce_iD_X^{s'}[\rho_i^2\sigma_i^{4+\alpha_H}
 +\rho_i^3\sigma_i^6]\).
For example, the product of the affine part and its error, divided
by \(\alpha_i^2\), gives exactly these two powers.  Squares of the
error and cubic terms are bounded by the same expression, since
their normalized sizes are bounded by \(C(T\sigma_i^{\alpha_H}+T^2)\)
and \(CT^3\).  If \(T\) is not small, whole-core bounds for the
member and its own action, and the comparison affine quadratic
value, give \(Ce_i(1+T^2)\le Ce_i C_TT^3\).  The fixed powers
of \(D_X\) are summable against the tree decay. A fixed finite set
of young ages failing the central half-datum-ball test is also
covered by the whole-core bound: their \(\sigma_i\) have a fixed
positive lower bound. The separate remote-datum admission is
retained. This is precisely the reason to use the ordinary
full-radius ratio \(\rho_i\ge1\) here.

At one base point, after its domain sum, the bound is therefore
\begin{equation}
 Ce_i\left[\rho_i^2\sigma_i^{4+\alpha_H}
                         +\rho_i^3\sigma_i^6\right].
 \label{eq:local-formation-OA14}
\end{equation}
A current window contains at most \(C\sigma_i^{-4}\) old base
points.  Its value bound is at most
\(Ce_i\rho_i^3\sigma_i^{\alpha_H}\).
Terms extending beyond that window retain an exponentially small
tree-length tail.  Outer and removed terms retain their actual
response factors.  Even the cruder amplitude
\(Ce_i\sigma_i^{-4}e^{-c/\sigma_i}\), when only a long-domain
value bound is used, is bounded by
\(Ce_i\rho_i^3\sigma_i^{\alpha_H}\).
This does not give such a tail a square of the actual small
curvature; it is a full analytic-radius estimate.

It is important where the inner estimate is localized.  At each
old base point form one window function from the near pointed sum,
its own action counterterm, and one common critical box map.  This
entire function has one fixed current-scale base, of uniformly
bounded size.  Only then apply the fine brackets of the window
formation rule.  Thus the cancellation in \eqref{eq:local-formation-OA14} occurs before a
single common bracket.  No sup norm of a sum is used to bound a
sum of differently based bracket operators.

On the independent fine discs the bracket with \(m\) differenced
cubes costs at most
\((e^{\bar\kappa}-1)^{-m}\), with no factorial.  The derivative
at \(\nu=0\) costs the reciprocal of its fixed disc, hence only
a fixed constant.  For outer origins retaining their original
base, split the old input rate into the domain-sum, long-tail
and retained parts.  If \(Y\) is the resulting current domain,
the hull and rooting geometry gives
\(d_k(Y)\le C_h(d_i(X)+m+1)\); for a bounded window base it gives
\(d_k(Y)\le C_h(m+1)\).
Choose the retained rates and \(\bar\kappa-1\) larger than
\(C_h(\lambda+\varepsilon)\), after spending the old rooted
entropy and diameter powers.  The remaining \(\varepsilon\)
absorbs the polynomial number of window anchors at a fixed \(Y\).
This proves the following pointwise analytic bound for the
actual chosen window filing:
\begin{equation}
 |V_{k,i}^{\rm old}(Y;y)|
       \le Ce_i\omega_{k+1-i}e^{-\lambda d_k(Y)},
 \qquad \|y\|_\infty<R_k,
 \quad
 \omega_{a+1}=C(1+a)^{3/2}L^{-\alpha_Ha}.
 \label{eq:local-formation-OA15}
\end{equation}
The current origin \(i=k\) uses its direct whole-core Cauchy
bound and \(\omega_1=O(1)\); it needs no age improvement.
The proof was carried out on the larger ball using a physical
fluctuation amplitude at most \(c_*\alpha_k\), rather than the
tested amplitude \(\delta_k\).  It does not use the final
potential bound with a Gaussian-growth factor
\(1+M^{-4}\sum_{b\subset Y}|y_b|^2\).  Using that weaker
envelope on a radius \(R_k\) would introduce \(R_k^2\) and
erase the desired Cauchy gain.

Let
\begin{equation}
 \mathcal E_k=\sum_{i\le k}\omega_{k+1-i}e_i
           \le C E_0 .
 \label{eq:local-formation-OA16}
\end{equation}
Absolute convergence is uniform on every smaller ball in \eqref{eq:local-formation-OA11},
by \eqref{eq:local-formation-OA15}.  It justifies summing old levels before taking two
coordinate derivatives.  All estimates can first be made with
finitely many old levels and on a finite torus.  Their constants
do not change in either limit.

\paragraph*{Cauchy, rooted counting and the weighted precision gap}

Write
\(P_{k,I}^{\rm old}(y)=\sum_Y\sum_{i\le k}V_{k,i}^{\rm old}(Y;y)\),
with exterior variables fixed.  Additive terms depending only on
the exterior can be omitted.  By two coordinate Cauchy formulas,
including the factor \(2!\) for a repeated coordinate, \eqref{eq:local-formation-OA15} gives
\begin{equation}
 \left|\partial_{y_b^a}\partial_{y_{b'}^{a'}}
            \sum_{i\le k}V_{k,i}^{\rm old}(Y;y)\right|
       \le C\mathcal E_k R_k^{-2}e^{-\lambda d_k(Y)}
 \label{eq:local-formation-OA17}
\end{equation}
throughout the real product window.  The derivative is zero unless
both bonds belong to the declared domain, including its fixed
collar.  For such a domain,
\(d(b,b')\le C_M(1+d_k(Y))\).
Choose \(\mu>0\) so small that
\(\lambda-\mu C_M>\kappa_D\), retaining a strict gap.  Root the
domains at a cube containing \(b\); the fixed number of possible
root cubes costs only a constant.  Then
\begin{equation}
 \sum_{Y\ni b,b'}e^{-\lambda d_k(Y)}
 \le C e^{-\mu d(b,b')}
       \sum_{m\ge0}e^{-(\lambda-\mu C_M-\kappa_D)m}
 \le C'e^{-\mu d(b,b')}.
 \label{eq:local-formation-OA18}
\end{equation}
Consequently the full old-activity Hessian, in the actual Gaussian
coordinates, satisfies
\begin{equation}
 \big\|D_bD_{b'}P_{k,I}^{\rm old}(y)\big\|
 \le C\mathcal E_k\left(\frac{g_k}{\alpha_k}\right)^2
                   e^{-\mu d(b,b')}
 \le C E_0(\log\theta_k)^{-2q}e^{-\mu d(b,b')} .
 \label{oa:localized-hessian}
\end{equation}
Color matrix norms cost a fixed factor only.  The count of bonds at
distance in \([r,r+1)\) is at most \(C(1+r)^3\).  Thus for every
\(0<\mu'<\mu\), both weighted row and column sums of the Hessian
are at most \(C\mathcal E_k(g_k/\alpha_k)^2\).  In particular,
Schur's inequality gives
\begin{equation}
 \|T_\phi D^2P_{k,I}^{\rm old}(y)T_\phi^{-1}\|_{2\to2}
       \le C\mathcal E_k(g_k/\alpha_k)^2,
 \qquad
 |\phi(b)-\phi(b')|\le\mu' d(b,b'),
 \label{eq:local-formation-OA20}
\end{equation}
where \(T_\phi\) multiplies the three coordinates at \(b\) by
\(e^{\phi(b)}\).  This estimate is uniform in \(I\).

For completeness, suppose the actual free precision \(A\), after
the prescribed constraint and gauge fixing, has a gap
\(A\ge a_-\mathbf1\) and uniform exponentially summable entries.
This is a separate free-background input, concerning the constrained
pullback Hessian in the coordinates \(\mathsf C_W\), not the ambient
curl Hessian.  Its Dirichlet compression \(A_{OO}\) has the same
gap.  For small Lipschitz weights its conjugation difference has
norm at most \(C\mu'\).  Hence the conjugated inverse of \(A_{OO}\)
has norm at most \(2/a_-\) when \(C\mu'<a_-/2\), by the Neumann
series.  Applying this to
\(\phi(b)=\mu' d(b,b')\) bounds each inverse entry by
\((2/a_-)e^{-\mu'd(b,b')}\).  This proof is uniform in \(O\).
It follows by two exponentially summable convolutions that
\[
 K_I=((A^{-1})_{II})^{-1}
       =A_{II}-A_{IO}A_{OO}^{-1}A_{OI}
\]
has exponentially summable entries at a possibly smaller rate.
Furthermore,
\(\langle v,K_Iv\rangle
 =\inf_w\langle(v,w),A(v,w)\rangle\ge a_-\|v\|^2\).
Thus the inverse marginal covariance, with its actual Schur
correction, retains the free gap.  Its conjugation difference
also has norm at most \(C\mu'\): use
\(|e^{\phi(b)-\phi(b')}-1|
 \le\mu'd(b,b')e^{\mu'd(b,b')}\) and the row and column sums.
Choose \(\mu'\) so that this norm is at most \(a_-/4\).
Then
\begin{equation}
 \operatorname{Re}\langle v,
      T_\phi(K_I-D^2P_{k,I}^{\rm old})T_\phi^{-1}v\rangle
 \ge\left(\frac{3a_-}{4}
       -C\mathcal E_k(g_k/\alpha_k)^2\right)\|v\|^2.
 \label{eq:local-formation-OA21}
\end{equation}
For a sufficiently large coupling floor the old activities cost
at most \(a_-/4\), leaving a weighted gap of at least \(a_-/2\).
Other intrinsic or remainder Hessians must be charged separately.
The reference precision is held fixed while differentiating in
\(y\); no derivative of a covariance with respect to \(y\) is
introduced.

\paragraph*{Remainders and the weaker finite-box alternative}

A holomorphic aged remainder with no fixed color spectators has
zero linear term at the trivial pair, but has no action
counterterm.  Its vacuum-only window bound has weight \(\rho_i^2\)
instead of \(\omega_{k+1-i}\).  The same proof gives \eqref{oa:localized-hessian} with
\(\mathcal E_k\) augmented by \(\sum_i\rho_i^2 r_i\), provided its
actual maps have the same domains and locality.  If
\(r_i\le C\theta_i^{-s}\), \(s>2\), and
\(\theta_i\ge\theta_k+c(k-i)-C_0\), use the actual running
radii to bound \(\rho_i^2\le C\theta_i/\theta_k\).
Indeed the squared radius ratio is bounded by this expression
times \((\log\theta_k/\log\theta_i)^{2q}\), which stays
bounded on the admitted flow. The maximum with one has
the same bound. Thus
\begin{equation}
 \sum_i\rho_i^2r_i\le C\theta_k^{-1}
       \sum_{a\ge0}(\theta_k+ca-C_0)^{1-s}
                  \le C_s\theta_k^{1-s} .
 \label{eq:local-formation-OA22}
\end{equation}
For the last inequality use the integral comparison with
\(\theta_k-C_0\ge\theta_k/2\).  This optional conclusion does not
cover a boundary term retaining stored color variables or a term
for which only an integrated norm is known.

If only a bound on the whole finite-box analytic function is
available, a useful but weaker statement remains.  Suppose
\(|P(y)|\le C E_0(r+D)^A\) on a Gaussian polydisc of radius
\(c\alpha_k/[g_k(r+D)^B]\).  The same coordinate Cauchy estimate
and \(m_I\le C(r+D)^4\) give
\begin{equation}
 \|D^2P\|_{2\to2}
       \le C E_0(r+D)^{A+2B+4}(g_k/\alpha_k)^2 .
 \label{eq:local-formation-OA23}
\end{equation}
Indeed each matrix entry has the bound without the last power
\((r+D)^4\), and the row and column contain at most \(m_I\)
entries.  For a direct volume bound \(A=4,B=0\), this is
\(CE_0(r+D)^8(\log\theta_k)^{-2q}\).  When
\(r+D\le C(\log\theta_k)^{2p}\), it tends to zero if \(q>8p\).
It preserves an unweighted gap.  It does not imply \eqref{eq:local-formation-OA20} at a
fixed positive spatial weight: conjugation could cost
\(e^{\mu'\operatorname{diam}I}\).  The rooted localization in
\eqref{eq:local-formation-OA15}--\eqref{eq:local-formation-OA18} proves the stronger weighted conclusion;
counting window coordinates alone would not suffice.

\paragraph*{The old-activity Hessian estimate and its scope}

Under the stated ordinary analytic map hypotheses and the reviewed
pointed and window inputs, the old regular-activity contribution to
the localized Hessian is bounded by
\(CE_0(\log\theta_k)^{-2q}\), uniformly over old ages and box sizes.
The analytic radius is in the physical-to-Gaussian coordinate
conversion in \eqref{eq:local-formation-OA11}.  A bare operational cutoff domain, or only a
Gaussian-growth potential bound on that domain, would not prove it.
The two required primitive verifications are the uniform unscaled
pivot estimates in \eqref{eq:local-formation-OA5}--\eqref{eq:local-formation-OA6} and the all-domain first-response and
membership estimates in \eqref{eq:local-formation-OA7}, including towers and outer substitutions.
They belong to the actual ordinary background analysis and have not
been proved for the Yang--Mills maps merely by this deduction.

The statement has no automatic consequence for the change in force
when exterior data are changed.  For that one needs, separately,
uniform estimates such as
\(\|D_b(P_e-P_{e_0})\|\le C\mathfrak b_k e^{-c d(b,O)}\)
on the actual compatible exterior and stored data.  A Hessian bound
in the variables of \(I\) does not bound derivatives with respect
to the exterior parameters.  Nor does it identify a clean-class
restriction with a Gaussian fibre kernel.  Field-independent
remote or holonomy bounds without the analytic domain assertion
used above remain separate tails; they receive neither an invented
curvature factor nor an unjustified Cauchy estimate.


\paragraph*{An analytic criterion for the old exterior force}

The larger fluctuation ball controls derivatives in its free
coordinates.  A useful further reduction states precisely what
function-level locality would also control a change of exterior
datum.  This reduction concerns regular old members, not a
conditional class probability or a boundary measure.

Let \(O_{e,Y}(y)\) be the actual old-origin member on a fixed
run window, where \(e\) is a compatible exterior datum.
Assume both \(e\) and a comparison datum \(e_0\) give admitted
functions on the same full complex ball
\(\|y\|_\infty<R\), with the same local filing.
Let \(O\) be the exterior region where they differ.
Suppose the actual background/boundary comparison yields
the amplitude difference
\begin{equation}
 |O_{e,Y}(y)-O_{e_0,Y}(y)|
       \le\varepsilon_\partial e^{-\lambda d(Y)}
                          e^{-c\,{\rm dist}(Y,O)}
              \quad(\|y\|_\infty<R).
 \label{fm:exterior-amplitude-premise}
\end{equation}
This is a uniform function bound on the entire analytic
ball.  A comparison on the real cutoff box alone is not
this premise.

Assume the actual supports have diameter at most
\(C_h(1+d(Y))\), their free-coordinate counts are at most
\(C_h(1+d(Y))^q\), and the number containing a fixed
coordinate with \(m\le d(Y)<m+1\) is at most
\(C_De^{\kappa_Dm}\).  When \(b\) belongs to \(Y\),
\begin{equation}
 {\rm dist}(b,O)
       \le C_h(1+d(Y))+{\rm dist}(Y,O).
 \label{fm:exterior-root-distance}
\end{equation}
Choose \(0<c'<c\), \(c''>0\) and a strict \(\eta>0\)
so that
\(\lambda-C_h(c'+c'')>\kappa_D+\eta\).
Such choices exist when \(\lambda>\kappa_D\).
On \(\|y\|_\infty<R/2\), coordinate Cauchy in the
common larger ball gives
\begin{equation}
 \left|D_b\sum_Y
            [O_{e,Y}-O_{e_0,Y}](y)\right|
       \le {C\varepsilon_\partial\over R}
                              e^{-c'{\rm dist}(b,O)} .
 \label{fm:exterior-force-reduction}
\end{equation}
Indeed only supports containing \(b\) contribute.
For each such support Cauchy costs \(2/R\).
Using \eqref{fm:exterior-root-distance}, its exponential
weight is at most
\(e^{c'C_h}e^{-(\lambda-c'C_h)d(Y)}
e^{-c'{\rm dist}(b,O)}\).
The rooted count is then a convergent geometric sum.
The estimates are uniform in finite volume, so the
absolutely convergent sum and its derivative have the
same limit on every smaller ball.

Two coordinate derivatives, with the same reasoning,
also give the weighted row bound
\begin{equation}
 \sum_{b'}e^{c''{\rm dist}(b,b')}
 \left|D_bD_{b'}\sum_Y
            [O_{e,Y}-O_{e_0,Y}](y)\right|
       \le {C\varepsilon_\partial\over R^2}
                              e^{-c'{\rm dist}(b,O)} .
 \label{fm:exterior-Hessian-reduction}
\end{equation}
For coincident coordinates the second Cauchy estimate
has the factor \(2!\); for different coordinates use
two independent discs.  Both are bounded by a fixed
constant times \(R^{-2}\).  For a given \(Y\), summing
\(b'\) costs at most its polynomial coordinate count
times \(e^{c''C_h(1+d(Y))}\).  The strict rate inequality
above leaves \(e^{-\eta d(Y)}\), which absorbs that
polynomial.  This proves the asserted row bound, with
constants independent of the number of window variables.

For the run-window radius \(R=c\alpha_k/g_k\), these
are explicit Cauchy costs in that same radius.
If only \(R/r^4\) has been admitted, the costs become
\(r^4/R\) and \(r^8/R^2\); they are not discarded.
The remaining model question for
\eqref{fm:exterior-force-reduction} is the actual
common-domain amplitude comparison in \eqref{fm:exterior-amplitude-premise}, with its
compatible exterior, stored fields, gauges and
relative boundary data.  The interior Hessian
estimate alone does not furnish that comparison.
Even this regular old-origin force bound would
leave the other terms of a clean conditional
exponent, its boundary contributions and its
normalized class law to be identified.



\subsubsection{Finite histories and the local formation norm}

The local estimates above are conditional on all admitted stored and
boundary data.  A finite history operation is an additional integral,
not an extra name for the local Gaussian functional.  We retain its
read variables and its inner Wilson weight explicitly.

Let \(\omega\) range over the discrete finite history records.  Write
\(\xi\) for the stored variables integrated by this history operation
and \(P\) for the remaining admitted read data.  Suppose the operation
has been expressed over one
parameter-independent positive reference measure \(d\mu_\omega(\xi)\).
Its complex integrand is \(K_\omega(\xi;P)\).  If a source derivative
also inserts a local function \(A_\omega(\xi;P)\), the relevant estimate
is on the product \(K_\omega A_\omega\).  One cannot first replace
\(K_\omega\) by its partition function and then take a supremum of
\(A_\omega\) on different fibres.

For \(c\ge0\), fixed history volume \(v_\omega\), and fixed read cores
\({\cal C}_\omega\), define
\begin{equation}
 \|K\|_c=
 \sum_\omega\int^{*}d\mu_\omega(\xi)\,
       e^{cv_\omega}
       \sup_{P\in{\cal C}_\omega}|K_\omega(\xi;P)|.
 \label{fm:history-local-norm}
\end{equation}
The upper integral permits a supremum whose measurability has not been
established; the same core and reference operation are used on both
sides of every estimate.  If \(M_\omega(\xi)\) is a common positive
envelope and
\begin{equation}
 \sup_{P\in{\cal C}_\omega}
       |K_\omega(\xi;P)A_\omega(\xi;P)|
 \le M_\omega(\xi)\,\varepsilon
       e^{-\lambda d(X)},\qquad
 \sum_\omega\int d\mu_\omega\,e^{cv_\omega}M_\omega\le C_H,
 \label{fm:history-envelope-transfer}
\end{equation}
then the norm of the localized insertion is at most
\(C_H\varepsilon e^{-\lambda d(X)}\).  This follows directly from
monotonicity of the upper integral and the sum over records.  It
requires no exchange of a supremum and an integral.

The inner Wilson factor belongs inside \(K_\omega\).  On a real fibre
it may supply positive action suppression, while its complex
continuation must be controlled by the kept precision and the actual
action reserve.  It is not replaced by Haar measure alone.  The
weighted history argument supplies a route to the envelope in
\eqref{fm:history-envelope-transfer}; the unmarked Gaussian polymer
bound, by itself, supplies only a local bound after the history has
been fixed.

For a finite list of records and a finite-dimensional source parameter,
one may derive derivatives under this operation as follows.  Suppose
each kernel is pointwise holomorphic on a source disc, the operations
and read cores are fixed, and a jointly measurable common envelope
bounds its field supremum on every smaller source circle.  Cauchy's
formula, Tonelli's theorem for that envelope, and dominated convergence
then give the kernel derivative and its norm bound.  For an infinite
record sum one needs the corresponding summable envelope or Banach-norm
holomorphy in \eqref{fm:history-local-norm}.  Pointwise partition
holomorphy alone does not establish that assertion.

Thus the finite local formation proof and the history-mass proof have
distinct conclusions.  The former gives the derivative in a small
potential direction, with the actual bond rooting and own
counterterm.  The latter must pay the norm over stored fields and the
sum of history records.  Their composition is justified by a common
envelope such as \eqref{fm:history-envelope-transfer}, not by dropping
the retained Wilson factor or replacing an integral of a supremum by
a pointwise polymer estimate.



% END MODULE chapters/local_formation_program.tex

% BEGIN LAYERED EXPANSION formation/actual_gaussian_mismatch.tex
\subsection{The actual Gaussian mismatch and its zero coefficient}
\label{layered:actual-gaussian-mismatch}

The action weight and the reference gauge weight need not coincide.
Their ratio contributes a Gaussian term of order \(|\tau-1|\),
which can be much larger than \(g\). This term has an exact total
at a physical constrained-critical background. The argument uses
the actual contour gauge and the complete Hessian on the nonlinear
averaging fibre. It explains why this Gaussian term is retained
separately in an absolute-radius formation estimate.

\subsubsection{The averaging fibre and its residual gauges}

Consider one step on a periodic four-dimensional lattice, with
odd block size \(L\ge3\). Let \(n\) be the number of coarse
sites. There are \(L^4n\) fine sites. The symmetric contour
average \(\mathcal M\) is the map in
\eqref{eq:rg-objects-R3}. Its means and logarithms are taken on
their admitted small branches. Fix a coarse field \(W\).

Write \(\mathcal S(V)\) for the physical action being expanded
on the fine-step field \(V\). At the first step it is the Wilson
action. At a homogeneous later step it is the actual effective
background action \(A(U_k(V))\). In the latter case its
gauge covariance, differentiability, and the asserted critical
background are hypotheses concerning that actual background;
the algebra below does not construct that background.

Let \(V_0\) be a constrained-critical point on
\(\mathcal M(V)=W\). Use the actual nonpivot chart
\begin{equation}
 V(u)=e^{i\Phi(u)}V_0,
 \qquad \mathcal M(V(u))=W,
 \qquad \Phi(0)=0,\quad D\Phi(0)=C,
 \label{agm:actual-chart}
\end{equation}
where \(u\in E\cong\mathbb R^d\) is the list of free
nonpivot colour coordinates. One three-colour pivot is eliminated
for each coarse bond. Hence
\begin{equation}
 d=12(L^4-1)n.
 \label{agm:free-dimension}
\end{equation}
The chart is assumed to cover a neighborhood of the critical
point, with the actual branch fixed. No assertion about a large
history chart is made here.

A residual site gauge is a function \(v\) equal to one at every
coarse node. Gauge covariance of the contour formula gives
\(\mathcal M(V^v)=\mathcal M(V)\), so these gauges preserve
the actual fibre, including its nonlinear constraints.
Their infinitesimal parameters form
\begin{equation}
 \mathfrak h=\{\lambda:\lambda(y)=0\text{ at every coarse node }y\},
 \qquad r:=\dim\mathfrak h=3(L^4-1)n.
 \label{agm:gauge-dimension}
\end{equation}

For a site \(x\in B(y)\), let \(T_{y,x}(V)\) be the actual
symmetric contour transport from the block node \(y\) to \(x\).
It obeys
\begin{equation}
 T_{y,x}(V^v)=v(y)T_{y,x}(V)v(x)^{-1}.
 \label{agm:contour-gauge-law}
\end{equation}
Indeed every path has this endpoint transformation, and uniqueness
of the admitted symmetric mean gives the same law for that mean.
Choose the block-axial representative with
\(T_{y,x}(V_0)=1\). Such a representative is obtained by the
single residual gauge \(v(x)=T_{y,x}(V_0)\), with \(v(y)=1\).
Define the actual local gauge coordinate
\begin{equation}
 \ell_x(V)=\frac1i\log T_{y,x}(V),\qquad x\ne y,
 \qquad \ell(V_0)=0.
 \label{agm:contour-gauge-coordinate}
\end{equation}
These contours lie within their blocks and use no pivot crossing
a block interface.

Let \(T:\mathfrak h\to E\) be the tangent to the residual gauge
orbit in the free coordinates of \eqref{agm:actual-chart}, and
let \(G=D_u\ell(V(u))|_0\). At
\(v_t(x)=e^{it\lambda(x)}\),
\eqref{agm:contour-gauge-law} gives
\(T_{y,x}(V_0^{v_t})=e^{-it\lambda(x)}\). Therefore the exact
linear identity is
\begin{equation}
                         GT=-I_{\mathfrak h}.
 \label{agm:actual-gauge-splitting}
\end{equation}
It is independent of the depth of the homogeneous background and
does not require an estimate for an inverse Hessian.

\subsubsection{The complete physical Hessian}

Put \(S(u)=\mathcal S(V(u))\). Constrained criticality means
\(DS(0)=0\). The physical action Hessian is
\begin{equation}
 K_A=D_u^2S(0).
 \label{agm:physical-action-Hessian}
\end{equation}
It contains the second derivative of the constraint chart. To
make this explicit, write
\(q(a)=\frac1i\log[\mathcal M(e^{ia}V_0)W^{-1}]\),
\(P=Dq(0)\), and take the actual pivot right inverse \(h\),
with \(Ph=1\). Differentiating \(q(\Phi(u))=0\) gives
\begin{equation}
 D^2\Phi(0)[u,v]=-hD^2q(0)[Cu,Cv].
 \label{agm:constraint-second-derivative}
\end{equation}
The ambient action gradient satisfies
\(D\mathcal S(0)=P^{\mathsf T}\mu\), since it annihilates
\(\ker P\). Thus
\begin{equation}
 K_A[u,v]=D^2\mathcal S(0)[Cu,Cv]
                       -\langle\mu,D^2q(0)[Cu,Cv]\rangle.
 \label{agm:constraint-curvature-Hessian}
\end{equation}
The second term is part of the action Hessian, rather than a
small potential to be discarded before the Gaussian split.

\begin{lemma}[Gauge null directions of the constrained Hessian]
\label{lem:agm-physical-gauge-null}
The complete Hessian satisfies \(K_AT=0\). Let the quadratic
reference gauge fixing be
\(K_G=G^{\mathsf T}aG\), where \(a>0\) on
\(\mathfrak h\). Then
\begin{equation}
 E=\ker G\oplus\operatorname{ran}T,
 \qquad K_A\operatorname{ran}T=0,
 \qquad K_G\ker G=0.
 \label{agm:orthogonal-form-sectors}
\end{equation}
The two summands need not be orthogonal in the original free
Euclidean coordinates; they are orthogonal for both quadratic
forms in the displayed sense.
\end{lemma}

\begin{proof}
The chart transports the residual gauge action to a local action
\(u\mapsto h_\lambda(u)\) on the exact fibre. Gauge invariance
gives \(S(h_{t\lambda}(u))=S(u)\). Differentiate first in
\(t\), then in a free direction \(v\), at \(t=0,u=0\).
The term involving the derivative of the gauge tangent is
multiplied by \(DS(0)=0\). The remaining term is
\(K_A[v,T\lambda]=0\). This uses the Hessian of \(S\),
including \eqref{agm:constraint-curvature-Hessian}.

For any \(u\in E\), write
\(u=(u+TGu)-TGu\). The first term lies in \(\ker G\)
by \eqref{agm:actual-gauge-splitting}; the second is a gauge
tangent. If \(T\lambda\in\ker G\), then
\(0=GT\lambda=-\lambda\), so the sum is direct.
The assertions about \(K_G\) follow from its Gram form.
\end{proof}

When \(V_0\) is a real constrained minimum, \(K_A\ge0\).
The reference Gaussian requires
\begin{equation}
                    K:=K_A+K_G>0.
 \label{agm:positive-complete-Hessian}
\end{equation}
This nondegeneracy is retained explicitly. The gauge argument
does not prove positivity on the physical sector \(\ker G\).
The actual reference gauge-fixing function has zero first
derivative at \(V_0\), because all its contour variables equal
one. Its Hessian is exactly \(G^{\mathsf T}aG\); no chart
curvature term is missing from that gauge Hessian.

\subsubsection{The determinant and the compensating expectation}

Let \(\tau=x/\theta\) with \(\operatorname{Re}\tau>0\),
and retain the reference gauge weight \(\theta\). The complete
quadratic exponent is
\begin{equation}
 Q_\tau=\tau K_A+K_G,
 \qquad
 -\tfrac12u^{\mathsf T}Q_\tau u
 =-\tfrac\tau2u^{\mathsf T}Ku
                +\tfrac{\tau-1}{2}u^{\mathsf T}K_Gu.
 \label{agm:complete-tau-quadratic}
\end{equation}
All determinant roots and \(\log\tau\) below are continued
from \(\tau=1\) in the right half-plane.

\begin{proposition}[The actual mismatch normalizer]
\label{prop:agm-exact-mismatch}
Under the preceding physical hypotheses,
\begin{align}
 \det Q_\tau&=\tau^{d-r}\det K,
 \label{agm:determinant-power}\\
 \int e^{-u^{\mathsf T}Q_\tau u/2}\,du
       &=\tau^{-(d-r)/2}\int e^{-u^{\mathsf T}Ku/2}\,du,
 \label{agm:unnormalized-tau-integral}\\
 \int\exp\left(\tfrac{\tau-1}{2}u^{\mathsf T}K_Gu\right)
               d\mu_{(\tau K)^{-1}}(u)
       &=\tau^{r/2}.
 \label{agm:compensating-Gaussian-integral}
\end{align}
Here \(d\mu_{(\tau K)^{-1}}\) is the normalized full
Gaussian on \(E\). It has no fluctuation cutoff. The two
normalizing contributions are separately
\(-\frac d2\log\tau\) and \(\frac r2\log\tau\);
they do not cancel.
\end{proposition}

\begin{proof}
Choose bases in the two sectors of
\eqref{agm:orthogonal-form-sectors}, and combine them into
one invertible change of coordinates \(S\), independent of
\(\tau\). In these coordinates
\[
 S^{\mathsf T}K_AS=\begin{pmatrix}A&0\\0&0\end{pmatrix},
 \qquad
 S^{\mathsf T}K_GS=\begin{pmatrix}0&0\\0&B\end{pmatrix},
 \qquad A>0,\ B>0.
\]
The dimensions are \(d-r\) and \(r\). The factors
\((\det S)^2\) cancel in the determinant ratio, giving
\eqref{agm:determinant-power}. Real Gaussian integration
gives \eqref{agm:unnormalized-tau-integral} for real
\(\tau>0\). For complex \(\tau\) in the right
half-plane, \(\operatorname{Re}Q_\tau\) is positive:
\[
 \operatorname{Re}Q_\tau
       =(\operatorname{Re}\tau)K_A+K_G
       \ge\min(\operatorname{Re}\tau,1)K.
\]
This gives absolute convergence and holomorphic integration
locally uniformly in that half-plane, so the same identity
continues there. Equivalently integrate the two displayed
blocks directly. Normalizing the Gaussian with precision
\(\tau K\) contributes \(\tau^{d/2}\); multiplication
by \eqref{agm:unnormalized-tau-integral} gives
\eqref{agm:compensating-Gaussian-integral}.
\end{proof}

\begin{corollary}[The mismatch has zero physical coefficient]
\label{cor:agm-zero-mismatch-coefficient}
Suppose the localized mismatch family \(\mathbf G_\tau\)
is the actual free-potential leg of the Gaussian expansion, so
its sum on a complete physical pair is the logarithm of
\eqref{agm:compensating-Gaussian-integral}. Then
\begin{equation}
 \sum_X\sum_{z\in X}\mathbf G_\tau(X,\mathbf p_W,z)
              =\tfrac r2\log\tau.
 \label{agm:physical-mismatch-total}
\end{equation}
This total is independent of \(W\). When the local Hessians
converge normally and use the actual physical response and
pointing, their bulk polarization and coefficient satisfy
\par\noindent
\begin{minipage}{\linewidth}
\begin{equation}
                 \Pi^{\mathbf G_\tau}=0,
                 \qquad b[\widehat{\mathbf G}_\tau]=0.
 \label{agm:zero-physical-polarization}
\end{equation}
\popQED\par
\end{minipage}
\end{corollary}

\begin{proof}
The right side of \eqref{agm:compensating-Gaussian-integral}
is positive at real \(\tau>0\). Hence the logarithm selected
from that real value is \(r\log\tau/2\), which gives the
sum rule. Differentiate it twice along the physical background
map. The normally convergent sum of local second derivatives
equals the second derivative of the constant total. At the
trivial pair a first derivative of a scalar gauge-invariant
family vanishes by constant colour rotations. Thus the extra
first-derivative term from differentiating the background map
vanishes, and the second derivative is precisely the bulk
polarization computed with its physical linear response.
It is zero. The second-moment readout defining the coefficient
therefore gives zero as well. Compatible finite-volume
localizations pass this equality to the companion family by
normal convergence; no estimate of an individual local summand
by zero is asserted.
\end{proof}

The sum rule extends to admitted complex physical data by
holomorphic continuation of the actual background map from its
real domain. It does not make the total constant on the whole
independent-current pair core. That core contains directions
transverse to the physical background manifold.

\subsubsection{Two essential qualifications}

First, the nonlinear constraint term in
\eqref{agm:constraint-curvature-Hessian} cannot be removed from
this argument. Even on a circle \(q(p)=|p|^2-1=0\), the
invariant action \(\mathcal S(p)=|p|^2\) is constant on
the rotation orbit. Its constrained action Hessian is zero,
while its ambient Hessian is \(2I\). The constraint curvature
subtracts exactly that tangent contribution. The same chain
rule accounts for the gauge null directions on the averaging
fibre.

Second, replacing the physical current by an independent
current changes that constrained Hessian. For example take
\[
 K_A=\operatorname{diag}(1,0),\qquad
 K_G=\operatorname{diag}(0,1),\qquad
 \widetilde K_A=\operatorname{diag}(1,e).
\]
The physical determinant has the power in
\eqref{agm:determinant-power}, while the changed form has
\(\det(\tau\widetilde K_A+K_G)=\tau(1+\tau e)\).
Its compensating expectation is
\[
       \tau^{1/2}\left(\frac{1+e}{1+\tau e}\right)^{1/2},
\]
which depends on the defect \(e\). This shows why the actual
independent-current completion and its defect must remain in
the local Gaussian family, even though its physical total has
zero coupling coefficient.

Finally, inserting sharp fluctuation tests into
\eqref{agm:compensating-Gaussian-integral} changes the identity.
The test contribution belongs to the cutoff leg of the
non-Gaussian formation remainder. The exact Gaussian mismatch
is defined before those tests are turned on.

% END LAYERED EXPANSION formation/actual_gaussian_mismatch.tex

% BEGIN LAYERED EXPANSION formation/first_step_absolute_formation.tex
\subsection{Absolute-radius formation for the first homogeneous step}
\label{layered:first-step-absolute-formation}

We prove an absolute-radius formation estimate for the first
homogeneous Wilson step. The precision, contour gauge, nonlinear
pivot chart, coarea density and sharp fluctuation cutoff are the
actual ones. A local covariance construction below supplies the
enlarged cube-parameter discs required by the local expansion.
Thus no decoupling-domain assertion is assumed for this step.
The construction does not extend a later nonlocal effective
Hessian by treating it as a finite-range Wilson matrix.

\subsubsection{The range of fields and the exact completion}

Fix the odd block size \(L\). Work on periodic lattices which
admit the \(L\)-blocks and the localization grids chosen below;
their side lengths are at least ten localization cubes. All
constants can depend on \(L\), the requested output rate and
the fixed localization choices, and are uniform in the volume.

Group delta integrals here mean the regular coarea
disintegration of the actual analytic averaging map.
Fix a sufficiently small absolute \(\epsilon_{\rm reg}\)
and restrict this disintegration to fields whose fine
plaquettes in every double block are below that ceiling.
The logarithms are the regular branches of the pivot
construction. No property of an arbitrary Borel
completion outside this regular set is used.

The actual first-step chart of
\cref{ch:first-step-nonflat-precision} solves
\begin{equation}
 q_U(a)=\frac1i\log[\mathcal M(e^{ia}U)\mathcal M(U)^{-1}]=0,
 \qquad a=\Phi_U(u),\qquad \Phi_U(u)|_{\rm nonpivot}=u.
 \label{afs:actual-pivot-chart}
\end{equation}
Its unscaled free radius is \(r_u=c_uL^{-9}\). Its linear
part is \(C_U\), and its nonlinear pivot values use only their
defining double blocks. On a fixed smaller radius,
\(\|\Phi_U(u)\|_\infty\le C_L\|u\|_\infty\),
and all local derivatives through the orders used below have
uniform bounds depending only on \(L\). These are the explicit
joint chart estimates proved in that section, not estimates for
an unspecified many-level constraint.

Here is a gauge-covariant current completion on the entire local
tube used in this proof. Put \(J_A(U)=DA_U(0)\), and define
\begin{align}
 S_{U,J}(u)={}&A(e^{i\Phi_U(u)}U)-A(U)
       +\langle J-J_A(U),\Phi_U(u)\rangle
       -\langle J,C_Uu\rangle,
 \label{afs:centered-action-completion}\\
 G_U(u)={}&\sum_{x\ne y(x)}a_x
   \left[1-\tfrac12\operatorname{Tr}
       \{T_{y(x),x}(e^{i\Phi_U(u)}U)
                         T_{y(x),x}(U)^{-1}\}\right],
 \qquad a_x=1.
 \label{afs:relative-contour-gauge}
\end{align}
These are the actual unit contour-gauge weights in
canonical Gaussian units, held fixed as \(\tau\) varies.
The relative contour product is holomorphic and
conjugates at the block node under a simultaneous gauge change.
Its value at zero is one. It uses no crossing pivot bond.
Consequently both expressions have zero linear part, and
\begin{equation}
 K_A(U,J)=D^2S_{U,J}(0),\qquad K_G(U)=D^2G_U(0),
 \qquad K(U,J)=K_A(U,J)+K_G(U).
 \label{afs:completed-precision}
\end{equation}
On the real block-axial locus this is exactly the first-step
matrix in \eqref{nfp:real-gap}, with its nonlinear constraint
curvature term. The relative gauge has the same Hessian there
as the absolute contour gauge. Away from that locus
\eqref{afs:relative-contour-gauge} specifies the completion;
no equality with a differently defined completion is presumed.

Choose fixed positive \(\epsilon_0,j_0\) strictly inside the
real smallness bounds of \eqref{nfp:real-gap}. At a real
reference pair suppose
\begin{equation}
 \sup_p|1-\partial U(p)|<\epsilon_0,
 \qquad \|J\|_\infty<j_0.
 \label{afs:real-reference-range}
\end{equation}
A node-fixed unitary contour gauge makes \(U\) block axial.
The relative completion is covariant under this change, so the
actual first-step estimates give, in the free coordinates,
\begin{equation}
 \gamma I\le K\le BI,\quad
 \operatorname{range}K\le R_K,
 \quad \|K\|_{\rm Sch}\le B_{\mathrm{Sch}},
 \qquad \gamma>0.
 \label{afs:real-matrix-data}
\end{equation}
Increase \(R_K=C_LL\) so it also bounds the radius of the
background and current variables read by each coefficient.
The Schur norm is the maximum of the absolute row and column
sums of the three-by-three colour blocks. These constants
depend only on \(L\) and the chosen strict real ceilings.

For \(\mathbf U=e^{iE}U\), \(\mathbf J=J+Z\), the same
finite local chart gives the explicit perturbation estimate
\begin{equation}
 \|K(\mathbf U,\mathbf J)-K(U,J)\|_{\rm Sch}
             \le C_L(\|E\|_\infty+\|Z\|_\infty).
 \label{afs:absolute-tube-perturbation}
\end{equation}
For completeness the relative-gauge version has this estimate
because it is a fixed finite product of the admitted means,
inverses and logarithms. In local unitary frames its arguments
lie in the joint chart of radius \(cL^{-5}\); a coefficient
reads at most \(C_L\) variables and has bounded derivatives
there. Summation over their bounded incidence gives the Schur
estimate. Thus the radius is absolute, rather than divided by
the number of lattice sites.

For \(\tau=x/\theta\), let
\(Q_\tau=\tau K_A+K_G\). Taking
\(|\tau-1|\) sufficiently small adds at most
\(C_L|\tau-1|\) to \eqref{afs:absolute-tube-perturbation}.
We shrink all these radii explicitly in the construction below.
They remain independent of \(g=\theta^{-1/2}\).

At an actual real constrained-critical background, with
\(J=J_A(U)\) and \(T_{y,x}(U)=1\), one has
\begin{equation}
 S_{U,J}(u)=A(e^{i\Phi_U(u)}U)-A(U),
 \qquad G_U(u)=\mathrm{GF}(e^{i\Phi_U(u)}U).
 \label{afs:physical-identification}
\end{equation}
The final linear term in \eqref{afs:centered-action-completion}
vanishes because \(C_U^{\mathsf T}J_A(U)=0\). At nonphysical
independent currents it is an auxiliary completion, which was
written explicitly to avoid an unaccounted term of order
\(g^{-1}\) in the linear expansion.

\subsubsection{A local covariance construction from the actual gap}

We give the construction with \(Q_\tau\) denoted by \(K_c\),
and its chosen real reference in
\eqref{afs:real-matrix-data} by \(K_r\). Choose a real
\(a\ge2B\), and write
\begin{equation}
 H_r=I-K_r/a,\quad H_c=I-K_c/a,
 \qquad r=1-\gamma/a<1.
 \label{afs:H-definition}
\end{equation}
The spectrum of \(H_r\) lies in \([0,r]\). Finite
propagation and Cauchy--Schwarz in each matrix row give
\begin{equation}
 \|H_r^j\|_{\rm Sch}\le C_L(1+j)^2r^j.
 \label{afs:power-Schur-bound}
\end{equation}
Indeed a row of \(H_r^j\) has at most
\(C_L(1+j)^4\) nonzero blocks, and its squared row norm is
bounded by a fixed colour factor times \(r^{2j}\). The
column estimate is the same. This estimate retains the
cancellations inside each matrix power.

We will also use the local version of the real gap.
There is a fixed radius \(R_{\rm gap}=C_L\), independent
of volume and of the set \(D\), with the following
property. If the real pair is given and obeys the
strict ceilings on the \(R_{\rm gap}\)-neighborhood of
\(D\), its actual Dirichlet coefficient matrix \(K_{r,D}\)
on free bonds of \(D\) satisfies
\(\gamma I_D\le K_{r,D}\le BI_D\).
To check this, repeat the partition-of-unity proof of
\eqref{nfp:real-gap} for a field supported on \(D\).
Only the fixed-size frame patches meeting its support
are used. Taking \(R_{\rm gap}\) larger than their
diameters and the coefficient margin makes all these
patches available in the given data. Axialize the real
reference by its internal means first; the relative
completion makes this a unitary change of the free
coordinates. The subsequent block-constant frames,
IMS error and flat comparison are then precisely those
in that proof. No extension of a real field through
the holes of \(D\) or across its exterior is required.

The estimate \eqref{afs:power-Schur-bound} concerns one common
real pair. The rooted activity norm below takes each local
coefficient's data supremum before summing its endpoints and
carriers. We retain a positive majorant for that stronger norm.
For an entry of the already combined power, the same local
compression and its spectral gap give
\(\|(H_r^j)(b,b')\|\le r^j\), uniformly over the real
data on its complete carrier. There are at most
\(C_L(1+j)^4\) possible endpoints in a row or column. Put
\begin{equation}
 p_j^+=C_L(1+j)^4r^j,\qquad
 P_*^+=\sum_{j\ge0}p_j^+<\infty .
 \label{afs:positive-power-majorants}
\end{equation}
These are sums of coefficientwise local-data suprema.
The smaller polynomial in \eqref{afs:power-Schur-bound} cannot
be used after moving separate suprema through the endpoint sum.

Fix a desired cube-parameter radius \(S>2\). Choose an
integer \(m\), to be made large, and put
\begin{equation}
                         T=H_c^m.
 \label{afs:short-step}
\end{equation}
One entry of \(T\) depends only on the matrix coefficients
within distance \((m+1)R_K\) of its two endpoints. Put
\(R_c=(m+1)R_K+R_{\rm gap}\). Assign
it the symmetric carrier consisting of both radius-\(R_c\) balls and
all localization cubes they meet. Choose the fine cube size
\(M=M(m)\) to be the smallest power of \(L\) at
least \(16R_c\). Then
\(16R_c\le M<16LR_c\), so it still grows only
linearly in \(m\), is odd, and is compatible with
the centered grids and power-of-\(L\) tori.
For a root set \(W\),
its guard \(W^\sharp\) is the union of all coarse
cubes containing a fine cube at fine-cube wall distance
at most eight from a cube incident to \(W\).
This definition preserves unions and is used in every
restriction and factorization below. Every short prefix
\(H_c^j\), \(j<m\), starting at a root bond, together
with its real-gap data collar, is contained in that
root guard. The number
\(c_s\) of coarse cubes in a short-step carrier is then
bounded independently of \(m\), volume and field.

Set \(\eta_m^+=p_m^+\). We also need a positive bound for
the finite complex perturbation. For each actual finite-range
coefficient, take the supremum of its norm over its common
local real-reference and complex-pair domain. The bounded
number of coefficients meeting a bond, and the finite local
chart estimate \eqref{afs:absolute-tube-perturbation}, give
a fixed positive Schur envelope \(B_+\) for both \(K_r\)
and \(K_c\) on an initial smaller tube, and a difference
envelope of Schur norm
\[
 \delta_K^+\le C_L\bigl(\|E\|_\infty+\|Z\|_\infty
                                      +|\tau-1|\bigr).
\]
These envelopes take the supremum separately for each raw
entry. They use only its full finite carrier. Define
\(C_H=1+B_+/a\). The telescoping polynomial identity,
with these positive entry envelopes multiplied and summed,
gives for every fixed \(j\ge1\)
\begin{equation}
 \|\mathbf{H_c^j-H_r^j}\|_{\mathrm{Sch},\sup}
 \le\Delta_j^+:=(\delta_K^+/a)j C_H^{j-1}.
 \label{afs:polynomial-perturbation}
\end{equation}
Set \(\Delta_0^+=0\).
Here the left side means the positive word majorant after
taking each local coefficient supremum. Thus the estimate
is stronger than a Schur bound at one common pair.
Choose the fixed absolute tube and \(\tau\) disc small
enough that \(\Delta_m^+\le r^m\) and
\begin{equation}
 \sum_{j=1}^{m-1}\Delta_j^+\le r^m .
 \label{afs:prefix-perturbation}
\end{equation}
Both are quantitative finite inequalities, imposed after
\(m,a,B_+\) have been fixed. Their radii are positive and
independent of \(g\), volume and the number of sites.

For the combined power, the real entry bound and the
positive difference envelope give a coefficientwise
Schur majorant at most \(p_j^++\Delta_j^+\).
In particular the short step has bound \(2\eta_m^+\),
and the sum of all complex prefixes has bound
\(P_*^++1\). The sum of their difference envelopes is
at most \(r^m\le\eta_m^+\). These bounds may be used
before every carrier sum even when different coefficients
are evaluated at different admitted local data.

Attach one parameter \(s_\Delta\) to each coarse cube.
Multiply an entire combined entry \(T(b,b')\) by
\(\prod_{\Delta\in\operatorname{car}(b,b')}s_\Delta\),
and call the resulting matrix \(T_s\). We combine the
polynomial entry before inserting this parameter, rather
than separating its mutually cancelling fine paths. For
\(|s_\Delta|<S\),
\begin{equation}
 \|T_s\|_{\rm Sch}\le2S^{c_s}\eta_m^+=:q_m^+.
 \label{afs:weighted-step-smallness}
\end{equation}
All its coefficients are holomorphic local functions of the
actual pair. Its carrier contains every background variable
used by the entry.

Let \(c_j=4^{-j}\binom{2j}{j}\), the coefficients of
\((1-z)^{-1/2}\). For a finite root set of free bonds
\(W\), hold the parameters at one on its fixed guard
\(W^\sharp\), and set
\begin{equation}
 R_W(s)=a^{-1/2}\chi_W
       \sum_{j=0}^{m-1}H_c^j
                      \sum_{k\ge0}c_{km+j}T_s^k.
 \label{afs:decoupled-inverse-root}
\end{equation}
The prefix is local inside the root guard. The tail is a
chain of the short-step carriers. Since \(0<c_j\le1\),
the positive prefix bound \(P_*^++1\) and
\eqref{afs:weighted-step-smallness} show that this series
converges in Schur norm as soon as \(q_m^+<1\).

At all parameters equal to one, grouping the absolutely
convergent binomial series by \(j\bmod m\) gives
\begin{equation}
 R_W(1)=\chi_W K_c^{-1/2},
 \qquad D_W(s):=R_W(s)R_W(s)^{\mathsf T},
 \qquad D_W(1)=\chi_WK_c^{-1}\chi_W^{\mathsf T}.
 \label{afs:complete-marginal}
\end{equation}
Transposition is bilinear; no conjugation is inserted in this
analytic covariance. The binomial identity is legitimate in
the convergent coefficient algebra. In particular its square
is \(K_c^{-1}\), and its transpose is itself before the
cube parameters are inserted.

Define the real short-prefix reference
\[
 R_{0,r}=a^{-1/2}\sum_{j<m}c_jH_r^j,
 \qquad D_{0,W}=\chi_WR_{0,r}^2\chi_W^{\mathsf T}.
\]
Since \(H_r\ge0\),
\begin{equation}
             a^{-1}I_W\le D_{0,W}\le\gamma^{-1}I_W.
 \label{afs:local-reference-gap}
\end{equation}
Its definition uses only the fixed guard of \(W\).
Equations \eqref{afs:prefix-perturbation} and
\eqref{afs:decoupled-inverse-root} give, uniformly on all
the enlarged parameter discs,
\begin{equation}
 \|R_W(s)-\chi_WR_{0,r}\|_{2\to2}
 \le a^{-1/2}\left[\eta_m^++
               (P_*^++1)\frac{q_m^+}{1-q_m^+}\right]=:\zeta_m^+.
 \label{afs:inverse-root-close}
\end{equation}
Consequently \(\|D_W(s)-D_{0,W}\|\le C_{a,\gamma}\zeta_m^++
(\zeta_m^+)^2\), independently of the local and auxiliary
dimensions. As \(m\to\infty\), this tends to zero
exponentially, while \(M(m)\) grows only linearly. We may
therefore choose \(m\) so large that
\begin{equation}
 q_m^+<1/4,\qquad
 C_{a,\gamma}\zeta_m^++(\zeta_m^+)^2
              \le c_*(LM(m))^{-4}.
 \label{afs:Gaussian-volume-budget}
\end{equation}
Here \(c_*\) is any prescribed small Gaussian comparison
constant. The order is: fix the desired rates and \(S\);
choose \(m\) and \(M(m)\); then choose the absolute tube
and \(\tau\) disc by
\eqref{afs:polynomial-perturbation}--\eqref{afs:prefix-perturbation}.

\begin{proposition}[Admitted local Gaussians for the actual first step]
\label{prop:afs-local-Gaussian}
On the resulting fixed absolute tube, the functions
\begin{equation}
 \mathcal J_W(s)[F]=\int_{\mathbb R^{3\#W}}
                            F(b)\,d\mu_{D_W(s)}(b)
 \label{afs:local-Gaussian-functional}
\end{equation}
are normalized holomorphic local Gaussian functionals for
\(|s_\Delta|<S\). They obey the local Gaussian comparison
of \cref{thm:localized-complex-Gaussian}, with perturbation
at most \(c_*(LM)^{-4}\) and with \(\Gamma=0\).
At complete physical data and \(s=1\), they are the actual
marginals of the first-step Gaussian with precision
\(Q_\tau\). When parameters vanish outside a union \(Z\)
containing \(W^\sharp\), they read no pair data outside
\(Z\), and they factor over the wall components of \(Z\).
\end{proposition}

\begin{proof}
The covariance is symmetric and within the bound
\eqref{afs:Gaussian-volume-budget} of the real positive
matrix \(D_{0,W}\). In particular its real part is
uniformly positive after decreasing \(c_*\). The Gaussian
integral and its determinant branch are therefore defined
on the whole indicated polydisc, are normalized, and are
holomorphic with measurable sharp functions allowed in
\(F\). Apply \cref{thm:localized-complex-Gaussian} with
\(C=D_W\), \(C_0=D_{0,W}\), \(\Gamma=\Gamma_0=0\).
Only \(3\#W\) coordinates occur in this comparison.
The possibly much larger column space of \(R_W\) is
not charged as a Gaussian dimension. Completeness is
\eqref{afs:complete-marginal}.

If parameters outside \(Z\) are zero, each surviving
short step has its full carrier in \(Z\). Every prefix
from a root in \(W\) is already in its guard. Thus every
coefficient and intermediate column used by that row lies
in the same wall component as the root. A short-step
carrier is a connected cube union containing both endpoints;
it cannot join two different components without reading a
zero parameter. Prefix carriers are connected and lie in
the guard. Rows rooted in different components therefore
have disjoint column supports. Their bilinear covariance
blocks are zero, proving Gaussian factorization. The same
description proves dependence only on pair data in \(Z\)
and the required hereditary restriction identity for
the fixed root support and its held guard: computing
from the restricted pair in \(Z\) agrees with any
larger ambient computation whose parameters are zero
outside \(Z\). At complete parameters the usual
nested marginal identities follow from
\eqref{afs:complete-marginal}.
\end{proof}

This is a localization of the actual covariance, fixed once
for all tori in the admitted class. It need not equal another
walk expansion term by term. Its full Gaussian marginals,
nonlinear integration variables, and complete physical
integral are exactly the original ones.

The estimates and local recipes hold also for pairs
given only on the output domain. To see the domain
point explicitly, retain only entries whose full
carrier is contained in that domain. If a retained
entry has first endpoint \(b\), the real data on its
radius-\(R_c\) ball are present. Compress the actual
matrix to the radius-\(mR_K\) ball about \(b\).
The preceding local gap applies to this compression.
The row at \(b\) of its \(m\)-th \(H\) power is the
same polynomial row as the uncompressed one: none
of its length-\(m\) paths can leave that ball.
This proves the pointwise estimate
\eqref{afs:power-Schur-bound} for that entire common-data
row before any inadmissible entries are deleted. For
coefficientwise local-data suprema, apply the same compression
to each entry separately and count its possible endpoints;
this gives \eqref{afs:positive-power-majorants}.
The positive finite-range perturbation envelopes give
\eqref{afs:polynomial-perturbation} on these local domains
as well. Deleting entries can only decrease either positive
majorant, and symmetry gives the corresponding column bounds. For the root
prefix and \(D_{0,W}\), use one common Dirichlet
compression to the \(mR_K\)-neighborhood of \(W\);
its needed real-data collar is inside \(W^\sharp\).
This preserves \eqref{afs:local-reference-gap}.
Thus the local estimates are on the actual local
pair domains with these collars, rather than only
on restrictions of globally extendible real fields.

The same construction localizes the Gaussian logarithms.
For the actual matrix \(K_c\), the identity
\begin{equation}
 \log K_c=(\log a)I-\sum_{\ell\ge1}H_c^\ell/\ell
 \label{afs:actual-matrix-log}
\end{equation}
converges in Schur norm after grouping
\(\ell=km+j\), \(0\le j<m\), and omitting \(\ell=0\).
Its diagonal trace at a bond is expanded as the finite
prefix \(H_c^j\) followed by \(k\) actual short steps.
The carrier is the union of their full carriers. It is
connected and has tree length at most \(C_0+c_sk\).
Multiplication by \(e^{\kappa d_c(X)}\) therefore costs
at most \(e^{\kappa C_0}e^{\kappa c_sk}\).
Choose \(S\) and \(m\) above with this additional rate
reserve. Each word is bounded by the product of the positive
prefix and short-step envelopes before taking the supremum
on its own union carrier. Summing those positive words proves
that the coefficientwise absolute rooted sum is bounded by
\[
 C\sum_{j<m}(p_j^++\Delta_j^+)
     \sum_{k\ge0}(2\eta_m^+ e^{\kappa c_s})^k
 \le C(P_*^++1),
\]
with the finitely many \(k=0,j>0\) terms included.
This proves locality and an absolute, volume-independent
rooted norm for the actual Gaussian log normalizer.
The value \(\log J_U(0)\) adds its finite local pivot
terms and fixed numerical powers; its nonconstant part
has the same type of bound by the joint pivot chart.
\par{}Use the same localization on a \(\tau\) disc four times
larger than the disc retained in the conclusion. Applying
it to
\(-\tfrac12\log\det Q_\tau+\tfrac12\log\det K
+\tfrac d2\log\tau\) gives the mismatch family.
Its terms vanish at \(\tau=1\), and the Cauchy bound
in this actual larger disc proves their rooted bound
\(C|\tau-1|\). Its physical total remains exactly
\eqref{agm:physical-mismatch-total}. The construction
does not classify that term as a \(g\)-small vertex.

The Gaussian normalizer is pointed at each diagonal
free bond in \eqref{afs:actual-matrix-log}; each
zero-field pivot factor is pointed at that pivot.
Colour traces are invariant under the local adjoint
frames. Their midpoint tents give the pointed
normalizer, with the same rate reserve.

For the mismatch use its actual gauge-quadratic
interpolation, which also fixes its bond density:
\begin{equation}
 \begin{split}
 Q_{\tau,u}&=\tau K-u(\tau-1)K_G,\qquad 0\le u\le1,\\
 g_{\tau,b}(P)&=\frac{\tau-1}{2}\int_0^1
             \operatorname{tr}[K_GQ_{\tau,u}^{-1}](b,b)\,du.
 \end{split}
 \label{afs:actual-mismatch-pointing}
\end{equation}
Indeed \(Q_{\tau,0}=\tau K\) and
\(Q_{\tau,1}=Q_\tau\). Differentiating its Gaussian
determinant gives
\(\sum_bg_{\tau,b}=\mathcal G_\tau\), with the
same normalization as above. The integrand is the
expectation of the actual insertion
\((\tau-1)B(b)\cdot(K_GB)(b)/2\).
All \(Q_{\tau,u}\) stay in the same admitted local
tube after a fixed decrease of the retained
\(\tau\) disc. Their inverse series uses
\(a^{-1}\sum_{\ell\ge0}H_{\tau,u}^\ell\), grouped
by the same short powers. Marking its one finite-range
\(K_G\) factor preserves the carrier bound. Hence
\eqref{afs:actual-mismatch-pointing} has local
pointed terms of norm \(C|\tau-1|\). This identifies
the pointing by the defining interpolation, including
in the noncommuting curved case.

\subsubsection{The actual coarea density and the intrinsic vertices}

Write \(g=\theta^{-1/2}\),
\(\mathfrak b=A_1(\log\theta)^p\ge1\), and keep the
cutoff \(\chi=\prod_b\mathbf1_{\{|B(b)|<\mathfrak b\}}\)
fixed at this real reference. Choose \(g\) so small that
\begin{equation}
                   16C_Lg\mathfrak b<r_u.
 \label{afs:nonlinear-chart-admission}
\end{equation}
Every argument \(\Phi_U(tgB)\), \(|t|\le4\), and every
local old-member substitution below then lies in the actual
unscaled graph chart. Cube parameters do not change these
local graphs: their complete double-block support is kept
in each selected potential's guard. They change the local
Gaussian functional outside that guard by the construction
just proved.

Let \(p_c\) be the three pivot coordinates of row \(c\).
Eliminating the exact constraint in these coordinates gives
the relative density
\begin{equation}
 J_U(u)=
 \prod_{b}j(\Phi_U(u)(b))
 \prod_c\det\!\left[D_{p_c}q_U(\Phi_U(u))\right]^{-1},
 \qquad j(a)=\left(\frac{\sin\sqrt{a\cdot a}}
                            {\sqrt{a\cdot a}}\right)^2.
 \label{afs:actual-coarea-density}
\end{equation}
The factor \(g^d\) and the fixed Haar and gauge
normalizers are independent of \(B\). We use the raw
logarithmic constraint \(q_U\) throughout. A per-length
constraint would require its compensating factor
\(L^{-3}\) for every coarse bond; it must not be
inserted a second time into the present raw convention.
At real arguments the pivot
determinant has its positive identity sign. Its reciprocal
in \eqref{afs:actual-coarea-density} is the holomorphic
continuation, with no complex absolute determinant.
Each pivot factor reads only its double block.

On the uniform unscaled radius, the logarithm of each
factor of \(J_U(u)/J_U(0)\) is holomorphic, vanishes at
zero and has derivative bounded by \(C_L\). For the
pivot factor this follows by differentiating
\(-\operatorname{Tr}\log[D_pq(u)(D_pq(0))^{-1}]\),
using the actual pivot lower bound and joint chart radius.
For Haar it follows from \eqref{sf:haar-ratio}. Hence its
local logarithmic vertices satisfy a linear bound
\(C_Lg\sum_{b\in S}|B(b)|\) on their fixed supports.

Subtract the full quadratic parts, and define
\begin{align}
 P^A_g(B)&=-g^{-2}\{S_{U,J}(gB)
                         -\tfrac12g^2B^{\mathsf T}K_AB\},\notag\\
 P^G_g(B)&=-g^{-2}\{G_U(gB)
                         -\tfrac12g^2B^{\mathsf T}K_GB\},\notag\\
 P^J_g(B)&=\log[J_U(gB)/J_U(0)].
 \label{afs:actual-intrinsic-potentials}
\end{align}
Every Wilson plaquette, relative contour gauge term and
current-completion term is a holomorphic function of a
fixed finite list of free coordinates, on an unscaled
radius. Its first three derivatives are uniformly bounded
there by the actual chart estimates. Taylor's integral
formula consequently gives, for each such local term,
\begin{equation}
 |g^{-2}R_3(gB)|
    \le C_Lg\left(\sum_{b\in S}|B(b)|\right)^3
    \le C_Lg\mathfrak b\sum_{b\in S}|B(b)|^2.
 \label{afs:local-g-Taylor}
\end{equation}
The second inequality uses the fixed bound on \(\#S\)
and the actual cutoff \(|B(b)|\le\mathfrak b\). The
coefficients are uniform on the absolute external pair
tube, not just at \(U=1\). The current-completion
correction is included before this Taylor expansion.
The same estimates, with a fixed increase in \(C_L\),
hold on nested complex fluctuation boxes
\(|B_b^a|<\nu(1+\mathfrak b)\), for a fixed
\(1<\nu<9/8\), and on their fixed radial margins.
The factor sixteen in
\eqref{afs:nonlinear-chart-admission} leaves this
room for the independent-bond pointing construction.

Group these terms into their actual fine localization
cube carriers, enlarged by the fixed double-block collar.
For \(Q_Y(B)=1+M^{-4}\sum_{b\subset Y}|B(b)|^2\), the
finite support and incidence counts give
\begin{equation}
 |V^{\rm int}_g(Y;P,B)|
 \le C_{L,M}g\mathfrak b\,e^{-\lambda_Vd_f(Y)}Q_Y(B),
 \qquad V^{\rm int}_g=\tau P^A_g+P^G_g+P^J_g.
 \label{afs:intrinsic-potential-norm}
\end{equation}
Indeed each elementary support meets a bounded number
of fine cubes, so its exponential carrier cost is fixed;
a fine cube contains \(O(M^4)\) elementary terms, and a
free coordinate enters only \(C_L\) of them. All these
finite factors are included in \(C_{L,M}\). The same
bound holds for \(|\tau-1|\) on its fixed admitted disc.

\subsubsection{Admitted old functions and ordinary source profiles}

Let a local input \(S_{\rm old}=\sum_Xf_X\) be holomorphic
on fixed pair balls of radii \(a_{H,U},a_{H,J}>0\)
around the indicated field and current, with bounds
\(B_X\) and
\begin{equation}
 \sup_b\sum_{X\ni b}e^{\lambda_Hd_f(X)}B_X\le H_0.
 \label{afs:old-absolute-input}
\end{equation}
Use the same physical input function before and after
the chart substitution. To include its current
argument, define the following explicit pair map:
\begin{equation}
 \begin{split}
 U_P(u)&=e^{i\Phi_U(u)}U,\\
 \Psi_P(u)&=\big(U_P(u),
       J+J_A(U_P(u))-J_A(U)\big),\qquad P=(U,J).
 \end{split}
 \label{afs:actual-pair-substitution}
\end{equation}
The fine Wilson gradient reads only its adjacent
plaquettes. Its local derivative is uniformly
bounded on these fixed unscaled tubes. Hence
\(\|J_A(U_P(u))-J_A(U)\|_\infty
\le C_L\|u\|_\infty\). This is the actual gradient
variation, rather than a current held fixed during
a physical field variation. The map is jointly
gauge covariant, equals \(P\) at zero, and at a
physical pair maps exactly to
\((U_P(u),J_A(U_P(u)))\).
Put \(a_H=\min(a_{H,U},a_{H,J})\), and require
\(16C_Lg\mathfrak b<a_H\) as well as
\eqref{afs:nonlinear-chart-admission}. The exact old
potential is
\begin{equation}
 V_X^{\rm old}(B)
   =f_X(\Psi_P(gB))-f_X(P).
 \label{afs:actual-old-difference}
\end{equation}
Its second value stays in the carried background action.
The local substitution reads the fixed double-block
enlargement of \(X\), including the adjacent plaquettes
needed for its current arguments. The analytic path
\(t\mapsto f_X(\Psi_P(tgB))\) is admitted up to
\(|t|<c\min(r_u,a_H)/(g\mathfrak b)\). The one-variable
Cauchy difference on that actual path gives
\[
 |V_X^{\rm old}(B)|
       \le C_L\frac{g\mathfrak b}{\min(r_u,a_H)}B_X.
\]
The enlargement of a domain has at most a fixed multiple
of its tree length plus one. Choose \(\lambda_H\)
larger than that multiple of \(\lambda_V\), with a
strict incidence/entropy reserve. Refiling by that
enlargement and summing \eqref{afs:old-absolute-input}
then gives the actual potential budget
\begin{equation}
 \|V^{\rm old}\|_{\lambda_V}
       \le C_{L,M,a_H}H_0g\mathfrak b.
 \label{afs:old-potential-norm}
\end{equation}
All input origins with the same carrier are summed before
this estimate; their number is not counted as a number
of independent Gaussian variables or potential types.

This includes the ordinary profile presentation of a
Wilson source \(\sum_p w(p)a_p\), with \(w\) on fixed
complex coefficient discs. Each term is a finite local
holomorphic word, and \eqref{afs:old-absolute-input}
holds with \(H_0\le C\|w\|_\infty\). The Gaussian,
chart, reference gauge weight and sharp cutoff stay fixed
as this source is varied. A constant-shift presentation
which instead changes \(\tau\) uses its separate
Gaussian mismatch and numerical factor.

\subsubsection{The three output families and the exact first-step law}

Let \(\mathcal N^\circ\) be the localization just
constructed of the vacuum-subtracted coarea and
Gaussian normalizer at \(\tau=1\):
\begin{equation}
 \mathcal N^\circ(P)=
 \log\frac{J_U(0)}{J_1(0)}
 -\frac12\log\frac{\det K(U,J)}{\det K(1,0)}.
 \label{afs:actual-N-family}
\end{equation}
Its vacuum factors stay in the exact numerical
normalization. Let \(\mathcal G_\tau\) be the free Gaussian
mismatch defined before cutoffs, as in
\cref{prop:agm-exact-mismatch}. Define the last family
from the actual normalized integral
\begin{equation}
 \mathcal F_g(P)=
 \log\int\chi(B)
       \exp\{V_g^{\rm int}(P,B)+V^{\rm old}(P,B)\}
                         \,d\mu_{Q_\tau^{-1}}(B).
 \label{afs:actual-F-integral}
\end{equation}
This is a definition of the output, not an assumed bound.

\begin{proposition}[Actual first-step absolute-radius formation]
\label{prop:afs-formation}
Choose the rates and localization sizes as above. There
is a coupling floor \(\theta_*\), independent of volume,
such that \eqref{afs:actual-F-integral} is nonzero and has
an actual union-filed logarithm \(F_g(X,P)\) on the full
fixed absolute pair tube. Its terms satisfy
\begin{equation}
 |F_g(X,P)|\le C\left[(1+H_0)g\mathfrak b
                          +e^{-c\mathfrak b^2}\right]
                                  e^{-\kappa d_c(X)}.
 \label{afs:first-step-F-bound}
\end{equation}
The same estimate, after a fixed decrease of the spatial
rate, holds for the constructed bond-rooted pointing and
its midpoint tent sum. The constants include the fixed
absolute radii, but contain no inverse power of \(g\).
The output splits into the actual Gaussian family,
the Gaussian mismatch, and this remainder. In particular
\(F_g=O(g(\log\theta)^p)\).
\end{proposition}

\begin{proof}
The preceding construction supplies every Gaussian
matrix, enlarged coarse-parameter disc, normalization,
hereditary restriction and wall factorization used in
the unmarked local expansion of
\cref{ch:local-formation-program}. The local coordinate
set for an index is taken to contain all free bonds
of its full fine support \(Z_0\), not merely the
bonds on which its integrand changes. Thus background
supports of the selected potentials, including a
background factor constant along the fibre, lie in
the held guard as well. Adding these integrated
coordinates does not change a Gaussian expectation
of that integrand. The local coordinate
count is at most \(C(LM)^4\#Z\) for a polymer \(Z\).
There is no auxiliary-volume factor. The potentials are
the literal local functions in
\eqref{afs:actual-intrinsic-potentials} and
\eqref{afs:actual-old-difference}, with their actual
cutoff supports. Equations
\eqref{afs:intrinsic-potential-norm} and
\eqref{afs:old-potential-norm} give
\(a_g\le C(1+H_0)g\mathfrak b\).

Apply the proved finite Mayer, cutoff and coarse
interpolation estimates to these data. Their activity
bound has amplitude \(C(a_g+p_*)\), where
\(p_*\le C(LM)^4e^{-c\mathfrak b^2}\). The rooted
connected-cluster bound and union filing then give
\eqref{afs:first-step-F-bound}. All its numerical
smallness conditions hold after increasing \(\theta_*\),
because \(g\mathfrak b\to0\) and
\(e^{-c\mathfrak b^2}\to0\). This uses that
\(A_1(\log\theta)^p\) tends to infinity; choose
\(p>1/2\), so its cutoff tail is also
\(O(g\mathfrak b)\).

The Jacobian and old functions are rooted along
\(B\mapsto tB\). For the cubic action and gauge
remainders use precisely the normalized path
\(t^{-2}P^{A,G}_g(tB)\), with the covariant shares
and independent bond derivatives of
\eqref{fm:explicit-intrinsic-pointing}. In particular
the bond insertion includes the derivative of
\(t^{-2}\); it is not the unnormalized Euler
insertion. The admitted chart has a fixed radial
margin by \eqref{afs:nonlinear-chart-admission}.
The order-three vanishing and the local unscaled
derivative estimates used in
\eqref{afs:local-g-Taylor}, followed by
\eqref{fm:constructed-pointing-integrability},
give the rooted budgets with the same
\(g\mathfrak b\) factor, uniformly down to \(t=0\).
The cutoff leg
is the actual moving-threshold boundary term of
\eqref{fm:rooted-cutoff-logarithm}. It supplies its
Gaussian tail and retains its bond root. Thus the
previously proved rooted expansion applies to these
literal legs; no pointing is inferred merely from an
unpointed norm. Summing their midpoint tents costs a
polynomial domain factor, absorbed by the reserved rate.
\end{proof}

Here is the exact physical normalization. Let \(U\) be
the actual first-step constrained-critical background,
in its block-axial representative. In this formula
\(S_{\rm old}(U)\) denotes the physical restriction
\(S_{\rm old}(U,J_A(U))\). Apart from the
external coarse indicator, the bare small-field integral
equals
\begin{equation}
 e^{-xA(U)+S_{\rm old}(U)}\,
 C_{\rm num}(\theta)g^dJ_U(0)(2\pi)^{d/2}
 (\det K)^{-1/2}\tau^{-d/2}
              \exp\{\mathcal G_\tau+\mathcal F_g\}.
 \label{afs:exact-physical-step}
\end{equation}
More explicitly, let \(n_f,n_c\) be the numbers of
positive fine and coarse bonds, and let
\(n_s=N_f-N_c\) be the number of node-fixed gauge
sites. In the raw convention of
\eqref{afs:actual-pivot-chart},
\begin{equation}
 C_{\rm num}(\theta)=\sigma_0^{n_f-n_c}
                      z(\theta,\varepsilon_\theta)^{-n_s},
 \qquad d=3(n_f-n_c),
 \label{afs:exact-numerical-prefactor}
\end{equation}
where \(\sigma_0\) is Haar density at zero and
\(z\) is the normalized group integral of the
single-site contour gauge with its retained test.
Indeed the fine Haar densities contribute
\(\sigma_0^{n_f}\), and each coarse group delta,
written in raw logarithms, contributes
\(\sigma_0^{-1}\). At the trivial field the raw
pivot derivative is \(L^{-3}I_3\), so
\(J_1(0)=L^{9n_c}\). Equivalently one may use the
per-length pivot derivative \(L^{-4}I_3\), its
factor \(L^{12n_c}\), and the delta conversion
\(L^{-3n_c}\); their product is the same.
Equation \eqref{afs:actual-coarea-density} proves the
coordinate identity; the centered action and gauge
agree with the physical ones by
\eqref{afs:physical-identification}. The quadratic
normalization is
\cref{prop:agm-exact-mismatch}. Every factor depending
on \(U\) is included, in particular \(J_U(0)\).

The first-step formation budget can therefore be
closed in the prescribed order. The construction
gives constants \(C_N,C_G,C_F\), independent of
the allowed old-input size \(N_*\), such that
\[
 \|\mathbf N^\circ\|_+\le C_N,\qquad
 \|\mathbf G_\tau\|_+\le C_G|\tau-1|,\qquad
 \|\mathbf F_g\|_+\le
       C_F\{(1+N_*)g\mathfrak b+e^{-c\mathfrak b^2}\}.
\]
Choose \(E_0\ge2C_N\), then the old-input ceiling
\(N_*\); decrease \(|\tau-1|\) so that its term is
at most \(E_0/4\); finally increase the coupling
floor so that the last term is at most \(E_0/4\).
This gives \(\|\mathbf N^\circ+\mathbf G_\tau+
\mathbf F_g\|_+\le E_0\) for the proved first-step
range, without putting \(E_0\) or \(N_*\) into
the preceding Gaussian or nonlinear chart choices.

\begin{corollary}[Ordinary source jets in the proved first-step range]
\label{cor:afs-ordinary-source-jets}
Insert the literal profile \(\sum_pw(p)a_p\) on the
fixed source ball \(\|w\|_\infty<\rho\). For
\(\|w\|_\infty\le\rho/2\), \(n\ge1\), and arbitrary profiles
\(h_1,\ldots,h_n\),
\begin{equation}
 |D^n_wF_g(X;w)[h_1,\ldots,h_n]|
 \le C_\rho g\mathfrak b\,n!
       (2e/\rho)^n e^{-\kappa d_c(X)}
                         \prod_{i=1}^n\|h_i\|_\infty.
 \label{afs:ordinary-source-jet-bound}
\end{equation}
The same estimate holds for the constructed pointing
with its reserved spatial rate. Only the source
plaquette vertices whose complete carriers are filed
inside \(X\) occur in that derivative.
\end{corollary}

\begin{proof}
The source potential is the actual admitted old
difference already considered; its norm on the full
source ball is at most \(C\rho g\mathfrak b\).
All Gaussian, chart and cutoff arguments stay fixed.
The Mayer integrals and convergent cluster sum are
therefore holomorphic in these source coefficients,
with \eqref{afs:first-step-F-bound} uniformly on
that ball. For nonzero \(h_i\), use the independent
complex directions \(w+\sum_i z_ih_i\) on
\(|z_i|<\rho/(2n\|h_i\|_\infty)\).
They lie in the admitted source ball. The
multivariate Cauchy formula bounds the mixed
derivative by
\(C_\rho g\mathfrak b(2n/\rho)^n
\prod_i\|h_i\|_\infty e^{-\kappa d_c(X)}\).
Use \(n^n\le e^nn!\). A zero direction gives zero.
The same argument applies to the analytic rooted
terms. Their source support is literal: an activity
filed inside \(X\) uses only selected potential
carriers inside \(X\), and its cluster union stays
there. No source variable outside that set is read.
\end{proof}

\subsubsection{Physical tests and the largest proved range}

The preceding fixed absolute domain is the analytic
domain of the filed functions. It is not the support
of every original external small-field indicator.
We first check that this regular coarea fibre contains
exactly the graph used above on the free fluctuation
box. Every middle pivot belongs to an elementary
plaquette whose other three bonds are free; choose a
fixed perpendicular direction. Odd \(L\ge3\) keeps
every other middle pivot out of that plaquette. Let
\(V\) be in the regular coarea set, with average
\(W=\mathcal M(U)\), and with its free bond logarithms
relative to \(U\) at most \(\delta=g\mathfrak b\).
Solving that plaquette product for its pivot bond,
and subtracting the corresponding expression for
\(U\), gives
\begin{equation}
 \max_{\rm pivots}\|V(b)U(b)^{-1}-I\|
       \le 2\epsilon_{\rm reg}+C\delta.
 \label{afs:regular-pivot-anchor}
\end{equation}
Here the other three factors are unitary and their
differences are at most \(C\delta\); the two
plaquette products differ by at most
\(2\epsilon_{\rm reg}\). These estimates are
unchanged by their left and right transports.
The pivot fixed-point map has derivative at most
\(C_L\|a\|_\infty\) in a fixed full-field ball
\(r_p=c_pL^{-5}\), as follows from its quadratic
constraint remainder and pivot inverse. Choose
\(\epsilon_{\rm reg}\) so small that
\eqref{afs:regular-pivot-anchor} puts all full
logarithms in that uniqueness ball, and then take
\(\delta\) below its fixed small floor.
The equation \(q_U(a)=0\), with its fixed free
values, has exactly the chart solution there.
Thus \(V=e^{i\Phi_U(gB)}U\) and its pivot values
obey the sharper \(C_Lg\mathfrak b\) bound.
Conversely the graph stays in the regular coarea
set when the background curvature is below
\(\epsilon_{\rm reg}/2\) and \(g\mathfrak b\)
is small enough. This proves the full change of
variables on the stated regular fibre, rather
than just an injective parametrization of a subset.

For the homogeneous tame operator itself, take real
coarse data with \(\|1-\partial W\|_\infty<\epsilon_W\),
where the fixed absolute \(\epsilon_W>0\) is small
enough for \cref{ch:first-step-background-expanded}
and \eqref{afs:real-reference-range}. No shrinking
curvature condition is needed for that operator.
Its additional tests are the contour-gauge tests
at threshold \(\varepsilon_\theta=
A_0g(\log\theta)^p\). The background contour means
are exactly one, and those means read only internal
free bonds. The product/mean derivative estimate
therefore gives their deviations at most
\(C_Lg\mathfrak b\), independently of the background
curvature. Choose \(A_0>2C_LA_1\). These indicators
are then identically one. The nonpivot fluctuation
field in \eqref{afs:actual-pivot-chart} is exactly
\(gB\), so its prescribed test is exactly \(\chi\).
This proves \eqref{afs:exact-physical-step} on the
fixed absolute physical background range.

If this factor is used in a history with extra
plaquette indicators, a sufficient common-test
range is smaller. Namely require
\(\|1-\partial W\|_\infty<c_L\varepsilon_\theta\).
The actual background fine curvature is then at most
\(C_Lc_L\varepsilon_\theta\), and the graph changes
fine plaquettes by at most \(C_Lg\mathfrak b\).
Taking \(c_L\) and \(A_1/A_0\) sufficiently small
puts their sum below half the indicated thresholds.
Those additional tests are identically one on the
same box. Their identity has not been asserted on
the larger absolute physical range.

Any such coarse external indicator remains outside
this formula. It cannot be analytically continued as an
indicator on a fixed absolute ball: a fixed nonzero
real curvature eventually exceeds
\(A_0g(\log\theta)^p\). The absolute-radius output
is instead the analytic extension of the displayed
fluctuation formula, with that external indicator
factored out. The physical background continuation
in the first-step theorem supplies its composition
on a fixed sufficiently small complex coarse-datum
ball. Its current is \(J_A(U(W))\), so it lies in
the physical completion just used. Source profiles
vary on their independent fixed discs while the
tests keep their real reference values.

The proof covers this actual first bare homogeneous
step, arbitrary compatible volume and real flat
holonomy, and its specified absolute complex tube.
For a later step the effective action
\(A(U_k(V))\) is nonlocal. Its Hessian does not have
the finite propagation used in
\eqref{afs:power-Schur-bound} and the short-step
carrier construction. Extending this result to all
homogeneous levels requires a proved localized
effective-precision expansion and the corresponding
unscaled nonlinear response domains. Those are the
earliest additional operator and domain inputs;
they are not replaced by assuming the desired
activity norm.

% END LAYERED EXPANSION formation/first_step_absolute_formation.tex

% BEGIN LAYERED EXPANSION formation/homogeneous_effective_action_trees.tex
\subsection{Connected trees for the actual homogeneous effective action}
\label{layered:homogeneous-effective-action-trees}

The homogeneous background of
\cref{prop:hb-actual-critical-orbit} has an expansion
whose spatial rate is independent of its Taylor order.
We derive it from the actual Wilson and averaging
equations. This supplies local analytic terms for
\(S_k(V)=A(U_k(V))\); localization by a single large
ball would not give the connected carriers below.

The coarse spacing is one and \(h=L^{-k}\). Use
the parallel frame of the fixed flat reference
\(W_0\), and coarse coordinates \(V_B=e^{iB}V_0\).
Fine pairings use \(h^4\) times counting measure.
The constants in this section are independent of
\(h\), compatible volume and flat holonomy. The
localization size, chosen after a requested rate,
is allowed to enter them. No independent fine-current
core or hereditary nonlinear tower is asserted.

Gauge invariance holds for the complete physical
sum \(S_k\). Its individual fixed-reference
Taylor terms need not be gauge invariant.
An application requiring individually invariant
core members needs a gauge-covariant localization
in addition to the analytic filing proved here.

\subsubsection{The actual cell equation and its analytic vertices}

Write \(x=(A,\omega)\), with the strong norm
\(\|A\|_{X_\alpha}+\|\omega\|_\infty\) of
\eqref{hb:norms}. The actual equation is
\begin{equation}
 x=\mathcal L_0^{-1}\{\iota B+\mathcal R_h(x)\},
 \qquad
 \mathcal R_h(A,\omega)=
 \binom{-N_h(A)+(q_h'(A)-Q)^\sharp\omega}
       {-[q_h(A)-QA]}.
 \label{het:actual-cell-equation}
\end{equation}
Here \(\iota B=(0,B)\), not a fine point-mass
insertion. Each component in the first row is an
actual fine forcing with the physical density
appropriate to \(\sharp\). The second row is a
coarse bond field. In particular
\(\mathcal R_h(0)=D\mathcal R_h(0)=0\).

Choose a smooth partition \(\sum_y\chi_y=1\) on
physical unit cells, with a fixed support enlargement
and uniformly bounded derivatives through order
three. Sample it on the fine lattice. The discrete
product rule, including the local H\"older difference,
gives
\begin{equation}
 \|\chi_yA\|_{X_\alpha}\le C
                \|A\|_{X_\alpha(\operatorname{star}y)}.
 \label{het:stable-cell-localization}
\end{equation}
Indeed \(D(\chi A)=(D\chi)A+\chi(\cdot+h)DA\);
one further difference uses only bounded
\(\chi,D\chi,D^2\chi\), \(A,DA\) and the
local H\"older seminorm of \(DA\). All distances
are physical, so there is no negative power of
\(h\) in this multiplication bound.

Let \(E_yx\) consist of \(\chi_yA\) and the
coarse multiplier components assigned to cell \(y\).
Then \(x=\sum_yE_yx\), with bounded local overlap.
Let \(I_z\) inject the fine forcing on cell \(z\)
and its assigned coarse forcing. The fine input
norm here is \(L^\infty\), so characteristic cell
restriction of that forcing needs no derivative.
The actual block propagator
\begin{equation}
              \mathcal G_{yz}=E_y\mathcal L_0^{-1}I_z
 \label{het:actual-cell-propagator}
\end{equation}
satisfies
\begin{equation}
 \|\mathcal G_{yz}\|\le C_Ge^{-s_0d(y,z)},
 \qquad
 G_s:=\sup_y\sum_z e^{sd(y,z)}
                    \|\mathcal G_{yz}\|<\infty
       \quad(0<s<s_0).
 \label{het:propagator-cell-sum}
\end{equation}
This follows directly from the weighted joint
inverse \cref{prop:hb-joint-strong-inverse}, with
the weight centered on the input cell. The constants
include the fixed support of \(E_y\). The inverse
retains both the field and the multiplier.

There is a fixed physical radius \(R_0=C_L\)
such that the output \(\mathcal R_{h,z}\) reads
only the components \(E_yx\) with
\(d(y,z)\le R_0\). This assertion concerns the
actual nonlinear functions in
\eqref{het:actual-cell-equation}. For the Wilson
term it follows from its fine plaquette stencil
and its exact normal-coordinate differential.
The factor \(h\) in that differential cancels
the possible \(h^{-1}\) from a second fine
difference, as in \eqref{hb:normal-current-differential}.
For \(q_h\), the successive carrier radii are
bounded by \(C_L\sum_{j<k}hL^j\le C_L\).
In its physical transpose only coarse rows whose
carriers meet the output fine bond occur. The
fixed-leg column density \(Ch^4\) in
\eqref{hb:actual-constraint-moments} cancels
the \(h^{-4}\) in \(\sharp\). Thus the same
locality and uniform norm statement holds for
\((q_h'-Q)^\sharp\omega\).

On a fixed local strong ball the resulting Taylor
expansion is
\begin{equation}
 \mathcal R_{h,z}(x)=
 \sum_{m\ge2}\ 
 \sum_{y_1,\ldots,y_m\in\mathcal N(z)}
 T_{h;z;y_1\ldots y_m}
                    [x_{y_1},\ldots,x_{y_m}],
 \label{het:actual-local-Taylor}
\end{equation}
where \(\mathcal N(z)=\{y:d(y,z)\le R_0\}\).
The coefficient includes \(1/m!\). Uniform
Banach Cauchy differentiation on the larger
admitted local ball, followed by polarization,
gives
\begin{equation}
 \sum_{y_1,\ldots,y_m\in\mathcal N(z)}
        \|T_{h;z;y_1\ldots y_m}\|
                 \le C_Rr_0^{\,2-m}.
 \label{het:actual-vertex-bound}
\end{equation}
To see the constants explicitly, the number of
possible cells in each slot is at most
\(D_0=\#\mathcal N(z)\), independent of \(h\).
The norm of \(m!^{-1}D^m\mathcal R\) is at
most \(C(e/r)^m\) on a smaller ball. Absorb
\(D_0^m\), the localization multiplication
bound and the fixed two-component choices into
\(r_0\). No fine-coordinate count is made.
This proves \eqref{het:actual-vertex-bound} from
the actual local residual, rather than postulating
an analytic vertex family for the desired output.

\subsubsection{The rooted expansion and its connected carriers}

Substitute \eqref{het:actual-local-Taylor} into
\eqref{het:actual-cell-equation} repeatedly. A leaf
is the actual coarse input \(B(c)\) propagated
by \(\mathcal G\). An internal vertex has arity
\(m\ge2\), coefficient \(T_h\), and an outgoing
propagator. Order its children in their Taylor
slots. The resulting ordered rooted trees are
exactly the formal substitution expansion; the
Taylor factorial is already in \(T_h\) and
must not be inserted again.

Assign a carrier to each term as follows. An edge
\(\mathcal G_{yz}\) is carried by a fixed
coordinate-order shortest unit-cell path from
\(y\) to \(z\). A local vertex is carried by
its radius-\(R_0\) cell star, including the
short paths to all its input components. The
root support and leaf coarse bonds are included.
The resulting carrier is a connected cube union
containing every input value read by the term.
Its physical tree length is bounded by
\begin{equation}
 \ell(\mathcal T)\le
 C_0\sum_{\text{propagator edges }e}(1+d_e)
                         +C_1\#\mathcal T.
 \label{het:tree-carrier-length}
\end{equation}
All local stars have fixed size, so \(C_0,C_1\)
are independent of arities, mesh and depth.

Choose a positive \(t\) so small that the
propagator distance cost in
\eqref{het:tree-carrier-length} uses at most
\(s_0/2\) of its decay. The remaining constant
per edge and per input slot is absorbed by a
fixed decrease of \(r_0\) and increase of
\(G_s,C_R\). With these modified constants,
the absolute sum of trees rooted at one cell,
charged by \(e^{t\ell(\mathcal T)}\), is
bounded by the small solution of
\begin{equation}
 U=G_*\|B\|_\infty
            +G_*C_*\frac{U^2}{1-U/r_*}.
 \label{het:actual-tree-majorant}
\end{equation}
For example choose \(u_*<r_*/4\) and
\(8G_*C_*u_*<1\), and then
\(G_*\rho_*<u_*/4\). The right side maps
\([0,u_*/2]\) into itself and is contractive
there. Its solution is at most
\(2G_*\|B\|_\infty\) for
\(\|B\|_\infty<\rho_*\).

This scalar estimate bounds the actual tree sum:
at a fixed root sum each outgoing propagator
over its endpoint using
\eqref{het:propagator-cell-sum}; then sum its
local child-cell choices using
\eqref{het:actual-vertex-bound}; finally repeat
from its children. The resulting recursion is
exactly dominated by
\eqref{het:actual-tree-majorant}. All sums are
absolute. A tree with \(n\) leaves has at most
\(n-1\) internal vertices, since every internal
arity is at least two. Hence its bounded-star
and zero-length-edge costs introduce only a
fixed exponential per leaf, already included
in that majorant.

The summed tree series solves the actual fixed
point equation and lies in the uniqueness ball
of \cref{prop:hb-actual-critical-orbit}.
It therefore equals its constructed background.
The spatial rate \(t\) has not been divided
by the number of leaves.

For a torus, expand each flat propagator in its
lifted, twisted infinite-lattice entries. These
entries have the same exponential cell bound.
Their absolutely convergent periodization solves
the periodic joint system; uniqueness identifies
it with \eqref{het:actual-cell-propagator}.
This can also be obtained by taking the local
limit of its uniformly exponentially bounded
finite-volume inverse and then periodizing.
The nonlocal linear site projector is included
in that inverse, not in a purported local
nonlinear vertex. Each lifted propagator is
given its path on the cover before projection
to the torus. Summing its endpoint and period
image is the single infinite-lattice sum in
\eqref{het:propagator-cell-sum}. Thus the
carrier and entropy bounds are uniform in
volume and flat twists, including wrapping
trees.

\subsubsection{The actual Wilson action on these backgrounds}

Let \(a_{h,z}(A)\) be the sum of the actual
fine plaquette actions assigned to physical
cell \(z\), each plaquette assigned once.
The exact plaquette logarithm has the form
\(U_A(p)=e^{ih^2F_p(A)}\), with
\(\|F\|_\infty\le C(\|DA\|_\infty+
\|A\|_\infty^2)\). Hence on a fixed local
strong ball
\begin{equation}
 |a_{h,z}(A)|\le C\|A\|_{X_\alpha(\operatorname{star}z)}^2.
 \label{het:actual-cell-action-bound}
\end{equation}
Indeed each plaquette contributes at most
\(Ch^4\|F\|^2\), and the number assigned to
one physical cell is at most \(Ch^{-4}\).
The same bound holds on the complex ball by
the analytic cosine estimate, without a real
part. The local action has no constant or
linear Taylor coefficient.

Expand this actual cell action in its local
strong variables, and substitute the convergent
background trees in every slot. The outer
action vertex has arity at least two. Its
Banach Cauchy coefficients have the same form
of bound as \eqref{het:actual-vertex-bound}.
Consequently, on a smaller fixed ball, the
sum of resulting action trees rooted at \(z\)
satisfies
\begin{equation}
 \sum_{\mathcal T\ \mathrm{rooted\ at\ }z}
 e^{t\ell(\mathcal T)}
       \sup_{\|B\|_\infty<\rho}|s_{h,\mathcal T}(B)|
                         \le C\rho^2,
 \qquad 0<\rho\le\rho_*/2.
 \label{het:actual-action-tree-bound}
\end{equation}
The literal fine gauge of the homogeneous
background preserves Wilson action exactly.
Therefore this series is the actual effective
action \(S_k(V_B)\), rather than the action
of an approximate or separately chosen lift.

Fix any desired \(\kappa>0\). Choose an odd
localization size \(M\), a power of \(L\),
large enough that a path of length \(d\)
meets at most \(C(1+d/M)\) localization
cubes and \(C\kappa/M<t/4\). For each action
tree file its full carrier into its actual
union \(X\) of these cubes. Constants per
tree edge cost \(e^{C\kappa}\) and are
absorbed by a further fixed decrease of
\(\rho_*\), since the number of edges is
bounded by a fixed multiple of its leaves.
Write the resulting local function as
\(s_{k,M}(X;B_X)\).

\begin{proposition}[Uniform local analytic effective action]
\label{prop:het-actual-effective-action}
There is a fixed positive \(\rho_\kappa\),
independent of homogeneous depth and compatible
volume, such that on
\(\|B\|_\infty<\rho_\kappa\)
\begin{equation}
 S_k(e^{iB}V_0)=\sum_Xs_{k,M}(X;B_X),
 \qquad
 \sup_Q\sum_{X\ni Q}e^{\kappa d_M(X)}
       \sup_{\|B\|_\infty<\rho}|s_{k,M}(X;B_X)|
                         \le C_M\rho^2
 \label{het:actual-action-family}
\end{equation}
for \(0<\rho\le\rho_\kappa\). Each term
is holomorphic only in its indicated coarse
bond coordinates and is zero to order two at
the flat datum. The representation is formed
from the actual fixed-point coefficients and
their connected propagator carriers.
\end{proposition}

\begin{proof}
Only incidence at a cube different from the
outer action root needs an additional check.
If the carrier reaches that cube, its charged
length is at least the physical distance from
the action root to that cube, up to a fixed
cell margin. Reserve half the exponential
weight in \eqref{het:actual-action-tree-bound}
to pay this distance. Summing action roots
then costs at most
\(C\sum_z e^{-ctd(z,Q)}\le C_M\),
including the \(O(M^4)\) possible roots
inside \(Q\). The remaining weight pays
\(e^{\kappa d_M(X)}\) and the fixed tree
costs just described. This proves the bound.
Absolute normal convergence proves holomorphy
and the sum identity. Every leaf is an actual
input coarse bond in the filed carrier, so
the indicated local coordinate dependence is
literal. Each action tree has at least two
leaves. Its zero of order two survives filing.
\end{proof}

\subsubsection{Marked leaves and the actual effective Hessian}

Taylor coefficients can be marked before
filing. Fix a coarse bond \(b\) and mark one
leaf there. Transfer the root to that leaf
by reserving an exponential weight for the
path to the action root. Sum the remaining
root positions as in the preceding proof.
Marking a second leaf at \(b'\) charges its
distance to \(b\) by the same connected
tree. The number of ordered choices of
marked leaves is at most \(n(n-1)\) at
order \(n\); its polynomial factor is
summable on a fixed smaller analytic ball.
Since each coarse cell has a fixed finite
number of bond and colour coordinates,
these are absolute kernel sums, not only
multilinear operator-norm estimates.

In particular the actual Hessian admits
a local symmetric family \(H_{k,X}(B)\),
obtained by differentiating the filed action
terms and, if necessary, symmetrizing each
carrier term. For a fixed smaller rate and
radius,
\begin{align}
 S_k''(B)&=\sum_XH_{k,X}(B),\notag\\
 \sup_b\sum_{X\ni b}e^{\kappa d_M(X)}
       \sum_{b'}\|H_{k,X}(B)(b,b')\|&\le C_M,\notag\\
 \sup_b\sum_{X\ni b}e^{\kappa d_M(X)}
       \sum_{b'}\|[H_{k,X}(B)-H_{k,X}(0)](b,b')\|
                                      &\le C_M\|B\|_\infty.
 \label{het:actual-filed-Hessian}
\end{align}
All rows and field variables read by a term
lie in its full carrier. The column estimates
are identical by symmetry. More finely one
can keep the physical tree length in these
estimates, with a fixed positive rate, by
reserving it before passing to the \(M\)-grid.

This conclusion uses the all-order rooted
expansion of the actual residual. Fixed-order
max-distance response estimates alone would
not justify the absolute marked-leaf sums
or the unaltered spatial rate here.

% END LAYERED EXPANSION formation/homogeneous_effective_action_trees.tex

% BEGIN LAYERED EXPANSION formation/homogeneous_absolute_formation.tex
\subsection{Absolute-radius formation at homogeneous depth}
\label{layered:homogeneous-absolute-formation}

We now use the actual effective-action expansion of
\cref{prop:het-actual-effective-action}. The external
variables in this section are coarse bond logarithms
in its fixed ball about a specified flat reference.
This proves a homogeneous physical range and its
local analytic extension in those coordinates.
Independent fine-current cores and proper-region
background towers are separate domains.
Choose the frozen flat coarse reference
\(V_0\) in its outer block-axial
representative, so its actual internal
contour means are one. Every flat
holonomy class has such a representative:
the node-fixed gauge \(T_{y,x}(V_0)\)
sets those means to one and preserves
the next coarse field. This is a fixed
reference choice before the Taylor
expansion. Consequently the later
axializing gauge of a nearby field is
small in that same fixed coordinate
ball. A parallel-frame calculation
alone would not replace this condition
on the actual absolute contour gauge.
The norms below are analytic spatial and
pointing bounds. The fixed-reference
filing does not establish individual
gauge invariance of its terms, so it
does not by itself give the termwise
PH marginal cancellation or transport
of older gauge-invariant core families.
The torus periods are divisible by the chosen
localization sizes and exceed ten such cubes.
The estimates are uniform on this compatible
volume class.

\subsubsection{The uniform flat fluctuation gap}

Let \(Q_h\) be the actual flat straight row in
\eqref{hb:actual-flat-rows}, with \(h=1/n\).
Its coarse curl is an actual positive surface
average of the fine curl:
\begin{equation}
 \begin{split}
 (\mathrm{d}_cQ_ha)_{\mu\nu}(y)
 &=h^6\sum_{x\in B(y)}\sum_{r,s=0}^{n-1}
       (\mathrm{d}_ha)_{\mu\nu}(x+rh e_\mu+sh e_\nu)\\
 &=:(T_hd_ha)_{\mu\nu}(y),\qquad
                         \|T_hF\|_2\le\|F\|_{2,h}.
 \end{split}
 \label{haf:actual-flat-curl-contraction}
\end{equation}
The identity follows by telescoping the difference
of the two straight rows. The sum of its positive
weights in one coarse row is
\(h^6n^6=1\). A fixed fine plaquette has exactly
\(n^2\) possible choices of \((r,s)\); for each
choice its starting point belongs to exactly one
block. Thus its column sum is \(h^6n^2=h^4\).
Jensen's inequality proves the contraction, with
the physical fine measure. Parallel unitary
transports in a flat reference preserve this
argument, including arbitrary flat twists.

Let \(\mathcal H_k=S_k''(0)\). The derivative
of the actual background at zero minimizes the
flat curl energy subject to \(Q_ha=v\).
Indeed its physical stationarity after removal
of the gauge penalty is
\(\mathrm{d}_h^\sharp \mathrm{d}_ha=Q_h^\sharp\omega\), so
the energy cross term against any other lift
with the same row value is zero. The literal
fine gauge does not change that energy.
Equation \eqref{haf:actual-flat-curl-contraction}
and the smooth exact row seed
\cref{lem:hb-row-seeds} consequently give
\begin{equation}
 \|\mathrm{d}_cv\|_2^2\le\langle v,\mathcal H_kv\rangle
      =\inf_{Q_ha=v}\|\mathrm{d}_ha\|_{2,h}^2
                             \le C\|v\|_2^2.
 \label{haf:actual-flat-effective-form}
\end{equation}
For the upper bound insert its exact lift
\(E_Qv\), whose physical \(H^1\) energy is
bounded uniformly in \(h\). No inverse of a
coarse point-mass readout is used.

Let \(C_0\) be the actual one-step pivot
inclusion at the coarse flat field, using
the plain outer average and its middle pivots.
Let \(K_{G,0}\) be its actual contour-gauge
quadratic form. Therefore
\begin{equation}
 K_{k,0}:=C_0^{\mathsf T}\mathcal H_kC_0+K_{G,0}
 \ge C_0^{\mathsf T}\mathrm{d}_c^\sharp \mathrm{d}_cC_0+K_{G,0}
 \ge\gamma_0 I,\qquad \|K_{k,0}\|\le B_0.
 \label{haf:uniform-flat-fluctuation-gap}
\end{equation}
The last lower bound is the proved actual first
flat precision gap; the outer constraint here
has not been replaced by the straight row.
The bounds on \(C_0,K_{G,0}\) depend only on
\(L\). Thus \(\gamma_0,B_0\) are independent
of the inner homogeneous depth.

\subsubsection{The complete nonlinear effective precision}

For a coarse background \(V=e^{iB_0}V_0\)
inside a fixed smaller ball, let
\(V(u)=e^{i\Phi_V(u)}V\) be the actual one-step
outer pivot graph. All its unscaled radii and
support bounds are those already used in
\cref{layered:first-step-absolute-formation}.
Define the centered composed action
\begin{equation}
 \begin{split}
 \mathscr S_{k,V}(u)&=S_k(V(u))-S_k(V)
                       -D_u[S_k(V(u))]_{u=0}u,\\
 K_{A,k}(B_0)&=D^2\mathscr S_{k,V}(0),\qquad
 K_k(B_0)=K_{A,k}(B_0)+K_G(V).
 \end{split}
 \label{haf:actual-composed-precision}
\end{equation}
Use the explicit relative contour gauge of
\eqref{afs:relative-contour-gauge} for its
auxiliary extension. At an actual axial
constrained-critical \(V\), the subtracted
linear functional is zero, so this is
precisely the physical fluctuation action.

The composition in
\eqref{haf:actual-composed-precision} includes
both inner response terms in \(D^2S_k\)
and the outer constraint-curvature term
\(DS_k\,D^2\Phi_V\). To control it, substitute
the finite local holomorphic map
\[
 B(u)=i^{-1}\log[V(u)V_0^{-1}]
\]
in the actual action trees. Its coefficients
have a fixed physical star and uniform
Cauchy bounds depending only on \(L\).
The marked-leaf proof of
\eqref{het:actual-filed-Hessian} therefore
gives a symmetric local operator family
\(K_{k,X}(B_0)\), with every matrix endpoint
and every background input in its carrier.
For a fixed \(t_0>0\), its absolute
physical-length Schur sums obey
\begin{align}
 K_k(B_0)&=\sum_XK_{k,X}(B_0),\notag\\
 \sup_b\sum_{X\ni b}e^{t_0\ell(X)}
             \sum_{b'}\|K_{k,X}(B_0)(b,b')\|&\le C_K,\notag\\
 \sup_b\sum_{X\ni b}e^{t_0\ell(X)}
  \sum_{b'}\|[K_{k,X}(B_0)-K_{k,X}(0)](b,b')\|
                                      &\le C_K\|B_0\|_\infty.
 \label{haf:actual-precision-carriers}
\end{align}
Here \(X\) is retained as a connected physical
unit-cell carrier, before the larger filing
grid is introduced. Retain its lifted cover
tree as a tag, and use its charged length
as \(\ell(X)\); tags with the same cubical
union can be combined after the norm sum.
This also retains winding lengths before
projection. The transpose column
sums satisfy the same bounds. Local contour
gauge terms are included as additional
fixed-star terms.

These inequalities follow from the actual
tree coefficients, not from a hypothesis
on the desired precision. A marked second
derivative contributes at most a polynomial
in its leaf number; summability follows on
the smaller ball. The first-derivative
factor in the outer curvature term vanishes
at the flat field and is \(O(\|B_0\|)\).
Thus after decreasing the fixed real ball,
\begin{equation}
            \gamma I\le K_k(B_0)\le BI
 \label{haf:actual-real-gap}
\end{equation}
uniformly in depth. Complex differences
have the corresponding Schur Lipschitz
bound on a fixed smaller ball. The same
carrier estimates hold for
\(Q_{k,\tau}=\tau K_{A,k}+K_G\), for
\(\tau\) in a fixed disc about one.

\subsubsection{Admitting cube parameters for the nonlocal precision}

We construct the local Gaussians rather than
declaring the nonlocal matrix finite-range.
Set \(H=I-Q_{k,\tau}/a\), with \(a\ge2B\).
The identity has a one-cell carrier. Take
the actual carrier expansion induced by
\eqref{haf:actual-precision-carriers}, and
write
\begin{equation}
 H=H_{\le R}+H_{>R},
 \qquad H_{\le R}=\sum_{\ell(X)\le R}H_X.
 \label{haf:actual-near-tail}
\end{equation}
Every entry of \(H_{\le R}\) reads only
fields within radius \(R+C_L\) of its
endpoints and has that finite matrix range.
For real data symmetrize each carrier term
if necessary. Its tail has Schur norm at
most \(Ce^{-t_0R}\). For large \(R\),
\eqref{haf:actual-real-gap} therefore implies
\begin{equation}
 0\le H_{\le R}\le r'I,
 \qquad r'=1-\gamma/(2a)<1.
 \label{haf:near-spectral-bound}
\end{equation}
The lower bound uses
\(I-Q/a\ge I/2\) and a tail smaller than
\(1/4\). All constants are uniform.

Fix the requested coarse parameter radius
\(S>2\). For a connected physical carrier,
round it to localization cubes of side
\(M\). The geometric covering-tree bound
gives
\[
 S^{\#X_M}\le S^{C_0}
          \exp\{C_0(\log S)\ell(X)/M\}.
\]
Choose \(M\) large enough that the rate in
this exponent is below \(t_0/4\). In the
operator-family norm charging this cube
weight, \eqref{haf:actual-precision-carriers}
then gives constants \(C_H,c>0\) such that
\begin{equation}
 \|H\|_{\rm car},\ \|H_{\le R}\|_{\rm car}\le C_H,
 \qquad \|H_{>R}\|_{\rm car}\le C_He^{-cR}.
 \label{haf:actual-weighted-tail}
\end{equation}
The norm is the maximum absolute row/column
sum of all its carrier terms, with the
displayed weight. Operator composition
files into the union of two carriers
sharing their intermediate bond. Thus
these unions are connected, their weights
are submultiplicative, and this norm is
submultiplicative as well.

The Gaussian parameters themselves are
placed on the coarse \(LM\)-cube grid
of the UM expansion. The displayed
\(M\)-cube carrier weight is stronger
than the needed \(LM\)-cube weight.
On the fixed held root guard use
\(\widehat s_\Delta=1\), and use
\(\widehat s_\Delta=s_\Delta\) elsewhere
in every carrier weight below.

Choose an integer \(m\), and then
\(R=Dm\), where the fixed \(D\) is large
enough that \(cD>2\log(1+C_H)+2\).
Use the exact identities
\begin{align}
 H^m&=(H_{\le R})^m+E_m,\notag\\
 E_j&=H^j-(H_{\le R})^j
   =\sum_{i=0}^{j-1}(H_{\le R})^iH_{>R}H^{j-1-i}.
 \label{haf:exact-short-power-tail}
\end{align}
The second line defines an actual filed
carrier series. It gives, simultaneously
for \(j\le m\),
\begin{equation}
 \|E_j\|_{\rm car}\le j C_H^je^{-cR}.
 \label{haf:short-power-tail-bound}
\end{equation}
Combine all paths inside the entries of
\((H_{\le R})^m\) before assigning their
parameters. At one common real background the
Hilbert-row estimate of
\eqref{afs:power-Schur-bound} gives a
pointwise Schur bound with the factor
\((1+m(R+C_L))^2\). For a local coefficient
family, take the supremum separately for each
entry over its full real read set. Denote the
maximum of this row sum and its transpose
column sum by \(\|\cdot\|_{{\rm Sch},+}\).
The appropriate bound is
\begin{equation}
 \|(H_{\le R})^m\|_{{\rm Sch},+}
       \le C(1+m(R+C_L))^4(r')^m.
 \label{haf:combined-near-short-power}
\end{equation}
Indeed each matrix entry is at most \((r')^m\)
by \eqref{haf:near-spectral-bound}, and a row
has at most \(C(1+m(R+C_L))^4\) possible
endpoints. Every permitted local real read set
extends by zero relative to \(V_0\) to admitted
global real data. Thus the entry estimate remains
valid when its supremum is taken independently
of the other entries. A supremum outside the
whole row sum would not establish this bound.

For the uncombined families \(H_X\), their tails
and the words defining \(E_j\), use the sums of
individual carrier suprema on a common larger
complex ball. The positive Taylor-tree
majorants proving
\eqref{het:actual-action-tree-bound} and
\eqref{het:actual-filed-Hessian} give precisely
these stronger bounds. Their products are
bounded by the positive products of those
majorants. Thus \(C_H\) in
\eqref{haf:actual-weighted-tail} and
\eqref{haf:short-power-tail-bound} can be
chosen for these norms as well.

Take \(M\) to be the smallest power of
\(L\) at least \(C_Lm(R+C_L)\). The
combined near entries then have connected
ball carriers meeting a bounded number
of \(M\)-cubes. Their cube weights cost
only \(S^{C_L}\). Weight the already filed
\(E_m\) terms individually by their full
carriers, and weight each combined near
entry by its full ball carrier. Denote
their sum by \(T_s\), symmetrizing its
carrier terms if needed. At \(s=1\),
\(T_1=H^m\) exactly. On all independent
\(|s_\Delta|<S\),
\begin{equation}
 \|T_s\|_{\rm Sch}\le
 C S^{C_L}(1+m(R+C_L))^4(r')^m
                    +mC_H^me^{-cR}=:q_m\longrightarrow0.
 \label{haf:actual-admitted-short-power}
\end{equation}
Here \(M=O_L(m^2)\). Its appearance in
the volume budget is polynomial, whereas
both terms on the right decrease
exponentially.

For each prefix \(j<m\), similarly put
\(A_{j,s}=(H_{\le R})^j+(E_j)_s\).
Only the near prefix is left unweighted;
its row from a root is contained in the
fixed held root guard. The tail prefix
is weighted by its actual full carriers.
Choose the guard to contain every
radius-\(m(R+C_L)\) root ball, using the
same fixed eight-cell collar convention
as before, with constants increased once.
With \(c_\ell=4^{-\ell}\binom{2\ell}{\ell}\), define
\begin{equation}
 R_W(s)=a^{-1/2}\chi_W\sum_{j<m} A_{j,s}
                         \sum_{v\ge0}c_{vm+j}T_s^v,
 \qquad D_W(s)=R_W(s)R_W(s)^{\mathsf T}.
 \label{haf:actual-decoupled-root}
\end{equation}
At complete parameters this is exactly
\(\chi_WQ_{k,\tau}^{-1/2}\), by the
grouped binomial identity. Thus its
covariance is the actual Gaussian marginal.
Every prefix-tail and short-step carrier
contains its endpoints and every external
coordinate read by that coefficient.
Zero outside a wall domain therefore
gives literal local dependence and wall
factorization, with fixed-root heredity.

For the real comparison use
\(R_{0,r}=a^{-1/2}\sum_{j<m}c_j
(H_{\le R,r})^j\). Its compression has
covariance at least \(a^{-1}I\), by
\eqref{haf:near-spectral-bound}. The
sum of the Schur norms of its prefixes
is at most a polynomial in \(R\).
Equations \eqref{haf:short-power-tail-bound}
and \eqref{haf:actual-admitted-short-power}
therefore make \(D_W(s)\) uniformly
\(o(M^{-4})\)-close to that real covariance
by increasing \(m\). Finally shrink the
fixed complex coarse-field and \(\tau\)
balls, after \(m,R,M\) have been fixed.
To justify this step coefficientwise, let
\(\delta_H^+\) be the positive carrier bound
for the difference between the complex near
matrix and the matrix at its real data.
Nested-ball Cauchy estimates on the actual
Taylor-tree majorants give
\(\delta_H^+\le C\delta\), where \(\delta\)
is the tube size. A positive word expansion
of the telescoping identity gives, for \(1\le j\le m\),
\[
 \Delta_j^+\le j C_H^{j-1}\delta_H^+.
\]
This is the same finite-power estimate as
\eqref{afs:polynomial-perturbation}; enlarge
\(C_H\ge1\) once if necessary. Choose the
tube so that
\(\sum_{j=1}^m\Delta_j^+\le(r')^m\).
This pays the whole prefix sum and the
\(m\)-th power perturbation. Increasing the
fixed constant in
\eqref{haf:actual-admitted-short-power}
absorbs the latter term, and the covariance
comparison remains \(o(M^{-4})\).
No bound polynomial in \(m\) for this
coefficientwise perturbation is required.
These radii are positive and independent
of depth and coupling. Local data in
this coarse product ball extend by zero
relative to \(V_0\), so no extra
curvature-extension premise is needed.
This proves the local Gaussian comparison,
enlarged independent cube discs and exact
structural identities needed by the
homogeneous unmarked expansion for the
actual nonlocal effective precision.

\subsubsection{The actual intrinsic and ordinary-source vertices}

Fix \(g=\theta^{-1/2}\) and
\(\mathfrak b=A_1(\log\theta)^p\), with
\(p>1/2\). Keep the real product cutoff
\(\chi(B)=\prod_b\mathbf1_{|B(b)|<\mathfrak b}\)
fixed throughout. Take \(\theta\) so large
that \(16C_Lg\mathfrak b\) is less than
both the outer pivot radius and the margin
in the actual coarse-data ball just proved.
Then \(V(gB)\), its normalized radial paths
and the needed complex fluctuation boxes
all remain in that ball. In particular
\(U_k(V(gB))\) is the actual admitted
homogeneous background on the entire box;
its existence has not been replaced by
an activity-norm assumption.

Compose each of the actual action terms
in \eqref{het:actual-action-family} with
the outer graph. Subtract its constant,
linear and quadratic Taylor terms in
\(u\), before summing. The global linear
sum is the centering term in
\eqref{haf:actual-composed-precision},
and the global quadratic sum is the
complete \(K_{A,k}\). Each composed
member is holomorphic on a fixed
\(\ell^\infty\) free-coordinate ball
and reads only the fixed chart collar
of its filed carrier. If its admitted
bound is \(B_X\), directional Cauchy
Taylor remainder gives
\begin{equation}
 |g^{-2}R_{3,X}(gB)|
 \le C B_Xg\max_{b\in X^+}|B(b)|^3
 \le C B_Xg\mathfrak b
                         \sum_{b\in X^+}|B(b)|^2.
 \label{haf:actual-effective-cubic-vertex}
\end{equation}
The first inequality includes the fixed
analytic-radius cost; the maximum is
over its actual chart support. The
second inequality is valid on the kept
cutoff and on its fixed complex margin
after a fixed change of constants.
Thus the actual Gaussian-growth
potential norm has amplitude
\(C_Mg\mathfrak b\), uniformly in depth.
No derivative count proportional to
the fine lattice size is present.

The outer relative contour gauge, Haar
and pivot coarea factors are the literal
one-step factors of
\eqref{afs:relative-contour-gauge} and
\eqref{afs:actual-coarea-density}, on
the present coarse lattice. Their
supports and unscaled derivative bounds
depend only on \(L\), so their intrinsic
remainders have the same amplitude.
The zero-field coarea value \(J_V(0)\)
stays in the Gaussian normalizer.

There is an actual all-depth ordinary
source class in this range. Take fine
plaquette profiles \(w(p)\), with
\(\|w\|_\infty<\rho_w\) on independent
fixed coefficient discs, and put
\begin{equation}
 H_{k,w}(V)=\sum_{p>0}w(p)a_p(U_k(V)).
 \label{haf:actual-ordinary-source}
\end{equation}
Its unit-cell action bound is
\(C\rho_w\|A\|_{X_\alpha}^2\), by
the same \(h^4\) plaquette count as
\eqref{het:actual-cell-action-bound}.
Apply the actual action-tree construction
with these weights at its outer action
vertices. It produces a genuinely local
analytic coarse-source family with
norm \(C_M\rho_w\), uniformly in depth
and in the number of fine plaquettes.
The literal fine gauge preserves each
plaquette trace separately, so this
family represents precisely
\eqref{haf:actual-ordinary-source}.

Its interaction is the exact difference
\(H_{k,w}(V(gB))-H_{k,w}(V)\).
Cauchy differentiation on the already
admitted coarse-data margin gives its
potential amplitude
\(C_M\rho_wg\mathfrak b\).
The background, Gaussian, gauge weight
and sharp cutoff have not been moved
when inserting this ordinary source.
Additional inputs already supplied as
local holomorphic coarse-data families
with incidence norm \(H_0\) obey the
same difference bound \(C_MH_0g\mathfrak b\).
This last optional statement does not
assert that arbitrary older fine-core
families have already been transported
into that class.

\subsubsection{The formed homogeneous families and their physical range}

\begin{proposition}[Actual homogeneous formation near a flat reference]
\label{prop:haf-actual-homogeneous-formation}
There are fixed positive coarse-data
and \(\tau\) radii and a coupling floor,
independent of homogeneous depth and
compatible volume, for which the actual
regular homogeneous step with action
\(xS_k\), contour-gauge weight \(\theta\)
and the ordinary sources
\eqref{haf:actual-ordinary-source}
has an analytic filed decomposition
\(\mathbf N_k^\circ+\mathbf G_{k,\tau}
+\mathbf F_k\). For every fixed requested
output rate its constructed pointing
satisfies
\begin{equation}
 \|\mathbf N_k^\circ\|_+\le C_N,\qquad
 \|\mathbf G_{k,\tau}\|_+\le C_G|\tau-1|,\qquad
 \|\mathbf F_k\|_+\le
 C_F\{(1+\rho_w+H_0)g\mathfrak b
                           +e^{-c\mathfrak b^2}\}.
 \label{haf:actual-homogeneous-F-bound}
\end{equation}
Here \(H_0=0\) unless the optional already
filed coarse-input class is used.
The constants are independent of depth.
The normalizer and mismatch are retained
as Gaussian terms, not counted in
\(\mathbf F_k\).
\end{proposition}

\begin{proof}
The local covariance construction
\eqref{haf:actual-decoupled-root} supplies
all enlarged cube discs, Gaussian
comparison, normalization, fixed-support
heredity and wall factorization needed
in the actual UM expansion. In each
functional take all free bonds of its
full fine support, so its held guard
also contains any background inputs
constant along that fibre. Its scalar
dimension is at most \(C(LM)^4\#Z\).
The comparison error can be chosen
below \(c_*(LM)^{-4}\), without an
auxiliary-volume factor.

The actual interaction estimates just
proved give the UM potential budget
\(C_M(1+\rho_w+H_0)g\mathfrak b\).
The unchanged sharp cutoff gives the
Gaussian tail \(Ce^{-c\mathfrak b^2}\).
The already proved finite Mayer,
connected cluster and union-filing
theorems therefore give the stated
\(\mathbf F_k\) bound. The normalized
cubic, ordinary radial and moving-
threshold pointings are exactly those
used in the first-step proof; their
analytic margins and budgets are now
uniform in depth.

For the Gaussian normalizer, use the
actual carrier expansion of
\(\log Q\), grouping
\(\log Q=(\log a)I-
\sum_{j<m,v\ge0,vm+j>0}H^j(H^m)^v/(vm+j)\).
The near combined powers and the filed
tails have precisely the summable cube
weights proved above. Root the trace
at its actual diagonal bond and add
the local zero-field pivot terms.
This proves the \(\mathbf N_k^\circ\)
bound after vacuum subtraction.
The same construction on a larger
\(\tau\) disc gives the mismatch bound;
its pointing is its actual resolvent
insertion as in
\eqref{afs:actual-mismatch-pointing}.
\end{proof}

The actual source functions are
holomorphic on independent fixed fine
profile discs. The same argument as
\cref{cor:afs-ordinary-source-jets}
therefore gives
\[
 |D_w^nF_k(X;w)[h_1,\ldots,h_n]|
 \le Cg\mathfrak b\,n!(2e/\rho_w)^n
       e^{-\kappa d_c(X)}\prod_j\|h_j\|_\infty
\]
on the half-radius source ball, with
the reserved rate for pointed terms.
There is no factor counting fine
plaquettes or source profiles in this
bound.

\subsubsection{The actual nested physical step}
\label{haf:physical-nesting}

To identify the complete physical step,
take \(W=e^{iB_+}V_{0,+}\), where
\(V_{0,+}=M(V_0)\) and \(\|B_+\|_\infty\)
is sufficiently small, and form its actual
depth-\(k+1\) background \(U_+(W)\).
Choose the flat fine extensions
compatibly at the two depths, so the
intermediate flat reference is the
partial average of that same flat
fine field.
The partial averaging recursion and
its raw bound \eqref{hb:uniform-raw-domain}
put \(V=M^kU_+(W)\) in the coarse-data
ball above. Its node-fixed one-step
axial gauge is also there after a
fixed decrease of the outer radius.
The required gauge extension is the
bounded local node extension from the
homogeneous background construction.
More explicitly, truncating that
recursion at its intermediate spacing
\(1/L\) gives
\(\|a_k\|_\infty+\|v_k-I\|_\infty
\le C_L\|B_+\|_\infty/L\).
On rescaling this spacing to one,
the fine potential is \(A/L\) and
its first difference is \(DA/L^2\).
The one-step axial node values have
size \(C_L\|B_+\|_\infty\); their
actual smooth fine extension has
the same bound in the required strong
gauge norm. Thus the intermediate
field and its axial representative
remain within \(C_L\|B_+\|_\infty\)
of the fixed strong orbit chart,
uniformly in \(k\).

This intermediate field is an actual
constrained-critical point for \(S_k\)
under \(M(V)=W\). Indeed every fine
variation with \(D(M^k)\delta U=0\)
also has \(D(M^{k+1})\delta U=0\).
Thus \(U_+(W)\) is critical on its
\(M^k\)-fibre. Both it and the
constructed \(U_k(V)\) are in the
fixed strong regularity class, after
rescaling the intermediate coarse
spacing by the fixed factor \(L\).
Orbit uniqueness in
\cref{prop:hb-actual-critical-orbit}
identifies their actions. To check
outer stationarity, lift any admitted
curve \(V(t)\) with \(M(V(t))=W\)
by its actual background \(U_k(V(t))\).
Then
\(M^{k+1}U_k(V(t))=W\) exactly.
At \(t=0\) its background is the
same critical orbit as \(U_+(W)\),
so
\(\partial_t A(U_k(V(t)))|_{t=0}=0\).
This proves that \(DS_k(V)\) annihilates
the outer constraint tangent. The linear
centering in
\eqref{haf:actual-composed-precision}
therefore vanishes physically.

On \(|B|<\mathfrak b\), the actual
\(V(gB)\) stays in the same fixed
coarse-data ball, so all inner
backgrounds exist on one admitted
chart. The outer regular-fibre anchor,
coarea density and contour-gauge test
removal are the literal one-step
proofs already given; their constants
do not involve \(k\). The exact
physical normalization is therefore
\eqref{afs:exact-physical-step} with
\(A(U)\) replaced by \(S_k(V)\), the
present complete precision, and its
actual zero-field pivot factor.
At these critical axial pairs gauge
invariance of the actual effective
action supplies the orbit-null action
sector. Thus the physical mismatch
total is again \((r/2)\log\tau\),
with the explicit numerical
\(\tau^{-d/2}\) retained separately.

This physical range is a fixed bond
ball about a frozen flat reference.
Extra history plaquette indicators
still require their own threshold
margin. The proof does not extend the
local functions across arbitrary
flat-reference changes, to all
curvature-only fields, to independent
fine currents, or to proper-region
and history towers. Those conclusions
are not prerequisites for the
homogeneous near-flat step proved here.

% END LAYERED EXPANSION formation/homogeneous_absolute_formation.tex

% BEGIN LAYERED EXPANSION formation/homogeneous_physical_pointing.tex
\subsection{Gauge invariance of the complete homogeneous pointing}
\label{layered:homogeneous-physical-pointing}

The fixed-reference Taylor terms in
\cref{prop:het-actual-effective-action} need not
be individually gauge invariant. A complete
physical pointed density can nevertheless be
constructed invariantly: specify its insertions
at the original fine plaquettes before making
the Taylor-tree filing. This distinction is
needed for its quadratic Ward identity.
The result below concerns the homogeneous
near-flat range already proved, with the
literal fine plaquette sources and no
untransported older fine-core inputs.
Here \(V\) is the input coarse
background of that one-step integral.
The actual output datum \(W\) enters
through the critical intermediate
background \(V=p_k(W)\), considered
explicitly below.

\subsubsection{Invariant functions before filing}

For each original fine plaquette put
\begin{equation}
 f_{k,p}(V)=a_p(U_k(V)),\qquad
 f_{k,p;V}(u)=f_{k,p}(e^{i\Phi_V(u)}V).
 \label{hpt:actual-plaquette-functions}
\end{equation}
These are the actual functions, before
decomposition into connected carriers.
For every sufficiently small coarse
site gauge \(g\) whose two data lie
in the admitted ball, the literal
node extension and orbit uniqueness
of \cref{prop:hb-actual-critical-orbit}
give
\begin{equation}
 f_{k,p}(V^g)=f_{k,p}(V).
 \label{hpt:individual-physical-plaquette-invariance}
\end{equation}
Indeed a fine gauge with those node
values takes the critical field to
the same regular orbit as
\(U_k(V^g)\). A fine gauge leaves the
trace of each plaquette, at its fixed
position \(p\), unchanged. This is
an assertion about the nonlocal
function \(f_{k,p}\), not its local
Taylor pieces.

Write \((O_gu)(b)=\operatorname{Ad}(g(b_-))u(b)\)
on outer free bonds. The actual pivot
equation, its relative coarse logarithm
and its free-identity convention give
\begin{equation}
 \Phi_{V^g}(O_gu)=O_g\Phi_V(u),\qquad
                f_{k,p;V^g}(O_gu)=f_{k,p;V}(u).
 \label{hpt:joint-chart-invariance}
\end{equation}
The first identity follows by covariance
and uniqueness of the same small
pivot graph. At real unitary \(g\),
\(O_g\) is a real block-orthogonal
map on the actual integration space.
It preserves each retained cutoff.

Let \(R_{3,p;V}(u)\) be the third-order
remainder of \(f_{k,p;V}\), subtracting
its own constant, linear and quadratic
jets. Each homogeneous jet transforms
under the same linear map \(O_g\).
Consequently
\(R_{3,p;V^g}(O_gu)=R_{3,p;V}(u)\).
This gives invariant cubic remainders
before their carrier expansion.

The other functions have the same
property directly. Every relative
contour-gauge trace conjugates at
its node. The Haar density is a
function of the invariant colour
pairing of \(\Phi\). A pivot matrix
transforms on its two sides by
adjoint matrices of determinant
one; its determinant ratio and
chosen logarithm are invariant.
The ordinary source
\(\sum_pw(p)f_{k,p}\) is invariant
with its fine profile held fixed.

\subsubsection{The actual normalization shares and insertions}

For each fine plaquette \(p\), choose
nonnegative fixed numbers
\(\vartheta_{p,b}\) on outer free
bonds with
\begin{equation}
 \sum_b\vartheta_{p,b}=1,\qquad
 \vartheta_{p,b}=0\quad\hbox{if }
             d(\operatorname{mid}p,\operatorname{mid}b)>C_L.
 \label{hpt:physical-plaquette-shares}
\end{equation}
One choice is equal shares on the
nearest free-bond midpoints. Their
distance is uniformly bounded because
each outer \(L\)-block contains free
bonds. This choice is covariant under
the geometric symmetries preserving
the actual pivot pattern. It is
independent of the field and its
Taylor-carrier presentation.

Put \(P_{p,V}(B)=-\tau g^{-2}R_{3,p;V}(gB)\).
The normalized intrinsic interpolation
and its actual bond insertion are
\begin{equation}
 \begin{split}
 P_{p,V;t}(B)&=t^{-2}P_{p,V}(tB),\\
 a_{p,t,b}(V;B)&=
    t^{-2}B(b)\cdot D_bP_{p,V}(tB)
       -2\vartheta_{p,b}t^{-3}P_{p,V}(tB).
 \end{split}
 \label{hpt:actual-invariant-intrinsic-insertion}
\end{equation}
The shares are specified at the
original plaquette, before any
carrier filing. For independent
parameters the defining function is
\[
 \left(\sum_b\vartheta_{p,b}t_b\right)^{-2}
                    P_{p,V}((t_bB(b))_b).
\]
Its derivative at \(t_b=t\) is
exactly the displayed insertion.
Thus its sum over \(b\) is the
derivative of the normalized diagonal
path, including its normalization
term. Nonnegative shares keep the
denominator at least \(t/2\) on
\(\max_b|t_b-t|<ct\), \(c<1/2\).
All the field arguments remain in
the proved common chart.

Equation \eqref{hpt:joint-chart-invariance}
and invariance of the colour pairing
show, for every bond separately,
\begin{equation}
 a_{p,t,b}(V^g;O_gB)=a_{p,t,b}(V;B).
 \label{hpt:bondwise-insertion-invariance}
\end{equation}
The same construction is applied
to the original relative gauge
remainders with shares on their
fixed local free supports. The
Haar, pivot and ordinary-source
legs use their ordinary radial
insertions \(B(b)\cdot D_bV(tB)\),
which satisfy the same identity.

There is a uniform rooted bound
for these actual insertions.
Retain each \(p\) as the outer
action-tree tag. Its local analytic
bound is \(Ch^4\), before summing
the fine plaquettes. The proof of
\eqref{haf:actual-effective-cubic-vertex}
then gives its intrinsic amplitude
\(Ch^4g\mathfrak b\). A derivative
rooted at \(b\) marks a leaf in
that actual tree. The exponential
root-transfer estimate sums these
terms over the plaquette's physical
cell positions. The normalization
root in \eqref{hpt:actual-invariant-intrinsic-insertion}
is within \(C_L\) of \(p\) by
\eqref{hpt:physical-plaquette-shares}.
It has the same carrier budget.
Finally there are \(O(h^{-4})\)
fine plaquettes per physical cell,
cancelling their \(h^4\) bound.
The order-three radial vanishings
and \eqref{fm:constructed-pointing-integrability}
therefore give the same
\(Cg\mathfrak b\) integrable rooted
norm as the homogeneous formation
theorem, with a fixed rate reserve.
This step uses the absolute marked-
tree sums, rather than differentiating
only an unpointed output estimate.

\subsubsection{Invariant normalized expectations}

The complete precision of the actual
composed action and relative contour
gauge transforms on the shared domain
as \(Q_\tau(V^g)=O_gQ_\tau(V)O_g^{\mathsf T}\).
For real unitary gauges, Gaussian
change of variables by \(O_g\),
the fixed product cutoff and
\eqref{hpt:bondwise-insertion-invariance}
therefore give invariance of each
normalized bond insertion expectation.
Every interpolation partition function
is zero-free in the ordinary cluster
domain proved above. Its canonical
logarithm and expectation are thus
unambiguous there. The identities
are holomorphic in the external
data and extend on their common
complex domain.

The cutoff leg is invariant as well.
Its one-bond moving threshold and
its boundary surface measure depend
only on \(|B(b)|\), which the real
orthogonal map \(O_g(b)\) preserves.
Its rooted Gaussian tail estimate
is the already proved one. The free
mismatch leg has its actual density
\((\tau-1)\int\operatorname{tr}(K_GQ_{\tau,u}^{-1})(b,b)du/2\);
the diagonal colour trace is invariant.

Let \(\mathcal E_b^F(V)\) be the
integral of the cutoff and interaction
insertions just specified. The exact
path identities give
\(\sum_b\mathcal E_b^F=\mathcal F_k\).
Use the fixed midpoint tents for
the output base point:
\begin{equation}
 F_{k,z}^{\rm phys}(V)=
                  \sum_bh_z(\operatorname{mid}b)\mathcal E_b^F(V).
 \label{hpt:actual-physical-pointed-density}
\end{equation}
The tents \(h_z\) are on the outer
output-node lattice, of spacing
\(L\) in the present input units.
After rescaling, the same \(z\)
is the output unit-lattice basepoint
in the composition below.
This complete physical density is
gauge invariant. Its local analytic
filing obtained from the actual
insertions has the uniform bound
\(Cg\mathfrak b\) of
\eqref{haf:actual-homogeneous-F-bound}.
The individual filed functions still
need not be gauge invariant.

\begin{proposition}[The physical quadratic Ward identity]
\label{prop:hpt-physical-quadratic-Ward}
At a trivial flat reference, the
complete density
\eqref{hpt:actual-physical-pointed-density}
has zero first derivative in coarse
bond logarithms. Its quadratic
kernel annihilates a coarse gradient
in either argument:
\begin{equation}
 D^2F_{k,z}^{\rm phys}(0)[\mathrm{d}_c\lambda,a]=0.
 \label{hpt:actual-quadratic-Ward}
\end{equation}
Its second spatial moments are
bounded by \(Cg\mathfrak b\),
uniformly in homogeneous depth.
\end{proposition}

\begin{proof}
A constant colour conjugation fixes
the trivial field. An invariant
linear functional in one adjoint
colour is zero, so
\(DF_{k,z}^{\rm phys}(0)=0\).
Differentiate the proved local gauge
identity in a small gauge parameter
and then in the field. The linear
gauge variation at zero is
\(\mathrm{d}_c\lambda\). The nonlinear
variation multiplies the vanished
first derivative. This proves
\eqref{hpt:actual-quadratic-Ward};
symmetry gives the second argument.
One may first use finite-support
gauges. The exponential marked-tree
bounds then extend the identity to
the stated finite-moment class.

The filed insertion bounds are
holomorphic on a fixed external
coarse-coordinate ball. Two Cauchy
derivatives cost its fixed inverse
radius squared. A carrier has at
most \(C(LM)^4(1+d_c(X))\) outer
coarse bond coordinates. Summing
all derivative pairs and multiplying
by their squared spatial distances
therefore costs only a fixed
polynomial in \(d_c(X)\), with an
\(M\)-dependent constant and no fine
depth factor. The reserved spatial
rate makes that polynomial summable.
Thus the second moments have
the asserted \(g\mathfrak b\) bound.
In particular an affine curvature
probe may be inserted in this
quadratic kernel by its absolutely
convergent spatial moments.
\end{proof}

\subsubsection{Composition with the actual critical background}

Let \(p_k(W)\) be the intermediate
coarse field \(M^kU_{k+1}(W)\), in
its actual one-step axial gauge,
as constructed in the physical-range
argument above. Define the complete
output density, now in output coarse
units, by
\begin{equation}
 \mathscr F_{k,z}(W)
             =F_{k,z}^{\rm phys}(p_k(W)).
 \label{hpt:actual-output-composition}
\end{equation}
This is the density for the actual
output datum \(W\). Gauge invariance
makes it independent of the admitted
intermediate representative.

\begin{corollary}[Physical homogeneous marginal bound]
\label{cor:hpt-physical-marginal-bound}
At the trivial flat reference and
zero ordinary profile, the complete
output densities
\eqref{hpt:actual-output-composition}
retain the quadratic Ward identity
and the uniform rooted second-moment
bound \(Cg\mathfrak b\).
Their infinite-lattice companion has
the physical Yang--Mills quadratic
coefficient
\begin{equation}
                 |b_k^{\rm phys}|\le Cg\mathfrak b,
 \label{hpt:actual-physical-marginal-bound}
\end{equation}
with the Wilson normalization of
\eqref{ph:action-normalization}.
This coefficient belongs to the
complete physical pointed density.
\end{corollary}

\begin{proof}
The map \(p_k\) is analytic on the
fixed output ball. Its actual first
response \(H_k=Dp_k(I)\), written in
potential coordinates, has uniform
exponential row and column moments.
Indeed it is the localized response
\eqref{hb:localized-background-response}
followed by a partial actual average
and its local axial gauge. Each
partial row has bounded physical
carrier and bounded directional
norm by \eqref{hb:uniform-raw-domain}.
There are only \(L^4\) intermediate
cells per output cell. Thus summing
its responses costs a fixed \(L\)
factor, not a fine-depth factor.

Write \(\mathcal K_{k,z}\) for the
quadratic kernel already proved for
\(F_{k,z}^{\rm phys}\). The first
jet vanishes, so the second-order
chain rule is exactly
\begin{equation}
 D^2\mathscr F_{k,z}(I)[a,b]
       =\mathcal K_{k,z}[H_ka,H_kb].
 \label{hpt:actual-composed-quadratic-kernel}
\end{equation}
No \(D^2p_k\) term remains.
The kernel moments on the right
are bounded by composing the two
exponentially localized responses.
For example the squared distance
of an output bond from \(z\) is
bounded by a constant times the
squared distance to its intermediate
bond plus that intermediate bond's
squared distance from \(z\).
Sum the former with the exponential
response moment and the latter
with the rooted moment of
\(\mathcal K_{k,z}\). This is the
pointed root-transfer bound and
retains \(Cg\mathfrak b\).

If \(W\) is changed by a sufficiently
small output site gauge, actual
orbit covariance of \(U_{k+1}\)
shows that the two intermediate
fields are gauge equivalent, with
the prescribed output node values.
The proved invariance of
\(F_{k,z}^{\rm phys}\) therefore
makes \(\mathscr F_{k,z}\) invariant.
Its first jet is zero and the
same differentiated gauge identity
gives the output Ward identity.

We check the physical affine
normalization using the actual
successive map. By
\eqref{abc:successive-affine-response},
its first response to an affine
output field of constant curl is
the intermediate affine potential
modulo a lattice gradient. That
identity follows from the actual
nesting in \cref{haf:physical-nesting},
the full Wilson affine response,
and then its literal partial mean.
The literal and straightened means
are kept distinct; their frame
correction is an intermediate
gradient. The vanished first jet
removes the nonlinear response,
and the Ward identity removes both
gradient legs. Thus
\eqref{abc:successive-coefficient-preserved}
applies to this complete invariant
density. Its physical Wilson
normalization is unchanged by
the intermediate change of units.

The original plaquette functions,
the nearest-midpoint shares and
the midpoint tents are covariant
under translations and the full
hypercubic group, including
reflections. Orbit uniqueness
preserves this covariance in the
complete composed density. The
quadratic tensor classification
\eqref{ph:Maxwell-T}--\eqref{ph:extractor}
therefore applies to this invariant
kernel with its proved moments.
It is applied to the complete
kernel, using the Ward and symmetry
identities just established; no
termwise invariant-core hypothesis
is inferred from its Taylor filing.
The extractor is a fixed linear
second-moment functional, so the
bound in
\eqref{hpt:actual-physical-marginal-bound}
follows from the composed moment
bound. Absolute lifted-carrier
and cluster convergence gives
the infinite-lattice companion
as the limit of compatible tori,
with these derivatives and moments.
\end{proof}

The quadratic identity belongs to
the complete invariant density.
It is not applied separately to
the fixed-flat Taylor terms. A
termwise invariant core family,
an older-family transport theorem
or a full fine-current marginal
identification still requires its
corresponding additional argument.

% END LAYERED EXPANSION formation/homogeneous_physical_pointing.tex

% BEGIN LAYERED EXPANSION formation_even_coupling.tex
\subsection{An even-coupling improvement of the intrinsic formation bound}
\label{fec:section}

The first absolute formation estimate uses the size
of the cubic fluctuation vertex. The normalized
integral has an additional cancellation: changing
the sign of the fluctuation coupling can be
compensated by changing the sign of every
Gaussian variable. We use that exact symmetry
before estimating the resulting local activities.
It improves their intrinsic smallness by one
power of the coupling.

Keep the homogeneous law and its actual
Gaussian structure from
\cref{prop:haf-actual-homogeneous-formation}.
In this subsection the cutoff radius
\(\mathfrak b\) is fixed independently of a
complex parameter \(g\). The parameters
\(\tau\), the external background and any
admitted ordinary coarse input are also
held fixed. Only at the end do we set
\(g=\theta^{-1/2}\) and
\(\mathfrak b=A_1(\log\theta)^p\).
The explicit numerical Jacobian and
Gaussian normalizer have already been
factored as in
\eqref{afs:exact-physical-step};
their logarithms are not included in the
intrinsic function considered here.

\subsubsection{The analytic deformation at zero coupling}

Write the exact normalized fluctuation
integral in the form
\begin{equation}\label{fec:intrinsic-integral}
 \mathcal Z(g;V,H,\mathfrak b)
      =\int \chi_{\mathfrak b}(B)
                 e^{P(g;V,H,B)}\,d\mu_{Q_\tau(V)^{-1}}(B),
 \qquad \mathcal F(g)=\log\mathcal Z(g).
\end{equation}
The Haar and pivot ratio in \(P\) is the
actual one. The action and relative gauge
sectors are their exact centred functions
minus their quadratic Gaussian parts.
The ordinary input is
\(H(V(gB))-H(V)\). The free coordinate
graph \(V(u)\) and all precision
localizations are independent of \(g\).

Every divided remainder in this formula
is analytic at \(g=0\). For example a
centred action remainder of order three
has the convergent expansion
\[
 g^{-2}R_3(gB)
       =\sum_{n\ge3}\frac{g^{n-2}}{n!}
                                      D^nR_3(0)[B,\ldots,B].
\]
The ordinary density and source terms
are analytic functions of \(gB\), with
their constant values subtracted.
Thus \(P(0;V,H,B)=0\).
The cutoff is a fixed real even function
of each \(B(b)\).

The absolute formation proof remains
valid on the complex disc
\begin{equation}\label{fec:coupling-disc}
                         |g|<r_g=c_*/\mathfrak b.
\end{equation}
Here \(c_*>0\) is fixed below all
intrinsic-vertex and analytic-chart
ceilings. Its proof estimates the
vertices by \(C|g|\mathfrak b\);
positivity of \(g\) was not used.
All Gaussian covariances are held
fixed on their already admitted
complex background domain.
The cluster series is therefore a
holomorphic, zero-free logarithm on
this disc, with the same fixed
spatial rate reserve. This gives
holomorphy in the norm of the local
activity family, rather than only
pointwise finite-volume holomorphy.

At zero coupling,
\[
             \mathcal F(0)
                   =\log\mu_{Q_\tau(V)^{-1}}(\chi_{\mathfrak b}).
\]
Its actual wall expansion is exponentially
small in \(\mathfrak b^2\). With a fixed
slightly smaller spatial rate, its
unpointed and canonical pointed norms
are bounded by
\begin{equation}\label{fec:zero-coupling-tail}
                           C e^{-c\mathfrak b^2}.
\end{equation}
Polynomial wall factors have been
absorbed by decreasing \(c\).
This is the wall estimate of the
formation proof, and does not bound
a whole-volume cutoff probability
away from zero.

\subsubsection{Exact sign symmetry of the filed functions}

The exact integrand satisfies
\begin{equation}\label{fec:integrand-parity}
              P(-g;V,H,-B)=P(g;V,H,B).
\end{equation}
Indeed every occurrence of the actual
group field has argument \(gB\),
which is unchanged. A prefactor
\(g^{-2}\) is also unchanged.
The subtracted quadratic Gaussian
form is even in \(B\). The ordinary
input depends on the same group
field, and the Haar and pivot
functions have that same argument.
The even Gaussian density and cutoff
therefore give
\begin{equation}\label{fec:total-parity}
                         \mathcal F(-g)=\mathcal F(g).
\end{equation}
The equality uses the zero-free branch
through \(g=0\), so there is no
undetermined logarithm constant.

We need this identity for the local
family, where an arbitrary refiling
could obscure cancellation. Use the
actual Taylor and carrier filing of
the homogeneous formation theorem.
That filing commutes with the
involution \((g,B)\mapsto(-g,-B)\):
each Taylor monomial keeps its
degree in \(g\) and its degree in
the fluctuation variables, and
its carrier is independent of their
signs. An intrinsic degree-\(n\)
monomial has the factor
\(g^{n-2}B^n\), whose sign is
\((-1)^{2n-2}=1\).
An ordinary degree-\(n\) term has
\(g^nB^n\), also with sign one.

Mayer exponentiation, multiplication,
the fixed Gaussian weakening,
connected extraction and union filing
all commute with this involution.
Any partial Gaussian integral has
the same symmetry under negation of
all its Gaussian coordinates.
Consequently every resulting filed
activity can be chosen even in \(g\).
This statement concerns sign parity;
it does not assert gauge invariance
of individual fixed-reference terms.

The canonical bond pointings have
the same parity. For an intrinsic
plaquette function \(P_g(B)\), its
normalized insertion is
\[
 t^{-2}B(b)\cdot D_bP_g(tB)
                      -2\vartheta_{p,b}t^{-3}P_g(tB).
\]
Equation \eqref{fec:integrand-parity}
changes both the vector and its
derivative sign in the first term;
the second term is unchanged.
Ordinary radial insertions satisfy
the same identity. A moving cutoff
threshold and its surface measure
are even in \(B(b)\). Thus the
actual pointed expectation, its
integral in the interpolation
parameter, and its carrier filing
are even as well.

\subsubsection{The improved absolute bound}

\begin{proposition}[Quadratic intrinsic formation smallness]
\label{fec:quadratic-formation}
Let \(\mathbf F(g)\) denote the actual
unpointed or canonical pointed homogeneous
intrinsic activity family, with the
spatial rate just fixed. For
\(|g|\le c_*/(2\mathfrak b)\),
\begin{equation}\label{fec:quadratic-formation-bound}
 \|\mathbf F(g)\|
       \le C\{\,|g|^2\mathfrak b^2
                                  +e^{-c\mathfrak b^2}\,\}.
\end{equation}
The constants are independent of the
inner fine depth and volume.
The same statement holds on the
smaller fixed external and ordinary
source discs of the formation theorem.
\end{proposition}

\begin{proof}
The activity-valued function
\(\mathbf F(g)-\mathbf F(0)\)
is holomorphic and even on
\eqref{fec:coupling-disc}. On a
fixed slightly smaller concentric
disc its norm is bounded by a
constant, since
the original estimate is
\(C|g|\mathfrak b\) plus the wall
tail. Use the Cauchy circle of
radius \(R=3r_g/4\). For the
stated \(|g|\le r_g/2\),
\(|g|/R\le2/3\).

Its Taylor coefficients of odd
order vanish, and its constant
coefficient is zero. Cauchy's
estimate in the activity norm gives
\[
 \|\mathbf F(g)-\mathbf F(0)\|
       \le C\sum_{n\ge1}(|g|/R)^{2n}
                              \le C|g|^2\mathfrak b^2.
\]
Add \eqref{fec:zero-coupling-tail}.
The same proof applies to each
pointed norm because the preceding
parity argument held before those
norms were estimated.
\end{proof}

For the physical choice
\(g=\theta^{-1/2}\),
\(\mathfrak b=A_1(\log\theta)^p\),
\(p>1/2\), this reads
\begin{equation}\label{fec:physical-smallness}
 \|\mathbf F\|
   \le C\theta^{-1}(\log\theta)^{2p}
                       +C\exp[-c(\log\theta)^{2p}].
\end{equation}
The tail is smaller than any fixed
inverse power of \(\theta\) after
increasing its threshold, because
\(2p>1\). This is a stronger bound
than estimating the cubic vertex
alone. It is still a bound on the
specified homogeneous model.

Two external field derivatives
cost a fixed inverse radius squared.
The reserved spatial rate absorbs
their coordinate incidence and
second spatial moments, exactly
as in
\cref{prop:hpt-physical-quadratic-Ward}.
For zero ordinary profile, or an
input already proved to satisfy the
same complete physical gauge identities
and the Euclidean symmetries required
for a scalar coefficient in
\cref{cor:hpt-physical-marginal-bound},
that complete physical quadratic
coefficient consequently satisfies
\begin{equation}\label{fec:physical-marginal}
 |b_k^{\rm phys}|
        \le C\{\theta^{-1}(\log\theta)^{2p}
                           +e^{-c(\log\theta)^{2p}}\}.
\end{equation}
This uses the complete invariant
pointed density and its actual
critical substitution. It is not a
coefficient readout from separate
noninvariant Taylor pieces.

\subsubsection{The first nonzero fluctuation coefficient}

The same argument identifies the
algebraic coefficient without a
numerical two-loop claim.
Write the actual exponent as
\[
                 P(g;B)=gP_1(B)+g^2P_2(B)+O(g^3).
\]
Then \(P_1\) is odd and \(P_2\)
is even. Let
\(\langle\cdot\rangle_{\chi}\)
be the normalized Gaussian
expectation with its retained
even cutoff. Expansion of the
normalized integral gives
\begin{equation}\label{fec:first-coefficient}
 \mathcal F(g)=\mathcal F(0)
   +g^2\left\{\langle P_2\rangle_{\chi}
                     +\tfrac12\langle P_1;P_1\rangle_{\chi}\right\}
                                                +O(g^4).
\end{equation}
Here the semicolon is the connected
two-variable expectation. Its
subtracted product is zero because
\(\langle P_1\rangle_\chi=0\).
The \(O(g^4)\) is in the smaller
activity norm, bounded by
\(C|g|^4\mathfrak b^4\).

For precision, let \(\mathscr S_V(u)\)
be the complete centred composed
action on the actual pivot graph,
and let \(\mathscr G_V(u)\) be
the actual relative contour gauge.
Their quadratic parts are the two
parts of \(Q_\tau(V)\).
Let \(\mathscr H_V(u)=H(V(u))\)
and \(\mathscr J_V(u)=\log J_V(u)\).
All derivatives below are at zero.
The two polynomials in
\eqref{fec:first-coefficient} are
\begin{equation}\label{fec:literal-first-two-vertices}
 \begin{split}
 P_1(B)={}&-\frac{\tau}{6}D^3\mathscr S_V[B^3]
            -\frac16D^3\mathscr G_V[B^3]
            +D\mathscr J_V[B]+D\mathscr H_V[B],\\
 P_2(B)={}&-\frac{\tau}{24}D^4\mathscr S_V[B^4]
            -\frac1{24}D^4\mathscr G_V[B^4]
            +\frac12D^2\mathscr J_V[B^2]
            +\frac12D^2\mathscr H_V[B^2].
 \end{split}
\end{equation}
Repeated arguments are denoted by
\(B^n\). In particular the action
vertices contain derivatives of
the nonlinear pivot constraint;
the relative gauge vertices are
those of the literal trace gauge.
Neither can be replaced by an
ambient unconstrained action or
a squared logarithmic gauge.

Without the cutoff,
\eqref{fec:first-coefficient} is
evaluated by Gaussian pairings.
Each contribution of
\(\langle P_1;P_1\rangle\)
has a line joining the two
vertices: their separate odd
expectations vanish. Thus its
connected carrier is rooted
without a second volume factor.
The cutoff correction has the
already proved wall smallness.
Equation \eqref{fec:first-coefficient}
therefore specifies the next
fluctuation coefficient as an
actual connected expression, but
does not assign it a universal
numerical beta-function value.

% END LAYERED EXPANSION formation_even_coupling.tex

% BEGIN LAYERED EXPANSION formation/homogeneous_gaussian_comparison.tex
\subsection{Gaussian comparison for the classical-pullback homogeneous law}
\label{layered:homogeneous-Gaussian-comparison}

We compare the actual regular homogeneous
one-step measure with its quadratic Gaussian
reference at the trivial flat background.
The observables are fine Wilson plaquettes
evaluated on the classical background
\(U_k(V)\). Their quadratic parts are
resummed into the Gaussian precision.
The source derivatives are then rooted at
their original fine plaquettes. This gives
an error proportional to source support
mass, rather than the whole lattice volume.

This measure is a classical-pullback
one-step model. It does not include an
arbitrary older-core Gibbs weight, history
conditioning or independent fine quantum
fluctuations. The Gaussian observables
below are quadratic functions of a
Gaussian field; their joint law need
not itself be a multivariate Gaussian.

\subsubsection{The actual normalized fibre measure}

Fix the compatible periodic geometry of
\cref{prop:haf-actual-homogeneous-formation},
inner depth \(k\), and \(h=L^{-k}\).
Let \(\tau\) be real in the proved
fixed interval about one, and put
\(g=\theta^{-1/2}\),
\(\mathfrak b=A_1(\log\theta)^p\),
\(p>1/2\). Work on the regular coarea
branch of the actual outer fibre
\(M(V)=1\). The background is exactly
one. Its actual graph is
\begin{equation}
 V(gB)=e^{i\Phi_1(gB)},\qquad
                 \Phi_1(gB)|_{\rm free}=gB,
 \qquad \chi(B)=\prod_b\mathbf1_{|B(b)|<\mathfrak b}.
 \label{hgc:actual-flat-fibre}
\end{equation}
The entire regular fibre on this
free box is this graph, by the
proved pivot anchor and uniqueness.
The contour-gauge indicators equal
one there. Their reference weights,
the gauge normalizers and the real
box remain fixed as sources vary.

Let \(J_1(u)\) be the literal Haar
and raw pivot coarea density in
\eqref{afs:actual-coarea-density}.
The normalized law \(\nu_{k,g}\)
is the pullback to \(B\) of
\begin{equation}
 \frac1{Z_{k,g}}
 \mathbf1_{\rm retained\ tests}(V)
 e^{-\tau S_k(V)/g^2-\mathrm{GF}(V)/g^2}
                          \delta_G(M(V))\,dV.
 \label{hgc:actual-group-law}
\end{equation}
The delta denotes the stated regular
Haar coarea disintegration. All
field-independent gauge, Haar and
\(g^d\) factors cancel in this
normalized measure; no determinant
depending on a source is discarded.

Write
\(K_0=Q_{k,\tau}=\tau K_{A,k}(0)+K_G(1)\),
\(C_0=K_0^{-1}\), and
\(\mu_0=\mu_{C_0}\). For this
subsection only, write \(\mathsf C\)
for the free-to-full pivot tangent
\(D\Phi_1(0)\), to distinguish it
from the covariance \(C_0\).
Define the exact remainder
\begin{align}
 P_g(B)={}&-g^{-2}
 \left\{\tau S_k(V(gB))+\mathrm{GF}(V(gB))
                      -\tfrac12g^2 B^{\mathsf T}K_0B\right\}
                  +\log\frac{J_1(gB)}{J_1(0)}.
 \label{hgc:actual-remainder}
\end{align}
There is no linear term: the flat
action, contour gauge and their
first derivatives vanish. Consequently
the literal normalized identity is
\begin{equation}
 d\nu_{k,g}(B)=
    \frac{\chi(B)e^{P_g(B)}}{\mathcal Z_g(0)}\,d\mu_0(B),
 \qquad \mathcal Z_g(0)=\mu_0[\chi e^{P_g}]>0.
 \label{hgc:exact-reference-law}
\end{equation}
The preceding formation proof supplies
the actual local vertices of \(P_g\),
with norm \(Cg\mathfrak b\), and
the complete admitted Gaussian structure.
All constants are uniform in \(k,h\)
and compatible volume.

\subsubsection{Actual quadratic source operators}

For a fine plaquette \(p\), put
\(f_{k,p}(V)=a_p(U_k(V))\).
It has value and first derivative
zero at one. Let \(\mathcal H_k\)
be the actual linear fine background
response in physical potential units.
Its literal gauge correction is a
gradient and does not change its
linear curvature. Define
\[
 \mathcal D_{k,p}B
       =(\mathrm{d}_h\mathcal H_k\mathsf C B)(p).
\]
The fine Wilson normalization then gives
\begin{equation}
 \frac12D^2[f_{k,p}(V(u))]_{u=0}[B,B]
       =B^{\mathsf T}A_pB,
 \qquad A_p=\frac{h^4}{2}\mathcal D_{k,p}^{\mathsf T}
                                      \mathcal D_{k,p}.
 \label{hgc:actual-plaquette-quadratic}
\end{equation}
The nonlinear outer chart contributes
no extra quadratic term, because
\(Df_{k,p}(1)=0\).
The second derivative of the inner
background is likewise multiplied
by the vanished first fine
plaquette variation. The remaining
jet is exactly the displayed
linear-curvature Gram matrix.

The actual weighted strong response
\eqref{hb:localized-background-response}
gives exponentially summable coarse
bond kernels for \(\mathcal D_{k,p}\).
In particular, for a fixed \(s>0\),
\begin{equation}
 \sum_{b,b'}e^{s(d(b,p)+d(b',p))}
                        \|A_p(b,b')\|\le Ch^4.
 \label{hgc:actual-quadratic-source-mass}
\end{equation}
Keep the original \(p\) and its two
propagator paths as its carrier tag.
This is an actual source-rooted
operator family: each carrier contains
\(p\) and its matrix endpoints.
The same bound holds with a fixed
connected-carrier exponential weight,
by the marked-tree construction.
There is no fine-coordinate dimension
factor in \eqref{hgc:actual-quadratic-source-mass}.

For profiles \(w_i\), \(1\le i\le r\le4\),
with \(\|w_i\|_\infty\le1\), put
\begin{equation}
 \begin{aligned}
 H_i(V)&=\sum_pw_i(p)f_{k,p}(V),&
 X_i(B)&=g^{-2}H_i(V(gB)),\\
 A_i&=\sum_pw_i(p)A_p,&
 Q_i(B)&=B^{\mathsf T}A_iB.
 \end{aligned}
 \label{hgc:actual-observables}
\end{equation}
Equivalently
\(Q_i=\tfrac12H_i''(1)[\mathsf CB,\mathsf CB]\).
Summing \eqref{hgc:actual-quadratic-source-mass}
over fine plaquettes near a fixed
coarse bond cancels the physical
\(h^4\) with their \(O(h^{-4})\)
density. It proves a uniform weighted
Schur bound \(\|A_i\|_{\rm Sch,s}\le C\).
Summing its absolute entries instead
gives \(\|A_i\|_{\rm tr}\le
Ch^4\sum_p|w_i(p)|\).

The exact source remainder is
\begin{equation}
 R_{p,g}(B)=g^{-2}f_{k,p}(V(gB))-B^{\mathsf T}A_pB,
 \qquad X_i=Q_i+R_{i,g},\quad
 R_{i,g}=\sum_pw_i(p)R_{p,g}.
 \label{hgc:actual-source-split}
\end{equation}
Retaining its original fine tag in
the action-tree proof gives a
rooted potential norm at most
\(Ch^4g\mathfrak b\) for each
\(R_{p,g}\), with a fixed reserve
in the carrier rate. Its zero of
order three follows from the own
value, first jet and quadratic
subtraction. Its chart is the
same fixed admitted chart as
\eqref{hgc:actual-flat-fibre}.
Thus the large factor \(g^{-2}\)
in the observable has been treated
by its actual Taylor decomposition;
it is not inserted into a bounded-
source theorem as an old coefficient.

\subsubsection{Exact quadratic resummation}

For \(t=(t_1,\ldots,t_r)\), define
\begin{equation}
 K_t=K_0-2\sum_i t_iA_i,
 \qquad C_t=K_t^{-1},\qquad \mu_t=\mu_{C_t},
 \qquad R_g(t)=\sum_i t_iR_{i,g}.
 \label{hgc:resummed-precision}
\end{equation}
The uniform Schur bounds and the
actual gap of \(K_0\) give a fixed
complex polydisc \(|t_i|<4\rho\)
on which the real part of \(K_t\)
is positive and its actual carrier
norms are uniform. Decrease \(\rho\)
further, if necessary, in the
near/tail and combined-power
construction of
\eqref{haf:actual-decoupled-root}.
It supplies precisely the local
Gaussian structures for every
\(K_t\), including enlarged independent
cube discs, normalization and walls.
The radius \(\rho>0\) is independent
of \(h,k,g\), volume and profiles
satisfying the displayed supremum
bound. At most four profiles are
used, so no shrinking radius in
their number is hidden here.

Let
\begin{equation}
 \mathcal Z_g(t)=\mu_t[\chi e^{P_g+R_g(t)}],
 \qquad E_g(t)=\log\mathcal Z_g(t),\qquad
 L_G(t)=-\frac12\log\det\left(I-2C_0\sum_it_iA_i\right).
 \label{hgc:normalized-remainder-log}
\end{equation}
Choose the determinant branch zero
at \(t=0\), and the cluster-log
branch for \(E_g\). The actual
nonlinear source vertices remain
of norm \(Cg\mathfrak b\) on this
disc. Formation therefore makes
\(\mathcal Z_g(t)\) zero-free.
Gaussian completion, with all its
normalizers retained, proves
\begin{equation}
 \log\mathbb E_{\nu_{k,g}}e^{\sum_it_iX_i}
     =L_G(t)+E_g(t)-E_g(0),\qquad
 \log\mathbb E_{\mu_0}e^{\sum_it_iQ_i}=L_G(t).
 \label{hgc:exact-log-ratio}
\end{equation}
The determinant here is essential.
It is the source-dependent ratio
of the normalized Gaussians before
the residual interaction is expanded.

\subsubsection{Source-rooted differentiation of the actual expansion}

We prove a source-mass bound, not
an unpointed error times the volume.
Introduce independent fine coefficients
\(\zeta_p\), with
\(K_\zeta=K_0-2\sum_p\zeta_pA_p\)
and source remainder
\(\sum_p\zeta_pR_{p,g}\).
The preceding argument is uniform
on a fixed product ball
\(\|\zeta\|_\infty<\zeta_*\):
the uniform Schur bound uses the
\(h^4\) density sum already proved.
At each filed carrier combine
the fine-plaquette contributions
into one source-potential type.
Keep \(p\) as the mark when
differentiating. The type count
is then uniform in the fine mesh.
Choose the retained \(\rho\) so
that \(16\rho<\zeta_*\); the
four-profile polydisc is then
inside this product ball.
For each original fine plaquette,
the complete covariance has the
exact derivative
\begin{equation}
              \partial_{\zeta_p}C_\zeta
                            =2C_\zeta A_pC_\zeta.
 \label{hgc:exact-source-resolvent-derivative}
\end{equation}
Its source tag and its \(h^4\)
mass are retained in every term.
At complete parameters the exact
normalized differentiation identity
is
\[
 \partial_{\zeta_p}E_g
   =\left\langle R_{p,g}+B^{\mathsf T}A_pB\right\rangle_{g,\zeta}
                              -\operatorname{Tr}(C_\zeta A_p),
\]
where the brackets have density
\(\mathcal Z_g(\zeta)^{-1}\chi
e^{P_g+\sum_q\zeta_qR_{q,g}}\,d\mu_\zeta\).
Thus the Gaussian quadratic
normalizer is subtracted exactly.

This tag can also be retained
through the local covariance
construction. Truncate by the
full connected carrier, including
the original \(p\), both matrix
endpoints and their propagator
paths. A term with nearby matrix
endpoints and a distant source
root belongs to the long-carrier
tail. Endpoint distance alone
does not define this truncation.
Symmetrize each real carrier
contribution under transpose,
retaining its \(p\) tag. The
source-rooted exponential entry
bound above and the physical
\(h^4\) density sum give a
uniformly small Schur tail.
Thus the short-carrier real
precision differs from the full
gapped precision by that small
tail; the near-power spectral
argument still applies. More
explicitly, reserve two copies
of a small length weight in
the tagged source family. Since
its carrier joins \(b\) to
\(p\), the tail then obeys
\[
 \sum_{b',\,\gamma:\ell(X_\gamma)>R}
       \|A_{p,\gamma}(b,b')\|
        \le Ch^4e^{-cR}e^{-c d(b,p)}.
\]
Summing \(p\) uses
\(h^4\sum_pe^{-c d(b,p)}\le C\),
so the full source tail has
Schur norm \(Ce^{-cR}\),
uniformly on the source product
ball. Transposition gives the
column estimate. In this
full-carrier algebra,
differentiating a near
power inserts exactly one
\(2A_p/a\); differentiating its
tail telescoping formula does
the same. For example
\[
 \partial_{\zeta_p}H_{\le R}^m
   =\sum_{j<m}H_{\le R}^j
               (2A_{p,\le R}/a)H_{\le R}^{m-1-j}.
\]
Its full combined ball carrier
contains the original \(p\).
The global absolute entry sum of
that marked matrix is \(Ch^4\),
with constants allowed to depend
on the fixed \(m,R,M\).
For products, this follows from
\(\|UE_pV\|_{\rm entry,1}
\le\|U\|_{\rm Sch}\|E_p\|_{\rm entry,1}
\|V\|_{\rm Sch}\).
The marked geometric series costs
\(\sum_{v\ge1}v q^{v-1}<\infty\).
The prefix derivatives have the
same estimate. Consequently the
filed local covariance \(D_{W,Z}\),
with root coordinate set \(W\),
held root guard, and all cube
parameters zero outside its
wall domain \(Z\), has a
source derivative \(\dot D_{W,Z,p}\)
with
\begin{equation}
 \|\dot D_{W,Z,p}\|_{\rm entry,1}\le Ch^4,
 \qquad \|D_{W,Z}^{-1}\|\le C,
 \label{hgc:actual-local-source-matrix-bound}
\end{equation}
uniformly on the full coarse
parameter discs. If the local
wall domain does not contain
the source root \(p\), this
derivative is zero: every
retained carrier contains that
root. The same statements hold
for the explicit source-potential
derivative \(R_{p,g}\).

At complete parameters the actual
covariance restriction instead
has derivative
\[
 \partial_{\zeta_p}(\chi_WC_\zeta\chi_W^{\mathsf T})
      =2\chi_WC_\zeta A_pC_\zeta\chi_W^{\mathsf T},
\]
which can be nonzero for a
source outside \(W\). The
vanishing above concerns the
filed \(D_{W,Z}\) for
\(p\notin Z\); it does not
assert vanishing of this
complete restriction for
\(p\notin W\).

Here the entry norm is the sum
of colour-block norms. It bounds
trace norm up to the fixed colour
factor. Thus the derivative of
the normalized marginal Gaussian
density is its density times
\begin{equation}
 \frac12\{B^{\mathsf T}D_{W,Z}^{-1}\dot D_{W,Z,p}D_{W,Z}^{-1}B
                      -\operatorname{Tr}(D_{W,Z}^{-1}\dot D_{W,Z,p})\},
 \label{hgc:normalized-local-Gaussian-score}
\end{equation}
whose absolute value is at most
\(Ch^4(1+|B|^2)\).
There is no factor equal to the
local or exterior Gaussian
dimension in this bound. Its
expectation against the normalized
constant functional is exactly
zero. In particular an empty
selection cannot contribute a
source-dependent vacuum activity.

Use a fixed small extra Gaussian
tilt \(\delta\asymp(LM)^{-4}\)
to bound \(1+|B|^2\le C_\delta
e^{\delta|B|^2}\). The Gaussian
comparison charge is then only
a fixed constant per polymer
cube, absorbed by the reserved
spatial rate and cube-Cauchy
radius. It retains the same
small amplitude as the selected
integrand: every nonempty selection
has either an actual remainder
potential of size \(Cg\mathfrak b\)
or a failed-cutoff tail of size
\(Ce^{-c\mathfrak b^2}\).
Polynomial boundary factors can
be absorbed by a fixed decrease
of \(c\). Differentiating a
potential instead inserts the
actual \(R_{p,g}\), whose norm
is already \(Ch^4g\mathfrak b\).
Apply the proved inserted-potential
estimate to that literal norm.

It follows, by the same finite
Mayer and coarse-interpolation
formulas, that the actual activity
derivative with its source root
satisfies
\begin{equation}
 |\partial_{\zeta_p}H_g(Z;\zeta)|
 \le Ch^4\eta(g)e^{-\lambda d_c(Z)}
                    \mathbf1_{\{p\in Z\}},
 \qquad \eta(g)=g\mathfrak b+e^{-c\mathfrak b^2}.
 \label{hgc:actual-source-rooted-activity}
\end{equation}
Membership includes the fixed
assigned carrier collar. A fixed
boundary multiplicity is absorbed
in \(C\). The extra tilt used
above is chosen in the ordinary
Gaussian margin before the final
spatial rates; it does not depend
on \(g,h,k\) or source support.

Differentiate the convergent
connected-cluster logarithm.
Mark the differentiated activity
and retain its original \(p\).
The rooted cluster theorem bounds
all polymers attached to it.
Summing its own polymers containing
the fixed coarse cube of \(p\)
uses their reserved exponential
rate. Hence
\begin{equation}
 |\partial_{\zeta_p}E_g(\zeta)|
                            \le Ch^4\eta(g).
 \label{hgc:actual-source-rooted-log}
\end{equation}
This is the source-rooted argument;
the estimate is not obtained by
first bounding \(|E_g|\) by a
volume term and then differentiating.

For \(\zeta_p=\sum_it_iw_i(p)\),
\eqref{hgc:actual-source-rooted-log}
therefore proves
\begin{equation}
 |\partial_{t_i}E_g(t)|
         \le C\eta(g)h^4\sum_p|w_i(p)|.
 \label{hgc:actual-source-mass-bound}
\end{equation}
Integrating on the straight path
from zero to \(t\) also gives
\[
 |E_g(t)-E_g(0)|\le C\eta(g)
              \sum_i|t_i|h^4\sum_p|w_i(p)|.
\]

\subsubsection{Cumulants and their Gaussian leading terms}

\begin{proposition}[Uniform comparison for classical-pullback observables]
\label{prop:hgc-actual-cumulant-comparison}
For \(1\le n\le4\) and the profiles
of \eqref{hgc:actual-observables},
\begin{equation}
 \left|\kappa_{\nu_{k,g}}(X_{i_1},\ldots,X_{i_n})
              -\kappa_{\mu_0}(Q_{i_1},\ldots,Q_{i_n})\right|
           \le C_n\eta(g)\,m_h(w),
 \qquad m_h(w)=h^4\sum_{p,i}|w_i(p)|.
 \label{hgc:actual-cumulant-error}
\end{equation}
The constants and the analytic
source disc are independent of
homogeneous depth, mesh, volume
and source support. For the
unscaled actual observables
\(H_i(V)\), the corresponding
error is \(C_ng^{2n}\eta(g)m_h(w)\).
\end{proposition}

\begin{proof}
Use the exact normalized identity
\eqref{hgc:exact-log-ratio}.
Bound one source derivative of
its error by
\eqref{hgc:actual-source-mass-bound}.
The remaining at most three
derivatives follow from Cauchy
on the fixed larger source disc.
Repeated indices give only their
fixed factorials. This proves
\eqref{hgc:actual-cumulant-error}.
Finally \(H_i=g^2X_i\) exactly,
so multilinearity of cumulants
gives the factor \(g^{2n}\).
\end{proof}

\subsubsection{The source-rooted even-coupling refinement}

Keep the physical retained threshold
\(\mathfrak b\) fixed and replace
the coupling in the remainder by
an auxiliary complex \(\gamma\).
In particular
\(R_{p,\gamma}(B)=\gamma^{-2}
f_{k,p}(V(\gamma B))-B^{\mathsf T}A_pB\),
with its removable value at zero.
Write \(E_{\gamma,\mathfrak b}(\zeta)\)
for the resulting logarithm. The
Gaussian precision \(K_\zeta\),
the profiles and the real cutoff
are held fixed. The divided
remainders have removable values
at \(\gamma=0\). The actual chart
and local potential bounds hold
on \(|\gamma|<c_*/\mathfrak b\),
with \(c_*\) fixed below their
uniform ceilings, exactly as in
\eqref{fec:coupling-disc}.
This is an analytic deformation
of the retained-cutoff integral;
it makes no change to the
original real tests at the
physical value \(\gamma=g\).

The source-rooted proof above
therefore holds throughout that
disc, with
\[
 |\partial_{\zeta_p}E_{\gamma,\mathfrak b}|
       \le Ch^4\{|\gamma|\mathfrak b+e^{-c\mathfrak b^2}\}.
\]
Its \(h^4\) tag is unchanged.
At zero coupling only failed-cutoff
activities remain, so
\( |\partial_{\zeta_p}E_{0,\mathfrak b}|
\le Ch^4e^{-c\mathfrak b^2}\).

Moreover the actual source
remainders satisfy
\(R_{p,-\gamma}(-B)=R_{p,\gamma}(B)\).
The Gaussian, cutoff and the
matrix construction are independent
of \(\gamma\). The exact sign
argument of \eqref{fec:integrand-parity}
thus applies with every original
source tag retained. The chosen
cluster logarithm, and its
\(\zeta_p\) derivative, are even
in \(\gamma\). Subtract the
zero-coupling derivative. This
holomorphic function has a double
zero at the origin. Cauchy on
a fixed smaller concentric
circle, followed by the order-two
Schwarz bound, gives
\begin{equation}
 |\partial_{\zeta_p}E_{g,\mathfrak b}|
    \le Ch^4\eta_2(g),\qquad
 \eta_2(g)=(g\mathfrak b)^2+e^{-c\mathfrak b^2},
 \qquad g\mathfrak b<c_*/4.
 \label{hgc:even-source-rooted-bound}
\end{equation}
This step uses the uniform
source-rooted bound, so it
retains \(h^4\) rather than
introducing a whole-volume factor.
Profile summation and the same
fixed source-disc Cauchy argument
show that the log-MGF comparison
and \eqref{hgc:actual-cumulant-error}
hold with \(\eta_2\) in place
of \(\eta\). The unscaled
observables retain their exact
factor \(g^{2n}\). The preceding
\(\eta\) estimate is its weaker
corollary when \(g\mathfrak b<1\).

For clarity the leading Gaussian
law is entirely explicit:
\begin{align}
 \mathbb E_{\mu_0}Q_i&=\operatorname{Tr}(C_0A_i),\notag\\
 \operatorname{Cov}_{\mu_0}(Q_i,Q_j)
                   &=2\operatorname{Tr}(C_0A_iC_0A_j),\notag\\
 L_G(t)&=\sum_{n\ge1}\frac{2^{n-1}}n
                   \operatorname{Tr}\left[\left(C_0\sum_it_iA_i\right)^n\right].
 \label{hgc:explicit-Gaussian-leading-law}
\end{align}
Equivalently, the actual linear
fine-curvature covariance is
\(\Gamma_k(p,q)=\mathcal D_{k,p}C_0\mathcal D_{k,q}^{\mathsf T}\).
The first two observable formulas
are
\[
 \mathbb E Q_i=\frac{h^4}{2}\sum_pw_i(p)\operatorname{Tr}\Gamma_k(p,p),
 \quad
 \operatorname{Cov}(Q_i,Q_j)=\frac{h^8}{2}
       \sum_{p,q}w_i(p)w_j(q)
              \operatorname{Tr}[\Gamma_k(p,q)\Gamma_k(q,p)].
\]
Thus a mixed Gaussian quadratic
cumulant is
\(\frac{2^{n-1}}n\sum_{\sigma\in S_n}
\operatorname{Tr}(C_0A_{i_{\sigma(1)}}C_0A_{i_{\sigma(2)}}
\cdots C_0A_{i_{\sigma(n)}})\).
For coincident observables this
is \(2^{n-1}(n-1)!\operatorname{Tr}[(C_0A_i)^n]\).
These are quadratic Gaussian
cumulants, including their
generally nonzero higher orders.

Since \(p>1/2\), \(\eta_2(g)\to0\)
and the cutoff term is eventually
\(O((g\mathfrak b)^2)\). The log-MGF
and cumulant errors tend to zero
whenever \(\eta_2(g)m_h(w)\to0\),
including bounded source masses
and masses growing by a fixed
power of \(\log(1/g)\). Their
reference is the actual depth-
\(k\) quadratic Gaussian law.
All bounds retain the actual
constraint, Haar/coarea density
and fixed tests of this concrete
classical-pullback model.

% END LAYERED EXPANSION formation/homogeneous_gaussian_comparison.tex

% BEGIN LAYERED EXPANSION formation/regular_cutoff_symmetry.tex
\subsection{Symmetry of the cutoff fibre integral}
\label{layered:regular-cutoff-symmetry}

The cutoff is imposed in free bond coordinates.
Its lattice symmetry therefore requires more than
symmetry of the group average. For the middle
pivots used here, an isometry permutes the free
bonds and acts orthogonally on their colour
coordinates. We prove this and the covariance
of the resulting integral.

Let \(L=2m+1\), and use the centered blocks and
the pivot policy of \cref{ch:block-average-pivot-chart}.
For a positive coarse bond \(c=(y,y+Le_\mu)\),
its pivot is
\begin{equation}
 b_0(c)=(y+me_\mu,y+(m+1)e_\mu).
 \label{rcs:middle-pivot}
\end{equation}
Write \(\Pi\) for these bonds and \(E\) for their
complement among positive bonds, with all three
colour coordinates on each bond. A real fluctuation
cutoff in this subsection means
\begin{equation}
 \chi_{\mathfrak b}(B)=
        \prod_{b\in E}\mathbf1_{\{\langle B(b),B(b)\rangle
                                                    <\mathfrak b^2\}} .
 \label{rcs:real-product-cutoff}
\end{equation}
Thus each factor is a colour ball. Complex
polydiscs used for estimates are separate
analytic domains.

Let \(\mathsf r(x)=Rx+a\), where \(R\) is a
signed permutation matrix and \(a\) is a coarse
lattice translation. On oriented bonds define
\((\mathsf rU)(\mathsf rb)=U(b)\). We consider
isometries between compatible tori, or of the
infinite lattice. On a cubic torus this includes
the full hypercubic group.

\begin{lemma}[The action on free coordinates]
\label{lem:rcs-free-coordinate-action}
The isometry preserves the pivot set, with
positive orientations restored. For a background
\(U\), define the linear map \(\mathsf r_U\)
on bond coordinates as follows. Let \(b^{\mathsf r}\)
be the positive representative of the image
unoriented bond, and let
\(\varepsilon_{\mathsf r}(b)\) be its image
orientation sign. Put
\begin{equation}
 (\mathsf r_U a)(b^{\mathsf r})=
 \begin{cases}
 a(b),&\varepsilon_{\mathsf r}(b)=1,\\
 -\operatorname{Ad}(U(b)^{-1})a(b),
       &\varepsilon_{\mathsf r}(b)=-1.
 \end{cases}
 \label{rcs:twisted-coordinate-action}
\end{equation}
Then
\begin{equation}
        \mathsf r(e^{ia}U)
                     =e^{i\mathsf r_Ua}\mathsf rU.
 \label{rcs:exact-coordinate-identity}
\end{equation}
For real \(U\), \(\mathsf r_U\) permutes bonds
and acts by an orthogonal matrix on each
colour space. Its restrictions
\(\mathsf r_U^E,\mathsf r_U^\Pi\) are well defined,
and
\(\chi_{\mathfrak b}(\mathsf r_U^E B)
=\chi_{\mathfrak b}(B)\).
\end{lemma}

\begin{proof}
If \(Re_\mu=e_\nu\), the image of
\eqref{rcs:middle-pivot} is the middle
pivot of the image coarse bond.
If \(Re_\mu=-e_\nu\), the new positive
coarse bond starts at \(\mathsf ry-Le_\nu\).
Its middle pivot has endpoints
\(\mathsf ry-(m+1)e_\nu,\mathsf ry-me_\nu\),
the reversed endpoints of the image pivot.
Thus \(\Pi\), and hence \(E\), are preserved.

On a reversed bond,
\[
 (e^{ia(b)}U(b))^{-1}
  =U(b)^{-1}e^{-ia(b)}
  =e^{-i\operatorname{Ad}(U(b)^{-1})a(b)}U(b)^{-1}.
\]
This proves \eqref{rcs:exact-coordinate-identity}.
For real \(U\), its adjoint action preserves
the colour pairing; the minus sign also
preserves it. The product cutoff is
therefore unchanged. The same assertion
would not follow for a box in fixed
individual colour components.
\end{proof}

\begin{lemma}[Covariance of the mean and graph]
\label{lem:rcs-mean-and-graph}
The actual one-step mean satisfies
\(\mathcal M(\mathsf rU)=\mathsf r\mathcal M(U)\).
For its admitted graph,
\begin{equation}
 \Phi_{\mathsf rU}(\mathsf r_U^E u)
                        =\mathsf r_U\Phi_U(u).
 \label{rcs:joint-graph-covariance}
\end{equation}
The identity is holomorphic wherever both
graph arguments lie in the joint chart.
\end{lemma}

\begin{proof}
The isometry sends the coordinate-order path
multiset from \(y\) to \(x\) to the one from
\(\mathsf ry\) to \(\mathsf rx\), preserving
its weights. The internal implicit mean
therefore has the same covariance, by
uniqueness.

The outer operation is the exponential of
the average of closed logarithms, followed
by multiplication by the central path.
It also respects reversal. To check this
directly, write its central transport as
\(C\) and its closed inputs as \(Z_x\), so
\(\mathcal M(U)(c)=e^{\overline{\log Z}}C\).
For the reversed coarse bond the closed
inputs are \(C^{-1}Z_x^{-1}C\).
Their average logarithm is
\(-\operatorname{Ad}(C^{-1})\overline{\log Z}\);
multiplication by \(C^{-1}\) gives
\((e^{\overline{\log Z}}C)^{-1}\).
This uses the actual outer operation.

Put \(W=\mathcal M(U)\). The relative raw
constraint \(q_U\) transforms as
\[
           q_{\mathsf rU}(\mathsf r_Ua)
                              =\mathsf r_W q_U(a).
\]
For a reversed coarse bond this is
\(-\operatorname{Ad}(W(c)^{-1})q_U(a)(c)\);
it follows by conjugating the inverse
relative coarse product before taking
its admitted logarithm.
The right side of
\eqref{rcs:joint-graph-covariance} thus
solves the transformed constraint with
the transformed free values.
Its small-graph uniqueness proves the
identity for real data. Joint analyticity
extends it on the common complex chart.
\end{proof}

\begin{proposition}[Covariance of the regular transfer]
\label{prop:rcs-cutoff-transfer-covariance}
Use the actual regular coarea component,
unit contour-gauge weights, the middle
pivots \eqref{rcs:middle-pivot}, and the
cutoff \eqref{rcs:real-product-cutoff}.
Assume the action and input pointings
are covariant under \(\mathsf r\).
Then the transfer and its complete
pointings are covariant under
\(\mathsf r\). This applies to the
canonical homogeneous step with action
\(S_j\), and to an ordinary input with
the same covariance. In particular the
scalar quadratic readout is available
for that model on the infinite companion.
\end{proposition}

\begin{proof}
For the physical background action,
reflecting a critical orbit gives a
critical orbit with reflected datum.
The strong norm, regularity bounds and
averaging constraint are preserved.
Orbit uniqueness in
\cref{prop:hb-actual-critical-orbit}
identifies the two orbits. Each original
plaquette trace, and hence \(S_j\), has
the required covariance. Any difference
between chosen literal representatives
is a node-fixed gauge and is removed
at the level of the complete function.

For a fixed real background, use
\(u'=\mathsf r_U^E u\) as free coordinates.
This is an orthogonal change with
absolute determinant one. Equation
\eqref{rcs:joint-graph-covariance} maps
the graph and the product colour balls
exactly. Regular plaquette bounds and
the coarea fibre are preserved.

The coordinate coarea factor is preserved
as well. Its pivot derivative transforms
by \(\mathsf r_W\) on its rows and
\((\mathsf r_U^\Pi)^{-1}\) on its columns.
For each reversed coarse row the
corresponding pivot is also reversed:
the two factors of determinant minus
one cancel. Their adjoint factors have
determinant one. Row and pivot
permutations cancel in the same way.
The reciprocal pivot determinant on its
identity branch is therefore unchanged.
The Haar density \(j(a)\) depends only
on the invariant colour pairing.
Consequently the full coarea density
in \eqref{afs:actual-coarea-density}
is carried to itself.

The internal transports and their
reference products are permuted by
the isometry. Their relative contour
trace, tame test and unit-weight sum
are unchanged. The centered action
and its full quadratic Hessian transform
by the same free orthogonal map.
Thus the Gaussian precision transforms
as \(K'=\mathsf r_U^E K(\mathsf r_U^E)^{-1}\),
and its normalized real Gaussian
measure is carried to the new one.
This proves covariance of the integral,
including its actual cutoff, Haar
factor, pivot factor and numerical
Gaussian normalization.

The block axial choice also commutes
with the isometry. If
\(s_U(x)=T_{y,x}(U)\) and \(s_U(y)=1\),
then
\(s_{\mathsf rU}(\mathsf rx)=s_U(x)\).
Hence
\((\mathsf rU)^{s_{\mathsf rU}}
=\mathsf r(U^{s_U})\).
This uses the symmetric contour mean,
not a choice of one ordered axial path.

For the pointings, choose their
insertions before carrier filing.
Original plaquettes and old anchors
are permuted, and midpoint tents are
covariant. On a reversed free bond,
the vector and its derivative are
contracted by the same orthogonal
map, so the radial and normalized
cubic insertions are preserved.
The cutoff surface and its colour
surface measure are also preserved.
For a Gaussian trace, each diagonal
colour block is conjugated and its
free bond is permuted. The pivot
normalizer factors have the covariance
already proved. Normalized expectations
and their interpolation integrals
therefore give covariant complete
pointed functions.

The real identities extend to the
admitted complex external-field domain
by holomorphy. This extension does not
change real integration contours by
a complex orthogonal map.
The auxiliary localization grid need
not be preserved by every translation;
it files an already defined complete
pointing and does not change its sum.
Its individual pieces need not have
the full covariance.
\end{proof}

% END LAYERED EXPANSION formation/regular_cutoff_symmetry.tex

% BEGIN LAYERED EXPANSION formation/homogeneous_irrelevant_transport.tex
\subsection{Physical transport of homogeneous invariant densities}
\label{layered:homogeneous-irrelevant-transport}

The subtraction appropriate to a classical pullback is a multiple
of its actual effective action. We prove a local analytic
irrelevance estimate for the remaining complete physical density.
The proof uses the literal intermediate mean, including its node
frames. Its localized analytic representatives are not required
to be separately gauge invariant.

We work on the fixed coarse bond ball about the trivial flat
field, with compatible periodic volumes, and on its infinite
lattice companion. All assertions about the companion concern
normally convergent pointed analytic functions. They do not
require an infinite volume probability measure.
Fix \(0<\beta<\alpha<1\). Constants below may depend on
\(L,\alpha,\beta\), a fixed analytic input radius, and the
localization size. They are independent of the two depths
and the compatible volume.

\subsubsection{Literal intermediate coordinates}

Use final coarse spacing one. For \(0\le s\le k\) put
\[
 h=L^{-k},\qquad \xi=hL^s=L^{s-k},\qquad
 W=e^{iw},\qquad U_k(W)=e^{ih\widehat A_k(w)}.
\]
Here \(U_k\) is the literal critical representative of
\cref{prop:hb-actual-critical-orbit}. Define
\begin{equation}
 V_{s,k}(W)=M^sU_k(W),\qquad
 B_{\xi,b}(w)=\frac1i\log V_{s,k}(W)(b).
 \label{hit:actual-intermediate-map}
\end{equation}
The bond \(b\) on the left is indexed in old coarse units;
its physical midpoint in final units is \(\xi m_b\).
The logarithm will be shown to stay in its actual chart.
This definition does not replace \(M^s\) by the straightened
linear average.

We first isolate a bound for the actual mean acting on an
arbitrary small fine field. Write
\[
 \mathscr P_{h,s}(A)=\frac1i\log M^s(e^{ihA}).
\]
For a homogeneous polynomial row, its absolute row norm
is the sum over its fine input bonds of the norms of its
colour multilinear coefficients. Ordered coefficients
include their Taylor factorial once.

\begin{lemma}[Absolute rows of the full plain mean]
\label{lem:hit-plain-mean-rows}
There are \(c,C>0\) such that the degree \(m\) row
coefficient \(T_{h,s;b}^{(m)}\) of \(\mathscr P_{h,s}\)
satisfies
\begin{equation}
 \sum_{b_1,\ldots,b_m}
       \|T_{h,s;b;b_1\ldots b_m}^{(m)}\|
                         \le C^m\xi^m,\qquad m\ge1.
 \label{hit:plain-mean-absolute-Taylor}
\end{equation}
Its fine inputs lie within physical distance \(C_L\xi\)
of the output bond. On \(\|A\|_\infty<c\),
\begin{equation}
 \|D\mathscr P_{h,s}(A)\|_{\rm row}\le C\xi,\qquad
 \|D^2\mathscr P_{h,s}(A)\|_{\rm row}\le C\xi^2.
 \label{hit:plain-mean-two-derivatives}
\end{equation}
The same conclusions, with zero constant term, hold
for the logarithm of the node frame in the partial
factorization.
For the row assertion alone \(C,c\) can be chosen
uniformly over \(L\ge3\), as in
\eqref{hfm:uniform-L-plain-rows}; the nonlinear
fixed-leg estimate has the different dependence
proved in \cref{prop:hfm-fixed-leg-density}.
\end{lemma}

\begin{proof}
Use the exact recursion and factorization
\eqref{hb:actual-recursion}--\eqref{hb:actual-plain-factorization},
stopped at \(s\). The one-step stencil is finite for fixed
\(L\). Its actual analytic chart therefore gives absolute
Taylor row bounds \(K_L^m\) for the terms of degree \(m\ge2\).
The derivative of the straightened row at zero must first
be combined: its row norm is exactly \(L\), by
\eqref{hb:one-step-bounds}. Taking the absolute values of
its separate contour and node-frame contributions before
this cancellation would give an incorrect depth factor.

The positive formal series majorizing the raw angle rows
can consequently be chosen to satisfy
\begin{equation}
 b_0(z)=hz,\qquad
 b_{t+1}(z)=Lb_t(z)+\frac{K_L^2b_t(z)^2}{1-K_Lb_t(z)}.
 \label{hit:actual-row-majorant}
\end{equation}
Increasing \(K_L\) also majorizes the zero-constant
one-step node increment by \(K_Lb_t/(1-K_Lb_t)\).
Let \(\xi_t=hL^t\), and evaluate these positive series
at \(z=c_0/\xi\). A first-exit argument gives
\[
 b_t(z)\le2\xi_tz\quad(0\le t\le s),\qquad
 \sum_{t<s}\xi_t\le\frac{\xi}{L-1}.
\]
Indeed, below this bound the quotient \(b_{t+1}/(L b_t)\)
is at most \(1+C_L b_t\). The product of these quotients
is at most \(\exp(C_Lz\sum_{t<s}\xi_t)\), less than two
when \(c_0\) is fixed sufficiently small. All inputs
then stay strictly in the one-step analytic chart.

The product defining \(v_s\), and the product in reverse
order defining \(v_s^{-1}\), have positive row majorants
bounded by
\[
             \prod_{t<s}(1-K_Lb_t(z))^{-1}.
\]
They differ from one by at most \(C_Lc_0\).
Apply the convergent exponential and logarithm series
to the exact product
\[
 M^s(e^{ihA})(b)
       =v_s(b_-)\,e^{ia_s(b)}\,v_s(b_+)^{-1}.
\]
Its zero-constant absolute logarithm majorant at
\(z=c_0/\xi\) is at most \(C_Lc_0\). Positivity of
the majorant gives the coefficient bound
\(C_Lc_0^{1-m}\xi^m\), which implies
\eqref{hit:plain-mean-absolute-Taylor}. This is an
absolute coefficient sum over fine inputs; it is
stronger than a Cauchy bound for a function of an
unspecified number of fine coordinates.

The one-step support enlargement at stage \(t\)
has physical size at most \(C_L\xi_t\). Its total
size is at most \(C_L\sum_{t<s}\xi_t\le C_L\xi\).
Differentiating the absolutely convergent majorant
on a fixed smaller supremum ball proves
\eqref{hit:plain-mean-two-derivatives}. The proof
for \(\log v_s\) uses the same product majorant.
Restriction of a node row to a coarser node set
does not increase any of these row sums.
\end{proof}
The literal representative also has a uniform connected
expansion. This step is needed because the action trees
can use a mean-slice representative by gauge invariance,
whereas the intermediate field in
\eqref{hit:actual-intermediate-map} uses the literal one.

\begin{lemma}[Connected expansion of the literal field]
\label{lem:hit-literal-field-carriers}
On a fixed coarse ball, \(\widehat A_k(w)\) has uniform
strong-cell norm \(C\|w\|_\infty\) and an analytic
connected-tree family rooted at each output fine cell,
with a fixed positive physical length rate. After
choosing \(M\) large enough for a requested cube rate
and decreasing the field radius, its rooted \(M\)-cube
norm is at most \(Cr\) on \(\|w\|_\infty<r\).
The same family bound holds for the entire field
vector on one localization cube, with its maximum
strong-cell norm. For a fixed requested rate its
amplitude constant can be chosen independent of
\(M\) and \(L\), with an allowed radius that may
depend on \(L\).
\end{lemma}

\begin{proof}
The mean-slice field \(A_k(w)\) has the strong-cell
trees of \eqref{het:actual-cell-equation} and
\eqref{het:propagator-cell-sum}. The prescribed literal
field is \eqref{hb:literal-representative}. Its node
function \(v_h(A)\) reads only a fixed physical star:
the support sum in the preceding proof is bounded
at full spacing one. That proof also gives its
uniform analytic supremum bounds on a fine
supremum ball.

The actual extension \(E_n\) in
\eqref{hb:node-extension} is local and bounded from
coarse supremum norm to \(Y_\alpha\). Hence
\(\lambda(A)=E_n(i^{-1}\log v_h(A))\) is a local
analytic map into \(Y_\alpha\) with norm
\(C\|A\|_\infty\). A fixed physical star contains
only a fixed number of strong cell slots. Banach
Cauchy in those slots gives coefficients \(C^m\);
it does not introduce the number of fine sites.

Finally use the uniformly analytic finite-gauge map
\eqref{hb:finite-gauge-BCH}. Its apparent \(h^{-1}\)
is removed before expanding by writing
\(\lambda(x+h e_\mu)=\lambda(x)+hD_\mu\lambda(x)\).
The logarithm of the remaining product vanishes
at \(h=0\), so division by \(h\) is analytic and
bounded in the stated strong cell norms.
Substitution of the field trees into these local
analytic maps attaches fixed physical stars and
vertices of arity at least one. Their number is
bounded by a fixed multiple of the leaves after
combining the linear maps. Their coefficients and
the fixed per-leaf carrier costs are absorbed by
a smaller input radius. The physical exponential
length reserve in the field trees pays the stars.
This proves both the connected family and the
strong bound for the actual literal representative.

For the vector assertion perform the same actual
fixed-point expansion with macroblocks from the outset.
Let \(\mathcal E_Q\) be a smooth macroblock
partition, whose physical derivatives are bounded
uniformly for \(M\ge1\), supported in the
\(M/10\) enlargement of its block. The product
slots are supported strong fields on these
enlargements; their bounded overlap reconstructs
the global field. Let \(\mathcal I_{Q'}\)
restrict the fine and coarse forcings to a macroblock.
The target norm is the maximum strong unit-cell
norm on the support of the output slot.
Work first with the infinite-cover inverse.
The unit-cell inverse
estimate gives
\begin{equation}
 \|\mathcal E_Q\mathcal L_{0,\infty}^{-1}\mathcal I_{Q'}\|
       \le C e^{-s_0\operatorname{dist}(Q,Q')/2},
 \qquad
 \sup_Q\sum_{Q'}e^{\kappa d_M(Q,Q')}
 \|\mathcal E_Q\mathcal L_{0,\infty}^{-1}\mathcal I_{Q'}\|
       \le C_\kappa .
 \label{hit:macroblock-joint-inverse}
\end{equation}
For the first estimate take the supremum over
target cells after summing the exponential unit-cell
kernel over source cells in \(Q'\). There is no
factor \(M^4\). For the second, touching macroblocks
give only a fixed number of terms; other terms
have physical separation at least
\(cM(d_\infty(Q,Q')-C)_+\). Choose \(M\) larger
than \(C\kappa/s_0\).
On a torus periodize these cover blocks,
retaining the image \(Q'+nP\) as a separate
label before any carrier assignment. Give
each such block a fixed path on the cover
joining its two macroblocks. The sum over
source blocks and period images is exactly
the cover sum just bounded. Short endpoint
distance in the quotient does not remove
a long image path. Projection of this
normally convergent tree construction solves
the actual periodic fixed-point equation;
contraction uniqueness identifies its sum
with the finite-volume background.

The actual residual in \eqref{het:actual-cell-equation},
restricted to \(Q\), reads only a fixed number of
neighboring macroblocks when \(M\) exceeds its
physical carrier radius. Its target is a forcing
supremum norm. Its analytic bound is \(C_L\|x\|^2\),
independent of \(M\): pointwise locality and the
physical transpose bound are used before taking
the supremum. Banach Cauchy counts these neighboring
macroblocks, not their fine coordinates.
The actual ordered-tree majorant therefore gives
a macroblock field family of size \(2C_\kappa r\)
after reducing the radius below its nonlinear
ceiling. The free constant in
\eqref{hit:macroblock-joint-inverse} is uniform in
\(L\), since its direct rows depend on \(N=L^k\),
as explained after
\eqref{hfm:uniform-L-plain-rows}. Growth in \(C_L\)
is paid only by the allowed radius.
The actual node extension and finite-gauge map
are local analytic maps on the same macroblock
spaces. Their uniform path-row bounds give the
stated literal-field vector family. Its amplitude
constant is independent of \(M,L\).
\end{proof}

\begin{proposition}[Actual smooth intermediate fields]
\label{prop:hit-actual-intermediate-fields}
There is a fixed \(r_*>0\) such that
\eqref{hit:actual-intermediate-map} is admitted on
\(\|w\|_\infty<r_*\), at every pair of depths.
For \(\|w\|_\infty<r\le r_*\) there is a continuous
connection \(a_\xi(w)\), holomorphic in \(w\), such that
\begin{equation}
 B_{\xi,b}(w)=\xi\,a_\xi(w)(\xi m_b),\qquad
              \|a_\xi(w)\|_{C^{1,\alpha}}\le Cr.
 \label{hit:actual-intermediate-C1alpha}
\end{equation}
For any fixed requested cube rate, with \(M\) and
the radius chosen accordingly, each raw coordinate
also has an actual local analytic family
\begin{equation}
 B_{\xi,b}(w)=\sum_{X\ni q(\xi m_b)}B_{\xi,b,X}(w_X),
 \qquad
 \sum_X e^{\kappa_0d_M(X)}
                 \|B_{\xi,b,X}\|_r\le C\xi r .
 \label{hit:actual-intermediate-family}
\end{equation}
The norm uses the final \(M\)-cube lattice; \(q(x)\)
is a cube containing \(x\).
There is also a Banach-valued version of
\eqref{hit:actual-intermediate-family}: all old bond
coordinates in one old \(M\)-cube can be grouped
into a single vector with its supremum norm.
Its rooted family has the same bound \(C\xi r\).
\end{proposition}

\begin{proof}
Apply \cref{lem:hit-plain-mean-rows} to the literal
field of \cref{lem:hit-literal-field-carriers}.
The raw angle and all intermediate raw inputs
are bounded by \(C\xi r\) and \(Cr\xi_t\).
The node products remain in their actual charts.
A fixed reduction of \(r_*\) proves admission.

For the spatial assertion write \(A=\widehat A_k(w)\).
Translation by \(\xi e_\nu\) preserves every grid in
the partial hierarchy. The actual mean commutes with
this translation. Put \(\delta_\nu A=\tau_{\xi e_\nu}A-A\).
The first two derivatives of the full plain mean give
\begin{equation}
 D_\nu^\xi(B_\xi/\xi)
   =\xi^{-2}\int_0^1
     D\mathscr P_{h,s}(A+t\delta_\nu A)
                              [\delta_\nu A]\,dt .
 \label{hit:intermediate-first-difference}
\end{equation}
The fine physical strong norm gives
\(\|\delta_\nu A\|_\infty\le C\xi r\), so the
right side is bounded by \(Cr\).
At an old-grid displacement \(y\) of physical size
\(d\le1\),
\[
 \|\tau_y\delta_\nu A-\delta_\nu A\|_\infty
                  \le C\xi r d^\alpha .
\]
The points on the two interpolation segments satisfy
\(\|\tau_y(A+t\delta_\nu A)-(A+t\delta_\nu A)\|_\infty
\le Crd\): use their convex-combination expression.
Equation \eqref{hit:plain-mean-two-derivatives} therefore
bounds the difference of \eqref{hit:intermediate-first-difference}
by
\[
                 Crd^\alpha+C\xi r^2d\le Crd^\alpha.
\]
The zeroth supremum of \(B_\xi/\xi\) is \(Cr\).
The exact staggered interpolation
\cref{wmc:interpolation}, at mesh \(\xi\), now gives
\eqref{hit:actual-intermediate-C1alpha}. It is a
fixed linear interpolation and hence preserves
holomorphic parameter dependence.

To prove the family assertion, keep the absolute
degree \(m\) raw row before attaching external-field
carriers. Its fine inputs lie within \(C_L\xi\)
of the root. Join those \(m\) points to the root
by a physical tree of length at most \(C_Lm\xi\),
and substitute their actual literal-field trees.
The coefficient sum is bounded by \(C^m\xi^m\);
the field trees contribute \((Cr)^m\).
Charging the join costs only \(e^{C\kappa_0m}\)
after passage to \(M\)-cubes, and a physical
length cost bounded by \(e^{Cm\xi}\).
These are per-leaf costs, absorbed by a smaller
fixed radius. Sum \(\sum_{m\ge1}(C\xi r)^m\).
This proves \eqref{hit:actual-intermediate-family}.
There is no separate rounding cost at each of
the arbitrarily many intermediate averaging stages.

For the grouped assertion, the image of an old
\(M\)-cube has side \(\xi M\le M\). Its raw output
vector reads the fine field in a fixed enlargement
of that image. Apply Banach Cauchy to the actual
plain-mean map on the fixed number of neighboring
macroblock strong spaces. The map is analytic
for an input supremum \(c/\xi\), and its whole
output vector has supremum \(Cc\) there, by
\eqref{hfm:uniform-L-plain-rows}. Its multilinear
coefficients are consequently bounded
\(C^m\xi^m\), with only a fixed macroblock-slot
factor. Compose the whole-cube literal-field
families of the preceding lemma. This proves
the vector bound without counting the old bond
coordinates in that cube.
\end{proof}

\begin{lemma}[All-depth physical nesting]
\label{lem:hit-all-depth-nesting}
On a fixed smaller ball, in common physical units,
\begin{equation}
 U_s(V_{s,k}(W))\sim U_k(W),\qquad
                      S_s(V_{s,k}(W))=S_k(W).
 \label{hit:actual-all-depth-action-identity}
\end{equation}
The first equivalence is a fine gauge equal to one
at the old coarse nodes. The old-unit fine mesh is
\(L^{-s}\).
\end{lemma}

\begin{proof}
Every variation tangent to an \(M^s\)-fibre is tangent
to the \(M^k\)-fibre, since
\(M^k=M^{k-s}\circ M^s\). Thus \(U_k(W)\) is critical
on the former fibre. In old physical units its
literal potential is \(\xi\widehat A_k\), its first
derivative is \(\xi^2D\widehat A_k\), and its local
H\"older derivative has the corresponding additional
\(\xi^\alpha\). Its strong norm is at most \(C\xi r\).
The old datum is admitted by
\cref{prop:hit-actual-intermediate-fields}.
Orbit uniqueness in \cref{prop:hb-actual-critical-orbit},
including its competing-multiplier bound
\eqref{hb:competing-multiplier-bound}, identifies
the two orbits for real data. Analytic continuation
gives the equality of their plaquette traces and
actions for complex data. Wilson action in four
dimensions is unchanged by this change of physical
units. This proves \eqref{hit:actual-all-depth-action-identity}
directly; no product of one-step uniqueness constants
is used.
\end{proof}

\subsubsection{Complete pointed inputs and their quadratic coefficient}

An old pointed family is a normally convergent representation
\begin{equation}
 f_z(B)=\sum_{X\ni q(z)} f_{z,X}(B_X),\qquad
 \sup_z\sum_X e^{\kappa d_M(X)}
                         \|f_{z,X}\|_\rho\le\epsilon .
 \label{hit:old-pointed-input-norm}
\end{equation}
Here \(X\) is a connected set of old \(M\)-cubes,
\(d_M(X)=|X|-1\), and \(\|\cdot\|_\rho\) is the
supremum on the old bond product ball
\(\|B\|_\infty<\rho\). Assign each positive bond
to its starting-site cube. The complete \(f_z\)
is normalized by \(f_z(0)=0\), invariant under
all sufficiently small admissible site gauges,
and covariant under translations of the anchor.
The individual \(f_{z,X}\) need not have these
invariances.

On finite tori we use the canonical physical family
and its infinite-lattice companion, with the
coupling parameters held fixed during the comparison.
Write \(\Pi_P f^\infty_z(B)=f^\infty_z(B^{\rm per})\)
for evaluation at the periodic lift of \(B\).
This is a function of the finite bond variables;
it is not a quantum expectation on a larger torus.
The compatibility property needed below is
\begin{equation}
 f_z^P-\Pi_P f_z^\infty=\mathcal W_z^P,\qquad
 \|\mathcal W_z^P\|_{\rho,\kappa/2}
                  \le C\epsilon e^{-c\kappa P/M}.
 \label{hit:finite-period-compatibility}
\end{equation}
Here \(P\) is the shortest old physical period.
The notation on the second line means an actual
rooted analytic family with the stated norm.
Its terms retain a finite-period-tail label.
For the canonical families this property is
proved below from their lifted operator and
cluster carriers. Gauge invariance on one finite
torus alone would not imply it.

\begin{lemma}[Ward moments of the complete density]
\label{lem:hit-complete-Ward-moments}
The infinite companion has zero first jet.
Its Hessian \(K_z=D^2f_z^\infty(0)\) has, for
some \(m>0\), an absolute two-leg moment
\begin{equation}
 \sum_{b,c}e^{m(|m_b-z|+|m_c-z|)}
                    \|K_z(b,c)\|\le C\epsilon,
 \qquad K_z[\mathrm{d}\lambda,a]=K_z[a,\mathrm{d}\lambda]=0.
 \label{hit:complete-Ward-moments}
\end{equation}
The Ward identity extends to polynomial
\(\lambda,a\). Define
\begin{equation}
 a^F_\mu(x)=\frac12\sum_\nu F_{\nu\mu}(x_\nu-z_\nu),
 \qquad q_f(F)=\frac12K_z[a^F,a^F].
 \label{hit:curvature-quadratic-readout}
\end{equation}
If the complete pointing is covariant under the
full hypercubic group, including reflections,
then
\begin{equation}
 q_f(F)=\frac{b_f}{2}\sum_{\mu<\nu}
                         \langle F_{\mu\nu},F_{\mu\nu}\rangle,
 \qquad |b_f|\le C\epsilon .
 \label{hit:scalar-curvature-coefficient}
\end{equation}
The colour pairing is bilinear on the complexified
algebra.
\end{lemma}

\begin{proof}
The first derivative is a colour vector at each
bond. Invariance under constant colour rotations
forces each vector to be zero. Differentiate
the local gauge identity once in the gauge
parameter and once in a bond direction.
The term containing a derivative of the
gauge-coordinate transformation is multiplied
by this vanished first jet. The remaining
identity is the Ward identity for the Hessian.

A carrier contains at most \(C M^4(1+d_M(X))\)
old bond and colour coordinates. Two Cauchy
derivatives cost this coordinate-pair count
and a fixed inverse power of \(\rho\).
Every coordinate in a rooted carrier is
within \(CM(1+d_M(X))\) of its anchor.
Choose \(m\) small relative to \(\kappa/M\);
the exponential and polynomial costs are paid
by a reserved fraction of the rate in
\eqref{hit:old-pointed-input-norm}.
This proves the moment bound.
Cut off polynomial gauge parameters and
arguments and pass to their limit using this
absolute moment. No unbounded polynomial
gauge transformation is applied to the
nonlinear function.

Constant colour rotations make the quadratic
colour pairing scalar. Independent axis
reflections kill terms joining different
curvature planes, since the two planes have
different reflection characters. Axis
permutations make the remaining six coefficients
equal. The definition of \(a^F\) has no
diagonal term \(F_{\mu\mu}\); using a bond
midpoint instead of its initial site therefore
does not change it. The moment bound proves
the estimate for its scalar coefficient.
\end{proof}

\subsubsection{An actual classical marginal representative}

Let
\(\theta_z(X)=\prod_{\lambda=1}^4(1-|X_\lambda-z_\lambda|)_+\),
periodized on a torus. On the old lattice with
fine spacing \(h_s=L^{-s}\), put
\begin{equation}
 e_{s,z}(B)=\sum_{p>0}\theta_z(m_p)
                              a_p(U_s(e^{iB})).
 \label{hit:actual-classical-pointed-density}
\end{equation}
The point \(m_p\) is the original fine plaquette
midpoint in old physical units.

\begin{lemma}[Exact representative and normalization]
\label{lem:hit-classical-marginal-representative}
The complete \(e_{s,z}\) is gauge invariant and
has an analytic rooted family with a bound
uniform in \(s\). It has zero value and first
jet at the flat datum, and
\begin{equation}
 \sum_z e_{s,z}(B)=S_s(e^{iB}),\qquad
 q_e(F)=\frac12\sum_{\mu<\nu}
                              \langle F_{\mu\nu},F_{\mu\nu}\rangle.
 \label{hit:exact-classical-representative}
\end{equation}
Moreover
\begin{equation}
 \sum_z e_{s,z}(B_\xi(w))=S_k(e^{iw})
 \label{hit:exact-classical-counterterm}
\end{equation}
on the actual common ball. Its finite periodic
point function is exactly the evaluation of its
classical infinite companion at periodic data.
\end{lemma}

\begin{proof}
The first identity follows from
\(\sum_z\theta_z=1\). Each original fine plaquette
trace is gauge invariant as a function of its
actual critical orbit, by
\eqref{hpt:individual-physical-plaquette-invariance}.
The action trees keep the factor \(h_s^4\)
at that original plaquette. Summing the
plaquettes in the fixed support of \(\theta_z\)
therefore gives a bounded rooted family,
by \eqref{het:actual-action-tree-bound};
there is no fine-count loss.

For every plaquette orientation the exact identity
\begin{equation}
                  h_s^4\sum_{p\ {\rm in\ that\ orientation}}
                                      \theta_z(m_p)=1
 \label{hit:exact-fine-tent-quadrature}
\end{equation}
holds. To see it in one coordinate, write the
fine index as \(Nn+r\), \(0\le r<N=h_s^{-1}\).
For each fractional shift, including the half
shift of a plaquette midpoint, the sum over
integer translates of the one-dimensional
tent is one. The \(N\) residue classes
therefore sum to \(N\). Take the product of
the four identities.

The actual affine first response is the sampled
affine fine potential modulo a fine gradient,
by \eqref{abc:literal-affine-response}.
Every plaquette has zero first variation at
the flat field, so no second background response
contributes to its quadratic jet. The fine
quadratic action is
\(h_s^4\langle F_{\mu\nu},F_{\mu\nu}\rangle/2\).
Equation \eqref{hit:exact-fine-tent-quadrature}
gives the second identity in
\eqref{hit:exact-classical-representative}.

Classical critical fields lift to periodic
covers: the lifted field and multiplier solve
the same local critical equations and constraints
and stay in the strong uniqueness class.
Orbit uniqueness identifies them with the
constructed background. Each fixed plaquette
trace is consequently unchanged by that lift.
The compact tent, with its periodic sum,
then gives the last assertion. This is a
statement about a classical function, not
about a Gaussian determinant or a quantum
cutoff integral.
Finally use \eqref{hit:actual-all-depth-action-identity}
to obtain \eqref{hit:exact-classical-counterterm}.
\end{proof}

For a scalar input set
\begin{equation}
          r_{s,z}=f_{s,z}-b_s e_{s,z},
                   \qquad b_s=b[f_s^\infty].
 \label{hit:physical-remainder-definition}
\end{equation}
It has complete first jet zero, curvature readout
zero, and input norm at most \(C\epsilon_s\).
The subtracted coefficient multiplies \(S_k\)
exactly after nesting. It is not replaced by a
coarse Wilson action that only approximates \(S_k\).

\subsubsection{A large analytic carrier bound from the actual map}

\begin{lemma}[Coarse block composition]
\label{lem:hit-big-composition}
Suppose \(f_z(0)=Df_z(0)=0\) and
\eqref{hit:old-pointed-input-norm} holds.
For a fixed range \(\xi\le\xi_{\max}\le1\), choose
a target rate \(\kappa_0>0\) and a field radius
\(r>0\) satisfying
\begin{equation}
 C\kappa_0\xi_{\max}
       +C_{\kappa_0}\xi_{\max}r/\rho<\kappa/4,
                  \qquad C_{\kappa_0}\xi_{\max}r/\rho<1/4.
 \label{hit:big-composition-rate-ceiling}
\end{equation}
The actual composition \(f_z(B_\xi(w))\) then
has a rooted final-scale analytic family with
\begin{equation}
                 \|f_z\circ B_\xi\|_{r,\kappa_0}
                              \le C\epsilon r^2\xi^2 .
 \label{hit:actual-big-carrier-bound}
\end{equation}
All constants are independent of the depths.
Using whole-cube Banach groups, the conditions
in \eqref{hit:big-composition-rate-ceiling} contain
no coordinate-count factor \(M^4\).
\end{lemma}

\begin{proof}
Replace each local representative by
\[
 f_{z,X}(B_X)-f_{z,X}(0)-Df_{z,X}(0)[B_X].
\]
Its complete sum is unchanged. The removed
constant and linear sums vanish; their absolute
convergence follows from Cauchy and a rate reserve.

Treat the variables in each old \(M\)-cube as
one complex Banach vector with its supremum
norm. There are \(n=|X|=1+d_M(X)\) such groups.
Multivariate Cauchy and polarization give a
coefficient bound
\(\|f_{z,X}\|_\rho(e/\rho)^m\) at total degree
\(m\), for each distribution of degrees among
the groups. To verify the polarization factor,
insert independent scalars in the \(j\) slots
of a homogeneous polynomial. Cauchy on circles
of radius \(1/j\) bounds its multilinear norm
by \(j^j/j!\le e^j\) times its polynomial norm.
Multiplying over the groups gives the stated
factor, independently of their dimensions.

The grouped actual intermediate family in
\cref{prop:hit-actual-intermediate-fields} has
norm \(d=C_{\kappa_0}\xi r\) for each group.
The Taylor terms of degree at least two are
therefore bounded by
\begin{equation}
 C\|f_{z,X}\|_\rho\,
       n^2(\xi r/\rho)^2
                 \exp(C_{\kappa_0}\xi r n/\rho).
 \label{hit:Banach-block-Cauchy-tail}
\end{equation}
Indeed sum the degree distributions with the
product generating function \((1-ed/\rho)^{-n}\),
subtract its degrees zero and one, and use
Taylor's formula for that scalar function.

The image of the connected old cube carrier
can be joined in the final cube lattice with
at most \(C(1+\xi d_M(X))\) edges. Scale a
spanning curve of length at most \(CMd_M(X)\)
and cover its \(O(\xi M)\) neighborhood by
final \(M\)-cubes; a curve of length \(\ell\)
needs at most \(C(1+\ell/M)\) such cubes.
This joins all grouped output roots at once,
including all coordinate positions in each
old cube image.
Union this join with their actual intermediate
families. Extra fixed costs per Taylor slot
are absorbed in \(C_{\kappa_0}\).
The geometric rate charge is
\(\exp(C\kappa_0(1+\xi d_M(X)))\).

The two exponential costs in this expression
and \eqref{hit:Banach-block-Cauchy-tail} fit in
the input rate by
\eqref{hit:big-composition-rate-ceiling}.
The remaining polynomial \(n^2\) uses another
fixed rate reserve. Summing the old family
gives \eqref{hit:actual-big-carrier-bound}.
This filing is an exact analytic substitution;
its pieces read only the final bond variables
in their indicated connected carriers.
\end{proof}

\subsubsection{Canonical physical inputs and finite-period tails}

\begin{lemma}[The actual canonical input class]
\label{lem:hit-actual-canonical-inputs}
The complete physical normalizer, mismatch and
interacting pointings obtained from
\cref{prop:haf-actual-homogeneous-formation} belong
to the class \eqref{hit:old-pointed-input-norm} on
a fixed output coarse bond ball, after subtracting
their vacuum values. Their rooted norm bounds can
be taken respectively as
\begin{equation}
 C,\qquad C|\tau-1|,\qquad
 C\{(g\mathfrak b)^2+e^{-c\mathfrak b^2}\}.
 \label{hit:actual-canonical-input-sizes}
\end{equation}
For the interacting term this assertion concerns
the homogeneous classical-pullback model with
zero ordinary profile and the physical pointing
of \cref{cor:hpt-physical-marginal-bound}.
These canonical families satisfy
\eqref{hit:finite-period-compatibility}.
\end{lemma}

\begin{proof}
The intrinsic functions are local analytic families
by \cref{prop:haf-actual-homogeneous-formation}.
Their physical arguments are the actual intermediate
data \(p_j(W)\), not independent fine currents.
The \(p_j\) in the physical pointing is in its
one-step axial gauge. On a fixed sufficiently
small output ball both it and
\(p_j^{\rm lit}(W)=M^{j-1}U_j(W)\) are admitted,
and the prescribed node-fixed gauge connects
them. Apply gauge invariance to the complete
pointed function first:
\[
       f_z^{\rm phys}(p_j(W))
                    =f_z^{\rm phys}(p_j^{\rm lit}(W)).
\]
The identity extends holomorphically on this
admitted overlap by the joint chart identity.
The local representatives below are composed
with \(p_j^{\rm lit}\); they are not individually
replaced by gauge invariance.
To see a full analytic carrier bound after this
substitution, group each intrinsic localization cube
as one Banach vector. In output units its side is
\(M\). The actual one-level-before-output map has
whole-cube family norm \(C\rho/L\), by
\cref{prop:hit-actual-intermediate-fields}.
Apply the block Cauchy argument in
\cref{lem:hit-big-composition}, retaining the
constant and linear degrees when necessary.
The old intrinsic cube has side \(LM\) in its
input units; after output rescaling the geometric
join costs \(C(1+d_M(X))\), with only fixed \(L\)
constants. Request a larger intrinsic rate first,
then choose a fixed smaller output radius so the
block Cauchy charge fits that rate. The bound is
linear in the intrinsic rooted norm. There is
no count of its \((LM)^4\) individual coordinates.

The complete interacting pointing is gauge
invariant by
\eqref{hpt:joint-chart-invariance} and
\cref{cor:hpt-physical-marginal-bound}.
For the normalizer use its exact definition
\eqref{afs:actual-N-family}. Point
\(\log J_U(0)\) at its pivot factors and
\(-\tfrac12\operatorname{Tr}\log K\) at its
diagonal free bonds, with the midpoint tents
and the stated vacuum subtraction.
Under a real coarse gauge the actual free
coordinates transform by the block-orthogonal
matrix \(O_g\) of
\eqref{hpt:joint-chart-invariance}. Thus
\(\log K\) conjugates by \(O_g\), and the
colour trace of each diagonal block is invariant.
The pivot determinant transforms on both sides
by adjoint matrices of determinant one.
This proves invariance of the complete pointed
normalizer before any carrier filing.
For the mismatch use its exact pointing
\eqref{afs:actual-mismatch-pointing}.
The diagonal colour trace of
\(K_GQ_{\tau,u}^{-1}\) has the same covariance.
These arguments also give the vanished first
jets by global colour rotations.

The matrix-log localization is
\eqref{afs:actual-matrix-log}, with the actual
nonlocal near/tail construction
\eqref{haf:exact-short-power-tail} and
\eqref{haf:actual-decoupled-root} in place of a
finite-range precision. Group powers before
attaching their parameters and retain every
precision carrier in the diagonal trace.
The absolute rooted logarithm series gives the
normalizer bound. The mismatch interpolation
gives its factor \(|\tau-1|\).
The interacting estimate is the physical
even-coupling bound \eqref{fec:physical-smallness}.
Taking a vacuum difference changes these
bounds by at most a fixed factor.
The original midpoint pointing, the symmetric
middle pivot pattern and the actual orbit map
preserve translations and the full hypercubic
group. The scalar readout is therefore applicable.

We give the finite-period argument explicitly.
Retain the lifted propagator paths in the
background trees, all precision carriers in
the grouped matrix powers, and the full assigned
carriers of the normalized local Gaussian
functionals and Mayer clusters.
In particular, the precision cutoff is by the
full lifted connected carrier \(\Gamma\),
containing both matrix ends and every external
background input, \emph{before} periodic
projection. Its charged length is \(\ell(\Gamma)\).
The stronger input estimate is already
\eqref{haf:actual-precision-carriers}:
the absolute Schur sum carries
\(e^{t_0\ell(\Gamma)}\), by the HET
connected-tree estimate and marked leaves.
Symmetrize a real term with its transposed
lifted carrier, and cut \(\ell(\Gamma)\le R\).
A long winding carrier never enters the
short precision because its projected
endpoints happen to be close.

The covering estimate used in HAF also holds
for cubes on the cover:
\[
 S^{\#\Gamma_M^{\rm cover}}
       \le S^C\exp\{C(\log S)\ell(\Gamma)/M\}.
\]
Thus its stronger lifted weight is paid by
the displayed exponential input norm.
The exact telescoping tails retain that
weight. Combine each short power on the
cover before projection; its finite-range
Hilbert-row proof is unchanged, since the
real cover gap follows by local limits
of the uniform quadratic-form gap.
The resulting short-step ball includes
all its external background support.
For the finite recipe lift the auxiliary
cube parameters and the fixed held root
guards periodically: parameters of period
images are identified, while their full
lifted carriers remain separately charged.
The cover operator is then deck covariant;
its row and column Schur sums descend
without a sheet factor. A nonwrapping
object sees the same single guard and
wall recipe as its companion.
The grouped root and the matrix-log series
therefore have their absolute lifted
carrier bounds as well. Subsequent
absolute Gaussian and cluster estimates
use submultiplicative carrier weights and
an upper bound on the number of integration
variables. The lifted weight is still
submultiplicative, and identifying period
images only reduces that number of variables.
Retaining these labels throughout the
absolute estimates gives the same rooted
canonical-family bound with a reserved
lifted rate.

For this comparison the support of a common
object includes its entire fixed held
root guard and the data and trace collars,
also before projection. These are fixed-layer
enlargements of the connected cube support.
Their size is at most a fixed multiple of
its size. Request that larger rate in the
formation estimate first and reserve the
needed fraction for this enlargement.
A rooted object contained in a common embedded
old neighborhood of radius \(P/8\) is calculated
by exactly the same local formulas on the
torus and on its lift. For a Gaussian trace,
a nonzero endpoint period image is a long
object and is not put in this common part.
For a local Gaussian marginal the comparison
uses the same fixed root, guard and wall
domain; it is not a comparison between
different choices of held root guards.
The normalized finite-dimensional integral,
cutoff, source pointing and cluster coefficient
of a common object are consequently identical.

Objects not paired in this way have connected
carriers reaching distance \(cP\) from their
root, or carrying a period winding. Their
charged lifted cube size is at least \(cP/M\).
The stronger weight just described pays
this size even when projection identifies
some of its cubes.
Reserve half of the exponential input rate.
The absolute remaining rooted sum is bounded
by \(C\epsilon e^{-c\kappa P/M}\).
Projecting a lifted carrier to the torus only
decreases the required cube-tree length.
This constructs the difference family in
\eqref{hit:finite-period-compatibility}, including
its finite-period-tail label. Cauchy derivatives
and any fixed rooted spatial moments use a
further rate reserve and obey the same
exponential comparison with a smaller \(c\).

This proves a comparison of pointed analytic
functions and their local expansions.
The finite determinant and cutoff integral
are not identified with those on a larger
cover, which would introduce extra quantum
fluctuation modes.
\end{proof}

\subsubsection{A slow-field estimate from a small local gauge}

\begin{lemma}[Complete slow-field expansion]
\label{lem:hit-slow-field-sup}
Let \(f_z^\infty\) obey the input norm and
complete gauge identity above. For a continuous
connection with \(\|a\|_{C^{1,\alpha}}\le r\)
put
\[
                 B_\xi(b)=\xi a(\xi m_b).
\]
On a fixed smaller radius, uniformly in
\(0<\xi\le1\),
\begin{equation}
 \left|f_z^\infty(B_\xi)
              -\xi^4q_f(F_a(\xi z))\right|
                         \le C\epsilon r^2\xi^{4+\alpha},
 \quad
 F_{\nu\mu}=\partial_\nu a_\mu-\partial_\mu a_\nu
                                      +i[a_\nu,a_\mu].
 \label{hit:complete-slow-field-bound}
\end{equation}
The same estimate holds uniformly on holomorphic
external-field balls of such connections.
\end{lemma}

\begin{proof}
For sufficiently small \(\xi\), take
\(R=R_0+D\log(1/\xi)\), with \(D\) large compared
with the inverse rooted spatial rates. In old
coordinates \(y=x-z\), let a smooth cutoff
equal one on \(|y|\le R\) and zero on
\(|y|\ge2R\). Set
\[
 A_\mu=a_\mu(\xi z),\qquad
 g(z+y)=\exp\!\left(i\xi\chi_R(y)
                                   \sum_\mu A_\mu y_\mu\right).
\]
Its exponent is bounded by \(C\xi rR\), and
its old first difference by \(C\xi r\).
For \(\xi\) below a fixed threshold,
\(\xi R\le1\). The original and transformed
bond logarithms both have supremum at most
\(C\xi r\), strictly inside the original
analytic ball. For real connections the
gauge is unitary. The equality
\(f_z^\infty(B_\xi^g)=f_z^\infty(B_\xi)\)
extends jointly and holomorphically while
both arguments remain in that ball.
In particular it applies to a holomorphic
connection \(a(w)\). No arbitrary complex
gauge saturation is used.

Within the region where the cutoff is one,
second-order BCH, with the stated gauge
convention, gives
\begin{align}
 B_{\xi,\mu}^g(z+y)
  =\xi^2\bigg\{&
   \sum_\nu\partial_\nu a_\mu(\xi z)
                         (y_\nu+\tfrac12\delta_{\nu\mu})
       +\frac i2\sum_\nu y_\nu[A_\nu,A_\mu]\bigg\}
       +E_\mu(y),\notag\\
 |E_\mu(y)|&\le Cr\xi^{2+\alpha}(1+|y|)^3 .
 \label{hit:nonabelian-local-gauge-expansion}
\end{align}
To check the commutator, the three leading
exponents are \(i\xi Y,i\xi A_\mu,
-i\xi(Y+A_\mu)\), where \(Y=\sum A_\nu y_\nu\).
Their BCH bracket is
\(\tfrac12\xi^2[Y,A_\mu]_{\rm H}\).
The continuous first-order Taylor remainder
costs \(r\xi^{2+\alpha}(1+|y|)^{1+\alpha}\);
the cubic BCH terms are bounded by
\(Cr^3\xi^3(1+|y|)^3\), which fits the
displayed error since \(\alpha<1\).

The leading polynomial is
\(\xi^2(a^F+\mathrm{d}\chi)\), where
\begin{equation}
 \chi(y)=\frac14\sum_{\mu,\nu}
       (\partial_\nu a_\mu+\partial_\mu a_\nu)(\xi z)
                                                   y_\mu y_\nu .
 \label{hit:staggered-gradient-polynomial}
\end{equation}
Its exact forward difference contributes
\(\tfrac12\partial_\mu a_\mu\), so this
identity includes the staggered midpoint
term. The Ward identity removes \(\mathrm{d}\chi\)
in both Hessian arguments.

Expand each local representative through
degree two at zero, and sum. The complete
constant and first terms vanish. For a
carrier contained in the inner patch the
transformed angle is bounded by
\(C\xi^2r(1+Md_M(X))\).
Block Cauchy bounds its cubic remainder
by \(\|f_{z,X}\|_\rho\xi^6r^3\) times a
fixed polynomial in \(1+d_M(X)\).
The exponential Taylor charge can be bounded
using the global supremum \(C\xi r\);
it is paid by a fixed rate reserve.
Summing such carriers gives \(C\epsilon\xi^6r^3\).
A carrier exiting the patch has
\(d_M(X)\ge cR/M\). Its cubic remainder,
estimated using the global angle bound,
is \(C\epsilon\xi^3r^3e^{-cR/M}\) after
reserving rate and coordinate-count polynomials.
Choose \(D\) so that this is also bounded
by \(C\epsilon\xi^6r^3\).

For the quadratic term use
\eqref{hit:complete-Ward-moments}.
Inside the patch, replacing either argument
by the error in
\eqref{hit:nonabelian-local-gauge-expansion}
costs \(C\epsilon r^2\xi^{4+\alpha}\);
the fixed spatial polynomials are summable
against the two-leg exponential moment.
Outside the patch, compare the globally
bounded transformed field with the leading
polynomial using the exponential tail of
that moment. Increasing \(D\) pays this
error as well. The remaining quadratic
term is exactly \(\xi^4q_f(F_a(\xi z))\).
Since \(\xi^6r^3\le r^2\xi^{4+\alpha}\),
the estimate follows for small \(\xi\).

For \(\xi\) above the fixed threshold the
ordinary degree-two Cauchy bound and
\(|q_f(F)|\le C\epsilon r^2\) prove the
same statement with a larger fixed constant.
All estimates hold in the complex norm,
so they are uniform on the indicated
holomorphic external-field ball.
\end{proof}

Apply this lemma to the actual interpolant in
\eqref{hit:actual-intermediate-C1alpha}. For the
remainder \eqref{hit:physical-remainder-definition}
the curvature readout is zero. It has the
complete bound
\begin{equation}
 \|r_{s,z}^\infty(B_\xi(\,\cdot\,))\|_r
                     \le C\epsilon_s r^2\xi^{4+\alpha}.
 \label{hit:actual-small-sup-bound}
\end{equation}
On a finite torus its old shortest period
satisfies \(P_{\rm old}\ge2/\xi\). By
\eqref{hit:finite-period-compatibility} and
\cref{lem:hit-big-composition}, the difference
from this periodic companion has a large
carrier norm bounded by
\begin{equation}
 C\epsilon_s r^2\xi^2e^{-cP_{\rm old}}.
 \label{hit:composed-finite-period-big-bound}
\end{equation}
For small \(\xi\) this is at most the right
side of \eqref{hit:actual-small-sup-bound}.
For \(\xi\ge\xi_0>0\), the large bound
\(C\epsilon_s r^2\xi^2\) suffices, with
a constant depending on \(\xi_0\).
The finite composed remainder therefore
has the same small supremum estimate.
Equivalently the logarithmic gauge patch can
be used in the finite old torus whenever
\(8R<P_{\rm old}\), which holds below a
fixed \(\xi_0\). No nonperiodic polynomial
gauge is applied to that finite torus.

\subsubsection{Block M\"obius interpolation without a volume loss}

The small complete supremum in
\eqref{hit:actual-small-sup-bound} is stronger
than the large norm
\eqref{hit:actual-big-carrier-bound}, but is
not itself a local-family estimate. We now
construct one by an exact block decomposition.

\begin{lemma}[Rooted block interpolation]
\label{lem:hit-block-interpolation}
Let \(F(0)=0\) be holomorphic on a bond ball
\(\|w\|_\infty<r\). Suppose it has a connected
family rooted at a cube \(Q_0\), with norm
at most \(E\) at rate \(\kappa_0>0\), and
\(\|F\|_r\le\delta\). Fix \(0<\vartheta<1\).
There are \(t>0,a>0\), independent of the
number of blocks, for which \(F\) has an
actual rooted local family on
\(\|w\|_\infty<tr\), satisfying
\begin{equation}
                    \|F\|_{tr,a}
                               \le C\delta^\vartheta E^{1-\vartheta}.
 \label{hit:block-interpolation-bound}
\end{equation}
The construction is linear in \(F\).
It preserves holomorphic auxiliary parameters.
\end{lemma}

\begin{proof}
For a finite set \(S\) of final \(M\)-cube
groups let \(w_T\) keep the groups in
\(T\subset S\) and set the others to zero.
Define
\begin{equation}
 F_S(w_S)=\sum_{T\subset S}
                      (-1)^{|S\setminus T|}F(w_T).
 \label{hit:canonical-block-Mobius}
\end{equation}
On a finite torus inclusion-exclusion gives
\(F=\sum_SF_S\) exactly; \(F_\varnothing=0\).
Each \(F_S\) reads only the indicated groups.

Put \(u=t/(1-t)\). For a fixed \(w_S\) of
norm at most \(tr\), write each group as
\(t v_Q\). The holomorphic function
\(F(\sum_{Q\in S}z_Qv_Q)\) has supremum
\(\delta\) on \(|z_Q|<1\), since the
groups are disjoint. The alternating sum
in \eqref{hit:canonical-block-Mobius} keeps
only the Taylor terms of positive degree
in every group. Cauchy in these scalars
therefore gives
\begin{equation}
                b_S:=\|F_S\|_{tr}\le\delta u^{|S|}.
 \label{hit:Mobius-small-sup}
\end{equation}
There is one scalar per block, not one
per bond or colour coordinate.

For a local representative \(G_X\) of the
given large family, its \(S\)-projection
vanishes unless \(S\subset X\), and obeys
the same estimate with \(\|G_X\|_r\)
in place of \(\delta\).
Let \(d_S\) be the length of a minimal
Steiner cube tree joining \(S\) and \(Q_0\).
Then \(d_S\le d_M(X)\) whenever \(S\subset X\).
Choose \(0<\kappa'<\kappa_0\) and \(t\)
small enough that
\(\log(1+u)+\kappa'<\kappa_0\).
Summing the projections of the large family
gives
\begin{equation}
 \sum_S e^{\kappa'd_S}b_S
 \le\sum_X\|G_X\|_r e^{\kappa'd_M(X)}
                              (1+u)^{|X|}
 \le C E .
 \label{hit:Mobius-large-weighted-sum}
\end{equation}
Only triangle inequalities are used here;
cancellations between representatives are
used solely in the complete bound
\eqref{hit:Mobius-small-sup}.

Choose \(a<(1-\vartheta)\kappa'\), and put
\(\lambda=((1-\vartheta)\kappa'-a)/\vartheta>0\).
H\"older's inequality, applied to
\(b_S^{1-\vartheta}b_S^\vartheta\), gives
\begin{equation}
 \sum_S e^{ad_S}b_S
 \le \delta^\vartheta(CE)^{1-\vartheta}
       \left\{\sum_Su^{|S|}e^{-\lambda d_S}\right\}^{\vartheta}.
 \label{hit:Mobius-Holder}
\end{equation}
The last sum is finite for a fixed sufficiently
small \(t\), at every positive \(\lambda\).
Here is a bound that does not assume a
large enough rate for an unweighted animal
count.

For \(n=|S|\), the minimum metric spanning
tree on \(S\cup\{Q_0\}\) has length at
most \(2d_S\): traverse a Steiner tree and
shortcut the resulting walk. Consequently
\[
 e^{-\lambda d_S}
 \le\sum_{\mathcal T\ {\rm on}\ \{0,\ldots,n\}}
       \prod_{\{i,j\}\in\mathcal T}
                e^{-\lambda |Q_i-Q_j|_1/2}.
\]
Allow repeated positions and divide by \(n!\)
to overcount the sets. For a fixed labelled
tree, sum its leaves successively. Each
leaf costs at most
\(B_\lambda=\sum_{x\in\mathbb Z^4}e^{-\lambda|x|_1/2}<\infty\).
The number of labelled trees is
\((n+1)^{n-1}\). Thus
\begin{equation}
 \sum_Su^{|S|}e^{-\lambda d_S}
 \le1+\sum_{n\ge1}\frac{u^n}{n!}
                       (n+1)^{n-1}B_\lambda^n<\infty
                    \quad\hbox{if } e uB_\lambda<1 .
 \label{hit:rooted-set-partition}
\end{equation}
On a periodic cube lattice the sum of the
metric exponential is no larger than this
infinite-lattice sum, by choosing nearest
representatives of the period classes.
The same constant is therefore valid.

File \(F_S\) on a fixed minimal Steiner
tree joining \(S\) to \(Q_0\). A fixed
geometric tie rule suffices; it depends
on no field or coupling parameter.
Equations \eqref{hit:Mobius-Holder} and
\eqref{hit:rooted-set-partition} prove
\eqref{hit:block-interpolation-bound}.
For the infinite lattice, the projections
can be interchanged with the given large
family by the absolute bound
\eqref{hit:Mobius-large-weighted-sum}.
Each local representative has the finite
inclusion-exclusion identity, so their
sum still reconstructs \(F\).
This proves both normal convergence and
the exact identity there. Every operation
is linear and the estimates are uniform
on smaller auxiliary polydiscs, proving
the last assertion.
\end{proof}

For the actual remainder take
\begin{equation}
                  \vartheta=\frac{2+\beta}{2+\alpha}.
 \label{hit:interpolation-exponent}
\end{equation}
Combining \eqref{hit:actual-small-sup-bound}
and \eqref{hit:actual-big-carrier-bound} gives,
per old anchor,
\begin{equation}
 \|r_{s,z}^\infty\circ B_\xi\|_{tr,a}
 \le C(\epsilon_s r^2\xi^{4+\alpha})^\vartheta
          (\epsilon_s r^2\xi^2)^{1-\vartheta}
             =C\epsilon_s r^2\xi^{4+\beta}.
 \label{hit:per-old-anchor-irrelevance}
\end{equation}
The equality of exponents uses
\(2+(2+\alpha)\vartheta=4+\beta\).
The reconstructed complete function remains
gauge invariant. This conclusion does not
assert invariance of its M\"obius pieces.

\subsubsection{A compatible finite-volume filing}

We retain the finite-period information when
making the improved filing. Merely replacing
every finite-volume coefficient by a short
minimal carrier could hide its winding
dependence.

First perform the construction in
\cref{lem:hit-block-interpolation} on the
infinite companion. Evaluate its normally
convergent local family at periodic final
data and project its full carrier to the
torus. Its analytic supremum does not grow:
identifying some of its input variables
still gives values inside the same product
ball. Its required cube length does not
grow either. Keep its lifted carrier label
before projection.

Classical lift compatibility identifies
the actual finite \(B_\xi\) with the
periodic restriction of its infinite
counterpart. The difference of the two
transported quantum densities is therefore
the composed old finite-period difference.
By \eqref{hit:composed-finite-period-big-bound}
it has large norm
\(C\epsilon_s r^2\xi^2e^{-c_wP_{\rm old}}\).
Put \(\gamma_{\rm per}=Mc_w\); the prescribed
period exponent is measured in \(M\)-cube
units, as \(e^{-\gamma_{\rm per}P_{\rm old}/M}\).
Both complete remainders have the small
supremum in
\eqref{hit:actual-small-sup-bound}, so
their difference has it as well.
Apply block interpolation to this difference.
It gives
\begin{equation}
 C\epsilon_s r^2\xi^{4+\beta}
                         e^{-(1-\vartheta)c_wP_{\rm old}}
 \label{hit:interpolated-finite-period-tail}
\end{equation}
for its minimal-carrier family.

Retain its finite-period-tail label and
enlarge each of these carriers by a fixed
shortest period cycle through the root.
This is a legitimate enlargement of a
carrier: it alters no function and adds
no field variables. If \(P_{\rm fin}=\xi
P_{\rm old}\) is the shortest final period,
the extra length is at most
\(C(1+P_{\rm fin}/M)\).
Choose \(a>0\) also to satisfy
\begin{equation}
          Ca\xi_{\max}<(1-\vartheta)\gamma_{\rm per}/2 .
 \label{hit:finite-period-rate-reserve}
\end{equation}
The added weight is paid by
\eqref{hit:interpolated-finite-period-tail}.
The remaining rooted norm is bounded by
\[
 C\epsilon_s r^2\xi^{4+\beta}
                       e^{-c'_wP_{\rm old}}
       = C\epsilon_s r^2\xi^{4+\beta}
                       e^{-(c'_w/\xi)P_{\rm fin}} .
\]
Thus all coefficients with genuine finite-period
dependence remain explicitly in long carriers.
The short part is the projected companion
construction. The output again satisfies a
compatibility estimate of the form
\eqref{hit:finite-period-compatibility}, now
with the improved size.
This argument compares functions through
their absolutely convergent local expansions;
it does not add or remove quantum modes.

\subsubsection{Summing old anchors and the exact counterterm}

\begin{proposition}[Actual homogeneous physical irrelevance]
\label{prop:hit-physical-irrelevant-transport}
Let \(f_{s,z}\) be complete gauge-invariant,
translation-covariant physical pointed densities
on the old coarse ball, with full hypercubic
covariance, input bound \(\epsilon_s\), and
the compatible companion property above.
This includes the actual canonical inputs of
\cref{lem:hit-actual-canonical-inputs}.
Set \(b_s=b[f_s^\infty]\) and
\(r_{s,z}=f_{s,z}-b_s e_{s,z}\).
For every fixed \(L\) there are fixed positive
output radius \(r_{\rm out}\), rate \(a\), and
constant \(C\) such that
\begin{equation}
 \sum_z f_{s,z}(B_\xi(w))
       =b_sS_k(e^{iw})+\sum_y\mathcal R_{s\to k,y}(w),
 \qquad
 \sup_y\|\mathcal R_{s\to k,y}\|_{r_{\rm out},a}
             \le C\epsilon_s r^2L^{-\beta(k-s)} ,
 \label{hit:actual-physical-transport}
\end{equation}
where \(r_{\rm out}=tr\), with the fixed
large-ball radius \(r\) chosen in
\eqref{hit:big-composition-rate-ceiling}.
Each complete output density is gauge invariant
and has zero first jet and zero companion
curvature coefficient. It has an actual
compatible local analytic representation;
its individual representatives need not be
gauge invariant.
\end{proposition}

\begin{proof}
With the final unit tents define the complete
output density by the exact formula
\begin{equation}
 \mathcal R_{s\to k,y}(w)
       =\sum_{z\ {\rm old}}\theta_y(\xi z)\,
                                      r_{s,z}(B_\xi(w)).
 \label{hit:actual-transferred-pointing}
\end{equation}
Only old anchors within fixed final distance
of \(y\) occur. Transfer their roots to \(y\)
by bounded cube joins. By
\eqref{hit:per-old-anchor-irrelevance} and
the compatible finite-period construction,
the rooted norm is at most
\[
 C\epsilon_s r^2\xi^{4+\beta}
                             \sum_z\theta_y(\xi z).
\]
Because \(\xi^{-1}\) is an integer,
the same residue-class tent computation as
\eqref{hit:exact-fine-tent-quadrature} gives
\begin{equation}
                    \xi^4\sum_z\theta_y(\xi z)=1 .
 \label{hit:exact-old-anchor-density}
\end{equation}
This proves the bound in
\eqref{hit:actual-physical-transport}.
It also proves the uniform finite-period-tail
bound after the old-anchor sum, with
exponential rate in \(P_{\rm fin}\).

The tents sum to one in \(y\).
Sum \eqref{hit:actual-transferred-pointing}
over \(y\), substitute the definition of
\(r_{s,z}\), and apply
\eqref{hit:exact-classical-counterterm}.
This proves the exact identity in
\eqref{hit:actual-physical-transport}.

Under a small final coarse gauge, actual
orbit covariance of \(U_k\) changes
\(V_{s,k}\) by an old-lattice site gauge.
The complete old functions are invariant,
and the tents are geometric scalars.
Thus each complete output density is
invariant. Its first jet is zero, either
by the chain rule or global colour rotations.

For the companion curvature coefficient
test a final affine potential of constant
curvature. The actual affine response
\eqref{abc:literal-affine-response},
followed by the literal partial mean,
is the old affine potential with its
length factor, modulo an old gradient.
The node-frame correction is included
as in \eqref{abc:successive-affine-response}.
The old constant term is another
polynomial gradient. The remaining
old curvature is \(\xi^2F\).
The second chain rule contains no second
background response because \(Dr_{s,z}(0)=0\).
The two-leg Ward identity removes all
gradient terms and gives
\[
 \frac12D^2(r_{s,z}\circ B_\xi)(0)[a^F,a^F]
                                  =\xi^4q_{r_s}(F)=0.
\]
All polynomial substitutions are justified
by the exponential moments of the actual
response and old Hessian. Multiplication
by the fixed root tents and their finite
anchor sum preserves zero. Translation
and hypercubic covariance of the complete
density are preserved since
\(\xi^{-1}\) is an integer.
The compatible filing was constructed
above and is an exact representation
of this same complete function.
\end{proof}

\begin{corollary}[Uniform sum of proper old scales]
\label{cor:hit-all-old-scales}
If \(0\le s<k\) and \(\epsilon_s\le E_0\)
on common parameter domains, their transported
remainders satisfy
\begin{equation}
 \sup_y\left\|\sum_{s<k}\mathcal R_{s\to k,y}
                              \right\|_{r_{\rm out},a}
 \le Cr^2\sum_{s<k}\epsilon_sL^{-\beta(k-s)}
                      \le \frac{Cr^2E_0}{L^\beta-1}.
 \label{hit:summed-old-scale-bound}
\end{equation}
Their complete sum is invariant, has zero
companion curvature coefficient, and has
the same finite-period compatibility.
The total classical counterterm is exactly
\((\sum_{s<k}b_s)S_k\).
\end{corollary}

\begin{proof}
Sum the positive rooted bounds. The age is
at least one, so the geometric series is
\(\sum_{n\ge1}L^{-\beta n}=(L^\beta-1)^{-1}\).
Normal convergence preserves the identities
and holomorphy. The finite-period comparison
has at least the final-period exponential
at every old scale, so its sum has the
same property.
\end{proof}

\subsubsection{Rate and analytic-radius bookkeeping}

For arbitrary fixed odd \(L\), the proof gives
a fixed positive output rate and one fixed
smaller field ball. The loss is independent
of the number of old scales.
One can prescribe an output rate \(a>0\)
in advance for proper old ages.
Choose \(\kappa_0,\kappa'\) with
\((1-\vartheta)\kappa'>a\), then choose \(t\)
using \eqref{hit:rooted-set-partition}.
Since \(\xi_{\max}=1/L\), the conditions
\eqref{hit:big-composition-rate-ceiling}
and \eqref{hit:finite-period-rate-reserve}
are met for all proper old ages by taking
\(L\) sufficiently large relative to these
fixed rates and a fixed ratio \(r/\rho\).
The old-carrier geometric charge contains
\(\xi d_M(X)\), and the finite-period
padding charge contains \(P_{\rm fin}=\xi
P_{\rm old}\). Both improve under this
coarsening. Write \(\gamma'_{\rm per}=Mc'_w\).
The remaining finite-period tail is
\(e^{-\gamma'_{\rm per}P_{\rm fin}/(M\xi)}\).
To retain a prescribed compatibility exponent
\(\gamma_{\rm per,target}\) in \(M\)-cube
units, impose
\(L\gamma'_{\rm per}\ge\gamma_{\rm per,target}\),
in addition to the padding reserve.
These conditions contain no \(M\): the
common \(M^{-1}\) length conversion cancels.
They compensate the fixed \(1-\vartheta\)
interpolation loss for every proper age.
Thus a chosen carrier rate and a chosen
finite-period compatibility exponent can
both be retained under sequential proper
transport when the blocking integer is
large enough. Choose \(M\) afterwards to
realize the requested macroblock field rate.

The grouped family bound in
\cref{prop:hit-actual-intermediate-fields}
has an amplitude constant independent of
\(M,L\), while its allowed radius can
depend on \(L\). Thus
\eqref{hfm:uniform-amplitude-coarse-contraction}
gives the necessary buffer in literal coordinates:
an outer data ball of radius \(R\)
maps inside a working input ball of
radius \(tR/4\) if \(C/L<t/4\).
Fluctuation substitution additionally
requires \(g\mathfrak b<tR/4\), with
the fixed chart constants included.
These inequalities can be imposed in
the order: rates and \(t\), blocking
integer, localization size, common
small physical radius, then coupling
floor. The nonlinear radius need
not be uniform as \(L\) grows.

The theorem proves the physical homogeneous
transport estimate from actual intermediate
maps and explicit canonical inputs.
A complete RG trajectory additionally
requires formation and substitution on
the buffered domains at every step.
No transport theorem for arbitrary
independent fine-current cores, proper
history domains or changing nonlinear
towers is asserted here.

% END LAYERED EXPANSION formation/homogeneous_irrelevant_transport.tex

% BEGIN LAYERED EXPANSION formation/literal_centre_physical_formation.tex
\subsection{Formation at the literal centre}
\label{layered:literal-centre-formation}

The literal intermediate map decreases bond angles
by a factor \(L^{-1}\). We use that decrease to
evaluate the old interaction on its original
analytic ball while allowing a larger ball of
output data. The relative contour gauge gives
an exact change to the axial step on their real
overlap. Holomorphy on the larger ball will
follow from the literal integral itself.

\subsubsection{The input and parameter domains}

Fix the middle pivots and real product colour
balls of \cref{prop:rcs-cutoff-transfer-covariance}.
Choose a carrier rate \(\kappa>0\).
An input \(h\) consists of pointed functions
\begin{equation}
 h_z(B)=\sum_{X\ni q(z)}h_{z,X}(B_X),\qquad
 \|h\|_{r,\kappa}
    =\inf_{\{h_{z,X}\}}\sup_z
             \sum_Xe^{\kappa d_M(X)}\|h_{z,X}\|_r .
 \label{lcf:input-norm}
\end{equation}
Here \(B=i^{-1}\log V\) is in old coarse units,
and \(X\) is a connected set of old \(M\)-cubes.
The complete \(h_z\) is gauge invariant, normalized
at \(B=0\), and covariant under old-lattice
translations and the hypercubic group.
Its companion coefficient
\eqref{hit:scalar-curvature-coefficient} is zero.
Take the Banach quotient of local analytic
families by representations whose complete
function is zero at every anchor.
The infimum in \eqref{lcf:input-norm} is
its quotient norm. Write
\(\mathfrak H_{r,\kappa}\) for the closed
subspace with the stated identities, and put
\begin{equation}
                   H(V)=\sum_z h_z(i^{-1}\log V).
 \label{lcf:total-old-interaction}
\end{equation}
The sum is finite on a torus.
We also write \(h_z(V)=h_z(i^{-1}\log V)\).
The conditions on the complete functions
are closed and linear: evaluation, the
Ward moment extractor and covariance
identities are continuous in this norm.
Complex variations of \(h\) below are
in this space. Reality is required only
when a real transfer is asserted.

For families over compatible periods,
include an infinite companion and take
the supremum of these spatial norms.
Fix \(\gamma_{\rm per}>0\) in cube units
and also require
\begin{equation}
 \|h\|_{r,\kappa,\gamma_{\rm per}}
 =\max\left\{\sup_P\|h^P\|_{r,\kappa},
       \|h^\infty\|_{r,\kappa},
       \sup_P e^{\gamma_{\rm per}P/M}
       \|h^P-\Pi_Ph^\infty\|_{r,\kappa/2}\right\}<\infty .
 \label{lcf:compatible-family-norm}
\end{equation}
The notation \(\Pi_P\) is periodic
evaluation, as in
\eqref{hit:finite-period-compatibility}.
Write \(\mathfrak H_{r,\kappa,\gamma_{\rm per}}\)
for this space. The same norm notation
will be used for output pointings before
their scalar coefficient is subtracted.

Let \(t\in(0,1/8)\) be the fixed shrink in
\cref{lem:hit-block-interpolation}. We shall
use the radii
\begin{equation}
                            r=tR .
 \label{lcf:two-radii}
\end{equation}
Choose an intrinsic output rate \(\kappa_{\rm int}\)
larger than a fixed geometric multiple of
\(\kappa+\gamma_{\rm per}\). The whole-cube intermediate bound
in \cref{prop:hit-actual-intermediate-fields}
has an amplitude \(C_p\), independent of \(L,M\)
on an allowed smaller ball. Choose an odd \(L\)
so large that
\begin{equation}
 \frac{64C_p}{Lt}<\min(1,\kappa_{\rm int}),\qquad
 \frac{C\kappa_{\rm int}}L<\frac{\kappa}{8}.
 \label{lcf:blocking-buffer}
\end{equation}
The fixed constants include the rate and padding
reserves in the transport theorem.
Choose \(M\) afterwards for those rates and for
the covariance construction. Then choose \(R>0\)
inside the common background and graph radii.
These choices are independent of depth and volume.
The nonlinear radius may depend on \(L\).

Fix a real reference coupling \(\theta_0\), and set
\begin{equation}
 g_0=\theta_0^{-1/2},\qquad
 \mathfrak b_0=A_1(\log\theta_0)^p,\quad p>1/2,
 \qquad x=\theta_0\tau .
 \label{lcf:frozen-reference}
\end{equation}
Choose \(\delta>0\) so that
\(|\tau-1|<8\delta\) lies in the admitted
Gaussian disc, and retain \(|\tau-1|<\delta\)
for the Lipschitz conclusions.
The reference gauge weight \(\theta_0\), all
tame thresholds, and the real cutoff
\(\chi_{\mathfrak b_0}\) stay fixed as \(x\)
varies. Numerical powers of \(g_0\) and
Gaussian determinants are separated below.

\subsubsection{The literal graph and old arguments}

At stage \(j\), define
\begin{equation}
 V_*(W)=\mathsf p_j(W):=M^jU_{j+1}(W),
                 \qquad W=e^{iw},\quad \|w\|_\infty<R .
 \label{lcf:literal-centre}
\end{equation}
The output spacing is one; the old input spacing
is \(1/L\) in these units. Thus \(\mathsf p_j\)
is the map \(p_{j+1}^{\rm lit}\) in
\eqref{hfm:uniform-amplitude-coarse-contraction}.
It is not put in its one-step axial gauge.
The actual graph about \(V=V_*(W)\) is
\begin{equation}
 Z_V(u)=e^{i\Phi_V(u)}V,\qquad
 M(Z_V(u))=W,\qquad \Phi_V(u)|_E=u,
 \qquad \Lambda_V(u)=i^{-1}\log Z_V(u).
 \label{lcf:literal-graph}
\end{equation}
All logarithms here are on the admitted branches.

\begin{lemma}[Admission on the two radii]
\label{lem:lcf-domain-admission}
For \(R\) as above,
\[
 \|\log V_*(W)\|_\infty\le C_pR/L<r/64.
\]
There is \(C_{\rm dom}=C_{\rm dom}(L)\) such
that the actual graph and logarithm satisfy
\[
 \|\Lambda_V(u)-\Lambda_V(0)\|_\infty
                  +\|\Phi_V(u)\|_\infty
                              \le C_{\rm dom}\|u\|_\infty .
\]
Consequently every argument of \(h\) in
\eqref{lcf:literal-graph} stays in its old
\(r\)-ball for all nested fluctuation boxes
used by the formation theorem, provided
\begin{equation}
       64C_{\rm dom}|\gamma|\mathfrak b_0<r
 \label{lcf:old-source-admission}
\end{equation}
and the fixed graph ceiling holds. These
domains are holomorphic in \(w,\gamma\).
At the real physical value \(\gamma=g_0\),
the retained regular fibre is precisely this
graph on the free box.
\end{lemma}

\begin{proof}
The first assertion is the literal mean bound
and \eqref{lcf:blocking-buffer}.
The graph estimate is the joint local pivot
estimate in \eqref{afs:actual-pivot-chart},
followed by the logarithm of
\(e^{i\Phi_V(u)}V\). Its constants do not
involve \(j\). Increase \(C_{\rm dom}\)
to include the fixed nested box factors
and radial margins of that theorem.
The two inequalities leave a fixed margin
inside the old \(r\)-ball. In particular
the same old function is evaluated before
and after substitution.

At real \(W\), the background curvature is
below the fixed regular ceiling when
\(R\) is small. The actual three-free-edge
pivot anchor \eqref{afs:regular-pivot-anchor}
puts every other point of that regular
coarea fibre, with its free values in the
box, inside the pivot uniqueness ball.
It therefore has the graph value. Conversely
the graph remains in the regular coarea
set by the same curvature and free-angle
bounds. The argument uses its fixed absolute
regular ceiling, not an extra history
plaquette threshold.
The background, graph and logarithm are
jointly holomorphic, and every argument of
the old functions remains in their declared
ball. This proves the stated complex domain
directly.
\end{proof}

The contour tests can be removed on the physical
box. Set
\(\varepsilon_0=A_0g_0(\log\theta_0)^p\) with
the same retained test convention as the
homogeneous step. The relative argument is
\(T_{y,x}(Z_V(g_0B))T_{y,x}(V)^{-1}\).
Internal contours use no middle pivot.
The normalized path/mean derivative bound
therefore puts this argument within
\(C_Lg_0\mathfrak b_0\) of one, uniformly in
the background. Choose \(A_0>2C_LA_1\).
The tests are then identically one on
\(\chi_{\mathfrak b_0}\), exactly as in
the proof of \eqref{afs:exact-physical-step}.
The analytic deformation in \(\gamma\) below
is defined after this removal. The remaining
real integration box does not vary.

\subsubsection{The change to axial gauge}

For a real \(V\), define
\[
 s_V(x)=T_{y,x}(V)\quad(x\in B(y)),\qquad s_V(y)=1,
                     \qquad V^{\rm ax}=V^{s_V}.
\]
This is a unitary gauge equal to one at
the output nodes. Put
\begin{equation}
 G_V(u)=\sum_{x\ne y(x)}
 \left[1-\frac12\operatorname{tr}
   \{T_{y(x),x}(Z_V(u))T_{y(x),x}(V)^{-1}\}\right].
 \label{lcf:relative-gauge}
\end{equation}

\begin{lemma}[Exact axial comparison on the real overlap]
\label{lem:lcf-real-axial-identity}
Whenever both old arguments are admitted,
the literal-centre integral with gauge \(G_V\)
equals the axial physical step.
The free-coordinate change is
\((O_Vu)(b)=\operatorname{Ad}(s_V(b_-))u(b)\).
It preserves the real product colour balls.
\end{lemma}

\begin{proof}
The endpoint gauge law gives the exact identities
\begin{equation}
 T_{y,x}(V^{s_V})=1,\qquad
 T_{y,x}(Z^{s_V})
                 =T_{y,x}(Z)T_{y,x}(V)^{-1}.
 \label{lcf:exact-gauge-matching}
\end{equation}
Thus the absolute contour weight and its
tame test at \(Z^{s_V}\) are the relative
ones at \(Z\), including their numerical
normalizer.
The actual graph uniqueness and joint gauge
identity give
\begin{equation}
 \Phi_{V^{\rm ax}}(O_Vu)=O_V\Phi_V(u),\qquad
                  Z_{V^{\rm ax}}(O_Vu)=Z_V(u)^{s_V}.
 \label{lcf:joint-axial-chart}
\end{equation}
For real \(V\), \(O_V\) is block orthogonal
and has determinant one. It preserves the
free box, Haar/coarea density and regular
plaquette bounds. It is independent of
the integration variable \(u\).

Gauge invariance gives
\(S_j(Z^{s_V})=S_j(Z)\) and
\(h_z(Z^{s_V})=h_z(Z)\) on their admitted
overlap. Apply these identities to the
complete functions \(h_z\), before local
filing. Equations
\eqref{lcf:exact-gauge-matching} and
\eqref{lcf:joint-axial-chart} now identify
the two coarea integrals by a real
orthogonal change of free coordinates.
No old function is evaluated outside its
own ball in this comparison.
\end{proof}

For fixed \(L\) this overlap contains a
neighborhood of \(W=1\). On the full
\(R\)-ball the literal integral below is
defined without evaluating \(h\) on the
axial representative. Its holomorphy
there will not require complex gauge
saturation of the old domain.

\subsubsection{The integral and its Gaussian factors}

Let \(d\lambda_W\) denote regular coarea
disintegration of normalized Haar measure
on \(M(Z)=W\). Restrict it to the retained
free box about \(V_*(W)\), and use the
contour tests just removed. Define
\begin{equation}
 \mathcal I_j(W;x,h)=
 z(\theta_0,\varepsilon_0)^{-n_s}
 \int_{M(Z)=W}^{\rm reg}
 \chi_{\mathfrak b_0}(u/g_0)\,
 e^{-xS_j(Z)+H(Z)-\theta_0G_V(u)}
                                      \,d\lambda_W(Z),
 \quad V=V_*(W).
 \label{lcf:literal-integral}
\end{equation}
The notation for \(\chi\) means the graph
free coordinates of \eqref{lcf:literal-graph}.
The omitted contour tests equal one on
this real support. Here \(n_s\) is the
number of old sites other than output nodes.
This definition uses only the retained
regular component.
For complex data we use its literal graph
pullback displayed below, with the same
fixed real integration variables \(B\).

Criticality and nesting give
\begin{equation}
 S_j(V_*(W))=S_{j+1}(W),\qquad
 D_u S_j(Z_V(u))|_{u=0}=0 .
 \label{lcf:physical-centering}
\end{equation}
Define the matrices
\begin{equation}
 K_A(V)=D_u^2S_j(Z_V(u))|_0,\quad
 K_G(V)=D^2G_V(0),\quad
 K_\tau(V)=\tau K_A(V)+K_G(V),\quad K(V)=K_1(V).
 \label{lcf:Gaussian-precision}
\end{equation}
They include the complete inner background
and outer constraint-curvature derivatives.
The gap and full lifted carrier bounds
are those of
\eqref{haf:actual-composed-precision} and
\eqref{haf:actual-precision-carriers}.
We do not replace a precision carrier by
the distance between its matrix endpoints.

Let \(J_V(u)\) be the exact coarea density
\eqref{afs:actual-coarea-density}. Put
\begin{align}
 R^A_V(u)&=S_j(Z_V(u))-S_j(V)
                                    -\tfrac12u^{\mathsf T}K_A(V)u,\notag\\
 R^G_V(u)&=G_V(u)-\tfrac12u^{\mathsf T}K_G(V)u,\notag\\
 P_{\gamma,\tau,h}(V,B)
  &=-\tau\gamma^{-2}R^A_V(\gamma B)
     -\gamma^{-2}R^G_V(\gamma B)
     +\log\frac{J_V(\gamma B)}{J_V(0)}
                         +H(Z_V(\gamma B))-H(V).
 \label{lcf:intrinsic-exponent}
\end{align}
The divided remainders have removable values
at \(\gamma=0\). The graph, precision and
cutoff are independent of \(\gamma\).
Define
\begin{equation}
 Z_j(\gamma;V,\tau,h)=
 \int\chi_{\mathfrak b_0}(B)
       e^{P_{\gamma,\tau,h}(V,B)}
                               \,d\mu_{K_\tau(V)^{-1}}(B).
 \label{lcf:normalized-integral}
\end{equation}
For complex data the normalized Gaussian
uses the determinant branch from positive
real data. The local expansion below fixes
the zero-free logarithm.

\begin{lemma}[Exact normalization at the literal centre]
\label{lem:lcf-exact-normalization}
Let \(n_f,n_c\) count positive old and output
bonds, \(d=3(n_f-n_c)\), and \(r_{\rm g}=3n_s\).
At \(V=V_*(W)\),
\begin{equation}
 \det K_\tau(V)=\tau^{d-r_{\rm g}}\det K(V).
 \label{lcf:physical-determinant-power}
\end{equation}
For the physical value \(\gamma=g_0\), the
coarea integral factors exactly as
\begin{align}
 \mathcal I_j(W;x,h)
  ={}& C_{\rm num}(\theta_0)g_0^d(2\pi)^{d/2}
       J_V(0)(\det K_\tau(V))^{-1/2}\notag\\
 &\qquad{}\times
       e^{-xS_{j+1}(W)+H(V)}
                        Z_j(g_0;V,\tau,h),\notag\\
 C_{\rm num}(\theta_0)
  ={}&\sigma_0^{\,n_f-n_c}
                   z(\theta_0,\varepsilon_0)^{-n_s}.
 \label{lcf:exact-factorization}
\end{align}
Here \(\sigma_0\) is Haar density at zero.
All determinant branches start at positive
real data and \(\tau=1\).
\end{lemma}

\begin{proof}
The graph change of variables has exactly
the density \eqref{afs:actual-coarea-density}.
Rescaling \(u=g_0B\) contributes \(g_0^d\).
The fine Haar and coarse delta densities
give \(\sigma_0^{n_f-n_c}\), as in
\eqref{afs:exact-numerical-prefactor}.
The action is centered by
\eqref{lcf:physical-centering}. Subtract its
full quadratic part and that of the relative
gauge. Gaussian integration with precision
\(K_\tau\) gives
\eqref{lcf:exact-factorization}, including
the extracted value \(H(V)\).

The action Hessian annihilates the tangent
to node-fixed gauge orbits, by constrained
criticality and gauge invariance.
For the relative gauge at a nonaxial \(V\),
the orbit derivative of its contour logarithm
at an internal site is
\(-\operatorname{Ad}(T_{y,x}(V))\lambda(x)\).
This is an invertible real orthogonal
change of the gauge parameter.
It gives the same form decomposition as
\eqref{agm:orthogonal-form-sectors}, with
gauge dimension \(r_{\rm g}\).
The complete precision is positive on
real data by HAF. Applying
\cref{prop:agm-exact-mismatch} after this
parameter change proves
\eqref{lcf:physical-determinant-power}.
Equivalently conjugate to the axial graph
using \eqref{lcf:joint-axial-chart}.
Both arguments retain the full constraint
curvature in \(K_A\).
The identities continue to the admitted
complex domain.
\end{proof}

Define
\begin{equation}
 N_j(W)=
 \log\frac{J_{V_*(W)}(0)}{J_1(0)}
       -\frac12\log\frac{\det K(V_*(W))}{\det K(1)} .
 \label{lcf:Gaussian-stage}
\end{equation}
It is independent of \(x,h\).
Equation \eqref{lcf:physical-determinant-power}
puts the entire remaining power of \(\tau\)
in a scalar factor. At the identity the
raw pivot derivative is \(L^{-3}I_3\), so
\(J_1(0)=L^{9n_c}\).
Consequently the factor outside the
field-dependent terms is
\begin{equation}
 C_j^0(\theta_0,x)=
 C_{\rm num}(\theta_0)g_0^d(2\pi)^{d/2}
 J_1(0)(\det K(1))^{-1/2}
                               \tau^{-(d-r_{\rm g})/2}.
 \label{lcf:scalar-prefactor}
\end{equation}
In particular no logarithm of this factor
is included in the intrinsic remainder.

\subsubsection{Admission and filing of the old source}

For an intrinsic construction before composing
with \(V_*(W)\), use the centered extension
\begin{equation}
 s_V(u)=S_j(Z_V(u))-S_j(V)
                           -DS_j(V)[D\Phi_V(0)u].
 \label{lcf:centred-extension}
\end{equation}
It is the coarse version of the completion
in \eqref{afs:centered-action-completion}.
It has zero first jet. Its Gaussian Hessian
is the full composed Hessian, including
the outer constraint curvature. At
\(V=V_*(W)\), its last term is zero
and it equals the action difference in
\eqref{lcf:intrinsic-exponent}.
This extension supplies local functions
on an old background ball; it does not
alter the physical integral.

\begin{lemma}[The ordinary source budget]
\label{lem:lcf-old-source-budget}
Suppose \(\|h\|_{r,\kappa}\le H_*\).
On the intrinsic background ball
\(\|i^{-1}\log V\|_\infty<r/8\), its
source in \eqref{lcf:intrinsic-exponent}
has an actual potential family with
budget
\begin{equation}
       \|\mathbf V^h_\gamma\|_{\kappa_{\rm int}}
                      \le C_{L,M,r}H_*|\gamma|\mathfrak b_0 .
 \label{lcf:old-source-potential-bound}
\end{equation}
The input functions and their arguments
are the ones in \eqref{lcf:input-norm};
no extension of their domain is used.
\end{lemma}

\begin{proof}
For each representative use exactly
\begin{equation}
 V^h_{z,X}(V,B)
       =h_{z,X}(\Lambda_V(\gamma B))
                              -h_{z,X}(\Lambda_V(0)).
 \label{lcf:actual-old-difference}
\end{equation}
On the nested cutoff boxes the path
\(\zeta\mapsto h_{z,X}(\Lambda_V(\zeta\gamma B))\)
is admitted for
\(|\zeta|<cr/(C_{\rm dom}|\gamma|\mathfrak b_0)\).
Its supremum is bounded by
\(\|h_{z,X}\|_r\).
One-variable Cauchy, with the constant
value subtracted, gives
\[
 |V^h_{z,X}|
       \le C\frac{C_{\rm dom}|\gamma|\mathfrak b_0}{r}
                                        \|h_{z,X}\|_r .
 \]
This bound treats all the coordinates
of \(X\) at once.

The local substitution reads only the
actual double-block enlargement of \(X\).
After changing to output units, its
carrier has at most \(C(1+d_M(X)/L)\)
output \(M\)-cube edges. Cover the
scaled old spanning curve and its
fixed output-size collar, as in
\cref{lem:hit-big-composition}.
Thus the requested rate cost fits the
old one by \eqref{lcf:blocking-buffer}.

We also convert the pointed norm to
an incidence norm before applying HAF.
If an input carrier contains a bond
in cube \(Q\), its size is at least
the cube distance from \(q(z)\) to
\(Q\). Reserve part of the old rate
and sum the anchors:
\[
 \sup_Q\sum_{z,X\ni Q}
       e^{\kappa d_M(X)/2}\|h_{z,X}\|_r
       \le C_{\kappa,M}H_* .
 \]
Indeed each old cube contains \(O(M^4)\)
unit anchors and the remaining sum is
a convergent metric exponential.
This factor affects the source budget,
not its analytic domain.
Combine all origins with the same
carrier before estimating its potential.
The source has no additional Gaussian
coordinates and introduces no number
of potential types proportional to
the volume.
These estimates prove
\eqref{lcf:old-source-potential-bound}.
\end{proof}

Choose a fixed ceiling \(4H_*+1\) for the
norm of \(h\). The actual interaction
budgets of HAF and the preceding lemma
are small on a disc
\begin{equation}
                   |\gamma|<c_{H_*}/\mathfrak b_0 .
 \label{lcf:auxiliary-coupling-disc}
\end{equation}
The positive constant \(c_{H_*}\) is
below the graph and old-source ceilings
and the UM potential ceiling.
It can depend on \(H_*,L,M,r\).
The Gaussian structure and its domains
were chosen before \(H_*\); only this
coupling ceiling depends on that norm.
Increasing the reference coupling floor
makes \(g_0\) lie in one half of this
disc.

\subsubsection{Formation on the larger output ball}

\begin{lemma}[A local family on the two radii]
\label{lem:lcf-large-ball-family}
The normalized integral
\eqref{lcf:normalized-integral}, with the
centered extension for intrinsic filing,
is zero free. Its logarithm is a holomorphic
family on \eqref{lcf:auxiliary-coupling-disc},
the larger \(\tau\)-disc, and the
\(4H_*+1\) input ball. After composing
with \(V_*(W)\), it has a local family
on \(\|w\|_\infty<R\), with any
fixed output rate chosen above.
The constants are independent of \(j\)
and volume.
\end{lemma}

\begin{proof}
Use the actual covariance construction
\eqref{haf:actual-decoupled-root}.
It truncates full lifted precision
carriers, including external background
inputs, before projection. Held root
guards contain the entire support of
each graph, potential and required
data collar. The matrix dimension of
a local functional remains at most
\(C(LM)^4\#Z\).
The old source changes its potential
but not this Gaussian structure.

The action, relative gauge and coarea
budgets are those proved in HAF.
Equation \eqref{lcf:old-source-potential-bound}
supplies the ordinary input budget on
its actual domain. The UM, Mayer and
connected cluster proof therefore gives
the intrinsic analytic family on
\eqref{lcf:auxiliary-coupling-disc}.
This is the existing proof applied
with the specified source, rather than
an assumption on its desired output.
It fixes a logarithm by its convergent
local series. No whole-volume lower
bound for a cutoff probability is used.

The intrinsic carrier is on output
localization cubes, of side \(LM\)
in old input units. Group its external
background coordinates as one Banach
vector per such cube. Its image in
output units has side \(M\).
Use the Banach-valued proof for all old
bond midpoints in this final cube:
the target is their supremum norm,
the pure mean row bound holds for
all of them, and only the same fixed
neighboring macroblocks are read.
Their number does not enter that proof.
The literal whole-cube family has
norm \(C_pR/L\), by
\cref{prop:hit-actual-intermediate-fields}.
Its ratio to the intrinsic background
radius \(r/8\) is at most \(8C_p/(Lt)\).
The block Cauchy argument of
\cref{lem:hit-big-composition}, retaining
all degrees, therefore gives a charge
\(\exp(C C_p|X|/(Lt))\).
The intrinsic rate reserve pays it
by \eqref{lcf:blocking-buffer}.
This counts cube groups, not their
\((LM)^4\) individual background
coordinates. The bounded root joins
and rate reserve give a bounded
composition operator on the local
family norm.

The map \(V_*(W)\) is independent of
\(\gamma,\tau,h\), so this composition
preserves their holomorphy and any
norm-valued derivative estimate.
On physical data the centered extension
equals the action difference by
\eqref{lcf:physical-centering}.
Thus the resulting family is the
logarithm of the literal integral
already defined, after its stated
Gaussian factors have been removed.
\end{proof}

\subsubsection{Pointing the old source}

Write \(\langle A\rangle_{\gamma,V,\tau,h}\)
for the normalized expectation obtained
from \eqref{lcf:normalized-integral}
by inserting \(A\). For one old anchor set
\begin{equation}
 \Delta h_{z,V,\gamma}(B)
                  =h_z(Z_V(\gamma B))-h_z(V).
 \label{lcf:complete-old-insertion}
\end{equation}
This definition precedes every local
representation of \(h_z\).
Let \(\mathscr F^0_{j,y}(W;\gamma,\tau)\)
be the HPT pointing with \(h=0\),
composed with the literal centre using
complete gauge invariance on the real
overlap. Extend it to the \(R\)-ball
by the literal formula. Define
\begin{equation}
 \begin{split}
 \mathscr F^{\rm raw}_{j,y}(W;\gamma,\tau,h)
   ={}&\mathscr F^0_{j,y}(W;\gamma,\tau)\\
   &+\sum_{z\ {\rm old}}\theta_y(z/L)
       \int_0^1
        \langle\Delta h_{z,V,\gamma}\rangle_
                              {\gamma,V,\tau,\ell h}\,d\ell ,
       \qquad V=V_*(W).
 \end{split}
 \label{lcf:complete-output-pointing}
\end{equation}
The tents use output units.
The sum over \(y\) is
\(\log Z_j(\gamma;V,\tau,h)\):
differentiate that logarithm along
\(\ell h\) and use the partition of
unity of the tents.

\begin{lemma}[A marked source family]
\label{lem:lcf-invariant-source-pointing}
The pointings
\eqref{lcf:complete-output-pointing}
have local analytic representations
on the \(R\)-ball. Their complete
functions are gauge invariant and
hypercubic covariant. They are even
in \(\gamma\). The part depending on
\(h\) vanishes at \(\gamma=0\).
\end{lemma}

\begin{proof}
Introduce independent ordinary source
parameters \(\zeta_z\), multiplying
\(\Delta h_{z,V,\gamma}\).
The incidence estimate in
\cref{lem:lcf-old-source-budget} is
uniform on a fixed product disc in
these parameters: their number adds
no Gaussian coordinates or potential
types. Differentiate the convergent
connected cluster expansion in
\(\zeta_z\). Each resulting object
contains one old-source insertion
rooted at \(z\), with its entire
original carrier and the graph
collar retained. The marked vertex
has budget \(C|\gamma|\mathfrak b_0\)
times the old rooted norm.
Sum its unmarked branches by the
same convergent cluster bound.
This constructs the local family
of the expectation in
\eqref{lcf:complete-output-pointing}
without a whole-volume factor.
The old origins are combined on
their assigned carriers before
forming the local Gaussian functionals.
Root transfer by \(\theta_y(z/L)\)
costs only fixed \(L,M\) constants.

Under a small real site gauge the
actual graph changes by the same
linear adjoint map in its free
variables as in
\eqref{hpt:joint-chart-invariance}.
Both \(Z_V(\gamma B)\) and \(V\)
remain in the old ball. Gauge
invariance of the complete \(h_z\)
therefore preserves
\eqref{lcf:complete-old-insertion}.
The Gaussian, cutoff and coarea
change by the corresponding real
orthogonal coordinates.
The normalized expectation is
unchanged. The same is true for
the \(h=0\) pointing, so each
complete function in
\eqref{lcf:complete-output-pointing}
is invariant. At the output,
orbit covariance of the actual
background induces exactly such
an old site gauge between the
two centres. The argument uses
complete functions before local
representatives are composed.
Joint holomorphy extends the
identity on the admitted complex
domain.

Lattice covariance follows from
\cref{prop:rcs-cutoff-transfer-covariance}.
The middle-pivot set and real colour
balls in that proposition are the
ones fixed here. The old insertions
are covariant by the input assumption,
and their tents transform with
their anchors. Thus this assertion
concerns the actual integral,
not a subsequent average of
pointings.

Finally every occurrence of the
group field has argument \(\gamma B\).
Hence the exact transformation
\((\gamma,B)\mapsto(-\gamma,-B)\)
preserves the exponent and all
insertions, including their
Taylor subtractions.
The centered Gaussian and fixed
cutoff are even. The HAF filing
commutes with this transformation;
each local Gaussian integration
and connected extraction therefore
gives an even function of
\(\gamma\), as in FEC.
At zero coupling
\(\Delta h_{z,V,0}=0\) exactly.
The normalized empty Gaussian
selection has zero derivative
in an old source. These facts
prove the last assertion at
the level of the marked family.
\end{proof}

\begin{lemma}[Even-coupling bound on the output family]
\label{lem:lcf-quadratic-smallness}
On a fixed smaller \(\tau\)-disc,
\(\|h\|_{r,\kappa}\le4H_*+1\), and
\(|\gamma|\le c_{H_*}/(2\mathfrak b_0)\),
the pointed logarithm satisfies
\begin{equation}
 \|\boldsymbol{\mathscr F}^{\rm raw}_j
                     (\gamma,\tau,h)\|_{R,\kappa}
        \le C_{H_*}\{
             |\gamma|^2\mathfrak b_0^2
                            +e^{-c\mathfrak b_0^2}\}.
 \label{lcf:quadratic-formation-bound}
\end{equation}
The constants are uniform in depth and volume.
\end{lemma}

\begin{proof}
The original HAF estimate and
\cref{lem:lcf-old-source-budget} give
the norm bound
\(C_{H_*}|\gamma|\mathfrak b_0+
Ce^{-c\mathfrak b_0^2}\)
on \eqref{lcf:auxiliary-coupling-disc},
including the marked family of
\cref{lem:lcf-invariant-source-pointing}.
The composition on the \(R\)-ball is
a bounded operator on this norm by
\cref{lem:lcf-large-ball-family}.
At zero coupling the logarithm is
the cutoff wall contribution
\(\log\mu_{K_\tau(V)^{-1}}(\chi_{\mathfrak b_0})\).
Its pointed norm is bounded by
\(Ce^{-c\mathfrak b_0^2}\), and it
is independent of \(h\).

Subtract this zero-coupling value.
The remaining activity-valued
function is analytic and even,
with a double zero at the origin.
The norm-valued Cauchy argument of
\cref{fec:quadratic-formation} on a
fixed inner circle of
\eqref{lcf:auxiliary-coupling-disc}
gives \(C_{H_*}|\gamma|^2\mathfrak b_0^2\).
All these bounds are for a rooted
local family; no extensive error
is differentiated.
Adding the cutoff wall term proves
\eqref{lcf:quadratic-formation-bound}.
\end{proof}

\subsubsection{Companions and the quotient norm}

In \eqref{lcf:compatible-family-norm},
the difference representations retain
their lifted carriers and finite-period-tail
labels, as in the construction following
\eqref{hit:interpolated-finite-period-tail}.
The infimum is over compatible
representations with those labels.
The formula compares complete functions;
it does not discard the labels when
an infimum is approached.

\begin{lemma}[Preservation of the period norm]
\label{lem:lcf-period-compatibility}
Fix \(\gamma_{\rm per}>0\) in cube units.
The rate choices and then \(L\) can be
made so that the bound
\eqref{lcf:quadratic-formation-bound}
also holds in the full norm
\(\|\cdot\|_{R,\kappa,\gamma_{\rm per}}\)
for input norm at most \(4H_*+1\).
The same exponent is retained at
every step.
\end{lemma}

\begin{proof}
The old shortest period is \(LP\)
when the output period is \(P\).
Decompose the old input as
\[
                 h^{LP}=\Pi_{LP}h^\infty+w^{LP}.
 \]
Retain the entire old carriers
and the finite-period-tail labels
of \(w^{LP}\) in the potential
\eqref{lcf:actual-old-difference}.
The source incidence estimate is
linear in its norm and gives
its extra factor
\(e^{-\gamma_{\rm per}LP/M}\).
Marking this source in the
connected expansion retains that
factor. The normalized empty
selection has zero old-source
derivative, and parity makes
the resulting difference even
in \(\gamma\) with value zero
at the origin. Cauchy on
\eqref{lcf:auxiliary-coupling-disc}
therefore supplies its
\(|\gamma|^2\mathfrak b_0^2\)
factor as well.

For the remaining difference
between the finite integral and
its companion, compare the full
lifted operator, Gaussian and
cluster objects before projection,
as in
\cref{lem:hit-actual-canonical-inputs}.
The precision has the
\(e^{t_0\ell(\Gamma)}\) bound
\eqref{haf:actual-precision-carriers}.
The near cutoff is by its whole
lifted carrier. Finite auxiliary
parameters and held guards are
lifted periodically with image
parameters identified.
Include all guards and data
collars in the common-object
support. Nonwrapping objects
therefore have identical
coefficients. Every unmatched
primitive Gaussian object or
connected cluster pays a fixed
cube-unit period exponential.
Request an intrinsic rate large
enough to make that exponent
larger than the prescribed
\(\gamma_{\rm per}\).

These estimates hold uniformly
on the auxiliary coupling disc.
At zero coupling the only
wrapping terms are cutoff
wall terms, with their
\(e^{-c\mathfrak b_0^2}\) factor.
The nonconstant difference in
\(\gamma\) is even and has
a double zero. The same
norm-valued Cauchy argument
therefore supplies
\(|\gamma|^2\mathfrak b_0^2\)
with the period factor retained.

In refiling the old source tail,
its period \(LP\) leaves an
exponential reserve for fixed
carrier and guard padding at
output period \(P\).
Choose \(L\) so that the
remaining exponent is at least
\(\gamma_{\rm per}\).
All exponents are in \(M\)-cube
units; \(M\) cancels from
this choice. The same estimates
apply to the vacuum value
at \(W=1\). This proves
the full-norm assertion.
\end{proof}

Define the vacuum-subtracted pointings
at the physical coupling by
\begin{equation}
 F_{j,y}(W;x,h)=
 \mathscr F^{\rm raw}_{j,y}(W;g_0,\tau,h)
                  -\mathscr F^{\rm raw}_{j,y}(1;g_0,\tau,h),
 \qquad \mathbf F_j(x,h)=(F_{j,y})_y .
 \label{lcf:formation-map}
\end{equation}
Their sum is the field-dependent
remainder. The subtracted total is
part of a scalar vacuum factor.

The definition
\eqref{lcf:complete-output-pointing}
depends only on the complete \(h_z\).
Two local representations of the
same pointings give the same
\(\mathbf F_j\) in the quotient.
All estimates above can be applied
to representatives whose norms
approach the input quotient norm.
On any finite-dimensional affine
subspace of inputs, choose local
representatives for the base and
directions. The cluster construction
is holomorphic in their coefficients.
Uniform local bounds, Cauchy
estimates and polarization give
a norm-convergent Taylor series
on a smaller input ball.
Thus \(\mathbf F_j\) is a
holomorphic map of the quotient
Banach spaces, including their
period component.

\subsubsection{The formation estimate and its derivatives}

\begin{proposition}[Formation with a fixed radius buffer]
\label{prop:lcf-buffered-formation}
Choose the rates, \(t,L,M,R,r=tR\)
as above. Fix \(H_*<\infty\).
There is a coupling floor such that
for every reference \(\theta_0\)
above that floor, the literal
regular step is defined on
\(\|w\|_\infty<R\),
\(|x/\theta_0-1|<\delta\), and
\(\|h\|_{r,\kappa,\gamma_{\rm per}}\le H_*\).
Its nonconstant Gaussian term is
\(N_j\) from \eqref{lcf:Gaussian-stage},
independent of \(x,h\).
The map \(\mathbf F_j\) is
holomorphic on fixed larger
relative coupling and profile
domains and satisfies
\begin{align}
 \|\mathbf F_j(x,h)\|_{R,\kappa,\gamma_{\rm per}}
                &\le C_{H_*}\epsilon(\theta_0),\notag\\
 \|D_h\mathbf F_j(x,h)[a]\|_{R,\kappa,\gamma_{\rm per}}
                &\le C_{H_*}\epsilon(\theta_0)
                              \|a\|_{r,\kappa,\gamma_{\rm per}},\notag\\
 \|\partial_x\mathbf F_j(x,h)\|_{R,\kappa,\gamma_{\rm per}}
                &\le C_{H_*}\epsilon(\theta_0)/\theta_0,\qquad
 \epsilon(\theta_0)=\theta_0^{-1}(\log\theta_0)^{2p}.
 \label{lcf:formation-API}
\end{align}
Its complete pointings are gauge invariant,
translation covariant and hypercubic covariant,
and are normalized at the flat field.
For real data the map is real.
The constants are independent of depth
and compatible volume.
\end{proposition}

\begin{proof}
At \(\gamma=g_0\),
\eqref{lcf:quadratic-formation-bound}
and \cref{lem:lcf-period-compatibility}
give the first bound, including
the vacuum subtraction.
The cutoff tail is absorbed in
\(\epsilon(\theta_0)\) after increasing
the floor, because \(2p>1\).
No cutoff varies in this estimate.

The same bound holds on
\(\|h\|<4H_*+1\) and the larger
fixed \(\tau\)-disc.
The quotient construction permits
Cauchy in any input direction
on a fixed norm margin.
It gives the \(h\) derivative
with its direction norm retained.
Cauchy on the larger \(\tau\)-disc
gives \(C_{H_*}\epsilon(\theta_0)\)
for \(\partial_\tau\mathbf F_j\).
Since \(\tau=x/\theta_0\),
\(\partial_x=\theta_0^{-1}\partial_\tau\).
This proves all three estimates.
Neither an indicator nor the
input \(h\) is rescaled by \(x\).

Gauge invariance and lattice
covariance are
\cref{lem:lcf-invariant-source-pointing};
the cutoff symmetry is
\cref{prop:rcs-cutoff-transfer-covariance}.
Their definitions on complete
pointings make the identities
independent of the chosen local
representatives.
For real data the integral is
positive and its coarea density
has the identity sign. Its
logarithm and source interpolation
are real on the selected branch.

Finally \eqref{lcf:exact-factorization}
and \eqref{lcf:physical-determinant-power}
give \(N_j\) with no remaining
field-dependent \(\tau\) power.
The scalar factors are outside
the remainder. Every domain and
norm provider is uniform in \(j\)
and volume, as required.
\end{proof}

For two inputs on the retained margins,
integrating the derivative bounds gives
\begin{equation}
 \|\mathbf F_j(x,h)-\mathbf F_j(\widetilde x,\widetilde h)\|
 \le C_{H_*}\epsilon(\theta_0)
       \left\{|x-\widetilde x|/\theta_0
                       +\|h-\widetilde h\|_{r,\kappa,\gamma_{\rm per}}\right\}.
 \label{lcf:joint-Lipschitz}
\end{equation}
The norm on the left is the output norm
in \eqref{lcf:formation-API}.
It controls the local interaction
rather than an extensive logarithm.

\subsubsection{Transport on the same field ball}

For \(i<k\), let \(\mathsf p_{i,k}(W)\)
be the literal partial mean of the
actual orbit at depth \(k\), expressed
on the old unit lattice. Put
\(\xi=L^{-(k-i)}\), and define the
linear transport of complete pointings by
\begin{equation}
 (T_{i,k}h)_y(W)
       =\sum_z\theta_y(\xi z)\,
                           h_z(\mathsf p_{i,k}(W)).
 \label{lcf:transport-definition}
\end{equation}
The sum means the periodized sum
on a torus. Its lifted carriers
are retained in the norm.

\begin{corollary}[A common field ball for transport]
\label{cor:lcf-common-radius-transport}
Choose \(0<\beta<\alpha<1\).
After increasing the fixed \(L\),
the rate and radius choices above
give
\begin{equation}
 \|T_{i,k}h\|_{r,\kappa,\gamma_{\rm per}}
 \le C_T L^{-\beta(k-i)}
                       \|h\|_{r,\kappa,\gamma_{\rm per}},
 \qquad i<k .
 \label{lcf:common-radius-transport}
\end{equation}
The constants are uniform in the depths
and compatible periods.
The transports obey
\(T_{j,k}T_{i,j}=T_{i,k}\).
\end{corollary}

\begin{proof}
Apply
\cref{prop:hit-physical-irrelevant-transport}
with its large test ball equal
to \(R\), its input radius
equal to \(r=tR\), and its
output radius equal to \(tR=r\).
The literal intermediate map has
size \(C_{\mathsf p}\xi R\)
on the large test ball, by
\cref{prop:hfm-fixed-leg-density}
and the actual intermediate-field
construction in the transport proof.
For every proper age,
\(\xi\le L^{-1}\).
The inequalities
\eqref{lcf:blocking-buffer}
therefore place the entire
intermediate vector family
inside the old \(r\)-ball.
They also pay the big-composition
Cauchy and carrier-rate costs.
Thus no shrink is iterated
with \(k-i\).
The Ward--Möbius estimate
and old-anchor count give
\eqref{lcf:common-radius-transport}.
The infimum over local
representations passes through
this linear bound.

The period estimate uses
\(P_{\rm old}=P_{\rm new}/\xi\).
The dimensionless retained
exponent from
\eqref{hit:interpolated-finite-period-tail}
is multiplied by \(\xi^{-1}\)
in output cube units.
The large-\(L\) choice retains
\(\gamma_{\rm per}\), including
the carrier padding reserve.
It is uniform for every
proper age.

For the semigroup identity,
actual nesting identifies
successive literal partial
means up to a gauge on the
old lattice. Apply complete
gauge invariance before
choosing local representatives.
For integer \(\xi^{-1}\), the
nodal hat identity is
\begin{equation}
 \sum_z\theta_y(\xi z)\theta_z(s)
                         =\theta_y(\xi s).
 \label{lcf:hat-refinement}
\end{equation}
Indeed both sides are multilinear
on each old unit cube, and
they agree at its vertices;
the knots of the left-hand
coarse hat are aligned with
old vertices. Periodization
preserves the identity.
Substitution in
\eqref{lcf:transport-definition}
now proves the semigroup law.
It is an identity of complete
functions, not of filed pieces.
\end{proof}

\subsubsection{The exact physical step}

Let \(e_{j+1,y}\) be the action pointing
from the actual fine-plaquette hats.
It has sum \(S_{j+1}\), companion
coefficient one and a uniform
norm on the fixed field ball.
For a compatible family \(a\),
write \(b[a]\) for the scalar
quadratic coefficient of its
infinite-volume companion,
defined by its Ward kernel.
This is a bounded linear
functional in the norms used
here: Cauchy controls the
quadratic jet, and the reserved
spatial rate controls its
second moment.

Choose the complete Gaussian pointings
\(\mathbf N_j\) with sum \(N_j\)
as in the homogeneous Gaussian
construction. Define
\begin{align}
 b_j&=b[\mathbf N_j],&
       n_j&=\mathbf N_j-b_j e_{j+1},\notag\\
 r_j(x,h)&=b[\mathbf F_j(x,h)],&
       f_j(x,h)&=\mathbf F_j(x,h)-r_j(x,h)e_{j+1}.
 \label{lcf:scalar-and-profile-projections}
\end{align}
All pointings in these differences
are normalized at the flat field.
The infinite-volume scalar
coefficient of \(n_j,f_j\) is zero.
The norm of \(n_j\) is bounded
uniformly in \(j\), and
\(b_j,n_j\) are independent
of \(x,h\). For \(r_j,f_j\),
the three bounds
\eqref{lcf:formation-API}
hold with a changed fixed
constant, both on the larger
output ball and after restriction
to the working ball.

\begin{corollary}[Closure of one regular physical step]
\label{cor:lcf-exact-profile-step}
On the domains of
\cref{prop:lcf-buffered-formation},
there is a field-independent,
holomorphic nonzero factor
\(C_j(x,h)\) such that
\begin{equation}
 \mathcal I_j(W;x,h)
   = C_j(x,h)
       \exp\left\{-x_+S_{j+1}(W)
                         +\sum_y(h_+)_y(W)\right\},
 \qquad
 \begin{cases}
 x_+=x-b_j-r_j(x,h),\\
 h_+=T_{j,j+1}h+n_j+f_j(x,h).
 \end{cases}
 \label{lcf:exact-profile-step}
\end{equation}
The new profile belongs to the
same quotient space on radius \(r\),
with the same spatial and
period exponents.
Its norm satisfies
\par\noindent
\begin{minipage}{\linewidth}
\begin{equation}
 \|h_+\|_{r,\kappa,\gamma_{\rm per}}
 \le C_TL^{-\beta}\|h\|_{r,\kappa,\gamma_{\rm per}}
                    +C_N+C_{H_*}\epsilon(\theta_0).
 \label{lcf:one-step-profile-bound}
\end{equation}
\popQED\par
\end{minipage}
\end{corollary}

\begin{proof}
By the partition of unity
for the hats,
\[
 H(\mathsf p_j(W))
                  =\sum_y(T_{j,j+1}h)_y(W).
 \]
The exact factorization
\eqref{lcf:exact-factorization},
\eqref{lcf:Gaussian-stage} and
\eqref{lcf:formation-map}
therefore give
\[
 \mathcal I_j(W;x,h)=C_j(x,h)
   \exp\{-xS_{j+1}(W)
             +H(\mathsf p_j(W))+N_j(W)
                              +\sum_yF_{j,y}(W;x,h)\}.
 \]
Here \(C_j\) is the scalar
\eqref{lcf:scalar-prefactor}
times
\(\exp\{\sum_y\mathscr F^{\rm raw}_{j,y}
                       (1;g_0,\tau,h)\}\).
The zero-free normalized integral
makes this factor nonzero
on the retained domains.
For real data it is positive.

Insert
\eqref{lcf:scalar-and-profile-projections}
and \(\sum_y e_{j+1,y}=S_{j+1}\).
This proves
\eqref{lcf:exact-profile-step}
with its stated signs.
Gauge invariance and hypercubic
covariance hold for every complete
term. The exact action
counterterm and the Ward readout
give zero companion scalar
coefficient for \(h_+\).
The period estimates and the
common-radius transport give
membership and
\eqref{lcf:one-step-profile-bound}.
\end{proof}

The integral in this result is the
regular step with action \(xS_j\),
the literal Haar and coarea densities,
the relative contour gauge and the
fixed free-bond colour-ball cutoff.
Its old input is a complete
coarse-field interaction \(h\).
The result supplies a fixed-radius
producer for this model.
It makes no identification with
an arbitrary transported interaction
on independent fine-current cores
or a conditional history measure.

\subsubsection{A derivative of constant diagonal weights}

Let \(a=(a_1,\ldots,a_4)\) range over a
fixed complex polydisc about \(1\), and set
\begin{equation}
 S_j^a(V)=\mathscr A_{h_j}^a(U_j^a(V)),\qquad
       \mathsf p_j^a(W)=M^jU_{j+1}^a(W).
 \label{lcf:weighted-action-and-centre}
\end{equation}
The action, actual orbit and literal
representative are those of
\cref{prop:ach-actual-anisotropic-background}.
The mean, middle pivots, relative contour
gauge and frozen free cutoff are unchanged.
For fixed \(a\), the required lattice
symmetries are translations and independent
coordinate reflections. Axis permutations
act jointly on \(a\) and the fields.

Write \(\mathfrak H^{\rm ref}_{r,\kappa,\gamma_{\rm per}}\)
for the quotient space with gauge invariance,
these reflection and translation identities,
flat-field normalization, and zero
coefficient in each of the six curvature
planes. Its norm is
\eqref{lcf:compatible-family-norm}.
At \(a=1\), the hypercubic subspace with
zero scalar coefficient is contained
in this space. The six readouts are
the \(b_{\mu\nu}\) of
\cref{ch:agc-anisotropic-Gaussian}.
Independent reflections exclude cross-plane
quadratic curvature terms, so these
readouts exhaust the quadratic curvature
part. They are continuous by the same
Cauchy and Ward moment bounds used above.

Put
\[
 K_\tau^a(V)=
       \tau D_u^2S_j^a(Z_V(u))|_0+D^2G_V(0).
 \]
Define \(\mathbf F_j^a\) by
\eqref{lcf:formation-map} with the
weighted action and centre. Define
\(N_j^a\) by
\eqref{lcf:Gaussian-stage} with
\(K_1^a,\mathsf p_j^a\), and use its
complete Gaussian pointings
\(\mathbf N_j^a\) as above.

\begin{proposition}[The first weight jet of formation]
\label{prop:lcf-constant-weight-jet}
On a smaller fixed weight polydisc, the
literal formation construction is jointly
holomorphic in \(a,x,h\), with input
\(\|h\|_{\mathfrak H^{\rm ref}}\le H_*\).
Its formation \(\mathbf F_j^a(x,h)\)
satisfies the three bounds
\eqref{lcf:formation-API} uniformly in \(a\),
and also
\begin{equation}
 \|D_a\mathbf F_j^a(x,h)[e]\|_{R,\kappa,\gamma_{\rm per}}
                 \le C_{H_*}\epsilon(\theta_0)|e|.
 \label{lcf:weight-formation-API}
\end{equation}
The nonconstant Gaussian \(\mathbf N_j^a\)
is independent of \(x,h\).
Its norm and first weight derivative
are uniformly bounded.
Covariance under axis permutations
is a joint identity in \(a,h\);
full hypercubic covariance at
\(a=1\) is recovered for
hypercubic \(h\).
\end{proposition}

\begin{proof}
We first verify the providers needed
uniformly on a weight disc.
The principal change in the actual
joint equation is the constant local
operator \(K_a-K_1\) of
\eqref{ach:principal-perturbation}.
The exponentially weighted
\(\mathcal W_p\) joint inverse in
\cref{lem:ach-Lp-joint-inverse}
therefore gives an \(a\)-dependent
inverse by a convergent Neumann
series on a fixed weight disc.
Each factor has a fixed physical
carrier; retain the connected
chain of those carriers in the
series. Its exponentially weighted
cell sum is uniform on a larger
weight disc.

Smooth unit-cell localization is
bounded on \(\mathcal W_p\) by the
product rule through two differences.
The remaining actual nonlinear
vertices have arity at least two
and uniform Banach Cauchy bounds,
by \eqref{ach:residual-bounds}
and the unchanged constraint
recursion. The ordered-tree proof
of \cref{prop:het-actual-effective-action}
thus applies with this inverse.
Its scalar majorant and its
connected-tree length charge
are uniform in \(a\).
This is an all-arity carrier
bound, not just a separate
bound for each fixed jet.
Multiplying each original
fine-plaquette density by
\(a_\mu a_\nu\) preserves
its \(h_j^4\) bound and the
cancellation of the fine
plaquette count. Consequently
the full lifted precision
family, including every
background leg and winding
label, has the uniform
bound required in
\eqref{haf:actual-precision-carriers}.

The actual free form has
the lower and upper bounds
\eqref{ach:actual-free-coercivity}.
The unchanged contour gauge
gives the complementary gauge
gap. Equivalently use the
relative-form perturbation
argument in
\cref{lem:agc-weighted-gap}
for the physical Gaussian.
The full composed constraint
curvature is retained.
On a smaller complex product
domain the normalized inverse
is a Neumann continuation of
the real gap. The HAF
near/tail construction therefore
works uniformly, always cutting
by the whole lifted carrier
before projection.

The free inverse and its
local \(L^p\) upgrade have
constants for the odd mesh
parameter \(N\), independent
of its factorization \(N=L^j\).
Their linear amplitude is
therefore uniform in \(L\).
The nonlinear radius may
shrink with \(L\).
Use this amplitude in
\eqref{lcf:blocking-buffer}
when choosing \(L\), and
decrease the common \(R\)
afterward. The HFM plain-row
bound then gives
\(\|\mathsf p_j^a(W)\|\le C R/L\)
uniformly on the larger
weight disc. This verifies
the literal old-source
admission, rather than
assuming it for the
deformed orbit.

All algebraic identities for
the relative gauge, the joint
graph, Haar and coarea densities
are independent of the action.
Weighted criticality and nesting
give \eqref{lcf:physical-centering}
with \(S_j^a,\mathsf p_j^a\).
The gauge-null Hessian and
relative gauge rank give the
same scalar power
\(\tau^{d-r_{\rm g}}\).
Thus the Gaussian defined
with \(K_1^a\), after vacuum
subtraction, remains independent
of \(x,h\).

The old-source budget uses
analytic norms and gauge
invariance, not axis permutations
or scalar isotropy.
It holds on
\(\mathfrak H^{\rm ref}\).
The intrinsic exponent remains
invariant under
\((\gamma,B)\mapsto(-\gamma,-B)\).
All activity, source-rooting
and period estimates above
are uniform on the larger
weight disc, giving
\eqref{lcf:formation-API}.
Cauchy in the fixed weight
margin gives
\eqref{lcf:weight-formation-API}.
It also bounds the first
weight jet of \(\mathbf N_j^a\).

Finally the middle-pivot and
colour-ball identity of
\cref{prop:rcs-cutoff-transfer-covariance}
commutes with reflections and
with simultaneous permutation
of \(a\). The actual weighted
orbit has that covariance by
uniqueness. Source interpolation
before filing gives the stated
covariance of the complete
pointings. No invariance of
an individual filed term
is used.
\end{proof}

For a fixed old input \(h\) in
\(\mathfrak H^{\rm ref}\), define
\(T_j^a h\) by
\eqref{lcf:transport-definition}
with \(\xi=L^{-1}\) and
\(\mathsf p_j^a\).
The big-composition estimate
already gives a uniform norm
on the larger weight disc:
the whole old-coordinate
family is \(CR/L<r/64\);
the coarsened old carrier
rate pays the output rate;
and the one-step old-anchor
count is the fixed \(L^4\).
The constant may depend
on this fixed \(L\).
Keeping the old period labels
uses \(P_{\rm old}=LP\)
exactly as above. Hence
Cauchy in \(a\) gives
\begin{equation}
 \left\|\left.D_aT_j^a h[e]\right|_{a=1}
                       \right\|_{r,\kappa,\gamma_{\rm per}}
                   \le C_L|e|\,
                            \|h\|_{r,\kappa,\gamma_{\rm per}}.
 \label{lcf:transport-weight-jet}
\end{equation}
Every plane readout of \(T_j^a h\)
is zero on this weight domain:
the actual affine response
\eqref{ach:affine-response}
is independent of \(a\) modulo
an actual gradient. The complete
first jet vanishes, and Ward
annihilates the gradient in
the quadratic jet. Differentiating
this exact identity proves that
the profile in
\eqref{lcf:transport-weight-jet}
also has six zero readouts.
It is a bounded profile;
no power of \(g_0\) is
claimed for it.

At the isotropic point, take
\(\sum_\mu e_\mu=0\).
Let \(e^{\rm str}_{j+1,e}\) be the
complete fine-action pointing
\begin{equation}
 (e^{\rm str}_{j+1,e})_y(W)
 =\sum_{p\parallel\mu\nu}
      (e_\mu+e_\nu)\theta_y(x_p)\,
                                  a_p(U_{j+1}(W)).
 \label{lcf:stress-action-pointing}
\end{equation}
Its six coefficients are \(e_\mu+e_\nu\),
and its sum is
\(\left.D_aS_{j+1}^a(W)[e]\right|_{a=1}\).
Indeed the derivative of the orbit
makes no contribution to the
total first variation at fixed
mean, by actual constrained
criticality. The coefficient
claim follows from the exact
fine-hat quadrature and the
affine response.

Use the transfer index \(j\) and put
\(b_{j,\mu\nu}^a=b_{\mu\nu}[\mathbf N_j^a]\).
This is the increment from depth \(j\)
to \(j+1\), indexed \(j+1\) in the
cumulative sum in the Gaussian comparison.
Joint permutation covariance gives numbers
\(u_j,v_j\) such that
\(\partial_{a_\lambda}b_{j,\mu\nu}^a|_1=u_j\)
for \(\lambda\in\{\mu,\nu\}\) and
equals \(v_j\) otherwise. Put
\(m_j=u_j-v_j\). Then
\begin{equation}
 \left.D_a\mathbf N_j^a[e]\right|_{a=1}
       =m_j e^{\rm str}_{j+1,e}+q_{j,e},
 \qquad
 \|q_{j,e}\|_{r,\kappa,\gamma_{\rm per}}\le C|e|,
 \qquad b_{\mu\nu}[q_{j,e}]=0 .
 \label{lcf:Gaussian-weight-jet}
\end{equation}
The normalization of this Gaussian
is the canonical contour and coarea
normalization of
\cref{prop:agc-actual-anisotropic-increment}.
That proposition identifies
\[
 m_j=\frac{11}{3\pi^2}\log L+O_\sigma(L^{-\sigma j})
                  \qquad(0<\sigma<1).
 \]
The uniform norm in
\eqref{lcf:Gaussian-weight-jet}
follows from weight Cauchy
and the uniform action-pointing
norm. The six-plane readouts
of its difference vanish by
the definition of \(m_j\).
Formation and its first weight
jet can likewise be projected
onto the six readouts and
their zero-readout complement,
retaining the
\(C_{H_*}\epsilon(\theta_0)\)
bounds.

These statements differentiate a constant
diagonal deformation on its actual
common domain. They do not assert
closure of a finite anisotropic
trajectory or admit a spatially
varying weight.

% END LAYERED EXPANSION formation/literal_centre_physical_formation.tex

% BEGIN LAYERED EXPANSION formation/frozen_centre_weight_response.tex
\subsection{The first weight derivative with a frozen centre}
\label{fcw:section}

The weighted formation step moves its critical centre.
A stress insertion in a fixed regular kernel does not.
We compare their first derivatives on the actual fibre.
A node-fixed gauge aligns the relative contour factors
exactly. The remaining difference is a surface source
on the fixed free colour balls. Its Gaussian tail is
small in the rooted local norm, including normalization
and the finite-period companions.

This is a first-derivative result. A fixed-centre
weighted kernel need not have the uniform small
formation disc of the moving-centre construction:
its off-centre linear force can have size
\(O((a-1)/g_0)\). No such disc is asserted here.

\subsubsection{The two actual integrals}

Use the homogeneous regular model and parameter
choices of \cref{prop:lcf-buffered-formation}.
In particular
\[
 g_0=\theta_0^{-1/2},\qquad
 b_0=\mathfrak b_0=A_1(\log\theta_0)^p,
 \quad p>\tfrac12,\qquad x=\theta_0\tau.
\]
The reference \(\theta_0\), colour-ball radius,
middle pivots and all tame thresholds are fixed.
The old input satisfies \(\|h\|\le H_*\) in the
compatible quotient norm of
\eqref{lcf:compatible-family-norm}. For the analytic
estimates retain the larger bounded \(h\) ball and
\(\tau\) disc used in the formation proof.
All old interactions below are complete gauge
invariant functions on their original field ball.

Fix real output data \(W\) first, and put
\[
 V_a=\mathsf p_j^a(W),\qquad V_0=V_1,
 \qquad Z_0(B)=Z_{V_0}(g_0B).
\]
The index zero on \(V_0\) means the frozen
isotropic centre. It is not the depth index.
Let \(\mathcal I^{\rm mov}(a)\) be the actual
weighted LCF integral with centre \(V_a\), action
\(S_j^a\), relative contour factor \(G_{V_a}\),
and the product colour balls in its own free
coordinates. Define the frozen-centre integral
on the same Haar-disintegrated fibre by
\begin{equation}\label{fcw:frozen-integral}
 \mathcal I^{\rm fr}(a)=
 z(\theta_0,\varepsilon_0)^{-n_s}
 \int_{M(Z)=W}^{\rm reg}
 \chi_{b_0}(B_0(Z))
     e^{-xS_j^a(Z)+H(Z)-\theta_0G_{V_0}(Z)}\,d\lambda_W(Z),
 \qquad Z=Z_0(B_0(Z)).
\end{equation}
Here the old lattice spacing is one and the raw
free angle is \(g_0B_0\). The measure and graph in
\eqref{fcw:frozen-integral} are frozen. The
regular/tame tests have their strict admission
margin and equal one on the retained support.
For a real weight near one both charts lie in
this same regular component. The size of that
real weight interval may depend on \(g_0b_0\).
Only its derivative at one will be used.

At \(a=1\), write their common value as
\(\mathcal I_0\) and the corresponding normalized
expectation as \(\langle\cdot\rangle_0\).
The envelope identity for the actual critical
action gives
\begin{equation}\label{fcw:frozen-stress-insertion}
 D_a\log\mathcal I^{\rm fr}(1)[e]
        =-x\langle E_{j,e}\rangle_0,
 \qquad
 E_{j,e}(Z)=\left.D_aS_j^a(Z)[e]\right|_1
                         =\sum_z E_{j,e,z}(Z).
\end{equation}
The pointings on the right are the original
fine-plaquette stress pointings, evaluated on
the isotropic actual orbit. The derivative of
that orbit pairs to zero by constrained
criticality at fixed literal mean.

\subsubsection{Exact contour alignment and the surface velocity}

At an old site \(x\) in the output block rooted
at \(y\), write
\(T_a(x)=T_{y,x}(V_a)\), \(T_0(x)=T_{y,x}(V_0)\).
Define
\begin{equation}\label{fcw:aligning-gauge}
 s_a(x)=T_0(x)^{-1}T_a(x),\qquad
       \widetilde V_a=V_a^{s_a},\qquad
 C_a(b)=i^{-1}\log(\widetilde V_a(b)V_0(b)^{-1}).
\end{equation}
Then \(s_a(y)=1\), \(C_1=0\), and
\(T_{y,x}(\widetilde V_a)=T_0(x)\).
Put \(c_{e,b}=D_aC_a(b)|_1[e]\).
The logarithm in this definition is only needed
near its zero at \(a=1\).

\begin{lemma}[The exact frozen-coordinate cutoff]
\label{lem:fcw-exact-cutoff-change}
For real admitted arguments, the change
\(Z\mapsto Z^{s_a}\) identifies every factor
in \(\mathcal I^{\rm mov}(a)\) with its
frozen-centre counterpart except its free
colour-ball indicator. In frozen coordinates
that indicator is
\begin{equation}\label{fcw:moved-cutoff-map}
 \prod_{b\in\mathcal F}
       \mathbf1_{\{|F_{a,b}(B)|\le b_0\}},\qquad
 F_{a,b}(B)=\frac1{g_0i}
               \log(e^{ig_0B_b}e^{-iC_a(b)}).
\end{equation}
Here \(\mathcal F\) is the actual set of free
middle-pivot coordinates. At \(|B_b|=b_0\),
with \(n_b=B_b/b_0\),
\par\noindent
\begin{minipage}{\linewidth}
\begin{equation}\label{fcw:radial-velocity}
                 n_b\cdot D_aF_{a,b}(B)|_1[e]
                                  =-g_0^{-1}n_b\cdot c_{e,b}.
\end{equation}
\popQED\par
\end{minipage}
\end{lemma}
\begin{proof}
The endpoint law for a node-fixed old gauge is
\(T_{y,x}(Z^s)=T_{y,x}(Z)s(x)^{-1}\).
Consequently
\[
 T_{y,x}(Z^{s_a})T_0(x)^{-1}
                         =T_{y,x}(Z)T_a(x)^{-1}.
\]
This proves \(G_{V_0}(Z^{s_a})=G_{V_a}(Z)\),
including the same numerical contour normalization
and the removed tame tests. Both the action
\(S_j^a\) and the complete old interaction are
gauge invariant on the real admitted overlap.
Haar measure and its coarea disintegration on
\(M(Z)=W\) are unchanged, because the gauge is
one at the output nodes. Thus no graph-Jacobian
or coarea term is left over. Equivalently, the
two graph densities transform as the pullbacks
of this same intrinsic measure.

The joint graph identity sends the old free
angle \(u_b\) to \(\operatorname{Ad}(s_a(b_-))u_b\)
relative to \(\widetilde V_a\). This is a real
orthogonal colour change with determinant one,
so it preserves each product ball. A field in
the frozen graph has free link
\(e^{ig_0B_b}V_0(b)\). Its relative free angle
about \(\widetilde V_a\) is exactly the logarithm
in \eqref{fcw:moved-cutoff-map}. Pivot coordinates
follow by the uniqueness of the actual graph.

For the last assertion differentiate
\(i^{-1}\log(e^{iX}e^{-itY})\) at \(t=0\).
Its derivative is \(-Y\) plus terms containing
commutators with \(X\). Pairing with \(X\)
kills every such term by invariance of the
colour inner product. The convergent logarithm
series therefore gives
\(\langle X,D_ti^{-1}\log(e^{iX}e^{-itY})|_0\rangle
=-\langle X,Y\rangle\).
Take \(X=g_0B_b\), \(Y=c_{e,b}\), and divide
the paired identity by \(g_0^2b_0\). This proves the exact normal
velocity, with no polynomial approximation
to the group chart.
\end{proof}

The centre jet has a compatible local analytic
family with
\begin{equation}\label{fcw:centre-jet-bound}
                       \|\mathbf c_e\|_{R,\kappa_+,
                             \gamma_{{\rm per},+}}\le C|e|,
                         \qquad c_{e,b}(1)=0.
\end{equation}
Indeed the weighted literal intermediate fields
have their all-degree connected bounds on a
larger weight disc. Cauchy gives the centre
derivative. The finite \(L\)-step contours in
\eqref{fcw:aligning-gauge} attach fixed local
stars and point products. Their constants can
depend on the fixed \(L,M,R\), but not on the
inner depth. The identity at \(W=1\) follows
because every weighted trivial background is
the identity. This construction uses only
the background matrices. It never evaluates
\(h\) on a complex gauge transform outside
its old ball.

Define a radial surface functional at one free
bond by integrating its three colour coordinates
over \(S_{b_0}^2\), with ordinary area measure,
and all other coordinates over their retained
balls. In the Gaussian pullback of LCF put
\begin{equation}\label{fcw:boundary-moment}
 \mathbf m_b(W;x,h)=
 \frac{\displaystyle
    \int_{|B_b|=b_0}\!n_b
       \prod_{d\ne b}\chi_{b_0}(B_d)
       e^{P_{g_0,\tau,h}(V_0,B)}
                           d\mu_{K_\tau(V_0)^{-1}}(B)}
 {\displaystyle Z_j(g_0;V_0,\tau,h)}.
\end{equation}
The numerator notation means the surface
functional just defined, not a product of
an ambiguous boundary value of \(\chi_{b_0}\)
with a delta distribution. The full coarea
and Haar densities are included in \(P\).
The scalar boundary source is
\begin{equation}\label{fcw:boundary-source}
                              \beta_{e,b}=g_0^{-1}c_{e,b}\cdot\mathbf m_b.
\end{equation}

\begin{proposition}[Exact total first-jet comparison]
\label{prop:fcw-total-jet-comparison}
For real admitted data,
\begin{equation}\label{fcw:total-jet-comparison}
 D_a\log\mathcal I^{\rm mov}(1)[e]
       -D_a\log\mathcal I^{\rm fr}(1)[e]
                              =\sum_{b\in\mathcal F}\beta_{e,b}.
\end{equation}
The right side has the local analytic continuation
and bounds proved below. This continuation does
not differentiate a sharp indicator at a complex
weight.
\end{proposition}
\begin{proof}
Use \cref{lem:fcw-exact-cutoff-change} to write
both integrals in the frozen fibre measure.
Their action derivative is identical, and
their old interaction, relative gauge and
measure are the same. The product of colour
balls has first shape derivative
\[
 \sum_{b\in\mathcal F}g_0^{-1}
      (n_b\cdot c_{e,b})\,
       \delta(|B_b|-b_0)\prod_{d\ne b}\chi_{b_0}(B_d).
\]
Intersections of two sphere boundaries have
codimension at least two and contribute no
first shape derivative. The formula follows
by the one-variable thin-shell change of
variables in each three-dimensional colour
group. The finite graph and strict chart
margin justify differentiation under these
compact integrals. Divide by their common
value \(\mathcal I_0\). The exact LCF Gaussian
factorization cancels its scalar prefactors,
giving \eqref{fcw:boundary-moment} and
\eqref{fcw:total-jet-comparison}.
\end{proof}

\subsubsection{A normalized rooted sphere estimate}

\begin{lemma}[One surface mark in the actual local expansion]
\label{lem:fcw-marked-sphere-bound}
There are fixed \(c,C,K>0\) such that the
complete sources \eqref{fcw:boundary-source}
have compatible holomorphic local families,
rooted at their marked free bonds, with
\begin{equation}\label{fcw:rooted-sphere-bound}
          \|\boldsymbol\beta_e\|_{R,\kappa,
                                    \gamma_{\rm per}}
             \le C_{H_*}|e|g_0^{-1}(1+b_0)^K e^{-cb_0^2}.
\end{equation}
The norm includes the original centre-jet
carriers, old-source carriers, coarea terms,
normalization and the period companions.
It is independent of total volume and inner
depth. It holds on the smaller complex
\(\tau,h,W\) domains retained from LCF.
\end{lemma}
\begin{proof}
We verify the additional local functional
estimate and then its normalization.
Use the actual filed Gaussian matrices of
\eqref{haf:actual-decoupled-root}. Use its
finer comparison error \(\varepsilon_D\) and
Gaussian tilt \(\delta_1\), both at most
\(c_*(LM)^{-4}\). Its scalar dimension is
\(d\le C(LM)^4\#Z\). Relative to the positive
real comparison covariance \(D_0\), the determinant
and inverse estimates give
\[
 |d\mu_D(B)|\le e^{C\varepsilon_Dd}
                   e^{C\varepsilon_D|B|^2}\,d\mu_{D_0}(B).
\]
Reweighting this positive Gaussian by
\(e^{(\delta_1+C\varepsilon_D)|B|^2}\)
changes its normalization by at most
\(e^{C(\delta_1+\varepsilon_D)d}\le e^{C'\#Z}\)
and leaves its covariance between fixed
positive constants. The constants are uniform
on the enlarged covariance-parameter discs
and local complex background domains.
Replacing the three-dimensional integral
at \(b\) by its sphere now uses the positive
Gaussian's three-coordinate marginal and gives
\begin{equation}\label{fcw:local-surface-estimate}
 \int_{|B_b|=b_0}\!|n_b|
           e^{\delta_1|B|^2}|d\mu_D(B)|
              \le e^{C\#Z}(1+b_0)^2e^{-c_2b_0^2}.
\end{equation}
The factor is the area of a three-colour
sphere. Remaining Gaussian moments and any
fixed local parameter derivatives give
only further polynomial factors, or a
fixed exponential per wall cube. They do
not introduce Gaussian coordinates from
the inner fine hierarchy.

The integration set of this marked local
functional is fixed. Gaussian densities,
graph vertices and source functions are
holomorphic in their actual parameters
and have the same tilted domination.
Thus the marked functional is holomorphic,
and the usual covariance-parameter Cauchy
estimates apply. Its full guard includes
the marked coordinate, every potential
support, the centre-jet coefficient carrier
and all background read collars. A wall
factorization is made only after these
complete guards are included. If separated
components remain, the marked functional
is the marked root component times the
ordinary functionals of the others.

The actual action/coarea interaction and
bounded old source have the small budgets
of \cref{lem:lcf-old-source-budget,lem:lcf-large-ball-family}
at the physical \(g_0\). Their existing
connected expansion uses the tilted Gaussian
bound, parameter holomorphy and this wall
factorization. The estimate
\eqref{fcw:local-surface-estimate} verifies
these same requirements with one marked
functional. It supplies an additional
root factor
\(\eta_\partial=Cg_0^{-1}(1+b_0)^Ke^{-c_2b_0^2}\).
The centre coefficient is expanded first
using \eqref{fcw:centre-jet-bound}; its
entire carrier is attached to that root.
The reserved intrinsic rate pays its
fixed guard enlargement. The sum over old
free bonds in an output cube is at most
\(C(LM)^4\), a fixed local factor.

For completeness normalization is obtained
by a single formal source, not a global
lower bound for an integral. Give the
boundary functional a parameter \(\zeta\)
with \(\zeta^2=0\). The finite partition
expression is \(Z+\zeta Z_\partial\).
The coefficient of \(\zeta\) in its
connected logarithm is exactly
\(Z_\partial/Z\). In the actual polymer
expansion every such connected term has
one marked object and otherwise ordinary
objects. At the marked coordinate the
surface functional replaces its physical
cutoff: an ordinary cutoff factor at that
same coordinate is excluded. Larger nested
auxiliary cutoffs equal one on the marked
sphere by the retained admission margin.
Thus no product of a radial delta with
an unspecified boundary value of its own
indicator is formed. This exclusion does
not enlarge any absolute activity bound.
The ordinary small activity
majorant is \(\varepsilon\), within the
retained cluster ceiling. Its rooted
tree sum is bounded by
\[
                  C\eta_\partial
                         \sum_{n\ge0}(C\varepsilon)^n
                                    \le C'\eta_\partial.
\]
All original source origins are combined
on their full carriers before this sum,
exactly as in LCF. This proves the local
bound for the normalized expectation,
without estimating an extensive
\(1/|Z|\). Since the unmarked logarithm
is the selected LCF branch, the formal
degree-one identity equals the actual
ratio in \eqref{fcw:boundary-moment}.
There is no product of two delta
distributions at one colour coordinate.

The intrinsic output is composed with
the actual literal centre on the large
\(R\)-ball by the same whole-cube Banach
composition of \cref{lem:lcf-large-ball-family}.
It is a bounded operator on the rooted
family norm; its rate cost and the old
anchor incidence count are already paid
by the blocking buffer. The surface
root creates no extra background
coordinates. The old input quotient is
respected because the complete integral
and marked expectation depend only on
the complete \(h_z\); take the infimum
over their compatible representations.

Finally retain lifted carriers and
finite-period labels before projection.
In a companion difference, telescope
the primitive matrices, centre factor,
vertices and old-source families. One
factor carries the existing period
tail or a retained winding carrier.
The marked Gaussian derivative still
obeys \eqref{fcw:local-surface-estimate}:
\(D\log d\mu_D\) is a trace term and
a quadratic form in \(B\), whose moments
add only a fixed polynomial in \(b_0\)
and wall size. A slight decrease of
\(c_2\) pays the former; the spatial
rate reserve pays the latter. Thus the
period interpolation and padding of
\cref{lem:lcf-period-compatibility} retain
\(\gamma_{\rm per}\). Construct all
Gaussian and centre families on one
common large cover, with the full
periodic held guards, before descent.
No comparison of quantum partition
functions on different volumes is made.
These estimates prove the full
compatible norm in
\eqref{fcw:rooted-sphere-bound}.
\end{proof}

Before multiplication by \(c_{e,b}\) and
\(g_0^{-1}\), the same marked construction
gives a compatible local vector family for
\(\mathbf m_b\), with the corresponding
sphere-tail bound and fixed \(W\) jets.
Its exponential moments justify the affine
product-rule probes used below; scalar
gauge invariance is required only after
pairing with the centre jet.

The construction also proves bounds for
any fixed number of \(\tau\) and bounded
old-input derivatives, by Cauchy on their
retained larger domains. Differentiation
in \(x\) adds \(\theta_0^{-1}\) per
derivative. No derivative of the frozen
reference cutoff radius is taken.
Because \(2p>1\), for every fixed \(N\)
the right side of
\eqref{fcw:rooted-sphere-bound} is at
most \(C_{H_*,N}\theta_0^{-N}|e|\)
after increasing the reference floor.

\subsubsection{Complete pointings and physical normalization}

Choose output hats \(\theta_y\) in output
units and define the complete pointing
\begin{equation}\label{fcw:pointed-boundary-source}
           (\mathcal B_{j,e})_y(W)
              =\sum_{b\in\mathcal F}
                    \theta_y(m_b/L)\,\beta_{e,b}(W).
\end{equation}
Here \(m_b\) is the old bond midpoint.
The finite
local anchor count changes only constants.
The hats sum to one, so this pointing
has total \(\sum_b\beta_{e,b}\).

\begin{lemma}[Symmetry and the flat value]
\label{lem:fcw-pointing-symmetry}
The complete functions in
\eqref{fcw:pointed-boundary-source} are
gauge invariant and jointly covariant
under lattice symmetries acting on the
weight direction \(e\), old input \(h\)
and fields. They
vanish at \(W=1\). These assertions
precede local filing.
\end{lemma}
\begin{proof}
Let an output gauge induce old gauges
\(k_a\) on \(V_a\) and \(k_0\) on
\(V_0\). Their output-node values are
the same. The contour endpoint law gives
\(s_a'=k_0s_ak_a^{-1}\). Hence
\(\widetilde V_a'=\widetilde V_a^{k_0}\):
alignment to the frozen contour frame
removes the weight-dependent gauge.
Thus \(C_a(b)\) and its first weight
jet transform by the same frozen
\(\operatorname{Ad}(k_0(b_-))\) as the
free colour variable. The sphere and
its normal transform by this real
orthogonal map. The complete old
interaction, action and coarea factors
are invariant on their admitted overlap.
Each complete scalar surface source
is therefore invariant. Holomorphy
extends this identity on the stated
complex domain.

Translations and reflections use the
fixed middle-pivot identity of
\cref{prop:rcs-cutoff-transfer-covariance}.
For a reversed stored bond the centre
jet and free colour variable have the
same sign and adjoint twist. Joint
axis permutation acts on \(e,h\) and
the weighted orbit simultaneously.
Their sphere pairing and the hats
therefore have the stated covariance.
At \(W=1\), \(c_{e,b}=0\), so the
complete source vanishes exactly.
\end{proof}

\begin{proposition}[A frozen-kernel stress jet with a rooted tail]
\label{prop:fcw-frozen-weight-jet}
The moving weighted logarithmic first jet
minus \(\mathcal B_{j,e}\) is a compatible
pointed representation of the actual
frozen-kernel first jet
\eqref{fcw:frozen-stress-insertion}, after
the same flat-field normalization.
Moreover
\begin{equation}\label{fcw:boundary-family-final}
 \|\mathcal B_{j,e}\|_{R,\kappa,\gamma_{\rm per}}
       \le C_{H_*}|e|\theta_0^{1/2}
                           (1+b_0)^Ke^{-cb_0^2}.
\end{equation}
The same type of bound holds for any
fixed \(x,h\) jet on smaller retained
domains. Constants are uniform in
depth and compatible volume.
\end{proposition}
\begin{proof}
The total identity is
\cref{prop:fcw-total-jet-comparison};
the local family and normalization are
\cref{lem:fcw-marked-sphere-bound}.
The symmetry lemma supplies the complete
function identities required by the
profile space. The flat value of the
error is zero, so vacuum subtraction
does not alter it. Select the frozen
output pointings by subtracting this
rooted family from the known moving
weight-jet pointings. Their total is
exactly the fixed-kernel first jet.
This selection does not assert equality
to every independently chosen per-anchor
filing of a stress insertion.
\end{proof}

\subsubsection{The affine readout and the regular stress equation}

For the present homogeneous literal convention
the boundary source has a further cancellation.
It is not needed for its superpolynomial norm bound.

\begin{corollary}[Zero boundary curvature readouts]
\label{cor:fcw-zero-curvature-readout}
All six quadratic curvature readouts of
\(\mathcal B_{j,e}\) vanish.
\end{corollary}
\begin{proof}
At \(W=1\), the frozen normalized law is
invariant under global colour conjugation,
also for complex \(h\) in the complete
gauge invariant input class. Its vector
surface moment in
\eqref{fcw:boundary-moment} consequently
vanishes. We already know \(c_{e,b}(1)=0\).

The exact affine response of
\eqref{ach:affine-response} is independent
of the constant weight modulo a gradient.
Both literal means are the same prescribed
affine datum, so that gradient can be
chosen zero at the final output nodes:
its node values are constant and the
flat constant can be subtracted.
Taking the actual partial mean gives
an old-site gauge equal to one at the
outer nodes. Contour alignment in
\eqref{fcw:aligning-gauge} is unique on
such an orbit. It therefore removes
this gradient and gives
\[
            D_Wc_{e,b}(1)[B^F]=0
\]
for every affine curvature probe \(B^F\).
These polynomial probes are interpreted
in the exponentially summable infinite
companion kernels, as in the affine
response proof.

Differentiate the exact product
\(\beta_{e,b}=g_0^{-1}c_{e,b}\cdot\mathbf m_b\)
twice in two affine directions. Its
constant factors vanish at the flat
field, and each remaining mixed product
contains one of the displayed zero
derivatives. Thus its entire quadratic
affine form is zero. The absolute local
family bound permits the hat and
companion sums. Complete gauge invariance
and the Ward readout then give all six
stated zero coefficients.
\end{proof}

In a response \(-tE_{j,e}+v\) of the
fixed isotropic kernel, the weight
direction is \(a=1+(t/x)e\). Hence the
actual output jet is the moving-weight
one minus \((t/x)\mathcal B_{j,e}\).
Using the selected pointings above,
the stress recurrence of
\eqref{rds:actual-linear-response} has
the same scalar leading coefficient;
its profile forcing is changed from
\(A_j[e]\) to \(A_j[e]-\mathcal B_{j,e}\).
The change is in the six-zero-readout
space and is superpolynomially smaller
than the permitted formation error.
Even without the affine cancellation,
its readouts are bounded by
\eqref{fcw:boundary-family-final} and
give only a superpolynomial scalar
correction after division by \(x\).
The coupling-weighted Volterra estimates
therefore retain their bounds.

This comparison proves only the first
constant-weight jet of the stated
homogeneous regular kernel with bounded
old input. It neither assigns holomorphy
to a sharp indicator nor establishes
small formation on a fixed-centre
weight disc. The frame, regular fibre,
cutoff and normalization in the
frozen response remain the isotropic ones.

% END LAYERED EXPANSION formation/frozen_centre_weight_response.tex


% BEGIN MODULE chapters/05_large_field.tex
\subsection{Large-field histories and induction}
\label{candidate:large_field}
\subsubsection{Large-field regions, their weights, and induction}

Convergence of the regular expansion in \cref{ch:local-formation-program} requires that the plaquette variables of the backgrounds lie within $\varepsilon_k\eta_k^2$ of $\mathbf{1}$, with $\varepsilon_k=g_kp_0(g_k)$, and that the fluctuation fields stay below $\delta_k=(A_1/A_0)\varepsilon_k$. The two thresholds carry the same power of the logarithm and differ by a small fixed ratio. Nothing guarantees these conditions. For this reason each step inserts decompositions of unity $1=\sum\prod\chi\prod\chi^c$ built from local tests. At a place where a test fails, no expansion is made and the fields there are stored. The weights present at that place then supply a factor $e^{-c\,x_k\varepsilon_k^2}=e^{-c\,p_0(g_k)^2}$. For a failed plaquette test this weight is the Wilson weight. When an approximate fluctuation test fails, the weights present are the Wilson weight together with the kept gauge-fixing exponents. When a fluctuation test, a collar test or a rim test fails, the weight is only the Gaussian form of the stored fluctuation field. (The failure of one of the sharper tests of the large-field operation $\mathbf{R}$ of \cref{sec:large-field-stored-variables-and-integration-in-a-large-field-patch} yields only $e^{-p_0(g_k)}$: 2a, 2c.) In this way the density becomes a sum over histories $h$, collected according to their labels $\mathcal{S}=(\Omega_j,\Lambda_j,S_j)_{j\le k}$, which consist of the regions of all levels. Beside the regions, a label carries the marks left by the operations $\mathbf{R}$, one for each island and for each drop-out that stands \cref{sec:large-field-the-tests-the-nested-regions-and-the-content-of-a}; the labels are thus labels with marks. The exponential form given below holds for sums of histories and not for an individual history (2a, 2d).

\begin{proposition}[{inductive description; ultraviolet stability}]\label{prop:large-field-2-1}
Assume the primitive hypotheses \(\mathcal H_{\mathrm{step}}\) of
\hyperref[rg:primitive-closure-premises]{the joint step}, on the actual
histories and domains used below. Suppose that the constants are chosen
in the prescribed order, beginning with the exponents and a large $L$ and ending with a small $\gamma$ \cref{sec:large-field-the-induction-closed-the-order-of-the-constants-the-sum}. Suppose also that the torus is admissible up to level $k$, and up to $K$ for the stability bounds; that is, $\ell L^{K-j}\ge2L^{11}M\check{R}_j$ for $j\le k$, respectively for $j\le K$. This is a lower bound on $\ell$, which is chosen after $M$ and $g$. When $\ell\in\mathfrak{L}(g)$ and the datum is a first-passage datum, the condition follows from the upper flow bound up to $K$ \cref{sec:flow-a-common-companion-and-the-torus-floor}. Assume further that $g_j\le\gamma$ and $\lambda_\sharp(j-i)-c_5\le x_i-x_j\le B_\sharp(j-i)$ for $i\le j\le k$ (up to $K$ for the stability bounds). When $g$ is small, \cref{sec:flow-the-run-its-coefficient-and-the-analytic-requirements} provides these conditions for any first-passage datum; they are conditions on how the couplings flow: $\beta$ has no upper bound, and $K=K(\beta,g)$ is determined together with it. Then $\rho_k=(\mathbf{R}\mathbf{T})^k\rho_0$ is defined, and

\begin{gather*}
\FitMath{\rho_k(V_k)=\sum_{\mathcal{S}}\mathbf{T}'_k(\mathcal{S})\big[\chi_{k,\mathcal{S}}\exp A_{k,\mathcal{S}}(U_k(\mathbf{V}))\big],} \\[3pt]
\FitMath{A_{k,\mathcal{S}}=-A(x_k(\cdot),U)+\mathbf{E}_k+\mathbf{R}_k+\mathbf{B}_k-\mathcal{E}_{k,\mathcal{S}},}
\end{gather*}


\begin{gather*}
\FitMath{\|E^{(j)}\|_{j,\kappa}\le E_0,} \\[3pt]
\FitMath{\|R^{(j)}\|_{j,\kappa}\le g_j^{\kappa_0},} \\[3pt]
\FitMath{|B^{(j)}(X,\cdot)|\le\mathsf{B}\,e^{-\kappa d^\natural_j(\mathcal{S};X)},} \\[3pt]
\FitMath{\mathsf{M}(i,X)\le e^{-(2-c')p_0(g_i)},}
\end{gather*}

where $-\mathcal{E}_{k,\mathcal{S}}$ is the constant term of the action of the label and $\mathsf{B}$ is a constant fixed before $\gamma$ (2d, 2e); $d^\natural_j$ is the hull length of \cref{candidate:large_field}, and $\mathsf{M}(i,X)$ is the mass of a component $X$ that, at some level $i\le k$, has passed all those tests of the operation $\mathbf{R}$ that come before the last one. Such a component is either healed at that level or an island, which remains in the label. The mass is defined as follows. Let $e^{\Phi_\omega}$ be the integrand of the operation $\mathbf{T}'_i(X)$ belonging to the component. This integrand contains the Wilson factor of the core (2b, 2c), and its exponent has the form $\Phi_\omega(\xi;\mathbf{p})=\Phi^\circ_\omega(\xi)+\sum_f\Phi_{\omega,f}(\xi;\mathbf{p})$, which is Borel in $\xi$, analytic in $\mathbf{p}$ and real at real pairs; each factor $\Phi_{\omega,f}$ depends on $\xi$ and on $\mathbf{p}$ only on a support $R_f\subset X^\sim$ \cref{candidate:large_field}. Let $\omega$ run over the local histories of $X$, each with a frame in the hypercubic group $W_4$. Let $\xi$ denote the stored variables and $\mu_\omega$ their measure. Let $\mathbf{p}=(\mathbf{U},\mathbf{J})$ denote the complex pairs of the space $U^c_i(X^\sim,\alpha_{0,i},\alpha_{1,i})$, which contains the pairs at which a term evaluates the factor. Let $\mathsf{v}(\omega)$ be the space-time volume of the patch, which counts $\log x_j$ for each stored site of level $j$ \cref{sec:large-field-rare-patches-gains-per-creation-and-costs-per-site}. Then

\begin{gather*}
\FitMath{\mathsf{M}(i,X):=\frac1{|W_4|}\sum_\omega e^{c\,\mathsf{v}(\omega)}\int\sup_{\mathbf{p}}\big|e^{\Phi_\omega(\xi;\mathbf{p})}\big|\,\mu_\omega(d\xi).}
\end{gather*}

The operation of a drop-out contains no Wilson factor of its own. For a drop-out the same bound holds in a modified form, with two changes. First, the Wilson weight over $X$ of the carrying term is included with the operation; the carrying term is a term of some level $k\ge i$ in which $X$ stands. Second, the supremum is taken over the evaluations at $\xi$. These are the real configurations that are compatible with $\xi$, regular layer by layer, and at which every test listed by $\omega$, with the exception of the test whose failure made $X$ a drop-out, takes its recorded value (the full list of these conditions is given in \cref{sec:large-field-rare-patches-gains-per-creation-and-costs-per-site}). This form follows from the induction and is not one of its clauses. The mass bound holds for every $c\ge0$ in the weight $e^{c\,\mathsf{v}}$ and for every small $c'>0$ in its exponent, provided the ceiling on the couplings is chosen after both $c$ and $c'$; the reserve constant $c_{\mathrm{r}}$ of \cref{sec:large-field-the-induction-closed-the-order-of-the-constants-the-sum} is one such $c$. Apart from the core factors of the islands in $Z$ \cref{sec:large-field-stored-variables-and-integration-in-a-large-field-patch}, the operation $\mathbf{T}_k(Z)$ contains no Wilson weight. The factor coming from a failed plaquette test sits in $e^{-A(x_k(\cdot),U_k)}$, which is the function the operation acts on. In the sum over the labels, every component $C$ of $Z_k$ keeps a factor $e^{-c\,p_0(g_k)-\kappa_1|C|_k}$. The constant $\kappa_1>0$ outweighs the number $C^n$ of connected sets of $n$ cells that contain a given cell \cref{sec:large-field-rare-patches-gains-per-creation-and-costs-per-site}. Each term is a non-negative measure. Neither $E^{(j)}$ nor $R^{(j)}$ carries a label index. On non-wrapping domains with matching step data they are the values, at the lifts of the domains in those units, of a single family on $\mathbb{Z}^4$. (The constant $\mathcal{E}_{k,\mathcal{S}}$, by contrast, depends on the label and on the torus.) Finally, let $\chi_k=\chi_k(T_\ell)$ denote the product, taken over every cube of the torus, of the tests $\chi_k(\square)$. Then, uniformly, there is a constant $E_-$ such that

\begin{gather*}
\FitMath{\chi_k\,e^{-x_kA(U_k(V_k))-E_-|\mathbb{T}^{(k)}|}\le\bar\rho_k,} \\[3pt]
\FitMath{\int\bar\rho_k\,dV_k\le e^{C(1+\log x_k)|\mathbb{T}^{(k)}|},} \\[3pt]
\FitMath{\int\rho_k\,dV_k=Z_{K,\ell,\beta}.}\qedhere
\end{gather*}


\end{proposition}

The theorem is the case $\theta=x(\beta)$ of an induction carried out along an arbitrary frame $\theta$ (the frame setting of \cref{candidate:flow}). That induction runs on all admissible tori simultaneously. The tube is a clause of the class $\mathfrak{C}_k(\theta)$ of \cref{sec:large-field-the-induction-closed-the-order-of-the-constants-the-sum}, and the step by itself does not yield this clause again (\cref{prop:large-field-2-11}). For complex field arguments, each term depends on a triple $(U,A',\mathbf{J})$, in which $U$ is $G$-valued and $A',\mathbf{J}$ are complex. For a term of level $j$, the pair of a triple is $\mathbf{p}=(\mathbf{U},\mathbf{J})$ with $\mathbf{U}=e^{i\eta_jA'}U$, and the triple is called a splitting of its pair. The domains of analyticity are defined by conditions on triples. One of the statements propagated by the induction is that a term takes the same value at any two splittings of one pair. The splitting is therefore a coordinate and not an argument, and a term is a function of the pair. (In the regular setting of \cref{ch:rg-objects-v2} the values coincide on a complex gauge orbit, with a margin left in the radii. Here what is asserted concerns two splittings of a single pair, and it holds on the domain itself.)

The core of the theorem is due to Ba{\l}aban: Comm. Math. Phys. 119 (1988) 243, Theorem 1 and Corollary 3, supplemented by 122 (1989) 175 and 122 (1989) 355, Theorem 1. For the idea of $\mathbf{R}$, he cites Gallavotti and co-workers as the closest earlier work. He states, for one torus, that neither the regular terms nor the part of the remainders produced by the fluctuation integral depend on the region $\Lambda_{k+1}$ of the step (119, pp. 278--279). He also names a second choice of the thresholds and sets it aside; in that choice only the logarithm is evaluated on a bound of the flow known a priori (p. 246). The upper bound he gives is pointwise, $\varrho_k(V_k)\le e^{E_+|T_\eta|}$, stated for his density and his volume ((2.50)), and \cref{prop:large-field-2-1} gives the integrated form in its place. Besides what is quoted here from his papers, this step contains the following:
\begin{itemize}
\item a proof of that independence and a range for it, covering the remainders in their entirety, simultaneously on all admissible tori and cutoffs, by means of one family on $\mathbb{Z}^4$ (for the counterterm, the paper 119, p. 281, uses ``$\mathbf{E}^{(j)}(U_j,z)$ defined on the whole lattice $L^{-j}Z^d$'');
\item the positivity of each term, as a clause of the statement (in 122 (1989) 175, p. 176, Ba{\l}aban announces that the densities $\rho(Z'',V)$ are positive);
\item measure and formula as two separate layers;
\item the class taken as a set of presentations along a frame;
\item the flow inequalities, hypotheses of \cref{prop:large-field-2-1}, obtained as conclusions of the induction, carried out jointly with the flow \cref{sec:flow-the-run-its-coefficient-and-the-analytic-requirements};
\item covariance of the $R$-terms under translations only, term by term, and covariance of their totals by way of measures;
\item the mass, with the complex supremum inside the integral and with a reserve weight;
\item the hull length;
\item the three possible fates of a component, two of which are recorded as marks of the label (in 122 (1989) 175, pp. 181--183, Ba{\l}aban removes from the region on which $\mathbf{R}$ acts any component in which a test of some stage fails, and in 122 (1989) 355, (1.102), he retains the operation $\mathbf{T}_k(Y)$ where the last test fails), and the two sequences of regions attached to a label;
\item an order of the constants that contains no circle;
\item exactness for each partition of the histories.
\end{itemize}

The route also makes use of the following:
\begin{itemize}
\item the collar and the rim of \cref{sec:large-field-the-tests-the-nested-regions-and-the-content-of-a}, which ensure that in the interior the mean of the conditional Gaussian measure stays under the thresholds set by the tests;
\item the factor depending on two levels in the sum over the scales of the $R$-terms \cref{sec:large-field-the-induction-closed-the-order-of-the-constants-the-sum};
\item the terms $V(Y)$, which are not of small amplitude, with the consequent removal of the empty family before the exponentiation \cref{candidate:large_field};
\item in \cref{sec:large-field-rare-patches-gains-per-creation-and-costs-per-site}, a single count over the genealogy of a component, which replaces the double induction of 122 (1989) 355, pp. 384--387.
\end{itemize}

In that induction, the case in which two regions join rests on a sentence on p. 387. That sentence states that after $n_1$ further hull operations, where $n_1$ is a fairly small number, for instance $n_1<10$, the domain obtained from the region whose wait is the longer of the two contains the domain obtained from the other region. The count never compares two such regions. A joining of regions costs it no separate payment, because the count charges the joining to a created cell \cref{sec:large-field-rare-patches-gains-per-creation-and-costs-per-site}. The later steps rely on the theorem and on its sourced form \cref{sec:sources-the-curvature-source-with-fixed-reference-data}, which contains the insertion formula.

\paragraph{The tests, the nested regions, and the content of a history}
\label{sec:large-field-the-tests-the-nested-regions-and-the-content-of-a}

\textbf{The tests of step $t=k+1$} are applied to a term of $\rho_k$ whose live region is $Y=\mathcal{Z}_k$. First, outside a margin of $Y$, the new field $V_t$ undergoes the plaquette test $\chi_t(\square)$, and the cubes where this test fails make up $P_t$. The factor $\mathfrak{G}_k$ is inserted away from $P_t\cup Y$. By the Faddeev--Popov argument, its integral over the gauge group of the blocks equals $1$. Next comes the approximate fluctuation test $\chi'_k(\square)=\mathbf{1}\{|V_k(\mathcal{M}^kU_{t,\square}(V_t))^{-1}-\mathbf{1}|<2\delta_k\}$. This test is evaluated at the bonds of $\square^{\sim2}$ other than the pivot bonds \cref{ch:local-formation-program}, and the third test below is evaluated at the same bonds. Here one takes from the orbit $U_{t,\square}(V_t)$ an element whose $k$-th average lies in the block axial gauge, and every such element gives the same tested number. The cubes where this test fails make up $Q_t$. These sets determine $\Omega_t$, the region on which $U_t(\mathbf{V})$ and $\mathcal{A}_k$ exist. The backgrounds and the propagators on sequences of regions are taken from \cref{ch:multilevel-backgrounds}, \cref{sec:propagators-propagators-in-a-background-field-exponential-decay-produced-by-the}, together with their regularity and their decay in the local scale. A label's sequence starts at $\Omega_0=\Omega_1^\sim$, the first region enlarged by one layer of level-$1$ cells, and not at $T_\ell$; this is required by \cref{prop:source-domains-4-13,lem:source-domains-4-15} in \cref{candidate:sources,ch:all-state-source-transfer}. After that the test $\chi^{(k)}(\square)$ is made, and its failures make up $R_t$. Hulls ($'$, $\widehat{\cdot}$) and layers are taken in level-$t$ cells.
At \(t=1\), \(P_0'\) is the level-one hull of failures of the
bare density test. Its indicator on a level-zero density
cube \(\square\) requires
\(\sup_{p\subset\square^\sim}|\partial U(p)-\mathbf1|<\varepsilon_0\).
Use the indicator and its complement in the positive
partition by \(P_0\). The first plaquette tests avoid \(P_0'^{\sim1}\)
and the first gauge draw avoids \(P_1'^{\sim1}\cup P_0'^{\sim2}\). The boxing map $\mathsf{bx}$ replaces each component contained in a cube of $O(1)$ cells by its bounding box, and repeats this until nothing changes. With these conventions,

\begin{gather*}
\FitMath{\Omega_t^c=Q_t'^{\sim2}\cup P_t'^{\sim4}\cup H_t,\qquad
 H_t=\begin{cases}\widehat{Y}^{\sim8},&t\ge2,\\
                   P_0'^{\sim5},&t=1,\end{cases}} \\[3pt]
\FitMath{Z^{[0]}=\mathsf{bx}\big((\Omega_t^c)^{\sim2}\cup R_t'^{\sim}\big),} \\[3pt]
\FitMath{Z^{[n+1]}=\mathsf{bx}\big(Z^{[n]}\cup(\text{cells of the failures of round }n)^{\sim}\big),}
\end{gather*}

where the rounds are those of the recursion described below. The set $Z_t=\Lambda_t^c$ is the final $Z^{[n]}$, and $S_t=(\Omega_t\setminus\Omega_t^{\sim-1})\cup R_t$. The numbers of layers are those of Ba{\l}aban. If the boxing map is omitted, $Z^{[0]}$ coincides with his $\Lambda_t^c$ in the form that the display (3.20) gives it, that is, before he replaces the small components by rectangular parallelepipeds (Comm. Math. Phys. 119 (1988) 243, (3.16), (3.20), p. 269). The order of these numbers matters: an older set receives more layers, and each move is made one layer inside the region of the move before it. The creation set $\mathcal{W}^\wedge_t$ is the union of the hulls of every failure of the step: $P_t'$, $Q_t'$, $R_t'$, the cells of the collar bonds that fail, and the rim cells that fail.
At \(t=1\) it also includes \(P_0'\).
If nothing is created, the step reduces to the free evolution

\begin{gather*}
\FitMath{\Phi_t(Y)=\mathsf{bx}\big(\widehat{Y}^{\sim10}\big);} \\[3pt]
\FitMath{\text{in general}} \\[3pt]
\FitMath{Z_t\subset\widehat{Y}^{\sim10}\cup(\mathcal{W}^\wedge_t)^{\sim O(1)}} \\[3pt]
\FitMath{\text{up to boxing}}
\end{gather*}

That is, one coarsens, adds a margin and boxes; this is the $S(Z)$ of \cref{sec:large-field-rare-patches-gains-per-creation-and-costs-per-site}. The size $\check{R}_k$ is the smallest power of $L$ that is not less than $(\log\min_{i\le k}x_i)^r$, with $r\ge2$. Because it is a power of $L$, all partitions are nested. Because it is formed from the running minimum, it is non-increasing (Ba{\l}aban forms it with $g_k$ itself), and therefore $t\mapsto t-\mathsf{N}_t$ is strictly increasing. Moreover $\check{R}_k\le L\check{R}_{k+1}$, and two consecutive drops never occur, because $x$ changes by $O(1)$ per step and $\gamma$ is small. A test cube $\square$ has side of order $M_2\check{R}_t$. Here $M_2$ is a power of $L$ with $L^2M_2\le M$ and $c\,\delta_0M_2\ge1$. The last condition is a floor on $M_2$, and $\delta_0$ is fixed together with $L$, which comes earlier. Hence a test cube lies well inside a cell of its own level.

\textbf{What factorises, and why.} The tests $\chi_t$ and $\chi'_k$ are built from localised backgrounds. They depend on $V_t,V_k$ only within $O(1)$ test cubes. The third test, and also the collar and rim tests, depend on $\mathcal{A}_k=\tfrac1i\log(V_k(V^{\mathrm{bg}}_k)^{-1})$, which is formed with the true averaged background on $\Omega_t$. These tests are therefore local in the new variable $\mathcal{A}_k$, but not in $(V_k,V_t)$. Once $\mathcal{A}_k$ has become the variable, the surviving $\chi'_k$ (on $S_t$) compare $e^{i\mathcal{A}_k}V^{\mathrm{bg}}_k$ with $\mathcal{M}^kU_{t,\square}$. They therefore depend on $V_t$ over the whole of $\Lambda_t$. Four facts together ensure that everything attached to $Z_t$ factorises over its components:
\begin{itemize}
\item The backgrounds are localised.
\item The substitution on the strip \cref{sec:large-field-stored-variables-and-integration-in-a-large-field-patch} is made. After it, every indicator kept in $\Omega_t\cap\Lambda_t^c$ depends on $V_t$ in $\Lambda_t^c$ alone, apart from the factors that equal $1$ on the supports.
\item The conditional quadratic form is cut \cref{sec:large-field-stored-variables-and-integration-in-a-large-field-patch}.
\item The pivot values are dealt with as follows. An old test $\chi_k(\square)$ of a cube $\square\subset\Omega_t^c$ adjacent to $\Omega_t$ also depends on $V_k$ at pivot bonds \cref{ch:local-formation-program} inside $\Omega_t$. There $V_k$ is not a stored variable, and through $e^{i\mathcal{A}_k}V^{\mathrm{bg}}_k$ a test of one component depends on the fields stored in another. At such a bond, $V_k$ is given the value that solves, for the pivot bond, the averaging constraint $\mathcal{M}(V_k)(c)=V_t(c)$ of the single coarse bond $c$. Ba{\l}aban (122 (1989) 175, (1.3)) retains only those old tests whose cubes lie outside $\Omega_t^{\sim}$; the tests that he leaves out equal $1$ on the supports by \cref{lem:large-field-2-2} (a), whose margin is this layer.
\end{itemize}

Inside $\mathbf{R}$, the tests are built from backgrounds localised in the component, never from the background of the torus. The induction carries the following property: each indicator stored in $\mathbf{T}_k(Z)$ depends only on those field variables that are localised in $Z$.

\textbf{Totality.} By definition, a test is the indicator function of the set where all objects occurring in it are defined and where its inequalities hold. These objects are the localised minimal orbit, which must exist and be a single orbit, the contour variables, and the logarithms. Consequently $\chi^c=1$ wherever some object is undefined. A factor of $\mathfrak{G}_k$ equals $1$ where $V(y,x)$ is undefined, by a separate convention for that factor. Hence $\int\mathfrak{G}_k(V^u)\,du=1$ for every $V$, and each insertion $1=\sum\prod\chi\prod\chi^c$ is an identity that needs no assumption. The tests are not empty conditions. The averaging constraint transports smallness: $|\partial V_k-\mathbf{1}|<2\varepsilon_k$ implies $|\partial V_t-\mathbf{1}|<CL^2\varepsilon_k$. Therefore $U_{t,\square}(V_t)$ exists on the carried supports, and the new test further restricts a field that is regular already (Comm. Math. Phys. 119 (1988) 243, p. 246 at the first step, p. 264 at a general step). The cost is that $\chi^c$ contains configurations that are not large in any quantitative sense. These configurations must also produce the small factor \cref{sec:large-field-rare-patches-gains-per-creation-and-costs-per-site}.

For the convention just stated, normalization can be checked one site
at a time. Let $D_{y,x}$ be the set where the prescribed mean
$V(y,x)$ exists. Uniqueness of the mean and its biunitary covariance
show that $D_{y,x}$ is gauge invariant and that, when the block centers
are fixed,
$V^u(y,x)=V(y,x)u(x)^{-1}$ on $D_{y,x}$. On this set the integral of
the normalized site factor over $u(x)$ is $1$ by Haar invariance.
Off $D_{y,x}$ the site factor is identically $1$, so its integral is
again $1$. Each noncentral site belongs to one block, and the Haar
variables $u(x)$ are independent. Integrating the product therefore
gives $\int\mathfrak G_k(V^u)\,du=1$ for every $V$. This uses the
sitewise extension: a defined site retains its prescribed weight even
if a different site's mean is undefined.

\textbf{Collar and rim}: this device is an addition to the step of Comm. Math. Phys. 119 (1988) 243, Sect. 3. The stored part of the fluctuation field moves the mean of the conditional Gaussian measure of \cref{sec:large-field-stored-variables-and-integration-in-a-large-field-patch}. The thresholds $\delta_k$ are stated in the unscaled normalisation, and in that normalisation $\mathcal{A}_k|_{\Lambda_t}$ has the mean

\begin{gather*}
\FitMath{\mathbf{m}=-\,C^{(k)}(\Lambda_t)\,\Delta^{(k)}\,\mathcal{A}_k\big|_{\Omega_t\setminus\Lambda_t}}
\end{gather*}

This map is linear, so in terms of the scaled fields $A$, $A'$ of \cref{sec:large-field-stored-variables-and-integration-in-a-large-field-patch} it is the same map, with each threshold divided by $g_k$; only $|\mathbf{m}|$ enters. This shift must remain well below the least threshold among the inner tests; that threshold is $\theta_c\delta_k$ on bonds that cross $\partial\Lambda_t$, and $\delta_k$ further in. A recursion therefore keeps enlarging the region until, along its boundary, $\mathcal{A}_k$ satisfies two sharper tests:
\begin{itemize}
\item $|\mathcal{A}_k|<\delta_k/(4\Theta_1)$ on the collar, the set of stored bonds whose distance from $\Lambda_t$, measured in level-$k$ units, is at most $\varpi_T$. With a qualifying word added, the same name serves for three further objects. In \cref{sec:large-field-stored-variables-and-integration-in-a-large-field-patch}, the collar of a stage means the shell over which that stage integrates. In \cref{sec:large-field-rare-patches-gains-per-creation-and-costs-per-site}, the waiting collar is the set that keeps $X\cap\Omega_{m_\flat}$ apart from the floor. In \cref{candidate:large_field}, \cref{sec:large-field-the-induction-closed-the-order-of-the-constants-the-sum}, the collar of a component $X$ means the enlarged component $X^\sim$, that is, $X$ together with one layer of $M$-cubes of its level.
\item $|\mathcal{A}_k|<\theta_c\delta_k$ on the rim, the set of cells of level $t$ of $Z_t$ that touch $\Lambda_t$.
\end{itemize}

At greater distance, all that is known is $|\mathcal{A}_k|<C\delta_k$. Every round examines bonds and cubes that no earlier round examined. Hence the recursion terminates, and the admissible outcomes it produces constitute a decomposition of unity. The kernel of $\mathbf{m}$ is bounded, with exponential decay. The contributions of the collar, the rim and the remaining part give

\begin{gather*}
\FitMath{|\mathbf{m}|\le C\big(\Theta_1^{-1}+\theta_ce^{-c\varpi_T}+e^{-cM\check{R}_t}\big)\delta_k<\tfrac12\theta_c\delta_k} \\[3pt]
\FitMath{\text{if }\theta_c\Theta_1\ge C,\,\varpi_T\ge C,\,\gamma\text{ small}.}
\end{gather*}

The rim fraction $\theta_c\in(0,1)$, the collar width $\varpi_T$ and the collar factor $\Theta_1$ serve this purpose.

\begin{lemma}[{the new tests imply the old; the form reproduces itself}]\label{lem:large-field-2-2} (a) On $\Omega_t^\sim$, that is, on $\Omega_t$ enlarged by a margin of one layer of cells, the indicators in $\mathfrak{G}_k$ and the tests $\chi_k(\square)$ that the term carries all equal $1$. This holds on the support of the new tests together with the tests and averaging constraints that the term still carries alongside them. (b) The test $\chi^{(k)}$ is made only on cubes $\square\subset\Omega_t^{\sim-1}$. For such a cube, if $\square\not\subset R_t$, then $\chi^{(k)}(\square)$ implies $\chi'_k(\square)$. Consequently, within $\Omega_t$, the test $\chi'_k$ remains only on $S_t$. (c) After the step, the tests that stand outside the integral operations are exactly $\chi_{t,\mathcal{S}}:=\chi_t(\Omega_t)=\prod_{\square\subset\Omega_t}\chi_t(\square)$. They are formed with the action sequence defined below, and with $\chi^c_t$ on an island less a margin.
\end{lemma}

For (a): up to a gauge transformation, 

\begin{gather*}
\FitMath{U_{k,\square}=\exp\big[i\eta_k\mathbf{H}_{k,\square}\big(\tfrac1i\log[V_k(\mathcal{M}^kU_{t,\square'})^{-1}]\big)\big]U_{t,\square'}}
\end{gather*}

. Here $\mathbf{H}_{k,\square}$ is the localised response $\mathbf{H}_k$ of \cref{ch:multilevel-backgrounds}, so the argument has the form displayed only off a boundary layer of the cube; inside that layer the argument is built from the lift of $V_k$ and is bounded by $C\varepsilon_k$. It follows that

\begin{gather*}
\FitMath{|\partial U_{k,\square}-\mathbf{1}|\le(1+c)\varepsilon_t\eta_t^2+C\big(\delta_k+e^{-c\,\delta_0M_2\check{R}_k}\varepsilon_k\big)\eta_k^2\le\big(CL^{-2}+CA_1/A_0+Ce^{-\check{R}_k}\big)\varepsilon_k\eta_k^2<\varepsilon_k\eta_k^2 .}
\end{gather*}

Three mechanisms appear repeatedly:
\begin{itemize}
\item Curvature scales by the factor $\eta_t^2/\eta_k^2=L^{-2}$.
\item The ratio $\delta_k/\varepsilon_k=A_1/A_0$ is small.
\item The artificial boundary conditions of localised minimisers act only through $e^{-c\delta_0\operatorname{dist}}$ (1b, 1c). On cubes whose side is proportional to $\check{R}_k\approx(\log x_k)^r$, this quantity is smaller than every power of $g_k$, because $r>1$. In the present estimate it needs to beat only constants (Ba{\l}aban, Comm. Math. Phys. 119 (1988) 243, (3.7)); on the strip \cref{sec:large-field-stored-variables-and-integration-in-a-large-field-patch} it must beat a negative power of $g_k$ (ibid., p. 269, below (3.22), where on $S_{k+1}$ the bound $O(1)\varepsilon_k\exp(-R_k)$ is obtained, ``hence by any positive power of $g_k$'').
\end{itemize}

For (b): 

\begin{gather*}
\FitMath{|V^{\mathrm{bg}}_k(\mathcal{M}^kU_{t,\square})^{-1}-\mathbf{1}|\le C\tfrac{A_0}{A_1}e^{-c\check{R}_k}\delta_k<\delta_k}
\end{gather*}

 ((3.17)--(3.19) of the same paper). The factor $2$ between the two thresholds gives exactly the room needed for this. For (c): a test cube of step $t$ serves also as a test cube of the density of level $t$, carrying the same background and the same threshold. The old tests inside $\Omega_t^c$ pass into the integral operation. The same ladder of implications holds inside $\mathbf{R}$ \cref{sec:large-field-stored-variables-and-integration-in-a-large-field-patch}. On the support of the tests of a stage, the following tests equal $1$: those of the preceding stage, and those of the single stored step whose fields the stage integrates (the plaquette tests of its two levels, the gauge indicators, the approximate fluctuation tests). Further, consider any point at which a term is evaluated: the triple $(U,0,J(U))$, where $U$ is the background of the label, at fields where no characteristic function of the term vanishes. At such a point, $\chi_q(\square)(V_q)=1$ holds on every cube $\square$ of level $q$ that lies in the layer $\Omega_{q+1}^\sim\setminus\Omega_{q+1}$ and not in a drop-out. Also at such a point, the current of the bare field, formed from the plaquettes inside $\Omega_0$ only, stays within the bound that the domains of analyticity impose on their current (the size $|\mathbf{J}|$ is smaller than the first radius of the domain, measured in the units of the local scale). This bound holds also at the bonds of $\Omega_0$ that have a plaquette leaving $\Omega_0$.

\begin{proposition}[{geometry}]\label{prop:large-field-2-3} The regions are nested, $T_\ell\supset\Omega_1\supset\Lambda_1\supset\dots\supset\Lambda_k$, with gaps of at least $cM\check{R}_j$ in level-$j$ units. Every large-field region and every creation set is a union of cells of its own level. Both statements hold for the regions that the steps produce (119, p. 256). Each operation $\mathbf{R}$ has stage sequences and prepared sequences (namely the regions that the label has after $m$ of its stages, and after all $\mathsf{N}_k$ of its stages; see \cref{sec:large-field-stored-variables-and-integration-in-a-large-field-patch}), on which the backgrounds of \cref{sec:large-field-stored-variables-and-integration-in-a-large-field-patch}, \cref{candidate:large_field} are formed, and drop-outs keep nests. In these sequences and nests the boxes are only unions of $M$-cubes, and the only lower bound on the gaps is one $M$-cube of the later level:

\begin{gather*}
\FitMath{\operatorname{dist}_\infty(Z_j,\Omega_{j+1})\ge LM,} \\[3pt]
\FitMath{\Omega_{j+1}\subset\Lambda_j^{\sim-1},} \\[3pt]
\FitMath{Z_j^\sim\cap\Omega_{j+2}=\emptyset} \\[3pt]
\FitMath{(\text{level-}j\text{ units; layers of }M\text{-cubes of level }j),}
\end{gather*}

the first inequality is an equality at the active indices
\(k_0(t)+2\le j<t\) of the upper part of a nest prepared at level
\(t\), where \(\Lambda_j=\Omega_j\). Here \(k_0(t)\) is the
change-of-scale index in \eqref{eq:elab-history-prepared-boxes}.
At \(j=k_0(t)+1\) the gap is instead \(w_j-w_{j+1}\) in bare
units, which still gives the displayed lower bound but can be larger; the reason is that a step places eight layers around the old region, and the preparation removes most of them again. Consequently the estimates of \cref{ch:multilevel-backgrounds}, \cref{sec:propagators-propagators-in-a-background-field-exponential-decay-produced-by-the} must hold with a separation of only one $M$-cube. A smallness $e^{-c\check{R}_j}$ from a separation is available only at regions produced by steps. In each decomposition exactly one outcome equals $1$, and hence $\rho_k\,dV_k=\sum_h\tau_h$ holds exactly. A history records every outcome: $P_t$, $Q_t$, $R_t$, the collars and rims, and the outcomes inside $\mathbf{R}$. It determines the label, the creations and the ages. There are two kinds of creation:
\begin{itemize}
\item Creations by step $t$ are the cells $\mathfrak{c}$ of $\mathcal{W}^\wedge_t$, which are $M\check{R}_t$-cubes of level $t$. They are counted one cell at a time, and each carries one factor $e^{-c\,p_0(g_t)^2}$.
\item Creations by an operation are its marked components (next paragraph); the tests of an operation are decompositions of unity component by component. Each marked component counts as one creation and carries one factor $e^{-p_0(g_t)}$, which must pay for $\mathsf{N}_t+1$ levels of one small box.
\end{itemize}

All later statements of locality rest on this.
\end{proposition}

\textbf{Three fates, marked components, two sequences of regions} (Ba{\l}aban declines ``a systematic description of all possible cases'': 122 (1989) 175, p. 183). We set $t_\flat:=t-\mathsf{N}_t$. In each component $X$ of its class \cref{sec:large-field-stored-variables-and-integration-in-a-large-field-patch}, the operation $\mathbf{R}$ of level $t$ prepares a nest $Z''_l(X)$, $t_\flat<l\le t$. The nest consists of increasing boxes in $\Omega_l^c\cap X$. Each box is what remains of the large-field region of scale $l$ after the stage that integrates the stored fields of level $l-1$. A component has one of three fates:
\begin{itemize}
\item \emph{Healed} (every test passed): $\Omega_l,\Lambda_l,S_l\mapsto\Omega_l\cup X,\Lambda_l\cup X,S_l\setminus X$ for $l\le t$, and $X$ leaves the label.
\item \emph{Island}, also called second class (only the last test fails, namely $\chi_t$ of the new field on $X$ less a margin): the same replacement is made for $l<t$ and for $\Omega_t$, $S_t$, but $\Lambda_t$ is left unchanged. Thus $X\subset Z_t$ remains, as a component without inner regions. On it, $\mathbf{T}'_t(X)$ is kept together with its variables, its measure and its factors \cref{candidate:large_field}.
\item \emph{Drop-out} (some test fails after $n\le\mathsf{N}_t$ stages): $X$ remains in $Z_t$ and retains, as its integral operation \cref{sec:large-field-stored-variables-and-integration-in-a-large-field-patch}, the operations carried out up to that point. Moreover

\begin{gather*}
\FitMath{\Omega_l\cap X=X\setminus Z''_l(X)} \\[3pt]
\FitMath{(t_\flat<l\le t_\flat+n),} \\[3pt]
\FitMath{\Lambda_l\cap X=X\setminus Z''_l(X),} \\[3pt]
\FitMath{S_l\cap X=\emptyset} \\[3pt]
\FitMath{(\text{the same }l,\,l<t),} \\[3pt]
\FitMath{\text{all other entries unchanged, so }\Lambda_t\cap X=\emptyset .}
\end{gather*}

(For $n=\mathsf{N}_t$, then, the entry $S_t\cap X$ keeps its previous value; on the support of the tests of the last stage, the tests retained on it equal $1$, by the ladder of implications inside $\mathbf{R}$ that is stated in the argument for \cref{lem:large-field-2-2}.)
\end{itemize}

Islands and drop-outs are the \emph{marked components}. A mark persists until a later healed component or island swallows it. The word also denotes two other objects: the mark $\Theta^{(k)}$ below, which is a weight on histories, and a marked term of a cluster expansion \cref{sec:large-field-the-induction-closed-the-order-of-the-constants-the-sum}. A drop-out's mark keeps the regions $\Omega^\circ_l=\Omega_l\cap X$, and likewise $\Lambda^\circ_l,S^\circ_l$, which $X$ possessed before the operation. For $0\le\boldsymbol{\nu}(X)\le n(X)$ we set

\begin{gather*}
\FitMath{\Omega^{\boldsymbol{\nu}}_l\cap X=X\setminus Z''_l(X)\,\,(t_\flat<l\le t_\flat+\boldsymbol{\nu}(X)),} \\[3pt]
\FitMath{\Omega^{\boldsymbol{\nu}}_l\cap X=\Omega^\circ_l\,\,(t_\flat+\boldsymbol{\nu}(X)<l\le t_\flat+n(X)).}
\end{gather*}

The choice $\boldsymbol{\nu}=n$ gives the \emph{determining sequence}, which describes the field in its present state. It governs the determining sets, $U_k(\mathbf{V})$ and the propagators; the analyticity spaces involve all the $\Omega^{\boldsymbol{\nu}}$. The choice $\boldsymbol{\nu}=0$ gives the \emph{action sequence}, the regions for which the terms were constructed. It governs the index sets of $\mathbf{E}_k,\mathbf{R}_k,\mathbf{B}_k$, the profiles $\varphi_j$ and hence $x_k(\cdot)$, the carried tests, and the constants $\mathcal{E}_{k,\mathcal{S}}$. No single sequence can play both roles. A stored region in a determining position requires variables that have already been integrated. A stage region in the sums re-sorts terms that were made for other regions. Three facts make the device workable. \emph{Marks wait}: a mark of level $s$ lying inside a class component of level $t$ satisfies $s<t_\flat$. Hence standing marks that share a part are nested, their windows $(s_\flat,s]$ are disjoint, and at any one place and level at most one drop-out mark is active. Next, each $\Omega^{\boldsymbol{\nu}}$ is the geometric label obtained by stopping the operations after $\boldsymbol{\nu}$ stages. Its steps, classes and nests coincide with those of the history at hand, and so \cref{prop:large-field-2-3} applies to it: \emph{the later geometry does not depend on how far an interrupted operation had progressed}. Finally, a term created at a moment $\mu$ (following the step, or the operation, of a level $j\le t$) retains permanently the space in which it was made:

\begin{gather*}
\FitMath{\Omega^{\boldsymbol{\nu}_\mu}_l(\mathcal{S}_t)=\Omega_l(\mathcal{S}^\mu)\cup Z_{[j,t]}} \\[3pt]
\FitMath{(l\le j),}
\end{gather*}

where $\mathcal{S}_t$ denotes the level-$t$ label as it stands once the operation $\mathbf{R}$ has been performed, $\boldsymbol{\nu}_\mu=n$ on each mark older than $\mu$ and is $0$ on each later mark, $\mathcal{S}^\mu$ is the label at that moment, and $Z_{[j,t]}$ is the union of all components that were healed or became islands at levels $j,\dots,t$.

\paragraph{Changing stage indices.}
Keep the recorded step data and the three fates fixed, and change
the stage indices on the drop-outs within their allowed ranges.
To form the new top regions of level \(s\), the recursion uses
the recorded step data and the top stored region from level
\(s-1\). Changing the stage index leaves that region fixed. At
each level the class and the nests for the next operation are
also fixed. Indeed, a difference at an earlier index \(l\) in a
candidate component can only be carried by a standing drop-out
of level \(s\) with \(l<s<t\). If \(l\) is in the eligibility
window of level \(t\), then \(s\) is in that window as well.
The drop-out lies in the candidate component by nesting, so that
component is ineligible in both labels. All eligible components
therefore have the same traces through their windows. Their nests
are determined by these traces and the grids.

The dependence on the indices can be checked one cell at a time.
Fix a level \(l\) in a later label and take \(q\) in a common cell refinement of the regions being compared. Look for a
standing drop-out \(X\) of level \(s\) with \(q\subset X\) and
\(s_\flat<l\le s\). There is at most one such drop-out because
standing marks that meet have disjoint windows. If there is none,
membership of \(q\) in the level-\(l\) regions is unchanged.
Otherwise let \(m\) be the replacement stage index of \(X\).
For the \(\Omega\) entry, the membership is given by
\begin{equation}\label{eq:elab-history-reindexed-cell}
 q\subset\Omega_l^{\mathrm{new}}
 \quad\Longleftrightarrow\quad
 \begin{cases}
  q\not\subset Z''_l(X),&l\le s_\flat+m,\\
  q\subset\Omega_l(\mathcal S'_s),&l>s_\flat+m.
 \end{cases}
\end{equation}
For \(l<s\) the same rule holds for \(\Lambda_l\) with
\(\Lambda_l(\mathcal S'_s)\) in the second line. At \(l=s\) the
drop-out keeps \(q\) outside \(\Lambda_s\).

These assertions follow together by induction over the operations.
At a step, the earlier entries and standing marks are copied.
At an operation, a healed component or an island overwrites the
older entries on its region, and any mark it swallows ceases to
stand. Outside the class no entries change. On a drop-out the
new entries through stage \(m\) are the complements of the
corresponding prepared boxes, and the earlier rule supplies all
entries outside its window. The eligibility check above shows at
each induction step that the recorded components and their fates
still define the same choices of regions. This proves
\eqref{eq:elab-history-reindexed-cell}, including the case of
zero stages. It also explains why the action and determining
sequences can use different indices on the same records. In
using this identity for an integral, one must still use the
measures, variables and test factors carried by that integral.

The resulting labels obey the stage version of the nesting
rules. A fully prepared box need not be a component of the
stored region, and at the current operation one uses the
prepared margins in place of the ten-layer margin of an
ordinary step.

\textbf{Locality in the history.} If the large-field regions of two histories never touch, the two histories combine into a single history $h^1\cup h^2$. There are two reasons for this. (i) \emph{A margin beats a reach}: two regions made of level-$t$ cells that do not touch are separated by at least one cell. One cell exceeds the largest range of a bond test, of a pivot bond and of a collar (in the case of the collar, this is the floor $M\ge\varpi_T$ of \cref{sec:large-field-the-induction-closed-the-order-of-the-constants-the-sum}). Taking hulls and adding layers commute with forming unions, and boxing depends on a single component only. A step therefore acts on each component separately. (ii) The class of $\mathbf{R}$ depends on the past: it depends on every object whose creation falls in the window $(t_\flat,t]$, or which carries a mark at a level in $[t_\flat,t)$. Moreover, the recursion is not monotone: a later region may cover a healed place once more. \emph{Window lemma}: no creation within the window of $t$ has been healed. Indeed, healing at a level $i<t$ something created at step $s$ requires $s\le i-\mathsf{N}_i$, and for something marked at level $s$ it requires $s<i-\mathsf{N}_i$. On the other hand $i-\mathsf{N}_i<t-\mathsf{N}_t\le s$, since $\mathsf{N}_t=\check{R}_t$ is non-increasing. Within the window, therefore, the creations and the marks of each history still lie in that history's own region. It follows that each history is the superposition of its clusters, and that clusters may be interchanged between histories. Each cluster has a root. The root is of one of two kinds: a component that remains live at the last level, or a healed component, considered at the level where its healing took place. The cluster is made up of the components of earlier levels that flow into the root, together with everything the history records in them. Each component, at whatever level, lies in the cluster of precisely one root, namely the last of its successors; thus nothing recorded about the healed components is lost.

\textbf{A hard-core gas; the standing part.} Erasing clusters of healed components from a history yields a history with the same label. The \emph{standing part} $h^\flat$ of $h$ is what remains after the clusters of all its healed components are removed. For a string $s$ of outcomes of the next step one has $(h,s)^\flat=(h^\flat,s)^\flat$. Let $\bar h$ be a history of level $k$ in which no component is healed at any level. In the rest of this paragraph each history is taken without the clusters healed below level $k$. It stands for the class of histories that yield it once these clusters are erased. Consider the histories that reduce to $\bar h$ once one erases the clusters belonging to those components that were healed at level $k$. They are in bijection with the pairs $(\mathfrak{X},(\varsigma_X)_{X\in\mathfrak{X}})$. Here $\mathfrak{X}$ is a finite collection of possible healed components which touch neither one another nor $Z_k(\bar h)$. Further, $\varsigma_X\in\mathcal{H}^1(X)$ ranges over the histories of $X$ alone in which $X$ is healed and nothing is healed before it, and the choices for different components are independent. All these histories share one label. An individual $\tau_h$ is a measure produced by the moves, and it has no exponent of the form in \cref{prop:large-field-2-1}. The part of $R^{(k)}$ produced by $\mathbf{R}$ arises only through the sum over $(\mathfrak{X},\varsigma)$, a sum which $\mathbf{R}$ itself performs (under \hyperref[hyp:C1]{the exact-operation hypothesis} and \hyperref[hyp:C2]{the locality hypothesis}). This sum is a polymer sum. The reasons are the gas property and two facts of locality in the history: $\mathbf{T}'_k(X)$ is the same in all the summands containing $X$, and a healed patch leaves no trace outside itself \cref{candidate:large_field}.

\textbf{Traces and charts.} Let $W$ be a region, and let $W^+_i$, the \emph{window} of level $i$, denote the enlargement of this region by a single layer of level-$i$ cells. The trace of a label is

\begin{gather*}
\FitMath{\operatorname{tr}_W\mathcal{S}:=\big(\Omega_i^c\cap W^+_i,\,Z_i\cap W^+_i,\,S_i\cap W^+_i\big)_{i\le k},}
\end{gather*}

with two additions. The first addition is the marks of the label. They enter collectively rather than individually. For every level, every kind of mark and, for drop-outs, every number of stages carried out, the trace records the union of all components carrying such a mark, together with the unions of their stored regions and of their nests. Hence the trace does not show which pieces in $W^+_i$ belong to the same component. The second addition consists of three separate sets of stored bonds of level $i-1$ contained in $W^+_i$: those in the collar, those at a rim cube, and those in the cut-off zone of the strip. The cut-off zone consists of those bonds of the strip that lie near the set where the fluctuation field is large and carry no kept small-field condition of that field; there the substitution of \cref{sec:large-field-stored-variables-and-integration-in-a-large-field-patch} uses the cut-off $\vartheta_{\mathrm{co}}$. Where a bond carries a kept condition, the variables are left unchanged, while on the remaining bonds of the strip the substitution proceeds with no cut-off. The terms depend on these three sets, and the regions in $W^+_i$ do not fix them: to decide if a cube is a rim cube one has to look at its neighbours, and $\Omega_i^c$ enters up to two layers away. The trace $\operatorname{tr}_Wh$ of a history is formed in the same way from everything that $h$ records, including the labels of its earlier moments. It is finer: knowing the trace of a label does not fix the traces of the histories that carry it. A \emph{chart} $\iota:W\to W'$ is a translation whose vector belongs to the lattice of level-$t$ cells. It acts either within one torus or, through a common lift of a non-wrapping region, from one torus and cutoff to another. It leaves invariant the lattices $\mathbb{T}^{(i)}$, $i\le t$, and also the partitions of all levels up to $t$ into cubes and cells. Neither a shift by a single site nor a mere isometry of the windows does so. Suppose that $\iota\operatorname{tr}_Wh=\operatorname{tr}_{W'}h'$, that the increments and the earlier terms in $W$ correspond, and that the bare weight corresponds on the window $W^+_t$ (this window contains $W^{\sim O(1)}$, since the cells of level $t$ are $M\check{R}_t$-cubes). Then the two histories produce corresponding terms with domains in $W$, including their natural domains of definition. Here a term is regarded as a partial map on the fields on $W$, and one of two corresponding terms is defined exactly where the other is. The integration of $V_k$ over the box is the move (3) of \cref{sec:large-field-stored-variables-and-integration-in-a-large-field-patch} and forms part of the main step of $\mathbf{R}$. Consider the part of the regular terms on a domain of level $m$ that this integration takes. It depends only on the trace through the levels $i\ge m$, whatever the values of $k$, of $i-m$, of the size of the domain and of the volume. The standing operations, that is the $\mathbf{T}'_i(X)$ of the islands carried by a term, are transported likewise, measures included and one factor at a time \cref{sec:large-field-the-induction-closed-the-order-of-the-constants-the-sum}. For the analyticity spaces more is needed. The spaces of two corresponding terms are mapped onto one another if two further conditions hold. First, the bonds of the bare lattice carrying their fields must form corresponding sets. This is needed because whether a bond leaving the domain ends in $\Omega_1^\sim$ is decided outside the window; for a domain in $\Omega_1$ there is no condition. Second, at each level, two unit cubes of the bare lattice lying in the domain must belong to a common drop-out mark of that level if and only if their images do. This is needed because a domain that is not convex may cut a mark into two pieces; where no mark meets the domain there is no condition. On the infinite lattice no sum over labels is formed. The only terms formed there are those of the never-large label, the label at which no region stands; domain by domain, these terms contain the sums over healed components that $\mathbf{R}$ performs (under \hyperref[hyp:C1]{the exact-operation hypothesis} and \hyperref[hyp:C2]{the locality hypothesis}). A term whose window $W^+_t$ lies in $\Lambda_t$ coincides with the corresponding term of that label (the terms belong to the standing part $h^\flat$, and everything that this part records lies in $Z_t$: 2e). For a non-wrapping domain and corresponding step data under the chart, the lifted term is the same function, and in this way $E^{(j)}$ and $b_j$ are computed \cref{sec:flow-a-common-companion-and-the-torus-floor}. It is this that links the geometry to \cref{sec:flow-a-common-companion-and-the-torus-floor} and \cref{sec:sources-analytic-dependence-on-the-source-locality-and-decay-in-tame}.

\textbf{The net form of the exponent.} Each function $\mathfrak{t}$ of the exponent has a fixed place, given by its level, its domain and the way it is made. The label enters only through the sorting into $\mathbf{E}_k,\mathbf{R}_k,\mathbf{B}_k$, the profile and one number:

\begin{gather*}
\FitMath{A_{k,\mathcal{S}}(U)=-\beta A(U)+\sum_{\mathfrak{t}}\mathfrak{t}(U)+\mathrm{const}(\mathcal{S}),} \\[3pt]
\FitMath{\text{because}} \\[3pt]
\FitMath{x_k(\cdot)+\sum_{j\le k}b_j\varphi_j(\cdot)=\beta}
\end{gather*}

and $c\mapsto A(c,U)$ is linear. Consequently a change of label changes no local term.

\textbf{The law of the histories.} The set $\mathfrak{H}_k$ of histories of level $k$ is finite. A run has the datum $\mathfrak{d}=(\beta,m,g)$ and takes place on $T_{\ell_m}$, $\ell_m=2L^m$. The mass of a history $h\in\mathfrak{H}_k$ is $\mathsf{t}(h):=\tau_h(\mathcal{V}_k)\ge0$. Since the steps are exact \cref{sec:large-field-the-induction-closed-the-order-of-the-constants-the-sum}, $\sum_{h\in\mathfrak{H}_k}\mathsf{t}(h)=Z_{K,\ell,\beta}$ at each level. The law of the histories of the last level, and in the same way the law of the histories of any level, is

\begin{gather*}
\FitMath{\tilde\nu_{K,\ell,\beta}(\mathsf{E}):=Z_{K,\ell,\beta}^{-1}\sum_{h\in\mathsf{E}}\mathsf{t}(h),}
\end{gather*}

where events of histories are written in sans-serif letters, $\mathsf{E}_k\subset\mathfrak{H}_k$. The creation set of step $j$ is $\mathcal{W}^\wedge_j(h)$, as before. Without the superscript, $\mathcal{W}_K(g)$ denotes the first-passage window. The live region $\mathcal{Z}_k(h)$ is $Z_k(h)$ with the components healed at step $k$ removed; it is the input region of step $k+1$. The event that the cell $\mathfrak{c}$ is live at level $k$, that is $\mathfrak{c}\subset\mathcal{Z}_k$, is written $\mathfrak{lv}_k(\mathfrak{c})$. For a cell $z$ of level $n$, the survival event $\mathcal{E}^{(n)}_z(k)$ is the event that $z\subset\mathcal{Z}_i$ for $n\le i\le k$. For the flow and history statements below, we continue the run past $K$. For these statements we fix an anchor coupling $\gamma_\star\in(0,\gamma_J]$, which is determined by the constants up to $c_a$ and by $\gamma_J$, and put $X_\star=\gamma_\star^{-2}$. The anchor level is $K_\star=K(\beta,\gamma_\star)$; it goes to infinity as $\beta$ does. In these statements the reference coupling obeys the ceiling $g\le\gamma_\star$, and the bounds on the flow hold up to $K_\star$. Also, at a given $\beta$, the run passes the window of $g$ no later than the window of $\gamma_\star$, $K(\beta,g)\le K_\star$; this follows from $g\le\gamma_\star$ and the monotonicity of the first passage.

\textbf{Lineages and marks.} Let $X$ be a component of $Z_k(h)$, at any level, that is either live at level $k$ or healed at step $k$; a component live at $k=K$ is a terminal component. Its lineage regions are $\mathsf{Z}_k(X):=X$ and, going down the levels, $\mathsf{Z}_j(X):=$ the union of all components that are live at level $j$ and contained in $\mathsf{Z}_{j+1}(X)$. Its lineage consists of these components $(j,C)$ together with the creations $(j,W)$ inside the $\mathsf{Z}_j(X)$. The weight of a level-$j$ cell is $\mathsf{r}_j:=M^4\check{R}_j^5$. We write $|C|_j$ for the number of level-$j$ cells in $C$. For comparison, $\#_j(Y)$ counts the sites of $\mathbb{T}^{(j)}$ in $Y$, and a cell contains $(M\check{R}_j)^4$ of them. Finally, for a component $W$ of the set that step $j$ creates in the lineage, $d'_j(W)$ is its tree length measured in these cells; the marks set by the operations $\mathbf{R}$ do not enter $\operatorname{Dep}$. Then

\begin{gather*}
\FitMath{\operatorname{Vol}_k(X;h):=\sum_{(j,C)\text{ in the lineage}}2\,\mathsf{r}_j\,|C|_j,} \\[3pt]
\FitMath{\operatorname{Dep}_k(X;h):=\sum_{(j,W)\text{ created in the lineage}}c\,p_0(g_j)^2\big(d'_j(W)+1\big).}
\end{gather*}

The mark at level $k$ has two dials, $\mathsf{w}\in[0,1)$ and $c_a\ge0$, the second also called the age constant. The mark and the end mark $\Theta:=\Theta^{(K)}$ are

\begin{gather*}
\FitMath{\Theta^{(k)}(h):=\sum_{X\text{ live at level }k}\big[\mathsf{w}\operatorname{Dep}_k(X;h)+c_a\operatorname{Vol}_k(X;h)\big]\ge0 .}
\end{gather*}

The mark keeps back the fraction $\mathsf{w}$ of the small factor of each creation in such a lineage. It also charges $2c_a$ for every cell, multiplied by the weight of the cell, and for every level. The constant $c$ in $\operatorname{Dep}$ is a single number. It is chosen after the fractions of the fluctuation threshold, namely those at which a step evaluates its tests: $1/(4\Theta_1)$ on the collar and $\theta_c$ on the rim. It is subject to $c\le c'\min\{\Theta_1^{-2},\theta_c^2\}$, where $c'$ is chosen after $(A_1,A_0)$ and contains the ratio $(A_1/A_0)^2$ of $\delta_k^2/g_k^2$ to $p_0(g_k)^2$. No ceiling on $g$ can take the place of this choice, since such a test has a deposit and a gain that are both of degree $2\mathfrak{p}_0$ in $\log x_j$.

\textbf{Three ages} are attached to a component $C$ that is live at level $k$. (1) The age in the label is determined by the oldest region of the label within $C$: $\operatorname{age}^{\mathrm{lab}}(C)=k-\min\{t:C\cap Z_t\ne\emptyset\}$, where $Z_t$ denotes the large-field region of the label $\mathcal{S}(h)$ at scale $t$. (2) The age of the lineage is $\operatorname{age}_k(C;h)=k-j_0$, with $j_0$ the earliest level that has a non-empty lineage region $\mathsf{Z}_j(C)$. The events $\mathfrak{lv}^{=a}_k(\mathfrak{c})$ and $\mathfrak{lv}^{\ge a}_k(\mathfrak{c})$ say that at level $k$ the cell lies in a live component whose lineage has age $a$, respectively at least $a$. (3) The seed age is $k$ minus the seed level. The seed level is the least $t$ such that $C$ contains, as a point set, some region created at step $t$. The corresponding events are written $\mathfrak{sd}^{=a}_k(\mathfrak{c})$. The first age is at most the second, and the second is at most the third. The second and the third can differ. Suppose that a region was healed at some earlier step and that a later live region again covers its place. Such a region counts for the seed level but is not part of the lineage. If the seed level is defined using the regions that no step before $k$ has healed, the two ages coincide. Payment is made for the age of the lineage alone \cref{sec:large-field-the-induction-closed-the-order-of-the-constants-the-sum}. Age has to be paid for with creations. In a lineage with $n$ creations, those made by operations included, a creation made at step $t$ satisfies $k-t\le Cn(\mathsf{N}_t+\log n)$. Here a creation is a created cell or a mark. By the waiting bound and the count of \cref{sec:large-field-rare-patches-gains-per-creation-and-costs-per-site}, this bound holds for every lineage; hence an old lineage contains many creations. These objects enter \cref{sec:large-field-the-induction-closed-the-order-of-the-constants-the-sum} and \cref{candidate:sources,ch:all-state-source-transfer}. Through the age series of \cref{candidate:sources,ch:all-state-source-transfer} they underlie both the two-point law \cref{sec:short-distance-the-short-distance-two-point-law-and-its-asymptotic-freedom}--\cref{sec:short-distance-the-law-passed-to-the-limits} and the four-point law \cref{thm:continuum-main}(c)(5), namely clauses (c)(3), (c)(4), (c)(5).


\paragraph*{History data, effective densities, and large-field operations}

\textbf{[1.6] The sequences of a label.} Let $\boldsymbol\nu'\le\boldsymbol\nu$. Then $\Omega^{\boldsymbol\nu}_l$ is a region for $\pi_l$. The chain $\Omega_0\supset\Omega^{\boldsymbol\nu}_1\supset\dots\supset\Omega^{\boldsymbol\nu}_s$ holds, and any two consecutive boundaries in it lie at least one $M$-cube of the coarser level apart, which is the separation that minimisers and propagators require. Further, $\Omega^{\boldsymbol\nu'}_l\subset\Omega^{\boldsymbol\nu}_l$ and $\Omega^n_l=\Omega_l$, so that \emph{the entries make up the determining sequence}. The action regions are unions of coarse cubes,

\begin{gather*}
\FitMath{T_\ell\supset\Omega^{\mathsf a}_1\supset\Lambda^{\mathsf a+}_1\supset\Lambda^{\mathsf a}_1\supset\Omega^{\mathsf a}_2\supset\dots\supset\Lambda^{\mathsf a}_s,} \\[3pt]
\FitMath{\Omega^{\mathsf a}_{l+1}\subset(\Lambda^{\mathsf a}_l)^{\sim-1},} \\[3pt]
\FitMath{Z^{\mathsf a}_s=Z_s,}
\end{gather*}

and the correspondence (entries, stored regions) $\leftrightarrow$ (action regions, nest) gives two presentations of one and the same datum.

\paragraph{Stored variables and integration in a large-field patch}
\label{sec:large-field-stored-variables-and-integration-in-a-large-field-patch}

\textbf{Measure and formula: two layers.} At the level of measures, a step acts as a finite sequence of moves on states $\sigma=(h,\xi)$, where $h$ records the outcomes obtained up to that point and $\xi$ the variables stored up to that point. Each move is given by positive kernels $\Pi^{(o)}$, finitely many and one for each outcome $o$, such that $\sum_o\Pi^{(o)}(\sigma,\cdot)$ has total mass one. A move is one of four kinds. It may be a test, that is, an indicator whose outcome is added to $h$. It may be a map. It may be a draw, that is, the image of a probability law, as for gauge fixings, hypercubic frames $r\in W_4$ and reference laws. Or it may delete a stored variable, which amounts to an integration. Everything expressed by a formula, such as Gaussian weights, Jacobians, expansions, resummations or the replacement of a test by $1$, only recasts an integrand in another form. A gauge fixing is a draw as a move. For it to be the factor that the formula writes, the measure must be gauge invariant. Let $\nu_0$ be any complex bare measure, let $\mu_0\ge0$, and let $\mathfrak{w}$ be any bounded Borel function of the bare field. Denote by $\mathsf{K}_h$ the kernel of the history $h$, obtained by composing the kernels $\Pi^{(o)}$ of its moves along its outcomes. Then

\begin{gather*}
\FitMath{\tau_h[\nu_0]=\int\nu_0(dU)\,\mathsf{K}_h(U,\cdot),} \\[3pt]
\FitMath{\sum_h\mathsf{K}_h(U,\mathcal{V}_k)=1\,\,\forall U,} \\[3pt]
\FitMath{\sum_h|\tau_h|(\mathcal{V}_k)\le|\nu_0|(\mathcal{V}_0),} \\[3pt]
\FitMath{\big|\tau_h[\mathfrak{w}\mu_0]\big|\le\sup|\mathfrak{w}|\,\tau_h[\mu_0].}
\end{gather*}

By totality (2a; a reference law whose normalisation does not lie in $(0,\infty)$ is not drawn, and in that case the component drops out; this case does not arise on the supports, see \cref{sec:large-field-the-induction-closed-the-order-of-the-constants-the-sum}, ``Reality and positivity''), every kernel is defined at every field, and no smallness condition is needed. They depend on the lattices and the constants. Of the frame $\theta$ (2e: the real sequence from which the choices are formed; in the present case $\theta=x$, and it is only through this frame that the kernels then vary with $\beta$; it is unrelated to the hypercubic frames $r$), they use only the gauge-fixing law, the coupling of the reference law, the cube sizes and the thresholds. Apart from their dependence through the frame, the kernels do not depend on any measure, on $\beta$, on a source, on any term of an effective action, or on whether some series converges. Exactness and positivity therefore come at no cost, because a move does nothing but split mass. The sourced densities come from $\nu_0=\rho^w_0\,dU$ with the same kernels \cref{sec:sources-the-curvature-source-with-fixed-reference-data}, and runs whose choices are frozen are runs of that same procedure \cref{sec:flow-frozen-choices-slopes-and-tuned-bare-couplings}. The estimates of \cref{sec:large-field-rare-patches-gains-per-creation-and-costs-per-site}--\cref{sec:large-field-the-induction-closed-the-order-of-the-constants-the-sum} add only one further statement: the formulas are formed on the supports and there coincide with the measures. \emph{Passengers}: divide the stored variables into those, $\mathfrak{v}$, that a move uses and the remaining ones, $\mathfrak{q}$. In step $t=k+1$ the operation $\mathbf{T}$ uses the fields of level $k$ and, from the lower levels, only the data of its backgrounds; the operation $\mathbf{R}$ uses variables only in the components of its class and next to them. Any variable not used by a move of the step is carried along unchanged, 

\begin{gather*}
\FitMath{\Pi^{(o)}((h,(\mathfrak{v},\mathfrak{q})),\cdot)=\Pi^{\flat(o)}((h,\mathfrak{v}),\cdot)\otimes\delta_{\mathfrak{q}}}
\end{gather*}

, in which $\Pi^{\flat(o)}((h,\mathfrak{v}),\cdot)$ is a measure on the states $((h,o),\mathfrak{v}')$. This is the locality of a step, in measure-theoretic form.

\textbf{$\mathbf{T}$ on a label} (Ba{\l}aban, Comm. Math. Phys. 119 (1988) 243, Sect. 3) comprises the following, carried out on the regions fixed by the tests of \cref{sec:large-field-the-tests-the-nested-regions-and-the-content-of-a}. First the multi-scale background $U_{j+1}(\mathbf{V})$ is formed, on the determining set $\mathbf{B}[\Omega_\bullet]$ \cref{ch:multilevel-backgrounds}; the third test of \cref{sec:large-field-the-tests-the-nested-regions-and-the-content-of-a} is formed with this background. Next, a change of variables is made on the strip. Then a conditional Gaussian integral is taken over $\Lambda_{j+1}$. Finally, the logarithm of this integral is split into $E$, $R$ and $B$. The operation sits between \cref{ch:local-formation-program}, in which no regions occur, and the operation $\mathbf{R}$.

\textbf{The strip.} The tests kept on $S_{j+1}\subset\Omega_{j+1}\setminus\Lambda_{j+1}$ depend on $V_{j+1}$ non-locally, by way of $V^{\mathrm{bg}}_j$. For $j\ge1$ use the joined background on $Z_{j+1}\sqcap\Lambda_j$ with $Z_{j+1}=\Lambda_{j+1}^c$, and denote its average by $V^{\mathrm{loc}}_j$. At the first step use the one-part localization. Then

\begin{gather*}
\FitMath{V^{\mathrm{bg}}_j(V^{\mathrm{loc}}_j)^{-1}=e^{i\mathbf{H}^{(j)}},} \\[3pt]
\FitMath{|\mathbf{H}^{(j)}|\le C\varepsilon_je^{-\check{R}_j};} \\[3pt]
\FitMath{V_j(b)=W'e^{-i\vartheta_{\mathrm{co}}(W')\mathbf{H}^{(j)}(b)}e^{i\mathbf{H}^{(j)}(b)}\,V^{\mathrm{loc}}_j(b),}
\end{gather*}

The last line is a bondwise change of variables with $W'=e^{i\mathcal{A}'(b)}$. In it, $\vartheta_{\mathrm{co}}$ is a smooth cut-off in $|\log W'|$. It equals zero for angles at most \(4\delta_j/3\) and one for angles at least \(5\delta_j/3\). This cut-off is applied on the cut-off zone of the strip (\cref{sec:large-field-the-tests-the-nested-regions-and-the-content-of-a}); at the remaining bonds where the substitution is made, $\vartheta_{\mathrm{co}}\equiv1$. Consequently, on the cut-off zone, a test whose threshold does not exceed $\delta_j$ is the same function of the new variable as of the old one, while a test at $2\delta_j$ is now made against the localised background. Because $|\mathbf{H}^{(j)}|\ll\delta_j$, a contraction argument shows that the map on a bond is a diffeomorphism of $G$. Its Jacobian equals $1$ wherever $\vartheta_{\mathrm{co}}$ is locally constant, and it contributes to the action as a boundary term. Once the substitution has been made, the factors of the kept tests at those bonds next to $\Lambda_{j+1}$ where no substitution takes place are omitted. On the supports these factors equal $1$, and where $\vartheta_{\mathrm{co}}<1$, the $2\delta_j$ check relative to either strip background passes by the angle bound below. Moreover, at a pivot bond of $\Omega_{j+1}$ the value of $V_j$ is taken from the constraint $\mathcal{M}(V_j)=V_{j+1}$. The strip tests after substitution have the support described below. Factorization of the integral operations also uses the locality of the integration domain and the full kernel.

Write \(W=V_j(b)V_j^{\mathrm{bg}}(b)^{-1}\) and
\(H=\mathbf H^{(j)}(b)\). In the strip change of variables the backgrounds
are held fixed while the bond variables are changed. The displayed
substitution gives
\[
 W=W'\exp\bigl(-\mathrm i\vartheta_{\mathrm{co}}(W')H\bigr).
\]
For \(W\in\mathrm{SU}(2)\) let \(\varphi(W)\in[0,\pi]\) be the angle
of its eigenvalues \(e^{\pm\mathrm i\varphi(W)}\). In particular,
\(\varphi(-\mathbf1)=\pi\). Off \(-\mathbf1\) the angle is the norm of
\((1/\mathrm i)\log W\) with the intrinsic logarithm.
On the cut-off zone the cut-off vanishes for
\(\varphi(W')\le4\delta_j/3\) and equals one for
\(\varphi(W')\ge5\delta_j/3\).
Here \(\delta_j\le1\), and we use the strip bound with
\(|H|\le\delta_j/8\). The triangle inequality for the angle gives
\(\varphi(W)\ge\varphi(W')-|H|\). If \(\varphi(W')<4\delta_j/3\),
the substitution leaves \(W=W'\). Otherwise
\[
                 \varphi(W)\ge(4/3-1/8)\delta_j
                                =29\delta_j/24 .
\]
Consequently, for \(0<a\le29/24\),
\(\varphi(W)<a\delta_j\) and \(\varphi(W')<a\delta_j\) are equivalent.
If the test passes, \(W=W'\ne-\mathbf1\). The strict inequality gives
the same conclusion at the upper allowed threshold. This proves
the assertion for the products and differences of sharp bond tests
whose bonds are all either in the cut-off zone or unchanged. At
unchanged bonds the test is literally the old one.
On the cutoff zone where \(\vartheta_{\mathrm{co}}<1\),
the upper triangle inequality gives
\(\varphi(W)<(5/3+1/8)\delta_j<2\delta_j\).
Also
\(V_j(V_j^{\mathrm{loc}})^{-1}
 =W'e^{i(1-\vartheta_{\mathrm{co}})H}\), so its angle
has the same upper bound. For an approximate test based
on the test cube's own background, use the comparison with
the localized strip background below.

\paragraph{The strip Jacobian.}
For fixed $H$ write $\vartheta=\vartheta_{\mathrm{co}}$ and
$\Psi(W')=W'\exp(-i\vartheta(W')H)$. The map is smooth on the
whole group because the cutoff is constant near both ends of
the angle chart. For given $W$ its inverse solves
$W'=W\exp(i\vartheta(W')H)$. With the bi-invariant metric
this is a contraction when $\|d\vartheta\|_\infty |H|<1$.
The cutoff derivative costs $C/\delta_j$, so the real shift
bound is used with $C|H|/\delta_j\le1/2$.

%
Measurability of the partial inverse holds on
the full unique-solution set. The graph
\(\Psi_H(W')=W\) is closed in the variables
\((H,W,W')\), and the last variable ranges
over the compact group. Its projection is
closed. The parameters with two different
solutions are the union, over \(n\ge1\), of
the projections with two solutions at
distance at least \(1/n\). Each of these
projections is closed by compactness of
the two group variables. The unique-solution
set is therefore Borel. Intersecting the
graph with \(W'\) in a closed subset of the
group shows that the recovered value is Borel
there. A fixed failure output on its
complement is measurable.

When a change of measure requires a
diffeomorphism, also require injectivity,
surjectivity and nonvanishing Jacobian
everywhere. Injectivity is the complement
of the same countable family of collision
projections, now over \(H\). The image
is compact, so surjectivity is equivalent
to containing a countable dense set of
target values. Nonvanishing of the
Jacobian everywhere is open by compactness.
This gives a Borel diffeomorphism domain.
The inverse is jointly continuous on it,
since any converging sequence of parameters
and target values has inverse subsequential
limits solving the same unique equation.

In left-trivialized tangent coordinates the differential is
\[
 X\longmapsto\operatorname{Ad}(\exp(i\vartheta H))X
                      -iH\,d\vartheta_{W'}(W'X).
\]
Factor out the adjoint, which has determinant one and fixes
$iH$. The remaining rank-one determinant gives the Haar
Jacobian for $dW$ in the variable $W'$ is
\[
 \mathcal J_b(W')=1-d\vartheta_{W'}(W'iH)>0.
\]
Its logarithm vanishes off the cutoff shell. At fixed real
$W'$ it is holomorphic in $H$ on the smaller complex ball where
$C|H|/\delta_j<1$. Here the real differential in the last display
is extended complex linearly.

For the local expansion take a decomposition
$H=\sum_Y H_Y$ in which the carriers all contain the bond read
and meet a boundary cube. Write $\ell(Y)$ for the size of a
shortest connecting tree in the cell grid, including the fixed
read. If $\sum_Y e^{\mu\ell(Y)}|H_Y|$ is available and its
product with $C/\delta_j$ is less than $1/2$, expand
$\log(1-u)=-\sum_{n\ge1}u^n/n$ with
$u=d\vartheta_{W'}(W'iH)$. Put a product on the union of its
carriers. Because they share the read, the size of a tree for
their union is no more than the sum of their tree sizes.
The absolute sum at weight $e^{\mu\ell}$ is then at most
$2(C/\delta_j)\sum_Y e^{\mu\ell(Y)}|H_Y|$.
The result still has a boundary cube on every carrier.
It uses the holomorphic shift expansion on the admitted pair
domain, while the cutoff variable itself stays real.

Where \(\vartheta_{\mathrm{co}}=1\), the substitution gives
\(V_j(b)=W'V_j^{\mathrm{loc}}(b)\). In an approximate fluctuation
test at such a bond, insert this expression for \(V_j(b)\).
The localized datum of a test cube \(\square\) uses only the
level-\((j+1)\) bonds with both ends in \(\square^{\sim4}\).
Agreement on those bonds fixes both whether the local background is
defined and its value in the test. To see that these bonds lie in
the retained region, put \(t=j+1\) and use level-\(j\) units.
The test cube has side \(L^2M_2\check R_t\), its layers have width
\(LM_2\check R_t\), and a level-\(t\) cell has side
\(LM\check R_t\). A bond being substituted has an end \(z\)
in the strip \((\Omega_t\setminus\Omega_t^{\sim-1})\cup R_t\).
If the bond occurs in the test on \(\square\), its other end, or
\(z\) itself, lies in \(\square^{\sim2}\). Each point of
\(\square^{\sim4}\) is thus at distance at most
\((L^2+6L)M_2\check R_t+1<LM\check R_t\) from \(z\), using
\(L\ge4\) and \(L^2M_2\le M\).
If \(z\) lies in the outer layer of \(\Omega_t\), its one-cell
enlargement lies in \((\Omega_t^c)^{\sim2}\).
If it lies in \(R_t\), that enlargement lies in the one-cell
enlargement of its hull. Both regions belong to the completed
stored region \(Z_t\), so \(\square^{\sim4}\subset Z_t\).
The strip construction here requires
\(V_j^{\mathrm{loc}}\) to depend on \(V_t\) only through
the bonds meeting \(\Omega_t\cap Z_t\). These restrictions determine the transformed test
including its zero value when the local background is undefined.
The statement concerns this change of variables. Independence
under a later integration requires the corresponding property
of that integration's domain and kernel.

\paragraph{The local strip background.}
For a step \(t=j+1\ge2\) use the joined region
\(W=Z_t\sqcap\Lambda_j\), with \(Z_t\) a part at level \(t\)
and \(\Lambda_j\) a part at level \(j\). Here we use the
step geometry in which their four-layer collars of
\(M_1\)-cubes lie in the corresponding ambient layers,
with the separation of the collars required for the join.
The one-cell enlargement of the proposed strip is in \(Z_t\),
and the top torus has at least \(6\rho\) sites in each direction.
Use either admitted first member of the joined sequence.
If the level-\(j\) curvature in \(W\) is at most
\(\varepsilon'\) and \(V_t=\mathcal M(V_j)\) on the
level-\(t\) bonds of \(W\), the joined test bound is
\(\max\{2L^2,\mathsf k_L\}\varepsilon'\) for sufficiently small
\(\varepsilon'\). At level \(j\) the tested holonomy is
the one of \(V_j\). At level \(t\) both types of bond
in a test give the same values of \(\mathcal M(V_j)\),
using the exact constraint for the given bonds and
the same average for the mixed ones. Its plaquette
deviation is at most \(L^2\varepsilon'+C_L(\varepsilon')^2\).
The face lift transfers this test bound to the joined datum.
Thus the joined datum has a real regular critical orbit
under the background smallness and torus conditions.

Normalize this orbit by setting its averaged contour
transports from the strip block centers equal to one at
the level-\(j\) sites that are free to be gauged. A site
that is a node of the joined determining set already
has its gauge fixed. The centers themselves are nodes.
The residual gauge relating two normalized representatives
therefore equals one at the free sites whose contour
transports were normalized as well. The level-\(j\)
averages on the strip bonds, whose blocks lie in the
joined first member and whose endpoints are such sites
or nodes, are consequently independent of the
representative. This gives \(V_j^{\mathrm{loc}}\) from
the retained real datum. For a union of components of \(W\)
the same argument restricts to that union. The determining
bonds and contour blocks used here belong to individual
components.

\paragraph{The shift in the block axial frames.}
For a strip bond $b$, include in its read set the double block
of the level-$j$ average and the contour blocks used to fix the
gauges of its two endpoints. Let $D_b$ be the least distance to
$W^c$ on the bare averaging carriers of this set, in the ambient
weighted distance, also including the order-$q$ blocks of the fixed
sites at their own levels $q\in\{j,t\}$. Consider a read on which the full datum is
in the real comparison class of \eqref{mb:real-collar-norm},
with test radius $\varepsilon_1$. Use a read set inside the common
interior of the full and joined sequences. The two sequences use
the same block centres and fix the same endpoints there. For each
of these fixed sites, use the node comparison of the real
collar estimate. In particular
these sites must be outside the wide collars.

Write the restriction of the full field as $E^w$, where
$E=e^{i\eta f}U^\circ_W$ in the notation of the collar comparison.
The level-$j$ fields $\mathcal M^jE$ and $\mathcal M^jU^\circ_W$
differ on the read set by at most
\[
 e_b=C\varepsilon_1 e^{-\delta_0 D_b/2}.
\]
Indeed the bare field has plaquette bound $C\varepsilon_1 L^{-2j}$
on these blocks, since both layers are of level at least $j$.
The relative perturbation, in the units of level $j$, is bounded
by $C\varepsilon_1 e^{-\delta_0 D_b/2}$ by
\eqref{mb:real-collar-norm}. Apply the local average estimate of
\cref{prop:averaging-1a-1-6} to the full double blocks, using the
bound for both the straightened angle and its endpoint factors.
The same estimate at the fixed sites gives
$|w(z)-1|\le C e_b$ there, by $w(z)=v_q[f](z)^{-1}$.
The order-$q$ spacing is at most $L$ times that of the level-$j$
blocks. The bound for $f$ is measured at the corresponding tops,
at weighted distance $O(1)$ from these blocks. The contour means
of the two fields in the comparison frame differ
by at most $C_L e_b$, since their transports are products of
bounded length and the derivative of the local mean is bounded
on its small-diameter chart.

For a free site $x$ with centre $y$ write $T_E(y,x)$ for the
contour mean of $\mathcal M^jE$. The normalization of the full
field uses $w(y)T_E(y,x)w(x)^{-1}$ at $x$, by gauge covariance.
Its composition with $w$ is therefore $w(y)T_E(y,x)$. The joined
field uses $T_{U^\circ_W}(y,x)$. The two factors differ by at
most $C_L e_b$ since $y$ is fixed and $w(y)$ is already close to
one. At a fixed endpoint the only comparison factor is $w(x)$.
Taking the product at the two endpoints of $b$ proves
\[
 \big|V_j^{\rm bg}(b)V_j^{\rm loc}(b)^{-1}-1\big|\le C_L e_b.
\]
The intrinsic logarithm has the same bound up to a constant.
There is no need to estimate the unconstrained values of $w$.

On the strip read sets with a full layer of level-$t$ cells
between their bond and $Z_t^c$, the read carriers use only
$O(L^t)$ of this margin. For the part of $W^c$ coming from
$\Lambda_j^c$, use the nested margin of $\Omega_t$ inside
$\Lambda_j$, with the same allowance for the read carriers.
The largest layer in the ambient sequence is $t$, so a physical
margin $cM\check R_t L^t$ gives $D_b\ge cM\check R_t$ in the
weighted distance. This yields the small real shift for such
reads, with the prefactor given by the real datum radius. Taking
the fixed $M$ large enough gives the cutoff bound at small
coupling. For each use we must check the full carriers and the
common fixed-site clause. The real endpoint estimate does not
supply the holomorphic expansion of the shift on an independent
current domain used to transfer the boundary terms.

The datum in the response equation must be distinguished from
the physical change of the averages. In the real collar comparison
write $\widehat U=(U|_W)^g$ for the aligned restriction and
$\widehat V=M_{\mathbf B^W}(\widehat U)$. The alignment gives
\[
 \widehat V=e^{iB_1}V^W,\qquad
 B_1(c)=\frac1i\log\!\left[
       \mathcal M^q(U^{v_a})(c)V^W_q(c)^{-1}\right]
\]
on the positively oriented collar bonds, and $B_1=0$ on the
remaining determining bonds. Here $c$ has level $q$ and belongs
to the collar of the part $a$. The physical increment that changes
$\widehat V$ into $V^W$ is therefore $b=-B_1$. For the response
built at $\widehat U$, its row input is $B'=b_\natural$, with
the conversion formed using $\widehat V$. At an interior row
this is $b$; at an outward row it includes the endpoint factor
of that conversion. In particular, the support of $B'$ may
include the outward rows whose endpoint blocks meet the support
of $b$.

This conversion is holomorphic wherever its stated pair and
logarithm bounds hold, and satisfies
\[
 |b_\natural|_y\le C_{\natural,\mathrm s}
             \max_{y'\in N_L(y)}|b|_{y'}.
\]
It identifies the row input once an admitted physical defect
$b$ is given. To apply the complex response theorem uniformly,
the aligned pair must lie in its joined local class and the
converted input must satisfy
\[
 \|b_\natural\|_\infty
       <\varepsilon_{\mathrm D,\mathrm s}e^{-32\varkappa}.
\]
The real collar alignment supplies a real defect of size
$C\varepsilon_1$.

\paragraph{The input on a fixed real collar chart.}
Fix the real configuration $U$ used in the collar comparison.
Let $g_0$ be its real alignment, equal to the supplied $v_a$ on
each wide collar $Q_a$ and to one elsewhere. Keep $g_0$ fixed
when the complex perturbation varies. On a part of level $j_a=j$,
use level-$j$ units and write
\[
 \mathbf U=e^{i\eta_j A'}U,\qquad
 \|A'\|_\infty\le\alpha_1,\qquad\eta_j=L^{-j}.
\]
The full averaging carriers considered below lie in the collar
region where the real plaquette bound is
$|\partial U-1|\le C\varepsilon_1\eta_j^2$.
Take $C\varepsilon_1$ and $\alpha_1$ below the fixed averaging
ceilings. Since $g_0$ is real, the same bounds hold after its
action, with $A'$ replaced by $\operatorname{Ad}(g_0)A'$.

The complex averaging formula, applied with real reference
$U^{g_0}$, gives, for $0\le q\le j$,
\[
 \mathcal M^q(\mathbf U^{g_0})(c)
 =v_q(c_-)e^{ia_q(c)}\mathcal M^q(U^{g_0})(c)v_q(c_+)^{-1},
\]
where the functions are holomorphic in $A'$ and
\[
 |a_q|\le2L^{q-j}\alpha_1,\qquad
 |v_q^{\pm1}-1|\le\frac8{Lr_*}L^{q-j}\alpha_1.
\]
For $q=0$ the site factors equal one. The central average is
unitary. Subtracting it in the preceding formula and using the
three displayed bounds therefore gives
\[
 \big|\mathcal M^q(\mathbf U^{g_0})(c)
                -\mathcal M^q(U^{g_0})(c)\big|
       \le C_L L^{q-j}\alpha_1.
\]
This estimate includes the endpoint factors. Its constant is
independent of $j$ and $q$.

Form the joined datum $V^W[\mathbf U]$ by the same rule as at
the real center: use the unchanged common rows and the face
lifts of the top averages on the collar rows. At a collar row,
the lift is either one or the value of a single top averaged
bond. The estimate at $q=j$ consequently gives
\[
 |V^W[\mathbf U](c)-V^W[U](c)|\le C_L\alpha_1.
\]
Set $\widehat V[\mathbf U]=M_{\mathbf B^W}
((\mathbf U|_W)^{g_0})$. On the collar rows define
\[
 b[\mathbf U](c)=\frac1i\log\!\left[
       V^W[\mathbf U](c)\widehat V[\mathbf U](c)^{-1}\right].
\]
At the center, $\widehat V[U]=e^{iB_1}V^W[U]$, so this
logarithm is $-B_1$. For sufficiently small
$C\varepsilon_1+\alpha_1$, its argument stays in a fixed
series-logarithm ball around one, and
\[
 \|b[\mathbf U]\|_\infty\le C_L(\varepsilon_1+\alpha_1),
 \qquad b[U]=-B_1.
\]
Indeed both matrix factors differ by $O(\alpha_1)$ from their
unitary center values, and the center product is
$e^{-iB_1}$. Inversion and the series logarithm are uniformly
bounded and holomorphic on this smaller ball.

On a common row the alignment is one at both endpoints, by
the collar geometry. Gauge covariance under the fixed real
$g_0$ thus gives $\widehat V[\mathbf U]=V^W[\mathbf U]$
there. Set $b=0$ on those rows. This proves exact support on
the collar rows and the identity
\[
           e^{ib[\mathbf U]}\widehat V[\mathbf U]
                         =V^W[\mathbf U].
\]
For a closed set of response rows, the converted datum
$B'=b_\natural$, formed at $\widehat V[\mathbf U]$, is
holomorphic and has norm at most
$C_{\natural,\mathrm s}C_L(\varepsilon_1+\alpha_1)$.
Its carriers include the endpoint blocks read by this conversion.
The response radius is met when this bound is smaller than
$\varepsilon_{\mathrm D,\mathrm s}e^{-32\varkappa}$.

The current remains an independent variable. First form the
restricted pair by keeping $\mathbf J$ at the bonds whose
incident plaquettes all remain inside the joined first member,
and by using the current of the restricted action only on its
free boundary. Then apply the fixed real gauge $g_0$ to both
members of this pair. The defect above does not depend on
$\mathbf J$. Thus its holomorphy imposes no equation
$\mathbf J=J(\mathbf U)$ in the interior. The response still
requires this restricted and gauged pair to belong to its
joined local class.

\paragraph{A target defined by the complex configuration.}\hypertarget{hi:strip-canonical-target}{}
The tower transports give a choice of target that depends
only on the complex configuration. The response can remain
based at the original restricted pair.

On a wide collar $Q_a$ of level $j$, write
$\mathbf W_q=\mathcal M^q\mathbf U$. For a bare site $x$
let $x_q$ be its ancestor at level $q$, so that
$x\in B^q(x_q)$, and put $y=x_j$. Define
\[
 g[\mathbf U](x)
  =\mathsf T_{\mathbf U}(y,x)
  :=\mathbf W_{j-1}(y,x_{j-1})
       \mathbf W_{j-2}(x_{j-1},x_{j-2})\cdots
       \mathbf W_0(x_1,x).
\]
Set $g[\mathbf U]=1$ off the collars and at their top sites;
for $j=0$ the product is empty. This is the supplied tower
transport, using the averaged contour variables of the
same averaging scheme. Its factors depend only on the
configuration in $B^j(y)$. On their averaging domain they
are holomorphic functions of $\mathbf U$ itself.
At a real configuration $U$ this is exactly
the descending collar alignment $g_0$.

The tower estimate is uniform in its depth. In the local
representation $\mathbf U=e^{i\eta_jA'}U$ used above,
with $\|A'\|_\infty\le\alpha_1$, the supplied product
formula writes
\[
 E(x):=g[\mathbf U](x)g_0(x)^{-1}
                 =\prod_{q=0}^{j-1}e^{Z_q(x)},
 \qquad
 \sum_q|Z_q(x)|\le
 z:=2\kappa_6 C''_{\mathrm{mav}}L^{-1}\alpha_1
                       \le\tfrac18.
\]
The order of the factors is the one in the tower product.
Thus $|E^{\pm1}-1|\le e^z-1\le C_L\alpha_1$.
The real matrices $g_0$ are unitary, so these bounds also
give uniform bounds on $g[\mathbf U]$ and its inverse.
Every full averaging carrier used here is included in the
collar carrier, as in the preceding chart estimate.

Let
\[
 V_*=M_{\mathbf B^W}(\mathbf U|_W),\qquad
 T=V^W[\mathbf U],\qquad
 R[\mathbf U]=T^{g[\mathbf U]^{-1}}.
\]
The last expression is just the action on the endpoints
of the determining rows. Define the physical input at the
original restricted pair by
\[
 b_{\mathrm{can}}(c)=\frac1i\log
           \bigl(R[\mathbf U](c)V_*(c)^{-1}\bigr).
\]
This expression is holomorphic and has the same uniform
size as the previous chart input. To check this, set
$\widehat V_0=V_*^{g_0}$. At a row $c$, its logarithm's
argument equals
\[
 g_0(c_-)^{-1}
 \bigl[E(c_-)^{-1}T(c)E(c_+)
                          \widehat V_0(c)^{-1}\bigr]
 g_0(c_-).
\]
The preceding chart calculation bounds
$T\widehat V_0^{-1}-1$ by
$C_L(\varepsilon_1+\alpha_1)$. The bounds on $E^{\pm1}$,
together with the bounded matrix factors in that chart,
give the same bound for the bracket minus one. On the
fixed logarithm ball this proves
\[
 \|b_{\mathrm{can}}\|_\infty
       \le C_L(\varepsilon_1+\alpha_1),\qquad
 e^{ib_{\mathrm{can}}}V_*=R[\mathbf U].
\]
At both endpoints of a common row the tower gauge is one,
by the same collar geometry used for the real alignment.
There $T=V_*$ and $b_{\mathrm{can}}=0$. Its possible
nonzero rows are therefore collar rows. Form
$B'_{\mathrm{can}}=(b_{\mathrm{can}})_\natural$ at the
original restricted pair. On a closed determining row set,
the conversion is holomorphic and gives
\[
 \|B'_{\mathrm{can}}\|_\infty
   \le C_{\natural,\mathrm s}C_L
                      (\varepsilon_1+\alpha_1).
\]
At a collar row $c$ of level $q<j$, let $y_\pm$ be the
level-$j$ ancestors of $c_\pm$. Its physical input is
determined by the bare configuration in
$B^j(y_-)\cup B^j(y_+)$: these are the two tower blocks,
the lower averaging blocks lie inside them, and the face
lift reads either one or their common top bond. Both
blocks lie in $Q_a$. The converted input also reads the
neighboring physical rows and the endpoint blocks
specified by the conversion. For a localized term these
complete carriers must lie in its admitted enlarged
domain. The response is defined when the displayed norm
fits its radius and the original restricted pair is
admitted to the joined class. The response is evaluated
directly at that pair.

This target also gives the required gauge transformation
law. For a real gauge $u$, define $\bar u(x)=u(y)$ in a
top block $B^j(y)$ of a collar and $\bar u=u$ elsewhere.
Covariance of the averaged contour variables cancels the
intermediate factors in the tower product and gives
\[
 g[\mathbf U^u]=\bar u\,g[\mathbf U]u^{-1}.
\]
The face lift is one inside a top block and takes the
value of one top bond across its face. Hence its target
transforms by $T[\mathbf U^u]=T[\mathbf U]^{\bar u}$.
The same identity holds on common rows, whose endpoints
in a collar are top sites. Consequently
\[
 R[\mathbf U^u]=R[\mathbf U]^u,\qquad
 b_{\mathrm{can}}[\mathbf U^u]
          =\operatorname{Ad}(u)b_{\mathrm{can}}[\mathbf U].
\]
Covariance of the conversion and of the response applies
at the unchanged restricted pair, with every interior
current entry still independent. The resulting response
is thus a holomorphic function of the pair and is
covariant for all admitted complex cube parameters.
Its definition is identical in overlapping real
factorizations, since every factor in the input is a
function of $\mathbf U$ itself.

At a real center, $g[U]=g_0$ and
$\widehat V_0=e^{iB_1}V^W[U]$. It follows that
\[
 b_{\mathrm{can}}[U]
          =\operatorname{Ad}(g_0^{-1})(-B_1).
\]
The converted inputs obey the same identity, with the
conversion on each side taken at its own base pair.
Real-gauge covariance of the response therefore identifies
this response with the earlier aligned response, conjugated
by $g_0^{-1}$. On a full read star outside the collars,
$g_0=1$, so their covariantly formed dressed reads agree
in the same fixed frame. For other read frames the explicit
frame factors are retained.

\paragraph{The rows and their complete carriers.}
Use the rows prescribed by the decoupled response theorem:
the positively oriented interior and outward bonds of
the joined determining set $(\mathbf B^W)^\flat$.
These rows form a closed set. Indeed, if a row of level
$q$ has outer endpoint $z$, every positive level-$(q-1)$
bond joining sites of $B(z)$ is an interior row of the
joined set. These are exactly the lower rows required
by the conversion $b\mapsto b_\natural$.

Let $\mathfrak k$ be an active wall-component of the
parameter partition $\mathsf P=\pi[\mathsf w]$, and put
$Y=\bigcup\mathfrak k$. Work in level-$t$ units. Every
collar top level satisfies $j\le t$, so its block has
side $\xi_j=L^{j-t}\le1\le\mathsf w$. Both side lengths
are powers of $L$ on the same centered grids. The nesting
of standard partitions has the following consequence:
if a bare lattice site $x$ lies in $Y$, then the entire
level-$j$ block containing $x$ lies in $Y$. To see this,
the site belongs to one parameter cube, since lattice
sites do not lie on its boundary walls. That cube is a
union of complete level-$j$ blocks. Its unique such
block containing $x$ is therefore contained in it.

Write $\mathbb Y^\circ_{\mathfrak k}$ for the sites
carried by the component's walk terms. The response
uses its source only at rows at
$N_L(\mathbb Y^\circ_{\mathfrak k})$. The conversion at
such a row uses physical inputs only at rows at
$N_L^2(\mathbb Y^\circ_{\mathfrak k})$. The endpoints
of these rows belong to
$N_L^3(\mathbb Y^\circ_{\mathfrak k})$: for an interior
row the neighboring endpoint has an adjacent block;
for an outward row this is the outer-end assertion
$B(z)\subset N_L(y_c)$.
The margin lemma places the blocks of all sites in
$N_L^4(\mathbb Y^\circ_{\mathfrak k})$ in the interior
of $Y$.

Consequently every endpoint just listed lies in $Y$,
and the partition argument places its complete collar
top block in $Y$. Thus every tower transport and every
top average used by a required physical input is read
inside $Y$. The lower averaging carriers and the
endpoint-block carriers of the conversion are inside
$Y$ by the same argument. The complete canonical input
carriers thus lie in the active component even when a
top block is larger than the joined row's own scale.

The original collar geometry is retained in this
argument. At a noncommon row, its two endpoints lie
in the wide collar of one part, in the original layer
$j$, and the complete ancestor blocks lie in that
collar. The original and joined region traces therefore
specify the common-row classification, this original
level and the tower blocks used in the source. When
these traces agree on $Y$, the canonical inputs read
by the component agree there as well. The margin
comparison for the response then applies. This uses
the original traces in addition to the joined traces;
the latter alone need not identify the collar top
level.

\paragraph{The source radius.}
Let $K_B$ denote the constant in the canonical input
bound, including the conversion and the fixed changes
of local units, so that
\[
 \|B'_{\mathrm{can}}\|_\infty
                   \le K_B(\alpha_0+\alpha_1).
\]
Suppose that the joined pair has the fixed configuration
radii required by the response theorem and current
radius at most $K_J\alpha_0$. In the strip domain
$\alpha_0+\alpha_1\le c e^{-C\varkappa}$, choose
$C\ge32$ and
\[
 0<c<\min\left\{c_0,
       \frac{\varepsilon_{\mathrm D,\mathrm s}}
                         {2\max\{K_B,K_J,1\}}\right\},
\]
where $c_0$ is smaller than the averaging and physical
logarithm ceilings used in the canonical construction.
Then both the converted source and the current radius
are smaller than
$\frac12\varepsilon_{\mathrm D,\mathrm s}e^{-32\varkappa}$.
The constants are chosen after the geometric scales
and before the final small-coupling threshold, as in
the strip domain. Both radii enter the boundary size.
Taking the infimum over the admissible radii gives
\[
 \|B'_{\mathrm{can}}\|_\infty\le K_B\mathsf n_W(P).
\]
The input has no cube-parameter dependence. Its radius
and carrier bounds consequently hold throughout the
parameter polydisc of the response theorem.

\paragraph{The retained current and the free boundary.}
The joined sequence is finer than the original sequence.
At a site where its layer is $q$ and the original layer
is $j$, one has $q\le j$ and $\xi_q\le\xi_j$.
Consequently the original pointwise bounds on curvature,
the complex potential and its covariant difference,
and the independent current imply the corresponding
joined pointwise bounds at every retained stencil.
This assertion concerns the pointwise clauses; local
gauges are a separate clause of the pair class.

At a free-boundary bond replace the current by the first
variation of the restricted plaquette action, as prescribed
above. For any set of retained plaquettes the supplied
first-variation formula is
\[
 J^\flat(\mathbf U)=D_{\mathbf U}^{\mathsf T}
                         F_{\mathbf U},\qquad
 F_{\mathbf U}(p)=\frac{\partial\mathbf U(p)
                 -\partial\mathbf U(p)^{-1}}{2i\xi^2}.
\]
Each bond receives at most six plaquette contributions.
The scaled curvature bound implies
$|F_{\mathbf U}(p)|\le C\alpha_0\xi^{-2}$, and the
one-bond transports in $D_{\mathbf U}^{\mathsf T}$ are
uniformly bounded by the complex-potential ceiling.
It follows that $|J^\flat(\mathbf U)(b)|
\le C\alpha_0\xi^{-3}$.
A free-boundary bond has a missing incident plaquette
within two bare spacings. Separation of the first two
joined regions therefore places its initial site in
joined layer zero. Since $\xi_0=\xi$, the bound reads
\[
                  \xi_0^3|J^\flat(\mathbf U)(b)|
                                      \le C\alpha_0.
\]
Together with the retained-current bound this gives
the strip current radius used above.

\paragraph{The local gauges on the joined cubes.}
For a triple in the structured domain on the full carrier,
the pointwise and good-gauge clauses in
\cref{rg:structured-local-gauge} give the original enlarged
local class. Restricting to the active parameter union keeps
the same splitting and all the inner stencils. The gauges of
the joined class are then obtained from the original gauges. We use the enlarged class
$\mathfrak P_t^\natural$, whose gauge cubes of index $q$
are all the cubes of side at most $12\mathfrak M\xi_q$
that are unions of $\mathfrak M$-cubes at that level and
lie in $\mathcal L_q\cup\mathcal L_{q+1}$. Such a cube
need not meet $\mathcal L_q$.

First recall the supplied truncation argument. Removing
the last member from a separated sequence carries
$\mathfrak P_{k+1}^\natural$ into
$\mathfrak P_k^\natural$ with the same radii and the
same decomposition $\mathbf U=e^{i\xi A'}U$.
All scaled inequalities can be read in the same original
units throughout. For the gauge clause, only a cube
$\mathsf q$ of index $k-1$ which meets the former last
region $\Omega_{k+1}$ needs attention. Its side satisfies
\[
       s(\mathsf q)\le12\mathfrak M\xi_{k-1}
                              <\mathfrak M\xi_k.
\]
Cover it by a cube $\mathsf q'$ of side
$2\mathfrak M\xi_k$, made of $2^4$ cubes at scale
$\mathfrak M\xi_k$. If $x\in\mathsf q\cap\Omega_{k+1}$,
every point of $\mathsf q'$ is within
$2\mathfrak M\xi_k$ of $x$. Separation therefore puts
$\mathsf q'$ in $\Omega_k=\mathcal L_k\cup\mathcal L_{k+1}$.
The original gauge clause at $\mathsf q'$ gives the
gauge on $\mathsf q$, with
\[
 \xi_{k-1}|A|\le L^{-1}a_0,
 \qquad
 \xi_{k-1}^2|\nabla^\xi A|\le L^{-2}a_0,
\]
where $a_0$ is its original gauge radius. Thus no loss
accumulates under successive truncations.

The same argument applies to the local class on an active
union $Y$ of parameter cubes. Here it is necessary to
choose $\mathsf q'$ inside $Y$. This is possible on the
actual nested grids. Put $m=\mathfrak M\xi_k$; then
$m$ divides $\mathsf w$ and $\mathsf w\ge2m$.
The parameter cubes whose interiors meet $\mathsf q$
form a rectangular block of cubes, all belonging to
$Y$. In each coordinate the interval of $\mathsf q$,
whose length is less than $m$, meets at most two
$m$-grid intervals. If it meets only one, include an
adjacent interval within that rectangular block.
The block has at least two such intervals in every
coordinate. The product of the chosen two-interval
unions is the required cube $\mathsf q'\subset Y$.
It still has side $2m$, so the separation argument is
unchanged. The original local gauge clause is therefore
applicable to $\mathsf q'$.

Now let $\widetilde{\mathsf q}$ be any gauge cube of
index $q$ for the joined sequence, contained in $Y$.
Fineness gives
\[
 \widetilde{\mathsf q}
   \subset\mathcal L_q^W\cup\mathcal L_{q+1}^W
   \subset\Omega_q^W\subset\Omega_q.
\]
Truncate the original sequence at $q$. In that sequence
$\Omega_q$ is the last layer, so
$\widetilde{\mathsf q}$ is one of its gauge cubes of
index $q$. The preceding local truncation argument
supplies a real gauge on a cube containing
$\widetilde{\mathsf q}$, with
$\xi_q|A|,\xi_q^2|\nabla^\xi A|\le a_0$.
Restrict this gauge to $\widetilde{\mathsf q}$ itself.
This is an allowed choice of the cube in the joined
gauge clause; it lies in $\Omega_0^W\cap Y$, and every
derivative stencil retained there is unchanged.
Thus the joined local-gauge clause holds with the
original radius.

Together with the pointwise and free-boundary-current
estimates above, this proves joined-pair admission from
the original enlarged local class on $\Omega_0\cap Y$.
The argument uses the actual standard parameter cubes
and the original local-class decomposition. It does not
require a holomorphic choice of the real local gauges.

\paragraph{Localization with the bond read fixed.}
Suppose the strip datum $B'[P]$ has been supplied on the joined
rows, holomorphically on the required pair domain, with
$\|B'[P]\|_\infty\le\beta<\varepsilon_{\mathrm D,\mathrm s}
e^{-32\varkappa}$. Its values are fixed while the cube parameters
are varied. Include the full local masks of this datum in its
carriers, and let $\mathfrak B$ be the cubes carrying its possible
nonzero rows. Thus these are boundary cubes, with the prescribed
fixed enlargement when the row conversion uses an endpoint block.
Use the admitted joined pair throughout.

Let $N_b$ be a wall-connected set of parameter cubes, where
adjacency means sharing a face. Include enough local margin
that the full averaging and frame read at $b$, with every read
stencil, lies in the interior of its union. Hold the parameters
on $N_b$ at one and vary those
on $S=\mathsf P\setminus N_b$. Write
\[
 F_b(s;P)=\mathsf{Rd}_b\bigl(P;
          \mathbf H(s;B'[P])\bigr),\qquad R=e^{\varkappa}.
\]
Here the read vanishes at zero response and is holomorphic on
the response ball. If its bound in the local response norm is
$C_{\rm rd}$, the response estimate gives
\[
 \sup_{|s_\Delta|<R}|F_b(s;P)|\le A,
 \qquad A=C_{\rm rd}C_{\mathrm D,\mathrm s}
                              e^{16\varkappa}\beta.
\]
All suprema below use the same admitted pair domain.

For $K\subset S$ define the mixed difference
\[
 H_b[K;P]=\sum_{A_0\subset K}(-1)^{|K|-|A_0|}
                            F_b(\mathbf1_{A_0};P).
\]
The fixed parameters on $N_b$ are understood in every summand.
Finite M\"obius inversion gives
$F_b(\mathbf1_S;P)=\sum_{K\subset S}H_b[K;P]$.
The response restricted to a wall component uses only the pair
and datum in that component. Since $N_b$ is wall-connected, the read
uses its one component. If $K\cup N_b$ has another component,
pair the terms in the mixed difference by switching one cube
of that component; they cancel. If the component of $N_b$ meets
no cube of $\mathfrak B$, its datum is zero and its response is
zero. Therefore
\[
 H_b[K;P]=0\quad\text{unless }K\cup N_b
 \text{ is wall-connected and meets }\mathfrak B.
\]
This reasoning uses the full component locality of the response,
not just decay of its kernel.

For $1<r<R$, the exact Cauchy formula for the difference is
\[
 H_b[K;P]=\frac1{(2\pi i)^{|K|}}
 \int_{|z_\Delta|=r,\,\Delta\in K}
 F_b(z_K,0_{S\setminus K};P)
                 \prod_{\Delta\in K}
                    \frac{dz_\Delta}{z_\Delta(z_\Delta-1)}.
\]
For $K=\emptyset$ this means $H_b[\emptyset;P]=F_b(0;P)$.
Taking $r\uparrow R$ gives
\[
                  \|H_b[K]\|\le A(R-1)^{-|K|}.
\]
Each parameter is differenced once, so there is no factorial.
The norm here is a supremum before summation.

Let $W_K$ contain $K\cup N_b$ and all its local read and datum
masks. Use a fixed bounded enlargement, and write $\ell(W)$
for the tree length on the output cube grid. Keep the geometric
comparison explicit:
\[
 \ell(W_K)\le c_h(|K|+n_b),\qquad n_b=|N_b|.
\]
The constant $c_h$ includes the change between the parameter
grid and the output grid. The rooted cube count gives constants
$C_b,c_{\rm e}$, independent of the torus and the number of
levels, such that at most $C_be^{c_{\rm e}n}$ wall-connected sets
$K\cup N_b$ have $|K|=n$. Here $C_b$ depends on the bounded
read star. Put
\[
 n_*:=\min\{|K|:K\cup N_b\text{ is wall-connected and meets }
                            \mathfrak B\},\qquad
 q_a:=\frac{e^{c_{\rm e}+a c_h}}{R-1}.
\]
If no such set exists, all terms vanish. Otherwise, for $q_a<1$,
\[
 \sum_K e^{a\ell(W_K)}\|H_b[K]\|
 \le A C_b e^{a c_hn_b}\frac{q_a^{n_*}}{1-q_a}.
\]
Group the terms having the same $W_K$. Their sum depends only
on the fields in $W_K$ and retains the boundary cube and the
whole read star. On each local pair domain where the component
equations and the supplied datum have been admitted with the
same bounds, this sum is holomorphic and obeys the same absolute
estimate. Locality alone would not extend a bound to local data
outside those domains. Reality and covariance of the datum,
response and read pass through the mixed differences.

For $\varkappa\ge1$, one may use
$\log(R-1)\ge\varkappa-1$. For example, when
$\varkappa>1+c_{\rm e}$, the choice
$a=(\varkappa-1-c_{\rm e})/(2c_h)$ gives
$q_a\le e^{-(\varkappa-1-c_{\rm e})/2}$.
The strip margin supplies the lower bound on $n_*$; it must be
measured after the read and datum masks have been included.
The factor $e^{16\varkappa}$ in $A$ and the fixed read-star cost
remain in the displayed estimate and can then be absorbed by
that margin. If the parameter partition is the level-$t$
$M$-cube partition, the response floor is
$\mathring\delta_0M\ge4352\varkappa$, in addition to its other
size conditions. Choosing $\varkappa$ proportional to
$\mathring\delta_0M$ within this floor gives a tree rate of
that order. This final choice does not supply the missing
holomorphic datum or its norm bound.

\paragraph{The background of an approximate test on the cutoff zone.}
Let $b$ be a bond of the cutoff zone that occurs in the
approximate test of $\square'$. Use bare length units and set
$w=M_2\check R_t L^t$, so that the test cubes have side $Lw$
and each of their enlargement layers has width $w$. By the
definition of the cutoff zone, an endpoint of $b$ is within
$2w+L^j$ of a test cube in $R$. Since $b$ is tested at
$\square'$, every point of $C=\square'^{\sim4}$ is within
$(L+8)w+2L^j$ of $R$. The region $R$ lies inside the one-cell
inset of $\Omega_t$. The one-cell enlargement of its hull
lies in $Z_t$. The side of a level-$t$ cell is
$M\check R_tL^t$. Choose the fixed $M$ with enough room that
the preceding distance leaves at least four $M_1$-layers in
both of these margins. For $t\ge2$ also use the one-cell
margin of $\Omega_t$ inside $\Lambda_j$. This is a margin in
level-$j$ cells, with side $M\check R_jL^j$. The cube sizes
at these consecutive levels are comparable, so the same
fixed choice of $M$ leaves room in this margin. For $t=1$
only the part $Z_1$ is used. Thus $C$ lies in the top layer of
the joined strip sequence.

In that top layer $C$ is a one-part collared region at level
$t$, with collar width $M_1$ and the full first member. Its joined sequence is the
tower of $C$, since all the members of the strip sequence
contain $C$. The top rows on $C$ carry the prescribed $V_t$.
The new collar rows are formed by the same lifts of those top
rows as in the definition of the test background. This identifies
the joined datum with the datum of $U_{t,\square'}$. In
particular its regular orbit is defined when the strip datum
belongs to the real comparison class.

The test reads $b$ inside $\square'^{\sim3}$, where its
endpoint contour blocks are wholly in the top layer of the
tower. In both sequences their centres are nodes and the
other level-$j$ sites of these blocks are free. The block
axial frames therefore fix the same sites in the comparison.
The read carriers have distance at least $w-CL^t$ from $C^c$,
which exceeds $3M_1L^t$. The top spacing in bare units is $L^t$.
Apply the real
collar estimate \eqref{mb:real-collar-norm} and the endpoint
calculation for the strip shift to the strip background and
its localization to $C$. It follows that
\[
 \varphi\bigl(V_j^{\mathrm{loc}}(b)
                    V_{j,\square'}^{\mathrm{bg}}(b)^{-1}\bigr)
       \le C\varepsilon_1 e^{-c M_2\check R_t},
\]
with $\varepsilon_1$ the real test radius of the joined strip
datum. On the strip support it is bounded by $C_L\varepsilon_j$.
The fixed scale floors and then the coupling ceiling give a
bound $\delta_j/8$ for this angle. Use also $|H|\le\delta_j/8$
for the strip shift. Where $\vartheta_{\mathrm{co}}<1$,
the angles of both $V_j(b)V_{j,\square'}^{\mathrm{bg}}(b)^{-1}$
and $W'V_j^{\mathrm{loc}}(b)V_{j,\square'}^{\mathrm{bg}}(b)^{-1}$
are at most $(5/3+1/8+1/8)\delta_j<2\delta_j$. Their group
distances from one are no larger than their angles. Where
$\vartheta_{\mathrm{co}}=1$ the two products coincide. Hence
on the cutoff zone the approximate bond factor can be written
through $W'$ and the two localized backgrounds on this support.

\textbf{What is frozen and what is integrated.} On $\Omega_{j+1}^c$, step $j+1$ stores $V_j$ together with its constraint. On $\Omega_{j+1}$ the action is expanded around $U_{j+1}(\mathbf{V})$. The only variable integrated there is $A=\mathcal{A}^{\mathrm{s}}_j|_{\Lambda_{j+1}}$, conditionally on the stored field $A'=\mathcal{A}^{\mathrm{s}}_j|_{\Omega_{j+1}\setminus\Lambda_{j+1}}$. The measure is Gaussian with covariance $C^{(j)}(\Lambda_{j+1})$, which is the inverse of the restriction of the form $\Delta^{(j)}$ to the bonds of $\Lambda_{j+1}$, and its linear term is $-\langle A,\Delta^{(j)}A'\rangle$. The coupling flows only where this integration is carried out, $\varphi_{j+1}\in C_0^\infty(\Lambda_{j+1})$, and it stops on the whole of $Z_{j+1}$, including $\Omega_{j+1}\setminus\Lambda_{j+1}$. Once the square is completed, the conditional weight appears as a factor inside $\mathbf{T}_k(Z)=\prod_{j<k}\mathbf{T}^{(j)}$, with the regions cut by the component $Z$ ((2.21), (3.23) there; this operation is the product shown for a component on which $\mathbf{R}$ has never acted, and in general the operation is defined by a recursion over the marks):

\begin{gather*}
\FitMath{\mathbf{T}^{(j)}=\int dV_j\big|_{\Omega_{j+1}^c}\delta\big(\mathcal{M}(V_j)V_{j+1}^{-1}\big)\,\mathfrak{G}_j\,\chi\chi^c\,z^{(j)}\!\int dA'\,\chi\chi^c\,e^{-\frac12\langle A',[\Delta^{(j)}-\Delta^{(j)}C^{(j)}(\Lambda_{j+1})\Delta^{(j)}]A'\rangle}.}
\end{gather*}

This operation contains no Wilson weight. That weight belongs to $\exp A_{k,\mathcal{S}}$, the function on which the operation acts; $z^{(j)}$ is explained below. The old $R$-terms, $\mathbf{R}'_j=\sum_{i\le j_1}\mathbf{R}^{(i)}$, appear alongside $\mathbf{E}_j$ in the regular action $\widetilde{\mathbf{E}}_j$ of the step \cref{sec:large-field-the-induction-closed-the-order-of-the-constants-the-sum}. Here $\mathbf{R}^{(i)}=\sum_X[R^{(i)}(X,\cdot)-R^{(i)}(X,\mathbf{1})]$ is the summand of level $i$ of $\mathbf{R}_j$, and $j_1$ is the index $k_1$ of \cref{candidate:large_field} at level $j$. The young ones, $\mathbf{R}''_j=\sum_{j_1<i\le j}\mathbf{R}^{(i)}$, are handled separately. Put $\delta F:=F(U_j)-F(U_{j+1})$. Attach a parameter $t$ to the terms present only next to $Z_j$, and a parameter $s$ to $\mathbf{R}''_j$ ((3.23)--(3.28) and pp. 270--271 there). Let $z^{(j)}=\prod_cz^{(j)}(c)$ be the factors that come from eliminating the delta functions. In $\mathbf{T}^{(j)}$ their product runs over those coarse bonds $c$ of $\Omega_{j+1}$ that lie outside $\Lambda_{j+1}$; under the logarithm it runs over those of $\Lambda_{j+1}$. Then

\begin{gather*}
\FitMath{\log\Big[z^{(j)}\!\!\int\!dA\,\chi^{(j)}e^{-\frac12\langle A,\Delta^{(j)}A\rangle-\langle A,\Delta^{(j)}A'\rangle+\cdots}\Big]-\tfrac12\big\langle A',\Delta^{(j)}C^{(j)}(\Lambda_{j+1})\Delta^{(j)}A'\big\rangle} \\[3pt]
\FitMath{=\int_0^1\!dt\,\big\langle\delta\mathbf{B}_j-\delta A(x_j(\cdot)-x_j,\cdot)+\cdots\big\rangle_{s=1,t}+\int_0^1\!ds\,\langle\delta\mathbf{R}''_j\rangle_{s,t=0}+\Big(\log[\cdots]-\tfrac12\big\langle A',\Delta^{(j)}C^{(j)}(\Lambda_{j+1})\Delta^{(j)}A'\big\rangle\Big)_{s=t=0}.}
\end{gather*}

In this identity $\langle\,\cdot\,\rangle_{s,t}$ denotes the normalised mean with respect to the measure that the parameters $s$, $t$ interpolate. The left-hand side is the unnormalised integral. Its second term is present since the square, once completed, was moved into $\mathbf{T}^{(j)}$. Inside the integral over $t$, the dots stand for two kinds of terms. Those of the first kind are the terms of the strip, namely the Jacobians together with the differences of all quantities that depend on a shifted variable. Those of the second kind are the differences, taken between the two backgrounds, of the exponents of the integral operations carried by the term. The coupling difference carries the sign it has in $A_{j,\mathcal{S}}=-x_jA-A(x_j(\cdot)-x_j,\cdot)+\dots$. A cluster expansion as in \cref{ch:local-formation-program} is applied to each part. The Gaussian normaliser is contained in the last part; it is that part of $E^{(j+1)}$ which carries $b^0_{j+1}$. On a label it is a determinant: that of a multi-scale operator, compressed to $\Lambda_{j+1}$ and taken in the background of that label; so by definition it is non-local. It is converted into local terms,

\begin{gather*}
\FitMath{\sum_c\log z^{(j)}(c)-\tfrac12\log\det\big(\mathsf{C}^{\mathsf{T}}\Delta^{(j)}\mathsf{C}\big|_{\Lambda_{j+1}}\big)=\text{number}+\sum_X\mathbf{N}(X,U_{j+1}),} \\[3pt]
\FitMath{|\mathbf{N}(X)|\le Ce^{-cMd_{j+1}(X)},}
\end{gather*}

which Ba{\l}aban does by means of a resolvent integral and walks \cref{ch:local-formation-program}, whereas here it is done by sums with alternating signs, taken over the sub-regions $Y\subset X$, whose terms are the diagonal entries of the logarithm of the operator built on $Y$. M{\"o}bius inversion yields the sum rule, and the decay is due to a sign-reversing involution, which cancels everything except the chains of walks that visit all of $X$. Only the last two parts give terms with domain $X\subset\Lambda_{j+1}$. These terms are built from walks and fields inside $X$, so the regions do not affect them. They are of two kinds. The first is $R^{(j+1)}$; it arises from the part with $s$, to which what $\mathbf{R}$ leaves is added \cref{candidate:large_field}. The second is $E^{(j+1)}$, which comes from the last part, where $\mathbf{E}_j$ stands with the old $R$-terms. These terms are the same for every label. For the regular terms this holds for every such $X$. For the remainders it holds for $X$ one cell inside $\Lambda_{j+1}$, the only place where a label shows them, and this is the margin of \cref{sec:large-field-the-induction-closed-the-order-of-the-constants-the-sum}. Thus the young $R$-terms go into $R^{(j+1)}$, and the old ones go into $E^{(j+1)}$ and $b_{j+1}$ (2d, 2e). Everything else makes up $B^{(j+1)}$; among it is the part of the quadratic form that is rejected (next paragraph). It amounts to $O(1)$ for each $M$-cube of level $j+1$ next to $Z_{j+1}$; this cost is covered by \cref{sec:large-field-rare-patches-gains-per-creation-and-costs-per-site}, and for that reason no renormalisation of it is required.

\textbf{What a one-step operation keeps.} The conditional form is a Schur complement on the free coordinates. In the following display the factors of \(\Delta^{(j)}\) act as its free-coordinate precision \(\mathsf C^{\mathsf T}\Delta^{(j)}\mathsf C\). A chain is either a precision walk or a product of a precision walk, a walk of the inverse of the integrated compression, and a precision walk. These chains have connected carriers. A chain is kept if the cells of level $j+1$ that it meets intersect $\Omega_{j+1}\cap Z_{j+1}$ in a non-empty connected set and remain within one layer of that set; a kept chain is then counted with the component $\mathfrak{k}$ of $\Omega_{j+1}\cap Z_{j+1}$ that contains this set:

\begin{gather*}
\FitMath{\big\langle A',[\Delta^{(j)}-\Delta^{(j)}C^{(j)}(\Lambda_{j+1})\Delta^{(j)}]A'\big\rangle=\sum_{\mathfrak{k}}\langle A',\Delta^{\mathrm{loc}}_{\mathfrak{k}}A'\rangle-2\sum_X\mathsf{B}^\times(X),} \\[3pt]
\FitMath{|\mathsf{B}^\times(X)|\le Cp_0(g_j)^2e^{-cM\check{R}_{j+1}}e^{-\kappa d_{j+1}(X)},} \\[3pt]
\FitMath{\langle A',\Delta^{\mathrm{loc}}_{\mathfrak{k}}A'\rangle\ge c\sum_{b\text{ in }\mathfrak{k}}|A'(b)|^2:}
\end{gather*}

(the bound on $\mathsf{B}^\times$ is for stored fields in range, $|A'(b)|\le Cp_0(g_j)$, which they are on the support of the kept tests) a rejected chain has length at least a fixed multiple of the cell side in fluctuation units, up to a fixed additive constant, and the size $p_0(g_j)^2$ of the stored field squared is outweighed by $e^{-cM\check{R}_{j+1}}$. The rule of Ba{\l}aban, by which the form ``couples only the fields $A_j$ in the same component'' (119, p. 258), fails to be local: whether two bonds belong to a common component is settled at a distance from them, and two histories having the same trace can differ on it. The kept form is not a Schur complement but a truncated series. Its lower bound is therefore not inherited and has to be proved. The proof localises to small cubes (Ismagilov--Morgan--Simon) and compares with a genuine Schur complement. \emph{Carriage}: in an integral operation that a term carries, the tests, the gauge factors and the delta functions depend on no background, whereas the forms $\Delta^{\mathrm{loc}}_{\mathfrak{k}}$ and the factors $z^{(j)}$ do. These forms and factors are given as functions of a pair $(\mathbf{U},\mathbf{J})$ near $\mathfrak{k}$, and they are constructed once and for all. They are evaluated at the current background, which in a term of level $k$ is $U_k(\mathbf{V})$. The background changes with every later step and with every stage of $\mathbf{R}$. Each change produces differences 

\begin{gather*}
\FitMath{-\tfrac12\langle A',[\Delta^{\mathrm{loc}}_{\mathfrak{k}}(U^\sharp)-\Delta^{\mathrm{loc}}_{\mathfrak{k}}(U^\flat)]A'\rangle}
\end{gather*}

, where $U^\flat$ is the background after the move and $U^\sharp$ the one before it, both written in the new variables. Localisation turns these differences into boundary terms. For a piece of a kept form whose carrier is $R$, the oscillation with respect to the pair is bounded by $Cp_0(g_j)^2e^{-\kappa d^\natural_{j+1}(R)}$, measured in the hull length of the carrier; this bound is large, being the square of a field at its threshold. The difference of such a piece at the two backgrounds of a stage of level $i$ is localised as well, and each of its localised pieces, on a domain $Y$, has modulus at most $C(\log x_i)^{-(q-3\mathfrak{p}_0-Cr)}e^{-c\,d^\natural_k(Y)}$. This is a negative power, because $q$ is chosen after $\mathfrak{p}_0$ and $r$, with $q>3\mathfrak{p}_0+Cr$.

\paragraph{The comparison at an admitted real reference.}
Fix the step and use level-\(j\) units. Let \(\mathsf A\)
be the complete free-variable precision of the step after
eliminating the pivot bonds, including the projected gauge
and current terms. The free-variable space is the space
on the nonpivot bonds meeting \(\Omega_{j+1}\).
The integration set \(I\) consists of the free bonds meeting
\(\Lambda_{j+1}\). Equivalently it is the set of free bonds
meeting the blocks of the level-\((j+1)\) sites in
\(\Lambda_{j+1}\). The stored set \(E\) is its complement
among the free bonds of the step. In particular a free
bond with one end in the integration region is assigned
to \(I\).

Take a real reference in the admitted class of the
fluctuation theorem, on the entire determining sequence,
with these conditioning sets.
The complete precision has a lower bound \(\gamma\,\mathrm{Id}\)
there, with \(\gamma>0\) fixed. The Schur complement
\[
 S=\mathsf A_{EE}
            -\mathsf A_{EI}(\mathsf A_{II})^{-1}\mathsf A_{IE}
\]
then satisfies \(S\ge\gamma\,\mathrm{Id}\). Indeed, its quadratic
form is the minimum of the full quadratic form over the
integrated coordinates with the stored coordinates fixed.
For a stored component \(D\) the compression of \(S\) to
its free bonds therefore has the same bound. When the
integration set is empty the second term of \(S\) is zero.

We can also control the difference between this
compression and the chain sum assigned to \(D\).
For the precision and for the covariance
\((\mathsf A_{II})^{-1}\), the absolute sum of the
length-\(m\) walk kernels is bounded entrywise by
\[
              C\,2^{-m-1}e^{-\delta|b_--b'_-|_1}.
\]
Their absolute row and column sums are at most
\(C'2^{-m-1}\), since the endpoints range over the unit
lattice in these units. For the three factors of a
Schur chain, sum first over the internal bond indices.
The resulting row and column bounds are at most the
products of the three bounds. At total length \(n\)
the number of choices of the three lengths is
\((n+1)(n+2)/2\). The tails above length \(r\) thus
have absolute row and column sums at most \(C''e^{-cr}\).
The same bound for a single precision walk follows
directly from its length estimate.

Each walk cell has a carrier of bounded diameter in
these units, with a bound fixed by the walk construction.
The carrier of a nonzero chain of total length \(n\)
has diameter at most \(C_{\rm cell}(n+3)\), since successive
carriers meet at their common bond. Let \(W\) be the
side of a level-\((j+1)\) large-field cell, measured in
level-\(j\) units. Use the closed cells for the retention
test, with a common corner counting as contact. Disjoint
closed cells of this grid are at sup distance at least
\(W\). Thus if the carrier has sup diameter less than
\(W\), all the cells it meets touch one another. A chain
contributing to the compression on \(D\) meets \(D\),
since its carrier contains both endpoints of each of its
exterior bonds. The cells it meets in the stored strip
then form a connected nonempty set in \(D\), and the
other cells it meets are at most one touching layer
away. Such a chain is retained for \(D\).
Thus every omitted term in this compression belongs
to the length tail \(n\ge W/C_{\rm cell}-3\).
The tail estimate and the matrix Schur test give
\begin{equation}\label{hi:retained-global-tail}
 \bigl\|\Delta_D^{\mathrm{loc}}-S|_D\bigr\|_{2\to2}
                  \le C e^{-cW/C_{\rm cell}} .
\end{equation}
When this error is at most \(\gamma/2\), the kept
form at this reference is bounded below by
\((\gamma/2)\,\mathrm{Id}\). This argument applies in the full
real reference class and on the literal integration
and stored coordinate sets just described.

To use this bound when the pair is given only near
the stored component, we must compare on small cubes
with admissible real references on which the same
short chain terms are defined. Such a comparison must
match the region traces, the integrated free bonds,
and the real pair on the carriers, in a common frame
for the retained coordinates. The walk locality
identities then match the short terms. The estimate
\eqref{hi:retained-global-tail} controls length tails
in an admitted comparison. If \(K\) is the kept matrix in the retained coordinates and
\(\sum_\nu\chi_\nu^2=1\), the kernel of
\(K-\sum_\nu\chi_\nu K\chi_\nu\) is
\[
 \tfrac12 K(b,b')\sum_\nu(\chi_\nu(b)-\chi_\nu(b'))^2.
\]
For cutoffs of bounded overlap varying on length \(h\), the
factor multiplying \(K(b,b')\) is at most
\(C h^{-2}d(b,b')^2\). The exponential kernel bounds give
uniform row and column sums with this second distance moment.
The matrix Schur test bounds the error by \(C h^{-2}\).
Thus patch lower bounds with a fixed margin give the complete
lower bound once the comparisons on the actual localized
histories have been constructed in the same coordinates.

\paragraph{Comparison in complex coordinates.}
The precision and inverse walk terms are holomorphic in the
background and current restricted to their carriers, on the
piecewise local classes of the fluctuation expansion. Two
realizations give the same coefficient for a decorated walk
with the same cell families and parent map if the region and
fluctuation-structure traces agree on its common padded hull
\([\![\omega]\!]\), at every level. Match the actual free-coordinate
and conditioning masks on all their reads in that hull.
The first-member trace then gives the same field carrier
\([\omega]=[\![\omega]\!]\cap\Omega_0\). The pair on this carrier
must be the same, in the same coordinates, and belong to both
local classes. This applies factor by factor to a Schur chain.

Suppose the actual later read and its comparison are expressed
in one chart about a coercive real reference. Normalize the
chart to radius \(1\) in the background and current norms.
If every carrier restriction belongs to its local class
throughout the chart and the absolute walk sums obey the uniform
row and column bounds with summable length tails there, the kept matrix \(Q(z)\) is
holomorphic and uniformly bounded in operator norm. Cauchy's
estimate along the radial line to \(z\) then gives
\(\|Q(z)-Q(0)\|_{2\to2}\le C\theta\) for \(\|z\|\le\theta\le
\frac12\). If \(Q(0)\ge c\,\mathrm{Id}\), it follows that
\(\operatorname{Re}(a^{\mathsf T}Q(z)a)\ge
(c-C\theta)\|a\|_2^2\) on the real retained coordinates.
For the localized histories we must construct the real
comparisons in these coordinates and verify the common chart
with a uniform fraction of the available radius for the
prepared and substituted data, including the independent
current. The existence of each coefficient on its open
carrier class alone does not give that fraction.

\paragraph{A real comparison by cutting off the potential.}
Keep the separated determining sequence of the step which
formed the strip, including its original free-coordinate
set and its integration mask. Work in level-\(j\) units.
On an enlarged comparison cube about a retained bond
suppose there is a real gauge in which the background is
\(e^{i\eta_j A}\), with \(|A|\) and
\(|\nabla^{\eta_j} A|\) at most \(a\), and with the
transported independent real current bounded by \(a\) in
absolute norm in level-\(j\) units. Take the cube and its
whole stencil collar strictly inside \(\Omega_j\), so that
only layers \(j,j+1\) occur there in the forming sequence.
The retained form and its comparison must use the same
pivots and the background-relative contour quadratic
term of the fluctuation operator. In that prescription the real gauge
acts on free coordinates by the adjoint at the initial
endpoint, an orthogonal change in the counting pairing.

Choose a real cutoff \(\chi\) equal to one on the middle
cube and zero in a neighborhood of the outer boundary,
with first difference bounded by \(C/h\) for a cube
of side \(O(h)\). Extend \(\chi A\) and \(\chi J\) by
zero in this gauge on the whole original torus.
The first difference of \(\chi A\) is at most
\(Ca(1+h^{-1})\). In a plaquette product the linear
sum of its bond logarithms is \(\eta_j^2\) times
the discrete curl of \(\chi A\), so its norm is bounded
by \(C\eta_j^2a(1+h^{-1})\). The product remainder is
bounded by \(C\eta_j^2a^2\). Thus the real reference
\((e^{i\eta_j\chi A},\chi J)\) has plaquette defects
at most \(C\eta_j^2(a+a^2)\), including at the cutoff.
It has a single small logarithm on the torus. The scaled
potential and its first difference cost at most a further
factor \(L^2\) on the adjacent upper layer. The independent
current costs at most \(L^3\) there. Outside the collar the
reference is flat with zero current. Choosing \(a\) small
after the comparison width makes this pair satisfy the
real curvature, gauge and current conditions of the
fluctuation theorem on the original geometry. The current
can be cut independently because that theorem takes a real
current as a separate argument.

Extend the local site gauge arbitrarily on the torus and
transform the reference back. The real pair conditions
are invariant under this gauge action, so the reference
is still admitted and agrees with the actual read on the
middle cube in its original frame. The global reference uses the same
region traces and conditioning masks there. Every chain
whose endpoints are tested in the inner cube and whose
length is less than its buffer to the middle boundary
divided by \(C_{\rm cell}\), up to the fixed additive
constant, has its carrier in the common region. Its terms
agree by carrier locality. On the local read the absolute
walk estimates must hold on all carriers of the kept form,
including the long ones. With these estimates the remaining
length tails in both forms have operator norm at most
\(C\exp(-ch/C_{\rm cell})\). Choose \(h\) to make these
errors and the partition error smaller than the fixed
reference gap, with the cutoff support fitting inside
the original collar. Thus this cutoff supplies a
positive real comparison on the patch once its gauge,
mask and current conditions hold. The partition estimate
above combines such comparisons with constants uniform
in the depth. For a later read, the two-layer bounds must
be verified in the forming sequence and in its retained
frames, including the enlarged cubes about the external
bonds. These comparisons use the reference-weighted
form. A change in the ratio of physical and reference
weights uses the gauge compensation in
\eqref{hi:compensated-quadratic}.

\paragraph{Reading a retained strip in its original domain.}
\hypertarget{hi:retained-strip-domain}{}
Let $D$ be a component of $\Omega_t\cap Z_t$ as formed by step
$t=j+1$. A later record retains this strip only if no intervening
healing, island operation or re-made level of a drop-out has
replaced its field. Write $D^{\sim}$ and $D^{\oplus}$ for its
one- and two-cell enlargements on the level-$t$ partition.
For a retained strip the trace on $D^{\oplus}$ is unchanged in
the following sense. The level-$t$ region is the one formed by
step $t$ and the regions of levels $1,\ldots,t-1$ contain $D^{\oplus}$.
This applies to the determining and action readings of level $t$,
with every stage index. The local structure has no radius lift
there and uses only the indices $t-1,t$. The retained factors use the pivots, integration
mask, and gauge of step $t$. A factor of the quadratic form reads
the field $\mathcal A^{\rm s}_{t-1}$ in these same coordinates.
The associated elimination factor depends only on the background.

Take $P=(U,A',\mathbf J)$ in the structured space of the
level-$t$ reading on an address containing the full carriers
of the read. Use the pointwise and gauge clauses of
\cref{rg:structured-local-gauge}, including layer zero at
the first strip. At its formation this is the reading on either side
of the step-$t$ operation that retains the strip. For a later label
use the level-$t$ reading for the retained factor on its original
carrier. Use
the larger of the running radii at $t-1$ and $t$ for each
quantity and set $\bar a_i$ to their uniform upper bounds over
the chosen coupling range, for $i=0,1,2$. The third radius is
for the independent current. The pointwise bounds on curvature,
potential and current in these two layers cost at most
$L^3(1+\beta_{\rm s})\bar a_i$ in the fixed strip domain. Indeed
the pointwise clause of the layered triple assigns a bond,
plaquette or difference by its base point. The strip domain may
assign it by another of its sites. Consecutive boundaries are
separated by more than the diameter of such a stencil, so the
two local spacings differ by at most $L$.
The largest power in these norms is the third, for the current.

In level-$(t-1)$ units the level-$t$ cells have side
$LM\check R_t$. The gauges required in the strip domain are
those on the cubes of the walk expansion, of width of order
$M_1$ in their local units, and on the enlarged tiny cubes,
of width of order $M/L^2$ in level-$(t-1)$ units. The latter
are required at the endpoints of the external bonds as well as
of the integration bonds. A required cube contained in the upper
region fits in a gauge cube of width of order $M$ in level-$t$
units inside the upper trace. One can place the containing cube
within the block intervals met by its coordinate projections.
For a required cube reaching the lower side of the boundary,
the containing cube has width of order $M$ in level-$(t-1)$
units. Its collar lies in $D^{\oplus}$ and hence in the
level-$(t-1)$ region. The stipulated hierarchy of the walk
and cell widths gives room for both choices. These are cubes
of the gauge clause of the layered triple, whose bound for
$U=e^{i\eta A}$ in that gauge is
$\xi_n|A|,\xi_n^2|\nabla A|\le CM\bar a_0$. Converting the
gauge estimate to the other layer costs at most $L^2$.
At $t=1$ the level-zero lattice has unit spacing in these units.
Axial transport in its local cube gives $|A|\le C\,\mathrm{side}\,\bar a_0$
from the plaquette bound, and the unit-spacing first difference
is bounded by twice the potential bound.

Choose the fixed radii $a_i^{\rm loc}$ inside the
positive strip class $\mathcal O^+$ of the original
kept operator. Its gauge radius includes the fixed
small radius of the walk cubes. The remaining
contributions to the radii tend to zero with the
coupling ceiling. It suffices to impose
\[
  L^3(1+\beta_{\rm s})\bar a_i<a_i^{\rm loc}\quad(0\le i\le2),
  \qquad CL^2M\bar a_0<a_0^{\rm loc}.
\]
For an external elimination bond also require the small
averaging chart on its full read set. The two-layer potential
and gauge bounds give a chart size $C(L,M)(\bar a_0+\bar a_1)$
in the upper units, so the separate condition is
$4C(L,M)(\bar a_0+\bar a_1)<r_{\rm P}(L)$, the pivot chart radius.
The floors on $M$, including that required by the decay rate
of the kept chains, precede the choice of these coupling ceilings.

These checks admit the read to the fixed domain of the original
kept operator and of its elimination factors. On the positive
strip class the kept-form estimate bounds the real part in its
original free coordinates, before multiplication by a coupling
ratio. In a real comparison
the potential may now be cut off in the gauges just obtained.
Choose the comparison width $h$ with its interaction collar inside
the two-cell trace. If $\gamma$ is the original real reference gap,
take the error in \eqref{hi:retained-global-tail} at most $\gamma/2$
and the sum $C e^{-ch/C_{\rm cell}}+C h^{-2}$ of the comparison
and partition errors at most $\gamma/4$. The free bonds and masks
have not changed, so this leaves a lower bound $\gamma/4$ for the
old conditioned form in its original coordinates. On its admitted
complex domain, the walk and elimination coefficients are
holomorphic. For the actual substitutions we must verify the
stated layered bounds on these carriers. For a normalized Gaussian
with a changed coupling use the compensation in
\eqref{hi:compensated-quadratic}.

\paragraph{The amplified substitution on carried strip factors.}
Let $n$ be the level of the main step and $j+1<n$ the level
at which a retained strip was formed. Consider a quadratic or
elimination factor of that strip which is carried through step
$n$. The carriage rule puts its carrier $X$ outside the one-layer
enlargement of the region being healed. The carrier is in the
exterior region of the healing, and the structure in which the
factor is read counts no stage in a component meeting its sites.
These are the served targets of the main step. The strip
coordinates and the pivots on $X$ are those at its formation.

Let $P=(U,A',\mathbf J)$ belong to the complex space of the
output structure on a saturated localization $Y$ containing the
enlarged carrier and the response collars. For the admitted
inner variables $\xi$ use the formed inner data
$\varsigma[P;\xi]$ of the main step. The belt map gives the
reduced increment $\beta=\mathfrak b^Y(s;P,\varsigma[P;\xi])$
of the row averages on the belts, with
$\|\beta\|_\infty\le\varepsilon_\beta/2$ throughout the complex
cube polydisc $|s(\Delta)|<e^{\bar\kappa}$. For each wall
component $w$ of the localization use the response with cube
parameters $s\mathbf 1_w$ and form
\[
 (\mathbf A,\mathbf K)=
 \sum_w\lambda_w(\mathbf H,\mathsf j)(s\mathbf 1_w;\beta)
                         [\mathfrak p^{\rm e}(P)].
\]
Here $\mathfrak p^{\rm e}(P)$ is the restriction to the local
exterior response class, with the current of that class.
Level this shift on the exterior pillars with the pins at the
sites of the stored data. Write $\mathfrak h^\delta$ for the
result restricted to the bonds of $X$. The argument of the
retained factor is $(P|_X)\dotplus\mathfrak h^\delta$, with
the potential composed by the group product and the current
increment added. Its stored fluctuation argument is held fixed.

The response estimate gives decay from the belts for the
potential, its first difference and the current in their local
units. The response of a wall is supported inside that wall,
so independent multipliers cost $\max_w|\lambda_w|$. Pinning
increases the bound by a fixed factor $C_\delta$. The wall
carriers lie strictly inside $Y$, so the shift vanishes on
the target boundary outside the interior of $Y$. Put
$\alpha_{*,n}=\min_{\iota=0,1}\alpha_{\iota,n}$ and let
$\Sigma$ denote the belts. The bound in the profile used by
the target clauses is
\[
 |\mathfrak h^\delta|_b\le
 c_{\rm sh}\alpha_{*,n}\,C_\delta a(X)\max_w|\lambda_w|
             e^{-\delta_1d_{\bar{\mathbf B}}(x_b,\Sigma)}.
\]
Here $|\cdot|_b$ takes the maximum of the potential, its
first difference and the current, multiplied by the first,
second and third powers of the local spacing, respectively.
The constant $c_{\rm sh}$ is the fixed profile allowance in
the substitution bound. The gap of a lower clause with
index $m'$ absorbs the profile in the units of that clause.
In particular, if $m$ is the exterior index at the base point,
the gap comparison for $m'\le m$ is
\[
 L^{-(m-m')}\alpha_{*,n}
 e^{-\delta_1d_{\bar{\mathbf B}}(x_b,\Sigma)}
 \le C e^{-c\delta_1M\check R_n}\mathsf{gp}(m'),
\]
where $\mathsf{gp}(m')$ is the gap supplied by the descent of
the layer bounds and the fixed gap constant is included in $C$.
With this profile the substitution into the served target
preserves its tower, potential and one-scale clauses whenever
$C_\delta a(X)\max_w|\lambda_w|<1$. The free rim and wide
shell used in that comparison hold for both classes of the
main step. The stored sites of a served target are outside
the boxes, so its pinned transports are among those admitted
by the comparison.

To bound $a(X)$, split the response decay between the
pointwise profile and the distance from the belts to $X$.
A path crosses a distance at least $cM\check R_n$ in
level-$n$ units before it can reach the exterior. It then
crosses the intervening layers, each of scaled width at least
$\rho$. If $n_X$ is the largest level for which $X$ meets
the interior of the healed exterior region $\bar\Omega_{n_X}$,
this gives
\[
 2a(X)\le
 \frac{C e^{C\bar\kappa}\varepsilon_\beta}{\alpha_{*,n}}
 \exp\bigl[-c\delta_0M\check R_n
           -c\delta_0\rho(n-1-n_X)_+\bigr].
\]
Thus for $a_X=2C_\delta a(X)$ the whole disc
$|\lambda|<2/a_X$ is available when every $\lambda_w$ equals
$\lambda$. Since $\check R_n$ grows at least quadratically in
the logarithm of the coupling parameter, the belt factor
absorbs $\alpha_{*,n}^{-1}\lesssim g_n^{-1}$ uniformly in $n$.
In particular $a_X<1$ at small coupling.

On these discs the substituted triple lies in the layered
space of the factor on its original carrier. The strip-domain
comparison above then applies to its holomorphic walk
coefficients and elimination factors in the original masks.
At $s=0$ or $\lambda=0$ the argument is $P|_X$. At the physical
point $s=\mathbf1$, $\lambda=1$ the reduced belt datum is the
true one. The pinned response gives the pair of a single
representative of the prepared critical orbit, restricted to
each served carrier. Its transports agree with those of the
healed background at the stored sites. This proves the
amplified read for the carried strip factors. Factors read
inside the healing require the localized inner substitutions.

\textbf{The idea of $\mathbf{R}$} (Comm. Math. Phys. 122 (1989) 175, (0.3)--(0.4)). We split $Z=Z'\cup Z''$. Following that paper (p. 176), $Z''$ is the union of the components ``for which the small factors connected with large field control some number of next steps'', and $Z'$ is the union of all the other components (the rule by which the present construction sorts the components into these two sets is given under ``Where it acts'', below). The density is a sum $\rho(V)=\sum_Z\rho(Z,V)$ over the large-field regions $Z$ ((0.2) of that paper), and $\rho(Z'',V)$ denotes the density belonging to the region $Z''$. Then

\begin{gather*}
\FitMath{(\mathbf{R}\rho)(V)=\sum_Z\rho(Z'',V)\,\frac{\int dV|_{Z'}\,\rho(Z,V)}{\int dV|_{Z'}\,\rho(Z'',V)},} \\[3pt]
\FitMath{\int\mathbf{R}\rho=\int\rho .}
\end{gather*}

In $Z'$ the law of the field is replaced by the small-field law, and the mass is preserved. This idealised form is not linear in $\rho$. The form actually used is linear, because its reference law contains the Wilson weight and nothing else.

\textbf{Where it acts.} Let $k_\flat=k-\mathsf{N}_k\ge1$ (\emph{depth}); this restriction excludes nothing at $k=\mathsf{N}_k$, since by the second condition below a member of the class is then a component of $Z_k$ in which no step has created anything, and no such component exists. A component $X$ of $Z_k$ belongs to the class when three conditions hold. First, for $k_\flat\le i\le k$ the set $X\cap Z_i$ is a single component of $Z_i$ and is small, that is, contained in a cube of $O(1)$ level-$i$ cells (\emph{single boxes}, and at $i=k$ \emph{size}). Second, for $k_\flat<i\le k$, step $i$ creates nothing in it, and it does not absorb a second component, $X\cap Z_i=\Phi_i(X\cap Z_{i-1})$ (\emph{quiet}). Third, it has not been acted on by any operation whose level lies in $[k_\flat,k)$, the lower end point included (\emph{quarantine}). Ba{\l}aban imposes two conditions, that the component be small and quiet (ibid., p. 177); to these two, the exclusion of a joining is added here as a separate condition, named explicitly. Two small quiet boxes that join within the window form a small component without creation, yet its trace one scale lower has two components: a joining does not count as a creation. Let $H_i$ denote the hull, in level-$i$ cells, of $X\cap Z_{i-1}$, and let $B^{[r]}$ be the box $B$ with each face pushed outward by $r$ such cells. Then

\begin{gather*}
\FitMath{\Omega_i^c\cap X=H_i^{[8]},} \\[3pt]
\FitMath{Z_i\cap X=H_i^{[10]}} \\[3pt]
\FitMath{(k_\flat<i\le k):}
\end{gather*}

the boxes are nested, the two of each level being concentric, and the number of cells along a side of $X$ is comparable to the fixed point of the free evolution \cref{sec:large-field-rare-patches-gains-per-creation-and-costs-per-site}. As long as $k\le\mathsf{N}_k$ the class is empty, so $\mathbf{R}$ does nothing at the first levels. For $k>\mathsf{N}_k$ the class may be non-empty: the component obtained by free evolution from a single isolated cell created at step $k-\mathsf{N}_k$, with no intervening creation, joining or operation before level $k$, belongs to it. Up to level $k$ the torus is admissible, $\tfrac12\ell L^{K-i}\ge CM\check{R}_i$ for $i\le k$, as the opening of \cref{prop:large-field-2-1} requires. For a first-passage datum, $x_K\le g^{-2}+B_\sharp$, this condition follows from $\ell\in\mathfrak{L}(g)$ by way of $\check{R}_i\le L^{K-i}\check{R}_K$, a consequence of the upper flow bound; it is a lower bound on $\ell$. Hence no cell, no box and no neighbourhood of a box wraps around the torus.

\paragraph{Restriction to a component.}
Let \(C\) be a union of components of the final region \(Z_i\)
of a step, and let \(Y\) be its input region. The hull, layer,
and boxing operations in this step restrict to \(C\). Indeed,
each cube of the hull of \(Y\) meets an input component. Their
union is connected and belongs to \(Z_i\), so the cube and the
input component belong to the same final component. The enlargement
of a cube by any of the layers used in the recursion is also
connected and lies in \(Z_i\). It therefore belongs to \(C\) precisely
when that cube belongs to \(C\). Finally, when a small component
is replaced by its box, the box lies in \(Z_i\) and is connected
to the component it replaces. The replacement can thus be made
before or after restricting to \(C\). Applying this argument at
each boxing pass and test round shows that the restricted output is
the one obtained from \(Y\cap C\) and the recorded failures in \(C\).
In particular, if \(C\) contains no creation of that step, the output
on \(C\) is given by the eight-layer hull and by boxing the ten-layer
hull of \(Y\cap C\), respectively.

\paragraph{Tracing an eligible component.}
Write \(\mathcal S'_s\) and \(\mathcal S_s\) for the labels just
before and after the operation of level \(s\). We use the nesting
of \cref{prop:large-field-2-3} on these labels. At a fixed earlier
index \(i\), the regions \(Z_i\) decrease as these labels advance
from \(\mathcal S'_i\). A step leaves earlier entries alone, and
the three rules for an operation only enlarge their integrated
regions. Fix \(h=k_\flat\le i\le k\) and a cell of
\(G_i=Z_i(\mathcal S'_k)\cap X\). For \(i\le s<k\) the cell
lies in a component of \(Z_s(\mathcal S'_s)\). Healing that component
would remove the cell from \(Z_i\). Making it an island would
have the same effect if \(i<s\). In the remaining cases an operation
on that component would give it a mark. If the mark ceased to stand,
a later healing or island would have integrated the cell at level \(i\).
If it still stands, the marked component lies in the component \(X\)
that contains the cell. This contradicts the exclusion of marks at
levels \(h,\dots,k-1\). Thus no intermediate class component
contains the cell. Its memberships in the regions at \(i\), and
at \(i-1\) when \(i\ge2\), are the same in \(\mathcal S'_k\)
as in \(\mathcal S'_i\).

The entire component of \(Z_i(\mathcal S'_i)\) through that cell
survives in \(Z_i(\mathcal S'_k)\). To see this, at each intermediate
operation the whole component still lies in the same component of
the top region as the chosen cell, by nesting and connectedness.
The operation leaves that top component alone. Induction proves the
assertion, and the surviving component is contained in \(X\).
Hence \(G_i\) is a union of components of the original output of
step \(i\). For \(i>h\), the step's input restricted to \(G_i\)
is \(G_{i-1}\).
The absence of creations for \(h<i\le k\) now lets us use the
preceding restriction argument at each of those steps. All these
steps have \(i\ge2\), since \(h\ge1\). The argument allows the
first step's creation set to be defined separately.

For clarity, measure lengths in bare units and set
\(a_i=M\check R_iL^i\), the side of a level-\(i\) cell.
For a box \(B\), write \(B^{[s]}\) for its enlargement by a
distance \(s\) in the sup norm in its lift. A negative \(s\)
means moving each face inward by \(|s|\). The oldest component \(G_h\) is a nonempty box by the
small-component boxing rule and the survival just proved.
Inductively its successor is a box. Its hull \(H_i\)
on the next grid satisfies
\(G_{i-1}\subset H_i\subset G_{i-1}^{[a_i]}\), by rounding the
endpoints in each coordinate to the level-\(i\) grid. A box
is fixed by the boxing operation. Thus the two restricted
outputs have the form
\begin{equation}\label{eq:elab-history-quiet-boxes}
 O_i:=\Omega_i^c\cap X=H_i^{[8a_i]},
 \qquad
 G_i=H_i^{[10a_i]},
 \qquad G_k=X.
\end{equation}
Here and below \(\Omega_i,Z_i\) without a label refer to
\(\mathcal S'_k\). Each \(H_i\) is nonempty by eligibility,
which requires a predecessor through the waiting window.
In particular, \(X\) has at least \(21\) cells along each side.
With the small-box cutoff of \(100\) cells per side, it has
at most \(100\).

\paragraph{The boxes used in preparation.}
Let \(u_i=ML^i\) be the side of an \(M\)-cube at level \(i\) in
bare units. In the scale choice used for the preparation, let
\(k_0+1\) be the last index with \(a_{k_0+1}=u_k\), with
\(h<k_0\le k-2\). Thus \(a_i\le u_k\) up to \(k_0+1\), and
\(a_i\ge Lu_k\) thereafter. The construction uses the following
distances from the original small-field region. For \(h<i\le k_0\)
take distance \(5a_i\) from \(\Omega_i\). For \(k_0<i\le k\)
take distance \(5u_k+w_i\) from \(\Omega_{k_0+1}\), where
\[
 w_{k_0+1}=2u_k,\qquad w_k=0,\qquad
 w_i=\sum_{q=i+1}^k u_q\quad (k_0+2\le i<k).
\]
In each case \(Z''_i(X)\) is the union of cells of \(X\) all of whose
points have at least the indicated distance. The change of scale
allows the later boxes to stay on the scale \(u_k\).

For \(0<s\le8a_i\), the points of \(X\) at distance at least \(s\)
from \(\Omega_i\) form \(O_i^{[-s]}\). The interior of \(O_i\)
is disjoint from \(\Omega_i\). Every face of \(O_i\) lies in the
closure of the part of \(\Omega_i\) between \(O_i\) and \(G_i\),
since \(G_i\) extends two cells beyond \(O_i\) in every direction.
Thus the distance inside \(O_i\) is the least distance to one of
its faces, which gives the erosion formula. It follows from
\eqref{eq:elab-history-quiet-boxes} that
\begin{equation}\label{eq:elab-history-prepared-boxes}
 Z''_i(X)=
 \begin{cases}
  H_i^{[3a_i]},& h<i\le k_0,\\
  H_{k_0+1}^{[3u_k-w_i]},& k_0<i\le k .
 \end{cases}
\end{equation}
The first box is a union of level-\(i\) cells. For the second,
the faces of \(H_{k_0+1}\) lie on the grid with side \(u_k\).
Since \(L\) is odd, those faces also lie on the grid with side
\(u_i\), and \(3u_k-w_i\) is a multiple of \(u_i\).
The second box is therefore a union of \(M\)-cubes of level \(i\).
These grid facts justify passing from the set of points at the
specified distance to the union of cells in the definition.

The earlier boxes satisfy \(H_i\subset Z''_i(X)\subset O_i
\subset G_i\). The next hull contains \(G_i\), so the nest increases
through the change of scale. At the later levels the centre is
fixed, and \(w_i\) decreases from \(2u_k\) to zero. The boxes remain
inside \(O_{k_0+1}\), which itself is contained in all later \(O_i\).
Thus \(G_h\subset Z''_{h+1}(X)\subset\cdots\subset Z''_k(X)\).
The separation from \(G_h\) to the exterior of \(Z''_{h+1}(X)\)
in \(X\) is at least \(3a_{h+1}\). At early stages, \(Z''_i(X)\)
is separated from the exterior of \(Z''_{i+1}(X)\) in \(X\)
by at least \(7a_i\), since \(G_i=H_i^{[10a_i]}\) is already
inside the next box. For \(k_0<i<k\), the separation is
\(w_i-w_{i+1}\). In particular the late separations have the
scale of the \(M\)-cubes, which explains the smaller margins
used for the prepared backgrounds.

With the cutoff of \(100\) cells per side,
\(G_{k_0+1}=H_{k_0+1}^{[10u_k]}\) has sides at most \(100u_k\).
Hence the sides of \(H_{k_0+1}\) are at most \(80u_k\), and
those of \(Z''_k(X)\) are at most \(86u_k\). The central
integration below adds one more layer of \(M\)-cubes.
Its box therefore has sides at most \(88M\) in level-\(k\)
units.

\textbf{The operation.} (0) \emph{The frame sum}, described first although it belongs inside (3). In each component the identity $1=|W_4|^{-1}\sum_{r\in W_4}1$ is inserted after the integration over the box and before the tree gauge of (3). Consequently no frame sum is attached to a component that leaves at (2) or during the stages. In the summand of $r$, every object of (3), (4) and \cref{candidate:large_field} that belongs to $X$ is formed in the frame $r$: the reference law with its denominator, the two trees, the test on $B$, and the last-shell test. If a component leaves after the tree gauge, its frame sum stays with it, and an island keeps, in each summand, the cut-off and the delta function belonging to the frame of that summand. Formed instead in a single frame $r_0$, these make the bracket of $X$ integrate to $|W_4|^{-1}\sum_r\mathfrak{a}_r\,\mathsf{D}^{(r_0)}_X/\mathsf{D}^{(r)}_X$ instead of $|W_4|^{-1}\sum_r\mathfrak{a}_r$. Here $\mathfrak{a}_r$ denotes that part of the summand of $r$ which is independent of the box, and $\mathsf{D}^{(r)}_X$ denotes the denominator of $\mathfrak{s}_X$ formed in the frame $r$. The reason for keeping the frame sum is that neither the test on $B$ nor the last-shell test is gauge invariant. Each of them enters only by way of decompositions of unity performed within a single gauge-fixed measure. Let $\mathsf{m}$ be the measure before the tree gauge, and let $\mathsf{a}_r,\mathsf{b}_r$ be its parts that pass, respectively fail, these tests in the frame $r$. Then $\mathsf{a}_r+\mathsf{b}_r=\mathsf{m}$ for every $r$, and hence $|W_4|^{-1}\sum_r(\mathsf{a}_r+\mathsf{b}_r)=\mathsf{m}$. With the drop-outs kept instead in a single fixed frame $r_0$, the total is $\mathsf{m}+(|W_4|^{-1}\sum_r\mathsf{a}_r-\mathsf{a}_{r_0})$. The average is taken so as to obtain covariance under the lattice isometries; trees break this covariance. (1) \emph{Stages}: for $j=k_\flat,\dots,k-1$, one value after another, the stored variables of the outer shell of $Z_{j+1}\cap X$ are integrated out as in a small-field step. Before this, tests at the sharper threshold $\delta'_j:=g_jA_1p_1(g_j)$ are made, formed with backgrounds localised in $X$. What remains is a box $\Lambda^{\mathrm{pr}}$ of side at most $CM$: it is the set $Z''_k(X)$ of the nest of \cref{sec:large-field-the-tests-the-nested-regions-and-the-content-of-a} enlarged by one more layer of $M$-cubes, and this layer is the shell of level $k$. In it the levels $k_\flat,\dots,k$ are stacked in thin shells, the shell of level $j$ carrying the field $V_j$. Together these fields form the datum $V''$ attached to the prepared determining set $\mathbf{B}''$; its background is $U''_k$. At the innermost position lies the core $C_{k_\flat}=X\cap Z_{k_\flat}$, whose side, measured in the unit of level $k$, is at most $CM\check{R}_{k_\flat}L^{-\mathsf{N}_k}$. (2) The test asks whether the exterior field $V_k$ admits an extension into $\Lambda^{\mathrm{pr}}$ whose curvature is less than $2\varepsilon_k$. If an exterior field satisfies $|\partial V_k-\mathbf{1}|<\varepsilon$, then it does admit an extension, one with $|\partial V_k-\mathbf{1}|<CM^2\varepsilon$. To see this, set the field equal to $\mathbf{1}$ inside after choosing an axial gauge on the surface $\partial^+\Lambda^{\mathrm{pr}}$. Hence ``no regular extension'' means that some exterior plaquette has $|\partial V_k-\mathbf{1}|\ge cM^{-2}\varepsilon_k$, and this pays as in \cref{sec:large-field-rare-patches-gains-per-creation-and-costs-per-site} (p. 193 there). (3) $V_k$ is integrated over $\Lambda^{\mathrm{pr}}$. Taken together, the parts from the test (2) up to the division into the two classes \cref{candidate:large_field} form the \emph{main step} of the operation; the stages that precede it and the exponentiation that follows it are counted separately. This integration is the first move to alter the density. Integrals against any bounded $F$ that is independent of $V_k$ on the box are left unchanged by it. Level by level, from the top down, it resolves the bare chain $\prod_j\int dV_j\,\delta(\mathcal{M}(V_j)V_{j+1}^{-1})$, $\int dV_i(c)\,\delta(\mathcal{M}(V_{i-1})(c)V_i(c)^{-1})=1$. What remains is a free Haar integral,

\begin{gather*}
\FitMath{\chi_k\,\chi_{k,\Lambda}\int dV_{k_\flat}\big|_{C_{k_\flat}}\chi\;\mathbf{T}_{k_\flat}(C_{k_\flat})\int dV''\big|_{\mathbf{B}''\cap(\Lambda^{\mathrm{pr}}\setminus C_{k_\flat})}\,\chi''_k\,e^{A''_k(U''_k)}}
\end{gather*}

Here $\chi_{k,\Lambda}$ is the test of (2); $\chi$ stands for the plaquette tests of level $k_\flat$ carried by the term on $\Omega_{k_\flat}\cap C_{k_\flat}$; and $\chi''_k$, $A''_k$ are the tests and the action that (1) leaves. The integral is regarded as a density that is constant in $V_k$ on the box. The field that is expanded is $V''$. It is expanded around the averages of the minimiser $U_0$ below, each average taken at its own level:

\begin{gather*}
\FitMath{V''=e^{ig_kB}V_0\,\text{ on the fluctuation bonds }\mathbb{B}_0=\mathbf{B}''\cap(\text{shells}),} \\[3pt]
\FitMath{V_0=\big(\mathcal{M}^jU_0\text{ on the shell of level }j\big)_{k_\flat\le j\le k}.}
\end{gather*}

Thus $B$ is a multi-scale field on bonds of every level from $k_\flat$ to $k$, with a single scaling $g_k$ and a single test $|B|<\delta'_k/g_k$. This is the field $B$ in \hyperref[nav:prepared-operation]{the prepared-component operation}. The gauge of this field is fixed along a tree $T_0$ in $\mathbb{B}_0$, whose bonds come from all these lattices, one lattice for each shell. For two nested boxes $P_1\supset P_2$, the tree is made up of axial contours which leave from a single initial site and go around $P_2$. The family of trees has two properties (pp. 195--196 there). First, the family is stable. Suppose one enlarges $P_2$ to a box $P'_2\subset P_1$ and deletes the contour of each site of $P'_2\setminus P_2$; then what is left is the tree of $P_1\setminus P'_2$. Ba{\l}aban describes this property as important, and boxes are needed because of it. Second, the tree of a shell uses up the gauge freedom at every site of the shell except one. For this reason, for each inner box on which a gauge is fixed, one external bond crossing its boundary is added, and one more is added for the outermost box. Through these extra bonds the trees of the shells become a single tree, and it fixes the gauge entirely, with Jacobian one. The core operation $\mathbf{T}_{k_\flat}(C_{k_\flat})$, which contains the failed tests, is at no point unfolded. One further test, the \emph{last-shell test}, compares with the vacuum and constrains $V_{k_\flat}$ between the core and $\mathbb{B}_0$. It requires 

\begin{gather*}
\FitMath{|\partial U_{k_\flat,\square}((\mathbf{1},V_{k_\flat}))-\mathbf{1}|<\tfrac12\varepsilon_{k_\flat}\eta_{k_\flat}^2}
\end{gather*}

 on the test cubes of that region, where $(\mathbf{1},V_{k_\flat})$ equals $\mathbf{1}$ on the side of $\mathbb{B}_0$. (4) In the run with its own choices, \(\theta_k=x_k\).
For each term $\tau$, its mass $m_\tau=\int\tau\,dV_k|_{\Lambda^{\mathrm{pr}}}$ multiplies the reference law. Of the two factors of this product, the level-$k$ variable $V'_k:=V_kV_\Lambda^{-1}$, where $V_\Lambda$ is the element of the critical orbit of \cref{prop:large-field-2-4} that equals $\mathbf{1}$ on the tree of $\delta_{\mathrm{ax}}$, occurs only in the reference law: in its numerator as the argument, and in its denominator as the variable of integration. Let $U_{k,X}$ denote the background localised in $X$, and let $\chi(\Lambda^{\mathrm{pr}})$ be the cut-off $\mathbf{1}\{|\tfrac1i\log V'(b)|<M_0\varepsilon_k$ on the bonds of the box$\}$, in which $M_0$ is a constant chosen large after $M$. Let $\delta_{\mathrm{ax}}$ be the tree gauge of $\mathbb{T}^{(k)}\cap\Lambda^{\mathrm{pr}}$, taken on a tree that contains all the sites of this set together with the external bond of $T_0$ into the box; that bond is part of the tree because, without it, a constant gauge transformation of the box is left unfixed. Let $\phi_X$ be a cut-off function that equals $1$ well inside $X$. Then

\begin{gather*}
\FitMath{\mathfrak{s}_X(V_k)=\frac{\delta_{\mathrm{ax}}(V'_k)\,\chi(\Lambda^{\mathrm{pr}})(V'_k)\,e^{-x_kA(\phi_X,U_{k,X}(V'_kV_\Lambda))}}{\int dV'|_{\Lambda^{\mathrm{pr}}}\,\delta_{\mathrm{ax}}(V')\,\chi(\Lambda^{\mathrm{pr}})(V')\,e^{-x_kA(\phi_X,U_{k,X}(V'V_\Lambda))}}.}
\end{gather*}

If one of the tests of (1)--(3) fails in a component, that component drops out \cref{sec:large-field-the-tests-the-nested-regions-and-the-content-of-a}: it remains live, and the failure counts as a creation. This can happen to the same component repeatedly. Each occurrence counts as a creation and is followed by a wait of at most $\mathsf{N}_k+1$ steps.

\textbf{The purpose of the stages, and the reason for a sharper threshold.} Continuing the field by least action costs a constant factor: for exterior data with $|\partial V_k-\mathbf{1}|<\varepsilon$ one has $|\partial V_\Lambda-\mathbf{1}|\le C(M)\varepsilon$, and $|\partial U_0-\mathbf{1}|\le C(M)\varepsilon_k\eta_k^2$ near the box. For every value of the coupling, the stages reduce the component from side $M\check{R}_k$ to a box of side $O(M)$. Hence $C(M)$ is a constant fixed before $\gamma$. On the shell of level $j$ it is overcome by the factor $L^{-(k-j)}$ by which a level-$k$ field scales when viewed from level $j$. On the outermost shells it is overcome by $\check{R}_k^{-c}$, because $1+\log_L\check{R}_i$ levels are compressed into the box, namely the levels from the last level $i$ with $L^{k-i}=\check{R}_i$ up to level $k$. Consequently $U_0$ belongs to the domain of the carried tests. It also belongs to the analyticity domains of the terms, insofar as the domain of a term lies inside the component or at a distance from the box; the remaining terms are treated in \cref{candidate:large_field}. In the words of 122 (1989) 175, p. 195, ``it is the reason why we have done the preliminary integrations''. As in \cref{lem:large-field-2-2}, the tests made at a stage have to imply the earlier ones. Here this must happen $\mathsf{N}_k=\check{R}_k$ times in succession, and the losses together amount to a fixed fraction. If two backgrounds satisfy the same constraints on a region, their difference is the response to the difference of their data, and this response decays. Let $U^{(n)}_{k,X}$ denote the background after $n$ stages, localised in $X$. Put $j=k_\flat+n$, and let $\beta_{\mathrm{s}}$ be a fixed small fraction, written $\beta$ by Ba{\l}aban. Then, on the support of the tests of stage $n+1$ and on the shell of level $i\le j$,

\begin{gather*}
\FitMath{|\partial U^{(n)}_{k,X}-\mathbf{1}|\le|\partial U^{(n+1)}_{k,X}-\mathbf{1}|\big(1+c\beta_{\mathrm{s}}2^{-(j-i)}\big)+c\beta_{\mathrm{s}}2^{-(j-i)}\varepsilon_i\eta_i^2,}
\end{gather*}

where the error term is 

\begin{gather*}
\FitMath{C\delta'_je^{-\delta_0\operatorname{dist}}\eta_i^2\le C\,\tfrac{\delta'_j}{\varepsilon_j}(1+(j-i)^{1/2})e^{-c(M/M_1)(j-i)}\varepsilon_i\eta_i^2}
\end{gather*}

. For $M$ large, the decay $e^{-c(M/M_1)(j-i)}$ across the shells dominates $1+(j-i)^{1/2}$, and $2^{-(j-i)}$ remains. The ratio $\delta'_j/\varepsilon_j=(A_1/A_0)(\log x_j)^{-(\mathfrak{p}_0-\mathfrak{p}_1)}$ dominates the fixed constants, and at this point smallness of $A_1/A_0$ suffices as well. After $n$ stages the test requires $|\partial U^{(n)}_{k,X}-\mathbf{1}|<\big(1-\beta_{\mathrm{s}}(1-2^{-(j-i+1)})\big)\varepsilon_i\eta_i^2$ on that shell. On the outermost shells, where the collars of the stages overlap, this requirement is relaxed by a factor $L_\flat^2$ for each overlap, where $L_\flat$ is an integer with $2\le L_\flat<\frac12L$ that is chosen together with $L$; such a choice requires $L\ge5$. Let $k_\sharp+1$ be the last of the levels at which an $M\check{R}$-cube, measured in units of level $k$, has side precisely $M$; in formulas, $M\check{R}_{k_\sharp+1}L^{k_\sharp+1-k}=M$. The collars of the stages with $j\le k_\sharp$ are disjoint, and the collars of the later stages overlap. On the shell $i$ the factor is $L_\flat^{2(i-k_\sharp-1)_+}$. It depends on the shell only, not on the stage, and it becomes larger with each successive shell until it reaches $\check{R}_k^{\,c}$. In the region not covered by the shells, the bound is expressed in the scale $j$ of the stage and equals $(1-\tfrac12\beta_{\mathrm{s}})\,c_l\,L_\flat^{2e}\varepsilon_j\eta_j^2$, where $c_l=3$ adjacent to $\Omega_k$ at the last stage and $c_l=1$ elsewhere. Far from the box, the curvature of the least-action continuation $U_0$ can be as large as $(2+c)\varepsilon_k\eta_k^2$, and the test of the last stage must be implied there as well. Here $e=0$ for $j\le k_\sharp$. For $j>k_\sharp$ one takes $e=(j-l-1)_+$ on $\Omega_l\setminus\Omega_{l+1}$, $k_\sharp<l\le j$, and $e=j-k_\sharp-1$ between the shells and $\Omega_{k_\sharp+1}$. This is the form of 122 (1989) 175, (1.24), where the level is denoted $k_0$ and the base $L_0$. On a layer $l<j$ this bound remains below $\varepsilon_l\eta_l^2$, the scale of that layer. In passing from the scale $j+1$ to the scale $j$ one gains $L^{-2}$ and pays $L_\flat^2$. On plaquettes read by both stages, the later threshold leaves a margin for the geometrically decaying response. The crossing plaquettes of the prepared interface require the comparison below. The sharper threshold enters in two places. In the first, the new fluctuation test has to imply the old fluctuation tests of step $j+1$, at $\delta_j$ and $2\delta_j$. This holds because $CL\delta'_j=CL(\log x_j)^{-(\mathfrak{p}_0-\mathfrak{p}_1)}\delta_j\le\tfrac12\delta_j$ for $\gamma$ small, and it puts no new condition on the constants chosen earlier. In the second, at the main step, $\delta'_k$ has to remain a small fraction of $\varepsilon_{k_\flat}$, which is the curvature threshold of the core level. In this estimate the thresholds of two levels may differ by $\varepsilon_j/\varepsilon_i\le C(1+(j-i)^{1/2})$, the bound valid for every pair of levels \cref{sec:large-field-the-induction-closed-the-order-of-the-constants-the-sum}. Restrict to the common support of the carried tests $\chi_k$, the tests of (2) and (3), and the tests and delta functions that stand under the core operation; the latter make the data stored in the core regular. On this set, and provided $C(M)\check{R}_k^{-c}\le c$ on the outermost shells, the final regularity threshold follows on its inset, with the seam implications discussed below. The equality does not reach $\partial X$: adjacent to $X^c$, the background localised in $X$ is constrained bond by bond on the bare lattice. There the test therefore demands $|\partial V_k-\mathbf{1}|<C\varepsilon_k\eta_k^2$ of the block field itself, while the carried tests provide only $2\varepsilon_k$. All of this holds under the condition

\begin{gather*}
\FitMath{C\,\mathsf{N}_k^{1/2}\,\frac{p_1(g_{k_\flat})}{p_0(g_{k_\flat})}+C(M)\,L^{-\mathsf{N}_k}\mathsf{N}_k^{1/2}\le c:}
\end{gather*}

the first term is $\delta'_k$ measured relative to $\varepsilon_{k_\flat}$, including the allowance $\mathsf{N}_k^{1/2}$. Since $\mathsf{N}_k=\check{R}_k\approx(\log x_k)^r$, it is a condition on the exponents, $\mathfrak{p}_1<\mathfrak{p}_0-Cr$, together with a ceiling on $g_k$. The second term is a level-$k$ field viewed from level $k_\flat$, so $\mathsf{N}_k$ is taken large.

\paragraph{The prefinal reads at the old wall.}
Let \(q=j+1<m\) and \(a_q=M\check R_qL^q\) in
bare units. Write \(S_X\) for the level-\(j\)
average of the prospective background localized
to \(X\), in the frame of the approximate collar
test. In an uncompressed stage the new
prepared face is five level-\(q\) cells from the
old \(\Omega_q\) wall. At \(q=k_0+1\) it is
seven cells away. For \(q>k_0+1\) the prepared
box is outside \(\Omega_{k_0+1}\), hence
outside \(\Omega_j\); the ordinary separation
from \(\Omega_j^c\) to \(\Omega_q\) is at least
eight level-\(q\) cells. At every prefinal
stage this last separation also bounds the
clearance to the older determining layers
on the exterior side. The clearance on the
interior side from the wall to \(\Omega_{q+1}\)
is at least eight level-\((q+1)\) cells.
Since \(a_{q+1}\ge a_q\), the four-cell wall
neighborhood lies in \(\Omega_{q+1}^c\cap X\),
which is inside the two-cell top inset of \(X\).
The two groups of ordinary
approximate tests have bases in
\((\Omega_q^{\sim2}\setminus\Omega_q)\cap X\)
and \((\Omega_q\setminus\Lambda_q^\sim)\cap X\).
Their bases lie within two level-\(q\) cells of the
wall, since in this quiet component \(Z_q\) is the
two-cell enlargement of \(\Omega_q^c\cap X\).
The four test layers of the localization, with
the endpoint blocks of every bond it reads, add
less than one level-\(q\) cell. This follows from
\(L^2M_2\le M\) and the fixed block reach. Thus the
complete reads lie within three cells of the wall,
with clearance from the new prepared face, the
older core, the next ordinary layer and \(X^c\).

Use this region with an additional fixed comparison
band on each side, of total reach less than four
cells from the wall. The ordinary global center
at step \(q\) and the preceding stage center
have the same rows there: the level-\(q\) rows on
the inside and the level-\(j\) rows on the outside.
Earlier promotions change only rows beyond the
original \(\Omega_j\) wall, with the associated
one-step completions there. The comparison band
misses those rows and their blocks by the
ordinary separation.
The two problems agree on the corresponding
localized datum with these two scales. The
joined comparison and averaging on the full
endpoint blocks bound the difference of their
axial readings by \(C_L\varepsilon_\star
e^{-c M\check R_q}\). Use the ordinary strip
shift on the bonds on which that substitution is
made. Its physical real bound on those bonds
is \(C_L\varepsilon_j e^{-c\check R_q}\).

The prospective stage has level \(q\) throughout
these test cubes and their enlargements. The
localized test-cube datum is the lift of the
same \(V_q\) there. Its background has the same
axial read at the common nodes. Comparing it with
the prospective background on the full
\(\square'^{\sim4}\) box gives
\[
 |S_X(b)(V^{\rm bg}_{j,\square'}(b))^{-1}
                                     -\mathbf1|
 \le C_L\varepsilon_\star e^{-c M_2\check R_q}.
\]
The distance from the endpoint blocks to the
comparison boundary is a fixed fraction of a
test layer, in level-\(q\) units. This is the
local comparison for the actual test centre,
with \(V_q\) fixed on its whole halo.

At a bond with neither endpoint in \(\Omega_q\)
the relative collar test gives \(2\delta'_j\)
against \(S_X\). For a bond read by the
completed ordinary strip test, compare instead
the ordinary center to the preceding-stage center
and then to \(S_X\) by the real fixed-target path
of \cref{hi:real-finer-collar-path}.
Across the old wall the coarse row joining the
two centres is an outward row of the old member
and a retained row after the promotion. Its value
is \(V_q(c)\) throughout the path. The two centres
are old nodes of levels \(j\) and \(q\). The
physical increment of this outward row is zero.
Its reduced increment still reads the changing
internal links of the exterior block. Use that
increment in the response and restore both node
values as in the two-block calculation below.
The restoration of the node gauges in the one-step
average leaves the identity at the old level-\(j\)
node and \(\mathsf v_j^{-1}\) at the level-\(q\)
node, with \(|\mathsf v_j-\mathbf1|\le C_L\delta'_j\).
The additional exterior contour of the ordinary
strip frame uses only internal bonds of the
exterior block. This block is entirely in the
new exterior collar, and those nonpivots are
tested there. The regular mean therefore bounds
that contour change by \(C_L\delta'_j\).
This gives the comparison in the inherited frame
with both endpoint factors included. At an old-wall
crossing the physical link itself is controlled by
the completed ordinary test at \(\delta_j\),
including the original real strip substitution.
The new collar test need not read that crossing.
Any adjacent factor on an old ordinary
Gaussian bond was removed by the ordinary
whole-term identity, with its original
cutoff and the retained strip factors.
That bond is absent from the free variables
of the prepared collar.
The sum of the center errors is at most
\(C_L\delta'_j+C_L\varepsilon_\star
e^{-c M_2\check R_q}=o(\delta_j)\), so both
groups of ordinary approximate factors equal
one on their complete reads. The argument uses
the direct exterior-only collar mask.

At the first stage put \(h=m_\flat\). The
surviving level-\(h\) density tests at
this first stage have bases in \(Z_h\cap\Omega_h
\cap X\). The stage-zero guard recovered in
\cref{hi:preceding-guard-support} implies the
other bottom density tests outside
\(\Omega_q^\sim\), by localization on
their own complete four-layer boxes. The bases
outside \(Z_h\) lie in \(\Lambda_h\), two
level-\(h\) cells from \(\Omega_h^c\), and
the exclusion of \(\Omega_q^\sim\) keeps
the full boxes in \(\Omega_h\setminus\Omega_q\).
The comparison with the pure level-\(h\)
test tower there costs at most
\(C\varepsilon_\star e^{-c\check R_h}\)
in level-\(h\) units, within the
\(\beta_{\mathrm s}/2\) margin of the
stage-zero bound. The retained bottom tests
read the frozen core: its gap to the first
prepared face is at least \(3a_q\) and their
halos reach less than \(a_q\).

\paragraph{The comparison of the centers at the last stage.}
Put \(j=m-1\), let \(Y=Z''_m(X)\), and write
\[
 K=(\Omega_m^c\setminus Y)\cap X .
\]
At the last stage the approximate collar test compares the
actual \(V_j\) to the \(j\)-fold average \(S\) of
\(U^{(\mathsf N_m)}_{m,X}\) on its nonpivot mask.
In the direct graded rule the mask consists of bonds with
an endpoint in \(K\) and neither endpoint in \(\Omega_m\).
The incident-bond convention additionally tests crossings
into \(\Omega_m\). In the direct rule those crossings use
the completed ordinary strip test.
Denote the corresponding preceding-stage average by \(S^-\),
the ordinary strip center by \(V^{\rm bg,loc}_j\), and the
ordinary test-cube center by \(V^{\rm bg}_{j,\square'}\).
We use the actual averaging constraints and keep the
preceding-stage guard during these comparisons.

The ordinary approximate tests in the fifth and sixth
groups have bases in
\((\Omega_m^{\sim2}\setminus\Omega_m)\cap X\)
and \((\Omega_m\setminus\Lambda_m^\sim)\cap X\),
respectively.
Their tested bonds are those of \(\square'^{\sim2}\) other than
the pivots. Their centers are formed on \(\square'^{\sim4}\).
The interior bases are at least one level-\(m\) cell from
the exterior of \(X\). The exterior bases lie in \(O_m\)
and within two level-\(m\) cells of \(\Omega_m\).
Each added test layer has width \(M_2\check R_m\)
in level-\(m\) units. The complete reads therefore remain at
distance at least \(cM\check R_m\) from \(\partial X\), including
the blocks of their end points. The prepared core \(Y\) lies in
\(\Omega_{k_0+1}^c\cap X\), which lies in the old live region
\(Z_{m-1}\cap X\). By \eqref{eq:elab-history-quiet-boxes} this
old region lies at least eight cells inside \(\Omega_m^c\).
Thus a test-cube box can meet the old wall and stays clear
of the prepared seam and the tower collar at \(\partial X\). The
prospective joined sequence has top level \(m\) throughout
the test-cube box.

On that box the actual top rows of \(U^{(\mathsf N_m)}_{m,X}\)
have values \(V_m\). Its joined datum for the test-cube
localisation is therefore exactly the one induced by the
lift of \(V_m\). The localised background is the ordinary
test-cube background. For a real datum below the ceiling
of the joined comparison, \eqref{mb:real-collar-datum}
and \eqref{mb:real-collar-norm} compare these backgrounds
in their common block axial reading. In the units of level
\(m\) the read is at distance \(cM_2\check R_m\) from the
boundary of its test-cube box. Consequently
\[
 |S(b)(V^{\rm bg}_{j,\square'}(b))^{-1}-\mathbf1|
 \le C\varepsilon_\star e^{-cM_2\check R_m}
\]
on the stated footprint, with \(\varepsilon_\star\) a
bound on the small regular datum. This application reads
the ordinary coarse density test and the retained rows.
It does not require the discarded sharp bound at the
artificial bare rim of \(X\).

For this real datum bound, the preceding-stage support
already controls the tested holonomies before the last
enlargement. Levels below \(j\) keep their rows. The
determining set of level \(j\) loses rows. Its second-kind
rows still use the unchanged layers \(j-1,j\), so the
remaining tested plaquettes have their old holonomies.
For a tested plaquette of level \(m\) with a corner in a
carried density cube, that test bounds its actual \(V_m\)
holonomy. The density cubes on \(\Omega_m^{\sim4}\) cover
the plaquettes near the old seam. In the other case the averaging
blocks of a newly tested plaquette are in the preceding
member of level \(j\), by the last prepared separation,
and outside the old \(\Omega_m\). Here the fine plaquettes
are formed from the old determining rows, and the exact
one-step average gives the coarse bound. At the exterior
tower a tested holonomy is either the identity or the
lifted top holonomy. These are bounds at the respective
levels of the datum. They do not divide the top plaquette
deviation by the square of the fine spacing.

To compare \(S^-\) with \(V^{\rm bg,loc}_j\), use the region
between an inner and an outer box. For the outer box take
\(X\) inset by a quarter of a level-\(m\) cell, rounded on
the grid of \(M_1\)-cubes at that level. The old exterior
inside \(X\) is a box separated from \(X^c\) by two cells
by \eqref{eq:elab-history-quiet-boxes}. For the excluded
inner box take \((\Omega_m^{\sim2})^c\cap X\), with this
enlargement in level-\(m\) cells. The outer boundary band
is in the old \(\Omega_m\), and the inner boundary band
is in \(\Lambda_j\setminus\Omega_m\). The earlier stages
alter rows only beyond the inner band, towards \(Z_j\).
The new gauge draws of the last stage start outside
\(\Omega_m^{\sim3}\). On the annulus both backgrounds
therefore have the same regions and the same determining
rows, including the top data used on its two boundary
bands. Join at level \(m\) on the outer band and at level
\(j\) on the inner band. Both localisations give the same
datum. For a read of either group with both ends in
\(\Omega_m\), use the localized strip center only on
the actual strip read and substitution mask. For a read
unaffected by that substitution use the ordinary global
center. It has the same determining rows on the annulus.
The block axial frame on the full read blocks is common,
and the distance to the annulus boundary is of order
\(M\check R_m\). With \(L_j(b)\) denoting the applicable
ordinary center, the comparison gives
\[
 |L_j(b)(S^-(b))^{-1}-\mathbf1|
 \le C\varepsilon_\star e^{-c\check R_m}.
\]

For the change from \(S^-\) to \(S\), use the finer
determining set at the averages of the prospective
background, then change the freed level-\(j\) rows to
their actual values. The two constraint sets are nested
at the induced datum. A critical point for the larger
constraint set is critical for the smaller set there.
The common rows, including the exterior tower rows, do
not change. On a nonpivot freed row the approximate
collar test bounds the increment by \(2\delta'_j\).
For a freed pivot, both blocks of its coarse row are
outside the old \(\Omega_m\). Each nonpivot in the
double block is controlled by the collar test or belongs
to the common fine rows. This includes the rows crossing
the prepared seam, whose outside block is in \(K\).
The two fine fields have the same \(V_m\) average on
that coarse row. The contour argument in
\cref{hi:real-finer-collar-path} places the actual pivot
in the small branch, and inversion of
the averaging equation bounds its increment by
\(C(L)\delta'_j\). A pivot crossing the old
\(\Omega_m\) seam is absent from the finer determining
set in the convention excluding inward bonds. The
resulting plain datum increment is supported in the
collar. The real response supplies the reduced datum
and the node factors, including any factors at a seam.
On the common axial reads inside the old \(\Omega_m\)
the response bounds \(S^-(b)S(b)^{-1}-\mathbf1\) by
\(C(L)\delta'_j\). This argument is made after the
approximate collar test and before the sharper collar
split.

On a read of either group in the old \(\Omega_m\), add these
three bounds to the carried or completed fluctuation bound
\(\delta_j\), with the ordinary strip shift on its actual mask.
The
exponentially small terms and \(C(L)\delta'_j\) fit in the
remaining \(\delta_j\) of the approximate threshold
\(2\delta_j\). A nonpivot read with an end point outside
the old \(\Omega_m\) has that end point in \(K\).
For a bond with neither endpoint in \(\Omega_m\), use
the collar test and the test-cube comparison to obtain
\(2\delta'_j+C\varepsilon_\star e^{-cM_2\check R_m}
<2\delta_j\). Both tests exclude the geometric
pivots on this overlap.

For a crossing of the old wall in the direct graded
rule, the full exterior endpoint block lies in \(K\),
and the interior endpoint is in the first stored
ordinary strip layer. The completed strip test there
gives the bound at \(\delta_j\) in the original plain
coordinate, without changing the physical link or
its soft gauge factor. Compare only the centres
\(L_j\) and \(S\) on the full two
blocks by the mean-contour axial reading used below
for the global collar chart. The exterior contour of
the ordinary frame is determined by the internal
bonds of the exterior block. The direct collar test
does read these internal bonds. Its bound controls
that contour without assuming an additional collar
bound on the crossing itself. On the common retained
centre pins the comparison gives
\[
 |L_j(b)S(b)^{-1}-\mathbf1|
 \le C_L\delta'_j+C_L\varepsilon_\star e^{-c\check R_m}
\]
in the inherited frame. Together with the test-cube
error and the original strip bound this is less than
\(2\delta_j\). The two-block comparison below uses
the actual retained centre nodes and the original
exterior contour.
When using the incident-bond convention and the
crossing itself lies in its mask, the direct collar
bound applies in that convention's prescribed frame.

For the inset split the finer path supplies the local
collar response at a fixed admitted prospective datum.
The clearances above place the complete footprints of
the two groups in its comparison region, including the
interior endpoints of tests based outside \(\Omega_m\).
The inherited strip quadratic form uses the original
strip variables.

\paragraph{The preceding guard on the prepared seam.}
\label{hi:preceding-guard-support}
The stage-zero plaquettes lie in
$\Omega_h\setminus\Omega_{h+1}$. At any stage with
$h<j<m$, every prepared shell lies in $\Omega_m^c$.
The deep exterior clause is in $\Omega_{i_0+1}^c$,
where $i_0+1\le j$, and the other exterior clauses are
in $\Omega_l\setminus\Omega_{l+1}$ with $l\le j$.
Consequently every earlier guard reads only
$X\cap\Omega_m^c$, which is two level-$m$ cells inside
$X$ by \eqref{eq:elab-history-quiet-boxes}. The
artificial bare rim occurs only in the final guard.

Consider the implication from stage $m$ to $j=m-1$
on the ordinary real collar chart. On a shell of index
$i<j$ already present at $j$, the geometric factor and
curvature scale are unchanged. The later coefficient
is smaller by $\beta_{\mathrm s}2^{-(j-i+2)}$. The
weighted response from the changed collar rows is at
most
\[
 C_L\frac{\delta'_j}{\varepsilon_j}
 (1+\sqrt{j-i}) e^{-c(M/M_1)(j-i)}
 \varepsilon_i L^{-2i}.
\]
Here the ratio of thresholds contributes
\(1+\sqrt{j-i}\) and the shell separation gives the
exponential. Choose \(M/M_1\) so their product is
bounded by a constant times \(2^{-(j-i)}\), then use
\(\delta'_j/\varepsilon_j\) to fit the reserve. Any
overlap factor on this fixed shell is the same at
both stages and is at least one. On a plaquette wholly in the new shell
$Z''_m\setminus Z''_j$ the factor is unchanged and the
coefficient decreases from $1-\beta_{\mathrm s}/2$ to
$1-3\beta_{\mathrm s}/4$. Outside both nests the ratio
of the two thresholds is at most
\[
 \frac{\varepsilon_m}{\varepsilon_j}
                    \frac{\max\{L_\flat^2,3\}}{L^2}.
\]
The factor $3$ occurs at $l=j$. The adjacent reference
couplings differ by a bounded increment, so their
threshold ratio tends to one in the small-coupling
range. These exterior comparisons have a fixed
fractional margin. On admitted real data the collar
response fits these reserves, by its weighted bound on
the old shells and the small ratio $\delta'_j/\varepsilon_j$
in the exterior. All these reads remain in the inset.

A plaquette can cross a face of $Z''_m$ while lying
wholly in $\Omega_{k_0+1}^c\setminus Z''_j$. Choose a
normal and a tangent edge away from the corners of the
face, using the half-spacing faces of the centered
blocks. The separation of the prepared boxes puts the
whole plaquette outside $Z''_j$, and the clearance from
$\Omega_{k_0+1}$ puts it inside the old deep exterior
region. The preceding guard reads it. The final shell
and exterior clauses both require containment, so neither
reads it. It is also outside the later original layers
used by the other exterior clauses.

At such a plaquette write $T_j(p)$ for the preceding
threshold and $E_j(p)$ for the bound on the curvature
change from the collar response, in bare plaquette units.
We must show that the prospective curvature is below
$T_j(p)-E_j(p)$ on the support on which the preceding
guard is removed. The small-paying-energy estimate
below gives this comparison on the part of the support
to which it applies. In the two joined sequences the
plaquette has the old and new scales $j$ and $m$ under
the largest-corner convention. The weighted profile
at $p$ is bounded, in either sequence, by
\[
 C\left(\varsigma_*+\check R_m^2\varepsilon_m
                    e^{-c\delta_{\mathrm c}M\check R_m}\right)
                         =o(\varepsilon_j).
\]
Indeed the old rows are read from the compatible field,
and the new rows use its retained averages. The distance
to the floor uses the same prepared-window separation.
The two bare curvature bounds gain respectively $L^{-2j}$
and $L^{-2m}$. Their sum is $o(T_j(p))$ and bounds the
change in any common real frame. This proves the preceding
seam test on the low-energy branch.

In an inset split let $\chi_{\rm in}$ be the literal final
clauses restricted to plaquettes whose corners are all in
$X$ with one full level-$m$ cell layer removed. The preceding
reads omitted by these clauses are the deep exterior
plaquettes straddling $\partial Z''_m$. Put
$F_j=L_\flat^{2(j-k_0-1)}$
and let $\kappa_\times$ test each of these plaquettes at
$(1-3\beta_{\mathrm s}/4)F_j\varepsilon_jL^{-2j}$ in the
prospective background. The datum for that test is held fixed
in the collar integration. If $\chi_{\rm in}$ is the final
inset test, the complementary factor of their product is
\[
 1-\chi_{\rm in}\kappa_\times
  =(1-\chi_{\rm in})+\chi_{\rm in}(1-\kappa_\times).
\]
Both terms belong to one regularity stop before the collar
integration. Under the payment hypotheses a failure in the
second term is impossible at small paying energy by the
profile estimate above. The threshold leaves a reserve
$\beta_{\mathrm s}F_j\varepsilon_j L^{-2j}/4$ for the
old-chart response. On the passing branch the
fixed-target response estimate below fits that reserve
throughout the admitted real collar box, so the
preceding seam test is one there. This implication
uses the reserve and the collar response. The
small-paying-energy argument concerns failure of
\(\kappa_\times\).

%
There is also a direct bound for these crossings on the
admitted real inset chart at the final stage. Translate a
crossing plaquette \(p\) a bounded number of bare spacings
into an exterior cell at one of its outside corners,
keeping its orientation. The translated plaquette lies in
the deep exterior and in the inset. The separation from
\(Z''_j\), the old wall, and the outer boundary lets both
plaquettes lie in the smaller part of a fixed-size
level-\(j\) regularity cube. Its enlargement is in the
prospective member of level \(j\). The local gauge on this
cube applies across the level-\(m\) face.
In that common gauge the difference of the linear plaquette
curls is at most \(C\varepsilon_\star L^{-(2+\alpha)j}\)
by the first-difference H\"older bound. The quadratic and
higher terms in the products of the four links vary by at
most \(C\varepsilon_\star^2 L^{-3j}\), by the bound on
their first differences. The inset test bounds the defect of
the translated plaquette by
\[
 (1-\tfrac12\beta_{\mathrm s})L_\flat^2 F_j
                    \varepsilon_m L^{-2m}.
\]
Divide the three bounds by \(F_j\varepsilon_j L^{-2j}\).
The last displayed bound has ratio at most
\((1-\beta_{\mathrm s}/2)(L_\flat/L)^2(1+o(1))\).
The two errors tend to zero, since \(j=m-1\) and
\(m\ge\check R_m\), with \(\check R_m\) growing at least as
the square of the logarithm of the inverse coupling.
This gives the reserve of \(\kappa_\times\) once the real
prospective datum has been admitted. The fixed-target
response then transfers it to the preceding guard.

%
An earlier transition has its own crossing reads. In
particular, with \(h=m_\flat\), the first prepared face
lies strictly in \(\Omega_h\setminus\Omega_{h+1}\).
A bare plaquette astride that face is tested at stage
zero. Neither containment clause of the next guard
tests it. The final-stage H\"older estimate uses
\(j=m-1\) to make its translation error small compared
with the tested threshold. At the first face that
argument does not supply the reserve.

For a transition from \(j\) to \(j+1\), use the actual
list of preceding reads. At a newly omitted crossing
the preceding coefficient is
\((1-\beta_{\mathrm s}/2)F_j\varepsilon_jL^{-2j}\),
with \(F_j=1\) in the uncompressed range and
\(F_j=L_\flat^{2(j-k_0-1)_+}\) thereafter. In a split
with an added crossing test, impose
\((1-3\beta_{\mathrm s}/4)F_j\varepsilon_jL^{-2j}\)
on the fixed prospective datum. This leaves
\(\beta_{\mathrm s}F_j\varepsilon_jL^{-2j}/4\) for
the old response. The first face has the separation
from the bottom core and the old wall needed for the
local paying-energy comparison. On the complete
crossing carriers the small-energy argument bounds
the prospective defect by
\(C(\sigma+\check R_m^2\varepsilon_m e^{-cM\check R_m})L^{-2j}\),
with \(\sigma=c_*\delta'_j\check R_m^{-4}\).
This fits the crossing test. Its failure on the
otherwise passing prospective branch belongs to the
same regularity stop on the old prefix. On its passing
branch \eqref{hi:real-crossing-response} fits the reserve
on the whole admitted real collar chart and gives the
preceding inequality there.

%
%
%
At the first transition put \(h=m_\flat\), \(q=h+1\)
and \(Y=Z''_q\). The original stage-zero guard tests
plaquettes contained in \(\Omega_h\setminus\Omega_q\).
All of these lie in the inset of \(X\). For one contained
in \(Y\), the first prepared shell tests it at
\((1-3\beta_{\mathrm s}/4)\varepsilon_h L^{-2h}\).
For one contained in \(Y^c\), use the new deep exterior
bound \((1-\beta_{\mathrm s}/2)\varepsilon_q L^{-2q}\).
Here \(h<k_0\), so \(i_0=h\) at the new stage and
the exterior power of \(L_\flat\) is zero. A plaquette
with corners on both sides uses the added crossing test
at the first shell threshold. In all three cases the
new bound leaves at least
\(\beta_{\mathrm s}\varepsilon_hL^{-2h}/4\)
below the old bound. For the exterior case use
\(\varepsilon_q/\varepsilon_h=1+o(1)\) and \(L^{-2}\).

The passing prospective guard fixes an admitted real
orbit at the outgoing datum. Hold that datum fixed and
use \cref{hi:real-finer-collar-path} with target \(V_q\)
on the full approximate collar ball. The old curvature
can change by at most \(C_L\delta'_h L^{-2h}\) on the
first face and on the adjacent layer-\(h\) stencils.
On the covered layer the same weighted comparison
applies to its full stencils, using the old units there.
This error fits the reserve. Thus the old stage-zero
guard equals one on the real chart for the passing
graded branch. The earlier value of that guard is
still used in assigning the incoming stage-zero
outcome. The removal in the collar integral follows
from the comparison on its whole real chart.

For an operation using these additional tests through
the whole window, assign a newly added crossing of
\(\partial Z''_{s+1}\) to shell \(s\). Retain it in
the subsequent guards with the same coefficient as
that shell. Its first coefficient is
\(1-3\beta_{\mathrm s}/4\). From the transition
\(j\) to \(j+1\) with \(j>s\), its coefficient drops
by \(\beta_{\mathrm s}2^{-(j-s+2)}\). The stencil
at this old face has scale at least \(L^s\), and the
separated faces give the same weighted distance to
the newly changed rows as for the adjoining shell,
up to a bounded stencil enlargement. For each
intervening layer the path crosses a width of at
least a fixed multiple of \(M\) in its own units.
The decaying form of \eqref{hi:real-crossing-response}
therefore gives a defect change bounded by
\[
 C_L\delta'_j L^{-2s}e^{-c(M/M_1)(j-s-1)}.
\]
The adjacent frame comparison costs at most
\(C(1+\sqrt{j-s})\) in comparing \(\varepsilon_j\)
with \(\varepsilon_s\). Choose the shell separation
and then the coupling floor as in the original ladder.
The change fits the decrease of the coefficient at
every transition. Thus the already assigned crossings
are carried with the earlier shell tests until the main
integration.

For comparison of the last-stage change alone with the
literal final test, let $r$ be the
product of its omitted indicators on the same prospective
datum. The literal continuation is $\chi_{\rm in}r$,
and the inset continuation is $\chi_{\rm in}\kappa_\times$.
Their two differences are the sectors
$\chi_{\rm in}r(1-\kappa_\times)$ and
$\chi_{\rm in}(1-r)\kappa_\times$. One may refine each
literal branch by multiplying it by $\kappa_\times$ and
by $1-\kappa_\times$, leaving that branch's kernel and
stage count on both pieces. Summing the two pieces recovers
the original labelled measure. To change the continuation
rule itself one must compare the two kernels on the
difference sectors. The later implication for
$\kappa_\times$ is on the main-step support. An earlier
stop belongs to a different branch.

The inset test can also define the split of a separate
operation from its first formation. Put
$G=\chi_{\rm in}\kappa_\times$ on the real pre-collar
fibre, counting an undefined tested background as a
failure. Use the old prefix kernel on $G=0$ and the
admitted continuation moves on $G=1$. Both are positive
kernels with unit mass when written as the real moves,
before their weights are expressed by the collar
integral. Thus the two branches exhaust the incoming
measure. A bare insertion multiplies the incoming
bare measure before applying the same kernels. For
example a curvature source replaces \(\rho_0(U)dU\)
by \(e^{A(w,U)}\rho_0(U)dU\) and uses the tests and
reference laws of the unsourced run. The stage count
on $G=0$ is the
preceding count, and failure of the product makes a
single mark. The mark's own mass omits its failing test.
A later containing fibre can use that failure for one
payment when the complete witness is available there.
The move identities give an exact weighted integral
with labels assigned by \(G\). To identify its
positive branch with the localized stage integral,
we must account for every retained test on its
actual chart. An indicator on the fixed outgoing
datum can stay outside the collar logarithm, as
the final top-density factor does. An indicator on
an integration variable can stay in the expansion
if it is fixed in the real integration coordinates
throughout the parameter domain, has a fixed carrier,
and accepts the smaller coordinate ball used in its
Gaussian tail estimate.
For example, on a plain strip bond with reference
\(S_{\rm o}\), an approximate factor with cube
centre \(D\) tests \(|e^{ig_j a_b}S_{\rm o}D^{-1}-\mathbf1|<2\delta_j\).
If the centre error is at most \(c\delta_j\) with
\(c<1\), the factor accepts the ordinary ball
\(|g_j a_b|<\delta_j\). This is a one-bond factor
when the centres are fixed on that fibre. A
test \(I_S\) with active slots \(S\) that accepts
the product ball of radius \(h\) has
\[
  1-I_S\le \sum_{b\in S}\mathbf 1_{\{|z_b|\ge h\}} .
\]
The absolute Gaussian tilt on its carrier therefore
gives a tail factor \(|S|e^{-\gamma_2 h^2}\) times
the normalization and mean costs of the tilt.
For several fixed selectors \(I_\alpha\), use their
complete carriers in every binary cut, and let \(S_\alpha\)
be their sets of active coordinates. Suppose that at
most \(\nu\) of the sets \(S_\alpha\) contain any one
bond, counting each selector only once. For a family
\(A\) of failed selectors choose one bond from each
\(S_\alpha\). Then
\[
 \prod_{\alpha\in A}(1-I_\alpha)
 \le \sum_{\substack{b_\alpha\in S_\alpha\\ \alpha\in A}}
 \prod_{b\in\{b_\alpha:\alpha\in A\}}
                    \mathbf 1_{\{|z_b|\ge h\}}.
\]
The set of bonds on the right has size at least
\(|A|/\nu\). Apply the stage Gaussian tilt once per
distinct bond. On a stage chart with the real-contour
Gaussian bound this gives the tail
\[
 \left(\prod_{\alpha\in A}|S_\alpha|\right)
          e^{-c h^2|A|/\nu}
\]
in addition to the Gaussian normalization cost
and the bounds on the analytic vertices. In the
tilt the positive quadratic form is diagonal on
the selected coordinates, so its size per
coordinate is independent of the multiplicity.
If no active coordinate enters a selector, carry
its real value outside the Gaussian estimate.

For a bond selector attached to an original test
cube, the sets of tested bonds lie in the two
test layers of that cube. A bond can belong to only
a fixed number of such enlarged cubes. Restrictions
to the integration mask can decrease this number.
The product of a repeated selector with itself
can be counted once. For a single stage, the
cardinality of each such set grows at most as a
fixed power of \(\check R_m\). At the stage
fluctuation threshold \(h\) is a fixed fraction
of \(\delta'_j/g_j\), so the factor
\(\log|S_\alpha|\) is absorbed by \(c h^2/\nu\)
at weak coupling. This count uses only the original
test-cube reads. If the inverse chart introduces
new active coordinates, use the resulting
supports and prove the overlap estimate on them.

Fixed Borel selectors in these coordinates require
no continuation as functions of \(z\). Keep them
on the real contour, and use the locally uniform
integrable Gaussian majorant for holomorphy in
the analytic parameters. Their compatibility with
the factorization uses their full carriers. For a
preceding curvature guard on the active collar,
we must either show it holds on the complete
chart or put it in this fixed-domain form and
prove that it accepts the reference ball. A bound
only on the support of the literal prospective
guard gives neither property for the inset rule.
If its domain moves with the analytic parameters,
the real tail estimate leaves the transported
integral and its bound to be proved.

For the remaining inset-positive evaluations use the real
collar path at fixed prospective datum. The seam estimate
concerns the prepared interface inside $X$. The face-lift
calculation at $\partial X$ supplies its own restriction.

%
%
%
\paragraph{The lower density and contour tests on the inset.}
%
%
%
%
%
%
In the ordinary step \(t=j+1\) use the query rule
in which a collar bond
has both ends in the explored region of its round.
Since that region only grows, a bond crossing the
final small-field boundary \(\partial\Lambda_t\)
has not been queried in the collar. Its inner
threshold in the ordinary mean estimate is the
rim threshold \(\theta_c\delta_j\). In a rule that
queries incident crossings as well, the accepted
ball must fit all the resulting tests, including
\(\delta_j^\flat\) on a queried pass. The stored
coordinates in either case are the complement
of the actual integrated nonpivots.

The ordinary labelled fluctuation split retains
the signs of the queried outcomes. For explored
bonds \(E^{\rm T}\), failure set \(I\), explored
rim cells \(E^{\rm H}\), and failed rim set
\(\mathsf R^{\rm bord}\), its factors include
\[
 \prod_{b\in I}\chi_b^1
 \prod_{b\in E^{\rm T}\setminus I}\chi_b^0
 \prod_{c\in\mathsf R^{\rm bord}}(1-\operatorname{lo}(c))
 \prod_{c\in E^{\rm H}\setminus\mathsf R^{\rm bord}}
       \operatorname{lo}(c).
\]
Keep also the failed-cube factors and the unexplored
bond factors of the ordinary split. An eligible
component of the collar operation misses every
ordinary creation set at the unfolded levels. A
failed queried bond puts its endpoint cells in that
set. A failed rim test puts its rim cell there, and
a failed cube puts its coarse hull there. Their
incident reads and the cube test halos lie in the
same enlarged ordinary component as the creation.
By nesting, such a component reaching the stage
component would put its creation inside it.
Thus only the passing queried factors occur on
the eligible collar carriers. The failed factors
in the frozen core keep their original signs.

Apply the test implications on the successful
real support in their original order. Use the
prospective regularity and approximate collar
tests, the exact retained averages, and the
ordinary completed fluctuation factors. These last
factors control the inherited strip variables read
by the belt tests, including the old-wall crossings.
The sharp exterior mask concerns a different stock
of nonpivots. Remove the preceding regularity,
ordinary density, contour, and two approximate belt
groups on the full reads where the corresponding
implications have been proved. Keep the final top
density, the soft gauge density, the carried
fluctuation factors, and the separate collar and rim
factors. Keep any other positive guard on its inverse
read until its own implication is established. An
active guard belongs inside the exact real collar
integral, while a fixed-datum factor can be taken
outside. The sharp collar split follows these
comparisons.

To apply the ordinary fluctuation factors, let
\(L_b\) and \(G_b\) be the localized and global
ordinary strip centres and put \(e^{iH_b}=G_bL_b^{-1}\).
Write \(W'=e^{iA'_j(b)}\) for the real strip
coordinate after the ordinary cutoff substitution.
The physical link on a substituted bond is
\[
 V_j(b)=W'e^{i(1-\vartheta(W'))H_b}L_b,\qquad
 V_j(b)G_b^{-1}=W'e^{-i\vartheta(W')H_b}.
\]
In the unchanged zone use the original coordinate.
On the plain strip \(\vartheta=1\). In the cutoff
zone \(\vartheta=0\) through angle \(4\delta_j/3\)
and \(\vartheta=1\) from \(5\delta_j/3\) onward.
If \(|H_b|\le\delta_j/8\), the hard angle tests
at thresholds at most \(29\delta_j/24\) are
unchanged. They agree below the start of the cutoff,
and above it the change of angle is at most
\(\delta_j/8\), so both fail. For an approximate
test with cube centre \(D\), put \(X_b=L_bD^{-1}\).
When \(\vartheta<1\) and
\(|H_b|+|X_b-\mathbf1|<\delta_j/3\), both
\(W'X_b\) and \(W'e^{i(1-\vartheta)H_b}X_b\)
have chord distance less than \(2\delta_j\) from
the identity. When \(\vartheta=1\) the two
arguments are equal. This proves the approximate
test identity on the cutoff mask with its
centre reserve.

%
%
%
Let \(h<j<m\) and \(q=j+1\). The lower
density cubes have bases in
\((\Omega_j^{\sim4}\setminus\Omega_q^\sim)\cap X\).
The full four-layer halo of a base extends
by less than one level-\(j\) cell. The
eight-cell gap to \(\Omega_{j-1}^c\) and
the exclusion of \(\Omega_q^\sim\) put this
halo in \(\Omega_{j-1}\setminus\Omega_q\).
It lies in the top inset of \(X\) because
it lies in \(\Omega_q^c\cap X\). It misses
the old prepared box \(Z''_j\). For \(j\le k_0\)
this follows from its five-cell distance
to \(\Omega_j\). At \(j=k_0+1\) the
distance is seven cells. At later levels
\(Z''_j\) is outside \(\Omega_{j-1}\).
Hence the whole density box is in the old
prepared layer \(j\). Use that layer's
actual determining rows at the old target.

For \(q\le k_0+1\), the halo can meet the
new prepared box \(Z''_q\). A plaquette
wholly inside that box is in the new shell
and is tested at
\((1-3\beta_{\mathrm s}/4)\varepsilon_j L^{-2j}\).
This includes any such plaquette crossing
the original \(\Omega_j\) wall. A plaquette
wholly outside the box uses the deep exterior
bound \((1-\beta_{\mathrm s}/2)\varepsilon_q L^{-2q}\),
since \(i_0=j\). A plaquette astride the
new face is wholly in \(\Omega_j\setminus\Omega_q\).
Indeed that face is five level-\(q\) cells
from the old \(\Omega_q\) wall, or seven
when \(q=k_0+1\), and the gap to
\(\Omega_j^c\) is eight cells. The preceding
guard therefore reads this plaquette and
the new crossing test supplies the shell
bound on it. In all three cases the new
bound leaves a fixed fractional margin
below \(\varepsilon_j L^{-2j}\).

For \(q>k_0+1\), the halo misses both
prepared boxes. When \(j=k_0+1\) both
faces are at least five level-\(j\) cells
from \(\Omega_j\), and for \(j>k_0+1\)
both boxes are outside \(\Omega_{j-1}\).
The halo is now in the prospective layer
\(q\). On plaquettes wholly in one of the
two original layers at the \(\Omega_j\)
wall the exterior bounds give
\[
 |\partial U^+(p)-\mathbf1|
 \le (1-\beta_{\mathrm s}/2)
      \max\{3,L_\flat^2\}\varepsilon_q L^{-2q}.
\]
For an astride plaquette translate it a
bounded number of bare spacings into the
original layer \(j\) in the same prospective
comparison cube. In a common real gauge
\eqref{mb:local-regularity-conclusion} and
\eqref{mb:holder-reduction} bound the
change in the curl by
\(C\varepsilon_\star L^{-(2+\alpha)q}\).
The variation of the terms of order at
least two in the plaquette logarithm is
\(O(L^{-3q})\). In the compressed range,
\(q\ge m-O(\log\check R_m)\), while
\(m\ge\check R_m\). With the frame
comparison in the waiting window, both
errors are \(o(\varepsilon_jL^{-2j})\).
Since \(L_\flat<L/2\), the resulting
bound has a fixed fractional margin here
as well.

The fixed-target real response transfers
the resulting bound to the preceding
background at a cost \(C_L\delta'_j L^{-2j}\).
On the full density localization box the
preceding background has the actual
level-\(j\) determining datum. Its
localization with the lifted \(V_j\)
on the pure level-\(j\) tower is precisely
the background used by the density test.
The inner plaquettes of the test have the
three-small-cube margin to the boundary
of that localization. The joined
comparison bounds the additional error
in level-\(j\) units by
\(C\varepsilon_\star e^{-c\check R_j}\).
Both errors fit the fractional reserve,
so the original lower density test passes
on the inset-positive real collar chart.

The contour indicators of the old exterior
belt use the order-\(q\) blocks based between
the first and third layers outside \(\Omega_q\).
Those blocks lie entirely in the collar. The
new gauge insertion starts beyond the third
layer, on the remaining exterior collar
blocks outside \(Z''_q\). All their internal
links are tested exterior nonpivots. There
are no finer determining nodes in these
blocks after promotion, so the prospective
axial frame has trivial mean contours.
On the local approximate collar ball the
links differ from their prospective axial
reference by \(O(\delta'_j)\), also after
changing the stage centre. The mean contour
map on the full block has a bounded
derivative, so these indicators pass at
their original \(\varepsilon_j\) threshold.
The soft density on the same contours
stays in the integral.

At a prefinal stage the level-\(q\) density
test has bases in
\((\Omega_q^{\sim4}\setminus\Omega_{q+1}^\sim)\cap X\).
Its small-cube halo and the fixed row
stencils fit in \(\Omega_q^{\sim5}\) and
miss \(\Omega_{q+1}\), since
\(a_{q+1}\ge a_q\). They are in the
promoted level-\(q\) member. Thus the
factor reads the fixed \(V_q\) and may
stand outside that collar integration. At the
following executed stage of the same component
its bases are precisely the lower density bases,
on the actual halo just treated. On an
intervening stop it keeps its real value
on the recorded outgoing graph. The final top
density remains with the fixed top field.

Thus on the passing graded collar branch
the old lower density and contour factors
can be removed by identities of real
integrands. The ordinary completed
fluctuation factors keep their original
strip coordinates, and the new sharp
factor uses the free exterior mask.
The bottom-core and fixed-outgoing
selectors stay on their real arguments.
This reduction uses the original chart
and the crossing clauses of the graded
rule. The next estimates must admit the
remaining analytic weights on the
complex collar section with these masks.

At the final stage set \(j=m-1\).
For an old approximate cube test whose
halo crosses \(\Omega_m\), distinguish the
actual collar mask and the inherited strip mask
as in the two belt comparisons. For the
remaining old approximate tests meeting
the small-field integration region,
apply the implication inside the original
ordinary-step integral, using its Gaussian
cutoffs on the variables integrated there.
The later stage has no independent copy of
those variables. If an earlier integral already
filed a factor involving those tests, use its
actual resulting factor until it has been
identified on the retained coordinates.

\paragraph{The real finer datum and the collar path.}\label{hi:real-finer-collar-path}
At a transition from level \(j\) to \(q=j+1\), put
\(Y=Z''_q\) and let \(K\) be its exterior collar in
\(\Omega_q^c\cap X\). Fix an admitted prospective real datum and its regular critical
configuration \(U^+\) in the frame of the approximate collar
test. Choose the common stored rows in that frame, which stays
fixed in the comparison. The prospective datum is held fixed
while the finer variables vary in the collar.
Use the fixed small regularity ceiling of the background
theorem, with room for constants depending on \(L\).
Let \(V^0\) be the datum obtained by reading \(U^+\) on
the preceding determining set. All common rows have their
prescribed values. The two joined sequences have the same bare
domain and agree except for the enlargement at level \(q\).
An old level-\(j\) inward bond has an end in the old \(\Omega_q\),
hence also in the enlarged member. It is absent from the new
no-inward determining set. The lower inward masks agree.
This is the nested restriction condition for the no-inward
determining sets. The old curvature and current scales are no
larger, so \(U^+\) is regular also in the old units.
A tested plaquette of the
old datum at level \(q\ge1\) has a corner in the old
member at level \(q\). Its full order-\(q\) carrier lies
in the member at level \(q-1\) by the gap between the
members. Its holonomy is the plaquette of \(\mathcal M^qU^+\),
including the one-step completions of the second-kind
bonds. The average bound gives an old datum radius of at
most a fixed multiple of the regularity ceiling of \(U^+\).
At level zero use the bare curvature bound. Also, the
constraints of the old set at \(V^0\) imply those of the
new set, since the new rows are obtained by composition
on their old fine blocks. It follows that \(U^+\) is
critical for the old set at \(V^0\).

Keep the inherited coarse rows at their prescribed values.
For each newly imposed coarse row \(c\), both endpoint blocks
are outside the old \(\Omega_q\), by alignment on the
level-\(q\) grid. They are inside the old member of level
\(j\). At the first stage the new face is at least three
level-\(q\) cells from the bottom core. At subsequent stages
the gap between the prepared faces exceeds the width of
the full double block.
Every nonpivot in the double block either meets \(K\) and
is read by the approximate collar test, or belongs to
the common fine rows on the other side of \(\partial Y\).
The actual fine field and the induced fine field
\(\mathcal M^jU^+\) therefore differ by at most
\(2\delta'_j\) on the nonpivots in the frame of that test.
Their common coarse value is \(W=V_q(c)\).

To check the actual pivot's branch, use the straight
corridor product \(q\) as its coordinate, with every other
fine bond fixed. In each of the two logarithmic cases of
the real averaging formula the coarse value is
\[
  F(q)=\exp_q\!\left(\sum_a\alpha_a\log_q H_a\right),
  \qquad \sum_a\alpha_a=1-L^{-3}.
\]
Here \(\exp_q,\log_q\) are the round exponential and
principal logarithm on \(\mathrm{SU}(2)\), with the group
logarithm transported by right translation from the identity.
The anchors \(H_a\) are the transports avoiding the
middle bond. On the regular branch their end pieces are
the averaged block contours. On the other logarithmic
branch use the literal contours. The \(L\) straight
transports through the middle bond give zero in the
logarithmic sum, which accounts for \(L^{-3}\).
If neither logarithmic case applies, \(F(q)=q\).

Write \(e\) for the regularity radius of \(U^+\).
The level-\(j\) plaquette defects of its induced field
on the double block are \(O_L(e)\), by the average
bound in the member of level \(j\). Sweeping the fixed
contour stencils shows that their transports lie within
\(C(L)e\) of the central transport, and hence within
\(C(L)e\) of \(W\) by the regular averaging formula.
Changing the nonpivots by \(2\delta'_j\) moves each literal
transport by at most \(C(L)\delta'_j\). The averaged
end contours have the same bound, since each stays in
its regular mean chart and reads no pivot. Thus
\(d(H_a,W)\le C(L)(e+\delta'_j)\) in either formula.
If \(F(q)=W\) in a logarithmic case, the tangent sum has
norm less than \(\pi\). Its norm therefore equals
\(d(q,W)\), and the triangle inequality gives
\[
  d(q,W)\le(1-L^{-3})
           \bigl(d(q,W)+C(L)(e+\delta'_j)\bigr).
\]
The remaining case gives \(q=W\) directly. This puts the
actual pivot within \(C(L)(e+\delta'_j)\) of the induced
pivot, since the other bonds of the corridor are already
close. Taking the fixed regularity ceiling sufficiently
small puts the full double block in the regular case of
the averaging formula. In a fixed double-block gauge of
the induced field, the derivative in the middle bond is
then close to its flat value \(L^{-3}I\).
The regular one-row inverse gives the sharper bound
\(C(L)\delta'_j\) for the pivot increment using the
prescribed coarse value and the tested nonpivots.

Interpolate the nonpivot logarithms and solve the pivot
equation along the interpolation, keeping \(W\) fixed.
Each double block contains only its own pivot. The
regular inverse is uniform on the small one-row chart,
so the row logarithms and their derivatives along the
path are \(O_L(\delta'_j)\). The endpoints are the
induced and the actual old datum. The common rows stay
fixed. This path lies in the real continuation
polydisc about \(V^0\). The old critical orbit therefore
exists along the path, with initial orbit that of \(U^+\).
This construction starts at the fixed prospective orbit
and applies to the full real product ball on which the
released nonpivots satisfy the approximate collar bound.
For a point of the ball use the regular implicit inverse
at the induced field. Its uniform radius and derivative
bound give the required pivot values on the whole ball.
This constructs the old datum without using the
preceding guard on that ball.

Write the old orbit on this path as
\((e^{i\mathcal H(B)}U^+)^{u_B}\) in bare units, with
\(B\) the plain determining logarithms relative to \(V^0\).
The real Landau response to a row at its top \(y_c\)
has value and covariant first-difference bounds with
factors \(\xi_y^{-1}\) and \(\xi_y^{-2}\), respectively,
times \(C e^{-a d_{\mathbf B}(y,y_c)}\). Here the metric
and the tops belong to the preceding joined sequence.
The derivative is taken in the plain row logarithm,
so its input is the velocity just bounded by
\(C_L\delta'_j\). The top sum at this rate is bounded
uniformly on the separated sequences, and at most eight
rows use any one top. Thus on the old level-\(j\)
plaquette carriers at the new prepared face,
\[
 \xi_j\|\mathcal H\|+\xi_j^2\|\nabla_{U^+}\mathcal H\|
                          \le C_L\delta'_j,
 \qquad \xi_j=L^j.
\]
This estimate sums in the actual metric, with the same
bound at the first stage and at the last stage.
If the affected stencil is a distance \(D\) from the
changed row tops, keep half the rate in the sum. The
bound at a stencil whose local scale is at least
\(\xi_s\) is then \(C_L\delta'_j e^{-aD/2}\) in the
norm with weights \(\xi_s,\xi_s^2\).

The scaled plaquette estimate on this complete stencil
gives the following bound for the difference of defects
between the old field \(U^-\) and the prospective field.
The real node gauge preserves the defect norm.
\begin{equation}\label{hi:real-crossing-response}
 \left|\,|\partial U^-(p)-\mathbf1|
             -|\partial U^+(p)-\mathbf1|\,\right|
                        \le C_L\delta'_jL^{-2j}.
\end{equation}
For small coupling this is less than
\(\beta_{\rm s}F_j\varepsilon_j L^{-2j}/4\), the reserve
left by the prospective crossing test.

\paragraph{The global real collar coordinates.}
%
%
%
Fix the prospective real datum with its global and local
critical orbits admitted. Set \(q=j+1\). Write
\(S_{\rm g}=\mathcal M^jU_m^{(n+1)}\) for the averaged global
background of the integral, and \(S_X=\mathcal M^jU_{m,X}^{(n+1)}\)
for the center of its approximate collar test. The two joined
problems have the same determining rows and data in the
interior of \(X\); the difference comes from the outer joining
collar. Use the prescribed block axial readings with the
same retained node data. On the full blocks these readings
include the node factors of the real response.

In bare units let \(a_m=M\check R_mL^m\). The collar
\(K\) lies in \(\Omega_q^c\cap X\subset O_m\) of
\eqref{eq:elab-history-quiet-boxes},
which is two level-\(m\) cells inside \(X\). Include the
order-\(j\) average carriers of the fine rows, the full
double blocks of the new coarse rows, and those of any
inherited coarse row whose averaging stencil reads a changed
fine bond. These finite stencils add only \(C L^m\) to
the distance from \(K\). The old groups of tests at the
original \(\Omega_q\) seam use their full four-layer
test-cube enlargements. Before the final stage the
preceding wall comparison gives their clearance.
At the last stage use the corresponding level-\(m\)
calculation. For the collar coordinates themselves
the full endpoint stencils have the stated clearance
at every stage by the prepared-face separations.
Their prospective scales are \(q\) and, across
\(\partial Y\), \(j\). The separation from the
preceding prepared face, or from \(Z_h\) at the first
stage, keeps these carriers out of finer layers.
Weighted lengths to the joining collar are at least
\(cM\check R_m\), since no layer has spacing greater
than \(L^m\). The real joined comparison, followed
by averaging on these blocks, gives
\[
 \sup_b\left|
   S_{\rm g}(b)S_X(b)^{-1}-\mathbf1\right|
      \le C\varepsilon_\star e^{-cM\check R_m}
      =o(\delta'_j)
\]
on the collar read set. The last equality uses the frame
drift and \(\check R_m\ge(\log(\theta_m-c_5))^r\), with
\(r\ge2\). The same bound is smaller than \(\delta'_j\)
times any fixed negative power of \(\log\theta_j\).

For the same actual nonpivot value the two coordinates
are related by the exact product
\[
  V_jS_X^{-1}=(V_jS_{\rm g}^{-1})(S_{\rm g}S_X^{-1}).
\]
Consequently the global sharp cutoff at \(\delta'_j\)
implies the local approximate cutoff at \(2\delta'_j\).
On the positive local approximate branch the global
nonpivot increment is bounded by \(3\delta'_j\), which
fits the same regular chart before the sharp split.
The actual middle bond in each double block
is the same in both charts, since the other physical
bonds are fixed and the unique regular pivot solves the
same constraint with target \(V_q(c)\).

\paragraph{The sharp guard in the two charts.}
On the nonpivot bonds put \(H=(1/i)\log(S_{\rm g}S_X^{-1})\)
and write \(V_j=e^{ig_j B_{\rm g}}S_{\rm g}=e^{ig_j B_X}S_X\).
Then
\[
 e^{ig_jB_{\rm g}}=e^{ig_jB_X}e^{-iH}.
\]
In each chart let \(h\) be the density of Haar measure in
exponential coordinates, with the same normalization. Right
translation gives on the free bonds
\[
 dB_{\rm g}=\prod_b
    \frac{h(g_jB_X(b))}{h(g_jB_{\rm g}(b))}\,dB_X .
\]
The scale factors \(g_j^{\dim G}\) cancel bond by bond.
The pivot densities are evaluated on the same solved
physical fields. They multiply both sides of the Haar
identity, since the retained row targets and the graph of
the regular inverse are the same in the two charts.

On the free collar variables the global sharp guard in
the first chart is the fixed product ball
\(\max_b|g_j B_{\rm g}(b)|<\delta'_j\).
In the second chart it is the indicator of this same
inequality after substitution. The local sharp guard tests
\(\max_b|g_j B_X(b)|<\delta'_j\). Either guard implies
the local approximate test on these bonds by the centre
comparison.
To compare the two sharp guards, use the uniform derivative
of the logarithm on the regular small chart. For
\(h_X=C\sup_b|H(b)|=o(\delta'_j)\),
\[
 \big|\max_b|g_j B_{\rm g}(b)|-
             \max_b|g_j B_X(b)|\big|\le h_X.
\]
In particular, where the guards disagree both maxima
lie within \(h_X\) of \(\delta'_j\). This band already
has a bond of size at least \(\delta'_j/2\) in either
chart, the size used in the sharp-test energy estimate.
On the common real fibre the four products of the two
guards and their complements sum to one. In changing
charts each product travels with its density. A comparison
of the labelled operations must also account for the two
disagreement terms and the variables their branches keep.

On the failed global branch the cutoff substitution
can be checked bond by bond. Write \(W'\) for its new
coordinate, and use the cutoff \(\vartheta\) at scale
\(\delta'_j\). The physical local and global readings
are respectively
\[
 W'e^{i(1-\vartheta)H}
 \quad\text{and}\quad
 W'e^{-i\vartheta H}.
\]
Here \(H\) is the fixed shift of the two prospective
centres on the real fibre. The joined estimate gives
\(|H|\le\delta'_j/8\). Below the start of the cutoff,
at \(4\delta'_j/3\), the global reading equals \(W'\).
Above that start its angle is at least
\(4\delta'_j/3-|H|>\delta'_j\). Thus its sharp
failure is exactly the complement of the sharp ball
in \(W'\). For the approximate local reading the two
chord tests in the old and new coordinates agree
where \(\vartheta=1\). Where \(\vartheta<1\) both pass
by \(5\delta'_j/3+|H|<2\delta'_j\).
Apply this identity on the actual transformed mask,
with the derivative of the cutoff map in its Haar
factor and with unchanged coarse targets for the pivots.

The domain of a test and the set of independent collar
coordinates must be distinguished. The direct graded
rule tests a sharp ball on its free exterior nonpivots.
The inherited strip coordinates, including the old-wall
crossings, keep their ordinary tests in their original
chart. Its local sharp failure keeps that prefix before
the shift. Under the incident-bond convention, the
sharp mask can also include a bond crossing into the
old \(\Omega_m\). Such a bond belongs to the inherited
ordinary chart. Write \(S_{\rm o}\)
for its reference. The global sharp factor there tests
the logarithm of
\((V_jS_{\rm o}^{-1})(S_{\rm o}S_{\rm g}^{-1})\), with
\(V_j\) read in the inherited chart. There is no new
integration coordinate for that bond. The whole sharp
indicator is the free-ball indicator times the product of
the sharp readings on the other bonds of its mask. The
ordinary strip bounds at scale \(\delta_j\) do not by
themselves imply these readings at scale \(\delta'_j\).

For the final collar, on a crossing of the old
\(\Omega_m\) wall this inherited
bond lies in the ordinary plain strip. The bases of both
passed and failed ordinary cube tests are inset by a full
level-\(m\) cell. In level-\(j\) units this cell has side
\(LM\check R_m\). The read reaches at most
\(2LM_2\check R_m+1\) beyond the base, including the
incident endpoint convention. This is less than one
cell by the choice of \(M\). Hence these masks do not
read the crossing bond. Its inside endpoint is in the
stored boundary strip. The ordinary bond map on this
plain part is just the expression of the physical link
relative to the local strip centre, with Haar measure.
The explored collar factors belong to the ordinary
interior test mask, and the ordinary rim test excludes
the stored strip bonds. The newly completed fluctuation
test and the inherited approximate test can still read
the crossing on their stated masks.

There is a direct way to keep an additional sharp read
on an inherited free bond. Let \(\rho_b\) be the smallest
retained ordinary upper threshold at that bond, including
its factors in the cube and rim tests. Suppose that on
its actual tested fibre
\[
 \left|\tfrac1i\log(S_{\rm g}S_{\rm o}^{-1})\right|
                      +\delta'_j<\rho_b .
\]
On the global sharp ball the triangle inequality for
the group angle makes all these ordinary upper factors
at \(b\) equal to one. On this plain strip use the
physical group variable with its Haar density. With the ordinary centre
fixed in the inherited coordinates being changed, holding
the fresh variables fixed makes the re-centering a right
translation bond by bond. Thus the new sharp read can be
written as a fixed ball in that inherited coordinate,
with the exact change in its conditional density. The
completed tests on the other bonds remain factors of the
original integral. For application we must check the
displayed centre bound on the complete crossing stencils,
the fixed-centre property on this joint fibre, and the
analytic bounds of the resulting density on the strip
discs. In particular the difference of the ordinary and
prospective centres, rather than the difference of the
two prospective centres, enters this bound.

On the ordinary reads with both ends inside \(\Omega_m\)
the preceding comparison gives this difference at order
\(O_L(\delta'_j)\) in a common axial frame. That size
fits any of the retained ordinary thresholds, which are
fixed fractions of \(\delta_j\). The estimate for the
approximate test on a crossing used a direct test of
\(V_j\) against the prospective centre. It leaves a frame
comparison to check for the present re-centering at a
crossing. The node factors at both ends of that bond must
be those of the inherited coordinate.

Let \(x_-=b_-\), \(x_+=b_+\), and let \(y_-,y_+\) be
the centres of their two complete level-\(m\) blocks.
Write \(T_W(y,x)\) for the one-step mean contour of a
level-\(j\) field \(W\) in such a block. Its axial reading
of the crossing is
\[
 \mathcal R_b(W)=T_W(y_-,x_-)W(b)T_W(y_+,x_+)^{-1}.
\]
For any gauge \(u\) the covariance of the mean contour gives
\[
 \mathcal R_b(W^u)
     =u(y_-)\mathcal R_b(W)u(y_+)^{-1}.
\]
Thus the comparison can be made through the factors at
the two block centres. If the two complete read blocks
lie in the common joined annulus with its matching rows,
the joined comparison controls the real potentials there.
At the old wall of the quiet component the two blocks
and their order-\(j\) carriers are within \(C L^m\) of
the wall. The outer band is separated from that wall
by the quiet-box margin, and the inner boundary lies
past the enlargement by two cells of the old member.
These margins place the two blocks in the common
interior of the annular comparison. The old exterior
block lies in the old member of level \(j\), and the
clearance from the prepared box puts both blocks in
the new member of level \(m\). Orient the crossing
from the exterior block to the interior block and
write \(c=\langle y_-,y_+\rangle\) for the coarse
bond between their centres. It is an outward row
of the old member at level \(m\), so it is retained
even when the fine inward rows are omitted.
It becomes an interior row after promotion of the
exterior block and has the same value \(V_m(c)\).
The old exterior centre is a node of layer \(j\)
by its internal level-\(j\) rows. The old interior
centre is a node of layer \(m\).
Both complete blocks lie inside the common part
of the two joined sequences away from their joining
bands. Their rows and the levels of these nodes
are therefore unchanged in the annular comparison.
The average factorization on the full two-block carrier
controls the level-\(j\) links and the factors at the two
centres. At the interior centre of the old sequence the
restoring gauge has value \(w(y)=v_{j+1}(y)^{-1}\).
In the averaged comparison its product with the
level-\(j\) factor is
\(w(y)v_j(y)=\mathsf v_j(y)^{-1}\), since
\(v_{j+1}(y)=v_j(y)\mathsf v_j(y)\). The exterior
centre lies in the old layer \(j\), even though it lies
on the level-\(m\) lattice. There the product is the
identity, as it is at the exterior fine node before
auxiliary axialization. The one-step means have uniformly bounded
derivatives on the same regular blocks. Applying the
last identity to both comparisons with the joined
reference shows that the axial readings of the ordinary
and preceding backgrounds differ by at most
\(C_L\varepsilon_\star e^{-cM\check R_m}\) on that
carrier. For the comparison with \(S_{\rm g}\),
start the real response on the old determining set
at the datum induced by the prospective field, in
its prescribed frame. The physical row increment
\(\beta(c)\) is zero throughout the interpolation.
The reduced increment \(\beta_\natural(c)\) can
be nonzero because its outward-row conversion
reads the changing internal links of the exterior
block. Use the response to \(\beta_\natural\)
and restore at each old node \(y\) by \(v(y)^{-1}\).
The datum of the uncorrected response is the desired
physical datum turned by \(v\), so this restoration
fixes the prescribed value of \(c\) and returns both
node frames to the inherited ones. At the two
centres the average factorization now leaves the
identity at the old layer-\(j\) node and
\(\mathsf v_j^{-1}\) at the old layer-\(m\) node,
as in the annular comparison. The latter is
\(O_L(\delta'_j)\) from the identity by the local
response on the full carrier. This gives the
same bound for the difference of the two axial
readings. Use a fixed residual gauge for the
response to start at the chosen prospective
representative. Both restored node values are
then prescribed for the whole path.

The original strip frame has a further contour at
the exterior endpoint. The ordinary joined strip
has nodes at all level-\(j\) sites of that block.
Its level-\(j\) internal bonds keep the actual
values of \(V_j\), which determine the old contour
from the block centre to the endpoint. The axial
condition is imposed at the nonnodes of the
ordinary interior blocks. An internal
path has no pivot, since each pivot lies in a corridor
between distinct blocks. The
interior endpoint is axial already. On the prospective
block the contour is the identity. The difference of
the exterior contour from the identity is therefore
bounded by \(C_L\max_e|V_j(e)S_{\rm g}(e)^{-1}-\mathbf1|\),
where \(e\) runs only over the internal bonds of that
exterior block. For the native collar mask the changed
exterior block next to the old wall lies wholly in the
active collar. The members consist of whole blocks, and
the clearance from the old wall to the prepared face is
greater than the width of this block, also at compressed
levels. Every internal contour bond of the exterior block
is therefore in the passed approximate test. That test
and the comparison of the two prospective centres bound
the maximum by \(C_L\delta'_j\). This uses the internal
paths even when the sharp test on the crossing fails.
Undoing the exterior contour gives the desired bound in
the inherited frame on the old-wall crossings just
described. If using a sharp mask with additional
inherited reads, we must check their complete carriers
and admit the resulting conditional density in the
complex stage space.
In this auxiliary axial comparison the stored fields keep
their original values. In a physical re-centering of a
crossing nonpivot, solve its old coarse constraint for
the pivot at the same fixed target. The native exterior
block bound and the inherited ordinary bounds put the
two-block read in the original regular pivot chart,
where the change of the solved pivot is bounded by a
constant times the change of the nonpivots. Keep its
elimination factor and the Jacobian of the original
strip substitution. For a crossing in the ordinary plain strip the
old rim factor excludes that strip bond. A completed
ordinary fluctuation test can read it, at threshold
\(\delta_j\); its factor at the crossing is then one on
the re-centred sharp ball by the centre bound above.

The induced datum argument also applies to the global
critical background itself. Read the global prospective
configuration on the global old determining set. Its old
units are no larger than its prospective units. As in the
localized argument, the complete tested carriers give an
admitted old datum and the nested constraints give
criticality at that datum. Interpolate the nonpivots of
the global chart and solve the corresponding pivots.
The modified rows are at level \(j\); the point response
sum in level-\(m\) units is bounded uniformly in the number
of components and in the torus. Thus the real global
first background in the collar integral is the regular
old orbit on this path. The localized old guard can be
checked at each of these physical row values by the
localized path estimate. Its value was not used to
continue the old orbit along the chart.

At the final stage the inherited strip is kept in the ordinary coordinates.
Some of its rows lie close to \(\partial X\), where the
estimate comparing \(S_{\rm g}\) with \(S_X\) need not
hold. An inherited double block that reads a changed
collar bond was already included in the interior
comparison. The ordinary and preceding global backgrounds
on the rest of the strip can be compared on a common
annulus. Use an outer box obtained by enlarging \(X\) by
half a level-\(m\) cell and rounding on the top tower grid.
Remove the inner box
\((\Omega_m^{\sim2})^c\cap X\).
The outer band is in the unchanged top layer, including
the neighboring cells outside \(X\). The inner band is in
the old layer \(j\). The earlier stages change rows over
the old components inside \(G_{m-1}\subset H_m\). Even
their full row carriers stay inside the removed box by
the six-cell gap. The two boundary
bands have the widths required for the joined comparison.
At the physical collar value, the ordinary and preceding
global problems read the same rows and data throughout
this annulus. Its interior contains the ordinary strip
rows touching the old \(\Omega_m\) and their full double
blocks, including those reaching across \(\partial X\).
This gives the comparison of the ordinary orbit in its
own gauge on those stencils. The inherited Schur complement
and the retained rim and collar indicators are evaluated
in the original variables of the strip.

\paragraph{The real endpoint triples.}
\hypertarget{hi:prepared-real-endpoint}{}
At a stage of the depth-\(m\) problem set \(q=j+1\)
and put \(e_l=\check R_l^2\varepsilon_l\). This allows
the complete averaging footprints at the original
exterior walls of a compressed stage. The tests with
sharper thresholds keep those thresholds.

%
%
%
The input at a point of a history includes the fields
of the operations still carried by it. A datum
plaquette uses the actual determining links and,
for a completed bond, the average of the lower
field on its full double block. These values agree
with the fields in the tests by the retained
averaging constraints. In the convention excluding
inward rows these plaquettes exhaust the datum tests.
For each test place in a stopped component use its
least-level departure. If the component is untouched,
or the full footprint misses the inset of its stage
test region, the ordinary density test supplies the
bound. The same ordinary test is used at untouched
higher levels and on the unchanged action top.
An inner footprint in a tested shell uses the last
positive stage test at that departure. At a restored
core the bound comes from the retained reference
cutoff and the surrounding ordinary tests.

For a departure at the sharp collar split after
\(n\) integrations, put \(j'=i_\flat+n\), where
\(i\) is the level of that operation. The retained
positive stage test is \(\chi_i^{(n+1)}\). A
level-\(j'\) test in its inset with a corner in
the collar uses also the retained approximate
bound at \(2\delta'_{j'}\). In the translated
coordinates this bound is on the reconstructed
field of the localized collar chart. The
failed sharp indicator is carried as a separate
factor. The ordinary test supplies a footprint
outside the inset. At the promoted level the
retained coarse row is fixed by its averaging
constraint. In the direct exterior collar
rule, a retained determining link of level \(j'\)
has no endpoint in the next old member. A
completed bond uses both full blocks outside the
corresponding next member. Thus the determining
reads use the collar bound on its exterior slots
and the retained values on the other slots.
The original strip tests remain the suppliers for
the old-wall fluctuation reads.

These choices use the tests carried at the relevant
moment. Before the first stage use the untouched
datum. While a positive stage test is carried it
supplies the ordinary assignment. After removal
and before integration the collar assignment uses
the new positive stage test and the approximate
collar test on the old collar variables. The new
ordinary assignment uses the new positive stage
test on the promoted datum. A standing departure
below the current waiting window keeps its earlier
determining reads and averaging completions.
Gauge moves conjugate the holonomies. A refresh
with the inner datum retained leaves its test
values in that inner slot. Where restoration
replaces the field, the tests of the restored
datum provide the new bound.

The witnesses on the complete footprints give
\(2\psi_l\varepsilon_l\) for the datum defect.
Here \(\psi_l\) is \(O(1)\) for an ordinary
witness or a restored core and is bounded by
\(C(\log\theta_l)^r\) on the remade shells.
The latter bound uses \(L_\flat^2<L\) and the
compressed interval of length at most
\(\log_L\check R_{k_0+1}\), with the frame
comparison of the running scales. Thus each
test gives a bound \(C e_l\) directly. Fix an
entering point on these retained tests and on
the constraints with coarse target \(V_q\).
This is the input for varying the collar. To
remove the preceding guard on an enlarged
support, use the separate comparison in
\cref{hi:preceding-guard-support}.

Let \(U_{\rm in}\) denote the regular global old
configuration at that input. The datum defects are
below the fixed real existence ceiling at weak
coupling. Apply the
\hyperlink{hi:real-datum-profile}{real profile estimate}
in its layer form.
In bare units it controls
\(L^{2\lambda(p)}|\partial U(p)-\mathbf 1|\) and
\(L^{3\lambda(b)}|J^0(U)(b)|\) at each corner or
endpoint by the profile there, where \(\lambda\)
of a stencil is its largest corner layer.
For the datum bounds just obtained the profile
has the form
\[
 C\sup_l e_l
   \exp\bigl[-\delta\,d_{\mathbf B}(x,\mathcal L_l)\bigr],
 \qquad \mathcal L_l=\Omega_l\setminus\Omega_{l+1}.
\]
The place of a tested level-\(l\) plaquette is in
\(\mathcal L_l\), and its corners lie in layers
\(l-1,l\). From layer \(s\) the path to this read
crosses the intervening layers, giving distance
at least \(c(|l-s|-2)_+\). The frame bounds make
\(e_l/e_s\) grow at most polynomially in
\(|l-s|\). Hence the profile in layer \(s\) is
at most \(C'e_s\). The zeroth layer uses the bare
datum tests. Apply the same profile on both stage
sequences.

In the prospective datum all holonomies below level
\(j\) are unchanged. At level \(j\), the new
first-kind row set is a subset of the old one, and
the second-kind completions still use the same
level-\((j-1)\) layer. At level zero there is no
completion from a finer layer. Above level \(q\)
the original row values and masks are retained.
For a completion at level \(q+1\), its full
two-block read stays in the old \(\Omega_q\),
by the old gap from \(\Omega_{q+1}\) to
\(\Omega_q^c\). Thus the enlargement at \(q\)
does not change that completion. At higher
levels neither layer of the completion changes.

Every tested holonomy at level \(q\) equals the
plaquette of \(\mathcal M^q U_{\rm in}\) on the
relevant complete stencil. A determining row
that met the old \(\Omega_q\) keeps its value.
Each newly imposed row has both blocks outside
the old \(\Omega_q\) and within the old
level-\(j\) member. Its retained constraint reads
the old determining \(V_j\) on these two blocks.
For a second-kind bond the two blocks are
outside the prospective member at \(q\) and
in the preceding layer, so they also use the
old level-\(j\) values. In either case the
identity with the average follows from the
actual retained constraint or its one-step
completion. For each tested plaquette the order-\(q\)
carrier stays inside the unchanged level-\(j\) member
by the separation of the prospective sequence. The
level-\(j\) profile and the adjacent-scale conversion give
bare curvature at most \(C e_jL^{-2j}\) on that carrier.
Averaging there gives
\(C(L)e_j\le C'(L)e_q\) for the datum defect at
the promoted level. The artificial face lift of
the localized problem does not enter this global
calculation. We have thus
obtained the levelwise datum estimate for the fixed
prospective global orbit and, by the weighted
maximum, its normalized curvature and current bounds
with \(C e_s\) in layer \(s\).

For the varied old datum, begin at the averages of
this prospective configuration on the old determining
set. For an old level-\(l\) tested plaquette, \(l\ge1\),
its order-\(l\) carrier is inside the old member of
level \(l-1\). The prospective scale on this carrier
is at least \(l-1\). In changing its units to \(l-1\)
the powers \(L^{-2(s-l+1)}\) absorb the frame
variation in \(e_s/e_{l-1}\), so averaging gives a
defect bounded by \(C(L)e_{l-1}\le C'(L)e_l\).
At level zero the bare profile gives the estimate
directly. Now vary the actual old collar coordinates
in the regular fibre chart. Only level-\(j\) rows
move, by \(O_L(\delta'_j)\) including the reconstructed
pivots. Only tests of levels \(j\) and \(q\) can
read these rows, either directly or through a one-step
completion. On their full double blocks the regular
average and the product of the four group factors
change by \(O_L(\delta'_j)\). This is smaller than the
datum allowance \(C e_l\) at these adjacent levels.
The old orbit on the varied fibre has the same
layerwise bound from the weighted maximum.
These are bounds for the two endpoints from the
entering datum profile of the history. In using them
through the stages, take the profile from the actual
retained tests at each entrance, as in the levelwise
datum bound. Repeatedly multiplying the adjacent
comparison constant would not give a uniform profile.

Consider now any stage clause of scale \(s\) for either
real endpoint. In a prepared layer and in an untouched
higher layer its reading set is in the actual member at
scale \(s\). In the compressed clause of an original
layer \(i\), the assigned scale is \(s=j\) for the old
stage and \(s=q\) for the prospective one. On the remade
component the actual member
at that scale contains the reading set, since the prepared
box is outside \(\Omega_{k_0+1}\). For a reading on an
unchanged part, use at least its original scale \(i\).
The loss of scale \(s-i\) is at most
\(\log_L\check R_{k_0+1}\), by the compression
identity. Thus one can use the actual member at
scale \(a\le s\) with \(a=s\) except for this
possible compression loss. Use the same calculation at an
earlier address whose determining regions are contained in
the endpoint sequence, with the compression interval of
that address.

For the pointwise bounds at scale \(s\), changing
from the scale \(a\) costs at most \(L^{3(s-a)}\).
This accounts for the physical current, whose scale
power is three, as well as the curvature, whose power
is two. At the plain triple \(A'=0\), so the two
potential clauses hold identically. For each of the
tower clauses, the actual level-\(\mathsf p\) tower
on a localization domain is formed by intersecting the
stage sequence, truncated at \(\mathsf p\), with its
localization tower. Hence it is subordinate, even if
\(\mathsf p\) exceeds the layer of the read. The
restriction identity \eqref{mb:restriction-identity}
identifies its real background at the own averages with
a gauge copy of the restriction. It therefore has the
same curvature and the same current at every full
current stencil. In level-\(\mathsf p\) units the
current gains \(L^{3\mathsf p}\) relative to bare
units. This gives the scale required in the tower
clause, since its allowed factor
\(L^{s-\min(\mathsf p,s)}\) is at least one. The
curvature factor in that clause simplifies to
\(L^{-2s}\).

For a local gauge cube of side \(C_{\mathrm g}ML^s\) wholly
in the actual member of level \(a\), use the gauge estimate
at \(b=\max\{0,a-1\}\). The halo of thickness \(8\rho L^b\)
lies in \(\Omega_b\) by separation, or in the ambient bare
torus at \(a=0\). The powers \(L^{-2(l-b)}\) and
\(L^{-3(l-b)}\) for its actual layers \(l\) absorb the
frame drift. Thus the curvature and current in those units
are bounded by \(C e_a\) on their complete halo stencils.
The side in level-\(b\) units costs \(O(M L^{s-b})\), and
conversion of the potential and its first differences costs
at most \(L^{2(s-b)}\). The normalized gauge bound is
\(C M L^{3(s-b)}e_a\). Use the torus side condition
\(2C_{\mathrm g}ML^s+16\rho L^b<\mathsf s_T\). The fixed
curvature ceiling and the product of the gauge radius in
level-\(b\) units with \(Ce_a\) fit the smallness
requirements at weak coupling. Restrict
the resulting gauge to the prescribed read set.

The frame bounds across the waiting window give
\(\check R_{k_0+1}\le C(\log\theta_s)^r\) on the
compressed interval, and \(e_a\le C e_s\) there.
Consequently the pointwise and tower bounds, normalized
at the nominal scale, are at most
\(C g_s(\log\theta_s)^{\mathfrak p_0+5r}\), and the
gauge bounds are at most the same quantity times \(M\).
The stage choice \(q_{\rm rad}\ge\mathfrak p_0+7r+1\) for
\(\alpha_{0,s}=C_0g_s(\log\theta_s)^{q_{\rm rad}}\), and the
positive coefficient floor \(1-4\beta_{\rm s}\) for
the choice \(0<\beta_{\rm s}<1/4\) in the allowed window,
put the plain triples inside half the radii for small
coupling. At earlier hereditary addresses use their wide
subordinate sequences, the full gauge cubes, and their fixed
positive coefficient floor and nonnegative lifts. The same
bounds fit the smaller adjacent running radius, with its
bare-level value at level zero, and the half radius of the
fixed local regularity clause. This verifies the real
clauses on the specified carriers.

\paragraph{The complex substitutions in the collar integral.}
Keep the real factor \(U\) fixed and use independent complex
potential and current coordinates. The source for the old
critical problem is the nonlinear datum increment on the full
double blocks of the collar. Solve the regular one-row
inverse at the applicable background. For a free field
\(B\) this gives the completed logarithmic datum
\(B'=g_j\mathsf C B-\mathsf h\widetilde{\mathsf D}
(g_j\mathsf C B)\). Take its corresponding reduced datum
for the reduced response. On the complex chart
the inverse gives \(\|B'\|_\infty\le C_Lg_j\|B\|_\infty\),
so the source in the box \(\|B\|_\infty<Cp_0(g_j)\) is
within the small response bound. Using the nonlinear row is
necessary for the equality at the physical endpoint.

%
%
For the current increment the reduced first variation
has a local linear term in addition to its supported
parts. Write \(D_{\rm lin}B'=\hat Q^{\mathsf T}\mathsf aB'\)
for this row term. The supported bracket of the raw
response is \(\mathbf F^\flat+D_{\rm lin}B'\).
Gate each row of the linear term by the product of
the switches in a connected set containing its
complete source and output stencils. Include the
plain-data and parent carriers in this set. The
neighbor margins and the one-step row charts give
bounded size on the aligned standard grain. Denote
the resulting linear term by \(D_{\rm lin,s}B'\).
It agrees with \(D_{\rm lin}B'\) at full switches,
and a nonzero row term lies in one active wall
component. In the current formula use
\[
 \mathbf J_s=
 \mathsf f(-i\xi\,\operatorname{ad}_{\mathbf H})^{-1}
 \bigl[\mathbf J+
       (\mathbf F^\flat+D_{\rm lin}B')-D_{\rm lin,s}B'\bigr].
\]
The raw formula gives exactly the current of the
perturbed real configuration at full switches.
At a cut the supported bracket and the gated
linear term use only their respective components.
The inverse \(\mathsf f^{-1}\) acts bond by bond,
so the increment from \(\mathbf J\) is supported
on those components. The product of the bounded number of
switches costs a fixed factor in the current
bound. Use this continuation to form the
component increments \(\mathsf j_c\) alongside
the localized potential increments.
In the weighted current estimate write
\(D_{\rm lin,s}=D_{\rm lin}G_s\), where \(G_s\)
multiplies each source row by its carrier product.
The increment from \(\mathbf J\) is the inverse
\(\mathsf f\) applied to
\(\mathbf F^\flat+D_{\rm lin}(I-G_s)B'\), plus
\((\mathsf f^{-1}-I)\mathbf J\).
The local row bound for \(\hat Q^{\mathsf T}\)
and the fixed switch cost bound the added row
term in cubic current units. The small
exponential chart bounds the last term by the
weighted response times the current radius.
The gauge turn has the same estimate for the
additional factor \(\operatorname{Ad}(u)-I\).

%
For a wall component \(\mathfrak w\), evaluate the
response with cube switches \(s\mathbf1_{\mathfrak w}\).
Call its potential \(\mathbf H_{\mathfrak w}\) and its
current increment \(\mathsf j_{\mathfrak w}\), where the
increment is the current in the displayed formula minus
the input \(\mathbf J\). Use the same component parameter
in both output sums,
\[
 \mathbf A=\sum_{\mathfrak w}\lambda_{\mathfrak w}
                         \mathbf H_{\mathfrak w},\qquad
 \mathbf K=\sum_{\mathfrak w}\lambda_{\mathfrak w}
                         \mathsf j_{\mathfrak w}.
\]
The independent input current is kept in its original
slot. The component parameter multiplies the response
after solving the equations. This is the interpolation
used in the common radius budget.

Use one levelled triple on the union of the carriers
of the inherited reads, restricting it to each
reader's address. If \(W=e^{i\eta_t\mathbf A}\mathbf U\)
is the unlevelled configuration, take
\(u(y)=v[P;\mathbf A](y)^{-1}\) at the pillar top \(y\),
using the same switched potential \(\mathbf A\) that
defines \(W\). At a stored site \(a\) of this pillar take
\[
 u(a)=\tau[\mathbf U](a)^{-1}u(y)\tau[W](a).
\]
Indeed \(\tau[W^u](a)=u(y)\tau[W](a)u(a)^{-1}
=\tau[\mathbf U](a)\). Write \(\mathbf A^\circ\) for
the relative logarithm in \(W^u=e^{i\eta_t\mathbf A^\circ}\mathbf U\).
Use the same analytic gauge in the current substitution
\[
 \mathbf J+\mathbf K^\circ=\operatorname{Ad}(u)
 \left(\mathbf J+\mathbf K\right).
\]
At the physical endpoint the corrected configuration and
current agree with those of the actual old critical orbit
in its adapted frame. The transport identity keeps the
spectators of each earlier term in their prescribed frames.
Each inherited term uses its own attachment. The new
crossing test reads the fixed prospective datum and adds
no pin.

On the integrated strip fields the substitution in an
inherited potential uses the scalar multiplier
\(\lambda_{\mathfrak w}s(\square_b)\). It multiplies the
far argument and serves as the switch for an unchanged
near argument. The real cutoffs act in their integration
charts. Keep the inherited
Schur form and take the difference of its two evaluations
as a potential on these discs. Its two evaluations use the
same stored variables. The shift of the local free coordinate
has a separate source radius of the form
\(c(\log\theta_j)^{-e_y}\delta'_j\). In the local sharp
rule, keep the free collar coordinate \(B_X\) in its
original real product ball \(|g_jB_X|<\delta'_j\). For a
free shift \(\mathsf y\) use the global increment
\[
 g_jB_{\rm g}(B_X,\mathsf y)
  =\frac1i\log\bigl(e^{ig_jB_X}e^{-i\mathsf y}\bigr)
\]
on each released nonpivot, and then complete the pivot at
the fixed coarse target. For the physical value of
\(\mathsf y\), take \(e^{-i\mathsf y}=V_XV_{\rm g}^{-1}\)
with \(V_X,V_{\rm g}\) the local and global centres in the
same frame. This is the orientation for left link logarithms.
Its bound is \(C\varepsilon_t(\log\theta_t)^C
e^{-cM\check R_t}\) in a component of level \(t\).
In its waiting window \(\theta_j\) and \(\theta_t\) are
comparable with one logarithm. The exponential therefore
puts the actual shift inside the source disc at every
stage, including the bottom. On a complex section its
source is the localized shift map formed on that section.

On the real free-source disc the trial global angle is
at most \(\delta'_j+|\mathsf y|\) in this same ball.
Combining with the actual centre comparison gives the
approximate bound \(2\delta'_j\) on its full free mask.
The old completion and rim tests act on the original
strip coordinates. At fixed admitted prospective datum
the guard implications, including the crossing reserve,
apply on this entire real ball. This permits the use of
the same integration domain for all the real trial shifts.
At the physical shift it gives the original local sharp
integral. The source is then continued in the integrand
while retaining those real integration coordinates.

%
%
On the passing local-sharp branch the smooth
stage Haar factor has a separate calculation.
Write \(z=g_jB_X\) and
\(z_{\mathsf y}=g_jB_{\rm g}(B_X,\mathsf y)\).
For the Haar density \(h\) normalized at zero,
the change from \(d z_{\mathsf y}\) to \(d z\)
has logarithm
\[
 L(z,\mathsf y)=\log h(z)-\log h(z_{\mathsf y}).
\]
Indeed right translation preserves Haar. The
logarithms are taken on the fixed small chart
where \(h\) is nonvanishing. Its evenness gives
\(L(0,\mathsf y)=-\log h(-\mathsf y)=O(|\mathsf y|^2)\).
The difference \(L(z,\mathsf y)-L(0,\mathsf y)\)
vanishes at either zero argument. Integrate the
mixed derivative along the two segments in the
product-logarithm chart to bound it by \(C|z||\mathsf y|\).
On the fixed real ball \(|z|<\delta'_j\) and the
independent shift disc this is at most
\(C(\delta'_j)^2(\log\theta_j)^{-e_y}\).
On an enlarged complex section use the actual
bound \(g_jB_*\) for \(|z|\) on that section.

When two moves use the same free link, let
\(z_0,z_1,z_2\) be its successive logarithms
in their actual order. The two smooth log
factors add to \(\log h(z_0)-\log h(z_2)\).
Subtract its centre at \(z_0=0\), with the
intermediate value \(z_1(0)\) induced by the first
map. Keep the centre once on the shift carrier.
After substituting the localized shifts its series
has two or more shift factors for a single
zero-centred translation, while the difference
has at least one. For the composition use the
ordered product-logarithm chart and the same
rooted union of the shift carriers. A cutoff
Jacobian between the smooth changes is carried
as its own factor.

For the global sharp branch keep \(B_{\rm g}\) as the
variable of the fixed sharp ball on the free nonpivots.
Its real local approximate test on those bonds is then
identically one. In comparing with the
analytic integral having a free shift source, the candidate
for this fixed ball is the source value zero with its
integration variable identified with \(B_{\rm g}\).
The physical integral with a local sharp guard instead
uses \(B_X\) as the variable of the fixed ball and the
actual shift in the analytic factors. The cutoff stays
fixed in both representations on the free coordinates.
In the direct graded rule the sharp mask is confined to
the free exterior nonpivots. The inherited crossing
coordinates remain in the ordinary strip integral with
their completed tests. A test on the broader incident
mask gives a separate real kernel.
If the full sharp mask also reads an inherited bond, we
must show that its exact pulled-back indicator is among
the fixed real factors of the analytic integral, or prove
it equals one on the support being integrated. In the
expansion with Gaussian tails each real cutoff uses its
prescribed carrier and accepts the corresponding smaller
ball. Freezing the real datum alone leaves those carrier
and tail checks to be made. The failed
global branch uses its own real strip map with the
transported tests.
Where that map has a cutoff, its derivative enters the
density in addition to the Haar factors.

\paragraph{Substituting the physical shift in the stage logarithm.}
Keep the shift $\mathsf y$ independent while forming the
analytic logarithm of the collar integral on the fixed real
domain. On the admitted graded section the physical value at
an exterior nonpivot $b$ is
\[
 H[P](b)=\sum_W H(b;W;P|_W).
\]
Each summand reads a decoupled response to the localized
difference of data. Its full carrier $W$ joins $b$ to the
boundary of the localization. At an ordinary strip the
resulting lower bound on the length is $d_j(W)+1\ge\check R_j$
in the strip units. In a stage it is $d_t(W)+1\ge\check R_t$
in the component units. In the respective units the pieces
have bound $C(\alpha_{0,j}+\alpha_{1,j})e^{-cM(d(W)+1)}$.
This construction requires the actual endpoint centres and
their decoupled reads on the common graded section. The
scalar annular bound on their difference alone would give
no such carrier assignment.

Let $F_X(\mathsf y)$ be a carrier term of the free-shift
logarithm. For a shift piece indexed by $a=(b,W)$ write
$h_a=H(b;W;P|_W)e_b$. The contribution using a finite
nonempty set $A$ of shift pieces is
\[
 \sum_{J\subset A}(-1)^{|A|-|J|}
                       F_X\left(\sum_{a\in J}h_a\right).
\]
Assign it to $X\cup\bigcup_{(b,W)\in A}W$, with the fixed
enlargements of the stage reads. Each nonzero difference
uses only the shift slots read by $F_X$. If the full
deformation $\mathsf y=\sum_{a\in A}t_a h_a$ is admitted
for $|t_a|<R_a$, $R_a>1$, Cauchy's estimate bounds this
mixed difference by the bound for $F_X$ on the deformation
disc times $\prod_{a\in A}(R_a-1)^{-1}$. Indeed the
alternating sum selects exactly the Taylor terms of positive
degree in every $t_a$.

The exponential bounds on the response pieces and their
lower length allow part of the decay to be spent on the
larger union carrier. In the stage sum use the rooted
closing load
\[
 \sup_\Delta\sum_{a=(b,W):\,c(b)=\Delta}
              \frac{e^{C r_\natural\varrho_{\rm a}\# W}}{R_a-1},
\]
where $c(b)$ is the anchor cell of the shift slot and the
exponent is the graded carrier price. Denote this load by
$\varepsilon_{\rm cl}$. The dilation radii must fit in the
free-shift disc and satisfy $C N^*\varepsilon_{\rm cl}\le1$,
where $N^*$ bounds the number of anchor cells of one component.
Use the native stage collars and the island hulls as exempt boxes,
with the stage separation, bounded hull multiplicity and adjacent
grades differing by at most one. The separated-component packing
estimate applies to the tree of $X$ with its prescribed trace
enlargement. Include the full anchor reads in checking this
enlargement. If the tree
$\Gamma$ is counted at level $m$ and the operation is at level
$t$, the weighted count of anchors, with weight
$\epsilon^{(D-1)_+}$ at wall distance $D$, is bounded by
\[
 N^*\left(1+\frac{4|\Gamma|L^{m-t}}{\mathsf N}
                   +\frac{|\Gamma|+1}{L\mathsf N}\right),
 \qquad 0<\epsilon\le c.
\]
The weight is one for a directly read anchor at wall distance
at most one. A fixed larger wall distance costs only a fixed
inverse power of $\epsilon$. In component units this gives
$C N^*(1+d(X)/\mathsf N)$ for the directly read anchors.
For graded trees keep the indicated scale conversion when
charging the length. For fixed $X$ the sum over subsets of
shift pieces is bounded by the product of one plus their
weighted Cauchy costs. Its logarithm is at most the sum of
those costs. The anchor count thus gives a bounded prefactor
and a loss $O(1/\mathsf N)$ in the spatial rate. In summing
over terms refiled at the same enlarged carrier, use the
anchored tree attachment count and the summed-to-pinned norm
estimate of the stage expansion. The spatial rate losses add
over the $\mathsf N$ stages, whereas the bounded prefactors
multiply individual stage amplitudes before those amplitudes
are summed.

Keep $F_X(0)$ on $X$. The nonempty anchored differences sum
to $F_X(H[P])-F_X(0)$ by telescoping, first for finitely many
pieces and then by absolute convergence in the closing norm.
The resulting terms,
with these enlarged carriers, are the ones passed to the
next stage and to the source oscillation norm. The
integration domain and the real factors in the stage law
stay those used in forming the free-shift integral.

If the inherited crossing coordinates are re-centred,
write \(\mathbf z\) for the old combined strip and collar
coordinates and \(\widetilde{\mathbf z}\) for the coordinates
with fixed sharp balls. The ordinary centre can depend on
the new collar variables even when it is fixed in the
inherited variables being translated. Its first derivative
in those collar variables contributes to the linear part
\(\mathsf T\) of \(\mathbf z\) in the new coordinates. For
old precision \(D\), the corresponding quadratic becomes
\(\mathsf T^{\mathsf T}D\mathsf T\), with any constant
shift accounted for in the mean and the scalar. If
\(D\ge\gamma I\) and \(\mathsf T\) has a bounded real
inverse, the new real floor is at least
\(\gamma\|\mathsf T^{-1}\|^{-2}\). Its covariance is
\(\mathsf T^{-1}D^{-1}\mathsf T^{-\mathsf T}\).
These identities show where the crossing correction
enters the Gaussian estimates. In particular the linear
mixing need not carry a small coupling factor. We must
bound its local matrix entries and the nonlinear remainder
on the declared sections, including the cut and switched
sections, before using the stage expansion for this
re-centred density.

Use the ordinary endpoint frames for the substitution in
the strip coordinate. In physical logarithmic units write
\(h(\beta)=\frac1i\log(S_{\rm g}S_{\rm o}(\beta)^{-1})\),
where \(\beta=g_j B\). The response in reduced row
coordinates has a coupling-independent analytic radius.
To use that radius here, suppose the completed row map,
the ordinary strip reference and the two endpoint
readouts give this same fixed-radius property for \(h\)
on the section, say \(\|\beta\|_\infty<R\), with
\(\|h(\beta)-h(0)\|_\infty\le C R\) there. Cauchy's
estimate along the line through \(\beta\) then gives
\[
 \|h(\beta)-h(0)-\mathsf L\beta\|_\infty
          \le 2C R^{-1}\|\beta\|_\infty^2,
 \qquad \|\beta\|_\infty\le R/2,
\]
where \(\mathsf L=Dh(0)\). The logarithm of
\(e^{ig_j a_{\mathrm{new}}}e^{ih(g_jB)}\) gives the old
strip coordinate. Bound all the combined coordinates by
\(T\), including the ordinary slots left untouched.
Take \(g_jT\le R/2\) and the product logarithm in the
small chart. If \(\|h(0)\|_\infty\le g_j d\), then after
subtracting \(a_{\mathrm{new}}+\mathsf LB+h(0)/g_j\) the
remaining scaled coordinate \(r\) satisfies
\[
 \|r\|_\infty\le C g_j\bigl(R^{-1}T^2+T(T+d)\bigr).
\]
The first term is the Cauchy remainder and the second
comes from the product logarithm. The remainder is
supported on the re-centred slots. If \(D\) has bounded
absolute row sums, the change from this remainder in
the quadratic is bounded by
\(C N_{\rm cr}[(T+d)\|r\|_\infty+\|r\|_\infty^2]\),
where \(N_{\rm cr}\) counts those slots. The log of the
Haar quotient costs at most \(C N_{\rm cr}g_j^2(T+d)^2\)
in the small exponential chart, since the log density
has zero first derivative at the identity. For fixed
\(R\) and logarithmic \(T,d\) these remainders are small
per crossing. Localization still requires the complete
carriers and decay of the coefficient maps.

For example, a uniform derivative bound
\(|\partial_c h_b|\le C e^{-\mu d(b,c)}\) on the fixed
source disc gives a second derivative bound on its
half-disc by Cauchy in one coordinate. Interchanging the
coordinates gives both exponentials, so that their
geometric mean bounds the same derivative:
\[
 |\partial_{c'}\partial_c h_b|
 \le \frac{2C}{R}
           e^{-\mu[d(b,c)+d(b,c')]/2}.
\]
The second-order remainder of \(h_b\) is bounded by
\(C/R\) times the square of the corresponding weighted
sum of \(g_j|B_c|\). Divide by \(g_j\) in the scaled
coordinate and include the product logarithm as above.
The source radius is needed here for
the centre map and its derivatives. The old entries are
still evaluated within their original strip discs.

On a fixed real ordinary chart the input to this centre
map can be formed in physical units. Solve the one-row
pivot equations on their fixed regular charts. On the
joined sequence a row reads a retained row, or the lift
of a named parent row. The latter is the identity or
the oriented parent link, so it is holomorphic in the
parent link on the same fixed geometry. Take logarithms
relative to the ordinary datum at zero and form the
reduced row input. These local operations preserve a
fixed small disc where the regular row and lift charts
have uniform room. Use the holomorphic critical family of
the ordinary joined sequence at its admitted real datum.
Read the averaged endpoint contours in the ordinary
strip frame. This gives the centre germ at that fixed
real reference, using the real identification of the
critical family at its physical data.
The annular comparison supplies its value at zero. To
use this construction on cuts, each parent row and the
complete endpoint stencils must remain in the prescribed
carrier, and the ordinary reference itself must be
provided by the compatible cut rule. Exponential real
comparison of two references does not give the exact
local dependence required of that rule.

For the ordinary localized branch at depth \(q=j+1\)
the strip join is formed on \(W_q=Z_q\cap\Lambda_j\),
with the corresponding cell intersections. Use \(W_1=Z_1\)
at the first ordinary step. For \(q>1\) this region
lies in the original \(\Omega_j\). The earlier
preparations change members at levels at most \(j\)
only within the old live sets contained in \(Z_j\).
Those changes miss \(W_q\). The preceding sequence
has yet to change the members at levels \(q\) and
above. Thus the joined row masks agree, including
which collar rows are lifts. Their values on the
real fibre come from the actual inherited and
completed physical datum. With that datum and
the same first-member action they give the same
localized critical orbit.

Form the input at the supplied global prospective pair
restricted to the first member of the ordinary join.
The ordinary join is finer than the prospective
sequence: each joined member is contained in the
original member of its level, and the stage enlarges
the latter. An inward bond omitted from the joined
set has an endpoint in the next joined member.
That endpoint lies in the next prospective member,
so the bond is also absent from the prospective
no-inward set. The removed-row condition in the
restriction argument is therefore satisfied.
Extend variations at fixed joined averages by the
unchanged exterior configuration. The original
averages remain fixed, and every interior bond of
a plaquette straddling the true first member is
fixed by a joined bare row. The action difference
of the extension equals the difference of the
restricted action. Thus on the real critical
slice the restriction is critical at its own
averaged datum. The complete boundary blocks use
their original collar levels on the prospective pair.

Use the \hyperlink{hi:strip-canonical-target}{tower transport
of the complex configuration} on these alignment collars.
Let \(D_*\) be the averaged datum of the unaligned
restriction on the ordinary joined rows, and let \(D_o\)
be the target ordinary joined datum in its aligned
frame, formed from the physical collar variables and
their retained rows. With \(g=g[\mathbf U]\), the source
at the unaligned response pair is
\[
 b_o(c)=\frac1i\log\bigl(D_o^{g^{-1}}(c)D_*(c)^{-1}\bigr).
\]
The restricted pair used in the response keeps the
independent current on complete stencils and the
restricted-action variation on its true incomplete
stencils. On the fixed regular
physical row charts the inverse pivot map changes the
prescribed rows by at most \(C_L\|\beta\|_\infty\).
The aligned zero-increment comparison and the bound on
the complex tower transport then give
\(\|b_o\|_\infty\le C_L(\varepsilon_1+\alpha_1+\|\beta\|_\infty)\)
on those charts. Here \(\varepsilon_1\) bounds the real
datum and \(\alpha_1\) the scaled complex potential on
the alignment collars. Where the ordinary row targets
are formed on this same chart, a fixed small physical
source radius and small coupling put this input in the
reduced response disc. The real response with the
endpoint factors identifies the ordinary localized
centre at full parameter values.

In the physical row conversion use the nonlinear
one-row correction at \(V=\mathcal M^j\mathbf U\).
It is holomorphic at complex \(V\) in the admitted
double-block chart and keeps the coarse value
\(\mathcal M(V)\) exactly. Only the rows released by the
stage are replaced by that completed field. The inherited
inward crossing and its ordinary pivot do not add
coordinates to the no-inward ordinary datum. Its other
actual rows and its lifted parents retain their prescribed
reads. If \(r_{\rm piv}\) is the physical radius of the
regular inverse, take a smaller fixed radius inside
\(r_{\rm piv}\) so that the bound for \(b_o\), multiplied
by the constant of reduction, is below half the native
response radius. The terms \(\varepsilon_1+\alpha_1\)
fit the other half of this allowance at weak coupling.
The fixed disc here is for the centre map. The inherited
actions are evaluated on their own running discs.

%
%
\hypertarget{hi:ordinary-strip-joined-section}{}
On the prospective pair take the running bounds on
the full ordinary read and on the original parent
blocks of the artificial rows. Use \(\varepsilon_*\)
for the real plaquette radius in layer units,
\(\alpha_*\) for the complex potential and first
difference, and \(\gamma_*\) for the independent
current. The full tower averaging estimate puts
each parent and endpoint average within \(C_L\alpha_*\)
of its unitary value in the same frame, uniformly
in the tower depth. After real gauge on a full
one-step block, the logarithms of the real reference
are \(O_L(\varepsilon_*)\). The complex change
and a fixed reduction of the physical source
disc therefore keep the pivot and contour maps
inside the one-step analytic charts. In particular
the complete pivot derivative is near its flat
invertible value on these discs. This supplies the
uniform chart room for the fixed-size collar
source. It also bounds the derivative of the
endpoint mean on the output of the response, after
reducing the source disc by the fixed response
constant.
For the complex gauge identities use the one-step
covariance on each complete double block. In
its small reference gauge let \(t\) bound the
link logarithms and let \(\varkappa_{\rm end}\)
bound the condition numbers of the gauge at
the coarse endpoints. The condition is
\[
 \varkappa_{\rm end}\bigl(e^{4(L-1)t}-1\bigr)<\omega_* .
\]
The tower transport bound keeps
\(\varkappa_{\rm end}\) bounded. The running
radii and a fixed reduction of the source disc
put \(t\) below the resulting threshold.
This checks covariance for the actual complex
contours and pivots. For the collared first
member form the alignment and parent averages
on their full original collars before restricting
the gauge to the chosen first member.

Put \(b_*=\varepsilon_{\mathrm D,\mathrm s}
e^{-32\varkappa}\), the reduced source and current
radius of the switched response. Choose the
physical collar radius so that its contribution to
\(\|(b_o)_\natural\|_\infty\) is at most \(b_*/4\).
In addition, take it small enough that the
response estimate times the derivative bound
for the endpoint read fits in half the
fixed contour and product-logarithm chart.
The source from the base is bounded by
\(C_L(\varepsilon_*+\alpha_*)\), including the
conversion. The restricted pair has current radius
at most \(C_L(\gamma_*+\varepsilon_*+\alpha_*)\):
use the independent current on its full stencils
and the variation of the restricted plaquettes on
the true boundary. For the latter estimate the
complex curvature bound and the bare-layer units
give the three powers of the bare spacing in the
current norm. At sufficiently small running
radii the base part and the current thus fit below
\(b_*/4\) and \(b_*/2\), respectively. The joined
gauge condition follows by the local truncation
and restriction of the enlarged gauge cubes.
These are bounds on one fixed-real-factor
section, with independent current. They allow the
whole fixed collar source disc before introducing
the old running argument multipliers.

For the ordinary joined class the pointwise bounds on
the prospective pair pass to the smaller layer spacings.
Use the complete prospective gauge cubes as in the
ordinary strip restriction. If a cube of the join lies
in a coarser prospective layer, cover it by the allowed
larger cube in the two adjacent prospective layers and
restrict its real gauge. This uses the enlarged family
of \cref{rg:structured-local-gauge} on the prospective
carrier.
For the ordinary response at depth \(q=j+1\)
use the join of the strip and the original
sequence on \(Z_q\cap\Lambda_{q-1}\), with parts
of levels \(q,q-1\). At \(q=1\) the join has the
single part \(Z_1\). Use the first member of the
actual ordinary operation. The full first member
and the version trimmed to the collared first
member give different boundary actions.
On a fixed joined section use the independent current
on complete stencils. At an incomplete stencil of the
first member use the first variation of the restricted
action. There are at most six retained plaquettes at a
bond, and such a boundary bond is in the bare layer of
the separated join. Thus its bare current is bounded by
the curvature radius there. This is the same restricted
current prescription as in the ordinary strip comparison.
The joined response uses the torus and native walk
floors at this depth. The ordinary kept strip
quadratic uses the separate conditioned one-step
walks. The full-depth collar form uses the native
full-depth operators.

%
When a comparison changes the first member itself,
the change of its boundary current is local. For
nested first members \(D^-\subset D^+\), let
\(J_{p,\mathbf U}\) denote the contribution of a
single retained plaquette to the first variation in
bare units. At a bond of \(D^-\) the physical currents
differ by the sum of \(J_{p,\mathbf U}\) over plaquettes
incident to that bond and retained only in \(D^+\).
This is an identity for the first variations of the
two actions. It extends holomorphically on the full
plaquette stencils. When transporting the independent
increment as well, subtract this sum from the
current on \(D^+\); the increment relative to the
physical current is preserved. In the bare current
norm a single subtracted plaquette contributes at
most \(C|\partial\mathbf U(p)-\mathbf 1|\). Indeed
the first-variation formula is the traceless part
of the holonomy transported to the bond, with the
bare factor \(\xi^{-3}\) canceled by the current norm.
If the admitted bound for this plaquette is
\(a_\ell\) in the prospective layer-\(\ell\)
curvature units, this contribution is at most
\(C a_\ell L^{-2\ell}\). Sum over the actual omitted
plaquettes, at most six at any bond. The layer-zero
bound suffices for a changed stencil in that layer;
the factor \(L^{-2}\) is available when its bound
comes from layer one. The running curvature ceiling
therefore gives the same small fixed current radius
for either first-member choice. On a fixed first
member, binary switches preserve its true boundary
current and its original plaquette list.

Use the fixed joined sequence in the switches of the
reduced response. The plain datum is formed once by
the row formula above. The cut removes operator terms
on this geometry and keeps the original boundary
current. In an active wall component the response
reads the reduced datum at the first neighborhood
of its response sites. Plain-data conversion reads
through the second neighborhood and the outer
ends of those rows through the third. The four
neighborhood margin contains the complete stencils.
For the walk terms use the closed padded hull
of the trail. It includes the output blocks and
their four-neighborhood with a one-bare-spacing
stencil margin. The switch carrier includes every
closed cell meeting that hull.
Thus a true-boundary current read on a surviving
term uses only the incident retained plaquettes
of the original join inside the active cells.
Only these used current entries enter the trace
comparison.

%
%
For this comparison use an aligned standard switch
grid at depth \(q\), with side a multiple of the
order-\(q\) block and with the native response
floors. The comparison uses the same grid and
native tags in both settings. If the guard is
filed on a coarser nested stage grid, take the parent
image of those tags. The lift is the identity
when the fine endpoints have the same parent.
Otherwise it is the directed link between their
two parents. The fine endpoints lie in the full
blocks of those parents. A released inner parent
that is a pivot lies in the corridor between two
order-\(q\) blocks. These are the blocks containing
the fine endpoints of its lift. The four-neighborhood
margin of a surviving row includes both endpoint
cells, hence those complete blocks by grid nesting.
For a held parent keep its actual prescribed link
and its carrier. The outer parents are level-\(q\)
links and keep their coarse targets. A parent
block can extend outside the joined first member.
Use the original prospective pair on that part
of the active cell. Thus the trace comparison
uses the pair and the row and parent masks on
the full active carrier in addition to the joined
traces on the first member.

%
%
On a generating carrier use the fixed-point
equations on the surviving walk terms. For a
component \(\mathfrak k\) let \(Y^\circ_{\mathfrak k}\)
be the sites over its walk terms. The potential
response is supported over \(Y^\circ_{\mathfrak k}\)
and the shifted multiplier
\(\lambda+\mathsf aB'\) at rows of its first
neighborhood. In the component equations one
can set \(B'\) to zero at the unused rows:
the margin statement gives the same response and
the same supported first variation. Every local
nonlinearity and every weighted difference in
the contraction proof then uses the stencils
in the four-neighborhood margin. The surviving
walk bounds are bounds on these carriers. Thus
the contraction runs on the supplied local
pair and its complete stencils. This comparison
requires the enlarged gauges on the carrier of
each surviving term.

Keep the cells of both endpoint blocks of the crossing
at switch value one. They are the same cell or adjacent
cells, so the read belongs to one active component.
For the perturbed averaged field \(T\), let \(v\) be
the accumulated node factor. At a retained node put
\(u(y)=g(y)v(y)^{-1}\). The order of the factors restores
the aligned datum from the pulled-back datum turned by
\(v\). For an auxiliary axial read at a nonnode \(x\)
in the block of \(y\), set \(u(x)=u(y)T(y,x)\).
This makes the contour of \(T^u\) the identity.
If the read instead has a stored transport
\(T_0(y,x)\) in the aligned inherited frame, use
\[
 u(x)=T_0(y,x)^{-1}u(y)T(y,x).
\]
Indeed \(T^u(y,x)=u(y)T(y,x)u(x)^{-1}=T_0(y,x)\).
When the data have been aligned by \(g\), use the
aligned transport \(T_0(y,x)=g(y)T_{0,\mathrm{old}}(y,x)g(x)^{-1}\).
At a retained exterior fine node use its own
node correction, since it is a top of the old
sequence. The ordinary exterior contour can enter
as an auxiliary link read and is undone on return
to the inherited chart. These
formulas are local on the two complete endpoint blocks.
To compare two cuts, match the restricted pair
on the surviving walk hulls and their four-neighborhood
margin. Match also the ambient parent and pivot reads
there. The latter equality gives the same plain source
on the rows of the conversion and hence the same
reduced source on the rows read by the fixed-point
equations. The true-boundary plaquettes are those
of the original first member on both cuts. Thus the
two component equations have the same coefficients,
source and restricted current. Uniqueness in the
small response ball identifies the resulting read.
At full switches the real response and the same
restoration identify the ordinary centre. The pair,
plain source and endpoint maps above therefore give
the cut rule on the fixed standard partition. On the
common graded cells we use the native walk construction
below, with the same complete row traces.

%
%
%
Distinguish the fields integrated by the stage from the
datum that sets the prospective base. The stage keeps
the prospective average fixed while the collar nonpivots
and their completed rows vary. These motions enter the
response as source coordinates at a fixed pair. On its
admitted section a variation of one physical collar
coordinate \(\beta_c\) changes a bounded set
\(S_c\) of plain rows. This includes a lift if its
named parent uses that released coordinate. Let
\(S_c^+\) include the finite row neighborhood used
by reduction. The regular inverse and the one-step
seam conversion give a reduced increment of size
\(C_L|\Delta\beta_c|\) supported on \(S_c^+\).
For a crossing \(b\), take its two complete endpoint
blocks as the read set \(K_b\). Apply the weighted
response difference bound with weight
\(e^{-\mu d(\,\cdot\,,K_b)}\). It bounds the increment
on \(K_b\) by
\(C_L e^{C\varkappa}e^{-\mu d(K_b,S_c^+)}
 |\Delta\beta_c|\).
The mean contour and relative logarithm are
holomorphic on fixed smaller charts there.
Dividing by the increment gives
\[
 |\partial_{\beta_c}h_b|\le C e^{-\mu d(y_b,y_c)}.
\]
Both endpoint blocks lie in the adjacent-scale geometry.
The endpoint change of units and the number of such
coordinates at a pillar are bounded in terms of \(L\).
The symmetric distance sum gives weighted row and column
bounds at a reduced rate. The switch product for the
two endpoint cells costs at most a fixed factor.
In the scaled logarithm the factor \(g_j\) from
\(\beta_c=g_jB_c\) cancels the output factor \(1/g_j\).
The constants in this derivative bound are uniform
as the prospective pair ranges over the joined
section used for the
\hyperlink{hi:ordinary-strip-joined-section}{ordinary joined response}.
The retained determining values used in that section
are the prospective averages of the same complex pair.
If another entry uses those values in a different
role, verify its admission on its original domain.

At the real prospective datum with the free increment
zero, the actual rows of \(D_o\) and \(D_*^g\) agree.
This is an analytic chart anchor. The sharp and other
retained tests remain in the real integral on their
original arguments.
On a released row, this uses the zero nonpivot increment
and uniqueness of the pivot at the prospective coarse
value. The plain logarithmic difference is then carried
by the original artificial boundary rows. In particular
if the alignment writes \(D_*^g=e^{iB_1}D_o\), the source
at the unaligned pair is
\(-\operatorname{Ad}(g(c_-)^{-1})B_1(c)\).
The finite neighbourhood used by reduction
keeps it in the boundary collar. The crossing reads
lie inside the original quiet window away from this
collar. The weighted response bound therefore gives
exponential smallness of the zero-increment change
there, uniformly in the switches with the endpoint
guard fixed. This uses the prospective frame for the
zero response, where the exterior contour of the
induced datum is the prospective axial contour.
%
%
%
For the complex version let \(W_*\) be the level-\(j\)
average of the supplied prospective configuration before
one-step axialization. The prospective datum on the
real fibre is the stock of retained determining rows,
including the top coarse targets. The prospective
background has precisely these averages. Define the
target on every common row of level \(q\) by the
corresponding average \(\mathcal M^q\mathbf U\).
At a released row use \(W_*\) in the regular pivot
inverse, with coarse target \(\mathcal M(W_*)\).
At an artificial row use the prescribed lift of
the same parent value. Put \(q=j+1\).
At the inner face of the ordinary join the
level-\(j\) parent carriers lie near
\(\partial\Lambda_j\). For \(q\le k_0\)
the prepared box begins five coarse cells
inside the old exterior; for \(q=k_0+1\)
it begins seven. The quiet exterior extends
eight cells beyond the hull of \(G_j\).
Hence the full inner alignment collar and
its parent blocks lie inside \(Z''_q\).
The margin before inserting this band is
at least one coarse cell.
Their lifts use the held level-\(j\) rows.
For \(q>k_0+1\), the prepared box is
inside the old exterior of level \(k_0+1\).
The carried box \(G_j\) includes the
two-cell enlargement of that exterior.
The inner band at \(\partial\Lambda_j\)
is now in the released collar. Its lifts
use the completed level-\(j\) field there.
At depth one there is no inner band.
The parents on the outer face
are level-\(q\) rows and use the fixed coarse
averages. If a free
increment is expressed in the prospective axial frame,
pull it back by the adjoint of the inverse gauge factor
at the start of its bond. The gauge fixes the coarse
centres, so the one-step covariance gives the same
coarse target in both frames on the admitted double
block. The one-row inverse is zero at zero free
increment. This gives the induced rows there in the
same frame on the entire section.

These are all the rows of the joined target. On the
real fibre the common values agree with the actual
retained datum, the released values are the completed
collar datum, and the lifts give the ordinary joined
datum by its defining rule. The complete averaging
blocks of the common rows stay in the join.
A common row meeting a join collar of level \(r\)
has that collar's level. If its endpoint in the
original layer is in the full collar band, that
band is in the original layer \(r\). An endpoint
beyond the band is at distance at least
\(\mu_{\rm col}L^r\) from the collar endpoint,
where \(\mu_{\rm col}\ge L^2\) is the joining
width. The row would then have level at least
\(r+2\). The separation of the original members
would put its collar endpoint in the member at
level \(r+1\), which is disjoint from the collar.
The alignment fixes both endpoints of a row at
level \(r\).
This uses the same geometric test for complex
and real fields.

A held determining row in the analytic target
is read as the prospective average. That value
equals its prescribed group datum on the real
fibre. A stored value can also occur as a spectator
of an old factor. Keep that separate occurrence
in its original slot, including when it is Borel.
Thus the common rows agree at zero collar increment,
even as the prospective pair varies. Only the
artificial boundary rows can supply a zero-increment
source. The finite reduction stencil enlarges that
support once. The current on the original restricted
action stays in the base of the response and is
not a separate input increment.

For a crossing oriented from the old exterior into the
old interior, let \(\tau_-^*,\tau_+^*\) be the two
prospective mean contours of \(W_*\), from the block
centres to the fine endpoints. Use
\(S_{\rm g}=\tau_-^*W_*(b)(\tau_+^*)^{-1}\) as the
prospective analytic link. If the response is zero,
the ordinary read with inverse node factors gives
\(W_*(b)(\tau_+^*)^{-1}\), since the exterior fine
site is an old node. Express the ordinary read in
the common frame by multiplying on the left by the
same \(\tau_-^*\), with the inverse prescription for
the reversed bond. At zero response it is now exactly
\(S_{\rm g}\). On the real prospective axial chart
\(\tau_-^*=1\), so the frame choice leaves the
physical ordinary read unchanged. This continuation
uses only the two complete blocks. It is applied to
the link read and does not change the independent
current in the response. With targets formed from the
same complex pair and the same stored data, the tower
transport, pivot completion, response and mean contours
are functions of that pair. Two splittings of the pair
give the same read by uniqueness of the regular
response on the admitted domain.

With this common-frame prescription and the row
identification, the zero-increment change on complex
sections comes only from the original join boundary.
The uniform decay bound for the switched response gives
\(\|h(0)\|_\infty\le
C(\varepsilon_*+\alpha_*)e^{-cM\check R_q}\)
at the guarded crossings, where the complex bound
uses both the real curvature and complex potential
radii. The old-wall read is at distance of order
\(a_q=M\check R_qL^q\) from either ordinary
alignment band, by the eight-cell inner clearance
and the two-cell outer clearance. At the fixed
depth \(q\) this gives the indicated weighted
distance after including the complete blocks.
Here \(\check R_q\) grows like
\((\log\theta_j)^r\), with \(r>1\), and
\(g_j^{-1}=\theta_j^{1/2}\). After division by
\(g_j\) the error is still smaller than any fixed
inverse power of \(\log\theta_j\). Thus the offset fits the smaller
crossing tier in the Gaussian mean.

%
%
%
%
%
For the substitution in the inherited density, pull the
prospective link back to the ordinary frame. For the
orientation above this gives
\(S_{\rm g}^{\rm o}=(\tau_-^*)^{-1}S_{\rm g}\).
The ordinary link is \(S_{\rm o}^{\rm o}\), before the
multiplication by \(\tau_-^*\). Reverse either link by
inversion for the opposite orientation. In the displayed
orientation the relative logarithm in the ordinary frame
is conjugate to the common-frame logarithm by
\((\tau_-^*)^{-1}\). This
factor depends on the background pair and is fixed as the
independent collar coordinates vary. A restored
interior transport common to both centres cancels
from their relative group element. The new coordinate
\(a_{\rm new}\) is in the ordinary frame. If a coordinate
was first introduced in the common frame, conjugate
it by \((\tau_-^*)^{-1}\) when pulling back the
physical link. This pullback is the identity at
the real axial anchor. The old coordinate is given
by the small logarithm of
\[
 e^{ig_j a_{\rm old}}
  =e^{ig_j a_{\rm new}}S_{\rm g}^{\rm o}
                      (S_{\rm o}^{\rm o})^{-1}.
\]
Here the reconstructed nonpivot values are used in the
inverse pivot map at the original coarse target. At
fixed collar and coarse data the Lebesgue Jacobian
is \(\mathfrak j(g_j a_{\rm new})/\mathfrak j(g_j a_{\rm old})\),
where \(\mathfrak j\) is the Haar density in exponential
coordinates. Keep it with the reconstructed row in
the inherited integral.

The old-wall crossing is in the far part of the plain
strip. At its formation, with \(q=j+1\), the collar of \(\Lambda_q\)
is separated from the old wall by the ordinary
coarse layers. The ordinary rim masks exclude the
strip. The failed-cube cutoff halos have their
bases inset by one coarse cell and stay inside
\(\Omega_q\) even with the two smaller halo layers
and a bond reach. Thus none of the three birth near masks
contains the crossing. The condition uses those
original masks even if a later stage tests the
crossing sharply. Its endpoints belong to adjacent level-\(q\)
blocks, one on either side of the old wall. A level-\(m\)
row can read this fine bond only if its double block
contains both endpoints. The only such row joins the
two block centres. Its inverse pivot uses the full
double block and the prescribed coarse target. On the
quiet geometry both blocks are clear of the old
integration collar and of the cut-off test enlargement.
Their stored bonds are in the ordinary boundary strip,
away from the rim tier. The fine exterior links are
ordinary determining values. Thus the local inverse
pivot and its Jacobian on this double block read far
stored variables and prescribed group values. A direct boundary read of
this coordinate and its occurrence in the conditional
polynomial of the strip use the old far analytic argument.
Evaluate both at \(a_{\rm old}\). The near strip fields
and the real island arguments keep their original
coordinates. The ordinary completion includes this
nonpivot bond with its stored threshold
\(\mathbf1_{\{|g_j a_{\rm old}|<\delta_j\}}\).
The direct sharp mask tests the free exterior
collar nonpivots and leaves this crossing in the
ordinary strip integral. If an enlarged mask also
tests the crossing, its new factor is
\(\mathbf1_{\{|g_j a_{\rm new}|<\delta'_j\}}\).
On its passing branch the triangle inequality for the
real angle gives \(|g_j a_{\rm old}|\le
|g_j a_{\rm new}|+|h|\), with \(h\) in the ordinary
frame. Hence the old factor is one if
\(\|h\|_\infty+\delta'_j<\delta_j\). Use this
inequality on the whole fixed real collar box. The
uniform row derivative of \(h\), integrated along
the segment from \(0\) to \(g_jB\), gives
\(\|h\|_\infty\le C_L\delta'_j+\|h(0)\|_\infty\)
on that box. The ratio \(\delta'_j/\delta_j\)
and the boundary estimate give the inequality at
weak coupling, including on the edge of the box.

%
%
For the simultaneous change, hold the collar
coordinates and the coarse targets fixed. The
ordinary centre depends on the exterior determining
rows and the coarse targets. The prospective
centre has the same fixed targets, and its finer
determining set excludes the old inward crossings.
Thus neither centre uses a crossing nonpivot being
changed. Translate all those nonpivots by their
respective right factors. For each coarse row solve
one pivot with all the nonpivots of its double block
at their new inverse values. Two rows can share
internal links, which are free or prescribed. Each
row uses only its own pivot. This gives the
simultaneous inverse and the product density.

The new ordinary completion tests the stored
coordinate. In the component-local run, on a plain
strip bond this is the original logarithm relative
to its local strip centre. Let \(d_{\rm lg}\) bound
the angle between this centre and the old global
centre on the actual test domain. The global
relative angle is at most \(\delta_j+d_{\rm lg}\)
on a local pass and at least \(\delta_j-d_{\rm lg}\)
on a local failure. Use this reserve when removing
the ordinary approximate bond test: the additional
angle from the global centre to the approximate
centre, together with \(d_{\rm lg}\), must be less
than \(\delta_j\). A failure-energy estimate using
the global angle can use the lower bound
\(\delta_j-d_{\rm lg}\). A factor with the literal
global sharp threshold keeps its own indicator
where it is not implied. In the enlarged mask,
on the new sharp ball its original relative
angle is at most \(\delta'_j+\|h\|_\infty+d_{\rm lg}\).
If this is less than \(\delta_j\), both old
completion indicators are one on that passing
support. The bound uses the ordered product of
the two centre shifts in the same real frame.
Carry the actual local completion indicator on
both branches of the component-local operation.
For the direct mask keep the crossing in the
native strip slot \(a_{\rm old}\) with its ordinary
ball. The sharp balls then apply only to the exterior
collar variables. If a crossing outside the sharp
mask is translated for a chart calculation, use the
inverse translation inside the old indicator. The
new radius gives no bound on that untested coordinate.
Only in the enlarged mask, where the crossing also
passes the new sharp ball, is the old completion
redundant by the displayed inequality. On the
passed ordinary cube mask, the approximate
test uses two different supplies of the \(\delta_j\)
bound on its nonpivot reads. A bond with no
endpoint in \(\Lambda_{j+1}\) has the bound from
its recorded split indicator. A bond with an
endpoint there has it from the cutoff of the
ordinary Gaussian integral. On this mask the
strip change is the ordinary one. The centre
reserve above puts both groups below the
approximate threshold \(2\delta_j\). The second
group is outside the large-field operation.
Thus removal of the approximate factor is an
identity on the support of the whole term,
with those exterior cutoffs present. Earlier
regularity tests are removed where the corresponding
stage comparison gives their value one. This product uses
one choice of sharp centre. A comparison of
both sharp masks also retains their two
disagreement sectors. A version
of the operation retaining any of those tests
still carries its original inverse argument.
On a local sharp failure retain the untransformed
prefix with its local indicator. On a global sharp
failure use the transported real tests and density.

%
%
With \(q=j+1\), the inherited free mask comes
from \(\Omega_q^{(j)*}\). The stage test uses
\((\Omega_q^c\setminus Z''_q)^{(j)*}\cap Z\).
The star uses incident bonds and deletes the
pivots for the corresponding region. The overlap
therefore consists of crossings between \(\Omega_q\)
and \(\Omega_q^c\setminus Z''_q\), subject to the
stage's restriction to \(Z\) and both pivot
exclusions. Keep only the crossings present
in the selected stored strip stock. These are
the additional inherited reads in the broader
sharp mask.

For a quiet inherited crossing in the broader
incident mask, the approximate pass also removes
the ordinary upper cutoff on the sharp-failure
branch. The approximate chord test on that mask
gives angle less than \(2\arcsin(\delta'_j)\) in
the local stage frame. The quiet-crossing estimate
bounds the separation of the local and ordinary
centres by \(C_L\delta'_j\) using the internal
approximate passes. It applies on both outcomes
of the sharp test on this crossing. Thus its
ordinary angle is less than
\(2\arcsin(\delta'_j)+C_L\delta'_j<\delta_j\)
at weak coupling. The plain ordinary upper
indicator is one throughout the approximate-pass
support for this crossing. In an ordinary cube
test this removes the crossing's factor from
the product. Any complement of the whole cube
test still uses the other factors.

For the global failure change of variables,
the raw identity transports the approximate
pass to the chord ball in each new relative
coordinate \(W'_b\) and keeps the sharp
threshold at angle \(\delta'_j\).
Order the finite tested mask and partition
the complement of the sharp pass by its
first failing bond. In each part the earlier
bonds lie in the sharp balls, the failing
bond has angle in
\([\delta'_j,2\arcsin(\delta'_j))\), and the
later bonds lie in the approximate balls.
These are fixed real factors on this mask.
The free \(V_j\) collar stock uses the
determining bonds in the exterior, which have
no endpoint in \(\Omega_q\). The crossing has
one integration slot in the inherited stock
even when it enters two tests. Its cutoff
Jacobian and the solved pivot of its row
remain in the same density. This partition concerns the
approximate-pass part of the global failure
kernel. The other real factors retain their
native arguments.

%
%
%
The density correction here is taken at the actual enlarged
radii. Let \(A\) and \(B_*\) bound the scaled strip and
collar arguments, including their multipliers when they
occur in the substitution. Write \(H_*=C(B_*+d)\), so
\(|a_{\rm old}-a_{\rm new}|\le H_*\). The normalized
Haar density has a nonvanishing analytic extension on
the fixed small chart and its logarithm has zero linear
term. The log of the free-slot quotient is therefore at
most \(C g_j^2(A+H_*)H_*\) in absolute value.
For the pivot factor compare the full coarea densities at
the old and new nonpivot arguments with the same coarse
target, in the original row chart. In Haar frames each
is the inverse determinant of the derivative in the
solved pivot. On a reduced fixed physical disc this
derivative is invertible and its log determinant has
bounded first derivatives. Insert all translated
nonpivots of the double block before evaluating the
inverse. The change of the solved link, and the log
ratio of the row densities, are bounded by
\(C_Lg_jH_*\). The fixed number of fine links of a
double block is included in \(C_L\). Count the row
once. The unchanged density and the ratio give the
original coarea factor at its inverse argument.
If the density is written with the centre determinant
and the linearizing derivative, their product with
the pivot Haar factor represents this same coarea
density. To check the normalization, differentiate
\(\widetilde{\mathsf Q}_V(\psi_V(B))=LQ_VB\)
at fixed nonpivots in the unconstrained pivot
direction, then set \(Q_VB=0\). For one private
row \(c\) at its pivot \(p\), this gives
\[
 (D_p\widetilde{\mathsf Q}_V)(\psi_V(B))
                    D_p\psi_V(B)=L\mathsf m_c(V).
\]
On the fibre the output logarithm is at the
coarse target, so its derivative at the identity
is the identity. The first factor on the left
is the actual row derivative in Haar frames
composed with the pivot exponential derivative.
The latter has determinant \(h(B'_p)\), where
\(B'=\psi_V(B)\). Taking determinants identifies
the coarea factor as
\[
 |\det(L\mathsf m_c(V))|^{-1}
       |\det D_p\psi_V(B)|\,h(B'_p).
\]
The analytic continuation uses the determinant
sign fixed at the real anchor. In the convention
using \(\mathsf m_c\) for elimination, the factor
\(L^{-3}\) is the row's fixed output normalization.
It cancels in a ratio with the same row. The full
coarea factor is evaluated once at the fully
composed inverse argument.

The double block is the fixed root of this ratio.
Its crossing increments read the same two complete
blocks of the row. On a face with this root active,
the restrictions to its active component give the
same nonpivot list, target and pivot solve as the
full face calculation. For a parent read outside
the joined first member, use the completed collar
value for a released link and the original prospective
average for a held one. Include that parent carrier
and the neighboring rows of the naturalization in
the read before restricting the switches. These
parent values are common to the source and the
row chart. Form the response from the completed
prospective collar first. Its restriction to the
face uses the same source values at the rows of
the response. Then use its output in the pivot
solve at the inverse strip arguments, with the
target held fixed. When the crossing gate is off
the two evaluations of the determinant coincide.
The difference of their logs is one holomorphic
rooted read. Apply the rooted Cauchy estimate with
the double block fixed and attach the response
carriers at this root. The one-slot Haar quotients
use their own endpoint roots.

The recentered free crossing is plain. Its group map is
a right translation at fixed collar variables, so it
has unit Haar Jacobian before passing to exponential
coordinates. The cutoff derivative occurs on the
original ordinary transition shell. At a fixed real
shell argument it is the rank-one determinant in the
ordinary shift. Its log is of size
\(C(\varepsilon_j/\delta_j)e^{-c\check R_{j+1}}\)
with the original carrier decay. Its old switched
version holds the shell argument fixed and multiplies
the difference from the zero argument by the old
near-field switch. This costs its switch radius. A
varying smooth cutoff argument in another composed
map requires its own real treatment.

%
On a fixed real part of the global sharp-failure
domain the new cutoff has scale \(\delta'_j\).
Hold its transformed argument \(W'_b\) real and
write the inverse factor, up to the fixed right
translation, as
\(W'_b\exp[-i\vartheta(W'_b)H_b]\).
Its Haar Jacobian is
\[
 J_b=1-d\vartheta_{W'_b}(W'_b iH_b).
\]
In left-trivialized coordinates factor out the
adjoint of the exponential. It has determinant
one and fixes \(iH_b\), leaving a rank-one
perturbation. The derivative of the cutoff is
supported in the transition annulus and has norm
at most \(C/\delta'_j\). Thus the principal
\(\log J_b\) is holomorphic in the shift on the
independent shift disc once \(C|H_b|/\delta'_j<1/2\).
Its absolute bound there is \(2C|H_b|/\delta'_j\).
This continuation leaves the real cutoff argument
and the fixed selectors in the Haar integral.
In the inherited frame the old relative link is
\(W'_b e^{-i\vartheta(W'_b)H_b}S_{g,b}S_{o,b}^{-1}\),
where \(S_{g,b}\) is the prospective global centre
and \(S_{o,b}\) the old strip centre. If a frame
map is needed, transport \(W'_b,H_b\) together
and keep \(\vartheta\) evaluated at the original
real \(W'_b\). For fixed
\(W'_b\), its logarithm is holomorphic in \(H_b\)
and the admitted relative centre logarithm
\(h_b\) of \(S_{g,b}S_{o,b}^{-1}\).
The bound in angular units is
\(C(2\arcsin\delta'_j+|H_b|+|h_b|)\).
Use the independent shift margin for \(H_b\)
and the original strip margin for the other
arguments. The old pivot is then evaluated at
these inverse arguments on its original chart.

If the admitted shift is expanded in carriers
that all contain the same complete bond read,
take its weighted absolute tree sum
\(\|H_b\|_\mu\). Each product in the series
for \(\log J_b\) has a carrier given by the
union of its factors, whose tree length is
at most the sum of their lengths. The weighted
logarithm sum is therefore at most
\(2C\|H_b\|_\mu/\delta'_j\) when
\(C\|H_b\|_\mu/\delta'_j\le1/2\).
At the physical quiet crossing the carrier
of the shift also reaches the distant
boundary, so its spare decay pays this
inverse radius. At fixed prospective centres
the Jacobians multiply on the free slots.
If an ordinary crossing centre reads the
new collar coordinates, form the collar
change first. Its reconstruction does not
read those old crossing coordinates, and
the full derivative is triangular.

%
%
The fixed physical radius for the relative logarithm
gives \(\|a_{\rm old}\|_\infty\le C_L(T+d)\) on the
smaller section above. Here \(T=O(p_0(g_j))\) bounds
the unamplified strip and collar fields. Form the
inverse on this box. Its use in a switched far
argument must fit the old radius
\(\varrho_{\rm st}\mathfrak b_b\), with
\(\mathfrak b_b\) of order \(p_0(g_j)\) in these
units. The direct field and the displacement will
use separate fractions of this radius. The
response of the triple uses its own margin. The
reductions concern only slots integrated in the
stage. Each carried far slot has its full output
radius. The Gaussian operators use their own
cut parameters.

Let \(\Phi(s,\lambda)\) be a switched weight with the
inverse arguments. Its perturbation at fixed switches
is \(\Phi(s,\lambda)-\Phi(s,0)\), where zero sets
all amplitude parameters to zero. This is the
difference to which the rooted estimate with a zero
at vanishing amplitude applies. For one factor per
wall-component take a mixed difference using a
separate amplitude for each component. Cauchy's
formula then gains one inverse radius for each such
amplitude, once the common cut satisfies component
deletion. A component carrying only the fixed join
boundary can affect a difference in a nonlinear
weight even when it carries no collar amplitude.
In the rooted estimate include that boundary among
the source cells and use its own small weight for
such a component. Compare the zero-amplitude value
with the prescribed reference value in a separate
term. When comparing two far arguments at the same
pair and fixed remaining arguments, Cauchy's estimate
along their line bounds the difference by the old
envelope times the displacement divided by the
remaining radius. In particular a residual shift
\(h(0)/g_j\) in this comparison is paid by the
boundary offset, with any multiplier on its slot
included in the displacement.

%
%
%
On the aligned flat stage partition, let \(K_b\)
be the cells containing the two complete endpoint
blocks. It consists of one cell or two face
neighbors. Use the same \(K_b\) wherever \(b\)
is read. Write \(s^{K_b}\) for the parameters with
the switches of \(K_b\) set to one and put
\(\eta_b(s)=\prod_{\Delta\in K_b}s_\Delta\).
An interpolation of the far coordinate is
\[
 a_b(s)=\frac1{ig_j}\log\!\left(
      e^{ig_j a_{\rm new}(b)}
      e^{i\eta_b(s)h_b(s^{K_b},g_j B)}
                              \right).
\]
Here \(h_b\) is the admitted response read in the
ordinary frame, using the original row and parent
traces. At the full evaluation this is the inverse
coordinate. If a cell of \(K_b\) is inactive, the
correction vanishes. When both endpoint cells are
active, \(K_b\) lies in one wall-component, and the
guarded response reads that component. Thus no
correction uses a different component's parameters
or collar increments. The coarse pivot for this
crossing uses the same blocks. At fixed collar
increment the interpolated group map is a right
translation, so the same density quotient applies.
The bound on the product of the at most two switches
costs a fixed factor. The preceding argument for the
far radius absorbs it in the common switch margin.

Every inherited weight uses these common far
coordinates at the crossings it reads. To estimate
the inherited weights, split the field argument
into the original coordinate and its displacement.
For an entry \(F\) with actively integrated far
and near slots \(\mathsf f,\mathsf n\), put
\(\varsigma_b=\lambda_{\mathfrak w}s(\square_b)\)
on its active component and zero otherwise. On
an unchanged far slot let \(a_{\rm new}(b)\)
be its original coordinate and set \(\Delta_b=0\).
On a moved far slot put
\(\Delta_b=a_b(s)-a_{\rm new}(b)\).
Use the joint interpolation
\[
 F\bigl(P_s;
       (\varsigma_b a_{\rm new}(b)+
                          \tau_{\mathfrak w}\Delta_b)_{\mathsf f},
                (\varsigma_b)_{\mathsf n},
                (a_b)_{\mathsf n},\mathfrak q\bigr).
\]
Here \(a_b\) on a near slot is its real value
in the original chart. A moved crossing has a
far slot. When all parameters are one, the far
argument is exactly its value in the original
chart. On a binary cut the displacement uses
its endpoint gate, so is zero when either
endpoint cell is inactive. Its read includes
\(a_{\rm new}(b)\), even at \(\lambda_{\mathfrak w}=0\),
because this coordinate enters the product
logarithm. Choose the owner cell \(\square_b\) from
\(K_b\), so this read is also guarded by the
endpoint switches. The direct argument uses the
owner's switch. The argument
\(\mathfrak q\) contains all the spectators and
their inherited switches, which keep their
values and radii. At zero switch a near entry
is independent of its Borel field. At zero field
it is independent of that switch. Thus the empty
cut gives the original zero-slot value, also for
the near entries. Use the unmultiplied physical
coordinates for the integrated arguments of the
Gaussian reference and its coordinate change.
For the analytic interaction formed from a
difference of strip quadratics use the
corresponding far variables and the switched near
variables \(\mu_b a_b\) on all its spectator legs.
In a rooted read the original
domain contains an endpoint of every crossing it
reads. The extra endpoint cells therefore attach
as face neighbors to those root cells. This is the
attachment estimate on the flat grain, with at most
eight choices of extra cell at each root cell.
The old history trace is kept at the reader's
address. The stage profile below controls admission
of the pair, and the switched chain and double-block
estimates control the native operators with several
crossing slots. A newly composed map that varies an
old near-field argument must supply its own real
coordinate and near switch. The old switch property
applies to spectators in their original slots.

%
%
The displacement part of an inherited term can be
present at \(a_{\rm new}=0\), since the transformed
argument then contains the collar field. The value
at \(\tau=0\) depends on \(a_{\rm new}\) and is
estimated with the ordinary strip amplitudes. Let
\(\varrho_{\rm st}\mathfrak b_b\) be the
old far radius. Even on the smaller enlarged disc
\(|a_{\rm new}(b)|<c\varrho_{\rm st}\mathfrak b_b\),
the bound \(g_j\varrho_{\rm st}\mathfrak b_b=o(1)\)
keeps the product logarithm on its fixed physical disc.
Differentiating it in its second factor gives
\[
 |\Delta_b(s)|
 \le C_L|\eta_b(s)h_b(s^{K_b},g_j B)|/g_j
 \le C'_L\mathfrak b_b .
\]
Here the last inequality uses the collar field on its
stage box and the bound on the boundary offset.
For the finite difference estimate the physical free
strip box has \(|a_{\rm new}(b)|\le C_0\mathfrak b_b\).
Reserve at most a fixed fraction of the old radius
for each of \(C_0e^\varkappa|\lambda_{\mathfrak w}|\)
and \(C'_L|\tau_{\mathfrak w}|\), in units of
\(\mathfrak b_b\). In particular the two amplitude
discs are simultaneously available, after intersecting
the first with the pair-admission disc. For an
estimate with an enlarged \(a_{\rm new}\)-disc, the
unshifted term instead uses its bound
\(e^\varkappa|\lambda_{\mathfrak w}|
\,|a_{\rm new}(b)|/\mathfrak b_b\) in this same
radius budget. The large strip-field disc and the
large direct amplitude cannot both be spent on
the same slot. Spectator slots keep their own discs.

Suppose an inherited term is bounded on the old joint
disc by \(b_X\). For a set of \(k\) components take
the mixed difference in their displacement multipliers
between zero and one. The Taylor coefficients that
survive have positive order in every one of these
multipliers. Summing the Cauchy bounds in those
variables bounds the mixed difference by
\[
 b_X(c'\varrho_{\rm st}-1)^{-k}.
\]
This uses one supremum of the whole term, including
when it reads several bonds of a component. In a
coefficient whose displacement adds an endpoint cell,
that cell must be active because \(\eta_b\) contains
its switch. It attaches to the old support at the
original endpoint. A component sourced only by the
fixed boundary offset uses the offset estimate.

%
%
%
In the switched estimate keep the unchanged strip
arguments and their old real selectors fixed while
expanding in the new switches. Use the pair amplitude
and the displacement amplitude separately. Let \(G\)
be the difference between the evaluations with
displacements \((\tau_{\mathfrak w}\Delta_b)_{\mathsf f}\)
and with zero displacements. Both evaluations use the
same switched pair, the same direct strip argument
and the same near dials. The fixed joined pair is
part of the base. Any zero-increment displacement
from the join enters \(G\) with the same multiplier
\(\tau_{\mathfrak w}\) as the rest of \(\Delta_b\).

Keep the original rooted read \(N\) fixed. For new
switch support \(K\), use the finite difference
\[
 \sum_{J\subset K}(-1)^{|K|-|J|}G(1_{N\cup J}).
\]
The amplitude of a wall component is inherited by
the pieces obtained on cutting it, with the original
anchors retained. The finite difference is the
integral of a mixed first derivative in the switches
of \(K\), over their unit intervals. Estimate this
derivative on the spatial switch discs while keeping
the two families of amplitudes on their joint discs.
If both amplitudes of a new wall component vanish,
the binary cut can delete its new cells. The pair
increment is zero there and so are the direct
amplified arguments and the guarded displacements
of that component. Pairing the subsets with and
without a new cell of this component makes the
finite difference zero. The fixed read is kept in
every summand. The difference also vanishes when all
the displacement amplitudes are zero. Consequently
its Taylor series in the amplitudes has at least
one positive degree per new wall component and a
positive displacement degree somewhere. The unchanged
strip dependence is summed by the original strip
block expansion.

Let \(a=(R_{\rm pair}-1)^{-1}\) and
\(u=(R_\tau-1)^{-1}\), where the radii refer to the
two amplitude discs. Let \(z\) be the ratio from
Cauchy's estimate in the spatial switches, with the
exponential margin of their original discs. For \(k\)
wall components the sum of the absolute amplitude
coefficients of the finite difference at physical
amplitudes is bounded by
\[
 b_X z^{|K|}\,k u(1+a)
                  [a+u(1+a)]^{k-1}.
\]
For a component marked with a displacement the degree
sum is \(u(1+a)\). At every other component sum the
degrees with at least one of the two amplitudes present.
Mark one of the \(k\) displacement components for this
bound. It has an original crossing anchor in \(X\),
since the old reader contains a crossing endpoint and
its displacement vanishes unless the endpoint gate is
active. An endpoint added by the gate is one face step
from that anchor. The other components attach as arms.
Their connected sets, weighted by the spatial switch
factor, have a sum bounded by \(C e^{-\nu D(a_0,X)}\)
at a source anchor \(a_0\). Here \(D(a_0,X)\) is the
length of a shortest attaching cell chain. The
attachment length is included in the decay weight of
the new domain. The old packing bound on the sum over
source anchors then gives a factor at most
\[
 \exp\!\left(CN^*(a+u)\left[1+
            \frac{d^\natural(X)+1}{\mathsf N_i}\right]\right)
\]
for the unmarked components. The amplified pair
radius here uses the stage gas choice
\(q-\mathfrak p_0-5r>9r\). Together with the old strip
radius this makes \(N^*(a+u)\) small. Charge the
displayed factor and the distance estimate below
by reserving a further fixed multiple of
\(1/\mathsf N_i\) in the stage rate. Over the full
waiting window this is a fixed loss.

For a pin in the new domain, let \(C'\) be the
component of the marked original anchor. The old
anchored tree and its attached arms place the pin at
distance at most \(CM(\mathsf N_i+d^\natural(X)
+\sum_W|W|)\) from \(C'\), in stage units. The
same bound holds for the distance from \(C'\) to a
component \(C\) met by the new domain. Reserve
\(c/\mathsf N_i\) of the rate and part of the
spatial switch margin to extract the corresponding
exponential in the distance divided by \(M\mathsf N_i\).
At fixed \(C'\), sum the marked original anchors
using \(J_{\mathsf A}^{\Sigma}(C')\). Earlier stage
terms give the additional \(C N^*\mathsf N_i\)
from their old pinned bound. Lower-level strip reads
are counted at their original endpoints, where their
long-domain estimate still applies. Summing the
distance weight over \(C'\) by the packing of the
treated components gives, for both the new pinned
norm and the new sum over domains meeting \(C^+\),
the bound
\[
 C u\sup_{C'}\bigl(J_{\mathsf A}^{\Sigma}(C')
                                  + C N^*\mathsf N_i\bigr).
\]
This uses the old summed-anchor norm for the marked
component even when the new pin lies in the original
domain. Replacing that sum by a count of possible
anchors in a domain would cost a further waiting-depth
factor.

To enter the difference in the primary potentials,
subtract its value at \(B=0\) in the actual collar
integration variables, at fixed pair and unchanged
strip arguments. Every surviving block difference
then uses at least one new collar block, and the old
collar ratio \(\hat q\) applies. The factor \(\hat q N^*\)
in the component bound is small, so this insertion
fits the old \(\mathsf N_i u N^*\mathsf N_i\) bound
after summing the stages. At zero free source the
subtracted offset has the exponentially small bound
from the joined comparison. If the independent free
shift \(\mathsf y\) moves the physical zero, its
extra displacement on the free-source disc is at
most \(C|\mathsf y|/g_j\), by the same fixed physical
chart estimate. This part uses the free-source closure
bound, with the small ratio
\(|\mathsf y|/(g_j\mathfrak b_b)\).

%
%
%
Count hot blocks among the current integration
variables. At stage \(j+1\) these are the level-\(j\)
strip and collar nonpivots. Assign each coordinate
block to the level-\((j+1)\) \(M\)-cube of one of its
endpoints in the stage collar. If both endpoints are
available, use a fixed convention. In the direct
graded rule take the new sharp product on the free
exterior mask and keep the inherited crossings in
their original strip coordinates. For the broader
re-centred rule use the product on the passing branch
above, where the smaller crossing ball makes the
transported old test redundant. The two disagreement
sectors of a comparison of sharp masks keep their real
indicators and are treated separately.
An earlier stored spectator or a solved pivot adds
no block. Each owner cube has at most \(C(LM)^4\)
slots. The quiet box at level \(j+1\) contains at
most \(C\check R_{j+1}^4\le C N^*\) owner cubes in
its collar.

%
At fixed shift source the smooth Haar difference
at a free collar bond reads just its own variable
block. The centre term reads none. Thus the
alternating block expansion of the difference
has the usual single active-block factor
\(\hat q_v\) and at most \(e^{2a_1}\) from the
hot-block charge. The shifts formed from fixed
stage data add their rooted carriers to the
localization and no current integration block.
This potential reads the same shift source
as the translated field. Add its envelope to
the primary potentials while keeping the one
closing index for that source and carrier.
The bounded number of free slots of a cell
then adds a bounded amount to the hot-block
activity, using the smooth Haar bound above.
In the stage tail estimate the weight of order
\((\mathsf N^*)^{-1}\) at the sharp threshold
gives the factor
\(\exp[-c(\log\theta_t)^{2\mathfrak p_1-4r}]\),
which pays the existing activity of order
\((\log\theta_t)^{4r}\) under the exponent
choices of the step. A failed sharp branch
uses the separate cutoff Jacobian in its real
history kernel. On the quiet branch the product
over at most \(C\mathsf N^*\) transformed slots
has log bounded by the sum of their shift bounds
divided by \(\delta'_j\). The quiet collar decay
pays this sum as well as the factor \(\mathsf N^*\).
This costs only a bounded density multiplier.
A space-time weight on that history uses
the survival reserve.

Use this ownership for the gas and keep the graded
tree for the term's carrier. In the flat last stage
the two cell systems coincide. If a graded carrier
cell meets new hot slots from a treated component,
their number is at most the total \(C N^*\) for
that component. The separation of the treated boxes
puts the slots of at most one treated component at
such a cell, including the fixed endpoint
enlargement. Hence the new hot-block factor on the
added cells is bounded by
\(\exp(C a_1 N^*\sum_W |W|)\).
With \(a_1=\hat q\) it fits the spatial switch
margin, also in the earlier stages. The bound on the
ratio of adjacent grades gives bounded degree for
the graded cell chains. Its volume conversion
\(L^{4(\mu-j-1)}\) is used for an old coarse pin,
where the path to the anchor pays for it. Count
only the assigned integration blocks in the gas.
The Gaussian cost below uses the same count, with
the error from each guard weighted by the number
of slots it protects.

%
%
%
%
There is a compatible grading for the switched
walks of the stage response. On the part below
the next layer use the level-\((j+1)\) \(M\)-cubes.
Above it, on the prospective layer \(\mu>j+1\),
use the level-\(\mu\) \(M\)-cubes. The stage
regions are unions of these nested cells and
consecutive faces are separated by a full cell
of the larger level. Small relative
enlargements of the cells meet at most adjacent
grades. They have bounded overlap, by separation
of nonadjacent layers and by the cube count
within each grade.

For a walk use the union of the closed hulls
of its constituent cells. Let \(G\) be a
graded cell of side \(M\xi_\mu\) that meets
one such hull. The hull's diameter is at most
\(C_L\mathfrak M\xi_{n_*}\), with \(n_*\)
its largest constituent level. If
\(n_*>\mu+1\), its level-\(n_*\) constituent
would lie in the region at that level at
distance at least \(M\xi_{n_*}\) from \(G\).
For sufficiently large \(M\) this is
impossible. Thus the tree that visits the
hull meets a small relative enlargement of
\(G\). Unless that tree lies wholly in the
enlargement with twice the width, it spends
weighted length at least \(c_L M\) there:
before exiting it uses spacings at most
\(\xi_{\mu+1}\). Sum over these enlargements
by bounded overlap. The cells whose larger
enlargements contain the whole tree are
bounded in number by the same overlap.
The initial hull, when it is counted
separately from the visited hulls, meets only
a bounded number of cells by its diameter
and the same level comparison. Hence, for
the walk length \(|\Gamma|_{\mathbf B}\),
the number of distinct graded switches is
at most
\[
 C_L+C_L|\Gamma|_{\mathbf B}/M.
\]
The switch loss therefore fits in the
native walk decay after taking \(M\)
large compared with the switch radius.
Apply the weighted walk summation to the
individual scalar monomials and their partial
sums. The cost \(e^{C_L\varkappa}\) is fixed
and the rate loss is \(C_L\varkappa/M\).
In particular, if the native absolute sum is
bounded with weight \(e^{\delta |\Gamma|_{\mathbf B}}\),
the sum after switching at rate \(\delta'\) is
bounded by the same constant times
\(e^{C_L\varkappa}\) for
\(\delta'+C_L\varkappa/M\le\delta\). This follows
termwise from the distinct-tag count. Replacing
old tags by their parent set can only decrease
the number of factors on this common polydisc.
The prefactor \(e^{C_L\varkappa}\) changes
the fixed source and current ceilings
in the contraction. Reduce these ceilings
before choosing the small-coupling range.

%
%
%
Use an auxiliary standard partition at the top of
the native tower to check the torus floors for the
unswitched walk setting. On the actual graded grid
attach to each native term the product over all
closed graded cells meeting its padded hull. For a
test of a native kernel, apply the graded count to
each tree joining the two test sites and visiting
the trail, including the initial hull as above.
If the native test has rate \(\vartheta\), this gives
\[
 |s^\omega|e^{-\vartheta d_{\mathbf B}(\omega;y,y')}
 \le e^{C_L\varkappa}
        e^{-(\vartheta-C_L\varkappa/M)d_{\mathbf B}(y,y')}.
\]
Here the distance in the exponent on the left is the
infimum over the test trees. The native sum over terms
starting in the test window gives the same estimate
for every partial sum with the native window and mass
constants. Use this for \(\mathsf H_0,\tilde{\mathsf G},
\mathsf H,\hat{\mathsf S}\) and, when used for the
potential read, the walk datum of \(D(1-\mathsf R)D^{\mathsf T}\).
Take \(M\) large so that the rate loss fits inside
the native spare decay in each test. The layer-ratio
estimate and the weighted count of bonds over each
site then give the same operator types for the four
maps of the reduced system, including the transpose of
\(\mathsf H\), with a larger constant. The contraction
and its weighted source estimate use these operator
types and the unchanged local structure. This gives
the switched response on the smaller fixed source and
current discs, with holomorphy by uniform convergence
of the fixed-point iteration.

For a face of the switch polydisc, keep the
switches in the allowed cells complex and put
the others at zero. Define wall adjacency by
a common face of positive three-dimensional measure,
including a subface at a change of grade. At any
point of the padded hull the closed cells containing
that point form a wall-connected set. Indeed the
small coordinate orthants at the point belong to
these cells, and consecutive orthants across a
coordinate plane give either the same cell or a
common wall. The union of the cells containing the
point contains a neighborhood of the point. Since
the hull is connected, its full tag set is wall
connected: different wall components of the tag
set would partition the hull into disjoint closed
sets, since all tags containing any given point
belong to one component.

A term survives on the face only when all its
tags lie in a single allowed wall component.
The native margin puts the four-neighborhood of
each kernel site inside the padded hull. Thus, if
\(\mathbb Y_{\mathfrak k}\) is the set of sites
whose four-neighborhood cubes lie in the interior
of the union of \(\mathfrak k\), the surviving
terms and their transposes are carried by
\(\mathbb Y_{\mathfrak k}\). The interiors for
distinct components are disjoint. In particular
the three-neighborhood of \(\mathbb Y_{\mathfrak k}\)
misses every other carrier set. Apply the splitting
of the reduced system over separated carriers.
The potential and the supported current bracket split.
In the multiplier the diagonal term \(-\mathsf a B'\)
is kept separately, as in the splitting identity for
the bracket. Use the ambient row types and true-boundary
stencils throughout. In the native exterior collar the
two parent blocks of a completed level-\(j\) row are
in the old layer \(j\), at distance \(O(L\xi_j)\)
from the row. Include both parent blocks and the
pivot in the source read. A lifted row with an
additional parent outside the raw margin uses the
connected carrier of that parent stencil as well.
Gate the corresponding source contribution by the
tags of its complete stencil. The extra switches
cost a fixed factor in the source bound, and at
full switches the source is the original completed
source. At full switches the native
terms sum to the original four operators and the
full response is identified by uniqueness.

%
%
%
To compare two placements of a cut on the same
torus, use a common union \(R\) of the graded cells
in native settings with the same scales and walk
parameters. Take the layer traces \(\Omega_n\cap R\)
on the common finest native grid. Choose the graded
cell widths on the nested grids of the native setting.
Each graded cell, each layer member and each padded
native hull is then a union on the finest grid.
Equality of the traces on \(R\) therefore gives
equality on each padded hull contained there.
For a surviving atom the identity on its hull
is then a native shift. Locality in the sequence
identifies its index and coefficient at matching
pair data on that hull. Repeat the comparison
from the other setting, so the sums have exactly
the same indices. The common graded grid gives
the same tags to the matched hulls. For a marked
conditional walk include its conditioning set in
the trace and use the same compression mask on the
matched endpoints.

For the reduced equations the same comparison
identifies the tops, levels, row types and outward
endpoints throughout the padded read. A site or a
row is assigned by the layer memberships of its
sites, with both ends covered by the margin. Thus
the same actual averaging blocks, pivot reads and
retained plaquettes enter the local structure.
Compare the pair and the independent current on
these stencils, including the true first-member
plaquette lists. The reduced source agrees when
the plain source agrees on the full stencils
of its naturalization and completion. These are
the data needed to apply uniqueness to the
component equations on \(R\). The comparison
is an identity with the switches in \(R\) still
complex and those outside \(R\) set to zero.

%
%
%
In a preparatory integration the changed member at
level \(q=j+1\) removes the released fine rows
from the level-\(j\) determining set and adds the
coarse target rows. A retained coarse row on the
old face can change from outward to interior.
Both versions of such a row may carry the same
grade-\(q\) tag. The outward equation reads the
seam factor of its full outer block, even at zero
plain increment on that row. The interior equation
uses its plain increment. At the changed seam form
each response in its own native setting. In the
old problem the new coarse value is the fixed
target of the pivot equation, and the completed
fine rows supply the old plain source. In the
prospective problem use its actual new row set.
In the ordinary centre construction the restriction
to the fixed join uses the induced ordinary datum
and its own action. It follows from the critical
restriction on that join, together with the
endpoint restoration above. The ordinary depth-\(q\)
operators and the full-depth collar operators
retain their own native settings.

Along a fixed sequence of stages the
regions grow and the base level increases.
An old graded cell lies in a current
cell because the current grade there is
at least its old grade and all grids are
nested. For an inherited walk take the
parent image of its actual tag set, using
its original boundary ownership. On a
binary cut the old tags are all active
exactly when all their parents are active.
The full-switch coefficient is unchanged.
For a walk formed with a closed-hull tag
rule the same image is the set of current
cells meeting that hull, since each is
the union of its old cells. Apply this
identity with the original hull. Additional
endpoint or response stencils are filed as
new attachments. The new monomial uses
each graded cell once. The per-slot switches
of the carried fields keep their own
arguments. For the ordinary join, which
ends at level \(j+1\), this retagging starts
from its admissible standard partition.
The full-depth collar Hessian and its roots
use their native walk terms on the same
active cells.

For the shear \(\mathsf T=I+\mathsf S\),
\(\mathsf S(a,B)=(\mathsf LB,0)\), one has
\(\mathsf S^2=0\) and \(\mathsf T^{-1}=I-\mathsf S\).
On a full component preserved by this shear the Gaussian
determinant is unchanged, and the covariance is the old
covariance multiplied by \(I-\mathsf S\) and its
transpose. Weighted row and column bounds for
\(\mathsf L\) and the old covariance give the corresponding
decay for this product by the triangle inequality for
the distance. A compression to a set preserved by the
shear is a congruence of the compressed precision.
Use \(\mathsf Tz+d\) as the affine reference of the
inverse map, with \(z\) the new coordinates. Here
\(\mathsf L\) is the derivative in \(B\) at zero
and \(d\) is the value of the inverse at zero.
Write \(r(z)\) for the remaining part. It includes the
term linear in the new strip coordinate from its
product with the offset, bounded by
\(C g_j|a_{\rm new}(b)||d_b|\). If the old
quadratic is \(Q(v)=\frac12v^{\mathsf T}Dv-\ell^{\mathsf T}v+c\),
then its new precision and linear coefficient are
\[
 D'=\mathsf T^{\mathsf T}D\mathsf T,\qquad
 \ell'=\mathsf T^{\mathsf T}(\ell-Dd).
\]
The new constant is \(c+\frac12d^{\mathsf T}Dd-\ell^{\mathsf T}d\).
The rest of the change in \(Q\) is
\[
 r(z)^{\mathsf T}[D(\mathsf Tz+d)-\ell]
                    +\tfrac12r(z)^{\mathsf T}D r(z).
\]
Include the sum of the one-slot logarithmic Haar quotients
and one inverse-pivot Jacobian per changed coarse row.

Suppose the derivative \(\mathsf L\) has a walk
expansion on the same grain as the old precision, with
both endpoints in each term's carrier. In
\(D'=D+\mathsf S^{\mathsf T}D+D\mathsf S+
\mathsf S^{\mathsf T}D\mathsf S\), join successive
carriers at their common index. Their union is connected,
its length is bounded by the sum of the lengths, and
the weighted masses multiply. The same joining applies
to the quadratic remainder. For \(r^{\mathsf T}Dr\)
the two remainder carriers attach at the two ends of
the same old precision term. The weighted row sum
therefore bounds the mass per crossing. Use the actual
products of switches in both expressions. A shared cell
can appear more than once in such a product.

For the derivative carrier expansion, fix the complete
endpoint carrier of the crossing and expand
\(\mathsf L_b(s;c)\) in the other switches. To read a
coefficient with support \(Y\), set to zero every
switch outside \(Y\) and the endpoint carrier. The face
identity then makes the resulting function independent
of any switch separated from the endpoint carrier.
It is also independent of the source coordinate if
its complete stencil is outside that component. This
gives the required vanishing of the corresponding
coefficients. On a
uniform switch disc of radius \(R_s\), a bound
\(C e^{-\mu d(b,c)}\) on the derivative gives, by Cauchy,
the absolute sum at \(|s|\le r_s<R_s\) over monomials
with support exactly \(Y\) bounded by
\[
 C e^{-\mu d(b,c)}
       \left(\frac{r_s}{R_s-r_s}\right)^{|Y|}.
\]
Indeed each cube of \(Y\) has a positive degree, and
summing those degrees gives the displayed geometric
factor. Choose the raw response disc larger than the
disc for the expansion, and use its corresponding
source radius. The degree of the graded wall graph is
bounded in terms of \(L\), since neighboring grades
differ by at most one. A connected set attached to the
endpoint carrier is visited by a walk around a spanning
tree, of length at most twice its size plus a fixed
endpoint term. The resulting exponential count is paid
by the switch margin. Choose the margin before \(M\).
The same argument applies
to \(r_b\), since subtracting the affine reference preserves
the cut dependence and the bound of order \(g_j(T+d)^2\)
on the larger section. For \(d_b\) the terms that remain
connect the two-block read to the join boundary. Use
this same grain for the wider guards.

%
For the restricted inverse on \(W\), take an intermediate
set \(W_1\) of variables containing \(W\) with a buffer.
Use a guard beyond \(W_1\) on which the switches are
one, and the old real reference pair. Set the exterior
switches to zero for the reference operators. Write
\(P_W,P_1\) for the projections onto \(W,W_1\).
The old local-positivity bound supplies a floor on
the compression \(D_{1,0}\) to \(W_1\). Its congruence
by \(I+P_1\mathsf S_0P_1\) has a uniform floor,
since its off-diagonal part is square-zero.
To compare its \(W\)-compression with
\(A_{W,0}=P_W\mathsf T_0^{\mathsf T}D_0\mathsf T_0P_W\)
on \(W\), use
\[
 \mathsf T_0P_W=P_1\mathsf T_0P_W+(I-P_1)\mathsf S_0P_W.
\]
The difference of the two compressions contains the
last term at least once. Its norm is exponentially
small in the buffer by the weighted bounds. This
gives the real floor on \(W\) for a sufficient buffer,
also when \(W\) cuts across the shear.

At the same reference pair, a term in the complex
cut perturbation must reach from \(W\) outside the
guard. The joined carriers, including those of
\(\mathsf S\), bound this perturbation by
\(C e^{-\nu d_{\rm cell}(W,G^c)}\) in the weighted
row and column norm, with \(G\) the guard. The full
perturbation \(E_W\) adds the change from the admitted
real pair to the actual pair on this guard. At fixed
switches the change of \(\mathsf T^{\mathsf T}D\mathsf T\)
is bounded by \(C(\|\Delta D\|_\nu+\|\Delta\mathsf S\|_\nu)\),
where \(\|\cdot\|_\nu\) denotes the maximum of the
weighted row and column sums. Expand the difference
in the three factors and use the uniform norms on
each section. The real
floor and the exponential kernel bound give the
weighted inverse of \(A_{W,0}\) at a reduced rate by
the local-resolvent estimate. Use
\[
 (A_{W,0}+E_W)^{-1}
       =\sum_{q\ge0}(-A_{W,0}^{-1}E_W)^q A_{W,0}^{-1}
\]
when the weighted norm of the product is less than one.

%
%
%
Use the old real Gaussian reference, namely the pair
\((U,0)\) with the same real factor \(U\). On the
joined response carrier use the same true-boundary
prescription there, so at an incomplete stencil its
current is the restricted variation at \(U\). Compare
the shear on the section with coordinates
\(q=(A',\mathbf J)\), with norm \(\|q\|_*\) the
supremum of the layer-scaled potential, its first
covariant difference and the independent current.
The gauges on the complete joined cubes depend only
on \(U\). Fix their ceiling and a fixed upper bound
for the real curvature with room in the raw response
class. On a small fixed ball \(\|q\|_*<r_*\), the
plaquette perturbation estimate bounds the change
of scaled curvature by \(C_L(r_*+r_*^2)\), using
adjacent-layer units on the full stencils. The
independent current and the restricted variation on
the true boundary fit the raw current radius by the
same choice of ceilings. Reduce \(r_*\) to leave room
also for the source conversion and the endpoint
logarithm on the fixed physical collar disc.

On this section, for \(\|q\|_*<r_*/2\), apply Cauchy's
estimate to the line \((U,zA',z\mathbf J)\) at fixed
switches. The uniform weighted bound for the shear
coefficients on \(|z|<r_*/\|q\|_*\) gives
\[
 |\mathsf L_{bc}(U,A',\mathbf J)
             -\mathsf L_{bc}(U,0,0)|
 \le C\,\frac{\|q\|_*}{r_*}e^{-\mu d(y_b,y_c)}.
\]
Summing at a smaller weight gives the required
\(\Delta\mathsf S\) bound. The old Gaussian comparison
supplies \(\Delta D\) on its own domain. The running
radii fit inside this fixed response ball at weak
coupling. The targets along this line use the averages,
completed rows and parent lifts above. At \(z=0\)
these give real targets in the same fixed frame.
The old spectator occurrences keep their values.

Transport the old admitted decoupled covariance factor
\(R_0(\hat s)\) by \(R_1(\hat s)=(I-\mathsf S(\hat s))R_0(\hat s)\),
using the same guarded switches as in the precision.
For the compression to \(W\) write \(A_W\) for the
transformed precision on \(W\) and \(B_W\) for its
block from the complement into \(W\). Take
\(C_W=A_W^{-1}\) and \(\Gamma_W=B_W R_1\), extending
\(B_W\) by zero on \(W\). The product has the same
carrier joining and guard comparison. At full
evaluation \(R_1R_1^{\mathsf T}=(D')^{-1}\), and
block inversion gives its marginal on \(W\) as
\[
 C_W+C_W\Gamma_W\Gamma_W^{\mathsf T}C_W .
\]
The full normalizer uses \(\det D'=\det D\).

%
For the cut identities choose the guard as a
union of fixed guards \(G_b\) of individual
active slots. The connected cell set \(G_b\)
contains every cell meeting the closed owner
cell enlarged by a fixed fraction of its side
divided by \(L\). Add the prescribed endpoint
and parent cells. For two cells with touching
closures the grades differ by at most one.
Subdivision on the finer of those grids
shows that only \(O_L(1)\) cells meet this
enlargement. With \(M\) larger than the fixed
resolvent patch width, it contains the local
buffers of every cutoff that is nonzero at \(b\).
All masks come from the ambient stage. For
a set \(W\) of integrated slots put
\(G(W)=\bigcup_{b\in W}G_b\).
If \(Z\) contains \(G(W)\), every \(G_b\)
lies in the wall-component of \(Z\) containing
its owner. Thus, for each such component
\(Z_\alpha\), with \(W_\alpha\) the slots of
\(W\) owned there,
\[
 G(W_\alpha)=G(W)\cap Z_\alpha .
\]
Use the same ambient lattice partition and
the original slot mask in each component.
Every surviving precision and root atom
retains its full parent and patch reads.
Its indices therefore agree with the indices
in the component restriction. In particular
\(A_W\) splits into the blocks \(A_{W_\alpha}\)
at a binary cut, and its inverse is taken
block by block. The same split for
\(B_WR_1\) uses independent noise slots
in the distinct components.

With \(F=\Gamma_W\), use the normalized
joint Gaussian in the active variable \(x\)
and standard auxiliary noise \(y\) with
quadratic exponent
\[
 -\tfrac12(x-m-C_WFy)^{\mathsf T}
                 A_W(x-m-C_WFy)-\tfrac12 y^{\mathsf T}y .
\]
With \(n_x,n_y\) the real dimensions of
the two vectors, the density factor is
\[
 (2\pi)^{-(n_x+n_y)/2}
          \exp(\tfrac12\operatorname{Tr}\log A_W).
\]
Both factors multiply over the components.
Unused noise integrates to one.
Here \(\log A_W\) is the matrix logarithm of
the actual switched compression, taken on the
branch fixed by its positive real reference.
The guarded perturbation is smaller in norm
than the real floor, so \(A_W\) stays strictly
accretive on the cut polydisc. Its logarithm
there is holomorphic and agrees with the
blockwise logarithm on a binary cut. This
normalization uses the inverse \(A_W^{-1}\)
and its own determinant. The independently
switched full root supplies \(F=B_WR_1\).
Away from full switches the definition only
uses that factor and its carrier bounds.
Integrate out \(y\) to obtain the
normalized cut measure \(\mathcal J(W;s)\).
At full switches it is the actual Gaussian
marginal, by the identity for \(R_1\).

For a cell set \(\mathcal K\) outside
\(G(W)\), take the alternating sum
\[
 \nu(W,\mathcal K)=
  \sum_{\sigma\subset\mathcal K}
  (-1)^{|\mathcal K\setminus\sigma|}
          \mathcal J(W;\mathbf1_{G(W)\cup\sigma}).
\]
Summing over \(\mathcal K\) gives the
full marginal by inclusion and exclusion.
For fixed \(\mathcal K\), the preceding
guard equality factors the sum over
wall-components of \(G(W)\cup\mathcal K\).
These are the two cut identities used
by the gas. The determinant in
\(\mathcal J\) is a conditional normalization.
The direct normalizer of the full stage
is filed separately. The root identity
with the inverse precision is used at
full switches.

%
%
%
For the Gaussian cost on a graded cell \(\Delta\)
use its physical side \(\ell_\Delta\) and put
\(q=j+1\). The ordinary kept form and the centre
response of the strip use towers through level
\(q\). The new collar form uses the full
prospective stage sequence. Its walks start at
active slots in \(Z_q\), with a fixed endpoint
enlargement. Inside a quiet component the old
\(\Omega_{q+1}\) is separated from \(Z_q\) by
at least \(c a_q\), where \(a_q=M\check R_qL^q\).
This member is unchanged by the present stage.
For an earlier stage the active slots are also
far inside \(X\), clear of the top joining collar.
At the last stage the top level itself is \(q\).
Thus, until a walk has traveled a distance
\(c a_q\) from an active slot in an earlier
stage, it cannot use a layer with spacing
greater than \(L^q\).
If it reaches another treated box, use
the same estimate up to the first exit
from the inset of its own box.

Put the guard at one on the cells touching the
cells of the active set and include the fixed
carrier reach of the chosen switched root.
Use the local collar root below on the full
stage tower and form the kept-strip root on
its ordinary level-\(q\) grid. Take the parent
image of the strip's fixed guard in the common
cells.
Adjacent grades differ
by at most one, so around each active cell this
guard covers a physical neighborhood of width
at least \(c\ell_\Delta/L\). Reserve the fixed
one-step carrier reach inside this neighborhood.
For \(q\) below the last stage the next
original member stays unchanged on \(X\).
The active slots are in the carried
\(G_q\), which has the quiet gap of at
least \(8a_{q+1}\) to that member. Since
\(a_{q+1}\ge a_q\), this gives the upper
clearance through the compressed stages.
At the last stage there is no higher member.
On the lower side the prepared separation
gives room \(ML^q\) for the local buffers.
For a patch crossing the exterior of the
treated box use the separation of the
prospective member at level \(q\) from
the complement of the member at level \(q-1\).

Use the padded hulls in the escape estimate.
A native atom whose constituent cells are
all of levels at most \(q\) stays within
\(C_L\mathfrak M L^q\) of the tree point
that meets it. If it has a constituent of
higher level, follow its connected cell
family to the first such constituent.
The bounded family of earlier cells puts
the tree point within \(C_L\mathfrak M L^q\)
of the higher member. The initial cell
is anchored to the active index and has
the same carrier bound. After allowing
these fixed reaches, an off-guard term
must already have paid weighted length
\[
 d_\Delta=c_L L^{-q}\min\{\ell_\Delta,a_q\}.
\]
Use the native weighted lengths for the absolute
walk sums. The displayed bound extracts an
exponential for a term which leaves the guard,
with room left for the row and column estimates.
The number \(n_\Delta\)
of actual active slots is at most
\(C_L L^{-4q}\min\{\ell_\Delta,a_q\}^4\):
use the volume of the cell and, separately,
the size of the level-\(q\) quiet box.
In the native stage grid the active slots
and their local guards are at base grade
\(q\). If a later stage retags the formed
term, retain the estimate at the time of
its formation. For each cell with active
slots, \(u=L^{-q}\min\{\ell_\Delta,a_q\}\ge M\).
Consequently \(n_\Delta e^{-c d_\Delta}\le
C_Lu^4e^{-c'u}\) is uniformly small once
\(M\) is large.

%
%
For the shear, compare the response at \(s\) with
the response at \(s_G\), obtained by setting the
switches outside the guard to zero. Keep the same
ambient pair and stencils. Fix an active anchor
\(b\) in \(\Delta\) and use the weight
\(w_b(y)=e^{-\eta d_{\mathbf B}(y_b,y)}\), with
\(0<\eta\le\theta_0\) in the contraction estimate
and \(\eta\) small compared with the physical
source-decay rate.
In each difference of primitive operators, a
surviving summand visits a tag outside the guard.
Join the test tree for a summand read at \(y\)
to a path from \(y_b\) to \(y\). The first native
cell meets the test window at \(y\), so its
padding is covered by the same initial allowance.
The first-escape
bound applies to the joined tree with its padded
carrier. The weight pays for the path to \(y\);
a fixed part of the native rate pays for the rest
of the escape. The remaining rate gives the
ordinary site sums. Thus for each primitive map
\(\mathsf O\) and bounded input \(f\), in the raw
scaled norms,
\[
 \|(\mathsf O(s)-\mathsf O(s_G))f\|_{w_b}
       \le C e^{-c d_\Delta}\|f\|_\infty .
\]
The source prepared before switching is the same in
both equations. If a source contribution is gated,
its change has the same weighted bound, since its
complete stencil meets the exterior.
The change in the initial multiplier
\(\lambda^0=\mathsf K\mathbf J\) satisfies the
operator difference bound as well.

Subtract the two reduced fixed-point equations in
the coordinates \((T,\lambda-\lambda^0)\).
The first difference, between two iterates of the
same map, costs at most half their distance in the
weighted norm. At the second solution the change of
map costs \(C e^{-c d_\Delta}\), by the preceding
bound and the local weighted estimate for
\(\mathsf C'^{\,\mathsf T}_T\Delta\lambda^0\).
This proves the bound for the solution difference
in that norm. The local endpoint read at \(d\) can
lose the factor \(e^{\eta d_{\mathbf B}(y_b,y_d)}\).
Use Cauchy on the smaller fixed physical source
disc to bound the difference of its derivatives.
For the shear the rescaling of the input by \(g_j\)
cancels the division of the output by \(g_j\).
Independently, the weighted source estimate bounds
either derivative by \(C e^{-\mu d_{\mathbf B}(y_d,y_c)}\)
for an input at \(c\). The geometric mean of the two
bounds gives, before summing,
\[
 \bigl|\mathsf L_d(s;c)-\mathsf L_d(s_G;c)\bigr|
 \le C e^{-c_0 d_\Delta}
       w_b(y_d)^{-1/2} e^{-\frac{\mu}{2}d_{\mathbf B}(y_d,y_c)} .
\]
For the transposed kernel use
\(w_b(y_d)^{-1/2}\le
w_b(y_c)^{-1/2}e^{\frac{\eta}{2}d_{\mathbf B}(y_d,y_c)}\).
The estimate holds on the larger switch disc, so
the same escape factor can be retained when taking
Cauchy coefficients that use an exterior switch.
For a product difference, telescope the factors.
At the changed factor the preceding bound leaves
the inverse weight at its entry site. The triangle
inequality bounds this inverse weight by the
product of \(e^{\eta d_{\mathbf B}/2}\) along the
preceding legs from the anchor. The unused decay
on those legs pays this product, with room for the
row and column sums. The inverse root uses the
local-resolvent carriers below.

This also gives a comparison away from a changed
seam. Use a common \(R\) with matching traces and
complete source and endpoint stencils as above.
In each setting restrict the switches to \(R\);
the restricted read at \(b\) is the same in the
two settings. For each unrestricted read subtract
its own restricted read by the weighted estimate.
The error is the sum of the two escape bounds,
each in its own weighted distance. The derivative
comparison uses the same physical input coordinate
and the same Cauchy disc on both sides. A carrier
that reaches a changed layer, boundary action or
completion stencil can instead contribute an
unmatched term. In the conditional strip form use
the native kept and rejected chains at each
evaluation, including the crossing and coupling
terms in the interaction. At the collar use the
actual pivot and seam curvatures in the native
Hessian below.

To see the normalization cost, set \(A=C_W^{-1}\)
and \(F=\Gamma_W\). In the representation with
auxiliary standard Gaussian noise the joint
precision is
\[
 K=\begin{pmatrix}
    A & -F\\
    -F^{\mathsf T} & I+F^{\mathsf T}C_WF
   \end{pmatrix},\qquad \det K=\det A .
\]
At the real reference it has a uniform floor
by the floor of \(A\), the bounded norms, and
the triangular change of variables. Let \(e_b\)
be the sum of the absolute row perturbations
of \(A\) and \(F\) at active slot \(b\). With
bounded absolute row and column sums,
\(C_W-C_{W,0}=-C_W(A-A_0)C_{W,0}\) and the
three-factor expansion of
\(F^{\mathsf T}C_WF-F_0^{\mathsf T}C_{W,0}F_0\)
bound the entrywise sum of \(K-K_0\) by
\(C\sum_b e_b\). This also bounds its trace
norm. The uniform operator perturbation is
small on the admitted section. Let \(P\)
be the projection onto the \(x\) coordinates.
After multiplying the absolute density by
\(e^{\gamma_2|x|^2}\), the real precision is
\(\operatorname{Re}K-2\gamma_2P\). In the
centered case the logarithm of its total mass is
\[
 \tfrac12\operatorname{Re}\operatorname{Tr}\log A
    -\tfrac12\log\det(\operatorname{Re}K-2\gamma_2P).
\]
At the reference the two determinants agree.
For a positive reference \(H_0\) of uniform
floor and a small perturbation \(E\), the
series for \(\operatorname{Tr}\log(I+H_0^{-1}E)\)
is bounded in absolute value by \(C\|E\|_1\),
where \(\|\cdot\|_1\) is the trace norm. Apply
this to \(A\) and to the real precision.
The projection has trace \(n_x\), so the tilt
cost is at most \(C\gamma_2 n_x\).
For complex \(m\) first translate the real
\(x\) by \(\operatorname{Re}m\). In the real
exponent the imaginary mean then contributes
a constant bounded by \(C\|\operatorname{Im}m\|^2\)
and a real linear term with norm bounded by
\(C\|\operatorname{Im}m\|\). The tilt contributes
the constant \(\gamma_2\|\operatorname{Re}m\|^2\)
and a linear term of norm \(2\gamma_2\|\operatorname{Re}m\|\).
Completing the square using the uniform floor gives
the additional cost
\(C\|\operatorname{Im}m\|^2+C\gamma_2\|\operatorname{Re}m\|^2\).
The offset on the physical section is
exponentially small. A free source that
moves the mean can be restricted to a disc
on which \(N^*\|m\|_\infty^2\) is small,
leaving room for the physical shift.

%
%
The guard and pair comparisons give
\(e_b\le C(e^{-c d_\Delta}+\epsilon_{\rm pair})\)
for \(b\) in \(\Delta\), with the switch loss
included in the constant. The running pair
radii give \(N^*\epsilon_{\rm pair}=o(1)\)
on the joint Gaussian section. Use the reference
coupling of the current scale in this Gaussian.
The difference from the position-dependent
Wilson weight is in the interaction. In
particular the separate disc for ratios in
the incoming terms does not enter \(e_b\).
If the Gaussian itself is rescaled by a ratio
\(\tau\), the same argument charges
\(Cn_\Delta|\tau-1|\) in each occupied cell,
so that option requires the corresponding
occupancy bound on its radius.
Take
\(\gamma_2\) a fixed small multiple of
\(1/N^*\). After integrating out \(y\), the
tilted total variation of the cut marginal is
bounded by the tilted absolute mass of the joint
density. Thus in any cut component the logarithm
of this bound is at most
\[
 C\sum_{\Delta} n_\Delta
 \bigl(e^{-c d_\Delta}+\epsilon_{\rm pair}+\gamma_2
       +\|\operatorname{Im}m\|_{\infty,\Delta}^2
       +\gamma_2\|\operatorname{Re}m\|_{\infty,\Delta}^2\bigr),
\]
where the sum is over its occupied cells.
Each summand is bounded, and a hot coordinate
at scale \((\log\theta_j)^{\mathfrak p_1}\)
gains \(\exp[-c(\log\theta_j)^{2\mathfrak p_1-4r}]\).
This pays the stage bound
\(\exp(C(\log\theta_j)^{4r})\) because
\(2\mathfrak p_1>8r\). The determinant of
the full active reference still uses the old
stage normalizer. These bounds use the native
response and local-resolvent families in the
fixed ambient stage. For a change of ambient
tower the corresponding completed rows and
native carriers must agree on the component
where the comparison is made.

%
%
In the complete stage integral the reference is the
sum of the homogeneous retained strip quadratic and
the homogeneous collar quadratic. The collar linear
term in the constant-weight Wilson part vanishes by
stationarity on the fixed averaging fibre. The remaining interactions,
including cross-component strip reads and changes of
the background or coupling, are in the potentials.
Unintegrated near fields occur there with their
original real values and switches. Thus \(\ell=0\)
for this reference on the complete integration vector.
The new mean is \(-\mathsf T^{-1}d=-d\), since \(d\)
has only strip components. On a guarded read take the
continuation \(-P_Wd(\hat s)\). It has the same
carrier as the offset and fits the smaller sharp
threshold on the crossing slots. At full switches
it is the mean of the full Gaussian restricted to
\(W\). The covariance formula accounts for the complementary
coordinates. The real sharp tests still act in the
specified fixed integration charts. All Taylor terms
of the potentials, including any linear terms, remain
in the interaction.

%
On the admitted joint switched section write the
exact active inverse as
\(\mathsf T z+d+r(z)\). With \(T\) bounding all
arguments on their actual enlarged discs, the
nonlinear Gaussian correction is
\[
 -r(z)^{\mathsf T}D(\mathsf T z+d)
                 -\tfrac12 r(z)^{\mathsf T}Dr(z).
\]
The product-logarithm and fixed-radius source
estimates give \(\|r\|_\infty\le Cg_jT^2\).
With the same endpoint gates for the exact and
affine coordinates, \(r_b\) vanishes if either
endpoint cell is inactive, and a surviving cut
coefficient reads only the wall component of those
cells. Cauchy's bound on the switched disc sums
these coefficients with an exponential weight on
their attached cells. The weighted row sums of
\(D\) and \(\mathsf T\) then bound the displayed
terms, filed at a changed strip slot, by
\(C(g_j T^3+g_j^2T^4)\) in the weighted sum.
For the term with two remainders use the row bound
at one of its slots and join the other slot's
carrier along the same matrix chain. This requires
only one sum over the changed anchors. The Haar
and pivot terms have the rooted bounds just proved.
For the norm pinned at any cell of a joined carrier,
keep a portion of the native spatial switch and
walk decay. The distance from that cell to the
changed anchor is at most the number of cells
in the connected carrier, with a fixed enlargement
for the endpoint blocks. Use the reserved decay
to sum over anchor cells, whose face graph has
bounded degree. Multiplication by the slot count
of an owner cube, or by the logarithmic slot count
over the component, is harmless since \(g_j T^3\)
is smaller than any fixed inverse logarithmic
power.

%
%
For a mixed strip term with a fixed exterior
slot, substitute only its active argument and
keep the original switch on any near exterior
argument. Its affine change involves
\(\mathsf S z+d\), so estimate it using the old
chain bound. In the physical stage the selected
strip masks are whole kept components. The kept
form therefore has no mixed block against the
fixed exterior components. Rejected chains and
the ratio term are still boundary interactions.
Use their formation rule chain by chain. Under
the original local keeping predicate, rejection
costs a cell of travel. This gives
\(e^{-c\delta_\star W_t}\) on the enlarged strip box
before the remaining carrier decay. The ratio term
uses kept chains, possibly short ones, with a factor
\(|1-\tau_k|\). Both estimates absorb the two enlarged
field radii. A Gaussian cut inside a kept component
uses the full marginal formula for the omitted variables.

The stage bounds allocate separate margins for the
pointwise reads, the tower reads and the core. The
independent current increment uses the pointwise margin.
The tower current, formed by averaging and solving the
auxiliary problem for the substituted configuration, has
its own estimate. On the graded interior windows, \cref{atc:theorem,atc:localization}
apply to settled elements and control the changes by the weighted
potential and its first difference. For an address read to depth
$p$, its nonempty level-$(p-1)$ localization cube forces at least
$M/L\ge L^3\rho\ge6\rho$ sites per torus period at that depth.
This supplies the period floor in the ambient estimate. The full
current stencil and the complete row reads remain those of that
fixed address.

%
%
Here is the budget on a fixed inherited address.
Write \(\mu(x)\) for the layer of the preceding
response sequence and \(\ell(x)\le\mu(x)\)
for the layer of the earlier address, truncated
at its reading depth when required. Put \(q=j+1\),
and write \(D_n(x\mid l)=|l-q|\) inside a
treated component and \((t-q)+(t-l)\)
outside the treated components of this stage
of level \(t\). The decrease in the coefficient
of the radius at a fixed index \(l\) is
\[
 \frac{\mathsf g_t(l)}{2\mathsf N_t}
       +\beta_{\rm R}2^{-|l-q|}
                         \mathbf 1_{\mathcal Z_n}(x)
 \ \ge\ \gamma_\star 2^{-D_n(x\mid l)},\qquad
 \gamma_\star=\frac{\theta_{\rm s}\beta_{\rm s}}{2\mathsf N_t}.
\]
For a fixed tower address its left-hand read
and its lift are unchanged. Multiply this
coefficient gap by
\(L^{i-r}L_\flat^{2e_{\sigma''}(x)}
  \alpha^{\sigma''|\sigma_n}_{0,\ell(x)}\)
for the slack in normalized tower units.
Where the layer changes from \(j\)
to \(q\), the ratio of the two scale weights
gives a fixed fractional margin, including
the lift gained there. The stage transition
uses the same unfinished top radius.

%
%
%
At an inherited core the restriction also
keeps the left-hand read and its radius
number. The coefficient in the restricted
bound is the later coefficient at the
core's index. The gap at that index supplies
the core margin throughout the original
domain, even outside the treated component.
In the room estimate let \(\mathsf v(b)\)
be the pointwise gap profile in the input
length \(\mathsf L(b)\). For a core of scale
\(\xi_i\) the shift uses the profile
\(\xi_i\mathsf v(b)/\mathsf L(b)\). The
layered profile comparison bounds it by a
fixed multiple of the core margin, with the
core's radius profile included. This is the
core condition in the same room estimate
that handles the pointwise and tower clauses.
Thus the stage amplitude bound also gives
admission to the inherited graded cores.

The weighted path from the changed zone
to \(x\) must traverse the intervening layers.
On leaving a treated component it first
reaches the original top layer \(t\), since
the quiet exterior of level \(t\) is inset.
Thus, with \(d^-(x,\Sigma')\) the preceding
weighted distance, \(D_n(x\mid\mu(x))\) is
at most a fixed constant plus
\(d^-(x,\Sigma')/\varrho_0\). Here
\(\varrho_0\) is the layer traversal floor.
The pinned response, in its own normalized
units, has profile bounded by
\[
 s_1(x)=C g_jp_0(g_j)|\lambda_{\mathfrak w}|
                 |B|_{\mathfrak w}
                 e^{-\delta_1 d^-(x,\Sigma')}.
\]
The constant includes the turn and switch
costs, and \(\delta_1\le\delta_{\mathbf H}/2\).
At the earlier address the shift and its
first difference are bounded by
\(L^{-(\mu-\ell)}s_1\) in that address's
units. The pointwise plaquette and current
comparisons have at least this scale gain.

The settled tower estimate convolves this
last bound with the tower decay in the
earlier ambient distance \(d^{\rm a}\).
One has \(d^{\rm a}\ge d^-\). The difference
of the layer losses \(\mu-\ell\) at two
points is at most \(2+2d^{\rm a}/\varrho_0\),
by counting crossings in each sequence.
It follows that this convolution, in the
form of the weighted supremum in the
tower estimate, is at most a constant times
\(L^{-(\mu(x)-\ell(x))}s_1(x)\) when
\(\delta_1+2\log L/\varrho_0\le\delta_{\rm T}\).
The required joint tower-class admission
uses the same small shift on the full
localization as in the estimate below.

Divide by the gap at \(\ell(x)\).
The bound
\[
 D_n(x\mid\ell(x))\le
 D_n(x\mid\mu(x))+\mu(x)-\ell(x)
\]
leaves the factor \((2/L)^{\mu(x)-\ell(x)}\).
It pays the part of the adjacent-coupling
comparison coming from lowering the index
to \(\ell(x)\). In the running radius
\(C_\iota\theta_l^{-1/2}(\log\theta_l)^{q_{\rm rad}}\)
the function of \(\theta_l\) decreases at
weak coupling. For an earlier index the
logarithmic factor in the radius is at
least as large, and the bounded drift of
\(\theta_l\) costs at most the square root
of one plus the level difference. For the remaining comparison
the bounded running increments cost at most
\(C(1+\sqrt{|j-\mu(x)|})\). The factor
\(e^{-\delta_1 d^-/2}\) pays this factor and
the factor \(2^{D_n(x\mid\mu(x))}\), since
\(\delta_1\varrho_0\ge2\).
This leaves the distance factor at rate
\(\delta_1/2\).

For an inherited domain \(X\),
write \(\mathsf d_n(X)\) for its distance in cubes from
the changed layer and \(|B|_{\mathfrak w}\) for the
supremum on component \(\mathfrak w\) divided by the
box radius. The stage profile budgets the increment by
\[
 \mathsf a_n^\delta |\lambda_{\mathfrak w}|
 |B|_{\mathfrak w}<e^{\delta_\star\mathsf d_n(X)},
 \qquad
 \mathsf a_n^\delta
 =O\bigl((\log\theta_j)^{-(q_{\rm rad}-\mathfrak p_0-r)}\bigr).
\]
Here \(q_{\rm rad}\) is the exponent of the
running radii. The power comes from
\(\gamma_\star^{-1}=O(\mathsf N_t)\),
\(\mathsf N_t=O((\log\theta_j)^r)\), and the
ratio \(g_jp_0(g_j)/\alpha_{0,j}\).
The physical values \(\lambda_{\mathfrak w}=1\) satisfy
this inequality for small coupling. For the common
substitution take the intersection of these discs
over the actual readers. The uniform disc obtained
at \(\mathsf d_n=0\) is sufficient for joint admission.
On the readerwise larger discs the same common
formula is restricted to that reader. The fixed
analytic disc of the ordinary joined centre is
used for its own unamplified source. The stage
discs are also limited by
\(|\lambda_{\mathfrak w}|e^\varkappa<c_{L,\varkappa}\varrho_{\rm st}\),
where \(|s|<e^\varkappa\) and \(\varrho_{\rm st}\) is the
strip radius in units of its field threshold. The fixed
factor \(c_{L,\varkappa}\) leaves the margin for the inverse strip
arguments. The physical value remains interior for small
coupling.

The real layer allowance \(e_l=\check R_l^2\varepsilon_l\)
is used in admitting the base pair. The largest
compressed conversion for current or local gauges
costs \(O(\check R_l^3)\). Thus it requires the
\(\mathfrak p_0+5r\) logarithmic power of the real
endpoint estimate. The stage floor
\(q_{\rm rad}\ge\mathfrak p_0+7r+1\) leaves room also for a
sum over the \(\mathsf N_i\) stages. The response
amplitude in the displayed profile inequality
is determined by the free collar box and the
running radius. It uses the actual collar
increment, rather than the bound for the base
pair.

\paragraph{Domains of the earlier entries.}
%
%
Use the graded formation rule from the bare density
onward. A tower address consists of an
earlier structure \(\sigma''\), a localization \(Z\) of
level \(i\), and a depth \(1\le\mathsf p\le i\). At
level zero the direct bare reads use the pointwise
clause. The tower of an earlier address is
the meet of the localization tower with the sequence of
\(\sigma''\), through depth \(\mathsf p\). For
\(\mathsf p\le r\le i\) remove two layers of level-\((r-1)\)
\(M\)-cubes from \(Z\) to form \(Z^{\wr r}\). Read a
plaquette or a current stencil when all its corners lie in
this window. Its base point then lies in the layer of the
tower of index \(\min\{\mathsf p,\ell\}\), where
\(\ell=\min\{m_{\sigma''}(x),i\}\) is the nominal layer.
Indeed \(Z^{\wr r}\subset Z^{\wr\mathsf p}\). The
two-layer localization comparison for \(M\ge L^3M_1\)
gives agreement of the tower and earlier region
memberships through \(\mathsf p\) on that window, where
\(M_1\) is the cube size used in the localization.
In bare units the curvature and current bounds use
the respective factors \(L^{-2\ell}\) and
\(L^{-2\ell-\min\{\mathsf p,\ell\}}\). Both are multiplied
by the layer radius with its lift and ambient coefficient,
and by \(L^{i-r}\). The tower itself and the ambient
coefficient thus play different roles under restriction.

\paragraph{The allowance and the stage gap.}
\label{rtcg:stage-profile}
The profile hypothesis in \cref{rtcg:cut} follows from the
actual radius clauses at an inherited stage address. We give
the comparison, including its dependence on the stage. Fix
that address, its ambient structure $\sigma''$, its level $i$
and its window index $r$. Put
$m(z)=\min\{m_{\sigma''}(z),i\}$. This is the nominal
ambient layer, denoted by $\ell$ above; the native settled
tower index is $\min\{\mathsf p,m(z)\}$. The clause radius
and its stage gap are indexed by the nominal layer.
Abbreviate by $\alpha_m$ the radius
$\alpha^{\sigma''|\sigma_n}_{0,m}$ selected by the bare-level,
unfinished-top and adjacent-index conventions. Let
$e(z)=e_{\sigma''}(z)$. Choose a fixed $C_\alpha\ge1$ large
enough to include the ratio of the potential radius to the
curvature radius, the adjacent-index comparisons and the
change of base point within a complete stencil. Here we use
the actual running radius system of the stage
formation, whose components have the same logarithmic power,
or a fixed absolute system. The ratio of its two components
is fixed in either case; taking the minimum of neighboring
running radii preserves that ratio. Thus $C_\alpha$ is
independent of the stage, coupling and depth. Define
\begin{equation}\label{rtcg:stage-allowance}
 B(z)=(1+\beta_{\rm s})C_\alpha
                         L_\flat^{2e(z)}\alpha_{m(z)}.
\end{equation}
The pointwise clauses on this inherited address give
$A(z)\le B(z)$, with $A$ as in \eqref{rtcg:profile}.
The independent current does not enter $A$. At a real stage
endpoint the stronger tested allowance
$\check R_l^2\varepsilon_l$, including its compressed-scale
conversion, also fits this bound by the endpoint calculation.

Fix $0<\kappa<\delta/4$, where $\delta$ is a common rate
in \cref{rtcg:alignment,rtcg:cut}. Decrease the common stage
rate $\delta_0$ if necessary so that $0<\delta_0\le\delta$.
Set
\begin{equation}\label{rtcg:allowance-envelope}
 R_\kappa(x)=\max_z e^{-\kappa d^\flat(x,z)}B(z).
\end{equation}
This is a deterministic function of the structure and its
radii. The triangle inequality gives
$R_\kappa(z)\le e^{\kappa d^\flat(x,z)}R_\kappa(x)$,
and $A\le B\le R_\kappa$. Thus
\eqref{rtcg:allowance} holds with $C_0=C_1=1$ and
$\mathcal R=R_\kappa$, without an additional profile assumption.

For the stage profiles one also has
\begin{equation}\label{rtcg:envelope-comparison}
 R_\kappa(x)\le C_B B(x),
\end{equation}
where $C_B$ is fixed before the coupling ceiling. Here is the
layer comparison that proves this assertion. Let $d=d^\flat(x,z)$
and let $\varrho_0$ be the common traversal floor of the
inherited determining sequence and its original regularity
ancestor. The latter sequence is coarser; truncating the
inherited sequence at $i$ only increases its metric. Counting
the intervening separated layers in each sequence gives
\[
 |m(z)-m(x)|\le2+2d/\varrho_0,\qquad
 |e(z)-e(x)|\le4+4d/\varrho_0.
\]
Indeed a path between layers differing by $h$ crosses at
least $h-2$ complete intervening collars. Apply this to the
two layer functions whose difference is $e$. The endpoint
conventions for the radius move its selected index by at most
one at either end.

For the running function
$\theta^{-1/2}(\log\theta)^{q_{\rm rad}}$, choose the
coupling ceiling so that $\log\theta>2q_{\rm rad}$.
This function decreases with $\theta$. The frame bounds
\eqref{eq:rg-objects-R19} then give, for radius indices $a,b$,
$\alpha_a/\alpha_b\le C(1+\sqrt{|a-b|})$.
To check this in the direction in which the radius grows,
use $\theta_b\le\theta_a+\bar B|a-b|+c_5$; the square-root
factor bounds the power of $\theta$, while the logarithmic
factor decreases. The opposite direction costs only the
fixed correction $c_5$ at the weak-coupling ceiling.
The minimum of neighboring radii obeys the same estimate.
Consequently
\[
 \frac{B(z)}{B(x)}\le
 C L_\flat^8(1+\sqrt{4+2d/\varrho_0})
               e^{8(\log L_\flat)d/\varrho_0}
 \le C_B e^{(8\log L_\flat+1)d/\varrho_0}.
\]
Choose the fixed geometric floor so that
$(8\log L_\flat+1)/\varrho_0\le\kappa/2$.
Multiplication by $e^{-\kappa d}$ and maximization over $z$
prove \eqref{rtcg:envelope-comparison}. This choice enlarges
only the fixed geometric constants.

Write $\mathsf{gp}_n(x,m)$ for the deterministic margin
between the two stage bounds at the read. The coefficient
decrease computed above gives, at a fixed inherited index,
\begin{equation}\label{rtcg:stage-gap}
 \mathsf{gp}_n(x,m)\ge
 \gamma_\star2^{-D_n(x\mid m)}
              L^{i-r}L_\flat^{2e(x)}\alpha_m,
 \qquad
 \gamma_\star=\frac{\theta_{\rm s}\beta_{\rm s}}{2\mathsf N_t}.
\end{equation}
Define the number used in the cut estimate by
\begin{equation}\label{rtcg:stage-cut-factor}
 a_{\star,n,x}=
 \min\left\{1,\frac{\gamma_\star2^{-D_n(x\mid m)}}
                    {C_B(1+\beta_{\rm s})C_\alpha}\right\}.
\end{equation}
It is positive and depends only on the address, the stage
and the read. Combining \eqref{rtcg:stage-allowance},
\eqref{rtcg:envelope-comparison} and \eqref{rtcg:stage-gap}
gives
\[
 R_\kappa(x)\le a_{\star,n,x}^{-1}\mathsf{gp}_n(x,m).
\]
The radius and lift cancel, and $L^{i-r}\ge1$ improves
the bound. No lower bound for $a_{\star,n,x}$ independent
of the stage and the read is asserted.
By \cref{rtcg:cut}, a cut satisfying its complete-carrier
separation changes either settled read by at most
\begin{equation}\label{rtcg:stage-cut-bound}
 K_{\rm D}a_{\star,n,x}^{-1}e^{-c\delta_0D}
                                 \mathsf{gp}_n(x,m).
\end{equation}
In particular, to spend at most a fraction $\vartheta\in(0,1)$
of that margin it suffices to impose
\begin{equation}\label{rtcg:stage-cut-distance}
 D\ge\max\left\{K_{\rm cl}M,
       \frac{1}{c\delta_0}
          \log\frac{2K_{\rm D}}{\vartheta a_{\star,n,x}}\right\}.
\end{equation}
The same criterion applies simultaneously to finitely many
reads by taking the largest required distance. In particular
the dyadic coefficient loss is paid by the displayed distance
when such simultaneous margin preservation is required.

%
%
\paragraph{A geometric cut of a settled tower clause.}
In the units of the earlier address, cut \(Z\) by a union \(R\)
of level-\(t\) cubes and keep the component \(Z_R\) of
\(Z\cap R\) containing the base \(x\) of the read. Components
here join also at corners. Suppose that \(x\) is at distance
at least \(\mathsf D+3M_1\) from \(Z\setminus R\), with
\(\mathsf D\ge K_{\rm cl}M\) as in the localization bound.
The cut tower is finer by monotonicity of the layer
construction under enlargement of the complement.
The complete-carrier argument following \cref{rtcg:cut}
locates every changed datum read within \(C_{\rm cut}M_1\)
of the actual cut \(Z\setminus R\), including the old parent
pillars and outward endpoint blocks. It also treats discarded
components by first exit through the actual cut; it requires
no extra spacing assumption between such components.
Increase \(K_{\rm cl}\), once, so that
\((C_{\rm cut}-3)_+M_1\le\mathsf D/2\).
Every changed complete carrier is then at distance at least
\(\mathsf D/2\) from the read. The retained \(r\)-window,
complete stencil, free pillars, nominal layer \(m\) and native tower index
\(\min\{\mathsf p,m\}\) are unchanged there.

For the same admitted pair, \cref{rtcg:cut} and
\eqref{rtcg:stage-cut-bound} bound the change in the magnitude
of either settled read by
\[
 K_{\rm D}a_{\star,n,x}^{-1}e^{-c\delta_0\mathsf D}
                                  \mathsf{gp}_n(x,m).
\]
The positive factor \(a_{\star,n,x}\) is the explicit number
in \eqref{rtcg:stage-cut-factor}; it depends on the stage
and the read. The bound of the clause depends on the
location, indices and structures, so it has the same
right-hand side before and after this cut. An available
margin larger than the displayed error survives the
restriction. In particular \eqref{rtcg:stage-cut-distance}
is sufficient for preserving a prescribed fraction of
the stage gap. This proves the row and domain comparison
with its stated margin and distance condition. A composed
reader uses the stencils and switches of its actual presentation.

\paragraph{Shifting the field in a fixed tower.}
On these windows the tower comparison estimates the
recomputed current itself. Write \(\mathbf U=e^{i\eta A'}U\)
and \(W^u=e^{i\eta(A'\boxplus\mathbf A^\circ)}U\) using one real
factor \(U\) and the local logarithm of the shifted pair.
When both potentials lie in the admitted tower class,
\cref{atc:theorem,atc:localization} give the normalized current increment
\[
 \begin{split}
 L^{2\ell+\min\{\mathsf p,\ell\}}
 |J^0_{\mathsf p,Z}(W^u)(b)-J^0_{\mathsf p,Z}(\mathbf U)(b)|
 &\le C\max_{c\in\mathsf Z^+} e^{-\delta_{\rm T}d^\flat(y_b,y_c)}
 \max\{l(c)|v(c)|,l(c)^2|\nabla_U v|_{\mathsf Z}(c)\},\\
  v&=A'\boxplus\mathbf A^\circ-A'.
 \end{split}
\]
Here \(J^0\) is in bare units, and \(\mathsf Z^+\)
consists of the positively oriented inner bare
bonds of the localized first member. Each \(y_c\)
is the earlier ambient pillar top of the base
point of \(c\). Use that sequence's distance
\(d^\flat\) and local length \(l\). The
comparison allows the base point of the
full stencil within \(2\xi_{\mathsf p}\) of
\(Z^{\wr\mathsf p}\) when the whole stencil is in the
localization's bond set. The agreement of the region
memberships on that halo covers stencils that straddle
an ambient layer boundary. The same estimate holds
for the plaquette increment with weight \(L^{2\ell}\).

%
Use the real factor \(U\) of the input on the whole
localized first member. The pinning estimate bounds
the relative logarithm of \(W^u\mathbf U^{-1}\) and
its covariant first differences with this same real
factor. Use the resulting \(\mathbf A^\circ\) in the
product \(A'\boxplus\mathbf A^\circ\). The small product
chart gives a single holomorphic choice on each
fixed-real-factor section. In particular the real
factor of the recomputed tower stays \(U\).

The ambient lengths of the earlier address are at
most the lengths at the input, including at the base
points of the full stencils. The pointwise input
radii therefore bound \(\mathsf n_c(A')\) and the
normalized curvature of \(U\) on the entire first
member. This uses the full localization for the
radius test, while its inset windows specify the
output stencils. Let \(\mathsf u(c)\) be the maximum
of the earlier weighted value and first-difference
bounds for the relative shift, and let
\(\mathsf u^+(c)\) include the neighboring bonds of
the product estimate. Apply the common radius
budget on these complete carriers. The turn cost
is fixed in terms of \(L,M\), and the bound on
\(\mathsf u\) there is \(C L^{-(\mu-\ell)}s_1\).
The same bound, with the neighboring maximum, is
small uniformly as the coupling ceiling decreases.
Thus the base norms can be put below \(a_{\rm T}/2\)
and the shift satisfies
\[
 \mathsf u(c)\le(256L^2)^{-1},\qquad
 \mathsf n_c(A')+2\mathsf u^+(c)<a_{\rm T}
\]
throughout the section. The geometric floor of the
tower theorem is the separation floor for its fixed
localization, chosen before the coupling ceiling.
The shift estimate now gives the recomputed tower
on this same mask. Fit the resulting normalized
increment in the available window margin, with
\(L^{i-r}\) at window \(r\), by the preceding budget.

Give each boundary term, potential away from the strip,
stage correction and interaction term the structure of its
formation. Carriage retains its level, localization domain,
function and space. When a stored coordinate is integrated,
its occurrence in the prospective evaluation is set to zero
and its switch is discarded. Before the first stage the sum
of earlier stage corrections is zero. At each subsequent
stage keep the original domain and stored attachment of
each correction already formed.
The old evaluation of the sum uses the levelled triple
restricted to each of those domains. The prospective
evaluation uses its own triple and sets the strip field
being integrated to zero. The difference of the two
evaluations enters the new collar integral. Lifting a lower
boundary term to a larger support for the expansion leaves
its original argument domain in place.

%
%
On the common stage profile bound, the levelled
substitution sends the stage-\((n+1)\) space into
every earlier graded space. The prospective evaluation
uses restriction from the output structure. At each
slot the formation structure precedes its own level's
output structure, which in turn precedes the current
output. For a carried slot take the standard reduction of
the outgoing history on its original domain. This
reduction has the completed trace of the healed
label on the domain's windows, treats no component
on its enlarged read, and keeps the recorded
drop-outs there. If the enlarged read misses the
treated components, one may use the history itself.
The traces on the smaller windows agree for all
these choices. Restrict the reduction to the old
level before taking the term and its space. The
first members of its tower clauses are contained
in the original domain. Intersecting the same
read of the old sequence with the same localization
tower gives the same clause on each such first
member. Identify the
original first member, the determining rows including
the one-step completions, and the pairs of tower and
window indices before comparing radii. The order of
layers and the distance comparison in the profile
estimate are comparisons within that common labelled
sequence. For a stored region transported to the new
label they apply after the address has been matched:
on its full read the sequences agree, and the tower,
window, lifts and address radii persist. Pointwise
descent between two output labels by itself gives
only the pointwise part of this comparison. A term
whose old hull has been absorbed uses the new filed
term, with its own formation domain.

For example, a kept-chain summand of the ordinary
ratio correction that reads a far strip bond has
the form
\[
 -\tfrac12(\tau_l-1)
       \sum_{b,b'} a_b^{\mathsf T}D_\omega(P)_{bb'}a_{b'} .
\]
Here \(l\) is its original fluctuation level and
\(D_\omega\) is the kept chain in the ordinary
conditional form. On its full carrier in the
original two-cell strip trace, the regular-pair
comparison applies to the pinned pair with the
original free rows. It uses pointwise potential
and current estimates. Its far arguments are
the joint strip coordinates bounded above, and
its near arguments use the old multipliers.
The row estimate on the enlarged field disc is
\(C|\tau_l-1|(\varrho_{\rm st}\mathfrak b)^2\)
times the old carrier decay, within the tube
bound. A rejected long chain uses its own
pair domain on its complete carrier.

For an inherited cluster or strip-transfer
entry that uses an earlier activity, keep the two
blocks and parent transports of the strip at its
original step, as well as every tower address
\((\sigma'',i,Z;\mathsf p,r)\) of its formed
argument. For instance, in a boundary cluster
\(E^\sharp(X)\) a vertex made from an incoming
boundary term uses its ordinary potential on
\(Y\), with \(Y^\sharp\subset X\). Recover the
ordinary trace on the one-layer windows of \(X\)
at each level. It specifies the old cuts, tiers
and which input potential has this support.
The ordinary fluctuation substitution in that
potential uses the determining rows and gauge
cubes of \(Y^\sharp\). Its fluctuation variables
are those of the ordinary formation, distinct
from the free variables of the later collar.

%
%
%
%
At the original step \(k\), a vertex of species
\(V^B\) is made from localized differences
\(f(P_{\rm o};\mathfrak q)-f(P;\mathfrak q)\)
of an incoming boundary term on its original
domain, say \(D\). In both evaluations the
stored arguments \(\mathfrak q\) are the same.
The ordinary substitution in \(P_{\rm o}\)
uses the corrected ordinary fluctuation at
\(\mathcal M^k\mathbf U\). With the real
factor \(U\) fixed, the substituted link
and the independent current are
\[
 \begin{split}
 e^{i\eta A''(b)}U(b)
 &=u(b_-)\,e^{i\eta\lambda_{\rm o}H(b)}
                 \mathbf U(b)\,u(b_+)^{-1},\\
 \mathbf J''&=\operatorname{Ad}(u)
                 (\mathbf J+\lambda_{\rm o}\mathsf j).
 \end{split}
\]
Here \(u\) is the ordinary pinned levelling
map and \(\mathsf j\) the reduced current
increment of this response. The identity of
the second line with the physical current
holds at the full real arrival. The complex
tower clauses use the recomputed tower current
of the shifted potential on their own windows.

For admission at this inner substitution, the
ordinary step supplies the coefficient gap
\(c_0\,2^{-(k+1-n)}\) at the index \(n\) of
each earlier clause, including when the
reference radius is the smaller of two adjacent
radii. Use that gap times the clause's radius
as the margin in the weighted potential
estimate. For a tower clause the comparison
of profiles includes the length ratio between
the tower and the input sequence, and the
exponential in their distance from the stencil.
The separation of successive layers pays the
geometric gap ratio and the square-root ratio
of the running radii. This gives the pointwise,
tower and core room for the ordinary response.
Write \(r_D\) for this room after the reduction
for the pinned map. For an ordinary field box
bounded by \(\mathsf C_{\mathcal A}\delta_k\),
use the separate amplitude radius given by
\[
 R_{{\rm o},D}^{-1}
  =\frac{C}{r_D}\,\mathsf C_{\mathcal A}
       \frac{\delta_k}{\alpha_*}
       e^{C\varkappa-\delta_*\mathsf d_k(D)}.
\]
The bound applies to the whole original graded
section, with \(R_{{\rm o},D}>2\) on the
enlarged ordinary formation box at small
coupling. Here one uses its enlarged
\(\mathsf C_{\mathcal A}\) in the radius
formula. The far stored fields in the vertex
retain their original radii.

In the closed cluster the integration over
the original ordinary \(B_I\) still uses its
fixed real cutoff balls and its original
conditional Gaussian, with the mean determined
by \(B_E\). At a moved crossing it is only
the external far slot of \(B_E\) that takes
the reconstructed coordinate. When the
outer substitution lies in the original
graded formation section and these external
slots lie in its field box, the original
absolute majorant of the cluster applies on
the whole integration section. The outer
stage margin is used for these membership
tests. The subsequent ordinary substitution
uses the gap and amplitude of its original
step, uniformly over the section just reached.

For the first earlier-stage correction that
reads this strip, carry the closed first-stage
function, including the carriers of its
original source closing. The second stage
uses the common stage triple on that full
formation domain, including its inherited
cores. The old evaluation uses the reconstructed
far field, and the prospective evaluation sets
that integrated field to zero. The original
source closing is already part of the
carried function. In this correction the
far slot is a stored-field argument. Its
parent targets are read through the pair.

For a tower-current read use the common real
factor and the earlier weighted potential bounds
on all of \(Z\) for both the base and the
pinned input. The six incident plaquettes
and the halo of the stencil must be in the
graded read of the same address. The new
filing hull carries the resulting term while
these evaluations use their original tower.
In a newly composed operator we must also
check the joint read of the retained datum
if it enters the strip inverse in a separate
role. The one-step forms and elimination
factors use their separate pair domains.

%
%
%
%
%
%
\paragraph{Carrying the formation reads through a cut.}
A formed summand has a fixed base read. Keep
the part of its original structure, row masks
and stored attachments used there. An earlier
normalizer or closed cluster term enters once
with its formed function. The closed carrier
includes the attachments of the earlier source
closing. The current substitution acts on this
closed function. It changes the pair and the
currently integrated variables, on the rooted
arms of the current bracket. In particular
the zero-increment value at the fixed base
has no new spatial switch. A bracket using
the old summand keeps its complete base,
and its response differences are formed on
that base with the new source arms.

An old quadratic or pivot factor instead
uses its one-step trace. For the strip
Schur form keep the original chain
classification into kept and rejected
terms. The rejected chains and the
coupling correction retain their original
filing in the exponent. At each full
evaluation their sum with the kept
form is the original conditional
quadratic. Retag a kept chain by the
parent image of its actual old tags,
as in the grading argument. Its
binary-cut criterion is inherited
through that image. A solved row
uses the original target and its
complete double block, including
the prescribed group links. Keep
its determinant with that row.

At a stored spectator the base includes
its original site, top, and full pillar.
In an actual stage structure a stored endpoint
at level \(l\) is a top or has its parent at
level \(l+1\). The complete level-\((l+1)\)
block is separated from level \(l+2\).
Its membership at level \(l+1\), and hence its
pillar and transport, persist while this
field remains stored. In a genuinely finer
geometric localization the same site may
become a top. The transport in that new
localization cannot be used for the earlier
occurrence just from the depth bound.

Use the fixed stage geometry in the binary
cuts of the response. For each original top
\(y\) keep the complete dependence set of the
pin, including the averaged contour, the
level-one neighbor margin, and the stencils
of its turning cells. Use the padded hull of
each response walk, which contains the blocks
of the output tops and of their four successive
neighbors. The node factor and contour use
internal links of the original block. The
turn uses its child blocks. Include the actual
two-bare-step stencil at the exceptional
level-one pins. There the fixed ambient
eligibility condition puts the nearby pillars
at levels at most one. For an actual stored
field of level zero the full parent block
is separated from layer two by the storage
margin. The complete pin and turn stencils
are over the two-neighborhood of the original
top. If a response at an output top \(z\) meets
this read, the symmetry of the neighborhood
puts the whole read in the four-neighborhood
of \(z\). This is the neighborhood in the
padded hull.

Tag every closed graded cell meeting the
padded hull. At each point of a surviving
hull the cells that contain that point
contain a neighborhood in their union. Their
star is connected through positive-area faces
by nested grid refinement. The connected hull
is therefore inside one active wall component,
also at faces between grades. For inherited
walks use the parent image of their actual
old tags. Pulling back binary switches along
the parent map gives exactly the old binary
test, even when the deduplicated complex
monomial uses fewer factors. Keep the old
closed functions as filed. For a new operator
without an admitted native walk presentation
we must still verify this padded-hull
condition on the actual pin and output
stencils.

On a binary cut a nonzero increment on
a pin carrier then has that whole carrier
in one active wall component. Otherwise
the increment on the carrier is zero.
On the original pillar form
\[
 p_y(a)=\tau[\mathbf U](a)^{-1}
 v[P;\mathbf A](y)^{-1}
 \tau[e^{i\eta\mathbf A}\mathbf U](a).
\]
If the response is zero there, \(p_y(a)=1\).
Otherwise the formula uses only the unchanged
pair on the fixed carrier and the increment
of the one wall component. On the child
turning cell use \(u=e^{-i\mu}\) with
\[
 \mu(x)=\operatorname{bp}_a(x)T_a[U](x)
             \tfrac1i\log\bigl(p_y(a)^{-1}\bigr),
\]
using the original bump and real-factor
transport. The turn is
supported in the inner third of that cell.
The pinned response therefore splits into the
responses of the wall components, with disjoint
nontrivial gauge factors. The cut restriction
agrees on each active pin because its carrier,
data and response traces agree.

Choose the pinned tops in the ambient
geometry once. Work on sections where all
needed tops are levelled on the whole
polydisc. The local pin estimate is
\(|p_y^{\pm1}-1|\le c_{\rm m}(L)\varepsilon_y(\mathbf A)\).
The turned pair costs at most
\(\mathsf C_{\rm m}\|(\mathbf A,\mathbf K)\|^+_y\),
where the norm includes the neighboring
level-one pillars. Check this bound against
the room of every earlier reader on the
complete turn polydisc. On that polydisc
gauge invariance at real turns and holomorphy
give independence from the auxiliary bump or
real-factor representation when the values at
the stored sites agree. On a fixed real-factor
section the transport equality extends from its
open real slice to the admitted complex branch
because both contour means there are holomorphic.
Equality of pairs gives the same point pins.
To identify the full turned triples compare
their real frames on the turning cells as well,
or use the invariance of the term on the
admitted turn polydisc. Throughout this argument
the stored values are fixed. In particular a
Borel near value stays in its original slot
and an inherited scalar switch keeps its old
radius. Only an actively integrated far field
is multiplied by \(\lambda_{\mathfrak w}s\);
for an actively integrated near field this
multiplier is its scalar switch and the value
itself is fixed. The recentered crossings use the far
arguments already checked on their
two blocks. A new reader which uses
a different near argument must
supply the corresponding real
coordinate and switch rule.

For an island factor keep its standing
history, integration measure and whole
exempt hull. At fixed internal real
data take the old and new pair
evaluations of its analytic exponent.
The internal real weight is the same
in those evaluations. The island
hull is one fixed read, with new
response arms at its attachment.
The cut factorization is applied
at fixed island data, and the
original island integral is taken
once on its carrier. This uses the
old bounds for the analytic exponent
and the mass on that carrier.

Once the Gaussian and potential estimates on the
common discs are available, the localized logarithm
defines the new correction in its own formation
space. On a fixed support, the accumulated correction
is read from the last contributing stage by restriction
to the spaces of its summands. The construction of
the incoming terms begins with the bare density as
described below.

%
%
The ordinary boundary terms, the lower
boundary terms other than pointed sums and
the earlier stage terms use only their
declared strip slots. Such a slot is read
from a domain containing one of its endpoints.
A term formed at level \(l\) reads stored
fields of levels strictly below \(l\).
Contact with an endpoint at the integration
level can therefore create an attachment
through the pair response while leaving
the stored arguments as they were.

For the direct exterior-only mask consider
the complete blocks of a new collar pivot.
Both endpoint blocks lie outside the old
\(\Omega_q\), whereas an old strip fluctuation
slot at the integrated level has an endpoint
in \(\Omega_q\). An inherited pivot that
reads the changed collar uses its original
ordinary strip component, whose affected
blocks are in the old-wall comparison.
The kept strip reference is integrated by
whole components, and its mixed blocks
across distinct components vanish.
Keep separate strip factors and long-chain
corrections on their original carriers.
A determining value has its own row
coordinate, even when it forms a reference
on a block with a stored fluctuation.

The successful completed strip tests are
products of bond tests. Split each product
between its active and retained coordinates.
On an active bond use the smallest of its
passing thresholds. On the retained bonds
keep the real indicators in the carried
kernel. A failed rim product belongs to
the component of its failure, so it is
carried in the frozen part or the separate
stopped branch. These splits give the real
cutoffs of the passing conditional integral.

In the analytic difference of strip
quadratics, write each chain on its carrier
as
\[
 -\tfrac12 \sum_{b,b'}
   \eta_b^{\mathsf T}(D_Y^- - D_Y^+)_{bb'}\eta_{b'},
 \qquad
 \eta_b=
 \begin{cases}
  a_b,& b\text{ far},\\
  \mu_b a_b,& b\text{ near}.
 \end{cases}
\]
The multiplier belongs to the original
slot in each occurrence, including both
occurrences on the diagonal. This preserves
the carrier and its cut rule. Inherited
analytic functions keep the switches of
their own slots. At \(\mu_b=0\) these
functions and the new interaction forget
\(a_b\). At \(a_b=0\) they forget
\(\mu_b\). The Gaussian reference and
its determinant use the integrated
coordinates themselves. The retained real
kernel carries its own selectors.
The surviving far fields keep their full
original discs, since their occurrences
in analytic functions and in the new
polynomial are separate from the determining
rows that form the response. The switch
identities pass through the connected
expansion once its old analytic readers
are admitted on their formation domains.

\paragraph{The first stage and the physical density.}
%
%
At level zero the array contains only the bare inverse
coupling, and the density is \(e^{-\beta A(U)}dU\).
There are no carried analytic entries at this level.
At a step of level \(t\), form the boundary terms of
\(\mathbf T\) as the sum of the localized normalizer,
the cluster terms assigned to the boundary, the covariance
and weight corrections, and the localized strip shift.
Inside the small-field region their triple argument lies
in the enlarged level-\(t\) core with radii
\((1+\beta_{\rm s})\alpha_{\iota,t-1}\). On the other
supports it lies in the graded space of the completed
label just after \(\mathbf T\). This structure precedes
the structure at stage zero and has the same radius
numbers there. The boundary construction gives holomorphy
in the triple, the far fields and the switches on those
spaces. On the cutoff zone it uses the real stored field;
its Jacobian difference there is the difference from the
zero-field value multiplied by the switch.

%
%
%
%
For a newly external plain far bond let \(R_{\rm out}(b)\)
be its declared carried radius. Choose the ordinary
formation box on these bonds to contain the discs
of radii \(R_{\rm in}(b)=2R_{\rm out}(b)\). This is
a choice of the fixed prefactor \(C_1^{\rm fl}\) of
the current far-field box, made before the weak-coupling
ceiling. Keep the real collar and rim bounds and all
inherited stored-field discs at their prescribed sizes.

Indeed, in the guarded conditional Gaussian, the
external-to-mean operator is a sum of kernels with a
uniform exponential row bound. Newly external far
bonds are at distance at least \(cW_t-O(1)\) from
the integrated bonds, where \(W_t=LM\check R_t\).
Their contribution to the absolute row sum, divided
by the integrated threshold, is bounded by
\[
 C\theta_{\rm c}^{-1}
     \sup_{b\ {\rm far}}\frac{R_{\rm in}(b)}{\mathfrak b_k}
     e^{-c'W_t}.
\]
It tends to zero even on the larger discs. Thus the
far part of the mean leaves the margin supplied by the
real collar and rim. The enlarged-box potential bounds
apply with the chosen fixed prefactor. In particular
the ordinary fluctuation bounds include the cubic error
on the enlarged box and the ratio error bounded by a
constant times \(|1-\tau_k|\sup_b R_{\rm in}(b)^2\).
The same tube bound and the same choice of logarithmic
powers make these errors small for the cluster sum.
Apply the uniform activity majorant on this formation
box before making the strip substitution.

In the plain zone write \(H(b;W)\) for a piece of the
strip shift with carrier \(W\). It includes both endpoint
blocks and reaches the join boundary. On the pair
sections of its formation,
\[
 \sum_W \sup |H(b;W)|
 \le C(\alpha_{0,k}+\alpha_{1,k})e^{-c\check R_t}.
\]
The same bound with a reduced \(c>0\) holds with a fixed
prescribed carrier weight in the sum, by choosing \(M\)
large at that rate. For a simultaneous transfer put
\[
 a_{\rm pre}(b;\mathbf z)
 =\frac1{ig_k}\log\left(
     e^{ig_k a_{\rm post}(b)}
             e^{-i\sum_W z_WH(b;W)}\right).
\]
The physical map is recovered at \(z_W=1\). On fixed
amplitude discs \(|z_W|<R\), \(R>1\), the product
logarithm gives a uniform displacement bound
\(C R g_k^{-1}\sum_W\sup|H(b;W)|\). The exponential
in \(\check R_t\) makes this smaller than
\(R_{\rm in}(b)-R_{\rm out}(b)\). This holds on all
plain far slots together. Thus the transferred cluster
can be evaluated on the full output discs and the
transfer differences can be estimated on the \(z_W\)
discs. On the cutoff zone use the separate real-field
switch, with the smooth cutoff evaluated at the real
field of that zone.

At full switches the strip correction is the change
of the localized exponent under the strip substitution,
together with the logarithm of its Jacobian. To see
this, let \(F_X[S]\) denote the combined term on \(X\)
with the strip substitution turned on at the bonds of
\(S\). Include in \(F_X\) the retained quadratic
term, the cluster and covariance terms that read the
shifted coordinates, and the negative of the
localized Haar log term in the strip formula.
If \(W\) is the set of eligible switches, the sum of
the differences with nonempty index set is
\[
 \sum_{\emptyset\ne S\subset W}\;
 \sum_{T\subset S}(-1)^{|S|-|T|}F_X[T]
 =F_X[W]-F_X[\emptyset].
\]
File a difference on the union of the old domain
and the shift carriers. This changes the index of the
term in the sum and leaves its value there unchanged.
In real coordinates \(a\mapsto a'\) the pullback of
the incoming Haar density is the Haar density at
\(a'\) times the strip Jacobian. Keep the Haar
density at \(a'\) in the reference factors and sum
the log strip Jacobian separately on its carriers.
At the zero stored field the cutoff is
constant near zero, so the strip map is the identity
there and its log Jacobian is zero.

At a physical argument \(U\) the ordinary step
\(t=k+1\) gives the exponent for the completed label
\(\mathcal S^+\).
\[
 \begin{split}
 \mathsf A^+(U)={}&A_{k,\mathcal S}(U)+\mathsf n_k
  +\sum_{z\in\Lambda^0}
       \sum_{z\in X\subset\Lambda}E^{(k+1)}(X,U,z)\\
 &+\sum_{X\subset\Lambda^{\sim-1}}R''^{(k+1)}(X,U)
  +\sum_X\big[E^\partial(X,U)
       +B_{\mathbf T}^{(k+1)}(X;U,\underline{\mathfrak f})\big].
 \end{split}
\]
Indeed on \(X\subset\Lambda\) the bond pointing
weights sum to one. Points outside \(\Lambda^0\)
give \(E^\partial\), and remainder terms outside
\(\Lambda^{\sim-1}\) go to the boundary cluster
family. Outside \(\Lambda\) the normalizer and
cluster terms enter the boundary family directly.
The form corrections and the strip differences just
described are added once. Use the reference convention
with covariance \(\tau_k^{-1}A_{\rm II}^{-1}\) on the
integrated bonds. The change in the gauge quadratic
relative to this convention belongs to the interaction.
The scalar in this identity is
\[
 \mathsf n_k
 =n'(\theta_k)\#(\mathfrak X_{\rm I}\sqcup\mathfrak X_{\rm E})
  +n''(\theta_k)\#(\Omega\cap\mathbb T^{(k+1)})
  +\bigl(n'''-\tfrac32\log\tau_k\bigr)\#\mathfrak X_{\rm I}.
\]
The three counts refer to fluctuation bonds, gauge
blocks and integrated bonds, respectively. They give
the Haar value and field rescaling, the gauge
normalization, and the Gaussian constant including
the ratio correction. These counts depend on the
label. The localized normalizer
accounts for \(-\tfrac12\log\det A_{\rm II}\) and
the log inverse determinants of the interior
eliminated rows. The exterior row factors belong
to the retained operation. If the local normalizer
is centered at the trivial field, put the removed
constant in the scalar part before recentering
the vacuum energy. Rewriting the coupling and vacuum
energy repartitions \(\mathsf A^+\). The tests of
the completed label and the reference measures of its
operations multiply this exponent in the real measure.

%
%
For the ordinary output, evaluate a stored factor in an
adapted representative of the native determining set.
The geometry of an ordinary stored bond of level \(\iota\)
puts its endpoints at tops or at sites in a pillar with
top of level \(\iota+1\). Its parent block is separated
from the next higher layer. Use the endpoint at which
the stored factor is represented. If it is a top, the
residual gauge group fixes it. Otherwise let \(a\) be
the site, \(y(a)\) its native pillar top, and
\(\tau[U](a)\) its contour transport. A gauge that is
\(\tau[U](a)\) at each such stored site and the identity
elsewhere sets those transports to one and fixes the
nodes. This uses the original pillars at all the stored
reads. If \(U\) and \(U^w\) are two such representatives,
\[
 1=\tau[U^w](a)=w(y(a))\,\tau[U](a)\,w(a)^{-1}
 \quad\Longrightarrow\quad w(a)=1 .
\]
The background and its physical current transform together.
Covariance of the factor therefore gives the same value on
both representatives with its stored arguments held fixed.
For a kept operation that reads a root, the same conclusion
uses the determining-node condition on that root. When a
cut or a test tower changes a pillar, use the native
evaluation for the stored factor.

%
For an earlier island \((I,i)\) the root is a level-\(i\)
site in \(I^\sim\subset\Lambda_i^+\subset\Omega_i\), which
is disjoint from \(\Omega_{i+1}\). For \(i=k\) this uses
the separation of the newly appended \(\Omega_{k+1}\)
from the standing island. Thus the root lies in layer
\(i\) of the ordinary output. If preparation acts on the
component containing \(I\), quarantine gives \(i<h_t\)
and the prepared sequence agrees with the original up
through \(h_t\). This leaves both layer membership tests
at \(i\) and \(i+1\) unchanged. Outside the acted
components the sequence was left fixed. In the healed
comparison the surviving spectator roots are outside
the added region \(Z\), where their layers are unchanged.
A level-\(i\) site in layer \(i\) is a determining node.
For a current main operation, the entrance bond has its
outer endpoint in the component and outside the prepared
core, by the clearance of the main box. That endpoint is
a node at the current top level. These checks fix the
roots in the corresponding residual gauge groups.
A crossing fluctuation whose initial point is outside its
box transforms at that point. The exterior starts in an
earlier operation lie on the same single-layer shell
inside \(I^\sim\). In the current operation they lie on
the outer level-\(t\) shell outside \(Z''_t\). They are
fixed nodes by the same layer argument. The residual
action fixes the group-valued drop-out frame variables
directly.

Lower boundary terms
formed in the same manner keep their original graded
spaces. The stage-origin terms have the graded spaces of
their respective stages. The terms from exponentiating an
earlier large-field operation use its output space. They
are the interaction with no chosen component together
with the additional localized remainder. The boundary
parts of the pointed regular sums are formed on their
single-scale cores. In the next ordinary step take the
boundary and stored-operation potentials by subtracting
the two evaluations at a common pinned substitution, with
the stored arguments fixed. This subtraction is analytic on
the graded step domain and preserves the far-field and
switch radii of the entering terms. The regular, remainder
and Wilson potentials use the plain substitution on their
respective domains, since their arguments carry no stored
fields. These are the argument spaces of the formation
rules. When using a reduced history as a donor, take its
member with the space in which that member was formed.

For real evaluation, fix the outgoing and retained data
on an admitted labelled fibre of a stage. Write \(a\)
for the active old strip fields and \(B\) for the
free collar variables. Let \(F_-(a,B)\) denote the
old evaluation of the entries being transferred,
and let \(F_+(a)\) denote their prospective
evaluation at the same strip values. In both
evaluations the switch of each integrated strip field
is one and the spectator fields and switches are
preserved. Each earlier correction has already
closed its own response source against its stage
shift. The fixed exterior exponent is \(F_+(0)\),
with zero in each integrated strip slot in the
original coordinates and its switch removed. The
zero of a translated integration coordinate can
give a different stored value. For the remaining
measure \(d\Gamma\), containing the tests, chart
Jacobians and Gaussian factors, the exact identity is
\[
 e^{F_+(0)}
 \int\exp\bigl\{
    [F_-(a,B)-F_+(a)]+[F_+(a)-F_+(0)]
          \bigr\}\,d\Gamma
  =\int e^{F_-(a,B)}\,d\Gamma .
\]
The first bracket uses the same-strip response
potentials. The second is the active strip dependence
already present at the prospective background. It
stays inside the integral. Use the estimates for the
prospective strip entries at fixed background and
include them once. In the joint strip substitution
above, the empty cut gives \(F_+(0)\) and the full
cut gives \(F_-(a,B)\). Its difference therefore
includes both brackets. The kept strip quadratic has its
separate Gaussian decomposition below.
At this stage put \(q=j+1\). The active variables of \(d\Gamma\) are the level-\(j\)
stored fluctuation on the selected old strip and the
free nonpivots of the newly opened collar. The old
slots meet \(\Omega_q\), have no endpoint in
\(\Lambda_q\), and exclude the original pivots.
Take the whole stored mask of each selected
component of \(\Omega_q\cap Z_q\) contained in
the treated \(X\). Other stored variables are fixed,
including deep core variables whose factors can still
read the moving background. On the active strip
keep the ordinary completion, collar and rim
indicators on their original masks. Where several
apply to the same slot their conjunction tests
at the least of their thresholds. The free sharp
collar ball acts on the new nonpivots. An additional
sharp test in the inherited strip uses its own
transported chart as above.
The prospective regularity guard depends on the fixed
datum and remains outside the collar logarithm. This
includes the prospective crossing tests wherever they
are used. The tests inside the integral are in their
original real slots, with their thresholds set by the
frozen frame. Completing a Gaussian square changes the
reference exponent and its normalization. The domain
of a hard cutoff stays in its original coordinate.

On a common parameter section, write the integral as
\(\int 1_E(a,B)\,F(P,\mathsf y;a,B)\,d\mu(a,B)\),
where \(E\) contains the fixed real cutoffs and \(F\)
contains the chart and exponent factors. At each point
of \(E\), holomorphy of those factors gives holomorphy
of \(F\) in the section parameters. A common integrable
absolute majorant on smaller parameter discs then lets
the Cauchy formula pass through the integral. This
gives the analytic integral on that section, evaluated
on the real data selected by the external guard. Its
nonvanishing and localized logarithm use the connected
expansion and its separate estimates below.

To identify \(d\Gamma\) in the two representations, use
the same reference scale \(\theta_j\) for the Gaussian,
the gauge draw, the thresholds and the fluctuation units.
Put the difference from the running action weight in the
interaction.

%
%
%
%
For the collar form put \(q=j+1\) and use the
full preceding determining sequence with its
closed no-inward row set. The exterior collar is
\(K_q=(\Omega_q^c\setminus Z''_q)\cap X\).
The direct graded rule releases level-\(j\) rows
with an endpoint in \(K_q\) and neither
endpoint in the old \(\Omega_q\).
The added level-\(q\) coarse rows have
the same endpoint rule on their own grid.
Alignment by complete level-\(q\) blocks
puts both blocks of any such row outside
the old \(\Omega_q\). They lie inside the
preceding member of level \(j\). In fact for
\(q\le k_0\) the collar extends at most
five coarse cells into the old exterior
and for \(q=k_0+1\) it extends seven.
The eight-cell quiet gap keeps the blocks
away from the old \(G_j\), also at the first
stage. For \(q>k_0+1\) the separation
between successive prepared faces is at
least one \(M\)-cell of level \(q\).
This gives the same conclusion for the
remade preceding member. The parent
and double-block reaches are \(O(L^q)\)
and fit these margins. Thus the nonlinear
source is tight in the preceding metric.

In a double block meeting the prepared
face, the internal fine links on the
prepared side are held, and the links
in the active block and those crossing
that face are released. The regular inverse
solves the new coarse row for its unique
middle pivot. A row with no endpoint
in \(K_q\) cannot read any released link
in its full double block. A crossing
into the old \(\Omega_q\) belongs instead
to the ordinary strip. Any change of
its old pivot through the exterior block
is handled in the inherited double-block
chart with the old coarse target and
the old coarea factor.
At the overlapping stages use the actual
preceding determining datum. In particular
the level-\(j\) rows completed by the last
stage already have its averaged values.
Among the exterior coarse deltas, use
exactly the newly added rows at this stage.
The exterior stock has neither endpoint
in the old \(\Omega_q\), since the ordinary
step already eliminated rows meeting that
member. This leaves the old-wall coarse
row with one elimination factor and no
additional exterior delta.

The base of the collar form is the prospective real
background at the induced datum. This background
is critical on the preceding fibre since that
fibre at the induced datum is contained in the
prospective fibre. Write \(z\) for the physical
level-\(j\) nonpivot perturbation, \(b(z)\) for the
plain row increment in that sequence, and
\(\beta(z)=b(z)_\natural\) for its reduced
logarithm. The retained plain rows have increment
zero. On the free rows and pivots use the exact
one-step graph
\[
 \mathsf C_V z-\mathsf h_V\widetilde{\mathsf D}_V
                                      (\mathsf C_Vz).
\]
The pivot contribution to the second derivative
of this graph is
\(-2\mathsf h_V C_V^{(2)}(\mathsf C_V\cdot,
                         \mathsf C_V\cdot)\).
The moving level-\(j\) rows of the old set
are interior there. The naturalizer can also
change an old outward row \(c\) of level \(q\)
whose outer block reads the moving fine rows,
even though \(b(c)=0\). For its outer end \(x\)
write \(\ell_x=(1/i)\log\theta_x(b)\) for the
seam factor on the complete internal links
of that block. Then \(\beta(c)\) is
\(\operatorname{Ad}(V(c))\ell_x\) if \(x=c_+\),
and \(-\ell_x\) if \(x=c_-\).
No internal link in that block is a coarse
pivot, since each pivot is between two blocks.
Reuse the same factor at an outer end shared
by several rows, with the orientation of
each row. Thus the pivot curvature belongs
to the released level-\(j\) pivot rows and the
seam curvature to the old level-\(q\) rows
that read the moving block. A newly added
coarse row is used in the fixed-target pivot
equation and has no coordinate in \(\beta\).
The naturalizer reads plain increments once.
It does not propagate a reduced increment
into a further level.

Write \(F(A)\) for the Wilson action at
\(e^{i\xi A}U_0\) in the reference normalization,
and \(\mathcal Q(A)\) for the straightened
row logarithm. Let \(A(z)=\mathbf H(\mathbf 1;\beta(z))\)
be the real response of the preceding problem.
Its node gauge restores exactly the plain datum
and preserves \(F\). Put \(K=\mathcal Q'(0)\).
On the real critical pair \(F'(0)=K^{\mathsf T}\lambda_0\).
For \(P_B=A'(0)[B]\), differentiate
\(\mathcal Q(A(z))=\beta(z)\) twice to obtain
\[
 \begin{split}
 (F\circ A)''(0)[B,C]={}&
 (P_B,F''(0)P_C)_\xi
 -\bigl(\lambda_0,\mathcal Q''(0)[P_B,P_C]\bigr)_{\mathrm{rows}}
 +\bigl(\lambda_0,\beta''(0)[B,C]\bigr)_{\mathrm{rows}} .
 \end{split}
\]
Here the reduced row pairing is the counting
pairing, and the two terms with \(\lambda_0\)
are distinct. The inner constraint curvature
has the minus sign from eliminating \(A''\).
The outer derivative uses the signed pivot
and seam terms of \(\beta''\). These terms
remain even though the first variation on
the fixed-target collar fibre vanishes.

Put \(\nu(z)(c)=\xi_c^{-1}\beta(z)(c)\) and
write \((\cdot,\cdot)_{\mathbf B}\) for the
pairing with row weights \(\xi_c^4\). On the
real response slice the first two terms are
the constrained form with
\(\mathsf G_1^{-1}-Q^{\mathsf T}\mathsf aQ\).
Indeed the gauge projection term vanishes
on the slice and the current correction
of \(\mathsf G_1^{-1}\) is the quadratic
constraint term there. The linearized response
in these units is
\(P_B=\mathsf H_1\nu'(0)[B]\):
it solves the linear constrained variation
at the prescribed row \(\nu'(0)[B]\), with
the same slice condition. Using
\(Q\mathsf H_1=I\) and
\(\mathsf G_1^{-1}\mathsf H_1=Q^{\mathsf T}\mathsf S_1\),
where \(\mathsf S_1=(Q\mathsf G_1Q^{\mathsf T})^{-1}\),
gives the Wilson form
\[
 \bigl(\nu'(0)[B],(\mathsf S_1-\mathsf a)
                           \nu'(0)[C]\bigr)_{\mathbf B}
 +\bigl(\lambda_0,\beta''(0)[B,C]\bigr)_{\mathrm{rows}} .
\]
Multiply by \(\theta_j\) and insert \(g_jB,g_jC\)
for the two physical perturbations. The scalar
factor is \(\theta_jg_j^2=1\).

On the complex section with independent current
use this operator formula for the continuation.
For the outer term take
\(\lambda_0(c)=\xi_c^3(\mathsf H^{\mathsf T}
\mathbf J)(c)\), using the same linear right
inverse \(\mathsf H\) as in the inner current
correction. On the real critical slice
\(J=\hat Q^{\mathsf T}\lambda_0\) and
\(\hat Q=\operatorname{diag}(\xi_c)Q\).
In the row pairing with weights \(\xi_c^4\),
this gives \(\mathsf H^{\mathsf T}J=\xi_c^{-3}\lambda_0\),
and \(\mathsf H_1^{\mathsf T}J\) gives the same
value there. The uncorrected and corrected
right inverses can give different current
contractions off the critical slice. The two
attachments of \(\nu'(0)\) to the native walk
family of \(\mathsf S_1\) use bounded local
stencils at levels \(j,j+1\). The subtraction
\(\mathsf a\) is a local row term. In the last
summand attach the native multiplier walk
\(\mathsf H^{\mathsf T}\mathbf J\) at the row
of \(\beta''\). In each case the walk and the
complete local stencils meet at a common row.
The bounds are absolute in the weighted walk
norm and apply to the actual full-depth form.

%
%
The switched fixed-source derivative used for
a centre response is obtained from the reduced
equations. At zero source the multiplier is
\(\lambda_s^0=\mathsf H(s)^{\mathsf T}\mathbf J\),
in the counting coordinates of that system.
Writing \(\mathsf m(c)=\xi_c\), the counting
right inverse is the weighted right inverse
times \(\mathsf m^{-1}\). The Schur operator
in the multiplier equation is
\(\mathsf m^3\mathsf S\mathsf m^{-1}\) in
place of the weighted \(\mathsf S\).
The diagonal conversions keep each native
walk carrier.
With \(T_\lambda v=(\mathsf C''(0)[v,\cdot])^{\mathsf T}\lambda\),
the derivative of the potential there is
\[
 \partial_{B'}\mathbf H(s;0)
 =\bigl(I-\widetilde{\mathsf G}(s)T_{\lambda_s^0}\bigr)^{-1}
                  \mathsf H_0(s).
\]
Expand the inverse by the small-current bound
of the contraction. Each \(T_\lambda\) uses
the complete local row stencils and the walk
for the multiplier attached at that row. At
a nonzero admitted source, differentiate the
whole contraction and use its Neumann inverse
on the potential and multiplier increments.
The carrier is formed by joining the surviving
walks and local row stencils at shared entries.
At full switches on a real critical datum the
linearized response on the constraint slice
is \(B'\mapsto\mathsf H_1\mathsf m^{-1}B'\),
the response used above in weighted units. This
identification does not replace the derivative
on the switched section.

On the prospective axial blocks the soft density
adds the Gram of the linearized mean contours.
The carried gauge factors cover the old band.
The gauge draws beyond that band cover the
remaining active blocks. Their reference weight
\(\theta_j\) cancels the \(g_j^2\) of the two
scaled derivatives.
Write its derivative in a changed gauge as the
derivative of the relative contour \(T(z)T(0)^{-1}\).
A fixed real gauge conjugates this relative contour
at the block centre and hence preserves the Gram.
In the no-inward sequence a noncentral site of
a complete active prospective block is free:
the level-\(j\) rows from the preceding layer
into that block are the omitted inward rows.
If a density factor also reads a contour at a
retained node in another mask, keep its exact
factor in the interaction unless its variation
on the collar graph is established. Constancy
on a fixed critical orbit alone does not give
its constancy along the varying datum.

%
The flat floor uses the actual plain tangent
\(v_j\) of \(\mathcal M^j U^-\), after the row and
node conversion. Its value at an active nonpivot
is \(g_jB\), and its linear coarse row increment
is zero on the fixed targets. At the flat pair
the multiplier term and the inner constraint
correction vanish. The Wilson Hessian controls
the level-\(j\) curl of \(v_j\). Indeed the
linear one-step flux map on full carriers has
row and column sums \(L^2,L^{-2}\).
Its \(\ell^2\) bound is one in dimension four,
so one may iterate on the full backward union
without a depth loss.

On a complete double block, zero circulations
make a fine one-form a site difference. Zero
mean contours make its site function constant
on each block. A zero coarse row makes the
constants equal. The resulting finite-dimensional
bound controls the fine norm by the squares of
the circulations, mean contours and coarse row.
For a collar slot internal to one block use
the single-block part. For a collar crossing
use both endpoint blocks. At the prepared
inner face the other block lies in the preceding
layer \(j\) and its internal rows are held, so
their tangent is zero. The coarse row across
that face remains fixed. The crossing into
the old \(\Omega_{j+1}\) belongs to the
inherited strip. Thus all the active collar
slots have their full curl carriers and the
required contour terms. Sum the block bound
with bounded overlap and use \(\theta_jg_j^2=1\).
This gives the flat floor uniformly in the
full depth, and coordinate compressions
inherit it by extending the test vector by
zero.

%
For the real comparison use the full form on
the fixed collar geometry. Partition the active
slots by real multipliers \(\chi_\nu\) with
\(\sum_\nu\chi_\nu^2=1\), supported in cubes of
side \(O(RL^{j+1})\) with bounded overlap, for
a large fixed \(R\). For a kernel \(D(b,c)\) the
partition error has matrix entries proportional
to
\[
 D(b,c)\sum_\nu(\chi_\nu(b)-\chi_\nu(c))^2 .
\]
The row and column bounds give an error
\(C(R^{-2}+e^{-cR})\). To see this, use the
first differences of the cutoffs for distances
at most \(RL^{j+1}\), and use the first
escape length in the native metric at larger
distances. Near an active slot the quiet collar
and the preceding prepared separation leave a
cell of clearance for the full buffered patch
in layers \(j,j+1\). Fix the native walk grain
\(\mathfrak M\), take \(R\) large at that grain,
then take \(M\) large enough to put the buffer
and its full stencils inside the clearance.
Use one real gauge on the complete comparison
cube, which meets the preceding layer \(j\).
For \(j\ge1\) the quantitative gauge clause
gives, for \(U^u=e^{i\mathsf v}\),
\[
 L^j|\mathsf v|+L^{2j}|\delta\mathsf v|
 \le C M L_\flat^{2e}\alpha_{0,j}.
\]
The radius tends to zero after the fixed
geometric choices, also with the allowed
overlap factor. At \(j=0\) the patch has
fixed size in bare lattice units. A tree
gauge on this box bounds the links and their
first differences by a fixed constant times
the bare plaquette bound, by filling its
fundamental cycles inside the box.

Conjugate the whole compressed patch by the
same real gauge, acting orthogonally on the
active bond variables. Compare short terms
with their values at the trivial pair by
interpolating its link logarithms and its
independent current from zero. Evaluate
each short term on its full carrier, using
the fixed local admission radii and Cauchy's
estimate. Sum with the native weighted walk
majorant. The multiplier term vanishes at
the trivial pair and has the same small
bound. A term leaving the buffer pays the
first-escape cost even if its matrix indices
coincide. Until escape it takes at least
\(cR/\mathfrak M\) native cell steps, which
gives the exponential tail from the walk
majorant. The floor follows by perturbing
the flat form on each patch and applying
the partition bound to the active nonpivot
vectors. In applying the flat estimate to
a localized vector, lift that vector through
the collar graph at the fixed targets.
Cutting off the previously lifted tangent
need not preserve those constraints.
This comparison also gives a real floor
for the analytic operator at the admitted
small real pairs with independent current,
including the zero-current reference.

A carrier-restricted comparison uses the
complete double blocks of its active slots
and the same gauge and coarse target rows.
If the guard includes the full local stencils
and the surrounding buffer, the terms it
deletes have the exponential first-escape
bound. Their operator norm is below a
fixed fraction of the real floor by the
Schur estimate. This supplies local positivity
for the guarded precision on these carriers.

%
%
%
Use the local resolvent construction for the
inverse root. Take the partition above by
translates on the level-\(j\) grid, with
supports of side \(O(RL^{j+1})\), normalized
by the sum of the squares of the translated
bump. Its normalizing sum is fixed by the grid
and independent of the active mask. Around
each nonempty support take an enlarged active
patch and a further buffer, both of width
of order \(RL^{j+1}\). Write \(D_\nu\) for
the compression to the active slots of that
patch of the sum of the unswitched precision
terms with complete carriers in its buffer.
Use the same rule after any prescribed
compression of the strip. The real floor and the
small complex pair perturbation give the
local resolvent \(G_\nu(x)=(D_\nu+x)^{-1}\)
for \(x\ge0\). The weighted conjugation
estimate bounds its rows and columns by
\(C/(x+\gamma)\) at a reduced rate, with
\(\gamma>0\) fixed. For two active indices in a
patch the distance is at least a fixed multiple of
their physical separation in level-\((j+1)\)
units. A path staying in the two-layer neighborhood
uses at most that spacing, while a path leaving
it first pays a length larger than the diameter
of the patch in those units. Summing the
entrywise resolvent bound over the level-\(j\)
slots gives uniform row and column bounds,
including on the compressed patches.

Set
\[
 G_{\rm p}(x)=\sum_\nu\chi_\nu G_\nu(x)\chi_\nu,
 \qquad (D+x)G_{\rm p}(x)=I+E(x).
\]
With local matrices extended by zero the
error has the two families
\[
 [D,\chi_\nu]G_\nu(x)\chi_\nu,\qquad
 \chi_\nu(D-D_\nu)G_\nu(x)\chi_\nu .
\]
In the first use the absolute walk sum of
\(D\), placing the cutoff difference on
each term. The square-sum of the cutoffs
and the \(\ell^2_\nu\) bound on their
translated differences give \(C/R\) after
the weighted row sum. More explicitly,
\(\sum_\nu|\chi_\nu(d)|
|\chi_\nu(b)-\chi_\nu(c)|\)
is at most \(C\min\{1,|b-c|/(RL^{j+1})\}\)
by Cauchy--Schwarz on the translates.
Use the common entrywise resolvent bound
before this sum. In the second, a discarded term
with both endpoints in the local matrix
must leave the buffer, which costs
\(Ce^{-cR}\). These estimates use the
same complete carriers for \(D\) and its
local inverse. The local inverse itself
has the buffered patch as its carrier.
All patches with active slots fit inside
the base-layer neighborhood of the collar.
Choose \(M\) larger than their fixed
size in one-step units. Then a patch
meets only a bounded number of the
aligned base cells, so its switches
cost \(e^{C\varkappa}\). A small
exponential weight on the physical
patch diameter also costs a fixed
factor, by choosing its rate of order
\(1/R\). With the remaining native
walk decay for \(D\), the absolute
majorant of \(E(x)\) is bounded by
\[
 \frac{C e^{C\varkappa}}{x+\gamma}
                   (R^{-1}+e^{-cR}) .
\]
Use \(R\) sufficiently large for this
bound to be less than \(1/2\) uniformly
in \(x\), then use the indicated floor
on \(M\).

Expand \(G_{\rm p}(x)(I+E(x))^{-1}\)
in these terms. The products join
their carriers at common indices.
Tag each resulting term by the
distinct graded cells meeting the
joined carrier. At binary cuts the
surviving terms use only their
active components. At full switches
the sum is the resolvent of the
physical precision. The integral
with weight \(x^{-1/2}/\pi\) gives
the inverse-root family, since
\((x+\gamma)^{-1}\) is integrable
against this weight. The pair comparison
uses
\[
 G_\nu(p;x)-G_\nu(p_{\rm r};x)=
 -G_\nu(p;x)\bigl(D_\nu(p)-D_\nu(p_{\rm r})\bigr)
                                   G_\nu(p_{\rm r};x).
\]
It bounds the difference by
\(C\epsilon_{\rm pair}/(x+\gamma)^2\)
on the same buffered mask. Telescope
each resolvent product in one factor
at a time. The remaining geometric
sum is bounded by
\(\sum_{n\ge1}n\,2^{-(n-1)}\).
Integration gives the pair bound
\(C\epsilon_{\rm pair}\) for the
switched root on the same carriers.
The guard error uses the separate
first-escape bound on its deleted
carriers.

For the log
use the resolvent representation
with kernel
\((1+x)^{-1}(D-I)(D+x)^{-1}\);
this keeps the cancellation at large
\(x\) before integrating. Each inverse
or log uses the same locally formed
precision and actual retained masks.
The trace for the direct stage term is
taken on the full compressed precision
before the shear, since its determinant
is preserved by that shear. Attach each
diagonal of a log term at its actual
active slot, with the whole carrier of
the term. Subtract its value at the
trivial pair using the same mask and
buffer, and use the pair comparison
of the resolvent products for the
subtracted term. This trace supplies
the direct stage normalizer. The
normalized cut measure uses the
separate determinant of \(A_W\).
This defines a switched root by
its carrier terms. The full-switch
identity is the one used in the
Gaussian marginal formula.

% Earlier Part I, Lemma 4.4 and Lemma 4.24(b):
% gluing on a filled box and the field dependence
% of the subtracted log-walk trace.
On the stage carriers that admit a
common-scale box cover, the subtracted
trace can be summed over stages.
Write \(\ell=\log\theta_i\) at the
top of the waiting window and choose
\(H=A\log\ell\). Split the trace
expansion by carrier length in its
native units. The terms of length
greater than \(H\) spend part of the
rooted decay to gain \(e^{-cH}\).
In the native local-gauge walk
use the sum of its itinerary
length and its number of factors
for this split. The itinerary
exponential pays for the first
quantity and the spare
walk-factor decay pays for
the second. Use the actual
stage anchors in the split.

For a short term, suppose the complete
carrier is in a filled nonwrapping
box of side at most \(CH\) in the
stage units, with an admissible
cover by overlapping gauge boxes
of that scale. The local potential,
first difference and current have
size bounded by
\(\mathfrak a=C\theta_i^{-1/2}\ell^C\)
there, using the running radii in
the waiting window. At small coupling
\(H\mathfrak a\) is within the
gluing radius. Gluing the real
gauges across the cover costs the
side ratio in the potential and its
first difference. Rotate the
independent complex perturbation
and current in the same gauges.
This gives one gauge on the whole
carrier, with the complete pair
bounded by \(C H^C\mathfrak a\)
in its local coefficient norms.
The trace is unchanged by this
change of gauge.

Compare with the trivial pair on
the same buffer and masks. The
local resolvent comparison, or
the telescoping estimate on the
factors of the log walk, bounds
the subtracted diagonal by its
summable majorant times
\(C H^C\mathfrak a\). In the
walk formula the closing loop
stays in the filled box. Its
transport is also within
\(C H^C\mathfrak a\) of the
identity, by multiplying its
links in that gauge. A fixed
power of the walk length is
paid for by the rate reserve.
The same bound holds when a short
term fits in just one of the
original gauge boxes.

%
%
%
The uncompressed stages supply
such a region for their interior
roots on the real support. At
stage level \(q\le k_0\), the old
component of level \(q-1\) lies
inside the hull whose eight-cell
enlargement is \(\Omega_q^c\) in
the class box. The prepared
\(Z''_q\) removes five of those
layers. The roots of the collar
outside \(Z''_q\) therefore have
at least three level-\(q\) cells
to the older core, before the
fixed carrier padding. The lower
re-made regions lie in that
core. In the other direction
use the quiet gap to the unchanged
upper member. If the padded root
neighborhood stays in the class
box, a box of side \(c a_q\)
around it, with fixed additional
padding, uses only the two
neighboring units. The ordinary
current and curvature bounds on
that box give the small real
gauge by the large-box gauge
estimate, with loss bounded by
a fixed power of \(\check R_q\).
The two units cost a fixed factor
of \(L\). A carrier leaving this
neighborhood pays \(e^{-c\check R_q}\)
in the first-escape bound. Thus
these uncompressed roots can be
summed on the physical real
sections. On a complex section
the same gluing applies when its
real factor has the native small
gauge clauses on the padded
cover. The independent current
is then estimated in its own
coefficient norm.

% Earlier Part I, Lemma 4.24(b), Steps 3 and 7.
In the local-gauge walk expansion,
join consecutive transition sites
by a path \(\gamma_a\) inside the
padded localization \(X_a^{\sim}\)
of the factor. Close the path inside
\(X_0\). The localizations used here
are connected unions of a bounded
number of boxes, so each segment
uses a bounded number of
coordinate-monotone paths. The
comparison of the two gauges at an
overlap uses a path in the retained
overlap. This must be checked for
any overlap clipped by the stage
mask. The complete closing loop
\(\gamma\) lies in
\(X_0\cup\bigcup_{a\ge1}X_a^{\sim}\).

Let \(u_a\) be the gauge of the
\(a\)-th factor. After conjugating
the transition product by the root
gauge, its factors have the form
\(u_a(x_a)^{-1}u_a(x_{a+1})\).
Write this as \(W_a U(\gamma_a)\),
where \(W_a\) is a conjugate of
\(U^{u_a}(\gamma_a)^{-1}\). If
the clause bound holds along the
path, and \(s(b)\) is the local
scale of the link, then
\[
 \|W_a-I\|\le C\mathfrak a
                   \sum_{b\in\gamma_a}\frac{\eta}{s(b)} .
\]
Include the closing factor in
this estimate. Telescoping the
products of unitaries bounds the
difference from \(U(\gamma)\)
by the sum of these errors. The
adjoint closing factor differs
from the identity by at most
twice this sum plus
\(2\operatorname{dist}(U(\gamma),Z(G))\).
The box diameter bounds the
physical path length. For the
weighted sum use the smallest
scale met by the path. On a
complete connected cover of
\(O(H)\) graded cells with
adjacent grades differing by
at most one, the ratio of the
largest and smallest scales is
at most \(L^{CH}\). The errors
of the local gauges then cost
only a fixed power of \(\ell\)
times \(\mathfrak a\).

For the walk comparison one can also
use a chain of clause boxes in which
only consecutive boxes intersect.
Use this when the complete closing
path runs through such a chain and
the successive grades differ by at
most one. Adjust the constant gauges
in order along the chain so that the
two gauges match at one site of each
overlap. On an overlap the change of
the relative gauge along a path is
bounded by the sum of the link
deviations in the two local gauges.
The coordinate paths in a clause box
have physical length at most a fixed
multiple of its side. If \(\xi_{\min}\)
and \(\xi_{\max}\) are the smallest and
largest units of the chain, moving
the transition to the point of the
closing path costs at most
\(C\mathfrak a\,\xi_{\max}/\xi_{\min}\).
The paths within the boxes obey the
same bound for their transports.
Multiplying these factors bounds the
closing loop by a fixed power of
\(H\) times
\(\mathfrak a\,\xi_{\max}/\xi_{\min}\).
For a short term through \(O(H)\)
boxes, this size ratio is at most
\(L^{CH}\). The remaining log-walk
factors use their original local
gauge estimates. This proves the
short-term bound also for these
chains, with an extra factor at
most \(\ell^{C_A}\).

%
%
There is a second closing-loop
estimate for the compressed region.
Suppose its complete closing loop
is carried between the class box
and the strictly inset old box
of level \(k_0\). On that full
shell the prepared grades are at
least \(k_0\). With the real
curvature bound on the shell,
each bare plaquette defect is
at most \(C\mathfrak a\eta^2/\xi_{k_0}^2\).
Choose each coordinate path in
a box by moving in one coordinate
at a time. A walk with \(p\)
nonroot factors then has at most
\(C(p+1)\) straight pieces.
The local path errors also obey
the weighted estimate above with
\(s(b)\ge \xi_{k_0}\). A padded
box of diameter \(C_L\mathfrak M\xi_i\)
therefore costs at most
\(C\mathfrak a\mathfrak M\xi_i/\xi_{k_0}\)
per factor.
The rectangular shell loop estimate
bounds the holonomy defect by
the number of straight pieces
times the square of the outer
side in bare units times this
plaquette bound. The two bare
mesh factors cancel. The remaining
scale cost is bounded by
\[
 C(M\check R_i)^2
       (\xi_i/\xi_{k_0})^2
  \le C(M\check R_i)^2(L\check R_{k_0+1})^2 .
\]
Here \(\mathsf s_{k_0+1}=M\)
gives the last ratio. The result
is a fixed power of \(\ell\)
times \(\mathfrak a\), with the
walk-length factors paid as above.
Thus a short loop on this full
shell is also covered by the
subtracted trace estimate, even
when it has no thin-chain cover.
This use of the shell requires
the plaquette bound between the
two boxes, including its filling
plaquettes. It applies on the
real support and to restrictions
of stage pairs on the whole class
domain with this bound on their
real factor. For a supremum on a
locally defined pair space, one
must show that the admitted pair
has such an extension or bound
its closing transport on the
filed domain itself. The local
gauge clauses can hold on an
annular carrier for a flat
connection with noncentral
holonomy.

In the native weighted distance
there is a lower escape bound
for these compressed collars.
At \(q=k_0+1\), the prepared
box removes seven of the eight
outer cells from the old hull
enlargement. To reach the old
core from the active collar
costs at least one of these
cells. At \(q\ge k_0+2\), a
connection from outside
\(Z''_q\) to the old core
crosses the band between
\(Z''_{k_0+2}\) and
\(Z''_{k_0+1}\).
Its width is
\(2M-\mathsf w_{k_0+2}\)
in top units, at least \(cM\).
That band has grade \(k_0+1\).
By the definition of weighted
length, crossing it costs at
least \(cM\xi_i/\xi_{k_0+1}
=cM\check R_{k_0+1}\).
Fixed root padding can be
reserved inside this width.
This proves confinement away
from the old core for short
connections charged their
native weighted length. For a
native walk whose padded factors
are unions of a bounded number
of proper-scale localization
boxes, the same escape also
controls the closing paths. These
are the native localizations of
\cite[Section 3C and Theorem 3.10]{B85c}.
Indeed, with the native grain
\(\mathfrak M\) fixed, the physical
diameter of any one padded factor
is at most \(C_L\mathfrak M\xi_i\).
Choose the stage \(M\) large
enough that this is less than
a fixed fraction of the width
of the protecting band. If a
padded factor meets the old
core, every itinerary through
that factor must penetrate to
the inner part of the band or
cross it. The cost is again at
least \(cM\check R_{k_0+1}\),
including the first compressed
stage where the protecting
width is one cell. In the native
walk distance the itinerary
visits the localizations in
order, using the weighted contour
metric of \cite[(2.45)--(2.46)]{B84b}.
Each possible itinerary
pays this cost, even when the
two matrix indices coincide.
All closing paths stay in the
padded factors, so no additional
escape estimate on their lengths
is needed for this confinement.

%
%
For the directly compressed collar
precision this check can be made
in each factor of the composite
resolvent. Its active slots are on
the outer side of the prepared box,
up to the fixed one-step stencil.
Every local inverse reads only
its buffer, which meets such a
slot and has diameter \(CR\xi_q\).
With \(M\) chosen after \(R\),
the buffer cannot reach the old
core. Between local inverses
each precision factor is compressed
to active slots. If such a factor
reaches the core, its own native
walk or tree must pay the inner
escape from an active endpoint.
Use the native tree distance for
the Schur and multiplier kernels.
Its tree meets the constituent
localizations. The endpoint
condition supplies the root visit
when that localization is omitted
from the tree convention. The
outer slots are within fixed reach
of the first and last localizations.
Thus if an escape occurs on a
multiplier branch, it is still
charged from an active endpoint.
Keep the attached row stencils
in the fixed padding of this
test. The bounds
on the remaining factors sum by the
parametrix majorant, with the count
of choices of an escaping factor
paid by the geometric series.
This argument uses the collar mask
before the shear. A factor with
additional active slots inside
the old core would require a
different estimate.

For a buffered inverse use the
real gauge on its full buffer.
The pair comparison for the
local resolvent replaces its
kernel there by the kernel at the
trivial pair with the same mask.
The vacuum matrix block between
two bond slots is scalar on
colour, by constant gauge
invariance and irreducibility
of the adjoint representation.
When the block is written back
in the original gauge, its scalar
factor multiplies
\(\operatorname{Ad}_{u(b_-)^{-1}u(c_-)}\).
Join \(b_-\) to \(c_-\) by a
coordinate path in the buffer.
At shared slots the segments
join without additional paths.
In a diagonal trace they add
the buffered inverses' segments
to the native precision segments
in the closed path. The matrix
errors are summed by the
telescoping resolvent comparison.
For the remaining closed path
use the shell estimate when the
complete loop and its filling
are admitted there.

%
On the outer side, a
wall-connected chain from a
collar bond to a cell touching
the class boundary has at least
\(2\mathsf N-1\) cells. When
the padded carrier of the
diagonal is filed at such a
bond, reaching the boundary
costs \(Ce^{-c\mathsf N}\) at
the reduced rate, using part of
the filed cell decay. This
controls the outer escape for
these collar roots. For the
other active slots one must
check their own outer margin.

For these terms in the pinned sum the number of
stages and the normalized anchor
count are bounded by a fixed power
\(\ell^{C_*}\). Thus these short
terms and the long tail together
cost at most
\[
 C\ell^{C_*}\bigl(H^C\ell^{C_A}\mathfrak a+e^{-cH}\bigr)
\]
at the reduced rate. Choose \(A\)
after the anchor count. This
bounds their contribution to the
summed direct norm uniformly.
The scalar stage numbers remain
separate from this subtraction.
The compressed stages may have
anchors over the same cells.
The outer \(M\)-cell layer of
\(\Omega_{k_0+1}^c\cap X\) lies
outside each \(Z''_q\) in the
compressed range and is inside
\(\Omega_q^c\) by nesting. Thus
even a fixed pin can meet all
these stages, and the absolute
sum needs a gain beyond the
count of its anchor layers.
For other short carriers crossing
the graded interfaces we must
still control the closing transport
on the complete admitted reads.
A chain that closes around a hole
can require filling plaquettes
beyond those reads.

%
%
%
For the strip part use the kept
chains of the external Schur form
with their original keep rule.
The source supplies their positive
floor after the chain selection.
The original Schur complement
already eliminated the ordinary
inner integration variables. File
each kept chain with its original
strip component, including its
one-layer carrier margin. This
margin includes the full read of
the middle resolvent in a Schur
chain. In the class box \(X\),
let \(H_q\) be the level-\(q\)
cell hull of \(Z_{q-1}\cap X\).
The complement \(\Omega_q^c\cap X\)
is its eight-layer enlargement.
The strip component lies in
\(\Omega_q\cap Z_q\cap X\), so
its one-layer enlargement is
still separated from \(H_q\)
by at least seven level-\(q\)
cell widths. At the first
compressed stage this excludes
the old hull. At later stages
the first prepared box
\(Z''_{k_0+1}\cap X\) is already
contained in \(Z_{q-1}\cap X\).
A new buffered inverse at an
active strip slot has fixed
reach in level-\(q\) units
and stays in this clearance.
Thus the kept strip precision
and its buffered inverse avoid
the old core on their complete
physical reads.

%
%
At a nonterminal ordinary strip
root the same one-scale argument
uses the two quiet gaps. Let
\(a_q=M\check R_q\xi_q\).
The root lies in \(\Omega_q\cap
Z_q\cap X\), up to the fixed
endpoint reach. Its distance
to \(Z''_q\) is at least
\(5a_q\) before compression and
\(7a_q\) at the first compressed
stage, before subtracting that
reach. At a later compressed
stage the old inset is contained
in \(Z_{q-1}\cap X\), giving the
\(8a_q\) gap of the original
class geometry. In the other
direction \(Z_q\cap X\) is at
least \(8a_{q+1}\) from the
unchanged \(\Omega_{q+1}\).
These roots are also away from
the class wall, by the same
separation at the top level.
The native towers of the direct
strip operator stop at \(q\).
Thus the native itinerary and
bounded padding put a short
trace in a box of side
\(CH\xi_q\), inside the single
grade \(q\) of the prospective
physical sequence. The real
profile and the box gauge give
the same small-field estimate.
Use the prospective strip
quadratic for this direct
term. Its difference from the
old quadratic is carried in
the interaction.

For the terminal strip, this also
gives a one-scale box for the
short trace on the physical pair.
At \(q=i\) the active roots are
in the original \(\Omega_i\cap X\).
The prepared inset \(Z''_i\cap X\)
is contained in \(Z_{i-1}\cap X\),
which is separated from these roots
by eight top cell widths.
The exterior top coarse cells
of \(\mathcal P_i\) touching \(X\)
belong to the clean region and
hence to \(\Omega_i\). In the
global prepared sequence the
first exterior layer stays
unchanged. Thus, on either side of
the wall of \(X\), a neighborhood
of width \(cM\check R_i\xi_i\)
about each strip root is entirely
in the global grade \(i\).
For a short word of length at
most \(H\), the native itinerary,
its bounded padding per factor,
and any local inverse buffers
put the entire trace in a box
of side \(CH\xi_i\).
The real profile controls all
plaquettes of that box. The box
gauge estimate then bounds the
short trace by its majorant times
a fixed power of \(H\) times
\(\mathfrak a\), as above.
Use this on the global physical
pair, retaining the original
operator masks and native weights.
If the strip mask is compressed
further at an admitted guarded
binary switch, the
comparison still uses the fixed
physical pair and the surviving
carriers. The box filling estimates
their closing transport. Its
plaquettes are used in the bound
and the coefficient continues to
read its own carrier.

The chain functions and their
read sets are fixed at formation.
Evaluate them on the supplied
old or prospective pair in the
common stored coordinates. An
endpoint transport needed for
this evaluation fits in one
additional layer. At the stage
that consumes a crossing slot,
use its prescribed endpoint
conversion, since the preparation
can change its determining row.
A response used to construct a
prospective pair contributes its
own carriers when the coefficient
is composed with that response.
The original keep rule applies
to the coefficient's reads on
the supplied pair.

On the native stage grid these margins
remain below the next higher member,
by the quiet separation of the
carried \(Z_q\). Rejected chains
keep their old carriers and their
separate boundary terms.
Compress this kept precision to
the strip slots integrated in the
stage, carrying mixed terms with
fixed exterior slots in the interaction.
For a further guarded subset \(W\)
take the compression
\[
 D_W(s)=P_W
  \Bigl(\sum_{\varpi\ {\rm kept\ at\ formation}}
          s^{\operatorname{par}(\operatorname{tags}(\varpi))}
          D_\varpi\Bigr)P_W .
\]
Here \(\operatorname{par}\) takes each
original cell to its parent on the
current grid. This compression
retains the keep decision made for
the original Schur form.
Run the local resolvent construction
on this compression at the ordinary
level-\((j+1)\) grain, then map the
old tags to the common graded cells.
The root is formed after compression.
A row restriction of the ambient root
factors the marginal covariance. A
square compression of the ambient root
need not factor either the marginal or
the conditional covariance. Use the
root of the compressed precision here
and take the
direct sum of the two roots of the
reference, then apply the shear transport.

Keep the inherited conditional strip form at its prescribed base
and put its complete difference at the two backgrounds
in the interaction once. In common stored coordinates
let \(D_-\) and \(D_+\) be the old and prospective
kept precisions, and write \(a=(a_{\rm I},a_{\rm E})\)
for the active and retained slots. The exact split is
\[
 \begin{split}
 -\tfrac12\langle a,D_-a\rangle={}&
 -\tfrac12\langle a_{\rm I},D_{+,{\rm II}}a_{\rm I}\rangle
 -\tfrac12\langle a_{\rm E},D_{+,{\rm EE}}a_{\rm E}\rangle\\
 &-\langle a_{\rm I},D_{+,{\rm IE}}a_{\rm E}\rangle
 -\tfrac12\langle a,(D_--D_+)a\rangle .
 \end{split}
\]
Use the first term for the integrated reference
Gaussian and keep the second in the prospective
retained factor. The last two terms go in the
interaction. The rejected chains and coupling-ratio
terms are boundary entries. Evaluate their old and
prospective values with the other boundary entries.
In the physical stage the active strip is a union
of whole kept components. The kept prospective
mixed block against unintegrated components then
vanishes. A subsequent Gaussian cut uses the
marginal of the integrated reference. Chains
coupling distinct original components are included
among the rejected chains. Together with the kept
form and the ratio term they give the original
conditional weight at each background. Keep the Haar density,
the derivative of the linearizing correction, the row
elimination factor, and the scalar vacuum contribution in
the assembled weight. In particular the row elimination
factor is in the collar integral and the vacuum term is
in its prefactor. The fixed-ball identification requires
these identities and the same retained strip tests on the
same active bond sets.

This identity allows the graded expansion to form its own
localized logarithm. For equality of the physical kernels
one starts with the same incoming density on each labelled
fibre and verifies the same measure transformation there.
%
%
%
The incoming kept strip in the real kernel
has its unnormalized quadratic and its
inherited row factors. When the complete
reference is written as a probability Gaussian,
its full logarithmic prefactor in Lebesgue
coordinates is
\[
 n_{\rm free}\nu_{\mathfrak g}
       -\tfrac12\operatorname{Tr}\log D
       +\log z_{\rm col}.
\]
Here \(n_{\rm free}\) counts the free Lie-algebra
slots of both integrals and \(\nu_{\mathfrak g}
=(\dim\mathfrak g/2)\log(2\pi)\) for ordinary
Lebesgue measure on a slot. The factor
\(z_{\rm col}\) consists of the new collar
pivot factors in their declared pairing
convention. Keep the old strip row factors
in their original places. The unnormalized
quadratic becomes the normalized Gaussian
with this one restoration of its mass.
If a previously normalized quadratic is
rewritten with a new precision, the
determinant factor is the ratio of the
two masses. The same rule applies when
a coupling difference is transferred
between the precision and the potential.

For the Haar delta the derivative at the
reference pivot is \(L\mathsf m_c\). The full
coarea away from that point uses the nonlinear
chart factors. In the convention
that files \(|\det\mathsf m_c|^{-1}\) in
\(z_{\rm col}\), the factor \(L^{-\dim\mathfrak g}\)
per row is in the scalar of the same operation.
The flat pivot is
\(\mathsf m_c(\mathbf1)=L^{-4}I\),
so their product is \(L^{3\dim\mathfrak g}\)
per new row. This is the same factor as
the inverse determinant of \(L\mathsf m_c\)
in Haar frames. Count the scalar once at
the row's owner. The normalized cut
determinant in \(\mathcal J\) is separate
from this full prefactor.

The shear preserves the full determinant,
and the affine mean translation has no
extra scalar after completing the homogeneous
reference square. In the direct logarithm
subtract the value at the trivial pair from
the field-dependent determinant and new
pivot factors. File their trivial values
with the fixed Gaussian number of that
stage and the collar constant \(E_0^{(j)}\).

%
The new pivot logarithm has a
local small-field bound. At the
reference row centre its Haar
differential \(L\mathsf m_c\) is
the pivot coefficient in the
relative logarithmic frame.
Conjugations at the two endpoint
frames have determinant one.
In the small gauge on the complete
double block write this coefficient
as \(L^{-3}I+E_c\).
The pivot estimate gives
\(\|E_c\|\le CL^2\varepsilon\)
with \(\varepsilon\) the local
chart size. The subtraction at
the trivial field is therefore
\(-\operatorname{Tr}\log(I+L^3E_c)\),
whose modulus is at most
\(C_L\varepsilon\) on the small
branch. This estimate uses only
the complete double block. The
running physical row bound makes
\(\varepsilon\le C\theta_i^{-1/2}\ell^C\)
on the waiting window, so the
polynomial count of stages and
anchors gives a uniform summed
bound for these local pivot terms.
The same bound holds on admitted
complex sections with this row
bound. The inherited strip coarea
keeps its original carrier.

The remaining direct terms and the cluster
logarithm are functions on their own
formed carriers. A term meeting several
components occurs once, through the
maximum of their executed stage counts.
The constants count the slots and rows
of each executed component once.

In the convention with the Wilson-weight
adjustment in the stage term, the expression
also contains \(A(x^{(n+1)}-x^{(n)},\cdot)\).
For a component stopped after \(n\)
integrations all earlier terms are
reevaluated at the common background \(U\),
with integrated stored slots set to zero.
The compensations then combine with the
current Wilson exponent as
\[
 -A(x^{(n)},U)+\sum_{\nu=0}^{n-1}
       A(x^{(\nu+1)}-x^{(\nu)},U)=-A(x^{(0)},U).
\]
This is the Wilson exponent that the stopped component
keeps when its stage logarithms are filed separately. The
scalar parts are carried with the component factors.
%
%
%
At each prefix apply the real operations to the joint
measure in the variables retained by that prefix. Use
the same labelled tests and chart in the physical
integral and in its graded representation. In the
local sharp convention, the sharp collar test is made
at the local centre before the shift. Its failure keeps
the actual prefix with its failure indicator and the
still untransformed variables of that stage. If the
global sharp convention is used, the failure branch
instead uses the transported test and its transported
density. In either case the failed stage contributes no
collar integral or stage logarithm. These two choices
need separate labelled identifications wherever the
tests split the fibre differently.

For a fixed output standing part \(\mathfrak o'\), let
\(\mathrm{Str}(\mathfrak o')\) consist of the real strings
whose output has that standing part. For such a string
write \(K_\varsigma\) for the product of its real kernels
and \(\pi_{\mathfrak o'}\) for the marginal onto the
retained variables. The measure of the standing part is
\[
 \sum_{\varsigma\in\mathrm{Str}(\mathfrak o')}
    (\pi_{\mathfrak o'})_*
       \bigl(\hat\tau_{\mathfrak o(\varsigma)}K_\varsigma\bigr),
\]
where \(\hat\tau_{\mathfrak o(\varsigma)}\) is the incoming
measure already summed over its standing part. Once the
stage identity holds for each applied real kernel, it
holds under continuation by a later kernel and under
this marginal and sum.

%
%
The continuation uses the completed determining rows,
the field on them and the stored data of the components
that still stand. On a healed \(X\), the lower rows have
been replaced by top rows, including the top bonds
meeting \(X\). The restored field is read on these
rows just as it is elsewhere. The integration over the
healed history is the identity on all remaining
coordinates. At a fixed value of these coordinates
the same geometric rules give the next tests and
operations. The class window cannot read a birth or
mark erased by an earlier healing. In particular, the
next normalized reference draw is computed from the
current exterior field, the prescribed coupling and
frame, the tree slice and the local cutoffs and Wilson
weight. Its denominator is the integral of that same
local weight on the slice. No term in the incoming
healed bracket enters this denominator. Repeating this
argument gives the same positive continuation on the
common retained space. The healed terms remain part
of the incoming density to which this continuation is
applied. This permits summation over the healed
histories before continuation and exponentiation of
the resulting bracket. Keep the
measure, carrier and internal coordinates of an island
that still stands. An absorbed island is removed from
the list of standing islands at filing, after its real
operation has been included. Equality here is an
identity on the retained variables of the construction,
with any physical insertions carried by the corresponding
kernels. The coupling coefficient still comes from the
whole-lattice small-field construction of \(\mathbf T\),
with its own Taylor expansion.
The analytic bounds still use the declared arguments of
every graded term, including the evaluations at \(F_+\).

%
%
The domains are part of the data carried by the
induction. At a new level a boundary entry of
\(\mathbf T\) uses the space of the completed ordinary
label. A stage entry uses the spaces of its contributing
executed stages, restricted to its formed domain. A
main-step entry uses the output space. The regular
remainder and pointed entries use their single-scale
cores. For the exceptional top boundary read of a
pointed entry the restriction lies in the core with
radii \(\bar a_\iota\); the boundary function is given
on that core. This read uses its own core radius.
Kept forms and elimination logs use the regular pair
on their birth carriers. A standing island keeps its
birth core, its measure and its inner factors. When it
stands inside a new drop-out, restrict its outer
arguments to the admitted smaller space.

For an older entry the output structure is restricted
to its original level, carrier and stored argument
list. Each earlier tower address is still a test of the
output structure, with the right-hand radius for that
read. The inclusion into the entry's space uses the
radius comparison on that address, including on the
interior of a fixed-real-factor section. In a re-made
member the retained field is evaluated by its actual
row map, with its original bond kind where carried.
This re-reads the same function. The boundary and stage
entries are formed on the graded domains from their
first occurrence.

To verify the real induction on these domains, fix the
outer list and the retained real variables of the
outcome. The ordinary sum rule gives the incoming
exponent of the first stage. At each executed stage the
full-switch evaluation, with the same cutoffs and row
coordinates, gives the physical collar integral. At
the main step use the real restoration and class split
on that retained measure. Composition and the marginal
on the output variables give the sum over the real
strings with the specified standing part. In a healed
region the equality is for this sum. For a standing
island the same inner measure and factors are compared
at equal outer reads. These identities are evaluations
on the real history kernels. The analytic estimates
for the new entries use their formation spaces. The
prospective regularity guard is the bare plaquette
test of the real background, with the crossing
clauses of the chosen rule. Its split precedes the
evaluation of the analytic collar logarithm.

If instead one imports the ordinary boundary total
formed on the whole-layer tower domain, its first use
in a collar is already at the first active stage. A
stage correction imported with that domain first enters
the earlier-stage sum at the second active stage. Each
such read would require an extension of that entry
or bounds on its original tower stencils. An evaluation
using that broader condition
on a stencil outside the comparison window requires a
bound for
\[
 L^{2\ell+\min\{\mathsf p,\ell\}}
 \left|J^0_{\mathsf p,Z}(W^u)(b)
             -J^0_{\mathsf p,Z}(\mathbf U)(b)\right|
\]
by the available tower margin at nominal layer \(\ell\)
and depth \(\mathsf p\), uniformly on the complex sections
and strip discs in use. Here \(J^0_{\mathsf p,Z}\) is the
tower current in bare units. The exponent follows from
the conversion \(J^{[\mathsf p]}=L^{3\mathsf p}J^0\)
and the layer factor
\(L^{\ell-\min\{\mathsf p,\ell\}}(L^{\ell-\mathsf p})^{-3}\)
in the allowed current. The broader condition is a
requirement on the domains of those particular terms.
To see the importance of the domain for the current bound,
at a fringe where the full stencil is fixed by bare
determining rows, a one-link perturbation of size
\(c\alpha_{0,\ell}L^{-2\ell}\) at the identity has a bare
current of the same order. For \(\mathsf p\ge\ell\) its
normalized current is of order
\(c\alpha_{0,\ell}L^\ell\), even if the independent
current coordinate is zero. Using the tower's own local
layer \(a\) for the current estimate supplies the weight
\(L^{3a}\); converting it to the required one costs
\(L^{2\ell+\min\{\mathsf p,\ell\}-3a}\).
There is a direct comparison for some such reads. Suppose
the weighted value and first difference of the levelled
potential change are at most \(A e^{-\gamma d^-(b,\Sigma)}\),
where \(\Sigma\) is the source collar and \(d^-\) is the
ambient weighted distance. Use the tower itself as the
ambient sequence for the tower comparison at the common
real factor. If its distance dominates \(d^-\) on the
carriers, the comparison gives
\[
 C A L^{2\ell+\min\{\mathsf p,\ell\}-3a}
            e^{-\min\{\gamma,\delta_{\rm T}\}d^-(x,\Sigma)}.
\]
Indeed, taking the finer scales only decreases the weighted
input norm. In the maximum of the tower estimate the two
distance factors combine by the triangle inequality. This
bound fits the margin when the scale loss is covered by
the radius budget or by clearance from the source. At a
bare fringe \(a=0\) the loss can grow with the depth of the
ambient sequence, even at a fixed distance from the source
in ambient units. The response estimate for second differences
also allows a scale loss in its own ambient layer. Its
zero-layer estimate applies at zero in that sequence.
The auxiliary tower can have layer zero at a site in a
coarse layer of the response sequence. On such a bare
fringe the recomputed current is the Wilson current of the
shifted links on its full stencil. The physical endpoint
identity for this current uses the equations at the full
parameter values. Those equations cannot be substituted
term by term into the walk series with independent switch
parameters. At such a read the argument requires an
additional estimate on the decoupled response, or an
analytic construction on the settled domains.

\paragraph{The crossing test at the main integration.}
Use the exterior, extendibility, fluctuation and last-shell
tests of the main step, together with its regular core and seam
data. On this support the joined datum \(V''\) is
\(\varepsilon_\star\)-regular, so the regular critical orbit
\(U''_{m,X}\) exists by the real main-step admission
assertion \cite[(1.89)--(1.98)]{B89a}, which uses these inputs
before the crossing test is evaluated.

%
%
Put \(h=m_\flat\). The frozen core is a nonempty component
of \(Z_h\), while the input large region at level zero is
empty. Thus \(h\ge1\). Work for the moment in the original
step-\(h\) frame. In level-\((h-1)\) units a coarse cell of
this step has side \(LM\check R_h\). The side of the test
cube enlarged by four test margins, with both endpoint
blocks allowed for in its averaging reads, is at most
\[
 (L^2+8L)M_2\check R_h+4L<LM\check R_h.
\]
Here \(L^2M_2\le M\) and \(L\ge13\). If the enlargement
straddles \(\partial\Omega_h\), its entire read lies within
one coarse layer on either side. The pre-operation region
\(Z_h\) contains \((\Omega_h^c)^{\sim2}\), so the connected
read lies in the component that meets the exterior at the seam.
For the seam of the frozen component this component survives
the operation of level \(h\). A healed component would have
disappeared, and the class rule excludes a drop-out or a
second-class mark at \(h\) itself. Operations on disjoint
components preserve its entries. The preceding input region
is at least eight coarse layers from \(\Omega_h\) when \(h>1\).
For \(h=1\) the initial rule gives the same test buffers with
the empty input region and the five-layer margin at \(P'_0\).

Read the seam field of level \(h-1\) from the retained one-step
factor of \(\mathbf T_h\) on this core. In the exterior of
\(\Omega_h\) this factor integrates the level-\((h-1)\) links
with their averaging delta. In the interior it stores the
nonpivot fluctuations, with the pivots fixed by the averaging
equation. On the strip read the links through its substitution.
The resulting field \(\mathsf v\) satisfies
\(\mathcal M(\mathsf v)=V_h\) on the coarse reads. The lower
determining rows are unchanged by preparation, and the
level-\(h\) rows on the frozen core read the original \(V_h\).
Consequently the seam fields extend these rows in the original
frame of the retained operation.

The two-layer buffer around the approximate failures and the
three-layer gauge buffer give the passing approximate test
and the passing coarse plaquette test on the test bases in
\(\Omega_h^{\sim2}\). The ordinary step estimate on
\(\Omega_h^\sim\) therefore gives
\(|\partial\mathsf v-1|<\varepsilon_{h-1}+\varepsilon_h\)
and the contour bound \(CL\delta_{h-1}<\varepsilon_{h-1}\)
on the seam reads. These are statements about the retained
fields in their original frame. Under a simultaneous change
of gauge, lift the level-\(h\) transformation constantly on
its blocks. The average transforms at the block nodes, and
the contour and the relative link against the axial center
are conjugated. The latter uses uniqueness of the center on
the tested bonds in that axial gauge.

%
This gives an identity of the retained factors on an
ordinary seam carrier whose moved sites belong to the
joint integration. The complete lower double blocks lie
on the \(\Omega_{h-1}\) side of its outer boundary,
since the old large region \(Z_{h-1}\) is separated from
\(\partial\Omega_h\) by eight coarse layers. The
complete level-\(h\) blocks for the lift fit in that
margin. There are no determining rows below level
\(h-1\) on these blocks. The level-\((h-1)\) rows
are read through the retained Haar factors, the solved
pivots or the strip inverse, as above.

Fix the exterior of the main step. Its reference
transformation is then fixed. On a union of these
complete blocks on which every moved lower coordinate
is carried by the joint law, lift its restriction
constantly on each block. Change variables in the
retained Haar integral and conjugate the free
logarithmic coordinates at their initial sites. The
Haar and Lebesgue reference measures are preserved.
At a row \(c\), the argument of the old delta becomes
\[
 \mathcal M(\mathsf v^u)(c)(V_h^u(c))^{-1}
 =u(c_-)\,[\mathcal M(\mathsf v)(c)V_h(c)^{-1}]\,u(c_-)^{-1}.
\]
The delta is unchanged. Evaluate the old indicators
and weights on the inverse reads of this same change,
including the original strip and reconstructed pivots.
Thus their joint integral gives the transported factor
at the new \(V_h\) with the guards still in place.
This applies where the full changed carrier is within
the joint law and has no moved genuinely fixed lower
endpoint or older mark with independent stored data.

%
There is also a comparison of the density factors
that only reads the level-\(h\) field. In the first
prepared box the inner cut leaves one coarse
\((h+1)\)-layer about the hull of the old core.
For an old density cube in \(\Omega_h\cap Z_h\),
the fourfold enlargement fits inside that layer
because \(4M_2\check R_h L^h+L^h<
M\check R_{h+1}L^{h+1}\), including its bond
reads.
The last-shell artificial datum keeps the
level-\(h\) links on this whole test box, in its
compensating site gauge. Lifting that gauge to the
localized tower identifies the two backgrounds up
to gauge and gives the same tested plaquette norms.
Thus the passing half-threshold implies the old
density factor on these cubes. This implication
uses only the literal level-\(h\) test reads.
In the full last-shell split it also applies
to a tested box crossing \(\partial\Omega_h\):
inside the old core the level-\(h\) field is
still a target of the retained graph there.
When the split is restricted to the inset
whose whole test box lies in \(\Omega_h\), keep
the old density factors at the other cubes.
The small-energy estimate for the inset
continues to use that full-box condition.

One can keep the old frame on the ordinary seam
without changing the reads of the two outer tests.
The seam and its full level-\(h\) test halo lie
strictly inside the inner cut. All genuine rows
below level \(h\) of the frozen core have their
endpoints in \(Z_h\), by the no-inward convention.
Choose a site gauge equal to one on those endpoints
and the seam halo, and equal to the mixed
compensator off an interior neighborhood of them.
Use this only to change coordinates of the retained
datum. Every endpoint of a fluctuation row lies
outside the change, so the mixed reference and
the relative fluctuation reads on \(\mathbb B_0\)
keep their values. On the artificial last-shell
field the additional gauge is supported strictly
inside the kept side, with a bond margin to the
cut. The bonds set to one therefore remain one,
and the whole artificial datum changes by a
site gauge. The last-shell indicator is unchanged.
With this gauge the old lower rows and the seam
averages are read in their original frame.

%
For the guarded integral, let \(\tau\) be the
mixed compensator and \(u\) the gauge just chosen,
both fixed by the exterior and the frame. Write
\(Y=V_h^u\) for the new target. The old factor
reads \(V_h=Y^{u^{-1}}\), and a factor written
in the mixed frame reads \(V_h^\tau=Y^{\tau u^{-1}}\).
The change from the mixed to the reset target
uses only Haar links in the interior of the
inner cut. In the old core these are the
\(dV_h\) targets of its retained operation.
Outside the old core and inside this cut the
level-\(h\) links belong to the prepared level-\(h\)
member, since they are in \(\Omega_h\) and clear
of \(\Omega''_{h+1}\). They lie in the second
Haar integral before the split. At the common
face use the original ownership of the links,
so the two stocks remain disjoint.

Keep the lower integration variables \(\xi\) in their
original charts. A native factor \(F(V_h,\xi)\),
including its stored gauge factor, its contour,
its reconstructed pivot and its indicator, is
then read as \(F(Y^{u^{-1}},\xi)\).
The averaging delta in that factor is
\[
 \delta\big(\mathcal M(\mathsf v)(Y^{u^{-1}})^{-1}\big)
       =\delta\big((\mathcal M(\mathsf v))^u Y^{-1}\big)
\]
row by row, by conjugation at the initial endpoint.
This transports the output graph and leaves the
lower rows and gauge factors in their original
variables. Haar measure on the changed targets is
preserved at each fixed exterior. Factors that
read the mixed target keep the argument
\(Y^{\tau u^{-1}}\), with the same rule for all their
other reads. Thus this change is made with the
guards still in the integral.

In the prepared datum use the transformed value
at both ends of each retained row. The old lower
rows are fixed by \(u\) at both endpoints. The
level-\(h\) rows in the inner cut are the targets
just described, and all fluctuation rows and the
outer join retain their mixed-frame values.
Consequently the prepared datum is the original
prepared datum transformed by \(u\) on its full
row set. Gauge covariance of the averages and
of the regular critical orbit identifies the
prepared guard with its original read, including
its domain of definition. The change of variables
does not use a passing value of that guard.

At an earlier face \(\partial Z''_{s+1}\), with
\(m_\flat\le s\le m-2\), the adjoining shell scales are
\(s,s+1\). The first shell is \(Z''_{m_\flat+1}\cap\Omega_{m_\flat}\).
Consider a crossing whose corners
are all in the inset \(\Omega_{m_\flat}\cap X^\natural\) of
the completed main-step estimate. The separation of successive
prepared faces is more than the bare plaquette diameter, so
every corner has one of these two profiles. For \(s\le m-2\)
the completed coefficient on shell \(s\) satisfies
\[
 1-\beta_{\mathrm s}(1-2^{-(m-s+1)})
       \le 1-7\beta_{\mathrm s}/8.
\]
On shell \(s+1\), the unmultiplied threshold is at most
\((L_\flat/L)^2(\varepsilon_{s+1}/\varepsilon_s)\) times
\(F_s\varepsilon_s L^{-2s}\). The ratio of consecutive
radii is \(1+o(1)\), and \(L_\flat<L/2\), so this too lies
below \((1-3\beta_{\mathrm s}/4)F_s\varepsilon_sL^{-2s}\)
for small coupling. The completed main-step estimate on the
inset bounds the curvature by the maximum of these corner
thresholds, including for astride plaquettes. It therefore
implies the added crossing test there. This uses the completed
corner estimate. At the prefinal passing stage one still uses
the separate prospective test and the real collar response.

We use the frame of the main fluctuation test and put
\(h=m_\flat\), \(Y=Z''_m(X)\), and
\[
 I_h=\Lambda^{\mathrm{pr}}\cap
                  ((\Omega''_{h+1})^{\sim2})^c .
\]
The fluctuation bonds \(\mathbb B_0\) extend from the
outer box to this inner cutoff. They include the bonds
crossing \(\partial\Lambda^{\mathrm{pr}}\) and exclude
those crossing the boundary of \(I_h\) towards its interior.

We first identify the reads of the tested holonomies of \(V''\).
A determining row of level \(i<m\) in the prepared sequence
has both ends in \(Y\), since neither end can lie in the
prepared member of level \(m\). A row of level \(i\ge h+2\)
has an end outside \(Z''_i(X)\). Its length is \(L^i\) in
bare units, less than the separation of \(Z''_i(X)\) from
\(Z''_{i-1}(X)\). For the early separations this follows
from \(7a_{i-1}>L^i\), and for the compressed ones from
the gap of at least one \(M\)-cube of level \(i\).
Such a row cannot reach \(Z''_{h+1}(X)\).
A row of level \(h+1\) has one end outside
\(Z''_{h+1}(X)\). Its other end can enter that box by
at most \(L^{h+1}\), so it stays on the fluctuation side
of the two-cell enlargement. Thus a genuine row missed
at the inner cutoff has level at most \(h\).
A top row beyond \(\Lambda^{\mathrm{pr}}\) instead has
its exterior value \(V_m=V_\Lambda\).

At a tested plaquette of level \(q\), the reads are
the determining rows on that plaquette and, for each bond
completed by one-step averaging, the two full blocks of
level-\((q-1)\) sites at its ends. The latter blocks lie wholly in the
preceding layer in the convention excluding inward bonds.
Consequently a test that reads an inner row missed by
\(\mathbb B_0\) has \(q\le h+1\). All its corners lie
in \(Z''_{h+1}(X)\) or within \(C L^{h+1}\) of it.
The shell between \(Z''_{h+1}(X)\) and \(Z''_{h+2}(X)\)
has width at least \(7a_{h+1}\) and spacing \(L^{h+1}\).
Every path from this read to a crossing plaquette at
\(\partial Y\) must cross that shell. After the footprint
allowance its weighted length is at least
\(cM\check R_m\). These holonomies can therefore be
estimated by \(\varepsilon_\star\) with the weight
\(e^{-c\delta_{\mathrm c}M\check R_m}\) in the local
regularity estimate.

For the remaining interior holonomies compare with the
averages of the basic background \(U_0\) in its mixed
axial--Landau gauge. The restriction of this datum to
\(\mathbb B_0\) is the actual center \(V_0\) of the
main fluctuation test.
On \(\mathbb B_0\) the relative bond logarithms have norm
less than \(\delta'_m\), and on the exterior top rows
the relative bond equals the identity. All lower-level
plaquette carriers lie within \(C L^m\) of \(Y\).
The same holds for a top plaquette using a completed bond,
because that bond reads the layer \(m-1\).
The outer \(M\)-layer of \(\Lambda^{\mathrm{pr}}\) and the
clearance of this box from the joining collar put these
carriers in the top layer of the tower for \(U_0\).
By \cref{prop:large-field-2-4} the top datum \(V_\Lambda\)
has defect at most \(C(M)\varepsilon_m\), using the
unchanged exterior and its outer tests off the box.
The lift to the tower preserves this bound on the tested
holonomies. Regularity on that tower gives
\(C(M)\varepsilon_m L^{-2m}\) in bare plaquette units
on the full backward carriers in its top layer. The real
change to the mixed gauge preserves this curvature bound.
Averaging on a level-\(q\)
carrier gives the bound \(C(M)\varepsilon_m L^{-2(m-q)}\)
for the baseline holonomy. A top holonomy with only top
rows uses the plaquette estimate of \(V_\Lambda\) directly.
The one-step relative average on a completed bond changes
it by at most \(C(L)\delta'_m\). Multiplication of the
four unitary factors gives a defect of at most
\(C(M)\varepsilon_m+C(L)\delta'_m\) for each of these
actual tested holonomies. The remaining tests at the
artificial join are the identity or lifted top plaquette
holonomies. Their defects obey the top datum bound.

Use the test-by-test profile \eqref{mb:datum-profile}
in \eqref{mb:weighted-maximum} at a corner of a crossing
plaquette \(p\).
Here \(\check R_m\) is formed from the running minimum
of the frame, which is at least \(\theta_m-c_5\),
and the exponent of its logarithm is at least \(2\).
Thus
\(\varepsilon_\star e^{-c\delta_{\mathrm c}M\check R_m}
=o(\varepsilon_m)\). The normalized curvature at \(p\)
is at most \(C(M)\varepsilon_m\).
Since \(p\) has a corner in the prepared member of level
\(m\), its bare curvature satisfies
\[
 |\partial U''_{m,X}(p)-\mathbf1|
                    \le C(M)\varepsilon_m L^{-2m}.
\]
Finally, \(L^{m-k_0-1}=\check R_{k_0+1}\) gives
\(F_j=L_\flat^{-2}\check R_{k_0+1}^{2\log_L L_\flat}\to\infty\).
With \(\varepsilon_m/\varepsilon_j\to1\), the last bound is
strictly below the threshold of \(\kappa_\times\) for small
coupling. This removes the crossing factor in the main
integration on its stated real support, before the
reference-law restoration of the free top field.

\textbf{How a patch recovers its coupling.} Each stage gives back to the collar it integrates the increments that this collar did not receive. After $m$ stages the Wilson term is formed with

\begin{gather*}
\FitMath{x^{(m)}(\cdot):=\beta-\sum_{i\le k}b_i\,\varphi\big[\Lambda^{(m)}_i\big](\cdot),} \\[3pt]
\FitMath{\Lambda^{(m)}_i=\Lambda_i\cup(X\setminus Z''_i)\,\,(k_\flat<i\le k_\flat+m),} \\[3pt]
\FitMath{\Lambda^{(m)}_i=\Lambda_i\,\text{otherwise}}
\end{gather*}

where $\varphi[\Lambda]$ is the profile of the region $\Lambda$. The difference $A(x^{(m)}-x^{(0)},\cdot)$ is carried separately, as a Wilson piece \cref{candidate:large_field}. In this way, at each bare plaquette the weights of all the Wilson terms in the exponent still sum to $-\beta$ \cref{sec:large-field-the-tests-the-nested-regions-and-the-content-of-a}. Hence $x''_k=x^{(\mathsf{N}_k)}$ coincides with $x_k$ on the part of $X$ outside the box, and equals $x_k+O(1+k-j)$ on the shell of level $j$. This is what makes the three Wilson terms in \hyperref[nav:prepared-operation]{the prepared-component operation} cancel; two of them carry the constant $x_k$.

\phantomsection\label{nav:prepared-operation}
\textbf{What a patch retains.} A component that has passed every test carries the operation below. It is (1.71) of 122 (1989) 355, in which the summation extends over $r$ and over all sequences of regions contained in $X$, with the prefactor $1/(2^dd!)$ in front. The notation is that of \cref{candidate:large_field}, and $\chi'$ denotes the test $|B|<\delta'_k/g_k$.

\begin{gather*}
\FitMath{\mathbf{T}'_k(X)F=\frac1{|W_4|}\sum_{\text{histories in }X,\,r}\int dB\,\delta_{rT_0}(B)\,\chi'\!\int dV_{k_\flat}\,\chi\,\mathbf{T}_{k_\flat}(C_{k_\flat})\,e^{-\frac12\|\phi^{1/2}D\mathsf{H}B\|^2-A(x''_k\phi_{\mathrm{c}},U''(e^{ig_kB},V''))+V(X^\sim,U_k)+\text{numbers}}\,F .}
\end{gather*}

%
%
The final test is \(\chi_k(X^{\sim-6})\), evaluated
after the draw of the new \(V_k\). Write \(H\) for
the level-\(k\) hull of the preceding region in \(X\).
In the quiet boxed nest, \(\Omega_k^c\cap X=H^{\sim8}\)
and \(X=H^{\sim10}\), with the layers taken in the
level-\(k\) partition. Thus \(X^{\sim-6}=H^{\sim4}\)
supplies the test bases left by the carried
\(\chi_k(\Omega_k^{\sim4})\). On the passing branch
the two tests give the ordinary level-\(k\) density
bound on \(X\).
The localized-background comparison on the margin
gives the extendibility implication of Comm. Math.
Phys. 122 (1989), p.~378: the new \(V_k\) gives an
extension whose localized bare plaquettes on
\(\Omega_k^c\cap X\) are strictly below the \(2\varepsilon_k\eta_k^2\)
threshold of \(\chi_{k,\Lambda}\).

In the recorded frame both \(V_k\) and \(V_\Lambda\)
are the identity on the box tree and its entrance
bond. For each non-tree bond, close it by the tree
path. Its value in each configuration is determined
by the corresponding cycle product. Express the
cycle products by plaquettes with their transports.
The box has fixed size in level-\(k\) units, so the
number of plaquettes used for each bond is bounded
in terms of \(M\). The density and minimizer bounds
give \(\lvert\log(V_kV_\Lambda^{-1})(b)\rvert
\le C(M)\varepsilon_k<M_0\varepsilon_k\), by the
choice of \(M_0\). Hence \(\chi_{k,\Lambda}\) and
\(\chi(\Lambda^{\rm pr})\) both equal one on the
passing support. Merge the two positive density tests
into the test of the enlarged region.

For the complete invariant density, the tree
constraint can now be removed. Hold the exterior
and the entrance root fixed. Recursion along the
tree writes box Haar measure as the product of Haar
measure on the box gauge transformations and Haar
measure on the non-tree links of the slice. The
recursive substitutions on the tree and the left and
right translations on the other bonds preserve Haar
measure. The gauge transformations lift constantly on
each level-\(k\) pillar. They preserve the exterior,
the stored variables and the inner quadratic and
Wilson factors, and transform the outer background
covariantly. The localized interaction factors and
the remaining tests are invariant. The integral of
this invariant extension against box Haar therefore
equals the slice integral. If a field in the original
gauge is carried, transport it by the same Haar gauge
draw. The healed \(X\) is removed from the label and
its internal history is summed in the bracket
\cref{candidate:large_field}. For an island the final
test fails. The two positive cutoffs and the tree
constraint then remain, and \(X\) retains 

\begin{gather*}
\FitMath{\chi^c_k\,\chi_{k,\Lambda}\,\chi(\Lambda^{\mathrm{pr}})\,\delta_{\mathrm{ax}}(V_kV_\Lambda^{-1})\,\mathbf{T}'_k(X)}
\end{gather*}

 as a single block under its frame sum, for as long as it stands ((1.72), (1.104)). The cut-off $\chi(\Lambda^{\mathrm{pr}})$ is absent from (1.102). A drop-out after $n$ stages retains the following: the one-step operations of the levels that were not reached; the bare integrals $\int dV_l\,\delta(\mathcal{M}(V_l)V_{l+1}^{-1})$ of the levels already treated; the core; and the functions present together with the failing factor $1-\chi$. The terms produced up to that point are added to the boundary terms.

\begin{proposition}[{minimiser in a box; Comm. Math. Phys. 122 (1989) 175, Proposition 1}]\label{prop:large-field-2-4} There exist a radius $a_\Lambda>0$ and a constant $C(M)$, both chosen after $M$ and before $\gamma$, with the following property. Suppose the exterior data satisfy $|\partial V_k-\mathbf{1}|<\varepsilon$ with $C(M)\varepsilon\le a_\Lambda$. In the application $\varepsilon\le C\varepsilon_k$, so this is a ceiling on $\gamma$. Take the class of continuations with $|\partial V_k-\mathbf{1}|<a_\Lambda$ at those level-$k$ plaquettes whose corners all lie within distance $1$ of $\Lambda^{\mathrm{pr}}$. On this class the background $U_{k,X}(V_k)$ exists. Within the class, the map $V_k|_{\Lambda^{\mathrm{pr}}}\mapsto A(U_{k,X}(V_k))$, which is a function of $V_k$ on the level-$k$ bonds having an end point in $\Lambda^{\mathrm{pr}}$, has exactly one critical gauge orbit $V_\Lambda$. This orbit is the set of its minima in the class, it is analytic in the data, and it satisfies $|\partial V_\Lambda-\mathbf{1}|\le C(M)\varepsilon$. Ba{\l}aban states it for the function on fields of mild regularity and for $\varepsilon$ sufficiently small, with the bound $CM^5\varepsilon$. The orbit of this class is the one used in (3) and (4). Moreover, $U_0:=U_{k,X}(V_\Lambda)$ satisfies $J(U_0)=0$ in $\Lambda^{\mathrm{pr}}$. 
\end{proposition}

The proof is a contraction argument in a ball whose radius depends on $M$, which is why a class of continuations is fixed. In axial gauge the Hessian is bounded below by $c(M)>0$; Ba{\l}aban obtains the lower bound $cM^{-5}$. Indeed, the Hessian is at least $c\|\partial B\|^2$ less a small multiple of $\|B\|^2$. In a tree gauge one has the Poincar{\'e} inequality

\begin{gather*}
\FitMath{\sum_b|B(b)|^2\le C(M)\sum_p|\partial B(p)|^2}
\end{gather*}

(122 (1989) 355, (1.7)--(1.9), with $CM^5$). The corresponding inequality for a single bond of the nest, $|B(b)|^2\le C(M\check{R}_k)^C\sum_p|\partial B(p)|^2$, yields the factor of the test in (3). It is (1.78) there, with every bond and every plaquette taken in its own lattice. The power of $\check{R}_k$ has two origins: the side of a shell in its own lattice, which is at most $CML\check{R}_k$ sites, and the number $\mathsf{N}_k+1$ of shells, where $\mathsf{N}_k=\check{R}_k$. The first form serves the level-$k$ variable $V'_k$ and the uniqueness of $V_\Lambda$. The form for one bond serves the multi-scale field $B$.

\textbf{A gauge for the configuration generated in a healing.} The configuration generated in a healing must belong to the analyticity spaces. To this end it is put in the form $e^{i\eta_kA^{\mathrm{g}}}U$, where $A^{\mathrm{g}}$ and its covariant differences $\nabla_UA^{\mathrm{g}}$ ($\nabla^U_\mu$ as in the Notation) are smaller than $(1+c\beta_{\mathrm{s}})\alpha_{1,k}$ \cref{candidate:large_field}. On boxes of side $O(1)$ around the sites $y$ there are local gauges $\hat{u}_y$, with potentials $A^{(y)}$. They are glued together by a partition of unity $(\chi_y)$ taken in the quaternions,

\begin{gather*}
\FitMath{u^{\mathrm{g}}(x)=\mathcal{N}\Big(\sum_y\chi_y(x)\hat{u}_y(x)\Big),} \\[3pt]
\FitMath{\mathcal{N}(m)=(\det m)^{-1/2}m,}
\end{gather*}

and then $A^{\mathrm{g}}=\sum_y\chi_yA^{(y)}$ is a mean whose weights sum to $1$. This holds apart from quadratic terms and from terms comparable to the differences $\hat{u}_y\hat{u}_{y'}^{-1}-\mathbf{1}$ between local gauges at neighbouring sites. In $\nabla_UA^{\mathrm{g}}$, the term $\sum_y(\partial\chi_y)A^{(y)}$ has weights summing to $0$. It is therefore comparable in size to the differences between the potentials of neighbouring sites, which have a leading part in common. The square root is the branch on $|\det m-1|<1$ that takes the value $1$ at $\det m=1$. The mean in the display lies in this domain. It is the product of one local gauge $\hat{u}_y$ with $\mathbf{1}$ plus an average of the small differences $\hat{u}_{y'}\hat{u}_y^{-1}-\mathbf{1}$ between local gauges at neighbouring sites. The real linear span of $\operatorname{SU}(2)$ is $\mathbb{H}$. There $\det m=|m|^2$, and $\mathcal{N}(m)=m/|m|\in G$ for $0<|m|<\sqrt2$. Wherever the root is defined, $\mathcal{N}$ is holomorphic and bi-covariant: $\mathcal{N}(u_1mu_2)=u_1\mathcal{N}(m)u_2$ whenever $u_1,u_2\in G^c$. Hence the same construction applies to complex fields.


\textbf{[2.2] Finer sequences: nested constraint sets, nested critical configurations.}

Suppose that $\mathbf S'$ is finer than $\mathbf S$ and that $U^*$ is a real configuration on $\Omega_0$, and put $V:=M_{\mathbf B}(U^*)$ and $V':=M_{\mathbf B'}(U^*|_{\Omega'_0})$. Given a configuration $U'$ on $\Omega'_0$, we write $\widetilde{U'}$ for the configuration obtained by extending it by $U^*$.

(a) Whenever $M_{\mathbf B'}(U')=V'$, one has $M_{\mathbf B}(\widetilde{U'})=V$; moreover, the difference $A_{\Omega_0}(\widetilde{U'})-A_{\Omega'_0}(U')$ is independent of $U'$.

(b) Suppose that $M_{\mathbf B'}(U')=V''$, where $V''=V'$ holds on the bare member $\Gamma'_0$. Then $(\mathcal M^q\widetilde{U'})(c)=V''(c)$ holds for each bond $c\in\Gamma_q\cap\Gamma'_q$, and the difference of the two actions takes the same value.

(c) When $U^*$ is a regular critical configuration for $(\mathbf B,V)$, its restriction $U^*|_{\Omega'_0}$ is a regular critical configuration for $(\mathbf B',V')$.

\textbf{[2.4] Comparison of two backgrounds, at real and at complex parameters.}

Suppose that $\mathbf S,\mathbf S'$ agree on a region $W$ which is collared in $\mathbf S$. Let $p$ be a bare plaquette lying well inside $W$, in the unit cell $y$, and set $D:=d_{\mathbf B}(y,W^c)$. Let $f_1,f_2$ be holomorphic functions on $\{\|w\|<r\}$ whose modulus is at most $\frac12\mathsf m^2\eta^4\xi(p)^{-4}$. Suppose further that for each real $w$ there exist real data $V_w$ on $\mathbf B$ and $V'_w$ on $\mathbf B'$ that are both $\varepsilon_1$-regular, with $\varepsilon_1\le\mathsf m$, and equal on $W$, and for which $f_1(w)=a_p(U_{\mathbf B}(V_w))$ and $f_2(w)=a_p(U_{\mathbf B'}(V'_w))$. Then

\begin{gather*}
\FitMath{|f_1-f_2|\le C\varepsilon_1^2\,e^{-\delta_0D}\,\eta^4\xi(p)^{-4}\,\text{ at real }w,} \\[3pt]
\FitMath{|f_1-f_2|\le C\,\mathsf m^2\,e^{-\frac23\delta_0D}\,\eta^4\xi(p)^{-4}\,\text{ for }\|w\|\le\tfrac15r.}
\end{gather*}


\textbf{[2.5] Averages of a complex configuration near $\mathbf 1$.}

Let $\mathbf A$ be a bare $1$-form with values in $\mathfrak g^c$, and suppose that $s:=\sup|\mathbf A|$ is sufficiently small. Then, for $0\le n\le k$ and for each bond $c$ of $\mathbb T^{(n)}$, the complex average is defined, and

\begin{gather*}
\FitMath{\mathcal M^n(e^{i\eta\mathbf A})(c)=e^{i\mathcal Q_n(\mathbf 1;\mathbf A)(c)},} \\[3pt]
\FitMath{|\mathcal Q_n(\mathbf 1;\mathbf A)(c)|\le C\,L^n\eta\,s=C\,\xi_n\,s,}
\end{gather*}

where $\mathcal Q_n(\mathbf 1;\mathbf A)(c)$ is holomorphic in $\mathbf A$ and real when $\mathbf A$ is real. It depends only on the values of $\mathbf A$ in the two blocks $B^n(c_\pm)$, so the supremum may be taken over these blocks alone. (Here the map $\mathcal Q_n(\mathbf 1;\cdot)$ is formed with the single spacing $\eta=\eta_k$ for every $n\le k$.)

\textbf{[2.10] The basic and the glued background from the exterior field alone.}

Let $V^{\rm e}$ be a configuration on $\mathsf C(Z)\setminus E$ with values in $G$, with $\varepsilon$ small and with $|\partial V^{\rm e}(p')-\mathbf 1|<\varepsilon$ at every plaquette made up of bonds of this configuration. Let $\mathsf G$ be an arbitrary gauge tree.

\emph{First part (the mixed gauge).} The objects $V_\Lambda$, $U_0$, $U_0^\circ$, $B^\Lambda$, $\mathbb A_0$, $w$, $\tilde w$, $\mathsf t_0$, $\tilde U_0=(U_0^\circ)^{\tilde w^{-1}}$ and $V_0=M_{\mathbf B''}(\tilde U_0)|_{\mathbb B_0}$ are well defined and real. They are determined by $V^{\rm e}$ and $\mathsf G$, and they satisfy

\begin{gather*}
\FitMath{|\partial V_\Lambda-\mathbf 1|<C\varepsilon;} \\[3pt]
\FitMath{|\mathbb A_0|,|\nabla^\eta\mathbb A_0|,|D^*_{\mathbf 1}D_{\mathbf 1}\mathbb A_0|\le C\varepsilon\,\text{on }\Lambda^-;} \\[3pt]
\FitMath{|\mathsf t_0-\mathbf 1|\le C\varepsilon;} \\[3pt]
\FitMath{|V_0(c)-\mathbf 1|\le CL^{i-k}\varepsilon\,\,(c\in\mathbb B_0\text{ of level }i<k).}
\end{gather*}


\emph{Second part (the quadratic form).} Assume in addition that $(\mathbf 1,V^{\rm in})$ is $\varepsilon$-regular on $\mathbf B^{\rm in}$. Then the glued datum $V^0$ is $C\varepsilon$-regular, and $U^0$, $\Phi_{V^0}$, $\mathcal H_{\mathbf B^{\rm f},V^0}$ and $\mathsf H_1$ are defined. The operator $\mathsf H_1$ is complex-linear and takes real values on real increments. In the unit cell $y$ it obeys $\xi|\mathsf H_1\mathsf f|,\ \xi^2|\nabla_{U^0}\mathsf H_1\mathsf f|\le\mathsf h(y;\mathsf f)$. Finally, for real $\mathsf f$ one has $\mathsf q_0(\mathsf f)=\sum_p\eta^4\zeta(p)|(D_{U^0}\mathsf H_1\mathsf f)(p)|^2\ge0$.

\paragraph*{The minimiser on $\Lambda^{\rm pr}$ and the inner background}

\textbf{[1.6] Regular data.} (a) \emph{Stability.} Suppose that $|\partial V(p')-\mathbf 1|<\varepsilon$, that on the double blocks $B(b_-)\cup B(b_+)$, where $b$ runs over the outer bonds of $p'$, the configuration $V$ is moderately regular, and that $\mathsf b$ is real and $|\mathsf b|_\infty=s$ is small. Under these conditions $|\partial(e^{i\mathsf b}V)(p')-\mathbf 1|<\varepsilon+K_{\rm s}s$. (b) \emph{Averages of a regular configuration.} Let $\Omega'_\bullet$ be finer than $\Omega_\bullet$, and let $U\in\mathfrak U(\mathbf B;\varepsilon)$. Then $M_{\mathbf B'}(U)$ is $C\varepsilon$-regular. The local form is as follows: whenever $|\partial U(p)-\mathbf 1|\le\alpha\eta^2\xi_j^{-2}$ holds in the blocks $B^j(x)$ around the corners of a plaquette $p'$ at level $j$, one has $|\partial(M_{\mathbf B'}U)(p')-\mathbf 1|\le2\alpha$. (c) \emph{Lifts.} Each holonomy of $V^\downarrow=M_{\mathbf B_k(\Omega)}(Q^{s*}_kV)$ is either $\mathbf 1$ or a plaquette variable of $V$. Moreover $(e^{i\mathsf b}V)^\downarrow=e^{i\mathsf b^\downarrow}V^\downarrow$. Here $\mathsf b^\downarrow(b)=\mathsf b(c)$ when $b$ crosses from the block of $c_-$ into the block of $c_+$, and the lifted increment takes the value $0$ on bonds lying inside a block.

\textbf{[1.14] The minimiser over $\Lambda^{\rm pr}$.} Take a geometry of the main step, and let $0<\varepsilon\le\varepsilon'$ with $C\varepsilon'\le a_\Lambda$. Fix a configuration $V^{\rm e}$ on $\mathsf C(Z)\setminus E$ satisfying $|\partial V^{\rm e}(p')-\mathbf 1|<\varepsilon$ at the plaquettes made up of these bonds. Let $V'$ be an extension of it to $\mathsf C(Z)$ that satisfies $|\partial V'(p')-\mathbf 1|<\varepsilon'$ on $S(P)$. We denote by $\mathcal V(V^{\rm e})$ the set of extensions $V$ for which $|\partial V(p')-\mathbf 1|<a_\Lambda$ on $S(P)$.

(o) There is always an extension for which $\varepsilon'=4N_{\rm sh}\varepsilon$, with the inequality taken in its non-strict form.

(a) On $\mathcal V(V^{\rm e})$ the configuration $U_{k,Z}(V)$ is defined. The map $V\mapsto A_Z(U_{k,Z}(V))$ is real-analytic and is invariant under gauge transformations at the sites of $P$.

(b) The critical points of this map form a single orbit. The minimum value is attained on that orbit and nowhere else. The orbit contains exactly one element $V_\Lambda=V_\Lambda(V^{\rm e})$ that is equal to $\mathbf 1$ on $\mathsf{Tr}$.

(c) On $S(P)$ one has $|\partial V_\Lambda(p')-\mathbf 1|<(1+4K_\Lambda)\varepsilon'$. Every bare plaquette $p$ of $Z$, lying in the cell $y$ of $\mathbf B^Z$, satisfies

\begin{gather*}
\FitMath{\big|\partial U_{k,Z}(V_\Lambda)(p)-\mathbf 1\big|\le\big|\partial U_{k,Z}(V')(p)-\mathbf 1\big|+K_5\,\varepsilon'\,e^{-\delta_0d_{\mathbf B^Z}(y,\Lambda^{\rm pr})}\,\vartheta_y .}
\end{gather*}


(d) \emph{Continuation.} If $\varepsilon'$ is reduced further, a unique map $\mathsf b\mapsto B_\Lambda(\mathsf b)\in\mathcal X^{\mathbb C}$ exists with the following properties. It is defined on the complex increments $\mathsf b$ on $\mathsf C(Z)\setminus E$ with $|\mathsf b|_\infty<a_\Lambda$. It satisfies $|B_\Lambda|_\infty\le r_\Lambda$. At every such increment, its value $B_\Lambda(\mathsf b)$ is a critical point, in the directions of $\mathcal X^{\mathbb C}$, of the function 

\begin{gather*}
\FitMath{\delta\mapsto\sum_p\big[a_p(\mathsf W_{V_\Lambda^\downarrow}(\mathsf b^\downarrow+\delta))-a_p(\mathsf W_{V_\Lambda^\downarrow}(\mathsf b^\downarrow))\big]}
\end{gather*}

. This map is holomorphic and satisfies $B_\Lambda(0)=0$, and

\begin{gather*}
\FitMath{|B_\Lambda(\mathsf b)|_\infty\le2K_\Lambda|\mathsf b|_\infty,} \\[3pt]
\FitMath{V_\Lambda(e^{i\mathsf b}V^{\rm e})=e^{iB_\Lambda(\mathsf b)}V_\Lambda(V^{\rm e})\,\text{ on }E\,\text{ for real small }\mathsf b .}
\end{gather*}


\paragraph{Rare patches: gains per creation and costs per site}
\label{sec:large-field-rare-patches-gains-per-creation-and-costs-per-site}

\begin{lemma}[{one factor per creation}]\label{lem:large-field-2-5} The following holds on the support of a term, at step $t=k+1$. (i) When the plaquette test of step $t$ fails at $\square$, each real $U$ that has the stored averages satisfies $\sum_{p\text{ near }\square}a_p(U)\ge c\,\varepsilon_t^2$. (ii) When the approximate fluctuation test fails, this sum, increased by the part of $\mathrm{GF}_k$ located at the blocks near $\square$, is bounded below by $c\,\delta_k^2$. (In what follows, $\mathrm{GF}_j(X;V_j)$ denotes the part of $\mathrm{GF}_j(V_j)$ located at the blocks in $X$.) (iii) One has $\langle B,\Delta^{(k)}B\rangle\ge c\|B\|^2$, and the bound extends to the kept conditional forms \cref{sec:large-field-stored-variables-and-integration-in-a-large-field-patch}. Consequently a failed fluctuation test yields 

\begin{gather*}
\FitMath{\tfrac12\langle\mathcal{A}^{\mathrm{s}}_k,\Delta^{(k)}\mathcal{A}^{\mathrm{s}}_k\rangle\ge c\,\delta_k^2/g_k^2}
\end{gather*}

. The same is true of a failed collar test or rim test, because the threshold of either differs from $\delta_k$ by a constant factor, namely $1/(4\Theta_1)$ or $\theta_c$ \cref{sec:large-field-the-tests-the-nested-regions-and-the-content-of-a}. 
\end{lemma}

For (i) and (ii) the contrapositive also establishes existence of the objects read by the tests. For (iii), once the actual background is admitted, use the group angle, including the value \(\pi\) at the cut of the logarithm. Let \(d\) be the tested threshold, including the fixed collar or rim fraction where appropriate. On the cutoff zone an angle below \(d\) lies where the cutoff is zero, so the substitution leaves it unchanged. At the other substituted bonds take the strip comparison bound \(C\varepsilon_k e^{-\check R_k}\) below \(d/2\), using the fixed lower bound on \(d/\delta_k\). The triangle inequality then shows that an old angle at least \(d\) requires a new angle at least \(d/2\). This includes an undefined old logarithm. The energy payment uses the coercive form in the retained free coordinates of that test. Since the coupling at that place is $x_{t-1}$, every created cell produces the factor $e^{-c\,p_0(g_t)^2}$. Inside $\mathbf{R}$ the tests work with the threshold $\delta'_j$, and the gauge trees they use have length at most $C\check{R}_j^{C}$. When a bond variable is large in a tree gauge, all that follows is that some plaquette exceeds $c\check{R}_j^{-C}\delta'_j$. For this reason these tests give no more than $e^{-c\check{R}_j^{-C}p_1(g_j)^2}\le e^{-p_0(g_j)}$, and this inequality requires $2\mathfrak{p}_1-Cr>\mathfrak{p}_0$.

\textbf{The energy of the actual averages.}
The small-action alternative admits a depth-uniform estimate even when
the bare curvature does not satisfy \eqref{ne:bare-regularity}.
The additional ingredient is an \(\ell^1\) estimate for the
\emph{local quadratic remainder}, followed by its propagation through
the positive flux matrices. We give the argument with its exact carrier.
All plaquette sums count each independent orientation once. For a real
\(\operatorname{SU}(2)\) field write
\[
 a_p(U)=1-\tfrac12\operatorname{Tr}U_p
       =\tfrac12\|1-U_p\|_{\mathrm{op}}^2.
\]
All group logarithms used in the proof are the branches near the identity.
The actual total measurable average retains its prescribed fallback
outside those branches.

\begin{lemma}[Local quadratic flux remainder]
\label{lem:small-energy-local-flux}
Fix the odd block size \(L\).
There are \(d_L>0\) and \(C_L<\infty\) with the following property.
For a coarse plaquette \(C\), let \(D(C)\) be the rectangular
site box formed by its four centered fine endpoint blocks.
If the fine plaquette logarithms
\(F_p=-i\log U_p\), for plaquettes internal to \(D(C)\), satisfy
\(\max|F_p|\le d_L\), all averages used by \(C\) are on their
regular branches, and
\begin{equation}
 F_C^+=\sum_p T_{C,p}O_{C,p}F_p+R_C,\qquad
 |R_C|\le C_L\sum_{p\subset D(C)}|F_p|^2 .
 \label{eq:small-energy-local-flux}
\end{equation}
Here \(T\) is the positive scalar flux matrix of
\eqref{ne:positive-flux}, with row sum \(L^2\) and column sum
\(L^{-2}\), and the color factors are orthogonal.
After periodic identifications, a collapsed coefficient may be a sum
of these transported terms; its norm is still dominated by the
corresponding scalar coefficient with multiplicity.
\end{lemma}

\begin{proof}
The union of the four endpoint blocks is already a rectangular
site box. In the two directions of \(C\) it contains two adjacent
blocks, and in each transverse direction it contains one.
Every internal coordinate-order contour and every joining straight
segment in the four actual averages stays in this box.
The fine plaquettes in the sheet formula
\eqref{ne:positive-flux} also have all four corners there.

Choose the coordinate comb from the initial corner of \(C\).
Its paths stay in this rectangle. The cycle made by a remaining
fine edge and its two comb paths is filled by commuting coordinate
steps inside the same rectangle. Every inserted elementary
plaquette is therefore internal to \(D(C)\); no extra layer is needed.
The unitary product inequality and the small logarithm give, in
this real gauge, \(U_b=e^{ia_b}\) and
\begin{equation}
 \|a\|_{\ell^2(D(C))}
       \le C_L\|F\|_{\ell^2(D(C))}.
 \label{eq:small-energy-local-comb}
\end{equation}
Indeed each edge cycle uses a finite number of internal plaquettes,
bounded for fixed \(L\), and the resulting finite incidence matrix
has bounded row and column sums. This proves the estimate directly,
as in \eqref{ne:local-comb-estimate}.

Decrease \(d_L\) so that these finite lists of \(a_b\)'s belong
to one common small analytic ball. The local implicit mean and the
coarse plaquette logarithm are then defined there.
The fine logarithm has linear part \(\mathrm d a\), and the
coarse logarithm has linear part \(\mathrm d_cQa\).
Their second derivatives on the smaller ball are bounded finite
matrices. Their remainders are consequently bounded by
\(C_L\|a\|_2^2\). The identity
\(\mathrm d_cQ a=T\mathrm d a\) follows from
\eqref{ne:linear-contour-average}--\eqref{ne:positive-flux}:
the endpoint gradient cancels and the remaining sum is the
positive sheet flux. Substituting the fine Taylor remainder,
using \eqref{eq:small-energy-local-comb}, and returning to the
original real frame proves \eqref{eq:small-energy-local-flux}.
All frame changes are unitary, even when open transports or
global holonomies are large.

On a periodic lattice the calculation is made on the finite
lifted rectangle and then projected. Repeated images are retained
in the finite incidence count. The row and column sums of \(T\)
are exactly the periodic counts in \eqref{ne:row-column-sums}.
This is an identity of actual classical averaging functions,
not an equality of quantum measures on different covers.
\end{proof}

For any family of such rectangles, each fine plaquette occurs
in at most \(C_L\) of them, with the same bound for their lifted
multiplicities. Thus, if all fine plaquettes read by a selected
set of coarse plaquettes belong to a set \(P\), the restricted
transported flux matrix \(A\) and remainder satisfy
\begin{equation}
 \|A\|_{\ell^2\to\ell^2}\le1,\qquad
 \|A\|_{\ell^1\to\ell^1}\le\rho:=L^{-2}<1,\qquad
 \|R\|_{\ell^1}\le C_L\|F|_P\|_{\ell^2}^2.
 \label{eq:small-energy-flux-norms}
\end{equation}
For the first inequality use the row sum \(L^2\), the column
sum \(L^{-2}\), and Cauchy--Schwarz as in
\eqref{ne:transported-flux-contraction}.
Deleting rows or columns only decreases these positive majorants.
The last inequality sums the actual local quadratic bound.

We next specify the backward carriers.
Let \(S\) be a finite set of plaquettes at level \(t\).
For \(C\in S\) and \(0\le j\le t\), let \(D_j(C)\) be the
union of the four centered blocks of side \(L^{t-j}\) on the
level-\(j\) grid, centered at the corners of \(C\).
Put
\begin{equation}
 P_j(S)=\bigcup_{C\in S}
       \{p\text{ at level }j:
                    \text{ every corner of }p\text{ lies in }D_j(C)\}.
 \label{eq:small-energy-ancestor-masks}
\end{equation}
The union is taken without multiplicity. We write
\(\mathcal N_0(S)=P_0(S)\): a bare plaquette is included exactly
when all its corners lie in the four blocks \(B^t(y)\)
belonging to one member \(C\) of \(S\).
Also \(S\subset P_t(S)\).

These masks are closed under all the required one-step readouts.
If \(C'\in P_{j+1}(S)\), choose a final \(C\) whose
\(D_{j+1}(C)\) contains its four corners.
The centered odd-\(L\) block tiling decomposes \(D_j(C)\)
into the children of the sites in \(D_{j+1}(C)\).
Consequently the four child blocks at the corners of \(C'\)
lie in \(D_j(C)\), and every fine plaquette in their rectangle
belongs to \(P_j(S)\). The same assertion for a bond uses
its two child blocks. It proves locality and regular-branch
admission for the actual averages on the cone, including their
internal means, without any assumption on the exterior field.
Periodic masks can contain additional plaquettes after projection;
the inclusion and the restricted matrix bounds remain valid.
They use sets of projected plaquettes and the periodic scalar
matrix entries, rather than charging an energy once per image.

\begin{proposition}[Small total energy on an exact averaging cone]
\label{prop:small-energy-all-depth}
There is \(e_*(L)>0\), independent of depth, compatible volume,
real holonomy and the exterior field, such that every real \(U\)
and every finite \(S\) satisfy
\begin{equation}
 \sum_{p\in\mathcal N_0(S)}a_p(U)
 \ge
 \min\left\{
       \tfrac14\sum_{C\in S}
                \|1-\partial\mathcal M^tU(C)\|_{\mathrm{op}}^2,
       \tfrac12e_*
     \right\}.
 \label{eq:small-energy-min-certificate}
\end{equation}
When the left side is less than \(e_*/2\), all the actual
means required by this cone are on their regular branches.
\end{proposition}

\begin{proof}
The empty set is immediate. Write
\(E=\sum_{p\in P_0(S)}a_p(U)\).
Choose the local threshold \(d_L\le1\) in the lemma smaller
if necessary. For \(0\le s\le d_L\),
\[
 s^2\le\tfrac94(1-\cos s).
\]
For example \(1-\cos s\ge s^2/2-s^4/24
\ge(11/24)s^2\) proves this for \(s\le1\).
Set
\[
 b_L=\min\left\{\tfrac34d_L,\,
                      \frac{3(1-\rho)}{32C_L}\right\}.
\]
Choose \(E_L>0\) so small that
\[
 E_L<1-\cos d_L,\qquad
             \tfrac94E_L<b_L^2,
 \qquad e_*=2E_L .
\]
All constants are chosen before the depth or the volume.

Suppose first \(E<E_L\).
Each fine plaquette in \(P_0(S)\) has its small logarithm,
and, with \(F_0\) restricted to that mask,
\begin{equation}
 a_0^2:=\|F_0\|_2^2\le\tfrac94E<b_L^2.
 \label{eq:small-energy-initial-log}
\end{equation}
If \(a_0=0\), the local remainder and flux identities inductively
give zero logarithmic curvature on every cone mask.
The flat local means are admitted, so the result follows.

Otherwise construct successive actual averages on the masks.
While \(a_j=\|F_j|_{P_j(S)}\|_2\le(4/3)a_0\), the local
threshold is respected, so the next cone means and their
plaquette logarithms exist by the lemma.
Their exact equations are
\(F_{j+1}=A_jF_j+R_j\), with
\eqref{eq:small-energy-flux-norms}.
The color matrices may depend on the actual fields; only their
stated norms are used. Unrolling these identities gives
\begin{equation}
 F_n=A_{n-1}\cdots A_0F_0+
       \sum_{j=0}^{n-1}A_{n-1}\cdots A_{j+1}R_j .
 \label{eq:small-energy-duhamel}
\end{equation}
An empty product is the identity. The first term has
\(\ell^2\) norm at most \(a_0\). For a remainder propagated
through \(n-1-j\) matrices, use first its \(\ell^1\) norm:
\[
 \|A_{n-1}\cdots A_{j+1}R_j\|_2
 \le\|A_{n-1}\cdots A_{j+1}R_j\|_1
 \le\rho^{\,n-1-j}C_La_j^2 .
 \]
Thus the induction hypothesis gives
\begin{equation}
 a_n\le a_0+\frac{C_L}{1-\rho}
                          \left(\tfrac43a_0\right)^2
       \le\tfrac76a_0<\tfrac43a_0 .
 \label{eq:small-energy-bootstrap}
\end{equation}
The middle inequality uses the specified choice of \(b_L\).
This proves the strict bootstrap, and at the same time admits
every required local mean. No unregularized exterior plaquette
has entered the induction.

For a real small logarithm,
\(\|1-e^{iF}\|_{\mathrm{op}}\le|F|\). Since \(S\subset P_t(S)\),
\[
 \sum_{C\in S}\|1-\partial\mathcal M^tU(C)\|_{\mathrm{op}}^2
 \le a_t^2\le\tfrac{16}{9}a_0^2\le4E .
\]
This proves the first alternative.
If instead \(E\ge E_L=e_*/2\), the minimum in
\eqref{eq:small-energy-min-certificate} is at most \(E\)
without any regularity assumption. This proves the assertion
for all real configurations under the total measurable averaging
convention.
\end{proof}

\paragraph{The plaquette and approximate fluctuation witnesses.}
Fix a test cube \(\square\) of step \(t=k+1\), and put
\(D=\square^{\sim4}\). Use the separated \(M_1\)-tower of
\(D\), with the geometry and smallness floors of the
real tower estimates. The test partitions are aligned,
\(\check R_t\le\check R_k\le L\check R_t\) are powers
of \(L\), and there are at least \(6\rho\) top sites in
each torus direction. The third enlargement is inside
the top member \(D_t\).
In bare units one layer has width
\(w=M_2\check R_tL^t\). On the third enlargement the
distance from the complement of \(D\) is at least \(w\),
and the part of \(D\) omitted by \(D_t\) is within
\(2M_1L^t\) of that complement. A level-\(t\) site of the third
enlargement is therefore a node. The noncentral sites
of its level-\(k\) block are free to be gauged: any node
from a lower row is within \(2M_1L^t+L^{t-1}<w\) of the
complement of \(D\). The blocks themselves lie in the
third enlargement by alignment of the grids.

For the plaquette test, let \(S_t\) be the level-\(t\)
plaquettes with all corners in \(D\). Use the Wilson
action on \(\mathcal N_0(S_t)\) from
\cref{prop:small-energy-all-depth}. A real evaluation
used to pay for this test must reproduce \(V_t\) on the
bonds of \(S_t\). If the action there is less than
\(c\varepsilon_t^2<e_*/2\), then
\[
 |\partial V_t-\mathbf1|<2\sqrt c\,\varepsilon_t
                         \qquad\text{on }S_t .
\]
Indeed, since the action is less than \(e_*/2\), the
minimum in \eqref{eq:small-energy-min-certificate} must
be attained by its squared-defect term. The lifted
datum of the tower, obtained by lifting \(V_t\) and
then taking the determining averages, is
\(\varepsilon_1\)-regular for
\(\varepsilon_1=3\sqrt c\,\varepsilon_t\).
For the real orbit estimate use the test's regularity
radius \(\varepsilon_{\mathrm v}\), with
\(\varepsilon_1\le a_1\) and
\(B_3\varepsilon_1\le\varepsilon_{\mathrm v}\le a_0\).
The real orbit exists and its curvature in the top
member is less than \(B_3\varepsilon_1L^{-2t}\) in
bare units. The plaquette test on \(\square^\sim\)
passes when \(3B_3\sqrt c<1\). This also rules out
an undefined orbit under the same small-action bound.

For the approximate fluctuation test, also include the
level-\(k\) plaquettes in the full double blocks used
in the one-step rigidity estimate for each tested bond,
including the plaquettes crossing the common face.
Call their union \(S_k\). Let \(E\) be the Wilson action
on the union of the two backward carriers. Require that
the compatible real field reproduce \(V_k\) on the
bonds of \(S_k\) and that \(\mathcal M(V_k)=V_t\) at the
corresponding top rows. Let \(G_{\square}\) be the
unweighted part of the soft axial gauge action containing the
contour variables used by these comparisons. At birth the
approximate test is placed one coarse-cell layer inside
the region of the gauge factors. Since \(M\ge L^2M_2\),
three test-cube layers fit inside this margin, so the
block centers of the third enlargement carry those
factors. If
\(E+G_{\square}<c\delta_k^2<e_*/2\), the same energy
argument bounds the level-\(k\) and level-\(t\)
plaquette defects by \(2\sqrt c\,\delta_k\) on their
respective sets. The tower estimate now applies with
\(\varepsilon_1=3\sqrt c\,\delta_k\), again with
\(\varepsilon_1\le a_1\) and
\(B_3\varepsilon_1\le\varepsilon_{\mathrm v}\le a_0\).
The two endpoint
blocks needed for a bond in \(\square^{\sim2}\) lie
in the top region with their averaging carriers.
Averaging the bare background to level \(k\) bounds the
plaquette defects of \(\mathcal M^kU_{t,\square}\) on
these blocks by \(2B_3\varepsilon_1 L^{-2}\).
For \(B_3\varepsilon_1\le L^2c^{\mathrm{elem}}_2\) the
contour means of this background exist. The small
plaquette bound on \(V_k\) ensures the existence of
its corresponding contour means as well.

In the chosen block axial representative the background
contour means equal \(\mathbf1\). For the means of \(V_k\),
the identity
\(1-\tfrac12\operatorname{Re}\tr H
=\tfrac12|H-\mathbf1|^2\) bounds each contour defect
by \((2G_{\square})^{1/2}\).
On the top bond of a double block the background has
average \(V_t\) by the determining constraint and the
identity that averaging the lift recovers \(V_t\).
The block centers are fixed in the axial normalization.
The other field has average \(V_t\) there by the
recorded one-step constraint. By
\cref{lem:averaging-1a-1-9} their relative bond defect
is at most \(C\sqrt c\,\delta_k\) on the tested
nonpivot bonds. Choosing \(c\) smaller makes this
less than \(2\delta_k\), so the approximate test
passes. The contour existence just checked also makes
this deduction valid with the sitewise extension of the
gauge factor at undefined means.

These two contrapositives apply to every real evaluation
with the specified averaging and gauge readouts, with
constants independent of depth. To use them in a history
integral, the recorded compatible rows must cover these
carriers, the gauge action must be the one retained on
the required blocks, and the paying sets must have the
ownership and floor exclusion in \eqref{hi:event-energy}.
The true fluctuation, collar and rim tests additionally
use the coercivity and domination in the actual stored
coordinates.

\paragraph{Ordinary creation witnesses in a later drop-out.}
The paying region of an ordinary creation is separated from older
large fields when the creation occurs. In a sufficiently later
operation it belongs to the protected floor; the new payment of that
operation is taken outside the floor. These are different exclusions.

Consider a failed $P_t$- or $Q_t$-test on a cube $\square$, put
$j=t-1$, and let $\mathfrak c$ be the level-$t$ cell containing
$\square$. The birth geometry gives
\[
 \mathfrak c^\sim\subset\Omega_t^c\cap D_{\mathrm{birth}},
\]
where $D_{\mathrm{birth}}$ is its component of $Z_t$ and the
enlargement is by one layer of level-$t$ cells. Containment in
$\Omega_t^c$ uses the two layers around $Q_t'$ and the four around $P_t'$;
when $t>1$, the birth margin also places this neighbourhood
inside $\Lambda_j$.
It also separates it from the older live region and its marks.

The exact averaging cones used above fit in this neighbourhood.
Indeed, $M\ge L^2M_2$, the partitions are aligned, and $L\ge3$.
Thus $\square^{\sim4}\subset\mathfrak c^\sim$.
Both sets are unions of complete level-$t$ blocks. For a plaquette
whose corners lie in $\square^{\sim4}$, all four blocks in
\eqref{eq:small-energy-ancestor-masks} therefore lie in
$\mathfrak c^\sim$. The double blocks for the approximate test
lie in the third enlargement, so the same conclusion holds for
their level-$j$ ancestor cones. Every paying bare plaquette has all
its corners in $\mathfrak c^\sim$; no enlargement of the paying
field's domain is hidden in this use of the averaging estimate.

Suppose this creation is carried into a later drop-out $X$ without
being absorbed into a healed component or an island. Every ordinary
step leaves the older entries unchanged. If an operation of level
$m$ contains this creation, its quiet-window condition gives
$t\le m_\flat$. A drop-out from that operation re-makes only
entries of levels $l>m_\flat$. It therefore leaves both
$\Omega_j$ and $\Omega_t$, and hence the determining rows of
level $j$ on $\mathfrak c^\sim$, unchanged. Its stages integrate
stored fields only at levels at least $m_\flat$, so the birth
field $V_j$, its gauge factors, and its one-step averaging
constraints are retained. This argument applies at every later
drop-out on the same branch.

Consequently, for a compatible real evaluation on $X$, the required
identities are
\[
 \mathcal M^jU=V_j\quad\hbox{on the protected birth rows},
 \qquad
 \mathcal M(V_j)=V_t\quad\hbox{on the retained step-$t$ rows}.
\]
Together with locality and the semigroup identity, these give
$\mathcal M^tU=V_t$ on every top row used in the birth test.
At $t=1$, the first identity uses the stored bare field wherever
the action field is completed by that field. The recorded gauge
readouts in the $Q_t$ case are likewise the original ones.
The exception for an earlier main-step drop-out, whose later
constraint reads the density argument rather than its integrated
datum, does not replace these protected birth rows. A mark older
than the birth is outside its paying neighbourhood; a mark
containing the birth has its re-made levels strictly above $t$.

The floor containment has the corresponding precise meaning.
For a class component $u$ containing the old creation, the
geometric statement $\mathfrak c^\sim\subset\mathsf{Fl}(u)$
contains the exact cones just checked. It separates this old
Wilson payment from the fresh payment of $u$, which uses no
plaquette with a corner in $\mathsf{Fl}(u)$. The same floor
rule places each older marked component enlarged by two
layers of unit cubes at its own level inside the later floor,
so that a fresh payment
which uses that neighborhood in the older window cannot
spend the later window's complement. At a fixed birth level,
one recorded witness per creation cell has its carrier within
two layers of that cell. A fixed plaquette or retained
gauge-factor site can therefore enter only a bounded number
of the same-level witnesses. Divide the available Wilson
and gauge coefficients by that multiplicity. For distinct
untouched pieces the source of
the shares is the disjoint union of their paying
neighborhoods. The action sequence
still assigns the old payment its preceding-level weight at birth;
at later reads the stated lower bound is a fixed multiple of
$\theta_t$. The local history must retain the failed birth
indicator, the stated one-step constraints, and, for $Q_t$, the
specified gauge factors. The compatible evaluation must reproduce
these protected rows on the whole neighbourhood. Under these
explicit readout conditions, the two contrapositives above apply
uniformly to the later evaluating field. Once an island absorbs
the creation, its payment belongs instead to the island's mass;
this argument does not charge it again. Nor does it supply the
separate witness needed for the later operation's own failed test.

For an ordinary plaquette or approximate fluctuation witness, the
protected birth rows give the same readout in the three places where
its payment is used. At a point of a later term, they are defining
constraints of the current background. For a drop-out's mass they
are imposed by \eqref{hi:compatible-fibre}. The retained one-step
constraints then give the required next average in both cases.
The failed birth test is retained in either variant of the mass:
the excluded test is the later operation's own root test. At the
first step the convention is explicit: $V_0$ is the stored bare
field and $\Gamma_0$ consists of bare bonds with both endpoints
in $\Omega_1^c$. Hence the birth neighbourhood already lies in
this prescribed bare readout.

The third place is the real inner Wilson factor of a completed
main integration. Let its level be $i$. The quiet window puts every
ordinary birth belonging to its outer history at a level
$t\le i_\flat$. Its paying neighbourhood lies in the old core.
On that core the localised inner determining set agrees with the
prepared set. Its datum is the retained $V^{\rm in}$ on the inner
rows, whereas the new increment is supported on rows with no
endpoint in the core. Thus the protected level-$t-1$ rows still
carry $V_{t-1}$, and the retained step-$t$ constraints still
produce $V_t$. The first-step bare completion is understood as above.

Here the representative matters for literal averaging identities.
The real Landau form used to write the inner Wilson action is
gauge-equivalent to a member of the constrained orbit; an arbitrary
member of that gauge orbit need not reproduce the fixed stored
datum. Use the adapted constrained representative in the preceding
identities. This leaves the Wilson sum unchanged, since each
plaquette holonomy changes by conjugation, and the gauge-fixing
variables in the certificate remain their stored values.
The prescribed adaptation acts on the read, not on those stored
variables. No equality of every historical coarse field is needed:
the protected birth rows and their one-step constraints suffice.

For an island, this paying inner Wilson factor is a function of
$\xi$ alone. Taking the supremum over a complex exterior pair does
not replace it by a Wilson action at that pair. These observations
settle the averaging and recorded-value part of the ordinary
witness read. Admission of the read to an analytic core and the
one-sided bounds on its kept forms and elimination factors are
separate requirements in \eqref{hi:read-coverage}.

%
%
For a failed ordinary fluctuation cube, a collar creation or a failed
rim cube, use the kept Gaussian form of the birth step \(t=j+1\).
The split places the failed fluctuation cubes in \(R\) and adds
the cells of the collar creations and the failed rim cubes to
the exploration set. The enlargements in the split put the
test bonds in \(Z_t\). They are kept nonpivots of the birth
step, and the quiet-window condition \(j<t\le i_\flat\)
keeps them below the integrations of the containing operation.

Write \(a'=g_j\mathcal A_j^{\mathrm s}\) for the retained
coordinate in angular units. An explored collar bond is in the
union of the passed cube masks, where the strip substitution
leaves the coordinate fixed. The rim mask excludes the ordinary
strip \(\mathfrak S\). The failed fluctuation cubes have their
masks in the enlargement of \(R\), so on these masks the
substitution either leaves the coordinate fixed or uses the
cutoff \(\theta_j(a')\). The same alternatives hold on the
rim mask, since the plain zone is contained in the bonds
meeting \(\mathfrak S\). On the ball below the corresponding
test threshold, \(\theta_j(a')=0\), and the exact product
logarithm gives \(\mathcal A_j=a'\). The logarithm also exists
throughout that ball. Thus a failed cube or rim test, or a
collar creation with \(\chi_b^1=1\), has a bond with
\(|a'|\ge c\delta_j\), where \(c>0\) is the smallest
fixed threshold fraction. The coercive reserve of the kept
form pays \(c'\delta_j^2/g_j^2\) at that original bond.
\par{}
\textbf{No integration is undone.} For the superrenormalisable model \cite[Sect.~3.A, Prop.~3.1]{B82b}, the conditional integrations are reversed at the end. This recovers a single positive quadratic form over many scales, and the factors of all the characteristic functions are then extracted from it simultaneously. Here the operations $\mathbf{T}_k(Z)$ keep the integrals, and every factor is obtained at the place where it sits. It comes from the Wilson weight, through the display above; or from the Wilson weight together with the kept gauge-fixing exponents (\cref{lem:large-field-2-5} (ii)); or from the conditional form, whose lower bound is $c$. (Taken as a whole, the conditional form inherits the bound because it is a Schur complement; the kept part of this form is treated in \cref{sec:large-field-stored-variables-and-integration-in-a-large-field-patch}.) With this, the cost of the stored Gaussian integrals is $e^{C\#_j(Z_j\cap\Omega_j)}$ (122 (1989) 355, p. 383).

\textbf{One action and many claimants.} Gains are added up in the ledger that follows, and this addition is legitimate: nothing is spent twice and the integrals that remain converge. The Wilson, gauge and stored Gaussian terms are assigned once to the witnesses used in the ledger. A \emph{piece} is a pair $(t,D)$ in which $D$ is a component of the region produced by step $t$ on which no $\mathbf{R}$ has acted up to level $k$. We write $\mathsf{w}_{t,D}$ for the weight of its one-step operation. This weight consists of the gauge-fixing factors with their normalisers, the Jacobians that arise when the averaging constraint is solved for the pivot bonds, and the conditional Gaussian. We write $\xi_{t,D}$ for the variables that belong to the piece. Let the shares of the action be as defined in (1), and extend $U_k$ beyond $\Omega_0$ by means of the stored bare field. Then at each point of the support of the term of a history one has $A(x_k(\cdot),U_k)\ge A_{\mathrm{cr}}:=\sum_{(t,D)}x_{t-1}\sum_{p\in\text{share of }(t,D)}a_p(U_k)\ge0$, so that the only part of the action spent is the Wilson sums over the shares. Moreover

\begin{gather*}
\FitMath{e^{-A_{\mathrm{cr}}}\prod_{(t,D)}\mathsf{w}_{t,D}\le\exp\Big[-c\sum_{\text{created cells }(t,\mathfrak{c})\text{ in the pieces}}p_0(g_t)^2\Big]\prod_{(t,D)}\bar{\mathsf{w}}_{t,D}(\xi_{t,D}),} \\[3pt]
\FitMath{\int\bar{\mathsf{w}}_{t,D}\le e^{C\#_t(D)},}
\end{gather*}

where each $\bar{\mathsf{w}}_{t,D}$ is a Borel function of the variables of its own piece only (the background no longer enters it), and $C=C(L)$ carries no $\log x$ per site. (1) \emph{The action is divided among the pieces, not held in common.} The share of a piece consists of the bare plaquettes that lie within one cell of those of its cells at which a plaquette test or an approximate fluctuation test has failed. Distinct pieces have disjoint shares, because the tests of a step are carried out several cells away from the hull of the older region. On a share one has $x_k(\cdot)=x_{t-1}$, since the older profiles have completed their transition and the newer ones have not yet begun it. Every other plaquette contributes $\ge0$. (2) \emph{The tests concern averages of the configuration that appears in the action.} The tests were built from localised backgrounds of step $t$, whereas the Wilson term contains $U_k$. The two are connected as follows. The identity $\mathcal{M}^{t-1}U_k=V_{t-1}$ on $\Gamma_{t-1}$ is the constraint that defines the background, and on the same set $V_t=\mathcal{M}(V_{t-1})$ is imposed by the corresponding delta function. To apply the witnesses of \cref{lem:large-field-2-5} to \(U_k\), the retained rows at the read must contain the averaging and contour readouts on their full carriers. This coverage is needed in addition to the identities on the retained rows. (3) \emph{Each exponent is split into two halves}: one half is used for payment and the other is kept for the integration. The only normaliser that depends on the coupling is integrated against half of its own exponent,

\begin{gather*}
\FitMath{\int du\,e^{-\frac12x(1-\frac12\operatorname{Re}\operatorname{tr}u)}\le C\int du\,\mathbf{1}\{|u-\mathbf{1}|<\epsilon\}\,e^{-x(1-\frac12\operatorname{Re}\operatorname{tr}u)}} \\[3pt]
\FitMath{(\epsilon\le1,\,x\epsilon^2\ge1),}
\end{gather*}

and both sides are of order $x^{-3/2}$. For a gauge factor at level \(j\), keeping half its exponent leaves the full denominator \(z(\theta_j,\varepsilon_j)\) in place. Dropping the numerator cutoff bounds its integral by \(z_\infty(\theta_j/2)/z(\theta_j,\varepsilon_j)\), where \(z_\infty\) is the untruncated integral. The displayed estimate bounds this ratio uniformly. On the cut-off zone of the strip the shift \cref{sec:large-field-stored-variables-and-integration-in-a-large-field-patch} vanishes for small $\mathcal{A}'$; consequently a large tested field is a large variable of integration. (The approximate test requires $\mathrm{GF}$ for the following reason. Take a gauge transformation at a site of $\mathbb{T}^{(t-1)}\setminus\mathbb{T}^{(t)}$. It moves $V_{t-1}$ away from $\mathcal{M}^{t-1}U_{t,\square}(V_t)$, the average of the localised background \cref{sec:large-field-the-tests-the-nested-regions-and-the-content-of-a}, which is in the block axial gauge \cref{sec:averaging-block-averages-gauge-fixing-and-the-renormalisation-group-of-this}. It does so at no cost in action, since it alters neither the action nor $V_t$.) (4) \emph{Within a patch on which the $\mathbf{R}$ of level $m$ acts}, a test that fails may not be paid for by the large field from which the patch arose. Whatever is older than the wait was paid for by plaquettes that lie in the \emph{floor}. The floor is made up of those cells of level $m_\flat$ of the old core, $\Omega^c_{m_\flat}\cap X$, which do not touch $\Omega_{m_\flat}$; this old core is a part of the core $C_{m_\flat}=X\cap Z_{m_\flat}$ of \cref{sec:large-field-stored-variables-and-integration-in-a-large-field-patch}, because $\Lambda_i\subset\Omega_i$. The gauge-fixing factors and stored Gaussians that paid for it are all of levels below $m_\flat$. Let $X$ be a drop-out, and consider stored variables with two properties: each test carried by $X$ takes its recorded value, and the plaquette variables of the datum in the old core, taken layer by layer ($V_l$ on $\Gamma_l$), are within $C\check{R}_l^2\varepsilon_l$ of $\mathbf{1}$. (These two properties hold on the supports. They also hold at the points of evaluation that occur below in the mass of a later drop-out that contains $X$. Within the mass of $X$ itself, no use is made of the test whose failure turned it into a drop-out.) For every real $U$ compatible with stored variables of this kind, one has

\begin{gather*}
\FitMath{c\,x_m\!\!\sum_{p\subset X^+,\,\text{no corner in the floor}}\!\!a_p(U)+\sum_{j=m_\flat}^{m-1}x_j\mathrm{GF}_j(X;V_j)+c\sum_{j=m_\flat}^{m-1}\|\mathcal{A}^{\mathrm{s}}_j\|_X^2\,\ge\,(\log x_m)^{2\mathfrak{p}_1-Cr}\,\ge\,p_0(g_m),}
\end{gather*}

where $X^+$ denotes $X$ enlarged by two layers of unit cubes of level $m$ and $U$ is compatible on $X^+$. For the stops with the full witness calculation below, smallness of the left-hand side forces the tested quantity below its threshold. At a final regularity stop this calculation applies to the inset and its prepared-seam test. In every case the bridge runs along the same chain. It starts from the action of a compatible configuration and passes to the plaquettes of the stored fields (by the energy of the averages). From there it passes to the regularity of the background used by the test, and then to the bond variables. That last link comes either from the rigidity of a single averaging step or from loops in the shells. In a shell of side $s\le CM\check{R}_m$, a loop composed of $\mathsf{n}$ straight pieces has a deviation from $\mathbf{1}$ that is bounded by $C\mathsf{n}s^2$ times the plaquette bound, and this is where the loss $\check{R}_m^{-C}$ comes from. The waiting collar screens off the unknown large field in the floor; this collar is what separates the floor from the region $X\cap\Omega_{m_\flat}$ in which the operation acts, and its width is at least one cell in the local units. Where the operation acts, the background of any stage has curvature at most $C(\varsigma+\check{R}_m^2\varepsilon_me^{-\delta_0M\check{R}_m})\eta_j^2$ on the layer $j$, which is the local scale. Here $\varsigma$ is the size of the plaquettes of the stored fields that follows from the smallness of the left-hand side. A drop-out uses no plaquette that has a corner in its floor. By contrast, each plaquette used beneath it lies inside the floor, and each gauge-fixing factor and stored Gaussian used beneath it is of a level below $m_\flat$. The windows of nested operations are disjoint.

\paragraph{Filling a loop on one shell lattice.}
Consider two lattice boxes of side at most \(s\), in
units of their common lattice spacing, with the inner
box strictly inside the outer one in each direction.
Use the lattice sites of the outer box outside the
inner box, with their bonds and plaquettes. Suppose
each side of the inner box has at least one layer of
sites beyond it in the outer box. A closed path in
this shell with \(n\) straight pieces contracts across
at most \(Cns^2\) plaquettes of the shell, counted with
multiplicity.

Indeed, the shell is the union of the eight rectangular
slabs beyond the faces of the inner box. Each straight
piece lies in one slab. Consecutive pieces have a
common endpoint in the intersection of their two slabs.
Choose once a lattice point in each nonempty slab
intersection. Join every change of slab to the chosen
point for that pair by a coordinate path in the
intersection, inserting the path and its reverse.
Within each slab, replace the portion between successive
chosen points by fixed coordinate paths through one
fixed point of the slab. Each replacement is a loop
of length at most \(Cs\) in a rectangular box of side
at most \(s\). Sweeping its edges to the fixed point
one coordinate at a time fills it with at most \(Cs^2\)
plaquettes.

It remains to contract the loop of slab labels. The
nerve has the eight slabs as vertices and has a simplex
for each selection with at most one of the two sides
in any coordinate. It is the join of four pairs of
points, hence is simply connected. In this fixed
finite nerve, choose paths from one vertex to each
vertex and fix a contraction of the loop associated
with each edge and the two chosen paths. Any loop of
\(n\) edges contracts by at most \(Cn\) triangle moves.
To realize one triangle move in the shell, take a
point in the common intersection of its three slabs.
Connect it to the three pairwise reference points
inside their intersections. The triangle then splits
into three loops in the individual slabs. Each has
length \(O(s)\) and costs \(O(s^2)\) plaquettes by the
same rectangular sweep. This gives the stated area.

For a real unitary field each plaquette move changes
the based loop transport by a conjugate of a plaquette
transport or its inverse. Thus the loop defect is
bounded by \(Cns^2\) times the maximum plaquette defect
on the filling.

For the main-step bond test write
\(\mathsf W_l=\mathcal M^lU\) and
\(\mathsf W_l^0=\mathcal M^l\tilde U_0\) in the
prescribed shell gauge. They must give the tested datum
and its reference in the same frame at \(b\). Suppose
the fields agree on the shell's gauge tree and exterior
bonds. Set \(\varsigma_* = c_*\delta'_m\check R_m^{-4}\)
with \(c_*>0\) fixed. Suppose the stored datum bound
is at most \(\varsigma_*\). The screening term is at
most \(\varsigma_*\) for small coupling, since
\(\check R_m^6(\varepsilon_m/\delta'_m)e^{-\delta_0 M\check R_m}\)
tends to zero. For \(\mathsf W_l\) use the small paying
energy on the full backward carriers to bound its shell
plaquette defects by \(C\varsigma_*\). For \(\mathsf W_l^0\)
the basic reference needs its own curvature estimate on
these carriers. Where its bare curvature is bounded by
\(C\varsigma_* L^{-2l}\) there, the averaging estimate
gives the same \(C\varsigma_*\) bound. On its prescribed
top rows the corresponding estimate uses the top datum
of the reference. The stage screening estimate above
does not establish this reference bound on its own.
For a test and return on the same level, take a loop \(\ell_b\) that traverses \(b\) in its
chosen orientation and otherwise uses these common
bonds on the level-\(l\) lattice. Writing the common transports before and after
\(b\) as \(P\) and \(Q\), the relative loop holonomy is
\[
  (P\mathsf W_l(b)Q)(P\mathsf W_l^0(b)Q)^{-1}
     =P\mathsf W_l(b)\mathsf W_l^0(b)^{-1}P^{-1}.
\]
Thus the norm of the tested defect is bounded by the
sum of the two loop defects. If the chosen loop has
a bounded number of straight pieces in a shell of
side \(s\le CM\check R_m\) in level-\(l\) units,
the stated filling bound
gives \(CM^2\check R_m^2\varsigma_*\) for the defect
and, after changing the constant, for its small
logarithm. This is below \(\delta'_m\) when
\(\check R_m\) is sufficiently
large. In applying this estimate we must find the
filling in the charged shell, with its area counted
with multiplicity and its backward carriers outside
the old floor. The required rows and the common
transports must come from the actual drop-out datum.

\paragraph{The datum tested before a stage is integrated.}
Consider stage \(n+1\) of the operation of level \(m\).
Put \(j=m_\flat+n\). At that point the old determining set
is the stage-\(n\) set. The only member of the sequence that
changes is the one of index \(j+1\), which grows from
\(\Omega_{j+1}\) to the prepared member. The level-\(j\)
determining rows that remain are a subset of the old rows.
The margin from the next region leaves the higher row read sets
unchanged.
The level-\((j+1)\) rows with an end in the old region are
already old determining rows. The new rows have both ends
outside \(\Omega_{j+1}\) and at least one end in the part
being added. The two order-\((j+1)\) blocks of each such
row lie outside \(\Omega_{j+1}\) by grid alignment. The
prepared separation puts both blocks inside the old
member at level \(j\), so all their level-\(j\) bonds are
old determining rows. The averaging delta on the new row
is still present when the stage is tested.

Let \(U\) be the real field of a compatible evaluation
reproducing the old rows on these full carriers. Its
level-\((j+1)\) averages on the new rows give the
prospective datum by the retained delta and the composition
of the averages. The other prospective rows have their old
values. In particular the test can be compared with \(U\)
before the collar variables have been eliminated.

For the no-inward determining set, each tested holonomy
in the prospective sequence is a plaquette of
\(\mathcal M^qU\) at its test level \(q\). A bond that
is not a determining row has both ends in the layer
just below \(q\). Its two blocks lie in that layer, where
the determining rows supply the complete one-step
average. Thus this identification also holds at a seam.
Upon joining to the tower of the component, an interior
holonomy is unchanged. A holonomy on the joining collar
is trivial or equals a top plaquette holonomy of the
original datum. The top data of this collar are outside
the old floor.

If the bare carrier of a tested plaquette avoids the
floor, the paying energy bounds its defect by
\(C\varsigma_*\). Otherwise its level is less than
\(m_\flat\) and its blocks lie in the old core. The floor
does not touch \(\Omega_{m_\flat}\), and the adjacent blocks
at the outward rows of that level reach only the touching layer.
The old plaquette defect is bounded by the regularity of the
recorded old datum.
This bound is at most \(C\check R_m^2\varepsilon_m\),
using the comparison of the thresholds at earlier
levels. The joined datum is therefore regular at a radius
tending to zero. On the part of the stage test at
weighted distance at least \(cM\check R_m\) from the
old carriers, the local regularity estimate gives
\[
 |\partial U^{(n+1)}_{m,X}(p)-\mathbf 1|
 \le C L^{-2\lambda(p)}
 \left(\varsigma_*+
       \check R_m^2\varepsilon_m
       e^{-c\delta_{\mathrm c}M\check R_m}\right).
\]
Here \(\lambda(p)\) is the local level of the joined
stage sequence at \(p\). This first establishes existence
of the tested orbit and then bounds its curvature.

In the interior away from the joining collar, the old
prepared shells have their stated levels. Before the
compressed part of the nest the exterior region has at
least the level prescribed by the test. In the compressed
part the loss in level is at most \(1+\log_L(C\check R_m)\).
Thus changing the curvature units to those of any of the
stage thresholds costs at most \(C\check R_m^2\). The
factor \(\check R_m^{-4}\) in \(\varsigma_*\) and the
exponential in the old-carrier term suffice for these
interior thresholds. This is the inset used in the implication
for carried tests. In the outer joining collar the tower can
descend to the bare scale. There a bound in its local units
gives the top-scale threshold only with an additional bound on
the lifted top datum. When the stage tests that collar, we must
estimate its failure separately.

At the final stage the literal regularity test includes
bare plaquettes of \(X\cap\Omega_m\) with the bound
\(3(1-\tfrac12\beta_{\rm s})\varepsilon_m L^{-2m}\).
Near the outer boundary of the localized background, an
inner bare bond of its level-zero layer is a determining
bond and carries the face lift of \(V_m\). Choose a
plaquette tangent to a flat outer face, in that bare
layer, crossing two faces of the level-\(m\) blocks.
All four of its bonds are determining bonds. Its
plaquette variable is exactly the level-\(m\) plaquette
variable of \(V_m\) at the four block centers. The
first-member convention that keeps only a collar of the
level-one member still leaves a layer of bare determining
bonds. It does not give the factor \(L^{-2m}\) on this
lifted plaquette.

This mismatch occurs even for a top datum obtained by
averaging a low-action bare field. Work in level-\(m\)
units and put \(\eta_m=L^{-m}\). In a fixed box containing
the four order-\(m\) blocks, take an abelian potential
\(A_1(x)=a x_2 T\), \(A_\nu=0\) for \(\nu\ne1\), where
\(T\) has eigenvalues \(1,-1\). Continue the potential
smoothly to zero outside a slightly larger fixed box
at the rim, keeping it away from the old core. Sample
it on links by \(U(b)=\exp(i\eta_m A_\nu(x_b))\).
All bare plaquette angles are bounded by
\(C|a|\eta_m^2\), and their number in the support is
at most \(C\eta_m^{-4}\). Hence the bare Wilson action
on the support is at most \(Ca^2\).

In the affine region a centered averaging step keeps
the same affine potential with the next spacing.
The formula for the contour mean reduces to arithmetic
averaging of its logarithms on the common abelian arc.
On a bond in direction \(1\) the contours at opposite
ends are translates with equal logs, and the sum of
straight path logs averages \(x_2\) to the coordinate of
the block center. On a bond in direction \(2\), translating
the contour changes its logarithm by a fixed multiple
of \(x_1-y_1\), where \(y\) is its block center. That
term averages to zero. Translations in the other two
directions change none of the logs. The construction
is local on the four order-\(m\) blocks, so the same
calculation iterates to their top plaquette. Its defect
is \(2|\sin(a/2)|\).

Taking \(a=8\varepsilon_m L^{-2m}\) for large \(m\)
therefore fails the literal rim bound, although the
available Wilson exponent of this local perturbation
is only \(C\theta_m\varepsilon_m^2 L^{-4m}\).
Thus the Wilson energy alone cannot pay for this failure.
We must obtain the additional rim estimate from the
other factors of the retained operation.

The same scale discrepancy affects the proposed fresh payment
on the retained real support. Fix the top coupling in a plateau
of the rounded cube size that contains the whole waiting window,
and put $h=m-\mathsf N_m$. Both $\mathsf N_m$ and the sides of
the fresh carriers in level-$m$ units are then fixed as $m$
varies. Create one cell at $h$ by a pure gauge, constant on
each order-$(h-1)$ block. Give one block whose centre is
outside the level-$h$ lattice a gauge value $u_x$ with
\[
  2\delta_{h-1}<|u_x-1|<\varepsilon_{h-1},
\]
and give the other blocks the identity. Before step $h$ the
gauge is constant within each fixing block and the tested
fields are flat. At step $h$ a nonpivot bond from $x$ inside
its level-$h$ block fails the approximate fluctuation test.
Its gauge-fixing factor still has positive support, and its
cost belongs to level $h-1$. Place the failing footprint in
one birth cell. With all subsequent tests passing, its single
successors stay within the class size bound. Indeed if $a_j$
is the side of a cell at level $j$, their sides satisfy
$s_j\le s_{j-1}+22a_j$, with $a_j=La_{j-1}$ on the plateau.
The creation enters the class at $m$, at the lower end of
the waiting window.

Multiply this field by the commuting perturbation at the rim
above, with its support disjoint from the birth footprint.
Use its actual successive averages as the retained fields.
The lower creation is unchanged, and from level $h$ upward
the pure gauge restricts to the identity. The real averaging
bound gives plaquette defects at level $j$ of size
$C|a|L^{-2(m-j)}$. The joined regular backgrounds exist at
radius $C|a|$. On the full double blocks used for the retained
relative bond tests, the contour variables and coarse rows
in the prescribed frames differ by $O(|a|)$. The one-step
comparison and the actual strip and collar translations
therefore keep the retained fluctuation coordinates of size
$O(g_j^{-1}|a|)$ on the fresh carriers. At the birth the failed
bond is outside the strip. The earlier stage tests read away
from the outer joining collar and at joined layers at least
$h$, where the bare curvature is bounded by $C|a|L^{-2h}$.
Since $m-h$ is fixed, these tests pass for large $m$.

The fresh gauge and Gaussian payments are $O(|a|^2)$, since
there are boundedly many fresh variables at fixed top coupling.
Together with the Wilson bound above this tends to zero for
$a=8\varepsilon_m L^{-2m}$, while the last stage test fails.
The gauge payment at the birth is already owned below the
fresh window. Thus the pointwise certificate must use the
joined layer of the rim plaquette, which is zero here, even
though the plaquette lies in the original region $\Omega_m$.
The inset estimate applies on the inset, and the ordinary
rim fluctuation tests have their separate kept-form payments.
For the negative branch of the full stage test we must supply
a different payment or an accounting that carries this branch
without a fresh gain.

For an internal seam at level \(l>m_\flat\), the common return
traverses the two adjacent shell trees and their single entrance.
Indeed, contracting each shell tree gives the chain of nested shells.
When the tested bond is itself the entrance, its defect vanishes.
Otherwise the resulting cycle uses bonds only at levels \(l\) and
\(l-1\).

We use the following reduction at a compatible real evaluation.
The tested bond and the original common return must read the retained
averages of the same bare field \(U\), in the specified frame, and the
reference values must read the averages of \(\widetilde U_0\) in that
frame. The exterior top datum used to form \(\widetilde U_0\) must
also equal the top averages of \(U\) on the exterior bonds of the
component. These are the exterior averages used in the real
component evaluation. A restored current-stage datum after the main
integration does not identify an earlier kept determining bond with
the average of the evaluation field. At such a use one needs the
retained-row identity for that bond. The extra fine bonds introduced
below are computed averages of the evaluation fields on their full
carriers and need no stored constraints of their own.

At a descendant drop-out \((m,Y)\), the field used by the
main-step bond test is the datum retained in the frame of that
drop-out. The two fields carried on a top bond meeting the
central box have different roles. The free argument \(V_m\)
enters the averaging delta functions of step \(m+1\).
On a determining bond the earlier variable \(V''(b)\) is
still the datum of the determining set. The same prepared
datum is retained on the finer shell rows. Before the central
integration, a re-made level instead keeps its bare averaging
delta functions. If the drop-out leaves at extendibility,
there has been no central integration and these are the
constraints still present in the innermost box.

For a descendant whose bond test was formed, the original
shell rows are therefore read from the retained variables of
the frame operation. The later real Wilson background is
constrained at those rows. Indeed, the drop-out keeps its
determining regions at departure for as long as it stands.
If an intervening island has closed the branch, the outer
mass reads that island's factor instead. Otherwise the
enlarged component of the descendant
lies in the old core of a containing operation, whose stages
leave those determining rows in place. In the broader
waiting convention a full re-entry at the lower end of
a window has the same distinction:
the coarse bonds whose averaging reads the free box argument
are deep inside the first retained core, and are not data
of the first stage. The top exterior of the child's central
box remains on the prepared top grid, with its original
exterior datum. Thus the bare field of the compatible
evaluation reproduces \(V''\) on the tested bond and on
the return tree, and it reproduces the child's exterior
top datum on its exterior determining bonds. The tree
condition compares these values with \(V_0\) in the
child's recorded frame.

This check concerns the test of \(V''V_0^{-1}\).
A test or delta function that reads the free top argument
would need its own comparison. The determining datum
gives no equality between that argument and the kept
prepared value on the same bond. When a descendant's
failure is charged in a later mass, use the evaluation
on the containing component and the child's charged
carriers in \(Y^+\). An evaluation in the mass of the
child itself is only defined on \(Y\), and its own
failure is omitted in that version of the mass bound.

%
%
For the selected waiting rule the child's enlargement
lies in the containing floor. Its paying region
consists of the plaquettes of \(Y^+\) with no corner
in the child's own floor. In a completed ancestor the
inner Wilson cutoff is one there. At a stopped ancestor
use its carrying Wilson weight on the same region.
Use the level-\(m\) weight in
\(Y\) and the bound of the first ordinary successor
on the added collar. The latter weight is bounded
below by \(c\theta_{m+1}\ge c'\theta_m\).
The annular carriers in the shell comparison are
contained in this paying region. The averaging energy
bound, followed by the bond and last-shell comparisons,
forces both shell tests to pass if that Wilson share is
less than \(c\theta_m(\delta'_m\check R_m^{-4})^2\).
This argument uses the child's prepared rows as
reproduced by the ancestor's compatible field in
the child's original frame.

These Wilson shares are disjoint. The share of a
descendant lies in the floor excluded from the share
of a containing mark. For incomparable marks use
the separate components at the larger level. The
smaller enlargement is inside its ordinary successor
and hence inside its later component, and the larger
enlargement stays clear of that component by the
cell separation. At equal levels the enlarged
components are already disjoint.

At main-step entry itself, the prepared real background has
the averages $V''$ on $\mathbf B''$, so it gives the test and
tree reads in the prepared frame.

For the inner Wilson read, let the child $Y$ stand as a drop-out inside a component
$Z$ whose main step is at level $i$. The waiting rule gives
$m<i_\flat$, and $Y^+$ lies in the protected old core of $Z$.
The parent's inner determining set agrees there with its prepared
set. Its inner datum retains the child's determining datum;
the parent's new increment has no endpoint in that core. Hence,
whenever the parent's real inner background has been formed,
its adapted constrained representative reproduces every retained
child shell row whose full carrier lies in $Y^+$. In particular,
it reproduces the same $V''$, the same tree values, and the same
exterior $V^{\rm e}$ in the preceding main-step comparison.
If the child stopped during preparation, this assertion concerns
its actual stopped determining list. The unperformed stages do
not supply additional rows. As with the ordinary witnesses,
replacement of a real Landau form by this constrained
representative leaves its paying Wilson sum unchanged.

Write \(P_1=Z''_{l+1}\) for \(l<m\) and \(P_1=\Lambda^{\rm pr}\)
for \(l=m\). Write \(P_2=Z''_{l-1}\) for \(l-1>m_\flat\) and
\(P_2=I_2\) otherwise. These boxes are aligned on the level-\((l-1)\)
grid. A level-\(l\) bond of the outer shell with both ends outside
\(Z''_l\) has both order-\(l\) endpoint blocks in \(P_1\setminus Z''_l\).
At either crossing, the inner endpoint lies in the first row of
level-\(l\) sites of \(Z''_l\). Its block is contained in the
one-block skin of \(Z''_l\). The prepared margin between \(P_2\) and
the boundary of \(Z''_l\) exceeds one level-\(l\) block. This follows
from the early and late nesting margins, and from the lowest-shell
width when \(P_2=I_2\). Thus every endpoint double block needed for
a level-\(l\) bond of the cycle is in \(P_1\setminus P_2\).
By the same alignment, the bare backward carriers of the
level-\((l-1)\) plaquettes used in that annulus stay in the annulus.
In particular they avoid the floor in \(I_2\) and lie in the charged
part of \(X^+\). The small-paying-energy estimate therefore bounds
the curvature of the evaluation field throughout these carriers
after averaging to level \(l-1\). This use requires the energy
on the full charged region and an evaluation field on those carriers.
Control of the stored bond values alone would give neither.

For the reference field, take the box \(Z\) in the main-step
geometry to be the component \(X\) in the drop-out estimate.
The reference uses the exterior restriction of the top averages
on this component. The endpoint blocks of each exterior top
plaquette avoiding the free bonds are outside
\(\Lambda^{\rm pr}\). By alignment with the top grid
they lie in \(X\) if all its corners lie there. For a
reference read that reaches across \(\partial X\), use
the full blocks in \(X^+\) and their original exterior
values. These carriers lie in the charged region of
the containing evaluation. The paying energy
bounds these exterior top plaquettes by \(C(L,M)\varsigma_*\).
The extension and reference-minimizer estimates put the same
bound on the top plaquettes touching a free bond. A free bond has an end in
\(\Lambda^{\rm pr}\), so each of its incident top plaquettes
lies within one top unit of the box, in the minimizer's
test set. All other plaquettes are unchanged exterior
plaquettes. For \(l<m\) the annulus lies in \(\Lambda^-\),
where the basic-reference potential and first derivative
bounds give the corresponding bare curvature bound.
For \(l=m\) use the regular basic-reference tower on \(Z\).
The full carrier of the prepared box and its first
exterior layer lies inside its top inset. The potential
bound on \(\Lambda^-\) alone would only treat the part of
that carrier in \(\Lambda^-\). In both cases the local
averaging estimate now gives the required curvature bound
of \(C(L,M)\varsigma_*\) at level \(l-1\) throughout the annulus.

For either field, replace each oriented level-\(l\) bond \(c\) by
the straight level-\((l-1)\) path with the same endpoints. The
one-step averaging formula gives the level-\(l\) value of \(c\)
as the straight-path product multiplied by a group element at
distance at most \(C(L,M)\varsigma_*\) from the identity. This
estimate uses precisely the level-\((l-1)\) plaquettes of its
endpoint double block. Gauge covariance gives the comparison at
the same endpoints in the original frame, including the entrance
and tested crossing. For a reversed orientation use the inverse
product. Telescoping the unitary products bounds the error on a
cycle with \(O(s)\) bonds by \(C(L,M)s\varsigma_*\), where \(s\)
is the outer side length in level-\((l-1)\) units. Here
\(s\leq C(L,M)\check R_m\).

The resulting cycle lies in the single-grid annulus
\(P_1\setminus P_2\). In the prescribed shell tree a contour from
the root has at most five straight pieces. The return within one
shell is obtained by concatenating two such contours, reversing
one of them and removing backtracking. With the entrance and the
tested crossing this gives a bounded number of straight pieces
on the fine grid as well. The filling therefore has area at most
\(C(L,M)s^2\). The two original cycle products are each at distance
at most \(C(L,M)(s+s^2)\varsigma_*\) from the identity. The
comparison along the original common tree bounds \(V'(b)\)
and its small logarithm by \(C(L,M)c_*\check R_m^{-2}\delta'_m\).
This is the comparison in the retained frame. Equality of the
auxiliary straight-path products of the two fields is unnecessary.

For a bond crossing the outer boundary of \(\Lambda^{\rm pr}\),
the outermost tree return, the outer entrance, and the return in
the first exterior layer are already on the top grid. The
single-grid argument applies with this exterior layer included
in \(P_1\), provided the full top carriers lie in the charged
region and the regular reference inset as above. The prepared
test set excludes bonds with an endpoint in \(I_2\), so there is
no further crossing at the inner cut.

The last-shell test uses the field obtained by replacing
part of the lowest prepared shell by the identity. Put
\(h=m_\flat\), and let \(\square\) be one of its test cubes.
Its full reading region \(\square^{\sim4}\) lies in
\(Z''_{h+1}\cap\Omega_h\). All level-\(h\) bonds of this
region are therefore determining rows of the prepared set.
The artificial datum keeps the original value on bonds
with an end in \(I_2\) and sets the remaining bonds to
\(\mathbf 1\). Each replaced bond belongs to \(\mathbb B_0\)
and has level \(h\).

Under the small-paying-energy assumption, each original
top plaquette of the reading region has defect at most
\(C\varsigma_*\). Indeed, its full order-\(h\) carrier lies
in the region of the compatible evaluation and avoids
the floor in \(\Omega_h^c\), since the reading region
and the prepared regions are unions of order-\(h\) blocks.
The preceding test \(\chi'\) has passed. In the retained
frame, \(V''_h=e^{ig_m B}V_0\) on every replaced bond,
so its distance from \(\mathbf 1\) is at most
\(\delta'_m+C L^{h-m}\varepsilon_m\) by the basic-reference
bound. Replacing the factors of any plaquette one at a
time gives the uniform bound
\[
  C\varsigma_*+4\delta'_m+C L^{h-m}\varepsilon_m
\]
on the plaquettes of the artificial datum. The two ratios
\(\delta'_m/\varepsilon_h\) and
\(L^{h-m}\varepsilon_m/\varepsilon_h\) tend to zero with
the coupling under the threshold and waiting choices.
Linear regularity of the tower on \(\square^{\sim4}\)
then bounds the bare curvature in the top inset
\(\square^\sim\) by less than
\(\tfrac12\varepsilon_h L^{-2h}\).
Thus the last-shell test also passes under small paying
energy. This argument uses both the original averages and
the bond bound in the retained frame at the artificial cut.

The cutoff of the reference law is distinct from the test
on \(B\). After the central integration the term is constant
in the free top argument on the box. In each frame the
tree constraint and the two shell tests are imposed on the
retained variables before a reference field is drawn. The
cutoff in the subsequent draw restricts
\(V'_m=V_m V_\Lambda^{-1}\) to its reference ball. Its complement
carries zero mass under that law. There is no failure event for
that cutoff on the free Haar argument.

On the support reaching the draw, the exterior plaquette
bounds and extendibility give a regular minimizer \(V_\Lambda\).
The reference ball is in the regularity domain of
\(U_{m,X}(V'_m V_\Lambda)\) for small coupling. In fact the
plaquette bound at \(V_\Lambda\) is \(C\varepsilon_m\), and the
four factors of a plaquette increase this bound by at most
\(4M_0\varepsilon_m\) on the ball. The centre \(V'_m=\mathbf 1\)
belongs to the reference tree slice. The numerator is positive
on a relative neighbourhood of that centre in the slice,
which has positive product Haar measure on the free bonds.
It is bounded by \(1\) there and on the rest of its support
after the tree delta functions have been evaluated, since the
real Wilson action with its nonnegative cutoff is nonnegative.
Thus \(0<\mathsf D_X\le 1\) in the same frame as the numerator.
The provision for an invalid denominator defines the move off
this support and needs no mark payment on it.

%
At fixed frame use the same product Haar on
the reference tree slice for every exterior.
The numerator is jointly Borel in the exterior
and the free slice variables. Its integral,
and the integral with any bounded Borel
insertion, are Borel in the exterior.
For the denominator, the continuation of
\(V_\Lambda\) in the admitted exterior chart
and the regular background continuation give
pointwise continuity of the Wilson weight
at fixed \(V'_m\). The reference cutoff in
these coordinates is fixed and the weight
is bounded by one. Dominated convergence
gives relative continuity of \(\mathsf D_X\)
on this chart. Dividing by its positive
value defines the normalized Borel draw.

The later plaquette split does read the restored field.
Its failing branch is an island. The restored \(V_m\) is
the top datum of the healed determining set there. In a
containing mass its failed plaquette test is evaluated on
these top averages of the compatible field. The test cube
lies inside the eroded island, with its reading region inside
the island. By the linear regularity estimate, failure implies
a top plaquette defect at least \(c\varepsilon_m\) in that
reading region. The energy of the averages transfers this
defect to the bare carrier in the containing Wilson action.
With its local weight bounded below by \(c\theta_m\), the
resulting exponent is at least \(c\theta_m\varepsilon_m^2\).
This is the island event in the outer genealogy. The island's
own operation mass is formed before the later plaquette
split and does not count that event.

\textbf{The cost.} Carrying $Z_n$ through step $n$ incurs three costs. The first is a factor $e^{C\log x_{n-1}}$ for each site of $\mathbb{T}^{(n)}$ lying in $T_\ell\setminus\Omega_n$. The logarithm in question is $-\tfrac32\log x_{n-1}$ times the number of free bonds less the number of gauge-fixed sites. It is produced where the fluctuation field is introduced and rescaled, $\mathcal{A}_{n-1}=g_{n-1}\mathcal{A}^{\mathrm{s}}_{n-1}$, which takes place on the whole of $\Omega_n$, whereas the field is integrated only on the smaller set $\Lambda_n\subset\Omega_n$; it is also produced where the gauge is fixed, namely on the region in which the step's plaquette tests were passed, and this region contains $\Omega_n$. The normalisation removes it everywhere. The second is a factor $e^C$ for each site of $\mathbb{T}^{(n)}$ lying in the rest of $Z_n$ enlarged by a layer. It comes from $\mathbf{E}_n,\mathbf{R}_n,\mathbf{B}_n$ and from the stored Gaussians, which are written in the scaled variable. The third is the entropy of the subregions. Altogether, for each component $Z$ of $Z_n$, the cost is $\exp[CM^4\check{R}_n^5(d'_n(Z)+1)]$, where $d'_n$ is the tree length measured in cells. Of the five powers of $\check{R}_n$, four account for volume and one for $\log x_{n-1}\le C\check{R}_n$. In the count below, the number $|Z_n|_n$ of level-$n$ cells of the region takes the place of Ba{\l}aban's linear size $d'_n$, because a region that consists of a single cell has $d'_n=0$ and still incurs a cost.

\begin{proposition}[{the count, the ledger, the credit}]\label{prop:large-field-2-6} Let a history and one of its components be given. Then the count and the ledger displayed below are valid. If moreover the component is neither healed nor marked, the credit inequality displayed below is valid for it as well. In all of these the constants depend neither on the history nor on the number of steps. The ledger and the credit are statements of Ba{\l}aban (Comm. Math. Phys. 122 (1989) 355, (1.79)--(1.81), (1.89)), here with volumes counted in cells. 
\end{proposition}

The proof does not follow his induction. It rests instead on counting once, globally, along the whole genealogy of a component.

\textbf{Shrinking and waiting.} We denote by $S(Z)$ the free evolution of \cref{sec:large-field-the-tests-the-nested-regions-and-the-content-of-a}, that is, the hull of $Z$ in the next partition together with $O(1)$ further layers, boxed. Every step that is not a drop of the cube size $\check{R}$ divides linear sizes, counted in cells, by the factor $L$. When a step is a drop, the partitions before and after it coincide, so that nothing shrinks; and of two consecutive steps, at most one is a drop. Consequently

\begin{gather*}
\FitMath{d'_n(S^{n-j}Z)\le CL^{-\nu}d'_j(Z)+C,} \\[3pt]
\FitMath{\nu\ge\lfloor(n-j)/2\rfloor\,\text{ the number of the steps }j+1,\dots,n\text{ that are not drops}}
\end{gather*}

(this display replaces the scaling of the tree length in Ba{\l}aban, Comm. Math. Phys. 122 (1989) 355, pp. 384--385, which in the present notation, with $Z^{(n-j)}$ denoting the hull of $Z$ in level-$n$ cells, reads: $d'_n(Z^{(n-j)})\le L^{-\frac12(n-j)}d'_j(Z)\le2^{-(n-j)}d'_j(Z)$ over $n-j>1$ steps, and $d'_n(S^{n-j}Z)\le C\,d'_n(Z^{(n-j)})$, or $\le C$ when the right-hand side vanishes). The added layers do not pile up, because every change of scale contracts what the earlier levels added. It follows that within $\mathsf{k}\le\mathsf{N}_j+C\log(d'_j(Z)+2)$ steps a component $Z$ of $Z_j$ does one of three things, provided the history continues up to level $j+\mathsf{k}$: it meets a creation, it joins another component, or it enters the class of $\mathbf{R}$.

\textbf{Why a quiet component is bound to enter the class.} Shrinking makes the component small, but smallness alone does not place it in the class. Consider a quiet line, that is, the sequence of components $v_i$, of the levels $i$, into which a single component is carried as long as nothing is created in it and nothing joins it. Such a line develops deterministically, $v_i=\Phi_i(v_{i-1})$, where $\Phi_i$ is the rule of the step for a lone component, and its side in cells obeys $\alpha_{i+1}\le\alpha_i+C$, $\alpha_{i+2}\le\alpha_i/L+2C$. This is an affine contraction over two levels; its fixed point is less than half of the threshold of smallness, so that room is left for a drop. Apart from the depth $k_\flat\ge1$, which is a matter of the arithmetic of the cube sizes, a line of this kind acquires on its own precisely the conditions of the class \cref{sec:large-field-stored-variables-and-integration-in-a-large-field-patch}. Moreover, a condition that fails on account of an event in the window also fails at every intermediate level, because $n\mapsto n-\mathsf{N}_n$ is strictly increasing. Hence the carrier of the event cannot have been healed in between. The waiting bound is founded on this property of the class (in Ba{\l}aban, a single sentence, p. 385).

\textbf{The count.} An unhealed component of level $n-1$ is a predecessor of the component $(n,C)$ of $Z_n$ that contains its hull. For a component $u_*$, the genealogy $\operatorname{Gen}(u_*)$ is the tree formed by all $u=(n(u),C(u))$ that equal $u_*$ or are iterated predecessors of it. It is the same object as the two of \cref{sec:large-field-the-tests-the-nested-regions-and-the-content-of-a}. The cluster of a root is its genealogy together with everything that the history records within it. The lineage of a live component is made up of the members of its genealogy, regarded as regions $\mathsf{Z}_j$, together with the creations in them. None of these three contains a component that was healed earlier and whose place a later region covers again; such a component is the root of a cluster of its own. Write $\mathsf{n}^{\mathrm{cr}}_t(u_*)$ for the number of created cells of level $t$ in $\operatorname{Gen}(u_*)$. Write $\mathsf{n}^{\mathrm{mk}}_t(u_*)$ for the number of its marked components of level $t$ other than $u_*$ itself. These are the drop-outs and islands, that is, the creations by an operation of \cref{sec:large-field-the-tests-the-nested-regions-and-the-content-of-a}, and they are counted separately here because, for each of them, the count uses only the factor $e^{-p_0}$, which is all that a drop-out contributes. Set $\mathsf{T}_t:=C\check{R}_t$; this number bounds the length of a wait and does not increase with $t$. Then for non-increasing costs $\hat{\mathsf{r}}_1\ge\hat{\mathsf{r}}_2\ge\dots\ge0$,

\begin{gather*}
\FitMath{\sum_{u\in\operatorname{Gen}(u_*)}\hat{\mathsf{r}}_{n(u)}\,|C(u)|_{n(u)}\,\le\,C\sum_t\mathsf{n}^{\mathrm{cr}}_t(u_*)\,\mathsf{T}_t\hat{\mathsf{r}}_t+C\sum_t\mathsf{n}^{\mathrm{mk}}_t(u_*)\,\mathsf{T}_t\hat{\mathsf{r}}_t .}
\end{gather*}

In particular, for each history and without any bound on the number of steps,

\begin{gather*}
\FitMath{\sum_{n\le k}\#\{\text{cells of }Z_n\}\le C\sum_{\text{creations of both kinds, at steps }t}(\mathsf{N}_t+1).}
\end{gather*}

The proof consists of three moves, and nothing is carried through an induction over the levels. (1) Contraction. The cells are counted by way of the occupied cells of a component, which are the hulls of the created cells beneath it. These form a skeleton that commutes with taking hulls, and the component never lies more than $O(1)$ cells away from it. Over a fixed number $s_0$ of levels, at least half of the occupied cells of a component that has many of them disappear, whereas a light component (one with at most $C$ of them) retains at most $C$. Let $\mathsf{o}_n$ be the number of occupied cells of level $n$ in $\operatorname{Gen}(u_*)$, and let $l_n$ be the number of its light components of level $n$. Then

\begin{gather*}
\FitMath{\mathsf{o}_n\le\tfrac12\mathsf{o}_{n-s_0}+C\,l_{n-s_0}+\sum_{n-s_0<t\le n}\mathsf{n}^{\mathrm{cr}}_t} \\[3pt]
\FitMath{(\text{quantities of levels}\le0\text{ are }0),} \\[3pt]
\FitMath{\text{whence}} \\[3pt]
\FitMath{\sum_n\hat{\mathsf{r}}_n\mathsf{o}_n\le C\sum_n\hat{\mathsf{r}}_nl_n+C\sum_t\hat{\mathsf{r}}_t\mathsf{n}^{\mathrm{cr}}_t,}
\end{gather*}

where the second inequality is obtained from the first by multiplying by $\hat{\mathsf{r}}_n$ and summing, since the costs are non-increasing. (2) Stretches. A member of $\operatorname{Gen}(u_*)$ is called an event if it contains a created cell of its own level, or comes after a marked component, or has two or more predecessors (in which case it is a joining). The genealogy is a disjoint union of stretches. Each stretch starts at an event and continues through components that have a single predecessor and contain nothing new. A stretch that starts at level $t$ has at most $\mathsf{T}_t$ light members. Indeed, after $O(1)$ levels a light component is contained in a cube of $O(1)$ cells, and once a further $\mathsf{N}_t$ quiet levels have passed it has entered the class, in which it is healed or else marked. (3) For each of the three kinds, the events are paid for through an injective assignment, and no created cell is charged more than twice. An event that contains a created cell is paid by that cell, the successor of a mark is paid by that mark, and a joining is paid by a created cell at a leaf of the tree. For joinings, select at each member that has predecessors one of them as its first predecessor. To a joining $v$, assign the leaf reached by starting from a predecessor of $v$ other than the first and always moving to the first predecessor. No two joinings receive the same leaf, since on the way up from that leaf the first member entered otherwise than from its first predecessor is $v$ (the leaves of a tree outnumber its branchings). Moreover, the leaf lies at a lower level, at which $\mathsf{T}_t\hat{\mathsf{r}}_t$ is no smaller (because $\mathsf{T}_t=C\check{R}_t$ is non-increasing in $t$), and each leaf contains a created cell. Thus a joining requires no payment of its own.

\textbf{The ledger.} In the numerical ledger assign gain $\mathsf{g}_t=c\,p_0(g_t)^2$ to each created cell and gain $\mathsf{p}_t=p_0(g_t)$ to each mark. To apply the ledger on a history we must verify the witnesses of \cref{lem:large-field-2-5} on their paying carriers. Against these stands, for every cell and every level, the cost $\mathsf{r}^{\mathrm{c}}_n=C\log x_n\,(M\check{R}_n)^4$ found above, which is bounded by the non-increasing sequence $\hat{\mathsf{r}}_n=C\mathsf{r}_n=CM^4\check{R}_n^5$; here $\mathsf{r}_n$ is the weight of a cell, \cref{sec:large-field-the-tests-the-nested-regions-and-the-content-of-a}. With these choices,

\begin{gather*}
\FitMath{\mathsf{F}(u_*):=\sum_t\mathsf{g}_t\mathsf{n}^{\mathrm{cr}}_t(u_*)+\sum_t\mathsf{p}_t\mathsf{n}^{\mathrm{mk}}_t(u_*)-\sum_{u\in\operatorname{Gen}(u_*)}\mathsf{r}^{\mathrm{c}}_{n(u)}|C(u)|_{n(u)}\,} \\[3pt]
\FitMath{\ge\,2\sum_t\mathsf{p}_t\mathsf{n}^{\mathrm{cr}}_t(u_*)\,\ge\,(2-c)\,p_0(g_{n(u_*)}),}
\end{gather*}

provided that $\mathsf{g}_n\ge2\mathsf{p}_n+C\mathsf{T}_n\hat{\mathsf{r}}_n$ and $\mathsf{p}_n\ge C\mathsf{T}_n\hat{\mathsf{r}}_n$. An event costs $\mathsf{T}_n\hat{\mathsf{r}}_n\le CM^4\check{R}_n^6\approx(\log x_n)^{6r}$. This cost is made up of $\log x_n\le C\check{R}_n$ for each site, $(M\check{R}_n)^4$ sites in a cell and a wait of length $\check{R}_n$. Since $\mathfrak{p}_0\ge6r+1$, the quantity $p_0(g_n)=A_0(\log x_n)^{\mathfrak{p}_0}$ exceeds this cost, and the extra power of the logarithm absorbs the constants. The final inequality holds for the following reason. The genealogy has a leaf, and so it has a created cell at some level $t\le n(u_*)$. Moreover $p_0(g_t)\ge(1-c)p_0(g_{n(u_*)})$, because, up to a bounded error, the flow is monotone: $x_t\ge x_{n(u_*)}-C$. When volumes are counted in cells, $e^{-\mathsf{F}(u_*)}$ coincides, up to two modifications, with the quantity whose supremum is taken on the right-hand side of (1.79). First, the $2p_0(g_j)$ that Ba{\l}aban reserves for each created region is taken from $c\,p_0(g_j)^2$ (the first of the two conditions). Second, Ba{\l}aban's cost of a cell, which in the notation of this text reads $CM^4\check{R}_j^5$, is replaced by the finer form $C\log x_j\,(M\check{R}_j)^4$. Components that consist of a single cell are included: such a component costs $\mathsf{r}^{\mathrm{c}}_n$ at every level, and this cost is paid through the light members.

\textbf{The credit, as a corollary.} Let $z=(j,Z)$ be a component that is neither healed nor marked. We write its factor as $e^{-\mathsf{F}(z)}=e^{-\kappa_j(Z)-2p_0(g_{j(Z)})}$, where $j(Z)$ is the level of the first creation occurring in its genealogy. We extend the history of $z$ by the fictitious continuation $S(Z),\dots,S^{\mathsf{k}}(Z)$. Nothing is created in this continuation, and it runs at the levels $n\le K$, where the flow bounds are available; beyond $K$ it is cut off. We mark the last member of this continuation and apply the ledger to the extended history. The extended history has the same created cells as $z$, and apart from that last member it has the same marked components. The ledger therefore gives 

\begin{gather*}
\FitMath{\mathsf{F}(z)-\sum_n\mathsf{r}^{\mathrm{c}}_n|S^{n-j}Z|_n\ge2\sum_t\mathsf{p}_t\mathsf{n}^{\mathrm{cr}}_t(z)\ge2\mathsf{p}_{j(Z)},}
\end{gather*}

which is to say

\begin{gather*}
\FitMath{\kappa_j(Z)\,\ge\,\sum_{n=j+1}^{\min\{j+\mathsf{k},K\}}\mathsf{r}^{\mathrm{c}}_n\,\big|S^{n-j}Z\big|_n ,}
\end{gather*}

that is, (1.80), in which Ba{\l}aban has $K$ in place of $\mathsf{k}$ and the summand 
\begin{gather*}
\FitMath{O(1)M^dR_n^{d+1}d'_n(S^{n-j}(Z)),}
\end{gather*}
 where $R_n$ is his size, built from $g_n$ itself. Thus the credit $\kappa_j(Z)$ pays for the entire wait. By the shrinking, the sum is bounded by $CM^4L^5\check{R}_j^6(d'_j(Z)+1)$ ((1.81)).

\textbf{Who pays.} At a creation, the deposit $c\,p_0(g_j)^2(d'_j+1)$ has degree $2\mathfrak{p}_0$ in $\log x_j$. It exceeds the total cost $CM^4L^5\check{R}_j^6(d'_j+1)$, whose degree is $6r$, so that per cell $p_0^2$ covers $\check{R}_j^6$. Now suppose a test fails inside $\mathbf{R}$. All that is known of the old credit is $\kappa\ge0$, so the renewed wait of a small region, $CM^4\check{R}_j^6$, must be paid by the new factor $e^{-p_0(g_j)}$ on its own. For this reason that factor has to be at least as good as $e^{-p_0}$. A joining neither pays anything nor releases anything.

\textbf{Joining: Ba{\l}aban's induction and the count.} Ba{\l}aban proves (1.80) by induction on $j$ and, within a step, by induction on the number of regions that join (pp. 384--387). When $X$ and $Y$ join to form $Z$, he sets free the reserve $2p_0(g_{j(X)})+2p_0(g_{j(Y)})-2p_0(g_{j(Z)})\ge c\,p_0(g_j)$ to pay for the renewed wait. This step relies on the claim of p. 387 that, once $S$ has been applied $n_1$ further times, the domain built from $Y$ includes the domain built from $X$; there $n_1$ is called a rather small number, with $n_1<10$ given as an instance. The count uses no containment of this kind. For the hull operation of this part, whether such a containment holds depends on where the two regions lie relative to the grids. In a single coordinate the hull map is $\mathsf{h}(y)=\lfloor(y+\tfrac12(L-1))/L\rfloor$, and $S$, taken with ten layers, maps an upper end $y$ to $\mathsf{h}(y)+10$. Set $y_{m_*}=10$ and $y_m=L(y_{m+1}-10)+\tfrac12(L-1)$. Then

\begin{gather*}
\FitMath{\mathsf{h}(y_m)=y_{m+1}-10,} \\[3pt]
\FitMath{\mathsf{h}(y_m+l)=y_{m+1}-9} \\[3pt]
\FitMath{(0\le m<m_*,\,1\le l\le L).}
\end{gather*}

Suppose that $Y$ ends at $y_0$ and that $X$, a block of side $5$, begins at $y_0+1$. The two regions touch. The upper end of $S^m(Y)$ is $y_m$ and the upper end of $S^m(X)$ is $y_m+1$ (for $m=0$ it is $y_0+5$), so $S^m(X)\not\subset S^m(Y)$ for every $m\le m_*$, where $m_*$ is arbitrary. At each of these steps the two partitions have ratio $L$, so that none of the steps is a drop. The count never compares $S^m(X\cup Y)$ with $S^m(Y)$.

\begin{proposition}[{the mass bound; where the small factor stands}]\label{prop:large-field-2-7} Consider a component $X$ that has passed every test of the operation $\mathbf{R}$ of some level $i\le k$ before the last decomposition $1=\chi_i+\chi_i^c$ (a healed component or an island). The operation of \cref{sec:large-field-stored-variables-and-integration-in-a-large-field-patch} on it is written

\begin{gather*}
\FitMath{\mathbf{T}'_i(X)F\,(\mathbf{p})=\frac1{|W_4|}\sum_\omega\int e^{\Phi_\omega(\xi;\mathbf{p})}F\,\mu_\omega(d\xi),} \\[3pt]
\FitMath{\mathbf{p}=(\mathbf{U},\mathbf{J})\,\text{a pair on }X^\sim,}
\end{gather*}

where $\omega$ ranges over the local histories of $X$, that is, the results of every decomposition of unity inside $X$ together with a frame in $W_4$ of the gauge tree, $\xi$ ranges over the stored variables, and $e^{\Phi_\omega}$ contains the core Wilson factor $\exp[-A(x''_i\phi_{\mathrm{c}},U''(\xi))]$. (a) For a healed component or an island, in which case $\mathsf{k}=0$,

\begin{gather*}
\FitMath{\mathsf{M}(i,X)=\frac1{|W_4|}\sum_\omega e^{c\,\mathsf{v}(\omega)}\int\sup_{\mathbf{p}\in U^c_i(X^\sim,\alpha_{0,i},\alpha_{1,i})}\big|e^{\Phi_\omega(\xi;\mathbf{p})}\big|\,\mu_\omega(d\xi)\le e^{-(2-c')p_0(g_i)},} \\[3pt]
\FitMath{\mathsf{v}(\omega)=\sum_{u\text{ listed by }\omega}\log x_{n(u)}\,\#_{n(u)}(C(u)):}
\end{gather*}

here the coupling is taken at the level of the action, and the current coupling then follows from $p_0(g_i)\ge(1-c')p_0(g_k)$; the supremum stands inside the integral; and a weight growing exponentially with the space-time volume of the patch is paid for in advance. It is this bound that the induction carries forward. Ba{\l}aban has (1.89) instead, at each fixed pair ((1.73)). (b) Take a single history and a single component $Z$ of $Z_k$ on which no $\mathbf{R}$ has acted. For every value of the exterior variables (in total variation),

\begin{gather*}
\FitMath{\Big\|\mathbf{T}_k(Z)\big[\chi_{k,\mathcal{S}}\,e^{-A(x_k(\cdot),U_k)}\big]\Big\|\le\prod_{\text{created cells }(t,\mathfrak{c})\subset Z}e^{-c\,p_0(g_t)^2}\cdot\exp\Big[C\sum_{t\le k}\#_t(Z_t\cap Z)\Big],}
\end{gather*}

with each factor taken at the coupling of its own step. Bounds of this kind do not multiply over the components, since the background couples the components. The pointwise majorant given above does multiply (Tonelli). Consequently, after integration and for an entire label, there exists a history $h$ carrying the label $\mathcal{S}$ for which the bound below holds, where $(t,\mathfrak{c})$ ranges over the created cells of $h$ lying outside the marked $X$. (The created cells are not determined by the label: a row of three failing cells and that row with its middle cell removed yield the same regions.)

\begin{gather*}
\FitMath{\int\mathbf{T}'_k(\mathcal{S})\big[\chi_{k,\mathcal{S}}\,e^{-A(x_k(\cdot),U_k)}\big]dV_k\le e^{C\sum_j\#_j(Z_j)}\prod_{(t,\mathfrak{c})\text{ outside the marked }X}e^{-c\,p_0(g_t)^2}\prod_{(i,X)\text{ marked, outermost}}\mathsf{M}(i,X).}
\end{gather*}

(The bound holds a fortiori when each site carries $C(1+\log x_j)$ and each created cell carries $e^{-c\,p_0(g_t)}$. This logarithmically weighted form is required by \hyperref[hyp:C3]{the rarity hypothesis} and by the source construction.) (c) Consider the sum over the labels, with the costs paid by \cref{prop:large-field-2-6} and the nested sequences inside each component summed. In this sum every component $C$ of $Z_k$ keeps the factor $\exp[-c\,p_0(g_k)-\kappa_1\#_k(C)]$, in which the volume is counted in sites (and so, a fortiori, the factor with $\kappa_1|C|_k$, in which it is counted in cells; this is the form in the introduction to this step). Since $d_k(C)\le\#_k(C)$, this per-site factor includes the term $-\kappa_1d_k(C)$, which is what p. 387 asserts for an arbitrary region. The decay is one further cost $O(1)$ per site; like the other costs, it is paid for by the events (\cref{prop:large-field-2-6}). If failed plaquette tests have produced a component, then $\mathbf{T}_k(Z)1\le\exp[C\sum_{t\le k}\#_t(Z_t\cap Z)]$ (119, (2.21), and p. 264, where it is said of the Wilson term that it ``yields the small factors for large field characteristic functions'').
\end{proposition}

\textbf{The bound has to reproduce itself.} (i) \emph{Supremum under the integral.} In a term of level $k\ge i$, the factor remaining on $X$ is evaluated at $(U_k,J(U_k))|_{X^\sim}$, and $U_k$ depends on the integration variables of every component. Such an integral is not controlled by a bound at each fixed pair, and bounds that are merely uniform in the other variables do not multiply. For example, take $\mathbf{1}\{|\xi_1-\xi_2|<\epsilon\}$ on $[0,T]^2$: each integral in one variable is $\le2\epsilon$, whereas the double integral is about $2\epsilon T$. What does multiply is a majorant given by a Borel function that depends only on the variables of $X$. The integral of the supremum itself is understood as an upper integral. Such a majorant can be obtained because only two factors depend on the pair: $\Phi_\omega=\Phi^\circ_\omega(\xi)+[-\tfrac12q_X+V(X^\sim)](\xi;\mathbf{p})$, where $q_X=\|\phi^{1/2}D\mathsf{H}B\|^2\ge0$ is formed at a real background, and $|V(X^\sim)|\le C\,\#\{\text{stored sites}\}$, with every site counted at its own level. The real background here is the critical configuration of the prepared sequence, taken with $B=0$ and at the datum of the main step, and it is built with the $G$-valued factor of the triple. Through that factor, $q_X$ depends on the splitting, whereas the exponent $\Phi_\omega$, taken in its entirety, is determined by the pair (see clause (vi) of the class of \cref{sec:large-field-the-induction-closed-the-order-of-the-constants-the-sum}). (ii) \emph{A fixed set.} If the supremum is taken over the arguments that occur at level $k$, the bound is not stable, since the next step replaces $U_k$ by $U_{k+1}$. The supremum is therefore taken over a larger fixed set, built only from the data of the levels $\le i$. This set contains every later argument, since on the supports the following bounds hold, with nothing lost in the age $k-j$ of the layer $\Omega_j\setminus\Omega_{j+1}$,

\begin{gather*}
\FitMath{|\partial U_k-\mathbf{1}|<C\check{R}_j^2\varepsilon_j\eta_j^2,} \\[3pt]
\FitMath{|J(U_k)|<C\check{R}_j^2\varepsilon_j} \\[3pt]
\FitMath{\text{on }\Omega_j\setminus\Omega_{j+1}\,(\text{the current in level-}j\text{ units}),} \\[3pt]
\FitMath{C\check{R}_j^2\varepsilon_j<\alpha_{0,j}\,\text{ for all }j\,\,(q\text{ is chosen after }\mathfrak{p}_0\text{ and }r,\text{ the ceiling on the couplings after }L)}
\end{gather*}

(the last inequality holds since the quotient is a negative power of $\log x_j$, since $q>\mathfrak{p}_0+2r$, and since the couplings stay below a ceiling; moreover, a layer at index distance $a$ has an influence that decays like $e^{-(a-1)_+}$, while its threshold grows like $(1+a)^{1/2}$). Hence the mass seen at level $k$ does not exceed the mass of level $i$, and a single step $\mathbf{R}\mathbf{T}$ gives back the bound with unchanged constants. (iii) \emph{The weight} $e^{c\,\mathsf{v}(\omega)}$ has as its only cost a larger constant in the ledger. Since it carries $\log x_j$ per site, it absorbs every perturbation of size $O(1)$ per site of the patch, such as the factor $e^\Psi$ of the sources, \cref{sec:sources-analytic-dependence-on-the-source-locality-and-decay-in-tame}. In \hyperref[hyp:C3]{the weighted rarity hypothesis}, the weight covers whatever the label still exhibits inside a drop-out, since the small factors produced there have already been consumed by the mass; this is why the masses are required with the weight raised. The argument requires a volume that counts each pair (level, cell) only a bounded number of times. A volume that counts the past again at every drop-out may exceed the total of all the gains; an example is a row of cells all created at the first step, with a drop-out at each completed wait. (iv) \emph{Drop-outs.} The factors kept by a drop-out contain no Wilson action, and no inner background has been formed; the kept factors are evaluated at the background of whichever later term contains $X$. On the other hand, the configuration enters the Wilson term of the carrying term, and the failed plaquette tests listed by $\omega$, which are the tests of the creations in $X$, bound plaquettes formed from averages of this configuration. That weight, over the bare plaquettes all of whose corners lie in $X$, is therefore counted with the operation. The pair is replaced by an \emph{evaluation} at $\xi$. By this we mean a real $U$ on $X$ with four properties. First, it is compatible with $\xi$, that is, $\mathcal{M}^jU=V_j(\xi)$ on the bonds of $\Gamma_j$ in $X$. Second, it is regular layer by layer, in the same sense as on the supports: on the layer $j$ one has $|\partial U-\mathbf{1}|<C\check{R}_j^2\varepsilon_j\eta_j^2$, while the plaquette variables of the stored fields differ from $\mathbf{1}$ by less than $C\check{R}_j^2\varepsilon_j$. Third, at this configuration each test listed by $\omega$ takes its recorded value, with the exception of the test whose failure made $X$ a drop-out; here the fields not formed by $\xi$, including $V_i$ on $X$, are read as averages of $U$. Fourth, $(U,J(U))|_{Y^\sim}\in U^c_m(Y^\sim,\alpha_{0,m},\alpha_{1,m})$ for each island $(m,Y)$ in $X$. Together with this configuration, the supremum also runs over further free parameters subject to one-sided bounds. These are the weight $\upsilon$ of the Wilson term, with $\upsilon\ge c\,x_t$ on $Z_t$, and the kept forms and Jacobians, which are evaluated at a background \cref{sec:large-field-stored-variables-and-integration-in-a-large-field-patch}, with form $\ge c$ and Jacobian $\le C$:

\begin{gather*}
\FitMath{\frac1{|W_4|}\sum_\omega e^{c\,\mathsf{v}(\omega)}\int^*\sup_{\text{evaluations at }\xi}e^{-A(\upsilon\mathbf{1}_X,U)}\,(\text{kept factors})\,\mu_\omega(d\xi)\le e^{-(2-c')p_0(g_i)}} \\[3pt]
\FitMath{(\upsilon\ge c\,x_t\text{ on }Z_t).}
\end{gather*}

(Over the empty set the supremum is $0$.) Take a term of some level $k\ge i$ in which $X$ occurs, and a point of its support. Put $\upsilon=x_k(\cdot)$ and take the forms and Jacobians at the background of that term. To apply the mass bound we must verify membership of this later read in the compatible fibre, using its original rows, carried tests, and form and density bounds. On prepared rows the read uses the retained datum, with the restored free outgoing field in its separate slot. For a drop-out, the left-hand side is its mass $\mathsf{M}(i,X)$. Where $X$ has no sum over frames, nothing depends on the frame listed by $\omega$, by the identity $1=|W_4|^{-1}\sum_r1$. This bound does not use the test whose failure made $X$ a drop-out. The gain $p_0(g_i)$ of that test is reserved for the renewed wait, and it is drawn on once $X$ occurs as a mark within a subsequent component (\cref{prop:large-field-2-6}). The supremum is taken over a set fibred over $\xi$. Over a product set it selects a flat $U$, at which $e^{-A}=1$, and the factor is lost. In a toy case, take $\Phi=-T(1-\cos p)$ together with the test $|\xi|\ge\epsilon$: its mass is $e^{-cT\epsilon^2}$ on the graph $p=\xi$, but approximately $1$ over the product set. For this reason each inequality of \cref{lem:large-field-2-5} is required to hold at every configuration of this set, and not merely at backgrounds. These inequalities do hold for all real configurations having the prescribed averages. Nothing is assumed inductively at a drop-out; the inductive bound is the one for the islands. (v) \emph{Islands inside $X$.} The count and the ledger are carried out on the outer genealogy of $X$, that is, on the genealogy cut off at the islands $(m,Y)$ in $X$. Each such island is a minimal node, and all of its cells, of which there are $O(1)$, are counted as created cells of level $m$. These cells are paid for in the same way as created cells. The factor $\chi^c_m$ of the new field on $Y$ less a margin is present in the term. Consequently there is a test cube of level $m$ in which the test fails at $V_m(\xi)$; this is the test $\chi_m$ of a step, and not one of the sharper tests of the stages. \Cref{lem:large-field-2-5} (i) then extracts the gain $c\,p_0(g_m)^2$ from the core Wilson factor of $X$, on plaquettes that are used by no other gain; when $X$ is a drop-out, the gain is extracted instead from the Wilson weight that (iv) counts together with its operation, and this holds at every evaluation. This test does not enter the mass of $Y$, since that mass is formed before the last decomposition. One evaluates the factor of $Y$ at a real pair of $U^c_m(Y^\sim,\alpha_{0,m},\alpha_{1,m})$, where it is bounded by its majorant. Since $\mathsf{v}$ is additive, the sum over the local histories of $Y$ separates and is at most $\mathsf{M}(m,Y)\le1$. Every leaf of the cut genealogy contains a created cell, and the ledger leaves $2\sum_t\mathsf{p}_t\mathsf{n}^{\mathrm{cr}}_t$, just as it does on the full genealogy.


\paragraph{The induction closed: the order of the constants, the sum over histories, exactness}
\label{sec:large-field-the-induction-closed-the-order-of-the-constants-the-sum}

\textbf{The class.} Let a frame $\theta_0,\dots,\theta_k$ be fixed (see the frame setting of \cref{candidate:flow}). Each choice of level $j$ that depends on the coupling is made with $\theta_j$. These choices are the thresholds, the radii, the scaling $\mathcal{A}_j=\theta_j^{-1/2}\mathcal{A}^{\mathrm{s}}_j$, the gauge-fixing weight, the Gaussian forms of the stages of $\mathbf{R}$ and the reference laws. The cube size and the waiting depth $\check{R}_j=\mathsf{N}_j$ are made instead with the running minimum $\min_{i\le j}\theta_i$ (a frame can go up by as much as $c_5$, and taking the minimum ensures that $\check{R}_j$ does not increase). Consequently the tests, the regions, the histories, the domains of analyticity and the measures of the integral operations depend on $\theta$ only, and on no entry of the density. In the bounds of this part we put $g_j:=\theta_j^{-1/2}$, so that $g_j\le\gamma$ is the floor $\theta_j\ge\gamma^{-2}$ of a frame (at $\theta=x$ this is the $g_j=x_j^{-1/2}$ of \cref{prop:large-field-2-1}). In this way the circle ``regions depend on couplings, couplings on activities, activities on regions'' is broken. The elements of $\mathfrak{C}_k(\theta)$ are presentations rather than densities: the families are not determined by the density, and locality in the history is a property of how the terms are placed. Each element is an array, specified simultaneously on every admissible torus and cutoff,

\begin{gather*}
\FitMath{\mathfrak{P}=\Big((x_j);\,(E^{(j)},\widehat{E}^{(j)},b_j);\,(R^{(j)},\widehat{R}^{(j)});\,\big(B^{(j)}_{\mathcal{S}}(X)\big);\,(\Phi_\omega)\Big)_{j\le k},} \\[3pt]
\FitMath{b_j=x_{j-1}-x_j=b[\widehat{E}^{(j)}]} \\[3pt]
\FitMath{(1\le j\le k;\,\text{at }j=0\text{ the one entry }x_0),}
\end{gather*}

whose density $\rho[\mathfrak{P}]$ has the form given in \cref{prop:large-field-2-1}. In it, $\widehat{E}^{(j)},\widehat{R}^{(j)}$ denote single families on $\mathbb{Z}^4$ that coincide with the families on the torus on lifts of non-wrapping domains; $b[\cdot]$ denotes the coefficient of \cref{ch:pointed-hessian-cancellation}; and $\Phi_\omega=\Phi^\circ_\omega(\xi)+\sum_f\Phi_{\omega,f}(\xi;\mathbf{p}|_{R_f})$ are the exponents of the integral operations attached to the islands. Membership in the class requires the following.
(i) \emph{the tube}: $|\mathsf{z}_j|\le\mathsf{t}_j\approx c\,\theta_j(\log\theta_j)^{-C}$, where $\mathsf{z}_j=\theta_j-x_j$. This relates the frame, which is given, to the weights of the array, and the choice of the frame ensures it. The tube is wide in absolute terms, but narrow enough that the discrepancy between the weight and the choices, regarded as a potential, is small: $|\mathsf{z}_j|\alpha_{0,j}\delta_j\le c(\log\theta_j)^{-c}$.
(ii) The bounds of \cref{prop:large-field-2-1} on $E^{(j)},R^{(j)},B^{(j)}$ hold, together with analyticity, reality, and invariance under a single fixed gauge action on the stored data. A boundary term is analytic as well in the switches $\mu_b$ that belong to those of its stored bonds lying close to a test, and in those of its stored fields that lie far from every test. Any two splittings of one pair give each term the same value. The boundary terms, terms away from the small-field region and collar corrections are supplied on the graded spaces assigned to their slots, including the formation addresses used for their reads. A carried entry retains its formation space, carrier and stored arguments. The standing-island factors use their birth cores. If the tests of a term take the value $1$, then each boundary term and each exponent has in its domain the triple $(U,0,J(U))$ formed from the background $U$.
(iii) \emph{identification}: $E^{(j)}$ coincides with the family produced from the earlier entries, in each of its two chains (see below), by the small-field step at the pair $(\theta_{j-1},x_{j-1})$, carried out on the whole lattice.
(iv) \emph{lattice covariance} holds in the forms stated below. Above all, the pointed totals are covariant under every isometry of $\mathbb{T}^{(j)}$, unit translations included (119, (2.29)).
(v) \emph{coherence}: whenever a chart \cref{sec:large-field-the-tests-the-nested-regions-and-the-content-of-a} identifies, scale by scale, what two histories record within one cell of $X$, it identifies the terms at $X$. This single clause contains locality in the history, covariance under translations and independence of the torus. For a standing operation, the chart identifies $\mu_\omega$ together with the Borel set carrying it, as well as its set of factors and $\Phi^\circ_\omega$. It also identifies the factors one by one: $\Phi_{\iota\omega,\iota f}(\iota\xi;\iota\mathbf{p})=\Phi_{\omega,f}(\xi;\mathbf{p})$, for arbitrary components and histories, provided that the chart $\iota$ transports both the record of the history at $R_f$ and $\xi$ on $R_f$.
(vi) For an island, $\mu_\omega$ is a Borel measure, positive and defined on a Borel set; $\Phi_\omega$ depends in a Borel way on $\xi$, is analytic in $\mathbf{p}$, and takes real values whenever the pair is real. Moreover, gauge invariance holds pointwise in $\xi$ and for each factor separately. The bounds of \emph{mass} (\cref{prop:large-field-2-7}, with a reserve constant $c_{\mathrm{r}}\ge0$ in the weight) and of \emph{oscillation} \cref{candidate:large_field} hold, the coupling in them being taken at the level at which the island was made. The class imposes no bound on the exponent itself, which is large and negative, and no clause concerns a drop-out.
(vii) Each term, taken together with its stored variables, is a non-negative measure $\hat\tau_h$.
The form and the bounds alone do not give a class that a step preserves. The gain $L^{-c(n-j)}$ of the regular terms does not follow from $\|E^{(j)}\|_{j,\kappa}\le E_0$, and neither do $|\mathbf{E}_k|\le C\sum_n|\Gamma_n|+C\sum_j\#_j(Z_j^+)$ (below) and the input of the next step, which depend on that gain. Clause (iv) cannot be omitted, for the following reason. Take the tents $h^{(j)}_z$ of \cref{ch:local-formation-program}; then $\mathsf{Q}_z:=c\,\alpha_{0,j}^{-2}A(h^{(j)}_z,\cdot)$ defines an activity whose norm does not exceed $1$, and

\begin{gather*}
\FitMath{\sum_z\mathsf{Q}_z=c\,\alpha_{0,j}^{-2}A(U),} \\[3pt]
\FitMath{\alpha_{0,j}^{-2}=c\,\theta_j(\log\theta_j)^{-2q}:}
\end{gather*}

this is the inverse coupling $\theta_j$ itself, up to the factor $c(\log\theta_j)^{-2q}$. For a regular family, such a part is removed by the counterterm $-b_jA$ together with the covariance of the pointed totals. That covariance makes the second-order part at each site equal to the multiple $b_j$ of the action at that site. Take the format of the complex class below, in which each site has its own member. In that format one can find arrays which obey all the bounds of clause (ii) and which, at each level, add no less than the same multiple of the action, with no part of it decaying as the age grows. Such arrays are ruled out by covariance of the pointed totals, already when it is required only for unit translations.

\begin{proposition}[{the step preserves the class}]\label{prop:large-field-2-11} Suppose the following:
\begin{enumerate}
\item $\mathfrak{P}\in\mathfrak{C}_k(\theta)$;
\item $\theta_{k+1}$ prolongs the frame;
\item all $g_j\le\gamma$, and they lie below a ceiling that depends on the reserve constant $c_{\mathrm{r}}$ and does not increase with it;
\item the next weight $x_{k+1}:=x_k-b[\widehat{E}^{(k+1)}]$ lies in the tube (here $\widehat{E}^{(k+1)}$ is formed first, by clause (iii), at the pair $(\theta_k,x_k)$, and its formation involves no choice of level $k+1$).
\end{enumerate}

Then each move of step $k+1$ is defined. The identity $\mathbf{R}\mathbf{T}\rho[\mathfrak{P}]=\rho[\mathfrak{P}']$ holds separately for each standing part \cref{sec:large-field-the-tests-the-nested-regions-and-the-content-of-a}, every part keeping its stored variables, and $\mathfrak{P}'\in\mathfrak{C}_{k+1}(\theta)$, the constants being unchanged. Moreover, $\mathfrak{P}'$ prolongs $\mathfrak{P}$. An array belongs to the class of level $k+1$ exactly when its cut at $k$ belongs to the class of level $k$ and, in addition, it satisfies the conditions of level $k+1$. Finally, the array whose single entry is $x_0=\beta$, with $|\beta-\theta_0|\le\mathsf{t}_0\le c\,\theta_0(\log\theta_0)^{-C}$, belongs to $\mathfrak{C}_0$. \emph{Not returned}: the hypothesis that the next weight belongs to the tube. At every step this hypothesis is provided by how the frame is chosen. For $\theta=x$ the deviation is $0$, and the requirement is that $x_{k+1}$ prolong the frame, which is the content of the flow inequalities \cref{sec:flow-the-run-its-coefficient-and-the-analytic-requirements} and of the nesting of the cube sizes, a consequence of the upper flow inequality \cref{sec:flow-a-common-companion-and-the-torus-floor}. For the reference line it is supplied by the drift bound of \cref{sec:flow-frozen-choices-slopes-and-tuned-bare-couplings}, below. (In Ba{\l}aban, Comm. Math. Phys. 119 (1988) 243, the second remark after Theorem 1 concerns the space of ``all densities satisfying the conditions of the inductive assumption'', and these conditions include (2.29) and (2.32), together with (2.33), which follows from (2.32).)
\end{proposition}

\textbf{Why the induction closes.} (a) Room: fix a fraction $\beta_{\mathrm{s}}$. The terms produced by $\mathbf{T}$ \cref{sec:large-field-stored-variables-and-integration-in-a-large-field-patch} are defined on the radii $(1+\beta_{\mathrm{s}})\alpha_{i,k}$ and decay at the rate $(1+4\beta_{\mathrm{s}})\kappa$. Moreover, $e^{-\kappa d_j(X)}\le e^{-c\kappa d_j(X)}e^{-(1+3\beta_{\mathrm{s}})\kappa d_k(\widehat{X})}$ for $j<k$, where $\widehat{X}$ denotes the union of all level-$k$ cubes meeting $X$, taken without any added layer. The operation $\mathbf{R}$ uses up this room and nothing beyond it. (b) Drift: for $m<n$ one has $g_m\le(1+c)g_n\le C(n-m)^{1/2}g_m$. This does no harm over the $\mathsf{N}_k$ levels that $\mathbf{R}$ spans, because every threshold changes only by powers of $\log x_k$. (c) Order. Every constant is of one of four kinds:
\begin{itemize}
\item free upward: its range has no upper bound (the odd integers for $L$, the powers of $L$ for $M$);
\item free downward: its range is some $(0,c]$;
\item confined to a window;
\item determined by earlier constants: a function of those alone.
\end{itemize}

The text states finitely many conditions on the constants. Each condition is assigned to the last constant on which it depends, after every determined constant in it has been expressed through the earlier ones. For every admissible choice of the earlier constants, the condition is then, at that constant, of one of three kinds:
\begin{itemize}
\item a floor on a constant free upward, which holds from some value onward;
\item a ceiling on a constant free downward, which holds on some $(0,c]$;
\item an endpoint of a non-empty window.
\end{itemize}

A condition that follows from such conditions is not counted separately, with one exception: the requirement that a window be non-empty. That requirement is a floor or a ceiling on some earlier constant, and it counts at that earlier constant. Two further stipulations apply:
\begin{itemize}
\item No equality between thresholds is imposed, since every assertion made with a threshold remains true when the threshold is lowered.
\item Some families of conditions are indexed by a quantified number, for instance the floor $c_a\ge C(b)\log L$, $C(b)\to\infty$, with an integer $b\ge1$. No single $c_a$ satisfies this for all $b$. Such a family belongs to the statement in which the number is quantified, and it is imposed only at finitely many values of that number, which are fixed before the constant.
\end{itemize}

Thus, apart from the constant it bounds, each of these conditions depends only on earlier constants, in the order

\begin{gather*}
\FitMath{(r,\mathfrak{p}_0,\mathfrak{p}_1,q,\kappa_0,\beta_{\mathrm{s}},\theta_{\mathrm{s}},\theta_c)\prec L\prec(M_1,R_1)\prec\kappa\prec(\bar\kappa,\kappa_1)\prec M_1'\prec M_2\prec(\Theta_1,\varpi_T)\prec M} \\[3pt]
\FitMath{\prec\sigma_\star\prec a_\Lambda\prec(A_1,A_0,C_0,C_1,\alpha_1,\alpha_0,E_0,M_0,\alpha',\mathsf{B})\prec(\underline\lambda,c_5,\bar B)\prec\gamma\prec(\varrho,c_a)\prec(g,\ell,\beta,K)}
\end{gather*}

Some of these constants require comment:
\begin{itemize}
\item $\underline\lambda=\tfrac14c_0$ and $c_5=2c_2$, where $c_2$ is the constant in \eqref{eq:flow-RF10}, are fixed together with $L$; in the display they appear as members of the triple that defines the frames.
\item $\bar\kappa$ is the exponent of the decoupling parameters of \cref{ch:local-formation-program}. The factor $e^{-\bar\kappa}$ of a link of the cluster expansion comes from the decay of the covariance across an $M$-cube.
\item $\delta_0M\gg\bar\kappa\gg L\kappa$: rates are chosen before the cube sizes $M_1'$, $M_2$, $M$. These are floors on $\bar\kappa$ and then on $M$, both chosen after $L$ and $\kappa$.
\item $M_1$ depends on $L$ alone, and $R_1$ depends on $L$ and $M_1$.
\item $\Theta_1\ge C/\theta_c$ and $\varpi_T\ge C$, and both are free upward.
\item $M\ge\max\{\varpi_T,L^2M_2\}$ and $A_0\gg A_1$.
\item The later ceilings on $g$ may depend on $c_a$ and on the dial $\mathsf{w}\in[0,1)$. The dial is a fraction that depends on no constant and is fixed together with the exponents.
\item The reserve constant $c_{\mathrm{r}}\ge0$ of (vi) is free upward and comes before the ceiling that depends on it.
\end{itemize}

Induction along this order now gives two facts, because finitely many floors hold together from some value onward, finitely many ceilings hold together on some $(0,c]$, and the finitely many endpoints of a window leave it non-empty owing to the conditions that count at earlier constants. Every admissible choice made up to a given position can be extended, and no later condition places a restriction on a constant chosen earlier. Adding one more floor or ceiling of the same kind leaves both facts true. The choice is therefore not circular, and a condition arising in the construction fits the order if, at the last of the constants on which it depends, it is of the kind named. A bound may contain a number that a statement returns, or a number given as a function of a whole group of constants. Any such number must be fixed before the constant that it bounds; the cycle cut below is a case of this. A representative condition on $M$ or on $\gamma$ asserts that $e^{-cM}$, $e^{-c\check{R}_k}$ or a power of $\log x_k$ dominates constants fixed earlier. \emph{The supremum form}: a condition imposed on all couplings of all runs reduces to a single ceiling on $\gamma$, namely $\sup_{x\ge\gamma^{-2}}h(x)\le c$ with $h\to0$, because the scale appears in it solely by way of $x_k$ (122 (1989) 175, p. 186: the conditions ``do not depend on any scale''). \emph{The cycle that is cut}: $\lambda_\sharp$ and $B_\sharp$ depend on $\gamma$ through $\sup|\mathsf{q}|$ (here $\mathsf{q}_i$ is the part of the increment $b_i$ contributed by the remainders; see below), whereas the ceiling $\gamma$ depends on the flow constants. Frames are therefore defined with constants $(\underline\lambda,c_5,\bar B)$ fixed before $\gamma$: $\underline\lambda(j-i)-c_5\le\theta_i-\theta_j\le\bar B(j-i)$ for $i\le j$. Hence a frame with first member $\theta_0$ has length at most $(\theta_0-\gamma^{-2}+c_5)/\underline\lambda$, in the same way that $K$ is bounded in terms of $\beta$ in \cref{sec:flow-the-run-its-coefficient-and-the-analytic-requirements}. In addition, $\underline\lambda\le\lambda_\sharp$ and $B_\sharp\le\bar B$. \emph{No caps}: there is no upper bound on $L,M,\kappa,q,\mathfrak{p}_0,c_a$, on the cutoff, the volume, an age or a number of steps, and no lower bound on a coupling. \emph{The role of each exponent}:
\begin{enumerate}
\item $\mathfrak{p}_0-\mathfrak{p}_1>Cr$, for comparing the ratio $\delta'_j/\delta_j$ with the square root of the waiting depth, and for the accumulation of the stages (\cref{sec:large-field-stored-variables-and-integration-in-a-large-field-patch}, \cref{candidate:large_field});
\item $2\mathfrak{p}_1-Cr>\mathfrak{p}_0$, so that a failed stage test still gives $e^{-p_0}$;
\item $\mathfrak{p}_0>3Cr$, so that $p_0$ dominates $\kappa\check{R}^4$ and the cost $\check{R}^6$ of a wait, and so that the window left for $\mathfrak{p}_1$ by the two preceding conditions is non-empty;
\item $q-\mathfrak{p}_0>Cr$, so that $\mathsf{N}_k\mathsf{s}_k\to0$;
\item $q>C\mathfrak{p}_0$, so that the product of an amplitude and an oscillation of degree $2\mathfrak{p}_0$ remains bounded;
\item $\kappa_0>C$, for the sum over the scales.
\end{enumerate}

(d) \emph{No circle.} New terms are built from old ones, and the old ones are bounded with the same constants that the new ones must satisfy. The circle is avoided because an old term enters again in one of three ways:
\begin{itemize}
\item unchanged;
\item through a difference, which carries a small factor. In $\mathbf{T}$ this factor is $\mathsf{s}_k$, the factor that the Cauchy estimate gives for a first difference in the fluctuation field \cref{candidate:large_field}.
\item as a whole, in the main step of $\mathbf{R}$. In that case there are two possibilities. Either the domain of the term must stretch across $c\check{R}_kL^{k-j}$ cubes of the term's own level $j<k$, in which case its decay makes the term small. Alternatively, the domain lies inside an enlarged component (a collar $X^\sim$ of \cref{candidate:large_field}). There only a volume bound $C\sum_{j\le k}\#_j(Z_j\cap X)$ is available, each scale counted in its own unit, and its constant is chosen later. The list is given below under \emph{Boundary terms}.
\end{itemize}

\emph{Uniformity of the input.} A term of level $j$ is defined on the radii at which it was created, and step $n+1$ evaluates it up to the output radius of that step. Let $w_{n,j}$ be the square of the ratio of these two radii. Let $\mathsf{f}_\nu\approx\nu^CL^{-c\nu}$ be forgetting weights; in \cref{ch:local-formation-program} they are written $\omega_\nu$. Then

\begin{gather*}
\FitMath{w_{n,j}\approx\frac{\theta_j}{\theta_n}\Big(\frac{\log\theta_n}{\log\theta_j}\Big)^{2q}\le C(n+2-j),} \\[3pt]
\FitMath{\sum_{j\le n}\mathsf{f}_{n+1-j}\|E^{(j)}\|_{j,\kappa}+\sum_{j\le n}w_{n,j}\|R^{(j)}\|_{j,\kappa}\le N_\star}
\end{gather*}

uniformly in $n$ and in the frame. This is everything that the small-field step requires of its input. The two factors in the approximation above are never unfavourable at the same time: the upper flow bound gives $\theta_j/\theta_n\le1+(n-j)$, and the lower one gives $\theta_n\le\theta_j+c_5$. The regular terms are of size $O(1)$ at each level and gain $L^{-c(n-j)}$ through their two counterterms. The remainders gain nothing; they are summed by means of the coupling. \emph{Regular terms}: let $\tau_k=x_k/\theta_k$, and let $E_{\mathsf{N}}$ be the bound on the localised Gaussian normaliser. Then

\begin{gather*}
\FitMath{\|E^{(k+1)}\|_{k+1,\kappa}\le E_{\mathsf{N}}+C\big(|1-\tau_k|+\delta_k+\mathsf{s}_kN_\star+e^{-c\,p_0(g_k)^2}\big)\le E_0 ,}
\end{gather*}

so $E_0$ is determined by the one-loop normaliser alone, and $C$ depends neither on $E_0$ nor on $N_\star$.

%
\emph{The new regular remainder.}
At output level \(t\), separate the fluctuation
clusters which contain a vertex from a remainder
of origin \(j\) with \(M\check R_j>L^{t-j}\).
Charge each of these clusters to one such vertex.
In the absolute cluster bound with one
distinguished insertion, this vertex has size
\(C\mathsf s_{t-1}w_{t-1,j}g_j^{\kappa_0}\),
using the induction bound of the input.
The root sum over the other vertices has a
uniform constant at the rate of the fluctuation
move. The sum over \(j\) must include both the
number of origins and their radius losses.

Indeed, for \(r=t-j\) the flow gives
\(x_t-C\le x_j\le x_t+C(r+1)\).
Since \(\check R_j\) is polynomial in
\(\log x_j\), the inequality \(L^r<M\check R_j\)
implies
\[
 L^r\le C(1+\log x_t)^C(1+\log(1+r))^C .
\]
The last factor on the right is bounded by
\(C L^{r/2}\). Hence the contributing ages
are \(O(1+\log\log x_t)\).
On this range \(g_j^{\kappa_0}\le Cg_t^{\kappa_0}\)
and \(w_{t-1,j}\le C(r+1)\).
The fluctuation contribution in the regular
remainder range therefore has norm at most
\[
 C\mathsf s_{t-1}(1+\log\log x_t)^2g_t^{\kappa_0}
 \le \tfrac12 g_t^{\kappa_0}.
\]
The negative power of \(\log x_t\) in
\(\mathsf s_{t-1}\) pays for the sum.
Use the final rate \(\kappa\) on the
domains one cell layer inside the output
regular region, with the ordinary contribution
supplied on the regular range of the appropriate
donor when the filing changes. There the exponentiation
of healed histories supplies the other
contribution, of norm at most \(e^{-p_0(g_t)}\).
The domains outside this range are filed
as boundary terms. Since
\(e^{-p_0(g_t)}\le g_t^{\kappa_0}/2\),
the new regular remainder has the same
bound \(g_t^{\kappa_0}\) as the input
remainders at their respective levels.

\emph{Boundary terms}: $\mathsf{B}$ is the maximum of the bounds coming from their origins \cref{candidate:large_field}, and none of these bounds depends on $\mathsf{B}$. This holds because data bounded by $\mathsf{B}$ enter the step in four ways:
\begin{itemize}
\item In $\mathbf{T}$, they enter as differences in the fluctuation field, with $C\mathsf{B}\mathsf{s}_k\le1$.
\item In the stages, they enter as differences in the collar field, taken one block of bonds at a time. Each block that is switched off contributes a negative power of $\log\theta_k$, with an exponent larger than the one belonging to (number of stages)$\times$($M$-cubes of level $k$ of a component).
\item In the main step, outside the enlarged components, they enter either through the difference under the substitution, or as undifferenced functions of a level $j<k$. In the second case the domain must extend over $c\check{R}_kL^{k-j}$ cubes, and $\mathsf{B}e^{-c\check{R}_k}\le1$.
\item Inside an enlarged component there is no smallness, and they enter only through the volume bound on $V(X^\sim)$. The constant of this bound is chosen after $\mathsf{B}$, and the bound is paid by the ledger of \cref{sec:large-field-rare-patches-gains-per-creation-and-costs-per-site}.
\end{itemize}

(e) \emph{No growth with the level.} A carried term is literally unchanged: it has the same function, the same domain of analyticity and the same tree length. The bounds of a standing island refer to the level at which it was made and do not involve $k$. Indeed, the supremum in the mass bound is taken over a fixed set, built from the data of the levels up to that level, and this set contains the arguments of all later levels. Hence the mass at level $k$ does not exceed the mass at the level at which the island was made (\cref{sec:large-field-rare-patches-gains-per-creation-and-costs-per-site}, ``The bound has to reproduce itself'', (ii)). The sizes that are returned are maxima rather than sums.


% END MODULE chapters/05_large_field.tex

% BEGIN MODULE chapters/history_mass_program.tex
\subsection{Uniform weighted masses of histories}
\label{ch:history-mass-program}

A large-field factor is an integral operation with a memory.  Its
arguments include the fields that a later operation is allowed to read;
its integration variables include the fields that an earlier operation
has retained.  These two lists need not describe the same region.  We
give the lists separate types, specify the mass of the operation, and
prove the counting and domination implications needed to propagate that
mass.  In particular, a supremum over complex arguments is placed inside
the integral.  An estimate of a partition function at each fixed argument
does not establish this stronger assertion.

We prove the geometric counts, scalar reserve estimates and their
consequences. Their application to the interacting construction requires
the background domains, kept quadratic forms, changes of variables and
localized energy bounds stated below.  A flat
quadratic comparison supplies a special case of the matrix estimates;
it does not establish their current-dependent, layered versions.

\subsubsection{The records carried by an operation}

Fix a finite four-dimensional periodic lattice.  Work in bare coordinates:
the bonds of level $j$ have length $a_0L^j$, where $L$ is an odd integer.
The torus and all partitions are chosen compatibly.  A region is a union
of partition cells, with a fixed half-open convention when a unique cell
address is required.  Closed cells are used for collars.  Passing from
one convention to the other changes only a fixed number of neighboring
cells in each estimate below.

Let $\theta_j>1$ be the reference inverse couplings, and write
\begin{equation}
 \ell_j=\log\theta_j,\qquad g_j=\theta_j^{-1/2},\qquad
 R_j=\check R_j,\qquad b_j^{\rm cell}=a_0L^jMR_j.
 \label{hi:scales}
\end{equation}
Here $M$ and every auxiliary partition size are fixed powers of $L$.
The symbol $g_j$ in this section uses the reference weight; the
physical coupling is $x_j^{-1/2}$ when $x_j\ne\theta_j$.
The partitions have cell side $b_j^{\rm cell}$.  The size sequence has
$R_{j+1}/R_j\in\{1,L^{-1}\}$, and two size drops are not consecutive.
Thus a passage between levels either coarsens these cells by $L$, or
leaves their physical partition unchanged.  Bounds are required uniformly
in the height and volume, with the fixed geometric parameters held fixed.

A label of height $s$ contains regions
\begin{equation}
 \mathcal S=(\Omega_j,\Lambda_j)_{1\le j\le s},\qquad
 \Lambda_j\subset\Omega_j,\quad
 \Omega_{j+1}\subset\Lambda_j,\quad Z_j=\Lambda_j^c.
 \label{hi:nested-label}
\end{equation}
There is an initial collar $\Omega_0=\Omega_1^\sim$; there is no
independently generated large-field region of level zero.  The state
after a step and the state after its large-field operation are different
labels.  A history records their ordered succession, rather than a
single final sequence with an inferred past.  This distinction is
necessary when an operation heals a region and a later creation covers
the same physical location.

The operation also records a class of components and, in each component,
the stages that have actually been completed.  A component in the
dynamic class has small size, the required depth, a quiet preparation
window, quarantine from earlier marks in that window including its
lower endpoint, and the single-box traces required by the local
background construction.  The dynamic class contains every eligible
component.  A static representation may deliberately use a smaller
class.  Waiting estimates below concern the dynamic rule; they would
be false for an arbitrarily withheld class.

For an eligible component $X$ of level $t$, let $t_\flat$ be the bottom
of its preparation window and let $Z''_j(X)$ be its prepared cores.
After $n(X)$ stages, with some components $\mathcal F$ completely
removed, the determining entries have the form
\begin{equation}
 \begin{split}
 N_j^{n,\mathcal F}
  &=\bigcup\mathcal F\ \cup
    \bigcup_{\substack{X\notin\mathcal F\\
            t_\flat<j\le t_\flat+n(X)}}(X\setminus Z''_j(X)),\\
 \Omega_j^{n,\mathcal F}&=\Omega_j\cup N_j^{n,\mathcal F},\qquad
 \Lambda_j^{n,\mathcal F}=\Lambda_j\cup N_j^{n,\mathcal F}.
 \end{split}
 \label{hi:stage-entries}
\end{equation}
The full prepared sequence uses every permitted stage.  In the
construction its consecutive determining boundaries have a collar of
at least one coarse $M$-cube.  This separation, and the joint admission
of all interpolated data to the background cores, are geometric and
analytic conditions on these actual entries.  Merely naming the entries
in \eqref{hi:stage-entries} does not prove either condition.

There is also an action sequence.  At a level redone by a drop-out it
lists what is still stored, rather than what has just been prepared:
\begin{equation}
 T\supset\Omega^{\mathsf a}_1\supset
 \Lambda^{\mathsf a+}_1\supset\Lambda^{\mathsf a}_1\supset
 \Omega^{\mathsf a}_2\supset\cdots\supset\Lambda^{\mathsf a}_s,
 \qquad Z_s^{\mathsf a}=Z_s.
 \label{hi:action-sequence}
\end{equation}
The extra symbol $\Lambda^{\mathsf a+}_j$ denotes the specified collared
support of the action weights.  Determining and action sequences can
coincide at recent levels and differ at older ones.  Equality at the last
level does not permit replacing one by the other everywhere.

The rules for reading these records are as follows.  Class membership,
the traced old cores and the old integration regions use the stored
label entries $Z_j$.  Backgrounds, their data, determining sets,
measures and gauge nests use the determining entries.  Density tests
and the last-shell test use the action regions.  Coupling and vacuum
weights use their specified action supports.  The integral operation
uses the retained operations and variables actually listed by the mark.
These rules fix five distinct uses of the same geometry.  They are part
of the meaning of a read, not five interchangeable descriptions of a
region.

\paragraph{A full re-entry in the broader waiting convention.}
This example allows a mark at the lower endpoint of the
preparation window. The dynamic class used below includes
that endpoint in its quarantine.
The reason for this separation is visible without an estimate.  Denote
the old central box by $P^\circ$, its old remaining core by $N^\circ$,
the old stored domain by the complement of $O^\circ$, and put
$W^\circ=O^\circ\setminus N^\circ$.  In a full re-entry the field on
$W^\circ$ is the current datum of the old level.  It is not a field of
the preceding level: those preceding-level variables have already been
integrated.  If the earlier failure preceded central integration, the
field on $N^\circ$ is supplied by its still-present averaging delta
function.  If the earlier failure occurred during or after central
integration, a bond touching $P^\circ$ can instead carry a free Haar
argument and a retained datum used by the determining sets.

Substituting the stored region at a determining place would in the
first case ask for an integrated-out variable on $W^\circ$.  Substituting
the determining region at a test place would in general assert the
presence of tests that the term no longer carries.  Both errors can
occur before any smallness estimate is attempted.  Conversely, an
argument that preserves the five lists above can check a read by a
finite inspection of its variable and test lists.  Its remaining
analytic question is admission of those data to the required cores.

\paragraph{Gauge factors at a retained seam.}
Each retained gauge factor includes a tag specifying its gauge
prescription and its actual background.  At the retained inward seam
one uses the relative prescription; away from that seam one uses the
prescription specified by the corresponding operation.  The group
element in such a factor must be evaluated by that tag.  An absolute
gauge factor at a field and a relative factor at a field and background
are different functions.  A bound for one cannot identify the other.
Removing every gauge factor on a bond with an endpoint in the next
determining region is a possible different convention, but using it
requires a new exact operation identity.  No such removal is made by an
estimate below.

\subsubsection{Local histories and their two argument spaces}

Write $\mathcal F=W_4$ for the finite group of gauge-tree frames and
$h=|\mathcal F|$.  A component that reaches the final decomposition,
and hence can be healed or made an island, carries an operation
\begin{equation}
 \mathbf T'_i(X)F(\mathbf p)
   ={1\over h}\sum_\omega
       \int_{S_\omega}e^{\Phi_\omega(\xi;\mathbf p)}
                           F\,d\mu_\omega(\xi).
 \label{hi:local-operation}
\end{equation}
The local history $\omega$ is a discrete record.  It lists the outer
frame, the outcomes of the decompositions in its outer genealogy,
and recursively a local history for each inner island.  The variable
$\xi$ lists the variables left free after every surviving averaging
delta function has been used as the definition
$V_{j+1}(c)=\mathcal M(V_j)(c)$.  Its reference measure is a product
of normalized Haar measures and Lebesgue measures.  Gaussian,
chart, elimination and gauge normalizers remain in the integrand;
they are not silently changed into probability measures of mass one.

The support $S_\omega$ records the tests actually carried by that
summand.  It is a Borel set when the finite background maps used to
define the tests are Borel.  At this point Borel dependence of actual
multi-level backgrounds is required; continuity of a formal expression
in a named background is not a substitute for that property.

The external argument is a pair $\mathbf p=(\mathbf U,\mathbf J)$
on $X^\sim$, where $X^\sim$ is the fixed coarse collar, not a
rectangular bounding box.  Its admitted complex core is denoted
$\mathcal U_i(X^\sim)$.  It is the core with the actual current
normalization and substitution rules of the construction.  An arbitrary
complex current array is not declared admissible by this notation.

To describe the splitting used in the main step, present a pair by
\begin{equation}
 \mathfrak z=(U,a',\mathbf J),\qquad
       \mathbf U=e^{ia'}U,\qquad U\in G^{\rm bonds}.
 \label{hi:triple}
\end{equation}
Here $a'$ is the dimensionless bond logarithm.  If a rescaled potential
$A'$ is used instead, its lattice spacing factor must be included in
$a'$.  A typical layer bound for a bare bond is
$|a'|<\hat\alpha_{1,j}L^{-j}$ and that for a bare plaquette is
$|\partial\mathbf U-1|<\hat\alpha_{0,j}L^{-2j}$.
The additional current and local-gauge conditions are retained in the
admitted core.  These displayed layer inequalities alone do not replace
them.

The complete exponent has the special decomposition
\begin{equation}
 \Phi_\omega(\xi;\mathbf p)
  =\Phi^\circ_\omega(\xi)
       -\tfrac12q_X(\xi;\mathfrak z)
       +V_X(\xi;\mathfrak z),\qquad q_X\ge0.
 \label{hi:tied-exponent}
\end{equation}
The inequality \(q_X\ge0\) here is on the real stored
integration fibre. The glued background defining $q_X$ is built at the $G$-valued factor
$U$, and is real.  This is the reason for its nonnegativity, including
when the represented pair is complex.  The two last summands need not
be functions of the pair separately.  Their sum is required to take
the same value at any two admitted triples representing the same pair,
and to be holomorphic in that pair.  Cauchy estimates apply to that
sum, not to an arbitrarily selected real quadratic summand.

\paragraph{Continuing the main-step quadratic form.}
Keep the real quadratic term in \eqref{hi:tied-exponent}. Its
continuation as a function of the pair can be constructed at the
admitted glued datum. Let $V$ be the real datum of the joined set
of $X$, using the prescribed inner rows and the exterior averages,
and let $V_\gamma=e^{i\gamma}V$. Write the continuation of its
critical representative in the form $\mathbf L_V(\gamma)^{u_\gamma}$,
where $\mathbf L_V(\gamma)$ is in the Landau chart about the real
critical representative of $V$. Use the on-shell pair
$(\mathbf L,J(\mathbf L))$, $\mathbf L=\mathbf L_V(\gamma)$.
Let $N_{V_\gamma}$ be the differential at zero of the map from
raw row data to reduced data at $V_\gamma$. It is the identity on
interior rows. At an outward row $c$ it adds
$\operatorname{Ad}(V_\gamma(c))\dot v_{c_+}(f)$ when $c_+$ is the
outer endpoint and subtracts $\dot v_{c_-}(f)$ when $c_-$ is the
outer endpoint. Here $i\dot v_z(f)$ is the first derivative at zero
of the one-step node factor at $z$ with datum direction $f$ on its
block. For a raw datum direction $f$, put
$f^\sharp=\operatorname{Ad}(u_\gamma^{-1})N_{V_\gamma}f$, with
the node gauge acting at the initial endpoint of each row.

With the cutoff $\zeta_X^{\rm co}$ of the main step and the
walk expansion $H_1=\sum_\omega H_{1,\omega}$ at this pair, define
\[
 t_{\omega\omega'}(\gamma;V)[f]
   =\sum_p\eta^4\zeta_X^{\rm co}(p)
      \big\langle(D_{\mathbf L}H_{1,\omega}f^\sharp)(p),
                 (D_{\mathbf L}H_{1,\omega'}f^\sharp)(p)\big\rangle .
\]
Here $D_{\mathbf L}$ is the covariant exterior difference and the
pairing is bilinear. The ordered pair of walks is counted once.
This uses the fixed regular-pair radii of the walk expansion. If
$a_{\rm D}$ is their ceiling and $K_{\rm S}$ the profile-to-response
constant, choose the real regularity ceiling $\varepsilon_{\rm w}$
and the datum radius $c_{\rm w}$ so that
\[
 C(L,\varrho_0)M_1\varepsilon_{\rm w}\le a_{\rm D}/4,\qquad
 C K_{\rm S}c_{\rm w}\le a_{\rm D}/4,\qquad
 Cc_{\rm w}\le\varepsilon_{\rm w}/2.
\]
The layer gaps of the joined set must give the weighted summation
for these walks. In particular the exponential decay after the
spacing comparison must absorb $L^{d/\varrho_0}$ at distance $d$.
On $|\gamma|_\infty<c_{\rm w}$, the profile and response bounds
admit the pair on the enlarged walk carriers, and the node gauge
and reduced datum give $|f^\sharp|_\infty\le C|f|_\infty$.

At a plaquette based in a cell of spacing $\xi_y$, the walk
estimate gives parameter-independent weights for the
differentiated lifts. Summing their termwise suprema
with these weights gives
$C\xi_y^{-2}|f|_\infty$. Its square costs $C\xi_y^{-4}|f|_\infty^2$,
while the cell has at most $6(\xi_y/\eta)^4$ plaquettes based in it.
Thus, for the $\mathsf n$ cells carrying the cutoff, the sum of the
suprema of the ordered-pair terms is at most
$K_{\rm w}\mathsf n|f|_\infty^2$. The sum $q_{0,X}^{\rm w}$ is
holomorphic on the datum disc. At a real datum the covariance of
each walk term removes the gauge $u_\gamma$ and identifies the
sum with the real main-step form there. This uses the differential
of the physical response, namely $H_1N_{V_\gamma}$ in the reduced
row coordinates. Recenter at $e^{ib}V$ for real
$|b|_\infty<c_{\rm w}/4$ and $|\gamma|_\infty<c_{\rm w}/2$.
The two formulas agree at real data, so their holomorphic
continuations agree on the common datum disc after the logarithmic
change of center. This uses the same glued datum in both splittings,
including the node and inner-row conventions.

For an actual triple, the real glued datum must be
$\varepsilon_{\rm w}/4$-regular, and its increment from the exterior
pair must satisfy $|\gamma|_\infty\le C\alpha_{1,t}<c_{\rm w}/4$.
Schwarz's lemma on the complex line through $\gamma$ gives
\[
 \left|\tfrac12\theta_t(q_{0,X}^{\rm w}-q_{0,X})(g_tB)\right|
 \le \frac{C K_{\rm w}}{c_{\rm w}}\,
       \theta_t\mathsf n(\delta'_t)^2\alpha_{1,t} .
\]
Here $|g_tB|_\infty<\delta'_t$ and $q_X=\theta_t q_{0,X}(g_tB)$
in the real presentation. The coupling ceiling imposes
$C\max\{2\alpha_{0,t},\varepsilon_1\}\le\varepsilon_{\rm w}/4$,
where $\varepsilon_1$ bounds the input datum tests, and makes the
last bound at most $1$. The absolute sum in the exponent is at
most $\tfrac12K_{\rm w}\theta_t\mathsf n(\delta'_t)^2$.
The same ceiling requires this to be at most
$(\log\theta_t)^{2\mathfrak p_0}/10$ using
$\mathsf n\le C\check R_t^5$. If the analytic form is separated
from the exponent, the corresponding remainder is
\(V_X+\tfrac12\theta_t(q_{0,X}^{\rm w}-q_{0,X})(g_tB)\).
For the modulus
bound keep the original nonnegative real term. This comparison
concerns the form of the main step at its glued datum. The older
conditioned forms carried by the history have separate free
coordinates, row restrictions and elimination factors.

An exact consequence worth isolating is
\begin{equation}
 \sup_{\mathbf p\in\mathcal U_i(X^\sim)}
         |e^{\Phi_\omega(\xi;\mathbf p)}|
   \le e^{\Phi^\circ_\omega(\xi)+C_VN_0(\omega)}
 \quad\hbox{if}\quad
 \sup_{\mathfrak z}|V_X(\xi;\mathfrak z)|\le C_VN_0(\omega).
 \label{hi:pair-sup-envelope}
\end{equation}
All representatives in the supremum must obey the same bound.  The
function $N_0$ counts the stored sites at their own levels, with the
fixed finite multiplicities belonging to the operation.  Proving the
complete localized bound on $V_X$ is an analytic formation task.  Once
it is proved, \eqref{hi:pair-sup-envelope} is a pointwise envelope before integration; it is
stronger information than a pointwise polymer logarithm.

\paragraph{The retained inner factors in the complete exponent.}
In the fixed real reference system, take a component $X$ of the
main step at level $t$ and write its kept exponent in
\eqref{hi:tied-exponent} as $\sum_R T_R^1$. Its carriers are in
the enlargements of older islands or of strips of levels at
most $h_t=t-\mathsf N_t$, where $\mathsf N_t$ is the number of
preparatory stages. The stored bonds are in the inner core.
The enlargement costs at most one cell of level $h_t$, and the
next two shells have widths at least $M\check R_t\xi_{h_t}$
and $LM\check R_t\xi_{h_t}$. The scale comparison
$\check R_{h_t}\le L\check R_t$ therefore puts every carrier
inside $Z''_{h_t+2}$. Write $l(x)$ for the layer index of the
prepared sequence and put $n_R=\max_{x\in R}l(x)$ and
$u_R=(h_t-1-n_R)_+$. In particular $n_R\le h_t+1$.

Let $U_\star$ be the adapted real background of the datum
$\mathbf V_\star$ on the component determining set. This datum
equals $V_X^{\rm in}$ on the inner rows and the identity on
the other rows. Its regularity follows from the kept tests:
the half-threshold test at the top of the prepared nest and
the tests inside its last preparatory operation give the
levelwise bounds on the holonomies of $(\mathbf1,V_X^{\rm in})$.
The real data $V_X^{\rm in}$ remain in that range while the
box variable $B$ varies on $|g_tB|<\delta'_t$.
Write $q(\xi)$ for the pair of $U_\star$ and use the pinned
representative at the stored sites. For every kept factor
$\varphi_R$, use its original target space, with a bound
$b_R^\sharp$ on the admitted spectator range. In the reference
evaluation the spectator values are held fixed at their stored
values. The factors are invariant, have the tie on real triples
with the same pair and have their data at the prescribed adapted
sites. Their structures precede the prepared structure.

To form the new argument, apply the exterior row map
$E_X^\natural$ to the decoupled response to the formed belt
datum $\mathfrak b^Y(s;P,\varsigma[P;\xi])$. Here $P$ is in the
space of the output structure on the localization $Y$ and
$|s(\Delta)|<e^{\bar\kappa}$. On $X^\sim$ the output space
has the single level $t$, so the endpoint blocks of the tower
rows lie in its one-scale core. For a tower row $c$, write
$T_c(W)$ for the average of $W$ transported along its two
combs to the tops $y_\pm$. Put $U^+=e^{i\eta\mathbf A}\mathbf U$
for the exterior shift, $v_\pm=v[P;\mathbf A](y_\pm)$ for its
site factors and $u_-=\mathbf U(\mathrm{cb}_{c_-})$. Then
\[
 e^{i\gamma(c)}=u_-^{-1}v_-^{-1}T_c(U^+)v_+
                       T_c(\mathbf U)^{-1}u_- .
\]
The set read by this product consists of the two combs, the
blocks of their tops and the double block of $c$. It is at
distance at least $cM\check R_t$ from the belt. If $\phi$ is
the supremum of $|\mathbf A|$ on this set, the product with
$z\mathbf A$ is holomorphic and within $1/2$ of the identity
for $|z|<r_{\rm E}/\phi$, by the site factor and relative
average bounds. The product is the identity at zero.
Schwarz's lemma bounds its logarithm at $z=1$ by
$\phi/r_{\rm E}$. The response bound for $\phi$ therefore gives
$|\gamma|\le Ce^{C\bar\kappa}\varepsilon_\beta
e^{-c\delta_0M\ell_t^r}$, where $\ell_t=\log\theta_t$.
The logarithm and the composed maps are holomorphic at fixed
real factor throughout the given cube polydisc.

Form $\mathbf V$ on the component
rows using $e^{i\gamma}\mathcal M^m(\mathbf U)$ on the tower
rows of level $m$, $\mathcal M^t(\mathbf U)$ on the remaining
exterior rows, the formed box averages with their left factor
$e^{ig_tB}$ on the box rows, and $V_X^{\rm in}$ on the inner
rows. Put
\[
 B^\natural=\frac1i\log\bigl(\mathbf V^{\hat w}
                                      \mathbf V_\star^{-1}\bigr),
\]
where $\hat w$ is the comb of $\mathbf U$ at the exterior nodes
and the identity at the nodes of the core. It acts on the datum
at the nodes. The low box rows have both ends in the core and
there
\[
 e^{iB^\natural(c)}=e^{ig_tB(c)}e^{i\nu_P(c)},\qquad
 |\nu_P(c)|\le C_1'L^{m-t}a_C,
 \qquad a_C=\max_\iota\alpha_{\iota,t}.
\]
The factor $e^{i\nu_P}$ is the continuation of the box average
in the exterior datum. At an entrance row the closed contour
through the gauge tree and the exterior path has $O(M^2)$
plaquettes. Its curvature bound and the exterior comb bound
give the bound $\epsilon_{\rm out}=C_1'M^2\check R_t a_C$ on
the remaining rows, including the small complex correction to
the comb. On the inner rows $B^\natural=0$. Thus the whole
formed datum is in the continuation ball $\|B^\natural\|_\infty<r_{\rm B}$.

Use the Landau response $(\mathcal H_\star,\mathsf j_\star)$
at $\mathbf V_\star$, transported to the adapted anchor. Its
current increment is
$J(e^{i\eta\mathcal H_\star(B^\natural)}U_\star)-J(U_\star)$.
On the bonds of
$R$ define the actual argument by
\[
 \mathfrak S_R(\lambda;\mathfrak z,\xi)
 =\mathfrak p[U_\star]|_R\dotplus
    (\lambda\mathcal H_\star(B^\natural),
     \lambda\mathsf j_\star(B^\natural))^\delta|_R .
\]
The superscript denotes levelling with all the stored sites
pinned. Both increments are multiplied before levelling.
The normalized potential, first difference and current of the
unlevelled response are bounded by a constant times the datum
profile. The current bound follows from the continued current
difference and the deviation of the response gauge from the
identity, since the scaled current is bounded in the chart.

For $h_t\le m\le t$ let $D_m(x)$ be the prepared weighted
distance to the sites outside the innermost region $I_2$ with
layer index at least $m-1$. Convolving the row bounds with
the decaying response gives the slowly varying profile
\[
 s_\star(x)=C\left[\delta'_t e^{-\theta_\flat D_{h_t}(x)}
 +a_C\sum_{m=h_t}^{t-1}L^{m-t}e^{-\theta_\flat D_m(x)}
 +\epsilon_{\rm out}e^{-\theta_\flat D_t(x)}\right].
\]
Separation gives $D_m(x)\ge\varrho_0(m-2-l(x))_+$.
Sum the middle series using $Le^{-\theta_\flat\varrho_0}\le1/2$.
Since $\mathsf N_t\gtrsim\ell_t^r$, both
$a_CL^{3-\mathsf N_t}$ and
$\epsilon_{\rm out}e^{-\theta_\flat\varrho_0(\mathsf N_t-3)}$
are smaller than a fixed multiple of $\delta'_t$ at small
coupling. Thus for $l(x)\le h_t+1$ the profile is at most
$C\delta'_t e^{-\theta_\flat\varrho_0(h_t-1-l(x))_+}$.

The anchor has half the radius of each target clause. Its
scaled curvature and current follow from the profile of the
regular datum. Local gauges follow on the required cubes from
the curvature bound, using one finer layer for cubes touching
the boundary. On a subordinate tower the restricted anchor
is critical over its own averages. The tower therefore has the
same curvature and current up to gauge. This also applies to
the towers of the single-scale cores after truncation. These
checks use separation of the component sequence, which can
have holes near its exterior boundary.

The available shift is $\kappa_{\rm R}\min_\iota\alpha_{\iota,l(x)}$
in prepared units for a slowly varying profile. For a target
index $m\le l(x)$ the shift gains at least $L^{-(l(x)-m)}$ in
its own units. The radii change by at most a factor
$C(1+l(x)-m)^{1/2}$, so the rescaled profile spends only a fixed
small part of the pointwise radii. The gauges of the unchanged
real factor still apply. For each subordinate tower the local
Lipschitz bound on its curvature and current takes the supremum
of the rescaled profile with an exponential distance penalty.
The tower decay rate exceeds
$\theta_\flat+2\log L/\varrho_0$, which absorbs the variation of
both layer indices and of the profile. The same bounds apply to
the single-scale cores and their towers. This gives the stated
allowance for the pinned shift, including the fixed cost of
levelling. The allowance also puts its local potential below
the levelling threshold on the neighboring tops.

Finally, the running radii satisfy
$\min_\iota\alpha_{\iota,l}\ge
c(1+t-l)^{-1/2}\min_\iota\alpha_{\iota,t}$, and
$1+t-l\le(2+\mathsf N_t)(1+(h_t-1-l)_+)$. The ratio
$\delta'_t/\min_\iota\alpha_{\iota,t}$ is a constant times
$\ell_t^{\mathfrak p_1-q}$, where $q$ is the radius exponent.
Absorbing the remaining square root into the layer exponential
shows that the normalized pinned shift on $R$ costs at most
$|\lambda|a_R/2$ of the allowance, with
\[
 a_R=C_*\ell_t^{-(q-\mathfrak p_1-r/2)}
         e^{-\delta_\flat\varrho_0u_R}
    \le C_*\ell_t^{-(q-\mathfrak p_1-r/2)}(4L^8)^{-u_R}.
\]
In particular the argument stays in the target space for
$|\lambda|<2/a_R$ uniformly on the cube polydisc and the
admitted spectator range. It is holomorphic with the real
factors held fixed. It reads $q(\xi)|_R$
at zero. At the real full parameters the belt and the box
averages give the actual prepared datum. The pinned response
has its current and the prescribed transports on the stored
sites. The adapted representative on $R$ can be taken from
the prepared orbit because its pillars inside the core are
the same as the component pillars. Its plain triple satisfies
the target bounds by the real profile comparison, and the tie
identifies its value with the pinned value. Thus
$T_R^1=\varphi_R(\mathfrak S_R(1;\mathfrak z,\xi))$.

Set $\delta_R=T_R^1-\varphi_R(q(\xi)|_R)$.  Schwarz's lemma, applied
to the difference along this substitution disc, gives
$\sup|\delta_R|\le a_Rb_R^\sharp$. 

\hypertarget{hi:formed-collar-majorant}{}
For the absolute bound of a kept member, split off the part of
each island factor that depends only on its stored variables.
The remaining factor vanishes at the trivial triple, for
each value of those variables, so its supremum is bounded by
its oscillation. A strip factor consists of the quadratic
walk sum and the logarithms of its external pivot Jacobians.
On the enlarged spectator discs the quadratic sum costs the
square of the radius multiplier $(\log\theta_{j-1})^{6r+1}$
at the birth level $j$. For each pivot, the full logarithm
has modulus bounded by $9\log L+c_{\rm P}$ on the pivot
chart. In the convention $\log z^{(j-1)}=-\mathsf p_{c'}+9\log L$
the second term is the scalar from the pivot normalization.
Its carrier has one or two adjacent cubes
and there are at most $8M^4$ such pivots per carrier. These
bounds give, for each level-$j$ carrier,
\[
 b_R^\sharp\le B_\Phi^\sharp\widehat\ell_j^{\mathfrak p^\sharp}
                      e^{-\kappa d_j^{\rm in}(R)},\qquad
 \mathfrak p^\sharp=2\mathfrak p_0+12r+2,
 \qquad \widehat\ell_j\le\ell_t(1+\log(2+t-j)).
\]
Here $\widehat\ell_j=\max\{\log\theta_j,\log\theta_{j-1}\}$
and $B_\Phi^\sharp$ absorbs the strip and island constants.
If the member is divided into identical copies, the bounds
are divided by the number of copies, so their sum is still
bounded by $b_R^\sharp$.

Write $z_j(X)$ for the number of side-$M$ cubes in level-$j$
units in the action large-field region of $X$, and put $Q=4L^8$. The levels
$j\ge h_t-1$ contain at most $C\ell_t^{5r}$ such cubes in
total. A kept carrier has a cube in the enlarged action
large-field region of its level. Counting through that cube
with the entropy bound, including the bounded multiplicity
of exempt island hulls, controls the contribution from these
levels by $CB_\Phi^\sharp\ell_t^{5r+\mathfrak p^\sharp}$,
since the bounded frame increments give $\theta_{j-1}\le
\theta_t+C(\mathsf N_t+2)$ throughout this window.

For $p=h_t-j\ge2$, the possible rooting cubes are in the
level-$(j+1)$ action large-field region. It is contained in
the inner region with at most $C\ell_t^{4r}$ cubes of level
$h_t$, so the rooting count at level $j$ costs at most
$C\ell_t^{4r}L^{4p}$. The root lies outside a member of the
prepared sequence with index $i'+1\le h_t$. Here $i'$ is
$j$, $j+1$, or the level $i<h_t$ of a drop-out having
$i-N_0(i)\le j<i$, where $N_0(i)$ is the number in the
old drop-out index $k_0(i)=i-N_0(i)$.
If $n_R\ge i'+2$, the carrier must cross a prepared gap
and $d_j^{\rm in}(R)\ge\tfrac12 L^{n_R-j}$. The exempt
hulls stay on the lower side of the gap and can reduce
the length by at most two level-$j$ cube units. Summing
over $m=n_R-j$ in this case costs, after entropy,
\[
 \sum_{m\ge2}e^{-c L^m}Q^{-(p-1-m)_+}
 \le C Q^{-(p-1)}.
\]
Otherwise, in the two unmarked cases $n_R\le j+2$ and
$u_R\ge(p-3)_+$. In the drop-out case
$u_R\ge(p-2-N_0(i))_+$, so the extra cost is at most
$Q^{N_0(i)}\le L^9\rho_i^9$, where $\rho_i=\check R_{j'}$
at $j'=i-N_0(i)+1$. The frame bound
$\theta_{j'}\le\theta_t+C(t-j')$ and
$N_0(i)\le2+r\log_L\log\theta_{j'}$ give
$t-j'\le2(t-j)$ at small coupling, since
$t-j\ge\mathsf N_t+2$ absorbs the logarithmic delay.
Consequently
\[
 \rho_i\le C[\ell_t(1+\log(1+\mathsf N_t))
                         (1+\log(1+p))]^r .
\]
In all three cases the geometric factor from $L^{4p}$
and $Q^{-p}$ sums against every fixed power of
$1+\log(1+p)$. Since $\mathsf N_t\le C\ell_t^r$, summation
gives
\[
 \sum_R b_R^\sharp(4L^8)^{-u_R}
 \le K_A' B_\Phi^\sharp\ell_t^{\mathfrak p^\sharp+13r}
              (1+\log\ell_t)^{9r+\mathfrak p^\sharp}.
\]

It follows that
\[
 \sum_R\sup|\delta_R|
 \le C_*K_A'B_\Phi^\sharp\ell_t^{5r+E+12r+2}
              (1+\log\ell_t)^{9r+\mathfrak p^\sharp},
 \qquad E=2\mathfrak p_0+\mathfrak p_1+\tfrac{17}{2}r-q.
\]
With $E+12r+2<0$, the coupling ceiling and the geometric bound
$N_0(\omega)\ge2^{-5}M^4\ell_t^{5r}$ make this at most
$N_0(\omega)$.  All suprema here are over the common admitted range.
Use \(b_R^\sharp\) on this range, including the enlarged
discs of the one-step quadratic factors.

Let $V_X^{\rm col}$ be the collar sum before the change of
inner reads. It consists of the localization of the action
difference, the small Wilson sum $\mathfrak v(X^\sim)$ and
the Wilson weight correction $\mathsf W_X$. Use the count
$N_0(\omega)=M^4\sum_{j\le t}z_j(X)$ from the action regions.

In the action difference, cancel each carried trunk with its
partner on the new side before estimating the localized sum.
The localized differences of a carried entry of size $b_e$
and amplitude $a_e$ have total grouped bound $a_e b_e$ on a
collar. Indeed each difference over switched cubes is bounded
by $a_e b_e$ times exponential decay in the number of cubes
needed to connect its support to its base. Summing the
connected extensions uses the cube entropy and the chosen
decoupling radius. The same estimate for the nontrunk
pieces of an unpaired entry can be absorbed into its bound
$Cb_e$, since $a_e\le1$. An unpaired entry is counted with
its whole trunk, and a plain entry with its whole value.
The kept trunks have the separate estimate just proved.

At level $t$ there are at most six ordinary entries on a
domain and at most $3^4z_t(X)$ collar cubes through which to
root them. For $j<t$ the index condition on an ordinary
entry gives a cube in a one-layer enlargement of the
action large-field region of level $j$. A root of a domain
in this collar belongs to the enlargement of $X$'s own
trace, by separation of the class components. There are at
most $3^4z_j(X)$ such cubes. The level-$j$ island hulls
in the collar have at most $3^4\mathsf m z_j(X)$ cubes in
total, where $\mathsf m$ bounds their multiplicity. Counting
domains with their length decay therefore costs at most
$C(1+\bar c)z_j(X)$ for the ordinary entries. This count
includes the unpaired plain numbers of the gained domains,
namely the pointed values and the remainder values at the
trivial pair. The large quadratic strip terms of level $t$
are carried outside the collars. The older strip factors
with base in the collar occur among the kept entries.

For the regular members, let $m$ be their level and $n$
the last action layer that their carrier meets. A member
with $n<t-1$ is rooted in the action complement of level
$n+1$. That complement in $X$ offers at most
$L^{4(n+1-m)}z_{n+1}(X)$ possible level-$m$ roots. The
carrier has length at least $2L^{n-m-3}-2$, so after paying
the domain entropy and one power of the length for the
member norm, the remaining decay sums over $m\le n$
against the root count. The last layer uses the collar
cubes as roots instead. This gives another
$C(1+\bar c)\sum_j z_j(X)$. Thus the action difference
has the required absolute majorant, with every grouped
piece, unmatched trunk and unpaired plain number counted.

In the small Wilson sum, the unpaired functions are the
local correction (including the analytic reference
normalizer) and the interior Wilson comparison with the
actual exterior row increment. Their combined supremum
is at most $C(1+\bar c)\ell_t^{5r}$. The two unsuppressed
Wilson differences are supported in at most $C\check R_t^4$
and $C\check R_t^5$ unit cells, respectively. The difference
of weights times the squared field bound is $O(1+\bar c)$
on these cells. This follows from the tube width
$C\theta_t\ell_t^{-2q_{\rm t}}$, $q_{\rm t}\ge q$, and
the field size $C\theta_t^{-1/2}\ell_t^q$ on the prepared
window. For the localized Wilson forms, the regularity
of the real inner data is bounded by $C\ell_t^r\varepsilon_t$
on the retained tests, as required by the box construction.
The other comparisons have small factors
$e^{-c\delta_0M\check R_t}$ or $L^{-\mathsf N_t}$, or the
Taylor loss $\theta_t^{-1/2}$, which absorb the polynomial
logarithmic factors. For the comparison in the deep inner
region the cells are summed with their weighted-distance
decay. The layer gap absorbs the possible increase of the
weight there. The analytic normalizer is uniformly bounded
on its box. These estimates are on the complex increment
ball used in the formed Wilson comparison.

The remaining Wilson entries are matched exterior-cell
terms. On this collar their total localized difference is
bounded by $C(1+\bar c)\check R_t^4\ell_t^{2q+2r}$ times
the pinned exterior response amplitude, and that amplitude
is bounded by $C\theta_t^{1/2}e^{-c\delta_0M\check R_t}$.
This term is bounded by $C(1+\bar c)\ell_t^{5r}$ at small
coupling. There is at least one level-$j$ label cell in the trace
for each $h_t\le j\le t$. Each such cell contains
$\check R_j^4\ge 2^{-4}\ell_t^{4r}$ side-$M$ cubes, and
there are $\mathsf N_t+1\ge \ell_t^r/2$ levels. Hence
$\sum_jz_j(X)\ge2^{-5}\ell_t^{5r}$. Thus the small Wilson
sum is bounded by $C(1+\bar c)\sum_jz_j(X)$.

Finally, $\mathsf W_X$ is the sum over plaquettes in $X$
of the Wilson density times the difference between the
new weight and $\theta_t$. The density is bounded by
$C\eta_t^4\alpha_{0,t}^2$, and each level-$t$ cube has
at most $CM^4\eta_t^{-4}$ such plaquettes. The correction
therefore costs $CM^4z_t(X)\alpha_{0,t}^2
(|x_{t-1}-\theta_t|+\bar c)$. The adjacent-level flow and
the tube bound make the last product of the radius squared
and the weight difference $O(1+\bar c)$. This estimate
holds after splitting the correction into its local
plaquette carriers as well.

Adding the three majorants gives
$\sum_\lambda\sup|v_\lambda|\le C_{\rm col}(1+\bar c)N_0(\omega)$
for the collected collar terms on their common range. The
separate Gaussian counting constant keeps its assignment in
the reference normalization. The exact assignment is
\[
 V_X=V_X^{\rm col}+\sum_R\delta_R,
 \qquad
 \Phi^\circ_\omega=\Phi^{\rm rest}_\omega
                         +\sum_R\varphi_R(q(\xi)|_R).
\]
Here $\Phi^{\rm rest}_\omega$ contains the retained inner Wilson
factor and all other terms depending on $\xi$ alone.  This preserves
the exponent and proves $|V_X|\le[C_{\rm col}(1+\bar c)+1]N_0$.
The pair-dependent part of the exponent remains
\(-q_X/2+V_X\), with its given holomorphy and independence
from the representation of the pair. The bound gives
\eqref{hi:pair-sup-envelope}. The integral of
\(e^{\Phi^\circ_\omega}\) uses the real coordinate estimates below.

%
%
%
\paragraph{The pair reader on a formed island.}
\hypertarget{hi:formed-island-pair-read}{}
On the output core of the formed island kit, write
\(f=g_tB\) and \(\zeta_X=\zeta_{0,X}-\zeta_{1,X}\).
Use the generated Wilson configuration \(W^{\rm f}(b,f)\)
of the prepared component. Its datum is the lift of the
outer average of the pair, the perturbed collar datum
\(e^{if}V_0[d;V^{\rm e}]\), and the retained inner datum
\(V_X^{\rm in}\). At fixed real stored variables these
give the same complex datum for all splittings of the pair.
On the real slice the glued representatives are related
by a gauge equal to one at the stored pins. The Wilson
trace is invariant there, and the tied outer factors also
have equal values there. On the two prepared leaves use a
complex parameter \(z\) taking the real interpolations to
the desired splitting at \(z=i\). For each real collar
parameter the two holomorphic reads agree for real \(z\),
so they agree at \(i\) by the identity theorem on the
neighborhood of \([0,i]\). Continuing in the collar
parameters gives equality on the common formed core.
The prepared leaves use the same pinned data and margins.

Define the weight-free quadratic form
\[
 q^\natural_{0,X}(f;P)=\left.\frac{d^2}{ds^2}\right|_{s=0}
                 A(\zeta_X,W^{\rm f}(b,sf)).
\]
It reads the pair \(P\) at fixed stored variables. No entry
of the source system enters it. The holomorphic remainder
\[
 T_6(b,f)=-\theta_t\left[
  A(\zeta_X,W^{\rm f}(b,f))-A(\zeta_X,W^{\rm f}(b,0))
  -\tfrac12q_{0,X}(f)\right]
\]
is bounded by one on the admitted radius \(5\delta'_t\)
under the prepared Wilson estimate and coupling ceiling.
For \(\|f\|_\infty\le\delta'_t\), its second coefficient
along \(sf\), bounded by Cauchy on \(|s|=1\), gives
\[
 \left|\tfrac12\theta_t
             (q^\natural_{0,X}-q_{0,X})(f)\right|\le1 .
\]
Before transferring the kept trunks to \(q(\xi)\), denote
the collar interaction by \(V_X^{\rm col}\). The new pair
reader is
\[
 V_X^\natural=V_X^{\rm col}+\sum_R\delta_R+
             \tfrac12\theta_t(q^\natural_{0,X}-q_{0,X})(f).
\]
Thus \(-\theta_t q^\natural_{0,X}(f)/2+V_X^\natural\)
plus the trunks read at \(q(\xi)\) gives exactly the old
exponent. The formed-core bounds use the carrier \(X^\sim\),
which already contains the prepared box and the gauge tree.
They apply on the output island core to the individual pieces
with the stored variables fixed. In the sourced-minus-zero
difference the common quadratic correction cancels, leaving
the Wilson source on the inner field and
\(V_X^{\natural,w}-V_X^{\natural,0}\), with the inherited
exponents read at their admitted real pairs. This is
\eqref{hi:actual-source-recursion} on that core. The absolute
cost and the source oscillation use its carrier bounds.

%
%
\paragraph{The factor bound at the birth of an island.}
On the admitted formed core, split each
collar member into its value at the fixed
anchor and the change from that value,
using the prescribed carrier of the member.
In the old core the anchor reads the stored
variables alone. File its value in
\(\Phi^\circ_\omega\). The oscillation of
the change costs at most twice the
amplitude on the anchored substitution
disc. The factor \((4L^8)^{-u_R}\) at
depth \(u_R\) pays for the \(L^{4u_R}\)
root count and its fixed multiplicities.
The additional negative power of
\(\ell_t=\log\theta_t\) absorbs the
inherited sizes. Outside the old core,
the waiting window contains at most
\(\ell_t^{C r}\) local cells, counted at
their own levels. The regular terms use
the pointed sums with their two
subtractions, and the Wilson terms use
their net weights. The collar
oscillations therefore use the
single-term collar bound with constant
\(\mathsf b_V\), independently of the
inherited oscillation constant.

For the kept contributions use the
bound on \(\sum_R|\delta_R|\) on the
common formed range. Their total
oscillation is at most twice that
bound. Its logarithmic factor, by the
choice of \(q\), makes the cost at most
\(\mathsf B_\Phi\ell_t^{2\mathfrak p_0}/4\)
when the coupling ceiling is chosen
after the strip and island constants.
The new quadratic and its displayed
correction cost at most
\(C A_1^2\check R_t^5\ell_t^{2\mathfrak p_1}+C\).
Use the prescribed tied expression for
these pieces. The hierarchy between
\(\mathfrak p_0\), \(\mathfrak p_1\) and
the waiting exponent makes their
oscillation at most another quarter.
Use the whole carrier \(X^\sim\) for the
quadratic and the retained changes
formed through the island's full
reader. That carrier is exempt in the
output label and has hull length zero.
The smaller carriers of the local collar
members keep their own decay.

Choose \(\mathsf B_\Phi\ge2\mathsf b_V\)
and keep the preceding quarters as
reserves. The new island then has the
same real oscillation bound per
carrier as the incoming islands.
For an island that stands, its exponent,
stored variables and carriers remain those
of its birth operation. Restricting its
arguments preserves the old oscillation
bound. These two cases return a maximum
over the formed and carried entries.
The bound applies on the formed cores
and uses the admission and mass bounds
of the actual histories.

\paragraph{Why the old Wilson factor is retained.}
The main action expansion is made only on the small-field part
$\zeta_0(1-\zeta_1)$, about the real glued background.  On the old
core the factor
\begin{equation}
 \exp\{-A(x''\zeta_1\mathbf1_X,U''(\xi))\}
 \label{hi:inner-wilson}
\end{equation}
is left intact.  The inner background $U''(\xi)$ is real and is formed
from the prepared determining sequence, with its old innermost datum
and its new fluctuation datum.  It is distinct from the glued
background used in $q_X$.  The factor in \eqref{hi:inner-wilson} is independent of the
external pair.  Its plaquettes include the old failed tests that have
survived in the record.  Replacing it by a quadratic Taylor expression
throughout the old core would remove the positive reserve used below.

The other factors in $e^{\Phi^\circ}$ are chart densities, constants,
kept gauge and elimination factors, stored Gaussian weights and the
inner operations evaluated at real pairs.  Their normalization and
their domains will be used individually.  The phrase ``the complete
exponent'' always includes all of them.

The restored law is combined with the old Wilson action before
estimating its normalization.  Fix one frame and all exterior
data of a component $X$ at level $i$.  Write $s$ for its retained
old variables and $v$ for its restored box field, with the tree
condition imposed.  Set
\[
 \begin{split}
 S_{\rm old}(s)&=A(x'',U''_i(s)),\\
 S_{\rm ref}(v)&=\theta_i A(\zeta_0,U_{i,X}(v)),\qquad
 S_0=\theta_i A(\zeta_0,U_0).
 \end{split}
\]
The denominator has the form \(D_X^{-1}=e^{S_0+I_X-E_X}\),
where \(I_X\) is the negative logarithm of the centered
reference integral after extracting its counting factor
\(e^{E_X}\). To combine its center with the Wilson actions,
write \(w=x''\), \(\theta=\theta_i\), \(W''=U''_i\),
\(W^0=U_0\), \(W^Z=U_{i,X}(v)\), and let \(\bar W\)
be the completed background. In a succession of completed
components these are the backgrounds before and after the
current replacement. Use \(w_{\rm fr}=\theta\) on the
component and \(w_{\rm fr}=w\) off it. The replacements
telescope over the components.

%
The three parts of the old weight use \(1-\zeta_0\),
\(\zeta=\zeta_0-\zeta_1\), and \(\zeta_1\), with
\(\zeta_0=1\) on the support of \(\zeta_1\). Write
\(\Delta_a(P,Q)=A(a,P)-A(a,Q)\) for weighted action
differences. The exterior part of the remainder is
\[
 \Delta_{w(1-\zeta_0)}(\bar W,W'')
 + A((w_{\rm fr}-w)(1-\zeta_0),\bar W)
 + \theta\,\Delta_{\zeta_0}(\bar W,W^Z).
\]
On the middle part put \(W_{\rm g}^0=W^{\rm f}(b,0)\)
and \(W_{\rm g}^f=W^{\rm f}(b,f)\), the two glued
backgrounds with the same exterior increment \(b\) and
with fluctuation increments \(0\) and \(f=g_iB\).
Its remainder, including the center of the reference
law, is
\[
\begin{aligned}
 &\theta A(\zeta_1,W^0)
 +\theta\,\Delta_\zeta(W^0,W_{\rm g}^0)
 -\bigl(\theta\,\Delta_\zeta(W_{\rm g}^f,W_{\rm g}^0)-q_X/2\bigr)\\
 &\hspace{2em}
 +\theta\,\Delta_\zeta(W_{\rm g}^f,W'')
 -A((w-\theta)\zeta,W'').
\end{aligned}
\]
The bracket contains the Taylor remainder after the
quadratic form \(q_X/2\) is extracted. The inner part
keeps \(S_{\rm core}=A(w\zeta_1,W^{\rm in}(f))\).
For the difference on this part insert the inner
background \(W^{\rm in}(\gamma^\flat)\). Its datum on
the rows adjacent to the core is \(e^{if}e^{iv_0}\)
relative to the inner center, where \(v_0\) is the
continued basic datum with
\(\lvert v_0\rvert\le C L^{-\mathsf N_i}\mathsf s\).
Here \(\mathsf s\) bounds the fields and increments.
On the remaining rows, including the joining rows,
keep \(f\). The two errors are
\[
 \Delta_{w\zeta_1}(W^{\rm in}(f),W^{\rm in}(\gamma^\flat)),
 \qquad
 \Delta_{w\zeta_1}(W^{\rm in}(\gamma^\flat),W'').
\]
The graded row response in the first and the screened
background comparison in the second bound them by
\(C\|w\|_\sharp\mathsf N_i\mathsf s^2\mathsf n\)
times \(L^{-\mathsf N_i}\) and
\(e^{-c\delta_0 M\check R_i}\), respectively.
Here \(\mathsf n\) bounds the number of remade pillars
and \(\|w\|_\sharp\) divides the weight on layer \(j\)
by \(1+i-j\). This places the response to the complex
outer variables in the remainders and leaves the
retained core at its real configuration.

Let \(w'\) be the actual outgoing weight, so that
\(S_{\rm new}=A(w',\bar W)\). Add
\(\mathsf W_X=A(w'-w_{\rm fr},\bar W)\) to the
remainder when completing the action. This difference
of weights is supported in the replaced component.
It has the positive sign since
\(-A(w_{\rm fr},\bar W)=-A(w',\bar W)+\mathsf W_X\).
The profiles give \(w'=w\) off the components. At the
unsourced center, the weight in a completed component
is the current bulk weight before the next increment.
The preceding bounds on the localized terms include
this correction. Summing the three differences and
\(\mathsf W_X\) gives \(P_X\) and the identity
\[
 -S_{\rm ref}-S_{\rm old}+S_0
   =-S_{\rm new}-\tfrac12q_X-S_{\rm core}+P_X.
\]
The exact product of the old weight and the restored law is
\[
 e^{-S_{\rm old}}\frac{\chi_Xe^{-S_{\rm ref}}}{D_X}
   =e^{-S_{\rm new}}\,
      \chi_X e^{-q_X/2-S_{\rm core}}
                 e^{P_X+I_X-E_X}.
\]
This is an identity in $(s,v)$ before either variable is
integrated.  It remains valid with the actual subsequent
averaging graph attached, including at a prescribed terminal
raw field.  No integration over that terminal field is used
to cancel $S_0$.

The factor $e^{-S_{\rm new}}$ belongs to the new label's action.
The operation retains the core Wilson factor and $q_X$ together
with the produced terms, scalar constants and its other carried
factors.  Thus the original old action is used once in this
rewrite.  One must not also keep it as a separate reserve, or
append the old factor $D_X^{-1}$ to the rewritten operation.
The terms $P_X+I_X$ are included once in the complete remainder;
$E_X$ is included once in the scalar normalization at its original
levels.

For example, on an admitted real fibre suppose the complete
remainder bound gives $P_X+I_X\le C N_0$, the scalar count gives
$|E_X|\le C V^\circ$, and $q_X,S_{\rm core}\ge0$.
Then the factor left after $e^{-S_{\rm new}}$ has been assigned
to the new action is bounded by $e^{CV^\circ}$ times the other
retained reference factors.  The new action itself is
nonnegative on a real carrying configuration, so it can also
be retained or bounded above by one in an estimate of the
complete positive measure.  If some of the core energy is used
for a failed-test payment, its unspent portion is the reserve
in this estimate; the same energy is not available twice.

This handles the nonflat normalization within the complete
Wilson product whenever those actual identities and bounds
apply.  It supplies no upper bound of this size for the
isolated probability density $\chi_Xe^{-S_{\rm ref}}/D_X$.
The reference integration at fixed terminal output still
requires the retained Haar or chart coordinates, including
the treatment of tree-fixed pivots.  The cancellation changes
none of those coordinates or constraints.

\paragraph{Drop-outs use a fibre.}
A drop-out has no inner background formed for its own main operation.
Its kept factors have no new core Wilson term.  Its mass therefore
includes the Wilson factor from the carrying term, on the bare
plaquettes whose corners lie in $X$.  The argument is an evaluation
$\mathfrak e=(e,U,\upsilon,Q_1,Q_2,\ldots)$ in a set
$\mathcal K^\sharp(\omega,\xi)$, with the following meaning:
\begin{enumerate}
\item $U$ is real on $X$ and compatible with the stored variables.
The exterior datum $e$ supplies the portions of the carried test reads
outside $X$, in the same physical frame and on the same averaging graph;
\item every carried test has its recorded value on its complete read,
using the reconstructed fields of $\xi$, the prescribed averages of $U$
and the exterior datum where the test reaches beyond $X$;
\item the real pair restricted to every inner island lies in that
island's admitted real core;
\item the complete kept quadratic forms and elimination factors obey
their stated one-sided lower and upper bounds;
\item $\upsilon(p)\ge c\theta_j$ on the action layer that owns $p$.
\end{enumerate}
Compatibility means the actual identities
\begin{equation}
 \mathcal M^jU(b)=V_j(\xi)(b)
 \quad\hbox{on every determining bond whose two blocks lie in }X.
 \label{hi:compatible-fibre}
\end{equation}
At level zero these are identities with the stored bare bonds.  Layer
regularity is imposed both on $U$ and on the relevant stored averages.
It cannot be inferred merely from the realness of $U$.

There are two variants.  The mass of the drop-out itself excludes the
test whose failure made it a drop-out.  A later component may restrict the evaluation set by including that test.
It can charge the root when the full paying carrier supplies a witness.
The mass of the earlier drop-out leaves this root uncharged.

Write the evaluated factor as
\begin{equation}
 \mathsf E_\omega(\xi;\mathfrak e)
   =e^{-A(\upsilon\mathbf1_X,U)}
                         \mathsf K_\omega(\xi;\mathfrak e).
 \label{hi:drop-factor}
\end{equation}
For a fixed exterior the actual kernel keeps the exterior tests on
their inverse reads. For the upper integral one may take the supremum
over admitted exterior values after expressing their kernels in the
common free coordinates. The shell and restored tree charts use
relative Haar in a fixed recorded frame. The exterior changes the
reconstruction and its weight, and these enter the pointwise envelope.
The Wilson share in \eqref{hi:drop-factor} is still the share on the
plaquettes owned by $X$. The set of evaluations varies with $\xi$.
For example, take a probability measure supported on
$\epsilon\le|\xi|\le2\epsilon<1$, and a factor
$e^{-T(1-\cos p)}$.  On the compatible graph $p=\xi$ its integral is
at most $e^{-cT\epsilon^2}$.  If $p=0$ is admitted independently of
$\xi$, the supremum of the factor is one.  For $p=i\eta$ it is
$e^{T(\cosh\eta-1)}$.  Neither replacing the fibre by a product nor
complexifying the real Wilson argument preserves the same small factor.

\subsubsection{Upper integrals and the strong real induction norm}

For a nonnegative function $f$ on a Borel measure space define
\begin{equation}
 \int^* f\,d\mu
   =\inf\left\{\int F\,d\mu:
                 F\text{ Borel},\ 0\le f\le F\right\}.
 \label{hi:upper-integral}
\end{equation}
All inequalities are pointwise, including on sets on which a
parameter-dependent supremum has not been proved Borel.  The value of
the supremum over an empty evaluation fibre is zero.  Monotonicity,
positive homogeneity and subadditivity follow directly: majorize each
summand by a Borel function within an arbitrary positive error and add
the majorants.  For a countable sum use errors of size $2^{-n}\delta$
and monotone convergence of the sum of the Borel majorants.  These
facts do not authorize interchanging an upper integral with an
arbitrary parameter supremum.

For an ordinary component put
\begin{equation}
 m_\omega^{\mathbb C}(\xi)
   =1_{S_\omega}(\xi)
        \sup_{\mathbf p\in\mathcal U_i(X^\sim)}
                           |e^{\Phi_\omega(\xi;\mathbf p)}|,
 \qquad
 m_\omega^{\mathbb R}(\xi)
   =1_{S_\omega}(\xi)
        \sup_{\mathbf p\in\mathcal U_i(X^\sim)\ \mathrm{real}}
                              e^{\Phi_\omega(\xi;\mathbf p)}.
 \label{hi:real-complex-majorants}
\end{equation}
For a drop-out use the supremum of $\mathsf E_\omega$ over its
compatible fibre, including the specified variant of the tests.
The weighted mass, for a nonnegative function $v$, is
\begin{equation}
 \mathsf M_c(i,X;v)
     ={1\over h}\sum_\omega e^{cv(\omega)}
                       \int^*m_\omega\,d\mu_\omega .
 \label{hi:weighted-mass}
\end{equation}
An artificial average $h^{-1}\sum_r1=1$ is inserted when the operation
does not otherwise depend on a root frame.  This gives the same
normalization in both classes.  The integrals over inner components
have their own factors $h^{-1}$; one does not cancel them a second
time in the outer count.

When the core is a fixed finite-dimensional separable space, the
ordinary supremum in \eqref{hi:real-complex-majorants} is Borel under the stated separate
measurability and continuity assumptions.  Indeed, choose a countable
dense subset of the core.  Continuity in the pair identifies its
supremum with the supremum over that subset.  A countable supremum of
Borel functions is Borel.  The fibre of \eqref{hi:drop-factor} need not have such a
fixed dense set.  We will dominate it by an explicit Borel envelope
instead of assuming a measurable optimizing choice.

The input needed in the induction is the real version of \eqref{hi:weighted-mass}, with
the supremum already inside the integral.  Uniform fixed-argument
bounds are weaker.  On the finite probability space
$\{1,\ldots,N\}$ with uniform measure, let
$f(\xi,p)=N1_{\{\xi=p\}}$.  For every fixed $p$,
$\int f(\xi,p)d\mu=1$, while $\int\sup_p f(\xi,p)d\mu=N$.
The example is measurable and positive.  Thus the issue is not cured
by positivity, finite volume or measurability of a partition function.

There is a parallel analytic example.  On $\xi\in\{1,\ldots,N\}$,
let the parameter space be the disjoint union of $N$ unit discs and
set $f(\xi,z)=N$ on the disc indexed by $\xi$, zero on the others.
The functions are holomorphic on that parameter space.  Holomorphy
alone does not interchange supremum and integral.  In the application
one must use the explicit structure of the complete exponent, such as
\eqref{hi:tied-exponent}--\eqref{hi:pair-sup-envelope}, and a common integrable envelope.

A real argument of an inner operation can depend on the outer
variables.  If it also depends on a variable integrated within that
inner operation, its factor is nevertheless bounded pointwise by
$m^{\mathbb R}$, provided it remains in the admitted real core.  The
strong real mass norm is precisely what permits this domination.  A
mere fixed-read estimate would instead require a read-before-sampling
condition, absent from that weaker norm.  These are two different
ways of making a kernel estimate applicable; they must not be conflated.

Finally, upper integrals are used only after a Borel majorant has been
constructed on the full product of stored variables.  We will apply
ordinary Tonelli to that majorant.  Bounds for fibrewise upper
integrals alone are not a substitute for such a construction.

\subsubsection{Genealogies and a volume that does not recount the past}

A node is a component $u=(n,C)$ of the region made at level $n$.
Unless it is healed, it has one successor, the component at level
$n+1$ that contains its propagated hull.  Predecessors of a node are
disjoint components at the preceding level.  The genealogy
$\mathcal G(x)$ of $x=(i,X)$ consists of $x$ and all its iterated
predecessors.  A component healed earlier and covered again later is
the root of a different genealogy.  It is not put back into this tree
because its physical location has been reused.

Cut $\mathcal G(x)$ at every clearing island below $x$.  The cut tree
$\mathcal G^\circ(x)$ includes these cut nodes as leaves; their own
trees are the inner genealogies.  A keeping mark records a drop-out
without clearing its earlier history.  Its older nodes are not copied
every time that mark is subsequently carried.

Write $|C|_n$ for the number of level-$n$ cells in $C$ and define
\begin{equation}
 \begin{split}
 w_n&=\ell_n(MR_n)^4,\qquad
 \operatorname{vol}(u)=w_{n(u)}|C(u)|_{n(u)},\\
 V(x)&=\sum_{u\in\mathcal G(x)}\operatorname{vol}(u),\qquad
 V^\circ(x)=\sum_{u\in\mathcal G^\circ(x)}\operatorname{vol}(u),\\
 v^\natural(x)&=V^\circ(x)+\sum_{y\ \mathrm{inner}}v^\natural(y).
 \end{split}
 \label{hi:volumes}
\end{equation}
This is a sum of sites times levels, with the reference logarithm at
the level of each site.  It is not the volume of the final component
multiplied by its age.  Nor is it a sum of all successive presentations
of its accumulated past.

\begin{lemma}[{Bounded multiplicity}]\label{lem:history-bounded-multiplicity}
For every finite genealogy,
\begin{equation}
 \begin{gathered}
 v^\natural(x)=\sum_{u\preceq x}m_x(u)\operatorname{vol}(u),\\
 m_x(u)=
 \begin{cases}
 2,&u\ne x\text{ is a clearing node},\\
 1,&\text{otherwise},
 \end{cases}
 \qquad V(x)\le v^\natural(x)\le2V(x).
 \end{gathered}
 \label{hi:bounded-multiplicity}
\end{equation}
\end{lemma}

\emph{Proof.} Expand the recursion in \eqref{hi:volumes}.  A node not on a cut
boundary belongs to exactly one outer tree in that expansion.  A
clearing node below $x$ is the last node of its parent's outer tree
and the first node of its own outer tree.  No third tree meets it.
The root belongs only to its own outer tree.  This accounts for every
term in the finite expansion and proves the formula.\hfill$\square$

For site counts without the logarithm put
\begin{equation}
 N^\circ(x)=\sum_{u\ \mathrm{listed\ with\ multiplicity}}
                         (MR_{n(u)})^4|C(u)|_{n(u)}.
 \label{hi:unweighted-site-count}
\end{equation}
Thus $\ell_{\min}N^\circ\le v^\natural$, where
$\ell_{\min}=\min_{j\le i}\ell_j$.  A fixed-width read collar of
one coarse $M$-layer has at most twice the sites of its owning
$MR_j$-cells once $R_j$ exceeds a fixed dimensional threshold.  To
see this, enlarge each owning cell separately and sum: its volume is
multiplied by at most $(1+C/R_j)^4$, while overlaps only reduce the
union.  Lattice endpoint conventions change $C$, not the conclusion.
Finite bond directions and color components are included in the
per-site bound on a term.  Consequently the site budget for such
reads can be chosen so that
\begin{equation}
 n^\sharp(x)\le2N^\circ(x),\qquad
              n^\sharp(x)\le {2\over\ell_{\min}}v^\natural(x).
 \label{hi:read-site-budget}
\end{equation}
The ownership of each read collar is important.  Adding its whole
past again at each keeping mark would not satisfy this argument.

The part of the earlier label that remains inside a drop-out has
volume
\begin{equation}
 V_{\rm in}(i,X;\mathcal S)
       =\sum_{j<i}\ell_j\,\#_j(Z_j(\mathcal S)\cap X).
 \label{hi:label-interior-volume}
\end{equation}
The volume in the left member is the volume of the retained
label. On the compatible inner histories the volume rule gives
$V_{\rm in}\le v^\natural(\omega_X)$. The restored action and stored
regions also enter the load. Compare their larger volume to
$V_{\rm in}$ using the waiting window. At an uncompressed level its
retained box is either the old quiet box or the prepared box
$H_j^{[3a_j]}$ in \eqref{eq:elab-history-prepared-boxes}. It contains
at least $7^4$ cells of its level. The compressed end of the
window has only $O(\log_L\check R_i)$ levels. Thus a fixed
fraction of the window gives this lower count, while the
restored count is at most a fixed multiple of $100^4$ at
each level. The weights $M^4\check R_j^4\log\theta_j$ of those cells
are comparable up to $O(L^4)$ by bounded drift in the window
and rounding to powers of $L$. Outside the window the entries
coincide. This bounds the restored load by $C_L V_{\rm in}$,
while the compatible mass pays with the retained volume in
\eqref{hi:label-interior-volume}.

\paragraph{Why an arbitrary larger volume is not available.}
Bounded multiplicity is a substantive restriction on a useful weight.
For an algebraic example, let a first set of $N$ creations have gain
$NP^2$ and volume $Nw$.  Let $r$ later keeping marks each have gain
$P$.  A volume that recounts that first set at every mark is at least
$(r+1)Nw$.  Choose $N$ so that $cNw>2P$.  Then
\begin{equation}
 c(r+1)Nw-NP^2-rP\longrightarrow+\infty
                         \quad(r\longrightarrow\infty).
 \label{hi:recounting-obstruction}
\end{equation}
The initial gain is finite and cannot pay for this repeatedly copied
volume.  Additional new volume at the marks only worsens the problem.
This example concerns the inference from a gain budget to a volume
weight; it does not assert that a particular field realizes every
abstract record.  It already proves that the scalar ledger alone
cannot dominate an unbounded number of copies of the past.  The
correct comparison is \eqref{hi:bounded-multiplicity}, or another volume with a proved fixed
multiplicity bound.

\subsubsection{Contraction, light stretches and the weighted count}

We now prove the geometric counting statement used in the reserve
calculation.  We state its local hypotheses precisely so that its
application can be checked independently of the analysis of fields.
The cell graph at every level has bounded degree.  The propagation of
an unmarked component takes its next-scale hull, adds a fixed number
of layers, and boxes only components whose diameter in cells is
bounded by a fixed threshold.  A new component is a union of such
propagated components and its new creation cells, with these local
collars.  Equivalently, one can construct an occupied skeleton $O(u)$
from the creation cells below $u$, so that
\begin{equation}
 \begin{gathered}
 O(u)\subset C(u)\subset O(u)^{[e]},\qquad
 O(u)\text{ is }h_0\text{-connected},\\
 |C(u)|_{n(u)}\le C_e\nu(u),\quad \nu(u)=|O(u)|_{n(u)}.
 \end{gathered}
 \label{hi:skeleton-geometry}
\end{equation}
Here $e,h_0,C_e$ are fixed.  The relation between two levels is the
actual cell parent map, with the bounded collars just described.
There are no consecutive size drops.

For clarity, the skeleton uses unique cell addresses.  The actual
closed-cube hull may meet several cells at a boundary.  Include that
bounded ambiguity in $e$.  The sum of the radii of earlier fixed
collars, expressed in the current cell units, is bounded: over every
two levels the older radii contract by $L^{-1}$.  Its bound is a
geometric series.  This proves the fixed-collar part of \eqref{hi:skeleton-geometry} under
the stated local propagation rule.  Boxing a bounded-size component
adds only another fixed collar.  Taking unions does not increase
these constants.  Connectivity holds if a merger is made by the
specified bounded-distance touching rule.  A rectangular hull of a
large thin domain would not obey these properties and is not the
propagation used here.

\begin{lemma}[{Coarse images of connected sets}]\label{lem:history-coarse-connected-count}
Let $A$ be $h_0$-connected, with $N$ cells, in the fine cell grid.
At a coarsening by a factor $q\ge1$, its occupied coarse hull satisfies
\begin{equation}
        |\mathsf H_q A|\le C_0(1+h_0N/q).
 \label{hi:coarse-connected-count}
\end{equation}
The same bound holds with a fixed additional coarse collar by
increasing $C_0$.
\end{lemma}

\emph{Proof.} A spanning tree of the connected cell graph has $N-1$
edges.  Traverse every tree edge in both directions.  Join the cell
centres by line segments along this walk.  Its total length in the
sup norm is at most $2h_0(N-1)$.  Split it into consecutive pieces
of length at most $q$.  Each piece lies in a sup-norm ball of radius
$q$ about its first point.  Such a ball, together with the width of
one fine cell, meets a fixed number of coarse $q$-cells in four
dimensions.  There are at most $1+2h_0N/q$ pieces.  Their union
contains every visited cell centre and hence the required cell hull,
up to its one fixed collar.  Counting this collar changes the same
fixed dimensional number.\hfill$\square$

Choose a fixed number $s_0$ of levels so large that the net
coarsening, at least $L^{\lfloor s_0/2\rfloor}$, makes the
coefficient of $N$ in \eqref{hi:coarse-connected-count} at most $1/4$.  Choose a fixed
$\nu_\star\ge4C_0$.  Propagating a component with
$\nu(u)>\nu_\star$ through $s_0$ levels without new material
then leaves at most $\nu(u)/2$ occupied cells.  For a component
with at most $\nu_\star$ cells, it leaves at most a fixed constant.
Call such components light.  The estimate concerns occupied cells;
a large component need not become light in $s_0$ levels.

Let $o_n$ be the total number of occupied cells in a fixed genealogy
at level $n$, and let $l_n$ be its number of light nodes there.
Unions at mergers reduce the count.  A creation in the intervening
$s_0$ levels contributes at most a fixed number of cells to their
coarse image.  Thus
\begin{equation}
 o_n\le\tfrac12o_{n-s_0}+C_1l_{n-s_0}
              +C_1\sum_{n-s_0<t\le n}n_t^{\rm cr},
 \label{hi:occupied-recurrence}
\end{equation}
with quantities at nonpositive levels set to zero.  At a cut leaf
of an outer genealogy its bounded number of boundary seed cells is
included among $n_t^{\rm cr}$.  The reserve for those seeds is the
last test outside the inner mass, not a second use of the inner
component's old gain.

Let $r_n$ be any nonnegative non-increasing sequence.  Multiply
\eqref{hi:occupied-recurrence} by $r_n$ and sum over the finite levels.  The shifted first
sum is at most half of $\sum_nr_no_n$, because
$r_{m+s_0}\le r_m$.  The light-node sum is at most
$C_1\sum_nr_nl_n$.  A creation of level $t$ occurs in at most
$s_0$ summands, with weights no larger than $r_t$.  Rearrangement
therefore gives
\begin{equation}
 \sum_nr_no_n\le2C_1\sum_nr_nl_n
                   +2C_1s_0\sum_tr_tn_t^{\rm cr}.
 \label{hi:weighted-contraction-sum}
\end{equation}
There is no initial-height term.  This is the first reason the
estimate is uniform in the number of levels.

\paragraph{Waiting only for the light part.}
A quiet node contains no new creation, has one predecessor and does
not immediately follow a mark.  Along a quiet stretch the side lengths
in cells satisfy, for a fixed $A$,
\begin{equation}
 \alpha_{n+1}\le\alpha_n+A,\qquad
       \alpha_{n+2}\le L^{-1}\alpha_n+2A.
 \label{hi:side-contraction}
\end{equation}
Iteration on either parity gives
$\alpha_{n+2m}\le L^{-m}\alpha_n+2A/(1-L^{-1})$.
The smallness threshold in the dynamic class is chosen strictly above
this fixed point and the intervening one-step enlargement.  A light
node has uniformly bounded side length by \eqref{hi:skeleton-geometry}: an $h_0$-connected
set with at most $\nu_\star$ cells has a spanning path of bounded
diameter.  After a fixed further number $b_0$ of quiet levels, the
component is small and is a single box evolving by the prescribed
free rule.

The preparation depths satisfy $\bar N_n\le R_n+1$, with a
non-increasing majorant, and the lower preparation endpoints are
strictly increasing.  These are the admitted size-profile conditions
of the dynamic construction.  After at most another
$\bar N_t+3$ quiet levels from a light node of level $t$, the
entire preparation window lies after the event that began its
stretch.  Quietness, quarantine including the lower endpoint, depth
and single-box traces then all hold.  The full dynamic class acts on
the component.  Hence a stretch begun at level $t$ contains at most
\begin{equation}
             T_t=b_0+\bar N_t+3\le C_TR_t
 \label{hi:light-wait}
\end{equation}
light nodes.  The numbers $T_t$ can be taken non-increasing.
Before these light nodes a stretch can have many heavy nodes; they
have already been controlled by \eqref{hi:weighted-contraction-sum}.  Replacing this argument by
a uniform bound on the whole stretch would incorrectly discard the
time needed to shrink a large component.

There is a uniform terminal alternative on a small torus.  If a
light component or its fixed required collar has no embedded cell
chart, some periodic cell side is at most a fixed constant
$C_\star=C_\star(h_0,\nu_\star,e)$, enlarged for the other fixed
collars.  Otherwise the bounded-diameter connected lift of that
component and its collar would embed.  After $q$ further levels the
same periodic cell side is divided by at least
$L^{\lfloor q/2\rfloor}$.  It must remain at least one at every
admitted partition.  Thus at most
$2\lceil\log_L C_\star\rceil+1$ levels remain.  Increase the
fixed constant in \eqref{hi:light-wait} to include them.  The proof then stops at
the terminal level; it does not claim existence of a nonwrapping
box background for a wrapped component.

There are three ways a stretch begins: a node contains a creation
of its own level, it follows a keeping mark, or it is a merger with
at least two predecessors.  Charge the first kind to one of its
new creations and the second to its preceding mark.  To handle a
merger, choose at every non-leaf node a first predecessor.  At a
merger choose one non-first predecessor and descend through first
predecessors to a leaf.  Assign that leaf to the merger.

This assignment is injective.  If two merger nodes are incomparable,
their leaf sets are disjoint.  If one lies below the other, a path
from the assigned leaf to its merger first enters that merger from
a non-first predecessor.  A descent by first predecessors from the
higher node cannot follow that same branch through it.  Thus it
cannot reach the same assigned leaf.  A leaf contains a creation;
one of its creation cells pays for the assigned merger.  A creation
can pay for its own start and for at most one merger.  This fixed
multiplicity is the second reason the count has no dependence on
height.

Because $T_nr_n$ is non-increasing, paying at the lower level of
an assigned creation or mark only increases the available weight.
Summing the at most $T_t$ light nodes on every stretch and then
using \eqref{hi:skeleton-geometry} and \eqref{hi:weighted-contraction-sum} proves
\begin{equation}
 \boxed{
 \sum_{u\in\mathcal G(x)}r_{n(u)}|C(u)|_{n(u)}
   \le C_2\sum_tT_tr_t n_t^{\rm cr}(x)
          +C_2\sum_tT_tr_t n_t^{\rm mk}(x).}
 \label{hi:weighted-genealogy-count}
\end{equation}
The root's own mark is excluded from $n_t^{\rm mk}$.  It has not
yet supplied a renewed wait.  The constants depend on the fixed local
geometry and $L$, not on the number of levels, the component size or
the torus volume.

\subsubsection{The scalar reserve and its coupling threshold}

Let $G_t$ be the gain per creation, $P_t$ the gain per mark and
$c_t$ the cost per cell at level $t$.  Choose a non-increasing
majorant $r_t\ge c_t$.  Define
\begin{equation}
 \mathfrak F(x)=\sum_tG_tn_t^{\rm cr}(x)
        +\sum_tP_tn_t^{\rm mk}(x)
        -\sum_{u\in\mathcal G(x)}c_{n(u)}|C(u)|_{n(u)}.
 \label{hi:reserve-definition}
\end{equation}
If
\begin{equation}
 G_t\ge2P_t+C_2T_tr_t,\qquad P_t\ge C_2T_tr_t,
 \label{hi:scalar-reserve-premise}
\end{equation}
then \eqref{hi:weighted-genealogy-count} gives
\begin{equation}
 \mathfrak F(x)\ge2\sum_tP_tn_t^{\rm cr}(x).
 \label{hi:scalar-reserve-conclusion}
\end{equation}
Every nonempty genealogy has a leaf and therefore a creation.  If
$P_t\ge(1+\beta_0)^{-1}P_i$ for $t\le i=n(x)$, the last bound is
at least $2(1+\beta_0)^{-1}P_i$.  The reserve is taken from the
creation gain; a mark in this conservative calculation can spend all
its credit on the later wait.

Here are sufficient quantitative choices.  Assume $R_t$ is
comparable to $\ell_t^r$, with fixed $r\ge1$, and put
\begin{equation}
 P_t=p_0(g_t)=A_0\ell_t^{p},\qquad
 G_t=c_GP_t^2,\qquad
 r_t=C_3c_1M^4R_t^5,\qquad p\ge6r+1.
 \label{hi:polynomial-budget}
\end{equation}
Here $c_G>0$ is a fixed gain constant, distinct from the bare spacing
$a_0$ in \eqref{hi:scales}.  The actual cost $c_1\ell_t(MR_t)^4$ is bounded by this $r_t$,
since $\ell_t\le CR_t$.  Also
$C_2T_tr_t\le C_4c_1M^4\ell_t^{6r}$.
Thus \eqref{hi:scalar-reserve-premise} follows from a lower bound
\begin{equation}
       \ell_t\ge K(1+c_1),\qquad c_GA_0\ell_t\ge4,
 \label{hi:affine-log-threshold}
\end{equation}
after increasing a fixed $K$.  To verify this explicitly, divide
the cost bound by $P_t$.  The quotient is at most
$(C_4M^4/A_0)c_1\ell_t^{-(p-6r)}$, at most a prescribed small
constant once \eqref{hi:affine-log-threshold} holds.  Divide the remaining condition
$c_GP_t^2\ge2P_t+C_2T_tr_t$ by $P_t$ and use the second
condition, increasing $K$ if necessary.  Only one extra logarithmic
power is needed.  In particular the lower threshold may be affine
in the coefficient of the volume weight.

For reference sequences with $\theta_t\ge\theta_i-c_5$ for
$t\le i$, one has
$\log\theta_t\ge\log\theta_i+\log(1-c_5/\theta_i)$.
For a sufficiently high coupling floor the ratio of its $p$-th power
to $\ell_i^p$ is at least $(1+\beta_0)^{-1}$, uniformly in
$t\le i$.  This proves the comparison used after \eqref{hi:scalar-reserve-conclusion}, without
assuming exact monotonicity of the reference couplings.

There is a useful consequence for an unfinished component $z=(j,Z)$.
Here unfinished means that its root is neither healed nor marked.
Continue it freely for $q$ levels, stopping before a new creation,
merger or completed wait.  The genealogy and its gains are unchanged;
only the later cell costs are added.  Apply \eqref{hi:scalar-reserve-conclusion} to the extended
genealogy.  If $t_0$ is its first creation level, the sum of its
creation credits is at least $P_{t_0}$.  Therefore
\begin{equation}
 \kappa(z):=\mathfrak F(z)-2P_{t_0}
   \ge\sum_{n=j+1}^{j+q}c_n|S^{n-j}Z|_n\ge0.
 \label{hi:future-wait-credit}
\end{equation}
The extension is cut at the terminal level if necessary.  It is
important to stop before the completed wait: withholding an eligible
component would change the dynamic construction.  The formula shows
exactly what portion of the old reserve is available to pay for the
future quiet evolution.

For comparison, if a separately admitted continuation begins with a
just-marked root, that mark becomes a non-root mark of the extended
genealogy. Its gain is then added once: the corresponding available
expression is $\mathfrak F(z)+P_j-2P_{t_0}$, not the unmarked
expression in \eqref{hi:future-wait-credit}. This observation does not assert the admission
of such a continuation or spend the same mark a second time.

\subsubsection{The entropy of the outer records}

The outer indices include choices of regions and tests, not just a
choice of frame.  We give a counting argument that includes their
positions.  Make a graph whose vertices are cell addresses together
with their levels.  Join neighboring cells at one level and join each
cell to its parent at the next level.  The degree is bounded by
\begin{equation}
                   D\le L^4+3^4+2.
 \label{hi:spacetime-degree}
\end{equation}
At a size drop the parent has only one child in this address graph;
otherwise it has at most $L^4$.  On a torus some neighboring addresses
can coincide, which only decreases the bound.  The outer component is
rooted at a fixed cell of $X$.  Its active cell addresses form a
connected set by the hull and merger rules.  If a local convention
requires a bounded additional collar, include it in the active record
and its already specified volume.

In any graph of degree at most $D$, the number of connected vertex
sets of size $N$ containing a fixed root is at most
\begin{equation}
                         D^{2(N-1)}.
 \label{hi:animal-count}
\end{equation}
Indeed, put a fixed order on the neighbors.  For a connected set choose
its canonical rooted spanning tree, and walk around that tree in depth
first order, traversing each edge twice.  The walk has length
$2(N-1)$ and its visited set determines the original set.  Each step
has at most $D$ choices.  This count includes possible birth levels and
translations of an old creation within a coarse component; they are
encoded in the path down the address graph.  There is no extra factor
equal to the torus volume or to the total height.

For a fixed active set $A$, the remaining discrete record has at most
\begin{equation}
                       \exp\{bV^\circ(A)\}
 \label{hi:ancillary-code}
\end{equation}
possibilities, for a fixed $b$.  Here is the finite coding condition
behind this estimate.  A cell of level $n$ contains
$O((MR_n)^4)$ sites and bonds.  Each local decomposition has a
fixed finite list of outcomes, and a fixed finite number of such
tests is attached to each site or bond.  Specify these outcomes by
letters from a fixed finite alphabet.  The stage place and the local
failure type have at most a fixed power of $R_n+1$ possibilities per
node; their logarithm is at most $C\ell_n$ since $R_n\asymp\ell_n^r$.
Tree frames are treated by their normalized averages below.  Lists
of owned seam factors, node entrances and retained-variable types
use the same finite local addresses, not arbitrary new continuous
parameters.  These facts prove \eqref{hi:ancillary-code} by multiplying the local
counts.  A proposed history representation with another discrete
choice must include its code and its cost here before this estimate
is applied.

There are $1+f(\omega^\circ)$ independently recorded outer frames,
including the artificially inserted root frame when needed.  Their
normalized sum is at most one for each non-frame record.  Since
$V^\circ(A)\ge|A|$ for $\ell_n\ge1$ and $MR_n\ge1$, choose
$c_{\rm e}\ge b+2\log D+\log2$.  \eqref{hi:animal-count}--\eqref{hi:ancillary-code} give
\begin{equation}
 \begin{split}
 &\sum_{\omega^\circ}
       h^{-1-f(\omega^\circ)}e^{-c_{\rm e}V^\circ(\omega^\circ)}\\
 &\hspace{10mm}\le
       \sum_{N\ge1}D^{2(N-1)}
                e^{-(c_{\rm e}-b)N}
       \le\sum_{N\ge1}2^{-N}=1.
 \end{split}
 \label{hi:outer-entropy}
\end{equation}
The outer component is nonempty, which is why the sum starts at one.
An empty component instead represents the identity operation and has
no rare-mass assertion.  This exact upper bound of one will prevent
an unwanted constant from accumulating at every induction step.

The coding condition is a condition on the literal finite
construction.  It must count every tested bond and every stage
place.  The connected-set argument then proves the uniform entropy
bound; it is not an assumption that arbitrarily chosen discrete
histories already have small total mass.

\subsubsection{Normalizer bounds with their complex branches}

We next isolate scalar and finite-matrix estimates that enter the
integrable envelope.  This keeps them distinct from the analytic
problem of identifying the actual kept matrix.

\paragraph{The truncated gauge normalizer.}
Use the matrix operator norm and the sharp radial cutoff
$|u-1|<\epsilon$, $0<\epsilon\le1$. If the class angle is
$r\in[0,\pi]$, then $|u-1|=2\sin(r/2)$. Thus this cutoff is
exactly the monotone interval $0\le t<t_\epsilon$ in
$t=1-\cos r$, with $t_\epsilon=\epsilon^2/2$.
The same proof applies to a cutoff explicitly known to be such an
interval with $t_\epsilon\ge c\epsilon^2$. Put
\begin{equation}
 z(y,\epsilon)=\int_G1_{\{|u-1|<\epsilon\}}
             e^{-y(1-\frac12\operatorname{tr}u)}\,du .
 \label{hi:gauge-normalizer}
\end{equation}
Haar measure is normalized to one.  This is an entire function of
$y$.  For real $x\ge1$ and $x\epsilon^2\ge1$ it satisfies
$c x^{-3/2}\le z(x,\epsilon)\le Cx^{-3/2}\le C$; also
$z(x,\epsilon)\le1$.  The lower bound follows by integrating a
Lie-algebra ball of radius $c_0x^{-1/2}$: it lies in the cutoff,
the exponential is bounded below there, and the Haar chart density
is bounded below.  For the upper bound use
$1-\cos r\ge cr^2$ on $0\le r\le\pi$ and the radial Haar density.
Both estimates have constants uniform in the cutoff satisfying the
displayed condition.

The required complex branch can also be proved uniformly, without
assuming that an integral of positive real values remains nonzero
after complexification.  Write
$t=1-\frac12\operatorname{tr}u=1-\cos r$.  The radial Haar density is
a constant times $t^{1/2}(2-t)^{1/2}\,dt$, $0<t<2$; the cutoff
restricts this to $0<t<t_\epsilon$.  Thus at $y=a+ib$, $a>0$,
$z(y,\epsilon)/z(a,\epsilon)$ is the expectation of $e^{-ibt}$
under a probability density proportional to
$t^{1/2}e^{-at}w(t)$, where
$w(t)=(2-t)^{1/2}1_{\{t<t_\epsilon\}}$ is nonnegative and
non-increasing.

For completeness, decreasing reweighting lowers increasing moments.
If $\rho$ is a probability, $f$ increases and $w$ decreases, then
\[
 \operatorname{Cov}_\rho(f,w)
   =\tfrac12\iint(f(s)-f(t))(w(s)-w(t))\,d\rho(s)d\rho(t)\le0.
\]
Divide by $\rho(w)>0$ to obtain
$\rho(fw)/\rho(w)\le\rho(f)$.  Apply this with the gamma density
proportional to $t^{1/2}e^{-at}$ and $f(t)=t^2$.  Two integrations
by parts show its second moment is $15/(4a^2)$.  Since
$\cos v\ge1-v^2/2$ for every real $v$, it follows that
\[
 \operatorname{Re}{z(a+ib,\epsilon)\over z(a,\epsilon)}
       \ge1-{15b^2\over8a^2}.
\]
On $|y-\theta|\le\theta/2$, elementary geometry gives
$|b|/a\le1/\sqrt3$.  Indeed this inequality is equivalent to
$3b^2\le a^2$, whose worst ratio on that disc occurs at its
tangent from the origin.  It also follows by maximizing
$[(\theta/2)^2-(a-\theta)^2]/a^2$ for
$\theta/2\le a\le3\theta/2$.  Hence the last real part is at
least $3/8$.  In particular
$z(y,\epsilon)\ne0$ throughout that disc.  When
$\theta\epsilon^2\ge1$, the real bounds, with their threshold
changed by a fixed factor if necessary, give
$|z(y,\epsilon)|^{-1}\le C\theta^{3/2}$ there.

Here the lower floor can be seen directly.  The fixed norm gives
$t_\epsilon\ge c_{\rm cut}\epsilon^2$.  Choose a fixed
$0<c_0<\min(1,c_{\rm cut}/2)$.  The interval
$0<t<c_0/\theta$ lies inside the cutoff because
$\theta\epsilon^2\ge1$.  On it $(2-t)^{1/2}\ge1$ and
$e^{-at}\ge e^{-3c_0/2}$ for $a\le3\theta/2$.
Integration of $t^{1/2}$ on this interval proves
$z(a,\epsilon)\ge c\theta^{-3/2}$.  Combining it with the
$3/8$ real-part bound proves the asserted inverse estimate.
The cutoff is frozen at its real reference value throughout:
neither $\epsilon$ nor its characteristic function is varied with
$y$.  Decreasing reweighting alone would not supply this lower floor
for an arbitrarily vanishing cutoff.

There is a unique logarithm on the open disc agreeing with the real
logarithm at $\theta$.  The closed-disc bounds follow by continuity.
Thus a kept inverse gauge normalizer costs at most
$C+\frac32\ell_j$ in the logarithmic budget even on this complex
disc.  This proves the normalizer part of the parameter extension;
it does not assert a parameter extension for the other factors.

\paragraph{The pivot factor.}
On an admitted one-step contour chart the three-color pivot matrix
is the per-length averaging derivative $\mathsf m_c$. Its identity
value is $L^{-4}I$; the raw logarithmic pivot derivative instead has
identity value $L^{-3}I$. In the per-length convention it obeys
\begin{equation}
 |L^{12}\det\mathsf m_c-1|\le\tfrac12,\qquad
 \ell_c=12\log L-\log(L^{12}\det\mathsf m_c),\qquad
 e^{\ell_c}=(\det\mathsf m_c)^{-1}.
 \label{hi:pivot-factor}
\end{equation}
The logarithm is its power series about one, so
$|\ell_c-12\log L|\le C$.  On the real admitted chart the
determinant has the positive identity sign, and this equals the
inverse absolute determinant for the per-length averaging coordinates.
The actual real operation also retains the conversion factors from
raw averaging and Haar coordinates and its other chart Jacobians.
There is no holomorphic absolute-value determinant.  The local
averaging-chart theorem supplies \eqref{hi:pivot-factor} once its small-field and
embedding hypotheses hold.  Its use at every retained layer
additionally requires that the actual iterated averages exist and
fall in those local charts; a one-step statement alone does not
verify that multi-level premise.

Each scalar coordinate, bond direction and local gauge variable is
charged to its owning site.  The inverse gauge normalizers in \eqref{hi:gauge-normalizer},
the pivot factors in \eqref{hi:pivot-factor}, and the rescaling powers of
$g_j=\theta_j^{-1/2}$ therefore have combined modulus at most
$\exp(CV^\circ)$, provided their number at level $j$ is at most
a fixed multiple of the stored sites there.  Chart Jacobians with
uniform local upper and lower bounds have the same property.  The
identity value of the pivot determinant fixes its power $L^{12}$;
counting the same three colors again in this factor would be an error.

\paragraph{The normalization numbers at their own levels.}
The numerical parts in the collar sum are distinct from the real
normalization number retained by an operation.  In the notation of
the step, the latter is
$e_X=n_X-n_G(X)\log(\theta_t^{-3/2}\sigma_0)$, where
\(n_G(X)\) is the number of integration bonds of the
reference law and \(\sigma_0\) is the Haar density at the
chart origin.
Here \(n_X\) subtracts the normalization of the label with \(X\)
adjoined to its regions and all marks in \(X\) removed, including
its own, from the current label normalization. It includes the
stage and drop-out numbers. Its ordinary step part is
the sum, for $1\le j\le t$, of
$N^{(j-1)}(\Omega_j^{\mathsf a},\Lambda_j^{\mathsf a+})-
N^{(j-1)}(\Omega_j^{\mathsf a}\cup X,\Lambda_j^{\mathsf a+}\cup X)$,
with the action region and collared integration region of
\eqref{hi:action-sequence}. In this formula,
\[
 \begin{split}
 N^{(j)}(\Omega,\Lambda)
 ={}&\bigl(\log\sigma_0-\tfrac32\log\theta_j\bigr)
                       |\Omega^{(j)*}|\\
 &-(L^4-1)\log z(\theta_j)\,|\Omega|_{j+1}
 +\bigl(\nu_{\mathfrak g}-\tfrac32\log(x_j/\theta_j)\bigr)
                       |\Lambda^{(j)*}|.
 \end{split}
\]
The three counts are respectively the fluctuation bonds, the
level-$(j+1)$ sites and the nonpivot integration bonds of level $j$.
Here \(z\) is the single-site gauge normalizer and
\(\nu_{\mathfrak g}\) is the logarithm of the standard
unnormalized Gaussian integral on \(\mathfrak g\).
The real normalization stays in $\Phi^\circ_\omega$ or in the identical
normalization of the reference measure.  The retained powers of
$\theta_j$ are charged at level $j$ in the volume
\eqref{hi:volumes}.

For complex normalization weights $y_j$ with
$|y_j-x_j|\le\bar c$, while the bare and net Wilson weights stay
fixed, let $\mathsf g_j(X)$ count the integration bonds of level $j-1$
added when $X$ is joined to $\Lambda_j^{\mathsf a+}$.
The exact scalar difference is
\[
 e_X[y]-e_X[x]
 =\tfrac32\sum_{j=1}^t \mathsf g_j(X)
       \left(\log\frac{y_{j-1}}{\theta_{j-1}}
                    -\log\frac{x_{j-1}}{\theta_{j-1}}\right).
\]
Indeed the stage and drop-out numbers here use only reference
weights and the first two terms of \(N^{(j)}\). In the ordinary
step difference, only the last logarithm changes. The logarithms
use their common branch near the real positive
weights, with $|x_j-\theta_j|+\bar c<\theta_j/2$ at every level.
If the added bonds are covered by one layer around the
\(z_j(X)\) owned \(M\)-cubes, then
\(\mathsf g_j(X)\le4\cdot3^4M^4L^4z_j(X)\).
Each enlarged cube accounts for at most \(M^4L^4\) sites
of level \(j-1\), and there are four positive bonds per site.
Using the bound \(2\) for the logarithm derivative on
the common half-disc and \(N_0\ge M^4\sum_jz_j(X)\) gives
\[
 |e_X[y]-e_X[x]|
 \le 972 L^4\bar c\,
       \bigl(\min_{0\le j<t}\theta_j\bigr)^{-1}N_0(\omega).
\]
In the sourced comparison this difference is added to the new
excess in \eqref{hi:actual-source-recursion}. The real
\(e_X[x]\) can carry logarithms of the reference weights,
which are counted at their own levels in \eqref{hi:volumes}.

\paragraph{The complete quadratic form.}
The real physical weight $x$ and the reference weight $\theta$ have
different roles.  Write $\tau=x/\theta$.  If $K_A$ is the physical
action part and $K_G$ is the reference gauge part of the actual
free-coordinate Hessian, the complete quadratic exponent uses
\begin{equation}
 Q_\tau=\tau K_A+K_G,\qquad K=K_A+K_G,
 \qquad
 e^{-B^{\mathsf T}Q_\tau B/2}
    =e^{-\tau B^{\mathsf T}KB/2}
       e^{(\tau-1)B^{\mathsf T}K_GB/2}.
 \label{hi:compensated-quadratic}
\end{equation}
Both matrices are in the actual free coordinates of the constraint
chart. The last factor compensates for the use of \(\tau K\) in
the normalized reference Gaussian. For unnormalized input the
prefactor is the Gaussian mass of \(\tau K\). With a normalized
incoming precision \(D_0\) it is the ratio of that mass to the
Gaussian mass of \(D_0\). Using \(Q_\tau\) directly instead gives
its own Gaussian mass and omits the compensating potential.

Here is a sufficient full half-disc estimate.  Suppose $K_A,K_G$
are real symmetric, $K_G\ge0$, $K_A\ge-\varepsilon I$ and
$K\ge\gamma I$ with $\varepsilon\le\gamma$.  For
$|\tau-1|\le1/2$, put $s=\operatorname{Re}\tau$.  If $s\le1$,
$sK_A+K_G=sK+(1-s)K_G\ge\gamma I/2$.  If $s\ge1$,
$sK_A+K_G=K+(s-1)K_A\ge(\gamma-\varepsilon/2)I$.
Therefore
\begin{equation}
       \operatorname{Re}Q_\tau\ge\gamma I/2.
 \label{hi:complete-sector-bound}
\end{equation}
If the matrices have complex perturbations with
$\frac32\|\delta K_A\|+\|\delta K_G\|\le\gamma/4$, the
same argument gives $\operatorname{Re}Q_\tau\ge\gamma I/4$.
This is a matrix estimate to be checked on the actual joint core.
If instead a free evaluation directly specifies a complete complex
form $Q$ with $\operatorname{Re}Q\ge\gamma I$, that one-sided
bound is sufficient by itself.

One must not multiply two independently accretive complex objects
and retain the same conclusion.  In one dimension,
$K=1+100i$ and $\tau=1+i/2$ have respectively positive real part
and $|\tau-1|=1/2$, whereas $\operatorname{Re}(\tau K)=-49$.
Nor should the compensating potential in \eqref{hi:compensated-quadratic} be bounded by its
modulus while independently replacing $\operatorname{Re}\tau$
by its worst possible value: their real coefficients are correlated.
Keeping the complete form avoids both losses.

\subsubsection{The integrable reference envelope}

We spell out the additional coordinate statement needed to integrate
the stored Gaussian reserve.  For each outer index there are finitely
many real Lebesgue blocks $b_1,\ldots,b_s$, together with normalized
Haar variables.  A permitted reference change of coordinates has
triangular form
$A_j=F_j(b_j;b_{j+1},\ldots,b_s,U)$ in a chosen order.  For fixed
later variables it is injective on its integration chart, and its
inverse Jacobian is bounded by $J_0^{d_j}$, where $d_j$ is the
number of scalar coordinates in that block.  It can be a translation,
in which case this Jacobian is one.  The actual stored-coordinate
identity or its explicitly bounded chart Jacobian is required; the
mere use of the symbol $A_j$ does not provide it.

Suppose the unspent part of the complete forms gives at least
$\beta\sum_j|A_j|^2$ in the exponent, for a fixed $\beta>0$.
Dropping restrictions only enlarges a nonnegative integral.  Change
the variables one block at a time, starting with the block whose
later variables are held fixed.  Each integral is at most the full
Euclidean Gaussian integral.  Induction over the blocks gives
\begin{equation}
 \int \exp\{-\beta\sum_j|A_j|^2\}\,d\mu^{\circ}
       \le J_0^{\sum_jd_j}(\pi/\beta)^{\frac12\sum_jd_j}
       \le e^{C_5V^\circ}.
 \label{hi:reference-gaussian-integral}
\end{equation}
The last inequality uses the actual coordinate count
$\sum_jd_j\le C N_0^\circ\le C V^\circ$.
Haar variables contribute at most one after positive restrictions
are dropped.  If the coordinate image is only a subset, the same
upper bound holds.  Surjectivity onto a product is not needed.

For a more general actual reference map it is enough to construct
directly a nonnegative Borel function $g_{\omega^\circ}$ with this
integral bound and with the pointwise domination used below.  Such
a construction is a real integration estimate.  It cannot be
replaced by a statement that a complex Gaussian at each fixed
external argument has a small partition function.

For a completed main operation whose core contains no earlier marks,
the reference coordinates can be read directly from its integral.
The one-step operations in that core integrate the stored strip fields
$\mathcal A_j^{\rm s}$ against Lebesgue measure.  Changing the background
changes their quadratic forms and elimination factors, but does not
change these integration variables.  The main operation also integrates
a field $B$ on the shells, with the single scaling $g_t$.
The variable \(B\) stays fixed when the external pair changes,
even though the reconstruction of the group field uses the
background. Its tree gauge
sets the tree bonds to zero and has Jacobian one; on the remaining bonds
the carried test is
\[
 \chi'(B)=\mathbf 1_{\{|B|<r_t\}},\qquad
 r_t=\delta'_t/g_t=A_1p_1(g_t).
\]
Here the averaging constraints can be eliminated from the top down:
the outgoing coarse field is integrated in the main operation, and
integration of its delta function sets it equal to the average of the
preceding field.  The remaining group variables carry normalized Haar
measure.  Chart densities and normalizers remain among the factors
estimated separately above.

Suppose the complete kept forms leave the reserve
$\beta\sum_j|\mathcal A_j^{\rm s}|^2$ after the event gains have been
extracted.  If there are $n_A$ free strip bonds and $n_B$ free shell
bonds, a reference function is
\begin{equation}
 g_\omega(B,\mathcal A^{\rm s})
   =\chi'(B)\exp\Bigl(-\beta\sum_j|\mathcal A_j^{\rm s}|^2\Bigr),
 \qquad
 \int g_\omega\,d\mu^\circ
   \le [C(1+r_t)]^{3n_B}(\pi/\beta)^{3n_A/2}.
 \label{hi:reference-cutoff-gaussian}
\end{equation}
Indeed, enlarge the domain to the product of the cutoff balls for $B$,
the full Euclidean spaces for the strip fields, and the full compact
groups.  The three factors are respectively a finite volume, ordinary
Gaussian integrals, and a Haar mass of one.  This estimate is uniform
in the background and uses no change of coordinates for the strip
fields.

If \(N_P\) pivot-elimination factors \(J_c\) remain in the
integrand relative to this product measure, use the Borel
majorant \(J_*^{N_P}g_\omega\), provided the
actual pivot charts satisfy $|J_c|\le J_*$.  For the convention
$J_c=|\det\mathsf m_c|^{-1}$, the stated regular-chart estimate is
$J_*=2L^{12}$. On an admitted full row the least singular value
of \(\mathsf m_c(V)\) is at least
\(L^{-4}(1-\kappa_9L^3\varpi)\), where \(\varpi\) bounds
the plaquette defect of its chart. The determinant bound follows
for \(\kappa_9L^3\varpi\le1-2^{-1/3}\). In a fixed common
real frame change \(u\), the coefficient transforms as
\[
 \mathsf m_c(V^u)=
 \operatorname{Ad}(u(c_-))\,\mathsf m_c(V)\,
                         \operatorname{Ad}(u(b_0(c)_-))^{-1}.
\]
Both rotations preserve its absolute determinant. The physical
Haar derivative is \(L\mathsf m_c\), with inverse density
\(L^{-3}|\det\mathsf m_c|^{-1}\). The carried elimination
factor above uses its stated normalized convention. This
calculation keeps \(u\) fixed in the pivot derivative and
requires the same admitted row and pivot. The local count $N_P\le C V^\circ$
absorbs this additional factor.  Factors already estimated separately
are not inserted a second time.

To express the bound in terms of $V^\circ$, let $n_{B,j}$ count the free
shell bonds on the lattice of level $j$, so that $n_B=\sum_{j\le t}n_{B,j}$.
The local coordinate count is
$n_A+\sum_{j\le t}\ell_j n_{B,j}\le C V^\circ$; strip bonds of level $j$
belong to their one-step operation at level $j+1$.  The fixed lattice
ratio is included in this constant.  Since $r_t\le C\ell_t^{\mathfrak p_1}$
and $\ell_t\le2\ell_j$ for $j\le t$, the logarithm of the right-hand side
of \eqref{hi:reference-cutoff-gaussian} is at most $C V^\circ$.
Thus this class has an explicit Borel reference function. Keep
\(\chi'\) when discarding the main quadratic form. The event
payments in this bound use the retained Wilson and gauge weights
and the strip forms with their stated reserve. For a drop-out
with fixed outgoing field we must still identify the reference
coordinates with that constraint in place.

A drop-out that leaves after the central box integration has a
different reference measure from one that leaves during a stage or
at the extendibility test.  The central integration has already
integrated the incoming top field.  The fresh Haar field left in the
box is a new variable; holding that field fixed does not restore a
constraint on the old incoming field.  The frame operation retained
by a late drop-out includes this completed integration and its tree
gauge.  Its old averaging chain is therefore eliminated from the
top down, just as in the completed main operation.

Suppose the core contains no earlier marks, and retain the shell
variables in their original group coordinates.  The non-tree shell
bonds then carry normalized Haar measure, also when the test
\(\chi'\) fails.  The tree gauge has Jacobian one.  Thus this case
does not need the cutoff volume in
\eqref{hi:reference-cutoff-gaussian}.  Let \(n_j\) count the retained
gauge normalizers at level \(j\), let \(n_A\) count the free stored
strip bonds, and let \(N_P\) count any remaining pivot factors.
Under the actual reserve and chart bounds stated above, a Borel
majorant for the reference part is
\[
 J_*^{N_P}\prod_j\mathfrak z(\theta_j,\varepsilon_j)^{-n_j}
 \exp\Bigl(-\beta\sum_j|\mathcal A_j^{\rm s}|^2\Bigr).
\]
Indeed the gauge exponentials and the remaining characteristic
functions are at most one.  Integration in the old stored
coordinates gives the bound
\[
 J_*^{N_P}\prod_j\mathfrak z(\theta_j,\varepsilon_j)^{-n_j}
 (\pi/\beta)^{3n_A/2}.
\]
The scalar estimate
\(\mathfrak z(\theta_j,\varepsilon_j)^{-1}
 \le C\theta_j^{3/2}\), together with the own-level coordinate
count, bounds this by \(e^{CV^\circ}\).  Scalar factors already
included in the normalization estimate are counted only once.
This reference bound applies to the late exits at the shell-field
or last-shell tests.  Merely having completed all preparatory
stages is insufficient: an exit at extendibility still precedes
the central integration and retains its fixed-output constraints.
The complete interaction estimate, the energy payments and
admission of the actual evaluation remain separate requirements.

One fixed-outgoing-field case is covered by the regular fibre charts.
Consider a drop-out before any stage has re-made a level, with no
earlier marks.  Suppose every averaging constraint in its retained
operations is already consumed by the regular pivot elimination, so
that no factor
$dV_j\,\delta(\mathcal M(V_j)V_{j+1}^{-1})$ on
$\Omega_{j+1}^c$ remains.  The outgoing field $V_t$ is held fixed.
For each retained step the exact fibre chart solves
$\mathcal M(V_j)=V_{j+1}$ for the pivot variables; its free strip
coordinates are the stored $\mathcal A_j^{\rm s}$.  Thus these
constraints need not be removed by integrating $V_t$.  The reference
integral is already an iterated Lebesgue integral in the free strip
coordinates, with the coarea and chart factors in the integrand.

Suppose the actual pivot charts have the bound $|J_c|\le J_*$, the
remaining densities have the separately established scalar bound,
and the complete forms leave
$c_{\mathrm{res}}\sum_j|\mathcal A_j^{\rm s}|^2$ after the required
payments, with $c_{\mathrm{res}}>0$ independent of the outgoing field.
If there are $N_P$ retained pivot factors and $n_A$ free strip bonds,
the reference part is bounded by the Borel function
\[
 J_*^{N_P}\exp\Bigl(-c_{\mathrm{res}}\sum_j|\mathcal A_j^{\rm s}|^2\Bigr),
 \qquad
 \int J_*^{N_P}e^{-c_{\mathrm{res}}\sum_j|\mathcal A_j^{\rm s}|^2}
                 \,d\mathcal A^{\rm s}
   =J_*^{N_P}(\pi/c_{\mathrm{res}})^{3n_A/2}.
\]
Restricting this function to the actual chart only decreases its
integral.  The local coordinate count bounds the result by
$e^{CV^\circ}$, uniformly in the fixed outgoing field for which these
same chart and reserve bounds hold.  The estimate concerns the
reference part; the complete interaction bound and admission of
the actual evaluation remain separate requirements.

If an unresolved averaging factor survives, positivity alone does not
give this conclusion.  Its normalization integrates the outgoing
field, whereas the drop-out keeps that field fixed.  Nor does
disintegration into conditional probability measures remove the
problem: their base measure is the actual pushforward of the fine
measure, whose density relative to outgoing Haar measure has not
thereby been bounded.

A remaining Haar averaging factor can also be estimated when its
carried support lies in one of the ordinary regular fibre charts.
Keep its outgoing field $W$ fixed and write the exact chart in
unscaled free coordinates $a$ as
\[
 V_W(a)=e^{i\psi_W^{\rm lin}(\mathsf C_Wa)}p_W,
 \qquad \mathcal M(V_W(a))=W.
\]
The fixed-fibre coarea formula on this chart gives
\[
 \int_D F(V)\,\delta(\mathcal M(V)W^{-1})\,dV
   =\int_{D_W}F(V_W(a))\rho_W(a)\,da.
\]
The complete chart density $\rho_W$ includes the pivot
elimination factors, the determinant of the nonlinear linearising
map, the Haar densities, and the fixed powers of $L$ from the group
delta functions.  Suppose the chart is admitted on $|a(b)|<r_L$,
where $r_L$ is fixed, and its stated local estimates give
$\rho_W\le C_L^N$, with $N$ the number of fine bonds in these double
blocks.  There are at most $N$ free bonds, so
\[
 \int_D\delta(\mathcal M(V)W^{-1})\,dV\le C_L'^N.
\]
The bounded coordinate box supplies the volume bound for this Haar
factor. The variables here are unscaled. If scaled
variables are used instead, their scaling Jacobian must remain in
$\rho_W$.

This covers the averaging factors on the belt
$\Omega_t^\sim\setminus\Omega_t$ wherever the two retained step-$t$
tests pass and the step-$t$ averaging constraint is imposed on the
full double blocks. The plaquette and approximate fluctuation
tests put the free bonds in the chart of their prescribed
background. In that chart the constraint solves the pivot
with contraction defect $O(L^{-3})$. The bound uses these
tests at the birth of the factor, in its original gauges.
For the remaining Haar factors the contour formula gives a direct
bound whenever the averaging pivot is one of the integrated bonds.
A regular extension of the coarse datum only supplies a single
candidate fine field.

The total contour formula gives a useful description of the remaining
one-row problem.  Fix all fine bonds except the middle corridor bond
of a coarse bond \(c\), and put \(q=U[c]\).  Haar measure in the middle
bond becomes Haar measure in \(q\).  For each of the \(L\) sites on
the central coordinate axis of \(B(c_-)\), the contour family has a
single member, and cancellation of the backtracking portions gives
\[
 U(\Gamma)U([x,x'])U(\Gamma')^{-1}=U[c].
\]
These terms contribute zero to the logarithmic sum.  Every other
joining path avoids the middle bond, as do both endpoint contour
families.  Thus, on the literal-contour branch, the row has the form
\[
 F(q)=\exp\!\left(\sum_{r=1}^{N_c}
          \alpha_r\log(H_rq^{-1})\right)q,
 \qquad
 \alpha_r>0,\quad \sum_r\alpha_r=1-L^{-3},
 \qquad N_c\le(4!)^2(L^4-L),
\]
where the \(H_r\) depend only on the other fine bonds.  Its intrinsic
logarithms are defined except at the finitely many values
\(q=-H_r\).  Moreover the norm of the exponent is at most
\(\pi(1-L^{-3})\).  On the regular branch the same form holds with
the averaged endpoint contours in place of the literal contours and
at most \(L^4-L\) terms.  Their values are again independent of the
middle bond.

\paragraph{Haar measure for a single averaging pivot.}
Identify $G=\mathrm{SU}(2)$ with the unit round three-sphere.
Write $\exp_q$ and $\log_q$ for its Riemannian exponential and
principal logarithm. The formula at a fixed set of nonpivot
bonds is, on either logarithmic branch,
\[
 F(q)=\exp_q(w(q)),\qquad
 w(q)=\sum_i\alpha_i\log_q H_i,\qquad
 s=\sum_i\alpha_i=1-\beta,\qquad \beta=L^{-3}.
\]
Indeed right translation to the identity sends $\log_q H_i$
to the intrinsic group logarithm of $H_iq^{-1}$. Each $H_i$
and $\alpha_i$ is fixed while the pivot varies. Consider this
formula on the complement of all the antipodes of its anchors,
before restricting it to the set on which its branch is used.

For a tangent vector $v$ with $r=|v|<\pi$, set
$H(v)=P_v+r\cot r\,(1-P_v)$, with $H(0)=1$. In the polar
metric $dr^2+\sin^2r\,d\omega^2$, this is the Hessian at $q$
of the half squared distance to $\exp_q v$. The logarithm is
the negative gradient of this squared-distance function in
the base point. Differentiate $\log_q F(q)=w(q)$ to obtain
\[
 B(w)\,dF_q=H(w)-\sum_i\alpha_i H(v_i),\qquad
 v_i=\log_q H_i,\qquad B(w)=D_2\log_q(F(q)).
\]
Use parallel transport back to $q$ to compare the tangent
spaces. In these coordinates $B(w)$ is positive with radial
eigenvalue one and transverse eigenvalue $|w|/\sin|w|$,
as follows by differentiating the exponential in polar coordinates.

The matrix function $1-H(v)$ is convex in quadratic-form order.
The cotangent expansion \cite[formula 4.22.3]{DLMF-pivot} gives
\[
 1-H(v)=2\sum_{n\ge1}
         \frac{|v|^2 I-vv^{\mathsf T}}{\pi^2n^2-|v|^2}.
\]
For a unit vector $e$ each summand tested against $e$ is a
positive multiple of $|P_{e^\perp}v|^2/(\pi^2n^2-|v|^2)$.
The quadratic-over-linear function is convex and decreasing
in its positive denominator. Its composition with the
linear numerator vector and the concave denominator is
therefore convex. Jensen with weights $\alpha_i/s$ gives
\[
 H(w)-\sum_i\alpha_iH(v_i)\ \,\ge\,\ H(w)-sH(w/s).
\]
Writing $r=|w|$, the radial eigenvalue of the lower bound is
$\beta$. The transverse eigenvalue divided by the transverse
eigenvalue of $B(w)$ is $\sin(\beta r/s)/\sin(r/s)\ge\beta$,
using concavity of sine on $[0,\pi]$. At $r=0$ this holds
by continuity. Thus the determinant of $dF_q$ is positive
and at least $\beta^3=L^{-9}$.

The multiplicity is at most one for this formula with its
fixed anchors. Combine coinciding anchors and let $p$ be one
of their antipodes, with weight $a_p$. As $q$ approaches $p$
along the unit tangent direction $\omega$, the map tends to
\[
 \exp_p(b_p+\pi a_p\omega),\qquad
 b_p=\sum_{i:H_i\ne-p}\alpha_i\log_p H_i .
\]
The limit is uniform in $\omega$. The ball of vectors of
centre $b_p$ and radius $\pi a_p$ lies inside $|v|<\pi$,
since $|b_p|+\pi a_p\le\pi s$. Replace each puncture by its
sphere of directions and attach a ball mapping by
$u\mapsto\exp_p(b_p+\pi a_p u)$ for $|u|\le1$. Give the
attached ball the orientation of the original geodesic
coordinates about $p$. Its interior map preserves orientation.
The resulting continuous map of a three-sphere has degree one.
To see this, multiply every weight by $t\in[0,1]$ and use the
same construction. At $t=0$ the map collapses each attached
ball to its puncture and agrees with the identity outside
these collapses, hence has degree one. For a target off the
finitely many cap boundary spheres, the preimages are finite
and each has local degree $+1$. Degree additivity
\cite[Proposition~2.30]{Hatcher-pivot} therefore
bounds the number of original preimages by one. Two original
preimages of a target on a cap sphere would give two for an
open set of nearby targets, by the local diffeomorphism
property. This proves the same bound there.

The total row is a restriction of at most two such injective
maps on measurable branch sets and is the identity in the
corridor product on the remaining cut set. That set is finite
for fixed nonpivots, by the literal-contour formula. Thus the
pushforward of normalized pivot Haar measure has density at
most $2L^9$ relative to normalized outgoing Haar measure.
In particular, for every nonnegative normalized approximate
identity $\delta_\epsilon$ and every outgoing value $W$,
\[
 \int_G\delta_\epsilon(F(q)W^{-1})\,dq\le2L^9.
\]
This upper bound is uniform even at an exceptional output.
An exact Borel choice of fibre kernel is available. For each
smooth branch take the inverse pivot on the image of its
actual Borel branch set, with weight $|\det dF|^{-1}$. The
map from the pair of nonpivot data and pivot to the pair of
nonpivot data and outgoing value is a Borel injection on
each such set. The branch sets and the graph of this map
are definable, by the same finite formulas as for the
real averaging map. Projection gives the image and the
graph of the inverse, so both are Borel by cell decomposition.
The Jacobian weight is Borel by the displayed derivative formula.
Add the two kernels,
using the zero measure on an empty branch. The resulting
kernel has mass at most $2L^9$ for every target and gives
the exact joint measure by change of variables on each
branch. The finite fallback set has zero conditional
incoming Haar measure and so contributes zero to this
version of the kernel. For the version retaining the
pointwise fallback branch, write $D_c(a)$ for its finite
pivot set at nonpivots $a$. Since the row is the identity
in the corridor product there, add
$\mathbf 1_{D_c(a)}(W)\delta_W(dq)$ to the kernel. This
keeps the exact joint Haar measure and increases the
pointwise bound to $2L^9+1$. The choice of version has
to be identified if a prescribed target comes from a
singular lower graph. Equality of the integrated joint
measures does not identify their pointwise values on
that graph. On a later gauge slice one must also check
the actual order and gauge action of the stored operation.

The row kernel is covariant under the real site gauges of
the row. The corridor product, the anchors and the output
are all multiplied by the same two endpoint factors. The
logarithmic branches and their cut sets are preserved.
Left and right multiplication preserve Haar volume, so
the inverse Jacobian weights are unchanged. In particular
the kernel can be used before taking a gauge quotient when
that is the order of the operation and the quotient acts
on the corresponding joint measure.

%
%
\paragraph{Gauge draws with an outgoing value held fixed.}
\hypertarget{history:fixed-output-gauge-lift}{}
Before discarding any variables, realize one partial history
on the bare field, the frames and the fresh fields of the
reference draws. Keep the temporary Haar variables of the
reference boxes as separate coordinates until they are
overwritten. The history weight is the product of its tests,
gauge factors and reference factors, each evaluated on the
field at the time of insertion.

Let \(u\) be a real gauge admissible for the specified draw.
It is constant on blocks already gauged and equals one near
the retained drop-outs and the retained second-class slices.
At time \(s\), use \(u^{(s)}=1\) on the inner boxes that will
be remade after \(s\), and \(u^{(s)}=u\) elsewhere. On the
bare field use \(u^{(0)}\), and on a fresh field use
\(u^{(s)}\) for the time of its draw. The frames are unchanged.
The temporary Haar coordinates can be transformed separately,
since they carry no earlier test and are subsequently overwritten.
The lifted change of variables preserves the bare Wilson
measure, the fresh Haar factors and the retained slice measures.
At a retained slice the lift is the identity.

Induct over the moves using covariance of the averages.
Outside the boxes already remade, the current field transforms
with the mask at the time of the read. On replaced rows
later moves take the fresh field at its own level, and
earlier factors retain their original arguments.
Earlier blocks are either contained in a later inner box or
disjoint from it, so the mask is still constant on each
previously gauged block. The plaquette tests and the relative
bond tests keep their norms. The inner tests use the transported
tree or a slice on which the gauge is the identity. Reference
normalizers transform with their defining integrals. Thus
the lift preserves the history weight and intertwines the
gauge action on its stored variables.

In a preparatory stage first impose the next-background
regularity test. Then insert the new level-\(j\) gauge factors
at the coarse centers in
\(((\Omega_{j+1}^{\sim3})^c\setminus Z''_{j+1})\cap Z\)
whose blocks do not yet carry the factor. The other collar
factors are retained from their earlier steps. The lift uses
these actual block groups and the normalizer of each factor
at its insertion.

For the ordinary draw the gauges fix all coarse centers.
Its raw outgoing average \(W\) is therefore unchanged.
Insert any bounded Borel \(F(W)\) in the invariant history
measure. Changing variables in the normalized gauge draw
gives the same weighted integral as inserting the gauge factor
as a product. This proves the identity of outgoing measures,
also in the presence of any other prescribed coordinates fixed
by the draw. On the pass support the contour variables used in
the factor are defined. Outside that support the total kernel
uses its separate normalized convention at undefined contours.
A pointwise bound on a singular prescribed fibre still uses
the corresponding branch or graph estimate.

On a row with a free Haar pivot the ordinary draw has a
stronger fixed-output identity. Hold the nonpivot coordinates
fixed before the draw, and write \(q=V[c]\) for the straight
corridor product. Both coarse endpoints are fixed by the
gauge. Thus \(q\), the coarse-endpoint contour anchors in its
row map, and the criterion for the averaging branch are
unchanged under the draw. The contour density of the draw
uses only internal block links and is independent of \(q\).
At a prescribed output \(W\), the pivot inverse in \(q\)
therefore has the same branches and densities before and
after the draw, using a common version at exceptional targets.
Equivalently each normalized group approximation to the
averaging delta has identical arguments on the two sides.
The gauge integration in these coordinates is the normalized
draw on the nonpivot data. For a later ordinary draw the old
large region is excluded by the five-layer buffer. The
blocks read by these contours and the neighboring endpoint
blocks are in the old small-field region with its positive
density tests. Their inherited plaquette bound implies that
the contour means in the normalized density are defined.
Other prescribed coordinates moved by the draw remain in
the joint graph. An older reference-tree gauge that moves
coarse endpoints requires its own row calculation.

In a nonwrapping one-step domain a pivot crosses the face
between the two endpoint blocks of its row. No other coarse
double block contains both ends of that fine bond. Thus, if
every constrained row has a free integrated pivot, fixing
all nonpivots gives a product of single-pivot maps. The
preceding bound is then $(2L^9)^N$ for $N$ constrained rows,
uniformly in all the outgoing values. In a stack of such
integrals one may integrate the pivots from the finest level
upward, keeping all still unintegrated coarse variables fixed.
This supplies the Haar part of a positive reference bound
on those stacked domains. The actual kernel requires a check
that each pivot is present in its stated integration measure,
that no pivot is also used as a separately stored coordinate,
and that the stated charge counts the participating rows.

In the completed main-box operation the incoming averaging
constraints have already been integrated when the frame tree
is imposed. An older island can instead carry a restored field
with a tree factor into a later averaging step. An interior
middle bond can belong to that tree. For example, in the frame
of a box comb whose last coordinate change is in direction
one, the interior bonds in that direction are tree bonds.
For a later row with a pivot on such a bond we must identify
another free coordinate or work with an exact gauge quotient
of the joint averaging law. A tree gauge moving the two coarse
endpoints also acts on the prescribed output of that row.

The next ordinary block gauge is fixed at the coarse centers,
so by covariance its insertion preserves that raw average.
In addition it is applied outside the enlargement of the older
input region. To see the role of the endpoint condition, take
identity exterior and the identity box minimum. The restored
cutoff bounds every bond in an interior double block by
$C\varepsilon_i$ from identity. The average of its row is then
within $C_L\varepsilon_i$ of identity. A Haar draw of an
independent old gauge at one coarse endpoint instead gives that
row a Haar marginal. Thus integrating those old gauges would
change the raw law on this fibre. The Wilson weights can change
the density inside its support; an estimate for the raw output
still needs the variables of its actual chart.

\paragraph{Averaging the restored box in its tree slice.}
\label{hi:restored-box-rows}
Fix the frame of one restored box $P=\Lambda^{\mathrm{prop}}_i$
and hold its exterior fixed. The box is a union of $M$-cubes
of level $i$. Since $M$ is a positive power of $L$ on the nested
centered grids, $P$ is a union of whole level-$(i+1)$ blocks.
Use a lift of the box, so that each coordinate line of blocks in
it is a finite interval. Take direction one last in the chosen
frame. In the restored variables $V'_i=V_iV_\Lambda^{-1}$ the
tree fixes all interior direction-one bonds. It fixes just one
boundary-crossing bond, the entrance. The other restored bonds
with an endpoint in $P$ carry product Haar measure. A row of the
first ordinary average can read one of these variables only if
at least one of its endpoint blocks lies in $P$.

For the derivative at the identity use dimensionless fine bond
logarithms $A$. Let $e_\mu$ denote a fine unit vector and let
\[
 \varphi(y)=L^{-4}\sum_{x\in B(y)}\bar A(y,x),
\]
where $\bar A(y,x)$ is the average of the path integrals along
the symmetric contours in $B(y)$. For a row $c=(y,y')$ the
derivative in the coarse bond logarithm is
\[
 (\mathbf Q^{\mathrm{lin}} A)(c)
  =\varphi(y)-\varphi(y')
   +L^{-4}\sum_{x\in B(y)} A([x,x+y'-y]).
\]
In particular, when both blocks are in $P$ and $c$ has
direction one, the straight-path term vanishes on the tree
slice. Choose in each block $B(y)$ the fine bond
$b_y=(y+e_1,y+e_1+e_2)$. It lies strictly off the first
coordinate face supporting the tree's direction-two bonds.
Its coefficient $\omega_y$ in $\varphi(y)$ is positive.
Indeed, the contour to $y+e_1+e_2$ in order $1,2$ uses it
with positive sign, giving at least $(2L^4)^{-1}$, and every
other contour using that bond has the same sign. No other
block contour uses $b_y$.

On each direction-one line through $P$, leave the selected
bond in the first block among the fixed variables. For each
interior direction-one row, select $b_{y'}$ in its right block.
Order the rows along the line. Each such row reads the current
and at most the preceding selected bond. The diagonal derivative
is $-\omega_{y'}$ on each color. For each remaining row, select
a free fine bond crossing the face between its two blocks.
In an interior transverse row the middle crossing bond is off
the first coordinate face, hence free. In a boundary row the
face has $L^3$ crossing fine bonds, at most one of which is the
entrance. Every chosen face-crossing bond belongs to just that
row among the double blocks and is absent from the endpoint
contours. Its coefficient in the straight-path average is
$L^{-3}$. Order these rows after the interior direction-one
rows. The square derivative in the selected variables is thus
block lower triangular, with these nonzero diagonal blocks.

The same order gives a real chart on the restoration cutoff.
Fix all unselected coordinates. On the full double block of a
row take a real gauge of the fixed reference $V_\Lambda$ with
links close to identity, using the comb estimate in the admitted
small-plaquette domain. The full blocks reach at most $L$ fine
spacings outside $P$. The prepared box is at distance at least
$4M$ from $\Omega_{k_0+1}$ in level-$i$ units and is inside
$X\cap\Omega_i^c$, two coarse cells from $X^c$.
Thus every full double block lies inside the region of the
retained exterior tests and the extendibility comparison.
The box minimum bounds the plaquettes meeting $P$ by
$C(M)\varepsilon_i$; on the remaining plaquettes it agrees
with the prescribed exterior, whose retained density tests
give the small-plaquette bound. On the cutoff the restored
logarithms have size at most $M_0\varepsilon_i$. Conjugating
the logarithms at their initial sites preserves the balls and
fixes their zeros on the old tree. The output gauge acts by
left and right translations at the coarse endpoints. Both
changes preserve the corresponding Haar measures.

In the row's gauge the complete field on its double block
is uniformly near identity. The analytic averaging formula
therefore applies throughout the selected coordinate ball.
The derivative of the row logarithm in its selected variable
differs from its scalar diagonal derivative at the identity
by at most $C(L,M)\varepsilon_i$. Take the coupling ceiling
small enough that this error is less than half the absolute
diagonal coefficient. On a convex relative-logarithm ball,
integration of the derivative along the segment between two
inputs bounds the difference of their outputs below by half
the coefficient times the input distance. Thus the row map
in this coordinate is injective with inverse Jacobian bounded
by a fixed $C_L$ for Haar measures. Fixed site conjugations
and output translations return the bound in the original frame.

The row dependencies used for the triangular order are exact
double-block dependencies of the nonlinear map. An interior
direction-one row sees its two selected block links only.
The remaining rows can also read previously selected block
links, and each reads only its own selected face-crossing link.
Use the one-row change of variables in this order to reconstruct
the selected coordinates at a raw output. Equivalently, apply
change of variables in the resulting triangular chart. The
Jacobian product and the remaining normalized Haar coordinates
give the bound $C_L^N\,dW$ on the joint output of the $N$ rows
meeting the restored box. This is a bound for the Haar slice
restricted to the cutoff, before applying its density. The
inverse on the chart and zero mass off its image give a Borel
version, with the bound at every output. The joint graph over
the unselected coordinates and the exterior is Borel, and the
inverse is taken on that graph. All these changes
hold with the actual fixed exterior values in the fibre.

After combining the restoration law with its old Wilson
weight, use this chart with the core factor, the real form
and the scalar prices assigned as above. The bound covers
the whole first-step row stock of one restored box at fixed
exterior. Its use in a subsequent constraint requires the
joint-coordinate check below.

\paragraph{Disjoint row charts and their outgoing fields.}
At a fixed level $i$, let $P_a\subset X_a$ be the boxes
restored in the distinct second-class components of the label,
whose tree conditions remain in the output. Every
double block used by a row meeting $P_a$ lies inside $X_a$.
Indeed $P_a\subset X_a\cap\Omega_i^c$ with the two-cell
buffer of \eqref{eq:elab-history-quiet-boxes}, and a double
block reaches at most $L$ fine spacings past $P_a$.
The $X_a$ are distinct components on the same grid. Thus
their first outgoing row sets are disjoint. Conditioning on
the complement of the restored bonds fixes all the exterior
coordinates in each box chart. In the ordinary step $i+1$
the complement of the new determining member contains the
hull of the carried region $Z_i$ enlarged by eight layers.
The block-gauge draw is outside its five-layer enlargement.
All the first outgoing double blocks of $P_a$ lie inside
$X_a\subset Z_i$, so they lie on the actual raw averaging
rows and outside the next gauge draw. This applies even
when the successor merges with another component. The
waiting rule leaves these rows in place, including when
$i+1$ is the ordinary bottom of the later operation.
The remaining rows have both blocks outside all restored
boxes and use no restored coordinate.

This also identifies the input coordinates at the outer face
of each second-class component. The top inner member of the
prepared nest lies in $P_a$. Central integration consumes its
last bare averaging graph and the old top bonds meeting $P_a$.
After the draw the determining list contains every level-$i$
bond meeting $X_a$. The tree constraint is only on bonds
meeting $P_a$. Outside those bonds the output carries the
original level-$i$ Haar coordinates with the old weights;
there is no remaining averaging constraint there from inside
the island. In particular the physical exterior that supplies
the centre of the restored law is read in the top coordinates
of $X_a\setminus P_a$. The earlier core variables stay in the
island integral.

Since $X_a$ is a union of whole level-$(i+1)$ blocks, a row
that reads a fine bond meeting $X_a$ has an endpoint block
there. If its two blocks miss $P_a$, its middle fine bond
meets $X_a$ and avoids $P_a$. It is one of these ordinary
Haar inputs, even when the other block is outside $X_a$.
The next-step enlargement contains the complete double block
in the bare region. The one-cell separation of distinct
same-level components prevents any such row from having a
pivot in two components. The rows touching $P_a$ and the
remaining rows touching $X_a$ together number at most
$8\#_{i+1}(X_a)$. Charge them to the actual successor of the
island, which contains their full blocks. Use the strong real
bound for the old island integral at its admitted argument.
The tree condition on the fresh restored field is outside that
integral and belongs to the first outgoing chart.

On any additional bare row having an actual free Haar pivot,
use its middle face bond. That bond crosses between its own two
blocks, so a row of a box chart cannot read it. Conversely
an ordinary row with both blocks outside all boxes reads none
of their selected coordinates. First solve the ordinary rows
outside the boxes at fixed unselected coordinates. These
row maps do not use the fresh box fields. Reconstruct the
ordinary pivots and fix that resulting physical exterior.
Then use the simultaneous chart for each disjoint box stock.
The reference field of a box may depend on the reconstructed
exterior, so it is fixed at this second step. Fubini's theorem
gives the product of the inverse-density bounds in this order.
Any retained weights are evaluated at the resulting inverse
graphs. In each recorded frame write the fresh field as
$V'V_\Lambda(e)$ at its reconstructed exterior $e$.
The reference tree fixes $V'$ on the same bonds for every
such $e$ in that frame, and right multiplication preserves
the relative Haar measure on the non-tree coordinates.
Bound the outgoing density uniformly in $e$ on this
relative chart. If an outside input to one of the ordinary
rows comes from a lower graph, keep it in that joint
measure while reconstructing $e$.
The frame and denominator are those of the original
summand. The actual tests restrict the domain of the chart. The number
of rows touching a box is at most eight times its number of
coarse blocks. The successor contains the propagated component
and hence these row blocks. Charge the rows to its sites at
level $i+1$. The disjointness at one level and the waiting
rule for successive marks count each assigned row once. Thus
the product of the first-box chart constants has an exponential
cost in the corresponding site-level count of
\eqref{hi:volumes}.

There is a different first chart for a component that
stops at the shell-field or last-shell test, after central
integration and before reference restoration. Central
integration has consumed the old box field, including
the bonds crossing its boundary. The resulting term has
unrestricted fresh Haar coordinates on the bonds with an
end in the box. Its frame tree acts on the old shell
variables. After bounding the retained weights, the
normalized Haar integral over that tree is at most one
at fixed exterior root. Any older marked core retains
its own joint integral and reference bound.

At this fixed exterior the reference used to define the
old shell fluctuation depends only on the old exterior
of the prepared box. The tree coordinates on that
fluctuation have the rooted-tree Jacobian one. They are
disjoint from the fresh Haar pivots of the outgoing rows,
so the combined coordinate change has this unit tree
factor times the raw row Jacobian. This permits weights
to be evaluated on the reconstructed old shell without
imposing any tree condition on the fresh pivots.

The box is a union of full level-$(i+1)$ blocks. A first
outgoing row of level $i+1$ that reads the box has at
least one of its two blocks there, so its middle fine
bond has an end in the box. It is one of the free Haar
coordinates, even for a row crossing the box face.
At fixed exterior the corresponding one-row bounds
therefore apply jointly to all these rows. A row whose
two blocks lie outside the box reads none of those
coordinates. The same count of at most eight rows per
coarse block pays for the first outgoing density, using
the unrestricted Haar row bound instead of the regular
restored chart. The reference-law slice enters only on
the branch that reaches restoration.

%
%
%
At a shell or last-shell stop the same order recovers
the exterior of the prepared integral. A density cube based in the
component uses four test layers of width \(M_2\check R_i\)
in level-\(i\) units, and at most one further bond.
This reaches less than one level-\(i\) cell, of side
\(M\check R_i\), since \(L^2M_2\le M\).
The next full double blocks add at most \(2L\) in
these units. The hierarchy \(M\ge L^2M_2\) and
\(M_2=L^2M_1\) gives
\(4M_2\check R_i+1+2L<M\check R_i\), so the resulting
read still lies within one old cell of the component.
An old cell is no wider
than a new one by \(\check R_i\le L\check R_{i+1}\).
The five- and eight-layer enlargements in the next
ordinary step therefore put these rows in the raw
region outside the gauge draw. Solve the ordinary rows
with both blocks outside the box first. Their inverses
return the selected exterior pivots. Keep the other
exterior coordinates among the unselected variables or
in their lower joint graph. A value shared by two
tests is reconstructed once. At this recovered exterior
solve the remaining rows, each with its free box pivot.
The prepared datum and the frame variables are read
in the old integral at this exterior. Use the same
values for the original positive tests and for the
prepared-datum support condition.

Keep the exact inverse charts when using the resulting bound.
Write $t$ for the selected variables and $s$ for the unselected
ones. At fixed $s$ and raw output $w$, let $t=\psi(s,w)$ be
the reconstructed variables on the image, with zero density
off the image. If a one-row change has several branches, keep
the corresponding sum of inverse densities. A nonnegative
Borel test or weight depending on $t$ is evaluated at these
reconstructions. The bound by a product Haar measure on $w$
can be used with an envelope uniform over the reconstructions.
This applies in particular to the uniform upper bounds for the
complete Wilson product, after the old, new and core factors
have been assigned their stated roles.

For the next raw averaging constraint the input is the
current-level field. On its full double blocks this field is
assembled from $w$ and the other current-level coordinates
that actually remain fixed or integrated. A later bare average
can use $w$ as a dominated Haar coordinate when these other
arguments are fixed by $s$ or supplied by their own joint chart.
Older variables in test and energy factors are still evaluated
by the inverse charts. The ordinary and prepared averaging deltas
both take the physical field at their own level as input. The tree
condition at an older island remains in that island's chart. An
old fine variable read by a carried weight therefore stays in the
inverse kernel for its outgoing row. If an operation prescribes
such a fine value as a target in its own right, we must use its
joint graph with the outgoing row. The available pivot choices
are then those left free under both prescriptions.

The waiting rule protects the first outgoing scale of an old
island. A later operation containing it remakes only scales
above $i+1$, so its retained variables below the first outgoing
field remain in the old core. The averaging constraint at
$i+1$ reads the restored argument $V_i$ and produces the raw
field $V_{i+1}$; the older determining datum $V''_i$ is in a
separate slot. This assigns the first raw rows once, without
giving a second Haar variable to a stored datum. For the
continued stack use the row assignment in
\cref{hi:remade-row-ownership}, with the joint outgoing law
of the earlier operation on the actual read set. Along a
consecutive bare chain whose current-level rows satisfy the
joint chart condition, iterate the conditional estimate from
the first outgoing field to the last one. The Jacobian product
then bounds the chain, with its actual stored weights read at
the inverse charts.

The intermediate fields must be read from the retained operations.
In an unmarked core, the field of level $j$ has its stored values on
$\Gamma_j$; inside $\Omega_j^c$ its values are the averages of the
preceding field, imposed by the preceding averaging kernel.  They
are not additional independent Haar variables.  The same reading
applies when the earlier drop-out left before completing its
preparatory stages. In the broader waiting convention, at a full re-entry
the remade shell instead carries the old current-level datum;
the remaining core carries either the preceding average or the
fresh argument from the central integration, according to the
recorded place of departure.  A fresh argument and a retained
determining datum on the same bond remain distinct variables.

Successive graph kernels can be compressed exactly.  If $x$ denotes
the variables of an earlier operation, $u=A(x)$ its output, and $p$
the other retained arguments, then, for every nonnegative Borel
function $H$,
\[
 \int \delta_{A(x)}(du)\int\delta_{B(u,p)}(dw)H(x,u,w)
   =H\bigl(x,A(x),B(A(x),p)\bigr).
\]
Here $B$ uses the full field assembled from its stored and generated
entries.  Consequently the compressed terminal constraint is the
constraint on this assembled field.  It becomes
$\mathcal M^2(V_{j-1})=V_{j+1}$ only when all entries used by the
second average are supplied by the first.  This identity follows by
substitution in Dirac measures and preserves the original integration
variables and their measure.

For a fixed-output mass estimate, this substitution does not supply
Haar measure in a generated pivot.  The one-row estimate applies
directly only to a pivot that is actually integrated against Haar
measure, with the other row arguments held fixed.  To apply it to
a pivot supplied by an earlier operation one first needs a bound
for that operation's joint output density on the variables used.
Likewise, a pivot crossing a retained seam must be assigned to its
actual integration or reconstruction.

The prepared-stage Haar factor integrates the determining bonds of
level $j$ in the nest $Z''_{j+1}$. Those bonds lie in or meet the shell
$Z''_{j+1}\setminus Z''_j$. After the collar change the
remaining delta rows are the inner bonds of level $j+1$ of
$Z''_{j+1}$. A row crossing the prepared face belongs to the
collar chart below.
For a row deep in $Z''_j$ its lower inputs have to be read through
the preceding operation. At the outside seam, the middle fine
bond of a crossing row may itself cross into $\Omega''_{j+1}$.
The no-inward determining set then omits that bond. In an ordinary
one-step operation such a row, where it is in the admitted chart,
belongs to the pivot reconstruction for the fluctuation step.
At a re-made level the distinction is between the old wall and
the new prepared face. The following decomposition identifies
which reconstruction each uses.

Use the preparatory integral of \cite[(1.60), (1.71)--(1.72)]{B89a}.
At a preparatory stage put $Y=Z''_{j+1}$ and
$C=(\Omega_{j+1}^c\setminus Y)\cap Z$. Divide the rows of the
step in this component into three sets. Let $R_{\rm ord}$ be
the rows meeting the old $\Omega_{j+1}$, whose pivots are in the
ordinary fluctuation and strip reconstruction. Among the rows
with both endpoints outside that region, let $R_{\rm col}$ be
the ones meeting $C$, and $R_{\rm in}$ the ones with both endpoints
in $Y$. A row crossing $\partial Y$ belongs to $R_{\rm col}$
by the separation from the old $\Omega_{j+1}$. The regions are
unions of complete level-$(j+1)$ blocks. The full double block
of a row in $R_{\rm in}$ is contained in $Y$, and that of a row
in $R_{\rm col}$ lies outside $\Omega_{j+1}$. Such a row
meets $\Omega_{j+1}^c$, which is contained in the complement
of the top determining member. The two-layer buffer of that
complement in $Z$ keeps the full double block in the same
component. The old-wall and preceding-face margins put it
in the level-$j$ member of the preceding stage. All its links
are then in the non-inward level-$j$ datum before the promotion.
The double block misses the earlier inner face. At the first
stage the collar uses at most five of the eight layers
separating the old members. At a compressed stage the spacing
of the prepared faces is at least one level-$(j+1)$ $M$-cube,
which exceeds the double-block reach. Thus a middle bond
selected from this collar stock is a level-$j$
integration coordinate. It is either an old exterior Haar
datum or a target left with its carried density after the
preceding collar integration. The earlier inner graph has
its targets wholly in $Z''_j$ and does not solve this pivot.
The soft gauge insertion is a density on this stock.
Rows carried outside the component keep their own reconstruction.

Use the collar correction on $R_{\rm col}$,
\[
 V_j=\exp i\bigl[g_j\mathsf C B_j
              -\mathsf h\widetilde{\mathsf D}
                    (g_j\mathsf C B_j)\bigr]V^{(j)}.
\]
The map $\mathsf C$ keeps the nonpivot increments, and
$\mathsf h$ acts only on its assigned pivots. The pivot of a
crossing row lies on the common face of its endpoint blocks,
outside every double block for $R_{\rm in}$. A row of
$R_{\rm col}$ uses no pivot of $R_{\rm ord}$, since that
pivot lies within or crosses into the old $\Omega_{j+1}$.
An ordinary row can read the collar field in its exterior
block. Its middle bond enters $\Omega_{j+1}$ and is inward
also in the promoted determining list. Its value was solved
by the original strip reconstruction at the prescribed
coarse row. In a sharp failure, first recover the collar
variables in the branch where the approximate test passed.
At fixed prospective datum the two stage centres are
independent of the changed collar bonds, since those bonds
are omitted from the new determining list. The extra joined
datum of the localized centre lies over the old top member
and does not reinstate those bonds. The old strip centre
may read the collar field. Read it at the reconstructed
exterior and use the original inverse at the same prescribed
coarse target. The two sets of pivot
derivatives are triangular in this order. Keep the old
elimination and strip Jacobian factors, and evaluate each
surviving old indicator on its complete original carrier.
The collar chart supplies only the density for its own rows.
On inner fine bonds of $Y$ the field still consists of the
retained shell values and those of the earlier operation.

Here is the exact factorization on any carried support on which
this collar chart is valid.  Denote those remaining variables by
$s$, their measure by $d\nu(s)$, and all the actual collar
coordinates by $a$.  Write $\Psi_W(s,a)$ for the reconstructed
field at the prescribed coarse datum $W$.  Its full coarea
density, with the old densities and tests, is $J_W(s,a)$.
If $F$ is the remaining integrand, the collar integration gives
\[
 I_W(s)=\int F\bigl(\Psi_W(s,a)\bigr)J_W(s,a)\,da.
\]
The surviving operation is therefore, with unsolved constraints
still understood as their original graph kernels,
\[
 \int I_W(s)\,
       \delta\bigl(\mathcal M(V_j(s))W^{-1}\bigr)
                    \big|_{R_{\rm in}}\,d\nu(s).
\]
Indeed the local coarea formula consumes each collar row once,
and the remaining row constraints are independent of $a$ by
the double-block containment.  All dependence on the prescribed
crossing-row values stays in $I_W$, including the elimination
factors and nonlinear Jacobians.  No extra Haar variable has
been adjoined. The delta factor for each solved row has already
been integrated in the chart, together with its coarea factor.

\paragraph{The seam after successive stages.}\label{hi:remade-row-ownership}
At the first stage the prepared-box separation from
$Z_{m_\flat}\cap X$ puts the full row blocks near the new face
outside the bottom operation. At each subsequent stage
the preceding completed collar leaves its level-$j$ graph
constraints only on rows with both endpoints in $Z''_j$.
The separation from $Z''_j$ to the new face exceeds the
reach of the full row blocks at that face. Thus the input
field there is read from the level-$j$ coordinates of the
retained law or from the outgoing side of the preceding
chart. The earlier inner graph is integrated in its own
variables. A read through an inverse chart adds no second
constraint. The three-way split at this stage leaves the
next inner graph on the rows wholly in $Z''_{j+1}$.
This proves the row assignment for every prefix of
completed stages.

A failure of the regularity or approximate collar test,
or of the completed ordinary fluctuation test, precedes
the collar pivot change. A failure of the sharp collar
test occurs after it. In this branch retain the collar
nonpivots and the exact translated density, including the
factors for the eliminated pivots. The determining list
can still be the one before the stage. The stage has no
integrated normaliser because its collar integral has yet
to be performed.

The row assignment uses the collar chart on its actual
real support. We must verify its admission on each support
where it is used. A bound for the surviving inner integral
also needs the joint law of the old core arguments from
their retained operation.

\paragraph{The joint law of the retained core rows.}
Fix the history and its outgoing field \(W\). Unfold each
older drop-out into the ordinary and stopped-stage operations
it retains, including the remaining averaging graphs. Each
unfolding lowers the level of the mark. The evaluation in
the upper integral varies the weights and compatible real
backgrounds at the stored values. The graph equations in
those variables and their target \(W\) are fixed throughout
the supremum.

Write \(h=t-\mathsf N_t\) for the bottom of a class
window of level \(t\), and \(G=Z_h\cap X\) for its core.
Every mark in \(G\) has level strictly below \(h\).
For \(h>0\), the last move in the core is thus the
ordinary operation \(\mathbf T^{(h-1)}\) on the output
of the lower operation. It keeps the density test at
level \(h\). The unresolved raw targets at that step
have their two endpoints in \(\Omega_h^c\cap G\).
At its birth the component contains two cell layers
around \(\Omega_h^c\cap G\), so the complete double
blocks of these targets are in the ordinary core. In the
unmarked part of each block a level-\((h-1)\) bond with an endpoint
in \(\Omega_{h-1}\) is an original Haar datum. There
is no inward higher endpoint on this block. A bond
wholly in \(\Omega_{h-1}^c\) is a raw output of the
preceding operation. This supplies the ordinary bottom
row bound from the joint measure of that lower operation
and the original Haar slots. At \(h=0\) there is only
the bare Haar input with its Wilson weight and tests.

In the lower operation use the field belonging to each
older mark at its recorded departure. For a partial
stop, use the original top determining data and the
surviving lower graph for the other rows. A fully
prepared stop at extendibility has its new top datum
on the remade shell and the surviving bare graph in
the inner box. Central integration comes later. A stop
after that integration and before restoration has the
fresh Haar slots and free first pivots described above.
In a closed island the fresh field is on its reference
tree slice, so the field at that level can be singular
with respect to product Haar. The simultaneous first
outgoing chart supplies the bound for those rows. The
fixed exterior includes the actual retained arguments
of each such chart. On a fresh slot used only by the
next graph, a same-bond datum retained in the old
determining set is a separate variable.

The unrestricted fresh field of a late open stop is
carried upward by averaging. To see which later reads
can prescribe it, let the mark be born at level \(i\).
While it stands, it remains a component of its original
\(Z_i\) with all entries through level \(i\) unchanged.
In the input \(Z_{r-1}\) of a later ordinary step \(r>i\)
it lies in the old large region. The ordinary step places
\(\Omega_r\) at least eight coarse cells from that input
region. The new plaquette tests are inserted outside the
four-layer enlargement of its hull. The approximate tests
and their gauge reads are outside the five-layer buffer.
Even the complete four-margin test cube with its endpoint
blocks has reach less than one coarse cell. These tests, and
the density tests with bases in \(\Omega_r^{\sim4}\), miss
the fresh box. The immediately carried test of level \(i\)
uses the separate clearance of the central box from the
original level-\(i\) density reads. At subsequent ordinary
steps the preceding density tests use the same ordinary
buffer argument at their own level.

If a later eligible component contains this standing mark,
its preparation begins at a level strictly above \(i\).
Inside the mark the lower determining entries are therefore
the frozen ones. A retained row at a bond address shared by
the fresh field still reads its old slot. This argument uses
the eight-cell gap only on the original ordinary walls.
The prepared walls use their own, possibly compressed, gaps.
Outgoing averages can use the fresh slot, and a boundary
weight reading it stays in the joint integral. Thus these
generated fields remain in the joint measure up to the
next average. The bound for their reconstructed row law
still has to hold at the actual exterior datum.

These cases identify the measure used at each row of the
bottom operation. For an induction through older marks,
use unnormalized joint upper measures with these Haar coordinates
and reference slices before the first outgoing average.
Then use the outgoing row bound for the whole read set
of the next operation. The prescribed-output estimate
requires this domination with the actual fixed exterior
and with all weights read on the reconstructed fibres.
If a generated slot is later used as a spectator, keep
it in this joint bound unless a bound for a fixed value
of that slot has been proved separately. Dividing a
joint upper bound by the spectator's marginal gives no
uniform conditional bound without a lower estimate for
that marginal.
The scalar mass of an older mark, obtained by integrating
its outgoing field, does not give that domination.

Write \(A\) for the fluctuation coordinates with kept
Gaussian forms and \(a\) for any free nonpivots of a
translated, unintegrated sharp stop, in unscaled logarithmic
units. Let \(v\) be the group coordinates and \(\zeta\) the
old variables of the inner islands.

For an open subtree of pre-main stops, the stopped factor
consists of the Wilson and gauge weights, scalar stage and
orientation terms, the elimination factors, the kept
quadratics, and the old island factors. Completed stage
corrections are in the carrying boundary exponent. After
allocating the event payments, the residual Wilson and
gauge numerators are bounded by one. The gauge denominators
cost at most \(C\theta_l^{3/2}\) per normalized factor at
level \(l\). On the admitted real chart each ordinary
pivot determinant inverse is at most \(e L^{12}\), and the kept forms
give a common Gaussian reserve \(\beta>0\). At an admitted
real pair use a Borel upper majorant \(m_a(\zeta_a)\) for
the old island, taking its integral as close to the earlier
upper-integral norm as needed.

For a sharp stop the passed approximate test in the
translated chart gives \(|e^{ia_b}-\mathbf1|<2\delta'_{j(b)}\)
on its free collar bonds, including when the sharp test
fails. In the principal chart this gives
\(|a_b|<C\delta'_{j(b)}\). At fixed centre and coarse
datum the pre-translation coordinate is also
\(O(\delta'_{j(b)})\) on the admitted chart. The density
on this box keeps
the Haar factors and the exact derivative of the
linearizing correction. In the correction the range is
on the pivots. Ordered after the nonpivots, the pivot
blocks are independent because each double block uses
only its own pivot. The correction has zero first
derivative at the origin and uniformly bounded second
derivative in the regular one-row chart. Its Jacobian
therefore costs at most a constant per pivot on this
small box. Include the Jacobian of the cutoff translation
on its own bonds as well. The translated free coordinates
need only the volume of their compact box.

Thus on admitted inverse charts of this subtree the
complete reference weight has the pointwise bound
\[
 H_\omega(A,a,\zeta,v)\le C_\omega
         e^{-\beta\sum_b|A_b|^2}\prod_y m_y(\zeta_y)
         \prod_{b\text{ at sharp stops}}
             \mathbf1_{\{|a_b|<C\delta'_{j(b)}\}}.
\]
Here \(\log C_\omega\) is bounded by the sum of the
stage and orientation scalar costs and the gauge, pivot,
Haar and cutoff density costs just listed, counted at
their own levels.
Old chart factors are read at their reconstructed
arguments. The tree and cutoff of a fresh restored field
stay in its coordinate measure. For earlier stops in the
main step use the retained measures described below.

At an old island, a later raw row uses the fresh field in
the restored box. The old island variables occur in the
weights. The pair-uniform envelope allows us to integrate
those old variables first and keep the fresh field and its
exterior as inputs to the next average.

Start at the lowest remaining scale. Carry a joint upper
measure for the full current field read by the next double
blocks, adjoining its original stored Haar coordinates.
For a fresh restored box use its first-average chart on the
whole row stock at fixed exterior. On each other bare row
use its middle bond when this is a Haar coordinate of the
assembled field. A double block reads no other selected
middle bond at the same scale. Rows already solved in the
ordinary or preparatory collar charts keep their coarea
factors. Use the remade masks in \cref{hi:remade-row-ownership}.
A complete inner double block at a prepared level \(j+1\)
lies in \(Z''_{j+1}\). For \(j>h\) in the completed
prefix, each of its level-\(j\) bonds
with an endpoint outside \(Z''_j\) is an original datum
coordinate at level \(j\), including a bond straddling
that earlier face. Both of its endpoints are outside the
new determining member at level \(j+1\), so the inward
mask does not remove it. A bond with both endpoints in
\(Z''_j\) is an output of the surviving bare averaging
graph at the preceding level. These two sets give all
the inputs to the next double block. Once the preceding
graph has a joint bound on its whole output, adjoin the
original Haar datum coordinates in the same conditional
integral. This supplies the next row bound with no seam
coordinate left unaccounted for. The first graph still
requires the joint output of the bottom operation, and a
stopped level uses the determining mask at departure.

A variable still read by a weight is evaluated at its inverse
graph. Every additional argument of the next constraint must
be in the same joint measure or fixed among the unselected
coordinates of that chart.

More explicitly, fix the genuine exterior \(e\), let \(v\)
be the selected current Haar inputs and let \(s\) be the
unselected coordinates with reference measure \(\nu_e\).
Use the lower joint density \(k_e(v,s)\) against
\(dv\,\nu_e(ds)\). If \(\kappa_{e,s}(W,dv)\) is the
inverse kernel of the new row chart and \(f_e\) is its
remaining weight, the composed density is
\[
 \int \nu_e(ds)\int \kappa_{e,s}(W,dv)\,
                k_e(v,s) f_e(v,s,W).
\]
If the inverse kernel has mass at most \(B\) and the
product of the two densities on its branches is at most
\(E_e(s)\), this is bounded by \(B\int E_e\,d\nu_e\).
This calculation uses the branch versions defined by the
original graph equations. If a generated spectator \(z\)
is itself a prescribed raw target, include its generating
row in the same complete lower row stock as the surviving
outputs \(v\). If the enlarged list admits a joint inverse
chart with independent pivot coordinates at the complete
fixed exterior, the inverse-density sum gives a kernel
\(k_e(v,z,s)\,dv\,\nu_e(ds)\) at the specified target \(z\).
Use the chart's actual branch version and its density bound.
For example, when $v,z$ are distinct targets of one bare
averaging operation with their original middle links still
free under Haar, their complete double blocks do not read
one another's middle links. Hold the common nonpivots fixed
and take the product of the corresponding inverse measures.
%
There is an additional regular case when the spectator \(z\)
occupies the middle link of the outgoing bare row. Suppose
the lower joint law at prescribed \(z\) leaves another
link \(b'\) of the same corridor free under Haar. Restrict
to the part of this row with small plaquette defects on
the entire double block. Choose the comb starting at the
initial top and traversing the row direction first.
Every comb crossing the face crosses it at the fixed
middle link \(z\), so the frame is independent of \(b'\).
The comb estimate puts the links in a small logarithmic
box in this same frame for all the admissible \(b'\).
Write \(f(a)\) for the logarithm of the output in these
frames, where \(a\) is the logarithmic coordinate of \(b'\).
The derivative in \(a\) at the flat field is \(L^{-3}1\).
Indeed, the \(L^4\) straight paths of the average traverse
that link for precisely \(L\) choices of the starting
point, and the averaged contours are internal to each
endpoint block. At fixed \(L\), analyticity gives
\(\|D_af-L^{-3}1\|\le L^{-3}/2\)
throughout a sufficiently small box. Choose a convex
coordinate ball in that box containing the admissible
free-link values. Integration on a segment gives
\[
             |f(a)-f(a')|\ge \tfrac12L^{-3}|a-a'|.
\]
The row thus has one inverse branch on the image of the
ball and an inverse Haar density bounded by \(C_L\).
Use this density on the actual test subset, evaluated on
the inverse, with value zero off its image. The other
rows do not read \(b'\), since it joins only these two
endpoint blocks. Index the constraints by positive
orientations and read the reverse as the inverse on the
graph. This replaces just the conflicted pivot,
with a geometric choice of alternative and the same
reference Haar measure before taking the exterior
supremum. The lower joint density \(k_e(b',z,s)\) at \(z\)
remains part of the integral, with the old pivot recovered
for every old weight that reads it. The alternate solve
uses no division by a marginal at \(z\).
A tree-fixed alternative or the intermediate logarithmic
branch outside the regular box requires its own chart.
On the straight-path fallback the output is independent
of \(b'\). If a different link \(b''\) of the original
straight path still has a free Haar coordinate, hold the
other inputs fixed and solve that link instead. The row
is \(Pb''Q=W\), so its inverse is \(P^{-1}WQ^{-1}\)
with Haar Jacobian one. Read the original fallback
predicate at that inverse. Other rows can read \(b''\),
but the fallback row does not read their face pivots.
Solve the fallback row first and then the other pivots
at the reconstructed inputs. If the straight path is
already fixed, its output stays on the prescribed graph.
Fix the free-coordinate reference measure separately for
each geometry, frame and branch choice.

A relation between prescribed outputs belongs to the graph
measure. Adding its generating equation to the list gives
no second free output or independent Haar factor. Fix the pivot assignment
by the geometry and the frame. The reference measure of
the remaining free variables is then the same for the
different exterior values. Put the inverse-density
factors and the weights read at reconstructed values
inside the pointwise supremum on those free variables.
Rows whose targets are only used to reconstruct the
exterior keep their target variables in the integral.
If all inputs of a row are already prescribed, its
target remains on the deterministic graph. This argument
applies to the complete protected first-average stock of
a restored box.

%
%
For an ordinary successor of a collar-kind departure of
level \(i\), keep the old top density halo through the
first outgoing row changes. The old top link of the
stopped prefix is either a value of its original
determining list or the output of its surviving
\(\mathcal M(V_{i-1})\) constraint. Include the full
lower double block of the latter row, as well as the
first outgoing blocks. This read lies within the
two-layer enlargement of the old component on the
next grid, inside the ordinary successor. When
the component later enters a containing preparation
this read is in its floor. The new prepared rows
leave these old graphs in place.

For a halo link used only in evaluating a retained
test, use the value \(z=\zeta(v,s)\) of that old graph
in the weight of the joint row integral. In the
notation above the corresponding operation is
\[
 \int\nu_e(ds)\int\kappa_{e,s}(W,dv)\,
       k_e(v,s) f_e(v,s,W,\zeta(v,s)).
\]
Equivalently the conditional integrand has the graph
measure \(\delta_{\zeta(v,s)}(dz)\) in that read slot.
The test retains this value until it has entered the
weight. Its presence uses the original inverse rows
and carries no new Haar density or integration volume.
In particular, the translated collar variable of the
stopped history remains the variable in its recorded
reference measure, as does the birth fluctuation in the
kept Gaussian. If a new constraint solves one of
these variables, insert its inverse value and chart
density before taking the supremum and omit its
integration volume.

For a prescribed exterior slot outside that stock, classify
its generating row at the stage where it was formed, using
\cref{hi:remade-row-ownership}. A promoted collar row uses
the original collar chart on its complete double block. An
ordinary wall row reads its old inverse at its coarse datum.
A row of the surviving inner graph uses the lower joint law.
If the backward closure of the read leaves the stored row
sets already supplied, we must supply that joint exterior
read or keep its relation on the actual graph.

The geometric closure along consecutive bare rows has an
exact description. Start with a level-\(q\) row and stop
the closure at level \(u<q\). Each first inverse reads
fine links with both endpoints in the union of the two
endpoint blocks. For every such link, the next inverse
reads only the two blocks of its endpoints. The union of
the child blocks of all sites in a parent block is the
same parent block in the finer lattice. Induction thus
keeps the level-\(u\) reads in the two order-\((q-u)\)
blocks of the original row. This applies also to the
nonpivots required by the inverse, which lie in its full
double block. If one of those slots comes from a restored
box or a collar change, stop the bare closure there and
include the stock of that actual chart. Test and weight
carriers are included separately at their own levels.

%
%
For an ordinary row outside \(\Omega_q\), apply the
no-inward rule to its level-\((q-1)\) inputs. A link
meeting \(\Omega_{q-1}\) is a datum there, since both
endpoints are outside \(\Omega_q\). A link with both
endpoints outside \(\Omega_{q-1}\) uses its actual
lower raw graph if that graph is still carried. The
two-cell enlargement of \(\Omega_{q-1}^c\) in
\(Z_{q-1}\) keeps its full double block in the same
predecessor component. The next growth contains the
predecessor in its successor. If the upper row reads
several predecessors, keep their inputs together in
the same conditional integral. When the recursion reaches
the first outgoing step of a closed island, the read ends
at its physical field. Use the restored first-average
stock and the ordinary Haar inputs at its face. On a region
absorbed by healing, the new determining list has dropped
the old lower members. Read the current datum there and
include the healed contribution through its weight. For a
standing drop-out use its departure list and continue only
through the rows it retains.

On a carrier clear of older replaced boxes, compatibility
with a real bare field propagates through the same graph.
On every complete double block the determining inputs
agree by admission and the other inputs agree by the
lower rows. Applying the average proves agreement at the
new target. For a carrier meeting a box replaced at a late
stop, the branch uses its fresh top argument. The nested
block inclusion above locates this exception using the
whole backward carrier. The older integrated datum stays
in the inner weight of the box.

If a later weight reads an
older branch or any of its remaining variables, keep that
argument in the joint integral, or include the weight in
the lower envelope before taking the marginal. An earlier generated coordinate
can be supplied as a selected Haar input only after its
joint density has been obtained by integrating the lower
rows. The test and weight functions still use the inverse
values of any lower arguments they read. In the prepared
seam, first solve the inner raw rows, then the collar
rows, then the ordinary rows that can read solved collar
values. This order keeps the actual one-sided dependencies.
An inner pivot can enter the collar centre through a
retained determining row. Reconstruct it before using
the collar chart. In the no-inward convention the
ordinary pivot at the old member is absent from the
fine determining list of the collar centre, since it
has an endpoint in the enlarged member. Thus this pivot
is solved only at the last step of the order.

At fixed output the compatible set is read at the inverse
of each branch. The pointwise envelope bounds the upper
integral in these coordinates before the unselected variables
are integrated. Use the Borel inverse-density sum with zero
mass off the chart images. For the outputs admitted to the
independent charts, iteration on full read sets gives the
joint outgoing density bound
\[
 C_\omega(\pi/\beta)^{3n_A/2}
 \prod_{b\text{ with remaining }da_b}(C\delta'_{j(b)})^3\,
 \prod_y\int m_y\,d\mu_y\,
 (2L^9)^{N_{\rm raw}}(C_L+1)^{N_{\rm alt}}C_L^{N_{\rm box}} .
\]
Here \(n_A\) counts independent Gaussian coordinates
still carrying \(dA\) after the row solves. A fixed or
reconstructed slot has its actual chart weight and no
volume factor. For each pivot choice, count the remaining
volumes on that branch. Sum the resulting bounds if the
branches have different free coordinates. The counts \(N_{\rm raw}\),
\(N_{\rm alt}\) and \(N_{\rm box}\) are for bare rows using
the middle link, bare rows using the alternate charts
and first-average rows of restored boxes, respectively.
The alternate charge allows a regular branch of bound
\(C_L\) and a straight-path fallback of bound one.
Only rows still to be solved enter these counts.
Use $2L^9+1$ in place of $2L^9$ for a middle-link row with
the pointwise fallback version. Deterministic read slots
stay on their graph. The already consumed ordinary and collar
coarea factors are in \(C_\omega\). Fixed exterior
slots keep their prescribed values. The row assignment
charges each price at its own level.

%
%
%
%
%
For an earlier open operation, start with the full
blocks read by the selected outgoing rows and take the
closure under lower reads. A row in that closure is
either genuine exterior data, a retained Haar datum,
part of a restored box's first-average stock, or an
output of lower rows. Include the whole stock when
using a restored box chart. Keep any unused outputs
until all weights that read them have been included.
An earlier promoted shell coordinate was the target
integrated against Haar in the corresponding averaging
integral. Its density already includes the elimination
of the old pivot. Use that density in the joint
integral, with the old pivot recovered for any weight
that still reads it.

At an extendibility stop all preparatory stages have
been completed, and the bare averaging rows inside
the prepared members survive. The last such row has
its target in the old field \(V_i\) on \(Z''_i\).
This is the field to be used in the next ordinary
average. The extendibility indicator reads only the
part of \(V_i\) outside the central box, which contains
\(Z''_i\). The infimum in that test gives a predicate
on the exterior. Its trial extensions do not appear
as variables or constraints in the retained kernel.

When forming the density of this history, take its
old top field as the argument of the inner graph.
At the next raw average use the same top field in
the outgoing Haar integral. Where the old field is
a target of the inner graph, include its earlier
joint density in that integral. The other top
bonds have their original Haar factors. Thus the
inner graph is evaluated at its prescribed target
before that target is used as an input of the next
row. Lower coordinates read by the old tests stay
with their inverse branches. In particular the
negative extendibility indicator is evaluated on
the reconstructed top exterior in this joint change
of variables. The fine inputs retain their integration
and constraints from the earlier graph.

At a shell or last-shell stop the central integration
has eliminated the prepared chain. Its old terminal
variables and the variables of the old core are
disjoint. Indeed, write \(h\) for the bottom of
this component's waiting window. The shell gauge
ends at the two-cell enlargement of the prepared
\(\Omega_{h+1}\). The prepared inner box at \(h+1\)
has three cell layers around the hull of the bottom
core. After the two layers have been removed, the
remaining cell layer still contains the complete
bottom-core reads.
The frame tree therefore acts only on the retained
shell. At fixed reconstructed exterior, the change
from the prepared group variables to their relative
values preserves Haar measure. Solving the frame
tree leaves normalized Haar on the free shell bonds.
Use this reference measure and the separate fresh
box Haar measure. The shell and last-shell failures
are separated before the operation that assigns the
reference-law density. Their fresh output therefore
carries bare normalized Haar. For a component that
remains in the class, combine the reference-law
density with the old Wilson weight by the pointwise
identity following \eqref{hi:inner-wilson}. Assign the
new action to the carrying term and keep the core
Wilson weight for the ledger. The scalar and
remainder costs enter at their own levels. This
cancellation holds before imposing any new averaging
row, so it also holds on the prescribed-output
branch. The subsequent Haar bound still uses the
retained tree and averaging coordinates. Keep all
earlier variables still read by the weights in the
lower joint integral.

Consequently a late node adds a normalized shell
factor and the free first-outgoing pivots to the
reference measure, with the old lower charts and
branch tests retained. If the middle bond of an
ordinary bottom row is itself a lower graph output,
the joint outgoing bound at that lower level
supplies its Haar domination together with the
other inputs. At a closed island use the simultaneous
first-outgoing chart on its restored slice. The
pointwise envelope can be formed before those row
integrations.

%
%
%
The upper-integral fibre allows one-sided evaluation
coordinates for the quadratic forms and pivot factors.
Write \(q_j\) for the complete kept quadratic at step
\(j\), including its coupling ratio. In the stopped
weight this is \(\tau_j\langle A_j,K_j A_j\rangle\),
where \(K_j\) is the sum of the kept chain contributions
on the original free-coordinate space. Use
\(\operatorname{Re}q_j\ge\gamma_{\rm p}\|A_j\|^2\)
and the upper bound \(eL^{12}\) for a carried
elimination factor. For an actual field read in the
original positive strip domain \(\mathcal O^+\), the local kept form
has a real-part lower bound \(c\gamma_0\). With the
original masks and the same free coordinates, the
absolute walk sums give \(\|K_j\|\le C_{\rm w}\).
For the ordinary strip the ratio is formed at its
original level \(j\). On \eqref{hi:parameter-tube},
the radius bound \(t_j\le C\theta_j(\log\theta_j)^{-2q_{\rm t}}\)
gives
\[
 |\tau_j-1|\le C(\log\theta_j)^{-2q_{\rm t}}
                        +\bar c/\theta_j.
\]
The radius estimate follows from adjacent frame
comparability. At the last level use monotonicity
of the tube radii and the preceding bound, or the
single-level tube condition. Choose the inverse
coupling floor so that twice the right-hand side
times \(C_{\rm w}\) is at most \(c\gamma_0/2\).
This leaves at least \(c\gamma_0/2\) for the complete
kept form throughout the ordinary strip ratio disc.

The ordinary normalization may instead keep the
form at \(\tau_j=1\) and place its ratio correction
in the boundary potential. The gauge compensation
in \eqref{hi:compensated-quadratic} also has to
be accounted for in the complete exponent. On the
stored boxes the compensation costs at most
\(C|\tau_j-1|M^4 p_0(g_j)^2\) per cube. This stays
bounded on the tube because \(q_{\rm t}>\mathfrak p_0\).
On the enlarged stored-field discs the squared field
bound adds \((\log\theta_j)^{12r+2}\). The same cost
tends to zero there with the hierarchy
\(q_{\rm t}>\mathfrak p_0+6r+1\).
If a compensation is included in the form, check
the lower bound for the resulting complete form.
Preparatory and main restoration Gaussians use
\(\theta_j\). In those operations the change of
coupling belongs to the Wilson potential. An
abstract one-sided evaluation is admitted by the
lower bound on the complete \(q_j\) itself.

For the event payments, let \(S_{\mathrm W}\) be the
real carrying Wilson action on the bare plaquettes
with all corners in \(X\), and let \(G_j\) be the
sum of the kept nonnegative gauge costs at level
\(j\). In the prepared component the Wilson weight
is the running weight times the core cutoff. The
cutoff is one on the one-layer hull of the old core,
which contains the charged plaquette footprints.
The running weight there is a convex combination of
the stage weights. At an event of level \(t\) it
can run for at most \(\check R_t\) more levels before
freezing. The bounded comparison over this range
gives a coefficient at least \(c\theta_t\) on the
event's plaquettes. Indeed, an earlier vertex of
the running profile has reference weight at least
\(\theta_t-c_5\), and a later vertex can lose at most
\(\bar B\check R_t\) before freezing. At a final
preparatory stage the extra vertex is \(x_s+b_{s-1}\).
Its frame value is \(\theta_s+\theta_{s-2}-\theta_{s-1}\),
so the adjacent frame comparison gives the same lower
bound after allowing the three relative tube errors.
The cutoff leaves all these bounds intact on the charged
footprint. Set
\[
 E_{\rm pay}=a_\sharp\sum_t p_0(g_t)^2 n_t^{\rm cr}
              +\sum_t p_0(g_t)n_t^{\rm mk},
\]
where the counts concern the creations and marks
paid in this component. A drop-out's own mass
omits its failed root test. Charge it in a later
waiting window when that component supplies its
root witness on the compatible inverse branch.
With all the assigned certificates on that branch,
the event ledger gives
\[
 S_{\mathrm W}+\sum_j\theta_jG_j
       +\tfrac14\gamma_{\rm p}\sum_j\|A_j\|^2
 \ge E_{\rm pay}.
\]
Use the same real branch for the Wilson action and
the witnesses, and the same stored coordinates for
the gauge costs and the kept quadratics. The part of
the exponent involving these three weights satisfies
\[
 e^{-S_{\mathrm W}-\sum_j\theta_jG_j
                   -\frac12\sum_j\operatorname{Re}q_j}
 \le e^{-E_{\rm pay}}
       e^{-\frac14\gamma_{\rm p}\sum_j\|A_j\|^2}.
\]
The bound is pointwise before taking the supremum
over admitted evaluations. It leaves
\(\beta=\gamma_{\rm p}/4\) in the reference Gaussian
majorant above. The normalized shell measure needs
no further scalar factor. The old island factor
uses its pair-uniform Borel majorant at the admitted
inner core. Any reference-law denominator belongs
with the corresponding Wilson product and restoration
measure.

For the application to an actual later term, use its
current background at each factor. On the full
carrier of a retained strip, the
\hyperlink{hi:retained-strip-domain}{original-domain comparison}
gives the coercive bound using the original masks and
the two adjacent layer radii. A carried determinant
uses the full two blocks of the original elimination
row and the small-curvature bound there. Keep this
determinant separate from the inverse-density factor
of a raw row change. The restriction to each inner
island's complete core gives the argument for its
majorant.

For the layer estimate at a point of a birth carrier,
separate the datum plaquettes of the current background
according to their complete test footprints. The
retained ones with the prepared-datum support bound
have defect at most \(C e_q\) at level \(q\). Use
the original positive density bounds on the unchanged
exterior tests. Where these tests cover a weighted
neighborhood of radius \(D\) about the carrier, the
profile \eqref{mb:datum-profile}--\eqref{mb:weighted-maximum}
gives
\[
 C e_s+C\varepsilon_\star e^{-\delta D}
\]
in actual layer \(s\), with \(\varepsilon_\star\)
the fixed admission ceiling on the rest of the
datum. The first term follows by the exponential
decay across separated layers and the frame bound
on \(e_q/e_s\). When a waiting collar separates
the carrier from those other tests, its weighted
width of order \(M\check R\) absorbs the second
term. At carriers reaching the exterior halo we
must include the positive tests of that halo in
the read.

%
%
Assign the tests of the current carrying datum before
taking any marginal. At an unchanged ordinary place use
its original density cube and the full lower blocks for
each completed bond. A passed test supplies a local
configuration with the stated averages and a curvature
bound. The averaging estimate then bounds the holonomy
of the datum plaquette. For a plaquette in a stopped
component use its recorded departure place and the full
union of its corner blocks. When this union lies in the
tested inset, use the last applicable stage guard. On a
collar departure include the approximate test on the
collar links. A footprint extending outside that inset
uses the original density cube on its complete exterior
read. The ordinary defect bound is of order
\(\varepsilon_q\). For a remade shell it has the factor
\(C_L L_\flat^{2(q-k_0-1)_+}\), and a collar departure
has the factor \(c_{\rm f}L_\flat^{2(q-k_0)_+}\), with
\(k_0\) the cutoff between the uncompressed and
compressed stages. These powers are bounded by
\(C(\log\theta_q)^r\) by the frame and the stage widths.

Carry the resulting holonomy bound through the actual
later moves. A gauge draw, a frame change or a tree
gauge conjugates the holonomy. A late Haar refresh
leaves the old determining datum in the inner slot.
The fresh field is used by the outgoing rows. A
reference-law draw changes the datum itself. At an
island its reference has the exterior bounds of the
passed extendibility test and of the surrounding
density cubes. On the restored box, comparing the
four links with the reference and using the retained
cutoff gives the bound \(C(M)\varepsilon_i\). If a
determining plaquette is deleted, the newly made
layer uses its own regularity guard. In each case the
certificate concerns the data that enter the current
background.

%
In a quiet containing component that heals, the old
top test enlarged by four cells covers the outer six
cells of the component, since the complement of the
top member is inset by two cells. The final positive
test on the six-cell inset supplies the interior.
Their merger gives the top test of the new domain.
The reverse tree gauge preserves the resulting
plaquette bounds. The full original localized read
of each test is used even where it crosses the two
sets of test places.

On a joint inverse branch these tests act on
the assembled datum of the current background.
Their supports give the input of the weighted
profile in its current determining sequence.
Restrict that profile to each complete birth
carrier before applying the strip, pivot or
island bound. In using the upper integral for a
later term we must verify the full inverse reads
of these tests and the energy inequalities for
its recorded event witnesses.

The coordinate count is local.  A localized Gaussian marginal with
$m$ scalar coordinates has an $\exp(Cm)$ cost from its proved
comparison estimate, regardless of the number of exterior auxiliary
coordinates used to construct it.  Here
$m\le C(LM)^4N_c$, where $N_c$ counts the relevant coarse cubes.
If the count already denotes $m$, the factor $(LM)^4$ is not
inserted again.  In particular the reference cost does not acquire
an outside-volume multiplier.

The analytic and geometric conditions needed for an outer component
can now be stated without assuming its desired mass.  On every
carried support and on every argument in the supremum, require:
\begin{enumerate}
\item the exact exponent identity in \eqref{hi:tied-exponent}, or the evaluated identity in \eqref{hi:drop-factor}, with the correct retained Wilson weight;
\item the complete bound $|V_X|\le C_VN_0$ on the admitted triples,
including the actual source and counterterm normalization;
\item real admissible reads for the inner factors, whose differences
from substituted complex reads have been included in $V_X$;
\item the scalar normalizer estimates just proved and the actual
kept-form coercivity and reference-coordinate domination;
\item a disjoint ownership of the action, gauge and Gaussian pieces
used to pay for the event witnesses.
\end{enumerate}
The last requirement has the quantitative meaning
\begin{equation}
 A(\upsilon,U)+\sum_j\theta_j G_j(\xi)
      +\beta_1\sum_j|A_j|^2
       \ge\sum_{e}\Gamma_e,
 \qquad \Gamma_e=c_{\rm w}\theta_{t(e)}a_e^2,
 \label{hi:event-energy}
\end{equation}
leaving the positive reserve $\beta$ of \eqref{hi:reference-gaussian-integral} unspent.
For ordinary components $U$ is the real inner background; for a
drop-out the inequality must hold for every compatible real evaluation
in its fibre.  A bound only at one preferred background would not
justify the supremum in its mass.

These are finite-domain conditions, uniform in the admitted parameters
and in the outer index.  Each has specific mathematical content.  In
particular \eqref{hi:event-energy} requires actual witnesses at the stated levels,
while \eqref{hi:reference-gaussian-integral} requires an integral in the actual stored coordinates.
We next derive the mass from them; no mass estimate is included in
this list of conditions.

\paragraph{Which coupling pays for an old event.}
There is a useful profile check in \eqref{hi:event-energy}.  Suppose a plaquette first
enters the large-field region at level $t$.  The action rule stops
its running weight within $C_RR_t$ further levels.  On its owned
support the retained weight is a convex combination of inverse
couplings with indices no later than $t+C_RR_t$.  Earlier couplings
are at least $\theta_t-c_5$, and a later coupling in this range is
at least $\theta_t-C_6R_t-c_5$ by the bounded one-step drift.
For a sufficiently high floor, $R_t=O(\ell_t^r)=o(\theta_t)$,
and hence
\begin{equation}
                \upsilon(p)\ge c\theta_t.
 \label{hi:stopped-weight}
\end{equation}
This argument uses the action supports.  Continuing the weight all
the way to an arbitrarily later level could make it much smaller,
and the same conclusion would not follow from bounded drift alone.
The freeze rule is thus a local support statement, not an informal
appeal to approximate monotonicity.

\paragraph{Elementary witness estimates.}
For $W\in\operatorname{SU}(2)$ in the operator norm,
$1-\frac12\operatorname{tr}W=\frac12|W-1|^2$.  A failed
kept gauge test $|W-1|\ge a$ therefore supplies the gain
$\theta a^2/2$.  A failed physical fluctuation test of size $a$
has $|A^{\rm s}|\ge\theta^{1/2}a$, and its reserved Gaussian
form supplies $c\gamma\theta a^2$.

For a plaquette certificate, suppose the actual tested quantity
has the following bound on the failure event, with nonnegative
coefficients $T_{ep}$. Weighted Cauchy--Schwarz then proves
\begin{equation}
 \begin{gathered}
 a\le\sum_{p\in P_e}T_{ep}|\partial U(p)-1|,
 \qquad\Longrightarrow\\
 \sum_{p\in P_e}\upsilon(p)
              \bigl(1-\tfrac12\operatorname{tr}\partial U(p)\bigr)
   \ge {a^2\over2\sum_{p\in P_e}T_{ep}^2/\upsilon(p)}.
 \end{gathered}
 \label{hi:wilson-certificate}
\end{equation}
For an ordinary loop identity such a certificate follows by writing
the loop as a product of conjugated plaquettes and telescoping the
product; all other unitary factors have norm one.  For a nonlinear
block average the certificate and its uniform weighted row bound
must instead be proved for that actual average.  A flat linear row
estimate alone does not supply it.  Similarly, a failed existence
or extendibility test needs a quantitative implication saying that
all smaller local witnesses would put its data inside the admitted
background chart.  Its contrapositive then provides a witness.
Failure of an unquantified existence assertion is not itself an
energy lower bound.

%
%
%
For an outer mass, assign the Wilson part of a mark witness
on \(u\) to the plaquettes of \(C(u)^+\) with no corner in
its floor. Use this assignment when the outer Wilson action
contains that share. A stopped root can use this halo when
it is read in a containing operation.
For a mark at level \(m\), \(C(u)^+\) adds two layers of
\emph{unit cubes of level \(m\)}, of total thickness \(2L^m\)
in bare units. Distinct components at this level are separated
by a coarse cell of side \(a_m=M\check R_m L^m>4L^m\),
so their enlargements are disjoint. For incomparable marks
at different levels, follow the earlier enlargement in the
ordinary hull to the level of the later mark. It is contained
in the earlier mark's successor there. That component and the
later mark are separated by a coarse cell of the later level,
which exceeds the thickness of the later enlargement.

If a standing mark \(v\) of level \(m\) lies below a later
eligible mark \(u\) of floor level \(h\), then \(m<h\).
The ordinary step into level \(m+1\) contains \(C(v)^+\)
in the one-cell enlargement of the hull of \(Z_m\). If
\(h=m+1\), the eight-cell buffer of the ordinary step
therefore puts the full \(h\)-cell carrier of \(C(v)^+\)
in the floor. If \(h>m+1\), the neighborhood has already
entered \(Z_{h-1}\), whose hull again receives eight cells
before reaching \(\Omega_h\). The floor entries on the
containing component agree with those of this ordinary
step. In both cases the earlier share lies inside the floor.
The same argument applies to Wilson paying halos at birth
levels \(t<h\), which are contained in the ordinary \(Z_t\)
at birth. For a plaquette or approximate-fluctuation birth
at \(h\), the cell lies in \(P'_h\) or \(Q'_h\), with the
level-zero plaquette births also included in the first step.
Its one-cell paying halo leaves at least one coarse cell of
the original two-cell buffer to \(\Omega_h\). Thus it too
is in the floor. A fresh ordinary test outside a live
component and its margin has its Wilson share outside the
carried halo.

The stored gauge terms are assigned by their original
level and contour. Assign the Gaussian payment from the
scalar reserve \(\tfrac14\gamma_{\rm p}\sum_b|A_b|^2\)
of the complete form, using the actual failed coordinate.
Nested mark windows are disjoint by the quarantine at the
lower end of the waiting window. A creation at the first
level \(h\) of a later window uses gauge and fluctuation
coordinates of the step from \(h-1\) to \(h\), below the
coordinates \(j\ge h\) used by the mark. Incomparable nodes
have disjoint spatial shares in any common layer. Within
a one-step share, divide by the bounded multiplicity of
its cell witnesses before assigning the gains.

Cut the outer genealogy at its inner islands. Their
earlier witnesses have already been used in the inner
masses. For a still-standing second-class island \(Y\) of
level \(m\), the negative density test is on its restored
level-\(m\) field. Its determining datum includes every
top bond meeting \(Y\). The original entries through \(m\)
on \(Y\) survive, and any later eligible window containing
this island starts above \(m\). On the full interior carrier
of that test, the retained constraints therefore identify
the level-\(m\) averages of a compatible real evaluation
with the restored field.

%
Choose a density cube \(C\) in the six-cell inset of \(Y\)
where the test fails. The regularity criterion uses \(C\)
enlarged by four test layers, of side \(M_2\check R_m\) in
level-\(m\) units. Completing the averaging blocks of a
plaquette on that carrier costs at most one more such unit.
Since \(4M_2\check R_m+1<M\check R_m\), all these blocks
stay in \(Y\). If every level-\(m\) plaquette defect on
the carrier were less than \(c\varepsilon_m\), the local
regularity criterion would give the required orbit with
regularity radius \(B_3c\varepsilon_m\). Choose \(B_3c<1\)
and the weak-coupling floor so that \(c\varepsilon_m\le a_1\)
and \(B_3c\varepsilon_m\le\varepsilon_{\mathrm v}\le a_0\),
with the torus condition of the regularity criterion.
The curvature on the tested inset is then strictly within
its threshold. Thus a
failure, including undefinedness, gives a plaquette defect
of size at least \(c\varepsilon_m\). The four-block
certificate \eqref{eq:small-energy-min-certificate} gives
bare Wilson energy at least \(c'\varepsilon_m^2\) in \(Y\).
The action gives this region weight \(x_m\) at birth and
includes a still-standing second-class \(Y\) in the
level-\(m\) action region. The core taper is one on the
whole old-core carrier just located. Use the lower
stopped-weight comparison \eqref{hi:stopped-weight} there and
\(x_m\asymp\theta_m\). This gives
\(c_0\theta_m\varepsilon_m^2=c_0P_m^2\) in the carrying
Wilson action.

The inner core Wilson action is on its prepared inverse
field. These are distinct action factors in the exponent
identity. If the outer ledger
counts the island as \(b\) creation cells and also as
a mark, both credits come from the fresh-field witness.
Here \(b\) is uniformly bounded. From an energy
\(c_0P_m^2\) reserve \(a_\sharp bP_m^2\) for the creation
cells, with \(a_\sharp b\le c_0/2\), and use
\(P_m\le c_0P_m^2/2\) for the mark at weak coupling.

For the untouched one-step pieces, the tests and the averaging constraints give the witness estimates
for every compatible real evaluation.  Each created cell has a
witness within two cells, and a plaquette, gauge factor or
fluctuation coordinate is used by at most $5^4$ of these witnesses.
Divide the available energy by this multiplicity.  The plaquette
witnesses use the Wilson action on $N_{t,D}$; the approximate
fluctuation witnesses use that action and the gauge-fixing energy;
the remaining fluctuation, collar and rim witnesses use the scalar Gaussian reserve at their actual failed slots.  The threshold
comparisons therefore give
\[
 \frac12 A_{\mathrm{cr}}
 +\frac12\sum_{(t,D)}\theta_{t-1}G_{t,D}
 +\frac{\gamma_{\rm p}}4
        \sum_{(t,D)}\sum_{b\in\mathfrak X_{t,D}^{\rm s}}
                  |\mathcal A^{\rm s}_{t-1}(b)|^2
 \ \ge\ c\sum_{(t,\Delta)\in\operatorname{Cr}} P_t^2 .
\]
Here \(\mathfrak X_{t,D}^{\rm s}\) is the stored strip mask of
the piece and $A_{\mathrm{cr}}$ is the Wilson sum over the disjoint shares. The sums run over the selected untouched pieces,
$\operatorname{Cr}$ is their set of created cells, and the actual
kept-form coercivity and coupling comparisons are those used in
the one-step witness estimates.  The unused halves of the Wilson,
gauge and Gaussian energies remain available for the reference
integral.

The smallest stated mark witnesses can have size
$cR_t^{-4}\min_{j\ \mathrm{in\ window}}\delta'_j$.  With
$\delta'_j\asymp g_j\ell_j^{p_1}$ and the bounded comparison
of reference couplings over a window of $O(R_t)$ levels, this gives
\begin{equation}
       \Gamma_e\ge c\ell_t^{2p_1-8r}
       \quad\hbox{for a mark witness}.
 \label{hi:mark-gain-power}
\end{equation}
The comparison uses bounded drift over that window, not over the
entire history.  Assume $2p_1-8r>p$ and $p>6r$.  Then this
gain exceeds both a fixed multiple of the later cell costs and
$2P_t+1$, after the coupling floor is raised.  Creation witnesses
of size $cg_tP_t$ have gains bounded below by $cP_t^2$ and obey
the same conclusion.  For example the two strict exponent
inequalities are compatible with the other hierarchy
$p_1\le p-4r-1$ once $p$ is sufficiently large.  This arithmetic
does not establish the witness certificates or their ownership;
it proves what those certificates would pay for.

The root witness of a drop-out needs a separate comment.  Its
certificate can use $X^+$, whereas that drop-out's own evaluation
has a field only on $X$.  Its root gain cannot then be claimed in
its own mass.  A later containing component can use the gain if its record includes
the failed last test and its compatible field supplies the witness on
the full paying carrier in $X^+$.  Restricting the earlier mass to this additional
test can only decrease it.  Comparing the old creation reserve gives the sharper mass bound
below without an additional root gain.

\subsubsection{The pointwise envelope and the mass induction}

The reference-coordinate maps used here are fixed by the stored
record.  They may depend on other stored variables, including inner
variables, but not on a parameter still inside the supremum.  An
argument-dependent translation that makes each separate Gaussian
integral small would not bound the integral of the supremum.

Under the conditions stated above, split the positive outer action,
gauge and Gaussian energy into two parts.  Use the first to pay
$\sum_e\Gamma_e$ in \eqref{hi:event-energy}.  Use the second to construct the
nonnegative Borel reference envelope of \eqref{hi:reference-gaussian-integral}, or use \eqref{hi:reference-cutoff-gaussian} for the completed operation with unmarked core. Bound the scalar
normalizers and the complete local energy by $e^{C_7V^\circ}$.
For every inner operation replace its actual real admissible read
pointwise by its strong real majorant in \eqref{hi:real-complex-majorants}.  The result is
\begin{equation}
 \boxed{
 m_\omega(\xi)
  \le h^{-f(\omega^\circ)}
       e^{C_7V^\circ-\sum_e\Gamma_e}
       g_\omega(\xi^\circ;\xi_{\rm in})
          \prod_y {m^{\mathbb R}_{\omega'_y}(\xi'_y)\over h},
 \quad
 \sup_{\xi_{\rm in}}\int g_\omega\,d\mu^\circ
                              \le e^{C_5V^\circ}.}
 \label{hi:actual-pointwise-envelope}
\end{equation}
The same formula is required over the compatible evaluation fibre
in a drop-out.  The envelope is constructed before that supremum:
every evaluation satisfies the same energy certificates, and every
complete free form satisfies the same coercivity bound.  Its stored
coordinates, reference measure and normalizers have been identified
as just specified.  These facts, not a pointwise gas expansion at
each evaluation, prove \eqref{hi:actual-pointwise-envelope}.

The displayed uniform conditional integral is enough even if the
envelope depends on inner variables.  When the actual reference
construction makes it a function of outer variables alone, this
condition follows immediately.  In the more general case it follows
from the uniform triangular calculation in \eqref{hi:reference-gaussian-integral}.  Its constant must
be independent of all inner histories and all their variables.
An estimate only after averaging some other exterior law would not
be the estimate in \eqref{hi:actual-pointwise-envelope}.

We prove explicitly the induction step.  Fix the outer index
$\omega^\circ$.  Its inner islands, their levels and their domains
are fixed.  The remaining histories are the independent local
records $\omega'_y$.  The retained operations have the corresponding
product reference measures, as in \eqref{hi:local-operation}, and
$v^\natural(\omega)=V^\circ(\omega^\circ)+
\sum_yv^\natural(\omega'_y)$ exactly.  Suppose the earlier
strong real masses with coefficient $c$ are at most one.
Multiply \eqref{hi:actual-pointwise-envelope} by the weight and use its Borel envelope.  Tonelli
allows the outer variables to be integrated first.  The bound on
that integral is uniform in the inner variables.  The remaining
nonnegative integrals and sums separate into the earlier masses.
Each is at most one.  We obtain
\begin{equation}
 {1\over h}\sum_{\omega\ \mathrm{over}\ \omega^\circ}
      e^{cv^\natural(\omega)}\int^*m_\omega\,d\mu_\omega
 \le h^{-1-f(\omega^\circ)}
          e^{(C_5+C_7+c)V^\circ-\sum_e\Gamma_e}.
 \label{hi:inner-history-integration}
\end{equation}
No measurable choice of a maximizing pair or evaluation has entered
this calculation.  No interchange of supremum and integral was
made.  The earlier masses include their own frame factors and
their own histories, exactly once.

The scalar reserve is applied with
$c_1=C_5+C_7+c+c_{\rm e}$.  The counting estimate applies to
the outer tree cut at its inner islands.  Its cut leaves have their
bounded number of boundary seeds.  Their last tests lie outside the
inner masses and supply the assigned creation gains.  Under the
witness hierarchy in \eqref{hi:mark-gain-power}, or directly under the strengthened
scalar conditions, the coefficient of each event can pay its
allocated part of $c_1V^\circ$ and leave the specified reserve.
In particular it is sufficient to have
\begin{equation}
 \sum_e\Gamma_e-c_1V^\circ
        \ge\sum_{e\ \mathrm{creation}}(2P_{t(e)}+1).
 \label{hi:outer-event-reserve}
\end{equation}
For the stronger event-by-event form one leaves $2P_{t(e)}+1$
also at every mark whose witness is available on this outer tree.
This stronger form follows from the actual witness powers, not
from treating a conservative credit of size $P_t$ as if it were a
larger gain.  The weaker form in \eqref{hi:outer-event-reserve} already suffices for the mass.

Here is the exact comparison that removes a possible constant loss.
For $t\le i$ the bounded-error monotonicity of the reference flow
gives, by the mean value theorem,
\[
 P_t\ge P_i-C_8{\ell_i^{p-1}\over\theta_i}.
\]
The error tends to zero.  Raise the common floor so that it is at
most $1/4$.  The outer tree has a creation leaf, and its reserved
$2P_t+1$ is then at least $2P_i$.  Thus, for any fixed
$0<a_\star\le2$, \eqref{hi:inner-history-integration}--\eqref{hi:outer-event-reserve} and \eqref{hi:outer-entropy} yield
\begin{equation}
 \boxed{\mathsf M_c(i,X;v^\natural)
                     \le e^{-a_\star p_0(g_i)}.}
 \label{hi:weighted-mass-conclusion}
\end{equation}
This argument proves the same constants at the new level; it does
not produce a factor $e$ at every step.  It proves the complex-pair
mass from the previous strong real masses.  Restricting the pair
supremum to the real core recovers that induction hypothesis at
the new level.  At the first level there are no earlier islands,
so the same proof starts with an empty product equal to one.

For a drop-out the same proof holds with every complete form in
its evaluated fibre and with its retained Wilson weight.  The
version including its own failed last test is a restriction and
has no larger mass.  The reserve argument with the extra unit per
old creation also gives the exact $e^{-2P_i}$ bound when that value
is desired.  Its proof does not use an unavailable root witness on
the smaller domain $X$.

\paragraph{The induction has a stronger input than stability.}
For emphasis, the induction statement is
\begin{equation}
 {1\over h}\sum_\omega e^{cv^\natural(\omega)}
       \int^*\sup_{\text{admitted real reads}}
                     e^{\Phi_\omega}\,d\mu_\omega\le1,
 \label{hi:strong-real-induction}
\end{equation}
not a bound on an integral at each fixed read.  It is also not a
bound on a final Gibbs partition function.  The real statement is
propagated by the pointwise structural envelope, while its complex
counterpart is an output of the same step.  The proof uses the
strong norm exactly where the inner histories are separated in
\eqref{hi:inner-history-integration}.  Stating this quantifier prevents an otherwise hidden
strengthening in the induction.

\subsubsection{Perturbations paid for by the stored-site reserve}

The same volume weight pays for an explicitly bounded source excess.
Suppose, on the full domain used in the supremum,
\begin{equation}
            |\Psi_\omega|\le\bar f\,n^\sharp(\omega).
 \label{hi:source-excess-budget}
\end{equation}
If $c\ell_j\ge2\bar f$ for all relevant levels, \eqref{hi:read-site-budget} gives
$|\Psi_\omega|\le cv^\natural(\omega)$.  Therefore, for a
bounded insertion $F$,
\begin{equation}
 {1\over h}\sum_\omega\int^*
       \sup\bigl|F e^{\Phi_\omega+\Psi_\omega}\bigr|
                 \,d\mu_\omega
       \le e^{-a_\star p_0(g_i)}\sup|F|.
 \label{hi:perturbed-mass}
\end{equation}
This is pointwise domination by the already weighted mass.  It
requires the same normalizers, reference measures, test supports
and argument domains on both sides.  A source that moves a sharp
test or changes the reference measure must first be rewritten in
that fixed representation; it is not covered by \eqref{hi:source-excess-budget} by name.

The weight may also be changed.  If $0\le v'\le v^\natural$,
monotonicity gives the mass bound with $v'$.  If an owned window
volume satisfies $v^{\rm win}\le2v^\natural$, use the original
mass with coefficient $2c$ to bound its weight with coefficient
$c$.  This comparison must have a fixed multiplicity.  The formula in \eqref{hi:recounting-obstruction} explains why there is no analogous statement for arbitrary
repeated-past volumes.

One can verify \eqref{hi:source-excess-budget} from local carrier information.  The following
criterion records the constant terms as well as the oscillations.
At level $j$, let $\mathcal A_j(\omega)$ be its owned stored
sites, with the fixed read collars assigned as in \eqref{hi:read-site-budget}.  Let
$\psi_f(z;\cdot)$ have carrier $R_f$ of that level, and let a
reference argument $p_f^0$ in its admitted carrier core be fixed.
Suppose
\begin{equation}
 |\psi_f(z;p_f^0)|+
     \sup_{p\ \mathrm{admitted}}|\psi_f(z;p)-\psi_f(z;p_f^0)|
         \le b_f(z),\qquad
 \sup_{x\in\mathcal A_j}
          \sum_{f:x\in R_f}{b_f(z)\over |R_f|}\le\bar f.
 \label{hi:carrier-site-criterion}
\end{equation}
Every carrier here has a nonempty owned site set $R_f$; using a
coarser root convention instead changes the displayed density
bound by its proved local counting constant.  Sum first over the
sites of $R_f$, then over the carriers.  Absolute convergence or
finite sums gives
$\sum_f\sup|\psi_f|\le\bar f\sum_j|\mathcal A_j|$.
The total source excess consequently satisfies \eqref{hi:source-excess-budget}, with its
fixed dimensional factors included in $\bar f$.

For a usual rooted decay criterion, assume each carrier has a
chosen owned root and tree length $d(R)$, and its total rooted
weight is at most $B e^{-\kappa d(R)}$. Assume, at any fixed owned
site $x$, that at most $C_0C_1^m$ carrier/root choices of length
$m$ contribute with $x\in R_f$, and that $\kappa>\log C_1$.
This is an incidence count at $x$, not merely a count of carriers
whose original root was $x$. Then the per-site sum is bounded by
$BC_0/(1-C_1e^{-\kappa})$.  A polynomial root multiplicity is
absorbed by a strictly smaller retained exponent: for
$\varepsilon>0$, $m^q e^{-\varepsilon m}$ has a finite
maximum, obtained by differentiating $t^qe^{-\varepsilon t}$.
This proves the carrier criterion with explicit rate expenditure.
It does not verify the amplitudes of the actual formed members.

The constant term in \eqref{hi:carrier-site-criterion} cannot be omitted.  A function constant
in its field arguments has zero oscillation and can have arbitrarily
large parameter-dependent value.  In the actual source class its
value is fixed by the vacuum, scalar-action and normalizer
readouts.  Those readouts must be bounded along with the
oscillation.  Gauge invariance does not bound an arbitrary constant.

%
%
For the actual recursion let \(V_X^{\rm col}\) be the
interaction before re-evaluating the kept inner factors.
Write the exponent of an inner source perturbation as its
stored-variable part plus its carrier functions \(\psi_R\).
The two systems use the same response map for a kept
trunk. Its contribution to the difference is
\[
 T_R^1[\psi_R]=\psi_R(q(\xi)|_R)+\delta_R[\psi_R].
\]
The first term rejoins the stored-variable part of that
inner exponent at its admitted real pair. Keep the second
term in the new excess. The resulting identity is
\begin{equation}
 \begin{aligned}
 \Psi_\omega
   &=\Psi^{\rm new}_\omega
        +\sum_y\Psi_{\omega'_y}(\xi'_y;\mathbf q_y(\xi)),\\
 \Psi^{\rm new}_\omega
   &=A(\zeta''(z)\zeta_1\mathbf1_X,U''(\xi))
     +V_X^{{\rm col},z}-V_X^{{\rm col},0}
     +\sum_R\delta_R[\psi_R].
 \end{aligned}
 \label{hi:actual-source-recursion}
\end{equation}
The reference measures, the bare stored fluctuation slots
and the reference quadratic \(q_X\) are common to the two
systems. The source profile \(\zeta''\) multiplies the
core cutoff \(\zeta_1\). Where that cutoff is nonzero
the real inner datum gives curvature
\(C\check R_l^2\varepsilon_l\eta_l^2\) on layer \(l\).
At weak coupling the square of this bound is
\(O(\eta_l^4)\). The \(O(\eta_l^{-4})\) bare plaquettes
per site of that layer give a cost \(C_W\bar c N_0\)
when \(|\zeta''|\le\bar c\). The bare layer uses the
uniform Wilson bound.

Apply the \hyperlink{hi:formed-collar-majorant}{formed
collar bound} to each of \(V_X^{{\rm col},z}\) and
\(V_X^{{\rm col},0}\). Their absolute bounds include
the unpaired plain values at the trivial pair and give
\((c_V-1)(1+\bar c)N_0\) for each. This application
uses the full common field and stored-slot ranges of
the two systems. For \(\delta_R[\psi_R]\), subtract
the value at \(q(\xi)|_R\) and use the same response
disc as for the unsourced kept factor. Its bound is
the product of the response amplitude and the admitted
carrier bound for \(\psi_R\). A term constant in the
pair cancels in this difference. The weighted carrier
sum for the kept factors therefore applies to these
differences with the bound for the source carriers.
In particular the resulting sum is at most
\(2(1+\bar c)N_0\) once the weak-coupling ceiling
is chosen with that carrier bound fixed. The power
margin is the one in the kept-trunk estimate, before
division by \(N_0\ge 2^{-5}M^4\ell_t^{5r}\).
Hence
\[
 |\Psi^{\rm new}_\omega|
       \le [C_W\bar c+2c_V(1+\bar c)]N_0.
\]
Add the older complete exponents at their own sites.
The additivity of the fresh and older counts takes
the maximum of their coefficients. For the source
induction we must check the admitted carrier bounds
of the actual generated members on the enlarged
source lines, including the stored switches. The
argument uses these bounds and the full interaction
bound on the same domains.

\subsubsection{Complex parameters with the inner Wilson weight fixed}

The preceding proof has a parameter version, but its domain must be
specified before taking suprema.  Fix the real reference sequence
$\theta$ and the real physical weight profile $x$.  Freeze all
characteristic cutoffs, stored-coordinate charts, reference measures
and the inner Wilson weight at those real data.  Let the complex
parameters $y_j$ satisfy
\begin{equation}
 |x_j-\theta_j|\le t_j,\qquad |y_j-x_j|\le\bar c,
                         \qquad \tau_j=y_j/\theta_j.
 \label{hi:parameter-tube}
\end{equation}
Here $t_j$ is the radius of the frozen tube at level $j$.
The ordinary strip uses the ratio disc of radius
$2(t_j+\bar c)/\theta_j$ about $1$. This radius tends to zero
uniformly as the inverse-coupling floor increases. In particular
it lies within the half-disc used for the scalar logarithms.
The actual ordinary presentation determines where $y_j$ occurs.
In particular its complete free-coordinate quadratic form is the
one prescribed by \eqref{hi:compensated-quadratic}, including a compensation if a normalized
$\tau_jK_j$ Gaussian is used.  The weight in \eqref{hi:inner-wilson} is still the
specified real weight $x''$.  Putting a complex weight on a large
real Wilson variable and estimating it as if it were the same
positive reserve is a different operation.

For a complete matrix on a fixed real free-coordinate space
there is a determinant calculation on a half-disc. Under the
real hypotheses preceding \eqref{hi:complete-sector-bound}, define
$B=K^{-1/2}K_AK^{-1/2}$.  Since $K_G\ge0$, one has $B\le I$.
Since $K_A\ge-\gamma I$ and $K\ge\gamma I$, one has
$B\ge-\gamma K^{-1}\ge-I$.  Thus $\|B\|\le1$ and
\begin{equation}
 Q_\tau=K^{1/2}[I+(\tau-1)B]K^{1/2},\qquad
 \operatorname{Re}Q_\tau\ge K/2,\qquad
                       \|Q_\tau^{-1}\|\le2/\gamma.
 \label{hi:relative-kept-sector}
\end{equation}
The matrix logarithm of $I+(\tau-1)B$ is its uniformly convergent
power series on $|\tau-1|\le1/2$.  Consequently the determinant
and its square root have coherent analytic branches agreeing with
the positive branches at one.  In $d$ free scalar coordinates,
the absolute determinant ratio is between $2^{-d}$ and
$(3/2)^d$.  The normalized Gaussian at each fixed real $K$ has
total variation at most $3^{d/2}$: its determinant costs
$(3/2)^{d/2}$ and replacing the real precision $K$ by $K/2$ costs
$2^{d/2}$ on integration.

For a supremum over varying backgrounds, this last fixed-matrix
integral is not the argument.  If $\gamma I\le K\le\Gamma I$
uniformly on those backgrounds, \eqref{hi:relative-kept-sector} and the determinant upper
bound instead give the common pointwise envelope
$C^d\Gamma^{d/2}e^{-\gamma|b|^2/4}$ in their same stored
coordinates.  Its integral is at most
$(C\Gamma/\gamma)^{d/2}$.  Alternatively one uses the proved
local comparison with a fixed positive reference matrix.  This
additional uniformity is precisely what places a parameter
supremum inside \eqref{hi:reference-gaussian-integral}.  The classical flat gap is one realized
matrix case.  A one-step nonflat estimate can verify the same
criterion by bounding the negative part of $K_A$ in the actual
free coordinates; neither case alone verifies it on every
retained layer.

At complex parameters the pair-independent factor of the exponent
need no longer be real.  The proof uses its modulus, with the
real positive energy and the fixed reference envelope just
constructed.  Suppose the actual local formations and read maps
are jointly holomorphic on \eqref{hi:parameter-tube}, with their cores admitted
throughout, and their complete energy bound is
$C_V(1+\bar c)N_0$ there.  The scalar gauge branch is supplied
by \eqref{hi:gauge-normalizer}; the pivot branches by \eqref{hi:pivot-factor}; the complete Gaussian
forms by \eqref{hi:relative-kept-sector} where its real matrix
hypotheses apply, and by the original-strip estimate for the
kept chains on their admitted reads.
All their costs are then bounded in modulus by the same type of
$e^{CV^\circ}$ factor.

Apply the mass proof to the norm with the parameter supremum
included from the start.  Its real part refers to real \emph{field
reads}, while the parameters can be complex:
\begin{equation}
 m^{\mathbb R,\mathcal D}_\omega(\xi)
     =\sup_{y\in\mathcal D}
        \sup_{\mathbf p\ \mathrm{real\ admitted}}
                                |e^{\Phi_\omega(\xi;\mathbf p,y)}|.
 \label{hi:real-field-complex-parameter}
\end{equation}
The parameter domains at earlier levels are the corresponding
prefix domains of the same reference frame and initial weight.
This prefix condition is essential for using the earlier norm in
\eqref{hi:inner-history-integration}.  It is not a conclusion of separate holomorphic
presentations with different frozen frames.

Uniform primitive estimates give \eqref{hi:actual-pointwise-envelope} with these parameter
suprema, so the same induction proves \eqref{hi:weighted-mass-conclusion} uniformly in $y$.
The floor is raised for the new local-energy and Gaussian cost
constants.  If the former is $C_V(1+\bar c)$, its stated dependence
is preserved in the floor.  This is a proof from uniform finite
domination.  It is not analytic continuation of an inequality
known only on the real parameter slice.

The positivity used to pay for old failed plaquettes is exactly
why the inner Wilson weight was frozen.  In a variable-profile
source presentation, physical Wilson perturbations are instead
carried as specified older potential members.  The equality of
that presentation with an original numerator and denominator is
an exact formation identity with its own counterterm rule.  A
complex parameter in \eqref{hi:parameter-tube} cannot be declared a physical source
in every factor merely because both families agree at zero.

\paragraph{Holomorphy of the actual integrals.}
Let $\mathcal D_0$ be a compact parameter polydisc strictly inside
the admitted domain.  The pointwise integrands are holomorphic
there, all reference measures are parameter independent, and the
mass proof supplies one integrable Borel majorant on the full
product of stored variables.  Integrate the one-variable contour
identity of each integrand around a triangle in any coordinate;
Fubini is justified by that majorant.  The integrated function is
continuous by dominated convergence and has zero contour
integrals.  Morera's theorem gives holomorphy in that coordinate.
The same bound and the iterated Cauchy formula give joint
holomorphy and its power series.  Countable sums of histories
are handled by the summable common majorant.  It is the reference
measure and not just the notation $d\mu$ that must be independent
of the parameter in this argument.

For a polydisc with radii $r_1,\ldots,r_m$ and a bounded holomorphic
numerator insertion $F$, Cauchy's formula gives
\begin{equation}
 |\partial_1^{a_1}\cdots\partial_m^{a_m}I(0)|
   \le\bigl(\prod_{\nu=1}^m a_\nu!r_\nu^{-a_\nu}\bigr)
                        e^{-a_\star P_i}\sup|F|.
 \label{hi:mass-cauchy}
\end{equation}
This differentiates the complete numerator with its actual
normalizers.  A normalized expectation additionally differentiates
its denominator.  It is not obtained by differentiating only
the newly inserted local member.

For example, if an actual one-zone perturbation has the uniform
bound
\begin{equation}
 |\Psi_\omega(z)|\le C|z|p_n n^\sharp(\omega),
 \qquad |z|<r_n,\qquad
 r_n\le {c\ell_{\min}\over2Cp_n},
 \label{hi:one-zone-test}
\end{equation}
and all its fields remain in the same admitted domains, \eqref{hi:perturbed-mass}
and \eqref{hi:mass-cauchy} prove the corresponding complex mass and derivative
estimate.  This specifies a concrete formation estimate to test:
the coefficient $C$ must be uniform in the histories, stored
fields, levels, zones and allowed reads.  A bound for one chosen
background or a source function before it is formed is not
\eqref{hi:one-zone-test}.

For several perturbations, a per-site criterion can contain
$\sum_\nu |z_\nu|p_\nu$ or a smaller local sum, depending on
their supports.  The admitted product radii must satisfy that
actual criterion.  A one-source disc alone does not authorize
an equally large product disc for arbitrarily many overlapping
sources.  Disjoint local supports and a proved local carrier
bound can supply such a product disc; its justification is the
local bound, not one-variable Cauchy.

\subsubsection{Consequences for stored volume, ages and marked integrals}

Fix $c>0$ at which the mass bound holds.  Any subcollection of
histories with $v^\natural\ge T$ has unweighted mass at most
\begin{equation}
 e^{-a_\star P_i-cT}.
 \label{hi:volume-tail}
\end{equation}
This follows pointwise from $e^{cv^\natural}\ge e^{cT}$ and
monotonicity of the upper integral.  More generally the power
series of the exponential gives
\begin{equation}
 {1\over h}\sum_\omega\int^*
       (v^\natural(\omega))^q m_\omega\,d\mu_\omega
     \le q!c^{-q}e^{-a_\star P_i},\qquad q=0,1,\ldots.
 \label{hi:volume-moments}
\end{equation}
For $q=0$ use the unweighted mass directly.  For $q>0$ use
$(cv)^q/q!\le e^{cv}$.  Combining this with \eqref{hi:read-site-budget} gives the
site-count moments with $c^{-1}$ replaced by
$2/(c\ell_{\min})$.

A path of length $A$ in the recorded genealogy contains a node
at every one of those levels.  Since each such node has at least
one cell, $v^\natural\ge\ell_{\min}A$.  Therefore its
recorded age has exponential tails with exponent
$c\ell_{\min}$.  For any fixed $b>0$ the age-weighted sum
with factor $L^{bA}$ is finite once
$c\ell_{\min}>b\log L$, by summing the geometric series from
\eqref{hi:volume-tail}.  This age is the length inside the recorded genealogy.
The number of later levels for which an already formed mark is
carried is not part of its unchanged local record; that later
duration is controlled by the live-component reserve and the
global label comparison, not by silently increasing its local
volume at every read.

These bounds control stored counts, not every unbounded function
of a stored field.  Small mass alone would allow an integrand
concentrated at field magnitude $T$ with a second moment of
order $T^2$ times that mass.  Field moments need the Gaussian
reserve or their own pullback estimate.  If the pointwise proof
retains $e^{-\beta|A|^2}$, then
\begin{equation}
 |A_b|^q e^{-\beta|A|^2}
       \le\left({q\over e\beta}\right)^{q/2}
                    e^{-\beta|A|^2/2}
       \le\beta^{-q/2}\sqrt{q!}\,
                    e^{-\beta|A|^2/2}.
 \label{hi:gaussian-mark}
\end{equation}
The maximum of $t^qe^{-\beta t^2/2}$ occurs at
$t^2=q/\beta$; the factorial estimate follows from
$q^q\le e^qq!$.  Run the same reference calculation with
$\beta/2$.  Its extra factor $2^{d/2}$ is another local volume
cost, independent of $q$.  With the floor chosen for that
unspent reserve, the marked mass has the prefactor in \eqref{hi:gaussian-mark}.
For a physical observable the coefficients and the map from its
field to $A_b$ must still be tracked, including their powers of
the scale ratio.  A small bad-event probability alone does not
give this marked estimate.

These consequences also explain what the volume weight can
contribute to a large-field extraction remainder.  It supplies
arbitrary fixed exponential weights in the recorded age after
the corresponding floor is chosen.  It supplies polynomial
stored-coordinate moments when a spare Gaussian fraction is
kept.  It does not identify a resummed positive conditional
measure, its boundary terms, or its exterior force.  Those are
different estimates in an interacting covariance calculation.

\subsubsection{Formation, reading and loading at later levels}

We make the quantifiers of transfer explicit.  A factor is formed
at level $i$, with a definite local history, stored-coordinate
space, support and normalizer.  For each later level $k\ge i$
at which its mark stands in a term, let
$\operatorname{Read}_{k,i,X}$ denote the actual argument supplied
to that same factor.  The coverage assertion needed is
\begin{equation}
 \begin{cases}
 \xi_X\in S_{\omega_X},\quad
 \operatorname{Read}_{k,i,X}\in\mathcal U_i(X^\sim)
                      \text{ and is real},&o=\mathsf2,\\
 \operatorname{Read}_{k,i,X}\in
                  \mathcal K^\sharp(\omega_X,\xi_X),&o=\mathsf d,
 \end{cases}
 \label{hi:read-coverage}
\end{equation}
for every admissible later history and every configuration on
its carried support.  It is a pointwise statement, before taking
an exterior expectation.  In the drop-out case it includes every
average in \eqref{hi:compatible-fibre}, every recorded test, the prescribed variant
of the last test, the inner-island cores and all one-sided
matrix and Jacobian bounds.  In the island case its reality is
what allows the strong real induction in the new formation.

The finite geometric part can be checked as an invariant of the
construction.  A later operation outside the component does
not change its stored variables.  An operation containing it
acts above the levels protected by its waiting floor, and its
determining entries use the actual re-entry field described
after \eqref{hi:action-sequence}.  The action entries keep the tests with their
recorded ownership.  Restricting a history before that later
operation consequently retains the same local history and the
same operation.  This proves the record statement under the
specified local separation and window rules.  Admission of
the resulting background to the analytic core, and the actual
matrix bounds in \eqref{hi:read-coverage}, are the remaining analytic part of
coverage; the invariant does not prove them by itself.

\paragraph{A stop at the prospective regularity test.}
With $n$ stages completed, let $j=i_\flat+n$. A stop
before the collar gauge and the removal keeps the ordinary
prefix of index $n$. Where the preceding guard is already
present, it remains a passing test of this prefix. For a
datum plaquette of level $q$, take the union of the
endpoint blocks at that level. Read the ordinary density
test at its original cube if $q>j$, if the top action
region there is untouched, or if this footprint leaves
the region covered by the preceding stage test. On the
covered footprints use the passing guard $\chi_i^{(n)}$.
An untouched prefix reads its original density tests
throughout. The old core keeps its own tests and the
averaging constraints remain in the original rows.

The ordinary-prefix real admission applies with these
tests and with the old determining list. It gives the
layer profile of the old background, and the half-radius
bounds on the corresponding earlier spaces, including
those of carried drop-outs. This uses the tests of the
prefix without the newly failed test. The stage has added
neither a gauge draw nor a pivot elimination. Its
one-step factors remain in their birth coordinates and
there is no normalizer from the interrupted stage.
For a later evaluation of this stopped fibre, the full
footprints of these tests and the same original pivot
and walk carriers must occur in the compatible read.

%
%
%
\paragraph{The datum at an early prepared seam.}
A complete test footprint can meet both sides of
a prepared face, even when the bare guard uses
plaquettes contained in its individual shells.
Use the coarse datum at the preceding averaging
fibre to control such a footprint.
Consider promotion to level \(q=j+1\) in the
uncompressed range \(q\le k_0\). In the quiet
geometry \(G_j\subset H_q\), and the new face is
the boundary of \(H_q^{[3a_q]}\), with
\(a_q=M\check R_qL^q\). Every new top test based on
this promotion has a row outside the new inner box.
Its complete blocks extend only \(O(L^q)\) inward,
so the footprint remains outside \(G_j\).
A footprint meeting the original exterior wall
of \(\Omega_q\) is read on the full halo of
the original level-\(q\) density test, before
that test is removed. The other new footprints
lie in the old exterior
\(\Omega_q^c\setminus Z''_j\). At the first
promotion their full carrier is in the region
\(\Omega_{i_\flat}\setminus\Omega_{i_\flat+1}\)
of the stage-zero
test. At later promotions the preceding exterior
guard reads the whole carrier. The local preceding
background reproduces the old determining
level-\(j\) rows there. Thus the guard bounds
the curvature of that background by
\(C\varepsilon_j L^{-2j}\) on the full carrier.
The new coarse rows are its averages by the
original constraints. The plaquette estimate gives
\(C_L\varepsilon_j\le C'_L e_q\), including a
plaquette that straddles the new prepared face.
The surviving tests of level \(j\) use a subset of
the previous rows and the same lower completions.
The lower tests are unchanged. If level \(q+1\)
is present, each tested plaquette there has a first-kind edge
meeting the original \(\Omega_{q+1}\). Its whole
carrier is within \(O(L^{q+1})\) of that member.
The quiet gap of \(8a_{q+1}\) keeps this carrier
inside the original \(\Omega_q\). Its completion
therefore uses the same rows as before. Higher
tests read neither the changed member nor its
immediate completions. Each promotion
uses its own positive tests and a single averaging
bound. Carry the new coarse holonomy bound
as a support condition before marginalizing
the old fibre.

%
%
%
\paragraph{The datum at a compressed prepared seam.}
For \(q=j+1>k_0\), take a level-\(q\) test whose
complete footprint meets \(\partial Z''_q\).
The footprint extends at most \(2L^q\) from this
face. The separation from \(Z''_j\) is at least
\(ML^q\). The face is at least \(5ML^i\) from
the old \(\Omega_{k_0+1}\), and the original exterior
has its clearance from \(\partial X\). Thus the
entire footprint lies in
\[
 X\cap\bigl(\Omega_{k_0+1}^c\setminus Z''_j\bigr),
\]
in the layer of index \(j\) before promotion.
At the first compressed transition \(j=k_0\) this
is the single exterior region of the preceding guard.
At the other transitions it is its deep exterior region.
In both cases its bound is at most
\(F_j\varepsilon_jL^{-2j}\), with
\(F_j=L_\flat^{2(j-k_0-1)_+}\).

The local preceding background has the actual old
level-\(j\) datum on this footprint. Every
first-kind link of the new test is promoted there,
and every second-kind link uses the same retained
level-\(j\) blocks. Their common coarse values are
exactly the one-step averages. Hence the test
holonomy is the level-\(q\) plaquette of the
averages of this local preceding background.
The guard on the full bare carrier and the averaging
estimate give a defect at most \(C_L F_j\varepsilon_j\).
Since \(L_\flat<L\) and
\(L^{i-k_0-1}=\check R_{k_0+1}\), this is at most
\(C_L e_q\). The preceding guard is available
before deletion, so the bound can be carried in
the coarse support after the integration. This
also applies to the last prepared face before
the final collar integration.

At a new level-\(q\) footprint away from the
prepared face, use the same calculation when
it is contained in a single preceding exterior
region. The factors in those clauses are at
most \(F_j\). A remaining new footprint meets
an original face \(\partial\Omega_l\), with
\(k_0+1\le l\le q\). Consecutive original walls
of levels \(l,l+1\) are separated by at least
\(8a_{l+1}\), so the footprint meets only one
wall level.
Use the density test at that face before it is
removed.

%
%
%
%
Put \(u=ML^i\) in bare units. Start with a tube
of width \(3u/2\) about the original face of
level \(l\). Its level-\(l\) plaquettes are
covered by the original density test with its
full halo. Before deletion this gives
\(|\partial V_l-\mathbf1|\le C\varepsilon_l\) there.
For each subsequent promotion to level \(p>l\)
reduce the width by \(4L^p\) to keep complete
backward averaging carriers. The total width
lost is less than \(8L^i\), so more than \(u\)
remains. These tubes stay clear of every
prepared inner face and every other original
face. In a later promotion at level \(p\) the
tube is outside the original \(\Omega_p\) and
inside the preceding prepared member. Thus
its level-\(p\) field is the exact average of
the retained level-\((p-1)\) field there.

On the joint graph before eliminating these
averaging constraints, the level-\(q\) field
on the remaining tube is the iterated average
of the same \(V_l\). Any frame changes act at
the block endpoints and preserve its plaquette
defect. Apply the full-carrier estimate of
\cref{prop:averaging-1a-1-4} once at depth
\(q-l\). This gives a bound
\[
 C L^{2(q-l)}\varepsilon_l
              \le C'\check R_q^2\varepsilon_q.
\]
The inequality uses \(L^{i-k_0-1}=\check R_{k_0+1}\)
and the frame comparison within the compressed
interval. At weak coupling the coarsest scaled
input is inside the averaging radius. This
bound is a condition on the coarse holonomies
in the tube. Insert it on the joint graph and
carry it through the exact pushforwards. When
\(q=l\) the bound follows directly from the
original positive density test. In this way the
quantitative datum bound uses the original
density test at its actual level of formation.

At a final regularity stop, $n=\mathsf N_i-1$. If the
broader waiting convention admits the component at
the lower endpoint of a later window, it is a partial
re-entry. Its level-$i$ density
test uses its original action region at that level, which
agrees there with the determining region. The top field
is the old datum on $\Gamma_i$. For an unprescribed
level-$i$ bond $c$ with both ends in $\Omega_i^c$, the
surviving one-step delta gives
\[
 V_i(c)=\mathcal M(V_{i-1})(c).
\]
On a protected footprint containing the two full blocks
of $c$, the later read has the same level-$(i-1)$ rows.
Composition of the averages therefore gives this same
value of $V_i(c)$ in the later read. The ordinary density
tests and the local stage guard are evaluated from the
resulting old datum on their own footprints. Deeper
remade levels still use their separate determining lists.
For a strip operator this argument applies when the
later structured address also contains its complete birth
carriers. The footprint needed for an energy witness can
be smaller than that operator address.

On a fixed incoming prefix and outgoing field, the old
and inset versions of the early stop have the same fibre
after omitting the root's negative indicator in its own
mass. The gauge factors, averaging graphs, retained
measure and weight are those of the common prefix.
A later read with the new failure imposed restricts this
common fibre. Any absolute reference bound on the
root-omitted fibre thus also bounds the restricted read.
The uniform joint-core bound on the common fibre and
the fresh payment for the new failure are separate estimates.

\paragraph{A stop between removal and collar integration.}
Consider a component $X$ of the operation of level $i$ with
$n$ completed collar integrations, and put $j=i_\flat+n$.
Stage $n+1$ first imposes the localized regularity test
$\chi_i^{(n+1)}$, the collar gauge and the approximate
closeness test $\chi'_j$ at $2\delta'_j$. It completes the
fluctuation tests of step $j+1$. The removal uses the
retained averaging deltas and the density tests at the
top and the old core, together with the lower and higher
one-step tests left in the term. The sharper collar radius
$CL\delta'_j<\delta_j/2$ supplies the fluctuation implications
on the covered exterior bonds. The completed ordinary
tests supply the bounds on inherited strip bonds. For the
regularity implication the error on a shell of level $l\le j$
is bounded by
$C\delta'_je^{-c(M/M_1)(j-l)}$ in its curvature units,
which fits the gap $c\beta_{\rm s}2^{-(j-l)}\varepsilon_l$
in the ladder. For the missing crossing reads use the
comparison of \cref{hi:preceding-guard-support}. Remove
an old indicator only when it equals one on the complete
chart used for its removal. Any retained active indicator
is part of the real collar integral. The sharp collar split
at $\delta'_j$ follows the removal.

%
%
At the collar gauge insertion, the group is pinned at
the next coarse sites, so the coarse row remains fixed.
Use \cref{prop:averaging-1a-3-3} on the incoming law
before selecting the subsequent gauge-dependent outcomes.
The group is supported beyond the reach of the old
retained factors. For the prospective regularity test
use covariance of the critical orbit. This gives the
invariance needed to use the soft gauge factor against
the incoming law at this insertion. Its support is in
complete blocks inside the component's exterior buffer.
A link moved across the prepared inner face belongs to
the collar, while the exterior of the component is fixed.
The later tests are weights on the resulting variables.

In the row chart this can also be seen directly. At a
site $x$ with contour $V_j(y,x)$ from its coarse centre,
the draw of a new contour value $q_x$ uses
$u(x)=q_x^{-1}V_j(y,x)$. Its input contour lies in a
single block and reads no corridor pivot. At fixed
nonpivots and gauge draw, each selected pivot changes
by fixed group translations preserving Haar, and the
averaging row has its original target. Carry the
resulting contour values in the branch kernel for the
tests that read them. A prescribed contour value would
be a graph relation with the gauge draw.

When the component leaves at that split, the retained
determining list still has index $n$. Its tests include the
stage-$(n+1)$ regularity test and the approximate collar
test at the old level $j$, together with the failing sharp
test. In the local sharp convention this test is formed
at the localized next-stage background, so its failure
is a condition on the variables of this component. In the
global sharp convention use the exact transported test and
density of the failed branch. In the variant of the fibre that omits
the root's failing test the positive tests are still
required. The same list supplies the passing branch up to
integration. For comparison one may form the prospective
ordinary stage-$(n+1)$ datum with the same variables and the
retained averaging constraints. The two determining
sequences differ only at $j+1$, by the collar. The
prospective read introduces no additional stored variables.

For a datum plaquette use the smallest-level drop-out
containing its place. Take its full footprint, the union
of the endpoint blocks at its test level. An ordinary
density test is read at its original action cube if the
footprint leaves the region of the stage test, if its
level exceeds $j+1$, or if it belongs to the untouched
top action region. The old core also keeps its own tests.
Within the covered part use $\chi_i^{(n+1)}$, supplemented
by $\chi'_j$ for a level-$j$ plaquette touching the collar.
The density cube and the footprint are taken in the
prescribed partitions of the test. The current mask of
an old exponent does not set the place of this guard.

On footprints covered by the literal clauses,
the prospective guard gives the curvature bounds
used in the block averages. At a prepared or
original wall use the coarse holonomy bound
from the preceding promotion calculation.
The prospective datum has the layer bounds
\(C e_q\) on all these reads. The unchanged
older rows use the bounds retained in their
own prefix.
For the old datum on the collar, where the complete double
block is in this comparison, the retained deltas give the
pivots from the nonpivot values and the prospective coarse
row values. The pivot comparison at equal coarse values
costs a fixed factor depending on $L$. Thus the approximate
test changes the old level-$j$ holonomy bound by at most
$C_L\delta'_j$, by telescoping the plaquette product.
The changed lower rows leave their prescribed coarse
values fixed. In particular the corresponding level-\((j+1)\)
holonomies agree with their prospective values, also
when an old edge reads the complete one-step average.
For the no-inward determining sets the global and localized
next-stage averages differ on the collar by at most
\[
 C\mathsf F e_i
                  e^{-\delta_0 M\check R_i}\le\delta'_j/2.
\]
Apply this comparison on the complete collar
blocks inside the quiet component. The collar
lies in \(O_{j+1}\). If \(j+1<i\) this is
inside \(G_{i-1}\), whose clearance from
\(\partial X\) is \(10a_i\). At \(j+1=i\)
the clearance of \(O_i\) is \(2a_i\).
The endpoint blocks and the original small
test-cube halos together reach less than
\(a_i\) past this carrier by the cell-size
hierarchy. This gives the weighted margin
of order \(M\check R_i\) in the comparison.
On a passing sharp branch the two centre
comparisons fit \(2\delta'_j\). At a sharp stop
the retained local approximate test gives
\(2\delta'_j\) in local coordinates, hence at most
\(3\delta'_j\) in global coordinates. Use the
transported density for the global sharp stop.

%
%
For the energy estimates at the approximate and sharp collar
stops, consider a bond crossing the old \(\Omega_{j+1}\) wall.
Let \(y_-\) and \(y_+\)
be the centres of its exterior and interior blocks. The
interior block is in the first ordinary strip cell. Both
the passed and the failed ordinary cube masks start at least
one full cell inside the wall, and their enlarged reads reach
at most \(2LM_2\check R_{j+1}+1\) in level-\(j\) units.
The complete interior block, including its internal bonds,
therefore lies in the plain zone. Its internal bonds are
nonpivots and are kept by the original ordinary split. On
them the exact plain-zone map gives
\[
 V_j(e)=\exp\bigl(ig_j\mathcal A_j^{\mathrm s}(e)\bigr)
            S_{\mathrm{o,loc}}(e),
\]
where \(S_{\mathrm{o,loc}}\) is the original localized strip
centre. The full block lies in the joined interior, and its
fine sites other than \(y_+\) are free for the ordinary axial
gauge. Hence the mean contours of \(S_{\mathrm{o,loc}}\) from
\(y_+\) are the identity. If the kept Gaussian reserve gives
\(\max_e |g_j\mathcal A_j^{\mathrm s}(e)|\le\sigma\) on these
links, telescoping the contour products and using the derivative
bound of the mean gives
\(\max_{x\in B(y_+)}|V_j(y_+,x)-\mathbf1|\le C_L\sigma\).
At the exterior block the retained soft gauge weight gives the
same bound if its gauge exponent, divided by its inverse
coupling, is at most \(\sigma^2\).

On the complete double block, the curvature bounds below
give defects of size \(C\sigma\) for \(V_j\) and the
prospective local centre. Axialize \(V_j\) with the two mean contours,
fixing the two block centres. Its coarse row is still
\(V_{j+1}(\langle y_-,y_+\rangle)\), the target of the
original ordinary pivot constraint. The prospective centre is
axial on both blocks and has that same target. Double-block
rigidity compares each axialized crossing to the common coarse
value with error \(C_L\sigma\). Undoing the auxiliary gauge
costs another \(C_L\sigma\), by the two contour estimates.
Thus the principal angle on this crossing is bounded by
\(C'_L\sigma\) in the local chart. On a full history the
global angle is at most this bound plus \(C\sup_b|H(b)|\),
where \(H\) is the stage centre shift. Consequently either
failure indicator on the full history is bounded above by
the local indicator that some tested angle is at least
\(\delta'_j/2\). Use this upper bound in a local mass,
where the evaluation need only give the localized centre.
The labelled integral still uses the actual sharp test.
The crossing estimate uses the stored ordinary coordinates
on the interior block even when the guard itself reads
only the crossing.

%
%
%
Write \(q=j+1\), and take the one or two full order-\(q\)
blocks of a tested collar bond together with the backward carriers
for their averages. The bond has an endpoint in
\((\Omega_q^c\setminus Z''_q)\cap X\). The blocks and
their carriers reach only \(O(L^q)\) past this set.
The gap from \(Z''_q\) to \(Z''_j\) is at least \(ML^q\),
also in the compressed nest. At the first stage use the
gap from \(G_{i_\flat}\) to the first prepared face and
\(\Omega_{i_\flat}^c\cap X\subset G_{i_\flat}\). Thus all
these carriers have prospective layer at least \(j\).
The clearance of \(O_q\) from \(\partial X\) is at least
\(2a_i\). After the \(O(L^q)\) enlargement this leaves
the carriers in the unchanged interior of the join.
In particular the old \(\Omega_i\) wall at the last
stage is separated from the outer tower collar.

Every tested datum plaquette whose backward carrier
touches the floor is below level \(i_\flat\), and its
blocks are in the protected core. To reach it from a
sharp carrier one must cross the band from \(G_{i_\flat}\)
to the first prepared face. That band has thickness
of order \(a_{i_\flat+1}\) in layers at most \(i_\flat\),
which gives weighted distance at least \(cM\check R_i\)
even at a later compressed stage. The artificial tower
collar at \(\partial X\) is also this far in the weighted
metric. The datum tests inside this weighted neighborhood
use the exact old or promoted averages on their complete
carriers. Their bare Wilson shares are in the charged
enlargement and outside the floor, so small energy gives
defect \(C\sigma\) there. Use the real admission bound
\(C\check R_i^2\varepsilon_i\) on the other tests.
The weighted regularity estimate now gives bare defect
\[
 C\bigl(\sigma+\check R_i^2\varepsilon_i e^{-cM\check R_i}\bigr)
 L^{-2j}
\]
on the sharp carriers. The averaging estimate gives
the same bound without \(L^{-2j}\) for the prospective
level-\(j\) field. The exponential term and the stage
centre shift are \(o(\sigma)\) for
\(\sigma=c_*\delta'_j\check R_i^{-4}\), uniformly in
the waiting window.

%
%
For the actual \(V_j\) at an old wall crossing, the
interior free links belong to the old ordinary fibre.
Use the original localized strip centre on the
component of \(Z_q\cap\Lambda_j\) in \(X\), with the
one-part join on \(Z_1\) at the first step. The strip
admission uses the original density bounds and exact
averaging rows on this domain. Reproduce those rows
and the density tests from the retained birth data in
the compatible evaluation. This admits the joined
centre at the fixed ordinary radius without requiring
an extension of the evaluation to the torus.

A tube of width \(c a_q\) about \(\partial O_q\)
lies deep in this join, stays outside \(G_j\), and
misses the next original wall. Its layers are \(j\)
and \(q\). The level-\(j\) datum reads on its exterior
are old determining values, and the level-\(q\) reads
are the old coarse targets and the exact averages of
the exterior rows. Their full backward carriers lie in
the charged enlargement and avoid the floor. The
datum away from the tube has the ordinary admission
bound, and leaving the tube costs weighted distance
\(cM\check R_q\). The weighted estimate therefore
gives level-\(j\) plaquette defects of size \(C\sigma\)
for the localized strip centre on the full double
block.

On the interior and crossing nonpivots the exact
plain-zone formula and the small retained Gaussian
coordinates bound the change from this centre by
\(C\sigma\). The exterior internal links are
determining and are identical for the two fields.
The centre and the actual field have the same coarse
target there by the joined row identity. Use the
straight-corridor estimate for the actual middle bond
with the centre curvature and the nonpivot changes both
of order \(\sigma\). This puts the actual solution in
the regular branch, where the one-row inverse gives the
bound \(C_L\sigma\) for its pivot. If the stopped
branch has translated a retained slot, use its original
value from the recorded inverse map for the Gaussian
reserve.
Telescoping each plaquette gives the required bound
for the actual field, including plaquettes using the
pivot.

If the sharp bond does not cross the old wall, its
endpoint blocks are in the old exterior and their fine
links are preceding determining rows. The compatible
field supplies the full \(V_j\) plaquettes directly, so
the Wilson share gives their bound. At the inner prepared
face the inner block remains in the determining layer.
The two fields have identical internal links and hence
identical mean contours there. When the two auxiliary
axial gauges are undone, this common contour factor
cancels in the norm of the relative link. Only the
contour errors on the collar side enter. On every collar
block the soft weight is either inherited from the old
exterior belt or inserted on the remaining blocks of
this stage. This covers all its contours. For a bond
internal to a single block use the single-block clause
of rigidity. The coarse target for a crossing of two
blocks is the common retained target.

At the approximate split, use the test in the original
coordinates, against \(V^{(j)}_X\). At this moment the
collar gauge has been inserted and the old ordinary
strip still carries the coordinates in the preceding
calculation. The complete coarse average is its original
constraint. The completion of the ordinary fluctuation
tests takes place later. The bound on the two block
contours here follows from the small paying energy in
the old strip and the collar gauge. Together with the
curvature and rigidity estimates it gives
\[
 |V_j(b)(V^{(j)}_X(b))^{-1}-\mathbf 1|\le C_L\sigma
\]
on the original approximate mask. With
\(\sigma=c_*\delta'_j\check R_i^{-4}\) this lies
below \(2\delta'_j\). On a sharp stop use the original
values recovered from the recorded translation, and
apply the angle bound in the chart of its sharp test.

The paying Gaussian and gauge variables in these
two estimates are of level \(j\ge i_\flat\). Their weights
and the Wilson coefficients on the charged enlargement
are comparable to the corresponding ones at level \(i\).
If the sum of the fresh Wilson, gauge and kept Gaussian
shares is at most \(c\theta_i\sigma^2\), the individual
bounds used above follow after reducing \(c\). Hence a
failure at either collar test forces the opposite inequality. In a
containing mass these shares use the child's full
carriers in its enlargement outside its own floor and
its original level-\(j\) slots. For a containing operation
whose waiting window starts at \(h\), an older standing
mark of level \(m<h\) is in the retained \(Z_{h-1}\)
and in the innermost box \(G_h\). Its witness extends
only two unit layers of level \(m\) beyond the mark.
The hull and eight-layer enlargement at step \(h\)
still contain the witness after removing the cells in
contact with \(\Omega_h\). Thus the child's bare
plaquettes lie in the containing operation's floor.
The containing operation spends only plaquettes with
no corner in that floor. Its gauge and Gaussian shares
start at \(h\), above the corresponding slots used by
the child. This separates payments through nested
waiting windows.

%
%
The newly inserted ordinary completion uses the
stored fluctuation coordinate with threshold
\(\delta_j\). The ordinary change of variables has
already been made, and the symbol \(\mathcal A_j\)
in the completion denotes its resulting coordinate.
The masks extend from the first strip cell by the
two small test-cube layers and fit in the retained
two-cell strip. A failed completion therefore has
a stored bond with
\(|g_j\mathcal A_j^{\mathrm s}|\ge\delta_j\).
The kept Gaussian reserve pays for this event on
the local fibre. Tests inserted before the ordinary
change of variables keep their original pullbacks,
as in the ordinary birth payment.

These bounds give regularity on the two stage sequences.
On directly guarded old collar reads the
layer factor at level \(j\) is
\(O(1)L_\flat^{2(j-k_0(i))_+}\). On the full datum
use the bound \(C e_l\) at level \(l\).
The real profile estimate then bounds the scaled
curvature and current on the retained carriers.
The adjacent-layer conversions and the compression
loss are those of the real endpoint comparison,
so the bounds fit the half radii of the earlier
structures. Local gauges follow from the curvature
bounds on their full cubes.
The restricted real background is critical over its own
averages on subordinate towers, so the tower and core
conditions use the same bounds. In particular an inner
island leaf is read in its original real core.

The stopped operation keeps the old one-step integrals of
levels $l=j+1,\ldots,i$, the bare nest averaging integrals
of levels $l=i_\flat+1,\ldots,j$, and the old core. The
factors present at departure multiply this product, including
the current test and any collar gauge already inserted.
At a sharp-test stop in the global convention the first
kept one-step integral uses the transported collar
coordinates, with the corresponding pivot and Jacobian
factors. A sharp failure in the component-local convention
keeps the untransformed prefix with its local test. The
earlier stops use the original pivot representation.
Only completed integrations contribute stage Gaussian
normalizers and terms to the stage boundary sum. At each
domain $D$ that sum is filed once, with stages up to the
maximum completed index of the components meeting $D$ and
with already integrated stored fields set to zero. The
stopped action keeps the Wilson weight from before the
stages. The Wilson differences from completed stages cancel
when one returns to that weight and are excluded from this
boundary sum.

%
%
For its norm separate each completed
stage term into its direct normalizer
and its localized cluster term. At a
fixed carrier, the cluster bounds cost
at most \(2\mathsf N_i\epsilon_{\rm comp}\)
at the final rate. The conversion to
pinned norms and the closing of carriers
spend \(O(1/\mathsf N_i)\) in the rate
at each stage. The count uses the
separation of the stage components.
The total loss is therefore bounded
independently of \(\mathsf N_i\).
The direct part must use the sum of its
pinned normalizer norms, not a
single-stage supremum multiplied by
\(\mathsf N_i\). In particular the
coefficient for the retained stage
entry is bounded by
\(\bar D_1^\Sigma+2\mathsf N_i\epsilon_{\rm comp}\),
where \(\bar D_1^\Sigma\) is this summed
direct norm. This gives the stage
amplitude in the standing-list bound
once the direct sum is bounded on the
actual graded cells. The identification
of each Gaussian determinant with its
own stage integral supplies the direct
terms for that sum.

%
%
For a sourced excess taken at the
same evaluated pair, the direct
terms cancel before the pinned
estimate is used. Keep the stage
prefix, masks, stored data and
reference Gaussian the same in
the two systems. The direct term
is then the same function of
the pair, including its subtraction
at the trivial pair. All system
entries occur in the potentials and
the resulting closed clusters.
Thus the difference of the retained
stage sums uses only the cluster
majorants on that common chart,
including for graded cells. The
absolute real bound on the standing
term still uses the direct sum.

% Earlier Part I, Lemma 4.25 on Gaussian-piece oscillation
% and the replacement of the field clause in Section 4.94(D).
A later connected coefficient of a
stored Gaussian piece also admits
a comparison at varying pairs.
Suppose the piece has been formed
with its local holomorphic majorant
\(\nu_e\) on the larger fixed
pair domain. Along the interpolation
\(t\mapsto p(t)\), let \(\delta\)
be the size of the pair displacement
in the local coefficient norms,
divided by their admitted radii.
Restrict the piece to the complex
line of the displacement. On a
disc of radius \(c/\delta\) its
supremum is bounded by \(\nu_e\),
so its change up to \(t=1\) is
bounded by \(C\nu_e\delta\).
For a stored quadratic form
\(Q_e(y,y)\) read at \(y=\operatorname{Ad}_g x\),
a change from \(y_1\) to \(y_2\)
gives the two terms
\(Q_e(y_2-y_1,y_2)+Q_e(y_1,y_2-y_1)\).
The bilinear majorant and the
bounded adjoint actions give the
corresponding linear bound in
the size of the relative rotation.
When a connected coefficient reads
the piece through a mixed derivative
in a nonempty set of its switches,
that derivative kills the value at
the initial pair. Cauchy's estimate
then applies to the increment.
This argument applies to an
arrested Gaussian piece, or a
survivor piece meeting its
arrested collar, when its
formation gives the stated
local majorant. The undifferentiated
term in the standing sum still
requires the absolute bound.

For each completed stage the number of integration
bonds is bounded by $C(L)$ times the number of sites of its
collar and strip at that stage's scale. Distinct drop-outs
sharing a re-made level are disjoint, so this count charges
the completed stage numbers once at that level. Changing the
departure place within the preparatory stage at fixed $n$
keeps the regions, the integration list and the action carrier.
The particular tests and any inserted gauge are recorded in
the local history. There are at most $4\mathsf N_i+6$ possible
places for an operation at level $i$, including the second
class, while the marked label uses at most $\mathsf N_i+2$
indices. This is the stage-place part of the count in
\eqref{hi:ancillary-code}.

%
\paragraph{A stop at extendibility.}
At extendibility all preparatory integrations
have been completed. Use the prepared datum
with the holonomy bounds carried from those
integrations and the original exterior and
core tests. The real profile estimate therefore
applies on the prepared read. The negative
extendibility test depends on its exterior
argument. At this stop the central integration
is still pending. The old top values and the
bare averaging chain in the core keep their
original places. A later read replays the
prepared support from those values and
constraints. The frame operation and the
new box Haar field enter only at the subsequent
stops.

%
For its failure in a containing mass, use the
old exterior top values on the complete
carriers in the child's enlarged component.
These carriers are outside the central box
and its floor. If their Wilson share is \(E\),
the averaging energy bound controls the
exterior top plaquette defects by
\(C\sqrt{E/\theta_i}\). The axial extension
on the boundary of the central box gives a
top field with the same exterior and
defects at most \(C(M)\sqrt{E/\theta_i}\).
In the top inset the localized critical
configuration for this extension has bare
curvature bounded by this number times
\(CL^{-2i}\). The required region
\(\Omega_i^c\cap X\) lies in that inset by
the two-cell clearance. Thus the recorded
negative extendibility test forces
\(E\ge c(M)\theta_i\varepsilon_i^2\).
This witness uses the original exterior
argument of the test.

%
%
%
%
%
\paragraph{The tests retained at a late stop.}
A shell or last-shell failure occurs after the final
preparatory guard and the central integration. It
retains the positive tests preceding the failure. Use
the prepared datum for the regularity comparison and
the fresh outgoing field for the later averaging rows
that read it. The removal that uses the two passing
shell tests takes place only on the continuing branch.
In the old prepared estimate the final stage test is
used with the original top density tests and the old
core tests.

For the inset operation, a level-$q$ tested plaquette
of the prepared sequence in $X$, with $q<i$, has its
place in $Z''_i\subset O_i$ by nesting. Its footprint
is the union of the four order-$q$ blocks. From its
place this extends by at most $3L^q/2$, so the
clearance in \eqref{eq:elab-history-quiet-boxes} puts
the full footprint inside the spatial inset. On the
no-inward list every corner is outside the prepared
$\Omega_i$: the first-kind ends are outside
$\Omega_{q+1}$ and the second-kind ends are outside
$\Omega_q$. Thus all four corners are in $Z''_i$.
The final box is a union of aligned level-$i$ blocks,
and each order-$q$ block lies in its containing
order-$i$ block. The footprint is therefore in
$Z''_i$ in this geometry. The top tests in the original
top action region keep their halos, including any
reads outside $X$.

A lower footprint can still meet a seam of an earlier
shell. The literal shell bound is on bare plaquettes
wholly contained in that shell. The comparison test at
the last prepared seam supplies a different set of
plaquettes. For real half-radius admission on an
inverse branch, use the
\hyperlink{hi:prepared-real-endpoint}{real endpoint comparison}.
Before marginalizing the central integral, keep its
prepared variables. In the framed integral these are
\(V''=V'V_0\) on the changed coordinates, with \(V_0\)
formed from the retained exterior. This substitution
recovers the same datum in each frame summand. On its
admitted incoming fibre the old lower holonomy bounds
and the exact top averaging relation give the prepared
datum bounds \(C e_q\) of the endpoint comparison.
Carry these bounds as a support condition on the
prepared datum. For each actual exterior \(e\), the
indicator of the bounds depends only on \(e\) and the
prepared variables. It equals one before the old
averaging deltas are integrated, on the admitted
entering fibre. The pushforward therefore assigns full
mass to the set where this indicator is one. Insert the
indicator in the retained kernel to use this density
version at each exterior. The real layer
estimate and the subordinate tower criterion then
give the half-radius membership on its earlier read
sets.

On a later read, use the prepared first-kind links
directly, including the old top determining values
on bonds meeting the box. A second-kind edge is
completed from its retained lower double block.
This replays each prepared holonomy on the complete
read set. The promotion relation on the original
fine rows was used to prove its bound before
integration. Its consequence on the prepared datum
is the support condition just described. The fresh
box field continues to enter the outgoing averaging
rows. The original exterior density tests near
\(\Omega_i\) retain their own halo. In the quiet
window the central box lies in \(\Omega_{k_0+1}^c\),
with \(k_0+1\le i-1\), inside the predecessor's hull.
The eight-cell growth leaves at least four cells
between it and \(\Omega_i^{\sim4}\). A complete
density-test read uses less than one more cell.
These tests therefore read the retained top exterior
even after the central draw. The shell
test itself reads the retained framed fluctuation,
and the last-shell test reads the prescribed
configuration with the retained bottom shell field
on the core side of the cut and the identity on the
shell side. It keeps the original values on bonds
crossing the inner cut, according to the convention
in the last-shell comparison.

The first outgoing chart recovers this exterior before
it uses the free box pivots. Keep the old tests and
support indicator in that joint measure. If a later
row change makes an exterior argument a function of
its outputs and remaining coordinates, compose each
test with that function on the actual inverse branch.
An earlier exterior row is read from a fixed
coordinate, a retained variable, or an output of
its original graph. Take the closure of the test
reads under these graphs before taking a marginal.
This preserves the support indicator even at a
generated exterior. We must check the joint row
bound at the resulting read set. A later estimate
that uses a sharper threshold at an earlier shell
seam must supply that bound on the full footprint.

For the analytic read of the kept one-step factors use
$X$ enlarged by one level-$i$ cell layer. The walk carrier
of a kept strip from step $l<i$ lies inside $X$. A strip
still carried from step $i$ can have a carrier extending
to that full one-cell enlargement. Once all preparatory
stages are complete the top strip no longer remains and
the one-$M$-cube collar suffices. In either case the pivot
factor uses the two endpoint blocks in its original birth
partition. Their regularity follows from the corresponding
pointwise bounds and the adjacent-level radius comparison.
The walk and pivot bounds are then taken on their original
masks. The larger analytic read does not extend the paying
Wilson plaquettes or the integration coordinates of the
component. The fibre on $X$ receives only the resulting
values of the forms and Jacobians, with their uniform
one-sided bounds.

%
%
\paragraph{The same halo at later levels.}
Include the old top density reads and their first outgoing
double blocks in the read set \(B_{\rm read}\) of a mark of
level \(i\).
For a still-carried top strip include its one-cell analytic
read as well. The old cells are nested in the new grid and
are no wider than the new cells. Thus \(B_{\rm read}\) lies within
two layers of the level-\((i+1)\) hull of \(X\). This is
inside the eight-layer growth and away from the new gauge
draw. It lies in the ordinary successor of the mark. At
subsequent ordinary steps, a full double block of a row
reading the preceding component lies in its hull enlarged
by one cell. The same growth includes these rows in the
raw region of the successor. Their outputs use the carried
lower graph and the current ordinary inputs.

For a later containing preparation let \(h\) be its bottom
level. The window rule gives \(i<h\). If \(h=i+1\), the
initial two-layer bound puts \(B_{\rm read}\) strictly inside the
old core of the new operation, even after removing the
cells touching \(\Omega_h\). If \(h>i+1\), apply the
eight-layer inclusion at \(h\) to the successor at
\(h-1\), which still contains \(B_{\rm read}\). In both cases
the original reads are in the floor. The preparation
uses the old outputs in its bottom core and modifies
only higher layers. The carried chart and its inverse
factors therefore stay at their original levels. A
healing that absorbs the component instead supplies
the new determining datum on the healed part. The
old graphs there remain in the integrated weight of
the healing.

%
%
%
%
\paragraph{The next ordinary strip.}
The larger coarse bounds at the original walls
are attached to footprints within the carried
components. At the next ordinary step the
complement of \(\Omega_t\) contains the hull of
\(Z_{t-1}\) and its eight-layer enlargement.
The old wall tubes are inside that hull. The
ordinary test and strip carriers use the exterior
halo of this growth rule. On that halo the
ordinary positive tests retain their original
thresholds and use their complete averaging
footprints. For the one- or two-cell neighborhood
of \(\Omega_t\), the complete small test and
contour carriers cost less than one more cell.
The no-inward convention also puts the places
of all tests of level below \(t-1\) outside
\(\Omega_{t-1}\), hence inside \(Z_{t-1}\).
The near part of the datum profile uses levels
\(t-1\) and \(t\), with comparable thresholds.

In the real profile near the ordinary strip,
the part of the datum carrying the enlarged
wall bound is at weighted distance at least
\(c M\check R_t\). Indeed the path across the
guarding cells uses layers of spacing at most
\(L^t\), and the fixed read halos leave a cell
margin. By \eqref{mb:weighted-maximum} its
contribution to the normalized curvature and
current is at most
\(C\varepsilon_\star e^{-c\delta M\check R_t}\).
This is smaller than \(\varepsilon_{t-1}\)
at weak coupling. The ordinary tests in the
remaining part of the profile give the usual
\(C\varepsilon_{t-1}\) bound near the strip.
For the localized strip at \(t\ge2\), the join is
formed on \(Z_t\cap\Lambda_{t-1}\). This domain
is disjoint from \(Z_{t-1}\), which contains the
carried wall tubes. Use the two-field construction
of the joined datum there. Its top values satisfy
the retained \(\mathcal M(V_{t-1})=V_t\) constraints
on the complete blocks, and its lower values use
the original positive density cubes meeting the
join. Their bounds give the local datum radius
of order \(\varepsilon_{t-1}\) on the full
joined test carriers. At \(t=1\) the strip has
the one-part join and there are no earlier marks.

One can formulate the transfer without identifying two
unnecessarily large sets of physically realizable arguments.
Let $\mathcal W_k(i,X,o)$ be the admitted reads at level $k$,
obtained by imposing the creation rule and all retained prefix
constraints up to that level.  Then
\begin{equation}
 \mathcal W_k(i,X,o)\subset\mathcal W_{k'}(i,X,o)
                \subset\mathcal W_i(i,X,o),\qquad i\le k'\le k.
 \label{hi:nested-read-sets}
\end{equation}
Actual read images are contained in these sets by \eqref{hi:read-coverage}.
If one defines $\mathcal W_k$ instead as those images themselves,
the same nesting additionally needs the exact prefix projection
to commute with the local read.  Label truncation alone is not
a proof of that stronger statement.  The admitted-set version
already gives the required mass transfer.

Restriction of a supremum and reduction of either its volume
or its coefficient can only reduce its upper integral.  Hence
for $0\le c'\le c$ and $0\le v'\le v^\natural$,
\begin{equation}
 \mathsf M_{c'}^\sharp(\mathcal W_k;v')
       \le\mathsf M_c(i,X;v^\natural)
       \le e^{-a_\star P_i}.
 \label{hi:read-mass-transfer}
\end{equation}
There is no new factor for the number of later levels.  If an
inductive class permits a slack constant, for example
$e\,e^{-a_\star P_i}$, that same slack is preserved for its
old marks; it is not multiplied once per read.  Newly formed
marks have the sharper bound proved in \eqref{hi:weighted-mass-conclusion}.  A run started
with those formed bounds retains them at all later levels as
long as its weights, domains and presentations remain in their
admitted tubes.

If a mark is absorbed into a new outer component, its local
record occurs inside the new local history.  Its weighted mass
is then used in \eqref{hi:inner-history-integration}, not loaded again as a second independent
outer mark.  This is the operational version of the multiplicity
identity in \eqref{hi:bounded-multiplicity}.

\paragraph{Field readings and compensators.}
Gauge covariance acts on stored variables as well as external
triples.  Denote that action by $\rho_\omega(u)$.  It conjugates
the variables rooted in a tree by the transformation at its
root and acts on an outside-starting stored bond at its actual
starting site.  Let $R_\omega$ be this finite list of sites.
A permitted compensator $c[U]$ is Borel, obeys the actual chart
conditions, and is the identity on $R_\omega$.  Then
$\rho_\omega(c[U])$ is the identity on every stored variable.
The compensated tied expression
\begin{equation}
 \Phi^\diamond_\omega(\xi;\mathfrak z)
   =\Phi^\circ_\omega(\xi)
      -\tfrac12q_X(\xi;c[U]\star\mathfrak z)
                           +V_X^\diamond(\xi;\mathfrak z)
 \label{hi:compensated-reading}
\end{equation}
therefore needs no change of integration variable depending on
$\xi$.  Its formation identity must say that the whole tied
sum takes the same value for equivalent triples.  Separate
invariance of the two summands is not required.

On every admitted fixed $U$ chart the dependence on the complex
potential and current is holomorphic, and the same envelope
bounds the compensated read.  Borel dependence on $U$ follows
from its Borel compensator and background maps.  These claims
require the actual compensator with these fixed sites and core
admission.  Gauge covariance for a fixed global transformation
would not alone construct it.  This is why the read rule retains
the sites on which the stored-variable action depends.

%
On the chart of a healed label one can construct this
compensator on the ordinary towers and the rooted trees.
Let \(a(x)\) be the top of the tower of \(x\), or the root
of its box if that tower starts in a treated box.
For each \(x\), take \(c[\mathbf U](x)\) to be the product
of the averaged transports from \(a(x)\) to \(x\).
The path first follows the box tree to the top, when
present, and then the combs in the successive layers
of the tower. The product is empty at a top outside
the boxes and at every root. At the real points of the
chart the factors are measurable group variables, and
on the admitted complex chart they are holomorphic.
Endpoint covariance cancels the transformations at
the intervening sites, giving
\[
             c[\mathbf U^u](x)u(x)
                     =u(a(x))c[\mathbf U](x).
\]
The compensator gauges each comb and tree transport to
one. For a tree variable the action on the stored
coordinate is at the root. For an outside-starting
variable whose start is a tower top the action is at
that top. Both kinds of coordinates are therefore fixed
in this reading. A term that reads another stored site
requires that site to be checked separately in \(R_\omega\).

Apply \(c[U]\) to the real arguments, and \(c[\mathbf U]\)
to the pair \((\mathbf U,\mathbf J)\), with the current
conjugated. At fixed \(U\) the real change is fixed. The
compensated pair depends only on the original pair, so a
tied slice expression is still tied on its admitted reads
at fixed stored variables. Under a real gauge change the
two normalized presentations differ by \(u\circ a\).
The invariance of the slice expression under these
root and top rotations gives the joint invariance of
the read expression. For a local term the carrier must
contain the full paths just used. In particular a carrier
that reads a box tree must contain that tree as well as
the affected towers.

The analytic bound applies on the preimage of the
original chart. This restriction matters even for flat
complex links. Take a constant imaginary potential in
one direction on a local order-\(j\) tower, with bare
spacing \(L^{-j}\) in level-\(j\) units and with \(U=1\).
The real compensator is then one. The complex comb
puts the transport across two adjacent tops on the bare
bond joining their towers. A logarithmic potential of
size \(\alpha\) before this change has size \(L^j\alpha\)
on that bond afterward. Hence the small potential bound
in the original chart does not give the same bound on
every compensated pair. At a frozen seam we must also
check the endpoint constraints on the unchanged lower
rows. The rooted product on the healed label supplies
the gauge correction only on the admitted chart just
described.


\subsubsection{Causal restrictions and the product of healed histories}

We give a finite algebraic statement that explains when a sum of
histories is a gas. Spatial separation by itself is not its
hypothesis. The determining reads, action reads and reference
normalizers must all be included.

An elementary move has a discrete type, a finite list of input
slots, a finite list of output slots and an integral kernel for
the variables it removes. A test is a move with a finite outcome
list and no integral. Every slot has a level, region and type,
such as a group field, free Gaussian coordinate, scalar constant
or background chart. Slots are versioned. The dependency relation
directs an edge from a move producing a slot version to each later
move reading it. It also orders every write/write, remove/read and
priority conflict consistently with a legal execution. We require
every pair left incomparable to be independent in the sense below;
producer-to-reader edges alone do not imply this property.
The relation includes the discrete reads deciding class membership
and the background box in which a critical configuration is formed.

Let $R(\gamma)$ be the closed read support of a group of moves
$\gamma$: include their written variables, every input they read,
every relevant action plaquette and every wall and gauge factor
they use. If a move evaluates a background on a whole box, that
box is included. A carrier cannot be reduced to a short
random-walk segment merely because a displayed kernel term has
a short segment. In particular $q_X$ in \eqref{hi:tied-exponent} has carrier
$X^\sim$; its separate walk terms use a minimizer on all of $X$
and are not local functions of their short walk supports.

The exact locality hypotheses for the finite statement are as
follows. Local functions and reference kernels obey hereditary
restriction: forming a move on a larger admitted box and setting
outside decoupling parameters to zero gives the move formed on
the smaller box. Its normalizers and counterterm readout obey
the same identity. Two independent groups have disjoint closed
read supports and no common decoupling wall. Their incoming
exterior slots agree with the reference carried history, and
all parameters stay in the common admitted domains.

\begin{lemma}[{Order and restriction}]\label{lem:history-order-restriction}
Under these hypotheses incomparable independent moves commute,
and a hereditary subcollection has the same outputs on its slots
whether executed alone or inside the larger history.
\end{lemma}

\emph{Proof.} For two integral moves, the inputs held fixed by
one contain no variable written or removed by the other. Their
integrands multiply, their reference measures are product
measures and their output substitutions concern disjoint slots.
Tonelli proves commutation for positive real kernels. For
complex kernels the common Borel envelope in \eqref{hi:actual-pointwise-envelope} supplies the
absolute integrability for Fubini. A test has the same value
before and after an independent move, since every slot it
reads has been included in its closed support. This also
proves commutation for a test and an integral move.

Two linear orders of a finite dependency partial order differ
by adjacent swaps of incomparable entries. To see this, take
the first entry of the second order and move it forward in the
first. Every entry passed is incomparable: an ancestor must
precede it in both orders, and a descendant cannot precede it
in the first. Remove this entry and continue by induction.
The commutation just proved makes the operation independent
of the chosen linear order.

For restriction, take the first retained move. Every input it
needs is either fixed exterior data or produced by a retained
move, by heredity. It therefore has the same value as in the
larger history. Determinism of its test and formation formula
gives the same outcome and output. Induct on the remaining
moves. Hereditary consistency identifies their formation on
the restricted boxes with their restrictions from the larger
boxes. This proves the statement for functions and integral
operations.\hfill$\square$

A minimizer on a larger domain generally changes with exterior
data. Locality of the bare action does not make that minimizer
a function of an arbitrarily smaller carrier. The actual
localized background functions, interpolation endpoints and
hereditary identities must be used in this lemma. Naming a
carrier is not such an identity.

Fix the outer history $\bar H$ and the real data that the completed
operation leaves unintegrated. In a history over $\bar H$ the healed
components of the current level $t$ form a compatible family
$\mathfrak X$. Restricting to a subfamily deletes the entire ancestral
cluster of each omitted component. Conversely, disjoint admissible
component histories that do not touch the standing region can be
united with $\bar H$. This follows from the history restriction and
union rules in
\cref{sec:large-field-the-tests-the-nested-regions-and-the-content-of-a}.
The output label and marks are common to the resulting histories.
The sums over the histories and frames of $C\in\mathfrak X$ are part of
the same local integral operation $\mathbf T'_t(C)$ in every family
that contains $C$.

These are statements about the records. In the integral, each record
uses its actual test support. In particular, on the first-class support
$\chi_t(C^{\sim-6})=1$ the two reference-law cutoffs in $C$ are one.
For the cutoff of the reference fluctuation this uses the bound by a
fixed multiple of $M^6\varepsilon_t$ and the choice of its radius
$M_0\varepsilon_t$ above that bound. The exterior
small-field test incorporates $\chi_t(C^{\sim-6})$, and the reverse
gauge-fixing identity removes the reference tree factor relative to
the variables of $C$. The second-class support retains the complementary
test, both cutoffs and the tree factor. The operations for the standing
components thus act on the following sum, with their variables kept
fixed while summing over the healed components.

Write $\mathcal Y(\mathfrak X)$ for the domains, other than single
collars, that contain at least one collar of a healed or second-class
component and contain every collar that they meet in a cube. Here the
collar is the whole enlargement $C^\sim$. Write $V_{\mathfrak X}(Y)$
for the interaction on such a domain. The single-collar terms are
included in the component operations. The completed action is common
to the families over $\bar H$. To compare the interactions, take the
local lists before assigning terms to slots. Localisation is linear
within each of the pointed, unpointed and Wilson classes. The Wilson
pieces with profile-dependent weights first use one common cover, on
which the level-$t$ filing profile $\varphi_t$ cancels by the Wilson
identity. The terms made by $\mathbf T$ and the stages are compared
before filing.

%
%
For a complete localized carrier missing the enlarged collars of the deleted
first-class components, use the comparison obtained by
erasing their histories. It preserves the local comparison
before filing. The input members on the domain must
correspond, the strip shifts must be determined by the
zone and label, and the standard rules for the levels
must apply. Under these conditions the local operators,
fluctuation potentials and stage terms agree as partial
maps. In particular the stage region is read on the
enlarged domain and its preparatory entries and nests
are read in their declared level windows. The Wilson
pieces using the final filing profile were combined above.
The comparison keeps the common functions even when
the final mask changes their slots.

The same argument applies to a deleted second-class
component when the formal history space permits
reassigning its class for this comparison. The ordinary
and stage functions are formed before the final class
test. Give the component the first-class designation
in the comparison of those unfiled functions, apply the
erasure identity, then file under the actual output mask.
This uses the order of formation and the formal
admissibility of the reassigned history. The real
histories in the integral retain their actual tests.

More generally, suppose $\mathfrak X\subset\mathfrak X'$ and $Y$ is
admitted for both families and meets none of the added collars. Every
entry with base in $Y$ then has the same carrier, domain and value:
the earlier entries agree by coherence, the response evaluations use
the same components and pinned sites in $Y$, and the kept factors
have carriers inside the collars of the components that remain.
Thus $V_{\mathfrak X'}(Y)=V_{\mathfrak X}(Y)$ on its common domain.
The requirement that $Y$ be admitted for both families is used here.
Adding a component can change the admissible resummation domains.

For $Y\in\mathcal Y(\emptyset)$ let $Y^{\mathfrak X}$ be $Y$ together
with each healed collar sharing a cube with $Y$. The collars are
disjoint, so this is a domain in $\mathcal Y(\mathfrak X)$. Set
\[
 W_{\mathfrak X}(Y')=V_{\mathfrak X}(Y')
   -\sum_{Y\in\mathcal Y(\emptyset):\,Y^{\mathfrak X}=Y'}V_\emptyset(Y).
\]
If $Y'$ contains no healed collar, the only preimage is $Y'$ and
the difference vanishes by the preceding comparison. Each other
difference depends on the family only through the healed components
in $Y'$ and their local histories. The subtracted terms are
independent of the variables integrated by the healed operations.
Grouping by $Y^{\mathfrak X}$ and taking their common exponential
outside these operations gives
\begin{equation}
\begin{aligned}
 &\sum_{\mathfrak X}\prod_{C\in\mathfrak X}\mathbf T'_t(C)
     \exp\!\left(\sum_{Y\in\mathcal Y(\mathfrak X)}V_{\mathfrak X}(Y)\right)
 \\
 &\quad=\exp\!\left(\sum_{Y\in\mathcal Y(\emptyset)}V_\emptyset(Y)\right)
  \left[1+\sum_{\mathfrak X\ne\emptyset}
    \prod_{C\in\mathfrak X}\mathbf T'_t(C)
    \exp\!\left(\sum_{Y'\in\mathcal Y(\mathfrak X)}W_{\mathfrak X}(Y')\right)
  \right].
\end{aligned}
 \label{hi:history-gas-identity}
\end{equation}
The component operations act on the exponentials following them,
including the dependence on their histories and integration variables.
This identity requires no division by a history density.

Expand the exponential inside the square bracket as a product of
$1+(e^{W_{\mathfrak X}(Y')}-1)$. Make a graph with one vertex for each
healed collar and each selected interaction domain. Two vertices
are joined when their supports have coinciding or adjacent cubes
outside the hulls of the fixed outer label. Each healed collar lies
outside those hulls. The restriction and union rule for the histories
and the agreement of the differences on retained domains give the
product of the activities for the different connected components
of this graph. The standing variables remain parameters during this
factorisation. Each activity contains a healed component, since
every selected nonzero difference has a healed collar in its
domain. The local terms $V_\emptyset$ remain in the boundary exponent,
since their domains contain a second-class collar.

The identity is for the complete operation on its physical data.
Local activity formulas can also be estimated on their admitted
complex cores, including independent currents. Their estimates do
not extend the global physical identity to independently varied
currents. To take the logarithm of the square bracket and file its
terms by their unions, use the polymer criterion after the
integrable kernel majorants and the activity bounds have been
established.

\paragraph{A kernel norm for composition.}
For a local kernel $K(\xi;p)$ define
\[
 \|K\|_c=\int^*e^{cv(\xi)}\sup_p|K(\xi;p)|\,d\mu
\]
including its discrete sums and frame factors. It gives
\[
 \|\mathbf T_KF\|_\infty\le\|K\|_0\|F\|_\infty
\]
for bounded insertions on the full read space. If a
second kernel is evaluated at an admitted read
$q=q(\xi,p)$ and has stored variables $\eta$, then
the additive weight and a common Borel envelope give
\begin{equation}
 \int^*e^{c(v_1+v_2)}
       \sup_p|K_1(\xi;p)K_2(\eta;q(\xi,p))|
                         \,d\mu(\xi)d\nu(\eta)
                    \le\|K_1\|_c\|K_2\|_c.
 \label{hi:kernel-composition}
\end{equation}
Bound the factors by their strong majorants and use
Tonelli. The read must lie in the core of the second
norm at every point of the first support. A real
read needs only the strong real version of that
norm. This is the kernel mechanism of \eqref{hi:inner-history-integration}, and
shows why a fixed-read integral estimate is not
the right norm for general composition.

\subsubsection{The loaded label sum and its analytic price}

At final level $k$ let $N_k$ be the number of lattice sites.
In the site load put $D_m=\Omega_m\setminus\Omega_{m+1}$ for $m\ge1$.
At level zero use $D_0=T_\ell\setminus\Omega_1$, counted on the bare lattice.
This counts all bare sites of that layer, including the sites outside the
collared first member. The determining sets of \eqref{hi:compatible-fibre}
still use their own first member. Put
$Z_j^+=((\Lambda_j^{\mathsf a})^c)^\sim$. The required load is
\begin{equation}
 \mathscr L_c(\mathcal S)
   =c\sum_{m=0}^k|D_m|
      +\sum_{j=1}^k
       (C_{\log}\ell_{j-1}+C_{\rm loc}+c)\#_j(Z_j^+),
       \qquad C_{\log},C_{\rm loc}\ge0.
 \label{hi:label-load}
\end{equation}
Every count is at its indicated level. At the last level the
outer remaining layer has at most $N_k$ sites. Older layers
may have been redone; their counts are not a single
present-day volume.

The scalar site load has a direct estimate. Since $\Omega_{m+1}$
is aligned at level $m+1$, for $m<k$ one has
\[
 |D_m|\le\#_m(\Omega_{m+1}^c)
   = L^4\#_{m+1}(\Omega_{m+1}^c)
   \le L^4\#_{m+1}(Z_{m+1}^{\mathsf a}).
\]
This uses inclusion of the action members in the determining
members and applies also to the full bare count at $m=0$.
The last layer costs at most $N_k$. For the source share charge
each dilated cell of $Z_j^+$ to a cell of $Z_j^{\mathsf a}$;
the one-layer enlargement costs a fixed multiplicity. The bounded
drift gives $\log\theta_{j-1}\le C\log\theta_j$. Thus the total
scalar load is bounded by
$cN_k+B(c)\sum_{j\le k}\log\theta_j\,\#_j(Z_j^{\mathsf a})$
for a fixed coefficient $B(c)$. To use it in the product over
histories, retain the ownership of each local action load.
The geometric load comparison is checked by ownership of
level-site addresses. Charge an unaltered layer site to its
action cell, with one-step fine-to-coarse multiplicity at
most $L^4$. Charge a site in a redone determining layer to
the local history of the mark that remade it, retaining its
original level. Prepared means are charged to their owned
level cells, not to an integrated-out variable at a different
level. The re-entry rule after \eqref{hi:action-sequence} specifies the actual
datum read there. With nested marks, use the latest operation
owning the remade slot; its inner record already accounts
for the earlier history.

The precise finite geometric assertion is that each such
read belongs to its owning cell collar, with a bounded
covering multiplicity $C_{\rm d}$. It involves the actual
determining entries, action entries and retained nests
together. Once this assertion holds, summing the assigned
addresses and using \eqref{hi:bounded-multiplicity}--\eqref{hi:read-site-budget} bounds the load by
current-level sites, live action-cell volume and marked
volumes. No factor counts the number of later reads of a
mark. In particular a coefficient of the following form
suffices to pay the redone load:
\begin{equation}
 b(c)=L^4(c+C_{\rm sf})+C+c,\qquad
 c^\sharp(c)=c+CC_{\rm d}(b(c)+c_{\rm e}).
 \label{hi:load-weight}
\end{equation}
The fixed $C$ includes the proved collar, scalar-number
and normalizer prices, and $C_{\rm sf}$ bounds final
unmarked local energy per site. Both expressions are
affine in $c$. The terms $C_{\log}\ell_{j-1}$ and
$C_{\rm loc}$ in \eqref{hi:label-load} enter the cell price before
the reserve is applied. Another price in an actual
presentation must be inserted with its actual sign
and magnitude, not omitted by this notation.

Separate the surviving real coarse Wilson factor.
The unmarked localized energy costs \(C_{\rm sf}N_k\)
plus the assigned defect-site prices. Rescaling and the
remaining Gaussian and chart factors have the local prices
proved above. The vacuum subtraction uses the following count.

%
%
\paragraph{Subtracting the label vacuum number.}
The label density carries \(e^{-\mathcal E_k(\mathcal S)}\).
Let \(E_{(k)}=-\mathcal E_k(\mathcal S_\emptyset)\) be the
vacuum number of the all-small-field prefix. Its subtraction
leaves the scalar \(-[E_{(k)}+\mathcal E_k(\mathcal S)]\).
Use the same one-step constants for the label and the
all-small-field term, with the actual ranges in each.

For a region \(Y\) made up of level-\(j\) blocks, put
\[
 n_{j-1}(Y)=|Y^{(j-1)*}|-(L^4-1)\#_j(Y).
\]
The first count consists of the incident free bonds after
pivot removal. The second counts the sites of the normalized
block gauge factors. In bare gauge dimension three, the
Gaussian scale contributes \(-\tfrac32\log\theta_{j-1}\)
per free bond and the gauge denominator returns the same
amount per gauge site. The remaining normalizer terms,
including the bounded ratio \(x_{j-1}/\theta_{j-1}\),
are the regular part of the count. On the full torus
\(n_{j-1}(T_\ell)=3(L^4-1)N_j\).

To compare a region with the torus, let \(b_j(Y)\) count
the bonds of level \(j\) with one end in \(Y\).
Counting bond endpoints outside \(Y\) shows that the
number of bonds with both ends outside is
\(4\#_j(Y^c)-b_j(Y)/2\). The corresponding fine count
uses \(L^4\#_j(Y^c)\) sites and \(L^3 b_j(Y)\) crossing
bonds. Subtract the coarse pivots and then the gauge sites.
This gives
\[
 n_{j-1}(T_\ell)-n_{j-1}(Y)
 =3(L^4-1)\#_j(Y^c)-\tfrac12(L^3-1)b_j(Y).
\]
Since \(b_j(Y)\le8\#_j(Y^c)\), this difference is nonnegative
and is at most \(3(L^4-1)\#_j(Y^c)\).
For \(Y=\Omega_j\) the logarithmic term in
\(\mathcal E_k(\mathcal S)-\mathcal E_k(\mathcal S_\emptyset)\)
is \(-\tfrac32\log\theta_{j-1}\) times that difference.
Thus its sign in the density after vacuum subtraction is
positive.

For the bounded activity terms, charge a missing base point
in \(Z_j^+\) by the rooted domain sum. Every other missing
pointing lies in a domain meeting \(Z_j\). Root there and
sum the at most \(CM^4(1+d_j(X))\) pointings with the spare
exponential decay. A missing remainder meets \(Z_j^+\)
because its range excludes the inner boundary layer.
These sums cost \(C\#_j(Z_j^+)\). The analogous source
difference uses the same ranges, and the Gaussian and
gauge counting factors cancel in that difference.

For the stage scalars, sum by the remade level \(j\).
At a fixed remade level the drop-out windows have disjoint
spatial pieces. Their strip and collar bonds are counted
in \(Z_j\), at most \(CL^4\#_j(Z_j)\) fine bonds in total.
The normalized Gaussian spectrum at the trivial pair
gives a bounded log contribution per integration coordinate.
Each pivot contributes its fixed flat Jacobian, and the
change of scale costs \(C\log\theta_{j-1}\) per bond.
This gives the stage price
\((C\log\theta_{j-1}+C)\#_j(Z_j)\).
Together these counts give the local price in
\eqref{hi:label-load}, including its logarithmic share.

For a later normalization \(E_{(n)}\) of the same run,
the full-lattice count gives
\[
 E_{(n)}-E_{(k)}
 =-\tfrac92(L^4-1)\sum_{j=k}^{n-1}
                 (\log\theta_j)N_{j+1}+\mathfrak r_k,
 \qquad |\mathfrak r_k|\le CN_k .
\]
Here the bounded one-step costs sum geometrically because
\(N_{j+1}=L^{-4(j+1-k)}N_k\). The logarithmic sum is
nonnegative before its minus sign, and is bounded by
\(C(1+\log\theta_k)N_k\) using bounded adjacent frame drift.
In particular the prefix floor expressed with \(E_{(n)}\)
loses at most \(e^{-CN_k}\) compared with its expression
using \(E_{(k)}\).

Under these local prices, the modulus of each
literal label term is bounded, before its last
Haar integration, by its positive history
envelopes, defect-site prices and
$e^{C_{\rm sf}N_k}$. The remaining real Wilson
factor is at most one. Normalized Haar
integration is at most one after positive
restrictions are dropped. For a drop-out restrict the sum to inner histories compatible
with its marked label. The volume bound on these histories is
$\mathsf v\ge V_{\rm in}$, where $V_{\rm in}$ is the inner
volume of that label. If $c'=c^\sharp(c)$, then on these histories
\[
 e^{c\mathsf v}\le e^{c'\mathsf v}e^{-(c'-c)V_{\rm in}} .
\]
The restricted mass with weight $c'$ is bounded by the full mass
with that weight. The extra factor pays the redone load of the
label. The full mass at weight $c$ includes other histories with
smaller inner volume, so the restriction precedes this estimate.
For a second-class mark use its full estimate at $c'$ by
monotonicity. Unfinished live
lineages use \eqref{hi:future-wait-credit}, not an additional
fictional marked mass.

Here is the final field-free summation. Root
each connected carried component at a cell of
the fixed $M$ partition at level $k$. There are
$N_k/M^4$ root addresses. Earlier healed
components have already been resummed into
local energies and are not independent roots
in this carried forest. Counting them again
would incorrectly introduce a bare-volume
root factor.

At each root the connected histories have
the bounded-degree address encoding of
\eqref{hi:animal-count} and the local code in \eqref{hi:ancillary-code}. Increase
the cell price by a fixed amount and use the
reserve at the correspondingly raised floor.
The residual positive weights $a_\gamma$ of
carried components then have
\begin{equation}
                  \sum_{\gamma\ \mathrm{at}\ r}a_\gamma\le1.
 \label{hi:loaded-root-sum}
\end{equation}
This is the connected-set sum in \eqref{hi:outer-entropy},
with the specified load in the price and
the marked masses already integrated.
Histories leading to the same label are
retained; the code pays for them. No
one-to-one identification of labels with
histories has been assumed.

More explicitly, let $a_\gamma$ be the
integrated positive envelope of the
particular component record, after its
marked operations have been dominated
and its owned defect prices assigned.
The per-record bound used here is
$a_\gamma\le h^{-f_\gamma}e^{-\sigma V_\gamma}$,
with $\sigma$ larger than its ancillary
code constant plus $2\log D+\log2$.
It follows from the event reserve with
this additional cell price only when
the complete kept forms and stored
coordinates of that literal term have
the uniform domination already stated.
They give a product envelope in the
disjoint owned coordinates; bounded
reads are first replaced by their
strong majorants. Counting that bound
by \eqref{hi:animal-count}--\eqref{hi:outer-entropy} proves \eqref{hi:loaded-root-sum}.
Thus \eqref{hi:loaded-root-sum} and \eqref{hi:loaded-label-sum} require
the product-coordinate and analytic
price estimates for each component
record. Positivity of the final gas
alone does not give these estimates.

Allowed collections are a subset of all
finite collections of distinct component
records. Dropping compatibility in a
nonnegative sum proves
\begin{equation}
 \sum_{\Gamma\ \mathrm{allowed}}\prod_{\gamma\in\Gamma}a_\gamma
       \le\prod_\gamma(1+a_\gamma)
       \le\exp\Bigl(\sum_\gamma a_\gamma\Bigr)
       \le e^{N_k/M^4}.
 \label{hi:forest-sum}
\end{equation}
For countably many records first prove
the finite inequality, then take the
increasing limit. The loaded integral
and label sum consequently satisfy
\begin{equation}
 \sum_{\mathcal S}e^{\operatorname{Re}\mathcal E_k(\mathcal S)}
       \tau^{\rm abs}_{\mathcal S,c}(\mathcal V_k)
                       e^{\mathscr L_c(\mathcal S)}
       \le e^{(c+C_{\rm sf}+M^{-4})N_k}.
 \label{hi:loaded-label-sum}
\end{equation}
Here $\tau^{\rm abs}_{\mathcal S,c}$ is
the literal measure with moduli and
$e^{cv^\natural}$ at its outer marks.
The term $cN_k$ is the current-level
part of \eqref{hi:label-load}; older loads have been
paid by the action-cell reserve or
the excess marked weight.

Thus \eqref{hi:loaded-label-sum} is a conditional deduction
from the exact normalized label
representation, local analytic
prices, read ownership and mass
bounds. It is not a consequence of
a final coarse gas alone. History
summation forms the local energies
and carried measures to which the
prices apply. Their formation and
common vacuum normalization must
first be matched. Complex sources
need uniform modulus prices on
their admitted domains. For a real
system, positivity of the remaining
Wilson term is used exactly in
the final Haar comparison.

\subsubsection{A weighted analytic extension from a prepaid real norm}

There is a second way to extend a parameter family. It does not
require the pointwise excess estimate in \eqref{hi:one-zone-test}. Its hypotheses are
instead an ordinary complex bound in the actual kernel norm and
a more strongly weighted bound along a real interval. We prove
the statement, including its norm and domain requirements.

Fix the stored-variable spaces and allowed argument fibres
$D_\omega(\xi)$. They may be ordinary cores or compatible
evaluation fibres. They do not change with the parameter being
continued. For a complex kernel family $F=(F_\omega)$ put
\begin{equation}
 \|F\|_{\mathfrak E_c}
   ={1\over h}\sum_\omega
      \int^*e^{cv(\omega,\xi)}
          \sup_{p\in D_\omega(\xi)}|F_\omega(\xi;p)|\,d\mu_\omega,
                    \qquad c\ge0.
 \label{hi:kernel-banach-norm}
\end{equation}
The weight is a fixed finite nonnegative Borel function. The
history volume is its principal example; the proof does not
require it to be integer valued. Reference measures, fibres
and weight are fixed throughout the argument. Kernels with
zero norm are identified. No measurable maximizing argument
is required by the definition.

\begin{lemma}[{Completeness and truncation}]\label{lem:history-completeness-truncation}
The finite-norm quotient in \eqref{hi:kernel-banach-norm} is a complex Banach space.
Multiplication by a Borel scalar of modulus one is an isometry.
Multiplication by $1_{\{v\le N\}}$, denoted $P_N$, has norm at
most one, and $P_N:\mathfrak E_0\to\mathfrak E_c$ has norm
at most $e^{cN}$.
\end{lemma}

\emph{Proof.} Homogeneity and the triangle inequality follow
from their pointwise supremum versions and \eqref{hi:upper-integral}. Zero norm
has a concrete meaning. If its supremum is $f$, choose Borel
majorants $g_n\ge f$ with weighted integrals at most $2^{-n}$.
Fatou gives $\liminf g_n=0$ almost everywhere, so $f=0$
outside a Borel null set. Conversely, zero outside such a set
has zero upper integral. The null set is common to every
argument in the fibre. Thus the quotient is a normed space with one common equivalence
relation for all arguments.

For completeness take a Cauchy sequence and a subsequence
$F_{n_j}$ whose consecutive norm differences are at most
$2^{-2j}$. On the countable disjoint union of the stored
spaces choose Borel majorants $b_j$ for the fibre suprema of
these differences, with total weighted integrals at most
$2^{-j}$. Countably many errors in the separate history
integrals can be allocated a further summable amount when
choosing these majorants. Tonelli gives $\sum_jb_j<\infty$
outside one Borel null set. Remove also the null sets on
which one of the countably many finite-norm kernels is
unbounded on its fibre. On the complement the subsequence
is uniformly Cauchy on every fibre, so it has a uniform
pointwise limit $F_\omega(\xi;\cdot)$. Define it to be zero
on the removed null set. Its difference from $F_{n_j}$ is
majorized by $\sum_{l\ge j}b_l$, whose weighted integral
tends to zero. Thus the subsequence, and then the original
Cauchy sequence, converges in norm. Uniform limits preserve
continuous field dependence or holomorphy on fixed field
cores when those properties are imposed. These give closed
subspaces if desired. The multiplier and projection claims
follow pointwise from their absolute values.\hfill$\square$

We will also use continuity from below of this upper integral
when the integrals have a finite bound. If $0\le f_n\uparrow f$
and $\int^* f_n\le M$, choose Borel majorants $g_n\ge f_n$
with integrals at most $M+1/n$. The functions
$h_n=\inf_{k\ge n}g_k$ are Borel, increase, and satisfy
$h_n\ge f_n$. Each has integral at most $M$, since it is
bounded by every $g_k$ with $k\ge n$. Measurable monotone
convergence shows that $h=\lim h_n$ majorizes $f$ and has
integral at most $M$. Monotonicity gives the reverse
inequality needed for equality of the increasing upper
integral limit. This applies on the disjoint union as well.

\paragraph{A criterion for ordinary Banach holomorphy.}
Suppose the actual representatives $K_\omega(z;\xi;p)$
are holomorphic in $z\in D_R$ at each fixed $(\xi,p)$,
and their $\mathfrak E_0$ norms are at most $M$ for
every $z$. This is a kernel bound, not just a bound
on an integrated scalar partition function. Assume
on each interior source circle $|z|=R_0<R$ that
\begin{equation}
 a_\omega(\xi,t)=\sup_{p\in D_\omega(\xi)}
                 |K_\omega(R_0e^{it};\xi;p)|
 \label{hi:measurable-circle-supremum}
\end{equation}
is jointly Borel. The reference measure and fibres
are independent of $z$.

For a fixed separable field core the usual
Carath\'eodory hypotheses give this measurability.
For each fixed field argument, Borel dependence
on $\xi$ and continuous dependence on the
source circle give joint measurability by
piecewise-constant circle approximations.
A countable dense set in the field core
then makes its supremum a countable supremum
of these jointly Borel functions. A fibre
depending on $\xi$ need not have a fixed
dense set; in that case \eqref{hi:measurable-circle-supremum}, or a
direct common Borel envelope, must be
checked on the actual evaluation rule.

Tonelli and the uniform ordinary norm give
\[
 {1\over h}\sum_\omega\int
       {1\over2\pi}\int_0^{2\pi}a_\omega(\xi,t)\,dt\,d\mu_\omega
                  \le M.
\]
Outside a Borel null set the circle average
is finite. Cauchy's formula at each field
argument therefore gives, simultaneously
for every argument and $|z|\le r<R_0$,
\begin{equation}
 \sup_{|z|\le r}\sup_p|K_\omega(z;\xi;p)|
    \le {R_0\over R_0-r}
                    {1\over2\pi}\int_0^{2\pi}a_\omega(\xi,t)\,dt.
 \label{hi:cauchy-common-envelope}
\end{equation}
The right side is a common Borel majorant
with total integral at most $MR_0/(R_0-r)$.
It was obtained first from the Cauchy
kernel, and then from Tonelli; no
uncontrolled supremum was interchanged
with the stored-variable integral.

The coefficient formula
$A_m=(2\pi)^{-1}\int K(R_0e^{it})
R_0^{-m}e^{-imt}\,dt$ has fibre supremum
at most $R_0^{-m}$ times the same circle
average. Thus
$\|A_m\|_{\mathfrak E_0}\le MR_0^{-m}$.
The Taylor remainder on $|z|\le r$ has
norm at most
$M(r/R_0)^{n+1}/(1-r/R_0)$ after degree
$n$, by the geometric series in Cauchy's
formula. It proves actual norm
holomorphy and identifies its coefficients
with the differentiated kernels outside
one common Borel null set. The centre
can be any point with an interior
source circle, so the proof works
throughout the domain.

This criterion applies when the ordinary
complex mass is available before the
weighted one. The ordinary complex
kernel bound, pointwise kernel
holomorphy and the stated measurability
then supply the Banach premise of
the following theorem. Scalar integral
holomorphy alone supplies none of
these kernel assertions.

\paragraph{An explicit harmonic barrier.}
Let $0<a<R$, remove the segment $[-a,a]$ from the disc
$D_R=\{|z|<R\}$, and for $z$ outside the segment let $w(z)$
be the root of
$w+w^{-1}=2z/a$ with $|w|>1$. There is exactly one such
root: the two roots have product one, and a root on the
unit circle would put $z$ on the segment. The roots are
simple off its endpoints. The unique choice is therefore
holomorphic there by its local simple-root formula, and
these choices agree on overlaps. Put
\begin{equation}
 g(z)=\log|w(z)|,\qquad H=\operatorname{arcosh}(R/a),
                         \qquad h_*(z)=1-g(z)/H.
 \label{hi:slit-barrier}
\end{equation}
The function $g$ is positive and harmonic, locally the
real part of a logarithm of its nonzero holomorphic root.
It tends to zero on the segment. The identity for $w$ and
the triangle inequalities give
\begin{equation}
          a\sinh g(z)\le|z|\le a\cosh g(z).
 \label{hi:barrier-bounds}
\end{equation}
Thus $g\ge\operatorname{arcosh}(R/a)$ on the outer
circle and $g\le\operatorname{arsinh}(r/a)$ on $|z|\le r$.
In particular $h_*\le0$ on the outer boundary, $h_*=1$
on the slit, and
\begin{equation}
 h_*(z)\ge h_0:=1-
   {\operatorname{arsinh}(r/a)\over\operatorname{arcosh}(R/a)}
                         \quad(|z|\le r).
 \label{hi:interpolation-height}
\end{equation}
Take $a=R/2$ and $r=R/4$. Then $h_0>1/2$: indeed
$\cosh(2\operatorname{arsinh}(1/2))=3/2<2$.
No assertion about a general harmonic-measure theorem is
needed for these explicit bounds.

\paragraph{The weighted extension theorem.}
Suppose $F:D_R\to\mathfrak E_0$ is holomorphic in the
Banach norm and $\|F(z)\|_{\mathfrak E_0}\le M$ there.
The case $M=0$ is immediate; take $M>0$ in the proof.
Suppose also that, on the real segment $[-a,a]$, its
same kernel classes belong to $\mathfrak E_C$ and obey
$\|F(x)\|_{\mathfrak E_C}\le M$. If $Ch_0>c$, then
\begin{equation}
 F:D_r\longrightarrow\mathfrak E_c\text{ is holomorphic},
 \qquad \sup_{D_r}\|F(z)\|_{\mathfrak E_c}\le M.
 \label{hi:weighted-extension}
\end{equation}
The same $M$ is retained. In particular, the choices
$a=R/2$, $r=R/4$, $C=2(c+1)$ give the conclusion with
an additional positive weight margin. If the two
hypotheses have different constants, their maximum
is a valid $M$.

\emph{Proof.} First truncate to $v\le N$. On any simply
connected patch off the slit choose a holomorphic function
$b$ with $\operatorname{Re}b=h_*$. Then
$e^{Cb(z)v}P_NF(z)$ is a holomorphic $\mathfrak E_0$
vector. To check the multiplier analyticity, its power
series and all derivatives converge in operator norm
on the truncation, since $v\le N$. Its norm is
\begin{equation}
 u_N(z)=\bigl\|e^{Ch_*(z)v}P_NF(z)\bigr\|_{\mathfrak E_0}.
 \label{hi:harmonic-weight-truncation}
\end{equation}
Different local harmonic conjugates differ by imaginary
constants. Their multipliers differ by diagonal phases
of modulus one. The norm in \eqref{hi:harmonic-weight-truncation} is consequently
single valued even though a global harmonic conjugate
around the slit need not exist.

The logarithm of the norm of a holomorphic Banach vector
is subharmonic. A direct proof at a nonzero value chooses
a norm-one complex linear functional attaining the norm
there by Hahn--Banach, applies the scalar logarithmic
submean inequality, and bounds its values on the circle
by the vector norm. At a zero, use the value minus
infinity and the same upper semicontinuity argument.
Thus $\log u_N$ is globally subharmonic off the slit,
because its local definitions agree.

On the slit the fixed truncation is continuous in the
$\mathfrak E_0$ norm and $h_*$ tends to one. Its boundary
norm is at most the assumed real $\mathfrak E_C$ norm,
hence at most $M$. On the outer circle $h_*\le0$;
the ordinary bound and $v\ge0$ give boundary upper
limit at most $M$. For this fixed $N$ the function is
bounded above throughout the punctured disc, and the
subharmonic maximum principle gives $u_N(z)\le M$.
This proof uses the norm-continuous extension of
$F$ across the real interval in $\mathfrak E_0$.
Scalar partition holomorphy would not supply it.

At any point of $D_r$ off the slit, $Ch_*(z)\ge Ch_0$.
Monotonicity of the weight gives
$\|P_NF(z)\|_{\mathfrak E_{Ch_0}}\le M$.
The continuity-from-below argument proves that
$F(z)$ belongs to that weighted space, with the
same bound. On $D_r\cap[-a,a]$ it follows
directly from the assumed $\mathfrak E_C$ bound.
The identification with the original $\mathfrak E_0$
class is exact, since these spaces have the same
common-null-set identification and $P_N$ exhausts
the fixed finite weight.

Choose $\delta>0$ with $c+\delta\le Ch_0$. Uniformly
on $D_r$ one now has
\begin{equation}
 \|(1-P_N)F(z)\|_{\mathfrak E_c}
        \le e^{-\delta N}
                    \|F(z)\|_{\mathfrak E_{c+\delta}}
        \le M e^{-\delta N}.
 \label{hi:weighted-analytic-tail}
\end{equation}
Each $P_NF$ is holomorphic into $\mathfrak E_c$, because
$P_N:\mathfrak E_0\to\mathfrak E_c$ is bounded. Their
uniform convergence on the smaller disc proves
holomorphy of the limit there. One can see this directly
by taking their Banach-valued Cauchy integrals on any
interior circle: uniform norm convergence permits the
limit in that formula and its differentiated formula.
This proves \eqref{hi:weighted-extension}.\hfill$\square$

It follows in particular that
\begin{equation}
       \|F^{(m)}(0)\|_{\mathfrak E_c}
                        \le m!M(4/R)^m,
       \qquad m=0,1,\ldots,
 \label{hi:interpolated-cauchy}
\end{equation}
for the quarter-disc choice. There is no factor
growing with the number of histories, no reciprocal
survival probability in the radius, and no extra
constant in front of $M$. The cost is the explicitly
larger prepaid real weight and a fixed radius loss.
The coupling floor for the real mass remains affine
in the requested final coefficient $c$.

When the formed kernels have their pointwise
holomorphic representatives, these norm derivatives
are their actual kernel derivatives. Here is a way
to avoid separate exceptional sets at uncountably
many parameters. The norm Taylor coefficients on
an interior circle satisfy a summable geometric
bound on any smaller circle. Choose Borel majorants
for their fibre suprema with a summable total
integral. Outside one Borel null set, the resulting
power series converges uniformly on every smaller
parameter disc and every fibre. At a countable
dense set of parameters it agrees with the original
kernel classes outside a further countable union
of Borel null sets. Pointwise holomorphic
uniqueness then makes it agree for all parameters
there. Differentiating this series gives the
actual differentiated kernels on that same full
set. No point evaluation is asserted to be a
bounded functional on $\mathfrak E_0$.

\paragraph{What would make this an actual source extension.}
For a fixed cutoff, use exactly its admitted
field or evaluation fibres and the fixed
history volume. The required ordinary family
is the complete kernel, including its
normalizers, not just the scalar integral of
that kernel. Its $\mathfrak E_0$ holomorphy can
be established from jointly measurable,
pointwise holomorphic kernels with a common
integrable fibre-supremum envelope on each
compact parameter disc. The dominated
Cauchy argument used before \eqref{hi:mass-cauchy}, now in
this complete normed space, proves that
assertion. A proved norm power series is
another sufficient verification. Pointwise
holomorphy without this norm control is not
the premise of \eqref{hi:weighted-extension}.

The other premise is the strongly weighted
real family over the same field core,
uniformly for every real parameter in the
stated interval. Real field values only on
the physical-current graph would not be
that premise. The theorem neither enlarges
the current core nor proves admission of a
noncovariant source array to the ordinary
complex source class. A moving domain must
first be replaced by a single common
admitted domain with the two stated bounds.
Here ``real'' concerns the source parameter.
If the desired norm takes its supremum
over a complex field core, the real-source
premise already takes its supremum over
that same complex field core.

If these checks hold with
$M=e^{-a_\star P_i}$, the theorem proves
the same weighted complex mass and its
Cauchy bounds on the smaller source disc.
It is thus an alternative to proving the
pointwise one-zone estimate in \eqref{hi:one-zone-test}.
Each route has a finite, testable analytic
input. The local-excess route uses actual
carrier and read bounds. The interpolation
route uses an ordinary complex kernel
bound and arbitrarily prepaid strong real
weights on the same admitted core. Neither
route follows from compactness of source
parameters or scalar partition holomorphy.

\paragraph{A dimension-free real-form corollary.}
Let $\mathcal B$ be a complex Banach source space with
an isometric conjugation $\kappa$. Its real form is
$\mathcal B_{\mathbb R}=\{u:\kappa u=u\}$. Thus every
$h$ has $h=h_1+ih_2$ with real $h_1,h_2$ and
$\|h_j\|\le\|h\|$. Finite coordinate arrays with
their usual real-form supremum or sum norms have
this property, with constants independent of the
number of coordinates.

Suppose $F$ is Fr\'echet-holomorphic from
$\{\|h\|<R\}$ into $\mathfrak E_0$, bounded by $M$,
and for every real $u$ with $\|u\|<R$ it obeys
the $\mathfrak E_C$ bound $M$ on the same field
fibres. For any $c\ge0$ the prepaid coefficient
$C=2(c+2)$ suffices to give
\begin{equation}
 \sup_{\|h\|<R/8}\|F(h)\|_{\mathfrak E_{c+1}}\le M,
 \qquad
 F:\{\|h\|<R/8\}\longrightarrow\mathfrak E_c
                      \text{ is Fr\'echet-holomorphic}.
 \label{hi:real-form-extension}
\end{equation}
The number of source coordinates does not enter the constant.

\emph{Proof.} Fix $h$ in that smaller ball. If $h_2=0$,
the asserted weighted bound is part of the real premise.
Otherwise put $v=h_2/\|h_2\|$ and consider the single
translated complex line
$f(z)=F(h_1+zv)$. For $|z|<R/2$ its source norm is
less than $5R/8$, so its ordinary bound is $M$.
For real $|z|\le R/4$ it is a real source of norm
less than $3R/8$, and has the prepaid weighted
bound. Apply \eqref{hi:weighted-extension} with ordinary radius $R/2$,
real half-segment $R/4$ and target radius $R/8$.
Since $h_0>1/2$, the chosen $C$ is strictly more
than sufficient for weight $c+1$. Evaluation at
$z=i\|h_2\|$ gives the first assertion.

This estimate is uniform on the whole smaller
source ball, not just on one line. Therefore
$\|(1-P_N)F(h)\|_{\mathfrak E_c}\le M e^{-N}$
uniformly there. Each $P_NF$ is Fr\'echet-holomorphic
into $\mathfrak E_c$ by its bounded linear
projection from $\mathfrak E_0$. Their uniform
norm limit is Fr\'echet-holomorphic. One direct
verification uses the Cauchy formula on every
complex line in an interior source ball and
polarizes its homogeneous Taylor coefficients.
Uniform convergence gives convergence of those
coefficients; polarization gives convergence
in each multilinear norm, and the Cauchy
bounds make their Taylor series converge
uniformly in a smaller ball. Its sum is
the asserted limit.\hfill$\square$

The derivative constants in this assertion can
be made explicit as well. Write $\rho=R/8$.
The $m$-th homogeneous Taylor polynomial at
zero is bounded by $M\rho^{-m}$ on the unit
source ball, by scalar Cauchy on each line.
For its symmetric $m$-linear form, the
coefficient of $z_1\cdots z_m$ in its value
at $\sum z_jh_j$ is $m!$ times that form.
Cauchy's formula on $m$ unit circles bounds
this coefficient by $m^m$ times the
polynomial bound, since $\|\sum z_jh_j\|\le m$
when $\|h_j\|\le1$. Consequently
\begin{equation}
 \|D^mF(0)\|_{\mathcal B^m\to\mathfrak E_c}
       \le m^mM(8/R)^m
       \le m!M(8e/R)^m.
 \label{hi:source-frechet-cauchy}
\end{equation}
These are multilinear operator norms. No
coordinatewise entropy or factor equal to
the dimension of the source space has
been introduced.

For a normalized excess with $F(0)=0$,
scalar Schwarz on every source line and
Hahn--Banach scalarization give
$\|F(h)\|_{\mathfrak E_c}\le(8M/R)\|h\|$
in the smaller ball. For a general kernel
family the difference $F(h)-F(0)$ has the
corresponding bound with $2M$. This is a
weighted estimate on the complete kernel,
not a bound on a chosen logarithm of that
kernel. A logarithm requires its own
zero-free normalization.

The corollary is useful precisely because
its premises say which source arrays are
allowed. If the real source bound is
known only on a gauge-covariant subspace,
the translated real point $h_1$ and
direction $v$ must both lie in that
subspace. The theorem then extends its
complexification, not the whole ambient
array space. It does not create a
noncovariant extension from data on
the physical-current graph. If the
ordinary complex family is available
on a radius depending on an order
parameter $p_n$, that same radius,
with the fixed factor eight, appears
in \eqref{hi:source-frechet-cauchy}; no larger radius has
been supplied by interpolation.

For a genuine array extension one must
verify all three statements: the ordinary
kernel map is Fr\'echet-holomorphic and
bounded in $\mathfrak E_0$ on its actual
complex source ball; the real weighted
bound holds on the full real ball used
in the argument; and all those families
use the same field/read fibres and
fixed reference data. These are
formation and domain assertions. Once
they hold, the weighted extension,
its source-dimension independence and
its derivative bounds follow from
\eqref{hi:kernel-banach-norm}--\eqref{hi:source-frechet-cauchy}.

\subsubsection{The finite estimates to be matched}

The weighted-history assertion
requires the following estimates
for finite operations. Their quantifiers
are uniform in the height, compatible
torus, component, local history,
stored variables on its support,
and the full joint argument domain.
Compactness at a single cutoff is
not that uniform statement.

First, every stored-coordinate
formula must be the literal
operation in \eqref{hi:local-operation} or \eqref{hi:drop-factor}, with
its fixed measures, normalizers,
tests and frame factors. Later
reads must satisfy \eqref{hi:read-coverage}. The
five region lists and retained
seam prescriptions determine this
check. Full re-entry must use the
current datum on its remade shell,
not an earlier field already
integrated there.

Second, every prepared, localized
and substituted datum must belong
to its joint background/current
core. Its complete kept forms
must have the lower bound and
reference domination in \eqref{hi:reference-gaussian-integral}
or \eqref{hi:reference-cutoff-gaussian} on their respective
coordinate domains, with local dimension counted once.
Independent-current completion
needs its projected gauge and
other current terms. It does not
follow from setting a coarse
averaging constraint equal to
the identity. The flat precision
and one-step perturbation results
verify restricted cases, not the
whole all-layer assertion.

Third, the sum assigned to $V_X$
needs \eqref{hi:pair-sup-envelope} on those cores for
all supported stored fields. Its
counterterm, vacuum and boundary
conventions belong to this bound.
Oscillation of isolated terms
does not bound their constants.
A pointwise KP logarithm does
not supply the integrable
supremum without its common
envelope. The local formation
and Gaussian theorems address
parts of this check under
their primitive matrix and
background hypotheses.

Fourth, failed tests need their
certificates and owned energy
sets in \eqref{hi:event-energy}. The core
Wilson term keeps its stopped
real weight, and every compatible
drop-out evaluation must preserve
the certificates used under its
supremum. A last test requiring
an extra collar supplies its
renewed gain only in a later
read with that collar present.

For parameters these checks
hold on one fixed-frame domain.
The finite gauge-normalizer
extension and the complete
Gaussian half-disc criterion
have been proved here. For a second-class component formed by the island kit, the
\hyperlink{hi:formed-island-pair-read}{pair reader} gives the
exact source excess on the admitted output core. The bound
following \eqref{hi:actual-source-recursion} uses the formed
carrier estimates there. At other weighted source lines or drop-out fibres
we must check \eqref{hi:source-excess-budget} on the actual
read, for example by the carrier criterion
\eqref{hi:carrier-site-criterion}. \eqref{hi:one-zone-test} is a precise one-zone
test for a finer weighted
source extension. Ordinary
source compactness and
holomorphy alone do not prove
that test.

Alternatively the full kernel
family can meet the explicit
same-core ordinary complex and
prepaid real norm hypotheses
of \eqref{hi:weighted-extension} or \eqref{hi:real-form-extension}. Those
theorems then prove its
weighted extension without
the pointwise excess test.
They do not assert those
formation or domain premises
for a source-array class
outside the stated hypotheses.

The loaded global comparison
also needs the literal label
representation, common vacuum
price, local unmarked energy
price and finite read-ownership
map. These are more specific
identifications than the claim
that a density is a sum of
labels. Once matched,
\eqref{hi:volumes}--\eqref{hi:loaded-label-sum} give the
counting, reserve, integration
and summation with constants
uniform in height and volume.
They retain the field,
current and history types
needed to apply those
scalar consequences.


% END MODULE chapters/history_mass_program.tex

% BEGIN PROOF EXPANSION formation/stage_formation_candidate.tex
\subsection{Complex stage integrals with retained tests}
\label{candidate:stage-formation}

Preparatory integrations of the large-field operation integrate a
stored fluctuation on a strip and a new fluctuation on the adjacent
collar; see \cite{B89a,B89b}. The earlier stage terms enter later
integrals. The estimates below control this accumulation on one
component. They also give a branch of the logarithm when the actual
tested integral is close, pointwise in its density, to a positive
reference integral. All coordinate maps and all retained factors are
specified before this comparison is applied.

\subsubsection{The tested local integral}

Let \(b\in\mathbb R^m\) be the actual free real coordinates of a
finite stage chart. A parameter \(\sigma\) includes its geometry,
real background, real stored spectators and retained tests. Fix
\[
 S_\sigma\subset[-T,T]^m,\qquad T\ge1,
 \qquad Q_{0,\sigma}=Q_{0,\sigma}^{\mathsf T}>0.
\]
The set \(S_\sigma\) is Borel. It can be disconnected and can contain
arbitrary carried sharp tests. On each complex parameter domain used
below it is fixed. Put
\begin{equation}
 I_{0,\sigma}
 =c_{0,\sigma}\int_{S_\sigma}
       e^{-b^{\mathsf T}Q_{0,\sigma}b/2+V_{0,\sigma}(b)}\,db,
 \qquad 0<I_{0,\sigma}<\infty,
 \label{sf:positive-reference}
\end{equation}
where \(c_{0,\sigma}>0\) and \(V_{0,\sigma}\) is real. No uniform
lower bound on \(I_{0,\sigma}\) is assumed.

For \(z\) in a complex polydisc \(\mathcal D\), suppose the actual
integral has the exact representation
\begin{equation}
 I_\sigma(z)=c_{0,\sigma}\int_{S_\sigma}
 e^{-b^{\mathsf T}Q_{0,\sigma}b/2+V_{0,\sigma}(b)}
 e^{H_\sigma(z,b)}\,db.
 \label{sf:exact-density}
\end{equation}
Here \(H\) is Borel in the real variables and holomorphic in \(z\),
with a uniform bound on the polydisc. In an application it contains
\begin{equation}
 H=\ell-\tfrac12b^{\mathsf T}(Q-Q_0)b
       +(V-V_0)+\log(J/J_0).
 \label{sf:complete-change}
\end{equation}
The term \(\ell\) includes every change of a Gaussian normalizer.
The factor \(J/J_0\) includes the nonlinear constraint determinant,
Haar densities and other chart Jacobians. These terms are retained
even when the quadratic exponent alone is small. An analytic
reciprocal determinant is used at complex arguments; an absolute
determinant is used only at real arguments.

\begin{proposition}[A tested integral in a disc about its reference]
\label{prop:stage-sector}
If \(\sup_{\sigma,z,b\in S_\sigma}|H_\sigma(z,b)|\le\delta\le1/8\),
then
\begin{equation}
 \left|\frac{I_\sigma(z)}{I_{0,\sigma}}-1\right|
 \le e^\delta-1<1,
 \qquad
 \left|\log\frac{I_\sigma(z)}{I_{0,\sigma}}\right|\le4\delta.
 \label{sf:sector-bound}
\end{equation}
The logarithm is the analytic branch defined by the power series
about \(1\). For a polydisc with coordinate radii \(R_i\), centred
at \(z_*\), and a multi-index \(\alpha\ne0\),
\begin{equation}
 \left|\partial^\alpha_z\log I_\sigma(z_*)\right|
 \le4\delta\,\alpha!\prod_iR_i^{-\alpha_i}.
 \label{sf:log-derivatives}
\end{equation}
Let \(\nu_{0,\sigma}\) be the probability measure obtained by
dividing the integrand in \eqref{sf:positive-reference} by
\(I_{0,\sigma}\), and let \(\nu_{z,\sigma}\) be its complex
reweighting by \(e^H\). Then
\begin{equation}
 \int|d\nu_{z,\sigma}-d\nu_{0,\sigma}|
 \le8\delta.
 \label{sf:normalized-law}
\end{equation}
These bounds are uniform when the positive reference masses tend
to zero.
\end{proposition}

\begin{proof}
By \eqref{sf:exact-density}, the ratio is
\(w=\int e^H\,d\nu_0\). Thus \(|w-1|\le r:=e^\delta-1\).
For \(\delta\le1/8\), one has \(r<1/4\), and
\(|\log w|\le r/(1-r)\le4\delta\). Dominated integration gives
holomorphy of the ratio. Its values remain in the simply connected
disc \(|w-1|<1\); hence this logarithm is coherent throughout
\(\mathcal D\). Apply the several-variable Cauchy formula on
smaller concentric polydiscs, then increase their radii to \(R_i\).
The reference mass is independent of \(z\), so its logarithm has
no derivative in \eqref{sf:log-derivatives}.

The complex density relative to \(\nu_0\) is \(e^H/w\). Therefore
\[
 \int\left|\frac{e^H}{w}-1\right|d\nu_0
 \le\frac{\int|e^H-1|d\nu_0+|w-1|}{|w|}
 \le\frac{2r}{1-r}\le8\delta.
\]
No inverse bound on the unnormalized mass has entered this proof.
\end{proof}

\subsubsection{The complete Gaussian and coarea factors}

The comparison can be verified using local matrices before taking
any absolute values in an auxiliary Gaussian. In the notation of
\cref{thm:localized-complex-Gaussian}, the exact marginal has
covariance \(D=C+C\Gamma\Gamma^{\mathsf T}C\). Its auxiliary
dimension need not be local. The following bound uses only the
number \(m\) of free coordinates remaining in the marginal.

\begin{lemma}[The Gaussian density ratio on a tested box]
\label{lem:stage-gaussian-ratio}
Suppose
\[
 \|C_0\|+\|C_0^{-1}\|+\|\Gamma_0\|\le K,
 \qquad
 \|C-C_0\|+\|\Gamma-\Gamma_0\|\le\epsilon,
\]
where \(C_0\) is real symmetric positive definite and
\(\Gamma_0\) is real, and take
\(\epsilon\le\epsilon_*(K)\). Let
\(D_0=C_0+C_0\Gamma_0\Gamma_0^{\mathsf T}C_0\).
The normalized marginal density ratio has an analytic logarithm
\begin{equation}
 L_D(b)=-\tfrac12\log\frac{\det D}{\det D_0}
       -\tfrac12b^{\mathsf T}(D^{-1}-D_0^{-1})b,
 \qquad
 \sup_{|b|_\infty\le T}|L_D(b)|
       \le C_K\epsilon m(1+T^2).
 \label{sf:gaussian-ratio}
\end{equation}
The bound is independent of the auxiliary dimension.
\end{lemma}

\begin{proof}
Expanding the products that define \(D\) gives
\(\|D-D_0\|\le C_K\epsilon\), since all rectangular factors
are bounded in operator norm. The real matrix \(D_0\) has
\(D_0\ge C_0\ge K^{-1}I\). For
\(E=D_0^{-1/2}(D-D_0)D_0^{-1/2}\), decrease
\(\epsilon_*\) so that \(\|E\|\le1/2\). The convergent
matrix logarithm gives
\[
 \left|\log\frac{\det D}{\det D_0}\right|
 =|\operatorname{Tr}\log(I+E)|
 \le m\sum_{q\ge1}\frac{\|E\|^q}{q}
 \le C_Km\epsilon.
\]
The resolvent identity gives
\(\|D^{-1}-D_0^{-1}\|\le C_K\epsilon\). Finally
\(|b|^2\le mT^2\). The marginal identity in
\cref{thm:localized-complex-Gaussian} must be applied before this
calculation; a total-variation estimate on the unreduced joint
Gaussian would charge its auxiliary dimension.
\end{proof}

Here is the corresponding finite coarea calculation. Suppose the
actual constraints in a chart are solved for three-component pivot
coordinates, and their real derivative matrices are
\(M_{0,a}(b)\), \(1\le a\le k\), with positive determinant.
Suppose their complex continuations satisfy
\begin{equation}
 E_a=M_aM_{0,a}^{-1}-I,
 \qquad\sup_{z,b}\|E_a\|\le\epsilon_a<1/2.
 \label{sf:pivot-assumption}
\end{equation}
Then the logarithm of the ratio of the inverse pivot determinants is
\begin{equation}
 L_{\rm piv}(z,b)=-\sum_{a=1}^k\operatorname{Tr}\log(I+E_a),
 \qquad |L_{\rm piv}|\le6\sum_{a=1}^k\epsilon_a.
 \label{sf:pivot-ratio}
\end{equation}
Indeed the power series used in the preceding proof applies to
each three-by-three matrix. The identity powers of \(L\) in the
per-length pivot convention cancel in the ratio. Their real
normalization remains in \(I_0\). Equation
\eqref{sf:pivot-assumption} requires the actual nonlinear pivot
chart on the tested domain.

For \(\operatorname{SU}(2)\), write \(a\in\mathbb C^3\) in the orthonormal Pauli
basis, using the bilinear square \(a\cdot a\). Apart from its fixed
normalizing constant, the analytic Haar density in exponential
coordinates is
\[
 j(a)=\left(\frac{\sin\sqrt{a\cdot a}}
                       {\sqrt{a\cdot a}}\right)^2.
\]
The quotient at zero is defined by its power series. On
\(\|a\|_2,\|a_0\|_2\le1/4\), it has a nonzero analytic
logarithm and
\begin{equation}
 |\log j(a)-\log j(a_0)|
 \le C(\|a\|_2+\|a_0\|_2)\|a-a_0\|_2.
 \label{sf:haar-ratio}
\end{equation}
To see this, put \(f(s)=\sum_{q\ge0}(-1)^qs^q/(2q+1)!\).
For \(|s|\le1/16\), its power series bounds \(|f-1|\) away
from one and bounds \(|f'|\) by an absolute constant. Thus
\(2\log f(s)\) has bounded derivative. Integrate that derivative
along the segment between \(a_0\cdot a_0\) and \(a\cdot a\).
Summing \eqref{sf:haar-ratio} over the chart bonds accounts for all
Haar density ratios. If the coupling is fixed, the rescaling
\(a=gb\) has the same Jacobian in the two integrals. A change of
\(g\) contributes an additional explicit scalar factor.

For comparison, the bare Wilson Taylor remainder also has a direct
local bound. For unitary background factors \(V_i\) and complex
Lie algebra elements \(B_i\), put
\[
 F(t)=g^{-2}\left(1-\tfrac12\operatorname{Tr}
                    \prod_{i=1}^4e^{itgB_i}V_i\right),
 \qquad S_B=\sum_i\|B_i\|_{\rm op}.
\]
Three differentiations and the multinomial theorem give
\(|F'''(t)|\le CgS_B^3e^{gS_B}\) for \(0\le t\le1\).
The integral Taylor formula therefore yields
\begin{equation}
 |F(1)-F(0)-F'(0)-\tfrac12F''(0)|
       \le CgS_B^3e^{gS_B}.
 \label{sf:bare-Wilson-remainder}
\end{equation}
On a bounded stage box its sum costs at most
\(Cg\,mT^3e^{CgT}\) when each coordinate meets a bounded number
of plaquettes. Composition with a nonlinear constrained response
requires the first three derivatives of that actual response.

\begin{proposition}[Response of the constrained density]
\label{prop:stage-constrained-response}
Let the actual constraint chart solve
\(F(z,b,y(z,b))=0\), where \(F\) is holomorphic and
\(M=\partial_yF\) is invertible throughout the admitted domain.
Suppose the complete stage density on its fixed tested set is
\[
 a(z,b)\exp[-S(z,b,y(z,b))]\,(\det M(z,b))^{-1}.
\]
All factors are holomorphic and have the uniform bounds needed to
differentiate under the finite integral; \(a\) includes the Haar
densities and numerical normalizers. A local analytic determinant
branch agrees with the positive coarea density at the real
reference. With dots denoting differentiation at fixed \(b\),
\begin{align}
 \dot y={}&-M^{-1}F_z,\notag\\
 \ddot y={}&-M^{-1}
  \big(F_{zz}+2F_{zy}\dot y+F_{yy}[\dot y,\dot y]\big),
 \label{sf:pivot-response}\\
 L={}&\partial_z\log a-S_z-S_y\dot y
                         -\operatorname{Tr}(M^{-1}\dot M).
 \label{sf:complete-score}
\end{align}
On a nonzero domain supplied by \cref{prop:stage-sector}, the
normalized complex stage expectation satisfies
\begin{align}
 \partial_z\log I={}&\langle L\rangle_z,\notag\\
 \partial_z^2\log I={}&\langle\dot L\rangle_z
                  +\langle L^2\rangle_z-\langle L\rangle_z^2,
 \label{sf:constrained-response}\\
 \partial_z\langle A\rangle_z={}&
       \langle\dot A\rangle_z+
       \langle AL\rangle_z-\langle A\rangle_z\langle L\rangle_z.
 \label{sf:constrained-observable-response}
\end{align}
The second determinant contribution in \(\dot L\) is
\begin{equation}
 \operatorname{Tr}
   (M^{-1}\dot M M^{-1}\dot M-M^{-1}\ddot M).
 \label{sf:coarea-second-response}
\end{equation}
For successive triangular eliminations the traces in these
formulas are summed over their pivot blocks.
\end{proposition}

\begin{proof}
Differentiate the exact constraint at fixed free coordinates.
Its first derivative is \(F_z+M\dot y=0\); differentiating
again gives \eqref{sf:pivot-response}. Differentiate the complete
logarithmic density, including the reciprocal determinant, to
obtain \eqref{sf:complete-score}. Jacobi's formula and
\(\partial_zM^{-1}=-M^{-1}\dot M M^{-1}\) give
\eqref{sf:coarea-second-response}. The quotient rule after
differentiation under the integral gives
\eqref{sf:constrained-response} and
\eqref{sf:constrained-observable-response}. The tests do not
produce a boundary term because they have not been varied.
\end{proof}

For example, suppose \(\|M^{-1}\|\le K_M\),
\(\|F_z\|\le a\), \(\|F_{zz}\|\le b\),
\(\|F_{zy}\|\le c\) and \(\|F_{yy}\|\le d\), in the
specified Euclidean and multilinear operator norms. Then
\[
 \|\dot y\|\le K_Ma,\qquad
 \|\ddot y\|\le K_M(b+2cK_Ma+dK_M^2a^2).
\]
For a pivot block of dimension \(d_p\),
\(|\operatorname{Tr}(M^{-1}\dot M)|
 \le d_pK_M\|\dot M\|\). These estimates identify which
constraint derivatives a quantitative source bound needs. In
particular the term \(S_y\dot y\) is retained unless the actual
stationarity equations cancel it. A derivative of a Gaussian
potential at fixed chart alone would omit these contributions.

\subsubsection{Resetting the stored fields}

A near stored field is a real variable under a sharp test. Its
reset is represented by a complex switch. Let \(a_v\) be such a
field, and let \(\mu_v\) range over \(|\mu_v|<\Lambda_v\),
where \(\Lambda_v>1\). The term at \(\mu_v=0\) is required to
be independent of \(a_v\); at \(a_v=0\) it is independent of
\(\mu_v\). At a far stored field the same notation represents the
radial parameter in the holomorphic substitution
\(a_v\mapsto\mu_va_v\). The carried tests are kept at their real
arguments in both cases.

\begin{lemma}[Block differences]
\label{lem:stage-block-differences}
Let \(f\) be holomorphic and bounded by \(B\) on
\(\prod_{v\in T}\{|\mu_v|<\Lambda_v\}\), together with its
other holomorphic arguments. For \(A\subset T\), set
\[
 f_A=\sum_{E\subset A}(-1)^{|A|-|E|}f(\mathbf1_E,\mathbf0_{T\setminus E}),
 \qquad q_v=(\Lambda_v-1)^{-1}.
\]
Then
\begin{equation}
 f(\mathbf1_T)=\sum_{A\subset T}f_A,
 \qquad |f_A|\le B\prod_{v\in A}q_v,
 \qquad
 \sum_{A\ne\emptyset}|f_A|
 \le B\left(\exp\sum_{v\in T}q_v-1\right).
 \label{sf:block-difference-bound}
\end{equation}
If \(f\) is independent of \(a_v,\mu_v\), then \(f_A=0\)
whenever \(v\in A\). The differences preserve the full declared
radii of the remaining stored fields and switches.
\end{lemma}

\begin{proof}
The finite inclusion-exclusion identity gives the first equality.
For one variable and \(1<r<\Lambda_v\), Cauchy's formula gives
\[
 f(1)-f(0)=\frac1{2\pi i}\int_{|w|=r}
                         \frac{f(w)}{w(w-1)}\,dw.
\]
Apply this identity in each variable of \(A\), with all other
variables in \(T\) set to zero. The bound is
\(B\prod_{v\in A}(r_v-1)^{-1}\); increase \(r_v\) to
\(\Lambda_v\). Summing this bound gives
\(B[\prod_v(1+q_v)-1]\), which is bounded by the last expression
in \eqref{sf:block-difference-bound}. Variables absent from the
function have zero difference. The proof fixes the other arguments
at arbitrary points of their full domains, so it spends none of
their radii.
\end{proof}

For a family \((f_e)\), let \(T_e\) be its set of read blocks and
\(B_e\) its analytic bound. Its incidence norm on the integrated
blocks is
\begin{equation}
 J=\sup_{v\in T}\sum_{e:v\in T_e}B_e.
 \label{sf:incidence-norm}
\end{equation}
If \(\sum_{v\in T_e}q_v\le1\) for every \(e\), then
\begin{equation}
 \sum_e\sum_{A\ne\emptyset}|(f_e)_A|
 \le2\sum_{v\in T}q_v\sum_{e:v\in T_e}B_e
 \le2J\sum_{v\in T}q_v.
 \label{sf:incidence-reset}
\end{equation}
This counts blocks actually read by a term. It introduces no
Gaussian coordinates on the rest of that term's domain.

A background displacement has the same estimate in an auxiliary
parameter. If \(f_e(P_e(\lambda))\) is holomorphic and bounded by
\(B_e\) on \(|\lambda|<a_e^{-1}\), with \(a_e\le1/2\), then
\begin{equation}
 |f_e(P_e(1))-f_e(P_e(0))|\le2a_eB_e.
 \label{sf:response-difference}
\end{equation}
Consequently, if \(P_e(1)\) is the old background read and
\(P_e(0)\) is the new one, the exact split is
\begin{align}
 f_e(P_e(1);a)-f_e(P_e(0);0)
 ={}&f_e(P_e(1);a)-f_e(P_e(0);a)\notag\\
   &+f_e(P_e(0);a)-f_e(P_e(0);0).
 \label{sf:two-differences}
\end{align}
The first line is bounded with \eqref{sf:response-difference}; the
second uses \eqref{sf:incidence-reset}. A displacement that fixes
the stored data at their carrying sites is needed for this identity
to be the split of the actual operation.

\subsubsection{Uniform accumulation in a finite sequence of stages}

We give the quantitative conclusion in a form that can be applied
to a fixed chain of actual tested charts. This is a local assertion
for one component. Let \(\theta>e\), \(\ell=\log\theta\), and
\(r\ge1\). Fix constants independent of the stage number, and assume
\begin{equation}
 N\le C_N\ell^r,\qquad K_s\le C_s\ell^{4r},\qquad
 m_n\le C_m\ell^{4r},\qquad T_n\le C_T\ell^p.
 \label{sf:stage-counts}
\end{equation}
At most \(K_s\) stored blocks are reset at a stage. Put
\(\Lambda=\ell^{6r+1}\), and suppose every such block has this
switch radius. The bounds hold on the full domains of the switches
that remain after integration.

For each remaining stored datum, assume that setting its switch to
zero makes both the complete reference kernel and the inserted
integrand independent of that datum. Setting the datum to zero
makes them independent of its switch. This includes the retained
tests and any normalization of the reference kernel. Thus integration
preserves the two independence properties defining a switch.

Let \(J_0\) bound the incidence norm of the earlier terms. Suppose
that each earlier stage contributes at most one function at each
block, with norm at most one; its incidence norm before stage
\(n\) is then at most \(J_0+n\). The admitted background
displacement satisfies the separate budget
\begin{equation}
 \sum_ea_{n,e}B_{n,e}
       \le C_a\ell^{4r-e}(J_0+n),\qquad e\ge6r+1.
 \label{sf:stage-response-budget}
\end{equation}
This budget concerns the individual old functions and their response
amplitudes. It does not assume a logarithm of the stage integral.

Assume that the sum of the Gaussian density ratio, the coarea and
Haar ratios, and the intrinsic interaction remainder is bounded by
\begin{equation}
 G_n\le C_g\theta^{-1/2}\ell^A
 \label{sf:stage-intrinsic-budget}
\end{equation}
on the tested box. Here \(A\) is fixed. For example,
\eqref{sf:gaussian-ratio} gives this form when
\(\epsilon_n\le C\theta^{-1/2}\ell^a\); then its contribution
has \(A=a+4r+2p\). Equations \eqref{sf:pivot-ratio},
\eqref{sf:haar-ratio} and \eqref{sf:bare-Wilson-remainder} give
the same form under the corresponding chart bounds. All these
estimates must hold on the full remaining switch domains.

\begin{proposition}[The logarithmic balance of the stages]
\label{prop:stage-accumulation}
Suppose each stage has \eqref{sf:exact-density}, with the old terms
split by \eqref{sf:two-differences}, and the preceding bounds hold.
There is a floor \(\theta_*\), depending only on the displayed
constants, \(r,p,A\) and \(J_0\), such that all the normalized
stage integrals are nonzero for \(\theta\ge\theta_*\). Their
logarithms \(C_n=\log(I_n/I_{0,n})\) satisfy
\begin{equation}
 \sum_{n=0}^{N-1}\|C_n\|_\infty
 \le C\left[\theta^{-1/2}\ell^{A+r}
                +\ell^{-1}+J_0\ell^{-r-1}\right].
 \label{sf:stage-total-bound}
\end{equation}
In particular each \(\|C_n\|_\infty\le1\), so the stated
incidence bound closes by induction whenever the next tested chart
admits the preceding functions on their declared switch domains.
The norm in \eqref{sf:stage-total-bound} is a supremum over those
domains, at fixed real spectators.
\end{proposition}

\begin{proof}
Increase the floor so that
\(q=(\Lambda-1)^{-1}\le2\ell^{-(6r+1)}\) and
\(K_sq\le1\). The exact exponent change at stage \(n\) is
bounded, by \eqref{sf:incidence-reset} and
\eqref{sf:response-difference}, by
\begin{equation}
 \delta_n=C_g\theta^{-1/2}\ell^A
       +4C_s\ell^{-2r-1}(J_0+n)
       +2C_a\ell^{4r-e}(J_0+n).
 \label{sf:individual-stage-budget}
\end{equation}
For \(n<N\), this tends to zero uniformly as
\(\theta\to\infty\). Choose the floor so that every
\(\delta_n\le1/8\) and \(4\delta_n\le1\).
\Cref{prop:stage-sector} proves existence and the bound
\(\|C_n\|\le4\delta_n\). This supplies the bound one required
at the next stage, rather than using the final sum as an inductive
assumption.

Sum \eqref{sf:individual-stage-budget}. The stored-field term is
at most
\[
 C\ell^{-2r-1}(NJ_0+N^2)
 \le C\big[J_0\ell^{-r-1}+\ell^{-1}\big].
\]
The response term is at most
\[
 C\ell^{4r-e}(NJ_0+N^2)
 \le C\big[J_0\ell^{5r-e}+\ell^{6r-e}\big],
\]
which has the same bound since \(e\ge6r+1\). The intrinsic
term sums to \(C\theta^{-1/2}\ell^{A+r}\). These are
\eqref{sf:stage-total-bound}. The reference factors remain in the
exact equality
\(\prod_nI_n=(\prod_nI_{0,n})\exp(\sum_nC_n)\).
They can be absorbed into vacuum constants only when they are
independent of the retained pair.
\end{proof}

The power \(6r+1\) pays for two stage counts and one block count:
\[
 N\,\Lambda^{-1}\,K_s\,N
       \le C\ell^{r-(6r+1)+4r+r}=C\ell^{-1}.
\]
This is the role of the enlarged stored-field domain. Its admission
is a condition on the actual nonlinear averaging and background
maps. Since \(g=\theta^{-1/2}\), a fixed power of \(\ell\)
times \(g\) still tends to zero. This arithmetic alone does not
prove that an asserted enlarged chart exists.

\subsubsection{Preservation of switches under integration}

For the induction just used, holomorphy is preserved by the integral
and the logarithm under a common integrable bound. Independence at
a zero switch requires a further condition on the actual kernel.
For a remaining stored datum \(a_v\), suppose that, when
\(\mu_v=0\), the integration domain and the complete normalized
reference kernel are independent of \(a_v\), and that the
non-reference integrand has the same independence. Then the
normalized ratio \(I/I_0\) and its logarithm are independent of
\(a_v\) there. The unnormalized integral has this independence
when its reference mass \(I_0\) does as well. If \(a_v=0\) removes
\(\mu_v\) from the integration domain, the normalized reference
kernel and the inserted integrand, the ratio and its logarithm are independent of that
switch. No smaller switch disc is
needed: dominated holomorphic integration and the fixed logarithm
branch apply at arbitrary points of the original disc.

This condition cannot be inferred from freezing the tests only in
the complex variable. For example, on \(0<b<1\), let
\(S_a=(0,a)\), \(0<a<1\), and use normalized Lebesgue measure
on \(S_a\). The function \(F(b)=b\) does not read \(a\), but
its expectation is \(a/2\). Thus a spectator-dependent test can
create dependence on that spectator even when the inserted
function has none. In the stage application, any such dependence
must either stay outside the integral as a retained factor or be
included in the declared switch structure of the kernel.

\subsubsection{A common envelope and the scope of the local result}

The preceding non-vanishing proof is relative to a positive
reference integral. A mass estimate for an upper integral of a
supremum is a different consequence. Fix the geometry, the number
\(m\) of free coordinates and its site count \(N_0\). Here is a
sufficient condition in these same coordinates. Suppose, uniformly
in all real read parameters,
\begin{equation}
 Q_{0,\sigma}\ge\gamma I,
 \qquad c_{0,\sigma}\le C_0^m,
 \qquad V_{0,\sigma}(b)\le B_0N_0+\tfrac14\gamma|b|^2,
 \qquad m\le c_0N_0.
 \label{sf:envelope-data}
\end{equation}
The complete unnormalized density then has the Borel envelope
\begin{equation}
 \sup_{\sigma,z}\left|
 \mathbf1_{S_\sigma}c_{0,\sigma}
 e^{-b^{\mathsf T}Q_{0,\sigma}b/2+V_{0,\sigma}(b)+H_\sigma(z,b)}
 \right|
 \le C_0^m e^{B_0N_0+\delta}e^{-\gamma|b|^2/4}.
 \label{sf:common-envelope}
\end{equation}
Its integral is at most \(\exp(CN_0+\delta)\), by the
Euclidean Gaussian integral. The supremum need not itself be
measurable: the displayed Borel function bounds its upper integral.
The common coordinates, the coercive bound and the determinant
bound in \eqref{sf:envelope-data} must hold before the supremum is
taken. Separate parameter-dependent changes of variables do not
establish this envelope.

The local conclusions require the exact stage chart, the
Gaussian marginal identity, the enlarged switch and displacement
domains, and the stated kernel condition at zero switches. For
simultaneous components, their interaction and the spatial filing
of logarithmic terms require the local Gaussian and connected
expansion estimates of \cref{ch:local-formation-program}. The
finite-component bound here supplies neither their decoupling
matrix estimates nor a common chart for an arbitrary history.

\begin{proposition}[Sharp tests can force a zero without switches]
\label{prop:stage-sharp-test-zero}
Let \(w_0\) be integrable on a real measure space, let
\(W=\int w_0\ne0\), and let \(x\ge0\) have an atomless
push-forward of \(|w_0|\,d\mu\). Then there are \(c>0\) and
\(\lambda\in[-\pi,\pi]\) such that
\par\noindent
\begin{minipage}{\linewidth}
\begin{equation}
 \int w_0e^{i\lambda\mathbf1_{\{x>c\}}}\,d\mu=0.
 \label{sf:phase-zero}
\end{equation}
\popQED\par
\end{minipage}
\end{proposition}

\begin{proof}
Set \(A(c)=\int_{x\le c}w_0\) and \(B(c)=W-A(c)\).
Atomlessness makes both continuous in \(c\), with
\(A(0)=0\), and \(A(c)\to W\) as \(c\to\infty\).
The continuous function \(|A(c)|-|B(c)|\) is negative at zero
and positive for sufficiently large \(c\). At a zero it has
\(|A(c)|=|B(c)|>0\), since \(W\ne0\). Choose
\(e^{i\lambda}=-A(c)/B(c)\). This gives
\eqref{sf:phase-zero}.
\end{proof}

For a free three-color strip coordinate,
\(x=|\mathcal A^{\rm s}(b)|\) is atomless under an integrable
Lebesgue density. Its norm is invariant under the adjoint action.
The phase in \eqref{sf:phase-zero} can therefore be split among
four local terms, each of modulus at most \(\pi/4\), without
changing analyticity in the background pair. The switch hypothesis
rules out precisely this uncontrolled dependence on a tested
stored field.

% END PROOF EXPANSION formation/stage_formation_candidate.tex


% BEGIN MODULE chapters/ordinary_source_formation_audit_v6.tex
\subsection{Source formation and volume bounds}
\label{ch:ordinary-source-formation-audit}

A source bound requires three estimates: formation of local
interactions, convergence of the polymer logarithm, and a bound on
the resulting integral.  A convergent polymer expansion supplies
the second estimate.  To apply it to a renormalisation construction one must
first verify its activity bounds for the interactions that the construction
actually forms.  The elementary implications below specify the ordinary source
estimates needed for that verification.

\subsubsection*{Site-dependent members and the summation over old levels}

Fix an odd integer \(L\ge3\).  Let \((\theta_j)_{0\le j\le n}\) satisfy
\begin{equation}
 \theta_j\ge Y,\qquad
 \lambda(n-j)-c_5\le\theta_j-\theta_n\le B(n-j),
 \qquad Y\ge2(c_5+1),
 \label{eq:ordinary-source-formation-OSF1}
\end{equation}
where \(\lambda>0\) and \(c_5\ge0\).  A pointed family at level \(j\) is a collection
\(G(X,P,z)\), where \(X\) is a localisation domain, \(z\in X\) is a base
point and \(P\) belongs to a specified complex field domain.  Its norm is
\begin{equation}
 \|G\|_{j,\kappa}
   =\sup_{X,z,P}e^{\kappa d_j(X)}|G(X,P,z)|.
 \label{eq:ordinary-source-formation-OSF2}
\end{equation}
Any companion families used to define its quadratic coefficient are included
in this supremum.  The field domains and restrictions are part of the data;
holomorphy does not assert that a function extends beyond these domains.

At every site \(z\), permit a different family \(G_{j,z}\).  Each such family
comes with its own quadratic coefficient \(b(G_{j,z})\).  The hypothesis on an
individual member is covariance of its \emph{pointed} total, so that subtracting
its coefficient times the tent action cancels the marginal quadratic term at
that base point.  No covariance relating the chosen members at different
sites is required.  This distinction permits a scalar source to be frozen at
the value belonging to each site.

Here is a precise quantitative formulation of the input used in this
argument.  Put \(\epsilon_{jn}=L^{-(n-j)}\), and let \(\rho_{jn}\) be the
ratio between the field radius required at level \(n\) and the radius on
which the member of level \(j\) is given, enlarged to be at least one.  Suppose
\(\rho_{jn}^2\le C_\rho(1+n-j)\).  Assume a linear localisation rule with the
following estimates, including summability over its rooted domains.
For one base point of level \(j\), a normalised regular member contributes at
most
\begin{equation}
 C\,\|G_{j,z}\|_{j,\kappa}\,
       \rho_{jn}^{\,3}\epsilon_{jn}^{\,4+\alpha}
 \qquad(0<\alpha<1)
 \label{eq:ordinary-source-formation-OSF3}
\end{equation}
to a window of level \(n\), after its vacuum and quadratic subtractions.
A member with its vacuum value alone subtracted contributes at most
\begin{equation}
 C\,\|F_{j,z}\|_{j,\kappa}\,
       \rho_{jn}^{\,2}\epsilon_{jn}^{\,4}.
 \label{eq:ordinary-source-formation-OSF4}
\end{equation}
The constants include the permitted loss of spatial decay and the count of
domains through the base point.  The rule must evaluate every member within
its stated field domain for every allowed decoupling and window parameter.
This domain inclusion is a separate hypothesis of \eqref{eq:ordinary-source-formation-OSF3}--\eqref{eq:ordinary-source-formation-OSF4}.

The factor \(\epsilon_{jn}^4\) has a direct dimensional meaning.  A window of
level \(n\) contains at most \(C\epsilon_{jn}^{-4}\) base points of level
\(j\).  Thus vacuum subtraction balances the number of base points in four
dimensions, whereas the quadratic subtraction must supply the additional
\(\epsilon_{jn}^{\alpha}\) to make old regular members summable.

\begin{proposition}[Uniform ordinary input]
Suppose the regular members obey
\[\sup_z\|G_{j,z}\|_{j,\kappa}\le G_\star.\]
Suppose defects and remainders
obey, respectively,
\begin{equation}
 \|D_j\|_{j,\kappa}\le L^{-(n+1-j)},\qquad
 \|H_j\|_{j,\kappa}\le C_H\theta_j^{-p},\qquad p>2.
 \label{eq:ordinary-source-formation-OSF5}
\end{equation}
Under \eqref{eq:ordinary-source-formation-OSF3}--\eqref{eq:ordinary-source-formation-OSF4}, their combined ordinary input at level \(n\) is bounded
by
\begin{equation}
 C G_\star
 +C\sum_{a=0}^{n}(1+a)L^{-(a+1)}
 +C C_H\sum_{a=0}^{n}(1+a)(\theta_n+\lambda a-c_5)^{-p}
 \le C G_\star+C+C C_H\theta_n^{\,2-p}.
 \label{eq:ordinary-source-formation-OSF6}
\end{equation}
The constant is independent of \(n\), the number of sites and the choice of
the member at each site.
\end{proposition}

\begin{proof}
The sum in \eqref{eq:ordinary-source-formation-OSF3} over the base points in a window.  The inverse volume factor
cancels \(\epsilon_{jn}^4\), leaving
\[
 C G_\star\sum_{a\ge0}(1+a)^{3/2}L^{-\alpha a}<\infty.
\]
This uses only the supremum of the member norms, so choosing a different
covariant member at each site has no effect on the bound.  The same
cancellation in \eqref{eq:ordinary-source-formation-OSF4} gives the second sum in \eqref{eq:ordinary-source-formation-OSF6} for defects.  That sum
is bounded by a convergent geometric series and its first derivative.

For remainders use \eqref{eq:ordinary-source-formation-OSF1} and put \(y=\theta_n-c_5\ge\theta_n/2\).
For \(a\ge0\), \(1+a\le C(y+\lambda a)\), since \(y\ge1\).
Comparison with the integral of the decreasing function
\((y+\lambda t)^{1-p}\) gives
\begin{equation}
 \sum_{a\ge0}(1+a)(y+\lambda a)^{-p}
 \le C\left(y^{1-p}+\int_0^\infty(y+\lambda t)^{1-p}\,dt\right)
 \le C y^{2-p}.
 \label{eq:ordinary-source-formation-OSF7}
\end{equation}
This proves the stated, deliberately nonoptimal, bound.
\end{proof}

For \(\operatorname{SU}(2)\), a bounded scalar action weight supplies another admissible
ordinary input.  If \(W\in SL(2,\mathbb C)\), then
\begin{equation}
 1-\tfrac12\operatorname{tr}W=\tfrac12\det(1-W).
 \label{eq:ordinary-source-formation-OSF8}
\end{equation}
On a field domain where each plaquette satisfies
\(|W_p-1|\le\alpha_n\eta_n^2\), its $\operatorname{tr}(F^2)$ is bounded by
\(C\alpha_n^2\eta_n^4\).  A local cube contains at most
\(C M^4\eta_n^{-4}\) plaquettes.  Filing each plaquette in a domain containing
its corners therefore gives a weight member of norm at most
\(C M^4\alpha_n^2\sup|u|\).  This proof uses the determinant identity and the
count of plaquettes; it does not use covariance of \(u\).  With
\(\alpha_n=O(\theta_n^{-1/2}(\log\theta_n)^q)\), any fixed bound on
\(\sup|u|\) is eventually admissible.

These summations do not prove \eqref{eq:ordinary-source-formation-OSF3}.  That estimate contains the
renormalisation assertion: covariance of the pointed Hessian, the correct
counterterm, regularity of the background within the window and uniform
membership after substitution must be proved for the actual formation rule.
They cannot be replaced by a bound on the full total or by holomorphy in the
source alone.

\subsubsection*{Exponential reserves and generated local perturbations}

Let a local history \(\omega\) have a positive measure \(\mu_\omega\) on
\(S_\omega\), a reference exponent \(\Phi_\omega(\xi;P)\), and a site count
\(N_\omega\).  Suppose its reserve volume \(V_\omega\) is additive over
the older histories swallowed by this history and satisfies
\begin{equation}
 V_\omega\ge A N_\omega,\qquad
 \sum_\omega\int_{S_\omega}
      e^{c_\star V_\omega}\sup_P|e^{\Phi_\omega(\xi;P)}|
             \,d\mu_\omega(\xi)\le\varepsilon,\qquad c_\star\ge1.
 \label{eq:ordinary-source-formation-OSF9}
\end{equation}
The supremum in \eqref{eq:ordinary-source-formation-OSF9} is assumed measurable, or is replaced by a measurable
majorant throughout.  If a holomorphic perturbation satisfies
\(|\Psi_\omega|\le fN_\omega\) and \(A\ge f\), then for bounded \(F\)
\begin{equation}
 \left|\sum_\omega\int F e^{\Phi_\omega+\Psi_\omega}\,d\mu_\omega\right|
 \le\varepsilon\|F\|_\infty.
 \label{eq:ordinary-source-formation-OSF10}
\end{equation}
Indeed, \(|e^{\Psi_\omega}|\le e^{V_\omega}\), so take absolute values
inside every integral and apply \eqref{eq:ordinary-source-formation-OSF9}.  If each integrand is holomorphic
in the parameter and all reference measures and integration sets are
parameter independent, a common integrable majorant on compact parameter
sets also proves holomorphy by Fubini and Morera.
This is an ordinary complex mass estimate.  It requires no derivative
with respect to the number of source zones.

Suppose a formation rule produces a new excess exponent in the explicit form
\begin{equation}
 \Psi_\omega^{\rm new}
   = A(\widehat c\,\upsilon1_X,U''(\xi))
       +V^{\rm col}(X^\sim;\mathfrak L)-V^{\rm col}(X^\sim;\mathfrak L^0)
       +\sum_R\delta_R[\psi_R],
 \label{eq:ordinary-source-formation-OSF11}
\end{equation}
where \(U''(\xi)\) is real, \(0\le\upsilon\le1\) and
\(\sup|\widehat c|\le\bar c\).  To apply \eqref{eq:ordinary-source-formation-OSF10} it is sufficient to prove
the actual formation bounds
\begin{equation}
 |A(\widehat c\,\upsilon1_X,U''(\xi))|\le C_W\bar c\,N_\omega^0,\qquad
 |V^{\rm col}(X^\sim;\mathfrak L)|
       +|V^{\rm col}(X^\sim;\mathfrak L^0)|
       +\sum_R\sup|\delta_R[\psi_R]|
          \le2c_{\rm int}(1+\bar c)N_\omega^0.
 \label{eq:ordinary-source-formation-OSF12}
\end{equation}
If \(\Psi_\omega=\Psi_\omega^{\rm new}
+\sum_{\omega'}\Psi_{\omega'}\) and
\(N_\omega=N_\omega^0+\sum_{\omega'}N_{\omega'}\), these inequalities give
\begin{equation}
 |\Psi_\omega|
 \le\max\{C_W\bar c+2c_{\rm int}(1+\bar c),f\}\,N_\omega.
 \label{eq:ordinary-source-formation-OSF13}
\end{equation}
The maximum, rather than a sum of the old and new constants, is what permits
unbounded nesting.  Preserving the carrier oscillation norm needed in the
next formation step is an additional estimate; it does not follow from
\eqref{eq:ordinary-source-formation-OSF13}.

There is a genuine reason to generate local perturbations through this
formation rule.  A bound on an exponent and on the oscillation of each
carrier does not bound arbitrary partial sums of carrier factors.  On a
one-point integration space, write
\(\Psi^\circ(\lambda)=-N\lambda\) and
\(\psi_1(\lambda;P)=N\lambda\), with a second, unused carrier if desired.
Then \(\Psi^\circ+\psi_1=0\), both carrier oscillations in \(P\) are zero,
and the full mass is unchanged for every \(\lambda\).  Interpolating only
\(\psi_1\), while leaving \(\Psi^\circ\) at its value at zero, instead gives
the mass \(e^{N\lambda}\).  No bound uniform in \(N\) survives on a fixed
complex disc.  This example refutes that partial-interpolation argument,
not a construction which defines its localized exponents recursively by
\eqref{eq:ordinary-source-formation-OSF11} and proves \eqref{eq:ordinary-source-formation-OSF12} on the whole complex data box.

%
%
%
\subsubsection*{Local data lines in the source induction}
\phantomsection\label{source:local-data-lines}

Use the fixed real step and its source radius \(r\).
Construct the constant-source data first. If the
radial increment bounds through level \(k\) sum to
at most \(1/3\), the running constant shift
satisfies \(|\zeta_k(\zeta)|\le4|\zeta|/3\).
The reserve \(2r\) puts the shifts with
\(|\zeta|<r\) in the real weight tube. The doubled
disc \(|\zeta|<2r\) instead uses the larger analytic
range of the increment maps. The frame bound
\(\theta_{i+1}\le\theta_i+C\) carries the upper
bound on the tube width at the next scale back to
\(C'\theta_i(\log\theta_i)^{-C_1}\) at scale \(i\).
Choose the tube constant smaller than the constant
in the analytic range. The fixed extra reserve
\(2r/3\) is then available at weak coupling.
Thus the previously formed shifts and remainders
are admitted when computing \(b_{k+1}(\zeta)\) on
the doubled disc. The summable increment bound
keeps the same \(1/3\) bound for the enlarged
list. The regular constant-source step supplies
the covariant member with this coefficient, before
the position-dependent source is introduced.

At level \(k+1\) let \(W\) be a region containing the complete
input reads of the output under consideration. For a
source \(w\) and a constant comparison value \(\zeta_0\)
in \(D(0,r)\), let \(\mathsf m\) be the source deviation
on \(W\), including the scaled Lipschitz bound at level
\(k+1\). Form the data line at \(\zeta_0\) by adding
\(t\) times the difference from \(w\). Move the site
values and the entries with their complete one-layer
enlargements inside \(W\). For a boundary sum make this
choice at each of its summands. Keep the reference entries
outside \(W\). The linear consistency rules are checked
at these same indices.

The inherited slope bounds put the line in the enlarged
data box for \(|t|\mathsf m<r\). In detail, the site
source has modulus at most \(2r\). If
\(\mathsf d_{2r}[b_i]\) is the bound for
\((b_i(\zeta)-b_i(0))/\zeta\) on \(D(0,2r)\), the
increment has Lipschitz constant at most
\(2\mathsf d_{2r}[b_i]\) on \(D(0,r)\), by Cauchy's
estimate for \(b_i-b_i(0)\) on the doubled disc.
Thus on the line its deviation is at most
\(3r\mathsf d_{2r}[b_i]\). The sum of the
\(\mathsf d_{2r}[b_i]\) is at most \(1/3\), so the
site source and the accumulated shifts fit the
\(3r\) bound. For an earlier defect at level \(i\),
the gradient scale is smaller by \(L^{-(k+1-i)}\).
The product \(|t|\mathsf m/r\) is at most one.
The other rooted entries use their bound at the
constant source and \(|t|\) times their inherited
slope, on the same carrier.

Define the excess exponents on this line by applying
the formation operations to the line. At the bare
level there are no excess exponents. A region for
the new grid is a region for the preceding grid, the
new source class is contained in the preceding one,
and the deviation at the earlier scale is at most
\(\mathsf m\). Thus the earlier line is available on
the same disc. The still-standing exponents are
carried from it. At a new operation apply
\eqref{eq:ordinary-source-formation-OSF11} and
\eqref{eq:ordinary-source-formation-OSF12} to the
data on the line. The fresh count and inherited
counts add, so the maximum in
\eqref{eq:ordinary-source-formation-OSF13} keeps the
same total bound at every depth. This defines the
line also when \(W\) cuts a large-field component.
At \(t=0\) it gives the constant-source output.
When the disc contains \(t=1\), the line there
gives the sourced output on complete reads in \(W\),
by locality of the actual operations and of the
lines at the inherited operations.

%
%
%
The carrier oscillation has its own bound on the
data box. At a new island the common quadratic
factors contribute zero to the source difference.
The remaining factors are differences of the
formed interaction pieces, with the whole island
carrier \(C^\sim\). Its length after exempting
the island hull is zero. Normalize the sum of
oscillations by \(\ell_t^{2\mathfrak p_0}\), with
the length weight of the carrier. The local
pieces on \(C^\sim\) cost \(C(1+\bar c)\check R_t^4\)
before normalization, and the source-dependent
Wilson pieces have a uniform bound. These contributions
give a fixed \(b_V^\times(\bar c)\). The difference
of an inherited factor is evaluated by the kept
response at fixed stored variables. Subtract its
value at the common real pair before taking the
response difference. The weighted sum of these
oscillations contributes at most
\(C_V^\times\vartheta(\theta_t)s\), where \(s\) is
the admitted normalized bound of the old carriers
and
\[
 \vartheta(\theta_t)
 = (A_1/C_0)(\log\theta_t)^{-(q-\mathfrak p_0)} .
\]
This estimate sums over the individual pieces and
therefore also bounds their grouped oscillations.
Standing factors retain their carriers and restrict
their old functions.

Choose \(c_\times\ge2b_V^\times(\bar c)\).
On the enlarged data line admit \(s\le4c_\times\).
With \(8C_V^\times\vartheta(\theta_t)\le1\), the new
normalized bound is at most \(c_\times\). Thus the
new carriers have the same room on the next data
line. This argument concerns the oscillation at
fixed stored variables. For the full source excess
use the absolute estimate of
\eqref{hi:actual-source-recursion}, including
the parts constant in the pair.

For a holomorphic output bounded by \(B\) on the
line disc, Schwarz applied after subtracting its
value at zero bounds the difference of the endpoint
values by \(2B\mathsf m/r\). If \(\mathsf m\ge r\),
the separate bounds at the two endpoints give the
same inequality. If \(\mathsf m=0\), the two reads
coincide. For carrier oscillations first subtract the
common reference value in each factor. Suppose their
sum of suprema on one carrier is bounded by \(B\)
throughout the line. For any finite collection choose
evaluation points and phases to make the endpoint
differences nonnegative and arbitrarily close to
their individual suprema. The corresponding phased
sum is a single holomorphic function bounded by \(B\)
on the line. Apply the same Schwarz estimate to it
and take the suprema and then the increasing sums.
This preserves the carrier oscillation slope separately
from the slope of the total exponent.

For the regular pointed coefficient the defect needs
the bound with unit slope and a spare unit of spatial
decay. The Gaussian normalizer is common to the line
and the constant source. After removing it, the
cluster contribution and the correction for the real
weight relative to the frozen one have combined bound
\[
 \left(C_V\vartheta_k^\sharp+
       C_{\mathbf M}\frac{\mathsf t_k}{\theta_k}\right)
                   e^{-(\kappa+1)d_{k+1}(X)} .
\]
Here \(\vartheta_k^\sharp\) includes the small
fluctuation-radius ratio, and \(\mathsf t_k\) is the
real weight tube radius. The
\hyperlink{source:ordinary-absolute-line}{ordinary
absolute estimate} gives this bound on the direct
regular windows. For an interior pointing the hulls
of the input vertices and the complete Gaussian
guards are in the small-field region, so the bracket
uses its ordinary legs without a stored exterior
mean. The normalizer and the free ratio term use
the same real reference in both members of the
comparison. The constant source chosen first at
each base supplies its own pointed subtraction.
The local line uses the direct estimates on its
complete bases. Any other origin entering the regular
rule must satisfy its own carrier and admission
bounds on this line. Taking the parenthesis at
most \(1/2\) gives the unit slope by the preceding
Schwarz estimate. Compare with the constant
\(\zeta_0=w(p_z)\) at the pointing \(z\). The spare
unit of decay pays for the diameter of the source read,
leaving the factor \(L^{-(k+1-i)}\) in a defect of
level \(i\) when it is read at level \(k+1\). The
uniform box bound supplies the output slope. The
input slopes only supply admission to that box.

%
%
Form the torus regular output of \(\mathbf T\) before choosing
its new extension. At each pointing \(z\) choose the Euclidean
member at the constant value \(w(p_z)\), with marginal coefficient
\(b_{k+1}(w(p_z))\). The pointed slope bounds the difference between
the actual regular output and this member, admitting the defect
of the extension. For estimates along the complex data line,
interpolate the extension members and their coefficients
affinely between the constant and sourced endpoints. On
\(|t|\mathsf m<r\) the member norm is at most
\(E_0+|t|\mathsf m/r\le E_0+1\), and the summed coefficient
deviation and site source obey the same \(3r\) scalar bound.
For the extension defect use the two regular outputs with the
common normalizer removed and the affine change of the member.
The first two terms are bounded by \(\tfrac12 e^{-\kappa d}\)
each and the affine change by \(e^{-\kappa d}\). The defect
therefore fits the enlarged box on the entire line. Applying
\(\mathbf R\) there returns the slope bounds for the new boundary
and remainder terms by the same Schwarz estimate. Carried terms
keep their previous slopes on their native reads.

For the actual sourced endpoint set
\[
 c_{k+1}(p;w)=\varphi_{k+1}(p)\sum_y h^{(k+1)}_y(p)
                                  b_{k+1}(w(p_y)).
\]
Use the real-center update
\(x^\natural_{k+1}=x^\natural_k-b_{k+1}(0)\varphi_{k+1}\)
for the position-dependent weight, and update the running shift
by \(\zeta_{k+1}=\zeta_k+c_{k+1}-b_{k+1}(0)\varphi_{k+1}\).
Then at every plaquette
\[
 x^\natural_k-\zeta_k=x^\natural_{k+1}-\zeta_{k+1}+c_{k+1}.
\]
Thus the pointed subtraction \(-A(c_{k+1},U)\) pays exactly for
the change of the Wilson coefficient. The profiles, tests and
integration kernels come from the frozen step throughout. On
the admitted formation domains the exactness of \(\mathbf R\mathbf T\)
with this extension gives the sourced output in the resummed
standing classes with their retained variables, on the same
source class for the next step.

\subsubsection*{What a pointwise polymer logarithm proves}

Consider a finite hard-core polymer gas \(\Xi_y\), depending on an exterior
field \(y\).  Suppose its activities \(h_y(C)\) satisfy a uniform
Koteck\'y--Preiss criterion, including absolute convergence of the logarithmic
series and the bound
\begin{equation}
 |\log\Xi_y|\le\sum_C|h_y(C)|e^{a(C)}\le c\,N.
 \label{eq:ordinary-source-formation-OSF14}
\end{equation}
If the activities are real for real \(y\), then \(\Xi_y>0\).
One proof is that the absolutely convergent logarithmic series is real.
Equivalently, replace every activity by \(t h_y(C)\), \(0\le t\le1\).
The same criterion applies, so \(\Xi_y(t)\) is real, continuous and nonzero.
Its value at zero is one, proving positivity at one.  Individual activities
need not be positive.  Moreover \(e^{-cN}\le\Xi_y\le e^{cN}\).

If a final representation really has the form
\(Z=\int q(y)\Xi_y\,d\nu(y)\), with \(q\ge0\), then
\begin{equation}
 e^{-cN}\int q\,d\nu\le Z\le e^{cN}\int q\,d\nu.
 \label{eq:ordinary-source-formation-OSF15}
\end{equation}
Thus even a pointwise polymer logarithm requires control of the remaining
integral.  It does not turn a restricted Haar integral into a probability
integral.  For example, take the empty gas and let \(q\) be the indicator of
a Haar set of measure \(e^{-N^2}\).  The polymer criterion is automatic,
but a lower bound \(e^{-CN}\) for \(Z\) is false.

In a representation that sums over standing labels, the useful upper
estimate instead has the form
\begin{equation}
 \sum_{\mathcal S}
       \big|e^{\mathcal E_{\mathcal S}(w)}
                  \tau_{\mathcal S}^{\,w}\big|(\mathcal V_n)
       e^{W(\mathcal S)}
 \le e^{cN_n}.
 \label{eq:ordinary-source-formation-OSF16}
\end{equation}
Here the absolute value is total variation, \(W\) pays for missing vacuum
contributions and the measures inside the retained integral operations are
also majorised before they are integrated.  A pointwise bound for one gas
does not supply this sum: arbitrarily many positive labels can have the
same individually bounded gas.  The counting of labels and the exponential
reserve in their operations are essential.

For the lower bound at zero source, positivity permits retaining the
all-small-field label.  Suppose its density, after vacuum normalisation,
is at least
\(\chi(y)\exp[-x A(U(y))-C N_n]\).  Assume the ball in which every one of
the \(4N_n\) link variables has distance less than \(s\) from the identity
is contained in \(\{\chi=1\}\), and that
\(A(U(y))\le C s^2N_n\) there.  Normalized Haar measure on \(\operatorname{SU}(2)\) assigns
this link ball mass at least \((c s^3)^{4N_n}\).  It follows directly that
\[
 e^{-E}Z(0)\ge
  \exp[-(C+4\log(1/(c s^3))+Cxs^2)N_n].
\]
Taking \(s=c_0x^{-1/2}\) gives
\(e^{-E}Z(0)\ge\exp[-(C+6\log x)N_n]\).
The terminal background estimates needed here are proved in
\hyperref[oh:terminal-link-ball]{the calculation on a product of link balls}.
No zero-free statement for the sourced full partition function is needed.


% END MODULE chapters/ordinary_source_formation_audit_v6.tex

% BEGIN MODULE chapters/weighted_extension_audit_v2.tex
\subsection{Weighted domains and scale forgetting}
\label{ch:weighted-extension-audit}

A multiscale argument needs several analyticity statements with different meanings. Holomorphy of a finite-dimensional partition function is the weakest of them. Holomorphy of a renormalization step on a class of effective interactions is stronger. Geometric forgetting of old interactions is stronger still. The latter can follow from an analytic extension with appropriate domains and bounds, but the weights in that extension must themselves be justified.

There is an additional distinction for large-field regions. A small probability of retaining an old region can appear as a small factor in the bound on an analytic function. It does not in general permit the analytic input radius to grow by the reciprocal of that probability. This section defines a suitable input space and states the local inequalities needed for continuum reconstruction. Their validity for the Yang--Mills step remains to be proved.

\subsubsection{The ordinary input space}

Fix a cutoff, a reference sequence of inverse couplings $\theta_0,\ldots,\theta_k$, a complete set of labels and tests, and hereditary complex domains $\mathcal U_j(X;r_{0,j},r_{1,j})$. A point of such a domain is a pair $(\mathbf U,\mathbf J)$; the current $\mathbf J$ is an independent complex variable. A regular background gives the special pair $(U,J(U))$. An identity verified at these special pairs is not automatically an identity at all pairs in the complex domain.

One convenient Banach space of pointed local families has the norm

\begin{gather*}
\FitMath{\|E\|_{j,\kappa}
 =\sup_{X,z}\,e^{\kappa d_j(X)}
    \sup_{\mathbf p\in\mathcal U_j(X;r_{0,j},r_{1,j})}
       |E(X,\mathbf p,z)|.}
\end{gather*}

When families have infinite-lattice companions, include their norms in the maximum. Gauge invariance, the prescribed translation covariance, and compatibility with restriction are closed linear conditions. Holomorphy is preserved under convergence in this supremum norm on the declared domains. These facts specify a closed space of analytic families once the domains themselves have been fixed. The existence and uniform size of those domains remain analytic inputs.

A regular Euclidean member comprises a family $G$, its companion, and its scalar quadratic coefficient $b(G)$. It has covariant pointed totals on the prescribed regular backgrounds and carries its own counterterm. A spatially varying input can assign a different such member $G^{(j)}_z$ to every base point $z$. This does not require the array $z\mapsto G^{(j)}_z$ to be translation invariant. Its norm is $\sup_z\|G^{(j)}_z\|_{j,\kappa}$. Each assigned member must retain its own coefficient and the required pointed covariance; covariance of the sum of all base points is a weaker condition.

The complete input also contains a bare plaquette weight $u$, defect families $D^{(j)}$, remainders $H^{(j)}$, boundary terms $B$, and excess exponents $\Psi$ in stored large-field operations. A useful product box has inequalities of the following form:

\begin{gather*}
\FitMath{\begin{split}
 &\|u\|_\infty+
   \sum_{j\le k}\sup_z|b(G^{(j)}_z)-b_j^{\mathrm{ctr}}|\le C_u,
 \qquad \sup_z\|G^{(j)}_z\|_{j,\kappa}\le E_0,\\
 &\|D^{(j)}\|_{j,\kappa}\le L^{-(m-j)},
 \qquad\|H^{(j)}\|_{j,\kappa}\le C_H\theta_j^{-\kappa_0/2},
 \qquad m\in\{k,k+1\}.
\end{split}}
\end{gather*}

Here $b_j^{\mathrm{ctr}}$ denotes the coefficient of the chosen real centre. Boundary terms are measured against specified positive majorants $\overline B_{j,h}(X)$:

\begin{gather*}
\FitMath{\|B\|_{\mathrm{bdry}}
  =\sup_{j,h,X}|B_{j,h}(X)|/\overline B_{j,h}(X).}
\end{gather*}

For a stored history $\omega$ let $N(\omega)$ count its stored sites, each at its own scale. Excess exponents require both a bound on $|\Psi_\omega|/N(\omega)$ and a carrier-by-carrier oscillation bound. If $\Psi_\omega=\Psi_\omega^0+\sum_f\psi_f$, the latter can be expressed as

\begin{gather*}
\FitMath{\sup_{\omega,R}\overline M_{\omega}(R)^{-1}
       \sum_{f:R_f=R}\operatorname{osc}\psi_f<\infty.}
\end{gather*}

The carrier majorant and the history size are part of the definition, not consequences of holomorphy.

A source trajectory fits this format by freezing the source separately at each base point. Schematically,

\begin{gather*}
\FitMath{G^{(j)}_z=E^{(j),\zeta}\big|_{\zeta=w(p_z)},} \\[3pt]
\FitMath{D^{(j)}(X,\cdot,z)
   =E^{(j),w}(X,\cdot,z)-E^{(j),w(p_z)}(X,\cdot,z).}
\end{gather*}

For a Lipschitz source, a local freezing estimate can give
$\|D^{(j)}\|\le C L^{-(K-j)}\|w\|_{\mathrm{Lip}}$.
This smallness is a property of the source trajectory. An arbitrary defect in an ordinary box has only the size stipulated by that box.

This organization resolves a possible misconception. Global noncovariance does not by itself invalidate a construction with sources: Euclidean members can be frozen site by site. The remaining task is to prove the asserted bounds for the full array, including its noncovariant defects, its variable Wilson weight, and its stored-history operations. Scalar marginal cancellation for a globally covariant list alone does not complete that task.

\subsubsection{Three kinds of weights}

Let $e_j$ measure a difference of regular families and let $r_j$ measure a difference of remainders. The regular irrelevance weights have the form

\begin{gather*}
\FitMath{\omega_\nu\le C\nu^4L^{-\sigma(\nu-1)},} \\[3pt]
\FitMath{\sigma>0.}
\end{gather*}

They represent decay after the vacuum and the applicable marginal quadratic terms have been removed. In contrast, an unrenormalized remainder can enter with a squared ratio of analyticity radii,

\begin{gather*}
\FitMath{w_{n,j}=\left(\frac{r_{0,n}}{r_{0,j}}\right)^2
 \asymp\frac{\theta_j}{\theta_n}
    \left(\frac{\log\theta_n}{\log\theta_j}\right)^{2q}.}
\end{gather*}

These weights account for the small size of the remainder. Geometric forgetting and any further symmetry cancellation require separate estimates.

A typed tangent list may also have a marginal amplitude $z_j$ and an increment $b_j$. Define

\begin{gather*}
\FitMath{z_n^\sharp=\sum_{j\le n}L^{-(n-j)}(z_j+z_{j-1}),} \\[3pt]
\FitMath{z_{-1}=0,}
\end{gather*}

and a load

\begin{gather*}
\FitMath{p_n(\delta)=\sum_{j\le n}\omega_{n+1-j}e_j
  +\sum_{j\in\mathcal J_n}w_{n,j}r_j
  +\gamma_n\left(z_n^\sharp+
          \sum_{j\le n}L^{-c(n-j)}b_j\right),} \\[3pt]
\FitMath{\gamma_n\asymp\theta_n^{-1}(\log\theta_n)^C.}
\end{gather*}

The set $\mathcal J_n$ consists of remainder slots still present at level $n$ in the chosen filing convention. If a remainder is re-filed into a regular family, it must be removed from this set and counted in that family instead. A recent-level window can have a separate term without geometric age decay. This bookkeeping cannot be suppressed by changing notation.

\phantomsection\label{notdef:aged-load}
The survival age of a large-field zone has a different role. For $0<q_a<1$ the desired integrated bound involves

\begin{gather*}
\FitMath{\Sigma_k(\delta)=\sum_{n\le k}(k-n+1)
        q_a^{(k-n-1)_+}p_n(\delta).}
\end{gather*}

The factor $k-n+1$ counts the steps at which the zone can be used. The survival factor $q_a^{(k-n-1)_+}$ belongs to an output estimate on the relevant histories. It need not, and generally cannot, be put into a single enlarged analytic input radius.

\subsubsection{What weighted holomorphy actually proves}

\begin{proposition}[A segment estimate]
Let $X,Y$ be complex Banach spaces, let $C\subset X$ be convex, and let $p$ be a continuous seminorm. Suppose $F$ is holomorphic on

\begin{gather*}
\FitMath{\mathcal T=\{a+h:a\in C,\,p(h)<\epsilon\},} \\[3pt]
\FitMath{\epsilon>0,}
\end{gather*}

and, for every $a\in C$ and $p(h)<\epsilon$, one has
$\|F(a+h)-F(a)\|_Y\le B$. Then, for $a,b\in C$,

\begin{gather*}
\FitMath{\|F(b)-F(a)\|_Y\le (B/\epsilon)p(b-a).}\qedhere
\end{gather*}

\end{proposition}
\begin{proof}
Put $h=b-a$. At every $a_t=a+th$, $0\le t\le1$, the map
$z\mapsto F(a_t+zh)-F(a_t)$ is holomorphic on
$|z|<\epsilon/p(h)$ and has norm at most $B$. The Banach-valued Cauchy formula gives
$\|DF(a_t)h\|\le Bp(h)/\epsilon$. Integrating over $t$ proves the inequality. If $p(h)=0$, the line is entire and bounded in oscillation. Every bounded linear functional on $Y$ makes it a bounded scalar entire function, so Liouville's theorem and separation by the dual imply that it is constant.
\end{proof}

The seminorm in this proposition is an assumption. Applying it with
$p(h)=\sum_j\omega_{k+1-j}\|h_j\|$ produces the displayed weights because they already occur in the domain. It does not derive them from an ordinary product box.

A differential bound is a weaker sufficient input. If
$\|DF(a)h\|\le C p(h)$ holds on the segment joining two actual trajectories, integrating this inequality proves the same comparison without constructing a full weighted tube around every off-trajectory point. Thus a derivative estimate along these segments suffices.

\begin{proposition}[A marked slot]
Suppose $F(a,m)$ is jointly holomorphic when
$p(a-a_0)<r_a$ and $\|m\|<r_m$, is bounded there by $B$, and satisfies $F(a,0)=0$. For $p(a-a_0)\le r_a/2$ and $\|m\|\le r_m/2$,

\begin{gather*}
\FitMath{\|D_aF(a,m)h\|\le
       \frac{CB}{r_a r_m}\,p(h)\|m\|.}\qedhere
\end{gather*}

\end{proposition}
\begin{proof}
Cauchy's formula first in the complex line $a+zh$, and then in the line $tm$, bounds the mixed derivative by
$CBp(h)\|m\|/(r_a r_m)$ on the reduced domains. Integrate that derivative from $m=0$ to $m$. The derivative in $a$ vanishes at $m=0$ because $F(a,0)$ is identically zero. The argument uses two complex lines and introduces no factor equal to the dimension of either input space.
\end{proof}

This explains how a small marked-slot size can multiply the sensitivity to other data. It does not establish the required radii, the marked-slot vanishing identity, or a uniform bound on the enlarged domain.

\subsubsection{Exact examples separating the assertions}

\begin{proposition}[Ordinary analyticity does not imply forgetting]
Fix $0<q<1$ and any positive numbers $s_k$. There are entire scalar functions $F_k$ bounded by $s_k$ on the unit polydisc of $\mathbb C^{k+1}$ for which no constant $C$ independent of $k$ satisfies

\begin{gather*}
\FitMath{|F_k(x)-F_k(y)|\le Cs_k
      \sum_{j=0}^k q^{k-j}|x_j-y_j|}
\end{gather*}

on the polydisc of radius $1/4$.
\end{proposition}
\begin{proof}
Take $F_k(x)=s_kx_0$, $x=0$, and $y=(1/4,0,\ldots,0)$. The inequality would give $1\le Cq^k$. This is impossible as $k\to\infty$. The same example is holomorphic on every weighted tube, but its supremum on a tube of fixed weighted radius grows like $s_kq^{-k}$. Thus the missing issue can be a uniform bound, even when the extension itself is entire.
\end{proof}

This example is a logical separation, not a model satisfying the gauge-theory power-counting identities. It shows exactly why analyticity and a small one-step factor alone do not establish those identities.

\begin{proposition}[Rarity does not enlarge the complex radius]
Let $0<q<1$, let $a$ be a nonnegative integer, and let an event have probability $q^a$. For $\delta>0$ its excess generating function can be

\begin{gather*}
\FitMath{R_a(\lambda)=q^a(e^{\lambda\delta}-1).}
\end{gather*}

On $|\lambda|\delta<c$ it has a bound $C_cq^a$, and $R_a'(0)=q^a\delta$. There is no age-independent bound $C_cq^a$ on the larger disc $|\lambda|q^a\delta<c$.
\end{proposition}
\begin{proof}
The first bound follows from $|e^z-1|\le e^{|z|}+1$. At the positive point $\lambda=cq^{-a}/(2\delta)$ the ratio $R_a(\lambda)/q^a$ equals $e^{cq^{-a}/2}-1$, which tends to infinity. The event model is a genuine two-point probability space, and its generating function is entire.
\end{proof}

The correct large-field strategy is therefore a bounded per-zone radius together with a rare-history factor in the bound. A demand for a bounded extension on a single tube measured only by the aged load would be stronger than what the argument needs and can fail even in this elementary example.

There is a separate reason for the Euclidean-member and marginal-slot conditions. Let

\begin{gather*}
\FitMath{H=\begin{pmatrix}0&1\\1&0\end{pmatrix},} \\[3pt]
\FitMath{M_\varepsilon=I+\varepsilon H.}
\end{gather*}

For real $|\varepsilon|<1$, take independent centered Gaussian vectors $X_1,\ldots,X_m$ with precision matrix $M_\varepsilon$. The normalized block average $m^{-1/2}\sum_iX_i$ has exactly the same covariance $M_\varepsilon^{-1}$, hence exactly the same precision matrix. The traceless quadratic perturbation $\varepsilon x_1x_2$ survives this blocking unchanged, although its scalar trace coefficient is zero. The Gaussian integrals are holomorphic and uniformly coercive on each fixed complex disc $|\varepsilon|<r<1$.

This example explains why a scalar counterterm cannot remove an arbitrary anisotropic quadratic perturbation. It does not contradict the use of a scalar counterterm in the Euclidean class. It also explains why a typed extension must supply all of its prescribed tensor counter-insertions before claiming that the remaining part is irrelevant.

Finally, identities on backgrounds do not determine an arbitrary extension to independent currents. The polynomial $P(z,j)=j-z^2$ vanishes on every pair $(z,z^2)$ and is nonzero elsewhere. An analogous cancellation known only on a graph of background pairs needs an additional argument before it can be used throughout a domain where the current is independent.

\subsubsection{Two local inequalities to establish}

The first test concerns the small-field potentials before exponentiation. Let $\dot V_{k,j}$ be the variation of the actual localized fluctuation potential caused by one regular family at level $j$. Include its own scalar counterterm, or the complete tensor counter-insertions for a typed family. A suitable test is
\begin{equation}
 \|\dot V_{k,j}[\delta E]\|_{\mathrm{pot}}
 \le C\varkappa_k\omega_{k+1-j}\|\delta E\|_{j,\kappa}.
 \label{eq:wext-potential-test}
\end{equation}
The supremum in the potential norm must range over the complex pairs and decoupling parameters actually used inside the fluctuation integral, not just real backgrounds. One possible norm is the least $M$ such that

\begin{gather*}
\FitMath{|\dot V(Y;\mathbf p,B)|\le
 M e^{-\lambda d_k(Y)}
   \left(1+M_0^{-4}\sum_{b\subset Y}|B(b)|^2\right).}
\end{gather*}

Here $M_0$ is the fixed block-size parameter. The complex field domains, the allowed fluctuation cut, and the enlarged decoupling discs must be specified. Defects and remainders have their own test with squared-radius weights in place of $\omega$. A variable Wilson weight is a marginal slot; it does not receive an irrelevance weight merely by being written as a local family.

For ordinary site-frozen Euclidean members, a regular-output dependence estimate already has this intended structure after the cluster expansion. The additional typed claim requires the same cancellation at the complex pairs in the integrand, with the full tensor removed. A background extraction formula alone is insufficient for that extension. Proving \eqref{eq:wext-potential-test} for these typed inputs is a concrete local task.

The second test concerns a single zone in a large-field operation. Fix a zone $n$ and a step $k'$ at which it is used, where $n\le k'\le k$. Let $\lambda_{n,k'}$ multiply only that use of a direction of load $p_n(\delta)$. Isolate the normally convergent subseries of histories affected by this dial, and denote the resulting output by $R_{k;n,k'}(\lambda)$. The required axis estimate has the form
\begin{equation}
 \sup_{|\lambda|p_n(\delta)<c}
   \|R_{k;n,k'}(\lambda)\|_{k,\kappa}
 \le C\epsilon_k q_a^{(k-n-1)_+}.
 \label{eq:wext-zone-test}
\end{equation}
The constants $c,C$ are independent of the levels, the history, and the positions of the sources. The source is allowed to be complex on the declared source ball. The factor $\epsilon_k$ is the small mass assigned to the output under consideration. For an unintegrated boundary term it must be replaced by that term's actual majorant; its residual dependence is subsequently paid under the relevant mass, rather than being assigned a false termwise survival bound.

Cauchy's formula applied to \eqref{eq:wext-zone-test} gives

\begin{gather*}
\FitMath{\|\partial_\lambda R_{k;n,k'}(0)\|_{k,\kappa}
 \le (C/c)\epsilon_k q_a^{(k-n-1)_+}p_n(\delta).}
\end{gather*}

Summing over the $k-n+1$ uses of that zone, and then over $n$, gives $(C/c)\epsilon_k\Sigma_k(\delta)$. This proves the differentiated age estimate from a precise bounded complex-mass statement. The proof needs no age-dependent increase of the analytic radius.

A complete verification must establish \eqref{eq:wext-zone-test} under the moduli and suprema used by the operation, including its reserve weights. Undifferentiated rarity at the unperturbed source alone does not establish it: the dial changes the integrand on the rare histories. Conversely, a full off-trajectory analytic tube is not necessary if these axis estimates and the needed first and mixed derivatives have been proved directly on the actual trajectories.

\subsubsection{What remains conditional and what can proceed separately}

The ordinary product-box theorem, a weighted dependence estimate for its regular output, and the finer marked and zone-wise statements are different inputs. Coefficient readout uses the physical radii of the chain, as proved in \cref{ch:extraction-radius-normalization}: fixed radii give a uniform norm, and running radii carry the inverse-square factor already recorded by the radius weights. Near-zero core inclusion is supplied by the stated tower bounds. The remaining finer assertions concern actual marked and typed evaluations in complex windows, with their proper radius normalization. When a linearized class includes these properties among its hypotheses, contraction rows derived from that class remain conditional on the class hypothesis. Such a deduction does not turn an ordinary complex-box bound into the weighted or zone-wise estimates.

For a prescribed whole-sequence continuum limit, the smallest useful replacement is an actual comparison of the two tuned trajectories: a step inequality, domination of the observable evaluations by the comparison distance, and a terminal-integral inequality with summable defects. These three estimates can be proved directly. Establishing a large weighted analytic class is one route to them, not a logical necessity.

Finite-cutoff source identities do not require this finer statement. Fixed-spacing DLR compactness and closedness of reflection positivity do not require it either. Subsequential separated-correlation compactness follows from its own cutoff- and volume-uniform source bounds. Those bounds remain substantial analytic requirements; their volume-free proof can use undifferentiated age-resolved rarity without proving the differentiated zone estimates above. Similarly, tuning at one fixed cutoff can follow from its own holomorphic slope and target-crossing estimates without establishing a comparison of consecutive cutoffs.

% Keep the closing statement and argument together.
\clearpage
The proposed summable matched-cutoff comparison and the resulting whole-sequence tuned convergence use the stronger comparison route. The typed tangent contractions, scale-uniform memory estimates, and differentiated anisotropy method for rotations also require their corresponding finer inputs. Estimates for several insertions at nonzero source require joint bounds beyond a single first derivative at zero. An independent vanishing-flux theorem would provide a different route to rotations and should be assessed on its own hypotheses.

Thus the unresolved point is a specific family of uniform estimates on the declared inputs, not a general objection to sources or a proof that the construction fails. The potential test and the one-zone complex-mass test locate two places where a proof can be supplied, sharpened, or tested by an explicit counterexample.


% END MODULE chapters/weighted_extension_audit_v2.tex

% BEGIN MODULE chapters/zone_reserve_test_v3.tex
\subsection{A reserve estimate for one large-field zone}
\label{ch:zone-reserve-test}

The one-zone estimate must bound a complexly perturbed integral;
rarity at a real centre alone does not suffice. This section proves a reserve lemma, verifies the direct Wilson-weight channel on regular zones, and tests whether carrier oscillation and a prepaid end mark supply the remaining perturbation bound. The result is partial. The missing statement can be expressed as a norm bound for an explicit local evaluation map. The reserve argument alone does not prove that statement for every formed copy of a direction.

\subsubsection{The integral before taking a logarithm}

Keep the reference sequence, tests, history labels, changes of variables and reference kernels fixed. A local operation is then represented on a disjoint union of stored-variable spaces by

\begin{gather*}
\FitMath{I(\lambda,p)=\sum_\omega\int_{S_\omega}
       e^{\Phi_\omega(\lambda,\xi;p)}\,\mu_\omega(d\xi),}
\end{gather*}

where each $\mu_\omega$ is positive and independent of $\lambda$. A harmless finite frame average can be included in these measures. The parameter $p$ ranges over the external complex-pair domain. A single zone-use parameter has an index $(n,k')$: it multiplies the chosen direction where the stored zone $n$ is used at step $k'$. Only histories containing that use are affected.

Before re-exponentiation, the substitution of an input list by a line of lists is affine in the exponent when its evaluation maps have been frozen. Its perturbation has the concrete form
\begin{equation}
 L_{\omega,n,k'}[\delta](\xi,p)
   =\sum_{e\in\mathcal E(\omega,n,k')}\epsilon_e\,
         \delta F_e\bigl(\kappa_{\omega,e}(\xi,p)\bigr)
       +\delta c_\omega(\xi).
 \label{eq:zr-local-evaluation-map}
\end{equation}
The index $e$ records the entry, its counter-insertions, its carrier, and its particular occurrence. The signs or fixed coefficients are $\epsilon_e$. The formed argument $\kappa_{\omega,e}$ must lie in that entry's original analytic domain. The last term records scalar parts, including the values subtracted at the trivial field and their contributions to the vacuum bookkeeping. Constants cannot be omitted from this map merely because the observable ultimately uses a connected derivative.

For this affine use one has
$\Phi_\omega(\lambda)=\Phi_\omega(0)+\lambda L_{\omega,n,k'}[\delta]$.
After local histories or polymer clusters have been summed, the resulting effective families are generally nonlinear in $\lambda$. Their control requires a convergent marked expansion; their logarithms need not remain affine. If the formed arguments themselves are changed with the input, their derivatives are additional terms in \eqref{eq:zr-local-evaluation-map}; they require their own domain and norm bounds.

\subsubsection{What an exponential reserve proves}

\begin{proposition}[A one-zone reserve lemma]
Let the fixed affected event $E_a$ consist of histories of age at least $a$
in a measurable disjoint union of history spaces with positive reference
measure. Assume that the integrands are measurable in the stored variables
and holomorphic in $\lambda$ for almost every stored configuration.
Suppose $W_\omega(\xi)\ge0$ is measurable, and assume a measurable majorant $M_\omega$ satisfies

\begin{gather*}
\FitMath{\sup_p e^{\operatorname{Re}\Phi_\omega(0,\xi;p)}\le M_\omega(\xi),} \\[3pt]
\FitMath{\sum_\omega\int_{E_a}M_\omega e^{cW_\omega}\,d\mu_\omega
       \le B\varepsilon q^a,}
\end{gather*}

where $c,C>0$ and $0<q<1$. Put $R=c/(2Cp_n)$, interpreted as infinity
when $p_n=0$. Suppose the perturbation is defined and holomorphic on
$|\lambda|<R$, and throughout that disc
\begin{equation}
 \sup_p|\Phi_\omega(\lambda,\xi;p)-\Phi_\omega(0,\xi;p)|
       \le C|\lambda|p_n W_\omega(\xi).
 \label{eq:zr-phase-hypothesis}
\end{equation}
Then on $|\lambda|p_n<c/(2C)$ the integral over $E_a$ is bounded in modulus by $B\varepsilon q^a$, and is holomorphic. Its difference from its value at zero is bounded by $2B\varepsilon q^a$. Its first derivative at zero is at most

\begin{gather*}
\FitMath{(2CB/c)\varepsilon q^a p_n.}
\end{gather*}

The same conclusions hold with an additional insertion which is measurable
in the stored variables, holomorphic in $\lambda$ on the same disc, and
uniformly bounded there, multiplied by that uniform bound. A fixed bounded
insertion is a special case.
\end{proposition}
\begin{proof}
The modulus of the perturbed integrand is at most
$M_\omega e^{cW_\omega/2}$. The assumed positive majorant is integrable over the disjoint union of histories. Dominated integration, or Morera's theorem and Fubini, proves holomorphy on the open disc. The mass estimate follows from the same domination. The difference is bounded by the sum of the two masses. Cauchy's formula on circles of radius approaching $c/(2Cp_n)$ gives the derivative bound. If $p_n=0$, the hypothesis makes the perturbation independent of $\lambda$, and the derivative is zero.
\end{proof}

Finite upper integrals can replace the measurable majorants in this proposition: choose measurable majorants with integrals arbitrarily close to their upper integrals, apply the argument, and take the infimum. The supremum over external pairs must remain inside the history integral.

A reserve can also produce the age factor. If $W\ge a$ on $E_a$ and

\begin{gather*}
\FitMath{\sum_\omega\int M_\omega e^{(c_{\rm pert}+c_{\rm age})W}\,d\mu_\omega
       \le B\varepsilon,}
\end{gather*}

then a perturbation paid by $c_{\rm pert}W$ leaves the estimate
$B\varepsilon e^{-c_{\rm age}a}$. Thus a per-zone radius independent of the age is compatible with geometric forgetting. The perturbation budget and age budget must be allocated before the coupling ceiling is selected.

This argument is effective only after \eqref{eq:zr-phase-hypothesis} has been proved. An oscillation bound for the unperturbed factors does not establish it.

\paragraph{The ordinary boundary terms in the zone allowance.}
Fix a history and time. Count each live first-group
boundary term once at its stored level \(e\), with one
term per formation domain at that level. Its assigned
anchor is a level-\(e\) cube in the native anchor
region. For a term whose domain misses the live
strips, the dial majorant is
\[
 C_{99} e^{-p_0(g_e)}
           e^{-(1+\beta/2)\kappa d_e(X)} .
\]
At a fixed anchor \(Q\) the ordinary rooted
domain sum is bounded by \(K_0\) at rate
\(\kappa-1\). Since each domain \(X\) contains \(Q\),
the summed majorant there is at most
\(C_{99}K_0e^{-p_0(g_e)}\).

Write \(\Gamma_l^\natural\) for the native
level-\(l\) strata, and use
\(N_l(S)=\ell_l^{-4}\operatorname{vol}(S\cap\Gamma_l^\natural)\)
for the volume in units of one native lattice cell.
The anchor geometry puts \(Q\) entirely in strata
with \(l\le e\). Since its side is \(M\ell_e\),
its contribution to the anchor count is
\[
 1=M^{-4}\sum_{l\le e}L^{-4(e-l)}N_l(Q).
\]
The level-\(e\) anchor cubes are disjoint.
Thus for anchors contained in \(S\), the summed
majorant is at most
\[
 C_{99}K_0 M^{-4}e^{-p_0(g_e)}
     \sum_{l\le e}L^{-4(e-l)} N_l(S).
\]
This uses the native stratum of each cell even
when it is finer than the anchor. In a convention
counting native \(M\)-cubes, absorb their volume
factor into the fixed coefficient.

The lower flow window gives
\(x_e=g_e^{-2}\ge X-2c_2+\lambda(K-e)\) for a
run with \(x_K\ge X\), including \(e=K\).
For \(p_0(g_e)=A_0(\log x_e)^{p_0}\) with \(p_0>1\)
we obtain the history-independent bound
\[
 \sum_{e\le k} e^{-p_0(g_e)}
 \le \sum_{t\ge0}
       \exp\!\bigl(-A_0[\log(X-2c_2+\lambda t)]^{p_0}\bigr).
\]
For large \(X\) the series is dominated by the
summable inverse squares of \(X-2c_2+\lambda t\),
so it tends to zero with the coupling ceiling.
This gives the all-level allowance for ordinary
first-group terms. On domains meeting live strips
the size estimate and the counting norm are
different. To use a small sum over the levels
we must sum the relative estimates over the
actual levels of storage. The age estimate below
also has a version using a uniform level bound.

\paragraph{The relative count in the strip allowance.}
For the relative comparison family the
formation estimate supplies a size
\(c_1^\flat(g_e)\le Cg_e^{\kappa_0}\),
\(\kappa_0\ge7\), after removing the
relative length weight. Its level sum is
bounded by
\[
 C\sum_{t\ge0}(X-2c_2+\lambda t)^{-\kappa_0/2},
\]
which tends to zero with the ceiling.
This is a sum of weighted suprema over
production domains, before rooting.
At an exterior anchor, the rooted
relative count converts each supremum
to a sum of values. For an activity bound
the native site expansion must also supply
the rate margin converting those values
to activities. The site anchor must
stay outside the excused islands with
the prescribed separation.

There is a different count for the
volume bound in a region \(Y\). When
the filed domain is contained in \(Y\)
and contains an excused island, assign
it to the first whole island \(I\) it
contains. The relative rooted sum
at \(I\) is at most
\(1+3^dC_{\rm fib}N_e(I)\), with
\(N_e(I)\) the number of its level-\(e\)
cubes. Thus at each level its cost is
proportional to the number of native
determining cells of \(I\). Summing
over islands of that level uses their
disjointness, and summing the size
coefficients then uses the level sum.
This pays the volume bound on the
admitted whole-island reads. A
single interior anchor for an activity
would lack this whole-island volume.

For the admitted one-run strip terms
filed on their tree domains, use the
unweighted rooted sum for terms formed
by \(\mathbf R\), after subtracting their
declared reference values. Their formation decay
\((1+\beta/2)\kappa\) pays the count of
domains over each tree at rate \(\lambda_\theta\)
and leaves at least \(c_d+1\) for the
tree count. The resulting bound at a
tree root is \(Cc_1^\flat(g_e)\).
This gives a true-read bound at the
formation level \(e\).
A direct read at a level-\(e\) cube
uses a tree root within one layer.
A pocket read uses a root within two
cells of its strip. There are a bounded
number of possible roots for pockets
of side at most one. The disjoint
larger pockets contribute at most
\(909^d\) choices, each with at most
\((s_e+7)^d\) roots for side \(s_e\).
For a pocket born at \(m\le e\),
\(s_e\le101\check R_m L^{-(e-m)}
\le101\check R_e\), by the lower bound
on the radius after each step.
Thus the reading-cube sum for this
class of \(\mathbf R\)-formed terms
is at most \(C(1+\check R_e)^d c_1^\flat(g_e)\).
Since \(\check R_e\) is polylogarithmic
in inverse coupling and \(\kappa_0\ge7\),
the bound is at most \(Cx_e^{-2}\) at
large \(x_e\). The lower flow window
therefore also sums it over levels.

% Earlier Part I 66583--66619, 77719--77736,
% 77798--77812, 81772--81815; Part III 24717--24774,
% 26557 and 126048--126088.
For the uniform bound at a fixed level there
is a direct count for the dialled family.
Use the individual first-group strip
domains \(X\) before combining their
values. At birth such an \(X\) contains
a union of whole excused strips and has
a nonempty part outside them. Its filed
skeleton \(A(X)\) is chosen from a shortest
relative tree. It contains \(X\) outside
the strips, and each included strip
touches it. The term reads \(A(X)\) and
the attached strips, including the older
strips read through a touching cube.
The separation of the new strips puts
one of these off-strip cubes in the
regular part. The filing rule keeps a cube
of the skeleton in \((Z_e^{\boxplus})^{(1)}\)
of the current label for a term still
stored at level \(e\).

On the admitted age bidisc use the
first-group strip majorant
\(C_{99}\alpha^{1/3}
e^{-(1+\beta/2)\kappa d_{\rm rel}(X)}\).
The strip amplitude \(\alpha\) is fixed.
The skeleton has \(d_e(A)\le d_{\rm rel}(X)\),
and there are at most \(e^{\lambda_\theta(d_e(A)+1)}\)
constituents over one \(A\). Summing the
individual majorants over that fibre leaves
\[
 C_{99}e^{\lambda_\theta}\alpha^{1/3}
 e^{-(1+\beta/4)\kappa d_e(A)}
\]
when \(\beta\kappa/4\ge\lambda_\theta\).
This bound also applies to any selected
subfamily, so the witness marks may be
kept on the individual terms. At a filed
cube, the rooted ordinary tree sum gives
a pile bounded by \(C\alpha^{1/3}\).
On the admitted relative-format reads
of \(\mathbf T\), the preparatory move,
and the large-field move, the same filed
skeleton and attached strips are used for
marked and dialled terms. The native
cube-to-cell attribution charges this
pile to the volume and activity bounds.
This gives the uniform part of the age
estimate without a summable strip amplitude.

To use the count on every actual mixed
history we must admit its fluctuation and
preparatory readers on the enlarged age
bidisc with the increased history weight.
For an inherited crossing in the prepared
sharp mask, the quiet-crossing calculation
gives the real domains on the approximate
pass. On a broader reader we must still
account for each pulled-back indicator,
and show that the densities on the fixed
domains have the required holomorphic
majorants. The skeleton sum applies once
those readers and majorants have been identified.

\paragraph{Carrying the age bound through earlier terms.}
Compare two lists on the same lattice and scaffold.
Give each use of a list entry in a production formula
its own parameter. A parameter of class \(z=(n,k')\)
is bounded using the load of zone \(n\). The
diagonal holomorphy proof extends to these separate
parameters when it charges each direct use by its
absolute size, and reads a previously produced object
only through its value and size bound. Indeed the
size allowance of every direct input is the same on
its own disc. In the order of formation, the earlier
inputs retain their admitted suprema, so the same
composition estimate gives the next holomorphic
object. In this product of discs one may turn off a
set of uses without moving the other parameters.

An occurrence means a stored object or a copy of
its formula inside an object that read it. The
outer domain contains the domain of the copy.
For a boundary term of level \(e\) its witnessing
components are those of \(Z_e\) whose enlarged
sets meet that domain. A frozen term uses its
arrested component. A use with \(n<k'\) starts
with the geometric witness containing its anchor in
\(\Omega_n\setminus\Omega_{n+1}\subset Z_{n+1}\).
A top-zone use starts in the own set of the object
that created it. A fluctuation or an arrest
deletes no witness. At healing, a carrier of a
use with a deleted geometric witness meets the
operated set at its anchor. If an object loses
its last witnessing component, the carrier
meets the operated set through the domain of
that object. The absorption rule places these
uses in the newly produced terms. All other
uses keep their live witnesses, including those
read through copies. This gives a unique current
assignment along each occurrence, with witness
level \(n+1\) for a surviving geometric witness
and level \(e\) for a surviving owner of level
\(e\). An occurrence carrying a copy need not
have the same level as its witness.

In a healed component \(X_i\), the histories
that contain a direct old use of zone \(n\)
must have \(Z_{n+1}\cap X_i\ne\varnothing\).
For a non-top use this follows from the zone
read in the enlarged component and the filling
rule for its older cells. An inner top use at
level \(k'=n\) is supported on
\(Z_{k'+1}\cap X_i\). A deleted geometric
witness must itself be present in \(Z_{n+1}\)
of the component history. In all these cases the
first index of that history is at most \(n+1\).
For an absorbed leaf with witness at level \(e\),
the component history must contain the witness in
\(Z_e\). This also applies if the leaf occurs
inside a carrier, because that carrier's formation
history includes the leaf's. At operation \(j\)
these facts give ages at least \(j-n-1\) and
\(j-e\), respectively.

For a first-group boundary term or a
second-group remainder \(\varphi\) formed
by \(\mathbf R_j\), fix
a class \(z=(n,k')\) of the source uses, and compare
the value of \(\varphi\) before and after turning off
its own uses. These are the direct old uses of the
healed histories, the direct top uses at formation,
and the uses whose witnesses are lost at the operation.
Let \(\mathfrak B_\varphi\) be its holomorphic
majorant and \(0<q\le1/2\) the age rate. The
history decomposition partitions the uses into
one group of direct and low-level uses and groups
of nested leaves of levels \(n+2\le e<j\).
A nested leaf in such a group loses its last witness
inside a healed component of this operation.
The history-absorption rule is needed also for
occurrences inside another stored term: a carrier
containing such a leaf contains its witnessing
anchor and is absorbed with that component.

Turn off the level groups one at a time.
For a group of level \(e\), replace a leaf
value by its value with the selected uses off
plus \(\nu\) times its change. If the carried
bound at level \(e\) is
\(C_*\mathfrak B_\psi q^{(e-n-1)_+}\)
for each such leaf \(\psi\), take
\[
 |\nu|\le R_e
    =[C_*\varepsilon_e^{\rm all}q^{(e-n-1)_+}]^{-1},
\]
where \(\varepsilon_e^{\rm all}\) is the
per-cell pile for all such leaves. A zero pile
requires no interpolation. This choice makes
the changed part at most
\(\mathfrak B_\psi/\varepsilon_e^{\rm all}\).
The summed changed parts therefore cost one
additional unit of the per-cell allowance
at a direct read with the stored-family count.
An intermediate term, once formed, enters
subsequent steps by its own majorant.

At the same time, multiply the sub-sums of
healed histories of age at least \(j-e\)
by \(\mu\), leaving the younger histories
unchanged. The bound on the age tail permits
\(|\mu|\le q^{-(j-e)}\) at the price of
doubling the component input. At \(\mu=0\)
the level-\(e\) leaves in this group cannot
occur: the last-witness condition requires
age at least \(j-e\) in the component that
absorbs the leaf. The term on this face
is therefore independent of \(\nu\).
At \(\mu=1\) the difference between
\(\nu=1\) and \(\nu=0\) is the required
change for this group. On an admitted
bidisc with bound \(\mathfrak B_\varphi\),
subtract both axis values and add the value
at the origin. The resulting function
vanishes on both axes and has supremum
at most \(4\mathfrak B_\varphi\).
Schwarz's estimate in the two variables
therefore bounds the group change by
\[
 4\mathfrak B_\varphi\,q^{j-e}/R_e
 =4C_*\varepsilon_e^{\rm all}\mathfrak B_\varphi
   q^{j-n-1}.
\]
Here the native formation estimate admits both
parameters by two separate charges. First use
the enlarged constant in the exponential
volume bound for stored values and in the
site activity, allowing one more unit of the
per-cell stored-value bound at each direct
read. The history tail estimate is taken
at this enlarged constant, so it is uniform
in \(\nu\). Then split the sequence sum in
each healed component. At \(|\mu|\le q^{-(j-e)}\)
the two pieces together cost at most twice
the unsplit component majorant. This factor
is paid in the product over components.
At a component formed at step \(j\), the
native geometric factor is at most
\(e^{2\kappa(100\check R_j)^4}\).
For the native growth number
\(\beta_0^{\rm pr}\le1/7\), the component
history bound is at most \(e^{1-7p_0(g_j)/4}\)
before the split. Thus the doubled bound
leaves \(e^{-3p_0(g_j)/2}\) by the
native coupling ceiling, which implies
\[
 -p_0(g_j)/4+2+2\kappa(100\check R_j)^4\le0 .
\]
The next formation sum costs a factor
\(e^{p_0(g_j)/2}\) and retains
the original \(e^{-p_0(g_j)}\) coefficient
per component. This proves the budget on
one age disc and one leaf-level disc in
the same receiving cluster, with the
other interpolations at their endpoints.
It uses the admitted per-use formation
and the history tail bound at the larger
volume and activity allowances.

After the level groups have been turned
off, the remaining uses require a healed
history of age at least \(a=(j-n-1)_+\),
except for direct top-level uses when \(a=0\).
The single age split at \(a\) gives
\(2\mathfrak B_\varphi q^a\) by the
same zero-face argument, or by the
trivial difference bound at \(a=0\).
Summing the levels gives the factor
\(2+4C_*\varepsilon_{\rm tot}^{\rm all}\)
in front of \(\mathfrak B_\varphi q^a\).
Here \(\varepsilon_{\rm tot}^{\rm all}\) bounds
the sum of the level piles uniformly.
Thus \(C_*=4/(1-4\varepsilon_{\rm tot}^{\rm all})\)
closes the induction when the total pile
is at most \(1/16\).

\paragraph{The carry with a uniform bound at each level.}
There is a second way to close the age estimate.
Suppose the actual leaf families have level piles
\(\varepsilon_e^{\rm all}\le B\) on the enlarged
formation and history domains. Put \(A=1+4B\).
Choose \(0<q_0\le q\) with \(2Bq_0\le1/3\),
and set \(q'=q_0/A\). Use the history weight
increased by \(\log(q/q')\). A history of age at
least \(s\) has volume at least \(s\), so the
weighted mass bound gives the share \((q')^s\).
The weight is fixed before the coupling ceiling.
We must have the formation bounds and the same
level-pile bound \(B\) at this weight.

With \(a=(j-n-1)_+\) use the inductive bound
\(D_a\mathfrak B_\varphi(q')^a\), \(D_a=2A^a\).
At a nested leaf of age \(b=e-n-1\ge1\), the inverse
leaf radius is \(\varepsilon_e^{\rm all}D_b(q')^b\),
at most \(2Bq_0^b\le1/3\). The same two-parameter
estimate bounds its level group by
\(4\varepsilon_e^{\rm all}D_b\mathfrak B_\varphi(q')^a\).
The direct contribution is at most
\(2\mathfrak B_\varphi(q')^a\).
The geometric sum gives
\[
 2+4B\sum_{b=1}^{a-1}D_b\le D_a
\]
with the sum empty for \(a\le1\). This closes the
induction, since \(D_a(q')^a=2q_0^a\le2q^a\).

For a term of the second group, its
domain is disjoint from the surviving
large-field regions and the enlargements
that would make the domain first-group.
A class use with a
surviving geometric witness would put
its anchor in one of those regions.
A copied occurrence with a surviving
witness would put its anchor in the surviving
region or its prescribed enlargement,
which is excluded for a second-group
domain. Thus every class use of a second-group
term belongs to its own set, and the carried
bound controls the full change of that class.
At a native zero-index cluster containing
a live strip \(Y_h\) and no healed
constituent, the only own uses
are the direct top uses at formation.
There is then no age loss in bounding
their change by twice the strip majorant.
Uses copied from a surviving witness
are still assigned to that witness level.
This special case gives a leaf difference
bound and still requires the strip count
for summation.

\paragraph{Reading the transported uses at the later component.}
Let \(X\) be a component of \(Z_k\) on which
the operation acts, and put
\(G=\mathbf T'_k(X)F\) for a fixed bounded
insertion \(F\). In each history summand of \(G\),
the filling property takes an older anchor in
\(X^\sim\cap Z_{n+1}\) into \(X\).
A surviving geometric witness at that anchor
therefore requires the summand's \(Z_{n+1}\)
to meet \(X\). An earlier owned use
of witness level \(e<k\) instead has an
anchor of its owner in \(Z_{e+1}\cap X^\sim\).
The same filling property puts that anchor
in \(X\), and its live witnessing component
lies in the same surviving component as the
anchor. The history containing the occurrence
must contain its witness already at \(e\).
Hence a use of witness level \(e\) requires
age at least \(k-e\) in the summand.
Direct inner top uses and constants attached
to a zone satisfy the direct bound \(b=(k-n-1)_+\)
by the index of the region at which they occur.

To turn off class \(z=(n,k')\), first
telescope over the own sets of live leaves
of levels \(n+2\le e<k\). These sets are disjoint
as sets of uses, and the surviving witness
assignment is made in each summand's own
history. For a fixed level use the disc of
radius \(R_e\) above. Its change is supported
on histories of age at least \(k-e\).
The absolute tail majorant for this subfamily,
with the extra stored-value unit included,
is \(\mathrm e\,\mathfrak E_k q^{k-e}\sup|F|\),
where \(\mathfrak E_k
=\exp[-2p_0(g_k)/(1+\beta_0^{\rm pr})]\).
It applies even though different
summands use different points of the product
of discs, since the tail bound sums the
individual majorants. Schwarz's estimate in
the single leaf parameter gives
\(2\mathrm e\,\mathfrak E_k q^{k-e}\sup|F|/R_e\).
The powers \(q^{k-e}q^{e-n-1}\) combine to
\(q^{k-n-1}\).

The remaining direct uses, geometric uses,
and own sets of carriers or leaves at levels
at most \(n+1\) require age at least \(b\).
Their difference costs at most twice the
age-\(b\) tail bound. Adding the telescope
therefore bounds the change of class \(z\) by
\[
 2\mathrm e\,\mathfrak E_k q^b\sup|F|\,
       (1+C_*\varepsilon_{\rm tot}^{\rm all}).
\]
For the uniform level bound, repeat the telescope
with \(q'\) and the age-dependent leaf radii. Writing
\(a=(k-n-1)_+\), it gives
\[
  2\mathrm e\,\mathfrak E_k(q')^a\sup|F|
     \left(1+B\sum_{d=1}^{a-1}D_d\right)
  \le 2\mathrm e\,\mathfrak E_k q^a\sup|F|,
\]
because \(1+B\sum_{d=1}^{a-1}D_d\le A^a\).
On the diagonal of the product of discs this
is the difference from setting the class
parameter to zero. Apply the one-variable
disc estimate in each class in succession.
At zone \(n\) there are \(k-n+1\) classes
\((n,k')\), each with the inverse disc radius
given by the zone load divided by the
fixed load allowance. Thus the change
between the endpoints is controlled by the
load sum with weights
\((k-n+1)q^{(k-n-1)_+}\).

\subsubsection{A common reserve for many zone uses}

One-zone bounds do not by themselves give a common polydisc for an
arbitrarily large number of zone uses. The reserve also needs a capacity
bound. The following precise version distinguishes the two issues.

\begin{proposition}[A capacity criterion for joint parameters]
Let $\mathcal Z$ be a finite set of zone-use indices and let $i$ label
finitely many directions at each index. Their nonnegative loads are
$p_{z,i}$. Use the measurable stored spaces and fixed reference measures
of the preceding proposition, and assume that all phase coefficients are
measurable in the stored variables. Fix nonnegative measurable history charges $W_{\omega,z}$,
a finite family of nonnegative measurable action pieces $S_{\omega,b}$, and nonnegative
coefficients $a_{\omega,z,b}$, with

\begin{gather*}
\FitMath{\sum_{z\in\mathcal Z}W_{\omega,z}\le W_\omega,} \\[3pt]
\FitMath{\sum_{z\in\mathcal Z}a_{\omega,z,b}\le1} \\[3pt]
\FitMath{\hbox{for every }b.}
\end{gather*}

Suppose the affine perturbation is

\begin{gather*}
\FitMath{\Phi_\omega(\vec\lambda)-\Phi_\omega(0)
    =\sum_{z,i}\lambda_{z,i}L_{\omega,z,i},} \\[3pt]
\FitMath{|L_{\omega,z,i}|
 \le p_{z,i}\left(CW_{\omega,z}
                    +\sum_ba_{\omega,z,b}S_{\omega,b}\right).}
\end{gather*}

Assume a majorant at zero obeys, on a fixed event $E$,

\begin{gather*}
\FitMath{\sum_\omega\int_E M_\omega
       \exp\left(cW_\omega+\sigma\sum_bS_{\omega,b}\right)d\mu_\omega
     \le B_E,} \\[3pt]
\FitMath{\sup_p e^{\operatorname{Re}\Phi_\omega(0,\xi;p)}\le M_\omega,}
\end{gather*}

where $c,\sigma>0$. Choose
$0<r<\min\{c/(2C),\sigma/2\}$. Then the integral over $E$ is jointly
holomorphic and bounded by $B_E$ on the product of weighted balls

\begin{gather*}
\FitMath{\mathcal D_r=\left\{\vec\lambda:
       \sum_i|\lambda_{z,i}|p_{z,i}<r
                  \text{ for every }z\right\}.}
\end{gather*}

All statements use fixed reference measures, fixed events and fixed
evaluation maps; if additional parameter dependence is present, the
integrands must instead be given holomorphically on this domain with
the same phase bound.
\end{proposition}
\begin{proof}
For $\vec\lambda\in\mathcal D_r$, the triangle inequality and the two
capacity inequalities give

\begin{gather*}
\FitMath{|\Phi_\omega(\vec\lambda)-\Phi_\omega(0)|
 \le Cr\sum_zW_{\omega,z}
       +r\sum_b\left(\sum_za_{\omega,z,b}\right)S_{\omega,b}
 \le (c/2)W_\omega+(\sigma/2)\sum_bS_{\omega,b}.}
\end{gather*}

Thus a single integrable majorant works on the entire parameter domain,
independently of $|\mathcal Z|$ and the number of directions. Dominated
integration proves separate holomorphy and local boundedness, and the
finite-dimensional Hartogs theorem gives joint holomorphy. Equivalently,
one can apply Morera's theorem successively on each compact polydisc.
\end{proof}

If a designated parameter affects only an event $E_a$ whose majorant
has $B_{E_a}\le B\varepsilon q^a$, setting that parameter to zero
changes the whole integral by at most $2B\varepsilon q^a$, uniformly
in the other parameters on $\mathcal D_r$. At their common zero,
Cauchy's formula gives the first derivative bound
$B\varepsilon q^a p_{z,i}/r$. To see this without subtracting a large
unaffected part, differentiate only the integral on $E_a$; outside
$E_a$ the integrand is independent of the designated parameter.
For $m$ distinct directions at the same index, Cauchy's formula on
circles with radii tending to $r/(mp_{z,i})$ gives

\begin{gather*}
\FitMath{\left|\partial_{\lambda_{z,1}}\cdots
             \partial_{\lambda_{z,m}} I_E(0)\right|
   \le B_E(m/r)^m\prod_{i=1}^m p_{z,i}.}
\end{gather*}

Here zero loads mean that the corresponding affine perturbations vanish,
and those derivatives are zero. Higher mixed bounds require an appropriate
majorant for the history event on which all marked uses occur; they do not
acquire a product of rarity factors from holomorphy alone.

The capacity assumptions have content. Without them, let a single
history have $W=1$ and let each of $N$ indices contribute
$L_z=p_z=1$. Every separate perturbation is bounded by $|\lambda_z|W$,
but choosing all $\lambda_z=r/2$ gives a total exponent change $Nr/2$.
A fixed reserve in $W$ cannot pay this uniformly in $N$. A proof for the
actual formation must therefore assign all simultaneous uses to a common
site and action budget, in addition to bounding each separate direction.

\subsubsection{A channel which can be bounded directly}

For $U\in\operatorname{SL}(2,\mathbb C)$, the determinant identity

\begin{gather*}
\FitMath{\det(U-I)=\det U-\operatorname{tr}U+1
          =2-\operatorname{tr}U}
\end{gather*}

gives $a(U)=\frac12\det(U-I)$. With the operator matrix norm,
\begin{equation}
 |a(U)|\le\tfrac12\|U-I\|^2.
 \label{eq:zr-plaquette-square}
\end{equation}
For real $U\in\operatorname{SU}(2)$ this is an equality, since the two singular values of $U-I$ coincide.

Consider a regular zone expressed in its own units, with bare mesh $\eta_n$, containing $N_n$ unit cells. Suppose the configurations used in the actual substitutions satisfy
$\|U_p-I\|\le\alpha_n\eta_n^2$ on that zone. There are at most $C\eta_n^{-4}$ plaquettes per unit cell. Hence
\begin{equation}
 |A(\delta c,U)|
    \le C\alpha_n^2\|\delta c\|_\infty N_n.
 \label{eq:zr-regular-weight}
\end{equation}
Consequently, if the zone load includes
$p_n\ge\alpha_n^2\|\delta c\|_\infty$, this direct Wilson-weight contribution to \eqref{eq:zr-local-evaluation-map} is bounded by $Cp_nN_n$. If the reserve volume dominates the number of such stored cells, it satisfies \eqref{eq:zr-phase-hypothesis}.

The usual scale $\alpha_n^2\asymp\theta_n^{-1}(\log\theta_n)^{2q}$ has precisely the form of the coefficient multiplying a marginal amplitude in the load. This calculation uses only the stated curvature regularity at the actual formed arguments. It does not prove that every substitution has that regularity; that is a separate inclusion requirement.

\subsubsection{An unrestricted bare zone needs action reserve}

There is a genuine limitation to a reserve measured only by stored sites. On a bare region where the real field is unrestricted, a plaquette can have $U_p=-I$, and then $a_p=2$. A direction which changes a local Wilson weight by one has an exponent variation of order one at such a configuration. If its proposed load is

\begin{gather*}
\FitMath{p_0=\beta^{-1}(\log\beta)^C}
\end{gather*}

and the site reserve is of order $W=N\log\beta$ on a fixed finite region, then

\begin{gather*}
\FitMath{\frac{|\delta\Phi|}{p_0W}
     \ge c\,\frac{\beta}{(\log\beta)^{C+1}}}
\end{gather*}

for an appropriate fixed configuration and fixed $N$. This diverges as $\beta\to\infty$. Such configurations are obtained, for example, by setting one bond to $-I$ and the other bonds of a small patch to $I$. The number of altered plaquettes remains bounded. Thus \eqref{eq:zr-phase-hypothesis} with this site reserve cannot hold pointwise on the entire unrestricted bare-field domain.

This does not refute the complex mass estimate. The Wilson factor strongly suppresses those configurations. For a bare insertion $\mathfrak b$ satisfying a pointwise action bound

\begin{gather*}
\FitMath{|\mathfrak b|\le C_b\|k\|_\infty A_{\rm available}(U),}
\end{gather*}

one instead has

\begin{gather*}
\FitMath{|\lambda\mathfrak b|
  \le\frac{C_b|\lambda|\|k\|_\infty}{\beta}
           S_{\rm available}(U),} \\[3pt]
\FitMath{S_{\rm available}=\beta A_{\rm available}.}
\end{gather*}

At a radius $|\lambda|p_0<c$, with
$p_0\ge C_b\beta^{-1}(\log\beta)^C\|k\|_\infty$, this spends at most
$c(\log\beta)^{-C}$ of the available Wilson action. It is a fractional action estimate, not a site-volume estimate.

A correct reserve lemma can include both budgets. If the majorant has the stronger reserve
$e^{cV+\sigma S_{\rm available}}$ and

\begin{gather*}
\FitMath{|\Phi(\lambda)-\Phi(0)|
       \le C|\lambda|p_nV+\eta(\lambda)S_{\rm available},} \\[3pt]
\FitMath{C|\lambda|p_n\le c/2,} \\[3pt]
\FitMath{|\eta(\lambda)|\le\sigma/2,}
\end{gather*}

the mass proof above applies on the intersection of these parameter
domains. A derivative estimate proportional to $p_n$ requires this
intersection to contain a disc of radius $c' /p_n$ with a fixed $c'>0$.
One sufficient condition is $|\eta(\lambda)|\le C_{\rm act}|\lambda|p_n$;
otherwise the Cauchy radius is the smaller action-controlled radius.
To use this statement one must identify action genuinely left available
after the failed-test payments, gauge fixing, and other insertions.
It cannot be spent twice. For several active zone-use parameters, their
fractions must have a bounded sum on each shared action piece. The existence
of a site-volume reserve alone proves neither of these action-allocation facts.

\subsubsection{What the carrier bounds do not control}

Write a standing exponent as

\begin{gather*}
\FitMath{\Phi_\omega(\xi;p)=\Phi_\omega^0(\xi)
                 +\sum_f\phi_{\omega,f}(\xi;p|_{R_f}).}
\end{gather*}

The carrier condition bounds the oscillation of each factor in its external pair argument. Choose a permitted anchor $p_\omega^0(\xi)$ and split

\begin{gather*}
\FitMath{\phi_{\omega,f}(\xi;p)=
  \phi_{\omega,f}(\xi;p_\omega^0)
  +\{\phi_{\omega,f}(\xi;p)-\phi_{\omega,f}(\xi;p_\omega^0)\}.}
\end{gather*}

Oscillation controls the braced expression. It controls neither the scalar anchor value nor the dependence of either expression on an input-list parameter. Cauchy's estimate in the pair variables is not Cauchy's estimate in that new input parameter.

A simple exact example separates these statements. Let $\pi_a\downarrow0$ be a sequence of putative loads, let one stored history have volume $W_a\ge a$, and let its formed exponent be

\begin{gather*}
\FitMath{\Phi_a(x;p)=\log\varepsilon-(c+s+1)W_a+xW_a,} \\[3pt]
\FitMath{s>0.}
\end{gather*}

It is entire, independent of $p$, and its carrier oscillation is zero. At $x=0$ it has the prepaid mass bound
$e^{\Phi_a(0)+cW_a}\le\varepsilon e^{-sa}$. The same bound, with room, holds throughout a fixed ordinary disc $|x|\le1/2$. It fails on a disc $|x|\pi_a<c'$: at $x=c'/(2\pi_a)$ the positive exponent eventually overwhelms every fixed reserve coefficient.

This is a counterexample to deriving the enlarged-radius estimate from carrier oscillation and an ordinary reserve alone. It is not a counterexample to the Wilson operation. Even if the incoming direction has size $\pi_a$, the example can be produced by an abstract local formation map which multiplies that direction by $\pi_a^{-1}$. All input sizes are then small, while the output exponent has the displayed variation. Locality and holomorphy do not exclude this map; a uniform norm bound on formation must exclude it. Exact preservation of a represented integral, by itself, also supplies no such bound on a particular decomposition into exponents and vacuum factors.

\subsubsection{A sufficient structural bound on the actual formation map}

The needed additional statement is precise. For the actual index set and formed arguments in \eqref{eq:zr-local-evaluation-map}, prove
\begin{equation}
 |L_{\omega,n,k'}[\delta](\xi,p)|
      \le C p_n(\delta)V_\omega(\xi)
       \quad\text{for every admissible }(\xi,p),
 \label{eq:zr-formation-test}
\end{equation}
with a separate fractional-action term when the bare zone is unrestricted. The norm must include scalar anchors, not only oscillations. The constant must be independent of the age, the level, and the history. If a parameter changes the arguments as well, prove the corresponding derivative bound throughout its proposed disc and integrate it to obtain \eqref{eq:zr-phase-hypothesis}.

One verifiable route to \eqref{eq:zr-formation-test} is a carrier incidence estimate.
Let the charged stored cells be a finite set $\mathcal C_\omega(\xi)$
with $|\mathcal C_\omega(\xi)|\le V_\omega(\xi)$, and decompose the
scalar part as $\delta c_\omega=\sum_{c\in\mathcal C_\omega}\delta c_{\omega,c}$.
Assign each formed occurrence to these cells and give nonnegative weights
$w_{e,c}$ with $\sum_cw_{e,c}=1$. Suppose the actual evaluations stay in their original domains and

\begin{gather*}
\FitMath{\sum_e |\epsilon_e|w_{e,c}\,
       \|\delta F_e\|_{\mathrm{evaluation}}
       +|\delta c_{\omega,c}|\le Cp_n(\delta)}
\end{gather*}

at every stored cell $c$. Summing over the cells proves
\eqref{eq:zr-formation-test}. The evaluation norm here is the norm after the specified unit change and counter-insertions. Using a norm at an earlier level without those transformations can lose exactly the scale weights under investigation.

A condition on the sum of sizes of the original direction terms in a unit cell is helpful, but it is not yet this incidence statement. The proof must show that substitutions and repeated occurrences preserve that bound. The formation rules distinguish copies still anchored in a live component, copies whose last anchor is removed now, and copies inside older already formed terms. The first kind is not a small new remainder. The last kind must be summed with its own norm and history mass. Counting its occurrence does not permit its bound to be reused without payment. In particular a factor obtained only after summation of histories is not a pointwise bound on the unsummed exponent.

% Earlier outline, 45354--45364.
For the deletion of a class born at level \(n\), use
the weighted mass on the class disc with survival
ratio \(q_*\). Suppose a term formed at level \(e\)
is read in a component that remains live until
level \(j>e\). Its own deletion bound has age
factor \(q_*^{e-n-1}\) when \(e\ge n+2\).
The later mass is restricted to histories that
carry this anchor from \(e\) to \(j\), giving
the factor \(q_*^{j-e}\). At healing the anchor
is removed and the two factors combine to
\(q_*^{j-n-1}\). Keep the size of the older
term in this product. A term carried by a
component left standing at its production
costs a fixed amount per level. The genuine
small remainders also carry the small factors
of their producing levels, which sum along the
flow.

More precisely, let \(A_j\) be the constant
after extracting \(q_*^{(j-n-1)_+}\) at
level \(j\). The per-level sum over the older
reads, including the size \(\epsilon_e\) of
their terms, gives
\[
 A_j\le C+C\sum_{e=n+2}^{j-1}\epsilon_e A_e .
\]
Here \(\epsilon_e\le\epsilon+\sigma_e\),
with the fixed bound \(\epsilon\) for the
standing-component terms and with
\(\sum_e\sigma_e\le C'\) for the genuine
remainders. Iterating the positive inequality
bounds \(A_j\) by a constant times
\(\prod_{e=n+2}^{j-1}(1+C\epsilon_e)\),
which is at most a constant times
\((1+C\epsilon)^{(j-n-2)_+}\).
Taking \(q_*=q_a/(1+C\epsilon)\) in the
weighted mass, by adding \(\log(1+C\epsilon)\)
to the reserved exponent per age level,
recovers \(q_a^{(j-n-1)_+}\) with a fixed
constant. This calculation uses the mass
of the actual histories on the disc.
A small new remainder has its domain away
from the surviving region. Since a copied
term's domain is contained in the domain
that carries it, such a remainder contains
no copy with an anchor still live there.

For the fresh ordinary crossing terms there is a direct
incidence bound at fixed history and stored fields. Assign each
complete marked cluster to one visited step cube in the strip
$\Omega_t\setminus\Lambda_t$, when its crossing carrier has such
a free cube. Retain this assignment while summing the complete
clusters. The rooted domain count, with the spare decay in the
marked crossing estimate, bounds the sum at a fixed free cube
by a constant times the crossing factor and the supremum of
its local input loads. For a cluster assigned to an exempt hull
$H$, the corresponding rooted sum costs at most $C\#H$. If
the whole hull is charged, distribute this bound over its
$\#H$ cells. An exempt hull whose cells are not all charged
needs the actual refiling estimate. This argument counts one
fresh crossing occurrence once. An old member carried by
restriction retains its carrier bound, and repeated reads in
later substitutions still require their own incidence count.


If ``size in a unit cell'' is defined to mean the sum of the evaluated sizes of all formed occurrences, including their scalar parts and the indicated unit changes, then the incidence estimate is already part of that definition. Under that stronger interpretation the reserve implication is proved. The substantive remaining task is to show that directions generated by the actual renormalization step meet that definition with the proposed load. A condition on the unformed entries cannot be silently replaced by this stronger condition.

\subsubsection{Passing the bound through a logarithmic sum}

Suppose that a finite polymer system has activities $z_X(\lambda)$
with a common Koteck\'y--Preiss majorant $r_X$, and that the corresponding
pinned bound is $e^{a_X}$. Choose the logarithm defined by the connected
expansion, equal to zero at the zero activity vector. Then a useful marked estimate is

\begin{gather*}
\FitMath{|\log Z(z(\lambda))-\log Z(z(0))|
    \le\sum_Xe^{a_X}|z_X(\lambda)-z_X(0)|,}
\end{gather*}

To prove it, join the two activity vectors by their straight segment. Every point on that segment has modulus at most $r_X$ coordinatewise. The absolutely convergent connected expansion gives
$|\partial_{z_X}\log Z|\le e^{a_X}$ by the pinned bound. Differentiate along the segment and integrate. Thus a summable one-zone age bound on the changed activities passes to the logarithm. No logarithm of an individual possibly zero history integral is required.

For a countably infinite system, apply the same estimate first to finite
restrictions of the same activity and incompatibility system.
If $\sum_X e^{a_X}|z_X(\lambda)-z_X(0)|<\infty$, the integrated
pinned series defines a relative cluster logarithm and satisfies this
bound, independently of the finite exhaustion by dominated convergence
of the rooted expansions. Writing two separate infinite-volume logarithms
requires convergence of both unmarked global series; a common local
majorant alone does not ensure that. This deduction has two inputs: the
activity majorant must remain valid on the complex parameter disc, and
the changed activities must satisfy the marked bound after the actual
formation and copy bookkeeping. A convergence theorem for polymers does
not prove those inputs. The exponential reserve can pay a controlled
perturbation of an integrand; it cannot manufacture control of the
perturbation map.

A recursive treatment of old copied terms can be sufficient instead of a single global phase estimate.
Fix $0<q_0<q<1$ and $C>0$. For example, suppose an already justified local operation estimate gives nonnegative bounds, for $j\ge n+1$,

\begin{gather*}
\FitMath{A_{j,n}\le Cq_0^{j-n}
       +C\sum_{n<e<j}\epsilon_e q_0^{j-e}A_{e,n},} \\[3pt]
\FitMath{0\le\epsilon_e\le\epsilon.}
\end{gather*}

Divide by $q_0^{j-n}$. A discrete Gronwall induction gives

\begin{gather*}
\FitMath{A_{j,n}\le Cq_0^{j-n}
       \prod_{n<e<j}(1+C\epsilon_e)
 \le C\{q_0(1+C\epsilon)\}^{j-n}.}
\end{gather*}

The induction follows by expanding the product, or using
$\prod_{e<j}(1+C\epsilon_e)=1+\sum_{e<j}C\epsilon_e\prod_{r<e}(1+C\epsilon_r)$.
Thus choosing $q_0\le q/(1+C\epsilon)$ preserves the desired rate $q$. No small sum of the $\epsilon_e$ is needed. This pays growth of output bounds after local integration. It does not pay an exponentially amplified perturbation inside an unsummed exponent. To apply this alternative one must establish the local parameterized operation bound at each new loss of an anchor and the stated copying recurrence. The sufficient formation test above is not claimed to be the only possible proof route.

The present test therefore proves the reserve implication and the direct regular-zone Wilson channel. It shows why an unrestricted bare channel requires available action, and isolates the additional uniform formation estimate for the remaining entries. The specified carrier oscillation, ordinary presentation bounds, and end-mark definitions do not alone establish that estimate. A proof of it for the actual local substitutions, with their scalar parts and copy incidence, would complete the proposed one-zone argument.


% END MODULE chapters/zone_reserve_test_v3.tex

% BEGIN MODULE chapters/extraction_radius_audit_v3.tex
\subsection{Coordinates for quadratic extraction}
\label{ch:extraction-radius-normalization}

The quadratic coefficient of an effective interaction is defined in the
physical block field. The field used to normalize a fluctuation Gaussian
is a different coordinate. Keeping these coordinates distinct explains
both the fixed radii of the unmarked chain and the squared radius ratios
in the marked chain. The required inclusion near the trivial field follows
from the background and tower estimates already used to define the
analytic field domains. It is not an additional loss caused by the number
of blocking steps.

All constants below may depend on the fixed blocking and localization
parameters. They are independent of the level, torus size and number of
earlier levels. The non-Abelian background and tower estimates are
explicit hypotheses in the statements that use them.

\subsubsection{Physical fields, Gaussian fields and the coefficient}

At level $j$, write the configuration in the trivial-field chart as

\refstepcounter{equation}
\begin{gather*}
\FitMath{\mathbf U=e^{i\eta_j A'},} \\[3pt]
\FitMath{v=(A',J),} \\[3pt]
\FitMath{\|v\|_{1,j}=\max\{\|A'\|_\infty,
                    \|\nabla_j A'\|_\infty,\|J\|_\infty\}.}\tag{\theequation}\label{eq:extraction-radius-ER1}
\end{gather*}

Here $0<\eta_j\le1$ is the bare spacing in level-$j$ units and $\nabla_j$
is a difference quotient in those units. The current $J$ is an
independent analytic variable. In particular, $A'$ in this chart is a
physical potential; no coupling multiplies it by definition.

The block datum of the background is $V=e^{iB}$. Let $p_j(B)$ be its
configuration-current pair in a fixed gauge near zero. Derivatives of its
configuration entry are taken in the chart in \eqref{eq:extraction-radius-ER1}. Suppose
\begin{equation}
\FitMath{p_j(0)=(\mathbf1,0),\qquad
 Dp_j(0)B=\mathfrak h_jB
       =(H_jB,D^*D H_jB).}
 \label{eq:extraction-radius-ER2}
\end{equation}
The datum $B$ is again physical. Its continuation radius $r_{\rm B}$
is a radius in $e^{iB}$, chosen independently of the coupling.
For a gauge-invariant local interaction $E$, the first derivative at
the trivial pair vanishes: invariance under constant gauge
transformations makes it an invariant linear functional on the
semisimple Lie algebra, and that functional is zero. Consequently

\refstepcounter{equation}
\begin{gather*}
\FitMath{D^2(E\circ p_j)(0)[B_1,B_2]
 =\operatorname{Hess}E(0)[\mathfrak h_jB_1,
                         \mathfrak h_jB_2].}\tag{\theequation}\label{eq:extraction-radius-ER3}
\end{gather*}

The possible second derivative of $p_j$ contributes nothing because it
is multiplied by $DE(0)$.

The polarisation kernel is the kernel of \eqref{eq:extraction-radius-ER3}. Its second spatial
moments define the curvature tensor $T[E]$, and, in the invariant
scalar channel, the action coefficient $b(E)$. These definitions use
derivatives with respect to $B$. For the Wilson functional of the block
datum itself, use the normalization

\refstepcounter{equation}
\begin{gather*}
\FitMath{b(A_{\rm block})=1,}\tag{\theequation}\label{eq:extraction-radius-ER4}
\end{gather*}

whose quadratic part is $\frac12\|\mathrm{d}B\|^2$. For a Wilson family evaluated
after the background map, the analogous normalization requires the
linear-response identities in that map. Likewise
$T[\mathsf A[c]]=c$ for a constant allowed curvature tensor $c$ is
the prescribed counter-insertion normalization. Neither identity for
pulled-back families is proved merely by the coordinate change below.

In a fluctuation integral one instead sets
$\mathcal A_j=g_j\mathcal A_j^{\rm s}$. If
$\widetilde E_j(B^{\rm s})=E(p_j(g_jB^{\rm s}))$, the chain rule gives

\refstepcounter{equation}
\begin{gather*}
\FitMath{D^2\widetilde E_j(0)=g_j^2D^2(E\circ p_j)(0),} \\[3pt]
\FitMath{T_{\rm s}[E]=g_j^2T[E],} \\[3pt]
\FitMath{b_{\rm s}(E)=g_j^2b(E).}\tag{\theequation}\label{eq:extraction-radius-ER5}
\end{gather*}

Thus a fixed probe in the Gaussian variable may fit inside a shrinking
physical domain. It extracts the rescaled coefficient in \eqref{eq:extraction-radius-ER5}.
Recovering the coefficient of the physical action requires division by
$g_j^2$.

For example, impose the response-decay and localization hypotheses of
the amplitude proposition below, and suppose that the usable physical
radius is $a_j$ and the family norm is at most $M_j$ there. Its radius
in $B^{\rm s}$ is comparable to $a_j/g_j$. The Cauchy radius factor
in the amplitude bound then transforms as
\begin{equation}
\FitMath{|g_j^2 b(E)|\le C M_j(g_j/a_j)^2,
 \qquad |b(E)|\le C M_j a_j^{-2}.}
 \label{eq:extraction-radius-ER6}
\end{equation}
The coupling cancels from the coordinate conversion. The same
calculation holds if the physical radius is written $g_j^\gamma$:
the canonical radius is $g_j^{\gamma-1}$, and the physical coefficient
bound remains $CM_jg_j^{-2\gamma}$. The change of variables creates no
extra loss and removes no loss already present in the physical radius.

\subsubsection{A near-zero ball lies in the hereditary core}

We state precisely the background information needed for this local
inclusion. A core at radii $(a_0,a_1)$ is described by the following
simple conditions and their tower counterparts. A triple is
$(U,A',J)$, with $U$ real and $\mathbf U=e^{i\eta_jA'}U$.
The real and complex curvatures are smaller than $a_0\eta_j^2$;
the real factor has the prescribed local regular gauges;
$|A'|,|\nabla_UA'|<a_1$; and $|J|<a_0$.
The tower counterparts require the corresponding curvature and current
bounds for every allowed restricted tower inside the domain.

By \cref{htc:all-towers}, if the simple conditions
hold at sufficiently small radii $(\epsilon_0,\epsilon_1)$, the tower
configurations exist holomorphically. Their curvature and current,
in the local units of their layer, satisfy the defining tower bounds
with radius

\refstepcounter{equation}
\begin{gather*}
\FitMath{K_{\rm T}(\epsilon_0+\epsilon_1^2).}\tag{\theequation}\label{eq:extraction-radius-ER7}
\end{gather*}

The estimate includes the stated change-of-unit factors and holds
for both tower rules, every restricted domain and every depth.
Its common analytic ball and physical-current bound are proved in
\cref{htc:common-ball,htc:tower-bound}.

\begin{proposition}[Near-zero core inclusion]
For sufficiently small $a_0,a_1$ there is a constant
$c>0$, independent of $j$, such that

\par\noindent
\begin{minipage}{\linewidth}
\refstepcounter{equation}
\begin{gather*}
\FitMath{\|v\|_{1,j}<r_j,} \\[3pt]
\FitMath{r_j=c\min\{a_0,a_1,1\}} \\[3pt]
\FitMath{\Longrightarrow} \\[3pt]
\FitMath{(\mathbf1,A',J)\text{ belongs to the core at }(a_0,a_1).}\tag{\theequation}\label{eq:extraction-radius-ER8}
\end{gather*}
\popQED\par
\end{minipage}

\end{proposition}
\begin{proof}
The plaquette product in the exponential chart obeys

\refstepcounter{equation}
\begin{gather*}
\FitMath{|\partial e^{i\eta_jA'}-\mathbf1|
 \le C\eta_j^2
       (\|\nabla_jA'\|_\infty+\|A'\|_\infty^2)
 \le C\eta_j^2(r_j+r_j^2).}\tag{\theequation}\label{eq:extraction-radius-ER9}
\end{gather*}

This follows by expanding a plaquette product: the linear term is the
discrete curl and each nonlinear term contains at least two bond
potentials. On a fixed small chart the remaining exponential series is
bounded by a uniform multiple of the quadratic term. In particular,
there is no inverse power of $\eta_j$ in \eqref{eq:extraction-radius-ER9}. This is the reason
for including the first difference quotient in \eqref{eq:extraction-radius-ER1}.

Choose

\begin{gather*}
\FitMath{\epsilon_0=\frac{a_0}{4K_{\rm T}},} \\[3pt]
\FitMath{\epsilon_1=\min\{a_1/2,\sqrt{a_0/(4K_{\rm T})}\}.}
\end{gather*}

Taking $c$ smaller if necessary ensures
$C(r_j+r_j^2)<\epsilon_0$ and $r_j<\min\{\epsilon_0,\epsilon_1\}$.
The real configuration is $\mathbf1$, so its curvature and local
gauge potential vanish. All simple conditions therefore hold at these
reduced radii. Moreover,
$K_{\rm T}(\epsilon_0+\epsilon_1^2)\le a_0/2$.
The tower estimate gives every hereditary clause at the desired
radii, and $|J|<r_j<a_0$. This proves \eqref{eq:extraction-radius-ER8}.
\end{proof}

A bound only on $\|A'\|_\infty$ would be different: arbitrary
oscillation on the bare lattice can give
$\|\nabla_jA'\|_\infty$ of order
$\eta_j^{-1}\|A'\|_\infty$. The ball relevant to \eqref{eq:extraction-radius-ER8} is the
$C^1$ ball, not that weaker ball.

For completeness, suppose the continued background in \eqref{eq:extraction-radius-ER2} is
holomorphic for $\|B\|_\infty<r_{\rm B}$ and satisfies

\refstepcounter{equation}
\begin{gather*}
\FitMath{\|\mathcal H_j(B)\|_{C^1}
       +\|J(e^{i\eta_j\mathcal H_j(B)})\|_\infty
 \le C_{\rm B}\|B\|_\infty.}\tag{\theequation}\label{eq:extraction-radius-ER10}
\end{gather*}

Then its triple belongs to the core whenever

\refstepcounter{equation}
\begin{gather*}
\FitMath{\|B\|_\infty<
 \min\{r_{\rm B}/2,r_j/(2C_{\rm B})\}.}\tag{\theequation}\label{eq:extraction-radius-ER11}
\end{gather*}

The same conclusion for a first-order response uses the derivative of
\eqref{eq:extraction-radius-ER10}. \eqref{eq:extraction-radius-ER8}--\eqref{eq:extraction-radius-ER11} give a fixed datum radius when
$(a_0,a_1)$ are fixed, and a radius comparable to
$\min\{a_0,a_1\}$ when they run. The hypotheses are exactly the
background continuation, its uniform first-order bounds and the
hereditary tower estimate. No analyticity radius in the normalized
Gaussian variable is substituted for this physical datum radius.

\subsubsection{Cauchy's estimate, response decay and entropy}

Work first with the infinite-lattice companion. Its torus versions use
the uniform response bounds and the lift convergence supplied by the
background hypotheses. Consider pointed local families satisfying

\begin{gather*}
\FitMath{|E(X,p,z)|\le M e^{-\kappa d(X)}}
\end{gather*}

on the core, and assume that the near-zero ball there has usable
radius $a$. Assume also the uniform response decay

\refstepcounter{equation}
\begin{gather*}
\FitMath{\|\mathfrak h_j\delta_c\|_{1,j;X}
 \le C e^{-\delta\operatorname{dist}(c,X)}.}\tag{\theequation}\label{eq:extraction-radius-ER12}
\end{gather*}

The distances in this section are measured in the unit of level $j$.
For a fixed root cube, suppose the number of domains with tree length
between $m$ and $m+1$ is at most $C e^{\kappa_{\rm D}m}$, while the
number of pointed sites in such a domain is at most $C(m+1)^4$.
These are the usual localization entropy and volume bounds.

\begin{proposition}[The amplitude bound with its radius]
If $\kappa>\kappa_{\rm D}$ with a fixed positive margin,
the polarisation sums converge absolutely and

\refstepcounter{equation}
\begin{gather*}
\FitMath{|\Pi_{\mu\nu}(u)|\le C a^{-2}M e^{-\delta_1|u|},} \\[3pt]
\FitMath{|b(E)|+\|T[E]\|\le C a^{-2}M,}\tag{\theequation}\label{eq:extraction-radius-ER13}
\end{gather*}

whenever the scalar coefficient or the stated summed tensor is defined
by the prescribed covariance conditions. The constants are uniform in
the level and torus volume. A tensor at a fixed base point satisfies
the same bound before any isotropy assertion is made.
\end{proposition}
\begin{proof}
The two-variable Cauchy formula on a bidisc inside the ball gives

\begin{gather*}
\FitMath{|\operatorname{Hess}E(X,z)[v,w]|
 \le Ca^{-2}M e^{-\kappa d(X)}\|v\|_{1,j;X}\|w\|_{1,j;X}.}
\end{gather*}

Combining this with \eqref{eq:extraction-radius-ER12} bounds the contribution of one term by

\refstepcounter{equation}
\begin{gather*}
\FitMath{Ca^{-2}M e^{-\kappa d(X)}
 e^{-\delta(\operatorname{dist}(c,X)+\operatorname{dist}(c',X))}.}\tag{\theequation}\label{eq:extraction-radius-ER14}
\end{gather*}

If the localization cubes have fixed side $M_0$, then

\begin{gather*}
\FitMath{|c_--c'_-|\le\operatorname{dist}(c,X)
          +\operatorname{dist}(c',X)+C M_0(d(X)+1).}
\end{gather*}

A sufficiently small fixed fraction of the exponents in \eqref{eq:extraction-radius-ER14} supplies
$e^{-\delta_1|c_--c'_-|}$ for a fixed
$\delta_1>0$. Choose that fraction so that the remaining localization rate stays
strictly above $\kappa_{\rm D}$. The remainder is summable. Indeed, root a domain
at a cube closest to $c$. Summing its pointed sites costs at most a
polynomial in $m+1$, and summing its shapes costs
$e^{\kappa_{\rm D}m}$. These costs are absorbed by the remaining
$e^{-\kappa_{\rm rem}m}$, where $\kappa_{\rm rem}>\kappa_{\rm D}$.
Summing the root location is bounded by the
remaining exponential in its distance from $c$. This proves the first
inequality in \eqref{eq:extraction-radius-ER13} and absolute convergence.

The second moment defining $b(E)$ is bounded by
$Ca^{-2}M\sum_u|u|^2e^{-\delta_1|u|}$, a finite constant times
$a^{-2}M$. For a tensor at a fixed base point, the two moment factors
are bounded by polynomials in the domain diameter and the two
response distances. Those polynomials are absorbed by a further
fixed fraction of the same exponentials. This proves the tensor
bound. Identifying an invariant tensor with the scalar action
coefficient uses the separate Ward and lattice-symmetry identities;
it is not needed for the amplitude estimate.
\end{proof}

In particular, the extraction norm is uniform on fixed physical
radii. On the running radii

\begin{gather*}
\FitMath{a_j\asymp\theta_j^{-1/2}(\log\theta_j)^q,} \\[3pt]
\FitMath{q\ge0,}
\end{gather*}

the correct general bound has the factor
$a_j^{-2}\asymp\theta_j(\log\theta_j)^{-2q}$.
One can express uniformity instead as
$a_j^2|b(E)|\le C\|E\|$, provided that this normalization is retained
in all later estimates.

\subsubsection{What the squared radius ratios pay for}

Let a remainder born at level $j$ have size $r_j$ on physical radius
$a_j$ and have its vacuum value removed. There is no linear term.
Assume a gauge-invariant window estimate in which the quadratic curvature
factors each scale by $L^{-2(n-j)}$, while a unit volume at level $n$ contains
$L^{4(n-j)}$ sites of level $j$. After these geometric factors cancel,
evaluating the quadratic part on a background of radius $a_n$ costs

\refstepcounter{equation}
\begin{gather*}
\FitMath{w_{n,j}r_j,} \\[3pt]
\FitMath{w_{n,j}=(a_n/a_j)^2.}\tag{\theequation}\label{eq:extraction-radius-ER15}
\end{gather*}

This is the relative quadratic evaluation weight. Extracting the
physical coefficient from the resulting family at level $n$ costs
$a_n^{-2}$, so

\refstepcounter{equation}
\begin{gather*}
\FitMath{a_n^{-2}w_{n,j}r_j=a_j^{-2}r_j.}\tag{\theequation}\label{eq:extraction-radius-ER16}
\end{gather*}

Charging another factor $a_j^{-2}$ on top of the left side would
count the same shrinking input radius twice. Omitting the displayed
readout factor $a_n^{-2}$ would instead confuse an evaluated
potential with its physical quadratic coefficient.
The full nonlinear remainder estimate also requires its actual
backgrounds to lie in the original analytic domains and the separate
control of large domains. \eqref{eq:extraction-radius-ER15} records the radius factor
in that estimate; it provides no continuation outside those domains.

For example, let $r_j\le C\theta_j^{-\kappa_0/2}$,
$\kappa_0>4$, and suppose
$\theta_j\ge\theta_n+\lambda(n-j)-C_0$ with $\lambda>0$ and a
sufficiently large common lower bound on $\theta_n$. Integral
comparison gives

\refstepcounter{equation}
\begin{gather*}
\FitMath{\begin{aligned}
 \sum_{j\le n}w_{n,j}r_j
 &\le C a_n^2
    \sum_{j\le n}\theta_j^{1-\kappa_0/2}
                         (\log\theta_j)^{-2q}
 \le C\theta_n^{1-\kappa_0/2},\\
 a_n^{-2}\sum_{j\le n}w_{n,j}r_j
 &\le C\theta_n^{2-\kappa_0/2}
                       (\log\theta_n)^{-2q}.
 \end{aligned}}\tag{\theequation}\label{eq:extraction-radius-ER17}
\end{gather*}

The first loss of one power of $\theta$ comes from summing the older
levels; the second comes from quadratic extraction. This is the
usual $g^{\kappa_0-4}$ bound, with its logarithmic improvement.
The two losses are already present in \eqref{eq:extraction-radius-ER17}.

% Keep the closing scope paragraph together.
\clearpage
The unmarked chain uses fixed physical radii, so \eqref{eq:extraction-radius-ER13} preserves
its small interaction bound. The marked chain uses running radii,
and its stronger smallness pays the readout loss in \eqref{eq:extraction-radius-ER17}. These
facts establish the radius bookkeeping under the stated background,
response and entropy hypotheses. They do not establish every later
power-counting assertion. Cancellation in a complex window, removal
of every required tensor marginal, weighted sensitivity to earlier
members, and differentiated bounds for a large-field zone remain
their own estimates. Neither the physical-to-Gaussian coordinate
change nor the local core inclusion supplies those conclusions.


% END MODULE chapters/extraction_radius_audit_v3.tex

\section{Couplings, sources and continuum limits}\label{sec:continuum}
The coupling recursion selects the physical scale and the tuned sequence. Uniform source estimates then bound connected correlations at separated supports. We track their dependence on the number of insertions before taking subsequences, and extend the limiting functionals to Schwartz tests flat on the coincidence diagonals. This gives the quantitative test-space estimates needed for \cref{sec:reconstruction}.

% BEGIN MODULE chapters/flow_arithmetic_v2.tex
\subsection{First passage for a marginal flow}
\label{ch:flow-arithmetic}

We derive the numerical consequences of a coupling recursion. Its
realization by a renormalization step, consistency between tori and
cutoffs, and estimates with nonconstant sources remain analytic
hypotheses. An
individual increment of the recursion is allowed to be negative.

\phantomsection\label{notdef:flow-recursion}
\subsubsection{Hypotheses and first passage}

Fix numbers $c_0>0$, $C_0\ge0$, $h>0$, and
$0\le\epsilon<c_0/4$. For each initial value $\beta\ge h+B$ consider a real
run with $x_0=\beta$. Assume that a next step exists whenever $x_j>h$, and
that every such step has the form

\begin{gather*}
\FitMath{x_j=x_{j-1}-b_j,} \\[3pt]
\FitMath{b_j=b_j^0+r_j.}
\end{gather*}

For every interval of defined steps assume

\begin{gather*}
\FitMath{\left|\sum_{l=i+1}^j(b_l^0-c_0)\right|\le C_0,} \\[3pt]
\FitMath{|r_j|\le\epsilon} \\[3pt]
\FitMath{\hbox{if }x_{j-1}\ge h.}
\end{gather*}

All constants are independent of $\beta$ and of the length of the run. The
existence of the next step is an assumption: an inequality about hypothetical
increments cannot establish that an underlying analytic construction is
defined. Set

\begin{gather*}
\FitMath{\lambda=c_0-\epsilon>0,} \\[3pt]
\FitMath{\Lambda=c_0+\epsilon,} \\[3pt]
\FitMath{B\ge c_0+C_0+\epsilon.}
\end{gather*}

The one-step case of the interval estimate gives $b_j\le B$. This is an upper
bound, not an assertion $b_j\ge0$.

Define

\begin{gather*}
\FitMath{K(\beta)=\min\{k\ge0:x_k\le h+B\},} \\[3pt]
\FitMath{\mathcal W_k=\{\beta\ge h+B:K(\beta)=k\}.}
\end{gather*}

The minimum is initially a possible stopping index; its finiteness will be
proved. At the endpoint $\beta=h+B$ it equals zero. When $\beta>h+B$ it is
positive, so the last increment is an actual step of the run.

\begin{proposition}[Block drift and a protected terminal window]
The first-passage time is finite. For defined segments up to that time,

\begin{gather*}
\FitMath{\lambda(j-i)-C_0\le x_i-x_j
                    \le\Lambda(j-i)+C_0} \\[3pt]
\FitMath{(0\le i\le j\le K).}
\end{gather*}

If $K>0$, then $h<x_K\le h+B$. Moreover

\begin{gather*}
\FitMath{\max\left\{0,\frac{\beta-h-B-C_0}{\Lambda}\right\}
 \le K\le\frac{\beta-h+C_0}{\lambda}.}\qedhere
\end{gather*}

\end{proposition}

\begin{proof}
Summing the recursion over a segment gives

\begin{gather*}
\FitMath{x_i-x_j=c_0(j-i)
       +\sum_{l=i+1}^j(b_l^0-c_0)+\sum_{l=i+1}^jr_l.}
\end{gather*}

The two assumed bounds give the block drift estimate. If the run never reached
$h+B$, all its values would remain above $h$, so it could be extended for
every finite number of steps. The lower drift bound would then imply
$x_j\le\beta-\lambda j+C_0$, contradicting $x_j>h+B$ for sufficiently large
$j$. Thus first passage occurs.

For $K>0$, minimality gives $x_{K-1}>h+B$. Since $b_K\le B$,
$x_K=x_{K-1}-b_K>h$. The upper endpoint follows from the definition of $K$.
The extra width $B$ is exactly what prevents the final step from overshooting
the lower reference value $h$.

Finally $\beta-x_K\le\Lambda K+C_0$ and $x_K\le h+B$ give the lower bound
on $K$. Conversely $\lambda K-C_0\le\beta-x_K<\beta-h$ gives the stated
upper bound for $K>0$. The endpoint case $K=0$ also satisfies both inequalities.
\end{proof}

In particular $\beta\to\infty$ forces $K(\beta)\to\infty$. Conversely a
sequence with $K\to\infty$ has $\beta\to\infty$, by the lower drift estimate
and the terminal lower bound. Since these stopping times are integers, the
sets $\mathcal W_k$ are nonempty for infinitely many different $k$. Nothing
here says that every one is nonempty, that one is an interval, or that the
terminal values in it coincide. Such assertions would require additional
information about the dependence of the analytic recursion on $\beta$.

For example, $b_j^0=c_0+(-1)^jq$, with $q>c_0$ and $r_j=0$, satisfies the
interval hypothesis with $C_0=q$. Every other increment is negative.
Nevertheless the block drift is positive and the preceding proof applies.
Replacing block drift by a silently assumed monotonicity would discard this
perfectly admissible numerical example.

\subsubsection{Vanishing mean error and the leading logarithm}

Assume in addition that

\begin{gather*}
\FitMath{|r_j|\le\omega(x_{j-1}),} \\[3pt]
\FitMath{\omega_* (M):=\sup_{x\ge M}\omega(x)\longrightarrow0} \\[3pt]
\FitMath{(M\to\infty),}
\end{gather*}

where $\omega$ is nonnegative on $[h,\infty)$. The earlier global bound by
$\epsilon$ remains in force. No monotonicity of $\omega$ is needed.

\begin{lemma}[Time spent below a fixed inverse coupling]
For $M\ge h$, the number $N_M$ of steps $1\le j\le K$ with $x_{j-1}\le M$
is bounded by

\begin{gather*}
\FitMath{N_M\le1+\frac{M-h+C_0}{\lambda}.}
\end{gather*}

This bound does not depend on the initial value.
\end{lemma}

\begin{proof}
There is nothing to prove if no such step occurs. Otherwise let $i$ be the
smallest index in $\{0,\ldots,K-1\}$ with $x_i\le M$. Every counted step
has its preceding index between $i$ and $K-1$, so $N_M\le K-i$. The block
drift from $i$ to $K$ gives
$\lambda(K-i)-C_0\le x_i-x_K<M-h$. The displayed, slightly weaker bound
therefore follows. Excursions back above $M$ cause no problem: all subsequent
indices are included in this bound on the entire remaining time.
\end{proof}

\begin{proposition}[Leading first-passage asymptotics]
With $h$ and the other constants fixed,

\begin{gather*}
\FitMath{\frac1{K(\beta)}\sum_{j=1}^{K(\beta)}|r_j|\longrightarrow0,} \\[3pt]
\FitMath{\frac{\beta}{K(\beta)}\longrightarrow c_0} \\[3pt]
\FitMath{(\beta\to\infty).}
\end{gather*}

If $a_K=L^{-K}$ for a fixed $L>1$, then
$\beta/\log(a_K^{-1})\to c_0/\log L$.
\end{proposition}

\begin{proof}
Split the steps according to whether $x_{j-1}\le M$. The lemma gives

\begin{gather*}
\FitMath{\frac1K\sum_{j=1}^K|r_j|
 \le\frac{\epsilon}{K}
             \left(1+\frac{M-h+C_0}{\lambda}\right)+\omega_*(M).}
\end{gather*}

First let $\beta\to\infty$ at fixed $M$, using $K\to\infty$, and then let
$M\to\infty$. This proves the first limit. Summing the whole recursion gives

\begin{gather*}
\FitMath{\beta=x_K+c_0K+e_K+\sum_{j=1}^Kr_j,} \\[3pt]
\FitMath{|e_K|\le C_0.}
\end{gather*}

The terminal value belongs to a fixed bounded interval, so division by $K$
proves the second limit. Finally $\log(a_K^{-1})=K\log L$ exactly.
\end{proof}

If the analytic one-loop calculation establishes
$c_0=11N^2\log L/(12\pi^2)$ in the convention $\beta=g_0^{-2}$, this
arithmetic yields $\beta/\log(a_K^{-1})\to11N^2/(12\pi^2)$.
The one-loop coefficient is an input to this implication, not a consequence
of the stopping argument. The result is uniform over any family of runs
satisfying the same constants and error envelope.

\subsubsection{A quantitative optional error bound}

Suppose $h\ge e$, $p\ge0$, and
$\omega(x)\le Cx^{-1/2}(\log x)^p$. Then

\begin{gather*}
\FitMath{|\beta-c_0K|\le h+B+C_0+
 C'\sqrt{\beta+C_0+2}\,[\log(\beta+C_0+2)]^p.}
\end{gather*}

To prove it, partition $[h,\infty)$ into unit intervals. If two preceding
indices $i<j<K$ have values in the same interval, then $x_i-x_j<1$ and
the lower drift gives $j-i<(1+C_0)/\lambda$. A unit interval is therefore
visited at most $J=2+\lceil(1+C_0)/\lambda\rceil$ times. Also
$x_i\le\beta+C_0$ by the lower drift from zero. Summing the supremum of the
error envelope in each visited interval consequently gives

\begin{gather*}
\FitMath{\sum_{j=1}^K|r_j|
 \le CJ[\log(\beta+C_0+2)]^p
        \sum_{n=0}^{\lceil\beta+C_0-h\rceil}(h+n)^{-1/2}.}
\end{gather*}

The last sum is at most a constant times $\sqrt{\beta+C_0+2}$, by comparison
with the integral of $x^{-1/2}$. Substitute this bound in the summed recursion.
For fixed $h$ this is an $O(\sqrt\beta(\log\beta)^p)$ error. It supplies no
second logarithmic coefficient, convergence rate for observables, or
summable comparison between neighboring cutoffs.

\subsubsection{Depth envelopes and sums over scales}

The same block estimate gives useful bounds measured from the terminal
scale. For $0\le s\le K$, apply it between $K-s$ and $K$ to obtain

\begin{gather*}
\FitMath{h+\lambda s-C_0<x_{K-s}
             \le h+B+\Lambda s+C_0.}
\end{gather*}

This estimate records the possible bounded backward excursions without
requiring the inverse couplings to be monotone. It also distinguishes the
fixed reference value $h$ from the terminal value $x_K$, which is allowed to
range over its window.

For an explicit scale-summation corollary suppose
$h\ge\max\{1,2C_0\}$. Then the lower bound becomes
$x_{K-s}\ge h/2+\lambda s$. Write $g_j=x_j^{-1/2}$. For every $q>2$,
comparison with a decreasing integral gives

\begin{gather*}
\FitMath{\sum_{j=0}^K g_j^q
 \le\sum_{s=0}^\infty(h/2+\lambda s)^{-q/2}
 \le(h/2)^{-q/2}
   +\frac{(h/2)^{1-q/2}}{\lambda(q/2-1)}
 \le C(q,\lambda)h^{1-q/2}.}
\end{gather*}

The bound is uniform in the number of steps. The first term of the sum is
bounded separately, and the integral accounts for all later terms; this
explains both the strict condition $q>2$ and the displayed dependence on the
reference inverse coupling.

At the marginal exponent $q=2$ the integral instead grows logarithmically.
Using the two depth envelopes gives, for example,

\begin{gather*}
\FitMath{\frac1\Lambda
 \log\left(1+\frac{\Lambda(K+1)}{h+B+C_0}\right)
 \le\sum_{j=0}^K g_j^2
 \le\frac2h+\frac1\lambda
                     \log\left(1+\frac{2\lambda K}{h}\right).}
\end{gather*}

For the lower bound integrate the decreasing function
$(h+B+C_0+\Lambda s)^{-1}$ from zero to $K+1$. For the upper bound separate
its counterpart's first term and integrate from zero to $K$. These elementary
estimates are often the numerical part of a renormalization argument:
sufficiently high powers of the running coupling are summable, whereas the
marginal power is not. They cannot replace the analytic cancellation or
renormalization that produces those powers in a particular model.

The extra lower bound on $h$ belongs only to this corollary. For a smaller
fixed $h$, separate the finitely many terminal depths $s<2C_0/\lambda$, bound
them using $x_j>h$, and apply the same integral comparison to the remaining
depths. The resulting constants may depend on $h$ and $C_0$.

\phantomsection\label{notdef:lattice-units}
\subsubsection{An exact change of lattice units}

There is also an identity that requires no coupling-flow assumption. Let
$\Lambda_{K,\ell}=(L^{-K}\mathbb Z/\ell\mathbb Z)^4$ be a compatible finite
lattice and let $b$ be an integer with $K-b\ge0$. Dilation $D_bx=L^bx$ maps
it bijectively to $\Lambda_{K-b,L^b\ell}$. Transfer a bond variable by
$U'_{D_b e}=U_e$. The map bijects positively oriented bonds and plaquettes,
preserves their incidence, and satisfies $U'_{D_bp}=U_p$. Product normalized
Haar measure is therefore unchanged, as is the Wilson action
$A(U)=\sum_p[1-\tfrac12\operatorname{Re}\operatorname{tr}U_p]$.

For a test function on $T_\ell$, put $f^{[b]}(x')=f(L^{-b}x')$ on
$T_{L^b\ell}$. Plaquette barycenters obey $x(D_bp)=D_bx(p)$, so

\begin{gather*}
\FitMath{\mathcal O_{K-b,L^b\ell}(f^{[b]})(U')
 =\sum_p f(L^{-b}D_bx(p))a_p(U)
 =\mathcal O_{K,\ell}(f)(U).}
\end{gather*}

No spacing prefactor appears in this observable. Consequently the two Wilson
partition functions at the same $\beta$ are equal, and all joint moments and
cumulants satisfy the exact identity

\begin{gather*}
\FitMath{S^T_{n;K,\ell,\beta}(f_1,\ldots,f_n)
 =S^T_{n;K-b,L^b\ell,\beta}(f_1^{[b]},\ldots,f_n^{[b]}).}
\end{gather*}

This is one finite-dimensional integral written in two units, including the
same number of sites in each direction. It is not a comparison between
different bare measures.

For $\|f\|_\sharp=\max(\|f\|_\infty,c_M\operatorname{Lip}f)$ the exact
norm transformations are

\begin{gather*}
\FitMath{\|f^{[b]}\|_\infty=\|f\|_\infty,} \\[3pt]
\FitMath{\operatorname{Lip}f^{[b]}=L^{-b}\operatorname{Lip}f.}
\end{gather*}

Thus $\|f^{[b]}\|_\sharp\le\|f\|_\sharp$ for $b\ge0$, and it is at most
$L^{|b|}\|f\|_\sharp$ for $b<0$. All support distances and support radii
are multiplied by $L^b$. The same bond bijection carries a fixed-spacing DLR
specification to its dilated specification, since it carries each finite
integration set and all plaquettes meeting that set to their counterparts.
It therefore also transports infinite-lattice Gibbs states and their local
moment identities.

To identify this unit change with two initial portions of a renormalization
run requires more: the blocking rules, gauges, auxiliary choices and
extraction of increments must commute with the bond bijection. Likewise,
torus-independent increments require a separately defined common local
construction. Neither property follows from the scalar recursion or the
exact dilation identity proved here.


% END MODULE chapters/flow_arithmetic_v2.tex

% BEGIN MODULE chapters/06_flow.tex
\subsection{The running coupling and first-passage arithmetic}
\label{candidate:flow}
\subsubsection{The run, its coefficient, and the analytic requirements}
\label{sec:flow-the-run-its-coefficient-and-the-analytic-requirements}

The numerical running coupling and the auxiliary frame have different
roles. This distinction is needed even for a real bare coupling. A frame
fixes the tests, cubes, Gaussian references, coordinate choices and
large-field operations. The weight of the Wilson action is another entry
of the construction. A derivative with the frame fixed differentiates
that weight; it does not differentiate a characteristic function whose
threshold is moving. We first state the actual analytic requirements on
these objects. The arithmetic consequences below are conditional
theorems under those requirements.

Fix odd \(L>1\), the block parameter \(M\), and the remaining geometric
constants. In the ordinary-source application use its actual
\(M=L^{m_0}\), fixed collar layers and compatible cofinal tori;
admissibility below refers to those assigned lattices.
The group is \(\operatorname{SU}(2)\); in the coefficient formulas
we retain \(N=2\). The pointed families, hereditary complex
potential/current cores, midpoint tents, companions and Euclidean list
norms are those of \cref{ch:rg-objects-v2}. In particular a pointed
family of level \(j\) has terms \(E(X,P,z)\), where \(X\) is a connected
localization domain and \(z\) is a unit lattice site in \(X\). Its norm
is a supremum in \(z\), not a sum in \(z\). The inequality
\begin{equation}
 |E(X,P,z)|\le \mathcal A e^{-\kappa d_j(X)}
 \label{eq:flow-RF1}
\end{equation}
is required on its declared core. Summing domains requires the separate
rooted-domain entropy estimate. A companion agrees on every common
nonwrapping lift, with the required collar, for all admissible tori
with the same frame prefix and initial weight. Local analyticity does
not establish that compatibility.
We write \(E^{(j)}\) for a torus family and
\(\widehat E^{(j)}\) for its common infinite-lattice companion.
The coefficient of the run is extracted from the latter.
A regular list member \(G^{(j)}\) has a companion
\(\widehat G^{(j)}\); neither is the additional one-step Gaussian
family \(G_\tau\). The family created at step \(k+1\) has pointed
level \(k+1\), although its creation choices are those of level \(k\).

Write \(P_\circ=(\mathbf1,0)\), and let \(\mathfrak h_\infty\) be the
infinite-volume first response of the background potential/current pair
at \(P_\circ\). Its existence, gauge covariance, affine reproduction
modulo gauge and exponential decay belong to the background hypotheses.
For the composed physical Hessian, average over an orthonormal color
basis \(t_a\), \(1\le a\le3\), and set
\begin{equation}
 \begin{aligned}
 \Pi(x,y)&=\frac13\sum_{a=1}^3\sum_z\sum_{X\ni z}
  D^2E(X,P_\circ,z)
       [\mathfrak h_\infty\delta_x^a,\mathfrak h_\infty\delta_y^a],\\
 b(E)&=\sum_{u\in\mathbb Z^4}\Pi_{\mu\nu}(u)u_\mu u_\nu,
                  \qquad \mu\ne\nu .
 \end{aligned}\label{eq:flow-RF2}
\end{equation}
The second line is used only after translation covariance and independence
of the choice \(\mu\ne\nu\) have been justified. Hessians are in the pair
exponential chart, not in an unspecified coordinate system. Their
directions have norm
\(\max\{\|A\|_\infty,\|\nabla A\|_\infty,\|J\|_\infty\}\).
For \(a'=\min\{a_1,c,a_0/C\}\), a ball of radius \(ca'\) in this
norm must lie in the core. A ball controlled by \(\|A\|_\infty\) alone
does not supply the uniform mesh normalization.

The coefficient and pointed-cancellation theorem in
\cref{ch:pointed-hessian-cancellation} applies under its full
response, summability and symmetry hypotheses. Its consequences used
here are absolute convergence, vanishing zeroth and first moments, and
\begin{equation}
 \begin{aligned}
 \sum_u\Pi_{\mu\nu}(u)u_\rho u_\sigma
 &=b(E)(\delta_{\mu\rho}\delta_{\nu\sigma}
       +\delta_{\mu\sigma}\delta_{\nu\rho}
       -2\delta_{\mu\nu}\delta_{\rho\sigma}),\\
 |\Pi(x,y)|&\le C(a')^{-2}\mathcal A
                      e^{-c|y_--x_-|},
 &|b(E)|&\le C(a')^{-2}\mathcal A .
 \end{aligned}\label{eq:flow-RF3}
\end{equation}
The extraction is linear, holomorphic in admitted parameters and real
on real families. It depends on the total, so two coherent localizations
of the same physical functional have the same coefficient. A constant
total has coefficient zero. At absolute radii the factor \((a')^{-2}\)
is a constant depending on \(M\); at running radii it can cost
\(C x_j(\log x_j)^{-2q}\). Neither the kernel nor its coefficient
depends on which admitted radii were used to estimate it.

\paragraph*{Pointed tensors and their period averages.}
There is a useful distinction between the Hessian of the total and that
of a pointed total. Write \(K_{\hat z}(x,y)\) for the contribution of
the base point \(\hat z\) to \(\Pi(x,y)\), including the color average
and the sum of domains in \eqref{eq:flow-RF2}. Assume here the actual pointed Ward
identities and absolute moments through order two, uniformly over one
period, and covariance under simultaneous translations by \(M\mathbb Z^4\).
These are pointed response hypotheses; the bulk estimate in \eqref{eq:flow-RF3} alone
does not imply them. Set \(\mathcal C_M=\{0,\ldots,M-1\}^4\) and
\begin{equation}
 \begin{aligned}
 T_{\mu\nu,\rho\sigma}(\hat z)
  &=\sum_{\substack{x\ {\rm of\ direction}\ \mu\\
                    y\ {\rm of\ direction}\ \nu}}
       (x_--\hat z)_\rho(y_--\hat z)_\sigma K_{\hat z}(x,y),\\
 \overline T&=M^{-4}\sum_{\hat z\in\mathcal C_M}T(\hat z).
 \end{aligned}\label{eq:flow-RF3a}
\end{equation}
If only the total is invariant under all unit lattice isometries,
the conclusion is
\begin{equation}
 \overline T_{\mu\nu,\rho\sigma}
   =b(E)(\delta_{\mu\nu}\delta_{\rho\sigma}
                        -\delta_{\mu\sigma}\delta_{\nu\rho})
                +c_{\rm top}\epsilon_{\mu\rho\nu\sigma}.
 \label{eq:flow-RF3b}
\end{equation}
In particular an invariant total does not determine \(T(\hat z)\)
at an individual base point.

Here is the algebra behind this assertion. At the trivial field the
first derivative of each scalar pointed total vanishes by constant
color invariance. The twice differentiated gauge identity therefore
makes its Hessian vanish when either argument is a gradient.
One may insert linear and quadratic gauge polynomials by cutoff and
the assumed absolute moments. Linear polynomials give zero row and
column sums in each bond direction. A quadratic polynomial
\((x-\hat z)_\mu(x-\hat z)_\rho\) gives antisymmetry of \(T\) under
\(\mu\leftrightarrow\rho\); the column identity gives antisymmetry
under \(\nu\leftrightarrow\sigma\). For a repeated index the extra
constant from the discrete derivative vanishes by the linear identity.
These same zero sums show that replacing lower bond endpoints by
midpoints does not change \(T\).

To identify its symmetric part, sum the bulk second moment over one
period cell and reindex the \(M\)-translates of \(\hat z\).
Absolute moments justify every interchange. For each fixed relative
bond position there is exactly one such translate with its lower
endpoint in \(\mathcal C_M\). Translation invariance of the bulk
kernel consequently gives
\begin{equation}
 \sum_u\Pi_{\mu\nu}(u)u_\rho u_\sigma
       =-\overline T_{\mu\nu,\rho\sigma}
        -\overline T_{\mu\nu,\sigma\rho}.
 \label{eq:flow-RF3c}
\end{equation}
Indeed the \(x_\rho x_\sigma\) and \(y_\rho y_\sigma\) terms
in \((x-y)_\rho(x-y)_\sigma\) vanish by the row and column zero sums;
the two cross terms give the right-hand side. Subtract the Maxwell
tensor in \eqref{eq:flow-RF3b}, using \eqref{eq:flow-RF3}. The residual is antisymmetric under
the three transpositions \(\mu\leftrightarrow\rho\),
\(\nu\leftrightarrow\sigma\), and \(\rho\leftrightarrow\sigma\).
These generate all permutations of four positions, so it is a
fully alternating tensor. In four dimensions it is a multiple of
\(\epsilon_{\mu\rho\nu\sigma}\), proving \eqref{eq:flow-RF3b}.

If the pointed totals themselves are scalar covariant under every
unit translation and reflection, \(T(\hat z)\) is independent of
\(\hat z\), and one reflection kills its alternating coefficient.
This is the pointed-symmetry case used by the local cancellation
theorem. With only period covariance, a local subtraction instead
requires its actual period/window argument; it cannot be inferred
from the invariant total. The layered estimate below retains these
pointed hypotheses.

The pointed Wilson family assigns the action to the same tents as the
regular member. Its norm is \(Ca_0^2\), so subtracting its own extracted
coefficient preserves a constant multiple of the family norm. This is
the subtraction used below. The stronger statement that the remaining
translation-invariant quadratic kernel is divisible by three lattice
derivatives, with exponentially decaying quotient kernels on a fixed
polyannulus, is a separate analytic target when that representation is
used. The moment identity alone is not presented as a proof of such
division.

\paragraph*{Frame, weight, tube and box.}
The functions used for the auxiliary choices are
\begin{equation}
 \begin{aligned}
 \varepsilon(y)&=A_0y^{-1/2}(\log y)^{\mathfrak p_0},
 &\delta(y)&=A_1y^{-1/2}(\log y)^{\mathfrak p_0},\\
 \delta'(y)&=A_1y^{-1/2}(\log y)^{\mathfrak p_1},
 &\alpha_\iota(y)&=C_\iota y^{-1/2}(\log y)^q,\\
 R(y)&=\min\{L^b:L^b\ge(\log y)^r\},
 &q&>\mathfrak p_0+2r,\quad \mathfrak p_1<\mathfrak p_0 .
 \end{aligned}\label{eq:flow-RF4}
\end{equation}
The fluctuation scaling is \(y^{-1/2}\). The frame also fixes the
gauge-fixing weights, the scalings in the large-field stages and their
main step, the stage Wilson constants, the Gaussian forms, and both
the coefficient and cutoff in the reference laws. Every occurrence of
a coupling is assigned either to such a choice or to the action
weight, never to both. Cube sizes and waiting depths use the running
minimum of the frame.

For a real positive-weight run and a terminal reference \(g>0\),
put \(h=g^{-2}\). The three real scales are
\begin{equation}
 \mathfrak g_j=\theta_j^{-1/2},\qquad
 g_j=x_j^{-1/2},\qquad g=h^{-1/2}.
 \label{eq:flow-RF4a}
\end{equation}
The first scales a fluctuation in the fixed-frame step; the second
is the actual running coupling on a real positive-weight run; the
third is the terminal reference coupling. In a self-calibrated run
\(\theta_j=x_j\), the first two coincide. On an admitted narrow real
tube they are uniformly comparable. Bounds for the creation domains,
cutoffs and remainders use the frame scale; numerical sums along the
actual flow use \(g_j\).

A finite reference sequence \(\theta_0,\ldots,\theta_n\) is real,
above its floor, and obeys
\begin{equation}
 \lambda_\flat(j-i)-c_5\le\theta_i-\theta_j\le\bar B(j-i),
 \qquad
 \check R_j=R(\min_{i\le j}\theta_i),\quad
 \check R_j\le L\check R_{j+1}.
 \label{eq:flow-RF5}
\end{equation}
The lower slope \(\lambda_\flat\) is the same slope used in the
background and history estimates; in the weak triple it is \(c_0/4\).
The frame is a real reference sequence, not an independently varied
list of complex action weights.
In the run along this sequence,
\begin{equation}
 x_0^\theta=\beta,\qquad
 x_{j+1}^\theta=x_j^\theta-b_{j+1}^\theta,\qquad
 z_j=\theta_j-x_j^\theta,\qquad \tau_j=x_j^\theta/\theta_j .
 \label{eq:flow-RF6}
\end{equation}
The frame and weights are separate lists. Put
\(\mathsf t^\sharp(y)=\varpi y(\log y)^{-2q_{\rm t}}\).
A tube has positive, nonincreasing radii \(\mathsf t_j\), with
\(\mathsf t_j\le\mathsf t^\sharp(\theta_{j+1})\) for \(j<n\).
For a one-member sequence require
\(\mathsf t_0\le2\mathsf t^\sharp(\theta_0)\). The constants
\(\varpi,q_{\rm t}\) are fixed before the coupling ceiling, with
\(2\varpi\le\min\{\varpi_0,\sigma_\star\}\) and \(q_{\rm t}\ge q_0\).
The pair is alive through \(k\) if its coefficients through \(k\) are
defined and \(|z_j|<\mathsf t_j\). It has room \(r'\) when
\(|z_j|+r'\le\mathsf t_j\) at every required level.

The box of level \(k\) is the full inductive array: weights, regular
families with companions and coefficients, remainders, boundary terms
and standing large-field operations. It includes the bounds
\(|E^{(j)}|\le E_0e^{-\kappa d_j}\),
\(|R^{(j)}|\le\theta_j^{-\kappa_0/2}e^{-\kappa d_j}\),
the covariance and positivity clauses, identification with the two
chains below, and the drift bounds. Its reserve constant \(c_{\rm m}\),
also denoted \(c_{\rm r}\), enters the large-field mass bound; increasing
it shrinks the box and does not change the step. An admissible torus
must satisfy
\begin{equation}
 \tfrac12\ell L^{K-i}\ge C(L)M\check R_i\qquad(i\le k).
 \label{eq:flow-RF7}
\end{equation}
All box statements are simultaneous statements on these tori and on
their common companions. The positivity of a label measure
\(\tau_{\mathcal S}=\sum_{\mathcal S(h)=\mathcal S}\tau_h\)
must not be confused with the numerical ratio \(\tau_j\).

The history part of each output box has a stronger requirement than
small real masses. For the next operation write \(i=k+1\), and let \(X\)
be a component
in either of the two classes with a pair-core representation at that
level. Use its actual local histories
\(\omega\), positive reference integration laws \(\mu_\omega\),
Borel supports \(S_\omega\), and exponents \(\Phi_\omega\).
All level normalizations remain in these assigned laws. For a
nonnegative integrand \(f\), the upper integral is
\begin{equation}
 \int^*f\,d\mu
   =\inf\left\{\int F\,d\mu:
               F\text{ is Borel and }F\ge f\right\}.
 \label{eq:flow-RF7b}
\end{equation}
It is used because the core supremum need not have been proved Borel.
The required complex mass clause is
\begin{equation}
 \frac1{|W_4|}\sum_\omega e^{c_{\rm m}v^\natural(\omega)}
   \int_{S_\omega}^{*}
      \sup_{\mathbf p\in U_i^c(X^\sim)}
           |e^{\Phi_\omega(\xi;\mathbf p)}|\,
                    \mu_\omega(d\xi)
       \le e^{-a_\star p_0(\mathfrak g_i)} .
 \label{eq:flow-RF7a}
\end{equation}
Here \(W_4\) is the finite hypercubic symmetry group,
\(v^\natural\) is the assigned history volume, \(U_i^c(X^\sim)\)
is the declared complex pair core, and \(p_0\) is the fixed positive
large-field logarithmic penalty. The history measure and integration
coordinates are fixed before taking the parameter supremum.
This is an integral of a supremum over the core; a pointwise real bound
or a small pointwise polymer logarithm does not prove it.

The actual split is
\(\Phi_\omega=\Phi_\omega^\circ-\frac12q_X+V_X\).
At real weights \(\Phi_\omega^\circ\) is real, Borel and independent
of the pair.
The reference quadratic form \(q_X\) may depend on the real factor of
the triple presenting the pair. Only the combination
\(-\frac12q_X+V_X\) is a function of that pair. A proof must retain this
tie. An extra perturbation is absorbed by the reserve only after its
actual support estimate
\(|\Psi_\omega|\le\bar f\,n^\sharp(\omega)\) and
\(c_{\rm m}\mathsf l_j\ge2\bar f\) have been verified at every relevant
level; \(n^\sharp\) is the assigned site count and \(\mathsf l_j\)
its volume-comparison weight.

Drop-outs have a separate requirement, with their assigned supports
\(\Xi_\omega\) and compatible evaluation fibers
\(\mathfrak K^\sharp(\omega,\xi)\):
\begin{equation}
 \frac1{|W_4|}\sum_\omega e^{c_{\rm m}v^\natural(\omega)}
   \int_{\Xi_\omega}^{*}
       \sup_{\mathfrak e\in\mathfrak K^\sharp(\omega,\xi)}
          |\mathsf E_\omega(\xi;\mathfrak e)|\,
                          \mu_\omega(d\xi)
       \le e^{-a_\star p_0(\mathfrak g_i)} .
 \label{eq:flow-RF7c}
\end{equation}
The majorant \(\mathsf E_\omega\) is the one assigned to that
evaluation, not a new product-core extension. When its own
\(1-\chi\) test is counted, the required gain is instead
\(e^{-2p_0(\mathfrak g_i)}\), with its root counted as an additional
event. A later label must read an evaluation in this same fiber at
every point of its support. For a later occurrence of a second-class
component, its stored variables must lie in \(S_\omega\) and its
later real pair in \(U_i^c(X^\sim)\). Those read-out statements are
additional geometric requirements. Complex quadratic forms in a
drop-out evaluation are admitted only with the specified positive
lower bound on their real part.

Complex variation of Gaussian ratios and stage constants additionally
requires the actual inner Wilson factor to retain its assigned weight.
Discarding that factor, changing its weight without a new comparison,
or bounding the two tied summands separately does not establish \eqref{eq:flow-RF7a}.
The distinction is also relevant to the gauge convention at seams.
If inward seam bonds are retained, an absolute contour gauge is
invariant only under the corresponding seam-tied site gauges. A
relative contour gauge uses its declared boundary reference. An
omission convention that removes every bond with an endpoint in the
next region can eliminate these seams, but its exact operation and
compensators must then be checked. These are history-formation
requirements, separate from the scalar flow arithmetic.

The sourced construction also fixes what a read-out means. At a
formed level \(j=K-u\), set
\begin{equation}
 \begin{aligned}
 \mathsf D^w_{[u]}(Y)
    &=\sum_{z\in Y}E^{(j),w}(Y,P_\circ,z)
                            +R^{(j),w}(Y,P_\circ),\\
 e_{[u]}(Y;w)&=\mathsf D^w_{[u]}(Y)-\mathsf D^0_{[u]}(Y).
 \end{aligned}\label{eq:flow-RF7d}
\end{equation}
The zero-source run in this difference starts from the same real
center and frame. The terminal read-out is the actual number
\(\mathsf I_k(w)\) at \(k=K-s_\circ\), where \(s_\circ\)
is the stop depth, and the exact normalized terminal factor is
\(\widetilde Z^\circ(w)=\mathsf I_k(w)/\mathsf I_k(0)\).
The common vacuum convention counts the removed Gaussian numbers
once. Neither these read-out identities nor the terminal ratio
supplies the complex local-energy or history majorants.

\paragraph*{The Gaussian and one-loop requirements; statement 3.3.}
Let \(\Phi_l=\log Z^{(l-1)}-\log Z^{(l-1)}(\mathbf1)\).
The localized Gaussian family is required to satisfy
\begin{equation}
 \Phi_l(V)=\sum_X\sum_{z\in X}N^{(l)}(X,P_V,z),\qquad
 |N^{(l)}|\le E_{\mathsf N}e^{-\delta_{\mathsf N}Md_l(X)},
 \qquad b(N^{(l)})=b_l^0 .
 \label{eq:flow-RF8}
\end{equation}
The constants depend only on \(L\). No coupling occurs in these
terms. Nonwrapping terms agree on all common lifts, and their pointed
totals are covariant under every unit lattice isometry.

For precision, the localization meant here is the following actual
one. Set \(k=l-1\), let \(A_Y=C^{(k)}(Y)^{-1}\) be the supplied
fluctuation form constructed inside \(Y\) and compressed to free bonds
at depth \(cM\), and define
\begin{equation}
 \begin{aligned}
 f_b(Y)&=\operatorname{tr}(\log A_Y)(b,b),\\
 E_{\mathsf N}(X;b)
   &=\sum_{\substack{Y\subset X\\b\ {\rm at\ depth}\ cM\ {\rm in}\ Y}}
                      (-1)^{|X|-|Y|}f_b(Y).
 \end{aligned}\label{eq:flow-RF9}
\end{equation}
The free-bond contribution to \(N^{(l)}\) is
\(-[E_{\mathsf N}(X;P,b)-E_{\mathsf N}(X;P_\circ,b)]/2\),
smeared by the midpoint tents. The pivot contribution is the logarithm
of its actual elimination factor minus its value at \(P_\circ\), smeared
by the same tents. The required spectral gap and complex continuation
of \(A_Y\), the locality of its walk expansion and the \(c/M\) radius
ceiling are part of this Gaussian-localization input. Finite Gaussian
integration estimates alone do not prove them.

The distinction between an on-shell pair and a critical pair matters.
On shell, the explicit current form expresses gauge covariance.
Identifying it with the Hessian on a constrained fibre also requires
the multiplier equation and actual criticality. The four tempered
localization properties used in the one-loop analysis are its sum
formula, common lifts, exponential domain decay on the shrinking
supremum polydisc of radius \(c\eta_l^2\), and invariance under constant
gauge transformations. Uniform Hessian Cauchy estimates instead use
the pair-direction norm preceding \eqref{eq:flow-RF3}.
For the canonical proper fluctuation precision,
\cref{pgc:section} supplies the physical identification, uniform
complex spectral gap and complete local inverse and logarithm
families used in \eqref{eq:flow-RF9}. The operator is explicitly
\eqref{pgc:actual-precision}: its inner correction uses the first
lift, while its outer constraint curvature uses the second lift.
The two readouts agree at a physical critical pair. Its relative
contour quadratic form agrees with the absolute contour Hessian
in the averaged-contour axial representative. An independent
signed or counterterm-modified matrix requires its own comparison;
zero-fold averaging alone would not justify one.

Input B3 requires, with the fixed momentum convention
\((2\pi)^{-4}dq\),
\begin{equation}
 \begin{aligned}
 c_0&=\frac{11N^2}{12\pi^2}\log L,\qquad
 \left|\sum_{l=1}^k b_l^0-kc_0\right|\le c_2,\\
 |b_{l+1}^0-b_l^0|&\le C_\alpha L^{-\alpha l}
                         \quad(0<\alpha<1).
 \end{aligned}\label{eq:flow-RF10}
\end{equation}
Thus \(b_l^0\to c_0\). The standard convention is
\(\beta=2N/g_{\rm std}^2\), giving the coefficient
\(11N/(24\pi^2)\) for \(g_{\rm std}^{-2}\) per factor \(L\).
The numerical value in \eqref{eq:flow-RF10} requires the actual lattice trace
calculation. It is not deduced by quoting a continuum perturbative
coefficient.

The one-loop calculation has three steps: compose the actual Gaussian
integrals, compare their averaging constraints with a positive mass,
and evaluate the massive lattice coefficient. The proofs below concern
the constructed homogeneous physical background. The complex-current
localization requirements in \eqref{eq:flow-RF9} remain separate.

A tempered functional is real analytic near the trivial field on
arbitrarily large tori, gauge and lattice-isometry invariant,
normalized at that field, and has a plane-wave Hessian per unit volume
converging uniformly in momentum to a bounded holomorphic symbol on a
fixed strip. For such functionals the background substitution is
required to preserve the coefficient. If \(\widehat\pi^F\) is the
limiting bond Hessian symbol with lower-endpoint Fourier convention,
its midpoint form obeys
\begin{equation}
 \begin{aligned}
 e^{i(p_\mu-p_\nu)/2}\widehat\pi^F_{\mu\nu}(p)
  &=b[F](|p|^2\delta_{\mu\nu}-p_\mu p_\nu)+O(|p|^4),\\
 \widehat\pi^{F\circ p_k}(p)
  &=L^2\widehat\pi^F(p/L)+O(|p|^3).
 \end{aligned}\label{eq:flow-RF11a}
\end{equation}
The second formula uses the actual background response and uniform
analytic control of its aliased modes. Its factor
\(L^2=L^{-1}L^{-1}L^4\) records two form scalings and the volume ratio.
The first formula follows from the moment identity and midpoint
lattice symmetries.

For the actual Gaussian densities, coefficient preservation is proved
in \S\ref{abc:section}. The first jet of each invariant density
vanishes, so its pullback Hessian uses only the first background
response. That response reproduces an affine field modulo a gradient;
the two-leg Ward identity removes the gradient, and the physical
density normalization supplies the scale cancellation above.
The successive maps use the actual nested criticality identity;
see \eqref{abc:successive-coefficient-preserved}.

The constrained Gaussian logarithms satisfy
\(\Psi_k=\Phi_k+\Psi_{k-1}\circ p_k\) by
\cref{pgg:actual-stage-composition}. This identity uses the complete
constrained Hessian, including the curvature of the nonlinear
constraint, before integrating in stages. Coefficient preservation
then gives \(\sum_{l\le k}b_l^0=b[\Psi_k]\).
The finite Schur representation uses
\begin{equation}
 \mathbb S(X,Q)=
 \begin{pmatrix}X&Q^{\mathsf T}\\Q&0\end{pmatrix}.
 \label{eq:flow-RF11}
\end{equation}
Here \(Q\) is the covariant form average, \(R\) is the covariant site
average, \(D\) is the covariant gradient, \(\Delta=D^{\mathsf T}D\)
is the site Laplacian, \(\Delta^{\rm W}\) is the Wilson-action
Hessian on forms, and
\(\mathcal K=\Delta^{\rm W}+DD^{\mathsf T}\),
\(\Xi=\mathcal K-DR^{\mathsf T}RD^{\mathsf T}+Q^{\mathsf T}Q\).
On the admitted real critical backgrounds with nondegenerate physical
Gaussians, \cref{pgg:physical-saddle} gives the representation
\begin{equation}
 \begin{aligned}
 \Psi_k={}&-\tfrac12\log|\det\mathbb S(
       \Xi-Q^{\mathsf T}Q+\mathfrak j[J,\mu],Q)|\\
 &+\log|\det\mathbb S(\Delta,R)|
          -(\text{the same expression at }\mathbf1).
 \end{aligned}\label{eq:flow-RF12}
\end{equation}
Here \(U=U_k(V)\) is the actual critical background,
\(J=\nabla_{\rm fine}A(U)=J_{\eta_k}(U)\), and
\(\mu=\nabla_{\rm block}(A\circ U_k)(V)\).
The convention is the bilinear color pairing
\(\langle B,C\rangle=\frac12\operatorname{tr}(BC)\);
fine-bond inner products have weight \(\eta_k^4\), and
block-bond inner products have unit weight. Explicitly,
\begin{equation}
 \begin{aligned}
 \left.\frac d{dt}A(e^{i\eta_k tB}U)\right|_0
     &=\eta_k^4\sum_b\langle J(b),B(b)\rangle,\\
 \left.\frac d{dt}(A\circ U_k)(e^{itC}V)\right|_0
     &=\sum_c\langle\mu(c),C(c)\rangle .
 \end{aligned}\label{eq:flow-RF12a}
\end{equation}
All gradients and adjoints in \eqref{eq:flow-RF12} use those fixed conventions.
The form \(\mathfrak j[J,\mu]\) is fixed by second variation of the
actual constrained background problem. To identify it, let \(q(a)\)
be the recentered logarithmic form of the straightened constraint at
\(U\), with \(q(0)=0\) and \(Dq(0)=Q\). The constrained Hessian
on \(\ker Q\) is
\[
        A''(U)-\langle\mu,D^2q(0)\rangle;
\]
this is \eqref{pgg:actual-current-form}. It is this fibre Hessian that
enters the saddle determinant.

The third-order momentum statement concerns the determinant
correction. Let \(C_k\) be the difference of the normalized vector
logarithms with and without the multiplier term, evaluated at the
actual \(\mu=\nabla_{\rm block}(A\circ U_k)\). Then
\cref{pgg:current-zero-coefficient} proves
\[
          \widehat\pi^{C_k}(p)=O_k(|p|^3),\qquad b[C_k]=0 .
\]
In that proof the mixed field--multiplier Hessian has order \(|p|\)
by the gauge Ward identity, while the derivative of the actual
multiplier has order \(|p|^2\). Their composition has no quadratic
coefficient. This does not give a pointwise third-order bound for the
multiplier matrix. For complex continuation the real absolute
determinants in \eqref{eq:flow-RF12} are replaced by coherent logarithmic
branches on a domain where the constrained forms remain nondegenerate.
The physical identity does not define an independent-current
extension.

The mass calculation uses finite determinants before any
thermodynamic limit. Work here on the unit bare lattice. For an
expanding box \(\Lambda_R=[-R,R]^4\cap\mathbb Z^4\), let
\(\iota_R^0\) and \(\iota_R^1\) be zero extension from its sites
and its positive bonds with both endpoints in the box. Put
\begin{equation}
 \begin{aligned}
 U_f(x,1)&=\exp(-ifx_2t_3),&
 U_f(x,\nu)&=I\quad(\nu\ne1),&
 \tfrac12\operatorname{tr}(t_3^2)&=1,\\
 \Delta_{f,R}&=\iota_R^{0*}\Delta_{U_f}\iota_R^0,&
 \mathcal K_{f,R}&=\iota_R^{1*}\mathcal K_{U_f}\iota_R^1 .
 \end{aligned}\label{eq:flow-RF13a}
\end{equation}
These are Dirichlet compressions of the scalar and vector
operators appearing in \eqref{eq:flow-RF12}, before its averaged-constraint
corrections. The vector Wilson Hessian is defined by second
variation against finitely supported bond fields, so no
infinite total action is differentiated. The oriented
\(12\)-plaquette is \(\exp(if t_3)\). Adjoint charged colors
therefore carry fluxes \(2f\) and \(-2f\), while the third
color is neutral.

For each fixed \(z>0\) and finite box, take the real logarithms
near \(f=0\), where the massive matrices are positive, and set
\begin{equation}
 \begin{aligned}
 \mathcal D_R(f,z)
   &=\log\det(\Delta_{f,R}+zI)
           -\tfrac12\log\det(\mathcal K_{f,R}+zI),\\
 \widetilde\beta(z)
   &=\lim_{R\to\infty}\frac1{|\Lambda_R|}
           \left.\frac{\partial^2}{\partial f^2}
                         \mathcal D_R(f,z)\right|_{f=0}.
 \end{aligned}\label{eq:flow-RF13b}
\end{equation}
\Cref{mlc:main} proves existence of this limit and independence of the
lattice F\o lner exhaustion with the stated zero-extension compression.
The identification with the bulk action coefficient uses the actual
affine response and physical mass rescaling in
\eqref{abc:mass-coefficient}. The order of the limits is the derivative
at fixed positive mass, then volume, then \(z\downarrow0\); no infinite
determinant is differentiated.

The regulator comparison requires the finite-volume normalized Hessian
kernels of this determinant pair to converge to the translation-invariant
quadratic form denoted \(\mathcal L(z)\). The bounded-remainder comparison
concerns
\(\mathcal L^\sharp=\mathcal L(1)-\mathcal L'(1)+
\mathcal L''(1)/2\), uniformly on the prescribed external-momentum strip.
The resulting formulas for \(N=2\) are
\begin{equation}
 \begin{aligned}
 \widetilde\beta(z)&=\frac{11}{6\pi^2}\log(1/z)+O(1),\\
 |z\widetilde\beta'(z)|+|z^2\widetilde\beta''(z)|&\le C,\\
 \sum_{l=1}^k b_l^0
 &=\widetilde\beta(\eta_k^2)
   -\eta_k^2\widetilde\beta'(\eta_k^2)
   +\tfrac12\eta_k^4\widetilde\beta''(\eta_k^2)+O(1).
 \end{aligned}\label{eq:flow-RF13}
\end{equation}
All field expansions are at fixed positive mass before taking that
mass to zero. The vector form need not be nonnegative without its
mass; the small-field hypothesis must make the massive form coercive.

The first two lines follow from \cref{mlc:main}. The third is
\eqref{mrc:actual-RF13-comparison}, combined with Gaussian composition
and coefficient preservation. Its uniform \(O(1)\) term comes from
\cref{prop:mrc-uniform-regulator-comparison}, which includes the
averaged constraints and the gauge determinants. After the
zero-coefficient multiplier correction has been removed, the
difference of the two Hessian symbols is bounded on a fixed complex
momentum neighborhood. If its bound is \(C\) on
\(\{|p_\nu|<\delta\}\), Cauchy's estimate bounds each mixed second
derivative at zero by \(C\delta^{-2}\). Both \(C\) and \(\delta\) are
independent of the depth. Thus the error compares the full cumulative
normalizer with the massive determinant; it is not a sum of one
bounded error per scale.

The stronger geometric convergence uses the Taylor triple
\(\mathfrak J_k=(\Delta_k,S_3^{(k)},S_4^{(k)})\) of
\(\mathcal S_k=A\circ U_k\). With
\(\Lambda(B,\upsilon)=S(B)+\langle\upsilon,m(B)\rangle\),
\(\mathbb G=(D^2\Lambda(0))^{-1}\),
\(V_i=D^i\Lambda(0)\) and
\(\dot z=\mathbb G(0,C)=(HC,-\Delta_{k+1}C)\), the identity is
\begin{equation}
 \begin{aligned}
 \Phi_{k+1}''(C,C')={}&\tfrac23\langle C,C'\rangle
                 -\tfrac23\langle HC,HC'\rangle\\
 &-\tfrac12\operatorname{Tr}(\mathbb G V_4[\dot z,\dot z'])
 +\tfrac12\operatorname{Tr}
       (\mathbb G V_3[\dot z]\mathbb G V_3[\dot z']).
\end{aligned}\label{eq:flow-RF14}
\end{equation}
Here \(\eta_k=L^{-k}\) in the mass comparison,
\(S(B)=\mathcal S_k(e^{iB})+\mathrm{GF}_k(e^{iB})\),
\(m(B)=i^{-1}\log\mathcal M(e^{iB})\),
\(\Delta_k=\mathcal S_k''(\mathbf1)\), and
\(S_j^{(k)}\) is its degree-\(j\) Taylor form in the same chart.
The vector \(HC\) is the linear constrained background response.
\Cref{prop:gpl-actual-formula} derives \eqref{eq:flow-RF14} with the
actual constraint and gauge vertices.
\Cref{prop:cac-Taylor-triple-convergence} proves
\(\|\mathfrak J_{k+1}-\mathfrak J_k\|\le C_\alpha L^{-\alpha k}\)
in the stated weighted kernel norms, and
\cref{prop:gpl-polarization-Lipschitz} supplies the uniform continuity
of the Gaussian polarization in those norms. Their combination is
\cref{cor:gpl-geometric-increments}, which proves the second line of
\eqref{eq:flow-RF10} on the homogeneous physical germ. The bounded
cumulative law alone would not give convergence of individual
increments. These estimates preserve the geometric comparison used
at equal depth in the stronger matched-cutoff route.

\paragraph*{The two-coupling Gaussian factor.}
At level \(k\), write the fibre form as
\(\Delta=\Delta_A+\Delta_{\rm GF}\). The frame fixes the field scaling
and gauge-fixing weight; the action weight supplies
\(\tau\Delta_A+\Delta_{\rm GF}\). The actual finite-dimensional
hypothesis is that \(\Delta_A\) vanishes on the block gauge directions
and \(\Delta_{\rm GF}\) acts only on those directions, with a positive
nondegenerate form on each complementary integration space. Under this
splitting the two ordinary Gaussian integrals factor, giving
\begin{equation}
 \begin{aligned}
 \int dB\,\delta(QB)
 e^{-\langle B,(\tau\Delta_A+\Delta_{\rm GF})B\rangle/2}
     &=\tau^{-n_\perp/2}Z^{(k)}(V),\\
 \int d\mu_{\tau^{-1}C^{(k)}}(B)
 e^{-(1-\tau)\langle B,\widehat G B\rangle/2}
     &=\tau^{n_{\mathcal O}/2},\\
 n_{\mathcal O}&=3\#(\mathbb T^{(k)}\setminus\mathbb T^{(k+1)}),
 &n_\perp&=3|\mathfrak X_k|-n_{\mathcal O},\quad l=k+1 .
 \end{aligned}\label{eq:flow-RF15}
\end{equation}
Here \(Q\) is the one-step average on the fibre, not the \(k\)-fold
average, \(\mathfrak X_k\) is the free-bond set, and
\(\widehat G\) is the gauge-fixing form in those coordinates.
For positive real \(\tau\), the formula follows by rescaling only the
non-gauge coordinates. Analytic continuation to the right half-plane
uses the branch of \(\log\tau\) real on the positive axis. Both
dimensions are independent of the background. In particular
\(n_\perp\) is a dimension, while \(d_l(X)\) is a domain metric.
The two quantities are never denoted by the same unqualified symbol.

Consequently the total of the extra Gaussian family \(G_\tau\) is
\(n_{\mathcal O}\log\tau/2\). Its physical second derivatives vanish
and \(b(G_\tau)=0\). Its norm may nevertheless be
\(K_{\mathbf M}|1-\tau|e^{-(1+c)\kappa d}\), with no small-coupling
gain. This cancellation is needed to avoid a spurious accumulation of
slopes of order \(\theta_k^{-1}\). The formula in \eqref{eq:flow-RF15} presupposes the actual
gauge splitting; it is not a consequence of a bound for a general
complex Gaussian.

The complete regular family has the form
\(E^{(k+1)}=N^{(k+1)}+G_{\tau_k}+F\).
Its physical step functional, up to field-independent normalizations,
is
\begin{equation}
 \begin{aligned}
 \mathcal E^{(k+1)}[\theta_k;x;S](V)
 ={}&\log\!\left(
  \mathbf T_k^{\rm sf}[\theta_k]\,
       e^{-xA(U_k(\cdot))+S(U_k(\cdot))}\right)(V)\\
 &+xA(U_{k+1}(V))-S(U_{k+1}(V)).
 \end{aligned}\label{eq:flow-RF16}
\end{equation}
The operator uses the whole torus as its small-field region, and \(S\)
is the action of the incoming list. Equality between this physical
functional and localized totals is asserted on their admitted
backgrounds. It is not an identity of a global Gaussian at arbitrary
independent-current pairs. The local formulas on such pairs have their
own complex cores.

\paragraph*{Lists, the two chains, and the increment as a function.}
A list contains regular families \(G^{(j)}\) with their own action
counterterms and unpointed remainders \(H^{(j)}\). At the output radii
\(\bar a_\iota\) its required load is
\begin{equation}
 \begin{aligned}
 \mathcal N(\mathfrak f)
  &=\sum_{j=1}^k\omega_{k+1-j}\|G^{(j)}\|
    +\sum_{j\ {\rm aged}}[\bar a/a_j]^2\|H^{(j)}\|,\\
 \omega_n&=(n+1)^{3/2}L^{-\alpha_{\rm H}(n-1)},\\
 |F(X)|&\le K_{\mathbf M}
   \{\delta(\theta_k)+\mathsf s\,\mathcal N(\mathfrak f)
                 +e^{-c\theta_k\delta(\theta_k)^2}\}
                  e^{-(1+c)\kappa d(X)} .
 \end{aligned}\label{eq:flow-RF17}
\end{equation}
The radius ratio is the maximum of the two coordinate-radius ratios;
\(\mathsf s=\delta(\theta_k)/\bar a_0\).
The first, second and third terms pay respectively for the interaction
beyond quadratic order, old-family differences, and failed Gaussian
tests. The parameter choices used to infer a power bound for the last
term must actually satisfy that power bound.

The fluctuation variable \(B\) ranges over nonpivot bonds. Its actual
fibre map \(B\mapsto V'(B)\) satisfies \(\mathcal M(V'(B))=V\);
the full fluctuation exponent \(B'(B)\), including its pivot
completion, is analytic and vanishes at \(B=0\). It is not identified
with the free array \(B\). Old-family amplitude estimates have to hold
on this actual map and its declared joint complex domains.

Likewise, in a localized Gaussian deformation with protected collar
\(Z_0^\sharp\), the independent coarse parameters \(s_\Delta\) belong
only to cubes outside that collar. The collar is a fixed-width cover
and neighboring layers, not a rectangular box filled around a thin
domain. Removing the protection can destroy invertibility of the
compressed form. The unmarked potential-ball theorem applies after
its actual matrix comparison, enlarged parameter discs, hereditary
filing and smallness conditions are verified. It does not construct
those domains. Local bounds on independent-current cores do not
become global Gaussian identities away from complete physical pairs.

The regular part has a summable age weight because its own affine
quadratic action is removed before localization. A remainder has no
such counterterm, so its squared-radius load is not given an age decay.
\Cref{ch:pointed-window-remainder,ch:finite-source-insertion,ch:renewed-source-pointing,ch:window-potential-formation,ch:linear-source-bootstrap} prove the corresponding finite
implications under their actual background, bond-split, potential,
core and boundary hypotheses. Their general polynomial age factors
do not automatically give the exact exponent \(3/2\) in \eqref{eq:flow-RF17}.
One must either verify that sharper estimate or allocate a smaller
positive exponential rate and prove the resulting comparison.

A Euclidean list has one companion, the correct coefficients of its
regular members, pointed totals covariant under the admitted unit
lattice isometries, and invariant remainder totals on the declared
covariance class. Its radii of covariance obey
\begin{equation}
 e_{j'}\le L^{2(j'-j)}e_j\qquad(j\le j').
 \label{eq:flow-RF18}
\end{equation}
Closure means that the actual output again has these properties.
Restriction to smaller cores does not prove covariance. The
backgrounds of fluctuating fields must first lie in a common
covariance class on an open real fluctuation ball; continuation then
requires a connected analytic domain on which both expressions exist.
The actual domain and naturality hypotheses remain part of closure.
For an isometry \(\mathsf r\), the Gaussian covariance assertion uses
the actual orthogonal change of free coordinates \(\mathsf O_{\mathsf r}\)
with
\(V'_{\mathsf rW}(\mathsf O_{\mathsf r}B)=\mathsf rV'_W(B)\).
This is not a walk-expansion operator. Translation and cover locality
of the supply rules do not by themselves give the full isometry
covariance of the new family: that also uses the Euclidean covariance
of the entire incoming list and its aged remainders.

%
For the remainder total, use the all-small-field
term on its invariant open real window. The
positive step kernels carry the term to its
image under the lattice isometry. The localized
density format is continuous on this window,
so equality of the measures gives equality of
the densities at the admitted backgrounds.
Divide by the regular factors and inherited
remainder factors whose totals are already
invariant on the two background windows of
the step. The remaining exponential is that
of the new remainder total. It is real there,
so equality of the exponentials gives equality
of the totals. This is the induction at the
actual backgrounds for the unsourced run and
for invariant real source data. On an admitted
connected complex family with the same density
identity, the difference of the two totals is
continuous with values in \(2\pi i\mathbb Z\).
It vanishes at the real
anchor and hence throughout the family. For a
holomorphic constant source, equality on the
real source interval also gives this continuation.

At the age gate the next lattice's isometries
preserve the old partition, so the group
available for the total includes all the
isometries required at the insertion step.
The termwise translation covariance uses the
translations preserving the actual partition.
The companion and matching index sets are
the other inputs for inserting the family in
the regular list.

The unmarked chain uses absolute radii and starts with the empty list.
The marked chain uses running radii
\(\check\alpha_{\iota,j}=\min\{\alpha_\iota(\theta_{j-1}),
\alpha_\iota(\theta_j)\}\) and inserts a remainder only when
\(M\check R_j\le L^{m+1-j}\) at step \(m+1\).
This age is distinct from the waiting depth of a large-field operation.
Before that age the remainder lies outside the regular exponent and
contributes to the remainder and boundary outputs.

Restriction from absolute to running radii must restrict the same
actual functions. Under that compatibility the marked difference
cancels \(N\) and \(G_\tau\). The required recurrence is
\begin{equation}
 \begin{aligned}
 u_{m+1}&\le K_{\mathbf M}\mathsf s(\theta_m)
       \left(\sum_{i=1}^m\omega_{m+1-i}u_i+
                         C_{\mathfrak h}\theta_m^{-p}\right),\\
 p&=(\kappa_0-2)/2>1,\\
 \sum_{j\le m}\frac{\alpha_0(\theta_m)^2}
                       {\check\alpha_{0,j}^2}
              \theta_j^{-\kappa_0/2}
       &\le C_{\mathfrak h}\theta_m^{-p},\\
 C_{\mathfrak h}&\le C\left(1+
                         \frac1{\lambda_\flat(p-1)}\right).
 \end{aligned}\label{eq:flow-RF19}
\end{equation}
The lower flow bound pays one inverse-coupling power in this sum.
With a constant frame its load can grow with the number of levels.
Smallness of \(K_{\mathbf M}\mathsf s\), relative to the summable
kernel and \(C_{\mathfrak h}\), closes a weighted supremum induction.
This is an actual marked-chain requirement, not a conclusion of an
unmarked finite source derivative.

The resulting increment requirement, with independent deviations
before imposing their recursion, is
\begin{equation}
 \begin{aligned}
 b_{k+1}&=b_{k+1}^0+f_{k+1}(z_0,\ldots,z_k)+q_{k+1},\\
 \sup|f_{k+1}|&\le C\theta_k^{-1/2}
                            (\log\theta_k)^{\mathfrak p_0}
                           =:\epsilon_k,\\
 \sup|q_{k+1}|&\le C_v\theta_k^{-p_R}=:\epsilon'_k,
 &p_R&=(\kappa_0-4)/2 .
 \end{aligned}\label{eq:flow-RF20}
\end{equation}
For \(f\), the supremum is on the whole product
\(|z_i|<\sigma_\star\theta_i\), with
\(0<\sigma_\star\le1/2\), and requires neither a lower flow bound
nor remainders. For \(q\), the frame obeys the weak flow bounds,
\(|z_i|<\varpi_0\theta_i(\log\theta_i)^{-2q_0}\) with
\(\varpi_0\le\sigma_\star/2\), and the admitted aged remainder data
range over their stated holomorphic parameter domain. Their bounds are
\(4\theta_j^{-\kappa_0/2}e^{-\kappa d_j}\), with \eqref{eq:flow-RF18}.
The functions are real on real data. Constants are independent of
the frame length, torus, cutoff and admitted parameter; the weak
triple affects their threshold. These are bounds at every formed
step \(j\ge1\) of the box:
\begin{equation}
 \begin{aligned}
 |b_j^{\mathfrak A}-b_j^0|&\le\epsilon(\theta_{j-1}),&
 |b_j^{\mathfrak R}|&\le\epsilon'(\theta_{j-1}),\\
 \omega_*(y)&=\sup_{v\ge y}\{\epsilon(v)+\epsilon'(v)\}.
 \end{aligned}\label{eq:flow-RF20a}
\end{equation}
Here \(\epsilon,\epsilon'\) are the two functions in \eqref{eq:flow-RF20}.
The function \(\omega_*\) is their least nonincreasing majorant on
the admitted half-line. Calling a run alive does not replace the
separate drift bounds at its earlier levels by a bound only at its
last level.

There is also a weaker sufficient marked-multiplier estimate. If the
marked exponent has norm
\(\mathfrak n_k\le C\mathfrak g_k^{\kappa_0-5}\)
and the actual step is holomorphic in its multiplier on a disc
of radius comparable to \(1/\mathfrak n_k\), Schwarz yields
\begin{equation}
 |E^{(k+1),\mathfrak R}|
 \le C\mathsf s(x_k)\mathfrak n_k e^{-(1+c)\kappa d},
 \qquad |b_{k+1}^{\mathfrak R}|
                  \le C\mathfrak g_k^{\kappa_0-7}.
 \label{eq:flow-RF21}
\end{equation}
This preserves an alternative sufficient error for the real flow;
it does not prove the stronger exponent in \eqref{eq:flow-RF20}.
If this weaker error is used in a tuning argument, its actual power
must replace \(p_R\) in the marked bare-disc estimate, and the
summability condition \(p_R>1\) must be checked anew.

\paragraph*{Well-founded induction and statement 3.1.}
The substantive construction input says that, for an admitted box
and frame prefix through \(k\), \(E^{(k+1)}\) and its coefficient
are formed before any choices using \(\theta_{k+1}\). They use the
creation radii \((1+\beta_{\rm s})\alpha_\iota(\theta_k)\).
The terminal covariance radius is
\(e_{k+1}=c\varepsilon(\theta_k)\), so it also needs no next frame
entry. Once the computed next weight is admitted, the actual
\(\mathbf R\mathbf T\) operation closes the next box, consistently
on every admissible torus and its companion. The standing-history
kernel identity is
\((m(\kappa\otimes1))|_{\mathsf C}=(m|_{\mathsf A})\kappa\)
for a measure on \(\mathsf A\times\mathsf B\) and a kernel from
\(\mathsf A\) to \(\mathsf C\). Here \(\kappa\) denotes a kernel,
not a decay rate. This identity preserves the history sum; it does
not supply its quantitative mass or locality bounds.

Choose \(c_5=2c_2\), \(C_\beta\ge c_0+2c_2+1\), and first fix
\begin{equation}
 \lambda_\flat=c_0/4,\quad
 \bar B=C_\beta+1,\qquad
 \bar\varsigma=C'\gamma^{\kappa_0-8},\quad
 \lambda_\sharp=c_0/2-\bar\varsigma,\quad
 B_\sharp=C_\beta+\bar\varsigma .
 \label{eq:flow-RF22}
\end{equation}
The coupling ceiling is chosen so
\(\bar\varsigma\le\min\{c_0/4,1\}\).
The error exponent and other parameters must permit this ceiling;
\(\kappa_0>8\) is a sufficient choice for the displayed power of
\(\gamma\). An error estimate making that ceiling impossible
cannot be used in this parameter selection.
The marked-chain constants are computed at the weak triple before
this choice. Then choose \(Y_H=\gamma_H^{-2}\) above every required
floor, with
\(\epsilon(y)+\epsilon'(y)\le\min\{c_0/4,1\}\) and
\(\epsilon'(y)\le\bar\varsigma\) for \(y\ge Y_H\).
In particular \(\gamma_H\le
(\gamma^{-2}+c_5+1)^{-1/2}\). Increasing a reserve constant beyond
\([0,1]\) is allowed only after its stated floor is increased.

Put \(s_k=kc_0-\sum_{l\le k}b_l^0\); then \(|s_k|\le c_2\).
For any formed prefix in the weak class,
\begin{equation}
 \begin{aligned}
 x_i-x_j-(j-i)c_0
 &=s_i-s_j+\sum_{i<l\le j}(b_l-b_l^0),\\
 \tfrac34c_0(j-i)-c_5
 &\le x_i-x_j\le B_\sharp(j-i).
 \end{aligned}\label{eq:flow-RF23}
\end{equation}
The last upper bound uses the individual coefficient bound
\(|b^{\mathfrak A}|\le C_\beta\),
\(|b^{\mathfrak R}|\le\bar\varsigma\).
The lower bound is stronger than the strong triple, which in turn
implies the weak triple. Above a further floor, a bounded one-step
drop changes the smallest power \(R(y)\) by at most one factor \(L\).
Thus the computed next weight extends the frame while it remains
above the floor. This is the numerical part of breaking the induction
circle. The actual formation and box-closing input has not been proved
by that telescoping calculation.

For \(h=g^{-2}\ge Y_H\), define
\begin{equation}
 K(\beta,g)=\min\{k:x_k(\beta)\le h+B_\sharp\},\qquad
 \mathcal W_k(g)=\{\beta\ge h+B_\sharp:K(\beta,g)=k\}.
 \label{eq:flow-RF24}
\end{equation}
The conditional flow statement 3.1 is that the run exists through
this index, has the same coefficients on all admissible tori, and
\begin{equation}
 \begin{aligned}
 \lambda_\sharp(j-i)-c_5&\le x_i-x_j\le B_\sharp(j-i)
                         &&(i\le j\le K),\\
 h<x_K&\le h+B_\sharp,\\
 (\beta-h)/B_\sharp-1&\le K
 \le1+(\beta-h-B_\sharp+c_5)/\lambda_\sharp .
 \end{aligned}\label{eq:flow-RF25}
\end{equation}
At \(\beta=h+B_\sharp\), \(K=0\); the terminal inequality still
holds. For \(K>0\), the previous value exceeds \(h+B_\sharp\)
and the last decrement is at most \(B_\sharp\), proving its strict
lower endpoint. The lower drift forces finite exit. The two bounds
for \(K\) follow by applying the upper drift to \(K\) and the lower
drift to \(K-1\). Arbitrarily large \(\beta\) give arbitrarily large
\(K\), hence infinitely many nonempty \(\mathcal W_k\).
No stepwise positivity of \(b_k\) was used.

At depth \(s\), write \(x_{[s]}=x_{K-s}\). The exact envelopes are
\begin{equation}
 h+\max\{0,\lambda_\sharp s-c_5\}
 <x_{[s]}\le h+B_\sharp(s+1).
 \label{eq:flow-RF26}
\end{equation}
Put \(\underline x_s=h+\lambda_\sharp s-c_5\),
\(\bar x_s=h+B_\sharp(s+1)\), and
\(g_-(g)=(g^{-2}+B_\sharp)^{-1/2}\).
A diagonal subsequence of any first-passage sequence with \(K\to\infty\)
therefore converges at every fixed depth, to values in the closed
versions of these intervals. Its couplings converge as well, with
terminal value in \([(h+B_\sharp)^{-1/2},g]\).
A whole-sequence limit along tuned runs requires the later
matched-cutoff input; compactness alone gives no such conclusion.

The arithmetic section, \cref{ch:flow-arithmetic}, proves the
leading logarithm from vanishing mean errors, including bounded
backward excursions. In the present notation its uniform form is
\begin{equation}
 \begin{aligned}
 |x_{[s]}-x_{[0]}-sc_0|&\le\Delta_g(s),&
 \Delta_g(s)/s&\longrightarrow0,\\
 \left|\frac{\beta}{K\log L}
       -\frac{11N^2}{12\pi^2}\right|
 &\le\frac{\Delta_g(K)+h+B_\sharp}{K\log L}
                       \longrightarrow0 .
 \end{aligned}\label{eq:flow-RF27}
\end{equation}
Here \(K\ge1\), \(a_K=L^{-K}\). Denote the displayed relative
error bound by \(\delta_K\), distinct from the fluctuation
threshold \(\delta_j\). The function \(\Delta_g\) can be chosen
nondecreasing. Under the stated power/log error bound, with the
marked exponent sufficient for the same estimate, one may take
\(\Delta_g(s)\le C+C(h+c_0s/2)^{1/2}
(\log(h+c_0s/2))^{\mathfrak p_0}\).
The smaller ceiling \(\gamma_{\mathrm d}\le\gamma_H\)
imposes the monotonicity and envelope floors for this quantitative
bound; its constant is independent of \(g\) and \(s\).
For fixed \(g\) the relative error is
\(O_g(K^{-1/2}(\log K)^{\mathfrak p_0})\).
For example \(p_R>1\), used below, suffices for the marked sum.
The qualitative limit needs only an error tending to zero.

For the uniform depth statement, apply the arithmetic argument to the
segment from \(K-s\) to \(K\). Its shifted one-loop partial sums have
the same bound \(2c_2\). Use the envelope \(\omega_*\) of \eqref{eq:flow-RF20a}
on the self-calibrated run. One may explicitly choose
\begin{equation}
 \Delta_g(s)=2c_2+
 \sum_{u=0}^{s-1}
 \omega_*\!\left(h+\max\{0,\lambda_\sharp u-c_5\}\right).
 \label{eq:flow-RF27a}
\end{equation}
This is independent of the initial coupling and cutoff, nondecreasing,
and sublinear by Cesaro averaging. Integral comparison for the
square-root/log envelope gives the quantitative bound above.

The sharper two-sided logarithmic comparison retained in statement
3.1 is
\begin{equation}
 \begin{aligned}
 h-2c_2+\tfrac12c_0(K-k)
 &<x_k\le h+B_\sharp+2c_2+\tfrac32c_0(K-k),\\
 |x_k-x_K-c_0(K-k)|&\le\Delta_g(K-k).
 \end{aligned}\label{eq:flow-RF28}
\end{equation}
These are the logarithmic inequalities since
\(\log(L^ka_K)^{-1}=(K-k)\log L\). The upper slope in the first
line uses the uniform error bound \( |b_l-b_l^0|\le c_0/2\)
and the partial-sum bound, not an assumption on the sign of each
one-loop increment.

\paragraph*{Scale comparisons and layered evaluation.}
These estimates retain the distinction between a numerical flow sum
and an estimate on an actual field. Write
\(t_p(x)=x^{-1/2}(\log x)^p\), \(p\ge0\). If
\(-c\le x_m-x_n\le B(n-m)\) for \(m<n\), then, above a floor
depending on a freely chosen \(\beta_0>0\),
\begin{equation}
 \begin{aligned}
 t_p(x_n)&\le(1+\beta_0)(n-m)^{1/2}t_p(x_m),\\
 t_p(x_m)&\le(1+\beta_0)^2(n-m)^{\beta_0}t_p(x_n),\\
 \check R_j/\check R_{j+1}&\in\{1,L\},\\
 \sum_{j=0}^kg_j^p
 &\le C_p\left(1+\frac1{\lambda_\sharp(p-2)}\right)g_k^{p-2}
                       \qquad(p>2).
 \end{aligned}\label{eq:flow-RF29}
\end{equation}
The last line also assumes the lower drift bound. The first two
follow by comparing the square roots and logarithms above their
monotonicity floor; a bounded backward excursion is absorbed by the
chosen \(1+\beta_0\). The third uses the running minimum and the
bounded one-step drop. For the last, use
\(x_j\ge x_k-c_5+\lambda_\sharp(k-j)\), absorb \(c_5\) by the
floor, and compare the decreasing power with its integral.
The exponent \(p=2\) is excluded: that integral is logarithmic.

For the field estimate fix an admitted label through \(k\), with
layers \(\Gamma_n=\Omega_n\setminus\Omega_{n+1}\),
\(\Omega_0=T_\ell\) and \(\Omega_{k+1}=\emptyset\).
The determining sequence starts at \(\Omega_1^\sim\).
Let \(\Lambda_j^+\) be \(\Lambda_j\) together with its level-\(j\)
islands. Base points are the unit sites lying one layer of
\(M\check R_j\)-cells inside this set. Let \(\mathsf Y\) be any
set of these points and \(\mathsf Y_n=\mathsf Y\cap\Gamma_n\).
Here \(\Gamma_n\) is a geometric region difference; the bond record
of a determining sequence is a separate object. Write
\(\#_n(D)=\#(\mathbb T^{(n)}\cap D)\). This is a site count, equal
to volume in level-\(n\) units when \(D\) is a union of those unit
cubes. It is not a count of localization cells.
The cardinality \(\#\mathsf Y_n\), by contrast, counts the original
level-\(j\) base points. Remainder norms at a general fixed frame
use \(\mathfrak g_j^{\kappa_0}\); in the self-calibrated array
\(\mathfrak g_j=g_j\), as in the displayed layered bounds.

Assume the actual window/core and boundary hypotheses for these
backgrounds, Euclidean regular members of norm at most \(E_0\),
remainders of norm at most \(g_j^{\kappa_0}\), and the stated
boundary-domain decay. In particular the actual boundary terms,
indexed by their recorded origins \(o\), must satisfy
\begin{equation}
 |B_{\mathcal S}(j,X,o)|
       \le C_Be^{-\kappa d_j^\natural(\mathcal S;X)} .
 \label{eq:flow-RF30a}
\end{equation}
The metric is the one in \eqref{eq:rg-objects-R23} of
\cref{ch:rg-objects-v2}: the minimizing tree represents
each chosen standing-island hull by one member, rather than
filling it with all of its cubes. It is not ordinary
\(d_j\)-decay across an island.
The retained layered requirements are
\begin{equation}
 \begin{aligned}
 &\left|\sum_{z\in\mathsf Y}
   \left\{\sum_{\substack{X\ni z\\X\subset\Lambda_j^+}}
       [E^{(j)}(X,U_k,z)-E^{(j)}(X,\mathbf1,z)]
                      -b_jA(h_z,U_k)\right\}\right|\\
 &\hspace{12mm}\le E_1\sum_{n=j}^k
      (1+n-j)L^{-\beta_T(n-j)}L^{-4(n-j)}\#\mathsf Y_n,\\
 &\sum_{\substack{X\in\mathcal R_j\\X\cap\Omega\ne\emptyset}}
       |R^{(j)}(X,U_k)-R^{(j)}(X,\mathbf1)|\\
 &\hspace{12mm}\le R_1g_j^{\kappa_0}
             \sum_{n=j}^k(1+n-j)\#_n(\Gamma_n\cap\Omega^\sim) .
 \end{aligned}\label{eq:flow-RF30}
\end{equation}
Here \(0<\beta_T\le\alpha_{\rm H}<1\) is fixed,
\(\mathcal R_j\) consists of domains outside \(Z_j^+\), and
\(Z_j^+\) is the stored large-field region enlarged by one cell
layer. The set \(\Omega\) is a union of level-\(k\) localization
cubes and \(\Omega^\sim\) has one level-\(j\) localization layer
added. The constant \(R_1\) in this display is not the geometrical
gap parameter. Constants may depend on \(\beta_T,L,N\) and the
construction parameters, but not on \(j,k,K,\ell\), the label or
the positions.

For clarity, the local estimate behind the first line, when its
actual background hypotheses have been verified, is
\begin{equation}
 \begin{aligned}
 &\left|\sum_{X\in\mathcal X}
    [G(X,U,z)-G(X,\mathbf1,z)]-b(G)A(h_z,U)\right|\\
 &\qquad\le C\mathsf g
  \left\{(\varepsilon_z/a_0)^2
                \sigma^{4+\beta_T}
        +e^{-\kappa\min\{\mathsf D,c/\sigma\}/2}\right\},
 \qquad \sigma=L^{-(n-j)} .
 \end{aligned}\label{eq:flow-RF31}
\end{equation}
The member has norm \(\mathsf g\) on radii \(a_0\le a_1\), with
\(a_0\le1\); the datum is \(\varepsilon_z\)-regular near \(z\),
\(\varepsilon_z\le c\), and
\(\varepsilon_z\sigma\le ca_0\).
Every term is evaluated in its actual core. The collection
\(\mathcal X\) contains all domains within \(\mathsf D\) layers
of \(z\). With complete near coverage to distance \(c/\sigma\),
the right side is bounded by
\(C\mathsf g(1+(\varepsilon_z/a_0)^2)\sigma^{4+\beta_T}\).
The local cancellation and nonlinear remainder sections supply
the analytic implication under their full hypotheses, including
the boundary-response compatibility and the own tent counterterm.
Checking those hypotheses on this layered background remains
part of the actual construction.

The tail in \eqref{eq:flow-RF31} has no factor
\((\varepsilon_z/a_0)^2\). For example a wrapping holonomy term,
multiplied by its allowed exponentially small domain weight, can
have zero nonwrapping companion and coefficient zero while its
normalized value is nonzero on a flat torus field. A curvature
amplitude cannot control that topological contribution.

To pass to \eqref{eq:flow-RF30} one also needs the actual geometry:
\(\varepsilon_z\le C\check R_n^2\varepsilon_n\le
c\alpha_{0,n}\), the comparison
\(\alpha_{0,n}\le C(2+n-j)^{1/2}\check\alpha_{0,j}\), and the
available collar depth. The condition \(q>\mathfrak p_0+2r\)
pays the first comparison. The factor \(L^{-4(n-j)}\) is used
entirely by converting level-\(j\) sites to level-\(n\) sites.
For remainders the corresponding near-domain Cauchy bound is
\begin{equation}
 |R^{(j)}(X,U_k)-R^{(j)}(X,\mathbf1)|
 \le C(d_j(X)+3)^2L^{-4(n-j)}
  \left(\frac{\check\alpha_{0,n-1}}{\alpha_{0,j}}\right)^2
                    g_j^{\kappa_0}e^{-\kappa d_j(X)} .
 \label{eq:flow-RF32}
\end{equation}
Far domains use their norm, not this amplitude estimate.
The second line of \eqref{eq:flow-RF30} also holds with \(1+n-j\) replaced
by \(\varpi_{jn}=\max\{1,(x_j/x_n)^2\}\).
Both dominate the required squared-radius ratio up to fixed
constants. The two resulting summations give
\(|\mathbf R_k(U_k)|\le C\sum_n
g_n^{\kappa_0-6}\#_n(\Gamma_n)\).

The regular counterterm in \eqref{eq:flow-RF30} is the sum of its own tents.
An arbitrary smooth weight equal to one on \(\Omega\) differs
from that sum on its collar; it produces a new action term
without an age gain. For the actual profile \(\varphi_j\),
\(\varphi_jh_z=h_z\) at every base point. The remaining rim
terms obey \(0\le A(\varphi_jh_y,U_k)\le C\alpha_{0,j}^2\).
The needed geometrical count is
\(\sum_j\#_j(Z_j^+)\le C\sum_n\#_n(\Gamma_n)\).
It includes the stored collars of healed components, not just
the re-done small-field regions. It is an actual window/history
count and is not proved by a local Cauchy estimate.

Here \(\mathbf R_k\) denotes the sum of the normalized remainder
differences on the label. The net regular action
\(\mathbf E_k\) includes its profile counterterms and rim terms.
With those counts and boundary sums,
\(|\mathbf E_k(U_k)|\le C\sum_n\#_n(\Gamma_n)\).
For admitted complex perturbations of total coefficient size
\(\bar c\), regular norms at most \(E_0+1\), remainder norms
at most \(4g_j^{\kappa_0}\), and the same boundary hull bounds,
the sum of the elementary moduli is at most
\(C(1+\bar c)\sum_n\#_n(\Gamma_n)\).
This is a pointwise estimate on the admitted arrays; it is
not an assertion that arbitrary such arrays are produced by
the RG map. Its constants and evaluation floor need no new
dependence on \(\bar c\). The variable action coefficient is
\begin{equation}
 x_k(y)=\beta-\sum_{j\le k}b_j\varphi_j(y),\qquad
 |x_k(y)-x_k|\le B_\sharp
                  \#\{j\le k:\varphi_j(y)<1\}.
 \label{eq:flow-RF33}
\end{equation}
It equals \(x_k\) one coarse cell inside all the small-field
regions. Let \(A_{k,\mathcal S}\) be the ACTUAL label exponent
and \(\mathcal E_k\) its once-counted field-independent number.
The sign convention is \(\mathcal E_k=-E_{(k)}\), where
\(E_{(k)}\) is the accumulated log vacuum normalizer of \eqref{eq:comparison-capsule-CP1}.
Thus the constant part of the exponent is \(E_{(k)}\).
The exact density format, which is a construction input,
identifies the sum of the elementary regular, remainder, rim
and boundary terms with the expression on the left of
\begin{equation}
 |A_{k,\mathcal S}(U_k)+\mathcal E_k+
                       A(x_k(\cdot),U_k)|
 \le C(1+\bar c)\sum_n\#_n(\Gamma_n) .
 \label{eq:flow-RF33a}
\end{equation}
Their modulus sum proves this consequence on the admitted
backgrounds. Defining an exponent by a bound would not prove
the required exact density format.

\subsubsection{A common companion and the torus floor}
\label{sec:flow-a-common-companion-and-the-torus-floor}

\paragraph*{Statement 3.4 and its compatibility hypotheses.}
For each fixed frame and weight prefix there must be a single
family \(\widehat E^{(j)}(\widehat X,P,z)\) on the infinite
level-\(j\) lattice, and likewise a companion for \(R^{(j)}\).
For every admissible \((\ell,K')\), \(K'\ge j\), and every
label, it agrees with the torus term when the projection of
\(\widehat X\) is injective with its margin and lies one full
level-\(j\) cell layer inside \(\Lambda_j\). A sufficient margin
places \(\widehat X\) in a rectangle with sides less than half
the corresponding torus side. The admissibility used here is
\begin{equation}
 \ell L^{K'-i}\ge2L^{11}M\check R_i\qquad(i\le j).
 \label{eq:flow-RF34}
\end{equation}
The family has the same analytic bounds and cube-translation
covariance on every such lift. Its coefficient \(b_j\) is a
function only of the frame and weights through level \(j-1\).

Uniqueness follows because every bounded domain, collar included,
embeds in arbitrarily large admissible tori. Existence requires
more. First, the walk, propagator, localized background and gauge
choices must be made once by a rule natural under covering maps,
cube translations and the later dilations. Independently chosen
objects on separate tori cannot be compared. Second, each local
formation formula must use only its supplied neighborhood,
earlier terms restricted there and the common numbers \(x_i,b_i\).
The induction is simultaneous over tori of unbounded side; it
does not define \(b_j\) using only domains that happen to fit a
particular torus.

Third, formation must be local in the history trace on the
prescribed one-layer windows \(X_i^+\). A chart identifying
these traces, the bare weight, prefix coefficients and incoming
local terms identifies the newly formed term on the common
local core. In the one-layer interior the standing history has
empty trace: healed components have already been absorbed into
local terms. This actual standing-history property is what
makes that interior term independent of the label. It is not
implied by an inequality for the modulus of an activity.

The order \(E^{(j)},b_j,x_j,R^{(j)}\) avoids a circular
definition. In particular the new large-field weight \(x_j\)
does not enter the preceding extraction of \(b_j\).
This compatibility is an input to the full interacting
construction. The finite formulas proved elsewhere can verify
pieces of it; their holomorphy alone does not prove it.

\paragraph*{A locality-preserving choice of the supplied kernels.}
There is a conditional construction of a natural supply rule that
also explains a restriction on gauge representatives. Suppose the
real frame prefix and numerical parameters have first been fixed, and
the actual finite-world construction has supplied local coefficients
\(\mathsf K_\omega[\mathfrak s;P]\), the declared joint complex
charts and their uniform bounds, and its exact operator or functional
sum rule. A setting \(\mathfrak s\) records the geometric and regional
data; its continuous field variables remain in \(P\).
Assume that the finite-cover identity for the total and the actual
filing of lifted walks hold. Those are properties to check for the
supplied linear algebra, not conclusions about an arbitrary collection
of small functions.

Let \(\mathcal T_{\mathfrak w}\) be the finite translation group of
a world under which the class of settings is closed. Average by
\begin{equation}
 \mathsf K_\omega^\flat[\mathfrak s;P](z,z_1)
   =\frac1{|\mathcal T_{\mathfrak w}|}
       \sum_{a\in\mathcal T_{\mathfrak w}}
       \mathsf K_{\tau_a\omega}
           [\tau_a\mathfrak s;\tau_aP](z+a,z_1+a).
 \label{eq:flow-RF34a}
\end{equation}
The norm and decay bounds are convex, so averaging preserves them.
If two settings have the same trace on the supplied reading
neighborhood, each pair of summands reads the same translated trace;
the averaged terms therefore remain local. Covariance follows by
reindexing the finite group. A single representative of each
translation orbit would not do this: its translating displacement
can depend on data outside the local hull.

For gauge orbits, the representative is instead fixed by an actual
comb tree of a specified box and its entrance bond, with the root
outside the box. Put the gauge equal to the identity outside and
propagate its values from the root along that tree. Every tree bond
then becomes the identity. This determines the gauge uniquely.
Under a site gauge \(u\), each root-to-site transport changes to
\(u(o)\,U[o,x]\,u(x)^{-1}\). Consequently
\(\operatorname{rep}(\mathcal O^u)=
\operatorname{rep}(\mathcal O)^{u^\natural}\), where
\(u^\natural=u\) outside the box and \(u^\natural=u(o)\)
inside. The entrance bond and the exterior root are part of its
reading data. Selecting a representative by a global orbit ordering
would not prove this local law.

The margins needed for this construction are also part of the input.
Write \(\xi_l\) for lattice spacing and \(\mathsf w_l\) for the
assigned cell side at level \(l\). A region \(\Omega_i\) used to
form a cover or cutoff has level at least \(i\). Its cutoff makes
its whole transition within distance less than one of its own cell
sides from its complement. A walk cell of index \(n\) has size
\(O(M_1\xi_n)\); write \(\mathsf q^\boxplus\) for its enlargement
by \(5M_1\xi_n\) in the maximum norm. If \(n^*\) is the largest
index with \(\mathsf q^\boxplus\not\subset\Omega_0\), require
\(CM_1\xi_{n^*}<\mathsf w_{l(\Omega_0)}\); the condition is absent
when no such cell occurs. Each layer enlargement adds exactly one
layer of the assigned cubes, and inclusion, touching and component
rules read the region within its one-cell margin. Finally each
representative box, with its required neighborhood, fits in a closed
cube of side less than one quarter of the side of each compared world.
These conditions turn graded reading sets into the asserted local
ones. Averaging alone does not verify them.

Push the averaged coefficients from a cofinal sequence of covering
worlds to a fixed world and take a diagonal normal-family limit on
their common finite-dimensional complex charts. Discrete settings,
walks and matrix entries are countable; a countable compact exhaustion
of those charts suffices. The relevant bounds are uniform before the
limit. The cover filing may give a multiplicity at most
\(|\omega|+1\). Its first moment is summable when the rooted walk
entropy and \((CM_1^{-1/2})^{|\omega|}\) have a strict margin.
Dominated convergence then preserves the exact sum rule. This uses
the actual cover identity and the same declared cores throughout;
it is not an extension of a physical identity to independent currents.
Resolvents and square roots can subsequently be formed from the
limiting basic coefficients on their supplied spectral domains.
There is no need to diagonalize an uncountable list of resolvent
parameters. This argument constructs coherent choices conditional
on the uniform matrix, chart, reading and filing inputs just stated.

Under these inputs, the coefficient of the physical totals,
taken in the large-torus limit, equals the companion
coefficient in \eqref{eq:flow-RF2}. Use the all-small-field totals for their
unit-lattice symmetries; a truncated label total has no such
invariance in general. On nonwrapping domains the Hessians
agree by compatibility and the torus response converges to
\(\mathfrak h_\infty\). Fix first a probing length
\(R_{\rm pr}\), measured in level-\(j\) units and unrelated to
the cube-size function \(R(y)\). Take the torus side to infinity at
this fixed length; once it exceeds \(2R_{\rm pr}\), the relevant
bounded domains are lifts. Only afterwards send \(R_{\rm pr}\)
to infinity. Domains escaping the probe carry the supplied
tree-length tail, with rate after entropy
\((\kappa-\kappa_{\mathcal D})(R_{\rm pr}/(2M)-2)\).
 The rooted entropy is
absorbed by the chosen \(\kappa\); response decay controls
remote differentiations. Dominated convergence therefore
passes the differentiated invariant totals to the companion
kernel. This proves the sum-rule identification, including
that a constant total has coefficient zero. It requires the
uniform response convergence and domain bounds just stated.
The response limit need exist along a coherent admitted family of
representatives, not along every possible gauge choice on separate
tori. An oscillating fine gauge fixing the unit-lattice sites can
destroy such a derivative limit. If two admitted limits differ by
\((\mathrm{d}\lambda_B,0)\), the differentiated gauge Ward identity annihilates
that direction in the Hessian at \(P_\circ\); the coefficient is
therefore unchanged. This independence does not prove existence of
the all-level response limit.

Wrapping domains have
\(d_j(X)\ge c\ell L^{K-j}/M\). Their exponential decay is used
directly and does not define an additional renormalized
coefficient. For a self-calibrated run, the cube ratio shows
that \eqref{eq:flow-RF34} at the terminal level implies it at all preceding
levels. Since \(x_K\le h+B_\sharp\), it suffices that
\begin{equation}
 \ell\ge\ell_\star(g):=
               2L^{11}M R(g^{-2}+B_\sharp).
 \label{eq:flow-RF35}
\end{equation}
Write \(\mathfrak L(g)\) for the compatible torus sides above
this floor. The same floor works on the reference line,
whose terminal value is \(h\). The side is selected after
\(g\); the cutoff \(K\) remains unrestricted. This statement
is a uniform construction requirement, not an assertion of
infinite-volume uniqueness.

\subsubsection{Frozen choices, slopes, and tuned bare couplings}
\label{sec:flow-frozen-choices-slopes-and-tuned-bare-couplings}

The reference line for a run of length \(K\) is
\begin{equation}
 \theta_k=h+c_0(K-k),\qquad h=g^{-2}.
 \label{eq:flow-RF36}
\end{equation}
Equal depth in runs of lengths \(K\) and \(K+1\) gives equal frame
numbers, lattices, partitions, radii and cube sizes. It does not give
equal tests, since their backgrounds live over different bare lattices,
and it identifies no labels or histories between those runs.
The line beginning at depth \(s\) is the reference line of
\(g_{[s]}=(h+c_0s)^{-1/2}\).

For fixed choices, every inserted characteristic function is the same
real Borel function of its integration fields. The actual operations
act linearly on the initial density. If their required dominated
analytic estimates hold, applying them to
\(e^{-(\beta-\zeta)A}\) gives holomorphic measures and local families
in \(\zeta\). No derivative of a cutoff is taken. Continuity of the
real run follows from the actual fixed-frame input-list Lipschitz
estimate while the arrays retain their tube room. It is continuity,
not a contraction. When each bare coupling chooses its own frame,
the tests and domains can change discontinuously and this argument
does not apply.
For a fixed-frame remainder comparison, both real centers
\(\beta,\beta'\) must be alive through the compared prefix.
The interpolated aged remainder list is
\begin{equation}
 H_\upsilon^{(j)}
   =R_{\beta'}^{(j)}
          +\upsilon(R_\beta^{(j)}-R_{\beta'}^{(j)}).
 \label{eq:flow-RF36a}
\end{equation}
At their common frame domains, the individual norm bounds give
\(\|H_\upsilon^{(j)}\|\le4\theta_j^{-\kappa_0/2}\) for
\(|\upsilon|<3/2\). Its companion, covariance radius and invariance
properties must be those required by the admitted list slot.
The holomorphic list estimate then applies to this actual
interpolation; an unrelated real center or an unspecified change
of the earlier histories cannot be substituted for it.

The exact decompositions of unity and normalized gauge factors imply
that the bare measure, partition function and observable integrals
are independent of admitted auxiliary choices. Intermediate
densities after a large-field operation need not coincide.
Comparing their extracted couplings requires E4. The retained
self-calibrated versus reference-run target is
\begin{equation}
 \max_{k\le K}|x_k(\beta)-x_k^\natural(\beta)|
 \le Cg(\log g^{-2})^{3\mathfrak p_0+1}\le1 .
 \label{eq:flow-RF37}
\end{equation}
It is required for first-passage data with
\(\varrho\ge B_\sharp+1\) and
\(0<g\le\gamma_{U'}\le\gamma_U\). Threshold and large-field
changes contribute at most \(C\theta_k^{-(\kappa_0-4)/2}\)
per step, and gauge-fixing changes at most
\(C\theta_k^{-3/2}(\log\theta_k)^{3\mathfrak p_0+1}\).
Proving those actual comparisons remains necessary. Equality of
bare integrals does not prove \eqref{eq:flow-RF37}.

The small-field coupling polydisc has relative radius
\(\sigma_\star\theta_j\). The large-field tube is narrower.
On a level-\(j\) unit cube, an action-weight change has size
\begin{equation}
 |x-\theta_j|\alpha_{0,j}^2
 \le C\mathsf t_j\theta_j^{-1}(\log\theta_j)^{2q}.
 \label{eq:flow-RF38}
\end{equation}
A logarithmic tube factor at least \((\log\theta_j)^{-2q}\)
keeps this bounded. A fluctuation difference gains a further
\(\delta_j/\alpha_{0,j}\). This explains the different roles
of the full small-field polydisc and the large-field tube.
It does not justify extending the entire interacting history
construction to the larger polydisc. Complex bare couplings in
the latter construction are obtained near admitted real runs
through the fixed-source analytic domain.

\paragraph*{Statement 3.5: the weaker marked-disc slope route.}
For this paragraph require \(p_R>1\). The unmarked term has the
independent-variable bound in \eqref{eq:flow-RF20}. Along the actual run the
marked term is allowed a different analytic input: on a disc in
the bare variable about each admitted real centre, of radius at
least \(c_\beta\varrho\), it is holomorphic and bounded by
\(\epsilon'_k\). Bounds are uniform on centres whose required
prefixes retain the prescribed tube room. On a smaller concentric
disc, Cauchy therefore gives \(C\epsilon'_k/\varrho\).
This hypothesis includes the actual remainder dependence on that
bare variable. It does not assert that this dependence is a
function of independent deviations alone.

% Earlier Part II, 14167--14387 and 14415--14422.
On the admitted constant-source disc this bound has a native
frozen-step calculation. Form \(E^{(k+1)}\) on the whole lattice
with the real fluctuation Gaussian, its bond cut and the source
vertex. The input list consists of the renormalized regular
families and the vacuum-subtracted remainders that pass the age
test. Form \(P^{\flat(k+1)}\) with the same Gaussian, cut and
running shifts, keeping just the recursively unmarked regular
families. Put \(Q^{(k+1)}=E^{(k+1)}-P^{\flat(k+1)}\).

At the frozen shift \(\zeta_k\), write \(s=1-\zeta_k g_k^2\).
The quadratic part of the source vertex is
\(-\tfrac12(s-1)\langle B,\Delta B\rangle\); together with the
real Gaussian it gives the Gaussian of the unmarked recursion
at ratio \(s\). Its other vertices are the remaining terms of
the same rewritten interaction. On a common cut and common
radii the renormalized inputs agree by induction, up to field-free
constants, with the unmarked list evaluated at the running shifts.
Those constants do not enter the marginal coefficient. Hence the
unmarked output supplies the one-loop and \(f_{k+1}\) terms in
\eqref{eq:flow-RF20}, and the marked coefficient is
\(\beta[Q^{(k+1)}]\).

In the frozen integral, split each inherited regular group by
linearity of its own marginal subtraction. The groups from
\(Q^{(j)}\), together with the admitted remainders, are exactly the
marked inputs. The first-stage contour estimate bounds their total
weight by the small fluctuation factor times the sum of the
renormalized \(Q\) loads and the aged remainder load. Removing
these inputs leaves the unmarked integral. The second-stage
cluster comparison bounds \(\|Q^{(k+1)}\|\) by a fixed multiple
of that marked weight. In the supremum form of
\eqref{eq:flow-RF19}, choose the coupling small enough that this
multiple times the fluctuation factor and the total regular-group
kernel is at most \(1/2\), and the remainder load has coefficient
at most \(1/2\) as well. The remainder weight here is comparable
to \(x_j/x_{k+1}\); its sum uses the smallness at birth. Since both
the earlier supremum and that load are bounded by the current
envelope, the new \(Q\) lies in the same envelope, of size
\(C\theta_k^{-(\kappa_0-2)/2}\). Extracting the coefficient on
the admitted running radius costs \(C\theta_k\), giving the
required \(C\theta_k^{-p_R}\) on the bare-variable disc.
This calculation requires holomorphic inherited remainders and
the identification of the actual frozen production with these
lists, including its cut and radii.

%
%
%
%
For the inherited remainders, the native step uses two bounds on
the same constant-source disc. The contribution from summing healed
components has weighted norm at most \(C e^{-p_0(g_{k+1})}\) on
domains away from the one-layer enlargements of the live regions.
Separate the empty-family term before taking the polymer logarithm.
A polymer disjoint from all second-class strips contains a healed
operation.
The terms meeting those enlargements go into the boundary list
and use the relative hull bound there. The complete boundary slot
also retains its empty-family term with the ordinary amplitude.

The second contribution comes from the ordinary fluctuation
integral. Keep its normalized reference Gaussian and its cut
fixed as in \cite[(3.26)--(3.29)]{B88b}. Use two potential
parameters. The first multiplies
the boundary terms, including the changes of the stored
operations, the weight variation near the large-field region,
the bounded rim part of the quadratic gauge correction,
and the change-of-variables terms of the rim. The second
multiplies the recent remainders evaluated as differences
between the pre-integration and post-integration backgrounds.
For an origin \(j\), first subtract the vacuum value of
\(R^{(j)}\). The vertex in the interpolation is the difference
\(R^{(j)}(U_k)-R^{(j)}(U_{k+1})\), evaluated on the two actual
backgrounds in the variables of this integral. The boundary
expectation varies the first parameter with the second at
one. The recent-remainder expectation varies the second with
the first at zero. The logarithm with both parameters zero
uses the aged, vacuum-subtracted remainders with the regular
vertices. The Gaussian normalizer belongs to the regular
part of the step and the reference Gaussian used for these
expectations has unit mass before the cut.

The two changes of the logarithm are the integrals of the
corresponding normalized expectations. In the local expansion,
the second change selects precisely the clusters with at least
one marked recent vertex on the same cut. File its terms on
the one-layer inset as \(R''^{(k+1)}\). Its terms in the rim,
and the first expectation's terms, go in the boundary family.
The conditional Gaussian in the history can have a nonzero
mean. With precision blocks \(\mathsf A_{\rm II}\) and
\(\mathsf A_{\rm IE}\), and exterior field \(a_{\rm E}\),
the mean is
\(-\mathsf A_{\rm II}^{-1}\mathsf A_{\rm IE}a_{\rm E}\).
The restriction to the centered term takes place after
decoupling. In \eqref{eq:local-formation-SI7} every parameter
off the carrier \(Z\) is zero. A walk meeting a cube off
\(Z\) then has weight zero. The remaining Gaussian terms
use the restricted form, its resolvent and the exterior
columns carried by \(Z\). A nonzero column must carry the
whole closed bond segment of its exterior argument. For
\(Z\subset\Lambda_{k+1}\) there are no such exterior or
crossing bonds. The localized mean therefore vanishes and
every fluctuating bond of \(Z\) has the interior cutoff \(T_k\).
With the same reference scale and Gaussian factor the
localized operators are those of the constant sequence of
regions \cite[p.~278]{B88b}. In the logarithm each such \(Z\), and the full
collar of each selected potential, lie in the filing domain
\(X\). The marked coefficient is thus the centered
coefficient at the same local pair and with the same
restricted list of input potentials. Match the decoupling
guard to compare individual summands. Regrouping the
decoupling brackets gives the guard-independent coefficient.
A term reading the rim uses the conditional formula and is filed at the
boundary. The one-layer inset decides where the centered
recent coefficient is filed.

The comparison uses a common localization and source
presentation. For the presentation which instead puts the
constant shift into the Gaussian, the comparison of marginal
coefficients uses the identity of the totals up to the
field-independent Gaussian normalization.

The recent input band is
\(J_{\rm rec}(k)=\{j\le k:M\check R_j>L^{k+1-j}\}\), restricted
to existing remainder families. In particular it is defined
at the input step, and includes the level \(k\) family when
it fails the ageing test. A localized vertex from \(j\) is
estimated first in its own box. In comb gauge the scaled
bond logarithm there is bounded by \(C D_Y t^2\) times the
output radius, where \(D_Y\) is the diameter of the input
domain plus the box margin in its own units. The first
localization is at input level \(k\), so \(t=L^{j-k}\).
Scale these bond logarithms by one complex parameter and
use the admitted background map on the resulting disc.
The first derivative at the trivial pair is an invariant
linear functional for the adjoint action of the constant
gauge group and vanishes by semisimplicity. The second-order
Schwarz estimate bounds the vacuum-subtracted value by the input size times
\(C D_Y^2t^4\) times the squared radius ratio. The difference
of the two background values has the same bound up to a
fixed factor. If the scaling disc has radius at most one,
use the input bound at the admitted endpoints. In the rooted
sum the output tree length is at most \(t\) times the input
tree length. The spare decay sums the diameter factor. A
domain extending beyond the containing box has input tree
length at least of order \(t^{-1}\). For \(j<k\), extract
\(t^4\) from its spare exponential tree decay and keep a
positive tree rate for summation. At \(j=k\) the input and
root weights already leave a positive tree rate and \(t^4=1\).
In both cases the \(O(t^{-4})\) fine anchors in a root cell
are paid by the \(t^4\) gain.

In the first localization of these marked vertices, fix a
real fluctuation \(B\) on the cut \(\sup|B|\le T_k\).
The response estimates of \cite[(1.17)--(1.22)]{B88a} on the
cut-parameter contours bound
the undilated response on the read set by \(K_Rg_k\sup|B|\)
and the part localized at a block by
\(K_R\eta g_k\sup|B|\). Here \(\eta\le1\) is the decay
factor from the block to the read set and \(K_R\) is uniform
on the output tube. The estimates apply also to complex \(B\).
If \(a\) is the amplitude of the localized part, the sufficient
containment inequality is
\[
 (1+\eta|a|)K_Rg_k\sup|B|\le\alpha_{*,k},
\]
where \(\alpha_{*,k}\) is the response radius. The top layer
uses the base and dilated parts separately within its
analyticity margin. At older layers the shift gains the
decaying transfer factors and uses the corresponding layer
room. With \(\theta_k=K_Rg_kT_k/\alpha_{*,k}\le1/3\), the
disc in the rescaled amplitude \(\eta a\) has radius
\(\theta_k^{-1}-1\ge2\). If the vertex differs from its
vacuum value there by at most \(N\), its change from the
undilated vertex is bounded by \(2N\). The Schwarz estimate
on the subdisc \(|\eta a|\le2\) bounds the change by
\(4N/(\theta_k^{-1}-1)\le6\theta_k N\).
Each differentiated cut parameter, integrated
over \([0,1]\), contributes at most
\(e^{\kappa_1}(e^{\kappa_1}-1)^{-2}\) by Cauchy's formula
on \(|\sigma|=e^{\kappa_1}\). The rooted sum therefore has
the factor \(\sigma_k\le C(\log x_k)^{\mathfrak p_0-q}\)
on the fixed real cut.

The second localization at zero exterior fluctuation keeps
this cut in the Mayer factors. In
\eqref{eq:local-formation-SI6}, \(Y_0\) contains the free
bonds read by each selected potential that requires a field
box, including each marked one. The failed-cut bonds lie
outside \(Y_0\). For the conditioned Gaussian
formula of \cite[(2.5)--(2.6), (2.14)]{B88a}, let \(C_\sigma\)
be the conditional covariance on the integration region and
\(\Gamma_\sigma X\) the exterior linear term, written in real
auxiliary Gaussian coordinates \(X\). The inner integral uses
the real variable \(B\), with Gaussian density of covariance
\(C_\sigma\) and the multiplier
\(\exp(-\langle B,\Gamma_\sigma X\rangle)\).
The cut-parameter contours complexify \(C_\sigma\) and
\(\Gamma_\sigma\). The potentials in the integrand are still
evaluated at \(B\) with the factor \(\chi_{Y_0}\), and each
Mayer contour multiplies its potential by a scalar. Take the
potential bound on the support of \(\chi_{Y_0}\) before
enlarging the Gaussian domain for the modulus estimate.
The determinant and Gaussian bounds then justify continuation
on the same real integration space. Thus the small factor
from the first localization is available on the contours
used to form the ordinary Mayer activities.

Scale the marked activities together by a complex parameter.
If their norm is \(\varepsilon_m\), the unused activity margin
gives a disc of radius proportional to \(\varepsilon_m^{-1}\)
on which the localized output has the usual uniform decay
bound. Subtract its value at zero and apply the Schwarz
estimate at one. This makes the output proportional to the
marked activity norm in the same weighted family. The bound
for this contribution is therefore
\[
 C\sigma_k\sum_{j\in J_{\rm rec}(k)}
        w_{k,j}g_j^{\kappa_0},
\]
where \(w_{k,j}\) is the squared ratio with the enlarged
radius of this step in the numerator. The amplitude bound
uses the remainder's own analyticity domain for each read.
The Mayer bound uses the same cut on every marked read and
the activity margin uniformly on the admitted source disc.
With the running-radius estimate \(w_{k,j}\le C(k+2-j)\), the sum of
these ratios is at most \(C(1+\#J_{\rm rec}(k))^2\). To bound
the band length, take
\(u=\lceil C(1+\log\log x_k)\rceil\) with the constant large enough.
The flow gives \(x_{k-u}\le x_k+C'(1+u)\) when \(k\ge u\).
Then \(M\check R_{k-u}\le L^{u+1}\) by the logarithmic bound
on \(\check R\). Thus \(k-u\) has already reached the ageing test.
The grid rule \(\check R_i\le L\check R_{i+1}\) implies
\(L^i\check R_i\le L^{k-u}\check R_{k-u}\) for \(i\le k-u\),
so those levels have reached it as well.
If \(k<u\), the whole list has length at most \(u\). The
reference flow then gives \(g_j^{\kappa_0}\le Cg_{k+1}^{\kappa_0}\)
throughout this band. For the native reference radii the sharper
comparison is
\[
 w_{k,j}=\tfrac94\,\frac{x_j}{x_k}
            \left(\frac{\log x_k}{\log x_j}\right)^{2q} .
\]
Here the numerator uses the radius enlarged by \(1+2\beta_s\)
at input level \(k\). The adjacent-scale room puts the actual
output space with radii \((1+\beta_{\rm pr})\alpha_{k+1}\) inside
this space, using
\((1+\beta_{\rm pr})\alpha_{k+1}\le(1+\beta_s)\alpha_{k+1}
\le\hat\alpha_k\) componentwise. Since \(x_j=x_k+O(1+k-j)\) in the band,
\(w_{k,j}\le3\) at small coupling, so this reference estimate
costs only a single power of the band length. For
\(q>\mathfrak p_0\), even the running-radius bound with the
squared count tends uniformly to zero after multiplication
by the fluctuation factor.

Each contribution is at most \(g_{k+1}^{\kappa_0}/2\) at weak coupling.
Both bounds use the native enlarged remainder charts. In the
recent term the running shifts are the functions of the single
bare constant source on its disc. The marking parameter is
independent of that source. The healed term is holomorphic by the majorized
resummation, and the recent term by the fixed-domain fluctuation
expansion. Restrict their output to the running remainder radii
when taking the inherited norm. The sum has norm at most
\(g_{k+1}^{\kappa_0}\) on the constant-source disc, so the next
frozen integral can use it as inherited data.
For this application the actual formation and norm domains must
agree with the native ones.

%
The inherited family also needs the common-lift and symmetry
properties. Compare the partial maps on the same local field
arguments, using the same reference sequence, local operator
terms, weights, source, increments and inherited lists.
For coherent input, the earlier lists are determined on
their own read sets by the standing history. In this history
the healed components have already been integrated out.
The remaining large-field regions are nested into the
current complement, so the standing trace on an interior
window is empty. The trace uses one layer of the label cells
at each relevant level. On a carrier
with this window inside \(\Lambda_{k+1}\), the local data
are those of the all-small-field history. This identifies
the inherited coefficients as functions on the common
domain. An evaluation in the actual history uses its own
local field argument. A common non-wrapping lift transports
the walk terms, cutoffs and inputs by a translation
preserving their grids. The cluster logarithm and the
filing therefore give the same coefficient on the lift.
At a real constant source, evaluate on the backgrounds of
the common hereditary window of very regular block fields.
For isometries preserving the level-\((k+1)\) label cubes,
the density obtained by summing the histories with the
all-small-field label and the regular integral with its aged
inputs are invariant. In the quotient of these densities the
carried background factors are invariant too, since the new
label cubes are unions of the older ones. Remove those factors
and normalize at the trivial field. The quotient is now the
exponential of the vacuum-subtracted new total. Its exponents
at a field and at its isometric image are real and therefore
equal.
Holomorphy extends that identity to the connected constant-source
disc on the same backgrounds.

When a remainder of origin \(j\) first enters the regular
integral at step \(k\to k+1\), the ageing rule gives
\(M\check R_j\le L^{k+1-j}\). Both lengths are powers of
\(L\) on the nested grids. Thus the origin cube side
divides the spacing of a level-\((k+1)\) block in origin
units. Coarse unit translations and the coordinate
isometries preserve the origin cube partition. The
invariance of the earlier total therefore suffices for its
use in the regular integral of this step.

% Earlier Part III, 3542--3565, 3778--3983, 4119--4129 and 4193--4266.
% Earlier Part II, 14834--14892, 15644--15669, 15818--15901,
% 15921--15968 and 16515--16869.
% Earlier Part II, 7150--7179, 7491--7504, 7590--7664 and 7709--7738
% (size-table weights), 16306--16340, 16870--16906, 21913--21936 and
% 24997--25022 (rooted values).
% Earlier Part I, 65765--65816, 66319--66334 and 137393--137474.
% Earlier Part I, 142661--142777 (interrupted change of variables).
% Earlier Part I, 75254--75786 and earlier Part II, 17840--17998,
% 24235--24327 and 24860--25054 (cut-family estimates).
% Earlier Part II, 23257--23420 and 25107--25250 (amplitude pieces).
% Earlier Part I, 75805--76106, 76667--76850, 80689--81196,
% 81200--81622 and 81722--81764 (wrapped terms).
% Earlier Part I, 102457--102690 and earlier Part II, 16013--16256
% (Wilson remainders); earlier Part II, 14732--14905, 15858--15899
% and 18033--18172 (volume estimates).
For the strip estimate in the healed contribution, measure lengths
in localization cubes of side \(M\). Work on the class of boxed
components of \(Z\) separated by at least \(\check R_k\) cubes.

First consider the plain values in the input of the strip estimate.
For a value unit \(v\) with carrier \(K_v\), adjoin the full components
of \(Z\) sharing a cube with \(K_v\) to obtain its filing domain
\(Y_1\). This leaves \(Y_1\setminus Z=K_v\setminus Z\). A tree for
\(K_v\) meets every cube of \(Y_1\setminus Z\), so
\(d_Z(Y_1)\le d_k(K_v)\). For an exterior \(Y_1\), fix a cube
\(\Delta\subset Y_1\setminus Z\). Every carrier in the fibre contains
\(\Delta\), giving the bound
\[
 e^{a d_Z(Y_1)}\left\|\sum_{v:\,Y_1(v)=Y_1}\mathrm{val}_v\right\|
 \le\sum_{v:\,K_v\supset\Delta}e^{a d_k(K_v)}\|\mathrm{val}_v\|,
 \qquad a=31\kappa/20 .
\]
Thus one needs a weighted sum at a cube outside \(Z\). This first
filing spends no decay and introduces no factor for the size of a
component of \(Z\).

The current chain values filed at \(m<k\) are confined to a single
component of \(\Lambda_k^c\), by connectedness and
\(\Omega_{m+1}^c\subset\Omega_k^c\subset\Lambda_k^c\). A carrier of
this kind meeting \(Z\) stays in \(Z\). For those filed at \(k\),
the stage norm at \(8\kappa/5\) bounds the column at rate \(a\) by
\(C_0^\sharp\) per stage. The current chain therefore costs at most
\(N(g_k)C_0^\sharp\).

For the plain top-level values left by an earlier arrest, put
\(j=k-s\), \(s\ge1\). Their native carrier \(X\) meets
\(\Omega_k^c\). If its coarse closure \(\bar X\) meets both \(Z\)
and its complement, it crosses the collar between
\(\Omega_k^c\) and \(\Lambda_k\). Consequently
\[
 d_j(X)\ge(2\check R_k-1)L^s,\qquad
 d_k(\bar X)\le L^{-s}d_j(X)+c_{\rm tr} .
\]
Coarsening spends at most \(31\kappa/60\) of the native rate
\(8\kappa/5\), leaving \(13\kappa/12\). Use the resulting factor
\(e^{-\kappa\check R_k L^s}\) and sum the native weighted sizes over
the \(L^{4s}\) fine cubes in \(\Delta\). The stored values are a
subsum of their birth-stage norm, which already sums the formed
terms over the components of that stage. In this count each formed
term is taken once. The bound over all stages and ages is thus
\[
 C_0^\sharp e^{ac_{\rm tr}}
 \sum_{s\ge1}N(g_{k-s})L^{4s}e^{-\kappa\check R_k L^s},
\]
where the sum uses only existing earlier levels. The reference
flow gives \(N(g_{k-s})\le C(\log x_k+c_s)^{r_{\rm pr}}\) with
\(c_s=O(\log(2+s))\). The width bound
\(\check R_k\ge(\log x_k-1)^{r_{\rm pr}}\) bounds this count by
\(C(2+c_s)^{r_{\rm pr}}\check R_k\) for \(\log x_k\ge2\).
Since \(\check R_k\ge1\), one has
\(\check R_k e^{-\kappa\check R_k L^s}\le
2(e\kappa L^s)^{-1}e^{-\kappa L^s/2}\). This bounds the contribution
uniformly in \(k\) and the coupling. At the final stage, a
lower-filed term whose carrier reaches
bonds adjacent to an arrested piece must also satisfy the collar
test if it enters the exterior column. When its native carrier
passes this test at level \(\ell<j\), use the birth-stage norm at
\(4\kappa/5\) and coarsen from \(\ell\). The spare rate is at
least \(17\kappa/60\), so its stage contribution is bounded by
\(C L^{4t}e^{-17\kappa\check R_k L^t/60}\), with \(t=k-\ell\).
At this filing level there are at most \(t\) earlier birth steps
and at most \(C(2+c_t)^{r_{\rm pr}}\check R_k\) stages at each
step. The collar factor sums these terms by the same absorption.

The change-of-variables terms of an earlier arrest use their own
filing level \(\ell<k\). Set \(t=k-\ell\). Their connected exterior
domains \(X\) meet their stage band inside \(\Omega_k^c\), so
\(d_\ell(X)\ge(2\check R_k-1)L^t\). After coarsening, a native
pointwise rate of \(4\kappa/5\) leaves \(17\kappa/60\). In using
this pointwise estimate we must count the birth blocks even when
several terms have the same localization domain.

For terms carrying their stage-cover anchor in \(X\) on the
filing grid, there is at most one term per cover block, and the
bands of simultaneous arrests are disjoint. At a fixed birth
step and stage the number of such blocks for \(X\) is therefore
at most a fixed multiple of \(1+d_\ell(X)\), by the cube count
for a tree meeting \(X\). The same bound applies for anchors in
a fixed-width enlargement on this grid. The share \(\kappa/20\)
of the spare rate sums the rooted domains with this additional
linear factor. The remaining \(14\kappa/60\) pays the collar,
and the \(L^{4t}\) fine roots in the coarse cube give a stage
bound of the form
\[
 C L^{4t}\exp\bigl(-\tfrac{14}{60}\kappa\check R_k L^t\bigr).
\]
There are at most \(t\) possible earlier birth steps at this
filing level. Their stage counts are bounded by
\(C(2+c_t)^{r_{\rm pr}}\check R_k\). Thus the sum over \(t\ge1\)
is uniform by the same exponential absorption as for the
top-level values. For a cover on a different grid we must show
that its localization still has this anchor count in filing
units, or use its own rooted norm before summing the scales. For
the bond-indexed terms of an interrupted change of variables,
the bound on each bond term gives its size on the admitted chart.
We must also identify the localized carriers and show the required
tree-length decay on them.

The two kinds of stored boundary terms use different carriers.
The boundary term formed in the \(\mathbf T\)-operation has a
pointwise bound at rate at least \(2\kappa\) on its native domain.
At \(j=k\), the spare rate \(2\kappa-a=9\kappa/20\) sums the domains
through the fixed root. At an earlier level \(j=k-s\), the same
coarsening bound leaves \(89\kappa/60\). Spend \(\kappa/20\) on the
rooted native domain count and use the collar to bound the rest.
With \(\kappa\ge10\log L\) and \(L\ge3\), the remaining
\(86\kappa/60\) pays the \(L^{4s}\) fine roots and leaves
\(e^{-\kappa\check R_k L^s}\) for the sum over ages.

The boundary term formed in the \(\mathbf R\)-operation is stored
on its relative label, whose coarse image serves as the filing
carrier. At the filing site the term is from an
earlier level \(j=k-s\), since the new one is formed after filing.
Its label meets an earlier strip or zone inside \(\Omega_k^c\).
For an exterior value the label therefore crosses the same collar.
On the label the stored rate is at least \(\kappa\), and coarsening
spends at most \(31\kappa/60\). Of the remaining \(29\kappa/60\),
use \(\kappa/20\) for the rooted label count and split the rest
into two equal parts. The first pays the fine roots since
\(\frac{13}{60}\kappa L^s\ge4s\log L\). The second gives the
summable factor \(e^{-13\kappa\check R_k L^s/60}\). The skeleton
fibres are summed when the stored terms are formed. This estimate
uses their resulting bounded prefactor and the decay on the
labels.

For a chain term with a nonempty decoupling bracket, the original
filing level supplies the rate before the bracket enlargement. At
the current level the rate \(8\kappa/5\) leaves \(\kappa/20\) after
allocating \(a\) to the original carrier. At a lower level \(j<k\),
coarsening leaves \(4\kappa/5-31\kappa/60=17\kappa/60\) in
level-\(j\) units. The fixed prefactor \(e^{ac_{\rm tr}}\) goes into
the constant paid by the inverse localization parameter of the
bracket estimate. The added decoupling carriers are charged in
that estimate.

For the terminal, renormalized, tail and boundary rows of the
size table, keep the point anchor when taking the weighted
domain sum. In a zone \(n\), set \(u=L^{j-n}\) for an anchor of
level \(j\le n\). The \(O(u^{-4})\) fine points in a zone cell
use up the \(u^4\) in a short piece or a local Wilson term.
For a short boundary piece, the anchor must lie in a slab of
thickness \(O(1+n-j)\) fine units around a zone face. The slab
count is \(O((2+n-j)u^{-3})\), leaving \((2+n-j)u\) after the
size is inserted. For an old tail the lower bound on its
diameter supplies \(u^5\) from the spare decay, so the same
fine-point count leaves \(u\). The frozen core has the extra
power \(u^{\hat\alpha}\) with \(\hat\alpha>0\), and the defect
row has \(L^{-(k-j)/2}\). The renormalized remainder row uses
its admitted weighted sum over levels after the \(u^4\) and
fine-point count cancel.

In the rooted column the weight \(7\kappa/4\) spends at most
\(7\kappa/12\) in level-\(j\) units when \(j<k\). At \(j=k\)
the ordinary rows present there have the improved rate
\(2\kappa\). After the rooted domain count and diameter factors,
keep \(e^{-\delta\kappa d_j(X)/4}\) for the point anchors. For
anchors in the current zone, use the per-cell scale sums and
then sum over cells. There are at most \((2i+3)^4\) cells in
a current cube shell of radius \(i\), and a domain reaching the
root from this shell has length at least \((i-1)_+\) in its
native units.
For a unit born at \(j=k-s<k\) with anchor in an old zone,
the domain must cross the collar as well. One can split the
remaining exponential into two halves and use
\(d_j(X)\ge(2\check R_k-1)L^s\) in one and the distance from the
anchor to the root in the other. In current units there are at
most \(C(ML^s)^4(2i+3)^4\) fine anchors in a shell of radius
\(i\), and the distance bound in fine units is at least
\(L^s(i-1)_+\). The shell sum is bounded by a constant times
\(L^{4s}\), which is summable in the ages against the retained
collar factor. These estimates use the assigned sizes and their
native anchor domains, including the domains of the subtracted
copies.

The short singles and the Wilson terms across the boundary use
an enlarged filing label \(\operatorname{lab}_1=K_v\cup\mathfrak h_v\),
where \(\mathfrak h_v\) covers the additional reads. Write \(c\) for
the fixed bound on \(d_k(\operatorname{lab}_1)-d_k(K_v)\). For these
terms split the filed value into the uncut value, which is charged
at the size in the table, and its difference from the cut value.
The carrier estimate bounds the absolute difference by
\(\mathfrak N^{\flat}(v)\) on the admitted cut chart, including when
all decoupling families are empty. To use this estimate we must
verify that the auxiliary consumer lies in its filing domain and
is admitted on the cut chart.

For the exposed units in this exterior column, their point lies in
a level-\(k\) collar cell or in a level-\(k\) cell next to the treated
boundary on the \(\Omega_k\) side. The cell weight there is \(3\).
For such a cell \(\Delta'\), the sum of
\(\mathfrak N^{\flat}(v)e^{(a+9)d_k(K_v)}\) over its point anchors is
bounded by \(27\mathsf K^{\flat}\) by the per-cell estimate at
\(a+9\le 8\kappa/5\). If the enlarged label contains the exterior
root \(\Delta\), its connectedness gives
\(d_k(K_v)\ge(i-1-c)_+\), where \(i\) is the cube distance between
\(\Delta'\) and \(\Delta\). Therefore even at the activity weight
\(a+8\) the cost of summing the difference over the anchors is at
most
\[
 27\mathsf K^{\flat} e^{(a+8)c}
 \sum_{i\ge0}(2i+3)^4 e^{-(i-1-c)_+} .
\]
The constant is uniform in the coupling and in the level. This
sum uses the location of the exterior point anchors. At an anchor
beside a foreign arrest window the per-cell estimate uses its
window weight. The amplitude Cauchy estimate for a pointed piece
supplies the extra factor there, even for its continuation with
all subsequent cut families empty.

For the terms with cut families, keep a connected label \(A\)
and a separate read set \(\mathcal R\) at each link. If \(\Delta\)
is a cut cell, let \(\rho(\Delta)\) be the smallest size of a
connected cut family joining it to a source cell. Here source
cells are the designated cells adjoining the exempt boxes. The
first link for plain interior terms cuts only cells disjoint from
\(\Lambda^{\sim}\), and links two to four leave cells meeting
\(\Lambda\) uncut. Each contributing
connected component of a cut family contains a source cell and
meets the reads, with a source cell at a box meeting the reads
also counted. The latter condition allows the cell's current
parent to be at distance up to \(c_\rho\rho(\Delta)\) from \(A\).
For a read already in the label this distance is at most one.
At an exempt box it is bounded by the box diameter. At the rim
of a component already met by the label it is at most
\(100\check R_k+1\), and the nearest source box is at distance
at least \(9\check R_k\) from the rim. At the first plain interior
link the read set is \(X_v\cup\hat X_v\) in the treated component.
Cells eligible there are current cells. Their parents are
within one cube of \(K_v\), or within the hub diameter plus
one for a source cell seeing the reads across a hub. These
bounds apply on the separated geometry. At each link we must
check them for the actual reads, including reads along a live
strip.

Apply the fundamental theorem of calculus separately in each
nonexempt cut variable. The term with all cut families empty
is the value at the zero cuts. In every other term use the mixed
derivative in precisely the selected variables, integrated over
their real intervals. Cauchy on the common product polydisc of
radius \(e^{\kappa_1}\) gives a factor \(e^{-(\kappa_1-1)}\) per
selected variable, since the circle at a point of \([0,1]\) can
have radius up to \(e^{\kappa_1}-1\). For a boundary-type term
use its supremum on this chart. For a pointed term one can
subtract the comparison value, which is constant in the cut
variables, and use the smaller difference bound. For an interior
term with reads deep in the core, the depth estimate gives the
additional factor \(C e^{-\delta_P\hat D_v/4}\), where \(\hat D_v\)
is the depth of the read hull. The mixed Cauchy argument requires
one product domain for the whole cut family and the source. In
the first two background comparisons of the terms with reads off
\(Z\), we must check this domain by the lower-layer membership
induction on the actual auxiliary carrier, including any portion
inside the treated component. The sealed cut discs for the later
comparisons apply at their own background layers.

At a link with family \(\mathsf y\), the new label is the old
label joined to the current parents of the selected cells. Add
the connecting cubes to the label and across the exempt boxes
used by the family. A connected part with \(n\) cells reaches
the old label by a corridor of length at most \(c_\rho n\),
and the bridges across its boxes cost a bounded number of
cubes per cell. Thus the new tree length is at most
\(d_k(A)+c_l|\mathsf y|\) with \(c_l\) fixed. To sum at weight
\(r\), spend \(c_l r\) per selected cell, leaving
\(\lambda_r=\kappa_1-1-c_l r\) for counting. A connected family
through a cell at depth \(t\) in the grid has at least \(t+1\)
members if it reaches a source. If it sees the reads from a
current parent at distance \(x\), it also has at least
\(x/c_\rho\) members. Thus the exponential reserve pays both the
\(L^{4t}\) fine cells in a parent and the polynomial shell count
in \(x\), as well as the exponential count of connected families
through each cell. With the lower bound on \(\lambda_r\), the
sum for one connected part is at most a small constant times
the number of cubes of \(A\). The sum over unordered sets of
parts is bounded by the exponential of that sum. The constant
can be chosen so the resulting factor is at most
\(e^{d_k(A)+1}\).

For a sum rooted at a prescribed cube, distinguish whether
that cube is in \(A\) or in the addition. In the first case,
after summing the other parts, one more unit of the old-label
weight supplies the distance from the root to the parent of
the cell seeing the reads. The reach correction costs at most
\(c_\rho\rho(\Delta)\), which is paid by the connected-family
reserve. In the second case, mark a connected part of size
\(n\) responsible for the root. A parent of one of its cells
lies within \(c_\rho n\) of the root. There are at most
\(O((1+n)^4)L^{4n}\) choices of parent and fine cell, and
exponentially many connected parts through that cell. One of
its \(n\) cells sees the reads. The factor \(e^{\rho(\Delta)}\)
to recover the reader sum costs at most \(e^n\), so the same
reserve sums over the marked parts. Finally the readers seen
from \(\Delta\) have a label cube among \(O((1+\rho(\Delta))^4)\)
candidates. The factor \(e^{-\rho(\Delta)}\) in their sum absorbs
this count and reduces it to the ordinary rooted label norm.
Write \(Q_r\) for the supremum over a root cube of the sum of
the assigned sizes times \(e^{r d_k(A)}\) over the labels \(A\)
through that cube. Thus a link, including the empty family,
satisfies
\[
 Q_r(\text{after the link})\le (1+C_\Sigma C_e)\,
 Q_{r+2}(\text{before the link}) ,
\]
where \(C_\Sigma\) is the count for the two root positions and
\(C_e\) bounds the number of nearby label cubes with its
\(e^{-\rho}\) weight.

After the last link, fill in the components of \(Z\) that meet
the final label. This adds no cube off \(Z\), and a tree for the
label bounds \(d_Z(Y_4)\). A nonempty cut term filed off \(Z\)
has both a source cell at distance at least \(9\check R_k-1\)
from \(Z^c\) and an exterior cube in its connected label. Its
length is therefore at least \(\check R_k-2\). Spend one unit of
the final weight to extract \(e^{-(\check R_k-2)}\). Summing in
reverse link order gives the rooted bound at weight \(a\) from
the input rooted norm at \(a+9\) for the four links of the plain
interior terms, and at \(a+7\) for the three cut links of the
plain exterior terms with no amplitude. Choose the cut radius
with \(\kappa_1-1\ge c_l(a+9)+\lambda_*\), where
\(\lambda_*\) the counting threshold. The plain interior input
uses the shallow column and the depth-weighted column in the
core. For the exterior input, the root of the intermediate
sum may lie inside \(Z\). Choose an exterior cube on the same
initial label and spend one more unit of weight to sum its
distance from that root. This uses the size column at \(a+8\),
including the exposed difference estimate on its enlarged label.
The two bounds are a polynomial in \(\log x_k\) times the collar
gain.

For a compact point-anchored amplitude piece \(p\), the input
to the later links has label \(K_p\cup\mathfrak h_v\) and reads
the domain \(Y_0(p)\). Keep its amplitude disc of radius
\(\rho_A=\mathsf w_L/(\eta_p\mathsf f_v)^{1/2}\) throughout these
links, where \(\eta_p\) is the decay to the amplitude block and
\(\mathsf f_v\) the decay from the point's window to the source
region. On the common product chart, amplitude Cauchy contributes
\(2/\rho_A\) to the dilated size, and the cut variables have their
own factors \(e^{-(\kappa_1-1)}\). The factor
\(\mathsf f_v^{1/2}\) pays the cubic window weight in the pointed
size column. The factor \(\eta_p^{1/2}\), with a reserve from
the first cut family, sums the possible amplitude blocks in shells
around the cells already reached. This sum is made over the pieces
behind one coarse root, retaining their actual zone and cut-cell
levels. It gives a bound proportional to \(\mathsf w_L^{-1}\) for
the input rooted norm. Three further links and the transfer to
an exterior input root use at most seven more units of weight.
The small factor is already in the amplitude column, including
its empty continuation.

The chain and long terms grouped into bundles use the sum over
blocks and fine cells in the term's own grain before the three
later links. For a bundle the comparison disc acts on that sum.
The small factor belongs to the resulting rooted bundle norm.
For a wrapped term of level \(j\), let \(A_v\) be its connected
label and \(K_v=[A_v]_k\). Its true read consists of \(A_v\) and
the attached older strips, each touching \(A_v\). An attached
strip healed at this step lies in the core of a treated component,
and the adjacent label cube lies in that component's inner box.
A strip that stays live is separated from the treated components.
Thus the true read meets \(Z\) exactly when \(K_v\) does. When
filing a value, adjoin the treated components met by \(K_v\).
This already includes all the healed strips read by the term.
For analytic evaluation also include the live strips touching
the label in the enlarged background domain. In the integration
over a treated component, the fields on the older live strips
are parameters. Two polymer factors that share only these
parameters factorize in that integral. Hence the older live
strips are background reads and the connected label is what
enters the length and compatibility tests.

In the cut expansion, keep the actual read at each link. It
consists of the original reads, the previously selected cells
and their touched boxes, with the rim of the preceding domain
as an overestimate. A cell on a live strip is eligible for a
cut under the same exemption rule as the other cells. Its
parent is within \(101\check R_k+1\) of the old label, while a
family from that strip to a source has size at least
\(10\check R_k\). The analogous rim estimate allows one more
cube. This pays the corridor from the cell to the label by
\(c_l|\mathsf y|\). The rest of the live strip need only be
adjoined to the background read domain. We must check that the
global chain stays in the joined chart for the stored wrapped
term, using inherited tubes on live strips and containment on
healed strips. Keep the twist in the common admitted source
slot and take the stored norm on the complex orbit domain,
since contour gauges can be complex. On that product chart the
same mixed Cauchy bound applies to each tuple of the composed
layer, current and gauge readers, including one whose only cut
is at the first link. At the third layer the support test applies
to the primitive current entering its inverse.

For a term with the whole read in a core, if all cut cells
stay in \(Z\), each connected family seeing the reads stays in
the same component. Its parents and connecting cubes stay in
that box. Induction over the four links then keeps the final
label inside the component. Thus an exterior tuple of such a
term triggers the core depth estimate, and its assigned size in
the exterior sum is the stored size times
\(C_{3F}e^{-\delta_P d^{\wedge}(X_v,\mathsf C_C^c)/2}\), with
\(X_v\) the true read and \(\mathsf C_C\) the core. Use the
unmodified stored size for the other wrapped terms. The depth
factor is assigned only in the majorant for exterior tuples.

There is a separate count for the initial wrapped reads.
Let \(\mathsf b_k(s)\) bound the weighted label column at native
level \(k-s\), with the spare native rate for labels finer than
the probing cell. If the read through a cell of level \(n\)
comes from an attached strip of level \(m\), count labels of
level \(j\ge m\) by a cube touching that strip. At level \(j\)
there are at most \((s_j+3)^4\) such cubes, where \(s_j\) is the
strip side in those units. Since \(s_j\le 101\check R_m\), this
count has no factor for the fine volume of the strip. A native
probe meets at most \(2^4\) live strips, all with \(m\ge n\).
For a read through the label itself, use the native meeting
convention for unlike levels, which requires a common interior
point. The levels \(j\ge n\) have a bounded number of anchor
cubes per level. At \(j<n\) the label
must reach from the probe to an older strip. With \(t=n-j\)
the spare rate bounds the fine-anchor count by a constant
times \(L^{4t}e^{-c(L^t/2-2)_+}\). These are the two parts of
the read count for a fixed probe. On the native cut geometry,
\(k-n+1\le\rho(\Delta)\). For a cell contained in a live strip,
a minimal path to a source must leave the birth-scale collar
of width \(5\check R_m/2\). Before its first source it takes no
hub jump, and in that collar each cell spans at most one birth
unit. In particular \(\check R_m\le2\rho(\Delta)+2\). Multiplying the read count
by \(e^{-\rho(\Delta)}\) absorbs its powers of the age and strip
side and bounds the reception sum by \(C_{\mathsf a}(1)\mathsf b_k(0)\).
For a source cell seeing the read through an exempt box, its
contribution is instead bounded by \(h_\sharp^4Q^\star_{a+9}\),
where \(h_\sharp\) bounds the box side and \(Q^\star\) is the
rooted norm for terms meeting \(Z\) with the assigned sizes
above.

The first-link family includes an empty choice for a term
meeting \(Z\). A term disjoint from \(Z\) enters through one of
four successive background differences, each with a nonempty
first family. Use the reception bound for the nonempty first
families, with at most four copies. The next three links use
the two-unit weight estimate. For an exterior final label the
collar length is at least \(\check R_k-2\). For a label reaching
a live strip this follows from its reach to \(Z\), supplied by
a source cell if the original term was disjoint. Otherwise the
label reaches from an inner box of a treated component to an
exterior cube. Spending one unit of weight yields the exterior
rooted bound
\[
 e^{-(\check R_k-2)}(1+C_\Sigma C_e)^3
 \bigl[(1+4h_\sharp^4 C_\Sigma)Q^\star_{a+9}
       +4C_\Sigma C_{\mathsf a}(1)\mathsf b_k(0)\bigr].
\]
For \(Q^\star\), labels joining a live strip to the treated
region already have a collar length in their own units, which
pays the number of fine roots. Shallow labels attached only
to healed strips are counted through the core on the native
grids, using the collar for levels finer than its rim level.
Deep labels use the depth factor when summing the core cells.
For the deep count, assign the depth to a cell whose closure
contains a point realizing the least distance to the core rim
and whose interior meets the read. Use half of its depth weight
to pay the probe count, and sum the remaining half by the
core shell estimate. Its prefactor is the area of the core
rim box, with side at most \(104L\check R_{h+1}+2\) in level-\(h\)
units. Here \(h\) is the level of the cells on the inner side
of the core rim. These counts bound \(Q^\star_{a+9}\) by a
polynomial in the stage count and the radius at the core rim.
The weighted column, depth estimate and joined chart must
hold for the actual stored terms in this format before using
the bound.

In the total pre-filed family, separate the ordinary values,
the brackets, and the source-free remainder of the Wilson
expansion. The all-empty term of an ordinary exterior unit is
a value. A plain interior term whose domain stays inside its
component goes into the interior volume sum. The first-link
amplitude term and its later nonempty continuations count in
the same amplitude budget, with each tuple charged once. The
nonempty wrapped terms have their own budget. The Wilson family
here consists of the action remainders after removing the explicit
source-weight term at fixed background. The background arguments
remain those of the chain. File the treated-component cell terms
at their native anchors, the terms depending on the whole
component at that component, and the terms in \(\Omega_k\) off
\(Z\) at their exterior cubes. Each resulting cut tuple belongs
to the Wilson family alone. Use the absolute exterior bound
and the summed strip bound \(1+c_I M^5\) for this family from
the Wilson localization estimate in Ba{\l}aban, Comm. Math.
Phys. 122 (1989), pp.~374--375, (1.66)--(1.68). The bound for
the denominator term \(I(U_0)\) in (1.11) is charged once per
component. To use these bounds, the localized Wilson family must
agree with the one in that estimate on the admitted backgrounds.

For the Wilson cell terms of a treated component, the local
bound follows from
\[
 a(W)=-\tfrac12\operatorname{tr}_{\rm norm}((W-1)^2W^{-1}),
\]
where \(\operatorname{tr}_{\rm norm}\) is the normalized trace.
If the plaquette increment at a level-\(n\) cell is at most
\(\widehat\theta(\eta/\ell_n)^2\), the plaquette count cancels
the fourth power of \(\eta/\ell_n\) in the density. Here \(\eta\)
is the fine spacing and \(\ell_n\) the spacing of the cell's
grid. The bound on a cell with action coefficients at most
\(x_n\) is \(6M^4x_n\widehat\theta^2\). In the core subtract the
value at zero cut parameters with the real rim datum fixed.
Plaquettes entirely in zone zero cancel because the two
backgrounds have the same datum there. For the other plaquettes
the linearized holonomy difference is a covariant first
difference of the displacement around the plaquette, with a
term proportional to \(1-\operatorname{Ad}_W\). The latter has
an extra curvature factor. The commutators in the product of
the four bond exponentials are quadratic in the displacement.
Thus if the scaled displacement bound is \(B\) at the cell,
the plaquette difference is at most \(C B(\eta/\ell_n)^2\).
The quadratic density identity then bounds the cell oscillation
by \(C M^4x_n(\widehat\theta B+B^2)\). The admitted exponential
depth estimate for \(B\) makes this sum proportional to the
core boundary count. Off the core use twice the cell bound
and count the preparatory shells. These two counts give a
polynomial in \(\log x_k\) for the Wilson cell descendants.
Four links give those descendants leaving \(Z\) the same
collar gain as other plain interior terms. This estimate for
the cell descendants supplements the absolute bound for the
full Wilson family.

Thus the exterior input norm is bounded by the sum of the value
bound, the bracket bound, and the Wilson bound. Once their admission has been
verified, the plain interior and exterior cut estimates take
one half of the bracket budget, the amplitude terms take one
half, and the wrapped terms take one, for a total bound of two.

Use the prescribed disjoint strips \(C^\sharp\) of width
\(m=\lfloor(\check R_k-1)/2\rfloor\)
about the separated components \(C\) of \(Z\). A pre-filed domain
\(Y_4\) is wall-connected and meets \(Z\) in a nonempty union of
full components. File its term on \(Y\) obtained by adjoining
every strip that shares a cube
with \(Y_4\). This partitions the terms into fibres and preserves
the sum. The true reads are inherited from the terms in each fibre,
including the enlargement of any wrapped read.

For an exterior output \(Y\), filing adds no outside-strip cubes,
so all preimages share a root cube there and a relative tree for
\(Y_4\) also meets every cube of \(Y\setminus Z^\sharp\). Thus
\(d_\sharp(Y)\le d_Z(Y_4)\). A wall path in \(Y_4\) from \(Z\) to
the outside of \(Z^\sharp\) visits an exterior cube touching \(Z\)
and an exterior cube at distance \(m\) from \(Z\). A relative tree
meets both cubes. The distance to \(Z\) changes by at least \(m-1\)
between the two points, so the tree has length at least \(m-1\).
With \(b=3\kappa/2\) and \(\lambda=\kappa/40\), the input rate
\(a=31\kappa/20=b+2\lambda\) therefore gives
\[
 e^{-a d_Z(Y_4)}\le e^{-b d_\sharp(Y)}e^{-\lambda(m-1)}
                         e^{-\lambda d_Z(Y_4)}.
\]
To count preimages of length at most \(n\), traverse a relative tree
twice and cut the walk into arcs of length at most one. Each arc
meets at most \(3^4\) cubes, so the cubes met by the tree form a
touching-connected family of \(O(n+1)\) cubes through the root. A
depth-first walk on this family gives exponentially many choices
in \(n+1\). Select the outside cubes among them, then select the full
components of \(Z\) incident to the selected cubes. Every component
in the preimage is incident to an outside cube, by wall
connectedness and separation of the components. Each outside cube
is incident to at most \(3^4\) components. This count sums the last
factor in the fibre when \(\lambda\) exceeds its entropy rate. Thus
the exterior strip amplitude is at most
\(C\|\mathfrak v\|_a e^{-\lambda(m-1)}\), where
\(\|\mathfrak v\|_a\) is the weighted exterior input norm. The
interior fibre of \(C^\sharp\) is the sum of the inputs contained
in that strip and uses a volume bound. With no prescribed
exterior root the free-family sum at a single link costs
\(e^{d_k(A)+1}\). Applied backward four times starting at weight
zero, it gives \(e^{4(d_k(K_v)+1)}\) times the assigned plain
size. The bound for a wrapped unit, with its value included,
has the additional factor four for the possible first brackets.
The four units of tree-length weight are available at the
unit's native rate, with a fixed coarsening offset absorbed
in the constant.

For the running chain choose one anchor in the filing-level
stage region of each term. At a given stage and anchor the
weighted sum is the corresponding stage-norm sum. At a cell of
zone \(n\), the bound on the number of such stages grows as
\(1+\log(1+\log x_n)\), using the running stage count and
\(x_n\ge x_k-2c_2\). Charge the other ordinary terms at the old
determining cells of their zones. A window cell coarsened by
\(e\) levels has at least
\(L^{4e}\) distinct old cells to charge. This absorbs its cubic
window weight, bounded by \(27\mathsf L_0^{6(e+1)}\), since
\(\mathsf L_0^2\le L/3\). For a frozen chain make this charge at
its own pre-arrest cells. A level-\(m\) stage anchor meeting its
band contains a cell of the band grid, and distinct anchors
at that level contain distinct band cells. If \(j+1\) is the
band level, the number of band cells is at most a constant
times \(\check R_{j+1}^4\sum_X(d'_{j+1}(X)+1)\), where \(X\)
runs over the band-containing components. The number of
possible filing levels is bounded by a constant times
\(\check R_{j+1}\). Thus the frozen volume allowance is a sum
of \(\check R_{j+1}^5(d'_{j+1}(X)+1)\) over the stages and
components of that arrest. The count uses the actual native
anchor and band inclusion for the stored terms.

An exterior bracket filed wholly inside \(C^\sharp\) reaches
a cube of \(C^\sharp\setminus C\). Charge it there using the
exterior rooted sum, which is at most one per cube for the
plain exterior brackets. Add the ordinary value allowance
on those strip cells and one Wilson allowance for the entire
component. This bounds the sum over all inputs contained in
the strip. The finite filing preserves holomorphy on the common
admitted source tube. The exterior gain absorbs the polynomial
logarithmic factors in the pre-filing bound, including those
from the Wilson vertices.
Since the width grows at least as a fixed positive multiple of
\((\log x_k)^{r_{\rm pr}}\) with \(r_{\rm pr}>1\), its gain beats
every fixed inverse power of \(x_k\), also after taking the cube
root of the amplitude.

% Earlier Part I, 23293--23417 and 24942--25544; earlier Part II,
% 17236--17377; earlier Part III, 4480--4799 (strip polymer sum).
% Earlier Part II, 12651--12843 and 14418--14669 (component modulus).
% Earlier Part II, 2663--2680 (unmatched counterterms); earlier Part III,
% 33222--33270 and 34285--34409 (creation deposit and charge unrolling).
To pass from the filed potential to the healed remainder, fix
the standing part of the history. Let \(W\) be the union of the
strips of components that stay as second-class regions in this
operation. Their variables are parameters in the healed sum.
Keep their tests, group delta factors and component integrals
outside that sum. A term on a single strip is local to its own
component and is included in that component's operation. For
the other domains, write the exponential as
\(1+(e^{V^\sharp(Y)}-1)\) at each domain \(Y\).
In a summand the integrations run over the chosen healed
components. Each interaction factor reads the integration
variables of a chosen component only if its domain includes
that component's strip. A connected class consists of the
touching parts off \(W\), also joined whenever one interaction
domain reads more than one of those parts. Each chosen healed
component belongs to exactly one class. The class factors
therefore multiply at fixed exterior data. Conversely, classes
with disjoint and nontouching parts off \(W\) can be assembled
by taking their union, by the locality of the history format.
This gives a hard-core gas whose empty family has value one.
For a nonempty healed set the term with no exterior interactions
still contains the component masses.
The class activity sums the compatible internal choices of
healed components and interaction domains with that class's
domain. Use the component operation of \cref{candidate:large_field}
with its original measure factors, tests and internal history
sum. Its localized action on the indicated reads is needed for
the factorization.

For the component mass in this argument, keep the stored
integration coordinates real on the supports of their original
tests. In the inner Wilson factor the field assembled from
these coordinates is unitary, while the source changes its
weight. Let \(c_j\) be the unshifted nonnegative coefficient
including its cutoff, and \(c_j^w\) the shifted coefficient.
On the admitted source tube \(\operatorname{Re}c_j^w\ge
(1-\theta_W)c_j\) with \(\theta_W\le1/4\). Plaquette by
plaquette, \(a_p(U)\ge0\) for the real stored field, so
\[
 \big|e^{-A(c_j^w,U)}\big|
 = e^{-A(\operatorname{Re}c_j^w,U)}
 \le e^{-(1-\theta_W)A(c_j,U)} .
\]
This comparison also applies inside the nested component
operations. In particular each Wilson test gain retains the
fraction \(1-\theta_W\) of its real exponent. The exterior
pair can be complex. The quadratic form depending on it uses
its source-independent complex-form estimate. The other
quadratics, tree constraints, Jacobian factors and tests keep
their real integration data. The sourced strip interaction
is bounded by the exponential of its supremum on the same
chart. Its ordinary value budget is charged to the original
cells, and the frozen contribution to the stage cells of its
arrest. The strip collar adds a fixed volume factor. The
additional weights of the inner real fluctuation integrals
use their per-cell bounds. For an unmatched counterterm on
the regular part of its cutoff collar, the action on a native
cell has modulus at most \(C\varepsilon(\gamma)^2\). The source
bound on its coefficient is \(b_++2r_{
m src}/3\) on the disc
of radius \(2r_{
m src}\). Thus its modulus costs at most
\(e\) per native cell once
\((b_++2r_{
m src}/3)C\varepsilon(\gamma)^2\le1\). Include
that cost in the existing per-cell allowance. This yields a
majorant before integrating the stored coordinates.

For a new region detected by a refused rim bond at level
\(j\), keep the corresponding gain from the constrained
quadratic form. If the full form has lower bound \(\gamma_0\),
so does its Schur complement on the retained coordinates:
minimizing \(\langle(y,x),\mathcal A(y,x)\rangle\) over the
discarded coordinate \(y\) gives at least \(\gamma_0\|x\|^2\).
With \(\mathfrak b_j=g_j^{-1}\delta_j\), the rim test forces
\(|x_b|\ge\mathfrak b_j/(4\Theta_1)\) at each refused bond \(b\).
Spend one quarter of the quadratic form on the test gains
and retain the other quarter in the density. The resulting
factor at each such bond is
\(e^{-\gamma_0\mathfrak b_j^2/(64\Theta_1^2)}\). Its first
square root pays the sum over refused-bond sets. At fixed
ordinary test region each set is counted once and the sum is
bounded by a product of
\(1+e^{-\gamma_0\mathfrak b_j^2/(128\Theta_1^2)}\) over rim
bonds. The other square root gives an exponent of at least
\(\gamma_0\mathfrak b_j^2/(256\Theta_1^2)\) per covered cube,
since each bond covers at most two. The ordinary fluctuation
tests use disjoint witnessing coordinates in this split. The
rim gain therefore joins the Wilson and ordinary fluctuation
gains in the cascade bound, and its cost is charged at its own
level.

The cascade charge estimate uses the cover by the coarsened
predecessor regions and the current test witnesses, including
the cubes of the refused bonds. Write \(\mathsf P_n=p_0(g_n)\)
and \(\mathsf Q_n\) for the common gain per created cell after
setting aside the refused-bond count. For a created region
\(H\) of level \(n\), the cell count satisfies
\(d'_n(H)+1\le\sqrt{4}\,|H|_n\), where \(|H|_n\) is its cell
count. Split \(\mathsf Q_n|H|_n\) into halves. One half gives
the deposit \(\mathsf Q_n(d'_n(H)+1)/(2\sqrt{4})\) and the
other leaves \(2\mathsf P_n\) when \(\mathsf Q_n\ge4\mathsf P_n\).
In unrolling the charge recursion, a quiet successor spends
its current charge from the balance. A re-created successor
adds its new test gain \(\mathsf P_n\) before spending its
charge. A join collects the predecessor balances and the
birth deposits and releases the older creation reserves
other than the one at the earliest birth of the joined
region. Thus each member keeps the balance and one oldest
reserve of \(2\mathsf P_{j_0}\). With the positive balance on
the admitted family and the internal history count paid in
the per-level charges, the exponent retains that reserve at
the final component. Compare
\(p_0(g_{j_0})\ge p_0(g_k)/(1+\beta_0)\)
along the reference flow. In the complex modulus estimate there
is one remaining factor \(e\), giving the bound
\[
 |\mathbf T'_k(X;w)F|
 \le e\,e^{-2p_0(g_k)/(1+\beta_0)}\sup|F| .
\]
The supremum is on the internal support at fixed external data.
For a product of compatible healed component operations use
this bound successively, provided the admitted chart and the
support bounds are valid with the other integration variables
fixed. The same integrable majorant permits holomorphic
integration of the analytic factors. To apply this estimate
we must identify the actual real Wilson arguments, witnesses,
charge family and common complex chart for the component
integral. These checks concern the stored histories that
enter the gas.

Let \(A\) denote one activity domain and \(d_{\rm out}(A)=d_W(A)\).
The tree for \(A\) off \(W\) is formed from the relative trees
of its chosen interaction domains and a path through the cube
centers of each chosen healed strip. A path in a rectangular
strip of \(N_h\) cubes has length at most \(N_h\) in cube units.
Join these trees and paths along a spanning tree of the class
graph, at cost at most \(2\sqrt{4}\) per join. The number of
cubes off \(W\) is bounded by the sum of the \(N_h\)'s and a
constant times \(\sum_Y(1+d_\sharp(Y))\). Thus in the input
rate \((1+2\beta_s)\kappa\) one can allocate
\((1+\beta_s)\kappa\) to the output length, half of
\(\beta_s\kappa\) to the cube count and the other half to
counting the interaction domains. If \(\alpha^\sharp\) is the
exterior strip amplitude, spend one third of its logarithm
per chosen interaction to pay the fixed join and cube-count
costs. For a healed component the modulus bound pays for the
strip, whose cost is at most \(2\kappa(100\check R_k)^4\) per
box. The positive coefficient
\(2/(1+\beta_0)-3/2\) at \(\beta_0\le1/7\) gives the room for
this cost and the sourced modulus prefactor. With the degree
ceiling the remaining factor is \(e^{-3p_0(g_k)/2}\) per
chosen healed component.

At fixed healed components the remaining interaction weights
have the form \(u_Y=(\alpha^\sharp)^{2/3}
e^{-\beta_s\kappa d_\sharp(Y)/2}\). Every interaction domain
contains an exterior cube off the strips used by that summand.
The rooted domain count bounds \(\sum_Y u_Y\) by a constant
times \((\alpha^\sharp)^{2/3}\) times the number of cubes off
\(W\) in \(A\). If the interaction family is nonempty,
\(\prod_Y(1+u_Y)-1\le(\sum_Y u_Y)e^{\sum_Y u_Y}\). Keep a
factor \((\alpha^\sharp)^{1/3}\) and bound the remaining
prefactor by an exponential in the cube count. The same count
pays the exponential from \(e^{V^\sharp(Y)}-1\). The resulting
cube penalty is small compared with the allocated cube-count
decay. A class containing a strip in \(W\) necessarily has a
nonempty interaction family, so it has the cube-root factor.

A class disjoint from \(W\) contains at least one healed
component, since each interaction meets a component of that
summand. This includes the class consisting of one healed
strip and an empty interaction family. Extract \(e^{-p_0(g_k)}\)
from the product of the component factors and use the remaining
\(e^{-p_0(g_k)/2}\) per component for counting. Each candidate
box is specified by its lowest cell and four bounded integer
side lengths. At fixed \(A\), the number of unordered sets of
\(q\) such boxes is bounded by the \(q\)-th power of the
candidate count divided by \(q!\). The resulting exponential
in the number of cubes is paid by the remaining volume decay.
Thus the activities have rate \((1+\beta_s)\kappa\) in
\(d_{\rm out}\), with amplitude \(e^{-p_0(g_k)}\) off \(W\)
and \((\alpha^\sharp)^{1/3}\) when they meet \(W\).

For the logarithm only connected incompatibility graphs
contribute. On a connected cluster the relative trees join
at the touching cubes off \(W\), with at most \(2\sqrt{4}\)
of extra length per edge of a spanning tree. Reweight each
activity by
\(e^{(1+\beta_s/2)\kappa(d_{\rm out}(A)+2\sqrt{4})}\).
The remaining \(\beta_s\kappa/2\) in the length pays the
exponential count of domains through a cube and the
exponential of their cube count in the polymer criterion
\cite{KP86}. Indeed, each incompatible domain contains an
exterior cube touching one of the cubes of the given domain
off \(W\), and for each such root the weighted series is bounded
by a constant times the maximum activity amplitude. For weak
coupling that constant times the maximum is small, giving
the criterion with size equal to the number of cubes off
\(W\). To bound the logarithmic coefficient at a fixed union,
mark a cube of the union off \(W\). Each cluster has at least
one polymer covering it. Summing the mark over the polymer
indices changes the factor \(1/n!\) into \(1/(n-1)!\) with
the marked polymer first. The rooted cluster majorant then
sums the other polymers and leaves the weighted rooted sum
of this first polymer. In particular, for a union disjoint
from \(W\) its amplitude is \(e^{-p_0(g_k)}\). The logarithm is
absolutely convergent on the admitted complex domain and the
output decay rate is \((1+\beta_s/2)\kappa\).

Classify the resulting domains using the one-layer enlargements
of the standing large-field region and the newly retained
second-class components. Away from them the output length is
\(d_k\) and the prefactor is \(C e^{-p_0(g_k)}\). The spare
\(\beta_s\kappa/2\) pays the rooted sum at rate \(\kappa\),
giving the healed contribution to the inherited remainder
bound. The other terms are new boundary terms. Their bound
uses the relative length off \(W\); the prefactor can be taken
as a fixed multiple of the larger of the two activity
amplitudes. To store them on relative labels, sum the domains
having a given skeleton. The count contributes
\(e^{\lambda_\theta(d_k(A)+1)}\) for the skeleton \(A\), and
\(\lambda_\theta\le\beta_s\kappa/4\) leaves the stored rate
\((1+\beta_s/4)\kappa\). This last step uses the skeleton
length comparison and the relative-format bound for the
actual complex charts.

The dimension-free polydisc lemma,
\cref{lem:polydisc-dimension-free-v7}, gives on the half
polydisc
\begin{equation}
 \sum_i r_i|\partial_iF|\le\tfrac43\sup|F|,
 \qquad
 |F(z)-F(z')|\le\tfrac43\sup|F|
                         \max_i|z_i-z_i'|/r_i .
 \label{eq:flow-RF39}
\end{equation}
One proves it by sending a single unit disc through the point
with coordinate disc automorphisms, choosing their phases to align
the derivative terms, and applying the one-variable Schwarz
inequality. Integrating the derivative bound on the segment gives
the second formula. Neither bound contains the number of earlier
levels. Neither implies age forgetting: the bounded function
\(F(z)=Bz_0/(\sigma_\star\theta_0)\) already has a nondecaying
weighted derivative in its first coordinate.

Here is the needed forced derivative estimate, including the
marked term. Let \(b_i'=d b_i/d\beta\),
\(M_i=\max_{0\le j\le i}|x_j'|\), with \(x_0'=1\).
The frame is fixed, so \(z_j'=-x_j'\).
The tube constants and coupling floor put these evaluation points
in the half polydisc required in \eqref{eq:flow-RF39}.
Since \(\min_{j<i}\theta_j\ge\theta_{i-1}-c_5
\ge\theta_{i-1}/2\) above the floor, \eqref{eq:flow-RF39} and the separate
bare-disc Cauchy estimate imply
\begin{equation}
 \begin{aligned}
 |b_i'|&\le a_iM_{i-1}+d_i,\\
 a_i&\le C\theta_{i-1}^{-3/2}
                      (\log\theta_{i-1})^{\mathfrak p_0},
 &d_i&\le C\varrho^{-1}\theta_{i-1}^{-p_R}.
 \end{aligned}\label{eq:flow-RF40}
\end{equation}
The \(\theta^{-1}\) factor in \(a_i\) is the unmarked
independent-polydisc gain. The factor \(1/\varrho\) in \(d_i\)
comes from the actual bare-variable disc and must not be
replaced by \(1/\theta_i\).

Put \(V_i=1+\sum_{l\le i}|b_l'|\). The differentiated recursion
gives \(M_i\le V_i\), and hence
\begin{equation}
 V_i\le(1+a_i)V_{i-1}+d_i,\qquad
 \sum_{i=1}^K|b_i'|
 \le e^{A_K}(1+D_K)-1 ,
 \quad A_K=\sum_{i=1}^Ka_i,\quad D_K=\sum_{i=1}^Kd_i .
 \label{fc:forced-gronwall}
\end{equation}
Indeed iterating the scalar inequality gives
\(V_i\le\prod_{l\le i}(1+a_l)+
\sum_{l\le i}d_l\prod_{m=l+1}^i(1+a_m)\);
each product is at most \(e^{A_i}\).
No triangular estimate for the physical field response is hidden
in this one-dimensional calculation.

On the reference line, integral comparison yields, uniformly in
\(K\),
\begin{equation}
 A_K\le Ch^{-1/2}(\log h)^{\mathfrak p_0},
 \qquad D_K\le C\varrho^{-1}h^{1-p_R}.
 \label{eq:flow-RF42}
\end{equation}
The condition \(p_R>1\) is exactly what makes the second bound
uniform over all depths for a fixed bare-disc radius.
For the original full Lipschitz numbers one may now take
\(\mathsf a_{[K-i]}=(4/3)a_i+d_i\), once the uniform slope
bound below is imposed. Thus
\begin{equation}
 \sum_{s<K}\mathsf a_{[s]}
 \le C\{h^{-1/2}(\log h)^{\mathfrak p_0}
                   +\varrho^{-1}h^{1-p_R}\}.
 \label{eq:flow-RF42a}
\end{equation}
The sharper original \(O(g(\log g^{-2})^{\mathfrak p_0})\)
rate follows when \(p_R\ge3/2\), in particular in the
intended large-\(\kappa_0\) regime. For merely \(p_R>1\),
retain the general forced rate displayed here.
Choose \(h\) large enough that
\(e^{A_K}(1+D_K)-1\le1/3\). Then
\(|x_j'-1|\le1/3\) at every admitted prefix. On real intervals,
\(2/3\le x_j'\le4/3\).
The statement is uniform only where the analytic inputs and
room conditions hold. It does not identify different frames.

The stronger independent marked-coupling polydisc route in the
comparison appendix remains available separately. Under its
actual hypotheses the marked row has order
\(\theta^{-p_R-1}\log^C\theta\); it is not the premise used
in \eqref{eq:flow-RF40}--\eqref{eq:flow-RF42}. Both routes require their actual domains.
In particular a fixed physical-source disc alone does not
supply the unmarked independent-variable gain.

\paragraph*{Canonical tube and conditional tuning theorem 3.2.}
For a general finite reference sequence of length \(n\ge1\),
the largest permitted tube is
\begin{equation}
 \mathsf t_j^*=
 \min_{1\le t\le\min\{j+1,n\}}\mathsf t^\sharp(\theta_t).
 \label{eq:flow-RF43}
\end{equation}
This is just the combination of all preceding upper bounds and
the requirement that tube radii do not increase with the level.
For \(n=0\) the largest radius is instead
\(2\mathsf t^\sharp(\theta_0)\). The exponent \(q_{\rm t}\)
and the coefficient \(\varpi\) are the ordinary tube constants,
not an insertion rate. The floor is chosen with
\(h\ge e^{2q_{\rm t}}\) and \(c_5/h\le1\) when the terminal
comparison is used.
A constant-tube criterion follows by setting
\(\mathsf K(a)=\max\{1,(4a/e)^a\}\), with
\(\mathsf K(0)=1\). Maximizing \(t^ae^{-t/4}\) gives
\((\log y)^a\le\mathsf K(a)y^{1/4}\) for \(y\ge1\).
In particular
\(\mathsf K(2q_{\rm t})=\max\{1,(8q_{\rm t}/e)^{2q_{\rm t}}\}\).
If the frame floor is \(x_\star\ge e\) and
\(\mathsf K(2q_{\rm t})t\le\varpi x_\star^{3/4}\), then
\begin{equation}
 t\le\frac{\varpi y^{3/4}}{\mathsf K(2q_{\rm t})}
       \le\varpi y(\log y)^{-2q_{\rm t}}
       \qquad(y\ge x_\star).
 \label{eq:flow-RF43a}
\end{equation}
Thus the constant sequence of radius \(t\) is a tube; its
criterion uses the floor rather than the frame length.
On the line in \eqref{eq:flow-RF36}, set
\begin{equation}
 \begin{aligned}
 \iota_k&=\epsilon_k+\epsilon'_k,\qquad
 T_K=\varrho+c_2,\qquad
 T_k=T_K+\sum_{j=k}^{K-1}\iota_j,\\
 F_k(\beta)&=\theta_k-x_k^\natural(\beta)+s_k,\qquad
 \widehat{\mathsf t}_k=T_k+c_2+2R',
 \qquad R'=4\varrho .
 \end{aligned}\label{eq:flow-RF44}
\end{equation}
This is the canonical tube: in depth notation
\(T_{[s]}=T_{[0]}+\sum_{u<s}\iota_{[u]}\).
The correction by \(s_k\) removes the bounded one-loop
oscillation, giving
\(|F_{k+1}-F_k|\le\iota_k\).
The accumulated drift satisfies
\(\sum_{j=k}^{K-1}\iota_j
\le C\theta_k^{1/2}(\log\theta_k)^{\mathfrak p_0}\)
in the stated power regime. For sufficiently large \(h\),
uniformly over \(K\) and \(k\), this and the fixed additive
terms are smaller than
\(\varpi\theta_{k+1}(\log\theta_{k+1})^{-2q_{\rm t}}\).
The same holds at the terminal convention. Thus
\(\widehat{\mathsf t}\) is an admitted tube.
The choice is made after \(\varrho\), not by shrinking the
prescribed offset range as \(K\) grows.

We spell out the domain premise needed by the interval argument.
At a prefix satisfying \(|F_j|\le T_j\), the actual array is
defined with room \(2R'=8\varrho\). The next coefficient is
formed using that prefix before the new large-field operation.
The uniform bounds and the marked bare-disc hypothesis above
apply on a neighborhood of every such prefix. Whenever the
computed next value satisfies \(|F_{k+1}|\le T_{k+1}\), the
next array closes with the same room. These are the actual
prefix formation/closure hypotheses, not consequences of
the scalar intermediate value theorem.

Define \(I_0=\{\beta:|F_0(\beta)|\le T_0\}\).
Since \(F_0=\theta_0-\beta\), this is a compact interval
mapped decreasingly onto \([-T_0,T_0]\). Choose the floor so
all its initial weights are admitted. Inductively suppose
\(I_{k-1}\) has this property for \(F_{k-1}\), and every
earlier prefix lies in its required room. The preceding
domain premise defines \(b_k\), hence \(F_k\), on a
neighborhood of \(I_{k-1}\), even before restricting its
new output to the target window. \eqref{eq:flow-RF40}--\eqref{eq:flow-RF42}
make \(F_k'=-x_k'\) strictly negative there.

At the left endpoint,
\(F_{k-1}=T_{k-1}=T_k+\iota_{k-1}\), so
\(F_k\ge T_k\). At the right endpoint the corresponding
inequality is \(F_k\le-T_k\). Continuity supplies unique
points with values \(T_k\) and \(-T_k\), and strict
monotonicity orders them. Their intervening compact interval
\begin{equation}
 I_k=\{\beta\in I_{k-1}:|F_k(\beta)|\le T_k\}
 \label{eq:flow-RF45}
\end{equation}
is mapped decreasingly onto the whole target interval.
On it
\(|\theta_j-x_j|\le T_j+c_2\) for \(j\le k\), leaving
room \(2R'\) in \eqref{eq:flow-RF44}. The actual closure premise therefore
permits the induction to continue. This proves existence and connectedness of the interval
from the closure premise.

At \(K\), \(x_K=h+s_K-F_K\). Since \(|s_K|\le c_2\),
its image contains every \(x\) with \(|x-h|\le\varrho\).
Each such value has exactly one preimage in \(I_K\).
Let \(\gamma_U\le\gamma_J\le\gamma_H\) impose the preceding
conditions for the fixed \(\varrho\), with \(\gamma_J\) also
imposing the admitted source analytic domain. For every \(K\)
and \(0<g\le\gamma_U\), the conditional tuning theorem is
\begin{equation}
 x_K^\natural(\beta_K^\natural)=g^{-2},\qquad
 \tfrac12c_0K-c_5\le\beta_K^\natural-g^{-2}
                       \le\tfrac32c_0K+c_5 .
 \label{fc:tuned-coupling}
\end{equation}
The common companion makes this the same number on every
admissible torus. Uniqueness has been proved in the admitted
canonical interval, not among arbitrary disconnected
presentations at all bare couplings.
Also
\(|\beta_K^\natural-h-c_0K|\le2c_2+\sum_k\iota_k\).
The single-run spacing estimate is
\begin{equation}
 \tfrac12c_0(j-i)-c_5
 \le x_i^\natural-x_j^\natural
 \le\tfrac32c_0(j-i)+c_5\qquad(0\le i\le j\le K).
 \label{eq:flow-RF46a}
\end{equation}
It follows by summing \(|b_l-b_l^0|\le c_0/2\).
It gives \eqref{fc:tuned-coupling}, but gives no
bound between the tuned numbers for two different cutoffs.

The bare offset \(t\) and terminal coordinate \(x\) are
distinct:
\(x_K^\natural(\beta_K^\natural-t)=h-t+O(c|t|)\)
when the slope error is at most \(c\). The reference frame
is independent of \(\beta\); the self-calibrated frame
\(\theta_j=x_j(\beta)\) and its frozen shifted run are
other applications of the uniform prefix induction.
For a constant source \(\zeta\), the frozen shifted
weights are \(x_j(\beta)-\zeta_j\), and their slope
normalization is \(x_j'(\beta)\). Identifying this derivative
with a particular F3 witness still requires that witness's
admitted run and comparison hypotheses.

\paragraph*{A finite nontriangular comparison lemma.}
The stronger two-run route uses an estimate different from
\eqref{fc:forced-gronwall}. Fix integers \(n\ge1\), \(0\le r_0<n\), numbers
\(0<\rho<1\), \(\Pi,\mathsf M\ge0\), nonnegative
\(y_w\) for \(r_0<w\le n\), and nonnegative \(D_u\),
\(k_{u,w}\), \(h_u,h'\) for the finite indicated ranges.
Put
\begin{equation}
 S_w=\frac{\rho^{\,n-w+1}}{1-\rho}+(w-r_0)\Pi .
 \label{eq:flow-RF47}
\end{equation}
Suppose, for \(r_0\le s<n\) and \(r_0\le u<n\),
\begin{equation}
 \begin{aligned}
 y_{s+1}&\le\sum_{u=r_0}^sD_u,\\
 D_u&\le\mathsf M(\rho^{\,n-u}+\Pi)
                           +\sum_{w=u+1}^n k_{u,w}y_w,\\
 \sum_{w=u+1}^nk_{u,w}S_w
       &\le h_uS_{u+1}+h'\Pi,\qquad
 \sum_{u=r_0}^{n-1}h_u+h'\le\tfrac12 .
 \end{aligned}\label{eq:flow-RF48}
\end{equation}
Then \(y_{s+1}\le2\mathsf M S_{s+1}\).
Indeed \(S_w>0\) and increases with \(w\). Let
\(Z=\max_{r_0<w\le n}y_w/S_w\), a finite number.
The forcing sum is at most \(\mathsf M S_{s+1}\).
The kernel sum is at most
\begin{equation}
 Z\left(\sum_{u=r_0}^s h_uS_{u+1}
                 +(s-r_0+1)h'\Pi\right)
 \le\tfrac12 Z S_{s+1}.
 \label{eq:flow-RF49}
\end{equation}
Taking the maximum gives \(Z\le\mathsf M+Z/2\), as
claimed. The sums couple a difference to later depths,
so this is weighted-max absorption, not a triangular
recursion. Its actual profile-kernel bound in \eqref{eq:flow-RF48}
remains a substantive two-run input. Bounds on increments
alone do not establish that input.

\paragraph*{Statement 3.6: real presentations of the same measure.}
A real presentation over a frame is an actual inductive array
whose weights remain in the tube. At each level it carries one
Euclidean member, with one scalar coefficient common to all
sites and no defect, and one shift \(\mathsf c_j\). These shifts
obey the admitted \(\bar c\) bound, with the frame floor chosen
in terms of \(\bar c\). Members, companions and coefficients
are independent of the admissible torus and real at real fields.
The all-small-field label is the label whose regions fill the
torus. In its exponent the action coefficient is
\(x_j-\mathsf c_j\). The weight actually used in the Gaussian
step is \(x_j\), and \(\mathsf c_jA\) is carried in the local
action. A spatially varying source array with defects is not
a presentation of this special kind.

Write \(b_j=x_{j-1}-x_j\), let \(\widehat b_j\) be the
coefficient actually extracted at step \(j\), and set
\(u=\mathsf c_0\). Coefficient bookkeeping gives
\begin{equation}
 \mathsf c_j=u+\sum_{l=1}^j(\widehat b_l-b_l),\qquad
 x_j-\mathsf c_j=x_{j-1}-\mathsf c_{j-1}-\widehat b_j .
 \label{eq:flow-RF50}
\end{equation}
In particular \(\widehat b_j\) need not equal the planned
decrement \(b_j\) when the shifts vary with the level.

The comparison hypothesis is equality of the all-small-field
measures at every level, on every admissible torus, for two
real presentations over the same frame. Require hereditary
open windows \(\mathfrak O_m\ni\mathbf1\) of block fields:
their backgrounds lie in every needed core, all the relevant
tests equal one, and
\(\mathcal M^iU_m(V)\in\mathfrak O_i\) when \(V\in\mathfrak O_m\)
and \(i<m\). Also require the common connected analytic
fluctuation domain on which continuation from such a window
to the next actual integral is used. These are properties of
the actual maps and tests.

Under these hypotheses the conclusions are
\begin{equation}
 \begin{aligned}
 x_j-\mathsf c_j&=x'_j-\mathsf c'_j,\\
 \widehat b_j&=\widehat b'_j,\\
 \sum_X(E^{(j)}-E'^{(j)})(X,U)
   &=\tfrac32|\mathfrak X_{j-1}|
                         \log(x_{j-1}/x'_{j-1}).
 \end{aligned}\label{eq:flow-RF51}
\end{equation}
The first line holds through level \(n\), and the coefficient
equality through \(n+1\) only when both presentations admit
that next step. The total identity is evaluated at backgrounds
\(U=U_m(V)\), \(j\le m\le n\), of the common windows.
On those backgrounds the remainder totals agree, and the
vacuum constants agree. When the total suppresses \(z\),
\(E^{(j)}(X,U)=\sum_{z\in X}E^{(j)}(X,P_U,z)\).
No assertion identifies individual localized terms of the two
presentations.

Here is the induction. At the bare level equality identifies
the effective weight \(\beta-u\), since the action is not a
constant function. On an open small-field window the positive
continuous densities agree everywhere if they agree almost
everywhere, because product Haar measure has full support on the
open window. Their logarithmic exponents are real, so their
equality involves no integer ambiguity. Analytic continuation
on the admitted fluctuation domain therefore identifies the
next unnormalized integral. The normalized Gaussian in each
presentation uses that presentation's own weight. Its
normalization has the term
\(-3|\mathfrak X_{j-1}|\log(x_{j-1}/\theta_{j-1})/2\).
It follows that the next regular totals differ by the
field-independent constant displayed in \eqref{eq:flow-RF51}.
The common-companion sum rule makes their extracted
coefficients equal. \eqref{eq:flow-RF50} then identifies the next
effective weights. The same physical formation identities on
the hereditary windows identify remainder totals and the
once-counted vacuum constants. This completes the induction.
It uses all-torus compatibility to extract a coefficient
from a constant total; one finite torus alone is insufficient.

Apply the result to a real \(\beta\) whose prefix along
\(\theta\) has room \(2\varrho\), and a real
\(|\zeta|<\varrho\) for which both constructed presentations
are admitted. One presentation starts at \(\beta-\zeta\)
with no shift; the other starts at \(\beta\) with constant
source \(\zeta\) and running shift \(\zeta_k\). The exact
identity is
\begin{equation}
 x_k^\theta(\beta-\zeta)
                  =x_k^\theta(\beta)-\zeta_k .
 \label{eq:flow-RF52}
\end{equation}
The slope bound keeps the running shift at most
\(4|\zeta|/3\) in its stated range. For
\(\theta=x(\beta)\), this describes the frozen shifted run,
not the self-calibrated run of \(\beta-\zeta\).
For the reference line it identifies the two presentations
of \(x_k^\natural(\beta-\zeta)\).

Complex gluing also uses one frame. First restrict the real
centres to the canonical interval \(I_K\) constructed above,
or to one of its admitted prefix intervals. Each source disc
is supplied by the actual analytic source construction.
Where a real point in that disc also belongs to the canonical
interval, \eqref{eq:flow-RF52} identifies its continuation with the real run.
If two discs about such centres overlap, their real overlap
contains an open interval in \(I_K\): the segment between
the centres lies in \(I_K\), and the open real disc intervals
overlap along that segment. For coincident centres at an
endpoint, a one-sided open interval inside the nondegenerate
\(I_K\) suffices. The identity theorem therefore identifies
the coefficient functions on the full complex overlap.

The union carries a single holomorphic bare-coupling
function. Remainder totals on common windows glue likewise;
regular totals glue after accounting for their explicit
Gaussian constant. This proves uniqueness of the supplied
continuations, not their original bounds or radii.

For arbitrary admitted real centres, every overlap used for gluing
must contain a common open interval of admitted real runs, or the
prefix/source-disc construction must prove real admission there.
This condition requires a separate proof when the admitted real set
has disconnected components.
Geometric intersection of two discs alone does not settle
it. With the stated admission property the same identity
theorem proof applies. Nothing here identifies runs made
with different frames, individual localized terms about
different centres, or functions on independent-current
cores from equality only on a physical graph.

\subsubsection{Units, common prefixes and change of reference}
\label{sec:flow-units-common-prefixes-and-change-of-reference}

\paragraph*{Statement 3.7: the finite dilation identity.}
For every integer \(b\) with \(K-b\ge0\) and \(L^b\ell\)
in the compatible torus family, the map \(D_bx=L^bx\)
identifies the bare lattices at \((K,\ell)\) and
\((K-b,L^b\ell)\). For negative \(b\) use the inverse of the
corresponding positive relabeling. The exact result of
\cref{ch:flow-arithmetic} is
\begin{equation}
 \begin{aligned}
 Z_{K,\ell,\beta}(w)
     &=Z_{K-b,L^b\ell,\beta}(w^{[b]}),\\
 S^T_{n;K,\ell,\beta}(\underline f)
     &=S^T_{n;K-b,L^b\ell,\beta}(\underline f^{[b]}),
 &f^{[b]}(x)&=f(L^{-b}x).
 \end{aligned}\label{eq:flow-RF53}
\end{equation}
There is no spacing prefactor in
\(\mathcal O(f)=\sum_pf(x(p))a_p\).
The map transports bonds, plaquettes and Haar variables
bijectively; both sides are the same finite integral.
The partition identity is entire in its finite source
variables and bare coupling. Logarithmic derivatives use
its nonzero domain. The same transport identifies periodic
Gibbs states on the two infinite lattices.
Radii and distances multiply by \(L^b\),
\(\operatorname{Lip}f^{[b]}=L^{-b}\operatorname{Lip}f\),
and
\(\|f^{[b]}\|_\sharp\le\max\{1,L^{-b}\}\|f\|_\sharp\).
The free bilinear form is invariant. This proves neither
convergence nor agreement of two RG constructions.

\paragraph*{Statement 3.8: what the step reads.}
Construction naturality is the additional assertion that
step \(k\) reads only its actual tower through \(k\),
partitions, frame and weight prefixes, earlier outputs
and fixed supply rules. Those rules are made for the
lattice, not for its unit of length. For two presentations
of the same bare lattice, their images under \(D_b\)
therefore identify the histories and their measures,
regions, healing data and local nonwrapping terms through
level \(K-b\), provided the prefixes agree. This is required
for both fixed-frame and self-calibrated runs. Separate
arbitrary choices on the two lattices would not establish it.

For a self-calibrated run this permits continuation from
first passage at \(g\) to first passage at \(g'\), provided
the prolonged run remains below the anchor ceiling and
the rescaled torus is admissible. Its history law through
the earlier level is unchanged. For reference runs,
\(\theta_{[r+s]}(g)=\theta_{[r]}(g_{[s]})\);
the prefix at \(g\), down to depth \(s\), is the
reference run at \(g_{[s]}\) of length \(K-s\), on
\(T_{L^s\ell}\). Its deviations, increments and tube
are the same. The implication
\(\ell\in\mathfrak L(g)\Rightarrow
L^s\ell\in\mathfrak L(g_{[s]})\) uses the stated torus
floor and parameter choices. Here the needed scalar comparison is
\(R(y+c_0)\le LR(y)\) for \(y\ge h\) above the chosen floor;
it is a comparison of the function \(R\), whereas a general run uses
its running-minimum values \(\check R_j\).

Keep the beyond-first-passage scope separate. If
\(\beta\in\mathcal W_{K'}(g)\) and \(K=K'+q\), \(q\ge0\),
then \eqref{eq:flow-RF53} identifies the cutoff-\(K\) measure with the
first-passage measure on \(T_{L^q\ell}\).
Estimates for it require \(L^q\ell\in\mathfrak L(g)\).
A torus volume factor becomes
\(\exp[nc_V(L^q\ell)^4]\), whereas a per-ball estimate
has its radius and separation changed to \(L^qR,L^qd\)
and no side dependence in its constant:
\begin{equation}
 |S^T_{n;K,\ell,\beta}(\underline f)|
 \le A(L^qd;L^qR)B(L^qR)^n(n!)^{\varkappa_\star}
                         \prod_i\|f_i\|_\sharp .
 \label{eq:flow-RF54}
\end{equation}
The stated lower torus floor for that bound also applies.
For \(q<0\), use the identity in reverse, require
\(L^q\ell\in\mathfrak L(g)\), and allow a norm cost
at most \(L^{|q|}\) per test. Applying this to tuned
offsets requires the separate comparison
\(|K(\beta_K^\natural-t,g)-K|\le C(\varrho)\).
It does not follow from tuning at a frozen frame.

\paragraph*{Statement 3.9: all four reference-change targets.}
Let
\begin{equation}
 S_\ell^{[t]}[g]=
   \lim_KS^T_{\cdot;K,\ell,\beta_K^\natural(g)-t},
 \qquad \mathsf J=[-\varrho/4,\varrho/4].
 \label{eq:flow-RF55}
\end{equation}
Existence along the whole sequence, its common complex
offset strip and moving-offset continuity require E3 and
the terminal-coordinate comparison in
\cref{ch:comparison-capsule-v7}.
The strip width and the cutoff floor may depend on the
fixed torus. The coordinates are
\begin{equation}
 \Xi_K(t)=x_K^\natural(\beta_K^\natural-t),\qquad
 \Xi_{\infty,s}(t)=
   \lim_Kx_{K-s}^\natural(\beta_K^\natural-t).
 \label{eq:flow-RF55a}
\end{equation}
The depth limit requires the matched comparison. Terminal
matching uses \(n_K=\Xi_{K+1}^{-1}\circ\Xi_K\) on its
admitted common strip. Its near-identity and summable-error
properties remain E3 hypotheses, not consequences of the
existence of individual tuning inverses. First choose \(\varrho\) large compared with
\(c_2,B_\sharp+1\) and a bound for the tuned increments;
then choose \(g\) small. The bound for those increments
is a two-run comparison requirement fixed before
\(\varrho\), not a consequence of one-run spacing.

(i) At fixed depth \(s\), the same bare coupling has two
offset labels:
\begin{equation}
 \begin{aligned}
 \widehat t_{K,s}
   &=\beta^\natural_{K-s}(g_{[s]})-\beta^\natural_K(g),\\
 \beta^\natural_K(g)-t
   &=\beta^\natural_{K-s}(g_{[s]})
                              -(t+\widehat t_{K,s}).
 \end{aligned}\label{eq:flow-RF56}
\end{equation}
Require the actual two-run estimate
\(|\widehat t_{K,s}|\le
C(c_2+\sum_{K-s\le j<K}\iota_j)\), its summable
cutoff increments, and the comparison domain on which
both presentations exist. The initial comparison requires
\(2c_2+\sum_{K-s\le j<K}\iota_j\le\varrho\); a smaller
fixed fraction of that bound keeps the offsets in the inner window.
For \(g\) small in terms of
the fixed \(s\), these offsets lie in
\(\frac12\mathsf J\) and converge to \(\widehat t_s\).
The whole-sequence moving-offset limit then gives
\begin{equation}
 \begin{aligned}
 S_\ell^{[t]}[g](\underline f)
 &=S_{L^s\ell}^{[t+\widehat t_s]}[g_{[s]}]
                                  (\underline f^{[s]}),\\
 \widehat t_{s+s'}(g)
 &=\widehat t_s(g)+\widehat t_{s'}(g_{[s]}).
 \end{aligned}\label{eq:flow-RF57}
\end{equation}
Here \(t\in\frac12\mathsf J\); each depth condition
must hold at its displayed reference coupling.
The condition on \(g\) is an upper bound, never a
lower bound. The terminal coordinate on the right is
the limiting depth coordinate \(\Xi_{\infty,s}(t)\).
Its existence belongs to the matched comparison.

(ii) For arbitrary sufficiently small \(g,g'\) and
\(t\in\mathsf J\), retain the stronger target that some
\(t'\in\mathsf J\) and integer \(b\) satisfy
\begin{equation}
 S_\ell^{[t]}[g](\underline f)
       =S_{L^b\ell}^{[t']}[g'](\underline f^{[b]}).
 \label{eq:flow-RF58}
\end{equation}
Its torus conditions are both
\(\ell\ge\ell_\star(g)\) and
\(\ell\ge
L^{(|g'^{-2}-g^{-2}|+C)/\lambda_\sharp}
\ell_\star(g')\).
The sign of \(b\) is that of
\(g'^{-2}-g^{-2}\), up to a bounded ambiguity.
The number of intervening passage levels lies between
\(|g'^{-2}-g^{-2}|/B_\sharp-1\) and
\((|g'^{-2}-g^{-2}|+C)/\lambda_\sharp\);
\(|b|\) differs from it by \(O(1)\).
This assertion needs recognition of the same bare run
in both calibrated families, not merely \eqref{eq:flow-RF53}.

(iii) Choose once an anchor
\(0<\gamma_\star\le\gamma_J\), sufficiently small in
terms of the fixed large-field window dials \(\mathsf w,c_a\). For any proposed
\(\mathbb R^4\) limit at \(g\le\gamma_\star\), retain
the prolonged-run target: after a subsequence,
\(K(\beta_j,\gamma_\star)-K_j=s_\star\) is constant,
of order \(g^{-2}-\gamma_\star^{-2}\), and
\begin{equation}
 S_n(\underline f)=S'_n(\underline f^{[-s_\star]}).
 \label{eq:flow-RF59}
\end{equation}
The anchor limit uses the same bare-coupling sequence.
\eqref{eq:flow-RF59} is the exact image under the point map
\(x\mapsto L^{s_\star}x\), expressed by the test
\(f^{[-s_\star]}(x)=f(L^{s_\star}x)\). No extra spacing
factor is inserted in the observable.
This preserves the target for arbitrary defining
sequences carrying the required construction and
recognition hypotheses. It does not give an automatic
mirror/cone theorem for unrestricted growing-torus
limits. Dilation-invariant properties need be checked
at only one reference coupling once this target is
established.

(iv) The bare tuned increments
\(\alpha_K=\beta^\natural_{K+1}-\beta^\natural_K\)
are a new notation, unrelated to the field radii.
The proposed conclusion is
\begin{equation}
 \alpha_K\longrightarrow c_0,\qquad
 S_\ell^{[t]}(\underline f)
     =S_{L^b\ell}^{[t-bc_0]}(\underline f^{[b]})
 \quad(b\ge1,\ t,t-bc_0\in\mathsf J).
 \label{eq:flow-RF60}
\end{equation}
For completeness the shape convention is
\(\mathfrak T_{\mathbb R^4}(d;\vartheta)\): both profiles are
nonnegative nonzero Lipschitz functions supported in radius-\(d\)
balls, with support distance in \([d,6d]\), and each obeys
\(\int f\ge\vartheta d^4\|f\|_\infty\) and
\(\operatorname{Lip}f\le\|f\|_\infty/(\vartheta d)\).
Clause (iv) retains the condition
\(\mathfrak T_{\mathbb R^4}(1;\vartheta)\ne\emptyset\) and
\(g\le\gamma_W(\vartheta)\). The proof also needs the actual
uniform smooth-profile envelopes on the pairs used in its
iteration; class nonemptiness by itself is not an envelope estimate.
The comparison appendix proves this implication with
the actual two-run depth/profile estimates, summable
comparison errors and uniform two-point envelopes.
Those envelopes are uniform in the offsets and tori
needed for repeated blocking; the fixed shape class
must be nonempty. Two distinct cluster increments
would give an offset period by analyticity, and
repeated blocking would contradict the uniform
short-distance envelopes. The unique increment
limit is then \(c_0\) by the averaged one-loop law.
The same argument yields
\(\beta^\natural_{K+b}-\beta^\natural_K\to bc_0\)
for every fixed integer \(b\).
The exact input hypotheses remain in force; finite
dilation and weak numerical inequalities do not
establish them.

\subsubsection*{Vacuum constants and the interfaces used later}

The vacuum recursion is \eqref{eq:comparison-capsule-CP1} of
\cref{ch:comparison-capsule-v7}.
Each assigned step number, Gaussian local number and
regular/remainder value at \(P_\circ\) is counted once.
Using the normalized \(N\) requires adding its removed
local numbers; using \(N^\circ\) does not. At fixed
calibration these numbers have zero source derivative.
They are not additional freely chosen normalizations.

The further vacuum estimates required are
\begin{equation}
 \begin{aligned}
 \log(\mathsf c_k(\theta_k)Z^{(k)}(\mathbf1))
      &=\mathsf N^{(k)}
               +\sum_X\sum_{z\in X}\mathsf n^{(k+1)}(X,z),\\
 \mathsf n_k^{\rm tot}
      &=36N_{k+1}\log L
        -\tfrac12\log\det(\mathsf C^{\mathsf T}
                             \Delta^{(k)}\mathsf C)|_{\mathbf1},
 &|\mathsf n_k^{\rm tot}|&\le CN_{k+1},\\
 \mathsf n_0^{\rm tot}
      &\le-3(L-1)^4N_1\log2<0 .
 \end{aligned}\label{eq:flow-RF61}
\end{equation}
Here \(N_j\) counts sites, \(\mathsf c_k\) is the
actual change-of-variable normalization and
\(\mathsf N^{(k)}\) is the free-bond/coarse-site
count, with \(-3\log\theta_k/2\) per free bond.
The map \(\mathsf C\) is the actual linear embedding
of nonpivot coordinates into full constrained
fluctuations; this determinant is on the free color
coordinates. It is not the covariance \(C^{(k)}\).
The local numbers are Gaussian values before
normalization, real and coupling independent,
with common nonwrapping lifts and the Gaussian
domain bound. Their coefficients vanish.
The count alone is not the entire step number.
The determinant estimates and first-step sign in
\eqref{eq:flow-RF61} remain actual Gaussian targets, beyond
the once-only accounting identity.

The explicit leading vacuum requirement is
\begin{equation}
 \left|E_{K,\ell,\beta}
   +\tfrac32(N^2-1)(L^4-1)
        \sum_{j=0}^{K-1}(\log x_j)N_{j+1}\right|
                  \le CN_0 ,
 \label{eq:flow-RF62}
\end{equation}
uniform in the admitted runs, tori and cutoffs.
It is not proved merely by naming a field-independent
constant.

Choose the constants in the following order. Fix the geometric
exponents and \(L\), with the Gaussian constants and
one-loop data; then choose \(\kappa\) above the domain
entropy, the required lower block scales and \(M\).
Choose \(\sigma_\star\), absolute radii below \(c/M\),
\(K_{\mathbf M}\), and
\(E_0\ge\max\{2E_{\mathsf N},4K_{\mathbf M}\sigma_\star\}\).
Next choose \(N_\star\), the ceiling on \(\mathsf s\),
the unmarked constants, \(C_\beta\) and the weak
triple. Compute the marked-load and tube constants
at that triple before choosing \(\gamma\), the
strong triple, \(\gamma_H\) and the smaller drift
ceiling. The source radius and comparison windows
have their stated place before the final smallness
of \(g\). Only then are \(\ell,\beta,K\) varied.
Each claimed uniformity depends on respecting this
order.

For the germ comparisons, let \(\mathfrak P_n(d;\vartheta)\)
consist of real Lipschitz tuples with disjoint supports in one ball
of radius \(5d\), satisfying
\(\operatorname{Lip}f_i\le\|f_i\|_\infty/(\vartheta d)\).
For \(n=2,3\), a modulus \((C,r_0)\) for \(\sigma\approx_n\sigma'\)
means
\begin{equation}
 |(\sigma-\sigma')(\underline f)|
       \le C r^{4-n}\prod_i\|f_i\|_\infty
 \quad(0<r\le r_0,\quad
       \underline f\in\mathfrak P_n(r;\vartheta)\cap C^2).
 \label{eq:flow-RF62a}
\end{equation}
A comparison for every \(C^1\) tuple is stronger and implies this
one on the \(C^2\) subclass. The converse requires its own density
argument with the stated shape bounds. Transitivity adds the
constants and takes the smaller radius. If
\((\mathsf T_b\sigma)(\underline f)=
\sigma(\underline f^{[b]})\), the modulus becomes
\((L^{(4-n)b}C,L^{-b}r_0)\), directly by substitution.
This accounts for the factor in an anchor comparison and keeps
the smoothness quantifiers of the short-distance input explicit.

Later applications use the following hypotheses.
Beyond-first-passage data use \eqref{eq:flow-RF54}; prolonged runs
use statement 3.8; fixed-depth reference prefixes
and terminal coordinates use \eqref{eq:flow-RF57}.
The germ/rigidity route in the comparison appendix
also carries its own shape and second-difference
hypotheses. In the anchor unit its correlation
constants acquire \(L^{-(4-n)s_\star}\).
Blocking and distinct-member comparisons use
\eqref{eq:flow-RF60}, not just a relabeling of \(K\).
The sub-unit cells of the angular route keep
their separate moment, mirror and ultraviolet
interfaces. The ordinary interacting source construction,
clean-class two-point extraction, cited reconstruction
theorem and full gauge-theory identification
retain their stated hypotheses.


% END MODULE chapters/06_flow.tex

% BEGIN LAYERED EXPANSION physical_gaussian_gauge_change.tex
\subsection{Gauge changes in the physical Gaussian integral}
\label{pgg:section}

The Gaussian normalizer on an averaging fibre contains a constraint
Jacobian and a gauge Jacobian. Both must be retained when a contour
gauge is replaced by a background gauge. We prove the finite-dimensional
identity with its compensating factor, then apply it to the actual
constraint and the complete constrained Hessian. The hypotheses needed
for this application are a differentiable physical background graph,
its local gauge covariance, and a nondegenerate physical Gaussian.
The identity does not construct that graph or an independent-current
extension of its Hessian.

\subsubsection{Constrained Lebesgue integrals and saddle determinants}

All spaces in this subsection are finite-dimensional real inner-product
spaces. Lebesgue measures and determinants are defined by their
orthonormal coordinates. Thus the formulas also apply to physical
lattice measures, after those measures have been put into the inner
products. Let \(Q:E\to F\) be onto, put \(n=\dim E\), \(m=\dim F\),
and let \(T=T^*\) be positive on \(\ker Q\). Define
\begin{equation}\label{pgg:constrained-integral}
 I(T,Q)=\int_E\delta_F(Qa)e^{-\langle a,Ta\rangle/2}\,da.
\end{equation}
The delta function is normalized by the Lebesgue measure of \(F\).
This convention is important when the constraint is changed.

\begin{lemma}\label{pgg:saddle-coarea}
With
\[
 \mathbb S(T,Q)=
       \begin{pmatrix}T&Q^*\\Q&0\end{pmatrix},
\]
one has
\begin{equation}\label{pgg:saddle-integral}
 I(T,Q)=(2\pi)^{(n-m)/2}
                         |\det\mathbb S(T,Q)|^{-1/2}.
\end{equation}
For every symmetric \(H:F\to F\), the saddle determinant and
the integral are unchanged by \(T\mapsto T+Q^*HQ\).
\end{lemma}

\begin{proof}
Split \(E=E_0\oplus E_1\) orthogonally, where \(E_0=\ker Q\),
and write \(Q=(0,Q_1)\). The square matrix \(Q_1:E_1\to F\)
is invertible and
\[
 \delta_F(Qa)\,da=|\det Q_1|^{-1}\,da_0\,\delta_0(da_1).
\]
Consequently the integral equals
\((2\pi)^{(n-m)/2}
(\det T_{00})^{-1/2}(\det QQ^*)^{-1/2}\).
Eliminating the \(T_{00}\) block in the saddle matrix gives
\[
 \det\mathbb S(T,Q)
       =(-1)^m\det T_{00}\det(QQ^*).
\]
The remaining two-by-two block has determinant
\((-1)^m|\det Q_1|^2\), independently of its upper-left
entry. This proves both assertions. In particular no positivity
of \(T\) on all of \(E\) was used.
\end{proof}

\subsubsection{The relative gauge determinant}

Let \(H\) be a fine site space and \(K\) a coarse site space.
Write \(D:H\to E\), \(\mathrm{d}:K\to F\) for the two infinitesimal
gauge actions. Assume initially that \(D\) and \(\mathrm{d}\) are
injective. Let \(R:H\to K\) and \(Q:E\to F\) be onto and satisfy
\begin{equation}\label{pgg:gauge-chain}
                              QD=\mathrm{d}R.
\end{equation}
Let \(B=B^*\) satisfy \(BD=0\). Assume that its restriction
to \(\ker Q\) is nonnegative and that its kernel there is
exactly \(D(\ker R)\). Put \(h=\dim H\), \(k=\dim K\), and
\(H_0=\ker R\). A residual gauge coordinate is a linear map
\(G:E\to Z\), where \(\dim Z=h-k\), such that
\[
                     GD_0:H_0\longrightarrow Z
               \quad\hbox{is invertible},\qquad D_0=D|_{H_0}.
\]
It follows that \(B+G^*G\) is positive on \(\ker Q\), so its
constrained integral is finite. Define
\begin{equation}\label{pgg:relative-invariant}
 \mathcal J(B,Q,R,G)
    =\sqrt{\det(RR^*)}\,|\det(GD_0)|
                           I(B+G^*G,Q).
\end{equation}
The determinant of \(GD_0\) uses the induced inner product
on \(H_0\); it is not a coordinate determinant in an
unnormalized list of site values.

\begin{proposition}[The gauge and constraint change identity]
\label{pgg:relative-gauge-identity}
The quantity \(\mathcal J\) is independent of \(G\).
If \(T:E\to K\) and
\begin{equation}\label{pgg:constraint-change}
       Q'=Q+\mathrm{d}T,\qquad R'=R+TD
\end{equation}
are onto, with the same nondegeneracy on the physical quotient,
then \(\mathcal J(B,Q',R',G')=\mathcal J(B,Q,R,G)\) for every
admitted \(G'\). The factor \(\sqrt{\det(RR^*)}\) in
\eqref{pgg:relative-invariant} is necessary with the delta
normalization \eqref{pgg:constrained-integral}.
\end{proposition}

\begin{proof}
Choose any full gauge coordinate \(C:E\to H\) for which
\(CD\) is invertible, and consider
\begin{equation}\label{pgg:full-gauge-integral}
 \mathcal Z(C,Q)=|\det(CD)|
 \int_{E\times K}\delta_F(Qa-\mathrm{d}\xi)
       e^{-\langle a,Ba\rangle/2-\|Ca\|^2/2}\,da\,d\xi.
\end{equation}
This is a finite Gaussian integral. Indeed, if \(Qa=\mathrm{d}\xi\),
choose \(v\in H\) with \(Rv=\xi\). Then \(a-Dv\in\ker Q\),
and the value of \(B\) is unchanged by this subtraction.
Its only null directions on that constrained space are gauge
directions. The full coordinate \(C\) controls them, while
injectivity of \(\mathrm{d}\) controls the remaining \(\xi\) variable.

To prove independence of \(C\), choose a fixed complement
\(E=E_{\rm tr}\oplus\operatorname{ran}D\) and write
\(a=b+Dv\). The Jacobian of this fixed decomposition is a
constant independent of \(C\) and \(Q\). In the other variable
replace \(\xi\) by \(\zeta+Rv\). Then
\[
 Qa-\mathrm{d}\xi=Qb-\mathrm{d}\zeta,\qquad
 \langle a,Ba\rangle=\langle b,Bb\rangle.
\]
The integral in \(v\) is
\[
 \int_H e^{-\|Cb+CDv\|^2/2}\,dv
                =(2\pi)^{h/2}|\det(CD)|^{-1}.
\]
It cancels the prefactor in \eqref{pgg:full-gauge-integral}
and is independent of \(b\). Thus \(\mathcal Z\) is independent
of \(C\).

It is also unchanged by \(Q\mapsto Q+\mathrm{d}T\). In
\eqref{pgg:full-gauge-integral}, make the triangular change
\(\xi\mapsto\xi+Ta\), holding \(a\) fixed. Its Jacobian is one,
and it removes \(\mathrm{d}Ta\) from the argument of the delta function.
The action and gauge weight are unchanged.

It remains to express \(\mathcal Z\) by a residual gauge
coordinate, retaining all measure factors. Choose \(c:F\to K\)
such that \(cd\) is invertible, and take
\[
                         Ca=(Ga,cQa)\in Z\oplus K.
\]
Its target has dimension \(h\). In the orthogonal splitting
of \(H\), use the right inverse
\(R^\dagger=R^*(RR^*)^{-1}\) and write
\[
                  v=v_0+R^\dagger\eta,\qquad v_0\in H_0.
\]
The Jacobian of \((v_0,\eta)\mapsto v\) is
\((\det RR^*)^{-1/2}\). The matrix of \(CD\) in these
coordinates is block triangular:
\[
 \begin{pmatrix}
      GD_0&GDR^\dagger\\0&cd
 \end{pmatrix}.
\]
Therefore
\begin{equation}\label{pgg:full-fp-factor}
 |\det(CD)|
   =\sqrt{\det(RR^*)}\,|\det(GD_0)|\,|\det(cd)|.
\end{equation}

For a fixed \(\xi\), the constrained integral with
\(Qa=\mathrm{d}\xi\), action \(B\), and gauge weight \(\|Ga\|^2/2\)
equals the one with \(Qa=0\). To see this choose
\[
 v_\xi=R^\dagger\xi
           -(GD_0)^{-1}GDR^\dagger\xi,
\]
where the second term is regarded as a vector in \(H_0\).
Then \(Rv_\xi=\xi\), \(GDv_\xi=0\), and translation
\(a\mapsto a+Dv_\xi\) preserves the action and residual
gauge weight. On the constraint, \(cQa=cd\xi\). Hence the
remaining integral in \(\xi\) in
\eqref{pgg:full-gauge-integral} is
\((2\pi)^{k/2}|\det(cd)|^{-1}\).
Combining it with \eqref{pgg:full-fp-factor} yields
\[
                    \mathcal Z(C,Q)
                        =(2\pi)^{k/2}\mathcal J(B,Q,R,G).
\]
The two independence assertions for \(\mathcal Z\) now give
the proposition.
\end{proof}

The constraint change in \eqref{pgg:constraint-change}
alters only the coarse gauge component of \(Qa\). Its effect
on the residual fine gauge group is \(R\mapsto R+TD\).
Omitting the resulting coarea factor would make the
invariance statement false, even though the quotient of
physical variations is the same.

\subsubsection{The background-gauge formula}

Let \(\Delta=D^*D\). Since \(D\) is injective on \(H_0\), its
compression \(\Delta_{00}\) to \(H_0\) is positive. Write
\[
 P_0=1-R^*(RR^*)^{-1}R
\]
for the orthogonal projection onto \(H_0\), and choose the
residual coordinate \(G=P_0D^*:E\to H_0\). Then
\[
 GD_0=\Delta_{00},\qquad G^*G=DP_0D^*.
\]
The same saddle calculation as in
\cref{pgg:saddle-coarea} gives
\[
        |\det\mathbb S(\Delta,R)|
                    =\det(RR^*)\det\Delta_{00}.
 \]
Consequently
\begin{equation}\label{pgg:background-gauge-product}
 \mathcal J
   =\frac{|\det\mathbb S(\Delta,R)|}{\sqrt{\det(RR^*)}}\,
                          I(B+DP_0D^*,Q).
\end{equation}
In particular, if \(RR^*=1\), its logarithm is
\begin{equation}\label{pgg:saddle-log-product}
 \log\mathcal J
   =-\tfrac12\log|\det\mathbb S(B+D(1-R^*R)D^*,Q)|
       +\log|\det\mathbb S(\Delta,R)|
       +\tfrac{n-m}{2}\log(2\pi).
\end{equation}
The site operator in the second determinant is the constrained
site Laplacian. The gauge form in the vector determinant is
\(D(1-R^*R)D^*\). It need not be the projected Landau gauge
form used to construct the physical background; the
gauge-change identity justifies this change in the integral.

\subsubsection{The actual constrained Hessian}

We now specialize the algebra. Let \(U(V)\) be an actual
differentiable critical background on a neighborhood of a
real coarse field \(V_0\), satisfying the prescribed homogeneous
constraint \(\mathcal M^jU(V)=V\). Assume its orbit is locally
unique, is gauge covariant, and has a nondegenerate Gaussian
on its physical fibre. Put \(U_0=U(V_0)\) and
\[
                         S_c(V)=A(U(V)).
\]
The action \(S_c\) is gauge invariant. All derivatives in what
follows use the actual logarithmic fine and coarse coordinates
at \(U_0,V_0\), with the physical measures declared for the
construction.

Let \(q_0(a)\) be the logarithm of the plain constraint:
\[
   \mathcal M^j(e^{i\eta a}U_0)=e^{iq_0(a)}V_0,
   \qquad q_0(0)=0,\quad Dq_0(0)=Q_0.
\]
At a constrained-critical point,
\begin{equation}\label{pgg:physical-multiplier}
     DA(U_0)=Q_0^*\mu,\qquad DS_c(V_0)=\mu.
\end{equation}
The second equality follows from the first by differentiating
the constraint along \(U(V)\). The derivative of that graph
is a right inverse of \(Q_0\), so no inverse of \(Q_0^*Q_0\)
is implicit here.

Let \(q(a)\) be the actual straightened constraint with its
actual node factor \(v(a)\), as in
\eqref{mb:node-diagram}. Thus the plain coarse field and
\(e^{iq(a)}V_0\) are related by that node gauge. Gauge
invariance gives the exact nonlinear identity
\begin{equation}\label{pgg:action-node-identity}
        S_c(\mathcal M^j(e^{i\eta a}U_0))
                          =S_c(e^{iq(a)}V_0).
\end{equation}
For the convention in which a site tangent acts by \(D\),
differentiation of the node factorization gives
\begin{equation}\label{pgg:actual-Q-shift}
 Q=Q_0+\mathrm{d}T,\qquad R=R_0+TD,\qquad QD=\mathrm{d}R.
\end{equation}
Here \(Q=Dq(0)\), \(R_0\) is evaluation of a fine gauge
parameter at the coarse nodes, and \(T\) is minus the
derivative of the node factor in that convention. Equivalently
the first identity defines its sign. The second identity
also follows by differentiating the same factorization along
a fine gauge transformation. The resulting \(R\) is the actual
transported site average, not an independently chosen map.
Its formula is the own-block average in the field
intertwining identity.

Consider the gauge-invariant fine function
\begin{equation}\label{pgg:subtracted-action}
  \widetilde A(W)=A(W)-S_c(\mathcal M^jW).
\end{equation}
Its first derivative vanishes at \(U_0\), by
\eqref{pgg:physical-multiplier}. Let \(B\) be its Hessian
there. Gauge invariance at a critical point gives \(BD=0\).
The chain rule in \eqref{pgg:action-node-identity} gives the
complete formula
\begin{equation}\label{pgg:actual-complete-Hessian}
 B=A''(U_0)-\langle\mu,D^2q(0)\rangle-Q^*H_cQ,
              \qquad H_c=D^2S_c(V_0).
\end{equation}
The middle term denotes the symmetric form
\[
 (a,b)\longmapsto
                  \langle\mu,D^2q(0)[a,b]\rangle.
\]
Thus, on the linearized fibre \(\ker Q\), the actual
quadratic action is
\begin{equation}\label{pgg:actual-current-form}
             B_{\rm fib}=A''(U_0)
                              -\langle\mu,D^2q(0)\rangle .
\end{equation}
It contains the curvature of the nonlinear constraint.
The last term in \eqref{pgg:actual-complete-Hessian}
does not change its saddle determinant, by
\cref{pgg:saddle-coarea}. This derivation defines the physical
current and multiplier correction from the actual constraint;
it is not a form selected to obtain a determinant identity.

The same construction with \(q_0\) gives the plain-fibre
Hessian. The full \(B\) is the same for both constraints,
because it is the Hessian of the single function
\eqref{pgg:subtracted-action}. This is the reason the
constraint-change part of
\cref{pgg:relative-gauge-identity} applies.

\subsubsection{The two actual site Jacobians}

In fine physical measure \(\eta^4\), node evaluation has
\begin{equation}\label{pgg:node-coarea}
                     R_0R_0^*=\eta^{-4}I.
\end{equation}
This is a field-independent factor. By contrast the
transported site average is a coisometry:
\begin{equation}\label{pgg:site-coisometry}
                             RR^*=I.
\end{equation}
To verify the latter at one step, write its value as the
normalized sum of the transported fine site values in a
block. Each real adjoint transport is orthogonal, distinct
blocks have disjoint site sets, and the fine measure divided
by the coarse measure is precisely the reciprocal of the
number of sites in a block. Multiplying the average by its
physical adjoint therefore gives the identity. For a
homogeneous composition, write the site average as the
composition of its one-step averages. Each is a coisometry
between its own physical site spaces, so their composition
is a coisometry as well. This calculation does not require
flatness.
Here the straightened constraint is re-centered at the actual
\(U_0\), so its relative contour logarithms vanish at \(a=0\).
Equation \eqref{pgg:site-coisometry} concerns that derivative.
It does not identify the effective gauge derivative of a
fixed-flat-reference nonlinear chart at a nonzero field
with a coisometry.

The actual contour gauge on the plain fibre supplies a
residual coordinate \(G_{\rm ax}\). Its derivative along a
gauge parameter which vanishes at the coarse nodes is
triangular in the levels, with field-independent determinant.
Here is the local verification. A contour from a block node
\(y\) to \(x\) transforms by its two endpoint gauges. At the
axial representative its value is one, so its logarithmic
derivative along \(v_t=e^{it\lambda}\) is
\(\lambda(y)-\lambda(x)\), with the fixed sign convention.
Group the noncoarse sites by their first parent in the
nested block hierarchy. The map from site values to these
parent differences is triangular with diagonal entries
\(\pm I_3\), since the final coarse-node values are zero.
Its determinant has absolute value one in these unweighted
site coordinates. Transported frames insert adjoint
rotations of determinant one. Physical measures and fixed
gauge weights insert only constants depending on the mesh
and the numbers of sites. Therefore
\begin{equation}\label{pgg:axial-fp-constant}
 \sqrt{\det(R_0R_0^*)}
          |\det(G_{\rm ax}D|_{\ker R_0})|=c_{\eta,j},
\end{equation}
where \(c_{\eta,j}>0\) is independent of the background.
For one step the endpoint derivative is exactly
\(GT=-I\) from \eqref{agm:actual-gauge-splitting};
the preceding triangular calculation is its composition
over the hierarchy.

Let \(\Psi_j(V)\) be the logarithm of the actual physical
Gaussian normalizer on the plain fibre, minus its value at
the trivial background. Its Haar factors at the centre are
field independent in left exponential coordinates. The
nonpivot elimination factor is precisely the constraint
coarea factor in \eqref{pgg:constrained-integral}. The
reference contour-gauge normalizations and the constant
\eqref{pgg:axial-fp-constant} also disappear from this
normalized difference.

\begin{proposition}[The physical saddle representation]
\label{pgg:physical-saddle}
On every admitted actual background domain described above,
\begin{equation}\label{pgg:physical-RF12}
 \begin{split}
 \Psi_j(V)={}&-\tfrac12\log\left|
   \det\mathbb S\left(
    A''(U_0)-\langle\mu,D^2q(0)\rangle
                       +D(1-R^*R)D^*,\,Q\right)\right|\\
 &+\log|\det\mathbb S(D^*D,R)|
                          -(\text{the same at }V=1).
 \end{split}
\end{equation}
The multiplier and the constraint curvature use their
physical coordinate normalizations. The formula is a
physical-pair identity. It does not specify how either
determinant is extended off that manifold to an independent
current.
\end{proposition}

\begin{proof}
At a coarse field with no infinitesimal stabilizer, \(\mathrm{d}\)
is injective. The fine gradient is injective as well:
\(D\lambda=0\) implies \(\mathrm{d}R_0\lambda=0\), hence
\(R_0\lambda=0\); a parallel fine field which is zero at a
coarse node is zero on its component. The same observation
makes \(D\) injective on \(\ker R\) without the generic
coarse-field restriction.

The physical Gaussian is nondegenerate on the plain fibre
modulo its residual gauges. Because the two constraints
differ by a coarse gauge component, their physical
quotients are isomorphic: given a vector on one constraint
space, correct its coarse gauge value by \(D\lambda\),
using surjectivity of the other site map. The value of
\(B\) is unchanged. Thus the required quotient
nondegeneracy holds for both \(Q_0,R_0\) and \(Q,R\).
The same argument shows \(Q\) is onto: if \(Q_0a=f\),
then \(Qa=f+\mathrm{d}Ta\); choose \(\lambda\) with \(R\lambda=-Ta\)
and replace \(a\) by \(a+D\lambda\).

Apply \cref{pgg:relative-gauge-identity} to the common
gauge-invariant Hessian \(B\). On the plain side use
\(G_{\rm ax}\) and \eqref{pgg:axial-fp-constant}.
On the straightened side use the background gauge and
\eqref{pgg:site-coisometry}. Formula
\eqref{pgg:saddle-log-product} gives the two logarithms.
Replace \(B\) by \(B_{\rm fib}\) using
\eqref{pgg:actual-complete-Hessian} and
\cref{pgg:saddle-coarea}. All remaining constants cancel
in the difference at \(V=1\), giving
\eqref{pgg:physical-RF12}.

The auxiliary integral \eqref{pgg:full-gauge-integral}
is singular when \(\mathrm{d}\) has a kernel, so it is not used
at such a field. After its coarse-gauge factors have
been canceled, both sides of
\eqref{pgg:physical-RF12} are finite constrained
Gaussian or saddle-determinant expressions. Their
nondegenerate fibre forms vary continuously on the
admitted background domain. Irreducible coarse
connections are dense near the trivial field: two small
loop holonomies with distinct colour axes have no
common nonzero adjoint fixed vector. Perturbations of
links can create such loops on the connected periodic
four-dimensional lattice. Thus the finite final
identity extends to the trivial field and to fields
with stabilizers by continuity. No continuity of the
uncanceled auxiliary integral is asserted.
\end{proof}

On an admitted complex background domain with coherent
logarithm branches, use
bilinear physical transposes and the logarithm branches
continued from the real positive Gaussians. The determinant
identity then continues holomorphically wherever the
constrained forms remain nondegenerate. Absolute values
belong only to the displayed real formula. In particular
one cannot insert them in a holomorphic localization.

\subsubsection{Gaussian composition with nonlinear constraints}

We also record why Gaussian integration in stages uses the
complete constrained Hessians. This identifies the algebraic
composition needed before a one-loop coefficient is extracted.
Let \(M:E\to F\) and \(N:F\to G\) be onto. Let \(T=T^*\)
be positive on \(\ker(NM)\), and let \(H:F\to E\) be its
constrained minimizer:
\[
 MH=1,\qquad
 \langle u,THb\rangle=0\quad(u\in\ker M).
\]
Put \(T_2=H^*TH\). The change \(a=u+Hb\), \(u\in\ker M\),
has Jacobian \((\det MM^*)^{-1/2}\); this follows first
for the orthogonal right inverse of \(M\), and then for
\(H\), whose difference from it is a shear into \(\ker M\).
The constraint becomes \(Nb=0\), and the exponent splits.
Hence exactly
\begin{equation}\label{pgg:gaussian-gluing}
                         I(T,NM)=I(T,M)I(T_2,N).
\end{equation}
Only positivity on the indicated constraint spaces is
needed. There is no additional determinant of \(N\) or
\(M\): both are already included by the two delta functions
with their declared measures.

For the nonlinear application let \(m_1,m_2\) be the
actual successive constraint maps in their local
coordinates. Let \(x=x(y)\) be the first physical
critical graph and \(y\) a critical point of
\(S_1(y)=F(x(y))\) under \(m_2(y)=z\). Write
\[
 M=Dm_1(x),\qquad N=Dm_2(y),\qquad
 DF(x)=M^*\mu_1,\qquad DS_1(y)=N^*\mu_2.
\]
Differentiation of the first graph gives
\(\mu_1=N^*\mu_2\). The complete first-stage Hessian is
\[
 B_1=F''(x)-\langle\mu_1,m_1''(x)\rangle.
\]
With \(H=Dx(y)\), differentiation of the criticality
equation gives
\begin{equation}\label{pgg:nonlinear-schur}
 MH=1,\qquad
 \langle u,B_1Hb\rangle=0\quad(u\in\ker M),\qquad
 S_1''(y)=H^*B_1H.
\end{equation}
For the last equality, differentiate \(m_1(x(y))=y\)
twice and pair with \(\mu_1\). Its acceleration term
cancels the corresponding acceleration term in
\(D^2(F\circ x)\). This calculation uses the
multiplier term in \(B_1\).

The total constraint has second derivative
\[
 D^2(m_2\circ m_1)[a,a]
     =m_2''[Ma,Ma]+Nm_1''[a,a].
\]
Therefore its complete Hessian is
\begin{equation}\label{pgg:composite-complete-Hessian}
 B_{\rm tot}=B_1-M^*\langle\mu_2,m_2''\rangle M.
\end{equation}
In the coordinates \(a=u+Hb\), this is the direct sum of
\(B_1|_{\ker M}\) and
\(S_1''-\langle\mu_2,m_2''\rangle\).

Gauge fixing does not spoil this decomposition. Choose
the actual first-stage axial representative throughout
the first background graph. Its first gauge coordinate
vanishes along that graph, so its derivative \(G_1\)
satisfies \(G_1H=0\). The second-stage gauge coordinate
acts on \(b=Ma\). The total quadratic gauge weight is
then \(\|G_1u\|^2+\|G_2b\|^2\), precisely the two
gauge weights used by the hierarchical Gaussian
integral. Apply \eqref{pgg:gaussian-gluing} to these
nondegenerate gauge-fixed forms, explicitly taking
\[
 T=B_{\rm tot}+G_1^*G_1+M^*G_2^*G_2M.
\]

\begin{corollary}\label{pgg:actual-stage-composition}
On an admitted homogeneous physical background graph,
the normalized Gaussian logarithms obey
\begin{equation}\label{pgg:actual-Psi-recursion}
             \Psi_j=\Phi_j+\Psi_{j-1}\circ p_j,
\end{equation}
where \(p_j\) is the actual preceding critical background
map and \(\Phi_j\) is the actual one-step Gaussian
logarithm. All constraint-curvature terms are included
before the Gaussian integration. Field-independent
measure constants cancel in the normalization.
\end{corollary}

\begin{proof}
Use the preceding Hessian decomposition at each stage
and integrate first in \(u\), then in \(b\).
The first factor in \eqref{pgg:gaussian-gluing} is the
earlier Gaussian normalizer evaluated at \(p_j(V)\);
the second is the new step normalizer. Normalize both
at the trivial coarse field and take their real
logarithms. This gives
\eqref{pgg:actual-Psi-recursion}.
\end{proof}

\subsubsection{The physical multiplier correction has zero coefficient}

The statement that a current correction is of third order in
external momentum concerns its determinant functional. The
matrix \(-\langle\mu,q''\rangle\) itself need not have that
order. We prove the required functional statement on the
actual background domain.

Fix a homogeneous depth \(j\), and work near a flat reference
with trivial coarse field for the coefficient extraction.
For the following argument assume the actual critical
background graph and its first two derivatives are
exponentially local on this neighborhood. The weighted
construction in \cref{prop:hb-actual-critical-orbit}, with
\eqref{hb:localized-background-response} and
\eqref{hb:localized-background-jets},
supplies this property on the constructed flat-reference
neighborhood; no estimate uniform in \(j\)
for the determinants in this argument will be required.
Let
\[
 T(V)=A''(U(V))+D(1-R^*R)D^*.
\]
Introduce an independent coarse bond covector \(\nu\),
with the same coordinate normalization as \(\mu\), and put
\begin{equation}\label{pgg:multiplier-log-difference}
 \begin{split}
 L_j(V,\nu)={}&-\tfrac12
       \log\det\mathbb S\bigl(T(V)-\langle\nu,q''_V\rangle,Q_V\bigr)\\
       &+\tfrac12\log\det\mathbb S(T(V),Q_V).
 \end{split}
\end{equation}
Here and below the logarithms are coherent local branches;
their difference is zero at \(\nu=0\). The actual
multiplier correction to the physical saddle logarithm is
\[
                        C_j(V)=L_j(V,DS_c(V)).
\]

\begin{proposition}\label{pgg:current-zero-coefficient}
The bulk Hessian of \(C_j\) has symbol
\begin{equation}\label{pgg:current-cubic-momentum}
                        \widehat\pi^{C_j}(p)=O_j(|p|^3)
                       \qquad(p\to0).
\end{equation}
Consequently \(b[C_j]=0\). The statement is exact at every
admitted depth. It does not assert a pointwise
\(O(|p|^3)\) bound for the multiplier matrix.
\end{proposition}

\begin{proof}
We first justify the local trace jets used here. At a fixed
fine mesh, the flat background-gauge vector form restricted
to \(\ker Q\) is nondegenerate. A zero of its sum of
nonnegative squares would satisfy
\[
             \mathrm{d}v=0,\qquad Qv=0,\qquad P_0D^*v=0,
                        \quad P_0=1-q^*q.
\]
In each Bloch fibre a closed field is a gradient, apart
from the parallel constant modes at zero frequency. The
actual average detects those constant vector modes, so
they vanish. For the gradient part \(v=D\lambda\),
the identity \(QD=\mathrm{d}_cq\) gives \(\mathrm{d}_cq\lambda=0\).
Subtract a parallel constant from \(\lambda\) if needed;
it then has \(q\lambda=0\). Pairing
\(P_0D^*D\lambda=0\) with \(\lambda\) gives \(D\lambda=0\).
Thus \(v=0\).

The finite Bloch matrices and their constraint maps are
continuous on the compact Brillouin torus. Surjectivity
of \(Q\), supplied by the actual row seeds, and the
preceding kernel argument imply invertibility of the
full saddle matrix in every fibre. At this fixed mesh,
its inverse is uniformly bounded over that torus.
It extends to a small complex momentum strip by the
resolvent identity. Fourier contour shifting therefore
gives exponential decay of its coarse-cell kernels,
with constants which may depend on \(j\).

All other coefficients in
\eqref{pgg:multiplier-log-difference} have fixed physical
carriers before the background is substituted. The
exponentially local background derivatives and the
resolvent identity show that its two trace derivatives
have absolutely summable exponential moments. For
example, differentiating a logarithm twice gives
\(\operatorname{Tr}(S^{-1}S''-
S^{-1}S'S^{-1}S')\); each differentiated coefficient
has its declared local carrier or the exponentially
decreasing background response. Products of these
kernels retain an exponential moment after decreasing
its rate. There are finitely many fine coordinates per
cell at fixed \(j\), so their internal trace is finite.
The same estimates make finite-volume winding terms
vanish exponentially. This proves existence of the
bulk Hessian symbols and their analyticity in a
momentum neighborhood of zero. It also justifies the
Ward identities below by first using finitely supported
gauge parameters and then taking these absolutely
convergent Fourier sums.

The function \(L_j\) is invariant under simultaneous
coarse gauge transformation of \(V\) and the covector
\(\nu\). Indeed, the background orbit is gauge
covariant, all fine and coarse tangent spaces transform
by their adjoint rotations, and the actual constraint
second derivative transforms with its covector
pairing. Both saddle determinants are conjugated
by these rotations, whose determinants are one.
In addition,
\begin{equation}\label{pgg:joint-current-first-jets}
 L_j(V,0)=0,\qquad DL_j(1,0)=0.
\end{equation}
The first equality is its definition. For the derivative
in \(\nu\), invariance under constant colour rotations
makes each coefficient an invariant linear functional
on the adjoint colour space, hence zero. The derivative
in \(V\) vanishes by the first equality.

Let \(K_j\) be the mixed \(V,\nu\) Hessian kernel at
\((1,0)\). Differentiate the joint gauge identity in a
gauge parameter and in \(\nu\). The term from
transforming \(\nu\) is multiplied by the zero first
derivative in \eqref{pgg:joint-current-first-jets}.
Thus
\[
              D^2_{V,\nu}L_j(1,0)[\mathrm{d}_c\chi,\nu]=0.
\]
Translation covariance gives, in Fourier coordinates,
\begin{equation}\label{pgg:mixed-current-Ward}
                 \widehat{\mathrm{d}}_c(-p)^{\mathsf T}\widehat K_j(p)=0.
\end{equation}
Since \(\widehat{\mathrm{d}}_{c,\mu}(p)=-ip_\mu+O(|p|^2)\),
its linear Taylor term forces
\(\widehat K_j(0)=0\): every column is annihilated by
every row vector \(p^{\mathsf T}\). Analyticity then gives
\begin{equation}\label{pgg:mixed-current-order}
                             \widehat K_j(p)=O_j(|p|).
\end{equation}

The actual multiplier is \(\mu(V)=DS_c(V)\), and
\(\mu(1)=0\). Let \(P_j\) be the symbol of its derivative,
that is the Hessian of \(S_c\) at the trivial field.
Gauge invariance gives \(P_j(0)=0\) by the same
zero-momentum Ward argument. In midpoint coordinates,
inversion sends both bond variations to their negatives
and \(p\) to \(-p\). The Hessian symbol in those
coordinates is therefore even. Its linear Taylor term
vanishes, and the endpoint-to-midpoint phase does not
change this conclusion after \(P_j(0)=0\). Hence
\begin{equation}\label{pgg:physical-multiplier-order}
                              P_j(p)=O_j(|p|^2).
\end{equation}
These are the actual symmetries of the symmetric
contour construction and its critical orbit.

Apply the second-order chain rule to \(C_j\).
The pure \(V,V\) Hessian of \(L_j(V,0)\) is zero.
The term involving \(D^2\mu\) is multiplied by
\(D_\nu L_j(1,0)=0\). The remaining terms are the
two mixed kernels composed with \(P_j\), and the
\(\nu,\nu\) kernel composed with \(P_j\) on both
sides. Their orders are respectively at least
\(|p|^3\) and \(|p|^4\), by
\eqref{pgg:mixed-current-order} and
\eqref{pgg:physical-multiplier-order}.
This proves \eqref{pgg:current-cubic-momentum}.
The quadratic Maxwell coefficient is therefore zero.
\end{proof}

The coefficient identity
\(\sum_{l\le j}b_l^0=b[\Psi_j]\) additionally uses
invariance of the coefficient under the actual map
\(p_j\). The saddle formula and the recursion alone
do not supply that response estimate. Likewise,
identification of the physical Hessian here does
not give localization or common complex-current
domains for a chosen off-shell completion. These
analytic requirements remain distinct from the
finite gauge-change and composition identities.

% END LAYERED EXPANSION physical_gaussian_gauge_change.tex

% BEGIN LAYERED EXPANSION massive_lattice_coefficient.tex
\subsection{The massive Wilson determinant in a constant field}
\label{mlc:section}

We compute the logarithmic coefficient directly on the unit lattice.
This is a component of the small-field Gaussian renormalization
method of \cite{B87}.
This calculation concerns the Wilson Hessian and the covariant site
Laplacian, before an averaging constraint is introduced. It includes
the neutral colour, the charged colours, and the order of the volume,
field, and mass limits. The comparison with a constrained determinant
is a separate step.

Write \(t_a\) for the Pauli matrices, so that
\(\tfrac12\operatorname{tr}(t_at_b)=\delta_{ab}\). A real bond
variation is written \(e^{it a_b}U_b\), with
\(a_b=\sum_a a_b^a t_a\). The Wilson action and its Hessian have the
normalization
\[
 A(U)=\sum_p\bigl(1-\tfrac12\operatorname{tr}U_p\bigr),\qquad
 A''(1)[a,a]=\sum_p|(\mathrm{d}a)_p|^2.
\]
Adjoints below use counting measure on the unit lattice and this
Euclidean colour inner product. Put
\[
 (D_U\lambda)_\mu(x)
   =\lambda(x)-\operatorname{Ad}_{U_\mu(x)}\lambda(x+e_\mu),\qquad
 \Delta_U=D_U^*D_U,\qquad K_U=A''(U)+D_UD_U^*.
\]
The opposite sign for both occurrences of \(D_U\) gives the same
operators. We use a star for the real Hilbert adjoint and for its
complexification in this calculation.

For \(f\in\mathbb R\), fix the Abelian field
\begin{equation}
 U_{f,1}(x)=e^{-ifx_2t_3},\qquad U_{f,\mu}(x)=1\quad(\mu\ne1).
 \label{mlc:background}
\end{equation}
Its positively oriented \((1,2)\)-plaquettes are \(e^{ift_3}\);
all its other plaquettes are the identity. The connection itself
grows in a logarithmic coordinate as \(|x_2|\), whereas every
plaquette is constant. We therefore define the determinant density
by finite compressions, rather than by placing arbitrary flux on
a periodic torus.

Let \(\Lambda_R=[-R,R]^4\cap\mathbb Z^4\). Site operators are
compressed to sites in \(\Lambda_R\); vector operators are
compressed to positive bonds with both endpoints there.
Thus, for example, \(K_{f,R}=P_RK_{U_f}P_R^*\). Compression is
taken after the full lattice operator has been formed. This is
the Hessian for variations extended by zero outside the indicated
coordinates. It does not discard plaquettes meeting the boundary.
For \(z>0\), define
\begin{equation}
 \begin{split}
 \mathcal L_R(f,z)&=\log\det(\Delta_{f,R}+z)
                       -\tfrac12\log\det(K_{f,R}+z),\\
 \mathcal D_R(f,z)&=|\Lambda_R|^{-1}
                          \bigl(\mathcal L_R(f,z)-\mathcal L_R(0,z)\bigr).
 \end{split}
 \label{mlc:finite-density}
\end{equation}
The determinants here are over all three real colours. For each
fixed positive \(z\), the matrices in this formula are positive
when \(|f|\) is sufficiently small in terms of \(z\).

\begin{proposition}\label{mlc:main}
The functions \(\mathcal D_R(\cdot,z)\) and their first three
derivatives at zero have limits as \(R\to\infty\), locally uniformly
for \(z\) in compact subsets of \((0,\infty)\). The limiting density
\(\mathcal D(f,z)\) is a real, even, smooth function of real \(f\)
near zero. Let
\[
 \widetilde\beta(z)=\partial_f^2\mathcal D(0,z).
\]
Then
\begin{equation}
 \widetilde\beta(z)=\frac{11}{6\pi^2}\log\frac1z+C_*+E(z),\qquad
 |E(z)|+|zE'(z)|+|z^2E''(z)|\le C\sqrt z,
 \quad 0<z\le1,
 \label{mlc:asymptotic}
\end{equation}
for finite constants \(C_*\) and \(C\). In particular,
\begin{equation}
 \sup_{0<z\le1}\left\{
 \left|\widetilde\beta(z)-\frac{11}{6\pi^2}\log\frac1z\right|
 +|z\widetilde\beta'(z)|+|z^2\widetilde\beta''(z)|\right\}<\infty.
 \label{mlc:uniform-mass-bound}
\end{equation}
The same limits hold for any lattice F\o lner exhaustion, using
the same zero-extension compression convention.
\end{proposition}

The proof will give an exact Brillouin-zone formula at every positive
mass. The coefficient is therefore obtained from the actual lattice
Hessian, without replacing it by a continuum vector operator.

\subsubsection{Charged coordinates and magnetic translations}

The neutral coordinate is \(a^3\). For the other two real
coordinates put \(w=a^1-ia^2\). Direct multiplication of Pauli
matrices gives
\[
 \operatorname{Ad}_{e^{i\theta t_3}}w=e^{2i\theta}w.
\]
Consequently the charged covariant shifts in the background
\eqref{mlc:background} are
\[
 (u_1w)(x)=e^{-i\varphi x_2}w(x+e_1),\qquad
 (u_\mu w)(x)=w(x+e_\mu)\quad(\mu\ne1),\qquad \varphi=2f.
\]
They are unitary, and
\begin{equation}
 u_1u_2=e^{i\varphi}u_2u_1,
 \qquad u_\mu u_\nu=u_\nu u_\mu
 \quad\text{if }\{\mu,\nu\}\ne\{1,2\}.
 \label{mlc:commutation}
\end{equation}
For \(m\in\mathbb Z^4\) define
\[
 W_m=e^{-i\varphi m_1m_2/2}
              u_1^{m_1}u_2^{m_2}u_3^{m_3}u_4^{m_4}.
\]
The commutation relation proves, including negative powers,
\begin{equation}
 W_mW_n=e^{i\varphi(m_1n_2-m_2n_1)/2}W_{m+n},
 \qquad W_m^*=W_{-m}.
 \label{mlc:weyl-product}
\end{equation}
If \(a(q)=\sum_m a_me^{im\cdot q}\), write
\(\operatorname{Op}_\varphi(a)=\sum_m a_mW_m\). This definition
also applies to matrix coefficients, with the matrix acting on
the vector index. Absolute Fourier convergence is sufficient.

Let \(\mathbb B=[-\pi,\pi]^4\), with opposite faces identified,
and write
\[
 \int_{\mathbb B}h(q)\,d\bar q
       =(2\pi)^{-4}\int_{[-\pi,\pi]^4}h(q)\,dq.
\]
The normalized trace of a finite Weyl polynomial is
\[
 \tau\bigl(\operatorname{Op}_\varphi(a)\bigr)
       =\operatorname{tr}a_0
       =\int_{\mathbb B}\operatorname{tr}a(q)\,d\bar q.
\]
Indeed, \(W_m\) has a zero diagonal unless \(m=0\), and
\(W_0=1\). The same formula holds for absolutely summable
coefficients by approximation. For two such symbols the zero
Fourier coefficient of their product involves only \(n=-m\).
The phase in \eqref{mlc:weyl-product} is then one. Hence
\begin{equation}
 \tau\bigl(\operatorname{Op}_\varphi(a)
                  \operatorname{Op}_\varphi(b)\bigr)
   =\int_{\mathbb B}\operatorname{tr}(ab)\,d\bar q.
 \label{mlc:trace-product}
\end{equation}
In particular this trace is cyclic, also for matrix-valued
symbols. It is a trace per lattice site, with the ordinary matrix
trace retained on the fibre.

For later remainder estimates, use
\[
 \|a\|_s=\sum_{m\in\mathbb Z^4}(1+|m|)^s\|a_m\|,
 \qquad s=0,1,\ldots.
\]
For real \(\varphi\), the phase in
\eqref{mlc:weyl-product} has modulus one, so
\(\|a\star_\varphi b\|_0\le\|a\|_0\|b\|_0\).
The same argument and
\((1+|m+n|)^s\le(1+|m|)^s(1+|n|)^s\) give the corresponding
weighted product bound. The operator norm is at most \(\|a\|_0\).
These estimates are uniform in the real magnetic field. They do
not assert holomorphy in a complex constant flux.

Taylor expansion of the phase in
\eqref{mlc:weyl-product} gives
\begin{equation}
 a\star_\varphi b
 =ab-\frac{i\varphi}{2}\{a,b\}
          -\frac{\varphi^2}{8}\{\!\{a,b\}\!\}
          +\mathcal R_3(a,b;\varphi),
 \label{mlc:moyal}
\end{equation}
where subscripts denote ordinary momentum derivatives and
\[
 \{a,b\}=a_1b_2-a_2b_1,\qquad
 \{\!\{a,b\}\!\}=a_{11}b_{22}-2a_{12}b_{12}+a_{22}b_{11}.
\]
For example,
\[
 \|\mathcal R_3(a,b;\varphi)\|_0
       \le C|\varphi|^3\|a\|_3\|b\|_3.
\]
To check the signs, the product of \(e^{im\cdot q}\) and
\(e^{in\cdot q}\) acquires the positive phase in
\eqref{mlc:weyl-product}, while each pair of momentum derivatives
contributes \(i^2=-1\). The remainder bound follows from the
third-order Taylor remainder of a real phase and
\(|m_1n_2-m_2n_1|\le |m||n|\). Weighted versions follow by the
same summation.

\subsubsection{The exact charged Wilson Hessian}

We first derive the finite Weyl polynomial that is needed in the
calculation. This also fixes the spin normalization. Let
\(p=(x;\mu,\nu)\) be a positive plaquette and let its background
holonomy be \(P=e^{if_{\mu\nu}t_3}\), where
\(f_{12}=f\) and the other \(f_{\mu\nu}\) vanish. Put
\(\lambda=e^{2if_{\mu\nu}}\).
Moving the four varied factors to the base point gives
\[
 P(t)=e^{itX_1}e^{itX_2}e^{itX_3}e^{itX_4}P,
\]
where
\[
 \begin{split}
 X_1&=a_\mu(x),\\
 X_2&=\operatorname{Ad}_{U_\mu(x)}a_\nu(x+e_\mu),\\
 X_3&=-\operatorname{Ad}_{P U_\nu(x)}a_\mu(x+e_\nu),\\
 X_4&=-\operatorname{Ad}_{P}a_\nu(x).
 \end{split}
\]
The inverses on the third and fourth edges are responsible for
both minus signs and both occurrences of \(P\). In charged
coordinates the four rows acting on \((w_1,\ldots,w_4)\) are
\begin{equation}
 C_1w=w_\mu,\qquad C_2w=u_\mu w_\nu,\qquad
 C_3w=-\lambda u_\nu w_\mu,\qquad C_4w=-\lambda w_\nu.
 \label{mlc:four-insertions}
\end{equation}
These are rows of operators; their adjoints reverse the order of
operator multiplication.

For real three-vectors \(x,y\), the Pauli product is
\( (x\cdot t)(y\cdot t)=(x\cdot y)1+i(x\times y)\cdot t\).
Differentiating \(1-\tfrac12\operatorname{tr}P(t)\) twice gives
\[
 \cos f_{\mu\nu}\left|\sum_r X_r\right|^2
           -2\sin f_{\mu\nu}\sum_{r<s}(X_r\times X_s)_3.
\]
For charged vectors,
\((x\times y)_3=-\operatorname{Im}(\overline w_xw_y)\).
Thus the charged plaquette Hessian is the Hermitian form
\begin{equation}
 H_{\mu\nu}
 =\cos f_{\mu\nu}\,S_{\mu\nu}^*S_{\mu\nu}
   +\frac{\sin f_{\mu\nu}}{i}
       \sum_{r<s}(C_r^*C_s-C_s^*C_r),\qquad
 S_{\mu\nu}=\sum_{r=1}^4 C_r.
 \label{mlc:exact-plaquette-hessian}
\end{equation}
For example, the second summand evaluated on \(w\) is
\(2\sin f_{\mu\nu}\sum_{r<s}
\operatorname{Im}\langle C_rw,C_sw\rangle\), as required.
There is no continuum approximation in
\eqref{mlc:exact-plaquette-hessian}.

Put \(D_\mu=1-u_\mu\). The exact charged vector operator is
\begin{equation}
 K_\varphi=DD^*+\sum_{\mu<\nu}H_{\mu\nu}
             =\operatorname{Op}_\varphi(\sigma_\varphi),
 \qquad
 \Delta_\varphi=\sum_\mu D_\mu^*D_\mu
             =\operatorname{Op}_\varphi(\psi),
 \label{mlc:exact-vector}
\end{equation}
where
\[
 \psi(q)=\sum_{\mu=1}^4(2-2\cos q_\mu).
\]
Each entry of \(\sigma_\varphi\) is a finite sum of terms
\(c e^{i\varphi p/2}e^{im\cdot q}\), with integer \(p\).
At zero flux the ordinary curl--divergence identity gives
\begin{equation}
 \sigma_0(q)=\psi(q)I_4.
 \label{mlc:flat-symbol}
\end{equation}
This identity can also be checked entry by entry: the diagonal
curl terms supply all \(D_\nu^*D_\nu\) with \(\nu\ne\mu\),
the divergence term supplies the missing direction, and their
off-diagonal entries cancel when the shifts commute.

We record the short computation that determines the first flux
coefficient. Only the \((1,2)\) plane can change the flat
cancellation. To evaluate the Weyl symbol at \(q=0\), one must
first put each product into the order
\eqref{mlc:weyl-product}; setting the shifts to one beforehand
would lose the magnetic phase. For example the full
off-diagonal entry is
\[
 \begin{split}
 \sigma_\varphi^{12}(q)
 ={}&1-e^{-iq_2}-e^{iq_1}-e^{i(q_1-q_2)}\\
   &+e^{-i\varphi/2}
       \bigl(e^{i(q_1-q_2)}+e^{iq_1}+e^{-iq_2}\bigr)
       -e^{i\varphi/2}.
 \end{split}
\]
The phases in the second line arise from putting products of
the two noncommuting shifts into Weyl order. Its combined
coefficient at \(q=0\) is
\[
 \begin{array}{c|rrr}
 \text{half-phase }p&-1&0&1\\ \hline
 \sum_m c_{m,p}^{12}&3&-2&-1 .
 \end{array}
\]
Here the contribution of \(DD^*\) is combined with the two
terms in \eqref{mlc:exact-plaquette-hessian}. These three
numbers can be obtained by substituting the four rows
\eqref{mlc:four-insertions}, using
\(\lambda=e^{i\varphi}\), and applying
\(W_mW_n=e^{i\varphi(m_1n_2-m_2n_1)/2}W_{m+n}\) to each
pair. The diagonal sums vanish, and the opposite off-diagonal
entry is the complex conjugate. Hence
\begin{equation}
 \sigma_\varphi(0)
 =\bigl(3e^{-i\varphi/2}-e^{i\varphi/2}-2\bigr)E_{12}
  +\bigl(3e^{i\varphi/2}-e^{-i\varphi/2}-2\bigr)E_{21}.
 \label{mlc:zero-momentum-symbol}
\end{equation}
The entries with indices outside \(\{1,2\}\) are zero in this
formula. Write
\begin{equation}
 \sigma_\varphi=\psi I_4+\varphi\Sigma_1
                       +\varphi^2\Sigma_2+O(\varphi^3).
 \label{mlc:symbol-jets}
\end{equation}
The remainder and all its momentum derivatives are bounded by a
constant times \(|\varphi|^3\). In particular,
\begin{equation}
 \Sigma_1(0)=-2i(E_{12}-E_{21}),\qquad
 \operatorname{tr}\Sigma_1(0)^2=8,
 \qquad \Sigma_2(0)=-\tfrac14(E_{12}+E_{21}).
 \label{mlc:spin-value}
\end{equation}
The finite Fourier expressions also give
\begin{equation}
 \|\Sigma_1(q)-\Sigma_1(0)\|\le C|q|,
 \qquad \|\Sigma_2(q)\|\le C,\qquad q\in\mathbb B.
 \label{mlc:jet-bounds}
\end{equation}
Representatives in \([-\pi,\pi]^4\) are used in the first bound.

There is a useful independent identity for this normalization.
The field \eqref{mlc:background} is critical under compactly
supported variations: on each bond the two adjacent
\((1,2)\)-plaquettes contribute opposite copies of
\(\sin f\,t_3\), and the remaining plaquettes contribute zero.
Gauge invariance, differentiated at this critical field, gives
\(A''(U_f)D_{U_f}=0\). Consequently
\begin{equation}
 K_\varphi D=D\Delta_\varphi.
 \label{mlc:critical-ward}
\end{equation}
It also follows by multiplication of the finite Weyl polynomials
above. The spin value in \eqref{mlc:spin-value} was computed
from the Hessian itself; it was not inferred from the Ward
identity alone.

\subsubsection{The neutral sector and positivity at fixed mass}

In the third colour all adjoint phases disappear. The four
insertions commute, so the only changed plaquette Hessian is
the factor \(\cos f\) in the \((1,2)\) curl. With ordinary
commuting shifts let
\[
 c_{12}(q)a=(1-e^{iq_2})a_1+(e^{iq_1}-1)a_2.
\]
The neutral site and vector symbols are exactly
\begin{equation}
 \Delta^0=\psi,\qquad
 K^0_\varphi=\psi I_4+
                  (\cos(\varphi/2)-1)c_{12}^*c_{12}
        =\psi I_4+\varphi^2N_2+O(\varphi^4),
 \quad N_2=-\tfrac18c_{12}^*c_{12}.
 \label{mlc:neutral}
\end{equation}
In particular
\begin{equation}
 \operatorname{tr}N_2(q)
   =-\tfrac18\bigl(|1-e^{iq_1}|^2+|1-e^{iq_2}|^2\bigr).
 \label{mlc:neutral-trace}
\end{equation}
This term is bounded and vanishes quadratically at zero momentum.
It is retained in the exact coefficient formula below.

We next justify that the logarithms used in the calculation
exist. By \eqref{mlc:flat-symbol} and the finite Fourier
description, the coefficient norm of
\(\sigma_\varphi-\psi I_4\) is at most \(C|\varphi|\).
This is a comparison with the magnetic scalar Laplacian, not
with the nonmagnetic operator in a globally small potential.
Since \(\Delta_\varphi\ge0\),
\begin{equation}
 K_\varphi\ge-C|\varphi|,
 \qquad K_\varphi+z\ge z/2
             \quad\text{if }|\varphi|\le c z.
 \label{mlc:mass-gap}
\end{equation}
The same lower bound holds for every compression. In the
neutral sector, for \(\cos f\ge0\), the sum of the weighted
curl squares and the divergence square is nonnegative. Thus
\(K^0_\varphi+z\ge z\). These observations prove positivity
of all finite determinants in
\eqref{mlc:finite-density} in a common real interval about zero,
for each fixed \(z>0\).

\subsubsection{Volume limits and differentiation in the field}

Here is a trace argument that deals with the linearly growing
potential. It will also justify differentiation of the determinant
density without a complex constant-field continuation.

\begin{lemma}\label{mlc:walk-limit}
Let \(T_\varphi=\operatorname{Op}_\varphi(t_\varphi)\) be a
Hermitian finite Weyl polynomial on \(d\) components. Suppose
\(t_0=\psi I_d\), and its finite Fourier coefficients are smooth
functions of real \(\varphi\). For fixed \(z>0\), the normalized
logarithms of the zero-extension compressions of
\(T_\varphi+z\) converge near \(\varphi=0\), together with any
fixed finite number of real flux derivatives. Their limits
are the cyclic trace of the logarithm defined by its convergent
power series. The convergence is locally uniform in positive
\(z\).
\end{lemma}

\begin{proof}
First restrict \(z\) to a compact positive interval. Choose
\(a>8+z\) throughout that interval and put
\[
 B_\varphi=I-a^{-1}(T_\varphi+z).
\]
At zero flux the sum of the norms of its combined Fourier
coefficients is
\[
 1-(8+z)/a+8/a=1-z/a<1.
\]
Coefficients with the same lattice shift are combined before
this norm is taken. Their smoothness and their finite number
give, on a sufficiently small real interval,
\(\|b_\varphi\|_0\le r<1\), uniformly in the restricted masses.
Consequently
\begin{equation}
 \log(T_\varphi+z)=(\log a)I
           -\sum_{n\ge1}\frac{B_\varphi^n}{n}
 \label{mlc:log-series}
\end{equation}
converges in operator norm and in the absolutely summable
coefficient algebra. The same series holds for each compressed
matrix, with the same norm bound.

Expand the trace of a power as a sum over walks. A term has
steps \(m_1,\ldots,m_n\) in a fixed finite set. It contributes
to the trace only if \(\sum_jm_j=0\). For such a closed word,
the product formula \eqref{mlc:weyl-product} gives a scalar
phase depending on the word and \(\varphi\), but not on its
starting site. Its phase exponent has absolute value at most
\(C n^2\). A finite compression keeps precisely the words whose
intermediate coordinates remain in the compression. A fixed
word is therefore lost only when its starting site lies within
distance \(Cn\) of the boundary. The proportion of such sites
tends to zero in a F\o lner exhaustion. The normalized trace
of each fixed power tends to its zero-coefficient trace.
For vector compressions with both endpoints in the box, the
omitted boundary bonds also have zero density. Their omission
changes only the same boundary set for each fixed word. The
constant term in the normalized logarithm tends to \(d\log a\),
because the number of retained vector coordinates divided by
the number of sites tends to \(d\).

For completeness, the flux derivative estimates must use the
combined coefficients of \(b_\varphi\); expanding separately
terms which cancel at zero flux would lose the bound \(r<1\).
For a derivative of order \(j\), differentiating coefficient
factors gives at most \(j\) exceptional factors, with uniformly
bounded coefficient norms. Differentiating the closed-word
phase gives a factor at most \(C_j n^{2j}\). Summing the other
factors in absolute value gives, after increasing the constant,
\[
 \left|\frac1{|\Lambda_R|}\partial_\varphi^j
            \operatorname{Tr}(P_RB_\varphi P_R^*)^n\right|
       \le C_j(1+n)^{3j}r^{\max\{n-j,0\}}.
\]
The same bound holds for the limiting zero-coefficient trace.
The sum of these bounds divided by \(n\) converges. It
justifies termwise differentiation, passage to the volume
limit, and their interchange. The argument is uniform on
compact positive mass intervals. It proves the assertion
for each finite number of derivatives.
\end{proof}

The proof applies to the charged vector polynomial by
\eqref{mlc:flat-symbol}, and to the charged scalar polynomial
directly. The neutral polynomial is an ordinary finite Fourier
polynomial with the same flat symbol, so it satisfies the same
conclusion. A constant colour rotation by \(it_1\) takes
\(U_f\) to \(U_{-f}\), preserves the compression coordinates,
and conjugates both real operators orthogonally. Hence every
finite density \(\mathcal D_R(f,z)\), and its limit, is even.

The cyclic trace in \eqref{mlc:log-series} will be denoted by
\(\tau\log(T_\varphi+z)\). This notation means the convergent
series with its scalar \(d\log a\) term. It is independent of
\(a\), since it is the limit of the finite logarithms. It also
satisfies
\begin{equation}
 \partial_z\tau\log(T_\varphi+z)
               =\tau(T_\varphi+z)^{-1}.
 \label{mlc:mass-jacobi}
\end{equation}
One may differentiate the locally uniformly convergent series
to prove this identity; cyclicity combines the \(n\) positions
of a differentiated factor. The same argument proves the
Jacobi identity for any bounded differentiable perturbation
which preserves a positive gap.

\subsubsection{The scalar orbital coefficient}

Let
\[
 u(q,z)=(\psi(q)+z)^{-1},\qquad
 a(q)=\psi_{11}(q)=2\cos q_1,\qquad
 b(q)=\psi_{22}(q)=2\cos q_2.
\]
For each fixed \(z>0\), \(u\) and all its momentum derivatives
have absolutely summable Fourier coefficients. For instance,
\(\psi+z\) stays away from zero on a sufficiently small complex
momentum strip, so the Fourier coefficients decrease
exponentially by shifting each integration contour. The strip
may depend on \(z\). Uniformity of that strip as \(z\downarrow0\)
is neither asserted nor needed here.

\begin{lemma}\label{mlc:scalar-coefficient}
If \(s_2(z)\) is the coefficient of \(\varphi^2\) in
\(\tau\log(\Delta_\varphi+z)\), then exactly
\par\noindent
\begin{minipage}{\linewidth}
\begin{equation}
 s_2(z)=\frac1{24}\int_{\mathbb B}a(q)b(q)u(q,z)^2\,d\bar q.
 \label{mlc:scalar-exact}
\end{equation}
\popQED\par
\end{minipage}
\end{lemma}

\begin{proof}
The Poisson bracket \(\{\psi,u\}\) is zero, because \(u\) is a
function of \(\psi\). The product expansion
\eqref{mlc:moyal} therefore gives a second-order parametrix for
the scalar resolvent with symbols
\[
 r_0=u,\qquad r_1=0,\qquad
 r_2=\frac18u\{\!\{\psi,u\}\!\}.
\]
More explicitly,
\begin{equation}
 (\psi+z)\star_\varphi(u+\varphi^2r_2)
                       =1+O_z(|\varphi|^3)
 \label{mlc:parametrix-residual}
\end{equation}
in the absolutely summable coefficient norm. All required
weighted norms are finite at fixed \(z\), so the remainder
estimate follows from \eqref{mlc:moyal}.

This parametrix determines the actual inverse coefficient.
Indeed, \(\Delta_\varphi+z\) has a convergent inverse series in
the same coefficient algebra, obtained from
\eqref{mlc:log-series} by the geometric series. Its coefficient
norm is locally bounded at fixed \(z\). Multiplying
\eqref{mlc:parametrix-residual} by that inverse shows that the
symbol of the actual resolvent differs from
\(u+\varphi^2r_2\) by \(O_z(|\varphi|^3)\). Thus no formal
inverse is being substituted for the lattice resolvent.

Since \(\psi_{12}=0\), differentiation of \(u\) gives
\[
 r_2=\frac14\left\{
 u^4(a\psi_2^2+b\psi_1^2)-u^3ab\right\}.
\]
Integration by parts on the momentum torus yields
\[
 0=\int_{\mathbb B}\partial_2(a\psi_2u^3)\,d\bar q
   =\int_{\mathbb B}(ab u^3-3a\psi_2^2u^4)\,d\bar q,
\]
and the same formula with indices exchanged. Consequently
\begin{equation}
 \tau\operatorname{Op}_\varphi(r_2)
       =\int_{\mathbb B}r_2\,d\bar q
       =-\frac1{12}\int_{\mathbb B}ab u^3\,d\bar q.
 \label{mlc:scalar-resolvent-coefficient}
\end{equation}
By \eqref{mlc:mass-jacobi}, the right side is \(s_2'(z)\).
The derivative in \(z\) of the right side of
\eqref{mlc:scalar-exact} is the same expression.

Both candidate coefficients tend to zero as \(z\to\infty\).
For the exact logarithm this follows from the series
\(\log(\Delta_\varphi+z)=\log z+
\log(1+z^{-1}\Delta_\varphi)\), whose flux derivatives converge
uniformly for large \(z\). For the integral it follows from
boundedness of \(a,b\) and \(u\le z^{-1}\). Equality of their
derivatives and of their limits at infinity proves
\eqref{mlc:scalar-exact}.
\end{proof}

The term \(s_2\) is the orbital contribution of one charged
scalar determinant. In particular it is already nonzero
before the vector spin term is included. Its sign and factor
come from the second magnetic bracket, not from a pointwise
Taylor expansion of the symbol \(\psi\).

\subsubsection{The vector logarithm and colour multiplicities}

We give the trace calculation with its noncommuting factors.
It prevents a scalar logarithm expansion from inadvertently
discarding the magnetic contribution. Put
\[
 B_\varphi=(\Delta_\varphi+z)I_4,\qquad
 \delta_\varphi=K_\varphi-\Delta_\varphi I_4
       =\operatorname{Op}_\varphi
          (\varphi\Sigma_1+\varphi^2\Sigma_2+O(\varphi^3)).
\]
At fixed positive \(z\), the segment
\(B_\varphi+t\delta_\varphi\), \(0\le t\le1\), has a
positive gap for small real \(\varphi\). The trace Jacobi
identity gives
\[
 \tau\log(B_\varphi+\delta_\varphi)-\tau\log B_\varphi
       =\int_0^1\tau\bigl((B_\varphi+t\delta_\varphi)^{-1}
                                  \delta_\varphi\bigr)\,dt.
\]
The exact twice-iterated resolvent identity and cyclicity then
give
\begin{equation}
 \tau\log(B_\varphi+\delta_\varphi)-\tau\log B_\varphi
 =\tau(B_\varphi^{-1}\delta_\varphi)
  -\frac12\tau(B_\varphi^{-1}\delta_\varphi
                        B_\varphi^{-1}\delta_\varphi)
  +O_z(|\varphi|^3).
 \label{mlc:trace-log-variation}
\end{equation}
For example, the remaining resolvent term contains three
copies of \(\delta_\varphi\), bounded inverse factors, and a
trace on four components per site, so its absolute value is
at most \(C_z\|\delta_\varphi\|_0^3\).

The scalar inverse has first flux coefficient zero by
\cref{mlc:scalar-coefficient}. The exact two-factor identity
\eqref{mlc:trace-product} therefore shows that the quadratic
coefficient of the first term on the right of
\eqref{mlc:trace-log-variation} is
\(\int u\operatorname{tr}\Sigma_2\,d\bar q\).
The quadratic coefficient of the second term is obtained at
zero flux, and equals
\(-\tfrac12\int u^2\operatorname{tr}\Sigma_1^2\,d\bar q\).
No first magnetic bracket contributes an additional cross
term: the scalar inverse has no linear symbol correction,
and the trace of a two-factor star product is the ordinary
integral exactly. If \(v_2(z)\) denotes the quadratic flux
coefficient of the charged vector logarithm, we have proved
\begin{equation}
 v_2(z)=4s_2(z)
       +\int_{\mathbb B}u\operatorname{tr}\Sigma_2\,d\bar q
       -\frac12\int_{\mathbb B}u^2
                                    \operatorname{tr}\Sigma_1^2\,d\bar q.
 \label{mlc:vector-coefficient}
\end{equation}
For the neutral vector logarithm the corresponding coefficient
is simply
\begin{equation}
 v_2^0(z)=\int_{\mathbb B}u\operatorname{tr}N_2\,d\bar q.
 \label{mlc:neutral-coefficient}
\end{equation}

There are two real charged colours. Equivalently, complexifying
the real two-colour space gives sectors with fluxes
\(\varphi\) and \(-\varphi\). The real determinant of a
positive complex-linear charged matrix is the square of its
complex determinant. Thus the scalar charged contribution to
the real logarithm is twice the one-sector scalar logarithm,
and likewise for the charged vector. The neutral scalar
logarithm is independent of \(f\). The coefficient of
\(\varphi^2\) in the combination
\(\log\det\Delta-\tfrac12\log\det K\) is therefore
\[
 2s_2-v_2-\tfrac12v_2^0
   =-2s_2-\int u\operatorname{tr}\Sigma_2\,d\bar q
         +\tfrac12\int u^2\operatorname{tr}\Sigma_1^2\,d\bar q
         -\tfrac12\int u\operatorname{tr}N_2\,d\bar q.
\]
Finally \(\varphi=2f\), and a quadratic \(\varphi\) coefficient
is multiplied by eight when it is differentiated twice in
\(f\). Combining this fact with
\eqref{mlc:scalar-exact} proves the exact formula
\begin{equation}
 \boxed{\displaystyle
 \widetilde\beta(z)=8\int_{\mathbb B}
 \left\{\frac12u^2\operatorname{tr}\Sigma_1^2
       -\frac1{12}ab u^2
       -u\operatorname{tr}\Sigma_2
       -\frac12u\operatorname{tr}N_2\right\}\,d\bar q .}
 \label{mlc:exact-beta}
\end{equation}
All four functions of momentum in this formula are independent
of \(z\). The last term is the neutral-colour correction.

\subsubsection{Infrared integrals and mass derivatives}

We now estimate \eqref{mlc:exact-beta} uniformly as the mass
decreases. On the chosen representatives of the momentum
torus,
\begin{equation}
 c|q|^2\le\psi(q)\le|q|^2,
 \qquad \psi(q)=|q|^2+O(|q|^4)\quad(q\to0).
 \label{mlc:dispersion}
\end{equation}
The lower bound follows, for example, by applying
\(|\sin(t/2)|\ge |t|/\pi\) on \([-\pi,\pi]\) to each term
\(2-2\cos t\).

\begin{lemma}\label{mlc:infrared-lemma}
Put \(I_2(z)=\int_{\mathbb B}(\psi+z)^{-2}\,d\bar q\).
There is a finite \(c_2\) such that
\begin{equation}
 I_2(z)=\frac1{16\pi^2}\log\frac1z+c_2+e_2(z),\qquad
 |e_2(z)|+|ze_2'(z)|+|z^2e_2''(z)|\le C\sqrt z.
 \label{mlc:I2-expansion}
\end{equation}
If \(h\) is smooth on \(\mathbb B\) and \(h(0)=0\), then
\(R_h(z)=\int h(\psi+z)^{-2}\,d\bar q\) has a finite
limit at zero and
\begin{equation}
 |R_h(z)-R_h(0)|+|zR_h'(z)|+|z^2R_h''(z)|\le C_h\sqrt z.
 \label{mlc:vanishing-numerator}
\end{equation}
For bounded smooth \(h\), the same conclusion holds for
\(\int h(\psi+z)^{-1}\,d\bar q\).
\end{lemma}

\begin{proof}
Choose a fixed ball \(|q|<\rho<\pi\). Outside this ball the
integrands and their first two mass derivatives are uniformly
bounded. Inside the ball, subtract
\((|q|^2+z)^{-2}\) from \((\psi+z)^{-2}\). The difference is
bounded by
\[
 C\frac{|q|^4}{(|q|^2+z)^3}.
\]
Its first and second mass derivatives are bounded respectively
by
\(C|q|^4(|q|^2+z)^{-4}\) and
\(C|q|^4(|q|^2+z)^{-5}\). Radial integration in four
dimensions gives an integrable limit for the difference,
a bound \(C(1+|\log z|)\) for the integral of its first
derivative, and \(C/z\) for that of its second derivative.
For instance the first derivative is bounded by a constant
times
\(\int_0^\rho r^7(r^2+z)^{-4}\,dr\); splitting at
\(r=\sqrt z\) gives the asserted logarithm. The second is
handled by the same split. Integrating the first derivative
from zero to \(z\) bounds the difference from its limiting
constant by \(Cz(1+|\log z|)\).

The remaining continuum radial integral is exact:
\[
 \int_{|q|<\rho}(|q|^2+z)^{-2}\,d\bar q
 =\frac1{16\pi^2}
   \left\{\log\frac{\rho^2+z}{z}
                   +\frac{z}{\rho^2+z}-1\right\}.
\]
Its nonsingular part has a finite limit and bounded first
two derivatives on a sufficiently small positive interval.
This proves \eqref{mlc:I2-expansion}, with the slightly
weaker but convenient error \(C\sqrt z\), since
\(z(1+|\log z|)\le C\sqrt z\) for \(0<z\le1\).

For the second assertion, smoothness and \(h(0)=0\) give
\(|h(q)|\le C_h|q|\) near zero. This makes the zero-mass
integral absolutely convergent. The first derivative has
absolute integral at most
\[
 C_h\int_0^\rho \frac{r^4}{(r^2+z)^3}\,dr\le C_hz^{-1/2},
\]
and the second at most
\(C_h\int_0^\rho r^4(r^2+z)^{-4}\,dr\le C_hz^{-3/2}\).
The estimates follow by scaling \(r=\sqrt z\,s\), extending
the resulting integrals to infinity, and using their
integrable large-\(s\) tails. Their weighted bounds are
therefore \(C_h\sqrt z\). Integration of the first derivative
from zero to \(z\) proves the remaining difference bound.

For the last assertion the zero-mass integrand is bounded
by \(C|q|^{-2}\), which is integrable in four dimensions.
The first derivative has integral at most
\(C(1+|\log z|)\), and the second at most \(C/z\).
The same argument again gives the stated, weaker
\(C\sqrt z\) error and weighted derivative bounds.
\end{proof}

\begin{proof}[Proof of \cref{mlc:main}]
The volume and field-differentiation assertions follow from
\cref{mlc:walk-limit}. Positivity and evenness were established
above. Formula \eqref{mlc:exact-beta} identifies the second
field derivative exactly.

By \eqref{mlc:spin-value},
\(\tfrac12\operatorname{tr}\Sigma_1(0)^2=4\), whereas
\(a(0)b(0)/12=1/3\). Thus the coefficient of the only
logarithmically divergent integral in
\eqref{mlc:exact-beta} is
\[
 8(4-\tfrac13)I_2(z)=\frac{88}{3}I_2(z).
\]
The remaining numerator in front of \(u^2\) vanishes at
\(q=0\), by \eqref{mlc:jet-bounds} and continuity of
\(a,b\). Its contribution is covered by
\eqref{mlc:vanishing-numerator}. The terms with one factor
of \(u\) have bounded smooth numerators and are covered by
the last assertion of \cref{mlc:infrared-lemma}. This
includes the neutral term; discarding it would change the
finite constant, though not the logarithm.

All these remainders have finite zero-mass limits and the
weighted derivative bounds in that lemma. Since
\[
 \frac{88}{3}\frac1{16\pi^2}=\frac{11}{6\pi^2},
\]
equation \eqref{mlc:I2-expansion} proves
\eqref{mlc:asymptotic}. Differentiating its logarithmic term
and using the weighted error estimates proves
\eqref{mlc:uniform-mass-bound}.
\end{proof}

\subsubsection{The second-order mass subtraction}

The mass comparison uses a Taylor polynomial evaluated at zero
mass, not merely the determinant at the comparison mass. The
preceding proof controls the two derivative terms as well.
Define
\[
 \widetilde\beta^{\sharp}(z)
   =\widetilde\beta(z)-z\widetilde\beta'(z)
                         +\tfrac12z^2\widetilde\beta''(z).
\]
\begin{corollary}\label{mlc:mass-subtraction}
For \(c=11/(6\pi^2)\),
\begin{equation}
 \widetilde\beta^{\sharp}(z)
       =c\log(1/z)+C_*+\tfrac32c+O(\sqrt z).
 \label{mlc:sharp-asymptotic}
\end{equation}
In particular, for fixed \(L>1\),
\begin{equation}
 \widetilde\beta^{\sharp}(L^{-2k})
       =\frac{11}{3\pi^2}k\log L+O(1),
 \qquad k\ge0.
 \label{mlc:scale-slope}
\end{equation}
The successive differences of these bare massive coefficients
converge to \(11\log L/(3\pi^2)\) with error \(O(L^{-k})\).
\end{corollary}

\begin{proof}
For \(c\log(1/z)\) the first derivative is \(-c/z\) and
the second is \(c/z^2\). Its two subtraction terms thus
add \(c+c/2\). The constant contributes only itself, and
\eqref{mlc:asymptotic} bounds the error and both subtracted
derivatives by \(C\sqrt z\). This proves
\eqref{mlc:sharp-asymptotic}. Substitution of \(z=L^{-2k}\)
proves \eqref{mlc:scale-slope}; subtraction at consecutive
values of \(k\) gives the last assertion.
\end{proof}

The successive differences in this corollary are defined by
the bare massive determinant. To identify them, up to a
uniformly bounded total error, with a sum of constrained
Gaussian increments requires the gauge-change identity, the
Schur-complement composition, and the uniform comparison
between the averaging regulator and the scalar mass regulator.
The present calculation supplies the lattice logarithm and
the mass-derivative estimates in that comparison. It does
not replace those three identifications. Nor does it imply
geometric convergence of the actual renormalization increments
from a bounded cumulative comparison alone.

% END LAYERED EXPANSION massive_lattice_coefficient.tex

% BEGIN LAYERED EXPANSION anisotropic_massive_coefficient.tex
\subsection{The diagonal anisotropy of the massive Wilson coefficient}
\label{amc:section}

The scalar logarithm does not determine the
response to a traceless change of the four
coordinate weights. We compute that response
from the weighted lattice Hessian. The computation
retains the lattice symbol until its infrared
estimate. In particular it differentiates an
actual family of Gaussian operators, rather
than differentiating a previously extracted
scalar asymptotic without a parameter bound.

Let \(a=(a_1,\ldots,a_4)\), initially real
and close to \((1,1,1,1)\), with all \(a_\mu>0\).
Write \(\mathsf A=\operatorname{diag}(a_\mu)\)
on the vector index, and put
\begin{equation}
 A_a(U)=\sum_{\mu<\nu}\sum_x
        a_\mu a_\nu
       \left(1-\tfrac12\operatorname{tr}U_{\mu\nu}(x)\right),
 \qquad
 \Delta_{a,U}=D_U^\sharp\mathsf A D_U .
 \label{amc:weighted-action}
\end{equation}
The colour inner product, covariant shifts,
and physical normalization are those in
\cref{mlc:section}. The symbol \(A_a\)
denotes an action; \(\mathsf A\) is its
constant diagonal weight.
Set
\begin{equation}
 \begin{split}
 K_{a,U}
   &=A_a''(U)+\mathsf A D_UD_U^\sharp\mathsf A,\\
 \widetilde K_{a,U}
   &=\mathsf A^{-1/2}K_{a,U}\mathsf A^{-1/2},
 \qquad \widetilde D_{a,U}=\mathsf A^{1/2}D_U .
 \end{split}
 \label{amc:weighted-vector}
\end{equation}
Square roots are positive for real \(a\);
near \(a=1\) their analytic branches are
fixed by that choice.

For a chosen coordinate plane \(\{\mu,\nu\}\),
use the constant Abelian field of
\eqref{mlc:background}, with the axes
renamed if necessary. Compress the full
operators after they have been formed,
exactly as in \cref{mlc:main}.
Define the vacuum-subtracted determinant
density from
\begin{equation}
 \mathcal L_{a,R}(f,z)=
 \log\det(\Delta_{a,f,R}+z)
       -\tfrac12\log\det(\widetilde K_{a,f,R}+z),
 \qquad z>0,
 \label{amc:determinant-density}
\end{equation}
and let \(\widetilde\beta_a^{\mu\nu}(z)\)
be its second \(f\)-derivative per lattice
site in the exhaustion limit.
All three real colours are included.

\begin{proposition}\label{amc:weighted-massive-law}
There is \(\delta>0\) such that
\(\widetilde\beta_a^{\mu\nu}(z)\) extends
holomorphically to
\(\max_\lambda|a_\lambda-1|<\delta\).
On every fixed smaller polydisc,
\begin{equation}
 \widetilde\beta_a^{\mu\nu}(z)
 =c\,d(a)a_\mu a_\nu\log(1/z)
                +C_{\mu\nu}(a)+E_{\mu\nu}(a,z),
 \qquad
 c=\frac{11}{6\pi^2},\quad
 d(a)=\prod_{\lambda=1}^4a_\lambda^{-1/2},
 \label{amc:weighted-logarithm}
\end{equation}
where \(C_{\mu\nu}\) is holomorphic.
For every fixed multi-index \(I\),
\begin{equation}
 \sup_a
 \sum_{j=0}^2
   \left|z^j\partial_z^j
                       \partial_a^I E_{\mu\nu}(a,z)\right|
                       \le C_I\sqrt z,\qquad 0<z\le1.
 \label{amc:uniform-parameter-error}
\end{equation}
For real \(a\), this is the coefficient
of the actual real weighted Wilson determinant.
The field variable \(f\) is differentiated
as a real variable at fixed \(z>0\).
No holomorphic constant-flux family is asserted.
\end{proposition}

\subsubsection{The gauge and the field-coordinate change}

The additional term in
\eqref{amc:weighted-vector} is the
square of the weighted divergence
\(D_U^\sharp\mathsf A v\).
Its Faddeev--Popov operator is
\(D_U^\sharp\mathsf A D_U=\Delta_{a,U}\).
For real \(a\), change vector coordinates by
\(v=\mathsf A^{-1/2}w\).
The Hessian and its gauge square become
\begin{equation}
 \widetilde K_{a,U}
  =\mathsf A^{-1/2}A_a''(U)\mathsf A^{-1/2}
                    +\widetilde D_{a,U}
                                      \widetilde D_{a,U}^\sharp .
 \label{amc:congruent-gauge}
\end{equation}
The constant Jacobian of this field-coordinate
change is independent of \(U\). It disappears
from the vacuum-subtracted density.
Equivalently the vector mass in the original
coordinates is \(z\mathsf A\):
\[
 K_{a,U}+z\mathsf A
   =\mathsf A^{1/2}
        (\widetilde K_{a,U}+z)\mathsf A^{1/2}.
\]
This equality also holds for the declared
zero-extension compressions.
It specifies the regulator used here;
comparison with an averaging regulator must
retain this choice.

Every constant Abelian field used here is
critical for \(A_a\) under compactly
supported variations. On a bond, the two
adjacent plaquettes in each plane contribute
opposite copies of the same weighted
\(\sin f_{\mu\nu}\,t_3\).
Gauge invariance therefore gives
\begin{equation}
 A_a''(U_f)D_{U_f}=0,\qquad
 \widetilde K_{a,U_f}\widetilde D_{a,U_f}
       =\widetilde D_{a,U_f}\Delta_{a,U_f}.
 \label{amc:weighted-Ward}
\end{equation}
The second identity is also obtained
directly from \eqref{amc:congruent-gauge}.
The spin coefficient below is computed from
the Hessian; the Ward identity is an
independent check on the chosen gauge.

\subsubsection{An exact weighted symbol}

It suffices to use the plane \((1,2)\).
Keep the magnetic shifts and the Weyl product
of \eqref{mlc:weyl-product}, with \(\varphi=2f\).
Let \(H_{\rho\sigma}\) be the exact charged
plaquette Hessian in
\eqref{mlc:exact-plaquette-hessian}.
The actual charged operators are
\begin{equation}
 \begin{split}
 \widetilde K_{a,\varphi}
  &=\mathsf A^{-1/2}
      \left(\sum_{\rho<\sigma}
                 a_\rho a_\sigma H_{\rho\sigma}\right)
                              \mathsf A^{-1/2}
       +\mathsf A^{1/2}DD^\sharp\mathsf A^{1/2},\\
 \Delta_{a,\varphi}
  &=\operatorname{Op}_\varphi(\psi_a),\qquad
 \psi_a(q)=\sum_{\rho=1}^4a_\rho(2-2\cos q_\rho).
 \end{split}
 \label{amc:exact-weighted-symbol}
\end{equation}
Both sides are finite Weyl polynomials.
At zero flux, their curl and divergence
terms cancel off the diagonal and supply
the same missing diagonal direction. Thus
\begin{equation}
 \widetilde\sigma_{a,0}(q)=\psi_a(q)I_4 .
 \label{amc:flat-weighted-symbol}
\end{equation}
For example the \(\mu\)-component of
the weighted curl square contains
\(\sum_{\rho\ne\mu}a_\rho D_\rho^\sharp D_\rho\);
the divergence square supplies
\(a_\mu D_\mu D_\mu^\sharp\).
This verifies the diagonal identity as
well as the off-diagonal cancellation.

Write the actual flux expansion as
\begin{equation}
 \widetilde\sigma_{a,\varphi}
  =\psi_a I_4+\varphi\Sigma_{a,1}
                    +\varphi^2\Sigma_{a,2}+O(\varphi^3).
 \label{amc:weighted-symbol-jets}
\end{equation}
Its coefficients are holomorphic in \(a\),
with uniformly bounded momentum derivatives
on a small complex polydisc.
Only the \((1,2)\) plaquette can change
the flat cancellation. Relative to the
unweighted calculation, its diagonal
corrections in the \(11\) and \(22\)
entries have multipliers \(a_2\) and \(a_1\);
its off-diagonal corrections have multiplier
\(\sqrt{a_1a_2}\).
All other corrections vanish.
This follows entry by entry from
\eqref{amc:exact-weighted-symbol}:
congruence divides the \(11\) entry of
\(a_1a_2H_{12}\) by \(a_1\), and its
\(12\) entry by \(\sqrt{a_1a_2}\).
The weighted divergence has exactly the
same factors in its corresponding entries.
Consequently
\begin{equation}
 \Sigma_{a,1}(0)
       =-2i\sqrt{a_1a_2}(E_{12}-E_{21}),\qquad
 \operatorname{tr}\Sigma_{a,1}(0)^2=8a_1a_2,
 \label{amc:weighted-spin}
\end{equation}
and
\begin{equation}
 \|\Sigma_{a,1}(q)-\Sigma_{a,1}(0)\|\le C|q|,
 \qquad \|\Sigma_{a,2}(q)\|\le C .
 \label{amc:weighted-jet-bounds}
\end{equation}
The same bounds hold after any fixed
number of \(a\)-derivatives.
The neutral colour has expansion
\[
 \widetilde K_{a,\varphi}^{\,0}
       =\psi_a I_4+\varphi^2N_{a,2}+O(\varphi^4),
\]
where \(N_{a,2}\) is a bounded finite
Fourier polynomial, holomorphic in \(a\).
Its \(11\) and \(22\) corrections have
the same respective weights \(a_2,a_1\).
This neutral term is retained in the
finite part of the coefficient.

\subsubsection{The weighted orbital and spin terms}

For \(z>0\), put
\[
 u_a(q,z)=(\psi_a(q)+z)^{-1},\qquad
 A_1(q)=2a_1\cos q_1,\quad A_2(q)=2a_2\cos q_2.
\]
The scalar Weyl calculation in
\cref{mlc:scalar-coefficient} applies
with \(\psi\) replaced by \(\psi_a\).
Its mixed second derivative is still zero.
The integration by parts has, for example,
the exact form
\[
 0=\int_{\mathbb B}
        \partial_2(A_1\,\partial_2\psi_a\,u_a^3)\,d\bar q
   =\int_{\mathbb B}
       (A_1A_2u_a^3
           -3A_1(\partial_2\psi_a)^2u_a^4)\,d\bar q .
\]
There is no change of magnetic phase,
and no approximation of the lattice
dispersion in this identity.
The scalar quadratic flux coefficient is
therefore
\begin{equation}
 s_{a,2}(z)=\frac1{24}
                 \int_{\mathbb B}A_1A_2u_a^2\,d\bar q .
 \label{amc:weighted-orbital}
\end{equation}

The trace logarithm is handled by
the same twice-iterated resolvent identity
as \eqref{mlc:trace-log-variation}.
At fixed real positive \(a,z\), the
segment between the scalar diagonal
operator and the full vector operator
is positive for sufficiently small
real flux. Its quadratic term is
\[
 4s_{a,2}
    +\int u_a\operatorname{tr}\Sigma_{a,2}\,d\bar q
    -\tfrac12\int u_a^2
                  \operatorname{tr}\Sigma_{a,1}^2\,d\bar q .
\]
The scalar inverse has no first flux
coefficient. The exact trace of a
two-factor Weyl product is its ordinary
momentum integral, so no first magnetic
bracket has been omitted here.
The two real charged colours and
\(\varphi=2f\) give the same multiplicity
eight as in \eqref{mlc:exact-beta}.
We obtain the exact formula
\begin{equation}
 \begin{split}
 \widetilde\beta_a^{12}(z)
   =8\int_{\mathbb B}\bigl\{&
      \tfrac12u_a^2\operatorname{tr}\Sigma_{a,1}^2
         -\tfrac1{12}A_1A_2u_a^2\\
      &-u_a\operatorname{tr}\Sigma_{a,2}
         -\tfrac12u_a\operatorname{tr}N_{a,2}\bigr\}\,d\bar q .
 \end{split}
 \label{amc:exact-weighted-beta}
\end{equation}
Its existence as a real exhaustion
limit follows from the closed-walk
proof in \cref{mlc:walk-limit}.
All step ranges remain fixed.
The coefficients and their parameter
derivatives are uniformly bounded
on compact sets of positive \(a,z\).
Every fixed number of flux and
parameter derivatives therefore has
the same dominated closed-walk limit.
For example a derivative of a
finite power inserts at most a
polynomial in its length, while the
fixed-mass resolvent series has
geometric decay.
Boundary walks occupy a vanishing
fraction of any declared F\o lner exhaustion.
This justifies taking the real field
derivative before the exhaustion limit.

For complex \(a\) close to one,
\(\operatorname{Re}\psi_a(q)\ge c|q|^2\)
on the standard momentum representatives.
The right side of
\eqref{amc:exact-weighted-beta}
is consequently holomorphic in \(a\)
at each \(z>0\). It defines the claimed
extension from its real values.
This extension concerns the weight
parameters; it does not complexify
the constant flux used in the
closed-walk argument.

\subsubsection{A uniform anisotropic infrared integral}

We prove the parameter estimate
needed to differentiate the logarithm.
For \(a\) in a small complex polydisc
put
\[
 I_{2,a}(z)=\int_{\mathbb B}(\psi_a+z)^{-2}\,d\bar q,
 \qquad
 H_a(t)=\int_{\mathbb B}e^{-t\psi_a(q)}\,d\bar q .
\]
For \(t\ge1\), uniformly on that polydisc,
\begin{equation}
 H_a(t)=\frac{d(a)}{16\pi^2t^2}+O(t^{-3}).
 \label{amc:weighted-heat-density}
\end{equation}
To verify the remainder, take a fixed
small momentum ball. On it,
\[
 \psi_a(q)=Q_a(q)+O(|q|^4),
 \quad Q_a(q)=\sum_\lambda a_\lambda q_\lambda^2,
 \quad \operatorname{Re}Q_a(q)\ge c|q|^2 .
\]
The difference of the two exponentials
is bounded by
\(Ct|q|^4e^{-ct|q|^2}\).
Its four-dimensional integral is
\(O(t^{-3})\).
Outside the ball the lattice integral
is exponentially small, and extending
the Gaussian integral to
\(\mathbb R^4\) costs the same type
of error. The latter integral is
\[
 (2\pi)^{-4}\int_{\mathbb R^4}e^{-tQ_a(q)}\,dq
                =\frac{d(a)}{16\pi^2t^2}.
\]
For complex \(a\), this identity follows
from the one-dimensional Gaussian
integrals on their branches near
the positive real axis. Every estimate
is uniform on a slightly larger
polydisc than the one on which
derivatives will be taken.

Since \(\operatorname{Re}(\psi_a+z)>0\),
Fubini gives
\begin{equation}
 I_{2,a}(z)=\int_0^\infty t e^{-tz}H_a(t)\,dt .
 \label{amc:heat-resolvent}
\end{equation}
Let \(c_a=d(a)/(16\pi^2)\).
Subtract \(c_at^{-2}\) only for \(t\ge1\).
The remaining integrand \(tH_a(t)\)
is bounded near zero and is
\(O(t^{-2})\) at infinity.
Thus
\begin{equation}
 I_{2,a}(z)=c_a\log(1/z)+C_I(a)+R_I(a,z),
 \qquad
 \sum_{j=0}^2|z^j\partial_z^jR_I(a,z)|
                         \le C\sqrt z .
 \label{amc:weighted-I2}
\end{equation}
Here is an explicit control of the
limit and its derivatives.
For the subtracted remainder \(r_a(t)\),
\(|r_a(t)|\le C\min(t,t^{-2})\).
Replacing \(e^{-zt}\) by one costs
\[
 \int_0^\infty
       |e^{-zt}-1|\,|r_a(t)|\,dt
                 \le Cz\{1+\log(1/z)\}.
\]
After one or two \(z\)-derivatives,
the corresponding estimates are
\[
 z\int_0^\infty t e^{-zt}|r_a(t)|\,dt
       \le Cz\{1+\log(1/z)\},\qquad
 z^2\int_0^\infty t^2e^{-zt}|r_a(t)|\,dt\le Cz .
\]
The subtracted leading term uses
\(\int_1^\infty e^{-zt}dt/t\).
Its difference from \(\log(1/z)\)
has a finite limit, an error \(O(z)\),
and weighted first and second
derivative errors \(O(z)\).
These estimates prove
\eqref{amc:weighted-I2}.
The constant \(C_I(a)\) is a
normally convergent integral plus
a holomorphic multiple of a fixed
number, and is holomorphic.
Cauchy's estimate on the larger
polydisc proves all fixed
\(a\)-derivative versions of
\eqref{amc:weighted-I2}.

Two elementary remainder estimates
complete the calculation.
If \(h_a(q)\) is holomorphic in \(a\)
and \(|h_a(q)|\le C|q|\), then
\begin{equation}
 \int_{\mathbb B}h_a(q)u_a(q,z)^2\,d\bar q
       =C_h(a)+R_h(a,z),\qquad
 \sum_{j=0}^2|z^j\partial_z^jR_h(a,z)|\le C\sqrt z.
 \label{amc:weighted-vanishing-numerator}
\end{equation}
The limit at zero is integrable.
Split its difference from the
positive-mass integral at
\(|q|=\sqrt z\).
The inner radial bound is
\(C\int_0^{\sqrt z}dr=C\sqrt z\);
the outer bound is
\(Cz\int_{\sqrt z}^1r^{-2}dr\le C\sqrt z\).
For the weighted first derivative,
use \(Cz\int_0^1r^4(r^2+z)^{-3}dr\);
for the second use
\(Cz^2\int_0^1r^4(r^2+z)^{-4}dr\).
Both are \(O(\sqrt z)\).
The region outside the unit ball
has a uniform positive dispersion
and obeys the same bound.
Normal convergence makes \(C_h\)
holomorphic. Cauchy again gives
its parameter derivative versions.

If the numerator is merely bounded
and there is one factor \(u_a\),
the zero-mass limit is also finite.
Its error is bounded by
\(Cz\{1+\log(1/z)\}\), its weighted
first derivative by the same
quantity, and its weighted second
derivative by \(Cz\).
These follow from the radial
integrals with numerator \(r^3\).
They too imply the bound \(C\sqrt z\).
This includes the neutral term
and the second vector-symbol jet.

\begin{proof}[Proof of \cref{amc:weighted-massive-law}]
At \(q=0\), the combined numerator
of \(u_a^2\) in
\eqref{amc:exact-weighted-beta} is
\[
 8\left(4a_1a_2-\frac13a_1a_2\right)
                            =\frac{88}{3}a_1a_2 .
\]
After this constant is subtracted,
the numerator vanishes at least
linearly in \(|q|\), by
\eqref{amc:weighted-jet-bounds}.
Apply \eqref{amc:weighted-I2}
to its constant part,
\eqref{amc:weighted-vanishing-numerator}
to its vanishing part, and the
single-resolvent estimate to the
remaining two terms.
The only logarithm is therefore
\[
 \frac{88}{3}a_1a_2
       \frac{d(a)}{16\pi^2}\log(1/z)
                  =c\,d(a)a_1a_2\log(1/z).
\]
All finite constants are
holomorphic in \(a\), and all
remainders have the stated
mass and parameter derivative
bounds. Renaming the two axes
proves every plane case.
\end{proof}

\subsubsection{The actual tracefree diagonal response}

Let \(e=(e_1,\ldots,e_4)\), and put
\(a_\lambda(t)=1+te_\lambda\).
The weight of the \((\mu,\nu)\)
plaquette is
\[
 a_\mu(t)a_\nu(t)
       =1+t(e_\mu+e_\nu)+O(t^2).
\]
Thus its derivative at zero is
exactly the prescribed diagonal
anisotropy; the harmless quadratic
choice of the weight family does
not change that derivative.
Define
\[
 \dot\beta_e^{\mu\nu}(z)
     =\left.\frac{d}{dt}\right|_{t=0}
                  \widetilde\beta_{a(t)}^{\mu\nu}(z).
\]
The derivative of the logarithmic
weight is
\begin{equation}
 \left.\frac{d}{dt}\right|_0
                  d(a(t))a_\mu(t)a_\nu(t)
       =e_\mu+e_\nu-\frac12\sum_{\lambda=1}^4e_\lambda .
 \label{amc:metric-weight-derivative}
\end{equation}
Consequently, if \(\sum_\lambda e_\lambda=0\),
\begin{equation}
 \dot\beta_e^{\mu\nu}(z)
       =c(e_\mu+e_\nu)\log(1/z)
                    +\dot C_e^{\mu\nu}+O(|e|\sqrt z),
 \label{amc:tracefree-diagonal-law}
\end{equation}
with the same two weighted mass
derivative estimates for the error.
For an isotropic direction
\(e_\lambda=s\), the logarithmic
derivative vanishes. This verifies
directly that the tracefree
diagonal channel and the scalar
weight channel have different
one-loop responses.

Apply the second-order mass
subtraction to the weighted family:
\[
 \widetilde\beta_a^{\mu\nu,\sharp}(z)
  =\left(1-z\partial_z+\tfrac12z^2\partial_z^2\right)
                              \widetilde\beta_a^{\mu\nu}(z).
\]
Parameter differentiation commutes
with this finite differential
operator. By
\eqref{amc:uniform-parameter-error},
\begin{equation}
 \widetilde\beta_a^{\mu\nu,\sharp}(z)
 =c\,d(a)a_\mu a_\nu\log(1/z)
       +C_{\mu\nu}(a)+\tfrac32c\,d(a)a_\mu a_\nu
       +O_a(\sqrt z),
 \label{amc:weighted-sharp-law}
\end{equation}
uniformly with any fixed number
of parameter derivatives.

\begin{corollary}\label{amc:bare-channel-slope}
Let \(h_k=L^{-k}\), with the fixed
blocking integer \(L\).
In the tracefree diagonal direction,
\begin{equation}
 \left.\frac{d}{dt}\right|_0
    \widetilde\beta_{1+te}^{\mu\nu,\sharp}(h_k^2)
   =\frac{11}{3\pi^2}(e_\mu+e_\nu)k\log L
                  +C_e^{\mu\nu}+O(|e|L^{-k}).
 \label{amc:bare-tracefree-cumulative}
\end{equation}
Its consecutive differences converge
geometrically to
\(11(e_\mu+e_\nu)\log L/(3\pi^2)\).
The analogous cumulative coefficient
in the isotropic weight direction
is bounded.
\end{corollary}

\begin{proof}
Substitute \(z=h_k^2\) in
\eqref{amc:weighted-sharp-law},
differentiate, and use
\eqref{amc:metric-weight-derivative}.
The two errors at consecutive
depths are each bounded by a
constant times \(L^{-k}\).
The coefficient of the logarithm
is zero in the isotropic direction,
so only its finite limit and
the bounded error remain there.
\end{proof}

This corollary concerns the bare
massive determinants just defined.
For the actual blocked anisotropic
increment one must also compare the
weighted constrained determinant,
its classical background, and its
regulator uniformly in the weight
parameter. The uniform parameter
estimate here is the input that
permits that comparison to be
differentiated. A bounded scalar
comparison without such a parameter
domain would not suffice.

% END LAYERED EXPANSION anisotropic_massive_coefficient.tex

% BEGIN LAYERED EXPANSION affine_background_coefficient.tex
\subsection{Affine response and preservation of the coupling coefficient}
\label{abc:section}

The coefficient of the action can be read from an affine field of
constant curvature. The nonlinear background map is defined on a
bounded field ball, so this readout uses the polynomial extension
of its linear response. We construct that extension from the actual
exponentially localized inverse and compute its value. No nonlinear
background is evaluated at an unbounded affine field.

Use the homogeneous construction of
\cref{prop:hb-actual-critical-orbit} at the trivial flat reference.
The coarse spacing is one and \(h=L^{-j}\), with \(j\ge0\).
The fixed flat fine gradient is \(D\), its curl is \(d\), and
the actual straightened derivative is \(Q\). The derivative
of the constructed strong-gauge representative will be denoted
by \(H_h\). The derivative of the literal physical representative
will be denoted by \(\widehat H_h\). Gauge covariance relates
them, but they are not identified as coordinate maps.

\subsubsection{Polynomial extension of the actual inverse}

Group fine coordinates into unit coarse cells. At a fixed
depth, each cell has finitely many fine coordinates. The
exponential kernel estimates for the joint inverse and its
coarse multiplier imply that its matrix entries have every
polynomial spatial moment. The following observation makes
the uniqueness argument precise.

\begin{lemma}\label{abc:polynomial-inverse}
Suppose a cell operator \(\mathcal L\) and a two-sided inverse
\(\mathcal G\) have absolutely summable exponentially
weighted row kernels, uniformly in the row. Then both act on fields of any fixed
polynomial growth, their products agree with their usual
matrix products there, and
\(\mathcal L\mathcal G=\mathcal G\mathcal L=1\) on that
space. In particular a polynomial-growth solution of
\(\mathcal Lu=0\) is zero.
\end{lemma}

\begin{proof}
For a field with \(|u(z)|\le C(1+|z|)^m\), use
\[
 (1+|z|)^m\le(1+|x|)^m(1+|z-x|)^m
\]
in the row defining the output at \(x\).
The exponential moment dominates the last polynomial,
so the row converges absolutely. In a product of two
kernels the corresponding double sum converges absolutely
by the same inequality applied twice. Fubini therefore
identifies it with the matrix product, whose identity
was already known on finitely supported fields.
\end{proof}

Apply this lemma to the actual flat joint operator
\[
 \mathcal L_h(A,\omega)
       =\bigl((\mathrm{d}^\sharp \mathrm{d}+DRD^\sharp)A-Q^\sharp\omega,\,QA\bigr).
\]
The site smoothing in \(R\) has exponential cell kernels,
as do its first derivatives; the remaining coefficients
have fixed physical carriers. The operator and its inverse
therefore satisfy the lemma at each fixed \(h\).
Their physical weighted estimates also show that the
extensions of their fixed-degree jets have constants
independent of depth. The finite-dimensional coordinate
count inside a cell is not used as a uniform bound.

\subsubsection{The actual response to constant curvature}

Fix a real antisymmetric tensor
\(F=(F_{\mu\nu})_{\mu,\nu=1}^4\), and one constant colour
\(t\), with \(\tfrac12\operatorname{tr}t^2=1\).
On the physical fine and coarse lattices set
\begin{equation}\label{abc:affine-fields}
 a_\mu(x)=\tfrac12\sum_\nu F_{\nu\mu}x_\nu t,\qquad
 b_\mu(y)=\tfrac12\sum_\nu F_{\nu\mu}y_\nu t.
\end{equation}
Their fine curl is the constant tensor \(F_{\mu\nu}t\).
Their fine divergence is zero, since
\(F_{\mu\mu}=0\). In particular
\[
                         \mathrm{d}^\sharp \mathrm{d}a=0,\qquad D^\sharp a=0.
\]

\begin{lemma}\label{abc:actual-affine-average}
For every homogeneous depth, the actual straightened
linear average satisfies \(Qa=b\).
\end{lemma}

\begin{proof}
Write \(n=h^{-1}\), which is odd. In the parallel frame,
the actual row is \eqref{hb:actual-flat-rows}:
\[
 (Qa)_\mu(y)=h^5\sum_{x\in B_n(y)}
                  \sum_{\ell=0}^{n-1}a_\mu(x+\ell he_\mu).
\]
The centred block has coordinate mean \(y\).
The line average shifts only the \(\mu\)-coordinate,
by \((1-h)/2\). Its coefficient in \(a_\mu\) is
\(F_{\mu\mu}/2=0\). The normalization of the block
and line weights is \(h^5n^5=1\). These facts give
\(Qa=b\) exactly.

The formula for the row also follows directly by
composing the actual one-step flat derivatives.
Balanced base-\(L\) digits run once through the centred
block of \(n\) sites in each direction, while ordinary
digits \(0,\ldots,L-1\) run once through the \(n\)
successive longitudinal steps. Thus no replacement of
the composed average by a continuum mean has occurred.
\end{proof}

\begin{proposition}\label{abc:actual-affine-response}
The polynomial extension of the actual strong-gauge
linear response satisfies
\begin{equation}\label{abc:affine-response}
                         H_hb=a,\qquad D\omega(0)b=0.
\end{equation}
The literal physical response has the form
\begin{equation}\label{abc:literal-affine-response}
                      \widehat H_hb=a+D\lambda_b,
\end{equation}
for a polynomial-growth fine gauge parameter
\(\lambda_b\). Its degree is bounded independently
of depth.
\end{proposition}

\begin{proof}
Differentiate the actual joint fixed-point equations
at zero. Their nonlinear residuals have zero first
derivative, so \((H_hb,D\omega(0)b)\) solves
\[
 (\mathrm{d}^\sharp \mathrm{d}+DRD^\sharp)H_hb-Q^\sharp D\omega(0)b=0,
                        \qquad QH_hb=b.
\]
The pair \((a,0)\) solves these same equations by
\cref{abc:actual-affine-average} and the zero
divergence and curl-divergence identities above.
It also satisfies the fixed Landau condition.
Both pairs have polynomial growth. Their difference
vanishes by \cref{abc:polynomial-inverse}, which
proves \eqref{abc:affine-response}.

The actual literal representative is obtained from
the strong-gauge one by the specified fine gauge in
\eqref{hb:literal-representative}. Differentiate that
identity at zero. The derivative of a fine gauge
is a fine gradient, giving
\eqref{abc:literal-affine-response}. Its gauge
parameter is the composition of the derivative of
the actual node factor, the local node extension,
and the response \(H_h\). The node factor has a
fixed physical carrier and bounded analytic
derivatives, while \(H_h\) has exponential cell
kernels. These operators map an affine input to
a polynomial-growth output of bounded degree.
This argument extends the linear jets only.
\end{proof}

For a single \((1,2)\) curvature component, the symmetric
potential in \eqref{abc:affine-fields} differs from the
potential \(a_1=-F_{12}x_2t,\ a_2=0\) by the gradient
of a quadratic polynomial, with sign determined by
the convention for \(D\). Thus either potential gives
the same gauge-invariant quadratic readout.

\subsubsection{The coefficient readout after substitution}

Use the pointed density convention of
\eqref{ph:pointed-total} and
\eqref{ph:exponential-moment}. Thus write
\(\mathcal F_h=\sum_z\mathcal F_{h,z}\), where each
localized scalar density is gauge invariant and is
normalized at the trivial field. Its Hessian kernel
at an anchor \(z\) satisfies, at each fixed mesh,
\begin{equation}\label{abc:pointed-two-leg-moment}
 \sum_{b,c}
   e^{m(\operatorname{dist}(b,z)+\operatorname{dist}(c,z))}
 \left\|D^2\mathcal F_{h,z}(1)[\delta_b,\delta_c]\right\|
                    <\infty
\end{equation}
for some \(m>0\). The anchors, physical volume
normalization and compatible pointing are part of
this convention. A sum of already localized
members with the same absolute moment has the
same meaning here.

Constant colour rotations force the first
derivative of each density to vanish. Its Hessian
annihilates fine gradients by the differentiated
local gauge Ward identity. The two-leg moment
\eqref{abc:pointed-two-leg-moment} extends that
identity to polynomial gauge parameters and
polynomial second arguments. In particular the
affine readout is taken in a density anchored at
\(z\), then summed with the prescribed density
of anchors per physical volume. It is not the
formal extensive pairing \(\langle a,Ka\rangle\)
of a translation-invariant bulk kernel with an
unbounded affine field.

For each localized density, the second derivative of
\(\mathcal F_{h,z}(\widehat U(B))\) has no term involving
the second derivative of the background: that term
is multiplied by \(D\mathcal F_{h,z}(1)=0\). Its affine
quadratic readout is consequently the rooted Hessian
applied to
\(\widehat H_hb\) twice. Equation
\eqref{abc:literal-affine-response} and the polynomial
Ward identity replace both arguments by \(a\).
Sum the resulting density identities with the
actual anchor weights; their partition of unity
preserves the physical density.
The action itself has density
\(\frac12\sum_{\mu<\nu}|F_{\mu\nu}|^2\) in this
normalization. We have proved
\begin{equation}\label{abc:coefficient-preserved}
       b[\mathcal F_h\circ\widehat U]
                     =b_h[\mathcal F_h].
\end{equation}
The subscript on the right specifies physical fine
spacing \(h\). In four dimensions this coefficient
has no extra power of the blocking scale.

For a successive map this also requires the actual
nested criticality identity, rather than an identification
of that map with a direct Wilson lift. We give the
argument explicitly. Work first in common physical
units. Write \(\widehat U^{(j)}\) for the full
Wilson critical lift to the original fine lattice,
\(\widehat U^{(j-1)}\) for the preceding lift, and
\(p_j\) for the one-step critical map of the actual
effective action \(S_{j-1}\). The common-chart
semigroup identity proved in
\cref{haf:physical-nesting} says
\[
       \widehat U^{(j)}(B)
          \sim \widehat U^{(j-1)}(p_j(B)),
\]
where the equivalence is a fine gauge equal to
one at the final coarse nodes. Apply the literal
\((j-1)\)-step mean. The right side is \(p_j(B)\).
The gauge on the left induces a gauge on the
intermediate coarse lattice. Consequently, at
the trivial field,
\begin{equation}\label{abc:successive-affine-response}
 Dp_j(0)b=Q_{\rm lit}^{(j-1)}\widehat H^{(j)}b
                         +\mathrm{d}\chi
                  =a_{\rm intermediate}+d\widetilde\chi.
\end{equation}
Here \(b\) is an affine final coarse potential
with constant antisymmetric curvature.
The symbol \(a_{\rm intermediate}\) denotes
its link variation in the same convention
as \(p_j\): it is \(\xi a_{\rm phys}\) at
intermediate bond centres for raw link-log
coordinates, where \(\xi=1/L\) in final
coarse units, and it is \(a_{\rm phys}\)
for variations normalized per unit length.
For the second equality, the direct fine response
is the affine fine potential modulo a gradient
by \eqref{abc:literal-affine-response}.
The straightened block-line mean reproduces
that affine potential at its intermediate bond
centres, with this length factor; see also
\eqref{wmc:mean-low-vertices}. Its literal mean differs by the
coarse frame gauge of the exact plain-mean
recursion. The linearized gauge covariance of
the literal mean takes a fine gradient to an
intermediate gradient. These observations
give exactly the last term in
\eqref{abc:successive-affine-response}, without
claiming that the literal derivative equals
the straightened average.

The same polynomial-growth extension used in
the direct affine response justifies these
linear identities on the infinite lattice.
In particular, for any compatible complete
pointed gauge-invariant density \(\mathcal F\),
its vanished first jet removes \(D^2p_j\),
and its two-leg Ward identity removes both
gradients in \eqref{abc:successive-affine-response}.
Thus
\begin{equation}\label{abc:successive-coefficient-preserved}
                           b[\mathcal F\circ p_j]=b[\mathcal F].
\end{equation}
This applies to the physical cumulative Gaussian
density in its successive-depth identity.
It concerns complete invariant densities, rather
than arbitrary fixed-reference Taylor pieces.

Finally change physical units between the
successive lattices.
Two factors \(L^{-1}\) from the one-form variation
and a factor \(L^4\) from the volume convert its
quadratic Hessian to \(L^2\widehat\pi(p/L)\);
the two powers of momentum remove that \(L^2\).
Equation \eqref{abc:coefficient-preserved} is the
coordinate-free affine version of this cancellation.
It establishes preservation of the coefficient;
it does not require equality of full polarization
symbols at nonzero momentum.

\subsubsection{Physical mass and the bare lattice formula}

For a physical mass parameter \(s>0\), let
\[
 \mathcal L_h(s,U)
   =\log\det(\Delta_{h,U}+s)
                   -\tfrac12\log\det(K_{h,U}+s),
\]
with its trivial-field value subtracted and with
the bulk trace convention of \cref{mlc:section}.
These are unconstrained site and Wilson-vector
operators. The fine physical variations and
inner products give
\begin{equation}\label{abc:physical-mass-scaling}
 \Delta_{h,U}=h^{-2}\Delta_{{\rm bare},U},\qquad
 K_{h,U}=h^{-2}K_{{\rm bare},U},\qquad z=h^2s.
\end{equation}
The scalar factors \(h^{-2}\) in their determinants
are field independent and disappear from the
normalized logarithm.

In a constant physical curvature \(tF\), the angle
of a fine plaquette in bare units is \(h^2tF\).
There are \(h^{-4}\) bare sites per unit physical
volume. Hence the coefficient per physical volume
is obtained by multiplying
\(\frac12\widetilde\beta(h^2s)(h^2tF)^2\)
by \(h^{-4}\). The two powers of plaquette area
cancel this density exactly. Thus
\begin{equation}\label{abc:mass-coefficient}
 b[\mathcal L_h(s,\widehat U(\cdot))]
                         =\widetilde\beta(h^2s).
\end{equation}
For the constant field, the symmetric and linear
potentials are related by an exact Abelian gauge.
Each finite compression of the massive bare
operators is conjugated by its unitary adjoint
matrices. Their determinant densities therefore
agree. Equivalently their quadratic jets agree
by the polynomial Ward identity.

At a fixed positive mass, the massive resolvent
and its mass derivatives have exponential kernels,
locally uniformly in that mass. Differentiation
therefore commutes with the second-moment
coefficient readout. Applying it once and twice
to \eqref{abc:mass-coefficient}, define
\[
 \mathcal L_h^\sharp
   =\mathcal L_h(1)-\partial_s\mathcal L_h(1)
                         +\tfrac12\partial_s^2\mathcal L_h(1).
\]
Then
\begin{equation}\label{abc:mass-Taylor-coefficient}
 \begin{split}
 b[\mathcal L_h^\sharp\circ\widehat U]
   ={}&\widetilde\beta(h^2)
           -h^2\widetilde\beta'(h^2)
           +\tfrac12h^4\widetilde\beta''(h^2).
 \end{split}
\end{equation}
This identity supplies the scale and normalization needed to use
\cref{mlc:mass-subtraction} in the physical
Gaussian comparison.

% END LAYERED EXPANSION affine_background_coefficient.tex

% BEGIN LAYERED EXPANSION background/mass_regulator_comparison.tex
\subsection{Uniform comparison of the averaging and mass regulators}
\label{ch:mrc-actual-comparison}

The massive coefficient and the finite gauge-change identity have
different regulators. We now compare them at the actual homogeneous
background, uniformly in the fine depth. The comparison includes the
constraint Schur determinants. Its ultraviolet remainder contains
three free resolvents, which are trace summable in four dimensions.
The few averaging directions are controlled by rank, rather than by
pretending that the averaging regulator is a scalar mass.

We use the background of \cref{prop:hb-actual-critical-orbit} on a
smaller fixed coarse ball. Fine scalar and vector spaces have physical
measure $h^4$, where $h=L^{-j}$; coarse spaces have counting measure.
At each actual background $U=\widehat U(B)$, re-center the fine
coordinate as $e^{iha}U$. Write $D=D_U$, $\Delta=D^\sharp D$,
and let $q=q_U$, $Q=Q_U$ be the actual site and straightened field
derivatives at this re-centered zero. Thus $qq^\sharp=1$ and
$QD=\mathrm{d}_cq$. They are not the effective $\tau_A$ of the fixed-flat
chart at a nonzero $A$. Let $A_W''(U)$ denote the actual physical
Wilson Hessian and put
\begin{equation}
 K=A_W''(U)+DD^\sharp,\qquad
 V_s=q^\sharp q,\qquad V_v=Q^\sharp Q-Dq^\sharp qD^\sharp.
 \label{mrc:actual-operators}
\end{equation}
The two averaging precisions are $\Delta+V_s$ and $K+V_v$.
The latter is the background-gauge precision with its field-row
penalty, rather than the Landau precision used to construct the
background. Its gap therefore requires a separate proof.

\subsubsection{A uniform gap for the flat background gauge}

All operators in this subsection are at a flat reference. Put
$P_0=1-q^\sharp q$, and retain $R_L$ for the previously proved
Landau projector. Define
\begin{equation}
 \Xi_0=\mathrm{d}^\sharp \mathrm{d}+DP_0D^\sharp+Q^\sharp Q.
 \label{mrc:flat-BG-precision}
\end{equation}

\begin{lemma}[The projected scalar Poisson equation]
\label{lem:mrc-projected-Poisson}
For $f\in\ker q$, the equations
\begin{equation}
                   q\lambda=0,\qquad P_0\Delta\lambda=f
 \label{mrc:projected-Poisson}
\end{equation}
have a unique solution and
\begin{equation}
 \|\lambda\|_2+\|D\lambda\|_2+\|D^2\lambda\|_2\le C\|f\|_2.
 \label{mrc:projected-Poisson-bound}
\end{equation}
The constants are uniform in mesh, volume and real flat holonomy.
\end{lemma}
\begin{proof}
The actual site Poincar\'e inequality on $\ker q$ makes
$\langle D\lambda,D\lambda\rangle$ coercive there. Its finite
variational equation gives existence, uniqueness and
$\|\lambda\|_{H^1}\le C\|f\|_2$. The missing site component is
a multiplier: $\Delta\lambda=f+q^\sharp\zeta$.
Read it with the smooth exact site lift $E_q$ of
\cref{lem:hb-row-seeds}. For every coarse $b$,
\[
 \langle\zeta,b\rangle
   =\langle D\lambda,DE_qb\rangle-\langle f,E_qb\rangle.
\]
The physical seed energy yields $\|\zeta\|_2\le C\|f\|_2$.
Since $q^\sharp$ is an isometry, $\|\Delta\lambda\|_2\le C\|f\|_2$.
In the flat parallel frame the Fourier identity
$\sum_{\mu,\nu}\|D_\mu D_\nu\lambda\|_2^2
=\|\Delta\lambda\|_2^2$ proves the last estimate. Twisted
periodic phases do not change this identity. No point readout of a
fine energy function is used.
\end{proof}

\begin{proposition}[Uniform background-gauge gap]
\label{prop:mrc-BG-gap}
There are $c,C>0$ such that
\begin{equation}
 c\|v\|_{H^1}^2\le\langle v,\Xi_0v\rangle\le C\|v\|_{H^1}^2.
 \label{mrc:flat-BG-gap}
\end{equation}
On a fixed smaller ball of real actual backgrounds the same comparison
holds for $K+V_v$ and for $\Delta+V_s$, relative to their free
physical $H^1$ norms. All constants are independent of $j,M$.
On its smaller complex continuation, the normalized precisions have
uniform relative-form Neumann inverses; no Hermitian inequality is
asserted at a complex background.
\end{proposition}
\begin{proof}
Use the proved flat Landau restoration to write
$v=s+D\lambda$, where $R_LD^\sharp s=0$ and $q\lambda=0$.
Its intertwining gives $Qs=Qv$ and its flat curl gives $\mathrm{d}s=\mathrm{d}v$.
The Landau gap therefore implies
\[
                 \|s\|_{H^1}\le C(\|\mathrm{d}v\|_2+\|Qv\|_2).
\]
The remaining scalar function satisfies
$P_0\Delta\lambda=P_0D^\sharp v-P_0D^\sharp s$.
Apply \cref{lem:mrc-projected-Poisson} to obtain
\[
 \|D\lambda\|_{H^1}
 \le C(\|P_0D^\sharp v\|_2+\|s\|_{H^1}).
\]
Together these inequalities prove the lower bound. The upper
bound follows from the actual flat row bounds and $\|P_0\|\le1$.

For the actual background, its literal normal coordinate has fixed
small $X_\alpha$ norm. The actual coefficient comparisons give
$D_U-D_0=O(r)$ as a zero-order $L^2$ operator and
$q_U-q_0,Q_U-Q_0=O(r)$ in physical $L^2$ row norms.
For the site map this follows directly by summing its disjoint
transported pillar averages; path lengths are bounded in physical
units and its transport errors are $O(r)$. For the field map use
the actual row/column density estimates. The exact Wilson Hessian
comparison bounds its form error by $Cr\|v\|_{H^1}^2$.
Expanding $D_U(1-q_U^\sharp q_U)D_U^\sharp$ as a form gives the
same error. Thus $K+V_v$ is a small form perturbation of $\Xi_0$.
The scalar statement is the actual site gap, or the identical
comparison with $\Delta_0+q_0^\sharp q_0$. A fixed smaller $r$
preserves both gaps uniformly.
\end{proof}

The gaps hold throughout the interpolation
\begin{equation}
 M=K_*+1,\qquad S_t=M+t(V-1),\qquad0\le t\le1,
 \quad (K_*,V)=(\Delta,V_s)\text{ or }(K,V_v).
 \label{mrc:interpolation}
\end{equation}
Indeed its form is the convex combination of $K_*+1$ and $K_*+V$.
At the flat reference both are uniformly comparable to $L+1$;
the just proved coefficient perturbation preserves this fact.
This observation also supplies coherent logarithm branches near
the real positive precisions.

\subsubsection{Fixed free normalization and the actual low-rank factors}

Set $G_0=(L+1)^{-1}$ on the flat fine scalar space, with identity
on colours and, for vectors, on components. All normalizations below
use this fixed free operator, independent of the background.
Define
\begin{equation}
 B=G_0^{1/2}MG_0^{1/2},\quad
 B_t=G_0^{1/2}S_tG_0^{1/2},\quad
 A=G_0^{1/2}VG_0^{1/2},\quad X=A-G_0.
 \label{mrc:fixed-normalization}
\end{equation}
The gaps prove uniform bounds for $B,B_t$ and their inverses;
they continue to a fixed smaller complex physical background ball
by their relative form Neumann series.

The actual rank factorization is explicit. In the scalar case
$A=(qG_0^{1/2})^\sharp(qG_0^{1/2})$. In the vector case it is
\begin{equation}
 A=(QG_0^{1/2})^\sharp(QG_0^{1/2})
       -(qD^\sharp G_0^{1/2})^\sharp(qD^\sharp G_0^{1/2}).
 \label{mrc:actual-rank-factorization}
\end{equation}
At a flat coarse Bloch fibre its rank is at most $3$ in the scalar
case and $12+3=15$ in the vector case. Its norm is uniformly
bounded. Indeed $Q,q$ have the actual physical row bounds, and
$D_U^\sharp G_0^{1/2}$ is bounded since its flat part is a
normalized derivative and its remaining part is zero order $O(r)$.
The negative sign in the vector regulator causes no problem for
this norm/rank assertion or for the interpolation gap.

We record why these bounds retain two actual background derivatives.
The site transport has physical path length bounded by a geometric
sum of scales. Its complex adjoint matrices and all fixed parameter
derivatives have bounded norm. Summation over disjoint pillars gives
bounded physical $q$ row norms on a fixed complex fine potential ball.
The actual straightened recursion gives bounded physical $Q$ row
norms on the same joint reference/increment ball.
To check that this joint ball is uniform after re-centering, truncate
the construction at scale $s$. Its formulas give both
$\|a_s\|_\infty\le Cr\xi_s$ and
$\|v_s-1\|_\infty\le Cr\sum_{r'<s}\xi_{r'}\le Cr\xi_s$.
Thus the plain reference $W_s=M^sU$ has raw logarithm of size
$Cr\xi_s$ in the flat frame. The joint one-step reference/increment
estimates therefore have summable errors at every remaining scale.
Their actual row/column products give the physical $Q$ bound on a
fixed ball, even for a complex reference. This does not assume that
the fixed-reference effective gauge derivative is a site average.
Banach-valued
Cauchy differentiation in the reference potential therefore bounds
the first two derivatives of these row operators; it does not infer
absolute entry sums from a general bilinear infinity-norm estimate.
The fine derivative $D_U$ has first parameter variation of zero
order $O(\|\delta A\|_\infty)$ and second variation
$O(h\|\delta A_1\|_\infty\|\delta A_2\|_\infty)$, directly from
the exponential link.

The physical Wilson Hessian can be checked without an unspecified
coefficient theorem. In one plaquette put
\[
 \mathcal C_Ua(p)=h^{-1}\sum_{b\in\partial p}
                 \epsilon_{p,b}\operatorname{Ad}(T_{p,b}(U))a_b,
\]
where $T_{p,b}$ is the ordered prefix which moves that link insertion
to the plaquette base. In a reversed edge its inverse-link insertion
is included in that prefix. Thus
$\delta U(p)U(p)^{-1}=ih^2\mathcal C_Ua(p)$ and
$\mathcal C_0=d$. Differentiating the four link factors a second
time splits their Hessian into a quadratic form in
$\mathcal C_Ua,\mathcal C_Ub$, with bounded analytic plaquette
coefficient, and a remainder with physical bound
$Ch^4|F_p|\sum_{b\in\partial p}|a_b|\sum_{b\in\partial p}|b_b|$.
To see the power, the two link insertions supply $h^2$; the part
not in their squared total insertion pairs with the antisymmetric
plaquette defect, of size $h^2F_p$. At the identity this part is
zero. The symmetric total-insertion coefficient is the bounded
matrix obtained from $U(p)+U(p)^{-1}$.

Every prefix contains at most four links. Consequently
$\mathcal C_U-d$ is zero order of size $C\|A\|_\infty$,
its first parameter derivative is zero order of size
$C\|\delta A\|_\infty$, and its second derivative has bound
$Ch\|\delta A_1\|_\infty\|\delta A_2\|_\infty$.
The curvature formula $F=\mathrm{d}A+E_h(A)$ bounds its first two parameter
jets by the corresponding $X_\alpha$ background norms. Differentiate
the preceding finite plaquette decomposition. Since
$\mathrm{d}G_0^{1/2}$ is bounded, it proves uniform bounds on the first two
derivatives of $G_0^{1/2}A_W''(U)G_0^{1/2}$. The same statement
for $DD^\sharp$ follows from the already stated zero-order
parameter derivatives of $D$. If these calculations are instead
performed in the fixed $A$ chart, the extra second coordinate
differential contains $hJ$; it is also bounded because
$h\|\mathrm{d}^\sharp \mathrm{d}\delta A\|_\infty\le C\|D\delta A\|_\infty$
and the complete Wilson residual is uniformly analytic.
We have therefore obtained the first two normalized $M$ jets from
the actual Wilson formula, without assuming the desired inverse.
The localized response estimates
\eqref{hb:localized-background-response} and
\eqref{hb:localized-background-jets} supply exactly those directions
after the physical background is substituted.

\subsubsection{The exact resolvent subtraction}

The following identity is finite dimensional and retains the full
regulator $V$, which need not be positive. Put $G=M^{-1}$ and
$W=V-1$. The resolvent identity applied twice gives
\[
 S_t^{-1}=G-tGWG+t^2GWGWS_t^{-1}.
\]
Integrate $\partial_t\log\det S_t=\operatorname{Tr}(S_t^{-1}W)$,
then expand $W=V-1$. One obtains exactly
\begin{equation}
 \begin{split}
 &\log\det(K_*+V)-\log\det(K_*+1)
                  +\operatorname{Tr}G+\tfrac12\operatorname{Tr}G^2\\
 &\quad=\operatorname{Tr}(GV)+\operatorname{Tr}(G^2V)
       -\tfrac12\operatorname{Tr}(GVGV)\\
 &\qquad\quad+\int_0^1t^2\operatorname{Tr}
                         (GWGWS_t^{-1}W)\,dt.
 \end{split}
 \label{mrc:exact-resolvent-Taylor}
\end{equation}
With the fixed normalization \eqref{mrc:fixed-normalization}, its
last trace is
\begin{equation}
 \operatorname{Tr}(B^{-1}XB^{-1}XB_t^{-1}X).
 \label{mrc:normalized-remainder}
\end{equation}
The three earlier terms become respectively
$\operatorname{Tr}(B^{-1}A)$,
$\operatorname{Tr}(B^{-1}G_0B^{-1}A)$ and
$\operatorname{Tr}(B^{-1}AB^{-1}A)$.
Every such term contains an actual finite coarse-rank factor.
Expand the three $X=A-G_0$ factors in
\eqref{mrc:normalized-remainder}. All terms except the pure
$G_0$ term again contain $A$, and are bounded by its rank times
the product of the already bounded operator norms. The remaining
term is bounded by Schatten H\"older:
\begin{equation}
 |\operatorname{Tr}(B^{-1}G_0B^{-1}G_0B_t^{-1}G_0)|
                                   \le C\|G_0\|_{\mathcal S_3}^3.
 \label{mrc:three-resolvent-bound}
\end{equation}
One must apply this reasoning per physical coarse cell, not with
the rank of a full finite-volume matrix. The next argument does so
and proves the two external momentum derivatives needed for the
coefficient.

\subsubsection{Two external jets without an unbounded complex plane wave}

Work at the trivial background for the bulk Hessian. Coarse Bloch
momentum is $\kappa\in[-\pi,\pi]^4$. At fine spacing $h=1/n$ the
free aliases in one cell have momenta $q=\kappa+2\pi m$, with
$m$ in $n$ consecutive centered integer values per coordinate.
Their Laplacian symbol is
\begin{equation}
                 \lambda_h(q)=4h^{-2}\sum_\mu\sin^2(hq_\mu/2).
 \label{mrc:free-alias-symbol}
\end{equation}
At real $\kappa$, $\lambda_h(q)\ge c|q|^2$ on the fine Brillouin
cube. A fixed small complex shift $z$ obeys
\begin{equation}
 \operatorname{Re}(1+\lambda_h(q+z))
                      \ge c(1+|q|^2),\qquad |z|<s_1,
 \label{mrc:complex-free-bound}
\end{equation}
after a uniform decrease of $s_1$. To verify it, the imaginary
shift changes the real symbol by at most $C|\operatorname{Im}z|^2$;
the real shift remains in a fixed slightly enlarged Brillouin cube,
where the same sine lower bound holds. Thus the principal square
root of $G_0(\kappa+z)$ is analytic on a fixed strip. In particular
\begin{equation}
 \int_{[-\pi,\pi]^4}\|G_0(\kappa+z)\|_{\mathcal S_3}^3
                         \frac{d\kappa}{(2\pi)^4}
 \le C\int_{\mathbb R^4}(1+c|q|^2)^{-3}\,dq<\infty.
 \label{mrc:uniform-Schatten-three}
\end{equation}
The sum over the aliases simply partitions the fine Brillouin cube.
The fixed colour/component multiplicities are included in $C$.

Let two external coarse variations have momenta $p$ and $-p$.
At first take $p$ real, and use the formal coefficient ring
$\mathbb C[\varepsilon_1,\varepsilon_2]/(\varepsilon_1^2,
\varepsilon_2^2)$. Substituting
$B=\varepsilon_1be^{ipx}+\varepsilon_2ce^{-ipx}$ records the
zeroth, two first and mixed second jets of each actual operator.
This is a jet computation at zero, not an assertion that a complex
plane-wave field lies in a globally bounded background domain.

Represent the jets on four graded fine-cell spaces indexed by
$\varnothing,\{1\},\{2\},\{1,2\}$. Their momenta are
$\kappa,\kappa+p,\kappa-p,\kappa$. A jet block from grade $T$
to grade $S$ is zero unless $T\subset S$, and otherwise is the
derivative with the external indices in $S\setminus T$.
Operator multiplication and inversion of these finite triangular
matrices are exactly the product rule and inverse rule for the
mixed two-jet. The free matrix $\mathbf G_0$ is diagonal in these
four grades, with the four shifted free symbols just specified.
Use it on both sides to normalize every jet matrix.

All normalized jet blocks of $\mathbf B,\mathbf B_t,\mathbf A$
are bounded uniformly in $h,\kappa,t$, on a fixed complex
neighborhood of $p=0$. Here are the details of this assertion.
The localized first and second background jets yield absolutely
convergent Fourier sums for their cell profiles when
$|\operatorname{Im}p|<s/4$. Their strong norms are uniformly
bounded by \eqref{hb:localized-background-jets}: the exponential
kernel bound dominates both external factors, after decreasing
its rate. Applying the finite-carrier transport and plaquette
formulas proved above gives bounded normalized $q,Q,D,K$ jets.
Every factor comparing two fine momenta is also bounded by
\eqref{mrc:complex-free-bound}. Thus all off-diagonal normalized
blocks are bounded without a fine-volume factor.

For the diagonal inverses, the actual flat precision gaps give a
uniform relative form inverse. Conjugation by a small physical
exponential changes each local derivative form by $O(s)$ in the
$H^1,H^{-1}$ pairing, and changes each averaging row form by
$O(s)$ because its physical carrier is bounded. Consequently the
flat interpolation inverse remains bounded for the shifted complex
Bloch momenta in a fixed strip. This proves uniform bounds for
the diagonal blocks of $\mathbf B_t^{-1}$; the corresponding
statement for $\mathbf B^{-1}$ is immediate from the free massive
operator. Inverting a four-grade triangular matrix then uses a
finite sum with at most two nontrivial jet steps. The inverses of
the full jet matrices are therefore bounded uniformly as well.

The rank statement also survives with no derivative-count problem.
Apply the jet functor to \eqref{mrc:actual-rank-factorization}.
It still factors through four copies of the coarse site/field
spaces. The rank of $\mathbf A$ is at most $12$ in the scalar case
and $60$ in the vector case. In this statement the jet of $Q^\sharp$
is the analytic derivative of that operator; it is not the ordinary
transpose of the four-grade jet matrix, which would reverse grades.
The actual factorization, and hence this finite-rank bound, remains
valid for complex external momentum.

Apply \eqref{mrc:exact-resolvent-Taylor} in the jet algebra. The
mixed trace is the trace of its $(\{1,2\},\varnothing)$ block,
not the ordinary trace of the entire triangular matrix. For real
$p$, cyclic permutation of trace factors is valid after integrating
in $\kappa$: it only reindexes $\kappa$ by one of the real external
momenta, and the coarse Brillouin torus and its fibre identifications
are periodic. Thus the displayed normalized trace identities give
the desired mixed bulk Hessian.

Each term containing $\mathbf A$ has uniformly bounded trace norm
by its rank and the operator bounds. For the pure free term apply
Schatten H\"older to its three diagonal free factors, and H\"older
once more in $\kappa$ if their momenta are different. Equation
\eqref{mrc:uniform-Schatten-three} bounds the resulting integral
uniformly. Projection onto one graded block cannot increase trace
norm, so it also bounds the required mixed trace. Integration over
$0\le t\le1$ preserves that bound.

Every expression just obtained is holomorphic in a fixed complex
neighborhood of $p=0$, by the same integrable bounds and analytic
inverse formulas. The equalities first proved for real $p$ extend
there by the identity theorem. This orders the steps correctly:
real cyclic reindexing is performed before analytic continuation;
there is no unproved change of a complex Bloch integration contour.
We have proved a uniform holomorphic bound for the mixed bulk
polarization of the entire remainder in
\eqref{mrc:exact-resolvent-Taylor}. Cauchy's formula bounds each
second external momentum derivative by a uniform constant.

For completeness the same construction justifies the bulk limit.
At fixed $h$ the gapped fine and coarse Bloch inverses have
exponential cell kernels. Finite torus kernels are their periodic
sums, and every differentiated trace above is a sum of products
with absolutely convergent spatial moments. Terms winding around
the torus vanish when its side tends to infinity. The bulk
polarization is the Bloch integral used above. Its bounds are
uniform in $h$, because rank and
\eqref{mrc:uniform-Schatten-three} were uniform before taking this
limit. A full finite-volume trace was never bounded by a fixed rank.

\subsubsection{The constraint Schur determinants}

Let $T_v=K-Dq^\sharp qD^\sharp$ and $T_s=\Delta$. For either
$C=Q$ or $C=q$, elementary block elimination gives
\begin{equation}
 |\det\mathbb S(T_*,C)|
   =\det(T_*+C^\sharp C)\,
          \det\{C(T_*+C^\sharp C)^{-1}C^\sharp\}.
 \label{mrc:actual-saddle-factorization}
\end{equation}
It can also be checked in orthogonal $\ker C$ coordinates. The
field/site constraint is onto, and the first precision has the
uniform gap proved above. Its coarse Schur matrix has uniform
positive lower and upper bounds. For the lower bound use the
actual exact smooth row lift $E_C$ and its physical energy:
\[
 \|b\|_2^2
 =\langle (T_*+C^\sharp C)(T_*+C^\sharp C)^{-1}C^\sharp b,E_Cb\rangle
 \le C\langle b,C(T_*+C^\sharp C)^{-1}C^\sharp b\rangle^{1/2}\|b\|_2.
\]
The upper bound follows from the gap and the bounded row form.
The actual nonflat lift is obtained from the flat smooth seeds by
the row comparison and its coarse Neumann inverse; its energy
bound is uniform on the present physical strong ball.

The Schur matrix and its inverse have a fixed holomorphic momentum
strip. This follows either by weighted relative conjugation of the
precision and its fixed-carrier rows, or directly by its normalized
jet formula. Its mixed logarithmic trace derivative is
$\operatorname{Tr}(S^{-1}S_{12}-S^{-1}S_1S^{-1}S_2)$, with the
appropriate shifted momenta and integrated cyclic convention.
The fine inverses in these coefficients have already been bounded.
There are only $12$ coarse vector colours and $3$ coarse site
colours per cell. Hence the mixed Schur-log polarization is
bounded holomorphically, uniformly in $h$. Cauchy's formula again
bounds its second momentum coefficient. Coherent branches are
fixed by the positive real Schur matrices; only the displayed real
saddle identity uses absolute values.

\subsubsection{The uniform physical coefficient comparison}

Let
\begin{equation}
 \mathcal L_h(s,U)=\log\det(\Delta_U+s)
                       -\tfrac12\log\det(K_U+s),\qquad
 \mathcal L_h^\sharp=\mathcal L_h(1)-\partial_s\mathcal L_h(1)
                                    +\tfrac12\partial_s^2\mathcal L_h(1),
 \label{mrc:mass-Taylor-functional}
\end{equation}
with trivial-field values subtracted and bulk trace per physical
coarse volume. Let $\Psi_j$ be the actual physical homogeneous
Gaussian logarithm of \eqref{pgg:physical-RF12}.

\begin{proposition}[Averaging versus mass with uniform two-jet error]
\label{prop:mrc-uniform-regulator-comparison}
On the constructed homogeneous background domain,
\begin{equation}
 \left|b[\Psi_j]-b[\mathcal L_h^\sharp\circ\widehat U]\right|\le C,
                  \qquad h=L^{-j},\quad j\ge0,
 \label{mrc:uniform-coefficient-comparison}
\end{equation}
where $C$ is independent of $j$. More precisely, the difference
after removal of the physical multiplier term has a bulk Hessian
bounded holomorphically on a fixed external momentum neighborhood.
\end{proposition}
\begin{proof}
Apply \eqref{mrc:exact-resolvent-Taylor} to the scalar and vector
precisions, with coefficients $1$ and $-1/2$. Their left sides
combine as the averaging logarithm minus
$\mathcal L_h(1)-\partial_s\mathcal L_h(1)
                    +\tfrac12\partial_s^2\mathcal L_h(1)$.
Indeed the two added trace terms on that left side combine to
$\partial_s\mathcal L_h(1)-\tfrac12\partial_s^2\mathcal L_h(1)$,
since the first log-determinant derivative is the resolvent trace
and the second is minus the square trace. Moving them to the
right gives precisely \eqref{mrc:mass-Taylor-functional}. The uniform jet
argument bounds the second-momentum coefficient of both remainder
terms. Equation \eqref{mrc:actual-saddle-factorization} adds only
the two coarse Schur logarithms, whose coefficients were uniformly
bounded. Finally \cref{pgg:current-zero-coefficient} proves that
the actual physical multiplier correction has coefficient zero.
This is a statement about its determinant functional; no false
pointwise third-order claim for the matrix $-\mu q''$ is used.
The finite-dimensional coefficient readout $b$ is a fixed linear
combination of these second momentum derivatives. The Cauchy bounds
therefore prove \eqref{mrc:uniform-coefficient-comparison}.
\end{proof}

The actual affine-response identity and physical mass rescaling in
\eqref{abc:mass-Taylor-coefficient} identify the massive term:
\begin{equation}
 b[\Psi_j]=\widetilde\beta(h^2)-h^2\widetilde\beta'(h^2)
                  +\tfrac12h^4\widetilde\beta''(h^2)+O(1).
 \label{mrc:actual-RF13-comparison}
\end{equation}
This uses the actual coefficient-preserving background pullback,
not an arbitrary smooth coarse interpolation. Combining the
massive lattice coefficient \eqref{mlc:scale-slope} yields the
cumulative physical law
\begin{equation}
                   b[\Psi_j]=\frac{11}{3\pi^2}j\log L+O(1).
 \label{mrc:cumulative-physical-law}
\end{equation}
By the exact physical Gaussian stage composition and coefficient
preservation \eqref{abc:coefficient-preserved}, its one-step
increments $b_\ell^0=b[\Phi_\ell]$ therefore satisfy
\begin{equation}
              \sum_{\ell=1}^j b_\ell^0
                          =\frac{11}{3\pi^2}j\log L+O(1).
 \label{mrc:physical-increment-sum}
\end{equation}
The bounded error here is a cumulative comparison. It does not
imply geometric convergence, or even convergence, of the actual
individual increments: a difference of two bounded comparison
errors need not tend to zero. The geometric error in the bare
massive coefficient is a separate, stronger property of that
particular regulator.

% END LAYERED EXPANSION background/mass_regulator_comparison.tex

% BEGIN LAYERED EXPANSION classical_mesh_consistency.tex
\subsection{Interpolation and consistency of the actual lattice mean}
\label{wmc:section}

To compare the classical effective actions at successive fine depths,
we need trial fields on both meshes with controlled curvature and
controlled coarse mean. The trial field must use the actual group
average. This subsection supplies the interpolation and consistency
estimates for that comparison.

Work at the trivial flat reference on the periodic physical torus
\(\mathbb T_M^4\), \(M\ge2\). This is the reference used to extract the
Taylor coefficients. Put \(h=L^{-j}\), with odd \(L\ge3\), and fix
\(0<\beta<1\). Constants can depend on \(L,\beta\) and the fixed
small coordinate ball, but not on \(j\) or \(M\). A fine bond field
\(a_\mu(x)\) is regarded as data at its physical midpoint
\(x+\frac h2e_\mu\). Write \(\|a\|_{X_\beta}\) for the physical
first-difference and local H\"older norm in
\eqref{hb:norms}. At the trivial reference all differences are
ordinary differences. Every continuous connection below is periodic.

\subsubsection{An exact interpolation with one H\"older derivative}

We first give the scalar interpolation; apply it to each colour
component on each of the four shifted midpoint grids. On a grid
of spacing \(h\), let \(a_x\) be the values and put
\[
 g_x=(D_1^+a(x),\ldots,D_4^+a(x)).
\]
Suppose \(\|a\|_\infty+\|g\|_\infty+[g]_{\beta,\rm loc}\le K\).
For grid points \(x,y\) with \(h\le|x-y|\le1\), a coordinate
path and telescoping give
\begin{equation}\label{wmc:discrete-Whitney-jets}
 |g_x-g_y|\le CK|x-y|^\beta,\qquad
 |a_y-a_x-g_x\cdot(y-x)|\le CK|x-y|^{1+\beta}.
\end{equation}
For the second inequality, each directed path increment is its
forward grid difference. Compare that difference with its value
at \(x\); all intervening points are within \(C|x-y|\) of \(x\).
There are at most \(C|x-y|/h\) increments, each with an error
\(CKh|x-y|^\beta\). Negative increments use the forward
difference at their other endpoint and obey the same estimate.
For fixed multiples of the unit radius the supremum terms
give the analogous bound after increasing \(C\).

Choose a smooth nonnegative product bump \(\rho\), supported
in \((-1,1)^4\), positive on \([-3/4,3/4]^4\), with
\(\rho(0)=1\). On the periodic lift of the grid define
\[
 \eta_x(X)=
   \frac{\rho((X-x)/h)}
        {\sum_z\rho((X-z)/h)},\qquad
 P_x(X)=a_x+g_x\cdot(X-x),\qquad
                         \mathcal I_ha(X)=\sum_x\eta_x(X)P_x(X).
\]
Only a fixed number of terms occur at each \(X\). The denominator
has positive lower and finite upper bounds, because a nearest
grid point is at coordinate distance at most \(h/2\).
Consequently
\[
 \sum_x\eta_x=1,\qquad
 |\partial^\ell\eta_x|\le C_\ell h^{-|\ell|}.
\]
At a grid point \(y\), the only nonzero bump is the one at \(y\).
All derivatives of bumps centered at its neighboring grid points
vanish there, since their arguments are on the boundary of the
open support. Hence
\begin{equation}\label{wmc:cardinal-interpolation}
             \mathcal I_ha(y)=a_y,\qquad
                          \nabla\mathcal I_ha(y)=g_y.
\end{equation}
The periodic lift makes these definitions independent of a
fundamental cube. A componentwise shifted grid is treated by
the same formula.

\begin{proposition}[Uniform interpolation of the fine connection]
\label{wmc:interpolation}
There is a linear periodic interpolation \(\mathcal I_h\)
from midpoint bond data to smooth continuous connections
such that
\begin{equation}\label{wmc:interpolation-bound}
 (\mathcal I_ha)_\mu(x+\tfrac h2e_\mu)=a_\mu(x),\qquad
               \|\mathcal I_ha\|_{C^{1,\beta}}\le C_\beta\|a\|_{X_\beta}.
\end{equation}
For a continuous \(C^{1,\beta}\) connection \(A\), its midpoint
samples
\[
                (\mathcal S_hA)_\mu(x)=A_\mu(x+\tfrac h2e_\mu)
\]
satisfy
\begin{equation}\label{wmc:sampling-bound}
                         \|\mathcal S_hA\|_{X_\beta}
                                      \le C_\beta\|A\|_{C^{1,\beta}}.
\end{equation}
Both constants are independent of the physical volume and
the number of fine sites in a unit cell.
\end{proposition}

\begin{proof}
For \(X\) within \(Ch\) of a grid point \(x\), compare each
polynomial in the interpolation formula with \(P_x\).
Equation \eqref{wmc:discrete-Whitney-jets} gives
\[
 |P_y(X)-P_x(X)|\le CK h^{1+\beta},\qquad
 |g_y-g_x|\le CK h^\beta
\]
whenever \(\eta_y(X)\ne0\). Differentiating the partition
formula once and twice therefore gives
\begin{equation}\label{wmc:local-interpolation-estimates}
 |\nabla\mathcal I_ha(X)-g_x|\le CK h^\beta,\qquad
                 |\nabla^2\mathcal I_ha(X)|\le CK h^{\beta-1}.
\end{equation}
The second bound is not a uniform second-derivative estimate;
only its consequence for the first-derivative H\"older norm
will be used.

If \(|X-Y|\le h\), integrate the second bound to obtain
\[
 |\nabla\mathcal I_ha(X)-\nabla\mathcal I_ha(Y)|
       \le CK h^{\beta-1}|X-Y|\le CK|X-Y|^\beta.
\]
If \(h\le|X-Y|\le1\), compare the two gradients with \(g_x,g_y\)
at their nearest grid points and use
\eqref{wmc:discrete-Whitney-jets}. The same H\"older bound
follows. Zeroth and first suprema follow directly from
the local formula, proving \eqref{wmc:interpolation-bound}.
For vector fields use the appropriate shifted grid for
each component. This preserves reality of the colour
coordinates.

For sampling, a forward difference is an average of the
corresponding continuous derivative on a segment of length
\(h\). Its supremum and its translated H\"older difference
are bounded by the continuous \(C^{1,\beta}\) norm. This
proves \eqref{wmc:sampling-bound}.
\end{proof}

\subsubsection{Wilson action consistency at this regularity}

For a continuous Hermitian connection \(A\), put
\[
 U_h(A)_\mu(x)=
        \exp\bigl(ihA_\mu(x+\tfrac h2e_\mu)\bigr),\qquad
 F_{\mu\nu}(A)=\partial_\mu A_\nu-\partial_\nu A_\mu
                                             +i[A_\mu,A_\nu].
\]
The norm of a Hermitian colour matrix is the one with
\(|X|^2=\tfrac12\operatorname{tr}X^2\). Let
\[
 \mathscr E(A)=\frac12
         \int_{\mathbb T_M^4}\sum_{\mu<\nu}|F_{\mu\nu}(A)(x)|^2\,dx.
\]

\begin{proposition}[Consistency of the Wilson energy]
\label{wmc:energy-consistency}
On a fixed sufficiently small \(C^{1,\beta}\) ball,
\begin{equation}\label{wmc:energy-consistency-bound}
               |A_h(U_h(A))-\mathscr E(A)|
                                  \le C_\beta h^\beta M^4.
\end{equation}
Here \(A_h\) is the actual Wilson sum over fine positive
plaquettes. In particular, for the same interpolated
connection,
\begin{equation}\label{wmc:two-mesh-energy}
               |A_{h/L}(U_{h/L}(A))-A_h(U_h(A))|
                                  \le C_\beta h^\beta M^4.
\end{equation}
The constants are uniform in the two meshes.
\end{proposition}

\begin{proof}
Let \(c_p\) be the plaquette center. Write the four edge
values of \(A\) as their affine Taylor polynomials at \(c_p\)
plus errors bounded by \(C_\beta h^{1+\beta}\).
The signed sum of the four link exponents is then
\[
 ih^2(\partial_\mu A_\nu-\partial_\nu A_\mu)(c_p)
                                  +O(h^{2+\beta}).
\]
The quadratic Baker--Campbell--Hausdorff term for constant
edge values is \(-h^2[A_\mu,A_\nu]\). Replacing those
constant values by the edge Taylor values changes that
term by \(O(h^3)\). The cubic and higher terms are also
\(O(h^3)\), uniformly on the fixed connection ball.
All four-factor logarithms are on the small branch.
Consequently, with the actual plaquette logarithm defined
by \(U_h(A)(p)=\exp(ih^2F_{h,p})\),
\begin{equation}\label{wmc:plaquette-consistency}
                         |F_{h,p}-F(A)(c_p)|\le C_\beta h^\beta.
\end{equation}
This argument uses one H\"older derivative, not a bounded
second derivative of \(A\).

The exact Pauli trace gives
\[
 1-\tfrac12\operatorname{tr}e^{ih^2F_{h,p}}
        =\tfrac12h^4|F_{h,p}|^2+O(h^8|F_{h,p}|^4).
\]
The plaquette fields are uniformly bounded.
Equation \eqref{wmc:plaquette-consistency} changes this
value by at most \(C_\beta h^{4+\beta}\) when
\(F_{h,p}\) is replaced by \(F(A)(c_p)\).
There are \(6M^4h^{-4}\) plaquettes. Finally
\(|F(A)|^2\) has a uniform \(C^\beta\) modulus, so its
midpoint Riemann sum has error at most
\(C_\beta h^\beta M^4\). These estimates prove
\eqref{wmc:energy-consistency-bound}. Apply it on both
meshes and subtract to obtain
\eqref{wmc:two-mesh-energy}.
\end{proof}

\subsubsection{The first two actual mean vertices}

Consider one step from fine spacing \(H/L\) to spacing
\(H\), with \(0<H\le1\). At the identity write its output
raw bond logarithm as
\[
       m_L(a)=m_{L,1}a+m_{L,2}[a,a]+R_{L,3}(a),\qquad
                         |R_{L,3}(a)|\le C_L\|a\|_\infty^3.
\]
The formula is on the finite double-block carrier of an
output row. Its coefficients and the remainder bounds
are the actual analytic coefficients of the symmetric
contour mean. The finite dimension depends on \(L\);
there is no fine-depth dependence in this one-step
statement.

\begin{lemma}[Midpoint affine reproduction and quadratic cancellation]
\label{wmc:mean-low-vertices}
The actual linear mean applied to midpoint samples of
an affine connection is its coarse midpoint value
multiplied by \(H\). The quadratic map \(m_{L,2}\)
vanishes on every constant connection.
\end{lemma}

\begin{proof}
At the identity, the derivative of a group mean is the
arithmetic mean of its Lie-algebra inputs. Let
\(T_{y,x}\) be the internal transport from the coarse
node \(y\) to a fine starting point \(x\) in its centered
block, and let \(t_{y,x}\) be its first logarithmic
variation. The positively oriented row has linear
variation
\begin{equation}\label{wmc:actual-plain-linear-row}
 \frac1{L^4}\sum_{x\in B(y)}
 \left\{t_{y,x}
       +\sum_{\ell=0}^{L-1}a_\mu(x+\ell H e_\mu/L)
       -t_{y+He_\mu,x+He_\mu}\right\}.
\end{equation}
The entries \(a_\mu\) here are fine link angles, sampled
at their fine bond midpoints.

For an affine continuous connection, midpoint sums along
each coordinate path integrate that affine connection
exactly. Translating such a path by \(He_\mu\) changes
its integral by
\[
       H\sum_\nu(\partial_\mu A_\nu)(x_\nu-y_\nu),
\]
independently of the order of its coordinate edges.
The mean of \(x-y\) over the centered odd block is zero.
Thus the two internal transport contributions in
\eqref{wmc:actual-plain-linear-row} cancel after averaging.
The remaining straight line samples have mean position
\(y+\frac H2e_\mu\): the starting block is centered at
\(y\), and the average of its \(L\) fine bond midpoints
is the midpoint of the coarse edge. They therefore give
exactly \(HA_\mu(y+\frac H2e_\mu)\). This proves the
linear assertion, including the plain transport factors.

For the second assertion set \(a_\mu\) to a constant
matrix in each direction. The input and output
connections are then translation invariant. Central
inversion of the lattice takes these constant fine
links to their inverses. The actual symmetric contour
mean commutes with this inversion, including reversal
of the coarse bond. Its admitted output logarithm
therefore satisfies
\[
                            m_L(-a)=-m_L(a).
\]
Its quadratic coefficient on a constant connection is
zero. Polarization gives zero for two different
constant connections as well. This is an exact
symmetry of the actual mean, not a continuum
replacement of its quadratic vertex.
\end{proof}

\subsubsection{One-step consistency for the nonlinear mean}

\begin{proposition}\label{wmc:one-step-consistency}
For a continuous connection \(A\) in a fixed small
\(C^{1,\beta}\) ball, one actual averaging step obeys
\begin{equation}\label{wmc:one-step-bound}
 \left\|
 \frac1i\log\mathcal M(U_{H/L}(A))
                  -H\mathcal S_HA\right\|_\infty
                                      \le C_\beta H^{2+\beta}.
\end{equation}
The logarithm is the small raw logarithm at the
identity. The estimate is local on each actual
double block and uniform in the volume.
\end{proposition}

\begin{proof}
Fix a coarse bond. Expand the continuous connection
to first order at its midpoint. Over its double
block, the remainder in \(A\) is \(O(H^{1+\beta})\);
the remainder in a fine link angle is consequently
\(O(H^{2+\beta})\), with \(L\) fixed.
By \cref{wmc:mean-low-vertices}, the linear vertex
applied to the affine part gives precisely the
coarse midpoint angle. It has a bounded finite
row sum, so its Taylor remainder has the asserted
size.

Split the fine angles as \(a=a_{\rm const}+a_{\rm var}\),
where \(\|a_{\rm const}\|\le C_LH\|A\|_\infty\) and
\(\|a_{\rm var}\|\le C_LH^2\|\nabla A\|_\infty\).
The constant--constant quadratic term is zero.
All other quadratic terms have size \(O(H^3)\).
The cubic analytic remainder has size \(O(H^3)\)
as well. Since \(H^3\le H^{2+\beta}\) for
\(0<H\le1\), these estimates prove
\eqref{wmc:one-step-bound}. The fixed small
connection ball keeps every path mean and raw
logarithm on its admitted branch.
\end{proof}

\subsubsection{Uniform propagation through the remaining levels}

The plain derivative of one step contains internal
transport terms. Iterating a crude bound on that
derivative could lose a factor at each level. We
instead use the exact node factorization to control
the full plain mean in fine physical coordinates.

\begin{lemma}\label{wmc:full-plain-Lipschitz}
There are \(r_0,C>0\), independent of \(h=L^{-j}\)
and volume, such that
\[
 \mathscr P_h(a)=\frac1i\log\mathcal M^j(e^{iha})
\]
is holomorphic on \(\|a\|_\infty<r_0\) and
\begin{equation}\label{wmc:full-plain-Lipschitz-bound}
 \|\mathscr P_h(a)-\mathscr P_h(\widetilde a)\|_\infty
                             \le C\|a-\widetilde a\|_\infty
\end{equation}
on its half ball. This statement uses no derivative
bound for the fine potentials.
\end{lemma}

\begin{proof}
Use the actual recursion and node factors in
\eqref{hb:actual-recursion}. Its first-exit
estimate uses only \(\|a\|_\infty\), and gives
\[
 \|a_s\|_\infty\le C\xi_s\|a\|_\infty,\qquad
 \|v_s-1\|_\infty\le C\xi_s\|a\|_\infty,
                 \quad \xi_s=hL^s.
\]
The geometric sum of the intermediate \(\xi_s\)'s
bounds the second expression. At the final level,
\eqref{hb:actual-plain-factorization} writes the
plain mean as a product of \(e^{iq_h(a)}\) and
the two node factors at its endpoints. Both
have norm and distance from the identity bounded
uniformly on this fine supremum ball.
Reduce \(r_0\) so their output has its small
logarithm. This proves a uniform bound for the
holomorphic map \(\mathscr P_h\) on a slightly
larger fixed ball.

For a fine direction \(b\), apply the one-variable
Banach-valued Cauchy estimate to
\(\mathscr P_h(a+zb)\) on that larger ball.
It gives a dimension-independent bound
\(\|D\mathscr P_h(a)b\|_\infty\le C\|b\|_\infty\)
on the half ball. Integrate it along the segment
from \(a\) to \(\widetilde a\) to obtain
\eqref{wmc:full-plain-Lipschitz-bound}.
\end{proof}

\begin{proposition}[Consistency of the full actual coarse mean]
\label{wmc:full-mean-consistency}
For a continuous \(C^{1,\beta}\) connection in
a fixed sufficiently small ball,
\begin{equation}\label{wmc:full-mean-bound}
 \left\|
 \frac1i\log\mathcal M^{j+1}(U_{h/L}(A))
          -\frac1i\log\mathcal M^j(U_h(A))
 \right\|_\infty
                                  \le C_\beta h^{1+\beta},
                   \qquad h=L^{-j}.
\end{equation}
The two coarse outputs have spacing one, and both
logarithms are raw logarithms near the trivial
coarse field.
\end{proposition}

\begin{proof}
After the first step from \(h/L\) to \(h\), write
its actual output as \(e^{ih\widetilde a_h}\).
Equation \eqref{wmc:one-step-bound} gives
\[
             \|\widetilde a_h-\mathcal S_hA\|_\infty
                                      \le C_\beta h^{1+\beta}.
\]
Both potentials lie in the half ball of
\cref{wmc:full-plain-Lipschitz} after the initial
continuous ball is reduced once. The remaining
\(j\) steps are exactly the same actual mean
on the two fine-\(h\) fields. Applying
\eqref{wmc:full-plain-Lipschitz-bound} proves
\eqref{wmc:full-mean-bound}. For \(j=0\) the
remaining mean is the identity and the argument
is unchanged.
\end{proof}

Together, \cref{wmc:interpolation,wmc:energy-consistency,wmc:full-mean-consistency}
give actual real trial
connections for comparing successive effective
actions. They require only the first H\"older
derivative supplied by the homogeneous background
construction. The exact repair of their small
coarse-mean error, the variational comparison,
and the passage from real action differences to
localized Taylor kernels are the next steps.

% END LAYERED EXPANSION classical_mesh_consistency.tex

% BEGIN LAYERED EXPANSION background/classical_action_convergence.tex
\subsection{Geometric convergence of the actual classical Taylor vertices}
\label{ch:cac-classical-convergence}

The comparison in this section concerns the actual minimized Wilson
action, rather than a postulated continuum minimizer. Interpolation
is used only to transfer a real trial field between two adjacent
meshes. Analyticity then transfers a real energy comparison to the
second, third and fourth Taylor vertices. This last passage needs a
real-to-complex argument: Cauchy's inequality alone cannot turn a
small bound on a real cube into a small complex derivative.

Work near the trivial flat reference, at physical coarse spacing
one. Put $h_j=L^{-j}$, let $\widehat U_j(B)$ be the actual literal
background of \cref{prop:hb-actual-critical-orbit}, and write
\begin{equation}
 S_j(B)=A_{h_j}(\widehat U_j(B)),\qquad
                    \mathcal V_{j,m}=D^mS_j(0),\quad m=2,3,4.
 \label{cac:actual-actions}
\end{equation}
The coarse torus has side $M\ge2$ and physical volume
$\mathcal V=M^4$. Its real coarse field coordinates are traceless
Hermitian bond logarithms. A complex ball always means their
holomorphic continuation. Fix an exponent $0<\beta<1$ and use
the preceding background construction with $X_\beta$.
Its coarse radius may be decreased by a fixed amount below.

We use the proved mesh-consistency facts in
\cref{wmc:interpolation,wmc:energy-consistency,wmc:one-step-consistency,wmc:full-mean-consistency}:
a periodic fine potential in $X_\beta$ admits a
periodic $C^{1,\beta}$ interpolation with the same midpoint values
and a uniform norm bound; its sampled Wilson energy differs from
$\tfrac12\int\sum_{\mu<\nu}|F_{\mu\nu}|^2$ by
$C h^\beta\mathcal V$; and the actual one-step mean applied to
the refined midpoint samples differs from the coarser midpoint
link by $C h^{2+\beta}$ in raw relative logarithm. All these
constants are uniform in $M,j$ on the fixed small norm ball.
These are consistency estimates for the actual mean and Wilson
action, not assumptions about a continuum critical field.

\subsubsection{Repairing a transferred trial field without losing derivatives}

\begin{lemma}[An exact repair of a small coarse mean error]
\label{lem:cac-mean-repair}
Let $U=e^{ihA}$ be a real fine trial with small $X_\beta$ norm,
and suppose its literal mean is $e^{i(B+e)}$, where
$\|e\|_\infty$ is small. There is a literal trial $U^{\mathrm{rep}}$
with mean $e^{iB}$ and
\begin{equation}
 |A_h(U^{\mathrm{rep}})-A_h(U)|
                                  \le C\|e\|_\infty\mathcal V.
 \label{cac:repair-energy}
\end{equation}
The trial and its repair stay in a fixed small strong chart.
All constants are uniform in mesh and volume.
\end{lemma}
\begin{proof}
Apply the actual mean-gauge slice map
\eqref{hb:mean-slice-solution} to the literal trial. Its result
$\overline A$ has $q_h(\overline A)=B+e$ and
$RD^\sharp\overline A=0$, with
$\|\overline A\|_{X_\beta}\le C\|A\|_{X_\beta}$.
It has exactly the same action as $U$.

Use the fixed smooth exact field seed $E_Q$ and solve
\begin{equation}
 z=-e-\{C_h(\overline A+E_Qz)-C_h(\overline A)\},
                         \qquad C_h=q_h-Q.
 \label{cac:repair-equation}
\end{equation}
The actual nonlinear row Lipschitz estimate makes its right side
a contraction on a ball of radius $2\|e\|_\infty$ when the
strong chart is fixed sufficiently small. It gives
$q_h(\overline A+E_Qz)=B$ and
$\|E_Qz\|_{Y_\beta}\le C\|e\|_\infty$.
The specified node-gauge extension in
\eqref{hb:literal-representative} converts this into a literal
field $U^{\mathrm{rep}}$ with the desired mean, preserving action.

To estimate the change, integrate the first variation along
$\overline A+tE_Qz$. Its curvature is uniformly bounded on the
small strong chart. Equation \eqref{hb:first-variation-readout}
uses only that curvature bound and the physical $W^{1,1}$ norm
of $E_Qz$, not a second derivative of $\overline A$.
The smooth seeds have bounded overlap, amplitude and first
differences, so
$\|E_Qz\|_{1,h}+\|DE_Qz\|_{1,h}
 \le C\mathcal V\|z\|_\infty$.
The normal-to-left differential has the same estimate, by its
explicit bounded coefficients and their first differences.
This proves \eqref{cac:repair-energy}. The actual gauge and seed
bounds also keep the repaired field inside the strong chart.
\end{proof}

The correction $E_Qz$ need not be horizontal. When comparing its
energy with the minimum, apply the literal-to-mean slice once more.
It preserves the action and puts the trial in $q_h=B$ and
$RD^\sharp A=0$, where the proved chord minimum applies.

\begin{proposition}[Real adjacent-mesh action comparison]
\label{prop:cac-real-action-comparison}
On a fixed sufficiently small real coarse ball,
\begin{equation}
 \sup_{\|B\|_\infty<\rho}
       \frac{|S_{j+1}(B)-S_j(B)|}{\mathcal V}
                                            \le C_\beta h_j^\beta.
 \label{cac:real-action-comparison}
\end{equation}
The constants are independent of $M,j$.
\end{proposition}
\begin{proof}
Interpolate the literal real background potential at spacing
$h=h_j$ by the proved periodic $C^{1,\beta}$ interpolation,
and sample that same potential at spacing $h/L$. The refined
trial has a uniformly small $X_\beta$ norm. The two Wilson
energies differ by at most $C h^\beta\mathcal V$, since each
differs from the energy of this common continuous interpolation
by that bound.

The one-step mean of the refined trial differs from the original
spacing-$h$ field by $C h^{2+\beta}$ in raw link logarithm.
In physical spacing-$h$ potential coordinates this is
$C h^{1+\beta}$. Apply
\cref{wmc:full-plain-Lipschitz} to the remaining $j$ actual
mean steps. Their final logarithms therefore differ by at most
$C h^{1+\beta}$. Thus the refined trial has literal final mean
$e^{i(B+e)}$ with $\|e\|_\infty\le C h^{1+\beta}$.
Repair it by \cref{lem:cac-mean-repair} and compare its action
with the actual spacing-$h/L$ minimum. This gives
$S_{j+1}(B)\le S_j(B)+C h^\beta\mathcal V$.

For the reverse inequality interpolate the spacing-$h/L$
background, and sample its interpolation on the spacing-$h$
grid. One-step consistency compares its actual fine first mean
with that coarse sample. The remaining $j$ means again produce
an error $C h^{1+\beta}$. The same energy and repair estimates
give $S_j(B)\le S_{j+1}(B)+C h^\beta\mathcal V$.

All transferred potentials have norm at most $C_\beta\rho$.
The plain-mean Lipschitz bound at zero also gives
$\|e\|_\infty\le C_\beta\rho$, in addition to its mesh-dependent
bound. This ensures uniform admission even for the initial coarse
meshes, where $h^{1+\beta}$ need not be small by itself.
Choose the real ball once so this bound, and its slice and repair
images, lie inside the fixed minimum chart. The local minimum is
then applicable in both directions; no global minimizer or
continuum critical field is invoked.
\end{proof}

The uniform complex bound needed below is simpler. For the actual
holomorphic background, $\|F\|_\infty\le C\rho$ and a plaquette
action is at most $Ch^4\rho^2$ in absolute value. There are
$O(h^{-4}\mathcal V)$ plaquettes. Consequently on a fixed
complex coarse ball,
\begin{equation}
                |S_j(B)|/\mathcal V\le C,\qquad j\ge0.
 \label{cac:uniform-complex-action-bound}
\end{equation}

\subsubsection{Small real values imply small fixed-order Taylor jets}

\begin{lemma}[A quantitative real-to-jet estimate]
\label{lem:cac-Chebyshev}
Fix $m\ge1$ and a complex polydisc. Suppose $f$ is holomorphic
there with $|f|\le C$, and $|f|\le\varepsilon$ on a fixed smaller
real cube centered at zero. Then
\begin{equation}
 |\partial_1\cdots\partial_m f(0)|
       \le C_m\varepsilon(1+\log(1/\varepsilon))^{2m},
                              \qquad0<\varepsilon<\tfrac12.
 \label{cac:real-to-jet-bound}
\end{equation}
The constant depends only on the two radii, $m$ and $C$.
\end{lemma}
\begin{proof}
Rescale the real cube to $[-1,1]^m$ and choose $R>1$ so that
the product of closed Bernstein ellipses of parameter $R$ lies
inside the holomorphy domain. Expand in product Chebyshev
polynomials: $f(t)=\sum_{n\in\mathbb N^m}c_n\prod_iT_{n_i}(t_i)$.
Real orthogonality gives $|c_n|\le2^m\varepsilon$.
Analyticity on the ellipses gives $|c_n|\le C_mR^{-|n|_1}$:
write $t_i=(w_i+w_i^{-1})/2$ and apply the Laurent coefficient
estimate on $|w_i|=R$ in each variable.

The derivative at zero has weight
$\prod_i|T_{n_i}'(0)|\le\prod_i n_i\le(1+|n|_1)^m$.
For $|n|_1\le N$, the real bound therefore gives at most
$C_m\varepsilon(1+N)^{2m}$ after counting the multi-indices.
The tail is at most
$C_mR^{-N}(1+N)^{2m-1}$, by summing the geometric series with
this fixed polynomial weight. Choose $N$ to be the smallest
integer at least $\log(1/\varepsilon)/\log R$.
Restore the fixed rescaling of the cube to prove the lemma.
The differentiated series is absolutely convergent. No Cauchy
estimate has been applied to a real-only supremum bound.
\end{proof}

Take real coarse profiles $b_1,\ldots,b_m$ with supremum norm
at most one and apply the lemma to
\[
 f(z)=\mathcal V^{-1}(S_{j+1}-S_j)
                              \left(\sum_{i=1}^mz_ib_i\right).
\]
The complex polydisc and real cube can be chosen independently
of profiles, depth and volume, by taking $m\max_i|z_i|$ smaller
than the fixed action ball. Equations
\eqref{cac:real-action-comparison} and
\eqref{cac:uniform-complex-action-bound} give, for $2\le m\le4$,
\begin{equation}
 \frac{|D^m(S_{j+1}-S_j)(0)[b_1,\ldots,b_m]|}{\mathcal V}
              \le C_m h_j^\beta(1+j)^{2m}.
 \label{cac:profile-Taylor-comparison}
\end{equation}
The finitely many initial $j$ for which the smallness parameter
exceeds $1/2$ are covered by the uniform complex Cauchy bound.
Complex Fourier profiles are permitted as well: expand each
$e^{ipx}$ into its real cosine and sine profiles and use complex
multilinearity, at a fixed cost at most $2^m$.

\subsubsection{Uniform analytic strips for the actual action vertices}

We spell out the locality assertion needed to pass from a small
symbol to a small kernel. Resolve a coarse bond into its cell
location and one of its twelve component-colour indices. In the
infinite-volume translation-covariant notation write the $m$-vertex
as $V_{j,m}(x_1,\ldots,x_m)$, retaining these finite indices.
Fix the first location at zero and define its star-weighted norm by
\begin{equation}
 \|V\|_{m,\sigma}
  =\max_{\text{component-colour indices}}
       \sum_{x_2,\ldots,x_m}
        e^{\sigma\sum_{i=2}^m|x_i|_1}|V(0,x_2,\ldots,x_m)|.
 \label{cac:weighted-vertex-norm}
\end{equation}
Here infinite-volume notation refers to kernels of compactly
supported variations, obtained from their local periodic limits.
No finite total action on a nondecaying infinite field is assigned.
One may sum the finite index maximum instead, changing only a fixed
constant. On a torus use periodic distance. This norm also controls
any weight given by the length of a shortest tree connecting the
inputs, since that length is no larger than the displayed star.

\begin{lemma}[Uniform exponential moments at fixed Taylor order]
\label{lem:cac-vertex-strip}
For $m=2,3,4$ there are $\sigma_0>0$ and $C_m$, independent of
$j,M$, such that
\begin{equation}
                         \|V_{j,m}\|_{m,\sigma_0}\le C_m.
 \label{cac:uniform-vertex-moment}
\end{equation}
The infinite-volume kernels exist, and finite-torus kernels are
their periodic sums in the $m-1$ relative input variables.
Their symbols are consequently bounded holomorphic functions
on a fixed polystrip in $4(m-1)$ complex momentum variables.
\end{lemma}
\begin{proof}
Split the Wilson action into unit-cell densities, assigning a
plaquette by its base point. Each density is uniformly analytic
in a small local $X_\beta$ ball and has bounded analytic norm:
the fine plaquette factor is $O(h^4\|A\|_{X_\beta}^2)$, and
only $O(h^{-4})$ plaquettes lie in one physical cell.
Its derivatives of every fixed order are therefore bounded
multilinear local maps, uniformly in $h$. Their carriers have
bounded physical diameter.

For a background jet with unit inputs at $y_1,\ldots,y_a$,
\eqref{hb:localized-background-jets} bounds its cell-$X_\beta$
norm at $z$ by $C_a e^{-s\max_i|z-y_i|_1}$. This holds also
for the literal background, as proved there. In the chain rule
for an order-$m$ cell-action jet, each partition of the $m$
inputs contributes a product of such bounds. For each block of
size $a$, its maximum is at least $a^{-1}$ times its sum.
Since $a\le m$, every term is bounded by
$C_m e^{-(s/m)\sum_i|z-y_i|_1}$, with a fixed change of the
constant for the bounded action carrier. Summing action cells gives
\begin{equation}
 |V_{j,m}(y_1,\ldots,y_m)|
       \le C_m\sum_z e^{-(s/m)\sum_i|z-y_i|_1}.
 \label{cac:vertex-cell-bound}
\end{equation}
With $y_1=0$, the triangle inequality implies
$\sum_{i=2}^m|y_i|_1\le(m-1)\sum_i|z-y_i|_1$.
Split the exponent into two halves. One half pays a star weight
of any size below $s/(2m(m-1))$; the other is summable in $z$
and all relative inputs after a further fixed decrease of that
weight. For example choose $\sigma_0<s/(4m(m-1))$ and sum
first in each $y_i$ and then in $z$. This proves the moment
bound, uniformly on periodic tori as well.

We justify the infinite-volume and periodization assertions at
the level of these actual jets. At zero field the joint inverse
has its translation-covariant Bloch representation and exponential
cell kernels. Its finite periodic kernel is the periodic sum of
the infinite one. Differentiate the actual fixed-point equation
at order $a$. Its order-$a$ unknown is the same linear joint
inverse applied to a finite sum of products of lower-order jets
and the finite-carrier derivatives of the actual residual.
Induction for $a\le m$ shows that each such jet is an absolutely
convergent rooted-tree kernel expression. Its leaves are the
distinct external differentiation slots. Once a lift of an
internal vertex is chosen, the images of the input slots in each
child subtree are independent. Expanding a periodic propagator
chooses the lift of its child vertex; continuing down the tree
then gives the independent periodic images of all external slots.
There is no closed internal cycle imposing an extra image
constraint. Thus periodic summation commutes with these
multilinear contractions; this is not the generally false
identity that periodization commutes with an ordinary pointwise
product in one shared argument. The exponential bounds justify
all changes of summation order.

Apply the same finite-order chain rule to the cell action and
sum its cells. Its action root joins those child trees, still
without a closed cycle. Fixing the first external input lift
removes the common translation; the remaining $m-1$ external
images are independent. This proves existence of the
infinite-volume action vertex and exact periodization in each
relative input, not just a subsequential compactness claim.

Finally Fourier series are absolutely convergent when the imaginary
parts of the $m-1$ independent external momenta are sufficiently
small. The moment bound gives a uniform bound on a smaller closed
polystrip. Cell positions rather than bond midpoints can be used
in this Fourier convention; changing that convention introduces
only fixed bounded analytic phases.
\end{proof}

Translation covariance now turns
\eqref{cac:profile-Taylor-comparison} into a pointwise bound for
the difference of vertex symbols. Choose complex unit component-
colour profiles with allowed momenta $p_1,\ldots,p_m$ whose sum
is zero. The volume-normalized $m$-linear functional is exactly
its vertex symbol at those momenta. It follows that
\begin{equation}
 \sup_{p_2,\ldots,p_m\text{ real}}
 |\widehat V_{j+1,m}(p_2,\ldots,p_m)
                 -\widehat V_{j,m}(p_2,\ldots,p_m)|
                       \le C_mh_j^\beta(1+j)^{2m}.
 \label{cac:real-symbol-comparison}
\end{equation}
Initially the momenta are on a finite-torus grid. Pass to arbitrarily
large tori using the just proved periodization, then use continuity
and density of the grids. This establishes the bound for every
real momentum, with no volume-dependent Fourier inversion loss.

\subsubsection{Interpolation into a smaller complex momentum strip}

We need an elementary several-variable refinement of the three-lines
principle. Let $F$ be periodic and holomorphic in
$|\operatorname{Im}p_i|<a_0$, $1\le i\le d$, and bounded by $C$
there. Suppose $|F|\le\varepsilon$ on the real torus. Choose
$0<a<a_0'<a_0$. Apply the three-lines bound in one variable,
keeping the remaining variables real, between its real axis and
the line of imaginary part $a_0'$ or $-a_0'$. This gives
$C^{1-\theta}\varepsilon^\theta$ for $|\operatorname{Im}p_1|\le a$,
where $\theta=1-a/a_0'$. Periodicity permits the argument on a
bounded cylinder, or equivalently by the annulus maximum principle.
Keeping this first variable in its reduced strip, repeat in the
second variable. Its real-axis bound is now the improved preceding
bound, whereas its outer-line bound is still $C$.
After all $d$ variables,
\begin{equation}
 \sup_{|\operatorname{Im}p_i|\le a}|F(p)|
       \le C^{1-\Theta}\varepsilon^\Theta,
                       \qquad\Theta=(1-a/a_0')^d.
 \label{cac:iterated-three-lines}
\end{equation}
This argument uses the uniform bound on the full larger polystrip
at each step. Smallness on the real torus alone is not assumed
to imply the same smallness throughout that larger strip.

Apply \eqref{cac:iterated-three-lines} to the vertex-symbol
difference. The larger-strip bound is supplied by
\cref{lem:cac-vertex-strip}; the real bound is
\eqref{cac:real-symbol-comparison}. Since $d=4(m-1)$,
\begin{equation}
 \sup_{|\operatorname{Im}p_i|\le a}|\Delta\widehat V_{j,m}(p)|
     \le C_m\{h_j^\beta(1+j)^{2m}\}^{(1-a/a_0')^{4(m-1)}}.
 \label{cac:reduced-strip-comparison}
\end{equation}
Again the finitely many initial indices are absorbed into $C_m$.
Fourier contour shifting on this reduced polystrip gives
\begin{equation}
 |\Delta V_{j,m}(0,x_2,\ldots,x_m)|
  \le C_m\{h_j^\beta(1+j)^{2m}\}^{\Theta_m}
                       e^{-a\sum_{i=2}^m|x_i|_1},
 \quad\Theta_m=(1-a/a_0')^{4(m-1)}.
 \label{cac:kernel-coefficient-decay}
\end{equation}
Shift each scalar momentum coordinate with the sign prescribed by
the corresponding relative lattice coordinate. Periodicity makes
the opposite contour faces cancel. The fixed analytic midpoint
phases, if used, change only $C_m$.

\begin{proposition}[Convergence of the actual classical Taylor triple]
\label{prop:cac-Taylor-triple-convergence}
For every $0<\alpha<1$ there are $\sigma_\alpha>0$ and $C_\alpha$
such that, for $m=2,3,4$,
\begin{equation}
 \|V_{j+1,m}-V_{j,m}\|_{m,\sigma_\alpha}
                                 \le C_\alpha L^{-\alpha j}.
 \label{cac:Taylor-triple-step}
\end{equation}
The limits $V_{\infty,m}$ exist in these weighted kernel norms and
\begin{equation}
 \|V_{\infty,m}-V_{j,m}\|_{m,\sigma_\alpha}
                                 \le C_\alpha L^{-\alpha j}.
 \label{cac:Taylor-triple-limit}
\end{equation}
The bounds hold uniformly on compatible finite tori by periodic
summation. They concern the actual Wilson background Taylor
vertices; no continuum minimizing field has been assumed.
\end{proposition}
\begin{proof}
Choose $\alpha<\beta<1$, and then choose the reduced strip $a>0$
so small that
$\beta(1-a/a_0')^{12}>\alpha$.
This gives $\beta\Theta_m>\alpha$ for every $m=2,3,4$.
Take $\sigma_\alpha=a/2$. Summing
\eqref{cac:kernel-coefficient-decay} with this weight costs only
the convergent geometric lattice sum
$\sum_{x_2,\ldots,x_m}e^{-(a/2)\sum|x_i|_1}$.
Its remaining polynomial in $j$ is absorbed by the strict
exponent margin $\beta\Theta_m-\alpha>0$. This proves
\eqref{cac:Taylor-triple-step}. The geometric series in $j$
converges in the Banach kernel norm, proving
\eqref{cac:Taylor-triple-limit}.

The periodization proved above preserves the weighted bound:
periodic relative distances are no larger than their lifted
distances, and the absolute value of a periodic sum is bounded
by the sum of absolute values. Thus no finite-volume factor
enters the result.
\end{proof}

The classical Taylor triple now has a genuine geometric mesh
comparison. Passing this estimate through a Gaussian polarization
requires a uniform Lipschitz estimate for that polarization in the
same kernel norms. That is a separate analytic step. The cumulative
one-loop comparison alone did not prove it, and no convergence of
individual Gaussian increments is asserted here.

% END LAYERED EXPANSION background/classical_action_convergence.tex

% BEGIN LAYERED EXPANSION inverse/gaussian_polarization_lipschitz.tex
\subsection{The actual one-step Gaussian polarization}
\label{ch:gaussian-polarization-lipschitz}

The Gaussian coefficient of a step depends only on the second,
third and fourth Taylor forms of its input effective action.
Here the constraint is the actual plain group average and the
gauge is the actual symmetric contour gauge. We derive this
dependence, including the Haar and coarea factors, and prove a
volume-independent Lipschitz bound. The resulting implication
from convergence of the actual Taylor triples to convergence
of the actual increments is quantitative.

Fix an odd step size \(L\). The input lattice has unit spacing;
it is not the original fine lattice of spacing \(L^{-k}\).
There are finitely many input coordinates, depending only on
\(L\), in one outer coarse cell. Use orthonormal Pauli colour
coordinates and product Lebesgue measures for group logarithms.
Write
\[
 S(B)=S_k(e^{iB}),\qquad m(B)=i^{-1}\log\mathcal M(e^{iB}),
             \qquad m(0)=0,\quad Q=Dm(0).
\]
The background for all Taylor forms in this section is the
trivial field. Let \(\ell_x(B)\) be the actual contour
coordinate from the node of a block to its nonnode site
\(x\), and put
\begin{equation}\label{gpl:actual-gauge}
 \begin{split}
 F_g(B)&=\sum_x\left(1-\tfrac12\operatorname{tr}T_x(e^{iB})\right)\\
       &=\sum_x\left(1-\cos|\ell_x(B)|\right).
 \end{split}
\end{equation}
This is the canonical \(\tau=1\) Gaussian gauge convention.
The proof also permits fixed gauge weights in a fixed
positive compact interval. Coupling powers and
field-independent reference-gauge constants have been
factored out; the two-coupling mismatch is a separate term.
On complex coordinates the trace formula is its analytic
continuation; the cosine expression is its power series
in the bilinear colour square of \(\ell_x\).

\subsubsection{The actual free effective action supplies the gap}

We first record why the free input of every homogeneous
depth has the lower bound used below. In the physical
fine spacing \(h=1/n\), the actual straight derivative is
\[
 (Q_ha)_\mu(y)=h^5\sum_{x\in B_n(y)}
                              \sum_{r=0}^{n-1}a_\mu(x+rh e_\mu).
\]
Telescoping the two coarse differences gives
\begin{equation}\label{gpl:actual-flux-contraction}
 (\mathrm{d}_cQ_ha)_{\mu\nu}(y)
 =h^6\sum_{x\in B_n(y)}\sum_{r,s=0}^{n-1}
             (\mathrm{d}_ha)_{\mu\nu}(x+rh e_\mu+sh e_\nu)
                         =:(T_hd_ha)_{\mu\nu}(y).
\end{equation}
These weights are positive. A row sum is \(h^6n^6=1\).
A fixed fine plaquette has \(n^2\) choices of \((r,s)\),
each with exactly one starting block, so its column sum
is \(h^6n^2=h^4\). Jensen's inequality therefore proves
\(\|T_hF\|_2\le\|F\|_{2,h}\), with constant one.

The plain derivative differs from this straight derivative
by the actual coarse gradient of the linear node factor.
Their curls agree. The physical derivative of the actual
background minimizes fine curl energy at its fixed mean:
its stationarity gives zero energy cross term against every
zero-mean variation. A literal fine node gauge preserves
this energy. The exact smooth row lift of
\cref{lem:hb-row-seeds}, followed by the
specified node gauge when passing to the plain mean, has
uniform physical gradient energy. Consequently
\begin{equation}\label{gpl:actual-effective-free-bounds}
             \mathrm{d}^\sharp \mathrm{d}\ \le T_2:=S_k''(0)\ \le C I .
\end{equation}
The lower bound is the actual flux contraction applied
to the minimizing lift. The upper bound inserts the exact
smooth trial. It is uniform in the inner homogeneous depth,
including depth zero. No continuum coercivity estimate is
substituted for these actual averaging identities.

\begin{lemma}[The fixed one-step contour gap]
\label{lem:gpl-contour-gap}
Let \(K_g=D^2F_g(0)\). On compatible tori, uniformly in
volume,
\begin{equation}\label{gpl:contour-gap}
           \mathrm{d}^\sharp \mathrm{d}+K_g+Q^\sharp Q\ge c_L I .
\end{equation}
The actual plain row \(Q\) is onto with a bounded right
inverse depending only on \(L\). Thus, whenever
\(\mathrm{d}^\sharp \mathrm{d}\le T_2\) and \(\|T_2\|\le M\),
\begin{equation}\label{gpl:actual-flat-saddle}
       \mathbb L(T_2)=
          \begin{pmatrix}T_2+K_g&Q^\sharp\\Q&0\end{pmatrix}
\end{equation}
has a uniformly bounded inverse. Only positivity on
\(\ker Q\) is required.
\end{lemma}

\begin{proof}
The actual pivot chart supplies the onto row and its
bounded local pivot right inverse at the trivial field.
All its coefficients and inverse pivot coefficients
belong to one fixed \(L\)-block geometry.
For the gap use the Bloch decomposition by these blocks.
The finite matrices of \(d,K_g,Q\) are continuous on
the compact Bloch torus.

Suppose the three squared terms vanish in one fibre.
A closed fine field is a gradient, except for constant
vector modes at zero Bloch frequency. This follows
also from the fine Fourier modes
\(p=(\theta+2\pi r)/L\): at every nonzero mode the
closedness equation says its vector is parallel to
the fine gradient symbol. At zero frequency write
the field as \(D\lambda+c\). The actual plain average
of its gradient is the coarse gradient of the node
values, and its average of \(c\) is \(Lc\).
In that zero-frequency block the former gradient is
zero, so \(Qu=0\) gives \(c=0\).

For the remaining gradient, contour covariance gives
\[
                  D\ell(0)D\lambda(x)=\lambda(y)-\lambda(x).
\]
The contour term therefore forces \(\lambda\) to be
constant inside each block. Its plain row is the
coarse gradient of that block value. Zero coarse
gradient makes the field itself zero, at both zero
and nonzero Bloch frequencies. Every Bloch matrix
has zero kernel, so compactness proves
\eqref{gpl:contour-gap}. In particular \(T_2+K_g\)
is uniformly coercive on \(\ker Q\). The right inverse
and the kernel equation prove the bounded full saddle
inverse, by lifting the datum and recovering the
multiplier from the remaining functional.
\end{proof}

\subsubsection{The global group-coordinate normalizer}

Define the actual Lagrangian in full input group
coordinates and coarse multiplier coordinates by
\begin{equation}\label{gpl:actual-Lagrangian}
 \Lambda_C(B,\upsilon)=S(B)+F_g(B)
                                  +\langle\upsilon,m(B)-C\rangle .
\end{equation}
Its stationary branch is \(z(C)=(B(C),\upsilon(C))\),
with \(z(0)=0\). Its Hessian at zero is
\(\mathbb L(T_2)\). The implicit function theorem
therefore gives the actual local branch; the
already constructed homogeneous background is this
branch after its specified residual contour gauge.
One can see the gauge removal directly: along a
fine gauge equal to one at coarse nodes both
\(S\) and \(m\) are constant. The derivative of
\(\ell_x\) in those gauge directions is an invertible
endpoint map near zero. Stationarity of
\eqref{gpl:actual-Lagrangian} thus implies
\(\ell(B(C))=0\): the derivative of the actual trace
weight in its \(\ell_x\) variable is
\((\sin|\ell_x|/|\ell_x|)\ell_x\), whose scalar
factor is nonzero on the chosen small branch.
Hence \(F_g=DF_g=0\) at the
centre, and this branch is physically constrained
critical for \(S\).
At every such centre its gauge Hessian is
\(D\ell^\sharp D\ell\), also the Hessian of the
quadratic logarithmic gauge. Thus their Gaussian
log normalizers, after subtraction at zero, agree
along this physical branch.
All vertices below use the literal trace function
\eqref{gpl:actual-gauge}.

Let \(j_f(B)\) and \(j_c(C)\) be the products of the
actual \(\operatorname{SU}(2)\) exponential Haar densities.
Factor out the classical weight \(e^{-S(B(C))}\).
The Gaussian log normalizer in these global
coordinates is
\begin{equation}\label{gpl:actual-normalizer}
 \begin{split}
 \Phi(C)={}&\log j_f(B(C))-\log j_c(C)\\
 &-\tfrac12\log\left|
                 \det D_z^2\Lambda_C(z(C))\right|
                    -(\text{the same at }C=0).
 \end{split}
\end{equation}
The omitted Gaussian powers and contour normalization
are independent of \(C\).
At real \(C\) the saddle determinant has a fixed nonzero
sign. Its local analytic logarithm means the logarithm
of its ratio to the value at zero, continued from one.

Here is why this formula is the actual constrained
Gaussian. In global logarithms the exact conditional
measure has density
\[
                  j_f(B)\,j_c(C)^{-1}\delta(m(B)-C)\,dB.
\]
Its tangent constraint is \(Dm(B(C))v=0\).
The quadratic form is the complete Lagrangian
Hessian, including \(\upsilon(C)D^2m\), plus the
actual contour Hessian. The constrained Gaussian
identity \cref{pgg:saddle-coarea} replaces its
quadratic integral and its coarea determinant by
the full saddle determinant in
\eqref{gpl:actual-normalizer}. In particular a
coordinate-dependent coarea factor has not been
discarded. Using left exponential fluctuation
coordinates instead would transfer the centre Haar
factor into this determinant; retaining just one
of these changes would alter the answer.

\subsubsection{All vertices and the colour cancellation}

Put
\begin{equation}\label{gpl:actual-vertex-definitions}
 \mathbb G=\mathbb L(T_2)^{-1},\quad
 ZC=\mathbb G(0,C),\quad HC=(ZC)_B,\quad
              V_r=D_z^r\Lambda_0(0)\quad(r=3,4).
\end{equation}
Differentiating stationarity gives
\[
                    z'(0)C=ZC=(HC,-\Delta_{\rm next}C).
\]
The sign of the multiplier follows from the plus sign
in \eqref{gpl:actual-Lagrangian}. More explicitly,
set \(S_{\rm next}(C)=S(B(C))\).
Its Hessian is \(\Delta_{\rm next}=H^\sharp T_2H\).
Differentiating the zero contour coordinate gives
\(K_gH=0\); pairing the linear stationary equation
with \(HC'\) then gives
\(\upsilon'(0)C=-\Delta_{\rm next}C\).

For clarity, if \(z_i=(a_i,\alpha_i)\), the full vertices
are
\begin{equation}\label{gpl:full-vertices}
 \begin{split}
 V_3[z_1,z_2,z_3]
   &=(D^3S+D^3F_g)[a_1,a_2,a_3]
         +\sum_{i=1}^3\langle\alpha_i,
                                   D^2m[a_j,a_l]\rangle,\\
 V_4[z_1,z_2,z_3,z_4]
   &=(D^4S+D^4F_g)[a_1,a_2,a_3,a_4]
         +\sum_{i=1}^4\langle\alpha_i,
                            D^3m[a_1,\ldots,\widehat a_i,\ldots,a_4]\rangle .
 \end{split}
\end{equation}
In the first sum \(\{j,l\}\) is the complement of
\(\{i\}\). The constraint and gauge terms are fixed
one-step terms; they are present before taking a
Gaussian. There is no multiplier at zero, so no
additional \(D^4m\) term occurs.

The scalar
\[
                 z\longmapsto\log|\det D_z^2\Lambda_0(z)|
\]
is invariant under every constant colour rotation.
Indeed \(S,F_g,m\) have their actual covariance,
and the simultaneous adjoint rotation of all field
and multiplier coordinates is orthogonal with
determinant one. Its differential at zero is an
invariant linear functional on a direct sum of
adjoint colour copies. Each copy has no invariant
linear functional. Therefore, in every field and
multiplier direction \(v\),
\begin{equation}\label{gpl:full-gradient-zero}
                         \operatorname{Tr}(\mathbb G V_3[v])=0.
\end{equation}
This is a finite-volume identity, before any trace
per volume or infinite-volume limit.

\begin{proposition}[Actual Gaussian Hessian formula]
\label{prop:gpl-actual-formula}
The Hessian of \eqref{gpl:actual-normalizer} is
\begin{equation}\label{gpl:actual-Phi-Hessian}
 \begin{split}
 \Phi''(0)[C,C']={}&\tfrac23\langle C,C'\rangle
                           -\tfrac23\langle HC,HC'\rangle\\
 &-\tfrac12\operatorname{Tr}
                     \bigl(\mathbb G V_4[ZC,ZC']\bigr)\\
 &+\tfrac12\operatorname{Tr}
             \bigl(\mathbb G V_3[ZC]\mathbb G V_3[ZC']\bigr).
 \end{split}
\end{equation}
It depends on \(S\) only through
\((D^2S(0),D^3S(0),D^4S(0))\).
\end{proposition}

\begin{proof}
The second derivative of a matrix logarithm is
the trace of its inverse times its second
derivative, minus the trace of the two first
derivatives separated by inverse factors.
Along \(z(C)\), the first matrix derivative is
\(V_3[ZC]\); the second is
\(V_4[ZC,ZC']+V_3[z''(0)[C,C']]\).
The last term has zero trace by
\eqref{gpl:full-gradient-zero}, in its full
augmented direction. Thus no unknown second
background response survives.

In the present Pauli normalization the exact
one-bond Haar density is
\[
 j(B)=\left(\frac{\sin |B|}{|B|}\right)^2,\qquad
                     \log j(B)=-|B|^2/3+O(|B|^4).
\]
Its first derivative is zero and its second
derivative is \(-2/3\) times the colour inner
product. This gives the two Haar terms in
\eqref{gpl:actual-Phi-Hessian}, including their
signs. The complete saddle contains all the
constraint and contour contributions in
\eqref{gpl:full-vertices}. Combining the terms
proves the formula.
\end{proof}

\subsubsection{A norm for the actual Taylor triple}

Assign input bonds, contour coordinates and
coarse multipliers to their outer \(L\)-cells.
The number of coordinates per cell is \(N_L<\infty\).
For a scalar symmetric \(r\)-tensor \(W\), define
\begin{equation}\label{gpl:vertex-norm}
 \|W\|_{s,r}
 =\max_{1\le j\le r}\ \sup_{i_j}
       \sum_{(i_l)_{l\ne j}}
 e^{s\,\mathfrak t(i_1,\ldots,i_r)}
                       |W(i_1,\ldots,i_r)|.
\end{equation}
Here \(\mathfrak t\) is the least cell-tree length
joining the indicated cells, and sums include
bond and colour indices. For \(r=2\) this is
the usual maximum absolute weighted row and
column Schur norm. On a torus the least torus
tree length is used; retained lifted kernels
can instead use their cover tree lengths.

Use the derivative triple
\(\mathfrak J=(T_2,T_3,T_4)\), with \(T_r=D^rS(0)\),
and norm
\begin{equation}\label{gpl:triple-norm}
                    \|\mathfrak J\|_s=\sum_{r=2}^4\|T_r\|_{s,r}.
\end{equation}
If \(S_3,S_4\) denote the homogeneous Taylor
forms, their polarizations have the additional
fixed factors \(3!\) and \(4!\) in \(T_3,T_4\).
This fixes the convention without changing
the asserted Lipschitz property.

Assume \(\|\mathfrak J\|_s\le M\).
The actual action-tree expansion supplies
such norms on its smaller fixed ball.
The fixed one-step constraint and gauge
vertices have finite carriers and bounded
norms of every fixed order. Thus the
augmented \(V_3,V_4\) satisfy the same
kind of bounds, depending only on \(M,L,s\).

\begin{lemma}[Inverse kernels uniformly over the input triples]
\label{lem:gpl-weighted-inverse}
On the class
\(\|\mathfrak J\|_s\le M\), \(T_2\ge \mathrm{d}^\sharp \mathrm{d}\),
there is a fixed \(t>0\) such that
\(\mathbb G\) and \(Z\) have bounded weighted
Schur norms at rate \(t\). For two such triples,
\begin{equation}\label{gpl:inverse-Lipschitz}
 \|\mathbb G-\mathbb G'\|_{t,2}
                  +\|Z-Z'\|_{t,2}\le C\|T_2-T_2'\|_{s,2}.
\end{equation}
Constants are independent of the torus volume
and the inner homogeneous depth.
\end{lemma}

\begin{proof}
The unweighted inverse is uniform by
\cref{lem:gpl-contour-gap}. For a bounded
one-Lipschitz cell function \(\psi\), the
weighted Schur norm of the coefficients
gives
\[
 \|e^{a\psi}\mathbb L(T_2)e^{-a\psi}
                           -\mathbb L(T_2)\|\le C(M,L,s)|a|
\]
for \(|a|\le s/2\). A normalized Neumann
series gives uniformly bounded inverses
for a fixed nonzero \(a\).
Truncated distance weights yield exponential
cell entries of \(\mathbb G\). At a fixed
smaller rate their row and column sums
are bounded: cell growth is polynomial
of degree four, with \(N_L\) internal
coordinates. This also proves the bounds
on \(Z\), a fixed column inclusion of
\(\mathbb G\).

Only the field-field \(T_2\) block differs
between the two full saddle matrices.
The exact resolvent identity is
\[
 \mathbb G-\mathbb G'
       =-\mathbb G
          \begin{pmatrix}T_2-T_2'&0\\0&0\end{pmatrix}
                                                     \mathbb G'.
\]
Weighted Schur norms are submultiplicative.
This proves \eqref{gpl:inverse-Lipschitz}
and its column statement for \(Z\).
\end{proof}

\subsubsection{Rooted trace contractions have no volume cost}

Let \(\Pi(\mathfrak J)\) be the coarse kernel
of \eqref{gpl:actual-Phi-Hessian}.
All statements below concern absolute
coordinate sums; an operator norm alone
would not suffice.

For the quartic trace, fix one external
coarse coordinate \(c\). Sum first its
response column \(Z_{ic}\), then the
quartic vertex rooted at \(i\), then the
second external response row. The internal
loop kernel is bounded by its Schur norm.
The distance from \(c\) to the other
external coordinate is no greater than
the first response distance, the vertex
tree length and the last response distance
combined. Hence, at a common smaller rate,
\begin{equation}\label{gpl:quartic-contraction}
 \left\|\operatorname{Tr}
                   (\mathbb G V_4[Z\,\cdot,Z\,\cdot])\right\|_{t,2}
                \le \|\mathbb G\|_{t,2}\|V_4\|_{t,4}\|Z\|_{t,2}^2 .
\end{equation}
For the pair of cubic vertices, root the
first vertex at that first response index.
Sum one intervening Green kernel, then
the second cubic vertex rooted at its
endpoint, then the second external
response. The other loop kernel is
bounded pointwise by its Schur norm.
This gives
\begin{equation}\label{gpl:cubic-contraction}
 \left\|\operatorname{Tr}
       (\mathbb G V_3[Z\,\cdot]\mathbb G V_3[Z\,\cdot])\right\|_{t,2}
          \le \|\mathbb G\|_{t,2}^2\|V_3\|_{t,3}^2\|Z\|_{t,2}^2 .
\end{equation}
The scalar tensor norms already sum all
internal coordinate indices. No free
vacuum component remains once the first
external coordinate is fixed. In particular
neither estimate has a volume factor.
The Haar contraction is a direct weighted
matrix product of \(H^\sharp\) and \(H\).

\begin{proposition}[Lipschitz Gaussian polarization]
\label{prop:gpl-polarization-Lipschitz}
For actual input actions whose Taylor
triples satisfy the preceding boundedness
and lower bound, there are \(t>0\) and
\(C<\infty\), independent of depth and
volume, such that
\begin{equation}\label{gpl:polarization-Lipschitz}
             \|\Pi(\mathfrak J)-\Pi(\mathfrak J')\|_{t,2}
                                      \le C\|\mathfrak J-\mathfrak J'\|_s.
\end{equation}
Each polarization itself has bounded
weighted norm. The full constraint and
gauge vertices in the formula are fixed
actual one-step vertices.
\end{proposition}

\begin{proof}
The augmented differences of \(V_3,V_4\)
are precisely the field vertices
\(T_3-T_3'\), \(T_4-T_4'\). All other
vertex terms are fixed.
Subtract the two formulas
\eqref{gpl:actual-Phi-Hessian} by replacing
one factor at a time. For a propagator or
response difference use
\eqref{gpl:inverse-Lipschitz}. For a
vertex difference use its rooted tensor
norm. Each remaining contraction is
bounded by \eqref{gpl:quartic-contraction}
or \eqref{gpl:cubic-contraction}.
The Haar difference uses the same
weighted product bound. These are
finitely many terms with fixed degrees,
which proves the assertion.
\end{proof}

These estimates also pass to compatible
bulk kernels: every rooted contraction
has an absolute exponential majorant.
Period images are summed with that same
majorant, and winding contributions tend
to zero along a growing compatible
volume. Thus \eqref{gpl:polarization-Lipschitz}
holds for the common infinite-lattice
companions whenever the input triples
have those actual companions.

\subsubsection{The implication for actual increments}

For the invariant physical total,
the colour-averaged second-moment
readout of the coefficient is continuous
in this norm. Explicitly its standard
off-diagonal form is
\[
 b(\Pi)=\sum_{u\in\mathbb Z^4}
                  u_\mu u_\nu\,\overline\Pi_{\mu\nu}(u),
                  \qquad \mu\ne\nu ,
\]
after the actual Ward and lattice
symmetries have identified this form.
The bar here denotes the orthonormal
colour average, not complex conjugation.
Since \(|u|^2\le C_t e^{t|u|}\), absolute
moments give
\begin{equation}\label{gpl:coefficient-Lipschitz}
 |b(\Pi(\mathfrak J))-b(\Pi(\mathfrak J'))|
                                \le C_t\|\mathfrak J-\mathfrak J'\|_s .
\end{equation}
The coefficient convention is the physical
one, with the outer coarse lattice rescaled
to unit spacing.

\begin{corollary}\label{cor:gpl-geometric-increments}
The actual homogeneous Taylor triples
have common bulk companions and, by
\cref{prop:cac-Taylor-triple-convergence},
for each \(0<\alpha<1\) obey the following
bound at some fixed positive rate
\(s_\alpha\):
\begin{equation}\label{gpl:actual-triple-convergence}
 \|\mathfrak J_{k+1}-\mathfrak J_k\|_{s_\alpha}
                         \le C_\alpha L^{-\alpha k},
                          \qquad 0<\alpha<1 .
\end{equation}
Then their actual Gaussian increments
\(b_{k+1}^0=b[\Phi_{k+1}]\) satisfy, with
a fixed index shift absorbed into the
constant,
\par\noindent
\begin{minipage}{\linewidth}
\begin{equation}\label{gpl:geometric-actual-increments}
 |b_{k+1}^0-b_k^0|\le C_\alpha L^{-\alpha k},
 \qquad
 |b_k^0-c_0|\le C_\alpha L^{-\alpha k},\qquad
                         c_0=\frac{11}{3\pi^2}\log L .
\end{equation}
\popQED\par
\end{minipage}
\end{corollary}

\begin{proof}
The difference bound follows from
\eqref{gpl:coefficient-Lipschitz}.
Its tail is a geometric series, so the
increments converge with the displayed
rate to some constant. The exact physical
stage composition and the proved massive
lattice comparison give
\eqref{mrc:physical-increment-sum}.
Dividing that cumulative identity by
the number of steps and taking the
Cesaro limit identifies the constant
as \(11\log L/(3\pi^2)\).
This last argument uses the lattice
coefficient, not a quoted continuum
beta coefficient.
\end{proof}

The actual formula and its uniform
Taylor-triple Lipschitz estimate complete
the geometric increment argument on the
homogeneous physical germ. The rooted star norm of
\cref{prop:cac-Taylor-triple-convergence}
dominates \eqref{gpl:vertex-norm}, after the
fixed \(L\)-cell grouping and the fixed
Taylor-factor conventions, and therefore supplies
\eqref{gpl:actual-triple-convergence}.
A formal response expansion alone does
not establish that mesh comparison. A bounded
cumulative determinant error alone
would likewise not give geometric
successive-increment convergence.

% END LAYERED EXPANSION inverse/gaussian_polarization_lipschitz.tex

% BEGIN LAYERED EXPANSION nonflat_effective_gaussian.tex
\subsection{The nonflat effective curl bound and the outer Gaussian gap}
\label{neg:section}

The positive massive row Schur matrix does not,
by itself, give a gap after its artificial
row mass is subtracted. We prove the
needed physical comparison. It gives
the actual outer contour-gauge gap
near a regular nonflat critical reference,
uniformly in the inner fine depth.

This subsection uses a homogeneous inner
torus, with fine spacing \(h=L^{-k}\)
and intermediate coarse spacing one.
The real fine reference \(U\) is an
actual constrained-critical field of
literal mean \(V=M^kU\), with
\begin{equation}\label{neg:reference-bounds}
             \|F(U)\|_\infty\le\kappa,\qquad
             \|j_U\|_\infty\le\gamma .
\end{equation}
The current is the physical Wilson
current in these units. The actual
hierarchical means are on their
regular branches. All ceilings below
are fixed sufficiently small.
Nontrivial holonomies are allowed.
We construct statements on the
critical germ provided by
\cref{prop:pnc-critical-continuation}.
We do not infer existence of this
reference from coarse curvature alone.

Let \(Q\) be its re-centered,
straightened, per-length field
derivative. Let \(\omega\) be its
physical multiplier, so
\(j_U=Q^\sharp\omega\).
The actual readout gives
\(\|\omega\|_\infty\le C\gamma\).
Write \(S\) for the local physical
critical action as a function of
the intermediate coarse field.
The straightened datum and literal
datum differ by their admitted node
gauge. The action is invariant under
that gauge, so their scalar action
germs coincide in the declared
coarse normal coordinates.

\subsubsection{The flat flux row with its physical norm}

We first give the exact flat identity
used in the comparison. Put \(N=h^{-1}\).
For fine translations \(\tau_\mu\),
define
\[
 A_h=h^4\sum_{x\in B_h(0)}\tau_x,
       \qquad S_{\mu,h}=h\sum_{t=0}^{N-1}\tau_\mu^t .
\]
The symbol \(A_h\) here is a block
average, rather than an action.
The set \(B_h(0)\) is the actual
centered \(N^4\)-site block; in
one coordinate its indices are
\(-(N-1)/2,\ldots,(N-1)/2\)
for the odd \(N=L^k\) used here.
The flat field row is
\(Q_\mu^0=A_hS_{\mu,h}\), sampled
at coarse bond origins.
The coarse difference obeys
\[
                    \tau_\mu^N-1=S_{\mu,h}D_{\mu,h}^+.
\]
All these flat translations commute.
Consequently
\begin{equation}\label{neg:flat-flux-identity}
 (\mathrm{d}_cQ^0v)_{\mu\nu}
       =A_hS_{\mu,h}S_{\nu,h}(\mathrm{d}_hv)_{\mu\nu}
                         =:Q^2_{\mu\nu}(\mathrm{d}_hv)_{\mu\nu}.
\end{equation}
The actual fixed-\(L\) hierarchy has
this same linear row: multiplying
the block averages combines the
base-\(L\) digits of the averaging
sites and line positions, giving
the displayed \(N\)-site formula.

Every weight of \(Q^2_{\mu\nu}\)
is nonnegative, its row sum is one,
and its counting column sum is
\(h^4\). The latter identity follows
by summing over all translated
coarse blocks; each fine site has
total block-average weight \(h^4\),
and each \(S_{\mu,h}\) preserves
that total weight. Jensen's
inequality therefore gives
\begin{equation}\label{neg:flat-flux-contraction}
                         \|Q^2f\|_{2,c}\le\|f\|_{2,h}.
\end{equation}
Thus the flat flux row has its
physical norm without a factor
counting fine plaquettes.
The identity is local. It may
be used in a flat parallel frame
on a cover without choosing a
logarithm of a global holonomy.

\subsubsection{A curvature-small covariant flux comparison}

\begin{lemma}[Actual coarse curl controlled by fine curl]
\label{neg:covariant-flux}
For the actual rows at the real
reference \eqref{neg:reference-bounds},
\begin{equation}\label{neg:covariant-flux-bound}
                   \|\mathrm{d}_VQv\|_{2,c}
                 \le C\|\mathrm{d}_Uv\|_{2,h}+C\kappa\|v\|_{2,h}.
\end{equation}
The constants are uniform in \(k\)
and volume.
\end{lemma}

\begin{proof}
For each coarse plaquette choose a
fixed surrounding physical cube
containing all four complete row
carriers. The current-free local
gauge construction gives a fine
potential of size \(C\kappa\)
on this cube, after decreasing the
curvature ceiling by its fixed
radius factor. Its literal coarse
mean has bond angles at most
\(C_L\kappa\): the actual path
lengths and normalized mean
majorants are bounded in these
physical units. Both fine and
coarse comparisons are made in
this same real frame.

The actual coefficient comparison
\eqref{npl:coefficient-comparison}
then gives \(Q=Q^0+\mathcal R_Q\)
on these rows, with
\(\|\mathcal R_Qv\|_{2,c}
 \le C_L\kappa\|v\|_{2,h}\)
on their fixed enlarged carriers.
There is no derivative of this
row error at the fine scale:
the following difference acts
on the coarse unit lattice.
Its \(L^2\) norm is bounded by
a fixed constant. Also
\(\mathrm{d}_V-\mathrm{d}_c=O_L(\kappa)\) as a
coarse matrix, while
\((\mathrm{d}_U-\mathrm{d}_h)v=O_L(\kappa)v\)
on the fine cube in this frame.
Insert these three comparisons
in \eqref{neg:flat-flux-identity}
and use
\eqref{neg:flat-flux-contraction}.
They give the asserted bound
for that plaquette and its
fixed neighbours.

The gauge transformations are
unitary, so these local norm
bounds are covariant. A fine
cell occurs in only a fixed
number of the surrounding
coarse cubes. Summing their
squared estimates proves
\eqref{neg:covariant-flux-bound}.
No single global flat frame
has been selected.
\end{proof}

\subsubsection{The effective Hessian, with its constraint curvature}

Let
\[
 H_L=\mathscr D_U-\Delta_U^{(2)}[\omega],
 \qquad
 G_1^{-1}=H_L+D_URD_U^\sharp+Q^\sharp a_fQ.
\]
This is the actual physical second
precision by
\eqref{acc:physical-second-precision}.
Its row lift is
\[
         \mathcal H=G_1Q^\sharp(QG_1Q^\sharp)^{-1}.
\]
It is horizontal, \(Q\mathcal H=1\),
and \(\|\mathcal Ha\|_{\mathcal E}
\le C\|a\|_{2,c}\).
The reference coarse spacing is
one, so all row weights \(a_f\)
are bounded above and below.

\begin{proposition}[Physical nonflat effective curl bound]
\label{neg:effective-curl}
The actual critical action germ satisfies
\begin{equation}\label{neg:effective-curl-bound}
 \begin{split}
 D^2S(V)[a,a]&\ge
            c\|\mathrm{d}_Va\|_{2,c}^2
                       -C(\kappa+\gamma)\|a\|_{2,c}^2,\\
 |D^2S(V)[a,b]|&\le C\|a\|_{2,c}\|b\|_{2,c}.
 \end{split}
\end{equation}
Its exact Hessian is
\begin{equation}\label{neg:effective-Hessian}
 D^2S(V)=\mathcal H^\sharp H_L\mathcal H
                       =(QG_1Q^\sharp)^{-1}-a_f.
\end{equation}
These are local statements at the
given physical critical reference.
\end{proposition}

\begin{proof}
Differentiate the actual critical
graph in its fixed reference
Landau slice. Its first derivative
is \(\mathcal H\), by the physical
row-saddle identity. Twice
differentiating its constraint
gives
\[
      QA''[a,b]+q''(0)[\mathcal Ha,\mathcal Hb]=0.
\]
The first derivative of the fine
action is \(Q^\sharp\omega\).
The second chain rule for the
action therefore gives
\(D^2S=\mathcal H^\sharp H_L\mathcal H\).
Equivalently the differentiated
stationarity equation is
\[
                  H_L\mathcal H=Q^\sharp D^2S.
\]
The horizontality of \(\mathcal H\)
annihilates its gauge penalty.
After adding the row mass,
multiplication by \(G_1\), then
by \(Q\), proves the second
identity in
\eqref{neg:effective-Hessian}.
Thus the artificial mass is
subtracted in the stated order.

The actual Wilson form estimate
and the constraint-curvature
estimate give
\[
 \langle v,H_Lv\rangle
       \ge(1-C\kappa)\|\mathrm{d}_Uv\|_2^2
                         -C(\kappa+\gamma)\|v\|_{\mathcal E}^2.
\]
Insert \(v=\mathcal Ha\).
The lift bound controls the
last term by \(C(\kappa+\gamma)\|a\|_2^2\).
Since \(Qv=a\),
\cref{neg:covariant-flux} implies
\[
             \|\mathrm{d}_Uv\|_2^2
                   \ge c\|\mathrm{d}_Va\|_2^2-C\kappa^2\|v\|_2^2.
\]
Absorb its last term by the same
lift bound. This proves the
lower inequality.
The fine Lagrangian form is
bounded in energy and the lift
is bounded, which proves the
upper inequality and its
polarized form. All its bounds
use the actual fine Hessian
and multiplier; an ambient
unconstrained Hessian has not
been substituted.
\end{proof}

\subsubsection{The actual outer plain-fibre precision}

Now use the actual one-step outer
pivot chart based at \(V\):
\[
        V(u)=e^{i\Phi_V(u)}V,\qquad
        \mathsf C_V=D\Phi_V(0),\qquad
        \Phi_V(u)|_{\rm free}=u.
\]
Its dimensionless outer block
size is the fixed \(L\). The
first-step chart estimates give
\begin{equation}\label{neg:outer-chart-bounds}
 \|\mathsf C_V\|_{2\to2}\le C_L,\qquad
 \|\mathsf C_Vu\|_2\ge\|u\|_2,\qquad
 \|\langle\mu,D^2\Phi_V(0)\rangle\|_{2\to2}
                                      \le C_L\|\mu\|_\infty .
\end{equation}
These are finite one-step bounds,
independent of the inner depth.
Let \(K_G(V)\) be the actual
relative trace-contour Hessian
of \eqref{afs:relative-contour-gauge}.
On a block-axial real reference
it is the absolute contour
Hessian.

At this physical reference the
coarse gradient is \(\omega\).
The full centred composed-action
Hessian on free outer coordinates
is therefore
\begin{equation}\label{neg:actual-outer-action-Hessian}
 K_A(V)=\mathsf C_V^\sharp D^2S(V)\mathsf C_V
                           +\langle\omega,D^2\Phi_V(0)\rangle.
\end{equation}
The second term cannot be omitted.
It is the nonlinear plain-constraint
contribution to the outer fibre
Hessian. Subtracting the linear
centering of the action has no
effect on this equation.

\begin{lemma}[The geometric outer curl-gauge gap]
\label{neg:outer-geometric-gap}
On the real block-axial, small-curvature
outer reference,
\par\noindent
\begin{minipage}{\linewidth}
\begin{equation}\label{neg:outer-geometric-gap-bound}
      \|\mathrm{d}_V\mathsf C_Vu\|_2^2+\langle u,K_G(V)u\rangle
                                          \ge c_L\|u\|_2^2.
\end{equation}
\popQED\par
\end{minipage}
\end{lemma}

\begin{proof}
Apply the proved actual first-step
gap \eqref{nfp:real-gap} with zero
independent current. Its action
part is
\(\mathsf C_V^\sharp\mathscr D_V^{\rm Wilson}\mathsf C_V\).
The exact plaquette upper estimate
bounds it by
\[
 C_L\|\mathrm{d}_V\mathsf C_Vu\|_2^2
                   +C_L\|F(V)\|_\infty\|\mathsf C_Vu\|_2^2.
\]
The outer curvature is \(O_L(\kappa)\)
by the local mean comparison in
\cref{neg:covariant-flux}. Use
\eqref{neg:outer-chart-bounds} and
decrease \(\kappa\) to absorb this
last mass term into half the
first-step gap. The resulting
inequality is
\eqref{neg:outer-geometric-gap-bound}.
This uses the first-step geometry
and gauge, rather than treating
the nonlocal effective action as
a finite-range Wilson action.
\end{proof}

\begin{proposition}[Nonflat late-depth outer Gaussian gap]
\label{neg:outer-Gaussian-gap}
For fixed sufficiently small
\(\kappa,\gamma\), the actual
outer precision satisfies
\begin{equation}\label{neg:outer-Gaussian-gap-bound}
          c_LI\le \tau K_A(V)+K_G(V)\le C_LI
\end{equation}
for real \(\tau\) in a fixed
interval about one.
The constants are independent
of the inner depth and volume.
\end{proposition}

\begin{proof}
Insert \(\mathsf C_Vu\) in
\eqref{neg:effective-curl-bound}.
Equations
\eqref{neg:outer-chart-bounds}
and \eqref{neg:actual-outer-action-Hessian}
give
\[
 \langle u,K_A(V)u\rangle
       \ge c\|\mathrm{d}_V\mathsf C_Vu\|_2^2
                         -C_L(\kappa+\gamma)\|u\|_2^2.
\]
Add the gauge term and use
\cref{neg:outer-geometric-gap}.
Choose the two ceilings small
enough to absorb the error.
The upper bound follows from
the effective Hessian and the
finite outer chart and gauge
bounds. Finally changing \(\tau\)
near one costs at most
\(C_L|\tau-1|\|u\|_2^2\),
so a fixed interval preserves
both inequalities.
\end{proof}

\subsubsection{The explicit independent-current completion}

We specify the operator to which
the complex assertion applies.
For a field-current pair
\(P=(\mathbf U,\mathbf J)\)
in a small tube about the preceding
real physical pair, form the
canonical \(G,H,G_1\) in
\eqref{acc:named-variants}, in
that order, and put
\begin{equation}\label{neg:off-current-outer-completion}
 \begin{split}
 \omega_P&=H(P)^\sharp\mathbf J,\qquad
 \Gamma_P=(Q_{\mathbf U}G_1(P)Q_{\mathbf U}^\sharp)^{-1}-a_f,\\
 \mathcal K_A(P)&=
      \mathsf C_{\mathbf V}^\sharp\Gamma_P\mathsf C_{\mathbf V}
                   +\langle\omega_P,D^2\Phi_{\mathbf V}(0)\rangle,\\
 \mathcal K_\tau(P)&=\tau\mathcal K_A(P)+K_G(\mathbf V),
                       \qquad \mathbf V=M^k\mathbf U .
 \end{split}
\end{equation}
The lift inside \(\omega_P\) is
the first lift. The row covariance
inside \(\Gamma_P\) uses the
second inverse.
At the physical pair,
\eqref{neg:effective-Hessian}
and \eqref{neg:actual-outer-action-Hessian}
identify this operator exactly
with the physical precision just
estimated.

The actual current inverse,
normalized row inverse and
readout bounds show
\begin{equation}\label{neg:outer-complex-perturbation}
 \|\mathcal K_\tau(P)-\mathcal K_1(U,j_U)\|_{2\to2}
 \le C_L\{\,\varepsilon+\gamma+\gamma_{\mathbf J}+|\tau-1|\,\}.
\end{equation}
Here \(\varepsilon\) is the
weighted \(C^1\) tube radius
about the real regular fine
reference and \(\gamma_{\mathbf J}\)
its independent-current amplitude.
To check this estimate, perturb
the actual inverse and its
normalized row trace in energy
norms, then invert the row matrix
using its real coarse gap.
The row weights are bounded
in the present unit coarse
spacing. The first-lift readout
has its already proved
supremum bound. Finally
the actual plain-mean, pivot
and relative gauge coefficients
have the uniform one-step
derivative bounds
\eqref{neg:outer-chart-bounds}.
Their background changes are
\(O_L(\varepsilon)\) by the
admitted hierarchical row and
node-frame majorants.
These facts give each term in
\eqref{neg:outer-complex-perturbation};
there is no fine operator-norm
comparison at scale \(h^{-2}\).

After reducing the fixed radii,
the right side is less than half
the real gap. Thus
\(\mathcal K_\tau(P)\) is
holomorphic and invertible on
a common complex tube, with
uniform inverse norm.
For real symmetric pairs in
that tube it has the same
positive lower bound. No
Hermitian positivity of the
complex matrix is asserted.

This proves a matched explicit
outer Gaussian precision on
physical regular critical germs
and the stated off-current
extension. Away from the
block-axial locus the relative
gauge in
\eqref{afs:relative-contour-gauge}
specifies the extension. An
independently prescribed
absolute-gauge completion,
signed profile, enlarged cube
domain or history-conditioned
matrix still requires its
own exact comparison.

% END LAYERED EXPANSION nonflat_effective_gaussian.tex

% BEGIN LAYERED EXPANSION actual_coarse_gaussian_admission.tex
\subsection{Local Gaussian domains for the canonical nonflat coarse form}
\label{acg:section}

We now admit independent cube parameters
for the actual canonical coarse Gaussian
near a real physical critical pair.
The real spectral gap and the local
coefficient expansion are separate inputs.
Here they are supplied by
\cref{neg:outer-Gaussian-gap,prop:ncw-named-families}.
The construction below also records the
background domain on which its local
coefficients are holomorphic.

Use a homogeneous inner hierarchy of
arbitrary depth, with inner coarse spacing
one. The outer step has the fixed blocking
integer \(L\). Fix a real physical critical
pair \(P_*=(U_*,j(U_*))\) satisfying the
small curvature and current requirements
of \cref{neg:outer-Gaussian-gap}.
Use its actual coarse datum
\(V_*=M^kU_*\), its one-step plain pivot
chart, and the stated relative contour
gauge. Choose a block axial representative
when identifying the physical absolute
gauge factor.
All real transposes are extended
bilinearly to the complex domains.

For \(P=(U,J)\) near \(P_*\), let
\(\Gamma_P,\omega_P\) and \(K_\tau(P)\)
be the explicit canonical extension in
\eqref{neg:off-current-outer-completion}.
In particular \(\omega_P\) uses the first
lift and \(\Gamma_P\) uses the second
inverse. We retain the
\(\langle\omega_P,\Phi_V''(0)\rangle\)
term in the outer action Hessian.
The following construction uses this
specified matrix. An independent-current
extension of its local Gaussians does
not identify a separately prescribed
history or signed-profile precision.

\subsubsection{The two uniform premises}

Assign the free outer fine-bond coordinates
to their starting unit cells. A cell has
at most a fixed number \(d_L\) of real
coordinates; this number does not involve
the inner fine mesh. The plain pivot map
and its first two derivatives read a
fixed \(L\)-dependent neighborhood.
After an aligned reblocking, the named
families therefore give an exact expansion
\begin{equation}
 K_\tau(P)=\sum_C K_C(P,\tau),\qquad
 \sup_y\sum_{z,C}
       e^{\kappa_0\sigma_0(C)}
                          \|1_yK_C1_z\|
   +\sup_z\sum_{y,C}
       e^{\kappa_0\sigma_0(C)}
                          \|1_yK_C1_z\|\le B_0 .
 \label{acg:actual-precision-family}
\end{equation}
Here \(C\) is a connected carrier of
fixed cubes of side \(M_0\), with its
full determining data, endpoints and
short joins. The shifted length
\(\sigma_0\) is the carrier length of
the positive-walk construction.
The constant \(B_0\) and \(M_0\) are
fixed before depth and volume.
The exponent \(\kappa_0\) can be
chosen at any required fixed finite
value by making the choices in
\cref{prop:ncw-named-families} first.
The permitted current ceiling and
complex tube are retained from
those choices.

At real physical data and
\(\tau\) near one, the geometric
gap gives
\begin{equation}
                  c_L I\le K_\tau(P_*)\le C_L I .
 \label{acg:actual-positive-gap}
\end{equation}
After a fixed decrease of the
background and current tube, the
same real bounds, with changed
constants, hold for every real
admitted \(P\). This follows from
the actual family perturbation bound;
it does not require the nearby
independent current to be a Wilson
gradient. Complex perturbations
have the same operator-norm bound.
The variable \(\tau\) ranges in a
fixed disc about one.

For completeness the locality
in \eqref{acg:actual-precision-family}
includes the row inverse in
\(\Gamma_P\). It is supplied by
\eqref{ncw:mass-removed-coarse-family},
with its actual mass subtraction.
The outer chart factors and
\(\Phi_V''(0)\) enlarge the carrier
by only a fixed star.
The rooted family for \(\omega_P\)
retains its fine-current input
bond by \cref{prop:ncw-rooted-multiplier}.
Thus the current term has a
genuine local carrier, as does
the row-Schur term. A completed
kernel decay bound alone would
not prove \eqref{acg:actual-precision-family}.

\subsubsection{An admissible local background tube}

The background tube is a weighted
\(C^1\) tube in the actual fine
relative field \(A\), with an
independent weighted current ball
for \(J-J_*\). For the homogeneous
geometry the assigned physical
cell scale is one.
Given a finite carrier, retain a
fixed complete data collar.
A local \(C^1\) field on that
collar can be extended to a global
relative field by a smooth cutoff
which equals one on the read set.
The cutoff is formed at the
unit-cell scale, rather than the
inner fine mesh.
It satisfies
\[
             |\chi|+|D\chi|\le C_L .
\]
Consequently
\begin{equation}
       \|\chi A\|_{C^1}\le C_L\|A\|_{C^1(\text{collar})}.
 \label{acg:local-tube-extension}
\end{equation}
To construct such a cutoff, take
the distance in the unit-cell
graph from the read set, truncate
it in a collar of fixed width,
and interpolate its values by
the fixed smooth partition on
unit cells. Bounded overlap
gives the displayed derivative
bound. Make it zero before the
outer edge of the supplied
collar. Extension by zero then
adds no fine-mesh jump.
Multiplication of the independent
current by a bounded cutoff
preserves its current norm and
takes no derivative of that current.

Choose the local radius smaller
than the global tube radius
divided by the constant in
\eqref{acg:local-tube-extension}.
Every such local datum has an
admitted extension relative to
the fixed real reference.
Exact determining-data locality
in \cref{prop:ncw-named-families}
makes the coefficient independent
of which extension is used.
This proves a local holomorphic
domain for each coefficient.
These are \(C^1\) field domains;
they are not independent scalar
fine-bond discs with a mesh-free
radius. The cube parameters
introduced below are independent
of these field variables.
The geometry is the fixed
homogeneous geometry in this
statement. No extension theorem
for an arbitrary partly specified
proper history is used.

\subsubsection{Full-carrier truncation and a combined short power}

Fix a requested cube-parameter
radius \(S>2\). Choose
\(b>2(C_L+1)\), and set
\[
                       H=I-K_\tau/b .
\]
Symmetrize each carrier coefficient
by its transpose. Its endpoints
and determining data are unchanged,
its row/column family norm does
not increase, and its sum is still
the symmetric matrix \(H\).
For a positive integer \(R\), put
\begin{equation}
 H_<=\sum_{\sigma_0(C)\le R}H_C,\qquad
 H_>=H-H_< .
 \label{acg:full-carrier-truncation}
\end{equation}
The tail has Schur norm at most
\(Ce^{-\kappa_0R}\).
For large \(R\), real data satisfy
\begin{equation}
                0\le H_<\le qI,\qquad q<1 ,
 \label{acg:near-positive-contraction}
\end{equation}
by \eqref{acg:actual-positive-gap}.
The lower bound uses \(H\ge I/2\)
and a tail smaller than \(1/4\).
Shrink the common complex tube
once more if necessary.
Then \(\|H_<\|_{2\to2}\le q'<1\)
throughout it, for a fixed \(q'\).

Every carrier in \(H_<\) has
diameter at most \(CM_0(R+1)\).
For \(j\ge1\), combine the matrix
entries of \(H_<^j\) before taking
absolute values. Its matrix range
is at most
\[
                     B_j=CM_0j(R+1).
\]
The number of unit-cell coordinates
in a ball of radius \(B_j\) is
at most \(C_L(1+B_j)^4\).
A Hilbert row has norm at most
\((q')^j\); Cauchy--Schwarz on
that row, and then on columns,
therefore gives
\begin{equation}
                  \|H_<^j\|_{\rm Sch}
                      \le C_L(1+B_j)^2(q')^j .
 \label{acg:combined-power}
\end{equation}
The data read by this combined
entry are retained in the two
radius-\(B_j\) endpoint balls.
Each short word has its whole
carrier in those balls: its
successive carriers are connected
and have total diameter at most
the sum of their declared ranges.
The balls can be joined by a
fixed short path. This is the
reason for truncating by the
full carrier in
\eqref{acg:full-carrier-truncation}.
Truncation by endpoint distance
would not justify this data claim.

The exact long-word remainder is
\begin{equation}
 E_j=H^j-H_<^j
       =\sum_{i=0}^{j-1}H_<^iH_>H^{j-1-i}.
 \label{acg:exact-long-word}
\end{equation}
Keep its actual joined-carrier
family. At a fixed reserved
carrier rate its norm is bounded
by
\begin{equation}
                  \|\mathbf E_j\|
                         \le CjB_H^j e^{-cR}.
 \label{acg:long-word-bound}
\end{equation}
The constants are independent
of the inner depth.
Choose \(R=Dm\), with
\(cD>2\log(1+B_H)+2\).
Next choose an aligned cube
side \(M_{\rm G}\) larger than
a fixed multiple of \(B_m\).
It may be taken to be a power
of \(L\). The combined short
entries now meet a bounded
number of these cubes, with
a bound independent of \(m\).

For the long terms, reblocking
does not lose their rate.
An old connected carrier can
be visited by a curve of length
at most \(CM_0\sigma_0(C)\).
Covering that curve and its
\(CM_0\) collar by
\(M_{\rm G}\)-cubes gives
\begin{equation}
 \# C_{\rm G}\le
       C_L\{1+(M_0/M_{\rm G})\sigma_0(C)\}.
 \label{acg:reblocked-parameter-cost}
\end{equation}
Choose \(M_{\rm G}\) also so
that the coefficient of
\(\sigma_0(C)\) in
\((\log S)\#C_{\rm G}\)
fits in the reserved old
carrier rate. Thus weighting
a full carrier by independent
parameters of modulus at most
\(S\) preserves
\eqref{acg:long-word-bound}.
No fine-background coordinate
count enters
\eqref{acg:combined-power}:
the coordinates counted there
are the outer free Gaussian
coordinates at spacing one.
The fine background is read
as determining data.

\subsubsection{The actual local covariance}

Let \(\mathbf T\) be the family
obtained by combining and filing
\(H_<^m\) as above and adding
the long-word family \(E_m\).
Its sum is \(H^m\). Multiply
each term by the product of
its declared cube parameters,
and denote the resulting matrix
by \(T_s\).
On the held root guard replace
each parameter by one.
For \(j<m\), similarly put
\[
          A_{j,s}=H_<^j+(E_j)_s .
\]
The near prefix is held without
parameters. Its rows from a
root wall read only its fixed
radius-\(B_m\) guard.
The long prefix keeps its
complete carrier parameters.
The preceding bounds give
\begin{equation}
 \sup_{|s_\Delta|<S}\|T_s\|_{\rm Sch}
 \le C_LS^{C_L}(1+B_m)^2(q')^m
                   +CmB_H^me^{-cDm}=:q_m\longrightarrow0 .
 \label{acg:small-admitted-power}
\end{equation}
Its convergence is exponential
up to a polynomial in \(m\).

For a finite root coordinate
set \(W\), let \(\chi_W\)
denote coordinate restriction,
not a fluctuation cutoff.
With \(c_n=4^{-n}\binom{2n}{n}\),
define
\begin{equation}
 \begin{split}
 R_W(s)&=b^{-1/2}\chi_W
        \sum_{j<m}A_{j,s}
                   \sum_{v\ge0}c_{mv+j}T_s^v,\\
 D_W(s)&=R_W(s)R_W(s)^{\mathsf T}.
 \end{split}
 \label{acg:local-Gaussian-covariance}
\end{equation}
The series converges normally
in the weighted row/column
norm after \(m\) is fixed
large enough. Prefix norms
grow only polynomially in
\(R\), by
\eqref{acg:combined-power}
and \eqref{acg:long-word-bound}.
The binomial coefficients
satisfy \(0<c_n\le1\).
At complete parameters the
ordinary matrix identity is
\begin{equation}
 R_W(1)=\chi_WK_\tau(P)^{-1/2},\qquad
 D_W(1)=\chi_WK_\tau(P)^{-1}\chi_W^{\mathsf T}.
 \label{acg:actual-marginal-identity}
\end{equation}
Indeed \(H^{mv+j}=H^j(H^m)^v\),
so grouping the convergent
binomial series by its residue
class modulo \(m\) gives
the first identity.
Its square gives the second.
This identity is asserted for
the complete matrices. Their
filed summands need not commute.

For real comparison data put
\[
 R_{0,*}=b^{-1/2}
           \sum_{j<m}c_jH_{<,*}^{\,j},\qquad
 D_{0,W}=\chi_WR_{0,*}^2\chi_W^{\mathsf T}.
\]
By \eqref{acg:near-positive-contraction},
\begin{equation}
 b^{-1}I\le D_{0,W}\le C_L I
 \label{acg:held-prefix-gap}
\end{equation}
for every \(W\).
The lower bound follows by
zero-extending a vector in
\(W\): \(R_{0,*}\ge b^{-1/2}I\).
There is no Dirichlet inverse
of an indefinite saddle here.

The long prefixes and the
\(v\ge1\) terms in
\eqref{acg:local-Gaussian-covariance}
are bounded by a polynomial
in \(m,R\) times \(q_m+
mB_H^me^{-cDm}\).
Their norms can therefore
be made smaller than any
prescribed fixed inverse
power of \(M_{\rm G}\),
since \(M_{\rm G}=O_L(M_0m^2)\).
Finally reduce the common
complex field, current and
\(\tau\) tubes so that the
finite held prefixes have
the same perturbation budget.
Their coefficient Lipschitz
bounds follow from Cauchy
on the local tube in
\eqref{acg:local-tube-extension}.
All choices are uniform in
depth and volume.
For example they can ensure
\begin{equation}
 \sup_{|s_\Delta|<S}
 \|D_{0,W}^{-1/2}(D_W(s)-D_{0,W})
                                      D_{0,W}^{-1/2}\|
                  <\frac14
 \label{acg:complex-marginal-sector}
\end{equation}
simultaneously for every
root coordinate set \(W\).
The finer \(o(M_{\rm G}^{-4})\)
budget is available for
per-cube determinant estimates.

\begin{proposition}\label{acg:canonical-Gaussian-admission}
For every fixed \(S>2\), the
canonical nonflat outer precision
specified above has the local
Gaussian extension
\eqref{acg:local-Gaussian-covariance}.
It is holomorphic on independent
cube discs \(|s_\Delta|<S\) and
on the common local \(C^1\)
background, independent-current
and \(\tau\) tubes.
Its complete covariance is the
actual marginal
\eqref{acg:actual-marginal-identity}.
Every local covariance has
the uniform complex sector
\eqref{acg:complex-marginal-sector}.
Its coefficients read their
full declared data carriers.
Setting parameters outside a
wall to zero gives literal
wall dependence and fixed-root
heredity, with the held guard
retained as part of that wall.
These assertions are uniform
in the homogeneous inner depth
and compatible volume.
\end{proposition}

\begin{proof}
The preceding construction proves
the identities, bounds and normal
convergence. A term which survives
zero parameters outside a wall
has its whole parameter carrier
there, in addition to the fixed
held guard. Its local coefficient
reads only those data.
Two wall systems which agree on
these complete read sets therefore
give the same coefficient.
This proves heredity.
Take components after adding
their complete held guards and
read sets, merging components
whose guarded sets meet.
Separated components then use
disjoint coordinate sets and
read sets; no surviving joined
carrier connects them.
Their Gaussian covariance is
block diagonal, and their
normalized Gaussian integrals
factor.

On small tori perform the finite
construction on one common
sufficiently large cover.
Retain all lifted data carriers
before truncation and all short
power endpoint balls before
projection. Average families
over deck shifts before descent.
Periodic auxiliary parameters
are identified on their images,
and held guards are their full
deck-periodic inverse images.
Their full lifted carrier
weights remain charged.
The row/column Schur bounds
descend without a sheet factor,
as in \cref{cor:cci-row-inverse}.
Descend the unrestricted root
matrix before restricting its
coordinates to \(W\).
Indeed let \(I\) be isometric
pullback from the base to the
periodic sector of the cover.
Periodicity of parameters,
guards and averaged families
gives \(R_{\rm cov}I=IR_{\rm base}\),
and the periodic sector is
reducing for these matrices
and their transposes. Thus
\[
 I^{\mathsf T}R_{\rm cov}R_{\rm cov}^{\mathsf T}I
                    =R_{\rm base}R_{\rm base}^{\mathsf T}.
\]
Only then restrict the base
coordinates to \(W\).
At complete parameters the
classical operators and their
positive matrix functions
restrict to this sector, so
this covariance is the actual
finite matrix marginal.
A covariance restricted to one
chosen lift of \(W\) would not
establish that assertion.
This is an operator identity,
not a comparison of quantum
partition functions on
different volumes.
\end{proof}

\subsubsection{Normalized complex Gaussian integrals}

The sector just proved also
specifies the determinant branch
and the normalization used in
local Gaussian estimates.
For \(d=\dim W\), define
\begin{equation}
 d\mu_{D_W}(u)=
 (2\pi)^{-d/2}\det(D_W)^{-1/2}
                  e^{-u^{\mathsf T}D_W^{-1}u/2}\,du,
 \qquad u\in\mathbb R^d .
 \label{acg:normalized-local-Gaussian}
\end{equation}
Continue the square root from
the positive real comparison
matrix \(D_{0,W}\).
Writing
\(D_W=D_{0,W}^{1/2}(I+E)D_{0,W}^{1/2}\),
with \(\|E\|<1/4\), gives
\[
 |\det(I+E)|^{-1/2}\le(4/3)^{d/2},
 \qquad
 \operatorname{Re}
        u^{\mathsf T}D_W^{-1}u
                 \ge c_L|u|^2
                 \quad(u\in\mathbb R^d).
\]
The inverse estimate follows
from the norm-convergent series
for \((I+E)^{-1}\) and
\eqref{acg:held-prefix-gap}.
Consequently
\begin{equation}
       \int_{\mathbb R^d}|d\mu_{D_W}|\le C_L^d,
       \qquad \int_{\mathbb R^d}d\mu_{D_W}=1 .
 \label{acg:Gaussian-normalization-bound}
\end{equation}
The second equality is first
the real Gaussian identity.
The same dominating Gaussian
on every smaller sector permits
holomorphic continuation to
the declared domain.
Polynomial moments and fixed
parameter derivatives follow
by differentiating that
dominated integral, or by
Cauchy on a smaller sector.
The empty selected expectation
is exactly one.

The factor \(C_L^d\) is local:
there are at most
\(C_LM_{\rm G}^4\) free
coordinates per localization
cube of a wall. Thus it costs
a fixed exponential per wall
cube, with no factor involving
the total torus or the inner
fine coordinate count.
This is the normalization
needed when a local Gaussian
is used in a connected
polymer integral.
It does not assert zero-freeness
of every sharply truncated
whole-volume Gaussian integral.
The normalized cutoff expansion
and its interacting vertices
must still be formed with
their actual frozen cutoffs.

The result proves the Gaussian parameter-domain statement
for the matrix and reference domain stated here.
\Cref{pgc:section} proves the proper-hierarchy version for
\eqref{pgc:actual-precision}, retaining its second-lift outer
current readout. The two off-current formulas agree on physical
critical pairs; they are not identified elsewhere.
Nonlinear core membership is established by the actual tower
and substitution maps in \cref{htc:all-towers,htc:parameter-core}.

% END LAYERED EXPANSION actual_coarse_gaussian_admission.tex

% BEGIN LAYERED EXPANSION inverse/proper_outer_gaussian_closure.tex
\subsection{The fluctuation precision on a proper inner hierarchy}
\label{pgc:section}

The mass removed from a row Schur matrix is essential for the
fluctuation covariance. We prove its lower bound on the actual
outer constraint, including the second derivative of that constraint.
The inner hierarchy may have proper regions and mixed scales.
The current remains an independent variable until the final physical
identification.

Work in level-$k$ units. The inner sequence ends at $\Omega_k$;
its top row length is one. Retain an additional union
$\Omega_{k+1}$ of complete $L$-blocks, with its prescribed separation
from $\Omega_k^c$. Put $\Theta'=\Omega_{k+1}\cap\mathbb T^{(k+1)}$
for the block centres and $\Theta=B(\Theta')$ for their level-$k$
sites. Let $X$ be the level-$k$ bonds touching $\Theta$, and
$\check X$ their nonpivot bonds. All pivot double blocks and
contour paths are retained. The separation floor below puts this
entire stock inside the interior top rows of the inner sequence.
Write $\iota$ for extension from $X$ by zero to its independent
row space. Thus $\iota^\sharp\iota=1$ in counting coordinates;
there is no mixed-scale measure factor on these length-one rows.
The field and row pairings elsewhere retain their physical weights.
The marked geometry may retain inward bonds, but $Q$ and its row
space use only the independent interior and outward rows.

Use a real regular reference and the common weighted $C^1$
complex tube of \cref{prop:ncw-named-families,dnr:linear-families}.
The independent current satisfies its weighted supremum bound.
The enlarged regularity clause includes every complete cube needed
in the two adjacent layers, including cubes lying wholly in the
upper member. Choose all proof tiles first, then a common larger
separation floor. The ceilings are independent of the fine mesh,
depth, volume, conditioning set and all larger admitted geometries.

\subsubsection{The precise form and its two multiplier readouts}

Write $V=\mathcal M^k\mathbf U$ on the complete outer stock and
let $q_V$ be its raw one-step constraint logarithm. The actual pivot
graph has
\[
 q_V(\Phi_V(u))=0,\qquad
 \Phi_V(u)|_{\check X}=u,\qquad C_V=D\Phi_V(0).
\]
Let $h_V$ be the pivot right inverse of $q_V'(0)$ and write
$C_V^{(2)}=\tfrac12q_V''(0)$. This superscript distinguishes the
quadratic row map from the pivot lift $C_V$.
The graph derivative satisfies
\begin{equation}
 \Phi_V''(0)[u,v]=-2h_VC_V^{(2)}(C_Vu,C_Vv).
 \label{pgc:pivot-second-jet}
\end{equation}
Indeed the free coordinates have zero second derivative, and twice
differentiating the constraint determines all pivot coordinates
by its invertible pivot block. Every map in this identity reads
its complete one-step double block and has a bounded row and
column norm depending only on $L$.

Let $b_Vw(y,x)$ be the derivative of the relative contour logarithm
\[
 \frac1i\log\{(e^{iw}V)(y,x)V(y,x)^{-1}\}.
\]
Set $G_V^{(2)}=b_V^\sharp b_V$, with one unit weight for every
noncentral site in a retained block. This is the relative contour
quadratic form. It is the Hessian of the absolute trace-contour
functional when the reference averaged contours equal one.

The ordered inner operators are exactly
\eqref{acc:named-variants}. In particular put
\begin{equation}
 \begin{aligned}
 \Sigma_1&=(QG_1Q^\sharp)^{-1},\\
 \omega_*&=H^\sharp\mathbf J,\\
 \omega_{1*}&=H_1^\sharp\mathbf J,\\
 E_V^{(2)}[\omega](w,z)
       &=2\langle\omega,\iota h_VC_V^{(2)}(w,z)\rangle,\\
 \Delta_{\rm out}
       &=\iota^\sharp\Sigma_1\iota-a_{\rm top}
                                      -E_V^{(2)}[\omega_{1*}],\\
 A_Y&=e_Y^\sharp C_V^\sharp
                 (\Delta_{\rm out}+G_V^{(2)})C_Ve_Y .
 \end{aligned}
 \label{pgc:actual-precision}
\end{equation}
Here $e_Y$ extends the free coordinates on a conditioning set
$Y=\check X(B(\Theta''))$, $\Theta''\subset\Theta'$, by zero.
Its retained blocks and collars are those of the conditioning
geometry. The constant top row mass is $a_{\rm top}=a_0I$;
for the unit mass convention it is $I$.
The term inside $G_1$ uses $\omega_*$; the outer constraint term
uses $\omega_{1*}$. These are distinct away from physical currents.

\subsubsection{A lower bound before the outer gauge is added}

\begin{lemma}[The top-supported Schur form]
\label{pgc:top-Schur}
At every admitted real pair, for a top row field $w$ supported
on $X$,
\begin{equation}
 \langle w,(\iota^\sharp\Sigma_1\iota-a_{\rm top})w\rangle
 \ge c\|\mathrm d_Vw\|_{2,\mathcal P_X}^2
                   -C(\alpha_0+\gamma_0)\|w\|_2^2 .
 \label{pgc:effective-curl}
\end{equation}
Here $w$ is extended by zero off $X$, and $\mathcal P_X$ is the
complete finite stock of level-$k$ plaquettes used by the outer
block curl estimates. The form on the left has operator norm
at most a fixed constant.
\end{lemma}
\begin{proof}
Put $v=H_1\iota w$. The actual proper row lift gives
$Qv=\iota w$ and $\|v\|_{\mathcal E}\le C\|w\|_2$.
The equality $G_1^{-1}v=Q^\sharp\Sigma_1\iota w$ gives
\begin{equation}
 \begin{split}
 \langle w,(\iota^\sharp\Sigma_1\iota-a_{\rm top})w\rangle
  =\langle v,(&\mathscr D_U-\Delta'_{\pi,U}[J]
       -S^\sharp\Delta_U^{(2)}[H^\sharp J]S
                                      +DRD^\sharp)v\rangle .
 \end{split}
 \label{pgc:Schur-identity}
\end{equation}
This is the variational identity for the complete Schur matrix,
followed by its exact row-mass subtraction. It uses neither
criticality nor horizontality of $v$.

At the real reference $R$ is an orthogonal projection, so the
last term is nonnegative. The Wilson plaquette form and the
current form estimates of \cref{ch:layered-current-precision}
and \eqref{acc:explicit-current-map} bound the right side of
\eqref{pgc:Schur-identity} below by
\[
 c\|\mathrm d_Uv\|_2^2
                   -C(\alpha_0+\gamma_0)\|v\|_{\mathcal E}^2.
\]
The multiplier $H^\sharp J$ has the required row-weighted
supremum bound. The second constraint jet and $S$ have their
bounded energy forms, so its term has the stated $C\gamma_0$
bound. No derivative of the independent current occurs.

Every plaquette in $\mathcal P_X$ and every carrier of its four
top rows lies wholly in $\Omega_k$. There the assigned scale is
one. The local proof of \cref{neg:covariant-flux} therefore gives
\[
 \|\mathrm d_V Qv\|_{2,\mathcal P_X}
        \le C\|\mathrm d_Uv\|_2+C\alpha_0\|v\|_{\mathcal E}.
\]
To see why the proper exterior changes nothing, choose a fixed
physical cube around each of these complete plaquettes. In its
real small gauge the actual top row differs from the flat flux
row by $O(\alpha_0)$ in both physical row and column norms.
The flat identity \eqref{neg:flat-flux-identity} and its Jensen
bound have no depth factor. The coarse derivative has unit
spacing. The cubes have bounded overlap; all norms outside
them on the right may simply be kept. Thus the estimate does
not replace the proper inner inverse by a homogeneous inverse.
Since $Qv=\iota w$, it proves \eqref{pgc:effective-curl}.
The upper bound follows from the proper row inverse and the
constant mass on the supported top rows.
\end{proof}

\begin{lemma}[The outer geometric inequality]
\label{pgc:outer-geometry}
For every retained block family and its full pivot stock,
\begin{equation}
 \|u\|_2^2\le C_L\{
       \|\mathrm d_V C_Ve_Yu\|_{2,\mathcal P_X}^2
                     +\|b_VC_Ve_Yu\|_2^2\}.
 \label{pgc:geometric-gap}
\end{equation}
On the free subspace whose reconstructed field vanishes on
each retained box tree, the same inequality holds with the
second term omitted.
\end{lemma}
\begin{proof}
First take a flat frame. In a block, compare a path to a site
with each coordinate-order path to that site. Successive
plaquette swaps express their difference as a sum of at most
$C_L$ plaquette curls. Hence an internal bond is its mean
contour-potential difference plus a sum of these curls.
Squaring and summing bounds all internal bonds by the two
terms on the right. In the tree version use the unique tree
path instead; its potential is zero.

Across a block face, a plaquette with one transverse step
compares neighboring crossing bonds. Its remaining two
bonds are internal and already bounded. The transverse
face grid is connected and has at most $L^3$ sites. Thus all
crossing values differ from a common face value by controlled
curls and internal bonds. The plain mean of that face is its
common crossing value plus controlled internal path terms.
It vanishes because $q_V'(0)C_Ve_Yu=0$, including each boundary
pivot row. This controls the common value. Outside retained
blocks the extended field has zero internal bonds, so the
same argument applies to a boundary face. A fixed bond or
plaquette enters at most $C_L$ such block comparisons.
The free restriction of $C_Ve_Yu$ is $e_Yu$, proving the claim
at a flat reference.

For a real small-curvature reference choose the regular real
gauge on each block and its neighboring double blocks. Their
sizes depend only on $L$. The plain row, pivot lift, relative
contour derivative and coarse curl differ from their flat
counterparts by $C_L\alpha_0$ in the finite row and column
norms. The gauge changes are unitary. Absorb this perturbation
in half the flat inequality and sum the same bounded-overlap
comparisons. Vanishing on a box tree is preserved by these
linear unitary coordinate changes.
\end{proof}

\begin{proposition}[Proper fluctuation covariance]
\label{pgc:proper-covariance}
For fixed sufficiently small ceilings, the exact matrix
\eqref{pgc:actual-precision} satisfies
\begin{equation}
                      c_LI\le A_Y\le C_LI
 \label{pgc:real-gap}
\end{equation}
at every admitted real pair. The same bounds hold on the
strict tree-zero subspace when the gauge term is omitted.
The inverses extend holomorphically to a common weighted
$C^1$ complex background tube and independent-current ball,
with uniform inverse norms. All conclusions are uniform in
the proper inner hierarchy and the conditioning set.
\end{proposition}
\begin{proof}
The outer second-jet bound and the proper readout estimate give
$\|E_V^{(2)}[H_1^\sharp J]\|\le C_L\gamma_0$.
Insert $w=C_Ve_Yu$ in \eqref{pgc:effective-curl}, add the
relative gauge, and use \cref{pgc:outer-geometry}. The pivot
lift is bounded and its free restriction is the identity.
Decreasing $\alpha_0+\gamma_0$ absorbs the negative mass
term. The same proof on the tree-zero subspace omits the
gauge and uses the second part of that lemma. All upper
bounds follow from the bounded factors in
\eqref{pgc:actual-precision}.

The proper energy and normalized row inverse bounds, together
with their local coefficient families, imply
\[
 \|A_Y(\mathbf U,\mathbf J)-A_Y(U,J_0)\|_{2\to2}
          \le C_L(\|A'\|_{C^1_{\rm weighted}}
                    +\|\mathbf J-J_0\|_{\rm current}).
\]
The same holds on the tree-zero free coordinates. For the row
Schur factor use its normalized inverse difference identity;
for the two readouts retain their different lifts. The outer
pivot and contour coefficients have fixed finite ranges.
Restriction and zero extension have norm one, so the bound
is independent of $Y$. It is a relative energy comparison,
not a fine-lattice operator comparison at scale $\eta^{-2}$.
A perturbation below half the real gap proves the assertion.
\end{proof}

\subsubsection{Local inverse families and logarithms}

The locality assertion requires more than the completed
operator gap. We record how the existing carrier construction
applies to this particular matrix. The mass-removed family
\eqref{ncw:mass-removed-coarse-family}, the rooted family of
$H_1^\sharp J$, and the finite outer factors give an exact
connected-carrier expansion of $A_Y$. For this second-lift
readout repeat the proof of \cref{prop:ncw-rooted-multiplier}
with the named family of $H_1$ in place of $H$: its physical
transpose pairs a unit top-row seed with the fine current,
and the family column bound sums the cellwise energy
pairings. It gives $C\gamma_0$ with no whole-volume current
norm. The bare restriction identity follows from $Q_0H_1=P_0$,
so that its separate bare term is retained exactly.
The precision family's row and column
sums, charged by a fixed exponential carrier rate, are
bounded. Each term retains the complete inner regions,
outer blocks, pivot rows and the trace of $Y$ on its read set.
These are precisely the data needed to multiply two terms.

The inverse construction of \cref{cor:cci-row-inverse} can
now be applied on the outer free coordinates, whose spacing
is one and whose multiplicity per cell is bounded by $C_L$.
Here are its domain and counting checks in this setting.
Choose $b>2C_L$ and write $H_Y=I-A_Y/b$. On a smaller common
tube it has Hilbert norm at most $q<1$. Truncate its family
by full carrier length at $R$, giving $H_<+H_>$. The omitted
family has norm $Ce^{-cR}$. Combine $H_<^m$ before taking
absolute values. Its matrix range is $O(mR)$ in physical
units, and hence
\begin{equation}
 \|H_<^m\|_{\rm Sch}
                  \le C_L(1+mR)^2q_1^m,
 \qquad q_1<1.
 \label{pgc:short-power}
\end{equation}
The complete data of every contributing short word lie in
the corresponding endpoint balls. Enlarge the separation
floor so that these balls, including their fixed carriers,
remain inside $\Omega_k$. This is possible after $m,R$ have
been fixed, independently of depth. In these balls the inner
assigned scale is one; in the extended outer sequence the
only additional scale is $L$. Consequently the parameter
count in a reblocked endpoint ball is bounded independently
of the number of finer layers. The arbitrary boundary of
$Y$ changes its recorded zero-extension mask, not that count.

The same collar gives the required local background domains.
Relative to the fixed global real reference, extend a local
complex potential from a short-word read ball by multiplying
it by a smooth cutoff equal to one on all reads and zero
before the edge of its supplied collar. Its first difference
is bounded in unit scale. The entire cutoff collar lies in
$\Omega_k$, so all its weights have scale one; multiplication
preserves the weighted $C^1$ norm within a fixed factor.
Extend by zero outside the collar. A bounded cutoff preserves
the independent-current norm without taking a current
derivative. Decrease the local complex radii by that factor.
This supplies an admitted global extension for estimating
each combined short coefficient. Its value is independent
of the extension by exact determining-data locality.
Long coefficients retain the local domains of their original
complete NCW carriers. This does not require an extension
of an arbitrary partially specified proper history.

Keep the long-word identity
\[
 H_Y^m-H_<^m
       =\sum_{i=0}^{m-1}H_<^iH_>H_Y^{m-1-i}
\]
as a family with its full joined mixed-scale carriers. Its
norm is at most $CmB_H^me^{-cR}$. Take $R=Dm$ with $D$ fixed
large, then $m$ large. Reblock after these choices. The
short-power term meets only a bounded number of new cubes;
the long carriers pay for their cube parameters from a
reserved part of their old carrier rate. Thus, for every
fixed requested parameter radius $S>2$, the filed $m$th
power has family norm less than $1/2$ on $|s_\Delta|<S$.
The finitely many prefixes are bounded. The exact identity
\[
 A_Y^{-1}=b^{-1}\Big(\sum_{j=0}^{m-1}H_Y^j\Big)
                                      (I-H_Y^m)^{-1}
\]
then gives a normally convergent local inverse family.
This is the inverse of the compressed precision, not a
marginal of an inverse taken before conditioning.

For the logarithm use the same grouped words in
\begin{equation}
 \log A_Y=(\log b)I-
       \sum_{j=1}^{m}H_Y^j
          \sum_{v\ge0}\frac{(H_Y^m)^v}{mv+j} .
 \label{pgc:local-log}
\end{equation}
At complete parameters this is the convergent real spectral
logarithm and its holomorphic continuation. At other parameters
the formula specifies its filed local extension. Its diagonal
trace at a bond is a sum of closed joined carriers containing
that bond. The same family norm bounds it, so the bound per
outer cell is uniform. Differences between two regions or
conditioning sets cancel term by term until a carrier meets
their differing read sets. Parameter walls and fixed-root
heredity follow by setting to zero every term whose complete
carrier leaves the wall. No coefficient is evaluated on an
unrecorded extension of the background.

For normalized Gaussian integrals use the binomial grouped
series \eqref{acg:local-Gaussian-covariance}, with $A_Y$ in
place of the homogeneous matrix. The two premises of that
construction have now been proved: \eqref{pgc:real-gap} and
the complete carrier family just described. Its combined
short words read the same top-scale endpoint balls; its
long words retain the same mixed-scale carriers. Its held
real prefix is at least $b^{-1/2}I$, and its tail can be
made arbitrarily small before decreasing the common complex
tube. Consequently it gives the uniform local complex sector
\eqref{acg:complex-marginal-sector} and the normalized
Gaussian identity \eqref{acg:Gaussian-normalization-bound}
for every coordinate selection inside $Y$. At complete
parameters its covariance is the selected marginal of
$A_Y^{-1}$. The conditioning by $Y$ has already occurred.
The determinant branch is continued from the same positive
real precision. These assertions do not identify two
independently chosen off-parameter extensions of a logarithm.
For small periods perform the construction on one common
compatible cover and retain all lifted read balls before
projection. Average the filed families over deck shifts,
then descend the complete operator before any coordinate
selection. The row and column sums descend without a sheet
factor, as in \cref{cor:cci-row-inverse}; the reducing-sector
calculation in the proof of
\cref{acg:canonical-Gaussian-admission} identifies the
descended covariance with the actual base matrix inverse.

\subsubsection{Identification at a physical critical pair}

Suppose now that $U$ is a regular constrained-critical field
of the inner hierarchy and $J=j_U=Q^\sharp\omega$. Then
$H^\sharp J=H_1^\sharp J=\omega$, since both lifts are right
inverses of $Q$. The physical Ward identity makes $H_1$ the
horizontal derivative of the actual critical graph, and
the differentiated constraint gives its Lagrangian Hessian.
Thus the second derivative of the critical action with
respect to the supported top datum is
$\iota^\sharp\Sigma_1\iota-a_{\rm top}$, as in
\eqref{neg:effective-Hessian}, now on the proper hierarchy.
Composing with the actual outer pivot graph adds
$\langle\omega,\iota\Phi_V''(0)\rangle$. By
\eqref{pgc:pivot-second-jet} this is exactly the negative
$E_V^{(2)}[H_1^\sharp J]$ term in
\eqref{pgc:actual-precision}. Adding the specified relative
contour Hessian and restricting free coordinates gives
$A_Y$ exactly. In an averaged-contour axial representative
it is also the absolute-contour physical Gaussian precision.
The equality continues on the corresponding physical
holomorphic chart. No identification with an independently
prescribed signed or counterterm-modified precision is used.

% END LAYERED EXPANSION inverse/proper_outer_gaussian_closure.tex

% BEGIN LAYERED EXPANSION background/anisotropic_classical_hierarchy.tex
\subsection{An actual diagonal deformation of the homogeneous backgrounds}
\label{ch:ach-anisotropic-classical}

A change of a plaquette weight changes the principal part of the
background equation. Its derivative contains second differences of
the potential. We therefore construct the deformed background in a
finite-$p$ second-order space before differentiating the hierarchy.
The parameter in this section is a constant diagonal weight, not a
spatially varying source.

Fix the odd blocking integer $L$, coarse spacing one, fine spacing
$h=L^{-j}$ and a compatible periodic volume. Real flat references
with arbitrary holonomy are permitted. Use their parallel frames and
physical pairings as in \cref{ch:hb-strong-inverse}. For
$a=(a_1,\ldots,a_4)$ near $1$ set
\begin{equation}\label{ach:actual-weighted-action}
 \mathscr A_h^a(U)=\sum_{\mu<\nu}a_\mu a_\nu
       \sum_x\left(1-\tfrac12\operatorname{tr}U_{\mu\nu}(x)\right).
\end{equation}
No factor depending on $h$ is inserted into this four-dimensional
Wilson action. Its normal-coordinate physical gradient is denoted
$J_h^a(A)$ when $U=e^{ihA}W_0$. The actual group mean, its
straightening $q_h$, and its node factor $v_h$ are unchanged.

\subsubsection{A second-order inverse with physical cell norms}

Let $C_y$ be unit cells, with a fixed enlargement $C_y^+$ for
difference stencils. For $p>4$ put
\begin{equation}\label{ach:physical-Lp-spaces}
 \begin{split}
 \|f\|_{\mathcal F_p}&=\sup_y\|f\|_{L^p(C_y^+)},\\
 \|A\|_{\mathcal W_p}&=\sup_y
       \sum_{r=0}^2\|D^r A\|_{L^p(C_y^+)}.
 \end{split}
\end{equation}
All mixed forward differences and the physical measure $h^4$ are
included. Any of the fixed larger cells gives an equivalent norm.
Uniform discrete extension and Morrey estimates imply
\begin{equation}\label{ach:Morrey}
 \|A\|_{X_\alpha}\le C_{p,\alpha}\|A\|_{\mathcal W_p},
                       \qquad 0<\alpha<1-4/p.
\end{equation}
One direct verification is to use the piecewise multilinear extension
of $DA$ on each shifted grid. Its $W^{1,p}$ norm is bounded by the
grid $W^{1,p}$ norm, and its nodal values are exact. The continuum
Morrey estimate on an enlarged unit cube bounds the nodal increments.
Applying the same argument to $A$, and summing the finitely many
staggered components, proves \eqref{ach:Morrey}. It has no mesh or
volume factor. Local frames on a fixed periodic cover deal with cells
that cross the fundamental boundary; unitary deck factors preserve
the estimates.

\begin{lemma}[The actual joint inverse on local $L^p$ forces]
\label{lem:ach-Lp-joint-inverse}
For the actual flat rows and $H_0=\mathrm{d}^\sharp \mathrm{d}+DRD^\sharp$, the system
\begin{equation}\label{ach:Lp-joint-system}
              H_0 A-Q^\sharp\omega=f,\qquad QA=b
\end{equation}
has a unique solution satisfying
\begin{equation}\label{ach:Lp-joint-bound}
 \|A\|_{\mathcal W_p}+\|\omega\|_\infty
          \le C_p(\|f\|_{\mathcal F_p}+\|b\|_\infty).
\end{equation}
The same bound holds with a sufficiently small exponential weight on
the cell norms. Constants are independent of depth, volume and real
flat holonomy.
\end{lemma}
\begin{proof}
Retain $K_0=H_0+Q^\sharp Q$ and
$S=QK_0^{-1}Q^\sharp$. The energy cell kernels of $K_0^{-1}$ and
the exponential coarse inverse of $S$ were proved in
\cref{prop:hb-joint-strong-inverse}. A force in $\mathcal F_p$,
$p>2$, has bounded local $L^2$ norm. Resolvent energy bounds between
unit cells and their exponential sums therefore give
\[
 \|K_0^{-1}f\|_{2,\mathrm{cell},\infty}
                 +\|QK_0^{-1}f\|_\infty\le C_p\|f\|_{\mathcal F_p}.
\]
For the second term, each physical row is supported in a fixed unit
star and has uniformly bounded $L^2$ density. There is no point
evaluation of an energy function. The exact formulas
\[
 \sigma=S^{-1}(b-QK_0^{-1}f),\quad
 A=K_0^{-1}(f+Q^\sharp\sigma),\quad\omega=\sigma-b
\]
bound the multiplier and the cell $L^2$ norm of $A$.

The component scalar equation is exactly
\begin{equation}\label{ach:scalar-equation}
 LA=f+DPD^\sharp A-Q^\sharp QA+Q^\sharp\sigma,
                   \qquad P=1-R.
\end{equation}
The actual factorization
$DPD^\sharp=(DF)M_s^{-1}(DF)^\sharp$ in
\eqref{hb:projector-strong} maps cell $L^2$ to $L^\infty$.
This uses the scalar smoothing of its columns already proved there.
The actual row density gives
$\|Q^\sharp c\|_\infty\le C\|c\|_\infty$; the division by $h^4$
in the physical transpose cancels the row coefficient $O(h^4)$.
Thus all terms in \eqref{ach:scalar-equation} other than $f$ are
bounded pointwise by the asserted right side of
\eqref{ach:Lp-joint-bound}. Here the local finite-$p$ estimate
\cref{dcz:interior-Lp} has outer $W^{1,p}$ terms, so an explicit
nested bootstrap is needed. The energy kernels first bound the
outer $H^1$ norm. The scalar equation and its local $p=2$ estimate
give $W^{2,2}$; the strict four-dimensional embeddings give
$DA\in L^3$ and $A$ in every fixed finite $L^q$. Its next $p=3$
estimate gives $W^{2,3}$, hence $DA\in L^8$ and
$A\in L^\infty$. If $p\le8$ the final $p$ estimate now applies.
If $p>8$, first apply the $p=8$ estimate to obtain
$DA\in L^\infty$, and then apply the desired finite $p$ estimate.
Choose a fixed enlargement large enough for all these successive
cutoffs. All other sources in the scalar equation were already
bounded; $f\in L^p$ belongs to each lower exponent used on these
fixed cells. This proves the field bound without an endpoint
singular-integral estimate.

For the weighted assertion use the existing exponential energy
kernels first and the scalar estimate separately on each cell.
Weights on neighboring enlarged cells differ by a fixed factor.
No derivative is taken of a truncated distance weight, and no
pointwise estimate is inferred directly from a cell $L^2$ estimate.
Uniqueness is the same $K_0,S$ energy argument as before.
\end{proof}

Let $K_a=\mathrm{d}^\sharp W_a \mathrm{d}$, where $(W_aF)_{\mu\nu}=a_\mu a_\nu
F_{\mu\nu}$. Constant weights give the elementary bound
\begin{equation}\label{ach:principal-perturbation}
 \|(K_a-K_1)A\|_{\mathcal F_p}
            \le C|a-1|\|A\|_{\mathcal W_p}.
\end{equation}
Indeed this is a sum of finitely many second differences with
coefficients $a_\mu a_\nu-1$. This is the place where a mere
$X_\alpha$ solution space would not suffice.

\subsubsection{The weighted equation is solved, rather than postulated}

The complete normal Wilson gradient has the decomposition
\begin{equation}\label{ach:complete-residual}
                   J_h^a(A)=K_aA+N_h^a(A).
\end{equation}
On a fixed complex $X_\alpha$ ball and a fixed weight polydisc,
\begin{equation}\label{ach:residual-bounds}
 \|N_h^a(A)\|_\infty\le C\|A\|_{X_\alpha}^2,
 \quad
 \|N_h^a(A)-N_h^a(\widetilde A)\|_\infty
                          \le Cr\|A-\widetilde A\|_{X_\alpha}.
\end{equation}
These are consequences of the actual plaquette proof of
\cref{lem:hb-Wilson-residual}, applied separately to each of the six
orientations before summing. To make the parameter dependence
explicit, the four-factor logarithm is
$F=\mathrm{d}A+E_h(A)$ with $DE_h=O(A\,DA)$; the sine remainder has rough
one-difference bound $O(h^3\|F\|_\infty^3)$. Each is multiplied by
the constant $a_\mu a_\nu$. The normal-to-left differential is
$1+O(hA)$, and its product with $K_aA$ is bounded by
$Ch\|A\|_\infty h^{-1}\|DA\|_\infty$.
Thus no unbounded second difference occurs in the nonlinear
remainder. These finite expressions are jointly holomorphic in
$A,a$ on a larger fixed ball; Cauchy estimates there give the
Lipschitz bound in \eqref{ach:residual-bounds}.

Set $C_h=q_h-Q$, and let $\mathcal L_0^{-1}$ be the inverse of
\cref{lem:ach-Lp-joint-inverse}. The actual penalized equation is
equivalent to the fixed point
\begin{equation}\label{ach:actual-fixed-point}
 \binom A\omega
 =\mathcal L_0^{-1}
 \binom{-(K_a-K_1)A-N_h^a(A)+(q_h'(A)-Q)^\sharp\omega}
             {B-C_h(A)}.
\end{equation}
The constraint in this equation is the entire actual recursion.
In particular its physical transpose uses the fixed-input-leg
$C_Lh^4$ moment of \eqref{hfm:hierarchical-fixed-leg-density}; no directional Hessian
estimate is substituted for a column-density estimate. Its
nonlinear terms have the bounds \eqref{hb:nonlinear-row-Lipschitz}.

\begin{proposition}[Holomorphic anisotropic critical backgrounds]
\label{prop:ach-actual-anisotropic-background}
Fix $p>4$ and $0<\alpha<1-4/p$. There are $\varepsilon,r,\rho>0$,
depending on these choices and $L$ but not on $j$ or the volume,
such that for $|a-1|<\varepsilon$, $\|B\|_\infty<\rho$, the actual
mean fibre has a holomorphic critical orbit $U_j^a(B)$. A unique
mean-slice representative and its full multiplier satisfy
\begin{equation}\label{ach:background-bound}
 \|A_j^a(B)\|_{\mathcal W_p}+\|\omega_j^a(B)\|_\infty
                                       \le C_p\|B\|_\infty.
\end{equation}
The prescribed literal representative has actual mean $e^{iB}V_0$.
For real $a,B$ the field is real and is the local constrained
minimum in a fixed strong chart, modulo gauges equal to one at the
coarse nodes. The orbit and total action descend across the
overlapping real-flat charts of \cref{prop:frd-flat-atlas-descent}.
\end{proposition}
\begin{proof}
In the joint norm $\|A\|_{\mathcal W_p}+\|\omega\|_\infty$, the
Lipschitz constant of \eqref{ach:actual-fixed-point} on radius $r$
is at most $C_p|a-1|+C_{p,L}r$, by
\eqref{ach:Morrey} and \eqref{ach:principal-perturbation} and \eqref{ach:residual-bounds}.
Choose $\varepsilon$ first so the first term is at most $1/8$,
then $r$ so the second term is at most $1/8$, and finally
$C_p\rho<r/4$. The fixed point is invariant and contractive.
The iteration from zero converges locally uniformly on the fixed
complex product ball. It constructs the jointly holomorphic pair
and proves \eqref{ach:background-bound}; it does not assume the
desired background or its inverse.

The fixed point first solves the equation with its gauge penalty
retained. The actual identities
\eqref{hb:effective-coarse-gauge} and \eqref{hb:exact-nonlinear-intertwining}
still hold because the group mean has not changed. For
$\tau_A\lambda=0$, gauge invariance of the weighted action gives
$\langle J_h^a(A),\Gamma_A\lambda\rangle=0$, while the constraint
force also vanishes in that test. The site construction used in
the penalty-removal proof of \cref{prop:hb-actual-critical-orbit}
makes $\lambda\mapsto RD^\sharp\Gamma_A\lambda$ onto
$\operatorname{ran}R\cap Z_\alpha$. Its correction uses the actual
$\tau_A$, the unchanged site inverse $G_sR$, and the smooth $E_q$.
Our field lies in its required $X_\alpha$ ball by
\eqref{ach:Morrey}. Therefore the remaining penalty annihilates
every vector in that range and is zero. For complex fields this
is a range and bilinear nondegeneracy argument, not an inference
from a possibly zero complex square. All tests are finite-array
gauge directions. They need not satisfy a uniform
$\mathcal W_p$ norm for this exact identity to apply.

The literal node correction uses sampled smooth profiles with
uniform bounds of every fixed order, as in
\eqref{hb:node-extension}. Its node values are $O(B)$ and its
logarithm has bounded local $W^{3,p}$ norm. The discrete product
and exponential rules hence preserve the field $W^{2,p}$ bound.
It gives the claimed actual literal mean by the exact
plain/straightened factorization, not by equating their derivatives.

For real weights in this smaller ball,
$W_a\ge cI$. On the linear slice and $\ker Q$, the weighted
quadratic form plus the unchanged gauge penalty is a relative
$O(|a-1|)$ perturbation of the proved positive flat form. The
actual nonlinear Hessian and multiplier terms add a relative
$O(r)$ form error. The constrained chord argument in the
isotropic proof therefore retains a positive lower bound. This
proves the local minimum and local orbit uniqueness.

For clarity, uniqueness is also valid against an admitted
$X_\alpha$ competing critical field. Put it into the actual mean
slice by the existing site construction. Its finite-array equation,
\eqref{ach:scalar-equation} and the constant principal perturbation
give
$\|A\|_{\mathcal W_p}\le C_p(\|B\|_\infty+r^2)
 +C_p\varepsilon\|A\|_{\mathcal W_p}$.
The competing multiplier is bounded by the same smooth-column
weak readout used in \eqref{hb:competing-multiplier-bound}; the
weighted Wilson form changes its constant only by a fixed factor.
Absorb the last term, shrink $r$, and use the contraction's larger
uniqueness ball. Thus no extra $W^{2,p}$ assumption on competing
orbits is hidden in this argument. The geometric flat-reference
transitions preserve the equation and the original strong bounds.
The same upgrade and uniqueness identify their total orbits on
overlaps, proving the asserted descent.
\end{proof}

\subsubsection{Localization, affine probes and actual free coercivity}

The preceding inverse also gives a localized analytic response,
including a derivative of the weight. Here the weight is one global
parameter: no decay in a fictitious location of $a$ is claimed.
The flat inverse with $K_a$ in place of $K_1$ is its Neumann
series in the bounded operator $(K_a-K_1):\mathcal W_p\to
\mathcal F_p$. The weighted cell bound in
\cref{lem:ach-Lp-joint-inverse} makes that series convergent at a
fixed positive spatial rate. Smooth multiplication by the unit-cell
partition is bounded on $\mathcal W_p$ by the discrete product
rule through order two. The actual Wilson and constraint residuals
are supported in fixed physical stars and have the same uniform
Banach Cauchy bounds as in \eqref{het:actual-cell-equation}.
Expanding this equation with the $a$-dependent inverse gives the
same ordered rooted trees as there. Their propagator sums are
uniform on a slightly larger weight polydisc. The scalar majorant
for arity at least two converges after decreasing $\rho$ once.
Consequently every fixed-order background jet, and its first $a$
derivative on a smaller polydisc, has exponential cell moments
uniform in $j$. This proves the required response localization
from the actual equation, rather than assigning a small global
inverse to an arbitrary short carrier.

At the trivial flat reference let
$A^F_\mu(x)=\tfrac12\sum_\nu F_{\nu\mu}x_\nu$, with constant
antisymmetric $F$. Then
\begin{equation}\label{ach:affine-response}
 K_aA^F=0,\qquad D^\sharp A^F=0,\qquad Q_hA^F=B^F.
\end{equation}
The first identity differentiates a constant weighted curl; the
second is antisymmetry, and the third is the exact centroid count
of the actual straight row. Exponential inverse kernels act on
polynomial-growth data, so uniqueness of the extended joint inverse
identifies its affine response with $A^F$ and multiplier zero.
The literal response adds only the actual node-frame gradient.
This remains true after one $a$ derivative: the transverse affine
response is independent of $a$. Complete invariant densities with
zero first jet therefore have coefficient-preserving direct and
successive pullbacks, by the argument of
\eqref{abc:successive-affine-response} and \eqref{abc:successive-coefficient-preserved}.
Individual fixed-reference Taylor pieces are not asserted invariant.

For the Taylor vertices and Fourier coefficient assertions below,
take the trivial flat reference, as in the isotropic classical
comparison. The existence and descent theorem above continues to
allow arbitrary real flat holonomy.
Put $S_j^a(B)=\mathscr A_{h_j}^a(U_j^a(B))$ and
$T_{j,m}^a=D_B^mS_j^a(0)$. Positivity of the actual flux averages
in \eqref{gpl:actual-flux-contraction} gives, for real $a$,
\begin{equation}\label{ach:actual-free-coercivity}
 c\,\mathrm{d}^\sharp \mathrm{d}\le T_{j,2}^a\le CI,
\end{equation}
uniformly in depth. Indeed $W_a\ge cI$ bounds the fine curl energy
below, the literal and straight means have the same curl, and the
smooth exact trial lift bounds it above. The weighted actual
background minimizes that fine energy at its mean. This establishes
the free gap needed for the Gaussian comparison below without
assuming an effective precision.

\subsubsection{Mesh comparison with a genuine parameter derivative}

\begin{proposition}[Anisotropic Taylor convergence and its weight jet]
\label{prop:ach-Taylor-convergence}
For each $0<\sigma<1$ there are a sufficiently small real weight
cube $\mathcal A_\sigma$ about $1$ and $s_\sigma>0$ such that,
for $m=2,3,4$,
\begin{equation}\label{ach:weighted-Taylor-step}
 \sup_{a\in\mathcal A_\sigma}
 \left\{\|T_{j+1,m}^a-T_{j,m}^a\|_{m,s_\sigma}
 +\sum_{\lambda=1}^4
   \|\partial_{a_\lambda}(T_{j+1,m}^a-T_{j,m}^a)\|_{m,s_\sigma}\right\}
                                      \le C_\sigma L^{-\sigma j}.
\end{equation}
The norm is the rooted exponential kernel norm
\eqref{cac:weighted-vertex-norm}. The actual bulk companions and
their first weight jets thus converge in that norm. Their finite
torus versions are the corresponding periodic multilinear kernels.
\end{proposition}
\begin{proof}
Choose $\beta<1$ close enough to one and $p$ so large that
$\beta<1-4/p$. The preceding construction and literal correction
give a uniform $X_\beta$ bound. For real weights the interpolation
and actual mean estimates
\cref{wmc:interpolation,wmc:full-mean-consistency} are unchanged.
The energy estimate is simply the sum of its six orientation
estimates multiplied by $a_\mu a_\nu$. The exact row-seed repair
of \cref{lem:cac-mean-repair} is unchanged as a field construction;
its first-variation bound is multiplied by bounded weights. The
weighted local minimum proved above is applicable in both
directions of the two-mesh trial comparison. Hence, uniformly for
real $a$ in a fixed cube and real $B$ in a fixed smaller ball,
\begin{equation}\label{ach:real-mesh-comparison}
       |S_{j+1}^a(B)-S_j^a(B)|/\mathcal V\le C_\beta h_j^\beta.
\end{equation}
The complex actions on the larger product ball are uniformly
bounded per volume: each plaquette contributes $Ch_j^4$, and
there are $O(h_j^{-4}\mathcal V)$ plaquettes.

Fix $m$ real source profiles of supremum at most one, a real weight
centre $a_*$ in a smaller cube, and one coordinate $\lambda$.
Apply \cref{lem:cac-Chebyshev} to the $m+1$ variables of
\[
 f(t,z)=\mathcal V^{-1}
       (S_{j+1}^{a_*+t e_\lambda}-S_j^{a_*+t e_\lambda})
                         \left(\sum_{i=1}^m z_i b_i\right).
\]
All its real values on the fixed cube are small by
\eqref{ach:real-mesh-comparison}, and the enclosing complex
polydisc is uniformly bounded. The resulting mixed derivative
is bounded by $C h_j^\beta(1+j)^{2(m+1)}$. The version without
$t$ gives the same statement for the undifferentiated vertex.
This is a joint real-to-jet argument; Cauchy's inequality has not
been applied to a real-only bound on $a$.

Cosine and sine profiles give the same bound for the vertex
symbols at real external momenta. The localized background
trees above, the physical plaquette density $h_j^4$, and the
$O(h_j^{-4})$ plaquettes per unit cell give a uniform larger
holomorphic momentum strip for these symbols. Its first weight
jet has the same bound by Cauchy on the larger weight polydisc.
For a fixed $m$ each external distance is dominated by a fixed
multiple of the root distances of the background jets; hence
the exponential moments used here require no all-arity claim.

Apply the iterated strip interpolation argument of
\cref{prop:cac-Taylor-triple-convergence} to each symbol and its
weight derivative. On a sufficiently smaller strip it gives
the power $\theta=(1-s/s_0)^{4(m-1)}$ of the real-momentum bound.
Fourier contour shifts on a slightly larger strip than the
output one bound its absolute exponentially weighted kernel sum.
Choose $\beta$ and then $s>0$ so that $\beta\theta>\sigma$;
the fixed powers of $1+j$ are absorbed in that strict margin.
This proves \eqref{ach:weighted-Taylor-step} for all three orders
at their common smaller rate. The parameter cube and the original
holomorphic radii may depend on these choices, but not on depth.

Finally the differentiated fixed-point trees commute with periodic
summation in their distinct external slots. After choosing a lift
of an internal vertex, each child input image is independently
summed. There is no extra closed-loop image constraint in these
classical rooted trees. Thus the periodic multilinear kernels
are the exact periodizations used above; no assertion that
periodization commutes with ordinary pointwise products is needed.
\end{proof}

The conclusion supplies an actual differentiated classical hierarchy.
It does not construct a proper history from arbitrary curvature-only
coarse data, admit independent enlarged cube parameters, or produce
a spatially varying anisotropy. Those are different domain questions.

% END LAYERED EXPANSION background/anisotropic_classical_hierarchy.tex

% BEGIN LAYERED EXPANSION background/anisotropic_gaussian_comparison.tex
\subsection{The actual anisotropic Gaussian increment}
\label{ch:agc-anisotropic-Gaussian}

We compare the constrained Gaussian of the weighted classical
hierarchy with the weighted massive lattice determinant. Both the
comparison and the convergence of successive increments retain a
weight derivative. These are separate steps: a uniformly bounded
cumulative regulator error alone does not imply convergence of its
successive differences.

Keep the homogeneous domain of
\cref{prop:ach-actual-anisotropic-background}, using the trivial
flat reference for the Taylor coefficients and their bulk symbols. Put
$\mathsf A=\operatorname{diag}(a_\mu)$, acting on the vector index.
For a complete invariant pointed functional, let $b_{\mu\nu}$ be
its quadratic coefficient in the single plane $\{\mu,\nu\}$,
in the physical convention of \cref{amc:weighted-massive-law}.
At fixed diagonal weights there are six such coefficients. Axis
permutations also permute the weights; scalar isotropy is not
assumed at $a\ne1$.

\subsubsection{The weighted physical gauge and its uniform gap}

At the actual background $U=U_j^a(B)$, re-center group coordinates
at zero and let $q=q_U,Q=Q_U,D=D_U$ be the actual re-centered rows
and gradient. Thus $qq^\sharp=1$ and $QD=\mathrm{d}_cq$. These site rows
are not the effective fixed-reference map $\tau_A$ away from zero.
Define
\begin{equation}\label{agc:weighted-precisions}
 \begin{aligned}
 \Delta_a&=D^\sharp\mathsf A D,&
 K_a&=(\mathscr A_h^a)''(U)+\mathsf A DD^\sharp\mathsf A,\\
 \widetilde K_a&=\mathsf A^{-1/2}K_a\mathsf A^{-1/2},&
 V_s&=q^\sharp q,\\
 V_v&=\mathsf A^{-1/2}Q^\sharp Q\mathsf A^{-1/2}
       -\mathsf A^{1/2}Dq^\sharp qD^\sharp\mathsf A^{1/2}.
 \end{aligned}
\end{equation}
The scalar averaging precision is $\Delta_a+V_s$. In original
vector coordinates it is
$(\mathscr A_h^a)''+\mathsf A D(1-q^\sharp q)D^\sharp\mathsf A
+Q^\sharp Q$; its congruent version is $\widetilde K_a+V_v$.
The gauge is $(1-q^\sharp q)D^\sharp\mathsf A$, with scalar
Faddeev--Popov operator $\Delta_a$ on the corresponding constrained
site variables. These choices are tied to one another.

The real positive branches of $\mathsf A^{1/2}$ extend
holomorphically to a fixed complex polydisc about $1$. Their
coordinate Jacobian is field independent and disappears on vacuum
subtraction. In the original vector coordinates the massive
regulator is $s\mathsf A$, not $sI$. After congruence it is $sI$,
as in \eqref{amc:congruent-gauge}.

\begin{lemma}[Uniform weighted averaging and interpolation inverses]
\label{lem:agc-weighted-gap}
For real $a$ and real actual backgrounds on a fixed sufficiently
small domain, both averaging precisions are uniformly comparable
to the free physical $H^1$ form. The same holds for
\begin{equation}\label{agc:interpolation}
 M=K_*+1,\qquad S_t=M+t(V-1),\quad0\le t\le1,
 \quad(K_*,V)=(\Delta_a,V_s)\ \hbox{or}\ (\widetilde K_a,V_v).
\end{equation}
On a smaller complex weight/background domain these normalized
operators have uniform relative-form inverses. Their weighted
Bloch inverses are holomorphic in a fixed external momentum strip.
No positivity assertion is made at complex parameters.
\end{lemma}
\begin{proof}
At $a=1$ and a flat reference these are the background-gauge and
site gaps of \cref{prop:mrc-BG-gap}. Their upper forms have the
same $H^1$ bound. Since the weights are constant,
\[
 |\langle v,(K_a-K_1)v\rangle|
             \le C|a-1|\|v\|_{H^1}^2
\]
at a flat reference; the corresponding statement for $\Delta_a$
is immediate. The $Q$ and $qD^\sharp$ terms have bounded relative
forms, so their weighted changes satisfy the same estimate.
The actual background has uniformly small $X_\alpha$ norm by
\eqref{ach:background-bound} and \eqref{ach:Morrey}. Exact plaquette and
transport comparisons give an additional relative $C\rho$
error, with the actual row densities retained. Taking
$|a-1|+\rho$ sufficiently small preserves both gaps. The massive
endpoint has the same gap; every real interpolation is their
convex combination.

For complex parameters normalize by the fixed free
$(L_h+1)^{-1/2}$. The error from the real flat operator is small
in operator norm, so a Neumann series proves the stated inverse
bound. A small exponential physical cell conjugation changes
each derivative form and each finite-carrier row form by $O(s)$.
Its normalized inverse is again a Neumann series, uniformly in
depth. This is the fixed-strip assertion. It uses relative-form
estimates and the actual re-centered rows, not a complex
orthogonal decomposition.
\end{proof}

\subsubsection{The regulator remainder is bounded with its weight jet}

Let $G=M^{-1}$ and $W=V-1$. The exact integrated resolvent Taylor
identity of \eqref{mrc:exact-resolvent-Taylor} reads
\begin{equation}\label{agc:exact-Taylor}
 \begin{split}
 \log\det(K_*+V)-\log\det(K_*+1)
       +\operatorname{Tr}G+\tfrac12\operatorname{Tr}G^2
  ={}&\operatorname{Tr}GV+\operatorname{Tr}G^2V
                    -\tfrac12\operatorname{Tr}GVGV\\
    &+\int_0^1t^2\operatorname{Tr}
                  (GWGW S_t^{-1}W)\,dt.
 \end{split}
\end{equation}
It holds first in finite dimensions, with coherent log branches;
the polarized bulk traces below then give its volume limit.
Use the \emph{fixed isotropic} free comparison
$G_0=(L_h+1)^{-1}$, with the identity in colour and vector indices,
and put
\begin{equation}\label{agc:fixed-normalization}
 B=G_0^{1/2}MG_0^{1/2},\quad
 B_t=G_0^{1/2}S_tG_0^{1/2},\quad
 C=G_0^{1/2}VG_0^{1/2},\quad X=C-G_0.
\end{equation}
The integral trace is then
$\operatorname{Tr}(B^{-1}XB^{-1}XB_t^{-1}X)$.
In particular anisotropy has not been hidden in a parameter-dependent
free reference whose differentiated trace might have a large rank.

For the vector case set
$E=Q\mathsf A^{-1/2}G_0^{1/2}$ and
$F=qD^\sharp\mathsf A^{1/2}G_0^{1/2}$. Then exactly
\begin{equation}\label{agc:actual-rank}
                      C=E^\sharp E-F^\sharp F.
\end{equation}
The scalar factorization is
$C=(qG_0^{1/2})^\sharp(qG_0^{1/2})$.
At a real Bloch momentum there are $12$ vector rows and $3$ site
rows per physical coarse cell. The normalized vector $C$ has rank
at most $15$, and the scalar one at most $3$. The physical
transposes and the constant congruence are part of these formulas.

All factors in \eqref{agc:fixed-normalization} and \eqref{agc:actual-rank},
and their first two background jets, are uniformly bounded on a
fixed complex weight polydisc and momentum strip. Here are the
estimates behind that assertion. The exact Wilson Hessian is a
sum of finitely many derivative forms with bounded plaquette
coefficients. Their two background derivatives are controlled
by the $X_\alpha$ bounds and localized jets established above.
Multiplication by $a_\mu a_\nu$ does not change their fine
difference orders. The actual $q,Q$ transports have uniformly
bounded physical path length; their complete column densities
and two differentiated moments are those of
\eqref{hfm:hierarchical-fixed-leg-density}. Also $D'$ is zeroth
order. Normalizing every derivative form on each side by
$G_0^{1/2}$ therefore bounds all these jets without a fine-volume
factor. Holomorphy on a larger weight polydisc gives the same
statements after one $a$ derivative. The diagonal inverse bounds
come from \cref{lem:agc-weighted-gap}.

To bound the polarized trace, use the four-grade algebra recording
the two external momenta $p,-p$, as in the proof of
\eqref{mrc:uniform-Schatten-three}. Its diagonal fine momenta are
$\kappa,\kappa+p,\kappa-p,\kappa$. The normalized matrix inverses
are finite triangular inverses with at most two strict jet steps.
The jet of \eqref{agc:actual-rank} still factors through four
copies of the coarse rows, hence has rank at most $60$ in the
vector case and $12$ in the scalar case. Analytic transposes are
taken before the jet construction; they do not reverse the four
grades. Each term with a $C$ factor thus has bounded trace norm
by rank times the bounded operator norms.

The only remaining terms in the last line of
\eqref{agc:exact-Taylor} have three $G_0$ factors. For a fixed
small complex shift $z$,
\begin{equation}\label{agc:three-resolvents}
 \int_{[-\pi,\pi]^4}\|G_0(\kappa+z)\|_{\mathcal S_3}^3
                        \frac{d\kappa}{(2\pi)^4}
       \le C\int_{\mathbb R^4}(1+c|q|^2)^{-3}\,dq<\infty.
\end{equation}
This is the exact alias partition of the fine Brillouin cube.
Schatten H\"older, followed by H\"older in $\kappa$ for differently
shifted factors, bounds their mixed traces. The bounded normalized
anisotropic factors $B^{-1},B_t^{-1}$ do not affect this estimate.
For real external $p$, cyclic trace identities are applied after
integrating $\kappa$, where their shifts are legitimate torus
reindexings. Holomorphic continuation then proves the identities
and bounds for complex $p$. No complex contour reindexing is
assumed. The bounds are uniform on a larger weight polydisc, so
Cauchy's inequality legitimately bounds one weight derivative
of the whole remainder.

The actual saddle factors must also be retained. For an onto
constraint $C_0$ and $T_*$ equal to the unconstrained site or
background-gauge precision, block elimination gives
\begin{equation}\label{agc:saddle-factors}
 |\det\mathbb S(T_*,C_0)|
  =\det(T_*+C_0^\sharp C_0)
       \det\{C_0(T_*+C_0^\sharp C_0)^{-1}C_0^\sharp\}
\end{equation}
for real parameters. A constant congruence gives the vector
version of \eqref{agc:weighted-precisions}. The actual smooth
row lifts and their $H^1$ energy bound give uniform upper and
lower bounds on these coarse Schur matrices. They have $12$ or
$3$ rows per cell, and the preceding weighted inverse estimates
bound their mixed log derivatives on the same strip. Their
branches are fixed by positivity on the real domain and then
continued. Thus their plane coefficients and first weight jets
are bounded as well. These are coarse finite ranks per fibre;
no assertion of fixed rank for the full finite-volume trace is used.

\subsubsection{Current corrections and coefficient-preserving pullbacks}

The finite physical Gaussian gauge-change identity
\eqref{pgg:physical-RF12} is algebraic for any invariant action.
Apply it to $\mathscr A_h^a$ with the weighted divergence just
used. Its actual multiplier correction is the Hessian term
$-\langle\mu,q''\rangle$ on the straightened fibre. It is not
assumed to be third order in momentum as a matrix.
As in that finite identity, coarse stabilizer factors are canceled
before taking the flat limit or a holomorphic continuation. The
remaining fine and coarse Schur expressions are the gapped ones
estimated above. Thus no logarithm of an uncanceled singular
coarse Faddeev--Popov factor is continued through the flat field.

For completeness its coefficient vanishes for every weight.
Write the joint determinant ratio as $L_a(V,\nu)$, with
$L_a(V,0)=0$. Global colour covariance gives the full first
$\nu$ derivative zero at $(1,0)$. The mixed $V,\nu$ Hessian
$K_a(p)$ obeys its gauge Ward identity
$d(p)^\sharp K_a(p)=0$; analytic expansion at $p=0$ gives
$K_a(0)=0$. The actual multiplier is the derivative of the
classical coarse action. Its linear response $P_a(p)$ is
$O(p^2)$: gauge covariance gives its zero constant term, and
independent coordinate reflections in midpoint coordinates
remove its odd linear term. Constant diagonal weights preserve
these reflections. The quadratic jet of $L_a(V,\nu_a(V))$
is consequently $O(p^3)$, since its terms are a product of
$K_a=O(p)$ and $P_a=O(p^2)$, or two $P_a$ factors.
This proves $b_{\mu\nu}[L_a\circ\nu_a]=0$ for every plane.
The exact identity on the weight polydisc also has zero first
weight derivative. It is the functional argument of
\cref{pgg:current-zero-coefficient}, with its hypotheses now
verified for the weighted action.

The direct affine-response identity
\eqref{ach:affine-response} proves coefficient preservation
for the weighted critical pullback. The successive version
uses actual nesting, not identification with a direct bare
Wilson lift. A full-depth weighted critical field is critical
on every smaller actual mean fibre. Its partial literal mean
lies in the common small chart by the same mean-domain estimates.
Local orbit uniqueness identifies it with the successive
critical construction. Differentiating this identity on an
affine probe gives the intermediate affine field, with its
raw link-length factor, plus an actual gradient. The complete
functional's vanished first jet removes the second derivative
of the critical map; its two-leg Ward identity removes that
gradient. This is precisely the literal/straightened distinction
in \eqref{abc:successive-affine-response} and is unchanged by $a$.

Define the vacuum-subtracted weighted massive physical logarithm
\begin{equation}\label{agc:mass-functional}
 \mathcal L_h^a(s,U)
   =\log\det(\Delta_{a,U}+s)
           -\tfrac12\log\det(\widetilde K_{a,U}+s),\qquad
 (\mathcal L_h^a)^\sharp
    =\mathcal L_h^a(1)-\partial_s\mathcal L_h^a(1)
                         +\tfrac12\partial_s^2\mathcal L_h^a(1).
\end{equation}
Let $\Psi_j^a$ be the actual cumulative physical Gaussian
logarithm of this weighted hierarchy in the canonical contour
gauge. Gauge-change invariance makes it equal to the weighted
background-gauge expression after the physical coarea factors
are included. The preceding trace and Schur estimates give
\begin{equation}\label{agc:uniform-RF13}
 b_{\mu\nu}[\Psi_j^a]
       =\widetilde\beta_a^{\mu\nu,\sharp}(h_j^2)
                                +E_j^{\mu\nu}(a),\qquad
 \sup_{j,a}\bigl(|E_j^{\mu\nu}(a)|
                      +|\partial_aE_j^{\mu\nu}(a)|\bigr)\le C.
\end{equation}
The error is holomorphic on a fixed smaller weight polydisc.
The mass substitution is $z=h_j^2s$ even in the weighted case:
both scalar and congruent vector physical operators are $h_j^{-2}$
times the bare ones. The plaquette area $h_j^2$ squared cancels
the site density $h_j^{-4}$ in its plane coefficient.
The left side of \eqref{agc:exact-Taylor} combines to the
averaging logarithm minus
$\mathcal L_h^a(1)-\partial_s\mathcal L_h^a(1)
+\tfrac12\partial_s^2\mathcal L_h^a(1)$; the signs here follow
from $(\log\det)'=\operatorname{Tr}G$ and
$(\log\det)''=-\operatorname{Tr}G^2$.

\subsubsection{The differentiated Gaussian polarization is continuous}

The actual one-step contour gauge remains
$\sum_x(1-\tfrac12\operatorname{tr}T_x)$.
Its cubic and quartic vertices, the actual plain mean vertices,
Haar terms and coarea convention in
\eqref{gpl:actual-Phi-Hessian} are unchanged. Only the input
Taylor triple $\mathfrak J_j^a=(T_{j,2}^a,T_{j,3}^a,T_{j,4}^a)$
has changed. The real gap \eqref{ach:actual-free-coercivity}
and the fixed contour gap retain uniform full saddle inverses;
replacing the lower constant one by $c>0$ only changes their
constants. The rooted trace proof of
\cref{prop:gpl-polarization-Lipschitz} therefore applies on this
weighted class and its smaller complex neighborhood.

Here is the extra derivative estimate needed, beyond continuity
of the undifferentiated formula. Let $\Pi_j^a$ be the actual
Gaussian polarization. For a weight direction put a dot over
its derivative. If $\mathbb G$ is its full field/multiplier
saddle inverse and $Z$ its actual response column, then
\begin{equation}\label{agc:differentiated-saddle}
 \dot{\mathbb G}=-\mathbb G
       \begin{pmatrix}\dot T_2&0\\0&0\end{pmatrix}\mathbb G,
 \qquad \dot Z=\dot{\mathbb G}\,I_{\rm datum}.
\end{equation}
The augmented $\dot V_3,\dot V_4$ are exactly the differentiated
action vertices. The constraint and literal trace-gauge vertices
are fixed. Differentiate the two rooted contractions in
\eqref{gpl:quartic-contraction} and \eqref{gpl:cubic-contraction}, using
\eqref{agc:differentiated-saddle}, and differentiate the Haar
matrix product. Each resulting expression has one dotted
factor, bounded in the same smaller weighted Schur norm.
Subtracting the formulas for two depths, one factor at a time,
gives
\begin{equation}\label{agc:jet-Lipschitz}
 \begin{split}
 \|\Pi_{j+1}^a-\Pi_j^a\|_{t,2}
 +\|\dot\Pi_{j+1}^a-\dot\Pi_j^a\|_{t,2}
 \le C\bigl(&\|\mathfrak J_{j+1}^a-\mathfrak J_j^a\|_s\\
       +\|\dot{\mathfrak J}_{j+1}^a-\dot{\mathfrak J}_j^a\|_s\bigr).
 \end{split}
\end{equation}
Uniform bounds for the undifferenced dotted triples follow from
their holomorphic weight ball. Rooting the first external leg
before summing the loop contractions removes any volume factor,
exactly as in the original polarization proof. The complete
field and multiplier first derivative of the log determinant
still vanishes by colour, so no omitted $V_3[z'']$ term is
introduced by the weight derivative.

\begin{proposition}[The actual diagonal one-loop law]
\label{prop:agc-actual-anisotropic-increment}
Let $b_{j,\mu\nu}^a$ be the plane coefficient of the actual
successive homogeneous Gaussian increment, in the same
normalization as \eqref{mrc:physical-increment-sum}.
For every $0<\sigma<1$, on a fixed sufficiently small real
weight cube,
\begin{equation}\label{agc:actual-increment-law}
 \begin{split}
 b_{j,\mu\nu}^a
     &=\frac{11}{3\pi^2}
            d(a)a_\mu a_\nu\log L+O_\sigma(L^{-\sigma j}),\\
 \partial_{a_\lambda}b_{j,\mu\nu}^a
     &=\frac{11}{3\pi^2}
       \partial_{a_\lambda}\{d(a)a_\mu a_\nu\}\log L
                                      +O_\sigma(L^{-\sigma j}).
 \end{split}
\end{equation}
In particular, for $a(t)=1+te$ with $\sum_\lambda e_\lambda=0$,
the actual anisotropic increment obeys
\begin{equation}\label{agc:tracefree-actual-law}
 \left.\frac{d}{dt}\right|_0 b_{j,\mu\nu}^{1+te}
      =\frac{11}{3\pi^2}(e_\mu+e_\nu)\log L
                                  +O_\sigma(|e|L^{-\sigma j}).
\end{equation}
The leading derivative in the isotropic weight direction is zero.
\end{proposition}
\begin{proof}
Apply \eqref{agc:jet-Lipschitz} to
\cref{prop:ach-Taylor-convergence}. The plane readout is a
fixed second-momentum derivative, or equivalently a second
absolute moment, and is continuous in the resulting exponential
kernel norm. Hence $b_{j,\mu\nu}^a$ and each first weight jet
have successive differences $O_\sigma(L^{-\sigma j})$.
Their geometric tails converge at that rate to limits, uniformly
on the smaller real weight cube.

Exact physical Gaussian composition and the weighted successive
affine coefficient preservation proved above give
\[
 \sum_{\ell=1}^j b_{\ell,\mu\nu}^a
                         =b_{\mu\nu}[\Psi_j^a].
\]
It can be differentiated at finite $j$ on the common weight
domain. The weighted massive law
\eqref{amc:weighted-sharp-law} and
\eqref{agc:uniform-RF13} show that these two cumulative quantities
are respectively
$2c\,d(a)a_\mu a_\nu j\log L+O(1)$ and its weight derivative
with an $O(1)$ error. Dividing by $j$ and taking Cesaro limits
identifies the already proved successive limits and their jets.
Since $2c=11/(3\pi^2)$, this proves
\eqref{agc:actual-increment-law}. Finally
$\partial_t[d(1+te)(1+te_\mu)(1+te_\nu)]_{t=0}
=e_\mu+e_\nu-\tfrac12\sum e_\lambda$, which gives
\eqref{agc:tracefree-actual-law} and the isotropic statement.
\end{proof}

The result concerns the actual homogeneous Gaussian increments
and their constant diagonal anisotropy. It uses neither an
unproved interacting cutoff symmetry nor a termwise invariant
local activity decomposition. It supplies no independent fine-core
or history-tower admission and no arbitrary proper-domain extension.

% END LAYERED EXPANSION background/anisotropic_gaussian_comparison.tex

% BEGIN LAYERED EXPANSION regular_homogeneous_trajectory.tex
\subsection{A tuned trajectory for the homogeneous regular steps}
\label{hrt:section}

The step estimates have to be solved with
conditions at opposite ends of the scale
interval. The ultraviolet profile is zero,
whereas the coupling is prescribed at the
infrared end. We solve this problem without
requiring the sum of the profile formation
bounds to be small. That sum need not
converge. The summable quantity is the
coupling derivative of the formation map.

Throughout this subsection the integral
operations are the homogeneous regular
steps on the literal buffered domains.
The frozen cutoffs, relative contour gauge,
physical pointings and Gaussian
normalization are the ones in
\cref{prop:lcf-buffered-formation}.
The large-field complementary operations
are not included in the trajectory
constructed here.

\subsubsection{The spaces and exact transport}

Fix the blocking integer and localization
cubes after the rate choices in
\cref{prop:hit-physical-irrelevant-transport}.
Use the same field radius \(r\) and
carrier rate at every scale, with the
larger formation ball and its literal
buffer retained.
Let \(\mathscr B_j\) be the space
of complete pointed analytic families
\(h=(h_z)\) at scale \(j\) which have
zero value, zero first jet and zero
companion curvature coefficient.
They are gauge invariant and covariant
under translations and the hypercubic
group. The norm is the infimum of the
rooted analytic carrier norms over their
compatible local representations.
Include the fixed finite-period
compatibility norm. All families have
their infinite companions and their
finite torus versions on the declared
compatible volumes.

This infimum is used to identify
representations of the same complete
pointed functions. The space is
completed in that norm. An absolutely
summable sequence of representations
gives a representation of its sum;
hence this completion consists of
normally convergent pointed analytic
functions, with the same identities.
In particular no invariance of each
individual local representative is
required.
Write \(H=\sum_z h_z\) for its total
finite-volume function.

Let \(T_{i,k}\) be the physical
transport of
\eqref{hit:actual-transferred-pointing},
with no new Gaussian integration.
For \(i=k\), set \(T_{i,i}=1\).
The common-radius construction gives
fixed \(C_T\ge1\) and \(0<\lambda<1\)
such that
\begin{equation}
       \|T_{i,k}\|_{\mathscr B_i\to\mathscr B_k}
                   \le C_T\lambda^{k-i},\qquad i\le k .
 \label{hrt:transport-bound}
\end{equation}
One may take \(\lambda=L^{-\beta}\)
with the fixed exponent in the
transport theorem. The same bound
holds for the compatibility component
of the norm, with its fixed cube-unit
period exponent.

The exact semigroup identity is
\begin{equation}
                       T_{j,k}T_{i,j}=T_{i,k}.
 \label{hrt:exact-semigroup}
\end{equation}
There are two points in this assertion.
First, all-depth classical nesting
identifies the two literal intermediate
fields up to an old site gauge.
The complete \(h_z\) is invariant,
so its two compositions agree.
Second, the pointing weights themselves
have an exact refinement identity.
For \(\xi=L^{j-k}\) and a point \(t\)
in the intermediate coordinate system,
\begin{equation}
            \sum_z\theta_y(\xi z)\theta_z(t)
                                  =\theta_y(\xi t).
 \label{hrt:hat-refinement}
\end{equation}
In one dimension the right side is
piecewise linear with all knots at
integers, since \(\xi^{-1}\) is an
integer. It equals its linear
interpolant from those nodal values.
Taking the product in four coordinates
proves \eqref{hrt:hat-refinement}.
The periodic sums give the same
identity on compatible tori.
The absolute local bounds justify
interchanging these sums.
Thus \eqref{hrt:exact-semigroup}
is an identity of complete pointed
functions, even if two chosen
filings differ.

Fix a terminal scale \(N\) and
choose \(\lambda<q<1\). Define
an equivalent norm at scale \(j\) by
\begin{equation}
 \|h\|_{j,*}
       =\sup_{0\le n\le N-j}
                        q^{-n}\|T_{j,j+n}h\|_{\mathscr B_{j+n}}.
 \label{hrt:adapted-norm}
\end{equation}
Then
\begin{equation}
 \|h\|_{\mathscr B_j}\le\|h\|_{j,*}
                    \le C_T\|h\|_{\mathscr B_j},\qquad
       \|T_{j,j+1}h\|_{j+1,*}\le q\|h\|_{j,*}.
 \label{hrt:one-step-contraction}
\end{equation}
The last inequality follows directly
by shifting \(n\) to \(n+1\) in
\eqref{hrt:adapted-norm} and using
\eqref{hrt:exact-semigroup}.
Its constants do not depend on
the horizon \(N\). This removes
the harmless prefactor \(C_T\)
from a one-step induction.

\subsubsection{The exact step and its two formation derivatives}

Let \(\mathcal N_j\) be the
vacuum-subtracted physical Gaussian
pointing for the step from \(j\)
to \(j+1\). Put
\[
 b_j=b[\mathcal N_j^\infty],\qquad
 n_j=\mathcal N_j-b_je_{j+1},
\]
where \(e_{j+1}\) is the actual
classical pointing in
\eqref{hit:actual-classical-pointed-density}.
The Gaussian comparison and the
canonical input bounds give
\begin{equation}
 \|n_j\|_{j+1,*}\le B_n,\qquad
 |b_j-b_\infty|\le B_b\lambda_b^j,\qquad
 b_\infty=\frac{11}{3\pi^2}\log L>0,\quad 0<\lambda_b<1 .
 \label{hrt:Gaussian-data}
\end{equation}
In particular the individual \(b_j\)
are bounded. Positivity of every
initial \(b_j\) is unnecessary.

The scalar power in the physical
Gaussian mismatch is kept in the
vacuum normalization. Consequently
the nonconstant \(\mathcal N_j\)
and \(b_j\) do not acquire a
coupling derivative from that
power. The remaining dependence
is in the formation map
\(\mathcal F_j(x,h)\).
For a fixed frame value
\(\theta_j\), its cutoff and
fluctuation scaling are frozen.
Write
\[
       \epsilon(t)=t^{-1}(\log t)^m,\qquad
                    m=2p>1,
\]
after increasing the coupling floor.
The Gaussian cutoff tail is absorbed
in this bound; \(p>1/2\) makes
\(\exp[-c(\log t)^{2p}]\)
smaller than every fixed negative
power of \(t\).

For every fixed profile bound \(H_*\),
the buffered formation theorem gives,
on slightly larger profile and
coupling domains,
\begin{equation}
 \begin{split}
 \|\mathcal F_j(x,h)\|&\le C_{H_*}\epsilon(\theta_j),\\
 \|D_h\mathcal F_j(x,h)\|&\le C_{H_*}\epsilon(\theta_j),\\
 \|\partial_x\mathcal F_j(x,h)\|
                      &\le C_{H_*}\epsilon(\theta_j)/\theta_j .
 \end{split}
 \label{hrt:formation-bounds}
\end{equation}
Here \(|x/\theta_j-1|\) is smaller
than a fixed constant and
\(\|h\|_{j,*}<H_*+1\).
All norms on the right concern
the complete compatible family.
The constants include the fixed
equivalence in
\eqref{hrt:one-step-contraction}.
They are uniform in \(j,N\)
and volume.

Let \(\mathcal P_{j+1}\) subtract
the companion curvature coefficient
times \(e_{j+1}\), and put
\[
 r_j(x,h)=b[\mathcal F_j(x,h)^\infty],
       \qquad f_j(x,h)=\mathcal P_{j+1}\mathcal F_j(x,h).
\]
The coefficient readout and this
projection are bounded on the
fixed-radius spaces, by the
complete two-leg Ward moment.
Thus \eqref{hrt:formation-bounds}
also holds for \(r_j,f_j\),
with changed constants.

Up to its positive scalar vacuum
factor, the actual regular step is
\begin{equation}
 e^{-xS_j+H}\ \longmapsto\
 e^{-xS_{j+1}+H\circ p_j^{\rm lit}
                         +\mathcal N_j+\mathcal F_j(x,h)} .
 \label{hrt:exact-regular-step}
\end{equation}
Extracting the two curvature
coefficients gives the exact
recursion
\begin{equation}
 \begin{split}
 x_{j+1}&=x_j-b_j-r_j(x_j,h_j),\\
 h_{j+1}&=T_{j,j+1}h_j+n_j+f_j(x_j,h_j).
 \end{split}
 \label{hrt:exact-recursion}
\end{equation}
The coupling is independent of
the torus because its readout
uses the infinite companion.
The finite-period corrections
remain in \(h_{j+1}\).
No finite quantum determinant
has been replaced by the
determinant on a periodic cover.

\subsubsection{Forward stability for a prescribed coupling sequence}

Choose
\begin{equation}
                  H_*\ge 1+\frac{2B_n}{1-q}.
 \label{hrt:profile-floor}
\end{equation}
Fix a sufficiently large real
number \(X\), and define the
nominal frame
\begin{equation}
             a_j=X+\sum_{\ell=j}^{N-1}b_\ell,\qquad a_N=X .
 \label{hrt:nominal-frame}
\end{equation}
The summable errors in
\eqref{hrt:Gaussian-data} imply
\begin{equation}
 |a_j-\{X+b_\infty(N-j)\}|\le B_a,\qquad
       c(X+N-j)\le a_j\le C(X+N-j)
 \label{hrt:nominal-envelope}
\end{equation}
once \(X\ge2B_a+2\).
Freeze \(\theta_j=a_j\)
throughout the construction.
For a fixed small \(\delta>0\)
consider coupling sequences in
\[
          \mathcal D_N
          =\{x:\ |x_j-a_j|\le\delta a_j,\quad 0\le j\le N\}.
\]
Use its slightly enlarged open
polydisc for holomorphic estimates.
The finite forward recursion
in \eqref{hrt:exact-recursion},
with \(h_0=0\), defines \(h_j(x)\).
For \(X\) large enough,
\begin{equation}
                         \|h_j(x)\|_{j,*}\le H_*
                         \quad(0\le j\le N).
 \label{hrt:profile-invariance}
\end{equation}
Indeed the next norm is at most
\(qH_*+B_n+C_{H_*}\epsilon(a_j)\).
By \eqref{hrt:profile-floor},
the first two terms leave at
least \((1-q)H_*/2\) of room,
and the last term fits there
when the coupling floor is fixed.
This is an induction in the
complete compatible family,
so the gauge, symmetry and
periodic identities are retained.

The dependence on the coupling
sequence has a stronger bound
than the profile size itself.
Put
\[
                  k_j=\epsilon(a_j)/a_j,\qquad
                         D=\|x-\widetilde x\|_\infty .
\]
For the two forward solutions,
\eqref{hrt:formation-bounds} gives
\[
 d_{j+1}\le
       \{q+C_{H_*}\epsilon(a_j)\}d_j+C_{H_*}k_jD,
       \qquad d_j=\|h_j(x)-h_j(\widetilde x)\|_{j,*}.
\]
Choose \(X\) so the bracket is
at most some \(q_1<1\).
For earlier \(i<j\), the envelope
in \eqref{hrt:nominal-envelope}
and monotonic decrease of
\((\log t)^m/t^2\) for large \(t\)
give \(k_i\le Ck_j\).
The bounded error \(B_a\) in
that envelope changes only
the fixed constant. Hence
\begin{equation}
 d_j\le C D\sum_{i<j}q_1^{j-1-i}k_i
                              \le Ck_jD .
 \label{hrt:profile-coupling-slope}
\end{equation}
No sum of \(\epsilon(a_j)\)
over all scales appears in
this stability estimate.

\subsubsection{The backward coupling contraction}

The terminal condition is \(x_N=X\).
Summing the first recursion in
\eqref{hrt:exact-recursion}
backward suggests the map
\begin{equation}
  (\Phi x)_j=a_j+
       \sum_{\ell=j}^{N-1}
                         r_\ell(x_\ell,h_\ell(x)),\qquad
                         (\Phi x)_N=X .
 \label{hrt:backward-map}
\end{equation}
We prove it is an invariant
contraction on \(\mathcal D_N\).

The following two elementary
estimates are uniform in \(N,j\):
\begin{equation}
 \begin{split}
 \sum_{\ell=j}^{N-1}\epsilon(a_\ell)
                    &\le C\{1+(\log a_j)^{m+1}\},\\
 \sum_{\ell=0}^{N-1}k_\ell
                    &\le C\frac{(\log X)^m}{X}.
 \end{split}
 \label{hrt:two-scale-sums}
\end{equation}
To prove them, use
\eqref{hrt:nominal-envelope}
to compare with the integer
sequence \(X+c n\), allowing
fixed constant changes in its
argument. The first integral
is bounded by
\(\int_X^{Ca_j}(\log t)^m\,dt/t\).
For the second, integration
by parts bounds
\(\int_{cX}^\infty(\log t)^m\,dt/t^2\)
by \(C(\log X)^m/X\).
For a noninteger \(m>1\), the
same integration by parts
uses \(m(\log t)^{m-1}/t^2\);
choosing \(X\) large absorbs
it in half of the original
integral. The finite first
terms are covered by the
same bounds.

The first line and
\eqref{hrt:formation-bounds}
give
\[
          |(\Phi x)_j-a_j|
               \le C_{H_*}\{1+(\log a_j)^{m+1}\}.
\]
Since \((1+\log^{m+1}t)/t\)
tends to zero, this is at
most \(\delta a_j/2\) for
all \(a_j\ge cX\), once
\(X\) is fixed large enough.
Thus the backward map is
strictly inside its relative box.

For two sequences, the
coupling and profile slopes give
\[
 \begin{split}
 |r_\ell(x_\ell,h_\ell(x))
       -r_\ell(\widetilde x_\ell,h_\ell(\widetilde x))|
 &\le Ck_\ell D+
           C\epsilon(a_\ell)d_\ell\\
 &\le Ck_\ell D .
 \end{split}
\]
The second line of
\eqref{hrt:two-scale-sums}
therefore proves
\begin{equation}
 \|\Phi x-\Phi\widetilde x\|_\infty
       \le K(X)\|x-\widetilde x\|_\infty,\qquad
       K(X)\le C\frac{(\log X)^m}{X}<\frac12 .
 \label{hrt:backward-contraction}
\end{equation}
For each finite \(N\), the
relative box is a complete
closed subset of the finite
supremum-norm space.
Its diameter may grow with
\(N\); neither its contraction
constant nor the floor does.

\begin{proposition}\label{prop:hrt-regular-trajectory}
For every sufficiently large
infrared inverse coupling \(X\)
and every finite horizon \(N\),
the actual homogeneous regular
recursion
\eqref{hrt:exact-recursion}
has a unique trajectory in
the displayed weak-coupling
box with
\[
                         h_0=0,\qquad x_N=X .
\]
It is real for real data,
has \(\|h_j\|_{j,*}\le H_*\),
and satisfies
\begin{equation}
 \begin{split}
 c(X+N-j)&\le x_j\le C(X+N-j),\\
 x_0&=X+b_\infty N
                    +O\{1+(\log(X+N))^{m+1}\}.
 \end{split}
 \label{hrt:tuned-coupling-law}
\end{equation}
All constants and the coupling
floor are independent of \(N\)
and compatible volume.
The finite-volume profiles
are the actual profiles of
the stated integral operations,
with their compatible companions.
\end{proposition}

\begin{proof}
The invariant contraction has
one fixed point. Its forward
profiles satisfy
\eqref{hrt:profile-invariance}.
Conversely every trajectory
with the two stated end
conditions has those forward
profiles and satisfies
\eqref{hrt:backward-map}.
This proves uniqueness in
the specified box.
The real maps preserve real
sequences, so iteration from
the real nominal frame gives
real trajectories.
The relative box and
\eqref{hrt:nominal-envelope}
give the first line of
\eqref{hrt:tuned-coupling-law}.
The sum of \(b_j-b_\infty\)
is bounded by
\eqref{hrt:Gaussian-data};
the first line of
\eqref{hrt:two-scale-sums}
bounds the sum of the
interaction corrections.
This gives its second line.
Finally the exact identity
\eqref{hrt:exact-regular-step}
at each step identifies
the recursively formed density
and its positive real
normalizer. This identification
uses the actual formation
map, not only the abstract
inequalities of a hypothetical
flow.
\end{proof}

\subsubsection{Analytic tuning in a frozen presentation}

Keep the frame and all its
cutoffs fixed at a real
reference endpoint \(X_0\).
Replace the terminal condition
by \(x_N=X_0+\zeta\), where
\(|\zeta|<cX_0\) for a fixed
small \(c>0\).
The backward map becomes
\[
      (\Phi^\zeta x)_j
         =a_j(X_0)+\zeta+
                  \sum_{\ell=j}^{N-1}r_\ell(x_\ell,h_\ell(x)).
\]
For small enough \(c\),
the same relative box and
slightly larger coupling
discs are invariant.
Its contraction constant
is unchanged.
Iteration converges normally
in \(\zeta\), and the forward
profile maps are holomorphic.
Thus the tuned trajectory is
holomorphic on that disc.

Let a prime denote its endpoint
derivative. The proof of
\eqref{hrt:profile-coupling-slope}
gives
\[
                \|h_j'\|_{j,*}
                          \le Ck_j\sup_\ell|x_\ell'|.
\]
Differentiate the backward
equation and use the same
summable slopes. Then
\[
 \sup_j|x_j'|\le1+K(X_0)\sup_j|x_j'|,
 \qquad
 \sup_j|x_j'-1|
                \le\frac{K(X_0)}{1-K(X_0)} .
\]
After increasing the floor,
\begin{equation}
 \frac34\le x_j'\le\frac54
 \quad\hbox{on a smaller real interval},\qquad
 \|h_j'\|_{j,*}\le Ck_j .
 \label{hrt:endpoint-slope}
\end{equation}
The bare tuning is consequently
strictly increasing in that
real interval.
These derivatives concern one
frozen presentation. They do
not differentiate its sharp
cutoff or recompute it from
the moving endpoint.

The complex bound for
\(x_j'-1\) on a fixed smaller
disc also yields, for every
integer \(s\ge2\),
\begin{equation}
        |x_j^{(s)}|
            \le C_s\frac{(\log X_0)^m}{X_0^s},
 \label{hrt:higher-endpoint-jets}
\end{equation}
uniformly in the horizon.
Apply Cauchy's estimate to
\(x_j'-1\), with disc radius
a fixed multiple of \(X_0\).
Its supremum is
\(O((\log X_0)^m/X_0)\),
and each additional derivative
costs one factor \(X_0^{-1}\).
This proves
\eqref{hrt:higher-endpoint-jets}.

The construction controls the
regular operations on the
stated common domains.
It is a coupling and profile
trajectory, not a bound for
the mass or source moments
of all complementary histories.
Those quantities require the
positive history estimates
used in the full construction.

% END LAYERED EXPANSION regular_homogeneous_trajectory.tex

% BEGIN LAYERED EXPANSION background/regular_trajectory_stability.tex
\subsection{Universality for bounded ultraviolet profiles at a fixed horizon}
\label{rts:section}

We keep one finite horizon, one family of regular maps and one
frozen cutoff presentation. The ultraviolet irrelevant profile
may now be any member of a bounded ball. Retuning the bare
coupling to the same infrared value makes the final profiles
exponentially close. This is a comparison within that same
finite problem; it does not compare different horizons or
construct an infinite-horizon flow.

\subsubsection{The fixed presentation and a bounded initial ball}

Use the complete compatible quotient spaces and physical
regular maps of \cref{prop:hrt-regular-trajectory}. Fix
\begin{equation}\label{rts:rate-order}
                \lambda<q<q_1<q_2<1.
\end{equation}
The adapted norm \eqref{hrt:adapted-norm} makes the exact
transport have norm at most $q$. Its equivalence constant
$C_T$ is independent of the horizon. The ultraviolet space
$\mathscr B_0$ consists of the same gauge-invariant,
hypercubic-covariant, normalized complete pointings, with
zero first jet and zero companion curvature coefficient.
All finite-period versions and their companions are included.
An input with a nonzero marginal coefficient is not part
of the irrelevant ball considered here.

Fix $R_{\rm uv}<\infty$. We shall allow
$\|h_0\|_{\mathscr B_0}<R_{\rm uv}+2$ and prove the
estimates on the smaller ball $\|h_0\|\le R_{\rm uv}$.
Take the common profile ceiling $H_*$ large enough that
\begin{equation}\label{rts:profile-ceiling}
 H_*\ge 2C_T(R_{\rm uv}+2)+1+\frac{2B_n}{1-q}.
\end{equation}
The enlarged formation domains are chosen with this
ceiling. All constants below may depend on this fixed
ball and on the fixed rates, but not on depth, horizon
or compatible volume.

Choose a real reference endpoint $X_0$ above the coupling
floor, and fix the entire presentation at
\begin{equation}\label{rts:frozen-frame}
       a_j=X_0+\sum_{\ell=j}^{N-1}b_\ell,\qquad
                        \theta_j=a_j,\quad a_N=X_0.
\end{equation}
These are scalar inverse-coupling frames, not diagonal
weight parameters. In particular every cutoff,
fluctuation scaling, gauge normalization and map stays
fixed when $h_0$ or the endpoint is varied. Set
\begin{equation}\label{rts:weights}
 e_j=\frac{(\log a_j)^m}{a_j},\qquad k_j=e_j/a_j,
              \qquad e_*=\frac{(\log X_0)^m}{X_0}.
\end{equation}
Increase the floor when necessary. The HRT envelopes
give $a_j\asymp X_0+N-j$ and
\begin{equation}\label{rts:weight-sums}
 \sum_{j<N}k_j\le C\frac{(\log X_0)^m}{X_0},\qquad
 \sum_{j<N}e_jk_j\le Ce_*^2,
 \qquad \sup_j e_jk_j\le Ce_*^2/X_0.
\end{equation}
For the second estimate compare with
$\int_{cX_0}^{\infty}(\log t)^{2m}t^{-3}\,dt$.
One integration by parts and a larger floor bound
this integral by $C(\log X_0)^{2m}/X_0^2$.
The last estimate is the corresponding supremum
bound. We will also use
\begin{equation}\label{rts:weight-comparison}
 k_i\le Ck_j\ (i<j),\qquad
 e_\ell\le Ce_j\{1+(\ell-j)/X_0\}\ (j\le\ell\le N).
\end{equation}
Indeed the bounded Gaussian increments give
$a_j\le a_\ell+C(\ell-j)$, while
$a_\ell\ge cX_0$. For large arguments the logarithmic
ratio is bounded when $a_\ell\le a_j+O(1)$; the
bounded error in the nominal envelope covers the
exceptional initial signs. Thus
$e_\ell/e_j\le C a_j/a_\ell$, which proves the
second inequality. The same envelope proves the
first one by monotonic decrease of $(\log t)^m/t^2$.

On the enlarged fixed domains the exact regular
recursion is
\begin{equation}\label{rts:actual-recursion}
 x_{j+1}=x_j-b_j-r_j(x_j,h_j),\qquad
 h_{j+1}=T_jh_j+n_j+f_j(x_j,h_j).
\end{equation}
The input ceiling above and \cref{prop:lcf-buffered-formation}
give, uniformly there,
\begin{equation}\label{rts:derivative-bounds}
 |r_j|+\|f_j\|\le Ce_j,\quad
 |\partial_xr_j|+\|\partial_xf_j\|\le Ck_j,\quad
 \|D_hr_j\|+\|D_hf_j\|\le Ce_j.
\end{equation}
Norms are the appropriate adapted norms. Fixed
equivalence constants have been included in $C$.
The scalar coefficients use the infinite companion,
so the coupling sequence is independent of the
torus. Finite-period differences remain in the profiles.

\begin{proposition}[Analytic tuning with a nonzero initial profile]
\label{prop:rts-bounded-UV-tuning}
For every fixed $R_{\rm uv}$ there is a sufficiently
large floor $X_0$, independent of $N$, such that
\eqref{rts:actual-recursion} has a unique trajectory
in the HRT relative coupling box with
\begin{equation}\label{rts:mixed-boundary-data}
                    h_0=H,\qquad x_N=X_0+\zeta,
 \quad \|H\|<R_{\rm uv}+2,\quad |\zeta|<cX_0.
\end{equation}
It is jointly holomorphic in $H,\zeta$, real for
real data, and has $\|h_j\|_{j,*}\le H_*$.
At a real endpoint in this presentation its
coupling envelope and bare law are those of HRT,
with constants uniform on the input ball.
\end{proposition}
\begin{proof}
For a prescribed coupling sequence in the fixed
relative box, solve the profile equation forward
starting from $H$. Its first norm is at most
$C_T(R_{\rm uv}+2)$. The induction
\[
 \|h_{j+1}\|_{j+1,*}\le qH_*+B_n+Ce_j\le H_*
\]
is valid by \eqref{rts:profile-ceiling} and a larger
floor. The map is holomorphic on the enlarged
coupling and input domains. For two sequences with
the same initial $H$, its difference obeys
\[
 d_{j+1}\le q_1d_j+Ck_j\|x-\widetilde x\|_\infty,
                          \qquad d_0=0,
\]
after making $Ce_j\le q_1-q$. Therefore
$d_j\le Ck_j\|x-\widetilde x\|_\infty$, exactly
as in \eqref{hrt:profile-coupling-slope}.

The backward map is
\[
 (\Phi^{H,\zeta}x)_j
       =a_j+\zeta+\sum_{\ell=j}^{N-1}
                          r_\ell(x_\ell,h_\ell(x,H)),
             \qquad (\Phi^{H,\zeta}x)_N=X_0+\zeta.
\]
The same sum of $k_j$ as in HRT makes its
supremum-norm contraction constant less than
$1/2$, uniformly in $H$. Its correction from
$a_j$ is at most
$|\zeta|+C\{1+(\log a_j)^{m+1}\}$.
Choose $c$ small and then the floor large so this
is strictly inside the relative box at every $j$.
The fixed-point iteration is invariant and
locally normally convergent in $H,\zeta$. It
constructs the claimed holomorphic tuning map.
No implicit inverse without a uniform bound is
being assumed. Every other trajectory with these
end data has the same forward profile and hence
is the same fixed point. Reality follows by
iteration. The bounded initial profile changes
only the common formation constants in the
proof of \eqref{hrt:tuned-coupling-law}.
\end{proof}

\subsubsection{The exact mixed-end linear inverse}

Linearize along any trajectory in the smaller
complex domains. Write $u_j=\delta x_j$ and
$v_j=\delta h_j$, and allow additive step sources
$\sigma_j\in\mathbb C$, $\tau_j\in\mathscr B_{j+1}$.
The exact linear system is
\begin{equation}\label{rts:linearized-system}
 \begin{split}
 u_{j+1}&=A_ju_j-R_jv_j+\sigma_j,\\
 v_{j+1}&=D_jv_j+F_ju_j+\tau_j,\qquad
                         u_N=\xi,\quad v_0=\eta,
 \end{split}
\end{equation}
where
\[
 A_j=1-\partial_xr_j,\quad R_j=D_hr_j,\quad
 D_j=T_j+D_hf_j,\quad F_j=\partial_xf_j.
\]
Thus $|A_j-1|,\|F_j\|\le Ck_j$,
$\|R_j\|\le Ce_j$, and $\|D_j\|\le q_1$.
All products below use the actual operators;
the symbols $D_j$ here are scale-profile maps,
not lattice derivatives.

For $j\le\ell$ put
\begin{equation}\label{rts:propagators}
 P_{j\ell}=\prod_{t=j}^{\ell}A_t^{-1},\qquad
 U_{ji}=D_{j-1}\cdots D_{i+1}\quad(i<j),
 \qquad U_{i+1,i}=1.
\end{equation}
The empty backward product is one. By
\eqref{rts:weight-sums}, the $A_j$ do not vanish
and $|P_{j\ell}|\le C$. Also
$\|U_{ji}\|\le q_1^{j-1-i}$. Solving one equation
backward and the other forward gives exactly
\begin{equation}\label{rts:mixed-factorization}
 \begin{aligned}
 u_j&=u_j^0+(\mathcal Bv)_j,&
 u_j^0&=P_{j,N-1}\xi-
                  \sum_{\ell=j}^{N-1}P_{j\ell}\sigma_\ell,\\
 v_j&=v_j^0+(\mathcal Ku)_j,&
 v_j^0&=D_{j-1}\cdots D_0\eta+
                         \sum_{i<j}U_{ji}\tau_i,\\
 (\mathcal Bv)_j&=\sum_{\ell\ge j}P_{j\ell}R_\ell v_\ell,&
 (\mathcal Ku)_j&=\sum_{i<j}U_{ji}F_iu_i.
 \end{aligned}
\end{equation}
In the sums with $\ell\ge j$ the upper limit is
$N-1$; at $j=N$ those sums vanish. The first
homogeneous forward product at $j=0$ is one.

For later use define positive scalar kernels
\begin{equation}\label{rts:positive-kernels}
 B_{j\ell}^+=Ce_\ell\mathbf1_{j\le\ell<N},\qquad
 K_{ji}^+=Ck_iq_1^{j-1-i}\mathbf1_{i<j}.
\end{equation}
They majorize the two exact operator kernels.
In unweighted coupling supremum norm their
product satisfies
\begin{equation}\label{rts:unweighted-loop}
 \|B^+K^+\|_{\infty\to\infty}
  \le C\sup_j\sum_{\ell\ge j}e_\ell
                    \sum_{i<\ell}q_1^{\ell-1-i}k_i
  \le C\sum_\ell e_\ell k_\ell
  \le Ce_*^2=:\vartheta<\tfrac12.
\end{equation}
The middle step uses $k_i\le Ck_\ell$.
This is the mixed-end loop. It contains both
formation derivatives. A sum of $e_j$ alone,
which could grow with the horizon, has not
been used as a smallness parameter.

Consequently the exact positive Green envelopes are
\begin{equation}\label{rts:Green-envelope}
 \begin{split}
 |u|&\le \sum_{n\ge0}(B^+K^+)^n
                                 (|u^0|+B^+|v^0|),\\
 |v|&\le |v^0|+K^+
            \sum_{n\ge0}(B^+K^+)^n
                                 (|u^0|+B^+|v^0|).
 \end{split}
\end{equation}
The notation $|v_j|$ means its adapted norm.
Every term is nonnegative and each intermediate
scale is explicitly summed. The series converges
absolutely in the coupling supremum norm; the
profile output then follows from the forward
kernel. This also proves uniqueness of
\eqref{rts:linearized-system} without a sign or
self-adjointness assumption on its complex
coefficients.

These formulas give useful estimates for arbitrary
finite step sources. In fact
\begin{equation}\label{rts:general-source-bound}
 \begin{split}
 \|u\|_\infty\le C\bigg\{&|\xi|+
      \sum_{i<N}|\sigma_i|+e_0\|\eta\|
                         +\sum_{i<N}e_{i+1}\|\tau_i\|\bigg\},\\
 \|v_j\|\le{}&q_1^j\|\eta\|+Ck_j\|u\|_\infty
                          +\sum_{i<j}q_1^{j-1-i}\|\tau_i\|.
 \end{split}
\end{equation}
For the first line, a profile impulse at step $i$
contributes at most
$C\sum_{\ell\ge i+1}e_\ell q_1^{\ell-1-i}
 \le Ce_{i+1}$ to $B^+|v^0|$.
The initial profile contributes $Ce_0$, by
\eqref{rts:weight-comparison}. The scalar direct
sources have exactly their backward sums in
\eqref{rts:mixed-factorization}. Apply
$(1-\vartheta)^{-1}$ to finish. The second line
is the forward equation and $k_i\le Ck_j$.

For example a scalar impulse at $i$ has direct
envelope $C\mathbf1_{j\le i}$; a profile impulse
at $i$ has direct envelope
$q_1^{j-1-i}\mathbf1_{j\ge i+1}$. Their complete
responses are those vectors propagated by
\eqref{rts:Green-envelope}. This explicitly retains
the backward leakage caused by retuning the
terminal coupling. No symmetric exponential
bound for an arbitrary scale impulse is claimed.

\subsubsection{Exponential dependence on the initial irrelevant profile}

For the initial-profile response set
$\xi=0$, $\sigma=\tau=0$ in
\eqref{rts:linearized-system}. Introduce
\begin{equation}\label{rts:exponential-response-norms}
 \|u\|_{E}=\max_j\frac{|u_j|}{e_jq_2^j},\qquad
 \|v\|_{Q}=\max_jq_2^{-j}\|v_j\|.
\end{equation}
The terminal $u_N=0$ causes no difficulty in
this norm. The backward kernel has the bound
\begin{equation}\label{rts:weighted-backward}
 \|\mathcal Bv\|_E
     \le C\|v\|_Q
            \sup_j\sum_{n\ge0}
                 q_2^n\frac{e_{j+n}}{e_j}
     \le C_B\|v\|_Q,
\end{equation}
where sums stop at the finite horizon.
The last inequality follows from
$e_{j+n}/e_j\le C(1+n/X_0)$, not from
a false monotonicity in $j$.
For the forward kernel,
\begin{equation}\label{rts:weighted-forward}
 \|\mathcal Ku\|_Q
     \le C\|u\|_E
        \sup_j\sum_{i<j}q_2^{-1}
                  (q_1/q_2)^{j-1-i}k_ie_i
     \le C_K e_*^2\|u\|_E.
\end{equation}
Here \eqref{rts:weight-sums} even gives an
additional factor $X_0^{-1}$, which is not
needed. The loop norm $C_BC_Ke_*^2$ is less
than $1/2$ after the same fixed floor increase.
Since $\|v^0\|_Q\le\|\eta\|$, solving the two
inequalities proves
\begin{equation}\label{rts:initial-derivative-decay}
             |u_j|\le Ce_jq_2^j\|\eta\|,
             \qquad \|v_j\|_{j,*}\le Cq_2^j\|\eta\|.
\end{equation}
The proof applies uniformly to complex
trajectories on the smaller product domains.
It is an inverse estimate with opposite end
conditions, rather than an estimate obtained
by evolving the scalar coupling forward.

\begin{proposition}[Finite-horizon exponential universality]
\label{prop:rts-exponential-universality}
Let $H,\widetilde H$ lie in the smaller bounded
initial-profile ball, and use the same endpoint,
horizon and frozen presentation. Then
\begin{equation}\label{rts:two-input-universality}
 \begin{split}
 |x_j(H)-x_j(\widetilde H)|
                   &\le Ce_jq_2^j\|H-\widetilde H\|_{\mathscr B_0},\\
 \|h_j(H)-h_j(\widetilde H)\|_{j,*}
                   &\le Cq_2^j\|H-\widetilde H\|_{\mathscr B_0}.
 \end{split}
\end{equation}
In particular the analytic bare tuning map has
$\|D_Hx_0\|\le Ce_0$, and
$\|D_Hh_N\|\le Cq_2^N$.
These constants are independent of the horizon
and compatible volume.
\end{proposition}
\begin{proof}
The jointly holomorphic trajectory can be
differentiated in an arbitrary Banach direction
$\eta$ of the initial profile. Its derivative
solves exactly \eqref{rts:linearized-system}
with the initial-response data just used.
The initial adapted norm costs at most
$C_T\|\eta\|_{\mathscr B_0}$. Apply
\eqref{rts:initial-derivative-decay} and
integrate this derivative along the straight
segment between $H$ and $\widetilde H$.
That segment stays in the retained Banach
ball. This proves the nonlinear difference
bounds, not only stability of a formal
linear recurrence.
\end{proof}

\subsubsection{All initial-profile jets in one frozen endpoint disc}

The derivative estimates above hold on a fixed
larger $H$ ball and a fixed larger disc
$|\zeta|<cX_0$. Banach-valued Cauchy and
polarization on strict inner margins therefore
give, for integers $r\ge1$, $s\ge0$,
\begin{equation}\label{rts:mixed-jets}
 \begin{split}
 \|D_\zeta^sD_H^r x_j\|
            &\le C_{r,s}e_jq_2^j X_0^{-s},\\
 \|D_\zeta^sD_H^r h_j\|_{j,*}
            &\le C_{r,s}q_2^j X_0^{-s}.
 \end{split}
\end{equation}
The constants include the fixed input-ball
margin and the usual finite-order factorials.
One way to obtain them is to apply Cauchy
to the already bounded first $H$ derivative,
taking $r-1$ further profile derivatives
and $s$ endpoint derivatives. Thus every
profile jet retains the exponential factor.
There is no count of its microscopic input
coordinates.

For completeness, pure endpoint derivatives
have the HRT bounds uniformly on this initial
ball. The direct scalar product differs from
one by at most $C\sum k_j$; the profile
feedback contributes at most $Ce_*^2$ by
\eqref{rts:unweighted-loop}. For endpoint
variation $\xi=1$, $\eta=\sigma=\tau=0$,
\eqref{rts:general-source-bound} consequently
gives
\begin{equation}\label{rts:endpoint-first-jets}
 \sup_j|\partial_\zeta x_j-1|
                  \le C\frac{(\log X_0)^m}{X_0},\qquad
                 \|\partial_\zeta h_j\|_{j,*}\le Ck_j.
\end{equation}
On a smaller real interval the first scalar
derivative lies between $3/4$ and $5/4$.
For $s\ge2$, Cauchy applied to
$\partial_\zeta x_j-1$ yields
\begin{equation}\label{rts:endpoint-higher-jets}
 |\partial_\zeta^s x_j|
               \le C_s(\log X_0)^m X_0^{-s},\qquad
 \|\partial_\zeta^s h_j\|_{j,*}
               \le C_sk_jX_0^{1-s}.
\end{equation}
These are derivatives of the single presentation
\eqref{rts:frozen-frame}. They do not differentiate
a sharp cutoff or replace its reference frame
by one recomputed from the moving endpoint.

\subsubsection{The actual finite-volume interaction and law}

For real profiles and the same real endpoint
$X=x_N$, let $\mathcal V$ be a fixed final
coarse volume. The complete finite-volume
interaction is
$\mathcal H_N^H(V)=\sum_z h_{N,z}^H(V)$.
There are $O(\mathcal V)$ unit anchors.
The local analytic norm controls the supremum
of each complete pointing on the working
field ball. Thus
\begin{equation}\label{rts:finite-volume-interaction}
 \sup_V|\mathcal H_N^H(V)-\mathcal H_N^{\widetilde H}(V)|
       \le C\mathcal V q_2^N
                          \|H-\widetilde H\|_{\mathscr B_0}.
\end{equation}
The exact regular-step identity of LCF and
the tuning construction identify these with
the actual finite-volume effective interactions
of the stated regular integral operations.
Scalar vacuum factors may differ and do not
affect a normalized final law.

To state such a comparison precisely, fix a
common real final integration domain $\mathcal D$
inside the working ball and a fixed positive
finite measure $\nu$ there, with nonzero total
mass. Neither depends on the ultraviolet
profile. Product Haar measure restricted to
a fixed smaller real bond ball is one choice.
Define
\begin{equation}\label{rts:IR-probabilities}
 d\mathbb P_H(V)=Z_H^{-1}
       e^{-X S_N(V)+\mathcal H_N^H(V)}
                           \mathbf1_{\mathcal D}(V)\,d\nu(V).
\end{equation}
The real continuous densities have finite
positive normalizers. If
$\delta=C\mathcal V q_2^N\|H-\widetilde H\|$,
then their likelihood ratio is between
$e^{-2\delta}$ and $e^{2\delta}$: the
unnormalized ratio is between $e^{-\delta}$
and $e^\delta$, as is the ratio of the
normalizers. Therefore, with the convention
$\|\mathbb P-\mathbb Q\|_{\rm TV}
=\tfrac12\int|d\mathbb P-d\mathbb Q|$,
\begin{equation}\label{rts:TV-bound}
 \|\mathbb P_H-\mathbb P_{\widetilde H}\|_{\rm TV}
 \le\min\left\{1,\quad
    \exp\bigl(C\mathcal V q_2^N
                  \|H-\widetilde H\|_{\mathscr B_0}\bigr)-1\right\}.
\end{equation}
The fixed multiplicative constant absorbs the
factor two. This is a finite-volume bound on
one common final domain. It is not uniform
over growing volumes, and it proves neither
clustering nor consistency between horizons.
No complementary-history measure has been
introduced in this comparison.

% END LAYERED EXPANSION background/regular_trajectory_stability.tex

% BEGIN LAYERED EXPANSION regular_diagonal_stress.tex
\subsection{Decay of the diagonal stress response in the regular trajectory}
\label{rds:section}

The diagonal one-loop coefficient is useful only
after the irrelevant part of the response has
been included. We derive that coupled response
from the weighted regular step and prove its
decay. The memory term is estimated with the
running coupling weight. An unweighted smallness
bound on that term would not give the result.

Use the actual regular trajectory
\((x_j,h_j)_{0\le j\le N}\) of
\cref{prop:hrt-regular-trajectory}, with
terminal coupling \(x_N=X\).
Its frame, real cutoff and all integral
conventions remain fixed during the response.
Write \(b_j\) for its Gaussian scalar
increment and \(m_j\) for the tracefree
diagonal Gaussian increment in
\eqref{lcf:Gaussian-weight-jet}.
The index in both symbols is the
transfer from \(j\) to \(j+1\).
The actual comparison proves
\begin{equation}
 |b_j-b_\infty|+|m_j-b_\infty|
                  \le C\lambda_0^j,\qquad
 b_\infty=\frac{11}{3\pi^2}\log L>0,\quad 0<\lambda_0<1 .
 \label{rds:two-Gaussian-increments}
\end{equation}

\subsubsection{The stress representative and its exact transport}

Let \(\mathcal E=\{e\in\mathbb R^4:
\sum_\mu e_\mu=0\}\), with its
Euclidean norm. For \(e\in\mathcal E\)
define the complete pointing
\begin{equation}
 E_{j,e,z}(V)=
       \sum_{p\parallel\mu\nu}
          (e_\mu+e_\nu)\theta_z(m_p)
                                      a_p(U_j(V)).
 \label{rds:stress-representative}
\end{equation}
The sum uses the original fine
plaquettes in scale-\(j\) units.
The actual critical orbit here
is isotropic.
These pointings are linear in
\(e\), gauge invariant, and
covariant under simultaneous
axis permutations and permutation
of the entries of \(e\).
They have uniform local norm
and plane readouts
\[
                   b_{\mu\nu}[E_{j,e}]=e_\mu+e_\nu .
 \]
These assertions follow from
the fine density and exact
hat quadrature, with the
actual affine response.

Their total is the derivative
of the weighted classical action:
\begin{equation}
 \sum_z E_{j,e,z}(V)
       =\left.D_aS_j^a(V)[e]\right|_{a=1}.
 \label{rds:total-action-derivative}
\end{equation}
Indeed differentiation of the
critical orbit at fixed literal
mean is tangent to its constraint
fibre. Pairing that derivative
with the physical action gradient
gives zero. Only differentiation
of the plaquette weights remains.
This is a total-action identity.
It does not identify the derivative
of every separately chosen
weighted action pointing.

The particular choice
\eqref{rds:stress-representative}
has exact isotropic transport:
\begin{equation}
                         T_{j,j+1}E_{j,e}=E_{j+1,e}.
 \label{rds:exact-stress-transport}
\end{equation}
All-depth nesting identifies the
original fine plaquette traces
in the two expressions.
The scale change and
\eqref{hrt:hat-refinement}
identify their pointing weights.
Both equalities are applied to
the complete functions before
any local filing. This avoids
an unnecessary irrelevant term
of size equal to the stress
coefficient itself.

\subsubsection{The zero-readout response space}

Use the reflection-invariant
profile space of the weighted
formation theorem, with all six
plane readouts zero.
Let \(\mathscr V_j\) consist
of bounded linear maps
\(v:\mathcal E\to\mathfrak H_j^{\rm ref}\)
which are covariant under the
axis permutation action.
Its norm is the operator norm
in the same compatible quotient
norm as HRT.
The identity
\[
             \mathcal E\ni e\longmapsto
                        (e_\mu+e_\nu)_{\mu<\nu}
 \]
is the stress embedding.

The transport bound used in HRT
also holds on the reflection
space with six zero readouts.
To see this, repeat the actual
slow-field estimate in
\cref{lem:hit-slow-field-sup}.
Reflections kill every cross-plane
curvature pairing. The six
zero readouts therefore make its
entire quadratic curvature form
zero. The remaining proof uses
only that cancellation, complete
gauge invariance, analytic carrier
bounds and the companion property.
Its block interpolation and
anchor sum are unchanged.
Thus the same-radius estimate
and exact complete-function
semigroup hold, with fixed
changed constants.
Use the horizon-adapted norm
on this space as in
\eqref{hrt:adapted-norm}.
We may therefore fix
\begin{equation}
                         \|T_j\|_{j,*\to j+1,*}\le q<1,
                    \qquad T_j=T_{j,j+1}.
 \label{rds:profile-contraction}
\end{equation}
The constants do not depend
on \(N\).

For clarity, the stress readout
of a permutation-covariant linear
response is one scalar.
The stabilizer of the plane
\(\{\mu,\nu\}\) makes the coefficient
of \(e_\mu,e_\nu\) equal, and
makes the two complementary
coefficients equal. The condition
\(\sum e_\mu=0\) consequently
reduces its readout to
\(c(e_\mu+e_\nu)\).
Axis permutations make \(c\)
the same for all planes.
Reflections exclude cross-plane
terms. A scalar isotropic
readout of this response is
zero. This elementary argument
also shows that the scalar
coupling has no first-order
variation in this channel.

\subsubsection{The response equations from the actual step}

At scale \(j\), perturb the exponent
by
\[
                   -t_j\sum_z E_{j,e,z}
                              +\sum_z v_j(e)_z .
 \]
Here \(t_j\) is scalar and
\(v_j\in\mathscr V_j\).
The leading-action perturbation
is obtained by differentiating
the constant weight
\[
                         a=1+(t_j/x_j)e .
 \]
Equation
\eqref{rds:total-action-derivative}
makes this exactly the stated
first-order perturbation of the
actual exponent.
The output leading action is
represented by \(E_{j+1,e}\)
using the same identity.
All subsequent zero-readout
terms are assigned to \(v_{j+1}\).
Thus no equality of individual
weighted action pointings is
being imposed.

Let \(\mathcal P^{\rm ref}\)
subtract the six plane readouts
times the corresponding original
fine-action plane pointings.
It is a bounded linear map on
the fixed-radius spaces.
Differentiate the exact weighted
regular step at \(a=1\) and
the actual pair \((x_j,h_j)\).
Its moving centre is compared with
the fixed isotropic kernel by
\cref{prop:fcw-frozen-weight-jet}.
The actual fixed-kernel jet is the
moving jet minus the rooted sphere
family \(\mathcal B_{j,e}\).
By \cref{cor:fcw-zero-curvature-readout}
this correction has six zero readouts.
The pointings below are the selected
fixed-total pointings of that comparison;
no independent per-anchor filing is
identified from equality of totals.
Equations
\eqref{lcf:weight-formation-API}
and \eqref{lcf:Gaussian-weight-jet}
give the exact linear system
\begin{equation}
 \begin{split}
 t_{j+1}
   &=\left(1-\frac{m_j}{x_j}+d_j\right)t_j+L_jv_j,\\
 v_{j+1}
   &=T_jv_j+\mathcal E_jv_j+\frac{t_j}{x_j}A_j .
 \end{split}
 \label{rds:actual-linear-response}
\end{equation}
The operators on the right
are derivatives of the specified
formation map, not independent
assumptions about a scalar recursion.
Their definitions are as follows.
If the stress readout of
\(D_a\mathbf F_j^a[e]|_1\)
is \(\mu_j(e_\mu+e_\nu)\),
then \(d_j=-\mu_j/x_j\).
The map \(L_j\) is minus the
stress readout of \(D_h\mathbf F_j[v]\).
The map \(\mathcal E_j\) is
\(\mathcal P^{\rm ref}D_h\mathbf F_j\).
Finally
\begin{equation}
 \begin{split}
 A_j[e]={}&\left.D_aT_j^a h_j[e]\right|_1+q_{j,e}\\
         &+\mathcal P^{\rm ref}
                     \left.D_a\mathbf F_j^a[e]\right|_1
                         -\mathcal B_{j,e}.
 \end{split}
 \label{rds:actual-profile-forcing}
\end{equation}
The first term has six zero
readouts by
\eqref{lcf:transport-weight-jet}.
The middle term is the actual
Gaussian remainder in
\eqref{lcf:Gaussian-weight-jet}.
This identifies the full forcing
and its normalization.
The final term is the actual sharp-cutoff
boundary correction. Its compatible
norm is superpolynomially small in the
frozen coupling reference, so it preserves
the stated forcing bound.
The signs in the first line
come from extracting the positive
Gaussian logarithm from a
negative action exponent.

Write
\[
             \epsilon_j=\frac{(\log x_j)^m}{x_j},
                    \qquad k_j=\epsilon_j/x_j .
 \]
The fixed-frame envelopes in HRT
make these quantities comparable
to their versions with \(a_j\).
The actual derivative bounds give
\begin{equation}
 |d_j|\le Ck_j,\qquad
 \|L_j\|+\|\mathcal E_j\|\le C\epsilon_j,\qquad
                         \|A_j\|\le C_{H_*}.
 \label{rds:actual-response-bounds}
\end{equation}
In particular the old profile's
weight derivative is only bounded;
it has no extra power of the
fluctuation coupling. Its factor
\(t_j/x_j\) in
\eqref{rds:actual-linear-response}
is retained.
No spatially varying weight is
used in these identities.

\subsubsection{Eliminating the irrelevant response}

Start with the pure bare stress
perturbation \(v_0=0\).
Choose the infrared floor so
\(q+C\sup_j\epsilon_j\le q_1<1\).
The second equation and its actual
operator bounds give the causal
estimate
\begin{equation}
 \|v_j\|_{j,*}
       \le C\sum_{i<j}q_1^{\,j-1-i}\frac{|t_i|}{x_i}.
 \label{rds:actual-short-memory}
\end{equation}
This follows by successively
substituting that equation and
bounding each product of its
profile operators by \(q_1\)
per step. The constant is
independent of horizon.
Substitution into the first
equation exhibits the memory
bound
\[
        C\epsilon_j
          \sum_{i<j}q_1^{\,j-1-i}\frac{|t_i|}{x_i}.
\]
The denominator here belongs
to the old index \(i\).
It cannot be replaced by a
claim of smallness for the
unweighted memory sum.

\subsubsection{The coupling weight and the diagonal product}

Set
\[
                 w_j=x_j/x_0,\qquad Y_j=t_j/w_j .
 \]
The actual scalar recursion is
\(x_{j+1}=x_j-b_j-r_j\),
with \(|r_j|\le C\epsilon_j\).
Consequently
\begin{equation}
 \frac{1-m_j/x_j+d_j}{w_{j+1}/w_j}
    =1+\frac{b_j-m_j+r_j+x_jd_j}{x_{j+1}} .
 \label{rds:exact-diagonal-ratio}
\end{equation}
This exact identity uses no
monotonicity of the trajectory.
For a large enough floor,
the bounded Gaussian increments
and small formation correction
make \(x_{j+1}\) comparable
to \(x_j\).
By HRT and
\eqref{rds:two-Gaussian-increments},
\begin{equation}
 \sum_{j<N}
   \frac{|b_j-m_j|+\epsilon_j}{x_{j+1}}
          \le \frac C X+
                       C\frac{(\log X)^m}{X}.
 \label{rds:diagonal-sum}
\end{equation}
The geometric numerator in
the first part has bounded
sum, and all \(x_j\ge cX\).
For the second part use
the second scale sum in
\eqref{hrt:two-scale-sums}.
Thus, for \(i<k\), the actual
diagonal product obeys
\begin{equation}
 \prod_{i\le j<k}\left|1-\frac{m_j}{x_j}+d_j\right|
             \le C\,\frac{x_k}{x_i}.
 \label{rds:diagonal-product}
\end{equation}
Take absolute values in
\eqref{rds:exact-diagonal-ratio},
use \(1+u\le e^u\) for
nonnegative \(u\), and
multiply. The factors \(w\)
telescope. Early increments
of either sign are harmless;
they were only required to
be bounded.

\subsubsection{The weighted memory estimate}

The identity
\[
                              \frac{w_i}{x_i}=\frac1{x_0}
\]
is the decisive cancellation.
Divide the first response
equation by \(w_{j+1}\),
insert
\eqref{rds:actual-short-memory},
and use that identity.
It gives
\begin{equation}
 |Y_{j+1}|\le(1+\eta_j)|Y_j|
       +C\frac{\epsilon_j}{x_{j+1}}
                    \sum_{i<j}q_1^{\,j-1-i}|Y_i|,
 \qquad
 \eta_j=C\frac{|b_j-m_j|+\epsilon_j}{x_{j+1}} .
 \label{rds:normalized-memory}
\end{equation}
The sum of the coefficients
\(\eta_j+\epsilon_j/x_{j+1}\)
is bounded by
\eqref{rds:diagonal-sum}
and becomes small with the
infrared floor.

For completeness put
\(M_j=\max_{i\le j}|Y_i|\).
The geometric age sum in
\eqref{rds:normalized-memory}
is at most \(M_j/(1-q_1)\).
Hence
\[
 M_{j+1}\le M_j\left\{
       1+\eta_j+
          \frac{C\epsilon_j}{(1-q_1)x_{j+1}}\right\}.
\]
Multiplication and
\eqref{rds:diagonal-sum}
give \(M_j\le C|t_0|\)
uniformly in \(j,N\).
Returning to
\eqref{rds:actual-short-memory}
then gives
\[
 \|v_j\|_{j,*}\le
       C\frac{|t_0|}{x_0}
                    \sum_{i<j}q_1^{\,j-1-i}
                              \le C\frac{|t_0|}{x_0}.
\]
Both bounds use the coupling
weight throughout the memory
sum.

\begin{proposition}\label{prop:rds-regular-stress-decay}
The first tracefree diagonal
response of the actual regular
trajectory, with bare coefficient
\(t_0\) and no bare irrelevant
response, satisfies
\begin{equation}
              |t_j|\le C\frac{x_j}{x_0}|t_0|,
       \qquad \|v_j\|_{j,*}\le C\frac{|t_0|}{x_0}.
 \label{rds:regular-stress-decay}
\end{equation}
The constants and infrared
floor are independent of
horizon and compatible volume.
For \(|t_0|\le C_0\sqrt{x_0}\)
and fixed sufficiently large
\(X=x_N\), its final
nonconstant insertion tends
to zero with bound
\begin{equation}
           |t_N|+\|v_N\|_{N,*}
                     \le C_X x_0^{-1/2}.
 \label{rds:square-root-input-decay}
\end{equation}
Since \(x_0=X+b_\infty N+
O(\log^{m+1}(X+N))\), this
is \(O_X(N^{-1/2})\).
\end{proposition}

\begin{proof}
The preceding causal response
construction and weighted
estimates prove
\eqref{rds:regular-stress-decay}.
At \(j=N\), \(x_N=X\).
Insert the stated bound
on \(t_0\) and the actual
bare coupling law from HRT.
The classical stress pointing
has uniformly bounded norm,
so these coefficient and
profile bounds also control
the resulting nonconstant
insertion on the final
working field ball.
\end{proof}

This proves the diagonal
stress decay for the stated
regular integral operations.
It uses the actual weighted
Gaussian law and the actual
formation derivatives.
It is a first-response
statement; it does not claim
a finite anisotropic trajectory
or a variable-weight domain.
The complementary-history
insertions in a full rotation
argument require their own
source and mass estimates.

% END LAYERED EXPANSION regular_diagonal_stress.tex

% BEGIN LAYERED EXPANSION formation/regular_scalar_insertion_nontriviality.tex
\subsection{A scalar insertion and nonvanishing in the regular sector}
\label{rsi:section}

Endpoint tuning gives an exact insertion identity when
the entire regular presentation is held fixed. We use
that identity to prove a positive variance bound which
is uniform in the ultraviolet depth on a fixed terminal
torus. The measure is the positive product of the
specified regular kernels, capped in one common coarse
field domain. It is not identified with the full Wilson
measure.

The proof has two ingredients. The flat flux contraction
makes every effective classical action nonconstant on
the same two coarse patches. The endpoint derivative
of the terminal profile is small, so the transferred
scalar insertion retains this separation. Positivity
then gives both patches a depth-independent probability.

\subsubsection{One frozen presentation and its exact density}

Fix a terminal torus in the compatible volume class of
\cref{layered:homogeneous-absolute-formation,prop:lcf-buffered-formation}.
Its integer period \(\ell\ge2\) is divisible by the chosen
localization sizes and exceeds ten cubes at each of those sizes.
These sizes and the volume floor are fixed before the endpoint
coupling is varied. Fix a real endpoint \(X\) above the coupling floor. For
each ultraviolet depth \(N\), take the HRT presentation
at this reference endpoint. Keep every frame
\(\theta_j\), fluctuation scale, relative contour
factor, free colour-ball cutoff, branch and domain
fixed. Different depths have different finite
lattices, but the estimates are uniform in \(N\).
No derivative below compares different presentations.

Vary only the terminal coupling to \(X+t\), with
\(|t|<cX\) in its analytic tuning disc. Write
\begin{equation}\label{rsi:endpoint-notation}
 \alpha_N(t)=x_0^{(N)}(X+t),\qquad
 H_N(t,V)=\sum_z h_{N,z}^{(N)}(X+t,V),
 \qquad S_N(V)=A(U_N(V)).
\end{equation}
The ultraviolet profile is zero. The classical
action \(S_N\) is independent of \(t\). All
intermediate coupling and profile dependences are
formed densities; they are not extra parameters
of the frozen kernels.

The exact positive composition of these kernels
has, relative to normalized terminal product Haar
measure, the density
\begin{equation}\label{rsi:exact-pushforward}
 \mathcal K_N(t,V)
       =e^{C_N(t)}\exp\{-(X+t)S_N(V)+H_N(t,V)\}.
\end{equation}
The scalar \(C_N\) includes every numerical,
Gaussian-vacuum and profile-vacuum factor in
the exact regular steps. It is real at real
\(t\), with its analytic logarithm fixed there.
It need not be bounded uniformly in \(N\).

Let \(\mathbb Q_N\) be the positive chain measure
formed from those frozen kernels and their
Haar/coarea measures, before the ultraviolet
factor \(e^{-\alpha_N(t)S_0}\) is inserted.
It is independent of \(t\). Equivalently, use
the successive deterministic mean constraints
and their positive regular weights. This
description keeps precisely the retained regular
component and its contour factors. On every
finite lattice the measure is finite and all
moments of the initial Wilson action \(S_0\)
are finite.

Choose below one common real raw-coarse domain
\(\Omega\) in the terminal field ball, independent
of \(N,t\). Define the capped regular probability
\begin{equation}\label{rsi:regular-measure}
 d\mathbb P_{N,t}^{\rm reg}
    =\frac{\mathbf1_{\{V\in\Omega\}}
                    e^{-\alpha_N(t)S_0}}{
                       \mathcal Z_N^{\rm reg}(t)}\,d\mathbb Q_N .
\end{equation}
Its terminal marginal is
\begin{equation}\label{rsi:infrared-measure}
 d\nu_{N,t}(V)=\frac{\mathbf1_\Omega(V)}{Z_N(t)}
                e^{-(X+t)S_N(V)+H_N(t,V)}\,dV,
 \qquad \mathcal Z_N^{\rm reg}(t)=e^{C_N(t)}Z_N(t).
\end{equation}
For real \(t\) near zero these are positive.
Conditional expectations at a specified \(V\)
mean the canonical coarea fibre ratios. This
version is defined throughout the admitted
coarse chart, including \(V=1\).

The HRT bounds, on fixed smaller tuning discs,
give
\begin{equation}\label{rsi:endpoint-estimates}
 \frac34\le\alpha_N'(0)\le\frac54,\qquad
 \|h_N(0)\|\le H_*,\qquad
 \|\partial_t h_N(0)\|\le C\frac{\epsilon(X)}X,
 \quad \epsilon(X)=X^{-1}(\log X)^m,
                         \quad m>1.
\end{equation}
For every fixed \(s\ge2\), Cauchy on the same
frozen discs also gives
\begin{equation}\label{rsi:higher-endpoint-estimates}
 |\alpha_N^{(s)}(0)|\le C_s\frac{(\log X)^m}{X^s},
 \qquad
 \|\partial_t^s h_N(0)\|
                       \le C_s\frac{\epsilon(X)}{X^s}.
\end{equation}
The profile assertion follows by applying Cauchy
to its first derivative on a fixed smaller disc.
These are exactly endpoint derivatives in
\cref{hrt:endpoint-slope,hrt:higher-endpoint-jets};
none differentiates a sharp cutoff.

\subsubsection{The scalar conditional identity and centering}

Put
\begin{equation}\label{rsi:transferred-observable}
 T_N=\alpha_N'(0)S_0,\qquad
 O_N(V)=S_N(V)-\partial_tH_N(0,V).
\end{equation}

\begin{proposition}[Exact scalar action insertion]
\label{prop:rsi-scalar-insertion}
For the capped regular measure,
\begin{equation}\label{rsi:conditional-action-identity}
 \mathbb E_{N,0}^{\rm reg}[T_N\mid V]
                         =O_N(V)-C_N'(0).
\end{equation}
Consequently its centered form is
\begin{equation}\label{rsi:centered-insertion-identity}
 \mathbb E_{N,0}^{\rm reg}
       [T_N-\mathbb E_{N,0}^{\rm reg}T_N\mid V]
                  =O_N(V)-\nu_{N,0}(O_N).
\end{equation}
This is a constant scalar action insertion.
It makes no assertion about a spatially varying
local source.
\end{proposition}
\begin{proof}
At fixed \(V\), differentiate the actual fibre
composition defining \(\mathcal K_N\). Its only
endpoint-dependent ultraviolet factor is
\(e^{-\alpha_N(t)S_0}\), since \(h_0=0\)
and \(\mathbb Q_N\) is frozen. Finite-dimensional
dominated differentiation gives
\[
 \partial_t\log\mathcal K_N(0,V)
                 =-\alpha_N'(0)\mathbb E[S_0\mid V].
\]
Differentiating \eqref{rsi:exact-pushforward}
gives \(C_N'(0)-S_N+\partial_tH_N(0)\).
These identities prove
\eqref{rsi:conditional-action-identity}.
The cap does not change a fibre conditional
at an admitted \(V\in\Omega\). Integrate
that equality against the terminal marginal
and subtract it. This cancels the complete
numerical vacuum derivative \(C_N'(0)\),
proving \eqref{rsi:centered-insertion-identity}.
The argument did not remove a vacuum term
from the uncentered identity.
\end{proof}

\begin{corollary}[Conditional variance and scalar information]
\label{cor:rsi-conditional-variance}
The exact decomposition is
\begin{equation}\label{rsi:variance-decomposition}
 \operatorname{Var}_{\mathbb P_{N,0}^{\rm reg}}(T_N)
  =\nu_{N,0}\{\operatorname{Var}(T_N\mid V)\}
                        +\operatorname{Var}_{\nu_{N,0}}(O_N).
\end{equation}
The normalized endpoint scores are respectively
\(-T_N+\mathbb E T_N\) and
\(-O_N+\nu_{N,0}O_N\). Thus the two variances
are the corresponding scalar Fisher informations,
and coarse conditioning cannot increase them.
\end{corollary}
\begin{proof}
Differentiate the normalized densities in
\eqref{rsi:regular-measure} and
\eqref{rsi:infrared-measure}. The centered
conditional identity identifies their scores.
Apply the elementary conditional-variance
formula. All variables have finite moments
on each finite regular graph.
\end{proof}

\subsubsection{The same two action patches at every depth}

Let \(n_b=4\ell^4\) be the number of stored
positive terminal bonds. In the normalization
\(\tfrac12\operatorname{tr}T^2=1\), fix a real
Abelian colour direction \(T\), and the coarse
field
\begin{equation}\label{rsi:cosine-probe}
 B_2(x)=T\cos(2\pi x_1/\ell),\qquad
                       B_\mu(x)=0\quad(\mu\ne2).
\end{equation}
Its supremum is one and its curl is nonzero. The
period-two case in the next formula is a kinematic
check; the regular-measure conclusions use only the
compatible torus fixed above. Write
\begin{equation}\label{rsi:curl-mass}
 q_\ell=\|\mathrm{d}_cB\|_2^2=
 \begin{cases}
       4\ell^4,&\ell=2,\\
       2\ell^4\sin^2(\pi/\ell),&\ell\ge3.
 \end{cases}
\end{equation}
The difference acts on stored oriented positive
bonds, so the two parallel links on a period-two
circle remain distinct lattice variables.

\begin{lemma}[A uniform positive quadratic direction]
\label{lem:rsi-uniform-action-direction}
There are \(R_0>0\) and \(M_\ell<\infty\),
independent of \(N\), such that
\begin{equation}\label{rsi:uniform-action-bounds}
 S_N(1)=DS_N(1)=0,\qquad
 D^2S_N(1)[B,B]\ge q_\ell,\qquad
 \sup_{\|A\|_\infty<R_0}|S_N(e^{iA})|\le M_\ell.
\end{equation}
For real admitted fields \(S_N\ge0\).
\end{lemma}
\begin{proof}
At the identity the actual fine response to
coarse datum \(B\) satisfies the literal
linear mean constraint. The difference between
the plain and straightened rows is a coarse
gradient, so their curls agree. The exact flat
identity and contraction
\eqref{neg:flat-flux-identity},
\eqref{neg:flat-flux-contraction} give
\[
                  \|\mathrm{d}_cB\|_{2,c}\le\|\mathrm{d}_hA^{\rm lin}\|_{2,h}.
\]
The fine Wilson Hessian in this direction is exactly
\(\|\mathrm{d}_hA^{\rm lin}\|_{2,h}^2\) in the displayed
colour normalization. Its first jet is zero,
so the second derivative of the background
map adds no acceleration term. This proves
the middle assertion in
\eqref{rsi:uniform-action-bounds}. No positivity
of a mass-added coarse Schur matrix is being
used for this conclusion.

The actual local action trees of
\cref{prop:het-actual-effective-action} give
a normally convergent analytic action on
a fixed complex coarse ball. Each original
fine density has its physical \(h^4\) weight;
the fine plaquette count therefore cancels
in a unit coarse cell. Summing its rooted
pointings on the terminal torus bounds
the total by \(C_A\ell^4\) on a smaller
fixed radius \(R_0\). Take this as
\(M_\ell\), enlarging it to include \(N=0\)
if necessary. The field is real at real
data, and its actual Wilson plaquette
terms are nonnegative. Their sum proves
the last assertion.
\end{proof}

Choose once
\begin{equation}\label{rsi:fixed-omega}
 0<r_\Omega\le\min\{r/4,R_0/4,\pi/4\},\qquad
 \Omega=\{V=e^{iA}:\max_b|A_b|<r_\Omega\}.
\end{equation}
The radius is fixed before \(N,X\) vary.
Let
\begin{equation}\label{rsi:patch-constants}
 \begin{split}
 0<a&\le\min\{r_\Omega/4,R_0/2,
                              q_\ell R_0^3/(8M_\ell)\},\\
 d_\ell&=q_\ell a^2/4,\qquad L_\ell=2M_\ell/R_0,\\
 0<\delta_\ell&\le\min\{r_\Omega/4,a/8,
                                           d_\ell/(8L_\ell)\}.
 \end{split}
\end{equation}
These constants are positive and do not depend
on ultraviolet depth or endpoint coupling.
Define product patches in raw coordinates by
\begin{equation}\label{rsi:two-patches}
 P_0=\{e^{iA}:\max_b|A_b|<\delta_\ell\},\qquad
 P_1=\{e^{iA}:\max_b|A_b-aB_b|<\delta_\ell\}.
\end{equation}
Both lie in \(\Omega\), and they are disjoint.

\begin{lemma}[Uniform separation on fixed patches]
\label{lem:rsi-action-patch-gap}
For every \(N\),
\par\noindent
\begin{minipage}{\linewidth}
\begin{equation}\label{rsi:action-patch-gap}
           \inf_{P_1}S_N-\sup_{P_0}S_N\ge\tfrac34d_\ell.
\end{equation}
\popQED\par
\end{minipage}
\end{lemma}
\begin{proof}
The scalar function \(z\mapsto S_N(e^{izB})\)
is bounded by \(M_\ell\) on \(|z|<R_0\).
Its divided coefficients have bounds
\(M_\ell R_0^{-m}\). For \(a\le R_0/2\),
the terms of degree at least three have
absolute sum at most
\(2M_\ell a^3/R_0^3\le q_\ell a^2/4\).
Its quadratic term is at least
\(q_\ell a^2/2\). Hence
\(S_N(e^{iaB})\ge d_\ell\), uniformly in depth.

On the real \(r_\Omega\)-ball, Cauchy in any
unit coarse Banach direction gives
\(\|DS_N\|\le L_\ell\): a complex disc of
radius \(R_0/2\) stays inside the action
ball. The line between a patch point and
its centre is in that real ball. The
variation of \(S_N\) on either patch is
therefore at most \(L_\ell\delta_\ell
\le d_\ell/8\). Subtracting the two
errors proves \eqref{rsi:action-patch-gap}.
\end{proof}

For normalized Haar measure on
\(\operatorname{SU}(2)\), the log-coordinate
density is
\[
              j(A)=(2\pi^2)^{-1}(\sin|A|/|A|)^2.
\]
On \(|A|\le\pi/2\) it is at least
\(2/\pi^4\). Thus the same two patches have
the explicit product-Haar lower bound
\begin{equation}\label{rsi:Haar-patch-mass}
             \operatorname{Haar}(P_i)\ge
       m_\ell:=\left(\frac8{3\pi^3}\delta_\ell^3\right)^{n_b}>0.
\end{equation}
Their dimension is \(3n_b=12\ell^4\), not an
ultraviolet dimension. The common cap has
Haar mass
\[
       m_\Omega=
       \left\{\frac{r_\Omega-\tfrac12\sin(2r_\Omega)}{\pi}\right\}^{n_b}>0.
\]

\subsubsection{Uniform positive probabilities and variance}

The terminal pointing bounds imply
\begin{equation}\label{rsi:total-profile-bounds}
 |H_N(0,V)|\le\ell^4H_*=:H_\ell,\qquad
 |\partial_tH_N(0,V)|\le
                      C\ell^4\frac{\epsilon(X)}X
                              \quad(V\in\Omega).
\end{equation}
Choose \(X_\ell\) above the HRT floor so that,
for all \(X\ge X_\ell\), the second quantity
is at most \(d_\ell/8\). This is possible
because \(\epsilon(X)/X\to0\).
Then the transferred observable has separation
\begin{equation}\label{rsi:observable-patch-gap}
                     \inf_{P_1}O_N-\sup_{P_0}O_N\ge d_\ell/2.
\end{equation}
All choices of patches precede this coupling
threshold.

Since \(0\le S_N\le M_\ell\) on the cap,
its terminal normalizer obeys
\begin{equation}\label{rsi:terminal-normalizer-bounds}
 m_\Omega e^{-XM_\ell-H_\ell}
                  \le Z_N(0)\le m_\Omega e^{H_\ell}.
\end{equation}
In particular each patch probability is at
least
\begin{equation}\label{rsi:patch-probability}
 \nu_{N,0}(P_i)\ge
            p_{X,\ell}:=m_\ell e^{-XM_\ell-2H_\ell}>0.
\end{equation}
These estimates do not use or bound the
possibly extensive vacuum factor \(e^{C_N}\).

\begin{proposition}[Finite-volume regular nonvanishing]
\label{prop:rsi-regular-nontriviality}
For the fixed compatible terminal torus above and \(X\ge X_\ell\),
uniformly over all ultraviolet depths,
\begin{equation}\label{rsi:positive-variance}
 \operatorname{Var}_{\nu_{N,0}}(O_N)
                \ge c_{X,\ell}:=p_{X,\ell}^2d_\ell^2/4>0,
 \qquad
 \operatorname{Var}_{\mathbb P_{N,0}^{\rm reg}}(T_N)
                                      \ge c_{X,\ell}.
\end{equation}
Consequently
\(\operatorname{Var}_{\mathbb P_{N,0}^{\rm reg}}(S_0)
\ge(16/25)c_{X,\ell}\).
\end{proposition}
\begin{proof}
For independent terminal samples \(V,V'\),
\(\operatorname{Var}(O_N)
=\tfrac12\mathbb E[(O_N(V)-O_N(V'))^2]\).
Restrict to \(P_0\times P_1\) and
\(P_1\times P_0\). The two probabilities
and \eqref{rsi:observable-patch-gap} give
the first lower bound. The second follows
from \eqref{rsi:variance-decomposition}.
Finally \(T_N=\alpha_N'(0)S_0\) and
\(\alpha_N'(0)\le5/4\). This proves the
last bound. The explicit constant may be
small as \(X\) or \(\ell\) changes; only
uniformity in ultraviolet depth is asserted.
\end{proof}

\subsubsection{Generating functions and endpoint contact terms}

The fixed presentation also gives exact higher
scalar identities. For sufficiently small \(t\),
the conditional coarea kernels satisfy
\begin{equation}\label{rsi:conditional-generating-identity}
 \begin{split}
 \log\mathbb E_{N,0}^{\rm reg}
       [e^{-(\alpha_N(t)-\alpha_N(0))S_0}\mid V]
   ={}&C_N(t)-C_N(0)-tS_N(V)\\
     &+H_N(t,V)-H_N(0,V).
 \end{split}
\end{equation}
Every logarithm is its germ at \(t=0\).
The global identity is
\begin{equation}\label{rsi:global-generating-identity}
 \mathbb E_{N,0}^{\rm reg}
       e^{-(\alpha_N(t)-\alpha_N(0))S_0}
              =e^{C_N(t)-C_N(0)}\frac{Z_N(t)}{Z_N(0)}.
\end{equation}
Both follow by changing only the ultraviolet
exponential in the same finite integral.
In particular no moving-frame or cutoff
derivative is hidden in these formulas.

For a genuine bare scalar source \(s\),
let \(t=\psi_N(s)\) solve
\(\alpha_N(t)=\alpha_N(0)+s\). On a common
smaller complex disc this inverse exists:
the HRT complex derivative is uniformly
close to one. Substitution in
\eqref{rsi:conditional-generating-identity}
is the conditional cumulant-generating
function of \(-S_0\) for that source.

Let \(\kappa_r(S_0\mid V)\) denote its
ordinary conditional cumulants. With
\(B_{q,r}\) the exponential partial Bell
polynomials, differentiation gives
\begin{equation}\label{rsi:Bell-contact-identity}
 \sum_{r=1}^q(-1)^r\kappa_r(S_0\mid V)
 B_{q,r}(\alpha_N'(0),\ldots,
                         \alpha_N^{(q-r+1)}(0))
   =C_N^{(q)}(0)-\mathbf1_{\{q=1\}}S_N(V)
                              +\partial_t^qH_N(0,V).
\end{equation}
The convention is determined by
\(\partial_t^q[m(t)^r/r!]_{t=0}
=B_{q,r}(m'(0),\ldots)\).
For example the second and third identities are
\begin{align}
 (\alpha_N')^2\operatorname{Var}(S_0\mid V)
           -\alpha_N''\mathbb E[S_0\mid V]
                     &=C_N''+\partial_t^2H_N,\notag\\
 -(\alpha_N')^3\kappa_3(S_0\mid V)
 +3\alpha_N'\alpha_N''\operatorname{Var}(S_0\mid V)
           -\alpha_N'''\mathbb E[S_0\mid V]
                     &=C_N'''+\partial_t^3H_N.
 \label{rsi:second-third-contacts}
\end{align}
All quantities here are evaluated at \(t=0\).
The contact terms involving higher bare
tuning derivatives cannot be dropped.

There is a useful vacuum-free version.
The profiles and action vanish at the
flat field for every endpoint. Subtract
\eqref{rsi:conditional-generating-identity}
at \(V=1\) to obtain
\begin{equation}\label{rsi:vacuum-free-generating-identity}
 \log\frac{\mathbb E[e^{-sS_0}\mid V]}
                {\mathbb E[e^{-sS_0}\mid1]}
       =-\psi_N(s)S_N(V)
                      +H_N(\psi_N(s),V)-H_N(0,V).
\end{equation}
The reference fibre is the canonical coarea
version, not an arbitrary value assigned to
a conditional expectation on a null event.
This exact ratio removes \(C_N\) at every
order. For instance
\begin{equation}\label{rsi:conditional-variance-variation}
 \operatorname{Var}(T_N\mid V)-\operatorname{Var}(T_N\mid1)
       =\partial_t^2H_N(0,V)
                          +\frac{\alpha_N''(0)}{\alpha_N'(0)}O_N(V).
\end{equation}
The endpoint bounds therefore control this
field-dependent difference by
\(C_\ell(\log X)^m/X^2\), uniformly in
depth. They do not bound its vacuum value.
More generally Cauchy for the inverse source
and \eqref{rsi:higher-endpoint-estimates}
give, for each fixed \(r\ge2\),
\begin{equation}\label{rsi:higher-conditional-variation}
 \sup_{V\in\Omega}
   |\kappa_r(S_0\mid V)-\kappa_r(S_0\mid1)|
                         \le C_{\ell,r}\frac{(\log X)^m}{X^r}.
\end{equation}
Indeed \(\psi_N'(0)=1/\alpha_N'(0)\),
while \(|\psi_N^{(r)}(0)|\le
C_r(\log X)^m/X^r\) for \(r\ge2\).
The latter follows by inverting the near-identity
analytic tuning on a disc of radius \(cX\)
and applying Cauchy to its derivative minus
its value at zero. Composing the profile
first-derivative bound gives the smaller
profile contribution. Thus the estimate
concerns coarse-field variation of the
conditional cumulant, rather than a
uniform ultraviolet cumulant bound.

Finally, direct differentiation of the total
normalizer gives the consistency identities
\[
 \begin{aligned}
 \partial_t\log\mathcal Z_N^{\rm reg}(0)&=-\mathbb E T_N,\\
 \partial_t^2\log\mathcal Z_N^{\rm reg}(0)
      &=\operatorname{Var}(T_N)-\alpha_N''(0)\mathbb E S_0,\\
 \partial_t^2\log Z_N(0)
      &=\operatorname{Var}_{\nu_{N,0}}(O_N)
                                  +\nu_{N,0}(\partial_t^2H_N).
 \end{aligned}
\]
They retain the same endpoint contacts as
\eqref{rsi:Bell-contact-identity} and agree
with the conditional-variance decomposition.

\subsubsection{A nondegenerate fixed-coarse subsequence}

The common terminal domain also permits an
ultraviolet subsequence without changing the
coarse cutoff or volume. Only finite-dimensional
coarse functions and probabilities are extracted.

Reduce the common endpoint disc once, and
write its radius as \(T=c_1X\). The complex
HRT slope bound gives
\begin{equation}\label{rsi:uniform-complex-tuning}
 \sup_{|t|<T}|\alpha_N'(t)-1|\le K(X)<\tfrac14,
 \qquad
 \sup_{|t|<T}\|\partial_th_N(t)\|
                           \le C\epsilon(X)/X.
\end{equation}
After decreasing \(c_1\) or increasing the
coupling floor, integration along the endpoint
segment gives \(\|h_N(t)\|\le H_*+1\).
Put
\(\delta\alpha_N(t)=\alpha_N(t)-\alpha_N(0)\).
Then \(|\delta\alpha_N(t)|\le(5/4)|t|\).
The bare value \(\alpha_N(0)\) itself is not
bounded in this argument.

Use the finite colour coordinates
\(A=(A_{b,c})\in\mathbb C^{3n_b}\), and the
polydisc
\begin{equation}\label{rsi:finite-polydisc}
 \mathcal D=\{(t,A):|t|<T,
                       \ |A_{b,c}|<2r_\Omega\ \hbox{for all }b,c\}.
\end{equation}
Its field coordinates obey
\(\max_b|A_b|<2\sqrt3r_\Omega<\min(r,R_0)\).
The closure of the real cap \(\Omega\) is
strictly inside this field domain. Thus
\(S_N(e^{iA})\), \(H_N(t,e^{iA})\),
\(\partial_tH_N(t,e^{iA})\), and
\(\delta\alpha_N(t)\) have uniform holomorphic
bounds on \(\mathcal D\), after using a
slightly smaller endpoint disc for the
displayed derivative. Their bounds depend
on the fixed \(\ell,X\), not on \(N\).

\begin{lemma}[Finite-dimensional compact extraction]
\label{lem:rsi-finite-dimensional-extraction}
There is a common subsequence \(N_k\to\infty\)
and holomorphic limits \(S_\infty,H_\infty,
\delta\alpha_\infty\), locally uniform on
\(\mathcal D\), with
\begin{equation}\label{rsi:analytic-subsequence}
 S_{N_k}\longrightarrow S_\infty,\quad
 H_{N_k}\longrightarrow H_\infty,\quad
 \partial_tH_{N_k}\longrightarrow\partial_tH_\infty,
 \quad\delta\alpha_{N_k}\longrightarrow\delta\alpha_\infty.
\end{equation}
They are real on real arguments,
\(\delta\alpha_\infty(0)=0\), and
\(3/4\le\delta\alpha_\infty'(0)\le5/4\).
\end{lemma}
\begin{proof}
This is a normal-family argument on the
\(1+3n_b\) complex coordinates just displayed.
On every smaller closed polydisc, the uniform
supremum bounds and coordinate Cauchy formulas
bound all first coordinate derivatives.
The scalar functions are therefore uniformly
bounded and equicontinuous on that compact
finite-dimensional set. Apply Arzel\`a--Ascoli
to the finite tuple of functions, then use
a countable nested exhaustion of \(\mathcal D\)
and a diagonal subsequence. The uniform limits
are holomorphic by the finite-dimensional
Cauchy integral formula. That formula also
identifies the limit of \(\partial_tH_N\)
with \(\partial_tH_\infty\).

Reality passes under uniform convergence.
The tuning vanishes at zero, and its first
derivative converges by the one-variable
Cauchy formula. The real slope interval
\eqref{rsi:endpoint-estimates} passes to
the limit. No normal-family theorem on
an infinite-dimensional field space is
being invoked.
\end{proof}

\begin{proposition}[A nondegenerate coarse ultraviolet subsequence]
\label{prop:rsi-fixed-coarse-subsequence}
On this subsequence, the terminal regular
probabilities converge in total variation to
\begin{equation}\label{rsi:limiting-coarse-measure}
 d\nu_\infty(V)=\frac{\mathbf1_\Omega(V)}{Z_\infty}
                    e^{-XS_\infty(V)+H_\infty(0,V)}\,dV.
\end{equation}
The observable
\(O_\infty=S_\infty-\partial_tH_\infty(0)\)
is real analytic near the cap and satisfies
\begin{equation}\label{rsi:limiting-nontriviality}
 \inf_{P_1}O_\infty-\sup_{P_0}O_\infty\ge d_\ell/2,
 \qquad
 \operatorname{Var}_{\nu_\infty}(O_\infty)
                         \ge c_{X,\ell}>0.
\end{equation}
In fact \(O_{N_k}\to O_\infty\) uniformly
on \(\overline\Omega\), and the variances
converge to the displayed positive variance.
\end{proposition}
\begin{proof}
The field closure of the cap lies in a
compact subpolydisc. Hence both the density
exponents and \(O_{N_k}\) converge uniformly
there. The positive normalizer bounds
\eqref{rsi:terminal-normalizer-bounds} pass
to \(Z_\infty\), so its value is nonzero.
The unnormalized real densities converge
uniformly; division by the convergent
positive normalizers proves total-variation
convergence. The cap is the same for every
depth, with fixed product Haar measure.
Its boundary has Haar measure zero.

The patch inequalities pass under uniform
convergence, and both limiting patch masses
are at least \(p_{X,\ell}\). The same
two-sample variance proof yields
\eqref{rsi:limiting-nontriviality}.
Alternatively bounded uniform convergence
of \(O_{N_k}\), together with total variation,
gives convergence of its first two moments
and hence of its variance. Analyticity and
reality are those of the extracted functions.
\end{proof}

There is also a common complex endpoint disc
on which normalized coarse source functionals
converge. Set
\[
 D_\ell=C\ell^4\epsilon(X)/X,\qquad
 t_\ell=\min\{T/2,\log(3/2)/(M_\ell+D_\ell)\}.
\]
For \(|t|<t_\ell\), the density relative
to its positive value at zero differs from
one by at most \(1/2\), since
\[
 |-tS_N(V)+H_N(t,V)-H_N(0,V)|
                           \le |t|(M_\ell+D_\ell).
\]
Integrating against \(\nu_{N,0}\) gives
\begin{equation}\label{rsi:common-zero-free-endpoint-disc}
                      |Z_N(t)/Z_N(0)-1|<1/2.
\end{equation}
Thus these capped normalizers are zero free
on this smaller, possibly \(\ell\)-dependent
disc. Uniform convergence on compact subsets
gives locally uniform convergence of normalized
coarse expectations of bounded fixed probes,
and of their endpoint derivatives. The same
holds for bounded holomorphic probes depending
on \(t\). This argument does not claim zero
freeness on the full tuning disc \(|t|<T\).

\begin{corollary}[Limit of the vacuum-free conditional source ratios]
\label{cor:rsi-source-ratio-subsequence}
On a common smaller bare-source disc,
\(\psi_{N_k}\to\psi_\infty\) locally uniformly,
where \(\delta\alpha_\infty(\psi_\infty(s))=s\).
The logarithms in
\eqref{rsi:vacuum-free-generating-identity}
converge locally uniformly in \((s,V)\) to
\begin{equation}\label{rsi:limiting-source-ratio}
 -\psi_\infty(s)S_\infty(V)
       +H_\infty(\psi_\infty(s),V)-H_\infty(0,V).
\end{equation}
Their derivatives give the corresponding
limiting coarse-field variations of conditional
cumulants. No convergence of \(C_N\), its
derivatives, or the uncentered conditional
source functions is asserted.
\end{corollary}
\begin{proof}
Write \(\delta\alpha_N(t)=t+r_N(t)\).
On the smaller common disc,
\(r_N(0)=0\) and \(\sup|r_N'|\le K(X)<1/4\).
For \(|s|<T/2\), the map \(t\mapsto s-r_N(t)\)
is an invariant contraction on
\(|t|\le2T/3\). It gives the common inverse
\(\psi_N\). The same applies to the limit.
Uniform convergence of \(r_{N_k}\) on that
closed smaller disc and the contraction
bound imply locally uniform convergence
of the inverses. Substitute in the exact
vacuum-free identity. Its analytic functions
and their derivatives converge by Cauchy
on compact subdomains, proving the result.
\end{proof}

This is a subsequential limit at one fixed
coarse cutoff and finite volume. The terminal
coarse observable is nondegenerate in that
limit. The vacuum factor is neither needed
nor extracted. No consistency between different
terminal volumes, horizons or coarse cutoffs
has been proved by this compactness argument.

\subsubsection{What a positive history mixture would add}

Suppose a separately established positive
decomposition on the same ultraviolet
observable space contains this capped
regular probability with weight \(p_N\):
\[
                  \mu_N=p_N\mathbb P_{N,0}^{\rm reg}
                                      +(1-p_N)\rho_N.
\]
For any common square-integrable observable,
the variance formula for its component flag
implies \(\operatorname{Var}_{\mu_N}(T_N)
\ge p_N\operatorname{Var}_{\mathbb P_{N,0}^{\rm reg}}(T_N)\).
Hence a genuine history floor \(p_N\ge p_*>0\)
would give \(\operatorname{Var}_{\mu_N}(T_N)
\ge p_*c_{X,\ell}\). The corresponding
\(S_0\) bound has the factor \(16/25\).

No such history floor is proved here.
The regular kernel product, its gauge weights
and its terminal cap are part of the measure
in this subsection. There is no asserted
full Wilson-law identification, volume-uniform
lower bound, continuum Schwinger field or
local-source transfer identity. The results
are the exact scalar normalization, uniform
finite-volume regular nonvanishing, and the
nondegenerate fixed-coarse subsequence just
constructed. They do not construct a limit
of the full ultraviolet measures.

% END LAYERED EXPANSION formation/regular_scalar_insertion_nontriviality.tex


% BEGIN MODULE chapters/07_sources.tex
\subsection{Sources, local energies, and separated correlation bounds}
\label{candidate:sources}
% BEGIN MODULE chapters/source_domains_v2.tex
\subsubsection*{Source domains and the exact running-shift identity}

This section fixes the source variables used in the conditional estimates.
The identities are finite-dimensional; the multiscale bounds for the
effective energies stated later require separate estimates.

\phantomsection\label{notdef:level-depth}
Fix $K$, a torus $T_\ell$, an odd $L>1$, and positive constants $M,c_M$ with $c_M\ge M$.
The depth $s$ corresponds to physical spacing $L^{-s}$; level $j$ of a
run of length $K$ corresponds to depth $K-j$. \phantomsection\label{notdef:source-triples}
A source triple is

\begin{gather*}
\FitMath{w=(c,(f_1,\ldots,f_n),(z_1,\ldots,z_n)),} \\[3pt]
\FitMath{w(x)=c+\sum_i z_if_i(x),}
\end{gather*}

where $f_i$ are real Lipschitz functions, $c,z_i$ are complex, and the
\phantomsection\label{notdef:source-norms}
family may be empty. Put

\begin{gather*}
\FitMath{\|f\|_{\sharp,s}=\max\{\|f\|_\infty,
                         c_ML^{-s}\operatorname{Lip}f\},} \\[3pt]
\FitMath{\|w\|_{\flat,s}=|c|+\sum_i|z_i|\|f_i\|_{\sharp,s}.}
\end{gather*}

The class $\mathcal W_r^{(s)}$ means the triples with
$\|w\|_{\flat,s}<r$. It is a union of finite-dimensional source balls,
one for each fixed real family. Holomorphy on this class means holomorphy
in $(c,z)$ on each such ball. The number $\|w\|_{\flat,s}$ belongs to the
chosen triple and can change when its function is represented differently.
The physical partition function and every source-local energy are required
to agree for two triples that have the same values on all bare plaquettes
on which that quantity depends. This representation-independence is a
separate part of the source format.

Always $\|w\|_\infty\le\|w\|_{\flat,s}$ and
$c_ML^{-s}\operatorname{Lip}w\le\|w\|_{\flat,s}$. Thus, if
$m=K-s$ and $j\le m$,

\begin{gather*}
\FitMath{c_Ma_{K-j}\operatorname{Lip}w
   =L^{-(m-j)}c_Ma_{K-m}\operatorname{Lip}w
   \le L^{-(m-j)}\|w\|_{\flat,s},} \\[3pt]
\FitMath{a_t=L^{-t}.}
\end{gather*}

This is the exact scale gain used when a source is frozen on an earlier
block. It follows from the norm definition, not from an analyticity theorem.

\phantomsection\label{notdef:source-profiles}
Here is an explicit algebraic version of the running profiles. Let $P_K$
be the finite set of bare plaquettes, with barycentres $x(p)$. For each
level $j$ choose a finite set $Y_j$ of base points, anchors $p_y\in P_K$,
and nonnegative functions $h_y^{(j)}$ on $P_K$ such that

\begin{gather*}
\FitMath{\sum_{y\in Y_j}h_y^{(j)}(p)=1,} \\[3pt]
\FitMath{h_y^{(j)}(p)\ne0\,\Longrightarrow\
       |x(p)-x(p_y)|\le C_hMa_{K-j}.}
\end{gather*}

Choose also profiles $0\le\varphi_j(p)\le1$. The grid, anchors, partition
and profiles are fixed independently of the source throughout the analytic
family. They are supplied by the geometric construction; the identities
below hold for every finite choice satisfying these conditions.

Let $b_j$ be holomorphic on $D(0,\rho)$ and define

\begin{gather*}
\FitMath{d_j=\sup_{0<|z|<\rho}\frac{|b_j(z)-b_j(0)|}{|z|},} \\[3pt]
\FitMath{\sum_{j\le k}d_j\le\frac13.}
\end{gather*}

For triples whose function has $\|w\|_\infty<\rho$, set

\begin{gather*}
\FitMath{x_k(p)=\beta-\sum_{j\le k}b_j(0)\varphi_j(p),}
\end{gather*}


\begin{gather*}
\FitMath{c_j(p;w)=\varphi_j(p)\sum_{y\in Y_j}
                   h_y^{(j)}(p)b_j(w(x(p_y))),}
\end{gather*}


\begin{gather*}
\FitMath{\zeta_k(p;w)=w(x(p))+
   \sum_{j\le k}\varphi_j(p)\sum_{y\in Y_j}h_y^{(j)}(p)
             [b_j(w(x(p_y)))-b_j(0)].}
\end{gather*}

Then the exact bookkeeping identity is

\begin{gather*}
\FitMath{x_k(p)-\zeta_k(p;w)
      =\beta-w(x(p))-\sum_{j\le k}c_j(p;w).}
\end{gather*}

To prove it, expand both sides and use $\sum_yh_y^{(j)}=1$ in each
constant term $b_j(0)\varphi_j$. Consequently the Wilson weights in a
sourced exponent can be added or subtracted without changing their total
bare coefficient. This assertion concerns that total coefficient; it does
not assert equality of two effective measures unless the exact integral
transport has also been proved.

The same definitions give

\begin{gather*}
\FitMath{|\zeta_k(p;w)-w(x(p))|\le\tfrac13\|w\|_\infty,} \\[3pt]
\FitMath{|\zeta_k(p;w)|\le\tfrac43\|w\|_\infty.}
\end{gather*}

Each value is holomorphic on every source ball with supremum norm below
$\rho$, because only finitely many evaluations and compositions occur.
It depends on the source only inside distance $C_hMa_{K-k}$ from $p$,
including the value at $p$ itself; here $a_{K-j}\le a_{K-k}$ for $j\le k$.

For a constant source $w\equiv z$, in a region where every
$\varphi_j(p)=1$, this reduces to

\begin{gather*}
\FitMath{\zeta_k(z)=z+\sum_{j\le k}[b_j(z)-b_j(0)],} \\[3pt]
\FitMath{\zeta_0(z)=z,} \\[3pt]
\FitMath{\zeta_{k+1}(z)=\zeta_k(z)+b_{k+1}(z)-b_{k+1}(0).}
\end{gather*}

Differentiating at zero gives
$\zeta_k'(0)=1+\sum_{j\le k}b_j'(0)$ and
$|\zeta_k'(0)-1|\le1/3$. The derivative bound follows from the definition
of $d_j$ by taking the limit at zero. For spatially varying sources, freezing
the function $w$ itself costs at most $\operatorname{Lip}w$ times the anchor
distance. To freeze $b_j\circ w$ one must work on a smaller source disc.
Indeed $h_j(z)=[b_j(z)-b_j(0)]/z$ extends holomorphically to zero and is
bounded by $d_j$. For $|z|\le\rho'<\rho$, Cauchy's estimate gives

\begin{gather*}
\FitMath{|b_j'(z)|=|h_j(z)+zh_j'(z)|
       \le d_j\frac{\rho}{\rho-\rho'}.}
\end{gather*}

Thus on the half-disc the additional factor is at most $2d_j$.
Apply the mean value integral on the segment between the two source values
to bound their composed difference. The displayed depth identity then supplies
the factor $L^{-(m-j)}$, with the fixed anchor constant $C_hM/c_M$.
The radial bound $d_j$ alone is not a Lipschitz bound on the whole open disc.

All smallness, radius and representation-independence conditions used here
must be retained when invoking the sourced RG format. They do not follow
from the entire dependence of the bare lattice partition function alone.

% END MODULE chapters/source_domains_v2.tex

\subsubsection{Transporting the observables: sources, analyticity, and bounds independent of the cutoff and of the volume}

The measure $\mu_{K,\ell,\beta}$ is described at the unit scale in \cref{candidate:averaging,candidate:large_field,candidate:flow}, and the present step gives the corresponding description of its correlations. Because $\mathcal{O}(f)=A(f,U)$, a source $w=c+\sum_iz_if_i$ acts as a complex shift of the bare inverse coupling that depends on position:

\begin{gather*}
\FitMath{\rho^w_0=e^{-A(\beta-w,U)},} \\[3pt]
\FitMath{Z(w)=\int\rho^w_0\,dU,} \\[3pt]
\FitMath{S^T_{n;K,\ell,\beta}(f_1,\dots,f_n)=\partial_{z_1}\cdots\partial_{z_n}\log Z(w)\big|_{w=0}.}
\end{gather*}

\phantomsection\label{notdef:local-source-energy}
The renormalisation steps are applied to $\rho^w_0$. Whatever the source leaves behind is gathered into local numbers, the source energies $e_j(X;w)$ of \cref{candidate:sources}. We set $\mathcal{E}(w):=\sum_{1\le j\le K}\sum_{X\in\mathcal{D}_j}e_j(X;w)$, $\widetilde{Z}(w):=e^{-E-\mathcal{E}(w)}Z(w)$ and $C_g:=C(1+\log g^{-2})$.

\begin{proposition}[{The construction with sources}]\label{prop:source-domains-4-1} Let $0<g\le\gamma_J$. For each first-passage datum $(K,\beta)$, each $\ell\in\mathfrak{L}(g)$ and each source satisfying $\|w\|_\flat<2\varrho$, the pair $(E,(e_j(X;\cdot)))$ forms a system of localised source energies. That is, each $e_j(X;w)$ is holomorphic in $(c,z)$ (analyticity). If the $f_i$ are real, it takes real values whenever $(c,z)$ is real. It depends on $w$ only through its values on $X^{\sim1}$, and it is zero at $w=0$ (locality). Moreover, uniformly,

\begin{gather*}
\FitMath{|e_j(X;w)|\le Ce^{-\kappa_ed_j(X)},} \\[3pt]
\FitMath{|\widetilde{Z}(w)|\le e^{C_g\ell^4},} \\[3pt]
\FitMath{\widetilde{Z}(0)\ge e^{-C_g\ell^4}}
\end{gather*}

(these are, in order: the exponential bound, whose rate $\kappa_e:=c\kappa$, $c<1$, exceeds the entropy rate of the localisation domains through a cube; upper stability; the floor). The same conclusions hold along a reference run, with every choice frozen at $\theta_k$, as long as $x^\natural_k$ remains in a tube around $\theta_k$ with a margin of $C\varrho$. For each $K$ this requirement is a condition on $\beta$ only. When $g\le\gamma_U$, it is met on the interval $I_K$ of \cref{sec:flow-frozen-choices-slopes-and-tuned-bare-couplings}, an interval that contains the tuned $\beta^\natural_K$. The conclusions also hold for the run stopped at any level $k$, where the sums are stopped at $j\le k$, the vacuum constant $E_{(k)}$ of the first $k$ steps replaces $E$, and $C(1+\log x_k)$ per site of $\mathbb{T}^{(k)}$ replaces $C_g\ell^4$ \cref{sec:sources-analytic-dependence-on-the-source-locality-and-decay-in-tame}.
\end{proposition}

%
%
%
\paragraph*{The localized energies in the source transfer.}
At the trivial pair \(P_0\), let \(\Delta_w\) denote the
difference from the unsourced term. The energy of a domain
at level \(j\) is
\[
 e_j(X;w)=
 \sum_{z\in X\cap\mathbb T^{(j)}}\!
       \Delta_w E^{(j)}(X,P_0,z)
       +\Delta_w R^{(j)}(X,P_0),\qquad e_0=0 .
\]
The normalizer and the ratio contribution computed with the
fixed real weight cancel in this difference. The source
dependence of the coupling shift is already in the sourced
regular term.
On the fixed source ball of radius $r$, apply the source-line
estimate with reference source zero. On the trivial pair the
regular difference is bounded by
$Cr^{-1}\|w\|_{\flat,K-j} e^{-(\kappa+1)d_j(X)}$ per pointing.
The remainder difference has bound
$Cr^{-1}\|w\|_{\flat,K-j}\theta_j^{-\kappa_0/2}e^{-\kappa d_j(X)}$.
The at most $CM^4(1+d_j(X))$ pointings use the extra unit of
decay, giving
\[
 |e_j(X;w)|\le CM^4r^{-1}\|w\|_{\flat,K-j}e^{-\kappa d_j(X)} .
\]
The retained source-energy rate can be chosen between the
rooted entropy rate and $\kappa$. At a witness stop the
separated-source bound uses only $j\le k_0$. Coarser origins
can read both supports and remain in the source derivatives
if the run is continued. If \(\mathcal S_\emptyset\) denotes the
all-small-field label, the vacuum bookkeeping gives
\[
 E+\sum_{j,X}e_j(X;w)
              =-\mathcal E_K(\mathcal S_\emptyset;w).
\]
Thus the same \(E\) and the sum of these local energies are
the constants removed in \(\widetilde Z\).
For a label \(\mathcal S\), write \(\mathsf T_j(\mathcal S;w)\)
for its share of these differences. The constant after
normalization is
\[
 -E-\sum_{j,X}e_j(X;w)-\mathcal E_K(\mathcal S;w)
 =-[E+\mathcal E_K(\mathcal S;0)]
  -\sum_j\bigl[\mathsf T_j(\mathcal S_\emptyset;w)
                       -\mathsf T_j(\mathcal S;w)\bigr].
\]
A missing pointing or remainder is rooted at a cube
of \(Z_j^+\) by its own domain or its base point.
Sum the trivial-pair difference with its rooted
decay, including the at most \(CM^4(d_j+1)\)
pointings per domain. The source share therefore
costs \(C\#_j(Z_j^+)\). The Gaussian, Haar, and
gauge constants at the fixed real weight give
the unsourced logarithmic price on the same
missing-site count. These are the scalar prices
in \eqref{hi:label-load}. The weighted product
majorant of the literal label terms is required
when summing the remaining factors.

Use \(s=K-j\) and put
\[
 V_{a,\ell}=E,\qquad U_{a,\ell}(W)=0,\qquad
 e_{a,\ell,s,X}(W)=e_{K-s}(X;w).
\]
Here \(w\) is the value of the physical source \(W\)
at the prescribed plaquette points. At \(s=K\)
the last entry is zero. The source appears through
these values and the values of the same function
at the running source points, so the entries are
functions of \(W\). The radius in the all-state
transfer is \(\varrho\) with the corresponding source
norm, or any smaller fixed radius under the norm
comparison of \cref{notdef:source-norms}. The two
stability estimates of \cref{prop:source-domains-4-1}
then give \eqref{eq:all-state-source-transfer-ST5}
with this assignment.

The actual domains at that scale are connected
unions of cubes of side \(Ma_s\). Their one-layer
enlargement contains every source read, with one
fixed additional reach for the barycentre of a
bare plaquette incident to that enlargement.
Use the compatible tori \(\ell=2L^m\) at sufficiently
large allowed sides. The cube count in each direction
is \(2L^{m-m_0+s}\), so the indexing exists at
all actual scales. The length \(d_j(X)\) of the
rooted estimate uses a cube-edge tree within \(X\).
The length \(m(X)\) used in
\cref{ch:graded-rooted-source} allows the tree to leave
\(X\), and hence \(m(X)\le d_j(X)\).
The bound for \(e_j\) therefore holds with \(m(X)\).
Choose the retained rate above the rooted
closed-walk count for these cube unions as in
\eqref{eq:graded-source-bound-GS1}. Each union
has its original single domain index. This gives
the indexing and source range needed for the
geometric estimates of D1. The uniform upper
bound and floor use the sourced history estimates
on the original source ball.

\begin{proposition}[{Correlation bounds: \cref{thm:continuum-main} (c)(1)--(2)}]\label{prop:source-domains-4-2} Take $g,K,\beta,\ell$ as above and $n\ge2$. Let the $f_i$ be real Lipschitz functions whose supports lie at mutual distances $\ge\mathsf{d}>0$. For the second and the third bound, let each support also be contained in a ball of radius $R$. Then there exists a side $\ell_G\in\mathfrak{L}(g)$, depending only on $g$, for which

\begin{gather*}
\FitMath{|S^T_n|\le C_n(\mathsf{d})e^{nc_V\ell^4}\prod_i\|f_i\|_\sharp;} \\[3pt]
\FitMath{|S^T_n|\le C^\infty_n(\mathsf{d};R)\prod_i\|f_i\|_\sharp\,\,(\ell\ge Cn(R+1));} \\[3pt]
\FitMath{|S^T_n|\le A(\mathsf{d};R)B(R)^n(n!)^{\varkappa_\star}\prod_i\|f_i\|_\sharp\,\,(\ell\ge\ell_G),}
\end{gather*}

uniformly, where every constant may also depend on $g$, $C^\infty_n(\mathsf{d};R)$ does not increase with $\mathsf{d}$, and one may take $A(\mathsf{d};R)=\exp(\max\{1,\mathsf{d}^{-5}\})$, after choosing $\varkappa_\star=k+1$ with the fixed graded exponent $k$ of \cref{thm:graded-rooted-source}. The sharper form proved there is $[2B(R)]^n(n!)^k\max\{1,\mathsf d^{-5n}\}\prod_i\|f_i\|_\sharp$. A polynomial separation prefactor independent of $n$ would require an additional estimate; it is not asserted here. The second and third bounds also hold for $S^T_n[\nu]$ in each periodic Gibbs state $\nu$ at $\beta$ on $\mathbb{L}_K$, with no condition on the volume.
\end{proposition}

\Cref{prop:source-domains-4-13} is the third main result. Under the law of the histories, the probability that a cell is live, or that it is created, is $\le Ce^{-cp_0(g_k)}$, uniformly in the volume. The short-distance laws rest on this result. \Cref{sec:short-distance-isolating-a-single-scale-coupling-and-normalisation-at-the-scale} treats the two-point law, clauses (c)(3), (c)(4). It uses these bounds on a single cell and, from their age-resolved form (\cref{prop:source-domains-4-14}), the age series at a rate $b\ge8$. The four-point argument for \cref{thm:continuum-main}(c)(5) requires the series at a rate $b>16$.

The same tools also give two further results that are used later. The first is the one-point transfer (\eqref{eq:all-state-source-transfer-ST14}), $|\nu(\mathcal{O}(f))-\nu'(\mathcal{O}(f))|\le C_1(R)\|f\|_\sharp$, which compares periodic Gibbs states taken at the same first-passage datum $(K,\beta)$. With it, one can centre a pure phase at the means of that periodic Gibbs state whose decomposition contains the phase, the coordinates of \cref{ch:phase-interface} become compact, and the one-point numbers $\mathfrak{c}_\xi$ of \cref{thm:phase-main}(3)(ii) are bounded. It also enters the growth constants of (b)(v), so that a single set of constants covers these components as well. The second consists of the insertion bound and the plate bound (\cref{thm:uniform-one-field-insertion,thm:plate-density-from-insertion}), which are the a priori estimates on which \cref{thm:all-orders-angular-closure} rests.

\Cref{thm:continuum-main}(a)(ii) uses only the bounds of \cref{prop:source-domains-4-2}. \Cref{thm:continuum-main}(a)(iii) (for the re-centring), \cref{ch:euclidean-hilbert-bridge}, \cref{ch:euclidean-hilbert-bridge} and \cref{ch:phase-interface}--\cref{ch:scalar-phase-hilbert-bridge} also use the one-point transfer in \eqref{eq:all-state-source-transfer-ST14}. The tuned couplings need not be first-passage data. \Cref{thm:terminal-coupling-limit-v7} and \cref{ch:comparison-capsule-v7} apply \cref{prop:source-domains-4-1} along reference runs, and this gives the tuned limits the torus bound only. The format underlying \cref{prop:source-domains-4-1} is the starting point of \cref{sec:short-distance-the-short-distance-two-point-law-and-its-asymptotic-freedom}--\cref{sec:short-distance-isolating-a-single-scale-coupling-and-normalisation-at-the-scale}. The difficulty lies in two places. The first is \cref{prop:source-domains-4-1}: the source term is marginal, since a step leaves it no smaller, just as it leaves the action no smaller; it takes complex values, so positivity is no longer available; and it must be carried through the large-field operations. The second is \cref{prop:source-domains-4-13}, in which the ratio of the stability bounds loses $e^{C_g\ell^4}$. Everything else (Cauchy's estimates, reflection positivity, the chessboard estimate) is classical, except for the bookkeeping in $n$.

\paragraph{The curvature source with fixed reference data}
\label{sec:sources-the-curvature-source-with-fixed-reference-data}

\textbf{Objects.} The norm $\|f\|_\sharp=\max\{\|f\|_\infty,CM\operatorname{Lip}f\}$ controls how much $f$ can vary within any $M$-cube of the unit lattice. The sourced densities $\rho^w_k$ arise from $\rho^w_0$ through the operations of \cref{candidate:averaging,candidate:large_field}, and every choice in them is the one made in the run at $w=0$. These choices are the tests $\chi$, the thresholds $\varepsilon_k,\delta_k$, the radii, the sizes $\check{R}_k$, the backgrounds $U_k(\mathbf{V})$ (which minimise the action with no source), the scaling $\mathcal{A}_k=g_k\mathcal{A}^{\mathrm{s}}_k$, the covariances $C^{(k)}$, the gauge-fixing weight and the reference laws of $\mathbf{R}$.

\begin{proposition}[{Kernels; same histories; exactness}]\label{prop:source-domains-4-3} To each history $h$ of level $k$ there corresponds a positive kernel $\mathsf{K}_h(U,dV_k)$, built from the choices of the run and from nothing else. For $\Phi$ bounded and gauge invariant, the operations carry $\Phi\rho_0$ to

\begin{gather*}
\FitMath{\sum_h\tau_h[\Phi],} \\[3pt]
\FitMath{\tau_h[\Phi]=\int\Phi(U)\,\rho_0(U)\,\mathsf{K}_h(U,\cdot)\,dU,} \\[3pt]
\FitMath{\sum_h\tau_h[\Phi](\mathcal{V}_k)=\int\Phi\rho_0\,dU.}
\end{gather*}

For the choice $\Phi=e^{A(w,U)}$ one obtains $\rho^w_k\,dV_k=\sum_h\tau^w_h$, where the sum runs over the histories of $\rho_k$. Each $\tau^w_h$ is a complex measure which, as a function of $(c,z)$, is entire, and $\tau^0_h=\tau_h\ge0$. Moreover $\int\rho^w_k\,dV_k=Z(w)$ for $k\le K$.
\end{proposition}

Each step consists of finitely many moves. The first is the insertion of $1=\int dV\,\delta(\mathcal{M}(U)V^{-1})$. The second is a family of decompositions of unity $1=\sum\prod\chi\prod\chi^c$. The third is a gauge-fixing factor, of total integral $1$ over the gauge group; its insertion is harmless since $\Phi\rho_0$, like each $a_p$, is gauge invariant. The fourth is a set of changes of variables. The last occurs inside $\mathbf{R}$: a term is replaced by the product of a fixed probability law and its integral over a component. Every move is linear in the density and leaves the total integral unchanged, and no move requires positivity. The tests depend on the fields only, so the regions, labels and histories are the same as at $w=0$.

\begin{proposition}[{Exact insertion of block observables}]\label{prop:source-domains-4-4} Given a history $h$ of level $k$, denote by $Y(h)$ the union of all components that $\mathbf{R}$ has acted upon at levels not exceeding $k$. A bond $c$ of $\mathbb{T}^{(k)}$ is said to be far from $Y(h)$ when neither the two blocks $B^k(c_\pm)$ nor the bare bonds that join them meet $Y(h)^\sim$. There exists a probability measure $\mathsf{P}_k$ on triples $(U,h,V)$, each made of a bare field, a history and a block field, with four properties. First, its marginal on $U$ is $\mu_{K,\ell,\beta}$. Second, under its marginal on $(h,V)$ each bounded gauge-invariant $\Phi(h,V)$ has integral $Z(0)^{-1}\sum_h\int\Phi(h,V)\,\tau_h(dV)$. Third, it is invariant under $(U,h,V)\mapsto(U^u,h,V^{u|})$, where $u$ is a gauge transformation of the bare lattice and $u|$ is its restriction to the sites of $\mathbb{T}^{(k)}$. Fourth, under it $V=\mathcal{M}^kU$ holds on every bond far from $Y(h)$. Consequently, setting $F^{(h)}(V):=\mathbb{E}_{\mathsf{P}_k}[F(\mathcal{M}^kU)\mid h,V]$, which is the conditional expectation with respect to $\mathsf{P}_k$, taken in a version that is invariant in $V$, one has for each bounded gauge-invariant $F$ on $\mathcal{V}_k$

\begin{gather*}
\FitMath{\big\langle F(\mathcal{M}^kU)\big\rangle_{K,\ell,\beta}=\frac{\sum_h\int F^{(h)}\,d\tau_h}{\sum_h\int d\tau_h},} \\[3pt]
\FitMath{|F^{(h)}|\le\sup|F|,} \\[3pt]
\FitMath{F^{(h)}=F\,\text{ if }F\text{ reads only bonds far from }Y(h);}
\end{gather*}

in particular, $F^{(h)}=F$ for any history in which $\mathbf{R}$ has acted on none of the components. In precisely this sense $\rho_k$ represents the block averages of the bare field. 
\end{proposition}

The mechanism has three parts. First, $\mathbf{T}$ integrates $F\circ\mathcal{M}$ exactly against the density one level lower. Second, on each component on which it acts, $\mathbf{R}$ puts a reference law in place of the law of the block field; this alters the joint distribution of $\mathcal{M}^kU$ and $V_k$ only in a neighbourhood of that component, and it leaves every mass as it was. Third, the steps are glued in succession, each being regarded as a transport of the gauge-averaged terms. All the assertions are identities between measures and involve neither a constant nor a smallness condition. The role of the proposition is to fix what $\rho_k$ means; nothing is derived from it. In the two-point law \cref{sec:short-distance-isolating-a-single-scale-coupling-and-normalisation-at-the-scale}, a single scale is not isolated by this coupling, because, inside the region it heals, the law of the block field is replaced by $\mathbf{R}$. It is isolated by the exactness with the sources read at depth $s$, that is, by $\sum_h\tau^w_h(\mathcal{V}_{K-s})=Z(w)$ (\cref{prop:source-domains-4-3}, together with \cref{sec:sources-analytic-dependence-on-the-source-locality-and-decay-in-tame}).

\textbf{Why the frozen choices still suffice.} Ba{\l}aban's theorems concern a single real coupling at each step (from these numbers the profiles build the function $g_k^2(x)$ of 119 (1988), (2.24), which coincides with $g_k^2$ on the last domain), a background that minimises the action appearing in the exponent, and a single coupling-constant counterterm in each step. For this reason the source is not built into that framework but is carried alongside it. In terms of the running shift $\zeta_k(\cdot;w)$ of \cref{sec:sources-analytic-dependence-on-the-source-locality-and-decay-in-tame},

\begin{gather*}
\FitMath{-A\big(x_k(\cdot)-\zeta_k(\cdot;w),U\big)=-A\big(x_k(\cdot),U\big)+\sum_{y\in\mathbb{T}^{(k)}}A\big(h^{(k)}_y\zeta_k(\cdot;w),U\big),}
\end{gather*}

where $h^{(k)}_y$ denote Ba{\l}aban's tent functions (Comm. Math. Phys. 119 (1988), (3.40), (3.65)). On the set where $|\partial\mathbf{U}-\mathbf{1}|<\alpha_{0,k}\eta_k^2$ one has

\begin{gather*}
\FitMath{\big|A(h^{(k)}_y\zeta_k(\cdot;w),\mathbf{U})\big|\le C\alpha_{0,k}^2\sup|\zeta_k|\le C\varrho\,g_k^2(\log x_k)^{2q}\ll1:}
\end{gather*}

it is merely an additional local term, small and analytic, once $g$ is chosen small after $\varrho$. When the fields are real, $|e^{-A(x_k(\cdot)-\zeta_k,U)}|\le e^{-(1-C\varrho g^2)A(x_k(\cdot),U)}$; hence a failed test still costs $e^{-cp_0(g_k)^2}$.

\begin{proposition}[{Reality and positivity at real sources}]\label{prop:source-domains-4-5} All the choices are real, and so are the $f_i$; consequently each move takes real data to real data. Hence, at real $(c,z)$, each of the following quantities is real: the boundary terms, the increments, the activities $E^{(j),w},R^{(j),w}$ (at real fields), the vacuum constants, the running shifts, and also the $e_j(X;w)$. The exponents are therefore real at real block fields. Suppose moreover that the source is a real constant with $|\zeta|<2\varrho$. Then, in the operations $\mathbf{T}_k(Z)$ and $\mathbf{T}'_i(X)$ of \cref{sec:sources-analytic-dependence-on-the-source-locality-and-decay-in-tame}, every exponential factor is strictly positive at real pairs, and

\begin{gather*}
\FitMath{x_k(\cdot)-\zeta_k(\cdot;\zeta)>0,} \\[3pt]
\FitMath{\tau^\zeta_h\ge0:}
\end{gather*}

thus the stability bounds of the run rely on a positivity that is lost at complex $w$ and is still available at real shifts. Furthermore, $\rho^\zeta_0$ coincides with the bare density at $\beta-\zeta$, written with the choices belonging to $\beta$: these are two presentations of a single measure \cref{sec:flow-frozen-choices-slopes-and-tuned-bare-couplings}. The choices of a reference run do not depend on $\beta$. Take a real $\beta$ whose run has room $2\varrho$ in the tube, and a real $\zeta$ with $|\zeta|<\varrho$; along such a run one then obtains $x^\natural_j(\beta-\zeta)=x^\natural_j(\beta)-\zeta_j(\zeta)$ and $b^\natural_j(\beta-\zeta)=b_j(\zeta)$. These facts enter the runs with real shifts in \cref{sec:flow-frozen-choices-slopes-and-tuned-bare-couplings} and in \cref{ch:comparison-capsule-v7}, as well as the monotone map that sends the bare coupling to the renormalised coupling \cref{candidate:flow,ch:comparison-capsule-v7}. The conjugated copies of \cref{candidate:sources,ch:all-state-source-transfer} and \cref{thm:all-state-source-transfer,thm:graded-rooted-source} make no use of these facts; holomorphy is all they require (with the common-normalizer condition in \eqref{eq:all-state-source-transfer-ST11}).
\end{proposition}

\paragraph*{Changing the real constant in the same frame.}
Let the reference run at $\beta$ have room $2r$ through the
prefix and take a real $|v|<r$. At the bare level the sourced
input is the unsourced one at $\beta-v$. All choices belong
to the fixed frame. Suppose the two descriptions agree
through level $k$. The ordinary formation at level $k$
can then be applied to both. The a priori bound by $r/3$
on each increment beyond its common one-loop term gives
\[
 \big|\theta_{k+1}-x^\natural_{k+1}(\beta-v)\big|
 \le \mathsf t_{k+1}-2r+|\zeta_k(v)|+2r/3
 <\mathsf t_{k+1} .
\]
Use $|\zeta_k(v)|\le4|v|/3$ here. This establishes the next
admission before comparing the output coefficients. The
source-independent kernels act on the same bare input in
the two descriptions. Apply the same-frame comparison of
presentations from
\cref{sec:flow-frozen-choices-slopes-and-tuned-bare-couplings}.
It compares the centered totals on the admitted hereditary
windows for the common companion. The pointed quadratic
readout of those totals uses the same physical convention.
This gives $b_{k+1}^\natural(\beta-v)=b_{k+1}(v)$ and, by
the update of the action weight,
$x_{k+1}^\natural(\beta-v)=x_{k+1}^\natural(\beta)-\zeta_{k+1}(v)$.

The same comparison identifies the unsourced vacuum constant
at $\beta-v$. In the all-small-field term, the two densities
are continuous on a common open neighbourhood of the identity
in the block field. They agree as measures there. The
normalized action and activity totals vanish at the trivial
background. Evaluate the two densities at the identity, where
all the tests pass, to obtain
\[
 E_{(k)}[\beta-v]=E_{(k)}[\beta]+\sum_{j\le k}\sum_Xe_j(X;v).
\]
This uses equality of the small-field term on the neighbourhood,
with the same Haar normalization. Equality of the total
partition functions would not determine the individual vacuum
constants. A comparison in a different frame needs its own
overlap and coupling estimate.

For the normalization at an intermediate witness $k_0$,
sum the constant-source increment bounds only through that
witness. The regular term costs
$C\theta_j^{-3/2}(\log\theta_j)^{\mathfrak p_0}$ at level $j$.
The marked term costs $Cr^{-1}\theta_j^{-\nu}$ by Cauchy's
estimate on the composed source disc. Along the reference
flow, $\theta_j\ge\theta_{k_0}+c(k_0-j)-C_0$ on this prefix.
The sums give, with $\nu=(\kappa_0-4)/2>1$,
\[
 |\zeta_{k_0}'(0)-1|\le C\left[
  \theta_{k_0}^{-1/2}(\log\theta_{k_0})^{\mathfrak p_0}
              +r^{-1}\theta_{k_0}^{1-\nu}\right] .
\]
This bound uses the actual prefix, at the same normalization
as the physical source derivatives.

\textbf{The scope of the mechanism.} Because $\mathcal O(f)=A(f,U)$ is a weighted action, its source is a coupling shift. A different bounded gauge-invariant local observable $\Phi$ need not be a weighted plaquette action. The exact insertion identity can transport $e^{z\Phi}$, but does not put it in the variable-Wilson format or supply its counterterms, source stability or uniform correlation estimates. Those require a separately closed source construction.


\paragraph*{Exact insertion of block observables: a coupling of bare field, history and block field}

\textbf{(i) The setting and the objects.}

\emph{Purpose of the block.} The effective density of level $k$ decomposes into a finite sum of positive measures indexed by the histories $h$ of level $k$: $\rho_k\,dV_k=\sum_h\tau_h$. A single application of $\mathbf T$ integrates $F\circ\mathcal M$ exactly against the full density one level lower, so after $k$ applications one has $\int F\,d(\mathbf T^k\rho_0)=\int F(\mathcal M^kU)\,\rho_0(U)\,dU$. In the case where $\mathbf R$ acts nowhere, $\rho_k\,dV_k$ is the image of $\rho_0\,dU$ under $\mathcal M^k$, and the assertion of this block then amounts to the definition of an image measure. Where $\mathbf R$ acts on a component, a reference law takes the place of the law of the block field inside that component. This modifies how $\mathcal M^kU$ and $V_k$ are jointly distributed on bonds close to that component, and no mass changes. What this block asserts is that the replacements fit together into a single probability measure, defined on triples (bare field, history, block field), with respect to which, away from the components acted upon, the block field coincides with the block average of the bare field. No small factor and no constant enters Statements 1 to 6.

\emph{Spaces.} Each $\mathcal V_j$ is a finite product of copies of the compact group $G$, and is therefore a compact metrisable space. Finite sets are given the discrete topology, and all measures are Borel measures.

\emph{Terms as the measures produced by moves.} In parts \cref{sec:sources-the-curvature-source-with-fixed-reference-data} and \cref{candidate:sources} the construction is regarded as a chain of moves, a move being a finite family of positive kernels with one kernel for each outcome. The outcomes of the moves of the steps $1,\dots,j$, listed in order, form a string, and such a string is called a history of level $j$. We write $\mathsf K_h(U,\cdot)$ for the kernel obtained by composing along the string $h$; it runs from $\mathcal V_0$ to the states reached after step $j$. The admissible histories of level $j$, namely the strings for which $\mathsf K_h\ne0$, form the finite set $\mathfrak H_j$, and $\mathfrak H_0=\{\varnothing\}$. The term of a history $h\in\mathfrak H_j$ is defined as the measure into which its string transforms the bare measure, taken as a measure in the block field. Its mass is $\mathsf t(h)$:

\begin{gather*}
\FitMath{\tau_h(\mathcal B):=\int\rho_0(U)\,\mathsf K_h\big(U,\{V_j\in\mathcal B\}\big)\,dU,} \\[3pt]
\FitMath{\tau_\varnothing=\rho_0\,dU,} \\[3pt]
\FitMath{\mathsf t(h):=\tau_h(\mathcal V_j).}
\end{gather*}

The statements and the proofs below use the terms only through this definition.

\emph{The two stages of a history.} Given $h'\in\mathfrak H_{j+1}$, we write $h'|_j\in\mathfrak H_j$ for the sub-list of those entries of $h'$ that come from the steps and from the large-field operations of the levels $\le j$. The strings at the stage reached once the moves of $\mathbf T$ of step $j+1$ have been made form a finite set, $\mathfrak H^{\mathbf T}_{j+1}$. Its elements are the lists $a=(h;P,Q,\Xi)$. Each is formed by some $h\in\mathfrak H_j$ together with the outcomes of those three moves of $\mathbf T$ which are built from tests:
\begin{itemize}
\item $P=P_{j+1}$, the outcome of the plaquette tests (at step $1$, the pair $(P_0,P_1)$);
\item $Q=Q_{j+1}$, the outcome of the approximate fluctuation tests;
\item $\Xi=\Xi_{j+1}$, the outcome of the fluctuation tests together with the recursion along the boundary, which determines $\Omega_{j+1},\Lambda_{j+1},S_{j+1}$.
\end{itemize}

For such an $a$ we set $a|_j:=h$. We let $\tau^{\mathbf T}_a$ be the measure produced by the string $a$, taken in the variable $V_{j+1}$, so that it is a measure on $\mathcal V_{j+1}$. The class of the large-field operation of level $j+1$ is denoted $\mathcal Q_{j+1}(a)$; it is a set of components whose argument is a string, and it is distinct from the constraint map $\mathcal Q_k(U;A)$ introduced in the Notation. It is determined by the entries of $a$ and consists of every component of $Z_{j+1}$ that satisfies the class conditions (part \cref{sec:large-field-stored-variables-and-integration-in-a-large-field-patch}). The union of its members is

\begin{gather*}
\FitMath{\mathsf Z(a):=\bigcup_{X\in\mathcal Q_{j+1}(a)}X\subset T_\ell .}
\end{gather*}

The components that leave the class in the course of the operation (the drop-outs) are contained in $\mathsf Z(a)$, and so are the second-class components, although both kinds remain live. For $h'\in\mathfrak H_{j+1}$, the sub-list $h'^{\,\mathbf T}\in\mathfrak H^{\mathbf T}_{j+1}$ consists of $h'|_j$ together with the outcomes $P,Q,\Xi$ of step $j+1$. Every other entry of $h'$ is an outcome or a class produced within the large-field operation of level $j+1$.

\emph{The region of action.} Let $h\in\mathfrak H_j$. We denote by $Y(h)\subset T_\ell$ the union of those components of the large-field regions of $h$ which have been acted on by the large-field operation at one of the levels $\le j$; equivalently, it is the union of the class components of these levels. In particular $Y(\varnothing)=\emptyset$, and

\begin{gather*}
\FitMath{Y(h')=Y(h'|_j)\cup\mathsf Z\big(h'^{\,\mathbf T}\big)} \\[3pt]
\FitMath{(h'\in\mathfrak H_{j+1}).}
\end{gather*}


\emph{Clear and far bonds.} To a bond $c$ of $\mathbb T^{(j)}$ we attach the set of bare sites $\mathrm{Bl}_j(c):=B^j(c_-)\cup B^j(c_+)$, where $B^0(y)=\{y\}$. We say that the bond $c$ is clear of a set $Y$ when $\mathrm{Bl}_j(c)\cap Y=\emptyset$. A bond $c$ of $\mathbb T^{(k)}$ is called far from $Y(h)$ when neither $\mathrm{Bl}_k(c)$ nor any bare bond joining its two blocks meets $Y(h)^\sim$. This set is the union of the enlarged sets $X^\sim$, one for each of the class components $X$ whose union is $Y(h)$. The enlargement adds to such a component one layer of cubes of $\pi_{j(X)}$, and $j(X)$ is the level at which the operation acted on $X$.

\emph{Gauge averages.} The $G$-valued gauge transformations of $\mathbb T^{(j)}$ form the group $\mathcal G_j$. It is a finite product of copies of $G$ and carries the Haar probability measure $dv$. It acts through $V\mapsto V^v$, with $(V^v)^{v_0}=V^{v_0v}$. Given $u\in\mathcal G_i$ with $i\le j$, we write $u|_j\in\mathcal G_j$ for the gauge transformation obtained from it by restriction to the sites of $\mathbb T^{(j)}$. A function on $\mathcal V_j$, or a measure on it, is called invariant if this action leaves it invariant. The same convention applies on the configurations over a set of bonds closed under reversal, and on $\mathfrak H\times\mathcal V_j$, where the group acts trivially on the finite set $\mathfrak H$. Given a bounded Borel function $F$ and a finite positive measure $\tau$ on $\mathcal V_j$, we define

\begin{gather*}
\FitMath{\bar F(V):=\int F(V^v)\,dv,} \\[3pt]
\FitMath{\int F\,d\bar\tau:=\int\bar F\,d\tau .}
\end{gather*}


\textbf{(ii) The statements.} Throughout, $\beta$ is real, $0\le s\le K$ and $k=K-s$. The run is assumed to satisfy, through level $k$, the inductive description of the effective densities (part \cref{sec:large-field-the-induction-closed-the-order-of-the-constants-the-sum}), with $x_0,\dots,x_k\ge\gamma^{-2}$, together with the two-sided flow bound (part \cref{sec:flow-the-run-its-coefficient-and-the-analytic-requirements}). Of all these data, the proofs use only the fact that the moves of the steps $1,\dots,k$ are formed, that is, they use nothing beyond a reference sequence and a torus that is large compared with the cubes of the levels $\le k$.

\textbf{1. The step, taken term by term.} Let $0\le j<k$ and $h\in\mathfrak H_j$. Then each $\tau^{\mathbf T}_a$ with $a|_j=h$ is a finite non-negative measure on $\mathcal V_{j+1}$. Moreover, every bounded Borel invariant $F$ on $\mathcal V_{j+1}$ satisfies the following identity (the transport identity of the stage $\mathbf T$):

\begin{gather*}
\FitMath{\sum_{a\in\mathfrak H^{\mathbf T}_{j+1}:\,a|_j=h}\int F\,d\tau^{\mathbf T}_a=\int(F\circ\mathcal M)\,d\tau_h .}
\end{gather*}

Now let $a$ be such that $a|_j=h$, and let $F$ be a bounded Borel invariant function on $\mathcal V_{j+1}$ which does not depend on the variables $V(b)$ for which $b$ is a bond of $\mathbb T^{(j+1)}$ having at least one end point in $\mathsf Z(a)$. Then (the transport identity of the stage $\mathbf R$)

\begin{gather*}
\FitMath{\sum_{h'\in\mathfrak H_{j+1}:\,h'^{\,\mathbf T}=a}\int F\,d\tau_{h'}=\int F\,d\tau^{\mathbf T}_a .}
\end{gather*}

After summation over $h\in\mathfrak H_j$, the first identity expresses the exactness of $\mathbf T$ on the whole density. Taking $F\equiv1$ in both identities gives

\begin{gather*}
\FitMath{\sum_{h'\in\mathfrak H_{j+1}:\,h'|_j=h}\mathsf t(h')=\mathsf t(h):}
\end{gather*}

thus the step divides the mass of a history among its continuations. The law of the histories (part \cref{candidate:sources,ch:all-state-source-transfer}) is based on this identity.

\textbf{2. Gauge averages.} Suppose a compact group $\mathcal G$ with Haar probability $dv$ acts continuously on a compact metrisable space $\mathcal X$. Define $\bar F(x)=\int F(v\cdot x)\,dv$ and $\int F\,d\bar\tau=\int\bar F\,d\tau$.
(a) The function $\bar F$ is bounded, Borel and invariant, with $\sup|\bar F|\le\sup|F|$, and $\bar F=F$ when $F$ is invariant. The measure $\bar\tau$ is finite, positive and invariant. It satisfies $\int F\,d\bar\tau=\int F\,d\tau$ for invariant $F$, and $\bar\tau(N)=\tau(N)$ for invariant Borel sets $N$.
(b) Every invariant finite positive measure $\nu$ satisfies $\int\bar F\,d\nu=\int F\,d\nu$. Hence, if two invariant finite positive measures give the same integral to every bounded Borel invariant function, they are equal.
(c) Let $r:\mathcal G\to\mathcal G'$ be surjective, and let $q:\mathcal X\to\mathcal X'$ be a Borel map with $q(v\cdot x)=r(v)\cdot q(x)$. Then the image under $q$ of any invariant measure is invariant, and $f\circ q$ is invariant whenever $f$ is invariant.

\textbf{3. Gluing.} Let $\mathcal X_1,\mathcal X_2,\mathcal Y$ be compact metrisable spaces and let $p_i:\mathcal X_i\to\mathcal Y$ be Borel maps. Let $\lambda_i$ be finite positive measures on $\mathcal X_i$ such that $(p_1)_*\lambda_1=(p_2)_*\lambda_2$. Then $\mathcal X_1\times\mathcal X_2$ carries a finite positive measure $\lambda$ with the following two properties: its marginals are $\lambda_1$ and $\lambda_2$, and $p_1(x_1)=p_2(x_2)$ for $\lambda$-almost every $(x_1,x_2)$.

\textbf{4. Clear bonds.} (a) One has $c_-,c_+\in\mathrm{Bl}_j(c)$. (b) Let $c$ be a bond of $\mathbb T^{(j+1)}$, and let $b$ be a bond of $\mathbb T^{(j)}$ whose two end points lie in the double block $B(c_-)\cup B(c_+)$. Then $\mathrm{Bl}_j(b)\subset\mathrm{Bl}_{j+1}(c)$. (c) Every bond of $\mathbb T^{(k)}$ that is far from $Y(h)$ is clear of $Y(h)$.

\textbf{5. A single step as a transport of the averaged terms.} Let $0\le j<k$. Then $(\mathfrak H_j\times\mathcal V_j)\times(\mathfrak H_{j+1}\times\mathcal V_{j+1})$ carries a finite positive measure $\mathsf m_{j+1}$ with the following properties:
(a) the marginal on $\mathfrak H_j\times\mathcal V_j$ equals $\sum_{h\in\mathfrak H_j}\delta_h\otimes\bar\tau_h$;
(b) the marginal on $\mathfrak H_{j+1}\times\mathcal V_{j+1}$ equals $\sum_{h'\in\mathfrak H_{j+1}}\delta_{h'}\otimes\bar\tau_{h'}$;
(c) for $\mathsf m_{j+1}$-almost every $((h,V),(h',V'))$ both of the following hold: $h'|_j=h$, and $V'(c)=\mathcal M(V)(c)$ for every bond $c$ of $\mathbb T^{(j+1)}$ that has no end point in $\mathsf Z(h'^{\,\mathbf T})$.

\textbf{6. Exact insertion of block observables.} We write $Z=Z_{K,\ell,\beta}=Z(0)$ and $\mu=\mu_{K,\ell,\beta}$, and $\mathbb E_{\mathsf P}$ denotes the expectation with respect to a probability measure $\mathsf P$.
(a) On $\mathcal V_0\times\mathfrak H_k\times\mathcal V_k$ there exists a probability measure $\mathsf P$ with the following four properties.
Its marginal on $\mathcal V_0$ equals $\mu$.
Its marginal on $\mathfrak H_k\times\mathcal V_k$ is built from the averaged terms:

\begin{gather*}
\FitMath{\bar\nu_k:=Z^{-1}\sum_{h\in\mathfrak H_k}\delta_h\otimes\bar\tau_h .}
\end{gather*}

Hence this marginal assigns to every bounded Borel function $\Phi(h,V)$ that is invariant in $V$ the integral $Z^{-1}\sum_h\int\Phi(h,V)\,\tau_h(dV)$, and the law of the histories of level $k$ is the marginal of $\mathsf P$ on $\mathfrak H_k$: for $\mathsf E_k\subset\mathfrak H_k$, $\mathsf P(h\in\mathsf E_k)=Z^{-1}\sum_{h\in\mathsf E_k}\mathsf t(h)=\tilde\nu(\mathsf E_k)$.
It is invariant under the maps $(U,h,V)\mapsto(U^u,h,V^{u|_k})$, $u\in\mathcal G_0$.
Conditionally on $(U,h)$, the third coordinate agrees with $\mathcal M^kU$ on every bond that is clear of $Y(h)$, and in particular on every bond far from $Y(h)$.
(b) For each bounded Borel invariant function $F$ on $\mathcal V_k$, let

\begin{gather*}
\FitMath{F^{(h)}(V):=\mathbb E_{\mathsf P}\big[F(\mathcal M^kU)\,\big|\,h,V\big]}
\end{gather*}

be chosen among the versions of the conditional expectation as one that is invariant in $V$. Versions with this invariance exist, and for each $h$ any two of them coincide $\tau_h$-almost everywhere. Each of them satisfies

\begin{gather*}
\FitMath{\big\langle F(\mathcal M^kU)\big\rangle_{K,\ell,\beta}=\frac{\sum_h\int F^{(h)}\,d\tau_h}{\sum_h\int d\tau_h},} \\[3pt]
\FitMath{|F^{(h)}|\le\sup|F|\,\,\tau_h\text{-almost everywhere},} \\[3pt]
\FitMath{\sum_h\int d\tau_h=Z .}
\end{gather*}

(c) Suppose $F$ depends only on the variables of bonds that are clear of $Y(h)$; this is the case in particular when its dependence is confined to bonds far from $Y(h)$. Then one has $F^{(h)}=F$ almost everywhere with respect to $\tau_h$. In particular, for every history in which the large-field operation has acted on none of the components, $F^{(h)}=F$. Without this restriction one has, in general, $F^{(h)}\ne F$.
Here the depth $s$ is arbitrary, $0\le s\le K$, and $k=K-s$.

\textbf{7. Bounded block observables.} Take $F$ as in Statement 6 (b) and level-$k$ cells $\mathfrak c_1,\dots,\mathfrak c_n$, and denote by $\mathsf E_k:=\bigcup_{i\le n}\mathfrak{lv}_k(\mathfrak c_i)\subset\mathfrak H_k$ the event that at least one of these cells is contained in the live region of level $k$. Then

\begin{gather*}
\FitMath{\Big|\big\langle F(\mathcal M^kU)\big\rangle_{K,\ell,\beta}-\frac{\sum_{h\notin\mathsf E_k}\int F^{(h)}\,d\tau_h}{\sum_h\int d\tau_h}\Big|=\Big|\frac{\sum_{h\in\mathsf E_k}\int F^{(h)}\,d\tau_h}{\sum_h\int d\tau_h}\Big|} \\[3pt]
\FitMath{\le\sup|F|\;\tilde\nu(\mathsf E_k)\le n\,C\,\sup|F|\,e^{-c_{\rm r}p_0(g_k)} .}
\end{gather*}

The final inequality is valid whenever the data satisfy the hypotheses of the one-cell bound of part \cref{candidate:sources,ch:all-state-source-transfer}, with $C$ and $c_{\rm r}$ the constants of that bound. It holds uniformly in $K$, $\beta$, $\ell$ and the positions of the cells. In this form the bound applies to an observable of bounded modulus that depends on the block variables of a bounded number of cells at its own scale. In the expectation of such an observable, the part contributed by the histories of $\mathsf E_k$ is smaller than any power of $g_k$.

\textbf{(iii) Proofs.}

\textbf{On 1.} All three assertions are contained in part \cref{candidate:sources}. There they are properties of the measures that the moves produce, taken one string at a time.
Positivity. All kernels of all moves are non-negative. Hence the measure into which a string turns $\rho_0\,dU\ge0$ is finite and non-negative.
The stage $\mathbf T$. In its first move, $\mathbf T$ sets $V_{j+1}:=\mathcal M(V_j)$, and no later move of $\mathbf T$ alters $V_{j+1}$. The reason is that the Faddeev--Popov draw uses gauge transformations supported on the sites $B(y)\setminus\{y\}$, and these act on $V_j$ without touching $V_{j+1}$. Hence at this stage the identity is valid for every bounded Borel $F$. As an equality of measures, $\sum_{a|_j=h}\tau^{\mathbf T}_a=\mathcal M_*\tau_h$.
The stage $\mathbf R$. Given a class component $X$, denote by $\Lambda^{\rm pr}\cap X$ its box. This box lies well inside $X$, and $\mathbf R$ integrates the block field over it. Denote by $\mathfrak b_X$ the set of bonds of $\mathbb T^{(j+1)}$ that have an end point in that box. The moves of $\mathbf R$ alter $V_{j+1}$ on the bonds of the sets $\mathfrak b_X$ and nowhere else. Hence at this stage the identity is valid for every bounded Borel $F$ that does not depend on these bonds. Each bond of $\mathfrak b_X$ has both end points in $X$. Therefore every function that depends on no bond with an end point in $\mathsf Z(a)$ belongs to this class. The property that matters is that $F$ does not depend on $\mathfrak b_X$. Gauge invariance by itself is not enough. A counterexample is the plaquette function that appears in two places: in the block of this part that treats the curvature observable as a source (in the proof concerning the steps as stochastic kernels: the two refinements), and in part \cref{candidate:sources} (stability; in the proof of Statement 2 on exactness: which functions are transported).
Why history by history. Each move, once summed over its outcomes, is a probability kernel at each state. For a move consisting of tests, exactly one outcome takes the value $1$. For the gauge fixing the property holds because the gauge fixing is carried out as a procedure, namely as a draw. Let $\mathfrak G\ge0$ be normalised by $\int\mathfrak G(V^u)\,du=1$. Then the measure

\begin{gather*}
\FitMath{(\tau\kappa_{\mathfrak G})(\mathcal B)=\int\tau(dV)\int du\,\mathfrak G(V^u)\,\mathbf 1_{\mathcal B}(V^u)}
\end{gather*}

has the same mass as $\tau$, for every measure $\tau$. The product $\mathfrak G\,\tau$, by contrast, has mass $\int\mathfrak G\,d\tau$. If $\tau$ is invariant under the group with respect to which $\mathfrak G$ is normalised, this coincides with the mass of $\tau$, and then $\tau\kappa_{\mathfrak G}=\mathfrak G\,\tau$. Suppose that the bare weight is invariant under gauge transformations and that the representatives transform covariantly. Then, history by history, the term on which a gauge-fixing move acts has this invariance (see the treatment of the steps as stochastic kernels at the beginning), and for such a term the draw coincides with the product. The invariance meant here is invariance under the group of the move. A term need not be invariant under the full gauge group of the level. For a measure without this invariance, replacing the draw by the product does not conserve the mass in general. The masses considered from here on are the ones produced by the draw, which preserves them for any bounded Borel weight. Likewise, the axial frames are drawn before the tests that use them. The test that follows the gauge fixing along the tree lacks gauge invariance. Therefore the mass $a_r$ that passes it depends on the frame $r$. Suppose the components that leave kept a single frame $e$ while the survivors were averaged over the frames. Then the mass after the reference laws would be off by $\bar a-a_e$, with $\bar a$ the average of the $a_r$, and already for $F\equiv1$ the identity of the stage $\mathbf R$ would not hold. Part \cref{candidate:sources} contains both points (exactness: frames first; the product form).
Why it does not suffice to have the identities in their form summed over the histories. Let $\mathcal V_0=\mathcal V_1=\mathcal V_2=\{0,1\}^2$, where bond $1$ is exterior and bond $2$ is the component. Let $\mathcal M=\mathrm{id}$, let the gauge group be trivial, and let $\mu$ be the uniform law on the diagonal. Level $1$ has two outcomes $a_1,a_2$, with $\tau^{\mathbf T}_{a_1}=\tau^{\mathbf T}_{a_2}=\mu/2$, and $\mathbf R$ acts on bond $2$ in $a_2$ alone. Thus $\tau_{a_1}=\mu/2$, while $\tau_{a_2}$ is the product of the marginal of $\mu/2$ on bond $1$ with the uniform law on bond $2$. At level $2$ no $\mathbf R$ acts, and every history has a single continuation, $a_i'$ for $a_i$. We set $\tau_{a_1'}:=\tau_{a_2}$, $\tau_{a_2'}:=\tau_{a_1}$. Every term is non-negative. The totals are transported exactly, and the masses are conserved history by history, $\mathsf t(a_i')=\mathsf t(a_i)$. Moreover, $\mathbf R$ replaces a term by the product of its exterior marginal with a probability law. On $\{h=a_1'\}$, however, $Y(a_1')=\emptyset$ forces $V=U$, which is a point of the diagonal. Yet $\tau_{a_1'}$ assigns mass $\frac14$ to the complement of the diagonal. Hence no $\mathsf P$ exists. The failure lies in the identity of the stage $\mathbf T$ applied to the one history $a_1$ by itself.
Below, the two identities are applied to invariant $F$ alone.

\textbf{Proof of 2.} (a) The map $(v,x)\mapsto F(v\cdot x)$ is bounded and Borel, so Fubini's theorem shows that $\bar F$ is Borel. Right invariance of Haar measure gives $\bar F(v_0\cdot x)=\int F(vv_0\cdot x)\,dv=\bar F(x)$. The functional $F\mapsto\int\bar F\,d\tau$ is positive and continuous under monotone limits. It therefore defines a measure $\bar\tau$. With $F_{v_0}(x)=F(v_0\cdot x)$, left invariance yields $\overline{F_{v_0}}=\bar F$, and this is the invariance of $\bar\tau$. The remaining assertions follow from $\bar F=F$ for invariant $F$, applied to indicators as well. (b) If $\nu$ is invariant, Fubini's theorem gives $\int\bar F\,d\nu=\int dv\int F(v\cdot x)\,\nu(dx)=\int F\,d\nu$. Hence $\nu$ is determined by the values it assigns to the integrals of the invariant functions $\bar F$. (c) This follows by a change of variables.
By this device, the proof converts the two identities of Statement 1 into equalities of measures, namely equalities between the gauge averages of the terms. The proof applies these identities to invariant functions only.

\textbf{Proof of 3.} We apply the disintegration theorem in the following form. Let $\mathcal X$, $\mathcal Y$ be compact metrisable, let $p:\mathcal X\to\mathcal Y$ be Borel, let $\lambda$ be a finite positive measure on $\mathcal X$, and let $\nu=p_*\lambda$. Then there exist probability measures $\kappa(y;\cdot)$ on $\mathcal X$, Borel in $y$, such that $\lambda(B)=\int\nu(dy)\,\kappa(y;B)$ and $\kappa(y;p^{-1}\{y\})=1$ for $\nu$-almost every $y$. Set $\nu=(p_i)_*\lambda_i$, and let $\kappa_i(y;\cdot)$ be a disintegration of this kind of $\lambda_i$ along $p_i$. Define

\begin{gather*}
\FitMath{\lambda(B):=\int\nu(dy)\,\big(\kappa_1(y;\cdot)\otimes\kappa_2(y;\cdot)\big)(B).}
\end{gather*}

The integrand is Borel in $y$, because the sets $B$ with this property form a Dynkin class that contains the products of Borel sets. We have $\lambda(B_1\times\mathcal X_2)=\int\nu(dy)\kappa_1(y;B_1)=\lambda_1(B_1)$, and the second marginal is treated in the same way. The set $\{p_1(x_1)=p_2(x_2)\}$ is Borel and contains $p_1^{-1}\{y\}\times p_2^{-1}\{y\}$. For $\nu$-almost every $y$, the latter set has full measure under $\kappa_1(y;\cdot)\otimes\kappa_2(y;\cdot)$.

\textbf{Proof of 4.} (a) Every site lies in its own block. (b) For a site $y$ of $\mathbb T^{(j+1)}$, its cube is the union, over the sites $x\in B(y)$ of $\mathbb T^{(j)}$, of their cubes. After intersection with the bare lattice this becomes $B^{j+1}(y)=\bigcup_{x\in B(y)}B^j(x)$. Consequently $B^j(b_\pm)\subset B^{j+1}(c_-)\cup B^{j+1}(c_+)$. (c) This follows from $Y(h)\subset Y(h)^\sim$.

\textbf{Proof of 5.} The measures $\tau_h$, $\tau_{h'}$, $\tau^{\mathbf T}_a$ are all finite and non-negative. Consequently the gauge average of each of them is a finite, positive and invariant measure.
Step 1 (the stage $\mathbf T$: Radon--Nikodym). Let $h\in\mathfrak H_j$ be fixed. The map $\mathcal M$ is covariant, $\mathcal M(V^v)=\mathcal M(V)^{v|_{j+1}}$, as well as Borel measurable, and $v\mapsto v|_{j+1}$ is a surjection of $\mathcal G_j$ onto $\mathcal G_{j+1}$. Therefore $\mathcal M_*\bar\tau_h$ is invariant, and $F\circ\mathcal M$ is invariant whenever $F$ is. Let $F$ be any bounded Borel function on $\mathcal V_{j+1}$. Applying the transport identity of the stage $\mathbf T$ to the invariant function $\bar F$, we obtain

\begin{gather*}
\FitMath{\sum_{a|_j=h}\int F\,d\bar\tau^{\mathbf T}_a=\sum_{a|_j=h}\int\bar F\,d\tau^{\mathbf T}_a=\int(\bar F\circ\mathcal M)\,d\tau_h=\int(\bar F\circ\mathcal M)\,d\bar\tau_h=\int\bar F\,d(\mathcal M_*\bar\tau_h)=\int F\,d(\mathcal M_*\bar\tau_h).}
\end{gather*}

Hence $\sum_{a|_j=h}\bar\tau^{\mathbf T}_a=\mathcal M_*\bar\tau_h$. In particular $0\le\bar\tau^{\mathbf T}_a\le\mathcal M_*\bar\tau_h$. Therefore $\bar\tau^{\mathbf T}_a=f_a\cdot\mathcal M_*\bar\tau_h$ for Borel functions $0\le f_a\le1$ with $\sum_{a|_j=h}f_a=1$ almost everywhere. Define $\mathsf m^{\mathbf T}_a$ as the image of $(f_a\circ\mathcal M)\,\bar\tau_h$ under $V\mapsto(V,\mathcal M(V))$. It is a measure on pairs $(V,W)\in\mathcal V_j\times\mathcal V_{j+1}$. The marginals of this measure are $(f_a\circ\mathcal M)\bar\tau_h$ and $\bar\tau^{\mathbf T}_a$, and almost surely $W=\mathcal M(V)$. What the identity transports is the sum over the outcomes and nothing more. It does not fix which portion of $\tau_h$ goes to which $a$, and the only splitting it supplies is the one by $f_a\circ\mathcal M$. Accordingly, under the coupling just constructed, the outcome $a$ is drawn with probability $f_a(\mathcal M(V))$, and it is not determined by the tests.
Step 2 (the stage $\mathbf R$: restriction off the class). Fix $a$. Let $\mathrm{Bd}(a)$ denote the set of bonds of $\mathbb T^{(j+1)}$ having no end point in $\mathsf Z(a)$. Let $\mathcal V(a)$ be the space of configurations on $\mathrm{Bd}(a)$, and let $\mathrm{res}_a:\mathcal V_{j+1}\to\mathcal V(a)$ be the restriction map, which commutes with the gauge action. Set $\lambda_a:=\sum_{h'^{\,\mathbf T}=a}\delta_{h'}\otimes\bar\tau_{h'}$. Let $f$ be a bounded Borel invariant function on $\mathcal V(a)$. Then $f\circ\mathrm{res}_a$ is invariant and depends on no bond with an end point in $\mathsf Z(a)$. The transport identity of the stage $\mathbf R$ therefore gives

\begin{gather*}
\FitMath{\sum_{h'^{\,\mathbf T}=a}\int(f\circ\mathrm{res}_a)\,d\bar\tau_{h'}=\sum_{h'^{\,\mathbf T}=a}\int(f\circ\mathrm{res}_a)\,d\tau_{h'}=\int(f\circ\mathrm{res}_a)\,d\tau^{\mathbf T}_a=\int(f\circ\mathrm{res}_a)\,d\bar\tau^{\mathbf T}_a .}
\end{gather*}

The measure $(\mathrm{res}_a)_*\bar\tau^{\mathbf T}_a$ and the image of $\lambda_a$ under $(h',V')\mapsto\mathrm{res}_a(V')$ are both invariant, so Statement 2 (b) shows that they are equal. The identity says nothing about the bonds with an end point in $\mathsf Z(a)$. For this reason the two measures are glued over $\mathcal V(a)$. The result is a measure $\mathsf m^{\mathbf R}_a$ on $\mathcal V_{j+1}\times(\mathfrak H_{j+1}\times\mathcal V_{j+1})$ with marginals $\bar\tau^{\mathbf T}_a$ and $\lambda_a$. Almost surely under it, $\mathrm{res}_a(V')=\mathrm{res}_a(W)$ and $h'^{\,\mathbf T}=a$.
Step 3 (composition). On the middle field $W$, the marginal of $\mathsf m^{\mathbf T}_a$ and the marginal of $\mathsf m^{\mathbf R}_a$ are both equal to $\bar\tau^{\mathbf T}_a$. We glue over $\mathcal V_{j+1}$ and discard the middle field. We then attach to $V$ the label $h=a|_j$ and sum over $a$. The resulting measure is $\mathsf m_{j+1}$. (a): $\sum_{a|_j=h}(f_a\circ\mathcal M)\bar\tau_h=\bar\tau_h$. (b): $\sum_a\lambda_a=\sum_{h'}\delta_{h'}\otimes\bar\tau_{h'}$, because each $h'$ has exactly one first stage. (c): almost surely $W=\mathcal M(V)$, $\mathrm{res}_a(V')=\mathrm{res}_a(W)$ and $h'^{\,\mathbf T}=a$, so that $h'|_j=a|_j=h$.

\textbf{Proof of 6.} For $0\le j\le k$ we define

\begin{gather*}
\FitMath{\mathsf A_j:=\big\{(U,h,V)\in\mathcal V_0\times\mathfrak H_j\times\mathcal V_j:\,V(b)=(\mathcal M^jU)(b)\text{ for every bond }b\text{ of }\mathbb T^{(j)}\text{ clear of }Y(h)\big\},}
\end{gather*}

which is a Borel set.
Step 1 (induction on the level). We show that, for each $0\le j\le k$, the space $\mathcal V_0\times\mathfrak H_j\times\mathcal V_j$ carries a probability measure $\mathsf P_j$ with the following three properties. Its projection to $\mathcal V_0$ is $\mu$. Its marginal on $\mathfrak H_j\times\mathcal V_j$ equals $Z^{-1}\sum_{h\in\mathfrak H_j}\delta_h\otimes\bar\tau_h$, which is the family of averaged terms. Finally, $\mathsf P_j(\mathsf A_j)=1$.
To begin the induction: the measure $\rho_0\,dU$ is invariant, so $\bar\tau_\varnothing=\tau_\varnothing$, and we take for $\mathsf P_0$ the push-forward of $\mu$ under $U\mapsto(U,\varnothing,U)$.
Passage from $j$ to $j+1$: by the second property together with Statement 5 (a), the measures $\mathsf P_j$ and $Z^{-1}\mathsf m_{j+1}$ induce the same marginal on $\mathfrak H_j\times\mathcal V_j$. We glue the two over this space and project the result onto $(U,h',V')$, which gives $\mathsf P_{j+1}$. The first two properties carry over, and Statement 5 (b) is what carries over the second. As for the third, with probability one we have $(U,h,V)\in\mathsf A_j$, $h'|_j=h$, and $V'(c)=\mathcal M(V)(c)$ outside $\mathsf Z(h'^{\,\mathbf T})$. Take a bond $c$ clear of $Y(h')=Y(h)\cup\mathsf Z(h'^{\,\mathbf T})$. Both of its end points belong to $\mathrm{Bl}_{j+1}(c)$, so neither lies in $\mathsf Z(h'^{\,\mathbf T})$, and therefore $V'(c)=\mathcal M(V)(c)$. For each bond $b$ of the double block of $c$ we have $\mathrm{Bl}_j(b)\subset\mathrm{Bl}_{j+1}(c)$. Hence such a bond is clear of $Y(h)$, and $V(b)=(\mathcal M^jU)(b)$. Because $\mathcal M(V)(c)$ depends on $V$ only through these bonds, we obtain

\begin{gather*}
\FitMath{V'(c)=\mathcal M(V)(c)=\mathcal M(\mathcal M^jU)(c)=(\mathcal M^{j+1}U)(c).}
\end{gather*}

Step 2 (symmetrisation). The group $\mathcal G_0$ acts on triples by $(U,h,V)\mapsto(U^u,h,V^{u|_k})$. Averaging $\mathsf P_k$ over this action gives a measure $\mathsf P$, and it is invariant under the action. Since $\mu$ and $\bar\nu_k=Z^{-1}\sum_h\delta_h\otimes\bar\tau_h$ are both invariant, $\mathsf P$ still has the first two properties. Covariance gives $\mathcal M^k(U^u)=(\mathcal M^kU)^{u|_k}$, so the condition $V(b)=(\mathcal M^kU)(b)$ holds exactly when $V^{u|_k}(b)=(\mathcal M^k(U^u))(b)$ does. Thus the set $\mathsf A_k$ is invariant, and $\mathsf P(\mathsf A_k)=\mathsf P_k(\mathsf A_k)=1$. Let $\Phi$ be a function and $N$ a Borel set, both invariant in $V$, with sections $N_h$. Then Statement 2 (a) gives

\begin{gather*}
\FitMath{\int\Phi\,d\bar\nu_k=Z^{-1}\sum_h\int\Phi(h,V)\,\tau_h(dV),} \\[3pt]
\FitMath{\bar\nu_k(N)=Z^{-1}\sum_h\tau_h(N_h).}
\end{gather*}

This is why the statement is restricted to invariant $\Phi$ and $F$. The marginal of $\mathsf P$ consists of the averaged terms, and the averaged terms coincide with the terms only on invariant functions and invariant sets.
Step 3 (assertion (a)). The marginals and the invariance were obtained in Step 2. A function of $h$ alone is invariant in $V$, and this yields the law of the histories. Because $\mathsf P(\mathsf A_k)=1$, for almost every $(U,h)$ a regular conditional distribution of $V$ given $(U,h)$ is concentrated on the configurations that agree with $\mathcal M^kU$ on every bond clear of $Y(h)$, and so in particular on every far bond.
Step 4 (invariant versions). Suppose that $F$ is invariant, and set $\mathsf f(U):=F(\mathcal M^kU)$. This function is invariant under $\mathcal G_0$. Take an arbitrary version $\Psi$ of $\mathbb E_{\mathsf P}[\mathsf f(U)|h,V]$ with $|\Psi|\le\sup|F|$, so that $\mathbb E_{\mathsf P}[\mathsf f(U)\phi(h,V)]=\int\Psi\phi\,d\bar\nu_k$ for every bounded Borel $\phi$. Each $v\in\mathcal G_k$ is of the form $u|_k$, where $u\in\mathcal G_0$ takes the value $v(y)$ on the block $B^k(y)$. Using the invariance of $\mathsf P$, of $\mathsf f$ and of $\bar\nu_k$, we get

\begin{gather*}
\FitMath{\mathbb E_{\mathsf P}[\mathsf f(U)\phi(h,V)]=\mathbb E_{\mathsf P}[\mathsf f(U^u)\phi(h,V^v)]=\mathbb E_{\mathsf P}[\mathsf f(U)\phi(h,V^v)]=\int\Psi(h,V)\phi(h,V^v)\,d\bar\nu_k=\int\Psi(h,V^{v^{-1}})\phi(h,V)\,d\bar\nu_k .}
\end{gather*}

The first equality holds because $\mathsf P$ is jointly invariant. We integrate over $v$, apply Fubini's theorem and use the fact that Haar measure is invariant under inversion. It follows that $F^{(h)}(V):=\int\Psi(h,V^v)\,dv$ is also a version, and that it is Borel, invariant in $V$ and bounded by $\sup|F|$. Any two invariant versions differ only on an invariant Borel set $N$ with $\bar\nu_k(N)=0$. Since the terms are non-negative, this gives $\tau_h(N_h)=0$ for every $h$.
Step 5 (assertion (b)). Take an invariant version. Then

\begin{gather*}
\FitMath{\big\langle F(\mathcal M^kU)\big\rangle=\int\mathsf f\,d\mu=\mathbb E_{\mathsf P}[\mathsf f(U)]=\int F^{(h)}(V)\,\bar\nu_k(dh,dV)=Z^{-1}\sum_h\int F^{(h)}\,d\tau_h .}
\end{gather*}

The last equality uses the invariance of $F^{(h)}$ in $V$. Taken by itself, the defining property of a conditional expectation yields the integral against $\bar\nu_k$ and nothing more. No use can be made at this point of a non-invariant version. Consider a term in which some component has been put into the second class by the large-field operation at level $k$. By part \cref{candidate:sources}, such a term carries, in the level-$k$ block field, a Dirac distribution on the gauge tree of the box. Consequently $\tau_h$ and $\bar\tau_h$ are mutually singular for such a term. A version is determined $\bar\tau_h$-almost everywhere, and in this case the integral against $\tau_h$ of a version that is not invariant is then arbitrary. Taking $\Phi\equiv1$ gives $\sum_h\int d\tau_h=Z$. For every $h$, any invariant version coincides $\tau_h$-almost everywhere with the version constructed in Step 4, and that version is bounded by $\sup|F|$.
Step 6 (assertion (c)). Fix $h_0$, and suppose that $F$ depends only on bonds clear of $Y(h_0)$. On $\mathsf A_k\cap\{h=h_0\}$ we then have $F(\mathcal M^kU)=F(V)$. It follows that 

\begin{gather*}
\FitMath{\mathbb E_{\mathsf P}[\mathsf f(U)\mathbf 1_{\{h=h_0\}}\phi(V)]=\mathbb E_{\mathsf P}[F(V)\mathbf 1_{\{h=h_0\}}\phi(V)]}
\end{gather*}

 for every bounded Borel $\phi$. In other words, $F^{(h_0)}=F$ $\bar\nu_k$-almost everywhere on $\{h_0\}\times\mathcal V_k$. The set on which this equality fails is invariant in $V$ and is therefore a $\tau_{h_0}$-null set. When the operation has acted on no component of $h_0$, we have $Y(h_0)=\emptyset$, and then every bond is clear.

\textbf{Proof of 7.} The equality comes from splitting the numerator of Statement 6 (b) into the part over $\mathsf E_k$ and the part over its complement. We have $\tau_h\ge0$, and $|F^{(h)}|\le\sup|F|$ holds $\tau_h$-almost everywhere. Hence $|\int F^{(h)}\,d\tau_h|\le\sup|F|\,\mathsf t(h)$. Summing this over $h\in\mathsf E_k$ and dividing by $Z=\sum_h\mathsf t(h)$ gives $\sup|F|\,\tilde\nu(\mathsf E_k)$. Subadditivity together with the one-cell bound $\tilde\nu(\mathfrak{lv}_k(\mathfrak c))\le C\,e^{-c_{\rm r}p_0(g_k)}$ of part \cref{candidate:sources,ch:all-state-source-transfer} yields $\tilde\nu(\mathsf E_k)\le n\,C\,e^{-c_{\rm r}p_0(g_k)}$. For a history outside $\mathsf E_k$, no one of the cells lies in a large-field region that is live at level $k$.

\textbf{(iv) What is used.}

\emph{From other parts.} The following results of part \cref{sec:averaging-block-averages-gauge-fixing-and-the-renormalisation-group-of-this} are used. The map $\mathcal M$ is Borel. It is gauge covariant: $\mathcal M(V^v)=\mathcal M(V)^{v'}$, where $v'$ is obtained from $v$ by restriction to the sites of the coarser lattice. It is local: $\mathcal M(V)(c)$ depends on $V$ only through the bonds of the double block of $c$. From part \cref{candidate:sources} we use three things. The first is the description of the construction as a chain of moves. The second is that every kernel is non-negative, so that every $\tau_h$ and every $\tau^{\mathbf T}_a$ is a finite non-negative measure. The third is the pair of transport identities of Statement 1. From part \cref{sec:large-field-stored-variables-and-integration-in-a-large-field-patch} we use the class of the large-field operation, which enters through the recursion $Y(h')=Y(h'|_j)\cup\mathsf Z(h'^{\,\mathbf T})$. Statements 1 to 6 use nothing further from the construction: no bound, no analyticity, and no form of a term. Only Statement 7 draws in addition on part \cref{candidate:sources,ch:all-state-source-transfer}, through the one-cell bound there on the law of the histories.

\emph{Where it is used.} Statement 7 is obtained by combining Statement 6 with the one-cell bound of part \cref{candidate:sources,ch:all-state-source-transfer}. The other parts draw on the identity $\sum_{h'|_j=h}\mathsf t(h')=\mathsf t(h)$, which is the last line of Statement 1. This identity makes the laws of the histories of the levels $j\le K$ a consistent family of probability measures. That family is the law $\tilde\nu$ of part \cref{candidate:sources,ch:all-state-source-transfer}. The bounds on large fields rest on it, and through them so do the laws at one scale. For $k\ge1$ the observables of the theorem do not have the form $F(\mathcal M^kU)$. Indeed, $\mathcal O(f)=\sum_pf(x(p))\,a_p(U)$ is a function of the bare plaquettes, and Statement 6 cannot be applied to it. These observables are carried as sources (first block of this part). A single scale is isolated by exactness with sources at depth $s$:

\begin{gather*}
\FitMath{\sum_{h\in\mathfrak H_{K-s}}\tau^w_h(\mathcal V_{K-s})=Z(w),} \\[3pt]
\FitMath{\text{whence, at }c=0,} \\[3pt]
\FitMath{\Big\langle\exp\sum_iz_i\,\mathcal O(f_i)\Big\rangle_{K,\ell,\beta}=\frac{Z(w)}{Z(0)}}
\end{gather*}

and this quantity is an expectation, taken over the histories, the block fields and the stored variables at level $K-s$, of the exponential of the sourced exponent minus the unsourced one.

\emph{Against Ba{\l}aban's papers.} Ba{\l}aban presents the operation $\mathbf R$ in a simplified form in the Introduction of Comm. Math. Phys. 122 (1989) 175, formula (0.3) on p. 176: 
\begin{gather*}
\FitMath{(\mathbf R\rho)(V)=\sum_Z\rho(Z'',V)\int dV|_{Z'}\,\rho(Z,V)\big/\int dV|_{Z'}\,\rho(Z'',V).}
\end{gather*}
 This form has the normalisation property $\int dV\,(\mathbf R\rho)(V)=\int dV\,\rho(V)$ ((0.4)). In the procedure as actually carried out, the denominators integrate parts of the densities that are localised near $Z'$ (p. 177), and the normalisation property takes the form (1.102), p. 201. The objects treated in Ba{\l}aban's papers are densities. On observables, he states that his Theorem 1 ``is the basis for many other applications, for example, for an analysis of expectation values of physical observables, like loop variables, averaged loop variables'', and calls these problems that ``deserve detailed analysis and further publication'' (Comm. Math. Phys. 122 (1989) 355, p. 356).


\paragraph{Analytic dependence on the source; locality and decay in tame and in wild regions}
\label{sec:sources-analytic-dependence-on-the-source-locality-and-decay-in-tame}

\textbf{Objects.} We write $w(p):=w(x(p))$. When the source is a constant shift $w\equiv\zeta$, the increments $b_j(\zeta)$ are computed on the infinite lattice in the same way as $b_j$. They therefore do not depend on the torus, and they are real for real $\zeta$. The running shift is $\zeta_k=\zeta+\sum_{j\le k}[b_j(\zeta)-b_j(0)]$, and the inverse couplings are $x_k-\zeta_k$. For a general source, let $p_y$ denote a bare plaquette at $y$. Then

\begin{gather*}
\FitMath{\zeta_k(p;w)=w(p)+\sum_{j\le k}\varphi_j(p)\sum_{y\in\mathbb{T}^{(j)}}h^{(j)}_y(p)\,\big[b_j(w(p_y))-b_j(0)\big]:}
\end{gather*}

in other words, the coupling counterterm of level $j$ is applied locally, and its coefficient is taken at the local value of the source.

\begin{proposition}[{The format with sources}]\label{prop:source-domains-4-6} The density $\rho^w_k$ has the form of \cref{sec:large-field-the-induction-closed-the-order-of-the-constants-the-sum}, with $x_k(\cdot)$ replaced by $x_k(\cdot)-\zeta_k(\cdot;w)$. Its activities $E^{(j),w},R^{(j),w},B^{(j),w}$ are holomorphic in $(c,z)$, depend on $w$ only through its values on $X^{\sim1}$, and obey the bounds valid at $w=0$ up to constants. Moreover $|\zeta_k(p;w)|\le\frac43\sup|w|$, and for every constant $|\zeta_*|<2\varrho$

\begin{gather*}
\FitMath{\big|E^{(j),w}(X,\cdot)-E^{(j),\zeta_*}(X,\cdot)\big|\le E_0\,\varrho^{-1}\mathsf{m}_j\,e^{-\kappa d_j(X)},} \\[3pt]
\FitMath{\mathsf{m}_j:=\max\Big\{\sup_{X^{\sim1}}|w-\zeta_*|,\,CML^{j-K}\operatorname{Lip}w\Big\}.}\qedhere
\end{gather*}



\end{proposition}

\textbf{Strategy: induction on the levels, on a complex polydisc of data.} On a complex polydisc of data (activities, boundary terms, shifts) centred at the real run, a single step (the cluster expansion of \cref{ch:local-formation-program} together with $\mathbf{R}$ of \cref{sec:large-field-stored-variables-and-integration-in-a-large-field-patch}--\cref{candidate:large_field}) is local, analytic and uniformly bounded. This is the step on the complex class of \cref{sec:large-field-the-induction-closed-the-order-of-the-constants-the-sum}. Its small-field core is the holomorphic map on local families required by \hyperref[hyp:B4]{the ordinary-step hypothesis}. The sourced data remain inside the polydisc, so holomorphy follows by composition. Locality follows because an activity at $X$ depends on the data, and on the outcomes of the history, only near $X$ (locality in the history, \cref{sec:large-field-the-tests-the-nested-regions-and-the-content-of-a}).

\emph{The marginal term.} When the source is constant, the step preserves the symmetries of \cref{ch:pointed-hessian-cancellation}, so once $b_{k+1}(\zeta)A$ has been extracted only irrelevant terms remain. The estimate is the one of \cref{sec:flow-frozen-choices-slopes-and-tuned-bare-couplings}. There is a remainder $\le Cg_k^{\kappa_0-4}$, coming from the $R^{(j)}$, which is holomorphic in $\zeta$ on a disc of radius $C\varrho$. Apart from this remainder, the increment $b_{k+1}$ is holomorphic in the shifts $\zeta_j$ of the earlier levels on discs of radius $cx_j$, where $c\le\sigma_\star$. These discs make up the polydisc of \cref{ch:local-formation-program}, taken along the run's own couplings. On these discs its non-constant part, less the remainder, is bounded by $Cg_k(\log x_k)^{\mathfrak{p}_0}$. Cauchy's estimate together with $x_k\ge g^{-2}+c(K-k)$ yields

\begin{gather*}
\FitMath{|b_{k+1}(\zeta)-b_{k+1}(0)|\le\mathsf{a}_{[K-k-1]}|\zeta|,} \\[3pt]
\FitMath{\sum_{s<K}\mathsf{a}_{[s]}\le C\sum_k\big[g_k^3(\log x_k)^{\mathfrak{p}_0}+\varrho^{-1}g_k^{\kappa_0-4}\big]\le\tfrac16,}
\end{gather*}

where the quantity $\mathsf{a}_{[s]}$ carries a depth index and denotes, for the step that arrives at depth $s$, the Lipschitz number of its increment (the same letter is used in \cref{sec:flow-frozen-choices-slopes-and-tuned-bare-couplings} and \cref{candidate:flow,ch:comparison-capsule-v7}). Consequently $|\zeta_k^{\prime}(0)-1|\le\frac16$. On the smaller real interval $|\zeta|\le\varrho$ one obtains $\zeta_k^{\prime}(\zeta)\in[\frac23,\frac43]$: on $|\zeta|<2\varrho$ the holomorphic quotient $[b_j(\zeta)-b_j(0)]/\zeta$ is bounded by $\mathsf a_{[K-j]}$, and Cauchy's estimate on the half-disc gives $|b_j^{\prime}(\zeta)|\le2\mathsf a_{[K-j]}$. The sum was chosen at most $1/6$ for this fixed loss. No derivative bound throughout the full open source disc follows from the radial inequality alone. On $X$ a general source is compared with $\zeta_*=w(p_*)$, where $p_*$ is a bare plaquette of $X^{\sim1}$. When $\mathsf{m}_j<\varrho$, the slope displayed above comes from the Schwarz lemma, applied along a line of data and not along a line of sources. This line is a family, holomorphic in $t$, of data of the level below; at $t=0$ it coincides with the data of $\zeta_*$, and at $t=1$ it coincides on $X^{\sim1}$ with those of $w$. For $|t|<\varrho/\mathsf{m}_j$ the family remains in the convex polydisc. When $\mathsf{m}_j\ge\varrho$, the slope is simply the sup bound. Subtract from each of the two activities the normaliser part $\mathbf{N}^{(j)}$ of \cref{ch:local-formation-program}, which is common to both; what is left of each is at most $\frac12E_0e^{-\kappa d_j(X)}$. The slope is thus a clause that the induction carries along. The defect estimate uses the sharper, normalizer-subtracted bound of \hyperref[source:local-data-lines]{the local data-line argument}. The induction carries similar slopes for $R^{(j)}$ and $B^{(j)}_{\mathcal S}$, and also, among the operations treated below, for the excess exponents $\Psi_\omega$. The slope does not follow, up to the boundary of the class of sources $\{\|w\|_\flat<2\varrho\}$, from holomorphy in $w$, from the bound, from locality and from the agreement at constant sources. Moreover, $\mathsf{m}_j\le e^{d_j(X)}L^{-(K-j)}\|w\|_\flat$. The part that no counterterm renormalises therefore carries a factor $L^{-(K-j)}$, which is summable over the scales. In short, a gradient of the source is irrelevant.

\emph{Wild regions.} Two kinds of operations act on the sourced exponent. The first kind consists of the operations $\mathbf{T}_k(Z)$ on the components on which $\mathbf{R}$ has not acted. These are built from the choices of the run only, namely delta functions, tests, gauge-fixing factors and conditional Gaussian weights, and they contain no Wilson weight \cref{sec:large-field-stored-variables-and-integration-in-a-large-field-patch}. On such components the source appears in the function on which these operations act, that is, in the Wilson weight and in the boundary terms. Taken on its own, $\mathbf{T}_k(Z)1$ carries no small factor: for a component produced by failed plaquette tests, this quantity is $\le\exp[C\sum_{t\le k}\#_t(Z_t\cap Z)]$. Three places carry the small factor \cref{sec:large-field-rare-patches-gains-per-creation-and-costs-per-site}: the Gaussian forms of the fluctuation fields, the kept gauge-fixing exponents and the Wilson weight. Only the last of these depends on the source; the first two are built from the choices of the run. The pointwise majorant of \cref{sec:large-field-rare-patches-gains-per-creation-and-costs-per-site} gives one factor $e^{-cp_0(g_t)^2}$ for each creation, at the coupling of that creation's own step, and one constant for each site. It is used in the form it has at $w=0$; the sourced exponent and the unsourced one are compared further below, on the supports of the tests.

\phantomsection\label{notdef:excess-source}
The second kind concerns a component $X$ on which $\mathbf{R}$ has acted at some level $i$ (whether it was healed, remained as an island, or dropped out because a test of a stage or of the main step failed). For such a component the source enters inside its operation $\mathbf{T}'_i(X)$. Except for a drop-out, this operation contains the Wilson factor of the core. When the component is healed or is an island, the operation coincides with that of Ba{\l}aban, Comm. Math. Phys. 122 (1989) 355, (1.71). Let $\omega$ run over the local histories of $X$, let $\xi$ be the stored variables and $\mu_\omega\ge0$ their measures, and let $\mathbf{p}=(\mathbf{U},\mathbf{J})$ be the pairs on $X^\sim$ at which a term evaluates the factor of $X$ \cref{sec:large-field-rare-patches-gains-per-creation-and-costs-per-site}. For each $\omega$, $e^{\Phi_\omega}$ is replaced by $e^{\Phi_\omega+\Psi_\omega}$, with an \emph{excess exponent}

\begin{gather*}
\FitMath{\Psi_\omega=A\big(\hat\zeta\,\phi_{\mathrm{c}}\mathbf{1}_X,U''(\xi)\big)+\big[V^w-V^0\big](X^\sim)+\sum_{\omega'}\Psi_{\omega'}:}
\end{gather*}

its summands are the Wilson term, weighted by the shift $\hat\zeta$ of the prepared regions ($|\hat\zeta|\le C\varrho$, and $0\le\phi_{\mathrm{c}}\le1$ is the cut-off on the old core, \cref{candidate:large_field}) and evaluated at the background $U''(\xi)$ of the inner region, which is constructed from the stored fields; the difference between the interaction terms $V$ of $\mathbf{R}$ with and without the source; and the exponents of the inner histories $\omega'$, namely those of the swallowed islands and drop-outs. The operation modified in this way is written $\mathbf{T}'_i(X;w)$. (In the case of a drop-out, the operation \cref{sec:large-field-rare-patches-gains-per-creation-and-costs-per-site} is taken to include the Wilson weight that the carrying term has over $X$. The operations that are kept are built from the choices of the run only, and none of their factors depends on the source, which appears in that Wilson weight and in the boundary terms. The exponent is then $\sum_{\omega'}\Psi_{\omega'}$ alone, the contribution of the swallowed islands, and it obeys the bound given below.) The exponent $\Psi_\omega$ is holomorphic in $(c,z)$, is real for real pairs and real sources, vanishes at $w=0$, and depends on $w$ only on $X^{\sim2}$. The background $U''$ is real, and its regularity holds only at the scale on which the respective layer was stored. Where $\phi_{\mathrm{c}}\ne0$, one has $|\partial U''-\mathbf{1}|\le C\check{R}_n^2\varepsilon_n\eta_n^2\le C\eta_n^2$ on $\Omega_n\setminus\Omega_{n+1}$, so the bound $a_p\le C\eta_n^4$ balances the $C\eta_n^{-4}$ plaquettes per site of level $n$ (on the bare layer $a_p\le2$, since $G$ is compact). Each layer therefore contributes $O(\varrho)$ for each site of its own scale:

\begin{gather*}
\FitMath{|\Psi_\omega|\le\bar{f}\,\#(\omega),} \\[3pt]
\FitMath{\bar{f}=C(1+\varrho),}
\end{gather*}

where, as in \cref{sec:large-field-the-induction-closed-the-order-of-the-constants-the-sum}, $\#(\omega)$ is the number of stored sites, each being counted at its own level. Since $\#(\omega)$ is additive over the inner histories, the constant after a step is the maximum of $\bar{f}$ and the constant of the new term, not their sum. Hence the constant does not increase with the level. The perturbation is extensive and must therefore be paid in advance. This is done by the mass bound of \cref{sec:large-field-rare-patches-gains-per-creation-and-costs-per-site}, \cref{sec:large-field-the-induction-closed-the-order-of-the-constants-the-sum}, which the induction maintains with a reserve weight $e^{c\mathsf{v}(\omega)}$, where $c=c_{\mathrm{r}}$ is the reserve constant of \cref{sec:large-field-the-induction-closed-the-order-of-the-constants-the-sum}, one fixed number (the same constant stands in every weight $e^{c\mathsf{v}}$ of this part except the raised one in the bound of the whole label, whereas in the exponents that contain $p_0$ the letter $c$ denotes a generic constant). Here the space-time volume $\mathsf{v}(\omega)$ counts $\log x_j$ for each stored site of level $j$:

\begin{gather*}
\FitMath{\mathsf{M}(i,X):=\frac1{|W_4|}\sum_\omega e^{c\mathsf{v}(\omega)}\int\sup_{\mathbf{p}}\big|e^{\Phi_\omega(\xi;\mathbf{p})}\big|\,\mu_\omega(d\xi)\le e^{-(2-c')p_0(g_i)},}
\end{gather*}

with the supremum over the complex pairs taken inside the integral and the sum. This form applies to a healed component and to an island. In the case of a drop-out, the integral contains the Wilson weight of the carrying term over $X$. The supremum is then taken over the evaluations at $\xi$, that is, over the real configurations which are compatible with the stored variables, are regular layer by layer, and give the listed tests their recorded values, with the exception of the test whose failure made $X$ a drop-out; these evaluations form a set fibred over $\xi$ \cref{sec:large-field-rare-patches-gains-per-creation-and-costs-per-site}. A product set contains a flat $\mathbf{U}$, at which $e^{-A}=1$, and the supremum over such a set removes the factor. Ba{\l}aban's bound is $\mathbf{T}'_k(X)1\le e^{-(2-c')p_0(g_k)}$ (ibid., (1.89), where Ba{\l}aban has $2(1+\beta_0)^{-1}$ in place of $2-c'$), for each fixed pair and without the weight. If $g$ is small enough once $\varrho$ has been fixed, then $c\log x_j\ge2\bar{f}$. Consequently $|e^{\Psi_\omega}|\le e^{c\mathsf{v}(\omega)}$, and for every bounded function $F$ of the stored variables and every evaluation of this kind

\begin{gather*}
\FitMath{\big|\mathbf{T}'_i(X;w)F\big|=\Big|\frac1{|W_4|}\sum_\omega\int F\,e^{\Phi_\omega+\Psi_\omega}\,d\mu_\omega\Big|\le\sup|F|\,\mathsf{M}(i,X)\le\sup|F|\,e^{-(2-c')p_0(g_i)}.}
\end{gather*}

Under a positive measure, the modulus of a complex factor serves as a bound for it; in the case of the small complex term produced by a complex pair, this is the device of Ba{\l}aban, Comm. Math. Phys. 122 (1989) 355, (1.73). The estimate just displayed is made on the moduli, and a sourced component retains the small factor of the unsourced one. That $\mathbf{T}'_i(X;w)F$ is holomorphic follows from dominated convergence together with the theorems of Morera and Fubini. For a healed component, the exponentiation of \cref{candidate:large_field} produces its local terms from this bound, each at most $e^{-p_0(g_i)}e^{-\kappa d_i(Y)}$ on a domain $Y$, just as at $w=0$.

\emph{The whole label.} Bounds on the individual components cannot be multiplied together, even though each is uniform in the variables outside its own component, because the background couples the components. (For $\mathbf{1}\{|\xi_1-\xi_2|<\epsilon\}$ on $[0,T]^2$, every integral in one variable is $\le2\epsilon$, while the double integral equals $2\epsilon T-\epsilon^2$.) By Tonelli, what does multiply is a product of functions of disjoint groups of variables; here the factors of this one product are the pieces of the pointwise majorant together with the integrands of the masses. For this reason the supremum is placed inside the integral. The result is the integrated bound of \cref{sec:large-field-rare-patches-gains-per-creation-and-costs-per-site}, with one factor for each creation and one mass for each outermost component on which $\mathbf{R}$ has acted. If every $e^{\Phi_\omega}$ in $\mathbf{T}'_k(\mathcal{S})$ is replaced by $e^{c\mathsf{v}(\omega)}|e^{\Phi_\omega}|$, then

\begin{gather*}
\FitMath{\int\mathbf{T}'_k(\mathcal{S})\Big[\chi_{k,\mathcal{S}}\,e^{-A(x_k(\cdot),U_k(\mathbf{V}))}\Big]dV_k\le\exp\Big[C\sum_{j\le k}(1+\log x_j)\,\#_j(Z_j)\Big]\prod_{(t,\mathfrak{c})}e^{-cp_0(g_t)}\prod_{(i,X)}\mathsf{M}(i,X),}
\end{gather*}

where the first product runs over those creation cells, outside these components, that belong to a single history with the label $\mathcal{S}$, namely the history for which this product is largest (the creation cells are not determined by the label, \cref{sec:large-field-rare-patches-gains-per-creation-and-costs-per-site}). The second runs only over the outermost ones among the components on which $\mathbf{R}$ has acted and which remain in the label. (In \cref{sec:large-field-rare-patches-gains-per-creation-and-costs-per-site} the bound carries $p_0(g_t)^2$ and has no logarithm. The sum over the labels needs only the present weaker form, because $CM^4\check{R}_t^5\log x_t\le\tfrac12c\,p_0(g_t)$ (the logarithmic cost allowed in \hyperref[hyp:C3]{the rarity hypothesis}): half of the exponent pays the cost and the other half is kept. Inside a drop-out the mass has consumed the factors of the creations, and whatever the label still shows there is paid for by the weight. The masses are required at a weight that exceeds the $e^{c\mathsf{v}(\omega)}$ spent by the sources by a constant; this constant is determined by the per-site cost of the part of the label visible inside the drop-out, and its choice comes after that of $\varrho$ and before the ceiling on $g$ is set (the mass bound is valid at every weight $c\ge0$, \cref{sec:large-field-rare-patches-gains-per-creation-and-costs-per-site}, and the masses $\mathsf{M}(i,X)$ in the display are understood at the raised weight). The surplus comes from that part of $\mathsf{M}(i,X)$ which is made up of the local histories admitted by the label; on these, $\mathsf{v}(\omega)$ is bounded below by the volume shown \cref{sec:large-field-rare-patches-gains-per-creation-and-costs-per-site}, \cref{sec:large-field-the-induction-closed-the-order-of-the-constants-the-sum}.) Once its vacuum constant is set aside, the sourced term of the label is dominated by the left-hand side multiplied by $e^{C(1+\varrho)}$ for each bond of the determining set, for three reasons. The moduli are taken inside the integrals. In the operations, $|e^{\Psi_\omega}|\le e^{c\mathsf{v}(\omega)}$. On the supports of the tests, what remains of the sourced exponent once the Wilson term $-A(x_k(\cdot),U_k)$ and the vacuum constant are set aside is at most $C(1+\varrho)$ for each bond of the determining set \cref{candidate:sources}. Consider the weighted sum over the labels of \cref{sec:large-field-the-induction-closed-the-order-of-the-constants-the-sum}, that is, the total masses summed over $\mathcal{S}$, where the credit pays the costs. In it every component $C$ of $Z_k$ keeps a factor

\begin{gather*}
\FitMath{\exp\big[-cp_0(g_k)-\kappa_1|C|_k\big],}
\end{gather*}

where $|C|_k$ is measured in cells. In \cref{sec:large-field-rare-patches-gains-per-creation-and-costs-per-site} the factor is taken per site, with $\#_k(C)\ge d_k(C)$ in place of $|C|_k$. This contains the form ``with the additional term $-\kappa_1d_k(X)$'', which Ba{\l}aban gives for a general region (ibid., p. 387).

\emph{The cost of the reserve.} A factor $e^{c\log x_j}$ for each site amounts to at most $e^{cM^4\check{R}_j^5}$ for each cell. Its form is that of the volume costs $CM^4\check{R}_n^5$ per cell and per step, which, until $\mathbf{R}$ acts on a component, are covered by the credit $\kappa_j(Z)$ of that component: an amount $p_0(g_j)^2$ for each cube, deposited at its creation, pays for $\check{R}_j^6$ for each cube (ibid., (1.79)--(1.81), in which both the deposit and the cost are measured by the linear size $d'_j(Z)+1$, counted in $MR_j$-cubes, and (1.89); these bound $\mathbf{T}'_k(X)1$ for the operation of (1.71), which contains the Wilson factor). Ba{\l}aban already pays $O(1)\log g_j^{-2}$ for each site of $Z_j$, on account of the normalisation constants (ibid., p. 380). The spare weight, whose constant $c$ is free, holds once $g$ is small enough, the constant $c$ being chosen first \cref{sec:large-field-rare-patches-gains-per-creation-and-costs-per-site}.

\begin{proposition}[{Sources along a frame}]\label{prop:source-domains-4-7} Let $(\theta_j)_{j\le n}$ be a frame such that $\theta_j\ge\gamma_J^{-2}$ throughout, and let $(x_j)$ be the run along it starting from a real $\beta$, as in the frame setting of \cref{candidate:flow}. The cube sizes are $\check R_j=R(\min_{i\le j}\theta_i)$. All the other auxiliary choices of level $j$ (the reference laws, the thresholds, the gauge-fixing weight and the scaling of the fluctuation field) are made with $\theta_j$, whereas the weight of the action obeys its own recursion. The reference line is a frame. A first-passage run, that is, the run starting from a first-passage datum $(K,\beta)$, is the run taken along its own couplings. The run is said to stay in the tube if $|\theta_j-x_j|\le\mathsf{t}_j$ for all $j$, and to have room $r$ if $|\theta_j-x_j|+r\le\mathsf{t}_j$. Suppose that the run has room $2r$ for $r=\varrho$ or for $r=2\varrho$. Then, for all sources with $\|w\|_\flat<r$ and for the steps of the frame, \cref{prop:source-domains-4-3,prop:source-domains-4-5,prop:source-domains-4-6,prop:source-domains-4-1} are valid, with upper stability and the floor holding with $C(1+\log\theta_n)$ for each site of $\mathbb{T}^{(n)}$. They hold on every admissible pair $(\ell,K)$ of a torus and a cutoff, admissibility meaning $K\ge n$ and $\ell L^{K-i}\ge CM\check{R}_i$ for $i\le n$, as in \cref{sec:flow-a-common-companion-and-the-torus-floor}; all of this is uniform in $\ell$, in $K$, in $n$, in $\beta$ and in the frame. Moreover

\begin{gather*}
\FitMath{\mathsf{z}_j(\zeta):=(\theta_j-x_j)+\zeta_j(\zeta)\in D(0,\mathsf{t}_j),} \\[3pt]
\FitMath{\sum_{j\le n}\,\sup_{0<|\zeta|<r}\frac{|b_j(\zeta)-b_j(0)|}{|\zeta|}\le\tfrac13.}\qedhere
\end{gather*}


\end{proposition}

The mechanism is as follows. A step depends separately on the choice $\theta_j$ and on the weight $x_j-\zeta_j$. For given remainders, its small-field part is holomorphic in their difference $\mathsf{z}_j(\zeta)$ on $D(0,\sigma_\star\theta_j)\supset D(0,\mathsf{t}_j)$ (the small-field holomorphy hypothesis in \hyperref[hyp:B4]{the ordinary step}, on the disc stated here); this difference is the deviation $\mathsf{z}_j=\theta_j-x_j$ displaced by the running shift. The full step is formed at a real deviation in the tube, together with a complex shift about it of modulus at most $\bar c$ (the complex class of \cref{sec:large-field-the-induction-closed-the-order-of-the-constants-the-sum}). Consequently, as far as the increments are concerned and with the remainders $R^{(i)}$ held fixed, a deviation from the frame and a constant source are the same variable. The remainders depend on the source as well, and the activities in the two presentations of \cref{prop:source-domains-4-5} need not coincide term by term. The room is what keeps the complex excursion $|\zeta_j|\le\frac43|\zeta|$ inside the disc. The objects constructed do not depend on the tube. A single statement therefore serves the first-passage runs, the comparison of two cutoffs in \cref{ch:comparison-capsule-v7,ch:weighted-matched-domains}, and the runs with frozen choices that occur in \cref{sec:flow-frozen-choices-slopes-and-tuned-bare-couplings} and \cref{thm:terminal-coupling-limit-v7}; it is used with $r=\varrho$, where the room is $2\varrho$, and with $r=2\varrho$, where the room is $4\varrho$ (which is the room $C\varrho$ of \cref{prop:source-domains-4-1}).

\emph{Reading at depth $s$.} For the run stopped at level $K-s$, everything above holds with the following changes. The norm $\|f\|_{\sharp,s}=\max\{\|f\|_\infty,CML^{-s}\operatorname{Lip}f\}$ replaces $\|f\|_\sharp$, the Lipschitz constant being measured at scale $L^{-s}$. The sums are restricted to $j\le K-s$. Either normaliser may be used, each with exponents of its own: $E$ of the whole run, or the vacuum constant $E_{(K-s)}$ of the first $K-s$ steps. The number of sites of $\mathbb{T}^{(K-s)}$ replaces $\ell^4$, and $L^{-(K-s-j)}$ replaces $L^{-(K-j)}$. The exponents of the two global bounds also depend on $s$ \cref{sec:short-distance-isolating-a-single-scale-coupling-and-normalisation-at-the-scale}.


\begin{proposition}[{Large fields are rare uniformly in the volume}]\label{prop:source-domains-4-13} (Live regions, creation sets, lineages, marks and ages, as well as the law $\tilde\nu$ of the histories, are defined in \cref{sec:large-field-the-tests-the-nested-regions-and-the-content-of-a}, and the prepaid bounds are stated in \cref{sec:large-field-the-induction-closed-the-order-of-the-constants-the-sum}. They have here the meaning given there.) For $g$ small, for every level $0\le k\le K$ (with $\mathcal Z_0=\emptyset$) and every level-$k$ cell $\mathfrak{c}$, one has uniformly

\begin{gather*}
\FitMath{\tilde\nu\big(\mathfrak{lv}_k(\mathfrak{c})\big)=\tilde\nu\big(\mathfrak{c}\subset\mathcal{Z}_k\big)\le Ce^{-cp_0(g_k)},} \\[3pt]
\FitMath{\tilde\nu\big(\mathfrak{c}\text{ is created at step }k\big)=\tilde\nu\big(\mathfrak{c}\subset\mathcal{W}^\wedge_k\big)\le Ce^{-cp_0(g_k)}.}\qedhere
\end{gather*}


\end{proposition}

The second bound ($k\ge1$) is a consequence of the first, because creation cells are live. Indeed, $\mathbf{R}$ touches a component only if no creation has occurred in it during the last $\mathsf{N}_k$ steps. Hence a component of $Z_k$ containing a creation cell of step $k$ is not healed at step $k$, which gives $\mathcal{W}^\wedge_k\subset\mathcal{Z}_k$. The same inclusion also bounds the event that $\mathfrak{c}$ lies in a live component in which a creation occurs at step $k$. The theorem relies on the following point. A label's determining set starts from a neighbourhood of its first region, $\Omega_0=\Omega_1^\sim$, and not from the whole torus. Its background therefore requires the bare field to be regular only on that neighbourhood (see under \cref{lem:source-domains-4-15}). If regularity is required at every bare plaquette of $T_\ell$, then a single rough plaquette, wherever it lies, leaves the step-$1$ background undefined. By totality \cref{sec:large-field-the-tests-the-nested-regions-and-the-content-of-a}, every test formed with that background then fails at every place, so that $Z_1$ equals the whole torus and is never healed. Moreover, at fixed $(\beta,g)$, the probability that such a plaquette is present tends to $1$ as $\ell\to\infty$. The theorem concerns only the law of the histories of the run; it gives no bound on the law that $\mu_{K,\ell,\beta}$ induces on the bare plaquettes inside a box of fixed physical size.

The proof is a first-moment argument. For a level-$n$ cell $\mathfrak{z}$, let $\mathcal{E}^{(n)}_{\mathfrak{z}}(k)$ be the survival event of \cref{sec:large-field-the-tests-the-nested-regions-and-the-content-of-a}, namely $\mathfrak{z}\subset\mathcal{Z}_i$ for $n\le i\le k$. The event $\mathcal{E}^{(k)}_{\mathfrak{c}}(K)\subset\mathfrak{lv}_k(\mathfrak{c})$ is the never-healed event. (i) If healing takes place at a level $j>k$, the relative
removal comparison of \cref{sec:aw-chessboard} gives
the factor $Ce^{-cp_0(g_k)-cp_0(g_j)}$, summable in $j$.
This use requires the outer-operation comparison on the
actual completed positive terms. (ii) We use $\operatorname{Vol}_i$, $\operatorname{Dep}_i$, the mark $\Theta^{(i)}$ with its dials $\mathsf{w}$, $c_a$, and the end mark $\Theta=\Theta^{(K)}$, all as in \cref{sec:large-field-the-tests-the-nested-regions-and-the-content-of-a}. For a component live at level $i$, along its lineage $\operatorname{Vol}_i$ charges every live cell, at every level $j\le i$, with twice the weight $\mathsf{r}_j=M^4\check{R}_j^5$, while $\operatorname{Dep}_i$ contributes $cp_0(g_j)^2(d'_j(W)+1)$ for every component $W$ that is created. By \cref{sec:large-field-the-induction-closed-the-order-of-the-constants-the-sum}, the upper bound of \cref{candidate:sources} stays valid once the mark is prepaid: for the run stopped at an arbitrary level when the dial is $\mathsf{w}=0$, and at the last level when the dial is $\mathsf{w}>0$. In the form needed here, the prepaid bound takes up the large-field bound one lineage at a time and consists of two statements, independent of each other. \emph{Volumes are paid by events.} Let $\mathsf r_j=M^4\check R_j^5$. For every history, and for every component $X$ of $Z_k$, whether live or healed at step $k$, one has

\begin{gather*}
\FitMath{\sum_{(j,C)}\mathsf r_j\,|C|_j\le c_G\sum_{(j,C)}\mathsf r_j\check R_j\,e(j,C),} \\[3pt]
\FitMath{e(j,C):=\big(|\operatorname{pred}|-1\big)^++\big|\operatorname{pred}\cap\operatorname{Act}_{j-1}\big|+\sum_W\big(d'_j(W)+1\big),}
\end{gather*}

where both sums run over the components of the lineage of $X$. Here $\operatorname{pred}$ denotes the set of components live at level $j-1$ that are contained in $C$; $\operatorname{Act}_{j-1}$ is the set of those components that are live even though $\mathbf R$ of level $j-1$ has acted on them; $W$ ranges over those components of the step-$j$ creation set that lie in $C$; and the constant $c_G$ is determined by $L$, $M$ and the remaining constants fixed ahead of the ceiling $\gamma$ imposed on the couplings, and by nothing else. The coupling $\gamma_\star$ is taken small. The inequality is a statement on the geometry of the regions alone; neither a field nor a measure enters. At levels where nothing happens, a component, measured in cells of its own level, shrinks geometrically until it becomes small; and once a component has stayed small over $\check R_j+1$ levels without any event, it belongs to the class of $\mathbf R$. \emph{Masses by lineage traces.} The lineage trace $\mathfrak l$ of $X$ is the list of the regions $\mathsf Z_j(X)$, of the sets created in them, and of the sets $\operatorname{Act}_j$. We set 

\begin{gather*}
\FitMath{\operatorname{wt}(\mathfrak l):=\exp\big[C_{\rm chg}\operatorname{Vol}_k-\operatorname{Dep}_k-\sum_{(j,W)}2p_0(g_j)-\sum_j|\operatorname{Act}_j|\,p_0(g_j)\big]}
\end{gather*}

, where $C_{\rm chg}$ is the charge constant; it is defined by the requirement that, for each cell of level $j$, the sums over the admissible regions, the Gaussian integrals, the normalisation constants and the vacuum-energy counterterms contribute a factor no larger than $e^{2C_{\rm chg}\mathsf r_j}$. This constant is determined by the parameters of the construction alone, and it is fixed before $\gamma_\star$. Let $\mathscr L$ be any set of lineage traces of level $k$, and sum over those histories of level $k$ whose live components carry the lineage traces $\mathscr L$. Then

\begin{gather*}
\FitMath{e^{-E_{(k)}}\sum_h\mathsf t(h)\le e^{C|\mathbb T^{(k)}|}\prod_{\mathfrak l\in\mathscr L}\operatorname{wt}(\mathfrak l):}
\end{gather*}

indeed, one unfolds the majorants of \cref{sec:large-field-rare-patches-gains-per-creation-and-costs-per-site}, passing through the components that $\mathbf R$ has acted on, as far as the creations, and one sums the entropy with the lineage trace held fixed. No factor arises from a joining, since a creation at a leaf of the lineage pays for it (\cref{prop:large-field-2-6}, move (3)). The relation $\mathsf r_j\check R_j=M^4\check R_j^6\ll p_0(g_j)$ is valid for all couplings not exceeding the anchor coupling $\gamma_\star$. This anchor coupling is fixed small only once $L$, $M$, the dials $\mathsf w$ and $c_a$, and the rate $c'$ introduced below (a constant of this proof only, different from the small $c'$ of the mass bound) have been fixed, and it does not exceed $\gamma_J$; $g$ is subsequently taken small depending on $\gamma_\star$. Consequently the volumes, the entropy of the lineage traces and the mark are all paid for by the events. Precisely, take the rate $c'\ge0$, which is fixed before $\gamma_\star$ and may be given any value; the value $c'=1$ already outweighs the number $C^{|X|_k}$ of connected sets $X$ of cells containing a given cell, because $\mathsf r_k\ge M^4$. Then for every set of lineage traces ending in $X$, 

\begin{gather*}
\FitMath{\sum_{\mathfrak l}\operatorname{wt}(\mathfrak l)\,e^{\mathsf w\operatorname{Dep}_k+c_a\operatorname{Vol}_k}\le e^{-c'\mathsf r_k|X|_k}\,e^{-2\min p_0(g_{j_0})}}
\end{gather*}

, where $j_0$ denotes the first level of a lineage. It is this bound that prepays the mark. There are two reasons for carrying out the argument at the last level, with the end mark. The first reason applies already at $\mathsf{w}=0$. The gain obtained from the first moment in (iii) is the quotient of two quantities: the exponent in the moment bound, equal to $\mathsf{B}(y)=C(1+\log y)$ for every site of the lattice on which the run stops, and the product of $c_a\mathsf{r}_k$ with the count of level-$k$ cells. When the run is stopped at level $k$, the quotient equals $\mathsf{B}(x_k)(M\check{R}_k)^4/(c_a\mathsf{r}_k)\approx\log x_k/(c_a\check{R}_k)$, and this carries no smallness of order $e^{-cp_0(g_k)}$. At the last level there are fewer sites, by the factor $L^{-4(K-k)}$. The second reason concerns the alternative of (iv) in which the torus shrinks. There a single deposit is weighed against the volume, and this requires $\mathsf{w}>0$. Whatever value the dial takes, the end mark registers only lineages of components still live at the last level, that is, only the never-healed lineages. This is why the healed part is split off in (i). Together with the floor one obtains $\mathbb{E}_{\tilde\nu}e^\Theta\le e^{C_g\ell^4}$, hence $\mathbb{E}_{\tilde\nu}\Theta\le C_g\ell^4$ by Jensen's inequality. On the other hand $\operatorname{Vol}_K$ charges every never-healed cell, so $\sum_{\mathfrak{c}}\mathbf{1}_{\mathcal{E}^{(k)}_{\mathfrak{c}}(K)}\le\Theta/(c_a\mathsf{r}_k)$, where $c_a>0$. (iii) Translations by level-$k$ cells commute with the steps up to level $k$; therefore $\tilde\nu(\mathfrak{lv}_k(\mathfrak{c}))$ is the same for every $\mathfrak{c}$. The event $\mathcal{E}^{(k)}_{\mathfrak{c}}(K)$, by contrast, depends on coarser partitions and is not covariant. Averaging gives

\begin{gather*}
\FitMath{\tilde\nu\big(\mathfrak{lv}_k(\mathfrak{c})\big)\le\max_{\mathfrak{c}'}\tilde\nu\big(\mathfrak{lv}_k(\mathfrak{c}')\setminus\mathcal{E}^{(k)}_{\mathfrak{c}'}(K)\big)+\frac{C_g\ell^4}{c_a\mathsf{r}_k\,\#\{\mathfrak{c}'\}}\le Ce^{-cp_0(g_k)}+C_gL^{-4(K-k)}.}
\end{gather*}

(iv) To make $K-k$ large, the run is extended on the same bare lattice up to a larger reference coupling, not exceeding the anchor coupling $\gamma_\star$. This is a fixed upper limit for the reference couplings. It is chosen small once the dials $\mathsf{w}$ and $c_a$ are fixed, in such a way that (i), (ii) and the comparison made in the first alternative below are valid at each reference coupling up to $\gamma_\star$; only afterwards is $g$ chosen small, depending on $\gamma_\star$. This continuation leaves the law up to level $k$ unchanged \cref{sec:flow-units-common-prefixes-and-change-of-reference}. Steps (i)--(iii) are applied to the continued run, with $\Theta$ the mark at its last level. 

%
%
Let \(K^\dagger\) be the largest admissible continued
first-passage level, \(g^\dagger\) its reference coupling,
and \(m^\dagger=m-(K^\dagger-K)\) its torus exponent.
Write \(K_\star\) for first passage at the anchor
coupling and \(X_\star=\gamma_\star^{-2}\).
Choose \(c_{\rm gap}=\lceil c_5/\lambda_\sharp\rceil+1\).
If \(K^\dagger+c_{\rm gap}\le K_\star\), the lower flow
bound produces a strict new record low within those
\(c_{\rm gap}\) steps. A strict record low is a first-passage
level for a reference threshold there. For a record \(q\),
use \(g'=(x_q-B_\sharp)^{-1/2}\), or the anchor if
\(x_q-B_\sharp<X_\star\), in which case \(q=K_\star\).
The new reference coupling is at least \(g^\dagger\)
and at most the anchor.
The admissible side floor \(m_\star(g')\) is nonincreasing
with \(g'\). Thus the record would still be admissible
if \(m^\dagger\ge m_\star(g^\dagger)+c_{\rm gap}\).
Maximality gives the opposite inequality. The terminal
site count is then at most \(C(1+\log (g^\dagger)^{-2})^{4r}\)
by the side floor. Every surviving lineage carries a
birth, with mark at least
\(\mathsf w c_{\mathsf D}(p_0(g_k)-1)^2\) by the drift.
For the chosen anchor ceiling this mark pays the terminal
moment exponent and leaves \(c p_0(g_k)\), since
\(2\mathfrak p_0>4r+1\).

Otherwise \(K^\dagger>K_\star-c_{\rm gap}\).
The upper flow bound and the first-passage windows give
\[
 p^\dagger:=K^\dagger-K
 > \frac{g^{-2}-X_\star-B_\sharp}{B_\sharp}-c_{\rm gap}.
\]
The continued exponent \(m^\dagger=m-p^\dagger\)
is nonnegative. In particular every side with
\(m\le m_V(g)=\lfloor(g^{-2}-X_\star-B_\sharp)/B_\sharp
\rfloor-c_{\rm gap}\) falls in the volume case.
This gives the horizon \(\ell_V(g)=2L^{m_V(g)}\).
Use the averaged never-healed bound at this terminal level.
Writing \(u=K-k\), division by the number of level-\(k\)
cells gives
\[
 \frac{\mathsf B(x_{K^\dagger})}{2c_a\check R_k}
 L^{-4(p^\dagger+u)} .
\]
Here \(x_{K^\dagger}\le g^{-2}+B_\sharp\) while
\(x_k\le (g^{-2}+B_\sharp)(1+u)\). Consequently the
required extra depth is at most
\(C(1+\log g^{-2})^{\mathfrak p_0}\), uniformly in \(u\):
the \(u\)-dependent part of \(c p_0(g_k)\) is bounded
by \(4u\log L+C\) because
\((\log(1+u))^{\mathfrak p_0}=o(u)\).
The available depth \(p^\dagger\) exceeds this bound for
small \(g\). The average is thus at most \(Ce^{-cp_0(g_k)}\).
Combine it with the pointwise healed estimate by averaging
the live-prefix events, whose probabilities are translation
invariant at level \(k\). This uses the common prefix of
the two runs.

\begin{proposition}[{Ages}]\label{prop:source-domains-4-14} Denote by $C$ the component live at level $k$ that contains $\mathfrak{c}$, by $\mathsf{Z}_j(C)$ its lineage regions, and by $j_0$ the least level at which these regions are non-empty. The three ages of $C$ are the age read in the label, the age of the lineage, $\operatorname{age}_k(C;h)=k-j_0$, and the seed age. These ages and the events $\mathfrak{lv}^{=a}_k(\mathfrak{c})$, $\mathfrak{lv}^{\ge a}_k(\mathfrak{c})$, $\mathfrak{sd}^{=a}_k(\mathfrak{c})$ are defined as in \cref{sec:large-field-the-tests-the-nested-regions-and-the-content-of-a}. The age read in the label refers to the oldest region of the label that appears in $C$: $\operatorname{age}^{\mathrm{lab}}(C)=k-\min\{t:C\cap Z_t\ne\emptyset\}$, where $Z_t=\Lambda_t^c$ is the large-field region of scale $t$ in the label $\mathcal{S}(h)$ of the history. The proposition states that the three ages are ordered,

\begin{gather*}
\FitMath{\operatorname{age}^{\mathrm{lab}}(C)\,\le\,\operatorname{age}_k(C;h)\,\le\,\text{seed age},}
\end{gather*}

and it gives a decay in the second of them. For the first inequality, note that every point of $C\cap Z_t$ belongs to the live region at each level between $t$ and $k$. It therefore belongs to $\mathsf{Z}_t(C)$, and so $t\ge j_0$. For the second inequality, no component of $\mathsf{Z}_{j_0}(C)$ contains a component that is live at the preceding level. Each such component is therefore created at step $j_0$, and it is contained in $C$. The second inequality may be strict. Suppose a region created at a step $t<j_0$ is contained in $C$. Then it was healed at some step before $k$, and the live region of some later level covers its place anew. As a point set it counts for the seed level. It contributes no factor, however, since no cell of it stays live through all the levels from $t$ to $k$. In terms of events, $\mathfrak{sd}^{=a}_k(\mathfrak{c})\subset\bigcup_{a'\le a}\mathfrak{lv}^{=a'}_k(\mathfrak{c})$. Let $b\ge1$, let $c_a\ge C(b)\log L$ (a floor on $c_a$, which is chosen after $L$), and let $g$ be small depending on $b$ and $c_a$. Then, uniformly, with constants $C$ on the right-hand sides that may depend on $b$ and $c_a$,

\begin{gather*}
\FitMath{\tilde\nu\big(\mathfrak{lv}^{=a}_k(\mathfrak{c})\big)\le Ce^{-cp_0(g_k)}(2L^b)^{-a},} \\[3pt]
\FitMath{\tilde\nu\big(\mathfrak{lv}^{=a}_k(\mathfrak{c})\cap\mathcal{E}^{(k)}_{\mathfrak{c}}(K)\big)\le C(2L^b)^{-a-(K-k)}.}
\end{gather*}

Both bounds decay in the age of the lineage, and so does the age series given further below. Suppose the seed level is formed only from regions that are healed at no step before $k$. Then the two bounds and the age series hold in the seed age as well, that is, for $\mathfrak{sd}^{=a}_k(\mathfrak{c})$: the seed level equals $j_0$, the two ages coincide, and $\mathfrak{sd}^{=a}_k(\mathfrak{c})=\mathfrak{lv}^{=a}_k(\mathfrak{c})$. When the seed level is formed from all regions, the price of survival yields no decay. Two facts show this. First, a level-$k$ cell contains of the order of $L^{4(k-t)}$ cells of level $t$, whereas the bound per creation has an exponent that is merely polylogarithmic in $k-t$. Second, take $k=K$ and $a=K-1$. Applied there, the bound gives at most $C(2L^b)^{-(K-1)}$, at fixed $g$ and for every $K$, for the probability that a fixed cell at the final level is live and occupies the place of a region created at step $1$. Yet the tests of the last step depend only on $x_{K-1},x_K\approx g^{-2}$. A first-moment argument yields only a density, and never a factor per level for a single cell. For this reason the chessboard estimate is now applied to the large fields, which is a second use of reflection positivity.

\phantomsection\label{notdef:age-survival}
\end{proposition}

\begin{lemma}[{Chessboard for the law of the histories}]\label{lem:source-domains-4-15} Put $q_a=e^{-c_a}$, where $c_a\ge c^{-1}\log(2L^4)$, equivalently $L^4q_a^c\le\frac12$. In this condition, in the two bounds below and wherever these bounds are quoted, the constant is $c=\frac1{16}$; it is fixed before $c_a$, and it arises as the root, taken over the tiles, of the bound for a pattern that is repeated in one tile out of $2^4$. Fix $\epsilon>0$ and take $g$ small depending on $c_a,\epsilon$. A \emph{tiling} is a translate, by a multiple of the side length of a level-$k$ cell, of the partition of the torus into cubes of side $\mathsf{l}_\square$, called \emph{tiles}, whose faces lie in mirror planes that pass through midpoints of bonds. In each tile, the cube with the same centre and with side $\frac34\mathsf{l}_\square$ is called its \emph{core}. A cell, of whatever level, is called \emph{deep} in a tile when the level-$k$ cell containing it lies in the core of that tile. Then for each level $k\le K$ there is a tile side $\mathsf{l}_\square\approx M\check{R}_k$, measured in level-$k$ units, which is a power of $L$ and spans at least sixteen level-$k$ cells, with the following property. Let $n\le k$, let $\mathfrak{z}$ be a level-$n$ cell situated deep inside a tile $\square_0$ of a tiling, and let $\square_1,\dots,\square_r$ be further tiles, all distinct. Then

\begin{gather*}
\FitMath{\tilde\nu\big(\mathcal{E}^{(n)}_{\mathfrak{z}}(k)\big)\le\mathsf{l}_\square^{-4}\,q_a^{c(k-n+1)},} \\[3pt]
\FitMath{\sum_{\mathfrak{c}_1,\dots,\mathfrak{c}_r}\tilde\nu\Big(\mathcal{E}^{(n)}_{\mathfrak{z}}(k)\cap\bigcap_{i\le r}\mathfrak{lv}_k(\mathfrak{c}_i)\Big)\le\mathsf{l}_\square^{-4}\,q_a^{c(k-n+1)}\,\epsilon^r,}
\end{gather*}

where the sum runs over the level-$k$ cells $\mathfrak{c}_i$ lying deep inside $\square_i$. The first bound is valid for every level-$n$ cell on every torus, including tori too small to carry a tiling with tiles of this side. The bounds are independent of how many tiles the torus contains. Survival costs a factor $q_a^c$ per level, and every further tile that contains a deep live cell costs a factor $\epsilon$. 
\end{lemma}

\subsubsection{Radial windows and bare-field witnesses}
\label{ch:age-window-witness}

This section supplies the conditional-density and witness steps in
\cref{lem:source-domains-4-15}. It concerns the actual positive law of a
finite history prefix, at zero source. The all-history hypotheses
\(\mathcal H_{\mathrm{step}}\) of
\cref{rg:primitive-closure-premises} remain in force. In particular, the
background and Gaussian estimates below are required on the actual
determining regions and stored-field fibres. A regular homogeneous
construction does not supply these hypotheses on its own.

The collar cutoff used here is the component-local cutoff specified in
\cref{sec:large-field-stored-variables-and-integration-in-a-large-field-patch}.
Its exponentially close global variant can support analytic estimates,
but it cannot replace this sharp cutoff in a comparison of two runs.
With the local cutoff, the comparison needs windows only at the ordinary
split of each step. No window at a preparatory completion or collar test
will be removed.

\paragraph{The finite prefix and its real estimates}
\label{sec:aw-prefix}

Fix \(k\), and measure lengths in level-\(k\) units. Write
\[
 s_j=M\check R_j L^{-(k-j)},\qquad
 r_j=M^4\check R_j^5,\qquad N_j=\check R_j,\qquad
 q_a=e^{-c_a},\qquad \widehat q=q_a^{1/16}.
\]
The cells are nested, \(s_j/s_{j-1}\in\{1,L\}\), and the flow satisfies
\[
 c(j-i)-C\le x_i-x_j\le C(j-i)+C,\qquad
 \check R_i=R\!\left(\min_{t\le i}x_t\right),\qquad
 (\log x)^r\le R(x)\le L\max\{1,(\log x)^r\}.
\]
All constants below may depend on the fixed construction parameters,
including \(L\), but not on the depth, the torus or the chosen history.

We record the real estimates that enter the argument. These are
formation and mass estimates, rather than the class-invariance target.
The background response on a regular collar has the form
\begin{equation}
 \left|V^{\rm bg}(b)V^{{\rm bg}\prime}(b)^{-1}-1\right|
 \le C_{\rm dec}e^{-\delta_{\rm dec}d(b,A^c)}
       \varepsilon_{\rm rim}.
 \label{aw:collar-response}
\end{equation}
The two determining sets and their data must coincide in \(A\), its rim
must lie in the common top domain, and the actual data must lie below
the regularity ceiling. Near the top domain,
\(\varepsilon_{\rm rim}\le C\varepsilon_j\). A factor
\(C(\log x_j)^r\) on rebuilt distant layers is allowed; two additional
cells of collar pay for it. The first determining region is
\(\Omega_0=\Omega_1^\sim\), rather than the whole torus.

In the literal nonpivot coordinates at a split, the smooth real
exponent satisfies
\begin{equation}
 \left|\partial_{a_b}\mathcal F\right|
 \le C_{\rm sl}p_0(g_j).
 \label{aw:stage-slope}
\end{equation}
This estimate includes the Haar, coarea and nonlinear terms. Its
quadratic contribution follows from an exponentially summable precision
row applied to fields bounded by \(Cp_0\).
%
For the ordinary remainder, the domains through the
tested bond have a bounded rooted sum with the ordinary
decay. A complex margin \(c/p_0(g_j)\) in the scaled
coordinate thus suffices by Cauchy's inequality. For
the preparatory remainder, use the coordinate
\(a_b=(\mathrm i g_i)^{-1}\log(V_i(b)C_b^{-1})\)
at its actual collar level \(i\), with \(C_b\) its
test centre. The stage domains meeting the
tested bond also meet the preparatory collar. Their
relative length does not charge the class box of side
\(O(n_{\rm sm}\check R_j)\) cells. Rooting at one of
its cells and summing the exterior continuations
therefore costs \(C(n_{\rm sm}\check R_j)^4\).
The complex margin in this collar coordinate is
\(c\delta'_i/g_i\), bounded below by a fixed constant.
Cauchy's inequality with the formed activity bounds
on that disc gives a contribution bounded by
\(C(n_{\rm sm}\check R_j)^4g_i/\delta'_i\).
The waiting-window bounds compare the scales, and
\(\mathfrak p_0>4r\) makes this at most \(Cp_0(g_j)\).
This uses the admitted collar margin with the relative
rooted sum, and the ordinary margin with the ordinary
rooted sum.

For an isolated box of free coordinates, write its conditional exponent
as
\begin{equation}
 H_e(y)=\tfrac12\langle y,K_{II}y\rangle+
        \langle y,K_{IO}A_O\rangle-P_e(y).
 \label{aw:isolated-exponent}
\end{equation}
Here \(I\) is the set of free bonds touching \(e^{+1}\), whereas the
marked bonds touch \(e\). The real Gaussian estimates and nonlinear
formation bounds give
\begin{gather}
 K\ge a_-I,\qquad |K(b,b')|\le a_+e^{-c|b-b'|},
 \qquad |A_O|\le Cp_0(g_{j-1}),                                      \label{aw:gaussian-input}\\
 \sup_I|\nabla P_e|\le \epsilon_{\rm nl},\qquad
 \|\nabla^2P_e\|_{\ell^2\to\ell^2}\le a_-/2,\qquad
 \epsilon_{\rm nl}\longrightarrow0.                                \label{aw:nonlinear-input}
\end{gather}
Distances in these formulas are in level-\((j-1)\) lattice units.
The bounds are needed after conditioning on the exterior datum.
Positivity of Haar measure alone supplies none of them.

Finally, let \(t(h)\ge0\) be the masses of the prefix and \(Z=\sum_ht(h)\).
The marked upper bound and the all-small lower bound, with the same
vacuum constant \(E_k\), are
\begin{equation}
 e^{-E_k}\sum_ht(h)e^{\Theta_k(h)}\le e^{C_+V_k},
 \qquad
 e^{-E_k}Z\ge e^{-C_-(x_k)V_k},\qquad
 C_++C_-(x)\le B(x):=C_B(1+\log x).
 \label{aw:prefix-masses}
\end{equation}
Here \(V_k=|\mathbb T^{(k)}|\), and
\[
 \Theta_k=c_a\sum_{X\ {\rm live\ at}\ k}\operatorname{Vol}(X),\qquad
 \operatorname{Vol}(X)=2\sum_{j\le k}r_j|\mathsf Z_j(X)|_j.
\]
The charge dial in the mass estimate has been increased before choosing
the coupling ceiling. Dividing the two inequalities in
\eqref{aw:prefix-masses} cancels \(E_k\) and gives
\begin{equation}
 \mathbb E_{\widetilde\nu}e^{\Theta_k}\le e^{B(x_k)V_k}.
 \label{aw:prefix-moment}
\end{equation}
In particular, if \(n_j\) distinct level-\(j\) cells are forced into
surviving lineages, their probability is at most
\(e^{B(x_k)V_k}q_a^{2\sum_jr_jn_j}\).

\paragraph{The coordinates read after a split}
\label{sec:aw-reads}

The independent split coordinates are the actual nonpivot link logs
\[
 a_b=g_i^{-1}i^{-1}\log\big(V_i(b)V_i^{\rm bg}(b)^{-1}\big),
 \qquad i=j-1.
\]
They are not coordinates obtained by a global nonlinear Landau
projection. The next field \(V_j\) is held fixed, and the literal
averaging constraint determines the pivots.

\begin{lemma}[{The split coordinate and the later decisions}]
\label{lem:aw-split-reads}
Fix all other stage coordinates. On either one-sided shell about any
of the three split thresholds, the conditional likelihood of a joint
survival event and of later ordinary-split no-window conditions is
constant along each radial ray. Existing sharp factors are constant
there as well. This statement uses the actual outward determining rows
and the margins of the ordinary step.
\end{lemma}

\begin{proof}
We describe the readsets, including the pivots. In a single averaging
step the pivot for a coarse bond is the middle fine bond of its central
line. Internal block contours contain no pivot. A pivot therefore occurs
in its own coarse constraint only. The local implicit-function theorem
solves that constraint, with the other nonpivot variables held fixed.
Changing \(a_b\) changes only \(b\) and the face pivots in its one-step
contour carrier. All constants in this finite-dimensional calculation
are \(C_L\); no lower bound uniform in \(L\) is being asserted.

The ordinary split reads cubes in \(\Omega_j^{\sim-1}\). Their full
one-step carriers stay inside \(\Omega_j\), with a cell margin.
Consequently none of the changed bonds belongs to the actual row
\(\Gamma_i\), whose bonds have no endpoint in \(\Omega_j\). The
background datum, the coarser field, and the older retained data are
unchanged. The same is true of data used by joined strip backgrounds:
their additional tower data are coarse lifts, and their inherited
fine rows are these same outward rows. Thus other stored nonpivot
coordinates acquire no hidden change through a background centre.

The plaquette tests depend on \(V_j\). The approximate fluctuation tests
exclude pivots, use the threshold \(2\delta_i\), and have exponentially
small centre offsets. They remain constant throughout the shells
about \(\delta_i^\flat,\theta_c\delta_i,\delta_i\).
The soft contour cutoffs are identically one there, since their
contours remain within \(C_L\delta_i\) and the parameter order gives
\(C_L\delta_i\ll\varepsilon_i\). For the density tests one needs the
quantitative new-test implication in \cref{lem:large-field-2-2},
not merely their previously recorded answer. With the new coarse
plaquette and approximate fluctuation tests passed, that implication
gives, on each old density cube in the one-cell enlargement,
\[
 |\partial U_{i,\square}(V_i)-1|
 \le\left(C_0L^{-2}+C_LA_1/A_0+
       Ce^{-c\delta_0M_2\check R_j}\right)\varepsilon_i\eta_i^2
 \le\tfrac14\varepsilon_i\eta_i^2.
 \]
The construction fixes these quantitative margins in this order:
the coarse-curvature floor on \(L\), the sufficiently small ratio
\(A_1/A_0\) after the finite-\(L\) response and pivot constants, and
the coupling ceiling after the collar constants. The shell enlarges
the fluctuation threshold by \(g_iW_\sigma=o(\delta_i)\), so the
same bound remains below \(\varepsilon_i\eta_i^2/2\).
Here \(\eta_i\) is the bare spacing in level-\(i\) units.
At \(i=0\), \(U_{0,\square}(V_0)=V_0\) on the tested plaquettes,
so these are the bare \(P_0\) tests on the full changed carrier.
Thus fixing the earlier \(P_0\) entry imposes no additional shell cut.
Old tests near a changed pivot are either removed in the strip
operation or identically one by its one-cell margin.

Now compare two states on the same ray, on the same side of a
threshold. All three interval flags agree. The ordinary recursion uses
only these flags, so its geometry, its creation entries and its
integration region agree. If the coordinate is integrated, its slot
is deleted and no subsequent move reads it.

If it is retained, it is still excluded from all later determining
data: the old domains only expand. Preparatory gauge reads lie in the
outer collars and use these datum rows and fixed coarser fields.
The other preparatory tests use backgrounds localised to the actual
component. A completion test can read the stored coordinate itself,
but then it tests the same previously stored modulus at the old
threshold. At a passed original cube test that answer is already one.
A failed original cube test is a creation; while its level remains
in the preparatory window, the class condition prohibits acting on its
component. Once it leaves that window, its old coordinate is below the
levels of those completion tests. No recomputation with a new centre
is permitted.

The remaining moves are fixed-law gauge or Haar draws, a frame draw,
coarse tree maps, or the coarse reference draw. The reference draw and
its normalizer read the exterior coarse datum and the fixed component
frame; they read no omitted fine row. Their read coordinates have
already been shown to agree. Its final density test therefore has
the same answer when the later coins agree. The old retained values
can remain as passengers; the assertion is about future decisions,
not about a full state kernel that carries such passengers unchanged.

Induction through the finite list of moves proves that the indicated
future events have equal indicators for equal later coins. Integration
over those independent coins gives the asserted constant likelihood.
This argument would fail if a retained fine staple were admitted into
a later datum row. Fixed coarse means by themselves would not exclude
that possibility.
\end{proof}

\paragraph{A deterministic family of split tests}
\label{sec:aw-diagnostics}

At step \(j\), after \(\Omega_j\) and its single global fluctuation field
have been formed, record the three interval flags for every candidate
nonpivot bond read by a cube in \(\Omega_j^{\sim-1}\). Then run the
original recursion using these flags. It uses exactly the flags it
would have queried before. Forgetting the unused flags recovers the
original history at every configuration, so this refinement changes
neither its law nor the moment \eqref{aw:prefix-moment}. It introduces
no new coin and no sum over possible recursion branches.

This precomputation is legitimate because the split uses one global
field, not a field recentered at a hypothetical explored component.
No such refinement is made for the context-dependent preparatory
collar cutoffs. Order the diagnostic split sites by step and then by
a fixed bond order. Each potential lattice bond now has one position
in this list; an inactive site has \(T_\sigma=0\).

Write \(p_i=p_0(g_i)\), and set
\[
 W_\sigma=(C_{\rm sl}p_i)^{-1},\qquad
 \mathfrak t_\sigma=
 \{\mathfrak b_i^\flat,\theta_c\mathfrak b_i,\mathfrak b_i\},
 \qquad \mathfrak b_i=\delta_i/g_i.
\]
Choose the ceiling so that every distinct threshold is at least
\(2W_\sigma\) from the other thresholds and is itself at least
\(2W_\sigma\). Let \(W_\sigma(u)\) denote the event that the site is
active and its modulus is within \(u\) of a threshold.

\begin{lemma}[{The conditional radial estimate}]
\label{lem:aw-radial}
Let \(E\) be a joint survival event, intersected with any set of
no-window conditions at diagnostic sites later in this deterministic
order. Then
\begin{equation}
 \mathbb P\big(E\cap W_\sigma(u)\big)
 \le C\,\frac{u}{W_\sigma}\,\mathbb P(E),
 \qquad 0\le u\le W_\sigma.
 \label{aw:relative-window}
\end{equation}
\end{lemma}

\begin{proof}
Disintegrate the stage law in the independent coordinates
\((s,a_\sigma)\). It has density
\(e^{\mathcal F_s(a)}\chi_s(a)\Pi(ds)\,da\), including all its
Haar and coarea factors. By \cref{lem:aw-split-reads}, the product of
\(\chi_s\) and the conditional future likelihood is constant on every
one-sided shell. Conditions at other diagnostic sites of the same
step depend on their fixed nonpivot coordinates, so they have the
same property. The stage slope is bounded by
\[
 C_{\rm sl}p_0(g_j)\le C_{\rm sl}(p_i+1)\le2/W_\sigma.
 \]
Along a ray in a shell of length \(W_\sigma\), the exponential density
changes by at most \(e^2\), and the radial Jacobian \(r^2\) by at most
four. Its initial portion of length \(u\) therefore carries at most
\(4e^2u/W_\sigma\) times the full shell mass. Apply this separately
above and below each threshold. The shells are disjoint. Integration
over the angle and \(s\) proves \eqref{aw:relative-window}.
\end{proof}

For a tile of side \(\ell_\square\), put
\[
 \begin{aligned}
 \bar p_j&=\max_{\max\{0,j-N_j-1\}\le i\le j}p_0(g_i),\\
 \delta_{\rm dec}D_j&=
 C+4\log\ell_\square+(4\log L+\log2)(k-j)\\
 &\quad+4(N_j+1)\log L+2\log\bar p_j,\\
 u_\sigma&=2C_{\rm dec}e^{-\delta_{\rm dec}\lceil D_j\rceil}
                  \varepsilon_j/g_{j-1}.
 \end{aligned}
 \]
The diagnostic sites of a tile satisfy the more economical count
\[
 |\mathfrak S_j^{\rm split}[\square]|
 \le C\ell_\square^4L^{4(k-j+1)}
 \le C\ell_\square^4L^{4(k-j+N_j+1)}.
 \]
Thus the original, stronger schedule still pays them. The flow over
\(N_j+1\) levels gives \(\varepsilon_j/g_i\le C\bar p_j\), and hence
\begin{equation}
 \sum_{\sigma\in\mathfrak S_j^{\rm split}[\square]}
       C u_\sigma/W_\sigma
 \le Ce^{-C}2^{-(k-j)}.
 \label{aw:window-budget}
\end{equation}
This displays the cancellation of all four count factors, including
the retained allowance \(L^{4(N_j+1)}\).
Before making any removal, enlarge the additive constant \(C\) so that
the same calculation gives
\(u_\sigma/W_\sigma\le C'e^{-C}\le
\min\{1,(2C_{\rm shell})^{-1}\}\) at every carried site, where
\(C_{\rm shell}\ge1\) is the constant in
\eqref{aw:relative-window}. Thus both \(u_\sigma\le W_\sigma\)
and \(\eta_\sigma=C_{\rm shell}u_\sigma/W_\sigma\le1/2\)
hold uniformly, including at large levels.

Let \({\rm NoWin}^{\rm split}(\square)\) exclude these diagnostic
windows. Remove them in reverse deterministic order. Applying
\eqref{aw:relative-window} at each removal and using
\(1-\eta\ge e^{-2\eta}\), with \(\eta\le1/2\), gives
\begin{equation}
 \mathbb P\!\left(E\cap\bigcap_{\ell=1}^r
          {\rm NoWin}^{\rm split}(\square_\ell)\right)
 \ge C_{\rm win}^{-r}\mathbb P(E)
 \label{aw:many-tile-windows}
\end{equation}
for distinct tiles. The constant in \(D_j\) can make \(C_{\rm win}>1\)
as close to one as desired. This multiplicative estimate uses the
conditional denominator in \eqref{aw:relative-window}; an
unconditional estimate on the probability of each window would not
give \eqref{aw:many-tile-windows}.

\paragraph{The bare-field witnesses}
\label{sec:aw-witness}

Realise every normalized measurable draw by an independent uniform
coin, in the fixed relative chart of its site or component. There are
finitely many potential draw sites in a finite prefix. Unused coins
are ignored. The resulting run is Borel, has the prescribed kernel
law, and, in the tiling regime, commutes with translations by
\(s_k\mathbb Z^4\). Here is the convention needed for the last
assertion. The torus side is at least \(2\ell_\square\), and
\(\ell_\square>\max\{128,2n_{\rm sm}\}s_k\). Every component on which
\(\mathbf R\) draws fits in a nonwrapping aligned box of side at most
\(n_{\rm sm}s_j\), less than half the torus side. Its bounding box
has a unique centre modulo the torus and has no nontrivial translation
stabilizer. Number its bonds by their offsets from that centre, in
the fixed frame. For each such shape and frame choose one Borel
coordinate isomorphism, and use the same isomorphism for every
translate. Translated states then have identical coordinate laws;
their quantile maps with the translated centre coin agree exactly.
Site and bond draws use the corresponding fixed local charts.

Induction over the moves proves covariance with this convention;
it does not follow from Borel quantile existence alone. On a smaller
torus an arbitrary Borel realization suffices for the direct moment
argument below, and no covariant realization for a component with a
nontrivial stabilizer is asserted. Integration is the deletion of a
slot; no later move may read a deleted slot.

Let \(m=dU\otimes d\xi\) be product Haar times Lebesgue, and
\(\mathbb P=\mu_{\rm Wilson}\otimes d\xi\).
Since the Wilson action is nonnegative, \(\mathbb P\le Z^{-1}m\).
The tile coordinates contain only bare bonds whose two endpoints are
strictly inside the tile and coins placed strictly inside it. For a
Borel event \(A\), define
\[
 F_\square(A)=
 \left\{x_\square:
 m^\square\{y:(x_\square,y)\notin A\}=0\right\},\qquad
 \phi_{\square,A}(U)=
 \int{\bf1}_{F_\square(A)}(U|_\square,\xi_\square)\,d\xi_\square.
 \]
Tonelli's theorem proves that these are measurable. The cylinder
over \(F_\square(A)\) is contained in \(A\) modulo a \(\mathbb P\)-null
set, and \(0\le\phi_{\square,A}\le1\).
For distinct tiles,
\begin{equation}
 \mathbb P\!\left(\bigcap_\ell
       {\rm pr}_{\square_\ell}^{-1}F_{\square_\ell}(A_\ell)\right)
 =\left\langle\prod_\ell\phi_{\square_\ell,A_\ell}\right\rangle_{\rm Wilson}.
 \label{aw:witness-product}
\end{equation}
The exterior measure in the definition is the product reference
measure, not a conditional Wilson measure.

These functions need not be gauge invariant. If the chosen statement
of chessboard positivity is restricted to gauge-invariant functions,
average each \(\phi\) over the gauge variables at all endpoints it
reads. Those site sets are disjoint between tiles. Wilson gauge
invariance makes the product expectation in
\eqref{aw:witness-product} unchanged. This averaging preserves the
support, the interval \([0,1]\), and translation covariance.
No reflection covariance of an individual coin realization is needed.

\paragraph{The isolated-cell payment}
\label{sec:aw-warm}

Suppose \(e^{+2}\cap Z_j=\varnothing\).
With \(I\) as in \eqref{aw:isolated-exponent}, varying its free
coordinates inside \(\prod_{b\in I}\{|y_b|<\mathfrak b_{j-1}\}\)
changes no explored collar or rim answer: every actual such query
touches a stage set \(Z^{[n]}\subset Z_j\), at index distance at least
three from \(e\), and therefore reads no bond of \(I\).
The original split queries that do read \(I\) are the full-threshold
cube tests, which stay passed. The precomputed lower-threshold
diagnostic flags may change, but none of them is read by this
recursion. Sum over these unused flags before taking the warm-cell
conditional law. The isolation factorization concerns the original
geometry and column event, not the artificially refined diagnostic
string. The approximate tests stay passed by their factor-two
margin and their centre offset. At \(j=1\), the bare density margin
used in \cref{lem:aw-split-reads} also applies. Thus isolation factors
as the indicator of this product window times an exterior indicator.
The region, the later run, and a column event are fixed by that
exterior datum. Earlier gauge coins must be integrated out: fixing
them would constrain the free coordinates to contour level sets.
Later coins are independent of the current stage.

Here is the quantitative conditional estimate. Fix a marked bond
\(b\) of \(e\), and let \(w(b')=e^{-\mu|b'-b|}\). For small fixed
\(\mu>0\), \eqref{aw:gaussian-input} gives
\[
 \left\|\tfrac12(wKw^{-1}+w^{-1}Kw)-K\right\|\le C\mu,
 \qquad
 \tfrac12(wKw^{-1}+w^{-1}Kw)\ge a_-I/2.
 \]
The same estimates hold for the full principal compression \(K_{II}\).
The exterior force has weighted norm at most
\(Cp_0(g_{j-1})e^{-cLM\check R_j}\); the nonlinear force has norm
at most \(C\epsilon_{\rm nl}\). In particular both together are
small, even though the unweighted exterior force need not be small.

For completeness, a product of convex windows admits the comparison
\[
 |\nu_{H_1}(f)-\nu_{H_0}(f)|
 \le \lambda^{-1}{\rm Lip}_w(f)
       \left\|w\sup|\nabla H_1-\nabla H_0|\right\|_2
 \]
when the conjugated symmetric part of \(H_0''\) is at least
\(\lambda I\). To see this, approximate the windows by separate
convex walls in each bond coordinate and couple the two overdamped
Langevin diffusions with the same Brownian motion. Weighted squared
distance contracts under the first drift; the remaining force is
bounded by the displayed weighted norm. The convex wall terms have
the same sign because the weight is scalar on each bond. Pass to
stationary couplings and then to hard walls. No convexity of \(H_1\)
is needed for this comparison.

Compare \eqref{aw:isolated-exponent} with the even centered quadratic
form on the same symmetric window. The comparison gives
\(|\nu_{H_e}(y_b^\alpha)|\le1\) after reducing the coupling ceiling.
By \eqref{aw:nonlinear-input}, \(H_e''\ge a_-I/2\). Comparing it
with \(H_e-t y_b^\alpha\) and integrating in \(t\) yields
\[
 \nu_{H_e}\left(|y_b^\alpha-\nu_{H_e}y_b^\alpha|\ge u\right)
 \le2e^{-a_-u^2/4}.
 \]
There are at most \(C(LM\check R_j)^4\) marked bonds. A union bound
therefore gives
\[
 \nu_{H_e}(e\ {\rm warm})
 \le C(LM\check R_j)^4
       e^{-c(\mathfrak b_{j-1}^{\flat\flat}/\sqrt3-1)^2}
 \le q_a^{r_j}.
 \]
The last inequality uses \(2\mathfrak p_0>5r\), with \(c_a\)
chosen before the coupling ceiling. Peeling disjoint boxes gives
the product of these factors, relative to any column event.
No derivative of the curvature or independent-current matching is
used in this concentration argument.

\paragraph{From the witnesses to the chessboard bound}
\label{sec:aw-chessboard}

The remaining geometric and mass ingredients are the following
ones of the actual history construction.
First, the boxing/cascade recursion is local in its interval flags.
Its pieces belong to the \(c_{\rm lay}\)-clusters of hot cells,
and each piece stays within \(c_h=n_{\rm sm}+1\) cells of its cluster.
This follows by keeping an original seed in every component: a
small bounding box is within \(n_{\rm sm}-1\) cells of that seed,
at every round. No monotonicity of the boxing map is required.
The same recursion on two common clusters gives the same labelled
piece even if distant clusters require more stages.

Second, retain the following exact representation premise on the
actual completed positive terms at a healing step. With the live
region and its internal histories fixed, the term for a family
\(\mathfrak I\) of healed components is
\[
 \mathbf O_{\mathfrak I}\big[W(\mathfrak I)\big],\qquad
 W(\mathfrak I)=
 \left[\prod_{X\in\mathfrak I}\mathbf T'_X\right]
                \exp\sum_Y V_{\mathfrak I}(Y).
 \]
%
Each \(\mathbf T'_X\) integrates against a nonnegative
measure. To see heredity in the native history, restrict
to the clusters of all roots other than \(X\). At each
level the separation of distinct clusters exceeds the reach
of the bond tests, pivots and collars. Hulls and enlargements
are taken on the separated regions, and boxing acts on one
component at a time. The class test also reads births in
\((j-\mathsf N_j,j]\) and marks in \([j-\mathsf N_j,j)\).
A birth at \(s\) can have been healed at an earlier level
\(i<j\) only if \(s\le i-\mathsf N_i\); for a mark the
inequality is strict. Since \(\mathsf N_i\ge\mathsf N_j\),
\(i-\mathsf N_i<j-\mathsf N_j\), so no item read by the
class test can have been erased in this way. These two
facts give the same tests and classes on the remaining
components. Each old component belongs to the cluster
of its last successor. For \(l\le j\), healing \(X\)
fills \(X\) into \(\Omega_l\) and \(\Lambda_l\) and
removes it from \(S_l\). Omitting the cluster leaves
these same completed label entries on \(X\).
The family parametrization over this label drops just
\(X\) and its isolated inner history.

Write
\(\mathbf O_{\mathfrak I}\) for the positive outer operation
on the common real exterior variables. For removal of \(X\)
we need
\[
 \mathbf O_{\mathfrak I}\le\mathbf O_{\mathfrak I\setminus\{X\}}
\]
on nonnegative inputs. In the completed-output representation,
fix the final determining sequence \(\bar{\boldsymbol\Omega}\),
its datum and the standing components, including the second-class
components \(Y_\iota\). The outer plaquette factor in the native
output formula is
\[
 \chi_j\!\left(
  \left(\bar\Omega_j\cap\bigcap_\iota(Y_\iota^c)^{\sim2}
  \right)^{\sim4}\right).
\]
%
%
On the healed branch the passed inset test is incorporated in
this output factor. The preparation's shell guards occur in
the island integral. In both terms of the comparison,
evaluate the output at the same retained datum, including
the restored top field. The completed determining rows
and the live-component operations are then common, as
are this outer factor and the action factored outside
the bracket. The localized interactions that read the
inner history remain inside \(W(\mathfrak I)\). In a
representation with a top mask
depending on the family, we must compare the preparation and
ordinary masks on the same level-\(j\) cubes at the actual
exterior field. The implication from the preparation mask to
the ordinary mask is sufficient, since multiplication by
these ordered indicators is composed with the same positive
maps. Resetting the coupling profile identifies the scale of
the tests. The comparison also needs their actual readouts
and supports.

%
Removal leaves the terms away from \(X\) unchanged and
independent of its integration variables. For the near terms
on the insertion side, follow a chain of cubes from the
enlarged strip of \(X\) to a cube outside the exempt strips.
Such a cube exists for the interaction domains of the
bracket. The first exit from the strip touches \(X\)'s
strip and is outside every other strip by their separation.
On the removal side, root the term at a cube of \(X\)'s
enlarged strip, now outside the remaining exempt strips.
In either case the interaction tree meets its root cube.
There are \(C\check R_j^4\) possible roots. The rooted
entropy bound and the strict decay reserve give a finite
geometric sum per root. Thus the sum of the absolute near
terms in the two brackets is bounded by \(C\check R_j^4\)
on their common admitted core. This uses the actual
localized interaction domains and the strip separation.

Consequently, pointwise inside the bracket, inserting
\(X\), with its internal history in
\(\mathscr H\), costs at most
\(t_X(\mathscr H)e^{C\check R_j^4}\), where
\(t_X(\mathscr H)=\sum_{\mathfrak h\in\mathscr H}\mathbf T'_{X,\mathfrak h}1\).
The trace mass bound of \cref{prop:large-field-2-6,prop:large-field-2-7},
with its charge dial increased by \(c_a\), gives
\[
 K_H\sup_X t_X(\mathscr H)e^{C\check R_j^4}
 \le\sup_{\mathfrak h\in\mathscr H}
        q_a^{\,2\sum_{i\le j}r_i|\mathsf Z_i(\mathfrak h)|_i}.
 \]
Here the unspent factor \(e^{-cp_0(g_j)}\) pays both the strip
exponent and the finite shape count \(K_H\).
The condition \(\mathfrak p_0\ge6r+1\) pays every increased lineage
charge, whose logarithmic degree is at most \(6r\).
The supremum in this estimate is on that same actual core and
is taken before the outer integration.

More precisely, let \(E_{c,\mathscr H}\) specify a healed member
at the cell \(c\), and let \(E_c\) allow any internal history there.
If \(B\) ignores earlier healed members and is unchanged by
deleting the member at \(c\), heredity and the pointwise comparison
give
\[
 \widetilde\nu(B\cap E_{c,\mathscr H})
 \le\left(\sup_{\mathfrak h\in\mathscr H}
               q_a^{\operatorname{Vol}(\mathfrak h)}\right)
          \widetilde\nu(B\setminus E_c).
 \]
Apply \(\mathbf O_{\mathfrak I}\) to the bracket
bound and dominate it on the reduced bracket by
\(\mathbf O_{\mathfrak I\setminus\{X\}}\). Then apply
the same subsequent positive kernels to the two sides. A joint survival
event has the required blindness: healing removes an entire
component, so its earlier lineage cannot also be the ancestor of
a component still live later. Peeling in increasing healing time
therefore proves the product bound for several exclusive healed
members. A bound on total polymer mass alone would not justify
this relative comparison.

Use the collar lengths
\[
 \begin{aligned}
 f_j&=\left\lceil\frac{M_1(\lceil D_j\rceil+C)}{M\check R_j}\right\rceil,\\
 c_{\rm adj}(j)&=2c_h+c_{\mathbf R}+f_j+C,\\
 m(j)&=1+2(k-j),\qquad
 m_0(j)=6(c_{\rm adj}(j)+1)m(j).
 \end{aligned}
 \]
At each level reserve \(m_0(j)+f_j+C(c_h+c_{\mathbf R})\)
additional cells, beyond the preceding depth expressed in the new
units. The additional two buffers have this same order. The total
physical width is bounded by
\[
 C M\check R_k+
 C M_1\sum_{n\ge0}(1+n)L^{-n}
     \big(\log\ell_\square+n+\check R_{k-n}+\log\bar p_{k-n}\big).
 \]
Since \(\check R_{k-n}\le C(1+n)\check R_k\), it is at most
\[
 M\check R_k
 \left[C_1+C_2\frac{\log\ell_\square+C+\log\log x_k}{\check R_k}\right].
 \]
Choose a power of \(L\) with
\(\ell_\square\asymp M\check R_k\), with its fixed
constant large enough for
\(\ell_\square>\max\{128,2n_{\rm sm}\}s_k\) and
\(\ell_\square\ge32\max_j B_j s_j\).
The width bound above allows this choice, since
\(\log\ell_\square/\check R_k\) tends to zero.
The core of side \(3\ell_\square/4\) leaves every
ancestor of a deep cell at depth at least \(4B_j\)
in its level-\(j\) partition.

%
The comparison of runs is an induction on the level.
Measure the preceding common depth in level-\(j\) cells
and round up to \(T_0\). The double blocks for a level-\(j\)
average stay in the cells of their centres, so the current
field is common on this inner lattice. Its inherited
live hull \(Z'\) is also common from \(T_0\) on. A cell
of this hull can be touched by a component of the preceding
live region from outside the inner lattice. The induction
therefore keeps the entire preceding components that meet
the common region, including those that only touch it.

At a step \(j\ge2\), a plaquette test is inserted outside
\(Z'^{\sim4}\). Its localized background, including the
existence check, uses only the full four-layer cube, which
lies in the cell and its neighbours. Thus the failure
set \(P_j\) is common from \(T_0+4\), its hull from
\(T_0+5\), and the gauge region
\((P_j'^{\sim1}\cup Z'^{\sim5})^c\) from \(T_0+6\).
The gauge draws use their own blocks before any draw and the
same tile coins at the drawn sites. For a tested bond \(b\)
both runs therefore obtain
\(u(b_-)V_{j-1}(b)u(b_+)^{-1}\) with the same factors.
The approximate test reads this field and the background
of its local test cube. Its one-cell neighbourhood puts
\(Q_j\) at depth \(T_0+7\) and \(Q_j'\) at \(T_0+8\).
The union defining \(\Omega_j^c\) is common from \(T_0+10\).

At \(j=1\) the inherited live hull is empty. Start at
\(T_0=1\) with \(P_0'\). A failed bare cube contributing
to its hull there has side at most \(L^{-1}s_1\).
Its density test adds at most \(L^{-2}s_1\) to the read,
so the bare plaquettes tested there are at distance at
least \((1-L^{-1}-L^{-2})s_1>0\) from the tile boundary.
The bare bonds of the tile determine this first hull.
With the first-step masks above, the depths for \(P_1\),
\(P_1'\), the gauge region, \(Q_1\) and \(Q_1'\) are
\(T_0+1,T_0+2,T_0+3,T_0+4,T_0+5\), respectively.
The first \(\Omega_1^c\) is therefore common from \(T_0+7\).
The bare regularity test supplies regularity on the
first determining region \(\Omega_0=\Omega_1^\sim\).

%
%
Call the common depth for \(\Omega_j^c\) \(A_0\).
The components of \(\Omega_j^c\) truncated to depths at
least \(A_0\) that touch its outer layer are common.
One with diameter exceeding \(m_0(j)\) is a first kind
of trigger. Otherwise these components lie at depths
at most \(A_0+m_0(j)\). From the region of depth
at least \(A_0\) remove their closed \(4M_1\)-neighbourhoods
and take the resulting union of \(M_1\)-cubes.
Its \(4M_1\) inner rim lies in \(\Omega_j\): a bad
cell meeting the rim would either touch the outer
layer or be adjacent to a removed cluster, because
\(10M_1<M\check R_j\) in level-\(j\) units. In either
case it would itself have been removed.
The region contains depth \(A_1=A_0+m_0(j)+2\).
At depth \(A_2=A_1+f_j\) the tested bonds have distance
at least \(f_jM\check R_j-5M_1\ge\lceil D_j\rceil\)
from its complement.

Use the common retained determining rows and data in
this region and the regular rim in
\eqref{aw:collar-response}. At the first step the
determining region includes \(\Omega_1^\sim\). A retained
bad cell in the comparison region is in a component
disjoint from the outer layer, so its touching cells
are still in the common \(\Omega_1\) data. In either
case the post-draw link is common. Bi-invariance of
the \(\operatorname{SU}(2)\) angle bounds the difference
of the two relative angle moduli by
\(\frac{\pi}{2}|V^{\rm bg}(V^{{\rm bg}\prime})^{-1}-1|/g_{j-1}\),
which is strictly less than the excluded window width
\(u_\sigma\). A relative group element at the cut locus
has angle \(\pi\), and the perturbed angle is still above
the largest small-field threshold. Thus the split
interval flags and the hot-cell clusters at depth
\(A_2\) are common. A long boundary cluster of these
hot cells gives the second kind of trigger.

%
Use the depth \(t\) in level-\(j\) cells. Set
\(h_j=m_0(j)+c_{\rm adj}(j)+c_h+3\) and
reserve the two depths \(A_\alpha=A_2+h_j\) and
\(B_j=A_\alpha+h_j\). These layers are included in the
width estimate above. If there is no trigger, a cluster
of hot cells at adjacency range \(c_{\rm adj}(j)\)
meeting \(t<\theta+c_{\rm adj}(j)+1\) has its
\(c_h\)-hull in \(t<\theta+h_j\), for
\(\theta=A_2,A_\alpha\). This follows by adding its
diameter budget \(m_0(j)\) and the hull enlargement.
In particular any cluster reaching the \(c_h\)-neighbourhood
of a cell at depth \(A_\alpha\) or greater stays beyond
\(A_2+c_{\rm adj}(j)\). Thus exterior hot cells cannot join
it at the smaller boxing adjacency \(c_{\rm lay}\).

Let \(\mathcal K\) be the union of the common hot clusters
meeting \(t<A_\alpha+c_{\rm adj}(j)+1\). Its \(c_h\)-hull
stays below \(B_j\). If a boxing cluster meets \(\mathcal K\)
or leaves the common region \(t\ge A_2\), its output pieces
stay in this hull or within \(c_h\) of \(t<A_2\).
The other clusters, their labels and their earlier-stage data
are common by the induction. They have separation greater
than \(c_{\rm adj}(j)\) from \(\mathcal K\).
An output piece is within \(c_h\) of its cluster, so its
distance from \(\mathcal K^{+(c_h+1)}\) is at least
\(c_{\rm adj}(j)+1-c_h-(c_h+1)=c_{\rm adj}(j)-2c_h\).
Choose the fixed additive constant in \(c_{\rm adj}(j)\)
so that this is at least \(c_{\mathbf R}+f_j+4\).
The same depth buffer separates the piece from the
unmatched outer data. Thus the \(c_{\mathbf R}\)-enlargement
used by a component-local preparatory read is inside
the common data. To apply this at a preparatory stage,
keep the fields before the operation as well as the
fields after each draw. The latter can overwrite a field
whose earlier value the next test reads.

%
%
%
For a component \(X\) in the common inner region,
the initial record consists of those incident fields,
the trace of the label, the current failed tests and
cascade cells, the earlier live components with their
records, and the coins assigned to \(X\). Append the
answers in their order of use. Equality of this record
compares physical link values and discrete decisions.
Relative coordinates formed with different global
centres can still differ. Their later uses must be
through an original recorded answer or a local read,
as in the completion test below. The waiting, quietness
and single-component tests are functions of this trace.
For an eligible component the boxes \(X\cap Z_i\) and
\(X\cap\Omega_i^c\) in the waiting window give the
prepared nest and every collar. Hence their row masks,
including the incident rows, are common.

At stage zero and at each successor stage, the
regularity test uses the localized background with
the joined datum of \(X\) and its prescribed lifts
of the top field there. Its determining rows and
their stored values are common. An undefined local
background gives the same failure on both sides.
For a stage integrating level \(\iota\), the gauge
draw reads complete blocks with centres in
\(((\Omega_{\iota+1}^{\sim3})^c\setminus
Z''_{\iota+1})\cap X\) that lack the earlier gauge
factor. That mask is common. Their contours are computed
from the original retained \(V_\iota\) links of these
blocks, which are stored as Haar links. The common
contour and the same normalized draw coin give the
same gauge transformation at each site. The
resulting incident links agree and the coarse
centres are fixed. The approximate test then reads
these links against the common local background.

%
Consider an ordinary completion test at a stage with
stored level \(\iota\), and put \(t=\iota+1\). The
completion reads a free bond with an endpoint in
\(\Omega_t\) from a test cube of the ordinary boundary
layer inside the quiet component \(X\). Its full bond
read lies in \(\operatorname{Box}_t(X)^{+2}\subset X\cap Z_t\),
where \(\operatorname{Box}_t(X)=\Omega_t^c\cap X\).
The waiting condition excludes creations there in the
stage window. The hull of the failed cubes at level \(t\)
is at distance at least five level-\(t\) cells from
\(\operatorname{Box}_t(X)\). Since the read of a failed
cube stays in its own cell and the touching cells,
the completion bond is outside the failed cube masks.
It therefore belongs to a previously passed mask or
to the plain strip. On the passed mask the recorded
ordinary bound makes the completion indicator one.

In the plain strip the test reads the stored relative
logarithm \(\frac1i\log(V_\iota(b)V_\iota^{\rm bg,loc}(b)^{-1})\).
The centre is the axial representative outside the
nodes of the joined problem on the corresponding
strip component of \(\Lambda_t^c\cap\Lambda_\iota\).
The strip nodes fix this representative at the test
bond. The joined problem uses the pre-operation
\(V_t,V_\iota\) of \(X\cap Z_t\).
The lineage in the waiting window stayed live through
each intermediate operation, so none of those operations
wrote its fields. The level-\(t+1\) gauge draw writes
outside the enlarged live region. Hence the inputs
have their stored values and are common in the two
runs. This gives the same completion answer even when
the test fails. It is used before testing the new
sharp ball. The latter uses the localized next-stage
centre and the post-draw link, both common. On each
outcome keep its actual prefix.

%
%
On a passing stage the deleted slots are the stored
nonpivot fluctuations on \(Z_{\iota+1}\cap
\Omega_{\iota+1}\cap X\) and the collar nonpivots on
\((\Omega_{\iota+1}^c\setminus Z''_{\iota+1})\cap X\).
The translation to the expansion centre is an integration
coordinate for the latter group. The real weight can use
the global expansion background here. The physical
retained state keeps the prescribed \(V_{\iota+1}\)
row and the surviving slots in their own charts.
On a failure of the local sharp test, the stored prefix
uses the untransformed real variables. Thus the two
runs agree on the retained physical record through
either branch. If a later read uses a translated
coordinate that is actually retained, its dependence
must be checked separately.

%
%
After the last collar, the extendibility test asks
for an extension of the common exterior coarse field
over the main box in the component problem. If this
passes, the Haar refresh uses the box coins and keeps
the old variables for the prepared datum. The same
component frame fixes both trees and the entrance
root, so the tree gauge gives the same shell variables.
The main fluctuation test reads the resulting field
in the prepared shells, and the last-shell test reads
the innermost stored field with its prescribed identity
data. The localized reference law has the same exterior,
frame, tree slice and cutoff in the two runs. Its
denominator is the integral of that same density on
the same slice, so the draw and any denominator failure
agree. The test of \(\chi_j(X^{\sim-6})\) is made
on the newly drawn field in the inset test cubes.
Its full local test neighbourhoods stay in \(X\).
Thus the healing decision agrees. The subsequent
deletions or inert declarations, the rotation on the
healed branch and the label update are functions of
the common record and the component coins.
Every bond in the final deep region either meets
one of these common components or is unchanged.
The common surviving labels specify which lower
slots remain. This is the induction output at level \(j\).

For the induction on the actual run we must verify that
any alternative stage or main-step chart still reads the
same incident rows and masks just identified. This is the
local-read input to the witness construction below.
A sharp test centred on a global determining set of
earlier outcomes requires a separate comparison.

If there is a trigger, take its first level \(J\).
The preceding survival record is common. At \(J\)
use the first kind of trigger if it occurs, since
the truncated bad components are already common.
Otherwise the split comparison gives the common
second kind of trigger. Along a cluster path
starting at its outer end, mark the first cells
reaching successive index radii
\(3(c_{\rm adj}(J)+1)(i-1)\), \(1\le i\le m(J)\).
From the initial cell one of the diameter endpoints
is farther than \(m_0(J)/2\). Take a path to it.
Its bounded steps give a pattern
of \(m(J)\) cells, mutually separated
by at least \(2c_{\rm adj}(J)+3\). Choose the cluster
and the path by a fixed order of cells in the tile. The
choice uses only the common data at the first trigger.
For a marked cell \(e\), a bad domain or a failed full
test puts a cell of \(Z_J\) in \(e^{+2}\). The remaining
hot kind is warm whenever \(e^{+2}\) contains no
\(Z_J\)-cell. Thus the pattern event is that each mark
has a cell of \(Z_J\) within two cells or is isolated
and warm. For a deep cell \(z\) whose lineage survives
from \(n\) to \(k\), every completion of the tile data
has the same survival decisions below \(J\) and the
marked alternative at \(J\). The status of an individual
mark as isolated may vary with the exterior.
With no trigger, the decisions agree through \(k\)
by the same induction on the verified local reads.
These assertions use the split windows only for the
original run. Under the local-read input the two types
of events give witness sets in the
sense of \cref{sec:aw-witness}. Apply the no-window
bound \eqref{aw:many-tile-windows}, then take the union
over witnesses in each tile and use
\eqref{aw:witness-product} for the resulting products.

%
%
Keep the tiles whose indices relative to the original
tile are even in all four coordinates. The marks in
their \(e^{+2}\) neighbourhoods are disjoint. In each
replica classify each marked cell as isolated and warm,
or choose a cell of \(Z_J\) within two cells of it.
The warm-cell bound can be peeled relative to the
event for these chosen cells and the survival columns.
The remaining event reads the regions and lineages
unchanged by varying the free variables of an isolated
box. Every warm mark thus pays \(q_a^{r_J}\).

Classify the chosen cells and the column, if \(J>n\),
by fate. Group objects absorbed by the same healed
component and specify its healing level \(h\ge J\).
Index the component by the level-\(h\) parent of its
first object. Its lineage at \(J\) must contain exactly
the level-\(J\) cells of that group among the chosen
objects. For a group carrying the column, also require
its lineage to contain every prescribed column cell
at levels \(n,\ldots,J-1\). These restrictions allow the relative healing
bound to be applied once per distinct component. A
group with \(d\) chosen marks pays \(r_Jd\) units of
the exponent. It has a cell at each level in \((J,h]\),
which pays a further \(h-J\) units. A column pays one
unit per level in \([n,J-1]\). The column's top is
used to specify its fate and is not charged as a
second mark if it is already a chosen cell.
No healed component serves different replicas. Even
after moving to nearby chosen cells and projecting to
the healing grid, their separation across an intervening
tile is at least \(\ell_\square-6s_k>n_{\rm sm}s_k\).
Apply
\eqref{aw:prefix-moment} once to the surviving marks
and columns in all replicas, so there is a single
volume factor. Each replica has at most \(m+1\)
objects. The choices of nearby cells and the group
labels cost at most \((C(m+1))^{m+1}\); summing the
healing times costs at most \((1-q_a)^{-(m+1)}\).
The cost per replica is at most
\[
 v_{J,\mathbf e}
   =\big(C(m+1)(1-q_a)^{-1}\big)^{m+1}
       q_a^{(J-n)_++r_Jm}.
 \]
This includes the sum over absorption times. The pure-survival
alternative has cost \(v_{\rm I}=q_a^{\sum_{j=n}^kr_j}\).
The whole replicated probability carries only the single factor
\(e^{B(x_k)V_k}\).

The tile grid has even side count in each coordinate.
In the chessboard dissemination, the copies in the
even-index class are pure translates, since two face
reflections restore the orientation in that coordinate.
Integrate the independent coins before dissemination.
On the kept tiles, restore the coin integral using
\eqref{aw:witness-product}. By translation covariance
of the run, each witness cylinder forces its translated
event almost surely. Dropping the other copies, which
lie in \([0,1]\), therefore bounds the product by the
probability of the replicated events.
The chessboard root is of order \(|\mathcal T|\) and the
even-index class has \(|\mathcal T|/16\) tiles. Consequently the
chessboard norm of a witness is at most
\[
 e^{B(x_k)\ell_\square^4}v_{\mathfrak w}^{1/16}.
 \]
There are at most
\(\ell_\square^4L^{4(k-J)}(m_0(J)+1)^{4(m-1)}\) patterns at level \(J\).
%
%
To bound the pattern count, put \(u=x_k+k-J+C_*\)
with \(C_*\) fixed. The flow gives \(N_J=\check R_J\le
C(\log u)^r\). The inverse coupling at any level in the
range defining \(\bar p_J\) is at most
\(C(u+(\log u)^r)\le C'u\), so
\(\log\bar p_J\le C(1+\log\log u)\).
Also \(\log\ell_\square\le C(1+\log\log u)\).
The displayed formula for \(D_J\) now gives
\(D_J\le C(1+k-J+(\log u)^r)\).
Using \(f_J\le C(1+D_J)\), we obtain
\[
 m_0(J)\le C(1+k-J)(1+k-J+(\log u)^r)\le C u^2 .
\]
The lower flow bound gives \(r_J\ge cM^4(\log u)^{5r}\)
after choosing the coupling ceiling. Apart from the
starting position factor \(\ell_\square^4L^{4(k-J)}\),
the logarithms of the enumeration and absorption
factors are at most \(Cm(J)\log u\). Half the surplus \((r_J-1)m\) pays this pattern count and its
absorption prefactor. The remaining half supplies the factor
\(\widehat q^{(r_k-1)/2}\). Also
\[
 (J-n)_++m\ge(k-n+1)+(k-J).
 \]
Choose \(c_a\ge16\log(2L^4)\). Then
\(\sum_{J\le k}(L^4\widehat q)^{k-J}\le2\).
Finally \(r\ge2\) gives \(5r>4r+1\), so, by taking the ceiling
small, the single-level surplus pays
\(\ell_\square^8e^{B(x_k)\ell_\square^4}\), \(C_{\rm win}\),
and any fixed requested extra factor \(\epsilon\).
The witness sum is therefore at most
\(\ell_\square^{-4}\widehat q^{\,k-n+1}\);
the sum for an additional tile's deep level-\(k\) live cell is
at most \(\epsilon\). Their product proves the two bounds of
\cref{lem:source-domains-4-15}.

For the first bound the anchor can be anywhere in the
torus. When the tiling exists, shift it by level-\(k\)
cells so the level-\(k\) parent of the anchor lies in
the central cell of its tile. The tile side and cell
side are powers of the odd integer \(L\), so the
central cell is well defined and the shift is allowed.
The side of the stopped lattice is twice a power of
\(L\), so whenever it is at least \(2\ell_\square\),
the tile count per side is even. If the torus is too
small to contain this tiling, then \(V_k\le(2\ell_\square)^4\). On survival,
\(\Theta_k\ge2c_ar_k(k-n+1)\). The same surplus comparison gives
\[
 B(x_k)V_k\le c_ar_k,\qquad
 e^{B(x_k)V_k-2c_ar_k(k-n+1)}
 \le\ell_\square^{-4}\widehat q^{\,k-n+1}.
 \]
Thus the first bound follows directly from the stopped moment on
every admitted prefix, without a tiling. Only the joint-tile
statement requires actual tiles. The argument concerns survival
and hence lineage age; it does not assign an age penalty to a
seed that was healed before the lineage began.
\par\emph{Proof of \cref{prop:source-domains-4-14}.} \emph{Anchor:} choose any cell $\mathfrak{z}$ of the lineage region $\mathsf{Z}_n(C)$ at the first level $n=j_0=k-a$ of the lineage. The lineage regions are nested and contained in the live regions. Hence $\mathfrak{z}$ stays live throughout the levels $n,\dots,k$, so that the event $\mathcal{E}^{(n)}_{\mathfrak{z}}(k)$ occurs, and the cell belongs to the component of $\mathfrak{c}$. This is the point at which the age of the lineage enters: a region healed in the meantime provides no such cell. \emph{Sum over anchors:} a box of side $\mathsf{l}_\square$ contains at most $\mathsf{l}_\square^4L^{4(k-n)}$ anchors, and each costs $\mathsf{l}_\square^{-4}q_a^{c(k-n+1)}$. For an anchor at a distance of $i\ge i_0$ boxes, the live component is connected and therefore passes through a chain of at least $i$ distinct adjacent boxes, each of which contains a live cell. These cells need not be deep. Consider the $2^4$ tilings whose face planes are displaced by $\pm\lfloor n_\square/4\rfloor$ cells from a corner of the level-$k$ cell containing $\mathfrak{z}$, where $n_\square\ge16$ is the number of level-$k$ cells along a side of a tile. Every level-$k$ cell is deep in a tile of one of these tilings, and $\mathfrak{z}$ is deep in all of them. By the pigeonhole principle a single tiling serves at least $i/16$ of the cells. These cells lie in at least $i/16^2$ distinct tiles, since a tile meets at most $2^4$ boxes, and at least $i/i_0$ of these tiles, where $i_0=2\cdot16^2$, are different from the tile of $\mathfrak{z}$. For a fixed reference chain, the choice of a shift, the
subset of its indices served by that shift and the choice
of one of at most \(2^4\) shifted tiles per index cost at
most \(16\cdot 2^i\cdot 2^{4i}\) choices. The chain itself
has at most \(3^{4i}\) choices. The second bound of the
lemma sums the deep-cell positions inside the selected
tiles. Hence with \(C_{\rm ent}=2^5 3^4\) and
\(\epsilon=(2C_{\rm ent})^{-i_0}\), the additional factor
for such an anchor is at most \(16\,C_{\rm ent}^{\,i}
\epsilon^{i/i_0}=16\cdot2^{-i}\), which is summable against the $(2i+3)^4$ boxes of a shell. Finally $L^4q_a^c\le(2L^b)^{-2}$, and this is exactly the floor $c_a\ge C(b)\log L$ imposed in the statement of the proposition. Therefore $\tilde\nu(\mathfrak{lv}^{=a}_k(\mathfrak{c}))\le C(2L^b)^{-2a}$. The geometric mean of this bound and that of \cref{prop:source-domains-4-13} is the first bound. Summing it gives the age series

\begin{gather*}
\FitMath{\sum_{a\ge0}L^{ba}\,\tilde\nu\big(\mathfrak{lv}^{\ge a}_k(\mathfrak{c})\big)\le Ce^{-cp_0(g_k)},}
\end{gather*}

which two later parts use. In \cref{sec:short-distance-isolating-a-single-scale-coupling-and-normalisation-at-the-scale} the series is needed with rate $b\ge8$, because for two points a live component of age $a$ carries a weight of up to $L^{8a}$; there the series underlies the two-point law, \cref{thm:continuum-main} (c)(3), (c)(4). In \cref{thm:continuum-main}(c)(5) the series enters with rate $b>16$, on each torus, since a weight $L^{16a}$ arises for four points; this concerns clause (c)(5). The expansions of \cref{sec:short-distance-isolating-a-single-scale-coupling-and-normalisation-at-the-scale} and \cref{thm:continuum-main}(c)(5) require the age up to which a live component stores fields. This is the age read in the label. For a component holding an island, the block of which retains the fields of its core \cref{sec:large-field-stored-variables-and-integration-in-a-large-field-patch}, it is at most the age of the lineage. Here, as in \cref{thm:continuum-main}(c)(5), the series is formulated with the age of the lineage, and this age is never smaller. The series therefore covers what these expansions require: 

\begin{gather*}
\FitMath{\{\mathfrak{c}\subset C,\,\operatorname{age}^{\mathrm{lab}}(C)\ge a\}\subset\mathfrak{lv}^{\ge a}_k(\mathfrak{c}).}
\end{gather*}

For the second bound, each old live cell is carried into its
whole parent cell at the next scale until its component is
healed. A merge retains its lineage through that level.
Thus if \(\mathfrak c\) has lineage age \(a\) at level \(k\)
and survives through \(K\), its parent at level \(K\)
has lineage age at least \(a+K-k\). Sum the first bound
at \(K\) over these terminal ages.

\emph{The two-term age bound.}
The prepaid trace bounds give a second age estimate on the tori
without using the tiles and windows of \cref{lem:source-domains-4-15}.
Decompose the event according to whether the lineage is healed at
some level \(j>k\) or stays live to the last level \(K'\) of the
run or its prolongation. Fix \(c_a>0\) and choose the coupling
ceiling small enough for the estimates below.

%
The event is empty if \(k-a<1\). For \(j\ge k\) write
\(\mathfrak c^{(j)}\) for the cell at level \(j\) containing
\(\mathfrak c\). For a top \(X\) at level \(j\), let \(\mathscr G_X\)
be the traces with first level \(1\le j_0\le k-a\).
If the component at \(\mathfrak c\) is carried into \(X\),
its lineage up to level \(k\) lies in the lineage of \(X\).
On \(\mathfrak{lv}^{\ge a}_k(\mathfrak c)\) this puts the trace
of \(X\) in \(\mathscr G_X\). Every level from \(j_0\) through
\(j\) contributes a cell to the lineage. At \(\mathsf w=0\)
the mark is therefore at least \(\bar c_a(a+j-k+1)\), where
\(\bar c_a=2c_aM^4\). The prepaid trace bound in the proof of
\cref{prop:source-domains-4-13} gives
\[
 \sum_{\mathfrak l\in\mathscr G_X}\operatorname{wt}(\mathfrak l)
 \le e^{-\bar c_a(a+j-k+1)}e^{-c_{\rm cell}|X|_j}
       e^{-2\min_{1\le i\le k-a}p_0(g_i)} .
\]
We fix \(c_{\rm cell}\) before choosing the coupling ceiling.

\emph{Survival to the terminal level.}
For each history in the survival event distinguish the unique
terminal component containing the cell
\(\mathfrak c^{(K')}\). Group the histories by their terminal
traces. Positivity and the mass bound at fixed traces bound the
unnormalized mass, after removing \(e^{E_{(K')}}\), by
\[
 e^{C|\mathbb T^{(K')}|}
 \left[\sum_{X\ni\mathfrak c^{(K')}}
       \sum_{\mathfrak l\in\mathscr G_X}\operatorname{wt}(\mathfrak l)\right]
 \left[\sum_{\mathscr C}\prod_{Y\in\mathscr C}
       \sum_{\mathfrak l:\,\mathrm{top}=Y}\operatorname{wt}(\mathfrak l)\right].
\]
Here \(X\) ranges over connected terminal cell unions, and
\(\mathscr C\) over collections of distinct such unions. For the
second bracket we drop the exclusion with the distinguished component.
Choose \(c_{\rm cell}\) so that the sum of
\(e^{-c_{\rm cell}|X|_{K'}}\) over connected sets containing a given
cell is at most one. The first bracket is then at most
\(e^{2-2p_0(g_k)-\bar c_a(a+K'-k+1)}\), by the drift
\(p_0(g_i)\ge p_0(g_k)-1\) for \(i\le k\). For the second
bracket the unrestricted trace bound gives total activity at
most one per cell. The product bound \(1+z\le e^z\) then gives
at most \(e^{|\mathbb T^{(K')}|}\), since the number of cells is
bounded by the number of sites. Divide once by the partition
function floor at \(K'\), using \(\mathsf B(x_{K'})\) to absorb
both volume charges and the floor. For a prolongation by \(p\)
levels on the same bare lattice, the number of terminal sites is
\((\ell L^{-p})^4\), and the history law through level \(k\)
is the same. A Markov bound on the global end mark alone would
lose the factor \(e^{-2p_0(g_k)}\).

\emph{Healing at level \(j>k\).}
Use the outer-operation domination of the relative
healing comparison above. Inside the bracket, deletion of the
healed member with top \(X\) costs at most its restricted trace
sum times \(e^{C_I\check R_j^4+C}\). A healed top is an aligned
box of at most \(n_{\rm sm}\) cells per side, so at most
\((\sum_{n=1}^{n_{\rm sm}}n)^4\) such boxes contain
\(\mathfrak c^{(j)}\). Discard the rooted size charge and split
the oldest creation payment using the drift to \(g_k\) and \(g_j\).
For \(0<c_{\rm r}<1/2\) the exponents allocated to these scales
are \(c_{\rm r}p_0(g_k)\) and \((2-c_{\rm r})p_0(g_j)\), up to
an additive constant. The ceiling gives
\(C_I\check R_j^4\le(1/2-c_{\rm r})p_0(g_j)\).
The probability at this healing level is bounded by
\[
 C e^{-\bar c_a(a+j-k)}e^{-c_{\rm r}p_0(g_k)}
     e^{-3p_0(g_j)/2}.
\]
The flow lower bound gives a bounded number of visits to any
\(x_j\)-interval of length \(\lambda_\sharp\), and the ceiling
gives \(e^{-3p_0(g_j)/2}\le x_j^{-3}\). Thus the last factor
is summable in \(j\), uniformly in the run. This yields the first
term of the bound without using the terminal partition function floor.

In particular, for each level-\(k\) cell \(\mathfrak c\) and \(a\ge0\),

\begin{gather*}
\FitMath{\tilde\nu\big(\mathfrak{lv}^{\ge a}_k(\mathfrak{c})\big)\le Ce^{-\bar c_aa}\,e^{-cp_0(g_k)}+C\exp\Big[\mathsf{B}(x_{K'})\,|\mathbb{T}^{(K')}|-2p_0(g_k)-\bar c_a\,(a+K'-k+1)\Big],}
\end{gather*}

where $K'\ge K$ is the last level of an arbitrary prolongation of the run \cref{sec:flow-units-common-prefixes-and-change-of-reference}, the run itself counting as one of its prolongations. The first term contains no volume, and no condition ties $k$, $a$, $K$ and $\ell$ together. The consequences of the bound follow from its second term, case by case; none of the cases is a hypothesis of any statement. (1) There is a side $\ell_V(g)$, tending to infinity as $g\to0$, with the following property. For every torus with $\ell\le\ell_V(g)$, the prolongation of (iv) terminates in the alternative where the torus is reduced to polylogarithmic size. In that case the gain $2p_0(g_k)$, of degree $\mathfrak{p}_0$ in the logarithm, dominates the volume term, of degree $4r+1$. The bound is then $Ce^{-\bar c_aa}e^{-cp_0(g_k)}$ at every level. (2) The last lattice of the reference torus $T_{\ell_\star(g)}$ of \cref{prop:short-distance-7-3} has polylogarithmic size even without prolongation. Hence that torus satisfies the same bound at all levels. (3) If $\ell>\ell_V(g)$, the torus satisfies the same bound at those levels at which $\bar c_a(K'-k+1)\ge\mathsf{B}(x_{K'})|\mathbb{T}^{(K')}|$. This is a floor on the depth that depends on the torus; once it is met at one level, it is met at every lower level. At all of these levels, if $\bar c_a\ge b\log L+\log2$, summing the bound gives the age series at the rate $b$. Thus this road alone supplies what \cref{sec:short-distance-isolating-a-single-scale-coupling-and-normalisation-at-the-scale} uses, namely the age series at the rate $b\ge8$, at all levels of the reference torus, as used in \cref{thm:continuum-main} (c)(3), (c)(4). For \cref{thm:continuum-main}(c)(5) it supplies the series at $b>16$ on each torus. On the tori of (1) it holds at all levels; on the tori of (3) it holds at the levels below one that depends on $\ell$. The additional four-point hypothesis in \hyperref[hyp:four-point]{the short-distance statement} has exactly this quantifier structure: $\mathsf{d}\le\mathsf{d}_4$ if $\ell\le\ell_4(g)$, and $\mathsf{d}\le\mathsf{d}_4(\ell)$ otherwise, where $\ell_4(g):=\max\{\ell_V(g),\ell_\star(g)\}$. The contribution of \cref{lem:source-domains-4-15}, through its tiles, windows and triggers, is decay at all levels and on all tori, with constants independent of $\ell$. This is uniformity as $\ell\to\infty$. The same age threshold therefore works on every torus. Passing a
four-point estimate to limits on \(\mathbb R^4\) additionally requires
the actual normalized interacting comparison and its \(o(g^8)\)
remainder, with the prescribed test factors, uniformly along the
volume and cutoff limits; see
\eqref{eq:four-point-unresolved-interacting-target}. These further
inputs are not deduced here. The torus and additional-hypothesis
restrictions in \cref{thm:continuum-main}(c)(5) remain in force.

\emph{The age constant.} In \cref{sec:large-field-the-induction-closed-the-order-of-the-constants-the-sum} two roles of $c_a$ are described; the one at work here is role (A), the dial of the mark. Its effect on the law of the histories is $C(2L^b)^{-a-(K-k)}$. Role (B), the age rate $q_a$ of the aged load \hyperref[notdef:aged-load]{$\Sigma_k$}, concerns the pairs of lists belonging to a class of the kind given there \cref{sec:large-field-the-induction-closed-the-order-of-the-constants-the-sum}, and it is used in \cref{ch:comparison-capsule-v7,ch:weighted-matched-domains}. The floors on $c_a$ listed in \cref{sec:large-field-the-induction-closed-the-order-of-the-constants-the-sum} concern histories, rotations and age-resolved rarity. \Cref{prop:source-domains-4-14} imposes the third of these.


\clearpage
 \textbf{Cutoffs beyond first passage.} The dilation identity of \cref{sec:flow-units-common-prefixes-and-change-of-reference} gives

\begin{gather*}
\FitMath{S^T_{n;K,\ell,\beta}(f_1,\dots,f_n)=S^T_{n;K-q,\,L^q\ell,\,\beta}\big(f^{[q]}_1,\dots,f^{[q]}_n\big),} \\[3pt]
\FitMath{\|f^{[q]}\|_\sharp\le\|f\|_\sharp,}
\end{gather*}

since $\mathcal{O}$ carries no power of the spacing. The same relation holds between $\mathbb{L}_K$ and $\mathbb{L}_{K-q}$. Let $\beta\ge g^{-2}+B_\sharp$ and $q\ge0$, and suppose that $K-q=K(\beta,g)$. The bounds of \cref{prop:source-domains-4-2} are then valid at the cutoff $K$, with $(L^q\mathsf{d};L^qR;L^q\ell)$ in place of $(\mathsf{d};R;\ell)$.

% END MODULE chapters/07_sources.tex

% BEGIN MODULE chapters/ordinary_stability_deduction_v6.tex
\subsection{Stability from a normalized history sum}
\label{ch:ordinary-history-stability}

This is a conditional deduction of upper stability and the unsourced
partition floor. It isolates the estimates a renormalization construction
must supply. In particular it does not replace those estimates by a
polymer convergence theorem. A representation with a sum over large-field
labels and label-dependent integral operations need not be a single
coarse Wilson integral multiplied by one uniformly small polymer gas.

Let the final four-dimensional periodic lattice have $N$ sites and
$4N$ independently oriented bonds. Haar measure on each
$\operatorname{SU}(2)$ bond is normalized. Let $w$ range over a specified
ordinary complex source neighborhood containing zero. Assume there
are complex measures $\tau_S(w)$, indexed by a finite set of labels,
such that
\begin{equation}
 Z(w)=\sum_S\tau_S(w)(\mathcal V),\qquad
                      \tau_S(0)\ge0.
 \label{oh:exactness}
\end{equation}
Assume all normalizers and scalar quantities below are well defined;
$E$ is the real vacuum constant of the unsourced all-small-field
label. Put
\[
 \mathcal E(w)=\sum_{j,X}e_j(X;w),\qquad
 \mathcal E(0)=0,\qquad
                  \widetilde Z(w)=e^{-E-\mathcal E(w)}Z(w).
\]
At each finite cutoff the sums are finite. The argument requires their locality and uniform
estimates; it uses no sum over cutoffs.

\subsubsection*{A total-variation criterion for upper stability}

Let $\mathcal E_S(w)$ denote the vacuum constant in the exponent of
the term carrying label $S$. Suppose there are nonnegative weights
$W_S$ and real label costs $q(S)$ satisfying
\begin{equation}
 \|e^{\mathcal E_S(w)}\tau_S(w)\|_{\mathrm{TV}}\le W_S,
 \qquad
 \operatorname{Re}[-E-\mathcal E(w)-\mathcal E_S(w)]
                                      \le C_0N+q(S),
 \label{oh:term-bound}
\end{equation}
and the weighted label bound
\begin{equation}
                     \sum_S e^{q(S)}W_S\le e^{C_1N}.
 \label{oh:label-sum}
\end{equation}
The constants are independent of cutoff, volume and $w$ in the
specified source neighborhood. Then
\begin{equation}
                         |\widetilde Z(w)|\le e^{(C_0+C_1)N}.
 \label{oh:upper}
\end{equation}
Indeed multiply each measure in \eqref{oh:exactness} by the common
normalization, insert and remove its own $e^{\mathcal E_S(w)}$,
take total variations, apply \eqref{oh:term-bound}, and sum by
\eqref{oh:label-sum}. This proof takes no logarithm of $Z(w)$ and
asserts no absence of its complex zeros.

In a history construction, $W_S$ is obtained by bounding the
unsourced integrations after inserting a reserve factor such as
$e^{c v(\omega)}$ on every local history. The ordinary source changes
an island exponent by an extensive amount,
$|\Psi_\omega(w)|\le C_r\,n(\omega)$, where sites are counted
in their own scales. A sufficient reserve is
$c v(\omega)\ge C_r n(\omega)$, uniformly in the history and
the exterior arguments. This is an ordinary mass reserve for one
presentation, and is distinct from estimates comparing marked or
weighted derivatives of two different presentations.

It is essential that the majorant used for each component be a
function of its own integration variables. A bound on its integral
at each fixed exterior field is insufficient in a coupled product.
For example, let both component factors be
$f_1(x,y)=f_2(x,y)=\mathbf1_{|x-y|<\epsilon}$ on $[0,T]^2$.
Each integral of a factor in its own variable is at most
$2\epsilon$, whereas the double integral of their product has
order $\epsilon T$, not $\epsilon^2$.
A bound with the supremum over the admissible exterior pair
\emph{inside} the local integral supplies a component majorant;
Tonelli then permits multiplication when the integration-variable
groups are disjoint. The construction must also specify which
Wilson factors are allotted to those majorants. An operation with
no Wilson factor of its own has no independent plaquette-tail gain.

Neither \eqref{oh:term-bound} nor \eqref{oh:label-sum} is a
consequence of the last triangle inequality. They encode ordinary
renormalization estimates: suppression of creation events,
control of their genealogies, the allocation of Wilson and Gaussian
factors without reuse, and payment of the missing vacuum normalizers
of every large-field label. In particular, estimates with the volume
of each fine lattice cannot simply be summed to obtain a bound
with the volume $N$ of the final lattice.

\subsubsection*{The floor from one positive label}

Let $x\ge1$ be the last inverse coupling and set
$r=cx^{-1/2}$ with a fixed $0<c\le1$. Write
\[
 \mathcal B_r=\{V:\|V_b-1\|_{\mathrm{op}}<r
                                      \hbox{ for every bond }b\}.
\]
Suppose the all-small-field label $S_0$ has a density with respect
to final product Haar measure satisfying
\begin{equation}
 d\tau_{S_0}(0)(V)
   \ge e^E\chi(V)\exp[-x\mathcal A(V)-C_IN]\,dV,
 \quad
 \chi=1\text{ on }\mathcal B_r,
 \quad \mathcal A(V)\le C_Ar^2N\text{ on }\mathcal B_r.
 \label{oh:small-field-floor-input}
\end{equation}
Here $\mathcal A$ may be the action of a minimizing background.
If so, inclusion in the characteristic set and the displayed action
bound are background estimates which must be verified; they do not
follow merely by renaming $\mathcal A$ as a coarse action.

For normalized Haar measure there is an absolute $c_H>0$ such that
the operator-norm ball of radius $r\le1$ about the identity has
measure at least $c_Hr^3$. To prove this write
$U=\cos\theta\,1+i\sin\theta\,n\cdot\sigma$, so
$\|U-1\|_{\mathrm{op}}=2\sin(\theta/2)$ and normalized Haar
measure is $(2\pi^2)^{-1}\sin^2\theta\,d\theta\,d\Omega(n)$.
The interval $0\le\theta\le r$ lies in this ball. Since
$\sin\theta\ge2\theta/\pi$ on $[0,1]$, its measure is at least
$8r^3/(3\pi^3)$. Decreasing the constant handles the irrelevant
choice of an open instead of a closed ball.

Positivity in \eqref{oh:exactness} permits restriction to $S_0$,
then to $\mathcal B_r$. Independence of the $4N$ Haar bonds gives
\begin{equation}
 \begin{split}
 e^{-E}Z(0)
  &\ge e^{-(C_I+C_Ac^2)N}(c_Hc^3x^{-3/2})^{4N}\\
  &\ge\exp[-(C+6\log x)N].
 \end{split}
 \label{oh:floor}
\end{equation}
Since $\mathcal E(0)=0$, this is the required floor for
$\widetilde Z(0)$. The constant depends on the fixed choices and
on $C_I,C_A$, but not on cutoff or volume.

\paragraph{The terminal tests on a product of link balls.}
\label{oh:terminal-link-ball}
For the all-small-field label at the terminal level $n=K$ the
determining sequence on the torus is constant. Its only data
are the final bonds $V$. Its final label has no standing regions.
The invariant form of the step gives the density
\[
 \chi_n(T_\ell)\,
       \exp\{E-x_nA(U_n(V))+F_n(U_n(V))\}\,dV .
\]
The equality of $E$ with minus the empty label's vacuum sum
holds at this depth. The pointed regular estimate on the all-small-field
label has only one layer, at level \(n\). For the level-\(j\)
family let \(h=n-j\). The tents at all its base points sum
to one. Subtract its vacuum value and the corresponding
tent-action term. The constant-strength cancellation gives
a bound per base point of
\(CE_0[(1+h)L^{-\beta_{\rm T}h}+L^{-h}]L^{-4h}\).
There are \(L^{4h}N_n\) base points. Thus these terms sum
to at most \(CN_n\).

The vacuum-subtracted remainder has bound
\(C(1+h)L^{-4h}\theta_j^{-\kappa_0/2}\) in the rooted sum
at a level-\(j\) site. Use the second-order estimate on the
near liftable domains, and the norm tail for the other
domains, including wrapping ones. The same site conversion
gives \(C(1+h)\theta_j^{-\kappa_0/2}N_n\). The lower flow
bound and comparison of the frame with the flow make the
sum over \(h\) bounded, since \(\kappa_0/2>2\).
On this empty standing label, the remainders include the
terms left by healed histories and there is no boundary term.
The resulting
\(F_n\ge-C_IN_n\) and positivity justify keeping just this
label in the lower bound.

We verify the remaining assertions of
\eqref{oh:small-field-floor-input} for this density. Write
$\eta_n=L^{-n}$ in level-$n$ units, and assume that
$|V(b)-\mathbf1|<r$ on every level-$n$ bond. Telescoping the
unitary products gives $|\partial V-\mathbf1|<4r$.
For a density cube $\square$ take its actual test tower in
$A=\square^{\sim4}$. Its determining datum is obtained by
lifting $V$ from the top. To see the size of this datum,
give a finer link the identity within a top block and the
corresponding oriented $V(c)$ across a top block face.
Each elementary plaquette of this lift is either trivial
or has the holonomy of the corresponding top plaquette.
The one-step average of this lift is its coarser lift:
the two endpoint contours are trivial within each block,
and every joining segment has the same single face value,
or the identity if it crosses no top face. This verifies
the same description on the averaged rows of the tower.

All determining plaquette holonomies, including those made
using averages across the interfaces, are holonomies of
the face lift. Thus the datum has radius $4r$. Apply the
linear regularity theorem on the prescribed separated density
tower, with its tied radii and with $4r$ below the datum ceiling.
The tested
plaquettes lie in $\square^{\sim}$, inside the top member
of the tower by the test-cube geometry. Their defects are
therefore at most $4B_3r\eta_n^2$. The constant-sequence
background of $V$ satisfies the same bound on the whole
torus. Take $r=c_0x_n^{-1/2}$. On the real tube $x_n$ and
$\theta_n$ are comparable, and $\varepsilon_n\sqrt{\theta_n}$
is bounded below. A fixed $c_0$ therefore makes
$4B_3r<\varepsilon_n$ and places the datum and background
below their respective regularity ceilings. Every factor of the terminal density test is
one on the product of balls. The earlier fluctuation tests
are accounted for in the integrated density and its regular
and remainder totals.

There are $6N_0$ bare plaquettes and $N_0=L^{4n}N_n$.
Using the Wilson normalization gives
\[
 A(U_n(V))\le\tfrac12(4B_3r)^2\eta_n^4\,6N_0
                    =48B_3^2r^2N_n .
\]
The action cost $x_nA(U_n(V))$ per final site is bounded.
The four Haar balls per site give the $6\log x_n$ term in
\eqref{oh:floor}. At $n=0$
the same conclusion follows directly by telescoping the bare
plaquette products, with no background construction.

\subsubsection*{Local energy bounds and the first-passage window}

Suppose the source families are actual functions of the source
values in a fixed-layer enlargement $X^+$, holomorphic in its
finite-dimensional source parameters and real at real arguments.
Suppose their differences at the trivial field obey
\[
 |\Delta E_j(X,z;w)|\le C_Ee^{-\kappa d_j(X)},\qquad
 |\Delta R_j(X;w)|\le C_Re^{-\kappa d_j(X)}.
\]
These are source-minus-unsourced differences: assume explicitly
$\Delta E_j(X,z;0)=\Delta R_j(X;0)=0$.
Let the number of possible base points $z\in X$ be at most
$C_M(d_j(X)+1)$, as follows for domains covered by the vertices
of a cubical spanning tree. For
$e_j(X;w)=\sum_{z\in X}\Delta E_j(X,z;w)+\Delta R_j(X;w)$,
\begin{equation}
 |e_j(X;w)|\le C_\delta
                     e^{-(\kappa-\delta)d_j(X)}\qquad(\delta>0).
 \label{oh:energy-bound}
\end{equation}
This uses only $(m+1)e^{-\delta m}\le C_\delta$. If the rooted
entropy rate is $\kappa_D<\kappa$, choose
$0<\delta<\kappa-\kappa_D$. The resulting energy decay still
exceeds entropy. Locality, holomorphy, realness and vanishing at
zero pass through the finite base-point sum. Dependence on the
source function itself is a structural hypothesis here; it is not
implied by a bound on a chosen list of source coefficients.

If a first-passage argument supplies
$g^{-2}<x\le g^{-2}+B_\sharp$, then the constants in
\eqref{oh:upper} and \eqref{oh:floor}, when bounded by
$C(1+\log x)$, are uniform over the cutoff data at this fixed
$g$. This is the final elementary step producing constants of
the form $C_g$. It supplies no proof that an admissible run reaches
that window.

\subsubsection*{What a pointwise polymer criterion would supply}

There is a different sufficient mechanism when a construction
actually has a pointwise hard-core gas. Let its polymers be nonempty
subsets of a set of $N$ sites. Suppose, uniformly in a real coarse
field $V$, its real activities obey
\begin{equation}
 \sum_{\gamma'\not\sim\gamma}|z_{\gamma'}(V)|e^{a(\gamma')}
       \le a(\gamma),\qquad
 \sup_x\sum_{\gamma\ni x}|z_\gamma(V)|e^{a(\gamma)}\le b.
 \label{oh:KP}
\end{equation}
These inequalities remain valid after multiplication of all
activities by any $t\in[0,1]$. The convergent cluster
expansion makes its partition function $\Xi_V(t)$ nonzero and
continuous on this interval. It is real there and
$\Xi_V(0)=1$; hence $\Xi_V(1)>0$. The absolute cluster theorem
bounds its logarithm by
$\sum_\gamma|z_\gamma|e^{a(\gamma)}\le bN$: sum the rooted
cluster bound and then use the nonempty site supports in
\eqref{oh:KP}. The branch is the one continued from $t=0$.
It is real on the interpolation, and therefore gives
\[
                     e^{-bN}\le\Xi_V(1)\le e^{bN}.
\]
For complex source activities, the corresponding uniform criterion
gives $|\Xi_V(w)|\le e^{bN}$. With a verified common Haar/Wilson
integral representation and its action estimates, these facts
could be used to prove stability and a floor.

Applying this lemma to a label/history representation requires an
identity with that integral. Local polymer exponentiations on small-field
regions, or exponentiation of the sum over healed components, do
not by themselves provide a pointwise gas over all final coarse
fields. For the history representation the concrete sufficient
conditions are \eqref{oh:term-bound}, \eqref{oh:label-sum}, and
\eqref{oh:small-field-floor-input}; verification of those ordinary
conditions remains necessary for the construction.


% END MODULE chapters/ordinary_stability_deduction_v6.tex

% BEGIN MODULE chapters/all_state_source_transfer_v3.tex
\subsection{Source bounds for periodic Wilson states}
\label{ch:all-state-source-transfer}

Reflection positivity transfers ordinary complex source bounds from
finite Wilson tori to arbitrary periodic states through their DLR
equations. The argument requires neither geometric forgetting for
derivatives of inherited activities nor moment bounds for Gibbs states
of a perturbed interaction.

Fix an odd integer $L\ge3$, $a=L^{-K}\le1$, a coupling $\beta>0$, and
$\mathcal O_a(f)=\sum_p f(x(p))a_p$. All sums defining this observable
are finite when $f$ has compact support. Although its supremum norm
may grow as $a^{-4}$, it is bounded at each fixed spacing. Every
periodic Wilson DLR state is invariant under all lattice isometries
and positive at all lattice mirrors by
\cref{thm:amrp-all-periodic}.

\subsubsection{An exact domination lemma, before source estimates}

Use a cubic board of cells of side $l=L^j\ge a$, whose face
planes cut links in half. Take cubic tori of side $\ell=2L^m$ with
$m\ge\max(0,j)$, having an even number of cells in every direction.
Every sequence of sides tending to infinity in this family eventually
accommodates any fixed coarser board of allowed side. All sequences in
this section are taken in this family. When localization cubes have
side $Ma_s$ with the construction convention $M=L^{m_0}$, restrict
to the cofinal subfamily $m\ge m_0$. Then
$\ell/(Ma_s)=2L^{m-m_0+s}$ is an even integer at every scale, and all
localization partitions are retained. The cell algebra consists of
bounded functions of bonds lying wholly inside the open cell; crossing
bonds belong to no cell algebra. Reflected copies are conjugated when
the reflecting isometry has negative determinant. For a cell function
$F$ define

\refstepcounter{equation}
\begin{gather*}
\FitMath{p_{l,\ell}(F)=
 \left\langle\prod_{C} F^{[C]}\right\rangle_{a,\ell,\beta}^{1/N},} \\[3pt]
\FitMath{N=(\ell/l)^4.}\tag{\theequation}\label{eq:all-state-source-transfer-ST1}
\end{gather*}

The number inside the root is nonnegative. To see this, split the
board at one face and the opposite face. The disseminated product is
$\Theta\Phi\,\Phi$ for the product of its copies in one half.
Finite-torus reflection positivity applies.

We will use, including for complex functions,

\refstepcounter{equation}
\begin{gather*}
\FitMath{\left|\left\langle\prod_C F_C\right\rangle\right|
 \le\prod_Cp_{l,\ell}(F_C),} \\[3pt]
\FitMath{p(F+G)\le p(F)+p(G),} \\[3pt]
\FitMath{p(\bar F)=p(F),} \\[3pt]
\FitMath{p(1)=1.}\tag{\theequation}\label{eq:all-state-source-transfer-ST2}
\end{gather*}

Here are details extending the finite chessboard argument to this
form. In one direction transport all slab profiles to one reference
slab, using conjugated reflection for odd copies. For an arbitrary
filling of the $n$ slabs, $n$ even, reflection Cauchy--Schwarz bounds
the squared absolute partition integral by the two integrals obtained
by retaining either half and reflecting it. Normalize label $j$ by
$(Z_j^{\rm dis}+\epsilon)^{1/n}$, where $Z_j^{\rm dis}\ge0$ is its
fully disseminated integral and $\epsilon>0$. The normalized filling
of maximum absolute value exists. Both reflected fillings also
maximize whenever that maximum is positive: each is nonnegative,
bounded by the same maximum, and their product bounds its square.
A constant run in a maximizing filling can be doubled by retaining
the half starting at that run. Once its length reaches $n/2$, a
constant filling results. Every constant normalized filling is at
most one. Thus the maximum is at most one; the zero maximum is
trivial. Letting $\epsilon\downarrow0$ proves the one-direction bound
even when some disseminated integral is zero. Applying this argument
successively to the four coordinate directions proves the first
inequality in \eqref{eq:all-state-source-transfer-ST2}: coordinate reflections commute, and each
successive dissemination copies the preceding pattern literally in
the other coordinates. The resulting exponent for each original cell
is $1/N$.

Expanding the disseminated product for $F+G$ and applying the first
inequality term by term bounds it by
$\sum_{k=0}^N\binom Nk p(F)^kp(G)^{N-k}$. Its nonnegative $N$th root
proves the triangle inequality. Exactly half the copies are
conjugated, giving $p(zF)=|z|p(F)$; conjugating the entire nonnegative
disseminated integral proves $p(\bar F)=p(F)$. In particular
$p(\operatorname{Re}(e^{-i\phi}F))\le p(F)$.

\paragraph{Domination.}
For any periodic Wilson DLR state $\rho$, any finite collection of
occupied cells, and any sequence of compatible tori
$\ell_k\to\infty$,

\refstepcounter{equation}
\begin{gather*}
\FitMath{\left|\rho\left(\prod_C F_C\right)\right|
 \le\liminf_{k\to\infty}\prod_Cp_{l,\ell_k}(F_C).}\tag{\theequation}\label{eq:all-state-source-transfer-ST3}
\end{gather*}

Only finitely many factors differ from $1$.

First consider a bounded real function $H$ of one cell. Sum its mirror
copies over a cube of side $\ell$. Reflection invariance of $\rho$
makes the expectation of the sum equal to $N\rho(H)$. Conditional
Jensen bounds that expectation by the supremum, over exterior
configurations, of the logarithm of its conditional exponential
moment. Replace the exterior condition by periodic bonds. After
adjoining dummy independent Haar variables if necessary, the two
integrals use the same variables. Only boundary plaquettes change;
each changes the Hamiltonian by at most $2\beta$, and there are
$O((\ell/a)^3)$ such plaquettes. Source factors do not cross the outer
cell faces and therefore agree in the two integrals. It follows that

\refstepcounter{equation}
\begin{gather*}
\FitMath{\rho(H)\le\log p_{l,\ell}(e^H)
                  +C\beta l^4a^{-3}\ell^{-1}.}\tag{\theequation}\label{eq:all-state-source-transfer-ST4}
\end{gather*}

This estimate is uniform in the exterior configuration, not in $a$;
the volume limit is taken with $a$ fixed.

For bounded real $G$, put $H=\log(1+tG)$ with
$0<t<(2\|G\|_\infty)^{-1}$. The elementary inequalities
$\log(1+tG)\ge tG-Ct^2\|G\|_\infty^2$ and
$\log p(1+tG)\le t p(G)$ turn \eqref{eq:all-state-source-transfer-ST4}, first as
$\ell_k\to\infty$ and then as $t\downarrow0$, into
$\rho(G)\le\liminf_kp_{l,\ell_k}(G)$. For complex $F$ choose $\phi$
so that $\rho(\operatorname{Re}(e^{-i\phi}F))=|\rho(F)|$ and use
the last consequence of \eqref{eq:all-state-source-transfer-ST2}.

Finally enclose the finitely many occupied fine cells strictly inside
one larger coarse board cell. The one-cell inequality applies to
$G=\prod_CF_C$. On every compatible torus the fine-board chessboard
inequality gives
$p_{\rm coarse,\ell}(G)\le\prod_Cp_{l,\ell}(F_C)$: disseminate the
coarse cell, apply \eqref{eq:all-state-source-transfer-ST2} on the fine board, and take the coarse root.
Choose the larger cell as a sufficiently large allowed integral power
of the fine side, so that both boards are compatible along the chosen
torus sequence. This proves \eqref{eq:all-state-source-transfer-ST3}. The lower limit along
\emph{any} sequence is essential below.

The perturbation $\log(1+tG)$ in this proof is arbitrary and need not
be an admissible linear source. No renormalization-group bound was
applied to it or to its Gibbs states. It enters only an exact Jensen
and boundary-comparison argument and the elementary seminorm
inequality.

\subsubsection*{The precise ordinary torus hypothesis}

Here is one sufficient source hypothesis, stated to expose the actual
input to the transfer. Fix constants $r,A,C_\pm>0$, an entropy
rate $\kappa_0$, and $\kappa>\kappa_0$, uniformly over the selected
cutoff--coupling data $(a,\beta)$ and allowed tori. Let
$\|f\|_\sharp=\max(\|f\|_\infty,c_M\|\nabla f\|_\infty)$.
For every finite family of real Lipschitz torus functions require the
ordinary complex source bound, on
$\sum_i|z_i|\|f_i\|_\sharp<2r$,

\refstepcounter{equation}
\begin{gather*}
\FitMath{|Z_{a,\ell}(W)|\le
 \exp\left[V_{a,\ell}+\operatorname{Re}U_{a,\ell}(W)
       +\operatorname{Re}\sum_{s=0}^K\sum_D e_{a,\ell,s,D}(W)
       +C_+\ell^4\right],} \\[3pt]
\FitMath{Z_{a,\ell}(0)\ge e^{V_{a,\ell}-C_-\ell^4}.}\tag{\theequation}\label{eq:all-state-source-transfer-ST5}
\end{gather*}

The terms in this expression have the following meaning and bounds.
The source-independent vacuum number $V_{a,\ell}$ is real.

The point counterterm is
$U_{a,\ell}(W)=\sum_p u_{a,\ell,p}(W(x(p)))$, where each $u$ is
holomorphic and zero at zero. For each fixed cutoff its norm on
smaller discs is bounded uniformly in $p$ and $\ell$; this bound may
depend on the cutoff. The domain term $e_{a,\ell,s,D}$ is holomorphic
in the ordinary source variables, is zero at zero, depends only on
$W|_{D^+}$, and satisfies
$|e_{a,\ell,s,D}|\le Ae^{-\kappa m(D)}$. Put $a_s=L^{-s}$ and require
$\operatorname{diam}D^+\le C_0a_s(m(D)+1)$.
The domain entropy bound, in a convenient localized form, is

\refstepcounter{equation}
\begin{gather*}
\FitMath{\sum_{D:\,D^+\cap B\ne\varnothing,\,m(D)\ge t}
                     e^{-\kappa m(D)}
 \le C(R+1)^4a_s^{-4}e^{-ct}}\tag{\theequation}\label{eq:all-state-source-transfer-ST6}
\end{gather*}

for every radius-$R$ ball $B$ and $t\ge0$, uniformly in its position.
A count per unit cube bounded by $e^{\kappa_0m}$ implies this form
after absorbing a fixed polynomial in $m$ into a smaller entropy gap.
All domain sums used are normally convergent on smaller source balls.
Source representation independence is required: local terms depend on
the source function on their indicated region, not on the particular
list representing that function.

The condition in \eqref{eq:all-state-source-transfer-ST5} is an upper bound and a floor, not a claim that
$Z(W)$ is zero-free on the whole source ball. Its ordinary source
class includes the global real Lipschitz profiles constructed by
mirror periodization. This last inclusion matters. For
$w=\sum_i z_if_i$ in one cell, the periodized source is

\refstepcounter{equation}
\begin{gather*}
\FitMath{W^{\rm ch}=\sum_i z_iH_i^++\sum_i\bar z_iH_i^- .}\tag{\theequation}\label{eq:all-state-source-transfer-ST7}
\end{gather*}

The even and odd copy profiles $H_i^\pm$ are real. If each support is
at distance $b> a$ from the faces, the profiles vanish near the faces
and their Lipschitz norms are at most $\|f_i\|_\sharp$. Hence
$\|W^{\rm ch}\|_\flat\le2\sum_i|z_i|\|f_i\|_\sharp$, independently
of the number of cells. The ordinary radius $2r$ in \eqref{eq:all-state-source-transfer-ST5} therefore
supplies the radius $r$ used below. One first regards the two lists of
coefficients as independent complex variables; evaluating one at the
conjugate of the other does not require a new analytic extension.

\subsubsection*{Holomorphic cell normalizers from the ordinary hypothesis}

Suppose $m$ radius-$R$ support balls lie at distance at least a fixed
$b>1$ from the faces of a cell $C$ of side $l$. For each torus form a
normalizer $E_{C,\ell}(z)$ as follows. Retain in $U+\sum e$ only the
point terms and domains whose complete dependence region lies in one
open board cell. Evaluate each retained term at the source copy in
that cell, conjugate its value if that cell has odd parity, sum over
cells, and divide by $N$. The result is holomorphic in $z$, even
though \eqref{eq:all-state-source-transfer-ST7} itself involves $\bar z$: on an odd cell the retained
term has the form $g(\bar z)$, and
$\overline{g(\bar z)}$ is holomorphic. Its real part is unchanged by
this outer conjugation. Also $E_{C,\ell}(0)=0$.

Any discarded source-dependent domain meets a support ball and
reaches a cell face. Its diameter is at least $b$, so
$m(D)\ge b/(C_0a_s)-1$. Counting once per ball, \eqref{eq:all-state-source-transfer-ST6} shows that the
discarded contribution divided by $N$ is at most

\refstepcounter{equation}
\begin{gather*}
\FitMath{B_{\rm face}(m,R,b)
 \le CAm(R+1)^4\sum_{s\ge0}a_s^{-4}e^{-c'b/a_s}<\infty.}\tag{\theequation}\label{eq:all-state-source-transfer-ST8}
\end{gather*}

The sum is independent of the cutoff. Overcounting domains meeting
several balls only enlarges this upper bound. Point counterterms
whose source is nonzero lie strictly inside cells, so none are
discarded. Combining \eqref{eq:all-state-source-transfer-ST1}, \eqref{eq:all-state-source-transfer-ST5} and \eqref{eq:all-state-source-transfer-ST8} proves

\refstepcounter{equation}
\begin{gather*}
\FitMath{p_{l,\ell}(e^{\mathcal O_a(w)})
 \le\exp[\operatorname{Re}E_{C,\ell}(z)+B_C],} \\[3pt]
\FitMath{B_C=(C_++C_-)l^4+B_{\rm face}(m,R,b).}\tag{\theequation}\label{eq:all-state-source-transfer-ST9}
\end{gather*}


At a fixed cutoff the normalizers are locally bounded uniformly in
$\ell$. For the point terms this follows from their fixed-cutoff
bound and $O(l^4a^{-4})$ plaquettes per cell. For the domain terms use
\eqref{eq:all-state-source-transfer-ST6} with $t=0$ and the finite sum of $a_s^{-4}$ over $s\le K$.
These bounds may grow with $K$, which is harmless for the next step.
Montel's theorem therefore gives a subsequence on which
$E_{C,\ell}$ converges locally uniformly to a holomorphic $E_C$.
Explicitly, Cauchy estimates on smaller source balls give
equicontinuity there; Arzela--Ascoli and a diagonal selection over a
countable compact exhaustion give the subsequence, and Cauchy's
formula proves holomorphy of its limit.
For finitely many occupied cells take one common subsequence. It is
chosen from the torus normalizers alone, independently of every
infinite-volume state.

There is a second uniform bound that does not grow with $K$. If
$n\ge2$ sources in the same cell have pairwise support distance at
least $d$, a point counterterm has zero mixed derivative, because a
point cannot belong to two of the supports. A contributing domain
must meet all supports; in particular its diameter is at least $d$.
Cauchy's formula and \eqref{eq:all-state-source-transfer-ST6}, with the first support ball as anchor,
give

\refstepcounter{equation}
\begin{gather*}
\FitMath{|\partial_1\cdots\partial_nE_C(0)|
 \le D_n(d;R)\prod_i\|f_i\|_\sharp,} \\[3pt]
\FitMath{D_n(d;R)=CA(R+1)^4\max(1,d^{-4})(n/r)^n.}\tag{\theequation}\label{eq:all-state-source-transfer-ST10}
\end{gather*}

Indeed the remaining scale sum is
$\sum_sa_s^{-4}e^{-c'd/a_s}\le C\max(1,d^{-4})$.
This follows by splitting at $a_s\approx d$; the tail becomes a
convergent sum $Cd^{-4}\sum_{j\ge0}L^{4j}e^{-c''L^j}$.
The same estimate holds before the Montel limit, and hence after it.

\subsubsection*{All-state generating functions, cumulants and mean differences}

\begin{proposition}[Ordinary torus source estimates imply all-state bounds]
\label{thm:all-state-source-transfer}
Assume the ordinary torus source hypothesis in \eqref{eq:all-state-source-transfer-ST5}--\eqref{eq:all-state-source-transfer-ST6}, including
the stated locality, point counterterms and uniformity, and the proved
lattice symmetries. Then \eqref{eq:all-state-source-transfer-ST13}, \eqref{eq:all-state-source-transfer-ST14} and \eqref{eq:all-state-source-transfer-ST15} hold for every
periodic Wilson DLR state, uniformly in the cutoff, the states and the
support positions. Their constants depend only on the displayed
source constants, the number of entries, the support radius and,
for higher cumulants, the support separation.
\end{proposition}

Let $f_1,\ldots,f_n$ have radius-$R$ supports at arbitrary positions.
A zero test makes every corresponding multilinear assertion trivial,
so such tests may be omitted.
Choose an allowed side $l\le C_Ln(R+b+1)$ and a board offset that
puts each support ball at distance at least $b$ from every face.
For each coordinate, the forbidden offsets occupy at most
$2n(R+b+a)$ of a period; choosing $l$ larger than a fixed multiple of
this number leaves an allowed link-plane offset. Thus at most $n$
cells are occupied. Put $w_C=\sum_{i:\,\operatorname{supp}f_i\subset C}z_if_i$.
Apply \eqref{eq:all-state-source-transfer-ST3} to $F_C=e^{\mathcal O_a(w_C)}$ along the common Montel
subsequence in \eqref{eq:all-state-source-transfer-ST9}. For every periodic DLR state $\rho$ this gives

\refstepcounter{equation}
\begin{gather*}
\FitMath{\left|e^{-E(z)}\rho\left(e^{\sum_i z_i\mathcal O_a(f_i)}\right)\right|
 \le Q,} \\[3pt]
\FitMath{E=\sum_CE_C,} \\[3pt]
\FitMath{Q=e^{\sum_CB_C},} \\[3pt]
\FitMath{\sum_i|z_i|\|f_i\|_\sharp<r.}\tag{\theequation}\label{eq:all-state-source-transfer-ST11}
\end{gather*}

The same $E$ and $Q$ serve all states at the datum. The number $Q$
depends on $n,R$ and the fixed source constants, but not on the
cutoff, state, or support positions. This conclusion uses the lower
limit along the selected sequence, rather than an upper limit over
all tori.

Write the normalized holomorphic function in \eqref{eq:all-state-source-transfer-ST11} as $h_\rho$.
It has $h_\rho(0)=1$ and $|h_\rho|\le Q$. A logarithm is needed only
near zero, where it exists by continuity. For any nonempty subset
$B\subset\{1,\ldots,n\}$, Cauchy's formula, with all other variables
zero, gives

\begin{gather*}
\FitMath{|\partial_Bh_\rho(0)|
 \le Q(|B|/r)^{|B|}\prod_{i\in B}\|f_i\|_\sharp.}
\end{gather*}

The finite derivative identity for a logarithm now implies

\refstepcounter{equation}
\begin{gather*}
\FitMath{|\partial_1\cdots\partial_n\log h_\rho(0)|
 \le r^{-n}\!\sum_{\pi}(|\pi|-1)!Q^{|\pi|}
                 \prod_{B\in\pi}|B|^{|B|}
                 \prod_i\|f_i\|_\sharp.}\tag{\theequation}\label{eq:all-state-source-transfer-ST12}
\end{gather*}

The sum is over set partitions. No uniform zero-free polydisc is
needed for this identity at zero.

For $n\ge2$, the cumulant is the sum of the mixed derivatives of $E$
and $\log h_\rho$. If more than one cell is occupied, the mixed
derivative of $E=\sum_CE_C$ is zero because every summand lacks a
variable. If one cell is occupied, use \eqref{eq:all-state-source-transfer-ST10}. \eqref{eq:all-state-source-transfer-ST10}--\eqref{eq:all-state-source-transfer-ST12}
prove the all-state separated cumulant bound

\refstepcounter{equation}
\begin{gather*}
\FitMath{|\kappa_n^\rho(\mathcal O_a(f_1),\ldots,\mathcal O_a(f_n))|
 \le C_n(d;R)\prod_i\|f_i\|_\sharp .}\tag{\theequation}\label{eq:all-state-source-transfer-ST13}
\end{gather*}

Here $C_n$ is explicit as \eqref{eq:all-state-source-transfer-ST12}, with the uniform choice of $Q$ above,
plus $D_n$. This proves a finite constant for each $n$; no factorial
growth assertion beyond these formulas is being made.

For one radius-$R$ test choose one cell with $l\le C_L(R+b+1)$ and
the corresponding common normalizer $E_f$. Cauchy directly gives

\begin{gather*}
\FitMath{|\rho(\mathcal O_a(f))-E_f'(0)|
       =|h_\rho'(0)|\le(Q_R/r)\|f\|_\sharp.}
\end{gather*}

Consequently, for every pair of periodic states at the same datum,

\refstepcounter{equation}
\begin{gather*}
\FitMath{|\rho(\mathcal O_a(f))-\nu(\mathcal O_a(f))|
 \le(2Q_R/r)\|f\|_\sharp .}\tag{\theequation}\label{eq:all-state-source-transfer-ST14}
\end{gather*}

The common derivative $E_f'(0)$ may diverge with the cutoff. The
estimate controls its differences from the means and hence differences
between the means. The argument also covers separated reflected
copies: their support radii, distances and test norms are unchanged.

If desired the one-point constant can be chosen to grow only as
$(R+1)^4$. Choose a fixed smooth partition of unity by translates of
compactly supported functions of uniformly bounded derivative. Its
members have supports in balls of a fixed radius $R_0$, only
$C(R+1)^4$ meet a radius-$R$ support, and
$\|f\psi_j\|_\sharp\le C\|f\|_\sharp$. Apply \eqref{eq:all-state-source-transfer-ST14} to these pieces
with the fixed constant $Q_{R_0}$ and add. Such a partition is obtained
by dividing a nonnegative compact bump, positive on a unit cube, by
the positive periodic sum of its integer translates.

%
\paragraph{Short-distance differences between states.}
At a common first-passage datum let \(D_n\) be the
difference of the connected \(n\)-point functions
of two periodic states. For tests in a common
unit ball, \eqref{eq:all-state-source-transfer-ST11}
uses the same energy \(E\) in both states. The
source bound admits overlapping supports. Thus
the mixed derivative of \(E\) cancels, and
\eqref{eq:all-state-source-transfer-ST12} gives
\[
 |D_n(F_1,\ldots,F_n)|\le C_n\prod_i\|F_i\|_\sharp .
\]
Take \(a_K\le\mathsf d\) and put the supports of
the \(f_i\) in a ball of radius \(5\mathsf d\).
For small enough \(\mathsf d\), pack \(N\ge
c\mathsf d^{-4}\) disjoint lattice translates
of this ball inside a unit ball, with gaps
between the copies. For each argument \(i<n\)
use an independent family of uniform signs
\(\epsilon_v^i\) on the translates. In the
last argument use \(\epsilon_v^n=\prod_{i<n}
\epsilon_v^i\). Write \(F_i=\sum_v\epsilon_v^i
\tau_v f_i\). Since every summand vanishes
at the boundary of its own ball, both the
supremum and Lipschitz bounds of \(f_i\)
hold for \(F_i\). In the multilinear expansion
the sign average keeps exactly the terms
whose translations all coincide. Translation
invariance makes that average \(ND_n(f_1,\ldots,f_n)\).
For \(n=1\) take the ordinary sum with all signs one.
Therefore
\[
 |D_n(f_1,\ldots,f_n)|\le
 C_n\mathsf d^4\prod_i\|f_i\|_\sharp
 \le C_{n,\vartheta}\mathsf d^{4-n}
                      \prod_i\|f_i\|_\infty ,
\]
where the last bound uses
\(\operatorname{Lip}f_i\le\|f_i\|_\infty/(\vartheta\mathsf d)\).

Finally put $b_1=C_1(R)$ and $b_m=C_m(d;R)$ for $m\ge2$. The
moment--cumulant formula and invariance of higher cumulants under
constant shifts give

\refstepcounter{equation}
\begin{gather*}
\FitMath{\left|\rho\prod_{i=1}^k
     [\mathcal O_a(f_i)-\nu(\mathcal O_a(f_i))]\right|
 \le\left(\sum_{\pi\in\operatorname{Part}\{1,\ldots,k\}}
                     \prod_{B\in\pi}b_{|B|}\right)
                     \prod_i\|f_i\|_\sharp .}\tag{\theequation}\label{eq:all-state-source-transfer-ST15}
\end{gather*}

Singleton blocks use \eqref{eq:all-state-source-transfer-ST14}; the other blocks use \eqref{eq:all-state-source-transfer-ST13}. This is the
centered moment estimate required when a component is centered at a
different state or mixture, with the same uniformity in both states.

Thus \eqref{eq:all-state-source-transfer-ST5}--\eqref{eq:all-state-source-transfer-ST6}, source locality, the stated ordinary holomorphy,
and the proved lattice symmetries imply the bounds used under G4.
The transfer requires no weighted analytic extension. The analytic construction must still prove the actual
ordinary source bound, entropy gap, upper stability and floor. A
torus statement controlling only separated mixed derivatives of a
logarithm does not state the cell estimate in \eqref{eq:all-state-source-transfer-ST9} or a common one-point
normalizer. Those concrete conditions must be retained when replacing
an all-state input by the transfer theorem.


% END MODULE chapters/all_state_source_transfer_v3.tex

% BEGIN MODULE chapters/graded_source_bound_v4.tex
\subsection{Factorial moment bounds from rooted domains}
\label{ch:graded-rooted-source}

We strengthen the spatial counting used in the ordinary-source transfer.
The resulting estimates have factorial growth in the number of entries
and a separation singularity of order proportional to that number.
Only the ordinary finite-torus source estimates are used. In particular,
neither weighted forgetting nor source bounds for perturbed Gibbs states
are added.

Retain the ordinary source upper bound and floor in \eqref{eq:all-state-source-transfer-ST5}, the point-local
counterterms, and the proved chessboard domination in \eqref{eq:all-state-source-transfer-ST3}. Let $L\ge3$
be odd, $a=L^{-K}$, and $a_s=L^{-s}$, $0\le s\le K$. All constants
below are uniform over the chosen cutoff--coupling data and all large
allowed tori. The physical source norm is
$\|f\|_\sharp=\max(\|f\|_\infty,c_M\|\nabla f\|_\infty)$.

\subsubsection*{The stronger count and its geometric origin}

The localization parameter obeys $M=L^{m_0}$ for a fixed integer
$m_0\ge0$. Use the cofinal family of original tori
$\ell=2L^m$, $m\ge m_0$. At every actual scale,
$\ell/(Ma_s)=2L^{m-m_0+s}$ is an even integer, so all localization
partitions tile the torus; no scale terms are omitted. Every fixed
power-of-$L$ board also tiles all sufficiently large tori in this
family, with even cell counts. At scale $s$ use a cubical partition
of side $\rho_s=M a_s$. A localization domain $X$ is a union of these cubes. It admits
a finite connected cubical edge tree meeting every cube of $X$, and
$m(X)$ is the minimum number of edges of such a tree. The source
dependence region is contained in a fixed number of neighbouring cube
layers around $X$. This includes source values at plaquette barycentres
whose plaquette has a corner in a one-layer enlargement of $X$,
because a bare plaquette has diameter $O(a)\le O(a_s)$.

For a fixed cube $Q$ the elementary rooted entropy estimate is

\refstepcounter{equation}
\begin{gather*}
\FitMath{\#\{X:Q\subset X,\,m(X)=m\}\le C e^{\kappa_Dm}.}\tag{\theequation}\label{eq:graded-source-bound-GS1}
\end{gather*}

Here is a proof for the stated domain indexing. A minimizing tree with
$m$ edges has at most $m+1$ vertices. Every cube meeting an edge of
the cubical grid contains an endpoint of that edge; hence every cube
of $X$ contains a tree vertex. Choose a tree vertex at a corner of $Q$
in at most $16$ ways. Traversing each tree edge twice gives a closed
walk of length $2m$ from that corner. There are at most $8^{2m}$ such
walks. At most $16(m+1)$ cubes contain a vertex of the walk, and a
subset of these specifies $X$. Thus

\begin{gather*}
\FitMath{\#\le16\,8^{2m}2^{16(m+1)}.}
\end{gather*}

This covers all admissible domains; restricting the domain class can
only reduce the count. The same upper bound applies on the finite
cubical quotient.

Suppose $|e_{s,X}(w)|\le A e^{-\kappa m(X)}$, where
$\kappa>\kappa_D$. Summing \eqref{eq:graded-source-bound-GS1} gives a positive entropy gap $c>0$.
If the dependence region sees a radius-$u$ ball, some cube of $X$ lies
within $C\rho_s$ of that ball. There are at most
$C(1+u/a_s)^4$ choices of this root cube. Consequently

\refstepcounter{equation}
\begin{gather*}
\FitMath{\sum_{X:\,X^+\text{ sees }B,\,m(X)\ge t}
                       e^{-\kappa m(X)}
 \le C(1+u/a_s)^4e^{-ct}.}\tag{\theequation}\label{eq:graded-source-bound-GS2}
\end{gather*}

There is no factor $a_s^{-4}$ when $u\ll a_s$. This improvement uses
the fixed-layer dependence geometry and rooted indexing; a bound only
by $C(u+1)^4a_s^{-4}$ would not suffice for the argument below.
Also, summing over every possible root cube, with harmless overcounting,
gives the global estimate

\refstepcounter{equation}
\begin{gather*}
\FitMath{\sum_X e^{-\kappa m(X)}\le C\ell^4a_s^{-4}.}\tag{\theequation}\label{eq:graded-source-bound-GS3}
\end{gather*}

The tree representation implies
$\operatorname{diam}X^+\le C a_s(m(X)+1)$, as needed in these estimates.

\subsubsection*{A polynomial bound for one finely supported cell source}

Let a board have side $\eta\le1$ and one real test $v$ supported in a
ball of radius $\delta\le\eta$, lying at distance at least $\delta$
from the faces. Assume $\delta\ge16a$. Periodize $zv$ by reflected
copies, conjugating odd copies. Its ordinary source norm is at most
$2|z|\|v\|_\sharp$. Thus ordinary source radius $2r$ allows
$|z|<r/\|v\|_\sharp$, independently of the number of cells.

Construct the interior-cell holomorphic normalizer $E_{v,\ell}(z)$
as in \eqref{eq:all-state-source-transfer-ST9}. The discarded terms are the source-dependent domains
whose full dependence region is not contained in one cell. Such a
domain sees a copy of the support ball and reaches a face, so
$m(X)\ge c_0\delta/a_s-C_0$. Let $N=(\ell/\eta)^4$.
The discarded sum divided by $N$ has three scale ranges.

For $a_s\ge\eta$, discard even the source restriction and use \eqref{eq:graded-source-bound-GS3}.
The cost is at most $CA(\eta/a_s)^4$ per scale and per cell; its sum
is bounded by a geometric series. For
$\delta\le a_s<\eta$, \eqref{eq:graded-source-bound-GS2} gives a bounded cost per scale, and there
are at most $C(1+\log(\eta/\delta))$ such scales. For $a_s<\delta$,
\eqref{eq:graded-source-bound-GS2} and the forced domain size give a cost at most

\begin{gather*}
\FitMath{CA(1+\delta/a_s)^4e^{-c'\delta/a_s}.}
\end{gather*}

The sum is bounded uniformly in $\delta$ by comparison with
$\sum_{j\ge0}(1+L^j)^4e^{-c''L^j}$. All ranges may be truncated at
the actual cutoff, which only decreases the bounds. The point-local
counterterms are retained: their nonzero sources are strictly inside
cells. Their possibly large first derivatives cause no face error.

The stability and floor cost per cell is
$(C_++C_-)\eta^4\le C_++C_-$. We obtain

\refstepcounter{equation}
\begin{gather*}
\FitMath{p_{\eta,\ell}(e^{z\mathcal O_a(v)})
 \le \exp[\operatorname{Re}E_{v,\ell}(z)]\,
             Q_0(\eta/\delta)^p}\tag{\theequation}\label{eq:graded-source-bound-GS4}
\end{gather*}

for constants $Q_0\ge1$, $p\ge0$ independent of $a,\eta,\delta$ and
the support position. Fixed-cutoff local boundedness and Montel
selection supply a limiting holomorphic univariate normalizer $E_v$.
For any finite list of pieces use one common subsequence. It is chosen
from the torus data and is independent of the infinite-volume states.

\subsubsection*{Board offsets and a genuine product source polydisc}

Consider $n$ tests supported in radius-$R$ balls, with distinct supports
at distance at least $d>0$. Put $\bar d=\min(d,1)$, and choose an
allowed board side, a power of $L$, such that

\refstepcounter{equation}
\begin{gather*}
\FitMath{\frac{\bar d}{16L}<\eta\le\frac{\bar d}{16},} \\[3pt]
\FitMath{\delta=\frac{\eta}{64n}.}\tag{\theequation}\label{eq:graded-source-bound-GS5}
\end{gather*}

First suppose $a\le\delta/16$; thus this board side is a multiple of
the mesh. A fixed smooth partition of unity at mesh proportional to
$\delta$ splits each $f_i$ into pieces $v_{i,k}$ with

\refstepcounter{equation}
\begin{gather*}
\FitMath{\#\{k:v_{i,k}\ne0\}\le C(1+R/\delta)^4,} \\[3pt]
\FitMath{\operatorname{supp}v_{i,k}\subset B_{i,k}(\delta),} \\[3pt]
\FitMath{\|v_{i,k}\|_\sharp\le C\delta^{-1}\|f_i\|_\sharp.}\tag{\theequation}\label{eq:graded-source-bound-GS6}
\end{gather*}

For example, scale a smooth compact partition obtained from translates
of a bump divided by their positive periodic sum. Scaling its bounded
first derivatives gives the last inequality. A zero piece contributes
nothing and is omitted.

Choose one piece from each original test. In any coordinate the
offsets for which one enclosing radius-$\delta$ ball comes within
$\delta$ of a face occupy an interval of length at most $4\delta$
modulo $\eta$. On the discrete link-plane offsets, one such interval
contains at most $4\delta/a+2$ candidates. The $n$ intervals exclude
at most $n(4\delta/a+2)<\eta/a$ candidates by \eqref{eq:graded-source-bound-GS5} and the fine-mesh
assumption. There is an allowed offset in each coordinate. Hence all
chosen balls have the required margin. Different original tests
occupy different cells: a four-dimensional side-$\eta$ cube has
Euclidean diameter $2\eta<d$, whereas their nonempty supports are
$d$ apart.

Write these chosen pieces as $v_1,\ldots,v_n$ and their different cells
as $C_1,\ldots,C_n$. By \eqref{eq:graded-source-bound-GS4} each has a univariate normalizer $E_i$
and a bound

\refstepcounter{equation}
\begin{gather*}
\FitMath{Q_i\le Cn^p,} \\[3pt]
\FitMath{r_i=r/\|v_i\|_\sharp.}\tag{\theequation}\label{eq:graded-source-bound-GS7}
\end{gather*}

Product domination in \eqref{eq:all-state-source-transfer-ST3} is applied to the separate cell factors.
It therefore gives, for every subset $S$ and every periodic state
$\rho$,

\refstepcounter{equation}
\begin{gather*}
\FitMath{\left|\exp\left[-\sum_{i\in S}E_i(z_i)\right]
      \rho\left(e^{\sum_{i\in S}z_i\mathcal O_a(v_i)}\right)\right|
 \le\prod_{i\in S}Q_i} \\[3pt]
\FitMath{(|z_i|<r_i, i\in S).}\tag{\theequation}\label{eq:graded-source-bound-GS8}
\end{gather*}

This is a product polydisc; no bound on its combined source norm is needed. Each
chessboard factor in \eqref{eq:all-state-source-transfer-ST3} uses only one periodized piece and its own
ordinary source disc. Inactive factors are exactly $1$, with
$p(1)=1$, so \eqref{eq:graded-source-bound-GS8} contains the product over the active subset,
not the product over all $n$ indices. The same normalizers work for
every state, including a second centering state $\nu$.

\subsubsection*{Common centering and factorial order}

Define $X_i=\mathcal O_a(v_i)-E_i'(0)$.
Differentiating the normalized function in \eqref{eq:graded-source-bound-GS8} once in each active
variable, and using product Cauchy, gives

\refstepcounter{equation}
\begin{gather*}
\FitMath{|\rho\prod_{i\in S}X_i|\le\prod_{i\in S}(Q_i/r_i).}\tag{\theequation}\label{eq:graded-source-bound-GS9}
\end{gather*}

Only the first derivative of each univariate $E_i$ occurs; no bound on
its higher derivatives is required. Applying the one-variable case
to $\nu$ gives $|\nu(\mathcal O_a(v_i))-E_i'(0)|\le Q_i/r_i$.
Expanding the product centered at $\nu$ over its active subsets and
using \eqref{eq:graded-source-bound-GS9} therefore proves
\begin{equation}
 \left|\rho\prod_{i=1}^n
       [\mathcal O_a(v_i)-\nu(\mathcal O_a(v_i))]\right|
 \le2^n\prod_i(Q_i/r_i).
 \label{eq:graded-source-bound-GS10}
\end{equation}
This explains both the necessity of subset bounds in \eqref{eq:graded-source-bound-GS8} and the
role of a common normalizer. Uniform moments for auxiliary perturbed
Gibbs states are not being presumed.

The sum in \eqref{eq:graded-source-bound-GS10} over the piece choices from \eqref{eq:graded-source-bound-GS6}, using linearity of
centering. The board may depend on that finite piece choice; every
constant is independent of that choice. Since $\delta\le1$,
\begin{equation}
\FitMath{\begin{aligned}
 \left|\rho\prod_{i=1}^n
       [\mathcal O_a(f_i)-\nu(\mathcal O_a(f_i))]\right|
 &\le\big[C(R+1)^4n^p\delta^{-5}\big]^n
                           \prod_i\|f_i\|_\sharp\\
 &\le B_0(R)^n n^{(p+5)n}\bar d^{-5n}
                           \prod_i\|f_i\|_\sharp.
 \end{aligned}}
 \label{eq:graded-source-bound-GS11}
\end{equation}

If $a>\delta/16$, dispense with the partition and the board. A
radius-$R$ ball contains at most $C(1+R/a)^4$ plaquette barycentres.
Since $0\le a_p\le2$,

\begin{gather*}
\FitMath{\|\mathcal O_a(f)-\nu(\mathcal O_a(f))\|_\infty
 \le C(R+1)^4a^{-4}\|f\|_\sharp
 \le C(R+1)^4\delta^{-4}\|f\|_\sharp.}
\end{gather*}

Its product bound is absorbed by \eqref{eq:graded-source-bound-GS11}, because $\delta$ is
comparable to $\bar d/n$, $p\ge0$ and $\bar d\le1$.
This handles every cutoff, including the case when the prospective
side $\eta$ is smaller than the mesh.

Choose the fixed integer $k=\lceil p\rceil+5$. The elementary inequality
$n!\ge(n/e)^n$, obtained by integrating $\log x$, converts \eqref{eq:graded-source-bound-GS11} to

\refstepcounter{equation}
\begin{gather*}
\FitMath{\left|\rho\prod_{i=1}^n
       [\mathcal O_a(f_i)-\nu(\mathcal O_a(f_i))]\right|
 \le B(R)^n(n!)^k\max(1,d^{-5n})
                           \prod_i\|f_i\|_\sharp.}\tag{\theequation}\label{eq:graded-source-bound-GS12}
\end{gather*}


\begin{proposition}[Graded moments and cumulants]
\label{thm:graded-rooted-source}
Under the ordinary torus source hypotheses and the rooted domain
geometry stated above, \eqref{eq:graded-source-bound-GS12} holds for every pair of periodic Wilson
DLR states at the same datum. The same moment and cumulant bounds hold
for finite Wilson tori in the compatible family above a fixed minimum
side, with the integrating and centering state equal to that torus law. For $n\ge2$ their cumulants satisfy

\refstepcounter{equation}
\begin{gather*}
\FitMath{|\kappa_n^\rho(\mathcal O_a(f_1),\ldots,\mathcal O_a(f_n))|
 \le [2B(R)]^n(n!)^k\max(1,d^{-5n})
                              \prod_i\|f_i\|_\sharp.}\tag{\theequation}\label{eq:graded-source-bound-GS13}
\end{gather*}

For the finite-torus assertion, use the univariate normalizers
$E_{i,\ell}$ directly. The finite chessboard inequality gives \eqref{eq:graded-source-bound-GS8}
with those normalizers and the same $Q_i$ on the product polydisc.
The remaining estimates use the torus distance for the support
separation. They are finite and use no volume limit. The
chosen sub-unit power-of-\(L\) board tiles every compatible
torus above the fixed minimum side.

The constants are independent of the cutoff, the states and support
positions. The exponent $k$ is one fixed finite integer.
One may take $B(R)=C(R+1)^4$ with a constant fixed by the source and
domain parameters.
\end{proposition}

To verify the cumulant assertion without an unnecessary extra
factorial, center at $\rho$ and apply the moment--cumulant formula.
For a partition with block sizes $n_1,\ldots,n_m$, use \eqref{eq:graded-source-bound-GS12} on each
block and $\prod_jn_j!\le n!$. The remaining combinatorial bound is

\begin{gather*}
\FitMath{\sum_\pi(|\pi|-1)!\prod_{D\in\pi}(|D|!)^k
 \le(n!)^{k-1}\sum_\pi(|\pi|-1)!\prod_{D\in\pi}|D|!
 \le2^{n-1}(n!)^k.}
\end{gather*}

For exactly $m$ blocks, the middle sum equals
$\frac{n!}{m}\binom{n-1}{m-1}$: order the blocks, sum over positive
compositions of $n$, and cancel their factorial denominators. Summing
over $m$ proves the inequality and \eqref{eq:graded-source-bound-GS13}.

% Keep the closing statement and argument together.
\clearpage
The separation singularity in this theorem is $d^{-5n}$, not a single
$d^{-5}$ prefactor independent of $n$. This is a factorial-growth
estimate at each fixed separation, with diagonal order proportional
to $n$. It is sufficient for extensions whose derivative order may
grow linearly with the number of entries. A stronger separation law
requires an additional argument. Likewise the theorem does not prove
the ordinary analytic source and stability hypotheses themselves.

One may also recover a presentation with a separation prefactor independent
of $n$, at the cost of one factorial power. Put $x=\max(1,d^{-5})$.
The exponential series gives $x^n\le n!e^x$, so \eqref{eq:graded-source-bound-GS12} implies

\refstepcounter{equation}
\begin{gather*}
\FitMath{\left|\rho\prod_{i=1}^n
       [\mathcal O_a(f_i)-\nu(\mathcal O_a(f_i))]\right|
 \le e^{\max(1,d^{-5})}B(R)^n(n!)^{k+1}
                            \prod_i\|f_i\|_\sharp.}\tag{\theequation}\label{eq:graded-source-bound-GS14}
\end{gather*}

Thus $A(d;R)=e^{\max(1,d^{-5})}$ is one possible prefactor in that
notation. Its separation behavior is much worse than a polynomial.
The quantitative extension uses the sharper graded form in \eqref{eq:graded-source-bound-GS12}, rather
than attempting to infer diagonal order from this rewritten bound.


% END MODULE chapters/graded_source_bound_v4.tex

% BEGIN MODULE chapters/continuum_compactness_v5.tex
\subsection{Compactness of separated correlations}
\label{ch:continuum-compactness}

A sequence of cutoff data with uniform source estimates has a
subsequence whose separated correlations converge. We also prove
existence of Wilson Gibbs states at each fixed datum. Constructing
the cutoff data and their uniform estimates, comparing different
cutoffs, tuning, and proving uniqueness of continuum limits require
further arguments.

\subsubsection*{The fixed-cutoff Gibbs measure argument}

Fix a spacing $a>0$, a finite inverse coupling $\beta\ge0$, and the
compact group $G=\operatorname{SU}(2)$. Choose one orientation of each nearest
neighbor bond of $a\mathbb Z^4$ and denote this countable set by $E$.
A configuration is $U\in\mathcal X=G^E$; the opposite bond has matrix
$U_e^{-1}$. The product topology is compact and metrizable. For example,
enumerating the bonds gives a metric
$\sum_{r\ge1}2^{-r}\min(1,\|U_{e_r}-V_{e_r}\|)$.
Continuous functions depending on finitely many bonds are uniformly
dense in $C(\mathcal X)$. Indeed, replacing all coordinates after the
first $N$ by a fixed configuration changes the metric by at most
$2^{-N}$, and uniform continuity proves this approximation.

Let $A_p(U)=1-\frac12\operatorname{Re}\operatorname{Tr}U_p\in[0,2]$
be the Wilson plaquette action. For a finite bond window $W$, let
$\mathcal P(W)$ be the plaquettes meeting $W$ and let $\partial W$
be their bonds outside $W$. Define
\begin{equation}
 \gamma_W F(U)=\frac{\int_{G^W}F(VU)
       \exp[-\beta\sum_{p\in\mathcal P(W)}A_p(VU)]\,dV}
      {\int_{G^W}\exp[-\beta\sum_{p\in\mathcal P(W)}A_p(VU)]\,dV}.
 \label{cc:specification}
\end{equation}
Here $VU$ replaces the variables in $W$, and Haar measure is
normalized. The denominator lies between
$e^{-2\beta|\mathcal P(W)|}$ and one. The normalized density is
therefore a continuous function of the finite variables
$(V,U_{\partial W})$. If $F$ is continuous local, $\gamma_WF$ is
continuous local, depending only on the remaining variables of $F$
and on $\partial W$. This proves the Feller property directly;
integration of a uniformly continuous function over the compact
finite product preserves continuity. By uniform approximation it
also holds for $F\in C(\mathcal X)$.

Take Wilson probability measures on periodic tori with integer periods
$N\to\infty$, and lift them periodically to $\mathcal X$. The finite
partition function is finite and strictly positive because the
integrand is continuous and positive on a compact space. The lifted
measures have a weakly convergent subsequence. One direct compactness
argument is as follows. The space $C(\mathcal X)$ is separable: on a
compact metric space finite approximations by distance functions with
rational coefficients form a countable uniformly dense family.
Successively extract convergence of the integrals on this family.
Uniform approximation then defines a bounded positive functional on
all of $C(\mathcal X)$, with value one at the constant one. Its Riesz
probability measure is the weak limit. This argument simultaneously
gives convergence against every continuous function, not merely the
chosen countable family.

The lifted torus measure is not itself an infinite-volume DLR state:
remote periodic copies can determine an inside variable. The identity
needed for passage to the limit is a local one. Fix $W$ and a
continuous local $F$. For all sufficiently large periods, the finite
bonds involved in $W$, $\partial W$ and $F$ embed injectively into
the torus, and the plaquettes meeting $W$ are precisely their usual
torus counterparts. Integrating the variables in that torus window
first gives exactly
\begin{equation}
       \widetilde\mu_N(\gamma_WF)=\widetilde\mu_N(F).
 \label{cc:torus-local-DLR}
\end{equation}
Both integrands are continuous local. Their weak limits satisfy
$\mu(\gamma_WF)=\mu(F)$. Cylinder approximation and
$\|\gamma_WF\|_\infty\le\|F\|_\infty$ extend this identity to all
continuous $F$. Equivalently the Borel probability measures
$\mu\gamma_W$ and $\mu$ agree. The identity consequently holds for
every bounded Borel $F$.

It is also the conditional-law formulation of DLR. For bounded
outside-measurable $g$, the kernel is proper and satisfies
$\gamma_W(gF)=g\gamma_WF$. Hence
$\mu(gF)=\mu(g\gamma_WF)$, and $\gamma_WF$ is outside measurable.
Thus \eqref{cc:specification} is the conditional expectation given
the bonds outside $W$. This holds for every finite window.

The torus lifts are invariant under lattice translations and the
hypercubic group. These transformations preserve Haar measure and the
Wilson action, including when a chosen bond orientation is reversed.
Weak convergence preserves their invariance, so $\mu$ is an invariant
Wilson DLR state. In particular it is periodic, with the whole lattice
as a period group. The same Feller argument makes the class of
translation-invariant DLR states weakly closed, hence compact, and
its defining identities make it convex. No compactness of the union
over arbitrary period groups is being asserted.

If a source argument uses a prescribed compatible cofinal family of
tori, take that family in this proof. No measure estimate beyond
compactness is needed. Reflection positivity and the stronger symmetry
theorems for every periodic DLR state are separate results; this
measure argument proves neither their all-state estimates nor any
uniformity as $a\downarrow0$.

%
%
\subsubsection*{Recognition of limits on sufficiently large tori}
\hypertarget{growing-torus-recognition}{}

Fix \((K,\beta)\) and a finite bond window \(W\). Use the same
cofinal family of admissible tori for the periodic lifts and
for their limit points in \(\mathcal G^{\mathrm{tor}}_{K,\beta}\).
Write \(h_W\)
for the density in \eqref{cc:specification}. If the torus is large
enough for \(W\cup\partial W\) to embed, the marginal on \(W\)
of its periodic lift has Haar density
\[
 m_{\widetilde\mu_N}(V)
   =\int h_W(V\mid U_{\partial W})\,d\widetilde\mu_N(U).
\]
Indeed, for functions of \(W\) the identity
\eqref{cc:torus-local-DLR} needs only this embedding.
The same formula gives the marginal \(m_\mu\) of any fixed-cutoff
torus limit point. All these densities are bounded below by
\(e^{-2\beta|\mathcal P(W)|}\). The functions
\(h_W(V\mid\cdot)\) of the exterior variables form a compact
family in the uniform norm as \(V\) ranges over \(G^W\).
A finite uniform net and weak convergence thus show that
the densities converge uniformly along any torus subsequence
converging weakly to \(\mu\). The lower bound on the densities
gives uniform convergence of their ratios to one. In particular
\[
 \inf_{\mu\in\mathcal G^{\mathrm{tor}}_{K,\beta}}
    \left\|m_\mu/m_{\widetilde\mu_N}-1\right\|_\infty
       \longrightarrow 0\qquad(N\to\infty).
\]
Otherwise a sequence of violating tori would have a weak limit
point in this very class and the ratios along that subsequence
would tend to one.

For each datum \((K,\beta)\) take \(W\) to be the bond window
of physical radius \(K+1\) in the max norm, centered at the
origin, and choose a physical side threshold
\(\ell_\flat(K,\beta)\). Choose it so that every larger admissible
torus contains the window without wrapping and the preceding
infimum is at most \(e^{-1/a_K}\). We may also take the threshold
at least \(2L^K\). Choose a limit state \(\mu\) with ratio error
\(\epsilon\le 2e^{-1/a_K}\). This single state serves for all
orders and all fixed compact supports eventually contained
in the window. In fact, if \(Y_i=\mathcal O_K(f_i)\) and the
test supports lie in a box of radius \(R\) within the window
with a margin for their plaquettes, then
\[
 \left|\widetilde\mu_N\!\left(\prod_{i\in I}Y_i\right)
             -\mu\!\left(\prod_{i\in I}Y_i\right)\right|
 \le \epsilon\prod_{i\in I}
      C(1+R/a_K)^4\|f_i\|_\infty .
\]
The density ratio gives this inequality directly by integration.
Expanding cumulants in moments gives an error at order \(n\)
bounded by \(C_n\epsilon(1+R/a_K)^{4n}\prod_i\|f_i\|_\infty\).
Expansion over subsets gives the same estimate for the products
centered in their respective states. Along first-passage data
at fixed \(g\) with \(K\to\infty\), these errors tend to zero
for every fixed \(R\) and \(n\). Thus each separated-correlation
limit from tori above these thresholds has the same separated
cumulants of order \(n\ge2\) and the same self-centered functional
as a volume-first limit. Conversely, for a volume-first sequence,
at each cutoff choose an arbitrarily large torus from the
subsequence realizing the state at that cutoff so its density
ratio satisfies the same tolerance. The same comparison realizes
its continuum limit by tori above the thresholds.

%
%
\paragraph{The diagonal reflection form on these tori.}
Let \(\Theta F=\overline{F\circ\theta^*}\) for a diagonal
lattice reflection through the origin. The window \(W\) is
invariant under \(\theta\). For a bounded Borel \(F\) of
the bonds on its positive side, \(\Theta F\cdot F\)
is a function of \(W\). The fixed-cutoff state \(\mu\)
is reflection positive by \cref{thm:amrp-all-periodic}.
The same density ratio as above gives
\[
 \widetilde\mu_N(\Theta F\cdot F)
 \ \ge\ -\epsilon\,\widetilde\mu_N(|\Theta F|\,|F|)
 \ \ge\ -\epsilon\,\widetilde\mu_N(|F|^2).
\]
The first expression is real by reflection invariance of
the torus law. The second inequality uses Cauchy--Schwarz
and that invariance. Thus the same side threshold controls
the reflection form simultaneously for every such \(F\).
Along growing tori above the threshold, the error tends
to zero for \(F=F_K\) whenever
\(\epsilon\widetilde\mu_N(|F_K|^2)\to0\).
In particular this holds for fixed polynomials in finitely
many compactly supported lattice observables, centered in
the torus law if desired, by the plaquette count above.

\subsubsection*{The continuum hypothesis}

Fix an odd $L\ge3$ and a sequence $a_j=L^{-K_j}\downarrow0$.
At the same datum $(a_j,\beta_j)$ let $\rho_j,\nu_j$ be periodic
Wilson DLR states. The all-mirror lattice theorem supplies the following
properties used here: both states are invariant under unit lattice
translations and the hypercubic group, and $\rho_j$ is reflection
positive at coordinate and diagonal lattice mirrors. The centering
state $\nu_j$ is invariant at those mirrors as well. These properties
are lattice inputs to the compactness argument.

The infinite-lattice case includes volume-first limits when
the states are torus limit points. Direct growing-torus
limits are treated below with the finite-torus bounds.
The \hyperlink{growing-torus-recognition}{density comparison above}
identifies limits from tori with sides at least \(\ell_\flat(K,\beta)\)
with volume-first limits. For general growing tori and for removal
of volume after a fixed-torus continuum limit, diagonal reflection
positivity requires a separate argument. Tuned fixed-torus limits,
when obtained by matched-cutoff estimates, remain separate from the
subsequence theorem.

For $f\in C_c^1(\mathbb R^4)$ write
\begin{equation}
 \mathcal O_j(f)=\sum_p f(x_p)A_p,
 \qquad X_j(f)=\mathcal O_j(f)-\nu_j(\mathcal O_j(f)),
 \qquad m_{j,n}(\mathbf f)=\rho_j\prod_{i=1}^n X_j(f_i).
 \label{cc:moments}
\end{equation}
There is one plaquette for each coordinate plane and lattice site,
so six orientations are counted once each; $x_p$ is its barycenter.
There is no power of $a_j$ in the observable. At a fixed cutoff these
are bounded continuous random variables with finite bond dependence.
Set $m_{j,0}=1$ and complexify multilinearly.

\phantomsection\label{notdef:separated-tests}
Let $\mathfrak S_n$ be the tuples of compact $C^1$ tests with pairwise
disjoint supports. A zero entry gives value zero and can be set aside.
The separation is the minimum distance of the supports, not that of
balls containing them. Every support may lie in its own radius-$R$
ball; these balls may overlap. Use the fixed physical norm
\[
 \|f\|_\sharp=\max(\|f\|_\infty,c_M\|\nabla f\|_\infty),
                      \qquad c_M>0\hbox{ fixed}.
\]
The ordinary source hypothesis is used only through the following
proved implications of it. At each fixed order there is a bound
\begin{equation}
 |m_{j,n}(\mathbf f)|\le C_n(d;R)\prod_i\|f_i\|_\sharp
 \label{cc:local-bound}
\end{equation}
when the supports have distance at least $d>0$ and the stated radius
bound. Bounds holding eventually on each fixed support neighborhood
suffice for extraction. The one-point version is included: for
different centering and integrating states it is a uniform bound on
their mean difference, not on either unsubtracted mean.

For the quantitative conclusion assume the graded bound, with one
finite integer $k\ge1$ and $B(R)<\infty$ independent of $j$ and the
two states,
\begin{equation}
 |m_{j,n}(\mathbf f)|
  \le B(R)^n(n!)^k\max(1,d^{-5n})\prod_i\|f_i\|_\sharp.
 \label{cc:graded-bound}
\end{equation}
For $n=1$ take $d=1$. A harmless change of $B$ includes complex
entries. These are all-state bounds at the indicated data and also
apply to the pair $(\rho_j,\rho_j)$ used for canonical recentering.
The ordinary-source theorem obtains this bound from D1 and
its actual rooted domain geometry, including the fixed-layer
dependence and strict entropy gap. A weaker counting assertion alone
does not supply this hypothesis. Those source hypotheses, and
existence of the given sequence of cutoff data with common constants,
are not proved by this section.

%
%
\paragraph{Compactness on growing tori.}
For a direct sequence of compatible tori of sides
\(\ell_j\to\infty\), use its Wilson law for both integration
and centering. For any fixed finite collection of bounded
support neighborhoods in \(\mathbb R^4\), choose a box
containing them and a fixed margin. Once its diameter
is less than \(\ell_j/4\), the quotient map to the torus
is injective and preserves distances on the box. Periodize
each test by summing its translates by the periods.
The copies are disjoint, so the sup and Lipschitz norms
are preserved. The torus distances between the projected
supports are the original distances. The finite-torus
part of \cref{thm:graded-rooted-source} thus gives
\eqref{cc:local-bound} and \eqref{cc:graded-bound}
eventually on each fixed support neighborhood.

The countable extraction below only uses eventual
bounds on each such neighborhood, so the same diagonal
sequence works for the tori. Each graded inequality
passes to the limit on its fixed tuple. The extension
to diagonal-flat Schwartz tests is then made on that
limit functional. Translations, hypercubic symmetries
and coordinate mirror positivity pass with the same
fixed margins. For coordinate reflection positivity
on a torus, the projected supports lie inside one half
of the torus cut by the rounded mirror and its antipodal
mirror for all large \(j\). The limiting diagonal mirror
is a separate case. The half-period is an integer number
of lattice spacings, so the antipodal mirror has the same
site or link type as the chosen mirror. The fixed support
margin accommodates the full diameter of every sampled
plaquette. Diagonal positivity uses the additional
recognition or another diagonal reflection-positivity input.

%
%
\subsubsection*{A fixed admissible torus}

For \cref{thm:continuum-main}(a)(ii), fix \(T_\ell\) and
take any sequence of first-passage data with \(K_j\to\infty\).
The first bound in \cref{prop:source-domains-4-2} applies
directly to the connected functions. On this torus the factor
\(e^{nc_V\ell^4}\) is fixed at each \(n\). Use a countable
basis of rational coordinate boxes on the torus. Every tuple
of separated compact supports is contained in finite unions
of such boxes whose closures remain separated. In each
union choose the smooth cutoff exhaustion used below and
multiply by rational real Fourier polynomials on \(T_\ell\).
Smoothing and the \(C^1\) convergence of the Fourier
approximants give a countable dense family with support
in that union. Diagonal extraction on this list of tuples
and all \(n\ge2\) gives one subsequence. Telescoping with
the torus bound on each product of the fixed support
neighborhoods, as in \eqref{cc:telescoping}, extends
convergence to every separated \(C^1\) tuple.
Each canonically centered moment is the finite sum of
products of connected functions over partitions without
singletons. The same subsequence therefore gives all
the centered moments. The bound and permutation and
hypercubic symmetries pass to the limit. Round a torus
translation to the lattice modulo \(\ell\) and use
the same telescoping estimate on fixed neighborhoods
of its translated supports. This gives all torus
translations in the limit.

On a product of separated coordinate neighborhoods
take smooth cutoffs \(\chi_i\) supported there and equal
to one near the support projections of a joint test \(F\).
For \(k_i\in\mathbb Z^4\) write
\(e_{k_i}(x)=e^{2\pi\mathrm i k_i\cdot x/\ell}\).
Define the distribution of the connected limit there by
\[
 F\longmapsto
 \sum_{k\in\mathbb Z^{4n}}\widehat F(k)
            S_n(\chi_1 e_{k_1},\ldots,\chi_n e_{k_n}).
\]
Each entry has \(\|\chi_i e_{k_i}\|_\sharp\le
C(1+|k_i|)\). The torus bound thus costs at most
\(C(1+|k|)^n\), and \(5n+1\) derivatives of \(F\)
give summable Fourier decay after this cost in
dimension \(4n\). The series agrees with \(S_n\)
on separated smooth tensors. The same Fourier
approximation with cutoffs gives density of these
tensors on each product neighborhood. Thus the local
definitions agree on overlaps and give a unique
distribution of locally finite order off the diagonals.
The same argument applies to the centered moments,
using their finite partition formula. Coordinate mirror
positivity passes to these limits by the moving-test
argument below, using supports compactly contained
in one of the two slabs between the rounded link
mirror and its antipodal mirror.

%
If these fixed-torus limits are taken along
\(\ell_j\to\infty\) and then converge on separated
tuples in \(\mathbb R^4\), the resulting limit is
also realized by one growing-torus cutoff sequence.
List the countable separated tuples below, projecting
them to \(T_{\ell_j}\) once their support neighborhoods
embed. At \(\ell_j\), choose one index from its defining
cutoff sequence with \(K\ge j\) and with error at most
\(1/j\) from the fixed-torus limit on the first \(j\)
listed tuples whose neighborhoods embed. These are
finitely many requirements on a convergent sequence.
For any fixed separated tuple, choose an approximant
from this list with the same separated neighborhoods.
The per-ball bound in \cref{prop:source-domains-4-2}
holds eventually on this product with one constant for
the finite cutoffs and the inner limits. Telescoping at
both ends gives twice the approximation error plus
\(1/j\) between the selected finite-torus values and
the inner limits. First let \(j\) tend to infinity,
then let the approximation error tend to zero. The
same indices give the centered moments by the finite
partition formula. This argument uses the uniform
per-ball bound on the growing family.

\subsubsection*{A countable support-adapted family}

Let $\mathcal U$ be the bounded finite unions of open boxes with
rational endpoints. This is countable. Every finite collection of
separated compact sets $K_i$ has neighborhoods $U_i\in\mathcal U$
such that $K_i\subset U_i$ and the closures $\overline U_i$ are
pairwise separated. Cover each compact set by sufficiently small
rational boxes inside a neighborhood of radius less than one fourth
of its distance from the other sets, then take a finite subcover.
The neighborhoods can also stay in any prescribed open support
regions.

For each $U\in\mathcal U$ we now exhibit a countable family dense
in $C_c^1(U)$ for the $C^1$ norm, with every approximant supported
in $U$. Fix a nonnegative smooth unit-mass bump $\varrho$ in the unit ball.
For integers $m\ge1$ let
\[
 \chi_{U,m}=
  \mathbf1_{\{x:\operatorname{dist}(x,U^c)>1/(2m)\}}
                                      *\varrho_{1/(8m)}.
\]
It is smooth, compactly supported in $U$, takes values in $[0,1]$,
and equals one where $\operatorname{dist}(x,U^c)\ge1/m$.
Use finite rational linear combinations of
\begin{equation}
 \chi_{U,m}(x)\cos(2\pi q\cdot x/N),\qquad
 \chi_{U,m}(x)\sin(2\pi q\cdot x/N),
                   \quad q\in\mathbb Z^4,\ N\in\mathbb N,
 \label{cc:dense-functions}
\end{equation}
and take their rational span over all indices. Call this real
countable family $D_U$.

To verify density, fix $f\in C_c^1(U)$ and choose $m$ such that
$\chi_{U,m}=1$ on a neighborhood of its support. Extend $f$ by
zero, smooth it on a sufficiently small scale, and view the resulting
smooth compact function as periodic on a large integer-side box.
Smoothing converges to $f$ in $C^1$ because $f$ and its gradient are
uniformly continuous. The Fourier partial sums of this smooth
periodic function converge with their first derivatives: repeated
integration by parts makes the differentiated Fourier series
absolutely summable. Approximate their finitely many real sine and
cosine coefficients by rationals and multiply by $\chi_{U,m}$.
The multiplication is bounded in $C^1$, and $\chi_{U,m}f=f$.
This produces the required approximants. Complexification gives the
countable dense family $D_U+iD_U$ for complex tests.

List all tuples $(g_1,\ldots,g_n)$, for every $n\ge1$ and every
ordered list of pairwise separated $\overline U_i$, with
$g_i\in D_{U_i}+iD_{U_i}$. This is one countable list. By
\eqref{cc:local-bound} the values of each coordinate form a bounded
sequence; finite initial values do not matter if a bound starts
eventually. Extract an infinite subsequence for convergence of the
first coordinate, then a subsequence of it for the second, and so
on. If $I_r$ is the $r$th nested infinite index set, choose increasing
$j_r\in I_r$. For each fixed coordinate its values along $(j_r)$
eventually lie in the convergent subsequence for that coordinate.
All listed values therefore converge along this single subsequence.

\subsubsection*{Extension to every separated tuple}

The following telescoping estimate explains why that one extraction
is enough. If $\operatorname{supp}f_i\cup\operatorname{supp}g_i$
lies in $\overline U_i$, let $d_U>0$ be a lower bound for their
pairwise distances, taking $d_U=1$ when $n=1$, and choose a common
finite radius bound $R_U$. Multilinearity
and \eqref{cc:local-bound} give, uniformly in the relevant tail,
\begin{equation}
 \begin{split}
 |m_{j,n}(\mathbf f)-m_{j,n}(\mathbf g)|
 \le C_n(d_U;R_U)\sum_{i=1}^n
   \|f_i-g_i\|_\sharp
   \prod_{h\ne i}\max(\|f_h\|_\sharp,\|g_h\|_\sharp).
 \end{split}
 \label{cc:telescoping}
\end{equation}
Every mixed tuple in this telescoping sum is separated because the
neighborhoods, rather than just the two endpoint tuples, carry the
separation.

Given any $\mathbf f\in\mathfrak S_n$, choose such rational
neighborhoods and approximate each entry in $C^1$ by the corresponding
countable family. The right side of \eqref{cc:telescoping} tends
to zero uniformly in $j$. For a fixed approximating tuple its values
converge along $(j_r)$. Comparing two indices through this tuple
shows that $m_{j_r,n}(\mathbf f)$ is Cauchy. Define its limit by
$\sigma_n(\mathbf f)$. This definition is independent of all
neighborhoods and approximants because it is the limit of the
original sequence at that tuple. The subsequence has not been chosen
again for $\mathbf f$; convergence now holds on every separated
compact $C^1$ tuple simultaneously. Set $\sigma_0=1$.

Passing in \eqref{cc:telescoping} proves the same local continuity
estimate for $\sigma_n$. On compact subsets of each product of the
fixed support spaces, convergence is uniform: choose a finite
$C^1$ net in that compact set, apply pointwise convergence at its
centers, and use \eqref{cc:telescoping} at both ends. Its constants
are uniform on a norm-bounded set.

Multilinearity passes whenever all tuples in the proposed identity
are admissible. In particular the union of the alternative supports
in a single entry must stay separated from the other entries.
Permutation symmetry, conjugation and reality pass directly from
their exact lattice identities. The bounds in \eqref{cc:local-bound} and \eqref{cc:graded-bound} pass with their
original constants at every fixed tuple, since absolute-value
inequalities are closed under convergence. Their radius and
separation parameters concern the original supports, not the larger
neighborhoods used to construct the limit.

\subsubsection*{Moving arguments, real symmetries, and rounded mirrors}

Suppose $f_{i,j}\to f_i$ in $C^1$ and all their supports lie in
fixed pairwise separated compact neighborhoods. \eqref{cc:telescoping} and fixed-tuple convergence give
\begin{equation}
 m_{j_r,n}(f_{1,j_r},\ldots,f_{n,j_r})
                   \longrightarrow\sigma_n(f_1,\ldots,f_n).
 \label{cc:moving-tests}
\end{equation}
This is also valid for every finite combination of tuples.
For translations the needed $C^1$ convergence has the explicit
modulus
\begin{equation}
 \|\tau_v f-\tau_w f\|_\sharp
 \le\max\{\|\nabla f\|_\infty|v-w|,
                       c_M\omega_{\nabla f}(|v-w|)\}\longrightarrow0,
 \label{cc:translation-modulus}
\end{equation}
where $\omega_{\nabla f}$ is the uniform modulus of continuity of
the compactly supported continuous gradient. No second derivative
is required. The chain rule gives the corresponding convergence
for a convergent sequence of Euclidean isometries.

For any fixed real translation $v$ choose lattice vectors
$v_j\in a_j\mathbb Z^4$ with $v_j\to v$. Both states are invariant
under $v_j$, so $m_{j,n}(\tau_{v_j}\mathbf f)=m_{j,n}(\mathbf f)$.
The moving-test statement proves real translation invariance of
$\sigma_n$. Hypercubic invariance is exact at every spacing and
passes without rounding.

For clarity the reflection limit also works at arbitrary affine
positions. A coordinate mirror $x_\mu=c$ can be approximated by link
mirrors with $c_j\in a_j(\mathbb Z+1/2)$ and $c_j\to c$.
A diagonal mirror $x_\mu\pm x_\lambda=c$ is approximated by
$c_j\in a_j\mathbb Z$. In unit-normal coordinates its offset grid
is $a_j\mathbb Z/\sqrt2$. These are different conventions for
the same diagonal mirror; no factor $\sqrt2$ is discarded.
The reflected tests converge in $C^1$: the linear reflection part
is fixed and the translation parts differ by $2(c_j-c)e_\mu$
in the coordinate case, and by
$(c_j-c)(e_\mu\pm e_\lambda)$ in the diagonal case.

Let $F$ be a finite complex combination of separated tuples whose
entries lie strictly on one open side of the limiting mirror.
Its compact support has positive distance from the mirror. For all
fine cutoffs the barycenters sampled by those entries and all the
bonds of their plaquettes lie on the same side of the rounded
mirror: a plaquette has diameter $O(a_j)$, while that distance is
fixed. Form its centered lattice polynomial $F_j$.
Reflection invariance of $\nu_j$ makes centering commute with the
lattice reflection. The moment identity gives
\[
       \rho_j(\overline{\theta_jF_j}F_j)
              =m_j(\theta_j\overline F\cdot F).
\]
The left side is nonnegative where the lattice mirror
has exact positivity. For a diagonal mirror through the
origin on tori above the local density threshold, its
negative part tends to zero by the diagonal reflection-form
bound above. Every tuple on the right has a common positive support separation
for large $j$: internal separations are fixed and the two sides
are separated by the mirror gap. Apply \eqref{cc:moving-tests}
to get
\begin{equation}
              \sigma(\theta\overline F\cdot F)\ge0.
 \label{cc:reflection-limit}
\end{equation}
Reflection invariance itself passes in the same way.
The argument applies at coordinate mirrors for both lattice
constructions and at diagonal mirrors for periodic states.
For tori above the
\hyperlink{growing-torus-recognition}{local density threshold},
the diagonal conclusion at the origin extends to affine
translates by translation invariance of the limit. A diagonal
torus limit outside this threshold requires its own mirror estimate.

\subsubsection*{The one-point limit and the connected functions}

The first moment is determined more precisely by the plaquette
means. Unit translations and hypercubic symmetry make
\[
 \delta_j=\rho_j(A_p)-\nu_j(A_p)
\]
independent of the plaquette position and its coordinate plane.
Define
\begin{equation}
 R_j(f)=\frac{a_j^4}{6}\sum_p f(x_p),\qquad
                         c_j=6a_j^{-4}\delta_j.
 \label{cc:Riemann-mean}
\end{equation}
Then $m_{j,1}(f)=c_jR_j(f)$. Each orientation samples a translate
of $a_j\mathbb Z^4$. Comparing its sample value to the integral
on the corresponding side-$a_j$ cube gives
\begin{equation}
 R_j(f)\longrightarrow\int f,
 \qquad |R_j(f)-\textstyle\int f|
             \le C_R a_j\|\nabla f\|_\infty
 \label{cc:Riemann-limit}
\end{equation}
for a test supported in a fixed radius-$R$ ball. Indeed the sample
and any point of its cube are at distance at most $2a_j$; only
$C(R+1)^4a_j^{-4}$ cubes can contribute. Average the six such
estimates. The same bound is independent of the center of that
ball.

Choose one real smooth compact bump $f_*$ with $\int f_*=1$.
Its limiting moment exists by the preceding extraction, and
$R_j(f_*)\to1$. Consequently along the extracted subsequence
\begin{equation}
 c_j=\frac{m_{j,1}(f_*)}{R_j(f_*)}\longrightarrow c\in\mathbb R,
             \qquad \sigma_1(f)=c\int f.
 \label{cc:first-moment}
\end{equation}
In particular the coefficients are bounded on that tail. This
bound is uniform over both states:
$|c_j|\le2B(R_*)\|f_*\|_\sharp$ once $R_j(f_*)\ge1/2$,
and $|c|\le B(R_*)\|f_*\|_\sharp$.
This
argument uses the uniform one-point comparison and the common
plaquette means; it is not a deduction from formal continuum
translation invariance alone. If $\rho_j=\nu_j$, then $c=0$.

All connected limits are finite algebraic consequences, requiring
no common disc of analyticity for an exponential. For a fixed
$n$ work in the algebra with $z_i^2=0$ and let
\[
 M_j(z)=1+\sum_{\varnothing\ne A\subset[n]}
                    m_{j,|A|}(\mathbf f_A)\prod_{i\in A}z_i.
\]
The logarithm is the finite series for $\log(1+(M_j-1))$.
Its coefficient of $\prod_{i=1}^nz_i$ is
\begin{equation}
 K_{j,n}(\mathbf f)=
 \sum_{\pi\in\operatorname{Part}[n]}
       (-1)^{|\pi|-1}(|\pi|-1)!
                            \prod_{A\in\pi}m_{j,|A|}(\mathbf f_A).
 \label{cc:cumulant}
\end{equation}
In fact the term of degree $r$ in the logarithm counts the $r!$
orders of each partition with $r$ blocks and divides by $r$,
giving $(r-1)!$. Every subtuple is separated. Hence these finite
polynomials converge to the same formula with $\sigma$ in place
of $m_j$; denote it by $S_n$. Conversely the finite exponential
series gives
$\sigma_n=\sum_\pi\prod_{A\in\pi}S_{|A|}$, since its factor
$1/r!$ cancels the orders of the blocks.

Adding constants to the $X_j(f_i)$ multiplies $M_j$ by
$\exp(\sum_i b_i z_i)$, so changes its logarithm only in degree
one. Thus for $n\ge2$, $K_{j,n}$ is the usual cumulant of the
raw variables $\mathcal O_j(f_i)$ under $\rho_j$, independent of
the centering state. The graded cumulant bound passes as well:
\begin{equation}
 |S_n(\mathbf f)|\le[2B(R)]^n(n!)^k
                  \max(1,d^{-5n})\prod_i\|f_i\|_\sharp,
                         \qquad n\ge2.
 \label{cc:connected-bound}
\end{equation}
One may obtain it either from the uniform lattice cumulant bound,
or directly from \eqref{cc:cumulant}: if block sizes are $n_h$,
$\prod_h n_h!\le n!$, and
\[
 \sum_\pi(|\pi|-1)!\prod_{A\in\pi}(|A|!)^k
       \le(n!)^{k-1}\sum_\pi(|\pi|-1)!\prod_{A\in\pi}|A|!
       \le2^{n-1}(n!)^k.
\]
For exactly $r$ blocks the last unweighted sum is
$\frac{n!}{r}\binom{n-1}{r-1}$, by ordered compositions of $n$.

There is a canonical centered functional distinct from the original
two-state moments. Center at $\rho_j$ instead of $\nu_j$; expanding
each product and using convergence of $m_{j,1}$ gives along the
same subsequence
\begin{equation}
 \sigma^0_n(\mathbf f)=
  \sum_{A\subset[n]}\sigma_{|A|}(\mathbf f_A)
                  \prod_{i\notin A}\bigl(-c\textstyle\int f_i\bigr)
  =\sum_{\pi\in\operatorname{Part}_{\ge2}[n]}
                                      \prod_{A\in\pi}S_{|A|}.
 \label{cc:centered-functional}
\end{equation}
It has empty value one, singleton zero, the same higher connected
functions, and all the bounds and positivity properties above.
For the bounds use the lattice pair $(\rho_j,\rho_j)$ before
passing to the limit, so no uncontrolled constant from recentering
is introduced.

\subsubsection*{The quantitative flat extension and closed inequalities}

For completeness the topology retained by this compactness statement
is explicit. Let $\Delta_n=\bigcup_{i<h}\{x_i=x_h\}$ and let
$\mathcal S_{\Delta,n}$ be the Schwartz functions whose every jet
vanishes on $\Delta_n$. Use
\[
 p_s(\Phi)=\max_{|\alpha|\le s}\sup_x
          (1+|x|^2)^{s/2}|\partial^\alpha\Phi(x)|.
\]
On unit-radius supports, \eqref{cc:graded-bound} and
$\|f\|_\sharp\le\max(1,c_M)\max(\|f\|_\infty,
\|\nabla f\|_\infty)$ give the quantitative flat-extension
hypothesis with separation exponent five. For either limit
the theorem gives a unique extension $\Sigma_n$ with
\begin{equation}
 |\Sigma_n(\Phi)|\le C^n(n!)^{k+32}p_{20n}(\Phi).
 \label{cc:flat-bound}
\end{equation}
The same bound holds for the centered and connected extensions,
after changing $C$. Separated smooth tensors are dense in
$\mathcal S_{\Delta,n}$.

For the infinite-lattice inputs the same uniform estimate
holds for the extensions $\Sigma_{j,n}$. In the direct torus
construction the graded bound is eventual on each fixed compact
tuple. Distant supports can project onto the same part of a torus
at finite $j$, so that argument gives the flat estimate on the
limit through the fixed compact tuples.

For the infinite-lattice sequence, given $\Phi$ approximate it
by a finite sum $T$ of smooth separated tensors in $p_{20n}$.
The common estimate controls both $\Sigma_{j,n}(\Phi-T)$
and $\Sigma_n(\Phi-T)$ uniformly in $j$.
On $T$ convergence has already been proved. Taking first the cutoff
limit and then the approximation error to zero proves
\begin{equation}
                \Sigma_{j_r,n}(\Phi)\longrightarrow\Sigma_n(\Phi).
 \label{cc:flat-convergence}
\end{equation}
The positive-time version of the density theorem uses tensors
strictly above the mirror. Their reflected products are flat on
all cross diagonals, including the time boundary. Thus the limit's
symmetry, reality, real translations and reflection positivity extend by this
density and the continuity of tensor products in Schwartz
seminorms. For each mirror with the lattice positivity input, use
its corresponding positive half-space. It follows by an orthogonal change
of variables in the flat-space density construction; it does not
assume invariance of the functional under that change of variables.

More generally any closed inequality involving finitely many fixed
tuple evaluations and continuous functions of them passes to the
limit. The statement allows moving tests as in
\eqref{cc:moving-tests}, convergent scalar coefficients and
eventual uniform bounds. Its proof is simply convergence of those
finitely many numbers followed by continuity and closedness of the
inequality. It does not turn strict inequalities into uniform
strict bounds or provide estimates whose constants diverge with
the cutoff.

The proposition records the infinite-lattice construction,
including its diagonal mirror input. Direct growing tori give the
compact-tuple convergence, limiting flat extension and coordinate
reflection positivity by the preceding argument.

\begin{proposition}[Subsequential continuum compactness]
\label{thm:continuum-compactness}
For every input sequence satisfying the lattice properties and
two-state uniform bounds above, there is one subsequence whose
moments converge on every separated compact $C^1$ tuple. The
limiting functional is admissibly multilinear, symmetric and real,
invariant under real translations and the hypercubic group, and
reflection positive at every coordinate and diagonal affine mirror.
It retains the graded bound and its quantitative diagonal-flat
extension. Its first moment is $c\int f$, its higher connected
functions are the limits of the integrating states' cumulants,
and \eqref{cc:centered-functional} is their canonical centered
functional. Moving arguments and fixed-tuple closed inequalities
pass as stated above.
\end{proposition}

The preceding sections prove every assertion. The fixed-cutoff
measure argument ensures that invariant Wilson DLR states are
available at each given datum. It supplies neither a sequence of
admissible renormalization data with $K_j\to\infty$ nor D1's
uniform source bounds. It also supplies no identification of two
different extracted limits, no matched-cutoff convergence rate,
and no claim of a full Yang--Mills observable theory.


% END MODULE chapters/continuum_compactness_v5.tex

% BEGIN MODULE chapters/quantitative_flat_extension_v4.tex
\subsection{Extension to diagonal-flat Schwartz tests}
\label{ch:quantitative-flat-extension}

A graded bound for compact separated moments implies an extension
to diagonal-flat Schwartz tests, with explicit dependence on the number
of arguments. The graded bound for the gauge theory remains a hypothesis.
We construct the local distributions by Fourier series on product boxes
and track the constants as the number of variables grows.

\subsubsection{The hypothesis and the quantitative conclusion}

Use $\|f\|_1=\max(\|f\|_\infty,\|\nabla f\|_\infty)$ on
$C_c^1(\mathbb R^4;\mathbb C)$. Suppose $S_n$, $n\ge1$, are complex
multilinear forms on tuples with pairwise separated compact supports.
Multilinearity is required whenever the supports in the identity remain
separated. Fix numbers $1\le b<\infty$, $\kappa\ge0$ and $\alpha\ge0$,
independent of $n$, such that

\refstepcounter{equation}
\begin{gather*}
\FitMath{|S_n(f_1,\ldots,f_n)|
 \le b^n(n!)^\kappa\max(1,d^{-\alpha n})
                         \prod_i\|f_i\|_1}\tag{\theequation}\label{eq:quantitative-flat-extension-QF1}
\end{gather*}

if every support lies in a ball of radius one, placed anywhere, and different
supports have Euclidean distance at least $d>0$. For $n=1$ only the radius
and norm conditions are imposed. A centered first moment equal to zero is
included. Real forms may first be complexified, changing $b$ by an absolute
factor. Neither invariance nor positivity is needed for the extension.

\phantomsection\label{notdef:diagonal-flat-space}
Put $A=\lceil\alpha\rceil$, $k_*=A+15$, and
$\Delta_n=\bigcup_{i<j}\{x_i=x_j\}$.
Let $\mathcal S_{\Delta,n}$ be the Schwartz functions on $\mathbb R^{4n}$
whose derivatives of every order vanish on $\Delta_n$. For $n=1$ this is
the ordinary Schwartz space. Write

\refstepcounter{equation}
\begin{gather*}
\FitMath{p_k(\varphi)=\max_{|\beta|\le k}\sup_x
       (1+|x|^2)^{k/2}|\partial^\beta\varphi(x)|.}\tag{\theequation}\label{eq:quantitative-flat-extension-QF2}
\end{gather*}


\begin{proposition}[Quantitative diagonal-flat extension]
\label{thm:quantitative-flat-extension}
There is $C<\infty$, depending on $b,\alpha$ and fixed numerical cutoff
choices but not on $n$, and a unique continuous linear functional
$\Sigma_n$ on $\mathcal S_{\Delta,n}$ agreeing with $S_n$ on smooth
separated tensors. It satisfies

\refstepcounter{equation}
\begin{gather*}
\FitMath{|\Sigma_n(\varphi)|
       \le C^n(n!)^{\kappa+32}p_{k_*n}(\varphi),} \\[3pt]
\FitMath{n\ge1.}\tag{\theequation}\label{eq:quantitative-flat-extension-QF3}
\end{gather*}

Finite sums of smooth compact separated tensors are dense in this
diagonal-flat space in its Schwartz topology. In its positive-time
subspace, the tensors may be chosen with supports compactly contained in
$\{x_4>0\}$. These are constrained density assertions, not density in
the unrestricted Schwartz space across the diagonals.
\end{proposition}

The conclusion is an H1 growth estimate with linear derivative order and
a fixed factorial exponent; its constants are intentionally not optimized.
For example, $\alpha=5$ gives the sufficient order $20n$.
The probability-theory normalization $\Sigma_0(1)=1$ may be added
separately.

The radius-one formulation loses no information about a hypothesis stated
as
\[
B(R)^n(n!)^\kappa\max(1,d^{-\alpha n}),
\]
with constants independent of support positions. Choose a fixed smooth unit-scale partition in each
slot. A function supported in a radius-$R$ ball is a sum of at most
$C(1+R)^4$ supported pieces of radius at most one, each with $C^1$ norm
at most $C\|f\|_1$. Supports of pieces from different slots retain the
original separation. Multilinearity and \eqref{eq:quantitative-flat-extension-QF1} therefore imply the same
bound with

\refstepcounter{equation}
\begin{gather*}
\FitMath{B_{\rm eff}(R)\le Cb(1+R)^4.}\tag{\theequation}\label{eq:quantitative-flat-extension-QF4}
\end{gather*}

Thus only $B(1)$ is used, even if the supplied large-radius upper bound
was exponential. This assertion concerns a separate ball for each entry,
not one ball enclosing the entire configuration.

\subsubsection{Cutoffs with controlled derivatives}

We fix all cutoffs once. They can be chosen in Gevrey class two, with

\refstepcounter{equation}
\begin{gather*}
\FitMath{0\le\eta\le1,} \\[3pt]
\FitMath{|\partial^\beta\eta|\le C^{|\beta|+1}(|\beta|!)^2.}\tag{\theequation}\label{eq:quantitative-flat-extension-QF5}
\end{gather*}

Here is an explicit construction and a justification of this bound. Set
$a(t)=e^{-1/t}$ for $t>0$ and $a(t)=0$ for $t\le0$.
On the complex circle $|z-t|=t/2$ one has
$\operatorname{Re}(1/z)\ge2/(9t)$. Cauchy's formula bounds
$|a^{(m)}(t)|$ by $m!(2/t)^m e^{-2/(9t)}$.
Maximizing $t^{-m}e^{-c/t}$ and using $(m/e)^m\le m!$ gives
$C^{m+1}(m!)^2$. The step

\begin{gather*}
\FitMath{s(t)=\frac{a(t)}{a(t)+a(1-t)}}
\end{gather*}

is set equal to zero for $t\le0$ and one for $t\ge1$.
Its denominator on $[0,1]$ is at least $e^{-2}$.
The same derivative bound holds for its reciprocal and hence for $s$.
For precision, differentiate $D(1/D)=1$ inductively; after division by
$(m!)^2$, the convolution coefficients are
$1/\binom mr\le1$. If $|D^{(r)}|\le C_D^{r+1}(r!)^2$,
a derivative constant $C_D(1+2e^2C_D)$ for the reciprocal closes
the resulting geometric sum. Products
obey the same bound by Leibniz's formula. This proves \eqref{eq:quantitative-flat-extension-QF5}, including
flatness at the join points.

Choose an even $\eta_0$ equal to one on $[-1,1]$, zero outside
$[-2,2]$, and nonincreasing on $[0,\infty)$, using $s(2-|t|)$ in
the transition. Put $\psi(t)=\eta_0(2t)$ and
$D(t)=\sum_{z\in\mathbb Z}\psi(t-z)$. At most three summands occur
locally and $D\ge1$. Then $\chi_1=\psi/D$ is supported in $[-1,1]$,
has \eqref{eq:quantitative-flat-extension-QF5}, and $\sum_z\chi_1(t-z)=1$. Define

\refstepcounter{equation}
\begin{gather*}
\FitMath{\chi(u)=\prod_{a=1}^4\chi_1(u_a),} \\[3pt]
\FitMath{\chi^+(u)=\prod_{a=1}^4\eta_0(u_a),} \\[3pt]
\FitMath{\zeta(u)=1-\chi^+(u).}\tag{\theequation}\label{eq:quantitative-flat-extension-QF6}
\end{gather*}

Thus $\sum_{z\in\mathbb Z^4}\chi(u-z)=1$,
$\chi^+=1$ on the support of $\chi$, and $\zeta$ is zero for
$|u|_\infty\le1$ and one for $|u|_\infty\ge2$.
All these functions lie between zero and one and have \eqref{eq:quantitative-flat-extension-QF5}, with
an absolute constant. In particular, undifferentiated factors cost one.

We will repeatedly use the following elementary product estimate.
For $N$ factors of this kind at scale $h$, every derivative of order
$q$ is bounded by

\refstepcounter{equation}
\begin{gather*}
\FitMath{(CN)^q(q!)^2h^{-q}.}\tag{\theequation}\label{eq:quantitative-flat-extension-QF7}
\end{gather*}

An absolute prefactor is allowed at $q=0$. Indeed, distribute the
multiindex among the factors by Leibniz. The sum of the multinomial
coefficients is $N^q$, and the product of the factorials of the
assigned derivative orders is at most $q!$. Derivatives of pair
differences $x_i-x_j$ introduce only signs. This argument produces no
factor $C^N$ when most factors are undifferentiated.

\subsubsection{An explicit local functional on a product box}

First, tensor relations cause no ambiguity. In a finite relation between
separated tensors, let $d$ be smaller than every separation. Localize all
entries by one finite smooth partition on sets of diameter less than
$d/8$. A product of slot sets with two sets closer than $d/2$ has a
zero factor in every original tensor: otherwise two original supports
would be closer than $d$. On every remaining product the slot sets are
separated, and the form is an ordinary multilinear map on their supported
function spaces. It factors through their algebraic tensor product.
The latter embeds injectively in functions on the Cartesian product:
for the finite-dimensional function spans in any relation, finitely many
point evaluations give invertible matrices in each slot. A zero tensor
therefore has zero value. Summing the localization proves independence
of the tensor presentation.

Here is a quantitative local extension. Let $0<h\le1/16$,
$\ell=h/64$, and let slot centers be $c_i\in\mathbb R^4$.
Suppose the cubes $Q_i^+=c_i+[-2\ell,2\ell]^4$ have pairwise distance
at least $h/2$. They lie in balls of radius one.
For a smooth function $u$ supported in
$\prod_i(c_i+[-\ell,\ell]^4)$, take its Fourier series on the product
of period-$8\ell$ cubes centered at these points. Write its coefficients
as $c_{k_1,\ldots,k_n}$, $k_i\in\mathbb Z^4$. Multiplying the series
by $\prod_i\chi^+((x_i-c_i)/\ell)$ gives a series of separated tensors
equal to $u$ everywhere. There is no boundary ambiguity: $u$ vanishes
near the boundaries of the period cubes.
Explicitly, with $Q=\prod_i(c_i+[-4\ell,4\ell]^4)$, the coefficients are

\begin{gather*}
\FitMath{c_{k_1,\ldots,k_n}=(8\ell)^{-4n}
    \int_Q u(x)\exp\!\left[-\frac{i\pi}{4\ell}
                  \sum_i k_i\cdot(x_i-c_i)\right]dx.}
\end{gather*}


In the scaled variables $y_i=(x_i-c_i)/(4\ell)$, apply
$(1-\pi^{-2}\Delta_i)^4$ in every
four-dimensional slot. Integration by parts gives

\refstepcounter{equation}
\begin{gather*}
\FitMath{|c_{k_1,\ldots,k_n}|
 \le C^n J_{8n,\ell}(u)
                    \prod_i(1+|k_i|^2)^{-4},} \\[3pt]
\FitMath{J_{q,\ell}(u)=\max_{|\beta|\le q}
                       \ell^{|\beta|}\|\partial^\beta u\|_\infty.}\tag{\theequation}\label{eq:quantitative-flat-extension-QF8}
\end{gather*}

The sum of absolute coefficients in the product differential operator
is at most $C^n$; no factor depending unspecifiedly on dimension is
hidden here. The localized exponential in slot $i$ has $C^1$ norm at
most $C\ell^{-1}(1+|k_i|)$. Since

\begin{gather*}
\FitMath{\sum_{k\in\mathbb Z^4}(1+|k|)(1+|k|^2)^{-4}<\infty,}
\end{gather*}

\eqref{eq:quantitative-flat-extension-QF1} makes the series of its moment values absolutely convergent and
gives a local functional $T(u)$ with

\refstepcounter{equation}
\begin{gather*}
\FitMath{|T(u)|\le (Cb)^n(n!)^\kappa
                 h^{-(\alpha+1)n}J_{8n,\ell}(u).}\tag{\theequation}\label{eq:quantitative-flat-extension-QF9}
\end{gather*}

For a tensor supported in the inner cubes it agrees with $S_n$:
its Fourier coefficients factor slot by slot, the partial sums converge
in $C^1$, and \eqref{eq:quantitative-flat-extension-QF1} permits passage through the multilinear form.

These local extensions glue to a unique functional $T$ on
$C_c^\infty(\mathbb R^{4n}\setminus\Delta_n)$. To check this explicitly,
in an overlap choose smaller product boxes enclosing the compact support
under consideration. Fourier expansions with cutoffs inside the overlap
converge with every fixed number of derivatives and consist of separated
tensors. Two local functionals agreeing on those tensors therefore agree
there, by their finite-order bounds in \eqref{eq:quantitative-flat-extension-QF9}. A finite partition of unity on
a compact off-diagonal set now defines $T$ independently of the cover.
For original separated tensors of larger radius, first partition each
entry into sufficiently small pieces and use multilinearity. This also
shows that the glued $T$ agrees with the given form on all smooth compact
separated tensors, not just on the boxes used in its construction.
This gluing argument is used for consistency, not for the numerical bound:
the latter uses only the fixed explicit grid and \eqref{eq:quantitative-flat-extension-QF9}.

\subsubsection{Dyadic configuration layers and the factorial accounting}

Set $h_j=2^{-j}/16$ and define

\refstepcounter{equation}
\begin{gather*}
\FitMath{P_h(x)=\prod_{i<l}\zeta((x_i-x_l)/h),} \\[3pt]
\FitMath{D_0=P_{h_0},} \\[3pt]
\FitMath{D_j=P_{h_j}-P_{h_{j-1}}} \\[3pt]
\FitMath{(j\ge1).}\tag{\theequation}\label{eq:quantitative-flat-extension-QF10}
\end{gather*}

For $n=1$ the empty product is one and $D_j=0$ for $j\ge1$.
Off the diagonals $\sum_{j\ge0}D_j=1$. A fine layer has support

\refstepcounter{equation}
\begin{gather*}
\FitMath{h_j\le\min_{i<l}|x_i-x_l|_\infty\le4h_j.}\tag{\theequation}\label{eq:quantitative-flat-extension-QF11}
\end{gather*}

The statements about support also hold for its derivatives. Monotonicity
of the one-variable cutoffs makes $0\le D_j\le1$, although positivity
of these differences is not needed for the estimates.

At level $j$ set $\ell_j=h_j/64$ and use the product grid partition

\refstepcounter{equation}
\begin{gather*}
\FitMath{\chi_{j,z}(x)=\prod_{i=1}^n\chi(x_i/\ell_j-z_i),} \\[3pt]
\FitMath{z\in(\mathbb Z^4)^n.}\tag{\theequation}\label{eq:quantitative-flat-extension-QF12}
\end{gather*}

Retain a box if $\chi_{j,z}D_j$ is not identically zero. At one point
of that box every pair has distance greater than $h_j$; its centers
therefore have pair distance at least $h_j-2\ell_j$.
The expanded cubes $\ell_jz_i+[-2\ell_j,2\ell_j]^4$ have pair
distance at least $h_j-6\ell_j>h_j/2$. Thus \eqref{eq:quantitative-flat-extension-QF9} applies to
$u_{j,z}=\chi_{j,z}D_j\varphi$. A box with only boundary contact can
be ignored if its smooth localized function is zero.

There are at most $Cn^2$ pair and grid factors in such a cutoff.
Using \eqref{eq:quantitative-flat-extension-QF7}, Leibniz's formula and $q=8n$, we obtain the uniform
scaled-derivative cost

\refstepcounter{equation}
\begin{gather*}
\FitMath{(Cn^2)^{8n}((8n)!)^2\le C^n(n!)^{32}.}\tag{\theequation}\label{eq:quantitative-flat-extension-QF13}
\end{gather*}

Here $n^{16n}\le e^{16n}(n!)^{16}$, while
$(8n)!\le8^{8n}(n!)^8$. This accounts for every dependence on the
number of pair cutoffs; the $\binom n2$ undifferentiated factors do
not create $C^{n^2}$.

On a fine layer choose a pair with $|x_i-x_l|_\infty\le4h_j$ and
replace $x_i$ by $x_l$. The resulting point $y$ lies on a diagonal
and $|x-y|\le8h_j$. Every derivative of $\varphi$ vanishes there.
Taylor's integral remainder along this segment gives, for $|\beta|\le8n$
and any positive integer $m$,

\refstepcounter{equation}
\begin{gather*}
\FitMath{|\partial^\beta\varphi(x)|
 \le\frac{(Ch_j)^m}{m!}
     C^{k_*n}(1+|x|^2)^{-k_*n/2}
                    p_{k_*n}(\varphi)}\tag{\theequation}\label{eq:quantitative-flat-extension-QF14}
\end{gather*}

when $8n+m\le k_*n$. The derivative of order $m$ along a vector
uses at most four coordinate directions; its multinomial sum is
bounded by the $m$-th power of that vector's $\ell^1$ norm.
Also $|x-y|\le1/2$, so the weight at every point of the segment is
comparable with that at $x$, at cost $C^{k_*n}$. These are exponential,
not factorial, costs in $n$.

Take

\refstepcounter{equation}
\begin{gather*}
\FitMath{m=(A+5)n+2,} \\[3pt]
\FitMath{8n+m=(A+13)n+2\le(A+15)n.}\tag{\theequation}\label{eq:quantitative-flat-extension-QF15}
\end{gather*}

In differentiating $u_{j,z}$, a derivative on a cutoff costs $h_j^{-1}$,
which its scaled-derivative factor in $J_{8n,\ell_j}$ cancels.
Derivatives remaining on $\varphi$ are bounded by \eqref{eq:quantitative-flat-extension-QF14}.
Since $h_j\le1$, unused positive powers of $h_j$ can be discarded.
\eqref{eq:quantitative-flat-extension-QF7}, \eqref{eq:quantitative-flat-extension-QF13}--\eqref{eq:quantitative-flat-extension-QF15} consequently imply

\refstepcounter{equation}
\begin{gather*}
\FitMath{J_{8n,\ell_j}(u_{j,z})
 \le C_A^n(n!)^{32}h_j^m p_{k_*n}(\varphi)
       \sup_{x\in\operatorname{supp}\chi_{j,z}}
                         (1+|x|^2)^{-k_*n/2}.}\tag{\theequation}\label{eq:quantitative-flat-extension-QF16}
\end{gather*}

For $j=0$ use the same estimate without $h_j^m$, by the ordinary
Schwartz derivative bound; all powers of the fixed $h_0$ are absorbed
in $C_A^n$.

The grid sum has exactly the needed size. Since
$1+\sum_i|x_i|^2\ge\prod_i(1+|x_i|^2)^{1/n}$, and grid points in
one slot differ from its box by a fixed distance less than one,

\refstepcounter{equation}
\begin{gather*}
\FitMath{\begin{aligned}
 \sum_z\sup_{x\in\operatorname{supp}\chi_{j,z}}
                  (1+|x|^2)^{-k_*n/2}
 &\le C_A^n\prod_{i=1}^n
       \sum_{z_i\in\mathbb Z^4}(1+|\ell_jz_i|^2)^{-k_*/2}\\
 &\le C_A^n h_j^{-4n}.
 \end{aligned}}\tag{\theequation}\label{eq:quantitative-flat-extension-QF17}
\end{gather*}

The one-slot sum is bounded by $C_A\ell_j^{-4}$ by comparison with
the integral of $(1+|u|^2)^{-k_*/2}$ on $\mathbb R^4$; it is finite
because $k_*>4$. No number of configurations depending exponentially
on $n^2$ is being counted.

Combine \eqref{eq:quantitative-flat-extension-QF9}, \eqref{eq:quantitative-flat-extension-QF16} and \eqref{eq:quantitative-flat-extension-QF17}. Since $\alpha\le A$, the sum
of absolute values at every fine level is at most
\begin{equation}
\FitMath{C^n(n!)^{\kappa+32}p_{k_*n}(\varphi)
                 h_j^{m-(\alpha+5)n}
 \le C^n(n!)^{\kappa+32}p_{k_*n}(\varphi)h_j^2.}
 \label{eq:quantitative-flat-extension-QF18}
\end{equation}
The coarse level has the same bound without $h_j^2$.
The dyadic sum $\sum_jh_j^2$ is finite. We may therefore define

\refstepcounter{equation}
\begin{gather*}
\FitMath{\Sigma_n(\varphi)=\sum_{j\ge0}\sum_zT(\chi_{j,z}D_j\varphi).}\tag{\theequation}\label{eq:quantitative-flat-extension-QF19}
\end{gather*}

This absolutely convergent sum proves \eqref{eq:quantitative-flat-extension-QF3} and continuity in its stated
seminorm. It agrees with $T$ on compact off-diagonal tests, because
there only finitely many layers and boxes occur and their partition
sums to one. In particular it has the required tensor values.

\subsubsection{Density, the time boundary, and independence of choices}

We give the density details also needed for uniqueness. First cut off
infinity by $Q_R(x)=\prod_i\chi^+(x_i/R)$.
For a fixed Schwartz seminorm $p_k$, an extra $L$ powers of decay give

\refstepcounter{equation}
\begin{gather*}
\FitMath{p_k((1-Q_R)\varphi)
       \le C_{n,k,L}R^{-L}p_{k+L}(\varphi),} \\[3pt]
\FitMath{R\ge1.}\tag{\theequation}\label{eq:quantitative-flat-extension-QF20}
\end{gather*}

This follows by Leibniz, since the nonzero terms have at least one
$|x_i|\ge R$ and cutoff derivatives cost $R^{-1}$.
Next cut out all pairwise diagonals with $P_h$.
Where a derivative of $1-P_h$ can contribute, some pair is within
$2h$. Applying the same Taylor argument with arbitrarily many extra
derivatives gives

\refstepcounter{equation}
\begin{gather*}
\FitMath{p_k((1-P_h)\varphi)
       \le C_{n,k,M}h^{M-k}p_{k+M}(\varphi).}\tag{\theequation}\label{eq:quantitative-flat-extension-QF21}
\end{gather*}

For each fixed $n,k$ choose $M>k$. Thus compact tests supported away
from every diagonal are dense in $\mathcal S_{\Delta,n}$.
For example both displayed density constants may be bounded by
$C^{k+M+L+1}(n^2+1)^k(k!)^2$, with unused parameters omitted.
This follows from \eqref{eq:quantitative-flat-extension-QF7} and the same directional Taylor remainder;
these constants are not used in place of the uniform estimate in \eqref{eq:quantitative-flat-extension-QF18}.
On each such compact set, use a finite product-box partition and the
Fourier expansions of \eqref{eq:quantitative-flat-extension-QF8}. Their partial sums converge with every
derivative and have a common compact support. They consequently converge
in the Schwartz topology and are finite sums of separated smooth tensors.
This proves the asserted constrained tensor density.
Choosing the truncations and Fourier partial sums successively for the
first $1,2,\ldots$ seminorms gives a single convergent sequence.

For positive time define $\mathcal S_{\Delta,+,n}$ by requiring all
derivatives to vanish outside
$\{x_{1,4}>0,\ldots,x_{n,4}>0\}\setminus\Delta_n$.
In addition to the preceding cutoffs, set
\begin{equation}
\FitMath{B_h(x)=\prod_{i=1}^n s(x_{i,4}/h-1).}
 \label{eq:quantitative-flat-extension-QF22}
\end{equation}
It is zero if some time is at most $h$ and one if all are at least
$2h$. Where $1-B_h$ contributes to a positive-time test, a time is
within $2h$ of zero. Taylor expansion to that forbidden hyperplane,
where every jet vanishes, proves the analogue of \eqref{eq:quantitative-flat-extension-QF21}.
Hence $Q_RP_hB_h\varphi$ approximates every positive-time test in
the full Schwartz topology. Product boxes for this approximation can
be expanded while staying strictly above time zero. Their tensor
approximants are therefore strictly positive-time separated tensors.

There is also a uniform quantitative version: replace $P_h$ throughout
\eqref{eq:quantitative-flat-extension-QF10} by $P_hB_h$. A fine layer then has all pair distances and all
times at least $h$, while one pair distance or one time is at most
$4h$. The expanded grid boxes retain separation and positive time.
There are still at most $Cn^2$ cutoff factors. Taylor to the nearest
forbidden pair diagonal or time hyperplane supplies \eqref{eq:quantitative-flat-extension-QF14}, so exactly
the same $k_*$ and factorial exponent apply. This verifies explicitly
that the time-boundary constraint introduces no hidden growth in $n$.

Uniqueness now follows without a choice of a particular partition.
Any two continuous extensions agreeing on separated tensors agree on
a dense subspace and hence on all of $\mathcal S_{\Delta,n}$.
Thus the Fourier boxes, cutoffs, dyadic grids and orders of their
absolutely convergent sums do not change the resulting functional.
The same statement holds on the positive-time subspace.

If the compact moments are symmetric, real or invariant under a specified
Euclidean motion, their extensions inherit that identity by density and
continuity of the test-function transformations. If the compact moments
are reflection positive, take finite positive-time Schwartz sequences
and approximate their components by the tensors just constructed.
For a negative-time and positive-time factor, their tensor product is
flat on all cross diagonals: away from the time boundary one factor
already vanishes, and at time zero its jets vanish. Internal diagonals
are handled by their own factors. The tensor product is continuous
in Schwartz seminorms. Thus the reflected quadratic forms converge,
and their nonnegativity passes to these full sequence domains.
This proves the distribution-space implication once \eqref{eq:quantitative-flat-extension-QF1} is supplied;
it is not a proof of the gauge-theory estimate in \eqref{eq:quantitative-flat-extension-QF1}.

\paragraph{The graded-source limit and phase consequence.}
Under \cref{thm:graded-rooted-source}, every pointwise limit on
separated tuples of its two-state centered moments retains \eqref{eq:graded-source-bound-GS12}: for
each fixed tuple the inequality is closed and its constants are uniform
in cutoff and in both states. The phase-interface theorem realizes each
of its component forms by such sequences of pure lattice states, centered
at the specified reference states. Therefore the same bound holds for
those components, including their possibly nonzero first moments; their
one-point shifts are covered by the all-state comparison bound.
The fixed equivalence between $\|\cdot\|_\sharp$ and $\|\cdot\|_1$
is absorbed in the radius-one constant $b$. Thus \eqref{eq:quantitative-flat-extension-QF1} applies with
$\alpha=5$ and the graded factorial exponent $k$. These limits and
components consequently have the unique diagonal-flat extensions with
order $20n$ and factorial exponent $k+32$, uniformly over the stated
phase realization. This uses that realization rather than assuming an
unrelated component decomposition has the same bounds. It supplies
neither contact terms across the diagonals nor full rotational covariance.

\subsubsection{Why the precise hypotheses matter}

Position uniformity cannot be replaced by unrestricted exponential
growth in an enclosing radius. For example,
$S_1(f)=\int e^{|x|^2}f(x)\,dx$ has finite bounds for compact tests
in every fixed ball about the origin, with an exponential radius cost,
but has no continuous extension to Schwartz space agreeing with those
values. Translated nonnegative compact bumps, damped by
$e^{-|x|^2/2}$, converge to zero in every Schwartz seminorm while their
values diverge. Their supports can be chosen in successively more distant
unit balls. This example violates the position-uniform radius-one
hypothesis. Under that actual hypothesis in \eqref{eq:quantitative-flat-extension-QF4} rules out this obstruction.

Nor may a prefactor such as $e^{c n^2}$ be silently called factorial
growth of a fixed order. A lognormal random constant field has moments
$S_n(f_1,\ldots,f_n)=e^{n^2/2}\prod_i\int f_i$.
It is translation invariant and reflection positive, but violates every
estimate in \eqref{eq:quantitative-flat-extension-QF3} with fixed linear seminorm order and fixed factorial
exponent. To see this, take unit-integral Gevrey-two bumps in $n$
separated unit balls with centers at distance $O(n)$ from the origin.
Their tensor is diagonal-flat. For every fixed $k$, its seminorm
$p_{kn}$ grows at most $C^n(n!)^{C_k}$, by \eqref{eq:quantitative-flat-extension-QF5} and the polynomial
location weights. Its moment grows like $e^{n^2/2}$, which is larger
than $C^n(n!)^q$ for every fixed $q$. Thus the graded, fixed-exponent
assumption in \eqref{eq:quantitative-flat-extension-QF1} has mathematical content; reflection positivity
alone does not supply it.


% END MODULE chapters/quantitative_flat_extension_v4.tex

\section{Short distances and Euclidean covariance}\label{sec:short-distance}
The short-distance calculation identifies the free-curvature coefficient after running the coupling to the separation scale. It requires a conditional interacting law, including the exterior configuration and the differentiated normalizer. We give these requirements explicitly before the comparison and insertion estimates. The angular arguments and the reflection-positive Hilbert-space construction provide the covariance needed in \cref{sec:reconstruction}.

% BEGIN MODULE chapters/08_short_distance.tex
\subsection{The short-distance coefficient and nontriviality}
\label{candidate:short_distance}
\noindent The short-distance conclusions below retain the interacting
extraction and state-transfer inputs of Group F. In particular the exact
conditional identities of \cref{ch:clean-conditional-law} require
the actual complete-law identification, connected force and boundary
estimates, a common reference for conditional means, and marked bad-event
moments with rates relative to the $g^4$ signal. The classical Gaussian
cutoff theorem supplies a reference estimate; the interacting
inputs require separate proof. The assertions of \cref{prop:short-distance-7-1,prop:short-distance-7-3} are proof-outline
conclusions conditional on that remaining package.\par\medskip

\subsubsection{Two-point short-distance asymptotics, torus four-point positivity, and universality}

At short distances, the connected two-point function behaves to leading order like that of $N^2-1$ free fields. On the tori the same holds for the connected four-point function, restricted to nearly collinear quadruples, which form the class $\mathfrak T^{(4)}$ introduced in \cref{thm:continuum-main}(c)(5). The strength of these free fields equals the running coupling evaluated at the observation scale; it is non-zero and tends to zero only logarithmically.

\textbf{Objects.} The test pairs are $(f,h)\in\mathfrak{T}_{\mathbb{E}}(\mathsf{d};\vartheta)$, that is, pairs of non-negative Lipschitz bumps, each supported in a ball of radius $\mathsf{d}$, placed at a distance from one another that lies between $\mathsf{d}$ and $6\mathsf{d}$ (\cref{sec:short-distance-the-short-distance-two-point-law-and-its-asymptotic-freedom}). The free kernel is $\mathcal{K}_{\mathrm{free}}(f,h)=\iint f(x)h(y)|x-y|^{-8}dx\,dy$. We take the coupling $g(\mathsf{d})=g_{K-s(\mathsf{d})}(\beta)$, as well as the normalisation $\zeta_K(\mathsf{d})=\zeta'_{K-s(\mathsf{d})}(\beta)$, at the level whose spacing is smaller than $\mathsf{d}$ by a power of $\log g(\mathsf{d})^{-2}$.

\begin{proposition}[{non-triviality; \cref{thm:continuum-main}(c)(3)-(5)}]\label{prop:short-distance-7-1} Assume $g\le\gamma_W(\vartheta)$. Then there exist an integer $n_0\ge1$, a threshold $\mathsf{d}_0=\mathsf{d}_0(g,\vartheta)\in(0,1)$, a function $\bar\eta:(0,\mathsf d_0]\to[0,\frac13]$ that is non-decreasing with $\bar\eta(\mathsf d)\to0$ as $\mathsf d\to0$, and, for each $\ell\in\mathfrak L(g)$, a function $\bar\varepsilon_\ell:(0,\mathsf d_0]\to[0,\frac12]$ that is non-decreasing with $\bar\varepsilon_\ell(\mathsf d)\to0$ as $\mathsf d\to0$, with the following properties whenever $\mathsf{d}\le\mathsf{d}_0$. (i) Take any first-passage datum with $K-s(\mathsf{d})\ge n_0$, any $\ell\in\mathfrak{L}(g)$ and any pair in $\mathfrak{T}_{T_\ell}(\mathsf{d};\vartheta)$. Then

\begin{gather*}
\FitMath{\langle\mathcal{O}(f);\mathcal{O}(h)\rangle_{K,\ell,\beta}=\zeta_K(\mathsf{d})^2\,\tfrac{3(N^2-1)}{\pi^4}\,g(\mathsf{d})^4\,\mathcal{K}_{\mathrm{free}}(f,h)\,(1+\varepsilon),} \\[3pt]
\FitMath{|\varepsilon|\le\tfrac12,}
\end{gather*}

and moreover $|\varepsilon|\le\bar\varepsilon_\ell(\mathsf d)$, the inequality $|\tfrac{11N^2}{12\pi^2}g(\mathsf{d})^2\ln(1/\mathsf{d})-1|\le\bar\eta(\mathsf d)$ holds, and $|\zeta_K(\mathsf{d})-1|$ is bounded by the same right-hand side; all of these bounds are uniform in $(K,\beta)$. The same identity is also valid in every periodic Gibbs state at $\beta$ on $\mathbb{L}_K$, with the same $\zeta_K(\mathsf{d})$ and $g(\mathsf{d})$ and with an error satisfying $|\varepsilon|\le\tfrac12$ (\cref{sec:short-distance-the-short-distance-two-point-law-and-its-asymptotic-freedom}). (ii) On pairs of class $C^1$, every continuum limit obtained from \cref{thm:continuum-main}(a)(ii) on the torus $T_\ell$, as well as every one obtained from \cref{thm:continuum-main}(a)(iii) on $\mathbb{R}^4$, obeys

\begin{gather*}
\FitMath{S_2(f,h)=\tfrac{3(N^2-1)}{\pi^4}\Big(\tfrac{12\pi^2}{11N^2}\Big)^2\frac{\mathcal{K}_{\mathrm{free}}(f,h)}{(\ln(1/\mathsf{d}))^2}\,(1+\eta),} \\[3pt]
\FitMath{c\le1+\eta\le C,}
\end{gather*}

and on each $T_\ell$ one has $\eta\to0$ as $\mathsf{d}\to0$. \phantomsection\label{hyp:four-point}
(iii) On the tori, let $\vartheta\le\vartheta_\square$, the number from \cref{thm:continuum-main}(c)(5), which depends only on the fixed data. For the additional four-point clause, assume there exist an integer $n_4\ge n_0$, a number $\mathsf d_4\in(0,\mathsf d_0]$, a side $\ell_4(g)\in\mathfrak L(g)$, together with, for each $\ell\in\mathfrak L(g)$ with $\ell>\ell_4(g)$, a threshold $\mathsf d_4(\ell)\in(0,\mathsf d_4]$, with the following property. Suppose $K-s(\mathsf{d})\ge n_4$. Suppose $\mathsf d\le\mathsf d_4$ when $\ell\le\ell_4(g)$, and $\mathsf d\le\mathsf d_4(\ell)$ otherwise. Then for every quadruple in $\mathfrak{T}^{(4)}_{T_\ell}(\mathsf{d};\vartheta)$ (\cref{thm:continuum-main}(c)(5)), the following bounds hold uniformly in $(K,\beta)$. Here $c,C$ are constants that are allowed to depend on $g$ as well as on the shape $\vartheta$.

\begin{gather*}
\FitMath{c\,(N^2-1)\,g(\mathsf{d})^8\prod_i\|f_i\|_\infty\le S^T_{4;K,\ell,\beta}(f_1,\dots,f_4)\le C\,(N^2-1)\,g(\mathsf{d})^8\prod_i\|f_i\|_\infty.}
\end{gather*}

Clause (iii) is a statement about the lattice functions on the tori. For four points, \cref{thm:continuum-main} asserts nothing about the limits, and nothing about $\mathbb{R}^4$.
\end{proposition}

Additional comparison estimates are required for four-point positivity on arbitrary shapes or continuum components and for rigidity of general short-distance germs.

\paragraph{The short-distance two-point law and its asymptotic-freedom normalisation}
\label{sec:short-distance-the-short-distance-two-point-law-and-its-asymptotic-freedom}

\textbf{Objects.} The class $\mathfrak{T}_{\mathbb{E}}(\mathsf{d};\vartheta)$ consists of the pairs of non-zero Lipschitz functions $f,h\ge0$ satisfying

\begin{gather*}
\FitMath{\operatorname{supp}f\subset B(x,\mathsf{d}),} \\[3pt]
\FitMath{\operatorname{supp}h\subset B(y,\mathsf{d}),} \\[3pt]
\FitMath{\mathsf{d}\le\operatorname{dist}(\operatorname{supp}f,\operatorname{supp}h)\le6\mathsf{d},} \\[3pt]
\FitMath{\textstyle\int f\ge\vartheta\|f\|_\infty\mathsf{d}^4,} \\[3pt]
\FitMath{\operatorname{Lip}f\le\tfrac{\|f\|_\infty}{\vartheta\mathsf{d}},}
\end{gather*}

where the final two conditions are required of $h$ as well. It follows that $c\vartheta^2\|f\|_\infty\|h\|_\infty\le\mathcal{K}_{\mathrm{free}}(f,h)\le C\|f\|_\infty\|h\|_\infty$. On $\Lambda^2\mathbb{R}^4$ the free curvature propagator is given by

\begin{gather*}
\FitMath{\mathbb{K}(u)=\pi^{-2}|u|^{-4}\,\Lambda^2\mathsf{I}(\hat{u}),} \\[3pt]
\FitMath{\mathsf{I}(\hat{u})=1-2\hat{u}\hat{u}^{\mathsf{T}},} \\[3pt]
\FitMath{\hat u=u/|u|,}
\end{gather*}

which is the two-point tensor of $\mathrm{d}A$ when $A$ is a one-form with covariance $\delta_{\mu\nu}(4\pi^2|x|^2)^{-1}$. Since $\Lambda^2\mathsf{I}$ is an involution, $\operatorname{tr}\mathbb{K}(u)^2=6\pi^{-4}|u|^{-8}$. On $\Lambda_{K,\ell}$, let $F=\mathrm{d}A$ be the free lattice curvature, where $\mathrm{d}$ denotes the lattice coboundary; this curvature is the standard Gaussian on exact $2$-cochains with a single colour and unit coupling, and $\mathbb{E}_{\mathrm{free}}$ denotes the expectation with respect to it. Its background part at depth $s$ is $F_B:=\mathbb{E}_{\mathrm{free}}[F\mid Q_{K-s}(\mathbf{1})A\text{ modulo closed forms}]$, and we set

\begin{gather*}
\FitMath{C^{\mathrm{free}}_K=C^{\mathrm{bg}}_{K,s}+C^{\mathrm{fl}}_{K,s},} \\[3pt]
\FitMath{C^{\mathrm{bg}}_{K,s}(p,q)=\mathbb{E}_{\mathrm{free}}\,F_B(p)F_B(q),} \\[3pt]
\FitMath{\mathcal{K}_{K,s}(f,h)=\sum_{p,q}f(x(p))h(x(q))\,C^{\mathrm{bg}}_{K,s}(p,q)^2.}
\end{gather*}

Let $Q^{\mathrm{pl}}_j$ be the average of plaquettes, that is, of $2$-cochains, over blocks of order $j$; it is the counterpart of the straight average of $1$-forms. Because $\mathrm{d}\,Q_j(\mathbf{1})=\mathrm{d}\,Q^{\mathrm{st}}_j(\mathbf{1})=Q^{\mathrm{pl}}_j\mathrm{d}$, the conditioning above is conditioning on $Q^{\mathrm{pl}}_{K-s}F$. The $C^{\mathrm{bg}}_{K,K-j}$ are therefore kernels of orthogonal projections whose ranges are nested, and the slices

\begin{gather*}
\FitMath{\mathcal{C}_{(j)}:=C^{\mathrm{bg}}_{K,K-j+1}-C^{\mathrm{bg}}_{K,K-j}} \\[3pt]
\FitMath{(C^{\mathrm{bg}}_{K,K}=C^{\mathrm{free}}_K),} \\[3pt]
\FitMath{C^{\mathrm{fl}}_{K,s}=\sum_{1\le j\le K-s}\mathcal{C}_{(j)},}
\end{gather*}

are the covariances of mutually independent pieces of $F$. For the free field, the renormalisation group thus reduces to an orthogonal decomposition.

\textbf{Statement.} The error in \cref{prop:short-distance-7-1}(i) satisfies $|\varepsilon|\le\frac12$ and, in addition, $|\varepsilon|\le\bar\varepsilon_\ell(\mathsf{d})\to0$ on each torus of the family, uniformly in the cutoff, the bare coupling, the pair and its position. The same identity, with the same $\zeta_K(\mathsf{d})$ and $g(\mathsf{d})$ and with $|\varepsilon|\le\frac12$, holds for $\langle\mathcal{O}(f);\mathcal{O}(h)\rangle_\nu$ in every periodic Gibbs state $\nu$ at $\beta$ on $\mathbb{L}_K$. The constant arises as follows: $a_p\approx\frac12|F(p)|^2$; for each colour, Gaussian half-squares have covariance $\frac12(\text{covariance})^2$; and $\frac{N^2-1}2\cdot\frac6{\pi^4}=\frac{3(N^2-1)}{\pi^4}$.

 \textbf{Two variants.} (i) \emph{The one-loop coupling.} Define $\mathsf{g}(\mathsf{d}):=\big(g^{-2}+\frac{11N^2}{12\pi^2}\ln\frac1{\mathsf{d}}\big)^{-1/2}$, which depends on $\mathsf{d}$ and $g$ only. \Cref{lem:short-distance-7-2}, together with $0\le s\log L-\ln\frac1{\mathsf{d}}-\ln C_W(s)<\log L$, gives

\begin{gather*}
\FitMath{\big|g(\mathsf{d})^{-2}-\mathsf{g}(\mathsf{d})^{-2}\big|\le o(s)+C+\tfrac{c_0}{\log L}\ln C_W(s)=o(s(\mathsf{d})),}
\end{gather*}

uniformly in $(K,\beta)$. Hence $g(\mathsf{d})/\mathsf{g}(\mathsf{d})\to1$, and if $\mathsf{g}(\mathsf{d})$ is put in place of the running coupling $g(\mathsf{d})$ of the construction, the law still holds below a different threshold, with error functions that differ but are of the same type. (ii) \emph{Cutoffs off first passage.} Let $K'=K(\beta,g)$, and let $b=K-K'$, which may be positive, zero or negative. The dilation identity of \cref{sec:flow-units-common-prefixes-and-change-of-reference} then takes the form

\begin{gather*}
\FitMath{\langle\mathcal{O}(f);\mathcal{O}(h)\rangle_{K,\ell,\beta}=\langle\mathcal{O}(f^{[b]});\mathcal{O}(h^{[b]})\rangle_{K',L^b\ell,\beta},}
\end{gather*}

and the map $f\mapsto f^{[b]}$ sends $\mathfrak{T}_{T_\ell}(\mathsf{d};\vartheta)$ into $\mathfrak{T}_{T_{L^b\ell}}(L^b\mathsf{d};\vartheta)$ and leaves $\mathcal{K}_{\mathrm{free}}$ invariant. Consequently, whenever $L^b\ell\in\mathfrak L(g)$, $L^b\mathsf{d}\le\mathsf{d}_0$ and $K'-s(L^b\mathsf{d})\ge n_0$, the law at $(K',L^b\ell,\beta)$ with separation $L^b\mathsf{d}$ determines the behaviour of $\mu_{K,\ell,\beta}$ at separation $\mathsf{d}$. For $K$ large and under the conditions of clause (d)(1), \cref{sec:flow-units-common-prefixes-and-change-of-reference} and \cref{candidate:flow,ch:comparison-capsule-v7} give $|b|\le C(\varrho)$ for the tuned data and for each of their offsets $t\in\mathsf J$. When $b<0$, the identity is to be read in the direction from $(K',L^b\ell)$ to $(K,\ell)$.

\textbf{Proof outline.} (1) The local analysis is carried out on a single torus, the smallest one, $T_{\ell_\star(g)}$. On that torus the law holds by \cref{prop:short-distance-7-3} of \cref{sec:short-distance-isolating-a-single-scale-coupling-and-normalisation-at-the-scale}, and its kernel there is $\frac{N^2-1}2\mathcal{K}_{K,s}$. Reflection positivity then carries it from this torus to the remaining volumes and to the Gibbs states. For the Gibbs states this is the only route available, because the expansion deals with torus measures. The local analysis can also be carried out on the other tori. \Cref{thm:continuum-main}(c)(5) does this, with four sources, on every torus of the family, but only below a threshold on $\mathsf{d}$ which may decrease with $\ell$. The transfer, by contrast, yields a single threshold $\mathsf{d}_0$ common to all of them. (2) Comparison of lattice and continuum. Let $u$ be the shortest vector representing $x(p)-x(q)$ on $T_\ell$, with $|u|\approx\mathsf{d}$, let $\sigma_{pq}:=a_K^4\pi^{-2}|u|^{-4}$, and let $\mathbb{K}(\cdot)_{pq}$ be the entry belonging to the planes of $p,q$. Then

\begin{gather*}
\FitMath{\big|C^{\mathrm{free}}_K(p,q)-a_K^4\,\mathbb{K}(u)_{pq}\big|\le C\Big[\Big(\frac{a_K}{|u|}\Big)^{1/2}+\Big(\frac{|u|}\ell\Big)^4\Big]\sigma_{pq}} \\[3pt]
\FitMath{(10a_K\le|u|\le\tfrac1{10}\ell),} \\[3pt]
\FitMath{\big|C^{\mathrm{fl}}_{K,s}(p,q)\big|\le o(1)\,\sigma_{pq}.}
\end{gather*}

The first estimate uses the exponent and range proved in
\cref{ch:lattice-curvature-comparison}.  Its range is compatible with
the witness requirements.  The fixed room floor gives
$C_W(s)\ge C_{W,0}A_0^2$ for the admitted small couplings.  Increase
$n_0$ so that $10L^{-n_0}\le C_{W,0}A_0^2$.  Then
\[
 \frac{a_K}{\mathsf d}
 =\frac{L^{-(K-s)}}{L^s\mathsf d}
 \le\frac{L^{-n_0}}{C_W(s)}\le\frac1{10}.
\]
For two supports in the stated shape class, every short difference has
$\mathsf d\le |u|\le10\mathsf d$: the upper bound is the minimum
support distance plus the two support diameters.  Decrease
$\mathsf d_0$ so that $\mathsf d_0\le\ell_\star(g)/100$.
Consequently $10a_K\le|u|\le\ell/10$ on every admitted torus.
The increase of $n_0$ uses only constants fixed before $g$; the decrease
of $\mathsf d_0$ is made after $g$, as in the witness construction.

The first inequality compares second differences of the Green's function on the lattice torus with those of the continuum Green's function. Because $a_K/\mathsf{d}\le L^{-n_0}/C_W(s)$, the first term on its right is small by the choice of $n_0$, and it tends to $0$ with $\mathsf{d}$ uniformly in $K$, since the room increases. On the reference torus the second term is at most $C(\mathsf{d}/\ell_\star)^4$, which is small with $\mathsf{d}_0$ and tends to $0$ with $\mathsf{d}$. The second inequality is \cref{sec:short-distance-isolating-a-single-scale-coupling-and-normalisation-at-the-scale}(c). Approximating by a Riemann sum, at a relative cost of $Ca_K/(\vartheta^3\mathsf{d})$, yields $\mathcal{K}_{K,s}=6\pi^{-4}\mathcal{K}_{\mathrm{free}}(1+o(1))$. (3) The common-normalizer transfer estimate in \eqref{eq:two-point-extraction-partial-ET7}, under its stated comparison hypotheses, transfers the law to every $T_\ell$ and every $\nu$ at a cost of $C_2\mathsf{d}^2\|f\|_\infty\|h\|_\infty$, in which $C_2=C_2(g,\vartheta)$ denotes the constant of that lemma at $n=2$. Take $\mathsf{d}\le\mathsf{d}_0(g,\vartheta)$; there \cref{lem:short-distance-7-2} gives $\tfrac{11N^2}{12\pi^2}g(\mathsf{d})^2\ln(1/\mathsf{d})\ge\tfrac23$. In this range the cost is a fraction of the main term. That fraction goes to $0$ as $\mathsf{d}$ does, and so do the $o(1)$ of (2) and the $R$ of \cref{prop:short-distance-7-3}. The threshold $\mathsf{d}_0$ is chosen at the point where these three errors combined give $|\varepsilon|\le\tfrac12$. The main term is $\ge c\vartheta^2(\ln(1/\mathsf{d}))^{-2}\|f\|_\infty\|h\|_\infty$. Consequently, on every torus, the error goes to $0$ as $\mathsf{d}$ does.

The same convergence extends to periodic Gibbs states and to $\mathbb{R}^4$ provided the reference-comparison and common-normalizer hypotheses of \eqref{eq:two-point-extraction-partial-ET7} hold uniformly along the defining sequence.

(4) The quantities $\zeta_K(\mathsf{d})$ and $g(\mathsf{d})$ are controlled by \cref{lem:short-distance-7-2}.




\paragraph{Isolating a single scale: coupling and normalisation at the scale $\mathsf d$}
\label{sec:short-distance-isolating-a-single-scale-coupling-and-normalisation-at-the-scale}

\textbf{Objects.} We set $s=s(\mathsf{d})$ and $k_0:=K-s$, the witness level. We also write $x:=x_{k_0}(\beta)=g(\mathsf{d})^{-2}$. The separation measured in level-$k_0$ units is $\hat{\mathsf{d}}:=L^s\mathsf{d}\in[C_W(s),LC_W(s))$. This inclusion holds for $\mathsf d<1$, in which case $s(\mathsf d)\ge1$, and it follows from the minimality of $s(\mathsf d)$ and the monotonicity of $C_W$. We write $\underline x_s:=g^{-2}+\lambda_\sharp s-c_5$ and $\bar x_s:=g^{-2}+B_\sharp(s+1)$ for the ends of the envelope of \cref{sec:flow-the-run-its-coefficient-and-the-analytic-requirements}; they do not depend on $(K,\beta)$, and they are non-decreasing in $s$. The room $C_W(s)=C_{W,0}\,p_0(\bar{x}_s^{-1/2})^2$ contains a constant $C_{W,0}$, chosen large after $L$, $M$ and $\vartheta$ and before $g$. The lower bounds placed on this constant are $\vartheta\,C_W(s)\ge CM$, which is used in the Strategy, and those required in (b) and (c) below. The room tends to infinity, but $\log C_W(s)=o(s)$.

\begin{lemma}[{coupling and normalisation}]\label{lem:short-distance-7-2} Uniformly in $(K,\beta)$, and with no torus entering \cref{sec:flow-a-common-companion-and-the-torus-floor}, one has

\begin{gather*}
\FitMath{|x_{K-s}-x_K-c_0s|\le o(s),} \\[3pt]
\FitMath{\zeta'_k=1+\sum_{j\le k}b_j'(0)\in[\tfrac23,\tfrac43],} \\[3pt]
\FitMath{|\zeta'_{K-s}-1|\le Cg_{K-s}\,p_0(g_{K-s}).}\qedhere
\end{gather*}


\end{lemma}

The first bound is the flow statement of \cref{sec:flow-the-run-its-coefficient-and-the-analytic-requirements}. Because $s(\mathsf{d})\log L/\ln(1/\mathsf{d})\to1$, it yields $g(\mathsf{d})^{-2}/\ln(1/\mathsf{d})\to c_0/\log L=\frac{11N^2}{12\pi^2}$. For the other two bounds, note that a constant source $w\equiv\zeta$ acts as a shift of $\beta$ in which every auxiliary choice is held frozen. The increment $b_{j+1}(\zeta)$ is holomorphic for $|\zeta|<2\varrho$, and its slope is $\le C\big[g_j^3(\log x_j)^{\mathfrak{p}_0}+\varrho^{-1}g_j^{\kappa_0-4}\big]$ \cref{sec:sources-analytic-dependence-on-the-source-locality-and-decay-in-tame}. Since $x_j\ge x_k+\lambda_\sharp(k-j)-C$, the slopes summed over the levels $j<k$ are at most $C\big[g_k(\log x_k)^{\mathfrak{p}_0}+\varrho^{-1}g_k^{\kappa_0-6}\big]$. This bound is a function of $x_k$ alone, it tends to $0$, and for $\kappa_0\ge7$ it is $\le Cg_k\,p_0(g_k)$.

\begin{proposition}[{the law on the reference torus}]\label{prop:short-distance-7-3} Let $g\le\gamma_W(\vartheta)$, let $(K,\beta)$ be a first-passage datum, and suppose that $\mathsf d\le\mathsf d_0$, $K-s(\mathsf d)\ge n_0$ and $(f,h)\in\mathfrak T_{T_{\ell_\star(g)}}(\mathsf d;\vartheta)$. Then on the torus $T_{\ell_\star(g)}$, uniformly in $(K,\beta)$ and in the pair, one has

\begin{gather*}
\FitMath{\langle\mathcal{O}(f);\mathcal{O}(h)\rangle=\zeta_{k_0}'^{\,2}\,\tfrac{N^2-1}2\,x^{-2}\,\mathcal{K}_{K,s}(f,h)\,(1+R),} \\[3pt]
\FitMath{|R|\le\tfrac14,} \\[3pt]
\FitMath{R\to0\,(s\to\infty).}\qedhere
\end{gather*}



\end{proposition}

\textbf{Strategy.} In this proof the pair $(f,h)$ is also written $(f_1,f_2)$. We set $k_1:=k_0+\mathsf n$ with $\mathsf n:=\lceil2\log_L\bar x_s\rceil$. Here $g$ is fixed first. Then $\mathsf{d}_0$ is chosen, small enough that $\mathsf n\le s(\mathsf d)$ whenever $\mathsf d\le\mathsf d_0$. This gives $k_1\le K$.
Reduce \(\mathsf d_0\) further to fit the complete clean read boxes
and their padded neighborhood into a single lattice patch on the
reference torus. The geometry of the two localized readouts below
bounds the physical diameter by
\[
 C(\mathsf r+L\mathsf D+M\check R_{k_0})L^{-(s-\mathsf n)}.
\]
Here \(L^{\mathsf n}\le L\bar x_s^2\), and \(\bar x_s\) is affine in
\(s\) at fixed \(g\). The factors \(\mathsf r,\mathsf D,\check R_{k_0}\)
grow at most polynomially in \(1+s\), uniformly in the first-passage
datum. Thus the displayed diameter tends uniformly to zero as
\(s\to\infty\). Taking it below \(\ell_\star(g)/4\) gives the common
lift used in the box comparisons. Moreover $x_j\approx x$ for $k_0\le j\le k_1$, because $B_\sharp\mathsf n\ll x$. The sourced construction \cref{sec:sources-the-curvature-source-with-fixed-reference-data} is carried out up to a level $k\le k_1$, with the sources measured at depth $K-k$. At depth $s$ one has $\|f\|_{\sharp,s}=\|f\|_\infty$, because $\vartheta\,C_W(s)\ge CM$ for $\mathsf d\le\mathsf d_0(g,\vartheta)$. Above the witness level, $\|f\|_{\sharp,K-k}\le L^{k-k_0}\|f\|_\infty$, so the disc in $z$ becomes smaller. The identities for the derivatives at $z=0$ remain valid, and in the bounds each derivative in a source costs a factor $L^{k-k_0}$. One then has $Z(z)/Z(0)=\mathbb{E}_{\mathsf{P}^{(k)}}\,e^{\Phi^{(k)}(z)}$. In what follows we write $\mathsf{P}=\mathsf{P}^{(k_0)}$, which is a law distinct from the law $\mathsf{P}_{k_0}$ of \cref{prop:source-domains-4-4} that appears further on, and $\Phi=\Phi^{(k_0)}$. Everything stated about these two objects in this paragraph holds at each level $k\le k_1$, with $k$ in place of $k_0$ and $K-k$ in place of $s$. $\mathsf{P}$ is a probability law on triples. Each triple consists of a class $[h]$ of histories of level $k_0$ sharing one standing part $h^\flat$ \cref{sec:large-field-the-tests-the-nested-regions-and-the-content-of-a}, a block field $V\in\mathcal{V}_{k_0}$, and the stored fields. The law is determined by the unsourced sums $\tau_{[h]}:=\sum_{h'\in[h]}\tau_{h'}$. Here $\tau_{[h]}(dV)$ is an integral over the stored fields, which are the fields left unintegrated at level $k_0$. Its integrand is the exponential of the exponent of the sum, integrated against a positive measure composed of the gauge-fixing factors, the characteristic functions, the delta functions of the averaging constraints and Haar measures. Dividing this integrand by $Z(0)$ yields the law on triples. $\Phi$ is the difference between the sourced and the unsourced exponent of $\tau_{[h]}$. An exponent is attached to such a sum and not to an individual $\tau_h$ \cref{sec:large-field-the-tests-the-nested-regions-and-the-content-of-a}, \cref{sec:large-field-the-induction-closed-the-order-of-the-constants-the-sum}. The marginal of $\mathsf{P}$ in $([h],V)$ equals $Z(0)^{-1}\tau_{[h]}(dV)$. Tested on gauge-invariant functions, this marginal agrees with the push-forward under $h\mapsto[h]$ of the $(h,V)$-marginal of $\mathsf{P}_{k_0}$, the law on triples (bare field, history, block field) of \cref{prop:source-domains-4-4} \cref{sec:sources-the-curvature-source-with-fixed-reference-data}. Projected further onto $[h]$, the marginal is the image of $\tilde\nu$. Both laws therefore extend one and the same law of $[h]$ and of the gauge orbit of $V$, but they add different variables. $\mathsf{P}_{k_0}$ adds the bare field and $h$, the particular member of the class. $\mathsf{P}$ adds the stored fields, on which the part of $\Phi$ attached to the live components depends. With the notation $\Phi^{(k)}_B:=\prod_{i\in B}\partial_{z_i}\Phi^{(k)}|_{z=0}$ we have

\begin{gather*}
\FitMath{\langle\mathcal{O}(f_1);\mathcal{O}(f_2)\rangle=\langle\Phi^{(k)}_1;\Phi^{(k)}_2\rangle_{\mathsf{P}^{(k)}}+\mathbb{E}_{\mathsf{P}^{(k)}}\Phi^{(k)}_{12}} \\[3pt]
\FitMath{(1\le k\le k_1).}
\end{gather*}

All of these are expressions for one and the same quantity. At $k=k_0$ the first term on the right is the main term and the second term is an error ((a), (b)). At $k=k_1$ the roles are exchanged, and the law is obtained at that level ((c)). Both $\Phi^{(k)}_i$ and $\Phi^{(k)}_{12}$ are derivatives with respect to the sources. The argument uses conservation of mass alone, $\sum_h\tau^w_h(\mathcal{V}_{k})=Z(w)$. In other words, it uses exactness with the sources measured at depth $K-k$, which is the depth $s$ when $k=k_0$ (\cref{prop:source-domains-4-3} of \cref{sec:sources-the-curvature-source-with-fixed-reference-data}, together with \cref{sec:sources-analytic-dependence-on-the-source-locality-and-decay-in-tame}). This is what isolates the scale; no conclusion is drawn from $\mathsf{P}_{k_0}$, the coupling with the bare field. The law of $\mathcal{M}^{k_0}U$ is not identified with the marginal of $\mathsf{P}$ in $V$, because inside what $\mathbf{R}$ heals it replaces the law of the block field. Under $\mathsf{P}_{k_0}$, the equality $V=\mathcal{M}^{k_0}U$ is valid only on bonds at a large distance from the components that $\mathbf{R}$ has acted on. For this reason the sources are carried through the steps, and no observable is conditioned.

(a) \emph{Locality.} The source enters the exponent in three ways. \emph{Pointed pieces:} away from the live components, the coupling term is a sum of functions, each of which depends on a single value of $w$; this holds because the running shift at a plaquette is a finite sum of functions of this kind. As $f_1f_2=0$, these pieces give no mixed derivative. \emph{Regular local terms:} these are indexed by $j\le k_0$ and $X\in\mathcal{D}_j$, and they depend on $w$ only on $X^\sim$ \cref{sec:sources-analytic-dependence-on-the-source-locality-and-decay-in-tame}. Every domain that meets both supports satisfies $d_j(X)\ge c\,\hat{\mathsf{d}}/M$. By Cauchy bounds on the source disc, the contribution of these terms to $\Phi_{12}$ is bounded pointwise by $C\hat{\mathsf{d}}^4e^{-c\kappa\hat{\mathsf{d}}/M}\prod\|f_i\|_\infty$. This is of order $e^{-c\,p_0^2}$, where, in this step, $p_0$ without an argument stands for $p_0(g_{k_0})$. It is therefore smaller than any power of $g(\mathsf{d})$. \emph{The part attached to the live components:} only a cruder bound is used for this part. Consider a class $[h]$ in which a live component is in contact with a support, that is, lies within distance $c\,\mathsf d$ of it. The label is common to the class, and hence each live component, its lineage and the age of that lineage are common to the class as well; the members of the class differ only in healed components. Let $a(h)$ denote the largest lineage age among such components (\cref{prop:source-domains-4-14} of \cref{candidate:sources,ch:all-state-source-transfer}). On such a class the mixed derivative of this part is bounded by

\begin{gather*}
\FitMath{(\log x)^C\Big[e^{-c\,\hat{\mathsf{d}}}+L^{C\check{R}_{k_0}}L^{8a(h)}\Big]\prod\|f_i\|_\infty.}
\end{gather*}

It is the mean that is small: $|\mathbb{E}_{\mathsf{P}}\Phi_{12}|\le o(1)\,x^{-2}\prod\|f_i\|_\infty$. This follows from the age series of \cref{candidate:sources,ch:all-state-source-transfer} (\cref{prop:source-domains-4-14} and its proof), taken with rate $b\ge8$. The growth of the amplitude and the decay of the series are governed by one and the same age, namely the age of the lineage. The age recorded in the label restarts wherever $\mathbf{R}$ leaves a component standing as an island, and it is not larger. The number of cells of level $k_0$ in contact with the supports is $\le C\hat{\mathsf{d}}^4$, so

\begin{gather*}
\FitMath{\mathbb{E}_{\mathsf{P}}\Big[\mathbf{1}\{\text{contact}\}\,L^{8a(h)}\Big]\le C\hat{\mathsf{d}}^4\sum_{a\ge0}L^{8a}\max_{\mathfrak{c}}\tilde\nu\big(\mathfrak{lv}^{\ge a}_{k_0}(\mathfrak{c})\big)\le C\hat{\mathsf{d}}^4e^{-c\,p_0(g(\mathsf{d}))},}
\end{gather*}

and this bound overcomes $(\log x)^CL^{C\check{R}_{k_0}}x^2$ because $\mathfrak{p}_0\gg r$. \emph{At level $k_0$ the two-point function draws on this series twice}, both times at a rate $b\ge8$. At level $k_1$, in (c), it draws on the series once, at the same rate. The first use is the present one, the bound on the mean of the live part's mixed derivative. The second occurs in (b), in the covariances of the live first derivatives; there the product of two such derivatives carries $L^{4a}\cdot L^{4a}$. At level $k_0$ these two uses are the source of the weight $L^{8a}$ for two points. At level $k_1$ its source is (c1). Of the growth, only one property enters: as a function of the age, it is no faster than geometric. A different exponent, or an additional factor polynomial in $a$, is handled by a larger rate $b$, that is, by a larger $c_a$, a constant that may be increased freely. Hence the theorem on the ages (\cref{prop:source-domains-4-14}) underlies \cref{thm:continuum-main} (c)(3) and (c)(4). It also underlies (c)(5), for which \cref{thm:continuum-main}(c)(5) employs this series with a larger rate.

(b) \emph{The marginal part.} Here the coupling carries the source. Let $h^{(j)}_y\ge0$, with $\sum_yh^{(j)}_y=1$, be the tent functions of \cref{sec:sources-the-curvature-source-with-fixed-reference-data}, \cref{sec:sources-analytic-dependence-on-the-source-locality-and-decay-in-tame}. These weights have range $CMa_{K-j}$ and serve to localise the counterterm of step $j$. At a fixed level these are the functions $h_z$ in \hyperref[notdef:source-profiles]{the source profiles}. Let $p_y$ be a bare plaquette at $y$. At points where every profile $\varphi_j$, $j\le k_0$, takes the value $1$, that is, off the live components, the running shift satisfies, by its formula in \cref{sec:sources-analytic-dependence-on-the-source-locality-and-decay-in-tame},

\begin{gather*}
\FitMath{\partial_{z_i}\zeta_{k_0}(p;w)\big|_0=f_i(x(p))+\sum_{j\le k_0}b_j'(0)\sum_{y\in\mathbb{T}^{(j)}}h^{(j)}_y(p)\,f_i(x(p_y))=:\zeta'_{k_0}\tilde{f}_i(p),} \\[3pt]
\FitMath{|\tilde{f}_i-f_i|\le\tfrac{CM}{\vartheta\hat{\mathsf{d}}}\sum_j|b_j'(0)|\,\|f_i\|_\infty;}
\end{gather*}

the factor $\zeta_K(\mathsf{d})^2$ arises here. Differentiating $-A(x_{k_0}(\cdot)-\zeta_{k_0}(\cdot\,;w),U_{k_0})$ gives

\begin{gather*}
\FitMath{\Phi_i=\text{const}+\zeta'_{k_0}X_i+(\text{irrelevant})+(\text{live})+(\text{far}),} \\[3pt]
\FitMath{X_i(V):=\sum_p\tilde{f}_i(p)\,\chi_{k_0}(\square(p))\,a_p\big(U_{k_0,\square(p)}(V)\big).}
\end{gather*}

In this formula $\square(p)$ is the test cube containing $x(p)$, and $U_{k_0,\square(p)}$ is the background localised about that cube, hence near $\operatorname{supp}f_i$. The test $\chi_{k_0}(\square(p))$ acts as a window: a plaquette contributes only if its localised background passes the plaquette test. Consequently $X_i$ depends on $V$ only near the support and is the same for every history. It also satisfies $|X_i|\le Cx^{-1}p_0^2\,\hat{\mathsf{d}}^4\|f_i\|_\infty$. Away from the large-field regions, the global background differs from the localised one by an amount exponentially small in the radius of localisation; this difference is the term (far). The size of the irrelevant terms is $\le C(\log x)^{3\mathfrak{p}_0-q}x^{-1}\hat{\mathsf{d}}^4\|f_i\|_\infty$. The physical bound uses the source-family coefficient and a quadratic estimate at the real field size, together with the Wilson subtraction at each grouped core. The factor $x^{-1}$ cannot be dispensed with. The law uses $p_0^2\,\hat{\mathsf{d}}^8(\log x)^{3\mathfrak{p}_0-q}\to0$ as the coefficient of $x^{-2}$, and Cauchy bounds over the source radius $\varrho$ yield no power of $x$. The factor comes from the bound on the regular family produced by a step (\cref{sec:large-field-the-induction-closed-the-order-of-the-constants-the-sum}, ``Dependence of the regular family on the array''). That bound also holds for shifts of the coupling that vary with position, and it weights a difference of shifts by $CM^4\alpha_{0,k}^2$. We apply it to the sourced and the unsourced arrays of \cref{sec:sources-analytic-dependence-on-the-source-locality-and-decay-in-tame}. The remainders enter through the slope towards the constants carried by the induction of \cref{sec:sources-analytic-dependence-on-the-source-locality-and-decay-in-tame}, and they are small because of their size. The result is a linear recursion for the slopes in the source, and the small factor in front of the bound closes this recursion over the levels. Its solution is $\approx x_j^{-1}(\log x_j)^{\mathfrak{p}_0+q}$, and this takes the place of $\varrho^{-1}$. The oscillating part of the source, namely what remains after subtracting from the source its value at a single point of the domain, has a slope bounded in the same way, with an extra factor $L^{-(k_0-j)}$. The live terms become larger as the age increases. If a live component has $\operatorname{age}^{\mathrm{lab}}=a$, the fields stored in it are those of every level from $k_0-a$ onwards, together with those of the next finer level. For a component holding an island, the block of which retains the fields of its core \cref{sec:large-field-stored-variables-and-integration-in-a-large-field-patch}, the stored fields begin at the first level of its lineage. Over one cube of level $k_0$, the plaquette sums of these fields contain at most $L^{4(a+1)}$ terms. For the live terms, only their first two moments on the events of contact enter, so that, as in (a), it suffices that their growth be geometric in the age. Let $h$ be a history in which a live component is in contact with $\operatorname{supp}f_i$. On such a history the live contribution to $\Phi_i$ is bounded by $(\log x)^C\big[e^{-c\,\hat{\mathsf{d}}}+L^{C\check{R}_{k_0}}(a(h)+2)^CL^{4a(h)}\big]\|f_i\|_\infty$. The power of the age in this bound is the one of \cref{thm:continuum-main}(c)(5). In the bound $Ce^{-cp_0(g_j)}(2L^b)^{-a}$ of \cref{prop:source-domains-4-14}, the factor $2^{-a}$ absorbs this power, so the age series is used at the same rate. This live contribution has two covariances. One is with the live contribution of the other factor, with weight $L^{8a(h)}$. The other is with those functions of $V$ that occur in the other factor, with weight $L^{4a(h)}$, taken against their supremum bound. By the age series of \cref{candidate:sources,ch:all-state-source-transfer}, as in (a), both of these covariances are $o(1)\,x^{-2}\prod\|f_i\|_\infty$.

\paragraph*{From the two witness insertions to the physical sources.}
Let \(\lambda_{k_0}\) be the marginal in \(V\) of the unsourced law
\(\mathsf P\) on the class, \(V\) and the stored variables. Write
\(\mathsf r=\mathsf d/a_s\), \(x=x_{k_0}\), and set
\(u=x^{-2}\|f_1\|_\infty\|f_2\|_\infty\).
The identity obtained by differentiating conservation of mass at the
witness is
\[
 S^T_2(f_1,f_2)=\mathbb E_{\mathsf P}\Phi_{12}
                    +\operatorname{Cov}_{\mathsf P}(\Phi_1,\Phi_2).
\]
The law on the right is the full unsourced law. In particular, integrating
out the stored variables before differentiating would add the conditional
covariance of their two source derivatives to the mixed term.

Use the decomposition
\(\Phi_i=c_i+\zeta'_{k_0}X_i+I_i+B_i\) at \(k_0\).
Here
\[
 c_i=\left.\partial_{z_i}\sum_{j\le k_0}\sum_X
                                  e_j(X;w_z)\right|_{z=0}.
\]
This is the same number for every class. The difference between
the actual share of source energy in a class and this sum is
included in \(B_i\). Count the missing shares of the actual
resummed label in the level units of its retained regions.
A healed region already removed by resummation enters through
its refiled activity and the compensated label. It is not a
present lineage contact of the witness.
The coefficient of \(X_i\) is obtained by summing the single-point source
shifts with their own tents. The sourced exponent on this ordinary carrier uses the running
weight $x_{k_0}-\zeta_{k_0}$ and the regular members with their
tent counterterms already subtracted. The difference of those
counterterms is the accumulated shift $\zeta_{k_0}-w$ there.
One can first cancel it against the shift in the weight and
then extract the Wilson part from the unsubtracted regular
derivative. Equivalently keep $\zeta'_{k_0}$ in the readout
and differentiate the already subtracted regular members.
In either form the term
\(b'_j(0)f_i(x(p_y))A(h_y^{(j)},\cdot)\) is counted once. There is no analogous mixed derivative of a pointed piece:
it is a function of one source value and \(f_1f_2=0\).

For \(I_i\), apply the regular formation estimate
\eqref{eq:local-formation-LB14} and the window estimates
\eqref{eq:local-formation-LB17}--\eqref{eq:local-formation-LB18}
at \(n=k_0\). The real regular windows used here have complete
retained data of size \(C\delta_n\). In the interior of the padded
small-field set, use \eqref{mb:real-layer-profile} and the local
gauge estimate on the full stencil of the localized background.
The input is the passed tests on the retained datum with its
regular background profile, rather than just the curvature cutoff
in \(X_i\). The fixed factors in the real gauge estimate fit the
regular analytic radius for small \(g\). The independent-coordinate
room in the window estimates is the room of the ordinary regular
formation rule.
Writing \(a=n-j\) and \(s_n=\delta_n/\alpha_n\), the near part per witness
window is bounded by
\[
 C\|f_i\|_\infty\sum_{j\le n} A_j
 \left[
  \left(\frac{\delta_n}{\alpha_j}\right)^2 L^{-\alpha_H a}
   +\left(\frac{\delta_n}{\alpha_j}\right)^3 L^{-2a}
 \right]
 + C\delta_n^2s_n\operatorname{Lip}(f_i)\sum_{j\le n}h_j.
\]
Indeed \(A_j/\alpha_j^2\le Cs_n\),
\(\delta_n/\alpha_j\le Cs_n(1+a)^{1/2}\), and \(h_j=h_n L^{-a}\).
Both geometric series, including the polynomial in \(a\) in the cubic
term, give \(C\delta_n^2s_n\|f_i\|_{\sharp,n}\).
Here \(\|f_i\|_{\sharp,n}\le C\|f_i\|_\infty\) by the lower bound
on \(\mathsf r\). The quadratic term that survives the Wilson subtraction
is included in this estimate. Its coefficient is bounded by
\(A_j/\alpha_j^2\), so the small factor \(s_n\) comes from the formed
source derivative. The cubic Cauchy estimate alone would leave this
quadratic term at the size of the insertion.

For the ordinary old remainder, the vacuum subtraction gives
\((\delta_n/\alpha_j)^2 L^{-4a}\) on the admitted core, which pays for
the fine-site count. Its derivative has the additional factor
\(x_j^{-\kappa_0/2}\) from the ordinary remainder bound and
\eqref{eq:local-formation-LB13}. Its contribution per witness window
is at most
\[
 C\|f_i\|_\infty\delta_n^2\sum_{j\le n}
                       x_j^{-\kappa_0/2}\alpha_j^{-2}
 \le C\|f_i\|_\infty x^{-1}p_0^2 x^{2-\kappa_0/2}.
\]
The last inequality uses \(\alpha_j^{-2}\le Cx_j\) and the lower
flow bound. It is useful to retain the small field ratio until after
the sum. The common-background, common-base and pointed assumptions of
the ordinary formation estimate are needed for these estimates. A bound
for a constant source on a larger disc cannot furnish the estimate
for the derivative in \(f_i\).

For a long ordinary domain entirely in a homogeneous nonwrapping
window, the same vacuum subtraction is performed in its axial chart.
At age \(a\) the complete field has potential at most
\(C(2+\operatorname{diam}_jX)\delta_n L^{-2a}\) in that chart, and
the normalized current has size at most \(C\delta_n L^{-3a}\).
The first derivative at the trivial pair is zero by constant colour
invariance. The quadratic Cauchy bound therefore contributes at most
\[
 C A_j\|f_i\|_\infty e^{-\kappa d_j(X)}
 (2+\operatorname{diam}_jX)^C
        (\delta_n/\alpha_j)^2 L^{-4a}.
\]
The bound is used only on the admitted core. If the polynomially
enlarged field ratio is at least the half-disc threshold, the
unexpanded ordinary core bound gives the same majorant after enlarging
the constant. For domains reaching beyond a fixed enlargement of the
witness window around their source anchor, a spare exponential in
\(d_j(X)\gtrsim L^a\) pays for the
polynomial and gives a summable gain after the fine-site count.
This argument concerns only the homogeneous domains. For an ordinary
domain depending on \(f_i\) whose enlargement reaches a live region
outside the contact neighborhood, the length is at least
\(c\mathsf rL^a/M-C\), since the determining set must also reach
\(\operatorname{supp}f_i\). For a winding domain the bound is
\(c(\ell/a_s)L^a/M-C\). The ordinary source-disc estimate and the
rooted count then put these terms into the exponentially small far
error. No small homogeneous field ratio is used on them. Inside
the contact neighborhood the live terms are estimated by the age
series. The formation and remainder bounds have to hold on the
actual regular domains of the unsourced histories.

Let \(i(x)=C[(\log x)^{3\mathfrak p_0-q}
       +p_0^2 x^{2-\kappa_0/2}]\), allowing a larger constant for the
ordinary far sum. The estimates give
\(\|I_i\|_\infty\le C\mathsf r^4x^{-1}i(x)\|f_i\|_\infty\)
for the regular part and
\(\|X_i\|_\infty\le C\mathsf r^4x^{-1}p_0^2\|f_i\|_\infty\).
% Earlier Part II, Section 8.5, source derivative, radius cancellation,
% shortened bound, long constituents and branch readings, 98771--98842,
% 99793--100296, 100416--100448 and 100661--100728.
\paragraph*{A bound using the physical readings.}
The correlation expansion gives a second way to obtain the
irrelevant-source bound. On the same bare torus it starts
with the product Haar measure, the same plaquette profiles,
and the action \(A(\beta-w,U)\). Its extra bare constant is
source independent and cancels in \(Z(w)/Z(0)\).
One can therefore use this calculation for the physical
generating function with its own standing histories and
running coefficients.

Use the vacuum-subtracted source derivatives of that
expansion on their original domains. Write \(a_j\) for
its tube radius and \(v_{j-1}\) for the running maximum
of its source weights. The
source-disc estimate gives the majorant \(C v_{j-1}a_j^2\)
before decay and the norm of the test function. The old
remainder has the separate majorant \(C r_{\rm s}^{-1}
g_j^{\kappa_0}\), with \(r_{\rm s}\) the source radius.
These bounds precede evaluation at a field.

At witness level \(n\), write \(t=L^{j-n}\), \(\nu=1+n-j\)
and \(D_X=\operatorname{diam}_j X+5M\). First group the
constant-source terms contained in the \(\nu\)-layer
enlargement of the anchor cube with their Wilson subtraction
at coefficient \(b'_j(0)\).
For the remaining constant-source terms and for the
nonconstant-source differences and old remainders, retain
a domain in the short sum when
\[
 D_Xt\le \rho/16,\qquad C_MD_Xt^2\mathfrak e/a_j\le1,
\]
where \(\rho\) is the witness separation in its lattice
units and \(\mathfrak e=C_{\rm w}\varepsilon_n\) is a
fixed bound for the real-field amplitude.
The native grouped core estimate uses its analytic
box extension on the field-rescaling disc. At the
real endpoint the extension must give the same field
and current as the read on the smaller box, up to gauge.
The grouped coefficient vanishes to second order at
the trivial field. Schwarz's estimate on this disc
gives the factor
\[
 C v_{j-1}a_j^2
 \left(\frac{\mathfrak e_Q}{a_n}\right)^2
 \left(\frac{a_n}{a_j}\right)^2
 \nu^3(1+a_n/a_j)t^{4+\hat\alpha},
 \qquad \mathfrak e_Q\le C\varepsilon_n .
\]
Here \(\hat\alpha>0\) is the exponent in the grouped core
estimate. The grouped core uses the same counterterm
coefficient as the running source. For the individual
terms in the rest of the short sum, the vacuum value
and the first field derivative vanish after subtraction
and constant colour invariance. In the real box gauge
the quadratic estimate costs \((C_MD_Xt^2\mathfrak e/a_j)^2\)
times the original majorant. Thus the two radii cancel
in the grouped core, and \(a_j^2\) cancels in each
individual term with majorant \(C v_{j-1}a_j^2\).

The factor \(t^4\) pays for the \(O(|S^+|t^{-4})\)
older anchors over the enlarged witness support \(S^+\).
In the core the remaining factor
\(\nu^3(1+a_n/a_j)t^{\hat\alpha}\) is summable by the
running radius comparison. The constant-source tails
outside the core have length at least \(\nu-1\), and
the absolute domain sum pays for \(D_X^2\) and gives
exponential decay in \(\nu\). The comparison of
nonconstant and constant sources uses the Lipschitz
profile and has the summable factor \(L^{-(n-j)/2}\).
For the old remainder the corresponding age sum is bounded
by its memory majorant \(C r_{\rm s}^{-1}\mathfrak m_n\).

Outside the grouped cores, if either short-domain
inequality fails, then
\[
 d_j(X)+5\ge c\,t^{-1}
                 \min\{\rho/M,\ a_n/(M\mathfrak e)\}.
\]
This uses \(D_X\le 4M(d_j(X)+5)\) and the bound on
\(t a_n/a_j\). Use twice the original analytic majorant
on these long terms. The length gain pays for the
anchor count and the ratio \(a_n^2/\mathfrak e^2\)
at weak coupling. Indeed \(a_n/\mathfrak e\) grows
as \((\log x)^{q-\mathfrak p_0}\) and the lower
witness bound for \(\rho\) grows as
\((\log x)^{2\mathfrak p_0}\). The exponential decay
at the smaller of these two scales pays for the
logarithmic ratio. With the kept part this gives
\[
 |I(f)|\le C|S^+|\,\|f\|_{\sharp,n}
 \bigl[\widehat v_{n-1}\varepsilon_n^2+
                 r_{\rm s}^{-1}\mathfrak m_n\bigr],
 \qquad \widehat v_{n-1}=\max_{j\le n}v_{j-1}.
\]
At the witness, \(\widehat v_{n-1}\le
C(\log x)^{\mathfrak p_0-q}\) and
\(\mathfrak m_n\le Cx^{-(\kappa_0-2)/2}\).
This gives the power \(3\mathfrak p_0-q\) in \(i(x)\).
The source bound on the analytic tube by itself has
the larger power \(\mathfrak p_0+q\).

The localized function of \(V\) evaluates only the
shortened sum at the prolonged box field and is zero
off its window. For the global field add the whole
long sum. On a history with no live label meeting the
neighborhood \(S_i^{+\rho/4}\) and with the window test
passed, evaluate the shortened sum at its real
background and add only those long domains in its
ordinary region \(\Lambda_j(h)\) at each level.
The buffer \(M\check R_jt\le\rho/8\) puts the source
anchors inside the ordinary region with their required
margin. The source profiles there are one. Use
equality with the whole-lattice constituents on
the admitted ordinary reads.
In the joint law multiply this irrelevant insertion
and the corresponding main insertion by the same
per-support indicator. The long terms use their own
original pair domains. The history bound is uniform
before averaging the integrated variables of a history.
The discarded indicators, the replacement by the
localized main insertion, and any history-dependent
share of the vacuum subtraction enter the separate
comparison terms.

% Earlier Part II 73715--73759 and 100592--100650;
% compare the manuscript's empty-trace recipe.
In a standing history, the geometric part of the reading
comparison uses the empty trace. The record sums out a
healed component and keeps its contribution in the ordinary
activities and constants. The live regions are those still
present at the witness level. On the event that none meets
\(N_i=S_i^{+\rho/4}\), the standing trace is empty throughout
\(N_i\). For the grouped core, \(\nu L^{j-n}\) is bounded
independently of \(j\), so its enlarged anchor cubes have
bounded size in witness units. The box-extension margins
of the core also have fixed size in those units, so its
reads fit inside \(S_i^{+3\rho/16}\). For the
other short terms the diameter bound \(D_Xt\le\rho/16\)
and the source support put their complete reads inside
\(S_i^{+3\rho/16}\). This leaves \(\rho/16\) for rounding
to the larger history cells. By \eqref{ccf:trace-padding},
this rounding costs at most \(C_{\rm tr}M\check R_n\)
in witness units. The lower witness bound for \(\rho\)
dominates this cost at weak coupling.
Thus all these short trace windows are inside \(N_i\).

For the history background on that window the preceding
levels obey \(x_j\ge x_n-C\). It follows that
\(g_j\le\sqrt2\,g_n\) when \(x_n\ge 2C\).
The degree function \(g(\log g^{-2})^{\mathfrak p_0}\)
is increasing for \(\log g^{-2}>2\mathfrak p_0\).
Take \(x_n\) large enough that this interval contains
\(\sqrt2\,g_n\). Then
\(\varepsilon_j\le\sqrt2\,\varepsilon_n\) throughout
the older levels. The native data bounds therefore use
one real-field amplitude at the witness. On each common
admitted pair domain,
\cref{lem:ccf-empty-standing-trace,prop:ccf-whole-lattice-identification}
identify the short history terms with their own
whole-lattice terms, for the same input list and
pointing. This identity also holds for their first
source derivatives. The long terms keep their
history-dependent domains and use the long-domain bound.

% Earlier Part II, Theorem 8.5(E), 101220--101271.
In this reading keep the variables integrated inside
each history in the unsourced joint probability law.
Let \(\mathcal G\) be the data fixed for its inner
integral. In the representation with a fixed
reference measure and a source-dependent exponent
\(A(z)\), the history's first jet is
\(\mathbb E[\partial_i A(0)\mid\mathcal G]\).
Its mixed jet is the conditional mean of
\(\partial_{12}A(0)\) plus the conditional covariance
of the first derivatives. Total covariance gives
\[
 \mathbb E\!\left[
  \operatorname{Cov}(\partial_1 A(0),\partial_2 A(0)
                    \mid\mathcal G)\right]
 +\operatorname{Cov}\!\left(
  \mathbb E[\partial_1 A(0)\mid\mathcal G],
  \mathbb E[\partial_2 A(0)\mid\mathcal G]\right)
 =
 \operatorname{Cov}(\partial_1 A(0),\partial_2 A(0)).
\]
Apply the uniform physical bound to the three
irrelevant covariances in this joint law. These pair
one irrelevant insertion with the other main
insertion, or the two irrelevant insertions with
each other. The two sources' factors give the
bound \(C\mathsf r^8x^{-2}i(x)(p_0^2+i(x))
\prod_i\|f_i\|_\infty\). This accounts for their
shares both inside and between histories, giving
the same covariance error in the bound below.
The bound on these three covariances holds under
any probability law with the stated physical
supremum bounds. To use it in the standing-class
decomposition here, it suffices to identify the
evaluated first-source sum, with its vacuum and
Wilson subtractions, on its admitted real pairs.
The empty trace comparison gives equality within
one whole-lattice recipe. For a transfer between
the two expansions one must also match their
formed whole-lattice sums and counterterms.
The mixed derivative of the exponent itself
uses the support and age estimates below.

In the mixed derivative the regular support must reach both original
supports, and at age \(a\) its length is at least \(c\mathsf r L^a/M-C\).
The Cauchy cost in the slow sources is uniform at the witness.
The anchor count times the spare exponential is summable in \(a\)
and gives \(Ce^{-c'\mathsf r}\prod_i\|f_i\|_\infty\).
For the remaining terms put
\(\epsilon_{\rm far}=e^{-c\mathsf r}+e^{-c\check R_{k_0}}\) and
\(\Upsilon=L^{C\check R_{k_0}}\). The localized-background error uses
both radii: the collar at the test cube and the distance to the
nearest live region. The bound on \(B_i\) has a factor
\(\epsilon_{\rm far}+\mathbf1_{\mathcal B_i}\Upsilon L^{4a_i}\);
the live mixed derivative has the corresponding \(L^{8a}\) bound.
Here the ages are lineage ages. Summing over the \(C\mathsf r^4\)
cells in contact with the supports and applying
\cref{prop:source-domains-4-14} at rate at least eight gives
\[
 \mathbb E_{\mathsf P}\big[\mathbf1_{\mathcal B}L^{8a}\big]
 \le C\mathsf r^4e^{-cp_0}.
\]
Since \(X_i\) depend only on \(V\), covariance bilinearity and
\(|\operatorname{Cov}(A,B)|\le 2\|A\|_\infty\|B\|_\infty\) for
the bounded regular terms give
\[
\begin{aligned}
 S^T_2(f_1,f_2)
   &= (\zeta'_{k_0})^2\operatorname{Cov}_{\lambda_{k_0}}(X_1,X_2)+E,\\
 |E|/u
   &\le C\mathsf r^8 i(x)(p_0^2+i(x))
       +x^2(\log x)^C\big[
              \epsilon_{\rm far}
             +\Upsilon^2\mathsf r^4 e^{-cp_0}\big].
\end{aligned}
\]
The second line uses the first and second contact moments for the
covariances involving \(B_i\), and the first moment of the mixed
derivative. In the witness window \(\mathsf r\le C(\log x)^{2\mathfrak p_0}\);
the first term tends to zero for \(q>21\mathfrak p_0\) and sufficiently
large \(\kappa_0\). The other terms tend to zero by the lower window
bound, \(r>1\), and \(\mathfrak p_0>\max\{1,r\}\).

The preceding error bound is uniform under the bare lattice translations
used in the witness covariance estimate. The physical left-hand side
is invariant under these translations, so average the identity over
the same translations whenever the witness estimate is an average.
The covariance estimate is applied with the widened profiles
\(\tilde f_i\). The widening is at most \(CMa_s\), and these supports
still have separation comparable to \(\mathsf d\). In the resulting
bare kernel sum, replacing \(\tilde f_i\) by \(f_i\) costs at most
\(C\|\tilde f_i-f_i\|_\infty\|f_{3-i}\|_\infty\) times the coupling
factor \(x^{-2}\). Indeed the squared curvature kernel on the
separated supports has a bounded absolute sum with the two plaquette
volume factors. By (b) the cost relative to \(u\) is \(O(M/(\vartheta
\mathsf r))\). This replacement precedes the Riemann estimate, which
uses the Lipschitz bounds of the original profiles.

(c) \emph{Gaussian extraction: the free template.} The covariance of the
bounded witness fields $X_i$ is obtained from
\hyperref[f2:witness-generating-function]{the generating function on the fluctuation segment}. Denote by $\mathcal{H}_jB$ the exact bare $2$-cochain of smallest norm that satisfies $Q^{\mathrm{pl}}_j\mathcal{H}_jB=\mathrm{d}B$. It is the curvature of the linear minimiser, namely $\mathrm{d}\,\mathsf{H}_j$ at $U=\mathbf{1}$, a point at which $J=0$ and the three Green's functions of \cref{sec:propagators-propagators-in-a-background-field-exponential-decay-produced-by-the} agree. We then have

\begin{gather*}
\FitMath{\mathcal{C}_{(j)}=\mathcal{H}_{j-1}C^{\mathrm{tr}}_j\mathcal{H}_{j-1}^*,} \\[3pt]
\FitMath{|\mathcal{H}_j(p,b)|\le C\eta_j^2e^{-\delta|x(p)-x_b|/a_{K-j}},} \\[3pt]
\FitMath{|\mathcal{C}_{(j)}(p,q)|\le C\eta_{j-1}^4e^{-\delta|x(p)-x(q)|/a_{K-j+1}}.}
\end{gather*}

The free propagator framework is developed in \cite{B84a}.  For the
centered blocks used here, the uniform pointwise estimate follows from
\cref{prop:hb-joint-strong-inverse} and its exponentially weighted version.
Rescale a level-$j$ block to side one and put $h=\eta_j$.  At the trivial
flat reference, let $H_hB$ be the solution with fine force zero and coarse
datum $B$.  In the notation of \eqref{hb:joint-inverse-formula},
\[
 H_h=K_0^{-1}Q^\sharp
                  (QK_0^{-1}Q^\sharp)^{-1}.
\]
This uses the straight row in \eqref{hb:actual-flat-rows}.  The gauge
penalty vanishes: testing the first saddle equation against $D\lambda$
with $q\lambda=0$ eliminates both its curl term and its multiplier term,
while \eqref{hb:site-gauge-inverse} makes the remaining gauge test
surjective.  Thus $R D^*H_hB=0$.  The first equation is consequently the
Euler equation for the curl energy on $QA=B$, and its curvature is the
least-norm exact curvature with datum $\mathrm{d}B$.  Indeed, if $f=\mathrm{d}v$ is
exact and $Q^{\mathrm{pl}}f=\mathrm{d}B$, then $Qv-B$ is a closed coarse one-form.
On the trivial periodic lattice it is a coarse gradient plus a constant.
The identities $qE_q=1$, $QD=\mathrm{d}_cq$ and $Q(\mathrm{constant})=
\mathrm{constant}$ change $v$ by a closed fine one-form to impose
$Qv=B$ without changing $f$.  This identifies the curvature with
$\mathcal H_j$ modulo closed forms.

Take $B$ to be a unit vector on the coarse bond $b$.  The weighted bound
\eqref{hb:joint-strong-bound} gives
\[
 |\mathrm{d}_h H_h(p,b)|\le C e^{-\delta\operatorname{dist}_{\rm phys}(p,b)}.
\]
Its constants are uniform in the fine subdivision and in the volume.
A raw bare plaquette curvature is $h^2\mathrm{d}_hH_h$; restoring the physical
block side therefore proves the factor $\eta_j^2$ and the distance scale
$a_{K-j}$ in the displayed estimate for $\mathcal H_j$.  This argument
uses the scalar equation and projector smoothing in
\cref{ch:hb-strong-inverse}, so it does not require a differentiated
resolvent to map $L^2$ into $L^8$.  The symbol $C^{\mathrm{tr}}_j$ denotes
the single-step covariance in a tree gauge,

\begin{gather*}
\FitMath{C^{\mathrm{tr}}_j:=\iota\,\big(\iota^*\mathcal{H}_{j-1}^*\mathcal{H}_{j-1}\,\iota\big)^{-1}\iota^*,} \\[3pt]
\FitMath{\mathrm{d}'\,C^{\mathrm{tr}}_j\,\mathrm{d}'^*=C^{\mathrm{curv}}_j,}
\end{gather*}

where $\iota$ embeds into the $1$-forms $B$ of level $j-1$ the subspace of those that satisfy $B=0$ on one contour for each site (a tree in each block, imposed sharply) together with $Q^{\mathrm{st}}B=0$. In this formula $Q^{\mathrm{st}}$ denotes the straight average over the blocks of the single step from level $j-1$ to level $j$, at $U=\mathbf 1$. It is thus the covariance, on that subspace, of the quadratic form $\langle B,\mathcal{H}_{j-1}^*\mathcal{H}_{j-1}B\rangle=\|\mathcal{H}_{j-1}B\|^2$, which equals twice the free action of the minimiser with datum $B$. On the subspace one has $\|\mathcal{H}_{j-1}B\|^2\ge c\|\mathrm{d}'B\|^2\ge c_L\|B\|^2$. It differs from the covariance $C^{(j-1)}$ of \cref{sec:propagators-propagators-in-a-background-field-exponential-decay-produced-by-the}-\cref{ch:local-formation-program} at the trivial background in two ways. The gauge differs: there it is imposed softly, on the means over all contours. The constraint also differs: the straight average is not the linearisation of $\mathcal{M}$. The two gauges induce the same curvature law.

%
\hypertarget{flat-curvature-gauge-comparison}{}
Let \(d'\) and \(d''\) be the fine and coarse coboundaries, and
\(\mathsf A\) the plaquette average for this step. The flat identities
\(d''Q=d''Q^{\rm st}=\mathsf A d'\) give the common image of the two
constraints under \(d'\). To see surjectivity, take \(h=d'B\) with
\(\mathsf Ah=0\). The coarse defect \(QB\) is closed. Remove its
constant harmonic part by a fine constant cochain, since the
average on constants is invertible. Remove its exact part by
a fine gradient whose scalar has the prescribed values at the
coarse centers. For \(Q^{\rm st}\) prescribe the corresponding
block means of the scalar instead. Both corrections preserve \(h\).

In \(QB=0\) the fibre of \(h\) consists of gradients with scalar
zero at the coarse centers after subtraction of one constant. A
fine harmonic difference would give a nonzero coarse harmonic
difference. On that fibre the contour mean to a site changes by
the scalar at the site minus its value at the center. Integration
of the normalized quadratic contour weight in these coordinates
has the same value on each fibre. In \(Q^{\rm st}B=0\), subtract
the tree path sums by a gradient with zero block means to set the
tree bonds to zero. A gradient vanishing on the trees has a scalar
constant on each block, and \(Q^{\rm st}d'\lambda=0\) then forces
these constants to agree. The tree condition therefore selects
one lift of each \(h\). All changes of coordinates are linear,
with constant measure factors canceled in normalization.

Consequently both flat quadratic laws induce the Gaussian with
curvature precision \(\mathsf V_{j-1}^{+}\) restricted to the exact
curvatures satisfying \(\mathsf Ah=0\), where \(\mathsf V_{j-1}\)
is the unconditioned covariance of \(Q^{\rm pl}_{j-1}F\). Its
covariance is \(C_j^{\rm curv}\). Transport by \(\mathcal H_{j-1}\),
which depends only on \(h\), gives the same bare slice
\(\mathcal C_{(j)}\) in either gauge. This calculation uses the
full quadratic laws. The comparison to the cut windows uses the
separate window estimates.

\hypertarget{witness-free-complement}{}Summing over $j\le k_0$, where the coarsest slice gives the dominant contribution, one obtains the following bound. Here $t=|x(p)-x(q)|/a_s\ge\frac12C_W(s)$, and $\delta_{\mathrm{fl}}>0$ is a rate formed from the rates $\delta$ of the slices, depending only on $L$ and the dimension.

\begin{gather*}
\FitMath{\big|C^{\mathrm{fl}}_{K,s}(p,q)\big|\le C\big(\tfrac{a_K}{a_s}\big)^4e^{-\delta_{\mathrm{fl}}t}=C\pi^2t^4e^{-\delta_{\mathrm{fl}}t}\,\sigma_{pq}:}
\end{gather*}

the fluctuation part is small because the curvature of the linear minimiser decays. Because $t^4e^{-\delta_{\mathrm{fl}}t}$ is decreasing only for $t\ge4/\delta_{\mathrm{fl}}$, the room is given the floor $C_{W,0}A_0^2\ge8/\delta_{\mathrm{fl}}$. Over $j>k_0$ the sum of the exponentials reconstitutes the power law: $\sum_iL^{-4i}e^{-tL^{-i}}\le Ct^{-4}$.

\emph{The interacting measure.} We apply the formula of the Strategy at level $k_1$, whose backgrounds are flat on the scale of the pair. Only the classes $[h]$ occur, and no finer class of histories: the sourced exponent of level $k_1$ is the one that \cref{sec:sources-analytic-dependence-on-the-source-locality-and-decay-in-tame} provides. (c1) \emph{The first member is an error term.} Consider the terms of the levels $k_0<j\le k_1$ and a source that is rough on their scale. Their slopes in such a source remain $\le Cx^{-1}(\log x)^{\mathfrak{p}_0+q}$, measured against $\|f_i\|_\infty$ (the norm $\|f_i\|_{\sharp,K-j}$ of the level of the term can exceed it by a factor of up to $L^{j-k_0}$). The recursion for these slopes closes since $\mathsf n\le C\log x$ (its size is set through $\bar x_s\le Cx$). Hence $\mathsf n^2(\log x)^{\mathfrak{p}_0-q}\le C(\log x)^{2+\mathfrak{p}_0-q}$ is small ($q-\mathfrak p_0>2$, $g$ small). Suppose first that no large-field region touches the supports at level $k_1$. Then $|\Phi^{(k_1)}_i-c_i|\le o(1)\,x^{-1}\|f_i\|_\infty$ for a constant $c_i$, where $o(1)$ refers to $\mathsf d\to0$, uniformly in the datum. In the other case $|\Phi^{(k_1)}_i-c_i|$ grows no faster than $L^{\mathsf n}L^{4a}\|f_i\|_\infty$ in the age $a$ of the touching region. The age series of \cref{candidate:sources,ch:all-state-source-transfer}, taken at the single level $k_1$ with $b\ge8$, beats this growth: $L^{\mathsf n}\le L\bar x_s^2$ is a power of $x$, while the series lies below every power. Consequently $|\langle\Phi^{(k_1)}_1;\Phi^{(k_1)}_2\rangle_{\mathsf{P}^{(k_1)}}|\le o(1)\,x^{-2}\prod\|f_i\|_\infty$. (c2) \emph{The second member becomes one number.} The quantity $\Phi^{(k_1)}_{12}$ is the sum of two parts. The first is the mixed derivative of the sum $\sum_{j\le k_1}\sum_Xe_j(X;w)$ of localised source energies from \cref{candidate:sources}. The second is a part that depends on the fields and is bounded only crudely, by a quantity of order $L^{2\mathsf n}$ for each bond of the determining set; by \cref{candidate:sources}, the size of that set is counted by the unit volumes of level $k_1$ together with the sites $\sum_{j\le k_1}\#_j(Z_j)$ of the large-field regions (every source derivative costs a factor $L^{\mathsf n}$). We translate the pair by the points $v$ of a grid $\mathfrak P\subset\Lambda_{s,\ell}$ whose spacing is comparable to $\mathsf{d}$, and we write $f^v_i:=f_i(\cdot-v)$ for the translated sources. The left-hand side of the formula is unchanged by this. The sum over $v$ of the mixed derivatives equals the mixed derivative evaluated at the two array sources $\sum_v\epsilon_vf^v_i$, whose norms are those of $f_i$, averaged over the signs $\epsilon\in\{\pm1\}^{\mathfrak P}$ (this is the sign-averaging argument proving \eqref{eq:two-point-extraction-partial-ET7}). The crude bound is therefore divided by the number $\approx L^{4\mathsf n}\hat{\mathsf{d}}^{-4}$ of pairs contained in a unit volume of level $k_1$. What survives is $L^{-2\mathsf n}\hat{\mathsf d}^4\le\bar x_s^{-4}\hat{\mathsf d}^4$, and this requirement fixes $\mathsf n$. There remains the number 

\begin{gather*}
\FitMath{c_{12}:=|\mathfrak P|^{-1}\sum_{v\in\mathfrak P}\partial_{z_1}\partial_{z_2}\sum_{j\le k_1}\sum_Xe_j(X;z_1f^v_1+z_2f^v_2)|_0}
\end{gather*}

. (c3) \emph{The number is evaluated.} The number $c_{12}$ is the value at a single point of the mean of $\Phi^{(k_1)}_{12}$ over $\mathfrak P$, namely of the term without large-field regions, taken at the trivial block field. Its exact single-step recursion runs from level $k$ to level $k+1$ for $k_0\le k<k_1$ and is evaluated at the flat terminal field. The measures there are positive, so a remainder is carried through a step by its supremum alone. For each step, two inequalities on a single fluctuation integral \cref{ch:local-formation-program} are then required. That integral belongs to the small-field step taken over the whole lattice, with a background that is real and regular and with coupling $x_j\approx x$. The first inequality concerns the covariance of two sums of background plaquette energies. The second concerns the mean value of a bilinear form whose arguments are the background plaquette angles. The expansion of \cref{ch:local-formation-program} gives terms for these comparisons. We must bound the interacting errors on the actual fluctuation integral, using Cauchy estimates on its admitted polydisc, and control the Taylor error in replacing the regular background by the flat one. At the flat background, the object that agrees with the free template is the curvature covariance $C^{\mathrm{curv}}_j$ introduced above. It determines $\mathcal{C}_{(j)}$, and $a_p$ depends on nothing else. 

\paragraph{The clean restriction at fixed exterior data.}

The clean restriction has to be made on a standing class with
its stored variables still present. Fix one such class and
keep its exterior datum fixed. At a step of level \(t\)
above the first level of the clean window, add a small
complex constant \(s\) to the input boundary term at
each level-\((t-1)\) live cube near the clean
neighbourhood. The modified boundary system belongs
to the complex box of the step identity for \(s\) in
a neighbourhood of zero. The real moves are
independent of the boundary system, so an input term
with \(N\) marked cubes acquires just \(e^{sN}\).
With the stored variables retained, the summed step
identity over the fixed output standing part therefore
equates two polynomials in \(z=e^s\) near \(z=1\).
The constant coefficient is the sum of the allowed
histories and, on the bracket side, the same bracket
with the marked inner histories omitted. The standing
input at the first level already enforces its live
condition. Take the coefficient algebraically, keeping
the original moves, shifts and counterterms. Healings
at or below the first level of the window have no
such mark. With
profiles redrawn after healing and compensated at each
plaquette, the coupling lag vanishes on the clean box.

To estimate the quotient, restrict the inner-history measure
of each bracket to the histories allowed by the mark. At the
same stored datum a restriction decreases its total variation.
An inadmissible collar gets weight zero. Keep the empty-family
factor, the relative geometry and the interaction potential on
each surviving tuple. Thus every absolute activity majorant of
the full bracket still applies to the restricted one, including
the small factor for each healed collar. The rooted logarithm
bound for the relative bracket is \(Ce^{-5p_0(g_t)/4}\) per
root, with the same field-independent majorant for both sums.
Only terms that meet a deleted inner-history choice or collar
can differ. Root a changed term in the healed collar
\(C^\sim\) of the deleted choice. In the fixed output
label \(S'\) we have \(C^{\sim2}\subset\Lambda_t(S')\).
Thus \(C^\sim\) is contained in the one-layer erosion of
\(\Lambda_t(S')\) and avoids the fixed hulls of the
standing components at all levels. The ordinary
rooted entropy sum applies there. The marked input
cube is in the padded box of the deleted choice.
Its possible roots lie within the operation's read
radius. The number of such roots per marked cube
is polynomial in \(\check R_t\), and that polynomial
is absorbed by decreasing the reserve in
\(e^{-5p_0(g_t)/4}\). Summing the two absolute bounds
over the resulting roots gives, with \(0<c<1/4\),
\[
 0\le\Delta_t
 :=\log K_t-\log K_t^{\rm cl}
 \le C c_t e^{-(1+c)p_0(g_t)}.
\]
The lower bound follows by positivity before exponentiation.
This comparison uses the standing-class identity with the
stored variables, the absolute bracket majorants and the full
carrier count. The restriction is of the coefficients of that
identity. The covariance and counterterms used to form it are
those of the full input.

Positive kernels preserve the resulting domination in later
moves. Thus, with \(\epsilon_j=C\sum_{t\le j}
c_t e^{-(1+c)p_0(g_t)}\) over the marked levels, the
clean and full incoming measures have ratio
\(r_{j-1}\in[e^{-\epsilon_{j-1}},1]\) on their common
standing variables. At step \(j\), use only the strings of
that step that remain clean on both sides. At a fixed
output datum, let \(\mathbb E_j^{\rm ref}\) be their positive
conditional expectation with the full incoming measure.
Disintegration and normalization give
\[
 \mathbb E_j^{\rm cl}F
   =\frac{\mathbb E_j^{\rm ref}(r_{j-1}F)}
          {\mathbb E_j^{\rm ref}r_{j-1}},
 \qquad
 e^{-\epsilon_{j-1}}\le
 \frac{r_{j-1}}{\mathbb E_j^{\rm ref}r_{j-1}}
 \le e^{\epsilon_{j-1}}.
\]
These statements are on the common positive output support,
up to null sets of its marginal. The lower bound splits off
a reference kernel with weight \(e^{-\epsilon_{j-1}}\).
Composing the splits leaves a reference branch of weight
\(P_0=\exp(-\sum_{k_0<j\le k_1}\epsilon_{j-1})\)
on the clean classes. The lost mass is at most the unclean
mass plus \(1-P_0\le\mathsf n\epsilon_{k_1}\).
This branch uses the original run weights and no counterterms
formed from a restricted exponent.

\paragraph{The product window on a clean interior box.}
At step $j=i+1$ let $I$ be the free bonds of the comparison
box. In input-level units let \(r_{\rm dec}\) bound
the diameter of the enlarged decision cube plus the
read of its affected pivot rows. Form \(I^+\) by enlarging
\(I\) by \(r_{\rm dec}\). Include every decision cube
whose enlarged stencil meets \(I^+\) in the clean
neighborhood. In the native box geometry a sufficient
collar width over \(I\), in these same units, is
\[
 \max\{100M\check R_i+4r_{\rm dec},\ 111LM\check R_j\}.
\]
The first width covers the complete decision stencils.
The second places the box on the ordinary side of the
live and creation buffers. The native cube count for
these enlargements costs a fixed factor bounded by
\(C L^8+(2L)^4\), which the rarity exponent absorbs.
%
%
Use the clean condition
before the integrations and healing of step $j$. The box and its
read margin are then inside $\Lambda_j$: no creation or live
component meets that margin and no operation of $\mathbf R$ at
this step acts there. In particular no rim, collar or strip test
reaches the changed bonds.

Write \(\widehat g_i=\theta_i^{-1/2}\) for the coordinate scale
of the reference frame. Fix $W=V_j$ and the exterior variables of the positive standing
class. Take the averaging fibre in the original block axial frame
at $V_* = V_i^{\rm bg}(W)$, with the older prescribed datum held
fixed. A change of the literal nonpivot coordinates $\widehat g_i y_b$ for
$b\in I$ changes only those bonds and the pivots in their one-step
read sets. Every such pivot solves its one coarse constraint on
the admitted chart. The change leaves the determining rows, the
coarse field $W$, and the older stored variables fixed. In
particular the coarse plaquette tests and the test backgrounds
computed from $W$ have the same values.

On the box the exact fluctuation test of each retained cube is
the full-threshold bound on the supremum of the free bond
logarithms that it reads. Its product over the cubes therefore
imposes exactly $|\widehat g_i y_b|<\delta_i$ for $b\in I$. Every such
bond is tested by a cube inside the ordinary split. The tests
involving both box and exterior bonds contribute these same
bounds and restrictions on the fixed exterior. There is no
additional pivot cutoff. The pivot change of variables fixes the
nonpivot coordinates and has a single solution on this support,
with pivot values of size at most $C_L\delta_i$.

To check the converse, keep the exterior of an admitted clean
term and vary the box coordinates throughout these balls. On the
affected approximate tests the fixed centers differ from $V_*$
by at most a small fraction of $\delta_i$ by the real collar
comparison on the test cube. Thus the exact threshold $\delta_i$
implies the approximate threshold $2\delta_i$ at every changed
nonpivot read. The contour cutoffs also pass, since their deviation
is at most $C_L\delta_i$ and the gauge cutoff uses $\varepsilon_i$.
The new-test implication of \cref{lem:large-field-2-2} applies
on the full old density carriers and gives the old threshold with
its fixed margin. At \(i=0\) check the bare plaquette test
on each affected plaquette. The complete ordinary read margin
contains all four bonds. At the retained background its plaquette
deviation is at most \(CL^{-2}\varepsilon_1\) in bare units,
by the curvature estimate of that lemma. Relative to this
background the nonpivots on the window and its fixed exterior
are within \(\delta_0\), and the solved pivots are within
\(C_L\delta_0\). Telescoping the four unitary factors bounds
the plaquette deviation by \(CL^{-2}\varepsilon_1+4C_L\delta_0\).
The adjacent frame radii bound \(\varepsilon_1/\varepsilon_0\).
Choose \(L\) for the curvature term, then use the ordinary
ceiling on \(\delta_0/\varepsilon_0\) to put
\(CL^{-2}(\varepsilon_1/\varepsilon_0)+4C_L\delta_0/\varepsilon_0\)
strictly below one on the whole window.
This check uses the estimates on the changed
fields and their read sets. The earlier recorded answers alone
would not show that they still pass. All other tests or decisions
of the step are exterior reads, or use only the fixed $W$.
The local geometry and the clean restriction of this step are
therefore constant throughout the box window. The gauge coins
inside the box have been integrated in the gauge-fixing factor.
%
%
Insert the normalized ordinary gauge weight on the summed
standing input before conditioning on the exterior. The
ordinary gauge group fixes the coarse centres and the
older held datum sites. The incoming density and insertion
are invariant on these orbits, and the normalized gauge
weight has orbit integral one. A
draw at an interior site can move an incident exterior
fine bond, so disintegrate at the exterior \emph{after}
this gauge integration. That exterior contains the retained
coarse field, the older stored data used by the standing
term, and the fine exterior free coordinates. A fine pivot
is reconstructed from its constraint
even if its location is outside \(I\). Include its full double
block when forming the test read. The exterior may also contain
coins of later moves. Integrate out the coins of the collar or
ordinary gauge draws made before the tests. At an interior site
such a coin fixes the contour and would impose an extra equation
on the free box coordinates. Its integrated contribution is the
gauge-fixing factor in the coarea integral. The ordinary
strip centre, its relative coordinate and Jacobian use the
held fine exterior rows and the coarse strip datum.
Thus a strip selector in its original coordinate is also
fixed in this box integral.

Check the clean margin from the large-field region \(Z_j\)
before healing, on the whole enlarged family of test reads, including the
reconstructed pivots. The collar and rim exploration reads
only the successive enlarging large-field sets, all contained in
\(Z_j\), the complement of the small-field integration region.
The selected exterior decisions therefore stay fixed as the
free box varies. When the earlier density indicator reads
the changed box, the prospective new-test implication above
proves its value on the full product window before it can
be kept among the fixed answers. Sum out the unused interior
diagnostic flags and apply the coarea formula to the term
summed over its standing part.

It follows that on the admitted clean reference class the
indicator in \eqref{eq:clean-conditional-law-CL2} is the product
$\prod_{b\in I}\mathbf1\{|y_b|<b_j\}$ with $b_j=\delta_i/\widehat g_i$,
times an exterior indicator. Use the exact coarea density of
\eqref{eq:clean-conditional-law-CL1} in both numerator and
denominator. The retained gauge-fixing exponential and the
Haar and pivot densities still occur in the weight. At fixed
exterior, freezing the other coordinates of the full quadratic
form gives the compression on $I$ and its linear cross force.
The exterior mixing measure is proportional to $Z_e m(de)$,
since the same box integral $Z_e$ normalizes every inserted
observable. This identifies the window and integration variables
for the matched conditional comparison.

%
%
The positive run also specifies how to impose later tests.
On the admitted box fibre let \(s_j(e)\) be the state
after deleting the integrated slots of step \(j\).
The retained field \(W\), the exterior fine slots, the
older stored slots and the completed label are fixed by
\(e\). The pivots which change with \(y\) are integrated
slots, since the full read margin separates the box
from the retained strip. Hence \(s_j(e)\) is constant
as \(y\) varies on this fibre. The state includes all
retained outcomes read by the subsequent moves.
The subsequent draws, tests and integrations form a positive kernel
\(\mathsf K_j(s_j(e),d\omega)\) on the later records.
Use the kernel conditional on any later coins already
included in \(e\), with the remaining coins integrated
in the kernel.
For a selection \(A\) of these records and the fixed
exterior, set \(k_A(e)=\mathsf K_j(s_j(e),A_e)\).
The conditional quotient with a box insertion \(F\)
is then
\[
 \frac{\displaystyle\int m(de)\,k_A(e) Z_e\,\nu_e(F)}
      {\displaystyle\int m(de)\,k_A(e) Z_e}.
\]
This follows by performing the box integral before the
later kernel integral, for a finite positive denominator.
For a standing class sum the positive terms on their
shared retained data, keeping each term's own \(Z_e\)
and \(k_A(e)\) in the mixture. A selection
on a field already deleted by the run can change the
conditional domain and must be kept in the joint integral.

%
\paragraph{Composition on the retained states.}
Let \(\widehat\tau_j\) be the measure on the disjoint union
of the level-\(j\) state spaces, with the block field,
all still stored slots and the labels used by the run.
Use the original positive operations of step \(j\) as a
map \(\mathcal T_j\) on measures, including the sums at
merges. For the carried insertion \(G\) put
\(\widehat\tau_{k_0}^G=G\widehat\tau_{k_0}\) and apply
the same operations to \(\widehat\tau_j^G\) and
\(\widehat\tau_j\). A test splits a measure by its
indicator and complement, an integration or change
of variables is a positive kernel, and a merge adds
the images from its incoming summands. Use gauge
fixing as a normalized orbit draw. The gauge-invariant
insertions used here can pass through this draw.
These operations conserve total mass after summing
the labels. At each output, the carried density obeys
\[
 G^{[j]}=\frac{d\widehat\tau_j^G}{d\widehat\tau_j}
         =\mathbb E_j G^{[j-1]},\qquad
 \mathbb E_j F=
   \frac{d\mathcal T_j(F\widehat\tau_{j-1})}
        {d\mathcal T_j(\widehat\tau_{j-1})}.
\]
Positivity gives \(|\mathbb E_jF|\le\|F\|_\infty\),
and \(F=1\) gives \(\mathbb E_j1=1\). At a merge
the two measures in this derivative are the sums
of the respective images. In particular the merge
uses the conditional weights of the original run.

For a terminal class retained by the run or by an
explicit refinement of it, restrict the terminal
measures to that class and disintegrate over
\(V_{k_1}=W\). The conditional expectation of \(G\)
is the conditional mean of \(G^{[k_1]}\) under the
uninserted terminal class measure. This follows by
testing the carried terminal measure against a
bounded function of \(W\). Thus the recursion
composes on complete states before projecting to
the block field. For a refined class that imposes
conditions on earlier deleted slots, retain those
conditions in the positive integrals used to form
the refinement. We must check the product window
and the box estimates for its actual conditional
law before applying the Gaussian comparison to
that class.

%
%
The deletion at an isolated ordinary cell
gives a useful case for later prepared and
healed exterior observations. At output level
\(j\), take \(I\) on the bonds of \(e^{+1}\) for
a cell \(e\), and use the event
\(e^{+2}\cap Z_j=\varnothing\).
The rim width is less than one cell. Each
test of the cascade reads within that width
of its current stage set. Run the cascade
from the data off \(I\), with the output field,
the older determining data and the coins fixed,
until a stage set first approaches the padded
cell. Up to that point the two cascades have
the same reads. Thus whether the approach
occurs is determined by this exterior.
On the event of no approach, the complete
post-step state depends only on that exterior.
The slots of \(I\) are integrated in the
ordinary product window.

Write \(\nu_\xi\) for their conditional law
at exterior \(\xi\). Later observations are
obtained by positive kernels from the retained
state. For a bounded insertion \(G\) in \(I\)
and a later test function \(F\), the integral
therefore has the form
\(\int \nu_\xi(G)\,K_\xi(F)\,d\pi(\xi)\),
where \(\pi\) is the pre-deletion exterior
measure on this event and \(K_\xi\) is the
later kernel. Disintegrating over a later
prepared or healed observation takes a positive
mixture of the \(\nu_\xi\). On a class with
a common reference law, a uniform normalized
density bound against that reference is
preserved by this mixture. A finer condition
on an earlier deleted coordinate instead
enters the original conditional integral.

For two bounded witness reads, also carry \(G_1G_2\)
as a single insertion. Put
\[
 T_j=(G_1G_2)^{[j]}-G_1^{[j]}G_2^{[j]} .
\]
Then \(T_{k_0}=0\) and the recursion on the retained
states gives
\[
 T_j=\mathbb E_j T_{j-1}
      +\operatorname{Cov}_j(G_1^{[j-1]},G_2^{[j-1]}) .
\]
Here \(\operatorname{Cov}_j(F,H)
=\mathbb E_j(FH)-\mathbb E_jF\,\mathbb E_jH\).
Indeed, expand the carried product at level \(j-1\),
apply \(\mathbb E_j\), and subtract the product of the
propagated means. Thus \(T_{k_1}\) is the sum of the
one-step covariances transported by the later
\(\mathbb E_j\)'s, with the weights of the actual merges.
In the terminal class, let \(\mathbb E^W\) denote
conditioning on \(V_{k_1}=W\) in the uninserted law.
The covariance under the terminal conditional law
\(\Pi\) of the original reads is
\[
 \operatorname{Cov}_\Pi(G_1,G_2)
   =\mathbb E^W T_{k_1}
      +\operatorname{Cov}^W(G_1^{[k_1]},G_2^{[k_1]}).
\]
Only the errors in the propagated means contribute to
the last covariance if their leading terms depend on
\(W\) alone. This is the decomposition used when
estimating one-step covariances separately.

\paragraph{The two one-step Gaussian comparisons.}
Fix \(k_0<j\le k_1\). Work in level-\((j-1)\) units on a
nonwrapping box containing the read sets in its interior. Write
\(V_*=V^{\mathrm{bg}}_{j-1}(V_j)\), and take derivatives in the
bond logarithms \(V=e^{iZ}V_*\). Let \(y\) have the centered
Gaussian law of the admitted reference box with its covariance
\(C^{(j-1)}\) on free coordinates. This is the covariance with
the reference frame scaling and the actual quadratic weight
\(\tau_{j-1}=x_{j-1}/\theta_{j-1}\) in the Wilson term.
The trace over bond or free coordinates includes the colour slots.
The first variation of the physical
bond logarithms is \(\widehat g_{j-1}\mathsf C y\), where \(\mathsf C\)
is the actual pivot tangent map. Set
\(C_*=\mathsf C C^{(j-1)}\mathsf C^{\mathsf T}\).
Use the uniform covariance bound and the derivative
bounds for the actual pivot map, including its higher
jets, in the norm of these readouts. The box dimension
and the real window radius are bounded by fixed powers
of \(\log x\).
In the following error bounds \(g=x^{-1/2}\asymp \widehat g_{j-1}\).

Use the normalized conditional law on the clean class, with its
actual product window, coarea factors and exterior mixture as in
\cref{ch:clean-conditional-law}. The comparison requires the
complete quadratic form and the same chart for the source
numerators and the denominator. Assume the exterior comparison on
the read set has density ratio between \(e^{-\Xi_j}\) and
\(e^{\Xi_j}\), with \(0\le\Xi_j\le1\), the Gaussian
window tail is \(0\le t_j\le1/2\), and the old
activity, marked remainder and any coupling lag contributions to
the reference box tilt have total bound \(\theta_j^{\mathrm E}\).
The intrinsic odd cubic and even higher terms have bounds
\(Cg n_j\mathsf b_j^3\) and \(Cg^2 n_j\mathsf b_j^4\),
respectively, where \(n_j\) counts the box coordinates and
\(\mathsf b_j\) is the window radius. These parity bounds use an intrinsic reference with every
exterior fluctuation argument set to zero. In particular they
apply to the homogeneous step with no stored exterior. To use
them on a box with stored exterior one must include the
change from those arguments in the exterior comparison.
The generating-function estimate with the stored values fixed
uses the full intrinsic oscillation instead, as in
\cref{cl:intrinsic-reference-hessian}. These estimates and the
gauge-equivariant means used below require a common admitted
analytic tube with room around the symmetric product window
and a small bound for the complete tilt.

Use the following formulas at a fixed compatible exterior,
with the matched reference having zero exterior fluctuation
arguments. For an exterior mixture they apply with a single
reference only after the density and readout comparisons to
that reference. If the reference varies with the exterior,
averaging the conditional covariances also requires the
covariance of the conditional means.

Let \(G,G_1,G_2\) be the gauge-invariant readouts, with common
tube bound \(\mathsf M\) and Taylor derivative bounds \(\mu_m\).
Put \(\bar\mu=\max_m\mu_m\) over the derivatives used here.
Write \(\kappa\) for the curvature size of \(V_*\) on the read
set, \(\kappa_i\) for its curvature size on the box of \(G_i\),
and \(\kappa_{\mathcal B}\) for the curvature size of \(V_j\)
on the window box. The curvatures are measured in their respective
levels, with \(\ell'\kappa\) small for a read box of side
\(\ell'\) in the common nonwrapping gauge. On the admitted clean boxes
use \(\kappa,\kappa_1,\kappa_2\le Cg p_0\). This permits the
variation of an old tilt to be bounded by its size times
the oscillation of the readout, using the first derivative
bound \(C\mu_2\kappa\) and the second-order Taylor bound.
The mean and covariance
comparisons have the form
\begin{align*}
 \mathbb E_jG
 &=G(V_*)+\tfrac12 \widehat g_{j-1}^2
              \operatorname{tr}(C_*D^2G(V_*))+E_1,\\
 \operatorname{Cov}_j(G_1,G_2)
 &=\widehat g_{j-1}^2\langle DG_1(V_*),C_*DG_2(V_*)\rangle\\
 &\quad+\tfrac12\widehat g_{j-1}^4\operatorname{tr}
              (D^2G_1(V_*)C_*D^2G_2(V_*)C_*)+E_2.
\end{align*}
Put \(s_i=g\mu_2\kappa_i+g^2\bar\mu\).
For readouts compared in the same coordinates the required errors
are bounded by
\begin{align*}
 |E_1|&\le C(\log x)^C
   \big\{(g^4+\theta_j^{\mathrm E}g^2)\bar\mu
      +(g^2+(\log x)^{-2q})\mu_2\kappa\kappa_{\mathcal B}\big\}
      +C\operatorname{osc}(G)\Xi_j+C(\log x)^C(\mathsf M+\bar\mu)\sqrt{t_j},\\
 |E_2|&\le C(\log x)^C
   \big\{(g^2+\theta_j^{\mathrm E})
           (g^4\bar\mu^2+g^2\mu_2^2\kappa_1\kappa_2)
       +(g+\theta_j^{\mathrm E})g^3\mu_2\bar\mu
           (\kappa_1+\kappa_2)\big\}\\
 &\quad+C(\log x)^C\Xi_j s_1s_2+C(\log x)^C(\mathsf M+\bar\mu)^2\sqrt{t_j}.
\end{align*}
If the exterior changes the readout map, let \(\eta_i\)
bound the supremum difference between its \(i\)-th
pullback and the reference pullback on the same window.
It costs \(\eta_i\) in the mean. For the covariance
subtract the reference constants \(G_i(V_*)\) first.
Writing \(B_i\) for the resulting reference suprema,
the pullback error is bounded by
\(2\eta_1(B_2+\eta_2)+2\eta_2B_1\).
Indeed, under the positive joint law of the exterior and
the box, expand the covariance in one pullback difference
at a time and use \(|\operatorname{Cov}(A,B)|
\le2\|A\|_\infty\|B\|_\infty\).
Here \(B_i\le C(\log x)^C s_i\) by the same
Taylor bounds as in the density comparison.

Expand the readouts in the actual pivot chart and compare with the
centered Gaussian on its symmetric window. The two linear
variations give the first covariance term. It vanishes at the
flat reference, since a gauge-invariant scalar has zero first
variation there. The two quadratic variations give the second
term by Wick contraction. For the mean, the trace is the quadratic
variation. The extra terms pairing \(DG\) with the quadratic pivot
increment, the odd cubic tilt, or the linear old-activity increment
use the vanishing of their gauge-equivariant means at flat data.

%
For the covariance, first work on the symmetric product
window at the matched reference. Write \(\mathbb E_0\) for
the normalized restricted Gaussian, \(P\) for its tilt
and \(A_i=G_i-\mathbb E_0G_i\). Put \(\delta=\|P\|_\infty\).
Expanding the normalized expectation at \(P=0\) gives
\(\mathbb E_P H=\mathbb E_0 H+
\operatorname{Cov}_0(H,P)+O(\delta^2\mathbb E_0|H-\mathbb E_0H|)\).
Indeed \(\mathbb E_0 e^P\ge e^{-\delta}\) for this real tilt,
and the linear term of its inverse subtracts \(\mathbb E_0P\).
Apply the identity to \(A_1A_2\) and then to each \(A_i\).
Subtracting the shifted means gives
\[
 \operatorname{Cov}_P(G_1,G_2)
 =\operatorname{Cov}_0(G_1,G_2)
   +\operatorname{Cov}_0(A_1A_2,P)
   -\operatorname{Cov}_0(G_1,P)\operatorname{Cov}_0(G_2,P)
   +O(\delta^2\|A_1\|_2\|A_2\|_2).
\]
The norms here are under the restricted Gaussian.
The product of shifted means has the same quadratic
size in \(\delta\).

The even part of \(A_i\) is \(O(g^2\bar\mu)\)
and the odd part is \(O(g\mu_2\kappa_i+g^3\bar\mu)\),
in the Gaussian norms, with the fixed logarithmic
factors from the window. The two same-parity products
pair only with \(P^+\), whose size is \(O(g^2+\theta_j^{\mathrm E})\)
with those factors. The mixed products pair only
with \(P^-\), of size \(O(g+\theta_j^{\mathrm E})\).
This gives the two displayed orders for the linear
tilt correction. Since \(\theta_j^{\mathrm E}\) is small,
\(\delta^2\) is bounded by a fixed logarithmic factor
times \(g^2+\theta_j^{\mathrm E}\). The quadratic
normalization terms have the same orders.
In the untilted covariance, the extra quadratic
pivot term is multiplied by the first derivative of
the readout. Its pairing with the main quadratic
term, and the pairing of a linear term with a cubic
term, cost \(O(g^4\mu_2\bar\mu(\kappa_1+\kappa_2))\).
The remaining higher-order pairings cost
\(O(g^6\bar\mu^2)\), with the logarithmic factors
included. In passing to the full Gaussian apply
Cauchy--Schwarz on the window complement to the
products and to each mean. This gives the
\(\sqrt{t_j}\) error in the comparison.

For a common pair of pullbacks the exterior density
ratio in the normalized comparison is bounded by
\(e^{\pm\Xi_j}\). The bound for the covariance from
the product of two centered suprema therefore
costs \(C(e^{\Xi_j}-1)\prod_i\|G_i-G_i(V_*)\|_\infty\).
The Taylor bounds on the window give
\(\|G_i-G_i(V_*)\|_\infty\le C(\log x)^C s_i\).
For \(\Xi_j\le1\) this is the stated exterior term.
The bound applies to the positive exterior mixture
relative to the same reference density.
For the first two means the argument uses fixed-radius
analytic tubes in the retained datum \(V_j\) and the
corresponding bounds for the mean on those tubes.
For the regular old origins the required joint chart
is obtained by combining the retained field and the
solved free field before taking the old averages.
In the matched local recipe write \(W=e^{iA}\),
let \(p_k(W)\) be its background link field, and
use \(e^{i\Phi_{p_k(W)}(u)}p_k(W)\) with the
physical free variable \(u\). At \(u=0\) the
induced datum agrees with the coarser background
by the nested constraint identity. The local
contour and pivot charts bound the logarithm of
the combined field by \(C(\|A\|_\infty+\|u\|_\infty)\).
On a joint disc \(\|A\|_\infty,\|u\|_\infty<c\alpha_k\)
we need the actual maps of the matched recipe to
enter the same window section as in
\eqref{eq:local-formation-OA12}. The constant \(c\)
must leave room for the sum of the two displacements.
Along the retained axial line its actual size is
at most \(C\operatorname{side}(\mathcal B)\kappa_{\mathcal B}\),
which must fit on the smaller disc.

The window cancellation uses this combined field:
form the old averages, use their common critical
window solution, and subtract the vacuum and its
own tent counterterm at each old pointing. Sum at
that pointing before applying the common bracket.
The full pointed quadratic tensor is zero. The
excluded domains give a long-domain tail, with
length at least \(2/\sigma_i-2\) for a window at
relative spacing \(\sigma_i=L^{-(k-i)}\).
The affine and Hölder estimates then give the
pointed remainder with the gain
\(\rho_i^2\sigma_i^{4+\alpha_H}+\rho_i^3\sigma_i^6\)
in units of \(e_i\). Here \(\rho_i\) is the ratio
of the output radius to the old radius. Counting
the \(\sigma_i^{-4}\) pointings in a current-scale
window gives the age weight
\(\omega_{k+1-i}\) of
\eqref{eq:local-formation-OA15}.
Once the actual old readers are identified on
this joint chart, the bound there is uniform in
\((A,u)\) for the same filed families.

Apply Cauchy's estimate in one retained direction
and one physical free coordinate. For a regular
term on \(Y\) it costs
\[
 C e_i\omega_{k+1-i}\alpha_k^{-2}
                    e^{-\lambda d_k(Y)} .
\]
If the free coordinate is outside the full carrier
its derivative is zero. The first free derivative
of the scalar regular sum vanishes at the matched
flat recipe with zero exterior free coordinates,
by constant colour invariance. Integrating the
mixed derivative along the retained axial line
and rooting at the free bond gives
\(C E_0\alpha_k^{-2}
  \operatorname{side}(\mathcal B)\kappa_{\mathcal B}\).
The box side costs a fixed power of \(\log x\).
This estimate concerns the regular part of the
old Wick vector. The vacuum-only remainders,
Wilson lag and boundary factors use their own
origin bounds. A factor with a fixed noncentral
spectator need not have zero free derivative
at the flat value of the varied datum.
The factor \(g^2\alpha^{-2}\) gives the inverse logarithmic
power in \(E_1\). These bounds, paired with the small
first derivative of \(G\), give the product
\(\kappa\kappa_{\mathcal B}\).

The three vectors in the mean correction can be computed
from the actual chart coefficients. Write \(\Phi_V\) for the
pivot graph in \eqref{afs:actual-pivot-chart}, with
\(\mathsf C=D\Phi_V(0)\), and put \(C=C^{(j-1)}\).
For two free directions \(u,v\), the quadratic pivot coefficient
is given by \eqref{agm:constraint-second-derivative}.
In particular, \(D^2\Phi_V(0)[u,v]=-hD^2q_V(0)[\mathsf C u,\mathsf C v]\)
for the pivot right inverse \(h\) of \(Dq_V(0)\).
Write the intrinsic first coefficient of the exponent as
\(P_1^{\mathrm{int}}(y)=\ell[y]+T[y,y,y]/6\), so that
the first-order exponent is \(\widehat g_{j-1}P_1^{\mathrm{int}}(y)\).
Here \(\ell\) is the first log-coarea derivative including Haar
density, and \(T\) is the signed third derivative of the complete
centered action and contour gauge, with the action weight retained
as in \eqref{fec:literal-first-two-vertices}. Write \(\ell_{\rm old}\)
for the first derivative of the actual old exponent after substitution
in the same graph. Take this derivative in the physical free
coordinate \(u=\widehat g_{j-1}y\), with the sign of the old term in
the exponent.

Use first the full Gaussian for these polynomial contractions.
Wick's formula gives the three vectors before multiplication by
\(\widehat g_{j-1}^2\):
\[
 \frac12D^2\Phi_V(0):C,\qquad
 \mathsf C C\bigl(\ell+\tfrac12 T:C\bigr),\qquad
 \mathsf C C\ell_{\rm old}.
\]
The contraction in \(T:C\) is over two of its three indices.
For example, the second identity follows by pairing \(y\) with
one of the three factors of the cubic. Each pairing leaves the
same contraction, hence the coefficient \(1/2\). The other
two identities use only the second Gaussian moment. For a
polynomial vector \(A\), passage back to the real Gaussian
conditioned on the window costs at most
\(C\sqrt{t_j}\|A\|_{L^2(\gamma)}\) when \(t_j\le1/2\).
This follows by subtracting the two normalized expectations and
using Cauchy--Schwarz on the complement of the window. The
polynomial moments are bounded using the same chart jets and
covariance. This tail comparison uses the real Gaussian.

For the homogeneous chart near the trivial datum, compose the
first two contractions with the actual intermediate field
\(p_k(W)\) of \eqref{hpt:actual-output-composition}.
The fixed output radius for that map, the outer chart radius
in \eqref{afs:actual-pivot-chart}, and the precision bounds
in \eqref{haf:actual-precision-carriers} give a common smaller
fixed complex radius for the Wick polynomials. The action family
in \eqref{het:actual-action-family} and the local coarea factors in
\eqref{afs:actual-coarea-density} bound the intrinsic coefficients
on the read coordinates. At the trivial datum each vector is fixed
by all constant colour rotations, so it vanishes. Cauchy's
estimate on this homogeneous datum ball bounds each of the
first two vectors by \(C(\log x)^C\|\log W\|_\infty\) on
a fixed smaller ball, with the polylogarithmic factors from the
box count. The third vector is bounded by the norm of the
summed old first derivative times the pivot and covariance
bounds.

%
%
For the conditional calculation first compare with its
matched local law with zero exterior free coordinates.
The density-ratio and readout comparisons account for
that change. For the Wick tensors of the local law,
keep its box, free and pivot masks, guard and weights
fixed. Compare its retained datum to the trivial datum
in the same chart. On the nonwrapping box choose the
coarse axial logarithm \(A_W\) with
\(\|A_W\|_\infty\le C\,\operatorname{side}(\mathcal B)
  \kappa_{\mathcal B}\)
in coarse units.
Use the line \(e^{itA_W}\) in the common admitted
datum ball. The local pivot jets and coarea derivative
use complete fixed stencils. For the cubic action use
the rooted coefficient bounds of
\eqref{het:actual-action-family}, and for the covariance
use the resolvent identity for the same precision and
mask. They give a holomorphic comparison on a smaller
fixed radius. Expand each contraction by replacing
one coefficient at a time along that line. The fixed
number of contracted indices costs at most a power
of the box count. Both tensors are zero at the
trivial datum since all the fixed data there
preserve constant colour rotations. The size of
\(\mathcal B\) is bounded by a power of \(\log x\),
so the first two vectors of the matched local law
are bounded by \(C(\log x)^C\kappa_{\mathcal B}\).
%
%
For the old regular near terms on this intrinsic
box, take the retained angular datum and the
physical free coordinate \(u\) as separate
analytic arguments. Use the jointly admitted
old-box substitution of \eqref{eq:local-formation-WF3},
with no fixed colour spectators. Keep the
box and its filing fixed. In its smaller ball,
the amplitude bound \eqref{eq:local-formation-OA15}
holds as the retained datum varies through
angular size \(c\alpha_k\), \(k=j-1\), while
\(\sup |u|\le c\alpha_k\).
Here the pointed counterterm is removed
before summing the old terms. The bound
uses the same age sum \(\mathcal E_k\)
on this joint ball.

To take the first derivative of this sum,
put the retained field on the line
\(e^{itA_W}\). For a free bond \(b\), let
\(F_b(t)\) be the \(u_b\)-derivative at
\(u=0\) of the summed regular exponent
in the chart. On the trivial retained datum
this derivative vanishes. Indeed the whole
intrinsic substitution commutes with constant
colour rotations, and their adjoint action
has no invariant covector. Cauchy's estimate
in \(u_b\), followed by the rooted sum,
bounds \(|F_b(t)|\) by \(C\mathcal E_k/\alpha_k\)
on \(|t|\|A_W\|_\infty<c\alpha_k\).
On the half datum ball, the zero at \(t=0\)
then gives
\[
 \|\ell_{\rm old}^{\rm reg}\|_\infty
   \le C\mathcal E_k \alpha_k^{-2}\|A_W\|_\infty .
\]
This takes the derivative in the physical
coordinate. After multiplication by
\(g_{j-1}^2\mathsf C C\), the summable
covariance and pivot rows give the bound
\(C\mathcal E_k(\log x)^{-2q}
\operatorname{side}(\mathcal B)\kappa_{\mathcal B}\)
for the regular part of the third vector.
Pairing it with the derivative of an invariant
readout supplies the second curvature factor,
by the same vanishing derivative at the trivial
field on the readout's own box.
For ordinary spectator-free remainders
with the same joint domain, replace
\(\mathcal E_k\) by \(\sum_{i\le k}\rho_i^2r_i\).
The vacuum-only bound supplies the
amplitude in the same calculation,
and \eqref{eq:local-formation-OA22}
sums these weights along the flow.

For ambient coefficients one must also
count carriers reaching the exterior.
Their exponential tails are the separate
boundary error. Terms with fixed colour
spectators or without the joint activity-radius
domain need their own first derivative bound.

%
%
The quadratic part of the first covariance term is a bilinear form
in the retained angles. Its mean at later steps requires the first
comparison for the carried readout, with its own derivative bounds.
On the admitted conditional kernels of the run, let \(G_i^{[j]}\)
be the propagated conditional mean of the \(i\)-th readout. Write
it as \(M_{i,j}+J_{i,j}\), with \(M_{i,j}\) the background read and
its accumulated second-order traces. The normalized conditional
expectation propagates \(J_{i,j}\) in supremum norm. In the
conditional covariance its mixed contribution is bounded by
\[
 \left|\operatorname{Cov}_j(M,J)\right|
    \le \|J\|_\infty\,\mathbb E_j|M-\mathbb E_jM|.
\]
For the scalar plaquette read on the clean boxes, the mean-value
bound for its background part uses the curvature of the retained
datum. With a fluctuation of size \(g\) times a power of \(\log x\),
the bound is a power of \(\log x\) times \(g\kappa_{\mathcal B}+g^2\).
The trace part already has a factor \(g^2\) and its fixed-radius
derivatives cost only powers of \(\log x\). The corresponding
contribution of two remainders is at most
\(2\|J_1\|_\infty\|J_2\|_\infty\). This is why the iteration
can use bounds on the values of the old remainders.

For the principal half-square term at \(p,q\), put
\(c_j=x_{j-1}^{-1}\mathcal C_{(j)}(p,q)\) at the flat reference.
At a fixed retained background with curvature readouts \(m_p,m_q\),
the part linear in each fluctuation has the flat-matrix term
\(c_j\langle m_p,m_q\rangle\). It has a quadratic zero at
the flat retained datum. In a later step \(l>j\), its leading
second-order trace increases its conditional mean by
\((N^2-1)c_jc_l\). The two symmetric terms in the Hessian of
the bilinear form cancel the \(1/2\) of the trace. This produces
the cross-level term in the principal covariance, in addition to
the single-step term \((N^2-1)c_j^2/2\).
The normalized comparison on the actual run must bound the error
in these transported traces on the common admitted chart. It also
must bound the covariance of the remaining terminal means. The
calculation uses the conditional kernels of the run, with the
standing data carried by them. Conditioning the initial law on a
plain composite block average defines a different law when the
operations have redrawn exterior fields.

%
%
\paragraph{Carrying a scalar plaquette read to a coarser datum.}
On the nested clean boxes let \(T_p(V)\) be the logarithm
of a level-\(k_0\) plaquette read from the background at
level \(k_0+i\), in the common chart. Use at each finer
step the actual datum induced by the coarser background.
The finer constraint fibre is contained in the coarser
one. Restriction and uniqueness of the regular orbit
identify the composed read with the read from the
single level-\((k_0+i)\) background. This applies on
the exact common constraints. When using two different
localization boxes, keep their collar comparison as a
separate error.

The fixed complex Landau tube for that background gives
\(\|T_p\|\le C L^{-2i}\) on the tube, with a constant
independent of \(i\). The real bound at the datum is
\(C L^{-2i}\kappa(V)\). Cauchy's estimate on a fixed
smaller datum ball gives
\(\|DT_p\|+\|D^2T_p\|\le C L^{-2i}\).
Let \(G\) be a scalar read of a fixed family of these
holonomies, gauge invariant on the nonwrapping box
with the other arguments fixed. At the trivial
field constant colour rotations leave the datum
fixed. Their adjoint action on each logarithm has
no invariant linear functional, so \(DG(0)=0\).
Thus in the read chart
\(\|DG(T)\|\le C\mu_2\|T\|\).
For the traced plaquette weight this cancellation
also follows by taking the trace of its traceless
linear logarithm. The two chain-rule terms give
\[
 \begin{split}
  \|D(G\circ T)(V)\|&\le C\mu_2 L^{-4i}\kappa(V),\\
  \|D^2(G\circ T)(V)\|&\le C\mu_2 L^{-4i}.
 \end{split}
\]
Indeed the first part of the second derivative
has two first derivatives of \(T\), while the
part with \(D^2T\) gains the other factor through
\(DG(T)\). With the bounded one-step Gaussian
covariance the trace of this latter part at depth
\(i\), including the fluctuation scale, costs
\(C(\log x)^C g_{k_0+i}^2\mu_2L^{-4i}\kappa(V)\).
It sums geometrically over the clean segment.
These are the derivative estimates
for this scalar readout family. A matrix plaquette
has only the \(L^{-2i}\) estimate for its derivatives.
The covariance kernels and the differentiated
old-origin sums have their own bounds in the
iteration.

\Cref{prop:hgc-actual-cumulant-comparison} proves the mean and
covariance replacement for the specified plaquette readouts of
the flat homogeneous classical-pullback fibre. For the actual
clean run, the bounds just stated use the normalized box law,
the full tilted exponent, exterior control, and uniform derivative
bounds for the carried readouts. These inputs must be verified
on the same chart.

\paragraph{The transported covariance on the fixed datum chart.}
At step $j$ use the local reference with its masks and all
other chart choices fixed. For a level-$k_0$ plaquette $p$ let
$T_p(u;W)$ be its logarithm in the common frame, after the
physical free displacement $u$ and the actual pivot substitution.
Write $m_p=T_p(0;W)$ and $R_{j,p}(W)=D_uT_p(0;W)$.
The physical matrix pairing between two plaquette reads is
\[
 K_{j,pq}(W)=\widehat g_{j-1}^{2}R_{j,p}(W)C^{(j-1)}(W)
                        R_{j,q}(W)^{\mathsf T} .
\]
All transposes here use the real coordinate inner products and
their bilinear continuations on the complex chart. The first
covariance term for the two half-squares is exactly
$\langle m_p,K_{j,pq}(W)m_q\rangle$. The second-order term
in the mean of one half-square is
\[
 \tfrac12\operatorname{tr} K_{j,pp}(W)
   +\tfrac12\widehat g_{j-1}^{2}
       \langle m_p,D_u^2T_p(0;W):C^{(j-1)}(W)\rangle .
\]
The trace is in the colour space. The second summand uses the
composed derivative and therefore includes the quadratic pivot
contribution to the trace. When using the separate bond Taylor
and pivot terms, count that contribution only once.

The fixed complex datum tube and the nested curvature bound
give $\|R_{j,p}\|_{\ell^\infty\to\mathfrak g}\le C\eta_{j-1}^2$
on a smaller tube. To see this, apply the plaquette bound
to the read with the perturbed datum on the tube and use
Cauchy's estimate in a physical free direction, with the
bounded pivot tangent included. In particular this bounds
the row in $\ell^2$. The covariance is uniformly bounded
on the same smaller tube by the precision estimates for
the fixed masks. Therefore $K_{j,pq}$ and any fixed number
of its datum derivatives have bounds of size
$C_m x_{j-1}^{-1}\eta_{j-1}^4$, allowing fixed powers of
$\log x$ from the box and chart constants. This follows by
applying Cauchy's estimate to the bounded product on the
larger common tube. The scalar half-square also has the
corresponding fixed-order derivative bound without the
factor $x_{j-1}^{-1}$, by applying Cauchy's estimate to the
bilinear square of the holomorphic logarithm. These bounds
use fixed readsets and the admitted common datum radius.

Along the coarse axial line, subtracting $K_{j,pq}(\mathbf 1)$
costs this same bound times the size of the coarse potential.
That size is bounded by the box side times its curvature.
The second summand of the mean trace contains $m_p$ itself
and so also gains the small curvature of the datum. At the
flat datum first compare the box covariance with the full
flat covariance restricted to the read set, at the same
relative weight $\tau_{j-1}$. The Schur correction in
\hyperref[f2:flat-box-marginal]{the flat marginal comparison}
bounds this discrepancy on the
read set by $Ce^{-c\mathsf D}$, where $\mathsf D$ is its
distance from the edge in the box units. The $\ell^1$
bounds on the two readout rows give a cost at most
$C x_{j-1}^{-1}\eta_{j-1}^4 e^{-c\mathsf D}$ in the
physical pairing. With this marginal the local read gives
the local unit-weight slice times $x_{j-1}^{-1}$ in each
colour, since $\widehat g_{j-1}^2\tau_{j-1}^{-1}=x_{j-1}^{-1}$.
Then compare the local and full readout rows by
\eqref{f2:local-flat-slice-error}. Thus the difference from
$c_j$ times the identity in colour is bounded by
\[
 C(\log x)^C x_{j-1}^{-1}\eta_{j-1}^4
 \big(\operatorname{side}(\mathcal B)\kappa_{\mathcal B}
     +e^{-c\mathsf D}+e^{-cd_{j-1,p}}+e^{-cd_{j-1,q}}\big).
\]
In the curved one-line covariance this bound is multiplied by
$|m_p||m_q|$. Its part coming from $K_{j,pq}(W)-K_{j,pq}(\mathbf 1)$
has zero quadratic jet at the flat datum, since both $m$'s
and the matrix difference vanish there. The flat matrix error
can contribute to the quadratic jet, with the exponential
losses in the displayed bound and the two factors from the
derivatives of the plaquette reads.

In the later mean step
the derivative of a trace already carries a factor of
$x_{j-1}^{-1}$. Taking another Gaussian trace at level $l>j$
costs at most $C(\log x)^C x_{l-1}^{-1}x_{j-1}^{-1}\eta_{j-1}^4$
for the mean. The sum of the $\eta_{j-1}^4$ factors is
geometric. The number of later steps is $O(\log x)$, so
these bounds retain a power of $x^{-1}$ beyond the first
trace. The principal quadratic term of the transported
covariance instead contributes the cross-level contraction
described above.

On an admitted clean terminal datum the background
half-square contributes at most
$C\|\phi_i\|_\infty x^{-1}\mathsf r^4 p_0^2 L^{-4\mathsf n}$
to the mean after summing its plaquettes. Here $\mathsf n$
is the number of transported steps and $\mathsf r$ is the
read box size at the initial scale. This follows from the
terminal curvature bound $Cp_0x^{-1/2}$ and the area factor
$L^{-2\mathsf n}$ in the nested read. Subtract the
deterministic reference mean
\[
 c_i=\frac{N^2-1}{2}\sum_p\phi_i(x(p))
       \sum_j x_{j-1}^{-1}\mathcal C_{(j)}(p,p).
\]
The trace differences above sum in $j$ by the same geometric
bound. The remaining one-step errors propagate by the
supremum norm. Thus in using total covariance on the clean
terminal law the covariance of the means costs at most
$2\|\mathbb E(G_1\mid W)-c_1\|_\infty
  \|\mathbb E(G_2\mid W)-c_2\|_\infty$, with the conditioning
index enlarged to include any class variable. The estimates
must be uniform over that index on the clean set. In the
generating-function argument the average is taken before
the logarithm, and this separate covariance estimate is
unnecessary.

\paragraph{The linear readout and the free slice.}
At level \(q=j-1\), write \(H_q^{\square}(p)\) for the
linearized curvature read of the localized problem for
\(p\), with its input row extended by zero to the
whole level-\(q\) torus. Write \(H_q(p)\) for the
linearized read of the regular critical problem on
that torus with every level-\(q\) bond prescribed.
Use the same averaging scheme in both problems.
Let \(d_{q,p}\) be the distance in level-\(q\) spacings
to the boundary of the localization, with the full
read stencil and the fixed collar width included
in measuring the margin.

For a real coarse \(1\)-form \(A\), take
\(V(t)(b)=\exp(it A(b))\) on the whole level-\(q\)
lattice. Its test holonomies have defect at most
\(C|t|\|A\|_\infty\), so the homogeneous torus
problem and the prescribed joined localization are
admitted for small \(|t|\|A\|_\infty\). Use the
restriction of \(V(t)\) on common rows and its
prescribed face lift on the new collar rows. By
\eqref{mb:real-collar-datum}, the discrepancy from
the induced datum is bounded by \(C|t|\|A\|_\infty\).
The weighted response \eqref{mb:real-collar-norm}
therefore bounds the difference of the plaquette
defects in the comparison frame by
\(C\eta_q^2|t|\|A\|_\infty e^{-cd_{q,p}}\).
The comparison can also be made in the ancestor
frame of the plaquette. The tower transport and the
normalized node factor equal one at the trivial field
and are continuous on the datum line. Consequently
their conjugations have identity linear part on a
defect vanishing at \(t=0\). Dividing by \(|t|\)
and letting \(t\) tend to zero gives
\[
 \bigl|(H_q^{\square}(p)-H_q(p))A\bigr|
   \le C\eta_q^2 e^{-cd_{q,p}}\|A\|_\infty .
\]
For one colour this is the \(\ell^1\) norm bound of
the row, by duality. In the same norm the two rows
are bounded by \(C\eta_q^2\), from the flat response
bound on the torus and the corresponding bound
on the localized datum ball. The constants are
independent of volume and depth. Here we differentiate
a real path with a fixed determining sequence in each
problem. The argument uses the actual full read domain
and gives no derivative estimate for changing history
labels or for an old activity.

Let \(\Gamma_j^0\) be the covariance of the complete
flat fluctuation law on the free bonds of the
level-\((j-1)\) torus in the step's contour gauge, at
unit relative Wilson weight. In the actual frame coordinates
the weight is \(\tau_{j-1}\). On curvature reads the actual
flat covariance is \(\tau_{j-1}^{-1}\) times the one at unit
weight, by the constrained curvature form and the separate
Gaussian integration along gauge directions. Consequently
\(\widehat g_{j-1}^2\tau_{j-1}^{-1}=x_{j-1}^{-1}\) is the
physical covariance factor multiplying the unit-weight slice.
Use one colour and put
\[
 R_{j,p}^{\square}=H_{j-1}^{\square}(p)\mathsf C^0,
 \qquad
 R_{j,p}=H_{j-1}(p)\mathsf C^0 .
\]
Here \(\mathsf C^0\) is the flat pivot inclusion of the
actual step. Its fixed finite range and bounded row
sums preserve the bounds for the readout rows.
All free bonds used by \(R_{j,p}^{\square}\) lie in
\(\mathcal S_j\). Thus after removing the windows its
covariance under the flat read-set marginal is exactly
\[
 \mathcal C_j^{\square}(p,p')
     =R_{j,p}^{\square}\Gamma_j^0(R_{j,p'}^{\square})^{\mathsf T}.
\]
The whole-torus covariance is
\(\mathcal C_{(j)}(p,p')
 =R_{j,p}\Gamma_j^0(R_{j,p'})^{\mathsf T}\), by the
identification of the curvature covariance in the
contour and tree gauges. It is this whole-torus
operator that is the orthogonal slice. Since
\(\|\Gamma_j^0\|_{2\to2}\le C\), subtract the two
products one row at a time and use the row bounds
to obtain
\begin{equation}
 \bigl|\mathcal C_j^{\square}(p,p')-\mathcal C_{(j)}(p,p')\bigr|
   \le C\eta_{j-1}^4
       \bigl(e^{-cd_{j-1,p}}+e^{-cd_{j-1,p'}}\bigr).
 \label{f2:local-flat-slice-error}
\end{equation}
The Gaussian used in both products is the same
complete flat Gaussian. The comparison of an
interacting box marginal with its flat restriction
is the preceding separate estimate.

\paragraph{The error on the generating function disc.}
Let \(d_*=\min_{p,j}d_{j-1,p}\) for the plaquettes
and levels in the segment. Since \(x_{j-1}\asymp x\)
and \(\eta_{j-1}=L^{-(j-1-k_0)}\eta_{k_0}\), the
difference of the covariances of the two sums of
independent slices is bounded entrywise by
\(Cx^{-1}\eta_{k_0}^4e^{-cd_*}\).
To estimate the determinant without a factor counting
the total number of covariance entries, use the same
standard Gaussian on the direct sum of the flat free
bond spaces for both sums. For the localized read
its row is
\[
 u_p=\bigoplus_{j=k_0+1}^{k_1}
        x_{j-1}^{-1/2} R_{j,p}^{\square}(\Gamma_j^0)^{1/2}.
\]
Let \(v_p\) denote the row with \(R_{j,p}\) in its
place. The geometric sum of the squared row bounds gives
\[
 \|u_p\|+\|v_p\|\le Cx^{-1/2}\eta_{k_0}^2,\qquad
 \|u_p-v_p\|\le Cx^{-1/2}\eta_{k_0}^2e^{-cd_*}.
\]
For \(z_p=\sum_a\lambda_a\phi_a(x(p))\), the
matrices of the quadratic Gaussian exponents are
\(K_{\square}=\sum_p z_p u_p^{\mathsf T}u_p\) and
\(K=\sum_p z_p v_p^{\mathsf T}v_p\), with the factor
\(1/2\) outside the quadratic form. The trace norm of
a rank-one term \(u^{\mathsf T}v\) is \(\|u\|\|v\|\).
Hence
\[
 \|K_{\square}-K\|_1
    \le C x^{-1}\eta_{k_0}^4 e^{-cd_*}\sum_p|z_p|,
 \qquad
 \|K_{\square}\|_1+\|K\|_1
    \le C x^{-1}\eta_{k_0}^4\sum_p|z_p|.
\]
For \(n\) profiles, the radii
\[
 |\lambda_a|\le
       \frac{c_* x}{n\mathsf r^4p_0^2\|\phi_a\|_\infty}
\]
and the count \(C\mathsf r^4\eta_{k_0}^{-4}\) of
plaquettes per profile give
\(x^{-1}\eta_{k_0}^4\sum_p|z_p|\le Cc_*p_0^{-2}\).
Zero profiles may be omitted. Take \(c_*\) small
enough that the two operator norms are at most \(1/2\).
The logarithms of the Gaussian generating functions
are then given by the convergent trace series, each
with the common colour factor \(3/2\) for
\(\mathrm{SU}(2)\). Telescoping each difference of
matrix powers and using the trace norm on its changed
factor bounds the difference of these logarithms by
\[
                    Cc_*p_0^{-2}e^{-cd_*}.
\]
The same bound, with another fixed constant, holds
for the generating functions. Cauchy's inequality
multiplies this cost by the product of the inverse
radii for the insertions being differentiated.
For the actual test box and the later boxes, the
native margins give
\(d_*\ge c_0\min\{\mathsf D,\check R_{k_0}\}-C_0\).
Taking \(\mathsf D=A\log x\) in this range therefore
makes the cost smaller than any prescribed inverse
power of \(\log x\) after a fixed number of
differentiations. The comparison allows the exact
whole-torus free projection and its ring sums to be
used for the Gaussian target. The normalized nonlinear
conditional comparison and the removal of the
fluctuation windows enter separately.

For the free calculation, write \(F_j\) for the background curvature
at level \(j\), with covariance \(G_j=C^{\mathrm{bg}}_{K,K-j}\).
The nested projections give \(F_{j-1}=F_j+Z_j\), where \(Z_j\)
is independent of \(F_j\) and has covariance \(\mathcal C_{(j)}\).
The recursion is evaluated at the zero terminal field \(F_{k_1}=0\).
In this Gaussian law \(F_j\) is the sum of the slices with
\(j<l\le k_1\), and its covariance is
\(G_j^\circ=G_j-G_{k_1}\). For one colour, the covariance of
the two half-squares at \(p,q\), conditional on \(F_j\), is
\[
 \frac12\mathcal C_{(j)}(p,q)^2
       +F_j(p)F_j(q)\mathcal C_{(j)}(p,q).
\]
The second term comes from the two terms linear in \(Z_j\) when the
half-squares are expanded at \(F_j\). Averaging over \(F_j\)
with the terminal field fixed gives
\[
 \frac12\left(G_{j-1}^\circ(p,q)^2-G_j^\circ(p,q)^2\right).
\]
Thus the free contributions telescope to
\(\tfrac12(G_{k_0}-G_{k_1})(p,q)^2\).
For separated supports \(|G_{k_0}(p,q)|\le Ca_K^4\mathsf d^{-4}\),
and the coarse covariance bound gives
\(|G_{k_1}(p,q)|\le CL^{-4k_1}\).
The error in replacing the square by \(G_{k_0}(p,q)^2\)
includes a product of these two kernels and the square of the
coarse kernel. Each support has at most \(C(\mathsf d/a_K)^4\)
plaquettes. After smearing, the error is therefore at most
\[
 C\left[(\hat{\mathsf d}L^{-\mathsf n})^4
         +(\hat{\mathsf d}L^{-\mathsf n})^8\right]
                  \|f_1\|_\infty\|f_2\|_\infty .
\]
This tends to zero at the stated choice of \(\mathsf n\).

Three facts are used here. The source arrives at level $k_0$ in the form $\zeta'_{k_0}\tilde{f}_i$ ((b)). Next, $a_p\approx\frac12|F(p)|^2$, where $F(p)$, an element of the Lie algebra, is the plaquette angle at $p$ of $U_{k_0,\square}(V)$ (the same angle is written $Y_p$ in \cref{thm:continuum-main}(c)(5)). Finally, half-squares of Gaussian variables have covariance $\frac12(\text{covariance})^2$ for each colour. With the two interacting comparisons and the earlier reductions, we obtain

\begin{gather*}
\FitMath{\Big|\langle\mathcal{O}(f_1);\mathcal{O}(f_2)\rangle-\zeta_{k_0}'^{\,2}\,\tfrac{N^2-1}2\,x^{-2}\,\mathcal{K}_{K,s}(\tilde{f}_1,\tilde{f}_2)\Big|\le o(1)\,x^{-2}\prod\|f_i\|_\infty.}
\end{gather*}

The formula at level $k_0$, combined with (a) and (b), gives the same estimate for $\zeta_{k_0}'^{\,2}\langle X_1;X_2\rangle_{\mathsf{P}}$ in the average over $\mathfrak P$.
\emph{Positions.} The cumulants are invariant under the translations of $\Lambda_{K,\ell}$. A translation $v\in\Lambda_{s,\ell}$ carries blocks onto blocks, which finer translations do not do. Hence

\begin{gather*}
\FitMath{C^{\mathrm{bg}}_{K,s}(p+v,q+v)=C^{\mathrm{bg}}_{K,s}(p,q),} \\[3pt]
\FitMath{\mathcal{K}_{K,s}(f,h)=\frac1{|\Lambda_{s,\ell}|}\sum_{v\in\Lambda_{s,\ell}}\mathcal{K}_{K,s}\big(f(\cdot-v),h(\cdot-v)\big):}
\end{gather*}

since both sides of the law are invariant, it is enough to make the comparison in the average over such positions.

(d) \emph{Wild histories.} For two points, the relevant event is that a large-field region touches the supports at the single level $k_1$ ((c1)). The weight of these histories is $\le C(M\check{R}_{k_1})^4e^{-c\,p_0(g_{k_1})}$, and no union over the levels is taken. In this bound the prefactor counts the sites of level $k_1$ in one cell of that level, and $O(1)$ cells touch the supports. For four points \cref{thm:continuum-main}(c)(5), the unclean histories are the ones that matter. At each level $j$ of the segment of \cref{thm:continuum-main}(c)(5), $k_0-O(L\check{R}_{k_0})\le j\le k_0+r_4$, the number of cells of level $j$ near the supports, that is, within $C\hat{\mathsf{d}}$ cells of level $k_0$, is $\le C\hat{\mathsf{d}}^4L^{4\max\{k_0-j,0\}}$; altogether there are $\le e^{CL\check{R}_{k_0}}$ such cells. A component on which $\mathbf{R}$ acts at a level $j$ of the segment contains a component that is live at level $j-1$. The one-cell bounds of \cref{candidate:sources,ch:all-state-source-transfer} give $\tilde\nu(\mathfrak{c}\subset\mathcal{Z}_j)\le Ce^{-c\,p_0(g_j)}$, and the same bound for creation at step $j$. Moreover $p_0(g_j)\ge c\,p_0(g(\mathsf{d}))$, because $x_j\ge x^{1/2}$. Hence $\tilde\nu(\text{unclean})\le e^{CL\check{R}_{k_0}-c\,p_0(g(\mathsf{d}))}$, and this is small because $\mathfrak{p}_0\gg r$. Even here these one-cell bounds must hold uniformly in the volume (prolonged run, translation covariance up to the level, first moment of the prepaid end mark: \cref{candidate:sources,ch:all-state-source-transfer}). The reason is that $T_{\ell_\star}$ contains $(\ell_\star/Ma_s)^4$ cubes of level $k_0$, and this number is unbounded.

\textbf{What must be overcome.} The natural size of a factor is $x^{-1}\|f_i\|_\infty$. The bound on $X_i$ is larger by $p_0^2\,\hat{\mathsf{d}}^4=(\log x)^{C\mathfrak{p}_0}$, and the irrelevant terms are controlled by sup bounds only. This is the reason for $q\gg\mathfrak{p}_0$ and for the super-polynomial smallness in (a), (d). The effective action is that of Ba{\l}aban (Comm. Math. Phys. 109 (1987) 249; 122 (1989) 175, 355); the sources and the extraction are added to it.

\paragraph*{The setting and the objects}

In \cref{sec:short-distance-the-short-distance-two-point-law-and-its-asymptotic-freedom} the two-point law is stated for two bumps $f,h$ whose size and mutual separation are both $\approx\mathsf d$. They are placed on a lattice of spacing $a_K\ll\mathsf d$ inside a volume of side $\ell\gg\mathsf d$. Isolating the scale $\mathsf d$ has two aspects. \emph{From below}: $K-s(\mathsf d)$ renormalisation steps integrate out the scales between $a_K$ and $\mathsf d$. In the leading term these scales survive only through two numbers of the run: $g(\mathsf d)$, the coupling, and $\zeta_K(\mathsf d)$, the normalisation. \emph{From above}: comparing $T_\ell$ with a single fixed torus eliminates the scales from $\mathsf d$ up to $\ell$, together with the volume itself. This comparison costs $O(\mathsf d^2)$, which is of lower order than the main term $\approx(\ln\frac1{\mathsf d})^{-2}$. This part defines the two numbers, bounds them uniformly in the cutoff and in the bare coupling, and carries out the comparison of volumes. The Gaussian main term is extracted in \cref{sec:short-distance-the-short-distance-two-point-law-and-its-asymptotic-freedom}, at the level $K-s(\mathsf d)$ on the fixed torus. There the cumulant is written in terms of a law on resummed histories and block fields (the history-resummation part of \hyperref[hyp:F2]{the conditional-law hypothesis}). Resummed histories are classes of histories that become equal once their entries inside healed components are deleted. The large-field operation $\mathbf R$ sums over the outcomes inside the components it heals, so an exponent is attached to such a class rather than to an individual history; write $\mathsf h$ for this history class.

\paragraph{The law passed to the limits}
\label{sec:short-distance-the-law-passed-to-the-limits}

\begin{corollary}[{\cref{thm:continuum-main}(c)(4)}]\label{cor:short-distance-7-4} Let $S$ be a continuum limit, either on $T_\ell$ of \cref{thm:continuum-main}(a)(ii) or on $\mathbb{R}^4$ of \cref{thm:continuum-main}(a)(iii). Let $\mathsf{d}\le\mathsf{d}_0$, and let $(f,h)\in\mathfrak{T}_{\mathbb{E}}(\mathsf{d};\vartheta)$ be a pair of class $C^1$. Then the conclusion of \cref{prop:short-distance-7-1}(ii) is valid with

\begin{gather*}
\FitMath{(1-\bar\varepsilon)\big(1-\bar\eta(\mathsf{d})\big)^4\le1+\eta\le(1+\bar\varepsilon)\big(1+\bar\eta(\mathsf{d})\big)^4,}
\end{gather*}

in which $\bar\eta$ and $\bar\varepsilon_\ell$ denote the functions of \cref{prop:short-distance-7-1}. Their values lie in $[0,\frac13]$ and $[0,\frac12]$ respectively, and both tend to $0$ with $\mathsf d$. The quantity $\bar\eta(\mathsf{d})$ is a common bound for $|\zeta_K(\mathsf{d})-1|$ and for $|\frac{11N^2}{12\pi^2}g(\mathsf{d})^2\ln\frac1{\mathsf{d}}-1|$. On $T_\ell$ one has $\bar\varepsilon=\bar\varepsilon_\ell(\mathsf{d})$, and on $\mathbb{R}^4$ one has $\bar\varepsilon=\frac12$. Consequently $S_2(f,h)>0$, and on $\mathbb{R}^4$ we have $\dim\mathcal{H}(S)\ge2$.
\end{corollary}

\textbf{Proof outline.} The argument follows from the uniformity in \cref{prop:short-distance-7-1}(i). In the lattice law, the only quantities that depend on $(K,\ell,\beta)$ are $\zeta_K(\mathsf{d})=1+\eta''$, $g(\mathsf{d})^2\propto(1+\eta')/\ln(1/\mathsf{d})$ and $\varepsilon$. Therefore

\begin{gather*}
\FitMath{\frac{S^T_{2;K,\ell,\beta}(f,h)\,(\ln(1/\mathsf{d}))^2}{\frac{3(N^2-1)}{\pi^4}(\frac{12\pi^2}{11N^2})^2\mathcal{K}_{\mathrm{free}}(f,h)}=(1+\eta'')^2(1+\eta')^2(1+\varepsilon)}
\end{gather*}

takes its value in the closed interval of the display above, for every datum. For fixed $\mathsf{d}$, we have $K_j-s(\mathsf{d})\ge n_0$ once $j$ is large. Every pair on $\mathbb{R}^4$ is contained in some ball of radius $5\mathsf{d}$. Hence, once transported to $T_{\ell_j}$, it belongs to the class and keeps the same free kernel. A closed condition is preserved in the limit. The bound on $\varepsilon$ that clause (c)(4) records on $\mathbb{R}^4$ is the only one in \cref{prop:short-distance-7-1}(i) that is free of any side, namely $|\varepsilon|\le\frac12$.

\textbf{The Hilbert space.} Let $f_{\mathsf{d}}$ be a $C^1$ bump whose support is the ball $B(2\mathsf{d}e_4,\mathsf{d})$. Then $(\theta f_{\mathsf{d}},f_{\mathsf{d}})$ is a test pair. Its shape $\vartheta_1$ is fixed and does not depend on $\mathsf{d}$, because the bumps are dilates of a single profile. When $\vartheta>\vartheta_1$, the pair may fail to belong to $\mathfrak{T}_{\mathbb{R}^4}(\mathsf{d};\vartheta)$. We take the ceiling $\gamma_W(\vartheta)$ smaller than the ceiling under which \cref{prop:short-distance-7-1}(i) is valid at the shape $\vartheta_1$. That theorem is applied there with the witness depth $s_1(\cdot)$, which is built from the room constant $C_{W,0}$ of that shape, and with the thresholds of the same shape. Once $\mathsf{d}$ is below its threshold, the three inequalities $\zeta'_{K-s_1(\mathsf{d})}\ge\frac23$, $1+\varepsilon\ge\frac12$ and $g_{K-s_1(\mathsf{d})}^{-2}\le\bar{x}_{s_1(\mathsf{d})}$ bound $S^T_{2;K_j,\ell_j,\beta_j}(\theta f_{\mathsf{d}},f_{\mathsf{d}})$ from below by a positive number that does not depend on $j$. Consequently \cref{ch:euclidean-hilbert-bridge} gives

\begin{gather*}
\FitMath{\big\|[(f_{\mathsf{d}})]\big\|^2=\sigma_S(\theta f_{\mathsf{d}},f_{\mathsf{d}})=S_2(\theta f_{\mathsf{d}},f_{\mathsf{d}})>0,} \\[3pt]
\FitMath{\langle\Omega,[(f_{\mathsf{d}})]\rangle=\sigma_S((f_{\mathsf{d}}))=0.}
\end{gather*}


\textbf{Periodic Gibbs states.} The expansion concerns torus measures. The law is carried over to these states by \eqref{eq:two-point-extraction-partial-ET7}, under its common-normalizer and reference-comparison hypotheses. That lemma relies on the periodic-state domination in \eqref{eq:all-state-source-transfer-ST3}, according to which the chessboard dominates the periodic Gibbs states. Precisely: let $F$ be a function of a single board cell of side $\mathsf{l}=O(1)$, and on $T_\ell$ let the chessboard seminorm be denoted $\|\cdot\|_{\mathrm{ch},\ell}$. Then

\begin{gather*}
\FitMath{|\nu(F)|\le\limsup_{\ell\to\infty}\|F\|_{\mathrm{ch},\ell}.}
\end{gather*}

This is a separate lemma. It is not the chessboard estimate, whose subject is torus measures. Its proof, taken in the form that involves all mirror copies, combines four ingredients: the reflections of the cells leave $\nu$ invariant \cref{ch:continuum-compactness,ch:all-mirror-periodic}; Jensen's inequality is used together with the Gibbs property; the boundary condition is changed at the cost of a surface term, which is weighed against a volume; and finally the seminorm enters.

\textbf{Ergodic components.} At $(K_j,\beta_j)$, each periodic Gibbs state obeys the law. A closed condition that all of these states satisfy is inherited by the components \cref{ch:phase-interface}. This yields \cref{thm:phase-main}(3)(ii): for a volume-first limit, each of its ergodic components satisfies the law in the form it takes on $\mathbb{R}^4$, with $\mathsf{d}_0$ and $\bar\eta$ unchanged.


% END MODULE chapters/08_short_distance.tex

% BEGIN MODULE chapters/two_point_extraction_partial_v6.tex
\subsection{Covariance bounds for interacting extraction}
\label{ch:two-point-extraction-tests}

The finite insertion, renewed pointing, window formation and linear
bootstrap sections derive the ordinary regular first response under
their explicit unmarked, background, domain and slope inputs.
They retain the potential normalization through local integration
and give the near bound \(C\delta^2s\) on complete nonwrapping fields
of size \(\delta\). The remainder below must also include the
separately bounded marked, remote and boundary contributions.
None of these estimates proves an interacting Gaussian law.

\subsubsection{Covariance bookkeeping with the retained source normalization}

A source-disc bound alone cannot replace the potential estimate.
Indeed \(H(z,B)=(z/\rho)(B/\alpha)^4\) is bounded by one on its
product unit balls, but its source derivative at \(B=\epsilon\)
is \(\rho^{-1}(\epsilon/\alpha)^4\). Relative to
\(\alpha^2(\epsilon/\alpha)^3\), this is
\(x(\log x)^{p-3q}/\rho\) when
\(\alpha=x^{-1/2}(\log x)^q\) and
\(\epsilon=x^{-1/2}(\log x)^p\); it diverges at fixed \(\rho\).
The finite construction supplies its response by differentiating
the inserted potential, with the stated unmarked hypotheses.

Here is the corresponding covariance bookkeeping. Let $\lambda$ be any
probability measure, put $n_i=\|f_i\|_\infty$, fix $p>0$, and suppose, uniformly in
$x\to\infty$,
\begin{equation}
 \begin{split}
 \|\mathsf O_i\|_\infty
       &\le C x^{-1}r^4(\log x)^{2p}n_i,\\
 \|\mathsf I_i\|_\infty
       &\le C x^{-1}r^4
                    [(\log x)^{3p-q}+x^{-a}]n_i,\qquad a>0,\\
 r&\le C_r(\log x)^{2p},\qquad q>21p.
 \end{split}
 \label{eq:two-point-extraction-partial-ET3}
\end{equation}
Then
\begin{equation}
 \begin{split}
 &|\operatorname{Cov}_\lambda(\mathsf O_1,\mathsf I_2)|
  +|\operatorname{Cov}_\lambda(\mathsf I_1,\mathsf O_2)|
  +|\operatorname{Cov}_\lambda(\mathsf I_1,\mathsf I_2)|\\
 &\hspace{1cm}\le
 C'x^{-2}\big[(\log x)^{21p-q}
                    +(\log x)^{18p}x^{-a}\big]n_1n_2
             =o(x^{-2}n_1n_2).
 \end{split}
 \label{eq:two-point-extraction-partial-ET4}
\end{equation}
To prove it use $|\operatorname{Cov}(A,B)|\le2\|A\|_\infty\|B\|_\infty$.
The product of the two support counts costs $r^8\le C(\log x)^{16p}$,
and the a priori insertion costs another $(\log x)^{2p}$.
For large $x$ the bracket defining $\mathsf I_i$ is at most one, so its
square is absorbed by itself. This gives \eqref{eq:two-point-extraction-partial-ET4}. The numerical inequality
$q>21p$ is a sufficient bound for this particular sup-norm step; it does
not claim to be the optimal requirement for the full extraction.


\subsubsection{A scalar transfer estimate with a genuinely common normalizer}

Let $\rho,\tau$ be two probability measures with real bounded observables
$\mathcal O(f)$ linear in a real compactly supported Lipschitz function
$f$. They may be measures on different configuration spaces. Write
\[
 D(f,h)=\operatorname{Cov}_\rho(\mathcal O(f),\mathcal O(h))
          -\operatorname{Cov}_\tau(\mathcal O(f),\mathcal O(h)).
\]
Assume that $D$ is invariant under joint translations by $a\mathbb Z^4$.
Set
\[
 \|f\|_\sharp=\max(\|f\|_\infty,c_M\operatorname{Lip}f)
\]
with $c_M$ fixed.

Suppose there are $r_0>0$ and $Q\ge1$ such that, for every pair
$F_1,F_2$ supported in one common unit ball, there is one holomorphic
function $E(z_1,z_2)$, $E(0)=0$, serving both measures, for which
\begin{equation}
 \left|e^{-E(z)}
       \omega\!\left(e^{z_1\mathcal O(F_1)+z_2\mathcal O(F_2)}\right)
       \right|\le Q,\qquad \omega\in\{\rho,\tau\},
 \label{eq:two-point-extraction-partial-ET5}
\end{equation}
on $|z_1|\|F_1\|_\sharp+|z_2|\|F_2\|_\sharp<r_0$.
The constants are uniform in the ball position and in the measures.
The tests in this hypothesis may overlap and have either sign.

Put $c_i=\|F_i\|_\sharp/r_0$ and denote the normalized holomorphic
function in \eqref{eq:two-point-extraction-partial-ET5} by $h_\omega$. It has $h_\omega(0)=1$.
Schwarz's lemma along rays gives
$|h_\omega(z)-1|\le(Q+1)\sum_i c_i|z_i|$.
On the bidisc with radii $[4(Q+1)c_i]^{-1}$ its principal logarithm has
modulus at most $\log2$. Thus Cauchy's formula gives
\begin{equation}
 |D(F_1,F_2)|
 \le 32\log2\,(Q+1)^2r_0^{-2}
                         \|F_1\|_\sharp\|F_2\|_\sharp.
 \label{eq:two-point-extraction-partial-ET6}
\end{equation}
The derivatives of $E$ cancel because it is the same function for both
measures. If an entry is zero the conclusion is immediate. No global
zero-free statement for the original generating functions is needed.

Suppose now $a\le d$ is sufficiently small, and $f,h$ have supports
in one common ball of radius $5d$, with
$\operatorname{Lip}f\le\|f\|_\infty/(\vartheta d)$ and the analogous
condition for $h$. Choose a lattice multiple $r'\in[5d,10d)$ and
$N\ge c d^{-4}$ translates of that ball, with centers in $4r'\mathbb Z^4$,
all lying inside a common unit ball. For each choice of signs
$\varepsilon_v=\pm1$ let
$F^\varepsilon=\sum_v\varepsilon_v\tau_v f$ and
$H^\varepsilon=\sum_v\varepsilon_v\tau_v h$.
The translated balls are disjoint. Since the functions vanish on their
ball boundaries, these signed sums do not increase either supremum or
Lipschitz norm. Independent signs with the same sign used in the two
arguments give
\[
 \mathbb E_\varepsilon D(F^\varepsilon,H^\varepsilon)=N D(f,h).
\]
Apply \eqref{eq:two-point-extraction-partial-ET6} and divide by $N$. The result is
\begin{equation}
 |D(f,h)|\le C(Q,r_0,c_M)\vartheta^{-2}
                         d^2\|f\|_\infty\|h\|_\infty.
 \label{eq:two-point-extraction-partial-ET7}
\end{equation}
The exponent is $2=4-2$: the number of copies gains four powers and
the two Lipschitz norms lose one power each.

If a reference law has a main term bounded below by
$c_\vartheta x^{-2}\|f\|_\infty\|h\|_\infty$, its relative transfer
error is at most $C_\vartheta d^2x^2$. In the actual short-distance
shape class, each nonnegative entry also has the explicit mass lower
bound $\int f\ge\vartheta d^4\|f\|_\infty$, and the two supports are
at distance between $d$ and $6d$. These requirements give
$\iint f(x)h(y)|x-y|^{-8}\,dx\,dy
\ge10^{-8}\vartheta^2\|f\|_\infty\|h\|_\infty$.
The mass requirement is additional to the Lipschitz bound in \eqref{eq:two-point-extraction-partial-ET7}.
At a witness depth with
$x\le\bar x_s=O(s)$ and
$d<L C_W(s)L^{-s}$, where $C_W$ is polylogarithmic in $s$, this tends
to zero uniformly. For a fixed $\vartheta=1/12000$ all constants remain
finite; their dependence on that fixed shape may affect the permitted
coupling and distance thresholds.

The common-normalizer hypothesis in \eqref{eq:two-point-extraction-partial-ET5} is not a consequence of having
unrelated bounded normalizers for the two measures. There are two
different ways of supplying it. If the scalar stability construction
already provides one common cell normalizer for every periodic
infinite-volume Gibbs state at the same datum, then \eqref{eq:two-point-extraction-partial-ET5} holds for any
two such states. This includes a reference state obtained by taking
finite-torus limits, provided its Gaussian estimate was uniform along
arbitrarily large tori for each fixed cutoff and pair. The source
$z_1F_1+z_2F_2$ is a single cell source, even when the supports overlap.
This route needs no comparison with a fixed finite torus.

If, instead, the reference Gaussian estimate has only been proved on
one prescribed finite torus, a comparison of that torus with the
infinite-volume states remains necessary. Compatible local energies
and vacuum-density/floor bounds give one scalar chessboard route to
\eqref{eq:two-point-extraction-partial-ET5}. Bounds valid on each torus with a distance threshold tending
to zero as its side grows do not by themselves furnish an infinite
reference with the needed fixed-pair estimate. In either route, a typed
Ward insertion estimate is not needed for \eqref{eq:two-point-extraction-partial-ET6}--\eqref{eq:two-point-extraction-partial-ET7}.

\paragraph*{A local companion criterion for vacuum-density compatibility}

For the finite-reference route just described, the following criterion
can remove a separate vacuum-pressure comparison assumption. It is not
needed for comparison of infinite-volume states when their common cell
normalizer is already available. It does not follow from upper stability
and a partition floor alone.
Fix $L>1$ and a cutoff $K$, and put $a_j=L^{j-K}$ for $0\le j\le K$.
Take compatible side lengths $\ell$ for which $\ell/a_j$ is integral.
Suppose the vacuum constant has the form
\begin{equation}
 E_{\ell,K}=\sum_{j=0}^K
 \left[(\ell/a_j)^4 c_{j,K}
       +\sum_{z\in(a_j\mathbb Z/\ell\mathbb Z)^4}
           \sum_{D\in\mathcal D_{\ell,j}(z)}v_{\ell,j}(D,z)\right].
 \label{eq:two-point-extraction-partial-ET8}
\end{equation}
The numbers $c_{j,K}$, the coupling sequence and every auxiliary choice
are the same for all $\ell$. There is a translation-covariant infinite
companion $v_{\infty,j}$ and an integer length $m(D)\ge0$. For both
families, assume that the number of domains of length $m$ at a root is
at most $C_D e^{\kappa_D m}$ and
$|v_{\ell,j}(D,z)|,|v_{\infty,j}(D,z)|
 \le C_Ee^{-\kappa_E m}$ with $\kappa_E>\kappa_D$.
The constants are uniform in $j,K,\ell$. Assume that rooted domains
with $m<\ell/(C_0a_j)$ correspond bijectively to companion domains,
with equal energy values. This is an explicit locality and nonwrapping
assumption, with $C_0$ fixed. It is stronger than simply asserting
that the two models use the same formula for their energies.

Absolute convergence defines
$\bar v_j=\sum_{D\in\mathcal D_{\infty,j}(0)}v_{\infty,j}(D,0)$ and
$A_K=\sum_{j=0}^K a_j^{-4}(c_{j,K}+\bar v_j)$.
At each root all terms below the stated threshold cancel in the
difference from the companion. Sum the two tails of the geometric
series with rate $\delta=\kappa_E-\kappa_D>0$. Division by the
physical volume therefore gives
\begin{equation}
 \left|\ell^{-4}E_{\ell,K}-A_K\right|
       \le C\sum_{j=0}^K a_j^{-4}e^{-c\ell/a_j}.
 \label{eq:two-point-extraction-partial-ET9}
\end{equation}
For every fixed $\ell_0>0$ and $\ell\ge\ell_0$, this is bounded by
$C_{\ell_0}e^{-c\ell/2}$, independently of $K$. Indeed, with $b=K-j$,
\begin{equation}
 \sum_{b\ge0}L^{4b}e^{-c\ell L^b}
 \le e^{-c\ell/2}
                \sum_{b\ge0}L^{4b}e^{-(c\ell_0/2)L^b}
 \le C_{\ell_0}e^{-c\ell/2}.
 \label{eq:two-point-extraction-partial-ET10}
\end{equation}
Thus two compatible tori obey the same bound with
$2C_{\ell_0}e^{-c\min(\ell,\ell')/2}$.
The divergent common number $A_K$ need not converge as the cutoff is
removed; it cancels exactly in the comparison.

In a localized source argument, equality of the corresponding
nonwrapping source energies also supplies the same interior-cell
counterterm for all tori. Bounds on the energies crossing the cell
faces and the appropriate partition floor are still needed to obtain
\eqref{eq:two-point-extraction-partial-ET5}. This lemma addresses its vacuum-density datum only. If the
coupling sequence or local normalizer numbers differ between tori,
or if nonwrapping compatibility has not been established for the actual
formation rule, \eqref{eq:two-point-extraction-partial-ET9} has not been proved for that construction.


\paragraph*{A common source normalizer from the reference torus.}
Fix one first-passage datum and let \(\ell_\star\) be the least
compatible side at its reference coupling.
Choose board cells of side \(\mathsf l_c=L^m\) such that a unit ball
can be placed in a cell with a fixed positive face clearance, and
\(\ell_\star/\mathsf l_c\) is even. This is possible with a fixed
\(m\) for the small reference couplings under consideration. Each
compatible side is an integer multiple of \(\ell_\star\).
Translate the two tests jointly by a bare lattice vector to the
interior of the designated cell. Write \(w_z=z_1F_1+z_2F_2\)
and \(N_\ell=(\ell/\mathsf l_c)^4\).
Disseminate this source by reflection and conjugation, using \(\bar z\)
in the odd cells, and call the resulting source \(W_z\).
Its supports have the same face clearance in each cell. The
dissemination of \(e^{\mathcal O(w_z)}\) is \(e^{\mathcal O(W_z)}\);
plaquettes in the support have all their corners in a single cell.
Joining the reflected profiles across faces leaves the Lipschitz
bound unchanged since each profile vanishes near the faces. Treating
even and odd copies separately gives
\(\|W_z\|_\flat\le2\sum_i |z_i|\|F_i\|_\sharp\).
In particular, \(\sum_i|z_i|\|F_i\|_\sharp<\varrho\) is within the
source domain of \cref{prop:source-domains-4-1}.

The numerator of the chessboard seminorm is nonnegative, and
\[
 \|e^{\mathcal O(w_z)}\|_{{\rm ch},\ell}
   =\left(\frac{Z_\ell(W_z)}{Z_\ell(0)}\right)^{1/N_\ell},
 \qquad
 Z_\ell(W_z)\le
   Z_{\ell_\star}(W_z)^{(\ell/\ell_\star)^4}.
\]
To see the second inequality, vary one side at a time, keeping the
other three fixed. Across a cell face there is no source. After
integrating the crossing bonds the Wilson transfer factor is the
product of the gauge-invariant projection and the convolution with
the product of the one-bond class functions
\(c(u)=e^{-\beta(1-\operatorname{tr}(u)/2)}\).
These operators commute. The character coefficients of \(c\) are
nonnegative: expand \(e^{(\beta/2)\chi_{1/2}}\) in powers and decompose
each tensor power of the defining representation. Hence the crossing
factor is positive, with self-adjoint square root \(S\). In a
cross-section there are finitely many bonds, and the convolution is
trace class since the nonnegative character series at the identity
converges. The intervening slab operator with the internal source
has the form \(A=SBS\), and its mirror copy is \(A^*\).
The trace of the shorter period is \(\operatorname{Tr}T\), where
\(T\) is the appropriate integer power of \(AA^*\).
Repeating the period \(b\) times gives \(\operatorname{Tr}T^b\),
at most \((\operatorname{Tr}T)^b\) by positivity. Applying this in
the four directions proves the comparison. The intermediate
rectangular tori here use the same Wilson integrals and need no
renormalisation construction.

The actual vacuum constants have the decomposition
\[
 E_{\ell,K}=\sum_{j<K}\left[
   \mathsf E^{(j)}_{\mathrm I}
  +\sum_y\sum_{X\ni y} E^{(j+1)}(X,\mathbf1,y)
  +\sum_X R^{(j+1)}(X,\mathbf1)\right].
\]
In \(\mathsf E^{(j)}_{\mathrm I}\), the Haar and Gaussian factor is
one scalar per nonpivot bond and the gauge-fixing factor is one
scalar per cell of the next lattice. The common coupling sequence
and fixed choice of gauges give the same two scalars for the two
tori, and both counts per unit volume are the same.
Point the remainder term by assigning \(R^{(j+1)}(X,\mathbf1)/|X|\)
to each root of \(X\). A domain whose complete window lifts to a
cube of side at most \(\min(\ell,\ell_\star)/4\) has the same
rooted value on either torus by the common companion rule of
\cref{sec:flow-a-common-companion-and-the-torus-floor}.
The density of these terms agrees, since the rooted family has the
periods of the smaller side. For the other domains the length at
depth \(u\) is at least \(c\ell_\star L^u/M-C\). Summing over their
roots and then over \(u\) gives
\[
 \left|E_{\ell,K}/\ell^4-E_{\ell_\star,K}/\ell_\star^4\right|
 \le C\sum_{u\ge0}L^{4u}e^{-c\ell_\star L^u/M}.
\]
Thus the common companion comparison applies to the vacuum constant
in the sourced construction. Neither vacuum density need have a
cutoff-independent limit.

Use upper stability of \cref{prop:source-domains-4-1} on
\(T_{\ell_\star}\) and the partition floor on \(T_\ell\).
Writing \(\mathcal E_\star(W_z)=\sum_{j,X}e_j(X;W_z)\) for the
reference source energy, the preceding estimates give
\[
 \|e^{\mathcal O(w_z)}\|_{{\rm ch},\ell}
 \le \exp\!\left[N_{\ell_\star}^{-1}\operatorname{Re}\mathcal E_\star(W_z)
              +C_g\mathsf l_c^4\right].
\]
Define \(E^\circ(z)\) on the reference torus as the average over
its cells of the sums of \(e_j(X;\cdot)\) with the complete
source footprint contained in one open cell. This includes the
barycentres of all bare plaquettes with a corner in \(X^{\sim1}\).
Evaluate the energy on
that cell's reflected source. In an odd cell, also conjugate the
answer. Such terms have the form \(\overline{h(\bar z)}\) with
\(h\) holomorphic, so \(E^\circ\) is holomorphic and \(E^\circ(0)=0\).
Its real part is exactly the corresponding interior portion of
\(N_{\ell_\star}^{-1}\operatorname{Re}\mathcal E_\star(W_z)\).
This identification uses locality on the complete source footprint
and compares the same energy function in each fixed domain.

A nonzero omitted term must meet a cell face and reach a plaquette
whose barycentre is in a source support. The face clearance implies
a lower bound \([cL^u/M-C]_+\) on its length at depth \(u\),
with the additional bare reach included in \(C\).
There are at most \(C(1+L^u/M)^4\) possible roots per cell near a
support. The length entropy is smaller than the rate \(\kappa_e\)
of \cref{prop:source-domains-4-1}. Thus the omitted sum per cell is
bounded independently of \(K\) by
\[
 C\sum_{u\ge0}(1+L^u/M)^4
   \exp\big(-(\kappa_e-\kappa_{\mathcal D})[cL^u/M-C]_+\big).
\]
The same \(E^\circ\) now gives \eqref{eq:two-point-extraction-partial-ET5}
for the reference torus and every compatible torus, by the one-cell
chessboard inequality. It gives that estimate for every periodic
Gibbs state of the theorem by \eqref{eq:all-state-source-transfer-ST3}.
The constants are uniform in the datum. This construction compares
with the prescribed finite reference side; it does not require
taking a subsequential torus limit.

By \cref{thm:amrp-all-periodic} and the composition of a site
reflection with the adjacent link reflection, the states in the
theorem are invariant under bare lattice translations. Therefore
\eqref{eq:two-point-extraction-partial-ET7} applies to their
differences from the reference. Its error is
\(C_{g,\vartheta}\mathsf d^2\|f_1\|_\infty\|f_2\|_\infty\);
relative to the nonnegative reference main term in the stated
shape class, this is bounded by \(C_{g,\vartheta}\mathsf d^2x^2\).
It tends uniformly to zero by the witness depth bounds. The
comparison extends the reference estimate whenever the physical
source remainder and the reference witness covariance have the
estimates used above.

% END MODULE chapters/two_point_extraction_partial_v6.tex

% BEGIN MODULE chapters/clean_conditional_law.tex
\subsection{The conditional law and its connected errors}
\label{ch:clean-conditional-law}

The short-distance coefficient concerns an interacting probability law.
Comparing quadratic kernels does not identify that law.  We specify the conditional measure and its complete exponent, then
derive the boundary terms, connected errors and exterior average.
Applying these identities requires the run-window identification and
uniform estimates stated at the end.  The primary route is the resummed clean
run; an averaging-fibre transfer is a distinct alternative.

\subsubsection*{Identify the measure before estimating its density}

Fix the exterior sigma-field of the run window, with value $e$.
The identity concerns that specified sigma-field and the resummed
positive term.  A finer history or interpolation resolution need
not be positive or carry the same conditional measure.

In an admitted real pivot chart write the raw logarithms as
$a=(p,u)$.  Suppose the constraints $q(p,u)=0$ have invertible
$D_pq$ and solution $p=\varphi(u)$.  If the original coordinate
density is $w(p,u)$, the constrained density is
\begin{equation}
 w(\varphi(u),u)|\det D_pq(\varphi(u),u)|^{-1}\,du .
 \label{eq:clean-conditional-law-CL1}
\end{equation}
This is the change of variables from $(p,u)$ to $(q(p,u),u)$,
followed by evaluation of the constraint delta function.
Haar factors, delta normalizations and the exact gauge-quotient
density belong in $w$.  A gauge-fixing factor requires the
separate identity inserting it.

Under a canonical linear rescaling $u=gCy$ retain its determinant.
For a variable-dependent map use its full Jacobian.  The domain
is the image of the actual integration set; a product window
is an identity to establish.

On a positive component the complete conditional density may
be written
\begin{equation}
 d\nu_e(y)=Z_e^{-1}1_{D_e}(y)
  \exp\{-\tfrac12 y^{\mathsf T}Q_e y-\ell_e^{\mathsf T}y+P_e(y)\}\,dy .
 \label{eq:clean-conditional-law-CL2}
\end{equation}
The actual extraction rule fixes $Q_e$.  The complete $P_e$
contains old activities, nonlinear action, coarea and Haar
logarithms, boundary and stored terms, and healed-origin
remainders.  A factor outside this exponential is carried in
every numerator, or absorbed where its logarithm exists.
Zeros require a domain or a decomposition.

An arbitrary choice of $Q_e$ would make \eqref{eq:clean-conditional-law-CL2} merely a rewriting.
The extraction rule, sigma-field, domain, normalizer and
observable numerators all belong to the required identification.
For the exact calculus below assume a bounded piecewise smooth
domain and a density in $C^1(\overline D)$, or in $W^{1,1}(D)$
with the indicated integrable traces and integration-by-parts identity.
All displayed force integrals are required to be integrable.
A rectangle suffices.  Unbounded domains require an exhaustion
and justified boundary limits.

\subsubsection*{The complete Hessian and the exterior force}

Set
\begin{equation}
 V_e=\tfrac12y^{\mathsf T}Q_e y+\ell_e^{\mathsf T}y-P_e(y),
       \qquad r_e=\nabla P_e-\ell_e .
 \label{eq:clean-conditional-law-CL3}
\end{equation}
The Hessian is $Q_e-P_e''$.  The old-activity theorem controls
its stated old-origin sum, not every other term in $P_e$.
Relative to a specified reference the exact force difference is
\begin{equation}
 \nabla V_e-\nabla V_0
  =(Q_e-Q_0)y+(\ell_e-\ell_0)-\nabla(P_e-P_0).
 \label{eq:clean-conditional-law-CL4}
\end{equation}
In the actual run this includes the marginal Schur force
$Q^{\rm sc}y$, the integrated exterior force $Q^e_{IO}y_O$,
the newly stored exterior force $\ell_E^e$, the interior
precision discrepancy and the interaction derivative.
Let \(\mathcal B_{\rm ext}\) consist of \(O\cup E\)
and the complement of the padded common region
on which the two operator geometries and coordinate
masks agree. The force bound required at an interior
coordinate is
\begin{equation}
 |f_b|\le Cb_j e^{-\mu d(b,\mathcal B_{\rm ext})}
 \label{eq:clean-conditional-law-CL5}
\end{equation}
uniformly in the admitted exterior data.

%
%
%
\paragraph{The determining data and the response to the free box.}
The restriction identity \eqref{mb:restriction-identity}
identifies the box background on its induced datum.
In the step from \(j\) to \(j+1\), fixing \(V_{j+1}\)
and the exterior free fluctuations still leaves
the retained lower determining data as arguments
of the background and the conditional precision.

One part of those arguments drops out
exactly. Let \(E_1\) be the fine bonds incident
to \(\Omega_1\), \(E_0\) their complement, and
\(C_0\subset E_0\) the bonds in a plaquette
touching \(E_1\). The bare datum fixes all of
\(E_0\). Enlarge \(C_0\) by the portions of \(E_0\)
read by the positive-level double blocks or the
fixed axial gauge stencil. In the outward-row
convention the additional averaging reads come
from the level-one outward rows, within the reach
of a level-one double block. Fix the nest and
axial frame. If two data agree at every positive
level and on this enlarged set, the constraints
on \(E_1\) and their chosen gauge slice coincide.
For any field on that slice,
the Wilson actions agree on plaquettes
touching \(E_1\), and the other plaquettes
give a scalar difference independent of
the field on \(E_1\). It follows that the
constrained axial solve and the constrained
Wilson Hessian on \(E_1\) agree exactly.
The scalar cancels in normalized box
expectations. In this argument the deep
bare fields may be rough. The other terms
of the density have their own reads.

On the native small-field support the retained
rows below \(j\) have the zone and interface
plaquette bounds of the determining nest. At the
top level, start with the bound \(2\varepsilon_j\)
on \(V_j\). Replacing an interior bond by the
average from the next step changes its factor by
at most \(c'\varepsilon_j\), with the same bound
in the reversed orientation. Telescoping the
product of four unitary factors gives at most
\(4c'\varepsilon_j\) in additional plaquette error.
This also covers mixed top plaquettes at the cut.
At the action collar in zone zero, the native
small-field guard gives the bare bound
\(\varepsilon_0\). These are bounds on the actual
base datum before varying the free box.

For a free box variation about one base
datum, extend its coordinate \(y\) by zero
on the other fluctuation bonds, writing
\(\widetilde y\) for the extension. In the
native top rows the variation is multiplication
by \(e^{iB(y)}\), where
\[
 B(y)=g_j\mathbf C\widetilde y
         -\mathbf h\,\widetilde D(g_j\mathbf C\widetilde y).
\]
Here \(\mathbf C\) is the linearized
constraint lift, \(\mathbf h\) the corridor
operator and \(\widetilde D\) the analytic
correction vanishing to second order.
On a tube of size \(Y\), its quadratic
bound gives \(\|B\|_\infty\le 2g_j c_C Y\),
below the datum radius by the native ceiling.
For a unit direction in one bond \(b\),
the derivative of \(B\) is \(O(g_j)\).
The linear lift reaches only \(b\) and the
pivots of its incident block pairs.
The corridor correction reads those
block pairs and their adjacent corridors.
Consequently this derivative reaches at
most \((4d+2)^2\) top bonds within
\(4L+2\) input-level units of \(b\).
The clean guard puts both endpoints of
every bond in that support inside
\(\Omega_{j+1}\), including the pivots.
All lower rows are left fixed.

On the admitted retained datum let
\(\mathcal H(B)\) be the analytic
response at its fixed base \(U_0\).
For real parameters the solution with
the varied top rows is gauge equivalent
to \(e^{i\eta\mathcal H(B(y))}U_0\).
The finite support of \(D_bB\) makes the
chain rule a sum over the bounded list of
top rows. The derivative bound of the
retained response applies on these carriers.
Use its comparison of distances to the top
rows and identify the derivative in its
plaquette estimate with the covariant
difference at \(U_0\). This gives, for
\(H(y)=\mathcal H(B(y))\),
\[
 \eta|D_bH(y)(x)|\le Cg_j L^{-n(x)}E(x,b),
 \qquad
 \eta^2|D_{U_0}D_bH(y)(p)|
       \le Cg_j L^{-2n(x)}E(x,b).
\]
Here \(n(x)\) is the retained zone and \(D_{U_0}\)
is the covariant lattice difference at the plaquette
\(p\) based at \(x\). The response weight is
\(E(x,b)=\min\{1,e^{-c\rho(x,b)+c_1}\}\),
where \(\rho\) is the retained response distance
with the top spacing as unit. The constant \(c_1\)
allows for the displacement of the rows of \(D_b B\)
from \(b\). On a smaller tube the
two orders of Cauchy differentiation
give a mixed bound with the minimum of
the two weights and the inverse tube
margin. This constructs the read and its
derivatives on a fixed labelled nest with
admitted restricted data and response.
The comparison of masks or of positive-level
data across different exteriors still
uses their matching estimates.

%
\paragraph{Terms outside the box that read the background.}
%
%
The ordinary pre-strip chart also allows nonzero
fixed exterior fluctuations. In that chart let
\(z\) be the other original free coordinates.
The native real windows give
\(\|z\|_\infty\le Cb_j\). The complete level-\(j\)
datum uses the multiplier \(e^{iB_{\rm full}(y,z)}\),
with the same lift and nonlinear correction as
above, applied to the full vector \((y,z)\).
Keep the outgoing target and all lower data fixed.
Apply the bound for the lift and its correction
with \(R=\max(Y,\|z\|_\infty)\). On the joint chart
the smallness of \(g_j c_C R\) gives
\(\|B_{\rm full}(y,z)\|_\infty\le2g_jc_C R\)
and \(\|D_yB_{\rm full}\|_{\infty\to\infty}
\le2g_jc_C\). The individual derivative \(D_b\)
uses the same rows within \(4L+2\) of \(b\).
These statements concern the original free slots,
including an exterior slot once its identification
with an original free coordinate has been made.

Choose \(y_*\) on the supported real fibre for
the fixed \(z\). Use the actual datum there as
the base. The relative logarithm for a variation is
\[
 A_{z,*}(y)=\frac1i\log\!
    \left(e^{iB_{\rm full}(y,z)}e^{-iB_{\rm full}(y_*,z)}\right).
\]
It is zero at \(y_*\). On the fixed exponential and
logarithm balls its norm is at most \(C_Lg_jY\)
and its derivative in one unit bond direction is at
most \(C_Lg_j\). For the size bound use the segment
from \(y_*\) to \(y\); for the derivative use the fixed
second factor of the relative product. Outside the
full top carriers of \(I\) the product is the identity.
The native response theorem is applied at the real
regular datum on the supported fibre, with \(A_{z,*}\)
in its input ball. Its chain rule gives the same bond
and plaquette estimates as above, with a fixed change
of constants. If the base is needed at zero, the
same segment varies only the clean top carriers.
Telescoping the four factors of each affected top
plaquette enlarges its real regularity bound by
\(C_Lg_j\|y_*\|_\infty\). The lower interface and
bare collar reads are unchanged when the full
carriers are in the padded top interior.

%
%
For the ordinary strip substitution this
identification is exact at fixed retained data.
The global centre is determined by the outgoing
block field, the level-\(j\) field off the
incident stock of \(\Omega_{j+1}\), and the
lower determining rows. The joined local
centre uses the outgoing rows and the old
retained rows of the output label. The latter
are fixed when the free box is varied. In
particular the old level-\(j\) determining
member used here has both endpoints outside
the released region. In the join
for the strip, the inner cut is in the top
layer and its new rows lift the fixed outgoing
rows. The outer cut is in the previous layer
outside \(\Omega_{j+1}\) and lifts the fixed
old rows there. The averaged local strip field
is the one in the axial frame of this joined
set. Neither centre reads the shifted
strip nonpivots. Thus both
centres, and their relative shift on a strip
bond, are fixed throughout this variation.
The real cutoff in the inverse strip map
is a function of that bond's stored coordinate.
Consequently a held ordinary post-strip
nonpivot yields a fixed pre-strip nonpivot,
including on the cutoff's transition shell.
The Jacobian of that held slot is also fixed.
The box coordinates themselves are unchanged
on the ordinary passing mask, and the
pivots are recomputed at the same target.
The ordinary response calculation therefore
applies to these conditional slots.

%
%
At the top preparatory stage of this same
output level the two centres stay fixed under
the original clean-box variation as well.
The stage member contains the old
\(\Omega_{j+1}\). Its old level-\(j\)
determining rows are the rows outside that
member, with neither endpoint inside it.
Thus they miss the box variation and its
complete pivot support inside the old
\(\Omega_{j+1}\). Keep the added coarse
targets and exterior data fixed. The joined
background uses the lift of the fixed output
field on the old inner region and the retained
datum on the outer region. The two stage
centres and the bondwise cutoff translation
are then fixed when the original box varies.
At a held real stage slot this again fixes
its inverse coordinate. A variation of the
released collar variables themselves must
use their stage chart.

%
There is also an exact fixed-coordinate case
for the healed reference draw. If its component
and full retained exterior rows are outside the
padded box read, fix those rows in the exterior
disintegration. The basic centre and the
normalized reference law depend only on that
retained exterior of the component. With its
relative reference draw held fixed, the restored
physical field \(V'V_\Lambda\) is then fixed,
and the reference density is scalar in \(y\).
The previous box variables have already been
integrated in this kernel. The completed global
background and the main-step density subsequently
read the replacement rows, and their comparison
uses those actual reads.

Consider a term \(f\) whose
remaining dependence on \(y\) is its gauge-invariant
evaluation on the responding background, restricted
to a bond set \(X_f\). Its other arguments are held
fixed. The response estimates give size
\(C_H\Theta(x)\) in each of the bond and plaquette
units above, where \(\Theta(x)=g_jY\sum_{b\in I}E(x,b)\)
for a tube of size \(Y\).
Integrate from the base, where the response
vanishes, along the segment to \(y\).
For the term's analytic
norm divide the two units by their radii \(\rho_f^{(1)}\)
and \(\rho_f^{(2)}\). Write \(r_f^{-1}=\max_i
(\rho_f^{(i)})^{-1}\). If
\(C_H r_f^{-1}\max_{x\in V(X_f)}\Theta(x)\le 1/2\),
the restricted field lies in the half ball for that
norm. A term bounded by \(M_f\) on the full ball
has the continuation \(\varphi_f(y)=f(W^{(y)}|_{X_f})\),
where \(W^{(y)}=e^{i\eta H(y)}U_0\) in the
response gauge at its base. Apply Cauchy's
estimate in the direction of the response derivative
inside the remaining half ball. With \(w_f=g_jM_f/r_f\),
this gives
\[
 |D_b\varphi_f(y)|\le 2C_H w_f
                      \max_{x\in V(X_f)}E(x,b).
\]
On a smaller tube with margin \(r\) the mixed
derivative has the bound \(2C_H w_f/r\) times
the smaller of the two endpoint weights for \(b,b'\).

Use the scaled path distance in \(E\). A path from
zone \(n\) to \(I\) crosses each intermediate zone.
If those zones have width at least \(M\) in their
own units, each contributes at least \(M\) to the
scaled distance. Split the response exponential
with rates \(\mu=\nu\) equal to half the rate in \(E\).
The count and the ordinary top-scale distance give,
for a level-\(j\) cell \(c\),
\[
 E(x,b)\le C e^{-\nu d_j(c,b)} e^{-\mu M(j-1-n)_+},
 \qquad x\in c\cap Z_n .
\]
Select a maximizing endpoint of each \(X_f\).
Suppose for any such choice the total \(w_f\) in
one cell and zone is at most \(C(1+j-n)^q L^{4(j-n)}\).
The age sum is finite by \(L^4 e^{-\mu M}\le1/2\).
If all the reads are at distance at least \(D_{\rm far}\ge1\)
from \(I\) in level-\(j\) units, the cell sum is at
most \(C D_{\rm far}^{d-1}e^{-\nu D_{\rm far}}\).
This bounds the sum of the absolute derivatives at
each \(b\in I\). It gives an entrywise mixed
derivative bound after division by the tube margin.

For the uncompensated frozen coupling, the far
term is the sum of Wilson plaquette terms with
coefficients \(g_j(x_p)^{-2}-g_j^{-2}\), fixed with
the exterior and bounded by \(C(1+j-n(p))\).
Their nonzero plaquette reads are outside the
collar of width \(D_{\rm far}\).
Plaquettes entirely in the fixed deep bare region
have zero derivative. At each responding plaquette
the source trace estimate gives
\(C g_j(\bar\varepsilon+\theta_*)L^{-4n(p)}
E(x_p,b)\) before multiplying by the coefficient.
Here \(\bar\varepsilon\) bounds the base curvature
in its zone units and \(\theta_*\) is the native trace
threshold for \(\Theta\) on the lag support.
On that support the unweighted top distance gives
\(\Theta(x)\le Cg_jY D_{\rm far}^{d-1}e^{-2\nu D_{\rm far}}\).
Use the collar width for which this is at most \(\theta_*\).
The trace derivative gains one factor of the
curvature and one from its variation.
There are at most \(CL^{4j}\) plaquettes in a
level-\(j\) cell, so the same cell and age sums
give the force bound
\[
 Cg_j(\bar\varepsilon+\theta_*)D_{\rm far}^{d-1}
                             e^{-\nu D_{\rm far}} .
\]
For the mixed Hessian of the lag, the
plaquette trace bound uses the weight
\[
 r^{-1}\min\{E(x_p,b),E(x_p,b')\}.
\]
In the cell sum the spatial exponential uses
the larger of the two distances from \(c\) to \(b,b'\).
This maximum is at least \(D_{\rm far}-C_d\),
at least \((d_j(b,b')-C_d)/2\), and at least
\(d_j(c,b)\). Use one third of the exponential
for each of these three bounds. The last bound
sums over cells, and the zone sum is unchanged.
Thus the lag Hessian entries are bounded by
\[
 C r^{-1}g_j(\bar\varepsilon+\theta_*)
        e^{-\nu D_{\rm far}/3}e^{-\nu d_j(b,b')/6}.
\]
The weighted row and column sums have the same
small prefactor at any smaller fixed rate.
At fixed exterior either force estimate bounds
the oscillation of the corresponding exponent on
the product window by its per-bond bound times
\(C|I|b_j\), by integrating along a segment in the
window. The ratio of the two normalized real
densities with and without these terms is between
the exponentials of minus and plus this oscillation.
Fixed spectator arguments require a compatible
gauge comparison of the partial read. A term
whose current or stored-colour arguments vary
with \(y\) must include their derivatives as well.

\paragraph{The exterior terms on a matched chart.}
Fix the retained datum and the coordinate frame of the read set.
Write \(I\) for the integrated bonds in the box and \(O\) for
the other integrated bonds of this step. Write \(E\) for the
newly external fluctuation slots of the same ordinary step,
with their fixed values \(a_E\) in scaled coordinates. Let
\(Q^e\) be the complete precision on \(I\sqcup O\) before
the exterior integrated coordinates are frozen. It includes
the ratio correction to the contour gauge quadratic. The
conditional mean of the ordinary fluctuation Gaussian gives
the additional linear term in the negative log-density
\[
 \ell_E^e=\tau_k\mathsf A^e_{I\sqcup O,E}a_E,
 \qquad k=j-1 .
\]
Indeed, for the ordinary Gaussian with precision
\(\tau_k\mathsf A^e_{I\sqcup O,I\sqcup O}\) its mean is
\(-(\tau_k\mathsf A^e_{I\sqcup O,I\sqcup O})^{-1}\ell_E^e\).
The gauge compensation has separate integrated and
exterior blocks. It changes the precision and leaves
this linear term unchanged.
Keep the \(E\)-slots fixed, including any real near values.
At a comparison datum \(e_0\), the centered Gaussian reference
on \(I\) has covariance \(\bigl((Q^{e_0})^{-1}\bigr)_{II}\).
Its precision is
\[
 K=Q^{e_0}_{II}-S,\qquad
 S=Q^{e_0}_{IO}(Q^{e_0}_{OO})^{-1}Q^{e_0}_{OI}.
\]
The reference exponent is \(-y^{\mathsf T}Ky/2+P_0(y)\).
The difference of the gradients of the two negative log-densities
gives the terms
\[
 Sy+Q^e_{IO}y_O+(\ell_E^e)|_I
     +(Q^e_{II}-Q^{e_0}_{II})y-\nabla(P_e-P_0)(y).
\]
All scalar factors depending only on \(e\) cancel in this comparison
at fixed exterior. This includes the square-completion
term of the ordinary move with two external fluctuation
arguments: its inverse chain may reach inside the box,
but both arguments are fixed. The Schur term \(S\) above
has two interior arguments and comes from taking the
Gaussian marginal of the reference.

The Schur and integrated exterior forces follow from the real
precision bounds. Suppose
\(Q^{e_0}\ge a_- I\), and that \(Q^{e_0}\) and \(Q^e\) have
uniform exponentially weighted row and column sums. Minimizing
the quadratic form over the exterior gives \(K\ge a_- I\).
For a weight \(e^\varphi\) with small Lipschitz constant \(\mu\),
conjugation changes \(Q^{e_0}_{OO}\) in operator norm by at most
\(C\mu\), by its weighted row and column bounds. Its inverse
therefore remains bounded in the conjugated norm. Taking
\(\varphi(b)=\mu d(b,b')\) gives exponential decay of its entries
and, at a smaller rate, summable weighted rows and columns.
For \(b\in I\), the first factor \(Q^{e_0}_{IO}(b,\cdot)\) in
\(S\) already crosses distance \(d(b,O)\). Summing the other
two factors gives
\[
 |(Sy)_b|+|(Q^e_{IO}y_O)_b|
       \le Cb_j e^{-c d(b,O)}
\]
on the stated windows \(|y|_\infty,|y_O|_\infty\le Cb_j\).
The old exterior tier bounds give \(\sup_E|a_E|\le Cb_j\)
on the real support. The weighted row bound for the
ordinary cross block therefore gives
\[
 |(\ell_E^e)_b|\le Cb_j e^{-c d(b,E)}.
\]
All three bounds hold at the read with the box padding,
since both \(O\) and \(E\) are outside the box. They
also hold with the distance to \(\mathcal B_{\rm ext}\).

For the change of the interior precision compare the full precision terms before
freezing the exterior. Use the local walk expansion with the
current, pivot, and relative contour gauge contributions in the
same chart.

Fix the outgoing datum \(V_j\) and take a collared box \(W\)
in the common clean top layer of two exteriors. Choose the same
collar grid in both settings, with the padding for all the
backward blocks and pivots. On \(W\) all the layer traces
through level \(j\) coincide. The only original determining
rows with both endpoints in \(W\) are top rows, which read
the fixed \(V_j\). The joined determining set therefore
coincides, and both joined data are formed by the same
face lift of those rows. In particular both ambient
backgrounds compare to one joined background by
\eqref{mb:real-collar-norm}. The common comparison box
gives one frame for the interior. If a lower read uses the
average of either ambient background, apply the local average
estimate on its entire double block in that frame. This gives
closeness of the two averaged data from the background estimate.

On a union of common standard cells with padding in \(W\),
the identical traces give identical indices, row types,
outward endpoints, and local carriers of the native walk
terms. For the ordinary compressed terms use the part
of this union on which both steps have the same integration,
pivot, and conditioning masks. A masked term uses this
identity only on its full carrier, including the endpoint
blocks of its masks.

% Earlier Part I 7267--7275, 11758--11825, 14925--14959.
For the native conditioned covariance, the local
resolvent at a cube \(q\) uses the sum of the form
terms with footprints inside \(q^\sharp\), compressed
to the kept bonds in \(Y'\cap q^*\). Thus one must
match the full \(q^\sharp\) read and the compression
on \(Y'\cap q^*\), even for an inverse entry whose
endpoints are well inside \(q^*\). On this common
support write \(R^a_q=(x_a+T^a_q)^{-1}\) for the
local resolvent of either datum \(a=e,e_0\). Then
\[
 R^e_q-R^{e_0}_q
 =-R^e_q\bigl[(x_e-x_{e_0})I+T^e_q-T^{e_0}_q\bigr]R^{e_0}_q .
\]
Equal data and equal scalar shifts give equality.
For nearby data the same identity uses the native
inverse bounds on the admitted domain.
An extra chain of a labeled reference
must have this same form family and compression, or its
own comparison.

Write \(d_b=d(b,\mathcal B_{\rm ext})\), with the cell and
stencil padding included in the common region. Compare the terms in row \(b\) whose
full carriers stay within distance \(d_b/2\) of \(b\).
The other terms have charged length at least \(c d_b-C\),
including the margins of their local read sets. The absolute
weighted row bound for each precision gives a cost
\(Ce^{-c'd_b}\) for this part, with the column variable summed.
This distance is bounded by the distance to the actual
cut even when the set of exterior integrated bonds has holes.

For the short carriers, compare the terms with common
masks in the frame of \(W\).
The real pair difference, in local units on a carrier, is bounded
by \(C\epsilon_0\sup_{z\text{ on the carrier}}e^{-c_0 d(z,W^c)}\).
The real estimates for the potential and its first differences
give the field bound, and the full-bond estimates give the current
bound. The adjoint rotations preserve the free-bond balls.

The ordinary relative contour square has a direct local
comparison. The paths in one block avoid the middle
crossing pivots. Thus \(G_Uh_U=0\) for those contour
rows, and their differential \(G_U C_U\) from free
coordinates is \(G_U\iota\) by
\eqref{eq:first-step-nonflat-precision-NFP3}.
At a common block it has the same path stencils and
free mask in both settings. The row coefficients in
\eqref{eq:explicit-relative-current-extension-NX4}
are analytic in the block links. Their difference is
bounded by the field difference on the block, in the
common node frame. The weight of the gauge quadratic
is one in \(Q_\tau=\tau K_A+K_G\). The difference of
the two Gram terms is the sum of a bounded row
times a row difference and its counterpart with the
difference in the first factor. The fixed number of
contours per block then gives the short row estimate
for this part of the complete precision.

On a common local operator domain with fixed room, Cauchy's
estimate on the segment between the two pairs bounds the difference
of a short term by its absolute analytic bound times $C$ times
the pair difference. Thus the absolute row sum of short
differences is at most $C\epsilon_0 e^{-c_0d_b/2}$. Together
with the long terms this gives
\[
 \sum_{b'\in I}|(Q^e_{II}-Q^{e_0}_{II})(b,b')|
                         \le Ce^{-c'd_b}.
\]
Here the row estimate uses the matching for each short summand
of the complete precision. For the reference used in the clean
mixture, we must file any additional boundary chains under the
same local comparison or bound their short contributions
separately. Multiplication by \( |y|_\infty\le Cb_j \)
gives the force bound at depth \(d_b\).

For the interaction force use the full carriers of the evaluated
members of \(P_e\). Write \(m(Y)\) for their length, with
\(\operatorname{diam}Y\le C_0(1+m(Y))\) and
\(\sum_{Y\ni b}e^{-\kappa m(Y)}\le C\).
For a carrier reaching \(\mathcal B_{\rm ext}\), a derivative bound
\(v_1e^{-2\kappa m(Y)}\) on each evaluation gives the needed
distance factor by spending half the decay. For a carrier
contained in the matched interior, use the matched derivative bound
\[
 \left|D_b(P_{e,Y}-P_{0,Y})\right|
     \le v_2e^{-2\kappa m(Y)}
                      \sup_{z\in Y}|U_e(z)-U_{e_0}(z)|
\]
in a common gauge. If the last difference is bounded by
\(\epsilon_0e^{-c_0d(z,\mathcal B_{\rm ext})}\) at \(z\), take
\(c'\le\min(c_0,\kappa/C_0)\) and use
\(d(b,\mathcal B_{\rm ext})\le d(z,\mathcal B_{\rm ext})+C_0(1+m(Y))\).
The rooted sum gives \(C(v_1+v_2\epsilon_0)e^{-c'd(b,\mathcal B_{\rm ext})}\).
The target in \eqref{eq:clean-conditional-law-CL5} follows
when \(v_1+v_2\epsilon_0\le Cb_j\).
For the regular old-origin members, the common-domain amplitude
estimate in \eqref{fm:exterior-amplitude-premise} also suffices,
by Cauchy's inequality and the same rooted sum. The force bound
must be checked for the other members of the complete exponent.

The readout must be compared in addition to the density.
On a common admitted pivot chart \eqref{afs:actual-pivot-chart},
the joint Lipschitz bound and finite double-block support of
\(\Phi_V\) bound the difference of the two evaluated link fields
by \(C\epsilon_0e^{-c'\mathsf D}\) on the read set, when
the complete pivot support has distance at least
\(\mathsf D\) from \(\mathcal B_{\rm ext}\).
Apply the background estimate to the coarse field
entering the graph, on its whole pivot support. This uses
the actual averaging estimate if the comparison was in fine
variables, and a fixed margin for the pivot support.
Thus a readout with Lipschitz constant \(L_G\) costs
\(\eta_G\le C L_G\epsilon_0 e^{-c'\mathsf D}\).
%
For the scalar insertions there is one more small
field factor. Put \(\Delta=C\epsilon_0e^{-c'\mathsf D}\)
for the difference of the two fluctuating link fields
at the same free coordinate, after common gauge transport
on the full pivot read. Let \(\kappa_*\) bound the
curvature of the reference \emph{evaluated} field on the
window, including its fluctuation. In the axial gauge
of this reference its link logarithms on a box of
side \(\ell_G\) are bounded by \(C\ell_G\kappa_*\).
The short linkwise interpolation to the other field
then has logarithms bounded by
\(C(\ell_G\kappa_*+\Delta)\).
The first derivative of the invariant scalar vanishes
at the flat field. On the small-field logarithm
chart, use the second derivative bound from zero
to each point of the path to estimate its first
derivative, then integrate along the path. This gives
\[
 \eta_G\le C\mu_2(\ell_G\kappa_*+\Delta)\Delta .
\]
For a Wilson plaquette in \(\mathrm{SU}(2)\) the same
bound follows by subtracting the squared norms of
the two plaquette deviations from the identity.

\paragraph{Intrinsic terms and the reference Hessian.}
\label{cl:intrinsic-reference-hessian}
Work in the canonical free coordinates \(y\), with \(g=\widehat g_{j-1}\)
and a tested radius \(b_j\ge1\). In the intrinsic exponent the
action and contour-gauge terms are composed with the actual pivot
graph before their quadratic parts are extracted. Their remainders
have the form \(g^{-2}R_{3,Y}(gy)\), with their signs and fixed
coupling weights. Here \(R_{3,Y}\) vanishes to order three, and the
complete quadratic terms are included in the precision. Any
linear term of the full action must be accounted for in the
choice of background or in the explicit linear force.
At the physical critical background, the tangent of the pivot
graph at fixed targets lies in the tangent space of the
constraint. The first variation of the Wilson action there
vanishes on that tangent space. The contour-gauge action
has zero first variation at the axial centre, where its
contours are the identity. Extract these Taylor coefficients
on the formed carriers before bounding the cubic remainders.
Freezing a nonzero exterior argument later still creates the
linear terms from the extracted quadratic.
The other intrinsic terms are
\(L_Y(gy)-L_Y(0)\), the local log Haar and pivot coarea factors
of \eqref{afs:actual-coarea-density}. In particular, the pivot
Jacobian is the derivative in the eliminated coordinates of the
actual constraint. The physical cubic and linear coefficients
are the ones in \eqref{fec:literal-first-two-vertices}.

On a fixed unscaled complex radius, let \(B_Y\) bound the
composed action and the coarea functions on their full
carriers. For the homogeneous action these bounds are supplied
by \eqref{het:actual-action-family} and its composition with
the outer graph in \eqref{haf:actual-effective-cubic-vertex}.
For the local Haar and pivot factors they follow from the
joint chart and \eqref{afs:actual-coarea-density}. Suppose that
\(g B_*\) stays within a fixed smaller radius on every
complete carrier. Here \(B_*\) includes \(2b_j\),
the real held bound \(C_O\widehat b_j\) on the other
integrated slots, and the full bound on the original
exterior slots \(a_E\). Leave a fixed fraction of
\(b_j\) beyond the real windows for the complex
coordinate variations, and use \(B_*\le C b_j\)
with the fixed native window constants.
Taylor's formula on this joint chart gives the
respective bounds \(CB_Ygb_j^3\) and \(CB_Ygb_j\).
For the coarea Hessian subtract its linear Taylor term before
estimating derivatives. The subtracted function is bounded by
\(CB_Yg^2b_j^2\) there. This subtraction is only for the
derivative estimate. The linear term remains in the density.
Two coordinate Cauchy estimates at a real point in the tested
box give the entry bound
\[
 |D_bD_{b'}P_{{\rm int},Y}|
       \le CB_Y(gb_j+g^2)\,
                  \mathbf1_{\{b,b'\in Y^+\}} ,
\]
where \(Y^+\) is the support of the term in the actual chart.
The radii in both derivatives are fixed fractions of \(b_j\).
For a repeated coordinate the factorial is \(2!\).
The rooted exponential bound for \(B_Y\) absorbs the polynomial
number of bonds in \(Y^+\) and the factor
\(e^{\mu\operatorname{diam}Y^+}\) for a smaller fixed \(\mu\).
The weighted row and column sums therefore give
\[
 \|T_\varphi P_{\rm int}''T_\varphi^{-1}\|_{2\to2}
      \le C(gb_j+g^2),\qquad
 |\varphi(b)-\varphi(b')|\le\mu d(b,b').
\]
This bound holds on the proved homogeneous chart, and on a
matched reference box with the same rooted intrinsic bounds.

%
%
The Haar and coarea part of this calculation
applies directly on the labeled clean double
blocks at nonzero held coordinates. Each such
term is the actual one-step Haar factor or the
inverse derivative in the private middle pivot.
In a clean complete block the map, the target
and the nonpivot and pivot masks are those of
the joint local chart. The regularity of the
whole double block and the held radii just
specified supply its analytic bounds. The
number and length of these carriers through
one bond are bounded. Two derivatives give
\(Cg^2\) in the weighted rows and columns of
this part, with the linear coarea term still
in the density.

% New Part II, Proposition 53.4(b), Proposition 53.5,
% especially 25347--25353, 25386--25392, 25473--25488,
% 25556--25569; new Part III, Lemmas 74.8 and 74.18.
The Wilson action on an ordinary determining set has
a direct weighted Hessian estimate on the same chart.
Use its exact response package at the physical critical
reference. Write \(F\) for the bare action in the
relative potential, \(K_0\) for the linearized rows,
\(\mathsf H\) for their right lift and
\(\widetilde{\mathsf G}\) for the constrained Green
operator. The nonlinear row linearization
\(A=T(A')\) of \cite{B85e} makes the
rows linear: \(\mathcal Q(T(A'))=K_0A'\). With
\(\mathsf L_0\) the constrained action Hessian, put
\[
 V_{\rm nl}(A')=F(T(A'))-F(0)
       -(F'(0),A')-\tfrac12(A',\mathsf L_0A').
\]
All pairings and adjoints in this calculation use the
native potential, source and row weights. The second
variation estimate gives, on a fixed absolute ball,
\[
 \|V_{\rm nl}''(A')w\|^\sharp_{\theta,y_0}
       \le C\|A'\|_\infty\|w\|_{\theta,y_0}.
\]
Here \(\theta\) is a spatial weight rate and can be a
fixed positive number below the
native response rate. This estimate includes the
second derivative of the nonlinear rows contracted
with the physical multiplier. The determining-set
response package supplies it with constants uniform
in the depth and the volume.

For a reduced row datum \(b\), its critical potential
is \(T(A'(b))\), where
\[
 \begin{gathered}
 A'(b)=\mathsf H b+\mathcal A_0(b),\qquad
                         K_0\mathcal A_0=0,\\
 \mathcal A_0+
     \widetilde{\mathsf G}V_{\rm nl}'(A')=0.
 \end{gathered}
\]
The gauge condition is
\(\mathsf RD^{\mathsf T}\mathcal A_0=0\).
Shrink the fixed absolute datum radius once so that
\(\|\widetilde{\mathsf G}V_{\rm nl}''(A')\|\le1/2\)
in these weighted spaces. This is possible since
\(\|A'(b)\|_\infty\le C\|b\|_\infty\).
Differentiating the stationary equation gives
\[
 DA'(b)=(I+
       \widetilde{\mathsf G}V_{\rm nl}''(A'))^{-1}
                  \mathsf H.
\]
Set \(\mathcal E(b)=F(T(A'(b)))\) and
\(\mathsf Z=\mathsf H^\sharp\mathsf L_0\mathsf H\).
The identities
\(\mathsf L_0\mathsf H=K_0^\sharp\mathsf Z\) and
\(\mathsf H^\sharp\mathsf L_0=\mathsf ZK_0\)
remove the cross terms containing \(\mathcal A_0\).
Stationarity removes its first variations. Thus
\[
 \begin{aligned}
 \mathcal E'(b)&=\omega_0+\mathsf Zb+
                      \mathsf H^\sharp V_{\rm nl}'(A'),\\
 \mathcal E''(b)-\mathsf Z
   &=\mathsf H^\sharp V_{\rm nl}''(A')
       (I+\widetilde{\mathsf G}V_{\rm nl}''(A'))^{-1}
          \mathsf H.
 \end{aligned}
\]
The inverse has norm at most two. The weighted
response bounds for \(\mathsf H\), its adjoint and
\(\widetilde{\mathsf G}\) therefore bound the last
operator by \(C\|b\|_\infty\). On the highest
active rows the scale weights are one. Apply the
estimate to a single row and Lie algebra component,
with the weight centred at its assigned top site.
The output entries then decay exponentially in
the distance between the two assigned sites.
There are boundedly many active row components
per highest-level site. The native exponential
sum, at any smaller fixed rate \(\mu<\theta\),
therefore gives absolute weighted row sums.
Symmetry gives the column sums, and the weighted
Schur estimate gives the required \(2\)-norm bound.

For this application the varying original datum
is confined to the highest level of the ordinary
determining set. In the reduced datum, an outward
node gauge reads only the interior datum one level
below. Those lower increments are zero in the
relative chart. There are no inward tying rows.
Consequently \(\mathsf b(B)=B\) on the full varying
top stock, as in \eqref{px:active-datum-flatness};
the preceding estimate is already in its physical
row coordinates.

Now compose with the actual one-step pivot graph
\(B=\Phi(v)\), written relative to that reference,
so \(\Phi(0)=0\). If \(\mathsf K(v)\) denotes
the Hessian of \(\mathcal E\circ\Phi\), the full
chain rule is
\[
 \mathsf K(v)=D\Phi(v)^\sharp
       \mathcal E''(\Phi(v))D\Phi(v)
                 +\mathcal E'(\Phi(v))\,D^2\Phi(v),
\]
where the last term denotes contraction in the
row coordinate. The joint pivot chart bounds its
first three derivatives on the complete double
blocks. Their supports have bounded diameter and
overlap. The physical multiplier satisfies
\(\|\omega_0\|_\infty\le C\alpha\).
To compare these Hessians, write
\(J(v)=D\Phi(v)\), \(H_2(v)=D^2\Phi(v)\) and
\(e(v)=\mathcal E'(\Phi(v))-\omega_0\).
The difference is the sum of
\[
 \begin{split}
 &J(v)^\sharp[\mathcal E''(\Phi(v))-\mathsf Z]J(v),\\
 &J(v)^\sharp\mathsf ZJ(v)-J(0)^\sharp\mathsf ZJ(0),\\
 &e(v)H_2(v),\qquad
             \omega_0[H_2(v)-H_2(0)].
 \end{split}
\]
The first term uses the reduced-action estimate.
The second and fourth use the local derivatives
of the pivot graph. For the third,
\(\|e(v)\|_\infty\le C\|v\|_\infty\) follows
by integrating its bounded Hessian along the
segment. Each term therefore gives
\[
 \|T_\varphi[\mathsf K(v)-\mathsf K(0)]
                        T_\varphi^{-1}\|_{2\to2}
                    \le C\|v\|_\infty.
\]
In particular, the change in the contraction
\(\mathcal E'(0)D^2\Phi\) is included here.
Its value at zero is part of the complete precision
\eqref{px:next-constrained-hessian}; it cannot be
discarded before subtracting that precision.
For \(v=g(y,y_O,a_E)\), the Hessian of the
\(g^{-2}\)-scaled action remainder is exactly
\(\mathsf K(v)-\mathsf K(0)\). This proves its
weighted bound \(Cgb_j\) throughout the full held
window, including nonzero original exterior
coordinates. The local contour action has the
same bound from its complete block chart after
its actual quadratic is subtracted. Together
with the Haar and coarea estimate above this
proves the intrinsic Hessian estimate on the
ordinary labeled clean stock. If lower-level
data are also released, keep their reduced-datum
derivatives and physical dual weights; the
highest-level identification just used does not
apply to those directions.

There is a separate parity condition in the value estimate.
When all exterior fluctuation arguments of the analytic
intrinsic reference are zero, its odd part is bounded by
\(Cn_jgb_j^3\).
The even parts of the action remainder and the coarea increment
start at orders four and two, respectively, so their sum is
bounded by \(Cn_jg^2b_j^4\). A newly stored \(a_E\) held
at a nonzero value can enter a cubic with two interior
slots, which is even in the interior. With these stored
values fixed use the full intrinsic oscillation bound
\(Cn_jgb_j^3\). On the joint analytic chart the segment
\((y,t y_O,a_E)\), \(0\le t\le1\), still gives the force
comparison for setting the other integrated fluctuations
to zero. At fixed background
and fixed quadratic extraction, the mixed Hessian bound
just proved gives
\[
 \left|D_b\{P_{\rm int}(y,y_O)-P_{\rm int}(y,0)\}\right|
 \le C(gb_j+g^2)b_j e^{-\mu d(b,O)}.
\]
Indeed, integrate the mixed derivatives along that segment
and sum the weighted row over \(O\). The physical exterior
variables remain in the actual law.

To reach the parity reference, this comparison must also
set the exterior slots of the current-step analytic intrinsic
terms to zero. Use the raw fluctuation coordinates before the
strip substitution, at a fixed pair and with the physical
tests and Jacobian held in the original density. The larger
analytic carrier boxes for these intrinsic terms allow a
segment in their current-step stored arguments. On the common
admitted box, integrating the mixed Hessian along
\((y,0,t a_E)\) bounds the extra change of interior force by
\(C(gb_j+g^2)b_j e^{-\mu d(b,E)}\) at bond \(b\).
This is a comparison of analytic evaluations. The strip
selectors and older stored fields are still present in the
physical law and require their own bounds. The frozen
quadratic's linear force is accounted for separately.

The regular old-origin contribution is controlled by
\eqref{oa:localized-hessian} and its weighted version
\eqref{eq:local-formation-OA20} once the matched reference
box reads the same subtracted family on its larger analytic
domain. The vacuum-only remainder uses the separate age sum
\eqref{eq:local-formation-OA22} on its admitted class.
When the profiles have been redrawn with compensation after
healing, the coupling lag vanishes in the clean box.
A term whose arguments, including its background read,
are fixed by the exterior contributes a scalar to the box.
It cancels from the normalized density at
that exterior and remains in the exterior mixing weight.
A boundary term that still depends on \(y\) requires its own
weighted Hessian bound in the reference. An estimate for the
difference of two exterior forces alone does not supply this.
Thus the proved intrinsic bound and the admitted old and
vacuum-only bounds can be spent against the Schur gap
in \eqref{eq:local-formation-OA21}. The remaining reference
terms must fit in the same gap before applying
\eqref{cl:weighted-force-comparison}.

\paragraph{Freezing only the exterior fluctuations.}
For a fixed full chart of the small-field step, the pivot graph
and the extracted precision are determined before freezing the
exterior free coordinates. Keep the retained datum, background,
stored spectators and coordinate sets of that chart fixed. In
the auxiliary reference put the other integrated fluctuations
equal to zero in the smooth exponent and keep the same full
precision for forming the centered Gaussian marginal on \(I\).
The newly external values \(a_E\) remain fixed in the smooth
interaction. The force from their cross quadratic is
\(\ell_E^e\) in the physical exponent, and dropping it for
the centered reference costs the estimate above. In this
comparison the interior precision difference vanishes. The
Schur correction and the force \(Q_{IO}y_O\) are still present.
The field readout on its declared read set is also unchanged:
each of its pivot values reads only the free bonds of its
defining double block, and those bonds are included in the
read set. This is the graph in \eqref{afs:actual-pivot-chart}.
The joint step box for the intrinsic action and coarea factors
contains the segment \((y,t y_O)\) with \(a_E\) held fixed.
The auxiliary zero is an evaluation on that analytic chart.
The exterior tests and the recorded exterior values in the
physical law are kept.

For any smooth exponent member \(P\) admitting the common
segment \((y,t y_O)\), its change of interior force is exactly
\[
 D_b\{P(y,y_O)-P(y,0)\}
   =\int_0^1\sum_{b'\in O}
            D_bD_{b'}P(y,t y_O)(y_O)_{b'}\,dt .
\]
A weighted Hessian row bound \(A\) on the full segment therefore
gives \(C A b_j e^{-\mu d(b,O)}\) on that exterior window.
This applies to the intrinsic estimate of
\cref{cl:intrinsic-reference-hessian}. It applies to the old
terms when their weighted estimate is available on the joint
free-coordinate domain, including the exterior slots of these
derivatives. A bound on interior derivatives at a fixed exterior
does not supply those slots. When the exterior changes the background, stored fields
or local masks, freezing its fluctuations first gives an
exterior-dependent reference. The
\hyperlink{f2-direct-flat-density}{direct flat comparison below}
uses the common rotated read window and the full row estimates
to obtain one reference density for the mixture.

\paragraph{Charging a boundary tail to the reference Hessian.}
Let \(W_I\) be the union of the current cubes covering the
entire free-coordinate box \(I\), with \(N_W\) cubes. Suppose
the surviving boundary members satisfy the majorants of
\cref{prop:ccf-old-boundary-tail} on one common complex ball
of free radius \(R_y\ge2b_j\), after their actual substitutions.
Write \(B_\partial\) for their sum. Terms independent of
the interior coordinates can be set aside as scalars at fixed
exterior. Then \eqref{ccf:old-boundary-tail} gives
\[
 \sup_{\|y\|_\infty<R_y}|B_\partial(y)-B_\partial(0)|
                  \le C B_*N_W e^{-cs}.
\]
Coordinate Cauchy on the real window gives bounds
\(CB_*N_WR_y^{-1}e^{-cs}\) for each first derivative and
\(CB_*N_WR_y^{-2}e^{-cs}\) for each Hessian entry.
If \(m_I\) counts the scalar coordinates and \(D_I\) is
the diameter of \(I\) in the metric of the weights, the
weighted Schur bound is
\[
 \|T_\varphi B_\partial''T_\varphi^{-1}\|_{2\to2}
       \le C B_*N_W m_I R_y^{-2}e^{-cs+\mu D_I},
 \qquad \operatorname{Lip}\varphi\le\mu .
\]
This uses the diameter of the box, so it does not require a
diameter estimate for an individual carrier in terms of its
length after exempt hulls have been removed.

For example, the additional margin
\(s\ge s_*+(\mu/c)D_I+C_1\log x\) makes this contribution
small when the coordinate counts and \(B_*\) grow
polylogarithmically. In this inequality the collar and box
distances use the units of their respective estimates, with
the fixed conversion included in the constants. The enlarged
collar still has polylogarithmic size when the original box
and the trace padding do. For two exteriors satisfying the
same bound, the first derivatives cost at most
\(CB_*N_WR_y^{-1}e^{-cs}\) in their difference.
As a uniform force error it adds at most
\(C n_j b_j B_*N_WR_y^{-1}e^{-cs}\) to the logarithmic
marginal error. This follows by using the summable weights in
\eqref{cl:weighted-force-comparison}. The bound requires the
larger complex ball for these evaluated boundary members and
the stated clean collar around \(W_I\). A real oscillation
bound alone gives no Hessian estimate.

%
%
%
For the current ordinary-step boundary
and kept potentials one can use the
native analytic domain directly. Fix
the actual exterior at an admitted support
point, with its adapted pair and stored
arguments. In the original scaled free
coordinates on the clean integrated bonds,
the complex step box has radius \(3b_j\).
Vary \(y\) on the ball of radius \(2b_j\)
and keep every other argument fixed.
The support has room in those other
coordinates, including the cut strip
evaluation. The pinned substitution
preserves the stored arguments. Hence
the boundary and kept potentials on
this ball use their original full
carriers and the inherited amplitudes
of the ordinary step.

Apply two coordinate Cauchy estimates
on each carrier \(Y\). A derivative
in the bonds \(b,b'\) can occur only
if the complete carrier reads both
of them. The two root cubes, with
the pivot padding included, are
nonexempt in the clean box. A hull
tree must meet both and spends their
separation in ordinary edge length,
including the part of the path that
passes through any exempt hull.
Splitting the rate between this
distance and the rooted sum
\eqref{ccf:hull-root-sum} gives
\[
 \begin{split}
 |D_bD_{b'}P_{\partial}(y)|
 &\le C_a b_j^{-2}\vartheta_k
   \bigl[\mathsf B(\log\theta_k)^{-a}
           +\mathsf B_\Phi e^{-L\check R_j}\bigr]\,
                          e^{-\mu_0 d(b,b')},
 \qquad k=j-1 .
 \end{split}
\]
Here \(P_\partial\) sums these
current-step substituted potentials.
Their inverse-logarithmic bound uses
the native crossing length. The
weighted Schur sum is uniform in the
number and diameter of the free bonds
for \(\operatorname{Lip}\varphi<\mu_0\)
with a fixed rate reserve. This
estimate is for interior derivatives
at fixed exterior and can be used
in the reference Hessian.

\subsubsection*{From the exterior force to the normalized marginal}

Use the same literal product of free-bond balls
\(D=\prod_{b\in I}D_b\) for the two real densities. Write
\[
 H_a(y)=\tfrac12y^{\mathsf T}A_a y+\ell_a^{\mathsf T}y-P_a(y),
 \qquad d\nu_a=Z_a^{-1}1_De^{-H_a}\,dy,\qquad a=0,1 .
\]
Here \(a=0\) denotes the reference with the Gaussian marginal
precision \(K\), and \(a=1\) the actual density at fixed exterior.
The two real potentials are \(C^2\) in the interior.
For \(T_\varphi y(b)=e^{\varphi(b)}y(b)\), suppose
\[
 \langle v,T_\varphi A_0T_\varphi^{-1}v\rangle
       \ge\lambda_A|v|^2,\qquad
 \|T_\varphi P_0''T_\varphi^{-1}\|\le\epsilon_0<\lambda_A .
\]
Use \(\lambda=\lambda_A-\epsilon_0\), and let \(\delta_b\)
bound the component at \(b\) of \(\nabla(H_1-H_0)\).
For a real function \(G\) with
\(|G(y)-G(y')|\le L_G|T_\varphi(y-y')|\), the common-momentum
comparison gives
\begin{equation}
 |\nu_1(G)-\nu_0(G)|
   \le \frac{L_G}{\lambda}
       \left(\sum_{b\in I}e^{2\varphi(b)}\delta_b^2\right)^{1/2}.
 \label{cl:weighted-force-comparison}
\end{equation}

Here is the argument including the hard boundary. First shrink
each ball to radius \(r_b\), and choose two intermediate radii
between \(r_b\) and the original radius. Smooth the metric
projection onto the larger intermediate ball by an even
mollifier. With support sufficiently small, the resulting map
\(r\), taken block by block, is the identity on the smaller
intermediate balls, takes values in the original domain and
has \(\|Dr\|\le1\). Replace \(P_a\) by \(P_a\circ r\) and add
the same wall \(\epsilon^{-1}\sum_b w(|y_b|^2-r_b^2)\) to both
energies. Take \(w\) smooth, convex and nondecreasing, zero on
the negative half-line, with positive derivative on the positive
half-line and affine for large arguments. The wall is smooth and
convex, with bounded Hessian for fixed \(\epsilon\). Its Hessian
has a positive lower bound off the smaller intermediate ball.
In the Hessian of \(P_0\circ r\), the term
\(Dr^{\mathsf T}P_0''(r)Dr\) has weighted norm at most \(\epsilon_0\)
because \(Dr\) commutes with the block weight. The remaining term
is block diagonal, vanishes on the smaller intermediate balls
in its respective blocks and is dominated by the wall Hessian
for small \(\epsilon\). Thus the weighted coercivity remains
\(\lambda\) on the whole space. The wall cancels from the
difference force \(f=\nabla(H_1^\epsilon-H_0^\epsilon)\).

For these smooth energies the Hamiltonian flow
\(\dot q=p,\ \dot p=-\nabla H_a^\epsilon(q)\) preserves the
product of its Gibbs law and the standard momentum Gaussian.
Energy conservation and zero phase-space divergence give this
identity. Start the two flows with the same Gaussian momentum.
For \(v=T_\varphi(q^1-q^0)\) their difference satisfies
\(\ddot v=-M(s)v-T_\varphi f(q^1(s))\) and \(\dot v(0)=0\).
The matrix \(M(s)\) is the conjugated average of the reference
Hessian on the segment between the two positions. It satisfies
\(\langle z,M(s)z\rangle\ge\lambda|z|^2\), and its norm is at most some
\(\Lambda\) at fixed \(\epsilon\).
The integral equation, for small \(t\), gives
\[
 |v(t)|\le(1-\lambda t^2/2+C\Lambda^2t^4)|v(0)|
  +(1+C\Lambda t^2)\int_0^t(t-s)|T_\varphi f(q^1(s))|\,ds .
\]
Indeed, the average of \(M(s)\) with weights \(2(t-s)/t^2\)
has the same lower bound and norm. Substituting \(v(0)\)
under the Hessian integral gives the displayed contraction
by squaring its norm. One further use of the integral equation
bounds the error from \(v(s)-v(0)\) by the fourth-order term
and the indicated increase of the force bound.

Refresh the common momentum every \(t\) units of time.
The marginal position laws stay invariant, so the mean of
the force integral equals \(t^2\nu_1^\epsilon(|T_\varphi f|)/2\).
Iterate the displayed contraction and use the resulting bound on
the mean weighted separation for a Lipschitz extension of \(G\).
Letting \(t\downarrow0\) at fixed \(\epsilon\) bounds the
observable difference by
\(L_G\nu_1^\epsilon(|T_\varphi f|)/\lambda\). As
\(\epsilon\downarrow0\), the smooth measures converge to the
Gibbs laws on the shrunken balls. The walls have quadratic
growth outside a bounded set, which gives an integrable
majorant for both densities and their first moments for all
sufficiently small \(\epsilon\). The force has at most linear
growth and equals the original force on the shrunken balls.
Dominated convergence therefore applies also to its mean.
Finally let the shrunken radii increase to the original radii.
The bounded reference Hessian and difference force give bounded
gradients on the original bounded convex window, so the last
limit is also dominated.
This proves \eqref{cl:weighted-force-comparison}. Only the
reference Hessian was used for coercivity.

Now assume the same coercivity bounds for all weights with
Lipschitz constant at most \(\mu\). Retain the read coordinates
\(u=y|_S\) and integrate over \(R=I\setminus S\). Write \(m_a\) for their marginal
densities and \(\nu_a^u\) for the conditional laws on \(R\).
The product window makes the integration domain independent
of \(u\), so for \(s\in S\)
\[
 \partial_s\log(m_1/m_0)
  =-\nu_1^u\bigl(\partial_s(H_1-H_0)\bigr)
    -(\nu_1^u-\nu_0^u)\bigl(\partial_s H_0\bigr).
\]
The first term costs \(\delta_s\). In the second, apply
\eqref{cl:weighted-force-comparison} to each unit colour
direction of \(\partial_sH_0\), with
\(\varphi(b)=-\mu d(b,s)\) on \(R\). The weighted
coercivity survives compression. The Lipschitz constant
is at most
\[
 L_s=\sup_{y\in D}
  \left(\sum_{b\in R}e^{2\mu d(b,s)}
                         \|H_0''(b,s)\|^2\right)^{1/2}.
\]
Thus the derivative norm is at most
\(\delta_s+\lambda^{-1}L_s
(\sum_{b\in R}e^{-2\mu d(b,s)}\delta_b^2)^{1/2}\).
Integrating along segments in the product of read balls
bounds the oscillation of \(\log(m_1/m_0)\) by the sum of
these bounds times the corresponding ball diameters.
Since both marginals have mass one, this oscillation
bound \(\Xi_j\) gives \(e^{-\Xi_j}\le m_1/m_0\le e^{\Xi_j}\).

If the forces obey \eqref{eq:clean-conditional-law-CL5}
and the weighted Hessian column bound makes \(L_s\) uniform,
the triangle inequality
\(d(b,s)+d(b,\mathcal B_{\rm ext})\ge d(s,\mathcal B_{\rm ext})
\ge\mathsf D\)
leaves a summable exponential
after extracting a smaller factor \(e^{-c'\mathsf D}\).
The derivative bound at a read coordinate is then
\(C b_j e^{-c'\mathsf D}\). Multiplying by the diameters
and the read-coordinate count gives
\(\Xi_j\le C n_j b_j^2e^{-c'\mathsf D}\).
The normalized density bound persists under the positive
exterior mixture when each exterior is compared with the
same reference density in the common read coordinates.
The readout discrepancy \(\eta_G\) remains a separate term.

\paragraph{The readout as the retained datum varies.}
At a fixed output of the step, the exterior mixture keeps the
same new field \(V_j\). The following comparison is uniform
as that field varies between outputs. Write \(\mathsf r=L^s\mathsf d\)
for the separation in witness-level units and \(p_0=p_0(g)\).
Put \(V_j\) in the axial
gauge of the nonwrapping comparison box and extend that gauge
constantly on each block to the previous level. Denote the
rotated free coordinates by \(\widetilde y\).
If the box has side \(l\) in its own units and plaquette
bound \(\kappa_j\), its logarithm satisfies
\(\|A_j\|_\infty\le Cl\kappa_j\). This follows by writing
the axial link as a product of at most \(Cl\) conjugated
plaquettes. For the intermediate background
\(V_*=\mathcal M^{j-1}U_j(V_j)\) in its block axial
representative, the contour variables are one, and the actual
one-step average is \(V_j\). Apply
\cref{lem:averaging-1a-1-9} against the trivial field on
the full read double blocks. The intermediate plaquettes
are bounded by \(C\kappa_j\) by the local averaging estimate.
Consequently its logarithm \(A_*\) in the same frame obeys
\(\|A_*\|_\infty\le C(l+1)\kappa_j\).

Let \(\mathsf C^0=D\Phi_1(0)\) be the flat pivot inclusion
on those blocks. On the joint chart
\eqref{afs:actual-pivot-chart}, \(D\Phi_{V_*}(0)-\mathsf C^0\)
has norm at most \(C\|A_*\|_\infty\) in that frame.
The chart remainder is quadratic in the free variables.
Multiplying the physical fluctuation by \(V_*\) and using
the local logarithm of the product gives
\[
 \|A_{j-1}-A_*-\widehat g_{j-1}\mathsf C^0\widetilde y\|_\infty
 \le C\bigl(\widehat g_{j-1}b_j\|A_*\|_\infty
                       +\widehat g_{j-1}^2b_j^2\bigr).
\]
Here \(A_{j-1}\) is the logarithm of the actual fluctuating
field in the chosen frame. The first error includes the
commutator from multiplying the two exponentials.

For a bare plaquette \(p\), denote the localized background
plaquette minus one at level \(j-1\) by \(P_p^{(j-1)}(A)\).
On its admitted fixed-radius datum ball this readout is
bounded by \(C\eta_{j-1}^2\). Its derivative at zero is
\(i\mathcal H_{j-1}^{\square}\), the linear curvature response
of that same localized background. Cauchy's estimate on a
smaller ball and the preceding inequality give
\[
 \begin{split}
 &\left|P_p^{(j-1)}(A_{j-1})-P_p^{(j-1)}(A_*)
       -i\widehat g_{j-1}(\mathcal H_{j-1}^{\square}
                             \mathsf C^0\widetilde y)(p)\right|\\
 &\hspace{16mm}\le C\eta_{j-1}^2
       \bigl(\widehat g_{j-1}b_j\|A_*\|_\infty+\widehat g_{j-1}^2b_j^2\bigr).
 \end{split}
\]
\paragraph{The two localized readouts of the retained field.}
The actual retained datum of step \(j\) includes \(V_j\) and
the lower rows that remain stored. On a fixed output, \(V_j\)
is common to the exterior mixture. The regular critical field
\(U\) at an exterior in that mixture still depends on the
exterior. In the interior the literal shift in the fluctuation
integral uses \(V_*=\mathcal M^{j-1}U\) in block axial gauge,
so \(\mathcal M(V_*)=V_j\) there. On the remaining
level-\((j-1)\) rows keep their stored values, and keep the
older rows. This is the datum induced on the truncated
sequence by \(U\). Indeed the retained rows already prescribe
their averages, and the newly exposed rows read the average
\(V_*\). Each equality uses the full double blocks.
A variation at fixed induced datum extends by \(U\) through
the frozen collar. The finer constraints imply the old
ones by induction in the averaging, and the fixed inner bonds
of crossing plaquettes make their contribution constant.
Thus \eqref{mb:restriction-identity} identifies the critical
field of the truncated problem with \(U\) up to residual gauge.
This argument keeps the same cut rows and averaging scheme as
the step.

For \(q=j-1,j\), let \(D_{q,p}\) be the domain of the
localized minimizer used to read \(p\). At \(q=k_0\) this is
the prescribed neighborhood of the test cube of \(p\),
including its four outer layers. At larger \(q\) take the
box around the two supports with the prescribed order
\(\mathsf D\) margin in level-\(q\) units, rounded out on
the collar grid. Include all top blocks of the plaquette
read and all the double blocks of the retained rows needed
by the pivots. The latter cost \(O(L)\) extra spacings on
the level-\((j-1)\) free-bond grid. The lower box and these
double blocks are in the read set of step \(j\), with the
extra \(\mathsf D\) margin to its integration-box edge.
The upper box reads \(V_j\), which is already fixed in this
conditional integral.

Use the clean neighborhood with enough room for these
fixed enlargements and their collars. In level-\(k_1\)
units the distance from the pair to the farthest of these
enlarged boxes is at most \(C(\mathsf r+\check R_{k_0})\):
the side at \(q\) is at most
\(C(\mathsf r L^{-(q-k_0)}+\mathsf D+\check R_{k_0})\)
in level-\(q\) units, and \(q\le k_1\). The same bound
allows a further margin of this size in level-\(k_1\)
units by increasing the fixed neighborhood constant.
The clean-history union bound counts cells meeting the
enlarged neighborhood in the same way. On this event
there are no standing regions or operations on these
boxes and their padded collars, at the two levels of the
step.
The padded collars lie in the ambient top layer, both
for the full problem and for the problem truncated at
level $j-1$.

On the padded boxes, the passed coarse tests bound the
plaquettes of the stored field $V_j$ by $Cgp_0$. For
$V_*$ use the corresponding bound on the full backward
cells of the ambient background and the one-step averaging
estimate. This is a bound on the averaged critical field
at zero fluctuation. In obtaining the local profile of
the ambient background, the fixed ceilings on data outside
the clean neighborhood enter with the exponential weight
from their distance to the padded boxes. The intervening
homogeneous layers have thickness at least
$c(\mathsf r+\check R_{k_0})$ in the native spacings.
Their contribution is bounded by $Ce^{-c(\mathsf r+\check R_{k_0})}$,
which is smaller than $gp_0$ in the short-distance range.

For each $q=j-1,j$, compare the local critical field with
the restriction of the same ambient field $U$.
The common rows are those of $V_*$ for the lower comparison
and of $V_j$ for the upper comparison. Each localization
uses the face lift of its own prescribed datum. The
alignment in \eqref{mb:real-collar-datum} is performed
on each collar using the $Cgp_0$ local profile. Thus the
collar source is of that size. In \eqref{mb:real-collar-norm}
the source size enters linearly, while the fixed regularity
ceilings admit the response path. At a plaquette $p$ the
result in the response gauges is bounded by
$C\eta_q^2gp_0e^{-cd_{q,p}}$, where $d_{q,p}$ is its
distance to the localization boundary in level-$q$ units,
with the fixed read stencils included in the margin.
Write $\Delta_j=\min_{q,p}d_{q,p}$. The native
test-cube and box margins give
\[
 \Delta_j\ge c_0\min\{\mathsf D,\check R_{k_0}\}-C_0.
\]

For a matrix comparison, express each plaquette defect
in the frame at its level-$q$ ancestor. This is conjugation
by the product of the averaged contours from the ancestor
to its base point. Use the coarse axial frame of the
common box as the input frame, extended in the block axial
representative for $V_*$. %
%
On the complete order-\(q\) double blocks of the ancestor
read, write the collar comparison in its aligned real
frame as \(U'=\exp(i\eta_q A)U\). The potential bound on
these blocks gives
\(a:=\|A\|_\infty\le Cgp_0e^{-c(d_{q,p}-C)}\).
The fixed enlargement in the exponent is the size of
the full blocks in level-\(q\) units. The straightened
logarithm and its node factor at level \(n\le q\) obey
\[
 \|a_n^{\rm str}\|_\infty+\|v_n^{\pm1}-1\|_\infty
       \le C L^n\eta_q a .
\]
Apply the one-step contour estimate on the complete
block of each contour from the next ancestor to the
current one. Covariance with \(v_n\) gives a change at
most \(CL^n\eta_q a\). The product of real unitary
transports changes by at most the sum of the changes,
which is bounded by \(Ca\). In the collar comparison
the gauge at the top is \(v_q^{-1}\), and this differs
from one by at most \(Ca\) as well. The other gauge
factors cancel between successive contours.

For a plaquette defect \(P\), changing the transport
to the top therefore changes its conjugation by at
most \(C|P|a\). The curvature bound gives
\(|P|\le C\eta_q^2 gp_0\), and the response bound
already controls the change of \(P\) itself. Together
they give \(C\eta_q^2gp_0e^{-c d_{q,p}}\) for the
change in the ancestor-frame defect.
On $U$ the two
ancestor-frame reads agree. The additional contour from
the level-$j$ ancestor to the level-$(j-1)$ ancestor is
the identity in the block axial representative used to
form $V_*$.

We use this comparison for the plaquettes in their
localized Landau representatives. On each actual local
datum ball in the common input frame, its logarithmic
datum, including the face lift, is bounded by $Cl\kappa_j$.
The normalized potential on the ancestor block has this
bound by the weighted response from the flat datum.
The same contour length sum makes its transport differ
from the identity by $Cl\kappa_j$. Thus each conjugation
changes a defect of size $C\eta_q^2\kappa_j$ by at most
$C\eta_q^2 l\kappa_j^2$. This argument uses only paths
in each localized tower, even when the initial test
cubes are disjoint. The common coarse axial gauge has
one residual constant rotation, which acts on both
plaquettes and leaves the flat Gaussian law invariant.
These errors, together with the two collar errors,
bound the difference of the lower and upper readouts
at zero fluctuation by
\[
 C\eta_{j-1}^2\bigl[gp_0e^{-c\Delta_j}
                          +(l\kappa_j)^2\bigr].
\]

Thus, with \(l\le C\mathsf r\), \(\kappa_j\le Cgp_0\)
and \(b_j\le Cp_0\), the remainder after retaining the
level-\(j\) background and the flat linear fluctuation is
bounded by
\[
 C\eta_{j-1}^2
       \bigl[(\mathsf r g p_0)^2+gp_0 e^{-c\Delta_j}\bigr].
\]
The term involving $\Delta_j$ uses both native margins.
For $\mathsf D=A\log x$, with $A$ fixed, the initial test
margin dominates $\mathsf D$ when $x$ is sufficiently large.

\paragraph{The common part of the two operator recipes.}
For a fixed step use the same tiling, masks and guard
prescription in the two boxes. When comparing a row
\(b\), include in a carrier every row of a local coefficient,
each endpoint double block of its pivot, its full current
stencil, the contour reads, and any held guard used in
forming its local Gaussian coefficient. A carrier
contained in the common interior uses the same functional
recipe on a common admitted pair domain. This follows
term by term: its masks are those of the empty standing
trace, its pivot is the same local analytic solution,
and its walks and guard parameters are specified by the
same kit. For the Gaussian root this is the fixed-root
restriction of \eqref{haf:actual-decoupled-root}. In
particular, one must preserve the held guard before
restricting a Gaussian functional. Independence of the
summed logarithm from auxiliary guard choices is an
identity after summation, which does not identify terms
filed with different guards.

With \(d_b\) the distance to the edge of the common
interior, take carriers lying within \(d_b/2\) of \(b\)
including their guards. They are therefore indexed in
both recipes. Any term left out of this matching either
reaches that edge or has a read or a guard reaching it.
Such a term has charged length at least \(c d_b-C\)
for the intrinsic and precision families; the fixed
guard widths only change \(C\). The absolute rooted
majorants give a row sum \(Ce^{-c'd_b}\) for the unmatched
terms. On each matched carrier use the actual difference
of its two potential and current reads before taking
Cauchy bounds. For comparison at the same retained
datum this difference is the exponential collar error.
For comparison with the flat datum it includes the
small background displacement \(C\mathsf r gp_0\).
This yields the short-carrier estimates in the real
precision and intrinsic force comparisons. Regular old
activities also require their own retained-datum radii
and origin estimates, and standing old boundary terms
use \eqref{ccf:old-boundary-tail}.

\paragraph{The Gaussian comparison in the rotated coordinates.}
\label{f2:flat-box-marginal}
%
%
The rotation of the actual free variables uses the
block-constant extension of the retained axial gauge.
In this frame let \(C_{\rm ax}\) be the old block-axial
centre field. Let \(u\) be the Landau gauge of the
comparison, restricted to the old sites and with
the same residual normalization at the tops.
By covariance of the average the old field in the
Landau frame is \(C_{\rm ax}^u\). If the Landau
potential and its first difference have size
\(\epsilon\) in the upper units,
the one-level average estimate bounds each of its
old link defects by \(CL^{-1}\epsilon\). For a site
\(x\) and its block centre \(y\), the symmetric mean
of the complete block contours satisfies
\[
 T_{y,x}(C_{\rm ax}^u)=u(y)u(x)^{-1}.
\]
This uses \eqref{agm:contour-gauge-law} and the axial
condition \(T_{y,x}(C_{\rm ax})=1\). Each path has at
most \(d(L-1)\) links. Telescoping the path products
and applying the local mean estimate bounds their
symmetric mean's distance from the identity by
\(C\epsilon\).
At the centre \(u(y)\) is the inherited node factor
of the localized response, or its inverse according
to the direction of the transfer. Its normalized
node estimate gives \(\lvert u(y)-1\rvert\le C\epsilon\).
Together these estimates give the same bound on
\(\lvert u(x)-1\rvert\) at every free site of the
complete block. Here \(\epsilon\le C\mathsf r g p_0\)
on the admitted localized core.

On the compression to the admitted blocks, or on
short carriers with all their endpoint blocks there,
the native precision and covariance transform
by the bondwise adjoint matrix \(R_u\) at the
initial sites. For either matrix \(A\), a spatial
weight \(T_\varphi\) commutes with \(R_u\), and
\[
 \bigl\|T_\varphi(R_u^{-1}AR_u-A)T_\varphi^{-1}\bigr\|
 \le 2\|R_u-1\|\,
              \bigl\|T_\varphi A T_\varphi^{-1}\bigr\|.
\]
The same estimate holds in the weighted row and
column sums, by applying the bound blockwise to
the two endpoint rotations. The return to the
actual read frame therefore costs \(C\epsilon\).
It uses the matched masks and full block contours.
The other carriers retain the separate exterior decay.
In the changed gauge use the relative contour
\eqref{eq:explicit-relative-current-extension-NX4}
at the transformed centre. It is conjugated at the
block node, so its quadratic weight is transported
with the same free-bond mask.

On short carriers inside the box, the rotated precision
and the flat one use the same local recipe and the same
ratio \(\tau_{j-1}=x_{j-1}/\theta_{j-1}\) in the Wilson
quadratic. The contour-gauge weight retains its original
normalization. At the flat pair the actual and straight
linearized averages give the same constrained exact
curvature space. In that space the action metric at
this level is \(\mathsf S_{j-1}^+\), obtained from the
energy of the flat minimizer. The soft contour weight
integrates to a constant on each curvature fibre,
since its contour values translate along the block
gauge orbit. Multiplying the curvature metric by
\(\tau_{j-1}\) gives the covariance factor
\(\tau_{j-1}^{-1}\). This is the argument of
\eqref{gcb:constrained-curvature-projection} with
the multiscale action metric. In physical units
\(\widehat g_{j-1}^2\tau_{j-1}^{-1}=x_{j-1}^{-1}\).

\paragraph{Matching the precision at a fixed exterior datum.}
Use the free stock of the step before the strip substitution,
with the middle pivot of every eliminated coarse constraint
removed. The interior integration stock includes the free
bonds with at least one endpoint in the small-field region.
In the native prepared chart, the operator at the physical
current is \(C_V^{\mathsf T}\Delta^{(j-1)}C_V\) on that
full stock \cite[(2.3)--(2.11)]{B87}. The identification
uses the perfect action on the actual lower determining
set. Its Hessian is the first quadratic term. The change of
variables in the plain outer constraint contracts its second
derivative with the current of that perfect action. The
relative contour square supplies the remaining term.

The same algebra identifies the symmetric-contour form on
an admitted compression of a physical graded branch. On
the active top datum, keep the finer arguments of the
outward node factor fixed. The relative datum then has
derivative equal to the active injection and second
derivative zero on those directions, as in
\eqref{px:active-datum-flatness}. Consequently the reduced
action Hessian is the compression of
\((QG_1Q^\sharp)^{-1}-a\) by
\eqref{px:physical-row-readout}. At the physical current
\(j_U=Q^\sharp\omega\), both right lifts have adjoint
readout \(\omega\), since they are right inverses of the
same \(Q\). Convert this Hessian and the coarse current to
the raw bond coordinates by \eqref{px:coarse-raw-duality}.
In the actual next pivot graph the constraint contraction
uses \(h_V^\sharp\widehat J\), where \(\widehat J\) is
that raw coarse current. Equation
\eqref{px:next-constrained-hessian} therefore gives the
outer action Hessian including its curvature term.
Adding the relative contour Gram in the same free
coordinates gives \eqref{px:complete-next-precision}.
This is the compressed physical value of the completed
form when the two constructions use the same branch,
active mask and contour gauge.

Freezing exterior fluctuation coordinates in this formed
quadratic leaves its principal interior block as the
conditional precision. If the full form is \(K\) and the
fixed coordinates are \(a_{\rm E}\), its mean is
\(-K_{\rm II}^{-1}K_{\rm IE}a_{\rm E}\).
The exterior Schur term is independent of the integrated
coordinates. For the comparison Gaussian, form the full
precision at the reference datum with the exterior
fluctuation set to zero. Its Gaussian marginal on the
box has covariance \((K^{-1})_{\rm II}\) and precision
\(K_{\rm II}-K_{\rm IE}K_{\rm EE}^{-1}K_{\rm EI}\), with
\(K\) now evaluated at that reference datum. This accounts
for the Schur correction in the force comparison. The
interactions and cutoffs enter the normalized comparison
separately.

At unit ratio use the inverse walks of this matched step
precision on each of the two nests. The native operator
theorem gives this identification and inversion on its
admitted prepared geometry. The homogeneous
symmetric-contour construction gives it by
\eqref{neg:off-current-outer-completion} and
\eqref{acg:actual-marginal-identity}. For a general
labeled nest in the symmetric construction we must admit
the same completed carrier families on the actual free
stock, including the separation and current bounds of
\cref{prop:ncw-named-families}. The local algebra above
identifies their physical matrix wherever this formation
is admitted. The two-level real representation gives, on the
padded top territory, a rotated potential of size
\(C\epsilon/L\), its first difference of size
\(C\epsilon/L^2\), and scaled current of size
\(C\epsilon\), all in input units. Let \(\epsilon_\star\)
bound the complete pair in the local operator norm on this
territory. Here \(\epsilon_\star\le C\mathsf r g p_0\).

The clean geometry identifies the top territory, with
its free and pivot indices, in both nests. A walk whose
enlarged carrier lies there is consequently assigned the
same kernel \(T_\omega\) in both expansions. In the local
extension use \(\chi_\omega=(1-(2L/M)d_{\rm w}(\,\cdot\,,[\omega]))_+\),
where \(d_{\rm w}\) is the operator's weighted distance and
\([\omega]\) is the full read. Thus the cutoff is one on
the read and its support lies in the padded territory.
If \((A',J')\) is the rotated pair based at the identity,
the local ball for the operator admits
\[
  z\longmapsto
  (e^{iz\eta_{j-1}\chi_\omega A'},\,z\chi_\omega J'),
  \qquad |z|<R_\star:=\frac{\varepsilon_1}{2\epsilon_\star},
\]
where \(\varepsilon_1\) is the fixed operator radius.
The current on this disc is the independent coordinate of
the operator. At \(z=1\) locality gives the physical kernel
on this walk, and at \(z=0\) it gives the flat kernel.
Weak coupling on the comparison window gives
\(R_\star\ge2\). If \(\epsilon_\star=0\) the short-walk
values coincide. Otherwise summing the Cauchy bounds for the
positive Taylor orders gives \(2/R_\star\) times the
walk majorant. The sum of the rooted majorants is bounded.
For a walk whose enlarged carrier exits the territory,
the tree from a read bond to that exit has length at least
the distance to the edge minus the fixed padding allowance
of the walk domains. The two outer sums are therefore
\(O(e^{-c\mathsf D})\). Returning to the read frame costs
\(C\epsilon\) by the endpoint estimate above. Thus at unit
ratio the entrywise comparison follows directly for the
full covariances restricted to the read set.

For the ratio tube use the affine precision
\(\tau(K_A+K_G)+(1-\tau)K_G\) from
\eqref{agm:complete-tau-quadratic}. Take a common
compression \(I\) on complete blocks between the read set
and the edge of the matched territory, with a fixed
fraction of the margin on each side. Retain part of the exponential walk rate
for the distance between the two endpoints in the
preceding estimates. The weighted row and column sums
then bound the difference of the compressed unit-ratio
inverses by \(C(\epsilon_\star+e^{-c\mathsf D})\).
The local relative gauge forms on \(I\) differ by
\(C\epsilon\) in the same norms. This follows from
the fixed-radius contour derivatives on the matched
complete blocks.

Write \(C_a\) for a unit-ratio inverse on \(I\), \(G_a\)
for the corresponding compressed gauge form, and
\(\lambda=(1-\tau)/\tau\). Its inverse at ratio \(\tau\)
is \(\tau^{-1}(1+\lambda C_aG_a)^{-1}C_a\).
The weighted row and column norms are submultiplicative.
The tube makes \(|\lambda|\|C_aG_a\|\) small there.
After taking out the common \(\tau^{-1}\), the
difference has a middle factor between
\((1+\lambda C_1G_1)^{-1}\) and
\((1+\lambda G_2C_2)^{-1}\). That middle factor is
\[
  C_1-C_2-\lambda C_1(G_1-G_2)C_2 .
\]
This gives the same bound on \(I\) at the actual ratio.
To return to the full covariance, use the corresponding
weighted inverse bound on each full free-bond stock.
For a full precision \(Q_\tau\) and its inverse
\(\Gamma_\tau\), block multiplication gives
\[
  (\Gamma_\tau)_{II}-(Q_\tau)_{II}^{-1}
    =-(Q_\tau)_{II}^{-1}(Q_\tau)_{I I^c}
                           (\Gamma_\tau)_{I^c I}.
\]
%
At a real admitted ratio, symmetry and the other
block equation give the factorization
\[
 (\Gamma_\tau)_{II}-(Q_\tau)_{II}^{-1}
       =B(\Gamma_\tau)_{I^c I^c}B^*,
 \qquad B=(Q_\tau)_{II}^{-1}(Q_\tau)_{I I^c}.
\]
For a linear read supported in \(I\), the weighted
inverse and precision bounds give an \(\ell^2(I^c)\)
tail for its row of \(B\). The tail spends its
exponential rate on the distance from the full
read to \(I^c\). Include the contour stencil and
the plaquette stencil before measuring this
distance. If the two reads have depths \(d_1,d_2\)
after that padding, the bounded exterior
covariance gives an error \(Ce^{-c(d_1+d_2)}\)
for unit read norms. This implies the required
\(Ce^{-c\mathsf D}\) bound on the buffered read set.
The endpoint rotations return this estimate to the read frame.
The ratio comparison thus requires the native
unit-ratio inverse and the affine gauge correction on
the actual free stock, with the same pivot elimination.
The resulting bound in scaled free coordinates is
\[
 \varpi_j\le C\bigl[\mathsf r g p_0+e^{-c\mathsf D}\bigr]
\]
against the flat covariance. This comparison uses the
complete precision, including the current and contour
gauge terms, in the common local recipe. In a physical
plaquette pairing each linear readout also carries its own
factor \(\widehat g_{j-1}\eta_{j-1}^2\).

Let $C_0,C_1$ be the flat and rotated Gaussian covariances on
the read coordinates, and put $C_t=(1-t)C_0+tC_1$.
The uniform precision bounds give common positive lower
and upper eigenvalue bounds for $C_t$, including after
compression to the read coordinates. If $h_t$ is its
centered Gaussian density and $D=C_1-C_0$, then
\[
 \partial_t h_t(y)=\tfrac12 h_t(y)
 \bigl(y^{\mathsf T} C_t^{-1}D C_t^{-1}y
                         -\operatorname{tr}(C_t^{-1}D)\bigr).
\]
The bracket has mean zero under $h_t$ and second moment
$2\|C_t^{-1/2}D C_t^{-1/2}\|_{\mathrm{HS}}^2$, by Wick's
formula. Integrating the $L^1$ bound on the derivative
gives a density difference at most $C\|D\|_{\mathrm{HS}}$.
On the common product window, normalization costs at
most another fixed factor when its Gaussian masses are
at least one half. The scalar read dimension is at most
$3n_j$, so the entrywise bound above costs $Cn_j\varpi_j$
in the normalized Gaussian comparison. The nonlinear
box tilt and exterior comparison enter separately.

\paragraph{The change in the intrinsic force.}
For the intrinsic terms of \cref{cl:intrinsic-reference-hessian},
take the same zero-exterior fluctuation argument for both
backgrounds. Put \(a_j=C\mathsf r g p_0\), which bounds the
rotated background potential on the short carriers. On
their common fixed-radius chart, the value bounds
\(CB_Y g b_j^3\) for the action remainder and
\(CB_Yg b_j\) for the coarea term also hold on the larger
fluctuation ball. Cauchy in a fluctuation variable costs
\(C/b_j\). Cauchy in the background direction on its fixed
radius, followed by integration to the flat datum, costs
\(Ca_j\). Thus the short-carrier part of the force difference
at bond \(b\) is at most \(C(gb_j^2+g)a_j\) after the
rooted sum. For the long carriers take the two force bounds
separately and use the distance to the edge as in the
precision comparison. The resulting bound is
\[
 C(gb_j^2+g)\bigl[a_j+e^{-c d(b,O)}\bigr].
\]
The additional comparison of nonzero exterior fluctuations
uses the mixed Hessian estimate on the admitted segment.
The old terms need their own retained-datum radii and
origin sums. The fixed-radius derivative used here applies
to the intrinsic terms.

%
%
\paragraph{A common flat reference for the exterior mixture.}
\hypertarget{f2-direct-flat-density}{}
Use the same flat precision on the common free-coordinate
box \(I\) and the same product window for every admitted
exterior. In \eqref{cl:weighted-force-comparison}, take the
reference potential to be the quadratic of this flat law.
The weighted gap of its precision gives the required
coercivity and the weighted column bound. The reference
has no nonlinear tilt or exterior linear force.

First set aside the old-origin exponent, whose oscillation
costs at most \(2\epsilon_j^{\rm old}\) in the normalized
log density. Keep the actual intrinsic action and coarea
terms in the other potential. Write \(a_j=C\mathsf r g p_0\)
as in the precision comparison. In the rotated coordinates
the difference of quadratic forces is bounded on a row by
the precision row estimate times \(b_j\). The Schur and
exterior forces have the row bounds at \(O\) and \(E\).
The intrinsic value bounds on the larger ball, followed by
one coordinate Cauchy estimate and the rooted sum, give
\(C(gb_j^2+g)\) for its force on that row at the actual
fixed exterior. Thus, with \(B=\partial I\cup O\cup E\),
the total force difference after the old-origin removal is
bounded by
\[
 \delta_b\le Cb_j\bigl[a_j+gb_j+e^{-c d(b,B)}\bigr]
 \qquad(b_j\ge1).
\]
This uses the complete precision row estimate, including
the matching of the gauge weight and the current stencil.

At a read coordinate \(s\in S\), the marginal log-density
estimate following \eqref{cl:weighted-force-comparison}
weights the remaining force components by
\(e^{-\mu d(b,s)}\). The two constant terms sum uniformly,
and the last term gains \(e^{-c'\mathsf D}\) when
\(d(s,B)\ge\mathsf D\). The read-window diameters are at most
\(Cb_j\). Restoring the old-origin exponent, the normalized
read marginal therefore has density ratio to the same flat
box-window marginal in \([e^{-\Xi},e^\Xi]\), where
\[
 \Xi\le2\epsilon_j^{\rm old}
        +C|S|b_j^2\bigl[a_j+gb_j+e^{-c'\mathsf D}\bigr].
\]
The pointwise ratio passes through the positive exterior
mixture in these common coordinates. This route only uses
the Hessian of the flat reference in the force lemma.
The actual chart, rotated readouts and selectors must
still be admitted on the common product window. Removing
the flat box window and comparing with the full flat
marginal use the Gaussian tail and Schur estimates above.

\paragraph{Admission on the real conditional window.}
For each positive string entering the exterior mixture,
use the completed label of its averaging step before
\(\mathbf R\) at that level. Keep the actual stored
fields, frames and cut rows of that label. In particular
a component that will be treated later by an exterior
operation is still counted in this exponent, at the
stage where the fluctuation integral is performed.

At a point of the support of this step, put
\(\mathcal A_k=\widehat g_k(y,y_O)\) on the nonpivot bonds that
touch \(\Omega_{k+1}\). In the strip use the corresponding
cut datum, background and shifted fluctuation. For every
carrier \(Y\) the step
field facts give
\[
 (U_{k+1},0,J(U_{k+1}))|_Y
        \in \widetilde{\mathcal U}^{1/2}_{\sigma^{\mathbf T}}(Y).
\]
Here \(\sigma^{\mathbf T}\) is the graded structure of
the completed label. In addition the free variable is
in the half-radius version of the step box
\[
 \left\{\begin{aligned}
 &\widehat g_k^{-1}|\mathcal A_k(b)|<3\mathfrak b_k
       &&\text{on free bonds touching }\Lambda_{k+1},\\
 &\widehat g_k^{-1}|\mathcal A_k(b)|<2C^{\mathrm{fl}}_1\mathfrak b_k
       &&\text{on the other free bonds of }\Omega_{k+1}.
 \end{aligned}\right.
\]
In this display the bond types are those of the step's
integration and exterior sets. The full product of this
box with the graded field domain is the domain used by
the old-origin substitution. The bound uses
\(|\mathcal A_k|<\delta_k=\widehat g_k\mathfrak b_k\) on the
integration bonds and the bound by
\(\mathsf c_{\mathrm{fl}}\widehat g_k\mathfrak b_k\) on the other
free bonds, or \(\mathsf c^\sharp_{\mathrm{fl}}\widehat g_k\mathfrak b_k\)
after the strip change. Choose \(C^{\mathrm{fl}}_1\)
above these two constants.

The first factor has four requirements. The passed
plaquette and fluctuation tests, and the retained tests
on the older or prepared rows, give the regularity of
the completed datum layer by layer. The bound includes
the prescribed shell factor at a prepared component.
Near \(\Omega_{k+1}\) that factor is one, and its far
contribution to the profile estimate is bounded by
\(C(\log x_k)^r e^{-cLM\check R_{k+1}}\).
The layer estimates of the regular critical field then
give the pointwise plaquette and current bounds in
their own units. The current is that of the restricted
action, including its prescribed boundary stencil.
The complex potential slot is zero. On each addressed
gauge cube the local transport estimate gives the
required real gauge with potential and first difference
bounded by the same layer size times the fixed
cube-size constant. For each addressed tower, restrict
the real critical configuration with its own averages.
The restriction identity and the inherited layer
bounds give the tower and its plaquette and current
estimates in the tower units. On a core the same
restriction with the core scale gives its two radius
bounds. These are the four clauses of the step
membership criterion. Their half-radius form uses
\(C\varepsilon_i^{[r]}\le c_1\check\alpha_{0,i}\) at
each layer. Here \(\varepsilon_i^{[r]}\) includes its
actual layer factor, \(\check\alpha_{0,i}\) is the
minimum of the two adjacent radii (the single radius
at level zero), and \(2c_1\) is below the fixed lower
bound on the graded coefficients. The test thresholds and radii are
chosen in that order. This also admits the structure
of any earlier member addressed by the label.

On the clean interior the exact product-window
identity supplies these support points when \(y\)
ranges over \(\mathcal W_I\). Its old-test implication
uses the same coarse datum on full double blocks,
the margin from \(\delta_k\) to the approximate test
at \(2\delta_k\), and the old plaquette bound from
the factors \(L^{-2}\) and \(\delta_k/\varepsilon_k\).
Exterior conditions are those of the chosen compatible
string. The reference branch retains the same step
strings and tests. At a fixed output its normalized
weight relative to the clean kernel is bounded above
and below by positive constants. Thus the test-support
properties carry to that branch under disintegration.
For a conditional oscillation estimate the resulting
almost-everywhere admission suffices. A bound for a
larger complex product of varying exterior data is a
different assertion.

Finally the actual step substitution on this domain
has the pivots of \(B'\) solved before it is used in
any older carrier. At complete parameters its averages
are \(e^{iB'}V_k^{\mathrm{bg}}\), its second member is
the actual current of \(U^\sharp\), and it preserves
the prescribed transports at sites with stored data.
The older datum induced by the literal fluctuation is
regular with half the margin on the support of the
tests. This is the regular evaluation clause for the
outgoing background of the fluctuation integral. The
identity by origins therefore applies on this domain
to both the old exponent and its local evaluated
potentials.

\paragraph{Which remainder is counted in a window.}
An ordinary \(R^{(i)}(X)\) enters the old remainder
sum only at its declared action-region place. The
remainder of the fluctuation integral is assigned
there when the entire domain \(X\) is at least one
level-\(i\) cell inside the corresponding small-field
region. The remainder produced by healing is assigned
there on the same one-cell interior of its output
action region. The term is a function of the core
pair at that place. A remainder of the fluctuation
integral in the rim, and a term of the healing
exponentiation outside its prescribed interior, are
filed as boundary terms. Membership in a larger clean
box at a later time does not change that filing. An
earlier term can enter the ordinary sum after healing
only through the actual identity for the refiled
term.

For an ordinary \(R^{(i)}(X)\), apply the window
function of its admitted old core, subtracting its
value at \((\mathbf 1,0)\). If the window is at scale
\(m\ge i\), write \(\sigma=L^{-(m-i)}\) and \(\rho\)
for the ratio of output to old radii. Its contribution
is bounded by
\[
 Cr_i e^{-\kappa d_i(X)}(2+d_i(X))^2 \rho^2\sigma^4.
\]
This is the vacuum subtraction estimate. At a fixed
old base the rooted domain sum absorbs the polynomial.
The \(C\sigma^{-4}\) possible old bases in the current
window cancel the scale factor. Thus the contribution
per window is \(C\rho^2 r_i\), with no additional
geometric gain in the age. The far members stay at
their original places with their rooted decay, as in
the step's remainder potential bound. Summation in
\(i\) is the flow sum \eqref{eq:local-formation-OA22}.
Neither a remainder at the rim nor an exponent with
stored colour data enters this sum. Their actual
boundary carriers, with the hull family of the label
at evaluation, have the separate boundary and
stored-factor bounds.

\paragraph{Removing the old interaction at fixed exterior data.}
For the generating function it suffices to bound the
oscillation of each old contribution on the real
conditional window. Use the reference branches of the clean
coupling, after removing the separate restriction
weights. Fix a compatible exterior \(e\)
of the output \(\omega\), including its retained datum
and all stored spectators. Write \(k=j-1\) and put
\[
 \vartheta_k=\frac{\delta_k}{\alpha_k}
       \asymp(\log x)^{\mathfrak p_0-q},
 \qquad \delta_k=\widehat g_k b_j.
\]
Use the complete evaluated potentials of the fluctuation
step. The input field in each potential is the pair of the
critical background and its current, with the physical
free coordinates \(\mathcal A=\widehat g_k(y,y_O(e))\).
Its output field \(U^\sharp\) is the actual solution at
the new datum \(e^{iB'}V^{\mathrm{bg}}_k\); \(B'\) includes
the solved pivot entries. At the complete parameters
the levelled substitution reads its actual current and
preserves the stored variables. Thus the sum over a given
origin equals the difference of that origin evaluated at
\(U^\sharp\) and at the unfluctuating \(U^\flat\).

For the regular family after the pointed counterterm
subtraction, the old remainders and the Wilson weight
difference, the step estimates on the admitted real box
give
\begin{equation}
 |V^o_e(Y;y)|\le C\vartheta_k\,S_k^o\,e^{-\lambda d_k(Y)} .
 \label{f2:old-real-potential-bound}
\end{equation}
The bound also holds for an ordinary defect if present,
with its separate size. In this formula the full fluctuation
box has width at most a fixed multiple of \(b_j\), including
the fixed exterior coordinates. The normalization constant
from \(U^\flat\) depends on \(e\) alone. The sizes are
\[
 \begin{split}
  S_k^{E}&=\sum_{i\le k}\omega_{k+1-i} e_i,\qquad
  S_k^{R}=\sum_{i\le k}\rho_i^2 r_i,\\
  S_k^{c}&\le CM^4\sup_m\alpha_{0,m}^2\,\mathsf b_\star .
 \end{split}
\]
Here \(c\) denotes the Wilson lag origin, \(\mathsf b_\star\)
bounds the increments of inverse coupling, and the source
in the generating function is carried as an insertion in
the unsourced step. The
regular sum is at most \(CE_0\) by \eqref{eq:local-formation-OA16}.
For the spectator-free remainders with \(r_i\le Cx_i^{-s}\)
and \(s>2\), \eqref{eq:local-formation-OA22} bounds
\(S_k^R\) by \(C x^{2-s}\). The Wilson size is bounded
uniformly. In its formation the level-\(m\) Wilson
members have size at most
\(C\alpha_{0,m}^2(k-m+1)\mathsf b_\star\); the
spatial and layer sum is already taken in the displayed
bound for \(S_k^c\). Replacing it by the global supremum
of \(x_k(\cdot)-x_k\) would lose uniformity.

The source of the lag in the unsourced density is
\(\sum_{i\le k} b_i(1-\varphi_i)\). When the profiles of
healed regions are redrawn, their changes are compensated
at each bare plaquette in the Wilson sum. The remaining
lag is supported in the standing large-field region.
On the clean box, including its trace window, every lag
potential whose full carrier stays there vanishes, by the
whole-lattice identification of the step potentials.
An exterior lag can still vary with \(y\) through the
background substitution. Use \eqref{f2:old-real-potential-bound}
for that localized origin. Its contributing carriers
reach the nonempty trace, so the spatial rate can also
pay for that distance. The bound without this extra
distance gain already suffices here. With profiles kept
before compensation, the same bound for the evaluated
lag origin is used wherever the step estimate is admitted.
A bound on the volume-weighted lag of an individual
history cannot be attributed to a resummed class unless
the lag is common to that class.

%
%
On the positive history resolution before this
compensation, the weighted lag can be controlled
without excluding all old creations under the box.
Let \(n_j\) be the size of the level-\(j\) window.
Count the older level-\(i\) cells under it whose
profile \(\varphi_i\) is less than one somewhere.
Each enters with weight
\((j-i+1)L^{-4(j-i)}\). There are at most
\(Cn_jL^{4(j-i)}\) cells to test, and a profile
can differ from one only next to a live or
newly created cell at that level. The one-cell
probability bound therefore gives the first moment
of the weighted count:
\[
 \mathbb E\,\mathfrak S_j^{\rm lag}
 \le C n_j\sum_{i<j}(j-i+1)e^{-c p_0(g_i)} .
\]
Below the witness level \(n\), the running lower
bound \(x_i\ge x_n+\lambda_\sharp(n-i)-C\)
makes this series summable. Use half the
exponent at \(x_n-C\), and use the other half to
dominate any fixed inverse power of \(x_n+n-i\).
Between \(n\) and the end of the fluctuation
segment the couplings are comparable to that
of the witness. Its logarithmic number of
levels contributes only a polynomial prefactor.
Markov's inequality at the threshold
\(\tau_\star n_j\), followed by the sum over
the segment, gives an exceptional weight
\(Ce^{-c' p_0(g_n)}\) for fixed \(\tau_\star>0\).
This uses the volume of each fine cell in the
window units. A count of all older test failures
with unit weight would lose that factor.
The bound concerns the positive refinement on
which the lag is recorded. In the compensated
class use the refiled exponent. A restriction by
the finer lag event must use its own joint law.

%
%
On this finer resolution, fix the exterior and
keep the original profiles in the Wilson
exponent. On the admitted complex box for
both matched fields, the curvature estimate
bounds the difference of their plaquette
actions per unit volume in window units by
\(C g^2 b_j^2\). A level-\(i\) cell occupies
\(L^{-4(j-i)}\) of a level-\(j\) cell.
The running bound on the increments of inverse
coupling bounds the missing profile coefficient
there by \(C(j-i+1)\) in the lag count.
Summing the action differences over the
lagging cells therefore gives
\[
 |\mathbb P_j^{\rm lag}|
 \le Cg^2 b_j^2\,\mathfrak S_j^{\rm lag}
 \le Cg^2 b_j^2\,\tau_\star n_j .
\]
On a smaller concentric box the Cauchy
estimate in two directions spends \(Cb_j^{-2}\),
so the Hessian has size \(Cg^2\tau_\star n_j\).
For its spatially weighted estimate the decay
of the two matched minimizer responses is used
on the same profile cells. These estimates read
the lag at fixed exterior on the finer
resolution.

At fixed exterior, subtract the value at one point
\(y_*\) of the product window. A term independent of
the interior coordinates then cancels. Each other full
carrier meets a current cube of the box, with its
prescribed fixed enlargement for pivot reads. After
summing its old origins, root that carrier at such a
cube. The strict rate reserve in
\eqref{f2:old-real-potential-bound} pays for the rooted
count. This gives
\[
 \sup_{y\in\mathcal W_I}|P_{e,\mathrm{reg}}(y)
                       -P_{e,\mathrm{reg}}(y_*)|
 \le C n_j\vartheta_k
              (S_k^E+S_k^R+S_k^c),
\]
with the ordinary defect size added if present. This is
a current-scale count. The old base points of a regular
activity have already been summed in its window, after
the counterterm of that base point has been subtracted.

%
%
For surviving boundary and kept factors, apply
the ordinary step estimates on the admitted
support just described. The actual support
point, including the cut strip point, lies
in the half-radius domain of the substitution.
The pinned levelling map uses the actual
output datum and its transformed current.
It preserves the stored arguments. Thus the
identities for the two evaluated exponents
use the same spectators, and their estimates
apply at the original fixed exterior.

For the filed boundary system, apply the
one-term-per-slot bound of Theorem 10(e)
by origin, with its stored-field and dial
domains. Its amplitude is
\(C\mathsf B\vartheta_k\) at the retained
hull rate.
Its carrier crosses the small-field boundary,
with hull length at least
\(c_2L\check R_j-2\).
For any fixed \(a\) the bound is
\(C_a\mathsf B\vartheta_k(\log\theta_k)^{-a}\)
times the hull decay at a smaller rate.
Spend a fixed part of the available rate
on the crossing length and retain a rate
above the rooted entropy.
For a kept factor use its inherited
carrier oscillation bound. The factor
form gives
\(C_a\mathsf B_\Phi\vartheta_k e^{-L\check R_j}\)
times the same decay. In the difference
at \(y\) and \(y_*\), all terms depending
only on the fixed exterior cancel.

The root cubes for the remaining terms
lie outside the alive exempt hulls by the
clean padding. The rooted sum
\eqref{ccf:hull-root-sum} thus gives, on
labels with these inherited bounds,
\[
 Cn_j\vartheta_k\bigl(
       S_k^E+S_k^R+S_k^c
      +C_a\mathsf B(\log\theta_k)^{-a}
      +C_a\mathsf B_\Phi e^{-L\check R_j}\bigr)
\]
for the old oscillation, including the
ordinary defect size when present.
This application is to the real support.
The induction of the inherited bounds
on all surviving histories is a separate
input to the exterior mixture.

\paragraph{Terms attached to healed regions.}
%
%
%
%
In the native positive history representation, a boundary
term produced at healing is filed at its healing level.
Its localization box can reach the read after the region
has disappeared, or when the region lies outside the
padded read. To use the estimate on attached terms, test
the terms themselves. At read level \(j\), an attached
term \(X\) filed at \(i'<j\) has weight
\(L^{-4\nu}x_{i'}^{2q+2}M_X d_{i'}(X)^2\), where
\(\nu=j-i'\), \(M_X\) bounds the term and \(d_{i'}(X)\)
is the size of its localization box in level-\(i'\) units.
Sum over boxes meeting the read. The native localization
estimate bounds the expectation per fine anchor by
\(C L^{-4\nu}x_{i'}^{2q+2}q_{i'}^{\mathrm{pk}}\), where
\(q_{i'}^{\mathrm{pk}}\) bounds the mass of the creation
or healing event that supplies a term at that position.
This uses the sum of \(M_X d_{i'}(X)^2\) at a fixed
anchor and the spatial decay for terms attached to
regions away from the read. The fine-anchor count
\(C n_j L^{4\nu}\) cancels the factor \(L^{-4\nu}\).
The per-bond version uses the same bound with the
localization overlap constant.

With \(q_{i'}^{\mathrm{pk}}\le C e^{-c_{\mathrm{cr}}p_0(g_{i'})}\),
the resulting level sum is bounded by
\(C n_j e^{-c_{\mathrm{cr}}p_0/2}\). For \(i'<k_0\)
use both running bounds on \(x_{i'}\). Half the
exponent controls the sum of the polynomial weights
in \(x+B_\sharp(k_0-i')\), since \(\mathfrak p_0>1\).
The levels of the fluctuation segment have comparable
coupling and logarithmic cardinality. Their polynomial
prefactor is absorbed by the reserve in the exponent
since \(p_0(g_{i'})/p_0\to1\) uniformly in that segment.
Markov's inequality
at \(e^{-c_{\mathrm{cr}}p_0/4}\) gives a failure mass
at most \(C n_j e^{-c_{\mathrm{cr}}p_0/4}\).
This estimates the budget with all filing levels.

For the weighted derivative estimate on the older list
\(i'<k_0-\mathfrak l\), use its axial-gauge representative
and the vanishing first jet at the flat field. The
Schwarz bound uses the fine fluctuation
\(C L^{-2\nu}g_jb_j\) and
\((g_j/R_{i'}^{\mathrm{an}})^2\le C(1+B_\sharp\nu)\).
Thus the term has second-order size at most
\(C M_X\kappa_{\mathrm{ax}}(X)^2 L^{-4\nu}
  (1+B_\sharp\nu)b_j^2\), where \(\kappa_{\mathrm{ax}}(X)\)
bounds the axial representative on the box. Apply
Cauchy's estimate on the admitted tube for the Hessian.
The weighted row sum contributes
\((d_{i'}(X)L^{-\nu}+C)^4e^{c_P(d_{i'}(X)L^{-\nu}+C)}\).
For the localized boundary terms the exponential uses
part of their decay. For the frozen-coupling term the
older-box bound \(d_{i'}(X)\le100M\check R_{i'}+C\)
and the choice of the lag depth bound
\(d_{i'}(X)L^{-\nu}\) uniformly. The remaining factors
are polynomial in the old inverse coupling, so the
same first-moment argument applies with an enlarged
fixed power. On this per-bond event the native gradient
bound is \(C\widehat b_j(b_j^2+p_0(g_j))e^{-c p_0}\),
with the pivot-map norm included in \(C\).
Integrating on the product window gives an oscillation
at most \(C n_j b_j\widehat b_j(b_j^2+p_0(g_j))e^{-c p_0}\),
which can be added to the old oscillation when this
decomposition is used. We must use the localization
and derivative estimates for the actual terms of the
conditional density. The all-filings mass estimate
alone does not extend the older-list derivative bound
to terms filed later.

%
For the native ordinary remainder, retain the
localization of each term in the gradient estimate.
The volume factor of a level-\(i'\) remainder
cancels the number of its fine placements. The
remaining old-level sum in the gradient is bounded
by \(C g_j x_j^{(3-\kappa_0)/2}
\operatorname{polylog}(x_j)\), with \(x_j=g_j^{-2}\).
At fixed background apply Cauchy's estimate in
the second coordinate, including a coordinate of the
exterior, on the joint window tube. Let \(m_j\) be
the available complex coordinate margin there.
Only terms whose read contains both coordinates
can contribute. Their localization reserve pays
the weighted rooted sum, giving
\[
 \sup_b\sum_{b'} e^{c\,d(b,b')}
 |\partial_b\partial_{b'}P^{\mathbf R}_j|
 \le \frac{C}{m_j}\,
       g_jx_j^{(3-\kappa_0)/2}\operatorname{polylog}(x_j).
\]
For the bound with \(b_j^{-1}\) one needs a
complex margin proportional to \(b_j\).
An absolute or inverse-logarithmic margin still
gives a small Hessian for the chosen
\(\kappa_0>4\). This estimate is for the
ordinary remainder with its native localization
and chart. The terms from marked or retained
spectators need their own mixed bounds.

%
The per-bond event for this comparison is imposed on
the incoming list of formed terms. At a compatible
exterior keep its filing levels, localization boxes
and recorded coupling lags, along with the spectators
and the history data that determine this list. The
geometric clean guard fixes the free-bond mask and
each dependence shell. The bound \(M_X\) is a
supremum on the admitted domain of the formed term,
including its recorded lag factor where present.
It is fixed as the real interior coordinate varies.
The axial factor and the distance and count weights
also depend only on this exterior record and the
fixed geometry. The pass indicator for the budget
is therefore an exterior factor on the whole product
window. At a merge sum the positive incoming terms
with their own list and indicator. A test on a list
filed after this step must be imposed at its own
incoming comparison or justified by the later-record
kernel on the actual retained state.

Let \(P_{e,\mathrm{old}}\) denote the sum of the
admitted contributions just bounded. Suppose the
resulting bound for
\(\sup_{\mathcal W_I}|P_{e,\mathrm{old}}(y)-
P_{e,\mathrm{old}}(y_*)|\) is \(\epsilon_j^{\mathrm{old}}\),
uniformly over the exterior data. For each \(e\),
keep the same domain, quadratic extraction, intrinsic
terms, stored variables and readout, and remove
\(P_{e,\mathrm{old}}\) only in an auxiliary normalized
interior law \(E_e^{\mathrm{int}}\).
Then
\[
 e^{-2\epsilon_j^{\mathrm{old}}}
 \le \frac{dE_e}{dE_e^{\mathrm{int}}}
 \le e^{2\epsilon_j^{\mathrm{old}}}.
\]
Indeed the ratio is the remaining exponential divided
by its mean under \(E_e^{\mathrm{int}}\).
The estimate costs
\((e^{2\epsilon_j^{\mathrm{old}}}-1)\sup|G|\)
for a bounded insertion, uniformly in \(\omega\) and
\(e\). Integrate this bound against the original
exterior mixing measure. The intrinsic
box comparison and the flat read-set approximation
are performed next, for each exterior, and have the same
flat target.

The loss \(C n_j(\log x)^{\mathfrak p_0-q}\) uses
one power of the window-to-radius ratio. It suffices
for the generating function route with any fixed
number of insertions by choosing \(q\) after the
polynomial logarithmic losses in the source radii.
A Hessian at fixed exterior cannot supply a derivative
with respect to a different retained datum or a stored
variable. Such a derivative would require one common
analytic domain for the two complete evaluations.
The normalized removal uses the real admission just
proved on the conditional support, with its retained
rows and margins. Each surviving boundary or stored
factor must also retain its declared carrier and
oscillation majorant at that evaluation.

Adding $a^{\mathsf T}y$ to $P_e$ leaves its Hessian unchanged
and changes its force by $a$.  It translates the Gaussian mean
even on all of $\mathbb R^n$.  Centering a quadratic observable
removes its scalar mean, but not the linear term caused by
that field translation.  We retain this term below.

\subsubsection*{Integration by parts with its boundary}

Suppress $e$, let $Q$ be real symmetric and invertible, and
set $C=Q^{-1}$.  Define
\begin{equation}
 {\cal B}_j(F)=Z^{-1}\int_{\partial D}
                       n_j(y)F(y)e^{-V(y)}\,dS(y).
 \label{eq:clean-conditional-law-CL6}
\end{equation}
The divergence theorem, with $\partial_jV=(Qy)_j-r_j$, gives
\begin{equation}
 {\mathbb E}[(QY)_jF]
       ={\mathbb E}[\partial_jF+r_jF]-{\cal B}_j(F).
 \label{eq:clean-conditional-law-CL7}
\end{equation}
On a rectangle this is the difference of its two face integrals.
The boundary has not been discarded.

The auxiliary linear-source function $M(z)={\mathbb E}e^{z^{\mathsf T}Y}$
is entire on a bounded window.  Since $M(0)=1$, its logarithm
$W(0)=0$ is defined near zero.  There put
\begin{equation}
 {\mathbb E}_z F={{\mathbb E}(Fe^{z^{\mathsf T}Y})\over M(z)},
       \qquad
 \beta_j(z)={{\cal B}_j(e^{z^{\mathsf T}Y})\over M(z)} .
 \label{eq:clean-conditional-law-CL8}
\end{equation}
Apply \eqref{eq:clean-conditional-law-CL7} to the exponential and multiply by $C$.  This proves
the exact connected identity
\begin{equation}
 \partial_iW(z)=\sum_jC_{ij}
                       \{z_j+{\mathbb E}_z r_j-\beta_j(z)\}.
 \label{eq:clean-conditional-law-CL9}
\end{equation}
The source here is an auxiliary linear field source.  It is not
an identification with a physical $\operatorname{tr}(F^2)$ numerator.

Let $m={\mathbb E}Y$, let $\Sigma$ be the covariance and let
$\kappa$ denote joint cumulants.  Differentiation gives
\refstepcounter{equation}
\begin{align*}
 m_i&=\sum_jC_{ij}\{{\mathbb E}r_j-\beta_j(0)\},\\
 \Sigma_{ia}-C_{ia}
 &=\sum_jC_{ij}\{\kappa(r_j,Y_a)-\partial_a\beta_j(0)\},\\
 \kappa(Y_i,Y_a,Y_b)
 &=\sum_jC_{ij}
       \{\kappa(r_j,Y_a,Y_b)-\partial_a\partial_b\beta_j(0)\},\\
 \kappa(Y_i,Y_a,Y_b,Y_c)
 &=\sum_jC_{ij}
       \{\kappa(r_j,Y_a,Y_b,Y_c)
                        -\partial_a\partial_b\partial_c\beta_j(0)\}.
 \tag{\theequation}\label{eq:clean-conditional-law-CL10}
\end{align*}
Each derivative differentiates the numerator and $1/M$.  The
first produces a covariance; the next remove the disconnected
partitions and give the mixed cumulants.  Equivalently, compare
the convergent local power series in \eqref{eq:clean-conditional-law-CL9}.

Put $Z_a=Y_a-m_a$.  The boundary derivatives are explicitly
\refstepcounter{equation}
\begin{align*}
 \beta_{j,a}(0)&={\cal B}_j(Z_a),\\
 \beta_{j,ab}(0)&={\cal B}_j(Z_aZ_b)-\Sigma_{ab}{\cal B}_j(1),\\
 \beta_{j,abc}(0)&={\cal B}_j(Z_aZ_bZ_c)
   -\Sigma_{ab}{\cal B}_j(Z_c)-\Sigma_{ac}{\cal B}_j(Z_b)\\
 &\hspace{10mm}-\Sigma_{bc}{\cal B}_j(Z_a)
                   -\kappa(Y_a,Y_b,Y_c){\cal B}_j(1).
 \tag{\theequation}\label{eq:clean-conditional-law-CL11}
\end{align*}
These follow by differentiating the quotient.  In particular
an unnormalized face-moment estimate is not automatically the
same estimate for a connected boundary derivative.

\subsubsection*{Replacing traces by weighted collar moments}

On the product window \(D=\prod_b K_b\), let \(b(j)\)
be the bond carrying the coordinate \(j\). Choose
\(\psi_j\) depending only on \(y_{b(j)}\), vanishing
near \(\partial K_{b(j)}\), equal to one off its collar
of width \(\delta_j\), with
\(|\partial_j\psi_j|\le c/\delta_j\).
Integration by parts in the \(j\) direction
has no boundary contribution for
\(\psi_j e^{z^{\mathsf T}y}\): on the
\(b(j)\) face the cutoff vanishes, and the
other faces have zero \(j\)-normal component.
Equivalently, use Fubini and integrate only
in the factor \(K_{b(j)}\). Comparing
it with \eqref{eq:clean-conditional-law-CL7} proves
\begin{equation}
 {\cal B}_j(e^{z^{\mathsf T}Y})
 =-{\mathbb E}\left[e^{z^{\mathsf T}Y}
   \{\partial_j\psi_j+
       (1-\psi_j)((QY)_j-r_j(Y)-z_j)\}\right].
 \label{eq:clean-conditional-law-CL12}
\end{equation}
Indeed the compactly supported formula equates
${\mathbb E}[\psi_j(QY-r)_je^{zY}]$ with
${\mathbb E}[(\partial_j\psi_j+z_j\psi_j)e^{zY}]$.
Adding the part with $1-\psi_j$ gives the sign and terms above.

Thus moments of the form
\begin{equation}
 {\mathbb E}\!\left[
  1_{\rm collar}(1+|Y|_{\rm loc})^q
   \{\delta_j^{-1}+|(QY)_j|+|r_j(Y)|+1\}\right]
 \label{eq:clean-conditional-law-CL13}
\end{equation}
control the corresponding numerator derivatives, with
the collar of the single bond \(b(j)\).
The local norm means the finite tested coordinate list
with its specified weights.
The derivatives of $1/M$ also use its ordinary moments.

Up to order three there are only finitely many product and
quotient terms.  If their ordinary moments are uniformly bounded
and the weighted collar moments are $O(\epsilon)$, the boundary
derivatives in \eqref{eq:clean-conditional-law-CL10} are $O(\epsilon)$ with the resulting
finite moment factors.  This follows directly by differentiating
\eqref{eq:clean-conditional-law-CL12} and \eqref{eq:clean-conditional-law-CL8}.  An unweighted collar probability does not
prove \eqref{eq:clean-conditional-law-CL13}; polynomial losses in the force and cutoff remain.

For a bond in the buffered interior the normalized
box comparison supplies one way to bound this moment.
Suppose the marginal at that bond has density ratio
at most \(e^\Xi\) relative to the reference on the
same product window. At the reference, the Gaussian is
centered, the coordinate variances are uniformly bounded,
the full window has Gaussian probability at least \(1/2\),
and the tilt has uniformly bounded oscillation. If the
bond window contains the Euclidean ball of radius
\(c_0 b\), a collar of Euclidean width at most \(c_0 b/2\) is
contained in \(\{|y_{b(j)}|\ge c_0 b/2\}\).
The Gaussian one-bond tail, divided by the full
window probability and multiplied by the exponential
of the tilt oscillation, is at most \(Ce^{-c b^2}\).
The density ratio transfers this bound at the same
bond with factor \(e^\Xi\).

On the product window, the uniform precision row
bound gives \(|(QY)_j|\le Cb\). The retained
linear force and the perturbation gradient use
their admitted row bounds. When these, the
tested coordinate sizes and the inverse collar width
are bounded by fixed powers of \(1+b\),
\eqref{eq:clean-conditional-law-CL13} is bounded
by a polynomial in \(1+b\) times \(e^{\Xi-cb^2}\).
This uses the density comparison at the differentiated
bond, with its full read margin. The bound is uniform
over exterior data for which these constants are
uniform, and passes to their positive mixture.

In \eqref{eq:clean-conditional-law-CL10}, the boundary
derivatives are contracted against a row of the Gaussian
inverse. Take a tested coordinate at depth at least
\(2D_0\) from the edge after accounting for its read.
For bonds at distance at most \(D_0\) from this coordinate,
use the collar bound with the common buffer. On the
remaining bonds use the deterministic polynomial bound
in \eqref{eq:clean-conditional-law-CL13} from the full
window and the uniform force and gradient bounds.
The weighted row estimate for the inverse makes the
contribution of these remaining bonds at most
\(C\operatorname{poly}(1+b)e^{-cD_0}\).
The quotient derivatives use only moments of the fixed
tested list, so they retain the same split.

When \(e\) varies, the bound must be uniform in \(e\) or integrated
in the order specified by the actual exterior law.  A
fixed-exterior bound is not a uniformly integrable envelope.
This matters for a clean restriction formed by healed components.

\subsubsection*{The covariance of quadratic insertions}

Let
\begin{equation}
 O_i=Y^{\mathsf T}A_iY+b_i^{\mathsf T}Y+c_i,\qquad
 A_i^{\mathsf T}=A_i,\qquad d_i=b_i+2A_i m .
 \label{eq:clean-conditional-law-CL14}
\end{equation}
Only the symmetric part contributes to a scalar quadratic form.
Expansion at $Y=m+Z$ gives
$O_i-{\mathbb E}O_i=Z^{\mathsf T}A_iZ-
\operatorname{tr}(A_i\Sigma)+d_i^{\mathsf T}Z$.
Multiplication yields
\refstepcounter{equation}
\begin{align*}
 \operatorname{Cov}(O_1,O_2)
 ={}&2\operatorname{tr}(A_1\Sigma A_2\Sigma)
                            +d_1^{\mathsf T}\Sigma d_2\\
 &+\sum_{a,b,c}(A_1)_{ab}(d_2)_c\kappa(Z_a,Z_b,Z_c)\\
 &+\sum_{a,b,c}(A_2)_{ab}(d_1)_c\kappa(Z_a,Z_b,Z_c)\\
 &+\sum_{a,b,c,d}(A_1)_{ab}(A_2)_{cd}
                                  \kappa(Z_a,Z_b,Z_c,Z_d).
 \tag{\theequation}\label{eq:clean-conditional-law-CL15}
\end{align*}
The fourth centred moment is its fourth cumulant plus three
pairings.  Subtracting the quadratic means removes
$\Sigma_{ab}\Sigma_{cd}$.  The other two pairings contract to
the trace above.  The mixed terms give the third cumulants,
and the two linear terms give $d_1^{\mathsf T}\Sigma d_2$.
This accounts for every term.

The centred Gaussian reference with covariance $C_0$ and purely
quadratic insertions is $2\operatorname{tr}(A_1C_0A_2C_0)$.
Its identification needs the mean term, covariance discrepancy
and third and fourth cumulants in the actual tested contractions.

For one sufficient dimension-explicit estimate put
\begin{equation}
 a_i=\sum_{a,b}|(A_i)_{ab}|,\quad
 v_i=\sum_a|(d_i)_a|,\qquad
 \|\Sigma\|_{\max},\|C_0\|_{\max}\le K .
 \label{eq:clean-conditional-law-CL16}
\end{equation}
If $\|\Sigma-C_0\|_{\max}\le\epsilon_2$ and the tested third
and fourth cumulants are bounded by $\epsilon_3,\epsilon_4$,
then
\refstepcounter{equation}
\begin{align*}
 &|\operatorname{Cov}(O_1,O_2)
                     -2\operatorname{tr}(A_1C_0A_2C_0)|\\
 &\quad\le4K\epsilon_2a_1a_2+Kv_1v_2
          +\epsilon_3(a_1v_2+a_2v_1)+\epsilon_4a_1a_2 .
 \tag{\theequation}\label{eq:clean-conditional-law-CL17}
\end{align*}
Expand one covariance factor at a time in the trace.  Each
product differs by at most $2K\epsilon_2$, and the prefactor
two gives the first term.  Absolute summation gives the others.
Weighted localized variants use the same actual contractions.

If $a_i\le Cr^4\|f_i\|_\infty$ and
$r\le C(\log x)^{p_r}$, the first and fourth errors cost
$(\log x)^{8p_r}$.  They require
$\epsilon_2,\epsilon_4=o((\log x)^{-8p_r})$ to be
$o(\|f_1\|_\infty\|f_2\|_\infty)$.  A claim that errors
merely tend to zero does not suffice.  The other rates are
read from the mean and third-moment terms.

Each physical quadratic insertion ordinarily carries $g^2$,
so the reference covariance carries $g^4$.  Its errors and
the physical remainders must be small relative to that signal.
A fixed source-disc estimate cannot replace a missing potential
normalization by an inverse power of $g$.

\subsubsection*{Exterior averaging adds the conditional means}

Suppose the resummed representation has a positive probability
realization with exterior variable $E$, clean event $A$ and
correct observable pullbacks.  Conditional expectation gives
\refstepcounter{equation}
\begin{align*}
 \operatorname{Cov}(O_1,O_2\mid A)
 ={}&{\mathbb E}\{\operatorname{Cov}(O_1,O_2\mid E,A)\mid A\}\\
 &+\operatorname{Cov}({\mathbb E}(O_1\mid E,A),
                      {\mathbb E}(O_2\mid E,A)\mid A).
 \tag{\theequation}\label{eq:clean-conditional-law-CL18}
\end{align*}
Add and subtract the product of conditional means to prove it.
It applies to a positive conditional law, not separately to
signed expansion terms.

Choose one deterministic number $a_i^0$ per observable and prove
\begin{equation}
 |{\mathbb E}(O_i\mid e,A)-a_i^0|\le\delta_i
                  \quad\hbox{for every admitted clean }e .
 \label{eq:clean-conditional-law-CL19}
\end{equation}
Then the last covariance in \eqref{eq:clean-conditional-law-CL18} has absolute value at most
$2\delta_1\delta_2$.  Choosing a different reference for every
exterior would not prove this.  The backgrounds, insertion
coefficients and conditional moments need one common reference.

Put $q=\mathbb P(A^c)$, assume $0<q<1$ and $|O_i|\le B_i$ on
the whole positive probability space.  Decomposing by $A$ gives
\refstepcounter{equation}
\begin{align*}
 \operatorname{Cov}(O_1,O_2)
 ={}&(1-q)\operatorname{Cov}(O_1,O_2\mid A)
       +q\operatorname{Cov}(O_1,O_2\mid A^c)\\
 &+q(1-q)(m_1^A-m_1^{A^c})(m_2^A-m_2^{A^c}).
 \tag{\theequation}\label{eq:clean-conditional-law-CL20}
\end{align*}
If every clean conditional covariance differs from a fixed
$\Gamma$ by at most $\epsilon$, this implies
\begin{equation}
 |\operatorname{Cov}(O_1,O_2)-\Gamma|
 \le\epsilon+2\delta_1\delta_2+q(6B_1B_2+|\Gamma|).
 \label{eq:clean-conditional-law-CL21}
\end{equation}
The bad covariance is at most $2B_1B_2$ and the product of
mean differences is at most $4B_1B_2$.  These prove the
constants.  At $q=0$ use the clean formula directly; at
$q=1$ there is no clean estimate to use.

For unbounded insertions retain the marked bad-event moments.
For instance Cauchy--Schwarz bounds
$|{\mathbb E}[1_{A^c}O_1O_2]|$ by
$({\mathbb E}[1_{A^c}O_1^2])^{1/2}
({\mathbb E}[1_{A^c}O_2^2])^{1/2}$.
Centering also requires the corresponding first moments.
A bad-event probability alone does not bound them.

An unnormalized clean sum has to be divided by its actual
mass.  Source derivatives of that scalar contribute unless
the exact numerator and denominator have already been
realized as this conditional probability.  A normalizer from
another volume or resolved class cannot be substituted.

\subsubsection*{Uniform shapes and the remaining model estimates}

The required positive shapes have $f_i\ge0$, support diameters
$O(d)$, separation between $d$ and $6d$, and
\begin{equation}
 \int f_i\ge\vartheta d^4\|f_i\|_\infty,\qquad
\|\nabla f_i\|_\infty\le C_{\rm shape}d^{-1}\|f_i\|_\infty .
 \label{eq:clean-conditional-law-CL22}
\end{equation}
The integral against $|x-y|^{-8}$ has a lower bound
$c_\vartheta\|f_1\|_\infty\|f_2\|_\infty$.  The maximum
distance is a fixed multiple of $d$ and the two masses
supply $d^8$.  The free comparison theorem gives the
lattice version at sufficiently fine compatible mesh, with this
uniform derivative bound. In fact, partition into mesh cells of side
$a\le d$. The oscillation on each cell is at most
$Ca\|\nabla f_i\|_\infty$. Only cells in the support enlarged by
$Ca$ contribute, whose total volume is $O(d^4)$ by the diameter
condition. Therefore the Riemann mass error is at most
$CC_{\rm shape}(a/d)d^4\|f_i\|_\infty$, uniformly in the
position of the support. Taking $a/d\le c\vartheta/C_{\rm shape}$
gives sampled mass at least
$(\vartheta/2)d^4\|f_i\|_\infty$. The actual source class permits
$C_{\rm shape}=\vartheta^{-1}$, in which case a sufficient floor is
$a/d\le c\vartheta^2$. Alternatively the sampled mass itself can
be imposed. Smoothness and continuum mass without a uniform derivative
bound do not supply this floor.

All errors in \eqref{eq:clean-conditional-law-CL21}, with physical remainders, must be
$o(g^4\|f_1\|_\infty\|f_2\|_\infty)$ uniformly on the
whole shape class, positions and required reference volumes.
The fixed $\vartheta=1/12000$ allows dependence on that
value followed by one uniform cutoff floor.

A reference estimate uniform along arbitrarily large tori
at each fixed datum and test pair passes to a reference
infinite-volume state and uses common-normalizer transfer.
An estimate only on one prescribed finite torus needs
the separate finite-reference and vacuum comparison.
These are alternatives; neither follows from a local kernel.

The first remaining task is the literal identity in \eqref{eq:clean-conditional-law-CL2}, with
the actual exterior sigma-field, domain, numerator and normalizer.
Next bound the complete Hessian and the individual forces.
Then bound the tested connected force and boundary quantities
in \eqref{eq:clean-conditional-law-CL10}, using either the trace or collar identity.
The common-reference mean comparison in \eqref{eq:clean-conditional-law-CL19} and marked bad
moments are separate.  Their rates must dominate the window
and coefficient losses in \eqref{eq:clean-conditional-law-CL17}.

These estimates would supply the interacting covariance on
the run route.  The proved free, remainder and state-transfer
implications then have their needed inputs.  The alternative
averaging-fibre route additionally requires its measure and
carried-weight identification; it is not an extra conjunct
once the run route closes.


% END MODULE chapters/clean_conditional_law.tex

% BEGIN LAYERED EXPANSION formation/clean_collar_filing.tex
\subsubsection{A clean collar and the actual local filing}
\label{layered:clean-collar-filing}

A local term made inside a sufficiently clean collar uses the
whole-lattice recipe. The collar must contain the actual trace
window, which includes a layer of cells of side
\(M\check R_j\xi_j\). A layer of \(M\xi_j\)-cubes is
smaller and does not suffice for that assertion. We first prove
the geometric and finite-filing statements, then sum the boundary
tails. The whole-lattice comparison keeps the same input list;
in particular it does not erase regular marked members produced
by healed histories.

\subsubsection{The two grids and the trace window}

Use physical spacings \(\xi_j=L^j\xi_0\). Let
\(\pi_j\) be the grid of cubes of side \(M\xi_j\), and
let \(\mathcal P_j\) be the cell grid of side
\(M\check R_j\xi_j\). The grids and all labels are the
ones fixed by the run. Assume the actual cube ratio
\begin{equation}
 \check R_j\le L\check R_{j+1},\qquad j<n.
 \label{ccf:cube-ratio}
\end{equation}
Thus the physical cell size is nondecreasing:
\(\check R_j\xi_j\le\check R_n\xi_n\).
Distances in current cube units mean physical distance divided
by \(M\xi_n\), using the periodic graph metric when needed.

For a region \(X\), its trace window \(X^{+1}\) rounds it
out, at each level \(j\), to all cells of \(\mathcal P_j\)
that meet it, and adds one cell layer. Include the fixed finite
additional layers used by the actual step in this definition.
Let \(c_{\rm tr}\) be a bound for their number. Rounding and
adding these layers gives the explicit inclusion
\begin{equation}
 X^{+1}\subset
 \{x:\operatorname{dist}_\infty(x,X)
             \le C_{\rm tr}M\check R_n\xi_n\},
 \qquad C_{\rm tr}=C(1+c_{\rm tr}).
 \label{ccf:trace-padding}
\end{equation}
Indeed rounding to one cell moves each coordinate by at most
one cell side; every additional layer costs at most one more.
Apply \eqref{ccf:cube-ratio} at each level. No sum over the
number of levels occurs, because the window is the union of
these levelwise neighborhoods, not their successive enlargement.

Let \(\mathcal S_n\) be a standing label and put
\(Z_n=\Lambda_n^c\). For \(W\), a set of current cubes,
write \(W_R\) for its \(R\)-neighborhood in current cube
units. The clean-collar condition used below is
\begin{equation}
 W_R\subset\Lambda_n,\qquad
 s:=R-C_{\rm tr}\check R_n-C_0>0.
 \label{ccf:clean-collar}
\end{equation}
Here \(C_0\) covers the fixed cube collars of the local
potential and integration formulas. They are specified as part
of the recipe, rather than changed after a bound is proved.

\subsubsection{What remains in a standing history}

The standing part deletes the entire cluster of each healed
component. All data which are still stored belong to an island,
a drop-out or an unhealed component. The following conclusion
uses this actual deletion rule.

\begin{lemma}[The standing trace in the clean region]
\label{lem:ccf-empty-standing-trace}
Every retained date, mark and stored variable of a standing
history of level \(n\) lies in \(Z_n\), with at least one
endpoint there for a stored bond. Outside \(Z_n\), the
action and determining sequences have no remade layers and
coincide with the all-small-field sequences. Consequently the
standing trace and its reading on \(X^{+1}\) are empty when
\(X^{+1}\subset\Lambda_n\).
\end{lemma}

\begin{proof}
Induct on the moves of a step. A new large-field outcome is
recorded only on the component where its test fails; it is
contained in that step's large-field region. At a subsequent
step, an old component and all its standing contents are either
carried, merged into another live component, or treated by the
large-field operation. Carrying and merging retain the contents
inside the new live component. An island or a drop-out retains
its contents inside its own component. Healing deletes that
component's entire cluster, including its older dates, fields,
frames and marks; none of those deleted variables remains as a
free integration variable in the standing part.

The only modifications of an earlier label sequence take place
inside a component of a later operation. If the component is
healed, its old regions are replaced by the healed output. If
it remains standing, it lies in \(Z_n\). The remade levels in
\eqref{eq:rg-objects-R21} are modified inside precisely such a
standing drop-out. Hence no remade layer survives outside
\(Z_n\). The nested unmodified label regions give the
all-small-field sequence there. This proves the induction.
An empty trace follows by restriction to \(X^{+1}\).
\end{proof}

An alive island \(I\) has a hull \([I]_j\), the union of the
closed level-\(j\) cubes meeting \(I\). At birth this gives the
one-layer enlargement. The nested partitions give
\begin{equation}
 ([I]_j)^{\uparrow n}=[I]_n,
 \qquad [I]_n\subset Z_n^{\sim c_I},
 \label{ccf:hull-coarsening}
\end{equation}
where \(c_I\) is the fixed hull collar. The first identity
uses the same standing island on both sides; an island absorbed
before level \(n\) is no longer in this family. The second
follows from seating of an alive island in the final live region
and from the definition of its hull. Thus an alive exempt hull
cannot enter the interior of a collar which excludes
\(Z_n^{\sim c_I}\). The possible failure of a finer place
condition inside a remade drop-out occurs in \(Z_n\) and does
not change this conclusion.

%
The hulls have a bounded corner multiplicity. Write \(u_j\) for
the physical side of a level-\(j\) localization cube and \(b_i\)
for the physical side of a level-\(i\) history cell. An island born
at \(i\) has sides between \(21b_i\) and \(100b_i\). Its level-\(j\)
hull has sides at most \(100b_i+2u_j\), hence at most
\(100\check R_j+2\) cubes per direction since \(b_i\le b_j\).
Use the torus margin \(204b_j<\ell\) for the following count.
If a hull contains a corner \(v\), its island meets the box of
radius \(u_j\) about \(v\), and both lift into the box of radius
\(102b_j\) about the lifted corner. When \(21b_i\ge u_j\), the
island contains a box of side \(u_j\) within the box of radius
\(2u_j\) about that corner. Distinct such islands are disjoint,
so volume comparison gives at most \(4^4\) of these boxes.
When \(21b_i<u_j\), each island side is shorter than \(5u_j\)
and each hull side contains at most six cubes. An aligned interval
of length \(n u_j\) containing the corner has at most \(n+1\)
positions. Summing over \(1\le n\le6\) gives at most \(27^4\)
distinct small hulls containing the corner. The exempt family
consists of geometric hulls, so coincident hulls count once here.
Thus \(m_H\le4^4+27^4\), also for the surviving subfamily or its
hulls on the next grid with the same torus margin.

\subsubsection{Equality of the filed recipes}

We use the actual local representation of the step: walk terms
are filed at their full carriers; pivot charts, means and
Jacobians carry their defining double blocks; Mayer selections
retain the whole selected potential domain; a cutoff selection
retains its bond collar; and a connected logarithmic term is
filed at the union of its polymer domains. The normalized
Gaussian functions are restricted by their prescribed hereditary
rules. These are the literal recipes of
\cref{ch:local-formation-program}, not merely numerical bounds
on a gas. Every dependency must be included in the carrier
before this convention is applied.

\begin{proposition}[The whole-lattice recipe on an empty trace]
\label{prop:ccf-whole-lattice-identification}
Fix the reference numbers, the physical action weights and one
coherent input list. Let \(E_h(X,\cdot)\) be a filed term made
from a standing history \(h\), and let
\(E_\emptyset(X,\cdot)\) be the term made by the whole-lattice
step with that same input. If
\(X^{+1}\subset\Lambda_n(\mathcal S_n)\), then
\begin{equation}
                 E_h(X,P)=E_\emptyset(X,P)
 \label{ccf:local-recipe-identity}
\end{equation}
on every common admitted local pair domain. The assertion
applies to the local Gaussian normalizer, the mismatch leg,
the interaction terms, and their declared pointing.
It asserts equality of functions at the same pair \(P\),
rather than equality of two generally different global
backgrounds.
\end{proposition}

\begin{proof}
By \cref{lem:ccf-empty-standing-trace}, the trace on the
entire set of cells used by a carrier in \(X\) is empty.
The local regions, tiers and integration-bond types therefore
coincide with the whole-lattice ones. In particular a carrier
contained in \(X\) has no external or crossing integration
bond. Its cutoff has the ordinary full threshold. A mean term
whose full carrier is contained in \(X\) cannot reach an
external bond, so that term is zero. The unlocalized global
Gaussian mean can still contain walks from the exterior;
their full carriers are not contained in \(X\).

Each local differential coefficient and each walk term depends
only on the pair, regions and bonds in its full carrier. At the
same \(P\), these data agree.

%
%
%
For the axial Gaussian, include the multiplier in this carrier
check. The second variation of the pivot row is local to its
double block, but it is paired at the pivot with the restricted
value of \(\mathsf H_1^{\mathsf T}\mathbf J\). Expand this
multiplier into its own walk terms before filing the pivot
contribution. Attach each such walk at the pivot and retain its
padded hull and current read. For the same indexed tree in
the two axial kernels, match every layer on that hull and the
admitted pair at its inner bonds. The active mask is the same on bonds
with both endpoints in the hull. The extensions by zero and
restrictions therefore give the same matrix entries in the two
settings, term by term. This gives equality for the native
axial covariance and precision terms after the free-coordinate
and mask identifications. If every term shorter than \(m_0\)
matches, the two remaining native tails give a kernel
difference bounded by
\[
 C2^{-m_0}e^{-\delta d_{\rm ax}(b,b')},
\]
where \(d_{\rm ax}\) is the distance between bond origins
in the axial lattice units. Indeed the absolute sum at length
\(m\) has the bound \(C2^{-m-1}\) times this spatial exponential.
Summing the geometric series and then the rows and columns
with a smaller spatial weight gives the weighted Schur bound
\(C2^{-m_0}\). An escaping multiplier branch is part of this
same tail.

 A local pivot chart and its
coarea factor agree by uniqueness of their admitted local
solution. The contour gauge and its normalizer use the same
contours and reference weight. Coherence of the input identifies
each old activity and its own counterterm used in \(X\).
Hence all individual potentials agree, including the fields
and current produced by their prescribed local substitutions.
This last assertion is on their common admitted domain; it
does not enlarge either domain.

Now apply the finite expansion. A Mayer term is a product of
the same selected potentials and the same cutoff indicators.
Its local Gaussian functional has the same coefficients and
the same real integration variables. Every coarse-parameter
derivative and integral in its polymer activity therefore
agrees. A logarithmic term on \(X\) uses only polymers
contained in \(X\), evaluated by hereditary restriction, so
every summand in its connected series agrees. Absolute
convergence, or equality at each finite truncation followed by
the same absolute majorant, proves
\eqref{ccf:local-recipe-identity}. Bond roots and midpoint
tents are the same in both recipes, so the conclusion also
passes to their actual pointing.
\end{proof}

The same proof shows that a nonzero difference between the two
recipes must have a carrier whose trace window meets the
nonempty standing trace. It cannot be filed on an interior
domain and then assigned a different, smaller carrier.

\subsubsection{Hull entropy rooted outside the exempt regions}

We supply the geometric sum needed for old boundary terms.
On the grid of level \(j\), let \(\mathsf H_j\) be the
hulls of the islands still alive at level \(n\). Their
multiplicity at a cube corner is at most \(m_H\). Define
\(d_j^\natural(X)\) by minimizing the length of a tree of
cube edges in \(X\), which meets every cube outside a
selected subfamily of complete hulls and meets every selected
hull. Edges inside a hull still cost their ordinary length;
only the requirement to visit all its cubes is removed.

\begin{lemma}[A root sum independent of island sizes]
\label{lem:ccf-hull-root-entropy}
If the cube \(q\) is contained in no exempt hull, then for
\(\kappa>c_E\),
\begin{equation}
 \sum_{X\ni q}e^{-\kappa d_j^\natural(X)}
 \le C_H\sum_{m\ge0}e^{-(\kappa-c_E)m},
 \qquad c_E=2\log8+(16+m_H)\log2.
 \label{ccf:hull-root-sum}
\end{equation}
The same count with a given root vertex $v$ holds for
domains whose chosen minimizing tree contains $v$, including
vertices in an exempt hull. It also bounds the sum restricted
to tree lengths $m\ge D$ by starting the series at $D$.
The constants are independent of the side lengths and volumes
of the hulls, the torus size and the number of levels.
\end{lemma}

\begin{proof}
Choose, for each \(X\), a minimizing tree and an exempt
subfamily by a fixed ordering. Its length in cube-edge units
is an integer \(m\). Since \(q\) is not exempt, the tree
meets one of its at most sixteen corners. Root it at the first
such corner. A tree with \(m\) edges has a depth-first tour
of length \(2m\); the lattice graph has degree at most eight.
There are at most \(16\cdot8^{2m}\) choices for this rooted
tree, with an overcount which only helps.

Every nonexempt cube of \(X\) meets a tree vertex, and at
most sixteen cubes meet one vertex. Every selected hull also
meets a tree vertex, and at most \(m_H\) hulls pass through
that vertex. Thus \(X\) is recovered by at most
\(16+m_H\) binary choices per vertex: choose the incident
cubes and the full hulls. There are at most \(m+1\)
vertices. This gives at most
\(16\cdot8^{2m}2^{(16+m_H)(m+1)}\) encoded domains.
The description is permitted to count disconnected or
inadmissible choices as well. Summation with
\(e^{-\kappa m}\) proves \eqref{ccf:hull-root-sum}.
For a specified vertex root the same encoding starts at that
vertex, and the factor sixteen for its choice is unnecessary.
Large hulls add one choice, rather than their volume.
\end{proof}

\begin{lemma}[A carrier crossing the clean collar]
\label{lem:ccf-crossing-length}
Suppose a level-\(j\) carrier meets \(W\), and its trace
window meets the standing trace. Under
\eqref{ccf:clean-collar} and \eqref{ccf:hull-coarsening},
after increasing \(C_0\) by the fixed hull collar,
\begin{equation}
                 d_j^\natural(X)\ge c_*L^{n-j}s-C_*.
 \label{ccf:crossing-length}
\end{equation}
The constants depend on the fixed grids and collars only.
\end{lemma}

\begin{proof}
The trace padding \eqref{ccf:trace-padding} shows that the
carrier reaches to current distance at least \(s\) from
\(W\). All current coarsenings of its exempt hulls lie
outside the clean interior by \eqref{ccf:hull-coarsening}.
An admissible tree must therefore meet a cube near \(W\),
which is not exempt, and either a second nonexempt cube at
distance at least \(s\), or a selected hull outside that
interior. In the latter case the tree still has to reach that
hull. Its connecting path pays ordinary geometric length,
including every portion inside a hull.

In physical units that path has length at least
\(cM\xi_ns-CM\xi_j\), allowing for the two touched cube
corners. Division by \(M\xi_j\) gives
\eqref{ccf:crossing-length}. The same argument on the
periodic graph uses its intrinsic distance. It does not
assume that the whole torus has one nonwrapping chart.
\end{proof}

\subsubsection{Summing old boundary origins in their actual units}

Let the observation region contain \(N_W\) current cubes.
It meets at most \(C L^{4(n-j)}N_W\) cubes at level \(j\).
Allow a fixed-origin multiplicity \(A_j\), and suppose the
actual local boundary majorants satisfy
\begin{equation}
 b_j(X)\le B_j e^{-\kappa d_j^\natural(X)},
 \qquad A_jB_j\le B_*(1+n-j)^p L^{q(n-j)}.
 \label{ccf:boundary-majorants}
\end{equation}
In this hypothesis \(d_j^\natural\) uses only hulls still
alive at the final level \(n\). This is an explicit admission
condition on the surviving boundary filing, not a consequence
of a creation-time bound using a larger exempt family. A term
whose carrier contains no island absorbed between \(j\) and
\(n\) has the same two lengths, so its old bound suffices.
At an absorbed hull use the refiling described below.

A simple example shows the necessity. Let \(H\) be a large
exempt box and \(q\) an adjacent cube. The old hull length
of \(H\cup q\) is zero. After \(H\) is removed from the
exempt family, a covering tree must visit its many cubes and
has length tending to infinity with their number. A bounded
nonzero function filed on \(H\cup q\) can satisfy the old
majorant and fail every proposed bound with the same amplitude
and the new length. Deleting the hull from the norm without
refiling the function is therefore invalid.

%
%
\paragraph{Boundary terms at an absorbed island.}
If a carrier \(X\) of origin level \(j\) contains the
level-\(j\) hull of an absorbed island, \(X\) meets the
treated region and enters the near list of that operation.
Consume the old member in the localized interaction.
Fill its slot from the reduced history obtained by erasing
the past of the treated components, and subtract that same
replacement in the localized interaction. The erased
history is a genuine history by the reduction rule. The
two changes preserve the sum of the exponents.

For an old boundary slot its donor has the same surviving
hulls in \(X\) as the completed label, so its hull length
equals the required output length. At the new level the
donor's length is at least the required length. Carry the
actual donor member with its birth core and its own level
and origin constant. Members made by the current operation
use their new bounds, and the pointed regular entries are
formed anew with the output base points. Each slot
holds one term, and at most four boundary origins
occur at one level and carrier. This transfers the decay
to the output length without using the old member's
discount for the absorbed hull.

%
%
At the new level \(t\) there is also a change
of slot for a domain disjoint from the treated
region. Put \(A=\Lambda_t(\mathcal S^+)\) and
write \(Z_1,Z_2\) for the unions of treated
components of the first and second classes.
The two output action regions are \(A\cup Z_1\)
and \(A\cup Z_1\cup Z_2\). The one-layer
window \(X^{+1}_t\) can reach one of the
treated components even when \(X\) does not.
The window criterion tests the new filing by
inclusion of \(X^{+1}_t\) in these regions.
For a former \(\mathbf T\)-boundary member
with \(X\) disjoint from the treated region,
if an outside cell of its window survives,
the member remains in the \(\mathbf T\) slot.
If all those cells are treated and at least
one is in \(Z_2\), the member goes into the
island remainder slot \(\mathsf H\). If all
are in \(Z_1\), it goes into the regular
remainder range.

The unpointed contribution supplied by
\(\mathbf T\) occupies exactly one of these
places. The change of filing gives no second
copy of the function. The rule that treats
a carrier as near whenever its window meets
the treated region supplies the same member
as the rule by slots, from the reduced
history. In particular a member in the
regular remainder range is supplied on that
range in the donor's pre-healing label.
On the narrower carry region we must
identify the passed function with that
member before applying its regular
remainder bound. Use the comparison
of the unfiled functions under erasure
of the deleted histories, with its
carrier and stage hypotheses. For the
pointed sum, split the old exterior base
points into the still exterior points and
the points newly in the one-layer erosion of
\(A\cup Z_1\cup Z_2\). The latter points enter
the regular sum once.

%
%
Apply the carriage bound to the standing
list of an output label. At a level and
carrier the bound is per origin. A carried
or donor member occupies that slot with
its own bound. A member formed at this
step occupies its new slot. The bounds
for the nonpointed origins are therefore
returned with amplitude
\[
 \max\{\mathsf B_{\rm in},\mathsf B^{\mathbf T},
                  \mathsf B^{\mathbf S},\mathsf B^{\mathbf R}\}
\]
in the output hull length. For the pointed
origin use the new sum of regular members
over its output base points. Its bound is
the pointed norm summed over those points.

For the ordinary formation the localized
Gaussian normalizer has amplitude
\(\mathsf B_{\mathsf N}\). The other three
boundary kinds have amplitude at most one
each on the declared step box once the
coupling ceiling makes their fluctuation
factor times the inherited sizes small.
Thus \(\mathsf B^{\mathbf T}\) can be taken
as \(\mathsf B_{\mathsf N}+3\), with the
ceiling chosen after the inherited bounds.
The stage and main-step origins require
their own formation amplitudes on their
actual domains.

%
For the interaction part of the stage amplitude,
factor the old term evaluated at the new stage
background with zero collar field out of the
stage integral. The integrated field's switch
is removed at that evaluation. The remaining
potential is its difference from that evaluation.
Each cluster of the input-dependent logarithm
therefore contains such a difference vertex.
The stage Cauchy bound gives it a factor
\(C B_*(\log\theta_t)^{-e}\), where \(B_*\)
bounds the inherited entries on the admitted
collar charts. Sum the rooted cluster bounds
over the performed stages. Even when the
collars overlap, the count is at most
\(C(\log\theta_t)^r\), and \(e>r\) gives a
uniform bound at the stage rate. The
input-independent normalizers use their own
summability estimate, with the field-independent
scalars kept outside the carrier logarithms.
For a fixed domain, the sum of the stage terms
formed on it occupies one stage slot. The Wilson
profile differences enter the action profile
sum separately.

The identity on the retained variables of
a standing list is the sum over the real
histories that return it. The resulting
measure uses these filed terms and keeps
the stored variables of that list. Hence
a uniform bound over the actual output
lists gives the bound at the next exterior
disintegration, which uses this positive
measure. There is no separate exponent of
each history in that sum.

%
%
The same filing applies at source \(w\) and at zero. A carried or
refilled member keeps its birth function and field space. The
output read uses the cut old list and the restriction of the
stored tuple to the carrier. At the fields already integrated by
a stage, the old function is read at zero and the associated
parameters are omitted. This restriction is independent of \(w\).

On the fixed source class of radius \(r_{\rm s}\) at a stop \(k\),
the sourced step bounds the difference of a boundary member from
its zero-source value by
\[
 2r_{\rm s}^{-1}\|w\|_{\flat,K-k}\,
 B_{j,o}e^{-\kappa d_j^\natural(\mathcal S;X)} .
\]
Here \(B_{j,o}\) is the boundary majorant's amplitude at origin
\(o\) and \(\mathcal S\) is the label at which the member is read.
The difference reads \(w\) on \(X^\sim\). At a refilled slot the
donor length is at least the required output length, so apply the
same bound to the restriction on the output list. Apply the
sourced formation estimate to each newly formed boundary member
on its own domain. The pointed origin uses its new sum of
pointings, with the spare decay paying their count. Repeating the
step yields the bound in the surviving hull norm at the witness,
on the same source class. This uses the sourced induction with
its frozen frame, tube room and constant-source data. The output
evaluation must lie in the native restricted field space.

A polynomial in \(\log\theta_n\) can be included in
\(B_*\). The factor \(A_j\) counts the actual origins and
slot multiplicities; a fine-origin incidence is not counted
as one coarse incidence. Suppose every term being summed has
a boundary carrier as in \cref{lem:ccf-crossing-length}.

\begin{proposition}[The old boundary tail]
\label{prop:ccf-old-boundary-tail}
If \(\kappa/2>c_E\), there are \(s_*,c,C>0\), depending
on the displayed rates and powers but not on \(n\), such that
for \(s\ge s_*\),
\begin{equation}
 \sum_{j\le n}A_j
     \sum_{X\text{ meeting }W\atop X\text{ has a boundary carrier}}
                         b_j(X)
                 \le CB_*N_W e^{-cs}.
 \label{ccf:old-boundary-tail}
\end{equation}
The same estimate bounds the oscillation when \(b_j(X)\)
is a bound on the complete evaluated local term on one
common chart; subtraction of two values costs a factor two.
\end{proposition}

\begin{proof}
Root a domain at a fine cube meeting \(W\). None of these
cubes is contained in an exempt hull. Spend half its decay
on \eqref{ccf:crossing-length}, and sum the remaining half
with \cref{lem:ccf-hull-root-entropy}. The level-\(j\)
contribution is bounded by
\[
 CB_jA_jN_WL^{4(n-j)}e^{-c\kappa L^{n-j}s}.
\]
Put \(a=n-j\). Equation \eqref{ccf:boundary-majorants}
leaves the scalar sum
\(\sum_{a\ge0}(1+a)^pL^{(q+4)a}e^{-c\kappa L^as}\).
Factor out \(e^{-c\kappa s/2}\). The remaining sum is
uniformly bounded for \(s\ge s_*\): its ratio at
successive integers is bounded by a fixed power times
\(e^{-c\kappa s(L-1)L^a/2}\), which is at most one half
once \(s_*\) is fixed sufficiently large. The finite initial
factor is harmless. This proves
\eqref{ccf:old-boundary-tail} without a cutoff-dependent
sum over old levels.
\end{proof}

For example, choose
\begin{equation}
 R\ge C_{\rm tr}\check R_n+C_0+s_*+C_1\log x.
 \label{ccf:logarithmic-clean-budget}
\end{equation}
Then the tail is bounded by \(CB_*N_Wx^{-cC_1}\).
If \(\check R_n\le C(\log x)^r\), the enlarged window
still has only \(O((\operatorname{diam}W+R)^4)\) current
cubes, hence a polylogarithmic number of retained coarse
link variables. The cell margin and the exponential margin
are paid separately in \eqref{ccf:logarithmic-clean-budget}.

\subsubsection{The input which the clean region retains}

Equation \eqref{ccf:local-recipe-identity} compares two
recipes with the same coherent input list. A regular marked
member generated from a previously healed region belongs to
that list even when the standing trace is empty. Its small
norm is estimated as a marked input, rather than set to zero.
Likewise, the complementary Gaussian cutoff factors used in
the finite expansion of a passed sharp test can have their
roots inside the clean region. Their size is the rooted
Gaussian tail of \eqref{fm:rooted-cutoff-logarithm}; they are
not boundary-history carriers.

Thus the local complete exponent separates into the
whole-lattice Gaussian and regular interaction recipe, the
retained small marked input, the cutoff correction, and the
boundary tail \eqref{ccf:old-boundary-tail}. To identify a
conditional law in retained coarse coordinates one must in
addition identify its physical background substitution and
its coarea/domain image on one admitted chart. A bound for
\(E_h(X,P)\) as a function of \(P\) does not prove that
two physical substitutions have the same arguments. The
geometric proof here establishes the local recipe and its
boundary-carrier tail before that separate analytic step.

% END LAYERED EXPANSION formation/clean_collar_filing.tex


% BEGIN PROOF EXPANSION positive_history_comparison.tex
\subsubsection*{The generating function on the fluctuation segment}
\phantomsection\label{f2:witness-generating-function}

Write \(\mathsf r=\hat{\mathsf d}\), \(g=x^{-1/2}\) and
\(\bar v=gp_0\). Let \(\phi_a=\widetilde f_a\), \(a=1,2\),
and let \(X_a\) be the two bounded, windowed observables
of the witness field. The tents move a support by at most
\(CMa_s\), so the enlarged supports are still separated
by at least \(\mathsf d/2\) when the room constant is
large. Use the union of their bare plaquette sets in the
readout construction. Each set has at most
\(C\mathsf r^4\eta_{k_0}^{-4}\) plaquettes.

For step \(j\), let \(\gamma_j^0\) be the flat Gaussian
marginal on the free read bonds and let \(\mathcal W_j\)
be their product of balls. Denote by \(T_j\) a uniform
bound for the change of expectation of a function of
these bonds of supremum at most one, from the rotated
actual law at a compatible exterior to
\(\gamma_j^0(\cdot\mid\mathcal W_j)\). The bounds above
give, on the admitted reference branches,
\[
 T_j\le C\left[
       e^{2\epsilon_j^{\mathrm{old}}}-1
       +e^{\Xi_j}-1
       +n_j \widehat g_{j-1}b_j^3+t_j+n_j\varpi_j
       \right],
 \qquad
 \Xi_j\le Cn_jb_j^2e^{-c\mathsf D}.
\]
Here \(t_j\le Cn_j e^{-cb_j^2}\) also bounds the
Gaussian error from the box and read-set cutoffs.
%
For this bounded generating function remove both
the old-origin tilt and the intrinsic remainder
at each original fixed exterior \(e\). If \(P\)
is a removed part, the density ratio on the
product window is \(e^P/\langle e^P\rangle\),
where the expectation uses the normalized
law after that removal. Thus the ratio lies
between \(e^{-\operatorname{osc}P}\) and
\(e^{\operatorname{osc}P}\). The old-origin
bound gives the first term in \(T_j\).
The full intrinsic oscillation at the original
stored values is at most
\(Cn_j\widehat g_{j-1}b_j^3\)
by \cref{cl:intrinsic-reference-hessian}.
This is the third term in \(T_j\) at weak
coupling. Both removals use the same product
window and the same physical background.

The resulting quadratic law on the box has
precision \(Q^e_{II}\) and linear term
\(Q^e_{IO}y_O+(\ell_E^e)|_I\).
Compare it with the centered law on the
same window whose precision is
\(K_e=(((Q^e)^{-1})_{II})^{-1}\).
The force difference is exactly
\[
 (Q^e_{II}-K_e)y+Q^e_{IO}y_O+(\ell_E^e)|_I .
\]
The first term is the Schur correction.
All three terms are estimated with the
quadratic weighted row bounds at \(e\).
The gap and marginal-density argument
give the term with \(\Xi_j\), without a
comparison of nonlinear evaluations at
different exterior values. Remove the
centered Gaussian box window, compare
its read covariance with the flat one,
and impose the flat read window. The two
window changes are included in \(t_j\).
Use the same Wilson ratio
\(x_{j-1}/\theta_{j-1}\) in the flat precision,
as in \hyperref[f2:flat-box-marginal]{the Gaussian comparison in the rotated coordinates}.
The final flat read law is the same for
every compatible exterior. Average the
normalized estimates against the original
positive exterior mixture.

Take the source radii
\[
 |\lambda_a|\le\Lambda_a
   =\frac{c_*x}{2\mathsf r^4p_0^2\|\phi_a\|_\infty}
 \quad(a=1,2),
 \qquad z_p=\sum_{a=1}^2\lambda_a\phi_a(x(p)).
\]
Zero profiles may be omitted. The constants below
hold with a fixed reserve in these radii. Let \(\xi_j\)
be independent flat read variables with the windowed
law just specified. Put
\[
 M_{\le j}(p)
   =i\sum_{k_0<l\le j}\widehat g_{l-1}
       \bigl(H^\square_{l-1}\mathsf C^0\xi_l\bigr)(p),
 \qquad
 Z_j(P)=\mathbb E\exp\!\left(
      \tfrac12\sum_p z_p|M_{\le j}(p)+P(p)|^2\right).
\]
The norm in the exponent is a norm of real-field
matrices. The complex variables are the source
coefficients. The row bounds of the linear readout,
\(b_l\le Cp_0\) and the geometric sum of \(\eta_{l-1}^2\)
give \(\sup_p|M_{\le j}(p)|\le C\eta_{k_0}^2\bar v\).
The same size holds for the remaining plaquette
backgrounds and the intermediate arguments in a
one-step replacement. Together with the count of
plaquettes this gives \(|Z_j|\le C\) and
\[
 |Z_j(P)-Z_j(P')|
    \le C\frac{\sup_p|P(p)-P'(p)|}{\eta_{k_0}^2\bar v}.
\]
Indeed one integrates the change of the exponent
along the segment between \(P\) and \(P'\). Both
bounds hold on the same source polydisc.

At each remaining level use one prescribed complete
axial tree on the common read box, with its root and
all paths included in the read. Gauge-transforming
the field before putting it in this fixed gauge
changes the resulting field by one constant rotation.
Use the tower frames of the localized readouts within
this gauge. Simultaneous rotation of all arguments
of \(Z_j\) leaves it unchanged, since the flat laws
and their balls are colour invariant and the flat
maps act separately on each colour. Hence
\[
                 Y_j(V_j)=Z_j(P^{(j)}(A^{\rm ax}(V_j)))
\]
is a bounded gauge-invariant function on the admitted
readout domain. It can be set to one outside that
domain. The read sets contain the entire dependence
on the tree and the old readouts. In particular an
operation on an exterior component of a clean string
leaves that read unchanged. At \(j=k_0\) take
\(Y_{k_0}=e^{\lambda_1X_1+\lambda_2X_2}\) on the
whole space. On the clean branch it is \(Z_{k_0}\)
of the localized plaquette defects. The equality uses
\(a_p(U)=\tfrac12|U(\partial p)-1|^2\) in the
normalized \(\mathrm{SU}(2)\) norm and the fact that
the witness tests pass there.

The one-step identity for the localized defects
replaces the argument of \(Z_{j-1}\), on the actual
fluctuation window, by
\[
 i \widehat g_{j-1}H^\square_{j-1}\mathsf C^0\widetilde y
                  +P^{(j)}(A^{\rm ax}(V_j)).
\]
Its pointwise error is bounded by
\(C\eta_{j-1}^2[(\mathsf r gp_0)^2+gp_0e^{-c\Delta_j}]\),
with \(\Delta_j\) the native readout margin. This
comparison uses one colour convention for both
supports. The residual constant rotation drops out
of \(Z_{j-1}\). The remaining expression depends
only on the free read coordinates and the retained
\(V_j\). The \(T_j\) bound and one integration with
\(\gamma_j^0(\cdot\mid\mathcal W_j)\) now give
\[
 \left|\mathbb E_j^{\mathrm{ref}}Y_{j-1}-Y_j\right|
   \le C\left[
     L^{-2(j-1-k_0)}
             \bigl(\mathsf r^2gp_0+e^{-c\Delta_j}\bigr)
       +T_j\right].
\]
For the unconditioned comparison below,
\(\mathbb E_j^{\mathrm{ref}}\) is the box law obtained from
the original adjacent-level law after summing the standing
parts and holding the compatible exterior fixed. Its good
current selector is then constant on the product window.
The comparison to \(\gamma_j^0\) is included in \(T_j\).
The last integration is exactly the independent
addition of the \(j\)-th slice in the definition of
\(Z_j\). No change to the exterior weights is made.

%
%
For the unconditioned two-point function use the
positive joint law of each adjacent pair of levels.
Let \(\mu_j\) be the standing law with its stored
variables retained, and let \(\pi_j\) be the joint
law of the original step from \(\mu_{j-1}\) to
\(\mu_j\), including the string of moves. Its
marginals give the exact telescoping identity
\[
 \mathbb E_{\mu_{k_0}}Y_{k_0}-\mathbb E_{\mu_{k_1}}Y_{k_1}
  =\sum_{j=k_0+1}^{k_1}
        \mathbb E_{\pi_j}(Y_{j-1}-Y_j).
\]
In the \(j\)-th term sum the earlier histories
over their standing parts in both the inserted
and uninserted measures. Use the good current
condition on the complete read box when it
selects an exterior weight of the resulting
integrals. At fixed exterior, the indicator
comparison for that step makes it constant
on the product window, with the redundant
interior flags summed out. Disintegrate at the
ordinary output before the large-field moves.
On the good current part those moves are exterior
to the read and preserve \(Y_j\), including
under the gauge changes. The coarea integral
with the insertion has the same weight as the
integral without it. After normalization the
box expectation is mixed with the exterior law
weighted by the uninserted box integral.
If \(d_j\) denotes the bound for the box
comparison and \(p_j\) the mass of the
excluded current part, the \(j\)-th term has
modulus at most \(d_j+2C p_j\), using the bounded
extensions of both \(Y\)'s. For a good condition
defined on finer strings, we must verify this
factorization after the standing-part sum.

Thus this comparison uses the summed one-step
exceptional mass. It allows the earlier healed
histories in the incoming standing law. On the
prescribed reference torus the age and creation
bounds control the probability of a bad current
part by summing the cells that meet the padded
read, with the lag estimate where the reference
requires it. A comparison conditioned on a
terminal exterior and on the full string uses the
same-exterior estimates separately. If a finer class
changes this box law by a normalized Radon--Nikodym
factor in \([e^{-\epsilon},e^{\epsilon}]\), the expectation
of the bounded insertion changes by at most
\((e^{\epsilon}-1)\|Y_{j-1}\|_\infty\). In that comparison
charge the class restriction at the step where it enters.

\paragraph{The top field and the Gaussian windows.}
%
%
\hypertarget{witness-top-native-tests}{}
To bound the last field, first use the native level-\(k_1\)
plaquette tests. The top read of the comparison uses its
own localized minimizer. Its box need not be a test cube.
Include the complete backward corner blocks for every top
plaquette in its determining read. Assign each coarse
plaquette to its native test cube by its barycenter. One
test layer is wider than the complete corner stencil.
At a passed native test
the tested background has bare plaquette defect at most
\(\varepsilon_{k_1}\eta_{k_1}^2\) on those blocks.
It has the prescribed level-\(k_1\) averages on the four
rows of the plaquette. The local averaging estimate of
\cref{prop:averaging-1a-1-4} therefore gives
\(\left|\partial V_{k_1}-1\right|\le2\varepsilon_{k_1}\)
there. The localized Landau curvature estimate on the top
comparison box then gives
\[
 \sup_p|P^{(k_1)}(p)|
 \le C\eta_{k_1}^2\varepsilon_{k_1}
 \le C\eta_{k_1}^2gp_0 .
\]
The last inequality uses the running coupling bound over the
segment. The transfer from the native test uses the scalar
curvature of the block datum.

A failed or undefined native test is a creation. A skipped
test lies in the margin of an earlier live component.
An operation at the top touching a complete read has a
live predecessor. Thus the terminal exception is covered
by the live and creation cells of levels \(k_1-1,k_1\)
meeting the padded reads and operation margins, before the
terminal field is formed. On the clean part, subsequent
exterior moves do not read or change the determining data
of these plaquettes. The same finite-band rarity estimate
bounds \(p_{\rm top}\), with the top stencils included in its
cell count.
On the part of the original terminal law where this
test fails, set the top argument \(P\) to zero. The
bounded extensions allow this replacement at cost
\(Cp_{\rm top}\), where \(p_{\rm top}\) is the mass
of that part. Keep the original terminal law in the
expectation of the resulting \(Z_{k_1}(P)\).

%
%
Write \(Z_\infty(P)\) for the unwindowed flat read
expectation and \(C_{\mathrm{loc}}\) for its one-colour
covariance. Set \(F_a=\operatorname{diag}(\phi_a(x(p)))\)
and \(h=\lambda_1F_1+\lambda_2F_2\). Decompose the real
matrix field \(P\) orthogonally into its component in the
three-colour Gaussian subspace, with coordinates \(p_c\),
and its perpendicular component \(P^\perp\). Then
\[
 \frac{Z_\infty(P)}{Z_\infty(0)}
  =\exp E_P(\lambda),\qquad
 E_P(\lambda)=\frac12\sum_{c=1}^3
    p_c^{\mathsf T} h(1-C_{\mathrm{loc}}h)^{-1}p_c
        +\frac12\sum_p h_{pp}|P^\perp(p)|^2.
\]
To obtain this identity, write each colour as
\(C_{\mathrm{loc}}^{1/2}\) times a standard Gaussian
and complete the square for real sources. This also
works for singular \(C_{\mathrm{loc}}\). The weighted
trace bound gives convergence and continuation to
the source polydisc. More explicitly, put
\(D_\lambda=\sum_a|\lambda_a||F_a|\) and write
\(h=D_\lambda^{1/2}S D_\lambda^{1/2}\) with
\(\|S\|\le1\). The trace of
\(D_\lambda^{1/2}C_{\mathrm{loc}}D_\lambda^{1/2}\)
is at most a fixed constant less than one, also for
the doubled radii after reducing \(c_*\).
The Neumann series therefore bounds \(|E_P|\) by
\(C\sum_p (D_\lambda)_{pp}|P(p)|^2\).
The plaquette count, the radii and the top-field
bound make this \(O(L^{-4\mathsf n})\).

Let \(t=\sum_j Cn_j e^{-cb_j^2}\). The joint flat
window conditions the full Gaussian on an event of
complement probability at most \(t\). At each fixed
\(P\) the determinant and the bound for \(E_P\), applied
to the doubled real sources, bound the second moment
of the insertion uniformly. Cauchy--Schwarz bounds
the deleted integral by \(C\sqrt t\), and the
conditioning denominator costs \(Ct\).
Consequently the terminal mixture
\[
  Z_{\mathrm{term}}
     =Z_\infty(0)\,\mathbb E_{\mathrm{term}} e^{E_P}
\]
approximates the windowed mixture within \(C\sqrt t\).
Here \(\mathbb E_{\mathrm{term}}\) is expectation under
the original terminal law with the argument \(P\)
defined above. The Gaussian integral is evaluated
at each fixed terminal datum before taking that
expectation.

Finally, the local/global flat comparison replaces
the centered determinant \(Z_\infty(0)\) by the
whole-torus one,
with error \(Cp_0^{-2}e^{-cd_*}\).
Write \(Z_{\mathrm{free}}\) for this determinant.
Its covariance in one colour is
\(C^x=\sum_{j=k_0+1}^{k_1}x_{j-1}^{-1}\mathcal C_{(j)}\).
The trace series gives a fixed positive lower
bound for \(|Z_{\mathrm{free}}|\) on the polydisc.

For the unconditioned law the generating function
differs from \(Z_{\mathrm{term}}\) by at most
\[
 C\left[
  \mathsf r^2gp_0+\max_j e^{-c\Delta_j}
   +\sum_j T_j+\sqrt t+\sum_jp_j+p_{\rm top}
 \right].
\]
The joint laws and their conditional kernels can
be chosen independently
of \(\lambda\). The bounds may first be imposed on
a countable dense subset of the polydisc and then
extended by analyticity of the bounded insertions.
This compares the generating functions before
taking logarithms.

\paragraph{Returning to the physical two-point function.}
The trace series bounds \(|Z_\infty(0)|\) below.
For large \(x\), the bound on \(E_P\) also bounds
\(|Z_{\mathrm{term}}|\) below. The quotients of the
witness generating function by \(Z_{\mathrm{term}}\)
and of \(Z_\infty(0)\) by \(Z_{\mathrm{free}}\) lie
in a fixed disc about one when the corresponding
errors are small. The two log differences then have
supremum \(C\delta\), with \(\delta\) the sum of the
generating-function error above and \(p_0^{-2}e^{-cd_*}\).
Mixed Cauchy differentiation costs
\(C\delta/(\Lambda_1\Lambda_2)\).
Relative to \(x^{-2}\|\phi_1\|_\infty\|\phi_2\|_\infty\)
this is the factor \(C(\mathsf r^4p_0^2)^2\) times
\(\delta\).

The remaining mixed derivative is that of
\(\log\mathbb E_{\mathrm{term}} e^{E_P}\) at zero.
Put
\[
 A_a(P)=\frac12\sum_p\phi_a(x(p))|P(p)|^2.
\]
Expansion of the resolvent to first order in \(h\)
gives this derivative as
\[
 \mathbb E_{\mathrm{term}}\sum_{c=1}^3
       p_c^{\mathsf T}F_1 C_{\mathrm{loc}}F_2p_c
       +\operatorname{Cov}_{\mathrm{term}}(A_1,A_2).
\]
Indeed, \(\partial_a E_P(0)=A_a(P)\) and the two
orders of \(F_1,C_{\mathrm{loc}},F_2\) give equal
quadratic forms by transposition. The covariance
comes from differentiating the logarithm after
averaging. On the union of the test supports
\(|C_{\mathrm{loc}}(p,p')|\le Cx^{-1}\eta_{k_0}^4\)
by the diagonal bound and positivity. With at most
\(C\mathsf r^4\eta_{k_0}^{-4}\) plaquettes and
\(|P|\le C\eta_{k_0}^2gp_0L^{-2\mathsf n}\),
the two terms are bounded at scale
\(x^{-2}\|\phi_1\|_\infty\|\phi_2\|_\infty\) by
\(C\mathsf r^8p_0^2L^{-4\mathsf n}\) and
\(C\mathsf r^8p_0^4L^{-8\mathsf n}\), respectively.

%
%
Take \(\mathsf D=A\log x\), with \(A\) fixed sufficiently large.
This fits the test margins since \(\check R_{k_0}/\log x\to\infty\).
The witness room gives \(\mathsf r\asymp p_0^2\), and
\(\mathsf D+\check R_{k_0}=O(\mathsf r)\) in the chosen range
of exponents. The box count gives \(n_j\le C\mathsf r^4\),
and \(\mathsf n=O(\log x)\). Thus the activity contribution to
the old oscillations, estimated with one power of the ratio of
radii, obeys
\[
 \sum_j C n_j(\log x)^{\mathfrak p_0-q}
 \le C(\log x)^{1+9\mathfrak p_0-q}.
\]
After the mixed Cauchy loss this is
\(O((\log x)^{1+29\mathfrak p_0-q})\), so take
\(q>1+29\mathfrak p_0\) as well as the earlier lower bounds.
On the admitted branches the other old terms use the power and
collar bounds in the old oscillation estimate. The intrinsic
and covariance terms involving \(x^{-1/2}\) have only fixed
logarithmic prefactors. The two collar distances and the box
margin give \(e^{-c\mathsf D}\) up to fixed constants, which
also absorbs any fixed logarithmic power. The window error
uses \(b_j\asymp p_0\), and \(L^{-4\mathsf n}\le Cx^{-8}\).

For the geometric exceptions, the union of padded reads fits
in the level-\(k_1\) neighborhood counted above. Its cover
at level \(j\) costs at most
\(C(\mathsf r+\mathsf D+\check R_{k_0})^4 L^{4(k_1-j)}\)
cells. Sum the one-cell bounds for live and creation events
over these covers, including the predecessors of the top
tests. This bounds their contribution to
\(p_{\rm top}+\sum_jp_j\) by
\[
 C(1+\mathsf n)(\mathsf r+\mathsf D+\check R_{k_0})^4
             L^{4\mathsf n}e^{-cp_0}.
\]
Here \(L^{4\mathsf n}=O(x^8)\) and \(\mathfrak p_0>1\), so
this too is negligible after the Cauchy loss. This count
uses the indicated geometric exceptions. When the
native decomposition with attached terms supplies the
conditional box estimate, charge the failure of its
per-bond budget separately. Its bound summed over
the segment is \(C\sum_j n_j e^{-c p_0}\), also
negligible after the Cauchy loss. For this native
budget split the adjacent law at the incoming list of
that step. The indicator is fixed at each exterior,
so the box normalization on the passing part is
unchanged. On its complement the bounded difference
of the two insertions costs at most a constant
times the failure mass. Other refinements must still
use the factorization of their actual history law.

The flat covariance contribution is
\(\frac32\sum_{p,p'}\phi_1(x(p))\phi_2(x(p'))C^x(p,p')^2\).
In the linear combination of the slices, replacing
\(x_{j-1}^{-1}\) by \(x^{-1}\) costs at most
\(C(B_\sharp\mathsf n+c_5)x^{-2}\) times the sum
of the absolute slice kernels. On separated supports
that sum is bounded by \(C\eta_{k_0}^4\mathsf r^{-4}\)
by the exponential slice bound and its geometric sum.
Multiplying by the sum of the two covariances to
compare their squares, and then summing the plaquettes,
gives a relative error \(C(B_\sharp\mathsf n+c_5)/x\)
in the smeared covariance at scale \(x^{-2}\).
The cross term and square from the
omitted coarse covariance give the displayed
coarse-complement error of the free calculation.

Apply this conclusion in the physical-source
identity at \(k_0\). There the small mixed-source
term and the errors in the two first derivatives
are the source estimates in (a) and (b).
The leading covariance is
\(\zeta'_{k_0}{}^2\langle X_1;X_2\rangle\).
This identifies the covariance of the bounded witness
observables in the physical-source identity. The
source remainder bounds at \(k_0\), the
nondegeneracy for the chosen class of tests,
and the reference-law transfer to the stated
Gibbs states keep their separate roles.

\subsection{Positive history sums and averaged removal of exceptional runs}
\label{ch:positive-history-comparison}

The passage from a run window to a connected correlation has two
normalizations. One first sums the positive histories that represent the
same physical event, and then normalizes their total mass. Conditioning
each history separately generally gives different weights. The purpose
of this section is to perform these operations in their correct order
and to estimate the error from exceptional runs without a uniform
lower bound on the clean probability at every exterior field.
The resulting estimates concern the second and fourth connected
correlations used in the short-distance argument.

All probability measures in this section are real and positive at zero
source. A differentiated polymer resolution can have signed or complex
terms; it must first be identified with the positive physical measure
before the statements below apply. No estimate on the absolute values
of a different resolution makes that identification automatically.

\subsubsection{Pushforward measures and the multiplicity of a history class}

Let $H$ be finite or countable. For $h\in H$ let $\sigma_h$ be a finite
positive measure on a standard Borel space $\Omega_h$, and let
$T_h:\Omega_h\longrightarrow\Omega$ be measurable, where $\Omega$ is
also standard Borel. Assume
\begin{equation}\label{phc:finite-total-mass}
 0<Z:=\sum_{h\in H}\sigma_h(\Omega_h)<\infty.
\end{equation}
The common physical measure is
\begin{equation}\label{phc:physical-pushforward}
 \mu=Z^{-1}\sum_{h\in H}(T_h)_*\sigma_h.
\end{equation}
For a nonnegative measurable $F$, its integral is the sum of the
integrals of $F\circ T_h$. This follows first for indicators from
the definition of a pushforward and then for general $F$ by monotone
convergence. For an integrable complex $F$, apply this result to its
absolute value before taking real and imaginary parts.

Suppose $H$ is partitioned into classes $\mathcal C$. Define the
unnormalized measure of a class by
\begin{equation}\label{phc:class-measure}
 \Sigma_C=\sum_{h\in C}(T_h)_*\sigma_h,
 \qquad C\in\mathcal C.
\end{equation}
Then $Z\mu=\sum_C\Sigma_C$. This is an equality of positive
measures, including all Jacobians and all tests present in the
individual histories. Rearrangement of the sum is justified by
positivity. It does not require a choice of one representative
configuration for the class.

\begin{proposition}[The density of a resummed class]
\label{prop:phc-class-density}
Suppose there is a positive measure $\nu_C$ such that
$(T_h)_*\sigma_h=a_h\nu_C$ for every $h\in C$, with $a_h\ge0$
measurable. Then
\begin{equation}\label{phc:class-density}
 \Sigma_C=a_C\nu_C,
 \qquad a_C=\sum_{h\in C}a_h.
\end{equation}
In particular the multiplicity can be taken outside an integral over
the class variables only if $a_C$ is constant almost everywhere in
those variables.
\end{proposition}
\begin{proof}
For every nonnegative $F$, monotone convergence gives
$\sum_h\int F a_h\,d\nu_C=\int F\sum_h a_h\,d\nu_C$.
This proves the equality of measures. A constant factor outside all
such integrals is equivalent to a constant Radon--Nikodym density.
\end{proof}

Equality of the coordinate domains, or independence of the histories
from the source parameter, does not imply constancy of $a_C$.
For example the two measures $dy$ and $(1+y^2)\,dy$ on $[-1,1]$
have the same domain and no source dependence, but their sum has
density $2+y^2$. Both densities contribute to a source integral
$\int e^{zX(y)}\,d\Sigma_C(y)$.
If an exact change of variables shows that all the pushforwards are
equal to the same measure, the finite cardinality of the class is a
constant multiplicity. For an infinite class of identical nonzero
positive finite measures, \eqref{phc:finite-total-mass} would fail.
Thus an infinite counting factor cannot be discarded using a
formal cancellation of infinite normalizers.

When $a_C>0$ is differentiable, $\log a_C$ belongs in the complete
conditional exponent of \cref{ch:clean-conditional-law}. There is
an additional derivative contribution even if every summand has
a simple Hessian. For a finite class write $a_h=e^{p_h}$ and
$w_h=e^{p_h}/\sum_j e^{p_j}$. Direct differentiation yields
\begin{align}
 \nabla\log a_C&=\sum_h w_h\nabla p_h,
 \label{phc:sum-force}\\
 \nabla^2\log a_C
 &=\sum_h w_h\nabla^2p_h
   +\sum_h w_h(\nabla p_h-\bar p)
                   (\nabla p_h-\bar p)^{\mathsf T},
 \quad \bar p=\sum_h w_h\nabla p_h.
 \label{phc:sum-hessian}
\end{align}
For a countable class, a sufficient condition is that, on every compact
subset \(K\) of the coordinate chart,
\[
 \sum_{h\in C}\sup_K
 a_h\big(1+|\nabla p_h|^2+\|\nabla^2p_h\|\big)<\infty,
 \qquad a_h=e^{p_h},
\]
and that \(a_C=\sum_h a_h\) stays positive. These bounds justify
differentiation and each separate weighted sum in
\eqref{phc:sum-force} and \eqref{phc:sum-hessian}. Uniform convergence of
\(\sum_h a_h\) and its first two derivative series alone does not
justify separating two divergent terms in the logarithmic Hessian.
\par{}The second term in \eqref{phc:sum-hessian} is the covariance of
the slopes of the summands. For $p_\pm(y)=\pm My$, each individual
Hessian vanishes, whereas $(\log(2\cosh My))''$ equals $M^2$
at zero. Control of individual Hessians consequently does not
control the Hessian after a positive history sum. The slopes or
the complete summed density must also be controlled.

\subsubsection{The exterior law at zero source}

Let $\pi:\Omega\longrightarrow E$ be the exterior variable and let
$\lambda=\pi_*\mu$. Standard Borel spaces admit a probability kernel
$e\mapsto\mu_e$ with
\begin{equation}\label{phc:disintegration}
 \int F\,d\mu=\int_E\left(\int F\,d\mu_e\right)\lambda(de)
\end{equation}
for every integrable $F$. Equalities involving this kernel are
understood outside a $\lambda$-null set, on which any fixed version
may be used. At a finite lattice cutoff all spaces used here are
finite products of compact Lie groups, Euclidean spaces and
discrete history sets, so this measurability setting applies.

The weights of separate histories in $\mu_e$ involve their
conditional masses at $e$. In particular they are not in general
the unconditional weights $\sigma_h(\Omega_h)/Z$.
One way to see the distinction is to retain the history label on
the disjoint union $\bigsqcup_h\Omega_h$, push it to the pair
$(h,e)$, and then disintegrate over $e$. The resulting conditional
law of $h$ supplies exactly the required weights. Summing first
and disintegrating first give the same joint integrals when these
conditional weights are retained, by \eqref{phc:disintegration}.

For bounded real physical insertions $X_1,\ldots,X_n$, put
\begin{equation}\label{phc:frozen-exterior-mgf}
 F(z)=\int_E\mu_e\!\left(e^{\sum_i z_iX_i}\right)\lambda(de).
\end{equation}
The exterior law $\lambda$ in this formula is the law at zero source
and stays fixed during differentiation. This follows by inserting
the exponential in \eqref{phc:disintegration}; no approximation is
being made. All changes of exterior weights induced by a source
are already present in the conditional generating functions.
Factoring out an $e$-dependent source normalizer from those
functions would require retaining that factor inside the outer
integral. Dividing it out before exterior averaging would change
$F$ and its connected derivatives.

\subsubsection{What an unconditional rarity estimate implies}

Let $C\subset\Omega$ be the clean event, measurable in the joint
physical and history space. Enlarge $\Omega$ to that joint space
when the event is not determined by the physical field alone.
Write
\begin{equation}\label{phc:clean-probability}
 q(e)=\mu_e(C),\qquad p=\mu(C^c)
                 =\int_E(1-q(e))\lambda(de).
\end{equation}
The number $p$ can be exponentially small although $q(e)$ is zero
on a set of exterior fields of positive measure. For example take
$E=\{0,1\}$, $\lambda(\{1\})=p$, and let cleanliness be exactly
the event $e=0$. Then $q(1)=0$. Thus no positive lower bound on
all the conditional clean probabilities follows from small $p$.

For $0<\eta<1$, define
\begin{equation}\label{phc:good-exterior}
 G_\eta=\{e:q(e)\ge\eta\}.
\end{equation}
Since $1-q>1-\eta$ on its complement, integration gives
\begin{equation}\label{phc:bad-exterior-probability}
 \lambda(G_\eta^c)\le\frac{p}{1-\eta}.
\end{equation}
More generally the set where $q<1-\delta$ has mass at most
$p/\delta$. The point is that the exceptional exteriors must be
retained in the error estimate; they do not disappear from a
conditional statement.

For $q(e)>0$, define the clean conditional kernel
\begin{equation}\label{phc:clean-kernel}
 \mu_e^C(F)=q(e)^{-1}\mu_e(F1_C).
\end{equation}
This kernel is a probability measure regardless of how small $q(e)$
is. Its existence says nothing about a uniform bound on its density
or its derivatives. Such bounds require the actual complete
conditional law. The next comparison uses only positivity.

\subsubsection{A resampling with the same exterior marginal}

It is convenient to compare laws on the feature space $\mathbb R^n$.
Let $X=(X_1,\ldots,X_n)$ be measurable and let
$\rho=X_*\mu$. Let $G\subset\{q>0\}$ be a measurable admitted
set of exteriors. For $e\in G$, use the law of $X$ under
$\mu_e^C$; for $e\notin G$, use the point mass at $0\in\mathbb R^n$.
Their mixture with respect to the unchanged $\lambda$ is a
probability measure $\nu_G$.

We use the total variation convention
$\|\rho-\nu\|_{\rm TV}=\sup_A|\rho(A)-\nu(A)|$ for probabilities.
For bounded $f$, the corresponding integral bound is
$|\rho(f)-\nu(f)|\le2\|f\|_\infty\|\rho-\nu\|_{\rm TV}$.
This convention is fixed to avoid an ambiguous factor of two.

\begin{proposition}[Averaged clean resampling]
\label{prop:phc-resampling}
With the preceding definitions,
\begin{equation}\label{phc:tv-resampling}
 \|\rho-\nu_G\|_{\rm TV}
       \le p+\lambda(G^c).
\end{equation}
If $G=\{q>0\}$, the right hand side can be improved to $p$.
If $G=G_\eta$, the bound $p+p/(1-\eta)$ is valid.
\end{proposition}
\begin{proof}
At a fixed $e$ with $q>0$, the full conditional probability is
$q\mu_e^C+(1-q)\mu_e^{C^c}$, with the second term omitted when
$q=1$. Its distance from $\mu_e^C$ is at most $1-q$.
Pushforward under $X$ does not increase total variation, since
preimages of feature events are physical events. At $e\notin G$
the distance from the chosen point mass is at most one. Integrate
these bounds and take the supremum over feature events to obtain
$\int_G(1-q)\,d\lambda+\lambda(G^c)\le p+\lambda(G^c)$.
When $G=\{q>0\}$, the integral of $1-q$ on $G^c$ is exactly
$\lambda(G^c)$, so the first displayed expression equals $p$.
The last assertion follows from \eqref{phc:bad-exterior-probability}.
\end{proof}

Global conditioning on $C$ gives another probability law,
$\mu(\,\cdot\mid C)$. Its exterior marginal is
$q(e)\lambda(de)/(1-p)$ when $p<1$. Thus global conditioning
and clean conditional resampling have different exterior marginals.
The latter is useful when a run-window estimate is to be integrated
against the original exterior law. This is the comparison that
\cref{prop:phc-resampling} supplies.

\subsubsection{Weighted exceptional moments}

A small exceptional probability alone does not control an unbounded
insertion. The estimate needed for a product of insertions can be
stated without global moment bounds on every physical field.
For each nonempty $S\subset\{1,\ldots,n\}$ let
\begin{equation}\label{phc:weighted-defect}
 d_S=\mu\!\left(1_{C^c}\prod_{i\in S}|X_i|\right),
 \qquad M_S=\prod_{i\in S}M_i,
\end{equation}
where $M_i>0$. Assume these numbers are finite and that
$|X_i|\le M_i$ on $C$, up to a $\mu$-null set. The clean bound
can use insertions from which fixed deterministic constants have
been subtracted. It cannot subtract an exterior-dependent random
constant without changing the global connected correlations.

\begin{proposition}[Mixed moment error with weighted tails]
\label{prop:phc-weighted-moments}
For every nonempty $S$,
\begin{align}
 \left|\rho\!\left(\prod_{i\in S}x_i\right)
       -\nu_G\!\left(\prod_{i\in S}x_i\right)\right|
 &\le d_S+[p+\lambda(G^c)]M_S,
 \label{phc:weighted-moment-error}\\
 \left|\rho\!\left(\prod_{i\in S}x_i\right)\right|
 &\le M_S+d_S,
 \qquad
 \left|\nu_G\!\left(\prod_{i\in S}x_i\right)\right|\le M_S.
 \label{phc:marginal-moment-bounds}
\end{align}
When $G=\{q>0\}$ the term $p+\lambda(G^c)$ in the first bound
can be replaced by $p$.
\end{proposition}
\begin{proof}
Write $X_S=\prod_{i\in S}X_i$. On $G$ the difference between a
full conditional moment and a clean conditional moment is
\begin{equation}\label{phc:conditional-moment-difference}
 \mu_e(X_S1_{C^c})-
       \frac{1-q(e)}{q(e)}\mu_e(X_S1_C).
\end{equation}
The absolute value of the second term is at most
$(1-q(e))M_S$, because $|\mu_e(X_S1_C)|\le q(e)M_S$.
This cancellation of $q(e)$ is the reason that a uniform clean
probability floor is unnecessary here. On $G^c$ the comparison
moment is zero. The full moment there has a bad part controlled
by the corresponding portion of $d_S$, and a clean part at most
$q(e)M_S$. Integrating these estimates gives
\begin{equation}\label{phc:sharper-weighted-error}
 d_S+M_S\left[\int_G(1-q)\,d\lambda
                         +\int_{G^c}q\,d\lambda\right].
\end{equation}
The bracket is at most $p+\lambda(G^c)$. If $G=\{q>0\}$,
its second integral is zero and its first is at most $p$.
Finally, the full absolute moment is at most $M_S+d_S$ and the
comparison law is supported in the clean feature box together
with the origin. This proves the remaining bounds.
\end{proof}

The defect $d_S$ is precisely where the Gaussian marks and unspent
reserves of \cref{ch:history-mass-program} are used. If the history
construction only bounds $p$, the weighted defects still require
proof. If it bounds the marked mass in the actual normalized
positive law, \cref{prop:phc-weighted-moments} converts that estimate
directly into a mixed moment error. An unnormalized mass estimate
must first be divided by a valid partition floor; the proposition
does not provide that floor.

For comparison, suppose $q_0>1$ and
$\|X_i\|_{L^{nq_0}(\mu)}\le B_i$. H\"older's inequality, first
on the exceptional event and then on the product, gives
\begin{equation}\label{phc:holder-defect}
 d_S\le p^{1-1/q_0}\prod_{i\in S}B_i,
 \qquad 1\le|S|\le n.
\end{equation}
Indeed the product has $L^{q_0}$ norm at most the product of the
$L^{|S|q_0}$ norms of its entries, and those norms are bounded by
the stated $L^{nq_0}$ norms on a probability space. In a fourth
connected correlation, uniform eighth moments therefore give a
square-root gain from rarity. A direct marked-mass estimate can
give a stronger gain and does not need these global eighth moments.

\subsubsection{Connected correlations as polynomials in mixed moments}

For nonempty $S$, write $m_S=\mathbb E\prod_{i\in S}X_i$.
The joint cumulant is the polynomial
\begin{equation}\label{phc:cumulant-polynomial}
 \kappa_S(X)=\sum_{\pi\in\Pi(S)}
       (-1)^{|\pi|-1}(|\pi|-1)!
                       \prod_{B\in\pi}m_B,
\end{equation}
where $\Pi(S)$ is the set of set partitions of $S$ into nonempty
blocks. The indices distinguish repeated insertions as well as
different ones. Consequently the same formula gives diagonal
moments and derivatives in repeated directions.

For bounded variables, \eqref{phc:cumulant-polynomial} follows by
expanding $\log(1+u)=\sum_{k\ge1}(-1)^{k-1}u^k/k$ in a
neighborhood of the origin and extracting the coefficient of
$\prod_{i\in S}z_i$ in
$\log\mathbb E e^{\sum z_iX_i}$. An unordered partition with
$k$ blocks arises in $k!$ orders in $u^k$, and the factor $1/k$
leaves $(k-1)!$. For variables with the finite mixed absolute
moments occurring in the formula, use this polynomial as the
definition; bounded truncations and dominated convergence recover
it from the usual moment definition. No moment generating
function is asserted to exist for an unbounded variable.

Let $\left\{\begin{smallmatrix}n\\k\end{smallmatrix}\right\}$
denote the number of partitions of $\{1,\ldots,n\}$ into $k$
blocks. Define
\begin{equation}\label{phc:stability-polynomial}
 b_n(D)=\sum_{k=1}^n k!
       \left\{\begin{matrix}n\\k\end{matrix}\right\}
                       (1+D)^{k-1},\qquad D\ge0.
\end{equation}
In particular
\begin{align}
 b_2(D)&=1+2(1+D),\label{phc:b-two}\\
 b_4(D)&=1+14(1+D)+36(1+D)^2+24(1+D)^3.
 \label{phc:b-four}
\end{align}
Thus $b_2(0)=3$ and $b_4(0)=75$.

\begin{proposition}[Connected error from weighted exceptional runs]
\label{prop:phc-connected-error}
Assume $d_S\le D M_S$ for every nonempty $S\subset\{1,\ldots,n\}$.
Put $r_G=p+\lambda(G^c)$. Then
\begin{equation}\label{phc:connected-error}
 |\kappa_\rho(x_1,\ldots,x_n)
        -\kappa_{\nu_G}(x_1,\ldots,x_n)|
 \le (D+r_G)b_n(D)\prod_{i=1}^n M_i.
\end{equation}
For $G=\{q>0\}$ one may use $r_G=p$.
\end{proposition}
\begin{proof}
For each partition with $k$ blocks, subtract the product of its
comparison moments from the product of its physical moments.
Use the telescoping product identity, replacing one factor at a
time. The replaced factor has error at most $(D+r_G)M_B$ by
\cref{prop:phc-weighted-moments}; every other factor has modulus
at most $(1+D)M_{B'}$. Since the blocks partition all indices,
the product of these bounds is
$k(D+r_G)(1+D)^{k-1}\prod_iM_i$.
Multiply by $(k-1)!$ from \eqref{phc:cumulant-polynomial}, sum
over partitions, and group them by $k$. This is exactly
\eqref{phc:stability-polynomial}.
\end{proof}

The same proof gives a useful comparison for two arbitrary bounded
feature laws with coordinate bounds $M_i$. If their total variation
distance is at most $\delta$, each mixed moment differs by at most
$2\delta M_S$, and all moments have modulus at most $M_S$.
Therefore
\begin{equation}\label{phc:bounded-tv-cumulants}
 |\kappa_\rho-\kappa_\nu|
              \le2\delta b_n(0)\prod_iM_i.
\end{equation}
The constants are $6\delta$ for a covariance and $150\delta$
for a fourth cumulant after the product of the bounds is factored
out. These constants are explicit sufficient bounds; no optimality
is needed in the application.

\subsubsection{Auxiliary source logarithms on the bounded feature box}

Bounded clean insertions also justify the connected interpretation
analytically. Let $\rho$ and $\nu$ be supported in
$\prod_i[-M_i,M_i]$. For
\begin{equation}\label{phc:auxiliary-source-region}
 t=\sum_i|z_i|M_i\le c<\log2,
\end{equation}
their generating functions satisfy
\begin{equation}\label{phc:zero-free-auxiliary}
 |F_\rho(z)-1|\le e^t-1<1,
 \qquad |F_\rho(z)|\ge2-e^c,
\end{equation}
and the same inequalities hold for $\nu$. The estimate follows
from $|e^w-1|\le e^{|w|}-1$ and integration. The principal
logarithm is therefore single valued throughout this region.

If the total variation distance is at most $\delta$, then
$|F_\rho-F_\nu|\le2\delta e^c$. The straight segment joining
their values remains in the disc $|w-1|\le e^c-1$, where
$|1/w|\le(2-e^c)^{-1}$. Integrating the derivative of the
logarithm along that segment gives
\begin{equation}\label{phc:log-comparison}
 |\log F_\rho(z)-\log F_\nu(z)|
       \le\frac{2\delta e^c}{2-e^c}.
\end{equation}
For a polydisc with radii $c/(nM_i)$, Cauchy's formula bounds
the derivative with one occurrence of every variable by the
right hand side times $n^n c^{-n}\prod_iM_i$. The polynomial
bound in \eqref{phc:bounded-tv-cumulants} is usually smaller for
$n=2,4$, but the logarithmic estimate explains the absence of
a normalization singularity.

These auxiliary sources are inserted into a fixed positive law at
zero source. Their region of holomorphy is not a claim that the
renormalization charts extend to sources of size $M_i^{-1}$.
One may use this elementary generating function to estimate a
real-law comparison without analytically continuing any history
map to that larger source domain. If the physical insertion itself
has been replaced by a derivative of a source-dependent effective
action, its higher source derivatives must be retained separately.

\subsubsection{The law of total cumulance}

The clean resampling still averages over the exterior field. The
conditional means cannot be discarded. The exact expression for
all the resulting terms is useful at order four.

Let $e\mapsto\xi_e$ be a probability kernel on $\mathbb R^n$,
let $\alpha$ be an exterior probability, and let
$\xi=\int\xi_e\,\alpha(de)$. For nonempty $S$ set
$c_S(e)=\kappa_{\xi_e}(x_i:i\in S)$. Suppose the coordinates
are bounded, or that all moments and integrals used below are
finite with an integrable dominating function. Then
\begin{equation}\label{phc:total-cumulance}
 \kappa_\xi(x_1,\ldots,x_n)
 =\sum_{\pi\in\Pi_n}
       \kappa_\alpha\big(c_B(e):B\in\pi\big).
\end{equation}
An outer cumulant of one entry means its expectation.

To prove the formula for bounded coordinates, write
$K_e(z)=\log\xi_e(e^{\sum_i z_ix_i})$ on a common small
polydisc. Its coefficient with distinct indices in $S$ is
$c_S(e)$. Expand
\begin{equation}\label{phc:total-cumulance-proof}
 \log\int e^{K_e(z)}\,\alpha(de)
  =\sum_{k\ge1}\frac1{k!}
       \kappa_\alpha(K_e(z),\ldots,K_e(z)).
\end{equation}
The series converges near zero because $K_e(0)=0$ and its
supremum tends to zero there. To obtain the coefficient of
$z_1\cdots z_n$, distribute the distinct indices among $k$
nonempty entries. These distributions are ordered partitions.
The factor $1/k!$ removes their ordering, giving
\eqref{phc:total-cumulance}. For finite moments one can instead
perform the same computation in the algebra of polynomials with
$z_i^2=0$. Every series truncates at degree $n$, so the proof
uses only the assumed moments and integrability.

At orders two and three the identity reads
\begin{align}
 \kappa_{12}&=\mathbb E_\alpha c_{12}
                   +\operatorname{Cov}_\alpha(c_1,c_2),
 \label{phc:total-two}\\
 \kappa_{123}&=\mathbb E_\alpha c_{123}
   +\operatorname{Cov}_\alpha(c_1,c_{23})
   +\operatorname{Cov}_\alpha(c_2,c_{13})
   +\operatorname{Cov}_\alpha(c_3,c_{12})
   +\kappa_\alpha(c_1,c_2,c_3).
 \label{phc:total-three}
\end{align}
At order four it reads
\begin{align}
 \kappa_{1234}={}&\mathbb E_\alpha c_{1234}
  +\sum_{i=1}^4\operatorname{Cov}_\alpha
                   (c_i,c_{\{1,2,3,4\}\setminus\{i\}})
       \notag\\
 &+\sum_{\{ij\}\{k\ell\}}
                   \operatorname{Cov}_\alpha(c_{ij},c_{k\ell})
       \notag\\
 &+\sum_{\{ij\},\{k,\ell\}=\{1,2,3,4\}\setminus\{i,j\}}
                   \kappa_\alpha(c_{ij},c_k,c_\ell)
       \notag\\
 &+\kappa_\alpha(c_1,c_2,c_3,c_4).
 \label{phc:total-four}
\end{align}
The second sum has the three pair partitions $12|34$, $13|24$,
$14|23$. The third has six choices of the pair $\{i,j\}$.
The displayed terms account respectively for partitions with
one, two, three and four blocks.

A mixture of centered Gaussian linear laws already illustrates
one essential term: its conditional fourth cumulants vanish,
but covariance of the conditional covariances can give a
nonzero global fourth cumulant. Similarly, for quadratic
Gaussian insertions the conditional one-point functions need
not be constant when the covariance depends on $e$. Subtracting
the global one-point function does not remove their exterior
variation. The mixed outer cumulants in
\cref{phc:total-two,phc:total-four} must be estimated.

\subsubsection{A fixed reference and quantitative exterior errors}

For $k\ge1$ define
\begin{equation}\label{phc:outer-cumulant-constant}
 a_k=\sum_{j=1}^k(j-1)!
              \left\{\begin{matrix}k\\j\end{matrix}\right\}.
\end{equation}
The first four values are $1,2,6,26$. The moment polynomial
implies $|\kappa(Y_1,\ldots,Y_k)|\le a_k\prod_i\|Y_i\|_\infty$.
For $k\ge2$ this cumulant is unchanged if a deterministic
constant is added to any one entry. For bounded variables the
invariance follows because such an addition changes the logarithm
of the generating function only by a linear term. The polynomial
identity extends it to finite moments.

\begin{proposition}[Exterior averaging around a deterministic reference]
\label{prop:phc-exterior-stability}
In the setting of \eqref{phc:total-cumulance}, suppose constants
$\gamma_S$ and nonnegative $\epsilon_S$ satisfy
\begin{equation}\label{phc:conditional-reference-errors}
 |c_S(e)-\gamma_S|\le\epsilon_S
\end{equation}
for every nonempty $S$, almost everywhere under $\alpha$. Then
\par\noindent
\begin{minipage}{\linewidth}
\begin{equation}\label{phc:exterior-error}
 |\kappa_\xi(x_1,\ldots,x_n)-\gamma_{\{1,\ldots,n\}}|
 \le\epsilon_{\{1,\ldots,n\}}
     +\sum_{\substack{\pi\in\Pi_n\\|\pi|\ge2}}
                       a_{|\pi|}\prod_{B\in\pi}\epsilon_B.
\end{equation}
\popQED\par
\end{minipage}
\end{proposition}
\begin{proof}
The one-block term of \eqref{phc:total-cumulance} has error at
most $\epsilon_{\{1,\ldots,n\}}$. In every other term subtract
$\gamma_B$ from each entry, using invariance of an outer
cumulant of order at least two under deterministic translations.
Its entries now have suprema at most $\epsilon_B$.
The bound following \eqref{phc:outer-cumulant-constant} and
the triangle inequality prove the result.
\end{proof}

For two insertions the bound is
\begin{equation}\label{phc:two-reference-error}
 \epsilon_{12}+2\epsilon_1\epsilon_2.
\end{equation}
For four insertions it is
\begin{align}
 \epsilon_{1234}
 &+2\sum_i\epsilon_i\epsilon_{\{1,2,3,4\}\setminus\{i\}}
  +2\sum_{\{ij\}\{k\ell\}}\epsilon_{ij}\epsilon_{k\ell}
       \notag\\
 &+6\sum_{\{ij\},\{k,\ell\}=\{1,2,3,4\}\setminus\{i,j\}}
                 \epsilon_{ij}\epsilon_k\epsilon_\ell
  +26\epsilon_1\epsilon_2\epsilon_3\epsilon_4.
 \label{phc:four-reference-error}
\end{align}
This form identifies exactly how accurately the conditional means
and the lower connected correlations must be controlled. It is
not necessary to demand the same absolute accuracy from quantities
that appear only as products. Conversely, a bound on the
conditional covariance alone is not enough for the fourth
connected correlation.

\subsubsection{Discarding unadmitted exteriors before reference comparison}

Suppose the conditional reference estimates hold only on $G$,
with $\lambda(G)>0$. Define the probability
$\widehat\nu_G=\lambda(G)^{-1}\int_G X_*\mu_e^C\,\lambda(de)$.
This removes the artificial point mass used in $\nu_G$.
Both measures are supported in the clean feature box, and
\begin{equation}\label{phc:discarded-exterior-tv}
 \|\nu_G-\widehat\nu_G\|_{\rm TV}\le\lambda(G^c).
\end{equation}
Indeed $\nu_G=\lambda(G)\widehat\nu_G+
\lambda(G^c)\delta_0$. The comparison with a deterministic
reference can now be made under the genuine exterior probability
$\lambda(\,\cdot\mid G)$.

Combining the preceding propositions gives the complete estimate
\begin{align}
 |\kappa_\rho-\gamma_{\{1,\ldots,n\}}|
 \le{}&(D+p+\lambda(G^c))b_n(D)\prod_iM_i
       \notag\\
 &+2\lambda(G^c)b_n(0)\prod_iM_i
       \notag\\
 &+\epsilon_{\{1,\ldots,n\}}
   +\sum_{\substack{\pi\in\Pi_n\\|\pi|\ge2}}
                         a_{|\pi|}\prod_{B\in\pi}\epsilon_B.
 \label{phc:complete-comparison}
\end{align}
Here $d_S\le DM_S$ and the conditional cumulants on $G$ are
within $\epsilon_S$ of $\gamma_S$. The first line compares
the physical law to clean resampling, the second removes the
exceptional exterior mass, and the remaining lines average the
conditional reference estimates. Every normalization in this
comparison is a normalization of a positive probability.

If a chart-admission set $A\subset E$ is also required, take
$G=A\cap G_\eta$. Its exceptional mass obeys
\begin{equation}\label{phc:admission-exception}
 \lambda(G^c)\le\lambda(A^c)+\frac{p}{1-\eta}.
\end{equation}
A model proof must supply the first term when admission is not
automatic. It cannot infer admission solely from the existence
of a conditional probability. The useful point is that a lower
bound on $q$ is only needed on the retained exteriors, and their
discarded mass is explicitly accounted for.

\subsubsection{Short-distance powers and the marked reserve}

We apply \eqref{phc:complete-comparison} with fixed order $n=2$
or $n=4$, inverse running coupling $x\longrightarrow\infty$,
and $N_i=\|f_i\|_\infty$. Suppose the admitted clean readouts
satisfy
\begin{equation}\label{phc:clean-readout-scales}
 M_i\le C_i x^{-1}(\log x)^{A_i}N_i,
 \qquad A_i\ge0.
\end{equation}
For the two-point readout bound used in
\cref{ch:two-point-extraction-tests}, one can take a fixed
exponent after substituting the permitted polylogarithmic
window radius. No optimization of that exponent is required;
it must be recorded before spending the rarity reserve.

Suppose the normalized marked mass gives, for fixed $c,b>0$,
\begin{equation}\label{phc:reserve-assumption}
 D+p+\lambda(A^c)\le C\exp[-c(\log x)^b].
\end{equation}
Then for fixed $\eta\in(0,1)$ the first two lines of
\eqref{phc:complete-comparison} are
\begin{equation}\label{phc:negligible-rare-error}
 o\!\left(x^{-n}\prod_iN_i\right).
\end{equation}
To check the assertion, $D\to0$, so $b_n(D)$ stays bounded.
The remaining factor, after $x^{-n}\prod_iN_i$ is removed,
is at most a constant times
$\exp[-c(\log x)^b+(\sum_iA_i)\log\log x]$.
Its exponent tends to $-\infty$ for every $b>0$.
The same conclusion follows from a power bound $O(x^{-a})$
with $a>0$. A logarithmic rarity bound
$O((\log x)^{-P})$ instead requires $P>\sum_iA_i$.

The estimate is a consequence of the normalized, marked hypothesis
\eqref{phc:reserve-assumption}. The unnormalized unmarked sum of
histories does not imply it. The distinction between a stored
volume mark and a field moment mark remains exactly the one in
\cref{ch:history-mass-program}; the present argument uses the
field moments in $d_S$ rather than replacing them by volume
moments.

For the conditional errors put $L_S=\prod_{i\in S}x^{-1}N_i$.
If the bounds on $G$ have the form
$\epsilon_S\le\delta_S(x)L_S$, then the last two lines of
\eqref{phc:complete-comparison}, divided by $L_{\{1,\ldots,n\}}$,
are at most
\begin{equation}\label{phc:relative-conditional-error}
 \delta_{\{1,\ldots,n\}}
 +\sum_{\substack{\pi\in\Pi_n\\|\pi|\ge2}}
                      a_{|\pi|}\prod_{B\in\pi}\delta_B.
\end{equation}
Thus it suffices that the displayed products tend to zero.
In particular conditional mean errors that are bounded at their
natural scale and actually tend to zero make every partition
containing a singleton negligible, provided the other conditional
errors remain bounded at their natural scales. A conditional two-point error
$o(x^{-2}N_iN_j)$ controls its pair-pair contribution to order
four, provided the second pair error stays bounded on its
natural scale. The one-block fourth-order error must itself be
$o(x^{-4}\prod_iN_i)$; averaging cannot improve it.

\subsubsection{The estimates still required from the interacting model}

The comparison just proved removes an unnecessary demand for a
uniform conditional clean-probability floor. It gives a rigorous
way to spend the positive history reserve in connected
correlations. Its remaining inputs are concrete.

First, the history sum and its class identifications must be
equalities for the zero-source physical measure, with their
pushforward densities retained as in
\cref{prop:phc-class-density}. Second, the complete conditional
exponent and the physical readout must be identified on the
admitted clean chart. Third, the marked normalized mass must
give the $d_S$ bounds. Fourth, the conditional cumulants,
including their means and their source contact terms, must
obey \eqref{phc:conditional-reference-errors} against one fixed
reference. These are the places where the background inverse,
formation and interacting force estimates enter.

For $n=2$, the reference coefficient is the free-curvature
two-point coefficient, and the errors in
\eqref{phc:two-reference-error} must be small relative to its
positive lower bound for the permitted shapes. For $n=4$, the
reference is the connected four-point kernel and the errors are
those in \eqref{phc:four-reference-error}. The free kernel and
its sign calculation do not prove the conditional interacting
estimate. The running-coupling conversion is applied only after
the comparison at scale $x^{-n}$ has been proved.

These observations also specify a useful order for closing the
short-distance argument. The conditional one-block errors and
the marked normalization come first. Quantitative shape and
grid errors are then compared with the established leading
coefficient. Exceptional exteriors and lower conditional
cumulants are controlled by the explicit products above.
No additional conditional clean-mass theorem is needed for
this averaged route.

% END PROOF EXPANSION positive_history_comparison.tex

% BEGIN PROOF EXPANSION effective_source_jets.tex
\subsection{Effective source derivatives and separated contact terms}
\label{ch:effective-source-jets}

Integrating a subset of the fields turns an affine physical source into
a nonlinear effective source action. Its mixed derivatives contribute
to connected correlations even at zero source. This section derives
the exact finite-order formula and estimates the mixed terms from
their spatial carriers. It explains when the covariance of the first
source responses suffices, and when additional contact insertions
must be retained.

\subsubsection{A fixed integration domain and its source derivatives}

Let $\mu$ be a probability on a measurable space. Suppose $A(z,y)$
is holomorphic in $z=(z_1,\ldots,z_n)$ near zero, measurable in $y$,
and $A(0,y)=0$. Assume that $e^{A(z,y)}$ and the source derivatives
used below have a common integrable bound on a smaller polydisc.
These hypotheses justify differentiation under the integral. Put
\begin{equation}\label{esj:generating-function}
 F(z)=\mu(e^{A(z,\cdot)}),\qquad W(z)=\log F(z),\quad W(0)=0.
\end{equation}
Since $F(0)=1$, the logarithm exists on a neighborhood of zero.
For nonempty $B\subset\{1,\ldots,n\}$ define the distinct-index jet
\begin{equation}\label{esj:source-jet}
 A_B(y)=\left.\prod_{i\in B}\partial_{z_i}A(z,y)\right|_{z=0}.
\end{equation}
The integral in \eqref{esj:generating-function} is on a fixed
measurable domain and with a fixed zero-source measure. A
source-dependent change of variables must first be written on
such a domain with its complete Jacobian. Differentiating a
sharp indicator of a moving domain is not justified by
holomorphy of the action alone.

\begin{proposition}[The partition formula for effective source jets]
\label{prop:esj-partition-formula}
Under these assumptions,
\begin{equation}\label{esj:partition-formula}
 \partial_{z_1}\cdots\partial_{z_n}W(0)
   =\sum_{\pi\in\Pi_n}\kappa_\mu(A_B:B\in\pi).
\end{equation}
An outer cumulant with one entry is its expectation.
\end{proposition}
\begin{proof}
For the coefficient with each index appearing once, work in
the commutative algebra in which $z_i^2=0$. Every power
series then truncates at degree $n$, so no interchange of
infinite series is required. In this algebra
$A(z,y)=\sum_{\varnothing\ne B}z_BA_B(y)$, where
$z_B=\prod_{i\in B}z_i$. The cumulant expansion of
$\log\mu(e^A)$ is
$\sum_{k=1}^n\kappa_\mu(A,\ldots,A)/k!$.
To obtain $z_1\cdots z_n$, the $k$ entries must use
disjoint nonempty index sets whose union is all indices.
These are the ordered partitions into $k$ blocks. The
factor $1/k!$ removes their ordering, proving the formula.
The dominated analytic assumptions identify this finite
coefficient with the indicated derivative of $W$.
\end{proof}

At order two,
\begin{equation}\label{esj:second-derivative}
 W_{12}(0)=\mu(A_{12})+\operatorname{Cov}_\mu(A_1,A_2).
\end{equation}
At order three,
\begin{align}
 W_{123}(0)={}&\mu(A_{123})
      +\operatorname{Cov}_\mu(A_{12},A_3)
      +\operatorname{Cov}_\mu(A_{13},A_2)\notag\\
 &+\operatorname{Cov}_\mu(A_{23},A_1)
      +\kappa_\mu(A_1,A_2,A_3).
 \label{esj:third-derivative}
\end{align}
At order four,
\begin{align}
 W_{1234}(0)={}&\mu(A_{1234})
      +\sum_i\operatorname{Cov}_\mu(A_i,A_{\{1,2,3,4\}\setminus\{i\}})
       \notag\\
 &+\sum_{\{ij\}\{k\ell\}}\operatorname{Cov}_\mu(A_{ij},A_{k\ell})
       \notag\\
 &+\sum_{\{ij\},\{k,\ell\}=\{1,2,3,4\}\setminus\{i,j\}}
                     \kappa_\mu(A_{ij},A_k,A_\ell)
       \notag\\
 &+\kappa_\mu(A_1,A_2,A_3,A_4).
 \label{esj:fourth-derivative}
\end{align}
The sums have respectively four, three and six terms.
Thus controlling only $A_i$ does not control either
$\mu(A_{12})$ or the additional fourth-order terms.

\subsubsection{A quantitative bound for all contact contributions}

Let $M_i>0$ and suppose $\|A_i\|_\infty\le M_i$.
For each block $B$ of size at least two suppose
\begin{equation}\label{esj:mixed-jet-bound}
 \|A_B\|_\infty\le\eta_B\prod_{i\in B}M_i.
\end{equation}
Using the cumulant constants $a_k$ defined in
\cref{phc:outer-cumulant-constant}, the partition formula gives
\begin{equation}\label{esj:contact-error}
 |W_{1\cdots n}(0)-\kappa_\mu(A_1,\ldots,A_n)|
 \le\prod_iM_i
     \sum_{\substack{\pi\in\Pi_n\\\pi\ne\{\{1\},\ldots,\{n\}\}}}
                 a_{|\pi|}\prod_{\substack{B\in\pi\\|B|\ge2}}\eta_B.
\end{equation}
Every term in the sum contains at least one mixed block.
The proof is to apply
$|\kappa(Y_1,\ldots,Y_k)|\le a_k\prod\|Y_j\|_\infty$
to each term of \eqref{esj:partition-formula}. It applies
also to complex-valued jets, since the moment polynomial
and its absolute-value bound are algebraic.

For example, the second-order error is at most
$\eta_{12}M_1M_2$. At order four it is at most
\begin{align}
 \prod_iM_i\bigg[&\eta_{1234}
      +2\sum_i\eta_{\{1,2,3,4\}\setminus\{i\}}
      +2\sum_{\{ij\}\{k\ell\}}\eta_{ij}\eta_{k\ell}
       \notag\\
      &+6\sum_{\{i,j\}\subset\{1,2,3,4\}}\eta_{ij}\bigg].
 \label{esj:fourth-contact-error}
\end{align}
In this formula the third-order jets occur paired with
one first-order jet. There is no isolated third-order
coefficient missing from the estimate.

If $A_i=Q_i+I_i$ with $\|Q_i\|_\infty\le M_i$ and
$\|I_i\|_\infty\le\rho_iM_i$, multilinearity gives
\begin{equation}\label{esj:first-jet-replacement}
 |\kappa(A_1,\ldots,A_n)-\kappa(Q_1,\ldots,Q_n)|
 \le a_n\left[\prod_i(1+\rho_i)-1\right]\prod_iM_i.
\end{equation}
This follows by expanding in the nonempty subset of entries
where $I_i$ is used. Combining
\cref{esj:contact-error,esj:first-jet-replacement} accounts
separately for the first response error and for every
mixed response. The latter cannot be absorbed into the
former without a bound on the mixed jets.

\subsubsection{Local carriers and a common source polydisc}

Let $\mathcal D$ be a family of finite domains in a graph
with cells $v$. Let $d(X)\ge0$ be a connected-size function.
For each $X$ suppose $H_X(z,y)$ is holomorphic on the
same closed source polydisc $|z_i|\le r_i$, measurable
in $y$ on one admitted field domain, and obeys
\begin{equation}\label{esj:local-holomorphic-bound}
 \sup_{z,y}|H_X(z,y)|\le h_X.
\end{equation}
Assume $H_X$ is independent of $z_i$ when its physical
carrier $X^{\sim}$ misses $\operatorname{supp}f_i$.
All source parameters here are held independently; a
single bound for a source in one direction would not
justify this mixed Cauchy estimate.

Suppose $H=\sum_XH_X$ converges normally on smaller
polydiscs. For a nonempty $B$ its mixed jet is
\begin{equation}\label{esj:local-jet-sum}
 H_B(y)=\sum_{X:\,X^{\sim}\cap\operatorname{supp}f_i\ne\varnothing
                          \text{ for all }i\in B}
                       \partial_BH_X(0,y).
\end{equation}
Multivariate Cauchy estimates give
\begin{equation}\label{esj:cauchy-local-jets}
 \|H_B\|_\infty
    \le\left(\prod_{i\in B}r_i^{-1}\right)
       \sum_{X\text{ meeting every support in }B}h_X.
\end{equation}
There is no factorial for distinct variables. Repeated
directions can be introduced as distinct source variables;
the usual factorial for a single-variable derivative
then follows from coefficient multiplicities.

The preceding statement requires a common field domain
for all $H_X$ after the actual substitution. If the norm
in \eqref{esj:local-holomorphic-bound} is established only
before substitution, the substituted family must be
shown to stay inside that domain. Cauchy's formula
does not repair a failure of domain admission.

\subsubsection{Exponential suppression of separated mixed jets}

Assume that a nonnegative local weight has the incidence bound
\begin{equation}\label{esj:incidence-bound}
 \sup_v\sum_{X\ni v}e^{\kappa d(X)}h_X\le E_0.
\end{equation}
Suppose every carrier that meets a selected support has
an assigned root in a finite set $W$, and that assigning
a root never reduces the estimate on the right: more
explicitly, the sum over such domains is bounded by
$\sum_{v\in W}\sum_{X\ni v}$ with the same $h_X$ and $d(X)$.
Assume also that if $X^{\sim}$ meets two of the separated
supports, then $d(X)\ge D_*$. These are the geometric
root and carrier properties needed in the estimate.
Then
\begin{equation}\label{esj:carrier-sum}
 \sum_{X\text{ meeting every support in }B}h_X
       \le |W|E_0e^{-\kappa D_*},\qquad |B|\ge2.
\end{equation}
Indeed factor $e^{-\kappa d(X)}$ out of the weighted
incidence norm and use $d(X)\ge D_*$ in each contributing
term. The root count is paid once by $|W|$.
This proves
\begin{equation}\label{esj:separated-jet-bound}
 \|H_B\|_\infty\le |W|E_0e^{-\kappa D_*}
                             \prod_{i\in B}r_i^{-1},
 \qquad |B|\ge2.
\end{equation}

A localization carrier may be an enlargement of $X$.
If its enlargement has radius $r_{\rm loc}$ and
the supports are separated by $d_*$, then the size
lower bound must use $d_*-2r_{\rm loc}$ in the
appropriate units. In particular it is not enough to
say that the original domains are disjoint. The
separation has to survive the actual enlargement used
in the source dependence.

If there are several origin levels, apply
\eqref{esj:separated-jet-bound} with their own incidence
bound, root count, carrier units and source radii,
and sum those numerical estimates. A coarse root count
cannot be used for fine-origin terms without the
filing estimate that transfers their incidence weights.
This is the same bookkeeping requirement as in
\cref{fwg:activity-oscillation}.

\subsubsection{The short-distance contact budget}

Let $x\to\infty$ and $N_i=\|f_i\|_\infty>0$.
Suppose the first responses satisfy
\begin{equation}\label{esj:first-response-scale}
 \|A_i\|_\infty\le Cx^{-1}(\log x)^P N_i,
\end{equation}
and all mixed jets of the regular effective source
arise from families satisfying
\eqref{esj:separated-jet-bound}. Suppose their summed
numerical bound is
\begin{equation}\label{esj:contact-numerical-bound}
 \|A_B\|_\infty\le
 C_n(\log x)^Qe^{-c(\log x)^b}
                         \prod_{i\in B}N_i,
 \qquad |B|\ge2,
\end{equation}
with $b>1$. This follows, for example, when $D_*\ge
c'(\log x)^b$, the root and norm factors are
polylogarithmic, and $r_i^{-1}\le C_nN_i$.

For a partition in \eqref{esj:partition-formula} with
$s$ singleton blocks and at least one mixed block,
its cumulant is bounded by
\begin{equation}\label{esj:partition-contact-scale}
 C_nx^{-s}(\log x)^{Q'}e^{-c'(\log x)^b}
                                      \prod_iN_i.
\end{equation}
For fixed $n\le4$ this is
$o(x^{-n}\prod_iN_i)$. In fact after division by
$x^{-n}\prod_iN_i$, the exponent has the form
$(n-s)\log x-c'(\log x)^b+Q'\log\log x$, which tends
to $-\infty$ because $b>1$. Summing the finitely many
partitions proves
\begin{equation}\label{esj:negligible-contacts}
 W_{1\cdots n}(0)
     =\kappa_\mu(A_1,\ldots,A_n)
          +o\!\left(x^{-n}\prod_iN_i\right),\qquad n=2,4.
\end{equation}
This establishes the regular separated contact estimate
with the actual source radius loss retained.

The condition $b>1$ in this argument is stronger than
the $b>0$ sufficient for rare-history errors in
\cref{phc:negligible-rare-error}. Here the mixed jets
were not assumed to have the natural factor
$x^{-|B|}$, so the exponential must overcome a fixed
positive power of $x$. If the geometric decay is only
$e^{-c\log x}=x^{-c}$, its coefficient must exceed
the missing source powers. If the mixed jets already
carry their full natural powers, $b>0$ suffices.
These are different uses of the same exponential
reserve and should not be conflated.

\subsubsection{Normalizer derivatives and the order of averaging}

The formulas also fix the effect of a chosen analytic
normalizer. If $\widetilde F(z)=e^{-E(z)}F(z)$ with
deterministic $E$, then
\begin{equation}\label{esj:normalizer-derivative}
 \partial_B\log F(0)=\partial_BE(0)
                              +\partial_B\log\widetilde F(0).
\end{equation}
A bound on $\widetilde F$ alone does not bound the
left side until the required mixed derivatives of
$E$ are estimated. If $E$ is a sum of terms each
depending on only one source variable, its mixed
distinct-index derivatives vanish for $|B|\ge2$.
For local terms involving several source variables,
the carrier estimate above supplies the needed bound.

For a conditional normalizer $E_e(z)$, the identity
must first be kept inside the exterior average:
\begin{equation}\label{esj:conditional-normalizer}
 F(z)=\int e^{E_e(z)}\widetilde F_e(z)\,\lambda(de).
\end{equation}
Replacing this by $\int\widetilde F_e(z)\lambda(de)$
would remove all the source derivatives of $E_e$ and
their exterior fluctuations. Apply the partition
formula to the full conditional logarithm, and then
the law of total cumulance of
\cref{ch:positive-history-comparison}. This order
retains both the source contact terms and the
conditional means.

At an actual clean run the remaining verification is
therefore precise. The original positive physical
source integral must equal the fixed-domain effective
integral. Its first responses must have the stated
readout and irrelevant remainder. Every mixed response
must retain its own carrier and source radius.
The marked history reserve controls the exceptional
contributions. Under these inputs,
\eqref{esj:negligible-contacts} and the averaged
comparison prove the connected extraction step
without discarding any differentiated normalizer.

\subsubsection*{Physical source derivatives at the four-support witness}

Use the sourced density of
\cref{sec:sources-the-curvature-source-with-fixed-reference-data}
at level $k_0=K-s$. The fixed-domain representation in the
two-point strategy applies with $w=\sum_{i=1}^4 z_i f_i$.
Use the exponent of each resummed standing class, including
the exponents of its integral operations. The measures of
the classes carry the unsourced gauge, averaging and test
factors on the joint space with stored fields. In particular,
an averaging delta can make the conditional measure in the
block field singular at fixed stored fields. Exactness of
both operations for the sourced total gives
$Z(w)/Z(0)=\mathbb E e^{\Phi(w)}$, with $\Phi$ the complete
exponent minus its zero-source value. If the representation
extracts deterministic local source energies as a separate
factor, include them in $\Phi$ in this identity.

The contact of a live component with a support and the age
of its lineage are already functions of the standing class.
For an event involving an individual history and the block
field together, let $\tau_h^0$ and $\tau_{[h]}^0$ denote
their zero-source block-field marginals before normalization.
Multiply the class joint measure by
$d\tau_h^0/d\tau_{[h]}^0(V)$ on the sheet for $h\in[h]$.
These weights are independent of $w$ and sum to one almost
everywhere. The lift gives the required individual zero-source
joint marginal and leaves the sum of the sourced integrals
unchanged. The exponent on a sheet is still the exponent of
the resummed class.

\paragraph{Mixed derivatives on island carriers.}
Separate the differences of regular activities, remainders
and boundary terms with no islands, together with the
label's pointed shares of the source energies. The boundary
terms involving islands form a second sum. In the remaining
sum keep the coupling pieces and the exponents of the integral
operations. Each nested operation exponent enters once, in
the operation that contains it. The carrying Wilson and
boundary terms enter in their own sums when a drop-out has
no exponent for them in its integral. All differences use
the same stored tuple at source $w$ and at zero.

Write $\rho_{\mathsf d}=\mathsf d L^s$ for the support
separation in witness lattice units. At origin level $j\le k_0$
the separation in localization cubes is
$n_j=(\rho_{\mathsf d}/M)L^{k_0-j}$. The rooms ensure
$n_j\ge C\check R_j$ and give source radii
$c\varrho/\|f_i\|_\infty$ at the witness. At fixed fields, the
running shift and each counterterm weight are sums of functions
of single plaquette source values, with fixed profiles and tents.
By linearity of the Wilson action in the weight, this also applies
to the separate carrying Wilson terms of a drop-out. Their mixed
derivatives vanish at disjoint supports regardless of the size
of the carrying region. An elementary operation whose source dependence
lies in its component of diameter $C\check R_j$ misses at
least one of any two of the separated supports. Apply this
criterion to the actual carrier of each sourced operation
exponent, including any subtraction from its normalizer.

For a formed operation the fixed-variable source formula
uses the twice-enlarged component, with the already formed
inner exponents included in the same read. At level $j$
the component at formation is in a cube of side
$100M\check R_j$ in its own units. Its enlarged physical
diameter is at most $CM\check R_j a_{K-j}$, below the
support separation by the witness room. Thus the mixed
derivative of this complete exponent vanishes. This applies
also to the executed operations retained inside a drop-out.
The Wilson and boundary factors never integrated by those
operations keep their own carriers. The calculation uses
the exponent at fixed stored fields and zero-source
reference measure throughout. In particular it involves
no division by a complex sourced operation mass.

For the regular activities and local energies, use the
enlarged carrier in \eqref{esj:carrier-sum}. An island boundary
term uses the tree length through the cells outside the
exempt hulls, with the tree also meeting each selected hull
as in \cref{lem:ccf-hull-root-entropy}. By the hull geometry at
\eqref{ccf:hull-coarsening}, every source contact in an enlarged
carrier is within $C(1+\check R_j)$ cells of a vertex of this
tree. The two contacts of an active carrier
therefore give length at least $n_j-C-C\check R_j$. Edges
inside a hull still count in the original lattice metric.
For the root pin the tree at a vertex near the first
contact. There are at most $Cn_j^4$ such vertices. The
rooted walk and hull incidence count in that lemma apply
at fixed history, before taking the essential supremum.
The exempt hulls are those in the surviving filing at the witness.
For a term whose old bound discounted an absorbed hull, apply the
donor refiling after \eqref{ccf:boundary-majorants} on the
admitted source box before using the incidence estimate.

For the boundary differences use the sourced bound on the retained
list at its actual field values. The witness room and monotonicity
of the depth norms give \(\|w\|_{\flat,K-j}\le\varrho/2\) for every
origin on \(|z_i|\le\varrho/(8\|f_i\|_\infty)\), after choosing
the fixed source radius \(\varrho\) within the sourced class.
Cauchy differentiation uses this common disc in the sources at
fixed stored fields. If an estimate holds almost everywhere for
each source separately, impose it first on a countable dense set
and extend it there by holomorphy of the actual read.
With the exponential incidence bound for the actual island
and regular terms, Cauchy's estimate thus costs
$Cn_j^4 e^{-c n_j}\prod_{i\in B}(\|f_i\|_\infty/\varrho)$
for a mixed derivative in $B$, $|B|\ge2$. Sum the levels
using $n_j=n_{k_0}L^{k_0-j}$ to obtain
\begin{equation}
 \|\Phi_B\|_\infty\le C(\rho_{\mathsf d}/M)^4
       e^{-c\rho_{\mathsf d}/M}\prod_{i\in B}
                             (\|f_i\|_\infty/\varrho).
 \label{eq:f3-relative-island-jets}
\end{equation}
This estimate uses the source localization on the resummed
standing classes, with the relative incidence bound on the
same source polydisc. For an operation that has a larger
source carrier one must include it in the mixed-jet estimate
with its actual size.

\paragraph{First derivatives on live histories.}
Let $B_i$ be the event that a live component at the witness
meets the fixed neighbourhood of support $i$. On this event
let $a_i$ be the greatest lineage age among those components.
Each of its $a_i+1$ stored levels can increase the spatial
count in the first source derivative. Write its upper weight
as $y_i=\mathbf 1_{B_i}(a_i+2)^C L^{4(a_i+1)}$. A renewal
of the label does not reset the count of retained levels.
For a nonempty set of indices $I$,
$\prod_{i\in I}y_i\le\sum_{i\in I}y_i^4$, since every
nonzero $y_i$ is at least one. On $B_i$ an attaining component
meets one of at most $C\rho_{\mathsf d}^4$ witness cells.
Summing over these cells and then the ages gives
\begin{equation}
 \mathbb E\prod_{i\in I}y_i
    \le C\rho_{\mathsf d}^4 Q_{17}(k_0),\qquad
 Q_b(k_0)=\sum_{a\ge0}L^{ba}
    \sup_{\mathfrak c}\mathbb P(\text{lineage age }a
                           \text{ at }\mathfrak c).
 \label{eq:f3-live-jet-moments}
\end{equation}
The four powers cost $L^{16a}$ and the remaining power
pays for the polynomial in the age. This is a normalized
history estimate, with the volume and depth range of the
age bound in force.

For a first derivative keep the witness, the class, the guards,
the background and the stored tuple fixed and restrict the
source to $w=t f_i$. This is the one-direction derivative of
the very same exponent used in \eqref{eq:f3-relative-island-jets}.
The source terms on the regular domain depend on the function
$w$, so restricting the other three parameters to zero changes
no integration domain or zero-source law. The derivative of
the full source energy of the empty standing label is a
deterministic constant $c_i$. For another standing label
\(\mathcal S\) the difference is the first derivative of
\(\sum_j[\mathsf T_j(\mathcal S;w)-\mathsf T_j(\mathcal S_\emptyset;w)]\),
with \(\mathsf T_j\) the label's share of the source energy.
Only a missing term near \(Z_j^+\) whose enlarged carrier meets
the source can contribute. This includes a missing one-point
gauge term if it is included in the source energy. The small-field
regions are nested in the reduced standing label, so the earlier
missing sites lie in the witness live footprint. The padding
at level \(j\) has physical width at most
\(CM\check R_j a_{K-j}\le CM\check R_{k_0}a_s\) by
\eqref{eq:rg-objects-R19}. Off \(B_i\) the padded missing sites
are therefore separated from the source by a fixed fraction of
\(\mathsf d\). A contributing carrier has
\(d_j(X)\ge c\rho_{\mathsf d}L^{k_0-j}/M\).
Rooting the sum near the source, the weighted carrier count
and Cauchy's estimate at the fixed source radius give a bound
\(C\|f_i\|_\infty\sum_{j\le k_0} m_j^4e^{-cm_j}\), where
\(m_j=\rho_{\mathsf d}L^{k_0-j}/M\). This is at most
\(C\|f_i\|_\infty\rho_{\mathsf d}^4e^{-c'\rho_{\mathsf d}/M}\).
On \(B_i\) the missing share belongs to the live error.

Off contact, the counterterm identity for the running weight
gives the profile of the insertion,
\[
 \left.\partial_{z_i}\zeta_{k_0}(p;w)\right|_0
   = f_i(x(p))+\sum_{j\le k_0} b_j'(0)
             \sum_y h_y^{(j)}(p)f_i(x(p_y))
   = \zeta'_{k_0}\widetilde f_i(p).
\]
Here $b'_j(0)$ differentiates the whole running increment
with respect to the input source. At constant input the
tent sum is one and the coefficient is $\zeta'_{k_0}$.
The tents have range at most $CMa_s$, so the bounded sum
of $|b'_j(0)|$ and the lower bound on $|\zeta'_{k_0}|$
give $|\widetilde f_i-f_i|\le CMa_s\operatorname{Lip}(f_i)$.
Use $\widetilde f_i$ as the profile of the bounded readout
$G_i$ in the four-insertion comparison below. Its window
tests all cubes needed to read the enlarged support. Include
these padded carriers in the neighborhood defining $B_i$.
%
On the actual zero-source joint law, the witness
characteristic functions imply passage of the individual
cube tests off \(B_i\). This is a statement about the
same retained fields and the localized witness chart.

On the admitted complete regular carriers the ordinary
one-direction bounds \eqref{eq:local-formation-LB14}
and \eqref{eq:local-formation-LB17}--\eqref{eq:local-formation-LB18}
apply to this restriction pointwise. After subtraction of the
vacuum value and the tent action, the quadratic coefficient
is bounded by $C\delta_{k_0}/\alpha_{k_0}$, using the
formed source slope. At age $a=k_0-j$, the quadratic
estimate retains a geometric factor $L^{-\alpha_H a}$
after the fine-site count. The cubic term retains $L^{-2a}$,
up to a polynomial in $a$. Summing gives
$C\delta_{k_0}^3/\alpha_{k_0}$ per witness window and per
unit source norm. For an old ordinary remainder the
quadratic field factor pays the fine-site count and the
extra factor $x_j^{-\kappa_0/2}$ is summed with
$\alpha_j^{-2}\le Cx_j$ by the lower flow bound.
This gives the additional $x^{2-\kappa_0/2}$ below.
%
For the localized readout off \(B_i\), the four-layer enlargement
of the test cube lies in the ordinary region. The history background
restricted there is the local minimizer at its own averages by the
nested-minimizer property. In a common gauge compare it with the test
minimizer on this enlargement. The logarithmic difference of their
averages is zero on the common interior rows and is bounded by
\(C\varepsilon_{k_0}\) on the outer layer. The stored fields enter
through these actual boundary averages. The collar response gives a
bound \(C\varepsilon_{k_0}e^{-c\check R_{k_0}}\) for the change of the
local potential and its first lattice derivatives in the plaquette
read. In bare plaquette units, with \(\eta=\eta_{k_0}\), the
plaquette change is bounded by
\(C\varepsilon_{k_0}\eta^2e^{-c\check R_{k_0}}\). Both plaquette
defects have size at most \(C\varepsilon_{k_0}\eta^2\), so
\[
 \tfrac12\bigl|\|U_p-\mathbf1\|^2-\|U'_p-\mathbf1\|^2\bigr|
 \le C\varepsilon_{k_0}^2\eta^4e^{-c\check R_{k_0}}.
\]
The separate contribution from the distant live region uses the
same curvature scale with decay \(e^{-c\rho_{\mathsf d}}\).
Summing with the witness profile cancels \(\eta^4\) against
the bare plaquette count and gives the action bound below.
These bounds use complete ordinary charts. Long or
wrapping carriers that leave those charts use the far
bound, or the live estimate if they reach a local contact.

\paragraph{Matching the windowed insertions.}
%
%
Let \(J_i\) denote the main Wilson response in the
physical derivative on \(B_i^c\), before replacing its
background by the witness minimizer. Extend
\(J_i\mathbf 1_{B_i^c}\) by zero on \(B_i\). On the same
joint law put \(O_i=\zeta'_{k_0}G_i(V)\), using the
witness window even on \(B_i\). Write \(F_i=\|f_i\|_\infty\).
On the witness window the localized minimizer passes its
own plaquette test and there are at most
\(C\rho_{\mathsf d}^4\eta_{k_0}^{-4}\) plaquettes in the
enlarged support. Off that window \(G_i=0\). Thus
\(\|O_i\|_\infty\le C x^{-1}\rho_{\mathsf d}^4p_0^2F_i\)
on the entire joint law. Then
\[
 J_i\mathbf 1_{B_i^c}-O_i
   =(J_i-O_i)\mathbf 1_{B_i^c}-O_i\mathbf 1_{B_i}.
\]
The two-radius comparison above bounds the first term
on the actual off-contact window, uniformly in the
stored fields. It has size at most
\(C x^{-1}\rho_{\mathsf d}^4p_0^2F_i
 (e^{-c\check R_{k_0}}+e^{-c\rho_{\mathsf d}})\).
The second term is bounded by the windowed supremum
times \(\mathbf 1_{B_i}\). This term occurs because the
witness window can pass on a live history.
Expand the difference of the fourth cumulants by
multilinearity. A term containing a deleted slot
\(O_i\mathbf 1_{B_i}\) costs its \(L^1\) bound times
the suprema of the other slots, by the moment formula
for the cumulant. With \(B=\bigcup_i B_i\), the total
deletion cost is at most
\[
 C\,\mathbb P(B)\,x^{-4}\rho_{\mathsf d}^{16}p_0^8
                   \prod_i F_i .
\]
Here \(\mathbb P(B)\le C\rho_{\mathsf d}^4 Q_{17}(k_0)\).
The off-contact defect costs the same prefactor with
\(\mathbb P(B)\) replaced by
\(e^{-c\check R_{k_0}}+e^{-c\rho_{\mathsf d}}\).
Thus the replacement of the physical main entries
uses the actual joint law and the localized comparison
on \(B_i^c\). The age-weighted bound below is still needed
for the live parts of the physical derivatives.

The off-contact ordinary error has size at most
$C x^{-1}\rho_{\mathsf d}^4 F_i$ times
$(\log x)^{3\mathfrak p_0-q}+p_0^2x^{2-\kappa_0/2}$, where the
last term pays the ordinary remainders. The far error costs a polynomial
in $\rho_{\mathsf d}$ times $F_i e^{-c\rho_{\mathsf d}}$.
%
For the live estimate, differentiate the outermost standing
operation exponents. An operation can contribute only when its
twice-enlarged carrier meets the source. By the diameter bound
above and \eqref{eq:rg-objects-R19}, its entire stored stock is
then in the witness source neighborhood of diameter \(C\mathsf d\).
Its lineage is among those counted by \(a_i\), and the oldest
retained level is at least \(k_0-a_i\). This also covers the
additional waiting-window stock of a drop-out, since its predecessor
already belongs to the lineage at the start of the window.
At a fixed level \(j\)
the stocks of distinct outermost operations are disjoint. Nested
stocks enter their outer exponent once. A site can occur once
more through a drop-out, since intersecting standing marks have
disjoint active level windows. The total site count at that level,
with this multiplicity, is bounded by
\(C\rho_{\mathsf d}^4L^{4(k_0-j)}\). Apply the Cauchy bound for
the sourced excess exponent, which is bounded by a fixed constant
times its stored site count on the fixed source disc. The derivative
sum is at most
\[
 C\varrho^{-1}F_i\rho_{\mathsf d}^4
       \sum_{j=k_0-a_i}^{k_0}L^{4(k_0-j)}
 \le C\varrho^{-1}F_i\rho_{\mathsf d}^4 L^{4a_i}.
\]
For the separate carrying Wilson terms, a level-\(j\) site has
\(O(\eta_j^{-4})\) bare plaquettes, each with action bounded by
\(C\eta_j^4\) on the corresponding regular layer. At the bare
layer use boundedness of the plaquette action. This gives a
fixed cost per site. The boundary terms use their own rooted
carrier estimates. With the source footprint and nested storage
bounds above, the live error has bound $s_i y_i$ with
$s_i=C(\log x)^C\rho_{\mathsf d}^8 L^{C\check R_{k_0}}F_i$.
This includes the differentiated operation exponent. The
joint law throughout is the one used for the mixed jets, and
the law of the four $G_i$ is its block-field marginal.

The lineage estimate of \cref{prop:source-domains-4-14}
with \(b=17\) gives \(Q_{17}(k_0)\le Ce^{-cp_0}\) on its
admitted runs. The fixed-torus alternative uses the same
age exponent at the witness depth where its terminal bound
applies. In particular the \(L^4\) error is at most
\(C s_i\rho_{\mathsf d}e^{-cp_0/4}\). Here
\(\check R_{k_0}=o(p_0)\) and \(p_0/\log x\to\infty\), so it is
\(o(x^{-1})F_i\). In every
partition with at least two blocks subtract the deterministic
constant from each first jet, using invariance of the
cumulant under those translations. The one-block partition
has only the four-source mixed jet. H\"older's
inequality in the moment formula bounds each of the fourteen
partitions containing a mixed source derivative by the
product of its mixed-jet bounds and the $L^4$ bounds for
its first jets. Relative to $x^{-4}\prod_iF_i$ this costs
at most $x^4(\log x)^C e^{-c\rho_{\mathsf d}/M}$.
In the partition into singletons expand the centered entries.
Each cumulant term with a live factor costs at most
$C\prod_i s_i\,\rho_{\mathsf d}^4Q_{17}(k_0)$ by
\eqref{eq:f3-live-jet-moments} in the corresponding moments.
No independence between contacts is needed. The relative
cost is bounded by
$C x^4\rho_{\mathsf d}^{36}(\log x)^C L^{C\check R_{k_0}}e^{-cp_0}$.
In the prescribed room $\check R_{k_0}=o(p_0)$ and
$p_0$ is of order $(\log x)^{\mathfrak p_0}$ with
$\mathfrak p_0>1$, so this tends to zero.
An ordinary error replacing one of the four main terms
costs $C\rho_{\mathsf d}^{16}(\log x)^{9\mathfrak p_0-q}$
relative to the same scale, apart from the exponentially
small pieces. The ordinary remainder has an extra inverse
power of $x$ that absorbs every fixed logarithmic power.
In the prescribed room
$\rho_{\mathsf d}\asymp(\log x)^{2\mathfrak p_0}$ this
requires $q>41\mathfrak p_0$. The far error is negligible
by the same separation as in \eqref{eq:f3-relative-island-jets}.
This reduces the physical fourth derivative to the fourth
cumulant of the bounded readouts, with the four running
source factors retained.

% Comprehensive Part II, 155900--156457, the own-scale
% four-source bounds and charge moments.
\paragraph{Counting the near reads on a finite torus.}
In the native four-source decomposition there is a second
bound for the near first derivatives. It uses the cells
actually read by each slot. For source \(i\) let \(T_i\)
count the cells of each stratum
\(\Omega_n\setminus\Omega_{n+1}\), \(n<k_0\), in
its enlarged support, measured in units of level \(n\).
Include the internal strata of pocket strips with their
own units. The action costs a constant on each such
cell and on the level-\(k_0\) portion, whose volume is
\(O(\rho_{\mathsf d}^4)\). In counting localization
and volume terms, a stratum of level \(n\) measured
at level \(j\ge n\) contributes \(L^{-4(j-n)}\)
times its own count. The sum over reading levels
therefore costs \(C(\rho_{\mathsf d}^4+T_i)\).
For the regular source terms the fine-anchor factor
\(L^{4(n-j)}\) cancels the same volume power in
their decay, leaving the summable \(L^{-\alpha(n-j)}\).
The boundary, potential and remainder bounds are
summed on their admitted native carriers. Together
these give, on the \(i\)-th live-contact event,
\[
 |\lambda_i^{\rm near}|
 \le C\mathfrak p_\flat(\bar x_s)(A_i+2)^{c_\flat}
               (\rho_{\mathsf d}^4+T_i)u_i ,
\]
where \(A_i\) is the greatest age read by that source
and \(u_i\) its source norm in this decomposition.

For a live component \(Z\) let \(D(Z)\) sum
\(M^4\check R_j^5d'_j(Z_j\cap Z)\) over its history,
including the internal pocket sequences. Here \(d'_j\)
is the geometric size in the history weight.
Pack the cells at level \(j\) by a maximal family
whose block indices are at least four apart in the
max norm. The family has at most \(d'_j\) members,
and each covers at most \(7^4\) cells of the original
region. The pocket strip costs a further
factor at most \(2^4\), and the level conversion
costs at most \(L^4\). Consequently \(T_i\) is bounded
by a fixed multiple of the sum of the charges of the
live components meeting that read. Each surviving
level contributes at least one, so \(D(Z)\ge
\operatorname{age}(Z)+1\).

Sum the charge of the live component of each read
block, counting it again if the block occurs in
another window. Call this sum \(D_{\rm read}\).
It dominates \(A_i+1\) on a live contact and controls
every \(T_i\). A product of at most four near
derivatives, divided by their deterministic factors,
is bounded on its contact event by \(C D_{\rm read}^{p}\)
for \(p\ge4(1+c_\flat)\). If each window meets at most
\(N_{\rm read}\) blocks, the power-mean inequality
bounds its expectation by \(C(4N_{\rm read})^p\)
times the largest rooted \(p\)-th charge moment.

On the prolonged finite torus the dialled history
bound pays \(e^{-c_a n}\) for charge at least \(n\).
If the component heals at \(k\), restrict the
healed sum to charge at least \(n+k-k_0\) and use
the healed-block bound. This yields a geometric sum
over \(k>k_0\). A component still live at \(\hat K\)
has charge at least \(n+\hat K-k_0\). Use the
terminal modulus and partition-function lower bound
with the charge dial. The additional decay in
\(\hat K-k_0\) pays the terminal volume under the
source's depth condition
\(c_a(\hat K-k_0)\ge \mathsf B(g^\dagger)\hat N^4\),
or the below-horizon bound applies. In this range
the resulting rooted moment is bounded by
\[
 C e^{-\theta_q p_0(g_{k_0})}
           \sum_{n\ge1}(n+1)^p e^{-c_a n}.
\]
This gives the near-slot moment estimate where the
actual four-source members have these own-scale
bounds and the dialled mass estimates apply. A
volume-uniform terminal charge estimate would be
needed for the same argument outside that range.

% END PROOF EXPANSION effective_source_jets.tex


% BEGIN MODULE chapters/gaussian_cutoff_clean_law.tex
\subsection{Sharp cutoff boundaries in the classical Gaussian law}
\label{ch:gaussian-cutoff-clean-law}

The boundary term in the conditional integration-by-parts identity need
not be small merely because a Gaussian precision has a gap.  The event
that every bond is small can have a probability tending to zero with
the volume.  Dividing a whole-space tail estimate by that probability
would therefore give no uniform conditional estimate.  In the centered
Gaussian reference law the symmetry and convexity of all the other
cutoffs give a different argument: their conditional probability
decreases on every radial ray.  A radial boundary density can then be
compared with an inner radial integral in the same direction.  The
normalization cancels before directions or bonds are summed.

This section proves that comparison for possibly singular physical
bond covariances.  It derives the resulting second through eighth
moment errors, including the covariance and third and fourth connected
errors of localized quadratic curvature observables.  The actual flat
pivot chart and contour gauge supply a concrete fixed-step application.
The measure here is classical Gaussian.  Haar and coarea factors, old
activities, nonlinear action weights and conditional exterior forces
are additional terms in the complete conditional law of
\cref{ch:clean-conditional-law}.

\subsubsection*{Physical blocks, actual covariance entries and the cutoff law}

Let \(Y\) be a centered Gaussian vector in a finite-dimensional real
space, with covariance \(G\ge0\), and let \(Z=CY\).  The blocks
\(Z_b\) belong to \(\mathbb R^3\); repeated or linearly dependent
blocks are allowed.  All covariances in this section are the actual
matrix entries
\begin{equation}
 \Gamma_{bc}=\mathbb E[Z_bZ_c^{\mathsf T}]=C_bGC_c^{\mathsf T},
 \qquad \Gamma=CGC^{\mathsf T}.
 \label{gcb:physical-covariance}
\end{equation}
For every tested block choose \(b_b>0\), put
\begin{equation}
 \mathcal A=\bigcap_b\{|Z_b|<b_b\},\qquad
 \nu(P)={\mathbb E[P(Z)1_{\mathcal A}]\over
                       \mathbb P(\mathcal A)},
 \qquad d_b=\operatorname{rank}\Gamma_{bb}\le3,
 \label{gcb:sharp-law}
\end{equation}
and choose an upper marginal scale \(\sigma_b>0\) with
\(\Gamma_{bb}\le\sigma_b^2 I\).  There is no lower eigenvalue
assumption.  A sufficiently small neighborhood of zero in the latent
Gaussian support is contained in \(\mathcal A\), so its probability
is positive.  If \(d_b=0\), then \(Z_b=0\) almost surely and that
constraint has no boundary contribution.  All blocks occurring in a
polynomial to be bounded below are included among the tested blocks.

For a finite volume no lower bound on \(\mathbb P(\mathcal A)\)
uniform in the volume is assumed or needed.  The estimates concern
these normalized finite-volume laws.  They do not define conditioning
on a simultaneous infinite collection of cutoff events, whose
intersection could have probability zero.

\subsubsection*{Marginal log-concavity from one-dimensional transport}

We first prove the integration fact needed for the conditional cutoff
probability.  Suppose nonnegative integrable functions \(f,g,h\)
on \(\mathbb R\), positive in the interiors of their interval
supports, satisfy, for \(0<t<1\),
\begin{equation}
 h((1-t)x+ty)\ge f(x)^{1-t}g(y)^t.
 \label{gcb:transport-hypothesis}
\end{equation}
It is enough initially that the functions be continuous and the
positive densities be smooth on their supports, as will hold in the
application.  Write \(F=\int f\), \(G_1=\int g\), and assume both
are positive.  The increasing quantile map \(T\) from density
\(f/F\) to \(g/G_1\) is continuously differentiable in the interior
and satisfies
\(T'(x)=f(x)G_1/[F g(T(x))]\).  The increasing map
\(T_t(x)=(1-t)x+tT(x)\) has derivative at least \(T'(x)^t\), by
the weighted arithmetic-geometric mean inequality.  Changing variables
on its image and using \eqref{gcb:transport-hypothesis} gives
\begin{equation}
 \begin{aligned}
 \int h&\ge\int h(T_t(x))T_t'(x)\,dx\\
 &\ge\int f(x)^{1-t}g(T(x))^t T'(x)^t\,dx
       =F^{1-t}G_1^t.
 \end{aligned}
 \label{gcb:transport-integral}
\end{equation}
One first integrates over compact subintervals of the positive support
and then lets them exhaust it.  All terms are nonnegative.  This also
handles unbounded interval endpoints and densities vanishing at the
boundary.  If either integral is zero the conclusion is immediate.

The same assertion holds in \(\mathbb R^q\).  Here is the induction
needed to avoid importing a stronger integration theorem.  Split
\(x=(x',u)\), \(y=(y',v)\).  For fixed \(x',y'\), apply the
one-dimensional assertion in \(u,v\), with the slice of \(h\)
at \((1-t)x'+ty'\).  The resulting marginals obey the same
inequality in \(q-1\) dimensions.  Apply the induction hypothesis
to those marginals and use Tonelli.  Slices of zero integral cause no
problem.  In the application all positive supports are open convex
sets and all slices have these interval supports.  Products of smooth
compact cutoffs with a Gaussian have integrable dominating functions
for their derivatives on every compact set of external parameters.
Their marginals are therefore continuous, smooth where required,
and positive on the interiors of their convex supports.  These facts
justify each induction step directly in this class of functions.

In particular, if \(k(z,w)\) is log-concave jointly in \(z,w\),
and its \(w\) slices are of the preceding integrable class, then
\begin{equation}
 w_0(z)=\int k(z,w)\,dw
 \quad\hbox{is log-concave in }z.
 \label{gcb:marginal-log-concavity}
\end{equation}
Indeed use the joint log-concavity inequality for the slices at
\(z,z'\) and \((1-t)z+tz'\), then integrate by
\eqref{gcb:transport-integral} and its induction.  The function
\(k\) need not be integrable over the external variable \(z\).
If \(k(-z,-w)=k(z,w)\), its marginal is even.  An even log-concave
function satisfies \(w_0(rn)\le w_0(sn)\) whenever \(0\le s\le r\):
write \(sn\) as a convex combination of \(rn\) and \(-rn\).
The argument below uses this ray monotonicity and requires no
volume estimate for the complete cutoff probability.

\subsubsection*{Smooth compact approximants with an inner plateau}

For \(0<\delta\le1/2\), construct a radial cutoff as follows.  Put
\(\phi(t)=0\) for \(t\le0\), and on \(0<t<1\) set
\begin{equation}
 \begin{aligned}
 \phi''(t)&={e^{-1/t^2}\over(1-t)^4},\qquad
 \phi(t)=\int_0^t(t-s){e^{-1/s^2}\over(1-s)^4}\,ds,
 \\
 \psi_{b,\delta}(z)&=
 \begin{cases}
 e^{-\phi((|z|/b-1+\delta)/\delta)},&|z|<b,\\
 0,&|z|\ge b.
 \end{cases}
 \end{aligned}
 \label{gcb:smooth-cutoff}
\end{equation}
The function \(\phi\) is convex, nondecreasing, flat at zero and
diverges at least as a positive multiple of \((1-t)^{-2}\) near
one.  Consequently \(\psi_{b,\delta}\) is smooth on all of
\(\mathbb R^3\), including the sphere \(|z|=b\), is even and
log-concave, and equals one on \(|z|\le(1-\delta)b\).
The smoothness at zero follows because it is constant near zero;
the smoothness at the outer sphere follows from the exponential
decay of \(e^{-\phi}\) faster than every inverse power of
\(1-t\).  Convexity of \(z\mapsto\phi((|z|/b-1+\delta)/\delta)\)
follows from convexity of the norm and monotonic convexity of \(\phi\),
extended by \(+\infty\) at and beyond one.

As \(\delta\downarrow0\), these cutoffs increase pointwise to
\(1_{\{|z|<b\}}\).  To check monotonicity inside the ball, the
argument of \(\phi\) is \(1-(1-|z|/b)/\delta\), which decreases
as \(\delta\) decreases.  Define
\begin{equation}
 \chi_\delta(Z)=\prod_b\psi_{b_b,\delta}(Z_b),\qquad
 \nu_\delta(P)={\mathbb E[P(Z)\chi_\delta(Z)]\over
                          \mathbb E\chi_\delta(Z)}.
 \label{gcb:smooth-law}
\end{equation}
Every denominator is positive, and Gaussian dominated convergence
gives \(\nu_\delta(P)\to\nu(P)\) for every polynomial \(P\).
No independence of the tested coordinates enters this statement.

The plateau in \eqref{gcb:smooth-cutoff} is essential to the estimate:
the derivative is supported in the outer annulus
\([(1-\delta)b,b]\).  A smooth cutoff whose derivative remained
nonzero close to zero would not by itself yield the exponential
boundary tail established next.

\subsubsection*{Conditional Gaussian support and monotonicity on rays}

Fix a block \(b\) of positive rank and let
\(V_b=\operatorname{ran}\Gamma_{bb}\), of dimension \(d=d_b\).
Orthogonal decomposition of a whitened latent Gaussian gives, when
conditioning on \(Z_b=z\in V_b\),
\begin{equation}
 Z_a=\Gamma_{ab}\Gamma_{bb}^{+}z+D_aW\quad(a\ne b),
 \qquad
 w_{b,\delta}(z)=
 \mathbb E_W\prod_{a\ne b}
 \psi_{b_a,\delta}(\Gamma_{ab}\Gamma_{bb}^{+}z+D_aW).
 \label{gcb:conditional-weight}
\end{equation}
Here \(W\) is a standard Gaussian independent of \(Z_b\), possibly
in dimension zero, and \(\Gamma_{bb}^{+}\) is the inverse on
\(V_b\), extended by zero.  For completeness, write the latent
standard Gaussian as its orthogonal projection onto the row space
of the map defining \(Z_b\) and the perpendicular projection.
The first projection is determined by \(z\), and the second is
an independent standard Gaussian on its subspace.  This proves
\eqref{gcb:conditional-weight}, including singular full covariances.
There is no bound asserted on \(\Gamma_{ab}\Gamma_{bb}^{+}\).

The integrand defining \(w_{b,\delta}\), multiplied by the Gaussian
density of \(W\), is jointly log-concave in \((z,W)\) and unchanged
under \((z,W)\mapsto(-z,-W)\).  The formula in \eqref{gcb:marginal-log-concavity} therefore proves that
\(w_{b,\delta}\) is even and log-concave on \(V_b\).  Its value
decreases along every ray.  It is positive near zero, although it
may vanish elsewhere because of exact dependencies between blocks.
Such zeros strengthen the inequalities below and require no separate
boundary trace convention.

\subsubsection*{A normalized boundary flux with no lower marginal variance}

Define the normalized smooth flux
\begin{equation}
 \beta_{b,\delta}=
 {\mathbb E[|\nabla\psi_{b_b,\delta}(Z_b)|
                 \prod_{a\ne b}\psi_{b_a,\delta}(Z_a)]
       \over\mathbb E\chi_\delta(Z)}.
 \label{gcb:smooth-flux}
\end{equation}
For \(d_b=0\), set \(\beta_b^*=0\); the flux is zero because
the cutoff is constant near zero.  In positive rank the radial gradient
lies in \(V_b\), so the
ambient Euclidean norm in \eqref{gcb:smooth-flux} is also its intrinsic
gradient norm.  Set \(B=b_b\), \(a=(1-\delta)B\), and use polar
coordinates \(z=rn\) in \(V_b\).  Up to a constant independent
of \(r,n\), the intrinsic Gaussian density is
\(\exp[-q(n)r^2/2]\), with
\(q(n)=n^{\mathsf T}(\Gamma_{bb}|_{V_b})^{-1}n
\ge\sigma_b^{-2}\).

For a fixed direction let
\begin{equation}
 \begin{aligned}
 D_n&=\int_0^B\psi_{B,\delta}(rn)w_{b,\delta}(rn)
                    r^{d-1}e^{-q(n)r^2/2}\,dr,\\
 I(q,a)&=\int_0^a s^{d-1}e^{-qs^2/2}\,ds.
 \end{aligned}
 \label{gcb:radial-denominator}
\end{equation}
For every \(r\in[a,B]\), ray monotonicity and the inner plateau
give \(D_n\ge w_{b,\delta}(rn)I(q(n),a)\).
The integral of the nonnegative radial derivative
\(-\partial_r\psi_{B,\delta}\) on \([a,B]\) is one.  Thus the
corresponding numerator in direction \(n\) is at most
\begin{equation}
 D_n\,{B^{d-1}e^{-q(n)a^2/2}\over I(q(n),a)}.
 \label{gcb:directional-flux}
\end{equation}
This inequality remains true if \(D_n=0\), because the preceding
pointwise inequality then forces the numerator to vanish.  In
particular it does not divide by a possibly zero directional weight.

There are constants depending only on \(d\le3\) such that
\(I(q,a)\ge c_d\min(a^d,q^{-d/2})\).  If
\(a\le q^{-1/2}\), integrate over \([0,a]\), where the exponential
is at least \(e^{-1/2}\).  Otherwise integrate over
\([0,q^{-1/2}]\).  Since \(B/2\le a\le B\), the ratio in
\eqref{gcb:directional-flux} is bounded by
\(C_d B^{-1}[1+(qB^2)^{d/2}]e^{-qB^2/8}\).
Taking the supremum over all \(q\ge\sigma_b^{-2}\) gives
\begin{equation}
 \boxed{\quad
 \beta_{b,\delta}\le\beta_b^*
 := {C_d\over b_b}
       \left[1+\left({b_b\over\sigma_b}\right)^d\right]
           e^{-c_d b_b^2/\sigma_b^2}\quad}
 \qquad(1\le d=d_b\le3).
 \label{gcb:flux-bound}
\end{equation}
To justify the last maximization explicitly, put
\(x=qB^2\), \(x_0=B^2/\sigma_b^2\).  The product
\((1+x^{d/2})e^{-x/16}\) is bounded on \([0,\infty)\).
For \(x\ge x_0\), the remaining factor is at most
\(e^{-x_0/16}\); increasing the constant gives the displayed
bound, with for example a smaller fixed \(c_d\).  Now integrate
\eqref{gcb:directional-flux} over directions.  The determinant
normalization of the intrinsic Gaussian and the total integral of
\(D_n\) cancel.  This proves \eqref{gcb:flux-bound} uniformly in
\(\delta\), the volume, the number of other constraints, their
dependencies and every lower marginal eigenvalue.

For \(b_b\ge\sigma_b\) the useful simpler consequence is
\begin{equation}
 \beta_b^*\le C_d b_b^{d-1}\sigma_b^{-d}
                  e^{-c_d b_b^2/\sigma_b^2}.
 \label{gcb:large-radius-flux}
\end{equation}
The qualification is necessary.  At small radius the normalized
radial flux can be of order \(d/b_b\); an estimate proportional to
\(b_b^{d-1}\sigma_b^{-d}\) alone would be false there.  The proof
has retained that regime in \eqref{gcb:flux-bound}.

\subsubsection*{Gaussian integration by parts before the sharp limit}

For a polynomial \(P\) in the ambient physical blocks, ordinary
Gaussian integration by parts, applied to a whitened latent vector,
gives
\begin{equation}
 \begin{aligned}
 \nu_\delta[Z_iP]&=
   \sum_j\Gamma_{ij}\nu_\delta[\nabla_jP]
       +\sum_b\Gamma_{ib}\mathcal L_{b,\delta}(P),\\
 \mathcal L_{b,\delta}(P)&=
 {\mathbb E[P(Z)\nabla\psi_{b_b,\delta}(Z_b)
                         \prod_{a\ne b}\psi_{b_a,\delta}(Z_a)]
        \over\mathbb E\chi_\delta(Z)}.
 \end{aligned}
 \label{gcb:smooth-ibp}
\end{equation}
The gradients are column vectors and \(\Gamma_{ib}\) are the
corresponding \(3\) by \(3\) blocks.  Degeneracy of \(\Gamma\)
does not alter the formula: in the latent variables its two
chain-rule factors are exactly the factors producing
\(CGC^{\mathsf T}\).  Gaussian domination justifies integration
by parts even when the cutoffs have no compact support in null
directions of \(C\).

Let \(\|P\|_{\mathcal B}\) be the supremum on the product of
closed balls \(\mathcal B=\prod_b\{|z_b|\le b_b\}\).  Then
\begin{equation}
 |\mathcal L_{b,\delta}(P)|
       \le\beta_b^*\|P\|_{\mathcal B}.
 \label{gcb:individual-defect-bound}
\end{equation}
For \(L=\sum_i\ell_i^{\mathsf T}Z_i\), define its canonical
sharp defect by
\begin{equation}
 \mathcal D_L(P)=\nu[LP]
       -\sum_{i,j}\ell_i^{\mathsf T}\Gamma_{ij}\nu[\nabla_jP].
 \quad
 |\mathcal D_L(P)|\le
 \left(\sum_{i,b}|\ell_i|\,\|\Gamma_{ib}\|\beta_b^*\right)
                   \|P\|_{\mathcal B}.
 \label{gcb:sharp-defect}
\end{equation}
Indeed the left side is the limit of the left side of
\eqref{gcb:smooth-ibp} contracted with \(\ell\), while
\eqref{gcb:individual-defect-bound} bounds the right side uniformly.
This proves the sharp identity and its estimate without assigning
a meaning to a product of a boundary delta function with another
sharp indicator on the same surface.

For repeated or dependent constraints the separate sharp limits of
\(\mathcal L_{b,\delta}\) need not be intrinsic: the allocation
between coincident constraint surfaces can depend on the smoothing
schedule.  The total defect in \eqref{gcb:sharp-defect} is intrinsic,
because its definition uses only sharp conditional polynomial
expectations.  This is the only boundary object needed for the moment
estimates.  When a classical smooth-boundary free-coordinate trace
exists, \(\nabla1_D=-n\,dS\), so this defect is the negative of
the covariance contracted with the outward boundary functional in
the conditional-law identity.  This also fixes the sign convention.

\subsubsection*{The actual weighted Schur accounting}

Choose positive amplitude scales \(a_b\) and positive cell weights
\(\mu_b\).  Assume the following explicit two-sided bound, using
the covariance entries of \eqref{gcb:physical-covariance}:
\begin{equation}
 S=\max\left\{
  \sup_i\sum_j\mu_j{\|\Gamma_{ij}\|\over a_i a_j},
  \sup_j\sum_i\mu_i{\|\Gamma_{ij}\|\over a_i a_j}
                         \right\}<\infty,
 \qquad
 \varepsilon_{\partial}
       =\sup_b{a_b\beta_b^*\over\mu_b}.
 \label{gcb:weighted-schur}
\end{equation}
The two sums give the usual bounded operator on weighted counting
\(\ell^2\) after normalizing by \(a_b\): weighted Cauchy in each
row and then the column bound give norm at most \(S\).
For the boundary estimate the more specific computation is
\begin{equation}
 {1\over a_i}\sum_b\|\Gamma_{ib}\|\beta_b^*
 =\sum_b\mu_b{\|\Gamma_{ib}\|\over a_i a_b}
                 {a_b\beta_b^*\over\mu_b}
 \le S\varepsilon_{\partial}.
 \label{gcb:weighted-boundary-row}
\end{equation}
Thus, with \(U(L)=\sum_i a_i|\ell_i|\),
\begin{equation}
 |\mathcal D_L(P)|\le
       \delta(L)\|P\|_{\mathcal B},\qquad
 \delta(L)=S\varepsilon_{\partial}U(L).
 \label{gcb:linear-defect-norm}
\end{equation}
There is no number of ambient coordinates in this estimate.
The factor \(\mu_b^{-1}\) in
\(\varepsilon_{\partial}\) cannot be dropped: coordinate
Gaussian integration by parts uses the ordinary partial derivative
with respect to that coordinate.  A weighted kernel convention
does not change that derivative.

For example, if a precision is written as a self-adjoint operator
\(K\) in pairing \(\sum_b\mu_b y_b\cdot z_b\), its coordinate
density has Euclidean precision \(M K\), where
\(M=\operatorname{diag}(\mu_b I_3)\).  Its covariance matrix is
\((MK)^{-1}\).  If
\((K^{-1}f)_i=\sum_j\mu_j\mathsf G(i,j)f_j\), then
\(K^{-1}=\mathsf G M\), and \((MK)^{-1}=\mathsf G\).
Accordingly \(\mathsf G(i,j)\), rather than
\(\mu_j\mathsf G(i,j)\), is the covariance entry used here.
This calculation prevents either suppressing or inserting a second
cell-volume factor during the physical coordinate conversion.

\subsubsection*{All linear marginals remain Gaussian dominated}

There is a useful moment bound independent of the cutoff radii.
For a real linear statistic \(L\), put
\(v(L)^2=\mathbb E L^2=\ell^{\mathsf T}\Gamma\ell\).
If this is zero then \(L=0\) almost surely.  Otherwise condition
the latent Gaussian on \(L=x\).  The argument of
\eqref{gcb:conditional-weight}, now including every cutoff, gives
an even log-concave weight \(w(x)\) for its one-dimensional
Gaussian density, first for \(\nu_\delta\) and then by the
sharp limit.  Therefore \(w\) decreases as \(|x|\) increases.
For any nonnegative increasing function \(f\) of \(|x|\),
\begin{equation}
 \nu[f(|L|)]\le
       \mathbb E f(|v(L)N|),\qquad N\sim N(0,1).
 \label{gcb:linear-moment-domination}
\end{equation}
To see the normalization explicitly, let \(R,R'\) be independent
absolute Gaussian variables of this variance.  The product
\((f(R)-f(R'))(w(R)-w(R'))\) is nonpositive.  Its expectation
is twice \(\mathbb E[fw]-\mathbb E[f]\mathbb E[w]\).
Divide by \(\mathbb E[w]>0\), and truncate \(f\) followed by
monotone convergence if necessary.  This proves the assertion
without a lower bound on the cutoff probability.

In particular the law is centered and, for \(m\ge1\),
\begin{equation}
 \nu[L^{2m}]\le(2m-1)!!\,v(L)^{2m},\qquad
 \nu[e^{tL}]\le e^{t^2v(L)^2/2},\qquad
 \operatorname{Cov}_\nu(Z)\le\Gamma.
 \label{gcb:linear-subgaussian}
\end{equation}
The exponential assertion follows by symmetry and the increasing
function \(\cosh(t|L|)\).  The covariance ordering follows by
testing every linear functional.  Every joint moment of an odd
number of linear statistics is zero, because the entire law is
invariant under \(Z\mapsto-Z\).

\subsubsection*{Covariance and fourth-order defects of linear statistics}

For linear statistics \(L_r\), write
\begin{equation}
 \begin{aligned}
 &\Gamma_{rs}^{L}=\mathbb E[L_rL_s],\quad
 \Sigma_{rs}=\nu[L_rL_s],\quad
 E_{rs}=\Sigma_{rs}-\Gamma_{rs}^{L},\\
 &B_r=\sum_i b_i|\ell_{r,i}|,\quad
 \delta_r=S\varepsilon_{\partial}U(L_r),\quad v_r=v(L_r).
 \end{aligned}
 \label{gcb:word-data}
\end{equation}
Here \(B_r\) bounds \(|L_r|\) on the entire cutoff box, whereas
\(v_r\) is its untruncated Gaussian standard deviation.  These
two quantities have different uses.  Taking \(P=L_s\) in
\eqref{gcb:sharp-defect} gives
\begin{equation}
 |E_{rs}|\le\min(\delta_r B_s,\delta_s B_r),\qquad
 |\Sigma_{rs}|+|\Gamma_{rs}^{L}|\le2v_rv_s.
 \label{gcb:covariance-error}
\end{equation}
The second estimate follows separately for its two terms from
\eqref{gcb:linear-subgaussian} and Gaussian Cauchy--Schwarz.
Every third joint moment of linear statistics is exactly zero.

Apply the same identity to \(P=L_2L_3L_4\).  The derivative
of this product produces exactly three pairings, and hence
\begin{equation}
 \nu[L_1L_2L_3L_4]
    =\sum_{j=2}^4\Gamma_{1j}^{L}\Sigma_{\{2,3,4\}\setminus\{j\}}
                    +\mathcal D_{L_1}(L_2L_3L_4).
 \label{gcb:fourth-recursion}
\end{equation}
In this formula \(\Sigma_{\{r,s\}}=\Sigma_{rs}\).  Subtracting
the three Gaussian pairings gives an error at most
\(\delta_1B_2B_3B_4+
\sum_{j=2}^4|\Gamma_{1j}^{L}||E_{\{2,3,4\}\setminus\{j\}}|\).
Subtracting instead the three pairings with the actual covariance
\(\Sigma\) gives the fourth connected moment and the explicit bound
\begin{equation}
 \begin{aligned}
 |\kappa_\nu(L_1,L_2,L_3,L_4)|
 \le\delta_1\big(&B_2B_3B_4+B_2v_3v_4\\
                         &+B_3v_2v_4+B_4v_2v_3\big).
 \end{aligned}
 \label{gcb:fourth-connected-error}
\end{equation}
Indeed replace \(\Gamma_{1j}^{L}-\Sigma_{1j}\) by
\(-E_{1j}\), use \(|E_{1j}|\le\delta_1 B_j\), and bound each
remaining \(\Sigma\) by its two Gaussian standard deviations.
Any of the four indices can be chosen as the first one.  The
identity proves the connected estimate without assuming a small
fourth cumulant.

For example,
\(\operatorname{Cov}_\nu(L_1L_2,L_3L_4)
=\kappa_\nu(L_1,L_2,L_3,L_4)+\Sigma_{13}\Sigma_{24}
+\Sigma_{14}\Sigma_{23}\).
Together with \eqref{gcb:covariance-error} this gives the covariance
error of two factored quadratic observables relative to their
untruncated Gaussian Wick contractions.  The higher degree estimates
below also control three and four such observables.

\subsubsection*{The moment recursion through degree eight}

Let a collection of localized linear statistics obey
\begin{equation}
 B_r\le R\,U_r,\qquad v_r\le v\,U_r,
 \qquad U_r=U(L_r),\qquad D=S\varepsilon_{\partial}.
 \label{gcb:local-radius-data}
\end{equation}
One may take \(R=\max_{i\in\mathcal S}b_i/a_i\) and
\(v=\max_{i\in\mathcal S}\sigma_i/a_i\), where \(\mathcal S\)
is the union of the supports of these statistics.  The variance
bound follows from
\(\|\Gamma_{ij}\|\le\sigma_i\sigma_j\), itself a consequence of
Gaussian Cauchy--Schwarz.  Thus neither \(R\) nor \(v\) requires
a maximum over coordinates outside their supports.

Let \(M_\nu(1,\ldots,n)=\nu[\prod_{r=1}^nL_r]\) and let
\(M_0\) denote the corresponding untruncated Gaussian moment.
For even \(n\), integration by parts gives the exact recurrence
\begin{equation}
 M_\nu(1,\ldots,n)-M_0(1,\ldots,n)
 =\sum_{j=2}^n\Gamma_{1j}^{L}
                  (M_\nu-M_0)(\{2,\ldots,n\}\setminus\{j\})
       +\mathcal D_{L_1}\Big(\prod_{r=2}^n L_r\Big).
 \label{gcb:general-moment-recursion}
\end{equation}
For \(n=2\) the empty moment difference is zero.  Odd moments
vanish for both laws.  Consequently
\begin{equation}
 \begin{aligned}
 |M_\nu(1,\ldots,2m)-M_0(1,\ldots,2m)|
             &\le d_{2m}\prod_{r=1}^{2m}U_r,\\
 d_2&=DR,\qquad d_n=DR^{n-1}+(n-1)v^2d_{n-2}.
 \end{aligned}
 \label{gcb:moment-error-recurrence}
\end{equation}
This is a direct induction: each derivative pairing contributes
at most \(v^2U_1U_j\), and the boundary defect contributes at
most \(D U_1\prod_{r=2}^n B_r\).  In particular
\begin{equation}
 \begin{aligned}
 d_2&=DR,\\
 d_4&=D(R^3+3v^2R),\\
 d_6&=D(R^5+5v^2R^3+15v^4R),\\
 d_8&=D(R^7+7v^2R^5+35v^4R^3+105v^6R).
 \end{aligned}
 \label{gcb:degree-eight-errors}
\end{equation}
All constants are independent of the total number of Gaussian
coordinates.  The factors \(U_r\) retain the actual localization
and scale of each inserted statistic.

Absolute mixed moments have a second bound useful when subtracting
products of moments.  H\"older with \(2m\) equal exponents and
\eqref{gcb:linear-moment-domination} gives, for either law,
\begin{equation}
 \left|\mathbb E_*\prod_{r=1}^{2m}L_r\right|
       \le(2m-1)!!\prod_{r=1}^{2m}v_r
       \le(2m-1)!!v^{2m}\prod_{r=1}^{2m}U_r,
 \qquad \mathbb E_*\in\{\nu,\mathbb E\}.
 \label{gcb:mixed-moment-bound}
\end{equation}
No independence of the \(L_r\) is used.  The simultaneous use
of the support bound in the defect and the variance bound in
unaffected moment factors avoids an unnecessary cutoff radius
in every term.

\subsubsection*{Localized quadratic observables and their connected errors}

Write each real quadratic observable in a finite factored form
\begin{equation}
 Q_j=\sum_\alpha c_{j,\alpha}
                      L_{j,\alpha,1}L_{j,\alpha,2},\qquad
 \mathcal A(Q_j)=\sum_\alpha |c_{j,\alpha}|
                        U(L_{j,\alpha,1})U(L_{j,\alpha,2}).
 \label{gcb:quadratic-coefficient-norm}
\end{equation}
Choose common \(R,v\) for this finite list as in
\eqref{gcb:local-radius-data}.  Every product of \(p\le4\)
quadratic observables is a sum of words of degree \(2p\).
\eqref{gcb:degree-eight-errors} and
\eqref{gcb:mixed-moment-bound} can therefore be summed with
absolute coefficients, with factor \(\prod_j\mathcal A(Q_j)\).
This step counts the inserted coefficients, not all coordinates
of the integration space.

For clarity define the joint cumulant by the finite partition formula
\begin{equation}
 \kappa_*(Q_1,\ldots,Q_p)=
 \sum_{\pi\in\mathfrak P_p}(-1)^{|\pi|-1}(|\pi|-1)!
                  \prod_{A\in\pi}\mathbb E_*\prod_{j\in A}Q_j.
 \label{gcb:quadratic-cumulant-definition}
\end{equation}
Write \(\Delta_p=\kappa_\nu(Q_1,\ldots,Q_p)
-\kappa_0(Q_1,\ldots,Q_p)\).  The following explicit estimates result:
\begin{equation}
 \begin{aligned}
 |\nu Q_1-\mathbb E Q_1|&\le d_2\mathcal A(Q_1),\\
 |\Delta_2|
  &\le(d_4+2v^2d_2)\mathcal A(Q_1)\mathcal A(Q_2),\\
 |\Delta_3|
  &\le(d_6+3v^2d_4+15v^4d_2)
                              \prod_{j=1}^3\mathcal A(Q_j),\\
 |\Delta_4|
  &\le(d_8+4v^2d_6+30v^4d_4+156v^6d_2)
                              \prod_{j=1}^4\mathcal A(Q_j).
 \end{aligned}
 \label{gcb:quadratic-connected-errors}
\end{equation}
Here the subscript zero always refers to the original untruncated
Gaussian with covariance \(\Gamma\).  In particular its quadratic
third and fourth cumulants are generally nonzero.  They are the
Gaussian Wick values, not quantities being set equal to zero.

We give the product bookkeeping in \eqref{gcb:quadratic-connected-errors}.
By multilinearity it is enough to use one factored monomial for each
\(Q_j\).  For covariance the full moment difference is bounded by
\(d_4\); replacing the product of its two means costs
\(2v^2d_2\).  At order three the three terms of partition type
\((2,1)\) cost at most
\(3(v^2d_4+3v^4d_2)\).  The product of three means, with
coefficient two, costs at most \(6v^4d_2\).  Together with
\(d_6\), these give the third line.  At order four, the four
partitions of type \((3,1)\) cost
\(4(v^2d_6+15v^6d_2)\); the three of type \((2,2)\) cost
\(18v^4d_4\).  The six of type \((2,1,1)\), with coefficient
two, cost \(12v^4d_4+72v^6d_2\).  The final product of four
means, with coefficient six, costs \(24v^6d_2\).
Their sum gives exactly the fourth line.  Each estimate uses the
telescoping difference of products, bounding unchanged factors by
\eqref{gcb:mixed-moment-bound}.  This establishes the constants
without an assumed connected-moment estimate.

Substituting \eqref{gcb:degree-eight-errors}, the three connected
error polynomials multiplying \(D\) are respectively
\begin{equation}
 \begin{aligned}
 \mathcal P_2(R,v)&=R^3+5v^2R,\\
 \mathcal P_3(R,v)&=R^5+8v^2R^3+39v^4R,\\
 \mathcal P_4(R,v)&=R^7+11v^2R^5+85v^4R^3+411v^6R.
 \end{aligned}
 \label{gcb:connected-error-polynomials}
\end{equation}
The same method covers mixed linear and quadratic observables:
expand each quadratic into its two factors, use
\eqref{gcb:general-moment-recursion} for each resulting word,
and then the finite partition formula.  An odd total linear degree
vanishes before taking absolute values.  No additional probabilistic
assumption is required for this mixed case.

The reference Gaussian Wick values themselves can be calculated
without a nonsingular physical covariance.  If
\(Q_j=Z^{\mathsf T}A_jZ\) with real symmetric \(A_j\), whitening
the latent Gaussian gives, for \(p\ge1\),
\begin{equation}
 \kappa_0(Q_1,\ldots,Q_p)
   ={2^{p-1}\over p}\sum_{\pi\in S_p}
       \operatorname{tr}(\Gamma A_{\pi(1)}\cdots
                                      \Gamma A_{\pi(p)}).
 \label{gcb:gaussian-quadratic-cycles}
\end{equation}
Indeed, near zero the logarithm of the Gaussian exponential
generating function is
\[
 \frac12\sum_{k\ge1}\frac{2^k}{k}
 \operatorname{tr}\left[\left(\Gamma\sum_jt_jA_j\right)^k\right].
\]
Taking the coefficient of \(t_1\cdots t_p\) gives
\eqref{gcb:gaussian-quadratic-cycles}.  This also identifies precisely
the contractions to which the sharp-cutoff cumulants are compared.

\subsubsection*{The fixed-step flat chart realizes the covariance hypothesis}

We now specialize to the actual first plain-average fibre at the
identity, with fixed odd block size \(L\), in current unit-spacing
coordinates.  Let \(P_L\) be the plain derivative of the nonlinear
average, \(\mathsf C_0\) its pivot inclusion with identity free
restriction, and \(\mathcal T\) the internal contour matrix.  The
definitions and the actual first-step precision are
\begin{equation}
 \begin{aligned}
 &\operatorname{ran}\mathsf C_0=\ker P_L,
   \qquad (\mathsf C_0u)|_{\mathrm{free}}=u,\\
 &\mathsf K_L=\mathsf C_0^{\mathsf T}
                  (\mathrm{d}_1^{\mathsf T}\mathrm{d}_1+
                         \mathcal T^{\mathsf T}\mathcal T)\mathsf C_0,
 \quad Y\sim N(0,\mathsf K_L^{-1}),\quad Z=\mathsf C_0Y.
 \end{aligned}
 \label{gcb:actual-flat-precision}
\end{equation}
The quadratic form here is exactly the first-step form, with no
previous finer-level minimization.  This is the classical reference with action weight
\(x=\theta\), so \(\tau=x/\theta=1\).  The canonical normalization
\(\theta g^2=1\) then gives coefficient one both to the Wilson curl
form and to the prescribed actual soft gauge.  In particular
\(\mathcal T\) is not a divergence gauge or a freely chosen
projector.

The pivot matrix has finite range with at most \(C_L\) nonzero
blocks in each row and column and entries bounded by \(C_L\).
The column assertion follows by the same double-block support as
the row assertion: a free bond belongs to only finitely many
double blocks when \(L\) is fixed.  The flat gap and the
exponential inverse estimate of
\cref{ch:flat-exponential-kernels} give
\(\|\mathsf K_L^{-1}(b,c)\|\le C_L e^{-\mu_Ld(b,c)}\), uniformly
on the admitted tori.  Consequently
\begin{equation}
 \Gamma=\mathsf C_0\mathsf K_L^{-1}\mathsf C_0^{\mathsf T},
 \quad
 \max\left(\sup_b\sum_c\|\Gamma_{bc}\|,
            \sup_c\sum_b\|\Gamma_{bc}\|\right)\le S_L,
 \quad \Gamma_{bb}\le\sigma_L^2 I.
 \label{gcb:realized-flat-schur}
\end{equation}
To check volume uniformity, a four-dimensional shell of distance
\(m\) has at most \(C_L(m+1)^3\) bond indices in block
coordinates.  Summing \(e^{-\mu_Lm}\) over these shells is
finite; periodic image sums have the same upper bound.  Multiplying
the finite row and column sums of \(\mathsf C_0\), its transpose
and the inverse proves the two sums in
\eqref{gcb:realized-flat-schur}.  The diagonal bound follows, for
example, from either sum.  Increase \(\sigma_L\) to at least
one if convenient.

Thus \eqref{gcb:weighted-schur} is realized with
\(a_b=\mu_b=1\), \(S=S_L\), independent of the number of blocks.
This is a statement at a fixed current spacing and fixed \(L\).
Pointwise decay in a multilevel physical metric would not alone
give its unweighted analogue: the number and normalization of
coordinates at different scales must still be charged in the
actual weighted norm in \eqref{gcb:weighted-schur}.

\subsubsection*{Why the actual contour gauge leaves no physical quadratic residue}

The following range calculation justifies integrating out this soft
gauge; it does not apply to an arbitrary soft gauge.  Let \(\mathcal S_0\)
be the real site fields vanishing at every block center.  For any
target internal contour data \(s=(s_{y,x})\), \(x\in B(y)\setminus
\{y\}\), define the site field \(R s\) by
\((Rs)(x)=s_{y,x}\) and \((Rs)(y)=0\).  Blocks partition the
sites, so this is well defined.  The plain-gradient identity and
the internal path telescoping identity give exactly
\begin{equation}
 \begin{aligned}
 &P_L\mathrm{d}_0Rs=0,\qquad \mathrm{d}_1\mathrm{d}_0Rs=0,\qquad
 (\mathcal T \mathrm{d}_0Rs)_{y,x}=s_{y,x},\\
 &\|Rs\|_2=\|s\|_2,\qquad \|\mathrm{d}_0Rs\|_2\le4\|s\|_2.
 \end{aligned}
 \label{gcb:actual-gauge-right-inverse}
\end{equation}
Here \(\mathrm{d}_0\) is the forward site gradient and \(\mathrm{d}_1\) the
plaquette circulation.  The last norm estimate follows by summing
\(|f(x+e_\mu)-f(x)|^2\le2|f(x+e_\mu)|^2+2|f(x)|^2\)
over four directions.  Internal contour means average paths, but
on a gradient every one of those paths gives the same difference
\(f(x)-f(y)\); hence no path count or \(L^{-4}\) factor occurs
in the third identity.  The constraint identity is
\(P_L\mathrm{d}_0f=\mathrm{d}_L(f|_{\mathrm{centers}})\), not a site-average
gradient identity.

In particular the image of \(\mathcal T\) on
\(\ker \mathrm{d}_1\cap\ker P_L\) is the entire internal target space.
It equals the image of \(\mathcal T\) on \(\ker P_L\).
This is the precise range condition that prevents a leftover
physical gauge penalty.  If an arbitrary alternative gauge failed
this condition, integrating its gauge directions could instead
leave a positive quadratic form in curvature.

Put \(\mathcal E_C=\mathrm{d}_1\ker P_L\).  The linear map
\begin{equation}
 \ker P_L\longrightarrow\mathcal E_C\oplus\mathcal S_0,
 \qquad a\longmapsto(\mathrm{d}_1a,\mathcal T a)
 \quad\hbox{is a bijection}.
 \label{gcb:curvature-gauge-coordinates}
\end{equation}
It is onto because any initial lift of a prescribed curvature in
\(\mathcal E_C\) can be adjusted to any prescribed contour data
by adding \(\mathrm{d}_0Rs\) from
\eqref{gcb:actual-gauge-right-inverse}.  It is one-to-one because
the flat gap on \(\ker P_L\) says that both \(\mathrm{d}_1a=0\) and
\(\mathcal T a=0\) force \(a=0\).  Equivalently use the actual
internal-bond and crossing-face zero-mode argument of
\cref{ch:flat-contour-precision-gap}.

Apply the constant linear change of variables
\(Y\mapsto(F,s)=(\mathrm{d}_1\mathsf C_0Y,\mathcal T\mathsf C_0Y)\).
Its Jacobian is a nonzero constant between the indicated Euclidean
spaces.  The quadratic exponent in
\eqref{gcb:actual-flat-precision} becomes exactly
\(\frac12\|F\|^2+\frac12\|s\|^2\).  The constant Jacobian and
the integral in \(s\) cancel from normalized curvature
expectations.  Therefore
\begin{equation}
 F=\mathrm{d}_1 Z\ \hbox{is a standard Gaussian on }\mathcal E_C,
 \qquad \mathrm{d}_1\Gamma \mathrm{d}_1^{\mathsf T}=P_{\mathcal E_C}.
 \label{gcb:constrained-curvature-projection}
\end{equation}
The same calculation with any prescribed positive quadratic form
in \(s\) changes only the gauge normalization.  It works here
because \eqref{gcb:actual-gauge-right-inverse} was proved for the
actual gauge.  It would be invalid merely to call a quadratic
penalty a gauge term and omit the range check.

If the action weight is instead \(x=\tau\theta\), \(\tau>0\),
multiply the curl term in \eqref{gcb:actual-flat-precision} by
\(\tau\), leaving the reference soft gauge coefficient one.
The same bijection gives exponent
\(\frac{\tau}{2}\|F\|^2+\frac12\|s\|^2\), so the corresponding
covariance satisfies
\(\mathrm{d}_1\Gamma_\tau \mathrm{d}_1^{\mathsf T}=\tau^{-1}P_{\mathcal E_C}\).
All numerical reference applications below use \(\tau=1\).
No uniform covariance bound as \(\tau\downarrow0\) is inferred.

\subsubsection*{The precise relation with the full free-curvature comparison}

Let \(\mathcal E=\operatorname{ran}\mathrm{d}_1\) on the same finite
periodic lattice, with each positive plaquette counted once.
Generally \(\mathcal E_C\) is a proper subspace of
\(\mathcal E\).  With \(\mathcal E_\perp=\mathcal E\cap
\mathcal E_C^\perp\), ordinary orthogonal decomposition gives
\begin{equation}
 P_{\rm ex}=P_{\mathcal E}=P_{\mathcal E_C}+P_{\mathcal E_\perp},
 \qquad F_{\rm full}=F_C+F_\perp,
 \label{gcb:free-curvature-decomposition}
\end{equation}
where the two summands on the right can be taken independent
standard Gaussians on their respective subspaces.  This proves
that the law of \(F_{\rm full}\) is the full exact-cochain
Gaussian used in the free-kernel comparison.  Conditional on
\(F_\perp\), its centered fluctuation has exactly the curvature
law in \eqref{gcb:constrained-curvature-projection}.  Conversely,
integrating this complementary Gaussian supplies the missing
curvature covariance.

The decomposition is an exact classical Gaussian identity.  It
does not identify the law of a coarse mean in an interacting run.
In particular it does not permit replacing
\(P_{\mathcal E_C}\) by \(P_{\rm ex}\) in a single fixed-fibre
Gaussian estimate.  Fine fluctuations alone omit the second
term of \eqref{gcb:free-curvature-decomposition}; exterior or
coarse averaging is a separate covariance contribution in the
complete conditional-law identity.

There is a finite-dimensional expression for this distinction
which needs no chosen coarse-coordinate parametrization.  Let
\(Q_C\) be any matrix whose kernel inside \(\mathcal E\)
is \(\mathcal E_C\).  Then
\begin{equation}
 P_{\mathcal E_C}=P_{\rm ex}
  -P_{\rm ex}Q_C^{\mathsf T}
       (Q_CP_{\rm ex}Q_C^{\mathsf T})^+
                         Q_CP_{\rm ex}.
 \label{gcb:projection-correction}
\end{equation}
To prove the formula, set \(B=Q_CP_{\rm ex}\) restricted to
\(\mathcal E\).  Singular-value decomposition shows that
\(B^{\mathsf T}(BB^{\mathsf T})^+B\) is the orthogonal
projection onto \((\ker B)^\perp\) within \(\mathcal E\).
Subtract it from the identity on \(\mathcal E\), then extend
by zero on \(\mathcal E^\perp\).  This yields exactly
\eqref{gcb:projection-correction}.  The identity is not an
estimate that the correction is small.  Such an estimate on a
particular physical window would require its own quantitative proof.

At fine physical spacing \(\eta\), write the canonical raw
curvature as \(F\).  The physical curvature of this canonically
normalized field is \(\eta^{-2}F\).  The full exact-cochain covariance is therefore
\(\eta^{-4}P_{\rm ex}\), precisely the normalization of
Requirement F1 in the dependency catalogue.
Its plaquette components are sampled at their own barycentres.
For three independent color copies and
\(Q_f=\frac12\sum_p f(x_p)|F_p|^2\), Gaussian Wick contraction
gives
\begin{equation}
 \operatorname{Cov}_0(Q_f,Q_h)
     ={1\over2}\sum_{p,q}f(x_p)h(x_q)
                    \|\mathbb E[F_pF_q^{\mathsf T}]\|_{\rm HS}^2.
 \label{gcb:curvature-wick-readout}
\end{equation}
The physical Riemann weight \(\eta^4\) multiplying the squared
physical curvature cancels its \(\eta^{-4}\) scaling in
\(Q_f\).  Thus the $\operatorname{tr}(F^2)$ readout is the unweighted
plaquette sum displayed here; inserting an additional
\(\eta^4\) into that sum would change the observable.

For the full law in \eqref{gcb:free-curvature-decomposition},
\eqref{gcb:curvature-wick-readout} has the F1 kernel and its
proved separated-smearing comparison.  For the first-step
fluctuation it has \(P_{\mathcal E_C}\) instead.  The cutoff
estimates compare each specified Gaussian reference with its own
Wick kernel.  They transfer F1 only after the reference used in
the application has actually been identified with, or compared
to, the full exact-cochain reference.  The decomposition records
how this works in the free Gaussian theory without asserting it
for a conditional interacting measure.

\subsubsection*{Uniform cutoff errors for the curvature polynomial}

Take all physical tangent bond radii equal to \(B\ge\sigma_L\)
in \eqref{gcb:actual-flat-precision}, and retain counting weights
and unit amplitudes.  From
\eqref{gcb:large-radius-flux}, including all ranks from one
through three and enlarging the fixed constants,
\begin{equation}
 \sup_b\beta_b^*\le\beta_L(B)
       :=C_L B^2e^{-c_LB^2},\qquad
 D\le S_L\beta_L(B),\qquad R=B,\quad v\le\sigma_L.
 \label{gcb:flat-exponential-error}
\end{equation}
The bounds remain valid if some physical blocks are deterministic.
For a sampled test function let
\(W_f=\sum_p|f(x_p)|\), and define the canonical quadratic
curvature polynomial
\begin{equation}
 Q_f(Z)={1\over2}\sum_p f(x_p)|\mathrm{d}_1Z(p)|^2,
 \qquad \mathcal A(Q_f)\le24W_f.
 \label{gcb:localized-curvature-norm}
\end{equation}
Indeed a fixed color component of a plaquette curl has four
coefficients of absolute value one, hence \(U\le4\).
There are three colors and the coefficient is \(f/2\), so
\((1/2)\cdot3\cdot4^2=24\).  Each positive direction pair is
already counted once in the sum over \(p\).  This calculation
requires no estimate on the number of unobserved bonds.

For \(p=2,3,4\), equations in \eqref{gcb:quadratic-connected-errors}--
\eqref{gcb:flat-exponential-error} now imply, with
\(\Delta_p=\kappa_\nu(Q_{f_1},\ldots,Q_{f_p})
-\kappa_0(Q_{f_1},\ldots,Q_{f_p})\),
\begin{equation}
 \begin{aligned}
 |\nu Q_f-\mathbb E Q_f|&\le C_L\beta_L(B)B W_f,\\
 |\Delta_p|
  &\le C_{L,p}\beta_L(B)
       (B^{2p-1}+B^{2p-3}+\cdots+B)
                        \prod_{j=1}^pW_{f_j}.
 \end{aligned}
 \label{gcb:curvature-connected-error}
\end{equation}
Here all expectations are normalized, including the sharp product
cutoff law.  The third and fourth lines implicit in the second
formula use sixth and eighth underlying linear moments.  The
vanishing of odd linear moments alone would not establish them.
The Gaussian values on the right are the actual cycle contractions in \eqref{gcb:gaussian-quadratic-cycles}, with the constrained curvature
projection of \eqref{gcb:constrained-curvature-projection} in this
particular application.

If each test is supported in a ball of current radius \(r\ge1\),
then \(W_f\le C(r+1)^4\|f\|_\infty\).  Thus the covariance
error has the expected two localized volume factors
\(C(r+1)^8\|f\|_\infty\|h\|_\infty\), rather than a factor
proportional to the volume of the torus.  At orders three and
four the factors are \((r+1)^{12}\) and \((r+1)^{16}\).
These are insertion-support counts; all constants in their
coefficients are uniform over torus size and support positions.

\subsubsection*{The nonlinear Wilson readout on the actual joint chart}

The Gaussian law need not be restricted to a merely formal observable.
At the flat background let \(a=\Phi_0(u)\) be the actual nonlinear
fibre chart of \cref{ch:block-average-pivot-chart}, so
\(u\) is the free raw logarithm and
\(D\Phi_0(0)=\mathsf C_0\).  Its domain contains
\(\|u\|_\infty<r_L\), with \(r_L=cL^{-9}\) after reducing
the fixed constant.  The finite local support and the uniform
holomorphic bounds on that joint domain give, on a smaller real ball,
\begin{equation}
 \begin{aligned}
 \|\Phi_0(u)-\mathsf C_0u\|_\infty
                      &\le C_L\|u\|_\infty^2,\\
 A_p(e^{i\Phi_0(u)})
       &={1\over2}|\mathrm{d}_1\mathsf C_0u(p)|^2+R_p(u),\\
 |R_p(u)|&\le C_L\|u\|_\infty^3.
 \end{aligned}
 \label{gcb:actual-readout-taylor}
\end{equation}
For the second assertion use the holomorphic extension
\(1-\frac12\operatorname{tr}U_p\), which equals the real Wilson
readout on \(\operatorname{SU}(2)\).  Its constant and linear terms
at the identity are zero.  Differentiating the four bond products
shows its quadratic term is
\(\frac12|\mathrm{d}_1\mathsf C_0u(p)|^2\), with Pauli pairing
\(|X|^2=\frac12\operatorname{tr}X^2\).  To bound the remainder,
apply one-variable Cauchy to the actual local function of \(tu\)
on a disc of radius proportional to
\(r_L/\|u\|_\infty\).  The joint chart bounds this function
uniformly; summing the Taylor terms of order at least three on
a smaller disc gives the displayed cubic estimate.  The first
assertion follows by the same argument at order two.  A plaquette
and its pivot dependencies occupy a fixed number of double
blocks, so neither remainder constant contains the lattice volume.

Set \(u=gY\).  Under the full physical tangent cutoff
\(|\mathsf C_0Y(b)|<B\), the identity free restriction implies
\(\|Y\|_\infty\le B\).  Thus \(gB<r_L/2\) places every
field of this cutoff event in the actual chart, and for
\(O_f=\sum_p f(x_p)A_p\),
\begin{equation}
 O_f(e^{i\Phi_0(gY)})=g^2Q_f(Z)+\mathcal R_f(Y),
 \quad |\mathcal R_f|\le C_Lg^3B^3W_f,
 \quad |Q_f|\le C_LB^2W_f.
 \label{gcb:nonlinear-readout-error}
\end{equation}
This uses the literal sampled $\operatorname{tr}(F^2)$ sum, with no extra
fine-cell volume factor.  The probability measure in this
identity is still the Gaussian cutoff measure, rather than the
nonlinear action measure.

The finite partition definition of a cumulant gives
\(|\kappa(X_1,\ldots,X_p)|\le C_p\prod_j\|X_j\|_\infty\).
Expand each slot of \eqref{gcb:nonlinear-readout-error}.  If
\(gB\le1\), the terms with at least one remainder have total
absolute value at most
\(C_{L,p}g^{2p+1}B^{2p+1}\prod_jW_{f_j}\).
Indeed a remainder replaces one factor \(g^2B^2\) by
\(g^3B^3\); the finite sum over subsets costs at most a
constant times \(gB(1+gB)^{p-1}\).
Together with \eqref{gcb:curvature-connected-error}, this proves
\begin{equation}
 \begin{aligned}
 &|\kappa_\nu(O_{f_1},\ldots,O_{f_p})
                    -g^{2p}\kappa_0(Q_{f_1},\ldots,Q_{f_p})|\\
 &\quad\le C_{L,p}g^{2p}\left[
       \beta_L(B)(B^{2p-1}+B^{2p-3}+\cdots+B)
                          +gB^{2p+1}\right]\prod_{j=1}^pW_{f_j},
 \qquad p=2,3,4.
 \end{aligned}
 \label{gcb:actual-readout-connected-bound}
\end{equation}
The same reasoning for one insertion gives mean error
\(C_Lg^2[\beta_L(B)B+gB^3]W_f\).
This is a concrete sharp-cutoff Gaussian clean-law estimate with
the actual nonlinear observable numerator.  It includes neither
the nonlinear change in its probability density nor any missing
coarse-mean contribution to the chosen free reference.

\subsubsection*{Literal radius conversions and comparison of cutoff sets}

At physical fine spacing \(\eta\), raw bond logarithm and physical
potential obey \(a_b=\eta A_b\).  In canonical tangent variables
\(a_b=g Z_b\).  Consequently the conversions are
\begin{equation}
 |a_b|<\alpha_b\ \Longleftrightarrow\
       |Z_b|<\alpha_b/g,\qquad
 |\ell_b A_b|<r_b\ \Longleftrightarrow\
       |Z_b|<{\eta r_b\over g\ell_b}.
 \label{gcb:literal-radius-conversion}
\end{equation}
The second relation refers to the tangent physical potential;
\(\ell_b\) is the actual scale multiplying it in that field
bound.  For the nonlinear chart use
\eqref{gcb:actual-readout-taylor}: its raw field differs from
\(gZ\) by at most \(C_Lg^2B^2\) on the event in question.
Any comparison with a cutoff on the nonlinear field must charge
this error and preserve a stated radius margin.

A free-coordinate box and the complete physical tangent product
box are not equal events.  With
\(c_L=\|\mathsf C_0\|_{\infty\to\infty}\ge1\), one has only
\begin{equation}
 \{\|Y\|_\infty<B/c_L\}
   \subset\{\|\mathsf C_0Y\|_\infty<B\}
   \subset\{\|Y\|_\infty<B\}.
 \label{gcb:free-physical-box-inclusions}
\end{equation}
The first inclusion uses the actual row norm of the pivot map;
the second uses its identity free restriction.  Inclusion alone
does not compare the corresponding normalized conditional
expectations.  The proof in this section works directly with
the middle, complete physical event, so requires no false
identification of these boxes.  If an actual integration rule
instead cuts only free coordinates, apply the abstract theorem
to those actual independent-coordinate blocks and carry the
observable coefficient map separately; it is a different cutoff
law with its own explicitly specified covariance and radii.

For growing canonical radius
\(B=c_*(\log\theta)^s\), \(s>0\), and
\(g=\theta^{-1/2}\), every fixed power of \(\log\theta\)
times \(\beta_L(B)\) tends to zero, as does every fixed such
power times \(g\).  The chart condition \(gB<r_L/2\) holds
eventually.  Thus for support radii bounded by any fixed power
of \(\log\theta\), the bracket and localized volume factors
in \eqref{gcb:actual-readout-connected-bound} tend to zero,
uniformly in torus volume.  This is an absolute error relative
to the displayed insertion norms.  Converting it into a relative
error for a specified two-point coefficient also requires a
proved lower bound for that coefficient in the same test class.

\subsubsection*{What passes to limits and what remains an interaction estimate}

The finite-volume cutoff laws have Gaussian-dominated moments
for every fixed local linear statistic by
\eqref{gcb:linear-moment-domination}.  In the fixed-step setting
their marginal variances are bounded uniformly by
\eqref{gcb:realized-flat-schur}.  These bounds give tightness on
each finite collection of bond coordinates.  A diagonal
subsequence on the countable lattice therefore has compatible
finite-dimensional limits.  Uniform higher local moments give
uniform integrability, so the local polynomial estimates above
pass to every such limit.  This conclusion does not assert
uniqueness of the limit, and does not identify it with conditioning
an infinite-volume Gaussian on a global event of positive
probability.  It only uses uniform finite-volume bounds.

The boundary proof also identifies its precise obstruction in
a general conditional law.  If in actual free coordinates, with
\(G>0\), its
complete density is
\(Z_e^{-1}\chi(CY)\exp[-\frac12Y^{\mathsf T}G^{-1}Y+R_e(Y)]\),
integration by parts has, in addition to the boundary defect,
the force term
\begin{equation}
 \mathbb E_{\nu_e}\!\left[
       P(CY)\,(CG\nabla R_e(Y))_i\right].
 \label{gcb:remaining-complete-force}
\end{equation}
This follows by differentiating the entire density, with its
actual normalizer; all integrability and chart-domain conditions
must hold.  An arbitrary exterior linear force can translate
the law toward the boundary, and an arbitrary old activity can
destroy the even log-concavity used in
\eqref{gcb:conditional-weight}.  Neither its small Hessian alone
nor the Gaussian boundary theorem bounds these two effects.

The same radial argument would remain valid for a further even
log-concave factor in the physical variables, because that factor
can be included in the joint log-concave integrand defining
\(w_{b,\delta}\).  Its new force term still has to be bounded
in the complete integration-by-parts identity.  No such convexity,
symmetry or force bound is asserted for the resummed interacting
conditional density.

Accordingly the result established here is the boundary and
connected-error estimate for the specified centered classical
Gaussian sector, with a realized fixed-step covariance and actual
observable coordinates.  Applying it to the interacting extraction
requires the complete measure, exterior averaging, old-activity
and force identifications of the conditional-law section, as well
as the correct full free reference for F1.  It supplies a proved
reference estimate for that boundary task, rather than assuming
the interacting extraction law or an all-depth unweighted
covariance bound.




% END MODULE chapters/gaussian_cutoff_clean_law.tex

% BEGIN PROOF EXPANSION finite_window_gaussian_law.tex
\subsection{A normalized Gaussian comparison on a finite run window}
\label{ch:finite-window-gaussian-law}

The conditional interacting law can be compared directly with a Gaussian
when the complete residual exponent has small oscillation on the tested
window. The oscillation is measured after removing a field-independent
constant, so an exterior-dependent vacuum energy causes no loss. This
route is useful when the number of retained coordinates and their tested
ranges grow only as powers of the logarithm of the inverse coupling.
It estimates the normalizer together with the numerator and avoids an
integration-by-parts argument at a sharp boundary.

The criterion below concerns the complete exponent. In particular it
includes the coarea density, Haar density, history multiplicity and all
old activities that vary with the retained coordinates. Bounds on the
nonlinear Wilson action alone are insufficient for its application.
We first prove the comparison, then give explicit estimates for several
of the factors that occur when the criterion is verified.

\subsubsection{The tested domain and the reference Gaussian}

Let $\gamma$ be a centered Gaussian on $\mathbb R^N$ with positive
definite covariance $C$. If the Gaussian is supported on a proper
linear subspace, choose orthonormal coordinates on that support
throughout; the statements then have exactly the same meaning.
Let $v=\max_jC_{jj}$. For $R>0$, put
\begin{equation}\label{fwg:box}
 B_R=\{y:\max_j|y_j|\le R\}.
\end{equation}
Let $D$ be measurable and contain $B_R$, up to a Gaussian null set.
The domain may be nonconvex. Define $t=\gamma(D^c)$.
The one-dimensional Gaussian bound and the union bound give
\begin{equation}\label{fwg:domain-tail}
 t\le2N\exp[-R^2/(2v)].
\end{equation}
Here and below a zero marginal variance can simply be omitted.
For completeness, for a marginal $Y$ with variance $\sigma^2$,
$\mathbb E e^{sY}=e^{s^2\sigma^2/2}$. Markov's inequality at
$s=R/\sigma^2$ bounds $\mathbb P(Y\ge R)$ by
$e^{-R^2/(2\sigma^2)}$; apply the same estimate to $-Y$.
This argument uses no independence between coordinates.

Suppose a real measurable function $P$ on $D$ has oscillation at most
$2\alpha$ in the following specified sense: there is a real constant
$c$ such that
\begin{equation}\label{fwg:oscillation}
 |P(y)-c|\le\alpha\quad\hbox{for $\gamma$-almost every $y\in D$}.
\end{equation}
The constant may depend on the exterior configuration. Set
\begin{equation}\label{fwg:interacting-law}
 d\nu=Z^{-1}1_De^{P-c}\,d\gamma,
 \qquad Z=\gamma(1_De^{P-c}).
\end{equation}
If $t<1$, this is a probability and
\begin{equation}\label{fwg:normalizer}
 e^{-\alpha}(1-t)\le Z\le e^\alpha(1-t).
\end{equation}
The sharper upper bound displayed here is useful: the clean domain
mass belongs to the same reference in both numerator and denominator.
Let $\gamma_D=\gamma(\,\cdot\mid D)$. Then
\begin{equation}\label{fwg:relative-density}
 e^{-2\alpha}\le\frac{d\nu}{d\gamma_D}\le e^{2\alpha}
 \quad\hbox{almost everywhere}.
\end{equation}
This proves non-vanishing of the real normalizer with its correct
domain factor. It does not assert a complex source-domain estimate.

\begin{proposition}[Weighted comparison of the normalized laws]
\label{prop:fwg-weighted-comparison}
For every $H\in L^2(\gamma)$,
\begin{align}
 |\nu(H)-\gamma(H)|
 &\le\frac{e^{2\alpha}-1+t}{1-t}\,\gamma(|H|)
       +\frac{\gamma(|H|1_{D^c})}{1-t}
       \label{fwg:weighted-comparison}\\
 &\le\frac{e^{2\alpha}-1+t+\sqrt t}{1-t}
                                \|H\|_{L^2(\gamma)}.
       \label{fwg:l-two-comparison}
\end{align}
Also $\nu(|H|)\le e^{2\alpha}(1-t)^{-1}\gamma(|H|)$
for every nonnegative integrable $|H|$.
\end{proposition}
\begin{proof}
The density estimate \eqref{fwg:relative-density} gives
$|\nu(H)-\gamma_D(H)|\le(e^{2\alpha}-1)\gamma_D(|H|)$.
On the other hand,
\[
 \gamma_D(H)-\gamma(H)
  =\frac{t\gamma(H)-\gamma(H1_{D^c})}{1-t}.
\]
Taking absolute values proves \eqref{fwg:weighted-comparison}.
Cauchy--Schwarz bounds its last tail integral by
$\sqrt t\|H\|_2$, and $\gamma(|H|)\le\|H\|_2$.
The last assertion follows by direct integration of
\eqref{fwg:relative-density}.
\end{proof}

The estimate only uses a Gaussian measure of the symmetric difference
from its full support through $t$. Thus one can replace a complicated
tested domain by a convenient inner box, provided that box actually
lies in the image of the physical tests. An upper bound on the size
of every retained coordinate does not prove this inner inclusion.
The latter is a separate domain-admission calculation.

\subsubsection{Dimension-explicit moments of quadratic readouts}

Let $A$ be real symmetric and set
\begin{equation}\label{fwg:quadratic-size}
 L(A)=\operatorname{tr}|C^{1/2}AC^{1/2}|,
 \qquad Q_A(y)=y^{\mathsf T}Ay-\operatorname{tr}(AC).
\end{equation}
The absolute value in $L(A)$ is the operator absolute value of a
real symmetric matrix. This definition also covers signed profiles.
It often gives a better bound than the matrix dimension times a
supremum of its entries.

For an integer $k\ge1$, diagonalize $C^{1/2}AC^{1/2}$ with eigenvalues
$\lambda_j$. Under $\gamma$, its quadratic form has the law
$\sum_j\lambda_j G_j^2$, where the $G_j$ are independent standard
Gaussians. The triangle inequality in $L^k$ yields
\begin{equation}\label{fwg:quadratic-moments}
 \|Q_A\|_{L^k(\gamma)}\le u_kL(A),
 \qquad u_k=((2k-1)!!)^{1/k}+1.
\end{equation}
Indeed $\|G_j^2\|_k=((2k-1)!!)^{1/k}$ and
$|\sum_j\lambda_j|\le\sum_j|\lambda_j|$.
The Gaussian even moment formula used here follows by integration
by parts: $\mathbb E G^{2k}=(2k-1)\mathbb E G^{2k-2}$ with
$\mathbb E G^0=1$. The estimate includes singular matrices $A$;
if $L(A)=0$, the quadratic readout is identically zero on the
Gaussian support.

For readouts $Q_i=Q_{A_i}$ and any nonempty set $S$ of $s$ indices,
H\"older gives
\begin{equation}\label{fwg:product-moments}
 \left\|\prod_{i\in S}Q_i\right\|_{L^2(\gamma)}
               \le u_{2s}^{\,s}\prod_{i\in S}L_i,
 \qquad L_i=L(A_i).
\end{equation}
There is no factor depending on $N$ apart from the actual sizes
$L_i$. This matters when an observable sums over a window: its
coefficient mass must be counted once in $L_i$, rather than once
for each moment estimate.

Let $0<g\le1$. Suppose the physical readouts on $D$ are
\begin{equation}\label{fwg:physical-readouts}
 X_i(y)=c_i+g^2[Q_i(y)+I_i(y)],
 \qquad |I_i(y)|\le\rho_iL_i.
\end{equation}
The constants $c_i$ are deterministic at the fixed exterior;
their variation with the exterior is retained when applying
the law of total cumulance. Assume $L_i>0$ and
$0\le\rho_i\le1$. An entry with $L_i=0$ may instead be handled
separately using a positive bound for its remainder.

\begin{proposition}[Mixed moments of the physical readouts]
\label{prop:fwg-readout-moments}
Assume $\alpha\le1/4$ and $t\le1/2$. For every nonempty $S$ of
size $s\le n$,
\begin{equation}\label{fwg:readout-moment-error}
 \left|\nu\!\left(\prod_{i\in S}(X_i-c_i)\right)
    -g^{2s}\gamma\!\left(\prod_{i\in S}Q_i\right)\right|
 \le K_n\left(\alpha+\sqrt t+\max_i\rho_i\right)
                         g^{2s}\prod_{i\in S}L_i,
\end{equation}
where $K_n$ depends only on $n$. The absolute mixed moments of
both displayed laws are at most $K_ng^{2s}\prod_{i\in S}L_i$.
\end{proposition}
\begin{proof}
First compare the product of $Q_i$ under $\nu$ and $\gamma$.
Apply \cref{prop:fwg-weighted-comparison} and
\eqref{fwg:product-moments}. Since $\alpha\le1/4$ and $t\le1/2$,
the coefficient there is at most a fixed multiple of
$\alpha+\sqrt t$.

Next expand $\prod(Q_i+I_i)-\prod Q_i$ by the nonempty subset
of positions at which the remainder $I_i$ is used. Its absolute
expectation is bounded by
\[
 \sum_{\varnothing\ne J\subset S}
    \prod_{i\in J}(\rho_iL_i)\,
           \nu\!\left(\prod_{i\in S\setminus J}|Q_i|\right).
\]
The final assertion of \cref{prop:fwg-weighted-comparison}
and H\"older bound the last moment by a constant depending only
on $n$ times $\prod_{i\in S\setminus J}L_i$. The sum of
the remainder coefficients is at most
$\prod_{i\in S}(1+\rho_i)-1\le(2^n-1)\max_i\rho_i$.
Multiply by $g^{2s}$ to obtain
\eqref{fwg:readout-moment-error}. The same estimates without
subtraction bound the absolute mixed moments.
\end{proof}

\subsubsection{The conditional connected coefficient}

The cumulant polynomial in
\cref{ch:positive-history-comparison} now proves, for $2\le n\le4$,
\begin{equation}\label{fwg:conditional-cumulant-error}
 \left|\kappa_\nu(X_1,\ldots,X_n)
       -g^{2n}\kappa_\gamma(Q_1,\ldots,Q_n)\right|
 \le K'_n\left(\alpha+\sqrt t+\max_i\rho_i\right)
                                      g^{2n}\prod_iL_i.
\end{equation}
Here the constants $c_i$ disappear by translation invariance of
a cumulant of order at least two. To see the stated dependence,
use \eqref{fwg:readout-moment-error} for the one replaced block
in the telescoping product proof of
\cref{prop:phc-connected-error}. The other blocks have the
moment bound of \cref{prop:fwg-readout-moments}. There are only
finitely many partitions at fixed $n$, giving $K'_n$.

For the one-point function the corresponding estimate is
\begin{equation}\label{fwg:conditional-mean-error}
 |\nu(X_i)-c_i|\le
       K_1(\alpha+\sqrt t+\rho_i)g^2L_i,
\end{equation}
because $\gamma(Q_i)=0$. When exteriors vary, the error from
$c_i(e)$ relative to a fixed reference must be added to this
bound; it is not removed by centering the global observable.

The reference quadratic cumulants have a finite trace formula.
Put $B(z)=\sum_i z_iC^{1/2}A_iC^{1/2}$. For
$2\|B(z)\|<1$, diagonalizing the real Gaussian integral first
for real $z$ near zero, and then continuing analytically, gives
\begin{align}
 \log\gamma\!\left(e^{\sum_i z_iQ_i}\right)
 &=-\operatorname{tr}B(z)-\tfrac12\log\det(I-2B(z))
       \notag\\
 &=\sum_{m\ge2}\frac{2^{m-1}}m\operatorname{tr}B(z)^m.
 \label{fwg:gaussian-trace-series}
\end{align}
The determinant branch equals one at zero. The logarithmic series
converges absolutely in operator norm on every smaller ball,
so coefficient extraction is legitimate. In particular
\begin{equation}\label{fwg:gaussian-covariance}
 \kappa_\gamma(Q_1,Q_2)=2\operatorname{tr}(A_1CA_2C).
\end{equation}
At order four it is the sum over closed four-cycles, with the
coefficient and colour multiplicity fixed in the free curvature
calculation. Thus \eqref{fwg:conditional-cumulant-error} controls
an interacting connected coefficient, rather than merely a
change in its quadratic precision.

\subsubsection{A finite-dimensional graph for the constraints}

We next estimate one contribution to $P$. Let $p\in\mathbb R^m$
be pivot variables and $y\in\mathbb R^N$ retained variables.
Suppose a real $C^1$ constraint is $q(p,y)=0$. On an open
neighborhood of the retained domain assume it has one admitted
solution $p=\varphi(y)$ and that $D_pq(\varphi(y),y)$ is invertible.
These assumptions concern the actual constraint and its chosen
branch; the implicit function theorem supplies only a local
solution until uniqueness and domain coverage are checked.

Let $Q$ be a fixed invertible pivot matrix and suppose
\begin{equation}\label{fwg:pivot-perturbation}
 D_pq(\varphi(y),y)=Q[I+E(y)],
 \qquad \|E(y)\|\le\delta<1.
\end{equation}
The operator norm here is Euclidean. Then
\begin{equation}\label{fwg:coarea-log-bound}
 \left|\log\frac{|\det D_pq(\varphi(y),y)|}{|\det Q|}\right|
                       \le\frac{m\delta}{1-\delta}.
\end{equation}
Indeed $I+sE$ is invertible for $0\le s\le1$. Its real
determinant remains positive, since it equals one at $s=0$.
The logarithmic derivative is
$\operatorname{tr}[(I+sE)^{-1}E]$, whose absolute value is at
most $m\delta/(1-s\delta)$. Integrating and using the weaker
constant $1/(1-\delta)$ proves the estimate. The same proof
works for a complex branch of the determinant logarithm on a
norm ball, but the probability comparison itself remains real.

Consequently the coarea factor
$|\det D_pq|^{-1}$ has a field-dependent logarithm bounded,
after subtraction of $-\log|\det Q|$, by the right hand side
of \eqref{fwg:coarea-log-bound}. The dimension $m$ is essential.
An operator perturbation $\delta=o(1)$ alone does not imply a
small total determinant logarithm if $m\delta$ is large.

The identity underlying this coarea factor is elementary but fixes
which determinant occurs. The map $(p,y)\mapsto(q(p,y),y)$ has
block-triangular derivative with determinant $\det D_pq$.
Changing variables in an integral with the constraint delta
function and then evaluating at $q=0$ gives
\begin{equation}\label{fwg:coarea-density}
 \int f(p,y)\delta(q(p,y))\,dp\,dy
  =\int f(\varphi(y),y)|\det D_pq(\varphi(y),y)|^{-1}\,dy
\end{equation}
on the admitted branch. With several branches the right side
is summed over them. A gauge quotient or a change of the
retained coordinates contributes its own factor and must be
included as well.

For a useful quantitative graph estimate, suppose
\begin{equation}\label{fwg:graph-equation}
 q(p,y)=Qp+Ry+g\mathcal N(p,y),
\end{equation}
where $\|Q^{-1}\|\le K$ and $\mathcal N$ is Lipschitz in $p$
with constant $L_p$ on a specified closed pivot ball for each
$y$ in the retained domain. Put $p_0(y)=-Q^{-1}Ry$.
If
\begin{equation}\label{fwg:graph-admission}
 gKL_p\le\tfrac12,\qquad
 gK\sup_y|\mathcal N(p_0(y),y)|\le b/2,
\end{equation}
and the radius-$b$ ball about every $p_0(y)$ lies in the
declared pivot domain, the map
$p\mapsto p_0(y)-gQ^{-1}\mathcal N(p,y)$ is a contraction
of that ball. Thus
\begin{equation}\label{fwg:graph-error}
 |\varphi(y)-p_0(y)|
       \le2gK|\mathcal N(p_0(y),y)|.
\end{equation}
The proof is the contraction inequality with its inhomogeneous
term moved to the left. If $\|D_p\mathcal N\|\le L_p$, then
\eqref{fwg:pivot-perturbation} holds with $\delta=gKL_p$.
This is a finite-window criterion with a specified radius.
For an actual multiscale constraint, a uniform $K$ and coverage
of the entire tested set still have to be proved.

\subsubsection{Haar factors and the action remainder}

The analytic Haar density and the bare Wilson remainder
are estimated in \eqref{sf:haar-ratio} and
\eqref{sf:bare-Wilson-remainder}. In particular, at a fixed
coupling, $M$ link logarithms $ga$ with $|a|\le R_1$ and
$gR_1\le1/4$ contribute at most $CMg^2R_1^2$ to the
Haar logarithm after subtraction of its value at zero.
The coordinate graph may change $a$ to a function of $y$;
the tested bound must then hold on that whole graph.
The fixed Jacobian powers of $g$ cancel in normalized
expectations. Source-dependent powers must be retained.

Let $S$ be a real analytic action on a finite coordinate chart
and write $a=gTy$, where $T$ is a specified linear map. Assume
the quadratic reference has been extracted exactly, so
\begin{equation}\label{fwg:action-taylor}
 g^{-2}S(gTy)=c_S+\ell^{\mathsf T}y
                 +\tfrac12y^{\mathsf T}Qy+R_S(y).
\end{equation}
If $S'(0)=0$ and the quadratic form is the actual Hessian,
then $\ell=0$. At a noncritical reference the linear term
must be retained or removed by a verified Gaussian shift.
Suppose that for every $y\in D$ the full segment
$a=tgTy$, $0\le t\le1$, lies in the declared analytic
chart, and that on that segment the third derivative obeys
\begin{equation}\label{fwg:third-action-bound}
 |D^3S(a)[Ty,Ty,Ty]|\le K_S(y).
\end{equation}
Taylor's integral remainder gives
\begin{equation}\label{fwg:action-remainder-bound}
 |R_S(y)|\le\tfrac16 gK_S(y).
\end{equation}
The power $g$ comes from the cubic numerator divided by
$g^2$. In a local action with $M$ terms, each depending on
a uniformly bounded list of coordinates of size at most
$R_1$, a uniform third-derivative bound gives
$K_S(y)\le CMR_1^3$. This assertion requires those local
derivative bounds; when $S$ is the Wilson action composed
with a constrained background, the derivatives of that
background are part of them.

If the graph of the constraints is nonlinear, it is equally
valid to take the third derivative of the composed action
on the graph. Replacing that derivative by the bare Wilson
third derivative would omit derivatives of $\varphi$.
The complete constraint-curvature term in the Hessian is
also part of the reference $Q$ in
\eqref{fwg:action-taylor}. The present comparison does not
justify replacing it by a flat Feynman precision.

\subsubsection{Old activities that meet the window}

Let $W$ be a finite set of retained cells. Suppose an old
interaction has a convergent representation
\begin{equation}\label{fwg:old-interaction}
 H(y)=\sum_XH_X(y),
\end{equation}
where a term is independent of $y$ if $X$ does not meet $W$.
Let $y_*$ be one admitted reference. If
\begin{equation}\label{fwg:activity-oscillation}
 \sup_{v\in W}\sum_{X\ni v}
       \sup_{y\in D}|H_X(y)-H_X(y_*)|\le\epsilon,
\end{equation}
then
\begin{equation}\label{fwg:old-total-oscillation}
 \sup_{y\in D}|H(y)-H(y_*)|\le|W|\epsilon.
\end{equation}
To prove this, bound the sum over terms that meet $W$ by
the sum over all incidences $(v,X)$ with $v\in W\cap X$.
Each contributing term is counted at least once. The
positive bound in \eqref{fwg:activity-oscillation} gives
absolute and uniform convergence of the varying part, so
the subtraction and sum are justified.

The norm in \eqref{fwg:activity-oscillation} is a norm of
the evaluated terms on one common admitted chart. A formal
norm at different centers, or a bound after a different
levelled substitution, cannot be inserted without an
estimate transferring it to this domain. Terms with old
fine-scale origins can meet the retained window in many
fine cells. Their incidences must be charged in their
actual units; the number $|W|$ cannot be taken to be the
coarse window volume unless the formation theorem has
already performed that accounting.

This qualification is one reason to state the criterion
in terms of the complete oscillation rather than claim
that every polylogarithmic coarse window automatically
has a small residual exponent. The primitive spatial
filing and the admitted substitution are what decide
whether \eqref{fwg:activity-oscillation} is available.

\subsubsection{Putting the complete exponent together}

Suppose the actual constrained positive density, after a
fixed change to its retained coordinates, has the form
\begin{equation}\label{fwg:complete-density}
 1_D(y)\exp\{-\tfrac12y^{\mathsf T}Qy+P(y)\}\,dy,
\end{equation}
with the same $Q$ as the covariance $C=Q^{-1}$. A Gaussian
shift may be made first if its determinant, domain image
and readout change are retained. Suppose there are
field-independent constants subtracting the following
parts of $P$, and that their residual absolute bounds
are respectively
\begin{equation}\label{fwg:component-bounds}
 \alpha_{\rm act},\quad\alpha_{\rm Haar},\quad
 \alpha_{\rm coarea},\quad\alpha_{\rm old},\quad
 \alpha_{\rm hist},\quad\alpha_{\rm bdry}.
\end{equation}
They denote the nonlinear action, Haar density, coarea and
coordinate Jacobians, old activities, resummed history
density, and the remaining boundary or stored-field terms.
Then \eqref{fwg:oscillation} holds with their sum
\begin{equation}\label{fwg:total-alpha}
 \alpha=\alpha_{\rm act}+\alpha_{\rm Haar}
       +\alpha_{\rm coarea}+\alpha_{\rm old}
       +\alpha_{\rm hist}+\alpha_{\rm bdry}.
\end{equation}
This follows from addition of logarithms on the specified
positive branch and the triangle inequality. A factor
that can vanish must be kept in the domain or resolved
into positive components; it cannot simply be assigned
a bounded logarithm.

The preceding estimates give, for example,
\begin{equation}\label{fwg:sample-alpha-budget}
 \alpha\le CMgR_1^3+CMg^2R_1^2
       +\frac{m\delta}{1-\delta}+|W|\epsilon
       +\alpha_{\rm hist}+\alpha_{\rm bdry}.
\end{equation}
Every quantity on the right is measured in the actual
retained coordinate chart. This is a sufficient estimate,
not a substitute for the actual determinant or old-core
calculation.

\subsubsection{A scale criterion for the second and fourth correlations}

Write $g=x^{-1/2}$, $x\to\infty$. Suppose all of the following
bounds hold uniformly on the admitted exteriors:
\begin{align}
 N+M+m+|W|&\le C(\log x)^A,
 &R_1&\le C(\log x)^B,\notag\\
 \delta+\epsilon+\alpha_{\rm hist}+\alpha_{\rm bdry}
       &\le Cx^{-a}(\log x)^E,
 &a&>0,\notag\\
 L_i&\le C(\log x)^{A_i}\|f_i\|_\infty,
 &\rho_i&\le Cx^{-a}(\log x)^E.
 \label{fwg:polylogarithmic-window}
\end{align}
Suppose also that $D$ contains a box for which
$R^2/v\ge c(\log x)^b$, $b>0$. Then
\cref{fwg:domain-tail,fwg:sample-alpha-budget} give
\begin{equation}\label{fwg:alpha-tail-limits}
 \alpha\le Cx^{-\min(a,1/2)}(\log x)^{E'},
 \qquad
 t\le C(\log x)^A e^{-c'(\log x)^b}.
\end{equation}
Thus \eqref{fwg:conditional-cumulant-error} yields
\begin{equation}\label{fwg:scaled-conditional-law}
 \kappa_\nu(X_1,\ldots,X_n)
  =x^{-n}\kappa_\gamma(Q_1,\ldots,Q_n)
           +o\!\left(x^{-n}\prod_i\|f_i\|_\infty\right),
 \qquad n=2,4.
\end{equation}
Indeed every power of $\log x$ in $\prod L_i$ is outweighed
both by a positive power of $x$ and by
$e^{-c'(\log x)^b/2}$. The same estimates give the lower
conditional errors relative to their conditional Gaussian values.
If $C$ or $A_i$ varies with the exterior, those Gaussian
cumulants must also be compared with one fixed reference before
applying the exterior averaging formula. If $c_i(e)$ varies,
its error relative to a deterministic reference must likewise
be added to \eqref{fwg:conditional-mean-error}.

The finite-window criterion supplies one explicit route to
the interacting law. It has no boundary-flux assumption and
does not require convexity of the image of the tests. Its
cost is the complete total-oscillation bound and an inner
Gaussian box on one actual chart. When fine-scale origins,
stored terms or class multiplicities prevent a small total
oscillation, the force method of
\cref{ch:clean-conditional-law} is the appropriate alternative;
one cannot infer \eqref{fwg:scaled-conditional-law} from a
small old-activity Hessian in that situation.

% END PROOF EXPANSION finite_window_gaussian_law.tex

% BEGIN PROOF EXPANSION four_point/four_point_kernel.tex
\subsection{The four-point curvature kernel: quantitative sign and lattice stability}
\label{ch:four-point-kernel-recovered}

The free curvature calculation fixes the normalization of every separated
Gaussian cycle.  Its collinear four-point value is positive, but a
continuity argument at one configuration does not give a uniform margin
on the class of test functions used in the short-distance assertion.
This subsection supplies that margin.  It also computes a negative
configuration, proves that its sign persists on an explicit open set,
and transfers both signs to the free lattice reference at a uniform mesh
floor.  These are statements about the leading kernel.  They do not
compare an interacting measure with its Gaussian reference.

We use the covariance convention of the free curvature chapter.  Put
$\mathcal V=\Lambda^2\mathbb R^4$, with orthonormal basis
$e_P=e_\mu\wedge e_\nu$, $P=(\mu,\nu)$, $\mu<\nu$.  Thus
$\dim\mathcal V=6$.  The position-space curvature kernel is denoted
by $\mathsf K$.  The lattice comparison used below is precisely
\eqref{eq:lcc-main}; no stronger lattice error exponent is assumed.

\subsubsection{Reflections and the closed-ring formula}

For a unit vector $u\in\mathbb R^4$ define
\begin{equation}
 s_u=I-2u\otimes u,\qquad R_u=\Lambda^2s_u,
 \qquad P_u\omega=u\wedge\iota_u\omega.
 \label{eq:four-point-reflections}
\end{equation}
Here $\iota_u$ denotes contraction.  The map $P_u$ is the orthogonal
projection onto $u\wedge u^\perp$, so $R_u=I-2P_u$.
In a basis $u,v_1,v_2,v_3$, $R_u$ is minus the identity on
$u\wedge v_i$ and plus the identity on $v_i\wedge v_j$.
Consequently it is a symmetric orthogonal involution, with three
eigenvalues of each sign.  For $x\ne0$ the normalization is
\begin{equation}
 \mathsf K(x)=\frac{R_{\widehat x}}{\pi^2|x|^4}.
 \label{eq:four-point-kernel-reflection}
\end{equation}
Indeed, differentiating $G(x)=(4\pi^2|x|^2)^{-1}$ gives
\begin{equation}
 \partial_\mu\partial_\rho G(x)
 =-\frac{\delta_{\mu\rho}-4\widehat x_\mu\widehat x_\rho}
              {2\pi^2|x|^4}.
 \label{eq:four-point-green-derivative}
\end{equation}
Substitution into the curvature covariance gives the entries of
$\pi^{-2}|x|^{-4}(I-2P_{\widehat x})$.  This verifies that the
reflection in \eqref{eq:four-point-reflections} acts first on
$\mathbb R^4$; its second exterior power, rather than a rank-one
reflection of $\mathcal V$, occurs in the curvature kernel.

Let $\mathcal R$ be the three unoriented cycles
$(1234)$, $(1243)$, $(1324)$.  For pairwise distinct points put
\begin{equation}
 k_4(x_1,x_2,x_3,x_4)
 =\sum_{\varsigma\in\mathcal R}
   \operatorname{tr}_{\mathcal V}
       \prod_{j=1}^4\mathsf K(x_{\varsigma(j)}-x_{\varsigma(j+1)}),
 \qquad \varsigma(5)=\varsigma(1).
 \label{eq:four-point-box-density}
\end{equation}
The products are ordered along the cycle.  A cyclic change of its
starting point preserves the trace.  Transposition reverses the
order, and $\mathsf K(-x)=\mathsf K(x)$, so reversing a cycle also
preserves its contribution.

\begin{proposition}[The trace of a closed ring]
\label{prop:four-point-closed-ring}
Let $w_1,\ldots,w_4\ne0$ and $\sum_jw_j=0$.  There is an
$\alpha\in[0,\pi]$ such that
\begin{equation}
 \operatorname{tr}_{\mathcal V}
     [\mathsf K(w_1)\mathsf K(w_2)\mathsf K(w_3)\mathsf K(w_4)]
 =\frac{2+4\cos\alpha}{\pi^8}\prod_{j=1}^4|w_j|^{-4}.
 \label{eq:four-point-closed-ring}
\end{equation}
In particular the absolute value of the angular factor is at most six.
\end{proposition}
\begin{proof}
Exterior powers preserve products.  Thus the angular matrix in the
left side is $\Lambda^2T$, where
$T=s_{\widehat w_1}\cdots s_{\widehat w_4}\in\operatorname{SO}(4)$.
The space $V=\operatorname{span}\{w_1,\ldots,w_4\}$ has dimension
at most three.  Each factor preserves $V$ and fixes $V^\perp$.
If $\dim V=3$, then $T|_V\in\operatorname{SO}(3)$ has an eigenvalue
one: the other nonreal eigenvalues come in conjugate pairs, while
the product of all eigenvalues is one.  If $\dim V\le2$, the
fixed space $V^\perp$ already has dimension at least two.
In all cases $T$ fixes a plane and rotates its orthogonal
complement.  Its eigenvalues are $1,1,e^{i\alpha},e^{-i\alpha}$.

For any four-dimensional matrix $T$, summing its diagonal
two-by-two minors gives
\begin{equation}
 \operatorname{tr}(\Lambda^2T)
 =\tfrac12\{(\operatorname{tr}T)^2-\operatorname{tr}(T^2)\}.
 \label{eq:four-point-exterior-trace}
\end{equation}
Equivalently this trace is the sum of the six pairwise products of
the eigenvalues.  For the eigenvalues just found it is
$2+4\cos\alpha$.  Multiplication by the four radial factors in
\eqref{eq:four-point-kernel-reflection} proves the formula.
\end{proof}

An orientation-preserving or orientation-reversing orthogonal
change of coordinates conjugates every angular matrix by the same
$\Lambda^2O$.  Hence $k_4$ is invariant under isometries and
homogeneous of degree $-16$.  Permuting its arguments permutes the
three cycles.  For separated integrable functions define
\begin{equation}
 \mathcal B_4(f_1,f_2,f_3,f_4)
 =\int_{(\mathbb R^4)^4}k_4(x_1,x_2,x_3,x_4)
                    \prod_{i=1}^4f_i(x_i)\,d x_1\cdots d x_4.
 \label{eq:four-point-integrated-box}
\end{equation}
If every two supports have distance at least $d>0$, absolute
convergence and \cref{prop:four-point-closed-ring} give
\begin{equation}
 |\mathcal B_4(\underline f)|
 \le18\pi^{-8}d^{-16}\prod_i\|f_i\|_{L^1}.
 \label{eq:four-point-absolute-bound}
\end{equation}
The sixteen powers of length in the measure cancel the homogeneity
of $k_4$.  Thus simultaneous dilations of the four arguments leave
$\mathcal B_4$ unchanged.

\subsubsection{A uniform positive cone of configurations}

For $0<\vartheta\le1$ consider nonzero nonnegative
$f_i\in C_c^1(\mathbb R^4)$ with $F_i=\|f_i\|_\infty$, each
supported in a ball of radius $d$, and satisfying, for $i\ne j$,
\begin{equation}
 \begin{split}
 d&\le\operatorname{dist}(\operatorname{supp}f_i,
                         \operatorname{supp}f_j)\le6d,\\
 \int f_i&\ge\vartheta d^4F_i,
       \qquad \operatorname{Lip}f_i\le F_i/(\vartheta d).
 \end{split}
 \label{eq:four-point-shape-class}
\end{equation}
The positive-cone condition is that one unit vector $e$ works for
every pair of points in distinct supports:
\begin{equation}
 |w-(w\cdot e)e|\le|w|/8,
 \qquad w=x_i-x_j,\quad i\ne j.
 \label{eq:four-point-cone}
\end{equation}

\begin{proposition}[Quantitative positive margin]
\label{prop:four-point-positive-cone}
Under \eqref{eq:four-point-shape-class} and
\eqref{eq:four-point-cone}, every ring has angular factor greater
than four, and
\par\noindent
\begin{minipage}{\linewidth}
\begin{equation}
 \mathcal B_4(\underline f)
 \ge c_4\vartheta^4\prod_i F_i>0,
 \qquad c_4=12\pi^{-8}10^{-16}.
 \label{eq:four-point-positive-margin}
\end{equation}
\popQED\par
\end{minipage}
\end{proposition}
\begin{proof}
Write $\omega=\arcsin(1/8)$.  The unoriented line of every edge
is at angle at most $\omega$ from $\mathbb Re$.  The angle
between the lines of any two edges is therefore at most
$2\omega$.  Two reflections with normal lines at angle $\psi$
give a rotation by $2\psi$ on their span and the identity on its
orthogonal complement.

For an orthogonal map $O$ put
$a(O)=\sup_{|v|=1}\angle(v,Ov)$.  The spherical triangle
inequality gives $a(O_1O_2)\le a(O_1)+a(O_2)$.  Grouping the
four reflections into two pairs therefore gives
$a(T)\le8\omega$.  For the closed ring in
\cref{prop:four-point-closed-ring}, $a(T)=\alpha$: in the rotating
plane the maximal displacement angle is exactly $\alpha$, and a
fixed-plane component cannot increase it.

Three applications of the double-angle identity yield
\begin{equation}
 \cos(2\omega)=\frac{31}{32},\qquad
 \cos(4\omega)=\frac{449}{512},\qquad
 \cos(8\omega)=\frac{70529}{131072}>\frac12.
 \label{eq:four-point-angle-margin}
\end{equation}
The first two positive cosines, starting from
$0<\omega<\pi/2$, put $8\omega$ in $(0,\pi)$.
Consequently $8\omega<\pi/3$ and
\begin{equation}
 2+4\cos\alpha\ge\frac{136065}{32768}>4.
 \label{eq:four-point-angular-factor-floor}
\end{equation}
Any two points in distinct supports have distance at most $10d$:
choose points realizing their minimum distance at most $6d$ and
add the two support diameters, each at most $2d$.
Each of the three rings thus contributes at least
$4\pi^{-8}(10d)^{-16}$.  Integrating the resulting bound and
using the four mass floors proves
\eqref{eq:four-point-positive-margin}.
\end{proof}

There is an independent matrix check on the positive margin.  If
$\psi$ is the angle between the lines of $u$ and $v$, then
$\|R_u-R_v\|_{\rm F}=4\sin\psi$.  To see this, choose
coordinates in which the two normals lie in a coordinate plane.
The product $s_us_v$ rotates that plane through $2\psi$;
on $\Lambda^2$ it has two fixed coordinates and two copies of that
rotation.  Taking the Frobenius norm gives the stated identity.
Under \eqref{eq:four-point-cone}, telescoping four orthogonal
factors gives $\|R_{u_1}\cdots R_{u_4}-I\|_{\rm F}\le2$.
For a real orthogonal six-dimensional matrix $A$,
$\operatorname{tr}A=6-\|A-I\|_{\rm F}^2/2$.
This gives the weaker but still uniform floor four without the
closed-ring eigenvalue calculation.

The nearly collinear class in the four-point assertion satisfies
the cone condition.  Indeed, if
$\operatorname{supp}f_i\subset B(c_i,d/16)$ and the $c_i$ lie on
$\mathbb Re$, the transverse part of $x_i-x_j$ has length at most
$d/8$; the separation in \eqref{eq:four-point-shape-class}
makes this at most $|x_i-x_j|/8$.
This class is nonempty with explicit constants.  Take
$\psi(x)=(1-|x|^2)_+^2$ and
$f_i(x)=\psi(16(x-c_i)/d)$, with consecutive centres separated
by $5d/4$ along the line.  Then
\begin{equation}
 \int\psi=\frac{\pi^2}{12},\qquad
 \operatorname{Lip}\psi=\frac8{3\sqrt3},\qquad
 0<\vartheta\le
 \min\left\{\frac{\pi^2}{12\cdot16^4},\frac{3\sqrt3}{128}\right\}
 \label{eq:four-point-explicit-shapes}
\end{equation}
suffices for all the shape conditions.  The first integral follows
by radial integration; the gradient norm is the maximum of
$4r(1-r^2)$ on $[0,1]$.  The smallest support separation is
$9d/8$, and the largest minimum separation is less than $6d$.

\subsubsection{A negative transverse configuration}

\begin{proposition}[A negative box with an explicit neighbourhood]
\label{prop:four-point-negative-box}
Let $h>0$ and let $e_1',e_2',e_3'$ be an orthonormal triple.
Put
\begin{equation}
 x_1^0=y,\quad x_2^0=y+he_1',\quad
 x_3^0=y+h(e_1'+e_2'),\quad
 x_4^0=y+h(e_1'+e_2'+e_3').
 \label{eq:four-point-negative-vertices}
\end{equation}
Then
\begin{equation}
 k_4(\underline x^0)=-\frac7{36\pi^8}h^{-16}.
 \label{eq:four-point-negative-value}
\end{equation}
If $|x_i-x_i^0|\le h/200000$ for all $i$, then
\begin{equation}
 k_4(\underline x)\le-\frac7{72\pi^8}h^{-16}.
 \label{eq:four-point-negative-neighbourhood}
\end{equation}
Thus nonzero nonnegative tests supported in these four balls have
a strictly negative integrated box.
\end{proposition}
\begin{proof}
Isometries and homogeneity reduce the calculation to $h=1$,
$y=0$ and the first three coordinate vectors.  Every reflection
fixes the fourth coordinate.  Direct multiplication on the
three-dimensional span gives the following exact values:
\begin{equation}
 \begin{array}{c|c|c|c}
 \text{cycle}& (|w_1|^2,|w_2|^2,|w_3|^2,|w_4|^2)
       &\operatorname{tr}_{\mathcal V}\Lambda^2T
       &\prod_j|w_j|^{-4}\\ \hline
 1234&(1,1,1,3)&-2&1/9\\
 1243&(1,2,1,2)&0&1/16\\
 1324&(2,1,2,3)&4&1/144
 \end{array}
 \label{eq:four-point-negative-table}
\end{equation}
For the first row, the first three reflections multiply to
$-I$ on that span, so the product has eigenvalues $1,-1,-1$
there.  For the second row the product is a three-cycle of the
coordinate axes.  For the third row its trace on
$\mathbb R^4$ is three, giving angular factor four by
\eqref{eq:four-point-exterior-trace}.  All entries may also be
obtained from the rational formula
$s_w=I-2ww^{\mathsf T}/(w^{\mathsf T}w)$, so no choice of square
root affects the calculation.  The sum is
$-2/9+4/144=-7/36$, proving
\eqref{eq:four-point-negative-value}.

For completeness the neighbourhood can be chosen without a
compactness argument.  The normal-form calculation above also
gives $\|R_u-R_v\|\le2\sin\psi$.  Differentiating
\eqref{eq:four-point-kernel-reflection} in a unit direction gives
\begin{equation}
 \|\mathsf K(z)\|\le\pi^{-2}|z|^{-4},\qquad
 \|\partial_v\mathsf K(z)\|\le6\pi^{-2}|z|^{-5}.
 \label{eq:four-point-kernel-derivative}
\end{equation}
If $|x_i-x_i^0|\le\rho h$ and $\rho\le1/20$, every edge
on the straight path between the configurations has length at
least $9h/10$.  On this path the kernel norm is less than
$2\pi^{-2}h^{-4}$ and its directional derivative norm less
than $11\pi^{-2}h^{-5}$.  Each edge moves by at most $2\rho h$.
Telescoping a product of four factors, using
$|\operatorname{tr}A|\le6\|A\|$, bounds the change in all three
rings by
\begin{equation}
 |k_4(\underline x)-k_4(\underline x^0)|
 \le12672\rho\pi^{-8}h^{-16}.
 \label{eq:four-point-negative-continuity}
\end{equation}
At $\rho=1/200000$ this is smaller than
$7\pi^{-8}h^{-16}/72$.  This proves
\eqref{eq:four-point-negative-neighbourhood}.  Integration against
$\prod_i f_i\ge0$ proves the last assertion, more precisely with
upper bound $-7h^{-16}\prod_i\int f_i/(72\pi^8)$.
\end{proof}

The restriction to a nearly collinear class therefore has
mathematical content.  Positivity of the covariance operator does
not imply positivity of every separated fourth cumulant of its
quadratic density.

\subsubsection{The finite Gaussian coefficient and the $\operatorname{SU}(2)$ normalization}

Let $\xi^1,\ldots,\xi^q$ be independent centred Gaussian vectors
with the same real positive semidefinite covariance $C$ on a
finite coordinate set.  For real diagonal matrices $D_i$ set
$Q_i=\frac12\sum_{a=1}^q(\xi^a)^{\mathsf T}D_i\xi^a$.
On a neighbourhood of the origin,
\begin{equation}
 \log\mathbb E\exp\Big(\sum_i z_iQ_i\Big)
 =\frac q2\sum_{m\ge1}\frac1m
       \operatorname{tr}\left(C\sum_i z_iD_i\right)^m.
 \label{eq:four-point-gaussian-log}
\end{equation}
The formula follows by diagonalizing
$C^{1/2}(\sum_i z_iD_i)C^{1/2}$ at real parameters and then
using the convergent matrix logarithm near zero; a singular $C$
is handled on its range.  Taking the coefficient of
$z_1z_2z_3z_4$, and grouping the 24 words by cyclic rotation and
reversal, gives
\begin{equation}
 \kappa(Q_1,Q_2,Q_3,Q_4)
 =q\sum_{\varsigma\in\mathcal R}
       \operatorname{tr}
        (CD_{\varsigma(1)}CD_{\varsigma(2)}
         CD_{\varsigma(3)}CD_{\varsigma(4)}).
 \label{eq:four-point-gaussian-cumulant}
\end{equation}
Each unoriented ring has eight words, canceling $2m=8$.
Replacing $C$ by $x^{-1}C$ multiplies this by $x^{-4}$.
Centering the $Q_i$ does not change this mixed cumulant.

For the Wilson observable of $\operatorname{SU}(2)$ use Pauli matrices
normalized by $\frac12\operatorname{tr}(\sigma^a\sigma^b)
=\delta_{ab}$.  For real $v\in\mathbb R^3$,
\begin{equation}
 1-\tfrac12\operatorname{tr}\exp(i v\cdot\sigma)
 =1-\cos|v|=\tfrac12|v|^2+O(|v|^4).
 \label{eq:four-point-su2-normalization}
\end{equation}
Hence the Gaussian quadratic reference has $q=3$, with exactly
the factor $1/2$ used above.  If each first source derivative has
coefficient $\zeta$, its four-point leading term is
\begin{equation}
              3\zeta^4x^{-4}\mathcal B_4(\underline f).
 \label{eq:four-point-su2-leading-term}
\end{equation}
There is no additional factor $1/2$ in this four-point formula.
The quartic Wilson Taylor term in
\eqref{eq:four-point-su2-normalization}, the interacting density,
and the actual source remainders have not been discarded by a
claim that their cumulants are small.  Estimating their combined
contribution is the separate comparison problem.

For comparison with three points, the Hodge star satisfies
$*R_u*=-R_u$.  Inserting $**=I$ between the factors and taking a
trace shows that every odd ring of the continuum kernel vanishes.
This is a kernel identity at separated points, rather than an
assertion that the quadratic density is a Gaussian random field.

\subsubsection{Uniform comparison with the free lattice box}

Let $a=\ell/N$, $N\in\mathbb N$, and let $P_{\rm ex}$ be the
orthogonal projection onto the exact lattice two-cochains.
Plaquettes $p=(y,P)$ have barycentres $x(p)=y+ac_P$, with
$c_P=(e_\mu+e_\nu)/2$ for $P=(\mu,\nu)$.
For any real symmetric covariance $C(p,q)$ define the lattice box
\begin{equation}
 \mathcal B_{4,a}[C](\underline f)
 =\sum_{\varsigma\in\mathcal R}\sum_{p_1,\ldots,p_4}
      \prod_i f_i(x(p_i))
      \prod_{j=1}^4 C(p_{\varsigma(j)},p_{\varsigma(j+1)}).
 \label{eq:four-point-lattice-box}
\end{equation}
In particular the fourth cumulant of the finite free Wilson
quadratic reference is $3\mathcal B_{4,a}[P_{\rm ex}]$.
There is no missing Riemann factor in this expression:
$P_{\rm ex}(p,q)=a^4\mathsf K_{a,\ell}(p,q)$, so a product of
four covariance entries contributes $a^{16}$.

\begin{proposition}[Uniform lattice error on the full shape class]
\label{prop:four-point-uniform-lattice}
Suppose \eqref{eq:four-point-shape-class} holds on a common short
lift of the four supports to $\mathbb R^4$, with $a\le d/100$
and $\ell\ge200d$.  There is a numerical $C$, depending only
on the constant in \eqref{eq:lcc-main}, such that
\begin{equation}
 \left|\mathcal B_{4,a}[P_{\rm ex}](\underline f)
                 -\mathcal B_4(\underline f)\right|
 \le C\left\{(a/d)^{1/2}+(d/\ell)^4
                    +\frac{a}{\vartheta d}\right\}\prod_iF_i.
 \label{eq:four-point-uniform-lattice-error}
\end{equation}
The bound is uniform in positions, directions, shapes and the
number of lattice sites.  The cone condition is unnecessary for
this error estimate.
\end{proposition}
\begin{proof}
Write the lattice box using four factors $a^4\mathsf K_{a,\ell}$.
On distinct supports every relevant distance lies between $d$
and $10d$.  The comparison in \eqref{eq:lcc-main}, entrywise
and hence in six-dimensional operator norm, gives
\begin{equation}
 \|\mathsf K_{a,\ell}(p,q)-\mathsf K(x(p)-x(q))\|
 \le C d^{-4}\{(a/d)^{1/2}+(d/\ell)^4\}.
 \label{eq:four-point-entry-error}
\end{equation}
All covariance factors are bounded by $Cd^{-4}$.
Replacing the four factors one at a time therefore costs at
most $C d^{-16}\{(a/d)^{1/2}+(d/\ell)^4\}$ per term,
before the $a^{16}$ factor and the four test functions.
For each orientation, the number of grid points in a ball of
radius $d$ is at most $C(d/a)^4$, by placing disjoint mesh cubes
about the points.  Summation in four variables and in the finitely
many orientations gives the first two error terms in
\eqref{eq:four-point-uniform-lattice-error}.

The remaining sum still samples the test functions and kernels at
orientation-dependent barycentres.  Move $y_i+ac_{P_i}$ to $y_i$,
one factor at a time.  Along these movements a nonzero test
function stays in its support enlarged by $Ca$, and distinct
supports remain at distance at least $d/2$.  The test-function
cost is at most $CaF_i/(\vartheta d)$, and
\eqref{eq:four-point-kernel-derivative} bounds each kernel cost by
$Ca d^{-5}$.  The same grid count bounds the entire change by
$Ca\prod_iF_i/(\vartheta d)$.

At the common grid sites $y_i$, summing the four orientation
indices gives precisely the matrix trace in
\eqref{eq:four-point-box-density}.  The result is
$a^{16}\sum_{y_1,\ldots,y_4}k_4(\underline y)\prod_i f_i(y_i)$.
This last Riemann sum has the same error bound.  In fact the
integrand has derivative at most
$C\vartheta^{-1}d^{-17}\prod_iF_i$ on the product of the
enlarged supports, whose sixteen-dimensional volume is
$Cd^{16}$.  Its oscillation in a mesh cell is at most the
derivative bound times $Ca$.  Summing cell errors gives
$Ca\prod_iF_i/(\vartheta d)$.
No diagonal singularity is crossed: the product of the four tests
vanishes in a neighbourhood of every coincidence diagonal, so the
integrand extends by zero there.  These three replacements prove
the claim.
\end{proof}

\begin{corollary}[A uniform lattice sign floor]
\label{cor:four-point-lattice-positive-floor}
There is a numerical $c>0$ such that the cone class satisfies
\begin{equation}
 \mathcal B_{4,a}[P_{\rm ex}](\underline f)
 \ge\tfrac12c_4\vartheta^4\prod_iF_i
 \label{eq:four-point-lattice-positive-floor}
\end{equation}
whenever $a/d\le c\vartheta^8$ and $d/\ell\le c\vartheta$.
\end{corollary}
\begin{proof}
Choose $c$ so that the right side of
\eqref{eq:four-point-uniform-lattice-error} is at most
$c_4\vartheta^4\prod_iF_i/2$.  This is possible because the
three displayed errors are then bounded by constants times
$c^{1/2}\vartheta^4$, $c^4\vartheta^4$ and
$c\vartheta^7$, respectively, and $\vartheta\le1$.
The same choice can enforce $a\le d/100$ and $\ell\ge200d$.
Apply \cref{prop:four-point-positive-cone}.
\end{proof}

For any fixed nonzero nonnegative $C^1$ tests supported in the
negative balls of \cref{prop:four-point-negative-box}, the same
comparison gives a strictly negative free lattice box at a
sufficiently fine mesh and sufficiently large torus.  To apply the
uniform version choose $d=9h/10$ and a fixed positive
$\vartheta$ smaller than all four mass and inverse Lipschitz
ratios.  The supports then meet the separation and size
conditions in \eqref{eq:four-point-shape-class}.  The continuum
negative margin is independent of the mesh, whereas the error in
\eqref{eq:four-point-uniform-lattice-error} tends to zero.

\subsubsection{Background covariances and perturbed profiles}

The continuum sign also applies to the background covariance once
its Gaussian complement has the required decay.  The exact
decomposition and the decay estimate are different facts.
Let $E$ be the space of exact two-cochains and let $A:E\to W$
be the actual linear coarse flux map.  For a standard Gaussian
$F$ on $E$ put $N=E\cap\ker A$ and $Y=E\ominus N$.
The restriction of $A$ to $Y$ is injective, so $AF$ determines
the $Y$ component of $F$.  Its $N$ component is an independent
standard Gaussian.  Thus
\begin{equation}
 \operatorname{Cov}(\mathbb E[F\mid AF])=P_Y,
 \quad \operatorname{Cov}(F-\mathbb E[F\mid AF])=P_N,
 \quad P_{\rm ex}=P_Y+P_N.
 \label{eq:four-point-background-projections}
\end{equation}
This finite-dimensional identity holds for the specified flux map;
it contains no spatial decay assertion.  To identify $P_Y$ with
the background covariance of a connection integral, one must also
verify that the actual constraint induces this flux map, that its
gauge integration has the prescribed normalization, and that the
closed and zero modes have been treated consistently.  The Gaussian
projection identity does not replace those identifications.

If, in addition, its fluctuation complement satisfies
\begin{equation}
 |P_N(p,q)|\le C(a/b)^4
                   e^{-\delta|x(p)-x(q)|/b},
 \quad |x(p)-x(q)|\ge b,
 \label{eq:four-point-needed-complement-tail}
\end{equation}
then the proof of \cref{prop:four-point-uniform-lattice} gives
\begin{equation}
 \begin{split}
 |\mathcal B_{4,a}[P_Y](\underline f)
                      -\mathcal B_4(\underline f)|
 &\le C\left\{(a/d)^{1/2}+(d/\ell)^4
               +\mathfrak e(d/b)+\frac a{\vartheta d}\right\}
               \prod_i F_i,\\
 \mathfrak e(t)&=\sup_{u\ge t}u^4e^{-\delta u},
 \qquad d\ge b.
 \end{split}
 \label{eq:four-point-background-comparison}
\end{equation}
Indeed $a^{-4}|P_N(p,q)|$ is bounded by
$C|x(p)-x(q)|^{-4}\mathfrak e(d/b)$ at the relevant distances.
Only the first covariance-replacement step changes. For the
actual flat witness, take the flux map $Q^{\mathrm{pl}}_{K-s}$
and $b=a_s$. The flat conditioning of
\cref{sec:short-distance-the-short-distance-two-point-law-and-its-asymptotic-freedom}
identifies $P_N$ with $C^{\mathrm{fl}}_{K,s}$. The
\hyperlink{witness-free-complement}{sum of the finer slices}
gives the tail at the separated distances used here, including
the enlarged profiles at distance at least $\mathsf d/2$.
The \hyperlink{flat-curvature-gauge-comparison}{flat gauge comparison}
identifies the curvature law of each slice with the flat step law.
Other flux maps require their own complement estimate.

The form is also stable under the small changes of profiles
produced by Gaussian extraction.  Retain $a\le d/100$.
Suppose a covariance satisfies
$|C(p,q)|\le C_0a^4d^{-4}$ between the distinct enlarged
supports.  Let $h_i(p)$ depend on the plaquette orientation if
necessary, with $|h_i(p)|\le\varepsilon F_i$, supported within
distance $d/8$ of $\operatorname{supp}f_i$, and
$0\le\varepsilon\le1$.  The enlarged supports still have
separation at least $3d/4$ and volume $O(d^4)$.  The absolute
grid-mass bound
$a^4\sum_p|h_i(p)|\le C\varepsilon F_i d^4$
uses only that support bound and the point count.  Expanding by
four-linearity, the fifteen nonempty choices of perturbed
arguments give
\begin{equation}
 |\mathcal B_{4,a}[C](f_1+h_1,\ldots,f_4+h_4)
                  -\mathcal B_{4,a}[C](\underline f)|
 \le C\{(1+\varepsilon)^4-1\}\prod_iF_i
 \le C'\varepsilon\prod_iF_i.
 \label{eq:four-point-profile-stability}
\end{equation}
No differentiability or positivity of the $h_i$ is needed.
Orientation-dependent perturbations therefore cause no
additional covariance or Riemann-sum assumption.

\paragraph{The normalized comparison of four bounded insertions.}
The four-point witness may require the larger room
\(C_W^\square(s)=c_W C_W(s)\), with a fixed \(c_W\ge1\) chosen
for the shape class. In this comparison write \(s\) for the depth
selected with that room and \(k_0=K-s\), \(x=x_{k_0}\),
\(\rho_{\mathsf d}=L^s\mathsf d\).
At the witness take the four windowed plaquette reads $G_i$,
with profiles $\phi_i$ supported in the allowed enlargements.
Use a segment of $r$ steps and a single nonwrapping box for
the four supports at each step. At the first input
level \(k_0\), include the full datum box of each
native localized witness minimizer, with its four
test layers. Their size in witness units is included
in \(C\rho_{\mathsf d}\). At a later input level
\(j-1\), use one localized minimizer cube around the
union of the four supports, with a margin of order
\(D\) in input units. Use one prescribed rooted axial
tree for the common frame. Its root and all its
paths belong to the read. A simultaneous change of
gauge leaves only the rotation at this common root
in the four transformed arguments.

%
For the input fluctuations, define \(S_j\) from
these full reads by adjoining all free coordinates
within the fixed double-block reach of the one-step
pivot and reconstruction maps. The exact dependence
on \(S_j\) follows from the finite range of those
maps. For the comparison integral \(I_j\), add a
further \(D\)-margin to a cube containing these
coordinates. With \(O_j\) the free exterior, this
gives \(S_j\subset I_j\), \(\#S_j\le n_j=\#I_j\),
and \(d(S_j,O_j)\ge D\) in input mesh units.
In particular
\[
 n_j\le C(\rho_{\mathsf d}L^{-(j-1-k_0)}+D+C_L)^4.
\]
These are the read-set conditions of the
normalized product-window estimate.

Use a separate neighborhood for the clean clause.
It contains \(I_j\), every native test whose reads
can reach it, their full row carriers, and the
operation margins. If an operation touches one
of these reads, its immediately preceding live
component is within the diameter of the operation
component. The enlargement by
\((200\check R_j+1)M\) in the appropriate level
units covers this possibility even when the old
live cells do not intersect \(I_j\). These layers
enter the clean-history count. The conditional
comparison integral retains its own box and
its distance \(D\) to the exterior.
The condition is imposed before that step's healing: no
creation or operation meets this enlargement, and no live
component meets it on either side of the step. Use the
compensated standing representation. The coupling lag
has no term supported wholly inside the box, since the
profiles of healed regions have been redrawn and the
standing large regions miss the padded neighborhood.
A lag term reaching the box through an exterior trace
uses the fixed-exterior old-origin oscillation estimate.

Fix the coarse field, the older stored slots and the
exterior free coordinates after the ordinary gauge draw.
Use the product-window calculation on the union of all
four reads. An affected pivot is reconstructed on its full
double block, including when its bond lies outside \(I_j\).
For every query that reads the changed coordinates the
interior window implication gives the approximate and
old density tests throughout the literal product of
balls. No operation or its guards meet these carriers.
For an exterior query the complete read is fixed.
Induction along the original order of the queries gives
the same exterior outcomes throughout the product,
including when the next query depends on an earlier
answer. Thus at fixed compatible exterior the selector
is constant on this product.

On a clean standing part the numerator with insertion
and the uninserted denominator therefore use the same
interior window and the same exterior measure.
Writing \(Z_{\mathsf e}\) for the uninserted interior
integral and \(m(d\mathsf e)\) for the exterior measure
with its retained weights, the normalized conditional law
on the box is averaged with measure proportional to
\(Z_{\mathsf e}\,m(d\mathsf e)\). A merge or an operation
on the exterior takes a further positive mixture of
these quotients. The pointwise clean quotient gives the
reference branch with the original weights for these
allowed strings. On this branch take the same cutoff
and the same flat slices for all four profiles. Use the
common-frame estimate for the nested localized readouts
on their union, including the separate first witness
minimizers. The hybrid through
the adjacent-level coupling is defined with the bounded
function $\exp(\sum_i\lambda_i G_i)$, with independent source
coordinates. Setting the top field to zero yields a single
Gaussian determinant independent of the terminal datum and
of the history. On the polydisc
\[
 |\lambda_i|\le\Lambda_i=
       \frac{c x}{4\rho_{\mathsf d}^4p_0^2\|\phi_i\|_\infty},
\]
the diagonal covariance bound and the count of plaquettes
give a uniform trace norm bound on the quadratic Gaussian
exponent. Zero profiles can be removed. The same bound on
a larger disc gives a uniform positive lower bound for the
determinant's modulus and the second moment needed to remove
the Gaussian windows.

Apply the one-step comparisons on this disc to the positive
law of the original segment. The nonlinear readout error
sums with the geometric factor $L^{-2(j-1-k_0)}$. The tilt,
covariance and exterior errors sum in the bounds for the
normalized box law. Charge a finer history restriction at
the steps where its density changes that law. The
\hyperlink{witness-top-native-tests}{native top tests}
give the bound on the top field for all four read sets
on the common admitted box. This field is present in the
Gaussian model as a shift; the linear
cross term is odd under simultaneous reflection of the
Gaussian variables. Pairing the reflected points makes its
effect quadratic in the shift. This and the top square
cost $O(L^{-4r})$ on the polydisc. In the cumulant
calculation these errors correspond to open chains with
two ends at the top curvature. The normalization by the
uninserted box mass is made at each step. In particular
decay of the free covariance alone does not give the
exterior error, which uses the decay of the response of
the tilted product-window law to the exterior force.

For an unconditional comparison, average the generating
functions before taking logarithms. A clean sum at terminal
datum $W$ can also use the normalized mixture of the clean
class laws, with weights $d\mathsf Q(h,\cdot)/d\mathsf Q_{\rm cl}(W)$
on the clean classes. Here $\mathsf Q_{\rm cl}$ is the sum of
their terminal measures. This gives a common determinant
for the clean mixture when the clean comparison is uniform
over its constituents. It uses a refinement for which
cleanness is a property of each class. The actual exponent
and the density of the refined classes have to be kept in
the normalized comparison. On an exceptional part the
$G_i$ still have their windowed sup bounds. Hence the
probability of that part, times the bounded insertion
estimate, suffices here. The physical first derivatives
instead use \eqref{eq:f3-live-jet-moments}.

For the unconditional comparison use the adjacent positive
laws \(\pi_j\) of
\hyperref[f2:witness-generating-function]{the two-insertion calculation},
now on the union of the four complete read sets. Take
\(Y_j\) to be the bounded generating function with all
slices through \(j\) replaced by flat ones. The same law
\(\pi_j\) supplies the two marginals in
\(\mathbb E_{\pi_j}(Y_{j-1}-Y_j)\). At a compatible good
exterior, where the selector equals one throughout the
product window, its conditional bound is
\[
 C\left[L^{-2(j-1-k_0)}
        \big(\rho_{\mathsf d}^2x^{-1/2}p_0+e^{-c\Delta_j}\big)+T_j\right].
\]
Here \(T_j\) is the normalized box error of the two-insertion
calculation, evaluated with the four-read coordinate count. The
bound for changing the plaquette field in the exponent has
the same constant up to a fixed factor because
\(\sum_i|\lambda_i|\|\phi_i\|_\infty\le c x/(\rho_{\mathsf d}^4p_0^2)\).
The sum over steps compares the two original endpoint
expectations. On a bad current part of mass \(p_j\), both
integrands have a bounded extension, so it contributes
at most \(Cp_j\). The top-field error, the deleted Gaussian
windows and the localized-to-flat comparison give, in all,
\[
 \delta\le C\left[\rho_{\mathsf d}^2x^{-1/2}p_0+
  \max_j e^{-c\Delta_j}+\sum_jT_j+L^{-4r}
  +\sqrt t+p_0^{-2}e^{-cd_*}+\sum_jp_j+p_{\rm top}\right].
\]
The notation \(t,d_*,p_{\rm top}\) has the same meaning as in
the two-insertion calculation. This argument uses the
uninserted exterior weights at each step. To use this
bound on the adjacent laws, we must estimate the box
errors on the actual four-read domains and the masses
of the incompatible exteriors.

There is also a comparison using the clean restriction of the
original run. On a clean class use the reference branch of
weight $P_0$ from the pointwise clean quotient, with the stored
variables and exterior datum kept in the disintegration.
The box and its padding are the same in the restriction and
the comparison. Thus each reference conditional expectation
uses the original positive weights for the allowed strings.
Compose the carried densities on the complete retained
states, and project to the terminal datum at the end.
Replace one slice at a time in these conditional expectations.
The remaining outer composition has sup norm at most one,
so the conditional errors add. For the fixed conditional laws,
prove the disc estimate on a countable dense set and extend
by analyticity, which gives a common exceptional null set
for the source parameters. Write $Z^0$ for the resulting
deterministic flat determinant and $\delta_{\rm ref}$ for
the uniform relative error of this reference branch.
If $\nu(\mathrm{cl})$ is the full law's mass of clean classes,
averaging the reference and remaining branches gives
\[
 \left|\frac{\mathbb E\exp(\sum_i\lambda_iG_i)}{Z^0}-1\right|
 \le C\bigl(\delta_{\rm ref}+1-P_0\nu(\mathrm{cl})\bigr).
\]
Indeed the reference weight is $P_0\nu(\mathrm{cl})$, and on
the complementary branch the uniform bounds on $G_i$ and
the lower bound on $|Z^0|$ control the ratio. This is a bound
on the unconditional generating function before taking its
logarithm. The hypotheses of this route are the support-adapted
clean restriction, the verified bracket majorants, and the
conditional comparisons on the common product windows.

For the clean class sum, the excluded histories are counted
by live cells meeting the padded box from the lower waiting
band through the end of the segment. A union bound over this
band and the $r$ steps costs at most
$C(1+r+\check R_{k_0})(\rho_{\mathsf d}+M\check R_{k_0})^4
 L^{4r+C\check R_{k_0}}$ times the normalized bound per cell.
Here the factor in $L$ converts the count from the terminal
scale to the finest scale in the band. The healed histories
outside this criterion are included in the standing sum.
For an estimate by adjacent-step selectors, the corresponding
count uses their actual current readsets. In either case
spend the rarity reserve after these counts and after the
source loss below.

If $\delta$ bounds the resulting relative error of the
averaged generating function on the polydisc and $\delta<1/2$,
the two logarithms differ by $C\delta$. Mixed Cauchy gives
an error at order four of at most
\[
 C\delta\prod_i\Lambda_i^{-1}
  \le C\delta\,x^{-4}\rho_{\mathsf d}^{16}p_0^8
                   \prod_i\|\phi_i\|_\infty .
\]
The comparison therefore requires that
$\delta\rho_{\mathsf d}^{16}p_0^8$ tend to zero. For instance,
if the old-origin removal has
\(\epsilon_j^{\mathrm{old}}\le C n_j(\log x)^{\mathfrak p_0-q}\)
on the four-read box and \(n_j\le C\rho_{\mathsf d}^4\), its
contribution after Cauchy is bounded by
\[
 C r\,\rho_{\mathsf d}^{20}p_0^8
              (\log x)^{\mathfrak p_0-q}.
\]
For \(r=O(\log\log x)\) and
\(\rho_{\mathsf d}\asymp(\log x)^{2\mathfrak p_0}\), this tends
to zero when \(q>49\mathfrak p_0\).

On the admitted clean reference windows take \(D=A\log x\)
and \(r=\lceil B\log_L\log x\rceil\). The native margins
give \(\Delta_j,d_*\ge c_0\min\{D,\check R_{k_0}\}-C_0\),
and \(\check R_{k_0}/\log x\) grows without bound.
Hence the geometric errors, including \(n_jb_j^2e^{-cD}\),
beat the Cauchy factor. In \(T_j\), the cubic density term
costs \(Cn_j x^{-1/2}p_0^3\), the covariance term is bounded
by \(Cn_j(\rho_{\mathsf d}x^{-1/2}p_0+e^{-cD})\),
and \(t_j\le Cn_j e^{-cb_j^2}\) with \(b_j\) comparable
to \(p_0\) on the segment. These and the square root of the
total window error also beat every required logarithmic
power. Removing the top field costs \(CL^{-4r}\), so take
\(4B>40\mathfrak p_0\) for
\(\rho_{\mathsf d}^{16}p_0^8L^{-4r}\to0\).
The loss of the clean reference branch contains both
\(1-\nu(\mathrm{cl})\) and \(1-P_0\le r\epsilon_{k_1}\).
The padded cell count and the rooted bound on the deleted
histories bound their sum by
\[
 C(1+r+\check R_{k_0})^C
      (\rho_{\mathsf d}+M\check R_{k_0})^4
      L^{4r+C\check R_{k_0}}e^{-cp_0}.
\]
Since \(\check R_{k_0}=o(p_0)\), this loss tends to zero
after multiplication by the Cauchy factor. The terminal
test reads are included in the clean neighborhood.
With the fixed-torus age estimate, use the unprolonged run and
write \(h=g^{-2}\). The last lattice has \(\ell^4\) sites and
\(x_K\le h+B_\sharp\). In the survival term of
\cref{prop:source-domains-4-14} a sufficient depth floor is
\[
 s_{\rm vol}(g,\ell)=\max\left\{0,
   \left\lceil \frac{\mathsf B(h+B_\sharp)\ell^4}{\bar c_a}\right\rceil-1
   \right\}.
\]
At \(k=k_0+r\), the condition \(s-r\ge s_{\rm vol}(g,\ell)\)
gives \(\mathsf B(x_K)\ell^4-\bar c_a(K-k+1)\le0\).
The survival term retains \(e^{-\bar c_aa-2p_0(g_k)}\), and the
healed term already has age decay without a volume factor.
For \(\bar c_a\ge17\log L+\log 2\) the weighted age sum is
geometric and gives the required bound for \(Q_{17}(k)\).
All earlier reads meet the same depth floor. At fixed \(g\),
the segment has \(r\le s/2\) above a fixed depth, so
\(s\ge2s_{\rm vol}(g,\ell)\) there suffices. The witness rule,
whose room is at least one, gives this depth for sufficiently
small separation. In particular a cutoff of the form
\(\mathsf d\le e^{-C_g(1+\ell^4)}\) suffices for this age
condition, in addition to the formation and clean-box conditions.
These choices require the stated nonwrapping room and
the native density and marginal estimates on the
common comparison boxes.

If the comparison
is also made at fixed clean class and datum, the same disc
controls the proper conditional cumulants relative to the
deterministic reference values. On the clean part the
multi-block terms in total cumulance contain products of at
least two conditional errors. They are absent when comparing
the averaged generating functions directly.

The determinant has covariance
$\sum_{j=k_0+1}^{k_0+r}x_{j-1}^{-1}\mathcal C_{(j)}$.
On separated supports the slice bound gives a summable
absolute kernel. Replacing $x_{j-1}^{-1}$ by $x^{-1}$
then costs $C(B_\sharp r+c_5)/x$ relative to the natural
smeared scale. Each term in the omitted ring sum has at
least one coarse edge, which costs $L^{-4r}$ in witness units.
Since $L^r\ge\rho_{\mathsf d}$ by the segment choice, any further
coarse edge costs at most the separated edge bound
$C\rho_{\mathsf d}^{-4}$. Counting the four supports gives
$C\rho_{\mathsf d}^4L^{-4r}$ at the same relative scale.
These two comparisons are made on the
ring itself and do not pay the source-polydisc loss. To
compare the enlarged witness profiles with the original
profiles in the ring, use \eqref{eq:four-point-profile-stability}
with relative perturbation \(CM/(\vartheta\rho_{\mathsf d})\).
Indeed the enlargement is at most \(CMa_s\), and
the slope bound is in units of \(\vartheta\mathsf d\),
with \(\mathsf d/a_s=\rho_{\mathsf d}\).
Choose the room so that \(CMa_s\le\mathsf d/8\)
and the relative perturbation is at most one.
The witness comparison and the ring use the same
enlarged profiles before this replacement. To
pay that loss for the $O(L^{-4r})$ term in the generating
function, take $L^{4r}$ larger than
$\rho_{\mathsf d}^{16}$ times the prescribed power of $\log x$.
A segment of length $O(\log\log x)$ suffices in the stated
room, with its constant chosen for the error budget. Fix the
reference coupling first. As the witness depth tends to infinity,
\(x\asymp\bar x_s\) and \(\bar x_s=O_g(1+s)\), so \(r\le s-1\)
and \(s-r\to\infty\). The padded neighborhoods have side at most
\(C(\rho_{\mathsf d}+D+M\check R_{k_0})L^{-(s-r)}\) in physical
units. Their polylogarithmic prefactor is dominated by the mesh
decay. Take the separation threshold small enough that this side
is below a fixed fraction of the torus side and that the stated
volume floor holds at \(s-r\). The resulting drift is \(o(x)\).
The box-law estimates must apply to the actual normalized laws
on these boxes.

%
%
The coupling in the four-point assertion uses the original
two-point depth, say \(s_2\). To compare, the room functions obey
\[
 \frac{C_W(u+t)}{C_W(u)}
 \le \bigl(1+\log(1+B_\sharp t)\bigr)^{2\mathfrak p_0}
 \qquad (u,t\ge0),
\]
since \(\bar x_u\ge e\) and \(\bar x_{u+t}=\bar x_u+B_\sharp t\).
Choose a fixed integer \(t_0\) so that \(L^{t_0}\) exceeds
\(c_W\) times the right side. By minimality of the witness
depths, \(s_2\le s\le s_2+t_0\): the original room condition
at \(s_2\) implies the enlarged one at \(s_2+t_0\).
Thus a required formation floor \(K-s\ge n\) follows from
\(K-s_2\ge n+t_0\). The two-sided flow bound gives
\(\lvert x_{K-s}-x_{K-s_2}\rvert\le B_\sharp t_0+c_5\).
Consequently
\[
 x_{K-s}^{-4}=g(\mathsf d)^8
     \bigl(1+O(x_{K-s_2}^{-1})\bigr).
\]
This transfers the leading size and the error scale of the
four-point comparison to the coupling used in the assertion.
The shape margin is taken from the ring at the chosen witness.

\subsubsection{The exact remaining interacting estimate}

This closes the sign, nondegeneracy and shape stability of the
continuum free box, and its comparison with the full free lattice
reference.  It supplies the corresponding background conclusion
whenever the explicitly named complement tail is proved.  With
the running law $x^{-1}\asymp(\log(1/d))^{-1}$ and
$\zeta\to1$, these bounds give size
$(\log(1/d))^{-4}$ for the leading four-point term.  They do
not prove that running law.

To obtain the interacting four-point assertion, the unresolved
comparison must identify the actual normalized history/block-field
law and establish, on its whole prescribed class,
\begin{equation}
 S_4^T(\underline f)
 =3\zeta^4x^{-4}\mathcal B_{4,a}[P_Y](\underline f)
            +\mathcal R_4(\underline f),\qquad
 |\mathcal R_4(\underline f)|
 \le\epsilon_4(d)\,x^{-4}\prod_iF_i,
 \quad \epsilon_4(d)\longrightarrow0.
 \label{eq:four-point-unresolved-interacting-target}
\end{equation}
For a relative error at most one half, the actual error must in
particular beat the fixed shape margin proportional to
$\vartheta^4$.  Every contribution belongs in
$\mathcal R_4$: nonquadratic Wilson terms, the interacting
conditional density, source jets, cutoff boundaries, fluctuations
not represented by $P_Y$, and the histories excluded from the
clean comparison.  The required constants must have the stated
cutoff, support-position, reference-volume and state uniformity;
the torus-only four-point claim must retain its torus scope.

The bounded windowed insertions use the probability of an
unclean history in the normalized comparison. The physical
first source derivatives can carry the age weight $L^{16a}$.
They use the weighted sum in \eqref{eq:f3-live-jet-moments}
or a direct bound for the four physical insertions.  Nor does a bound on
the interior Hessian determine the exterior force in the fixed
conditional density.  Those estimates concern the interacting
construction, and none is inferred from the sign of the free box.

% END PROOF EXPANSION four_point/four_point_kernel.tex


% BEGIN MODULE chapters/uniform_insertion_bound_v4.tex
\subsection{A uniform one-field insertion bound}
\label{ch:uniform-insertion}

An estimate for each correlation function does not automatically give
an estimate in the Hilbert norm of an arbitrary polynomial of fields.
The argument below obtains that stronger norm by repeated reflections
at a fixed lattice spacing, before taking a continuum limit. Its analytic hypotheses are the ordinary scalar source upper bound
and floor; no typed insertion estimate is required.

Assume the all-mirror invariance and reflection positivity of every
periodic Wilson DLR state, the ordinary torus source hypothesis in \eqref{eq:all-state-source-transfer-ST5},
and the fixed-layer rooted domain geometry in \eqref{eq:graded-source-bound-GS1}--\eqref{eq:graded-source-bound-GS3}. The
localization size is $M=L^{m_0}$, with $L$ odd. Torus limits use the
cofinal compatible family $\ell=2L^m$, $m\ge m_0$. Every scale cube
$Ma_s$ then tiles the torus, since
$\ell/(Ma_s)=2L^{m-m_0+s}$. Every fixed board of power-of-$L$ side,
and every fixed coarser such board, tiles all sufficiently large tori
in this family. No scale terms are omitted.

\subsubsection*{A raw-algebra reflection argument}

Fix the spacing $a=L^{-K}$ and a periodic state $\rho$. Take a board
cell $C$ of side $l=L^j\ge a$, centered at a lattice site, with lower
time face $t$. Its face planes are link mirrors: $l/a$ is odd.
Its middle time plane is a site mirror. Both are supplied by the
all-mirror theorem. For bounded finite-support functions above $t$ put
\begin{equation}
 \langle B,A\rangle_t=\rho(\overline{\theta_tB}A),
 \qquad \|A\|_t^2=\langle A,A\rangle_t.
 \label{eq:uniform-insertion-bound-UI1}
\end{equation}
For a bounded complex cell function $D$, let
$q_D=\limsup_kp_{l,\ell_k}(D)$ along a compatible cofinal sequence.
Define, on the raw half-space algebra,
\[
 T_DA=D\,\tau_{le_4}A,
 \qquad D^*=\overline{\theta_{t+l/2}D}.
\]
Reflection and translation invariance give the formal adjoint identity
$\langle B,T_DA\rangle_t=\langle T_{D^*}B,A\rangle_t$.
No bounded operator on a quotient has been assumed.

Put $K_D=T_{D^*}T_D$ and $m_j=\langle A,K_D^jA\rangle_t$.
Formal adjoints and reflection Cauchy--Schwarz imply
\begin{equation}
 m_{2r}=\|K_D^rA\|_t^2,\qquad
 m_{2r+1}=\|T_DK_D^rA\|_t^2,\qquad
 m_j^2\le m_{j-1}m_{j+1}.
 \label{eq:uniform-insertion-bound-UI2}
\end{equation}
For odd $j=2r+1$ the last inequality uses the two vectors
$K_D^rA,K_D^{r+1}A$; for even $j=2r\ge2$ it uses
$T_DK_D^{r-1}A,T_DK_D^rA$. Thus it is a raw semidefinite-form
calculation, not an application of a spectral theorem to an operator
whose boundedness is still to be proved. If $m_0=0$, Cauchy--Schwarz
already gives $m_1=0$. Otherwise log-convexity gives
\begin{equation}
 m_1^N\le m_0^{N-1}m_N. \label{eq:uniform-insertion-bound-UI3}
\end{equation}

The expression for $m_N$ consists of the left endpoint
$\overline{\theta_tA}$, a column of $2N$ alternating translated copies
of $D^*$ and $D$, and the right endpoint $\tau_{2Nle_4}A$.
Choose a coarse side $l'=L^rl$ large enough to contain $A$ in a
coarse half-space cell above $t$, spatially aligned with $C$.
The oddness of $L^r$ makes all coarse faces fine-board faces.
Take $N$ through sufficiently large multiples of $L^r$.

On this coarse grid retain fixed endpoint zones containing both
endpoint supports and, if convenient, a fixed number $c_A$ of the
adjacent column factors. Bound the complete endpoint functions by
their fixed-cutoff supremum norms. This yields a finite constant
$C(a,A,D)$ independent of $N$. In the middle, every occupied coarse
cell is a product of fine cell functions only. Fine-to-coarse
chessboard refinement and product domination in \eqref{eq:all-state-source-transfer-ST3} apply there.
Grid-preserving translates and reflected conjugates of $D$ have the
same fine seminorm. Therefore, for every $\epsilon>0$,
\begin{equation}
 m_N\le C(a,A,D)(q_D+\epsilon)^{2N-c_A}.
 \label{eq:uniform-insertion-bound-UI4}
\end{equation}
The volume limit is taken first for each fixed $N$. The estimate uses
no assertion $p(AD)\le\|A\|_\infty p(D)$: a mixed endpoint block,
if present, is bounded in its entirety by its supremum norm.
The endpoint zones contain only a bounded number of coarse cells;
their factor count and supremum bound stay fixed as the column grows.

Taking the $N$th root in \eqref{eq:uniform-insertion-bound-UI3}--\eqref{eq:uniform-insertion-bound-UI4}, then sending $N\to\infty$
and finally $\epsilon\downarrow0$, gives
\begin{equation}
 \|T_DA\|_t\le q_D\|A\|_t,
 \qquad
 |\rho(\overline{\theta_tB}D\tau_{le_4}A)|
       \le q_D\|B\|_t\|A\|_t.
 \label{eq:uniform-insertion-bound-UI5}
\end{equation}
The same order of limits proves this when $q_D=0$; no negative power
of $q_D$ is introduced. It also proves that null vectors are mapped
to null vectors. Only now may the map be regarded as a bounded
operator on the Hilbert quotient. Endpoint constants may depend on
the number of fields inside $A$, but they disappear before \eqref{eq:uniform-insertion-bound-UI5}.
The final constant does not depend on that number, on $B$, or on
their polynomial coefficients.

\subsubsection*{One ordinary source and a common centering}

Let $h$ be real and supported in the ball of radius $l/4$ around the
cell center. For $a\le l/64$, apply the three-range count in \eqref{eq:graded-source-bound-GS4}
with board side $\eta=l$, support radius $\delta=l/4$ and the same
face margin. Since $\eta/\delta=4$, its face cost is bounded
uniformly as $l\downarrow0$. Coarse scales use the global count
divided by the number of cells, not a separate root count at every
source copy. Upper stability and the floor cost only $Cl^4$.
Thus, for $l\le1$, there are holomorphic univariate normalizers with
\begin{equation}
 E_{h,\ell}(0)=0,\qquad
 p_{l,\ell}\big(e^{z\mathcal O_a(h)-E_{h,\ell}(z)}\big)\le Q_0,
 \qquad |z|\|h\|_\sharp<r_0,
 \label{eq:uniform-insertion-bound-UI6}
\end{equation}
where $Q_0,r_0>0$ are independent of the spacing and cell side.
They are also independent of the surrounding fields. At a fixed
cutoff these normalizers are locally bounded uniformly in the torus,
so a compatible Montel subsequence gives one $E_h$, independently of
every infinite-volume state.

The remaining coarse mesh ratios can also be included. If
$3\le l/a<64$, the ball and its contributing plaquettes lie wholly
inside the cell, and
$\|\mathcal O_a(h)\|_\infty\le C64^4\|h\|_\infty$.
Set $E_{h,\ell}=0$ and enlarge $Q_0$ in \eqref{eq:uniform-insertion-bound-UI6}. The only smaller
allowed ratio is $l/a=1$; then a radius-$a/4$ ball about a lattice
site contains no plaquette barycentre, so $\mathcal O_a(h)=0$.
Consequently \eqref{eq:uniform-insertion-bound-UI6} covers every allowed ratio, with the same constants.

For fixed $z$ put $D_z=e^{z\mathcal O_a(h)-E_h(z)}$.
Homogeneity of the seminorm and the locally uniform convergence of
the normalizers give $q_{D_z}\le Q_0$. Apply \eqref{eq:uniform-insertion-bound-UI5} and then Cauchy's
formula on the one source disc. The expectation is holomorphic,
because at fixed cutoff all observables are bounded. We obtain
\[
 \left|\rho\left(\overline{\theta_tB}
       [\mathcal O_a(h)-E_h'(0)]\tau_{le_4}A\right)\right|
 \le(Q_0/r_0)\|h\|_\sharp\|B\|_t\|A\|_t.
\]
Taking $A=B=1$ in a second arbitrary periodic state $\nu$ bounds
$|\nu(\mathcal O_a(h))-E_h'(0)|$ by the same expression without
the outside norms. The case $D=1$ of \eqref{eq:uniform-insertion-bound-UI5} handles the scalar error.
Hence
\begin{equation}
 \left|\rho\left(\overline{\theta_tB}
       [\mathcal O_a(h)-\nu(\mathcal O_a(h))]
                         \tau_{le_4}A\right)\right|
 \le C\|h\|_\sharp\|B\|_t\|A\|_t.
 \label{eq:uniform-insertion-bound-UI7}
\end{equation}
Both states are arbitrary. Large cutoff-dependent normalizer
derivatives cancel; uniform moments of perturbed Gibbs states have
not been assumed. In using \eqref{eq:uniform-insertion-bound-UI5} for $D_z$, endpoint supremum bounds
may depend on $z$ and the cutoff. They disappear pointwise before the
uniform bound $q_{D_z}\le Q_0$ is used in Cauchy's formula.

\subsubsection*{Continuum passage with fixed support gaps}

Suppose $a_j\downarrow0$ and two periodic states $\rho_j,\nu_j$ at
each datum give convergent centered separated moments
\[
 \sigma_j(f_1,\ldots,f_n)
 =\rho_j\prod_i[\mathcal O_{a_j}(f_i)
                         -\nu_j(\mathcal O_{a_j}(f_i))]
 \longrightarrow\sigma(f_1,\ldots,f_n).
\]
Use the uniform all-state local moment bounds already proved to obtain
equicontinuity under small $C^1$ test changes. In particular, on each
fixed separated tuple, $\sigma_j(\tau_{v_j}X)\to\sigma(X)$ whenever
$v_j\to0$. This follows directly by multilinear telescoping and
$\|\tau_vf-f\|_\sharp\to0$ for compact $C^1$ functions. The constants
in this passage may depend on the fixed tuple and its length; they
do not enter the constant in \eqref{eq:uniform-insertion-bound-UI7}.
Approximating a fixed real translation by lattice translations also
proves translation invariance of $\sigma$, since both states in each
datum are lattice-translation invariant.

Let $F,G$ be arbitrary finite complex linear combinations of field
products whose entries have compact supports strictly in positive
time and are separated within each product. Define
$\|[F]\|_\sigma^2=\sigma(\theta\overline F\cdot F)$.
This is nonnegative as a limit of the lattice reflection forms.
For a fixed power $l\le1$ and a real $h$ supported in
$\bar B((\vec y,0),l/4)$, choose the nearest spatial lattice point
$\vec y_j$ and put
\[
 v_j=(\vec y_j-\vec y,a_j/2),\quad t_j=a_j/2,
 \quad x_{*,j}=(\vec y_j,l/2+a_j/2).
\]
For all sufficiently fine cutoffs $l/a_j$ is odd. Therefore $x_{*,j}$
is a lattice site and the cell faces are link mirrors. The shifted
inserted test is explicitly
$h_j=\tau_{v_j+(l/2)e_4}h$, supported around that cell center.
The identity $\theta_{t_j}\tau_{v_j}=\tau_{v_j}\theta_0$ identifies
the lattice reflected norm tuples with the corresponding small
translations of the continuum norm tuples. No half-mesh translation
invariance of $\nu_j$ is being assumed.

Apply \eqref{eq:uniform-insertion-bound-UI7} to the centered lattice polynomials associated to
$\tau_{v_j}F$ and $\tau_{v_j}G$. All insertion and norm tuples have
fixed positive support gaps for large $j$. Their values converge by
the preceding equicontinuity, including when a limiting Hilbert norm
is zero. Thus
\begin{equation}
 |\sigma(\theta\overline G\cdot(\tau_{(l/2)e_4}h)
                                      \cdot\tau_{le_4}F)|
 \le C\|h\|_\sharp\|[G]\|_\sigma\|[F]\|_\sigma.
 \label{eq:uniform-insertion-bound-UI8}
\end{equation}
The constant is independent of the numbers of entries and the finite
polynomial coefficients of $F,G$. Arbitrary centering states $\nu_j$
are retained throughout.

For completeness time translations are contractions in these norms.
At fixed spacing, reflection Cauchy--Schwarz gives
$\|\tau_uA\|_t^2\le\|A\|_t\|\tau_{2u}A\|_t$ for every positive
lattice time $u$. Iteration against
$\|\tau_{2^Nu}A\|_t\le\|A\|_\infty$ proves contraction. The same
fixed-gap limit passage proves it for every positive real time in
the continuum. It uses no bounded insertion operator in advance.

\subsubsection*{A $C^1$ insertion estimate near the time mirror}

We record the form needed independently of any plate or angular
argument. Let $0<\eta\le1$, and let $f\in C_c^1$ be supported in
$\{\eta\le x_4\le3\eta\}$, with support diameter at most $R$.
Take the largest power $l$ of $L$ below $\eta$ and partition $f$
at mesh $cl$ into pieces in radius-$l/4$ balls. There are at most
$C_L(1+R)^3\eta^{-3}$ pieces: only the three spatial directions
have unbounded length in units of $\eta$. Since $f$ vanishes on both
time faces, $\|f\|_\infty\le\eta\|\partial_4f\|_\infty$.
This cancels the $1/l$ from the partition derivative, because
$\eta/l\le L$. Each piece has sharp norm at most
$C\|f\|_{C^1}$.

Translate a piece to the middle of its side-$l$ time slab and use
\eqref{eq:uniform-insertion-bound-UI8}. Its center time lies between $\eta-l/4$ and $3\eta+l/4$;
both outside factors therefore retain a positive margin of at least
$\eta-3l/4$. Time contraction absorbs those margins. Summing gives
\[
 |\sigma(\theta\overline G\cdot(f)\cdot\tau_{4\eta e_4}F)|
 \le C(1+R)^3\eta^{-3}\|f\|_{C^1}
                                  \|[G]\|_\sigma\|[F]\|_\sigma.
\]
Taking $G=(f)\cdot\tau_{4\eta e_4}F$ and dividing its norm proves
the field-operator estimate; splitting real and imaginary parts handles
complex $f$. We may weaken $\eta^{-3}$ to $\eta^{-4}$ and write
\begin{equation}
 \|[(f)\cdot\tau_{4\eta e_4}F]\|_\sigma
 \le C(R)\eta^{-4}\|f\|_{C^1}\|[F]\|_\sigma.
 \label{eq:uniform-insertion-bound-UI9}
\end{equation}

\begin{proposition}[Uniform one-field insertion]
\label{thm:uniform-one-field-insertion}
Under the stated ordinary source, rooted-domain and lattice-symmetry
hypotheses, every limit just described satisfies \eqref{eq:uniform-insertion-bound-UI8}--\eqref{eq:uniform-insertion-bound-UI9}.
Choose a fixed $0<c<1/64$. If a real $k\in C_c^1$ is supported in a
ball of radius $c\eta$ centered at $y$, $|y_4|\le\eta/8$, and all
entries of $F,G$ lie strictly above $\eta/4$, then
\begin{equation}
 |\sigma(\theta\overline F\cdot(k)\cdot G)|
 \le C\eta^{-4}\|k\|_\sharp
                  \|F\|_{(\eta/4)}\|G\|_{(\eta/4)},
 \quad
 \|F\|_{(\eta/4)}=\|[\tau_{-(\eta/4)e_4}F]\|_\sigma.
 \label{eq:uniform-insertion-bound-UI10}
\end{equation}
The constant depends only on the ordinary source and localization
constants and the fixed geometric choices. It is independent of the
defining sequence, the states, support positions, and the numbers of
fields in the surrounding finite complex polynomials.
\end{proposition}

Indeed put $\xi=c\eta$ and translate the insertion in time by
$2\xi-y_4$, placing it in $[\xi,3\xi]$. After this translation the
left and right polynomials are, respectively,
$\tau_{y_4-2\xi}F$ and $\tau_{-y_4-2\xi}G$, the latter followed by
$\tau_{4\xi e_4}$. Each is a positive time translate of its
$\eta/4$-downshift by at least $\eta/8-2\xi>0$.
Apply \eqref{eq:uniform-insertion-bound-UI9}, reflection Cauchy--Schwarz and contraction, absorbing
the fixed factor $c^{-4}$. This proves \eqref{eq:uniform-insertion-bound-UI10}.
The theorem asserts the $C^1$ insertion bound only. It does not assert
a plate density theorem, an $L^1$ insertion bound, or full rotational
covariance.


% END MODULE chapters/uniform_insertion_bound_v4.tex

% BEGIN MODULE chapters/plate_density_from_insertion_v5.tex
\subsection{Analytic insertion densities and the plate bound}
\label{ch:plate-density-from-insertion}

This section derives an $L^1$ insertion bound from the proved $C^1$
Hilbert insertion bound. All operators are regularized by positive
Euclidean time. No unsandwiched point field is presumed to exist.
The exact spectral cube cone is used, but rotational covariance is not.

Let $\sigma$ be a symmetric, real, time-reflection-invariant,
translation-invariant
functional on admissible separated tests, with its reflection-positive
Hilbert reconstruction. Write $\mathcal D$ for the dense classes of
finite complex positive-time polynomials and $\Omega=[1]$.
Use the probability normalization $\sigma(1)=1$, so $\Omega$ is a
unit vector fixed by all translations.
Positive time translation is $e^{-tH}$, spatial translation is
$U(\mathbf x)=e^{i\mathbf x\cdot\mathbf P}$, and these operators
strongly commute. Assume their joint spectral measure is carried by
\begin{equation}
 E\ge0,\qquad |q_j|\le E\quad(j=1,2,3).
 \label{eq:plate-density-from-insertion-PD1}
\end{equation}
Also assume local $C^1$ continuity of the separated moments, as proved
from the ordinary source bounds. This continuity is used for fixed
finite polynomials only; none of its field-number constants enter the
uniform bounds below.

The insertion input has fixed constants $C_I,c_I>0$, $p\ge0$, with
$c_I<1/64$. For $0<\eta\le1$, let
\[
 I(F;h;G)=\sigma(\theta\overline F\cdot(h)\cdot G),\qquad
 \|F\|_{(\eta/4)}=\|[\tau_{-(\eta/4)e_4}F]\|.
\]
Whenever the entries of $F,G$ lie strictly above $\eta/4$ and $h$ is
real with support in a radius-$c_I\eta$ ball centered at $y$,
$|y_4|\le\eta/8$, the input is
\begin{equation}
 |I(F;h;G)|\le C_I\eta^{-p}\|h\|_\sharp
                       \|F\|_{(\eta/4)}\|G\|_{(\eta/4)}.
 \label{eq:plate-density-from-insertion-PD2}
\end{equation}
The norm is $\|h\|_\sharp=\max(\|h\|_\infty,c_M\|\nabla h\|_\infty)$.
Splitting real and imaginary parts permits complex $h$, at a fixed
factor two. The preceding insertion theorem supplies $p=4$; keeping
$p$ general makes the additional five powers transparent.

\subsubsection*{Bounded forms, real translations and holomorphy}

For $h$ supported in $\bar B(0,c_I\eta)$, define the form on raw
positive-time polynomials by
\[
 b_{h,\eta}(A,B)=I(\tau_{(\eta/4)e_4}A;h;
                                      \tau_{(\eta/4)e_4}B).
\]
By \eqref{eq:plate-density-from-insertion-PD2} it descends through null vectors and extends to a bounded
sesquilinear form on the Hilbert space. Riesz representation gives
a unique bounded operator $B_{h,\eta}$, with
\begin{equation}
 \|B_{h,\eta}\|\le C\eta^{-p}\|h\|_\sharp,
 \quad
 B_{h,\eta}[A]=[(\tau_{(\eta/4)e_4}h)
                                \cdot\tau_{(\eta/2)e_4}A].
 \label{eq:plate-density-from-insertion-PD3}
\end{equation}
The vector identity follows by testing against the dense polynomial
classes and translating the insertion form. Its displayed product
is admissible: the insertion is strictly positive in time and is
separated from the shifted entries of $A$. Thus this identity proves
the action of the bounded operator; it does not postulate a point
field on the quotient.

For complex $z=(z_4,\mathbf z)$ put
\[
 T(z)=\int e^{-z_4E+i\mathbf z\cdot\mathbf q}\,d\Pi(E,\mathbf q).
\]
The spectral cone gives $\|T(z)\|\le1$ if
$\operatorname{Re}z_4\ge\sum_j|\operatorname{Im}z_j|$.
It is operator-norm holomorphic in the strict interior. Indeed, on
a neighbourhood with margin $\epsilon>0$, derivatives of total order
$m$ of the integrand are bounded by $E^m e^{-\epsilon E}$.
Its supremum is at most $m!\epsilon^{-m}$. The exponential Taylor
series consequently converges uniformly in spectral variables on
smaller complex neighbourhoods, so its spectral integrals converge
in operator norm and may be differentiated there.

Now take $h$ supported in $\bar B(0,c_I\eta/2)$, and set
$\Phi=[\tau_{-(\eta/4)e_4}F]$, $\Psi=[\tau_{-(\eta/4)e_4}G]$.
For real $x$, $|x_4|<\eta/8$, direct use of \eqref{eq:plate-density-from-insertion-PD3} gives
\begin{equation}
 I(F;\tau_xh;G)=\langle\Phi,M_{h,\eta}(x)\Psi\rangle,
 \quad M_{h,\eta}(z)=
 T(\tfrac\eta8e_4+z)B_{h,\eta/2}
 T(\tfrac\eta8e_4-z).
 \label{eq:plate-density-from-insertion-PD4}
\end{equation}
For clarity, the real time factors in this formula are
$e^{-(\eta/8+x_4)H}$ and $e^{-(\eta/8-x_4)H}$; the spatial
factors are $U(\mathbf x)$ and $U(-\mathbf x)$.
Substitution in \eqref{eq:plate-density-from-insertion-PD3}, followed by one common time translation,
recovers exactly $\theta\overline F\cdot(\tau_xh)\cdot G$.
Both factors in \eqref{eq:plate-density-from-insertion-PD4} are contractions on
\begin{equation}
 \mathcal O_\eta=
 \{z:|\operatorname{Re}z_4|+
                     \textstyle\sum_{j=1}^3|\operatorname{Im}z_j|
                                                    <\eta/8\}.
 \label{eq:plate-density-from-insertion-PD5}
\end{equation}
Hence $M_{h,\eta}$ is operator-norm holomorphic there and
$\|M_{h,\eta}(z)\|\le C\eta^{-p}\|h\|_\sharp$.
There is no restriction on real spatial translations.

\subsubsection*{A scaled compact parametrix, with its distribution proof}

In four dimensions let
$G(x)=-(64\pi^2)^{-1}|x|^2\log|x|$, with value zero at the origin.
It is $C^1$. For $x\ne0$,
\[
 \Delta(|x|^2\log|x|)=8\log|x|+6,
 \quad \Delta\log|x|=2|x|^{-2},
 \quad \Delta|x|^{-2}=0.
\]
The first two identities also hold as distributions: Green's formula
on $\{|x|>\epsilon\}$ has boundary errors tending to zero, respectively
of orders $O(\epsilon^4|\log\epsilon|)$ and $O(\epsilon^2)$,
up to smaller terms. In the third identity the boundary term is
nonzero. Against a compact smooth test $f$, the inner normal is
$-\partial_r$, so the term containing $f\partial_n|x|^{-2}$ tends
to $-2|S^3|f(0)=-4\pi^2f(0)$; the other boundary term is $O(\epsilon)$.
Thus
\begin{equation}
 \Delta|x|^{-2}=-4\pi^2\delta_0,
 \qquad \Delta^3G=\delta_0.
 \label{eq:plate-density-from-insertion-PD6}
\end{equation}

Choose real $\chi\in C_c^\infty(B(0,1))$ equal to one near zero.
Put $n=\chi G$ and $s=\delta_0-\Delta^3n$. Since $\chi=1$ near
the singularity, \eqref{eq:plate-density-from-insertion-PD6} cancels its delta term. Every remaining
Leibniz term contains a derivative of $\chi$ supported away from
zero, where $G$ is smooth. Therefore $s\in C_c^\infty(B(0,1))$.
For $0<r\le1$ define
\begin{equation}
 n_r(x)=r^2n(x/r)
       =-\frac{\chi(x/r)|x|^2\log(|x|/r)}{64\pi^2},
 \qquad s_r(x)=r^{-4}s(x/r).
 \label{eq:plate-density-from-insertion-PD7}
\end{equation}
Distributional scaling in four dimensions gives
$\Delta^3n_r=\delta_0-s_r$: six derivatives contribute $r^{-6}$,
the prefactor contributes $r^2$, and the scaled delta contributes
$r^4$. Ordinary differentiation gives
\begin{equation}
 \|n_r\|_\sharp\le Cr,\qquad
 \|s_r\|_\sharp\le Cr^{-5}.
 \label{eq:plate-density-from-insertion-PD8}
\end{equation}
Both supports lie in $\bar B(0,r)$. The logarithm is
$\log(|x|/r)$, not $\log|x|$ after simply cutting off $G$.
The subtracted quadratic polynomial removes a spurious $|\log r|$
in \eqref{eq:plate-density-from-insertion-PD8}.

For smooth compact $k$, convolution with \eqref{eq:plate-density-from-insertion-PD6}--\eqref{eq:plate-density-from-insertion-PD7} gives
\begin{equation}
 k=(\Delta^3k)*n_r+k*s_r.
 \label{eq:plate-density-from-insertion-PD9}
\end{equation}
This can be verified directly against compact smooth tests by Fubini
and distributional integration by parts. Alternatively, the compact
$C^1$ kernel $n_r$ is integrable, so all derivatives of its convolution
with $k$ may be placed on $k$ under the integral; the distribution
identity then yields \eqref{eq:plate-density-from-insertion-PD9}. Both sides are $C^1$, so their
distributional equality is also pointwise equality.

\subsubsection*{A bounded analytic density and the $L^1$ bound}

Fix $\epsilon=\min(c_I/8,1/1024)$ and put $r=\epsilon\eta$.
Both $n_r$ and $s_r$ fit the source ball used in \eqref{eq:plate-density-from-insertion-PD4}.
Define the operator-valued holomorphic function
\begin{equation}
 V_\eta(z)=\Delta_z^3M_{n_r,\eta}(z)+M_{s_r,\eta}(z).
 \label{eq:plate-density-from-insertion-PD10}
\end{equation}
On the smaller domain
$|\operatorname{Re}z_4|+\sum_j|\operatorname{Im}z_j|<\eta/32$,
a polydisc of radius $\eta/128$ lies in $\mathcal O_\eta$.
Cauchy's inequalities bound the six derivatives by
$6!(128/\eta)^6$ times the bound in \eqref{eq:plate-density-from-insertion-PD4}. The sum of coefficients
of $\Delta^3$ is $4^3$. \eqref{eq:plate-density-from-insertion-PD8} therefore give
\begin{equation}
 \|V_\eta(z)\|
 \le C\eta^{-p}(\eta^{-6}r+r^{-5})
 \le C\eta^{-(p+5)}.
 \label{eq:plate-density-from-insertion-PD11}
\end{equation}
Every constant is independent of the surrounding polynomials and
their numbers of entries.

Let $k\in C_c^\infty$ have support in $|x_4|<\eta/64$.
For a continuous compact $\phi$ and a compact $C^1$ kernel $q$,
Riemann sums for $\phi*q=\int\phi(x)\tau_xq\,dx$ converge in
$C^1$, with common compact support. Local $C^1$ continuity of the
insertion form consequently implies
$I(F;\phi*q;G)=\int\phi(x)I(F;\tau_xq;G)\,dx$.
All convolution supports here stay strictly between the two outside
time supports. Applying this observation to \eqref{eq:plate-density-from-insertion-PD9}, then integrating
six derivatives off $k$ onto the analytic matrix element in \eqref{eq:plate-density-from-insertion-PD4},
proves
\begin{equation}
 I(F;k;G)=\int k(x)\langle\Phi,V_\eta(x)\Psi\rangle\,dx.
 \label{eq:plate-density-from-insertion-PD12}
\end{equation}
The integrations by parts have no boundary terms, since $k$ is
compactly supported in the open strip. This proves the density
identity, not just an inequality.

Mollification carries \eqref{eq:plate-density-from-insertion-PD12} to compact $C^1$ tests in the strip:
the left side converges by local $C^1$ continuity and the right side
by $L^1$ convergence and \eqref{eq:plate-density-from-insertion-PD11}. Thus it has the unique continuous
extension to every $L^1$ function on that strip, with
\begin{equation}
 |I(F;k;G)|\le C\eta^{-(p+5)}\|k\|_{L^1}
                            \|F\|_{(\eta/4)}\|G\|_{(\eta/4)}.
 \label{eq:plate-density-from-insertion-PD13}
\end{equation}
For an $L^1$ test this formula defines the extended insertion; no
pointwise field interpretation is required. The operator weak
integral $\int k(x)V_\eta(x)\,dx$ is well-defined by its bounded
sesquilinear forms and has the same norm bound. For compact $C^1$
tests it agrees with the regularized field product in \eqref{eq:plate-density-from-insertion-PD3}.
If two choices of $\eta$ admit the same outside polynomials and
test strip, their forms agree first on compact smooth tests and then
on the common $L^1$ domain by density. Thus the extension is
compatible on intersecting pulled-down domains.

Different compact parametrices give the same density on the real
strip: \eqref{eq:plate-density-from-insertion-PD12} makes their continuous matrix-element differences zero
against every compact smooth test. They therefore agree pointwise.
Their holomorphic continuations agree on connected overlapping
domains, by uniqueness from an open set of real points, first in
scalar matrix elements and then as bounded operators.

\subsubsection*{Plates and compatible point-plate operators}

Split $x=(t,u,x_\perp)$, with $x_\perp\in\mathbb R^2$.
For $\varphi\in L^1(\mathbb R^2)$ set
\begin{equation}
 V_{\varphi,\eta}(t,u)
   =\int\varphi(x_\perp)V_\eta(t,u,x_\perp)\,dx_\perp.
 \label{eq:plate-density-from-insertion-PD14}
\end{equation}
The weak integral is a bounded operator, of norm at most
$C\eta^{-(p+5)}\|\varphi\|_1$. It is operator-norm holomorphic for
$|\operatorname{Re}t|+|\operatorname{Im}u|<\eta/32$.
Indeed the integrands are uniformly bounded and holomorphic on
smaller neighbourhoods, so their locally uniformly convergent Taylor
series may be integrated. For an $L^1$ plane profile $\alpha$ supported
in $|t|<\eta/64$, Tonelli and \eqref{eq:plate-density-from-insertion-PD12}--\eqref{eq:plate-density-from-insertion-PD13} give
\begin{equation}
 |I(F;\alpha\otimes\varphi;G)|
 \le C\|\varphi\|_1\eta^{-\wp}\|\alpha\|_1
                 \|F\|_{(\eta/4)}\|G\|_{(\eta/4)},
 \qquad \wp=p+5.
 \label{eq:plate-density-from-insertion-PD15}
\end{equation}
Profiles may extend arbitrarily far in the spatial coordinate $u$.
For smooth compact profiles and labels this is the original insertion;
all other stated $L^1$ cases are its continuous extension.

For $0<a\le1/4$ define the regularized point-plate operator
\begin{equation}
 B_{\varphi;a}=V_{\varphi,4a}(0,0),\qquad
 \|B_{\varphi;a}\|\le C a^{-\wp}\|\varphi\|_1.
 \label{eq:plate-density-from-insertion-PD16}
\end{equation}
It is the operator-norm limit of the regularized insertions of
$\alpha_\epsilon\otimes\varphi$, where $\alpha_\epsilon$ is a
unit-integral smooth plane approximate identity at zero. This follows
from the operator-norm continuity in \eqref{eq:plate-density-from-insertion-PD14}. For $0<a<b\le1/4$,
\begin{equation}
 B_{\varphi;b}=e^{-(b-a)H}B_{\varphi;a}e^{-(b-a)H}.
 \label{eq:plate-density-from-insertion-PD17}
\end{equation}
To prove it, first use a compact smooth plate supported in the central
strips of both scales. The vector formula in \eqref{eq:plate-density-from-insertion-PD3}, or its dense
matrix-element identity, proves the relation for its two regularized
operators. Approximate the plane point by $\alpha_\epsilon$ and pass
in operator norm. Finally approximate an $L^1$ label by smooth compact
labels using \eqref{eq:plate-density-from-insertion-PD15}. This proves compatibility on the stated domains without assuming
an unbounded field common to the scales.

The same calculation gives real translation compatibility. For
$0<a'<a\le1/4$ and $|t|<a/16$,
\begin{equation}
 V_{\varphi,4a}(t,u)=
 e^{-(a-a'+t)H}U(ue_1)B_{\varphi;a'}
 U(-ue_1)e^{-(a-a'-t)H}
 \label{eq:plate-density-from-insertion-PD18}
\end{equation}
whenever $a-a'>|t|$. One proves this first for a sufficiently narrow
smooth approximate identity, using the real translation calculation
of \eqref{eq:plate-density-from-insertion-PD4}, and then passes in operator norm. The right side also
supplies the corresponding holomorphic continuation. These identities
are suitable for composing independently regularized plate insertions.

\begin{proposition}[Analytic density and plate bound]
\label{thm:plate-density-from-insertion}
Under \eqref{eq:plate-density-from-insertion-PD1}--\eqref{eq:plate-density-from-insertion-PD2} and the stated local continuity, the insertion has
the analytic operator density in \eqref{eq:plate-density-from-insertion-PD10}--\eqref{eq:plate-density-from-insertion-PD12}, the uniform $L^1$ bound in \eqref{eq:plate-density-from-insertion-PD13}, and the plate bound in \eqref{eq:plate-density-from-insertion-PD15} with exponent $\wp=p+5$.
The point-plate operators in \eqref{eq:plate-density-from-insertion-PD16} obey the compatibility and translation
identities in \eqref{eq:plate-density-from-insertion-PD17}--\eqref{eq:plate-density-from-insertion-PD18}. All constants are independent of the numbers
of outside fields. No additional typed source estimate is used.
\end{proposition}

As a useful spectral consequence, suppose $\varphi$ is real. Reflection
of real, time-even approximate identities proves
$B_{\varphi;a}^*=B_{\varphi;a}$.
Let $v_{\varphi,a}=B_{\varphi;a}\Omega$ and let $\mu_{\varphi,a}$
be its finite positive joint spectral measure for $(H,P_1)$.
Since $H\Omega=0$, \eqref{eq:plate-density-from-insertion-PD17} gives
$\mu_{\varphi,b}=e^{-2(b-a)E}\mu_{\varphi,a}$.
The positive measure $\mu_\varphi=e^{2aE}\mu_{\varphi,a}$ is thus
independent of $a$, is finite on bounded-energy sets, and may have
infinite total mass. For $t>0$ choose $0<a<\min(t/2,1/4)$ and put
\begin{equation}
 K_\varphi(t,u)=
 \langle v_{\varphi,a},e^{-(t-2a)H+iuP_1}v_{\varphi,a}\rangle
 =\int e^{-tE+iuq}\,d\mu_\varphi(E,q).
 \label{eq:plate-density-from-insertion-PD19}
\end{equation}
It is independent of $a$, finite, and its spectral measure is supported
on $E\ge|q|$. The integral is holomorphic on
$\operatorname{Re}t>|\operatorname{Im}u|$: on smaller neighbourhoods
all differentiated integrands are bounded by a polynomial in $E$
times a strictly decaying exponential, integrable by \eqref{eq:plate-density-from-insertion-PD19} at a
smaller positive real time.

To identify this as the original two-plate kernel, take real plane
approximate identities even in time, first with a smooth compact real
label $\varphi$. The vectors
$[(\tau_{ae_4}\alpha_\epsilon)\otimes\varphi]$
converge to $v_{\varphi,a}$ by \eqref{eq:plate-density-from-insertion-PD16}. Their matrix elements in
\eqref{eq:plate-density-from-insertion-PD19} are precisely the two-plate moments at time separation $t$,
with the narrow profiles, by reflection and common translation.
For compact profiles with a fixed positive time gap, choose $a$
smaller than half that gap. Convolution approximates those profiles
in $C^1$, while the approximating matrix elements are uniformly
bounded by the two vector norm bounds. Fubini and dominated
convergence therefore recover the integral of $K_\varphi$ against
the original profiles. Smooth compact approximation of an $L^1$
label uses the uniform bounds above. The kernel here is the full
two-point moment kernel; for connected moments one applies the
construction to the centered functional.
This is a statement in the chosen coordinate frame. No diagonal-frame cone,
contact value at $t=0$, or rotational covariance is asserted here.


% END MODULE chapters/plate_density_from_insertion_v5.tex

% BEGIN MODULE chapters/angular_kernel_construction_v5.tex
\subsection{Analytic planar kernels and the finite angular band}
\label{ch:angular-kernel-construction}

A plate estimate uniform in the numbers of outside fields yields the
analytic planar kernels used in the angular argument. Rotation invariance
also requires the short-distance two-point estimate. Both estimates are
hypotheses in this section.

\subsubsection*{The assumptions and the meaning of a product}

Write a point of Euclidean space as $(y,z)$, where
$y=(y^0,y^1)=(x_4,x_1)$ and $z=(x_2,x_3)$. Put
$\theta(y^0,y^1)=(-y^0,y^1)$. A tuple is a finite ordered list of
compactly supported $C^1$ test functions with pairwise disjoint
supports. Products of tuples mean concatenation, whenever the
concatenation still has disjoint supports. We use finite complex
linear combinations of such tuples, including the empty tuple $1$.
The notation $(h)$ denotes the tuple with its single entry $h$.

Assume the following properties of a centered functional $\sigma$.
It is multilinear in the entries, symmetric, real on real entries,
and invariant under simultaneous real translations and the
hypercubic group. Its values satisfy $\sigma(1)=1$ and
$\sigma((h))=0$. On every fixed collection of separated compact
support sets it is continuous for the $C^1$ norms, with a bound
independent of the positions of those sets when their diameters and
mutual distances are fixed. No bound on the growth of these constants
with the number of entries is assumed in this section.

Let $\mathscr A_+$ denote the combinations of tuples supported
strictly in $\{x_4>0\}$. Reflection positivity and the reconstruction
give a Hilbert space $\mathcal H$, a dense class map
$F\mapsto[F]$, and $\Omega=[1]$, such that
\begin{equation}
 \langle[F],[G]\rangle=\sigma(\theta\overline F\cdot G),
 \qquad \|\Omega\|=1.
 \label{ak:hilbert}
\end{equation}
The inner product is antilinear in the first entry. The reconstructed
time and planar spatial translations have the actions
\begin{equation}
 e^{-tH}[F]=[\tau_{te_4}F],\qquad
 e^{ixP}[F]=[\tau_{xe_1}F],\qquad t\ge0,
 \label{ak:translations}
\end{equation}
where $H,P$ are strongly commuting self-adjoint operators,
$H\ge0$, and
\begin{equation}
 H\ge |P|,\qquad H\Omega=P\Omega=0.
 \label{ak:cone}
\end{equation}
Only this one planar consequence of the exact cube cone is used.

For $F$ supported above a height $a$ define the pulled-down norm by
\[
 \|F\|_{(a)}=\|[\tau_{-ae_4}F]\|.
\]
A label is a real $\varphi\in C_c^1(\mathbb R^2)$ in the transverse
variables. A profile is a complex $\alpha\in C_c^1(\mathbb R^2)$.
The corresponding plate is $\alpha\otimes\varphi$.
Fix constants $0<c<1/4$ and $\wp>0$. The plate input is that, for
each label, a constant $C_\varphi<\infty$ satisfies
\begin{equation}
 \begin{split}
 &|\sigma(\theta\overline F\cdot
                    (\alpha\otimes\varphi)\cdot G)|\\
 &\hspace{8mm}\le C_\varphi\eta^{-\wp}\|\alpha\|_1
                 \|F\|_{(\eta/4)}\|G\|_{(\eta/4)}
 \end{split}
 \label{ak:plate-input}
\end{equation}
for $0<\eta\le1$, profiles supported in $|y^0|<c\eta$, and
$F,G$ supported strictly above $\eta/4$. This includes arbitrary
finite combinations and arbitrary numbers of entries in $F,G$.
The constant $C_\varphi$ is the same for all these choices and all
$\eta$. Every concatenation in this inequality has separated
supports. The plate theorem obtained from the ordinary-source
insertion estimate supplies precisely this input, with a fixed small
$c$ and $\wp=p+5$.

\subsubsection*{Bounded plates and their padding}

For a profile in the strip of \eqref{ak:plate-input} put $u=\eta/4$.
On the dense classes define
\begin{equation}
 B_\varphi[\alpha;\eta][G]
   =\big[(\tau_{ue_4}(\alpha\otimes\varphi))
                          \cdot\tau_{2ue_4}G\big].
 \label{ak:plate-operator}
\end{equation}
The plate on the right lies at times strictly between $0$ and
$2u$; the other entries lie strictly above $2u$. Thus the right
side is an actual class in $\mathcal H$.

To prove that it depends only on $[G]$, test it against $[F]$ and
translate the full tuple down by $u$. The resulting expression is
\[
 \langle[F],B_\varphi[\alpha;\eta][G]\rangle
 =\sigma\big(\theta\overline{\tau_{ue_4}F}\cdot
          (\alpha\otimes\varphi)\cdot\tau_{ue_4}G\big).
\]
The two pulled-down norms on the right are $\|[F]\|$ and
$\|[G]\|$. By \eqref{ak:plate-input}, density and Hilbert duality,
\begin{equation}
 \|B_\varphi[\alpha;\eta][G]\|
 \le C_\varphi\eta^{-\wp}\|\alpha\|_1\|[G]\|.
 \label{ak:operator-bound}
\end{equation}
In particular null classes map to null classes. There is a unique
bounded extension to the whole Hilbert space, with the same norm
bound. No domain containing an unspecified number of fields remains
in this statement.

The definitions also fix all padding conventions. If $\eta'\le\eta$
and the profile is supported in $|y^0|<c\eta'$, then, with
$d=(\eta-\eta')/4$,
\begin{equation}
 B_\varphi[\alpha;\eta]
   =e^{-dH}B_\varphi[\alpha;\eta']e^{-dH}.
 \label{ak:padding}
\end{equation}
Indeed both sides applied to $[G]$ give the right side of
\eqref{ak:plate-operator}; density proves equality of the bounded
operators. Realness, symmetry, reflection invariance and a
translation by $2u$ similarly give
\begin{equation}
 B_\varphi[\alpha;\eta]^*
       =B_\varphi[\overline\alpha\circ\theta;\eta].
 \label{ak:plate-adjoint}
\end{equation}
For example, conjugate the expression defining
$\langle[G],B_\varphi[\alpha;\eta][F]\rangle$, reflect all its
entries, translate by $2u$, and permute the entries. The resulting
expression is
$\langle[F],B_\varphi[\overline\alpha\circ\theta;\eta][G]\rangle$.
The label remains unchanged because it is real and the reflection
acts only on the plane.

\subsubsection*{The tube and its bounded translations}

For $v=(v^0,v^1)\in\mathbb C^2$ use the linear coordinates
$v_\pm=v^0\pm iv^1$, and set
\begin{equation}
 \mathbb V=\{\operatorname{Re}v_+>0,
                   \operatorname{Re}v_->0\}
       =\{\operatorname{Re}v^0>|\operatorname{Im}v^1|\},
 \qquad v^*=(\overline{v^0},-\overline{v^1}).
 \label{ak:tube}
\end{equation}
This is an open convex cone. Its real points are the half-plane
$\mathbb H=\{v^0>0\}$. Both entrywise conjugation and the star
preserve it, and $\operatorname{Re}v\in\mathbb V$ whenever
$v\in\mathbb V$. These assertions follow directly from the two
strict real linear inequalities in \eqref{ak:tube}. Notice the
distinct identities
\[
 (v^*)_\pm=\overline{v_\pm},\qquad
 v^*=-\theta v\quad\hbox{for real }v.
\]

The joint spectral measure of $(H,P)$ is carried by
$\{(E,q):E\ge |q|\}$. Define on the closed tube
\begin{equation}
 T(v)=\int e^{-v^0E+iv^1q}\,d\mathsf P(E,q).
 \label{ak:T}
\end{equation}
Writing $E_\pm=(E\pm q)/2\ge0$ gives
$-v^0E+iv^1q=-v_+E_--v_-E_+$. Consequently
\begin{equation}
 \|T(v)\|\le1,\quad T(v)T(w)=T(v+w),\quad
 T(v)^*=T(v^*),\quad T(v)\Omega=\Omega.
 \label{ak:T-properties}
\end{equation}
The last assertion uses $H\Omega=P\Omega=0$. Dominated convergence
in the spectral measure of each vector proves strong continuity on
the closed tube.

The map $T$ is holomorphic in operator norm on $\mathbb V$.
Here is a direct proof, which also records why the exact cone is
appropriate. At a point with $\operatorname{Re}v_\pm\ge r>0$,
expand in increments $\Delta_\pm$. The coefficient of
$\Delta_+^j\Delta_-^k$ is the bounded functional-calculus operator
with multiplier
\[
 \frac{(-E_-)^j(-E_+)^k}{j!k!}
                    e^{-v_+E_--v_-E_+}.
\]
Its norm is at most $r^{-j-k}$, since
$x^je^{-x}/j!\le1$ for $x\ge0$. The resulting two power series
converge in operator norm when $|\Delta_\pm|<r$. For real $v$ with
$v^0\ge0$, the reconstruction gives
\begin{equation}
 T(v)[F]=[\tau_{v^0e_4+v^1e_1}F].
 \label{ak:T-real}
\end{equation}

Let $R_z$ be the complex rotation with the usual matrix
\[
 R_z=\begin{pmatrix}\cos z&-\sin z\\ \sin z&\cos z\end{pmatrix}.
\]
In light coordinates $(R_zv)_\pm=e^{\pm iz}v_\pm$. In particular
\begin{equation}
 R_{is}\mathbb V=\mathbb V,\qquad
 (R_{is}v)^*=R_{is}(v^*)\quad(s\in\mathbb R).
 \label{ak:boost-star}
\end{equation}
These are algebraic identities: the multipliers $e^{\mp s}$ are
positive real numbers. For a tuple $v=(v_i)$ in the tube define
\[
 d_{\mathbb V}(v)=
       \min\{1,\min_{i,\pm}\operatorname{Re}(v_i)_\pm\}.
\]
Then its exact transformation and a useful lower bound are
\begin{equation}
 d_{\mathbb V}(R_{is}v)
 =\min\{1,\min_i e^{-s}\operatorname{Re}(v_i)_+,
                   \min_i e^s\operatorname{Re}(v_i)_-\}
 \ge e^{-|s|}d_{\mathbb V}(v).
 \label{ak:margin}
\end{equation}
The same assertions hold in every quarter-turn frame
$\mathbb V_\nu=R_\nu\mathbb V$, $\nu\in(\pi/2)\mathbb Z$.
The rotations commute; $R_\nu$ is real; and the real points of
$\mathbb V_\nu$ are $\{e_\nu\cdot v>0\}$, where
$e_\nu=R_\nu(1,0)$.

We will use the following elementary uniqueness principle. A
holomorphic function on a connected open subset of $\mathbb C^N$
which vanishes on a nonempty relatively open subset of $\mathbb R^N$
vanishes identically. At an interior real point all real derivatives
are zero, so its complex Taylor series is zero; overlapping balls
then propagate the equality along a path. The same proof applies to
Hilbert-valued holomorphic functions after testing against each
vector. An antiholomorphic variable is treated by conjugating that
variable first. This last use will always take place on domains
stable under conjugation.

\subsubsection*{Products with smooth plane profiles}

For $n\ge1$ define the open convex domain
\begin{equation}
 \mathbb D^+_{n,\eta}=
 \{\mathbf x:x_1-\eta e_0/4\in\mathbb V,
       \ x_{i+1}-x_i-\eta e_0/2\in\mathbb V\ (i<n)\}.
 \label{ak:Deta}
\end{equation}
For profiles $\alpha_i$ in the strip $|y^0|<c\eta$ and labels
$\varphi_i$, abbreviate $B_i=B_{\varphi_i}[\alpha_i;\eta]$ and put
\begin{equation}
 \begin{split}
 \Psi_\eta[\vec\alpha](\mathbf x)
   ={}&T(x_1-\eta e_0/4)B_1
        T(x_2-x_1-\eta e_0/2)B_2\cdots\\
      &\hspace{24mm}T(x_n-x_{n-1}-\eta e_0/2)B_n\Omega.
 \end{split}
 \label{ak:smooth-vector}
\end{equation}
It is Hilbert-valued holomorphic, and
\begin{equation}
 \|\Psi_\eta[\vec\alpha](\mathbf x)\|
       \le\prod_{i=1}^n C_{\varphi_i}\eta^{-\wp}\|\alpha_i\|_1.
 \label{ak:smooth-vector-bound}
\end{equation}
The definitions remain valid on the closed domain using strong
continuity of $T$. When all $x_i$ are real and the displayed
non-strict time inequalities hold, direct substitution gives
\begin{equation}
 \Psi_\eta[\vec\alpha](\mathbf x)
       =\Big[\prod_{i=1}^n\tau_{x_i}
                                (\alpha_i\otimes\varphi_i)\Big].
 \label{ak:real-vector}
\end{equation}
For clarity this is an induction from the right. The rightmost
operator first puts its plate at time $\eta/4$. The preceding
translation moves it by $x_n-x_{n-1}-\eta e_0/2$; the next plate
operator translates that entry forward by $\eta e_0/2$ and inserts
its own plate at $\eta e_0/4$. The two centers are therefore
$\eta e_0/4$ and $\eta e_0/4+x_n-x_{n-1}$. Iteration and the
leftmost translation put all centers at the $x_i$. Spatial shifts
follow the same addition rule. Time gaps and $c<1/4$ ensure that
every tuple in this induction is separated.

Define the smoothed scalar kernel on
$U_\eta=(\eta e_0/2+\mathbb V)^{n-1}$ by
\begin{equation}
 K_{n,\eta}[\vec\alpha](v)
   =\langle\Omega,B_1T(v_1-\eta e_0/2)B_2\cdots
                   T(v_{n-1}-\eta e_0/2)B_n\Omega\rangle.
 \label{ak:smooth-kernel}
\end{equation}
It is holomorphic and bounded by the right side of
\eqref{ak:smooth-vector-bound}. If the $y_i$ are real and
$y_{i+1}^0-y_i^0\ge\eta/2$, then
\begin{equation}
 \sigma\Big(\big(\tau_{y_i}
                      (\alpha_i\otimes\varphi_i)\big)_{i=1}^n\Big)
       =K_{n,\eta}[\vec\alpha](v(\mathbf y)),\qquad
 v(\mathbf y)_i=y_{i+1}-y_i.
 \label{ak:real-smooth-kernel}
\end{equation}
Translate the first center to $(\eta/4,0)$ and use
\eqref{ak:real-vector}; taking its scalar product with $\Omega$
proves the formula. Translation invariance then restores the
original centers. Set $K_0=1$ and $K_1=0$, consistently with the
empty and centered one-point moments.

\subsubsection*{Removing the plane smoothing}

Fix $\eta$ and choose real radial nonnegative smooth bumps
$\beta_{i,j}$, each of mass one, whose support radii tend to zero.
For all large $j$ the bumps fit in $D(0,c\eta)$.
The kernels $K_{n,\eta}[\vec\beta_j]$ have a common bound
\begin{equation}
 C_\eta=\prod_i C_{\varphi_i}\eta^{-\wp n}
 \quad\hbox{on }U_\eta.
 \label{ak:normal-bound}
\end{equation}
The bumps need not have the same radii or shapes in different
entries. This freedom is needed later for mixed inner products.

Here are the details identifying every limit. Let $b_i$ be $C^1$
profiles whose successive support time gaps exceed $\eta$. For all
large $j$, their convolutions with $\beta_{i,j}$ have separated
supports in fixed compact sets. The identity
\[
 b_i*\beta_{i,j}=\int b_i(y)\tau_y\beta_{i,j}\,dy
\]
holds also as a $C^1$ Riemann integral. Indeed, for fixed $j$,
the integrand and its first derivatives in the test variable are
uniformly continuous on a common compact set. Riemann sums converge
there in $C^1$, with their supports in the same compact
neighborhood. Multilinearity, the assumed separated $C^1$
continuity, and successive Riemann summation in all entries give
\begin{equation}
 \begin{split}
 \sigma\big(((b_i*\beta_{i,j})\otimes\varphi_i)_i\big)
   =\int\prod_i b_i(y_i)\,
       K_{n,\eta}[\vec\beta_j](v(\mathbf y))\,d\mathbf y.
 \end{split}
 \label{ak:convolution-identity}
\end{equation}
The integrands are continuous and compactly supported; thus Fubini
and the successive integrals are legitimate.

A locally bounded family of holomorphic functions has locally
uniformly convergent subsequences here. One can see this directly:
Cauchy estimates on smaller polydiscs bound first derivatives,
Arzel\`a--Ascoli gives convergence there, and a diagonal choice on
a countable exhaustion gives convergence on all compact subsets.
The Cauchy integral formula shows that the limit is holomorphic.
Apply this observation to \eqref{ak:normal-bound}. Along a
subsequence let the limit be $K'$. Since
$b_i*\beta_{i,j}\to b_i$ in $C^1$, passage in
\eqref{ak:convolution-identity} gives
\begin{equation}
 \sigma\big((b_i\otimes\varphi_i)_i\big)
    =\int\prod_i b_i(y_i)K'(v(\mathbf y))\,d\mathbf y.
 \label{ak:eta-density}
\end{equation}
All differences in this integral lie in a compact subset of
$U_\eta$, so locally uniform convergence suffices.

There is at most one holomorphic $K'$ with this property. If two
continuous densities give the displayed integrals, choose the
$b_i$ to be unit-mass bumps in successively smaller discs around
any fixed real configuration with gaps $>\eta$. Their integrals
converge to the values of the two densities there. These values
coincide. Such differences contain the nonempty relatively open
set $\{v\hbox{ real}:v_i^0>\eta\}$ in the connected set $U_\eta$;
the uniqueness principle proves equality on $U_\eta$.
Consequently the entire smoothing family converges, not merely a
subsequence. Otherwise a sequence staying away on a compact set
would have a further subsequence with the same uniquely determined
limit. Denote it by $K_{n,\eta}$.

The argument proves a useful strengthened conclusion: for any
choice of unit-mass bumps in the different entries with maximal
support radius tending to zero, the corresponding kernels converge
to the same $K_{n,\eta}$, locally uniformly on $U_\eta$. No relation
between their rates of shrinkage is required.

If $\eta'<\eta$, the two limiting kernels agree on $U_\eta$.
Both satisfy \eqref{ak:eta-density} for gaps $>\eta$, and the same
uniqueness proof applies. Since $\bigcup_{\eta>0}U_\eta$
equals $\mathbb V^{n-1}$, this constructs a unique holomorphic
function
\[
 K_n(v;\vec\varphi)\quad\hbox{on }\mathbb V^{n-1}.
\]
It represents $\sigma$ on every time-ordered family of profiles:
\begin{equation}
 \sigma\big((\alpha_i\otimes\varphi_i)_i\big)
   =\int\prod_i\alpha_i(y_i)
                       K_n(v(\mathbf y);\vec\varphi)\,d\mathbf y.
 \label{ak:unsmeared-density}
\end{equation}
Choose $\eta$ below all the support gaps to obtain this formula.
It also proves reality at every real point of the tube, by testing
against real profiles concentrated near the chosen points.

The bound retains exactly one factor per field:
\begin{equation}
 |K_n(v;\vec\varphi)|
      \le\Big(\prod_i C_{\varphi_i}\Big)
                              d_{\mathbb V}(v)^{-\wp n}.
 \label{ak:kernel-bound}
\end{equation}
Indeed $v\in U_\eta$ whenever
$\eta<\min\{1,2\min_{i,\pm}\operatorname{Re}(v_i)_\pm\}$.
Take the limit up to this upper endpoint in
\eqref{ak:normal-bound}; that endpoint is at least
$d_{\mathbb V}(v)$. Neither a derivative of a shrinking bump nor
an estimate whose constant depends on outside field counts has
entered this step.

\subsubsection*{Real overlap identities and the four frames}

For a quarter turn $\nu$ define
$K_n^{(\nu)}(v;\vec\varphi)=K_n(R_{-\nu}v;\vec\varphi)$ on
$\mathbb V_\nu^{n-1}$. Hypercubic invariance and the change of
variables $y_i=R_\nu y_i'$ show that
\eqref{ak:unsmeared-density} holds with $K_n^{(\nu)}$ whenever
the profiles are ordered in the direction $e_\nu$. The rotations
act only on profiles and leave the labels unchanged.

It is useful to account for permutations as well. For a permutation
$\pi$ of $\{1,\ldots,n\}$ let
\begin{equation}
 \mathcal U_{\nu,\pi}=
 \{\mathbf z\in(\mathbb C^2)^n:
       R_{-\nu}(z_{\pi(i+1)}-z_{\pi(i)})\in\mathbb V\ (i<n)\}.
 \label{ak:ordered-domain}
\end{equation}
On it use $K_n$ with these ordered differences and the correspondingly
permuted labels. Each domain is open, convex, nonempty, invariant
under simultaneous complex translations, and satisfies
$\operatorname{Re}\mathbf z\in\mathcal U_{\nu,\pi}$ whenever
$\mathbf z\in\mathcal U_{\nu,\pi}$.

Two of these kernels agree on every nonempty intersection. The
intersection is convex and hence connected. It contains the real
part of each of its points, so contains a relatively open real
ball. At any real configuration in the intersection take small
discs about its points which preserve both strict orders.
The moment of profiles supported in these discs has both continuous
kernel densities, by \eqref{ak:unsmeared-density}, the frame
identity and symmetry of $\sigma$. Testing with unit-mass profiles
shrinking to the points shows equality of the two densities at that
configuration. The holomorphic uniqueness principle now gives
equality on the whole complex intersection. This proves the
overlap identity without assuming arbitrary-angle rotation symmetry.

Thus the ordered kernels glue to a single holomorphic function
$\mathbf K_n(\mathbf z;\vec\varphi)$ on
\begin{equation}
 \mathcal U_n^{\rm ord}=\bigcup_{\nu,\pi}\mathcal U_{\nu,\pi}.
 \label{ak:ordinary-domain}
\end{equation}
This union is connected. To check this explicitly, a real
configuration can be assigned any prescribed strict time ordering
and, independently, any prescribed strict abscissa ordering. Hence
each $0$-frame permutation domain intersects every $\pi/2$-frame
permutation domain. These convex domains form a connected union.
The other two frames give the same domains with their orderings
reversed, so add no disconnected components. Its real points form
\[
 \operatorname{Ord}_n=
 \{\mathbf y:\hbox{all times are distinct or all abscissae
                       are distinct}\}.
\]
The function is translation invariant, symmetric in the labelled
points, real on real configurations, and invariant under the
quarter turns. These properties follow first from the corresponding
moment densities on real open sets and then from uniqueness.
Imaginary rotations preserve each domain
$\mathcal U_{\nu,\pi}$ and hence preserve their union.

\subsubsection*{Point product vectors and their Gram kernel}

Put
\begin{equation}
 \mathbb D_n^+=\{\mathbf x:x_1\in\mathbb V,
                           \ x_{i+1}-x_i\in\mathbb V\ (i<n)\}.
 \label{ak:Dplus}
\end{equation}
This is the union of the nested $\mathbb D^+_{n,\eta}$ as
$\eta$ decreases to zero. For $n=0$ there is one empty
configuration and its vector will be $\Omega$.

We prove strong convergence of the smooth product vectors, rather
than analytically continuing an inequality. Fix $\eta$, choose a
real radial unit-mass bump $\beta_\epsilon$ supported in
$D(0,\epsilon)$, and take $\epsilon<c\eta$. Use it in every
entry of \eqref{ak:smooth-vector}, writing the result as
$\Psi_{\eta,\epsilon}(\mathbf x;\vec\varphi)$.

If $\mathbf x\in\mathbb D^+_{n,\eta}$ and
$\mathbf y\in\mathbb D^+_{m,\eta}$, define the list of differences
\begin{equation}
 J(\mathbf x,\mathbf y)=
 \big((x_n-x_{n-1})^*,\ldots,(x_2-x_1)^*,
             x_1^*+y_1,y_2-y_1,\ldots,y_m-y_{m-1}\big).
 \label{ak:J}
\end{equation}
If $n=0$, omit the left list and the central entry. If $m=0$,
omit the central entry and the right list. The labels of this list
are $(\varphi_n,\ldots,\varphi_1,\psi_1,\ldots,\psi_m)$.
Every entry belongs to $\eta e_0/2+\mathbb V$: for the central
one subtracting $\eta e_0/2$ gives the sum
$(x_1-\eta e_0/4)^*+(y_1-\eta e_0/4)$ in the open cone.
The other entries follow from star invariance.

For different radii $\epsilon,\epsilon'$ the exact identity is
\begin{equation}
 \begin{split}
 &\langle\Psi_{\eta,\epsilon}(\mathbf x;\vec\varphi),
                  \Psi_{\eta,\epsilon'}(\mathbf y;\vec\psi)\rangle\\
 &\quad=K_{n+m,\eta}[
       \underbrace{\beta_\epsilon,\ldots,\beta_\epsilon}_{n},
       \underbrace{\beta_{\epsilon'},\ldots,\beta_{\epsilon'}}_{m}]
                    (J(\mathbf x,\mathbf y)),
 \end{split}
 \label{ak:mixed-inner-product}
\end{equation}
with the reversed left labels specified above. For real positions
this is just \eqref{ak:real-vector}, the Hilbert inner product,
symmetry, and \eqref{ak:real-smooth-kernel}. Radiality and realness
are important: $\theta\overline{\tau_x\beta_\epsilon}
=\tau_{\theta x}\beta_\epsilon$ for real $x$. The left side
is antiholomorphic in $\mathbf x$ and holomorphic in $\mathbf y$;
the right side has exactly the same properties because the star is
conjugate linear. The product domains are convex, conjugation
stable and have real interior points. Mixed holomorphic uniqueness
therefore proves \eqref{ak:mixed-inner-product} for complex
positions too. For total length one both sides are zero; for
length zero both equal one.

The strengthened smoothing conclusion after
\eqref{ak:eta-density} shows that the right side converges to
$K_{n+m}(J(\mathbf x,\mathbf y))$ when both radii tend to zero,
with no restriction on their relative rate. Convergence is uniform
on compact sets of the position domains. In particular, for a
compact subset of $\mathbb D^+_{n,\eta}$, expand
\[
 \|\Psi_{\eta,\epsilon}(\mathbf x)
              -\Psi_{\eta,\epsilon'}(\mathbf x)\|^2
\]
as its four mixed inner products. All four tend uniformly to the
same diagonal kernel value; their signed sum tends to zero.
This proves a locally uniform strong Cauchy property. The limit
$\Psi_\eta(\mathbf x;\vec\varphi)$ is Hilbert-valued holomorphic:
uniform convergence on smaller polydiscs passes the vector Cauchy
formula and its power series to the limit.

These vectors are independent of $\eta$ on overlaps. For any
fixed sufficiently small physical bump and real positions in an
overlap, both definitions give the same class in
\eqref{ak:real-vector}. Vector holomorphic uniqueness on the
connected overlap proves equality for complex positions, and then
the strong limits agree. We have therefore constructed
\begin{equation}
 \Psi(\mathbf x;\vec\varphi)\in\mathcal H,
       \qquad\mathbf x\in\mathbb D_n^+,
 \label{ak:point-vector}
\end{equation}
holomorphically, with the bound
$\|\Psi(\mathbf x;\vec\varphi)\|
\le\prod_i C_{\varphi_i}\eta^{-\wp}$ for every admissible
$\eta$. On the domain where $v\in\overline{\mathbb V}$, the
semigroup rule also gives
\begin{equation}
 \Psi(\mathbf x+v;\vec\varphi)=T(v)\Psi(\mathbf x;\vec\varphi).
 \label{ak:point-vector-translation}
\end{equation}
It holds before removing the bumps and hence afterward.

For labelled positive configurations define
\[
 \operatorname{Gr}(\mathsf x,\mathsf y)
   =K_{n+m}(J(\mathbf x,\mathbf y);
                       \varphi_n,\ldots,\varphi_1,\vec\psi).
\]
Use the empty and one-point conventions already specified. The
construction proves the exact equality
\begin{equation}
 \operatorname{Gr}(\mathsf x,\mathsf y)
       =\langle\Psi(\mathsf x),\Psi(\mathsf y)\rangle.
 \label{ak:gram-vector}
\end{equation}
Thus this is a positive semidefinite kernel at complex positions,
for arbitrary finite families and different lengths, including the
empty configuration.

Finally $R_{is}\mathbb D_n^+=\mathbb D_n^+$, and
\eqref{ak:boost-star} gives
$J(R_{is}\mathbf x,R_{is}\mathbf y)=R_{is}J(\mathbf x,\mathbf y)$.
Consequently
\begin{equation}
 \begin{split}
 \operatorname{Gr}_s(\mathsf x,\mathsf y)
   &:=K_{n+m}(R_{is}J(\mathbf x,\mathbf y);\hbox{the same labels})\\
   &=\langle\Psi(R_{is}\mathsf x),\Psi(R_{is}\mathsf y)\rangle
 \end{split}
 \label{ak:boost-gram}
\end{equation}
is positive semidefinite for every real $s$. This is an actual
Gram representation; positivity has not been inferred from analytic
continuation of a positive scalar function.

\subsubsection*{Narrow pencils and the entire angle function}

Fix $0<\phi_0<\pi/4$ and let
\[
 \operatorname{Sec}=\{v\in\mathbb R^2\setminus\{0\}:
                            |\arg(v^0+iv^1)|<\phi_0\}.
\]
For $n\ge2$, a pencil is a configuration with every consecutive
difference in this sector. The cases $n=0,1$ have their separate
constant kernels fixed above. Write $r_{\mathbf y}=\min_i|y_{i+1}-y_i|>0$.
For a pencil with differences $v_i$ and labels $\vec\varphi$ put,
on the vertical strip
\[
 S_\nu=\{z=t+is:|t-\nu|<\pi/2-\phi_0\},
       \qquad\nu\in(\pi/2)\mathbb Z,
\]
the function
\begin{equation}
 a^{(\nu)}_{\mathbf y}(z)
     =K_n(R_{z-\nu}v;\vec\varphi).
 \label{ak:angle-patch}
\end{equation}
It is well defined and holomorphic. In fact, writing
$(v_i)_\pm=|v_i|e^{\pm i\phi_i}$ with $|\phi_i|<\phi_0$ gives
\begin{equation}
 \operatorname{Re}(R_{t+is-\nu}v_i)_\pm
       =|v_i|e^{\mp s}\cos(\phi_i+t-\nu)>0.
 \label{ak:pencil-margin}
\end{equation}

Only neighboring strips intersect. Their intersection is exactly
\[
 S_\nu\cap S_{\nu+\pi/2}
   =\{t+is:\nu+\phi_0<t<\nu+\pi/2-\phi_0\},
\]
which is nonempty and connected. At a real angle there, all the
rotated differences lie strictly in the quadrant between the two
frame normals. Small discs about the rotated points are ordered in
both frames. The proved real overlap identity therefore identifies
the two values in \eqref{ak:angle-patch} on this real interval.
The one-variable identity theorem identifies the holomorphic
functions throughout the intersection.

The strip width $\pi-2\phi_0$ exceeds their spacing $\pi/2$.
They cover the complex plane and glue to a unique entire function
$a_{\mathbf y}(z;\vec\varphi)$. The construction immediately gives
\begin{equation}
 a_{\mathbf y}(z+\pi/2)=a_{\mathbf y}(z).
 \label{ak:period}
\end{equation}
Indeed use the patch with center $\nu+\pi/2$ for the shifted
argument and the patch with center $\nu$ for the original argument;
the arguments of $K_n$ are identical. For real $t$ it equals
$\mathbf K_n(R_t\mathbf y;\vec\varphi)$, since a suitable frame
orders this rotated pencil. In particular it is real on the real
axis, and $a_{\mathbf y}(is)=K_n(R_{is}v;\vec\varphi)$.

For every real $t$ choose a nearest quarter turn $\nu$, so
$|t-\nu|\le\pi/4$. Set $c_0=\cos(\phi_0+\pi/4)>0$.
By \eqref{ak:pencil-margin}, every light-coordinate margin is at
least $c_0r_{\mathbf y}e^{-|s|}$. Hence
\begin{equation}
 |a_{\mathbf y}(t+is)|
 \le\Big(\prod_iC_{\varphi_i}\Big)
       \min\{1,c_0r_{\mathbf y}e^{-|s|}\}^{-\wp n}
 \le A_{\mathbf y,\vec\varphi} e^{\wp n|s|},
 \label{ak:angle-type}
\end{equation}
where
$A_{\mathbf y,\vec\varphi}
=\prod_iC_{\varphi_i}\max\{1,(c_0r_{\mathbf y})^{-\wp n}\}$.
This proves the claimed exponential type with its linear dependence
on $n$; an uncontrolled prefactor depending on $n$ would not change
that conclusion.

\subsubsection*{The finite band and analytic harmonic coefficients}

Define
\begin{equation}
 c_m(\mathbf y;\vec\varphi)=\frac2\pi
       \int_0^{\pi/2}a_{\mathbf y}(t;\vec\varphi)e^{-4imt}\,dt.
 \label{ak:Fourier}
\end{equation}
Cauchy's theorem on a rectangle of height $s$ and periodicity give
\[
 c_m=\frac2\pi\int_0^{\pi/2}
              a_{\mathbf y}(t+is)e^{-4im(t+is)}\,dt,
 \qquad |c_m|\le A_{\mathbf y,\vec\varphi}
                                      e^{4ms+\wp n|s|}.
\]
The vertical sides cancel exactly. Sending $s$ to $-\infty$ for
$4m>\wp n$, and to $+\infty$ for $4m<-\wp n$, proves
\begin{equation}
 c_m=0\quad\hbox{if }4|m|>\wp n,\qquad
 a_{\mathbf y}(z)=\sum_{4|m|\le\wp n}c_m e^{4imz}.
 \label{ak:finite-band}
\end{equation}
For the equality, periodicity makes $a(z)$ a single-valued
holomorphic function of $w=e^{4iz}$ on $\mathbb C\setminus\{0\}$.
Its Laurent coefficients are the displayed $c_m$, so its Laurent
series is the finite sum after the vanishing just proved.
Equivalently one may use the uniformly convergent Fourier series
of the smooth periodic function on the real axis and then the
identity theorem. Reality implies $c_{-m}=\overline{c_m}$.

We next construct coefficients beyond pencils, without assuming
continuity across different ordering charts. Put
$N=\lfloor\wp n/4\rfloor$ and on planar configurations introduce
the angular vector field
\[
 \mathcal L=\sum_{i=1}^n
             (-y_i^1\partial_{y_i^0}+y_i^0\partial_{y_i^1}).
\]
On translation-invariant functions it is the same angular vector
field in the consecutive-difference coordinates. Let
\[
 P(X)=\prod_{m=-N}^N(X-4im),\qquad
 p_m(X)=\prod_{\substack{-N\le j\le N\\j\ne m}}
                         \frac{X-4ij}{4im-4ij}.
\]
The finite-band identity gives $P(\mathcal L)\mathbf K_n=0$
on the relatively open real set of pencils. The left side is
holomorphic on the connected domain $\mathcal U_n^{\rm ord}$.
Uniqueness therefore proves this equation everywhere on that
domain. Define there
\begin{equation}
 c_m=p_m(\mathcal L)\mathbf K_n\quad(|m|\le N).
 \label{ak:projectors}
\end{equation}
On pencils these are the Fourier coefficients in
\eqref{ak:Fourier}: each polynomial $p_m$ selects its corresponding
distinct exponential. The polynomial identities
$\sum_m p_m=1$ and $(X-4im)p_m(X)=\text{constant}\,P(X)$ give
\begin{equation}
 \sum_m c_m=\mathbf K_n,\qquad
                    \mathcal Lc_m=4imc_m.
 \label{ak:eigen-coefficients}
\end{equation}
The coefficients are holomorphic on the ordered domain, symmetric
in labelled points and invariant under simultaneous translations.
On its real points $c_{-m}=\overline{c_m}$: the angular vector
field has real coefficients, the density is real there, and
$\overline{p_m(X)}=p_{-m}(X)$ for real $X$.
For translation invariance, note that applying $\mathcal L$ to a
translation-invariant function again gives one: the extra term
under a common shift is a linear combination of the total
derivatives $\sum_i\partial_{y_i^\mu}$, which vanish.

Since imaginary rotations preserve the domain, the last equation
in \eqref{ak:eigen-coefficients} can be integrated along them:
\begin{equation}
 c_m(R_{is}\mathbf z)=e^{-4ms}c_m(\mathbf z),\qquad
 \mathbf K_n(R_{is}\mathbf z)=\sum_m e^{-4ms}c_m(\mathbf z).
 \label{ak:boost-expansion}
\end{equation}
The same equation gives
$c_m(R_t\mathbf y)=e^{4imt}c_m(\mathbf y)$ on any real interval
whose rotated configurations remain in $\operatorname{Ord}_n$.

These coefficients extend uniquely and real-analytically to every
configuration of distinct real planar points. Here is the full
gluing argument. Let $\mathcal G_n$ consist of configurations with
distinct points and with no two nonzero difference vectors
perpendicular. This is open and dense: its complement, in addition
to the collisions, is a finite union of zero sets of nonzero real
polynomials. Each polynomial is nonzero, as can also be checked
by taking distinct points on one line. Every real rotation of a
configuration in $\mathcal G_n$ is ordinary. Failure of ordinariness
would produce a time tie and an abscissa tie, whose nonzero
difference vectors are perpendicular. Thus the preceding angular
ordinary differential equation proves covariance for every real
angle on $\mathcal G_n$.

For an arbitrary distinct configuration $\mathbf y$, choose an
angle $t$ at which all its rotated times are distinct. Only
finitely many angles are forbidden. In a neighborhood of
$\mathbf y$ define the extension by
\begin{equation}
 \widetilde c_m(\mathbf u)
           =e^{-4imt}c_m(R_t\mathbf u).
 \label{ak:global-extension}
\end{equation}
This is real-analytic. Two choices agree on the dense set
$\mathcal G_n$ in every overlap, by the covariance already proved
there; continuity makes them agree everywhere in the overlap.
They consequently glue, agree with the original coefficients on
ordinary configurations, and obey exact covariance
\begin{equation}
 \widetilde c_m(R_t\mathbf y)=e^{4imt}\widetilde c_m(\mathbf y)
 \quad(t\in\mathbb R).
 \label{ak:global-covariance}
\end{equation}
This last equality, permutation symmetry, translation invariance
and $\widetilde c_{-m}=\overline{\widetilde c_m}$ follow by
continuity from $\mathcal G_n$. Uniqueness follows because
$\operatorname{Ord}_n$ is dense in the collision-free set.
The collision-free set itself is connected: one may move its
points successively along polygonal paths avoiding the finitely
many positions occupied by the other points, with small preliminary
moves if necessary. Removing finitely many points does not
disconnect the plane.

Set $\overline m(0)=\overline m(1)=0$; for $n\ge2$ let
$\overline m(n)$ be the largest $|m|$ occurring for any pencil and
any real labels, or zero if none occurs. It is a finite integer and
\begin{equation}
 \overline m(n)\le\wp n/4.
 \label{ak:top-frequency}
\end{equation}
Coefficients with $|m|>\overline m(n)$ vanish identically on the
ordered domain by uniqueness, and then on every collision-free
configuration by \eqref{ak:global-extension}. If
$\overline m(n)>0$, some pencil and labels have
$c_{-\overline m(n)}\ne0$, by reality.

\subsubsection*{Boost positivity without time ordering and the top kernel}

For arbitrary labelled configurations of $k$ and $l$ distinct
points in the real open half-plane, define
\begin{equation}
 \widetilde{\operatorname{Gr}}_s(\mathsf x,\mathsf y)
  =\sum_m e^{-4ms}\widetilde c_m^{(k+l)}
       (\theta\mathbf x,\mathbf y;\vec\varphi,\vec\psi).
 \label{ak:global-gram}
\end{equation}
All points in this combined configuration are distinct: the two
groups lie on opposite sides of the mirror. Empty tuples have the
conventions $\widetilde c_m^{(0)}=\delta_{m0}$ and
$\widetilde c_m^{(1)}=0$.

This kernel is positive semidefinite for every real $s$.
First suppose each individual configuration has distinct times.
Sort those times. The combined reflected configuration is then
ordinary, and \eqref{ak:boost-expansion}, symmetry and
\eqref{ak:boost-gram} identify \eqref{ak:global-gram} with the
proved Gram kernel. For a finite family with time ties, perturb
the positions slightly inside the positive half-plane to remove
all such ties. The reflected combined configurations remain
collision free. Every coefficient is continuous there, so their
finite matrices converge entrywise to the desired matrix.
Positive semidefiniteness is closed under this limit. This proves
the assertion with time ties, arbitrary lengths and the empty
configuration included.

In particular, fix $k\ge1$, put $m_*=\overline m(2k)$, and set
$c_* =\widetilde c^{(2k)}_{-m_*}$. On $k$-point configurations
the kernel
\begin{equation}
 (\mathsf x,\mathsf y)\longmapsto
                c_*(\theta\mathbf x,\mathbf y;\vec\varphi,\vec\psi)
 \label{ak:top-positive}
\end{equation}
is positive semidefinite: multiply the matrix in
\eqref{ak:global-gram} by $e^{-4m_*s}$ and send $s\to+\infty$.
All the other finitely many terms tend to zero.

There is also an exact formulation across an arbitrary line
$\mathfrak l$. Let $b\in\mathfrak l$ and let $R_\nu e_0$ be
its normal pointing into a chosen side $\mathfrak l_+$. The map
$r(y)=R_{-\nu}(y-b)$ sends that side to $\mathbb H$ and
intertwines its reflection $\theta_{\mathfrak l}$ with $\theta$.
Using translation invariance and \eqref{ak:global-covariance},
\eqref{ak:top-positive} becomes positivity of
\begin{equation}
 (\mathsf x,\mathsf y)\longmapsto
       e^{4im_*\nu}c_*(\theta_{\mathfrak l}\mathbf x,
                                      \mathbf y;\vec\varphi,\vec\psi)
 \label{ak:top-all-lines}
\end{equation}
on $k$-point configurations in $\mathfrak l_+$. The compensating
unit phase is essential. This is a statement about the top
coefficient, and does not assume reflection positivity of the
original functional at every line.

\subsubsection*{The cut inequality, including the empty tuple}

The preceding constructions imply
\begin{equation}
 2\overline m(k+l)\le\overline m(2k)+\overline m(2l)
                  \qquad(k,l\ge0).
 \label{ak:convexity}
\end{equation}
If the left top frequency is zero there is nothing to prove.
Otherwise choose a pencil and labels whose negative top coefficient
is nonzero. When $k,l\ge1$, translate the midpoint between
$y_k,y_{k+1}$ to the origin, and put
\[
 x_i=\theta y_{k+1-i}\ (1\le i\le k),\qquad
 z_j=y_{k+j}\ (1\le j\le l),
\]
with the correspondingly reversed left labels. Both are positive
configurations. Their first points and all consecutive differences
are in $\operatorname{Sec}$: reflection followed by a minus sign
changes the sign of a real spatial coordinate and preserves this
sector. The first points are the two reflected halves of the cut
difference.

The combined reflected configuration for $\mathsf x,\mathsf z$
is exactly the original pencil. The doubled configurations for
$\mathsf x,\mathsf x$ and $\mathsf z,\mathsf z$ are also pencils.
Their central differences are respectively $2x_1^0e_0$ and
$2z_1^0e_0$, and all other differences are sector vectors or their
stars. The Gram Cauchy--Schwarz inequality gives
\[
 |a_{\mathbf y}(is)|^2
   \le a_{(\theta x_k,\ldots,\theta x_1,x_1,\ldots,x_k)}(is)
       a_{(\theta z_l,\ldots,\theta z_1,z_1,\ldots,z_l)}(is).
\]
The left side has leading rate
$e^{8\overline m(k+l)s}$ as $s\to+\infty$, with nonzero
coefficient. Each right factor has rate at most
$e^{4\overline m(2k)s}$ or $e^{4\overline m(2l)s}$ by its finite
band. Comparison proves \eqref{ak:convexity}.

If $k=0$, translate a nonzero pencil so that its first point is in
$\operatorname{Sec}$. Apply Gram Cauchy--Schwarz to the empty
vector $\Omega$ and its product vector. The empty diagonal entry
is one, so the same comparison gives
$2\overline m(l)\le\overline m(2l)$. This is exactly the required
inequality with $\overline m(0)=0$. The case $l=0$ follows by
symmetry, and the centered one-point and fully empty cases are
immediate. Thus the initial convexity information has not been
lost by omitting the vacuum vector.

\subsubsection*{The two-point interface and its positive plane measure}

For two points, $\operatorname{Ord}_2$ is the entire real
collision-free set: two distinct points cannot have both coordinates
equal. Translation invariance therefore gives a real-analytic
function $G_{\varphi,\psi}(r)$ for every real $r\ne0$, equal to
$K_2(r;\varphi,\psi)$ when $r^0>0$. Its global coefficients obey
\begin{equation}
 \begin{aligned}
 G_{\varphi,\psi}(r)&=\sum_m\widetilde c_m(r;\varphi,\psi),\\
 \widetilde c_m(R_t r)&=e^{4imt}\widetilde c_m(r),\\
 K_2(R_{is}r)&=\sum_m e^{-4ms}\widetilde c_m(r)
       \qquad(r^0>0).
 \end{aligned}
 \label{ak:two-point-interface}
\end{equation}
The last identity holds for every real $s$, not merely near zero;
it is \eqref{ak:boost-expansion}. Every sum is finite, with
$|m|\le\overline m(2)\le\wp/2$. The kernel in
\eqref{ak:global-gram} for one-point configurations is positive
semidefinite for every $s$. A quarter-turn frame with projection
at least $|r|/\sqrt2$ also gives the useful bound
\begin{equation}
 |G_{\varphi,\psi}(r)|
 \le C_\varphi C_\psi
              \min\{1,|r|/\sqrt2\}^{-2\wp}.
 \label{ak:global-two-point-bound}
\end{equation}

The density formula for two points also holds for arbitrary
disjoint compact plane-profile supports, even if those whole
supports have no strict coordinate ordering. Let their distance
be $d>0$ and partition each profile into finitely many pieces
supported in sets of diameter less than $d/(10\sqrt2)$. For each
pair of pieces, two chosen points in their supports differ in at
least one coordinate by $d/\sqrt2$ or more. The small diameters
then make the two entire piece supports strictly ordered in that
coordinate, in one of the four frames. Apply
\eqref{ak:unsmeared-density} and the frame identity to this pair.
Summing the finitely many identities by bilinearity proves
\[
 \sigma(\alpha\otimes\varphi,\beta\otimes\psi)
  =\iint\alpha(y)\beta(y')G_{\varphi,\psi}(y'-y)\,dy\,dy'
\]
for every such pair of profiles. All partitions can be smooth;
their products with the original $C^1$ profiles remain $C^1$ and
have supports in the original separated sets.

There is a positive spectral plane measure for equal real labels.
Let $\zeta_s=\Psi((s,0);\varphi)$, $s>0$. By
\eqref{ak:point-vector-translation},
$\zeta_{s'}=e^{-(s'-s)H}\zeta_s$ if $s'>s$.
Let $\mu_s(B)=\langle\zeta_s,\mathsf P(B)\zeta_s\rangle$,
using the joint spectral measure of $(H,P)$. The locally finite
positive measure
\begin{equation}
 d\mu_\varphi(E,q)=e^{2sE}\,d\mu_s(E,q)
 \label{ak:plane-measure}
\end{equation}
is independent of $s$, because the preceding vector identity
multiplies $\mu_s$ by $e^{-2(s'-s)E}$. It is carried by
$E\ge|q|$ and satisfies
\begin{equation}
 K_2(v;\varphi,\varphi)
     =\int e^{-v^0E+iv^1q}\,d\mu_\varphi(E,q)
                       \qquad(v\in\mathbb V).
 \label{ak:Laplace-Fourier}
\end{equation}
For real $v=(t,x)$ choose $0<s<t/2$ and use
\eqref{ak:gram-vector} and \eqref{ak:point-vector-translation}:
the left side is
$\langle\zeta_s,T(t-2s,x)\zeta_s\rangle$, which is the
displayed integral. In particular
$\int e^{-tE}\,d\mu_\varphi<\infty$ for every $t>0$.
On a compact sub-tube the multiplier is bounded by $e^{-dE}$
for some $d>0$. Its derivatives are controlled by a smaller
exponential weight because $|q|\le E$. The integral is therefore
holomorphic there, and uniqueness extends its equality from real
points to all of $\mathbb V$.

No uniqueness of the vacuum or absence of a zero-energy mass is
asserted here. A positive plane measure in a diagonal-reflection
picture requires the reflection-positivity input for that picture;
it is not inferred from the coordinate cone. Likewise, eliminating
the nonzero two-point harmonics requires a separate ultraviolet
estimate and the diagonal-frame argument. The polygon and integer
arguments can use \eqref{ak:top-all-lines} and
\eqref{ak:convexity}, but are not contained in this section.
Finally, the analytic extension of coefficients to distinct planar
points does not identify the original moment functional on every
exceptional support arrangement. That distributional passage, in
particular arrangements with coincident planar projections and
separated transverse supports, requires a separate proof.


% END MODULE chapters/angular_kernel_construction_v5.tex

% BEGIN MODULE chapters/angular_two_point_elimination_v5.tex
\subsection{Eliminating two-point angular harmonics}
\label{ch:angular-two-point-elimination}

This section proves an implication for the two-point function. Its inputs
are a continuous planar kernel with its axis-frame spectral realization,
the exact cone in that frame, one diagonal reflection, and a uniform
short-distance bound on a specified class of positive shapes. It constructs
the diagonal spectral measure and proves its cone support; a diagonal cone
and continuous rotation covariance are not assumed. The angular continuation
needed at order two is proved directly. In particular this argument does
not require an all-orders highest-harmonic theorem.

The kernel and Hilbert realization are analytic inputs to this section.
They are supplied by the planar-kernel construction when its insertion,
continuity and reconstruction hypotheses hold. This section does not prove
the interacting source estimates or the short-distance input used below.

\subsubsection{The precise two-point interface}

Write a point of $\mathbb R^4$ as $(y,z)$, with
$y=(y^0,y^1)=(x_4,x_1)$ and $z=(x_2,x_3)$. On the plane let
$e_0=(1,0)$, let $R_t$ be the rotation
\[
 R_t(u,v)=(u\cos t-v\sin t,u\sin t+v\cos t),
 \qquad \theta(u,v)=(-u,v).
\]
The same notation denotes these maps on $\mathbb R^4$, leaving $z$ fixed.
Put $n=R_{\pi/4}e_0$. The reflection at $n^\perp$ is
$\theta_n=R_{\pi/4}\theta R_{-\pi/4}$.

Let $S_2$ be a real symmetric bilinear form on pairs of real
$C_c^1(\mathbb R^4)$ tests with disjoint supports, complexified bilinearly.
It is translation- and $W_4$-invariant and satisfies coordinate reflection
positivity. It also satisfies
\begin{equation}
 S_2(\theta_n\overline f,f)\ge0
       \quad\bigl(f\in C_c^1(\{n\cdot y>0\};\mathbb C)\bigr).
 \label{eq:angular-two-point-elimination-TU1}
\end{equation}
Positivity at any one diagonal normal in four dimensions implies this
choice by hypercubic conjugation. All scalar products below are conjugate
linear in the first variable.

The kernel interface is the following explicit assumption. For real
$\varphi,\psi\in C_c^1(\mathbb R^2)$ there are continuous functions
$G_{\varphi,\psi}$ on $\mathbb R^2\setminus\{0\}$ such that
\begin{equation}
 S_2(\alpha\otimes\varphi,\beta\otimes\psi)
 =\iint\alpha(y)\beta(y')G_{\varphi,\psi}(y'-y)\,dy\,dy'
 \label{eq:angular-two-point-elimination-TU2}
\end{equation}
whenever the compact supports of $\alpha,\beta\in C_c^1(\mathbb R^2;\mathbb C)$
are disjoint. There is a Hilbert space with strongly commuting
self-adjoint operators $H,P$, with $H\ge |P|$, and vectors
$\zeta_{\varphi,a}$, $a>0$, such that
\begin{equation}
 \begin{split}
 \zeta_{\varphi,b}&=e^{-(b-a)H}\zeta_{\varphi,a}\quad(b>a),\\
 G_{\varphi,\psi}(t,x)
 &=\langle\zeta_{\varphi,a},
       e^{-(t-2a)H+ixP}\zeta_{\psi,a}\rangle\quad(t>2a).
 \end{split}
 \label{eq:angular-two-point-elimination-TU3}
\end{equation}
This is a point-kernel interface, not a bound on unsmeared field
operators. At order two every distinct planar pair is in an ordinary
coordinate frame. Consequently the two-point specialization of the planar
kernel theorem supplies \eqref{eq:angular-two-point-elimination-TU2}--\eqref{eq:angular-two-point-elimination-TU3} without an exceptional planar set.

When $S_2$ is the connected function of a functional with unit vacuum vector
$\Omega$, apply the interface to
$\zeta_{\varphi,a}-\langle\Omega,\zeta_{\varphi,a}\rangle\Omega$.
The vacuum is translation invariant, so its orthogonal projection commutes
with $H,P$; this subtraction preserves \eqref{eq:angular-two-point-elimination-TU3}, positivity and the cone. Thus
the interface here concerns the connected kernel when the ultraviolet
hypothesis concerns the connected function. No uniqueness of the vacuum
is required.

Kernel uniqueness follows by testing \eqref{eq:angular-two-point-elimination-TU2} with small nonnegative
approximate identities at two distinct points. The same test shows that
the kernels inherit reality, bilinearity in the labels, and all symmetries
of $S_2$ that preserve the plane. In particular, for
$G_\varphi=G_{\varphi,\varphi}$,
\begin{equation}
 G_\varphi(R_{\pi/2}y)=G_\varphi(y),\qquad
 G_\varphi(y^0,-y^1)=G_\varphi(y),
 \qquad G_\varphi(y)\in\mathbb R.
 \label{eq:angular-two-point-elimination-TU4}
\end{equation}

\subsubsection{The axis measure and the plate norm}

Fix a real label $\varphi$. Let $\mathsf P$ be the joint spectral measure
of $(H,P)$ and set
\begin{equation}
 \mu_{\varphi,a}(B)=
    \langle\zeta_{\varphi,a},\mathsf P(B)\zeta_{\varphi,a}\rangle,
 \qquad d\mu_\varphi(E,q)=e^{2aE}\,d\mu_{\varphi,a}(E,q).
 \label{eq:angular-two-point-elimination-TU5}
\end{equation}
The first measure is finite and positive. \eqref{eq:angular-two-point-elimination-TU3} gives
$d\mu_{\varphi,b}=e^{-2(b-a)E}d\mu_{\varphi,a}$, so the second measure
is independent of $a$. It is positive and carried by
\[
 \mathcal C=\{(E,q)\in\mathbb R^2:E\ge |q|\}.
\]
For every $t>0$ choose $a<t/2$ in \eqref{eq:angular-two-point-elimination-TU3}. Then
\begin{equation}
 G_\varphi(t,x)=\int e^{-tE+ixq}\,d\mu_\varphi(E,q),
 \qquad
 \Gamma_\varphi(t)=\int e^{-tE}\,d\mu_\varphi=G_\varphi(t,0)<\infty.
 \label{eq:angular-two-point-elimination-TU6}
\end{equation}
The measure is locally finite and $\sigma$-finite: on $\{E\le N\}$ its
mass is at most $e^{tN}\Gamma_\varphi(t)$. The function $\Gamma_\varphi$
is nonnegative and nonincreasing. By a quarter turn, any $y\ne0$ has a
positive first coordinate at least $|y|/\sqrt2$. Therefore
\begin{equation}
 |G_\varphi(y)|\le\Gamma_\varphi(|y|/\sqrt2).
 \label{eq:angular-two-point-elimination-TU7}
\end{equation}
This bound will give the diagonal reconstruction its bounded time orbits.
It uses no decay at large distance. Define also
\begin{equation}
 \mathbb V=\{v\in\mathbb C^2:
                   \operatorname{Re}v^0>|\operatorname{Im}v^1|\},
 \qquad K_\varphi(v)=\int e^{-v^0E+iv^1q}\,d\mu_\varphi.
 \label{eq:angular-two-point-elimination-TU7a}
\end{equation}
On a compact subset of this tube the absolute multiplier is at most
$e^{-cE}$ for some $c>0$. Its derivatives are dominated by a smaller
exponential because $|q|\le E$. Thus $K_\varphi$ is holomorphic and
agrees with $G_\varphi$ on the real positive-time half-plane.

On positive-time tests put
$\|[f]\|^2=S_2(\theta\overline f,f)$; reflection positivity and
polarization give the associated Hilbert seminorm and its triangle
inequality. For $\alpha$ supported in $\{y^0>0\}$, \eqref{eq:angular-two-point-elimination-TU2} and \eqref{eq:angular-two-point-elimination-TU6} give
\begin{equation}
 \|[\alpha\otimes\varphi]\|^2
 =\int\left|\iint\alpha(t,x)e^{-tE+ixq}\,dt\,dx\right|^2
                                                 d\mu_\varphi(E,q).
 \label{eq:angular-two-point-elimination-TU8}
\end{equation}
Indeed, if $t\ge\delta>0$ on the support, the absolute integral before
interchanging the variables is bounded by
$\|\alpha\|_{L^1}^2\Gamma_\varphi(2\delta)$. This proves Fubini's
applicability, including when the unweighted measure has infinite mass.

\subsubsection{The ultraviolet class and fixed shapes}

For $d>0$ and $\vartheta>0$, let $\mathfrak T(d;\vartheta)$ consist of
pairs of nonzero, nonnegative $C_c^1(\mathbb R^4)$ tests $f,h$, each
supported in a ball of radius $d$, satisfying
\begin{equation}
 \begin{split}
 d&\le\operatorname{dist}(\operatorname{supp}f,
                          \operatorname{supp}h)\le6d,\\
 \int f&\ge\vartheta\|f\|_\infty d^4,\qquad
 \operatorname{Lip}f\le\|f\|_\infty/(\vartheta d),
 \end{split}
 \label{eq:angular-two-point-elimination-TU9}
\end{equation}
and the analogous conditions for $h$. The balls may be placed anywhere.
The ultraviolet input is a finite nonnegative function $u$, with
$u(d)\to0$ as $d\downarrow0$, such that for $0<d\le d_0$,
\begin{equation}
 0\le S_2(f,h)\le u(d)\|f\|_\infty\|h\|_\infty
               \quad\bigl((f,h)\in\mathfrak T(d;\vartheta_*)\bigr).
 \label{eq:angular-two-point-elimination-TU10}
\end{equation}
Here the universal $\vartheta_*>0$ is chosen below, before $S_2$.
A hypothesis at a smaller shape parameter also suffices. A hypothesis
only on a narrower positive-shape class must be checked against this
choice; that inclusion is part of the ultraviolet premise.

The following explicit shapes fix all regularity constants. Choose
$a\in C_c^1((2,3))$, between zero and one, equal to one on
$[9/4,11/4]$, zero outside $[17/8,23/8]$, with cubic smoothstep ramps on
the two intervening intervals. Then
$\int a=5/8$ and $\|a'\|_\infty\le12$.
Let $I=[-1/4,1/4]$, $\tau=\mathbf1_I*\mathbf1_I$, and
\begin{equation}
 b=12(\tau*\tau),\qquad c_b=\int b=3/4.
 \label{eq:angular-two-point-elimination-TU11}
\end{equation}
This is a nonnegative $C_c^2$ spline, supported in $[-1,1]$, with maximum
one and $\|b'\|_\infty=2$. Explicitly it is
$1-6s^2+6|s|^3$ for $|s|\le1/2$, and $2(1-|s|)^3$ for
$1/2\le|s|\le1$. With the Fourier convention
$\widehat b(q)=\int b(s)e^{isq}\,ds$, convolution gives
\begin{equation}
 \widehat b(q)=c_b\left(\frac{\sin(q/4)}{q/4}\right)^4,
 \label{eq:angular-two-point-elimination-TU12}
\end{equation}
where the quotient has value one at zero.

For completeness a transverse partition can use
$q(s)=b(4s/3)$,
$D(s)=\sum_{j\in\mathbb Z}q(s-j)$ and $\chi(s)=q(s)/D(s)$.
There are at most two nonzero summands, $D\ge2/27$ by choosing a nearest
integer, and $\sum_j\chi(s-j)=1$.
The bounds $\|q'\|_\infty\le8/3$ and
$\|D'\|_\infty\le16/3$ give $\|\chi'\|_\infty\le1008$.
Put $\psi(z)=\chi(z_1)\chi(z_2)$. It takes values in $[0,1]$, has
support in $[-3/4,3/4]^2$, and its integer translates sum to one.
Let $\widetilde\psi(z)=\rho(z_1)\rho(z_2)$, where $\rho=1$ on
$[-1,1]$, $\rho=0$ outside $[-3/2,3/2]$, and on each intervening
interval use the cubic smoothstep. Thus
\[
 \int\widetilde\psi=25/4,\quad
 \|\nabla\widetilde\psi\|_\infty\le6,\quad
 \|\nabla\psi\|_\infty\le2016,\quad
 \psi\le\widetilde\psi.
\]
All these functions are globally $C^1$. The numerical choice
\begin{equation}
 \vartheta_* =1/12000
 \label{eq:angular-two-point-elimination-TU13}
\end{equation}
will suffice for all tests constructed next.

\subsubsection{Dilation removes the spline zeros}

For $h>0$ and $1\le\lambda\le2$ put
\[
 \alpha_h^\lambda(t,x)=a(t/h)b(4x/(\lambda h)),
 \qquad f_h^\lambda=\alpha_h^\lambda\otimes\varphi.
\]
The time transform and spatial transform are
\[
 A_h(E)=\int a(t/h)e^{-tE}\,dt,\qquad
 B_{h,\lambda}(q)=\frac{c_b\lambda h}{4}
       \left(\frac{\sin(\lambda hq/16)}{\lambda hq/16}\right)^4.
\]
For $E\ge0$, $A_h(E)\ge(h/2)e^{-3hE}$. There is an absolute $c>0$
such that
\begin{equation}
 \int_1^2|B_{h,\lambda}(q)|^2\,d\lambda
                  \ge c h^2(1+h^2q^2)^{-4}.
 \label{eq:angular-two-point-elimination-TU14}
\end{equation}
Here is a proof that includes every zero of the transform. Set
$x=h|q|/16$. For $0\le x\le1$, the sinc in the display is bounded
below on $0\le\lambda x\le2$, giving a fixed positive lower bound
after division by $h^2$. For $x\ge1$, that divided integral is at least
\[
 c x^{-8}\int_1^2\sin^8(\lambda x)\,d\lambda
       =c x^{-9}\int_x^{2x}\sin^8 v\,dv.
\]
The number
$c_1=\min_{r\in[0,\pi]}\int_r^{r+1}\sin^8v\,dv$ is positive:
the integral is continuous, periodic and never zero. An interval of length
$x\ge1$ contains $\lfloor x\rfloor\ge x/2$ successive intervals of
length one. The last display is therefore at least $c x^{-8}$.
Both regimes imply \eqref{eq:angular-two-point-elimination-TU14}, with constants independent of $h,q$.

The finite number
$C_*=\sup_{r\ge0}e^{-r}(1+r^2)^4$ and $|q|\le E$ imply
\[
 e^{-6hE}(1+h^2q^2)^{-4}\ge C_*^{-1}e^{-7hE}.
\]
Using \eqref{eq:angular-two-point-elimination-TU8}, \eqref{eq:angular-two-point-elimination-TU14} and nonnegative Fubini yields
\begin{equation}
 \int_1^2\|[f_h^\lambda]\|^2\,d\lambda
    \ge c h^4\Gamma_\varphi(7h).
 \label{eq:angular-two-point-elimination-TU15}
\end{equation}
The cone is essential here. Without it, arbitrarily large $|q|$ at small
$E$ could make the spline factor too small.

\subsubsection{Positive pedestals put the plate in the ultraviolet class}

Assume $\varphi\ne0$, choose $r_1>0$ with
$\operatorname{supp}\varphi\subset B(0,r_1)$, and take
$0<h\le\min(1,r_1,2d_0/5)$. Put
$A=\|\varphi\|_\infty+h\|\nabla\varphi\|_\infty>0$ and
\[
 \psi_i(z)=\psi(z/h-i),\quad
 \widetilde\psi_i(z)=\widetilde\psi(z/h-i),\quad
 \varphi_i^-=2A\widetilde\psi_i,\quad
 \varphi_i^+=2A\widetilde\psi_i+\psi_i\varphi.
\]
Only indices in
$\mathcal I=\{i:\psi_i\varphi\not\equiv0\}$ are needed. Its size is
at most $(2r_1/h+3)^2\le25r_1^2h^{-2}$. For a fixed $\lambda$ let
$h_i^\pm=\alpha_h^\lambda\otimes\varphi_i^\pm$. Then
\begin{equation}
 f_h^\lambda=\sum_{i\in\mathcal I}(h_i^+-h_i^-),\qquad
 A\widetilde\psi_i\le\varphi_i^\pm\le3A\widetilde\psi_i.
 \label{eq:angular-two-point-elimination-TU16}
\end{equation}
In particular these tests are nonnegative, nonzero, and have sup norm
between $A$ and $3A$.

To check the support geometry in \eqref{eq:angular-two-point-elimination-TU9}, consider each shape. Each support lies in the ball of radius $d=5h/2$ about
$(5h/2,0,hi)$: its four coordinate half-widths are at most
$h/2,h/2,3h/2,3h/2$, whose squared sum is $5h^2<d^2$.
The reflected ball also has radius $d$. The distance between the support
and its time reflection is at least $4h$ and at most $6h$, hence between
$d$ and $6d$.

Moreover, direct product integration gives
\[
 \int h_i^\pm\ge A h^4
    \frac{\int a\,\int b\,\int\widetilde\psi}{4}
       =\frac{375}{512}A h^4.
\]
Since $\|h_i^\pm\|_\infty\le3A$, the integral condition holds for
every $\vartheta\le1/160$. The product derivative rule gives
\[
 \operatorname{Lip}h_i^\pm\le\frac{A}{h}
 \left[3(12+4\cdot2)+2\cdot6+2016+1\right]
 =2089A/h.
\]
As $\|h_i^\pm\|_\infty\ge A$, \eqref{eq:angular-two-point-elimination-TU13} gives the Lipschitz condition:
$2089A/h\le\|h_i^\pm\|_\infty/(\vartheta_*d)$.
The reflected tests have the same integrals and norms. Thus
$(\theta h_i^\pm,h_i^\pm)\in\mathfrak T(5h/2;\vartheta_*)$.

By \eqref{eq:angular-two-point-elimination-TU10}, reflection positivity and the Hilbert triangle inequality,
\[
 \|[f_h^\lambda]\|
 \le\sum_i\bigl(\|[h_i^+]\|+\|[h_i^-]\|\bigr)
 \le6\,\#\mathcal I\,A\sqrt{u(5h/2)}.
\]
Combining its square with \eqref{eq:angular-two-point-elimination-TU15} proves
\begin{equation}
 \Gamma_\varphi(7h)
   \le C r_1^4
       (\|\varphi\|_\infty+\|\nabla\varphi\|_\infty)^2
       h^{-8}u(5h/2).
 \label{eq:angular-two-point-elimination-TU17}
\end{equation}
Consequently
\begin{equation}
                 t^8\Gamma_\varphi(t)\longrightarrow0
                         \quad(t\downarrow0).
 \label{eq:angular-two-point-elimination-TU18}
\end{equation}
The zero label is immediate. The two factors of $h^{-4}$ in \eqref{eq:angular-two-point-elimination-TU17}
have different origins: normalization in \eqref{eq:angular-two-point-elimination-TU15}, and the square of the
number of transverse cells. No uniform bound on a compact profile's
Fourier transform from below was used.

%
%
%

\subsubsection{A short-distance estimate from the ordinary sources}
\label{sec:angular-mirror-budget}

There is a second way to obtain \eqref{eq:angular-two-point-elimination-TU18}.
It uses reflection positivity at coordinate mirrors, hypercubic and
translation invariance, and the joint spectrum
\(E\ge |p_i|\), \(1\le i\le3\), of the Euclidean Hilbert evolution.
For the connected function use the one-field vectors projected off the
unit vector \(\Omega\). The projection commutes with the evolution and
spatial translations. The short-distance input in this argument is an
upper bound on a single family of positive mirror pairs.

\paragraph{The four-dimensional spectral measure.}
For a compact \(C^1\) test \(f\) in the positive half-space, put
\[
 \widetilde f(E,p)=\int f(x)e^{-x_4 E+i\vec x\cdot p}\,dx,
 \qquad m_f(B)=\langle[f],\mathsf P(B)[f]\rangle,
\]
where \(\mathsf P\) is the joint spectral measure and the brackets
refer to the connected form. Integration of positive-time translations
against another such test gives
\([h*f]=\widetilde h(H,P)[f]=\widetilde f(H,P)[h]\).
The integrals exist in the Hilbert space by strong continuity and compact
support. Consequently
\(|\widetilde h|^2m_f=|\widetilde f|^2m_h\).
Choose a nonnegative \(a\in C_c^1((2,3))\) with positive integral and let
\(v(s)=e^{-s}\mathbf 1_{[0,1]}(s)\), \(v^\vee(s)=v(-s)\), and
\(b=(v*v^\vee)^{*2}\). This is an even nonnegative \(C_c^1\) function
supported in \([-2,2]\). Its Fourier transform satisfies
\[
 \widehat b(q)=\frac{|1-e^{-1+iq}|^4}{(1+q^2)^2}
             \ge \frac{c}{(1+q^2)^2}.
\]
Set \(B_1(x)=a(x_4)\prod_{i=1}^3b(4x_i)\). Its transform is nonzero on
the entire spectral carrier, since the time factor is a positive Laplace
transform. Thus
\(d\rho^{(4)}=|\widetilde B_1|^{-2}\,dm_{B_1}\) is a positive measure
on that carrier, with \(dm_f=|\widetilde f|^2d\rho^{(4)}\) for each \(f\).
Polarization gives the analogous mixed formula. This construction uses
only one-field vectors, before taking any pointwise values of a kernel.

For \(h>0\) let \(B_h(x)=B_1(x/h)\). The time transform is bounded
below by \(c h e^{-3hE}\), and the spatial Fourier bound just proved
gives, on the carrier,
\[
 |\widetilde B_h(E,p)|^2
 \ge c h^8e^{-6hE}(1+h^2E^2)^{-12}
 \ge c' h^8e^{-C hE}.
\]
In particular \(s_{\rm a}(t):=\int e^{-tE}\,d\rho^{(4)}\) is finite
for every \(t>0\), by taking \(Ch\le t\) and bounding the integral by
\((c'h^8)^{-1}\|[B_h]\|^2\). It is nonnegative and nonincreasing.
Each support \(\operatorname{supp} B_h\) lies in a ball of radius
\(d=5h/2\) centered at \((5h/2)e_4\), since its four coordinate
half-widths are at most \(h/2\). Its distance to its reflected support
lies between \(4h\) and \(6h\). Its integral is a positive constant
times \(h^4\), its supremum is fixed, and its Lipschitz constant is
\(O(h^{-1})\). Thus \((\theta B_h,B_h)\) belongs to
\(\mathfrak T(5h/2;\vartheta_1)\) for some fixed \(\vartheta_1>0\).
If this family of pairs satisfies
\begin{equation}\label{eq:angular-mirror-budget-bound}
 S_2(\theta B_h,B_h)\le u(5h/2)\|B_h\|_\infty^2,
 \qquad d^2 u(d)\longrightarrow0\quad(d\downarrow0),
\end{equation}
then \(s_{\rm a}(t)\le Ct^{-8}u(ct)\) for small \(t\), by setting
\(h=t/C\) in the spectral estimate.

For \(t>0\) the continuous kernel on the half-space of displacement
vectors is
\[
 S(te_4+\vec x)=\int e^{-tE+i\vec x\cdot p}\,d\rho^{(4)}(E,p).
\]
Indeed the mixed spectral formula gives this kernel when the two test
supports lie on opposite sides of the coordinate mirror. By translation
it gives the same formula whenever their time projections are strictly
ordered. For each nonzero displacement there is a coordinate whose
positive or negative projection is nonzero. Use the hypercubic
symmetries and small boxes about the two points to obtain the kernel
on that neighborhood. Two such formulas agree on overlaps because
both represent the same bilinear form there, as can also be seen by
testing with approximate identities in the two boxes. Compact
separated test pairs are handled by partitions into finitely many
such boxes. Choosing a coordinate of largest modulus proves
\begin{equation}\label{eq:angular-mirror-budget-axis}
 |S(x)|\le s_{\rm a}(|x|_\infty)\qquad(x\ne0).
\end{equation}

\paragraph{The transverse integral.}
For a real transverse label \(\varphi\), the definition of the plane
kernel and uniqueness on separated planar tests give
\[
 \Gamma_\varphi(t)=\int_{\mathbb R^2}S(te_4+\zeta)
                 (\varphi*\varphi^\vee)(\zeta)\,d\zeta,
 \qquad \varphi^\vee(\zeta)=\varphi(-\zeta).
\]
The weight is bounded in modulus by \(\|\varphi\|_2^2\), by
Cauchy--Schwarz. If \(\varphi\) is supported in a ball of radius \(R\),
the weight vanishes for \(|\zeta|>2R\). Use
\eqref{eq:angular-mirror-budget-axis} on the square \(|\zeta|_\infty\le t\)
and on the annuli with inner radii \(2^it\) and outer radii
\(2^{i+1}t\), \(i\ge0\). The area of an annulus is bounded by
\(C(2^it)^2\). For \(0<t\le\delta\) with \(\delta\) sufficiently
small this gives
\begin{equation}\label{eq:angular-mirror-budget-plate}
 \Gamma_\varphi(t)\le C\|\varphi\|_2^2
 \left[t^{-6}\sum_{2^it\le\delta}2^{-6i}u(c2^it)
       +R^2s_{\rm a}(\delta)\right].
\end{equation}
In this sum the central square is included in the \(i=0\) term.
Annuli with inner radius greater than \(\delta\) are bounded using
monotonicity of \(s_{\rm a}\). After multiplication by \(t^8\), the
sum is bounded by \(C\sup_{0<d\le c\delta}d^2u(d)\), since its terms
are \(c^{-2}2^{-8i}(c2^it)^2u(c2^it)\). For fixed \(\delta\) the
remaining term vanishes as \(t\downarrow0\). Then let
\(\delta\downarrow0\). This proves \eqref{eq:angular-two-point-elimination-TU18}.

\paragraph{The stopped source bound on the mirror pair.}
For a periodic-state limit of first-passage data at \(g\), assume the
ordinary source systems of \cref{prop:source-domains-4-1} at their
first-passage stops, including the bounds in
\eqref{eq:all-state-source-transfer-ST5}--\eqref{eq:all-state-source-transfer-ST6}.
The radius, localized energy bound, entropy gap and geometric constants
must be uniform over the stops used here. The sum of the upper stability
and floor constants at a stop with inverse coupling \(x\) must be bounded
by \(C(1+\log x)\), as in the stopped estimate of
\cref{prop:source-domains-4-1}. Uniformity at a single reference coupling
would give no bound on this sum along the varying stops. The argument
also uses the flow inequalities and independence of the self-calibrated
prefix from its last level and the reference coupling, together with
the exact dilation \eqref{eq:flow-RF53} for the periodic Gibbs states.

Call \(j\) a record if \(x_j<x_i\) for every \(i<j\). The indices
\(0\) and \(K\) are records at first passage. For
\(j_0=1+\lceil c_5/\lambda_\sharp\rceil\), the lower bound
\(x_i-x_{i+j_0}\ge\lambda_\sharp j_0-c_5>0\) gives a new record
within \(j_0\) steps of any record \(i\le K-j_0\). When fewer steps
remain, \(K\) itself supplies the next record. For each fixed small
\(d>0\) and all sufficiently large \(K\), choose \(b\ge0\) so that
\(K-b\) is a record and
\[
 cL^{-j_0-1}\le L^bd\le c,
 \qquad b\le\log_L(c/d),
\]
where \(c\le1\) is fixed. In fact, if \(n=\lfloor\log_L(c/d)\rfloor\),
the depths of successive records have gaps at most \(j_0\), so we
may take the greatest record depth \(b\le n\). For \(b>0\), set
\(g_b^{-2}=x_{K-b}-B_\sharp\). Then \(g_b<g\) and \(K-b\) is the
first-passage cutoff at \(g_b\): every earlier inverse coupling exceeds
\(g_b^{-2}+B_\sharp=x_{K-b}\). The upper flow bound and the original
first-passage window give
\[
 x_{K-b}\le g^{-2}+B_\sharp(b+1).
\]
For \(b=0\) use the original reference coupling. At each stop the
ordinary source systems are needed on arbitrarily large compatible tori.
Their minimum side may depend on \(g_b\), since the all-state transfer
takes the side to infinity at fixed cutoff and coupling.

Dilation by \(L^b\) takes a periodic state at cutoff \(K\) to one at
cutoff \(K-b\). It preserves the two-source cumulant exactly, with
test functions replaced by \(f_i^{[b]}(x)=f_i(L^{-b}x)\).
For \((f_1,f_2)\in\mathfrak T(d;\vartheta_1)\), the rescaled tests
have support radii at most one, a fixed positive separation floor,
and \(\|f_i^{[b]}\|_\sharp\le C\|f_i\|_\infty\). The latter uses
\(\operatorname{Lip}f_i^{[b]}\le\|f_i\|_\infty/(\vartheta_1L^bd)\).
In \eqref{eq:all-state-source-transfer-ST9}--\eqref{eq:all-state-source-transfer-ST12}
take \(n=2\) and use a cell side bounded by a fixed constant. There
are at most two occupied cells. Therefore the normalized generating
function has bound \(Q\le \exp[C(1+\log x_{K-b})]\).
The mixed derivative of the holomorphic normalizer is bounded by
\eqref{eq:all-state-source-transfer-ST10} at the fixed separation floor.
For the remaining factor \(h_\rho\), use
\(\partial_{12}\log h_\rho(0)=\partial_{12}h_\rho(0)
-\partial_1h_\rho(0)\partial_2h_\rho(0)\). Cauchy's formula bounds
this expression by \(C(Q+Q^2)\prod_i\|f_i^{[b]}\|_\sharp\).
Returning to the original unit yields the all-state estimate
\begin{equation}\label{eq:angular-mirror-budget-running}
 |S^T_2[\rho](f_1,f_2)|
 \le C\bigl[1+g^{-2}+\log(1/d)\bigr]^C
        \|f_1\|_\infty\|f_2\|_\infty.
\end{equation}
The constants do not depend on the cutoff, the periodic state, or the
support positions. The estimate passes to the continuum limit at each
fixed \(d>0\). For the mirror pair this is
\eqref{eq:angular-mirror-budget-bound}, since \(d^2\) times any fixed
power of \(1+g^{-2}+\log(1/d)\) tends to zero. Thus the ordinary source
estimates at the running stops give
\eqref{eq:angular-two-point-elimination-TU18} for periodic-state limits.
The angular argument below uses this limit and the two mirror
representations. The coefficient law has a separate role in the
short-distance assertions.

\subsubsection{Entire angular continuation and the critical Fourier modes}

For $r>0$, define on $|\operatorname{Re}z|<\pi/2$
\begin{equation}
 g_r^{[0]}(z)=\int
       \exp[-rE\cos z+irq\sin z]\,d\mu_\varphi(E,q).
 \label{eq:angular-two-point-elimination-TU19}
\end{equation}
Writing $z=u+is$, the absolute integrand is
$\exp[-r\cos u(E\cosh s+q\sinh s)]$.
On the cone this is at most $\exp[-r\cos u\,e^{-|s|}E]$.
On a compact substrip its derivatives are dominated by a polynomial in
$E$ times a smaller exponential. Thus the integral is holomorphic.
Its real restriction is $G_\varphi(r\cos u,r\sin u)$.

On $|\operatorname{Re}z-k\pi/2|<\pi/2$ define
$g_r^{[k]}(z)=g_r^{[0]}(z-k\pi/2)$. On real overlaps these functions
agree by \eqref{eq:angular-two-point-elimination-TU4}. The complex overlaps are connected vertical strips with
real intervals, so the identity theorem glues them. Their union is the
complex plane. The resulting $g_r$ is entire and $\pi/2$-periodic.
For every real $u,s$ choose $k$ with $|u-k\pi/2|\le\pi/4$ to obtain
\begin{equation}
 |g_r(u+is)|\le\Gamma_\varphi(r e^{-|s|}/\sqrt2)
                         =o(e^{8|s|}).
 \label{eq:angular-two-point-elimination-TU20}
\end{equation}
The last little-o, uniform in $u$, is exactly \eqref{eq:angular-two-point-elimination-TU18}.

Define
$a_m(r)=\frac2\pi\int_0^{\pi/2}g_r(u)e^{-4imu}\,du$.
Cauchy's theorem on a rectangle of height $s$, with cancellation of its
vertical sides by periodicity, gives
\[
 |a_m(r)|\le e^{4ms}\sup_{u\in[0,\pi/2]}|g_r(u+is)|.
\]
For $m\ge2$ take $s\to-\infty$; for $m\le-2$ take
$s\to+\infty$. \eqref{eq:angular-two-point-elimination-TU20} gives $a_m(r)=0$ in both cases,
including $|m|=2$. A mere $O(t^{-8})$ bound on $\Gamma$ would not
eliminate that boundary mode. Fourier uniqueness for a continuous periodic
function now gives, first on the real line and then everywhere by the
identity theorem,
\begin{equation}
 g_r(z)=a_0(r)+a_1(r)e^{4iz}+a_{-1}(r)e^{-4iz}.
 \label{eq:angular-two-point-elimination-TU21}
\end{equation}
One elementary uniqueness proof uses the positive Fejer kernels: their
integrals are one, their mass outside every fixed neighbourhood of zero
tends to zero, and their convolution converges uniformly on continuous
periodic functions. A function with all Fourier coefficients zero has
all these convolutions zero. Thus no convergence assumption on a formal
Fourier series is implicit in \eqref{eq:angular-two-point-elimination-TU21}.

\subsubsection{The positive measure in the diagonal frame}

Put $D(t,x)=G_\varphi(R_{\pi/4}(t,x))$ for $t>0$.
We construct its positive measure without assuming a diagonal cone.
For $y=(t,x)$, $t>0$, introduce a formal vector $e_y$ with
\begin{equation}
 \langle e_y,e_{y'}\rangle_D=D(t+t',x'-x).
 \label{eq:angular-two-point-elimination-TU22}
\end{equation}
This kernel is positive semidefinite. To see it, approximate each planar
point $R_{\pi/4}y$ by a nonnegative radial smooth mollifier supported
strictly on the positive side of the diagonal mirror, tensor it with
$\varphi$, and apply \eqref{eq:angular-two-point-elimination-TU1} to an arbitrary complex linear combination.
The formula in \eqref{eq:angular-two-point-elimination-TU2} and continuity away from zero give \eqref{eq:angular-two-point-elimination-TU22} in the limit.
Thus its quotient and completion define a Hilbert space.

On finite sums define $T_s e_{(t,x)}=e_{(t+s,x)}$, $s\ge0$, and
$U_b e_{(t,x)}=e_{(t,x+b)}$, $b\in\mathbb R$. The $U_b$ preserve the
form and have inverse $U_{-b}$. The $T_s$ are symmetric and commute with
the $U_b$. If $v=\sum c_i e_{(t_i,x_i)}$ and
$\delta=2\min_i t_i>0$, \eqref{eq:angular-two-point-elimination-TU7} gives a finite bound on
$\|T_s v\|$ for every $s\ge0$, namely its square is at most
$\Gamma_\varphi(\delta/\sqrt2)(\sum|c_i|)^2$.
On the unquotiented vector space Cauchy--Schwarz gives
\[
 \|T_s v\|^2=\langle v,T_{2s}v\rangle_D
                         \le\|v\|\,\|T_{2s}v\|.
\]
If $\|v\|=0$ this preserves the null space. Otherwise iterate and use
the bounded orbit to get $\|T_s v\|\le\|v\|$. Therefore $T_s$ descend
to self-adjoint contractions. They are positive because
$T_s=T_{s/2}^2$. Continuity of $D$ implies strong continuity of both
families on point vectors, hence on their completion.

There are strongly commuting generators $H_D\ge0$ and $P_D$ with
$T_s=e^{-sH_D}$ and $U_b=e^{ibP_D}$. For specificity, $T_1$ is
injective: $T_1v=0$ implies $T_{1/2}v=0$, then all its dyadic roots
annihilate $v$, and strong continuity gives $v=0$.
Define $H_D=-\log T_1$ by spectral calculus. Uniqueness of positive
roots proves the semigroup identity first at rational times and then at
all times. The spectral theorem for the commuting unitary group supplies
$P_D$ and the joint spectral measure. This is the elementary Hilbert
reconstruction for \eqref{eq:angular-two-point-elimination-TU22}, not an OS analytic-continuation theorem.

Let $v_a=e_{(a,0)}$. If $\mathsf P_D$ is that joint spectral measure,
define
\begin{equation}
 \nu_a(B)=\langle v_a,\mathsf P_D(B)v_a\rangle_D,\qquad
 d\nu(E,q)=e^{2aE}\,d\nu_a(E,q).
 \label{eq:angular-two-point-elimination-TU23}
\end{equation}
As $v_b=T_{b-a}v_a$ for $b>a$, $\nu$ is independent of $a$.
It is a positive, locally finite, $\sigma$-finite measure on
$[0,\infty)\times\mathbb R$, and for $t>0$
\begin{equation}
 D(t,x)=\int e^{-tE+ixq}\,d\nu,\qquad
 \int e^{-tE}\,d\nu=D(t,0)<\infty.
 \label{eq:angular-two-point-elimination-TU24}
\end{equation}
Indeed choose $a<t/2$ and evaluate
$\langle v_a,T_{t-2a}U_xv_a\rangle_D$.
Nothing so far asserts $|q|\le E$ for $\nu$.

\subsubsection{Fourier--Laplace uniqueness and the boost identity}

We shall use the following uniqueness fact. Positive Borel measures
$\rho_1,\rho_2$ on $[0,\infty)\times\mathbb R$ with finite
$\int e^{-tE}d\rho_i$ for every $t>0$ are equal if their transforms
$\int e^{-tE+ixq}d\rho_i$ agree for all $t>0,x\in\mathbb R$.
For each real $x$ these transforms are holomorphic in $t$ on
$\operatorname{Re}t>0$, by exponential domination, so the identity
theorem gives equality at $t=a+i\tau$. These are precisely the joint
Fourier transforms of the finite measures $e^{-aE}\rho_i$.
Finite Fourier transforms determine their measures: convolve their
difference with a Gaussian, apply the Gaussian Fourier integral and
Fubini to get zero, and let the Gaussian width tend to zero against
compactly supported continuous functions. Hence the weighted measures,
and then the original measures, agree. The same proof applies to finite
positive linear combinations on either side of an identity.

For $s\in\mathbb R$ define the Lorentz map of spectral variables
\begin{equation}
 L_s(E,q)=(E\cosh s+q\sinh s,E\sinh s+q\cosh s),
 \qquad \mu_s=(L_s)_*\mu_\varphi.
 \label{eq:angular-two-point-elimination-TU25}
\end{equation}
This preserves $\mathcal C$ and $L_sL_t=L_{s+t}$. Its energy satisfies
$E_s\ge e^{-|s|}E$, so $\int e^{-tE}d\mu_s\le
\Gamma_\varphi(te^{-|s|})$. Define
$F_s(t,x)=\int e^{-tE+ixq}d\mu_s$ for $t>0$.
Direct multiplication of the rotation matrix by $i s$ gives
\begin{equation}
 F_s(v)=K_\varphi(R_{is}v)
       =g_r(\theta+is)
       \quad\bigl(v=R_\theta(re_0),\ v^0>0\bigr),
 \label{eq:angular-two-point-elimination-TU26}
\end{equation}
where the middle expression uses \eqref{eq:angular-two-point-elimination-TU7a}.
In fact $R_{is}(t,x)=(t\cosh s-ix\sinh s,
x\cosh s+it\sinh s)$, and its exponential multiplier is exactly
$e^{-tE_s+ixq_s}$. Thus \eqref{eq:angular-two-point-elimination-TU26} holds on the whole real half-plane.

Write $C=\cosh4>1$. The three terms in \eqref{eq:angular-two-point-elimination-TU21} obey the algebraic identity
\[
 (C-1)g_r(\theta+\pi/4)
       +g_r(\theta+i)+g_r(\theta-i)
                         =(C+1)g_r(\theta).
\]
Using \eqref{eq:angular-two-point-elimination-TU6}, \eqref{eq:angular-two-point-elimination-TU24} and \eqref{eq:angular-two-point-elimination-TU26} and then transform uniqueness gives the
identity of positive measures
\begin{equation}
 (C-1)\nu+\mu_1+\mu_{-1}=(C+1)\mu_\varphi.
 \label{eq:angular-two-point-elimination-TU27}
\end{equation}
Every measure has the required finite Laplace weights. In particular
\begin{equation}
       0\le\nu\le\frac{C+1}{C-1}\mu_\varphi.
 \label{eq:angular-two-point-elimination-TU28}
\end{equation}
The diagonal measure is therefore carried by the exact cone. This
conclusion was deduced from \eqref{eq:angular-two-point-elimination-TU27}, not inserted as an additional
reflection-positivity consequence.

For $r>0$ and real $s$ the function
$w_s(E,q)=\exp[-r(E\cosh s+q\sinh s)]$ is positive on the cone and
at most $e^{-r e^{-|s|}E}$. It is integrable against all four measures
in \eqref{eq:angular-two-point-elimination-TU27}; for $\mu_{\pm1}$ use $L_sL_{\pm1}=L_{s\pm1}$.
Integrating \eqref{eq:angular-two-point-elimination-TU27} against $w_s$ yields
\[
 (C-1)\int w_s\,d\nu+F_{s+1}(r,0)+F_{s-1}(r,0)
                                  =(C+1)F_s(r,0).
\]
Substitute \eqref{eq:angular-two-point-elimination-TU21} into this identity to obtain
\begin{equation}
 \begin{split}
 F_s(r,0)&=a_0(r)+a_1(r)e^{-4s}+a_{-1}(r)e^{4s}\ge0,\\
 \int w_s\,d\nu
       &=a_0(r)-a_1(r)e^{-4s}-a_{-1}(r)e^{4s}\ge0.
 \end{split}
 \label{eq:angular-two-point-elimination-TU29}
\end{equation}
The first positivity comes from the axis measure, the second from the
diagonal measure. Divide by $e^{4s}$ and let $s\to+\infty$:
the two limits give $a_{-1}(r)\ge0$ and $a_{-1}(r)\le0$.
Then let $s\to-\infty$ to get both signs for $a_1(r)$.
The limits also show these coefficients are real, if that was not already
deduced from \eqref{eq:angular-two-point-elimination-TU4}. Thus $a_1(r)=a_{-1}(r)=0$ for every $r>0$.

% Keep the closing statement and argument together.
\clearpage
\begin{proposition}[Two-point angular elimination]
\label{thm:angular-two-point-elimination}
Under \eqref{eq:angular-two-point-elimination-TU1}--\eqref{eq:angular-two-point-elimination-TU3} and the short-distance hypothesis in \eqref{eq:angular-two-point-elimination-TU9}--\eqref{eq:angular-two-point-elimination-TU13},
every equal-real-label plane kernel is radial:
\[
             G_\varphi(R_t y)=G_\varphi(y)
                      \quad(y\ne0,\ t\in\mathbb R).
\]
Every nonzero angular coefficient of the two-point kernel vanishes,
including coefficients with mixed transverse labels.
\end{proposition}

\begin{proof}
The preceding argument proves the first assertion. It uses the
short-distance hypothesis only through
\eqref{eq:angular-two-point-elimination-TU18}, so the same conclusion
holds under the other stated hypotheses when that limit is obtained
from \cref{sec:angular-mirror-budget}. Symmetry of $S_2$
gives $G_{\varphi,\psi}(y)=G_{\psi,\varphi}(-y)$; invariance under
the planar half turn gives $G_{\psi,\varphi}(-y)=G_{\psi,\varphi}(y)$.
Therefore the mixed kernel is symmetric in its real labels and
\[
 G_{\varphi,\psi}
       =\tfrac14(G_{\varphi+\psi}-G_{\varphi-\psi}).
\]
Both terms on the right are radial. Complex labels follow by bilinear
extension. Uniqueness of Fourier coefficients identifies this conclusion
with vanishing of every $m\ne0$ coefficient in any finite-angular-band
construction of the same two-point kernel. In its notation,
$\overline m(2)=0$.
\end{proof}

This is the order-two conclusion used by a later all-orders argument.
It does not prove that argument, nor does it replace the required
extension from planar plate kernels to arbitrary separated four-dimensional
tuples. The ultraviolet vanishing and the kernel/spectral interface remain
the stated premises. No rotational Ward identity, infrared decay,
diagonal cone assumption, or full $\operatorname{O}(4)$ covariance was used.


% END MODULE chapters/angular_two_point_elimination_v5.tex

% BEGIN MODULE chapters/angular_harmonic_closure_v5.tex
\subsection{Euclidean covariance from the two-point angular law}
\label{ch:angular-harmonic-closure}

The analytic kernels and boost Gram identities of
\cref{ch:angular-kernel-construction}, together with the two-point
elimination of \cref{ch:angular-two-point-elimination}, give
Euclidean covariance at every order and on distributions. The argument
requires the two-point ultraviolet estimate and uses no rotational Ward
identity.

Work first with the centered scalar hierarchy. Compact separated moments
are symmetric, translation invariant and hypercubic invariant, have position-uniform local
$C^1$ bounds, and are reflection positive in every coordinate picture.
In each such picture the Hilbert generators have the exact cube cone
$H_j\ge |P_{j,r}|$ for every tangential coordinate $r\ne j$.
Their off-diagonal distributions are those constructed in
\cref{ch:quantitative-flat-extension}.

\subsubsection{The harmonic data used below}

The chosen plane has coordinates $(x_4,x_1)$, and $R_u$ denotes its
rotation through $u$. A plate label is a real smooth compact function of
the two transverse coordinates. The analytic kernel construction supplies,
for each $n$, real-analytic functions
\begin{equation}
 c_m(y_1,\ldots,y_n;\varphi_1,\ldots,\varphi_n),
 \qquad y_i\in\mathbb R^2,\quad y_i\ne y_j,
 \label{eq:angular-harmonic-closure-HC1}
\end{equation}
multilinear in the labels, with only $|m|\le Cn$ possibly nonzero.
Their sum is the moment density on the ordinary ordering charts.
They are symmetric under simultaneous permutations of positions and labels,
translation invariant in the plane, and satisfy
\begin{equation}
 c_m(R_u\mathbf y;\vec\varphi)
       =e^{4imu}c_m(\mathbf y;\vec\varphi),
 \qquad c_{-m}=\overline{c_m}
 \quad\hbox{for real labels.}
 \label{eq:angular-harmonic-closure-HC2}
\end{equation}
An ordinary chart means that all time coordinates are distinct, or that
all abscissae are distinct; a fixed ordering is imposed on its connected
component. The functions in \eqref{eq:angular-harmonic-closure-HC1} already extend over the nongeneric planar
configurations with distinct points. This is the analytic-projector
conclusion of the kernel construction, not a density assertion about
distributions on the exceptional set.

\phantomsection\label{notdef:harmonic-degree}
Let $b_n$ be the greatest absolute index of a harmonic that is not
identically zero for some choice of real labels. Set $b_n=0$ if there is
no nonzero nonconstant harmonic, and put $b_0=0$ for the vacuum kernel.
Thus $b_n$ are nonnegative integers and $b_n\le Cn$. The boost Gram
identity gives
\begin{equation}
 2b_{k+l}\le b_{2k}+b_{2l},\qquad k,l\ge0.
 \label{eq:angular-harmonic-closure-HC3}
\end{equation}
For clarity, this inequality compares global degrees. A nonzero coefficient
can be tested on an open configuration set split into $k$ negative and
$l$ positive times: real analyticity and connectedness prevent a nonzero
coefficient from vanishing on that whole open set. Cauchy--Schwarz for the
boosted product vectors bounds its largest exponential by the two diagonal
ones. The empty product is the vacuum vector, so the cases $k=0$ or $l=0$
are included.

For $m_*=b_{2k}$ the coefficient $c_{-m_*}$ is a positive kernel on
reflected pairs of $k$-point configurations in a half-plane. This follows
by multiplying the positive boost Gram matrix by $e^{-4m_*s}$ and taking
$s\to+\infty$. The covariance in \eqref{eq:angular-harmonic-closure-HC2} transports it to every oriented line:
after multiplication by the fixed unit phase for that line it is again
positive. We use the zero-row consequence of this positivity, which is
unchanged by that phase. All these statements have explicit product-vector
proofs in the preceding kernel section; analytic continuation of a
nonnegative real-valued function alone would not prove them.

\subsubsection{The polygon divisibility argument, with label release}

Fix $k\ge1$, put $m_*=b_{2k}$, and abbreviate $H=c_{-m_*}$.
We prove
\begin{equation}
                    k\mid 2m_* . \label{eq:angular-harmonic-closure-HC4}
\end{equation}
Suppose otherwise. For any real plate label $\varphi$, place that same
label at every vertex of a regular $2k$-gon. Rotation through $\pi/k$
permutes the labelled vertices. Symmetry and \eqref{eq:angular-harmonic-closure-HC2} imply
\begin{equation}
 H(\mathbf y;\varphi,\ldots,\varphi)
 =e^{-4\pi i m_*/k}H(\mathbf y;\varphi,\ldots,\varphi)=0.
 \label{eq:angular-harmonic-closure-HC5}
\end{equation}
The polygon may be translated; rotation about its center has the same
effect by translation invariance. The assumption that \eqref{eq:angular-harmonic-closure-HC4} fails is
exactly what makes the phase in \eqref{eq:angular-harmonic-closure-HC5} different from one.

Call a labelled set $S$ a zero set if $H$ vanishes on every $2k$-point
extension of it, with all remaining real labels arbitrary. A full polygon
with the labels above is a zero set, since no completion is then needed.
Here is the reduction of such a set. Suppose it is contained, with its
labels, in $\theta Q\cup Q$, where $Q$ has $j\le k$ points strictly on
one side of a line, and the reflected labels agree. Fill $Q$ by any
$k-j$ additional points $A$ on that side, with arbitrary real labels.
The positive kernel has zero diagonal at $Q\cup A$, because its diagonal
configuration contains $S$. Its entire row is therefore zero. Indeed,
positivity of the $2\times2$ Gram matrix gives
$|K(X,Y)|^2\le K(X,X)K(Y,Y)$; a zero diagonal makes the right side zero.
Consequently
\begin{equation}
 H\big(\theta(Q\cup A)\cup B\big)=0
 \label{eq:angular-harmonic-closure-HC6}
\end{equation}
for every $k$-point $B$ on the other side of the reflected set, with
arbitrary real labels.

For each fixed choice of these free labels, \eqref{eq:angular-harmonic-closure-HC6} holds on a nonempty
open set of all the free positions: the reflected points of $A$ and the
points of $B$ can be varied independently in their respective half-planes.
The configuration space of these $2k-j$ free points, avoiding each other
and the fixed reflected $Q$, is connected. It is the complement of a
finite union of affine sets of codimension two. One elementary proof is
to join two configurations through a generic intermediate point; the
intermediate points whose segment meets one of these affine sets form a
finite union of measure-zero cones. Real analyticity and the identity
principle extend \eqref{eq:angular-harmonic-closure-HC6} to this whole configuration space. It follows that
$\theta Q$ itself is a zero set, with no restriction on the remaining
labels or half-planes. This is the label-release step; it is stronger than
just showing that a special diagonal configuration has zero value.

Apply this reduction to the polygon. Reflection in the perpendicular
bisector of an edge pairs its vertices and has no vertex on its line.
It first reduces the full polygon to a consecutive arc of $k$ vertices.
A consecutive arc of $j\le k$ vertices fits into the reflected union of
an arc of $\lceil j/2\rceil$ vertices. To check the odd case, index the
vertices by integers modulo $2k$. For the arc $0,\ldots,j-1$, reflect
indices by $i\mapsto j-i$ when $j$ is odd; the union has $j+1$ vertices,
paired without a fixed vertex. For even $j$, use $i\mapsto j-1-i$ and
the arc itself is paired. These are edge-bisector reflections, and the
selected half has the stated size. All its labels are still $\varphi$.
Repeated reduction reaches a single labelled point.

Translation covariance makes every point with label $\varphi$ a zero
set. Since $\varphi$ was arbitrary, every labelled $2k$-point
configuration has zero value: choose any one of its labels as that
single fixed label and use the unrestricted completion. Thus $H$ is
identically zero, contradicting the definition of $m_*$. This proves
\eqref{eq:angular-harmonic-closure-HC4}. When $m_*=0$ the divisibility is automatic; no contradiction
argument is needed.

\subsubsection{Discrete convexity reduces every order to two points}

Set $a_k=b_{2k}$, $a_0=0$. \eqref{eq:angular-harmonic-closure-HC3} with $k-1,k+1$ gives
\begin{equation}
 2a_k\le a_{k-1}+a_{k+1},\qquad 0\le a_k\le C'k.
 \label{eq:angular-harmonic-closure-HC7}
\end{equation}
The integer increments $d_k=a_k-a_{k-1}$ are nondecreasing, initially
nonnegative, and bounded. If one increment exceeded the linear-growth
constant, all later increments would do so, contradicting \eqref{eq:angular-harmonic-closure-HC7}.
Hence they are eventually a fixed integer $d$. For large $k$ write
\begin{equation}
 a_k=dk-D,\qquad D=\sum_{j\ge1}(d-d_j)\ge0,
 \label{eq:angular-harmonic-closure-HC8}
\end{equation}
where only finitely many summands are nonzero. By \eqref{eq:angular-harmonic-closure-HC4}, $k\mid2a_k$
for every $k$, so every sufficiently large $k$ divides $2D$. Thus
$D=0$. Each summand in \eqref{eq:angular-harmonic-closure-HC8} is nonnegative, and every increment equals
$d=a_1$. Therefore
\begin{equation}
 b_{2k}=k b_2,\qquad 2b_n\le b_{2n}=n b_2.
 \label{eq:angular-harmonic-closure-HC9}
\end{equation}
The second inequality is the empty-tuple case of \eqref{eq:angular-harmonic-closure-HC3}. Convexity and
linear growth without \eqref{eq:angular-harmonic-closure-HC4} would not suffice: $a_k=\max(0,k-1)$ is
a counterexample. The polygon argument is what removes that defect.

The two-point elimination theorem gives $b_2=0$ under its ultraviolet
shape estimate. \eqref{eq:angular-harmonic-closure-HC9} then give $b_n=0$ for every $n$.
Consequently the ordinary-chart kernel is invariant under rotations in
the chosen plane. The following argument upgrades this fact to the
actual off-diagonal distributions, including their exceptional supports.

\subsubsection{A local analytic translation lemma}

Fix a coordinate $j$ and divide the entries of a compact separated tuple
into two groups whose supports lie on opposite sides of a coordinate
hyperplane, with a positive gap. Their scalar moment is a matrix element
in the coordinate-$j$ Hilbert picture. After removing some unused gap,
it has the form
\begin{equation}
 \sigma(F_-\cdot\tau_z F_+)
   =\langle\Phi,e^{-(d+z_j)H_j}U_{z_\parallel}\Psi\rangle,
 \qquad d>0,
 \label{eq:angular-harmonic-closure-HC10}
\end{equation}
for real $z$ sufficiently small. Here the left group has been reflected
to define $\Phi$, both groups have been shifted away from their mirror,
and $d$ is the remaining normal separation. This is exactly the
compact-tuple translation dictionary, applied to products in each group.
Its vectors need not be point fields, and no coincident product is used.

Let $v$ be any real four-vector. The right side with $z=sv$ is
holomorphic for small complex $s$. The precise sufficient domain is
\begin{equation}
 d+v_j\operatorname{Re}s>
          \sum_{r\ne j}|v_r|\,|\operatorname{Im}s|.
 \label{eq:angular-harmonic-closure-HC11}
\end{equation}
Indeed the joint spectral cone $|P_{j,r}|\le H_j$ bounds the spectral
exponential by one there. On every smaller domain a positive damping
margin controls all powers of $H_j$ and $P_{j,r}$, so difference quotients
converge in operator norm. This proves holomorphy and local uniform
bounds, rather than formally translating test functions to complex
arguments. Translating the negative group instead gives the same
conclusion with the corresponding signs. In particular, any specified
half of a strict coordinate cut may be moved in any real direction,
with a holomorphic parameter near zero.

\subsubsection{Removing the exceptional distributional support}

For $n\ge2$ put
\[
 \mathcal C_n=\{(x_1,\ldots,x_n)\in(\mathbb R^4)^n:
                                      x_i\ne x_l\ (i\ne l)\}.
\]
Let $T_n$ be the off-diagonal moment distribution. For a fixed rotation
$R_u$ define its difference by
\begin{equation}
 D_u(f_1\otimes\cdots\otimes f_n)
   =\sigma(f_1,\ldots,f_n)-\sigma(R_uf_1,\ldots,R_uf_n),
 \qquad (R_uf)(x)=f(R_u^{-1}x).
 \label{eq:angular-harmonic-closure-HC12}
\end{equation}
Both terms are well-defined distributions on $\mathcal C_n$.

First suppose all original and all rotated time coordinates at a
configuration are distinct. Choose small product support boxes keeping
both time orderings strict. In those boxes approximate each smooth test
in $C^1$ by finite sums
$\alpha(x_4,x_1)\varphi(x_2,x_3)$, with supports staying inside the
boxes. For example, a product Fourier expansion with fixed compact
cutoffs gives these approximants with all first derivatives. The
rotation-invariant ordinary-chart kernel gives equality for the plate
approximants. The local moment bounds pass it to the original tests.
Local product-tensor density then says that $D_u$ vanishes on this
ordinary open set. This step alone would not rule out a distribution
concentrated at the remaining time ties.

Fix any $y\in\mathcal C_n$. For each coordinate take all threshold
hyperplanes halfway between successive distinct coordinate levels of
the finite set $\{y_1,\ldots,y_n\}$. Each threshold divides the indices
into two nonempty sets, and these finitely many cuts together separate
every pair of indices. No point lies on one of their hyperplanes.
There exist $u_0>0$ and a product neighborhood $W$ of $y$ such that,
for every $|u|<u_0$, each cut has the same index partition for $x$ and
for $R_ux$, for all $x\in W$, with a uniform positive gap in both
configurations. To see this, take the minimum distance of the finitely
many original points from these finitely many hyperplanes. Rotation is
continuous and the points are bounded, so their displacement and then
the size of $W$ can be made smaller than a quarter of that minimum.
Also take $u_0<\pi/3$. The cuts and their partitions are kept fixed
throughout the following local argument.

Fix $|u|<u_0$. At $x\in W$ let
\begin{equation}
 q(x)=\#\{i<l:x_{i,4}=x_{l,4}\}
       +\#\{i<l:(R_ux_i)_4=(R_ux_l)_4\}.
 \label{eq:angular-harmonic-closure-HC13}
\end{equation}
The set $\{q\le q_0\}$ is open: every inequality between distinct
time coordinates remains an inequality in a sufficiently small
neighborhood. We prove by induction that $D_u$ vanishes on this open
set for each integer $q_0$. The case $q_0=0$ is the ordinary-chart
argument above.

Take a point with $q=q_0>0$. Choose a tied pair in either frame and
one of the fixed common cuts that separates it. Let $A$ be the
positive side of that cut. Translate every entry in $A$ by $s e_4$.
Original cross-group time differences have slope $1$ or $-1$ in $s$;
rotated cross-group differences have slope $\cos u$ or $-\cos u$.
All slopes are nonzero. Internal differences in either group are
unchanged. For all sufficiently small nonzero real $s$, no previously
nonzero difference becomes zero, while every previously zero cross
difference becomes nonzero. At least the chosen tied pair is of this
kind, so the number in \eqref{eq:angular-harmonic-closure-HC13} decreases.

Choose one such $s_*$ and a small real interval $I$ about it. Shrink
the product neighborhood $V$ of the original point so that its translates
stay in $W$, all previously nonzero differences stay nonzero, and all
cross-group differences in both frames are nonzero for every $s\in I$.
Every configuration in the translated product boxes then has
$q\le q_0-1$. By induction, $D_u$ vanishes there as a distribution.

Now take arbitrary smooth slot tests supported in $V$. Their difference in \eqref{eq:angular-harmonic-closure-HC12}, after the group translation, is a holomorphic function of $s$
near the real segment joining zero to $I$. For the first moment term
this is \eqref{eq:angular-harmonic-closure-HC10}--\eqref{eq:angular-harmonic-closure-HC11} with $v=e_4$. For the rotated term the same
common cut separates the same two groups, and the direction is
$v=R_ue_4$. Strict gaps persist on the chosen segment; shrinking the
segment and its complex neighborhood supplies \eqref{eq:angular-harmonic-closure-HC11} in both cases.
This is why both configurations were required to share the cut. The
holomorphic function is zero on $I$, hence is zero at zero. Thus $D_u$
vanishes on every such product tensor in $V$. The explicit local
Fourier tensor density for distributions makes it vanish on all tests
in $V$.

This proves the induction step. Since $q\le n(n-1)$, the induction
terminates and $D_u$ vanishes on $W$. It excludes a singular part on
any of the exceptional time-tie sets, including configurations with
coincident plane projections but separated transverse positions. At
no point was the kernel integral assumed analytic across those sets.

\subsubsection{All angles, all planes, and the Schwartz domain}

We have proved local invariance for every sufficiently small $u$ on a
neighborhood of each point of $\mathcal C_n$. Let a compactly supported
off-diagonal test be given, and fix a finite angle $v$. The union of
the supports of its rotations between zero and $v$ is a compact subset
of $\mathcal C_n$. Cover it by finitely many neighborhoods just
constructed, and take the minimum of their positive angle bounds.
Partition $v$ into finitely many smaller steps. On each step the
difference distribution is zero on the current support, so composition
gives invariance through $v$.

Thus all separated smooth moments are invariant in the chosen plane.
Compact $C^1$ tests follow by smooth approximation on fixed separated
neighborhoods. The hypercubic group conjugates this circle of rotations
to rotations in every coordinate plane. Givens elimination of the
columns of an orthogonal matrix expresses every member of
$\operatorname{SO}(4)$ as a finite product of such plane rotations: successively
zero the entries below each diagonal position, using a rotation in the
corresponding two-coordinate plane, and finish with diagonal signs.
A coordinate reflection supplies the other component of $\operatorname{O}(4)$.

For $n=1$, the translation-invariant one-point functional is
$c\int f$ and is already invariant. Restoring such a mean to a centered
field uses a finite subset expansion on each compact tuple; this
commutes with every rotation because $\int Rf=\int f$.
Finally rotations act continuously on Schwartz space, preserve every
pairwise coincidence diagonal, and preserve flatness there. The
diagonal-flat extension and its constrained tensor density therefore
transfer the same invariance to that full test domain.

\begin{proposition}[All-orders angular closure]
\label{thm:all-orders-angular-closure}
Under the moment, cone, analytic-kernel and two-point hypotheses stated
above, the full scalar hierarchy is $\operatorname{O}(4)$ invariant, both on
separated compact $C^1$ tuples and on its unique diagonal-flat Schwartz
extension. The polygon, discrete-convexity and exceptional-support steps
require no further analytic rotation assumption.
\end{proposition}

This closes a conditional Euclidean-covariance implication. It neither
proves the ordinary interacting source estimates nor the ultraviolet
two-point law, and it does not identify the scalar cyclic theory with a
complete family of Yang--Mills observables.


% END MODULE chapters/angular_harmonic_closure_v5.tex

% BEGIN PROOF EXPANSION ward_exact/ward_displacement.tex
\subsection{Displacement calculus and the exact lattice Ward identity}
\label{ch:exact-displacement-ward}

This section derives the finite-lattice rotation identity, including
its contact terms and its shell flux. The identity is exact before
any renormalization or continuum limit. Its separate terms identify
which estimates a Ward proof of Euclidean covariance needs. In
particular the divergence of the displacement field contributes to
the flux and must be retained.

\subsubsection{Plaquette holonomies based at a bond}

Work on the periodic lattice $a\mathbb Z^4/\ell\mathbb Z^4$ with
$\ell/a$ an integer at least three. Write $\mathcal B_+$ for the
positive bonds. A configuration is a point of
$\operatorname{SU}(2)^{\mathcal B_+}$; for a reverse bond $\bar b$ set
$U_{\bar b}=U_b^{-1}$. Let $x_b$ be the midpoint of a bond and
$x_p$ the barycentre of a plaquette. A plaquette is counted once
in a sum over $p$, and a sum over its sides uses its four
positive bonds, irrespective of the orientation of its boundary.

If $b$ is either orientation of a side of $p$, let $W_{b,p}$
be the holonomy of the plaquette boundary that starts at $b_-$
and traverses $b$ first. Its first matrix factor is $U_b$.
Reversing the orientation of this boundary inverts the
holonomy and leaves its trace unchanged. Set
\begin{equation}\label{wd:quaternion-parts}
 W=w^0(W)I+i\mathbf w(W),\qquad
 w^0(W)=\tfrac12\operatorname{tr}W,\qquad
 \mathbf w(W)=\frac{W-W^*}{2i}.
\end{equation}
With the Pauli basis $t_\alpha$ and inner product
$\langle X,Y\rangle=\tfrac12\operatorname{tr}(XY)$,
$\mathbf w=\sum_{\alpha=1}^3w^\alpha t_\alpha$ and
$(w^0)^2+|\mathbf w|^2=1$. The plaquette density is
$a_p=1-w^0(W_{b,p})$, independent of the chosen side and
base point. Thus
\begin{equation}\label{wd:curvature-vector-bound}
 0\le a_p\le2,\qquad |\mathbf w(W_{b,p})|^2\le2a_p.
\end{equation}

For a positive bond, define $\partial_b^\alpha$ by differentiating
$U_b\mapsto e^{ist_\alpha}U_b$. The inverse bond changes by the
inverse of this matrix. The same definition for an oriented
negative bond is understood as a vector field on the positive
bond coordinates. Since the four bonds in a plaquette are
distinct, a derivative at one of its sides differentiates
exactly one factor in the based holonomy.

\begin{lemma}[The derivatives of the scalar and vector parts]
\label{lem:wd-holonomy-derivative}
For a side $b$ of $p$,
\begin{align}
 \partial_b^\alpha w^0_{b,p}&=-w^\alpha_{b,p},
 &\partial_b^\alpha a_p&=w^\alpha_{b,p},
       \label{wd:scalar-derivative}\\
 \partial_b^\alpha w^\gamma_{b,p}
 &=\delta_{\alpha\gamma}w^0_{b,p}
          +\sum_\delta\varepsilon_{\alpha\gamma\delta}
                                             w^\delta_{b,p},
 &\sum_\alpha\partial_b^\alpha w^\alpha_{b,p}
       &=3w^0_{b,p}.
       \label{wd:vector-derivative}
\end{align}
The derivatives of $a_p$ vanish on bonds not in its boundary.
\end{lemma}
\begin{proof}
Write $W_{b,p}=U_bR$, where $R$ has no dependence on that
bond coordinate. Its derivative is $it_\alpha W_{b,p}$.
Expand this product with
$t_\alpha t_\gamma=\delta_{\alpha\gamma}I+
i\sum_\delta\varepsilon_{\alpha\gamma\delta}t_\delta$.
Its scalar part is $-w^\alpha$ and its vector part is the
right side of \eqref{wd:vector-derivative}. The alternating
term vanishes when $\gamma=\alpha$, giving the final sum.
\end{proof}

\subsubsection{The displacement vector field}

For $\rho$ different from the direction of $b$, let
$p(b,\pm\rho)$ be the plaquette spanned by $b$ and
$\pm ae_\rho$, and abbreviate its based vector part by
$\mathbf w_{b,\pm\rho}$. Put
\begin{align}
 \mathbf u_{b,\rho}
     &=-\tfrac12(\mathbf w_{b,+\rho}-\mathbf w_{b,-\rho}),
 &\mathbf u_{b,\mu(b)}&=0,\notag\\
 \mathbf v_b[\eta]&=a^{-1}\sum_\rho\eta^\rho(x_b)
                                               \mathbf u_{b,\rho},
 &\mathsf D_\eta&=\sum_{b\in\mathcal B_+,\alpha}
                                      v_b^\alpha[\eta]\partial_b^\alpha.
 \label{wd:displacement}
\end{align}
Only the finitely many values of $\eta$ at bond midpoints
are used. No regularity in the spatial variable is needed
for this definition. Its coefficients are smooth functions
of the group variables, even if $\eta$ is merely a prescribed
list of numbers.

For $b$ a side of $p$ define the elementary density
\begin{equation}\label{wd:elementary-density}
 d_{b,p,\rho}=\langle\mathbf u_{b,\rho},\mathbf w_{b,p}\rangle.
\end{equation}
Then \cref{lem:wd-holonomy-derivative} gives
\begin{equation}\label{wd:action-on-plaquette}
 \mathsf D_\eta a_p
       =a^{-1}\sum_{b\subset\partial p,\rho}
                                  \eta^\rho(x_b)d_{b,p,\rho}.
\end{equation}
All densities here are bounded functions on a compact
configuration space, so every finite-cutoff derivative and
integral used below exists.

\subsubsection{Gauge covariance and bond reversal}

A site gauge $u$ acts by
$U_b\mapsto u(b_-)U_bu(b_+)^{-1}$. Its intermediate factors
cancel in a based loop, hence
\begin{equation}\label{wd:based-gauge-covariance}
 W_{b,p}(U^u)=u(b_-)W_{b,p}(U)u(b_-)^{-1}.
\end{equation}
The same conjugation acts on $\mathbf w$, $\mathbf u$ and
$\mathbf v$. Since the adjoint action preserves the colour
inner product, $d_{b,p,\rho}$ is gauge invariant. The
transformation rule for the bond derivative is the
corresponding adjoint transformation of its direction.
The two rules cancel in their contraction, so
\begin{equation}\label{wd:derivation-gauge-covariance}
 \mathsf D_\eta(F\circ(\,\cdot\,)^u)
                      =(\mathsf D_\eta F)\circ(\,\cdot\,)^u.
\end{equation}

The exact reversal relations are
\begin{align}
 W_{\bar b,p}&=U_b^{-1}W_{b,p}^{-1}U_b,
 &\mathbf w_{\bar b,p}
       &=-\operatorname{Ad}(U_b^{-1})\mathbf w_{b,p},\notag\\
 \mathbf u_{\bar b,\rho}
       &=-\operatorname{Ad}(U_b^{-1})\mathbf u_{b,\rho},
 &\mathbf v_{\bar b}[\eta]
       &=-\operatorname{Ad}(U_b^{-1})\mathbf v_b[\eta].
 \label{wd:reversal}
\end{align}
The first identity follows by following the inverse boundary
from the other end of the same bond; the others use inversion
and conjugation of the vector part. The midpoint is unchanged.
At a fixed configuration,
$\partial_{\bar b}^X=-\partial_b^{\operatorname{Ad}(U_b)X}$.
It follows that $\mathbf v_{\bar b}\cdot\partial_{\bar b}
=\mathbf v_b\cdot\partial_b$. Thus the vector field
$\mathsf D_\eta$ is independent of which orientation is
chosen for each unoriented bond, and
$d_{\bar b,p,\rho}=d_{b,p,\rho}$.

For a signed permutation $S$ of the spatial coordinates and
a lattice translation, the image of a transverse positive
direction can be negative. Extend $\mathbf u_b(\xi)
=\sum_\rho\xi^\rho\mathbf u_{b,\rho}$ linearly to every
$\xi\in\mathbb R^4$. Interchanging $+\rho$ and $-\rho$
changes its sign. Consequently, for the isometry
$\mathsf r(x)=Sx+v$,
\begin{equation}\label{wd:spatial-covariance}
 d_{\mathsf rb,\mathsf rp}(S\xi)(\mathsf rU)
                    =d_{b,p}(\xi)(U),\qquad
 \mathsf D_{\mathsf r\cdot\eta}(F^{\mathsf r})
                    =(\mathsf D_\eta F)^{\mathsf r},
\end{equation}
where $(\mathsf r\cdot\eta)(x)=S\eta(\mathsf r^{-1}x)$.
Bond reversals that arise in this formula are harmless by
\eqref{wd:reversal}. This verifies covariance under coordinate
reflections as well as orientation-preserving permutations.

\subsubsection{The exact Haar divergence}

The divergence of a vector field written in left-invariant
orthonormal frames is
$\operatorname{div}\mathsf D_\eta=
\sum_{b,\alpha}\partial_b^\alpha v_b^\alpha$.
By \eqref{wd:vector-derivative},
\begin{equation}\label{wd:divergence-bond-form}
 \operatorname{div}\mathsf D_\eta
  =\frac3{2a}\sum_{b\in\mathcal B_+,\rho\ne\mu(b)}
          \eta^\rho(x_b)[a_{p(b,+\rho)}-a_{p(b,-\rho)}].
\end{equation}
The coefficient three is the dimension of the colour algebra.
The sign follows from $w^0=1-a_p$ and the minus sign in
the definition of $\mathbf u$.

Define the midpoint central difference by
\begin{equation}\label{wd:midpoint-central-difference}
 \nabla^c_\rho\eta^\rho(x)
       =a^{-1}[\eta^\rho(x+\tfrac a2e_\rho)
                        -\eta^\rho(x-\tfrac a2e_\rho)].
\end{equation}
A plaquette in the $\mu\nu$ plane occurs twice in
\eqref{wd:divergence-bond-form} for each of its two directions.
For the direction $\nu$, its lower $\mu$ side has midpoint
$x_p-ae_\nu/2$ and sign plus; its upper $\mu$ side has
midpoint $x_p+ae_\nu/2$ and sign minus. The other direction
is identical. Regrouping the finite sum therefore gives
\begin{equation}\label{wd:divergence-plaquette-form}
 \operatorname{div}\mathsf D_\eta
 =-\tfrac32\sum_{p\text{ in plane }\mu\nu}
        a_p[\nabla^c_\mu\eta^\mu(x_p)
                             +\nabla^c_\nu\eta^\nu(x_p)].
\end{equation}
The notation in this formula sums over all plaquettes once;
the two displayed directions are the directions of that
plaquette and have no further implicit sum.

\begin{proposition}[Integration by parts with the sourced Wilson measure]
\label{prop:wd-haar-ibp}
For a smooth vector field $\mathsf D$ on the finite product
of groups and a smooth function $F$,
\begin{equation}\label{wd:haar-ibp}
 \int\mathsf DF\,dU=-\int F\operatorname{div}\mathsf D\,dU.
\end{equation}
For $A(c,U)=\sum_pc(p)a_p(U)$ and $Z(c)=\int e^{-A(c,U)}dU\ne0$,
the normalized, possibly complex, expectation satisfies
\par\noindent
\begin{minipage}{\linewidth}
\begin{equation}\label{wd:sourced-ibp}
 \langle\mathsf DA(c,\cdot)\rangle_c
                         =\langle\operatorname{div}\mathsf D\rangle_c.
\end{equation}
\popQED\par
\end{minipage}
\end{proposition}
\begin{proof}
Left multiplication preserves Haar measure, so
$\int\partial_b^\alpha H\,dU=0$ for smooth $H$. Apply it
to each product $v_b^\alpha F$, expand by the product rule,
and sum to prove \eqref{wd:haar-ibp}. Apply that identity
to $F=e^{-A}$ and use $\mathsf De^{-A}=-(\mathsf DA)e^{-A}$.
Division by the stipulated nonzero integral proves
\eqref{wd:sourced-ibp}. Compactness justifies differentiation
under each integral. The proof works for complex coefficients
by linearity; positivity is only needed later for estimates.
\end{proof}

\subsubsection{Orbital, tensor and remainder densities}

The four positive sides of a plaquette have lever arms
$x_b-x_p=\pm ae_\mu/2,\pm ae_\nu/2$. Set
\begin{align}
 t_\rho(p)&=\sum_{b\subset\partial p}d_{b,p,\rho},\notag\\
 s_{\lambda\rho}(p)&=\sum_{b\subset\partial p}
               a^{-1}(x_b-x_p)_\lambda d_{b,p,\rho},\notag\\
 r_\rho(p)&=t_\rho(p)-\tfrac12[a_{p+ae_\rho}-a_{p-ae_\rho}].
 \label{wd:orbital-spin-remainder}
\end{align}
These are local gauge-invariant polynomials. Under the
signed-permutation isometries, $s$ is a rank-two tensor
and $r$ a vector: this follows from
\eqref{wd:spatial-covariance}, the transformation of the
lever arm by $S$, and the signed transformation of the
central difference.

The sphere relation and $2ab\le a^2+b^2$ give
\begin{equation}\label{wd:elementary-energy-bound}
 |d_{b,p,\rho}|\le1,\qquad
 |d_{b,p,\rho}|
      \le\tfrac12[a_{p(b,+\rho)}+a_{p(b,-\rho)}]+a_p.
\end{equation}
Indeed $|\mathbf u|\le1$, while the product with
$|\mathbf w_{b,p}|$ is bounded by
$\tfrac14(|\mathbf w_+|^2+|\mathbf w_-|^2)
+\tfrac12|\mathbf w_{b,p}|^2$.
If $N(p)$ includes the plaquettes sharing a side with $p$
and the translates $p\pm ae_\rho$, these inequalities imply
\begin{equation}\label{wd:local-energy-bounds}
 |s_{\lambda\rho}(p)|+|t_\rho(p)|+|r_\rho(p)|
                            \le C\sum_{p'\in N(p)}a_{p'}.
\end{equation}
The constant is absolute and independent of $a,\ell$.
Each term depends only on bonds within $2a$ of the
plaquette barycentre in the sup norm. These bounds
identify the energy weight needed to estimate an
insertion; they do not themselves supply decay of a
connected correlation.

\subsubsection{A cutoff rotation and the product rule}

Let $\omega\in\mathfrak{so}(4)$. Choose a Lipschitz function
$\chi$ supported in $(-\ell/2+2a,\ell/2-2a)^4$ and read
it on the torus through representatives in
$[-\ell/2,\ell/2)^4$. Define $\eta(x)=\chi(x)\omega x$.
The support margin ensures that every shift used in a
nonzero term is represented by its ordinary Euclidean
shift. The linear field $\omega x$ alone is not periodic;
this support condition is essential for the next
identities.

Put
\begin{equation}\label{wd:cutoff-variation}
 h_{\omega,\chi}(p)=a^{-1}
   \sum_{b\subset\partial p,\rho}
       [\chi(x_b)-\chi(x_p)](\omega x_b)_\rho d_{b,p,\rho}.
\end{equation}
Using $\omega x_b=\omega x_p+\omega(x_b-x_p)$ in
\eqref{wd:action-on-plaquette} gives the exact product rule
\begin{align}
 \mathsf D_{\chi\omega x}a_p
  ={}&\chi(x_p)a^{-1}\sum_\rho(\omega x_p)_\rho t_\rho(p)
       \notag\\
    &+\chi(x_p)\sum_{\rho,\lambda}
                             \omega_{\rho\lambda}s_{\lambda\rho}(p)
      +h_{\omega,\chi}(p).
 \label{wd:cutoff-product-rule}
\end{align}
Only the antisymmetric part of $s$ contributes to its
contraction with $\omega$. This is an algebraic tensor
statement and does not assume rotation symmetry of the
quantum measure.

For a plaquette weight $f$ define
\begin{equation}\label{wd:rotation-central-difference}
 (\nabla^s_\omega f)(p)
   =\sum_\rho(\omega x_p)_\rho
                     \frac{f(p+ae_\rho)-f(p-ae_\rho)}{2a}.
\end{equation}
For any weight $c$, define
\begin{align}
 \mathcal S_\chi[c]&=\sum_pc(p)\chi(x_p)
                     \sum_{\rho,\lambda}
                                      \omega_{\rho\lambda}s_{\lambda\rho}(p),
       \notag\\
 \mathcal K_\chi[c]&=\sum_pc(p)\chi(x_p)a^{-1}
                             \sum_\rho(\omega x_p)_\rho r_\rho(p).
 \label{wd:spin-orbital-functionals}
\end{align}
Omit the subscript $\chi$ when its factor is absent.

The central-difference part of the first line of
\eqref{wd:cutoff-product-rule}, summed against $c$, is
\begin{equation}\label{wd:rotation-summation-parts}
 \sum_{p,\rho}(c\chi)(p)(\omega x_p)_\rho
                     \frac{a_{p+ae_\rho}-a_{p-ae_\rho}}{2a}
                   =-\mathcal O(\nabla^s_\omega(c\chi)).
\end{equation}
To verify it, shift the finite plaquette sum separately in
its plus and minus terms. In direction $\rho$, the
$\rho$ component of $\omega x$ is unchanged by that
shift because $\omega_{\rho\rho}=0$. The compact support
margin removes every seam ambiguity. This proves the
identity without a discrete product-rule approximation.

\subsubsection{The exact sourced identity}

Let $w$ be a complex plaquette weight and $c=\beta-w$,
with $\beta>0$. Assume $\chi=1$ on the Euclidean
$2a$-neighborhood of the representatives of all
plaquettes where $w\ne0$. This implies
$w\chi=w$ and $w(p)h_{\omega,\chi}(p)=0$.
Define the source-independent shell insertion
\begin{equation}\label{wd:flux-definition}
 \mathcal B_{\omega,\chi}
    =\beta\mathcal O(\nabla^s_\omega\chi)
       -\beta\sum_ph_{\omega,\chi}(p)
       +\operatorname{div}\mathsf D_{\chi\omega x}.
\end{equation}
For $Z(w)=\int e^{-A(\beta-w,U)}dU\ne0$, denote
the normalized expectation by $\langle\cdot\rangle_w$.

\begin{proposition}[Finite-lattice rotational Ward identity]
\label{prop:wd-ward-identity}
Under the stated support and normalization hypotheses,
\begin{equation}\label{wd:sourced-ward-identity}
 \langle\mathcal O(\nabla^s_\omega w)\rangle_w
   =-\langle\mathcal S_\chi[c]+\mathcal K_\chi[c]\rangle_w
                         +\langle\mathcal B_{\omega,\chi}\rangle_w.
\end{equation}
The divergence contribution has the explicit shell form
\par\noindent
\begin{minipage}{\linewidth}
\begin{equation}\label{wd:rotation-divergence}
 \operatorname{div}\mathsf D_{\chi\omega x}
 =-\tfrac32\sum_{p\text{ in plane }\mu\nu}a_p
       \sum_{\rho\in\{\mu,\nu\}}(\omega x_p)_\rho
                                              \nabla^c_\rho\chi(x_p).
\end{equation}
\popQED\par
\end{minipage}
\end{proposition}
\begin{proof}
Sum \eqref{wd:cutoff-product-rule} against $c$. Insert the
definition of $r$ and use
\eqref{wd:rotation-summation-parts}. Since
$c\chi=\beta\chi-w$ and $wh=0$, this gives
\[
 \mathsf D_{\chi\omega x}A(c,U)
  =\mathcal S_\chi[c]+\mathcal K_\chi[c]
      -\beta\mathcal O(\nabla^s_\omega\chi)
      +\mathcal O(\nabla^s_\omega w)
      +\beta\sum_ph_{\omega,\chi}(p).
\]
Take normalized expectations and apply
\eqref{wd:sourced-ibp}. Rearrangement proves
\eqref{wd:sourced-ward-identity}. For the divergence,
substitute $\eta=\chi\omega x$ in
\eqref{wd:divergence-plaquette-form}. Its component
$(\omega x)_\rho$ is constant along the $\rho$ direction,
again because $\omega_{\rho\rho}=0$. The remaining
central difference acts only on $\chi$, proving
\eqref{wd:rotation-divergence}.
\end{proof}

Every occurrence of $\chi$ in the flux is a difference
of its values at points at distance at most $a$.
In the first term of \eqref{wd:flux-definition}, split
$\chi(x_p+a e_\rho)-\chi(x_p-ae_\rho)$ at $x_p$.
The second term already has a bond-midpoint difference,
and the last is \eqref{wd:rotation-divergence}. Thus a
plaquette contributes nothing when $\chi$ is constant
on its $a$-neighborhood. This local assertion does not
license setting $\chi=1$ globally on a torus while
treating $\omega x$ as periodic.

\subsubsection{Connected derivatives and the contact terms}

Choose weights $f_1,\ldots,f_n$ and a single cutoff
$\chi$ equal to one on the prescribed neighborhood of
all their supports. Set $w(z)=\sum_i z_if_i$.
Since the group product is compact and $0\le a_p\le2$,
$Z(w(z))$ and every inserted numerator are entire
functions of $z\in\mathbb C^n$. Also $Z(0)>0$.
Continuity gives a common polydisc about zero on which
$Z$ is nonzero and all the preceding normalized
identities hold holomorphically.

For a bounded insertion $Y$ independent of $z$,
\begin{equation}\label{wd:inserted-cumulant-derivative}
 \partial_{z_1}\cdots\partial_{z_n}
       \langle Y\rangle_{w(z)}\big|_{z=0}
       =\langle Y;\mathcal O(f_1);\ldots;\mathcal O(f_n)\rangle.
\end{equation}
This follows by adding one more source $sY$ to $Z$ and
differentiating $\partial_s\log Z(s,z)$ at $s=0$.
Boundedness justifies all derivatives under the integral,
and local non-vanishing justifies the logarithm.

Differentiate \eqref{wd:sourced-ward-identity} once in
each of the $n$ source variables. The source on its
left gives $n$ terms: differentiating its explicit
factor $z_i$ leaves a cumulant with the remaining
$n-1$ sources. The same product rule applies to
$\mathcal S_\chi[\beta-w]+\mathcal K_\chi[\beta-w]$.
Since $\chi=1$ at the source plaquettes, its resulting
contact insertions are $\mathcal S[f_i]+\mathcal K[f_i]$.
Consequently
\begin{align}
 &\sum_{i=1}^nS_n^T(f_1,\ldots,\nabla^s_\omega f_i,\ldots,f_n)
       \notag\\
 &\quad=-\beta\langle(\mathcal S_\chi+\mathcal K_\chi)[1];
                          \mathcal O(f_1);\ldots;\mathcal O(f_n)\rangle
       \notag\\
 &\qquad\quad+\sum_{i=1}^n
       \langle(\mathcal S+\mathcal K)[f_i];
                                  (\mathcal O(f_j))_{j\ne i}\rangle
       \notag\\
 &\qquad\quad+\langle\mathcal B_{\omega,\chi};
                          \mathcal O(f_1);\ldots;\mathcal O(f_n)\rangle.
 \label{wd:connected-ward-identity}
\end{align}
For $n=1$ the contact term is an expectation. No
assumption of disjoint supports is used for this exact
finite-lattice identity. Such separation enters estimates
of its connected terms, rather than the algebraic
differentiation.

\subsubsection{The shell estimate and the continuum use}

For a cutoff $\chi_R$ equal to one on a ball of radius
$R$ and zero outside a ball of radius $2R$, assume
$\operatorname{Lip}\chi_R\le C/R$, $a\le R/10$, and
$2R+2a<\ell/2$, retaining the earlier torus support margin.
The difference form of the flux, together with
\eqref{wd:elementary-energy-bound}, gives
\begin{equation}\label{wd:flux-energy-bound}
 |\mathcal B_{\omega,\chi_R}|
    \le C'(\beta+1)|\omega|
       \sum_{p:\,R-3a\le|x_p|\le2R+3a}a_p.
\end{equation}
To check the factors, a difference of $\chi_R$ contributes
at most its point separation times $C/R$, cancelling
the displayed $a^{-1}$ in each term. In this shell
$|x_p|+a\le3R$, which cancels the remaining $R^{-1}$.
The bounded neighbor multiplicity in
\eqref{wd:elementary-energy-bound} only changes $C'$.
The divergence term has the same estimate with the
factor $\beta$ replaced by one.

This bound supplies integrability at fixed cutoff.
It does not make the connected shell insertion tend to
zero as $R\to\infty$: the shell contains an increasing
number of plaquettes. A Ward proof of continuum rotation
invariance must establish vanishing of the last line of
\eqref{wd:connected-ward-identity}, as well as the stated
ultraviolet estimates for the spin and orbital terms.
Clustering without a rate sufficient for the shell
sum is not such an estimate.

For $f\in C_c^3$ in a fixed bounded region, Taylor's
integral formula yields
\begin{equation}\label{wd:rotation-difference-limit}
 \|\nabla^s_\omega f-(\omega x)\cdot\nabla f\|_\infty
       \le Ca^2|\omega|\sup_{x\in\operatorname{supp}f+B_{2a}}
                                      (1+|x|)|D^3f(x)|.
\end{equation}
For \(f\in C_c^4\) supported in a fixed ball of radius \(R_f\),
the same calculation after one spatial derivative gives
\[
 \|\nabla^s_\omega f-(\omega x)\cdot\nabla f\|_{C^1}
 \le C a^2|\omega|(1+R_f+a)
                  \bigl(\|D^3f\|_\infty+\|D^4f\|_\infty\bigr).
\]
Indeed differentiation also acts on the linear coefficient
\((\omega x)_\rho\), producing a central-difference error for
\(f\); the other term is the same error for a first derivative
of \(f\). Both have the displayed bound. The error is supported
in \(\operatorname{supp}f+B_a\). On any fixed separated tuple
these enlarged supports remain separated for small \(a\).
The uniform correlation bound in the declared \(C^1\) test norm,
together with multilinearity, therefore passes the left side
of the Ward identity to the infinitesimal rotation derivative
for smooth tests.
If the spin, orbital and shell terms all vanish on
that same limiting route, the limiting functionals
satisfy the infinitesimal rotation identity. Integrating
along one-parameter rotations gives rotation covariance
on the smooth test class, and continuity extends it
to the declared test space. The exact finite-lattice
calculation identifies all terms to be controlled; it
does not replace their interacting estimates.

% END PROOF EXPANSION ward_exact/ward_displacement.tex

% BEGIN PROOF EXPANSION ward_classical/classical_invariants_wilson.tex
\subsection{Classical $\operatorname{SU}(2)$ invariants and the smooth Wilson action}
\label{ch:classical-invariants-wilson}

The rotational Ward calculation distinguishes a local density from its
integral, and a classical identity from an estimate in the lattice
measure.  We make those distinctions explicit.  First we classify the
local gauge-invariant polynomials through dimension six under the full
hypercubic group.  We prove the finite enumeration, its independence,
and the quotients by divergences and by the field equation.  We then
derive the smooth-field Wilson expansion with the same SU(2)
normalization as the lattice construction.  Its pointwise coefficient
contains a divergence which disappears only from an unweighted
compact-support integral.  None of these classical calculations bounds
that coefficient under a rough interacting lattice measure.

\subsubsection{Conventions and the local polynomial spaces}

The Lie algebra consists of traceless Hermitian matrices.  Write
$X=\vec x\cdot\sigma$, using Pauli matrices with
$\frac12\operatorname{tr}(\sigma^a\sigma^b)=\delta_{ab}$, and set
\begin{equation}
 \langle X,Y\rangle=\tfrac12\operatorname{tr}(XY),\qquad
 \langle X,Y,Z\rangle=\tfrac i2\langle X,[Y,Z]\rangle.
 \label{eq:classical-pauli-conventions}
\end{equation}
Multiplication of the three Pauli matrices gives
\begin{equation}
 XY=(\vec x\cdot\vec y)I+i(\vec x\times\vec y)\cdot\sigma,
 \quad i[X,Y]=-2(\vec x\times\vec y)\cdot\sigma,
 \quad \langle X,Y,Z\rangle=-\det(\vec x,\vec y,\vec z).
 \label{eq:classical-pauli-products}
\end{equation}
In particular the triple bracket is real and alternating.  This fixes
the sign of the cubic invariant below.

Let $A_\mu$ be a smooth connection on $\mathbb R^4$.  Define
\begin{equation}
 D_\alpha X=\partial_\alpha X+i[A_\alpha,X],\qquad
 F_{\mu\nu}=\partial_\mu A_\nu-\partial_\nu A_\mu+i[A_\mu,A_\nu],
 \qquad J^{\rm c}_\nu=\sum_\mu D_\mu F_{\mu\nu}.
 \label{eq:classical-calculus-definition}
\end{equation}
The superscript distinguishes this continuum field equation from a
current defined as a positive action gradient.  For
$S_0=\frac14\int\sum_{\mu,\nu}|F_{\mu\nu}|^2$, a
compactly supported variation gives
$\delta S_0=-\int\sum_\nu\langle J^{\rm c}_\nu,\delta A_\nu\rangle$.
Thus its sign is the opposite of that positive gradient; the equation
$J^{\rm c}=0$ is unaffected.

Commuting $\partial_\alpha+iA_\alpha$ as operators, and using
the Jacobi identity, proves
\begin{equation}
 \begin{aligned}
 [D_\alpha,D_\beta]X&=i[F_{\alpha\beta},X],\\
 D_\alpha F_{\mu\nu}+D_\mu F_{\nu\alpha}
                         +D_\nu F_{\alpha\mu}&=0,\\
 \partial_\alpha\langle X,Y\rangle
 &=\langle D_\alpha X,Y\rangle+\langle X,D_\alpha Y\rangle.
 \end{aligned}
 \label{eq:classical-calculus-identities}
\end{equation}
The same product rule holds for the triple bracket.  Under a site
gauge $u$, the connection is
$A^u_\alpha=uA_\alpha u^{-1}+i(\partial_\alpha u)u^{-1}$.
Conjugation of $\partial+iA$ shows that every iterated covariant
derivative of $F$ transforms by $X\mapsto uXu^{-1}$.

Give $D_{\alpha_1}\cdots D_{\alpha_k}F_{\mu\nu}$ dimension
$k+2$.  Let $\mathcal V_\delta$ be the real gauge-invariant
polynomials in these variables which are homogeneous of dimension
$\delta$, identified as functions of $(A,x)$.  Constant polynomials
form $\mathcal V_0$.  No naked connection variable is admitted in
these spaces.  A total derivative of dimension six will mean
\begin{equation}
 \partial\mathcal V_5
 =\left\{\sum_\rho\partial_\rho v_\rho:
                          v_\rho\in\mathcal V_5\right\}.
 \label{eq:classical-total-derivative-space}
\end{equation}
This is a quotient by local gauge-invariant currents, not a statement
that a divergence vanishes pointwise.  We keep the current ideal
separate: $\mathcal I_6$ is spanned by
\begin{equation}
 \langle J^{\rm c}_\nu,D_\alpha F_{\kappa\lambda}\rangle,
 \qquad \langle D_\alpha J^{\rm c}_\nu,F_{\kappa\lambda}\rangle.
 \label{eq:classical-current-ideal}
\end{equation}
It contains all dimension-six gauge-invariant multiples of the field
equation and its covariant derivatives in this polynomial algebra.
Indeed a factor $D^kJ^{\rm c}$ leaves dimension $3-k$ for its
cofactor.  For $k=0,1$ this gives precisely the two displayed
families; $k=2$ has no possible cofactor, and $k=3$ is a linear
colour expression, whose gauge average is zero.

An orthogonal map acts on a connection by
$(R\cdot A)_\mu(x)=\sum_\nu R_{\mu\nu}A_\nu(R^{-1}x)$.
All indices in an iterated derivative of $F$ are tensor indices.
The group $W_4$ consists of every signed coordinate permutation,
including orientation-reversing ones.  We write
$\mathcal V_\delta^{W_4}$ for scalar invariants.  For tensor
channels, $\operatorname{Hom}_{W_4}(M,\mathcal V_\delta)$ denotes
equivariant linear maps from the indicated real tensor space $M$.

\subsubsection{The finite invariant enumeration}

\begin{lemma}[Colour and index reduction]
\label{lem:classical-finite-enumeration}
The spaces below are spanned by the displayed families:
\begin{equation}
 \begin{aligned}
 \mathcal V_4&=\operatorname{span}\{\langle F_{\mu\nu},F_{\kappa\lambda}\rangle\},\\
 \mathcal V_5&=\operatorname{span}\{\langle F_{\mu\nu},D_\alpha F_{\kappa\lambda}\rangle\},\\
 \mathcal V_6&=\operatorname{span}\{\langle F_{\mu\nu},D_\alpha D_\beta F_{\kappa\lambda}\rangle,
 \langle D_\alpha F_{\mu\nu},D_\beta F_{\kappa\lambda}\rangle,
 \langle F_{\mu\nu},F_{\kappa\lambda},F_{\sigma\tau}\rangle\}.
 \end{aligned}
 \label{eq:classical-colour-spanning}
\end{equation}
For $1\le\delta\le3$, $\mathcal V_\delta=0$.
For a tensor monomial with index tuple $I$, averaging over all
coordinate signs kills it unless its set of odd-multiplicity indices
matches the set prescribed by the output channel.  A scalar requires
the empty set; a vector component $\rho$ requires $\{\rho\}$;
an off-diagonal two-tensor component $\rho\sigma$ requires
$\{\rho,\sigma\}$.
\end{lemma}
\begin{proof}
Each variable has dimension at least two, so through dimension six
a monomial has at most three colour vectors.  Average its coefficient
tensor over constant SU(2) gauges.  Conjugation acts as all of
$\operatorname{SO}(3)$ on the colour vectors.  Invariant coefficient
tensors of ranks one, two and three are respectively zero, multiples
of $\delta_{ab}$, and multiples of $\epsilon_{abc}$.  Here is an
elementary verification.  Rotations through $\pi$ about coordinate
axes kill a rank-one invariant and all off-diagonal rank-two entries;
quarter turns equate the remaining diagonal entries.  For rank three,
the same half turns require each of the three coordinate labels to
occur an odd number of times.  They therefore occur once each.
Proper signed permutations make the remaining coefficients
alternating, giving $\epsilon$.  These arguments classify coefficient
tensors before evaluating any differential identities, so additional
relations among the jets cannot enlarge the spanning space.
The possible dimensions now give exactly
\eqref{eq:classical-colour-spanning}.  Linear colour expressions,
including a single high derivative of $F$, average to zero.

A coordinate sign vector $\varepsilon$ multiplies a monomial by
$\prod_j\varepsilon_j^{n_j}$, where $n_j$ is its index count.
Averaging this character against the scalar, vector, or two-tensor
character gives the stated parity rule.  Averaging subsequently over
coordinate permutations gives orbit sums; for a fixed output
component use the subgroup fixing its distinguished indices.
This works with spanning families and does not assume that all
formal monomials are independent.
\end{proof}

The central element $-I$ acts on $\mathcal V_\delta$ by
$(-1)^\delta$.  Hence odd-dimensional scalar invariants vanish,
and the vector channel vanishes at even dimension.  At dimension
four the scalar parity patterns allow only squares of an individual
$F_{\mu\nu}$; permutation averaging gives
\begin{equation}
 \mathcal V_4^{W_4}=\mathbb R\sum_{\mu,\nu}|F_{\mu\nu}|^2,
 \qquad \mathcal V_5^{W_4}=0.
 \label{eq:classical-low-scalar-space}
\end{equation}

Define the following three symmetric two-tensors and vectors:
\begin{equation}
 \begin{aligned}
 T^1_{\alpha\beta}&=\sum_\kappa\langle F_{\alpha\kappa},F_{\beta\kappa}\rangle,
 &T^2_{\alpha\beta}&=\delta_{\alpha\beta}\sum_{\kappa,\lambda}|F_{\kappa\lambda}|^2,\\
 T^3_{\alpha\beta}&=\delta_{\alpha\beta}\sum_\kappa|F_{\alpha\kappa}|^2,
 &v^1_\rho&=\partial_\rho\sum_{\alpha,\beta}|F_{\alpha\beta}|^2,\\
 v^2_\rho&=\sum_\nu\langle F_{\rho\nu},J^{\rm c}_\nu\rangle,
 &v^3_\rho&=\partial_\rho\sum_\kappa|F_{\rho\kappa}|^2.
 \end{aligned}
 \label{eq:classical-tensors-vectors}
\end{equation}
There is no sum on a repeated free index in $T^3$ or $v^3$.

\begin{proposition}[The low-dimensional tensor channels]
\label{prop:classical-tensor-channels}
The symmetric two-tensor channel of dimension four has basis
$T^1,T^2,T^3$.  Its antisymmetric channel is zero.
The vector channel of dimension five has basis $v^1,v^2,v^3$,
and
\begin{equation}
 v^1_\rho=\sum_\alpha\partial_\alpha T^2_{\alpha\rho},\quad
 v^3_\rho=\sum_\alpha\partial_\alpha T^3_{\alpha\rho},\quad
 v^2_\rho=\sum_\alpha\partial_\alpha
              (T^1_{\alpha\rho}-\tfrac14T^2_{\alpha\rho}).
 \label{eq:classical-vector-divergences}
\end{equation}
Moreover $\Lambda^2\mathbb R^4$ is irreducible for the full $W_4$,
and has no equivariant image in $\mathbb R$, $\mathcal V_4$, or
$\mathcal V_5$.
\end{proposition}
\begin{proof}
For a diagonal output component $\rho\rho$, sign parity requires
even counts.  In a product of two $F$'s this leaves
$\sum_\kappa|F_{\rho\kappa}|^2$ and the sum of the squares whose
indices avoid $\rho$.  These are the same span as the diagonal
components of $T^2$ and $T^3$.  For $\rho\ne\sigma$, the only
nonzero parity pattern is $\rho\sigma\kappa\kappa$ with
$\kappa$ different from the two distinguished indices.  Its orbit
sum is $\sum_\kappa\langle F_{\rho\kappa},F_{\sigma\kappa}\rangle$.
It is symmetric in $\rho,\sigma$, so the antisymmetric channel is
zero, and the full two-tensor space is the three-dimensional span
claimed, subject to the independence verification below.

For the vector component $\rho$, a dimension-five monomial has
five indices.  The odd-set rule gives only multiplicities $(3,2)$
or $(1,2,2)$; $(5)$ is zero by antisymmetry of $F$.  Up to signs
and orbit multiplicities the resulting list is
\begin{equation}
 \begin{aligned}
 a_\rho&=\sum_\kappa\langle F_{\rho\kappa},D_\rho F_{\rho\kappa}\rangle,\\
 b^1_\rho&=\sum_{\sigma,\tau\ne\rho}
                         \langle F_{\sigma\tau},D_\rho F_{\sigma\tau}\rangle,\\
 b^2_\rho&=\sum_{\sigma,\tau\ne\rho,\,\sigma\ne\tau}
                         \langle F_{\rho\sigma},D_\tau F_{\tau\sigma}\rangle,\\
 b^3_\rho&=\sum_{\sigma,\tau\ne\rho,\,\sigma\ne\tau}
                         \langle F_{\sigma\tau},D_\sigma F_{\rho\tau}\rangle.
 \end{aligned}
 \label{eq:classical-vector-orbits}
\end{equation}
This list is exhaustive: in the $(3,2)$ pattern the derivative
has the distinguished index and both antisymmetric pairs contain
it; in the $(1,2,2)$ pattern it either occurs on the derivative,
or on one of the two curvature factors, giving the last three
entries.  The product rule and Bianchi identity give
\begin{equation}
 v^3=2a,\qquad v^1=2b^1+4a,\qquad v^2=a+b^2,
 \qquad \tfrac14v^1=a+b^3.
 \label{eq:classical-vector-orbit-reduction}
\end{equation}
For the last identity contract Bianchi with $F_{\kappa\lambda}$
and sum: the two terms with a derivative on an index of that
curvature are equal, while the remaining one is half the derivative
of $\sum|F|^2$.  Explicitly
\begin{equation}
 \sum_{\kappa,\lambda}\langle F_{\kappa\lambda},D_\kappa F_{\rho\lambda}\rangle
 =\tfrac12\sum_{\kappa,\lambda}\langle F_{\kappa\lambda},D_\rho F_{\kappa\lambda}\rangle
 =\tfrac14v^1_\rho.
 \label{eq:classical-contracted-bianchi}
\end{equation}
The left side is $a_\rho+b^3_\rho$.
The same contracted identity, applied to the divergence of $T^1$,
proves the third formula in \eqref{eq:classical-vector-divergences};
the other two are immediate product rules.

Independence can be checked on genuine connections rather than
formal jets.  Multiply all polynomials in the following tests by a
smooth cutoff equal to one near zero, and set $A_\mu=q_\mu\sigma^3$.
At zero, only the displayed $q_\mu$ are nonzero:
\begin{equation}
 \begin{array}{c|c}
 q&(v^1_1,v^2_1,v^3_1)\\ \hline
 q_2=x_1+x_1^2/2&(4,1,2)\\
 q_3=x_2+x_1x_2&(4,0,0)\\
 q_2=x_1+x_3^2/2&(0,1,0)
 \end{array}
 \label{eq:classical-vector-independence-tests}
\end{equation}
The determinant is eight.  For the tensor channel, the field
$q_3=x_1+x_2$ has $T^1_{12}=1$, while $T^2_{12}=T^3_{12}=0$.
The field $q_2=x_1$ has $(T^1_{33},T^2_{33},T^3_{33})=(0,2,0)$
and $(T^1_{11},T^2_{11},T^3_{11})=(1,2,1)$.
These three evaluations prove tensor independence.

Finally the sign subgroup acts on the six coordinate wedges by
six distinct characters $\varepsilon_\rho\varepsilon_\sigma$.
Averaging against one character projects onto its one-dimensional
coordinate wedge.  Any invariant subspace containing a nonzero
vector therefore contains a coordinate wedge.  Permutations act
transitively on the six wedges, so it contains them all.  This
proves irreducibility.  Its image in the trivial representation is
zero, its dimension-four image is the antisymmetric channel already
computed, and its dimension-five image is zero by the action of
the central element $-I$.
\end{proof}

\subsubsection{Dimension six as densities}

Define
\begin{equation}
 \begin{aligned}
 O_1&=\sum_{\mu,\nu}|D_\mu F_{\mu\nu}|^2,&
 O_2&=\sum_{\alpha,\mu,\nu}|D_\alpha F_{\mu\nu}|^2,\\
 O_3&=\sum_\nu|J^{\rm c}_\nu|^2,&
 O_F&=\sum_{\mu,\nu,\rho}
                       \langle F_{\mu\nu},F_{\nu\rho},F_{\rho\mu}\rangle,\\
 Q&=\sum_\alpha\partial_\alpha^2\sum_\mu|F_{\alpha\mu}|^2,
 &d_i&=\sum_\rho\partial_\rho v^i_\rho.
 \end{aligned}
 \label{eq:classical-dimension-six-operators}
\end{equation}
In $O_1$ the derivative index is tied to one curvature index before
squaring.  It is not $O_3$, which first sums the derivative
contributions.  In particular $d_3=Q$.

\begin{lemma}[An exact off-shell identity]
\label{lem:classical-offshell-identity}
For every smooth connection, pointwise,
\par\noindent
\begin{minipage}{\linewidth}
\begin{equation}
 O_2=2O_3-4O_F+\tfrac12d_1-2d_2.
 \label{eq:classical-offshell-relation}
\end{equation}
\popQED\par
\end{minipage}
\end{lemma}
\begin{proof}
Differentiate Bianchi and commute the two covariant derivatives:
\begin{equation}
 \sum_\rho D_\rho D_\rho F_{\alpha\beta}
 =D_\alpha J^{\rm c}_\beta-D_\beta J^{\rm c}_\alpha
                  +2i\sum_\rho[F_{\rho\alpha},F_{\rho\beta}].
 \label{eq:classical-curvature-laplacian}
\end{equation}
Pair with $F_{\alpha\beta}$ and sum.  Antisymmetry gives twice
$T=\sum_{\alpha,\beta}\langle F_{\alpha\beta},D_\alpha J^{\rm c}_\beta\rangle$
from the first two terms.  The commutator term is $4O_F$ by
\eqref{eq:classical-pauli-conventions} and relabeling the three
indices.  The product rules give
\begin{equation}
 d_1=2O_2+4T+8O_F,\qquad d_2=O_3+T.
 \label{eq:classical-divergence-identities}
\end{equation}
Eliminate $T$.  No integration or field equation was used.
\end{proof}

\begin{proposition}[Complete off-shell scalar basis]
\label{prop:classical-six-density-basis}
The space $\mathcal V_6^{W_4}$ has basis
\begin{equation}
                  O_1,\ O_2,\ O_3,\ d_1,\ d_2,\ Q.
 \label{eq:classical-six-density-basis}
\end{equation}
The cubic-in-$F$ invariant subspace is exactly $\mathbb RO_F$.
\end{proposition}
\begin{proof}
We first give an exhaustive reduction.  In a cubic curvature
monomial regard each antisymmetric pair as an edge between its
two coordinate labels.  The sign rule requires even degree at
every label.  Three non-loop edges can then form only a triangle;
the colour triple bracket gives its oriented contraction $O_F$.
Permutation averaging changes only its multiplicity and sign.
Thus the cubic subspace is at most one-dimensional.

For the $DF\,DF$ terms the six index counts are $(4,2)$ or
$(2,2,2)$; the all-equal case is zero.  The full orbit list, up
to signs and factors, is
\begin{equation}
 \begin{aligned}
 O_1,&\qquad S_1=\sum_{a,b,c}^{\ne}|D_aF_{bc}|^2,\\
 S_2&=\sum_{a,b,c}^{\ne}\langle D_aF_{ac},D_bF_{bc}\rangle,
 \qquad S_3=\sum_{a,b,c}^{\ne}\langle D_aF_{bc},D_bF_{ac}\rangle.
 \end{aligned}
 \label{eq:classical-six-index-orbits}
\end{equation}
Here the superscript means that the three labels are distinct.
The $(4,2)$ pattern gives $O_1$.  In $(2,2,2)$, either a
derivative repeats a curvature index in each factor, giving $S_2$,
or all three labels occur in each factor.  Equal derivative labels
then give $S_1$ and different ones give $S_3$.
Expanding the sums and contracting Bianchi yields
\begin{equation}
 O_2=S_1+2O_1,\qquad O_3=O_1+S_2,\qquad S_1=2S_3.
 \label{eq:classical-six-index-reduction}
\end{equation}
For the last equality multiply
$D_aF_{bc}=D_bF_{ac}-D_cF_{ab}$ by $D_aF_{bc}$ and sum
over distinct indices.  Interchanging $b,c$ makes the second
cross term the negative of the first, giving exactly $2S_3$.

It remains to reduce $F\,DDF$.  The pointwise identity
\begin{equation}
 \langle F_{\mu\nu},D_\alpha D_\beta F_{\kappa\lambda}\rangle
 =\partial_\alpha\langle F_{\mu\nu},D_\beta F_{\kappa\lambda}\rangle
       -\langle D_\alpha F_{\mu\nu},D_\beta F_{\kappa\lambda}\rangle
 \label{eq:classical-second-derivative-reduction}
\end{equation}
requires no commutation of $D_\alpha,D_\beta$.
After averaging over $W_4$, the first term is the divergence of
an equivariant vector of dimension five.  By
\cref{prop:classical-tensor-channels} it lies in
$\operatorname{span}\{d_1,d_2,Q\}$.  The second term was just
enumerated.  The scalar space is therefore spanned by
$O_1,O_2,O_3,O_F,d_1,d_2,Q$, and
\eqref{eq:classical-offshell-relation} reduces this to the six
members in \eqref{eq:classical-six-density-basis}.

We verify all six dimensions on actual compactly supported
connections.  Use the commuting cutoff connections of the previous
proof.  The following three choices give the indicated first three
columns; since $F(0)=0$, their last three columns are
$(d_1,d_2,Q)=(2O_2,O_3,2O_1)$:
\begin{equation}
 \begin{array}{c|ccc}
 q&O_1&O_2&O_3\\ \hline
 q_2=x_1^2/2&1&2&1\\
 q_3=x_1x_2&0&4&0\\
 q_2=(x_1^2-x_3^2)/2&2&4&0
 \end{array}
 \label{eq:classical-density-first-tests}
\end{equation}
Their three-by-three determinant is $-8$.
Two further commuting tests, $q_2=x_1+x_1^3/6$ and
$q_2=x_1+x_1x_3^2/2$, have $DF(0)=0$ and give respectively
$(d_1,d_2,Q)=(4,1,2)$ and $(4,1,0)$.
For the last test choose constant antisymmetric $f_{\mu\nu}$ with
$f_{12}=\sigma^1$, $f_{23}=\sigma^2$, $f_{31}=\sigma^3$ and
$f_{\mu4}=0$, and put
$A_\mu=\frac12\sum_\nu x_\nu f_{\nu\mu}$ near zero.
Then $F(0)=f$, $DF(0)=0$, and
$O_F(0)=-6$.  Moreover
$D_\mu^2F_{\mu\nu}(0)=0$: the ordinary second derivative of
$i[A_\mu,A_\nu]$ in the $\mu$ direction has a zero factor,
and $A(0)=\partial_\mu A_\mu(0)=0$.
Thus $Q(0)=0$.  The curvature Laplacian and
\eqref{eq:classical-divergence-identities} give
$(d_1,d_2,Q)=(-12,9,0)$.
These last three tests have $O_1=O_2=O_3=0$ and divergence
matrix
\begin{equation}
 \begin{pmatrix}4&1&2\\4&1&0\\-12&9&0\end{pmatrix},
 \qquad \det=96.
 \label{eq:classical-density-divergence-tests}
\end{equation}
The full six-by-six evaluation matrix is block triangular with
determinant $-768$.  This proves independence as functions of
connections, and also shows $O_F\ne0$.  The spanning result is
therefore a basis result, not merely a list of possible operators.
\end{proof}

\subsubsection{Rotational scalars, divergences and the field equation}

\begin{proposition}[The two density obstructions and the single integrated obstruction]
\label{prop:classical-rotation-quotients}
One has the direct sum
\begin{equation}
 \mathcal V_6^{W_4}
 =\operatorname{span}\{O_2,O_3,d_1,d_2\}
                    \oplus\mathbb RO_1\oplus\mathbb RQ,
 \qquad
 \mathcal V_6^{\operatorname{O}(4)}=\operatorname{span}\{O_2,O_3,d_1,d_2\}.
 \label{eq:classical-rotation-density-splitting}
\end{equation}
Modulo total derivatives the scalar space has basis $O_1,O_2,O_3$.
Modulo rotational scalars and total derivatives its dimension is one,
represented by $O_1$.  This last obstruction remains nonzero after
also quotienting by $\mathcal I_6$.
Modulo divergences and $\mathcal I_6$ alone, a basis is given by
the classes of $O_1,O_2$, with $O_2\equiv-4O_F$.
For the dimension-four symmetric tensor channel, a combination
$c_1T^1+c_2T^2+c_3T^3$ is $\operatorname{O}(4)$-equivariant exactly
when $c_3=0$.
\end{proposition}
\begin{proof}
Complete contraction of indices shows that $O_2,O_3,O_F$ are
rotational scalars and that $T^1,T^2,v^1,v^2$ have their stated
full orthogonal tensor types.  Their divergences $d_1,d_2$ are
therefore scalars.  To prove that no further linear combination
is a scalar, fix a radial cutoff equal to one near zero and use
\begin{equation}
 A_3(x)=\phi(p\cdot x)\sigma^3,\qquad A_\mu=0\quad(\mu\ne3),
 \qquad |p|=1,\ p_3=0
 \label{eq:classical-rotated-test-family}
\end{equation}
near zero.  Rotations fixing $e_3$ take this family from $p$ to
$Rp$.  At zero direct differentiation gives
\begin{equation}
 O_1=\phi''(0)^2\sum_\alpha p_\alpha^4,
 \qquad Q=2\{\phi''(0)^2+\phi'(0)\phi'''(0)\}
                                      \sum_\alpha p_\alpha^4.
 \label{eq:classical-rotated-test-values}
\end{equation}
The factor $\sum p_\alpha^4$ changes from one to one half under
a rotation taking $e_1$ to $(e_1+e_2)/\sqrt2$.
For $\phi(t)=t^2/2$, invariance of $aO_1+bQ$ would force
$a+2b=0$; for $\phi(t)=t+t^3/6$ it would force $b=0$.
Hence $a=b=0$.  Together with the six-dimensional basis this
proves \eqref{eq:classical-rotation-density-splitting}.

The $W_4$-scalar part of $\partial\mathcal V_5$ equals
$\operatorname{span}\{d_1,d_2,Q\}$.  Indeed average a current
over the finite group with its vector action; the divergence of
the average is the scalar average of its divergence.  The vector
classification then applies.  These three divergences are
independent by the density basis.  Quotienting gives the asserted
three-dimensional basis.  The off-shell identity also gives
$O_2\equiv2O_3-4O_F$ in this quotient.

The same sign and permutation enumeration in
\eqref{eq:classical-current-ideal} gives
\begin{equation}
 \mathcal I_6^{W_4}=\operatorname{span}\{O_3,d_2\}.
 \label{eq:classical-current-ideal-invariants}
\end{equation}
For $\langle J^{\rm c}_\nu,D_\alpha F_{\kappa\lambda}\rangle$
the only nonzero even index pattern contracts $\nu$ with one
curvature index and $\alpha$ with the other, yielding $O_3$.
For $\langle D_\alpha J^{\rm c}_\nu,F_{\kappa\lambda}\rangle$
the same pattern yields $d_2-O_3$ by the product rule.
The alternative pairing sets the two curvature indices equal and
is zero.  Both surviving invariants are rotational scalars.
If an invariant belonged to a sum of the rotation, divergence and
current subspaces, averaging that membership over $W_4$ would
give precisely these invariant subspaces.  Thus the extra quotient
cannot remove the class of $O_1$.  This proves the on-shell
quotient assertion without setting $J^{\rm c}$ to zero inside an
off-shell density identity.
The divergence subspace together with the current ideal has basis
$d_1,d_2,Q,O_3$, by the six-density basis.  Its quotient is
therefore two-dimensional, with basis $O_1,O_2$; the exact
off-shell relation reduces there to $O_2=-4O_F$.  In particular
$O_F$ is not a divergence, even after the field-equation quotient.

For the tensor assertion, $T^1,T^2$ already have full orthogonal
covariance.  In the test family with $\phi'(0)=1$,
$T^3_{\alpha\alpha}=p_\alpha^2+\delta_{\alpha3}$ and every
off-diagonal entry is zero.  At $p=e_1$, rotating the tensor
$\operatorname{diag}(1,0,1,0)$ by $\pi/4$ in the first two
coordinates produces a nonzero off-diagonal entry.  The tensor
$T^3$ evaluated on the rotated connection still has zero
off-diagonal entries.  It is therefore not equivariant, and
subtracting the equivariant $T^1,T^2$ proves the claim.
\end{proof}

The conclusion has three distinct meanings.  There are two
nonrotational directions as local densities, represented by $O_1$
and $Q$.  In the integral modulo the specified divergences only
$O_1$ remains.  In the further algebraic quotient by the field
equation it still remains.  None of these statements removes a
weighted divergence: for a compactly supported test $h$,
\begin{equation}
 \int hQ=\sum_\alpha\int(\partial_\alpha^2h)
                                  \sum_\mu|F_{\alpha\mu}|^2.
 \label{eq:classical-weighted-divergence}
\end{equation}
This term must be retained when the local Ward calculation uses a
nonconstant cutoff.

\subsubsection{Parallel transport and centred plaquettes}

For a piecewise $C^1$ path $\gamma:[0,T]\to\mathbb R^4$, define
its transport by the right-ordered differential equation
\begin{equation}
 P'(t)=iP(t)\eta(t),\qquad P(0)=I,
 \qquad \eta(t)=\sum_\mu A_\mu(\gamma(t))\dot\gamma_\mu(t).
 \label{eq:classical-transport-ode}
\end{equation}
Its value at $T$ is denoted $\Pi_\gamma[A]$.  Iterating the
integral equation gives the absolutely convergent series
\begin{equation}
 \Pi_\gamma=I+\sum_{n\ge1}i^n
  \int_{0<t_1<\cdots<t_n<T}\eta(t_1)\cdots\eta(t_n)
                                                 \,dt_1\cdots dt_n.
 \label{eq:classical-transport-series}
\end{equation}
The $n$th term is bounded by $(\int|\eta|)^n/n!$, proving
existence and allowing termwise differentiation for smooth finite
parameter families on compact parameter sets.  Uniqueness follows
by iterating the integral equation for a difference, with the same
factorial bound.  Since $\eta$ is Hermitian and traceless,
$(PP^*)'=0$ and $(\det P)'=i\operatorname{tr}(\eta)\det P=0$.
Thus the transport is in SU(2).

Concatenation multiplies transports in path order; reversing a path
inverts its transport.  These follow either from uniqueness or by
reversing the integral equation.  Differentiating
$u(\gamma(0))P(t)u(\gamma(t))^{-1}$ shows that
\begin{equation}
 \Pi_\gamma[A^u]=u(\gamma(0))\Pi_\gamma[A]u(\gamma(T))^{-1}.
 \label{eq:classical-transport-gauge-covariance}
\end{equation}
The sign in $A^u$ is thereby checked against the transport
ordering, rather than chosen independently.  The same differentiation
gives
\begin{equation}
 \frac d{dt}\{P(t)X(\gamma(t))P(t)^{-1}\}
 =P(t)\left(\sum_\mu\dot\gamma_\mu(t)D_\mu X(\gamma(t))\right)P(t)^{-1}.
 \label{eq:classical-transport-covariant-derivative}
\end{equation}

Fix $\mu<\nu$.  Let $c(t)$, $0\le t\le4$, traverse the
boundary of the unit square centred at zero in the $(\mu,\nu)$
plane, starting at $-(e_\mu+e_\nu)/2$ and following
$+e_\mu,+e_\nu,-e_\mu,-e_\nu$ at unit speed.  Extend it
periodically.  Thus $c(t+2)=-c(t)$.  Define
\begin{equation}
 W_{\mu\nu}(s,y)=\Pi_{y+sc}[A],\qquad
 G_{\mu\nu}(s,y)=1-\tfrac12\operatorname{tr}W_{\mu\nu}(s,y).
 \label{eq:classical-centred-plaquette}
\end{equation}
These are smooth functions of $(s,y)$ for real $s$, including
$s=0$.  Apply the ODE separately on the four sides to see the
smooth parameter dependence despite the corners.  The trace is real
for SU(2).  Gauge transformations conjugate $W$ at its base point,
so $G$ is gauge invariant.

The centred plaquette has the exact parity
\begin{equation}
                         G_{\mu\nu}(-s,y)=G_{\mu\nu}(s,y).
 \label{eq:classical-centred-parity}
\end{equation}
Indeed changing $s$ to $-s$ starts the same oriented square halfway
round its boundary.  If its two half-transports are $P_1,P_2$,
the two full transports are $P_1P_2$ and $P_2P_1$, whose traces
are equal.  This is parity of the centred traced plaquette; it
does not assert that a matrix-valued transport based at a fixed
corner is itself even.

\begin{lemma}[The pointwise plaquette coefficients]
\label{lem:classical-plaquette-expansion}
For every smooth connection,
\begin{equation}
 G_{\mu\nu}(s,y)
 =\frac{s^4}{2}|F_{\mu\nu}(y)|^2
  +\frac{s^6}{24}
       \langle F_{\mu\nu},(D_\mu^2+D_\nu^2)F_{\mu\nu}\rangle(y)
  +O(s^8).
 \label{eq:classical-pointwise-plaquette-expansion}
\end{equation}
On a compact set of $y$ the remainder is uniform.  For a compactly
supported connection it has a common compact support when $|s|\le1$.
\end{lemma}
\begin{proof}
Use radial gauge about the selected point $y$.  Let $u(y+z)$ be
the transport of $A$ along the straight segment from $y$ to
$y+z$, and put $B=A^u$.  By
\eqref{eq:classical-transport-gauge-covariance}, every radial
segment for $B$ has transport one.  Hence
\begin{equation}
 z\cdot B(y+z)=0,\qquad B(y)=0,\qquad
 \partial_\lambda B_\kappa(y)=\tfrac12F_{\lambda\kappa}(y).
 \label{eq:classical-radial-gauge-jets}
\end{equation}
The last two conclusions follow by differentiating the first
identity once and twice, and then using the definition of $F$.
This gauge need not be compactly supported; it is used only to
identify coefficients at $y$.  The uniform remainder will be
estimated using the original plaquette function.

Put $f=F_{\mu\nu}(y)$.  In the series
\eqref{eq:classical-transport-series}, the coefficient on the
square is
\begin{equation}
 \eta_s(t)=s\sum_\kappa B_\kappa(y+sc(t))\dot c_\kappa(t)
 =s^2\eta_0(t)+O(s^3),\qquad
 \eta_0(t)=\tfrac12(c_\mu\dot c_\nu-c_\nu\dot c_\mu)f.
 \label{eq:classical-radial-square-coefficient}
\end{equation}
All leading matrices are multiples of the single $f$.
Let $H=\int_0^4\eta_s(t)\,dt$ and let $W_n$ denote the
ordered integral of $n$ coefficient factors, without $i^n$.
The trace of $H$ is zero.  Symmetry of $\operatorname{tr}(XY)$
gives $\operatorname{tr}W_2=\frac12\operatorname{tr}H^2$.
The order-$s^6$ part of $\operatorname{tr}W_3$ is zero, because
it is proportional to $\operatorname{tr}f^3=0$; replacing one
leading coefficient by its $O(s^3)$ remainder costs $O(s^7)$.
All $n\ge4$ terms cost $O(s^8)$ by the factorial estimate.
Consequently
\begin{equation}
                         G_{\mu\nu}(s,y)=\tfrac12|H|^2+O(s^7).
 \label{eq:classical-square-trace-reduction}
\end{equation}
This is where the SU(2) trace algebra is used.

Ordinary Stokes' formula in the square gives
$H=\int_{[-s/2,s/2]^2}\psi(y+z_\mu e_\mu+z_\nu e_\nu)
\,dz_\mu dz_\nu$, where
$\psi=\partial_\mu B_\nu-\partial_\nu B_\mu$.
Taylor expansion on the symmetric square yields
\begin{equation}
 H=s^2\psi(y)+\frac{s^4}{24}
                       (\partial_\mu^2+\partial_\nu^2)\psi(y)+O(s^6).
 \label{eq:classical-square-integral}
\end{equation}
The odd moments vanish and each quadratic moment equals $s^4/12$.
At $y$ one has $\psi=f$.  Also, for $i=\mu$ or $\nu$,
$\partial_i^2[B_\mu,B_\nu](y)=0$: both potentials vanish,
and in the remaining product one factor is
$\partial_iB_i(y)=0$.  For the same reason
$D_i^2F_{\mu\nu}(y)=\partial_i^2F_{\mu\nu}(y)$.
Therefore
\begin{equation}
 H=s^2f+\frac{s^4}{24}(D_\mu^2+D_\nu^2)F_{\mu\nu}(y)+O(s^6).
 \label{eq:classical-square-covariant-coefficient}
\end{equation}
Insert this into \eqref{eq:classical-square-trace-reduction}.
It identifies the coefficients of orders four and six, with an
initial $O(s^7)$ remainder.  The original function $G(s,y)$ is
smooth and even.  Thus all odd Taylor coefficients vanish and
Taylor's theorem gives an $O(s^8)$ remainder, bounded by
$|s|^8\sup_{|t|\le1}|\partial_s^8G(t,y)|/8!$ on $|s|\le1$.
On compact sets this supremum is finite.  If $A$ has compact
support, any square of side at most one sufficiently far from
that support has all its transports equal to one.  The plaquette
function and all coefficients then vanish.  The required uniform
support and remainder follow on the fixed enlarged compact set.
\end{proof}

There is a direct noncommuting check on both coefficients.  For a
connection constant near the square, with $A_1=r\sigma^1$,
$A_2=r\sigma^2$ and the other components zero, multiplication of
the four exponentials gives
\begin{equation}
 G_{12}(s,y)=2\sin^4(sr)
           =2r^4s^4-\tfrac43r^6s^6+O(s^8).
 \label{eq:classical-constant-nonabelian-check}
\end{equation}
Here $F_{12}=-2r^2\sigma^3$ and
$D_1^2F_{12}=D_2^2F_{12}=-4r^2F_{12}$, giving exactly the
two coefficients in \eqref{eq:classical-pointwise-plaquette-expansion}.
Ordinary derivatives of $F$ vanish in this test, so replacing the
covariant derivatives by ordinary ones would fail the check.

\subsubsection{Summing over the lattice with a controlled remainder}

For a compactly supported connection let its support lie inside
$(-\ell/2+1,\ell/2-1)^4$.  Assume $0<a\le1$ and
$\ell/(2a)\in\mathbb N$.  Define the torus link variables by
transport along each positive bond.  Since $A$ vanishes in the
boundary collar, a plaquette crossing a torus face agrees with the
corresponding square in $\mathbb R^4$.  The complete unweighted
action is therefore exactly
\begin{equation}
 A(U)=\sum_{\mu<\nu}\sum_{x\in a\mathbb Z^4}
 G_{\mu\nu}\left(a,x+\tfrac a2(e_\mu+e_\nu)\right).
 \label{eq:classical-action-as-squares}
\end{equation}
Only finitely many summands are nonzero.  The same formula with a
weight at the barycentre defines $A(h,U)$.

We record the quadrature estimate so the error in the action is
quantitative.  For $\varphi\in C_c^\infty(\mathbb R^4)$ and
any offset $v\in\mathbb R^4$,
\begin{equation}
 \left|a^4\sum_{n\in\mathbb Z^4}\varphi(a(n+v))
                      -\int\varphi\right|
 \le Ca^5\max_i\|\partial_i^5\varphi\|_{L^1}.
 \label{eq:classical-shifted-quadrature}
\end{equation}
To prove it, periodize $r\mapsto\varphi(ar)$ on the unit torus.
Its Fourier coefficient at $k\in\mathbb Z^4$ is
$a^{-4}\widehat\varphi(2\pi k/a)$, with
$\widehat\varphi(\xi)=\int e^{-i\xi\cdot x}\varphi(x)\,dx$.
The Fourier series converges absolutely after five integrations by
parts.  Evaluating it at $v$ gives
\begin{equation}
 a^4\sum_n\varphi(a(n+v))
       =\sum_{k\in\mathbb Z^4}e^{2\pi ik\cdot v}
                                      \widehat\varphi(2\pi k/a).
 \label{eq:classical-shifted-poisson}
\end{equation}
For $k\ne0$ choose a coordinate with $|k_i|=|k|_\infty$ and
integrate five times in that coordinate.  The term is bounded by
$(a/(2\pi|k|_\infty))^5\|\partial_i^5\varphi\|_1$.
The number of points in the shell $|k|_\infty=m$ is $O(m^3)$,
so the nonzero-frequency sum converges and proves
\eqref{eq:classical-shifted-quadrature}.  The offset occurs only
in phases, giving the asserted uniformity.

\begin{proposition}[The weighted and unweighted smooth Wilson expansion]
\label{prop:classical-smooth-wilson-expansion}
For a fixed smooth compactly supported $A$ and a smooth weight $h$,
bounded with its derivatives near that support, one has
\begin{equation}
 A(h,U)=\frac14\int h\sum_{\mu,\nu}|F_{\mu\nu}|^2
            +a^2\int h\left(-\frac{O_1}{24}+\frac Q{48}\right)
            +\mathcal E_{A,h}(a),
 \qquad |\mathcal E_{A,h}(a)|\le C_{A,h}a^4.
 \label{eq:classical-weighted-wilson-expansion}
\end{equation}
The constant depends on finitely many smooth norms and a support
bound for $A,h$, and is independent of the torus side, the number
of sites and the coupling.  In particular,
\begin{equation}
 \beta A(U)=\frac\beta4\int\sum_{\mu,\nu}|F_{\mu\nu}|^2
            -\frac{a^2\beta}{24}\int O_1+O_A(\beta a^4).
 \label{eq:classical-unweighted-wilson-expansion}
\end{equation}
The divergence is removed only in this unweighted formula.
\end{proposition}
\begin{proof}
Denote the coefficients in
\eqref{eq:classical-pointwise-plaquette-expansion} by
$c_4^{\mu\nu}$ and $c_6^{\mu\nu}$.  They and the remainders
have support in one fixed compact set $K$.  There are at most
$C(\operatorname{diam}K+1)^4a^{-4}$ relevant mesh points.
The sum of the pointwise $O(a^8)$ remainders is consequently
$O_{A,h}(a^4)$.

Apply \eqref{eq:classical-shifted-quadrature} to
$h c_4^{\mu\nu}$ and $h c_6^{\mu\nu}$, with
$v=(e_\mu+e_\nu)/2$.  The errors in replacing their sums by
integrals are $O_{A,h}(a^5)$ and $O_{A,h}(a^7)$ after the
additional $a^2$ factor.  Both are bounded by $C_{A,h}a^4$.
The leading coefficient satisfies
$\sum_{\mu<\nu}c_4^{\mu\nu}=\frac14\sum_{\mu,\nu}|F_{\mu\nu}|^2$.
For the next coefficient set
\begin{equation}
 Y=\sum_{\mu,\nu}\langle F_{\mu\nu},D_\mu^2F_{\mu\nu}\rangle.
 \label{eq:classical-wilson-Y}
\end{equation}
Pairing the two ordered orientations gives
$\sum_{\mu<\nu}c_6^{\mu\nu}=Y/24$ exactly.
Applying the product rule twice gives $Q=2O_1+2Y$.
Thus the summed pointwise coefficient is
$-O_1/24+Q/48$, proving
\eqref{eq:classical-weighted-wilson-expansion}.
For $h=1$, the compact support makes $\int Q=0$.
Multiplication by $\beta$ proves
\eqref{eq:classical-unweighted-wilson-expansion} with its stated
dependence on $\beta$.
\end{proof}

Equivalently the weighted second coefficient is
\begin{equation}
 -\frac1{24}\int h O_1
  +\frac1{48}\sum_\mu\int(\partial_\mu^2h)
                                      \sum_\nu|F_{\mu\nu}|^2.
 \label{eq:classical-weighted-second-coefficient}
\end{equation}
This displays the source or shell contribution explicitly.  It is
an integration-by-parts identity for smooth fields, not a reason to
omit $Q$ before localizing the action.

The expansion may be differentiated for a specified smooth finite
parameter family, but not by differentiating a bare big-$O$
inequality.  Suppose $A(t,x),h(t,x)$ are smooth for $t$ in a
compact interval, have a common support bound for the connections,
and the derivatives used above are uniformly bounded.  Apply the
Taylor proof directly to $\partial_t^mG(s,y;t)$: this function is
still even in $s$, and its eighth $s$ derivative is uniformly
bounded on the compact parameter set.  The quadrature proof
applies to the differentiated coefficients as well.  It follows
that for each fixed $m$,
\begin{equation}
 \left|\partial_t^m\mathcal E_{A(t),h(t)}(a)\right|
                                      \le C_m a^4.
 \label{eq:classical-parameter-remainder}
\end{equation}
The hypotheses are smooth classical family bounds; they are not
uniform bounds on fields distributed by the lattice Gibbs law.

\subsubsection{What the classical calculation supplies to the Ward problem}

The classification excludes scalar, antisymmetric and vector partners
in the stated dimensions by finite colour and coordinate algebra.
It proves that the dimension-six rotational obstruction has two
directions as a density, one after the specified divergences are
removed, and still one after the field-equation quotient.  The
Wilson expansion identifies the coefficient of that last direction
as $-a^2\beta/24$, while retaining the weighted divergence with
coefficient $a^2\beta/48$.  These conclusions are off-shell
identities or controlled classical asymptotics in their declared
test spaces.

To make the bulk term of the interacting rotational Ward identity
vanish, one still needs the actual typed insertion estimates,
their source formation and normalization, the required treatment
of equation-of-motion insertions under the measure, and the shell
or volume transfer bounds.  A smooth-field expansion is not
uniform on the rough lattice fields in that measure.  The factors
$a^2$ and the absence of lower-dimensional partners do not, by
themselves, prove the quantum estimate.  The finite exact Ward
identity and the interacting bounds remain distinct parts of that
argument.

% END PROOF EXPANSION ward_classical/classical_invariants_wilson.tex


% BEGIN MODULE chapters/euclidean_hilbert_bridge_v3.tex
\subsection{Euclidean evolution from separated moments}
\label{ch:euclidean-hilbert-bridge}

This section proves a conditional reconstruction on compactly supported
separated test tuples. The time contraction follows from a bound which is
uniform in the positions of finitely many supports. No estimate on arbitrary
Schwartz sequences, and no bound uniform in the number of fields, enters
that argument. A further lattice approximation hypothesis gives a spectral
cone with a cubical spatial section. The round cone and continuous rotations
require additional input.

\subsubsection{The precise compact test domain}

Put $\mathcal D=C_c^1(\mathbb R^4;\mathbb C)$ and use the norm
\[
\|f\|_1=\max\{\|f\|_\infty,\|\nabla f\|_\infty\}.
\]
A separated tuple $X=(f_1,\ldots,f_n)$ has compact supports at pairwise
positive distance. The empty tuple is denoted by $1$. Write $\mathscr V$
for the complex vector space of finite formal linear combinations of these
tuples. Multilinear and permutation relations may be imposed at the outset;
equivalently they will be null relations in the forms below. Concatenation
$X\cdot Y$ is used only if the supports of all its entries remain separated.
There is no claim that $\mathscr V$ is an algebra under unrestricted
concatenation.

Our translation convention is
$\tau_vf(x)=f(x-v)$, so its support is translated by $v$.
For the time reflection $\theta(x_1,x_2,x_3,x_4)
=(x_1,x_2,x_3,-x_4)$ put $\theta f=f\circ\theta$.
These maps act entrywise on tuples and linearly on formal combinations.
The bar conjugates both the coefficients and the test functions.
Let $\mathscr V_+$ consist of combinations whose entries have supports
compactly contained in $\mathbb E_+=\{x_4>0\}$.

\phantomsection\label{notdef:moment-functional}
Assume normalized symmetric multilinear moments $S_n$ on this separated
domain. Write $\sigma(X)=S_n(f_1,\ldots,f_n)$, $\sigma(1)=1$,
and extend linearly. They obey the following hypotheses.
They preserve conjugation, are translation invariant, and are invariant under
the coordinate reflections. Their reflection-positive form satisfies

\refstepcounter{equation}
\begin{gather*}
\FitMath{q(F,F):=\sigma(\overline{\theta F}\cdot F)\ge0} \\[3pt]
\FitMath{(F\in\mathscr V_+).}\tag{\theequation}\label{eq:euclidean-hilbert-bridge-HB1}
\end{gather*}

This is positivity for every finite linear combination, including the empty
tuple, rather than just for individual monomials. The concatenations in
\eqref{eq:euclidean-hilbert-bridge-HB1} and its polarized form are separated: reflected supports lie strictly
below time zero, and the others strictly above it.

For every integer $n\ge1$, $R<\infty$ and $d>0$, assume a finite number
$M_n(d,R)$ such that

\refstepcounter{equation}
\begin{gather*}
\FitMath{|S_n(f_1,\ldots,f_n)|
       \le M_n(d,R)\prod_{i=1}^n\|f_i\|_1,}\tag{\theequation}\label{eq:euclidean-hilbert-bridge-HB2}
\end{gather*}

whenever each support is in a ball of radius $R$, centered anywhere, and
their pairwise distances are at least $d$. The centers need not lie in a
common bounded region. For one entry the separation condition is empty.
These are moment bounds, not merely cumulant bounds. Bounds on cumulants
give \eqref{eq:euclidean-hilbert-bridge-HB2} after the finite moment--cumulant sum, if the needed one-point
bound is included. There is no restriction here on the growth of $M_n$ as
$n$ increases or $d$ decreases.

\begin{proposition}[Compact-tuple Euclidean reconstruction]
Under \eqref{eq:euclidean-hilbert-bridge-HB1}--\eqref{eq:euclidean-hilbert-bridge-HB2} and the stated invariances, the completion of
$\mathscr V_+$ by its reflection form is a separable Hilbert space
$\mathcal H$. The empty tuple has class $\Omega$ of norm one. The formulas

\refstepcounter{equation}
\begin{gather*}
\FitMath{T_t[F]=[\tau_{te_4}F]} \\[3pt]
\FitMath{(t\ge0),} \\[3pt]
\FitMath{U_x[F]=[\tau_{(x,0)}F]} \\[3pt]
\FitMath{(x\in\mathbb R^3)}\tag{\theequation}\label{eq:euclidean-hilbert-bridge-HB3}
\end{gather*}

define, respectively, a strongly continuous semigroup of positive
self-adjoint contractions and a strongly continuous unitary representation.
They commute and fix $\Omega$. There is a unique nonnegative self-adjoint
operator $H$ with $T_t=e^{-tH}$. Each $T_t$, $t>0$, is injective.
The spatial generators $P_i$ and $H$ commute strongly.
\end{proposition}

\subsubsection{The quotient and the contraction proof}

Positivity of the quadratic form gives Hermiticity by polarization and the
Cauchy--Schwarz inequality
$|q(F,G)|^2\le q(F,F)q(G,G)$. Thus
$\mathscr N=\{F:q(F,F)=0\}$ is a linear subspace orthogonal to everything.
The quotient $\mathscr V_+/\mathscr N$, with inner product $q$, is
pre-Hilbert; its completion defines $\mathcal H$.
In particular $\|\Omega\|^2=\sigma(1)=1$ and
$\langle\Omega,[F]\rangle=\sigma(F)$.

Time and spatial translations in \eqref{eq:euclidean-hilbert-bridge-HB3} preserve their positive-time domain.
Translation invariance and $\theta\tau_{te_4}=\tau_{-te_4}\theta$ give

\refstepcounter{equation}
\begin{gather*}
\FitMath{q(F,\tau_{te_4}G)=q(\tau_{te_4}F,G),} \\[3pt]
\FitMath{q(\tau_{(x,0)}F,\tau_{(x,0)}G)=q(F,G).}\tag{\theequation}\label{eq:euclidean-hilbert-bridge-HB4}
\end{gather*}

The second identity proves that spatial translations preserve null vectors
and induce isometries with inverse $U_{-x}$.
For the first identity, boundedness of the induced time maps remains to be
proved; symmetry alone would not suffice.

Fix $F=\sum_\alpha c_\alpha X_\alpha\in\mathscr V_+$.
There are $\delta>0$, $d>0$ and $R<\infty$ such that every entry has
time coordinate at least $\delta$, every individual tuple has internal
separation at least $d$, and each entry has support in a ball of radius $R$.
For every $s\ge0$, every concatenation
$\overline{\theta\tau_{se_4}X_\alpha}\cdot
\tau_{se_4}X_\beta$ has internal separation at least
$d_*=\min\{d,2\delta\}>0$. Indeed its two groups lie below
$-\delta-s$ and above $\delta+s$. Moving the groups apart changes
their centers but not their radii, norms or internal separations.
Applying \eqref{eq:euclidean-hilbert-bridge-HB2} to the finitely many pairs gives

\refstepcounter{equation}
\begin{gather*}
\FitMath{\sup_{s\ge0}\|\tau_{se_4}F\|_q\le C_F<\infty.}\tag{\theequation}\label{eq:euclidean-hilbert-bridge-HB5}
\end{gather*}

The same assertion covers empty entries and constants by normalization.
Its essential hypothesis is that the bound is independent of support
positions. A bound with a growing enclosing-volume factor would not prove
\eqref{eq:euclidean-hilbert-bridge-HB5}.

Write $a(s)=\|\tau_{se_4}F\|_q$. Symmetry and Cauchy--Schwarz give

\begin{gather*}
\FitMath{a(t)^2=q(F,\tau_{2te_4}F)\le a(0)a(2t).}
\end{gather*}

Iteration yields

\refstepcounter{equation}
\begin{gather*}
\FitMath{a(t)\le a(0)^{1-2^{-k}}a(2^kt)^{2^{-k}}
       \le a(0)^{1-2^{-k}}C_F^{2^{-k}}.}\tag{\theequation}\label{eq:euclidean-hilbert-bridge-HB6}
\end{gather*}

If $a(0)=0$, the first inequality already proves $a(t)=0$.
Otherwise let $k\to\infty$ to obtain $a(t)\le a(0)$.
Consequently time translations preserve $\mathscr N$ and extend to
contractions on $\mathcal H$. Their symmetry in \eqref{eq:euclidean-hilbert-bridge-HB4} makes each bounded
extension self-adjoint. Moreover

\begin{gather*}
\FitMath{q(F,\tau_{te_4}F)=q(\tau_{(t/2)e_4}F,
                         \tau_{(t/2)e_4}F)\ge0,}
\end{gather*}

so every $T_t$ is positive. The semigroup law is inherited from the exact
translation law on the dense domain.

\subsubsection{Continuity, generators and elementary consequences}

For $f\in C_c^1$, uniform continuity of $\nabla f$ gives

\refstepcounter{equation}
\begin{gather*}
\FitMath{\|\tau_h f-f\|_1
 \le\max\left\{|h|\|\nabla f\|_\infty,
       \sup_{|x-y|\le|h|}|\nabla f(x)-\nabla f(y)|\right\}
       \longrightarrow0.}\tag{\theequation}\label{eq:euclidean-hilbert-bridge-HB7}
\end{gather*}

This statement need not hold in the Lipschitz norm for a merely Lipschitz
function. To apply it to a moment, enclose the old and slightly moved
supports of each entry in one common neighborhood. For sufficiently small
shifts these neighborhoods retain half their original positive separation.
The multilinear telescoping identity and \eqref{eq:euclidean-hilbert-bridge-HB2} then bound the change of a
moment by a constant times the sum of the $C^1$ changes of its entries.
Differences of functions are evaluated only on those fixed neighborhoods.

Expanding $\|[\tau_hF]-[F]\|^2$ into its finitely many moment terms and
using this observation proves convergence to zero for $h=te_4$, $t\downarrow0$,
and for spatial $h\to0$. The contraction and unitary bounds extend this
continuity from the dense quotient domain to all of $\mathcal H$.
The semigroup and group laws then give strong continuity at every parameter.

Here is a construction of the generator which also rules out a hidden
zero-transfer subspace. If $T_1v=0$, positivity and
$T_1=T_{1/2}^2$ give $T_{1/2}v=0$. Repeated halving and strong continuity
at zero imply $v=0$. Thus the spectral measure of $T_1$ has no atom at
zero. Define the densely defined self-adjoint operator
$H=-\log T_1$ by functional calculus; its spectrum is nonnegative.
Uniqueness of positive roots gives $T_{m/n}=T_1^{m/n}$ for nonnegative
rationals, and strong continuity gives $T_t=e^{-tH}$ for all $t\ge0$.
Functional calculus also proves uniqueness of $H$ and injectivity of every
$T_t$, $t>0$. Since translations fix the empty tuple,
$H\Omega=0$ and $U_x\Omega=\Omega$.

The commuting spatial unitaries have self-adjoint generators
$U_x=e^{i\sum_i x_iP_i}$. They commute with $T_1$ and hence with its
spectral projections and with $H$. The joint spectral theorem, equivalently
the spectral theorem for the strongly continuous abelian unitary group
$(s,x)\mapsto e^{isH}U_x$, supplies one projection-valued measure
$\mathsf E$ on $[0,\infty)\times\mathbb R^3$. Here and below strong
commutation means commutation of these spectral projections.
For $F,G\in\mathscr V_+$, $t\ge0$ and $x\in\mathbb R^3$ the dictionary is

\refstepcounter{equation}
\begin{gather*}
\FitMath{\begin{aligned}
 \sigma(\overline{\theta F}\,\tau_{(x,t)}G)
 &=\langle[F],e^{-tH}U_x[G]\rangle\\
 &=\int e^{-tE+ix\cdot p}
                  \langle[F],\mathsf E(dE\,dp)[G]\rangle.
 \end{aligned}}\tag{\theequation}\label{eq:euclidean-hilbert-bridge-HB7a}
\end{gather*}


For completeness, the quotient is separable. Cover compact positive-time
supports by finite unions of rational boxes with closures still in
$\mathbb E_+$; choose these unions separately for each entry so that they
remain separated. There are countably many such finite systems. On each
fixed bounded union, the $C^1$ space is separable: restriction of a function
and its four derivatives embeds it isometrically in a finite product of
separable continuous-function spaces. Choose countable dense subsets inside
the supported $C^1$ spaces, and rational complex coefficients for formal
combinations. Telescoping with \eqref{eq:euclidean-hilbert-bridge-HB2}, on the fixed separated neighborhoods,
shows that the resulting countable set is total in $\mathcal H$.

The formal description does not create extra vectors from zero tensors.
For a finite homogeneous relation $\sum_\alpha c_\alpha
f_{\alpha,1}\otimes\cdots\otimes f_{\alpha,n}=0$ whose tuples are
separated, choose $d>0$ below all their separations. Localize every entry
with one finite smooth partition covering their supports, whose sets have
diameter less than $d/8$. In a product of partition sets, either two sets
have distance less than $d/2$, in which case every original tensor has a
zero factor there, or all the slot sets remain separated. In the latter
case the moment is multilinear on the product of their supported function
spaces and hence is a linear functional on their algebraic tensor product.
The map from this tensor product into functions on the Cartesian product
is injective: on the finitely many function spans involved, independent
functions can be distinguished by finitely many point evaluations, and
the corresponding evaluation matrices invert each slot. Thus the localized
zero tensor has zero moment. Summing proves the assertion. Applying it
also after concatenation with a reflected test shows that a zero tensor
is a reflection-null vector. This is an algebraic compact-support fact,
distinct from the Schwartz density input below.

Any realization with a dense image of the same test domain and the same
scalar products is therefore canonically unitarily equivalent, by sending
each test class to its other realization; the generators intertwine by
density. The construction can use any coordinate as time if its RP
hypothesis is supplied. For $c\ge0$ and $F$ supported strictly above time
$c$, the exact identity is $[F]=T_c[\tau_{-ce_4}F]$.

Any further linear symmetry of the moments which preserves positive time
and commutes with $\theta$ acts isometrically on this quotient. In particular,
if the moments also have the spatial signed-permutation symmetry $W_3$,
then $V_w[F]=[wF]$ is unitary,
$V_wT_t=T_tV_w$ and $V_wU_xV_w^{-1}=U_{wx}$.
Uniqueness of the joint spectral measure gives the corresponding covariance
under $(E,p)\mapsto(E,wp)$ and hence $W_3$ symmetry of the joint spectrum.
This optional finite symmetry is not continuous rotation invariance.

Spectral dominated convergence gives $T_t\to P_0=1_{\{0\}}(H)$ strongly
as $t\to\infty$. Hence the large-time reflected correlations equal the
matrix elements of $P_0$. In particular, $\ker H=\mathbb C\Omega$ exactly
when their limits, or their Ces\`aro limits, equal
$\overline{\sigma(F)}\sigma(G)$ for every positive-time $F,G$.
If $\sigma(f)=0$ and $S_2(\theta\bar f,f)>0$ for one positive-time test,
then $[(f)]$ is a nonzero vector orthogonal to $\Omega$, so
$\dim\mathcal H\ge2$; this requires an actual positive two-point input.

A shift of the fields does not change this representation. More precisely,
let $m$ be a complex linear, translation- and reflection-invariant functional
on $\mathcal D$, satisfying $m(\bar f)=\overline{m(f)}$ and bounded on
each fixed compact support. Define

\begin{gather*}
\FitMath{C_m(f_1,\ldots,f_n)
  =\sum_{A\subset\{1,\ldots,n\}}
       \prod_{i\notin A}[-m(f_i)]\,(f_i)_{i\in A},} \\[3pt]
\FitMath{\sigma^{\rm c}(F)=\sigma(C_mF).}
\end{gather*}

Then $C_{-m}$ is the inverse of $C_m$, and on every admissible
concatenation $C_m(FG)=C_m(F)C_m(G)$. It commutes with the bar, reflections
and translations. Therefore
$[F]_{\sigma^{\rm c}}\mapsto[C_mF]_\sigma$ is a unitary intertwining
the two evolutions. Its finite expansion also transfers in \eqref{eq:euclidean-hilbert-bridge-HB2}, since local
boundedness and translation invariance bound $m$ uniformly on balls of each
fixed radius. This algebraic statement does not estimate the size of
the one-point difference between two cutoff states.

In the particular unweighted Wilson-density normalization, the one-point
calibration can be written explicitly. Let two states $\rho_j,\nu_j$ on
spacing $a_j$ be translation and hypercubic invariant, and put
$\mathcal O_j(f)=\sum_p f(x(p))a_p$. Their plaquette mean difference
$d_j=\rho_j(a_p)-\nu_j(a_p)$ is independent of $p$. Define

\begin{gather*}
\FitMath{R_j(f)=\frac{a_j^4}{6}\sum_p f(x(p)),} \\[3pt]
\FitMath{c_j=6a_j^{-4}d_j.}
\end{gather*}

Then
$\rho_j(\mathcal O_j(f))-\nu_j(\mathcal O_j(f))=c_jR_j(f)$
exactly, and the six shifted Riemann sums give $R_j(f)\to\int f$.
If these one-point differences converge to $m(f)$ on all compact $C^1$
tests, choose one $f_0$ with $\int f_0=1$. The ratio at $f_0$ shows
$c_j\to m(f_0)=c$, and consequently $m(f)=c\int f$ for every test.
A uniform one-point bound at $f_0$ bounds $|c|$; this bound is an input,
not a consequence of translation invariance alone. When the two states
coincide, $c_j=c=0$.

\subsubsection{Coarse lattice inequalities and the continuum cube cone}

The spectral cone needs additional information. Assume infinite-lattice
probability measures $\nu_j$ at spacings $a_j\downarrow0$, invariant
under lattice translations and reflection positive on the closed side of
all coordinate and diagonal lattice mirrors. Their moment forms will be
denoted by $S_n^{(j)}$. For periodic Wilson DLR states these hypotheses are proved in
\cref{thm:amrp-all-periodic}. For other lattice measures they
remain explicit assumptions. The normal transfer $T_j$
and tangential translation $V_{j,i}$ in their reflection Hilbert spaces
then have the carrier established in the plateau argument:

\refstepcounter{equation}
\begin{gather*}
\FitMath{\sinh(E/2)\ge|\sin(p/2)|,} \\[3pt]
\FitMath{T_j=e^{-E},} \\[3pt]
\FitMath{V_{j,i}=e^{ip}.}\tag{\theequation}\label{eq:euclidean-hilbert-bridge-HB8}
\end{gather*}

Here $E$ is dimensionless and may be infinity on $\ker T_j$, while $p$
is a circle variable. Only this lattice carrier, rather than a round
relativistic cone, is assumed from that argument.

One lattice distinction is useful. For spacing $a$, the site mirror at
zero and the link mirror at $-a/2$ have the same positive-side lattice
algebra. Their forms obey
$q_{\rm link}(A,B)=q_{\rm site}(A,T B)$, by translation invariance.
Consequently $[A]_{\rm link}\mapsto T^{1/2}[A]_{\rm site}$ is an
isometry onto the closure of $\operatorname{Ran}T^{1/2}$, namely
$(\ker T)^{\perp}$. It intertwines the two normal transfers, since this
holds on the common local domain. Thus the link transfer is injective,
even if the site transfer has a kernel. This proves the relation between
the two pictures without discarding site-kernel contributions in \eqref{eq:euclidean-hilbert-bridge-HB10}.

For every integer $k\ge1$, convexity of $\sinh$ and the elementary
inequality $|\sin(kv)|\le k|\sin v|$ give

\refstepcounter{equation}
\begin{gather*}
\FitMath{\sinh(kE/2)\ge k\sinh(E/2)
       \ge k|\sin(p/2)|\ge|\sin(kp/2)|.}\tag{\theequation}\label{eq:euclidean-hilbert-bridge-HB9}
\end{gather*}

The sine inequality follows by induction from the addition formula and
$|\cos v|\le1$. Thus functional calculus gives the bounded operator
inequality

\refstepcounter{equation}
\begin{gather*}
\FitMath{(I-T_j^k)^2\ge T_j^k(2I-V_{j,i}^k-V_{j,i}^{-k}).}\tag{\theequation}\label{eq:euclidean-hilbert-bridge-HB10}
\end{gather*}

At $T_j=0$ this inequality holds directly, so no division by a transfer
eigenvalue is required. Its quadratic form is a five-point inequality.
If $s=ka_j$ and
$\Phi_{j,F}(v)=\nu_j(\overline{\theta F_j}\,\tau_vF_j)$,
it reads

\refstepcounter{equation}
\begin{gather*}
\FitMath{\Phi_{j,F}(0)-4\Phi_{j,F}(se_4)+\Phi_{j,F}(2se_4)
       +\Phi_{j,F}(se_4+se_i)+\Phi_{j,F}(se_4-se_i)\ge0.}\tag{\theequation}\label{eq:euclidean-hilbert-bridge-HB11}
\end{gather*}


To pass to the limit, use the following approximation between cutoffs. Assume linear local lattice
bounded Borel representatives $\iota_j(f)$ for compact $C^1$ tests. Their supports lie
within distance $\epsilon_j\to0$ of $\operatorname{supp}f$. They commute
with coordinate reflection and with translations by $a_j\mathbb Z^4$.
Let $F_j$ be the product representative of each tuple, extended linearly,
and let $S_n^{(j)}$ be its moments. Assume convergence on every fixed
separated compact tuple to $S_n$, and assume \eqref{eq:euclidean-hilbert-bridge-HB2} for these lattice moments
with constants independent of $j$ once $j$ is sufficiently large for the
fixed $n,d,R$. For $\operatorname{tr}(F^2)$ representatives their supremum norms may
depend on $j$; the asserted uniform control is on these moments.

If $v_j\to v$, the supports of $\tau_{v_j}f$ and $\tau_vf$ lie in fixed
slightly enlarged sets. Whenever the limiting tuple is separated, these
sets retain a positive separation. Telescoping and the uniform moment
bound give

\refstepcounter{equation}
\begin{gather*}
\FitMath{S_n^{(j)}(f_{1,j},\ldots,f_{n,j})
       -S_n^{(j)}(f_1,\ldots,f_n)\longrightarrow0}\tag{\theequation}\label{eq:euclidean-hilbert-bridge-HB12}
\end{gather*}

for such translated tests, by \eqref{eq:euclidean-hilbert-bridge-HB7}. Together with fixed-tuple convergence,
this proves convergence at the moving arguments.

The same argument also explains inheritance of the Hilbert hypotheses from
these approximants, when it is needed. A continuum translation is approached
by lattice vectors $v_j$; exact lattice covariance and \eqref{eq:euclidean-hilbert-bridge-HB12} pass its
moment identity to the limit. Coordinate reflections through the origin
need no rounding. For any fixed $F\in\mathscr V_+$ the representatives
$F_j$ lie on the positive side for large $j$, so their finite quadratic
RP inequalities pass directly to \eqref{eq:euclidean-hilbert-bridge-HB1}. Normalization, conjugation,
multilinearity and the uniform moment bounds pass by their finite defining
relations. The lattice mirror hypotheses themselves are supplied by
\cref{thm:amrp-all-periodic} in the periodic Wilson application.

There is also a finite-torus scalar-product statement, with a narrower
conclusion. If fixed separated moments converge from growing periodic tori
and satisfy uniform $C^1$ bounds on the fixed neighborhoods under
consideration, RP at the link mirror $x_4=a_j/2$ gives

\begin{gather*}
\FitMath{q(F,G)=\lim_j\nu_j(\overline{\theta_{a_j/2}F_j}\,G_j).}
\end{gather*}

Indeed $\theta_{a_j/2}f=\tau_{a_je_4}\theta f$, whose difference from
$\theta f$ tends to zero in $C^1$, and the same telescoping estimate
applies. This compares scalar products as numbers; it asserts no maps
between the cutoff Hilbert spaces. Finite-torus positivity of this kind
does not by itself supply the global lattice semigroup or \eqref{eq:euclidean-hilbert-bridge-HB8}.

Fix $t>0$ and choose integers $k_j\ge1$ with $k_ja_j\to t$.
For positive-time $F$, its lattice representative is supported on the
positive side of the fixed site mirror for large $j$.
All five moment expressions in \eqref{eq:euclidean-hilbert-bridge-HB11} are separated uniformly: the negative
reflected group and every translated positive group retain their original
positive-time margin. Internal separations are unchanged. Therefore \eqref{eq:euclidean-hilbert-bridge-HB12}
passes \eqref{eq:euclidean-hilbert-bridge-HB11} to its continuum version. In $\mathcal H$ this is

\refstepcounter{equation}
\begin{gather*}
\FitMath{B_{t,i}:=(I-T_t)^2-T_t(2I-U_{te_i}-U_{-te_i})\ge0.}\tag{\theequation}\label{eq:euclidean-hilbert-bridge-HB13}
\end{gather*}

Initially its form is nonnegative on the dense compact-tuple domain; since
$B_{t,i}$ is bounded and self-adjoint, the inequality holds on all of
$\mathcal H$. No differentiation of a cutoff correlation at $t=0$ has
been used. The fixed-time coarse inequality is passed first.

In the joint functional calculus its multiplier is

\refstepcounter{equation}
\begin{gather*}
\FitMath{b_{t,i}(E,p)=1-4e^{-tE}+e^{-2tE}
                     +2e^{-tE}\cos(tp_i)
       =2e^{-tE}[\cosh(tE)+\cos(tp_i)-2].}\tag{\theequation}\label{eq:euclidean-hilbert-bridge-HB14}
\end{gather*}

Take a countable sequence $t_m\downarrow0$, for example $t_m=1/m$.
Positivity of \eqref{eq:euclidean-hilbert-bridge-HB13} means that the joint projection-valued measure gives
zero mass to $\{b_{t_m,i}<0\}$. Remove the countable union of those
sets for the three spatial indices. At every remaining spectral point,

\begin{gather*}
\FitMath{0\le\lim_{m\to\infty}
       \frac{\cosh(t_mE)+\cos(t_mp_i)-2}{t_m^2}
       =\tfrac12(E^2-p_i^2).}
\end{gather*}

Since $E\ge0$, this proves

\refstepcounter{equation}
\begin{gather*}
\FitMath{\mathsf E\{(E,p):E<\max_i|p_i|\}=0.}\tag{\theequation}\label{eq:euclidean-hilbert-bridge-HB15}
\end{gather*}

In particular $\operatorname{Dom}H\subset\operatorname{Dom}P_i$ and
$\|P_i\psi\|\le\|H\psi\|$ on that domain. The corresponding
form-domain statement follows by integrating $|p_i|\le E$ rather than
their squares. Zero-energy vectors are spatially invariant.
Consequently vectors fixed by all Euclidean time and spatial translations
are exactly $\ker H$. This statement concerns invariant Hilbert vectors;
identifying them with an ergodic decomposition of correlation states needs
the separate phase argument.
The cone in \eqref{eq:euclidean-hilbert-bridge-HB15} has a cubical spatial section.
Nothing here proves $E\ge|p|_2$ or supplies an $\operatorname{O}(4)$ action.
Finite periodic torus laws cannot simply be substituted for the infinite
lattice hypotheses at arbitrarily large separations; any required
infinite-volume reflection-positive state must first be obtained.

\subsubsection{What the full Schwartz sequence domain additionally needs}

The compact-tuple construction above does not establish the claimed
Schwartz-domain growth property. Here is a precise sufficient additional
input for that extension. For $n\ge2$ put
$\Delta_n=\bigcup_{i<j}\{x_i=x_j\}$ and let
$\mathcal S_{\Delta,n}$ be the Schwartz functions on $\mathbb R^{4n}$
whose derivatives of every order vanish on $\Delta_n$.
Use $\mathcal S_{\Delta,0}=\mathbb C$ and the usual Schwartz space for
$n=1$. Define the seminorms

\begin{gather*}
\FitMath{p_k(\varphi)=\max_{|\alpha|\le k}\sup_x
       (1+|x|^2)^{k/2}|\partial^\alpha\varphi(x)|.}
\end{gather*}

The positive-time subspace $\mathcal S_{\Delta,+,n}$ consists of functions
all of whose derivatives vanish outside
$\{x_{1,4}>0,\ldots,x_{n,4}>0\}\setminus\Delta_n$.
This includes flatness at the time boundary, as well as at the diagonals.

The separate continuity/growth input, called H1 here, supplies continuous
linear moment functionals $\Sigma_n$ on these full diagonal-flat spaces,
agreeing with the compact smooth separated tensors, and constants
$c,C<\infty$, $k\in\mathbb N$, $\varkappa<\infty$ such that

\refstepcounter{equation}
\begin{gather*}
\FitMath{|\Sigma_n(\varphi)|\le c C^n(n!)^{\varkappa}p_{kn}(\varphi).}\tag{\theequation}\label{eq:euclidean-hilbert-bridge-HB16}
\end{gather*}

This is a condition on moments, not just connected functions. It is not a
consequence of having a finite $M_n(d,R)$ for every fixed triple of
parameters. The Density lemma of \cref{ch:precise-reconstruction} proves that finite linear combinations of
smooth compact separated tensors are dense in $\mathcal S_{\Delta,n}$,
and those with strictly positive-time supports are dense in
$\mathcal S_{\Delta,+,n}$, in their Schwartz subspace topologies.
These assertions include the necessary cutoffs at infinity, at the time
boundary, and near the diagonals. They cannot be replaced by ordinary
tensor density without those constraints.

Using H1 and that density lemma, the finite Schwartz sequence space
$\bigoplus_{n\ge0}\mathcal S_{\Delta,+,n}$ has form
$\sum_{m,n}\Sigma_{m+n}(\overline{\theta F_m}\otimes G_n)$.
The tensor here lies in the diagonal-flat space: an internal coincidence
is killed by the corresponding factor, and a coincidence between a negative
and a positive time argument forces the common time to zero, where the
factors are flat. Tensor multiplication is continuous in the Schwartz
topologies. Thus H1 and tensor density transfer the compact reflection
positivity and show that its quotient completion is canonically the same
$\mathcal H$. Compact $C^1$ tuples themselves are approximated by compact
smooth ones inside the same separated positive-time neighborhoods, using
\eqref{eq:euclidean-hilbert-bridge-HB2}. The already proved contraction and unitary groups then extend by
this identification. The use of H1 here is its continuity on the enlarged
domain; \eqref{eq:euclidean-hilbert-bridge-HB16} also records the additional quantitative growth claim needed
when that OS hypothesis is asserted. The compact construction does not supply the factorial constants
in H1; the constrained tensor-density proof is given in the cited section.


% END MODULE chapters/euclidean_hilbert_bridge_v3.tex

\section{Reconstruction and canonical components}\label{sec:reconstruction}
This section passes from the Euclidean functionals to scalar quantum fields. We first construct the spectral and common-domain objects and compare their moments with the reconstructed theory. The compact decomposition then identifies the canonical components and their vacuum spaces. The final reconstruction statement records the precise growth, positivity, covariance and clustering hypotheses. The intervening observable-algebra results specify the information carried by the scalar field and the additional requirements for extending it to gauge observables.

% BEGIN MODULE chapters/reconstruction_spectral_bridge.tex
\subsection{Complex translations and reconstruction domains}
\label{ch:reconstruction-spectral-interface}

The diagonal-flat growth theorem controls the order of the moment
functionals. The Euclidean Hilbert construction gives the Hilbert space
and its evolution.
We match their domains here, prove the round spectral cone and a joint
complex translation theorem for that evolution, and specify the field
domain obtained once the Minkowski distributions have been constructed.
The translation theorem here has fixed-vector scope.  The actual
field-product continuation is constructed in
\cref{ch:h3-energy-band-fields}; the explicit reconstruction
norm estimates and the locality theorem are matched in
\cref{ch:quantitative-reconstruction-matching}.

\subsubsection*{The exact test-space match}

Let \(\mathcal S_{\Delta,n}\) be the Schwartz functions on
\(\mathbb R^{4n}\) whose every derivative vanishes whenever two
four-dimensional arguments coincide.  Write
\begin{equation}
 p_j(f)=\max_{|\alpha|\le j}\sup_x
            (1+|x|^2)^{j/2}|\partial^\alpha f(x)|.
 \label{eq:reconstruction-spectral-bridge-RHI1}
\end{equation}
The quantitative flat-extension theorem gives, under its actual source
hypotheses, normalized symmetric moments with
\begin{equation}
 |S_n(f)|\le A B^n(n!)^{k+32}p_{20n}(f).
 \label{eq:reconstruction-spectral-bridge-RHI2}
\end{equation}
This is the diagonal-flat version of the growth requirement.  If
\(B_+=\max(1,B)\), the elementary inequality \(n!\ge2^{n-1}\)
gives
\begin{equation}
 B^n\le B_+(n!)^{\log_2 B_+},\qquad
 \sigma_n=A B_+(n!)^{k+32+\lceil\log_2B_+\rceil}.
 \label{eq:reconstruction-spectral-bridge-RHI3}
\end{equation}
Thus \(|S_n(f)|\le\sigma_n p_{20n}(f)\), with one linear order
coefficient and factorial \(\sigma_n\).  No unquantified change of
seminorm order has occurred.

The applicable external result is the corrected scalar theorem
\(E'\!\to R'\) with \(E0'\), in
\cite[Section IV.1, pp.~287--289]{OS75}.  The exact diagonal-flat seminorm comparison is given in
\cref{thm:quantitative-os-seminorm-ledger}.  In particular the
proved weighted order \(20n\) and factorial exponent \(k+32\)
imply the safe total-monomial order \(40n\), with an explicit
additional factorial exponent twenty and exponential constant.
Thus the actual growth is sufficient without identifying two
seminorm conventions by an unspecified dimension-dependent constant.  We retain its Euclidean covariance, positivity,
symmetry, reality and clustering requirements.  The argument here does
not replace that result by an assertion about independently continuous
correlations or assign an ordinary Euclidean contact prescription.

The positive-time space already used in
\cref{ch:euclidean-hilbert-bridge} is denoted by
\(\mathcal S_{\Delta,+,n}\).  Its functions are flat outside the
region of distinct arguments with all times positive.  Define also
\begin{equation}
 \begin{gathered}
 O_n=\{0<t_1<t_2<\cdots<t_n\},\\
 \mathcal S_{{\rm ord},n}
  =\{f\in\mathcal S(\mathbb R^{4n}):
               \partial^\alpha f=0\text{ on }O_n^c
                \text{ for every }\alpha\}.
 \end{gathered}
 \label{eq:reconstruction-spectral-bridge-RHI4}
\end{equation}
The ordered spaces, used in the historical positivity convention
\cite[Sections 2--3]{OS73}, are subspaces of
\(\mathcal S_{\Delta,+,n}\).  Set both degree-zero spaces to
\(\mathbb C\), and use finite direct-sum sequences throughout.

Let \(\mathcal K\) be the completed full positive-time quotient,
with its proved form
\begin{equation}
 q(F,G)=\sum_{m,n}S_{m+n}
              (\Theta\overline{F_m}\otimes G_n).
 \label{eq:reconstruction-spectral-bridge-RHI5}
\end{equation}
The historical reflected form also reverses the first group of
arguments.  Scalar permutation symmetry makes the two forms equal
on the ordered subspace.  Restriction of \eqref{eq:reconstruction-spectral-bridge-RHI5} consequently gives an
isometry of its ordered quotient into \(\mathcal K\).  Denote its
closed range by \(\mathcal K_{\rm ord}\).

\begin{proposition}[The ordered sector and its inherited evolution]
\label{thm:ordered-sector-restriction}
The subspace \(\mathcal K_{\rm ord}\) contains the vacuum and
reduces the positive-time contractions \(T_t=e^{-tH}\) and all
spatial translations \(U_a=e^{ia\cdot P}\).  Their restrictions are
strongly continuous, have the restricted joint spectral measure,
and have the same nonnegative Hamiltonian in functional calculus.
Separability, spatial clustering and a one-dimensional kernel of
\(H\), when established on \(\mathcal K\), pass to this sector.
\end{proposition}

\emph{Proof.}
Simultaneous positive-time translation preserves strict time ordering,
as do all simultaneous spatial translations.  The corresponding
operators therefore preserve the ordered core and its closed range.
For a bounded self-adjoint \(T_t\), invariance of a closed subspace
\(M\) also makes \(M^\perp\) invariant: if \(x\perp M\) and
\(y\in M\), then \(\langle T_tx,y\rangle=\langle x,T_ty\rangle=0\).
Thus its projection commutes with every \(T_t\), and hence with the
spectral projections of \(H\).  Both signs of each spatial translation
preserve the range, so its projection commutes with the spatial
unitaries and their spectral projections as well.  This proves the
joint spectral restriction and the domain identity
\begin{equation}
 \operatorname{Dom}H_{\rm ord}
       =\operatorname{Dom}H\cap\mathcal K_{\rm ord},\qquad
 H_{\rm ord}=H|_{\operatorname{Dom}H_{\rm ord}}.
 \label{eq:reconstruction-spectral-bridge-RHI6}
\end{equation}
Strong continuity restricts.  The degree-zero vector belongs to the
ordered sector and still has norm one.  A closed subspace of a separable
Hilbert space is separable.  A clustering matrix-element limit restricts
to any pair of its vectors.  Finally, a vector annihilated by the
restricted Hamiltonian lies in \(\ker H\); if the latter is
\(\mathbb C\Omega\), the restricted kernel is exactly that line.
\hfill\(\square\)

Equality \(\mathcal K_{\rm ord}=\mathcal K\) is not needed in
this theorem or for these transfers to the reconstructed sector.
It is not a Schwartz density statement: evaluation at a configuration
with equal positive times and distinct spatial points vanishes on
every ordered test, while it need not vanish on a test in
\(\mathcal S_{\Delta,+,n}\).  A claim of equality of the completed
Hilbert spaces would require an additional argument in the Hilbert
seminorm.  The imported reconstruction supplies fields on its ordered
realization.  The theorem above does not extend their field operators
or their Lorentz boosts to a larger space.

\subsubsection*{Implementing spatial rotations on the actual quotient}

Suppose now that the established Euclidean functional has spatial
\(\operatorname{SO}(3)\) covariance.  This is an explicit input from the
full covariance route.  Neither the finite hypercubic group nor
reflection positivity alone supplies it.  For \(R\in\operatorname{SO}(3)\)
let
\begin{equation}
 (r_R f)(\mathbf x_1,t_1,\ldots,\mathbf x_n,t_n)
   =f(R^{-1}\mathbf x_1,t_1,\ldots,R^{-1}\mathbf x_n,t_n).
 \label{eq:reconstruction-spectral-bridge-RHI7}
\end{equation}
The pullback preserves all the flat and positive-time spaces, commutes
with reflection, and preserves the form in \eqref{eq:reconstruction-spectral-bridge-RHI5}.  It consequently
preserves its null space and gives an isometry
\(V_R[F]=[r_RF]\) with inverse \(V_{R^{-1}}\).  This defines
unitaries on the actual completed quotient.

We check their continuity.  For each fixed \(n,j\), the chain rule
and the orthogonality of \(R\) bound \(p_j(r_Rf)\) locally uniformly
by a constant times \(p_j(f)\).  Near identity write \(R=e^A\),
where \(A\) is skew and \(\|A\|\le C\|R-I\|\).
Differentiation along \(e^{sA}\) inserts
\(\sum_i(A\mathbf x_i)\cdot\nabla_{\mathbf x_i}\), up to the
pullback sign.  Its coordinate factors consume one weight and its
derivative consumes one derivative.  Thus
\begin{equation}
 p_j(r_Rf-f)\le C_{n,j}\|R-I\|p_{j+1}(f)
 \label{eq:reconstruction-spectral-bridge-RHI8}
\end{equation}
for \(R\) in a fixed neighbourhood of identity.  Continuity at every
rotation follows by composition.  Tensor multiplication is continuous
in the Schwartz topology, and each \(S_{m+n}\) is continuous on its
flat domain.  Expanding
\(\|[r_RF]-[F]\|^2=q(r_RF-F,r_RF-F)\) therefore gives convergence
to zero for each finite-sequence class.  These classes are dense;
unitarity extends the convergence to every vector of \(\mathcal K\).
Hence \(R\mapsto V_R\) is a strongly continuous unitary
representation.  The same proof applies to \(\mathcal K_{\rm ord}\),
since rotations preserve its core.

On the dense core, the exact coordinate identities are
\begin{equation}
 V_RT_t=T_tV_R,\qquad
 V_RU_aV_R^{-1}=U_{Ra}.
 \label{eq:reconstruction-spectral-bridge-RHI9}
\end{equation}
They extend to the Hilbert space by boundedness.  Uniqueness of the
self-adjoint generators gives the unbounded-operator statements
\begin{equation}
 \begin{gathered}
 V_R\operatorname{Dom}H=\operatorname{Dom}H,
       \qquad V_RHV_R^{-1}=H,\\
 V_R\operatorname{Dom}(a\cdot P)
       =\operatorname{Dom}((Ra)\cdot P),\qquad
 V_R(a\cdot P)V_R^{-1}=(Ra)\cdot P.
 \end{gathered}
 \label{eq:reconstruction-spectral-bridge-RHI10}
\end{equation}
The last form states the convention without an ambiguity about whether
a matrix acts on momentum components or on the translating vector.
It also identifies the transformed spectral projections.

\begin{proposition}[The round cone on the Euclidean Hilbert space]
\label{thm:round-cone-from-spatial-rotations}
Suppose the joint spectral measure \(\mathsf E\) of \((H,P)\)
has the established cube carrier
\(E\ge\max_i|p_i|\).  Under the spatial covariance just implemented,
\begin{equation}
 \mathsf E\{(E,p):E<|p|_2\}=0.
 \label{eq:reconstruction-spectral-bridge-RHI11}
\end{equation}
This holds on the full \(\mathcal K\) and on its reducing ordered
sector.  In particular
\(\operatorname{Dom}H\subset\operatorname{Dom}|P|\), and
\(\||P|\psi\|\le\|H\psi\|\) on that domain.
\end{proposition}

\emph{Proof.}
Rotating the first coordinate carrier with \eqref{eq:reconstruction-spectral-bridge-RHI10} gives
\(\mathsf E\{E<|n\cdot p|\}=0\) for every unit vector
\(n=Re_1\).  Choose a countable dense set \(n_j\) on the unit
sphere, and remove the countable union of these spectral null sets.
At every remaining point, \(|n_j\cdot p|\le E\) for every \(j\).
Continuity in \(n\) and
\(\sup_{|n|=1}|n\cdot p|=|p|_2\) prove \eqref{eq:reconstruction-spectral-bridge-RHI11}.  Integrating
\(|p|_2^2\le E^2\) against a vector's spectral measure gives the
stated domain inclusion and norm bound.  Integrating
\(|p|_2\le E\) also gives the corresponding form-domain statement.
Restriction to a reducing subspace restricts the same measure.
\hfill\(\square\)

This argument uses spatial rotations and the already constructed
joint spectrum.  It does not construct a Lorentz boost.

\subsubsection*{A joint tube of bounded complex translations}

Use the joint measure just proved, and define the open parameter tube
\begin{equation}
 \mathcal T=\{(z,w)\in\mathbb C\times\mathbb C^3:
                    \delta(z,w):=\operatorname{Re}z
                                      -|\operatorname{Im}w|_2>0\}.
 \label{eq:reconstruction-spectral-bridge-RHI12}
\end{equation}
Joint functional calculus defines
\begin{equation}
 \mathcal V(z,w)=\int e^{-zE+iw\cdot p}\,
                           \mathsf E(dE\,dp).
 \label{eq:reconstruction-spectral-bridge-RHI13}
\end{equation}
On the carrier \eqref{eq:reconstruction-spectral-bridge-RHI11} its multiplier has modulus at most
\(e^{-\delta(z,w)E}\).  Thus it is a bounded operator of norm at
most one.  The notation \(e^{-zH+iw\cdot P}\) will mean precisely
this bounded operator.

\begin{proposition}[Complex translations and their boundary]
\label{thm:hilbert-complex-translation-tube}
The map \(\mathcal V\) is jointly holomorphic in operator norm on
\(\mathcal T\).  For \(a\ge0\), a three-dimensional multiindex
\(\alpha\), and \(m=a+|\alpha|\),
\begin{equation}
 \|\partial_z^a\partial_w^\alpha\mathcal V(z,w)\|
       \le m!\,\delta(z,w)^{-m}.
 \label{eq:reconstruction-spectral-bridge-RHI14}
\end{equation}
It obeys the addition law and its unitary real boundary is
\begin{equation}
 \lim_{\substack{s,y\to0\\s>|y|_2}}
       \mathcal V(s+it,x+iy)
          =e^{-itH}U_x
          \quad\hbox{strongly}.
 \label{eq:reconstruction-spectral-bridge-RHI15}
\end{equation}
For every pair of Hilbert vectors its matrix coefficients have
tempered boundary values in \((t,x)\), with Fourier carrier
\(E\ge|p|_2\) in the convention \(e^{-itE+ix\cdot p}\).
Every assertion also holds on \(\mathcal K_{\rm ord}\).
\end{proposition}

\emph{Proof.}
The differentiated multiplier is
\((-E)^a(ip)^\alpha e^{-zE+iw\cdot p}\).  Its absolute value is
at most \(E^m e^{-\delta E}\), and
\begin{equation}
 E^m e^{-\delta E}\le m!\delta^{-m},
 \label{eq:reconstruction-spectral-bridge-RHI16}
\end{equation}
as follows from the \(m\)-th term of the exponential series.
This gives \eqref{eq:reconstruction-spectral-bridge-RHI14}.  It also justifies norm differentiation: on a
small parameter ball with \(\delta\ge\delta_0>0\), the scalar
exponential Taylor remainder after its linear term is bounded by
\(C|h|^2E^2e^{-\delta_0E/2}\).  Its spectral supremum tends to
zero quadratically in \(|h|\).  The same argument after any fixed
number of derivatives proves joint complex differentiability and
local norm-convergent Taylor series.

Multipliers multiply, so
\(\mathcal V(z,w)\mathcal V(z',w')=
\mathcal V(z+z',w+w')\); the sum stays in the tube.  For the boundary
in \eqref{eq:reconstruction-spectral-bridge-RHI15}, the multipliers converge pointwise on the spectrum and
remain bounded by one.  Spectral dominated convergence proves strong
convergence for every vector, along every such approach.  In particular
there is no choice of regulator path in this bounded-operator statement.

For fixed \(\xi,\eta\), let
\(\mu_{\xi,\eta}(A)=\langle\xi,\mathsf E(A)\eta\rangle\).
Its total variation is at most \(\|\xi\|\|\eta\|\): for a finite
disjoint partition, apply Cauchy--Schwarz to
\(\sum_j\|\mathsf E(A_j)\xi\|\|\mathsf E(A_j)\eta\|\).
Therefore
\begin{equation}
 \langle\xi,e^{-itH}U_x\eta\rangle
     =\int e^{-itE+ix\cdot p}\,d\mu_{\xi,\eta}(E,p)
 \label{eq:reconstruction-spectral-bridge-RHI17}
\end{equation}
is bounded and continuous.  For a Schwartz test \(g\) on
\(\mathbb R^4\), its distributional pairing has absolute value at
most
\begin{equation}
 \|\xi\|\|\eta\|\int|g(a)|\,da
       \le C\|\xi\|\|\eta\|p_5(g).
 \label{eq:reconstruction-spectral-bridge-RHI18}
\end{equation}
Dominated convergence with this integrable majorant proves the
boundary convergence in \(\mathcal S'\) as well.  The measure in
\eqref{eq:reconstruction-spectral-bridge-RHI17} is supported on the forward cone, proving its stated carrier.
The reducing-subspace result follows from the common restricted
spectral calculus.\hfill\(\square\)

There is also an explicit dense domain for the translation generators:
\begin{equation}
 \mathcal D_{\rm heat}=\bigcup_{s>0}\operatorname{Ran}e^{-sH}.
 \label{eq:reconstruction-spectral-bridge-RHI19}
\end{equation}
The ranges are nested, so this is a linear space; strong continuity
at zero makes it dense.  The cone and \eqref{eq:reconstruction-spectral-bridge-RHI16} give, for
\(v=e^{-sH}\eta\),
\begin{equation}
 v\in\operatorname{Dom}(H^aP^\alpha),\qquad
 \|H^aP^\alpha v\|\le m!s^{-m}\|\eta\|.
 \label{eq:reconstruction-spectral-bridge-RHI20}
\end{equation}
It is invariant under the real translations and spatial rotations.
It is also invariant under each generator: reserve half of the heat
factor and apply the bound to the other half.  Thus generator and
domain assertions in this bridge do not depend on a formal exponential
with an unspecified domain.  Field invariance of this heat domain is
not asserted.

\subsubsection*{The field-product reconstruction beyond this lemma}

The coefficients in \eqref{eq:reconstruction-spectral-bridge-RHI17} have a finite measure and hence a simple
boundary.  The fixed-vector translation theorem alone does not
construct arbitrary field products.  In particular the strict-gap
insertion bound in \eqref{eq:uniform-insertion-bound-UI9} does not permit globally undoing
\(e^{-sH}\) or putting \(s=0\) in an unbounded field formula.

\Cref{ch:h3-energy-band-fields} supplies the additional
construction under the actual UI/PD inputs.  It defines the bounded
sharp-time Riesz family, removes its heat factors only on finite
energy bands, and constructs the spacetime fields on a common
energy domain.  Its exact sewing, simultaneous forward-tube
continuation and full cone tempered boundary apply to arbitrary
field products.  It also derives Lorentz covariance from actual G2.
These results extend the fixed-vector statement in \eqref{eq:reconstruction-spectral-bridge-RHI14} to field
products.

\Cref{ch:quantitative-reconstruction-matching} checks the hypotheses of the
corrected reconstruction theorem by explicit seminorm estimates.
Analytic uniqueness identifies the constructed moments with those of
the theorem, giving local commutativity and the stronger linear
output-order estimate. The theorem
does not assert preservation of the Euclidean factorial exponent.
Nor does either construction assign ordinary Euclidean contact
values: the diagonal-flat input continues to have its stated domain.

\subsubsection*{The common field domain after a Minkowski boundary exists}

We state the operator hypotheses explicitly.  Suppose the
continuation step has produced tempered distributions \(W_n\) on
all of \(\mathcal S(\mathbb R^{4n})\), with \(W_0=1\), and the
positive Wightman form.  This is a new assumption in the rest of this
section.  Let
\begin{equation}
 \begin{gathered}
 \mathcal B=\bigoplus_{n\ge0}^{\rm finite}\mathcal S(\mathbb R^{4n}),
 \qquad (FG)_n=\sum_{r+s=n}F_r\otimes G_s,\\
 F_n^*(x_1,\ldots,x_n)
       =\overline{F_n(x_n,\ldots,x_1)},\qquad
 \omega(F)=\sum_nW_n(F_n).
 \end{gathered}
 \label{eq:reconstruction-spectral-bridge-RHI21}
\end{equation}
The empty component is the unit.  Positivity means
\(\omega(F^*F)\ge0\) for every finite sequence, with no separation
restriction.  Its polarized form is
\(\langle F,G\rangle=\omega(F^*G)\), and its null space is
\(\mathcal N_W=\{F:\omega(F^*F)=0\}\).

\begin{proposition}[The field-domain construction]
\label{thm:explicit-wightman-field-domain}
The completion of \(\mathcal B/\mathcal N_W\) has a common dense
invariant field domain \(\mathcal D=\mathcal B/\mathcal N_W\).
For a one-variable Schwartz test \(f\), left multiplication by its
degree-one sequence defines \(\phi(f)\) there.  It is linear in
\(f\), satisfies
\(\phi(f)\subset\phi(\bar f)^*\), and is closable.  For each
\(\eta\in\mathcal D\), the map \(f\mapsto\phi(f)\eta\) is
continuous from Schwartz space to the Hilbert space.  The vacuum is
cyclic for the field polynomials.  Distributional Poincare covariance
and spacelike commutativity of \(W\), when supplied, give the
corresponding operator identities on \(\mathcal D\).
\end{proposition}

\emph{Proof.}
Positivity gives Cauchy--Schwarz.  A null element \(A\) is therefore
orthogonal to every element.  For arbitrary \(B\),
\begin{equation}
 \|BA\|^2=\omega(A^*B^*BA)
              =\langle A,B^*BA\rangle=0.
 \label{eq:reconstruction-spectral-bridge-RHI22}
\end{equation}
All expressions are finite algebra products.  Thus \(\mathcal N_W\)
is a left ideal, which is exactly what is required for left
multiplication to be well-defined on the quotient.  Let \(\iota(f)\)
denote the degree-one sequence.  Then
\begin{equation}
 \phi(f)[G]=[\iota(f)G],\qquad
 \langle[F],\phi(f)[G]\rangle
             =\langle\phi(\bar f)[F],[G]\rangle.
 \label{eq:reconstruction-spectral-bridge-RHI23}
\end{equation}
The domain is dense by construction and invariant under every such
left multiplication.  The adjoint domain of \(\phi(f)\) contains
this same dense domain, by the second identity.  Hence \(\phi(f)\)
is closable.  For real \(f\) it is symmetric on \(\mathcal D\);
essential self-adjointness has not been inferred.

Fix \(\eta=[G]\), with finitely many nonzero degrees.  Its image
norm squared is a finite sum of terms
\begin{equation}
 W_{r+s+2}(G_r^*\otimes\bar f\otimes f\otimes G_s).
 \label{eq:reconstruction-spectral-bridge-RHI24}
\end{equation}
Temperedness of each distribution and continuity of tensor
multiplication give, for this fixed \(G\), an integer \(j\) and
constant \(C_G\) such that the absolute sum is at most
\(C_Gp_j(f)^2\).  Thus
\(\|\phi(f)\eta\|\le C_G^{1/2}p_j(f)\), proving the asserted
Hilbert-valued distribution continuity.  The order and constant here
may depend on the fixed vector; no degree-independent operator norm
has been postulated.

Finite sums of ordinary Schwartz tensors are dense in every
\(\mathcal S(\mathbb R^{4n})\).  This can be seen directly by
cutting off infinity and using Fourier series on a product box,
whose partial sums converge in every Schwartz seminorm.  The norm of
a homogeneous degree-\(n\) vector is continuous under this convergence,
because its squared norm uses the single tempered distribution
\(W_{2n}\) and a tensor product.  Consequently the polynomial vectors
\(\phi(f_1)\cdots\phi(f_n)\Omega\) have dense span.  Their span
is itself a common dense invariant domain if that smaller domain is
preferred.

A scalar Poincare pullback \(\alpha_g\) is an algebra automorphism
of \(\mathcal B\) commuting with the involution.  If
\(\omega\circ\alpha_g=\omega\), the prescription
\(\mathcal U_g[F]=[\alpha_gF]\) gives a unitary with its exact
inverse.  The usual Schwartz continuity of affine pullbacks and the
same finite form expansion prove strong continuity.  On the common
domain the algebra identity gives
\begin{equation}
 \mathcal U_g\phi(f)\mathcal U_g^{-1}
                        =\phi(\alpha_gf).
 \label{eq:reconstruction-spectral-bridge-RHI25}
\end{equation}
This derives the implementation once distributional covariance has
been supplied; it does not derive a Lorentz identity from the spatial
rotation argument.

Finally suppose the Minkowski distributions give the adjacent
commutator identity for tests \(f,h\) with spacelike-separated
supports and arbitrary spectator tests.  Its pairing with
\(F^*\) and \(G\) says
\begin{equation}
 \langle[F],[\phi(f),\phi(h)][G]\rangle=0
                       \quad(F,G\in\mathcal B).
 \label{eq:reconstruction-spectral-bridge-RHI26}
\end{equation}
The commutator maps the common domain into itself.  Density of the
test classes therefore makes it zero there.  This proves operator
locality from that distributional identity.\hfill\(\square\)



% END MODULE chapters/reconstruction_spectral_bridge.tex

% BEGIN MODULE chapters/h3_energy_band_fields.tex
\subsection{Tempered fields from analytic insertions}
\label{ch:h3-energy-band-fields}

We construct spacetime-smeared fields from the actual Euclidean
insertion estimates.  The construction starts with the bounded analytic
density of \cref{ch:plate-density-from-insertion}, derives a
compatible sharp-time Riesz family, and removes its heat factors only
on finite energy bands.  Temporal smearing then gives a common invariant
dense field domain and positive tempered vacuum moments.  Exact
Euclidean sewing also gives a joint positive-gap continuation and its
tempered real-time boundary.  No unrestricted Minkowski moment sequence
is assumed in this construction.

The hypotheses are those already declared for UI and PD: the actual
ordinary-source and symmetry inputs, the separated-moment continuity,
and the reflection-positive Hilbert evolution with its exact cube cone.
Write \(T_s=e^{-sH}\), \(U_x=e^{ix\cdot P}\), and let
\(\mathcal K_{\rm ord}\subset\mathcal K\) be the reducing ordered
sector of \cref{ch:reconstruction-spectral-interface}.
The round cone is available when the spatial rotation hypothesis of
that section is supplied.  All estimates below can first be made on
\(\mathcal K\) and then restricted to \(\mathcal K_{\rm ord}\).
The field's cyclic space will be identified with this ordered sector.

\subsubsection*{The sharp family is a defined bounded form}

The PD input is an actual operator-norm analytic function
\(V_\eta(t,x)\), uniform in real \(x\in\mathbb R^3\), with
\begin{equation}
 \|V_\eta(t,x)\|\le C\eta^{-N},\qquad
 N=p+5=9,
 \label{eq:h3-energy-band-fields-EB1}
\end{equation}
on its retained complex neighbourhood.  The value \(N=9\) uses the
proved UI exponent \(p=4\).  \eqref{eq:plate-density-from-insertion-PD12} identifies this density with the
original insertion form between pulled-down outside polynomials.
For \(0<s\le1/4\) and \(h\in L^1(\mathbb R^3)\), define
\begin{equation}
 \Phi_s(h)=\int h(x)V_{4s}(0,x)\,d^3x,
       \qquad \|\Phi_s(h)\|\le C s^{-N}\|h\|_1.
 \label{eq:h3-energy-band-fields-EB2}
\end{equation}
The weak integral is defined by bounded sesquilinear forms and Riesz
representation.  Norm continuity of the density also permits a Bochner
integral.  This definition uses no unsandwiched sharp-time operator.

We verify the identification and consistency needed later.  Choose a fixed nonnegative even real
\(\chi\in C_c^\infty(\mathbb R)\), with integral one, and set
\(\chi_\epsilon(t)=\epsilon^{-1}\chi(t/\epsilon)\).
This temporal approximate identity has
\(\|\chi_\epsilon\|_1=1\) and support shrinking to zero.  For compact smooth \(h\), the
actual regularized insertion operator for
\(k_\epsilon(t,x)=\chi_\epsilon(t)h(x)\) has core form
\begin{equation}
 \langle[A],\Phi_s(k_\epsilon)[B]\rangle
      =I(\tau_s A;k_\epsilon;\tau_s B),\qquad
 \Phi_s(k_\epsilon)[B]
      =[(\tau_s k_\epsilon)\cdot\tau_{2s}B].
 \label{eq:h3-energy-band-fields-EB3}
\end{equation}
Here \(\tau_s=\tau_{se_4}\), and the second identity is exactly the
\eqref{eq:plate-density-from-insertion-PD3} vector identity, extended to a compact spatial profile by its
proved density formula.  The two time supports in the displayed product
are separated for small \(\epsilon\).  \eqref{eq:plate-density-from-insertion-PD12} gives
\(\Phi_s(k_\epsilon)=\int\chi_\epsilon(t)h(x)V_{4s}(t,x)\,dt,dx\).
A temporal Cauchy derivative of \eqref{eq:h3-energy-band-fields-EB1} is bounded uniformly in real
\(x\), so
\begin{equation}
 \|\Phi_s(k_\epsilon)-\Phi_s(h)\|
       \le C s^{-N-1}\epsilon\|h\|_1.
 \label{eq:h3-energy-band-fields-EB4}
\end{equation}
Thus the sharp limit is an operator-norm limit of actual core
insertions.  Approximation in spatial \(L^1\) extends this conclusion
to the domain in \eqref{eq:h3-energy-band-fields-EB2}.

Testing \eqref{eq:h3-energy-band-fields-EB3} after an additional time shift of both outside
polynomials gives
\begin{equation}
 \Phi_{s+r}(h)=T_r\Phi_s(h)T_r
               \quad(0<s<s+r\le1/4).
 \label{eq:h3-energy-band-fields-EB5}
\end{equation}
Indeed \(\langle A,T_r\Phi_s(k)T_rB\rangle\) is the insertion
form with outside entries shifted by \(s+r\).  PD's compatibility
on the overlapping pulled-down domains identifies it with the
\(s+r\) form.  Pass to the norm limit in \eqref{eq:h3-energy-band-fields-EB4}, then to spatial
\(L^1\) labels.  Reflection with the even approximants and spatial
translation invariance give, by the same dense-form calculation,
\begin{equation}
 \Phi_s(h)^*=\Phi_s(\bar h),\qquad
 U_y\Phi_s(h)U_y^{-1}=\Phi_s(\tau_y h).
 \label{eq:h3-energy-band-fields-EB6}
\end{equation}
If spatial rotations are supplied, their actual quotient unitaries
similarly intertwine \(\Phi_s(h)\) and \(\Phi_s(r_Rh)\).

For larger \(s\), define \(\Phi_s\) using the right side of
\eqref{eq:h3-energy-band-fields-EB5} with any smaller fixed heat parameter.  The small-parameter
identity proves independence of that choice and then \eqref{eq:h3-energy-band-fields-EB5} for all
\(s,r>0\).  Its safe global estimate is
\(C\min(s,1/4)^{-N}\|h\|_1\).  Only the small-\(s\) polynomial
estimate is used below; no decay as \(s\to\infty\) is required.

The ordered sector is preserved.  To see the domain explicitly,
compact tests supported inside \(0<t_1<\cdots<t_n\) can be covered
by finitely many product boxes with strictly ordered time intervals.
A product-box Fourier approximation gives sums of ordered tensors.
For a general ordered Schwartz test, first cut off infinity and then
cut out the time boundary and adjacent equal-time hyperplanes.  Flatness
on those boundaries makes the cutoff errors tend to zero in every
Schwartz seminorm by Taylor's remainder.  Thus such compact ordered
tensors are dense in the ordered quotient norm.  For each of them,
\eqref{eq:h3-energy-band-fields-EB3} prepends an insertion near time \(s\) and shifts all old entries
by \(2s\), giving another ordered vector.  The bounded norm limit
therefore preserves \(\mathcal K_{\rm ord}\).  The adjoint in \eqref{eq:h3-energy-band-fields-EB6}
does the same, so this subspace reduces every \(\Phi_s(h)\).

\subsubsection*{Removing heat on finite bands}

Let \(P_m=\mathbf1_{[m,m+1)}(H)\), \(m=0,1,\ldots\).
Their orthogonal sum is identity.  The notation
\begin{equation}
 A_{mn}(h)=(e^{sH}P_m)\Phi_s(h)(P_ne^{sH})
       :P_n\mathcal K\longrightarrow P_m\mathcal K
 \label{eq:h3-energy-band-fields-EB7}
\end{equation}
uses bounded outer operators, of norms at most \(e^{s(m+1)}\) and
\(e^{s(n+1)}\).  If a larger heat parameter is used, \eqref{eq:h3-energy-band-fields-EB5} cancels
the extra factors on these same finite bands.  Hence \eqref{eq:h3-energy-band-fields-EB7} is
independent of \(s>0\).  Taking
\(s=[8(1+m+n)]^{-1}\) in \eqref{eq:h3-energy-band-fields-EB2} gives
\begin{equation}
 \|A_{mn}(h)\|\le C_A(1+m+n)^N\|h\|_1,
       \qquad A_{mn}(h)^*=A_{nm}(\bar h).
 \label{eq:h3-energy-band-fields-EB8}
\end{equation}
All spatial covariance identities restrict to these blocks because
the spatial unitaries commute with \(H\).  The same construction
can be made at a spatial point using \(\Phi_s(x)=V_{4s}(0,x)\).
Compatibility and the block bound hold pointwise in \(x\), by
continuity and \eqref{eq:h3-energy-band-fields-EB5} tested against spatial labels, and
\(A_{mn}(h)=\int h(x)A_{mn}(x)\,dx\).

The unregularized block matrix in \eqref{eq:h3-energy-band-fields-EB8} need not define a Hilbert-space
operator at a sharp time.  Its polynomially growing columns have not
been shown square summable.  We next define the spacetime-smeared
operator, for which the missing summability is supplied by time
integration.

\subsubsection*{General spacetime tests and the block estimate}

For \(g\in\mathcal S(\mathbb R\times\mathbb R^3)\), write
\(g_t(x)=g(t,x)\).  The bounded linear map
\(h\mapsto A_{mn}(h)\) on spatial \(L^1\) makes
\(t\mapsto A_{mn}(g_t)\) differentiable in operator norm, with
derivatives \(A_{mn}(\partial_t^jg_t)\).  Define directly
\begin{equation}
 B_{mn}(g)=\int_{\mathbb R}
           e^{itH}A_{mn}(g_t)e^{-itH}\,dt,
 \qquad
 \mathfrak t_q(g)=\sum_{j=0}^q\|\partial_t^jg\|_{L^1(\mathbb R^4)}.
 \label{eq:h3-energy-band-fields-EB9}
\end{equation}
This is a norm integral between the fixed bands.  It is already a
definition for arbitrary spacetime Schwartz tests, not just tensors.

Pull out the scalar phase \(e^{it(m-n)}\).  The remaining generators
\((H-m)|_{P_m\mathcal K}\) and \((H-n)|_{P_n\mathcal K}\) have
norm at most one.  A derivative of their conjugation costs at most
two.  Leibniz's rule followed by \(q\) integrations by parts therefore
gives, when \(|m-n|\ge1\), the bound with factor
\(3^q|m-n|^{-q}\mathfrak t_q(g)\).  Boundary terms vanish by the
Schwartz decay in temporal \(L^1\).  For \(m=n\) use the original
norm integral.  Together these estimates give the explicit safe bound
\begin{equation}
 \|B_{mn}(g)\|\le C_A6^q(1+m+n)^N
                       (1+|m-n|)^{-q}\mathfrak t_q(g).
 \label{eq:h3-energy-band-fields-EB10}
\end{equation}
The adjoint blocks satisfy \(B_{mn}(g)^*=B_{nm}(\bar g)\).
The proof used actual derivatives of a bounded map into fixed energy
bands and did not differentiate an unbounded sharp-time field.

Put \(W=1+H\).  On the \(m\)-th band, powers of \(W\) have
norm at most \((m+2)^k\), while \(\|W^{-l}P_n\|\le(n+1)^{-l}\).
The elementary comparison
\begin{equation}
 \frac{(m+2)^k(1+m+n)^N}{(n+1)^{N+k}}
       \le2^{N+k}(1+|m-n|)^{N+k}
 \label{eq:h3-energy-band-fields-EB11}
\end{equation}
shows exactly how the energy weights are paid.  Set
\(q=N+k+2\), and let
\(Z_2=2\sum_{d\ge0}(1+d)^{-2}<\infty\).
The scalar row and column sums of the weighted blocks are at most
\(C_A2^{N+k}6^{N+k+2}Z_2\mathfrak t_{N+k+2}(g)\).
Block Schur consequently gives
\begin{equation}
 \|W^k\phi(g)W^{-(N+k)}\|
       \le C_k\mathfrak t_{N+k+2}(g)
       \le C'_kp_{N+k+7}(g).
 \label{eq:h3-energy-band-fields-EB12}
\end{equation}
For the last inequality use the ordinary four-dimensional Schwartz
seminorm: its weight of order at least five is integrable, and it
controls all the indicated temporal derivatives.  These are actual
seminorm estimates for non-tensor spacetime tests.

Here is the operator and domain construction behind \eqref{eq:h3-energy-band-fields-EB12}.  For
\(k=0\), the block matrix \(B_{mn}(g)W^{-N}P_n\) defines a bounded
operator by Schur.  Multiply it on the right by \(W^N\) on
\(\operatorname{Dom}W^N\) to define \(\phi(g)\).
For finite-energy vectors its blocks are exactly \eqref{eq:h3-energy-band-fields-EB9}.
For \(\xi\in\operatorname{Dom}W^{N+k}\), approximate by
\(Q_R\xi\), where \(Q_R=\sum_{m=0}^RP_m\).
The weighted Schur estimate makes both \(\phi(g)Q_R\xi\) and
\(W^k\phi(g)Q_R\xi\) converge.  Closedness of \(W^k\) proves
the claimed domain inclusion and the operator identity in \eqref{eq:h3-energy-band-fields-EB12}.

It follows that
\begin{equation}
 \mathcal D_\infty=\bigcap_{k\ge0}\operatorname{Dom}H^k
 \quad\hbox{is dense and}\quad
 \phi(g)\mathcal D_\infty\subset\mathcal D_\infty.
 \label{eq:h3-energy-band-fields-EB13}
\end{equation}
Finite-energy vectors make it dense.  The adjoint-block identity,
first on such vectors and then by their graph-norm approximation,
gives
\begin{equation}
 \langle\xi,\phi(g)\eta\rangle
       =\langle\phi(\bar g)\xi,\eta\rangle
                 \quad(\xi,\eta\in\mathcal D_\infty).
 \label{eq:h3-energy-band-fields-EB14}
\end{equation}
Thus each field restricted to this common invariant domain is
closable, and it is symmetric for real tests.  \eqref{eq:h3-energy-band-fields-EB12} also
proves continuity into every graph seminorm for each fixed domain
vector.  Essential self-adjointness is not inferred.

All energy projections, sharp insertions and real translations reduce
\(\mathcal K_{\rm ord}\).  The same block limits therefore preserve
\(\mathcal D_\infty\cap\mathcal K_{\rm ord}\).  Spatial covariance
and time smearing give exact domain identities
\begin{equation}
 \begin{gathered}
 \mathscr U(t,y)=e^{itH}U_y,\qquad
 \mathscr U(t,y)\mathcal D_\infty=\mathcal D_\infty,\\
 \mathscr U(t,y)\phi(g)\mathscr U(t,y)^{-1}
                      =\phi(\tau_{(t,y)}g).
 \end{gathered}
 \label{eq:h3-energy-band-fields-EB15}
\end{equation}
Spatial rotations have the analogous covariance when supplied.  Our
time convention is fixed by the positive sign in \(e^{itH}\).
With the standard Minkowski pairing \(Et-p\cdot x\), its spatial
momentum is \(-P\); the round cone is unchanged by this sign.

\subsubsection*{Positive tempered moments with an explicit extension}

The vacuum belongs to \(\mathcal D_\infty\), and \eqref{eq:h3-energy-band-fields-EB12} can be
iterated.  Normalize \(W_0=\|\Omega\|^2=1\).
For each fixed \(n\),
\begin{equation}
 \|\phi(g_1)\cdots\phi(g_n)\Omega\|
       \le D_n\prod_{j=1}^np_{M_n}(g_j),\qquad M_n=nN+7.
 \label{eq:h3-energy-band-fields-EB16}
\end{equation}
Indeed the energy exponent needed at the next factor increases by
\(N\) at each step; the largest test order is \(nN+7\).
Define the vacuum moments on tensors by
\(W_n(g_1\otimes\cdots\otimes g_n)
 =\langle\Omega,\phi(g_1)\cdots\phi(g_n)\Omega\rangle\).
Multilinearity factors through the algebraic tensor product.  Its map
into functions is injective on every finite collection: choose point
evaluations distinguishing the finite-dimensional function span in
each slot and invert those evaluation matrices.  Tensor relations
therefore introduce no ambiguity.

We prove a fixed-order tempered estimate explicitly.  Take a smooth unit product-grid partition
and expand each localized test on period-eight boxes in every
four-dimensional slot.  Multiply each Fourier wave by a fixed
larger compact cutoff.  Its one-slot \(p_M\) norm is at most
\(C_M(1+|z|)^M(1+|l|)^M\), where \(z\in\mathbb Z^4\) is the
box center and \(l\in\mathbb Z^4\) its Fourier index.
Choose \(2r=2\lceil(M_n+6)/2\rceil\), and integrate
\((1-\Delta_i)^r\) in each slot when estimating the coefficients.
This gives the factor \(\prod_i(1+|l_i|^2)^{-r}\).
Derivatives of the fixed cutoffs cost a finite constant depending on
\(n,r\).  With
\begin{equation}
 K_n=n(M_n+8),
 \label{eq:h3-energy-band-fields-EB17}
\end{equation}
the seminorm \(p_{K_n}(F)\) pays both their total derivative order
\(2rn\le n(M_n+7)\) and the per-slot location weight
\((1+|z_i|)^{-(M_n+8)}\).  The latter follows from
\(1+\sum_i|x_i|^2\ge\prod_i(1+|x_i|^2)^{1/n}\), with fixed-box
weight comparisons.  After the tensor bound in \eqref{eq:h3-energy-band-fields-EB16}, the absolute
series is controlled by the products of
\begin{equation}
 \sum_{z\in\mathbb Z^4}(1+|z|)^{-8}<\infty,
 \qquad
 \sum_{l\in\mathbb Z^4}(1+|l|)^{M_n}(1+|l|^2)^{-r}<\infty.
 \label{eq:h3-energy-band-fields-EB18}
\end{equation}
The second sum converges since \(2r-M_n\ge6>4\).
The Fourier series of the localized smooth tests converges in every
required seminorm.  For a tensor its coefficients factor, so the
absolutely summed value agrees with its original multilinear value.
This constructs the unique continuous extension, with
\begin{equation}
 |W_n(F)|\le D'_np_{K_n}(F)
                     \quad(F\in\mathcal S(\mathbb R^{4n})).
 \label{eq:h3-energy-band-fields-EB19}
\end{equation}
The same absolute series, using the vector bound in \eqref{eq:h3-energy-band-fields-EB16}, constructs
the Hilbert-valued vacuum-product distribution.  Higher graph seminorms
are treated by starting the iteration of \eqref{eq:h3-energy-band-fields-EB12} at the corresponding
energy exponent.

Thus \(W_n\) is tempered at every fixed \(n\).  The displayed
order \(K_n=O(n^2)\) is sufficient for that assertion; no linear
order bound or preserved factorial moment exponent is claimed here.
For finite tensor sequences, \eqref{eq:h3-energy-band-fields-EB14} gives
\begin{equation}
 \sum_{r,s}W_{r+s}(F_r^*\otimes F_s)
       =\left\|\sum_s\phi(F_s)\Omega\right\|^2\ge0.
 \label{eq:h3-energy-band-fields-EB20}
\end{equation}
Here \(F_r^*\) reverses and conjugates its arguments, and for a
tensor \(\phi(F_s)\Omega\) means its field product.  The
Hilbert-valued extension just constructed defines the same vector for
a general Schwartz \(F_s\).  Ordinary tensor density and \eqref{eq:h3-energy-band-fields-EB19}
then prove \eqref{eq:h3-energy-band-fields-EB20} on all finite unrestricted Schwartz sequences.
This proves the required positive Minkowski moment functional from
the actual insertion input.

\subsubsection*{Exact Euclidean words and joint positive-gap analyticity}

Let \(z=(z_1,\ldots,z_{n-1})\) have positive real components and
take spatial labels \(h_j\in L^1(\mathbb R^3)\).  Choose small
\(a_j>0\) with \(a_j+a_{j+1}<\operatorname{Re}z_j\).  Define
\begin{equation}
 \begin{gathered}
 \mathcal E_n(z;h)=\langle\Omega,
        \Phi_{a_1}(h_1)\mathcal P(z,a;h)\Omega\rangle,\\
 \mathcal P(z,a;h)=\prod_{j=1}^{n-1}
       \left[T_{z_j-a_j-a_{j+1}}\Phi_{a_{j+1}}(h_{j+1})\right].
 \end{gathered}
 \label{eq:h3-energy-band-fields-EB21}
\end{equation}
The product is ordered by increasing \(j\) from left to right and
is identity for \(n=1\).
Here \(T_z=e^{-zH}\) is the actual bounded analytic semigroup.
Changing \(a_j\) to a smaller \(b_j\) inserts the two factors
\(T_{a_j-b_j}\) of \eqref{eq:h3-energy-band-fields-EB5}; they join their adjacent heat factors
and give precisely the formula with \(b_j\).  The endpoint heat
factors fix \(\Omega\).  Two arbitrary allowed choices are compared
through a common smaller choice.  Thus \eqref{eq:h3-energy-band-fields-EB21} is independent of every
auxiliary parameter.

On a neighbourhood of any positive-real-gap point, keep the \(a_j\)
fixed.  The formula is a finite product of bounded norm-holomorphic
operators.  The matching identities just proved glue these formulas
on overlaps, giving joint analyticity on
\(\operatorname{Re}z_j>0\) for every gap simultaneously.  With
\begin{equation}
 r=\min\{1/8,\tfrac14\min_j\operatorname{Re}z_j\},
 \qquad
 |\mathcal E_n(z;h)|\le C^nr^{-Nn}\prod_j\|h_j\|_1,
 \label{eq:h3-energy-band-fields-EB22}
\end{equation}
the bound follows by taking every \(a_j=r\).  When \(n=1\), fix
\(a_1=1/8\); its vacuum value is independent of that choice.

We verify that this is the actual Euclidean continuation.  For real
ordered times \(\tau_1<\cdots<\tau_n\), replace every sharp label
by \(\chi_\epsilon(t)h_j(x)\) in the allowed core formula.  Starting
at the rightmost factor, repeated use of \eqref{eq:h3-energy-band-fields-EB3} puts the first insertion
at \(a_1\) and the later ones at
\(a_1+(\tau_j-\tau_1)\), with the corresponding narrow profiles.
A common translation by \(\tau_1-a_1\) identifies its vacuum
expectation with the original Euclidean moment.  Operator-norm
convergence in \eqref{eq:h3-energy-band-fields-EB4} therefore identifies the point-time limit with
\(\mathcal E_n(\tau_{j+1}-\tau_j;h)\).

For compact temporal profiles in disjoint ordered intervals there is
a fixed positive gap.  Choose the auxiliary heat parameters uniformly
below that gap.  Approximate the profiles by their convolutions with
the narrow kernels.  The original moments converge by local \(C^1\)
continuity, while \eqref{eq:h3-energy-band-fields-EB22} permits Fubini and dominated convergence on
the operator side.  This recovers their exact original moment values.
Product-box Fourier approximation gives the same identification for
general compact ordered tests, not only time-space tensors.
No diagonal Euclidean contact prescription is involved.

The corresponding vector formula, when all absolute times are
positive, is
\begin{equation}
 T_{\tau_1-a_1}\Phi_{a_1}(h_1)
       \mathcal P((\tau_{j+1}-\tau_j)_j,a;h)\Omega,
 \label{eq:h3-energy-band-fields-EB23}
\end{equation}
with \(a_1<\tau_1\).  The same core argument and norm limits
identify its ordered-profile integrals with the actual ordered
Euclidean vectors.  It is a bounded formula throughout the retained
positive gaps.

\subsubsection*{A joint tempered boundary, not separate limits}

For positive gap regulators \(\epsilon_j\), define on tensor tests
\begin{equation}
 W_n^\epsilon(g_1\otimes\cdots\otimes g_n)
   =\langle\Omega,\phi(g_1)T_{\epsilon_1}\phi(g_2)
                   \cdots T_{\epsilon_{n-1}}\phi(g_n)\Omega\rangle.
 \label{eq:h3-energy-band-fields-EB24}
\end{equation}
Every heat contraction has norm at most one in every energy graph
norm.  The proof of \eqref{eq:h3-energy-band-fields-EB16}--\eqref{eq:h3-energy-band-fields-EB19} therefore gives, uniformly in all
\(\epsilon_j\ge0\),
\begin{equation}
 |W_n^\epsilon(F)|\le D'_np_{K_n}(F),\qquad
 W_n^\epsilon(F)\longrightarrow W_n(F)
        \quad\hbox{as all }\epsilon_j\downarrow0.
 \label{eq:h3-energy-band-fields-EB25}
\end{equation}
For tensors the convergence follows by replacing the finitely many
heat factors one at a time.  Strong convergence in every graph norm,
combined with \eqref{eq:h3-energy-band-fields-EB12}, controls each term.  For an arbitrary Schwartz
test, first approximate it by finite tensors in \(p_{K_n}\).
The uniform bound controls the two approximation errors, and tensor
convergence controls the fixed approximation.  This proves joint,
regulator-path-independent convergence in \(\mathcal S'\).

To identify it as the boundary of \eqref{eq:h3-energy-band-fields-EB21}, first insert the finite
total-energy cutoff \(Q_R\) at every field.  All generators and
field blocks are then bounded, so time integration gives exactly
\begin{equation}
 z_j=\epsilon_j-i(t_{j+1}-t_j)
 \label{eq:h3-energy-band-fields-EB26}
\end{equation}
in that formula.  For fixed positive regulators, its intermediate
energy-band sums are absolutely convergent: \eqref{eq:h3-energy-band-fields-EB8} contributes only
fixed polynomial factors, while every intermediate energy is damped
by its positive gap.  The convergence of the cutoff expressions is
uniform in the real times and bounded by an integrable product of
the spatial \(L^1\) norms of the time slices of the tests.
Fubini and dominated convergence therefore identify their integrals
with \eqref{eq:h3-energy-band-fields-EB24}.  On the field side, the cutoff products converge by the
graph bounds and \(Q_R\to I\) in every graph norm.
\eqref{eq:h3-energy-band-fields-EB25}--\eqref{eq:h3-energy-band-fields-EB26} consequently prove a jointly tempered
boundary of the actual sewn Euclidean functions.  Independent
one-variable boundary statements have not been substituted for it.

\subsubsection*{The full spectral carrier and the cyclic sector}

Assume now the supplied round cone.  The finite-energy cutoff point
field
\(A_R(x)=\sum_{m,n\le R}A_{mn}(x)\) is bounded, spatially covariant,
and has bounded translation generators on its cutoff range.
Its vacuum words have the usual difference-variable form with one
joint translation between each adjacent pair of fields.  If a test
in one difference variable has Fourier transform vanishing on the
forward carrier, integrating that translation gives zero by the joint
spectral theorem.  Other differences and spectator fields can be
held fixed during this integration.  On a compact Fourier set outside
the product carrier, choose a finite cover by product boxes, in each
of which one difference momentum is outside its carrier.  A subordinate
partition and product Fourier approximation reduce each localized
test to the just-proved zero identities.  This proves the full product
support condition.

The cutoff vacuum distributions converge to \(W_n\) by \eqref{eq:h3-energy-band-fields-EB12} and
the same uniform extension estimate as in \eqref{eq:h3-energy-band-fields-EB25}.
Their difference-variable limits are legitimate distributions:
fix a Schwartz test \(\rho\) of integral one in the common center
coordinate, and pair the full distributions with
\(\rho(x_1)\psi(x_2-x_1,\ldots,x_n-x_{n-1})\).
This is a continuous map from difference tests to full tests.
The cutoff identities show that the limiting difference functional
is independent of \(\rho\) and that its pairing with a full test
uses the center integral of that test.  Thus it is the actual
difference-variable distribution of \(W_n\).  Support in a fixed
closed set passes to this tempered limit.  Hence
\begin{equation}
 \operatorname{supp}\widehat{W_n^{\rm diff}}
                       \subset(\overline V_+)^{n-1},
 \qquad \overline V_+=\{(E,p):E\ge|p|_2\},
 \label{eq:h3-energy-band-fields-EB27}
\end{equation}
with the Fourier convention of \eqref{eq:h3-energy-band-fields-EB15} and \eqref{eq:h3-energy-band-fields-EB26}.  In unreduced
momenta this is the corresponding condition on all partial sums,
not merely a bound on total translations of one fixed vector.
Without the spatial rotation input the same proof retains the cube
carrier provided by \eqref{eq:plate-density-from-insertion-PD1}.

Let \(\mathcal H_\phi\) be the closed span of field-polynomial
vacuum vectors.  They lie in \(\mathcal K_{\rm ord}\), and \eqref{eq:h3-energy-band-fields-EB15}
makes their span invariant under both signs of real translations.
Thus \(\mathcal H_\phi\) reduces \(H\), its heat factors and its
energy projections.  The reverse inclusion, and hence cyclicity, follows as below.

First, polynomial vectors are graph-dense in
\(\mathcal H_\phi\cap\operatorname{Dom}W^l\) for every \(l\).
For a smooth compact spectral cutoff \(c_R(H)\), its inverse Fourier
transform is integrable.  Therefore \(c_R(H)\) is an integral of
real time translations.  Applied to a polynomial vector, that
integral is approximated in the \(W^l\) graph norm by finite sums
of translated polynomial vectors: the orbit is graph-norm continuous
and its graph norm is constant.  Given an arbitrary vector in the
stated domain, first cut its spectrum off and then approximate the
cutoff vector by these smoothed polynomial approximants.  Finally let
the cutoff tend to identity.  This proves graph density.
The energy estimate in \eqref{eq:h3-energy-band-fields-EB12} now shows that every \(\phi(g)\) maps
\(\mathcal H_\phi\cap\mathcal D_\infty\) into
\(\mathcal H_\phi\).

For compact smooth \(h\), the bounded operators
\(T_s\phi(\chi_\epsilon\otimes h)T_s\) converge in operator
norm to \(\Phi_s(h)\).  On finitely many bands this follows from
the temporal approximate identity and the bounded band generators.
The tails converge uniformly: their block norms are bounded by
\begin{equation}
 C(1+m+n)^Ne^{-s(m+n)}\|h\|_1,
 \label{eq:h3-energy-band-fields-EB28}
\end{equation}
independently of \(\epsilon\), since \(\|\chi_\epsilon\|_1=1\).
The row and column sums outside a growing finite square tend to
zero.  The limiting blocks are exactly those of \(\Phi_s(h)\)
by \eqref{eq:h3-energy-band-fields-EB7}.  Each approximating operator preserves
\(\mathcal H_\phi\): the right heat factor sends its vectors into
all energy domains, the preceding graph-density argument handles the
field, and the left factor preserves the subspace.  The norm limit
does so as well.  Spatial \(L^1\) approximation extends this to all
labels in \eqref{eq:h3-energy-band-fields-EB2}.

Every vector in the exact Euclidean sewing formula in \eqref{eq:h3-energy-band-fields-EB23} therefore
belongs to \(\mathcal H_\phi\).  Its ordered-profile integrals are
the dense ordered Euclidean vectors established above.  Consequently
\begin{equation}
 \mathcal H_\phi=\mathcal K_{\rm ord}.
 \label{eq:h3-energy-band-fields-EB29}
\end{equation}
This proves the actual cyclic realization and its inherited
separability, evolution and phase vacuum conclusions.  It makes no
assertion that \(\mathcal K_{\rm ord}=\mathcal K\).

\subsubsection*{The joint forward tube and its full cone boundary}

The preceding argument also supplies the simultaneous complex spatial
variables under the round-cone hypothesis; the following bounded formula proves joint continuation.
Work on \(\mathcal K_{\rm ord}\) and put
\begin{equation}
 S(z,w)=e^{-zH+iw\cdot P},\qquad
 \operatorname{Re}z>|\operatorname{Im}w|_2,
 \qquad \Phi_a^0=V_{4a}(0,0).
 \label{eq:h3-energy-band-fields-EB30}
\end{equation}
The joint spectral proof in \cref{ch:reconstruction-spectral-interface}
makes \(S\) a norm-holomorphic contraction on this tube.  It also
shows directly that derivatives of total order \(l\) on a smaller
tube with margin \(\delta\) have norm at most
\(l!\delta^{-l}\), because \(|P|_2\le H\).
The point operator \(\Phi_a^0\) is the norm limit of the spatial
approximate identities in \eqref{eq:h3-energy-band-fields-EB2}, by continuity of \(V_{4a}\).
It therefore preserves the ordered sector, has norm at most
\(Ca^{-N}\), and satisfies exactly the same heat consistency.

For \(\delta_j=\operatorname{Re}z_j-
|\operatorname{Im}w_j|_2>0\), choose
\(a_j+a_{j+1}<\delta_j\) and define
\begin{equation}
 \begin{gathered}
 \mathscr E_n(z,w)=\left\langle\Omega,\Phi_{a_1}^0
    \prod_{j=1}^{n-1}
       \left[S(z_j-a_j-a_{j+1},w_j)\Phi_{a_{j+1}}^0\right]
                                      \Omega\right\rangle,\\
 r=\min\{1/8,\tfrac14\min_j\delta_j\},\qquad
 |\mathscr E_n(z,w)|\le C^nr^{-Nn}.
 \end{gathered}
 \label{eq:h3-energy-band-fields-EB31}
\end{equation}
The product has the same order as \eqref{eq:h3-energy-band-fields-EB21}.  Decreasing the auxiliary
\(a_j\) and joining their extra heat factors to the adjacent joint
semigroups proves independence, exactly as before: the heat factors
commute with \(S\), and every resulting joint semigroup remains
inside its contraction tube.  Fixed auxiliary parameters on a
neighbourhood give joint norm holomorphy, and their independence
glues the neighbourhoods.  Taking every \(a_j=r\) proves the bound.
For \(n=1\) use the constant vacuum value with \(a_1=1/8\).

For real spatial gaps, spatial covariance of the point operators
identifies this formula with the spatial point version of \eqref{eq:h3-energy-band-fields-EB21}.
Spatial approximate identities converge in norm, so the exact
Euclidean sewing already proved identifies it with the original
ordered Euclidean density and with its spatially smeared values.
Thus \eqref{eq:h3-energy-band-fields-EB31} is a jointly defined function of all \(4(n-1)\)
complex difference variables, not merely a family indexed by real
spatial labels.

Let \(\epsilon_j>|y_j|_2\), and insert
\(R_j=S(\epsilon_j,iy_j)=e^{-\epsilon_jH-y_j\cdot P}\)
between the fields in \eqref{eq:h3-energy-band-fields-EB24}, obtaining \(W_n^{\epsilon,y}\).
The round cone gives \(\|R_j\|\le1\).  These operators commute
with \(W\), are contractions in every graph norm, and converge
strongly there to identity as \((\epsilon_j,y_j)\to0\) within
that cone.  The last assertion is dominated convergence in the joint
spectral measure: the multipliers are bounded by one and converge
pointwise to one.  The proof of \eqref{eq:h3-energy-band-fields-EB25} therefore gives
\begin{equation}
 |W_n^{\epsilon,y}(F)|\le D'_np_{K_n}(F),\qquad
 W_n^{\epsilon,y}(F)\longrightarrow W_n(F)
       \quad\hbox{for every }F\in\mathcal S(\mathbb R^{4n}).
 \label{eq:h3-energy-band-fields-EB32}
\end{equation}
This is a simultaneous limit along every cone-admissible approach.

Finite energy cutoffs identify the integrated function in \eqref{eq:h3-energy-band-fields-EB31} with
\(W_n^{\epsilon,y}\) at
\begin{equation}
 z_j=\epsilon_j-i(t_{j+1}-t_j),\qquad
 w_j=x_{j+1}-x_j+iy_j.
 \label{eq:h3-energy-band-fields-EB33}
\end{equation}
Indeed the inserted spectral multipliers are now bounded by
\(e^{-(\epsilon_j-|y_j|_2)H}\).  Hence each intermediate energy
sum has the same absolute polynomial-times-exponential bound as in
\eqref{eq:h3-energy-band-fields-EB26}, uniformly in all real coordinates.  Finite-cutoff Fubini,
then dominated convergence, proves the identity on tensor tests.
The bound in \eqref{eq:h3-energy-band-fields-EB31} makes its fixed-regulator integral continuous on
Schwartz tests, while \eqref{eq:h3-energy-band-fields-EB32} does the same on the field side; tensor
density proves equality on every such test.  \eqref{eq:h3-energy-band-fields-EB32} is
therefore the actual joint tempered boundary of \eqref{eq:h3-energy-band-fields-EB31}, with all
complex spatial and time regulators allowed by the round cone.
In the convention of \eqref{eq:h3-energy-band-fields-EB15}, these are difference points
\((t+i\epsilon,x+iy)\) with imaginary part in the forward cone.

\subsubsection*{Lorentz covariance from the actual Euclidean rotations}

For this step assume the full scalar Euclidean covariance G2 of the
actual diagonal-flat moments.  The finite hypercubic group alone does
not supply this hypothesis.  In particular G2 includes every real
rotation in a time--space coordinate plane.  No Lorentz covariance
of the field has been used in constructing the tube.

Put \(\zeta_j=(iz_j,w_j)\).  The domain of \eqref{eq:h3-energy-band-fields-EB31} is precisely
\(\operatorname{Im}\zeta_j\in V_+\) for every gap.  Denote the
same holomorphic function by \(\mathcal F_n(\zeta)\).  Its
Euclidean slice is \(\zeta_j=(i\tau_j,r_j)\), with
\(\tau_j>0\) and \(r_j\in\mathbb R^3\).  We first justify
pointwise use of G2 on this slice.  A compact neighbourhood of any
such Euclidean configuration has strictly ordered time coordinates.
The sewing formula represents its actual moment functional by the
continuous point density in \eqref{eq:h3-energy-band-fields-EB31}.  For a sufficiently small fixed
Euclidean rotation, its rotated neighbourhood still has strictly
ordered times.  Apply G2 to arbitrary smooth tests supported in this
neighbourhood.  The scalar change-of-variables identity, whose
Jacobian is one, makes the two continuous densities equal there.
Their difference is zero against every compact smooth test and hence
zero pointwise.  Thus the local rotation identity is an identity for
the actual sewn function, not an imposed extension across the diagonals.
Common Euclidean translations allow the first point to be set at zero,
so rotations act linearly on all the difference variables.

Choose a time--space rotation with infinitesimal action
\(\dot\tau_j=r_j^k\), \(\dot r_j^k=-\tau_j\), simultaneously
for every gap.  Differentiating the just-proved pointwise identity
and using \(\zeta_j^0=i\tau_j\) gives
\begin{equation}
 i\mathcal B_k\mathcal F_n=0\quad\hbox{on the Euclidean slice},
 \qquad
 \mathcal B_k=\sum_{j=1}^{n-1}
   \left(\zeta_j^k\partial_{\zeta_j^0}
                    +\zeta_j^0\partial_{\zeta_j^k}\right).
 \label{eq:h3-energy-band-fields-EB34}
\end{equation}
The derivative exists because the point density is holomorphic on
a complex neighbourhood of that slice.  The slice is maximal totally
real: the complex linear coordinate change
\((\tau_j,r_j)\mapsto(i\tau_j,r_j)\) sends real open boxes
to its open pieces.  A holomorphic function vanishing on such a box
vanishes on a complex neighbourhood, by applying the one-variable
identity theorem successively in each coordinate.  The product tube
is connected, indeed convex.  Apply this observation to the
holomorphic function \(\mathcal B_k\mathcal F_n\); it follows that
it vanishes on the entire tube.

Every real proper orthochronous boost preserves \(V_+\), so its
whole real orbit stays inside the tube.  Differentiating
\(s\mapsto\mathcal F_n(\Lambda_k(s)\zeta)\) now gives zero
at every \(s\).  The function is constant along that orbit.
Spatial rotations have the same invariance, either by G2 on the
slice and the same uniqueness argument, or by their already
implemented quotient unitaries.  These transformations generate
\(\mathrm{SO}_+(1,3)\).  Explicitly, for any \(\Lambda\) in
this group write \(\Lambda e_0=(\gamma,\gamma v)\), where
\(|v|<1\).  The boost in direction \(v\) sends \(e_0\) to
this vector.  Its inverse times \(\Lambda\) fixes \(e_0\) and
is a spatial proper rotation.  A directional boost is a spatial
rotation conjugate of a coordinate boost.  Therefore
\begin{equation}
 \mathcal F_n(\Lambda\zeta)=\mathcal F_n(\zeta)
              \quad(\Lambda\in\mathrm{SO}_+(1,3)).
 \label{eq:h3-energy-band-fields-EB35}
\end{equation}

The boundary passage uses the full cone limit proved above.  If
\(\eta_j\in V_+\), then also \(\Lambda^{-1}\eta_j\in V_+\).
Changing real variables \(\xi=\Lambda\chi\) in the integral
of \(\mathcal F_n(\xi+i\eta)\) against
\(f(\Lambda^{-1}\xi)\) changes its regulator to
\(\Lambda^{-1}\eta\), by \eqref{eq:h3-energy-band-fields-EB35}, and leaves the Jacobian one.
Both regulator families have the same joint tempered limit by
\eqref{eq:h3-energy-band-fields-EB32}.  Consequently
\begin{equation}
 W_n^{\rm diff}(f\circ\Lambda^{-1})=W_n^{\rm diff}(f),
 \qquad
 W_n(\alpha_\Lambda F)=W_n(F),
 \label{eq:h3-energy-band-fields-EB36}
\end{equation}
where \(\alpha_\Lambda\) pulls every argument back by
\(\Lambda^{-1}\).  The second identity follows from the center
integration defining the difference distribution, which commutes
with a determinant-one simultaneous linear change of coordinates.
The Schwartz pullback is continuous for every fixed \(\Lambda\).
This proves distributional Lorentz covariance without a separate
boundary-exchange assertion.

\subsubsection*{The covariant field representation on the ordered sector}

Let \(X(F)\) denote a finite sum of the vacuum-product vectors
constructed in \eqref{eq:h3-energy-band-fields-EB16}--\eqref{eq:h3-energy-band-fields-EB20}.  Their span is dense in
\(\mathcal K_{\rm ord}\) by \eqref{eq:h3-energy-band-fields-EB29}.  Define there
\begin{equation}
 U(\Lambda)X(F)=X(\alpha_\Lambda F),\qquad
 U(\Lambda)\Omega=\Omega.
 \label{eq:h3-energy-band-fields-EB37}
\end{equation}
\eqref{eq:h3-energy-band-fields-EB20} and \eqref{eq:h3-energy-band-fields-EB36}, applied by polarization to two finite
sequences, preserve their inner product.  In particular a zero
vector is sent to zero.  The map is thus well-defined and isometric
on the actual quotient, and its inverse is the same construction
for \(\Lambda^{-1}\).  It extends to a unitary on the ordered
sector.  The pullback composition law gives a representation
of \(\mathrm{SO}_+(1,3)\) on this sector.
On the spatial rotation subgroup it equals the quotient unitary
already constructed: both maps have the same action on every
field-polynomial vacuum vector, whose span is dense by \eqref{eq:h3-energy-band-fields-EB29}.

If \(\Lambda\to I\), its pullbacks converge in every Schwartz
seminorm.  This follows by the mean-value formula for the derivatives
and polynomial weights on a bounded matrix neighbourhood of identity;
one additional derivative and weight suffice for each fixed
seminorm.  The Hilbert-valued extension after \eqref{eq:h3-energy-band-fields-EB19} then makes
\(X(\alpha_\Lambda F)\to X(F)\).  Unitarity and density extend
this convergence to all vectors.  Hence \(U\) is strongly
continuous.  The translation implementation already satisfies
\(\mathscr U(a)X(F)=X(\tau_a F)\).  The exact pullback identity
therefore gives
\begin{equation}
 \begin{gathered}
 U(\Lambda)\mathscr U(a)U(\Lambda)^{-1}
                       =\mathscr U(\Lambda a),\\
 \mathscr U(a)U(\Lambda)\mathscr U(b)U(M)
                       =\mathscr U(a+\Lambda b)U(\Lambda M).
 \end{gathered}
 \label{eq:h3-energy-band-fields-EB38}
\end{equation}
The second line exhibits the strongly continuous proper
orthochronous Poincare representation, with the positive spectral
cone already established.  The notation \(M\) in this formula
denotes another Lorentz matrix.

Covariance also holds on the energy domain.  Stone's theorem applied to the
first identity in \eqref{eq:h3-energy-band-fields-EB38} gives the exact self-adjoint identity
\(U(\Lambda)HU(\Lambda)^{-1}=\gamma H+b\cdot P\), where
\((\gamma,b)=\Lambda e_0\), \(\gamma>|b|\).  The right side
is defined by the joint spectral calculus.  The round cone gives
\begin{equation}
 (\gamma-|b|)H\le\gamma H+b\cdot P
                   \le(\gamma+|b|)H.
 \label{eq:h3-energy-band-fields-EB39}
\end{equation}
These are pointwise inequalities in that calculus.  Their positive
lower and upper constants show that the domains of every positive
integer power of the two operators agree, with equivalent graph
norms.  Thus \(U(\Lambda)\mathcal D_\infty=
\mathcal D_\infty\) on the ordered sector.  The polynomial
intertwining identity follows directly from \eqref{eq:h3-energy-band-fields-EB37} and multiplication
of the field polynomials.  Polynomial vectors are graph-dense in
every ordered-sector energy domain by the proof preceding \eqref{eq:h3-energy-band-fields-EB28},
now using \eqref{eq:h3-energy-band-fields-EB29}.  The graph bound in \eqref{eq:h3-energy-band-fields-EB12}, and the equivalent graph
norms just proved, extend that identity to the common domain:
\begin{equation}
 U(\Lambda)\phi(g)U(\Lambda)^{-1}
                =\phi(\alpha_\Lambda g)
       \quad\hbox{on }\mathcal D_\infty\cap\mathcal K_{\rm ord}.
 \label{eq:h3-energy-band-fields-EB40}
\end{equation}
This is a common-domain covariant scalar field construction on the
actual cyclic space.  No action of these Lorentz unitaries on
\(\mathcal K\ominus\mathcal K_{\rm ord}\) is needed or asserted.

The construction has supplied common-domain Hermitian fields,
positive tempered moments, the full translation spectral carrier,
ordered Euclidean matching, a joint forward-tube tempered boundary,
and Lorentz covariance under G2.  Spacelike locality still requires
the permutation and complex-continuation argument.  It does not
follow from the block estimate or the real boost identities alone
and has not been asserted here.  All underlying ordinary-source and
ultraviolet gauge-theory hypotheses retain their declared conditional
scope.


% END MODULE chapters/h3_energy_band_fields.tex

% BEGIN MODULE chapters/quantitative_reconstruction_matching.tex
\subsection{Quantitative reconstruction and local commutativity}
\label{ch:quantitative-reconstruction-matching}

The fields of \cref{ch:h3-energy-band-fields} were constructed
before any local-commutativity conclusion was used.  Their common
domain, positive moments, ordered Euclidean values and joint tempered
boundary are therefore available for a comparison with the corrected
reconstruction theorem.  We make that comparison explicitly.  It
supplies locality on the actual ordered Hilbert sector, together with
the stronger distribution-order conclusion of the imported theorem.
The proof does not infer locality from a spectral cone alone.

Throughout this section the ordinary-source, ultraviolet and symmetry
hypotheses retain the scope declared in the preceding sections.  In
particular, full scalar Euclidean covariance means G2 on the actual
diagonal-flat moments.  It is a condition on the chosen state or
component, not a conclusion drawn merely from the symmetry of a
mixture.  Clustering is used only for the conclusions that explicitly
invoke a clustering component.

\subsubsection*{The distribution space and its numerical seminorms}

For \(d\ge1\), use the global Euclidean seminorm
\begin{equation}
 p_j^{(d)}(f)=\max_{|\alpha|\le j}\sup_{x\in\mathbb R^d}
       (1+|x|^2)^{j/2}|\partial^\alpha f(x)|.
 \label{eq:quantitative-reconstruction-matching-QRM1}
\end{equation}
The derivative order is the total multiindex order.  The weight is
one global weight in \(d\) variables.  Neither a separate weight
in every argument nor a derivative order separately bounded in every
coordinate is being substituted for it.  These choices matter when
\(d=4n\) and \(j\) itself grows with \(n\).

Let \(\Delta_n=\bigcup_{i<j}\{x_i=x_j\}\), and let
\(\mathcal S_{\Delta,n}\) consist of the functions in
\(\mathcal S(\mathbb R^{4n})\) whose every partial derivative
vanishes on \(\Delta_n\).  It has the induced Schwartz topology.
It is closed: evaluation of a fixed derivative at a fixed point of
\(\Delta_n\) is a continuous functional, and the space is the
intersection of their kernels.  No quotient by contact
distributions has been chosen in defining the test space.  Rather,
such distributions are invisible when a functional is restricted to
this space.

For comparison with derivative-sum conventions, put
\[
 s_j^{(d)}(f)=\sum_{|\alpha|\le j}\sup_x
       (1+|x|^2)^{j/2}|\partial^\alpha f(x)|.
\]
The exact count of multiindices is
\(\binom{d+j}{j}\).  One proof appends a nonnegative slack
coordinate so that \(|\alpha|\le j\) becomes a composition of
\(j\) into \(d+1\) parts; the usual placement of separators
then counts the compositions.  Consequently
\begin{equation}
 p_j^{(d)}(f)\le s_j^{(d)}(f)
       \le\binom{d+j}{j}p_j^{(d)}(f)
       \le2^{d+j}p_j^{(d)}(f).
 \label{eq:quantitative-reconstruction-matching-QRM2}
\end{equation}
At \(d=4n\), \(j=20n\), the cost is at most \(2^{24n}\).
Thus a switch between a maximum and a sum costs one exponential in
\(n\), not an unspecified equivalence constant depending on the
dimension.  The total derivative order remains \(20n\).

The global weight can also be compared with a coordinate maximum
weight, but that comparison has a different numerical cost.  Since
\(|x|_\infty\le|x|_2\le\sqrt d\,|x|_\infty\), the reverse
weight comparison costs \(d^{j/2}\).  For the displayed values
this is \((4n)^{10n}\), which is still factorial growth.  Indeed
\[
 \log(n!)=\sum_{r=1}^n\log r
       \ge\int_1^n\log t\,dt=n\log n-n+1
\]
implies \(n^n\le e^n n!\), and hence
\begin{equation}
 (4n)^{10n}\le(4e)^{10n}(n!)^{10}.
 \label{eq:quantitative-reconstruction-matching-QRM3}
\end{equation}
This estimate explains why a change of seminorm convention must
record its dimension cost.  We keep \eqref{eq:quantitative-reconstruction-matching-QRM1} for the proved extension
and give the following explicit comparison for the theorem interface.

\begin{proposition}[The quantitative seminorm ledger]
\label{thm:quantitative-os-seminorm-ledger}
For \(f\in\mathcal S(\mathbb R^d)\) and an integer \(r\ge0\), put
\begin{equation}
 \begin{gathered}
 q_{{\rm tot},r}(f)=\max_{|\alpha|+|\beta|\le r}
                      \sup_x|x^\alpha\partial^\beta f(x)|,\\
 q_{{\rm sep},r}(f)=\max_{|\alpha|\le r,\ |\beta|\le r}
                      \sup_x|x^\alpha\partial^\beta f(x)|.
 \end{gathered}
 \label{eq:quantitative-reconstruction-matching-QRM22}
\end{equation}
Then
\begin{equation}
 \begin{gathered}
 q_{{\rm tot},r}\le q_{{\rm sep},r}\le p_r,
 \qquad p_r\le(1+d)^r q_{{\rm sep},r},\\
 p_r\le(1+d)^r q_{{\rm tot},2r}.
 \end{gathered}
 \label{eq:quantitative-reconstruction-matching-QRM23}
\end{equation}
In particular, in dimension \(4n\),
\begin{equation}
 p_{20n}(f)\le[(5e)^{20}]^n(n!)^{20}q_{{\rm tot},40n}(f).
 \label{eq:quantitative-reconstruction-matching-QRM24}
\end{equation}
\end{proposition}

\emph{Proof.} A monomial of degree at most \(r\) is bounded by
\((1+|x|^2)^{r/2}\), which proves the first two comparisons.
For the reverse comparison use
\(\sqrt{1+|x|^2}\le1+\sum_{i=1}^d|x_i|\), and expand its
\(r\)-th power.  All monomials have total degree at most \(r\),
and the sum of their nonnegative coefficients is exactly
\((1+d)^r\).  This proves the separate-order bound.  The constraints
\(|\alpha|,|\beta|\le r\) imply
\(|\alpha|+|\beta|\le2r\), giving the last line.  Finally
\(1+4n\le5n\), and the elementary factorial estimate preceding
\eqref{eq:quantitative-reconstruction-matching-QRM3} gives \((5n)^{20n}\le[(5e)^{20}]^n(n!)^{20}\).
\hfill\(\square\)

The comparison distinguishes total combined order from separate
monomial and derivative orders.  It also proves a sufficient safe
match with a doubled linear order, independently of which of these
explicit seminorm conventions is used to state the standard theorem.
No appeal to a dimension-dependent topological equivalence is needed.

\subsubsection*{The exact factorial input}

For the actual scalar moments the quantitative extension has already
proved, with one fixed finite \(k\),
\begin{equation}
 S_0=1,\qquad
 |S_n(f)|\le A B^n(n!)^{k+32}p_{20n}^{(4n)}(f),
       \quad f\in\mathcal S_{\Delta,n}.
 \label{eq:quantitative-reconstruction-matching-QRM4}
\end{equation}
The coefficient \(20\) controls every order at once.  This is
stronger information than continuity at each fixed \(n\).  The
constant \(A\) can be replaced by \(\max(1,A)\), and similarly
put \(B_+=\max(1,B)\).  For \(n\ge1\), the factors
\(2,3,\ldots,n\) give \(n!\ge2^{n-1}\).  Therefore
\begin{equation}
 B^n\le B_+(n!)^{\log_2B_+},\qquad
 \sigma_n=A B_+(n!)^{k+32+\lceil\log_2B_+\rceil}
 \label{eq:quantitative-reconstruction-matching-QRM5}
\end{equation}
is one factorial-growth majorant, and
\(|S_n(f)|\le\sigma_n p_{20n}(f)\).
If the derivative-sum convention is used, its harmless exponential
cost in \eqref{eq:quantitative-reconstruction-matching-QRM2} is absorbed by replacing \(B\) by \(2^{24}B\)
before the same argument.  Moreover the seminorm comparison gives
\begin{equation}
 |S_n(f)|\le A\widetilde B^n(n!)^{k+52}
                         q_{{\rm tot},40n}(f),
 \qquad \widetilde B=B(5e)^{20}.
 \label{eq:quantitative-reconstruction-matching-QRM25}
\end{equation}
Absorbing \(\widetilde B^n\) as in \eqref{eq:quantitative-reconstruction-matching-QRM5} gives a factorial
majorant at the one fixed order \(40n\).  Since
\(q_{{\rm tot},40n}\le p_{40n}\), this also gives the weighted
version at that order.  Thus both the direct weighted route and the
safe total-monomial route meet a fixed linear-order hypothesis.
The increase by twenty in the factorial exponent is explicit, and
no order proportional to \(n^2\) has entered this Euclidean input.

The extension in \eqref{eq:quantitative-reconstruction-matching-QRM4} acts on precisely the coincidence-flat
Schwartz space.  A stronger assertion on unrestricted Euclidean
tests is not needed for the corrected linear-growth hypothesis.
In particular there is no step that assigns a distribution supported
on \(\Delta_n\), then invokes the theorem as if that assignment
had been fixed by the separated data.

It is useful to distinguish three estimates at this point.  The formula in \eqref{eq:quantitative-reconstruction-matching-QRM4} is a proved implication of the declared source assumptions.
The energy-band construction gives a possibly larger, quadratic in
\(n\), order for unrestricted Minkowski moments.  The corrected
reconstruction theorem below has its own stronger linear-order
Minkowski conclusion.  We will obtain that conclusion for the
constructed moments only after identifying the two constructions.
The identification does not use that stronger bound in advance.

\subsubsection*{The reflected products really belong to the flat domain}

Let \(\mathcal S_{\Delta,+,n}\) be the flat functions whose
jets vanish whenever some time coordinate is nonpositive.  Set
\(\theta(\mathbf x,t)=(\mathbf x,-t)\).  For
\(F_m\in\mathcal S_{\Delta,+,m}\) and
\(G_n\in\mathcal S_{\Delta,+,n}\), consider
\begin{equation}
 (\Theta\overline{F_m}\otimes G_n)(x_1,\ldots,x_{m+n})
   =\overline{F_m(\theta x_1,\ldots,\theta x_m)}
                         G_n(x_{m+1},\ldots,x_{m+n}).
 \label{eq:quantitative-reconstruction-matching-QRM6}
\end{equation}
Every derivative of this product vanishes on a coincidence within
the first group, because the corresponding derivative of \(F_m\)
does; the same applies within the second group.  At a cross-group
coincidence, simultaneous support of the two factors requires its
time coordinate to be both nonpositive and nonnegative.  It is
therefore zero.  At that time boundary every jet of at least one
factor vanishes.  Leibniz's rule proves vanishing of all derivatives
of the product there as well.  Thus \eqref{eq:quantitative-reconstruction-matching-QRM6} belongs to
\(\mathcal S_{\Delta,m+n}\), including the case in which one
degree is zero.

Tensor products and reflection are continuous for the stated
seminorms.  For example, for each fixed \(j,m,n\),
\begin{equation}
 p_j^{(4(m+n))}(\Theta\overline F\otimes G)
      \le2^{j/2}p_j^{(4m)}(F)p_j^{(4n)}(G).
 \label{eq:quantitative-reconstruction-matching-QRM7}
\end{equation}
The estimate follows by distributing a derivative between the two
groups and using
\(1+|x|^2+|y|^2\le2(1+|x|^2)(1+|y|^2)\).
Reflection changes no derivative magnitude or Euclidean length.
This proves joint continuity on any fixed pair of degrees.
Consequently the full finite-sequence reflected form
\begin{equation}
 q(F,G)=\sum_{m,n}S_{m+n}
                 (\Theta\overline{F_m}\otimes G_n)
 \label{eq:quantitative-reconstruction-matching-QRM8}
\end{equation}
is well-defined.  Its positivity is the actual full-sequence
positivity proved by the constrained density argument.  The sum is
finite; no completion in the number of arguments is being assumed.

The historical ordered space consists of the Schwartz functions
whose jets vanish outside \(0<t_1<\cdots<t_n\).  It is a subspace
of \(\mathcal S_{\Delta,+,n}\): a coincidence forces equal
times and hence lies outside the strict ordered region.  The
historical reflected product reverses the first argument group.
Scalar permutation symmetry makes its value under \(S_{m+n}\)
equal to \eqref{eq:quantitative-reconstruction-matching-QRM8}, by the permutation reversing just those \(m\)
slots.  This establishes the positivity convention required by the
imported theorem on its actual test space.

The isometry between ordered quotients follows directly.  Two
ordered sequences with zero difference in \eqref{eq:quantitative-reconstruction-matching-QRM8} have the same
image in the full quotient.  Conversely a zero image has zero
ordered norm because the two forms agree.  Complete this isometry
and denote its closed range by \(\mathcal K_{\rm ord}\).
The resulting Hilbert space is the same one on which the band fields
were shown cyclic.  An equality with the larger full positive-time
space is not part of this argument.

\subsubsection*{Symmetry, reality and the clustering clause}

Scalar permutation symmetry and reality are retained properties of
the actual Euclidean functional.  On separated tensors these are the
identities for products of the same real scalar observable; the
quantitative flat extension and constrained tensor density transfer
them to the entire flat test space.  G2 transfers in the same way.
Thus the theorem is applied to the chosen functional itself.  No
invariance of a component is inferred solely from invariance of its
barycenter.

For a clustering component, the compact-tuple cluster estimate gives
on its ordered compact core
\begin{equation}
 \langle[F],U_a[G]\rangle
       \longrightarrow
       \langle[F],\Omega\rangle\langle\Omega,[G]\rangle
                    \quad(|a|\longrightarrow\infty),
 \label{eq:quantitative-reconstruction-matching-QRM9}
\end{equation}
where \(a\) is spatial.  To extend this to any two ordered
Hilbert vectors \(\xi,\eta\), choose compact-core approximants
\(\xi_l,\eta_l\).  Unitarity bounds the matrix-element error,
uniformly in \(a\), by
\begin{equation}
 \|\xi-\xi_l\|\,\|\eta\|
       +\|\xi_l\|\,\|\eta-\eta_l\|.
 \label{eq:quantitative-reconstruction-matching-QRM10}
\end{equation}
The corresponding rank-one vacuum expression has the same bound,
since \(\|\Omega\|=1\).  First take the far-translation limit
at fixed \(l\), and then let \(l\) grow.  This proves \eqref{eq:quantitative-reconstruction-matching-QRM9}
on the whole sector, hence in particular on the historical ordered
Schwartz sequences.  It proves the exact cluster clause without
assuming translation-uniform Schwartz bounds for approximate tests.

Several components need not satisfy \eqref{eq:quantitative-reconstruction-matching-QRM9} after mixing.  The
cluster-free part of reconstruction still applies to their positive
mixture, whereas a unique-vacuum conclusion must retain the component
or rank-one cluster assumption.  A polynomial decay estimate on
compact tuples does not become the same numerical rate on all Hilbert
vectors by the preceding density proof.

\subsubsection*{One verified reconstruction import}

We use the corrected scalar theorem \(E'\!\to R'\) of
\cite[Section IV.1, pp.~287--289]{OS75}.  In the notation fixed
above, with the exact input match in
\cref{thm:quantitative-os-seminorm-ledger} and \eqref{eq:quantitative-reconstruction-matching-QRM25}, its
import is the following.  Normalized coincidence-flat
moments with a factorial majorant at one linear order, scalar
Euclidean covariance, ordered reflection positivity, reality and
permutation symmetry reconstruct positive local tempered scalar
moments with the prescribed ordered Euclidean restrictions, the
forward spectrum and relativistic covariance.  They obey
\begin{equation}
 |\widetilde W_n(F)|\le a b^{n^2}p_{wn}^{(4n)}(F)
       \quad(F\in\mathcal S(\mathbb R^{4n}))
 \label{eq:quantitative-reconstruction-matching-QRM11}
\end{equation}
for some fixed \(a,b,w\).  Adding the cluster clause gives the
clustering conclusion.  The axiom-subset form without clustering is
the clause specified in \cite[Introduction, Remark~6]{OS75}, with
\cite[Section~3, Remark~3(b)]{OS73}.

This is one standard theorem import.  Its locality conclusion
includes the permutation and complex-continuation step.  That step
is not an additional gauge-model estimate, nor is its
proof asserted to follow from the band estimate.  The previous
sections have matched every model-facing input and its domain.

We specify the time convention before comparing moments.  If the
imported formulation uses a negative-imaginary-time tube, let
\(D(t,x)=(-t,x)\).  Transport its moments by this real linear
involution.  The transported field has time action \(e^{+itH}\)
when the original one uses \(e^{-itH}\), and its Euclidean slice is
\((i\tau,r)\).  Positivity is unchanged, because pullback by
\(D\) commutes with reversal and complex conjugation.  Locality is
unchanged, because \(D\) preserves the spacelike relation.  Lorentz
covariance is transported by \(\Lambda\mapsto D\Lambda D\),
which again belongs to \(\mathrm{SO}_+(1,3)\).  Finally
\(p_j(F\circ D)=p_j(F)\), so \eqref{eq:quantitative-reconstruction-matching-QRM11} is unchanged.  Equivalently,
if \(\widetilde E_n(z;r)\) is the ordered-gap function, the two
conventions have boundaries
\begin{equation}
 \widetilde W_n^{(-)}: z_j=\epsilon_j+it_j,
 \qquad
 \widetilde W_n^{(+)}: z_j=\epsilon_j-it_j.
 \label{eq:quantitative-reconstruction-matching-QRM12}
\end{equation}
In the rest of this section \(\widetilde W\) denotes the second,
which is the convention of the constructed band fields.  This is a
coordinate convention, not a change of the Euclidean data.

\subsubsection*{Identifying the local moments without assuming locality}

Let \(E_n(z;r)\) be the point-insertion ordered-gap function from
the band construction, with the first Euclidean point fixed at zero.
Let \(\widetilde E_n(z;r)\) be the ordered-gap function of the
import.  Both are jointly holomorphic for all
\(\operatorname{Re}z_j>0\), and continuous in real spatial gaps
\(r\).  Their Euclidean restrictions represent the same actual
moment functional on compact strictly ordered configurations.  The
sewing proof gives that assertion for \(E_n\); the prescribed
Euclidean restriction is part of the reconstruction import for
\(\widetilde E_n\).

We verify equality of the densities.  On any compact ordered neighbourhood
the actual moment is represented by the continuous band density.
The imported function represents the same functional there.  Their
difference integrates to zero against every compact smooth test.
If it were nonzero at a point, continuity would give a smaller ball
and a fixed complex phase on which its real part had a strict sign;
a nonnegative test in that ball would contradict the zero integral.
Thus the two continuous densities agree pointwise.  After common
center integration this gives
\begin{equation}
 E_n(\tau;r)=\widetilde E_n(\tau;r)
        \quad\hbox{for every }\tau_j>0\hbox{ and real }r.
 \label{eq:quantitative-reconstruction-matching-QRM13}
\end{equation}

Fix \(r\).  For fixed positive real values of every variable except
\(z_1\), the difference is a holomorphic function of \(z_1\) in
the right half-plane and vanishes on its positive real interval.
The one-variable identity theorem makes it zero for every such
\(z_1\).  Keep this first variable at an arbitrary complex value
and apply the same argument to \(z_2\).  Continue through all the
gap variables.  It follows that
\begin{equation}
 E_n(z;r)=\widetilde E_n(z;r)
       \quad\hbox{whenever every }\operatorname{Re}z_j>0.
 \label{eq:quantitative-reconstruction-matching-QRM14}
\end{equation}
This induction uses functions already jointly holomorphic on one
product domain.  It does not replace a joint boundary
statement by separate tempered limits.

At every strictly positive regulator, \eqref{eq:quantitative-reconstruction-matching-QRM14} identifies the two
integrals against arbitrary spacetime Schwartz tests.  The band
integrals are well-defined and have their regulator-uniform seminorm
bound; the import has its actual tempered boundary in the convention in \eqref{eq:quantitative-reconstruction-matching-QRM12}.  Let all regulators tend to zero jointly.  The independently
proved band boundary therefore gives
\begin{equation}
 W_n(F)=\widetilde W_n(F)
       \quad\hbox{for every }F\in\mathcal S(\mathbb R^{4n}).
 \label{eq:quantitative-reconstruction-matching-QRM15}
\end{equation}
For \(n=1\) this is the same equality of constant vacuum means;
for \(n=0\) it is normalization.  The argument has never used
locality of \(W\).  Uniqueness among already local theories alone
would not have justified \eqref{eq:quantitative-reconstruction-matching-QRM15}; equality of the actual analytic
functions and of their unrestricted boundaries does justify it.

In particular \eqref{eq:quantitative-reconstruction-matching-QRM11} now holds for the constructed moments.  This
is an imported improvement of their growth bound, obtained after
identification.  Their original energy-band proof required only its
own explicit fixed-order estimates and was not circularly based on
this improvement.

\subsubsection*{Distributional locality and the full common domain}

For adjacent slots \(j,j+1\), let \(\pi_j\) swap those slots.
The imported locality and \eqref{eq:quantitative-reconstruction-matching-QRM15} say exactly
\begin{equation}
 \left.(W_n-W_n\circ\pi_j)\right|_{
     \{(x_j-x_{j+1})^2<0\}}=0.
 \label{eq:quantitative-reconstruction-matching-QRM16}
\end{equation}
This is a distributional identity on that open subset of all
\(4n\) real coordinates; spectator arguments are unrestricted.
It is not merely a commutation identity for a two-point function.
It applies first to compact smooth tests in the open subset.
Infinity cutoffs then extend it to a Schwartz test supported there:
the cutoffs converge in every Schwartz seminorm, and each cutoff
has compact support contained in the same open subset.

Take compact smooth \(g,h\) with spacelike separated supports.
Let \(\xi,\eta\) be any two finite field-polynomial vacuum
vectors.  The adjoint identity on the common domain writes
\(\langle\xi,[\phi(g),\phi(h)]\eta\rangle\) as a finite sum
of the expressions in \eqref{eq:quantitative-reconstruction-matching-QRM16}, with the reversed conjugate bra
labels before the adjacent \(g,h\) pair and arbitrary ket labels
after it.  Each is zero.  The polynomial bra vectors are dense in
\(\mathcal K_{\rm ord}\), and the commutator ket is an actual
Hilbert vector by invariance of the common domain.  Hence the
commutator is zero on every polynomial ket.

There is a further domain issue before asserting the same identity
on \(\mathcal D_\infty\).  Let \(W=1+H\) and retain the proved
energy exponent \(N=9\).  Applying the band bound first with
\(k=0\), then with \(k=N\), gives
\begin{equation}
 \begin{split}
 \|\phi(g)\phi(h)\eta\|
   &\le C_0C_Np_{2N+7}(g)p_{2N+7}(h)
                                      \|W^{2N}\eta\|,\\
 \|[\phi(g),\phi(h)]\eta\|
   &\le2C_0C_Np_{2N+7}(g)p_{2N+7}(h)
                                      \|W^{2N}\eta\|.
 \end{split}
 \label{eq:quantitative-reconstruction-matching-QRM17}
\end{equation}
Here \(\eta\in\operatorname{Dom}W^{2N}\): the first field
maps that domain into \(\operatorname{Dom}W^N\), where the
second is defined.  Thus both products are actual continuous maps
from this graph domain to the Hilbert space.  Polynomial vectors
are dense in that graph domain on the ordered sector, by the
spectral-smoothing argument and cyclic identification already
proved.  Approximate \(\eta\) there and use \eqref{eq:quantitative-reconstruction-matching-QRM17}.  We obtain
\begin{equation}
 [\phi(g),\phi(h)]\eta=0
 \quad\left(\eta\in\operatorname{Dom}W^{2N}
                            \cap\mathcal K_{\rm ord}\right),
 \label{eq:quantitative-reconstruction-matching-QRM18}
\end{equation}
in particular on the invariant common domain.

The same statement holds for Schwartz labels with spacelike
separated supports.  Multiply each label by smooth infinity
cutoffs; their supports remain spacelike separated, and they
converge in every seminorm appearing in \eqref{eq:quantitative-reconstruction-matching-QRM17}.  The two product
maps therefore converge in the stated graph-to-Hilbert norm.
This proves the extension without imposing a uniform spacelike
distance at infinity.

\eqref{eq:quantitative-reconstruction-matching-QRM18} is the local commutativity of the constructed
unbounded Wightman field.  It does not say that arbitrarily selected
self-adjoint extensions have commuting spectral projections.  Real
test fields have been proved symmetric and closable, while essential
self-adjointness and strong spectral commutation require their own
operator criteria.  No exponentials of such unspecified extensions
are needed for \eqref{eq:quantitative-reconstruction-matching-QRM18}.

\subsubsection*{The stronger vector and energy seminorm conclusions}

The scalar bound in \eqref{eq:quantitative-reconstruction-matching-QRM11} also improves the vacuum-product test-space
estimate.  For a general \(F\in\mathcal S(\mathbb R^{4n})\),
let \(X_n(F)\) be its actual Hilbert-valued vacuum product.  On
finite tensors positivity and the adjoint identity give
\(\|X_n(F)\|^2=W_{2n}(F^*\otimes F)\).  Both sides are
continuous under the already proved unrestricted extensions, so
this equality holds for every \(F\).  Tensor seminorm comparison
then gives
\begin{equation}
 \|X_n(F)\|\le A_n p_{2wn}^{(4n)}(F),
 \qquad A_n=\sqrt a\,b^{2n^2}2^{wn/2}.
 \label{eq:quantitative-reconstruction-matching-QRM19}
\end{equation}
One may increase \(a,b,w\) to make them at least one.  The harmless
power of two comes from \eqref{eq:quantitative-reconstruction-matching-QRM7}; the specific larger constant displayed
is sufficient.  Reversal and conjugation do not change the seminorm.
This estimate is compatible with the actual band-vector extension,
because their values on the dense tensor space already agree.

There is an explicit energy version of \eqref{eq:quantitative-reconstruction-matching-QRM19}.  Set
\(D_t=\sum_{j=1}^n\partial_{t_j}\).  Translation covariance and
continuity of the vector-valued test map imply
\begin{equation}
 H^lX_n(F)=i^lX_n(D_t^lF),\qquad
 p_j(D_t^lF)\le n^lp_{j+l}(F).
 \label{eq:quantitative-reconstruction-matching-QRM20}
\end{equation}
For the first identity, the simultaneous test translation is
\(F(t_1-s,\ldots,t_n-s)\); its derivative is \(-D_tF\),
while the vector translation is \(e^{isH}\).  Difference quotients
converge in the Schwartz topology and therefore in the Hilbert
space.  Stone's closed generator criterion proves the domain
identity for \(l=1\); induction proves every \(l\).  The second
estimate is the multinomial expansion, whose absolute coefficients
sum to \(n^l\), followed by the monotonicity of the weights.
Consequently
\begin{equation}
 \|W^lX_n(F)\|
       \le A_n(1+n)^lp_{2wn+l}(F).
 \label{eq:quantitative-reconstruction-matching-QRM21}
\end{equation}
This proves continuity of the vacuum-product distribution into every
energy graph domain with a linear test order in \(n\) at fixed
\(l\).  It does not give an essential-self-adjointness conclusion:
such a conclusion would need control of repeated products of the
same unbounded field on a suitable set of analytic vectors, which
is a different estimate.

\subsubsection*{The exact conditional conclusion}

\begin{proposition}[The scalar reconstruction on its actual cyclic sector]
\label{thm:scalar-reconstruction-on-ordered-sector}
Suppose the actual scalar moments satisfy \eqref{eq:quantitative-reconstruction-matching-QRM4}, the full positive
form in \eqref{eq:quantitative-reconstruction-matching-QRM8}, reality, permutation symmetry, G2 and the UI/PD insertion
hypotheses.  The energy-band field on
\(\mathcal D_\infty\cap\mathcal K_{\rm ord}\) is Hermitian,
tempered, Poincare covariant, local and cyclic, with the full forward
spectrum.  Its Euclidean correlations are exactly the original
diagonal-flat moments in their declared sense.  The conclusion uses
one verified corrected reconstruction theorem import for the
permutation and locality step and its stronger output-order bound.
For a component satisfying \eqref{eq:quantitative-reconstruction-matching-QRM9}, the corresponding cluster
conclusion holds.  No rank-one cluster assertion is required for
the cluster-free local conclusion.
\end{proposition}

\emph{Proof.} The field, domain, sewing, spectral and joint-boundary
arguments were proved before the import.  \eqref{eq:quantitative-reconstruction-matching-QRM13}--\eqref{eq:quantitative-reconstruction-matching-QRM15}
identify its moments; \eqref{eq:quantitative-reconstruction-matching-QRM16}--\eqref{eq:quantitative-reconstruction-matching-QRM18} prove locality on the actual
energy domain.  The remaining assertions are the already established
covariance, positivity and cyclicity.\hfill\(\square\)

The assertion concerns this scalar cyclic sector. Gauge-theory
existence, full observable identification and a positive mass gap
require separate arguments.


% END MODULE chapters/quantitative_reconstruction_matching.tex

% BEGIN MODULE chapters/reconstruction_testspace_estimates.tex
\subsection{Unrestricted test functions and joint boundary estimates}
\label{ch:reconstruction-testspace-estimates}

This section proves the test-space statements used in constructing the
band fields.  Its estimates apply to general spacetime tests and to
Hilbert-valued vacuum products.  They also justify simultaneous
regulator limits, the passage to difference variables and the spectral
support in every cut.  The arguments use explicit partitions and
Fourier coefficients.  No unrestricted Euclidean contact extension and
no unquantified tensor-product assertion enter the proof.

Write \(\mathcal S_d=\mathcal S(\mathbb R^d)\) and use
\(p_j^{(d)}\) from \eqref{eq:quantitative-reconstruction-matching-QRM1}.  Constants indexed by fixed integers
\(n,j,r\) below are finite constants for those integers.  A bound
at every fixed field order is sufficient for the unrestricted
temperedness constructed here.  The stronger linear-order conclusion
of the corrected reconstruction import has already been distinguished
from these constructive estimates.

\subsubsection*{A product partition with controlled support}

Choose \(b\in C_c^\infty(\mathbb R^4)\), nonnegative, supported
in \((-1,1)^4\) and positive on \([-3/4,3/4]^4\).  The function
\(d(x)=\sum_{z\in\mathbb Z^4}b(x-z)\) is positive, smooth and
periodic.  Only finitely many terms meet any point; the number is
uniformly bounded.  Positivity on one compact fundamental cell gives
a strictly positive lower bound for \(d\).  Therefore
\begin{equation}
 \chi_z(x)=\frac{b(x-z)}{d(x)},\qquad
 \sum_z\chi_z(x)=1,
 \qquad
 \sup_z\|\partial^\alpha\chi_z\|_\infty<\infty.
 \label{eq:reconstruction-testspace-estimates-RTE1}
\end{equation}
The derivative bound follows by repeatedly differentiating the
reciprocal of the smooth periodic function with its positive lower
bound.  Every \(\chi_z\) is supported in \(z+[-1,1]^4\).
For a fixed degree \(n\), use the product partition
\(\chi_Z(X)=\prod_{i=1}^n\chi_{z_i}(x_i)\), where
\(Z=(z_1,\ldots,z_n)\).  It still sums to one and is locally
finite.  A derivative of prescribed total order has a bound depending
only on that order and \(n\), with no dependence on the locations.

Fix also \(\eta\in C_c^\infty((-3,3)^4)\) equal to one on
\([-1,1]^4\).  Put \(\kappa=2\pi/8\).  On the period-eight
box in each slot, expand
\(\chi_Z(X)F(X)\) in its ordinary smooth Fourier series.  Its
coefficients are
\begin{equation}
 c_{Z,L}(F)=8^{-4n}\int_{[-4,4]^{4n}}
    \chi_Z(Z+Y)F(Z+Y)
                    e^{-i\kappa\sum_i l_i\cdot y_i}\,dY,
 \label{eq:reconstruction-testspace-estimates-RTE2}
\end{equation}
where \(L=(l_1,\ldots,l_n)\in(\mathbb Z^4)^n\).
The localized function vanishes in a neighbourhood of the periodic
boundary.  Its periodic extension is consequently smooth with all
derivatives matching.  There is no boundary term when derivatives are
integrated off it.

Define one-slot waves
\begin{equation}
 v_{z,l}(x)=\eta(x-z)e^{i\kappa l\cdot(x-z)}.
 \label{eq:reconstruction-testspace-estimates-RTE3}
\end{equation}
Multiplying the periodic Fourier series by the larger cutoffs gives
the global identity
\(\chi_ZF=\sum_Lc_{Z,L}\bigotimes_i v_{z_i,l_i}\).
Indeed on the support of \(\chi_ZF\) the cutoffs are one; off
that support inside the periodic box the represented periodic
function is zero; outside the larger cutoffs both sides are zero.
The reason for using two nested cutoffs is precisely to obtain this
global identity without leaving noncompact plane waves.

\subsubsection*{The coefficient estimate and the exact order budget}

On the support of \(\chi_Z\), the distances between \(z_i\)
and \(x_i\) are bounded by two in Euclidean norm.  Thus their
polynomial weights of any fixed order are comparable, uniformly in
the centers.  Also
\begin{equation}
 \prod_{i=1}^n(1+|x_i|^2)^{1/n}
          \le1+\sum_{i=1}^n|x_i|^2.
 \label{eq:reconstruction-testspace-estimates-RTE4}
\end{equation}
Each factor on the left is at most the common right-hand side before
taking the product and the \(n\)-th root.  Apply
\(\prod_i(1-\kappa^{-2}\Delta_{y_i})^r\) in \eqref{eq:reconstruction-testspace-estimates-RTE2}, where
\(r\) is a positive integer.  Its Fourier multiplier is
\(\prod_i(1+|l_i|^2)^r\).  The derivative order used in all
slots together is at most \(2rn\).  Leibniz's rule, \eqref{eq:reconstruction-testspace-estimates-RTE1} and
\eqref{eq:reconstruction-testspace-estimates-RTE4} therefore give, whenever \(J\ge2rn\),
\begin{equation}
 |c_{Z,L}(F)|\le C_{n,r,J}p_J^{(4n)}(F)
    \prod_{i=1}^n(1+|z_i|)^{-J/n}
                       (1+|l_i|^2)^{-r}.
 \label{eq:reconstruction-testspace-estimates-RTE5}
\end{equation}
The integration boxes have fixed volume, so no location-dependent
volume factor is present.  Derivatives of the fixed partition cost
only \(C_{n,r,J}\).  The single global weight gives decay \(J/n\) in each slot.

For every nonnegative integer \(M\), differentiating \eqref{eq:reconstruction-testspace-estimates-RTE3} gives
\begin{equation}
 p_M^{(4)}(v_{z,l})
         \le C_M(1+|z|)^M(1+|l|)^M.
 \label{eq:reconstruction-testspace-estimates-RTE6}
\end{equation}
Each derivative differentiates either the fixed cutoff or its
exponential.  The latter contributes a power of \(\kappa l\),
and the former has a fixed derivative bound.  On its support
\(1+|x|\le C(1+|z|)\), proving the weight factor.  Since there
are only finitely many derivative splittings up to order \(M\),
their sum is absorbed in \(C_M\).

Choose
\begin{equation}
 2r=2\left\lceil\frac{M+6}{2}\right\rceil,
 \qquad K=n(M+8).
 \label{eq:reconstruction-testspace-estimates-RTE7}
\end{equation}
Then \(M+6\le2r\le M+7\), so \(2rn\le K\).
Combining \eqref{eq:reconstruction-testspace-estimates-RTE5}--\eqref{eq:reconstruction-testspace-estimates-RTE7} yields
\begin{equation}
 \sum_{Z,L}|c_{Z,L}(F)|
        \prod_{i=1}^np_M^{(4)}(v_{z_i,l_i})
                     \le C_{n,M}p_K^{(4n)}(F).
 \label{eq:reconstruction-testspace-estimates-RTE8}
\end{equation}
Here the sums to be checked are exactly
\(\sum_z(1+|z|)^{-8}\) and
\(\sum_l(1+|l|)^M(1+|l|^2)^{-r}\), in four dimensions.
For completeness, the shell \(|z|_\infty=k\), \(k\ge1\),
has \((2k+1)^4-(2k-1)^4\le C(k+1)^3\) points.  Euclidean
and maximum norms in this fixed dimension differ by at most two.
The first shell sum is dominated by \(\sum(k+1)^{-5}\).
The second is dominated by \(\sum(k+1)^{3+M-2r}\), whose
exponent is at most \(-3\).  Both converge.  Taking their
\(n\)-fold products proves \eqref{eq:reconstruction-testspace-estimates-RTE8} with a finite constant at the
specified field order.

\subsubsection*{Actual Schwartz convergence of the tensor series}

The identity used in \eqref{eq:reconstruction-testspace-estimates-RTE8} also converges in the full Schwartz
topology, not merely in the one displayed seminorm.  Fix an arbitrary
target order \(j\).  A tensor in independent slots satisfies
\begin{equation}
 p_j^{(4n)}\left(\bigotimes_i f_i\right)
                         \le\prod_i p_j^{(4)}(f_i),
 \label{eq:reconstruction-testspace-estimates-RTE9}
\end{equation}
because a total derivative splits between the slots and
\(1+\sum|x_i|^2\le\prod_i(1+|x_i|^2)\).
Repeat \eqref{eq:reconstruction-testspace-estimates-RTE5} with \(2r\ge j+6\), and choose
\(J\ge\max\{2rn,n(j+8)\}\).  \eqref{eq:reconstruction-testspace-estimates-RTE6} and \eqref{eq:reconstruction-testspace-estimates-RTE9}
now make the tensor series absolutely summable in \(p_j^{(4n)}\).
Since \(F\) is Schwartz, its \(p_J\) in this argument is finite
for every such choice.  This proves absolute summability in every
target seminorm.

The Schwartz space is complete for this family: a Cauchy sequence
converges with each derivative uniformly on compact sets; the limits
are successive derivatives by the fundamental theorem of calculus,
and the weighted Cauchy bounds put the resulting function in every
seminorm.  Therefore the tensor series has a Schwartz sum.
Its pointwise sum is \(F\), by the locally finite partition and
the smooth periodic Fourier identities.  Its Schwartz sum is hence
also \(F\).  Finite sums of tensors are consequently dense in the
ordinary unrestricted Schwartz space.  Absolute convergence permits
any enumeration of the indices, with no conditional rearrangement.

There is also no algebraic tensor ambiguity.  For a finite-dimensional
space \(E\) of functions, the evaluation functionals span \(E^*\):
otherwise a nonzero vector of their annihilator would be a function
vanishing at every point.  Select a basis of evaluations and invert
its matrix on a basis of \(E\).  Applying these evaluation matrices
in each slot shows that the natural map from a finite tensor product
of such function spaces to functions on the product is injective.
Every finite tensor relation lies in one such product.  A multilinear
functional therefore has a well-defined linear value on the actual
finite tensor functions used above.

\subsubsection*{Scalar, vector and energy-graph extensions}

Let \(\mathcal E\) be a Banach space and suppose that an actual
\(n\)-linear map satisfies
\begin{equation}
 \|\mathcal L(f_1,\ldots,f_n)\|_{\mathcal E}
                           \le D\prod_i p_M^{(4)}(f_i).
 \label{eq:reconstruction-testspace-estimates-RTE10}
\end{equation}
Define its value on a general \(F\in\mathcal S_{4n}\) by
\begin{equation}
 \mathcal L(F)=\sum_{Z,L}c_{Z,L}(F)
                    \mathcal L(v_{z_1,l_1},\ldots,v_{z_n,l_n}).
 \label{eq:reconstruction-testspace-estimates-RTE11}
\end{equation}
The series is absolutely convergent in \(\mathcal E\) by \eqref{eq:reconstruction-testspace-estimates-RTE8},
and \(\|\mathcal L(F)\|\le DC_{n,M}p_{n(M+8)}(F)\).
Linearity follows from linearity of every coefficient and absolute
summation.  On a tensor \(F=\bigotimes_i f_i\), the localized
coefficients factor into their one-slot coefficients.  The one-slot
series converges in \(p_M\); expand the finite product and pass
to its absolutely convergent series using \eqref{eq:reconstruction-testspace-estimates-RTE10}.  This proves that
\eqref{eq:reconstruction-testspace-estimates-RTE11} agrees with the original multilinear value.  Ordinary tensor
density proves uniqueness of its continuous extension.  In
particular a different choice of cutoffs gives the same answer.

For the actual band fields put
\(\mathcal L(f_1,\ldots,f_n)=\phi(f_1)\cdots\phi(f_n)\Omega\).
Iteration of the weighted field bound gives, at every integer
\(l\ge0\),
\begin{equation}
 \begin{gathered}
 \|W^l\mathcal L(f_1,\ldots,f_n)\|
        \le D_{n,l}\prod_i p_{M_{n,l}}(f_i),\\
 M_{n,l}=nN+l+7,\qquad K_{n,l}=n(nN+l+15).
 \end{gathered}
 \label{eq:reconstruction-testspace-estimates-RTE12}
\end{equation}
At the first field the required right energy exponent is \(l+N\);
at the next it is \(l+2N\), and so on.  The largest temporal
test order in the band estimate is exactly \(nN+l+7\).
Apply \eqref{eq:reconstruction-testspace-estimates-RTE11} first in the Hilbert space, then in each energy graph
space with its \(\|W^l\cdot\|\) norm.  The partial tensor sums
are the same vectors.  They converge both in the Hilbert space and
after applying \(W^l\).  Closedness of \(W^l\) identifies the
two limits.  Hence the extension \(X_n(F)\) satisfies
\begin{equation}
 X_n(F)\in\mathcal D_\infty,\qquad
 \|W^lX_n(F)\|\le C_{n,l}p_{K_{n,l}}(F).
 \label{eq:reconstruction-testspace-estimates-RTE13}
\end{equation}
This constructs the vector-valued distribution in all its graph
seminorms.  Coherence between graph orders is a consequence of
closedness, not an assumed nuclear extension into an intersection
of unbounded-operator domains.

If a fixed field \(\phi(g)\) is applied to these vectors, the
energy bound makes it continuous from the graph norm of \(W^{l+N}\)
to that of \(W^l\).  Passing finite tensors to the limit proves
\begin{equation}
 \phi(g)X_n(F)=X_{n+1}(g\otimes F)
                         \quad\hbox{in every energy graph norm}.
 \label{eq:reconstruction-testspace-estimates-RTE14}
\end{equation}
This is the precise multiplication identity needed for the
unrestricted polynomial representation.  It remains valid for
non-tensor \(F\); no product of two Euclidean separated-tuple
symbols has been used to define it.

\subsubsection*{Uniform families and the topology of the boundary limit}

Suppose \(\mathcal L_\alpha\) obey \eqref{eq:reconstruction-testspace-estimates-RTE10} with common \(D,M\), and
converge to \(\mathcal L_0\) on every finite tensor function.  Their
extensions have a common \(p_K\) bound.  For fixed \(F\), choose
a finite tensor approximation in that seminorm.  The approximation
error is bounded uniformly in \(\alpha\); convergence handles
the one fixed finite approximation.  This proves
\(\mathcal L_\alpha(F)\to\mathcal L_0(F)\) for every unrestricted test.  It is
valid for simultaneous parameters and therefore does not impose an
order on their limits.

In fact the convergence is uniform on bounded Schwartz sets.
Here is an elementary compactness argument for exactly the
seminorm needed.  For a bounded set \(\mathcal B\subset\mathcal S_d\),
smooth infinity cutoffs satisfy
\begin{equation}
 \sup_{F\in\mathcal B}p_K((1-\chi_R)F)
             \le C_KR^{-1}\sup_{F\in\mathcal B}p_{K+1}(F)
                         \longrightarrow0.
 \label{eq:reconstruction-testspace-estimates-RTE15}
\end{equation}
The derivative bounds on \(\chi_R\) and the extra power of
decay prove this by Leibniz's rule.  At fixed \(R\), the cut-off
set has common compact support and uniformly bounded derivatives
through order \(K+1\).  The derivatives through order \(K\)
are uniformly equicontinuous by the mean-value formula.  Successive
Arzela--Ascoli extraction for their finitely many derivatives proves
precompactness in \(p_K\).  Combining this with \eqref{eq:reconstruction-testspace-estimates-RTE15} gives a
finite \(p_K\) net for \(\mathcal B\), with centers in the
Schwartz space.

Apply the common bound to the distance from a net center.  Convergence
at its finitely many centers then gives
\begin{equation}
 \sup_{F\in\mathcal B}\|\mathcal L_\alpha(F)-\mathcal L_0(F)\|_{\mathcal E}
                                    \longrightarrow0.
 \label{eq:reconstruction-testspace-estimates-RTE16}
\end{equation}
Thus scalar convergence is in the strong tempered-distribution
topology of uniform convergence on bounded Schwartz sets.  The
Banach-valued statement and the energy-graph version follow from the
same finite-net proof.  This strengthens the joint pointwise-in-test
boundary passage and records exactly which uniform estimate makes
the strengthening legitimate.

For the regulated field moments the intermediate operators are
\(e^{-\epsilon_jH-y_j\cdot P}\), with
\(\epsilon_j>|y_j|_2\).  The cone makes them contractions in
every energy graph norm, and the joint spectral theorem makes them
converge strongly there to identity.  The product field estimates
therefore have the same constants as without regulators.  Replacing
the finitely many intermediate contractions one at a time proves
convergence on tensors in every needed graph norm.  \eqref{eq:reconstruction-testspace-estimates-RTE15}--\eqref{eq:reconstruction-testspace-estimates-RTE16} then prove simultaneous strong distributional
convergence on unrestricted tests.  The construction does not rely
on a separately tempered Laplace transform in each gap.

\subsubsection*{Center integration with an explicit continuous right inverse}

For \(n\ge2\) let
\begin{equation}
 C_n(a,\xi)=(a,a+\xi_1,\ldots,a+\xi_1+\cdots+\xi_{n-1}),
 \qquad
 I_nF(\xi)=\int_{\mathbb R^4}F(C_n(a,\xi))\,da.
 \label{eq:reconstruction-testspace-estimates-RTE17}
\end{equation}
The linear coordinate change has determinant one.  Its inverse is
\(a=x_1\), \(\xi_j=x_{j+1}-x_j\), and
\[
 |a|^2+\sum_j|\xi_j|^2\le5\sum_i|x_i|^2,
 \qquad |C_n(a,\xi)|\le n|(a,\xi)|.
\]
The first estimate follows from
\(|x_{j+1}-x_j|^2\le2|x_{j+1}|^2+2|x_j|^2\).
For the second, each cumulative sum is bounded by Cauchy--Schwarz
and then the \(n\) slot bounds are added.

A derivative in a difference coordinate differentiates at most
\(n\) original slots, with coefficients of magnitude one.
Separate the center weight of order five from the remaining weight
of order \(j\).  The first preceding comparison shows that
\begin{equation}
 p_j^{(4(n-1))}(I_nF)
       \le C_4\,5^{(j+5)/2}n^jp_{j+5}^{(4n)}(F),
 \quad
 C_4=\int_{\mathbb R^4}(1+|a|^2)^{-5/2}\,da<\infty.
 \label{eq:reconstruction-testspace-estimates-RTE18}
\end{equation}
Derivatives may be taken under the integral by the same integrable
bound.  Thus center integration is an actual continuous map of
Schwartz spaces.

Choose \(\rho\in\mathcal S_4\) with integral one and put
\begin{equation}
 J_\rho\psi(X)=\rho(x_1)
                 \psi(x_2-x_1,\ldots,x_n-x_{n-1}).
 \label{eq:reconstruction-testspace-estimates-RTE19}
\end{equation}
Then \(I_nJ_\rho\psi=\psi\).  The coordinate comparison above,
the chain rule and the separate-variable weight inequality give
\begin{equation}
 p_j^{(4n)}(J_\rho\psi)
             \le(3n)^jp_j^{(4)}(\rho)
                              p_j^{(4(n-1))}(\psi).
 \label{eq:reconstruction-testspace-estimates-RTE20}
\end{equation}
There are at most three summands when an original coordinate
derivative hits the center factor or the adjacent difference
factors; iterating gives \(3^j\).  The factor \(n^j\) pays
the global weight under \(C_n\).  These safe constants suffice
without a dimension-independent assertion about the coordinate
change.  For \(n=1\), center integration is simply the integral
of a Schwartz function and the right inverse is \(c\mapsto c\rho\).

\subsubsection*{Factoring every translation-invariant distribution}

We prove the factorization, including independence of \(\rho\).
In center coordinates simultaneous spacetime translation changes only
\(a\).  A translation-invariant distribution annihilates each
\(\partial_{a_j}\Psi\), because translation difference
quotients converge in the Schwartz topology to that derivative.
It remains to identify the kernel of center integration.

Choose a one-dimensional Schwartz function \(\rho_1\) of
integral one.  For a Schwartz function \(f(a,\xi)\), define
\(P_jf=\rho_1(a_j)\int f\,da_j\).  Each \(P_j\) is
continuous on the full Schwartz space: an extra weight of order
two gives an integrable majorant in its one integration variable,
and the remaining weights are paid from the global weight.  More
explicitly \(p_j(P_if)\le C_{j,\rho_1}p_{j+2}(f)\).
The \(P_j\) commute.

If a Schwartz function \(g(t,z)\) has integral zero in \(t\)
for every parameter \(z\), set
\(Qg(t,z)=\int_{-\infty}^t g(s,z)\,ds\).
For \(t\ge0\) it also equals \(-\int_t^\infty g(s,z)\,ds\).
In either tail \(|s|\ge|t|\).  The extra weight of order two
therefore proves
\(p_j(Qg)\le(1+\pi)p_{j+2}(g)\).
Derivatives in \(z\) pass through the integral and still have
zero integral; positive \(t\) derivatives are derivatives of
\(g\).  Thus \(Q\) is a continuous Schwartz primitive and
\(\partial_tQg=g\).

The telescoping identity
\begin{equation}
 f-P_1P_2P_3P_4f
       =\sum_{j=1}^4P_1\cdots P_{j-1}(1-P_j)f
       =\sum_{j=1}^4\partial_{a_j}\Psi_j
 \label{eq:reconstruction-testspace-estimates-RTE21}
\end{equation}
uses this primitive in the \(j\)-th variable of each summand.
The summand has zero integral in that variable, and all preceding
\(P_i\) commute with its derivative.  Every \(\Psi_j\) is
Schwartz.  If \(\int f(a,\xi)\,da=0\), the product \(P_1P_2P_3P_4f\)
vanishes; this proves that the kernel of center integration is
exactly the sum of these derivative ranges.

It follows that every translation-invariant tempered distribution
\(T\) has the unique difference distribution
\begin{equation}
 t(\psi)=T(J_\rho\psi),\qquad T(F)=t(I_nF).
 \label{eq:reconstruction-testspace-estimates-RTE22}
\end{equation}
Both maps are continuous by \eqref{eq:reconstruction-testspace-estimates-RTE18}--\eqref{eq:reconstruction-testspace-estimates-RTE20}.  Independence of
\(\rho\) follows because two right inverses differ by an element
of the just-identified kernel.  Uniqueness follows from
\(I_nJ_\rho=1\).  This proof applies to the unrestricted
Minkowski distributions.  It does not require assigning a value to
an arbitrary Euclidean contact distribution.

\subsubsection*{Finite spectral cutoffs and their actual graph limits}

Let \(Q_R=\sum_{m=0}^RP_m\) for the energy bands used earlier,
and set \(\phi_R(g)=Q_R\phi(g)Q_R\).  These are bounded
operators with finite energy range, even though the individual
band subspaces can be infinite dimensional.  Their weighted field
bounds have the same constants as the original fields, since
\(Q_R\) commutes with \(W\) and has norm one.
For \(\xi\in\mathcal D_\infty\),
\begin{equation}
 \begin{split}
 \|W^l(\phi_R(g)-\phi(g))\xi\|
 &\le\|(Q_R-1)W^l\phi(g)\xi\|\\
 &\quad+C_l\mathfrak t_{N+l+2}(g)
                        \|W^{N+l}(Q_R-1)\xi\|
                \longrightarrow0.
 \end{split}
 \label{eq:reconstruction-testspace-estimates-RTE23}
\end{equation}
Both terms tend to zero by the spectral theorem.  Iterating the
graph estimate for products and replacing one factor at a time
proves convergence of every finite cutoff vacuum product.  A
uniform graph estimate is retained at each replacement, so the
vectors on which the remaining field factors act are in the exact
domains required.  Applying the extension and uniform-family lemma
gives convergence of the scalar and vector distributions on all
unrestricted tests, including strong convergence on bounded
Schwartz sets.

At a spatial point the bounded cutoff field is
\(A_R(x)=\sum_{m,n\le R}A_{mn}(x)\).  The generators of its
translations on \(Q_R\mathcal K_{\rm ord}\) are bounded by
the spectral cone and the bounded energy range.  Its spacetime
translate is therefore a norm-continuous bounded point field.  Its
vacuum products have an ordinary difference-variable expression,
with one joint translation between every adjacent pair of fields.
No finite-dimensionality of a band is needed for this assertion.

\subsubsection*{The carrier in every difference variable}

Write a difference variable as \(\xi=(t,x)\), and use physical
momentum \(p=-P\).  The joint translation is
\(\mathscr U(\xi)=e^{i(Et-p\cdot x)}\) in joint spectral
variables.  If the test transform
\(f^+(E,p)=\int f(t,x)e^{i(Et-p\cdot x)}\,dt,dx\)
vanishes on the forward cone, then
\begin{equation}
 \int f(\xi)\mathscr U(\xi)\,d\xi
                         =f^+(H,-P)=0.
 \label{eq:reconstruction-testspace-estimates-RTE24}
\end{equation}
This follows by Fubini against scalar spectral measures; the
integrable norm bound of \(f\) permits the operator weak integral.
Insert this identity at any one of the intermediate translations
of a cutoff vacuum word.  The other difference tests and the
bounded spectator operators can be held fixed during that
integration.  Thus its Fourier distribution has its carrier in
the cone in each cut.

To obtain product support rigorously, take a compact Fourier set
outside the product of the closed cones.  Around each point choose
a product box in which one designated momentum is outside its
cone.  A finite subcover and a smooth subordinate partition reduce
an arbitrary test on that compact set to finitely many tests on
such boxes.  The product Fourier construction, with cutoffs inside
each box, approximates each by finite tensors in the test topology.
Every tensor gives zero by \eqref{eq:reconstruction-testspace-estimates-RTE24}.  Hence the complete test gives
zero.  This proves the joint support statement in all momentum variables.

The difference distributions of the cutoffs converge to the actual
difference distribution, because \eqref{eq:reconstruction-testspace-estimates-RTE22} uses the continuous fixed
right inverse \(J_\rho\).  Fourier transformation is continuous
on Schwartz space: differentiating its integral introduces powers
of the original coordinates, while multiplication by a Fourier
coordinate is integration by parts.  An extra integrable weight
gives, for instance,
\(p_j^{(d)}(\widehat f)\le C_{d,j}p_{2j+d+2}^{(d)}(f)\).
Thus convergence passes through Fourier transformation.  A test
outside a fixed closed carrier was zero at every cutoff and remains
zero in the limit.  Consequently
\begin{equation}
 \operatorname{supp}\widehat W_n^{\rm diff}
                           \subset(\overline V_+)^{n-1}.
 \label{eq:reconstruction-testspace-estimates-RTE25}
\end{equation}
The unreduced momentum statement follows by the determinant-one
center-and-difference coordinate change.  It is the corresponding
condition on every partial sum.  The center Fourier variable is
zero by translation invariance.  Nothing here replaces these
multiple carrier conditions by the spectrum of a single vector.

\subsubsection*{An explicit absolute bound before removing the cutoffs}

There is a useful stronger justification of Fubini at a strictly
interior tube point.  In a vacuum word set the endpoint band
indices to zero, since the vacuum has energy zero.  Every field
block contributes at most
\(C(1+m_j+m_{j+1})^N\), and each intermediate cone regulator
with margin \(\delta_j>0\) contributes at most
\(e^{-\delta_jm_j}\).  Since
\(1+m+n\le(1+m)(1+n)\), the absolute sum is bounded by
\begin{equation}
 C^n\prod_{j=1}^{n-1}
            \sum_{m\ge0}(1+m)^{2N}e^{-\delta_jm}<\infty.
 \label{eq:reconstruction-testspace-estimates-RTE26}
\end{equation}
There is no band-dimension factor: each path is estimated in
operator norm before its band indices are summed.  For fixed
integer \(q\), comparison with the integral of
\((1+x)^qe^{-\delta x}\) gives
\(\sum_{m\ge0}(1+m)^qe^{-\delta m}
 \le C_q(1+\delta^{-q-1})\).
For \(0<\delta\le1\), expand \((1+x)^q\) and integrate its
monomials; for \(\delta\ge1\), dominate by the fixed sum with
\(\delta=1\).  This proves the asserted bound in both ranges.

The series in \eqref{eq:reconstruction-testspace-estimates-RTE26} is uniform in all real coordinate translations.
Its tails outside a growing finite band cube tend to zero by
absolute summability.  Integrating against the absolute value of a
Schwartz spacetime test is therefore legitimate, and dominated
convergence removes the cutoffs.  This proves the fixed-regulator
identity with the actual holomorphic point function.  The bound
blows up at the boundary and is used only in the interior.  The
regulator-uniform Schwartz estimate, which is required for the
joint boundary, comes separately from the field energy bound and
graph contractions.  The joint boundary therefore follows from the regulator-uniform
estimate without interchanging the energy cutoff and a singular
boundary limit.

\subsubsection*{Ordered cutoffs in every Schwartz seminorm}

We give the constrained density proof also needed for exact sewing.
Let \(O_n=\{0<t_1<\cdots<t_n\}\), and suppose that every jet
of \(F\) vanishes outside \(O_n\).  Put
\(\ell_0=t_1\), \(\ell_j=t_{j+1}-t_j\) for \(1\le j<n\).
Choose a smooth one-variable cutoff \(c\) equal to zero on
\(( -\infty,1]\) and one on \([2,\infty)\), and set
\(C_\delta(t)=\prod_{j=0}^{n-1}c(\ell_j/\delta)\).
Derivatives of total order \(r\) cost at most
\(C_{n,r}\delta^{-r}\).

On the support of \((1-C_\delta)F\), a nonzero value of a
derivative of \(F\) can occur only with all \(\ell_j>0\)
and some \(\ell_j\le2\delta\).  Move to that hyperplane
along its normal.  The distance is at most \(2\delta\); all
jets of \(F\) on the hyperplane are zero.  Taylor's formula
with integral remainder to order \(M\), applied also to every
derivative of \(F\), gives an extra factor \(C_M\delta^M\).
The normal segment has bounded length for \(\delta\le1\),
so the global weights are uniformly comparable along it.
Leibniz's rule consequently proves, for \(M>j\),
\begin{equation}
 p_j((1-C_\delta)F)
             \le C_{n,j,M}\delta^{M-j}p_{j+M}(F).
 \label{eq:reconstruction-testspace-estimates-RTE27}
\end{equation}
The same estimate covers points where a cutoff derivative is
nonzero, because they lie in one of these boundary strips.  A
finite sum over the possible hyperplanes costs only a constant at
fixed \(n,j,M\).  Thus \(C_\delta F\to F\) in every
Schwartz seminorm.

After cutting off infinity, these tests have compact support a
positive distance from the ordered boundary.  Choose a sufficiently
fine product grid in the full configuration space.  A product box
meeting that compact support has its one-slot time intervals
strictly ordered and strictly positive, after slightly enlarging
them while retaining their gaps.  Only finitely many boxes occur.
Localize and apply the product Fourier expansion with its larger
cutoffs still inside those enlarged intervals.  Every tensor term
then has precisely the required support order.  Its finite partial
sums converge in every Schwartz seminorm.  This proves density by
compact ordered tensors, including the time boundary and all
adjacent equal-time hyperplanes, rather than only a pointwise
ordering of generic configurations.

\subsubsection*{Exact sewing beyond tensor profiles}

For a compact ordered time configuration let
\(\rho=\min\{1,t_1,t_2-t_1,\ldots,t_n-t_{n-1}\}\).
The point vector sewing formula, with all auxiliary heat parameters
less than \(\rho/4\), gives a vector \(\mathcal V_n(X)\)
and the scalar Euclidean density \(\mathcal E_n(X)\), with
\begin{equation}
 \|\mathcal V_n(X)\|+|\mathcal E_n(X)|
                              \le C_n\rho^{-Nn}.
 \label{eq:reconstruction-testspace-estimates-RTE28}
\end{equation}
For the vector formula the first positive absolute time is needed
to make its outer heat operator bounded.  For the scalar formula
it is harmless to include that additional minimum.  The constants
are uniform in the real spatial coordinates.  On a compact set
with a fixed positive ordered gap these are bounded continuous
functions, by the actual bounded semigroup and insertion formulas.

For compact ordered tensors the narrow-time core insertion formula
proves their exact original Euclidean scalar and vector values.
The Fourier approximations just proved converge in the original
test-space seminorms, and also in \(L^1\) on this common compact
set.  The first convergence controls the original moments and
quotient vectors; the second, with \eqref{eq:reconstruction-testspace-estimates-RTE28}, controls their sewn
integrals.  Passing the finite tensor identity to the limit proves
exact sewing for every compact ordered test.

There is also an absolute-integral formulation for a general ordered
Schwartz test.  If \(\rho<1\), select its smallest boundary
distance and use Taylor's flatness along that hyperplane to order
\(M=Nn+2\).  If \(\rho=1\), the ordinary Schwartz bound
already suffices.  With \(J=4n+2\), in either case
\begin{equation}
 |F(X)|\le C_n\rho^{Nn+2}
             (1+|X|^2)^{-J/2}p_{J+Nn+2}(F).
 \label{eq:reconstruction-testspace-estimates-RTE29}
\end{equation}
The finite normal displacement in the Taylor argument preserves
the weight comparison.  The indicated seminorm pays both its
derivative order and its weight.  Multiplying by \eqref{eq:reconstruction-testspace-estimates-RTE28} leaves
\(C_n\rho^2(1+|X|^2)^{-J/2}\), which is integrable on
\(\mathbb R^{4n}\).  It follows that
\begin{equation}
 \int_{O_n}\bigl(\|\mathcal V_n(X)\|
             +|\mathcal E_n(X)|\bigr)|F(X)|\,dX
                    \le C_n p_{(N+4)n+4}(F).
 \label{eq:reconstruction-testspace-estimates-RTE30}
\end{equation}
For \(N=9\), this is at most a \(p_{17n}\) bound for
\(n\ge1\).  In particular it is compatible with the available
\(p_{20n}\) Euclidean continuity.  No factorial-in-\(n\)
assertion for the constants of this auxiliary ordered estimate is
needed or inferred.

The integral in \eqref{eq:reconstruction-testspace-estimates-RTE30} is a Bochner integral for the vector value:
the vector sewing function is norm continuous on compact ordered
sets, and the displayed majorant is integrable.  The ordered
cutoffs converge in Schwartz topology by \eqref{eq:reconstruction-testspace-estimates-RTE27}.  The formula in \eqref{eq:reconstruction-testspace-estimates-RTE30}
also controls the difference of their integrals.  Thus the compact
sewing identity extends to every ordered Schwartz test, both as a
scalar identity and as the actual ordered quotient vector.  This
proves the precise unrestricted and constrained test-space bridges
used in the field construction while leaving the ordinary
Euclidean contact domain unassigned.


% END MODULE chapters/reconstruction_testspace_estimates.tex

% BEGIN MODULE chapters/reconstruction_representation_domains.tex
\subsection{The scalar polynomial representation and its domains}
\label{ch:reconstruction-representation-domains}

The unrestricted moments now define a positive noncommutative
polynomial state.  We identify its quotient representation with the
actual band fields, including their generators, graph cores and
bounded heat-regularized insertions.  The construction is on
\(\mathcal K_{\rm ord}\).  The larger positive-time quotient
need not be identified with it.

The algebra in this section is an algebra of Minkowski test
sequences.  Its concatenations are available on all Schwartz
functions because the unrestricted joint boundary has been proved.
The earlier separated Euclidean tuple domain, on which arbitrary
concatenations need not be admissible, is not being declared an
algebra by analogy.

\subsubsection*{The full test-sequence algebra}

Let
\begin{equation}
 \mathfrak B=\bigoplus_{n\ge0}^{\rm finite}
                  \mathcal S(\mathbb R^{4n}),\qquad
 \mathcal S(\mathbb R^0)=\mathbb C.
 \label{eq:reconstruction-representation-domains-RD1}
\end{equation}
Every element has finitely many nonzero components.  Define
\begin{equation}
 \begin{gathered}
 (FG)_n=\sum_{r+s=n}F_r\otimes G_s,\qquad
 1=(1,0,\ldots),\\
 F_n^*(x_1,\ldots,x_n)
       =\overline{F_n(x_n,\ldots,x_1)}.
 \end{gathered}
 \label{eq:reconstruction-representation-domains-RD2}
\end{equation}
Associativity follows by grouping the same finite triple sum over
degrees; the unit is immediate.  Reversal interchanges the two
blocks of a tensor, so \((FG)^*=G^*F^*\), and a second reversal
and conjugation gives \((F^*)^*=F\).  These are identities of
actual Schwartz functions, including components that are not finite
tensor sums.

The topology is the locally convex direct-sum topology of the
Schwartz spaces.  A linear map from it to a locally convex space is continuous
precisely when its restriction to each degree is continuous.  On any prescribed finite
set of degrees, multiplication is jointly continuous: the tensor
seminorm estimate in \eqref{eq:reconstruction-testspace-estimates-RTE9} bounds each summand, and the degree sum is
finite.  No completion over infinitely many field degrees is needed
for the state or the representation below.

Set
\begin{equation}
 \omega(F)=\sum_nW_n(F_n),\qquad
 X(F)=\sum_nX_n(F_n).
 \label{eq:reconstruction-representation-domains-RD3}
\end{equation}
The scalar and vector sums are finite.  The proved unrestricted
positivity and tensor-density passage give
\begin{equation}
 \omega(F^*G)=\langle X(F),X(G)\rangle,
 \qquad \omega(1)=1.
 \label{eq:reconstruction-representation-domains-RD4}
\end{equation}
Thus the polynomial state has been obtained from the actual moments,
not imposed on an undefined Euclidean product.  It is continuous on
every degree.  Its topology and positivity do not assert a
\(C^*\)-norm bound for arbitrary polynomial multiplications.

\subsubsection*{Cauchy--Schwarz and the null left ideal}

For completeness the quotient construction can be made directly
from positivity.  The form \(q(F,G)=\omega(F^*G)\) is
Hermitian and positive semidefinite.  Hermiticity follows either
from \eqref{eq:reconstruction-representation-domains-RD4}, or by expanding the nonnegative real number
\(\omega((F+zG)^*(F+zG))\) for real and imaginary values of
\(z\).  For \(q(G,G)>0\), choose
\(z=-q(G,F)/q(G,G)\) in the quadratic form.  If
\(q(G,G)=0\), its linear coefficient must vanish for all complex
\(z\).  These two cases prove
\begin{equation}
 |q(F,G)|^2\le q(F,F)q(G,G).
 \label{eq:reconstruction-representation-domains-RD5}
\end{equation}
Consequently a null vector is orthogonal to every element.

Let \(\mathfrak N=\{F:q(F,F)=0\}\).  It is a linear
subspace by \eqref{eq:reconstruction-representation-domains-RD5}.  It is also a left ideal.  If \(F\in\mathfrak N\)
and \(A\in\mathfrak B\), put \(B=A^*AF\), which is another
finite algebra element.  Orthogonality gives
\begin{equation}
 \|AF\|_q^2=\omega(F^*A^*AF)=q(F,A^*AF)=0.
 \label{eq:reconstruction-representation-domains-RD6}
\end{equation}
All terms are finite algebra products; no norm bound on left
multiplication has been used.  The ideal need not be a right ideal,
and no tracial property has been assumed.

Complete \(\mathfrak B/\mathfrak N\) to a Hilbert space
\(\mathcal H_\omega\), and let
\(\mathcal D_{\rm alg}=\mathfrak B/\mathfrak N\) be its
dense quotient domain.  Define
\begin{equation}
 \pi(A)[F]=[AF],\qquad \Omega_\omega=[1].
 \label{eq:reconstruction-representation-domains-RD7}
\end{equation}
The left ideal makes this well-defined.  The domain is invariant
under every \(\pi(A)\), the algebra law is exact there, and
\(\langle[F],\pi(A)[G]\rangle
 =\langle\pi(A^*)[F],[G]\rangle\).
The adjoint's domain therefore contains the dense subspace
\(\mathcal D_{\rm alg}\), with
\(\pi(A^*)\subset\pi(A)^*\).  Every \(\pi(A)\) is
closable.  For \(A=A^*\) it is symmetric on this domain; its
closure is not thereby proved self-adjoint.

For a one-field label \(g\in\mathcal S_4\), let
\(\iota(g)=(0,g,0,\ldots)\).  The corresponding GNS field is
\(\phi_\omega(g)=\pi(\iota(g))\).  It depends linearly on
\(g\), and \(\phi_\omega(g)^*\) contains
\(\phi_\omega(\bar g)\) on the common quotient domain.
All of these assertions are about defined Hilbert-space operators.

\subsubsection*{The exact unitary identification with the ordered sector}

\eqref{eq:reconstruction-representation-domains-RD4} defines an isometry
\begin{equation}
 J:\mathcal D_{\rm alg}\longrightarrow\mathcal K_{\rm ord},
              \qquad J[F]=X(F).
 \label{eq:reconstruction-representation-domains-RD8}
\end{equation}
It is well-defined because \(\mathfrak N\) is exactly the
kernel of the actual vector map.  Conversely every vector in that
kernel is null, so no additional quotient has entered.  The range
contains every finite field-polynomial vacuum vector.  Their span
is dense in the ordered sector by the bounded sewing and cyclicity
proof.  Therefore the continuous extension of \(J\) is onto.
It is a unitary from \(\mathcal H_\omega\) to
\(\mathcal K_{\rm ord}\), sending \(\Omega_\omega\) to the
actual vacuum \(\Omega\).

For a general non-tensor component the multiplication identity in \eqref{eq:reconstruction-testspace-estimates-RTE14} says
\begin{equation}
 J\phi_\omega(g)[F]=\phi(g)J[F].
 \label{eq:reconstruction-representation-domains-RD9}
\end{equation}
The right side is defined because \(X(F)\in\mathcal D_\infty\).
For finite tensors it is the ordinary product identity.  Passing
to a general component in the energy graph norm of order \(N\)
proves exactly the same equality; passing only in Hilbert norm
would not have sufficed for the unbounded field on the right.

This identifies the field representation, not merely its numerical
vacuum moments.  In particular the cyclic Hilbert space in the
reconstruction theorem is the actual closed ordered range.  The
unitary \eqref{eq:reconstruction-representation-domains-RD8} supplies no fields on
\(\mathcal K\ominus\mathcal K_{\rm ord}\), and it does not
require that complement to vanish.

\subsubsection*{Translations, generators and spectral projections}

Let \(\tau_a\) shift every spacetime argument by the same real
vector in the active test convention.  Translation invariance of
\(W\) preserves \eqref{eq:reconstruction-representation-domains-RD4}, so
\(U_\omega(a)[F]=[\tau_aF]\) extends to a unitary on the GNS
space.  The vector-valued Schwartz continuity makes it strongly
continuous on the dense quotient domain; unitarity extends that
continuity to all vectors.  The actual band-field covariance gives
\begin{equation}
 JU_\omega(t,y)J^{-1}=e^{itH}U_y
                         \quad\hbox{on }\mathcal K_{\rm ord}.
 \label{eq:reconstruction-representation-domains-RD10}
\end{equation}
This equality of bounded groups is first verified on the dense
polynomial vectors and then extended.  Stone's uniqueness theorem
therefore identifies their self-adjoint generators with their exact
domains.  In particular
\begin{equation}
 \begin{gathered}
 J\operatorname{Dom}H_\omega
          =\operatorname{Dom}H\cap\mathcal K_{\rm ord},
 \qquad JH_\omega J^{-1}=H_{\rm ord},\\
 Je^{-sH_\omega}J^{-1}=T_s|_{\mathcal K_{\rm ord}},
 \qquad J\mathsf E_\omega(A)J^{-1}
                 =\mathsf E(A)|_{\mathcal K_{\rm ord}}.
 \end{gathered}
 \label{eq:reconstruction-representation-domains-RD11}
\end{equation}
The last assertion holds for every Borel set of joint energy and
momentum.  Indeed the separate strongly commuting groups in \eqref{eq:reconstruction-representation-domains-RD10}
have the same joint spectral representation after conjugation.
One can also obtain it first for their bounded Fourier transforms,
then for the generated spectral projections by the uniqueness part
of joint functional calculus.  Thus there is no new choice of
Hamiltonian or of momentum in the quotient representation.

For any integer \(l\), the same functional calculus gives
\(J\operatorname{Dom}(1+H_\omega)^l
 =\operatorname{Dom}W^l\cap\mathcal K_{\rm ord}\).
The common energy domain is transported exactly.  The signs in
\eqref{eq:reconstruction-representation-domains-RD10} are those of the band construction; physical spatial
momentum in \(Et-p\cdot x\) is \(-P\).  Its joint cone and
all-cut moment spectrum are consequently the already proved ones.

\subsubsection*{Integrated polynomial operators on the energy domain}

The vector extension can be strengthened to an operator extension.
For a finite product of \(m\) fields and \(l\ge0\), iteration
of the energy bound gives
\begin{equation}
 \left\|W^l\phi(g_1)\cdots\phi(g_m)
                              W^{-(mN+l)}\right\|
       \le C_{m,l}\prod_i p_{mN+l+7}(g_i).
 \label{eq:reconstruction-representation-domains-RD12}
\end{equation}
The product is initially defined on the invariant energy domain.
The bound gives its unique bounded extension between the indicated
weighted spaces.  Regard this as an \(m\)-linear map into the
Banach space of bounded operators from
\(\operatorname{Dom}W^{mN+l}\), with norm
\(\|W^{mN+l}\xi\|\), to
\(\operatorname{Dom}W^l\), with norm \(\|W^l\xi\|\).
The constructive Fourier lemma applies in that actual Banach space.
For \(F\in\mathcal S(\mathbb R^{4m})\) it defines
\(\Phi_m(F)\) with
\begin{equation}
 \|W^l\Phi_m(F)W^{-(mN+l)}\|
              \le C'_{m,l}p_{m(mN+l+15)}(F).
 \label{eq:reconstruction-representation-domains-RD13}
\end{equation}
The partial Fourier sums at different \(l\) are the same
polynomial operators.  On a vector in a higher energy domain their
Hilbert limits agree; closedness of \(W^l\) identifies their
graph limits.  Hence these definitions give one coherent operator
on \(\mathcal D_\infty\), mapping it into itself.

For finite tensors, concatenation and reversal are the ordinary
field-product rules.  For general \(F,G\), approximate each by
finite tensors in every seminorm needed by \eqref{eq:reconstruction-representation-domains-RD13}.  Before composing,
use the estimate with the larger graph exponent required by the
left factor.  Both compositions then converge on the same domain.
This proves
\begin{equation}
 \begin{gathered}
 \Phi_m(F)\Phi_n(G)=\Phi_{m+n}(F\otimes G)
                         \quad\hbox{on }\mathcal D_\infty,\\
 \langle\xi,\Phi_m(F)\eta\rangle
       =\langle\Phi_m(F^*)\xi,\eta\rangle
                         \quad(\xi,\eta\in\mathcal D_\infty).
 \end{gathered}
 \label{eq:reconstruction-representation-domains-RD14}
\end{equation}
The tensor approximations converge in graph norms, so this passage
does not apply a merely closable operator to an uncontrolled
Hilbert-space limit.  Each \(\Phi_m(F)\) restricted to the
common domain is closable, because its adjoint contains the densely
defined conjugate operator in \eqref{eq:reconstruction-representation-domains-RD14}.

The finite sum \(\Phi(A)=\sum_m\Phi_m(A_m)\) is now an actual
representation of \(\mathfrak B\) on the invariant common
domain.  It extends its GNS polynomial realization: on
\(J\mathcal D_{\rm alg}\), \eqref{eq:reconstruction-representation-domains-RD14} and vacuum cyclicity give
\(\Phi(A)J[F]=J[AF]\).  Thus the unrestricted algebra product agrees with the field product
on the energy domain.

\subsubsection*{Polynomial vectors form an energy graph core}

Let \(\mathcal P\) be the span of finite tensor field-polynomial
vacuum vectors.  It is dense in the ordered sector and is contained
in every energy domain.  It is invariant under real time
translations by field covariance.  We prove, for each fixed
\(l\ge0\),
\begin{equation}
 \overline{\mathcal P}^{\,\|W^l\cdot\|}
          =\operatorname{Dom}W^l\cap\mathcal K_{\rm ord}.
 \label{eq:reconstruction-representation-domains-RD15}
\end{equation}
Choose \(c_R\in C_c^\infty(\mathbb R)\), equal to one on
\([0,R]\), with \(0\le c_R\le1\) on the positive axis and
support there below \(2R+2\).  The inverse Fourier transform
is integrable, and the spectral theorem gives
\(c_R(H)=\int\check c_R(t)e^{itH}\,dt\) in the strong
sense, with a consistent Fourier normalization.

For \(\zeta\in\mathcal P\), this is also a Bochner integral
in the \(W^l\) graph norm.  The translated vector has constant
graph norm \(\|W^l\zeta\|\); its orbit is graph continuous,
because \(e^{itH}\) commutes with \(W^l\) and is strongly
continuous on the vector \(W^l\zeta\).  First truncate the
integrable Fourier tail.  On the remaining compact time interval,
uniform continuity gives finite simple-function or Riemann-sum
approximations in the graph norm.  Every sample is a translated
polynomial vector, hence remains in \(\mathcal P\).
Therefore \(c_R(H)\zeta\) belongs to its graph closure.

Now let \(\xi\in\operatorname{Dom}W^l\cap\mathcal K_{\rm ord}\).
At fixed \(R\), choose \(\zeta_j\in\mathcal P\) converging
to \(\xi\) in Hilbert norm.  Boundedness of \(W^lc_R(H)\)
makes \(c_R(H)\zeta_j\to c_R(H)\xi\) in graph norm.
The previous paragraph puts all these approximants in the graph
closure of \(\mathcal P\).  Finally
\(c_R(H)\xi\to\xi\) in that graph norm by dominated spectral
convergence.  This proves \eqref{eq:reconstruction-representation-domains-RD15}, including the order of the two
approximations and the bounded operator which justifies the first.

The dense GNS domain \(J\mathcal D_{\rm alg}\) contains
\(\mathcal P\), so it is also an energy graph core.  Consequently
the extension from the polynomial quotient to
\(\mathcal D_\infty\) does not create a different closure of a
fixed field.  For example,
\begin{equation}
 \begin{aligned}
 \overline{\phi(g)|_{\mathcal P}}
 &=\overline{\phi(g)|_{J\mathcal D_{\rm alg}}}\\
 &=\overline{\phi(g)|_{\mathcal D_\infty\cap\mathcal K_{\rm ord}}}\\
 &=\overline{\phi(g)|_{\operatorname{Dom}W^N\cap\mathcal K_{\rm ord}}}.
 \end{aligned}
 \label{eq:reconstruction-representation-domains-RD16}
\end{equation}
To prove the last equality, approximate a vector in
\(\operatorname{Dom}W^N\) by polynomial vectors in that graph
norm.  Their field images converge by the \(k=0\) energy bound.
Thus its graph lies in the closure of the polynomial graph.
The reverse graph inclusions follow from the domain inclusions.
All operators in \eqref{eq:reconstruction-representation-domains-RD16} are closable, by the conjugate-field adjoint
identity.  The same proof applies to a fixed integrated polynomial
of degree at most \(m\), using \(W^{mN}\).

\subsubsection*{Poincare implementation on both natural domains}

The reconstructed moments are Poincare covariant.  Pullback of tests
therefore gives a null-preserving isometry
\([F]\mapsto[\alpha_{a,\Lambda}F]\) on the quotient.  Its
inverse is the inverse group transformation.  It extends to a
unitary, with exact group law.  Schwartz continuity on every fixed
degree makes it strongly continuous on the quotient domain, and
unitarity extends this to the completion.  On polynomial vectors
it has exactly the band representation's action, so \eqref{eq:reconstruction-representation-domains-RD8}
intertwines the two representations on the entire Hilbert spaces.

The quotient domain is invariant under this action.  The energy
domain is invariant as well, for a different reason.  If
\((\gamma,b)=\Lambda e_0\), the exact translation group law
and Stone's theorem give
\(U(\Lambda)HU(\Lambda)^{-1}=\gamma H+b\cdot P\).
The round cone yields the positive comparisons
\begin{equation}
 (\gamma-|b|)H\le\gamma H+b\cdot P
                   \le(\gamma+|b|)H.
 \label{eq:reconstruction-representation-domains-RD17}
\end{equation}
These are inequalities in one joint spectral representation.
They imply equality of every integer-power domain and equivalence
of its graph norms, since \(\gamma>|b|\).
For a fixed group element, the action is consequently continuous
in each energy graph norm.  The field intertwining law on the
polynomial domain extends by \eqref{eq:reconstruction-representation-domains-RD15} and \eqref{eq:reconstruction-representation-domains-RD13} to the full common
energy domain, for both single fields and integrated words.

This verifies two distinct useful domains: the algebraically
invariant quotient domain and the complete intersection of energy
domains.  They need not be equal as sets.  They give the same
canonical field closures by \eqref{eq:reconstruction-representation-domains-RD16}, and the exact same Hilbert
representation by \eqref{eq:reconstruction-representation-domains-RD8}.  No choice of an unspecified larger
extension is hidden in the comparison.

\subsubsection*{Local polynomial algebras and the locality ideal}

For an open spacetime set \(O\), let
\(\mathfrak B_{\rm pol}(O)\) be the unital algebra generated by
the degree-one labels in \(C_c^\infty(O)\), and use the actual
representation \(\Phi\) on \(\mathcal D_\infty\).  Isotony
is immediate from inclusion of supports.  Poincare transformations
carry this algebra to that of the transformed region, by the
already proved common-domain covariance.

If \(O_1\) and \(O_2\) are spacelike separated, every generator
from the first commutes with every generator from the second on
the invariant domain.  Repeatedly move adjacent generators past
one another, applying the proved field commutator at each step.
Every intermediate vector remains in the common domain.  Hence
the represented polynomial algebras commute there:
\begin{equation}
 \begin{gathered}
 \relax[\Phi(A),\Phi(B)]\xi=0,\\
 A\in\mathfrak B_{\rm pol}(O_1),\quad
 B\in\mathfrak B_{\rm pol}(O_2),\quad
 \xi\in\mathcal D_\infty.
 \end{gathered}
 \label{eq:reconstruction-representation-domains-RD18}
\end{equation}
One may also use compact jointly smeared words supported in
\(O_i^m\).  Cover that compact support by product boxes whose
larger cutoffs still lie in \(O_i\), and apply the support-preserving
Fourier approximation.  The operator graph bound in \eqref{eq:reconstruction-representation-domains-RD13} then
passes the same commutation identity to those general smearing
functions.  This is a topological extension of the local polynomial
algebra, not an assertion that every such smearing is literally a
finite tensor sum.

Let \(\mathfrak I_{\rm loc}\) be the two-sided algebra ideal
generated by \(\iota(g)\iota(h)-\iota(h)\iota(g)\) with
spacelike separated supports.  Each generator acts as zero on
the GNS domain by locality.  Left and right multiplication
preserve that domain, so every element of the ideal acts as zero.
In particular
\begin{equation}
 \mathfrak I_{\rm loc}\subset\mathfrak N,
 \qquad \pi(\mathfrak I_{\rm loc})=0.
 \label{eq:reconstruction-representation-domains-RD19}
\end{equation}
The positive scalar state therefore factors through the local
polynomial quotient.  The ideal is invariant under the scalar
Poincare pullbacks because those preserve spacelike separation.
It is also invariant under star: the star of a generating commutator
is the negative commutator of the conjugate labels, with their
supports unchanged.
This provides an actual local polynomial observable structure.
It does not assert strong commutation of spectral projections of
closures of real fields, or select self-adjoint extensions of
them.  A bounded local operator net needs an additional operator
criterion if it is to be built through such spectral calculi.

\subsubsection*{Recovering the actual bounded insertion operators}

The preceding equivalence can be checked against the original
bounded Euclidean insertions.  On the ordered sector the
energy-band proof established the operator-norm limit
\begin{equation}
 T_s\phi(\chi_\epsilon\otimes h)T_s
                      \longrightarrow\Phi_s(h)
                  \quad(\epsilon\longrightarrow0).
 \label{eq:reconstruction-representation-domains-RD20}
\end{equation}
Here \(\chi_\epsilon\) is the fixed nonnegative scaled temporal
approximate identity chosen in the band construction, with
\(\|\chi_\epsilon\|_1=1\), and \(h\) is a compact smooth spatial
label.  On finitely many bands the limit follows
from the time approximate identity and bounded generators.  The
remaining block tails are bounded by
\(C(1+m+n)^Ne^{-s(m+n)}\|h\|_1\), with vanishing absolute
row and column tails uniformly in \(\epsilon\).  Thus \eqref{eq:reconstruction-representation-domains-RD20}
is a bounded-operator limit, not a formal sharp-time substitution.

The same construction on the GNS representation is conjugated by
\(J\) at every \(\epsilon\), by \eqref{eq:reconstruction-representation-domains-RD9}--\eqref{eq:reconstruction-representation-domains-RD11}.
The bounded limit is consequently the actual original Riesz
operator \(\Phi_s(h)\).  Spatial approximate identities give
the point operator \(\Phi_s(x)=V_{4s}(0,x)\) in operator norm,
because the analytic density is norm continuous in real space.
This identifies more than the vacuum state of the insertion:
it identifies the operator between every pair of Hilbert vectors.

There is an exact formula for its real Euclidean time translate.
For \(0<s\le1/4\) and \(|t|<s\), choose \(0<a<s-|t|\).  Then
\begin{equation}
 V_{4s}(t,x)
          =T_{s+t-a}\Phi_a(x)T_{s-t-a}.
 \label{eq:reconstruction-representation-domains-RD21}
\end{equation}
To verify it, first replace the point label by a narrow compact
spacetime approximation and test on ordered compact core vectors.
The right heat factor shifts the old entries by \(s-t-a\).
The insertion at scale \(a\) puts the new entry at time \(a\)
and shifts the old entries by another \(2a\).  The left heat
factor finally puts the new entry at \(s+t\) and all old entries
at their original times plus \(2s\).  This is precisely the
actual core formula with outside entries shifted by \(s\) and
the central insertion at \(t\).  Both sides are bounded.  Dense
core equality and the norm point limits prove \eqref{eq:reconstruction-representation-domains-RD21} wherever
the original density is defined.  The displayed bounded formula
also extends it to the larger interval \(|t|<s\).

For the jointly complex translate, write
\(S(z,w)=e^{-zH+iw\cdot P}\).  On the common domain where its
two factors are contractions, the formula is
\begin{equation}
 V_{4s}(z_4,z_{\rm sp})
    =S(s+z_4-a,z_{\rm sp})\Phi_a(0)
                           S(s-z_4-a,-z_{\rm sp}),
 \label{eq:reconstruction-representation-domains-RD22}
\end{equation}
with \(|\operatorname{Re}z_4|+
|\operatorname{Im}z_{\rm sp}|_2<s-a\).
At real arguments it agrees with \eqref{eq:reconstruction-representation-domains-RD21} and spatial covariance.
Both sides are operator-norm holomorphic on the intersection
with the original PD strip.  Testing matrix elements and applying
the several-variable identity theorem from a real open box proves
their equality on connected common neighbourhoods.  In this way
the reconstructed generators and field identify the actual PD
analytic density, with its branch fixed by the original insertion
form.  An arbitrary bounded coefficient function having the same
vacuum expectation would not have sufficed for this conclusion.



% END MODULE chapters/reconstruction_representation_domains.tex

% BEGIN MODULE chapters/compact_disintegration_v3.tex
\subsection{Conditional measures on compact metric spaces}
\label{ch:compact-disintegration}

The conditional laws used in the phase construction take values in the
Borel probability measures on a compact metric configuration space.
The conditioning field may be completed and need not be countably
generated as a literal collection of sets. It is the Borel field of the
\emph{target} space that is countably generated.

\begin{proposition}[Compact-metric conditional law]
\label{thm:compact-conditional-law}
Let $X$ be a nonempty compact metric space, let $\mu$ be a Borel
probability measure on $X$, and let $\overline{\mathcal B}_\mu$ be the
$\mu$-completion of its Borel field $\mathcal B(X)$. Suppose

\refstepcounter{equation}
\begin{gather*}
\FitMath{\mathcal N_\mu\subset\mathcal H\subset
       \overline{\mathcal B}_\mu,}\tag{\theequation}\label{eq:compact-disintegration-CD1}
\end{gather*}

where $\mathcal H$ is a $\sigma$-field and $\mathcal N_\mu$ denotes
all subsets of Borel $\mu$-null sets. There exists a map
$K:X\to\mathcal P(X)$, where $\mathcal P(X)$ carries its weak topology,
with the following properties.

The map $K$ is both Borel measurable and $\mathcal H$-measurable. For
every bounded Borel $f:X\to\mathbb C$,

\refstepcounter{equation}
\begin{gather*}
\FitMath{Kf(x):=\int_X f(y)\,K(x,dy)} \\*[3pt]
\FitMath{\text{is a version of }E_\mu[f\mid\mathcal H].}\tag{\theequation}\label{eq:compact-disintegration-CD2}
\end{gather*}

In particular $x\mapsto K(x,B)$ is $\mathcal H$-measurable for every
$B\in\mathcal B(X)$, and

\refstepcounter{equation}
\begin{gather*}
\FitMath{\int_H K(x,B)\,\mu(dx)=\mu(H\cap B)} \\*[3pt]
\FitMath{(H\in\mathcal H, B\in\mathcal B(X)).}\tag{\theequation}\label{eq:compact-disintegration-CD3}
\end{gather*}

The measure on the right is understood on the completion. If $K'$ is
another probability-valued kernel satisfying \eqref{eq:compact-disintegration-CD2}, then there is a
single $\mu$-null set $N$ such that

\refstepcounter{equation}
\begin{gather*}
\FitMath{K(x,\cdot)=K'(x,\cdot)} \\*[3pt]
\FitMath{(x\notin N)}\tag{\theequation}\label{eq:compact-disintegration-CD4}
\end{gather*}

as Borel probability measures. One may take $N$ to be a Borel null set.
Consequently the law $K_*\mu$ on $\mathcal P(X)$ is independent of the
choice of conditional kernel.
\end{proposition}

\subsubsection*{Borel representatives and countably many tests}

We first record why completion causes no obstacle to the Borel version
in this statement. Every bounded
$\overline{\mathcal B}_\mu$-measurable real function $u$ has a Borel
representative. Indeed, for each rational $r$ replace $\{u>r\}$ by a
Borel set differing from it inside a Borel null set. Outside the union
of those countably many null sets, the rational level cuts reconstruct
$u$. Their supremum, clipped to a bound for $u$, is a Borel function
equal to $u$ there. If $u$ was $\mathcal H$-measurable, this
representative remains $\mathcal H$-measurable: the change is confined
to a Borel null set, all of whose subsets belong to $\mathcal H$ by
\eqref{eq:compact-disintegration-CD1}. This observation will be applied only countably many times.

Choose a countable dense set $\{x_j\}$ in $X$. Let $D\subset C(X;\mathbb R)$
be the collection generated by the constant $1$ and the functions
$y\mapsto d(y,x_j)$ under rational linear combinations, multiplication,
pointwise minimum and pointwise maximum. It is countable, is a rational
vector lattice and algebra, and separates points: if $x\ne y$, a dense
point sufficiently near $x$ gives different distances from $x$ and $y$.
Its uniform closure is a real algebra containing the constants and
separating points, so the Stone--Weierstrass theorem makes $D$ uniformly
dense in $C(X;\mathbb R)$.

For each $f\in D$ take a conditional expectation $u_f$ of $f$ given
$\mathcal H$. These conditional expectations can be obtained as the
orthogonal projections of $f$ onto the closed subspace
$L^2(\mathcal H,\mu)$ of $L^2(\mu)$. The projection identity against
$1_H$, $H\in\mathcal H$, gives

\refstepcounter{equation}
\begin{gather*}
\FitMath{\int_Hu_f\,d\mu=\int_H f\,d\mu.}\tag{\theequation}\label{eq:compact-disintegration-CD5}
\end{gather*}

Testing the sets on which $u_f<0$, or $u_f$ exceeds an upper bound for
$f$, proves positivity and the usual $L^\infty$ bound. Linearity and
uniqueness follow from the same identity. Choose for each $u_f$ a
bounded representative that is both Borel and $\mathcal H$-measurable,
using the preceding observation.

There are only countably many assertions

\refstepcounter{equation}
\begin{gather*}
\FitMath{u_1=1,} \\*[3pt]
\FitMath{u_{rf+sg}=ru_f+su_g} \\*[3pt]
\FitMath{(f,g\in D, r,s\in\mathbb Q),} \\*[3pt]
\FitMath{u_f\ge0} \\*[3pt]
\FitMath{(f\in D, f\ge0\text{ on }X).}\tag{\theequation}\label{eq:compact-disintegration-CD6}
\end{gather*}

Each holds almost surely. Their exceptional sets therefore have a
single Borel, $\mathcal H$-measurable null union $N_0$. This is the
simultaneous exceptional set needed for the construction. No
intersection over all continuous or all Borel functions is taken.

\subsubsection*{Pointwise Riesz representation}

For $x\notin N_0$ define $L_x(f)=u_f(x)$ on $D$. It is rational linear,
positive on nonnegative elements of $D$, and $L_x(1)=1$. If rational
$r>\|f\|_\infty$, then $r\pm f\in D$ are nonnegative, so

\begin{gather*}
\FitMath{|L_x(f)|\le r.}
\end{gather*}

Letting rational $r$ decrease to $\|f\|_\infty$ proves
$|L_x(f)|\le\|f\|_\infty$. Uniform density consequently gives a unique
continuous extension to $C(X;\mathbb R)$. The extension is real linear:
approximate a real scalar by rationals and use the norm bound. It is
positive: if $f\ge0$ is continuous, choose $f_j\in D$ with
$\|f_j-f\|_\infty<1/j$; then $f_j+1/j\ge0$ and
$L_x(f_j)\ge-1/j$. It still takes the value one at $1$.

The Riesz representation theorem supplies a unique regular Borel
probability measure $K(x,\cdot)$ representing $L_x$. For $x\in N_0$
set $K(x,\cdot)=\delta_{x_*}$ at one fixed $x_*\in X$. Thus $K(x,\cdot)$
is a probability measure at \emph{every} point. Its values on the null
set are a choice, with no conditional interpretation required.

We verify measurability of the probability-valued map. Enumerate $D$ as
$\{d_j\}$ and define

\refstepcounter{equation}
\begin{gather*}
\FitMath{I:\mathcal P(X)\longrightarrow\mathbb R^{\mathbb N},} \\*[3pt]
\FitMath{I(\rho)=(\rho(d_j))_j.}\tag{\theequation}\label{eq:compact-disintegration-CD7}
\end{gather*}

The map is continuous and injective, by uniform density and uniqueness
in Riesz representation. The weak probability space $\mathcal P(X)$
is metrized by
$\sum_j2^{-j}|\rho(d_j)-\rho'(d_j)|/(1+2\|d_j\|_\infty)$:
uniform density identifies convergence in this metric with weak
convergence. It is compact. Indeed, compactness follows by taking a
diagonal subsequence of integrals on $D$, extending the resulting
positive normalized functional as above, and applying Riesz again.
Hence $I$ is a homeomorphism onto its compact image in the Hausdorff
space $\mathbb R^{\mathbb N}$.

Every coordinate of $I\circ K$ equals $u_{d_j}$ outside $N_0$ and the
constant $d_j(x_*)$ on $N_0$. It is both Borel and
$\mathcal H$-measurable. The countable product map is therefore
measurable for both fields; composing with the inverse of $I$ proves
the asserted measurability of $K$. This also supplies an explicit
countable determining family for uniqueness.

\subsubsection*{All bounded Borel tests and uniqueness}

For $f\in D$, the equality $Kf=u_f$ holds outside $N_0$; thus \eqref{eq:compact-disintegration-CD5}
holds with $Kf$ in place of $u_f$. Uniform approximation extends it to
every real continuous $f$, and $Kf$ is measurable as a uniform limit
of the corresponding measurable integrals.

Here are the remaining monotone-class details. The collection of
bounded real Borel functions $f$ for which $Kf$ is measurable for both
fields and

\refstepcounter{equation}
\begin{gather*}
\FitMath{\int_H Kf\,d\mu=\int_H f\,d\mu} \\*[3pt]
\FitMath{\text{for every }H\in\mathcal H}\tag{\theequation}\label{eq:compact-disintegration-CD8}
\end{gather*}

is a vector space containing $C(X;\mathbb R)$. It is closed under
uniformly bounded increasing limits, by monotone convergence for each
probability $K(x,\cdot)$ and for $\mu$; subtracting a constant reduces
to nonnegative functions. The bounded functional monotone-class
theorem therefore makes it all bounded Borel functions, since the
continuous algebra generates $\mathcal B(X)$. One can see the latter
directly from a countable dense family of distance functions and the
countable base of metric balls. Complex functions follow by their real
and imaginary parts. This proves \eqref{eq:compact-disintegration-CD2}, and indicators give \eqref{eq:compact-disintegration-CD3}.
In particular it also verifies directly that $K$ is a probability
kernel, including measurability of every Borel-set evaluation.

For uniqueness, apply uniqueness of conditional expectation to $Kd_j$
and $K'd_j$ for the countably many $d_j\in D$. Outside a single null
set all these integrals agree. Uniform density then gives equality on
every continuous function, and Riesz uniqueness gives \eqref{eq:compact-disintegration-CD4}. The null
set belongs to the completed Borel field and hence is contained in a
Borel null set, as asserted. Changing a map on a null set does not
change its probability-valued law. The law has barycenter $\mu$,
because \eqref{eq:compact-disintegration-CD8} with $H=X$ gives

\refstepcounter{equation}
\begin{gather*}
\FitMath{\mu(f)=\int_{\mathcal P(X)}\rho(f)\,(K_*\mu)(d\rho)} \\*[3pt]
\FitMath{(f\text{ bounded Borel}).}\tag{\theequation}\label{eq:compact-disintegration-CD9}
\end{gather*}

The integrand is Borel on $\mathcal P(X)$: this is immediate for
continuous $f$, and the same bounded monotone-class argument proves it
for every bounded Borel $f$.\hfill$\square$

\subsubsection*{The completed invariant field and the exact null-set scope}

For a countable group of Borel, $\mu$-preserving transformations,
the field of events in $\overline{\mathcal B}_\mu$ invariant modulo
$\mu$ under every group element is a $\sigma$-field containing
$\mathcal N_\mu$. Thus it satisfies \eqref{eq:compact-disintegration-CD1}. In particular the theorem
applies to the completed translation-invariant field of a compact
metric lattice configuration space. The conditional measures are Borel
probabilities on that configuration space; no countable generation of
the completed invariant field is needed.

The uniqueness assertion is simultaneous for the resulting probability
measures. The assertion that $Kf$ is a conditional-expectation version
holds separately for each bounded Borel $f$ and, of course, for any
preselected countable collection on one common conull set. It does not
assert equality outside one universal exceptional set to an arbitrary
previously selected version for \emph{every} Borel function.

Nor does the theorem assert pointwise properness
$K(x,H)=1_H(x)$ simultaneously for all conditioning events. Even when
those events are Borel, such a demand is false for completed fields.
For example let $\mu$ be nonatomic and let
$\mathcal H=\{H\in\overline{\mathcal B}_\mu:
\mu(H)\in\{0,1\}\}$ be the null/co-null field.
The conditional kernel is $K(x,\cdot)=\mu$ almost surely.
If simultaneous properness held outside one null set, then at any
point $x$ of that conull set the conditioning event $\{x\}\in\mathcal H$
would require $K(x,\{x\})=1$, whereas $\mu(\{x\})=0$. This does not
contradict \eqref{eq:compact-disintegration-CD2}: for a fixed Borel $H\in\mathcal H$, it gives
$K(x,H)=1_H(x)$ almost surely, with an exceptional set that may depend
on $H$. \eqref{eq:compact-disintegration-CD3} is the corresponding integral identity. The theorem concerns the countably generated Borel target
$\mathcal B(X)$; no regular-conditional-law assertion for an arbitrary
non-countably-generated target field is being made.


% END MODULE chapters/compact_disintegration_v3.tex

% BEGIN MODULE chapters/phase_interface_v2.tex
\subsection{Canonical phases from lattice estimates}
\label{ch:phase-interface}

This section gives a conditional implication with completely specified
probabilistic objects. The finite-spacing law G3 is proved in
\cref{ch:plateau-proof}; it is used before any continuum component
is constructed. The all-state moment and
one-point comparison estimate G4 is derived from the ordinary torus
source hypothesis D1 in \cref{ch:all-state-source-transfer}. Conditional probability measures are
constructed in \cref{ch:compact-disintegration}. The argument here
deduces the canonical phase measure from those precise ingredients.

In this section a pure state means an extreme point of the translation-invariant
DLR class $\mathcal G_j$, equivalently an ergodic state for lattice translations.
Extremality in the larger class of all, possibly noninvariant, DLR states is not
used and is not identified with this property without an additional argument.

\phantomsection\label{notdef:invariant-dlr}
Put $a_j=L^{-K_j}\downarrow0$, with $L$ an integer, and let
$\Omega_j=\operatorname{SU}(2)^{\mathcal B_j}$ be the space of bond configurations
on $a_j\mathbb Z^4$. The set $\mathcal B_j$ contains one orientation of each
bond. Its countability makes $\Omega_j$ compact and metrizable: enumerate the
bonds and use the sum of $2^{-m}$ times bounded metrics on the group factors.
The Wilson specification $\gamma_W$ integrates finitely many bonds $W$ with
density proportional to
$\exp[-\beta_j\sum_{p:p\cap W\ne\varnothing}a_p]$.
The denominator is positive, and compactness of the integration space shows
that $\gamma_W F$ is continuous and local whenever $F$ is. Define
$\mathcal G_j$ to consist of the translation-invariant probability measures
with $\rho(\gamma_WF)=\rho(F)$ for every finite $W$ and continuous local $F$.
The equations are weakly closed, so $\mathcal G_j$ is compact and convex.
Its nonemptiness follows from the fixed-spacing torus-limit construction,
and a centering state $\nu_j\in\mathcal G_j$ is fixed.
G1 supplies the lattice symmetries and mirror positivity used in G3.

Let $\mathfrak T$ be the set of finite tuples of real $C^1$ compactly supported
functions with pairwise disjoint supports, including the empty tuple.
Concatenation is used only when its resulting supports are disjoint. For
$X=(f_1,\ldots,f_k)$ define the bounded continuous local function

\begin{gather*}
\FitMath{X_j=\prod_{i=1}^k
       [\mathcal O_j(f_i)-\nu_j(\mathcal O_j(f_i))],} \\[3pt]
\FitMath{\varnothing_j=1.}
\end{gather*}

Assume $\nu_j(X_j)\to\sigma(X)$ on $\mathfrak T$.
For support radii $R$ and separation $d$, put $b_1=C_1(R)$ and
$b_n=C_n(d;R)$ for $n\ge2$, and define

\begin{gather*}
\FitMath{M_k(d;R)=\sum_{\pi\in\operatorname{Part}(\{1,\ldots,k\})}
                         \prod_{J\in\pi}b_{|J|},} \\[3pt]
\FitMath{M_0=1.}
\end{gather*}

The ordinary moment--cumulant identity, applied under any $\rho\in\mathcal G_j$,
gives

\begin{gather*}
\FitMath{|\rho(X_j)|\le M_k(d;R)\prod_i\|f_i\|_\sharp.}
\end{gather*}

Indeed, a singleton block contributes the difference between the means in
$\rho$ and $\nu_j$, and a larger block contributes the corresponding cumulant
under $\rho$. This is why a bound on one-point differences is needed.
Telescoping a difference of two products gives equicontinuity with the same
constants whenever all supports remain in fixed separated containers.

\subsubsection{Decomposition at one spacing}

Let $\mathcal I$ denote the completed invariant $\sigma$-field for a fixed
$\nu\in\mathcal G_j$. \Cref{thm:compact-conditional-law} gives
a probability kernel $\omega\mapsto\nu^\omega$ satisfying
$\nu^\omega(F)=E_\nu[F\mid\mathcal I](\omega)$ for bounded Borel $F$.
We verify that the components are states of the required class.

First, a bounded invariant density preserves the DLR equations. To prove it,
approximate such a density $h$ in $L^1(\nu)$ by bounded continuous local
functions $h_n$. This approximation is available because cylinder functions
are uniformly dense in continuous functions on the compact product space,
and continuous functions are dense in $L^1$ of a Borel probability measure.
Translate $h_n$ so its support misses $W$. Invariance of $h$ and $\nu$ leaves
the approximation error unchanged, while
$\gamma_W(h_nF)=h_n\gamma_WF$ for the translated function. Passing to the
limit proves $\nu(h\gamma_WF)=\nu(hF)$.

Consequently $\nu^\omega(\gamma_WF)=\nu^\omega(F)$ almost surely.
Likewise $\nu^\omega(\tau_vF)=\nu^\omega(F)$ almost surely for every lattice
translation $v$. These statements can be imposed on one common full-measure
set: finite bond sets and translations form countable families, and continuous
functions on the compact configuration space have a countable dense family
of continuous cylinder functions. Supremum-norm continuity then extends the
equations to all continuous local functions. A monotone-class argument extends
the integral identities to bounded Borel functions. Thus
$\nu^\omega\in\mathcal G_j$ almost surely.

The law G3 proves that almost every such state is ergodic. Fix a continuous
local $A$. Applied to $\nu^\omega$ and then integrated, its far-translation
limit gives

\begin{gather*}
\FitMath{\int\operatorname{Pl}_{\nu^\omega}(A,A)\,d\nu(\omega)
 =\operatorname{Pl}_\nu(A,A)
 =\int|\nu^\omega(A)|^2\,d\nu(\omega).}
\end{gather*}

The first passage to the limit uses the bound $\|A\|_\infty^2$ and dominated
convergence. The difference of the integrands on the two sides is the
nonnegative variance of the invariant conditional expectation of $A$ in
$\nu^\omega$. It is therefore zero almost surely. Impose this on a countable
dense family of continuous local $A$. Density in $L^2(\nu^\omega)$ then says
that every invariant $L^2$ function is constant, which is ergodicity.

\phantomsection\label{notdef:ergodic-decomposition}
Write $\mathcal E_j$ for these ergodic states. They are exactly the extreme
points of $\mathcal G_j$. A nontrivial invariant event decomposes a state
into two DLR states by the invariant-density argument. Conversely, the
Radon--Nikodym density in a convex decomposition of an ergodic invariant state
is invariant, and hence constant. Moreover $\mathcal E_j$ is closed. If
$\rho_n\in\mathcal E_j$ tends weakly to $\rho$, G3 gives a decay estimate
uniform in $n$. Passing first to this weak limit and then to far translations
identifies $\operatorname{Pl}_\rho(A,A)$ with $|\rho(A)|^2$. The same dense-set
argument proves that $\rho$ is ergodic. Thus $\mathcal E_j$ is compact.

Let $D_j$ be the law of $\nu_j^\omega$ under $\nu_j$. It is supported on
$\mathcal E_j$, has barycenter $\nu_j$, and does not depend on versions of the
kernel. For this last assertion choose one countable determining family of
continuous cylinder functions. Conditional expectations are unique almost
surely for that entire family simultaneously; equality on the family implies
equality of the probability measures. This also establishes measurability of
the probability-valued map and of every evaluation used below.

Uniqueness among all decompositions supported on $\mathcal E_j$ requires one
more observation. For any such law with barycenter $\nu_j$, the uniform law G3
identifies the integral of a product of finitely many state evaluations with
successive far-translation limits of $\nu_j$. These moments are therefore the
same as those of $D_j$. Evaluations of continuous local functions are continuous
and separate measures in the compact space $\mathcal E_j$. Their polynomial
algebra is dense by Stone--Weierstrass, so equal mixed moments imply equal laws.

\phantomsection\label{notdef:phase-coordinates}
\subsubsection{Continuum evaluation coordinates}

Choose countable dense test sets as follows. Let $\mathcal Q$ contain all finite
unions of closed boxes with rational corners. The space of $C^1$ functions
vanishing off $Q\in\mathcal Q$, with its $C^1$ norm, is separable: embed it
by its value and four first derivatives in a finite product of separable
continuous-function spaces. Choose a countable dense set in each space and
enlarge their union by the countable group of $L$-adic translations and by
$W_4$. Let $\mathfrak T^\circ$ consist of the separated tuples from this union.
It is countable. Every separated tuple can be approximated by such tuples
with supports in fixed separated containers: cover each original support by
a finite union of sufficiently small rational boxes inside its chosen open
container and use the dense set for that union.

For each $X\in\mathfrak T^\circ$ choose once and for all a moment bound $M_X$
from the preceding estimate. Define the compact metrizable product

\begin{gather*}
\FitMath{\mathsf K=\prod_{X\in\mathfrak T^\circ}[-M_X,M_X],} \\[3pt]
\FitMath{e_j(\rho)=(\rho(X_j))_{X\in\mathfrak T^\circ},} \\[3pt]
\FitMath{\lambda_j=(e_j)_*D_j.}
\end{gather*}

The evaluation map is continuous for each fixed $j$. In particular
$\lambda_j$ is a well-defined probability measure on $\mathsf K$.

\begin{proposition}[Conditional canonical decomposition]
Under G1, the finite-spacing law G3, the estimates G4, and the stated
convergence of $\nu_j(X_j)$, the measures $\lambda_j$ converge weakly to one
probability measure $\lambda$. Put $\mathsf X=\operatorname{supp}\lambda$.
Every $\xi\in\mathsf X$ defines a unique functional $\sigma_\xi$ on
$\mathfrak T$, obtained as a limit along pure states
$\rho_j\in\mathcal E_j$ centered at $\nu_j$. The maps
$\xi\mapsto\sigma_\xi(X)$ are continuous and separate points. Furthermore

\begin{gather*}
\FitMath{\sigma(X)=\int_{\mathsf X}\sigma_\xi(X)\,d\lambda(\xi).}
\end{gather*}

Their joint moments are the successive far-translation limits of $\sigma$.
More precisely, for $q\ge2$,

\begin{gather*}
\FitMath{\int_{\mathsf X}\prod_{i=1}^q\sigma_\xi(X^i)\,d\lambda(\xi)
 =\lim_{|v_q|\to\infty}\cdots\lim_{|v_2|\to\infty}
   \sigma(X^1\cdot\tau_{v_2}X^2\cdots\tau_{v_q}X^q).}
\end{gather*}

The innermost limit is taken first. At each stage the remaining groups must
have disjoint supports; every sufficiently distant next translation is then
admitted. Each of these intermediate limits exists.
These properties determine the compact space, measure and functionals uniquely
among compact metrizable full-support representations with continuous
point-separating evaluations and the displayed mixed-moment property, up to
the homeomorphism preserving all evaluations. Each component clusters
with a bound $C(X,Y)|v|^{-2}$ for sufficiently distant separated products,
where $C(X,Y)$ and the distance threshold are uniform in $\xi$.
\end{proposition}

\begin{proof}
The crucial uniformity comes from reflection sizes rather than supremum norms
of $X_j$. Choose an integer physical half-side $\varpi$ so that all supports
under consideration lie in $(-\varpi+1,\varpi-1)^4$. This also leaves room
for the corners of plaquettes whose barycenters are in those supports. Then
$r_j=\varpi/a_j$ is an integer. Each expression defining
$\mathcal N_{\rho,r_j}(X_j)$ is the moment of two separated reflected copies
of $X$. G4 bounds it uniformly in $j$ and $\rho$ by a constant determined by
$M_{2|X|}$ and the test norms. For a displacement $v$ common to the lattices,

\begin{gather*}
\FitMath{\frac{r_j}{|v|_\infty/a_j-2r_j}
 =\frac{\varpi}{|v|_\infty-2\varpi}.}
\end{gather*}

Thus G3 gives, in every pure state, factorization of the two separated products
with error at most $C(X,Y)|v|^{-2}$, uniformly in the cutoff. The norms of the
individual representatives, which can diverge, do not enter this bound.

For finitely many coordinates $X^1,\ldots,X^q$, integrate this factorization
over $D_j$. Choose $v_q$ first, then $v_{q-1}$, down to $v_2$, with each
new vector sufficiently far from the already fixed union. In the factorization
at that stage the enclosing support box and its constant may depend on the
fixed outer translations, but not on $j$. Repeating the estimate and multiplying
errors only by the uniform individual moment bounds makes the total error
smaller than any prescribed $\epsilon$, uniformly in $j$. At that one fixed
finite configuration the remaining single moment converges by the hypothesis
on $\nu_j$. The mixed coordinate moments of $\lambda_j$ are therefore Cauchy.
Their precise iterated-limit representation will follow after the components
and their uniform clustering have been constructed.
Polynomials in finitely many coordinates separate points of $\mathsf K$ and
contain constants. By Stone--Weierstrass they are uniformly dense in its
continuous functions. Compactness therefore supplies a weak limit, and the
convergent mixed moments make that limit unique along the full sequence.

If $\xi$ belongs to its support, every neighborhood of $\xi$ has positive
$\lambda$ measure. Portmanteau implies that its preimage has positive $D_j$
measure for all sufficiently large $j$. Use a nested countable neighborhood
base and a diagonal choice to obtain $\rho_j\in\mathcal E_j$ with
$e_j(\rho_j)\to\xi$. The uniform telescoping estimate extends convergence
from $\mathfrak T^\circ$ to every tuple, giving $\sigma_\xi$. It also proves
independence from the choice of $\rho_j$ and continuity in $\xi$, by uniformly
approximating a given evaluation with dense-set evaluations. Multilinearity,
symmetry and the moment bounds pass to this limit.

Barycenters and higher mixed moments pass to the limit first on the countable
set and then on all tuples by the same uniform approximation. The evaluations
separate points because their countable subfamily consists of the coordinates
of $\mathsf K$. If another compact full-support representation has the same
mixed moments, embed both spaces by \emph{all} their tuple evaluations into
$\mathbb R^{\mathfrak T}$. Their images have compact union. Polynomials in
finitely many of these coordinates are dense in the continuous functions on
that union, again by Stone--Weierstrass. Equality of mixed moments therefore
gives equality of the two image measures. Each compact image equals the
support of its image measure. The images consequently coincide, and
injectivity and compactness identify the two spaces by a homeomorphism
preserving the measure and every evaluation. This argument does not presume
that a countable subfamily separates an arbitrary competing representation.

Finally pass the pure-state factorization estimate to the limit for $L$-adic
translations. Continuity of $C^1$ tests in the norm $\|\cdot\|_\sharp$ extends
it to real translations, preserving the uniform constants after a fixed
enlargement of the support margin. This proves the asserted clustering.
Finally hold $v_3,\ldots,v_q$ fixed and let $|v_2|$ tend to infinity in the
barycenter identity. Uniform component clustering separates $X^2$ from the
fixed union of all the other groups. The uniform moment bounds justify passing
this limit through the integral. Then let $|v_3|$ tend to infinity, and continue
in that order. This proves existence of every intermediate limit and the
displayed mixed-moment formula, with the innermost limit taken first.
\end{proof}


%
%
\paragraph{Plaquette long-range order and the phase means.}
Let \(\mathcal I_j\) be the invariant sigma-field of
the lattice state \(\nu_j\). For a pure state \(\rho\)
define its rescaled plaquette mean difference by
\[
 C_j(\rho)=6a_j^{-4}\bigl[\rho(A_p)-\nu_j(A_p)\bigr].
\]
Choose a real singleton \(f_*\) from the determining
family with \(\int f_*\ne0\). By \eqref{cc:Riemann-mean},
its coordinate in \(e_j(\rho)\) is \(R_j(f_*)C_j(\rho)\),
with \(R_j(f_*)\to\int f_*\). The coordinate is uniformly
bounded. Weak convergence of \(\lambda_j\) therefore gives
convergence of \(\int C_j^2\,dD_j\) to
\(\int\mathfrak c_\xi^2\,d\lambda(\xi)\), where
\(\sigma_\xi((f))=\mathfrak c_\xi\int f\).
The barycenter has zero first moment, so
\(\int\mathfrak c_\xi\,d\lambda=0\).

The finite-spacing conditional law in G3 gives
\(\rho(A_p)=E_{\nu_j}[A_p\mid\mathcal I_j]\) under
the ergodic decomposition. Hence
\[
 \kappa_S:=\int\mathfrak c_\xi^2\,d\lambda(\xi)
   =\lim_j 36a_j^{-8}
       \operatorname{Var}_{\nu_j}
          \bigl(E_{\nu_j}[A_p\mid\mathcal I_j]\bigr).
\]
The lattice plateau estimate identifies the unscaled
variance with the far-translation limit of
\(\nu_j(A_p\,\tau_v A_p)-\nu_j(A_p)^2\) at that
fixed spacing. For separated \(f,h\), take the
barycenter of the two-point moment and subtract the
product of the means within each phase. This gives
\[
 S_2(f,h)=\int S_2^\xi(f,h)\,d\lambda(\xi)
                  +\kappa_S\int f\int h .
\]
In the zero-energy representation, the projection
of a singleton with nonzero integral has a nonconstant
phase coordinate when \(\kappa_S>0\). Its variance
tests the singleton sector of the vacuum criterion.

%
For separated \(C^1\) tests in one ball of radius
\(5\mathsf d\) with
\(\operatorname{Lip}f_i\le\|f_i\|_\infty/(\vartheta\mathsf d)\),
take the realizing states of two phases at the
same cutoffs and apply the short-distance
state comparison. Higher cumulants are invariant
under the centering shifts. At order one use the
same centering state for both means. On taking
limits the oscillation across phases of the cumulant
of a subtuple \(B\) is at most
\[
 C_{|B|,\vartheta}\mathsf d^{4-|B|}
                        \prod_{i\in B}\|f_i\|_\infty .
\]
The same comparison includes the mixture by using
its defining states. Apply \eqref{phc:total-four}
with the phase law \(\lambda\). In each outer
cumulant with at least two blocks subtract a
fixed reference phase value from every argument.
The bounded-variable cumulant estimate gives
\(O(\mathsf d^{4|\pi|-4}\prod_i\|f_i\|_\infty)\)
for a partition \(\pi\). Hence
\[
 \left|S_4(f_1,\ldots,f_4)
      -\int S_4^\xi(f_1,\ldots,f_4)\,d\lambda(\xi)\right|
 \le C_\vartheta\mathsf d^4\prod_i\|f_i\|_\infty .
\]
Where the interacting lower bound of order
\(\bar x_s^{-4}\prod_i\|f_i\|_\infty\) holds for
the mixture, \(\mathsf d^4\bar x_s^4\to0\) makes
the phase average positive for sufficiently small
\(\mathsf d\). By continuity the phases with
positive fourth cumulant on that tuple then
form an open set of positive \(\lambda\)-measure.

\paragraph{Isometries of the decomposition.}
\hypertarget{phase-isometries}{}
Let \(r(x)=Rx+c\) be an isometry preserving \(\sigma\).
The family \(\sigma_\xi\circ r^{-1}\) has the same mixed moments
as the original family. Indeed,
\(r^{-1}\tau_v=\tau_{R^{-1}v}r^{-1}\) on tuples, and the far-translation
limits above allow all directions of approach to infinity. Apply the
mixed-moment formula to the tuples \(r^{-1}X^i\) with translations
\(R^{-1}v_i\), then use the invariance of \(\sigma\).
For one evaluation use the barycenter formula.
The transformed evaluations are continuous and separate points
since \(r\) is a bijection of the test tuples. Uniqueness gives
a homeomorphism \(h_r\) with \((h_r)_*\lambda=\lambda\) and
\(\sigma_{h_r(\xi)}=\sigma_\xi\circ r^{-1}\). Uniqueness of the
point with those evaluations also gives \(h_{rs}=h_rh_s\).

Suppose now that \(\sigma\) is \(\operatorname{O}(4)\)-invariant.
Each realized component is fixed by the coordinate reflection
\(\theta_1\) by the lattice symmetry of its realizing states and
the continuum passage in G1. The centering state has that symmetry
as well. For \(R\in\operatorname{O}(4)\),
\(\sigma_\xi\circ R=\sigma_{h_{R^{-1}}(\xi)}\) is another such
component. It follows that
\[
 \sigma_\xi(R\theta_1R^{-1}X)
   =\sigma_{h_{R^{-1}}(\xi)}(\theta_1R^{-1}X)
   =\sigma_{h_{R^{-1}}(\xi)}(R^{-1}X)
   =\sigma_\xi(X).
\]
These hyperplane reflections generate \(\operatorname{O}(4)\).
Thus every component is invariant under this group. Invariance
under translations and rotations also transfers coordinate
reflection positivity to each affine mirror by conjugation.

% END MODULE chapters/phase_interface_v2.tex

% BEGIN MODULE chapters/phase_vacuum_v3.tex
\subsection{The invariant Hilbert space of scalar phases}
\label{ch:phase-vacuum}

The canonical decomposition is a statement about correlation forms. A density argument below identifies its zero-energy Hilbert space
without assuming that the linear span of phase evaluations is an algebra.

Let $X$ be a compact metrizable phase space, let $\lambda$ be a Borel probability
measure of full support on $X$, and let $\sigma_\xi$ be its component
correlation forms. They act on finite separated tuples of compactly supported
test functions and their complex
linear span, including the empty tuple $1$. Assume that

\refstepcounter{equation}
\begin{gather*}
\FitMath{f_F(\xi):=\sigma_\xi(F)\in C(X),} \\[3pt]
\FitMath{\sigma(F)=\int_X f_F\,d\lambda,} \\[3pt]
\FitMath{f_1=1.}\tag{\theequation}\label{eq:phase-vacuum-PV1}
\end{gather*}

Assume that the functions $f_F$ separate points of $X$, and that the
components preserve conjugation, time reflection, and translations.
Concatenation is used only when all its supports are separated.

The quantitative phase construction supplies uniform clustering: for
any fixed separated tuples $F,G$, once their translated supports are
disjoint,

\refstepcounter{equation}
\begin{gather*}
\FitMath{\sup_{\xi\in X}
 |\sigma_\xi(F\,\tau_{te_4}G)
                   -\sigma_\xi(F)\sigma_\xi(G)|\longrightarrow0} \\[3pt]
\FitMath{(t\longrightarrow\infty).}\tag{\theequation}\label{eq:phase-vacuum-PV2}
\end{gather*}

Finite linear combinations inherit this property by summation. Uniform
separated moment estimates ensure that all functions displayed in
\eqref{eq:phase-vacuum-PV1}--\eqref{eq:phase-vacuum-PV2} are bounded, for each fixed tuple. The phases need not have
the same one-point mean.

Assume that the reflection-positive quotient has been constructed and
completed to a Hilbert space $\mathcal H$, with dense classes $[F]$ of
separated tuples whose compact supports lie strictly in $x_4>0$, and their linear combinations. Suppose
positive time translations induce a strongly continuous contraction
semigroup $e^{-tH}$ with a nonnegative self-adjoint generator $H$ and

\refstepcounter{equation}
\begin{gather*}
\FitMath{\langle[F],e^{-tH}[G]\rangle
      =\sigma(\overline{\theta F}\,\tau_{te_4}G),} \\[3pt]
\FitMath{t>0.}\tag{\theequation}\label{eq:phase-vacuum-PV3}
\end{gather*}

These are the explicit Hilbert-space hypotheses; this section does not
derive them from qualitative convergence alone. Write
$P_0=1_{\{0\}}(H)$ for the zero-energy projection.

\begin{proposition}[The zero-energy representation]
Under \eqref{eq:phase-vacuum-PV1}--\eqref{eq:phase-vacuum-PV3}, there is a unique unitary

\refstepcounter{equation}
\begin{gather*}
\FitMath{U:\ker H\longrightarrow L^2(X,\lambda)} \\[3pt]
\FitMath{\text{such that}} \\[3pt]
\FitMath{U P_0[F]=f_F.}\tag{\theequation}\label{eq:phase-vacuum-PV4}
\end{gather*}

The vacuum vector $[1]$ maps to the constant function $1$. If $X$ has
exactly $m<\infty$ points, then $\dim\ker H=m$. If $X$ is infinite,
then $\ker H$ has countably infinite Hilbert dimension.
\end{proposition}

\begin{proof}
The spectral theorem and dominated convergence give
$e^{-tH}\to P_0$ strongly as $t\to\infty$. For fixed $F,G$, use
\eqref{eq:phase-vacuum-PV1}, reflection invariance, conjugation, and \eqref{eq:phase-vacuum-PV2} in \eqref{eq:phase-vacuum-PV3} to obtain

\refstepcounter{equation}
\begin{gather*}
\FitMath{\begin{aligned}
 \langle P_0[F],P_0[G]\rangle
  &=\lim_{t\to\infty}\int_X
      \sigma_\xi(\overline{\theta F}\,\tau_{te_4}G)
                                                      \,\lambda(d\xi)\\
  &=\int_X\overline{f_F(\xi)}f_G(\xi)\,\lambda(d\xi).
 \end{aligned}}\tag{\theequation}\label{eq:phase-vacuum-PV5}
\end{gather*}

The limit is justified by uniform clustering, applied to the finite
linear combinations as well as to individual tuples. In particular,
if $P_0[F]=0$, then $f_F=0$ in $L^2$. Thus the rule in \eqref{eq:phase-vacuum-PV4} is
well defined and isometric on the span of the displayed vectors.
That span is dense in $\ker H$, since the classes $[F]$ are dense
in $\mathcal H$ and $P_0$ is a bounded projection. Its extension is
an isometry with closed range. It remains to prove surjectivity.

Let $\mathcal V$ be the complex linear span in $C(X)$ of all the
phase evaluations $f_F$ with $F$ supported at positive times.
Any separated tuple can be moved into this half-space by a positive
time translation, which leaves its phase evaluations unchanged.
Consequently $\mathcal V$ still separates points of $X$. It contains
$1$ and is closed under complex conjugation.

For positive-time tuples $F,G$, the concatenation
$F\,\tau_{te_4}G$ is a positive-time separated tuple for all sufficiently
large $t$. By \eqref{eq:phase-vacuum-PV2},

\begin{gather*}
\FitMath{f_{F\,\tau_{te_4}G}\longrightarrow f_F f_G} \\[3pt]
\FitMath{\text{uniformly on }X.}
\end{gather*}

It follows first for generators, and then for their finite linear
combinations, that products of elements of $\mathcal V$ belong to
its uniform closure. The uniform closure is itself closed under
products: approximate each of two factors uniformly by elements of
$\mathcal V$ and use boundedness of the approximants. It is therefore
a unital conjugation-closed subalgebra separating points. The complex
Stone--Weierstrass theorem makes that closure equal to $C(X)$.
Since continuous functions are dense in $L^2(X,\lambda)$, the image
in \eqref{eq:phase-vacuum-PV4} is dense and closed, hence is the whole space.

The stated vector correspondence determines the unitary on a dense
set, proving uniqueness. It sends $[1]=P_0[1]$ to $1$, since
translation fixes the empty tuple.

For a finite phase space, full support gives positive mass to
each of its $m$ points, so $L^2$ has dimension $m$. For an infinite
compact metric space, $L^2$ is separable. One way to see this is to
take a countable uniformly dense subset of $C(X)$ and then its
rational complex linear span, which is dense in $L^2$. Its dimension
cannot be finite: for every integer $m$ choose $m$ distinct points
and disjoint open neighbourhoods of them. Full support gives each
neighbourhood positive measure, so their indicator functions are
$m$ nonzero mutually orthogonal vectors. The dimension is therefore
countably infinite. This proves all assertions.
\end{proof}

In particular, a single phase gives a one-dimensional zero-energy
space under these hypotheses. Several phases can produce several
vacua, even when each component has a unique vacuum. An infinite
phase set is recorded by the extended natural count $\infty$ in
dimension statements; its set-theoretic cardinality is not equated
with Hilbert dimension. Nothing in this identification gives a
strictly positive lower spectral bound on $(\ker H)^\perp$.


% END MODULE chapters/phase_vacuum_v3.tex

% BEGIN MODULE chapters/scalar_phase_hilbert_bridge.tex
\subsection{The Hilbert realization of scalar phases}
\label{ch:scalar-phase-hilbert-bridge}

The canonical phase decomposition first specifies scalar correlation
forms.  We now identify its ordered Hilbert realization, including
its common-domain fields and zero-energy subspace.  The proof constructs the measurable Hilbert field and shows that
the mixed cyclic space fills its direct integral by proving density
of the polynomial sections.

The phase space used here is the canonical \emph{scalar} phase
space.  Its evaluations are those of the specified scalar observable.
They separate this space by its construction.  This statement does
not say that a single scalar observable distinguishes every state
of a larger gauge-observable algebra.

\subsubsection*{The actual phase hypotheses and their uniform extension}

Let \(X\) be compact metrizable, let \(\lambda\) be a Borel
probability measure of full support, and write \(S^\xi\) for
the component scalar moments.  The canonical decomposition gives,
on separated compact tuples \(F\),
\begin{equation}
 f_F(\xi)=S^\xi(F)\in C(X),\qquad
 S(F)=\int_Xf_F(\xi)\,d\lambda(\xi),\qquad f_1=1.
 \label{eq:scalar-phase-hilbert-bridge-SPH1}
\end{equation}
These evaluations separate points.  Their reality, translations
and permutation identities hold componentwise.  We use the actual
component G2 covariance and the uniform component source bounds;
they are not inferred from symmetry of the mixture.  The scalar
phase construction supplies uniform clustering of each fixed pair
of separated compact tuples after sufficiently far translations.

The quantitative scalar route has constants uniform on this state
class.  In particular, on the coincidence-flat Schwartz domain,
\begin{equation}
 \sup_{\xi\in X}|S_n^\xi(F)|
       \le A B^n(n!)^{k+32}p_{20n}(F).
 \label{eq:scalar-phase-hilbert-bridge-SPH2}
\end{equation}
The UI/PD insertion constants and the energy exponent \(N=9\)
are uniform as well.  These are the state-uniform conclusions of
the declared ordinary-source hypotheses.  They retain those
hypotheses; no new uniform bound is extracted from compactness alone.

Constrained tensor density extends \eqref{eq:scalar-phase-hilbert-bridge-SPH1} to every flat Schwartz
test: approximate in \(p_{20n}\), use \eqref{eq:scalar-phase-hilbert-bridge-SPH2} to make the
approximation error uniform in \(\xi\), and take the same limit
in the barycenter identity.  It also makes
\(\xi\mapsto S_n^\xi(F)\) continuous, as a uniform limit of
the compact-tuple evaluations.  In particular the reflected forms
\(q_\xi(F,G)\), for any fixed finite ordered Schwartz sequences,
have continuous Gram entries and satisfy
\begin{equation}
 q(F,G)=\int_Xq_\xi(F,G)\,d\lambda(\xi).
 \label{eq:scalar-phase-hilbert-bridge-SPH3}
\end{equation}
The products in these forms are in the flat domain by their
proved time-boundary vanishing.  Their degree sums are finite.

If the original determining family was defined with compact
\(C^1\) labels, smooth labels still suffice.  Choose fixed
disjoint open containers around the original supports and mollify
inside them in the \(C^1\) norm.  Telescoping the product and
using the uniform separated moment bound gives
\begin{equation}
 \sup_\xi|f_{(h_1,\ldots,h_n)}(\xi)
               -f_{(h'_1,\ldots,h'_n)}(\xi)|
             \le C_{h,\mathcal U}\sum_i\|h_i-h'_i\|_{C^1}.
 \label{eq:scalar-phase-hilbert-bridge-SPH4}
\end{equation}
The approximating labels have a common norm bound and stay in
the fixed separated containers, so the constant is uniform along
the approximation.  This proves that the continuous scalar
coordinates remain determined by smooth compact labels.

\subsubsection*{Ordered evaluations still separate the phase space}

Let \(\mathcal V_{\rm ord}\) be the complex span of evaluations
of compact smooth tuples with strictly ordered positive-time
supports, including the empty tuple.  We prove that its uniform
closure contains every evaluation in \eqref{eq:scalar-phase-hilbert-bridge-SPH1}.

Take a compact configuration set away from all coincidence
diagonals.  At a fixed configuration \((x_1,\ldots,x_n)\),
choose a unit vector \(v\) for which the numbers \(v\cdot x_i\)
are all distinct.  Such a vector exists: the finite product
\(\prod_{i<j}v\cdot(x_i-x_j)\) is a nonzero homogeneous
polynomial.  If it vanished on the whole sphere, homogeneity
would make it zero everywhere, a contradiction.  Select a proper
Euclidean rotation whose new time projection is this \(v\),
permute the slots into increasing projected-time order, and
translate all times to be strictly positive.  These strict
inequalities persist in an open neighbourhood of the configuration.

A finite subcover of the compact set and a smooth subordinate
partition split any supported test into finitely many such pieces.
Componentwise G2, permutation symmetry and translation invariance
move each piece into a compact ordered positive-time test.
The support remains off the coincidence diagonals under all
these transformations.  Ordered support-preserving Fourier
approximation expresses that piece as a Schwartz limit of finite
ordered tensors.  The formula in \eqref{eq:scalar-phase-hilbert-bridge-SPH2} makes their evaluation errors
uniform on \(X\).  Applying this to a separated tensor, and
then using \eqref{eq:scalar-phase-hilbert-bridge-SPH4} for the original \(C^1\) family, proves
\begin{equation}
 \{f_F:F\text{ in the determining separated family}\}
                  \subset\overline{\mathcal V_{\rm ord}}^{\|\cdot\|_\infty}.
 \label{eq:scalar-phase-hilbert-bridge-SPH5}
\end{equation}
Thus ordered evaluations retain point separation.  This is a
statement about scalar phase evaluations and their uniform norm,
not a rotation of the reflection-positive Hilbert plane or a
claim that the full positive-time quotient equals its ordered
sector.

Uniform time clustering gives the multiplicative part.  For two
ordered compact tuples \(F,G\), the concatenation
\(F\cdot\tau_{te_4}G\) is again ordered for all sufficiently
large \(t\), and
\begin{equation}
 f_{F\cdot\tau_{te_4}G}\longrightarrow f_Ff_G
                              \quad\hbox{uniformly on }X.
 \label{eq:scalar-phase-hilbert-bridge-SPH6}
\end{equation}
Finite linear combinations obey the same assertion by summation.
The uniform closure of \(\mathcal V_{\rm ord}\) is therefore
closed under products: first multiply finite linear combinations
using \eqref{eq:scalar-phase-hilbert-bridge-SPH6}, then approximate each of two general closure elements
uniformly and use boundedness of the approximants.  Reality gives
closure under conjugation, and the empty tuple gives the unit.
It separates points by \eqref{eq:scalar-phase-hilbert-bridge-SPH5}.  The complex Stone--Weierstrass
theorem consequently yields
\begin{equation}
 \overline{\mathcal V_{\rm ord}}^{\|\cdot\|_\infty}=C(X).
 \label{eq:scalar-phase-hilbert-bridge-SPH7}
\end{equation}
The proof supplies the algebra property; it has not assumed that
the original linear span of phase evaluations was already an
algebra.

\subsubsection*{One countable core for all component Hilbert spaces}

Write \(\mathcal K_\xi\) for the component ordered Hilbert
sector.  Choose a countable family \((F_j)_{j\ge1}\) of finite
ordered compact tensor sequences, with \(F_1=1\), whose images
are dense in every \(\mathcal K_\xi\).  Here is an explicit
way to make the choice common.  Use rational boxes with strictly
ordered positive-time intervals, fixed nested compact cutoffs,
and the countable Fourier waves on rational grids.  Finite sums
with rational complex coefficients form a countable family.
For any ordered compact test the support-preserving Fourier
construction approximates it in every required Schwartz seminorm
using such boxes and waves; rational coefficients approximate the
finitely many retained coefficients.  Infinity and ordered-boundary
cutoffs handle general ordered Schwartz tests.  The uniform bound
in \eqref{eq:scalar-phase-hilbert-bridge-SPH2}, applied to the reflected form, makes these approximations
uniform in the component Hilbert norms.

Put
\begin{equation}
 v_j(\xi)=[F_j]_\xi,\qquad
 g_{ij}(\xi)=\langle v_i(\xi),v_j(\xi)\rangle
                       =q_\xi(F_i,F_j).
 \label{eq:scalar-phase-hilbert-bridge-SPH8}
\end{equation}
The Gram entries are continuous by \eqref{eq:scalar-phase-hilbert-bridge-SPH3}.  In particular
\(\|v_j(\xi)\|\) is uniformly bounded on the compact phase
space for each fixed \(j\).  Zero vectors at some phases cause
no difficulty; no constant rank is assumed.

Construct a measurable orthonormal family by Gram--Schmidt with
zeros retained.  Recursively set
\begin{equation}
 u_j=v_j-\sum_{i<j}e_i\langle e_i,v_j\rangle,
 \qquad
 e_j=\begin{cases}u_j/\|u_j\|,&\|u_j\|>0,\\0,&\|u_j\|=0.
                      \end{cases}
 \label{eq:scalar-phase-hilbert-bridge-SPH9}
\end{equation}
Every coefficient is a Borel function of the finite Gram data
already constructed.  The set \(E_j=\{\|u_j\|>0\}\) is
Borel.  At each phase the nonzero \(e_j\) form an orthonormal
basis, because their finite spans successively recover those of
the total fundamental family \(v_j\).  This explicit procedure
handles phase-dependent ranks and possible vanishing Gram
determinants.

\subsubsection*{A concrete direct-integral construction}

A section \(\psi\) is called measurable when its coordinates
\(\psi_j(\xi)=\langle e_j(\xi),\psi(\xi)\rangle\) are
Borel after null modifications.  Its fiber norm is
\(\sum_j|\psi_j(\xi)|^2\).  Define
\begin{equation}
 \begin{gathered}
 \mathcal D=\int_X^\oplus\mathcal K_\xi\,d\lambda(\xi),\\
 \|\psi\|_{\mathcal D}^2
       =\int_X\sum_j|\psi_j(\xi)|^2d\lambda(\xi)<\infty.
 \end{gathered}
 \label{eq:scalar-phase-hilbert-bridge-SPH10}
\end{equation}
Its elements are measurable sections with finite displayed norm,
modulo equality almost everywhere.  Coordinates vanish off \(E_j\), where the
corresponding basis vector is zero.  This is the Hilbert direct
sum \(\bigoplus_jL^2(E_j,\lambda)\), so completeness is
immediate.  Its scalar product is the integral of the fiber
scalar products.  The construction is independent up to its
natural fiberwise unitary of the Gram--Schmidt presentation:
another common total measurable family has measurable transition
coefficients, and finite orthogonal projections followed by
their strong limits identify the same sections and norms.

The elementary sections \(\mathbf1_Av_j\), with \(A\) Borel,
span a dense subspace.  If a section were orthogonal to all of
them, then
\(\int_A\langle v_j(\xi),\psi(\xi)\rangle d\lambda=0\)
for every \(A,j\).  The integrands are integrable, since the
fundamental section has bounded norm and \(\psi\) is square
integrable.  They vanish almost everywhere.  Countability gives
one full-measure set for all \(j\), and totality of the fibers
then makes the section zero.

Continuous coefficients already suffice:
\begin{equation}
 \operatorname{span}\{c(\xi)v_j(\xi):c\in C(X),\ j\ge1\}
                            \text{ is dense in }\mathcal D.
 \label{eq:scalar-phase-hilbert-bridge-SPH11}
\end{equation}
Indeed a Borel probability on a compact metric space is regular.
Approximate the indicator of a Borel set in \(L^2(\lambda)\)
by continuous functions bounded by one: choose a compact subset
and an open superset with arbitrarily small measure difference,
then use a continuous function equal to one on the compact set
and zero outside the open set.  The uniform bound on
\(\|v_j\|\) transfers that approximation to the section norm.
The preceding elementary-section density proves \eqref{eq:scalar-phase-hilbert-bridge-SPH11}.

\subsubsection*{The mixed quotient embeds isometrically}

Let \(\mathcal K_{\rm mix}\) be the ordered sector of the
barycenter functional \(S\).  On its common ordered core set
\begin{equation}
 J_0[F](\xi)=[F]_\xi.
 \label{eq:scalar-phase-hilbert-bridge-SPH12}
\end{equation}
For finite linear combinations of the \(F_j\) the section is
measurable by construction.  Any other ordered test is approximated
uniformly in the component Hilbert norms by the common family.
The section obtained in this way is measurable and agrees with
the actual quotient vector in every fiber.  Its norm is integrable
by the same uniform estimate.  The formula in \eqref{eq:scalar-phase-hilbert-bridge-SPH3} gives
\begin{equation}
 \langle J_0[F],J_0[G]\rangle_{\mathcal D}
                         =q(F,G).
 \label{eq:scalar-phase-hilbert-bridge-SPH13}
\end{equation}
Thus null vectors are sent to zero and the rule extends to an
isometry \(J:\mathcal K_{\rm mix}\to\mathcal D\).
Denote its closed range by \(M\).  At this stage no onto
statement has been used.

The section \(\Omega(\xi)=\Omega_\xi=[1]_\xi\) is the
image of the mixed vacuum and has norm one.  Each component
Hilbert space is the actual cyclic scalar space established
earlier.  Neither \(\mathcal K_{\rm mix}\) nor a fiber has
been replaced by the full positive-time quotient.

\subsubsection*{Measurable evolution and exact range invariance}

For a fixed \(s\ge0\), the component heat operator sends
\([F]_\xi\) to \([\tau_{se_4}F]_\xi\).  Its images of
the common core are measurable, by uniform test approximation.
Its norm is at most one.  It therefore defines a bounded
decomposable operator
\(T_s^\oplus\psi(\xi)=T_s^\xi\psi(\xi)\) on
\(\mathcal D\).  The same reasoning applies to both signs
of spatial translations.  Fiberwise strong continuity and the
bound \(2\|\psi(\xi)\|\) give strong continuity on
\(\mathcal D\) by dominated convergence.  They obey their
group and semigroup laws pointwise and hence on this space.

The exact common-core translation identity gives
\begin{equation}
 JT_s^{\rm mix}=T_s^\oplus J,\qquad
 JU_y^{\rm mix}=U_y^\oplus J.
 \label{eq:scalar-phase-hilbert-bridge-SPH14}
\end{equation}
Thus \(M\) is invariant under the heat operators and both
spatial translation signs.  Since each heat operator is
self-adjoint, invariance also makes \(M^\perp\) invariant.
Thus \(M\) reduces them and their spectral projections, including
the strong limit of the heat semigroup as \(s\to\infty\).

The insertion family is measurable as well.  For a narrow smooth
label \(k\), its actual core vector identity is
\begin{equation}
 \Phi_s^\xi(k)[F]_\xi
      =[(\tau_{se_4}k)\cdot\tau_{2se_4}F]_\xi.
 \label{eq:scalar-phase-hilbert-bridge-SPH15}
\end{equation}
The right side is a common ordered test when the central temporal
support is sufficiently narrow.  Its Gram entries are continuous
in \(\xi\).  The operator-norm approximation to the sharp
spatial insertion has error
\(Cs^{-N-1}\epsilon\|h\|_1\), uniformly in the phase.
Hence the sharp operators \(\Phi_s^\xi(h)\) form a measurable
bounded family with common norm \(Cs^{-N}\|h\|_1\) for
small \(s\), and the safe extended bound for larger \(s\).

\eqref{eq:scalar-phase-hilbert-bridge-SPH15} proves range invariance at the vector level.
Before taking a sharp limit, its right-hand test is independent
of the phase and is exactly the corresponding mixed test.  Thus
the decomposable insertion sends \(J[F]\) to the image under
\(J\) of the actual mixed inserted vector.  Pass to the
uniform norm sharp limit, and then extend from the dense core.
We obtain
\begin{equation}
 J\Phi_s^{\rm mix}(h)
                 =(\Phi_s(h))^\oplus J.
 \label{eq:scalar-phase-hilbert-bridge-SPH16}
\end{equation}
The conjugate insertion has the same property, so \(M\) reduces
these bounded operators.  A calculation of matrix elements only
between vectors in \(M\) would not by itself have proved this
invariance; the actual common test in \eqref{eq:scalar-phase-hilbert-bridge-SPH15} is the needed step.
Spatial approximate identities give the same assertion for the
point operator, with uniform norm convergence from the analytic
PD derivative bound.

\subsubsection*{Component zero energy is exactly the vacuum line}

Uniform compact-tuple component clustering first gives the
rank-one spatial matrix-element limit on the ordered compact
core.  The unitarity approximation argument extends it to all
vectors of \(\mathcal K_\xi\) for each fixed phase.
If \(\psi\in\ker H_\xi\), the cube cone already implies
that its spatial momentum is zero.  Indeed the joint spectral
carrier on \(E=0\) has \(p=0\), so \(U_y^\xi\psi=\psi\)
for every spatial \(y\).  Apply the rank-one cluster limit to
its own matrix element.  It gives
\(\|\psi\|^2=|\langle\Omega_\xi,\psi\rangle|^2\),
which forces \(\psi\) onto the vacuum line.  Conversely that
line has energy zero.  Therefore
\begin{equation}
 \ker H_\xi=\mathbb C\Omega_\xi,
 \qquad P_0^\xi[F]_\xi=f_F(\xi)\Omega_\xi.
 \label{eq:scalar-phase-hilbert-bridge-SPH17}
\end{equation}
This proof uses component clustering.  It does not apply that
rank-one assertion to a nontrivial mixture.

Strong spectral convergence gives
\(T_s^\xi\to P_0^\xi\) as \(s\to\infty\).
Dominated convergence of the squared section norms makes
\(T_s^\oplus\to(P_0)^\oplus\) strongly on \(\mathcal D\).
Since \(M\) reduces the heat family, it contains the projection
of every one of its vectors.  For an ordered compact \(F\),
\begin{equation}
 JP_0^{\rm mix}[F]
                      =f_F(\xi)\Omega_\xi\in M.
 \label{eq:scalar-phase-hilbert-bridge-SPH18}
\end{equation}
\eqref{eq:scalar-phase-hilbert-bridge-SPH7} and \(\|\Omega_\xi\|=1\) now give
\(c(\xi)\Omega_\xi\in M\) for every \(c\in C(X)\).
This is the supply of phase coefficients needed for surjectivity.

\subsubsection*{Bounded sewn words make the range onto}

Fix a compact ordered test.  Its times and adjacent gaps have a
strictly positive common margin.  Choose the auxiliary heat
parameters below a quarter of that margin.  Write the configuration
as \(X=((\mathbf x_i,t_i))_{i=1}^n\).
At every point of its compact support its exact vector is a bounded sewn word
\begin{equation}
 \begin{gathered}
 B_\xi(X)\Omega_\xi
   =T_{t_1-a_1}^\xi\Phi_{a_1}^\xi(\mathbf x_1)
                      \mathcal Q_\xi(X)\Omega_\xi,\\
 \mathcal Q_\xi(X)=\prod_{j=1}^{n-1}
       \left[T_{t_{j+1}-t_j-a_j-a_{j+1}}^\xi
                    \Phi_{a_{j+1}}^\xi(\mathbf x_{j+1})\right].
 \end{gathered}
 \label{eq:scalar-phase-hilbert-bridge-SPH19}
\end{equation}
The product is ordered from left to right with increasing \(j\),
and is identity when \(n=1\).
The operator norms are bounded uniformly in the phase and in the
support variables.  The positive heat margins also give norm
continuity of these words uniformly in the phase: for heat use
the bound on \(He^{-\delta H}\); for spatial insertions use
the uniform PD Cauchy derivative.  Thus their decomposable
operators have a norm-continuous compactly parameterized family.

Every factor preserves \(M\) by \eqref{eq:scalar-phase-hilbert-bridge-SPH14}--\eqref{eq:scalar-phase-hilbert-bridge-SPH16}.  Applying a word
to \(c\Omega\in M\) consequently gives the section
\(c(\xi)B_\xi(X)\Omega_\xi\in M\).  Integrate against the
fixed compact test.  The uniform norm bound permits a Bochner
integral, approximated in norm by finite sums.  The range is
closed, so the result is in \(M\).  Exact ordered sewing in
each fiber identifies that result with
\(c(\xi)[F]_\xi\).  For the fundamental core this proves
\begin{equation}
 c(\xi)v_j(\xi)\in M
                         \quad(c\in C(X),\ j\ge1).
 \label{eq:scalar-phase-hilbert-bridge-SPH20}
\end{equation}
One may use the ordered Schwartz Bochner estimate to make the
same assertion for all ordered tests, but the compact fundamental
family already suffices.

These sections are dense in \(\mathcal D\) by \eqref{eq:scalar-phase-hilbert-bridge-SPH11}.
Since \(M\) is closed, we have proved the onto identification
\begin{equation}
 J:\mathcal K_{\rm mix}
      \xrightarrow{\ \cong\ }
          \int_X^\oplus\mathcal K_\xi\,d\lambda(\xi).
 \label{eq:scalar-phase-hilbert-bridge-SPH21}
\end{equation}
No density of constant sections was assumed.  The zero-energy
projection first produced nonconstant phase coefficients; actual
bounded insertions and sewing then propagated them to a fiberwise
total family.  This proves precisely why the complete direct
integral is obtained for this scalar phase construction.

\subsubsection*{Exact generator domains in the direct integral}

The heat family on \(\mathcal D\) is a strongly continuous
self-adjoint contraction semigroup, hence has a nonnegative
self-adjoint generator \(H^\oplus\).  Its resolvent is
\begin{equation}
 (1+H^\oplus)^{-1}
       =\int_0^\infty e^{-s}T_s^\oplus\,ds,
 \qquad
 ((1+H^\oplus)^{-1}\psi)(\xi)
                   =(1+H_\xi)^{-1}\psi(\xi).
 \label{eq:scalar-phase-hilbert-bridge-SPH22}
\end{equation}
The integral is strongly convergent, and fiberwise integration
follows from the integrable bound \(e^{-s}\|\psi(\xi)\|\).
The resolvent identifies the generator domain as follows.

If \(\psi\) is in that domain, write it as the resolvent
applied to \(f\in\mathcal D\).  In each fiber it then lies
in \(\operatorname{Dom}H_\xi\), and
\(H_\xi\psi(\xi)=f(\xi)-\psi(\xi)\), a square-integrable
section.  Conversely, if the fiber domain condition holds and
that section is measurable and square integrable, put
\(f(\xi)=(1+H_\xi)\psi(\xi)\).  The pointwise resolvent
identity recovers \(\psi\), placing it in the global generator
domain.  Iterating the same argument with resolvent powers gives
\begin{equation}
 \begin{split}
 \operatorname{Dom}(1+H^\oplus)^l
   =\{\psi: &\psi(\xi)\in\operatorname{Dom}(1+H_\xi)^l
                                      \text{ a.e.},\\
             &\int_X\|(1+H_\xi)^l\psi(\xi)\|^2d\lambda<\infty\}.
 \end{split}
 \label{eq:scalar-phase-hilbert-bridge-SPH23}
\end{equation}
The domain includes the measurability of the displayed section;
equivalently it is the range of the measurable resolvent power.
The mixed generator is unitarily identified with this operator
by \eqref{eq:scalar-phase-hilbert-bridge-SPH14} and \eqref{eq:scalar-phase-hilbert-bridge-SPH21}, including its exact domain.

The common global energy domain requires \eqref{eq:scalar-phase-hilbert-bridge-SPH23} for \emph{every}
integer \(l\).  Membership in the fiberwise intersection alone
does not ensure these integrability conditions.  Since the
intersection is countable, a vector in the global domain is in
the fiber common domain on one common full-measure set.  This
distinction is retained in every field formula below.

Spatial unitary groups are also decomposable by \eqref{eq:scalar-phase-hilbert-bridge-SPH14}.
Stone's theorem identifies their generator domains in the same
way.  Their groups commute with the heat semigroup, so their
joint spectral calculus is the componentwise one.  One can see
this first for products of bounded resolvents and Fourier--Laplace
transforms, then extend to bounded Borel functions by bounded
pointwise spectral approximation and strong limits.  In particular
all band projections and bounded inverse-heat factors on a band
are measurable decomposable operators.

\subsubsection*{Measurable band fields and the global common domain}

For fixed \(m,n,s>0\), the component block
\begin{equation}
 A_{mn}^\xi(h)=(e^{sH_\xi}P_m^\xi)
             \Phi_s^\xi(h)(P_n^\xi e^{sH_\xi})
 \label{eq:scalar-phase-hilbert-bridge-SPH24}
\end{equation}
is a measurable bounded operator field: its outer factors are
bounded Borel functions of the measurable energy generator and
its middle factor was constructed above.  Its norm has the
uniform bound \(C(1+m+n)^N\|h\|_1\), by the same small-heat
optimization.  Time smearing is a norm integral between the
fixed bands.  Approximation by simple temporal sums, or scalar
dominated integration of its matrix entries, makes the resulting
blocks measurable.  The weighted block-Schur estimate has common
constants over all phases.

It follows that
\begin{equation}
 L_g^\xi=\phi_\xi(g)(1+H_\xi)^{-N}
 \quad\hbox{is measurable and}\quad
       \sup_\xi\|L_g^\xi\|\le C p_{N+7}(g).
 \label{eq:scalar-phase-hilbert-bridge-SPH25}
\end{equation}
To justify measurability of the full matrix, first retain finitely
many bands.  Those operators are measurable by the preceding
construction.  Their strong limit is the weighted field matrix:
the bounded Schur operator has those blocks, and the finite
orthogonal projections tend strongly to identity.  The uniform
operator bound passes to the limit.  This proves \eqref{eq:scalar-phase-hilbert-bridge-SPH25} without
postulating a globally unsandwiched sharp-time field.

Every step defining the blocks commutes with the unitary \eqref{eq:scalar-phase-hilbert-bridge-SPH21}:
the energy functional calculus, the actual bounded insertion,
the fixed-band inverse heat, the temporal integral and finally
the uniformly bounded strong matrix limit.  Therefore the actual
mixed field has the identity
\begin{equation}
 (J\phi_{\rm mix}(g)\psi)(\xi)
       =\phi_\xi(g)(J\psi)(\xi)
       \quad\left(\psi\in\mathcal D_\infty^{\rm mix}\right).
 \label{eq:scalar-phase-hilbert-bridge-SPH26}
\end{equation}
For its construction on \(\operatorname{Dom}(1+H^{\rm mix})^N\),
apply the bounded decomposable field \((L_g)^\oplus\) to the
global vector \((1+H^\oplus)^NJ\psi\).  The integrability
condition in \eqref{eq:scalar-phase-hilbert-bridge-SPH23} makes this defined.  The stronger uniform
weighted bounds show that it preserves the global common energy
domain.  Thus \eqref{eq:scalar-phase-hilbert-bridge-SPH26} is an exact common-domain formula, not a
claim about a maximal direct-integral domain for arbitrarily
chosen self-adjoint field extensions.

The same operator-valued Fourier construction gives this identity
for every integrated scalar polynomial.  Its graph bounds are
uniform because the input field energy estimates are uniform.
Consequently the full actual scalar polynomial representation,
with its common invariant domain, is decomposed by \eqref{eq:scalar-phase-hilbert-bridge-SPH21}.

\subsubsection*{Measurability of the reconstructed Poincare action}

We prove measurability for each component directly.  Fixed Minkowski polynomial vacuum vectors in a
component form measurable sections.  To verify it, start with
the measurable vacuum section and use finite band approximations
to each smeared field.  All are measurable.  Their products
converge in every required energy graph norm by the uniform
field bounds.  Pointwise limits of their measurable coordinates
are measurable.  The general joint-smearing statement follows
from the uniform vector Fourier extension.

Choose a countable dense family of Minkowski Schwartz labels
and their finite polynomial words.  These vectors are total in
every component by its cyclic scalar reconstruction.  Their
Gram entries are measurable by the preceding paragraph.  They
can therefore be used as another measurable fundamental family,
with a Gram--Schmidt construction exactly as in \eqref{eq:scalar-phase-hilbert-bridge-SPH9}.

For a fixed Poincare element \((a,\Lambda)\), the component
unitary sends every member of this total family to the same
word with transformed tests.  The result is a measurable
section, by the just-proved test continuity and field
measurability.  Finite orthogonal projections in that family
then prove measurability on the original fundamental sections:
expand each original section in its measurable orthonormal
coordinates, apply the unitary to a finite expansion, and take
the fiber norm limit.  The operator norm is one throughout.
This constructs the decomposable unitary
\begin{equation}
 (U^\oplus(a,\Lambda)\psi)(\xi)
                         =U_\xi(a,\Lambda)\psi(\xi).
 \label{eq:scalar-phase-hilbert-bridge-SPH27}
\end{equation}
The component group laws give its group law.  Fiberwise strong
continuity and the bound by \(2\|\psi(\xi)\|\) give strong
continuity after integration.  No continuous choice of a basis
in the fibers is necessary.

On a mixed Minkowski polynomial vacuum vector the intertwining
identity follows from \eqref{eq:scalar-phase-hilbert-bridge-SPH26}, the invariant component vacua and
the common transformed tests.  Those vectors are dense in the
mixed ordered sector, and both maps are bounded unitaries.
Thus
\begin{equation}
 JU_{\rm mix}(a,\Lambda)J^{-1}=U^\oplus(a,\Lambda).
 \label{eq:scalar-phase-hilbert-bridge-SPH28}
\end{equation}
All covariance statements here use actual component G2.  A
symmetry which only permuted components would instead require
an accompanying action on the phase variable; that different
case is not silently substituted for the componentwise
invariance assumed in this scalar route.

\subsubsection*{Phase multipliers and zero-energy multiplicity}

For \(c\in L^\infty(X,\lambda)\), multiplication of sections
defines a bounded operator \(M_c\) of norm at most
\(\|c\|_\infty\).  It commutes with the decomposable
Poincare operators and with energy functional calculus.
It preserves every global graph domain by \eqref{eq:scalar-phase-hilbert-bridge-SPH23}, and the
pointwise field identity gives
\begin{equation}
 M_c\phi(g)\psi=\phi(g)M_c\psi
                           \quad(\psi\in\mathcal D_\infty).
 \label{eq:scalar-phase-hilbert-bridge-SPH29}
\end{equation}
The same is true for integrated polynomial operators on that
domain.  This is a precise commuting multiplier structure for
the scalar representation.  It is not an assertion about the
center of an unspecified bounded local gauge-observable net.

By \eqref{eq:scalar-phase-hilbert-bridge-SPH17}, the global zero-energy sections are exactly
\(c(\xi)\Omega_\xi\) with \(c\in L^2(X,\lambda)\).
The unitary identification therefore gives
\begin{equation}
 \ker H^{\rm mix}\cong L^2(X,\lambda),\qquad
       P_0^{\rm mix}[F]\longmapsto f_F,
 \qquad\Omega^{\rm mix}\longmapsto1.
 \label{eq:scalar-phase-hilbert-bridge-SPH30}
\end{equation}
This recovers the phase evaluation rule on the actual ordered
sector.  Its onto property follows from \eqref{eq:scalar-phase-hilbert-bridge-SPH7}, or directly
from the full direct-integral identification, and does not
assume equality with the larger positive-time Hilbert space.
The same sections are invariant under the full reconstructed
Poincare group, because every component vacuum is invariant.
Conversely any Poincare-invariant vector has energy zero.

If \(X\) has exactly \(m\) points, full support gives positive
mass at every one and \(\dim\ker H^{\rm mix}=m\).  If \(X\)
is infinite, \(L^2(X,\lambda)\) is separable: continuous
functions have a countable uniformly dense family and are
dense in \(L^2\).  Its dimension is not finite, because for
any \(m\) one can choose \(m\) distinct points with disjoint
open neighbourhoods; full support gives positive measure to
each neighbourhood, whose normalized indicators are orthogonal.
The Hilbert dimension is therefore countably infinite.  It is
not the set-theoretic cardinality of the phase space.

\subsubsection*{The precise uniform spectral-gap criterion}

The direct integral also identifies exactly the uniform spectral
statement that a further estimate would have to prove.  Put
\(\Pi_m^\xi=\mathbf1_{(0,m)}(H_\xi)\) for \(m>0\).
These projections are measurable and decompose under \(J\).
One explicit verification uses the bounded resolvent
\(R_\xi=(1+H_\xi)^{-1}\).  Continuous functions of \(R_\xi\)
decompose by polynomial approximation and \eqref{eq:scalar-phase-hilbert-bridge-SPH22}.  Continuous
functions on \([0,1]\), bounded by one and increasing to the
indicator of \(((1+m)^{-1},1)\), then converge strongly to
\(\Pi_m^\xi\).  The same approximation holds for the global
resolvent, and dominated section convergence identifies the two
limits.  Hence
\begin{equation}
 \begin{gathered}
 J\mathbf1_{(0,m)}(H^{\rm mix})J^{-1}
       =\int_X^\oplus\Pi_m^\xi\,d\lambda(\xi),\\
 \mathbf1_{(0,m)}(H^{\rm mix})=0
            \ \Longleftrightarrow\ \Pi_m^\xi=0\text{ a.e.}
 \end{gathered}
 \label{eq:scalar-phase-hilbert-bridge-SPH31}
\end{equation}
For the reverse implication in the last equivalence integrate
the squared norm of a projected section.  For the forward
implication apply the zero global projection to every fundamental
section \(v_j\).  The integrals of their nonnegative squared
projected norms vanish.  Countability and fiber totality make
the component projection zero almost everywhere.

Define the extended component gap by
\(g(\xi)=\sup\{q\in\mathbb Q_{>0}:\Pi_q^\xi=0\}\),
with supremum zero for the empty set and infinity when there are
no positive energies.  It is measurable, since each zero-projection
event is the countable intersection of
\(\{\|\Pi_q^\xi v_j(\xi)\|=0\}\).
Using rational thresholds in \eqref{eq:scalar-phase-hilbert-bridge-SPH31} gives
\begin{equation}
 g_{\rm mix}=\operatorname*{ess\,inf}_{\xi\in X}g(\xi).
 \label{eq:scalar-phase-hilbert-bridge-SPH32}
\end{equation}
Thus positive component gaps almost everywhere are insufficient
unless their essential lower bound is positive.  No continuity
or positive lower bound of this measurable function follows
merely from the continuous finite Gram entries.

An elementary Hilbert-family example verifies that distinction.
Take \(X=[0,1]\) with Lebesgue measure, a vacuum vector in every
fiber, and one extra orthogonal unit vector \(e\) for \(x>0\).
There is no extra vector at \(x=0\).  Use fundamental sections
\(v_1=\Omega\) and \(v_2(x)=\sqrt{x}\,e\), with \(v_2(0)=0\),
and let \(H_xe=xe\).  Their nonconstant Gram entry is \(x\);
its heat-translated entry is \(xe^{-tx}\).  Both are continuous
at zero.  Every fiber has a one-dimensional zero-energy space.
Its gap is \(x\) for \(x>0\), and infinity at zero, while
the direct integral has gap zero.  This example establishes the
spectral point about Gram data; it is not asserted to be a
gauge-theory phase model.

The equivalent heat estimate is
\begin{equation}
 \begin{gathered}
 \mathbf1_{(0,m)}(H^{\rm mix})=0
 \quad\Longleftrightarrow\\
 \|e^{-tH^{\rm mix}}-P_0^{\rm mix}\|\le e^{-mt}
                         \quad(t\ge0).
 \end{gathered}
 \label{eq:scalar-phase-hilbert-bridge-SPH33}
\end{equation}
The forward implication is the spectral supremum on positive
energies.  Conversely a nonzero spectral projection in \((0,m)\)
has a nonzero subprojection in some compact interval strictly
below \(m\).  A unit vector in that subspace has a slower
exponential lower bound, contradicting the displayed norm
estimate.  Here \(P_0^{\rm mix}\) is the whole phase vacuum
space in \eqref{eq:scalar-phase-hilbert-bridge-SPH30}, not just the projection onto the one vector
\(\Omega^{\rm mix}\).

A dense-core version can be useful without an operator-norm
estimate as input.  Suppose one common \(m>0\) satisfies, for
every vector \(\zeta\) in a dense core, an estimate
\begin{equation}
 \begin{gathered}
 0\le\langle\zeta,(e^{-tH^{\rm mix}}-P_0^{\rm mix})\zeta\rangle
                       \le C_\zeta e^{-mt},\\
 t\text{ sufficiently large}.
 \end{gathered}
 \label{eq:scalar-phase-hilbert-bridge-SPH34}
\end{equation}
Each \(C_\zeta\) may depend on the vector.  Its positive
spectral measure cannot charge \((0,c]\) for any \(c<m\):
such a charge would give a lower bound proportional to
\(e^{-ct}\), contradicting \eqref{eq:scalar-phase-hilbert-bridge-SPH34}.  Thus the projection onto
\((0,m)\) kills the dense core and hence vanishes.  Conversely
\eqref{eq:scalar-phase-hilbert-bridge-SPH33} implies \eqref{eq:scalar-phase-hilbert-bridge-SPH34} with \(C_\zeta=\|\zeta\|^2\).
These are exact conditional criteria, not an estimate of a
positive \(m\).  The available compact-tuple clustering does
not itself supply \eqref{eq:scalar-phase-hilbert-bridge-SPH34}.

The identification concerns the scalar correlation phase space
and its scalar fields.  Extending it to a larger gauge-observable
family requires the additional observable and state-consistency
data specified separately.


% END MODULE chapters/scalar_phase_hilbert_bridge.tex

% BEGIN MODULE chapters/gauge_observable_algebra.tex
\subsection{The gauge-invariant cutoff algebra}
\label{ch:gauge-observable-algebra}

The scalar field constructed above has a precise polynomial
observable sector.  To compare it with a gauge theory one first
needs the larger gauge-invariant algebra whose observables would
have to be transported and identified.  We construct that algebra
at a cutoff, prove concrete dense families, and state the exact
measure-preservation consequence of testing them.  These are
algebraic statements about the actual bond variables.  A continuum
limit for the additional observable family is a further requirement.

\subsubsection*{Bonds, orientations and the full local gauge action}

Let \(\Gamma=(V,E)\) be a finite or countable locally finite
graph, and select one orientation \(E_+\) of its unoriented
bonds.  The hypercubic lattice and its admissible periodic
quotients are the intended cases.  Write
\begin{equation}
 K=\operatorname{SU}(2),\qquad \Omega_\Gamma=K^{E_+},
 \qquad U_{\bar e}=U_e^{-1}.
 \label{eq:gauge-observable-algebra-GIA1}
\end{equation}
For a finite composable oriented path
\(\gamma=e_1\cdots e_r\), use the ordered product
\(U_\gamma=U_{e_1}\cdots U_{e_r}\); the empty path has
holonomy identity.  With \(s(e),t(e)\) its endpoints, a site
gauge \(g=(g_v)_{v\in V}\) acts by
\begin{equation}
 (g\cdot U)_e=g_{s(e)}U_eg_{t(e)}^{-1},\qquad
 (g\cdot U)_\gamma=g_{s(\gamma)}U_\gamma
                                      g_{t(\gamma)}^{-1}.
 \label{eq:gauge-observable-algebra-GIA2}
\end{equation}
The path identity follows by cancelling every adjacent
\(g_v^{-1}g_v\).  It is compatible with reversed paths and
with insertion or deletion of immediate retracings.

For a closed path the normalized trace
\begin{equation}
 w_\gamma(U)=\tfrac12\operatorname{Tr}U_\gamma
 \label{eq:gauge-observable-algebra-GIA3}
\end{equation}
is real, lies in \([-1,1]\), and is gauge invariant.
Reality follows from the conjugate eigenvalues of an
\(\operatorname{SU}(2)\) matrix.  Cyclic changes of a loop's
basepoint and reversal do not change this trace.  The matrix
holonomy of an open path is instead gauge covariant at its two
endpoints; it is not a scalar gauge-invariant observable by itself.

The gauge group \(\mathcal G=K^V\) has its compact product
topology.  Its action on \(\Omega_\Gamma\) is continuous:
each output bond depends continuously on only that bond and its
two site variables.  For countable sets, both product spaces
are compact metrizable.  An action of all site gauges is meant
here.  Invariance under gauges fixing selected centers is a
different, generally weaker condition.

\subsubsection*{Continuous local functions and their uniform completion}

For finite \(F\subset E_+\), let
\(\mathcal A_F=C(K^F)\), embedded into
\(C(\Omega_\Gamma)\) by ignoring the remaining bonds.
If \(F\subset F'\), this is an injective unital star map
and preserves the supremum norm.  Any two local functions
belong to the algebra for their finite union of bonds.
The union is therefore a unital star algebra.

In the infinite case enumerate the bonds and choose the bounded
bi-invariant matrix metric \(d_K\le1\).  The product metric
\(d_\Omega(U,U')=\sum_j2^{-j}d_K(U_{e_j},U'_{e_j})\)
is compatible with the topology.  Fix a reference configuration.
Replacing all coordinates after the first \(N\) by their
reference values changes this metric by at most \(2^{-N}\).
For \(f\in C(\Omega_\Gamma)\), uniform continuity shows
that the corresponding cylinder functions converge uniformly
to \(f\).  Hence
\begin{equation}
 \mathcal A=\overline{\bigcup_{F\Subset E_+}\mathcal A_F}
                                  ^{\|\cdot\|_\infty}
                      =C(\Omega_\Gamma).
 \label{eq:gauge-observable-algebra-GIA4}
\end{equation}
This argument gives the concrete meaning of the quasilocal
completion in the compact bond model.  It does not include
an additional charged field at infinity or a choice of boundary
frame for an open path.

Let \(S(F)\) be the finite set of endpoints of the bonds in
\(F\).  A local function is invariant under every full site
gauge exactly when it is invariant under \(K^{S(F)}\), since
no other gauge coordinate enters its value.  Define
\begin{equation}
 \mathcal A_F^{\rm inv}
       =C(K^F)^{K^{S(F)}},\qquad
 \mathcal A^{\rm inv}
       =\{f\in C(\Omega_\Gamma):f(g\cdot U)=f(U)\}.
 \label{eq:gauge-observable-algebra-GIA5}
\end{equation}
These are unital closed star subalgebras.  The local embeddings
preserve invariance.  Finite products of Wilson traces belong
to the appropriate local invariant algebra.

\subsubsection*{Finite Haar averaging and the actual projection}

For \(f\in\mathcal A_F\), set
\begin{equation}
 (\mathcal E f)(U)=\int_{K^{S(F)}}f(g\cdot U)\,dg,
 \label{eq:gauge-observable-algebra-GIA6}
\end{equation}
with normalized product Haar measure and gauges outside \(S(F)\)
equal to identity.  The result is continuous and depends only
on \(F\).  If a larger endpoint set is used the additional
integrations do not change the integrand.  If the same cylinder
function is represented in a larger bond algebra, the answer
is again the same function.  Thus \eqref{eq:gauge-observable-algebra-GIA6} is one well-defined
linear map on the local union.

Haar invariance makes its range gauge invariant.  It is positive,
preserves the unit and star, and satisfies
\begin{equation}
 \|\mathcal E f\|_\infty\le\|f\|_\infty,
 \qquad \mathcal E^2=\mathcal E,
 \qquad \mathcal E b=b\quad(b\text{ invariant}).
 \label{eq:gauge-observable-algebra-GIA7}
\end{equation}
For idempotence, average twice and combine the two site gauges;
the normalized product Haar integral is unchanged.  All the
integrals are on finite compact groups at this stage.
The norm bound extends \(\mathcal E\) uniquely to
\(C(\Omega_\Gamma)\), still positive, unital and contractive.

The extended image is precisely \(\mathcal A^{\rm inv}\).
Images of local approximants are invariant, so their uniform
limits are invariant.  Conversely, if \(b\) is invariant
and \(f_N\to b\) are cylinder approximants, then
\(\|\mathcal E f_N-b\|_\infty\le\|f_N-b\|_\infty\):
inside \eqref{eq:gauge-observable-algebra-GIA6}, replace \(f_N(g\cdot U)\) by
\(b(g\cdot U)=b(U)\).  Thus \(\mathcal E b=b\).
In particular
\begin{equation}
 \mathcal A^{\rm inv}
       =\overline{\bigcup_F\mathcal A_F^{\rm inv}}
                                      ^{\|\cdot\|_\infty}.
 \label{eq:gauge-observable-algebra-GIA8}
\end{equation}
No assertion that arbitrary cylinder approximants are themselves
invariant was needed; averaging makes them so without increasing
their approximation error.

For invariant \(b,c\), the finite integral also gives
\(\mathcal E(bfc)=b(\mathcal Ef)c\), and uniform approximation
extends this to the full algebra.  The same pointwise argument
proves complete positivity: an average of positive matrices of
continuous functions remains positive when tested against every
constant coefficient vector.  Thus \(\mathcal E\) is an
actual conditional expectation with the stated range and norm.

\subsubsection*{The compact orbit space and invariant states}

The product metric can be chosen gauge invariant, since left
and right multiplication preserve the matrix metric on each
bond.  On gauge orbits define
\begin{equation}
 d_{\rm orb}([U],[V])
       =\min_{g\in\mathcal G}d_\Omega(U,g\cdot V).
 \label{eq:gauge-observable-algebra-GIA9}
\end{equation}
The minimum exists by compactness of \(\mathcal G\) and
continuity.  Symmetry follows by inversion of the gauge, and
the triangle inequality follows by composing two gauges and
using that they act by isometries.  If the distance is zero,
a minimizing gauge gives \(U=g\cdot V\), so the orbits
coincide.  Thus it is a metric.  The orbit map is continuous
because \(d_{\rm orb}([U],[V])\le d_\Omega(U,V)\).
The compact domain maps onto a Hausdorff metric space, so this
is its quotient topology.  In particular the orbit space is
compact metrizable and all orbits are closed.

Invariant continuous functions descend and lift through this
quotient.  Hence
\begin{equation}
 \mathcal A^{\rm inv}\cong C(\Omega_\Gamma/\mathcal G)
 \label{eq:gauge-observable-algebra-GIA10}
\end{equation}
isometrically as unital star algebras.  States of the invariant
algebra are therefore probability measures on this compact
orbit space.  There is no choice of a possibly discontinuous
gauge-fixed representative in this assertion.

Every state \(\omega\) of \(\mathcal A^{\rm inv}\) has the
gauge-invariant extension
\(\widehat\omega(f)=\omega(\mathcal E f)\) to \(\mathcal A\).
It is positive and normalized by \eqref{eq:gauge-observable-algebra-GIA7}.  Conversely a
gauge-invariant state \(\tau\) on \(\mathcal A\) satisfies
\(\tau(\mathcal Ef)=\tau(f)\).  For a cylinder function this
follows by integrating its state value over the finite gauge
group; the map into continuous functions is norm continuous,
so the bounded functional passes through the integral.  Uniform
approximation handles all \(f\).  Thus
\begin{equation}
 \begin{gathered}
 \mathrm{St}(\mathcal A^{\rm inv})
             \longleftrightarrow\mathrm{St}_{\mathcal G}(\mathcal A),\\
 \omega\longmapsto\omega\circ\mathcal E,\qquad
 \tau\longmapsto\tau|_{\mathcal A^{\rm inv}}.
 \end{gathered}
 \label{eq:gauge-observable-algebra-GIA11}
\end{equation}
is an affine bijection; \(\mathrm{St}\) denotes normalized positive
states and its subscript denotes gauge invariance.
By the Riesz representation theorem
these are also the corresponding gauge-invariant probability
measures on bond configurations.  In particular the \emph{full}
invariant state determines its unique gauge-invariant lift.  A
state on a proper scalar subalgebra has a different information
content, addressed below.

\subsubsection*{Gauge invariance of the actual cutoff Gibbs states}

For a plaquette \(p\), let
\(A_p=1-\frac12\operatorname{Re}\operatorname{Tr}U_p\).
It is gauge invariant by \eqref{eq:gauge-observable-algebra-GIA2}.  On a finite periodic lattice,
the Wilson density with respect to product bond Haar measure
is therefore gauge invariant, and the measure itself is
invariant by the left and right invariance of every bond
integral.  Its positive normalization does not affect the
identity.

The same fact holds for a Wilson DLR state on the infinite
locally finite graph with its local plaquette specification.
Let \(g\) be supported on finitely many sites and let \(W\)
contain every bond incident to those sites.  This is finite
by local finiteness.  All outside bonds are fixed under this
gauge.  In the conditional kernel \(\gamma_W\), change the
inside variables by that gauge.  Product Haar measure is
preserved, and each plaquette term meeting \(W\) has unchanged
trace, including the plaquettes containing fixed outside bonds.
Consequently
\begin{equation}
 \gamma_W(f\circ g)=\gamma_Wf,
 \qquad \mu(f\circ g)=\mu(\gamma_W(f\circ g))=\mu(f).
 \label{eq:gauge-observable-algebra-GIA12}
\end{equation}
The second identity is the DLR equation.  It applies to a local
continuous \(f\); its additional outside variables are simply
held fixed in the conditional integral.

For an arbitrary site gauge and a local \(f\), only the
endpoint values of the finitely many bonds in its support
matter.  Copy those values to a finitely supported gauge and
set the rest to identity.  It has exactly the same action on
\(f\), so \eqref{eq:gauge-observable-algebra-GIA12} proves invariance under the arbitrary gauge
on all local tests.  \eqref{eq:gauge-observable-algebra-GIA4} then gives invariance on
the whole continuous algebra.  Thus the cutoff Gibbs and DLR
states supply actual states of \(\mathcal A^{\rm inv}\),
not merely states of its $\operatorname{tr}(F^2)$ part.

This proof is at a fixed cutoff and does not establish joint
continuum limits for all these observables.  Gauge invariance
of a cutoff state and existence of a particular scalar scaling
limit have distinct analytical contents.

\subsubsection*{A concrete dense family of invariant polynomials}

On a finite bond set, let \(\mathcal P_F\) be the unital
star algebra generated by the matrix entries of each \(U_e\)
and their complex conjugates.  It separates points of \(K^F\)
because the defining matrix representation is faithful.  It is
closed under conjugation.  Stone--Weierstrass gives uniform
density in \(C(K^F)\).  Apply the contractive expectation:
\begin{equation}
 \overline{\mathcal E\mathcal P_F}^{\|\cdot\|_\infty}
                       =\mathcal A_F^{\rm inv}.
 \label{eq:gauge-observable-algebra-GIA13}
\end{equation}
The averaged polynomial is still a matrix-entry polynomial.
Indeed substitute \eqref{eq:gauge-observable-algebra-GIA2} in each monomial, expand its finitely
many indices, and integrate the gauge-dependent coefficients.
Only constant coefficients change.  It follows that
\(\mathcal E\mathcal P_F=\mathcal P_F\cap\mathcal A_F^{\rm inv}\),
which is itself a unital star algebra.  Its density is therefore
an algebraic density statement with a concrete averaging map,
not a formal span assumption.

The invariant tensors in this family can also be read directly.
At each site collect the fundamental, dual and conjugate
representation factors carried by the incident matrix entries
of a monomial.  They form one finite-dimensional unitary
representation \(R_v(g)\).  Haar averaging at that site is
\begin{equation}
 P_v=\int_KR_v(g)\,dg.
 \label{eq:gauge-observable-algebra-GIA14}
\end{equation}
It is an orthogonal projection onto the invariant tensors:
inversion of \(g\) proves \(P_v^*=P_v\), two Haar integrals
prove \(P_v^2=P_v\), and left invariance gives
\(R_v(h)P_v=P_v\).  Conversely it fixes every invariant
vector.  Choosing an orthonormal basis of that finite invariant
space writes \(P_v\) as a finite sum of rank-one tensor
contractions.  Taking the product over vertices shows exactly
how an averaged matrix monomial becomes a finite sum of
gauge-invariant tensor networks.

This argument uses tensor powers of the defining representation
and their invariant projections.  No classification of all
irreducible representations is required for the density proof.
In the \(\operatorname{SU}(2)\) case the next section gives a second
complete family using only closed Wilson trace words and an
explicit orbit-separation argument.

\subsubsection*{Based gauges and physical invariance}

The gauge group fixing one root in a connected graph is a
proper subgroup of the full site gauge group.  A rooted
holonomy can be invariant under the former and still transform
by conjugation under the remaining root value.  Likewise a
coordinate chart which fixes block centers may retain coarse
bond matrices with their endpoint covariance.

For example, under a gauge fixing all internal path vertices
but changing the two endpoints, an open holonomy still obeys
the two factors in \eqref{eq:gauge-observable-algebra-GIA2}.  Its individual matrix entries
therefore do not belong to \(\mathcal A^{\rm inv}\).
A gauge-fixed background or a free coordinate in a constrained
chart is useful geometric data, but its covariance must be
included when an observable is read out.  Full gauge-invariant
traces or invariant index contractions remove the remaining
endpoint action.  The distinction prevents a uniqueness
theorem modulo center-fixing gauges from being mistaken for
a complete specification of physical gauge observables.

Boundary frames can be kept deliberately, in which case the
relevant algebra is covariant under the boundary group rather
than fixed by it.  Its endpoints, transformation law and any
later invariant contraction must then be part of the observable
type.  The invariant algebra in \eqref{eq:gauge-observable-algebra-GIA5} uses no unspecified
boundary frame.

\subsubsection*{What a dense insertion test really proves}

Suppose a finite-cutoff operation gives a bounded linear readout
\(T\) from the invariant continuous algebra to another
continuous algebra, and probability states \(\mu,\nu\) on
the two sides.  If
\(\mu(P)=\nu(TP)\) for every member of a uniformly dense
invariant polynomial family, then for any invariant continuous
\(f\), choose \(P_l\to f\) uniformly and estimate
\begin{equation}
 |\mu(f)-\nu(Tf)|
       \le(1+\|T\|)\|f-P_l\|_\infty
                    +|\mu(P_l)-\nu(TP_l)|\longrightarrow0.
 \label{eq:gauge-observable-algebra-GIA15}
\end{equation}
For a positive unital Markov readout, \(\|T\|=1\); if a
different bounded norm is used it is the displayed norm which
must be checked.  Thus exactness on a concrete complete
invariant family can establish exactness on the entire
continuous gauge-observable algebra.

If the induced full measures are gauge invariant, equality of
their invariant states identifies their gauge-invariant lifts
by \eqref{eq:gauge-observable-algebra-GIA11}, hence identifies their Borel probability measures.
Their integrals then agree on all bounded Borel insertions.
For a change of variables or a probability kernel this is an
actual measure identity once its induced measure and gauge
covariance have been specified.  Without that covariance,
\eqref{eq:gauge-observable-algebra-GIA15} still states exactly the invariant continuous identity
that was proved and does not silently add a stronger lift.

Preserving only the normalization tests the single element
\(1\), not the dense family in \eqref{eq:gauge-observable-algebra-GIA13}.  Preserving scalar
$\operatorname{tr}(F^2)$ moments tests another subalgebra whose size
must be determined.  The next two sections make the finite
\(\operatorname{SU}(2)\) complete family and the scalar restriction
explicit, so that a proposed larger gauge-theory identification
can be checked against defined observables rather than inferred
from one sector's existence.


% END MODULE chapters/gauge_observable_algebra.tex

% BEGIN MODULE chapters/su2_wilson_orbit_coordinates.tex
\subsection{Complete Wilson coordinates for a finite $\operatorname{SU}(2)$ graph}
\label{ch:su2-wilson-orbit-coordinates}

We prove a concrete completeness result for the cutoff observable
algebra.  On each finite graph, finitely many Wilson traces of
words of length at most three in a fundamental loop family
separate every gauge orbit.  Their polynomials are uniformly
dense in all continuous gauge-invariant functions.  We also
describe their exact compact coordinate image.  The result is
kinematic; it does not assume a Gibbs measure or a continuum
limit of these coordinates.

\subsubsection*{Tree frames with a specified root}

First take a connected finite graph with vertex set \(V\) and
one orientation for each bond in \(E_+\).  Choose a root
\(o\) and a spanning tree \(T\).  For \(v\in V\), let
\(p_v\) be the tree path from \(o\) to \(v\), using inverse
matrices whenever a stored orientation is traversed backwards.
Set
\begin{equation}
 H_v(U)=U_{p_v},\qquad H_o=I,
 \qquad (\mathcal T U)_e=H_{s(e)}U_eH_{t(e)}^{-1}.
 \label{eq:su2-wilson-orbit-coordinates-GWO1}
\end{equation}
This is the actual site gauge with values \(g_v=H_v\).
Every tree bond becomes identity.  For a tree edge traversed
from parent \(v\) to child \(w\), the path identity is
\(H_w=H_vU_{vw}\), so the transformed bond is
\(H_vU_{vw}H_w^{-1}=I\).  The reversed stored orientation
gives the inverse identity and hence also identity.

Enumerate the non-tree bonds as \(e_1,\ldots,e_r\), where
\(r=|E_+|-|V|+1\), and put
\begin{equation}
 h_i(U)=H_{s(e_i)}U_{e_i}H_{t(e_i)}^{-1}
          =U_{p_{s(e_i)}e_i p_{t(e_i)}^{-1}}.
 \label{eq:su2-wilson-orbit-coordinates-GWO2}
\end{equation}
Thus \(h_i\) is the holonomy of a specified closed path based
at the root.  All these functions are continuous finite products
and inverses of actual bond variables.
Every tuple in \(K^r\) is obtained: set every tree bond to
identity and assign the desired tuple to the non-tree bonds.

For an arbitrary site gauge \(k\), the path transformation law
gives
\begin{equation}
 H_v(k\cdot U)=k_oH_v(U)k_v^{-1},\qquad
 h_i(k\cdot U)=k_oh_i(U)k_o^{-1}.
 \label{eq:su2-wilson-orbit-coordinates-GWO3}
\end{equation}
In particular the tuple is invariant under gauges fixing the
root.  If two configurations have the same tuple, their tree
gauges send them to the same tree-identity configuration.  The
composition of the inverse of one tree gauge and the other
is a gauge fixing the root and mapping one configuration to
the other.  This proves that based gauge orbits are exactly
the fibers of the tuple map.

Under the full site group the only residual action on this
tuple is simultaneous conjugation.  A constant site gauge
keeps all tree bonds identity and conjugates every non-tree
bond, so every such conjugation is realized.  Conversely
\eqref{eq:su2-wilson-orbit-coordinates-GWO3} shows that a full gauge transformation can do no more.
We have therefore proved a continuous bijection of compact
orbit spaces
\begin{equation}
 K^{E_+}/K^V\ \cong\ K^r/\operatorname{Ad}K.
 \label{eq:su2-wilson-orbit-coordinates-GWO4}
\end{equation}
Both quotient spaces are Hausdorff by the compact orbit metric
construction.  The continuous bijection is consequently a
homeomorphism.  A tree gauge is not a discarded formal
coordinate choice; its exact root covariance is what produces
the remaining conjugation quotient.

\subsubsection*{Several components and changing the tree}

For a graph with finitely many connected components choose a
root and spanning tree in each.  Its orbit space is the product
of the corresponding \(K^{r_c}/\operatorname{Ad}K\), with an
independent root conjugation in every component.  In particular
a component that is a tree has one orbit and contributes only
constant invariant functions.  Products of invariant functions
from the different components separate this product space.
No closed path joining distinct components is introduced.

The construction is independent of the chosen forest at the
level of observables.  In tree-identity gauge, any closed path
is a word in the non-tree bond matrices and their inverses;
tree traversals contribute identity.  New fundamental loops
from another tree are therefore continuous word functions of
the old tuple.  The converse is true by interchanging the
trees.  After quotienting by root conjugation these are the
same intrinsic orbit coordinates related by homeomorphisms.
We will prove completeness using one convenient forest,
without attributing physical meaning to its choice.

If a larger finite graph is used, the old closed paths still
define exactly the old functions on its bond configurations.
The embedding of local invariant algebras is the isometric
one proved earlier, not a prescription to identify two
different root frames.  This will make the final dense
Wilson family compatible with nested local windows.

\subsubsection*{Real quaternion coordinates and the product signs}

Let \(\sigma_1,\sigma_2,\sigma_3\) be the Pauli matrices,
with
\(\sigma_i\sigma_j=\delta_{ij}I+i\varepsilon_{ijk}\sigma_k\).
Every \(H\in\operatorname{SU}(2)\) has the unique form
\begin{equation}
 H=aI+i\,b\cdot\sigma,
 \qquad a\in\mathbb R,\quad b\in\mathbb R^3,
 \qquad a^2+|b|^2=1.
 \label{eq:su2-wilson-orbit-coordinates-GWO5}
\end{equation}
For example the usual matrix form
\(\left(\begin{smallmatrix}\alpha&\beta\\-\bar\beta&\bar\alpha
\end{smallmatrix}\right)\), \(|\alpha|^2+|\beta|^2=1\),
gives \(a=\operatorname{Re}\alpha\),
\(b_3=\operatorname{Im}\alpha\),
\(b_1=\operatorname{Im}\beta\),
\(b_2=\operatorname{Re}\beta\).
Conversely \eqref{eq:su2-wilson-orbit-coordinates-GWO5} is unitary with determinant one.  Its inverse
is \(aI-i b\cdot\sigma\).

In these conventions direct multiplication gives
\begin{equation}
 \begin{aligned}
 (aI+i b\cdot\sigma)(cI+i d\cdot\sigma)
 &=(ac-b\cdot d)I\\
 &\quad+i\,(ad+cb-b\times d)\cdot\sigma.
 \end{aligned}
 \label{eq:su2-wilson-orbit-coordinates-GWO6}
\end{equation}
The minus sign in the cross product is fixed by the two factors
of \(i\).  It is useful to retain that sign explicitly when
extracting the triple trace.

Write \(H_i=a_iI+i b_i\cdot\sigma\).  The normalized
traces of one, two and three matrices are
\begin{equation}
 \begin{aligned}
 a_i&=\tfrac12\operatorname{Tr}H_i,\\
 c_{ij}&=\tfrac12\operatorname{Tr}(H_iH_j)
                              =a_ia_j-b_i\cdot b_j,\\
 c_{ijk}&=\tfrac12\operatorname{Tr}(H_iH_jH_k)\\
 &=a_ia_ja_k-a_k b_i\cdot b_j-a_i b_j\cdot b_k
                  -a_j b_i\cdot b_k+(b_i\times b_j)\cdot b_k.
 \end{aligned}
 \label{eq:su2-wilson-orbit-coordinates-GWO7}
\end{equation}
This follows by taking the scalar part of \eqref{eq:su2-wilson-orbit-coordinates-GWO6}, then
multiplying by the third factor.  All the displayed traces
are real.  Thus the selected Wilson traces recover
\begin{equation}
 \begin{gathered}
 G_{ii}=1-a_i^2,
 \qquad G_{ij}=a_ia_j-c_{ij}=b_i\cdot b_j,\\
 D_{ijk}=c_{ijk}-a_ia_ja_k+a_kG_{ij}+a_iG_{jk}+a_jG_{ik}
                         =\det(b_i,b_j,b_k).
 \end{gathered}
 \label{eq:su2-wilson-orbit-coordinates-GWO8}
\end{equation}
Pairs with \(i<j\) and triples with \(i<j<k\) suffice;
the diagonal Gram entries are already fixed by \eqref{eq:su2-wilson-orbit-coordinates-GWO5}.

\subsubsection*{Conjugation is exactly proper spatial rotation}

Conjugation of the traceless Hermitian matrices
\(b\cdot\sigma\) by an \(\operatorname{SU}(2)\) matrix is a real
orthogonal map of \(\mathbb R^3\).  Indeed it preserves that
real three-dimensional space and its inner product
\(\frac12\operatorname{Tr}((b\cdot\sigma)(d\cdot\sigma))
 =b\cdot d\).
It has determinant one: the map depends continuously on the
connected group \(\operatorname{SU}(2)\), while its determinant
takes only the values \(\pm1\) and is one at identity.

Every proper rotation occurs.  For a unit axis \(n\) and
angle \(\theta\), choose
\(q=\cos(\theta/2)I-i\sin(\theta/2)n\cdot\sigma\).
Using the Pauli identity twice gives
\begin{equation}
 \begin{aligned}
 q(b\cdot\sigma)q^{-1}
 &=\bigl[\cos\theta\,b+\sin\theta\,(n\times b)\\
 &\qquad+(1-\cos\theta)(n\cdot b)n\bigr]\cdot\sigma.
 \end{aligned}
 \label{eq:su2-wilson-orbit-coordinates-GWO9}
\end{equation}
This is the rotation about \(n\).  Any proper real orthogonal
matrix in dimension three fixes some axis and acts as a
two-dimensional rotation on its orthogonal plane.  To see
the fixed axis, use its real eigenvalue and determinant one:
the remaining eigenvalues are either a conjugate unit pair
or real signs, which forces at least one eigenvalue one.
Its plane action has determinant one and hence is a rotation.
Thus \eqref{eq:su2-wilson-orbit-coordinates-GWO9} realizes every element of \(\operatorname{SO}(3)\).

It follows that simultaneous conjugation of the \(H_i\)
leaves the \(a_i\) fixed and rotates all the \(b_i\) by one
common proper rotation.  The completeness problem is now a
finite real linear-algebra problem, including degenerate
spans.

\subsubsection*{Gram data and orientation separate every orbit}

Suppose two tuples have equal selected data in \eqref{eq:su2-wilson-orbit-coordinates-GWO7}.  Their
scalar parts agree, and \eqref{eq:su2-wilson-orbit-coordinates-GWO8} gives the same Gram matrix and
the same oriented triple determinants.  Define on the span
of the first family
\(T(\sum_i t_i b_i)=\sum_i t_i b'_i\).
If the first sum is zero, its squared norm is \(t^TGt=0\);
the same Gram matrix makes the second sum zero.  Thus the
map is well-defined.  It preserves inner products and is an
isometry onto the second span.  Extend it by orthonormal
bases of the complements to an element of \(\operatorname{O}(3)\).

If the span has rank three, some triple determinant is
nonzero.  Its equality in the two families forces the
determinant of this isometry to be \(+1\).  If the rank is
at most two, the complement has a nonzero dimension.  Reverse
one vector in the chosen complementary image basis if
necessary; this changes the full determinant while leaving
all the prescribed images \(b_i\mapsto b'_i\) untouched.
In every rank there is therefore a proper rotation taking
the first family to the second.  Lift it by \eqref{eq:su2-wilson-orbit-coordinates-GWO9}.
We have proved
\begin{equation}
 \begin{gathered}
 \bigl(a_i,c_{ij},c_{ijk}\bigr)
     =\bigl(a'_i,c'_{ij},c'_{ijk}\bigr)\\
 \Longleftrightarrow\quad
 H'_i=qH_iq^{-1}\text{ for one }q\in\operatorname{SU}(2).
 \end{gathered}
 \label{eq:su2-wilson-orbit-coordinates-GWO10}
\end{equation}
The reverse implication is trace invariance.  The proof
includes central matrices, collinear families and coplanar
families; it has not divided by a possibly vanishing triple
determinant in these cases.

For rank three the orientation data cannot simply be omitted.
For example \(i\sigma_1,i\sigma_2,i\sigma_3\) and
\(i\sigma_1,i\sigma_2,-i\sigma_3\) have the same scalar
parts and pair Gram entries, but their triple traces have
opposite sign.  No common proper rotation maps the two
triples.  The triple coordinate is the part of the trace
family which records this distinction.

\subsubsection*{The exact compact coordinate image}

Let \(Z_r\) be the image in
\(\mathbb R^{M_r}\),
\(M_r=r+\binom r2+\binom r3\), of the selected data.  It is
compact as the continuous image of \(K^r\).  The data map
factors through the simultaneous-conjugation quotient and
is injective there by \eqref{eq:su2-wilson-orbit-coordinates-GWO10}.  Hence
\begin{equation}
 K^r/\operatorname{Ad}K\ \cong\ Z_r.
 \label{eq:su2-wilson-orbit-coordinates-GWO11}
\end{equation}
This is a homeomorphism of compact spaces, not just a
one-to-one generic-coordinate parametrization.

There is also an explicit test for membership in the image.
Given real \(a_i,c_{ij},c_{ijk}\), form the symmetric matrix
\(G\) and numbers \(D_I\) from \eqref{eq:su2-wilson-orbit-coordinates-GWO8}, with increasing
triples \(I\).  They are realizable precisely when
\begin{equation}
 \begin{gathered}
 G\ge0,\qquad \operatorname{rank}G\le3,
                   \qquad G_{ii}=1-a_i^2,\\
 D_ID_J=\det G_{I,J}\quad\hbox{for all increasing triples }I,J.
 \end{gathered}
 \label{eq:su2-wilson-orbit-coordinates-GWO12}
\end{equation}
Here \(G_{I,J}\) is the indicated three-by-three cross
Gram matrix.  When there are fewer than three indices the
second line has no conditions.

Necessity is immediate from a family of real vectors:
\(G=B^TB\) is positive with rank at most three, and
\(\det(B_I^TB_J)=\det B_I\det B_J\).
For sufficiency, the real spectral theorem supplies vectors
\(b_i\in\mathbb R^3\) with the given positive Gram
matrix.  Their lengths and the diagonal condition make
\(a_iI+i b_i\cdot\sigma\) elements of \(K\).
If the rank is below three, every three-by-three principal
Gram determinant is zero, so the conditions \(D_I^2=0\)
force all the desired triple determinants to vanish, as
they do for the reconstructed vectors.

If the rank is three, choose a triple \(I_0\) with positive
principal Gram determinant.  Its desired \(D_{I_0}\) has
the same nonzero absolute value as the determinant of those
three reconstructed vectors.  Apply one common orthogonal
reflection if necessary to arrange equality of their signs.
For every other triple \(J\), the cross-determinant
condition gives
\(D_{I_0}D_J=\det G_{I_0,J}\).  The reconstructed vectors
obey the same equality with their actual determinants.
Division by this one nonzero \(D_{I_0}\) forces all remaining
triple determinants to agree.  The formula in \eqref{eq:su2-wilson-orbit-coordinates-GWO7} now gives
every desired trace.  This proves sufficiency, including all
degenerate cases.

These are finite polynomial equalities and inequalities:
positive semidefiniteness is expressed by principal minors,
and rank at most three by the vanishing higher minors.
For a positive matrix, vanishing of all four-by-four
principal minors already suffices for that rank bound,
because four independent Gram vectors would have a positive
principal determinant.  The cross-determinant conditions
link the orientations of different triples.  Requiring only
their squared determinants would not impose one common
orientation across all indices.

\subsubsection*{Finite Wilson trace words are a complete family}

Return to the fundamental based loops \(\gamma_i\) in
\eqref{eq:su2-wilson-orbit-coordinates-GWO2}.  The coordinates in \eqref{eq:su2-wilson-orbit-coordinates-GWO7} are exactly the Wilson
traces of
\begin{equation}
 \gamma_i,\qquad
 \gamma_i\gamma_j\ (i<j),\qquad
 \gamma_i\gamma_j\gamma_k\ (i<j<k).
 \label{eq:su2-wilson-orbit-coordinates-GWO13}
\end{equation}
These are closed paths based at the same root; repetitions,
intersections and retracings are allowed.  Each is contained
in the original finite graph.  The coordinate family for
separate connected components is used componentwise and
then multiplied as scalar functions.

The unital complex polynomial algebra in these real-valued
coordinates is closed under conjugation and separates the
compact orbit space.  Stone--Weierstrass and \eqref{eq:su2-wilson-orbit-coordinates-GWO4}, \eqref{eq:su2-wilson-orbit-coordinates-GWO11}
therefore give
\begin{equation}
 \overline{\mathbb C[
       w_{\gamma_i},w_{\gamma_i\gamma_j},
       w_{\gamma_i\gamma_j\gamma_k}]}
                    ^{\|\cdot\|_\infty}
                           =C(K^{E_+})^{K^V}.
 \label{eq:su2-wilson-orbit-coordinates-GWO14}
\end{equation}
This proves uniform completeness of an actual finite
gauge-invariant observable family.  It requires neither
a generic full-rank assumption nor a separate trace-word
invariant-theory assumption.  The real linear algebra and
compactness above supplied separation in every rank.

The number of coordinates is finite, with
\(M_r\le(r+1)^3\).  Each fundamental loop has length at
most \(2(|V|-1)+1\); the longest selected word has at most
three times that length.  Thus the construction supplies
explicit finite paths within the original support window.
This is a statement about a complete test family, not a
uniform low-degree approximation rate for every continuous
observable.

\subsubsection*{Exact trace relations at the cutoff}

Some consistency relations can be written without the tree
coordinates.  For \(B\in\operatorname{SU}(2)\), its characteristic
equation gives \(B+B^{-1}=(\operatorname{Tr}B)I\).
Multiply by \(A\) and take the trace to obtain
\begin{equation}
 \begin{gathered}
 \operatorname{Tr}A\operatorname{Tr}B
          =\operatorname{Tr}(AB)+\operatorname{Tr}(AB^{-1}),\\
 w_Aw_B=\tfrac12(w_{AB}+w_{AB^{-1}}).
 \end{gathered}
 \label{eq:su2-wilson-orbit-coordinates-GWO15}
\end{equation}
For based closed paths this is an exact identity of
continuous bond functions.  Empty paths have trace one,
retracings can be deleted, reversal leaves the trace fixed,
and conjugating a loop by a connecting path leaves its trace
fixed.  These relations, together with the coordinate-image
conditions in \eqref{eq:su2-wilson-orbit-coordinates-GWO12}, give concrete finite compatibility tests
for a proposed bare loop readout.

They are identities of the bounded cutoff variables.  If a
different renormalized observable is defined by cutoff-dependent
normalizations or mixing, its transformed relations must be
derived from that definition.  The unchanged bare identities
are not automatically identities between unspecified
renormalized composite fields at coincident points.

\subsubsection*{From a finite complete family to the quasilocal algebra}

Every finite local gauge-invariant function is uniformly
approximated by the Wilson polynomials just constructed on
its own finite graph.  Combining this with \eqref{eq:gauge-observable-algebra-GIA8} gives
\begin{equation}
 \mathcal A^{\rm inv}
      =\overline{\mathbb C[w_\gamma:
               \gamma\text{ a finite closed edge path}]}
                                      ^{\|\cdot\|_\infty}.
 \label{eq:su2-wilson-orbit-coordinates-GWO16}
\end{equation}
A polynomial in finitely many traces is still local; its
support is contained in the finite union of the paths'
bonds.  Conversely no enlargement of a local graph was
needed to approximate its invariant functions.  The
statement is therefore compatible with local support and
with the isometric inclusions of finite windows.

On a countable graph the finite closed paths form a
countable family.  Polynomials with rational complex
coefficients in them are a countable uniformly dense
family of the invariant algebra.  A proposed bounded
observable transport can consequently be tested on an
explicit determining family, subject to the boundedness
and gauge-covariance conditions of \eqref{eq:gauge-observable-algebra-GIA15}.

At one finite cutoff a state of the invariant algebra is
determined by all joint polynomial moments of the selected
coordinates.  Indeed these moments give its values on
the dense polynomial algebra in \eqref{eq:su2-wilson-orbit-coordinates-GWO14}, and a state is
norm continuous.  Equivalently they determine a probability
measure on the compact image \(Z_r\).  It is the joint
moments which are required; separate one-variable laws
alone do not specify the correlations of the coordinates.

No continuum path limit has been taken in this argument.
The countability of lattice paths, exact finite graph
relations and compact bounded moment problem are the
cutoff kinematics.  Geometric rounding independence,
renormalized loop limits, their joint correlations with
the scalar field and their physical reconstruction remain
separate observable readout questions.


% END MODULE chapters/su2_wilson_orbit_coordinates.tex

% BEGIN MODULE chapters/scalar_observable_information.tex
\subsection{The scalar observable and its information}
\label{ch:scalar-observable-information}

The complete cutoff invariant algebra has now been described without
reference to a continuum limit.  We compare it with the actual scalar
observable used in the construction.  This comparison has two different
outcomes.  At a fixed cutoff the scalar algebra is a proper subalgebra,
and we exhibit the missing information explicitly.  In the continuum,
the scalar distributions do determine their own cyclic representation
and canonical field closures.  Neither statement supplies the joint
continuum law of other gauge observables.

\subsubsection*{The actual cutoff normalization}

Write $a>0$ for a fixed lattice spacing.  Count each elementary positive
coordinate-plane plaquette once, write $x_p$ for its barycenter in
physical coordinates, and let $U_p$ be its oriented holonomy.  The
observable is
\begin{equation}
 \begin{gathered}
 A_p(U)=1-\tfrac12\operatorname{ReTr}U_p\in[0,2],\\
 O_a(f)=\sum_p f(x_p)A_p,\qquad
 X_a(f)=O_a(f)-\rho_a(O_a(f)).
 \end{gathered}
 \label{eq:scalar-observable-information-GSI1}
\end{equation}
There is no power of $a$ multiplying this displayed sum.  Any proposed
different normalization or mixing is a different readout and requires
its own estimates.  For $f\in C_c^\infty(\mathbb R^4)$ the sum is finite;
for Schwartz $f$ it converges uniformly as a function of $U$, since the
plaquette barycenters form a finite union of translates of $a\mathbb
Z^4$.  In particular it is a bounded continuous invariant function at
this fixed $a$.

Let $\mathfrak A_a^{\rm inv}$ be the invariant algebra constructed
above, with the finite periodic version used on a torus.  Define
\begin{equation}
 \mathfrak A_{a,d}
  =C^*(1,A_p:\ p\text{ an elementary plaquette})
  =C^*(1,O_a(f):f\in C_c^\infty(\mathbb R^4))
  \ \subset\ \mathfrak A_a^{\rm inv}.
 \label{eq:scalar-observable-information-GSI2}
\end{equation}
For the second equality, the inclusion from right to left follows from
the finite sum in \eqref{eq:scalar-observable-information-GSI1}.  Conversely the barycenters of distinct
positive coordinate-plane plaquettes are distinct: their fractional
coordinate pattern records the plane, and their integer parts record
the translate.  Around any one barycenter choose a smooth bump which is
one there and zero at every other barycenter.  The corresponding
$O_a(f)$ is exactly $A_p$.  The same argument uses smooth periodic
functions on a finite torus.  Adding or subtracting the constants
$\rho_a(O_a(f))$ leaves the unital algebra unchanged.

This is an equality of algebras at a fixed cutoff.  The bump isolating a
single plaquette may depend on $a$.  For example
$f_{a,p}(x)=\chi((x-x_p)/(ca))$, with a fixed sufficiently small $c$,
has derivatives of order $r$ of size $O(a^{-r})$.  Consequently the
algebraic equality does not give a uniform estimate for the isolation
of a plaquette as $a\downarrow0$.  It cannot be substituted for the
fixed-test estimates used in a continuum limit.

\subsubsection*{The compact scalar quotient and its moments}

Put $Z_a=\Omega_a/\mathcal G_a$ and let $Y_a$ be the image of the
continuous map
\begin{equation}
 \pi_a:Z_a\longrightarrow[0,2]^{\mathcal P_a},\qquad
 [U]\longmapsto (A_p(U))_p.
 \label{eq:scalar-observable-information-GSI3}
\end{equation}
The countable product on the right is compact metrizable.  Its image
$Y_a$ is compact, and $\pi_a$ is a surjection onto that image.  The
unital polynomial algebra in the coordinate functions on $Y_a$ is
conjugation closed and separates its points.  Stone--Weierstrass and
the isometric pullback therefore identify
\begin{equation}
 \mathfrak A_{a,d}=\pi_a^*C(Y_a),\qquad
 \mathfrak A_a^{\rm inv}=C(Z_a).
 \label{eq:scalar-observable-information-GSI4}
\end{equation}
For an infinite graph this includes uniform limits of local scalar
functions.  It does not require that a scalar function have finite
support after taking that limit.

All joint polynomial moments of the $A_p$ determine a state on
$\mathfrak A_{a,d}$.  Indeed they determine its values on the
coordinate polynomials, and for any continuous $b$ and polynomial $P$
the difference between two states is bounded by
$2\|b-P\|_\infty$.  Taking a uniformly approximating sequence proves
uniqueness on $C(Y_a)$.  The same statement applies to all joint
moments of the $O_a(f)$ or the centered variables, provided the
centering constants are specified.  It is a statement about complete
bounded cutoff moment data.  A sequence of continuum separated-point
limits is not, by itself, this complete fixed-cutoff data set.

What is lost by $\pi_a$ can be seen without a dimension count or an
invariant-theory assumption.

\subsubsection*{An explicit small-curvature loss of relative color data}

Use the unit lattice $\mathbb Z^4$ and the Pauli matrices of \eqref{eq:su2-wilson-orbit-coordinates-GWO5}.
Let $0<\alpha<\pi/2$ and set
\begin{equation}
 A=e^{i\alpha\sigma_3},\qquad
 B_\parallel=e^{i\alpha\sigma_3},\qquad
 B_\perp=e^{i\alpha\sigma_1}.
 \label{eq:scalar-observable-information-GSI5}
\end{equation}
Consider two configurations.  In both, all positively oriented bonds
are identity except the bond $c_1=(0,e_1)$, which is $A$, and the bond
$c_2=(4e_4,4e_4+e_1)$.  The latter is $B_\parallel$ in the first
configuration and $B_\perp$ in the second.  Reversed bonds have inverse
values as always.

No elementary plaquette contains both special bonds.  Each elementary
plaquette holonomy is consequently identity, $A^{\pm1}$, or
$B^{\pm1}$.  The plaquettes meeting a special bond are the same in
the two configurations, and
\begin{equation}
 \tfrac12\operatorname{Tr}A^{\pm1}
  =\tfrac12\operatorname{Tr}B_\parallel^{\pm1}
  =\tfrac12\operatorname{Tr}B_\perp^{\pm1}
  =\cos\alpha.
 \label{eq:scalar-observable-information-GSI6}
\end{equation}
Every plaquette action is therefore equal in the two configurations.
Each nonidentity plaquette is within $O(\alpha)$ of identity in
matrix norm.  Thus the example has arbitrarily small plaquette curvature and
requires no large-field exceptional set.

Let $\gamma_1$ be the positive $12$-plaquette based at zero with first
edge $c_1$, and let $\gamma_2$ be its translate by $4e_4$, based at
$4e_4$.  Their holonomies are $A$ and $B$, respectively.  Let $P$ be
the straight path from zero to $4e_4$ along coordinate four.  All its
bonds are identity.  The closed path
$\gamma=\gamma_1P\gamma_2P^{-1}$ has holonomy $AB$, so that
\begin{equation}
 \begin{gathered}
 w_\gamma(U_\parallel)=\cos(2\alpha),\qquad
 w_\gamma(U_\perp)=\cos^2\alpha,\\
 w_\gamma(U_\parallel)-w_\gamma(U_\perp)=-\sin^2\alpha\ne0.
 \end{gathered}
 \label{eq:scalar-observable-information-GSI7}
\end{equation}
The first equality uses the same axis twice.  For the second, the
scalar part of the product in \eqref{eq:su2-wilson-orbit-coordinates-GWO6} is
$\cos^2\alpha-\sin^2\alpha\,e_3\cdot e_1=\cos^2\alpha$.
These configurations are not gauge equivalent, because the normalized
trace of every closed path is invariant under the full site gauge
group.

Let $\mathbb E_a$ be the Haar expectation of \eqref{eq:gauge-observable-algebra-GIA6}.  Evaluation at
either configuration followed by $\mathbb E_a$ is a normalized
positive, fully gauge-invariant state on $C(\Omega_a)$.  Its restriction
to $C(Z_a)$ is simply evaluation at that orbit.  The two states agree
on $\mathfrak A_{a,d}$, hence on every product of all the scalar
observables in \eqref{eq:scalar-observable-information-GSI1}, but disagree on $w_\gamma$.  This proves that
$\mathfrak A_{a,d}$ is a proper subalgebra, including in an arbitrarily
small-plaquette region.

These two states are algebraic cutoff examples.  No assertion that
they are Wilson Gibbs states, DLR states, reflection-positive states,
or scaling limits is involved.  Their role is exact: positivity,
gauge invariance, and the complete scalar marginal do not logically
determine the full invariant state.  Additional dynamical hypotheses
could select a unique extension within a smaller class; such a
selection would be a further theorem.

\subsubsection*{A separate loss of periodic holonomy data}

On a periodic four-dimensional lattice with side length $N\ge3$, put
$h_\theta=e^{i\theta\sigma_3/N}$ on every positive coordinate-one bond
and identity on every other positive bond.  All elementary plaquettes
are identity, whereas the closed path winding once around coordinate
one has
\begin{equation}
 A_p=0\quad\hbox{for every }p,\qquad
 w_{\rm wind}=\tfrac12\operatorname{Tr}h_\theta^N=\cos\theta.
 \label{eq:scalar-observable-information-GSI8}
\end{equation}
For different $\theta\in[0,\pi]$ these are different full gauge orbits
with exactly the same scalar plaquette data.  A periodic gauge
transformation changes the winding holonomy only by conjugation.  A
nonperiodic change of frame which removes $h_\theta$ on the covering
lattice is not a periodic gauge transformation and cannot identify
these torus orbits.

The complete graph coordinates of \eqref{eq:su2-wilson-orbit-coordinates-GWO13} retain such winding loops.
Discarding them would change the observable algebra.  This example
concerns different holonomy data; it does not contradict uniqueness
on a fixed nonlinear coarse fibre in the first-step background
theorem.  Together, \eqref{eq:scalar-observable-information-GSI7} and \eqref{eq:scalar-observable-information-GSI8} distinguish a local loss of
relative nonabelian orientation from a global loss of periodic
holonomy.

\subsubsection*{Exactly what abstract state extension gives}

For any state $\nu$ on $C(Y_a)$ consider the set
\begin{equation}
 \mathcal E_a(\nu)=
 \{\mu\in\operatorname{Prob}(Z_a): (\pi_a)_*\mu=\nu\}.
 \label{eq:scalar-observable-information-GSI9}
\end{equation}
We use the Riesz correspondence between states and probability
measures on compact spaces.  The set in \eqref{eq:scalar-observable-information-GSI9} is convex and weakly
closed in the compact state space of $C(Z_a)$.  It is nonempty for
every $\nu$.  Here is a proof which does not choose conditional
measures on individual fibres.

Regard $\nu$ as a norm-one unital functional on the closed subalgebra
$\pi_a^*C(Y_a)$.  The complex Hahn--Banach theorem extends it to a
functional $\ell$ on $C(Z_a)$ with
\begin{equation}
 \ell|_{\pi_a^*C(Y_a)}=\nu,\qquad
 \ell(1)=1,\qquad \|\ell\|=1.
 \label{eq:scalar-observable-information-GSI10}
\end{equation}
Such a functional is positive.  To check this directly, let $b$ be
real-valued with $\|b\|_\infty\le1$, and write $\ell(b)=u+iv$.
For every real $t$,
\begin{equation}
 |1+it\ell(b)|^2\le\|1+itb\|_\infty^2\le1+t^2.
 \label{eq:scalar-observable-information-GSI11}
\end{equation}
After expanding the left side and taking $t\to0$ with both signs we
obtain $v=0$.  Thus $\ell$ is real on real functions.  If $0\le b\le1$,
then $|\ell(b)|\le1$ and $|1-\ell(b)|\le1$, since also
$\|1-b\|_\infty\le1$.  These two real inequalities give
$0\le\ell(b)\le1$.  Scaling proves positivity for every nonnegative
$b$.  Riesz now supplies the measure $\mu$, and equality of the
restrictions is precisely the pushforward condition in \eqref{eq:scalar-observable-information-GSI9}.

Composing this extension with $\mathbb E_a$ gives a full
gauge-invariant probability law on the bond configuration space.  This
last lift is unique once $\mu$ on the \emph{full} orbit space has
been specified, by \eqref{eq:gauge-observable-algebra-GIA11}.  The Hahn--Banach step generally has no such
uniqueness.  Its conclusion concerns a positive bounded cutoff state;
it supplies no action, DLR specification, ultraviolet estimates,
reflection positivity, or continuum observable normalization.

The nonuniqueness can be located exactly for a point scalar state.
For $y\in Y_a$,
\begin{equation}
 \mathcal E_a(\delta_y)=\operatorname{Prob}(\pi_a^{-1}\{y\}),
 \qquad
 \mathcal E_a(\delta_y)\text{ is a singleton}
 \ \Longleftrightarrow\ \pi_a^{-1}\{y\}\text{ is a singleton}.
 \label{eq:scalar-observable-information-GSI12}
\end{equation}
Indeed a measure with pushforward $\delta_y$ gives mass one to the
closed fibre, and conversely any measure on that fibre has that
pushforward.  Distinct points of the fibre give distinct Dirac
measures.  The two orbit points in \eqref{eq:scalar-observable-information-GSI7} lie in one such nonsingleton
fibre.  We do not infer that every possible scalar state has multiple
extensions.

More globally, equality $\mathfrak A_{a,d}=C(Z_a)$ is equivalent to
injectivity of $\pi_a$.  If it is injective, compactness makes it a
homeomorphism onto $Y_a$, so its pullback is onto $C(Z_a)$.  Conversely,
if the pullback is all of $C(Z_a)$, equal images would force all
continuous functions to agree at two points; continuous functions on
a compact metric space separate points.  This criterion isolates the
precise algebraic reason that the scalar marginal does not determine
the full cutoff law.

\subsubsection*{Ordinary positivity does not give reflection positivity}

The distinction is already visible on a two-site classical system.
Let $s_-,s_+\in\{-1,1\}$ and give mass $1/2$ to each of
$(s_-,s_+)=(-1,1),(1,-1)$.  This is a positive normalized probability
measure invariant under the reflection exchanging the sites.  Take
the positive-time function $F=s_+$ and let $\Theta$ act by exchange
and complex conjugation.  Then
\begin{equation}
 \mu((\Theta F)F)=\mu(s_-s_+)=-1.
 \label{eq:scalar-observable-information-GSI13}
\end{equation}
Thus even reflection symmetry and ordinary probability positivity do
not imply OS positivity.  The latter is a distinct family of Gram
inequalities pairing reflected positive-time words.  An abstract
extension from \eqref{eq:scalar-observable-information-GSI10} is not an extension satisfying those
inequalities unless they too are proved.  This example is a logical
distinction between two positivity requirements, not a purported
counterexample to the Wilson construction.

\subsubsection*{The scalar continuum distributions determine a cyclic theory}

There is also a positive uniqueness statement, now for the actual
continuum scalar output.  Let two vacuum representations have the
same full Minkowski distributions $W_n$ on every
$\mathcal S((\mathbb R^4)^n)$, with their common invariant polynomial
domains generated from their vacua.  For finite scalar words define
\begin{equation}
 U\bigl(\varphi_1(f_1)\cdots\varphi_1(f_n)\Omega_1\bigr)
   =\varphi_2(f_1)\cdots\varphi_2(f_n)\Omega_2.
 \label{eq:scalar-observable-information-GSI14}
\end{equation}
The formula is extended linearly over finite sums, including the
empty word.  The squared norm of any such sum is a finite sum of
values of $W_{m+n}$ on the reversed, complex-conjugate first word and the
second word.  Equality of all distributions proves equality of these
Gram forms.  In particular a null sum on one side is null on the
other, and \eqref{eq:scalar-observable-information-GSI14} is well defined and isometric.  Its range is the
other polynomial core, so it extends to a unitary between the two
scalar cyclic Hilbert spaces.

Concatenating one further test function proves the field
intertwining on these polynomial cores.  Taking closures of the
graphs proves equality of the corresponding canonical closed
operators whenever these closures exist.  Vacuum covariance gives
intertwining of translations and proper orthochronous Lorentz
transformations first on the core and then on the Hilbert spaces.
For translations this also intertwines their joint spectral
measures, the energy operator, all its powers and their domains.
The common energy domains and canonical closures constructed in RD
are therefore part of this determined scalar representation.  This
argument does not add essential self-adjointness or a bounded local
observable net to the established conclusions.

Under their stated hypotheses, EB and QRM supply the full $W_n$ and
identify their cyclic space with $\mathcal K_{\rm ord}$.  The formula in \eqref{eq:scalar-observable-information-GSI14} thus gives uniqueness in the precise cyclic sense just
proved.  It makes no assertion that $\mathcal K_{\rm ord}$ exhausts
the entire Euclidean OS space, or that adding an independent gauge
observable leaves this cyclic space unchanged.  A larger joint
observable theory contains this scalar representation as the
closure of its scalar vacuum words; surjectivity onto the larger
space is an additional condition.

\subsubsection*{Bounded cutoff moments and unbounded continuum fields}

At a fixed cutoff one has the elementary bound
\begin{equation}
 \|O_a(f)\|_\infty\le2\sum_p|f(x_p)|,
 \qquad
 \|X_a(f)\|_\infty\le4\sum_p|f(x_p)|.
 \label{eq:scalar-observable-information-GSI15}
\end{equation}
For a fixed nonzero physical test function the right side can grow
as $a\downarrow0$.  Thus the compact moment argument for a fixed
$a$ does not pass to a common compact spectral interval in the
continuum.  Under those hypotheses, the continuum field has the graph bound,
with $W=1+H$,
\begin{equation}
 \|W^k\varphi(g)W^{-(N+k)}\|
       \le C_k\,\mathfrak t_{N+k+2}(g),\qquad N=9.
 \label{eq:scalar-observable-information-GSI16}
\end{equation}
Here $\mathfrak t_q$ is the finite spacetime test seminorm defined
in EB.  This gives a common invariant energy-smooth domain, products,
closability, and the canonical graph identifications of RD.  Repeated
use increases the energy order and the required test seminorm.  It
does not give a bound of the form $\|\varphi(g)^n\Omega\|\le C^n n!$
for every Schwartz $g$ with a fixed constant.  Such an analytic-vector
assertion has not been used to assign a spectral calculus to the
field.

A concrete scalar moment calculation shows why positivity and a
bound exponential in $n^2$ alone cannot replace compact moment
determinacy.  On $(0,\infty)$ let
\begin{equation}
 d\mu_\epsilon(x)=
  \frac{e^{-(\log x)^2/2}}{x\sqrt{2\pi}}
       \bigl(1+\epsilon\sin(2\pi\log x)\bigr)\,dx,
       \qquad -1\le\epsilon\le1.
 \label{eq:scalar-observable-information-GSI17}
\end{equation}
These are nonnegative measures.  With $z=\log x$, their sine
contribution to every nonnegative integer moment is
\begin{equation}
 \begin{split}
 &\frac1{\sqrt{2\pi}}\int_{\mathbb R}
        e^{-z^2/2+nz}\sin(2\pi z)\,dz\\
 &\hspace{1cm}=e^{(n^2-4\pi^2)/2}\sin(2\pi n)=0,
       \qquad n=0,1,2,\ldots.
 \end{split}
 \label{eq:scalar-observable-information-GSI18}
\end{equation}
For completeness the complex Gaussian identity used here follows
first for real $u$ by completing the square in
$\int e^{-z^2/2+uz}\,dz=\sqrt{2\pi}e^{u^2/2}$.
Both sides are entire in $u$ by dominated convergence on compact
complex sets, so the one-variable identity theorem permits
$u=n+2\pi i$ and taking imaginary parts.  The formula in \eqref{eq:scalar-observable-information-GSI18} at $n=0$
proves normalization.  For every $\epsilon$ the moments are
\begin{equation}
 \int x^n\,d\mu_\epsilon(x)=e^{n^2/2},\qquad n\ge0,
 \label{eq:scalar-observable-information-GSI19}
\end{equation}
while the measures differ when $\epsilon$ differs, since their
densities differ on an interval.  There are therefore bounded
continuous functions whose integrals differ, although all these
polynomial moments agree.

This calculation is not an alternative continuum model satisfying
the other hypotheses.  It proves the narrower implication needed
here: polynomial positivity and $e^{O(n^2)}$ moment growth alone do
not determine all bounded functions of an unbounded variable.
The imported OS output has an $a b^{n^2}$ type distribution bound;
no compact moment or analytic-vector conclusion has been inferred
from that form of bound.  The actual scalar Wightman functional and
its canonical cyclic graph representation remain determined by
\eqref{eq:scalar-observable-information-GSI14}.  The distinction concerns further spectral completions or
additional observables, not an ambiguity in those already
constructed distributions.

\subsubsection*{Scalar phase labels and full gauge-state labels}

The compact phase space used in SPH is separated by scalar
evaluation functions $f_F(\xi)$ of ordered Euclidean words.  Its
points distinguish those scalar functionals.  They are not, by this
definition, points of $Z_a$, nor are they automatically all extremal
states of a continuum gauge-observable algebra.  If two proposed
joint observable states have the same scalar restriction, all
$f_F$ agree for those restrictions even when a mixed or nonscalar
correlation differs.

Conversely a theorem specifying a physical class of extensions
could prove that its restriction map to scalar states is injective.
The finite-cutoff examples show that this would require the
additional class restrictions, not just gauge invariance and
positivity.  In a continuum construction the required conditions
include actual joint readouts, their compatible limits, their
positivity and covariance, and their relation to the scalar
normalization in \eqref{eq:scalar-observable-information-GSI1}.  The next section formulates these conditions
as observable statements and proves the algebraic and
representation consequences once they hold.


% END MODULE chapters/scalar_observable_information.tex

% BEGIN MODULE chapters/gauge_observable_extension_criteria.tex
\subsection{Joint gauge-observable extension and identification}
\label{ch:gauge-observable-extension-criteria}

We give precise conditions under which further observables extend the
constructed scalar theory.  They concern joint moments, not separate
one-observable laws.  We prove the resulting positivity, consistency,
domain and identification statements, and distinguish them from the
model estimates needed for a continuum readout.  The scalar
reconstruction already established in EB and QRM is retained
throughout; none of its locality conclusions is inferred from an
unconstructed larger observable family.

\subsubsection*{Typed words and the scalar restriction}

Let $I$ be a set of observable types with a distinguished scalar type
$0$.  Each type has a complex test space $E_\alpha$, an antilinear
involution $f\mapsto f^*$ into $E_{\alpha^*}$, and a specified
localization region for each local test.  For point fields take
$E_\alpha=\mathcal S(\mathbb R^4)$, with the appropriate finite tensor
components if necessary.  A loop type instead needs an explicit space
of geometric labels or smearing tests and a support convention; it is
not silently a point field.  Types are required to have compatible
involutions, so $(\alpha^*)^*=\alpha$ and $(f^*)^*=f$.

First form the unital free algebra of finite typed words, with
complex linearity in each test.  The star is
\begin{equation}
 (q_{\alpha_1}(f_1)\cdots q_{\alpha_n}(f_n))^*
   =q_{\alpha_n^*}(f_n^*)\cdots q_{\alpha_1^*}(f_1^*).
 \label{eq:gauge-observable-extension-criteria-GEX1}
\end{equation}
Let $\mathcal I$ be a declared two-sided star ideal of further
relations and let $\mathcal A$ be the quotient.  The relations may
encode exact trace identities, conjugate observables, or other
proved identifications.  A normalization-dependent continuum
identity is not included merely because the corresponding bare
cutoff variables satisfy it.  This distinction applies in particular
to \eqref{eq:su2-wilson-orbit-coordinates-GWO15}.

There is a specified unital star map $\iota$ from the scalar word
algebra into $\mathcal A$, sending its field generator to $q_0$.
If an algebraic embedding is required, its injectivity is an
additional algebraic property to check.  For the Hilbert-space
restriction below the necessary condition is equality of the
scalar functional:
\begin{equation}
 \omega(\iota(F))=\omega_d(F)
 \quad\hbox{for every scalar word sum }F,
 \label{eq:gauge-observable-extension-criteria-GEX2}
\end{equation}
where $\omega_d$ is the full scalar Minkowski functional already
constructed.  Any scalar relation killed by $\iota$ must consequently
be null for that functional.  The readout identifying $q_0$ is the
actual centered normalization in \eqref{eq:scalar-observable-information-GSI1}, not an independently chosen
scalar spectator.

\subsubsection*{The exact mixed positivity and relation tests}

A candidate joint moment family gives a complex linear functional
$\omega$ on finite word sums.  Its elementary requirements are
normalization, hermiticity and all finite Gram inequalities:
\begin{equation}
 \begin{gathered}
 \omega(1)=1,\qquad \omega(F^*)=\overline{\omega(F)},\\
 \sum_{i,j=1}^r\overline{c_i}c_j\,
                  \omega(F_i^*F_j)\ge0
 \quad(r<\infty, c_i\in\mathbb C).
 \end{gathered}
 \label{eq:gauge-observable-extension-criteria-GEX3}
\end{equation}
Here the $F_i$ may involve different types and different word
lengths.  Separate positivity for each type is insufficient.  The
inequality for all finite sums is equivalent to
$\omega(F^*F)\ge0$ for every element $F$ of the free algebra.

The ideal must also be annihilated.  For a star-closed set of
generators $R$ of $\mathcal I$, an exact equivalent test, under
\eqref{eq:gauge-observable-extension-criteria-GEX3}, is
\begin{equation}
 \begin{split}
 &\omega(A R B)=0\quad\hbox{for every }A,B\text{ and every }R;\\
 &\hspace{3mm}\Longleftrightarrow\quad
 \omega((RB)^*RB)=0\quad\hbox{for every }B\text{ and every }R.
 \end{split}
 \label{eq:gauge-observable-extension-criteria-GEX4}
\end{equation}
To prove the reverse implication, apply Cauchy--Schwarz to the
positive form $\langle C,D\rangle=\omega(C^*D)$ with $C=A^*$ and
$D=RB$.  Its second diagonal entry is zero, so
$\omega(A R B)=0$.  The forward implication uses $A=B^*R^*$.
Linearity then annihilates the entire generated ideal.  Both $R$
and $R^*$ are included in the declared generating set.  Merely
requiring $\omega(R)=0$ would not suffice: a centered real variable
can have zero mean and positive variance.

Even pairwise consistency of positive laws misses a genuine mixed
constraint.  Consider three proposed signs $A,B,C\in\{-1,1\}$ with
uniform individual marginals.  Give $(A,B)$ and $(B,C)$ the law of
two identical uniform signs, and $(A,C)$ the law of two opposite
uniform signs.  Every pair law is positive and all their one-sign
restrictions agree.  A common joint positive state would nevertheless
have
\begin{equation}
 \omega((A-B+C)^2)
  =3-2\omega(AB)+2\omega(AC)-2\omega(BC)=-3.
 \label{eq:gauge-observable-extension-criteria-GEX5}
\end{equation}
This contradicts \eqref{eq:gauge-observable-extension-criteria-GEX3}.  Thus a family of compatible two-observable
laws is not a substitute for positive matrices of all finite words.

\subsubsection*{Test continuity has a concrete meaning}

For a fixed word of point-field types, require a bound
\begin{equation}
 |\omega(q_{\alpha_1}(f_1)\cdots q_{\alpha_n}(f_n))|
    \le C_{\boldsymbol\alpha}\prod_{j=1}^n p_M(f_j)
 \label{eq:gauge-observable-extension-criteria-GEX6}
\end{equation}
for some finite $M$, which may depend on the word.  Different
orders in the factors can be replaced by their maximum.  This is
equivalent to continuity of that multilinear form on the product of
the Schwartz spaces: continuity at zero gives a product neighborhood
where its absolute value is at most one, and rescaling each nonzero
factor gives \eqref{eq:gauge-observable-extension-criteria-GEX6}.  The seminorms contain the supremum norm, so a
zero factor seminorm means the test is zero.

The explicit tensor expansion of RTE, applied with this $M$, gives
a unique full Schwartz distribution for the word, with
\begin{equation}
 |W_{\boldsymbol\alpha}(F)|
      \le C'_{\boldsymbol\alpha}p_{n(M+8)}(F),
       \qquad F\in\mathcal S((\mathbb R^4)^n).
 \label{eq:gauge-observable-extension-criteria-GEX7}
\end{equation}
This does not postulate a tensor theorem without a norm: the
absolutely summable grid/Fourier expansion proved in RTE bounds
$\sum|c_\nu|\prod_jp_M(f_{\nu,j})$ by that full-test seminorm.
The same construction works for vector-valued forms whenever the
right side of \eqref{eq:gauge-observable-extension-criteria-GEX6} bounds a norm in the stated complete target.
Conversely a full distribution bound implies a simple-tensor bound,
since $p_r(f_1\otimes\cdots\otimes f_n)\le\prod_jp_r(f_j)$.

Thus one may complete each fixed typed-word space to its full
Schwartz space.  Star reversal and concatenation are continuous.
Positivity on this completed finite-degree algebra follows by
approximating each of the finitely many full tests by tensor sums:
every entry of its Gram matrix converges, and finite positive
matrices are closed under entrywise limits.  The ideal test is
imposed on this completed algebra as well, or extended there by a
proved continuity bound for the specified relations.  Other test
types need their own explicitly stated continuity/completion
criterion; \eqref{eq:gauge-observable-extension-criteria-GEX7} does not assign them point-field test spaces.

These degreewise bounds do not by themselves verify the linear
growth hypothesis of a Euclidean reconstruction theorem.  Such a
theorem has a further quantitative hypothesis across all word
lengths.  QRM checked that hypothesis for the scalar family; it has
not checked it for arbitrary additional types.

\subsubsection*{The joint GNS construction and the scalar subspace}

Assume \eqref{eq:gauge-observable-extension-criteria-GEX2}--\eqref{eq:gauge-observable-extension-criteria-GEX4} and the specified test continuity.  On
$\mathcal A$ use the positive form
$\langle A,B\rangle=\omega(A^*B)$.  Its Cauchy--Schwarz inequality
follows by applying positivity to $A+zB$ for every complex $z$.
The null space $\mathcal N$ is a left ideal.  Indeed if
$\omega(B^*B)=0$, then
\begin{equation}
 \omega((AB)^*AB)
   =\langle B,A^*AB\rangle=0
 \label{eq:gauge-observable-extension-criteria-GEX8}
\end{equation}
by Cauchy--Schwarz.  All the words in the second argument are
elements of the algebra, so their norms are finite.  There is no
implicit bounded-operator estimate in this proof.

Complete $\mathcal D=\mathcal A/\mathcal N$ to $\mathcal H$, put
$\Omega=[1]$, and act by left multiplication.  It is well defined
by \eqref{eq:gauge-observable-extension-criteria-GEX8} and gives an invariant dense polynomial domain.  On that
domain the star identity gives
\begin{equation}
 \langle \pi(A^*)v,w\rangle=\langle v,\pi(A)w\rangle,
       \qquad v,w\in\mathcal D.
 \label{eq:gauge-observable-extension-criteria-GEX9}
\end{equation}
In particular each field is closable, because its adjoint contains
its declared star partner on a dense domain.  Positivity alone gives
neither boundedness nor essential self-adjointness of these fields.

The equality in \eqref{eq:gauge-observable-extension-criteria-GEX2} gives an isometry from the scalar polynomial space
to $\mathcal H$, with
\begin{equation}
 J[F]_d=[\iota(F)],\qquad
 \mathcal H_d=\overline{\pi(\iota(\mathcal A_d))\Omega},
       \qquad J:\mathcal K_{\rm ord}\longrightarrow\mathcal H_d
       \text{ unitary}.
 \label{eq:gauge-observable-extension-criteria-GEX10}
\end{equation}
The norm identity is exactly the scalar restriction applied to
$F^*G$; it proves that null vectors are respected.  Surjectivity is
onto the displayed closed subspace by its definition.  It is onto
the entire joint Hilbert space precisely when the scalar vacuum
words are dense there.  That density is an additional condition,
not a consequence of matching the scalar moments.

Already for bounded commutative algebras one can see the distinction.
The state $\nu\otimes\tfrac12(\delta_{-1}+\delta_1)$ on
$C(Y\times\{-1,1\})$ restricts to $\nu$ on $C(Y)$.
Its GNS scalar subspace is $L^2(\nu)\otimes\mathbb C1$, whereas
the sign observable supplies the orthogonal nonzero vector
$1\otimes s$.  This is an algebraic example of additional degrees
of freedom with exactly the same scalar state; it makes no claim
about a continuum gauge extension of the constructed model.

There is a useful exact compression criterion in this bounded case.
For a representation of a unital star algebra by bounded operators,
let $P$ project onto a closed scalar subspace.  Compression is
multiplicative on the full algebra if and only if that subspace is
reducing.  One direction is immediate.  For the other, a
multiplicative compression would give
\begin{equation}
 P\pi(A^*)(1-P)\pi(A)P
   =P\pi(A^*A)P-P\pi(A^*)P\pi(A)P=0.
 \label{eq:gauge-observable-extension-criteria-GEX11}
\end{equation}
The left side is the square of the norm of $(1-P)\pi(A)P$ in
quadratic-form language.  Hence this operator vanishes for every
$A$, and also for $A^*$, proving reduction.  If $\Omega$ is in the
subspace and is cyclic for the full algebra, reduction then forces
the subspace to be all of $\mathcal H$.  For unbounded fields an
arbitrary projection need not even preserve their domains, so no
unqualified compression representation is asserted.

\subsubsection*{Consistency across finite observable families}

Suppose a directed system of unital star algebras
$\mathcal A_J$ has compatible injective star maps for $J\subset K$.
Suppose normalized positive functionals $\omega_J$ satisfy
\begin{equation}
 \omega_K|_{\mathcal A_J}=\omega_J.
 \label{eq:gauge-observable-extension-criteria-GEX12}
\end{equation}
There is then a unique normalized positive functional on their
algebraic inductive limit.  Define it on an element by using any
stage containing that element.  A common larger stage and \eqref{eq:gauge-observable-extension-criteria-GEX12}
make the value independent of the choice.  Any finite word sum
belongs to a common stage, so its square has a nonnegative
expectation there; this proves positivity on the union.  The same
argument proves uniqueness.  For a system of isometric unital
$C^*$ embeddings the state has norm one on the union and extends
to its norm completion.  An unbounded field algebra has no such
automatic $C^*$ completion; its degreewise test continuity is
checked as in \eqref{eq:gauge-observable-extension-criteria-GEX6}--\eqref{eq:gauge-observable-extension-criteria-GEX7}.

If global relations are first specified in one free quotient, the
maps from smaller quotients need not be injective automatically.
One can instead prescribe all finite mixed moments in that global
quotient and apply \eqref{eq:gauge-observable-extension-criteria-GEX3}--\eqref{eq:gauge-observable-extension-criteria-GEX4} directly.  Adding a new type must
not silently impose a previously untested relation on an old type.

A practical sufficient route to \eqref{eq:gauge-observable-extension-criteria-GEX12} is uniqueness of a matched
finite-family reconstruction.  Suppose the same Euclidean joint
data on every overlap have reconstructions unique in the exact
analytic/boundary class being used.  Restrict a larger reconstruction
to the smaller family.  It has the smaller Euclidean data and the
same required class properties.  Uniqueness makes it the smaller
reconstruction.  Thus the extensions at successive stages are consistent.  Both the class
membership and the uniqueness theorem must be checked for the
actual new types; scalar uniqueness alone is not that theorem.

\subsubsection*{Covariance and its continuity on the joint space}

Assume a declared spacetime symmetry group acts by star
automorphisms $\alpha_g$ of $\mathcal A$, respecting the relations,
and that $\omega\circ\alpha_g=\omega$.  Then
\begin{equation}
 U_g[A]=[\alpha_g A]
 \label{eq:gauge-observable-extension-criteria-GEX13}
\end{equation}
is well defined and isometric; $\alpha_{g^{-1}}$ gives its inverse.
It implements the action on the common polynomial domain, fixes
$\Omega$, and satisfies the group law there and hence everywhere.
Strong continuity follows if, for every polynomial $A$,
$g\mapsto\omega(A^*\alpha_gA)$ is continuous at the identity,
because
\begin{equation}
 \|U_g[A]-[A]\|^2
   =2\omega(A^*A)-2\operatorname{Re}\omega(A^*\alpha_gA).
 \label{eq:gauge-observable-extension-criteria-GEX14}
\end{equation}
Density and the unitary norm then extend continuity to all vectors.
For point-field types this condition follows from the actual
continuous test action and the joint moment bounds.  For geometric
loop labels it is a condition on that declared label topology and
their joint moments, not a consequence of a name assigned to a
loop.

When the scalar inclusion intertwines these actions, its closed
cyclic subspace is invariant under $U_g$ and $U_{g^{-1}}$, hence
reducing.  The unitary in \eqref{eq:gauge-observable-extension-criteria-GEX10} intertwines the established scalar
symmetries with their restrictions.  It still says nothing about
additional joint states orthogonal to that subspace.

Use the physical translation convention
$U(a)=\exp(i(a^0H-\mathbf a\cdot\mathbf P))$; in the notation of
RHI the physical spatial momentum here is minus its spatial
generator.  Let $E$ be the joint spectral measure of the strongly
continuous translation group.  Positivity of energy-momentum has
the following exact dense test:
\begin{equation}
 E(\mathbb R^4\setminus\overline V_+)=0
 \quad\Longleftrightarrow\quad
 \operatorname{supp}\mu_A\subset\overline V_+
 \ \text{ for every }A\in\mathcal A,
 \qquad
 \mu_A(B)=\|E(B)[A]\|^2.
 \label{eq:gauge-observable-extension-criteria-GEX15}
\end{equation}
Indeed the forward implication is immediate.  Conversely the
spectral projection of the open complement kills every vector of
the dense polynomial domain, and hence is zero.  The measures
$\mu_A$ are the positive Fourier spectral measures of
$\omega(A^*\alpha_a A)$.  Checking only the vacuum coefficient
checks just the atom at zero, and checking only single-field
vectors checks only their span.  The criterion is for all words,
unless an independent density theorem reduces that test family.

\subsubsection*{Locality as an all-spectator moment identity}

For spacelike separated declared localization regions, put
$R=[q_\alpha(f),q_\beta(g)]$.  The exact joint locality criterion is
\begin{equation}
 \omega(A R B)=0\quad\hbox{for every polynomial }A,B.
 \label{eq:gauge-observable-extension-criteria-GEX16}
\end{equation}
It is equivalent to $\pi(R)=0$ on the common polynomial domain.
For if \eqref{eq:gauge-observable-extension-criteria-GEX16} holds, the vector $\pi(R)[B]$ has zero inner product
with every polynomial vector $[C]$, by taking $A=C^*$.  Density
forces the vector to vanish.  Conversely vanishing of those
vectors gives \eqref{eq:gauge-observable-extension-criteria-GEX16}.  Equivalently the norm-square test in
\eqref{eq:gauge-observable-extension-criteria-GEX4} can be used for these commutator relations.

Equality of only the two-point functions with exchanged entries
does not test \eqref{eq:gauge-observable-extension-criteria-GEX16}, since it leaves out the spectators.  Once
the elementary commutators vanish on an invariant common domain,
the Leibniz rule for commutators proves locality of their localized
polynomial algebras on that domain.  For extended geometric
observables the regions are their actual supports, not nominal
base points.  No spacelike commutativity is inferred from two
disjoint curves unless their complete declared localization
regions satisfy the required condition.

This is a criterion for the larger moment family.  For the scalar
family the necessary identity is already a conclusion of QRM's
matched corrected OS theorem.  For a new family it must follow
from that family's actual joint analytic/permutation argument or
a reconstruction theorem whose hypotheses have been checked for
it.  It is not used here as an assumption purporting to prove that
same joint continuation.

\subsubsection*{A precise common-domain criterion for added point fields}

We next prove how verified energy bounds upgrade the polynomial
construction.  Restrict this paragraph to point-field Schwartz
types, or to other test types with the same smooth translation
action and moment continuity.  Let the joint representation above
contain only types meeting that assumption in this domain argument;
otherwise the conclusion applies to the cyclic subspace of the
qualified subfamily, without a graph-core assertion for the larger
space.  Let that representation
have the forward spectral property.  Its polynomial test domain
$\mathcal D$ lies in $\bigcap_k\operatorname{Dom}H^k$.
For a full word test $F$ translation in time translates all its
arguments; continuity of the two-word norm form permits every
Schwartz difference quotient in Hilbert norm.  Stone's generator
identity gives
\begin{equation}
 H^k[F]=i^k[D_0^kF],\qquad
 D_0=\sum_j\partial_{x_j^0}.
 \label{eq:gauge-observable-extension-criteria-GEX17}
\end{equation}
Iterating proves the assertion for all $k$.  The sign agrees with
$f_a(x)=f(x-a)$ and $U(t)=e^{itH}$.

The same domain is a graph core for every $(1+H)^k$.  To see this,
choose a smooth compactly supported scalar cutoff $c_R$ equal to
one on $[0,R]$ and tending pointwise to one on $[0,\infty)$, with
$|c_R|\le1$.  For fixed $R$, $(1+H)^kc_R(H)$ is bounded, so dense
polynomial vectors approximate $c_R(H)\psi$ in graph norm after
applying $c_R(H)$.  For a polynomial vector $v$, Fourier inversion
expresses $c_R(H)v$ as an integral of $U(t)v$ with a Schwartz
coefficient.  This integral is approximated by finite Riemann sums
in graph norm: on compact time intervals the translated full test
is graph continuous by \eqref{eq:gauge-observable-extension-criteria-GEX17}, and the tail is bounded by the
$L^1$ tail times $\|(1+H)^kv\|$.  The Riemann sums are polynomial
vectors.  Finally $c_R(H)\psi\to\psi$ in graph norm for
$\psi\in\operatorname{Dom}(1+H)^k$ by spectral dominated
convergence.  This proves the core claim with both approximations
in their required order.

Set $\mathsf E=1+H$.  Suppose for each added type there is a
nonnegative integer $N_\alpha$ such that, for every integer $k\ge0$,
the \emph{actual} polynomial operator satisfies
\begin{equation}
 \|\mathsf E^k\pi(q_\alpha(f))v\|
   \le C_{\alpha,k}p_{M_{\alpha,k}}(f)
                         \|\mathsf E^{k+N_\alpha}v\|,
                 \qquad v\in\mathcal D.
 \label{eq:gauge-observable-extension-criteria-GEX18}
\end{equation}
These are bounds to prove on the stated joint space.  The scalar
case with $N_0=9$ is already available on its own scalar cyclic
space.  In a larger representation even the action of the scalar
field on vectors created by new observables requires the
corresponding mixed bound.  The scalar result proves neither that
larger-space bound nor \eqref{eq:gauge-observable-extension-criteria-GEX18} for a Wilson-loop or composite type.

The core claim and closedness of $\mathsf E^k$ extend the operator
coherently to
$\mathcal D_\infty=\bigcap_k\operatorname{Dom}H^k$, mapping it
into itself.  More explicitly choose polynomial $v_j$ approaching
$v\in\mathcal D_\infty$ in graph norms through order $j$.
For each fixed $k$, \eqref{eq:gauge-observable-extension-criteria-GEX18} makes the field images Cauchy in the
$\mathsf E^k$ graph norm.  Their Hilbert limits coincide for all
$k$, and closedness identifies the higher energy powers of that
same limit.  The extension is independent of the approximating
sequence.  The star identity extends by the same graph
approximations for both vectors, so the extended operators remain
closable.

For $k=0$ the operator also extends from
$\operatorname{Dom}\mathsf E^{N_\alpha}$ to the Hilbert space
with its stated graph bound.  Approximating in that graph norm
shows that this extension lies in the closure of the original
polynomial operator.  Conversely all three operators, on
$\mathcal D$, on $\mathcal D_\infty$, and on
$\operatorname{Dom}\mathsf E^{N_\alpha}$, have the same closure:
their graphs are nested, and every graph point of the largest
restriction is approximated by the smallest.  This is a canonical
closure statement, not essential self-adjointness.

For a word let $S_j=\sum_{i<j}N_{\alpha_i}$ and
$S=\sum_iN_{\alpha_i}$.  Inserting the appropriate energy powers
between factors gives, first on $\mathcal D_\infty$ and then as a
bounded weighted map,
\begin{equation}
 \bigl\|\mathsf E^l\pi(q_{\alpha_1}(f_1))\cdots
       \pi(q_{\alpha_n}(f_n))\mathsf E^{-(l+S)}\bigr\|
   \le\prod_{j=1}^n C_{\alpha_j,l+S_j}
                         p_{M_{\alpha_j,l+S_j}}(f_j).
 \label{eq:gauge-observable-extension-criteria-GEX19}
\end{equation}
Indeed the $j$th weighted factor is
$\mathsf E^{l+S_j}\pi(q_{\alpha_j}(f_j))
\mathsf E^{-(l+S_j+N_{\alpha_j})}$, and the adjacent powers
cancel on the energy-smooth domain.  Their bounded extensions
give the asserted norm inequality.  Applying the vector-valued
RTE expansion with the maximum of the displayed test orders
then gives a full non-tensor word operator with an explicit
finite Schwartz seminorm bound.  The vacuum vector has
$\mathsf E^r\Omega=\Omega$, so this also bounds every fixed
mixed vacuum distribution.

If \eqref{eq:gauge-observable-extension-criteria-GEX16} holds, the product estimate with energy order
$N_\alpha+N_\beta$ extends its commutator identity from the
polynomial domain to $\operatorname{Dom}
\mathsf E^{N_\alpha+N_\beta}$, using the graph core just proved.
In particular it holds on $\mathcal D_\infty$.  This extension is
operator locality on a declared common domain; strong commutation
of self-adjoint spectral calculi would require further hypotheses.

If the full joint Poincare representation is also established,
its forward cone implies
$c_\Lambda H\le U_\Lambda HU_\Lambda^{-1}\le C_\Lambda H$
for each fixed proper orthochronous Lorentz transformation, in the
joint spectral sense.  The time component of a transformed
forward momentum is bounded above and below by positive
constants times its energy.  It follows that
$U_\Lambda\mathcal D_\infty=\mathcal D_\infty$, with equivalent
energy graph norms.  Covariance on the polynomial domain then
extends by \eqref{eq:gauge-observable-extension-criteria-GEX18} to this common domain, with the specified
finite-component test transformations.  No unrestricted inverse
heat operator or all-energy band sum is needed for this domain
argument.

\subsubsection*{Actual cutoff readouts and renormalization}

A continuum label requires a concrete cutoff function on the same
bond configuration space and the same real measure as the scalar
construction.  Denote it by $Q_{a,\alpha}(f)$.  It must be fully
site-gauge invariant.  For compactly localized tests it also has a
declared finite support, including every connecting path, smearing
region and auxiliary loop.  A typical linear mixing prescription is
\begin{equation}
 Q_{a,\alpha}(f)
  =\sum_{\beta\in J_\alpha}Z_{a,\alpha\beta}
             B_{a,\beta}(D_{a,\alpha\beta}f)
       +c_{a,\alpha}(f)1,
 \label{eq:gauge-observable-extension-criteria-GEX20}
\end{equation}
where $J_\alpha$ is finite for this prescription, the bare functions
$B_{a,\beta}$ and lattice smearing maps $D_{a,\alpha\beta}$ are
specified, and the subtraction $c_{a,\alpha}$ is specified and
linear.  The formula in \eqref{eq:gauge-observable-extension-criteria-GEX20} is a template for a definition, not an
assertion that suitable coefficients or bounds exist.  A proposed
readout has to supply its actual coefficients, their normalization,
and the moment bounds after mixing.  Diverging $Z$ factors cannot
be ignored because each bare Wilson trace is bounded.

The star and Euclidean reflection operations must be matched by
that definition.  In particular the real cutoff identities are
\begin{equation}
 Q_{a,\alpha^*}(f^*)=\overline{Q_{a,\alpha}(f)},\qquad
 \Theta_aQ_{a,\alpha}(f)
       =Q_{a,\vartheta\alpha}(\vartheta f),
 \label{eq:gauge-observable-extension-criteria-GEX21}
\end{equation}
with the antilinear test/reflection action on the right specified
for the type.  For a tensor observable the reflection can act on
its components.  For a loop it acts on its entire oriented path.
The complex coefficients and counterterms in \eqref{eq:gauge-observable-extension-criteria-GEX20} must respect
these identities or their explicitly proved limiting versions.
The constant counterterm has empty localization support; the
nonconstant part has its actual support.

For a point-field readout a useful localization requirement is
\begin{equation}
 \operatorname{supp}Q_{a,\alpha}(f)
     \subset(\operatorname{supp}f)^{\epsilon_a},
                  \qquad \epsilon_a\downarrow0,
 \label{eq:gauge-observable-extension-criteria-GEX22}
\end{equation}
where the left side denotes the physical locations of all bonds
on which the function depends.  For an extended type replace
$\operatorname{supp}f$ by its declared geometric region.  It is
this condition, not the name of the observable, that allows a
fixed positive-time gap to survive at small cutoff.  It does not
by itself give a moment bound or independence of the lattice
rounding of a curve.

The Euclidean readout map takes products into the commutative
algebra of bond functions.  It is not a homomorphism from the
noncommutative Minkowski field algebra.  A Euclidean family must
first be reconstructed, with its ordered analytic continuation
and boundary conventions, before it can supply the joint
Minkowski functional used in \eqref{eq:gauge-observable-extension-criteria-GEX2}--\eqref{eq:gauge-observable-extension-criteria-GEX19}.  Ordinary positive
probability expectations of the commuting cutoff functions do
not alone prove Minkowski positivity and spectrum.

\subsubsection*{Passing the actual reflected Gram matrices to a limit}

Let $\rho_a$ be real probability measures with a proved lattice OS
reflection inequality.  Write $\Theta_a$ for its antilinear
reflection, including complex conjugation.  For positive-time
cutoff functions $F_{a,1},\ldots,F_{a,r}$ the actual inequality is
\begin{equation}
 \sum_{i,j}\overline{c_i}c_j\,
      \rho_a((\Theta_aF_{a,i})F_{a,j})\ge0.
 \label{eq:gauge-observable-extension-criteria-GEX23}
\end{equation}
Suppose the word tests defining these functions have a fixed
positive-time gap $\delta>0$ and their readouts satisfy \eqref{eq:gauge-observable-extension-criteria-GEX22}.
If the discrete reflection plane is displaced by $b_a\to0$ from
the continuum plane, then for all sufficiently small $a$,
$\epsilon_a+|b_a|<\delta$.  Every bond support is on the required
positive side.  Thus \eqref{eq:gauge-observable-extension-criteria-GEX23} applies to the actual readouts rather
than a formally localized replacement.

Suppose also that each entry in \eqref{eq:gauge-observable-extension-criteria-GEX23} has a joint limit given by
the proposed Euclidean data.  Passing a finite sum to the limit
gives a positive continuum reflected Gram matrix.  This is a
complete proof of that implication: for each fixed vector of
coefficients the quadratic form is a limit of nonnegative real
numbers.  One must nevertheless verify that the reflected word
test for each entry belongs to the space on which the limit is
known.  A separated or diagonal-flat distribution cannot be
evaluated on an arbitrary coincident-point test without this
check.

For the scalar family QRM performs the reflected-product
diagonal-flat check and the enlargement of the positive-time
quotient.  For a proposed new family one can start with compact
strictly ordered positive-time words away from all internal
diagonals, verify that their reflected products are in the
declared joint limit space, and then prove the density/insertion
estimate needed to enlarge that space.  This explicitly
identifies the further test-space task.  It does not invent
Euclidean contact terms to make an undefined Gram entry exist.

Likewise lattice symmetries pass to any joint limit when their
actual readouts intertwine the corresponding test actions and
the limiting distributions are continuous in the needed
topology.  Discrete lattice rotations alone do not prove
continuous Euclidean rotational covariance.  The scalar G2
argument supplies it for its own family.  An extension needs the
matching assertion for its mixed data if that assertion enters
its reconstruction theorem.

Real OS inequalities are used at real physical measures.  An
analytic continuation of coarse backgrounds or of external
currents is a separate analytic operation.  A complex source
functional is not a probability state merely because its
restriction at zero real source is positive.

\subsubsection*{A useful bounded loop compactness statement}

There is a limited joint compactness result available without any
renormalized insertion estimate.  Choose countably many labelled
closed paths $\gamma_{a,j}$ at each cutoff, and set
$Y_{a,j}=w_{\gamma_{a,j}}\in[-1,1]$.  The labels include their
chosen lattice rounding.  For any sequence $a_k\downarrow0$ there
is a subsequence on which all finite joint polynomial moments of
these bounded real coordinates converge.

Indeed the finite coordinate monomials form a countable set,
and each expectation is bounded by one.  Repeated subsequences
and a diagonal selection give a common subsequence.  Define $L$
on coordinate polynomials by these limits.  For
$K=[-1,1]^{\mathbb N}$,
\begin{equation}
 |L(P)|\le\|P\|_{C(K)},\qquad L(1)=1,
       \qquad L(P)\ge0\quad\hbox{if }P\ge0\text{ on }K.
 \label{eq:gauge-observable-extension-criteria-GEX24}
\end{equation}
The inequalities follow by applying them before taking the
limit, since every vector of readout values belongs to $K$.
Coordinate polynomials are conjugation closed and separate
points of the compact product.  Their uniform closure is
$C(K)$, so \eqref{eq:gauge-observable-extension-criteria-GEX24} extends $L$ uniquely to a positive state.
Riesz gives a probability measure on $K$ with exactly the
selected joint moments.  This proves full finite-coordinate
consistency, not just convergence of separate one-loop laws.

If a fixed finite-coordinate polynomial relation $R$ holds for
the readouts at every cutoff, then
\begin{equation}
 L(|R|^2)=0,
 \label{eq:gauge-observable-extension-criteria-GEX25}
\end{equation}
so $R=0$ almost surely in the limit.  More generally the same
conclusion holds if the coefficients of relations $R_a$ converge
to those of $R$, their finite coordinate set is fixed, and
$R_a(Y_a)=0$.  In fact $R_a\to R$ uniformly on $K$, whence
$\rho_a(|R(Y_a)|^2)\to0$.  This is a direct way to retain a
specified finite collection of trace relations, provided the
rounded labels preserve their concatenations and inverses.  If
those label identifications are not exact, a quantitative
replacement estimate is needed.

This compactness statement need not give a nontrivial limit:
the state concentrated at $Y_j=1$ is permitted.  It gives no
rounding independence, continuous spacetime covariance,
renormalized local composite limit, or mixed law with the
unbounded scalar distributions.  It also gives no uniform
control after multiplying these coordinates by diverging
renormalization constants.  Those are different assertions from
the bounded probability statement just proved.

\subsubsection*{What mixed bounds a bounded loop actually needs}

Let $Y_a$ be any cutoff loop polynomial with $|Y_a|\le1$ and
let $Q_a=\prod_{j=1}^nX_a(f_j)$.  At the same real cutoff measure
one has
\begin{equation}
 |\rho_a(Y_aQ_a)|
   \le\rho_a\Bigl(\prod_j|X_a(f_j)|\Bigr),\qquad
 |\rho_a(Y_aQ_a)|^2\le\rho_a(|Q_a|^2).
 \label{eq:gauge-observable-extension-criteria-GEX26}
\end{equation}
The first is the triangle inequality and the second is
Cauchy--Schwarz with $\rho_a(|Y_a|^2)\le1$.  For example a
sufficient further estimate is
\begin{equation}
 \sup_a\rho_a\Bigl(\prod_j|X_a(f_j)|^2\Bigr)
              \le C_n\prod_j p_M(f_j)^2.
 \label{eq:gauge-observable-extension-criteria-GEX27}
\end{equation}
It gives uniform boundedness and test continuity for the mixed
word in \eqref{eq:gauge-observable-extension-criteria-GEX26}.  A complete mixed subsequence argument would
then also require consistency of every such word, the actual
relation tests, and the reflected Gram limits where they are
used.  It would not yet infer its joint Minkowski continuation
without a matched reconstruction.

The separated signed scalar distribution bound does not itself
evaluate the right side of \eqref{eq:gauge-observable-extension-criteria-GEX27}.  That expression is an
absolute repeated product at one real cutoff measure.  Its
test is not automatically in a diagonal-flat Euclidean domain.
If a stronger insertion estimate is available, its exact
observables, repeated products and exterior variables must be
matched before using it here.  In particular a formal
arbitrary-current coefficient is not identified with the needed
physical mixed product merely by naming it $E^{(2)}$.  The coefficient/readout identification must be proved on its
stated domain wherever it is used.

For a fixed family, uniformly bounded absolute moments of
every required degree permit diagonal moment subsequences on a
countable dense test set.  Uniform multilinear bounds of the
form in \eqref{eq:gauge-observable-extension-criteria-GEX6} then extend the limits to all tests by its explicit
tensor argument.  Positivity and relation identities pass
through finite matrices.  The uniform bounds are essential in
this extension from a countable test set; convergence on a
dense set alone does not establish distributional convergence.

\subsubsection*{A quantitative path-sliding comparison}

There is a concrete geometric estimate useful in checking
rounding consistency.  Suppose two fine paths with the same
endpoints are connected by an explicit finite sequence of
elementary plaquette moves and retracings.  An elementary move
replaces two sides of a plaquette by the other two sides with
the same endpoints.  If $A$ and $B$ are the two path products,
$AB^{-1}$ is the plaquette holonomy with the appropriate
orientation and base point.  All other prefix and suffix
products are unitary.  Hence the norm of the change in the
whole path product is exactly $\|I-U_p\|$ for that move, up to
conjugation or inversion, both of which preserve this norm.
Telescoping along the supplied sequence gives
\begin{equation}
 \|U_\gamma-U_{\gamma'}\|_{\rm op}
       \le\sum_{\ell=1}^{M}\|I-U_{p_\ell}\|_{\rm op}.
 \label{eq:gauge-observable-extension-criteria-GEX28}
\end{equation}
Retracings give zero change.  Repeated use of a plaquette is
counted with its multiplicity.  For closed paths the same
bound applies to normalized traces, because
$|\operatorname{Tr}C|/2\le\|C\|_{\rm op}$ for two-by-two
matrices.

For $U\in\operatorname{SU}(2)$ its eigenvalues are $e^{\pm i\theta}$,
so the actual plaquette action satisfies
\begin{equation}
 \|I-U_p\|_{\rm op}^2
      =2\bigl(1-\tfrac12\operatorname{ReTr}U_p\bigr)=2A_p.
 \label{eq:gauge-observable-extension-criteria-GEX29}
\end{equation}
Taking $L^q(\rho_a)$ norms in \eqref{eq:gauge-observable-extension-criteria-GEX28}, with $q\ge2$, and using
the triangle inequality therefore gives the explicit sufficient
comparison
\begin{equation}
 \|w_\gamma-w_{\gamma'}\|_{L^q(\rho_a)}
    \le\sqrt2\sum_{\ell=1}^{M}
                   \|A_{p_\ell}\|_{L^{q/2}(\rho_a)}^{1/2}.
 \label{eq:gauge-observable-extension-criteria-GEX30}
\end{equation}
This is a true estimate for the real unitary cutoff variables.
If the right side tends to zero for the specified rounding
moves, it proves convergence of their trace difference in that
norm.  Mere proximity of two physical curves does not bound
the number of moves or this sum.  A proof must count the actual
filling and provide the needed plaquette norm rate.  A uniform
pointwise small-curvature estimate $\|I-U_p\|\le\varepsilon_a$
only yields $M\varepsilon_a$, which need not tend to zero.

No filling is asserted between paths of different winding
class on a torus.  Indeed \eqref{eq:scalar-observable-information-GSI8} rules out any estimate which
would remove nontrivial flat winding holonomies by such a
claim.  The supplied move sequence, endpoints and periodic
identifications are part of the hypotheses of \eqref{eq:gauge-observable-extension-criteria-GEX28}.

\subsubsection*{A controlled change of readout scheme}

Suppose two readout prescriptions $Q_{a,j}$ and $Q'_{a,j}$ for a
fixed word satisfy, for that word length $n$,
\begin{equation}
 \|Q_{a,j}-Q'_{a,j}\|_{L^n(\rho_a)}\longrightarrow0,
 \qquad
 \sup_a\bigl(\|Q_{a,j}\|_{L^n}+\|Q'_{a,j}\|_{L^n}\bigr)<\infty.
 \label{eq:gauge-observable-extension-criteria-GEX31}
\end{equation}
Then their joint word expectations have the same limit whenever
one exists.  Expand the difference of products by replacing one
factor at a time.  Holder's inequality with $n$ factors bounds
the $j$th term by
\begin{equation}
 \|Q_{a,j}-Q'_{a,j}\|_{L^n}
       \prod_{i<j}\|Q_{a,i}\|_{L^n}
       \prod_{i>j}\|Q'_{a,i}\|_{L^n},
 \label{eq:gauge-observable-extension-criteria-GEX32}
\end{equation}
which tends to zero.  The case $n=1$ is just its $L^1$ bound.
Versions with different Holder exponents or proved exterior
insertion norms give the corresponding same conclusion, but
their actual norms must be stated.  For bounded loops,
\eqref{eq:gauge-observable-extension-criteria-GEX30} can supply the difference norm in \eqref{eq:gauge-observable-extension-criteria-GEX31} once its
right side is controlled.  For renormalized fields the
differences include their coefficients and counterterms as
well.

There is also a useful two-limit density criterion.  Suppose
for each intended readout $Q_{a,j}$ there are polynomials
$P_{a,j}^{(r)}$ in a fixed listed family of loop readouts, with
coefficients and finite labels independent of $a$ for each
$r$.  Suppose all fixed-polynomial joint moments have limits,
the required $L^n$ norms are uniformly bounded, and
\begin{equation}
 \lim_{r\to\infty}\limsup_{a\downarrow0}
      \|Q_{a,j}-P_{a,j}^{(r)}\|_{L^n(\rho_a)}=0
 \label{eq:gauge-observable-extension-criteria-GEX33}
\end{equation}
for every fixed word length involved.  The product telescoping
bound gives a uniform limiting error $C_n\varepsilon_r$.
For each $r$ take the known polynomial-word limit.  Applying
the same error bound for $r$ and $s$ shows these limits are
Cauchy as $r,s\to\infty$.  The limiting value is then the
unique limit for the original word: its limsup distance from
the $r$th polynomial limit is at most $C_n\varepsilon_r$.
This proves the criterion without interchanging an
uncontrolled pair of limits.

The fixed-cutoff density theorem in \eqref{eq:su2-wilson-orbit-coordinates-GWO14} alone gives no such
criterion.  Its approximating degrees, loop labels and
coefficients may depend arbitrarily on $a$.  Convergence of
each fixed loop polynomial does not control a moving
polynomial with diverging coefficients.  The condition in \eqref{eq:gauge-observable-extension-criteria-GEX33}, or
an actual alternative uniform insertion norm, is the missing
quantitative link in such a use of cutoff completeness.

\subsubsection*{Exact laws and dense-observable transport}

At a fixed real cutoff, a genuine integration or change of
variables may induce a bounded observable map $T_a$ and an
identity for a determining class.  The dense-class norm
estimate in \eqref{eq:gauge-observable-algebra-GIA15}, applied to the complete Wilson polynomial
family in \eqref{eq:su2-wilson-orbit-coordinates-GWO16}, extends that identity to every continuous
invariant observable.  If the actual induced full bond
measure is gauge invariant, equality there determines the
full gauge-invariant probability law by \eqref{eq:gauge-observable-algebra-GIA11}, and hence
all bounded Borel readouts with that law.

The hypotheses include the map's actual observable type, its
boundedness at that cutoff, and the identity on a complete
class.  An identity only for the scalar partition function,
its scalar derivatives, or one special insertion is not the
determining-class identity for all Wilson polynomials.  A
based gauge-fixed calculation must also be transported to
the full-site invariant functions before the orbit-space
state statement applies.  None of these requirements is
replaced by preservation of the constant function alone.

Across different cutoffs the norms and readouts in this
argument may change.  A finite-cutoff equality is exact, but
any passage to a renormalized joint continuum family still
needs the continuity and limit estimates just described.
This separates an exact integral identity from the
independent estimates needed to read a continuum observable
through it.

\subsubsection*{Conditional joint reconstruction and physical identification}

The previous statements give a precise conditional extension
theorem.  Suppose a declared family of joint Minkowski word
distributions satisfies the scalar restriction
\eqref{eq:gauge-observable-extension-criteria-GEX2}, mixed positivity
\eqref{eq:gauge-observable-extension-criteria-GEX3}, the relation test
\eqref{eq:gauge-observable-extension-criteria-GEX4}, and the specified test
continuity (for point-field words,
\eqref{eq:gauge-observable-extension-criteria-GEX6}--\eqref{eq:gauge-observable-extension-criteria-GEX7}).  Suppose its symmetry
action obeys the invariance and continuity used in
\eqref{eq:gauge-observable-extension-criteria-GEX13}--\eqref{eq:gauge-observable-extension-criteria-GEX14}, its full polynomial spectral measures pass
\eqref{eq:gauge-observable-extension-criteria-GEX15}, and its spacelike moment identities pass \eqref{eq:gauge-observable-extension-criteria-GEX16}.
Then it has a normalized positive cyclic representation,
closable fields on a common invariant polynomial domain,
the stated strongly continuous covariance and forward
spectrum, and operator locality on that domain.  It contains
the constructed scalar representation as exactly the closed
subspace in \eqref{eq:gauge-observable-extension-criteria-GEX10}.  If the additional whole-space smoothness
and energy bounds in \eqref{eq:gauge-observable-extension-criteria-GEX18} are verified, it also has the
common energy-smooth domain and canonical closure statements
proved above.  All assertions follow from the explicit
arguments in \eqref{eq:gauge-observable-extension-criteria-GEX8}--\eqref{eq:gauge-observable-extension-criteria-GEX19}; no boundedness or domain property
of an added type is imported from its scalar restriction.

If the data are first supplied for finite families, \eqref{eq:gauge-observable-extension-criteria-GEX12}
gives the unique global functional once their actual
overlap consistency is proved.  A matched finite-family
Euclidean reconstruction with uniqueness can prove that
consistency as explained there.  Its theorem must match the
new types, the joint positive-time test domain, the growth
seminorms and the unrestricted boundary values.  For the
scalar family those checks are the established EB/QRM
arguments.  For a larger family they are explicit new
checks, not a generic consequence of citing an OS theorem.

For a gauge-theory identification this conditional theorem
must be attached to actual readouts of the same lattice
theory.  In addition to the compatible joint limits it must
identify their renormalizations, exact limiting relations,
localization, and relation to the scalar $\operatorname{tr}(F^2)$.
One must state why the chosen observable family is sufficient
for the intended continuum gauge-invariant theory.  The
finite-graph theorem supplies that sufficiency at a fixed
cutoff.  Passing it to the continuum needs a proved uniform
criterion such as \eqref{eq:gauge-observable-extension-criteria-GEX33}, or another stated reconstruction
of the target observable algebra.  Bare cutoff completeness
does not silently perform that passage.

Nor do scalar cluster or spectral assertions automatically
become assertions on the enlarged Hilbert space.  The scalar
subspace is reducing for the joint spacetime symmetries, so
its spectral measure is a restriction of the full one.
Restriction does not exclude additional spectrum or
translation-invariant vectors in its orthogonal complement.
For example let any constructed scalar theory be tensored
with $\mathbb C^2$, choose the unit vector
$\eta=(1,1)/\sqrt2$, and put
\begin{equation}
 \begin{gathered}
 \widehat\Omega=\Omega_d\otimes\eta,\qquad
 \widehat U(g)=U_d(g)\otimes I,\qquad
 \widehat\varphi(f)=\varphi_d(f)\otimes I,\\
 \chi(f)=\Bigl(\int f(x)\,dx\Bigr)I\otimes
                  \begin{pmatrix}1&0\\0&-1\end{pmatrix}.
 \end{gathered}
 \label{eq:gauge-observable-extension-criteria-GEX34}
\end{equation}
The added field is bounded for each test and central.  Its
moments are tempered, it is a scalar under proper
orthochronous Poincare transformations, and it is local.
All mixed moments are those of this tensor representation,
so they are positive.  The scalar restriction is unchanged.
The vectors $\eta$ and
$\operatorname{diag}(1,-1)\eta$ span $\mathbb C^2$, so the
full polynomial vacuum space is the full tensor space.
Both copies of any scalar zero-energy vector have zero
energy.  In particular a unique scalar vacuum need not
become a unique joint vacuum.  This is an explicit
extension illustrating the logical distinction; it is not
claimed to be an actual readout from the lattice gauge
measure.

A full joint cluster assertion would resolve this particular
vacuum issue.  If for every pair of polynomial vectors the
spatial translation coefficients cluster to their vacuum
product, density and unitarity extend that convergence to
the weak operator limit
\begin{equation}
 U(0,\mathbf a)\ \rightharpoonup\
       |\Omega\rangle\langle\Omega|
                 \qquad (|\mathbf a|\to\infty).
 \label{eq:gauge-observable-extension-criteria-GEX35}
\end{equation}
The extension uses fixed polynomial approximants for both
vectors and the uniform unitary norm, so no uniformity over
all polynomials is being assumed.  By the forward cone,
$H=0$ implies $\mathbf P=0$.  A zero-energy vector is
therefore invariant under every spatial translation.
Applying \eqref{eq:gauge-observable-extension-criteria-GEX35} to it proves
$\ker H=\mathbb C\Omega$.  This is a proof from a
\emph{full joint} cluster assertion; scalar clustering
only supplies the statement on the scalar cyclic
subspace.  The phase mixture discussed in SPH retains its
own explicit invariant subspace instead of assuming such
a joint purity statement.

\subsubsection*{Bounded Euclidean multiplication is not yet an OS field}

The ordinary bounded loop result in \eqref{eq:gauge-observable-extension-criteria-GEX24} and the measure
inequalities in \eqref{eq:gauge-observable-extension-criteria-GEX26} do not give an OS operator norm.  This
failure can be demonstrated at the level of a finite
reflection-positive system.  Let $s_-$ and $s_+$ be
independent uniform signs, with reflection exchanging them
and conjugating coefficients.  For positive-time functions
$F,G$ of $s_+$ the OS form is
\begin{equation}
 \langle[F],[G]\rangle_{\rm OS}
   =\overline{\mathbb E F(s_+)}\,\mathbb E G(s_+).
 \label{eq:gauge-observable-extension-criteria-GEX36}
\end{equation}
It is positive semidefinite, so this is genuinely a
reflection-positive probability system.  But $F=s_+$ is
null.  The positive-time function $Y=s_+$ has ordinary
supremum norm one, whereas $YF=1$ has OS norm one.  Thus
multiplication by $Y$ does not even preserve the OS null
space.  There can be no bound
$\|[YF]\|_{\rm OS}\le C\|[F]\|_{\rm OS}$ for this
insertion.  This example asserts only the insufficiency of
ordinary boundedness and reflection positivity; it is not
assigned the continuous-time smoothness hypotheses of the
scalar reconstruction.

An added sharp-time point type, with spatial test
$h\in\mathcal S(\mathbb R^3)$, instead needs its precise
time-shifted two-exterior form $B_{s,\alpha,h}(F,G)$,
defined from its actual Euclidean readout, with a bound
\begin{equation}
 |B_{s,\alpha,h}(F,G)|
    \le C_{\alpha,s}p_M(h)
                    \|[F]\|_{\rm OS}\|[G]\|_{\rm OS}.
 \label{eq:gauge-observable-extension-criteria-GEX37}
\end{equation}
Here $p_M$ is the weighted Schwartz seminorm on
$\mathbb R^3$.  The form is antilinear in $F$ and linear
in $G$ and $h$.  A different label space needs its own
specified continuous seminorm.
This bound makes it vanish if either exterior vector is
null, so it descends to the actual quotient.  Density and
Riesz then give a unique bounded operator
$\Phi_{s,\alpha}(h)$ on that specified exterior Hilbert
space, with norm at most $C_{\alpha,s}p_M(h)$.
The reflection adjoint, spatial covariance and heat
consistency must be verified from the same placements.
When that space has an established self-adjoint contraction
heat semigroup $T_r$,
an exact form identity
$B_{s+r}(F,G)=B_s(T_rF,T_rG)$ gives
$\Phi_{s+r}=T_r\Phi_sT_r$ by density and boundedness.

If the exterior space is an ordered sector inside a larger
OS space, matrix elements alone specify a compression.
The actual vector-sewing identity is needed to show that
the insertion maps ordered words to the ordered sector;
the star partner's sewing and adjoint identities then
give reduction.  A sharp-time
definition also needs its actual narrow-time limit or
density estimate.  If a proved bound in \eqref{eq:gauge-observable-extension-criteria-GEX37} has a
polynomial small-$s$ loss, the heat-consistent energy-band
construction of EB supplies the subsequent field method,
with its explicitly checked test continuity and domain
orders.  Point tubes and Lorentz continuation still need
their actual density and covariance checks.  For a new
type none of these model inputs follows
from $|w_\gamma|\le1$ or from the scalar insertion estimate.
This separates the bounded cutoff observable, the OS
Riesz operator, and the eventual Minkowski field with
their exact intervening requirements.

\subsubsection*{The resulting scope of the $\operatorname{tr}(F^2)$ construction}

The established scalar theory has its unrestricted
tempered distributions, positive cyclic representation,
forward spectrum, Lorentz covariance, local field
relations and the exact phase/Hilbert decomposition within
the stated scalar scope.  The finite-cutoff gauge analysis
also gives an actual complete Wilson observable family and
the precise information lost by taking only plaquette
actions.  These are substantial compatible conclusions.

A full Yang--Mills identification additionally needs the
sufficient joint gauge-observable family and its actual
continuum readouts, including mixed limits, relations,
positivity, covariance, locality and the stated relation
to the scalar density.  The conditional extension theorem
organizes those requirements and proves what follows
from them; it does not verify their model estimates.
Nontriviality of the scalar field guarantees a nontrivial
scalar subspace in any extension satisfying \eqref{eq:gauge-observable-extension-criteria-GEX2}, but
does not specify all its added degrees of freedom or their
dynamics.


% END MODULE chapters/gauge_observable_extension_criteria.tex

% BEGIN MODULE chapters/reconstruction_audit_chapter.tex
\subsection{Conditional scalar reconstruction}
\label{ch:precise-reconstruction}
\subsubsection{Reconstruction: hypotheses, consequences, and scope}

All growth constants below are positive, factorial exponents are
nonnegative, and derivative-order coefficients are positive integers.
D2 and H1 denote the conditional graded-moment and diagonal-flat growth
conclusions. G2 denotes scalar rotation covariance derived under the
ordinary-source, mirror, cone and ultraviolet hypotheses stated below.

The passage from Euclidean correlation functions to a quantum field
requires several logically separate inputs. Compactness of separated
correlations gives an off-diagonal distribution. Reflection positivity
gives a Hilbert space. Uniform growth in the number of arguments is
needed for the corrected reconstruction theorem. Full Euclidean
covariance and clustering have to be established in their own right.
Keeping these steps separate prevents a qualitative convergence
statement from carrying quantitative information it does not contain.

\paragraph*{The Euclidean test domain}

Write

\begin{gather*}
\FitMath{D_n=\{(x_1,\ldots,x_n)\in(\mathbb R^4)^n:
             x_i\ne x_j\text{ whenever }i\ne j\}.}
\end{gather*}

Let $\mathcal S_{\ne,n}$ be the closed subspace of the Schwartz space
$\mathcal S(\mathbb R^{4n})$ consisting of functions whose every derivative
vanishes on the complement of $D_n$. These functions are flat on all
coincidence diagonals. Let $\mathcal S_{+,n}$ consist of those elements
that are also flat outside the region in which all time coordinates
are positive. Set $\mathcal S_{\ne,0}=\mathcal S_{+,0}=\mathbb C$.
The time reflection is $\theta(x_1,x_2,x_3,x_4)
=(x_1,x_2,x_3,-x_4)$, and

\begin{gather*}
\FitMath{(\Theta\overline\varphi)(x_1,\ldots,x_n)
       =\overline{\varphi(\theta x_1,\ldots,\theta x_n)}.}
\end{gather*}

Symmetry of the functions makes reversal of the order of the arguments
immaterial in this scalar setting.

These are valid test domains for the historical
Osterwalder--Schrader reconstruction problem. They do not assert that
the Euclidean functions have a preferred extension to every Schwartz
test. A distribution supported on a coincidence diagonal vanishes on
$\mathcal S_{\ne,n}$. Such a distribution therefore cannot be recovered
from separated correlations. Any later use of contact terms, composite
operator identities at coincident points, or local Ward identities has
to specify the additional extension that it uses. This observation
does not obstruct reconstruction on the stated domain.

\begin{lemma}[Density of separated tensors]\label{lem:separated-tensor-density}
Finite sums of tensors $f_1\otimes\cdots\otimes f_n$, with smooth
compactly supported $f_i$ having pairwise disjoint supports, are dense
in $\mathcal S_{\ne,n}$. The corresponding assertion holds in
$\mathcal S_{+,n}$ when all supports lie at positive times.
\end{lemma}

The proof, including the simultaneous infinity, diagonal and positive-time cutoffs and the product Fourier approximation, is given in \cref{ch:quantitative-flat-extension}.

The lemma is a statement about the flat subspace. It gives no density
claim in the full Schwartz space, where evaluation on the diagonal
is a continuous functional that vanishes on every such tensor.

\paragraph*{What convergence preserves}

Let $\sigma_j$ be centred Euclidean moment functionals at finite
cutoffs. Suppose that for every separated smooth tuple $X$ the
numbers $\sigma_j(X)$ converge to $\sigma(X)$. Suppose also that each
$\sigma_j$ is reflection positive for finite families of internally
separated positive-time tuples. Different tuples in a family may
overlap. This last qualification is essential: reflection positivity
is a matrix condition, not a condition on one tuple at a time.

\begin{lemma}[Closed reflection positivity]\label{lem:closed-reflection-positivity}
For positive-time tuples $X_1,\ldots,X_r$ and complex coefficients
$c_1,\ldots,c_r$,

\begin{gather*}
\FitMath{\sum_{i,k=1}^r\overline c_i c_k\,
              \sigma(\theta\overline X_i\cdot X_k)\ge0.}
\qedhere
\end{gather*}



\end{lemma}

\emph{Proof.}
Each $X_i$ is internally separated. Its reflected entries have
negative times, so $\theta\overline X_i\cdot X_k$ is a separated tuple for
every pair $(i,k)$. The displayed expression is the limit of the
corresponding finite sum for $\sigma_j$, all of which are
nonnegative.\hfill$\square$

Suppose, in addition, that the limiting functionals extend
continuously to $\mathcal S_{\ne,n}$. Reflection positivity then holds
on finite Schwartz sequences $\Phi=(\varphi_0,\ldots,\varphi_N)$ with
$\varphi_n\in\mathcal S_{+,n}$:

\begin{gather*}
\FitMath{\sum_{n,m=0}^N
 S_{n+m}(\Theta\overline\varphi_n\otimes\varphi_m)\ge0.}
\end{gather*}

Indeed, the mixed tensor belongs to $\mathcal S_{\ne,n+m}$.
Coincidences within one factor are killed by its flatness; a
coincidence across the two factors has time coordinate zero or lies
outside the allowed half-space of one factor, so is killed there.
Approximate each $\varphi_n$ by the tensors in \cref{lem:separated-tensor-density},
using a seminorm strong enough for all the finitely many mixed
terms, and pass to the limit. Permutation symmetry and invariance
under any fixed Euclidean isometry pass in exactly the same way
once they hold on separated tensors.

There is one further step when translations are initially known
only for vectors in a dense subgroup, for example the union of
nested lattice translation groups. Given an arbitrary vector $v$,
choose subgroup vectors $v_k\to v$. Translation is continuous on
the smooth tests, and the local bounds permit passage from
$S_n(\tau_{v_k}\varphi)=S_n(\varphi)$ to
$S_n(\tau_v\varphi)=S_n(\varphi)$. This argument does not generate
rotations from the finite hypercubic group.

\paragraph*{The quantitative reconstruction hypothesis}

Fix the seminorms

\begin{gather*}
\FitMath{|\varphi|_k=
 \max_{|\alpha|\le k}\sup_x
       (1+|x|^2)^{k/2}|\partial^\alpha\varphi(x)|.}
\end{gather*}

\begin{theorem}[Scalar reconstruction]\label{thm:reconstruction-main}
For reconstruction, assume a sequence of symmetric functionals
$S_n\in\mathcal S_{\ne,n}'$, with $S_0=1$, and constants
$A,B,k_0,p<\infty$, independent of $n$, such that

\refstepcounter{equation}
\begin{gather*}
\FitMath{|S_n(\varphi)|\le A B^n(n!)^p|\varphi|_{k_0n}.}\tag{\theequation}\label{eq:reconstruction-audit-RC1}
\end{gather*}

Assume full Euclidean invariance, reflection positivity on the
domains just described, and the cluster property for separated
positive-time groups of arguments. Under these hypotheses the
corrected Osterwalder--Schrader reconstruction theorem gives a
Hermitian scalar Wightman field with positive energy, locality,
Poincare covariance, and a unique vacuum in its cyclic
representation.
\end{theorem}

Reality in the reconstruction import follows here from
reflection positivity. The sesquilinear form on the finite
positive-time sequences is Hermitian by polarization.
Its degree-zero and degree-\(n\) entries give
\(S_n(\Theta\bar\varphi)=\overline{S_n(\varphi)}\)
on positive-time tests. A compact separated tensor can be
translated into the positive half-space. Translation
invariance of \(S_n\) on the tensor and on its reflected
conjugate then gives the identity before the translation.
Separated tensor density and continuity give reality on
the entire diagonal-flat space.

We apply the corrected reconstruction theorem of Osterwalder and Schrader \cite[Sect.~IV.1, Thm.~$E'\!\to R'$, p.~287]{OS75}, with the test space and growth hypothesis just stated.

The exact all-order seminorm comparison is proved in
\cref{thm:quantitative-os-seminorm-ledger}.
\Cref{ch:h3-energy-band-fields} gives the direct
field, tube-boundary and covariance construction.
\Cref{ch:quantitative-reconstruction-matching}
identifies its moments with those supplied by the
corrected reconstruction theorem, which gives locality.
The representation is the ordered sector; its common
domains and scalar phase decomposition are developed
in the reconstruction sections. The cited theorem applies
by the hypothesis checks given there.

The condition in \eqref{eq:reconstruction-audit-RC1} must not be replaced by the weaker assertion that
every $S_n$ is separately continuous with unspecified constants.
The growth of the distribution order and the constants in $n$
matters to the existence of tempered boundary values in the
analytic continuation. Likewise, a local off-diagonal kernel
theorem alone does not establish \eqref{eq:reconstruction-audit-RC1}.

\paragraph{The explicit rooted-source and dyadic extension route.}
The rooted-source and dyadic extension estimates give a quantitative
proof of \eqref{eq:reconstruction-audit-RC1}, with controlled dependence on $n$. \Cref{thm:graded-rooted-source} derives
moments directly from the ordinary source hypothesis D1 and the proved
G1 chessboard domination, using the actual domain geometry in \eqref{eq:graded-source-bound-GS1}--\eqref{eq:graded-source-bound-GS3}.
At scale $a_s=L^{-s}$, domains are unions of $Ma_s$ cubes, indexed by
those unions, and their source dependence occupies only a fixed number of
neighboring cube layers. Their size is the minimum cubical edge-tree
length meeting all cubes, with source-energy decay above its rooted
entropy rate. Here $M=L^{m_0}$, and the cofinal tori $\ell=2L^m$,
$m\ge m_0$, tile every actual scale with even cube counts. No scale is
discarded for a failure to tile. The weak abstract \eqref{eq:all-state-source-transfer-ST6} count alone does
not furnish the sharper rooted and global counts needed for this proof.

For real compact $C^1$ tests, the theorem gives one fixed finite integer
$k$, $B(R)=C(R+1)^4$, and the two-state centered bound

\refstepcounter{equation}
\begin{gather*}
\FitMath{\left|\rho\prod_{i=1}^n
   [\mathcal O_a(f_i)-\nu(\mathcal O_a(f_i))]\right|
 \le B(R)^n(n!)^k\max\{1,d^{-5n}\}
                                 \prod_i\|f_i\|_\sharp.}\tag{\theequation}\label{eq:reconstruction-audit-RC2}
\end{gather*}

Each support lies in its own radius-$R$ ball, placed anywhere; the
separation is $d$. The same factorial exponent bounds higher cumulants.
Uniformity lets \eqref{eq:reconstruction-audit-RC2} pass to the separated continuum limits and their
component limits. This is requirement D2 conditional on the actual D1
source representation and G1, not a new independent analytic estimate.

\Cref{thm:quantitative-flat-extension} uses only the unit-radius
version, with the fixed equivalence of the sharp and $C^1$ norms absorbed
in $B(1)$. Its product Fourier construction uses $8n$ derivatives;
Gevrey-two pair/grid cutoffs cost at most $C^n(n!)^{32}$.
Dyadic Taylor flatness to order $10n+2$ pays the local $h^{-6n}$ cost
and the $h^{-4n}$ configuration grid count, leaving a summable dyadic
square. It yields \eqref{eq:reconstruction-audit-RC1} with

\begin{gather*}
\FitMath{k_0=20,} \\[3pt]
\FitMath{p=k+32.}
\end{gather*}

The proof includes independence of tensor presentations, overlap gluing,
cutoffs at infinity and every pairwise diagonal, and the positive-time
boundary. Its constrained tensor density supplies the full-sequence
reflection-positive form. Thus H1 is a proved implication of D2 here;
the ordinary analytic RG source estimates in D1 are still assumptions.
Neither separate continuity with arbitrary order constants nor weak counting in \eqref{eq:all-state-source-transfer-ST6} alone may be substituted for this GS/QF route.

\paragraph*{Hilbert space and clustering}

\Cref{ch:euclidean-hilbert-bridge} proves the quotient construction
on compact positive-time tuples from their invariant reflection-positive
moments and position-uniform per-ball bounds. Fixed separation of the
reflected supports bounds the translated orbits before iterated
Cauchy--Schwarz proves contraction. The resulting time semigroup is
strongly continuous, positive and self-adjoint, and spatial translations
are strongly continuous unitaries. The stated lattice approximation also
gives the cube spectral cone. The GS/QF growth and density argument then
identifies its completion with the full positive-time Schwartz sequence
quotient. These deductions retain their explicit hypotheses; positivity
alone is not being substituted for the continuity and growth checks.

For completeness, clustering on a dense class of vectors extends
to the whole Hilbert space. Let $U(v)$ denote the spatial unitary
translations, and suppose that for dense vectors $\xi,\eta$,

\begin{gather*}
\FitMath{\langle\xi,U(v)\eta\rangle\longrightarrow
       \langle\xi,\Omega\rangle\langle\Omega,\eta\rangle} \\[3pt]
\FitMath{(|v|\to\infty).}
\end{gather*}

Choose dense approximants $\xi_k\to\xi$ and $\eta_k\to\eta$.
Unitarity bounds the difference of the matrix elements, uniformly
in $v$, by

\begin{gather*}
\FitMath{\|\xi-\xi_k\|\,\|\eta\|
       +\|\xi_k\|\,\|\eta-\eta_k\|.}
\end{gather*}

The vacuum matrix elements obey the same estimate. First let
$|v|\to\infty$ at fixed $k$, then let $k\to\infty$.
The argument proves the qualitative cluster limit. It does not
preserve a numerical decay rate on arbitrary vectors without
further quantitative control of the approximations.

\paragraph*{Phases and full rotational covariance}

The phase-interface construction realizes the components with
the G1 mirror properties and the D1/GS/QF moment bounds. Once
the angular argument gives covariance of the mixture, the
\hyperlink{phase-isometries}{action on its canonical decomposition}
and coordinate reflection invariance give covariance of every
component. The state-uniform two-point extraction is used to
obtain the two-point law for the realizing states of each component.
The convergent cutoff one-point differences give the real
phase mean \(c_\xi\int f\), with \(\sup_\xi|c_\xi|\)
bounded by the one-point transfer. To restore it on the
Schwartz domains, let \(I_A\varphi\) be the integral of
\(\varphi\) over the variables with indices in \(A\).
Differentiation under the integral preserves flatness on
every remaining coincidence diagonal. If the original test
is supported at positive ordered times, then so is the
marginal on the retained arguments, with flatness at its
boundary. Indeed any violation of a retained order or
positivity constraint already violates the original one.

Take \(k_0>4\), increasing the seminorm order if needed.
Writing \(x=(y,z)\) with \(z=(z_i)_{i\in A}\), use
\[
 (1+|x|^2)^{k_0 n/2}
 \ge (1+|y|^2)^{k_0(n-|A|)/2}
       \prod_{i\in A}(1+|z_i|^2)^{k_0/2}.
\]
The derivatives of order at most \(k_0(n-|A|)\) in
\(y\) therefore give
\(\lvert I_A\varphi\rvert_{k_0(n-|A|)}
\le C_{\rm int}^{|A|}\lvert\varphi\rvert_{k_0 n}\),
where \(C_{\rm int}=\int_{\mathbb R^4}(1+|z|^2)^{-k_0/2}\,dz\).
For \(A\) the full set of indices this also bounds the
vacuum integral. The restored moment is
\[
 S_{\xi,n}(\varphi)=
 \sum_{A\subset\{1,\ldots,n\}} c_\xi^{|A|}
                  S^{\rm c}_{\xi,n-|A|}(I_A\varphi).
\]
It agrees with the finite subset expansion on separated
tensors, hence with the continuous extension by density.
For a centered bound with exponential base \(B\) and
nonnegative factorial exponent \(p\), the binomial sum gives
the same exponent and seminorm order with base
\(B+C_{\rm int}\sup_\xi|c_\xi|\), after including the
vacuum in the prefactor. On the ordered sequences the
same marginal maps define inverse shifts with \(c_\xi\)
and \(-c_\xi\). The reflected pairing identity on separated
tensors extends by continuity and gives the same
intertwining unitary as on compact tuples. On the cyclic
field domain the centered insertion subtracts
\(c_\xi\int f\) times the identity.

Similarly, invariance under the hypercubic group, coordinate
reflection positivity, and translation invariance do not by
themselves imply invariance under $\operatorname{O}(4)$. The finite group
is not dense in the rotation group. Any route through diagonal
mirrors, a sharp energy--momentum cone, bounded insertions, or
the exclusion of nonconstant angular harmonics needs an additional
theorem. The following chain now supplies that theorem conditionally,
before the external reconstruction theorem is invoked.

\Cref{thm:all-orders-angular-closure} completes the scalar
covariance implication. Its preceding sections prove, in order, the
uniform Hilbert insertion bound, the analytic $L^1$ density, the product
kernels and continued Gram identities, and the ultraviolet elimination
of the two-point harmonics. Polygon divisibility and common coordinate
cuts then remove all higher harmonics and exceptional distributional
support. The proof works with finite support boxes; density of generic
point configurations alone would not exclude a singular part.

The angular continuation uses
\eqref{eq:angular-two-point-elimination-TU18}. The dependency
catalogue obtains it through the positive-shape bound
\eqref{eq:angular-two-point-elimination-TU10}, deduced from F2--F3
on the stated all-state test class. The second route in
\cref{sec:angular-mirror-budget} uses the bound for a fixed
family of reflected bumps from the ordinary source estimates
at the running stops. Both deductions apply to periodic-state
limits with the coordinate and diagonal mirrors and the exact cone.
The \hyperlink{growing-torus-recognition}{finite-window density comparison}
includes growing tori above its side threshold. General growing-torus
limits require their own mirror and cone argument.

The application uses the interacting D1 estimates and the F2--F3
extraction and normalization lemmas in their full forms. These are the
model statements invoked to obtain the nontrivial scalar theory. The OS analytic-continuation theorem is still imported,
and the identification with full Yang--Mills observables remains separate.

\paragraph*{What the reconstructed field establishes}

When the limiting observable is the centred scalar $\operatorname{tr}(F^2)$,
the theorem above constructs its scalar observable sector. A
nonzero two-point function on separated tests excludes a field
that is a multiple of the identity. A nonzero higher cumulant
excludes Gaussianity of that scalar field. Neither conclusion
alone identifies a full Yang--Mills observable algebra. In
particular, a Wick square of a free Gaussian field already has
nonzero higher cumulants. Establishing a Yang--Mills theory in a
stronger sense requires the stated additional observables,
identities, and compatibility relations that define that sense.

Nor does polynomial clustering establish a mass gap. The positive
measure $E e^{-E}\,dE$ on $[0,\infty)$ has Laplace transform

\begin{gather*}
\FitMath{\int_0^\infty E e^{-E}e^{-tE}\,dE=(1+t)^{-2}.}
\end{gather*}

Its support reaches zero, although its transform decays
quadratically. A positive spectral gap requires an estimate that
excludes spectral support in an entire interval above the vacuum.
That question remains separate from the existence and
reconstruction of the scalar sector described here.


% END MODULE chapters/reconstruction_audit_chapter.tex
\appendix
\section{Auxiliary lattice and comparison estimates}\label{app:calculations}
These calculations supply the finite-dimensional normalization, kernel decay, polymer estimates and reflection-positive limits used in the main argument. Each result retains its own parameter range and boundary conditions. The comparison calculations at the end specify the domains needed when the coupling or cutoff changes.

% BEGIN MODULE chapters/block_normalization_v2.tex
\subsection{Block normalization and a uniform scalar gap}
\subsubsection{Block means, adjoints, and a uniform scalar Poincar\'e inequality}

An adjoint depends on the inner products of both its domain and its range.
This matters when a block mean maps a fine lattice to a coarse lattice. The
following finite-dimensional calculation fixes the convention and gives the
scalar estimate used in a gauge-fixed Gaussian comparison. It applies to real
or complex scalar functions; a finite number of color components are handled
by summing the same estimate component by component.

Fix an integer $m\ge1$ and put $\eta=m^{-1}$. A block consists of the sites

\begin{gather*}
\FitMath{B=\{0,\eta,\ldots,(m-1)\eta\}^{4}.}
\end{gather*}

Its associated physical cell has side $m\eta=1$ and volume one. Translation of
the block, including translation by half a fine spacing to use cell centers,
does not change the calculation. Let $\mathcal B$ be a finite disjoint family
of such blocks, for example the blocks of a periodic torus whose side is an
integer number of coarse cells. Write $\Lambda_f$ for their union and
$\Lambda_c$ for the set indexing the blocks. A coarse function has one value
per block. We count every fine bond once, with positive orientation. In the
block Neumann form only bonds whose two endpoints belong to the same block
are counted.

\phantomsection\label{notdef:block-inner-products}
The physical inner products are

\begin{gather*}
\FitMath{\langle f,g\rangle_f=\eta^4\sum_{x\in\Lambda_f}\overline{f(x)}g(x),} \\[3pt]
\FitMath{\langle u,v\rangle_c=\sum_{B\in\mathcal B}\overline{u(B)}v(B),} \\[3pt]
\FitMath{\langle a,b\rangle_e=\eta^4\sum_{e}\overline{a(e)}b(e).}
\end{gather*}

The last sum can mean all fine bonds of the periodic lattice or just the
internal bonds of the blocks; the operator being considered determines which
set is used. The same weight is used on fine sites and fine bonds. The physical
difference is
$\mathrm{d}_\eta f(x,x+\eta e_\mu)=\eta^{-1}[f(x+\eta e_\mu)-f(x)]$.
The unscaled combinatorial difference is $\mathrm{d}_0=\eta \mathrm{d}_\eta$. These two operators
must not be interchanged in an estimate with a mesh-independent constant.

\phantomsection\label{notdef:block-mean}
\begin{proposition}[The mean, its adjoint, and its projection]
Define the block mean $Q:H_f\to H_c$ by

\begin{gather*}
\FitMath{(Qf)(B)=\eta^4\sum_{x\in B}f(x)=m^{-4}\sum_{x\in B}f(x).}
\end{gather*}

For the physical inner products just specified, its adjoint is
$(Q^*u)(x)=u(B)$ for $x\in B$. Consequently $QQ^*=1$, and
$P=Q^*Q$ is the orthogonal projection onto the functions constant in each
block. In particular,

\begin{gather*}
\FitMath{\|f\|_f^2=\|(1-P)f\|_f^2+\|Qf\|_c^2.}\qedhere
\end{gather*}

\end{proposition}

\begin{proof}
Directly from the definitions,

\begin{gather*}
\FitMath{\langle Qf,u\rangle_c
 =\sum_B\eta^4\sum_{x\in B}\overline{f(x)}u(B)
 =\langle f,Q^*u\rangle_f.}
\end{gather*}

Since a block contains $m^4$ sites and $\eta^4m^4=1$, its mean applied to a
constant is that constant. Thus $QQ^*=1$. It follows that $P^*=P$ and
$P^2=Q^*(QQ^*)Q=P$. Its range is precisely the block-constant functions, and
the displayed norm identity is the orthogonal decomposition by that projection.
No limiting argument is involved.
\end{proof}

\begin{proposition}[A block Neumann estimate]
Let $\mathrm{d}_{\eta,N}$ denote the physical difference restricted to internal bonds
of each block. If $m\ge2$, then

\begin{gather*}
\FitMath{\|(1-P)f\|_f^2
 \le \frac{\eta^2}{4\sin^2(\pi\eta/2)}
                 \|\mathrm{d}_{\eta,N}f\|_e^2
 \le \tfrac14\|\mathrm{d}_{\eta,N}f\|_e^2.}
\end{gather*}

For $m=1$ the left side is zero. If $\mathrm{d}_\eta$ includes bonds between blocks,
the same bound holds with $\mathrm{d}_\eta$ on the right. With the combinatorial
difference the exact first bound is instead

\begin{gather*}
\FitMath{\|(1-P)f\|_f^2
 \le \frac{1}{4\sin^2(\pi/(2m))}\|\mathrm{d}_{0,N}f\|_e^2.}\qedhere
\end{gather*}

\end{proposition}

\begin{proof}
First consider a one-dimensional path with $m$ vertices indexed by
$j=0,\ldots,m-1$. Its Neumann graph Laplacian is the matrix associated with
the quadratic form $\sum_{j=0}^{m-2}|f_{j+1}-f_j|^2$. For an interior
vertex it acts as $2f_j-f_{j-1}-f_{j+1}$; at the two endpoints it acts as
$f_0-f_1$ and $f_{m-1}-f_{m-2}$. For $k=0,\ldots,m-1$, set

\begin{gather*}
\FitMath{v_k(j)=\cos\bigl(\pi k(j+\tfrac12)/m\bigr).}
\end{gather*}

The cosine addition formula gives, in the interior,

\begin{gather*}
\FitMath{2v_k(j)-v_k(j-1)-v_k(j+1)
       =4\sin^2(\pi k/(2m))v_k(j).}
\end{gather*}

At the left endpoint,

\begin{gather*}
\FitMath{\cos(\theta/2)-\cos(3\theta/2)
 =4\sin^2(\theta/2)\cos(\theta/2),} \\[3pt]
\FitMath{\theta=\pi k/m.}
\end{gather*}

The right endpoint follows by reflection. Thus these are eigenvectors with
eigenvalues $4\sin^2(\pi k/(2m))$. The eigenvalues are distinct because sine
is strictly increasing on $[0,\pi/2)$; the nonzero eigenvectors therefore form
an orthogonal basis. The $k=0$ vector is constant.

On the four-dimensional block the graph Laplacian is the sum of the four
one-dimensional matrices acting on separate tensor factors. Tensor products
of the $v_k$ therefore form an eigenbasis. The eigenvalues of the physical
Neumann Laplacian $\mathrm{d}_{\eta,N}^*\mathrm{d}_{\eta,N}$ are

\begin{gather*}
\FitMath{4\eta^{-2}\sum_{\mu=1}^4\sin^2(\pi k_\mu/(2m)),} \\[3pt]
\FitMath{0\le k_\mu<m.}
\end{gather*}

Its only zero mode is the constant vector, and the smallest positive eigenvalue
is $\lambda_N=4\eta^{-2}\sin^2(\pi/(2m))$. Multiplication of the site and
bond inner products by the common constant $\eta^4$ does not change this
matrix or these eigenvalues. Expanding $f-Pf$ in the eigenbasis gives
$\lambda_N\|f-Pf\|_f^2\le\|\mathrm{d}_{\eta,N}f\|_e^2$ on one block and hence on
their disjoint union.

Finally, concavity of sine on $[0,\pi/2]$ gives
$\sin t\ge2t/\pi$ there. At $t=\pi/(2m)$ this yields
$\sin(\pi/(2m))\ge1/m=\eta$, so $\lambda_N\ge4$.
Adding the nonnegative contributions of bonds between blocks can only increase
the Dirichlet form. Substitution of $\mathrm{d}_{0,N}=\eta \mathrm{d}_{\eta,N}$ gives the
combinatorial version.
\end{proof}

\phantomsection\label{notdef:scalar-gap}
These estimates immediately give a uniform scalar gap on a periodic fine
lattice. If $\mathcal L_\eta=\mathrm{d}_\eta^*\mathrm{d}_\eta$ and
$\mathcal A=\mathcal L_\eta+Q^*Q$, then

\begin{gather*}
\FitMath{\langle f,\mathcal A f\rangle_f
 =\|\mathrm{d}_\eta f\|_e^2+\|Qf\|_c^2
 \ge\|(1-P)f\|_f^2+\|Qf\|_c^2=\|f\|_f^2.}
\end{gather*}

Thus $\mathcal A$ is invertible and $\|\mathcal A^{-1}\|\le1$, independently
of the mesh and number of blocks. This proves a scalar block-constrained gap.
It supplies neither a vector-field coercivity theorem nor a mass gap for a
physical Hamiltonian; those are different assertions.

\phantomsection\label{notdef:scalar-projection}
\paragraph*{A uniform inverse for a coarse Gram matrix}

The same convention also makes a useful projection formula precise. On the
periodic lattice put $G'=\mathcal A^{-1}$, with $\mathcal A$ as above. There
is a numerical constant $c>0$, independent of $m$ and the number of blocks,
such that

\begin{gather*}
\FitMath{QG'^2Q^*\ge c1.}
\end{gather*}

Here is a finite-dimensional proof, including a choice of the test functions.
For $m\ge2$ define on a block indexed by $j=(j_1,\ldots,j_4)$

\begin{gather*}
\FitMath{\phi_m(j)=16\prod_{\mu=1}^4
             \sin^2\bigl(\pi(j_\mu+\tfrac12)/m\bigr).}
\end{gather*}

The sum of the $m$ roots $e^{2\pi i(j+1/2)/m}$ is zero. Taking real parts
in $\sin^2 t=(1-\cos 2t)/2$ shows that the one-dimensional arithmetic
mean of each sine-square factor is $1/2$. Thus the mean of $\phi_m$ on a
block is one. For a coarse function $u$, define $Tu(x)=u(B)\phi_m(x)$ on
the block $B$. Then $QT=1$ and $\|T\|\le16$.

There is also a bound $\|\mathcal L_\eta Tu\|_f\le C\|u\|_c$ uniform in
$m$. To see it, the second physical difference of
$\sin^2(\pi t)$ at interior sites is bounded by the supremum $2\pi^2$ of
its second derivative. At an end site, its value is
$\sin^2(\pi\eta/2)\le\pi^2\eta^2/4$. A second difference that crosses a
block boundary therefore contributes at most a constant times the absolute
coarse values in the two neighboring blocks: the factor $\eta^{-2}$ is
canceled by this endpoint bound. The other three sine-square factors are
bounded by one. After summing the four coordinate differences, a pointwise
bound is $C\sum_{B'\sim B}|u(B')|$, where the sum includes $B$ and its
face neighbors. Squaring, summing with weight $\eta^4$, and using the bounded
number of neighbors proves the norm estimate. For $m=1$, set $T=1$; the
unit-spacing Laplacian has norm at most $16$, so the same conclusion holds.
Consequently $\|\mathcal A Tu\|_f\le C\|u\|_c$ for all $m$.

Self-adjointness and $G'\mathcal A=1$ now give

\begin{gather*}
\FitMath{\|u\|_c^2=\langle Q^*u,Tu\rangle_f
          =\langle G'Q^*u,\mathcal A Tu\rangle_f
          \le C\|G'Q^*u\|_f\|u\|_c.}
\end{gather*}

This proves the claimed positive lower bound for $QG'^2Q^*$.

The map $V=G'Q^*$ is injective, its Gram matrix is $V^*V=QG'^2Q^*$,
and elementary linear algebra shows that

\begin{gather*}
\FitMath{R=1-G'Q^*(QG'^2Q^*)^{-1}QG'}
\end{gather*}

is the orthogonal projection onto $\ker(QG')$. This kernel equals
$\mathcal L_\eta\ker Q$. Indeed, if $Qv=0$ then
$\mathcal A v=\mathcal L_\eta v$, so $QG'\mathcal L_\eta v=0$.
Conversely, if $QG'h=0$, put $v=G'h$. Then $Qv=0$ and
$h=\mathcal A v=\mathcal L_\eta v$. The expression for $R$ is therefore
well defined with a uniformly invertible coarse Gram matrix. Every adjoint in
this argument uses the physical inner products fixed above.

\paragraph*{Translation to plain Euclidean coordinates}

There is an equivalent convention with unweighted coordinate inner products,
but the averaging matrix has to change. Let $S_f f=\eta^2 f$ on fine sites,
$S_e a=\eta^2 a$ on fine bonds, and $S_cu=u$ on coarse sites. These maps are
unitaries from the physical spaces just defined to the spaces with plain
coordinate sums. The correctly conjugated averaging matrix is

\begin{gather*}
\FitMath{\widetilde Q=S_cQS_f^{-1},} \\[3pt]
\FitMath{(\widetilde Q v)(B)=\eta^2\sum_{x\in B}v(x),} \\[3pt]
\FitMath{(\widetilde Q^*u)(x)=\eta^2u(B).}
\end{gather*}

Each row has $m^4$ entries of size $\eta^2=m^{-2}$, so its squared Euclidean
length is one. Rows have disjoint supports, proving
$\widetilde Q\widetilde Q^*=1$. The conjugated projection is
$\widetilde P=\widetilde Q^*\widetilde Q=S_fPS_f^{-1}$. Because sites and
bonds have the same physical weight,
$S_e \mathrm{d}_\eta S_f^{-1}$ is the same physical difference matrix
$\eta^{-1}(\text{shift}-1)$. Therefore the entire Poincar\'e estimate survives
unitary conjugation with its exact constant.

In contrast, suppose that one retains the arithmetic mean matrix
$\overline Qv=\eta^4\sum_{x\in B}v(x)$ while using plain inner products on
both lattices. Its plain adjoint is
$(\overline Q^*u)(x)=\eta^4u(B)$, and

\begin{gather*}
\FitMath{\overline Q\overline Q^*=\eta^4 1,} \\[3pt]
\FitMath{P=\eta^{-4}\overline Q^*\overline Q.}
\end{gather*}

Thus $\overline Q^*\overline Q$ itself is not the projection. For a nonzero
constant fine function and $m>1$, the expression
$\|f-\overline Q^*\overline Q f\|^2$ is positive whereas $\mathrm{d}_\eta f=0$.
This gives an immediate check against accidentally combining the arithmetic
mean with the wrong adjoint. Either the physical weighted convention or the
conjugated unweighted convention is consistent; mixing them is not.

\paragraph*{Consequences for cochains and Gaussian formulas}

When a Gaussian curvature is first defined using the plain cochain inner
product, its covariance is an orthogonal projection for that inner product.
After changing to physical coordinates, transform the random vector and the
covariance together. For a linear random vector $F$ with covariance $C$ in
plain coordinates and a linear coordinate change $T$, the covariance of $TF$
is $TCT^*$. An orthogonal projection is carried to an orthogonal projection
when the coordinate change is interpreted as the unitary between the two
specified Hilbert spaces; an arbitrary rescaling of its coordinate entries
without the corresponding inner product change has a different covariance.

For a Hilbert-space operator $B:H_1\to H_2$, unitary changes of coordinates
$S_1,S_2$ give

\begin{gather*}
\FitMath{\widetilde B=S_2BS_1^{-1},} \\[3pt]
\FitMath{\widetilde B^*=S_1B^*S_2^{-1}.}
\end{gather*}

In particular products, orthogonal projections, positive inverses and
Moore--Penrose inverses commute with this conjugation. A conditional Gaussian
covariance written as
$C-CB^*(BCB^*)^+BC$ remains the same Hilbert-space operator after every factor
is conjugated. This observation permits a proof to begin in plain finite
cochain coordinates and then use physical block coordinates. The change must
be announced before the first adjoint, constraint, or Poincar\'e estimate that
uses the latter convention. Every discrete derivative must also say whether
it is $\mathrm{d}_0$ or $\mathrm{d}_\eta$. These declarations prevent a missing power of the mesh
from being hidden inside an adjoint symbol.


% END MODULE chapters/block_normalization_v2.tex

% BEGIN MODULE chapters/block_decay_v2.tex
\subsubsection{A scalar block resolvent with uniform off-diagonal decay}

This finite-dimensional lemma supplies a scalar part of the propagator input.
All distances are physical distances. The argument uses a relative form bound
and a Neumann series; the exponentially conjugated operator is generally not
self-adjoint.

Let a four-dimensional periodic lattice of spacing $\eta=1/m$, $m\ge1$, be
partitioned into blocks of $m^4$ sites and physical side one. Its torus has an
integer number of blocks in each direction. On site functions use
$\langle u,v\rangle=\eta^4\sum_x\overline{u_x}v_x$. On positively oriented
fine bonds use the same weight. Put

\begin{gather*}
\FitMath{(\mathrm{d}_\eta u)(x,x+\eta e_\mu)
       =\eta^{-1}(u_{x+\eta e_\mu}-u_x),} \\[3pt]
\FitMath{\mathcal L_\eta=\mathrm{d}_\eta^*\mathrm{d}_\eta.}
\end{gather*}

The block mean $Q$ takes the arithmetic mean of the $m^4$ fine values. Its
coarse inner product is the unweighted sum, so $Q^*$ embeds a coarse value as
a constant on its block. Thus $P=Q^*Q$ is an orthogonal projection and its
matrix entries, as an operator acting on coordinate vectors, are
$P_{xy}=\eta^4$ when $x,y$ are in the same block and zero otherwise.

Write $\operatorname{dist}_\eta$ for shortest-path distance on the periodic
fine graph, with length $\eta$ per bond. The distance includes wrapping paths.
The diameter of a single block in this metric is at most $4$.
The block Poincar\'e inequality proves

\begin{gather*}
\FitMath{A_\eta:=\mathcal L_\eta+P\ge1,} \\[3pt]
\FitMath{\|u\|\le\|A_\eta^{1/2}u\|,} \\[3pt]
\FitMath{\|\mathrm{d}_\eta u\|\le\|A_\eta^{1/2}u\|.}
\end{gather*}

These statements also hold when $m=1$, since then $P=1$.

\begin{proposition}[A mesh-independent Combes--Thomas bound]
Let $w$ be a real function with $|w_x-w_y|\le\eta$ on every neighboring pair.
For real $t$ let $M_t$ multiply by $e^{tw_x}$ and set
$A_t=M_tA_\eta M_t^{-1}=A_\eta+E_t$. There is a numerical constant $C$,
independent of the mesh, volume and function $w$, such that, for $|t|\le1$,

\begin{gather*}
\FitMath{\|A_\eta^{-1/2}E_t A_\eta^{-1/2}\|
                  \le C(|t|+t^2).}
\end{gather*}

Consequently one may choose a fixed $t_0>0$, independent of those parameters,
so that $\|A_t^{-1}\|\le2$ for $|t|\le t_0$. For arbitrary nonempty sets of
sites $X,Y$ this gives

\begin{gather*}
\FitMath{\|1_X A_\eta^{-1}1_Y\|
 \le2\exp[-t_0\operatorname{dist}_\eta(X,Y)].}\qedhere
\end{gather*}

\end{proposition}

\begin{proof}
Set $A=A_\eta$. First consider the contribution $E_t^{\mathcal L}$ from the
Laplacian. On a positively oriented bond with endpoints $x,y$, put
$\delta=w_x-w_y$, $a=e^{t\delta}-1$ and $b=e^{-t\delta}-1$. The diagonal
entries of the conjugated edge Laplacian do not change. Its off-diagonal
contribution to $\langle u,E_t^{\mathcal L}v\rangle$ is therefore

\begin{gather*}
\FitMath{-\eta^4\eta^{-2}
             (a\overline{u_x}v_y+b\overline{u_y}v_x).}
\end{gather*}

The exact identity

\begin{gather*}
\FitMath{a\overline{u_x}v_y+b\overline{u_y}v_x
 =a(\overline{u_x}-\overline{u_y})v_y
   +b\overline{u_y}(v_x-v_y)+(a+b)\overline{u_y}v_y}
\end{gather*}

exposes a difference of $u$ or $v$ in the terms linear in $t$. Since
$|\delta|\le\eta\le1$ and $|t|\le1$,

\begin{gather*}
\FitMath{|a|,|b|\le e|t|\eta,} \\[3pt]
\FitMath{|a+b|=2|\cosh(t\delta)-1|\le e t^2\eta^2.}
\end{gather*}

Each vertex occurs as the head of at most four positively oriented bonds.
Cauchy--Schwarz applied to the first term, summed over the bonds, consequently
gives $2e|t|\|\mathrm{d}_\eta u\|\|v\|$. The second term is bounded by
$2e|t|\|u\|\|\mathrm{d}_\eta v\|$. The last term is bounded by
$4e t^2\|u\|\|v\|$. In particular,

\begin{gather*}
\FitMath{|\langle u,E_t^{\mathcal L}v\rangle|
 \le 2e|t|(\|\mathrm{d}_\eta u\|\|v\|+\|u\|\|\mathrm{d}_\eta v\|)
                         +4e t^2\|u\|\|v\|.}
\end{gather*}

The two powers of $\eta^{-1}$ in the Laplacian have been canceled by the
Lipschitz difference and the difference of a test vector. Estimating the
off-diagonal matrix entries directly would miss this cancellation.

For the block term, the neighbor condition implies
$|w_x-w_y|\le4$ for sites in one block. Its conjugation error has entries

\begin{gather*}
\FitMath{(E_t^P)_{xy}
 =\eta^4 1_{\{x,y\text{ in one block}\}}
                              (e^{t(w_x-w_y)}-1).}
\end{gather*}

Their absolute row sums and column sums are at most $4e^4|t|$, because a
block has $m^4$ sites and $m^4\eta^4=1$. The Schur test therefore yields
$\|E_t^P\|\le4e^4|t|$. The common site weight $\eta^4$ does not alter this
Schur estimate. Using the two energy bounds preceding the proposition gives

\begin{gather*}
\FitMath{|\langle u,E_t v\rangle|
 \le C(|t|+t^2)\|A^{1/2}u\|\|A^{1/2}v\|,}
\end{gather*}

where, for example, $C=8e+4e^4$ suffices. Substitute $u=A^{-1/2}f$ and
$v=A^{-1/2}g$ and take the supremum over unit vectors $f,g$. This proves
the asserted relative operator bound.

Choose $t_0=\min\{1,(4C)^{-1}\}$. Then
$B_t=A^{-1/2}E_t A^{-1/2}$ has norm at most $1/2$ when $|t|\le t_0$.
The exact factorization and its inverse are

\begin{gather*}
\FitMath{A_t=A^{1/2}(1+B_t)A^{1/2},} \\[3pt]
\FitMath{A_t^{-1}=A^{-1/2}(1+B_t)^{-1}A^{-1/2}.}
\end{gather*}

The Neumann series gives $\|(1+B_t)^{-1}\|\le2$ and hence
$\|A_t^{-1}\|\le2$. This step requires no assertion that $A_t$ is
self-adjoint or ordered below by a positive scalar.

Finally take $w_x=\operatorname{dist}_\eta(x,Y)$ and $t=t_0$. The triangle
inequality makes $w$ admissible, and $w=0$ on $Y$. If
$D=\operatorname{dist}_\eta(X,Y)$, then $w\ge D$ on $X$. Similarity of the
two inverses gives

\begin{gather*}
\FitMath{1_XA^{-1}1_Y=1_X M_t^{-1}A_t^{-1}M_t1_Y.}
\end{gather*}

The rightmost multiplication is the identity on $Y$, while the norm of the
leftmost multiplication is at most $e^{-t_0D}$. Combining these statements
with $\|A_t^{-1}\|\le2$ proves the final estimate.
\end{proof}

The estimate concerns operator norms between sets in the weighted lattice
Hilbert space. It does not assert a mesh-independent bound for a Green
function kernel normalized relative to the physical counting measure
$\eta^4\sum_x$. Converting a coordinate matrix to that kernel introduces its
own power of $\eta$. The scalar lemma likewise does not establish the
non-Abelian vector-field or multiscale background estimates. Those require
their additional gauge constraints and response bounds.

\paragraph*{Dirichlet restrictions and the distance being estimated}

The same estimate holds for a Dirichlet compression of this particular
operator, with constants independent of the shape of the region. To state it
without ambiguity, let $D$ be any set of sites of the periodic fine lattice.
Let $H_D$ consist of the functions supported in $D$, with the inherited weighted
inner product, and let $J_D:H_D\to H$ be extension by zero. Define

\begin{gather*}
\FitMath{A_D=J_D^*A_\eta J_D.}
\end{gather*}

In particular the block term is the compression of the \emph{original}
projection $P$. On a block intersected by the boundary it is not a new mean
taken over just the surviving sites. Replacing it by that new mean would
define a different operator and is not justified by this corollary.

For $u\in H_D$, the quadratic form is

\begin{gather*}
\FitMath{\langle u,A_Du\rangle
 =\|\mathrm{d}_\eta J_Du\|^2+\|QJ_Du\|_c^2\ge\|u\|^2.}
\end{gather*}

It controls the physical gradient of the zero extension, including boundary
bonds. Restrict an admissible ambient weight $w$ to $D$ and conjugate by its
exponential on $H_D$. Because multiplication by this exponential commutes with
the site indicator $1_D$, the new error form is exactly the restriction of
$E_t$ to zero-extended vectors. The estimates in the proof therefore give

\begin{gather*}
\FitMath{\|A_D^{-1/2}(M_tA_DM_t^{-1}-A_D)A_D^{-1/2}\|
 \le C(|t|+t^2).}
\end{gather*}

The same Neumann argument implies, for $X,Y\subset D$,

\begin{gather*}
\FitMath{\|1_XA_D^{-1}1_Y\|
 \le2e^{-t_0\operatorname{dist}_\eta(X,Y)}.}
\end{gather*}

Here the distance is computed in the ambient periodic lattice. It may use a
path outside $D$. No estimate in the possibly much larger intrinsic distance
of a region with holes has been asserted.

For comparison with continuum notation, the shortest-path distance is the
periodic physical Manhattan distance, and hence is no smaller than the
periodic Euclidean distance. The displayed bound therefore implies the same
bound with Euclidean distance in its exponent. One can alternatively take
$w$ to be Euclidean distance to $Y$ directly, since its change across a fine
bond is at most $\eta$. Both formulations concern physical distances; using
the number of bonds without the factor $\eta$ would incorrectly strengthen
the assertion as the mesh is refined.

The calculation explains the uniform mechanism: the scalar energy controls
one physical difference, the weight varies by only $O(t\eta)$ across an
edge, and its genuinely quadratic error is $O(t^2\eta^2)$. The block projection
has range at most one physical cell and row mass one. Consequently the
conjugation cost depends on the dimension and physical block size, with no
factor proportional to the number of fine sites in a block or the torus.


% END MODULE chapters/block_decay_v2.tex

% BEGIN MODULE chapters/polymer_theorem_v2.tex
\subsection{Convergence of the polymer logarithm}
\label{ch:polymer-exact}

The small-field integration must produce activities satisfying a summability
criterion. Once that estimate is available, the passage from activities to a
local logarithm is a combinatorial theorem. We prove that theorem here, so its
convergence mechanism is independent of the particular gauge-field
coordinates. The estimate on the activities remains a separate analytic task.

\subsubsection{Finite polymer systems}

A finite polymer system consists of a finite set $\mathcal P$ and a symmetric
incompatibility relation $\not\sim$ on it. Every polymer is incompatible with
itself. Give each $X\in\mathcal P$ a complex activity $z_X$. A compatible
family contains no two incompatible polymers. Its partition function is

\begin{gather*}
\FitMath{Z(z)=\sum_{\substack{\mathcal F\subset\mathcal P\\
                      \mathcal F\,\mathrm{compatible}}}
                              \prod_{X\in\mathcal F}z_X.}
\end{gather*}

The empty family contributes one. There are no questions of convergence in
this polynomial. The problem is to construct its logarithm and bound it in
terms of local activities without a loss exponential in $|\mathcal P|$.

For a list $\mathbf X=(X_1,\ldots,X_n)$, let $H(\mathbf X)$ be the simple
graph on the labeled vertices $1,\ldots,n$ whose edges are the pairs with
$X_i\not\sim X_j$. Repeated polymers are permitted in lists; their distinct
vertices are adjacent by self-incompatibility. Define

\begin{gather*}
\FitMath{\phi(\mathbf X)=
  \sum_{\substack{F\subset E(H(\mathbf X))\\([n],F)\,\mathrm{connected}}}
                       (-1)^{|F|},} \\[3pt]
\FitMath{n\ge2,} \\[3pt]
\FitMath{\phi(X_1)=1.}
\end{gather*}

If $H(\mathbf X)$ is disconnected the sum is zero.

\begin{lemma}[The connected graph identity]
As a formal power series at $z=0$,

\begin{gather*}
\FitMath{\log Z(z)=\sum_{n\ge1}\frac1{n!}
                  \sum_{X_1,\ldots,X_n\in\mathcal P}
                      \phi(\mathbf X)\prod_{i=1}^nz_{X_i}.}\qedhere
\end{gather*}

\end{lemma}
\begin{proof}
The ordered-list expression for $Z$ is

\begin{gather*}
\FitMath{Z(z)=\sum_{n\ge0}\frac1{n!}\sum_{\mathbf X\in\mathcal P^n}
  \prod_i z_{X_i}\prod_{i<j}
                  (1-\mathbf1_{\{X_i\not\sim X_j\}}).}
\end{gather*}

The product is one on pairwise compatible lists and zero otherwise. Such
lists contain no repetition, and their $n!$ orderings give exactly the
original polynomial. Expand the product as a sum over edge sets. Each graph
decomposes uniquely into its connected components. The weight factors over
the components, including the activity at each vertex. Expanding the
exponential of the proposed connected-graph series enumerates a partition
of the labeled vertices into these components. The factor for ordering the
components is canceled by the exponential's factorial, and the factorials
within the components cancel their vertex orderings. This reconstructs the
ordered-list expression coefficient by coefficient. Its constant term is
one, so its formal logarithm is the connected series.
\end{proof}

\subsubsection{The tree bound, including the cancellation}

\begin{lemma}[A hard-core tree inequality]
For a finite simple graph $H$ on $n\ge1$ labeled vertices,

\begin{gather*}
\FitMath{\left|\sum_{\substack{F\subset E(H)\\([n],F)\,\mathrm{connected}}}
                  (-1)^{|F|}\right|
       \le \#\{\text{spanning trees of }H\}.}
\end{gather*}

For $n=1$ the empty tree is counted once.
\end{lemma}
\begin{proof}
Fix a total order of the edges of $H$. Call an edge $e$ eligible for an edge
set $F$ if the endpoints of $e$ can be joined by a path in $F$ consisting
only of edges smaller than $e$. This definition does not require $e\in F$.
If at least one edge is eligible, choose the smallest eligible edge $e$ and
toggle its membership in $F$.

This is an involution. Eligibility of $e$ depends only on edges smaller than
$e$, which were not changed; the same holds for every smaller edge, so the
smallest eligible edge stays the same. Toggling $e$ preserves the connected
components, because its endpoints already have a path using smaller edges.
It reverses the parity of $|F|$. Thus all connected edge sets with an
eligible edge cancel in pairs in the alternating sum.

A surviving edge set has no cycle: if it contained one, the largest edge of
that cycle would be eligible. Consequently every surviving connected edge
set is a spanning tree. Each contributes the same sign $(-1)^{n-1}$, and
there can be no more survivors than spanning trees. This proves the bound.
The same argument also proves that the alternating sum has sign
$(-1)^{n-1}$ when it is nonzero.
\end{proof}

In particular,

\begin{gather*}
\FitMath{|\phi(X_1,\ldots,X_n)|\le
 \sum_{T\,\mathrm{tree\,on}\,[n]}
                       \prod_{\{i,j\}\in E(T)}
                             \mathbf1_{\{X_i\not\sim X_j\}}.}
\end{gather*}

The cancellation in this inequality is essential. Taking absolute values
before summing over connected graphs would count all connected subgraphs,
which need not admit the rooted-tree bound used next.

\subsubsection{The rooted-tree estimate}

Let $a_X\ge0$ be real numbers and suppose
\begin{equation}
 \sum_{Y:Y\not\sim X}|z_Y|e^{a_Y}\le a_X
                      \quad\text{for every }X\in\mathcal P.
 \label{eq:kp-explicit}
\end{equation}
This is the polymer criterion used here \cite{KP86}.
For a fixed root polymer $X$, let $\mathcal T_X$ be the total nonnegative
weight of all rooted labeled trees, with the root fixed to have polymer
$X$, every nonroot vertex of polymer $Y$ carrying weight $|z_Y|$, and every
edge requiring incompatibility of its endpoint polymers. A tree on $m$
vertices has the factor $1/(m-1)!$ for the nonroot vertex labels.

One can define this sum without assuming its convergence. Start with
$T_X^{(0)}=1$ and set

\begin{gather*}
\FitMath{T_X^{(h+1)}=\exp\left(\sum_{Y:Y\not\sim X}|z_Y|T_Y^{(h)}\right).}
\end{gather*}

The quantity $T_X^{(h)}$ sums trees of height at most $h$. To see this,
delete the root. Each child becomes the root of a tree of height at most
$h-1$, and an unordered collection of $r$ child trees contributes the
factor $1/r!$ in the exponential. Expanding those exponentials and assigning
the nonroot labels gives exactly $1/(m-1)!$: the multinomial coefficient
for assigning labels to the child subtrees cancels their factorials and
the ordering of the $r$ children. This also verifies the formula for every
fixed degree before any infinite sum is taken.

The iterates increase with $h$. Moreover $T_X^{(h)}\le e^{a_X}$ for every
$h$: this is true at $h=0$, and the induction step is precisely
\eqref{eq:kp-explicit}. Monotone convergence proves
\begin{equation}
 \mathcal T_X=\lim_{h\to\infty}T_X^{(h)}\le e^{a_X}.
 \label{eq:rooted-tree-bound}
\end{equation}

\begin{proposition}[Absolute convergence and a nonzero partition function]
Under \eqref{eq:kp-explicit} the connected series converges absolutely and
defines a number $L(z)$ with $e^{L(z)}=Z(z)$. In particular $Z(z)\ne0$.
For each root polymer $X$,

\begin{gather*}
\FitMath{\sum_{n\ge1}\frac1{(n-1)!}
 \sum_{X_2,\ldots,X_n}
 |\phi(X,X_2,\ldots,X_n)|\prod_{i=2}^n|z_{X_i}|
                       \le e^{a_X}.}
\end{gather*}

Also $|L(z)|\le\sum_X|z_X|e^{a_X}$.
If the criterion holds throughout a complex activity domain with a uniform
majorant, the series converges locally normally there and gives the branch
of $\log Z$ that equals zero at zero.
\end{proposition}
\begin{proof}
The tree inequality bounds the root series by $\mathcal T_X$, proving its
estimate by \eqref{eq:rooted-tree-bound}. Rooting each nonempty connected
list at its first vertex bounds the absolute unrooted series by
$\sum_X|z_X|e^{a_X}$: replacing $1/n!$ by $1/(n-1)!$ only enlarges it.

For $|t|<1$ the same majorant applies to the activities $tz_X$, and the
connected series is analytic in $t$ there. Near zero its exponential is
$Z(tz)$ by the formal identity and absolute convergence. The analytic
identity theorem gives $e^{L(tz)}=Z(tz)$ on the entire unit disc. The
absolute majorant is summable also at $|t|=1$, so the connected series is
continuous as $t$ tends radially to one. Taking that limit proves
$Z(z)=e^{L(z)}\ne0$.

On a domain with a uniform summable majorant, the Weierstrass criterion
gives normal convergence on smaller domains and hence holomorphy. The
branch statement follows by agreement near zero and continuation within
the connected component of that domain containing zero. This statement
does not assert that an arbitrary disconnected activity domain has a
single preferred logarithm.
\end{proof}

For a countable polymer set the same statements hold in a finite region
whenever $\sum_X|z_X|e^{a_X}<\infty$ there, by exhaustion and the displayed
absolute majorant. In an infinite lattice the total unrooted sum is
generally extensive and infinite. The rooted bound remains meaningful if
the incompatibility sums satisfy the criterion, and it is the form used
to define local connected expansions. No finite infinite-volume partition
function is asserted.

\subsubsection{Locality, weighted bounds, and analytic dependence}

If a polymer activity depends only on field variables in an enlargement
of that polymer, each connected term depends only on the union of those
enlargements. This is a direct consequence of its product formula. Extra
exponential weights can be incorporated in the activity majorant before
the criterion is checked. The root estimate then retains those weights.
Holomorphic dependence on external fields or sources passes to the
connected expansion on every domain where the same absolute majorant is
valid.

The theorem does not prove that gauge-field fluctuation activities satisfy
\eqref{eq:kp-explicit}. In a multiscale application their majorant must
already include the fluctuation covariance, the number of localization
domains, source radii, boundary costs, and any required derivatives.
Likewise, a tree bound for a small-field expansion does not by itself sum
large-field histories. Stating the activity estimate separately keeps the
proved combinatorial step from assuming its analytic conclusion.


% END MODULE chapters/polymer_theorem_v2.tex

% BEGIN MODULE chapters/free_curvature_v2.tex
\subsection{Correlations of the free curvature}
\label{ch:free-curvature-exact}

This section computes the noncoincident correlations of $\operatorname{tr}(F^2)$
for a free curvature in four Euclidean dimensions. The calculation fixes a
normalization and proves the cycle identities used in the short-distance
comparison. It does not identify an interacting lattice measure with this
model. That comparison, including its error, remains a separate estimate.

\subsubsection{A Gaussian field defined without a pointwise potential}

Fix a positive number $g$ and an integer $d_G\ge1$; for the adjoint color
space of $\operatorname{SU}(2)$, $d_G=3$. Coordinates in $\mathbb R^4$ are denoted
by $x_1,\ldots,x_4$. Our Fourier transform is

\begin{gather*}
\FitMath{\widehat f(k)=\int_{\mathbb R^4}e^{-ik\cdot x}f(x)\,dx,} \\[3pt]
\FitMath{f(x)=(2\pi)^{-4}\int e^{ik\cdot x}\widehat f(k)\,dk.}
\end{gather*}

Let $G$ be the tempered distribution with Fourier transform $|k|^{-2}$.
In four dimensions this multiplier is locally integrable at zero. Its
inverse transform, away from zero, is

\begin{gather*}
\FitMath{G(x)=\frac1{4\pi^2|x|^2}.}
\end{gather*}

One direct verification uses
$|k|^{-2}=\int_0^\infty e^{-t|k|^2}\,dt$ and the product Gaussian integral.
The inverse transform of $e^{-t|k|^2}$ is
$(4\pi t)^{-2}e^{-|x|^2/(4t)}$. Integrating in $t$, with
$u=|x|^2/(4t)$, gives $1/(4\pi^2|x|^2)$. The same identity holds as
tempered distributions by Fubini: the integral
$\int |\varphi(x)|\,|x|^{-2}dx$ is finite for every Schwartz function
$\varphi$, controlling the absolute integral of the heat-kernel formula.

For real Schwartz functions define

\begin{gather*}
\FitMath{(f,h)_G=(2\pi)^{-4}\int_{\mathbb R^4}
              \frac{\overline{\widehat f(k)}\widehat h(k)}{|k|^2}\,dk.}
\end{gather*}

This is a positive real bilinear form. The integral is finite near zero
because $\int_0^1r^3r^{-2}dr<\infty$, and is finite at infinity by Schwartz
decay. Quotient by its null space and complete to a real Hilbert space
$\mathcal H_G$. It is separable: the map $f\mapsto|k|^{-1}\widehat f(k)$
embeds it into a separable real form of $L^2(\mathbb R^4)$.

Choose an orthonormal basis $(e_j)$ and independent real standard normal
variables $(N_j)$ on their countable product probability space. For
$f\in\mathcal H_G$, the series

\begin{gather*}
\FitMath{W(f)=\sum_j(f,e_j)_G N_j}
\end{gather*}

converges in $L^2$ and defines a centered jointly Gaussian family with
$\mathbb E[W(f)W(h)]=(f,h)_G$. Finite partial sums prove the covariance
formula and joint Gaussianity; convergence of their characteristic
functions proves it for the limit. Independent copies give families
$A_\mu^a(f)=gW_\mu^a(f)$ for $1\le\mu\le4$, $1\le a\le d_G$.
Their covariance is

\begin{gather*}
\FitMath{\mathbb E[A_\mu^a(f)A_\nu^b(h)]
       =g^2\delta_{ab}\delta_{\mu\nu}(f,h)_G.}
\end{gather*}

No value $A_\mu^a(x)$ at a point is used in this definition.

Define the curvature as the Gaussian distribution

\begin{gather*}
\FitMath{F_{\mu\nu}^a(f)=-A_\nu^a(\partial_\mu f)
                         +A_\mu^a(\partial_\nu f),} \\[3pt]
\FitMath{\mu<\nu.}
\end{gather*}

There are six independent coordinate labels $P=(\mu,\nu)$, although the
six curvature components need not be independent random variables.
Differentiating the covariance as a distribution gives

\begin{gather*}
\FitMath{\mathbb E[F_{\mu\nu}^a(x)F_{\rho\sigma}^b(y)]
       =g^2\delta_{ab}K_{\mu\nu,\rho\sigma}(x-y),}
\end{gather*}

where the notation on the left denotes a distributional kernel and

\begin{gather*}
\FitMath{K_{\mu\nu,\rho\sigma}(x)=
 -\delta_{\nu\sigma}\partial_\mu\partial_\rho G(x)
 +\delta_{\nu\rho}\partial_\mu\partial_\sigma G(x)
 +\delta_{\mu\sigma}\partial_\nu\partial_\rho G(x)
 -\delta_{\mu\rho}\partial_\nu\partial_\sigma G(x).}
\end{gather*}

The signs follow from differentiating one argument at $x$ and the other
at $y$: $\partial_{y_\rho}G(x-y)=-\partial_{x_\rho}G(x-y)$.
At nonzero separation this is an ordinary smooth matrix-valued function.

\subsubsection{The six-dimensional curvature matrix}

Give $\Lambda^2\mathbb R^4$ the Euclidean inner product for which
$(e_\mu\wedge e_\nu)_{\mu<\nu}$ is an orthonormal basis. Put $r=|x|$,
$u=x/r$, and $R_u=I-2u\otimes u$. This $R_u$ is the orthogonal reflection
normal to $u$. Write $\Lambda^2R_u$ for its induced operator on two-forms.

\begin{proposition}[The curvature coefficient]
For $x\ne0$,

\begin{gather*}
\FitMath{K(x)=\frac1{\pi^2r^4}\Lambda^2R_u,} \\[3pt]
\FitMath{K(x)^\mathsf T K(x)=\frac1{\pi^4r^8}I_6,} \\[3pt]
\FitMath{\operatorname{tr}(K(x)^\mathsf T K(x))
                  =\frac6{\pi^4r^8}.}\qedhere
\end{gather*}

\end{proposition}
\begin{proof}
Twice differentiating $G$ gives

\begin{gather*}
\FitMath{\partial_\mu\partial_\rho G(x)
   =\frac1{2\pi^2r^4}(4u_\mu u_\rho-\delta_{\mu\rho}).}
\end{gather*}

At $u=e_4$ this Hessian is diagonal with entries
$(-1,-1,-1,3)/(2\pi^2r^4)$. Inserting it in the component formula for $K$
shows that $K$ is diagonal: it has eigenvalue $1/(\pi^2r^4)$ on
$e_1\wedge e_2,e_1\wedge e_3,e_2\wedge e_3$, and eigenvalue
$-1/(\pi^2r^4)$ on $e_1\wedge e_4,e_2\wedge e_4,e_3\wedge e_4$.
This is exactly the action of $\Lambda^2R_{e_4}$ times the stated scalar.
Both the differentiated radial kernel and the claimed expression transform
by conjugation under orthogonal changes of coordinates, proving the first
identity for every $u$. Since $\Lambda^2R_u$ is an orthogonal involution
on a six-dimensional space, its squared norm matrix is $I_6$, giving the
other two identities.
\end{proof}

Fix the standard orientation and let $J=*$ be the Hodge operator on
two-forms. It satisfies $J^2=I$. For an orthogonal map $R$,

\begin{gather*}
\FitMath{J(\Lambda^2R)=(\det R)(\Lambda^2R)J.}
\end{gather*}

To check the formula, use the defining identity
$\alpha\wedge *\beta=\langle\alpha,\beta\rangle\,dx_1\wedge\cdots\wedge dx_4$
and the fact that the top exterior power of $R$ multiplies the oriented
volume by $\det R$. Since $\det R_u=-1$, $JK(x)=-K(x)J$.
Every product of an odd number of matrices $K(x_i)$ therefore anticommutes
with $J$ and has trace zero: conjugation by $J$ preserves its trace while
changing its sign.

\subsubsection{Regularized densities and connected cycle contractions}

Take a real smooth even function $\rho$ supported in the unit ball with
$\int\rho=1$ and put $\rho_\varepsilon(x)=\varepsilon^{-4}\rho(x/\varepsilon)$.
Smear every curvature component with its translate to define
$F_{\varepsilon,P}^a(x)$. This is a Gaussian variable with a smooth
covariance kernel. Its centered quadratic density is

\begin{gather*}
\FitMath{q_\varepsilon(x)=\frac12\sum_{a=1}^{d_G}\sum_{P=1}^6
  \big(F_{\varepsilon,P}^a(x)^2
           -\mathbb E[F_{\varepsilon,P}^a(x)^2]\big).}
\end{gather*}

On any compact set its covariance is continuous and bounded. The integral
$q_\varepsilon(f)=\int f(x)q_\varepsilon(x)dx$ for compactly supported
smooth $f$ is defined by $L^2$ limits of Riemann sums. The continuous
covariance makes those sums Cauchy. The Gaussian moment formula gives
convergence of their moments of each fixed order as well: all the
integrands are finite sums of products of bounded continuous covariances
on a fixed compact set.

For clarity, first fix distinct points $x_1,\ldots,x_n$ at a fixed
$\varepsilon$. For one color let $C_{ij}$ be the $6\times6$ covariance block
of $F_\varepsilon(x_i)$ and $F_\varepsilon(x_j)$, including the factor $g^2$.
Put $Z=\operatorname{diag}(z_1I_6,\ldots,z_nI_6)$ and let $C$ be the full
block covariance matrix. The Gaussian integral yields, near $z=0$,

\begin{gather*}
\FitMath{\log\mathbb E\exp\big(\textstyle\sum_i z_iq_\varepsilon(x_i)\big)
   =-\frac{d_G}{2}\log\det(I-CZ)
                         -\frac{d_G}{2}\sum_i z_i\operatorname{tr}C_{ii}.}
\end{gather*}

Here $C$ is positive semidefinite; write the Gaussian vector as
$C^{1/2}N$ and diagonalize the real symmetric quadratic form for small
real $z$. The one-dimensional Gaussian integral proves the determinant
identity there. Holomorphic continuation proves it in a complex
neighborhood of zero. Sylvester's determinant identity identifies
$\det(I-C^{1/2}ZC^{1/2})$ with $\det(I-CZ)$, including singular $C$.

The last term only subtracts the first cumulants. Expand the logarithm of
the determinant by its convergent matrix power series. For $n\ge2$ the
joint cumulant is consequently

\begin{gather*}
\FitMath{\kappa_n(q_\varepsilon(x_1),\ldots,q_\varepsilon(x_n))
  =\frac{d_G}{2n}\sum_{\sigma\in S_n}
  \operatorname{tr}\big(C_{\sigma(1)\sigma(2)}\cdots
                      C_{\sigma(n)\sigma(1)}\big).}
\end{gather*}

Indeed only the term of degree $n$ in
$\frac{d_G}{2}\sum_{m\ge1}\operatorname{tr}((CZ)^m)/m$ contributes to the
mixed derivative. Each permutation records the order of the $n$ distinct
vertices of the cycle. The factor $1/n$ removes its cyclic repetitions;
the factor $1/2$ remains even when $n=2$.

\subsubsection{The noncoincident limit and its consequences}

Let the test functions $f_1,\ldots,f_n$ have pairwise disjoint compact
supports, and let $d>0$ be their minimum separation. The regularized
off-diagonal kernel is
$C_{ij}=g^2(K*\rho_\varepsilon*\rho_\varepsilon)(x_i-x_j)$.
For $\varepsilon<d/4$ the convolution samples $K$ only away from zero.
It converges uniformly, with every fixed derivative, on the compact set
of relevant differences. The cycle formula and ordinary dominated
convergence therefore give a finite limit of the smeared connected
correlations. Its kernel is

\begin{gather*}
\FitMath{C_n(x_1,\ldots,x_n)=\frac{d_Gg^{2n}}{2n}\sum_{\sigma\in S_n}
 \operatorname{tr}\big(K(x_{\sigma(1)}-x_{\sigma(2)})\cdots
 K(x_{\sigma(n)}-x_{\sigma(1)})\big),} \\[3pt]
\FitMath{x_i\ne x_j.}
\end{gather*}

This defines separated correlations. It does not assert convergence of
$q_\varepsilon(f)$ in $L^2$ when the same test is used twice; coincidence
singularities require a different distributional treatment.

\begin{corollary}[Two points, odd cycles, and a positive four-point kernel]
The limiting connected two-point kernel is

\begin{gather*}
\FitMath{C_2(x,y)=\frac{3d_Gg^4}{\pi^4|x-y|^8}.}
\end{gather*}

Every odd connected kernel $C_n$ with $n\ge3$ vanishes at noncoincident
arguments. If four distinct points are collinear, then

\begin{gather*}
\FitMath{C_4(x_1,x_2,x_3,x_4)=\frac{6d_Gg^8}{\pi^8}
 \sum_{\text{three unoriented cycles}}
       \prod_{(i,j)\text{ in the cycle}}|x_i-x_j|^{-4}>0.}
\end{gather*}

For four sufficiently small disjoint neighborhoods of such points the
four-point kernel remains strictly positive throughout their product.
\end{corollary}
\begin{proof}
For two points, the permutation formula is
$d_Gg^4\operatorname{tr}(K(x-y)K(y-x))/2$. The kernel is even and symmetric,
so the curvature coefficient proposition proves the stated value. For odd
$n$, every trace in the cycle sum vanishes by the Hodge anticommutation
identity. For collinear points all difference directions give the same
reflection $R_u=R_{-u}$. The product of four curvature matrices is the
positive scalar product of their radial factors times $(\Lambda^2R_u)^4=I$.
Every trace is therefore six times that scalar. The 24 permutations divide
into three unoriented cycles, each represented eight times, canceling
the factor $2n=8$. Positivity on small neighborhoods follows from
continuity and the strictly positive value at the chosen configuration.
\end{proof}

For nonnegative nonzero separated tests this gives

\begin{gather*}
\FitMath{C_2(f,h)=\frac{3d_Gg^4}{\pi^4}
                 \iint\frac{f(x)h(y)}{|x-y|^8}\,dx\,dy>0.}
\end{gather*}

Thus the coefficient $3(N^2-1)/\pi^4$ used for $\operatorname{SU}(N)$ is the
exact free coefficient in this normalization. For $\operatorname{SU}(2)$ it
is $9/\pi^4$. Recovering this coefficient from the interacting lattice
theory still requires its ultraviolet comparison and uniform remainder
estimate. The four-point calculation also shows why non-Gaussianity of
the scalar density alone cannot certify an interacting underlying gauge
theory: this density already has a positive connected four-point kernel
in an explicitly Gaussian model.


% END MODULE chapters/free_curvature_v2.tex

% BEGIN MODULE chapters/lattice_curvature_comparison_v2.tex
\subsection{The free lattice curvature comparison}
\label{ch:lattice-curvature-comparison}

This section proves a free-field input of the construction. The covariance of a Gaussian exact lattice two-form is an orthogonal projection. After multiplication by the physical factor $a^{-4}$, its entries converge away from coincidence to the continuum curvature kernel. We prove an error estimate uniform in the lattice spacing and the torus side, and compute the resulting separated quadratic-density coefficient. No interacting measure is compared with this Gaussian model here.

\phantomsection\label{notdef:curvature-kernels}
\subsubsection{Cochains, Fourier conventions, and the statement}

Let $a>0$, $\ell>0$, and $N=\ell/a\in\mathbb N$. The lattice is
$\Lambda_{a,\ell}=a\mathbb Z^4/\ell\mathbb Z^4$. A one-form has coordinates $A_\mu(x)$, one per positive bond based at $x$. A two-form has coordinates $F_{\mu\nu}(x)$ for $\mu<\nu$, based at the lower vertex of its positively oriented plaquette. Both cochain spaces use the plain coordinate inner products. Define

\begin{gather*}
\FitMath{(\mathrm{d}_1A)_{\mu\nu}(x)
 =a^{-1}\{A_\nu(x+ae_\mu)-A_\nu(x)
          -A_\mu(x+ae_\nu)+A_\mu(x)\}.}
\end{gather*}

The factor $a^{-1}$ does not affect the range. Let $P_{\mathrm{ex}}$ be the orthogonal projection onto $\operatorname{Ran}\mathrm{d}_1$. It is the covariance of a centered Gaussian vector which is standard on that subspace. Potential gauge directions need not be inverted to define this curvature covariance. The constant two-forms are excluded because they are not in $\operatorname{Ran}\mathrm{d}_1$.

Write $P=(\mu,\nu)$, $\mu<\nu$, and put $c_P=(e_\mu+e_\nu)/2$. The centre of the plaquette $(x,P)$ is $x+ac_P$. The physical lattice kernel is

\begin{gather*}
\FitMath{\mathsf K_{a,\ell}^{PQ}(x,y)
    =a^{-4}P_{\mathrm{ex}}((x,P),(y,Q)).}
\end{gather*}

For a displacement on the torus, its length means distance to $\ell\mathbb Z^4$. Representatives are chosen in $[-\ell/2,\ell/2]^4$ when needed.

The continuum Fourier transform on $\mathbb R^4$ is

\begin{gather*}
\FitMath{\widehat f(k)=\int e^{-ik\cdot x}f(x)\,dx,} \\[3pt]
\FitMath{f(x)=(2\pi)^{-4}\int e^{ik\cdot x}\widehat f(k)\,dk.}
\end{gather*}

Let $G(x)=1/(4\pi^2|x|^2)$ and define the continuum curvature distribution by
\begin{align*}
 \mathsf K_{\mu\nu,\rho\sigma}(x)
 ={}&-\delta_{\nu\sigma}\partial_\mu\partial_\rho G(x)
      +\delta_{\nu\rho}\partial_\mu\partial_\sigma G(x)\\
    &+\delta_{\mu\sigma}\partial_\nu\partial_\rho G(x)
      -\delta_{\mu\rho}\partial_\nu\partial_\sigma G(x).
\end{align*}
At $x\ne0$ this is an ordinary smooth matrix. Its contact terms at zero are retained in its Fourier definition and are not squared in the statements below.

\begin{proposition}[Free lattice curvature comparison]
There is a numerical constant $C$ such that, for every $a,\ell$ as above and every pair of plaquettes with centre displacement

\begin{gather*}
\FitMath{u=x-y+a(c_P-c_Q),} \\[3pt]
\FitMath{10a\le r=\operatorname{dist}(u,\ell\mathbb Z^4)\le\ell/10,}
\end{gather*}

one has
\begin{equation}
 \left|\mathsf K_{a,\ell}^{PQ}(x,y)-\mathsf K_{PQ}(u)\right|
 \le Cr^{-4}\left\{(a/r)^{1/2}+(r/\ell)^4\right\}.
 \label{eq:lcc-main}
\end{equation}
The continuum kernel on the right is evaluated at the short representative of $u$. The constants do not depend on the orientation labels, the lattice spacing, the number of sites, or the torus side.
\end{proposition}

The exponent $1/2$ is sufficient for convergence and is not claimed optimal. Keeping the barycentre phases in a sharper symbol expansion can improve the lattice power. The proof here deliberately keeps a periodic raw symbol and moves to barycentres only after the Fourier estimate.

\subsubsection{The exact projection symbol}

The torus frequencies are $k\in\delta\mathbb Z^4$, where $\delta=2\pi/\ell$. Choose one Brillouin representative of each frequency modulo $(2\pi/a)\mathbb Z^4$ and denote this set by $\mathcal Q_a$. Its coordinates lie between $-\pi/a$ and $\pi/a$, with one endpoint omitted when necessary. Define

\begin{gather*}
\FitMath{t_{a,\mu}(k)=a^{-1}(e^{iak_\mu}-1),} \\[3pt]
\FitMath{\omega_a(k)=\sum_\mu|t_{a,\mu}(k)|^2
        =4a^{-2}\sum_\mu\sin^2(ak_\mu/2).}
\end{gather*}

Let $B_a(k):\mathbb C^4\to\Lambda^2\mathbb C^4$ be

\begin{gather*}
\FitMath{(B_a(k)v)_{\mu\nu}=t_{a,\mu}(k)v_\nu-t_{a,\nu}(k)v_\mu.}
\end{gather*}

This is exactly the Fourier symbol of $\mathrm{d}_1$. Direct multiplication gives

\begin{gather*}
\FitMath{B_a^*B_a=\omega_a I-t_a t_a^*,} \\[3pt]
\FitMath{B_at_a=0.}
\end{gather*}

Consequently $(B_aB_a^*)^2=\omega_a B_aB_a^*$. For $k\ne0$ in the Brillouin zone,

\begin{gather*}
\FitMath{P_a(k)=B_a(k)B_a(k)^*/\omega_a(k)}
\end{gather*}

is the orthogonal projection onto $\operatorname{Ran}B_a(k)$. Set $P_a(0)=0$. Thus $\|P_a(k)\|\le1$ and its rank at nonzero frequency is three. At frequency zero the exact-form subspace is zero.

The unitary finite Fourier transform and its inverse are

\begin{gather*}
\FitMath{\widehat F(k)=N^{-2}\sum_{x\in\Lambda_{a,\ell}}e^{-ik\cdot x}F(x),} \\[3pt]
\FitMath{F(x)=N^{-2}\sum_{k\in\mathcal Q_a}e^{ik\cdot x}\widehat F(k).}
\end{gather*}

Therefore
\begin{equation}
 a^{-4}P_{\mathrm{ex}}((x,P),(y,Q))
  =\ell^{-4}\sum_{k\in\mathcal Q_a}
          e^{ik\cdot(x-y)}[P_a(k)]_{PQ}.
 \label{eq:lcc-fourier-lattice}
\end{equation}
Indeed the projection entry before multiplication by $a^{-4}$ has prefactor $N^{-4}$, and $a^{-4}N^{-4}=\ell^{-4}$.

For the continuum symbol replace $t_a$ by $t_0(k)=ik$. Put

\begin{gather*}
\FitMath{P_0(k)=B_0(k)B_0(k)^*/|k|^2} \\[3pt]
\FitMath{(k\ne0),} \\[3pt]
\FitMath{P_0(0)=0.}
\end{gather*}

It is again a norm-one rank-three projection. With the specified Fourier convention its inverse transform is precisely the distribution $\mathsf K$ defined above. On the torus define the distribution
\begin{equation}
 \mathsf K_\ell(z)=\ell^{-4}
        \sum_{k\in\delta\mathbb Z^4}e^{ik\cdot z}P_0(k).
 \label{eq:lcc-fourier-torus}
\end{equation}
The bounded Fourier coefficients define a periodic distribution. Below we show that it is smooth away from the origin and identify its finite-volume correction.

\subsubsection{Uniform derivative and counting estimates}

For each multi-index $\alpha$ there is a constant $C_\alpha$ such that, in the Brillouin zone away from zero,
\begin{equation}
 \|\partial^\alpha P_a(k)\|\le C_\alpha |k|^{-|\alpha|}.
 \label{eq:lcc-symbol-derivatives}
\end{equation}
The same estimate holds for $P_0$. To prove it, use
$2|\sin(t/2)|\ge2|t|/\pi$ for $|t|\le\pi$. Hence
$c|k|^2\le\omega_a(k)\le|k|^2$ throughout the Brillouin zone. Differentiating the exponential gives
$|\partial^\alpha t_a(k)|\le C_\alpha |k|^{1-|\alpha|}$; for positive derivative order this uses $a|k|\le2\pi$. The product rule gives the analogous bounds of order $2-|\alpha|$ for $\omega_a$ and $B_aB_a^*$. Induction in the identity $\omega_a\omega_a^{-1}=1$ bounds derivatives of the reciprocal. The quotient rule proves \eqref{eq:lcc-symbol-derivatives}. All functions involved are periodic; near a Brillouin face the same bounds hold in periodic coordinates, where $|k|$ is comparable to $a^{-1}$.

For $a|k|\le c_0$, with a fixed small $c_0>0$, Taylor's formula gives
$t_a(k)=ik+O(a|k|^2)$. The preceding denominator bounds imply
\begin{equation}
 \|P_a(k)-P_0(k)\|\le Ca|k|.
 \label{eq:lcc-low-symbol}
\end{equation}
For example, the numerator difference is $O(a|k|^3)$ and the denominator difference is $O(a|k|^3)$, while both denominators are comparable to $|k|^2$. Both symbols have been assigned zero at $k=0$, so that frequency contributes no error.

We will use two elementary grid bounds. If $R\ge C\delta$ and $M>4$, then
\begin{align}
 \ell^{-4}\#\{k\in\delta\mathbb Z^4:|k|\le2R\}
       &\le C R^4,\label{eq:lcc-count-low}\\
 \ell^{-4}\sum_{k\in\delta\mathbb Z^4}(R+|k|)^{-M}
       &\le C_M R^{4-M}.\label{eq:lcc-count-high}
\end{align}
For a proof, place a cube of side $\delta$ about each grid point. The cubes are disjoint up to boundaries. If $R\ge C\delta$, the values of $R+|k|$ and $R+|q|$ on such a cube differ by bounded factors. The second sum is bounded by a constant times $(2\pi)^{-4}\int(R+|q|)^{-M}dq$, which is $C_MR^{4-M}$ after scaling. The first bound follows by enclosing the cubes in a ball of radius $3R$.

\subsubsection{Summation by parts on the frequency grid}

Fix a smooth cutoff $\chi$ equal to one on $|k|\le1$ and zero on $|k|\ge2$. Suppose

\begin{gather*}
\FitMath{R\ge2M\delta,} \\[3pt]
\FitMath{2R<\pi/(2a).}
\end{gather*}

Extend $\chi(k/R)$ periodically on the frequency torus of side $2\pi/a$. It is smooth because its support stays away from the Brillouin faces. The high-frequency lattice symbol

\begin{gather*}
\FitMath{H_a(k)=(1-\chi(k/R))P_a(k)}
\end{gather*}

is smooth on that frequency torus, including at zero. Write

\begin{gather*}
\FitMath{\rho_a(k)=\operatorname{dist}(k,(2\pi/a)\mathbb Z^4).}
\end{gather*}

The derivative estimates and the cutoff product rule give

\begin{gather*}
\FitMath{\|\partial_\mu^M H_a(k)\|
       \le C_M(R+\rho_a(k))^{-M}.}
\end{gather*}

Let $\Delta_\mu H(k)=H(k+\delta e_\mu)-H(k)$. The iterated fundamental theorem of calculus yields

\begin{gather*}
\FitMath{\Delta_\mu^M H(k)=\delta^M
 \int_{[0,1]^M}\partial_\mu^M H
       (k+\delta(t_1+\cdots+t_M)e_\mu)\,dt.}
\end{gather*}

The torus distance $\rho_a$ changes by at most $M\delta$ along this segment. Since $R\ge2M\delta$, $R+\rho_a$ remains comparable to its value at $k$. Thus
\begin{equation}
 \|\Delta_\mu^M H_a(k)\|
       \le C_M\delta^M(R+\rho_a(k))^{-M}.
 \label{eq:lcc-differences}
\end{equation}

Now let $z\in a\mathbb Z^4/\ell\mathbb Z^4$, $z\ne0$, and choose a short representative. Its length $r_z$ is at most $\ell$, and some coordinate satisfies $|z_\mu|\ge r_z/2$. On $|z_\mu|\le\ell/2$,

\begin{gather*}
\FitMath{|e^{-i\delta z_\mu}-1|
 =2|\sin(\pi z_\mu/\ell)|\ge4|z_\mu|/\ell
 \ge2r_z/\ell.}
\end{gather*}

The phase $e^{ik\cdot z}$ is periodic on the frequency grid because $z/a$ is integral. Finite summation by parts therefore gives the exact identity

\begin{gather*}
\FitMath{(e^{-i\delta z_\mu}-1)^M
   \sum_{k\in\mathcal Q_a}e^{ik\cdot z}H_a(k)
 =\sum_{k\in\mathcal Q_a}e^{ik\cdot z}\Delta_\mu^M H_a(k).}
\end{gather*}

There is no Brillouin boundary term. Apply \eqref{eq:lcc-differences} and \eqref{eq:lcc-count-high}, using $\rho_a(k)=|k|$ on the chosen representatives. The result is
\begin{equation}
 \left\|\ell^{-4}\sum_{k\in\mathcal Q_a}
                 e^{ik\cdot z}H_a(k)\right\|
 \le C_M r_z^{-M}R^{4-M}.
 \label{eq:lcc-high-lattice}
\end{equation}
The powers of the torus side cancel: the inverse difference factor contributes $(\ell/r_z)^M$ and the frequency difference contributes $\delta^M=(2\pi/\ell)^M$.

The continuum torus tail obeys the same estimate. Here
$H_0(k)=(1-\chi(k/R))P_0(k)$ on the infinite reciprocal grid. Its derivatives satisfy the same bounds with $|k|$ in place of $\rho_a$. The $M$th differences are absolutely summable for $M>4$. Multiplication of a Fourier distribution by $(e^{-i\delta z_\mu}-1)^M$ shifts its Fourier coefficients and hence gives the same identity distributionally. The right side is an absolutely convergent continuous Fourier series. On the open set where the multiplying factor is nonzero, division by that smooth factor identifies the tail with an ordinary continuous function and proves
\begin{equation}
 \|\mathsf K_\ell^{\mathrm{high}}(z)\|
        \le C_M r_z^{-M}R^{4-M}.
 \label{eq:lcc-high-continuum}
\end{equation}
This distributional argument avoids assuming convergence of an undifferentiated order-zero Fourier series.

\subsubsection{Comparison with the continuum torus kernel}

On the common low-frequency support, the finite and infinite reciprocal grids agree. \eqref{eq:lcc-low-symbol} and \eqref{eq:lcc-count-low} give

\begin{gather*}
\FitMath{\left\|\ell^{-4}\sum_k e^{ik\cdot z}\chi(k/R)
                  (P_a(k)-P_0(k))\right\|\le CaR^5.}
\end{gather*}

Combining this with the two high-frequency bounds yields, for $M>4$,
\begin{equation}
 \|\mathsf K_{a,\ell}^{\mathrm{raw}}(z)-\mathsf K_\ell(z)\|
       \le C_M\{aR^5+r_z^{-M}R^{4-M}\}.
 \label{eq:lcc-split}
\end{equation}
Here $\mathsf K_{a,\ell}^{\mathrm{raw}}$ is the matrix Fourier sum in \eqref{eq:lcc-fourier-lattice} at the common plaquette-origin displacement $z$.

Take $M=9$ and choose

\begin{gather*}
\FitMath{R=\frac D{r_z}(r_z/a)^{1/10},}
\end{gather*}

where $D$ is a fixed sufficiently large numerical constant. Because $r_z\le\ell$, this ensures $R\ge18\delta$. For $r_z/a$ larger than a fixed numerical threshold $C_0$, it also ensures $2aR<c_0$ and $2R<\pi/(2a)$. Both terms in \eqref{eq:lcc-split} are then at most a constant times

\begin{gather*}
\FitMath{r_z^{-4}(a/r_z)^{1/2}.}
\end{gather*}

Indeed $aR^5=D^5r_z^{-4}(a/r_z)^{1/2}$ and
$r_z^{-9}R^{-5}=D^{-5}r_z^{-4}(a/r_z)^{1/2}$.

For $1\le r_z/a\le C_0$, the lattice Fourier sum has norm at most $a^{-4}$ because $\|P_a\|\le1$. The continuum bound
$\|\mathsf K_\ell(z)\|\le C r_z^{-4}$, established in the next section, shows that their difference is at most $C'a^{-4}$. Increasing the numerical constant by a factor depending only on $C_0$ bounds this by
$Cr_z^{-4}(a/r_z)^{1/2}$. Consequently, for every nonzero raw lattice displacement,
\begin{equation}
 \|\mathsf K_{a,\ell}^{\mathrm{raw}}(z)-\mathsf K_\ell(z)\|
       \le Cr_z^{-4}(a/r_z)^{1/2}.
 \label{eq:lcc-torus-comparison}
\end{equation}
All thresholds used here are fixed numbers. They impose no relation between the physical volume and the ultraviolet cutoff beyond $\ell/a\in\mathbb N$.

\subsubsection{The finite-volume correction and contact terms}

The continuum torus heat kernel is

\begin{gather*}
\FitMath{p_t^\ell(x)=\sum_{m\in\mathbb Z^4}
       (4\pi t)^{-2}e^{-|x+m\ell|^2/(4t)}
   =\ell^{-4}\sum_{k\in\delta\mathbb Z^4}e^{-t|k|^2}e^{ik\cdot x}.}
\end{gather*}

For completeness, the first expression is an absolutely convergent periodic sum for $t>0$. Its Fourier coefficient at $k$ is the ordinary Gaussian Fourier integral $\ell^{-4}e^{-t|k|^2}$, which proves the second expression. Define the mean-zero torus Green distribution by

\begin{gather*}
\FitMath{G_\ell(x)=\int_0^\infty(p_t^\ell(x)-\ell^{-4})\,dt.}
\end{gather*}

Its nonzero Fourier coefficient is $\ell^{-4}|k|^{-2}$ and its zero coefficient is zero. Applying the same second-derivative combination as in the definition of $\mathsf K$ gives \eqref{eq:lcc-fourier-torus}.

We record the uniform local correction. At $\ell=1$ and $|x|\le1/4$,
\begin{align*}
 G_1(x)-G(x)
 ={}&\int_0^1\sum_{m\ne0}(4\pi t)^{-2}
                 e^{-|x+m|^2/(4t)}\,dt-1\\
 &+\int_1^\infty(p_t^1(x)-1)\,dt
       -\int_1^\infty(4\pi t)^{-2}e^{-|x|^2/(4t)}\,dt.
\end{align*}
Every fixed derivative of the first integral is uniformly bounded: differentiating a Gaussian gives a polynomial times a Gaussian, while $|x+m|\ge c|m|$ and
$\sum_{m\ne0}e^{-c|m|^2/t}\le C e^{-c'/t}$ for $0<t\le1$. The resulting power of $t^{-1}$ is integrable against this last exponential. For the second integral, its differentiated Fourier series is bounded by a constant times $e^{-ct}$ for $t\ge1$. Derivatives of the last integral are bounded by $C_\alpha t^{-2-|\alpha|/2}$ on this ball and are integrable. Thus $G_1-G$ extends smoothly across zero with bounded derivatives of every fixed order.

Scaling gives $G_\ell(x)=\ell^{-2}G_1(x/\ell)$ and the corresponding homogeneity for $G$. In particular,
\begin{equation}
 \|\mathsf K_\ell(z)-\mathsf K(z)\|\le C\ell^{-4}
       \quad(0<|z|\le\ell/4).
 \label{eq:lcc-volume-error}
\end{equation}
The same argument gives an extra factor $\ell^{-1}$ for each additional derivative of the difference. On the remainder of the unit torus, away from zero, the heat-kernel formula and its derivatives converge uniformly on compact sets. Scaling and compactness therefore also give the global bounds

\begin{gather*}
\FitMath{\|\partial^\alpha\mathsf K_\ell(z)\|
       \le C_\alpha r_z^{-4-|\alpha|},} \\[3pt]
\FitMath{z\ne0.}
\end{gather*}

These justify the estimate used in the bounded $r_z/a$ regime above.

All these statements concern the kernels away from coincidence. The symbols have not discarded their order-zero distributional terms. For example, the trace of $P_0(k)$ is three for $k\ne0$. Accordingly the trace of the torus kernel as a distribution is
$3(\delta_{\mathrm{per}}-\ell^{-4})$; away from zero it is $-3\ell^{-4}$. The trace of the free noncoincident kernel is zero. This is consistent with \eqref{eq:lcc-volume-error}; it would be inconsistent to replace the full distribution by its noncoincident formula before taking Fourier coefficients.

\subsubsection{Moving to barycentres and proving the theorem}

Let $u$ be the centre displacement in \eqref{eq:lcc-main} and let
$z=u-a(c_P-c_Q)$ be its raw-origin displacement. The shift has length at most $a$. Choose representatives so $|u|=r\le\ell/10$ and $|z-u|\le a$. Then

\begin{gather*}
\FitMath{0.9r\le|z|\le1.1r<\ell/4.}
\end{gather*}

Use \eqref{eq:lcc-torus-comparison} for the raw kernel and \eqref{eq:lcc-volume-error} for the continuum torus kernel at $z$. The free kernel is homogeneous of degree $-4$ and smooth away from zero, so $\|\nabla\mathsf K(v)\|\le C|v|^{-5}$. The segment from $z$ to $u$ stays at distance at least $0.9r$ from zero. Hence

\begin{gather*}
\FitMath{|\mathsf K_{PQ}(z)-\mathsf K_{PQ}(u)|\le Ca r^{-5}.}
\end{gather*}

Combining the three differences proves

\begin{gather*}
\FitMath{|\mathsf K_{a,\ell}^{PQ}(x,y)-\mathsf K_{PQ}(u)|
 \le C\{r^{-4}(a/r)^{1/2}+\ell^{-4}+ar^{-5}\}.}
\end{gather*}

Since $a/r\le1$, the last term is absorbed by the first, and
$\ell^{-4}=r^{-4}(r/\ell)^4$. This proves \eqref{eq:lcc-main}.

The argument also proves convergence to $\mathsf K_\ell$ uniformly on any fixed separated subset of the torus, when plaquette centres are used in both arguments. In this statement use the global derivative bound for $\mathsf K_\ell$ to move from origins to centres. No assumption that a chosen separation is small relative to $\ell$ is needed for this fixed-torus convergence.

\subsubsection{The separated quadratic-density coefficient}

Let $\xi_p^b$, $1\le b\le d_G$, be independent color copies of the Gaussian exact two-form, each with covariance $g^2P_{\mathrm{ex}}$. For a test $f$ on the torus define the finite-dimensional centered density

\begin{gather*}
\FitMath{Q_a(f)=\frac12\sum_{b=1}^{d_G}\sum_p f(x_p)
             \{(\xi_p^b)^2-\mathbb E(\xi_p^b)^2\}.}
\end{gather*}

Here $x_p$ is the plaquette barycentre. Equivalently, $a^{-2}\xi_p^b$ is the physical curvature and the displayed sum includes its $a^4$ Riemann weight. For two tests of positive mutual separation, the Gaussian identity
$\operatorname{Cov}(X^2,Y^2)=2\operatorname{Cov}(X,Y)^2$ gives

\begin{gather*}
\FitMath{\operatorname{Cov}(Q_a(f),Q_a(h))
 =\frac{d_Gg^4}{2}\,a^8\sum_{p,q}
         f(x_p)h(x_q)\,|\mathsf K_{a,\ell}(p,q)|^2.}
\end{gather*}

For real kernels the square here is the ordinary real square of the scalar entry. The uniform separated convergence just proved, followed by ordinary shifted-grid Riemann sums, gives

\begin{gather*}
\FitMath{\lim_{a\downarrow0}\operatorname{Cov}(Q_a(f),Q_a(h))
 =\frac{d_Gg^4}{2}\iint f(x)h(y)
             \sum_{P,Q}\mathsf K_{\ell,PQ}(x-y)^2\,dx\,dy.}
\end{gather*}

For continuous compactly supported separated tests the Riemann-sum integrand is continuous after setting it to zero in a neighbourhood of the diagonal; for $C^1$ tests it is $C^1$ there. There is therefore no use of an integral of $|x-y|^{-8}$ through coincidence.

Now keep the tests in a fixed compact subset of $\mathbb R^4$ and let $\ell\to\infty$. The correction estimate makes $\mathsf K_\ell\to\mathsf K$ uniformly on the relevant separated differences. Direct differentiation of $G$ shows

\begin{gather*}
\FitMath{\mathsf K(x)=\frac1{\pi^2|x|^4}\Lambda^2(I-2\widehat x\otimes\widehat x).}
\end{gather*}

One can check the coefficient at $\widehat x=e_4$: the six eigenvalues are three copies of $1/(\pi^2|x|^4)$ and three copies of its negative. Orthogonal covariance proves the identity in every direction. The induced two-form reflection is orthogonal, and therefore

\begin{gather*}
\FitMath{\sum_{P,Q}\mathsf K_{PQ}(x)^2=\frac6{\pi^4|x|^8}.}
\end{gather*}

It follows that the infinite-volume separated covariance is exactly
\begin{equation}
 \frac{3d_Gg^4}{\pi^4}
       \iint f(x)h(y)|x-y|^{-8}\,dx\,dy.
 \label{eq:lcc-density-coefficient}
\end{equation}
For $\operatorname{SU}(2)$, $d_G=3$, so the coefficient is $9g^4/\pi^4$ in this convention. Higher separated Gaussian cumulants can likewise be passed through their finite cycle formulas using the same kernel convergence.

There is also a relative finite-volume statement for nonnegative pairs. Suppose every cross-distance between the supports lies between $d$ and $C_{\mathrm{sh}}d$, and $C_{\mathrm{sh}}d\le\ell/10$. The correction estimate and the Frobenius norm of the free kernel give pointwise

\begin{gather*}
\FitMath{\left|\sum_{P,Q}\mathsf K_{\ell,PQ}(v)^2
              -\frac6{\pi^4|v|^8}\right|
 \le C\left(\frac{|v|}{\ell}\right)^4|v|^{-8}.}
\end{gather*}

Indeed the difference of the squares is bounded by
$C|v|^{-4}\ell^{-4}+C\ell^{-8}$. Multiplication by $f(x)h(y)\ge0$ and integration show that the limiting torus covariance differs from \eqref{eq:lcc-density-coefficient} by a relative error at most $C_{\mathrm{sh}}'(d/\ell)^4$, provided both tests are nonzero. The mesh error tends to zero for each such fixed test pair by the uniform kernel estimate and the Riemann-sum argument above. A rate relative to a varying family of tests would require their norms and nondegeneracy to be specified.

This proves the free covariance comparison and the free density coefficient. It asserts neither an $L^2$ limit of $Q_a(f)$ with the same test repeated nor an interacting short-distance asymptotic. Establishing the latter requires a comparison of the interacting measure with this Gaussian model and a remainder smaller than the leading term.


% END MODULE chapters/lattice_curvature_comparison_v2.tex

% BEGIN MODULE chapters/gibbs_symmetry_v2.tex
\subsection{Link reflections of periodic Gibbs states}
\label{ch:periodic-link-symmetry}

This section concerns one fixed lattice and one fixed positive coupling.
No continuum limit or renormalization estimate is used. The conclusion
quantifies over every periodic DLR state. In particular, the state is not
assumed to arise as a limit of torus measures. The finite-dimensional
transfer calculation below, together with a local elimination identity,
is what permits that quantifier.

Set the lattice spacing equal to one, let $G=\operatorname{SU}(2)$, and use
normalized Haar measure. On one orientation of the bonds of
$\mathbb Z^4$ assign variables $U_b\in G$; reverse orientation means
$U_{\bar b}=U_b^{-1}$. For an elementary plaquette set

\begin{gather*}
\FitMath{a_p=1-\tfrac12\operatorname{tr}U_p,} \\[3pt]
\FitMath{0\le a_p\le2.}
\end{gather*}

The trace is real for $G=\operatorname{SU}(2)$. Fix $\beta>0$. A DLR state
is a probability measure whose conditional law on any finite bond set
$D$, given all other bonds, has density proportional to
$\exp[-\beta\sum_{p:p\cap D\ne\varnothing}a_p]$. A state is periodic
when it is invariant under a subgroup of finite index of
$\mathbb Z^4$. A link reflection is reflection in a plane
$\{x_\mu=i+\frac12\}$, with $i\in\mathbb Z$.

\begin{proposition}[Link symmetry of every periodic Wilson state]
Every periodic DLR state of the four-dimensional $\operatorname{SU}(2)$
Wilson interaction at $\beta>0$ is invariant under every link
reflection. Consequently every such state is invariant under
$2\mathbb Z^4$.
\end{proposition}

The proof uses positivity and strict positivity of different objects.
The one-bond convolution operator is positive as an operator, with
strictly positive Peter--Weyl eigenvalues. Its original kernel is
strictly positive pointwise. An auxiliary operator inserted between
two half-kernels will be a contraction, and need not be self-adjoint
or positive. Confusing these conditions would invalidate the
variational step.

\paragraph*{A lower tangent valid for an arbitrary periodic state}

Let $\phi$ be a real continuous local observable and let
$m\mathbb Z^4$ be a subgroup of periods of a DLR state $\nu$.
For torus sides $s$ divisible by $m$, write

\begin{gather*}
\FitMath{W_s=\sum_{v\in m\mathbb Z^4/s\mathbb Z^4}\tau_v\phi,} \\[3pt]
\FitMath{Z_s(\phi)=\int e^{W_s-\beta\sum_pa_p}\,dU.}
\end{gather*}

Here local observables are evaluated on the periodic extension of
the torus configuration. Suppose

\refstepcounter{equation}
\begin{gather*}
\FitMath{Z_s(\phi)\le Z_s(0)} \\[3pt]
\FitMath{\hbox{for all sufficiently large multiples }s\hbox{ of }m.}\tag{\theequation}\label{eq:gibbs-symmetry-GL1}
\end{gather*}

Then

\refstepcounter{equation}
\begin{gather*}
\FitMath{\nu(\phi)\le0.}\tag{\theequation}\label{eq:gibbs-symmetry-GL2}
\end{gather*}


To prove this implication, take a cube $D_s$ of sites of side $s$,
and let $B_s$ be the positively oriented bonds whose initial sites
belong to the cube. Choose the representatives of the source
anchors inside $D_s$, and denote their sum, evaluated with an
arbitrary outside configuration, by $\widetilde W_s$. This is a
finite local function; terms near the boundary can involve outside
bonds. The number of anchors is exactly $M_s=(s/m)^4$.
The DLR equation and Jensen's inequality give

\begin{gather*}
\FitMath{\log\gamma_{B_s}(e^{\widetilde W_s})(U')
       \ge\gamma_{B_s}(\widetilde W_s)(U').}
\end{gather*}

The conditional partition-function ratio on the left can be
compared with $Z_s(\phi)/Z_s(0)$ uniformly in $U'$. Interior terms
are identical under the identification of $B_s$ with the torus
bond coordinates. Only plaquettes and translates of $\phi$ in a
fixed-width boundary layer can differ. Their number is at most
$C_{\phi,m}s^3$, and each difference is bounded by
$C(\beta+\|\phi\|_\infty)$. Thus

\begin{gather*}
\FitMath{\log\frac{Z_s(\phi)}{Z_s(0)}
   \ge \nu(\widetilde W_s)-C_{\phi,m,\beta}s^3
   = M_s\nu(\phi)-C_{\phi,m,\beta}s^3.}
\end{gather*}

In the first inequality we integrated the conditional inequality
against $\nu$ and used the DLR equation. The equality uses only
$m\mathbb Z^4$ invariance. Divide by $M_s$ and let $s\to\infty$
through multiples of $m$. The assumption in \eqref{eq:gibbs-symmetry-GL1} proves \eqref{eq:gibbs-symmetry-GL2}.

This argument needs neither differentiability of an infinite-volume
pressure nor representation of $\nu$ by finite torus measures.
The local source is fixed during the limit $s\to\infty$. If the
source later depends on a parameter tending to zero, that parameter
is kept fixed until the surface error has disappeared. Constant
parts of $\phi$ cause no difficulty with the anchored definition of
$\widetilde W_s$: they occur with exactly the same multiplicity in
the box and torus source sums.

\paragraph*{The one-bond convolution and its smooth conjugations}

Define

\begin{gather*}
\FitMath{c(g)=e^{-\beta(1-\frac12\operatorname{tr}g)},} \\[3pt]
\FitMath{(CF)(w)=\int_G c(wv^{-1})F(v)\,dv.}
\end{gather*}

The kernel is real, symmetric, and bounded below by $e^{-2\beta}$.
Let $\Pi_l$ be the Peter--Weyl projection corresponding to the
irreducible representation of spin $l\in\frac12\mathbb N_0$ and
dimension $d_l=2l+1$. Schur's lemma and the Weyl integration
formula give

\refstepcounter{equation}
\begin{gather*}
\FitMath{C=\sum_l c_l\Pi_l,} \\[3pt]
\FitMath{c_l=\frac{2e^{-\beta}}{\beta}I_{2l+1}(\beta)>0.}\tag{\theequation}\label{eq:gibbs-symmetry-GL3}
\end{gather*}

For completeness, write $g$ with eigenvalues $e^{it},e^{-it}$ and
use $\chi_l(g)=\sin((2l+1)t)/\sin t$. The character coefficient is

\begin{gather*}
\FitMath{d_lc_l=\frac{2e^{-\beta}}\pi
  \int_0^\pi e^{\beta\cos t}\sin t\sin(d_lt)\,dt
  =e^{-\beta}(I_{2l}-I_{2l+2})(\beta).}
\end{gather*}

The Bessel recurrence gives \eqref{eq:gibbs-symmetry-GL3}. Its series

\begin{gather*}
\FitMath{I_n(\beta)=\sum_{k\ge0}
       \frac{(\beta/2)^{n+2k}}{k!(n+k)!}}
\end{gather*}

has positive terms and implies factorial decay in $n$.
Moreover,

\refstepcounter{equation}
\begin{gather*}
\FitMath{\frac{I_{n+1}}{I_n}\le\frac{\beta}{2(n+1)},} \\[3pt]
\FitMath{\frac{I_n}{I_{n+1}}
  =\frac{2(n+1)}\beta+\frac{I_{n+2}}{I_{n+1}}
  \le C_\beta(1+n).}\tag{\theequation}\label{eq:gibbs-symmetry-GL4}
\end{gather*}

Consequently both ratios between neighboring $c_l$ grow at most
polynomially in $l$. Factorial decay also shows that $C$ is trace
class and that $C^{1/2}$ has a real symmetric smooth kernel.

For a finite block $\mathcal B$ of spatial bonds put

\begin{gather*}
\FitMath{\mathcal C=C_{\mathcal B}=\bigotimes_{b\in\mathcal B}C,} \\[3pt]
\FitMath{c_{\boldsymbol l}=\prod_{b\in\mathcal B}c_{l_b},} \\[3pt]
\FitMath{\Pi_{\boldsymbol l}=\bigotimes_{b\in\mathcal B}\Pi_{l_b}.}
\end{gather*}

Its kernel is
$\mathcal C(u,v)=\prod_{b\in\mathcal B}c(u_bv_b^{-1})$.
In particular it is strictly positive on the compact product
space. It commutes with the simultaneous vertex gauge action on
its two arguments.

Let $f$ be a polynomial in matrix coordinates and their conjugates
on $G^{\mathcal B}$. Define on $C^\infty(G^{\mathcal B})$

\refstepcounter{equation}
\begin{gather*}
\FitMath{\widehat f F=
   \sum_{\boldsymbol l,\boldsymbol l'}
    \frac{c_{\boldsymbol l}}{c_{\boldsymbol l'}}
       \Pi_{\boldsymbol l'}M_f\Pi_{\boldsymbol l}F,}\tag{\theequation}\label{eq:gibbs-symmetry-GL5}
\end{gather*}

where $M_f$ denotes multiplication by $f$. This is a continuous
map of the smooth Frechet space into itself. Indeed, if the degree
in coordinate $b$ is $d_b$, the Clebsch--Gordan rule makes the
summand zero unless $|l_b-l_b'|\le d_b/2$. By \eqref{eq:gibbs-symmetry-GL4}, the remaining
ratios are bounded by a fixed polynomial in the indices. Smooth
functions have Peter--Weyl coefficients decreasing faster than
every polynomial, as follows by applying powers of the Casimir.
The finite-dimensional reproducing-kernel estimate bounds the
supremum norm on a block by a polynomial times its $L^2$ norm.
The same holds after any fixed number of derivatives. These facts
give convergence of \eqref{eq:gibbs-symmetry-GL5} with all derivatives and the asserted
continuity.

The notation $\widehat f=\mathcal C^{-1}M_f\mathcal C$ is therefore
legitimate on this smooth domain when interpreted by \eqref{eq:gibbs-symmetry-GL5}.
It does not assert that $\mathcal C^{-1}$ acts on every smooth
function or that $\widehat f$ is bounded on $L^2$. The formal
$L^2$ adjoint on smooth functions is given by the same finite-band
series with the coefficient ratios reversed and $f$ replaced by
$\overline f$. Those ratios also have polynomial growth, so this
formal adjoint maps smooth functions continuously to smooth
functions. In particular no unbounded inverse is applied outside
its justified spectral domain in the following identity:

\refstepcounter{equation}
\begin{gather*}
\FitMath{\mathcal C\widehat fF=M_f\mathcal CF.}\tag{\theequation}\label{eq:gibbs-symmetry-GL6}
\end{gather*}

It follows by multiplying \eqref{eq:gibbs-symmetry-GL5} by $\mathcal C$ and summing the
projections. In kernel form, for every $w$,

\refstepcounter{equation}
\begin{gather*}
\FitMath{f(w)\int\mathcal C(v,w)F(v)\,dv
       =\int\mathcal C(v,w)(\widehat fF)(v)\,dv.}\tag{\theequation}\label{eq:gibbs-symmetry-GL7}
\end{gather*}

The equality is pointwise, since the series and both integrals
define smooth functions. All these assertions also hold for
smooth families depending on extra variables on a finite compact
product: differentiation in those variables commutes with the
series. The operator $\widehat f$ commutes with multiplication by
functions of variables on which it does not act, and operators
acting on disjoint variable sets commute.

\paragraph*{A transfer contraction inequality}

Fix a direction, called time, and take a cubic torus of side $s$.
Let $\mathcal E$ be its spatial bonds in one layer. On
$L^2(G^{\mathcal E})$, let $P$ be vertex gauge averaging and let
$V$ be $\beta$ times the spatial plaquette action in that layer.
The operator

\begin{gather*}
\FitMath{L=P C_{\mathcal E}^{1/2}e^{-V}C_{\mathcal E}^{1/2}P}
\end{gather*}

is positive and trace class. The convolution, its square root,
and $e^{-V}$ all commute with $P$. Integrating the time-like bonds
in the torus integral inserts $P$ between successive layers.
Using cyclicity of the trace gives

\refstepcounter{equation}
\begin{gather*}
\FitMath{Z_s(0)=\operatorname{Tr}L^s.}\tag{\theequation}\label{eq:gibbs-symmetry-GL8}
\end{gather*}

Here normalized Haar measure makes the gauge average an
orthogonal projection without an additional volume factor.

Let $\Xi$ be a real bounded operator on $L^2(G^{\mathcal B})$
commuting with the block vertex gauge action, and suppose
$\|\Xi\|\le1$. Suppose the kernel

\begin{gather*}
\FitMath{k_\Xi=\text{kernel of }\mathcal C^{1/2}\Xi\mathcal C^{1/2}}
\end{gather*}

is continuous and strictly positive pointwise. In two adjacent
layers, use transported spatial variables $u,v$ and set

\refstepcounter{equation}
\begin{gather*}
\FitMath{\phi_\Xi(u,v)=\log\frac{k_\Xi(u,v)}{\mathcal C(u,v)}.}\tag{\theequation}\label{eq:gibbs-symmetry-GL9}
\end{gather*}

This is a real continuous gauge-invariant local observable.

Place translates of the block so far apart that blocks in a
single time step have disjoint bond sets. Multiplication by
$e^{\phi_\Xi}$ replaces the corresponding factor $\mathcal C$
by $\mathcal C^{1/2}\Xi\mathcal C^{1/2}$. Thus the full kernel in
step $j$ is
$C_{\mathcal E}^{1/2}\Xi_j C_{\mathcal E}^{1/2}$, where $\Xi_j$
is a tensor product of copies of $\Xi$ and identities and has
norm at most one. Gauge averaging the time-like bonds gives
the trace of the product of
$e^{-V}C_{\mathcal E}^{1/2}\Xi_j C_{\mathcal E}^{1/2}P$.
Moving the first half-kernel and layer weight cyclically to the
end, and using commutation with $P$, rewrites this trace as

\begin{gather*}
\FitMath{\operatorname{Tr}(\Xi_1L\Xi_2L\cdots\Xi_sL).}
\end{gather*}

The unperturbed steps have $\Xi_j=1$. Schatten Holder gives
\begin{equation}
\FitMath{\left|\operatorname{Tr}(\Xi_1L\cdots\Xi_sL)\right|
 \le\prod_{j=1}^s\|\Xi_jL\|_s
 \le\|L\|_s^s=\operatorname{Tr}L^s.}
 \label{eq:gibbs-symmetry-GL10}
\end{equation}
The last equality uses $L\ge0$. The modified integral is
positive because every replaced kernel is positive pointwise.
Therefore \eqref{eq:gibbs-symmetry-GL10} is precisely the partition-function inequality in \eqref{eq:gibbs-symmetry-GL1} for a sufficiently sparse periodic array of $\phi_\Xi$.

Given an arbitrary periodic state $\nu$, choose $m$ large enough
that $m\mathbb Z^4$ is a subgroup of its periods and the copies
are disjoint. \eqref{eq:gibbs-symmetry-GL1}--\eqref{eq:gibbs-symmetry-GL2} imply

\refstepcounter{equation}
\begin{gather*}
\FitMath{\nu(\phi_\Xi)\le0.}\tag{\theequation}\label{eq:gibbs-symmetry-GL11}
\end{gather*}

The norm contraction and the lower tangent are the only
ingredients of this conclusion. No self-adjointness of $\Xi$
and no differentiability of a limiting pressure has been used.

\paragraph*{Two adjacent layers cannot carry an antisymmetric mean}

Let $Y(u,v)$ be a real continuous function on
$G^{\mathcal B}\times G^{\mathcal B}$ that is invariant under
simultaneous vertex gauge transformations and antisymmetric:
$Y(v,u)=-Y(u,v)$. Then every periodic DLR state satisfies

\refstepcounter{equation}
\begin{gather*}
\FitMath{\nu(Y(U_i,U_{i+1}))=0.}\tag{\theequation}\label{eq:gibbs-symmetry-GL12}
\end{gather*}

To prove this, first suppose $K(u,v)=Y(u,v)\mathcal C(u,v)$ is
a polynomial. Its integral operator has finite Peter--Weyl
support. Antisymmetry makes $K$ a real skew-adjoint operator.
On that finite-dimensional support define

\begin{gather*}
\FitMath{D=\mathcal C^{-1/2}K\mathcal C^{-1/2}.}
\end{gather*}

Every eigenvalue $c_{\boldsymbol l}$ is strictly positive, so
the definition is valid. Extend $D$ by zero on the orthogonal
complement. It is a finite-rank real skew-adjoint operator,
commutes with gauge transformations, and
$\Xi_\varepsilon=e^{\varepsilon D}$ is a real orthogonal
operator for real $\varepsilon$ of either sign. In particular
$\|\Xi_\varepsilon\|=1$.

The kernel difference from $\mathcal C$ has finite
Peter--Weyl support and admits, uniformly in $(u,v)$,

\begin{gather*}
\FitMath{k_{\Xi_\varepsilon}(u,v)
   =\mathcal C(u,v)+\varepsilon K(u,v)+O(\varepsilon^2).}
\end{gather*}

Since $\mathcal C\ge e^{-2\beta|\mathcal B|}$, this kernel
is strictly positive for sufficiently small $|\varepsilon|$.
Consequently
$\phi_{\Xi_\varepsilon}=\varepsilon Y+O(\varepsilon^2)$
uniformly. Apply \eqref{eq:gibbs-symmetry-GL11}, first at each fixed nonzero
$\varepsilon$ after taking the thermodynamic limit in the
lower-tangent argument. For both signs it gives
$\varepsilon\nu(Y)\le C\varepsilon^2$. Letting
$\varepsilon\to0$ proves \eqref{eq:gibbs-symmetry-GL12} in this case.

For general $Y$, approximate $Y\mathcal C$ uniformly by
polynomials on the compact product group. Antisymmetrize each
approximant and average it over the simultaneous gauge group.
These operations preserve polynomiality, because a finite
matrix-coordinate degree spans a finite-dimensional
representation, and preserve the uniform approximation.
Divide by the strictly positive kernel $\mathcal C$.
The resulting functions $Y_n$ converge uniformly to $Y$ and
have the already proved property $\nu(Y_n)=0$. This proves
\eqref{eq:gibbs-symmetry-GL12}. The argument also specifies why injectivity of
$\mathcal C$, rather than merely its operator positivity,
is needed.

\paragraph*{The exact one-layer DLR elimination identity}

Fix a finite stack of layers $[a,b]$. Write $T_{j,y}$ for
the time-like bond from $(j,y)$ to $(j+1,y)$ and
$\mathring U_{j,e}$ for a spatial bond in layer $j$. Define

\refstepcounter{equation}
\begin{gather*}
\FitMath{g_{j,y}=T_{a,y}\cdots T_{j-1,y},} \\[3pt]
\FitMath{U_{j,\langle y,y'\rangle}
  =g_{j,y}\mathring U_{j,\langle y,y'\rangle}g_{j,y'}^{-1}.}\tag{\theequation}\label{eq:gibbs-symmetry-GL13}
\end{gather*}

These are temporal-gauge coordinates on the stack, with the
empty product at $j=a$. A gauge-invariant local function of
the stack is exactly a function of the $U_j$ invariant under
one common vertex gauge action in all layers. This follows
by setting the time-like bonds to the identity successively.
For any fixed finite observable only finitely many vertices
are involved, so only finite-support gauge transformations
are needed.

Let $a<j<b$ and resample the spatial bonds
$\{j\}\times\mathcal B$, leaving every other original bond
conditioned. The matrices $g_{j,y}$ use only time-like bonds
and hence are fixed during this resampling. The substitution
$\mathring U_j\mapsto v=U_j|_{\mathcal B}$ preserves Haar
measure. The two time-like plaquettes attached to each
resampled spatial bond produce the factors below; all
spatial plaquettes that meet the block produce the finite
potential $V_j(v;\eta)$, where $\eta$ denotes the other
spatial variables of layer $j$. Thus the conditional density
is exactly

\refstepcounter{equation}
\begin{gather*}
\FitMath{\frac1{z(w,w',\eta)}
 \mathcal C(w,v)e^{-V_j(v;\eta)}\mathcal C(v,w')\,dv,} \\[3pt]
\FitMath{w=U_{j-1}|_{\mathcal B},} \\[3pt]
\FitMath{w'=U_{j+1}|_{\mathcal B}.}\tag{\theequation}\label{eq:gibbs-symmetry-GL14}
\end{gather*}

No plaquette meeting the resampled bonds is omitted: it is
either spatial in layer $j$ or one of the two time-like
plaquettes of that bond. The potential can involve spatial
neighbors outside $\mathcal B$, but they lie in the same
layer and in a finite enlargement of the block.

Let $f$ be a real polynomial on $G^{\mathcal B}$ and let
$\Phi$ be smooth in finitely many transported variables
from layers $j$ and above, and independent of the spatial
variables of layer $j-1$. Define

\refstepcounter{equation}
\begin{gather*}
\FitMath{\Phi^\flat
  =\frac{e^{V_j(v;\eta)}}{\mathcal C(v,w')}
     \widehat f_v\!\left[
        e^{-V_j(v;\eta)}\mathcal C(v,w')\Phi\right].}\tag{\theequation}\label{eq:gibbs-symmetry-GL15}
\end{gather*}

All factors are smooth on a finite compact product, and
the denominator is strictly positive. \eqref{eq:gibbs-symmetry-GL5}
therefore proves that $\Phi^\flat$ is smooth. Its new
dependence is confined to layer $j$ and the already present
higher layers; it acquires no spatial variable in layer
$j-1$.

Every DLR state, periodic or not, satisfies the exact identity
\begin{equation}
\FitMath{\nu\big(f(U_{j-1}|_{\mathcal B})\Phi\big)
       =\nu(\Phi^\flat).}
 \label{eq:gibbs-symmetry-GL16}
\end{equation}
Indeed, condition on all bonds except the specified block
in layer $j$. The conditional numerator on the left is

\begin{gather*}
\FitMath{f(w)\int\mathcal C(v,w)
       e^{-V_j(v;\eta)}\mathcal C(v,w')\Phi\,dv.}
\end{gather*}

Apply \eqref{eq:gibbs-symmetry-GL7} to the smooth bracketed function. Substitution
of \eqref{eq:gibbs-symmetry-GL15} identifies the result with the numerator of the
conditional expectation on the right. The same normalizer
$z(w,w',\eta)$ occurs on both sides; it is not
differentiated or canceled through an operator. Integration
of this conditional identity gives \eqref{eq:gibbs-symmetry-GL16}. Reflection of
the construction yields the corresponding identity that
eliminates an upper outer layer.

\paragraph*{Removing two outer layers and preserving reflection}

Let $\theta$ be reflection about the middle of a stack with
$2k+2$ layers, numbered $i-k,\ldots,i+k+1$. Let $X$ be a
real gauge-invariant polynomial observable in that stack.
In temporal coordinates, a finite sum of products writes
it as

\begin{gather*}
\FitMath{X=\sum_\alpha f_\alpha(w)h_\alpha(z)\Phi_\alpha,}
\end{gather*}

where $w=U_{i-k}|_{\mathcal B}$ and
$z=U_{i+k+1}|_{\mathcal B}$ are the outer layers and
$\Phi_\alpha$ depends on the inner layers. Take one finite
block containing all the variables that occur in this
decomposition. The potentials used during elimination can
enlarge this block spatially, but only finitely and only in
the inner layers.

Use \eqref{eq:gibbs-symmetry-GL16} to remove $f_\alpha(w)$ through the first inner
layer $u=U_{i-k+1}|_{\mathcal B}$, and then its reflected
version to remove $h_\alpha(z)$ through the last inner
layer $v=U_{i+k}|_{\mathcal B}$. The result is

\refstepcounter{equation}
\begin{gather*}
\FitMath{\nu(X)=\nu(EX),} \\[3pt]
\FitMath{EX=\sum_\alpha \Upsilon^{-1}
       \widehat f_{\alpha,u}\widehat h_{\alpha,v}
          (\Upsilon\Phi_\alpha).}\tag{\theequation}\label{eq:gibbs-symmetry-GL17}
\end{gather*}

When the two inner layers are not adjacent, $\Upsilon$ is
the product of their two factors $e^{-V}$ and their bond
kernels toward the remaining interior. When $k=1$ they
are adjacent and
\begin{equation}
\FitMath{\Upsilon=e^{-V(u;\eta)}e^{-V(v;\eta')}
                    \mathcal C(u,v).}
 \label{eq:gibbs-symmetry-GL18}
\end{equation}
The common bond kernel occurs once, not twice. To see the
cancellation directly, the first elimination introduces
$e^{V(u)}/\mathcal C(u,v)$ outside its conjugated
operator. In the upper elimination the conditional
numerator contains $e^{-V(v)}\mathcal C(u,v)$, canceling
that denominator before the second operator acts. The
operators act on distinct variables and commute. Moving
the remaining independent factors through them gives
exactly \eqref{eq:gibbs-symmetry-GL18}. This calculation is also a check that
normalizing factors from different conditional laws have
not been interchanged.

Let $R$ denote literal reversal of the array of transported
layer coordinates. An actual reflection of the original
bond configuration induces $R$ followed by a common vertex
gauge transformation arising from the change of temporal-gauge
base. The spatial potential is the same function of its own
layer variables at both ends, and $\mathcal C(u,v)$ is
symmetric. Reflection interchanges $u$ and $v$, including
all their spatial neighbors, and preserves $\Upsilon$.
It also sends $\widehat f_u$ to $\widehat f_v$.
Consequently coordinate elimination commutes with literal
reversal:

\refstepcounter{equation}
\begin{gather*}
\FitMath{E(RX)=R(EX).}\tag{\theequation}\label{eq:gibbs-symmetry-GL19}
\end{gather*}

The output need not be a polynomial, but it is a smooth
function of finitely many inner-layer variables.

Every DLR state is invariant under finite-support vertex
gauge transformations. This follows by applying its
conditional law to the finite set of bonds incident to
the transformed vertices and changing variables using
Haar invariance and plaquette conjugacy. We may therefore
replace $EX$ in expectation by its average under the
common temporal-gauge action; denote it by $\mathcal A$.
This averaging commutes with $R$. Once averaged, the choice of the
lowest layer in \eqref{eq:gibbs-symmetry-GL13} disappears: changing the base from
$i-k$ to $i-k+1$ conjugates all remaining transported
variables by one common vertex gauge transformation.
The averaged function is invariant under that
transformation, even when its values are the actual
time-like transport matrices. This last statement is a
pointwise identity of invariant functions, and does not
invoke invariance of the state under a field-dependent
change of variables. Thus the averaged output is a
continuous gauge-invariant observable of the smaller
stack of $2k$ layers. Since $X$ is gauge invariant, its actual
reflection is represented by $RX$ in temporal coordinates.
After averaging the common gauge factor on the output also
disappears. Thus the identity needed on the original bond
space is $\mathcal A E(\theta X)=\theta\mathcal A EX$.

\paragraph*{Proof of the link-reflection theorem}

Fix the reflection $\theta$ in $x_\mu=i+\frac12$.
We prove by induction on $k$ that every real
gauge-invariant polynomial in the stack
$[i-k+1,i+k]$ has equal expectation to its reflection.

For $k=1$, temporal gauge writes the observable as
$\Psi(U_i,U_{i+1})$. Reflection interchanges the two
arguments, up to a common gauge transformation, which
does not affect $\Psi$. The difference from its
reflection is the function
$Y(u,v)=\Psi(u,v)-\Psi(v,u)$, to which \eqref{eq:gibbs-symmetry-GL12}
applies. Hence the expectations agree.

Assume the assertion for $k$. For a polynomial $X$ in
the stack with $2k+2$ layers, \eqref{eq:gibbs-symmetry-GL17}--\eqref{eq:gibbs-symmetry-GL19} and gauge
averaging reduce the comparison of $\nu(X)$ and
$\nu(\theta X)$ to the same comparison for a
continuous gauge-invariant function of the smaller
stack. Invariant polynomials are uniformly dense in
the invariant continuous functions on a finite product
of compact matrix groups: first apply Stone--Weierstrass
to matrix coordinates and their conjugates, then
average the approximants over the finite vertex gauge
group. The induction hypothesis therefore applies to
this continuous output. It gives
$\nu(\mathcal A EX)=\nu(\theta\mathcal A EX)$ and proves
the next step.

Every gauge-invariant continuous local function fits
in a finite symmetric stack and is uniformly
approximated as above. We have proved equality of
its expectation and that of its reflection.
For a general continuous local $F$, gauge average
over the finite set of vertices relevant to its
support. Gauge invariance of both $\nu$ and its
reflected state shows that both expectations of
$F$ equal the expectations of that gauge average.
Thus $\nu(F)=\nu(\theta F)$ for all continuous
local $F$, which determine the probability measure
on the compact product configuration space. The
state is reflection invariant. The direction $\mu$
and integer $i$ were arbitrary. Composing
reflections in two successive link planes gives
translation by $2e_\mu$, proving the stated
$2\mathbb Z^4$ invariance.\hfill$\square$

\paragraph*{Passage to all mirrors}

The common period is the input to the independent fold-and-separation
proof in \cref{ch:all-mirror-periodic}. That section supplies
adapted diagonal split tori, finite crossing kernels, the full
chessboard inequality, and the arbitrary-state separation argument.
The link theorem above does not require those later ingredients.


% END MODULE chapters/gibbs_symmetry_v2.tex

% BEGIN MODULE chapters/all_mirror_rp_audit_v3.tex
\subsection{Reflection positivity of periodic Wilson states}
\label{ch:all-mirror-periodic}

The result in this section concerns the complete algebra of bonds on a
closed side of a mirror. Its proof does not identify an arbitrary Gibbs
state with a limit of unperturbed torus measures. Instead, a pressure
inequality produces, for each separating observable, a suitable limit of
\emph{perturbed} torus measures. The common period supplied by link
reflection invariance is what makes the separation argument work.

We use the following result proved in the preceding link-symmetry
argument: every periodic DLR state of the four-dimensional
$\operatorname{SU}(2)$ Wilson specification, at $\beta>0$, is invariant under
all coordinate link reflections and consequently under
$P=2\mathbb Z^4$. No uniqueness or extremality is assumed here.

Let $\Omega=\operatorname{SU}(2)^{\mathcal B}$, with one group variable for each
positively oriented nearest-neighbour bond; reversing a bond inverts its
variable. Put

\begin{gather*}
\FitMath{a_p(U)=1-\tfrac12\operatorname{tr}U_p,} \\[3pt]
\FitMath{c(g)=\exp[-\beta(1-\tfrac12\operatorname{tr}g)].}
\end{gather*}

The trace is real for $\operatorname{SU}(2)$, and $0\le a_p\le2$.
The product topology makes $\Omega$ compact and metrizable. A state is
periodic if it is invariant under some finite-index subgroup of
$\mathbb Z^4$. Denote by $\mathcal G_\beta^P$ the compact convex set of
Wilson DLR states invariant under $P$.

The normals considered below are

\begin{gather*}
\FitMath{\mathcal N=\{e_\mu:1\le\mu\le4\}
 \cup\{e_\mu+e_\lambda,e_\mu-e_\lambda:1\le\mu<\lambda\le4\}.}
\end{gather*}

For $n\in\mathcal N$ write $q=|n|^2$ and

\begin{gather*}
\FitMath{\theta_c x=x-2(n\cdot x-c)n/q.}
\end{gather*}

The admissible positions are $c\in\frac12\mathbb Z$ when $q=1$ and
$c\in\mathbb Z$ when $q=2$. These and only these positions in the stated
families give lattice reflections. Let $\mathcal A_c^+$ be the bounded
Borel functions of bonds both of whose endpoints satisfy $n\cdot x\ge c$,
and define $\mathcal A_c^-$ similarly. The antilinear reflection is
$\Theta_cF=\overline{F\circ\theta_c^*}$.

\begin{proposition}[All-mirror positivity of periodic states]
\label{thm:amrp-all-periodic}
Every periodic Wilson DLR state at $\beta>0$ is invariant under every
admissible reflection $\theta_c$ in the sixteen families above. For every
such mirror,

\refstepcounter{equation}
\begin{gather*}
\FitMath{\nu(\Theta_cF\,F)\ge0} \\[3pt]
\FitMath{(F\in\mathcal A_c^+\text{ or }F\in\mathcal A_c^-).}\tag{\theequation}\label{eq:all-mirror-rp-AM1}
\end{gather*}

In particular, the conclusion applies to every bounded observable on
any finite stack contained in either closed half-space, regardless of
the number of layers in that stack.
\end{proposition}

Composing the reflections in $x_\mu=0$ and $x_\mu=1/2$ gives
translation by one lattice unit. Coordinate reflections give sign
changes, and reflection in $x_\mu-x_\lambda=0$ exchanges two
coordinates. Thus every such state is invariant under all lattice
translations and the full signed-permutation group. At $\beta=0$,
the DLR equations directly force every finite-bond marginal to be
product Haar; that unique product state has the same properties.

\subsubsection*{Positive crossing kernels on split tori}

Write an element of $\operatorname{SU}(2)$ as a unit quaternion
$s(g)\in S^3\subset\mathbb R^4$. Then

\begin{gather*}
\FitMath{\tfrac12\operatorname{tr}(gh^{-1})=s(g)\cdot s(h).}
\end{gather*}

Consequently

\refstepcounter{equation}
\begin{gather*}
\FitMath{c(gh^{-1})=e^{-\beta}\sum_{m=0}^\infty\frac{\beta^m}{m!}
 \sum_{i_1,\ldots,i_m=1}^4
 \prod_{j=1}^m s_{i_j}(g)\prod_{j=1}^m s_{i_j}(h).}\tag{\theequation}\label{eq:all-mirror-rp-AM2}
\end{gather*}

Every coefficient is nonnegative. Each feature has absolute value at
most one, and the sum of the coefficients is $e^{3\beta}$, so the
series converges absolutely and uniformly. Tensor products of finitely
many such expansions are again uniformly convergent positive Gram
expansions. The coarse total coefficient bound is immaterial; each
torus used below is finite.

A torus is called split for the mirror normal $n$ if its period lattice
has the form

\refstepcounter{equation}
\begin{gather*}
\FitMath{\Gamma=\Gamma_\parallel\oplus\mathbb Z(2z n/q),} \\[3pt]
\FitMath{\Gamma_\parallel\subset n^\perp,}\tag{\theequation}\label{eq:all-mirror-rp-AM3}
\end{gather*}

where $\Gamma_\parallel$ has rank three, $2zn/q$ is a lattice vector,
and $z$ is a sufficiently large integer. The reflection has two fixed
cuts, at heights $c$ and $c+z$, on this torus. The two intervening closed
halves are exchanged. All tori below have arbitrarily large shortest
nonzero period, so that no small-quotient ambiguity affects plaquettes.

We verify reflection positivity on these tori for all bounded
half-space functions. Partition the bond variables into the interiors
of the two halves, the bonds lying on the cuts, and, when present, the
bonds crossing a cut. A coordinate site cut has no plaquette with
vertices strictly on both sides. Conditional on its boundary bonds,
the action is a sum of a positive-side action, its reflected action,
and an action on the boundary. Its reflection form is thus an integral
of absolute squares.

At a diagonal cut, signed permutations reduce the geometry to
$n=e_1+e_2$, $c\in\mathbb Z$. A bond changes the height $x_1+x_2$ by
$0$ or $1$ in absolute value, so no bond crosses from a strictly
negative height to a strictly positive one. The only plaquettes meeting
both strict sides are in the $(1,2)$ plane, with heights
$c-1,c,c,c+1$. Each is cut into two two-bond paths between the two
vertices at height $c$. Orienting both paths in the same direction
expresses its Wilson factor as $c(P_+P_-^{-1})$. Reflection interchanges
the paths, and \eqref{eq:all-mirror-rp-AM2} supplies the required Gram expansion. Plaquettes
in all other planes lie in a closed half. The bonds on the mirror, in
directions $3$ and $4$, are fixed pointwise. Conditioning on them and
multiplying the crossing expansions again gives an integral of squares.
The same argument applies at the second cut.

For a coordinate link cut there are crossing bonds. Integrate them
before evaluating a half-space function, which cannot depend on them.
Let $X$ and $Y$ be the collections of bonds parallel to the cut in its
two adjacent layers. The product crossing kernel before integration is

\begin{gather*}
\FitMath{C(X,Y)=\prod_b c(X_bY_b^{-1}).}
\end{gather*}

It is a positive kernel by \eqref{eq:all-mirror-rp-AM2}. If $T=(T_x)$ is a collection of group
elements at the vertices of one boundary layer, set
$(T\cdot X)_{xy}=T_xX_{xy}T_y^{-1}$. Let $P_0$ be Haar averaging under
this action. Simultaneous gauge covariance gives
$C(T\cdot X,T\cdot Y)=C(X,Y)$, since each plaquette trace changes by
conjugation. Hence the associated integral operator commutes with the
orthogonal projection $P_0$. Integration of the crossing bonds gives
the kernel of

\refstepcounter{equation}
\begin{gather*}
\FitMath{P_0C=CP_0=P_0CP_0\ge0.}\tag{\theequation}\label{eq:all-mirror-rp-AM4}
\end{gather*}

Equivalently, in each Gram feature of the product expansion one may
Haar-average both its $X$ argument and its $Y$ argument. Absolute
uniform convergence justifies this operation. The crossing bonds at
the two cuts are disjoint, so the two constructions multiply. This
argument places no gauge-invariance condition on the test function.

In all three geometries the resulting reflection form has the form

\refstepcounter{equation}
\begin{gather*}
\FitMath{\int\Theta_cF' F\,e^{-\beta\sum_pa_p}\,dU
   =\sum_\gamma b_\gamma\int
       \overline{I_\gamma(F';Y)}I_\gamma(F;Y)\,d\lambda(Y),} \\[3pt]
\FitMath{b_\gamma\ge0,}\tag{\theequation}\label{eq:all-mirror-rp-AM5}
\end{gather*}

where $Y$ denotes the fixed boundary variables; in the link case $Y$
is absent. The functions $I_\gamma$ integrate over one half, with its
Boltzmann factor and the corresponding Gram feature. The boundary
Boltzmann factor is included in the positive measure $\lambda$.
Fubini's theorem follows from the coefficient summability above.
This proves the Cauchy--Schwarz inequality for the reflection form and
reflection positivity for all bounded half-space functions. Haar
invariance, orientation reversal, and invariance of the plaquette trace
also prove reflection invariance of the finite measure.

If $E$ is a bounded real function on one closed half, multiplying the
density by $e^{E+\theta_cE}$ preserves both properties: insert $Fe^E$
in \eqref{eq:all-mirror-rp-AM5}. This observation will apply to folded perturbations.

\subsubsection*{A finite one-direction chessboard proof}

Fix parallel slab faces at heights $c+kH$, where $H\ge2$ is an integer.
Use a split torus with normal height period $2MH$, $M\ge1$. Its $2M$
slabs are indexed by $\mathbb Z/(2M)$. The group generated by reflection
in the slab faces acts simply transitively on the slabs: its
translations change the index by an even integer and its reflections
send $k$ to $2j-1-k$. Thus a function on any one slab has a well-defined
copy on every other slab obtained using this group. Boundary bonds may
belong to two closed slabs; this causes no problem, since the products
are functions of the same boundary variables and \eqref{eq:all-mirror-rp-AM5} permits those
variables in both half-algebras.

For positive bounded slab functions $f_0,\ldots,f_{2M-1}$ with strictly
positive infima, denote by $Z^{\rm dis}(f_j)$ the partition function
obtained by copying $f_j$ into every slab. We claim

\refstepcounter{equation}
\begin{gather*}
\FitMath{\int\prod_{j=0}^{2M-1}f_j\,e^{-\beta\sum_pa_p}\,dU
 \le\prod_{j=0}^{2M-1}Z^{\rm dis}(f_j)^{1/(2M)}.}\tag{\theequation}\label{eq:all-mirror-rp-AM6}
\end{gather*}

Here is a proof valid for every $M$, without a power-of-two restriction.
Transport the $2M$ functions to one reference slab and retain their
labels. For any filling $\sigma:\mathbb Z/(2M)\to\{0,\ldots,2M-1\}$,
let $Z(\sigma)$ be the partition function obtained by placing profile
$\sigma(k)$ in slab $k$. Put

\begin{gather*}
\FitMath{w_j=Z^{\rm dis}(f_j)^{1/(2M)},} \\[3pt]
\FitMath{T(\sigma)=Z(\sigma)\big/\prod_kw_{\sigma(k)}.}
\end{gather*}

All these numbers are positive, and every constant filling has $T=1$.
Reflection Cauchy--Schwarz at any pair of opposite slab faces gives

\refstepcounter{equation}
\begin{gather*}
\FitMath{T(\sigma)^2\le T(\sigma_+)T(\sigma_-),}\tag{\theequation}\label{eq:all-mirror-rp-AM7}
\end{gather*}

where $\sigma_+$ retains one semicircle of labels and reflects it into
the other, and $\sigma_-$ does the reverse. Indeed, the unnormalized
inequality is \eqref{eq:all-mirror-rp-AM5}; the two denominators together contain each
original denominator factor twice.

There are finitely many fillings. If $\sigma$ maximizes $T$, \eqref{eq:all-mirror-rp-AM7}
forces both reflected fillings to maximize it as well. Take a cyclic
constant run of length $l$ in a maximizing filling and reflect at the
face at the start of that run. If $l\le M$, retaining the half beginning
there gives a maximizing filling with a run of length $2l$. If
$l\ge M$, retain a semicircle inside the run and reflect it; the result
is constant. Iterating therefore produces a constant maximizing
filling. The maximum is one. Applying this to the original filling
proves \eqref{eq:all-mirror-rp-AM6}.

This argument uses reflection positivity only for finite tori. It does
not infer the positivity of an arbitrary DLR state from the positivity
of one adjacent-layer kernel.

\subsubsection*{A lower tangent for an arbitrary periodic DLR state}

Let $Q$ be a fixed finite-index period subgroup, let $A$ be a real
continuous local observable, and form the locally finite interaction
array $\mathbb A=(\tau_vA)_{v\in Q}$. Write
$\kappa=[\mathbb Z^4:Q]^{-1}$. Let $\Gamma_i\subset Q$ be rectangular
adapted period lattices with every side tending to infinity. On the
corresponding tori put

\begin{gather*}
\FitMath{Z_i(t\mathbb A)=\int
   \exp\left[t\sum_{v\in Q/\Gamma_i}(\tau_vA)^{\Gamma_i}
                    -\beta\sum_pa_p\right]dU.}
\end{gather*}

For every $Q$-invariant Wilson DLR state $\nu$ and every fixed real $t$,

\refstepcounter{equation}
\begin{gather*}
\FitMath{\frac1{|D_i|}\log\frac{Z_i(t\mathbb A)}{Z_i(0)}
 \ge t\kappa\nu(A)-o(1).}\tag{\theequation}\label{eq:all-mirror-rp-AM8}
\end{gather*}

The error may depend on $A,Q,\beta,t$, but is uniform in $\nu$.

For completeness, choose a fundamental parallelepiped $D_i$ and
integrate the positively oriented bonds anchored in it, using arbitrary
external bonds as boundary conditions. Let $W_i$ sum the source copies
whose entire support is interior to this bond set. Conditional Jensen
and the DLR equation give

\begin{gather*}
\FitMath{\log\frac{Z_{D_i}^{t}(\eta)}{Z_{D_i}^{0}(\eta)}
 \ge t\gamma_{D_i}W_i(\eta),} \\[3pt]
\FitMath{\int\gamma_{D_i}W_i\,d\nu=\nu(W_i).}
\end{gather*}

The open and torus Hamiltonians differ only in a boundary layer of
fixed thickness, determined by the range of $A$ and the plaquettes.
Each energy term is bounded; each bond is met by only a fixed finite
number of array terms. Thus their logarithmic partition ratios differ
by at most
$C_A(\beta+|t|\|A\|)|\partial_R D_i|$, uniformly in $\eta$. The number
of omitted array anchors is at most $C_A|\partial_R D_i|$. Since the
remaining $Q$-translates all have expectation $\nu(A)$, integrating the
Jensen bound proves \eqref{eq:all-mirror-rp-AM8}, with an error bounded by a constant times
$|\partial_R D_i|/|D_i|$. This ratio tends to zero for the adapted
rectangular tori, including the diagonal ones. No differentiability of
infinite-volume pressure is used.

\subsubsection*{Fold domination with a fixed common period}

Fix one admissible mirror $(n,c)$. Choose an even positive integer $h$
and set

\refstepcounter{equation}
\begin{gather*}
\FitMath{w=hn\in P,} \\[3pt]
\FitMath{H=n\cdot w=hq,} \\[3pt]
\FitMath{P_\parallel=P\cap n^\perp,} \\[3pt]
\FitMath{Q=P_\parallel\oplus\mathbb Zw,} \\[3pt]
\FitMath{Q_f=P_\parallel\oplus\mathbb Z(2w).}\tag{\theequation}\label{eq:all-mirror-rp-AM9}
\end{gather*}

Both $Q$ and $Q_f$ are finite-index subgroups of $P$ and
$[Q:Q_f]=2$. The tangent lattice $P_\parallel$ has an orthogonal
integer basis: for a coordinate normal use the three other vectors
$2e_j$; for $e_1\pm e_2$ use $2(e_1\mp e_2),2e_3,2e_4$, with the
corresponding signed permutations for the other cases. Thus

\refstepcounter{equation}
\begin{gather*}
\FitMath{\Gamma_{M,J}=J P_\parallel\oplus\mathbb Z(2Mw)}\tag{\theequation}\label{eq:all-mirror-rp-AM10}
\end{gather*}

is an adapted rectangular split torus. Its shortest period tends to
infinity as $M,J\to\infty$.

Assume for the moment that $A$ is supported in the slab
$c\le n\cdot x\le c+H$. The unfolded array contains all its translates
by $Q$. The folded array contains translates of $A$ in the even slabs
and the reflected copies $\theta_{c+H}A$ in the odd slabs, with period
$Q_f$. It is invariant under reflection at every face $c+kH$.
In \eqref{eq:all-mirror-rp-AM6} take

\begin{gather*}
\FitMath{f_k=\exp\left[t\sum_{v\in P_\parallel/JP_\parallel}
                    (\tau_{v+kw}A)^{\Gamma_{M,J}}\right],} \\[3pt]
\FitMath{t>0.}
\end{gather*}

Their product is the unfolded perturbation. Their disseminated
partition functions are identical: normal translation by $kw$ carries
one dissemination to another, permutes the slab faces, and preserves
the unperturbed Wilson measure. Each is the folded partition function.
Hence

\refstepcounter{equation}
\begin{gather*}
\FitMath{Z_{M,J}(t\mathbb A^{\rm u})
 \le Z_{M,J}(t\mathbb A^{\rm f}).}\tag{\theequation}\label{eq:all-mirror-rp-AM11}
\end{gather*}


Let $\mu_{M,J,t}$ be the normalized folded measure, and define

\begin{gather*}
\FitMath{p_{M,J}(t)=|D_{M,J}|^{-1}
       \log[Z_{M,J}(t\mathbb A^{\rm f})/Z_{M,J}(0)].}
\end{gather*}

It is convex, $p_{M,J}(0)=0$, and therefore
$p_{M,J}(t)\le t p'_{M,J}(t)$ for $t>0$. The folded measure is invariant
under $Q_f$ and under $\theta_{c+H}$. All its source terms consequently
have the same expectation as $A$. There are $\kappa|D_{M,J}|$ terms,
where $\kappa=[\mathbb Z^4:Q]^{-1}$, so

\begin{gather*}
\FitMath{p'_{M,J}(t)=\kappa\mu_{M,J,t}(A).}
\end{gather*}

Combining this identity, \eqref{eq:all-mirror-rp-AM8}, and \eqref{eq:all-mirror-rp-AM11} gives, for every
$\nu\in\mathcal G_\beta^P$,

\refstepcounter{equation}
\begin{gather*}
\FitMath{t\kappa\nu(A)-o(1)
 \le p_{M,J}(t)\le t\kappa\mu_{M,J,t}(A).}\tag{\theequation}\label{eq:all-mirror-rp-AM12}
\end{gather*}


The finite folded measure is invariant and reflection positive at
every face $c+kH$. Indeed, its perturbation is $E+\theta_{c+kH}E$,
where $E$ is the sum of its copies lying in one closed half of the
torus. Copies on a boundary can be assigned half weight to each side.
The observation following \eqref{eq:all-mirror-rp-AM5} applies.

First fix $t>0$ and let $M,J\to\infty$. Compactness gives a limit
$\rho_t$. It is a DLR state for the folded, finite-range perturbed
interaction, is $Q_f$-periodic, is invariant and reflection positive at
all its slab faces, and satisfies $\nu(A)\le\rho_t(A)$. Here DLR passage
is local and exact: any fixed collection of bonds, its interaction
neighbourhood, and the support of a cylinder test embed injectively in
all sufficiently large tori. The finite torus conditional equation is
then exactly the required local DLR equation on that test. Its kernel
is continuous, so the equation passes to a weak limit. A finite test
on one half-space eventually lies in a single torus half; hence its
reflection inequality passes to the limit as well.

Next let $t\downarrow0$ along a subsequence for which $\rho_t$ converges
to $\rho_0$. On any fixed finite set, only finitely many source copies
meet that set. Their bounded total contribution is $O(t)$ uniformly in
the exterior configuration. The perturbed conditional densities thus
converge uniformly to the Wilson conditional density. It follows that
$\rho_0$ is a Wilson DLR state. It retains $Q_f$ periodicity, reflection
invariance and positivity at the slab faces, and
$\nu(A)\le\rho_0(A)$.

Now apply the previously proved link-symmetry theorem to $\rho_0$.
Its finite-index period $Q_f$ is enough for that theorem; the result
upgrades its period to $P$. We have proved the following domination:

\refstepcounter{equation}
\begin{gather*}
\FitMath{\begin{gathered}
 A\text{ supported in }[c,c+H],\quad \nu\in\mathcal G_\beta^P
 \quad\Longrightarrow\quad
 \exists\rho_0\in\mathcal G_\beta^P:
 \\
 \rho_0\text{ is invariant and positive at }(n,c),
 \qquad \nu(A)\le\rho_0(A).
 \end{gathered}}\tag{\theequation}\label{eq:all-mirror-rp-AM13}
\end{gather*}

The order of limits is essential: the surface error in \eqref{eq:all-mirror-rp-AM12} is
divided by $t\kappa$, so volume tends to infinity before $t$ tends to
zero.

\subsubsection*{Separation closes the full half-space statement}

Let $K=\mathcal G_\beta^P$ and let $K_c\subset K$ consist of the states
invariant under $\theta_c$ and positive on every continuous cylinder
function in the positive half-space. The set $K$ is compact and convex:
the DLR kernels preserve continuous cylinder functions, and all its
defining equations and period equations are closed affine conditions.
The set $K_c$ is also closed and convex. It is nonempty, for example
by the folded construction with $A=0$.

For any real continuous local $A$, choose $v\in P$ so that $\tau_vA$
is supported strictly above the fixed mirror $c$. Then choose the even
$h$ large enough that this support lies below $c+H$. Such a $v$ exists
because $n\cdot(2ke_j)=2k$ for a coordinate $j$ with $n_j=1$.
Apply \eqref{eq:all-mirror-rp-AM13} to $\tau_vA$. Both $\nu$ and the resulting $\rho_0$ are
$P$-invariant, so

\refstepcounter{equation}
\begin{gather*}
\FitMath{\nu(A)=\nu(\tau_vA)\le\rho_0(\tau_vA)=\rho_0(A),} \\[3pt]
\FitMath{\rho_0\in K_c.}\tag{\theequation}\label{eq:all-mirror-rp-AM14}
\end{gather*}

If some $\nu\in K$ were outside $K_c$, separation of a point from a
compact convex set of measures would supply a real continuous function
$A$ with $\nu(A)>\sup_{\rho\in K_c}\rho(A)$. Continuous cylinder
functions are uniformly dense in $C(\Omega)$, so $A$ could be chosen
local while retaining a strict inequality. This contradicts \eqref{eq:all-mirror-rp-AM14}.
Therefore $K=K_c$.

This step explains why the common period cannot be omitted. A state
produced by folding initially has only the normal period $2w$. Such a
large period need not move a test straddling the mirror into one slab.
The link-symmetry theorem supplies $2\mathbb Z^4$ for both states in
\eqref{eq:all-mirror-rp-AM14}, independently of the observable and slab width. No claim that
periodic DLR states are torus-selected is needed.

The mirror was arbitrary, so every periodic state is invariant and
positive at all the stated mirrors on continuous local half-space
functions. Finally, on the compact metrizable product of the positive
half-space bonds, continuous cylinder functions are dense in
$L^2(\nu)$. For bounded Borel $F$ choose such $F_j$ with
$\|F_j-F\|_2\to0$. Reflection invariance makes $\Theta_c$ an isometry
of $L^2(\nu)$, and Cauchy--Schwarz gives

\begin{gather*}
\FitMath{\left|\nu(\Theta_cF_jF_j)-\nu(\Theta_cFF)\right|
 \le\|F_j-F\|_2(\|F_j\|_2+\|F\|_2)\longrightarrow0.}
\end{gather*}

This proves \eqref{eq:all-mirror-rp-AM1} on the complete positive-half algebra. Replacing a
negative-half function by its reflected conjugate proves the other
half. In particular, arbitrary finite stacks and their bounded Borel
tests are covered.\hfill$\square$

As a further consequence, every periodic state is invariant under all
unit lattice translations: in direction $e_\mu$ the product of the
site reflection at $0$ and the link reflection at $1/2$ is translation
by $e_\mu$. It is also invariant under the hypercubic group. The
coordinate reflections at zero generate coordinate sign changes, and
the diagonal reflections with normal $e_\mu-e_\lambda$ interchange
coordinates. These transformations generate all signed permutations.

\subsubsection*{Scope of the implication}

Positive semidefinite contractions of an adjacent-layer kernel can
produce a one-layer positivity test, just as skew contractions produce
an antisymmetric invariance test. That observation alone supplies no
identity reducing arbitrary odd or even stacks to a single layer.
The proof above uses a different route: exact finite-torus positivity
on the full half-algebra, a finite chessboard inequality, a lower
tangent for each arbitrary periodic DLR state, and separation after
the common period is established. Thus the all-stack and diagonal
conclusions do not rest on a separate layer-elimination positivity
rule. The result is at a fixed lattice spacing and coupling; it makes
no continuum-limit or interacting renormalization-group assertion.


% END MODULE chapters/all_mirror_rp_audit_v3.tex

% BEGIN MODULE chapters/phase_plateau_audit_v3.tex
\subsection{A uniform lattice plateau estimate}
\label{ch:plateau-proof}

The estimate in this section is a consequence of translation invariance and
reflection positivity in coordinate and diagonal directions. No infrared bound,
mass gap, or differentiability of a pressure is assumed. The Gibbs equations
are not needed in the proof. Their role in the Wilson application is to place
every periodic Wilson Gibbs state in the reflection-positive class by
\cref{thm:amrp-all-periodic}, proved independently in
\cref{ch:all-mirror-periodic}. The state here is a measure on the
infinite lattice, rather than a finite-torus law lifted periodically.

By dilation it suffices initially to take lattice spacing one. Let $\nu$ be
a translation-invariant probability measure on lattice configurations. A local
function means a bounded Borel function of finitely many bonds. The algebra
$\mathfrak A(S)$ contains the functions of bonds with both endpoints in $S$.
Assume reflection positivity on the closed side of every coordinate mirror
$x_\mu=c$, $c\in\frac12\mathbb Z$, and every diagonal mirror
$x_\mu\pm x_\lambda=c$, $c\in\mathbb Z$. These are sixteen families.
Reflections act linearly on functions by pullback, without complex conjugation.
Write $\theta_\mu$ for reflection in $x_\mu=0$.

For $r\ge1$ an integer put $Q_r=[-r,r]^4$. For $A\in\mathfrak A(Q_r)$
define

\refstepcounter{equation}
\begin{gather*}
\FitMath{\mathcal N_{\nu,r}(A)
 =\max_{\mu,\gamma}
   \nu\bigl(\overline{\theta_\mu\gamma A}
                     \tau_{2re_\mu}\gamma A\bigr),}\tag{\theequation}\label{eq:phase-plateau-PL1}
\end{gather*}

where $\gamma$ ranges over the sixteen coordinate sign changes. The terms
in this maximum are nonnegative squared reflection norms. Let $\mathcal I$
be the completed invariant $\sigma$-field for all lattice translations and set

\refstepcounter{equation}
\begin{gather*}
\FitMath{\operatorname{Pl}_\nu(A,B)
   =\nu\bigl(\overline{E_\nu[A\mid\mathcal I]}
                           E_\nu[B\mid\mathcal I]\bigr).}\tag{\theequation}\label{eq:phase-plateau-PL2}
\end{gather*}


\begin{proposition}[Uniform plateau bound]
\label{thm:lattice-plateau}
There is a numerical constant $C$ such that, for the measures just specified,
$A,B\in\mathfrak A(Q_r)$, and $y\in\mathbb Z^4$ with
$n=|y|_\infty-2r\ge1$,

\refstepcounter{equation}
\begin{gather*}
\FitMath{\left|\nu(\overline A\tau_yB)-\operatorname{Pl}_\nu(A,B)\right|
 \le C\left[(r/n)^2+(r/n)^3\right]
          \mathcal N_{\nu,r}(A)^{1/2}
          \mathcal N_{\nu,r}(B)^{1/2}.}\tag{\theequation}\label{eq:phase-plateau-PL3}
\end{gather*}

The constant is independent of the measure, the support size, and any
coupling. At lattice spacing $a$, the same statement holds with
$Q_r=a[-r,r]^4$ and $n=|y|_\infty/a-2r$.
\end{proposition}

We use the spectral theorem for commuting bounded normal operators and the
mean ergodic theorem for a unitary operator. All estimates supplying the
power in \eqref{eq:phase-plateau-PL3}, including the momentum-density estimate, are proved below.

\subsubsection{Reflection Hilbert spaces and their spectral carriers}

Reflection positivity and translation invariance first imply invariance under
each of the indicated reflections. On one closed side, positivity makes
$\langle F,G\rangle_\theta=\nu(\overline{\theta F}G)$ Hermitian.
Testing it against $1$ shows $\nu(F)=\nu(\theta F)$ for functions supported
on that side. Any local function can be translated into the side by a
lattice vector normal to the mirror. Reflecting reverses that translation;
translation invariance therefore proves the equality for the original
function. The reflections consequently act unitarily on $L^2(\nu)$.

In the coordinate picture $\mu$, complete the quotient of local functions on
$x_\mu\ge0$ for the positive semidefinite form

\begin{gather*}
\FitMath{\langle F,G\rangle_\mu=\nu(\overline{\theta_\mu F}G).}
\end{gather*}

Denote the class of $F$ by $[F]_\mu$. Its norm is at most
$\|F\|_\infty$. Translation by $e_\mu$ induces a positive self-adjoint
contraction $T_\mu$, and tangential coordinate translations induce commuting
unitaries $U_{\mu,\lambda}$, $\lambda\ne\mu$.

Here are the details of this elementary reconstruction. The normal translation
is symmetric on the pre-Hilbert space by translation invariance and
$\theta_\mu\tau_{e_\mu}=\tau_{-e_\mu}\theta_\mu$. Its quadratic form is
nonnegative by reflection positivity at the mirror $x_\mu=-1/2$.
Repeated Cauchy--Schwarz gives

\begin{gather*}
\FitMath{\|T_\mu[F]\|\le
       \|[F]\|^{1-2^{-k}}\|T_\mu^{2^k}[F]\|^{2^{-k}}.}
\end{gather*}

The last norm is at most $\|F\|_\infty$. Letting $k\to\infty$ proves
contraction and preservation of null vectors. The bounded symmetric extension
is self-adjoint and positive. Tangential translations preserve the form and
have inverses of the same kind, so are unitaries. Their commutations on local
functions pass to the completion. In particular, for integers $s\ge0$ and tangential
$w\in\mathbb Z^4$,

\refstepcounter{equation}
\begin{gather*}
\FitMath{\langle[F]_\mu,T_\mu^s U_\mu^w[G]_\mu\rangle
       =\nu(\overline{\theta_\mu F}\tau_{se_\mu+w}G).}\tag{\theequation}\label{eq:phase-plateau-PL4}
\end{gather*}


There is a joint spectral measure with variables
$(E,\eta)\in[0,\infty]\times(-\pi,\pi]^3$ for which
$T_\mu=e^{-E}$ and $U_{\mu,\lambda}=e^{i\eta_\lambda}$.
The value $E=\infty$ represents a possible kernel of $T_\mu$.
The following carrier bound is essential:

\refstepcounter{equation}
\begin{gather*}
\FitMath{\cosh E+\cos\eta_\lambda\ge2} \\[3pt]
\FitMath{(\lambda\ne\mu),} \\[3pt]
\FitMath{E\ge\kappa|\eta|_\infty,} \\[3pt]
\FitMath{\kappa=\frac2\pi\log(1+\sqrt2).}\tag{\theequation}\label{eq:phase-plateau-PL5}
\end{gather*}

We now derive it from the diagonal mirrors.

\paragraph{A polynomial-measure lemma.}
Suppose a finite positive measure $m$ on the closed unit disc satisfies

\refstepcounter{equation}
\begin{gather*}
\FitMath{\int|p(z)|^2\,dm(z)
   \le C_0\sup_{[0,1]}|p|\sup_{|z|=1}|p|}\tag{\theequation}\label{eq:phase-plateau-PL6}
\end{gather*}

for every polynomial. Then it gives zero mass to
$\{|1+z|^2<4|z|\}$.

To prove this without invoking an extremal-function theorem, set

\begin{gather*}
\FitMath{L(\zeta)=\operatorname{Log}\frac{1+\zeta}{1-\zeta},} \\[3pt]
\FitMath{f_s(z)=\frac{\cos(sL(\sqrt z))}{\cosh(\pi s/2)},} \\[3pt]
\FitMath{s>0.}
\end{gather*}

The Cayley transform has positive real part in the disc. Thus $L$ is
holomorphic, odd, and satisfies $|\operatorname{Im}L|<\pi/2$.
The cosine is even in $\sqrt z$, so $f_s$ is a single-valued holomorphic
function of $z$ in the disc. Put
$\varphi(z)=|\operatorname{Im}L(\sqrt z)|$.
The identity $|\cos(u+iv)|^2=\cos^2u+\sinh^2v$ gives

\begin{gather*}
\FitMath{|f_s|\le1,} \\[3pt]
\FitMath{\sup_{[0,1)}|f_s|\le\cosh(\pi s/2)^{-1},} \\[3pt]
\FitMath{|f_s(z)|\ge\frac{\sinh(s\varphi(z))}{\cosh(\pi s/2)}.}
\end{gather*}

Apply \eqref{eq:phase-plateau-PL6} to Taylor approximants of $f_s(qz)$, $0<q<1$,
which converge uniformly on the closed disc. The resulting inequality,
followed by Fatou as $q\uparrow1$, is

\begin{gather*}
\FitMath{\int\sinh^2(s\widetilde\varphi(z))\,dm(z)
       \le C_0\cosh(\pi s/2).}
\end{gather*}

Here $\widetilde\varphi=\varphi$ inside the disc, equals $\pi/2$ on
the boundary except at $1$, and may be set to zero at $1$.
Indeed the radial limits of $\varphi$ at the other boundary points
equal $\pi/2$. For each $\epsilon>0$ this inequality forces
$m(\{\widetilde\varphi\ge\pi/4+\epsilon\})=0$ by letting
$s\to\infty$. Finally, if $z=\zeta^2$ inside the disc,

\begin{gather*}
\FitMath{\tan\varphi(z)=\frac{2|\operatorname{Im}\zeta|}{1-|\zeta|^2},} \\[3pt]
\FitMath{|1+\zeta^2|^2-4|\zeta|^2
       =(1-|\zeta|^2)^2-4(\operatorname{Im}\zeta)^2.}
\end{gather*}

Thus $\varphi>\pi/4$ is precisely the excluded set. The same
assertion holds on the boundary, where only $z=1$ is not excluded.
This proves the lemma.

Fix $\lambda\ne\mu$ and a local $A$ on $x_\mu\ge0$. Choose an
integer $R\ge1$ enclosing its support in
$0\le x_\mu\le R$, $|x_\lambda|\le R$. Let $\theta'$ be reflection
in $x_\mu+x_\lambda=R$, put
$B=\theta'\theta_\mu A$, $A_0=\tau_{2Re_\mu}A$, and
$d_\pm=e_\mu\pm e_\lambda$. Both $B$ and $A_0$ lie on the positive
side of this diagonal mirror. Reconstruction at that mirror gives a
positive contraction $T_+$ for $d_+$ and a unitary $V_-$ for $d_-$.
The preceding reconstruction proof applies: positivity of $T_+$ uses
the admissible diagonal mirror one height unit behind, and its
contraction uses the bounded local representative. For integers $k,l\ge0$,

\refstepcounter{equation}
\begin{gather*}
\FitMath{\langle[A]_\mu,T_\mu^{2R+k+l}U_{\mu,\lambda}^{k-l}[A]_\mu\rangle
       =\langle[B]',T_+^kV_-^l[A_0]'\rangle.}\tag{\theequation}\label{eq:phase-plateau-PL7}
\end{gather*}

Both sides are the same lattice correlation: their displacement is
$2Re_\mu+kd_++ld_-$. No analytic continuation is used here.

Push the spectral measure of $[A]_\mu$, multiplied by $t^{2R}$,
forward by $z=t\zeta$, where $t=e^{-E}$ and
$\zeta=e^{i\eta_\lambda}$. Expanding a polynomial and its conjugate
in \eqref{eq:phase-plateau-PL7} gives

\begin{gather*}
\FitMath{\int|p|^2\,dm
  =\langle[B]',p(T_+)\overline p(V_-)[A_0]'\rangle,} \\[3pt]
\FitMath{\overline p(z)=\sum_k\overline{p_k}z^k.}
\end{gather*}

The spectral theorem bounds this by the right side of \eqref{eq:phase-plateau-PL6}, with
$C_0=\|[B]'\|\|[A_0]'\|$. The lemma excludes
$|1+t\zeta|^2<4t$. Since $t^{2R}>0$ for $t>0$, the joint spectral
measure of $[A]_\mu$ itself vanishes there. The local classes are
dense, so the operator spectral measure vanishes there. At $t=0$
the equivalent operator inequality is automatically satisfied.
We have proved

\begin{gather*}
\FitMath{(1-T_\mu)^2\ge T_\mu(2-U_{\mu,\lambda}-U_{\mu,\lambda}^*).}
\end{gather*}

For $0<t\le1$ this is equivalent to the first assertion of \eqref{eq:phase-plateau-PL5}.
The curve is $E\ge2\operatorname{arcsinh}|\sin(\eta_\lambda/2)|$.
Concavity of $x\mapsto\operatorname{arcsinh}\sin x$ on $[0,\pi/2]$
puts it above its endpoint chord, giving the constant $\kappa$.
Intersect the three carriers to obtain \eqref{eq:phase-plateau-PL5}.

\subsubsection{The zero-energy plateau is the invariant conditional field}

Write $\Pi_\mu=1_{\{E=0\}}$; spectral dominated convergence gives
$T_\mu^s\to\Pi_\mu$ strongly as $s\to\infty$. By \eqref{eq:phase-plateau-PL5}, every
tangential unitary equals the identity on this zero-energy space.
Let $\mathcal I_\mu$ denote the invariant field for the single
translation $\tau_{e_\mu}$. The mean ergodic theorem on $L^2(\nu)$
and \eqref{eq:phase-plateau-PL4} give

\refstepcounter{equation}
\begin{gather*}
\FitMath{\langle[F]_\mu,\Pi_\mu[G]_\mu\rangle
       =\nu(\overline{\theta_\mu F}E_\nu[G\mid\mathcal I_\mu]).}\tag{\theequation}\label{eq:phase-plateau-PL8}
\end{gather*}


Reflection acts as a unitary involution on
$L^2(\mathcal I_\mu,\nu)$. That involution is positive.
Indeed, approximate an invariant function $J$ in $L^2$ by
bounded local $G_k$ and translate each approximant into
$x_\mu\ge0$. The error is unchanged because $J$ is invariant.
Reflection positivity then passes to the $L^2$ limit and gives
$\langle\theta_\mu J,J\rangle\ge0$. A positive self-adjoint
involution has only the eigenvalue $+1$, so reflection fixes $J$.
Consequently \eqref{eq:phase-plateau-PL8} equals

\begin{gather*}
\FitMath{\nu\bigl(\overline{E_\nu[F\mid\mathcal I_\mu]}
                       E_\nu[G\mid\mathcal I_\mu]\bigr).}
\end{gather*}

The map $\Pi_\mu[F]_\mu\mapsto E_\nu[F\mid\mathcal I_\mu]$
is therefore an isometry onto a dense subspace of
$L^2(\mathcal I_\mu,\nu)$. Density follows by the same shifted
local approximations. It intertwines the tangential translations.
Since those translations are the identity on $\Pi_\mu$, every
$\mathcal I_\mu$-measurable function is invariant under all four
coordinate translations. Hence $\mathcal I_\mu=\mathcal I$ modulo
$\nu$, and

\refstepcounter{equation}
\begin{gather*}
\FitMath{\langle[F]_\mu,\Pi_\mu[G]_\mu\rangle
       =\nu\bigl(\overline{E_\nu[F\mid\mathcal I]}
                           E_\nu[G\mid\mathcal I]\bigr).}\tag{\theequation}\label{eq:phase-plateau-PL9}
\end{gather*}

This also proves that $E_\nu[\theta_\mu F\mid\mathcal I]
=E_\nu[F\mid\mathcal I]$. The possible $T_\mu$ kernel has
energy infinity and does not affect any of these zero-energy identities.

\subsubsection{A positive time filter and a bounded momentum density}

For $A\in\mathfrak A(Q_r)$ define its cube vector
$\Phi_A^\mu=[\tau_{re_\mu}A]_\mu$. Translation invariance shows
$\|\Phi_A^\mu\|^2=\nu(\overline{\theta_\mu A}\tau_{2re_\mu}A)$.
Thus \eqref{eq:phase-plateau-PL1} is the maximum of the corresponding squared norms of
$\gamma A$. In particular $\mathcal N^{1/2}$ is a seminorm,
is invariant under the sign-change group, and is at most the
supremum norm.

Let $b$ be the nonnegative periodic triangle function supported on
$[-\pi/4,\pi/4]$, normalized to have mean one. Its Fourier
coefficients are

\begin{gather*}
\FitMath{b_m=\left(\frac{\sin(m\pi/8)}{m\pi/8}\right)^2,} \\[3pt]
\FitMath{b_0=1.}
\end{gather*}

This follows by writing the triangle as the convolution of two
interval indicators and normalizing their zeroth coefficients.
Let $z_m$ be the even period-four sequence $(0,1,-2,1)$ and set

\refstepcounter{equation}
\begin{gather*}
\FitMath{\widehat q(u)=\sum_{m\in\mathbb Z}z_mb_m e^{imu},} \\[3pt]
\FitMath{\Psi(q)=2\sum_{m\ge1}z_mb_mq^{m-1},} \\[3pt]
\FitMath{0\le q\le1.}\tag{\theequation}\label{eq:phase-plateau-PL10}
\end{gather*}

Both series converge absolutely. The identity
$z_m=\frac12i^m-(-1)^m+\frac12(-i)^m$ gives

\begin{gather*}
\FitMath{\widehat q(u)=\tfrac12b(u+\pi/2)-b(u+\pi)
                               +\tfrac12b(u-\pi/2).}
\end{gather*}

In particular $\widehat q$ vanishes whenever
$\operatorname{dist}(u,2\pi\mathbb Z)\le\pi/4$.
The crucial energy property is

\refstepcounter{equation}
\begin{gather*}
\FitMath{0\le\Psi\le C,} \\[3pt]
\FitMath{\Psi(1)=0,} \\[3pt]
\FitMath{\Psi(e^{-x})\ge c\tanh x} \\[3pt]
\FitMath{(x>0).}\tag{\theequation}\label{eq:phase-plateau-PL11}
\end{gather*}


To verify its nontrivial lower bound, use the Poisson kernel

\begin{gather*}
\FitMath{P_x(u)=\sum_me^{-|m|x}e^{imu}=\frac{\sinh x}{\cosh x-\cos u}.}
\end{gather*}

Since $z_0=0$, the Poisson integral yields

\begin{gather*}
\FitMath{\Psi(e^{-x})=\frac{e^x}{2\pi}\int_{|v|\le\pi/4}b(v)
  [\tfrac12P_x(v-\pi/2)+\tfrac12P_x(v+\pi/2)-P_x(v-\pi)]\,dv.}
\end{gather*}

Writing $C_x=\cosh x$, the bracket is

\begin{gather*}
\FitMath{\sinh x\left[
  \frac{C_x}{C_x^2-\sin^2v}-\frac1{C_x+\cos v}\right]
 \ge\frac{\tanh x}{\sqrt2(C_x+2^{-1/2})}.}
\end{gather*}

For this inequality first replace the first fraction by $1/C_x$,
then use $\cos v\ge2^{-1/2}$ and monotonicity of
$q/(C_x+q)$. The factor $e^x/(C_x+2^{-1/2})$ has a positive
lower bound, and $b$ has mean one. This proves the lower bound.
The upper bound follows from absolute convergence of \eqref{eq:phase-plateau-PL10}, and
$\Psi(1)=\widehat q(0)=0$. These statements extend continuously
to $q=0$, where $\Psi(0)=2b_1>0$.

First assume $\theta_\mu A=\epsilon A$, $\epsilon\in\{1,-1\}$,
and put $\Phi=\Phi_A^\mu$. For $s\in\mathbb Z$ and tangential
$w\in\mathbb Z^3$, define

\begin{gather*}
\FitMath{g_A(s,w)=\epsilon\nu(\overline A\tau_{se_\mu+w}A).}
\end{gather*}

For $s\ge2r$, the cube-crossing identity from \eqref{eq:phase-plateau-PL4} gives

\refstepcounter{equation}
\begin{gather*}
\FitMath{g_A(s,w)=\langle\Phi,T_\mu^{s-2r}U_\mu^w\Phi\rangle,} \\[3pt]
\FitMath{|g_A(s,w)|\le\mathcal N(A).}\tag{\theequation}\label{eq:phase-plateau-PL12}
\end{gather*}

Reflection in $x_\mu=0$ and translation invariance give
$g_A(-s,w)=g_A(s,w)$ and
$g_A(s,-w)=\overline{g_A(s,w)}$. Define

\begin{gather*}
\FitMath{K_A(w)=\sum_{m\in\mathbb Z}z_mb_mg_A(2rm,w).}
\end{gather*}

The term at $m=0$ vanishes. Every other time separation has
absolute value at least $2r$, so \eqref{eq:phase-plateau-PL12} bounds this sum by
$C\mathcal N(A)$, uniformly in $w$. By the same identities,

\refstepcounter{equation}
\begin{gather*}
\FitMath{K_A(w)=\langle\Phi,\Psi(T_\mu^{2r})U_\mu^w\Phi\rangle
       =\int e^{i\eta\cdot w}\,dM_A(\eta),}\tag{\theequation}\label{eq:phase-plateau-PL13}
\end{gather*}

where $M_A$ is the momentum pushforward of the positive measure

\begin{gather*}
\FitMath{\Psi(e^{-2rE})\,d\langle\Phi,\mathsf E_\mu(E,\eta)\Phi\rangle.}
\end{gather*}

Thus $M_A$ is positive; the signs in $z_m$ do not compromise
positivity because \eqref{eq:phase-plateau-PL11} controls their spectral sum.

For $|w|_\infty\ge2r$, choose $\lambda\ne\mu$ at a maximal
coordinate. Replace $w$ by $-w$ if necessary, which only conjugates
$K_A(w)$, and suppose $w_\lambda\ge2r$. In the coordinate picture
$\lambda$, for every integer $s$ the same correlation is

\begin{gather*}
\FitMath{g_A(s,w)=\epsilon\left\langle
 \Phi_{\theta_\lambda A}^\lambda,
 T_\lambda^{w_\lambda-2r}U_{\lambda,\mu}^s
       \prod_{\lambda'\ne\lambda,\mu}U_{\lambda,\lambda'}^{w_{\lambda'}}
 \Phi_A^\lambda\right\rangle.}
\end{gather*}

Sum this identity at $s=2rm$ against $z_mb_m$. In this picture
the filter becomes the bounded spectral function
$\widehat q(2r\eta'_\mu)$ of a unitary translation. Whenever it
is nonzero, $|\eta'_\mu|>\pi/(8r)$, since the gap around zero
modulo $2\pi$ includes the ordinary interval about zero. Aliases
at other multiples of $2\pi$ do not invalidate this implication.
The carrier \eqref{eq:phase-plateau-PL5} consequently gives
$E'\ge\alpha=\kappa\pi/(8r)$ on the filter support. Both cube
vectors have norm at most $\mathcal N(A)^{1/2}$. Hence

\refstepcounter{equation}
\begin{gather*}
\FitMath{|K_A(w)|\le C\mathcal N(A)
       e^{-\alpha(|w|_\infty-2r)}} \\[3pt]
\FitMath{(|w|_\infty\ge2r).}\tag{\theequation}\label{eq:phase-plateau-PL14}
\end{gather*}


There are $O(r^3)$ points with $|w|_\infty<2r$ and $O(k^2)$
points on the shell $|w|_\infty=k$. Since $\alpha$ is a fixed
constant times $r^{-1}$, \eqref{eq:phase-plateau-PL14} and the previous bound give

\begin{gather*}
\FitMath{\sum_w|K_A(w)|
 \le C\mathcal N(A)
       \left[r^3+\sum_{j\ge0}(r^2+j^2)e^{-\alpha j}\right]
 \le Cr^3\mathcal N(A).}
\end{gather*}

Therefore $h_A(\eta)=(2\pi)^{-3}\sum_wK_A(w)e^{-i\eta\cdot w}$
is continuous and bounded by $Cr^3\mathcal N(A)$.
It has the same Fourier coefficients as $M_A$. Uniqueness of a
finite measure from its Fourier coefficients follows by density
of trigonometric polynomials in the continuous functions on the
torus. Consequently

\refstepcounter{equation}
\begin{gather*}
\FitMath{M_A=h_A(\eta)\,d\eta,} \\[3pt]
\FitMath{0\le h_A\le Cr^3\mathcal N(A).}\tag{\theequation}\label{eq:phase-plateau-PL15}
\end{gather*}

This is a derived bound for an energy-weighted momentum measure.
It is not an assumed bound for the unweighted spectral measure.

\subsubsection{Removing the filter and proving the plateau estimate}

Let $m_\Phi=\langle\Phi,\mathsf E_\mu(\cdot)\Phi\rangle$.
Since $M_A$ has a density, it has no atom at $\eta=0$.
By \eqref{eq:phase-plateau-PL11}, its weight is strictly positive at every finite
$E>0$. Thus $m_\Phi$ gives zero mass to
$\{0<E<\infty,\eta=0\}$. The energy-zero contribution is precisely
the plateau and was not removed by this inference: its weight is zero.
The point $E=\infty$ contributes zero to $T_\mu^n$ for $n\ge1$.

For $E>0$ the function $e^{-nE}/\tanh(2rE)$ decreases in $E$.
Using \eqref{eq:phase-plateau-PL11}, then \eqref{eq:phase-plateau-PL5}, and then \eqref{eq:phase-plateau-PL15}, gives
\refstepcounter{equation}
\begin{align*}
 \langle\Phi,T_\mu^n(1-\Pi_\mu)\Phi\rangle
 &\le C\int_{\mathbb T^3}
       \frac{e^{-n\kappa|\eta|_\infty}}
            {\tanh(2r\kappa|\eta|_\infty)}h_A(\eta)\,d\eta\\
 &\le Cr^3\mathcal N(A)\int_0^\infty
       u^2e^{-n\kappa u}\left(1+\frac1{2r\kappa u}\right)du\\
 &\le C\mathcal N(A)\bigl[(r/n)^3+(r/n)^2\bigr].
 \tag{\theequation}\label{eq:phase-plateau-PL16}
\end{align*}
In the second line we used $1/\tanh v\le1+1/v$ and the
three-dimensional cube shell volume $24u^2du$. The apparent singularity
at $u=0$ is integrable. The pushforward step can include $E=\infty$
and other unused spectral mass, since the integrand is nonnegative;
it does not add an atom at $\eta=0$ by \eqref{eq:phase-plateau-PL15}.

For an arbitrary complex $A$, write
$A=A_++A_-$ with $A_\pm=(A\pm\theta_\mu A)/2$.
The seminorm property and sign-change invariance imply
$\mathcal N(A_\pm)\le\mathcal N(A)$.
The operator $D_n=T_\mu^n(1-\Pi_\mu)$ is positive, so its
quadratic form at a sum is at most twice the sum of its two
quadratic forms. Applying \eqref{eq:phase-plateau-PL16} to $A_+$ and $A_-$ proves
\eqref{eq:phase-plateau-PL16}, with an enlarged numerical constant, for every $A$.

Finally suppose $y_\mu=|y|_\infty>0$. The cube-crossing identity is

\begin{gather*}
\FitMath{\nu(\overline A\tau_yB)=
   \langle\Phi_{\theta_\mu A}^\mu,
        T_\mu^n U_\mu^{y_\perp}\Phi_B^\mu\rangle.}
\end{gather*}

The $\Pi_\mu$ term is \eqref{eq:phase-plateau-PL2} by \eqref{eq:phase-plateau-PL9}, conditional invariance under
translations and reflection, and $U_\mu^{y_\perp}\Pi_\mu=\Pi_\mu$.
For the remaining term, Cauchy--Schwarz in the positive form of
$D_n$ bounds the modulus by the geometric mean of the two quantities
in \eqref{eq:phase-plateau-PL16}. The tangential unitary commutes with $D_n$, so does not
change that form. Also
$\mathcal N(\theta_\mu A)=\mathcal N(A)$. This proves \eqref{eq:phase-plateau-PL3}.
For a negative maximal coordinate use
$\nu(\overline A\tau_{-y}B)
=\overline{\nu(\overline B\tau_yA)}$ and conjugate symmetry of
the plateau. Dilation restores spacing $a$ and replaces distances
by their numbers of lattice steps, proving the stated scaling.
\hfill$\square$

For an ergodic state the invariant conditional expectations are constants,
so \eqref{eq:phase-plateau-PL3} is clustering. For a mixture its generally nonzero invariant plateau
is retained. When local continuum representatives have large supremum norms,
the reflection sizes in \eqref{eq:phase-plateau-PL1} can still be bounded by separated reflected
moments. This is the reason \eqref{eq:phase-plateau-PL3}, rather than a supremum-norm substitute,
is the useful estimate for a continuum phase construction. No spectral gap
above the invariant subspace has been proved: the filter weight vanishes
linearly near zero energy, and removing it is exactly what produces the
inverse-square term.

\begin{corollary}[Closedness of the ergodic reflection-positive class]
At fixed spacing, the translation-invariant measures satisfying these
mirror-positivity hypotheses form a weakly closed set, and so does its
subset of ergodic measures.
\end{corollary}
\begin{proof}
Translation invariance and each positivity inequality on continuous
local functions pass to a weak limit. Reflection invariance passes
as well, and approximation in $L^2$ extends positivity to bounded
Borel local functions. Thus the full reflection-positive class is
closed.

Suppose now that $\nu_j$ are ergodic in this class and converge weakly
to $\nu$. For fixed continuous local $A,B$ in a common cube,
their plateaux under $\nu_j$ are
$\overline{\nu_j(A)}\nu_j(B)$. Passing to the limit in \eqref{eq:phase-plateau-PL3} gives

\begin{gather*}
\FitMath{|\nu(\overline A\tau_yB)-\overline{\nu(A)}\nu(B)|
 \le C[(r/n)^2+(r/n)^3]\|A\|_\infty\|B\|_\infty.}
\end{gather*}

In particular the correlation converges to the product along
$y=ke_\mu$, $k\to\infty$. Its Cesaro means, with $A=B$, identify
the squared $L^2$ norm of $E_\nu[A\mid\mathcal I_\mu]$ by the mean
ergodic theorem. That norm is $|\nu(A)|^2$, so the conditional
expectation is constant. Continuous local functions are dense in
$L^2(\nu)$; hence $\mathcal I_\mu$, and therefore $\mathcal I$, is
trivial. This is ergodicity of $\nu$.
\end{proof}


% END MODULE chapters/phase_plateau_audit_v3.tex

% BEGIN MODULE chapters/comparison_capsule_v7.tex
\subsection{Coupling coordinates and conditional comparisons}
\label{ch:comparison-capsule-v7}
\label{candidate:power_counting}
\label{candidate:uniqueness}
\label{candidate:rotations}

The coupling, matched-run and Ward estimates below imply their stated
conclusions by elementary analytic arguments. The interacting estimates
themselves remain hypotheses. The
pointed families, companions, physical Hessian extractor, midpoint tents,
list action and radius weights are defined in \cref{ch:rg-objects-v2}.
The pointed cancellation, nonlinear window estimate and finite formation
sections have their stated background and domain hypotheses. None computes
the one-loop trace or proves a matched-cutoff or typed Ward estimate.

\subsubsection*{Vacuum accounting and the stronger coupling inputs}

Write \(P_\circ=(\mathbf1,0)\). With \(E_{(0)}=0\), the accumulated
vacuum normalization is defined by
\begin{equation}
\begin{aligned}
 E_{(k+1)}&=E_{(k)}+\mathsf e_k,\qquad E=E_{(K)},\\
 \mathsf e_k&=\mathfrak n_k+
 \sum_X\left[\sum_{z\in X}E^{(k+1)}(X,P_\circ,z)
                         +R^{(k+1)}(X,P_\circ)\right].
\end{aligned}\label{eq:comparison-capsule-CP1}
\end{equation}
Here \(\mathfrak n_k\) is the assigned field-independent step number.
For \(\tau=x/\theta\), its Gaussian scaling part is
\(-3|\mathfrak X_k|\log\tau/2\); \(\mathfrak X_k\) is the set of
non-pivot bonds, with three real color coordinates per bond. Haar and
gauge-fixing counts are included in that assigned number.
The family in \eqref{eq:comparison-capsule-CP1} contains the unnormalized local Gaussian numbers:
\(\mathbf N^\circ=\mathbf N+\mathbf N^\circ(P_\circ)\).
Those numbers must not also be added to \(\mathfrak n_k\).
Equivalently, using \(\mathbf N\) requires adding the removed local
numbers once separately. This convention counts every constant once.
At fixed calibration, \(\mathfrak n_k,\mathbf N^\circ,\mathbf G_\tau\)
have zero source derivative. Source energies are the differences of
the displayed \(E\)- and \(R\)-values between sourced and unsourced runs.
More explicitly, the localized source energy of level \(j\) is
\[
 e_j(X;w)=\sum_{z\in X}\{E^{(j),w}(X,P_\circ,z)-E^{(j),0}(X,P_\circ,z)\}
          +R^{(j),w}(X,P_\circ)-R^{(j),0}(X,P_\circ).
\]
The source vacuum exponent is the sum of these values over levels and
local domains. This definition introduces no additional arbitrary
normalizer. The zero source derivative of the Gaussian pieces refers
to the frozen potential presentation. Moving the source into the
Gaussian action weight changes the determinant and is another
presentation, whose numbers must be compared by its own sum rule.

Requirement B3 retains an actual one-loop input:
\begin{equation}
 b_i=b_i^0+f_i+q_i,\qquad
 c_0=\frac{11N^2}{12\pi^2}\log L,\qquad
 \left|\sum_{i=1}^r(b_i^0-c_0)\right|\le C_0 .
 \label{eq:comparison-capsule-CP2}
\end{equation}
Here \(N=2\), and \(b_i^0\) is independent of the bare coupling and
admissible torus. A summable estimate
\(|b_i^0-c_0|\le CL^{-\sigma i}\), \(\sigma>0\), suffices.
The asserted normalization requires the actual lattice determinant/trace
calculation and its zero-mode convention. Gaussian integration alone
does not establish it.

Some later slope conclusions require more than \eqref{eq:comparison-capsule-CP2}. Fix a calibration
\(\theta_0,\ldots,\theta_K\), independent of the varied bare coupling.
Put \(\xi_j=\theta_j-x_j\). Require the actual unmarked \(f_i\) and
marked \(q_i\), for \(1\le i\le K\), to be jointly holomorphic on the
independent earlier-variable polydiscs with radii
\begin{equation}
\begin{aligned}
 r_j^{\rm u}&=c\theta_j,&
 \sup|f_i|&\le C\theta_i^{-1/2}(\log\theta_i)^p,\\
 r_j^{\rm m}&=c\theta_j(\log\theta_j)^{-C_*},&
 \sup|q_i|&\le C\theta_i^{-p_R},
 \qquad 0\le j<i .
\end{aligned}\label{eq:comparison-capsule-CP3}
\end{equation}
The constants \(p,C_*\ge0\), \(p_R>0\), and \(c>0\) are fixed.
The actual runs must remain in the half polydiscs, with common
localization, counterterms and admissible domains.
These are stronger actual coupling inputs, including the marked tube.
They are not inferred from holomorphy of the bare partition function,
nor from a running-radius activity bound followed by two field
derivatives. The marked exponent is often \(p_R=(\kappa_0-4)/2\),
but any later use must retain the actual proved loss and tube radius.
All norms and source Taylor expansions here use fixed declared domains;
shrinking the field domains while differentiating would be a different
claim. Compatibility of the companion and earlier cutoff prefix
also remains part of the construction.

\subsubsection*{A dimension-free polydisc lemma and its precise consequence}

\begin{lemma}[The bounded-polydisc estimate, code \(1e.3.7\)]
\label{lem:polydisc-dimension-free-v7}
If \(F\) is holomorphic on \(\prod_{j=1}^dD(0,r_j)\) and
\(|F|\le B\), then
\begin{equation}
 \sum_jr_j\left(1-\frac{|\xi_j|^2}{r_j^2}\right)
                   |\partial_jF(\xi)|\le B .
 \label{eq:comparison-capsule-CP4}
\end{equation}
On the half polydisc this implies
\begin{equation}
\begin{aligned}
 \left|\sum_jv_j\partial_jF(\xi)\right|
       &\le\tfrac43 B\max_j|v_j|/r_j,\\
 |F(\xi)-F(\xi')|&\le\tfrac43B
                         \max_j|\xi_j-\xi'_j|/r_j .
\end{aligned}\label{eq:comparison-capsule-CP5}
\end{equation}
No constant depends on \(d\).
\end{lemma}
\begin{proof}
The case \(B=0\) is immediate. Set \(a_j=\xi_j/r_j\) and
\(\phi_a(z)=(a+z)/(1+\bar a z)\). Choose unit complex numbers
\(\epsilon_j\) so the terms in the derivative at zero of
\[
 H(z)=F\big(r_j\phi_{a_j}(\epsilon_jz)\big)_{j=1}^d
\]
have the same argument. Since \(\phi_a'(0)=1-|a|^2\), the modulus
of that derivative is the left side of \eqref{eq:comparison-capsule-CP4}.
The one-variable Schwarz lemma, applied after a disc automorphism
sending \(H(0)/B\) to zero, gives
\(|H'(0)|\le B(1-|H(0)/B|^2)\le B\).
The Schwarz lemma itself follows by applying the maximum principle
to a holomorphic function vanishing at zero divided by \(z\), first
on smaller circles and then letting their radius increase to one.
An attained boundary value makes \(H\) constant.
On the half polydisc, \eqref{eq:comparison-capsule-CP4} gives
\(\sum_jr_j|\partial_jF|\le4B/3\).
This proves the first estimate of \eqref{eq:comparison-capsule-CP5}; integration along the segment
from \(\xi\) to \(\xi'\), which stays in that half polydisc, proves
the second.
\end{proof}

Assume \(\theta_j\ge\theta_i\) for \(j<i\), with all centers sufficiently
large that \(t/(\log t)^{C_*}\) is increasing.
Applying the lemma separately in \eqref{eq:comparison-capsule-CP3} bounds the absolute row sum by
\begin{equation}
 a_i:=\sum_{j<i}|\partial_{\xi_j}(f_i+q_i)|
 \le C\left[
 \theta_i^{-3/2}(\log\theta_i)^p+
 \theta_i^{-p_R-1}(\log\theta_i)^{C_*}\right].
 \label{eq:comparison-capsule-CP6}
\end{equation}
Suppose also that the calibration has backward linear drift
\(\theta_i\ge h+c_1(K-i)\), \(c_1>0\).
Comparison with integrals gives
\[
 A:=\sum_{i=1}^Ka_i
 \le C\left[h^{-1/2}(\log h)^p+
                       h^{-p_R}(\log h)^{C_*}\right]\longrightarrow0.
\]
Indeed substitute \(t=hu\) in each integral; the integrals of
\(u^{-1-\delta}(1+|\log u|)^r\) are finite for \(\delta>0\).
A drift with a fixed additive error is equivalent after increasing \(h\).

For the actual differentiable recursion
\(x_0(\beta)=\beta\),
\(x_i(\beta)=\beta-\sum_{r\le i}b_r(\beta)\),
put \(M_i=\max_{j\le i}|x_j'|\). Calibration is held fixed, so
\(|b_i'|\le a_iM_{i-1}\) and
\(M_i\le(1+a_i)M_{i-1}\). Consequently
\begin{equation}
 M_K\le\prod_i(1+a_i)\le e^A,\qquad
 \sum_i|b_i'|\le\prod_i(1+a_i)-1\le e^A-1 .
 \label{eq:comparison-capsule-CP7}
\end{equation}
The middle estimate telescopes the successive products.
In particular
\[
 |x_K'(\beta)-1|\le\sum_i|b_i'|
 \le e^A-1
 =O\!\left(h^{-1/2}(\log h)^p+h^{-p_R}(\log h)^{C_*}\right).
\]
On witness prefixes with terminal \(h\to\infty\), under the same
uniformly admitted hypotheses, this makes the constant-source
normalization tend to one with the displayed error. Identifying
that prefix and derivative with the actual F3 witness remains
an input; the estimate does not establish that identification.
At weak coupling \(A\le\log(4/3)\), hence the last sum is at most
\(1/3\). Constant sources varying \(\beta\) to \(\beta-\zeta\)
have the same absolute slopes. This supplies the slope premise in
\eqref{eq:local-formation-LB12} under \eqref{eq:comparison-capsule-CP3} and the admitted drift/domain conditions.
It does not prove those conditions.
Reality gives \(2/3\le x_K'(\beta)\le4/3\) on an admitted real
interval. Integrating the derivative proves the uniform terminal
bi-Lipschitz bound. If the same row estimates hold on a whole complex
offset disc, write \(\Xi_K(t)=h-t+R_K(t)\); then
\(R_K(0)=0\) and \(\sup|R_K'|\le1/3\).
One can construct the matching inverse without any global uniqueness
assertion. For \( |t|<\varrho\), solve
\[
 s=h-\Xi_K(t)+R_{K+1}(s)
\]
on the closed disc \( |s|\le2\varrho\), assuming the functions are
defined on a neighborhood of that disc. Its right side maps the disc
strictly inside itself and has Lipschitz constant at most \(1/3\).
Iteration converges geometrically to its unique solution, uniformly on
smaller \(t\)-discs. The iterates are holomorphic, so the solution
is holomorphic; reality gives a real solution at real \(t\).
The bounds \( |\Xi_K(t)-h|\le4|t|/3\) and
\( |\Xi_{K+1}(s)-h|\ge2|s|/3\) give
\( |\mathfrak n_K(t)|\le2|t|\).
This supplies the coordinate part of E2 when its actual full-disc
premises are available, rather than proving the existence of those
premises. Nor does a polydisc supremum
give forgetting in \(i-j\): \(F_i(\xi)=B\xi_0/r_0\) is bounded
by \(B\) for every \(i\), with unchanged first derivative.

\subsubsection*{The fixed-torus coordinate and the matched limit}

Fix a torus \(T_\ell\), \(h=g^{-2}\), and a separated tuple of
\(n\ge2\) compact Lipschitz tests, with the fixed norm
\(\|\cdot\|_\sharp\) of the source convention. Every bound below includes their product
norm; its constant may depend on order and positive separation.
Let \(\beta_K^\natural\) be its torus-independent tuned bare center
and write depth as \([s]=K-k\). Define
\begin{equation}
\begin{aligned}
 \mathsf J&=[-\varrho/4,\varrho/4],&
 \Xi_K(t)&=x_K^\natural(\beta_K^\natural-t),\\
 F_{K,n}(t;\mathbf f)
   &=S^T_{n;K,\ell,\beta_K^\natural-t}(\mathbf f),&
 \mathfrak n_K&=\Xi_{K+1}^{-1}\circ\Xi_K .
\end{aligned}\label{eq:comparison-capsule-CP8}
\end{equation}
The inverse is the specified holomorphic matching branch.
Inputs E2--E3 require its existence on a common neighborhood, reality
on the real axis, uniform upper and lower real Lipschitz bounds for
\(\Xi_K\), \(\Xi_K(0)=h\), and
\(|\Xi_K(t)-h+t|\le|t|/3\).
On a fixed smaller complex neighborhood of \(\mathsf J\), require
\begin{equation}
\begin{aligned}
 |\mathfrak n_K-t|+|\mathfrak n_K'-1|
       +|\Xi_{K+1}-\Xi_K|&\le C\eta_K,\\
 |F_{K,n}(t)-\mathfrak n_K'(t)^n
                    F_{K+1,n}(\mathfrak n_K(t))|
       &\le C_{n,\ell,\mathbf f}\epsilon_K,\qquad
 \sum_K(\eta_K+\epsilon_K)<\infty .
\end{aligned}\label{eq:comparison-capsule-CP9}
\end{equation}
The one-run estimates below supply the uniform holomorphic
bound for \(F_{K,n}\) on a slightly wider strip, for example
\(\mathcal N_\ell=\{|\operatorname{Re}t|<\varrho/2,\
|\operatorname{Im}t|<\upsilon_\ell\}\), where
\(\upsilon_\ell=c\varrho e^{-C(g)\ell^4}>0\).
Matching paths stay in that strip after shrinking it.
Constants and cutoff floors may depend on this fixed \(\ell\).
Often \(\eta_K=\Upsilon_K\le
CK^{-3/2}(\log(2+K))^{p+1}\); a separate remainder error
\(\Pi_K\) is part of \(\epsilon_K\), not a polarization tensor.

\subsubsection*{The scalar coordinate and the local energies}

The two estimates in \eqref{eq:comparison-capsule-CP9} use different
comparisons. For the coordinate, let $b^A_{[u]}(t)$ and $b^B_{[u]}(t)$
be the scalar increments at the same depth for the two tuned runs.
Compare them at the matched constant offsets $t$ and
$\mathfrak n_K(t)$, using the inverse branch already constructed.
They include their marked remainders. We must show on a common
offset disc, with fixed room about $|t|\le\varrho$, that
\begin{equation}
 \sup_t|b^B_{[u]}(\mathfrak n_K(t))-b^A_{[u]}(t)|
       \le C(\rho^{K-u}+\Pi_K),\qquad 0\le u<K .
 \label{eq:comparison-matched-scalar-step}
\end{equation}
Here $0<\rho<1$ and $\Pi_K$ is the error from the unmatched
large-field terms at the bare end. The derivative of each
increment, including the extra increment $b^B_{[K]}$, has the
one-run bound $Ca_{[u]}$, with
$a_{[u]}\le C\theta_{[u]}^{-3/2}(\log\theta_{[u]})^p$.
This uses the sensitivity estimate on the actual full disc as
in \eqref{eq:comparison-capsule-CP7}. The derivative loses one
power of $\theta_{[u]}$ from the holomorphy radius in the
bound for the part beyond one loop.

\paragraph{Scalar increments of partnered runs.}
Use the decomposition \eqref{eq:flow-RF20} into the one-loop
increment, the unmarked analytic coefficient $f$, and the
marked correction $q$. The one-loop difference has a geometric
bound at equal depth by \eqref{agc:actual-increment-law} at
equal isotropic weights. The unmarked coefficient at equal
deviations uses the transfer of the regular chain, with the
fixed physical extraction radius. Changing deviations within
one cutoff uses the weighted forgetting estimate on the
admitted product of offset discs. At the step where a changed
deviation enters, Cauchy's estimate divides the one-step
bound by a radius of order the inverse coupling. Subsequent
steps propagate the resulting difference by the regular age
weights. This gives the single-run forcing coefficient of order
$\theta_{[s+1]}^{-3/2}(\log\theta_{[s+1]})^p$ at depth $s$,
times a geometrically decreasing kernel in the depth lag.

For the marked correction, first change the deviations of
the longer run to those of the shorter run on the shared
depths, keeping its extra deviation and its own remainders.
This first difference uses the one-cutoff marked estimate
with fixed remainder data. Next compare the two cutoffs at
equal deviations. The latter difference is the difference
of the two marked coefficients minus the difference of the
two unmarked coefficients, with all four readouts on their
common running radii. The gate for inserting an old remainder
is the one of \eqref{eq:flow-RF19}. It admits the remainder
to the regular exponent at the age with the prescribed lattice
symmetries. Its coefficient estimate thereafter uses the
squared radius ratio. If $\chi_t$ is the distance of the two
remainder families of depth $t$, the contribution on reading
the coefficient at depth $s<t$ is at most
$C\theta_{[t+1]}\chi_t$. Indeed, the extraction costs the
inverse square of the running radius at the output, which
cancels that radius squared in the ratio. There is no
geometric age factor in this bound.

The regular comparison of the four chains bounds the
equal-deviation coefficient difference by
\begin{equation}
 C\left[\theta_{[s+1]}(1+K-s)\lambda^{K-s}
          +\Pi_K+\sum_{s<t<K}\theta_{[t+1]}\chi_t\right]
 \label{eq:comparison-marked-four-chains}
\end{equation}
for some $\lambda<1$. The first term pays for the running
extraction radius and the geometric mesh defects of the
partners. The bare-end remainder terms in $\Pi_K$ are
estimated by their sizes. The one-cutoff bound for each
marked correction also bounds this difference by
$C\theta_{[s+1]}^{-\nu}$, with $\nu=(\kappa_0-4)/2$.
For $b=(1+K-s)\lambda^{K-s}$,
$\min\{b\theta,\theta^{-\nu}\}\le b^{\nu/(\nu+1)}$.
Retain the last two terms of \eqref{eq:comparison-marked-four-chains}
and apply this inequality to the first. A slightly slower
geometric rate absorbs the polynomial in $K-s$.

We must bound the weighted sum of remainder distances for
the actual matched histories. Where $\theta_{[K]}\le
2\theta_{[s+1]}$, the two individual size bounds already
give $C\Pi_K$ for this sum. Away from that range the comparison
uses the two parts of the new remainder. The ordinary step
puts in one part from clusters that contain a recent remainder.
%
On the common admitted carriers let \(D_j\) be the
norm of the difference of the recent remainder inputs
of origin \(j\) before forming the localized vertices,
and \(A_j\) the larger of their two norms.
Let \(\eta\) be the size of the change of the other inputs
in their admitted units. For the ordinary contribution at
level \(t\), the difference of the cluster sums has bound
\[
 C\mathsf s_{t-1}
   \sum_{j:\,M\check R_j>L^{t-j}}
                 w_{t-1,j}\bigl(D_j+\eta A_j\bigr).
\]
Indeed, interpolate the recent inputs on the common
domain. Differentiating the logarithm expansion marks
one varied recent vertex. Its bound is
\(C\mathsf s_{t-1}w_{t-1,j}D_j\),
and the rooted sum over the other vertices is uniform
in that interpolation. At fixed recent inputs, every
cluster in the selected sum has a recent vertex,
including at complex values of the other inputs.
Apply the bound with that vertex marked on the
complex neighborhood of the path changing the other
inputs. Cauchy's estimate on discs of radius
proportional to \(\eta^{-1}\) along this path gives
the factor \(\eta\) without losing the recent
vertex size. The supports and age gate must be
common, with the recent-vertex formation bound
valid along the enlarged complex path for the
other inputs, including the operators. This path
uses a fixed reserve in the admitted norms.
The other part comes
from integrating histories in the marked operation. On a single
lattice, the weighted extension estimates that part by the rare
mass times the sum of the loads at level $n$, each with weight
$(k-n+1)q_a^{(k-n-1)_+}$ at level $k$. The exponent next to a
live component can read its retained layer repeatedly. Before
integration the factor $(k-n+1)$ counts these uses. A bounded
price per use is paid by strengthening the survival reserve
under the eventual integral. Thus the geometric age weight
for that part belongs to the integrated history. The direct
recent-input coefficient has no rare-event factor. Its window has length
$m_s=O(\log_L(M\check R_s))$. In the units of the remainder
size at depth $s$, the ratios of the sizes over this window
are bounded. On the profile $\mathsf S$ of
\eqref{eq:comparison-matched-profile}, the window can cost
$C m_s(1+m_s)\rho^{-m_s}$. This is a power of $\log\theta_{[s]}$
and is absorbed by the logarithmically small coefficient.
The integrated part requires the full comparison of weighted
histories, including the shell and availability estimates below.

%
%
Here is the summation required of the local
comparison. Put \(a_t=\theta_{[t+1]}\chi_t\) and
\(q_t=C\theta_{[t]}^{-\nu-1}(\log\theta_{[t]})^c\).
A sufficient estimate after separating the direct
recent-input contribution is
\[
 \begin{split}
 a_t\le{}&
 q_t\left\{\lambda^{K-t}+(1+\varrho^{-1})
 \left[\sum_{u=t+1}^K\lambda^{u-t-1}y_u
       +\mathsf h+\sum_{j=t+1}^{K-1}a_j\right]\right\}\\
 &+\sum_{j=t+1}^{K-1}P_{tj}a_j+P_{tK}\mathsf h ,
 \end{split}
\]
where \(\mathsf h\le C\theta_{[K]}^{-\nu-1}\)
is the extra bare-end input. The \(P_{tj}\) are
the coefficients of the ordinary recent
contribution in the weighted units. For fixed
\(j\), they vanish unless \(j>t\) and the
origin \(j\) is still recent at \(t\).
There are \(O(1+\log\log\theta_{[j]})\)
such depths. In this window
\(\theta_{[t+1]}/\theta_{[j+1]}\) is bounded,
and the radius loss is at most \(C(1+j-t)\).
The column sum is therefore bounded by
\(C\mathsf s^{\rm run}(\theta_{[j]})
 (1+\log\log\theta_{[j]})^2\).
The coupling ceiling makes its supremum
small. This charge has no rare-mass factor.

Sum over \(r<t<K\). The triangular term
with \(q_t a_j\) has total coefficient at most
\((1+\varrho^{-1})\sum_tq_t\), which is also
small at the coupling ceiling. Absorb these
two column bounds on the left. The extra
input contributes \(C\mathsf h\). For the
coefficient of \(y_u\), split the sum in \(t\)
at \((u+r)/2\). On the earlier portion the
geometric factor gives
\(C\theta_{[r+1]}^{-\nu}\lambda_H^{u-r-1}\)
for a fixed \(\lambda_H<1\) larger than \(\sqrt\lambda\).
On the later portion \(\theta_{[t]}\asymp\theta_{[u]}\)
and the geometric series gives
\(C\theta_{[u]}^{-\nu-1}(\log\theta_{[u]})^c\).
The term \(\sum_tq_t\lambda^{K-t}\) has the
same split with \(u=K\). These bounds give the
weighted tail estimate used below, with the
bare-end powers included in \(\Pi_K\).
The local estimate still requires the compared
history masses on their actual test domains.

\paragraph{Feedback from the terminal pin.}
Let $y_u=\sup|\Delta_u|$ be the deviation difference on the
matched offset disc and let $D_s$ be the supremum of the
increment difference in \eqref{eq:comparison-matched-scalar-step}.
The pin gives $y_u\le\sum_{s<u}D_s$. On this disc the
required remainder-distance estimate bounds the last sum in
\eqref{eq:comparison-marked-four-chains} by
\[
 C(\rho_0^{K-s}+\Pi_K)
  +C\sum_{u>s}\left[(1+\varrho^{-1})\theta_{[s+1]}^{-\nu}
               \rho_0^{u-s-1}+p_u\right]y_u ,
 \qquad p_u=C\theta_{[u]}^{-\nu}(\log\theta_{[u]})^b.
\]
With this estimate and the one-cutoff variation at fixed
remainders, the marked difference obeys the same bound, with
a rate $\rho_0<1$ large enough to include the interpolation
loss. Add the one-loop and the two unmarked differences.
For $r_s=\rho^{K-s}+\Pi_K$, with $\rho_0<\rho<1$, the
result is
\begin{equation}
 D_s\le C r_s+\sum_{u>s}(k_s\rho_0^{u-s-1}+p_u)y_u,
 \label{eq:comparison-scalar-feedback}
\end{equation}
where $k_s$ is a sum of the one-cutoff sensitivity and
$C(1+\varrho^{-1})\theta_{[s+1]}^{-\nu}$. The feedback
acts in the opposite direction to the pin. Use the finite
maximum $Y=\max_sD_s/r_s$. Then $y_u\le Y\mathsf S_u$.
Against this profile, the geometric kernel in
\eqref{eq:comparison-scalar-feedback} costs
$C k_s[\rho^{K-s}+(s+1)\Pi_K]$.
The factor $(s+1)k_s$ is small at weak coupling since
$s+1\le\theta_{[s+1]}/c_0$. The polynomial kernel costs
at most $\Pi_K\sum_u u p_u+C K^{-\nu}(\log(2+K))^b$,
by the convolution estimate for \eqref{eq:comparison-matched-profile}.
Use weak coupling for the first term and the independent
prefactor of $\Pi_K$ for the second. Thus the kernel acting
on $\mathsf S$ costs at most $r_s/2$, and $Y\le C+Y/2$.
This proves \eqref{eq:comparison-matched-scalar-step} once
the marked comparison has the stated bound.

% Earlier Part III, 1862--1889, 91705--91966 and 120565--120739.
\paragraph{The comparison families on the common scaffold.}
The coupled comparison supplies the matched increment
estimate when its response and comparison rows apply
to the actual pair of runs. Write \(r\) for depth
from the terminal lattice. At the unmatched bare
end \(r=K\), place the longer run's family on the
shorter run's scaffold by the transfer that preserves
its total. On the common depths, proceed in decreasing
\(r\), and in the order of the production blocks at
fixed \(r\). For each new block use the shorter
run's production with the comparison families already
constructed in the input slots. This list uses the
longer run's totals. It belongs to the domain of the
response estimate because every inserted family
obeys the bound in the units of its own component.

There are two representatives of the difference in
each component \(c\). The first adds the change
\(M_r^c\) caused by the changed inputs of the shorter
run's production and the change \(N_r^c\) from
that production to the longer run's. The middle
total cancels. The second subtracts the shorter
family directly from the longer family transferred
to its scaffold. Its norm is at most
\((C_V+1)u_r^c\), where \(C_V\) bounds the transfer
and \(u_r^c\) is the one-run unit in this component.
Keep the first representative if its norm is
within this bound and otherwise use the second.
The norm is a supremum on the entire admitted
parameter, source and torus domains. The choice
is therefore fixed across those domains. The
relations between the components that involve
only their totals survive the choice. The
source-defect components vanish at constant
sources in both representatives, and the local
and frozen parts of a remainder component are
chosen together. The resulting family has the
required local reads and is holomorphic on
the admitted domain. Its norm \(g_r^c\)
satisfies
\[
 g_r^c\le (C_V+1)u_r^c,
 \qquad
 g_r^c\le\|M_r^c\|+\|N_r^c\|.
\]
In the second case the second inequality follows
because the first representative exceeded the
threshold. This closes the construction of
the next input list.

Use the blocks \(P=(\mathsf P,Q)\),
\(A=(P^\flat,D)\) and \(R=R'+R''\) in
their prescribed reading units. Here \(D\)
is the difference between the local source
read and its value with the source frozen
at the base point. The pointed block \(A\)
is scaled by an inverse power of
\(\log\theta_{[r]}\), and \(Q,R\) carry
the inverse squared coupling in their units.
At the matched depth \(r_0\)
the running-coupling difference is zero.
The increment functional on the \(P\) block
then bounds the difference \(d_{s+1}\) of
the two running couplings by
\[
 d_{s+1}\le A_0\sum_{r_0\le u\le s}g_u^P,
 \qquad r_0\le s<K .
\]
Insert this sum into the response rows, including
their weighted coupling tails. The resulting
inequalities have nonnegative coefficients in
the block norms. They may involve both earlier
and later depths.

For the envelope put \(e_r=\rho^{K-r}+\Pi_K\).
The three block bounds to check are
\(C_E e_r\), \(4C_Ee_r\) and
\(4C_E\widehat r_r e_r\), respectively.
Here \(\widehat r_r\) is the one-run remainder
coefficient of order \(\theta_{[r]}^{1-\kappa_0/2}\),
enlarged by the window factors and one factor
of \(1+r\) from the coupling sum. Its
bare-end value sets the size of the floor
\(\Pi_K\). For the geometric response kernel
\(q^{s-r}\), use \(q\le\rho^2\) and
\(e_{r+j}\le\rho^{-j}e_r\).
The coupling sum up to \(r\) is at most
\(\rho^{K-r}/(1-\rho)+(r-r_0+1)\Pi_K\);
convolution with the geometric kernel pays
the linear factor in the lag. For the
unweighted remainder sum use
\(\sum_s\widehat r_s\le1\) and the
one-run comparison
\(\widehat r_s\le C_T\rho^{-(K-s)/2}
\widehat r_K\). This bounds
\(\sum_{s>r}\widehat r_s e_s\) by \(2\Pi_K\)
with the prescribed prefactor in the floor.
In the short remainder window of length
\(\ell_r\), the relative one-run bound
\(\widehat r_{r+j}\le 2(1+j)\widehat r_r\)
leaves at most \(C\ell_r^2\rho^{-\ell_r}\)
times the coefficient of that window.
The small-coupling bounds for the three
rows make the sum of the weighted envelope
terms at most half the envelope in each
row, and the direct forcing at most the
other half. The transferred bare-end
family satisfies the envelope there.

To apply these tests simultaneously, take
the maximum of the ratios of the actual
block norms to their proposed bounds. If
the maximum were \(t>1\), it would occur
in a response row. At that row the two
half-envelope bounds give \(t\le(t+1)/2\).
Thus every block satisfies its envelope.
In particular the coupling increment
difference is \(O(e_r)\) on the common
offset disc, and the running-coupling
difference has the preceding summed
bound. This application requires the
matched productions and the response
rows in the same running radii as the
increment functional.

% Earlier Part III, 122654--122727.
For a comparison of the pointed remainder
shares, make the volume allocation in each
actual run before comparing them. The
allocation depends on the original filing
domains. For example, file a fixed
analytic function first on a cube wholly
inside one filing cell, and then on a larger
domain containing that cube and meeting a
second cell. The total function is unchanged.
The second cell receives zero in the first
case and the positive volume fraction of
the function in the second. At a shared
filing domain the two actual allocations
have the same geometric weights, so their
difference is the weighted sum of the
differences of the individual terminal
reads. That sum requires the paired
terminal-read estimate on its actual
carriers.

% Earlier Part III, 122691--122872 and 164579--164655.
Let \(i''\) be the filing level of a remainder
from level \(i\), and \(t_i=L^{i-i''}\) the
mesh ratio. The share \(\mathrm{sh}(X,y)\)
at \(y\) is the fraction of the level-\(i\)
cubes of \(X\) in that filing cell. Evaluate
\(R^{(i)}(X)\) on the box background formed
from the actual level-\(i\) averages in the
padded box of \(X\), and subtract \(R^{(i)}(X,1)\).
The assigned domain \(Y(X,y)\) contains
that box and \(Y^+\) contains the old source
read. The box-background and restriction
estimates on the filing radii give
holomorphy. On admitted backgrounds
restriction of the critical orbit identifies
the box evaluation with the original one.
Summing the shares over \(y\) gives one for
each \(X\), so the sum is the vacuum-subtracted
total. At constant sources the identity
before summing \(y\), with boxes and shares
transported under the lattice symmetries,
gives covariance of the pointed total.

% Earlier outline, 44361--44365.
For the one-run compact estimate, take the
box background along the complex radial
line in the coarse datum. Its inverse
analytic radius is bounded by
\(C(d_i(X)+1)t_i^2 w_{i''-1,i}^{1/2}\).
The first derivative of the remainder
at the identity vanishes. The second-order
Schwarz estimate then gives the square of
this inverse radius for the vacuum-subtracted
value. This is the \(t_i^4\) gain in the
one-run filing norm.

Write \(S_i^+[R]\) for the difference norm
of the level-\(i\) remainder and
\(\widehat\vartheta_i\) for the decay in
the datum comparison. At the shared
filing slice the paired read bound gives
\(S_i^+[R]+\widehat\vartheta_i g_i^{\kappa_0}\)
times the size of the individual read.
For compact \(X\) of diameter \(D_X\),
that size has the factor \(D_X^2t_i^4\)
times \(w_{i''-1,i}e^{-\kappa_Td_i(X)}\),
with \(w_{i''-1,i}\) the crossing weight.
There are \(O(\check R_i^4)\) fine
anchors in a filing cell. The diameter
cost is bounded by \(CM^2(d_i(X)+5)^2\)
and \(\check R_i^4 M^2t_i^4\le CM^{-2}\).
The target domains of compact reads have
bounded length. For long \(X\), the
length is at least \(t_i^{-1}-2\).
The output weight spends a fixed fraction
of the input length decay, and another
fraction leaves \(Ce^{-cM\check R_i}\)
to pay the anchor count. The crossing
weight at this filing level is bounded
below. The rooted sums give
\[
 S_{i''}^+[\mathsf{RF}R]\le
 C w_{i''-1,i}
       \bigl(S_i^+[R]+\widehat\vartheta_i g_i^{\kappa_0}\bigr).
\]
This estimate uses the compact and long
terminal read hypotheses at the actual
shared slice.

Localize the difference at each \(y\)
by the alternating box sum on the
shorter run's scaffold. Summing the
boxes cancels the proper restrictions
and leaves the full-background difference
at that \(y\). This gives a pointed
comparison family on the admitted domains.
Subtract its frozen value at \(W(p_y)\)
to obtain the source-defect component,
which vanishes at constant sources.
The factor of two for that subtraction
is paid in changing from the \(Q\) unit
to the pointed \(A\) unit under
\(E_0 C_t g^{-2}\ge2\), with \(E_0,C_t\)
the unit constants. Add these two filed
differences to the comparisons of the
ordinary \(Q\) and \(D\) outputs. The
hybrid production uses the filed components
already in its input. The filing operation
is carried out on the two actual runs.

An older remainder is read as such
only in the finite band before filing.
After that band its share is part of
the pointed regular family. The \(P\)
and \(A\) inequalities gain a sum of
remainder distances over the filing
depths and a term from the relative
datum comparison. At most two depths
file at \(r\), each within the window
of length \(\ell_r\). The relative
datum term costs \(C\rho^{-\ell_r}
\mathfrak r_r\) times the geometric
envelope, where \(\mathfrak r_r\) is
the one-run remainder size in its
reading unit. The sum of the remainder
distances in that window costs
\(C\ell_r^2\rho^{-\ell_r}\widehat r_r\)
times the same envelope. These are
the additional terms in the three
row smallness test. When the finite-band
response and source rows hold, the
maximum argument then gives the
geometric envelope with the pin at
the terminal coupling of the run
with filed shares.

\paragraph{The change of offset.}
The two terminal couplings agree at the matched offsets.
Subtract this equality at $t=0$ from the equality at $t$.
The bare terms leave $t-\mathfrak n_K(t)$, so the matching
defect is, up to sign, the sum of the centered increment
differences and the centered extra increment. The one-run
derivative bounds and $|\mathfrak n_K(t)|\le2|t|$ bound
each centered difference by $C|t|a_{[u]}$. By the Schwarz lemma and
\eqref{eq:comparison-matched-scalar-step} it is also bounded by
$C|t|(\rho^{K-u}+\Pi_K)/\varrho$. The extra increment costs
$C|t|a_{[K]}$. Thus
\begin{equation}
 |\mathfrak n_K(t)-t|\le |t|\Upsilon_K,\qquad
 \Upsilon_K=C a_{[K]}+C\sum_{u<K}
    \min\{a_{[u]},(\rho^{K-u}+\Pi_K)/\varrho\}.
 \label{eq:comparison-matched-scalar-rate}
\end{equation}
For the $O(\log(2+K))$ depths nearest $K$, use $a_{[u]}$ and
$\theta_{[u]}\asymp K$ as $K\to\infty$. Farther away, sum the
geometric bound, choosing the multiple of $\log(2+K)$ so its
sum is $O(K^{-2})$. The other cost is $CK\Pi_K/\varrho$.
For $\Pi_K=O(K^{-5/2})$ this gives
$\Upsilon_K=O_\varrho(K^{-3/2}(\log(2+K))^{p+1})$ for $p\ge0$.
The bound on $\Xi'_{K+1}$ gives the same rate for
$|\Xi_{K+1}(t)-\Xi_K(t)|$ by matching. Cauchy's estimate
on a smaller disc bounds the derivative
of $\mathfrak n_K-\mathrm{id}$ by $C\Upsilon_K$. This gives
the first line of \eqref{eq:comparison-capsule-CP9}.
Writing $N_K=-\Xi'_K(0)$, differentiation of
$\Xi_{K+1}\circ\mathfrak n_K=\Xi_K$ gives
$N_{K+1}\mathfrak n'_K(0)=N_K$. Hence $N_K$ converges to a
number in $[2/3,4/3]$, with error bounded by
$C\sum_{k\ge K}\Upsilon_k$.

\paragraph{The comparison through the sourced run.}
For the full comparison, index the reference choices by the
depth from unit scale. In the two runs the cubes, thresholds,
waiting depths and radii then agree at common depths.
%
Cut the histories at a depth \(T\) proportional to \(K\),
with \(K-T\) greater than the waiting depth there by
the fixed geometric margin. On the subcollection
with no creation at \(u\ge T\), the bare-end rules lie
beyond the retained births. After matching the real
decisions, pair the labels with the same regions and
retained rows at common depths. Comparison of the
weights still uses their actual stored variables and
factors. Consider
a record formed with root at common depth \(s<T\),
containing a birth at depth \(u\ge T\). The path to that
birth occupies at least \(u-s\) levels of the record.
Each node on the path contributes at least \(M^4\)
sites with its own logarithmic weight, as in
\eqref{hi:unweighted-site-count} and
\eqref{hi:volume-tail}. On the common increasing
depth frame these logarithmic weights are at least
a fixed multiple of \(\log\theta_{[s]}\).
Thus \(v^\natural\ge c_*M^4\log\theta_{[s]}(u-s)\).
Use the mass bounds at \(c'>c\) on the same
raised floor and write
\[
 e^{cv^\natural}
   \le \theta_{[s]}^{-c_*(c'-c)M^4(T-s)} e^{c'v^\natural}
\]
on this subcollection. The inequality holds before the
upper integral, with the same localized root sum.
Thus for \(s\le\alpha'K<T=\alpha K+O(1)\), with
\(0<\alpha'<\alpha<1\), this unpaired part has an
exponential cutoff gain. There is no extra factor for
choosing a position or the first such birth because
the subcollection is counted once in the mass of the
whole rooted record. A birth that has been healed and
erased from that record is accounted for through the
terms it leaves in the recursion.
The duration for which a formed mark is later carried
belongs to the live-component reserve and the global
label comparison. Its stored genealogy is fixed.

For a visible creation on the outer tree one can also
retain its own reserve in \eqref{hi:outer-event-reserve}.
In the proof of \eqref{hi:weighted-mass-conclusion},
choose a creation at \(u\ge T\) on this subcollection
in place of the arbitrary leaf. The other entropy and
volume allocations are unchanged. The factor
\(\exp[-2p_0(g_{[u]})]\) gives a cutoff gain faster than
any power of \(K\) for \(T\asymp K\), even if this
visible outer record is short. This estimate uses the
actual creation reserve. A failure whose witness has
been left for a later operation, or an earlier birth
internal to an inner island, uses its own mass.

For the remainders formed at \(u\ge T\), use the
individual size \(C\theta_{[u]}^{-\kappa_0/2}\) and the
readout weight \(C\theta_{[u]}\), with the logarithmic
factors of the norm. On the weak-coupling reference
line \(\theta_{[u]}\asymp K\) for \(T\le u\le K\).
The sum over these \(O(K)\) depths, including the
unpaired last remainder, is at most \(\Pi_K\) of
order \(K^{-\nu}\) times a power of \(\log(2+K)\),
where \(\nu=(\kappa_0-4)/2\). This part enters the
comparison at lower depths without an age factor.
A first disagreement in a short history at fixed
common depth \(s\) can occur with no ancestral path
reaching \(T\). Its birth-scale bound depends on
\(\theta_{[s]}\), so one still needs the local
near-threshold comparison.

Write $\Delta_u(t)$ for the difference of the two deviations
at depth $u$ and at the matched offsets $t,\mathfrak n_K(t)$.
The matching identity sets $\Delta_0=0$. The deviation
recursion, with the same reference increment in both runs,
expresses $\Delta_u$ as a signed sum of the scalar increment
differences through depth $u-1$. With
\begin{equation}
 \mathsf S_u=\frac{\rho^{K+1-u}}{1-\rho}+u\Pi_K,
 \label{eq:comparison-matched-profile}
\end{equation}
the scalar estimate \eqref{eq:comparison-matched-scalar-step}
gives $\sup|\Delta_u|\le C\mathsf S_u$. If the field $v$
takes values in a smaller offset disc, Cauchy's estimate there
bounds the derivative of $\Delta_u$ by $C\mathsf S_u$ with a
fixed radius cost. Thus the mismatch of the pointwise deviation
fields is at most $C\mathsf S_u$ in the maximum of the
supremum norm and the depth-scaled Lipschitz seminorm of the
admitted source field.

We must show on the two sourced trajectories that the distance
$X_s$ of their outputs at depth $s$ satisfies
\[
 X_s\le C\mathsf S_{s+1}
     +C\sum_{u\ge s}\alpha_{s,u}\mathsf S_u
     +\sum_{u>s}\varkappa_{s,u}X_u .
\]
Use a distance that compares values on common fields, with
exponential localization and the units of the defect and
remainder families. The small multiple of the individual
sizes needed for information lost by the transfer goes in
the one-step defect. The unmatched family at depth $K$ also
goes in that defect. Its scalar deviation can still enter
the input sum.

%
%
For the sourced trajectory the size of its defect can be
read off at formation. Let \(i\le m\) and use the source
norm in the units of level \(m\). In the defect at a base
point \(z\), compare with the constant \(w(p_z)\).
A tree meeting all cubes of \(X\) has \(d_i(X)\) edges
of length \(ML^{i-K}\). The prescribed enlargement
adds a bounded number of cube lengths to the diameter.
Thus
\[
 \sup_{p\in X^\sim}|w(p)-w(p_z)|
 \le C M L^{i-K}(1+d_i(X))\operatorname{Lip}w .
\]
Use the constant in the scaled source norm so that the local
norm of this difference is at most \(e^{d_i(X)}\mathsf l_i\),
where \(\mathsf l_i\) is its Lipschitz part. This uses the
boundedness of \((1+d)e^{-d}\). For a source ball of radius
\(r\), the unit-slope estimate on the admitted line bounds
the defect by \(r^{-1}e^{d_i(X)}\mathsf l_i\) times the
rate-\(\kappa+1\) decay. Since
\(\mathsf l_i\le L^{-(m-i)}\|w\|_{\flat,K-m}\), this gives
\[
 |D^{(i),w}(X,\cdot,z)|
  \le r^{-1}\|w\|_{\flat,K-m}L^{-(m-i)}e^{-\kappa d_i(X)} .
\]
For a drive in the unit-scale norm take \(m=K\), giving
size \(L^{-u}\) at depth \(u=K-i\), from the moment the
entry is made. The old defect is passed on unchanged.
The stronger admission bound required by the next box
comes from this original estimate on the actual source
trajectory. A one-step estimate on free input arrays
must assume that stronger bound on the inherited entries.

%
For a newly formed paired defect, put the two
step maps on the same slope line in the source.
At depth \(s\), take a common base point \(y\)
of the enlarged read and set
\[
 w_\lambda=w(y)+\lambda(w-w(y)).
\]
The partnered source uses the pointwise matching
map applied to \(w_\lambda\). In each run subtract
the formation map at its own frozen value at the
base point of the family. Both defects vanish
at \(\lambda=0\), even if the two frozen points
differ. The Lipschitz bound on the full read
admits the line up to radius
\(R\asymp L^s/(1+d(X))\).
Suppose the unscaled two-run array comparison
on this enlarged disc bounds the difference
of the two formation defects by
\(\mathcal B_X e^{-(\kappa+\delta)d(X)}\)
with \(\delta>0\). For \(R\ge2\) Schwarz's
lemma gives at \(\lambda=1\)
\[
 C L^{-s}(1+d(X))\mathcal B_X
                       e^{-(\kappa+\delta)d(X)} .
\]
For \(R<2\), the same bound follows from
the unscaled comparison at the actual source,
since \(L^{-s}(1+d(X))\) is then bounded
below. The spare decay pays for \(1+d(X)\).
The output unit \(L^{-s}\) is therefore
recovered for the paired defect. To separate
the displacement of the frozen points,
its bound on the line is
\(C|\lambda|\operatorname{Lip}(w)\) times their
physical separation, using the derivative of
the matching map. The slope estimate gives
that same displacement cost at \(\lambda=1\).
The step still needs the unscaled array
comparison on the admitted line.

%
%
On the common all-small label the carried Wilson
weight can be compared separately. At matched levels
use the same physical tent grid. Write \(w_A=w\) and
\(w_B=\mathfrak n_K\circ w\), and let \(T_u^a\) denote
tent interpolation at depth \(u\), sampling at the
barycentres of the plaquettes \(p_y^a\) based at the
tent vertices. Set
\[
 f_u^a(x)=b_{[u]}^a(w_a(x))-b_{[u]}^a(0).
\]
Here the increments are written as functions of the
constant source in their respective runs. Denote the
constant-source shift at depth \(s\) by \(\zeta^a(t)\)
and the actual shift on plaquettes by \(c^a\). The
formula for the carried shift on this label gives,
at a common physical position,
\[
 c^a(x)-\zeta^a(w_a(x))
      =\sum_{u}(T_u^a-I) f_u^a(x),
\]
where the sum is over the increments already made.
The difference of the two errors is
\[
 \sum_{u\ {\rm common}}(T_u^A-I)(f_u^B-f_u^A)
 +\sum_{u\ {\rm common}}(T_u^B-T_u^A) f_u^B
 +(T_*^B-I) f_*^B .
\]
The star is the extra first increment in run \(B\).
For a common depth, the constant-source recursion
expresses the difference of the increments as the
difference of two successive matched coupling
mismatches. Subtracting its value at zero and using
Cauchy's estimate on a smaller source disc gives
\[
 \operatorname{Lip}(f_u^B-f_u^A)
 \le C\varrho^{-1}\operatorname{Lip}(w)
                    (\mathsf m_u+\mathsf m_{u+1}) .
\]
If \(\xi_u\asymp L^{-u}\) is the tent width,
\(T_u^A-I\) costs at most \(C\xi_u\) times that
Lipschitz bound. The ratio \(\xi_u/\xi_s\) gives
the geometric share in the source units at depth \(s\).
In the second sum only the sampling plaquettes move,
by \(O(a_K)\). The unmatched tent has width \(O(a_K)\).
These two terms together cost \(C a_K\operatorname{Lip}(w)\),
using the summable derivative bounds for the increments.
For later quadrature, first differences of the
multilinear tents give
\(\operatorname{Lip}(c^a)\le C\operatorname{Lip}(w)\).
Indeed the samples at neighboring vertices differ
by at most \(C\xi_u\operatorname{Lip}(f_u^a)\), which
pays for the \(\xi_u^{-1}\) in a tent derivative.

In the action the scalar coefficient is \(x^a-c^a\).
The difference of its frozen parts is the paired
constant-source coupling mismatch, and the two
tent errors give the additional difference just
estimated. This uses the scalar weight as a
potential in the frozen step. The Gaussian
structure and the independent analytic current
retain their own arguments.

For an inherited gauge-invariant defect, separate
the family difference on the transferred fine
argument from the change between the two
evaluation arguments of the flat family.
After subtraction of the vacuum values, the
first part uses the vacuum-only window estimate
of \cref{cor:vacuum-only-window}, with its
squared ratio of field radii. The size already
present in the inherited defect supplies
\(L^{-(u-s)}\) in the output units. The second
part uses the actual paired evaluation fields
on the complete reads. It is a separate term
in the unscaled array comparison.

The regular frozen families are compared one source value at
a time, with their own counterterm. A displaced barycenter costs the Lipschitz
norm of $v$ times its displacement. The source defect at
depth $u$ already has the size factor $L^{-u}$. Relative to
the output unit $L^{-s}$, this contributes $L^{-(u-s)}$
times the squared ratio of the field radii. Marked remainders
have the squared ratio and their own power bound. In the
recent window the term passed on by the ordinary operation
has a logarithmically small coefficient multiplying the
input difference. It must be included when closing the
full recursion. In the marked operation the dependence
before the mass estimate can instead cost a power of
$\log\theta_{[s]}$. There the weighted mass and the age sum
must supply the smallness.

The required kernels have a geometric part $c_{u-s}$ and a
remainder part $p_u$ of size
$C\theta_{[u]}^{-\nu}(\log\theta_{[u]})^b$ with a fixed
constant. For $q\ge0$ compare
$\mathsf S_{s+q+1}$ with $\mathsf S_{s+1}$ using the factors
$\rho^{-q}$ and $1+q$ on the two terms of
\eqref{eq:comparison-matched-profile}. Thus the geometric
part preserves the profile if both
$\sum_qc_q\rho^{-q}$ and $\sum_q(1+q)c_q$ are finite.
For the remainder kernel, $u(1+K-u)\ge K$ for $1\le u\le K$.
The polynomial size bound for $p_u$ and this inequality
give $\sum_{u\le K}p_u\rho^{K-u}\le C K^{-\nu}(\log(2+K))^b$.
The other part costs $\Pi_K\sum_u(u+1)p_u$. This sum
can be made small at weak coupling by $\nu>2$ and
$\theta_{[u]}=g^{-2}+c_0u$. In the geometric convolution
at the bare end, smallness relative to $\Pi_K$ instead
comes from choosing the prefactor of $\Pi_K$ large after
the kernel constants. The same bounds apply to the input
kernel on $\mathsf S_u$, where boundedness suffices.
With small weighted moments of $c_{u-s}$ in the feedback
kernel, its full action on $\mathsf S_{u+1}$ is at most
$\mathsf S_{s+1}/2$. Backward induction gives
$X_s\le C'\mathsf S_{s+1}$ uniformly in $K$.

The local energy comparison at $(\mathbf 1,0)$ uses this
profile. A carrier of size $d(X)$ has at most
$CM^4(1+d(X))$ base points. A fixed reserve in the exponential
decay pays for this count, giving the local estimate with
$C\mathsf S_{s+1}e^{-\kappa_e d(X)}$.
Since $\mathsf S_{s+1}\le C(1+s)(\rho^{K-s}+\Pi_K)$,
this is the rate used for the mixed derivatives. For the
terminal integral we must also compare the weights on the
common terminal fields, with the same scalar at the source
and at its real center. The corresponding terminal estimate
uses the kernels at depth zero and gives
$C(g,\ell)\mathsf S_1$ on the admitted source ball. Moving
the stop by a fixed number of depths restores the corresponding
local energies. The same profile bounds their differences
at this fixed stop.

% Earlier Part III, 54212--54269, 55400--55412, 55981ff., 56438--56494,
% 56967--56990 and 121119--121183.
\paragraph{The source-independent part of the terminal comparison.}
For a native healed component, form the two remainders on
the same prepared determining sets. Keep the real frozen
data, masks, Haar coordinates and weights of the common
geometry. With the cutoff \(\zeta_0\) of the component,
the last term is
\[
 K_5=-\log\int e^{-S(B')}\,d\mu(B'),
 \qquad
 S(B')=x_m\bigl[A(\zeta_0,U^{B'})-A(\zeta_0,U_0)\bigr].
\]
The common measure \(d\mu\) includes the Haar density, the
gauge constraint and the sharp integration domain. The
remaining terms of the remainder come from the exact Wilson
and quadratic split. In its notation the remainder is
\(\mathbf C_m=K_5+K_3-(E_{12}+E_3+E_4+K_1+K_2+K_4)\).
The Wilson terms use the same weights in both runs. For
\(K_5\), divide the two integrals before taking the logarithm.
On the real data the ratio is the expectation of
\(e^{-(S_B-S_A)}\) under the normalized \(A\) integral.
Thus its log difference costs at most
\(\sup|S_B-S_A|\) on the same cut. This comparison uses the
whole action difference in \(S\) with its own centre.

At real common data the layerwise bound on the paired
minimizers is obtained by telescoping their data defects.
The paired data at each occurrence have bounded depth.
The weighted response bound controls the change from the
defect at each step, and the common local index fixes the
rows. For a Wilson cell this first bounds the curvature
difference, after transporting to the same frame. The
action estimate multiplies it by the bound on the two
curvatures and sums the plaquettes of the cell. Before
multiplying by \(x_n\), the result is
\(CM^4\mathsf e_m^2(\vartheta_A^n+C_\chi L^{-n_\chi})\),
where \(n_\chi\) is the cutoff's scale exponent on the
cell. The native cutoffs have \(n_\chi\ge n-1\), so
the cutoff term has the same decay. Two-constants
interpolation on the fixed comparison discs gives the
square root of the rate on the common complex data.
The geometric mean with the one-run bound decaying
into the core then leaves \(\vartheta_A^{n/4}\) for
these Wilson cell terms. The normalization and the
cutoff block use their whole-term estimates.

In localizing \(E_{12}\) and \(E_3\), assign the
switch-independent subtraction to the smallest-box piece
of its own cell. For the other pieces it disappears under
the switch difference. The two-run estimate applies at
the natural binary vertices, where the boxes contain the
component and the required cell padding. The complex
polydisc bound uses the sum of the separate one-run
majorants. The interpolation estimate combines the two
bounds, with a rate reserve for the tuple count. This gives
the native class estimate
\[
 \sum_{Y\subset C^\sharp}
 \sup\bigl|\delta\mathfrak v^{\mathbf C,(n)}(Y)\bigr|
          \le C\vartheta_w^{n/8},\qquad 0<\vartheta_w<1 .
\]
The class occurs only at the prepared levels in the
component and at the two endpoint levels of the split.
There is no coupling-difference term from these common
weights. This supplies the source-independent load for the
matched native remainder. Its use in a terminal read
requires the same measure, masks and normalization above.
A remainder depending on the variable source uses the
separate excess comparison. A source-independent factor
that varies with the label still occurs in the weights of
both stop integrals.

% Earlier Part III, 89847--90040, 90095--90134,
% 90754--90854 and 90940--91102.
\paragraph{The exterior sum of the Wilson remainder.}
For the outer part of \(E_3\), the unit on a
cell \(\Delta\) is the weighted Wilson value at
the box field minus the value at the root, with
the same cutoff and plaquette base points.
Its core \(K_\Delta\) is the level-\(m\) cube
of \(\Delta\) with four layers of padding.
For a treated component \(C\) let \(H_C\) be
its protected block with one layer of padding
and put
\(\epsilon_\Delta=\sum_C e^{-\delta_P d^\wedge(\Delta,H_C)}\).
At the switch vertices the paired Wilson
estimate is bounded by
\(C M^4 x_l\mathsf e_m^2
\epsilon_\Delta^{1/2}\langle l\rangle^{1/4}\),
where \(l\) is the level of \(\Delta\) and
\(\langle l\rangle\) is the two-cutoff rate
at that level. On the contour the four
Wilson values are bounded by
\(C M^4x_l\mathsf e_m^2\), using their admitted
plaquette bounds on the complete core.
The switch estimate gives the square root
of the product and an exponential in the
number of switched cells. The switch sum
at an exterior cell \(Q\) spends this last
exponential to pay for the tree weight
and leaves \(Ce^{-\operatorname{dist}_m(Q,K_\Delta)}\).

Use subadditivity to distribute
\(\epsilon_\Delta^{1/4}\) over \(C\). For
each \(C\), split its decay into four equal
parts. The first pays to replace
\(x_l\langle l\rangle^{1/8}\) by
\(Cx_m\vartheta_\diamond^m\), using the
flow and the crossed zones, where
\(\vartheta_\diamond\) is the exterior
two-cutoff rate. The second leaves the
gain \(e^{-5\check R_m}\) from the
separation of the outer cells from \(H_C\).
The third joins the distance from \(Q\)
to \(K_\Delta\) to that from \(\Delta\)
to \(H_C\). The fourth sums the cells
relative to the boundary of \(H_C\).
The separated-block count sums the
resulting weights at \(Q\) over \(C\).
Thus the weighted norm is at most
\(Cx_m\mathsf e_m^2e^{-5\check R_m}
\vartheta_\diamond^m\). The factor
\(x_m\mathsf e_m^2\) is polynomial in
\(\log x_m\), and for a class \(n\le m\)
one has \(\vartheta_\diamond^m\le
\vartheta_\diamond^n\le C\Delta_n^c\).
The analogous
one-run vertex estimate starts with
\(\epsilon_\Delta\) and no two-cutoff
rate, leaving \(\epsilon_\Delta^{1/2}\)
after the switch estimate. Three of its
eight decay parts pay the lift, anchor
and cell count, and the remaining five
give \(e^{-25\check R_m}\).

For \(E_{12}\), first compare the field in a
box \(G\) with the restriction from a box
containing both \(G\) and the treated component.
The determining data differ only on the new
outer rows near the cut. Then compare the
restriction of the larger field to the field
of the component's own box. In this comparison
only the component's boundary shells supply
the mismatch. Both mismatches are outside
the core collar and are \(O(\mathsf e_m)\).
The two localized responses give the one-run
depth bound on the cell, which combines
with the paired vertex bound and the contour
majorant. In the exterior sum the switch
family must reach from a core inside the
treated component to an exterior cell.
Its minimum length leaves \(Ce^{-c\check R_m}\)
after paying for the tree weight and the
family count. The depth-weighted cell count
costs a polynomial in the logarithm of
\(x_m\). This gives a small coefficient times
the class rate. Retain the class of the
original cell or endpoint row. The
outer \(E_3\) terms at older levels use
the classes in the untreated live regions.
The full pieces confined to a treated
component use the value estimate at its
own box, rather than this exterior gain.

% Earlier Part III, 61090--61111, 61461--61467,
% 61633--61668 and 78937--78948.
For the source-free remainder on the aligned scaffold,
the exterior estimate is used on the filed domains
\(Y_4\). At an exterior root \(Q\), terms filed as a
value on one treated component do not occur. The label
\(\mathrm{lab}_4\) of each remaining piece is formed
from the four successive switch families. Filing adds
whole touched components of \(Z\), so \(Q\in Y_4\)
implies \(Q\in\mathrm{lab}_4\) and
\(d_{m,Z}(Y_4)\le d_m(\mathrm{lab}_4)\). For the
admitted outgoing pieces the label also reaches the
source cell of a nonempty switch family. The protected
margin between such a cell and \(Z^c\) gives
\(d_m(\mathrm{lab}_4)\ge\check R_m-2\).

With \(a=31\kappa/20\), spend one additional unit of
ordinary length weight to extract \(e^{-(\check R_m-2)}\).
The four successive rooted switch sums each spend at
most two more units of weight. The rate condition
\(\kappa_1^\natural\ge c_l(a+9)+\lambda_*\) pays these
losses, where \(c_l\) is the length allowance per switch
cell and \(\lambda_*\) pays the connected count. What
remains is the rooted sum of the original units of
class \(n\), whose supports have length zero. Its
weights are \((2\beta_v\mathfrak B_v)^{1/2}\), with
\(\beta_v\) the paired binary bound and \(\mathfrak B_v\)
the one-run contour majorant. At an interior cube only
units of its treated component contribute. At an
exterior cube there is just the outer unit there. The
classwise cell sum gives \(C\vartheta_{\mathbf C}^n
\mathsf N_{\mathbf C}(m)\), where
\(\mathsf N_{\mathbf C}(m)\) is polynomial in
\(\check R_m\). Thus the filed load satisfies
\[
 \sup_{Q\not\subset Z}\sum_{Y_4\ni Q}
 e^{a d_{m,Z}(Y_4)}
 \sup|\delta\mathfrak v^{\mathbf C,(n)}(Y_4)|
 \le C e^{-(\check R_m-2)}
        \vartheta_{\mathbf C}^{\,n}\mathsf N_{\mathbf C}(m).
\]
This sum uses the common switch scaffold and the native
piece bounds at its admitted reads.

% Earlier Part III, 38914--38979, 85546--85833,
% 121073--121114 and 121675--121763.
\paragraph{The reads on a live strip.}
Use two sets for a stored term on a live strip. Its true
domain \(X\) specifies its fields and dependencies. For
counting, choose an optimal tree for its relative length
and use the union \(X^\theta\) of the cells met by the tree.
Take the fat core of \(X^\theta\) for the ordinary
expansion. The thickening bound here uses the ordinary
tree length of \(X^\theta\). A strip block whose enlargement
reads \(X^+\) outside this core carries a separate parameter
for each response increment that reads it. The initial core
contains the tree and its prescribed padding. For a selected
set \(D\) of these parameters the new core adds only their
four-layer blocks and the corridors, with the fixed final
enlargement. This includes the complete carrier of the
coefficient selected by \(D\).

In the native kept term of a component \(Y\), subtract
the value at its declared reference. The true domain is
\(Y^\sharp\) with its assigned adjacent cell
\(\Delta_1(Y)\). Both lie in the one-layer enlargement
of the strip, and the tree domain is \(\Delta_1(Y)\)
alone. At a fixed level the strips born there are
separated, so at most one strip per face can choose
the same cell. This gives at most \(2d\) choices at
that cell. The reference value stays in the original
identity for that label. The nonempty switch difference
and the plain derivative both kill that constant.
On the admitted tube the resulting coefficient is
therefore bounded using the oscillation of the kept
term. At a fixed anchor the summed oscillation budget
is at most \(2d c_H^{\rm kp}C_\Sigma^U\check R_m^{p_\Sigma}\).
On the admitted strip discs
of radii \(r_a=r_0e^{E_a}\), a nonempty difference in the
parameters \(a\in D\) gains
\(\prod_{a\in D}(2/r_0)e^{-E_a}\).
Indeed its alternating vertex sum is the integral of
the mixed derivative over the unit cube. The remaining
radius at each real point is at least \(r_a/2\).
The domain and locality rules for the strip expansion
then treat this coefficient as an ordinary stored term
with its enlarged core. For a plain-block derivative
the source bound carries the separate decay in the
plain block and in the domain's distance to the shift.

In the kept family the source separation supplies
\(E\ge\varpi(\tfrac52\check R_m-8)\), with
\(\varpi=\delta_1M/8\), in the factor reserved for
the strip count. All read blocks lie within eight
level-\(m\) units of \(Y^\sharp\), and the source
is outside the collar of width \(5\check R_m/2\).
For a plain-block derivative split the domain factor
\(e^{-E/2}\) into two factors \(e^{-E/4}\).
The first pays the original level sum under
\(\varpi/8\ge2\log(2L^d)\). For a nonempty switch
family split each \(e^{-E_a}\) in halves. One half
per block pays the original block sum under
\(\varpi/2\ge(\bar a+2)a_d+1\) and
\(\varpi/4\ge4\log(2L^d)\), with \(\bar a,a_d\)
the exponents of the block count. A spare half
from any selected block pays the kept count.
The plain block's own decay is retained.
In either case a factor at least as small as
\(e^{-\varpi(\frac52\check R_m-8)/4}\) remains for
the polynomial anchor budget. In particular
their effective contribution is bounded by a fixed
constant times
\[
 C_\Sigma^U\sup_{R\ge1}
             R^{p_\Sigma}e^{-\varpi(\frac52R-8)/4}.
\]
This uses the source separation on the actual
response paths as well as the native strip-domain rule.

% Earlier Part III, 38621--38640, 38850--38910,
% 39266--39276, 39438--39489 and 39695--39724.
\paragraph{Room for the enlarged strip parameters.}
In one run, let \(\Omega^{\mathsf s}\) be the
source region and \(D(x)\) the distance to it in
the scaled path metric. For a datum \(b\) supported
there, write \(\mathsf P_b\) for the exponential
smearing of \(|b|\) at rate \(\delta_1\).
The geometric summation bound gives
\(\mathsf P_b(x)\le c_1\|b\|_\infty
e^{-\delta_1 D(x)/2}\). The strip letters grow
with distance at rate \(\delta_1/8\). On the
enlarged blocks their envelope is
\(\max\{c_0^\chi,r_0^\chi C_\Delta
e^{\delta_1 D(x)/8}\}\), with \(c_0^\chi\)
the bound for ordinary letters and \(C_\Delta\)
the fixed enlargement cost. In the partition
of unity, at most \(2^d\) cutoffs overlap,
and the ratio of the local scale to a cutoff's
scale is at most \(L\). The first and second
differences cost \(L\) and \(L^2\), respectively.
Use \(C_\chi=2^d C_\zeta L^2\) for both
derivative bounds, with \(C_\zeta\) the
partition bound, and take the radius floor
\(8C_\chi\) before rescaling the strip radii.

Multiplying the letter envelope by
\(C_\chi\mathsf P_b\) bounds the resulting
majorant by
\[
 C_\chi c_1\|b\|_\infty
 \max\{c_0^\chi e^{-\delta_1 D(x)/2},
        r_0^\chi C_\Delta e^{-3\delta_1 D(x)/8}\}.
\]
Dividing the strip radii by \(C_\chi\)
cancels that factor in the strip term.
The ordinary letter bound still carries it.
For the pointwise room, if
\(|b(y)|\le\varepsilon u(y)\) on the source
region and
\(\sum_y e^{-\delta_1 d(x,y)/2}u(y)\le C_T u(x)\)
in the same units, split the smearing exponent
in half. The identical profile estimate then
holds with \(c_1\|b\|_\infty\) replaced by
\(C_T\varepsilon u(x)\). To use the enlarged
discs in the actual box we must fit the first
bound under the absolute response ceiling and
the second under \(\ell_\iota(x)a(x)/4\),
where \(\ell_\iota\) is the label scale and
\(a\) is the local response majorant. This
requires comparing the native datum budget
with the local room on the whole box,
including its covered and inactive sites.

For the two-run reading difference at a foreign site,
hold the leaf function of the second run fixed and
subtract its evaluation at the first run's read with
the second run's common datum slots inserted. Write
\(h(t)\) for this difference in the strip parameters.
The native reading comparison uses the full one-run
discs, the relative row of the additional strip shift,
and the net datum comparison on the true hull.
It bounds \(h\) by \(C\vartheta^\sharp\|F\|\) on those
discs, with \(\vartheta^\sharp\) taken at the smaller
of the leaf level and the lowest zone of its true hull.
Cauchy's estimate as above bounds the mixed strip
piece by \(C\vartheta^\sharp\|F\|\prod_{a\in D}2/r_a\).
In the subsequent affine piece estimate use the smaller
radii \(r_a/4\). Both bounds in the geometric mean
carry the full product on these radii, so the square
root falls on \(\vartheta^\sharp\). A change of
the leaf function at a shared read uses the ordinary
one-run strip filing. At the component's own treatment
the terms touching its Gaussian core use the comparison
about the moved real centre.

% Earlier Part III, 15593--15966 and 16002--16170.
\paragraph{Counting stored domains away from live strips.}
At level \(\ell\) let \(W_\ell\) be the union of the
four-layer enlargements of the level-\(\ell\) covers
of the live strips. Its touching components are the
islands of this counting convention. File a domain by
including every island it meets. This changes its
index and preserves the true dependency set. The
relative length \(d_\ell^\circ(Y)\) is the least
length of a tree in the filed \(Y\) that meets every
cube outside \(W_\ell\). The length is measured
along the tree, including any path through an island.

For the blind filing used in this count, healing a
component absorbs the terms filed on it, including
terms indexed by a domain containing one of its
islands. If such a term survives until \(m\ge\ell\),
its excused islands at \(\ell\) survive with it. Their
births belong to one component after the youngest
birth, and healing that component would absorb
the stored term with the rest of its local terms.
The coarser cover of its excused set lies in \(W_m\).
Indeed, for \(m>\ell\) the four finer layers have
width \(4L^{-(m-\ell)}<1\) in level-\(m\) cells.
The cover thus adds at most one layer about the
coarser cover of the still-live strips. If a
level-\(m\) cube of the filed domain is outside
\(W_m\), it contains a level-\(\ell\) cube of \(Y\)
outside its excused set. The old relative tree meets
that cube, so it serves as a tree for the filed
domain at level \(m\) with length at most
\(L^{-(m-\ell)}d_\ell^\circ(Y)\).

At an exterior level-\(\ell\) root, the blind relative
domain count gives a bounded sum with rate
\(\lambda_0\ge\mathsf k_d+1\). In that count each
cube meets at most \(3^d\) other cubes, and a
neighbor belongs to at most one island. Now take an
exterior level-\(m\) root \(Q\) at distance at least
one from \(W_m\). Set \(y=L^{m-\ell}\). Its \(y^d\)
subcubes can serve as level-\(\ell\) roots. Islands
whose births have healed can be dropped for this
count since no surviving domain includes them.
If a surviving domain counted at \(Q\) contains
an island \(I\), the distance from a root subcube
in \(Y\cap Q\) to \(I\) is at least \(y\). A chain
in \(Y\) from \(I\) to that subcube contains a first
cube off \(I\). It touches \(I\) and lies outside
all the islands. The relative tree meets it and
the root subcube, so its length is at least \(y-3\).
At rate \(\lambda\ge\mathsf k_d+2\), spend one
unit of the decay on this distance. The remaining
rooted sum over each subcube costs at most
\(C_{\rm fib}\). The total is at most
\[
 C_{\rm fib} y^d e^{-(y-3)} ,
\]
uniformly bounded for \(y\ge1\). If the root is an
island \(I\) itself at level \(\ell\), the pure
domain \(Y=I\) contributes one. Every larger domain
containing \(I\) has an exterior cube touching \(I\).
Rooting at these cubes bounds the sum by
\(1+3^d C_{\rm fib}N_\ell(I)\), with \(N_\ell(I)\)
the number of cells of \(I\) in this partition.

An output of \(\mathbf T_\ell\) that meets an old
island also meets \(\Omega_\ell\). The old-strip
buffer gives distance at least \(3\check R_\ell-6\)
from that island to \(\Omega_\ell\). The same first
exterior cube argument gives
\(d_\ell^\circ(Y)\ge3\check R_\ell-8\).
One unit of the rate therefore pays the crossing
weight of this store. These are counts on the
filed domains. Analytic estimates use the true
dependency sets and the cores of their actual trees.

% Earlier Part III, Lemma 9.5 on density over strips,
% 82226--82343.
\paragraph{The one-run sum at a strip tree cell.}
For a term \(F\) formed at \(j<m\), write \(A=X^\theta\)
for its tree domain and \(\mathfrak D_F\) for the tube
on its true domain. Its size after the reference
subtraction is bounded on \(\mathfrak D_F\) by
\(2c_1^\flat e^{-(1+\beta/2)\kappa d_j(A)}\) for
\(\mathbf R\) and by \(x^{\mathbf T}e^{-\kappa_{\rm out}d_j(A)}\)
for \(\mathbf T\). The formation estimate for a core
at the reading level gives
\[
 d_m(K(A))\le c_H\bigl(1+L^{-(m-j)}d_j(A)\bigr).
\]
For \(r\le8\kappa/5\), the floor on \(L\) gives
\(rc_H/L\le a_\flat/4\). Thus the core weight
\(e^{r d_m(K(A))}\) spends at most \(a_\flat/4\)
of the old length decay, besides \(e^{rc_H}\).
The density norm spends a further \(a_\flat/4\).
The remaining rate for the healing term is
\((1+\beta/2)\kappa-a_\flat/2\ge c_d+1+\lambda_\theta\),
and the fluctuation term has at least the same
reserve. First count all true domains with the
same \(A\), at cost at most
\(e^{\lambda_\theta(d_j(A)+1)}\). Then sum the
trees through the specified root with the remaining
rate \(c_d+1\). The two families together cost
\[
 e^{rc_H+\lambda_\theta}K(d)
                         (2c_1^\flat+x^{\mathbf T}).
\]
Preparatory chain terms add their stage bounds in
the tree norm. At \(j=m\) use the stronger
creation rate for the terms that carry it,
since the core length has no scale reduction there.
This proves the stated rooted sum for the admitted
one-run terms. For the paired estimates the
coefficient and domain comparisons are needed
at each reading site.

% Earlier Part III, 81567--81726 and 83809--83834.
For the strip class, take an actual reading cell
\(\Delta\subset C\cap\Omega_n\) of level \(n\) in
the native grading of the treated component \(C\).
Every term in this class carries a pocket. Let
\(\mathsf b(m-j)\)
be the bound for the weighted sum at a tree root
of level \(j\), including the reserve
\(\varsigma_0 d_j(A)\), \(\varsigma_0=a_\flat/4\),
for \(j<m\). A true read reaches \(\Delta\) either
through the one-layer enlargement of \(A\) or
through a pocket strip enlarged by three layers.
Every attached pocket has a tree root within two
cells of its strip.

If \(j=n-u<n\), the pocket was born no later than
\(j\), and its strip satisfies
\(\operatorname{dist}_n(Y^\sharp,\Omega_n)\ge 5/2\).
Thus its
three-layer enlargement at level \(j\) cannot reach
\(\Delta\subset\Omega_n\). A direct read needs a
tree cube near \(\Delta\) as well as the cube near
the pocket. The tree length is at least
\((\tfrac52 L^u-6)_+\). There are at most
\((L^u+2)^d\) choices for the root near \(\Delta\).
The resulting count over \(u\ge1\) is bounded by
\(C^{<,\theta}\mathsf b(m-n)\), since
\[
 C^{<,\theta}=C_{\mathsf b}\sum_{u\ge1}(1+u)^{p_{\mathsf b}}
       (L^u+2)^d e^{-\varsigma_0(\frac52L^u-6)_+}
 <\infty ,
\]
using the polynomial lag bound for \(\mathsf b\).
For \(j\ge n\), a direct read or a pocket of side at
most one has a root in a fixed neighborhood of
\(\Delta\). A larger pocket has all sides at least
\(1/101\) in its level-\(j\) units. Its strip must
meet the seven-cell box about \(\Delta\), and the
disjoint strips then fit at most \(909^d\) such
pockets against that box. A pocket of largest side
\(s_j\) supplies at most \((s_j+7)^d\) roots.
If three level-\(j\) layers reach a quarter of
the pocket's birth radius, \(s_j\le1212\).
Otherwise use the full neighborhood of the birth
strip in the screened core. Its width at birth
level \(m'\) is \(5\check R_{m'}/2\) for \(m'<h\),
and \((\check R_h+1)/2\) for \(m'=h\), where \(h\)
is the healing level of the core. Its cells have
level at most \(m'\), so a contour leaving it
from \(\Delta\) spends at least \(w\) in contour
length per traversed birth-level unit, with \(w\)
the depth unit of the screened grading. Thus
\(\varrho(\Delta)\ge1+\check R_{m'}/4\), and the
root count is
at most \(C(1+\varrho(\Delta))^d\), with
\(\varrho(\Delta)\) its contour-depth variable.
Summing these levels costs
\(C(m-n+1)(1+\varrho(\Delta))^d\mathsf b(m-n)\).

On an outer rim cell \(\Delta\subset
(C^\sharp\setminus C)\cap\Omega_m\) of level \(m\),
consider terms
whose cores lie in \(C^\sharp\). A pocket-only
read at this cell is impossible. An own pocket
is at least \(9\check R_m\) from \(C^c\).
A foreign pocket is far enough from \(C^\sharp\)
that its root within two cells would put the
core outside \(C^\sharp\). The strip sum on
this rim consequently costs just
\((3^d+C^{<,\theta})\mathsf b(0)\).
The estimates use the native reading sets.
They supply the local count with no factor
depending on the age of a pocket.

% Earlier Part III, 81734--81886.
\paragraph{The strip sum in a healing component.}
Use the strip terms whose full admitting cores lie
in \(C\), so that their true reads \(X_F^{\mathsf r}\)
are contained in \(C\). Write \(\mathsf C\) for the
never-integrated core of \(C\), \(h=m-N\) for its
healing level, and \(W(F)\) for the nonnegative
weight in the cell estimate. If \(X_F^{\mathsf r}
\subset\mathsf C\), choose a cell of the true
read closest in contour distance to \(\mathsf C^c\).
The contour factor \(e^{-\delta_P\hat D_F/4}\),
with \(\hat D_F=d^\wedge(X_F^{\mathsf r},\mathsf C^c)\),
is at most the corresponding factor for this
cell. The rooted sum therefore reduces to
\[
 \sum_{\Delta\subset\mathsf C}
 e^{-\delta_P d^\wedge(\Delta,\mathsf C^c)/4}
 \sum_{F:\,X_F^{\mathsf r}\cap\Delta\ne\varnothing} W(F).
\]
At a cell of depth \(\varrho\) and level \(n\),
crossing the intervening zones gives
\(h-n\le\varrho\). Consequently
\(m-n+1\le (N+2)\varrho\), and the polynomial
lag bound for \(\mathsf b\) gives
\(\mathsf b(m-n)\le C(1+\varrho)^{p_{\mathsf b}}
\mathsf b(N+1)\). The strip cell estimate is
then at most \(C(N+2)\mathsf b(N+1)e^\varrho\).
Use half the screening exponent to pay \(e^\varrho\):
\(d^\wedge(\Delta,\mathsf C^c)\ge(\varrho-1)w\) and
\(\delta_P w/8\ge2\). The remaining exponential
sum is bounded by the contour cell-sum estimate
times the number \(\mathsf n_\partial(m)\) of
boundary cells of the closed core complement.
Thus the screened strip sum is bounded by
\(C\mathsf n_\partial(m)(N+2)\mathsf b(N+1)\).

If the true read leaves \(\mathsf C\), root it at a
cell in the closure of \(C\setminus\mathsf C\).
There \(\varrho=1\) and \(h\le n\le m\). The strip
cell estimate, summed over these cells, costs
\(C\mathsf n_{\rm sh}(m)(N+1)\mathsf b(N)\), with
\(\mathsf n_{\rm sh}(m)\) the cell count of this
region. This case includes a term whose tree
stays in \(\mathsf C\) while its true read leaves
it. These are one-run estimates for the strip
part of the healing sum.

% Earlier Part III, 117377--117423 and 117674--117848.
\paragraph{Oscillation of a stored term with framed slots.}
Consider the read of a stored term \(v\) with retained
variables that rotate at sites inside the blocks.
Use two admitted seeds with the same frozen component
datum. In the screened core the two reads can be
brought to one real centre \(U^s\) using the
real gauges of the frozen-centre construction.
Write their fields as \(q_i=e^{i\eta X_i}U^s\)
in that gauge, for \(i=0,1\). If \(g_i\) is
the change to this gauge and \(Q_i\) is the
product of the frame factors of the original
presentation, the physical field is \(q_i\)
acted on by \((g_iQ_i)^{-1}\). Hence both
physical reads can use the same value of the
retained variables.

Set \(X_\lambda=(1-\lambda)X_0+\lambda X_1\)
and interpolate the logarithms of the factors
\(c_i=g_iQ_i\) in the same way. With
\(q_\lambda=e^{i\eta X_\lambda}U^s\) and
\(c_\lambda\) the exponential of the interpolated
logarithm, compute the slice datum \(B[q_\lambda]\)
and its critical representative \(V(B[q_\lambda])\).
Let \(u[U]\) be the site map to the block frame
for a field \(U\) about \(U^s\). The chart
argument of the term is
\[
 u[V(B[q_\lambda])]\,u[q_\lambda]^{-1}c_\lambda .
\]
The two site maps convert the frame of the
configuration to that of the slice. They are
evaluated at the actual rotation sites, including
sites inside a block. At the endpoints the
physical-frame identity gives the two original
values with the retained variables held fixed.

At a rotation site of depth \(d_s\), the screened
field and frame estimates bound the log of this
product by the local rotation allowance times
\(C(1+2|\lambda|)e^{-\delta_1 d_s/4}\), with
the available margin in the constant. For a
deep term let \(D_v\) be the depth of its read,
so \(d_s\ge D_v-c\) on each determining stencil.
The core response bound keeps the potential
and current on the admitted slice, and the
site bounds keep the frame in its rotation
domain on the disc
\(\rho_v=e^{\delta_1D_v/8}\). If \(N_v\) bounds
the holomorphic term minus its fixed constant
on that domain, subtract its value at zero
and use Schwarz's inequality at one. The
oscillation is at most
\(2N_v/\rho_v\le 2N_v e^{-\delta_P\widehat D_v/2}\)
in the scaled depth \(\widehat D_v\). When
\(\rho_v=1\) or the term is outside the deep
class, use the two endpoint majorants,
giving \(2N_v\). These are the
two bounds for the oscillation sum at the
component's own healing.

% Earlier Part III, 81734--81854.
For the deep oscillation sum, assign a read to
one of its nearest core-boundary cells. The
depth rate is measured from the full read.
One part of the rate pays for bringing the
unit weight to that boundary, including the
finer cells in its layers. A second part pays
for the sum over read locations relative to
the boundary. This bounds the sum of the
screened oscillations by the boundary count
times the per-cell unit bound. The absolute
values of pieces filed inside the component
are charged at their own interior stage cells.
For those pieces the switch sum costs at most a
constant times \(e^{4d_m(K_v)}\), with \(K_v\)
the core of the stored term. The original carrier
tree length bounds \(d_m(K_v)\), so four units
of the original carrier decay cover this cost.
An arrested term whose label leaves the component
is counted in the exterior sum at its collar.

% Earlier Part III, 117977--117993 and 118083--118248.
\paragraph{The exact datum in the frame of a kept factor.}
For the kept factors whose frame sites are top
block nodes, let \(\Sigma\) be those sites,
\(V(B)\) the critical slice about the real centre,
and \(c\) the frame argument of the kept factor.
In the frame of the real inner variables the
datum is
\[
 [M(V(B))^{[c^{-1}]}]
     =[(e^{iB}D_c)^{[c^{-1}\tau_B]}].
\]
Here \(D_c\) is the real centre datum and
\(\tau_B\) is the node map of the slice.
Brackets identify data related by a node gauge
equal to the identity on \(\Sigma\). A further
frame \(d\) sends \(c\) to \(cd\), since
\(d^{-1}c^{-1}=(cd)^{-1}\). The same slice is
used. The logarithmic polar norms add as an
upper bound, so composition spends the
corresponding part of the frame radius.

At a leaf admitting the frame extension, write
\(c=e^Yw\) with \(w\) unitary. The regauging
map \(\rho[Y]\) takes the value \(e^{-Y}\)
at \(\Sigma\). If \(\widehat w\) is a real
extension of \(w\), the framed value equals
the unframed value at
\((V(B)^{\rho[Y]})^{\widehat w^{-1}}\)
with the original real inner variables.
This follows by cancelling the real rotation
of the variables against the real gauge.
The resulting field is critical, with the
exact datum displayed above. The polar
extension spends \(\|Y\|/(8\varepsilon_{\rm fr})\)
of the slice room, where
\(\varepsilon_{\rm fr}\) is the frame radius.
The real gauge spends none. Thus half of
the frame radius costs at most \(1/16\)
of the slice room. For admitted charts,
equality on overlapping parameter families
follows first at the real values from uniqueness
of the real orbit and invariance of the kept
factor under the datum stabilizer. Holomorphy
of the frame extension on the product of the
slice and frame domains then extends that equality.
The polar representative supplies a pointwise
tube witness. This argument
concerns frames at top nodes. The interior
rotation sites of the preceding estimate need
the full site map there.

For paired centres with common datum on the
rows at \(\Sigma\), use the same \(B\) on
those rows. Denote the additional collar datum
of the finer run by \(b_{\rm col}\). Matching
the exact datum in the two birth frames uses
the finer frame \(h(B)c\), where
\[
 h(B)=
 \left.
 \tau^{\mathsf B}_{B\oplus b_{\rm col}}
       (\tau^{\mathsf A}_B)^{-1}
 \right|_\Sigma .
\]
Indeed \((h(B)c)^{-1}\tau^{\mathsf B}\)
and \(c^{-1}\tau^{\mathsf A}\) agree on
\(\Sigma\). The far ends of the rows are
compared modulo the same trivial-on-\(\Sigma\)
gauges. A common chart then acts by right
multiplication on both frames and leaves
the collar variables alone. To apply this
formula at the mixed complex reads we must
construct the paired bases on the actual
joined rows, choose the finer lift, and fit
both \(c\) and \(h(B)c\) in their frame
budgets on the contour. The matching identity
uses the admitted charts. The extension on
each older factor must also be returned through
the nesting step before that factor is read again.

% Current Part III, 49345--49349 and 49660--49699.
On complete stencils the curvature and the
Wilson current give part of the room estimate
for a polar regauge. Each plaquette is
conjugated at its base point, and the six
plaquettes in a current stencil are transported
to the base of the bond. Consequently the
current is conjugated there in the same level
units. If \(z\) bounds the polar log of the
regauge, the two operator norms of the
regauge and its inverse are at most \(e^z\).
Each of these weighted entries therefore
grows by at most \(e^{2z}\), with a relative
increase bounded by \(2ze^{2z}\).
This controls the terms that transform by
conjugation. The chart potential also reads
the variation of the lifted gauge between
sites. Its first differences and the changes
of the joined boundary rows still have to
fit in the room for the mixed contour.

% Current Part III, Lemma 81.22(b,d) and Theorem 81.2(a,c).
On a matched closed set of determining rows
one can reconstruct a critical field for the
twisted datum locally. Take its physical row
logarithm \(b\) about the admitted real datum.
If the frame is near a nonidentity unitary
section, use the real reference transported
by that section. Convert \(b\) to the reduced
row logarithm \(b_\natural\) using the inverse
of the outward-row correction. Solve the
response equations with input \(b_\natural\)
and all switches at one. Its exact average
is the specified physical datum acted on
by the response's node factor. Removing
this factor by a fixed analytic extension
of the inverse node map gives the required
datum. The response and the local row maps
are holomorphic for small rows, and the
stationary equation extends from the real
data by holomorphy. This reconstruction
gives an exact local representative whenever
the physical row disc lies in the response
radius. The actual paired contour must still
give that disc on the closed rows and leave
the required margins for the potential and
the frame.

% Earlier Part III, 44858--45013 and 45094--45222.
\paragraph{The relative frame at an already admitted later read.}
Let \(Y\) be a live component born at step
\(m\), with padded strip \((Y^\sharp)^+\).
The strip remains in the unreblocked
large-field neighbourhood. Its separation
from the younger determining sets leaves
the read indexed by the birth level \(m\).
Use the actual paired data at this site.
Write \(D_0\) for the first run's datum
with the finer run's collar entries inserted,
\(d_{\rm net}\) for the net phase,
and \(g_{\rm cen}\) for the centre
map in the site's exact representation. Let
\(u^{\mathsf A}\) and \(u^{\mathsf B}\) be
the site frame corrections in the two reads
and \(c^{\mathsf A}\) the first frame.
The path for the second run's kept factor
uses the arguments
\[
 D_s=g_{\rm cen}^{s}e^{is d_{\rm net}}D_0g_{\rm cen}^{-s},
 \qquad c_s=e^{sZ}c^{\mathsf A},
 \qquad e^Z=u^{\mathsf B}(u^{\mathsf A})^{-1}.
\]
In the datum formula \(g_{\rm cen}^{s}\)
and its inverse act at the respective endpoints.
At \(s=0\) the arguments are the mixed
read, and at \(s=1\) they are the
second run's actual read. Hold the real
inner variables of \(Y\) fixed.

The bound for \(Z\) first compares the
fine site map to the one constructed on
the matched coarse inputs. It bounds the
logarithm of their quotient by the local
radius times the relative decay. A local
Lipschitz estimate then replaces the
coarse inputs by those of the actual
first run. This gives the smallness of
\(Z\) in units of the read margin. Use the
centre-map and net-phase bounds in the
same units. At a site with these room
estimates the path stays admitted for
\(|s|<(8C_{\rm read}\vartheta_Y^\sharp)^{-1}\),
where \(\vartheta_Y^\sharp\) is the
relative decay of the read at level \(m\).
Schwarz applied after subtracting the
value at zero bounds the reading change
by \(8C_{\rm read}\vartheta_Y^\sharp\)
times the two-point oscillation of the
kept factor at fixed inner variables.

For a nonempty switch coefficient subtract
the reading change at the empty switch
point. The contour bound of the resulting
difference is twice the oscillation:
each of the two read maps stays in the
admitted datum domain. On the vertices
use the reading estimate for both values.
The geometric mean leaves
\(C(\vartheta_Y^\sharp)^{1/2}\) times the
oscillation, in addition to the switch
coefficient decay. Thus the two kept
factors together have a reading bound
\(C\vartheta_f^{m/2}\check R_m^{p_\Sigma}\)
in their class, by the oscillation sum at
their birth. For the kept quadratic form
the actual datum read is the block
average on the exterior part of the
healing core. It has no independent
frame argument. This later-read
calculation uses the admitted old
contours and their comparison maps.

% Earlier Part III, 51948--52035, 69339--69686,
% 70122--70191 and 71105--71413.
\paragraph{The kept Gaussian at a component's own read.}
For a read contained in the Gaussian core, use the
supported real field \(U^s\) of the prepared determining
set. The moved centre on the padded read has the form
\((e^{i\eta\mathfrak X}U^s)^v\), with \(\mathfrak X\) real
and small in the local response norm. To test its admission
in the older domain, put a real cutoff on \(\mathfrak X\)
which is one on the full clause read and zero outside
the next padded set, still within the core. If the cutoff
width is \(u/8\), then on any older clause cube meeting it,
\(\xi_l|\nabla\chi|\le 8L/M\) since \(\xi_l\le L\xi_j\)
and \(u=M\xi_j\) in the level-\(j\) read. The product rule
preserves the bound on the potential and its first
difference in the clause units. In the real gauge of
\(U^s\) this is within the available radius of the older
clause. Extend the site gauge \(v\) off the padded read
and apply it after the cutoff. The resulting ambient
field agrees with the moved centre on the required
reads. The argument uses the margin from the read
to the artificial strata of the core.

On an aligned label at healing, consider a live
interrupted attempt \(W\) nested in the operated
component. Use the common frozen data and the common
collar variable in the two block-centre frames.
The stored quadratic \(q_W\) is evaluated at that
variable. Its native expansion is a sum of units
\(T_e\) on complete carriers. At an admitted real
vertex its sum is the retained quadratic on the
physical pair. The comparison uses the determining
set refined by the attempt at both vertices.
Add and subtract the second unit at the first read
expressed in the common datum. This separates the
change of the stored coefficient from the change of
its read. The reading difference vanishes if the two
data agree. For a positive displacement \(d_e\),
the line in the native unit tube has radius \(r_G/(4d_e)\).
With unit majorant \(G_e\), Schwarz's inequality bounds
the reading part at the vertex by \(4d_eG_e/r_G\).
The extension in the switch parameters uses the
geometric mean of this bound and the one-run majorant.
Together with the stored mark, the resulting size is
at most \(C\varpi(e)\Delta_{n_e}^c\), where \(\Delta_n^c\)
is the classical load of zone \(n\) and \(\varpi(e)\) is
the native unit weight.

The overlap bound sums these weights by the visited
collar cells. For all attempts nested in \(C\), at a fixed
zone \(n\), the absolute vertex differences thus sum to
at most \(C N_n^{\rm vis}(C)\Delta_n^c\), with
\(N_n^{\rm vis}(C)\) the number of visited cells of that
zone. Multiplying by a zone parameter of radius
\(E_1/\Delta_n\), with \(\Delta_n^c\le\Delta_n\), costs
at most \(CE_1\) per such cell. The rooted sum for the
analytic units also retains the spatial weights.
Their two tree lengths are bounded by a fixed multiple
of the number of constituent cells. Spend the matching
part of the native unit decay on these weights and sum
the remaining overlap weights at the chosen root.
This gives a bound \(C\Delta_n^c\) at that root.

% Earlier Part III, 5418--5457, 5777--5830,
% 6207--6310 and 6504--6737.
\paragraph{Composed reads through the preparatory stages.}
In one run let a chain enter at stage \(n_e\)
and descend through stages to a stored read at
\(\nu\le n_e\). Use the fixed real centre \(U_c\)
of the entry. On the padded leaf hull \(X^+\)
write \(\mathbf U=e^{i\eta A'}U_c\). Extract
the datum \(B[\mathbf U]\) by the chart's
nonlinear averaging map on the induced box
set \(\mathfrak S\). For this datum solve in the
Landau chart at \(U_c\) with its own response
\(H^\dagger J(U_c)\), obtaining \(V(B[\mathbf U])\).
The stored term \(F(X;\cdot)\) is evaluated at
\[
 S(\mathbf U)=
 \big(V(B[\mathbf U]),J(V(B[\mathbf U]))\big)|_X .
\]
The current in this pair is recomputed on the
slice. At real points the chart represents the
native leaf datum up to gauge. On the complex
disc the slice is its holomorphic continuation.
With no moves, the entry \(U_e\) is already
critical on the induced box set by restriction
from the entry chart. Uniqueness in that chart
gives \(S(U_e)=(U_e,J(U_e))|_X\).

For the radius check let \(\mathsf T_i^a\) be
the stage \(i\) bound on member \(E_a\) of the
field and current table. Put
\(\mathrm{rm}_i^a=\mathsf T_i^a-\mathsf T_{i+1}^a\)
and \(G^a=\mathsf T_\nu^a-\mathsf T_{n_e}^a\).
The stage-room bounds give \(\mathrm{rm}_i^a\ge0\).
Let \(a_i^0\) be the size share of the move at
stage \(i\), with the local stencil suprema
included by the slowness bounds. The classical
entry condition at the stage bounds
\(\sup(|h_i|+\ell_i|\nabla h_i|)\le a_i^0/K_o\)
on its local read. The
composite logarithm \(\mathcal Y\) of the moves
over \(U_e\) then has local size at most
\(2\ell\sum_i a_i^0/K_o\) by the small-product
estimate, with \(\ell\) the slice scale.
On the admitted chart, the change of the
box datum on each cube is bounded by four
times the supremum of \(|\mathcal Y|\)
on that cube. This follows from
the linear average bound and the remainder
at \(U_e\), using the classical ceiling
\(C_w|A_e|\le2\) there, where
\(U_e=e^{i\eta A_e}U_c\).
The local Landau response bounds and slowness
therefore control the change of the member with
scale power \(e_a\) by
\[
 \frac{16 B_{\natural}C_{\rm sl}}{K_o}\,
 \ell^{1-e_a}\sum_{\nu\le i<n_e}a_i^0.
\]
Here \(e_a=1,2,3\) for the potential, its
first difference and the current respectively.
For the plaquette member use the gradient
bound with its extra \(\eta^2\) factor. The
real-centre member and the cube condition
are inherited from the entry.

At a point in the units of the read, let \(i_1\)
be the first stage that raises its level, going
from \(\nu\) toward \(n_e\). Before \(i_1\)
the converted share at stage \(i\) is at most
\(80\mathsf r_1\,\mathrm{rm}_i^a\), where
\(\mathsf r_1=C_1/C_0\) is the ratio in the
stage table. The later shares cannot be charged
to their own stage rooms in the finer read units.
Instead the share sum from \(i_1\) onward is at
most \(C_\Sigma\bar u^0_{i_1}\), with \(\bar u^a_i\)
the local stage radii. The units at the first
raising stage are still the read units, and
\(\bar u^a_{i_1}\le40\,\mathrm{rm}_{i_1}^a\).
Charge this whole part to the first raising
room. The two parts spend disjoint rooms,
so the total converted share is at most
\(400\mathsf r_1G^a\). Without a raising stage
only the first bound is needed. The fixed
choice \(K_o\ge2^{18}B_{\natural}C_{\rm sl}\mathsf r_1\)
then makes the displacement at most \(G^a/8\).
Adding the entry size gives
\[
 E_a(S(\mathbf U))\le \mathsf T_{n_e}^a+G^a/8
                     =\mathsf T_\nu^a-7G^a/8.
\]

Where the point has been raised beyond
the leaf's own read level \(n_F(x)\),
convert the read table to that scale. If
the level \(n_0(x)\) before the chain exceeds
\(n_F=n_F(x)\), the factor in this
comparison is bounded by
\(K_u L^{-u}3^{-v}\sqrt{u+v}\le K_u/L\),
where \(u=n_0-n_F\ge1\) and \(v\ge0\)
is the excess of the current label level over
\(n_0\). At \(n_0=n_F\) the stage-table
bound \((9/4)3^{-v}\sqrt v\le3/4\) for
\(v\ge1\) gives the same allowance.
In the leaf's units the entry spends at most
\(3/4\) of its radius. For the displacement,
the conversion gains at least one power
\(\ell_{n_F}/\ell\), and the share sum at the
read stage is at most \(C_\Sigma\bar u_\nu^0\).
The unit comparison and the choice of \(K_o\)
then bound the displacement by \(1/200\)
of the leaf's radius. Thus the total is at
most \(3/4+1/200<4/5-1/40\). On the
label-equal part, use the scheduled table
containment
\(\mathsf T_\nu^a\le(c_F^+-\varrho_F)\mathsf u_F^a\),
where \(c_F^+\mathsf u_F^a\) is the leaf's
full radius and \(\varrho_F\mathsf u_F^a\)
is its reserve. Fresh terms of the top layer
of the healing step and terms left from a
preparatory chain use their separate containment
clauses, with the latter's own reserve. For
the composite plaquette and
current members with index at most the
birth index, restriction of the slice chart
identifies them with its own plaquette and
current members. These checks give the
one-run stored read on the admitted
preparatory chain. The paired Gaussian
hull comparison has its own admission.

% Earlier Part III, 7198--7278 and 7280--7475.
\paragraph{The common datum at a preparatory vertex.}
Take two aligned native vertex reads over a
preparatory seed. In run \(\mathsf B\) let \(v_0\)
be the common datum with its prescribed
\(\mathsf B\) completion. Write
\(v_s=e^{is d_{\rm net}}v_0\) for the interpolation
to the actual datum \(v_1\), after removing
the real block-centre frame factors from the
two ends of the read. The logarithm
\(d_{\rm net}\) is of the resulting relative
group product. In the fixed chart the
response bound controls the change from \(v_1\)
in each member by \(|1-s|D^a\), where
\(D^a=2B_{\natural}\ell^{-e_a}(|d_{\rm net}|)^\approx\)
uses the exponentially smeared modulus of
the mismatch. The class of the read is governed
by \(q=\vartheta^{(j\wedge n_X)/2}\), where \(j\)
is the formation level and \(n_X\) is a lower
bound for the own read levels on \(X\).

At a point where the label level equals
the leaf's own read level, split the
mismatch into cells where a move acts and cells
of the entry absent from the common data.
For a move cell the level \(l\) is the
label level and satisfies \(l\le j\).
The noncommon cells have \(l=j\) below
the label or \(l\) below both \(j\) and
the label. On a move cell of level \(l\) the share bound
is \(C\vartheta^{l/2}\) times its size share.
Here \(l\ge j\wedge n_X\), so the preceding
first-raising sum bounds this contribution
to \(D^a\) by \(CqG^a\). At a noncommon
cell the mismatch is at most \(C\vartheta^{l/2}\)
times its local datum unit. For the smeared
sum, split \(e^{-2\delta_1 d}\) at distance \(d\)
into rates \(\delta_1,\delta_1/4,\delta_1/4,
\delta_1/2\). The first converts the zone
factor \(\vartheta^{l/2}\) to \(Cq\) by
screening. The second pays the change of local
radii, the third yields the depth factor
\(C2^{-(j-n_F)}\) at read level \(n_F\), and
the fourth sums the rooted kernel. The depth
factor is proportional to the leaf's tube
fraction \(\varrho_F\), so these cells cost
\(Cq\varrho_F\mathsf u_F^a\), with \(\mathsf u_F^a\)
the leaf's member unit. At a raised point
convert to the earlier leaf units as above,
using the screened and smeared bound with
the depth factor omitted. With the constant
of the comparison enlarged
once, the two estimates are
\[
 D^a\le C_{TL}q\,(G^a+\varrho_F\mathsf u_F^a)
 \quad\text{at a label-equal point},\qquad
 D^a\le (C_{TL}/10)q\,\mathsf u_F^a
 \quad\text{at a raised point}.
\]
Use the datum comparison in the range
\(8C_{TL}q\le1\). This follows from the
coupling ceiling \(8C_{TL}\vartheta^{n_c/2}\le1\)
when both the formation level and the
chosen lower read level are at least \(n_c\).
Take \(R=(8C_{TL}q)^{-1}\ge1\). For \(|s|\le R\),
\(|1-s|\le2R\), so the first estimate spends at
most a quarter of \(G^a+\varrho_F\mathsf u_F^a\)
and the second spends at most \(\mathsf u_F^a/40\).
The one-run slice margins admit these
displacements in the quarter-room tube.
At a raised point the largest coefficient
used is \(3/4+1/200+1/40\), below the
quarter-room allowance \(4/5-3/(4\cdot40)\).
At \(v_0\) the factor \(|1-s|\) is one,
which gives the smaller eighth-room tube.
The datum radius itself is
checked before solving for the slice. On the
noncommon entries the local datum unit carries
the scale factor \(L^{l-\iota}\) relative to
the current level \(\iota\). The comparison
floor gives \(\vartheta^{l/2}L^{l-\iota}
\le\vartheta^{\iota/2}\le q\); on the move cells
use the share bound. Thus \(R\sup|d_{\rm net}|\)
is controlled by the input datum ceiling and
the fixed choice of \(C_{TL}\). Together with
the baseline datum budget this puts the full
path in the chart's datum ball, as required
for the box response estimate.

On this disc the difference
\(\widehat F^{\mathsf B}(v_s)-\widehat F^{\mathsf B}(v_0)\)
is holomorphic, vanishes at zero, and is bounded
by \(2\sup|\widehat F^{\mathsf B}|\) on the
admitted tube. Schwarz's lemma at \(s=1\)
gives \(16C_{TL}q\sup|\widehat F^{\mathsf B}|\).
This supplies the reading comparison on the
native preparatory vertices, using the native
share and noncommon datum estimates. For
a family with extra strip letters the same
argument requires these datum bounds on that
family.

% Earlier Part III, 84573--84995, and Earlier Part I,
% 80239--80353.
\paragraph{The kept-term sum in a preparatory integration.}
At preparatory stage \(n+1\), take a native
kept term \(v\) of a live region \(Y\) born
at \(m<k\). Its pointed and window slots are
empty. Its field read uses the padded strip
hull \((Y^\sharp)^+\). The active response
blocks are those whose enlargements meet
this hull. At zero block parameters the
move is the identity, so put this value of
\(v\) in the base. Expand the remaining
difference first in the nonempty active
block sets \(A\), then in cell switches for
the box at fixed \(A\).

For the cell expansion, leave the cells of
the enlarged core off the mandatory core
at switch value one. Partition the other
cells into grains, which are unions of
cells. Give all cells of a grain one switch
parameter. The restricted family is holomorphic
on the grain polydisc. On a zero face, deleting
the grain one cell at a time uses the
one-machine zero-face identity. At each grain
vertex the value is the sum of cell coefficients
whose selected cells lie in the enabled
grains and the fixed core cells. Sorting
by the set of grains met, then inverting,
shows that each grain coefficient is the
sum of all cell coefficients meeting exactly
those grains. The subset of fixed core cells
in this sum is unrestricted.

Use Cauchy's estimate on this grain coefficient
before taking absolute values of the cell
pieces. If \(\mathcal Y\) is the selected
grain set, its bound is the oscillation of
\(v\) on its admitted read, times
\[
 e^{-\kappa_1^{(1)}|\mathcal Y|}
       \prod_{\square\in A}(r_\square-1)^{-1}.
\]
Here \(r_\square=c_R/\vartheta(\square)\ge2\),
so each inverse radius costs at most
\(2\vartheta(\square)/c_R\). Piece locality
requires every wall component of the selected
grains to touch the enlarged core. This uses
the preparatory locality of the full padded
read of \(v\). The coefficient also reads
the moved base on the enlarged core and
selected grains and the stored variable
\(\Phi_Y\).

Let \(R_n\) be the moved region at the stage,
\(w=M/M_1\), and \(S_v\) the native one-cell
anchor of the strip. Use level \(h+1\) for its
corridors in a treated component and level
\(k\) in an untreated one. The native corridor
from a block \(\square\) to \(S_v\) has
fewer than \(d(x_\square+1)\) cubes, where
\(x_\square\) is the distance in its corridor
units. The treated and untreated corridor
bounds give \(x_\square+1\le110d^\wedge(\square,R_n)/w\).
Thus the length weight spends at most
\(\delta_P d^\wedge(\square,R_n)/4\) when
\(\delta_P w/4\ge110d(a_k+1)\).
The inverse block radius decays at rate
\(\delta_P/2\) in this distance. After paying the
corridor and the fixed flat-format costs,
the block weight \(y_{v,\square}\) is bounded
by \(K_2^\flat\vartheta e^{-\delta_P d^\wedge(\square,R_n)/4}\).
The weighted grain activity is bounded by
\(x^U=\exp[-\kappa_1^{(1)}+(a_k+1)+3^dc_\flat]\).
The flat grain floor
\(2e^2\widehat b_t^2 N^\flat x^U\le1\),
with \(\widehat b_t\) the adjacency bound
and \(N^\flat\) the core size bound per block,
gives the shape sum at most \(e^{|A|}\).
The factor \(e^{-1}\) per block reserved
by the flat-format bound pays this cost.
Put \(\Sigma_v=\sum_\square y_{v,\square}\)
over the blocks of the full read. The
confined or two-box block count gives
\(\Sigma_v\le2c_{\rm p}K_2^\flat\vartheta V'\),
where \(V'\) bounds the contour block sum
for a single component of \(R_n\). The
coupling ceiling gives
\(\Sigma_v\le1/(3N)\le1/2\), where \(N\) is
the number of young levels of the step.
The remaining block subset sum at a
specified anchor \(\square_0\) costs at most
\(e^{\Sigma_v}y_{v,\square_0}\).

At a target filing cell, the anchor can
reach the cell through its own core or
corridor, or through an attached wall
component of grains. In the first case
the cell is within \(103\check R_k\) of the
block in level-\(k\) units. In the second
case the positional grain bound supplies
the factor \(e^{-(\operatorname{dist}_k-103\check R_k)_+}\).
This gives the same position factor for
both cases. For the other root sum, a cell
meeting \(S_v\) is within that distance of
every eligible block of \(v\), and the
nonempty subset sum is bounded by
\(\Sigma_v e^{\Sigma_v}\).

The near-strip geometry gives
\(d^\wedge(\square,R_n)\ge w(5\check R_m/2-6)
\ge w(\check R_k-5)\) for every eligible
block. In an own chain it uses the
screened collar of \(Y^\sharp\); at \(m=h\)
use the width in level-\((h+1)\) units and
\(L\check R_{h+1}\ge\check R_h\). In a foreign
chain the collar is inside the untreated
component. The hull count gives at most
\(2^d\) regions at the same birth level per
block. Their polynomial oscillation bound
is absorbed at rate \(\delta_P/16\)
in \(d^\wedge\), leaving rate \(3\delta_P/16\)
in the block sum. Use the uniform bound for
\(y_{v,\square}\) before summing over the leaves
at the block, since their corridors can differ.

This last sum must use the ordinary old
blocks together with all blocks meeting the
padded kept hulls. The old-block summation
proof extends to this enlarged class.
Every added block lies in the same cover
and has the same lower bound
\(w(\check R_k-5)\) on its distance to
\(R_n\), hence to each component of \(R_n\).
In the component sum at rate \(\delta_P/8\),
half the exponent pays \(V'\). The near
component shells then gain \(e^{-(\check R_k-5)}\),
and the shells at distance \(l\check R_k\)
gain \(e^{-l\check R_k}\). The outer shells
around the target cell gain the position
factor from the anchor bound. The original
polynomial shell counts and tail floor
therefore bound the enlarged block sum
by the same \(V'\). The kept contribution
to each of the two rooted one-run sums
is at most \(C K_2^\flat\vartheta V'\)
times the kept oscillation constant. This
uses the native one-run tube and
preparatory piece locality.

% Earlier Part III, 67078--67204, 119152--119161,
% 119445--119568.
At a preparatory stage use the same unit identity
after admitting the stage's endpoint maps in the
local Gaussian tube. On each complete hull
align the parameter-independent axial centres
\(U_c^{\mathsf A},U_c^{\mathsf B}\) by a real gauge
\(\Gamma\), so that their coarse comparison is
\(\mathcal M((U_c^{\mathsf B})^\Gamma)
=e^{i\eta a_c}U_c^{\mathsf A}\).
The paired chart gives the corresponding bounds
for the centre displacement, the gauge at the
datum points and the displacement of the two chart
fields on the parameter disc. In the local length
units each is bounded by the local radius times the
chart decay. When comparing the actual bond datum,
take these bounds on both endpoint top blocks and on
the smeared averaging contour. In addition to the
chart decay there is the error \(L^{-m}\) of the
averaging contour at separation \(m\) in levels.
At the preparatory site both errors are bounded by
\(C\vartheta_2^{n_e/2}\) for the lowest zone \(n_e\)
of the unit hull. The affine estimate keeps the
radius factor, since it occurs in both bounds before
the geometric mean, and gives decay
\(\vartheta_2^{n_e/4}\).
The weighted pile over a finite set of cells
\(\mathcal A\) is at most
\(C|\mathcal A|\Pi_G^+(x_{k'})\Delta_n^c\)
for the class \((n,k')\).
Apply this to the cells meeting the actual
preparatory block and use its block count. This
gives the stage load
\(C\Pi_G^{\rm blk}(x_{k'})\Delta_n^c\), with
\(\Pi_G^{\rm blk}\) including that count.

% Earlier Part I, 118810--118897, 119037--119113,
% 119191--119254 and 120375--120381.
The small factor for a block uses the count
and the unit's local radius together. For
each unit meeting the padded response block,
choose a common site in its carrier and a
current cell at that site. The cell count
of \(\mathcal A\) is bounded by the square
of the window bound at this current level,
by the spread condition of the response block.
The current level and the level of the
largest chart radius of the unit differ
by at most \(N_e+3\), where \(N_e\) is two
plus its number of collar cells. The flow
then bounds both the window size and the
old collar polylogarithm by functions of the
coupling at the largest-radius cell, with
an increment \(O(N_e+1)\) in inverse
squared coupling. If \(\Theta_e\) is that
radius and \(\varrho_e^2\) the old
polylogarithm, the reserve \(e^{-N_e}\)
absorbs the extra dependence on \(N_e\)
and gives
\[
 |\mathcal A|\varrho_e^2\Theta_e e^{-N_e}
 \le C\Pi_G^{\rm blk}(x_{k'}).
\]
Use the factor \(1/|\mathcal A|\) from
this bound against the positive unit pile
over the cells of \(\mathcal A\).
For a nonempty response coefficient,
subtract the value of the unit at zero
before applying the anchored Cauchy
estimate. The mixed derivative kills that
constant. Its rim-anchor sets are empty
because the retained collar variable is
held fixed in this integration.

% Earlier Part III, 120215--120289.
The layer bound in the first comparison
uses the actual chain vertices at unit
dilation and the cut vertices with dilation
parameter between zero and one. In the
product rule for these reads the dilation
member can be evaluated at unit radius.
It costs at most \(16c_H\sup_p m_p\),
where \(m_p\) is the amplitude of the
source at the row \(p\) and \(c_H\) is the
product-rule constant. The source bound
gives \(\sup_p m_p\le C\alpha_{0,m}\).
This estimates the member directly from
the amplitude. On the outer layer the
available row room can include the
large window weight. That weight remains
in the room for the other members of the
layer estimate, with the coupling of
the same window.

These estimates concern the stored quadratic,
in its own collar variables. The new Gaussian
integral and its determinant have their separate
comparison. At healing, for each unit, machine,
filing domain and real hull centre, we must give
holomorphic bond configurations on the full padded
hull for all subsets of its switches. The
configurations must agree bond by bond when a
switch is set to zero and its membership in the
subset is changed. At the binary vertices they
must recover the physical axial reads in the
fixed chart. If the far read sees no operated
component, the unit must have its root value in
the corresponding source frame. The family must
stay in the hull tube and satisfy the deep spatial
bound. For far units one can use the weaker
contour estimate below. The vertex
comparison and the rooted estimates apply to
the unit values in this assembly.

% Earlier Part III, 87420--87432, 87517--87647,
% and 118851--118870.
For the all-one vertex take a fixed wall-connected
box \(G\) containing the complete read and the full
enlargements of the operated components it meets.
Use the restriction rule for the determining sets
and the matched frames at the glued seam. Put the
four stages of the construction on this same box.
The starting field is the restriction of the root.
On the outer rows the shell restriction
gives identical rows and levels in the successive
determining sets. Each increment and the intervening
gluing are supported away from these block pairs.
Consequently the next stage inherits the root
average on these rows. On the inner rows use the
prescribed prepared data and the frozen data on
their respective parts. At the glued stage the cut
data are the identity on one side and averages of
the matched top data on the other. Composition of
the averages gives the prescribed coarse rows there.
Thus the last stage has exactly the native datum of
the single box. At an admitted datum, uniqueness of
the critical orbit identifies the last field in the
fixed axial frame. If there are no operated components
in the box, there are only outer rows, and the last
field and the restricted root are in the same orbit.
For the Gaussian value this comparison uses the
retained variables transported in the same gauge.

% Earlier Part III, 118336--118339 and 118388--118512.
When the one-run and comparison presentations
are used at a natural vertex in the same cutoff,
their physical boundary functions must agree
at a fixed collar variable. Consider the two
admitted presentations as continuous families
on an open set of real core data, with the same
core datum and absolutely axial endpoints.
On a dense set of those data the endpoints
belong to one orbit at real parameters.
Remove the two accumulated frame maps to
obtain the relating gauge of the physical reads.
Axial transport makes this gauge constant
on each core block. On the graph of core data
rows the gauge at one end of a row is transported
by that row to the other. At a reference node
the gauge must commute with every loop holonomy.
If two core-loop holonomies are free coordinates
of the datum, their noncommuting values are dense.
For such a pair in \(\mathrm{SU}(2)\) the common
centralizer is \(\{1,-1\}\). The gauge is then the
same central element on the blocks containing
the hull base points, so it fixes the hull datum
and the collar variable. Continuity in the real
core data gives equality also at reducible
data. Holomorphy in the presentation parameters
then gives equality on the common parameter
disc at this vertex. This use of continuity
requires two admitted families over an open
set. At an isolated reducible datum a stabilizer
can rotate the collar variable.

% Earlier Part III, 117244--117375.
The chart source of a link is fixed by the handed
increment in the physical frame. If \(\mathfrak Q_{s-1}\)
is the accumulated frame map before the link, and
the chart updates the average at \(c\) by \(e^{ib_s(c)}\),
frame covariance conjugates that factor by
\(\mathfrak Q_{s-1}(c_-)^{-1}\) in the physical average.
To realize the prescribed increment \(\beta_s(c)\)
one therefore uses
\(b_s(c)=\operatorname{Ad}_{\mathfrak Q_{s-1}(c_-)}\beta_s(c)\).
This is a recursive holomorphic assignment on the
admitted tube, since the frame is known before
the link. It preserves the size at real data,
and the complex frame bound multiplies the
size by at most \(e^{2\cdot1.07 K_u\varrho_\Pi}\).
A row set to the previous object's own average
has zero increment in both frames. For a kept
Gaussian the increment may also enter its
quadratic in the carrier frame, so its literal
value must be kept in that quadratic.

% Earlier Part III, 115716--115795 and 115972--116110.
For the zero-face check, the local fixed-point
comparison uses switches on enlarged walk carriers.
In top-level units give a switch cell \(\Delta\)
a halo \(\mathcal N(\Delta)\) of width \(3LM_1\).
Multiply a walk and its transpose by the product
of the switches whose halos meet its enlarged
localization set, including switches of cells
outside the enabled box. A shell row of level
\(n\) depends on the box within \(3s_n\) of
the row, and \(3s_n\le3LM_1\). Thus it is the
halo condition that removes walks whose rows
could change when a cell is deleted.

If the walk has contour tree length \(\widehat d\),
the geometry floor places all its switch cells
within \(M/4\) of the tree. There are at most
\(C_d(1+\widehat d/M)\) such cells by the
covering count. For \(|s(\Delta)|\le e^{\kappa_1}\)
the product costs at most
\(e^{C_d\kappa_1}e^{\delta_0\widehat d/2}\)
under \(\delta_0M\ge2C_d\kappa_1\).
This leaves half of the original walk decay.

For one native response link with a fixed base,
set a switch to zero and compare the two boxes.
The surviving walk words have the same
intermediate local stencils, levels and base
values off the halo. The first kernel and its
local output stencil kill the output on the halo.
The last switched kernel reads only input
stencils away from the halo. Where the prescribed
data agree on those stencils, the two systems
therefore have the same equations on the common
part and the same weighted uniqueness ball for
the response extended by zero. For the ordinary
slot system this applies with the slot source on
core rows and the prescribed current slot of the
fixed base.
For the sourced slot system the two extra terms
are
\[
 -\mathcal G_0(1-\theta)Q^*\theta_1
 H_{\boldsymbol\Omega_1(W)}^\dagger J_1,\qquad
 -\mathcal G_0(\theta J_{1,2}),
\]
where \(W\) is the box of the link. Switch the
cut-geometry operator \(H^\dagger\) and the
outer kernel \(\mathcal G_0\) by the same rule.
The two \(J\) sources are fixed in the core, and
the operators between the switched kernels read
only local stencils. This makes the sourced terms
obey the same zero-face restriction. Rim data read
from a fixed configuration agree on the common
stencils as well. Uniqueness gives equality of
the two outputs on the common bonds and the base
on the halo. To apply this check to the
Gaussian family, we must verify those walk and row
identities for the boxes generated by the unit's
switch subsets, including the junction data, and
identify the reconstructed bonds in the fixed
physical frame on the full padded read.

% Earlier Part III, 58375--58512, 119266--119305,
% 119574--119605, 119733--119787 and 119873--119910.
\paragraph{The exterior sum of the Gaussian units.}
At healing, form the core \(K_e\) of a unit by padding
its hull by one layer of level-\(m\) cells and adjoining
the whole protected blocks it touches. For a set of
enabled exterior cells \(J\), read in the component
\(G_e(J)\) of this core after adding \(J\) and the
protected blocks joined to it. A unit is far if its
padded hull avoids the enlarged operated region.
For a far unit subtract the value at the root before
taking the alternating sum. This makes the vertex
value zero when the box sees no operated component.
For a deep unit use the full value, with its hull inside
the collar of one treated component. The other units
are shallow. The coefficient of a set \(B\) of cells is
the alternating sum of the values for \(J\subset B\).
The sum over \(B\) recovers the full vertex value,
with the root subtraction for a far unit.
If a cell of \(B\) is outside \(G_e(B)\), pairing subsets
which differ at that cell gives zero. Thus each
connected part of \(B\) must attach to the core,
counting cells sharing a protected block as adjacent.

At an admitted far read let \(b_e\) bound the
displacement of the hull datum from its root value
at every vertex and let \(r_e\) be the local chart
radius on the hull. To estimate \(b_e\), compare the
exact average of the last field with that of the
root, restoring both endpoint frames. The frame
corrections are products of the factors for the
intervening layers. Their logarithms are bounded
by the sum of the layer increments. The two
averages can be compared on the disc of the same
real centre, where Cauchy's estimate gives the
corresponding Lipschitz bound for their logarithms.
If the increments at the switch vertices decay
away from the enlarged operated blocks \(C^\sharp\),
uniformly on the admitted tube of hull data, this gives
\[
 b_e\le C\alpha_{0,m}\sum_C
 e^{-\delta_P d^\wedge(X_e^+,C^\sharp)} .
\]
Here \(\alpha_{0,m}\) is the step amplitude and
\(d^\wedge\) is distance measured along paths in
the local scale units. On the complex polydisc it
suffices that the total size of the increments is
bounded by the chart radius. Both estimates are
on the same switch family.

For a function of that family bounded by \(V\) on
the vertices and by \(H\) on the polydisc of radius
\(\rho_1\), the alternating coefficient for \(B\)
is bounded by
\[
 \left(\frac{4}{\rho_1-1}\right)^{|B|/2}(VH)^{1/2}.
\]
Indeed the vertex sum gives \(2^{|B|}V\).
Face agreement lets one replace the vertex value
in the alternating sum by its integrated mixed
derivative in all the switches of \(B\): the
terms on a zero face cancel in pairs. Cauchy's
estimate gives \((2/(\rho_1-1))^{|B|}H\).
Switches outside \(B\) stay at one.
Take the geometric mean of these two bounds.

Write \(\nu_e\) for the majorant of the unit and
\(\nu_e^\times\) for that of its common-datum
difference. With the fixed tube constants absorbed,
the one-run difference from the root has vertex
bound \(C\nu_e b_e\) and contour bound
\(C\nu_e r_e\). The change of the common-datum
part has the same pair of bounds with \(\nu_e\)
replaced by \(\nu_e^\times L^{-\theta j_e/4}\),
where \(j_e\) is the level of the arrest.
For the reading part of the two-run difference
from the root, use the better of the separate
one-run vertex estimates and the direct vertex
comparison. With \(n_{X_e}\) the lowest zone on
the hull this gives a vertex bound
\[
 C\nu_e\min\{b_e,\ r_e\vartheta_2^{n_{X_e}/2}\}
 \le C\nu_e(b_e r_e)^{1/2}\vartheta_2^{n_{X_e}/4}.
\]
Its contour bound is \(C\nu_e r_e\). Thus in
addition to the common switch factor the reading
coefficient has size at most
\(C\nu_e r_e^{3/4} b_e^{1/4}\vartheta_2^{n_{X_e}/8}\).
The other two coefficients have the factor
\((b_e r_e)^{1/2}\). The quarter-power loss in the
reading part therefore affects the spatial and
coupling gain as well as the zone rate. The
comparison uses the zone rate bound
\(\vartheta_2^{n/8}\le\vartheta_f^{n/2}
\le C\Delta_n^c\) for the reading class
\(n\le n_{X_e}\), with \(\vartheta_f\) the decay
rate at the later reading sites. Keep the
common-datum part in its own class, at \(j_e\)
or at the level of its filing cell.

For the sum at an exterior cell \(Q\), file the
piece for \((e,B)\) on the union of the core, the
cells of \(B\), the protected blocks they touch
and their attached treated components. Weight
it by \(e^{a d_{m,Z}}\), where \(a=31\kappa/20\)
and \(d_{m,Z}\) is its tree length outside the
large-field region. If the coefficient
is nonzero, this label connects \(K_e\) to \(Q\).
The length is at most \(C(N_e+|B|)\), where
\(N_e\) is two plus the number of collar cells
of the unit. Its ordinary length is at least
\(\operatorname{dist}_m(Q,K_e)-1\). Bound the
relative length by the ordinary one and use an
extra unit of ordinary length weight to factor
out \(e^{1-\operatorname{dist}_m(Q,K_e)}\),
with distance in level-\(m\) units. If \(c_l\)
is the length allowance per switch cell, take
the coefficient rate \(\kappa_1^\natural\) so
that \(\kappa_1^\natural-(a+1)c_l\) is at least
the rate required to sum connected families
at the boundary of the core. This sum costs at most
\(e^{CN_e}\), since the number of core cells is
\(O(N_e)\) and the adjacency has bounded degree.
The unit weights have a corresponding reserve
in \(N_e\). What remains to sum is the unit
weight times its spatial profile and the factor
\(e^{-\operatorname{dist}_m(Q,K_e)}\).

For each \(C\) in the far profile, use one fraction
of its spatial decay to continue that last distance
to \(C^\sharp\). Use a second fraction to move the
decay from the hull to the filing cell of \(e\).
The hull has \(O(N_e)\) touching cells, so this
uses another fixed exponential in \(N_e\).
The remaining weight at each filing cell is
summable over all live arrests at once. The
distance sum over filing cells is bounded by a
constant times the number of exterior boundary
cells of \(C^\sharp\). For this step the distance
rate times the width of a cell in its own units
must be large enough to pay for subdivision
across levels. Finally the separated blocks
satisfy
\(\sum_Ce^{-\operatorname{dist}_m(Q,C^\sharp)}
\le C\), by counting them in shells around \(Q\).
The boundary count costs at most a fixed power
of the separation scale.

The unit weight contains a polynomial in the
logarithm of the inverse squared coupling at
\(j_e\). For the common-datum comparison it
is absorbed in the classical load. For the
one-run estimate and the weaker reading estimate
use a further fraction of the spatial decay:
the path from the hull to the boundary crosses
the intermediate zones, giving exponential
decay in \((m-j_e-2)_+\) after the hull-size
cost. The flow bounds the logarithm at \(j_e\)
by its value at \(m\) plus \(C\log(2+m-j_e)\).
Thus the cost after summation is polynomial in
the logarithm at \(m\). Denote that cost by \(P_m\)
and the common upper bound on the boundary-cell
counts by \(N_{\partial,m}\). In the exterior norm
(the sum of the weighted suprema over the labels
containing \(Q\), followed by a supremum over
exterior \(Q\)), the far one-run bound is
\(C N_{\partial,m}P_m\alpha_{0,m}^{1/2}\).
The two-run bound for class \(n\) is
\[
 C N_{\partial,m}
 \bigl(\alpha_{0,m}^{1/2}
       +P_m\alpha_{0,m}^{1/4}\bigr)\Delta_n^c .
\]
There are enough spatial fractions in the weakest
rate \(\delta_P/4\): the cell-summation rate may
be chosen at most \(\delta_P/16\) in the
scaled distance.

% Earlier Part III, 119813--119860.
For the deep units, use a fixed vertex as the
reference. The reference has zero alternating
coefficient for nonempty \(B\). Only those \(B\)
enter the exterior norm because the core is
inside the operated region. On a deep hull use
the admitted contour bound on the datum increment
from this vertex. The preceding comparison then
uses the full contour gain
\(C\alpha_{0,m}e^{-\delta_P d^\wedge(X_e^+,\mathsf C_C^c)}\)
in place of \(b_e\), where \(\mathsf C_C\) is
the collar of \(C\). Cauchy's bound for one
run and the vertex/contour bound for the reading
difference retain, respectively, \(\alpha_{0,m}\)
and \(\alpha_{0,m}^{1/2}\). Sum at the exterior
boundary of the collar in its own level units.

The shallow units use the tube majorant even
when their hulls meet the boundary. If the arrest
is inside a treated collar but its hull reaches
outside, the collar-cell carrier reaches across the
collar margin. If the arrest is outside all
treated components but its hull meets an enlarged
one, the carrier spans the separation from that
arrest to the enlargement. Either way,
\(N_e\ge c\check R_m-C\), where \(\check R_m\)
is the separation scale. Spend the spare
exponential in \(N_e\) to get
\(Ce^{-c'\check R_m}\); in the one-run sum reserve
part of it for the polylogarithm at the arrest
level, since the carrier has a cell for each
intervening level. The same boundary sum now
applies with this carrier gain in place of the
small spatial amplitude.

All three estimates are small after the boundary
and logarithmic costs. For the two-run direction,
the dial radius cancels the classical load against
the full load. These are the exterior Gaussian
contributions to the terminal comparison.

% Earlier Part III, 85382--85522, 116250--116453,
% and 116850--117048.
\paragraph{Box pieces for a term at its own healing.}
Consider an interior unit \(v\) at the healing site
\(\mathbf L\). By the pocket-extinction rule, each
of its pockets dies inside the integration region
at this step. Use its full fat core \(K_v\) in
every box of the four links, along with the
mandatory hubs \(H_C^{(t)}\) of link \(t\). The
eligible cells at link \(t\) are the cells of its
partition that meet neither the interior of any
\(H_C^{(t)}\) nor \(K_v^\circ\). Thus a switch never
cuts the core. At zero switch on \(\Delta\), the
one-machine locality identity identifies the
restricted response and averaging operators on
\(K_v^\circ\) for boxes differing by \(\Delta\).
The moved point reads the operators through a
function supported in \(K_v^\circ\), so it agrees
for the two boxes. This applies to a shallow
interior term as well as a deep one. A kept term
reading a foreign live strip still uses the
strip parameters above.

For a tuple of finite eligible switch sets
\(\vec B=(B_1,\ldots,B_4)\), fix nested regions
\(Y_t\) before varying the vertices. Take \(Y_t\)
to be the union of the level-\(m\) cube cover of
\(K_v\) and of the cells of \(B_1,\ldots,B_t\),
enlarged by the components of \(Z\) it meets.
Thus \(Y_t\) depends on the full sets \(B_s\),
\(s\le t\). If \(J_t\subset B_t\) is the
enabled set, the box in \(Y_t\) consists of the
fixed core and hubs there together with the
cells of \(J_t\). Hubs outside \(Y_t\) are
separate wall components and are unread. Build
the chain from link \(4\) to link \(1\), using at
each new chart the parent readings on its layer
and the prescribed data at the original sites.

For the analytic handoff to the next link, let
\(c\) be a young site in that chart, with padded
read \(P^+(c)\). In a common frame write
\(b_t^{\mathrm{off}}(c)\) for the parent layer
increment before the move and
\(b_t^{\mathrm{nat}}(c)\) for its increment
after the move. Use the handed increment
\[
 b_t^{\mathrm{off}}(c)+
 \prod_{\Delta:\ \mathcal N(\Delta)\cap P^+(c)\ne\varnothing}
    s_t(\Delta)
       \bigl(b_t^{\mathrm{nat}}(c)-b_t^{\mathrm{off}}(c)\bigr),
\]
where the product uses the link's switch cells
and \(\mathcal N(\Delta)\) is the switch halo.
Compose its exponential, after frame transport,
with the baseline group datum. At original
sites the prescribed increment is zero. If a
zero switch meets the padded read, the layer
datum is the parent datum on both faces. If
its halo misses the read, the natural increment
is unchanged by the local response identity.
At all ones the handoff is the native layer
datum. In particular, the row \(2\) continuation
uses the switched datum even off its box.
Keeping the old exterior field without that
switch would give two different layer data on a
zero face when the toggled cube contains that
read. On the admitted reference tube the number
of halo factors per young site is bounded and
the layer estimate keeps the new datum in its
chart.

On the admitted switch discs the responses,
layer averages and chart maps are holomorphic.
Their composition at a fixed \(\vec B\) has
the same zero-face agreement: the regions are
fixed, the parent of the toggled link is fixed,
and the remaining links read its output only.
For the invariant unit evaluation \(F_{\vec B}\)
on this family let \(T_v(\vec B)\) be its
alternating sum over \(J_t\subset B_t\) at
unit switch values. For \(\vec B\ne\vec\varnothing\)
the mixed Cauchy estimate at radius \(\rho>1\)
gives
\[
 |T_v(\vec B)|\le
 \left(\frac{2}{\rho-1}\right)^{\sum_t|B_t|}
 \inf_{c_0}\sup|F_{\vec B}-c_0|.
\]
This estimate needs only the family for the
given full tuple.

To recover the endpoint, take \(D=C\) or
\(C^\sharp\), and sum over switch sets contained
in \(D\). Each \(Y_t\) is contained in \(D\)
because \(D\) contains \(K_v\) and contains every
component of \(Z\) it meets. In the sum over
\(B_4\), the regions \(Y_1,Y_2,Y_3\) are fixed.
The link \(4\) value for \(J_4\) depends only
on \(J_4\) and the hubs reached by the reads
from the fixed smaller regions or connected
to their wall components. Extra hubs included
only by \(B_4\setminus J_4\) are separate
summands. At the full enabled set every cube of
\(D\) outside the fixed core and hubs is enabled.
Any hub in \(D\) belongs to a component of \(Z\)
that meets the core or those enabled cells,
so the enlargement includes it. Since \(D\) is
wall connected and the other hubs are separated
from it, M\"obius inversion on its eligible cells
gives the all-enabled link \(4\) on \(D\).

The chart at \(Y_3\) is now the restriction of
that natural vertex. Indeed its block set is
the induced one, its layer data are the readings
of that vertex, and original sites of \(Y_3\)
are original sites of \(D\) with the same datum.
The restriction rule and uniqueness of the
regular critical orbit identify the two charts,
up to the complex gauge in the chart. The same
identification at complex data follows on the
admitted connected seed discs by holomorphy.
At fixed \(J_3\), the next links use this
restriction only on the enabled box and on the
reads passed to \(Y_2\). The native vertex
depends on the parent on its box, and its
output off the box is the parent. These reads
are independent of any extra part of \(Y_3\)
included through \(B_3\setminus J_3\); the box
determining set is the label restricted and
eroded at the enabled box itself. Thus the
sum over \(B_3\) gives the all-enabled link
\(3\) on \(D\). For link \(2\), use the
natural vertex with the prescribed data on
the two sides. Its tilted naturality gives
the critical chart with the original datum
passed to link \(1\). The same restriction and
inversion first at link \(2\) and then at link
\(1\) give
\[
 \sum_{\vec B:\, \bigcup B_t\subset D\ \forall t}
 T_v(\vec B)
  = F_v\bigl(X_v;(q^D,J(q^D))|_{X_v}\bigr).
\]
Here \(q^D\) is the field of the four natural
links on the full box \(D\). Restriction to
the hull followed by its invariant read gives
the last expression. This proves the endpoint
identity with the tuple-dependent regions,
using the native chart restriction and
locality on the admitted one-run tubes.

% Earlier Part III, 120292--120488.
\paragraph{Counting the four successive box reads.}
For a kept term whose construction uses four
successive links on boxes with a fixed core \(K\),
count the switched cells in their order of
dependence, starting with the link read by the
term. Let \(B_t\) be the switched cells of
link \(t\). Start with \(\mathcal R_0=K^+\),
the core and its touching cells. At link \(t\)
adjoin \(B_t\) and the protected blocks to \(K\).
Keep the wall components that meet
\(\mathcal R_{t-1}\), and add those components
and their touching cells to obtain
\(\mathcal R_t\). These sets contain the possible
reads of the first \(t\) links at any vertex
with smaller switch sets.

The alternating coefficient is zero if a
switch cell lies outside the retained wall
components of its link. Removing it leaves
every read unchanged at each vertex paired
with that removal. For a nonempty tuple the
coefficient is also zero if none of the
retained components meets a protected block.
In that case every link is at its base on
the required reads. A connected part of
\(B_t\) may attach to the already reached
region without carrying a source of its own.
Count such parts by their first cell touching
the reached region, directly or through a
protected block. Following a wall path to
\(\mathcal R_{t-1}\) gives such an attachment:
as long as the path uses switched cells and
their shared protected blocks it stays in one
connected part, and the first passage into
the core already meets \(K^+\).

On this geometry the cells at the core wall
are level-\(m\) cubes. Together with the
bounded size of a protected block and the
bounded adjacency between cells, this gives
\(C(1+d_m(K)+\sum_{s<t}|B_s|)\) choices of
an attachment cell. Here \(d_m\) is ordinary
tree length in level-\(m\) units. The filing
label is built by joining each part at its
attachment and including the parents of its
cells and its touched protected blocks.
This costs at most \(c_l|B_t|\) additional
tree length at link \(t\). This construction
uses the same filing rule for membership in
each box. The coefficient bound is used on
this larger set of allowable families.

At the chosen connected-family decay rate,
the sum for a single attachment cell is
bounded by the convergent geometric series
for connected cell sets through that cell.
With the rate margin for the cell geometry, the product
over possible attachments is at most
\(\exp\{(d_m(K)+1)/4+1+\sum_{s<t}|B_s|\}\).
Sum the four links in reverse order, starting
with link four. Each summation spends at
most one unit of switch decay at each
earlier link. After allowing for the tree
weight \(e^{r d_m}\), a margin of three units
above the connected-family rate suffices.
The four sums together cost
\(e^5 e^{(r+1)d_m(K)}\).

In the exterior estimate the label of a
nonzero piece reaches a protected block
inside the operated region and also the
exterior anchor \(Q\). Its ordinary length
is at least \(9\check R_m+m_\sharp-2\)
and at least the distance from \(Q\) to the
enlarged component containing \(K\).
Here \(m_\sharp\) is the width of the exterior
enlargement in level-\(m\) units. Spend one
unit of length decay on each lower bound.
The preceding count with \(r=a+2\) leaves
the weight \(e^{(a+3)d_m(K)}\) for the
ordinary unit size sum. The remaining
distance to \(Q\) is summable over separated
components. The factor
\(e^{-(9\check R_m+m_\sharp-2)}\) pays the
strip and shallow-cell counts. This uses
the four-link filing and its stated
locality at the vertices.

% Earlier Part III, 84998--85190, and Earlier Part I,
% 102846--102912.
\paragraph{Old kept terms at a later healing step.}
For a kept term of \(Y\) outside the region \(Z\)
being treated at step \(k\), take the four-link
expansion on its admitted native family.
The empty tuple reads no response: its boxes
are the hubs of \(Z\), separated from the
padded kept hull. The remaining pieces
are split by the first link with a nonempty
switch set. The oscillation on the admitted
tube bounds the piece for \(\vec B\) by
\(\mathfrak O(v)e^{-\kappa_1^{(1)}\sum_t|B_t|}\).
The wall path from the padded hull to a
source leaves the strip's collar. The
depth and the number of selected cells
are therefore at least \(5\check R_m/2-1\)
for a strip born at \(m\), as needed for
the kept-anchor pile. Its label connects
\(Q_Y(k)\) outside \(Z\)
to a source cell in \(Z\), so the native
separation bound gives length at least
\(\check R_k-2\). Spend one unit of the
length decay on this separation, starting
the remaining pile at rate \(a+1\).

For the rest of the sum use the native
weighted piles \(\mathsf Q_r[\mathfrak L^t]\),
with \(t\) the number of links and \(r\)
the rate of the label weight. The reading
sets give the distance condition at each
link. Summing the last link bounds the pile
by \((1+C_\Sigma^{\rm sum}C_e)
\mathsf Q_{r+2}[\mathfrak L^{t-1}]\)
for \(t\ge2\). At the first link the kept
anchor contributes \(4C_\Sigma^{\rm sum}
\mathsf T^*_{r+2}\), where \(\mathsf T^*_r\)
is the kept-anchor pile and the four counts
come from the first nonempty link. The
rate floor \(\kappa_1^{(1)}\ge c_l(a+9)+\lambda_*\)
permits the four sums and the one unit
already spent. Since the anchor pile
is at most \(\mathfrak n^{\rm kept,U}\),
the exterior norm at rate \(a\) costs
\[
 e^{-(\check R_k-2)}
 (1+C_\Sigma^{\rm sum}C_e)^3\,
 4C_\Sigma^{\rm sum}\mathfrak n^{\rm kept,U}.
\]

If \(Y\) lies inside a treated component \(C\),
its padded hull lies in the never-integrated
core \(\mathsf C\) of \(C\). On the admitted
switch polydisc, compare with a fixed shape
point over the same seed and retained datum
in the same frame. The framed dilation radius
is at least \(e^{\delta_P D'/2}\), with
\(D'=d^\wedge((Y^\sharp)^+,\mathsf C^c)\).
Schwarz's lemma bounds the difference by
\(\mathfrak O(v)e^{-\delta_P D'/2}\)
throughout the switch polydisc. Subtracting
the fixed value leaves every nonempty
mixed coefficient unchanged. The switch
Cauchy estimate then retains the same
switch-count decay, with the improved
oscillation weight.
For any exterior piece at least one selected
family leaves \(C\): if all families stay
inside, so do their corridors and hubs,
which lie in the smallest boxes containing
their endpoints. Assign \(Y\) a cell of
\(\mathsf C\) containing a hull point nearest
to \(\mathsf C^c\). The distance of that
cell to \(\mathsf C^c\) is at most \(D'\),
and at most \(2^d\) regions use a given
cell. Since \(D'\ge w(\check R_m-1)/2\)
for a region born at \(m\), one quarter of
the exponent absorbs the polynomial
oscillation at that birth level. The
other quarter sums by the boundary cell
estimate for \(\mathsf C\). Thus the deep
kept contribution is bounded by the
boundary cell count times
\(C\mathfrak n^{\rm kept,U}\).
Both bounds use the native kept tube,
the locality of its pieces and the one-run
geometry at healing.

% Earlier Part III, 120164--120172, 121189--121269
% and 121408--121481.
\paragraph{The exponent of a nested history.}
At an admitted root read, separate the exponent
of each node into its own list and quadratic
form, the kept members of its immediate
children, and the Gaussians of its live arrests.
A child's kept member is read
as its empty-switch value plus its nonempty
pieces at the parent. For each such member
write the two-run difference as the family
difference evaluated at a common datum plus
the change of read in one family. At the
evaluation of a treated root the two runs use
the same seed, so the choice of seed introduces
no reading mark. Retain the reading marks
of the individual stored terms,
including the read of a child at its parent's
vertex. In the child's family
difference, recurse only on the empty-switch
value, at the child's datum furnished by the
parent's empty-switch evaluation. Keep the
nonempty pieces at the parent. The levels
decrease along branches and the recursion is
an identity at every node. In the resulting
sum each own unit is charged once, with the
original class for each family mark, reading
mark, form and Gaussian.

Write \(N(Y)\) for the count of determining
sites by stratum in the padded region of
\(Y\) at its own step, and \(K_{\rm cell}(m_Y)\) for
the one-run per-cell bound at that node's
step \(m_Y\). The tree totals are
\[
 N_{\rm tree}=\sum_Y N(Y),\qquad
 V_{\rm tree}=\sum_Y K_{\rm cell}(m_Y)N(Y).
\]
The children inject into the parent's stage
cells. The live arrest collars at a node
also charge those cells, with disjoint arrest
regions. Hence the sum of the child count and
the arrest collar-cell count is at most
\(2N(Y)\). At a real root the retained Gaussian
value is nonpositive, and the one-run
oscillation bounds its increase on the
admitted tube. The node's own quadratic form
has bounded real part and costs a fixed
amount at its node. The screened nonempty pieces
of a child cost a fixed amount per child:
subtract its empty-switch value and apply
the depth estimate at fixed collar variable
before the switch sum. Summing by generation
gives the real-part bound \(V_{\rm tree}+C
N_{\rm tree}\) in either run.

For the comparison on a class dial, the
per-cell family bounds give a contribution
proportional to \(N_{\rm tree}\) after
cancelling the load. In the child-family
contribution of a nonempty switch set,
subtract the family mark at the empty-switch
datum before taking the alternating sum.
The comparison at fixed collar variable gives
the depth factor for this difference on the
contour. It pays the polynomial size of the
child's family pile. For a child the
change-of-read estimate at its parent's
vertex has factor \(\vartheta_f^n\).
For the nonempty pieces take the geometric
mean of the vertex and screened contour
bounds, leaving \(\vartheta_f^{n/2}\)
and enough depth decay to pay the child
sum. The dial radius is at most
\(C\vartheta_f^{-n/4}\), by the lower bound
on the full load. The absolute contributions
from reads are therefore bounded by fixed
multiples of \(V_{\rm tree}\) and \(N_{\rm tree}\).
The family part of the node's own form is
charged at its determining rows. To bound
the reading part, use the mixed-datum line
from the common datum at zero to the second
run's datum at one. At the node's step \(m_Y\)
it lies in the tube for \(|s|<R\), with
\(R=(8C_{\rm read}\vartheta_Y^\sharp)^{-1}\ge2\).
Here \(\vartheta_Y^\sharp\le\vartheta_f^{m_Y}\)
is the relative reading rate and \(C_{\rm read}\)
includes the frame allowance. Along the line
the imaginary part of the form is bounded by
its small one-run majorant \(\Sigma_Y\). For
\(\Sigma_Y>0\), mapping the strip
\(|\operatorname{Im}w|<\Sigma_Y\) to the disc
by \(w\mapsto\tanh(\pi w/(4\Sigma_Y))\)
and using Schwarz--Pick bounds the change by
\[
 \frac{2\Sigma_Y}{\pi}\log\frac{R+1}{R-1}
 \le \frac{2\Sigma_Y}{R}.
\]
Use a slightly larger strip and then its limit
if needed. This also gives zero when \(\Sigma_Y=0\).
The dial bound \(C\vartheta_f^{-m_Y/4}\)
makes the cost at most one per node for the
class \(m_Y\), by the smallness of \(\Sigma_Y\).
For the Gaussians use the dial bound per arrest
collar cell. The preceding count charges these
costs to \(N_{\rm tree}\). Thus the real-part
estimate on the interpolated exponent has
the same two tree totals, with the prescribed
enlargements of the coefficients.

% Earlier Part III, 41880--41956, 82753--82799
% and 120164--120172. Current Part III, Proposition 80.8,
% Lemmas 87.31--87.33, and 67189--67199.
For a first-class component \(C\) reached
from a larger residual domain, take the average
of the larger root over the determining set
of \((C^\sharp)^+\), including its boundary
rows. Use the same real inner variables
\(\Phi_C\) and construct the slice on
\(C^\sharp\) from this averaged datum.
On a finer separated box the real
restriction is critical. Indeed, extend
a competitor on the smaller fibre by the
original root outside the box. The finer
constraints preserve the larger averages.
Any plaquette meeting the cut has its
inside bonds on the frozen fine boundary
rows, so its action is unchanged. The
action difference is exactly the difference
of the two smaller-box actions. This also
identifies the restricted current on the
bonds whose six plaquettes are retained.
The inherited layer bounds and regularity
of the induced datum give the real orbit
identification when their bounds lie in the
box radii.

On complete carriers protected by the padding,
the restriction of the shell hierarchy preserves
the levels of the rows. Indeed, the shell at a
given level is obtained by intersecting its
ambient member with the box eroded at that
level, so it is unchanged on blocks away from
both cut collars. The cut can still create new
boundary rows. On each complete double block
the exact average is holomorphic about an
admitted real reference, with both endpoint
node factors and the bounded averaged
logarithm. The datum on the collar must be
computed by this read. Even on a matched real
collar the aligned datum can be \(e^{iB_1}V^W\),
where \(V^W\) is the native box datum and
\(B_1\) is the small collar correction.

To use this restriction in the residual
production, we must match the finer separated
sequence with the native padded determining
set of \(C\), including the boundary rows
where the averaging contours change. The
slice with the assigned current must land
in the half tube of the kept-factor bound.
For the paired and mixed reads we must also
fit the quarter-tube and eighth-tube
reserves and the allowance for the frame
change. Real restriction on the matched
rows gives the real endpoint of this test.
The second-class kept-member read uses its
own admitted box.

% Earlier Part I, 23032--23061; Part II, 14526--14572;
% Part III, 12646--12659, 13001--13027, 119174--119221
% and 124540--124587.
\paragraph{Arguments of the production formula.}
For an aligned label use the recorded direct
units, the read terms from earlier productions,
and the operator and normalizing entries.
The direct entries include the coupling and
source coefficients and the constants in the
fifth load. At a fluctuation step the operator
entries are the covariance, the walk kernel and
its composed quadratic form. The remaining
entries include the localized potentials,
Jacobians and the normalizer with its stored
boundary share. At a healing step use the
corresponding chain and kernel entries.

For the large-field integral, on its native
history tree the exponent consists of the
raw Wilson term, the inner integrals' quadratic
forms and source weights, and the normalization
constants at each node, together with the kept
exponent of the output. The latter is the fresh
form and kept potential, with the child terms
and arrested Gaussians inserted as in the
nested sum. Haar factors, tree constraints,
the tests and the inner integral kernels are
part of the production too. When a source is
held fixed, some of these factors still depend
on the cutoff through the coupling, density
or chart. In particular the old inner
quadratics have their own cutoff differences
even when independent of the later background.
The histories of an inner integral
also act on the outer potential terms whose
domains contain that integral's component.

Use the localized remainder after the Wilson
and quadratic cancellation. In the resulting
array its terms take the place of the raw
brackets they represent. The kept factors of
a pocket are separate entries until they
are read into a later exponent. For a live
arrest use the Gaussian piece, its shifted
value minus the base, its grain piece, or
its retained weight, according to the
reading site. Each is taken at the site and
in the variables in which it occurs. The
new integral's determinant is a separate
normalizing entry.

On a common chart for the integrated and
retained variables, the measures, constraints,
tests and reference factors must either be
identified by the alignment or entered with
their actual changes. With these identifications,
substitution of the two arrays into the same
integral formula gives the two productions.
The componentwise family and reading identities
then give \(x^{\mathsf B}-x^{\mathsf A}
=\sum_z W^z\) in the entries of that formula.
This check uses the full list of native
factors at the production, including the
terms carried by its inner histories.

% Earlier Part III, 14022--14056, 119539--119570,
% 119920--120121 and 120174--120179.
\paragraph{Bounds for the production on the class discs.}
At a fluctuation half-step the class derivative
splits into an operator direction and a marked
potential direction. The Gaussian units
enter the latter through their weighted
pile. The marked expansion at the convex
base is linear in the marked potential,
and bounds its derivative by the pile times
the native output decay. The one-cell pile
and the fluctuation weight therefore add
at most \(C\theta_i\Pi_G^+(x_i)\Delta_n^c\)
to the class-\(n\) input. Here \(\theta_i\)
bounds the inverse fluctuation weight.
Every filed Gaussian unit holds a retained
collar cube, so it enters the large-field
activity. The list of outputs wholly in the
clean region does not read it. The operator
direction keeps its own bound.

For a live arrest \(W\) outside the components
being treated, evaluate the retained quadratic
at the new root on the top storey. Write
\(D_W^z\) for the sum of its marks of class
\(z=(n,k')\) and \(N(W)\) for the number of
filing cells holding its collar. The vertex
bound gives \(|D_W^z|\le K_G N(W)\Delta_n^c\).
For \(0\le\lambda\le1\) and \(|\varepsilon|\le
\Lambda_z\), convexity of the real part gives
\[
 \left|\exp\big((1-\lambda)q_W^{\mathsf A}
       +\lambda q_W^{\mathsf B}
       +\varepsilon D_W^z\big)\right|
 \le
 \max_{\mathsf m}|e^{q_W^{\mathsf m}}|\,
 e^{K_G E_1^zN(W)} .
\]
The collar cells of distinct live arrests
are charged on their disjoint regions. Keep
the maximum separately for each arrest. The
exterior difference from the root to the
healed field is filed in the potential and
uses the exterior sum above. At the last
lattice the weight is read at the root
without intervening links. For an arrest
nested in a treated component \(C\), use
the collar-cell bound inside the exponent of
the healer. Its retained variable is
integrated there. For a class \(n\) present
in that arrest the birth index satisfies
\(j(C)\le n+1\), giving age at least \(m-n-1\)
in the estimate for the healed term.

Before applying the strip estimate, take the
convex bound on the one-run potentials and
add the class direction. The exterior
remainder and the exterior Gaussian each
contribute a one-run norm and a two-run norm
multiplied by the dial radius. The class
direction is bounded by using
\(\Lambda_z\Delta_n^c\le E_1^z\).
The four coefficients tend to zero under
the same coupling ceiling, so each may be
bounded by one before multiplying by the
native strip factor. In the rim direction
the two exterior marks enter the linear
potential with their original classical load
\(\Delta_n^c\). At a preparatory stage let
\(\nu_\zeta\) be the array of positive
unit bounds for the response direction
\(\zeta\), and \(\mu\) the moment bound
which controls the stage holomorphy.
The stage disc has radius
\(1/(4\mu[\nu_\zeta])\) for a nonzero
moment bound. Convexity gives the
one-run majorant on the base, and the
block estimate gives the additive
majorant on the disc. The staged leaf
estimate reads this direction bound
on the common chart. The Gaussian
collar coordinate stays fixed there.

The exterior bracket of the fifth
comparison row uses the one-run outer
cell estimate in the coefficient of its
own unknown. It therefore multiplies the
fifth load in that row.

These modulus bounds are for the unnormalized
activities. In the activity and age estimates
the new costs are the nested exponent, the
collar cells of the untreated arrests, and
the bounded remainder per treated component.
Charge them in the entropy constant and use
the age bound at each generation. For the
residual output with filing domain \(X'\),
the same production map \(\mathcal P_{X'}\)
is used at the two arrays of slot values.
Apply the modulus bounds to its values on
the disc \(x^\lambda+\varepsilon W^z\),
where \(x^\lambda\) is the convex array and
\(W^z\) is the class direction. Cauchy's
estimate can then bound the corresponding
first variation by the supremum on that
disc divided by \(\Lambda_z\). The own
classes use their healed exponent bound.
This is the passage from the component
estimates to the class directions used in
the terminal sum.

% Earlier Part III, 13365--13404, 14473--14524
% and 120110--120121.
\paragraph{Accumulation of the stored boundary loads.}
Fix the birth zone \(n\). Let \(u_j\) be
the partial sum of the class-\(n\) loads
of stored boundary terms after the first
\(j\) productions that create such a term.
The next production reads at most \(a_n+u_j\),
where \(a_n\) is the class input before this
feedback. A fluctuation step adds its store
with the crossing weight \(w_{\ell,n}\) of
its output. A healing step adds its residual
store with its residual comparison coefficient.
Thus the two estimates take the form
\[
 u_{j+1}\le (1+\mu_j)u_j+\mu_j a_n+b_j,
 \qquad u_0=0.
\]
The \(b_j\) are the weighted classical
inhomogeneous bounds at the fluctuation
steps. The first two crossing weights
are one. For \(\ell\ge n+2\), let \(t\)
be the depth of level \(\ell-1\) from
the terminal level. The weight is bounded
by a multiple of
\(\exp[-(2/L)(\log(x_*-2c_2+\lambda_{\rm flow}t))^{r_{\rm pr}}]\),
where \(x_*\) is the lower bound on the
inverse squared coupling and
\(\lambda_{\rm flow}>0\) is the flow lower
bound. The depths for different \(\ell\)
are distinct. Since \(r_{\rm pr}>1\),
the terms are bounded by a summable tail
uniformly for \(x_*\) above a fixed lower
bound. The tail tends to zero as \(x_*\)
increases. This gives \(\sum_\ell w_{\ell,n}\le3\)
at small coupling.

The fluctuation coefficients are bounded by
their uniform bound times \(w_{\ell,n}\), and
the residual coefficients already have a
summable bound. With the enlarged smallness
budget of the class estimates these give
\(\sum_j\mu_j\le1/40\). Iterating the
inequality and bounding the products of
\(1+\mu_j\) gives
\[
 u_j\le
 e^{1/40}\left(a_n\sum_{i<j}\mu_i+\sum_{i<j}b_i\right)
 \le a_n/38+e^{1/40}\sum_{i<j}b_i .
\]
The new classical marks are bounded by the
existing class load and so fit in its right
member. The fifth-row term stays in its own
coefficient. Only these crossing-weighted
stores and the residual stores enter the
accumulated load.

% Earlier Part III, 13743--13764, 13864--13911,
% 14181--14197, 14305--14318 and 17645--17682.
The zone of a slot is assigned before its load is
formed. A resolved slot inherits an anchor in the
zone of its formation. An ordinary zone-\(b\) anchor
is in \(\Omega_{b+1}^c\), and a prepared zone uses
the quiet window to give the same birth bound.
Trace the anchor through the later reads of that slot.
As long as the containing component survives, its
successors retain that birth and meet the read domain.
A second-class or arrested successor retains its record.
If the containing component is healed, the interactions
reading its inner history are summed in the healing
bracket. A slot whose zone occurs only in the summed
sequence is no longer a resolved read of that zone.
Its old age factor is spent in the own-output estimate
at production. File the resulting whole slot, with its
norm, at the later of its stored level and the block
zone. At subsequent reads use the live-incidence
condition of the storage rule for that new class.
For \(b\le k-2\), the terminal comparison requires
each sensitive label in class \(b\) to have a live
component born by \(b+1\). This is checked using the resolved
interpolation on the actual reads and the block
incidence for the own slots. The birth of an erased
component gives no terminal age.

% Earlier Part III, 12045--12099, 14402--14654,
% 15155--15321 and 17261--17495.
\paragraph{The envelope for the full zone load.}
The preceding recurrence sums the stores resolved in one
birth zone. The boundary terms produced in healing and
filed at their own level also enter the load. Let \(\ell_j\)
be the full load in zone \(j\), including the regular and
remainder reads, both boundary loads, the shift and
increment reads, and the classical index terms. Put
\(s=K-j\), \(e_s=\rho^{K-s}+\Pi_K\), and
\(\mathsf l_s=(\log\theta_{[s]})^{c_1}\), with the exponent
covering the logarithmic weights of the comparison.
For a lag \(m\) use
\[
 e_{s+m}\le\rho^{-m}e_s,\qquad
 \mathsf l_{s+m}\le\widehat{\mathsf l}_m\mathsf l_s ,
\]
where \(\widehat{\mathsf l}_m\) grows at most as a power
of \(1+\log(1+m)\). Use the constituent response
envelopes on the actual common slot reads. On a live
second-class strip we must verify the relative-format
version of those bounds on its stored reads, as well
as the ordinary version on the other slots. In the
weighted term norms \(S_i^+\), at \(i=j-m\) they give
\[
 S_i^+[E]\le C\mathsf l_{s+m}e_{s+m},\qquad
 S_i^+[R]\le C\theta_{[s+m]}^{-1}\widehat r_{s+m}e_{s+m}.
\]
The first bound adds the anchored, quadratic and
source-defect entries in their respective reading
units. The weight for the \(E\) terms is bounded by
\(C(m+1)^3q^m\). Since \(q/\rho<1\), its sum is
\(O(\mathsf l_s e_s)\). The \(R\) weight instead
contains \(\theta_{[s+m]}/\theta_{[s]}\).
It cancels the inverse coupling in the \(R\) reading
unit. The unweighted remainder tail
\(\sum_{r>s}\widehat r_r e_r\le2\Pi_K\) then gives
\(O(e_s)\) for the \(R\) load.

The gauge load convolves the shifts
\(O((1+s+m)e_{s+m})\) with \(L^{-m}\), and the
increments \(O(e_{s+m})\) with \(L^{-m/2}\).
The prefactor is bounded by
\(C\theta_{[s]}^{-1}(\log\theta_{[s]})^{2q_\flat}\).
Both convolutions are summable relative to \(e_s\).
The ratio \((1+s)/\theta_{[s]}\) is bounded by the
flow. Thus this load is \(O(\mathsf l_s^2 e_s)\).
For the classical \(E\) and boundary entries the
kernels are \(C(m+1)^3q^m\) and \(L^{-m/2}\),
respectively, times the class rate at \(i\), bounded
by \(C\rho^i\). The classical \(R\) entry with
\(g_i^{\kappa_0}\) has instead the coupling weight
of the \(R\) row. After its cancellation the sum
contains \(\theta_{[r]}^{1-\kappa_0/2}\rho^{K-r}\)
at depth \(r\). The one-run comparison of remainder
coefficients bounds its sum by \(C\Pi_K\).
The last classical entry costs \(C\mathsf l_s e_s\)
by the index rate and the logarithmic weight.
Together these give \(O(\mathsf l_s^2e_s)\) before
the two boundary loads are added.

At healing level \(r\), the own output bound uses
the load in each zone \(i\le r\) with its terminal
age weight. Reading that output in zone \(j\) costs
\(L^{-(j-r)/2}\). Interchanging the sums over \(r\le j\)
and \(i\le r\), and adding the resolved boundary
estimate above, gives
\[
 \ell_j\le a_j^0+\sum_{i\le j}\mathcal K(j,i)\ell_i,
 \qquad
 \mathcal K(j,i)=\frac{4(1+\epsilon_0)}{E_0}
 \sum_{i\le r\le j}L^{-(j-r)/2}c_1^\flat(g_r)
             (r-i+1)q_a^{(r-i-1)_+}.
\]
Here \(\epsilon_0\) is the resolved-load coefficient
of the own contribution. The direct and classical
bounds just proved enter \(a_j^0\) with fixed constants.
The additional own inhomogeneity at \(r\) is bounded
by \(C c_1^\flat(g_r)(q_a^{r-n_c}+\vartheta^r)\)
before the \(L^{-(j-r)/2}\) convolution, with
\(\vartheta\le\rho\). The factor \(q_a^{-n_c}\)
is bounded by a power of \(\log(X+CK)\), where
\(X=\theta_{[0]}\). Split that logarithm at depth
\(K-r\): the factor that grows with \(\theta_{[K-r]}\)
is absorbed by \(c_1^\flat(g_r)\le
C\theta_{[K-r]}^{-\kappa_0/2}\), and the factor
growing with \(r\) by \((q_a/\rho)^r\). The resulting
inhomogeneity after convolution is \(O(e_s)\).
Hence \(a_j^0\le C\omega_j\), where
\(\omega_j=\mathsf l_{K-j}^2 e_{K-j}\).

For the weighted supremum norm with weight \(\omega_j\),
write \(\xi=j-r\) and \(\eta=r-i\). The ratio
\(\omega_i/\omega_j\) is at most
\(\widehat{\mathsf l}_{\xi+\eta}^2\rho^{-(\xi+\eta)}\).
The two geometric factors \(L^{-\xi/2}\) and
\(q_a^{(\eta-1)_+}\) make the resulting double sum
finite, since \(L^{-1/2}<\rho\) and \(q_a<\rho\).
The norm of \(\mathcal K\) is therefore bounded by
\(C\sup_r c_1^\flat(g_r)<1\) at weak coupling.
The loads are finite by the a priori dial bound.
Taking the weighted norm in the displayed inequality
gives \(\ell_j\le C\omega_j\).
For the stop at depth \(s_\circ=K-k\), the same
lag comparison yields
\[
 \sum_{j\le k}(k-j+1)q_a^{(k-j-1)_+}\ell_j
 \le C\mathsf l_{s_\circ}^2 e_{s_\circ}
 \sum_{m\ge0}(m+1)q_a^{(m-1)_+}
             \widehat{\mathsf l}_m^2\rho^{-m}
 \le C'\mathsf l_{s_\circ}^2e_{s_\circ}.
\]
This uses the two-run fluctuation and healing bounds
on the admitted slots and the actual crossing weights.

% Earlier Part III, Proposition 9.2:prop-rest-term-tail,
% 26707--26746; Proposition 9.8d:prop-transport-O,
% 122081--122132; the whole-input allowance at
% 24709--24715 and 25837--25858.
\paragraph{Removing the old histories from a remainder.}
For a history-created second-group remainder
\(R'^{(k)}(X)\), retain in each operated
component only the sequences whose first
index exceeds \(n<k\). Write \(0<q_{\rm op}\le1/2\)
for their geometric forgetting rate.
Denote the resulting
formula by \(R'^{(k)}_{>n}(X)\). Make this
restriction on the full component operation,
including its last exponential, before
forming the receiving cluster. The following
estimate uses its admitted holomorphic
formation chart, with absolute subfamily
majorants and the doubled whole-component
input allowed by the age-dial construction.

Put \(a_0=k-n\), \(p_k=p_0(g_k)\),
\(P=p_0^\sharp(n)=\min_{j\le n}p_0(g_j)\),
and \(E_k=\exp[-2p_k/(1+\beta_0)]\).
Here \(0\le\beta_0<1\) is the degree-comparison
constant, and the age comparison gives
\(P\ge p_k/(1+\beta_0)\).
An age-\(a\) sub-sum in the removed part,
\(a\ge a_0\), has absolute majorant at most
\(\mathrm e\,e^{-2P}q_{\rm op}^{a+1}\sup|F|\).
Multiply all these old sub-sums, in every
operated component, by the same variable
\(\mu\). Use the disc
\[
 |\mu|\le R_\mu
       :=q_{\rm op}^{-a_0}
          \exp\bigl[2P-2p_k/(1+\beta_0)\bigr].
\]
The whole old input on this disc is bounded by
\[
 \mathrm e E_k\sup|F|
       \sum_{a\ge a_0}q_{\rm op}^{a+1-a_0}
 \le\mathrm e E_k\sup|F|,
 \qquad q_{\rm op}\le\tfrac12.
\]
The younger input has the same bound without
the dial. Thus their sum stays within the
allowed \(2\mathrm e E_k\sup|F|\) input.
This sum is made at the level of the
history-dependent absolute majorants, with
the last exponential split separately in
each history. No unchanged per-age bound
is needed after the old part is enlarged.
The receiving formation therefore remains
holomorphic and has its bound
\(\mathfrak B_X=e^{-p_k}e^{-\kappa d_k(X)}\)
on the disc. At \(\mu=1\) it is the original
remainder and at \(\mu=0\) its younger part.
Since \(R_\mu\ge2\), the one-variable
Schwarz estimate gives
\[
 \begin{split}
 |R'^{(k)}(X)-R'^{(k)}_{>n}(X)|
 &\le 2\mathfrak B_X/R_\mu\\
 &\le 2e^{-\kappa d_k(X)}
       e^{-(1-\beta_0)P}q_{\rm op}^{k-n}.
 \end{split}
\]
The second inequality follows from
\(p_k\le(1+\beta_0)P\). Multiplication
by the domain weight gives the same bound
in the flat remainder norm.

The restriction also identifies which old
uses can remain. A direct use of zone
\(n'<n\) requires an operated sequence
with first index at most \(n'+1\le n\),
so it has been removed. If every live
witness of a leaf on this reader belongs
to one of its operated components, a
live witness of level \(e\le n\) likewise
requires first index at most \(e\), and
has been removed. Old zone copies can
still occur inside an owned leaf formed
at \(e>n\). The carry bound gives them
the factor
\[
 C_*q_{\rm op}^{e-n'-1}
                      \le C_*q_{\rm op}^{n-n'}.
\]
Such a copy can occur only when the
containing operated sequence starts by
\(e\). Thus the younger remainder retains
these controlled old-zone reads. In using
the live-witness exclusion we must check
the stated witness containment on the
actual receiving carriers.

% Earlier Part III, Lemma 9.2:lem-aligned-parts,
% 26748--26890, and Proposition 9.8d:prop-transport-G,
% 122135--122221.
\paragraph{Aligning the nested stored terms.}
An aligned outer sequence can still read a stored
term whose own healed histories start below
\(n_c\). Restrict those histories before using
the comparison on zones at or above \(n_c-1\).
For a first-group stored leaf formula \(\psi\) formed at level
\(e\), define \(\psi^{\rm AL}\) recursively
in its formation level. Restrict the sequence
sum of each healed component in its cluster
to first index at least \(n_c\). Whenever a
nested occurrence loses its witness under
this restriction, replace it by its own aligned
formula. The recursion terminates since nested
formation levels decrease. This changes the
formula on its original variables and preserves
the full carriers used in its majorant.

First consider the branch with per-cell leaf
piles \(\epsilon_e^{\rm all}\) satisfying
\[
 \epsilon_{\rm tot}^{\rm all}
       :=\sum_e\epsilon_e^{\rm all}\le\tfrac1{16},
 \qquad
 C_{\rm AL}:=(1-4\epsilon_{\rm tot}^{\rm all})^{-1}
                                      \le\tfrac43.
\]
Use the same admitted per-use domains and the
two independent age and leaf dials as in the
carry estimate, with the operation rate
\(q_{\rm op}\). If \(\mathfrak B_\psi\) is the
original full-carrier majorant, then
\[
 |\psi-\psi^{\rm AL}|
       \le C_{\rm AL}\mathfrak B_\psi
                       q_{\rm op}^{(e-n_c)_+}.
\]
Indeed, removing its own unaligned histories
costs at most
\(\mathfrak B_\psi q_{\rm op}^{(e-n_c)_+}\) by geometric
forgetting. For nested occurrences formed at
\(r<e\), the induction hypothesis gives the
per-cell dial size
\(C_{\rm AL}\epsilon_r^{\rm all}
q_{\rm op}^{(r-n_c)_+}\). The age dial has radius
\(q_{\rm op}^{-(e-r)}\), and the leaf dial has the
reciprocal of that size as its radius.
The mixed difference estimate therefore bounds
this level's contribution by
\[
 4C_{\rm AL}\epsilon_r^{\rm all}
   \mathfrak B_\psi q_{\rm op}^{e-r+(r-n_c)_+}
 \le 4C_{\rm AL}\epsilon_r^{\rm all}
   \mathfrak B_\psi q_{\rm op}^{(e-n_c)_+}.
\]
Summing over \(r\) closes the induction because
\(1+4C_{\rm AL}\epsilon_{\rm tot}^{\rm all}
=C_{\rm AL}\). The constant is independent
of the formation level. If only a uniform
per-level pile bound is available, the
strengthened carry estimate above remains
available; the present constant specifically
uses the summed small pile.

Run the affine comparison on these fully
aligned formulas, including all their nested
occurrences. Every retained use is then in
the zones where the comparison input applies.
The discarded nested differences enter the
unmatched term of an operated component.
A reader \(\mathbf T'_k(X)F\) gives a
level-\(e\) difference the additional factor
\(q_{\rm op}^{k-e}\). Since
\[
 (k-e)+(e-n_c)_+\ge k-n_c,
\]
the transported difference has the same age
power as the own unaligned histories. The
reader production factors and full carriers
are the same in these two estimates. Sum
\(C_{\rm AL}\epsilon_e^{\rm all}\) over the
formation levels. The own unmatched bound is
multiplied by
\[
 C_{\rm align}:=1+C_{\rm AL}
                 \epsilon_{\rm tot}^{\rm all}
                               \le\tfrac{13}{12}.
\]
Keep the terminal rate \(q_a\) fixed and
choose the operation age reserve so that
\(q_{\rm op}\le\min\{q_a,1/2\}\). Take the admitted
dial and history bounds at this reserve
before fixing the coupling ceiling.
For \(k\ge n_c\), the operation error is
then bounded by a constant times
\(q_a^{k-n_c}\).
This construction uses
the common scaffold, admitted dial domains
and zone comparison of the two runs; it does
not supply those inputs on an unadmitted read.
% Earlier Part III, the terminal slot budget at
% 13794--13818, stored-family pile at 24762--24788,
% terminal modulus at 17584--17624, and absorption
% into both output groups at 25252--25255.
For a first-group term still stored at the
terminal step, the slot argument is more
direct. Let \(\Psi_u\) be the term of level
\(e\) occupying such a slot, and set
\(H_u^{(e)}=\Psi_u-\Psi_u^{\rm AL}\).
Count each outer stored slot once; its
aligned formula already includes the nested
replacements. On its admitted full reading
domain the preceding bound gives the
per-cell size
\[
 \rho_e=C_{\rm AL}\epsilon_e^{\rm all}
                  q_{\rm op}^{(e-n_c)_+}.
\]
For nonzero \(\rho_e\), use the slot path
\(x_u(\nu_e)=\Psi_u^{\rm AL}
                 +\nu_e H_u^{(e)}\),
with the other level slots held at their
current endpoints. Its values at zero
and one are the aligned and original
slots. Omit levels with zero pile.
The stored-family value and activity budgets
are linear in this pile. Thus a single
level-\(e\) parameter multiplying
\(H_u^{(e)}\) has a disc of radius a fixed
multiple of \(\rho_e^{-1}\) in the terminal
slot. Reserve the corresponding one-unit
value and activity allowance, with the
coupling ceiling making the radius at least
two. Keep the ramps, tests and real geometry
fixed. The terminal modulus proof uses
these value and activity bounds and the
same raw Wilson factors, so its bound on
this disc is \(CA^\sharp\).

A still-stored boundary term has an anchor
in the enlargement of the level-\(e\)
large-field region in the current label.
Its sensitive terminal labels therefore
have a live component of age at least
\((k-e-1)_+\). Subtract the insensitive
labels and use their age-weighted modulus
bound on the leaf disc. The Schwarz
estimate, with arbitrary labelwise Borel
phases, bounds the sum of the resulting
total variations by
\[
 CA^\sharp C_{\rm AL}\epsilon_e^{\rm all}
       q_{\rm op}^{(e-n_c)_+}
                         q_a^{(k-e-1)_+}.
\]
For \(k\ge n_c\), summation over \(e\)
gives
\(C(q_a)A^\sharp C_{\rm AL}
\epsilon_{\rm tot}^{\rm all}q_a^{k-n_c}\).
This argument uses the full leaf reading
domain, including its configuration tube.
We must check that the actual terminal
read lies in that tube and that its slot
has the reserved value and activity
allowance.

When a leaf has been absorbed by a healing
operation, its change can enter both a
younger boundary term and a second-group
regular remainder. For the former use the
operation-reader estimate and then the
surviving younger boundary slot. For the
latter we must identify the recomputed
list and terminal slots and verify their
comparison bounds. Such absorption does
not by itself supply an old live witness
at the terminal step. After these terminal transport
checks, write \(C_{\rm term}\) for the
fixed coefficient of the total unmatched
bound.

% Earlier Part III, Lemma on comparison of aligned terms,
% 17504--17734, and the terminal sum, 18118--18225.
\paragraph{Summing the changes of the aligned terminal terms.}
On the common terminal variables first compare the fully aligned
labels with the ramps and the stop-level production slots
of the ramp move fixed to those of run A. On a
standing second-class component the common variable is
\(V'\), with \(V=V'V_{\Lambda}^{\mathsf c}\) in run
\(\mathsf c\). On a removed component with stored
shell variables use \(V''=V'M_{\mathbf B''}(U_0^{\mathsf c})\).
Elsewhere keep the raw variables. In the order of the
reference construction each multiplier is determined
by the already fixed exterior. Conditional right
translation thus preserves the free Haar factors,
and the recorded tree constraints act on the common
relative variables. Both physical charts push forward
from the same reference space for this label and frame.
The two physical tree slices can have different images.
All total variations of aligned differences below are
computed on the common chart. Write \(\tau_h^{B\mid A}\)
for the B term with these A ramps and production slots,
and \(\tau_h\) for a term along the interpolation to it.
Test the family against arbitrary Borel
\(\varphi_h\) with \(|\varphi_h|\le1\). The first variation of
\(\sum_h\int\varphi_h\,d\tau_h\) is the sum of the directions
indexed by a charged birth zone \(b\) and a read level
\(t\in[b,k]\), where \(k\) is the stop level. Use the slot
comparison at the common endpoints, so the sum of these
directions is the actual interpolated derivative.

Write \(\ell_b\) for the load of a zone. Its direction
has an analytic disc of radius \(E_1/\ell_b\) in the
aligned comparison, with zero loads omitted. Terms that
read no slot of that direction are constant on the disc.
Subtract the sum over those insensitive labels before
applying Cauchy's inequality. Put \(a=(k-b-1)_+\)
and let \(m_a(h)\) count the live components of age
at least \(a\). If \(a>0\), the incidence check gives
\(m_a(h)\ge1\) on each remaining label. The age modulus bound on the
disc allows multiplication of each term by
\(\mu^{m_a(h)}\), with \(|\mu|=q_a^{-a}\). The age
reserve pays this multiplier in the sum of the
absolute masses, bounded by \(A^\sharp\). The sum over
the sensitive labels without the multiplier is therefore
at most \(A^\sharp q_a^{a}\). When \(a=0\) the
ordinary modulus bound suffices. The common integrable
majorant permits the holomorphic integration against the
tests. Cauchy's inequality, integration along the
interpolation and summation over \(t\) give
\[
 \left|\sum_h\int\varphi_h\,d(\tau_h^{B\mid A}-\tau_h^A)\right|
 \le \frac{A^\sharp}{E_1}
       \sum_{b\le k}(k-b+1)q_a^{(k-b-1)_+}\ell_b
\]
at the fixed ramps. For each label take the conjugate
polar phase of its difference measure as the test. The
left side becomes the sum of its total variations, so
the comparison holds before summing the label masses.

% Earlier Part III, ramp comparison, 17754--17907.
For the joint estimate of the ramp change and the recent
unaligned labels, keep the objects of run B at its real
center. On the aligned labels use the move that
interpolates the ramps \(t^{\mathsf A}+s(t^{\mathsf B}
-t^{\mathsf A})\) and the small production slots that read
only \(V_k\). On the common slot chart, each such
interpolated slot is bounded by
\(1+20 C_\chi\vartheta_\chi^k|s|\) times the enlarged
support indicator in run B. The indicator is fixed on
the real integration space. There are at most \(N_w^4\)
slots on the last lattice. Their product is bounded
by \(e\) times the enlarged support on the disc
\[
 |s|\le (20 C_\chi N_w^4\vartheta_\chi^k)^{-1}.
\]
In using this disc, the admitted fillings and the
labelwise modulus estimate must hold with doubled
component majorants. The weight bound pays this doubling
from the component exponent and sums the fillings.
It permits the insertion of the same arbitrary bounded
Borel phases as above. On the disc the integral with
these phases is bounded by a fixed multiple of
\(A^\sharp\). Cauchy's inequality bounds its change
from \(s=0\) to \(s=1\) by
\(C N_w^4\vartheta_\chi^k A^\sharp\). Here the disc
contains \(|s|\le2\) for sufficiently large \(k\).
The two endpoints have the B objects and, respectively,
the A and B ramps and stop-level slots.

For each run separately, multiply the common unaligned
labels born above \(n_c\) by \(\mu\). The admitted
multiplier disc has radius
\((4C_\chi^\sharp\vartheta_\chi^k)^{-1}\), with the same
modulus bound. The sum for this class is linear in \(\mu\),
so its coefficient has size at most
\(C C_\chi^\sharp\vartheta_\chi^k A^\sharp\). Taking
the supremum over the labelwise phases gives the
variation sum for these terms. This gives the joint
bound \(C_{\rm test}\vartheta_\chi^k A^\sharp\), using
the common-slot, support and modulus bounds of the
admitted ramp moves. Add it once to the fixed-ramp
comparison. The labels
born at or below the common-level threshold \(n_c\)
which are left unpaired contribute at most
\(C_{\rm term}q_a^{k-n_c} A^\sharp\),
including the nested remainders when the
terminal marked-leaf check just described
holds on the summed-pile branch. With no
nested terms take \(C_{\rm term}=4\).
Thus the total of the aligned differences and the
unaligned masses is bounded by
\[
 A^\sharp\left[
  E_1^{-1}\sum_{b\le k}(k-b+1)q_a^{(k-b-1)_+}\ell_b
  +C_{\rm term}q_a^{k-n_c}+C_{\rm test}\vartheta_\chi^k
 \right].
\]
At a fixed stop depth \(s_\circ=K-k\), the alignment level
satisfies \(n_c\le C+C\log\log(x_0+C)\) on the admitted runs.
Since \(x_0\le C_g(1+K)\), the factor \(q_a^{-n_c}\) grows
at most as a fixed power of \(\log(2+K)\). For \(q_a<\rho<1\)
this power is absorbed by \((q_a/\rho)^k\). Thus the old-label
term is \(O(\rho^{K-s_\circ})\), as is the joint ramp term if
\(\vartheta_\chi\le\rho\).
With the full load envelope, the three terms therefore
cost \(C_{s_\circ}A^\sharp(\rho^{K-s_\circ}+\Pi_K)\)
at a fixed stop, at both the source and its real centre.
This is the error for the scalar comparison below once the
actual terminal labels have the stated slot interpolation,
age incidence, marked-leaf transport and jointly controlled ramp.

\paragraph{The normalized integral at the stop.}
For the second line of \eqref{eq:comparison-capsule-CP9}, put
$v=t+\sum_i z_i f_i$. Use $v$ in run A and
$\mathfrak n_K\circ v$ in run B. Take the real centers
at \(t_1=\operatorname{Re}t\) and \(\mathfrak n_K(t_1)\).
The sources relative to the centers are \(v-t_1\) and
\(\mathfrak n_K\circ v-\mathfrak n_K(t_1)\), respectively.

% Earlier Part III, 1499--1549, admission of the local pair.
For the source-class comparison, let \(r_J\) be the
parameter-disc radius and take \(|t_1|\le r_J/4\).
Choose the residual source \(W=v-t_1\) in the class
of radius \(3r_J/8\). Its amplitude is at most
\(1.6(3r_J/8)=0.6r_J\), and its Lipschitz bound is
\(3/8\) of the bound in the radius-\(r_J\) class.
Thus the values of \(t_1+W\) lie in
\(\bar D(0,0.85r_J)\). Write \(\mathfrak n_K(z)=z+e(z)\),
with \(|e(z)|\le\Upsilon_K|z|\) on \(D(0,r_J)\).
Cauchy's estimate on the smaller disc gives
\(\sup|\mathfrak n'_K|\le1+(20/3)\Upsilon_K\le8/3\)
when \(\Upsilon_K\le1/4\). Integration along the
segment joining \(t_1+W(x)\) and \(t_1+W(y)\) gives
the Lipschitz bound of the radius-\(r_J\) class for
the residual source of run B. Its amplitude is at
most
\[
 \|W\|_\infty+\sup_x|e(t_1+W(x))|+|e(t_1)|
 \le (0.6+0.85/4+0.25/4)r_J=7r_J/8.
\]
The matched real center satisfies
\(|\mathfrak n_K(t_1)|\le(5/4)|t_1|\).
Choose the real center interval with room for this
change and the source radius in both runs. These
bounds admit the matched residual source on the
local pair, including its spatial variation.

%
For one run, take \(w\) relative to its real center and let \(Z(w)\)
be the corresponding bare sourced partition function. At level \(j\)
write \(\mathsf V_j(X;w)\) for the sum of the regular anchored
terms and the remainder produced at that level, evaluated at the
trivial field, together with the one-point gauge terms extracted
by the completed gauge steps. Group the latter on their single
\(M\)-cubes and include each at the step at which it is extracted.
The gauge term at a pin is
\(\log[\mathfrak z(a_j(y;w))/\mathfrak z(x_j)]\), with the
reference coefficient \(x_j\) of that run at its real center.
Set
\[
 e_j(X;w)=\mathsf V_j(X;w)-\mathsf V_j(X;0),
 \qquad E_u^{\rm src}(w)=\sum_X e_{K-u}(X;w).
\]
The values of these energies are independent of where the run is
stopped. For a stop at depth \(s\) put \(k=K-s\) and define
\[
 I_k(w)=
 \exp\!\left(-E_{(k)}-\sum_{j=1}^k\sum_X e_j(X;w)\right)
 \int\rho_k^w .
\]
Here \(E_{(k)}\) is the source-independent vacuum constant and the
integral includes the stored variables. Exact preservation of mass
under the first \(k\) operations gives \(\int\rho_k^w=Z(w)\).

% Earlier Part III, 11949--11967 and 18408--18424.
In the native history representation, the boundary energy is
extracted globally. A label \(\mathcal S\) produces its regular
share only on the carriers contained in its small-field regions
\(\Lambda_j(\mathcal S)\). Thus its term in \(I_k(w)\) contains
the compensating exponent
\[
 -\sum_{j\le k}
   \sum_{X\not\subset\Lambda_j(\mathcal S)} e_j(X;w).
\]
This includes the negative copy of the centered gauge term on
each omitted cube. The compensating exponent is the unit-field
evaluation of stored boundary families at their respective
levels, and their real-center vacuum constants are stored
separately. A common multiplicative normalization in a run is
chosen once at its real center and applies to every label and
source value in the comparison.

Consequently the normalized stop factor is
\[
 \widetilde Z^{[s]}(w)=\frac{I_k(w)}{I_k(0)}
   =\exp\!\left(-\sum_{u=s}^{K-1}E_u^{\rm src}(w)\right)
      \frac{Z(w)}{Z(0)}.
\]
Increasing the stop depth from \(s\) to \(s+1\) therefore
multiplies the stop factor by \(e^{E_s^{\rm src}(w)}\).
In the logarithms normalized at \(w=0\) this adds exactly
\(E_s^{\rm src}(w)\). For a fixed stop depth, its cost in the
bound for \(|\widetilde Z|\) is at most
\(\exp(C\ell^4\sum_{u<s}L^{4u})\), using the connected-domain
sum for each transferred energy. This cost is independent of the
cutoff at a fixed \(s\).

At a fixed stop depth \(s_\circ\), the one-run positive lower
bound and stability upper
bound give $\widetilde Z(0)=1$ and $|\widetilde Z|\le Q_\ell$
on the source ball, with $Q_\ell$ fixed at fixed torus. If
the radius is $r$, then on the ball of radius
$r/(2(Q_\ell+1))$ the Schwarz lemma on each complex line
gives $|\widetilde Z-1|\le1/2$. The logarithm there is
single-valued and Lipschitz on the range.

% Earlier Part III, 11907--11968 and 18287--18326.
To apply the history estimate to the stop integral, choose
one positive scalar \(a_{\mathsf c}\) for each run and real
center, in the one-run vacuum normalization. Let \(P,Q\)
be the resulting stop integrals \(a_{\mathsf B} I_k^{\mathsf B}\)
and \(a_{\mathsf A} I_k^{\mathsf A}\), at the matched
residual sources. Use these same scalars at zero residual
source. Let \(\tau_h^{\mathsf c}\) denote the normalized
history measure in the chart of its label. Its total mass
is unchanged by the chart substitution. With \(H_{\rm al}\)
the aligned labels, the difference of the scalar readouts is
\[
 P-Q=
 \sum_{h\in H_{\rm al}}(\tau_h^{\mathsf B}-\tau_h^{\mathsf A})(1)
 +\sum_{h\in H^{\mathsf B}\setminus H_{\rm al}}\tau_h^{\mathsf B}(1)
 -\sum_{h\in H^{\mathsf A}\setminus H_{\rm al}}\tau_h^{\mathsf A}(1).
\]
Pair the first sum in the common label charts and bound
each evaluation at \(1\) by the corresponding variation.
The sum estimate above therefore gives both
\(|P-Q|\le A^\sharp_\ell\varepsilon\) and
\(|P_0-Q_0|\le A^\sharp_\ell\varepsilon\), where
subscript \(0\) denotes the zero residual source at those
same real centers. In terms of the original two integrals
the common comparison scalar is \(a_{\mathsf A}/a_{\mathsf B}\).
The one-run lower bound gives \(P_0\ge B_\ell^{-1}\) in
this normalization. Subtracting the two ratios gives
\[
 \left|\frac{P}{P_0}-\frac{Q}{Q_0}\right|
 \le B_\ell\bigl(|P-Q|+|Q/Q_0|\,|P_0-Q_0|\bigr)
 \le B_\ell A^\sharp_\ell(1+Q_\ell)\varepsilon .
\]
The logarithms on the smaller ball have the same error
up to a fixed factor. This uses the history estimate for
the actual tests and stored variables at both arguments.

%
\paragraph{A uniform bound on the wider strip.}
The admitted interval \(I_K\) in \eqref{eq:flow-RF45}
leaves the stated room for the source. Its terminal
image contains \(h\pm\varrho\), and the upper slope
bound for \(x_K\) there is at most \(5/3\). The
preimages of these two values are each at least
\(3\varrho/5\) from \(\beta_K^\natural\). Thus the
one-run estimates apply at the real centers
\(|\operatorname{Re}t|<\varrho/2\).
Let \(r_\ell>0\) be the radius of the
smaller source ball just obtained from the Schwarz bound.
For \(|\operatorname{Im}t|<r_\ell/2\), take a product
disc in the \(z_i\) on which
\(\sum_i |z_i|\|f_i\|_\sharp<r_\ell/2\). The residual
source \(w=\mathrm i\operatorname{Im}t+\sum_i z_i f_i\)
lies in the smaller ball, with real center
\(\operatorname{Re}t\). There the stop factor is
nonzero and its logarithm has a uniform bound. The
local energy factors in the exact partition-function
identity are exponentials. Thus the original partition
function is nonzero on this disc and the mixed derivative
of its logarithm includes a uniformly bounded stop term
by Cauchy's inequality.

In each individual local energy the single-point frozen
share depends on at most one \(z_i\) at a plaquette.
For \(n\ge2\) it disappears from the mixed derivative.
The remaining terms contributing at depth \(s\) have
enlarged carriers meeting every support. The individual
energy bound and the separation of the supports then
give the summable majorant
\(C(\ell L^s)^4 e^{-c\mathsf d L^s/M}\)
after the connected-domain count. Cauchy's inequality on
the same fixed product disc bounds their mixed derivatives
uniformly in \(K\). Together with the stop term this
bounds \(F_{K,n}\) on the wider strip. At each finite
cutoff the bare integral is jointly entire in the bare
coefficient and the finitely many source parameters.
At nonzero values its mixed logarithmic derivative is
independent of the local logarithm branch and is
holomorphic in \(t\). This gives holomorphy across the
real-center slices used to obtain the bound.

\paragraph{Mixed derivatives and the separation of supports.}
At a plaquette barycenter at most one $f_i$ is nonzero.
Every mixed derivative in distinct $z_i$ of the pointwise
change $\mathfrak n_K\circ v$ thus vanishes from order two on.
The mixed derivative of the full logarithm for run B is
$\mathfrak n_K'(t)^nF_{K+1,n}(\mathfrak n_K(t))$. The same
argument kills the sitewise frozen pieces of the local
energies for $n\ge2$. In the other local terms the derivative
selects only enlarged carriers $X^\sim$ meeting every support.
If the separation is at least $\mathsf d>0$, this forces
$d(X)\ge \mathsf d L^s/M-C$ at depth $s$, where $d(X)$ is
the size used in the connected-domain sum.

On a shared carrier at depth \(s\), decompose the trivial-field
evaluation \(\mathsf V_{K-s}(X;w)\) into the sum of the frozen
regular terms at their separate base points and the evaluation
of the defect, remainder and gauge terms. For the frozen term at a
base point \(y\) the source enters only through \(w(p_y)\).
The full mixed derivative kills this share in each run separately.
Represent the difference of the remaining sums at the matched
sources by a local family \(G_s(X;v)\) on the common scaffold.
Its sum gives the difference by the total identities of the
comparison families. The centered difference uses
\(G_s(X;v)-G_s(X;t_1)\), with the same real centers as above.

% Earlier Part III, 11949--12058, 17333--17345, 92214--92281,
% 93020--93134 and 164537--164546.
The one-point part of \(G_s\) on a single \(M\)-cube is the sum
over its pins \(y\) of the paired log normalizer differences.
On the common scaffold these are
\[
 \log\mathfrak z\bigl(\theta_{[s]}-\zeta^{\rm B}_{(s)}(p_y)\bigr)
 -\log\mathfrak z\bigl(\theta_{[s]}-\zeta^{\rm A}_{(s)}(p_y)\bigr),
\]
up to constants independent of the source. Here \(\zeta_{(s)}\)
is the running shift at the gauge step being compared. For the
common gauge cutoff write \(\mathfrak z(a)=\int\chi e^{-aQ}\,dH\),
where \(dH\) is Haar measure and \(0\le Q\le\epsilon_s^2\) on the
cutoff support. On the strip \(|\operatorname{Im}a|\epsilon_s^2\le1\)
its real part is at least \(\cos(1)\mathfrak z(\operatorname{Re}a)\),
and differentiating under the integral gives
\(|(\log\mathfrak z)'(a)|\le\epsilon_s^2/\cos(1)\).
Apply this on the segment joining the admitted weights. The
pointwise shift bound gives \(C(1+s)(\rho^{K-s}+\Pi_K)\) for the
difference at a pin, and the fixed number of pins per cube
preserves this rate. On the plateau this is the pointwise bound
for the local pair. On a common transition layer the two shift
fields interpolate their consecutive-depth values with the same
coefficient \(\lambda(p_y)\in[0,1]\). Their difference is thus the
interpolation of the two paired differences. If the two adjacent
depths are \(s,s+1\), write \(e_s=\rho^{K-s}+\Pi_K\). Then
\(e_{s+1}\le\rho^{-1}e_s\) and \(s+2\le2(1+s)\), so the layer has
the same \(C(1+s)e_s\) bound. The consecutive reads in this argument
are those of the actual paired residual sources. At each pin we must
check the common profile and admission of both reads. For different
interpolation coefficients the additional term is their difference
times the change of one run's shift between the two depths.

%
For the compensating gauge copy in the terminal integral, the
standing aligned label supplies the same regions and cell
partition in the two runs. The fixed profile is a product over
rounded complement cells. At a pin only cells within one quarter
of their side length contribute a factor different from one.
Thus the common pin and the matching cells of the standing label
give identical profile values. In a restoration step the sequence
in force can differ from the retained label. A gauge read there
requires this check for the sequence that actually supplies it.

The argument of the normalizer on an omitted cube is the stored
gauge-weight entry at its own level \(j\), read at the terminal step.
The mark of this entry is charged in the running-weight load. Write
\(\mathfrak g_j=g_j^2(\log x_j)^{2q_\flat}\) for the load weight and
\(E_1\) for the bound on the product of this weight and the moved
slot along the dial. The change of the gauge exponent on a cube
is at most \(CE_1\epsilon_j^2/\mathfrak g_j\), by the strong
logarithmic derivative bound above. Since
\(\epsilon_j^2\le Cg_j^2(\log x_j)^{2\mathfrak p_0}\) and
\(q_\flat\ge\mathfrak p_0\), this is \(CE_1\) per cube on the
admitted dial. The negative sign of this copy is retained in the
exponent. The global gauge term uses the local energy comparison.

% Earlier Parts II, 25447--25455, and III, 92122--92176, 92889--92918.
An older-layer read also has a summable bound when the pointing
measures match, even if the active depths at the pin are far from
the nominal gauge depth. At each shared depth \(u\), the pointing
measure is the sum of the tents at the pin times the profile value,
with atoms at their source sampling points. Its mass is at most
one. Use the two real centers \(t_1,\mathfrak n_K(t_1)\) and put
\(W_A=v-t_1\), \(W_B=\mathfrak n_K(v)-\mathfrak n_K(t_1)\).
By the exact shift formula, the difference of the residual shifts
is \(W_B-W_A\) plus the sum of the paired centered increments
integrated against the common pointing measures, and the extra
finest increment of the longer run. At an atom with source value
\(v_*\), the shared increment difference obeys
\[
 |\delta b_u(v_*)-\delta b_u(t_1)|
 \le C\min\{a_{[u]},e_u/\varrho\}\,|v_*-t_1|.
\]
The two bounds follow from the one-run slopes and Cauchy's estimate
on the smaller disc, respectively. The extra increment costs
\(Ca_{[K]}\|W_A\|_\infty\). Since \(\mathfrak n_K'-1=O(\Upsilon_K)\),
integration along the source segment bounds \(W_B-W_A\) by
\(C\Upsilon_K\|W_A\|_\infty\). The same sum as in
\eqref{eq:comparison-matched-scalar-rate} then bounds the residual
shift difference by \(C\Upsilon_K\|W_A\|_\infty\). If the pointing
measures differ, their difference acts on a one-run centered
increment and requires a separate comparison.

For the centered logarithm, the full gauge arguments also contain
the real-center inverse couplings. Denote their difference at the
pin by \(d_s\). The two logarithm differences, at the source and
at its center, are bounded together by
\(C\epsilon_s^2(2d_s+\Upsilon_K\|W_A\|_\infty)\) on the same
admitted gauge strip. When the centers are the scalar centers,
their difference uses the paired scalar estimate.

The shift reads the source within two depth-\(s\) spacings of
the pin. The gauge term therefore uses the enlarged carrier test
for mixed derivatives. Since its carrier is a single cube,
once \(CM L^{-s}<\mathsf d\) it cannot meet two of the separated
supports and its mixed derivative vanishes.
The older-layer estimate thus adds only \(C_{n,\ell,\mathbf f}\Upsilon_K\)
to the mixed-derivative comparison when used on gauge pins with
the indicated locality and compared real centers.

% Earlier Part III, 92127--92151 and 164259--164265, 164493--164512.
For a sharper comparison, consider a pin \(p\) in a common native
zone of actual depth \(m<K\). Its profiles are zero for \(u<m\),
possibly partial at \(m\), and one for \(u>m\) among the available
inputs, including the extra finest input of the longer run. Use
this calculation where that pattern and the common tents apply
to the full gauge read, and the re-centered shift formula uses
the same increment functions as the frozen matching, as in
\eqref{eq:flow-RF52}. Let
\(D_u(z)=\delta b_u(z)-\delta b_u(t_1)\)
for the shared increments, and let \(D_K\) denote the centered
extra increment at the matched offset. Terminal matching at
\(v(p)\) and at \(t_1\) gives
\[
 W_B(p)-W_A(p)+\sum_{u=0}^K D_u(v(p))=0.
\]
Replace each tent sample by its value at the pin in the residual
shift difference. The identity cancels the full sum, leaving the
omitted increments for \(u<m\) and the omitted fraction at \(m\).
On the smaller source disc their total is at most
\(C(1+m)e_m\|W_A\|_\infty\), since \(e_u\le e_m\) for \(u\le m\).
For the error of this replacement, the sample distance is
\(O(L^{-u})\) and the increment difference has derivative bounded
by \(Ce_u/\varrho\) on the shared depths. Hence the error is at most
\[
 C\operatorname{Lip}(v)\left(\sum_{u=m}^{K-1}e_u L^{-u}
                                  +a_{[K]}L^{-K}\right)
 \le C\operatorname{Lip}(v)e_m L^{-m},
\]
using \(L^{-1}<\rho\) and bounded \(a_{[K]}\). The real-center
coefficient is compared separately as above. If the extra input
is not fully present, the omitted fraction of \(D_K(v(p))\)
must also be included.

The logarithm on a cube is a sum over individual pins. At one
pin the source halo has radius \(O(L^{-m})\). Once this is below
a fixed fraction of \(\mathsf d\), the halo misses at least one
support and the mixed derivative at that pin vanishes. On the
remaining pins with \(s\le m\), both depths are bounded in terms
of the separation.
The bound \((1+m)e_m=O_{\mathsf d}(\rho^K+\Pi_K)\) then gives
the sharper mixed-derivative rate after the finite cube count and
Cauchy's estimate. The identification of the actual zone and the
profile pattern is needed for this use of the rate.

For the remaining two parts use the total identities at the
unit field, where the regular background is the unit field
and the current is zero. The source defect \(D\) contributes
the sum over base points of its comparison representative
with weight \(E_0\mathsf l_s\). The remainder contributes one
representative per carrier with weight \((C_t\theta_{[s]})^{-1}\).
These are the weights of the native comparison units, with
\(\mathsf l_s\le C(1+s)\) and \(\theta_{[s]}\ge h\).
On the admitted source polydisc the \(A\) and \(R\) envelopes
therefore bound these two parts together by
\(C(1+s)(1+d(X))(\rho^{K-s}+\Pi_K)e^{-\kappa^c d(X)}\).
Here \(\kappa^c\) is the decay rate in the comparison norm,
and the factor \(1+d(X)\) counts the base points in the defect.
Absorb that factor by taking \(\kappa_e<\kappa^c\). This gives
the following bound for the regular parts, and for the complete
family at depths where the gauge pins satisfy the nominal-depth
estimate. Use it on the shared depths through \(K/2\):
\[
 |G_s(X;v)|\le
 C(1+s)(\rho^{K-s}+\Pi_K)e^{-\kappa_e d(X)} .
\]
We must verify the comparison envelopes and the pin estimates
for the actual paired trajectories on this polydisc. Summing by
an anchor cube and reserving exponential decay after the
connected-domain count costs at most
\(C(\ell L^s)^4e^{-c\mathsf d L^s/M}\).
This cost times \((1+s)\rho^{-s}\) is summable. Thus Cauchy's
estimate gives \(C_{n,\ell,\mathbf f}(\rho^K+\Pi_K)\)
for the mixed derivative from these depths. For deeper
shared depths, and for the one extra depth of run B, use
the two individual local energy bounds. The same carrier
count gives $C_{n,\ell,\mathbf f}e^{-c'\mathsf dL^{K/2}/M}$.
For the stop logarithms use Cauchy's estimate on the smaller
polydisc, with the fixed-torus radius above. In particular
the second line of \eqref{eq:comparison-capsule-CP9} follows
if the stop error $\varepsilon$ has a summable bound and the
shared-depth estimate holds. At $t=0$ the identity for
$N_K$ cancels the derivative factor in the comparison of
$N_K^{-n}F_{K,n}(0)$, with no additional error from the
matching derivative.

For summability one can stop the shared-depth comparison
earlier. For any $q>1$ take $s_K$ with
$c\mathsf d L^{s_K}/M\ge 2q\log K$ and
$s_K=O(\log_L\log K)$. The individual local bounds give a
summable tail above $s_K$, including the carrier count.
This leaves the comparison of the terminal normalized
integral at the fixed stop and the shared-depth estimate through
\(s_K\). One can also transfer the shallower totals into the stop
factor at depth \(s_K\) using the exact identity above. In that
form the stop factor must be compared at the moving depth.
For sources near their real centers, the Schwarz estimate
for each local energy and the connected-domain sum give
\(\left|\sum_{u<s} E_u^{\rm src}(w)\right|
 \le C_\ell L^{4s}\|w\|_\sharp\).
The stop factor at depth zero is close to one on a fixed
source ball. The exponential of the transferred energies
is also close to one on the ball of radius \(c_\ell L^{-4s}\).
At \(s=s_K\) the resulting mixed Cauchy estimate costs only
a fixed power of \(\log K\), provided the comparison of the
two moving-stop factors is established on that ball.

\begin{proposition}[The conditional terminal-coupling limit]
\label{thm:terminal-coupling-limit-v7}
These inputs imply uniform convergence \(F_{K,n}\to S^{[t]}_{\ell,n}\)
on the smaller strip and \(\Xi_K\to\Xi_\infty\).
The latter is a real bi-Lipschitz homeomorphism from \(\mathsf J\)
onto \(I\supset[h-\varrho/6,h+\varrho/6]\). Set
\[
 S^{(\ell,x)}=S^{[\Xi_\infty^{-1}(x)]}_\ell,\qquad x\in I .
\]
If \(t_K\in\mathsf J\) and \(\Xi_K(t_K)\to x\), then
\(F_{K,n}(t_K)\to S^{(\ell,x)}_n\), without a subsequence.
\end{proposition}
\begin{proof}
Cauchy's estimate on the wider strip bounds \(F_{K,n}'\) uniformly.
Thus replacing \(\mathfrak n_K(t)\) by \(t\) in \eqref{eq:comparison-capsule-CP9}, and its
derivative factor by one, costs \(C\eta_K\). The resulting increments
are summable, uniformly on the smaller strip. Their limit is
holomorphic and continuous on \(\mathsf J\). The coordinate increments
are summable too. Its real Lipschitz bounds pass to the limit;
the endpoint bounds give the asserted interval.
Uniform coordinate convergence gives \(\Xi_\infty(t_K)\to x\);
the lower Lipschitz bound forces
\(t_K\to\Xi_\infty^{-1}(x)\). Uniform convergence of \(F_K\)
and continuity of its limit finish the proof.
\end{proof}

This gives the terminal-coupling limit in part (d) of \cref{thm:continuum-main} under the stated comparison premises;
infinite-volume uniqueness requires a separate argument.
The depth coordinate
\(\Xi_{\infty,s}(t)=\lim_Kx_{K-s}^\natural(\beta_K^\natural-t)\)
requires the corresponding fixed-depth comparison. If its envelope is
\[
 \lambda s-C\le\Xi_{\infty,s}(t)-\Xi_{\infty,0}(t)
                         \le\Lambda s+C,
\]
then \(\Xi_{\infty,s}(t)^{-1/2}\) decreases in order as
\(s^{-1/2}\); at length \(L^{-s}\) this is logarithmic running.
The reference value \(g_{[s]}=\theta_{[s]}^{-1/2}\) is a calibration,
not this arbitrary member's running coupling. Common-prefix matching
must also control
\(\hat t_{K,s}=\beta_{K-s}^\natural(g_{[s]})-\beta_K^\natural(g)\).
If it converges to \(\hat t_s\), exact finite dilation passes to
\(S^{[t]}_\ell[g]=S^{[t+\hat t_s]}_{L^s\ell}[g_{[s]}]
(\mathbf f^{[s]})\), wherever both offsets are in their windows.
% Earlier Part III, 111457--111525 and 120741--120865.
For the first-passage application let \(K'\) be the stopping
index of the run whose choices follow the coupling, at bare
inverse coupling \(\beta\). The scaffold comparison in E4
must give a live frozen run of length \(K'\) from the same
bare value, in the admitted real tuning class. Write \(y\)
for its terminal inverse coupling. Write \(\omega(g)\le1\)
for the terminal cost in \eqref{eq:flow-RF37}. The
first-passage window then gives
\[
 h-\omega(g)\le y\le h+B_\sharp+\omega(g).
\]
Use the increasing inverse of the frozen terminal coupling
as a function of the bare value in this class. Its value at
\(h\) is \(\beta_{K'}^\natural\). The inverse tuning estimate
has Lipschitz constant at most \(3/2\), so
\[
 -\tfrac32\omega(g)
 \le\beta-\beta_{K'}^\natural
 \le\tfrac32\bigl(B_\sharp+\omega(g)\bigr).
\]
For \(\varrho\ge6(B_\sharp+1)\) this puts the offset
\(\beta_{K'}^\natural-\beta\) in \(\mathsf J\).
The identification uses uniqueness of the admitted
real frozen run at terminal value \(y\).

Conversely, for sufficiently large \(K\) and
\(\beta=\beta_K^\natural-t\) with \(t\in\mathsf J\),
the forward comparison at its
first-passage index \(K'\) gives
\(|\beta_{K'}^\natural-\beta_K^\natural|\le\varrho/2\).
The matched-depth estimates below give a lower spacing
bound between the tuned bare centres. To see this, compare
their couplings at depth \(\min(K,K')\) by telescoping in
the run length. The geometric error and the summed floor
give a fixed bound there. The remaining depth interval of
the longer run then contributes at least
\(\lambda|K-K'|-C\) by the flow bound. This gives
\(|K-K'|\le C(\varrho)\).
The comparison between the ordinary and the frozen
bookkeeping, including admission of the frozen run,
is the input to this argument.
Exact lattice dilation uses \(f^{[b]}(y)=f(L^{-b}y)\), without
an observable prefactor, and \(T_\ell\mapsto T_{L^b\ell}\).
Conditional offset limits and depth shifts pass through that identity
only when the offsets converge and both torus floors hold.
The change-of-reference claims retain these matching premises.
The tuned bare increments must also satisfy:
\begin{equation}
 \beta_{K+1}^\natural(g)-\beta_K^\natural(g)\longrightarrow c_0,
 \qquad
 \beta_{K+b}^\natural(g)-\beta_K^\natural(g)\longrightarrow bc_0
 \quad(b\in\mathbb Z\text{ fixed}).\label{eq:comparison-capsule-CP9a}
\end{equation}
The second assertion follows by finite telescoping from the first.
The following proof uses the stronger comparison hypotheses.
Write \(x^{[i]}_{[s]}\) for the tuned depth-\(s\) coupling in the run
of length \(i\), with \(x^{[i]}_{[0]}=h\) and
\(x^{[i]}_{[i]}=\beta_i^\natural\). Require
\(x^{[i]}_{[v]}-x^{[i]}_{[u]}\), for \(u<v\), to lie in
\([\lambda(v-u)-C,\Lambda(v-u)+C]\), and
\[
 |x^{[i+1]}_{[s]}-x^{[i]}_{[s]}|
 \le C\big(L^{-c(i-s)}+s\Pi_i\big),\qquad
 \sup_{s\ge1}s\sum_{i\ge s}\Pi_i<\infty .
\]
These constants and the last bound are fixed before \(\varrho\).
Comparing two bare centers at the same depth \(K'\), then telescoping
in the run length, bounds their extra discrepancy by
\(C\sum_{i\ge K'}L^{-c(i-K')}+CK'\sum_{i\ge K'}\Pi_i\le C'\).
The spacing bounds therefore give a uniform bound \(A\) on
\(\alpha_K=\beta_{K+1}^\natural-\beta_K^\natural\).
Choose \(\varrho>20A\).

Each cluster value \(\alpha\) of \(\alpha_K\) gives
\(S^{[t]}_\ell=S^{[t-\alpha]}_{L\ell}(\mathbf f^{[1]})\)
by exact finite dilation and uniform moving-offset convergence.
Two cluster values \(\alpha^-<\alpha^+\) give an offset period
\(\omega=\alpha^+-\alpha^-\) on an open interval. Holomorphy in the
offset extends that identity throughout the overlapping strips.
Repeated blocking can then restore every offset modulo \(\omega\)
into \([0,\omega)\); all intermediate offsets stay in \(\mathsf J\)
by the choice of \(\varrho\).
Require, additionally, one fixed positive shape \(\Psi\), a floor
\(\ell_0\), and \(r_0>0\), such that
\[
 0<\mathfrak m(r)\le S^{[t]}_{\ell,2}(\Psi_r)\le\mathfrak M(r),
 \quad \mathfrak M(r)\longrightarrow0,\qquad
 t\in\mathsf J,\ \ell\ge\ell_0,\ 0<r\le r_0 .
\]
Normalization by the two fixed supremum norms is understood.
The required uniformity in offset, shape placement and arbitrarily
large tori is part of the stated tuned two-point comparison.
For \(r_m=L^{-m}r_0\), blocking gives
\(\mathfrak m(r_0)\le\mathfrak M(r_m)\to0\), a contradiction.
Thus \(\alpha_K\) has one cluster value and converges.
Finally \eqref{eq:comparison-capsule-CP2}--\eqref{eq:comparison-capsule-CP3}, backward drift and \(x_K^\natural=h\) make the
average nonlinear increment error tend to zero, so
\(\beta_K^\natural/K\to c_0\). Cesaro averaging identifies the unique
cluster value with \(c_0\). Without these two-run and uniform-envelope
premises, \eqref{eq:comparison-capsule-CP9a} requires a separate comparison estimate.
With \eqref{eq:comparison-capsule-CP9a}, exact finite dilation and the moving-offset limit give
\[
 S^{[t]}_\ell(\mathbf f)
       =S^{[t-bc_0]}_{L^b\ell}(\mathbf f^{[b]})\qquad(b\ge1),
\]
when both offsets lie in their windows and the torus floors hold.
Indeed the same bare coupling has the second offset
\(t-[\beta_{K+b}^\natural-\beta_K^\natural]\), which converges to
\(t-bc_0\). Thus the original blocking target is preserved without
claiming it follows from one-step potential estimates.

Comparison of distinct offsets requires the following additional two-point assumption.
Fix a nonempty class of nonnegative \(C^2\) pairs: each test is supported
in a radius-\(r\) ball, the two supports have distance in \([r,6r]\),
and each test satisfies
\(\int f\ge\vartheta r^4\|f\|_\infty\) and
\(\operatorname{Lip}f\le\|f\|_\infty/(\vartheta r)\).
Define \(\sigma\approx_2\sigma'\) by an error
\(Cr^2\|f\|_\infty\|h\|_\infty\), uniformly on this class at all
sufficiently small \(r\). Set
\((D_\mu\sigma)_2(f,h)=\sigma_2(f(\cdot/\mu),h(\cdot/\mu))\).
Require \(\sigma_t\approx_2D_\mu\sigma_{t'}\), with
\(\mu=\exp((t-t')\log L/c_0)\), and uniform two-point upper and
lower envelopes of order \((\log(1/r))^{-2}\).
A germ cannot be fixed by \(D_\mu\), \(\mu\ne1\). For \(\mu>1\),
apply its error bound to one fixed rescaled shape, with normalized
value \(Q(r)\), at radii \(r\mu^{-a}\).
The upper envelope gives \(Q(r\mu^{-a})\to0\), so
\[
 Q(r)\le C\sum_{a\ge1}(r\mu^{-a})^2
                     =\frac{Cr^2}{\mu^2-1}.
\]
This contradicts the logarithmic lower envelope for arbitrarily
small \(r\). For \(\mu<1\), invert the dilation. Thus equality of two
different-offset limits would contradict rigidity. The comparison
retains its shape, placement, offset and torus-floor quantifiers;
it is an additional premise, not a consequence of matching alone.

\subsubsection*{Why a geometric relabelling needs another construction}

Take three test tuples of the same fixed correlation order, for example
three pairs, with observable values
\[
 F_K(t)=(1,\,1+t,\,1+\Pi_K)
       =\alpha+\beta t+\gamma\Pi_K,\qquad
 \Pi_K=K^{-q},\quad q>1 .
\]
Here \(\alpha=(1,1,1)\), \(\beta=(0,1,0)\),
\(\gamma=(0,0,1)\), with determinant one.
On \(|t|\le1/4\) the values are positive and uniformly bounded.
With \(\Xi_K(t)=h-t\) and \(\mathfrak n_K=\mathrm{id}\), all
summable matching inequalities hold; their increments are
\(O(K^{-q-1})\).
Suppose a reparametrization \(t\mapsto\tau_K(t)\) and one common
normalization \(a_K(t)\), bounded above and away from zero, made
their increments geometrically small. At a fixed \(t\), geometric
summability forces \(a_K\), the first component, to converge
geometrically to a nonzero limit. The third component minus the
first is \(a_K\Pi_K\), tends to zero, and would also be geometrically
small. This contradicts its polynomial lower bound.
The same argument allows arbitrary reparametrizations: their value
does not enter either the first component or the third-minus-first.
Only a common, nonvanishing field normalization is required. This is
a counterexample to the proposed implication from numerical inequalities,
not an assertion that these three numbers already form a physical
Yang--Mills state. Thus those inequalities alone do not supply the
stronger second-coordinate claim.
% Earlier Part III, 99409--99554 and 109230--109232.
The construction with filed remainders uses its
own running coordinate. At the filing step add the
vacuum-subtracted pointed shares to \(E\). Extract
their Wilson coefficient on the filing lattice
and subtract the corresponding counterterm.
The vacuum values go into the scalar normalization
and the shares still in the filing band remain
in the unfiled part. The shifted Wilson coefficient
and the gauge-fixing coefficient both use the new
running shift. The filing identity is applied
within the same admitted production. If \(x_K\)
is the coupling before this transfer, the two
zero-source coordinates at the same bare datum satisfy
\[
 \hat x_K=x_K-
 \sum_{i:\,i''\le K}
 \operatorname{coeff}_{A,i''}\bigl(\mathsf{RF}[R^{(i)}]\bigr).
\]
Here \(i''\) is the actual filing level and
\(\operatorname{coeff}_{A,i''}\) denotes the coefficient
of the vacuum-subtracted total at that level.
On a common admitted complex offset disc, denote
this sum by \(T_K(t)\). The filing norm and
coefficient estimate bound each summand by
\[
 C\theta_{i''}w_{i''-1,i}\|R^{(i)}\|
 \le C\theta_i^{\,1-\kappa_0/2}.
\]
The inequality uses the comparison of radii at
the filing depth and the one-run remainder bound.
Summing over the depths of the input remainders,
where the inverse coupling grows at least linearly,
gives
\[
 \sup_t |T_K(t)|\le C h^{\,2-\kappa_0/2}
 \qquad(\kappa_0>4).
\]
Cauchy's estimate gives the same small bound for
the derivative on a fixed smaller disc. Thus
the coordinate \(\Xi_K(t)-T_K(t)\) of this
presentation keeps a lower real Lipschitz bound
when \(\Xi_K\) has one and the coupling is small.
The ordinary pin at \(h\) is then at
\(h-T_K(0)\) in the filed coordinate at that
bare datum.

% Earlier Part III, 99275--99342, 122691--122773
% and 122793--122872.
For the two-cutoff estimate align the remainders
by birth depth \(s\) from the terminal lattice.
If the two runs use the same scaffold, the
filing delay \(a_s\) at that depth is fixed
by the scaffold. The remainder is included
precisely if \(a_s\le s\), so the inclusion
test is common on shared depths. Terms
at an extra bare depth can be bounded by
their one-run size. If the schedules differ,
their inclusion tests must be compared as well.

On the admitted paired paths write \(\chi_s\)
for the termwise difference norm of the two
remainders at birth depth \(s\), and \(\iota_s\)
for the relative datum error of their filed
read. The termwise filing estimate and the
Wilson coefficient bound give
\[
 |\delta T|
 \le C\sum_{\substack{s\ \mathrm{shared}\\a_s\le s}}
       \bigl(\theta_{[s]}\chi_s+
              \iota_s\theta_{[s]}^{\,1-\kappa_0/2}\bigr)
      +C\theta_{[K+1]}^{\,1-\kappa_0/2}.
\]
Here the coefficient estimate at the filing
level uses the crossing weight to return to
the unit at birth. Each filed remainder
occurs once in the sum. The termwise norm
uses the actual filed slices of the runs.

In the comparison pinned in the filed
coordinate, the filed response rows bound
\(\theta_{[s]}\chi_s\) by
\(C\widehat r_s\rho^{K-s}\). The datum term
has the same bound on the admitted paths.
Using \(\widehat r_s\le
C\rho^{-(K-s)/2}\widehat r_K\) and summing
the geometric factor gives an \(O(\widehat r_K)\)
bound for \(\delta T\), including the bare end.
Thus a geometric comparison at each fixed
depth may still yield a polynomial error
for the sum of the transferred coefficients.
If the paired reads are available only
through depth \(K/2\), their contribution
is bounded by a geometric factor
\(\rho^{K/2}\) times the summable one-run
weights. The other depths are bounded
separately by the one-run coefficient sum,
which on that depth interval is
\(O_g(K^{2-\kappa_0/2})\).
With \(\widehat r_K\) of order
\(\theta_{[K]}^{\,2-\kappa_0/2}\) times a
power of its logarithm, the bound is
summable for \(\kappa_0>6\).

On a common admitted terminal interval
for the filed runs, let \(u\) be the
filed terminal coupling and write
\(\mathcal T_K(u)\) for the sum \(T_K\)
at its corresponding bare datum. The
pre-transfer coordinate of this same
presentation is \(X_K(u)=u+\mathcal T_K(u)\).
On a smaller interval it has a uniform
positive lower Lipschitz bound by the
one-run derivative estimate and the
admitted inverse for \(u\).
For a pre-transfer target \(x\), put
\(u_K=X_K^{-1}(x)\). Comparing
\(X_{K+1}(u_K)\) with \(x\) gives
\[
 |u_{K+1}-u_K|\le
 C\sup_u|\mathcal T_{K+1}(u)-\mathcal T_K(u)|.
\]
This transfers the summable bound to
the change of pin within the same
presentations. Its use for the separately
constructed ordinary run requires the
matching of its production and source
normalization to that presentation.

\subsubsection*{The optional Ward target on arbitrary defining sequences}

In this section \(S_{j,n}\) and \(S_n\), \(n\ge2\), are
\emph{connected} forms. Require convergence on every smooth compact
separated tuple on \(\mathbb R^4\), along the specified defining
sequence. On every fixed collection of separated compact support
neighborhoods require uniform local multilinear \(C^1\) bounds.
Mollification inside slightly enlarged separated neighborhoods then
extends the forms and their convergence to compact \(C^1\) tuples:
first approximate each entry, use the uniform telescoping bound,
and then pass to the limit for the smooth approximants.
Assume full signed-permutation symmetry \(W_4\) of the limit.
The specified lattice probability laws realize these connected forms;
the mixed cumulants below use those laws.
For torus approximants a fixed \(R\) must fit in one common chart
for all sufficiently large \(j\).
For \(\omega\in\mathfrak{so}(4)\), set
\(\delta_\omega f=-(\omega x)\cdot\nabla f\) and
\[
 D^s_{j,\omega}f(x)=
 \sum_\rho(\omega x)_\rho
       \frac{f(x+a_je_\rho)-f(x-a_je_\rho)}{2a_j}.
\]
Here are actual additional Ward inputs, independent of the angular route.
Choose smooth \(\chi_R=1\) on \(B_R\), zero outside \(B_{2R}\),
with derivatives bounded by \(C_rR^{-r}\).
\phantomsection\label{notdef:ward-insertion}
For every \(j,\omega,R\), supply fixed annular local functions
\[
 \mathfrak B_{j,\omega,R}=\sum_{\Delta\in\mathcal A_{j,R}}
                              \mathfrak B^\Delta_{j,\omega,R}.
\]
There are at most \(CR^4\) physical unit boxes; each summand is a
specified gauge-invariant Borel function of links in a fixed enlargement
of its box, linear in \(\omega\), supported near
\(R-Ca_j\le|x|\le2R+Ca_j\).
The functions may contain the divergent Wilson weight.
They are supplied by the actual change of variables, independently of
the tests, not defined as their desired rotation defect.
Require the exact finite identity
\begin{equation}
 \sum_i S_{j,n}(f_1,\ldots,D^s_{j,\omega}f_i,\ldots,f_n)
  =\mathfrak U_{j,\omega,R}(\mathbf f)
   +\langle\mathfrak B_{j,\omega,R};
                  \mathcal O_j(f_1);\ldots;\mathcal O_j(f_n)\rangle^T .
 \label{eq:comparison-capsule-CP10}
\end{equation}
In \eqref{eq:comparison-capsule-CP10}, \(\mathcal O_j(\mathbf f)\) denotes the separate list
of cumulant arguments, not their coincident product. The bulk functional
\(\mathfrak U\) must also be the specified correlation from that
same Ward construction; it is not an arbitrary error chosen after the
boundary functions. \phantomsection\label{hyp:ward-bulk}
For every fixed sufficiently large \(R\), require
\(\mathfrak U_{j,\omega,R}\to0\) and the analogous sum with
\(D^s_{j,\omega}f_i+\delta_\omega f_i\) to tend to zero.
These include the actual typed anisotropy and volume-transfer estimates
where needed, and are not proved by scalar source bounds.

Multilinear telescoping and local \(C^1\) continuity differentiate the
limit's rotation orbit on smooth tuples. Applying the assumed convergence
at the fixed derivative tests in \eqref{eq:comparison-capsule-CP10}, without differentiating a
cutoff limit, gives
\begin{equation}
 \lim_j\langle\mathfrak B_{j,\omega,R};\mathcal O_j(\mathbf f)\rangle^T
     =-\delta_\omega S_n(\mathbf f).
 \label{eq:comparison-capsule-CP11}
\end{equation}
It is independent of \(R\).
Consequently the condition in part (e) of \cref{thm:continuum-main}
\begin{equation}
 \liminf_{R\to\infty}\limsup_{j\to\infty}
 \left|\langle\mathfrak B_{j,\omega,R};
                      \mathcal O_j(\mathbf f)\rangle^T\right|=0
 \label{eq:comparison-capsule-CP12}
\end{equation}
is equivalent to vanishing of this derivative.
The inner limit always has \(R\) fixed.

For completeness, one nonzero generator suffices.
Assume also the invariant one-point convention \(c\int f\) and
zeroth moment one, or use the centered functional. Let \(G\) be the
closed subgroup of orthogonal transformations preserving the connected
forms. Closedness follows from \(C^1\) continuity on separated tuples.
Equivalently, define \(G\) from the centered moment forms obtained
by the finite sum over partitions of these cumulants; the inverse
partition formula gives the same group. The invariant one-point
function \(c\int f\) makes the uncentered moment group the same.
Let \(\mathcal A\) consist of skew matrices whose entire one-parameter
group lies in \(G\). It is closed under real scalar multiplication
and addition: the finite-matrix product formula
\[
 e^{t(A+B)}=\lim_{r\to\infty}(e^{tA/r}e^{tB/r})^r
\]
follows by expanding each factor to second order and telescoping the
product error, which is \(O(r^{-1})\). Closedness of \(G\) then applies.
It is also stable under conjugation by \(W_4\).
For \(a<b\), sign matrices \(D_\epsilon=\operatorname{diag}(\epsilon)\)
give the projection
\[
 2^{-4}\sum_{\epsilon\in\{\pm1\}^4}
       \epsilon_a\epsilon_b D_\epsilon A D_\epsilon
                         =A_{ab}E_{ab}.
\]
Thus any nonzero element of \(\mathcal A\) supplies a coordinate-plane
generator; permutations supply all of them.
Successive plane rotations reduce any determinant-one orthogonal
matrix to the identity column by column. Hence \(G\) contains
\(\operatorname{SO}(4)\), and its coordinate reflections give all
\(\operatorname{O}(4)\).
The condition in \eqref{eq:comparison-capsule-CP12} for a nonzero \(\omega\), on every smooth tuple,
gives its whole one-parameter group by applying the zero derivative
also at each rotated tuple. Conversely full invariance makes \eqref{eq:comparison-capsule-CP11}
zero for every \(\omega\).
No Lie-subgroup theorem is needed. Indeed an infinite closed \(G\)
has nonidentity \(g_r\to I\); the normalized skew logarithms have a
nonzero subsequential limit \(A\).
Integer powers \(g_r^{\lfloor t/\|\log g_r\|\rfloor}\to e^{tA}\)
supply the required one-parameter group. Therefore such a
\(W_4\)-invariant group is finite or all of \(\operatorname{O}(4)\).

\clearpage
\subsubsection*{Infrared meaning and the Hilbert-space consequence}

A uniform mixed-insertion estimate
\(|\langle\mathfrak B^\Delta;\mathcal O_j(\mathbf f)\rangle^T|
\le C D^{-4-\delta}\), \(\delta>0\), at distance \(D\) from the
fixed supports, implies \eqref{eq:comparison-capsule-CP12}: its shell sum is at most
\(CR^4R^{-4-\delta}\).
Decay \(D^{-4}\) alone yields only a bounded sum.
It also suffices that the sum of absolute mixed insertions tend to
zero along one sequence \(R_l\to\infty\), after \(\limsup_j\).
Ordinary per-ball density bounds do not give either estimate for
these different, weight-bearing insertions.
If an older survival argument is used, its event
\(\mathcal L_{\mathfrak c,k,a,K}\) means that cell \(\mathfrak c\)
is live at level \(k\), its component lineage starts at \(k-a\),
and it remains live at every level \(k,\ldots,K\).
Its weighted law is an additional history input.

Assume the relevant centered or recentered moment functional also
satisfies the HB quotient hypotheses. Once full scalar rotation
invariance is supplied, spatial
\(R\in\operatorname{O}(3)\) acts on the HB quotient by
\(V_R[F]=[(R\oplus1)F]\).
It preserves the positive form and null space, has inverse \(V_{R^{-1}}\),
fixes \(\Omega\), commutes with time evolution and conjugates a spatial
translation by \(x\) to one by \(Rx\).
Local \(C^1\) continuity proves strong continuity on dense test vectors,
and unitarity extends it to the completion.
The joint spectral measure is consequently covariant under
\((E,p)\mapsto(E,Rp)\).
For limits for which HB's cube inequalities have been established,
they give \(E\ge|p|\): use a countable dense
set of rotations, remove their common null set, and approximate a
rotation carrying any given \(p\) to \(|p|e_1\).
Centering is the map \(C_m\) proved in
\cref{ch:euclidean-hilbert-bridge} and changes none of these
conclusions. In the alternative notation,
\(\mathsf{sh}_\xi=C_{-m_\xi}\),
\(m_\xi(f)=c_\xi\int f\), so its inverse is \(C_{m_\xi}\).
No operator for a general time-mixing Euclidean rotation is asserted.
External OS reconstruction, F4 identification of the full theory and
any mass-gap estimate remain separate. The periodic-state angular
proof retains its own mirror, cone and ultraviolet interface; the
optional Ward premises above do not silently enlarge that state class.


% END MODULE chapters/comparison_capsule_v7.tex

% BEGIN MODULE chapters/weighted_matched_domains.tex
\subsection{Weighted derivatives and matched presentations}
\label{ch:weighted-matched-domains}

This section formulates weighted derivative estimates and comparisons
of matched presentations on their specified domains. These stronger
estimates are separate from the ordinary scalar source argument and
its angular continuation.  We prove the analytic, convolution and normalization steps.
The remaining input is the corresponding extension of the actual
formation and history maps, with the weights and reserves specified
here.  No survival probability is used as an enlarged input radius.

\subsubsection*{A direction is a list, including its own coefficient}

Fix a real frame and all its reference operations.  A direction at old
level $j$ contains a pointed regular family $\dot E^{(j)}$, any
awaiting remainder $\dot R^{(j)}$, and the specified weight or vertex
amplitude.  Its regular family has the actual gauge covariance and
the covariance of its pointed total required by the chosen channel.
Its own coefficient is differentiated with that family.

For the scalar channel let $T$ be the fixed physical Hessian
readout.  If $E^\lambda$ is an admitted holomorphic family, put
$b^\lambda=T[E^\lambda]$.  With the frame fixed the map $T$ is
linear and independent of $\lambda$, so
\begin{equation}
 \dot b=T[\dot E],\qquad
 {d\over d\lambda}
       [E^\lambda-b^\lambda h_zA]\big|_0
                         =\dot E-\dot b\,h_zA.
 \label{eq:weighted-matched-domains-WM1}
\end{equation}
The tent $h_z$ and background convention are the assigned frozen ones.
If the readout itself changes between two frames, its difference is
an additional term; it is not covered by \eqref{eq:weighted-matched-domains-WM1}.
For a tensor channel the scalar coefficient is replaced by the full
assigned vector of counter-insertions.  Its variation is taken at the
same time as that of its own family.

This rule has analytic content at a single old base.  Marginal
cancellation is applied to the whole covariant member, including
its own counterterm, before the bracket interpolation.  A
noncovariant defect is a separate input with its own scale-small
norm.  Cauchy differentiation of a partition function cannot give
that defect the cancellation of another member.

The norms below are norms at the birth level, on the actual
hereditary cores.  Write $\mathsf e_j,\mathsf r_j$ for the regular
and remainder norms.  Let $\tau_j$ be the amplitude and $\mathsf k$ its
fixed generating test.  The norms also cover the associated
frozen families at every admitted constant source $\zeta$
and every constant generator $h_0$ in the channel with
$|h_0|\le\|\mathsf k\|_\infty$.  The amplitude norms below
are the corresponding suprema over $\zeta$.  This frozen
coverage is part of the weighted input class; it is not
inferred from a norm of only the varying test.  Put
\begin{equation}
 \mathsf z_j=\|\tau_j\|_\infty\|\mathsf k\|_\infty,\quad
 \mathsf b_j=\|\tau_j-\tau_{j-1}\|_\infty\|\mathsf k\|_\infty,\quad
 \mathsf z_n^\sharp=\sum_{j\le n}L^{-(n-j)}
                              (\mathsf z_j+\mathsf z_{j-1}).
 \label{eq:weighted-matched-domains-WM2}
\end{equation}
Use $\mathsf z_{-1}=0$ and the actual first-entry convention.
The symbols here measure this direction, not the deviations of the
running coupling.

Let $\varsigma_\nu$ be the actual forgetting weights and let
$w_{n,j}=(\alpha_n/\alpha_j)^2$ be the squared radius ratios.
If a remainder born at $j$ is refiled at $j''$, define
\begin{equation}
 \begin{aligned}
 {\cal L}_n={}&\sum_{j\le n}\varsigma_{n+1-j}\mathsf e_j
       +\sum_{j\le n\le j''}w_{n,j}\mathsf r_j\\
 &+\mathfrak g_n\left[
       \mathsf z_n^\sharp+
             \sum_{j\le n}L^{-c(n-j)}\mathsf b_j\right].
 \end{aligned}\label{eq:weighted-matched-domains-WM3}
\end{equation}
The second sum includes precisely the remainders still awaiting
refiling.  A bare stored insertion contributes its stated
additional size $C\mathfrak g_0N_{\rm bare}\|\mathsf k\|_\infty$ at
level zero.  The size $N_{\rm bare}$ is counted from that
insertion and is not silently set to one.

The vertex response factor is of order
$\mathfrak g_n=x_n^{-1}(\log x_n)^C$ on the
coupling-independent analytic route.  For a member analytic only
on a running field radius the corresponding gain is generally the
square of the fluctuation-to-radius ratio.  These two factors
are different.  The parity calculation below explains the distinction.

\subsubsection*{Weighted extension and the derivative rows}

Attach one dial $\lambda_{n,k'}$ to the entry when it is used in
zone $n$ at step $k'$.  A useful actual class requires a joint
extension with axis radii
\begin{equation}
 |\lambda_{n,k'}|<c/{\cal L}_n .
 \label{eq:weighted-matched-domains-WM4}
\end{equation}
When a load is zero its direction has no contribution and that
axis can be omitted.  The class also specifies the joint
multi-axis budget; individual axis discs alone do not give
their Cartesian product.

The required analytic assertion has two parts.  In the small-field
map the differentiated potential norm must retain
$C\varkappa_k{\cal L}_k$, where
$\varkappa_k$ is the actual fluctuation-to-radius ratio.
For a large-field output, setting one dial to zero must change
its assigned term by a bound proportional to the share of
histories retaining that zone, on the same radius in \eqref{eq:weighted-matched-domains-WM4}.
The small factor and the age bound are bounds of the extension,
rather than consequences of its holomorphy.

Here are the exact implications once that class has been proved.
The inserted finite-map theorem applied to the potential derivative
gives its regular potential row.  In addition the actual
refile operation must carry a remainder born at $j$ with
norm at most $Cw_{k,j}\mathsf r_j$, and the new birth has
norm $\sigma_{k+1}$.  These primitive refile bounds give
the complete regular row below.
For the amplitude row one must also use the actual channel update
and its normalization:
\begin{equation}
 {\cal T}_{k+1}
       [\dot E_{\zeta,h_0}^{(k+1),{\rm fr}}]
       =[\tau_{k+1}(\zeta)-\tau_k(\zeta)]h_0 .
 \label{eq:weighted-matched-domains-WM4a}
\end{equation}
Here ${\cal T}_{k+1}$ is the own scalar or tensor coefficient
readout in the fixed chosen channel.  Its field-radius
bound is $|{\cal T}_{k+1}F|\le
C\alpha_{k+1}^{-2}\|F\|$.
Apply it to every frozen constant generator; taking
$|h_0|=\|\mathsf k\|_\infty$, or the supremum over its
unit directions for a linear amplitude map, proves the
displayed amplitude-increment bound.  If
$\|\mathsf k\|_\infty=0$ that bound is immediate.
Without this actual update and frozen coverage the
argument bounds only the own coefficient
${\cal T}_{k+1}[\dot E]$, not the separately defined
amplitude increment.
A coefficient readout costs the inverse square of its
physical field radius:
\begin{equation}
 \mathsf e_{k+1}\le C\varkappa_k{\cal L}_k
          +C\sum_{j\ {\rm refiled\ at}\ k}w_{k,j}\mathsf r_j
          +\sigma_{k+1},\qquad
 \mathsf b_{k+1}\le C\alpha_{k+1}^{-2}\mathsf e_{k+1}.
 \label{eq:weighted-matched-domains-WM5}
\end{equation}
If a different normalized coefficient convention is used, its
explicit conversion replaces the last factor.  In Gaussian
coordinates that conversion includes the square of the coupling.
It cannot be charged once in $w_{k,j}$ and a second time as
another old-radius loss:
\begin{equation}
 \alpha_n^{-2}w_{n,j}=\alpha_j^{-2}.
 \label{eq:weighted-matched-domains-WM6}
\end{equation}

Suppose a single zone dial has holomorphic output difference
$F(\lambda)-F(0)$ bounded on \eqref{eq:weighted-matched-domains-WM4} by
$C\epsilon_k q_a^{(k-n-1)_+}$, with $0<q_a<1$.
Cauchy's formula gives
\begin{equation}
 |F'(0)|\le C'\epsilon_k
                   q_a^{(k-n-1)_+}{\cal L}_n .
 \label{eq:weighted-matched-domains-WM7}
\end{equation}
There are $k-n+1$ possible uses of this zone by step $k$.
Summing its actual dial derivatives gives the aged load
\begin{equation}
 {\cal A}_k=\sum_{n\le k}(k-n+1)
                         q_a^{(k-n-1)_+}{\cal L}_n .
 \label{eq:weighted-matched-domains-WM8}
\end{equation}
This derivation uses an already contracted direction at each
use, with its own radius.  It does not apply Cauchy again on
a radius chosen at the original birth.

The own part of a new boundary term can have this age-weighted
bound.  Dependence on an inner stored layer in a live component
is different: its boundary exponent can carry order-one
coefficients counted at each later use.  Its small factor
belongs to the component's mass integral.  That inner part
must be paid under that mass where the component enters.
Placing the same age factor in an exponent norm would claim
a stronger estimate than \eqref{eq:weighted-matched-domains-WM7}.

\subsubsection*{The exact convolution requirements}

Assume the actual forgetting kernel satisfies
$\varsigma_\nu\le C\nu^4L^{-c(\nu-1)}$.
For a fixed $\theta$ with $L^{-c}<\theta<1$ put
\begin{equation}
 S_\theta=\sum_{\nu\ge1}\varsigma_\nu\theta^{-\nu}<\infty .
 \label{eq:weighted-matched-domains-WM9}
\end{equation}
The convergence is geometric after a polynomial factor.
If the shifts agree, the unmarked row
$u_{k+1}\le\varkappa_k\sum_{j\le k}
\varsigma_{k+1-j}u_j$ propagates $u_j\le A\theta^j$
whenever $\sup_k\varkappa_k S_\theta\le1$.
Substitute the induction bound, factor $A\theta^{k+1}$,
and the remaining sum is at most $S_\theta$.
Thus $\theta$ is chosen first and the weak-coupling floor
after it.  The weight itself is not generated by this induction.

For the age kernel define $a_\nu=(\nu+1)q_a^{(\nu-1)_+}$.
On sequences decaying as $\theta^n$ its norm is
\begin{equation}
 A_\theta=\sum_{\nu\ge0}a_\nu\theta^{-\nu}<\infty
                   \quad\Longleftrightarrow\quad \theta>q_a .
 \label{eq:weighted-matched-domains-WM10}
\end{equation}
For $\nu\ge1$ the ratio in the geometric series is
$q_a/\theta$, with the extra factor $\nu+1$.
At equality the terms do not even tend to zero.
Consequently ${\cal L}_n\le B\theta^n$ implies
${\cal A}_k\le BA_\theta\theta^k$ under this strict inequality.
The first finitely many ages are included in the same finite
constant.  They are not discarded by an asymptotic argument.

Squared radius weights often grow polynomially in age under
the admitted frame.  If
$w_{n,j}\le C(n-j+2)^A$ and a refiled remainder has
$\mathsf r_j\le B\theta^j\eta^{n-j}$ for its actual allowed
range, then
\begin{equation}
 \sum_{j\le n}w_{n,j}\mathsf r_j
 \le CB\theta^n\sum_{a\ge0}(a+2)^A(\eta/\theta)^a .
 \label{eq:weighted-matched-domains-WM11}
\end{equation}
This is finite for $\eta<\theta$.  A remainder merely bounded
by $B\theta^j$ has no such conclusion for growing weights.
The recent-window remainders must instead be treated by
their own recursive small factor and finite window length.
These are distinct origin classes.

More generally if a nonnegative sequence satisfies
$e_{n+1}\le\epsilon\sum_{\nu\ge1}K_\nu e_{n+1-\nu}+s_{n+1}$,
with $e_j=0$ at negative indices, the weighted supremum norm
gives
\begin{equation}
 \sup_{n\ge0}\theta^{-n}e_n
 \le {e_0+\sup_{n\ge1}\theta^{-n}s_n
          \over1-\epsilon\sum_{\nu\ge1}K_\nu\theta^{-\nu}}
 \label{eq:weighted-matched-domains-WM12}
\end{equation}
when the denominator is positive.  Prove it first up to a
finite horizon, take the maximum, and move its coefficient
to the left.  The same bound is independent of that horizon.
This is the discrete contraction used after the actual
weighted rows are obtained; it does not establish those rows.

\subsubsection*{A single-zone incidence budget on its original mass}

A zone parameter can occur in several copied potentials,
counterterms and scalar anchors.  To exhibit its true domain,
list these occurrences before estimating a complex exponential.
Let $B_p(Y)\ge0$ be the specified reserve quantities of a
component, including any size reserve, and let $a_{pt}\ge0$
be the charge of occurrence $t$ against reserve $p$.
Suppose the actual exponent increment satisfies
\begin{equation}
 |\Psi(\lambda,Y,e)-\Psi(0,Y,e)|
       \le\sum_pB_p(Y)\sum_t a_{pt}|\lambda_t|
 \label{eq:weighted-matched-domains-WM13}
\end{equation}
on its admitted parameter domain.  All occurrences and their
scalar constants are included.  An oscillation norm alone
does not bound a scalar anchor.

The needed unsourced reserve is the weighted mass
\begin{equation}
 \int^{*} \sup_{e\in{\cal E}_\omega(Y)}|w_\omega^0(Y,e)|
             \exp\left\{\sum_p r_pB_p(Y)\right\}\,dY
                           \le M_\omega .
 \label{eq:weighted-matched-domains-WM14}
\end{equation}
The set ${\cal E}_\omega(Y)$ is the actual admissible
external fibre at the stored variable $Y$.  The upper
integral is the infimum of integrals of measurable
majorants; one may equivalently assume the displayed
common integrable envelope as part of the premise.
The parameter-dependent integrands are measurable in
$Y$ and pointwise holomorphic on their joint source
domain.  Fixed integration sets and this common
envelope justify the dominated Cauchy argument below.
The original inner Wilson factor is part of $w_\omega^0$
with its prescribed original weight.  This assertion does not
replace it by a freely reduced coupling.  The positive
reserve must have been established in the original mass
estimate with its exact allocation.

On the joint domain
\begin{equation}
 \sum_t a_{pt}|\lambda_t|<r_p\quad\hbox{for every }p,
 \label{eq:weighted-matched-domains-WM15}
\end{equation}
equations in \eqref{eq:weighted-matched-domains-WM13}--\eqref{eq:weighted-matched-domains-WM14} bound the perturbed component mass
by $M_\omega$.  The proof multiplies the absolute value
of the original integrand by the exponential majorant
and integrates.  The exterior supremum stays inside the
stored-variable integral.  For fixed integration sets,
this same domination on smaller parameter domains justifies
holomorphy of the integral by the Cauchy formula.

For polydisc radii $R_t$ it suffices to prove
$\sum_ta_{pt}R_t\le(1-\epsilon)r_p$ simultaneously.
Single-axis bounds $a_{pt}R_t\le r_p$ do not imply this
sum-budget condition.  This is the no-reuse incidence test:
every copy is charged once, and all copies sharing a
reserve are summed.

For one dial with load ${\cal L}_n$, an incidence estimate
$\sum_ta_{pt}\le C_p{\cal L}_n$ gives radius
$c/{\cal L}_n$, with $c$ chosen from the unused $r_p/C_p$.
Its radius does not increase as the zone ages.  Age rarity
is a bound on $M_\omega$ or its sum, and yields the factor
in \eqref{eq:weighted-matched-domains-WM7}.  The zone-reserve test section proves why placing
its reciprocal into the input radius fails.  The application
uses \eqref{eq:weighted-matched-domains-WM13}--\eqref{eq:weighted-matched-domains-WM14} for the complete
one-zone list and its assigned mass, including the original
inner Wilson factor and dropout fibres.

For $m$ source directions sharing one potential allowance,
radii proportional to $c/(m{\cal L}_i)$ give a valid
sum-budget disc.  Its mixed derivative costs at most
$M(m/c)^m\prod_i{\cal L}_i$.  Since
$m^m\le e^m m!$, this has a factorial majorant.
It is incorrect to use the two-direction spare-slot radii
for arbitrary $m$ without summing the shared charge.

\subsubsection*{The even-coupling vertex gain}

There is an exact analytic gain for a vertex admitted on
coupling-independent field radii.  Let $q$ be an auxiliary
coupling, distinct from any constraint derivative.
Suppose the actual finite integral and localized formation
are holomorphic for $|q|<R_q$ and $|u|<u_0$.  The
parameter $u$ multiplies the specified vertex.  Assume
the Gaussian and cutoffs are fixed and even, the potentials
are unchanged by $(q,Y)\mapsto(-q,-Y)$, and at $q=0$
the output is independent of $u$.  Write its normed output
as $\Phi(q,u)$ and suppose $\|\Phi\|\le M$.

Cauchy in $u$ bounds
$F(q)=\partial_u\Phi(q,0)$ by $2M/u_0$.
Changing the integration variable $Y\mapsto-Y$ proves
$F(-q)=F(q)$, term by term only when the actual
localization rule respects that change of variables.
Independence at zero gives $F(0)=0$.  Therefore
\begin{equation}
 \|F(g)\|\le {2M\over u_0}|g/R_q|^2,\qquad
 \|F(g)-g^2L\|\le C{M\over u_0}|g/R_q|^4,
 \quad L=\tfrac12F''(0),\quad |g|\le R_q/2 .
 \label{eq:weighted-matched-domains-WM16}
\end{equation}
For a proof expand $F(q)=\sum_{j\ge1}a_{2j}q^{2j}$.
Cauchy bounds $\|a_{2j}\|\le(2M/u_0)R_q^{-2j}$.
Alternatively the even function defines a holomorphic
function of $q^2$ and Schwarz proves the first bound.
The geometric tail starting with $j=2$ proves the second.

The radius $R_q$ is determined by the actual potential
and field cutoff.  If the cut field has size
$b_k\asymp(\log x_k)^{p_0}$, a cubic intrinsic remainder
$q^{-2}V_{\ge3}(qY)$ has removable value zero and
size $O(|q|b_k^3)$.  With the other potential and
core estimates admitting this whole disc, one may take
$R_q=c b_k^{-3}$.  Further proved polynomial logarithmic
losses replace the exponent three by their total exponent.
Then $(g/R_q)^2$ has the form
$x_k^{-1}(\log x_k)^C$.  A mere upper bound on a
possible radius would not prove this gain.
If the admitted disc is instead only
$R_q=c\alpha_k/b_k$ on a running field radius, the
ratio squared is $c^{-2}(\delta_k/\alpha_k)^2$,
where $\delta_k=g b_k$.  Parity alone does not replace
that smaller radius by a coupling-independent one.

The one-loop coefficient has a direct normalized formula.
Suppose the inserted vertex and other potentials have
$S=qS_1+q^2S_2+\cdots$ and $V=qV_1+\cdots$.
Differentiate
$\log{\mathbb E}_0 e^{V+uS}$ in $u$ and expand at zero.
Parity makes ${\mathbb E}_0S_1={\mathbb E}_0V_1=0$.
The coefficient of $q^2$ is consequently
\begin{equation}
 {\mathbb E}_0S_2+
          {\mathbb E}_0(S_1V_1)
                 -{\mathbb E}_0S_1\,{\mathbb E}_0V_1
       ={\mathbb E}_0S_2+
                         \operatorname{Cov}_0(S_1,V_1).
 \label{eq:weighted-matched-domains-WM17}
\end{equation}
The expectation includes the actual cutoff and its
normalizer.  Replacing it by an uncut Gaussian requires
the separate cutoff-tail estimate.  The formula in \eqref{eq:weighted-matched-domains-WM17}
does not compute its tensor eigenvalue or its sign.

\subsubsection*{Matched reference measures on a common domain}

Two local coefficient lists can be compared by
differentiating a common finite integral only after their
variables, sharp sets and numerators have been pulled to
the same reference domain.  Let
\begin{equation}
 Z_t(z)=\int_D e^{V_t(y,z)}\,d\mu_{A_t}(y),\qquad
 d\mu_{A_t}=(2\pi)^{-N/2}\det(A_t)^{1/2}
                     e^{-y^{\mathsf T}A_ty/2}\,dy .
 \label{eq:weighted-matched-domains-WM18}
\end{equation}
Assume the determinant branch, matrices and integrands
are holomorphic on a common neighbourhood of the
interpolation $0\le t\le1$, with fixed $D$ and
uniform domination there.  Assume the generating
function has a common nonzero neighbourhood of $z=0$.
Differentiation gives exactly
\begin{equation}
 {d\over dt}\log Z_t(z)
 =\tfrac12\operatorname{Tr}(A_t^{-1}\dot A_t)
       +\left\langle \dot V_t
                      -\tfrac12y^{\mathsf T}\dot A_ty\right\rangle_{t,z}.
 \label{eq:weighted-matched-domains-WM19}
\end{equation}
This follows by differentiating the determinant and
the quadratic exponent.  A change-of-variables
Jacobian belongs in $V_t$ with its own derivative.
If the cutoff moves, \eqref{eq:weighted-matched-domains-WM19} is not the complete
derivative unless the induced boundary term or an
exact fixed-domain pullback has been supplied.

For the normalized moment generator
$\log Z_t(z)-\log Z_t(0)$, the trace cancels if
$A_t$ has no source dependence.  With
$L_t=\dot V_t-\tfrac12y^{\mathsf T}\dot A_ty$ the
derivative is
\begin{equation}
 \langle L_t(z)\rangle_{t,z}
                           -\langle L_t(0)\rangle_{t,0}.
 \label{eq:weighted-matched-domains-WM20}
\end{equation}
A subtracted vacuum energy has its derivative as
well; it cancels between two presentations only
when it is the genuinely common energy.
If $V_t$ includes a transported observable,
$\dot V_t$ includes that observable's coordinate
derivative.  It is not just the difference of old
interaction coefficients.

Suppose the right side of \eqref{eq:weighted-matched-domains-WM20}, including every
required scalar energy term, has holomorphic supremum
at most $\epsilon_t$ on a common polydisc of source
radii $r_i$.  Integrating in $t$ and using coordinate
Cauchy gives for a mixed cumulant derivative
\begin{equation}
 |\Delta\kappa_{i_1,\ldots,i_m}|
       \le \gamma!\prod_i r_i^{-\gamma_i}
                                  \int_0^1\epsilon_t\,dt ,
 \label{eq:weighted-matched-domains-WM21}
\end{equation}
where $\gamma_i$ is the number of repetitions of
coordinate $i$.  For distinct coordinates there
is no $m!$.  A common source radius with repeated
derivatives has the indicated factorial.
This is a cumulant bound because it differentiates a logarithm.
For raw moments use the actual normalized generator
$M_t(z)=Z_t(z)/Z_t(0)$.  If $|M_t(z)|\le Q$ on the
same polydisc and $G_t=\partial_t\log M_t$ is bounded
by $\epsilon_t$, then
$M_1-M_0=\int_0^1M_tG_t\,dt$.
Coordinate Cauchy gives the raw-moment version of \eqref{eq:weighted-matched-domains-WM21}
with the additional factor $Q$.  Scalar energy
subtractions must be restored or included in the
identified target generator before it is called a
physical moment.
Thus a coefficient comparison gives a matched
moment estimate only after the actual common
domain, scalar normalization, source generators
and expectation bound have been proved.

For adjacent cutoff estimates the resulting
error sequence has to be summable in the
telescoping theorem.  A direct two-cutoff bound
by $\epsilon_{\min(K,K')}$ with $\epsilon_K\to0$
is another sufficient route.  Bounded derivatives
at each cutoff, with no comparison of the
generators across cutoffs, supply neither.
The comparison-capsule section gives the
resulting fixed-torus Cauchy limit in its exact scope.

\paragraph{Transfer between the two bare lattices.}
At a common physical depth transfer a regional pair $(U,J)$ by
blocking $U$ one step and taking the cell mean of the current $J$.
On common rows the composition of the straightened constraints
identifies the retained block field. The local Jacobian and the
block gauge factors agree when formed from this same field by the
same rule. The other quantities to compare are the regional
propagators and kernels, the localized Wilson response, the
localized Gaussian normalizer, and the fields used to evaluate
earlier activities. Define these on the two cutoffs by the same
regional rule before taking differences. A walk term indexed by
a finite set $S$ of cubes is obtained from the alternating sum
of the regional fluctuation operators of $Y\subset S$, each
compressed to the interior of $Y$. The switched precision is
linear in those walk terms. The normalizer uses the logarithm
of an operator built on the recursive tower of each region.
On the full torus the sequences are constant and the two
regional rules use the same uncut fluctuation operator
for their respective totals.

%
%
%
\paragraph{The analytic line on ordinary carriers.}
Take the sharp ordinary data \(U,A',\mathbf J\)
with radii \(a_0,a_1,c_0\), where \(a_1^2\le C a_0\).
Let \(h=\eta/L\) be the sharp spacing in the common
units. Use the line
\[
 A_z=A_{\rm R}+zA_{\rm I},\qquad
 \mathbf U_z=e^{ihA_z}U,\qquad
 \mathbf J_z=J_{\rm R}+zJ_{\rm I},
 \qquad |z|\le R ,
\]
where the real and imaginary parts are in the real
Lie algebra and \(R>1+\sqrt2\) is fixed. The target
is at \(z=i\). The line has a fixed real base \(U\).

For a plaquette \(p\), let \(d_U A'(p)\) be the
sum of its four signed increments transported to
the base of \(p\) in \(U\). The first-order expansion
of \(\partial(e^{ihA'}U)(p)(\partial U(p))^{-1}\)
is \(1+ih d_U A'(p)\), with error \(C h^2a_1^2\).
Both endpoint curvatures are bounded by \(a_0h^2\),
so \(|d_U A'|\le Ch(a_0+a_1^2)\).
The transports in \(U\) preserve the real Lie
algebra. Applying the same expansion to \(A_z\)
therefore bounds the curvature along the line
by \(C_R h^2(a_0+a_1^2)\).

The complete double blocks of a flat gauge cube
and their neighboring translates lie in its
\(4\eta\) halo. Use the one-layer sharp collar
when this halo is contained there. A cube whose
halo fits within the sharp chart size can use
the given gauge on the halo. Otherwise
start from the sharp gauge on the cube, rounding
inward by at most a fixed number of fine layers.
Work on the lattice lift of the cube, using the
covering charts on a short torus.

For the axial continuation add one whole face
layer at a time. In the signed outward direction
\(d\), let \(W(x)\) be the last gauged normal
link from \(x-hd\) to the face point \(x\).
With \(g\) the gauge already defined, set
\[
 g(x+hd)=W(x)^{-1}g(x)U_d(x).
\]
The new gauged normal link is then \(W(x)\).
For a tangential direction \(\mu\), the plaquette
identity for the gauged links \(V\) gives
\[
 V_\mu(x+hd)=W(x)^{-1}
        (\partial U^g)_{d\mu}(x)V_\mu(x)W(x+h e_\mu).
\]
The plaquette error and the tangential change of
\(W\) are both \(O(h^2a_0)\). Commuting \(W\)
past \(V_\mu\) costs \(O(h^2a_0^2)\), since both
links differ from one by \(O(ha_0)\). Thus each
new tangential link differs from its old face
value by \(O(h^2a_0)\). For tangential first
differences compare the two corresponding old
links and add these two errors. Normal links
inherit their tangential differences and agree
on successive links of each extended ray.
The small logarithm divided by \(h\) therefore
preserves the field and first-difference bounds,
with a fixed constant at each added layer.
Extend successive faces of the enlarged rectangle
to include the corners. The halo takes \(O(L)\)
fine layers in each direction, so the total
constant is independent of the cutoff.

In the resulting sharp chart write
\(U=e^{ih a}\). The fixed one-step logarithm map
has bounded first and second derivatives on a
small chart and is the same function on translated
stencils. Applying it to the real base gives a
coarse gauge potential and first difference bounded
by \(C a_0\). For the complex factor use
\[
 F(a,B)=\frac1{i\eta}
     \log\!\left(\mathcal M(e^{ihB}e^{ih a})
                               \mathcal M(e^{ih a})^{-1}\right).
\]
It vanishes when \(B=0\). The derivative in \(B\)
is bounded by \(C_L\), and the mixed derivative
bound gives \(|D_aF|\le C_Lh\sup|B|\).
Comparing the complete neighboring stencils
therefore bounds \(F(a,A_z)\) and its first
covariant difference by \(C_Ra_1\).
The coarse plaquette estimate uses exchanges
of adjacent steps in the averaging paths:
the \(O(L^2)\) resulting sharp plaquette errors
are transported by products of bounded length.
This bounds the two coarse curvatures by
\(C_{L,R}\eta^2(a_0+a_1^2)\).
The finite transport in the cell mean bounds the
independent current by \(C_{L,R}c_0\).

These are the single-scale clauses on any ordinary
domain covered by these charts. In either tower
rule, at an address of depth \(p\) below level \(j\),
put \(\lambda=L^{-(j-p)}\). If \(a'_0,a'_1\)
bound the resulting curvature and complex factor,
the tower theorem gives the bound
\(K_{\rm red}(\lambda a'_0+\lambda^2(a'_1)^2)\)
at that address. In its tower subcall the independent
current is set to zero. The current slot of the
evaluation pair keeps its own radius. Taking the
sharp radii smaller than the paired running
radii by one fixed factor gives, on both lines,
\(a'_i\le b_i\) for the corresponding field
radii, \(a'_0+a'_1\le a^\sharp\), and
\(K_{\rm red}(a'_0+(a'_1)^2)\le b_0\).
Thus the line and its transfer lie in their
cores at every tower address.
On the real diameter the line gives real pairs in
a fixed small radius, with real chart bound
\(C_{L,R}(a_0+a_1)\). This is enough for the real
regional estimates stated at a fixed small radius.
For a shifted or graded read the same argument
requires its own complete collar and bounds.

\paragraph{The regional row form.}
Write $\sharp$ for the finer bare lattice and $\flat$ for the
other lattice. For the switched kernel use scale sequences on
a fixed region $Y$. At physical spacing \(\xi_n\)
the direct rule keeps precisely the cubes of side
\(LM_1\xi_n\) in \(Y\) whose distance from \(Y^c\)
is at least \(\mathsf r_{\rm t}\xi_n\).
The tests defining a member at a common physical scale
therefore agree cube by cube:
\(\Omega^\sharp_{n+1}(Y)=\Omega^\flat_n(Y)\) for \(n\ge1\).
At the lower cut include the full double blocks and
the neighboring stencils of the first rows retained.
%
Use real pairs of small radius whose transfers are also in the
small-radius class. Extend flat row data by zero off the common
interior and outward rows. Let $\mathsf F$ be the fluctuation
form before the row penalty, $Q$ the row map, and set
\[
 \mathsf G=(\mathsf F+Q^{\mathsf T}\xi_c^{-2}Q)^{-1},\qquad
 \mathsf M=Q\mathsf G Q^{\mathsf T},\qquad
 \mathsf H=\mathsf GQ^{\mathsf T}\mathsf M^{-1} .
\]
Here $\xi_c$ denotes the local row scale in the diagonal
penalty. Let $\mathbb P$ be the cell mean and $q_1$ the
straightened linearized bond average for one step. For
$f=\mathsf H^\flat B_1$ and $h=\mathsf H^\sharp B_2^\uparrow$,
write \(\lambda_i=(\mathsf M^{-1}-a_{\rm row})B_i\) for the
row multiplier, with \(a_{\rm row}=\operatorname{diag}(\xi_c^{-2})\)
in that run. Then \(\mathsf F\mathsf H B_i=Q^{\mathsf T}\lambda_i\).
Set \(\mathsf d=Q^\sharp\mathbb P^{\mathsf T}f-B_1^\uparrow\).
Use \(\mathbb P^{\mathsf T}f\) in the sharp equation and
\(q_1h\) in the flat one. The row penalties agree on the
shared rows. If \(Q^\flat q_1h=B_2\), the identity is
\begin{equation}
\begin{aligned}
 &(B_1^\uparrow,(\mathsf M^\sharp)^{-1}B_2^\uparrow)_\sharp
      -(B_1,(\mathsf M^\flat)^{-1}B_2)_\flat \\
 &\quad=\mathfrak D(f,h)-(\mathsf d,\lambda_2),\qquad
 \mathfrak D(f,a)=(\mathbb P^{\mathsf T}f,\mathsf F^\sharp a)_\sharp
                     -(f,\mathsf F^\flat q_1a)_\flat .
\end{aligned}
\label{eq:comparison-regional-row-defect}
\end{equation}
On common rows \(Q^\flat q_1h=B_2\) by composition of
the straightened averages. For an unmatched bottom row use
the complete double block in the sharp bare rim. Its fine
links are rows of the sharp determining set with zero
prescribed values, so its average is zero. If a regional
cut changes this relation, keep the extra term
\((Q^\flat q_1h-B_2,\lambda_1^\flat)\) on the right of
\eqref{eq:comparison-regional-row-defect}. The prolongation
defect pairs with the sharp row multiplier. Its size is
\(O(\eta\nabla f)\) on the common interior, where \(\eta\)
is the coarser of the two bare meshes.

The consistency estimate for $\mathfrak D$ uses the fields cut
off at the unmatched rim. In scaled local first-order norms
$n_1(y),n_2(y)$, the second of which includes the H\"older
quotient of the gradient, its bound is
\[
 C(1+u_0)\sum_y (\eta/\xi_y)^\beta n_1(y)n_2(y),\qquad 0<\beta<1 .
\]
The row currents must satisfy $\xi_c^3|\mathsf j^\sharp(c)|\le1$
and $\xi_c^3|\mathsf j^\sharp(c)-\mathsf j^\flat(c)|
\le u_0(\eta/\xi_c)^\beta$ on the common rows. The term linear
in the prescribed current uses the exact transpose pairing
$(\mathbb P^{\mathsf T}f,J)_\sharp=(f,\mathbb PJ)_\flat$.
In the Landau projection use the straightened site and
field transfers \(q_0,q_1\), for which \(q_1D_U=D_Wq_0\)
on the complete carriers, with \(W=\mathcal M U\).
Write \(\mathbb P^{\rm f}\) for the covariant face mean.
On the full stencils the site Laplacian commutator is
\[
 (\Delta^{\circ\flat}q_0-q_0\Delta^{\circ\sharp})\lambda
 =D_W^{\mathsf T}(q_1-\mathbb P^{\rm f})D_U\lambda
   -(q_0D_U^{\mathsf T}-D_W^{\mathsf T}\mathbb P^{\rm f})D_U\lambda.
\]
This adds and subtracts \(D_W^{\mathsf T}\mathbb P^{\rm f}D_U\lambda\)
after intertwining. The last term is bounded by
\(C\eta^{-1}\varpi\sup |D_U\lambda|\) on its box, with
\(\varpi\) the plaquette defect there. The unit-scale curvature
bound gives \(\varpi\le C\eta^2\). The first term is the
divergence of the small face-mean difference.
After summation by parts the divergence term is
paired with a controlled first difference. The resulting
site resolvent term is controlled in cell $L^2$ and is used
in the final pairing. A second application of a propagator
to that term would need a different bound.

In screening the rim, crossing one layer pays for the growth
of the number of tops by at most $L^4$ and for the growth of
the consistency weight by $L^\beta$. The resulting sum is
\[
 \sum_y (\eta/\xi_y)^\beta e^{-\delta_0d(y,y')}
       \le C(\eta/\xi_{y'})^\beta .
\]
First work with the forms at prescribed row current zero.
For the bond force \(J^\sharp\), put
\[
 a_*=(1-\mathsf H^\sharp Q^\sharp)\mathsf G^\sharp J^\sharp,\qquad
 \mathsf j_*^\sharp=(\mathsf H^\sharp)^{\mathsf T}J^\sharp.
\]
Then \(Q^\sharp a_*=0\), and multiplying out the inverse
equations gives
\(\mathsf F^\sharp a_*=J^\sharp-Q^{\sharp\mathsf T}\mathsf j_*^\sharp\).
The exact transpose pairing for \(J^\flat=\mathbb PJ^\sharp\)
therefore gives
\[
 (B_1,\mathsf j_*^\flat)-(B_1^\uparrow,\mathsf j_*^\sharp)
   =\mathfrak D(f,a_*)+(\mathsf d,\mathsf j_*^\sharp)
                        +(Q^\flat q_1a_*,\lambda_1^\flat).
\]
The last term vanishes on the matched constraint carriers
by the same common-row and bare-rim calculation. The
weak form estimate with the weighted fields \(f,a_*\) and
the screening sum then give the induced current difference.
Insert these currents in
\eqref{eq:comparison-regional-row-defect}. At unit rows the
entrywise rate is $CL^{-\beta k}$ when the step has $k$ finer
levels. The consistency estimate supplies the power of the
mesh and the propagation bound screens the bottom layers.

For the switched kernel use the common regional alternating
sums with switches held at one on the prescribed collar.
The alternating sums have a real smallness bound with a cost
exponential in the number of cubes and an independent analytic
bound with exponential decay. To combine the two, write the
complex pair on the line through its real and imaginary parts
whose real arguments give real pairs. This line and its transfer
must be admitted on a disc of radius $R>1+\sqrt2$, with the
target at $i$. The two-constants estimate uses the bound on
the real diameter and the analytic bound on the disc. Retain
a sufficiently small positive power of the real smallness
so the remaining localization absorbs the growth of the
alternating sums. With the summable spatial decay this gives
a bound $CL^{-\alpha_K k}$ on the differences of the inverse
covariance and the auxiliary Gaussian kernel, for some
$\alpha_K>0$.

\paragraph{The normalization on the regional towers.}
The tower for the determinant is defined by the recursive
regional rule. To compare its operators, let
$\mathbb E^\sharp,\mathbb E^\flat$ be the inverses of the
two first-order systems with bilinear forms $\mathfrak K^\sharp$
and $\mathfrak K^\flat$. For right-hand sides $G^\sharp,G^\flat$,
put $\Phi^\sharp=\mathbb E^{\sharp\mathsf T}G^\sharp$ and
$\Psi^\flat=\mathbb E^\flat G^\flat$. If $\mathfrak i$ prolongs
the flat variables and $\mathfrak a$ averages the sharp variables,
the inverse identities give
\[
 (G^\sharp,\mathfrak i\Psi^\flat)_\sharp
 -(\mathfrak a\Phi^\sharp,G^\flat)_\flat
 =\mathfrak K^\sharp(\Phi^\sharp,\mathfrak i\Psi^\flat)
     -\mathfrak K^\flat(\mathfrak a\Phi^\sharp,\Psi^\flat) .
\]
Use multilinear prolongation $\mathfrak i_0$ and straightened
averaging $\mathfrak a_0$. Compare real pairs of small radius
with transfers of small radius. The flat test field $\Psi$
vanishes below a fixed margin above the first common layer.
Require $\xi_c^3|\mathsf j^\sharp(c)|\le1$ and
$\xi_c^3|\mathsf j^\flat(c)|\le1$ on every row of the two
systems, respectively.

%
For a common bond row the chain rule gives
\(Q^\sharp=Q^\flat q_1\), and for a common top it gives
\(Q'^\sharp=Q'^\flat q_0\). This cancels the terms
with the averaged fine field. For the reversed pairing
with a prolonged coarse field \(B\), the bond-row
error is \(Q^\flat(q_1\mathfrak i_0-1)B\).
On a complete carrier the flat means reproduce
constants. The local comb and the first difference
of \(B\) bound the \(L^2\) norm of the expression
before the row map by
\(C\eta\|\nabla B\|_2+C_L\varpi\|B\|_2\).
The raw bond-row coefficients give
\(|\xi_c Q^\flat(c)F|\le C\xi_c^{-1}\|F\|_2\)
on its full cells. This gives the mesh factor in
the pairing with the raw multiplier. The reversed
top pairing uses the estimate for
\(q_0\mathfrak i_0-1\) in its site scale.

For the top mass use
\(a'(y)=\kappa_n\xi_y^{-2}\), where
\(\kappa_0=\kappa_1=1\) and
\(\kappa_{n+1}^{-1}=1+L^{-2}\kappa_n^{-1}\)
for \(n\ge1\). The insertion of a fine level
changes \(n\) to \(n+1\) on a common top, so
\[
 |a'^\sharp(y)-a'^\flat(y)|
             \le \xi_y^{-2}(\eta/\xi_y)^2 .
\]
The local site row bound makes this a weak error
of the same order in the scaled cell norms.

For the site slot \(p\) the site equation gives
a first-difference bound.
With \(g_0,D^{\mathsf T}x\) the two parts of its
source, test on a complete collar against \(\chi^2p\),
where \(\chi=1\) on the cell and \(|D\chi|\le C/\xi_y\).
Covariant summation by parts and the local bound on
\(Q'\) give
\[
 \xi_y\|Dp\|_{2,y}\le
 C\bigl(\|p\|_2+\xi_y^2\|g_0-Q'^{\mathsf T}r\|_2
                            +\xi_y\|x\|_2\bigr),
\]
with the norms on the right taken on the collar.
The current in \(x\) is used in norm. The row-current
term in \(x\) is \(-\Delta^{(2)}[\mathsf j](A-Dv)\).
Its one-cell bound uses the raw amplitude
\(\xi_c^3|\mathsf j(c)|\), with no difference of
the multiplier. For the
Laplacian pairing with \(p\), apply the difference
of the two gradient means to the other slot.
In flat coordinates the derivative of the
multilinear prolongation is the prolongation of
each directional derivative, constant along its
coarse segment and multilinear transversely.
This map and its volume transpose reproduce
constants on complete incidences. Together with
\(D_W q_0=q_1D_U\) this moves the difference of
means onto the gradient of the other slot, costing
one mesh times its extra \(L^2\) difference.
In the local comb the covariant corrections cost
\(C_L\varpi/\eta\) on site values and \(C_L\varpi\)
on gradient means. Thus \(p\) requires only its
first-difference bound. The site equations for
\(v,w,q\) have \(L^2\) sources controlled by the
cell norms and first differences of the bond field.
In their admitted regular gauge, localizing to a cell and
using the flat identity
\(\sum_{\mu,\nu}\|\partial_\mu\partial_\nu f\|_2^2
=\|\Delta_0 f\|_2^2\) gives their interior
second-difference bounds. The cutoff terms
and the difference from the covariant site
Laplacian cost a bounded first difference and
value in cell units. In the scaled cell norms
of the first-order systems this gives
\[
 \big|\mathfrak K^\sharp[\mathsf j^\sharp](\Phi,\mathfrak i_0\Psi)
   -\mathfrak K^\flat[\mathsf j^\flat](\mathfrak a_0\Phi,\Psi)\big|
 \le C\sum_y\left(\frac{\eta}{\xi_y}+\mathsf u(y)\right)
           [\![\Phi]\!]_{(y)}^\vee [\![\Psi]\!]_{(y)},
\]
where $\mathsf u(y)=\max_{c\text{ at }y}\xi_c^3
|\mathsf j^\sharp(c)-\mathsf j^\flat(c)|$ on the common rows.
The unmatched layers are screened by the cutoff and the
exponential propagation bound. First use the estimate with
zero row currents. The screening sum with exponent one then
gives $\xi_c^3|\mathsf j_*^\sharp(c)-\mathsf j_*^\flat(c)|
\le C a\,\eta/\xi_c$ on common rows, where $a$ is the pair radius
and $\mathsf j_*=\mathsf H^{\mathsf T}J$. Using these induced
currents in the consistency estimate gives the real entrywise
rate $CL^{-k}$ on the common interior bonds. The exponential kernel estimate
on the larger complex domain and the two-constants estimate
give a Schur norm bound $CL^{-k/2}$ on the smaller domain.

%
\paragraph{Common layers of a recursive tower.}
Suppose the two calls use the same physical first
member \(Y\), with aligned grids, and the same
recursive rule. At the next grid of side \(s\),
assume that \(Y\) is a union of its cells. If the
previous member was obtained by removing
\(k\) layers of side \(s/L\) from \(Y\), a
new cell at distance \(ts\) from \(Y^c\) has
clearance \(ts-ks/L\) when it is contained in
that member. The distance index \(t\) is an
integer. This follows by thickening the union
of excluded closed cells in the sup metric,
so the same rule applies at faces and corners.

The rule that requires clearance at least \(s\)
therefore removes \(1+\lceil k/L\rceil\) new
layers. Starting from \(k=0\) its counts are
\(1,2,2,\ldots\). Since the finer call has one
more erosion before the first common scale, the
two members agree at the second common erosion.
This is the rule of the minimal determining
tower with \(M_1\)-cubes.

For the rule with \(\mathfrak M\)-cubes and
strict clearance \(>R_1s/L\), the new count is
\(1+\lfloor (R_1+k)/L\rfloor\). If \(R_1=L^a\),
the count after \(n\) erosions for \(1\le n\le a\)
is \(1+\sum_{b=1}^n L^{a-b}\).
The count at \(n=a\) is a fixed point of the
recursion. Hence the members agree after \(a\)
common erosions for this rule. The full torus
and empty regions agree directly.

These calculations require alignment of \(Y\)
through the indicated first common grids.
Once a member agrees, the same recursive rule
at the same physical grids gives agreement
thereafter. The common determining rows and tops
are obtained from the coincident members above
this fixed bottom margin. A row at level \(n\)
uses its predecessor member for the complete
double block and its successor member for the
row type. In the \(M_1\) tower take \(n\ge3\)
before allowing for the fixed stencil margin.
Use the same rule one level above synchronization
in the strict-clearance tower. Cut below the
complete carriers and include the full stencils
in the unit-error halo before the form comparison.

The fine tower is one erosion ahead at the
first common grid, so its member there is
contained in \(Y\). Monotonicity preserves this
inclusion at each common physical scale.
Consequently the local fine spacing is at most
the local coarse spacing. The one-run kernel
bounds can use the coarser weighted distance,
with a fixed reduction in the exponential rate.
Indeed, away from layer boundaries the length
weights are ordered pointwise. On a boundary,
use the smaller of the two adjacent spacings
as the weight convention, which is equivalent
up to a factor of two to the length counting both
adjacent layers: a boundary segment of length \(t\) contributes
\(t\max(\xi_1^{-1},\xi_2^{-1})\) under the first convention
and \(t(\xi_1^{-1}+\xi_2^{-1})\) under the second.
A coarse cell in the
omitted layers contains only a bounded number
of finer cells, since these layers span a fixed
number of meshes. Thus the same shell count
applies to the one-run terms in the unit-error
halo.

%
%
For the inner window \(\square_0=\square^{\sim5}\),
the raw determining tower uses the strict-clearance
rule. In units of the output level the original
\(\pi_k\)-cells have side \(M/L\). The inequalities
\(M\ge L^4M_1\) and \(M_1\ge R_1\mathfrak M\)
give alignment throughout this tower when the
two calls use the same physical window.
The inner and removed-window reads lie in
\(\square^{\sim2}\), and \(\square^{\sim3}\) lies
in the protected interior of \(\square_0\).
The native pillar bound puts its \(4L^n\)-halo in
the free level-\(n\) pillars in bare units. This
includes the full plaquette and bond-current
stencils for the inner reads. The inverse solves
use the whole tower, with its unmatched lower
data screened as above. If the depth is shorter
than the fixed margin needed for synchronization,
the individual bounds suffice with a larger
constant at those bounded depths.

%
For the inner transfer at old level \(1\le j\le k\),
use \(X^\dagger\), the enlargement of \(X\) by one
layer of \(\pi_{j-1}\), as the read collar.
In level-\(k\) units its width is
\(M L^{-(k-j+1)}\), at most \(M/L\) and at least
\(M L^{-k}\). It contains the one-step stencils
and keeps \(X^\dagger\) strictly inside
\(\square^{\sim3}\) for \(X\subset\square^{\sim2}\).
This places the transfer in the ordinary field
and current bounds on the protected collar.
For the inner input of \cref{htc:inner-substitution},
the same proof gives the old analytic core on
the collar itself. Its complete bare stencils
reach at most two bare steps from \(X^\dagger\)
and remain in \(\square^{\sim3}\), where the flat
representative has the bounds \eqref{htc:inner-readouts}.
For a tower address \((p,Z)\) of \(X^\dagger\),
\(Z\subset\square^{\sim3}\subset\square_{(j)}\),
so the exact change of gauge \eqref{htc:nested-identity} applies to
the restriction of the native response to \(Z\).
Its conjugation norm is less than \(6/5\).
For a truncation at depth \(n'\le p\), the curvature
and current bounds after this change are
\[
 \frac65\frac{a_0}{8}L^{-2j},
 \qquad \frac65\frac{a_0}{16}L^{-3(j-n')}.
\]
These are below
\((a_0/4)L^{-(j-p)}L^{-2p}\) and
\((a_0/4)L^{-(j-p)}L^{2(n'-p)}\), respectively,
since \(n'\le p\le j\). This gives both hereditary
rules for the triple on \(X^\dagger\), with core radii
\((a_0/4,a_1/2,a_0/4)\), for the admitted native
substitution. The pair-orbit evaluation is still
on the activity domain \(X\).
Use the distance with the thinner
read collar and a single sharp-to-flat radius loss
for the one bare transfer. The frame and shifted
inputs in the paired comparison must still be
identified on that collar.

On a common row the plain relative datum has
an exact transfer. First transfer the shifted
sharp field and its sharp base separately.
Their relative logarithm on a flat bond is
\(\mathsf Q_1^\sharp(\mathbf U;\mathsf a)\).
On the common small logarithm branches the
semigroup identity gives
\[
 \mathsf Q_{q+1}^\sharp(\mathbf U;\mathsf a)
  =\mathsf Q_q^\flat\bigl(\mathcal M\mathbf U;
                        \mathsf Q_1^\sharp(\mathbf U;\mathsf a)\bigr).
\]
Compare this input to the actual flat shifted
input in the common frame on the complete
carrier. At fixed base the row estimate
\eqref{htc:inner-row-difference} bounds the datum
difference by \(CL^q\) times the supremum
difference of the dimensionless input logs.
For the conjugation of the input potential,
only the frame value at the initial site
enters this bound. The change of base and the
node factors of the straightened chart must
also be matched. The baseline inner datum
uses the output radius \(a_{0,k+1}\).
Its ratio to the old radius is at most
\(C(n+1)^{1/2}\), without a scale gain for
that baseline.

The Gaussian compression uses the calculation
for its specified tower. For any other response
at a vertex, the same physical subregion and
cutoff must be used at both depths, and the
sequence of its core must be identified before
applying the tower calculation. If the cutoff
\(\zeta\) is applied after the solve, the extra
term in a common averaging row is
\[
 [q_1,\zeta]A(c)=
 \sum_b q_1(c,b)\bigl(\zeta(b_-)-\zeta(c_-)\bigr)A(b).
\]
Here the cutoff is evaluated at the initial
endpoints of the bonds. This term vanishes
on a plateau containing the whole carrier.
For a smooth cutoff it costs at most
\(C\eta\|\nabla\zeta\|_\infty\sup_{\operatorname{supp}q_1(c,\cdot)}|A|\)
for the complete one-step stencil of mesh
\(\eta\), by its bounded row sum. One must
check the cutoff on the actual vertex read,
including any derivative of this commutator
needed in its transfer estimate. A
characteristic cutoff on a straddling carrier
gives only the individual row bound.

For the raw switched response a protected case
is available. The field on an active wall component
is its own restricted solve, with support in
\(\mathbb Y^1\), and its reduced current has support
in \(\mathbb Y^2\). The margins of \(\mathbb Y^3\)
keep the full row, plaquette and first-difference
stencils that meet either support inside that component.
Its indicator therefore adds no commutator on
these native reads, by \cref{dnr:walls}.
In the specified raw sequence, if every
partition cube meeting a cell hull has parameter
one, the field, projected field and reduced
current equal the all-one fields on that
sequence. Indeed, every term of their operator
series uses only these cubes, and the local
structures agree throughout the sequence.
This is the all-one identification in \cref{dnr:walls}.
When it holds on both actual regional sequences, use their
undecoupled comparison with their own determining data.

For other switch assignments, the raw operator
at a vertex keeps the walk terms of the ambient
sequence supported on the active cubes. The
Möbius sum defining a named vertex selects the
undecoupled response of its designated region.
To identify the two, one must first identify
their linear lifts on the specified row data.
At the all-one vertex $QH_1=1$ on the rows.
For the proper mask the linearized constraint is
$QH_1(s)B'$, by \cref{dnr:walls}. The support property
alone therefore supplies no lift identity for a different
designated regional solve.

For a named vertex using the direct region rule
one can test the required top inset on its full read
carrier \(K\). If the required inset has width
\(r_{\rm in}\xi_{\rm top}\), it suffices that
\[
 \operatorname{dist}_\infty(K,Y^c)
 \ge(\mathsf r_{\rm t}+LM_1+r_{\rm in})\xi_{\rm top}.
\]
Every cube of the top grid meeting this neighborhood
of \(K\) then passes the direct clearance condition.
For an outer vertex the region and the reads used
in this check must be those of its own supply.

Write $\mathsf A^-_Y$ for the interior compression of the
tower operator and $T_Y=1-c_N^{-1}\mathsf A^-_Y$, with
$\|T_Y\|\le q<1$ on the matched domain. At a free bond $b$ the
root of the logarithm is
\[
 f_b(Y)=3\log c_N-\sum_{m\ge1}m^{-1}\operatorname{tr}(T_Y^m)_{bb} .
\]
Use the same alternating sum over regions $Y\subset X$ containing
the determining cubes of $b$. Subtract the value of each sum at
the trivial pair, weight by the tents of the base points, and
multiply by $-1/2$. The alternating sum and partition of unity
restore the normalized log determinant. The local pivot terms
agree on the common retained rows because their coefficients
are functions of the same retained field.

The difference of the roots is bounded by the operator norm
difference times $C\sum_{m\ge1}q^{m-1}$, since telescoping $T_Y^m$
produces $m$ terms and the logarithm divides by $m$. The sum over
regions thus costs at most $C2^{|X|}L^{-k/2}$, before localization.
Independently, a chain missing a cube cancels by switching that
cube in the alternating sum. Short chains meeting all cubes
pay the distance of the carrier. Long chains are summed using
$q^m$. The split at a small multiple of the carrier length in
units of the regional mesh retains exponential decay of the
localized sum. Interpolating the two summed bounds with a small
power of the norm difference absorbs the factor $2^{|X|}$ and
retains the required exponential decay together with a positive
power of $L^{-k}$. This also applies to the vacuum subtraction.

\paragraph{The regional Wilson increment.}
For each action carrier $X$ and each vertex subregion $Z$
containing $X$, let $H_Z$ be the response on the core of $Z$
times the prescribed cutoff. The localized interaction uses
the action increment
$\operatorname{Br}_{X,Z}=a_X(e^{i\eta H_Z}U)-a_X(U)$.
For $Y\supset X$ sum these increments over $Z=X\cup S$,
$S\subset Y\setminus X$, with sign $(-1)^{|Y\setminus X|-|S|}$,
and take the third-order Taylor remainder in the fluctuation.
Inversion at the full vertex restores the Wilson increment
with the same subtraction. The one-run bound with the
decoupling parameters in their larger polydisc provides decay
in the number of cubes of $Y\setminus X$.

For a matched real vertex write $U'$ for the perturbed fine
field, $W'$ for the perturbed coarse field, and
$W=\mathcal M U$. Put $D=A^\flat\circ\mathcal M-A^\sharp$.
Subtract the increments in the two actions to get
\begin{equation}
 \begin{aligned}
 &(A^\sharp(U')-A^\sharp(U))-(A^\flat(W')-A^\flat(W))\\
 &\qquad=-\big(D(U')-D(U)\big)
          +\big(A^\flat(\mathcal M U')-A^\flat(W')\big) .
 \end{aligned}
 \label{eq:comparison-regional-wilson-defect}
\end{equation}
The local version includes the plaquette shell of the carrier.
The second term compares configurations on the coarse lattice.
Use the response estimate on the matched real cores and
the continuity of the action there. The first term uses
consistency of the action under a regular perturbation of a
possibly rough background.

For this consistency, the exterior derivative commutes with
the linearized averaging of bond fields. In unscaled plaquette
logarithms the induced average on two-forms is $L^2$ times
a probability average of fine plaquettes. The difference of
the flat quadratic actions is therefore $-L^4/2$ times the
sum of the corresponding flux variances, plus the shell term.
If the perturbation changes the curvature from $F$ to
$F+\gamma$, the change of variance on a block is
$2\operatorname{Cov}(F,\gamma)+\operatorname{Var}(\gamma)$,
where the covariance uses the Lie algebra inner product.
Subtract the mean of $\gamma$ and use Cauchy--Schwarz. The
rough $F$ enters through its size, and the oscillation of
$\gamma$ supplies a positive power of the mesh from the
H\"older bound on the covariant gradient of the perturbation.
In the localized consistency estimate the shell and the
nonlinear errors from multiplication of plaquettes are kept
with this same carrier. This is an estimate of the variation
$D(U')-D(U)$, since the variance of $F$ itself need not be small.

%
The first variation of this quadratic defect in the chosen
cell frame is
\[
       -L^4\sum_P\operatorname{Cov}_P(F,\delta_A F)
                      +d(\text{shell})[A] .
\]
Subtract the block mean of \(\delta_A F\) in the covariance.
The H\"older bound on the covariant gradient of the test
variation then gives the mesh gain while the rough curvature
is used only in norm. Keep the differentiated nonlinear
plaquette errors with the same full carrier.

%
For the nonlinear force itself, keep the transports in
the plaquette chart fixed and separate the explicit
background holonomy. Use \(h\) for the bare spacing in
one run, and write
\(\Pi^0=c_0\mathbf 1+iH_0\) for the background plaquette.
The perturbed plaquette has trace
\(\operatorname{tr}(\Pi^0 e^{i\theta})\), where
\(\theta=d\gamma+n\) and the four \(\gamma\)'s are the
increments transported to the plaquette base. Put
\(\sigma=\max_i|\gamma_i|\) and \(w=|\Pi^0-1|\).
The chart bounds \(n\) by \(C\sigma^2\).
In the differentiated expansion the terms of order zero
and one are the background force and the Hessian.
The remaining explicit dependence on \(\Pi^0\) is in
the two oriented remainders of the plaquette force.
Replacing \((c_0,H_0)\) by \((1,0)\) in those remainders
changes them by \(C\sigma^2w\). Indeed \(H-H_0\)
and \(H'-H_0\) change by at most \(Cw|\theta|\) and
each is multiplied by an \(O(\sigma)\) factor. The other
changed terms have a factor \(n\), a second-order
exponential remainder, or \((\cos\sqrt{\langle
\theta,\theta\rangle}-1)H_0\). Each has the same
\(C\sigma^2w\) bound, since
\(|c_0-1|=w^2/2\) and \(|H_0|\le w\).
The transport of the output to its bond base is
kept in both terms. At the local scale \(\xi\) use
\(w\le a_0(h/\xi)^2\), \(\sigma\le h\sup|A|\), and
the factor \(h^{-3}\) in the bond force. The resulting
scaled force error is at most
\[
          C a_0(h/\xi)\,(\xi\sup|A|)^2 .
\]
Thus the two cutoffs can bound their explicit rough
holonomy remainders separately in the weak pairing.

To compare the remaining force, put
\(N(A)_{\mu\nu}=i[A_\mu,A_\nu]\).
The second-order term of
\(\log(e^{-i\gamma_3}e^{-i\gamma_4}
e^{i\gamma_1}e^{i\gamma_2})/i\) at
\(\gamma_3=\gamma_1,\ \gamma_2=\gamma_4\) is
\(i[\gamma_1,\gamma_4]\). On the full stencil the
opposite increments differ by \(h^2\nabla_U A\).
Consequently
\[
 h^{-2}\theta=D_U A+N(A)
       +O\big((h/\xi)\xi^{-2}(a^2+a^3)\big),
 \qquad a=\xi\sup|A|+\xi^2\sup|\nabla_U A|.
\]
Differentiating the same finite expansion in a
test \(\widehat A\) replaces \(a^2+a^3\) in the
error by \(C(a+a^2)b\), where
\(b=\xi\sup|\widehat A|+\xi^2\sup|\nabla_U\widehat A|\).
For the plaquette with the explicit holonomy removed
the action is \(\frac12\langle\theta,\theta\rangle+
O(|\theta|^4)\). Subtract the derivative of its
quadratic action. Dividing by the cell measure \(h^4\),
the leading nonlinear weak density is
\[
 \langle N(A),D_U\widehat A\rangle+
 \langle D_U A+N(A),N'_A\widehat A\rangle .
\]
The omitted density is bounded by
\(C(h/\xi)\xi^{-4}a^2b\) on a small first-order ball.

%
On complete common cells compare this density at
\((U,A,\widehat A)\) and
\((W,q_1A,q_1\widehat A)\), with \(W=\mathcal M U\).
Now take \(h\) and \(\eta=Lh\) to be the two meshes.
For a coarse \(\mu\nu\)-plaquette the flat per-area
mean \(Q_2^0\) assigns weight \(L^{-6}\) to the \(L^2\)
fine plaquettes of each translate based at a site in
its initial block of \(L^4\) sites. Stokes on each
translated square gives \(D_\eta Q_1^0=Q_2^0D_h\).
Each fine plaquette occurs \(L^2\) times over the
incident coarse plaquettes. The lattice volume
ratio is \(L^4\), so the volume transpose of \(Q_2^0\)
reproduces constants on complete incidences.

Use the transports of a local comb to carry the
fine plaquette values to the coarse base and call
this mean \(Q_2\). In a common comb on the enlarged
stencil the one-step row coefficients and these
transports differ from their flat values by
\(C_L\varpi\), with \(\varpi\) the maximum fine
plaquette defect there. The coarse links in this
frame are also within \(C_L\varpi\) of one.
Taking the difference of the two local operators
therefore gives
\[
 |D_Wq_1 B-Q_2D_U B|
      \le C_L\eta^{-1}\varpi\sup|B|
\]
on that stencil. The estimate compares the row
coefficients at each of the neighboring rows using
their sup bound. It requires no difference estimate
on the plaquette defect.

Use this identity for \(B=A,\widehat A\). In the
first pairing move \(Q_2\) onto the coarse
\(N(q_1A)\) by the volume transpose. In the
pairing with \(D_U A\), move it onto
\(N'_{q_1A}(q_1\widehat A)\). Constant reproduction
and the first-difference bounds compare these
adjoint means with the corresponding fine
products at cost \(C\eta\) times a product
difference and \(C_L\varpi\) times its size.
The purely algebraic pairing uses the same
volume quadrature. Since
\(\varpi\le C(\eta/\xi)^2\), all three errors
are bounded by \(C(\eta/\xi)\xi^{-4}a^2b\)
in density, hence by \(C(\eta/\xi)a^2b\) per
cell. Only first differences of the test are used.
If a common smooth cutoff multiplies the product,
its difference enters in that same quadrature
bound. At a cubical face with an incomplete
incidence strip of width \(O(\eta)\), the density
bound \(C\xi^{-4}a^2b\) and the strip volume
\(O(\eta\xi^3)\) give the same cell estimate.
Keep any wider transition or unmatched row in
its native cutoff estimate. This comparison is for
the physical Wilson remainder on its admitted
common carriers.

%
%
For the scalar weight carried by a sourced step,
let \(c_h,c_\eta\) be the two plaquette weights on
an admitted common carrier. Use \(S_h\) for the
fine Wilson increment density, with its background
subtracted, or its first variation. Include the
lattice volumes in the sums. On complete incidences
one has the exact identity
\[
 \sum_p h^4 c_h(p)S_h(p)
 =\sum_P\eta^4 c_\eta(P)(Q_2^0 S_h)(P)
  +\sum_p h^4
     \big[c_h(p)-((Q_2^0)^*c_\eta)(p)\big]S_h(p).
\]
The transpose is taken for these lattice volumes.
Extend the coarse weight to the common physical
stencil, using its tent formula for the actual
source. Every coarse barycentre contributing to
the transpose at \(p\) is within \(C\eta\) of
the barycentre of \(p\). Constant reproduction
therefore bounds the coefficient in brackets by
the paired weight difference at that barycentre
plus \(C\eta\operatorname{Lip}(c_\eta)\).
Compare the first sum to the coarse action
increment with the same weight \(c_\eta\).
In the variance calculation each block
difference is multiplied by \(c_\eta(P)\).
The oscillation of the curvature increment
still gives the gain against the rough background.
In the curl pairings one includes \(c_\eta\) in
the coarse test product before transposition.
Its first difference adds the same
\(C\eta\operatorname{Lip}(c_\eta)\) error, times
the bound on the corresponding pairing density.

For a plaquette filed to a single cube,
incomplete incidences lie in the strip along
its faces. A plaquette filed to two or four
cubes already lies in that strip, since it
crosses a face. On a fixed admitted carrier its
volume is \(O_M(\eta\xi^3)\). The bound for
the normalized increment or variation density
then gives the shell cost \(O_M(\eta/\xi)\) in
cell units, multiplied by the size of the
weight. The physical variation of the weight
member is the sum of the plaquette variations
with their individual \(c(p)\). It is used
inside this weak pairing. The current argument
of the response remains independent.

In particular let \(U_*\) be an actual fine critical
field on the uncut core, \(W_*=\mathcal M U_*\), and
\(P_*=d\mathcal M\) in compatible tangent coordinates.
Take \(A\) with its complete action and row stencils
where the vertex cutoff equals one and away from all
unmatched row readouts. The physical row
differentials then obey \(K^\sharp A=K^\flat P_*A\).
Write \(J^a\) for the inner-action gradient at
cutoff \(a\), which is the Wilson current on full
stencils. With the fine multiplier \(\omega_*\),
stationarity gives
\[
 \big\langle J^\flat(W_*)-K^{\flat\mathsf T}\omega_*,
                              P_*A\big\rangle_\flat
                =dD_{U_*}[A].
\]
If the carrier reaches an unmatched row or the action
shell, keep its term in the comparison. This identity
tests the physical coarse residual on transferred
variations. The transition shell of a cutoff response
retains its own force. With a general fine field the fine residual
\(\langle J^\sharp-K^{\sharp\mathsf T}\omega,A\rangle_\sharp\)
also enters. The full equation with an independent
current still requires the weak source estimate.

Use the common real-to-complex interpolation disc for the
vertex bounds and then three circles in the fluctuation
amplitude. A sufficiently small positive power of the real
bound retains the required carrier decay after the alternating
sum. On the resulting admitted disc the cubic remainder is
bounded in powers of the unscaled fluctuation
$\mathsf t^{-1/2}B$. Its contribution to the interaction
is multiplied by $\mathsf t$, so the cubic bound is
$C\mathsf t^{-1/2}\|B\|^3$ times the mesh gain and the carrier
decay. On $\mathsf t^{-1/2}\|B\|\le C\delta(\mathsf t)$ this is
$C\delta(\mathsf t)\|B\|^2$ times those same factors.
Compare the two vertex constructions with the same reference
weight $\mathsf t$ here. A difference of actual scalar action
weights is paid on the coupling disc in the step comparison.

\paragraph{Transport of earlier activities.}
Use the same regional rule to construct the response and its
reduced current at the two spacings. The row data are related by
composition of the relative averages. At vertices of the
decoupling parameters the response is the undecoupled response
of the selected subregion, multiplied by its cutoff. Thus the
comparison uses responses with specified determining sets.
Insert the averaged fine response into the coarser response
equation. Its error in the field equation is measured in a weak
pairing with one controlled difference on the test field. The
nonlinear constraint on a common row composes with the actual
one-step transfer. Write \(\eta\) for the coarser bare spacing,
\(h=\eta/L\), and \(\breve q(A)=\eta^{-1}\mathcal Q^1(U^\sharp,hA)\)
for that transfer in the straightened chart. If the comparison
uses its linearization \(q_1 A\), the Taylor bound on the full
carrier \(\mathsf P_c\) gives
\[
 |\breve q(A)(c)-q_1A(c)|
       \le C_L(h^2/\eta)\sup_{\mathsf P_c}|A|^2 .
\]
Thus in the row scale \(\xi_y\) the discrepancy is at most
\(C(\eta/\xi_y)(\xi_y\sup_{\mathsf P_c}|A|)^2\). Use either the
exact transfer or this small value defect in the constraint
comparison.
%
The same row estimate in a variation \(\widehat A\) gives
\[
 |(D\breve q(A)-q_1)\widehat A(c)|
     \le C_L(h^2/\eta)\sup_{\mathsf P_c}|A|
                                   \sup_{\mathsf P_c}|\widehat A|.
\]
Indeed the derivative of the dimensionless one-step
remainder is bounded by \(C_Lh\sup_{\mathsf P_c}|A|\),
and the input and output changes of units give \(h/\eta\).
On common rows the actual constraint derivatives obey
\(K^\sharp(A)=K^\flat(\breve q(A))D\breve q(A)\).
When a weak equation uses \(q_1\widehat A\) in place of
the exact tangent transfer, the constraint-multiplier
error is
\[
 \big\langle\omega,\,
        K^\flat(\breve q(A))(D\breve q(A)-q_1)\widehat A
                         \big\rangle .
\]
This error uses the row estimate on the full carriers
and the actual dual weights of the multiplier. The
Wilson and prescribed-current terms of the field
equation require their own comparisons.
The noncommon rows are in the lowest layers. For the
weighted contraction put \(e(y)=\eta/\xi_y\) above the
fixed bottom cut and \(e(y)=1\) below it. Define
\(\mathsf E(x)=\sup_y e^{-\mu'd(x,y)}e(y)\), where \(d\)
is the regional distance and \(\mu'\) is smaller than
the propagation rate. The triangle inequality gives
\(\mathsf E(y)/\mathsf E(x)\le e^{\mu'd(x,y)}\).
Thus a kernel decaying faster than this weight varies
retains a summable bound in the norm with weight
\(\mathsf E^{-\gamma}\). A layer \(r\) spacings below
the top can enlarge the bare-to-local mesh ratio
by at most \(C L^r\), and each intervening layer
has the fixed distance floor. Choosing the floor
so the exponential decay absorbs \(L^r\) bounds
\(\mathsf E\) at the top by \(CL^{-m}\) in a tower
with \(m\) levels.
The reduced current is estimated from the row multiplier, without
applying two differences to a response known only in a first-order
norm. The comparison of the reduced current is used
well inside the top layer.

%
Here is the passage from the cell estimate to the field norm
used at a vertex. In units of one common cell let $f$ be
the difference of the transferred fine field and the coarse
field, written on the common mesh of spacing $b$. Let $E$
bound its lattice $L^1$ norm on the cell and a fixed collar.
An $L^2$ estimate there gives this input by Cauchy--Schwarz.
Use the first-difference bound and the common $\beta_H$-H\"older
bound for the gradients. Both bounds are on the same chart
and cover the full neighboring stencils. At an interior site
$x$, average $f$ on a cube of side $r\ge 2b$ near $x$.
The number of sites is comparable to $(r/b)^4$. Thus the
value at $x$ has norm at most $C(Er^{-4}+r)$.
For a gradient, average the difference of $f$ between this
cube and its translate by a mesh multiple of size comparable
to $r$ in the gradient direction, divided by that translation
length. Telescoping the difference and using the H\"older
bound on every intervening stencil compares this average to
$\nabla f(x)$ with error $Cr^{\beta_H}$. The two cube sums
are bounded by $CEr^{-4}$, so
\[
 |\nabla f(x)|\le C(Er^{-5}+r^{\beta_H}).
\]
Take $r$ comparable to $\max\{b,E^{1/(5+\beta_H)}\}$ for
the gradient, and to $\max\{b,E^{1/5}\}$ for the value.
When either choice exceeds the collar radius the original
bounded field norm gives the result with a larger constant.
The resulting bounds are $C(E^{1/5}+b)$ and
$C(E^{\beta_H/(5+\beta_H)}+b^{\beta_H})$, respectively.
The same bounds hold at the edge of an aligned cell
when the cell supplies the weak input and all the
one-sided stencils. Choose the averaging cube inward
in every coordinate and translate it in the inward
sign of the derivative direction. Telescoping then
uses only the given Hölder bound in the cell. For a
backward difference the base of the stencil is one
mesh from the forward base, which is covered by the
same error. The small radii can be chosen below a
fixed fraction of the cell side, and larger radii
are covered by the original first-order bound.
Consequently a weak estimate with a power of the local mesh
ratio gives a positive power in the first-order field norm
on the common interior. At the unmatched bottom use its
unit bound in the screened weight. The multiplier comparison
used for the current is separate from this interpolation.

%
To obtain the weak input for that interpolation, write
\(\lambda_y=\eta/\xi_y\) and let \(f^\sharp,f^\flat\)
be the actual bond sources in the direct regional
comparison. In the scaled cell norms the source
condition is
\[
 \xi_y\|f^\flat\|_{L^2(\Delta_y)}\le N(y),\qquad
 |(f^\sharp,\widehat A)_\sharp-(f^\flat,q_1\widehat A)_\flat|
          \le\sum_y\lambda_y^\gamma N(y)\|\widehat A\|_y .
\]
Inject \(f^\sharp,f^\flat\) into the bond slots. Take
\(\Phi^\sharp=\mathbb E^\sharp f^\sharp\) and
\(\Psi^\flat=(\mathbb E^\flat)^{\mathsf T}g\), where
\(g\) tests the wanted flat output in its own slot.
Write the bilinear
form with its operator acting in the second argument.
Then the inverse identity is
\[
 \begin{split}
 (g,\mathfrak a_0\Phi^\sharp-\mathbb E^\flat f^\flat)_\flat
   ={}&\mathfrak K^\flat(\Psi^\flat,\mathfrak a_0\Phi^\sharp)
        -\mathfrak K^\sharp(\mathfrak i_0\Psi^\flat,\Phi^\sharp)\\
      &+(\mathfrak i_0\Psi^\flat,f^\sharp)_\sharp
        -(\Psi^\flat,f^\flat)_\flat .
 \end{split}
\]
%
To use the form estimate in the first line, write \(S\)
for the constant self-transpose map on the first-order
slots given by
\[
 S(A;v,w,\sigma;p,q,r)
           =(A;p+w,q+v,r-\sigma;v,w,\sigma).
\]
For the full system add the multiplier slot
\(\omega\mapsto-\omega\). The reduced system obeys
\(\mathbb K^{\mathsf T}S=S\mathbb K\).
In the full system the extra pairings in
\((S\Psi',\mathbb K\Psi)\) are
\(-(QA',\omega)-(\omega',QA)\), so the same identity
holds. The common transfer and cut use the same site
map on \(v,w,p,q\) and the same top map on \(\sigma,r\).
They commute with this slot change. Thus
\(\mathfrak K(\Psi,\Phi)=\mathfrak K(S\Phi,S^{-1}\Psi)\)
puts the first line in the order of the form estimate,
with the same support at the cut.
The source condition in the second line is tested at
the prolonged bond component \(\mathfrak i_0\Psi^\flat\).
It leaves the quadrature error pairing \(f^\flat\) with
the difference of \(q_1\mathfrak i_0\Psi^\flat\) and
\(\Psi^\flat\) on their bond slots. On complete common
cells the flat averaging and prolongation reproduce
constants, so the difference of the means costs one
mesh times one difference. The common comb frame
adds a term proportional to the plaquette radius.
Use the local \(L^2\) bound on the prolonged test and
its differences in this pairing.
At the fixed-margin cut use the weighted kernel decay
for the lower unmatched contributions. This gives the
cell-mean comparison for the fields and the separate
row comparison for the multipliers under the stated
source condition.

In the middle part of the projected inverse
\(G_1DRD^{\mathsf T}G_1\), use the first-order
system with the multiplier and its equation removed.
Its inverse contains the block \(G_1\). Put
\(u^a=R^aD^{a\mathsf T}G_1^a f^a\) for \(a=\sharp,\flat\).
For the actual force with its weighted supremum bound,
the first-difference Hölder bound for \(G_1 f\) and the
smoothing part of \(R\) give a scaled site Hölder bound
of some exponent \(\beta>0\) for \(u\). On a unit cell
the site mean error \(E\) then gives a supremum error
at most \(C(Er^{-4}+r^\beta)\) by averaging on a
cube of side \(r\) around the test site. Use \(r\) of
order \(E^{1/(4+\beta)}\) above the mesh floor.
This retains a positive mesh power in the site
supremum norm. Split the second solve as
\[
 \begin{split}
 q_1G_1^\sharp D_Uu^\sharp-G_1^\flat D_Wu^\flat
  ={}&q_1G_1^\sharp D_Uu^\sharp-G_1^\flat D_Wq_0u^\sharp\\
     &+G_1^\flat D_W(q_0u^\sharp-u^\flat).
 \end{split}
\]
The last term uses the single-run differentiated-source
kernel on the improved site norm. The first line uses
the weak comparison again, now with the actual
gradient source. On complete common stencils
\(D_Wq_0u^\sharp=q_1D_Uu^\sharp\).
For such sources the weak source defect is the pairing
of \(D_Uu^\sharp\) with
\((1-q_1^{\mathsf T}q_1)\widehat A\).

%
At the flat pair the per-length bond average has row
mass one. In the weighted transpose the mass is also
one, since a fine link of a given direction occurs
in \(L\) paths counted across the incident coarse rows
and the measure ratio is \(L^4\). Both means only read
a fixed number of bare blocks. Thus
\((1-q_1^{\mathsf T}q_1)\widehat A\) costs
\(C\eta\sup|\nabla\widehat A|\).
In a common comb on the full union of incident double
blocks the straightened row coefficients differ from
the flat ones by \(C_L\varpi\), where \(\varpi\) is the
plaquette bound there. This adds at most
\(C_L\varpi\sup|\widehat A|\). Since
\(\varpi\le C(\eta/\xi_y)^2\), the cell pairing
gains a mesh ratio using the scaled \(L^2\) bound of
the gradient source.

%
%
For that bound, use the bond-force entry of the reduced
first-order inverse on the admitted undecoupled regional
frame, including its independent current. For a force \(f_{y'}\) supported
at one source cell, put \(A=G_1f_{y'}\). The cell estimate
for this entry is
\[
 \xi_y^{-1}\|A\|_{L^2(\Delta_y)}
 +\|\nabla A\|_{L^2(\Delta_y)}
 +\xi_y\|\nabla\nabla A\|_{L^2(\Delta_y)}
 \le C e^{-\delta d(y,y')}\xi_{y'}\|f_{y'}\|_{L^2(\Delta_{y'})}.
\]
The stencils in this estimate carry the covariant
differences of that frame. In
\(D u=D D^{\mathsf T}A-D(1-R)D^{\mathsf T}A\),
the first term on full interior stencils is bounded by
the covariant second differences on the neighboring cells,
with the curvature times \(A\) if their order is exchanged.
For the second term, with \(\eta_a\) the bare mesh in
the run, the kernel from \(z\) to \(y\) is bounded by
\(C(\eta_a/\xi_z)^4\xi_y^{-2}e^{-\delta d(y,z)}\).
Cauchy--Schwarz on the source cell and the volume of the
output cell give the \(L^2\) contribution
\(C e^{-\delta d(y,z)}\xi_z^{-2}\|A\|_{L^2(\Delta_z)}\).
After multiplying by \(\xi_y\), absorb \(\xi_y/\xi_z\)
in the exponential by the scale-conversion bound.
Summing the cell estimates yields the same scaled
\(L^2\) force bound for \(D u\) with a reduced decay
rate. Thus the second gradient-source comparison uses
an off-shell one-run estimate on complete common cells.
To return its output to the first-order field norm,
use the differentiated-source Hölder bound as well.
On the complete stencils
\(D_\mu u=-\nabla_\mu^{\mathsf T}X_{\mu\mu}\), where
\(X_{\mu\mu}(x)=\operatorname{Ad}(U_\mu(x))u(x+h e_\mu)\)
on the lattice of mesh \(h\).
The bound for
\(\nabla G_1\nabla^{\mathsf T}\) applies with an output
Hölder exponent strictly below the Hölder exponent
of this transported site input. Together with the
mean comparison this gives the required
first-order interpolation for the second solve.
At the cut the cutoff-gradient and incomplete-row terms must be
kept in the lower-layer estimate. This second solve
uses the site estimate just interpolated. The
terminal \(L^2\) correction for the site Laplacian
commutator is only paired in the first comparison.

%
%
The lift, inverse row map and induced current used by
the nonlinear equation are compared at zero prescribed
row current, with the actual independent bond current
still present. Use \(m(c)=\xi_c\) to convert from
per-length rows to raw rows. If \(H,S\) are the
per-length lift and inverse row map, the raw quantities
are \(H m^{-1}\) and \(m^3S m^{-1}\). The row mass
converts by the same rule as \(S\); on a common row
it agrees at the two cutoffs. Hence the comparison of
\(P\) uses the difference of the converted \(S\)'s.

At a common source row \(c'\), inject \(m^{-1}e_{c'}\)
in the constraint slot of the full first-order system.
Its bond component is the raw lift applied to \(e_{c'}\).
Its multiplier component, multiplied by \(m^3\), is
the raw inverse row map applied to \(e_{c'}\).
For an output row \(c\), use \(m^{-1}e_c\) in the
dual row slot, so the pairing with its multiplier
has exactly the factor \(m(c)^3\). For the output
field, test in one cell at \(y\) by a bond force
\(g\) with \(\xi_y\|g\|_2\le1\). The source-row
terms agree in the two inverse identities because
the common source row and its normalized injection
are the same on the complete carrier. The form
difference then controls the inverse row entry and
the scaled cell \(L^2\) difference of the lift.
This comparison uses both directions of the form
pairing, with the quadrature errors above.

Before comparison the raw row and lift kernels
are bounded by \(C(\xi_y/\xi_{c'})e^{-\delta d(y,y(c'))}\)
in their respective scaled norms. The scale
conversion absorbs the displayed ratio in a
smaller exponential. In the form estimate use the
two localized solutions centered at \(y(c')\) and
\(y\). Summing their product against the mesh
weight gives \(C\mathsf E(y)^\beta e^{-\delta' d(y,y(c'))}\)
for some \(\beta>0\) and a smaller rate
\(\delta'\). This follows by the triangle inequality
for \(d\) and \(\mathsf E(z)/\mathsf E(y)\le
e^{\mu'd(y,z)}\), leaving a summable exponential
in the intermediate cell \(z\).
For the lift apply the cell interpolation to
the difference divided by its one-run exponential
envelope. The one-run gradient Hölder bound
and the scaled cell error give a positive power
of \(\mathsf E(y)\) in the first-order norm with
a remaining exponential in \(d(y,y(c'))\).
Thus, after reducing the exponent and the rate,
\[
 \|(\Delta H)e_{c'}\|_{X_y}
             +|\Delta P(c,c')|
       \le C\mathsf E(y)^\alpha e^{-\delta' d(y,y(c'))},
 \qquad y=y(c)
\]
for the row entry, and for any output cell \(y\)
in the lift term. Here \(X_y\) is the scaled
first-order field norm. A uniformly bounded
number of rows per top and the summable distance
bound give the weighted operator comparisons.
For a row input bounded by \(\mathsf E^\gamma\)
times its norm, the output is bounded by
\(\mathsf E^{\alpha+\gamma}\) times that norm.
Take \(\gamma\mu'\) below the remaining rate.

For the induced current use a bond-force input
\(\mathbf J\). Elimination gives the multiplier
component \(-H^{\mathsf T}\mathbf J\), so its
negative, multiplied by \(m^3\), is the raw
\(\phi\). If the transferred current is
\(\mathbf J^\flat=\mathbb P\mathbf J^\sharp\),
the source discrepancy in the same row test
pairs \(\mathbf J^\sharp\) with the difference
between the prolongation and \(\mathbb P^{\mathsf T}\)
of the bond test. The two means reproduce
constants in the flat frame. On their common
carrier the error costs \(\eta\) times one
difference and the small comb correction.
This gives the mesh gain in \(\phi\) without
differentiating the independent current.

Enlarge the lower cut by the full stencils of
the rows and their neighboring cells. For the
admitted local-box cut the first-order commutator
is supported on this enlarged transition and
has a uniform local weak \(L^2\) bound. The
extra constraint and multiplier terms are
bounded there by their own complete row
coefficients times differences of the cutoff.
Use \(e=1\) in the transition, and screen its
terms by the localized kernels. In particular
a source row \(c'\) in the cut gives only the
one-run exponential there, which is paid since
\(\mathsf E(y)\ge e^{-\mu'd(y,y(c'))}\).
This treats the admitted direct sequences on
common complete upper carriers. A complex
comparison uses a common admitted line and
the two-constants bound with a reserve in the
localization rate.

%
For the nonlinear field equation, compare the force
on the undecoupled solves at \(s=1\) before multiplying
by the vertex cutoff. The form of the force is
\eqref{eq:elab-reduced-force}. For a test \(\widehat A\)
put \(x=dX(T)[\widehat A]\), and use raw row pairings.
Moving the transposes to the test gives the identity
\[
 \begin{split}
 (\mathsf W_s(T),\widehat A)
 ={}&-(x-d\mathcal C_2(T)[\widehat A],\phi)
       -(\widehat Q\widehat A,P_sX)\\
    &-(x,P_s\widehat Q T)+(x,P_sX)
       +(\mathcal V(Y),\widehat A-H_sx).
 \end{split}
\]
The first four terms are row pairings. The last uses
the weak Wilson estimate on the same carriers, with
the adjusted test \(\widehat A-H_sx\).

On the common rows the exact nonlinear composition
identifies \(\mathcal C^\sharp(B)-\mathcal C^\flat(q_1B)\)
with the difference of the flat full row map at
\(\breve q(B)\) and \(q_1B\). The common linear part
cancels. Use the value and derivative estimates for
\(\breve q-q_1\) above, together with the local
Lipschitz estimate for the row derivative on the
admitted ball. This gives the small weak row error
and its derivative, in the full row norms. The
quadratic part has the same estimate by extracting
the degree-two term on a fluctuation disc.

In \(X=\mathcal C(T-H_sX)\), transfer \(T\) by \(q_1\).
On any common row stock where the weighted lift
comparison holds, the difference of the right sides
is the row error, plus the effect of the lift
difference on \(X=O(\|T\|^2)\), plus a term
bounded by \(C\|T\|\|H_s\|\) times the difference of
\(X\). Absorb the last term by the contraction margin.
The same comparison for \(dX\) follows from the
differentiated equation, or from Cauchy's estimate
in the fluctuation at the fixed admitted pair.
In the first pairing retain the cancellation
\(X(T)=\mathcal C_2(T)+O(T^3)\) at each cutoff.
On a smaller complex fluctuation disc the difference
of \(x-d\mathcal C_2(T)[\widehat A]\) then has
quadratic order in \(T\) and the same mesh gain.
The other row terms have quadratic order by
\(X=O(\|T\|^2)\) and \(x=O(\|T\|\|\widehat A\|)\).
These estimates use the row comparison for
\(P_s\) and for \(\phi=H_s^{\mathsf T}\mathbf J\)
at the actual rows. The latter is the induced
current comparison, with the transferred independent
current kept in the pair.

The lift comparison also controls the difference
between the two \(Y\)'s and between the two
adjusted tests in the first-order field norm.
Thus on the admitted direct regional charts the
Wilson pairing and the three row-operator comparisons
give the force defect in the norm of the
one-difference test. The estimate for the vertex
field uses the common smooth cutoff after the solves.
For each native vertex mask we must check its row
and stencil carriers against these estimates, or
apply the corresponding comparison on the actual
recursive tower. Keep the cutoff commutators and
incomplete stencils in the estimate at the lower
margin.

%
For the inner evaluation, anchor both frames at
the same base point. The mixed transport of a
quaternion from \(y\) to \(x\) is
\(P_{xy}v=E(x,y)v\mathbf U(x,y)^{-1}\).
At real data it is an isometry. In the frame
equation write \(D_b=\nabla_b+C_b S_b\) at
spacing \(b\), with \(C_b=(P-1)/b\) on positive
edges and \(S_b\) the shift to their terminal
sites. Thus \(\mathcal K_b=D_b^*D_b\), and the
forcing is \(F_b=(\theta+\Delta_b\theta)w_{0,b}\),
where \(\Delta_b=\nabla_b^*\nabla_b\).

Here is the comparison on a fixed box whose
Dirichlet faces are common nodal planes, with
zero values at their nodes on both meshes.
Use regular gauges for \(E,\mathbf U\) on
the box and all double blocks of the coarse
path and energy reads. Work in these
gauges with \(\|C_b\|_\infty+\|\nabla_b C_b\|_\infty
\le C\varepsilon_{\rm g}\). Let \(a_b\) be the
unit anchor and take the cutoff to vanish near
the boundary. The comb has bounded physical
length, so \(\|w_{0,b}-a_b\|_\infty
\le C\varepsilon_{\rm g}\). For the one-run
normalization use
\[
 (\mathcal K_b+1)(w_b-\theta a_b)
   =(\theta+\Delta_b\theta)(w_{0,b}-a_b)
                    -(\mathcal K_b-\Delta_b)(\theta a_b).
\]
Pairing forward and backward edges bounds the
last term by \(C\varepsilon_{\rm g}\).
In fact the sum of the two transport increments
in a coordinate direction is a difference of
the forward coefficients, times \(b\), plus
\(O(b^2\varepsilon_{\rm g}^2)\). The remaining
terms contain one difference of \(\theta\).
The resolvent has supremum norm at most one
by the equation at a site of maximal quaternion
norm. It follows that \(|w_b|\) is bounded below
on each fixed positive plateau of \(\theta\)
when \(\varepsilon_{\rm g}\) is small.

%
To compare the forcing, use the same comb order
on the box lift. Put
\(\epsilon_E=\eta^{-1}\sup|E^\flat-\mathcal M E^\sharp|\)
in the aligned gauges. Use the restriction
of the fine gauge for the coarse \(\mathbf U\),
so \(\mathbf U^\flat=\mathcal M\mathbf U^\sharp\)
in these frames.
Each coarse transport differs from its straight
fine product by \(O_L(h^2)+O(\eta\epsilon_E)\),
where \(h=\eta/L\), by the one-step path
comparison. Telescope the products from the
base to a coarse site. For a nearby fine site,
append the short cell path and put it in the
native comb order by interchanging neighboring
steps. There are at most \(C_L/h\) such exchanges,
each with curvature error \(O(h^2)\).
The path terms therefore differ by
\(C(\eta+\epsilon_E+|a_h-a_\eta|)\) on paired
cells. Away from the cells at the breaks of the
piecewise polynomial \(\theta\), the two second
differences of the cutoff agree up to \(C\eta\).
They are bounded in the break cells, of total
volume \(O(\eta)\). The forcing discrepancy
thus costs \(C(\eta^{1/2}+\epsilon_E+|a_h-a_\eta|)\)
in physical \(L^2\).

For the mixed energy let \(R\) denote restriction
to coarse sites and \(I\) multilinear prolongation
of a coarse test \(g\). The two forms are
\(\mathcal A_b(v,g)=(v,g)_b+(D_bv,D_bg)_b\)
in the physical-volume pairings. The
one-step path comparison and the first differences
in the good gauges bound the discrepancy of
the per-length coefficients on corresponding
stencils by \(C(\eta+\epsilon_E)\).
The flat gradient of \(Ig\) is an interpolation
of the parallel coarse gradients. Its volume
transpose preserves constants on each complete
incident stencil. In the flat gradient terms
replace this mean of \(\nabla_h w_h\) by the
mean along the straight segment giving
\(\nabla_\eta Rw_h\). Their difference costs
\(C\eta\|\nabla_h^2 w_h\|_2\) by telescoping
over the fixed number of fine edges in the
stencil. In terms with \(C_b S_b\), move
coefficient factors onto \(w_h\) or \(\nabla_h w_h\)
before this quadrature. Differences then cost
one difference of the coefficient and one of
the field, or the coefficient discrepancy above.
The mass terms use the constant-preserving
site transpose. On a common grid-aligned
Dirichlet box, boundary sites are zero and
the coarse edges paired with a nonzero test
gradient have their full transverse incidence.
This yields
\[
 \big|\mathcal A_\eta(Rw_h,g)-\mathcal A_h(w_h,Ig)\big|
 \le C(\eta+\epsilon_E)\|w_h\|_{H^2_h}\|g\|_{H^1_\eta}.
\]
The one-run \(H^2_h\) bound follows by writing
\((\Delta_h+1)w_h=F_h-(\mathcal K_h-\Delta_h)w_h\).
The last term is bounded in \(L^2\) by
\(C\varepsilon_{\rm g}(\|w_h\|_2+\|\nabla_h w_h\|_2)\).
Energy controls these first-order norms.
Odd reflection across the Dirichlet faces and
commutation of the flat differences bound the
second differences by the flat Laplacian norm.
Pair the frame equations with
\(g=Rw_h-w_\eta\), using \(Ig\) on the fine
lattice. The source bound and the same
site quadrature give
\[
 \|Rw_h-w_\eta\|_{H^1_\eta}
       \le C\bigl(\eta^{1/2}+\epsilon_E+|a_h-a_\eta|\bigr).
\]

For the Dirichlet realization one can also
remove a thin boundary strip before comparing
coefficients. Taking the norm in the mixed
equation gives
\((\Delta_b+1)|w_b|\le |F_b|\).
In a coordinate of width \(D\), a multiple
of \(s(D-s)\), with \(s\) the distance to one
face in that coordinate, dominates the bounded
forcing and is nonnegative on the boundary.
Comparison in each coordinate bounds \(|w_b|\)
by \(C\) times the distance to the boundary.
Setting zero boundary data on a common inset
at distance at most \(r+\eta\) changes the
solution there by at most \(C(r+\eta)\) in
supremum norm, by the homogeneous mixed
equation. The coefficient and forcing
comparison can then be made on that inset.

On an interior read the first differences are uniformly
bounded by the local bootstrap
\eqref{psr:scalar-point-bootstrap}. Its proof applies componentwise
to these orthogonal transports: the flat principal part is the
same, and in the regular gauge the first-order coefficient is
\(O(\varepsilon_{\rm g})\), while the zeroth coefficient and the
forcing have bounded finite \(L^p\) norms. The energy estimate
supplies its outer \(H^1\) hypothesis. All cutoffs fit in the fixed
read collar.
A maximum of size \(s\) for the difference
therefore gives an \(L^2\) contribution at
least \(cs^6\) in dimension four. Below mesh
radius use the single site with weight
\(\eta^4\) and \(s\lesssim\eta\).
The resulting supremum error has the power
\(1/3\) of the \(L^2\) error. The plateau
lower bound permits normalization in
\(\operatorname{SU}(2)\) with the same error. This supplies the
frame value for conjugating the inner datum on
the common carrier.

%
%
The native box has faces between sites and is a
union of complete paired cells. For the realization
with only inner edges, use its inner-edge scalar
Laplacian for \(\Delta_b\theta\) and take \(\theta\)
constant in the normal directions near the faces.
In particular a cutoff vanishing in a boundary
collar has the same forcing comparison as above.
Use the mixed cell mean \(P_0\) from the fine sites
to a coarse centre. Its transports are those of
the two averaged contours. Let \(Q_P\) average
fine bond fields along the straight \(L\)-step
paths from one cell to the neighboring cell,
transported to the initial centre. The weight
of each fine edge occurrence in the path sum is \(L^{-5}\).
Telescoping the path differences gives
\[
 D_\eta P_0=Q_P D_h+\mathcal R,\qquad
 \|\mathcal R\|_{2\to2}\le C(\eta+\epsilon_E).
\]
Indeed the mismatch of the two end transports
is \(O_L(h^2)+O(\eta\epsilon_E)\) on each path,
and the path difference is divided by \(\eta\).

For coarse \(v\), let \(I_Pv\) be the multilinear
prolongation, transporting the coarse corner values
by short paths to the centre of the cell containing
the fine site and from there to that site. The mass defect
\(I_Pv-P_0^*v\) costs \(C\eta\|D_\eta v\|_2\).
In a local comb gauge on a complete stencil the
gradients of \(I_Pv\) and \(Q_P^*D_\eta v\)
average the neighboring parallel slopes, with
total weight one in each case. In the latter
average this follows by counting the paths through
the fine bond in its own direction. Their difference
is bounded by \(C\eta\) times the second differences
of \(v\). Restoring the transports costs
\(C(\eta+\epsilon_E)\|v\|_2\), by the complete
stencil curvature bound and the same link
comparison. At a cell face reflect the field
and the tangential links, with identity transport
on the link to the mirror site. The straddling
plaquettes are then flat and the outward coarse
slope is zero. The reflected stencil therefore
gives the same estimate at the inner boundary.

For the second-difference bound put \(v=w_\eta\) and
\(\psi=D_\eta v\). If \(T_\mu\) denotes an orthogonal
transport followed by one translation, the difference
\(T_\mu T_\nu-T_\nu T_\mu\) is a two-step transport
times the plaquette holonomy minus the identity.
Consequently \([D_\mu,D_\nu]\), and the corresponding
forward-backward commutator, have bounded \(L^2\) operator
norm under the normalized curvature bound. Commuting one
difference through \(D^\sharp D\) gives, with the shifted
carriers retained,
\[
 D_\mu D_\nu^\sharp D_\nu v
 =D_\nu^\sharp D_\nu D_\mu v
  +[D_\mu,D_\nu^\sharp]D_\nu v
  +D_\nu^\sharp[D_\mu,D_\nu]v.
\]
Thus the equation for \(\psi\) has source \(D_\eta F_\eta\),
a divergence of a field with norm at most \(C\|v\|_2\),
and a term of norm at most \(C\|\psi\|_2\).
Pair with \(\psi\). Summing the source pairing over components
gives
\((D_\eta F_\eta,D_\eta v)_\eta
  =(F_\eta,F_\eta-v)_\eta\).
Move the divergence onto the differences of \(\psi\) and use
Young's inequality to absorb half their square norm. The site
energy estimate then yields
\(\|D_\eta\psi\|_2\le C\|F_\eta\|_2\).
For the inner-edge boundary convention, perform the calculation
after the preceding even reflection of the site field and
links. The added normal difference is zero and the straddling
plaquette is flat, so restriction gives the same estimate on
the original box without a boundary term.

The exact form difference is
\[
\begin{split}
 \mathcal A_h(I_Pv,f)-\mathcal A_\eta(v,P_0f)
  ={}&(I_Pv-P_0^*v,f)_h\\
    &+(D_h I_Pv-Q_P^*D_\eta v,D_h f)_h
       -(D_\eta v,\mathcal R f)_\eta .
\end{split}
\]
With \(f=w_h-I_Pv\), the forcing difference
and these defects give the same \(L^2\) error
as in the nodal-plane case. Use the regular
gauges on the read collar for the first-difference
bound and the supremum interpolation. In those
gauges the covariant prolongation differs from
the ordinary one by \(O(\eta)\) there.

%
%
For the native real box, use the common good
gauge of the input \(\mathbf U\), normalized
at \(y_0\) and restricted from fine to coarse.
Here \(a_0\) bounds the actual real potential
in that gauge, including any real increment.
The box and its one-step carrier fit the
ordinary gauge cube by the size hierarchy.
The actual box datum in this gauge is formed
by the full iterated averages of the gauged
input. It is gauge equivalent to the native
comb datum by a small node transformation.
Indeed the difference of the two gauges,
normalized at \(y_0\), is the inverse comb
transport in the good gauge, hence close to
one by the potential bound. The small normalized-gauge
and real-base change argument in the proof of
\cref{htc:inner-substitution}, on these complete box reads
and the common smaller logarithm ball, identifies the
two representatives of the box orbit.

The ordinary datum in the common good gauge
agrees on common rows by composition and
covariance of the averages. On a row whose
full preceding outward read is also common,
the reduced data agree by \cref{dnr:physical} and
\eqref{mb:reduced-datum}. At the trivial base,
\cref{dnr:flat-identification} identifies the Landau
potential of the box minimizer with
$\mathbf H(\mathbf 1;B_\natural)$ on the closed
noninward row stock, for its actual ordinary
logarithmic datum and reduced datum, on the
fixed small ball of that result. On a row of spacing $\xi_c$, the
averaging bound at the trivial base gives
\(|B(c)|\le C\xi_c a_0\) in the good gauge.
The same profile for \(B_\natural\) follows
from the local change of data, since its
row lengths are comparable. Integrating the
screened derivative of \cref{dnr:flat-identification}
from zero bounds the potential by \(Ca_0\)
and its first difference at a cell \(y\)
by \(Ca_0/\xi_y\). In the sum one absorbs the
ratio of the source and output row lengths
in the exponential kernel.
Thus on the protected upper reads the
regional response estimate compares the
two Landau potentials at these data.

The actual coarse average of the fine
minimizer uses the one-step node factors
as well as the straightened potential.
In the trivial-base chart the one-step
factor differs from one by \(C_Lh\sup|A|\).
Comparing the same holomorphic map on two
translated blocks bounds its change by
\(C_Lh\eta\sup|\nabla_h A|\).
In the normalized coarse link the two
endpoint factors therefore cost at most
\(C_Lh(\sup|\nabla_h A|+\sup|A|^2)\).
Here the quadratic term comes from
commuting the near-identity factor past
the coarse link. Include also the
nonlinear straightened transfer error.
This supplies the actual \(E\)-coefficient
estimate on the matched protected reads.

The anchor in a Landau representative
uses its gauge at $y_0$. By \cref{dnr:physical},
the inverse of that gauge factor there is $v_B(y_0)$.
The native comb node factor is one at
this site. For the accumulated factor,
separate the fine one-step factor from
the recursion. The remaining product
is the flat factor at the exactly
straightened fine input. The one-step
factor is \(1+O(h\sup|A|)\). On the
complete common top pillar the flat
factor has a bounded holomorphic chart
by \cref{dnr:flat-identification}, and Cauchy's
inequality bounds its change by the
supremum input error there. Hence the
protected response comparison also
compares the Landau anchor factors,
when the original anchors coincide in
the common input frame.

%
%
For the inner-edge realization the matched
interior comparison suffices, with a shrinking
boundary strip. In fixed box units let
\(C\eta<r<1\). The accumulated erosion of
the native tower through scale \(\xi\) is
at most \(C\xi\), by summing the successive
mesh widths. Thus stencils at distance at
least \(r\) from the faces have local scale
\(\xi_y\ge cr\). The fixed unmatched lower
layers lie in the strip. By \eqref{px:level-barrier},
the ratio of two row spacings is at most a fixed
factor times \(\exp((\log L)d/\rho)\), where \(d\)
is the regional distance.
Taking a fixed power \(\sigma\in(0,1]\)
small enough for the screening rate gives
\(\mathsf E(y)\le C(\eta/\xi_y)^\sigma\).
On the complete
common rows outside the strip the response
bound and the one-step node correction
give the per-length box coefficient error
\[
 \epsilon_E(r)
   \le C\bigl(\eta^\alpha r^{-1-\alpha}+\eta/r\bigr),
 \qquad \alpha>0.
\]
The one-run mixed coefficient is bounded
throughout the box by the potential
bound, and has first difference at most
\(C/r\) outside the strip.

For this application one need not bound
the normalized curvature throughout the
strip. Write the mixed equation as
\[
\begin{split}
 (\Delta_b+1)w_b={}&F_b
    -\nabla_b^*(C_bS_bw_b)\\
    &-S_b^*C_b^*\nabla_bw_b-S_b^*C_b^*C_bS_bw_b .
\end{split}
\]
Reflect the scalar site functions evenly
at each cell face. In the divergence term
reflect an edge component oddly across
the faces normal to its direction, with
zero flux across those faces, and evenly
across the transverse faces. Then the
extended flat equation uses the usual
lattice differences. Integrating the heat
kernel bound \eqref{eq:discrete-holder-estimates-DH5}
with the mass factor \(e^{-t}\) gives a uniform
\(L^\infty\) bound for the first derivative
of the flat resolvent. Its derivative
followed by a divergence costs at most
\(C(1+|\log b|)\), since the second derivative
of the heat kernel has summed bound
\(C\max(t,b^2)^{-1}\). The mixed maximum
principle bounds \(w_b\). The term linear
in \(\nabla_bw_b\) is absorbed by smallness
of the coefficient in the common gauge.
Hence
\[
 \|\nabla_bw_b\|_\infty\le C(1+|\log b|).
\]
The same logarithmic bound holds after
adding the bounded transport term.

Apply the nodal restriction and multilinear
test prolongation to the weak forms in
these gauges, using even extension for the
scalar quadrature on the doubled box.
Split the complete interpolation stencils
by whether they meet the strip, including
the \(O(\eta)\) stencil halo. The forms
on the strip of width \(Cr\) cost at most
\(Cr^{1/2}(1+|\log\eta|)\) times the test
energy norm, by the volume of the strip
and the preceding bounds. In the interior
the flat equation with a cutoff changing
on length \(r\) gives an \(L^2\) bound
\(Cr^{-2}\) on the second differences.
The interior quadrature therefore costs
\(C(\eta r^{-2}+\epsilon_E(r))\).
For the forcing the coordinate combs from
\(y_0\) to the points of the smaller
rectangle stay within that rectangle.
Use the per-length comparison on their
full carriers. On the rim use the uniform
bound on \(F_b\), which requires the
stated face property of the cutoff.
This yields the energy error
\[
 C\bigl(r^{1/2}(1+|\log\eta|)
       +\eta r^{-2}+\epsilon_E(r)+\eta^{1/2}
       +|a_h-a_\eta|\bigr).
\]
Take \(r=\eta^\gamma\) with
\(0<\gamma<\min\{1/2,\alpha/(1+\alpha)\}\).
The mesh terms are \(O(\eta^\nu)\) for
any \(\nu<\min\{\gamma/2,1-2\gamma,
\alpha-(1+\alpha)\gamma\}\). On the read
collar the one-run interior Lipschitz
estimate and plateau normalization give
the frame comparison with the cube root
of the energy error. The
plateau lower bound here also follows
by comparing the forms for \(w_b\) and
\(\theta a_b\) in energy, which uses only
the small supremum of the mixed coefficient;
interpolate the result on the interior
collar. To apply this argument to the
named smoothing rule use its actual
edge stock, cutoff at the faces and
anchored real input.

To compare regular activities use their own counterterms before
estimating the old evaluations. Write $G^\sharp,G^\flat$ for
partnered families, $e^\sharp,e^\flat$ for their scalar
coefficients, and $A^\sharp,A^\flat$ for the actions. With
$\widetilde G=G-eA$ in each run the exact decomposition is
\begin{equation}
\begin{aligned}
 \widetilde G^\sharp
 &=\big(G^\sharp-\mathbb M^*G^\flat
           -(e^\sharp-e^\flat)A^\sharp\big)\\
 &\quad+e^\flat(\mathbb M^*A^\flat-A^\sharp)
              +\mathbb M^*\widetilde G^\flat .
\end{aligned}
\label{eq:comparison-transport-counterterm}
\end{equation}
The second-order test on a linear potential of constant
curvature extracts $e$. The linearized averaging reproduces
that potential in the coarser units. Thus the first bracket
obeys the same cancellation as a regular counterterm-subtracted
family. The action defect in the second term obeys the same
quadratic cancellation. In a fluctuation insertion its value
at the background is removed by the fluctuation derivative.
Its variation is bounded by the regular perturbation estimate
for the plaquette defect. The two powers of the radius in that
estimate pay for extracting $e^\flat$ by Cauchy's inequality.

In the last term compare the two evaluations of the same
counterterm-subtracted function. Integrating one derivative in
each of $N$ decoupling parameters gives an alternating sum of
its $2^N$ vertex values. Use the regional response and
frame comparison on the reduced polydiscs where the
full vertex carrier admits it.
%
For an outer evaluation, use the domain field at \(B(t)\)
plus \(w\) times its window direction. Write
\(\rho_j=a_{0,k+1}/a_{0,j}\) and \(\sigma_j=L^{-(k-j)}\)
in the partnered units. The one-run field bound on
the full \(X^\sim\) gives a direction of size at most
\(C\sigma_j\delta(\mathsf t)\le C\mathsf s\rho_j\sigma_j a_{0,j}\).
The field at \(B(t)\) has size at most
\(C\mathsf s a_{0,j}\), since \(\sigma_j\rho_j\le C\).
After allowing for the fixed transfer loss in the
background radius, the disc
\(|w|\le c/(\mathsf s\rho_j\sigma_j)\) therefore fits
the old radius on each cutoff. Choose \(c\) smaller
than the available radius fraction divided by the
one-run bound. The same disc gives the Cauchy cost
\(C\mathsf s\rho_j\sigma_j\) for the derivative.

%
%
For the raw switched outer field, restrict to
the wall component supplying its read of \(X\).
If all determining rows of that component have
complete averaging carriers disjoint from the
window direction, its row datum is independent
of the window parameter. Indeed the parameter
enters the literal averages only through
\((t+t_\square\zeta_\square)H_k\) on those
carriers. The potential and assigned current
of the restricted response therefore both
have zero window derivative. When this holds
at both cutoffs the two activity derivatives
vanish at this vertex, even when their
baseline fields differ.

For the named supply the vertex reduction is
instead given by its alternating sum. Write
\(F_a\) for the cutoff times the undecoupled
field and current of the region \(Y_a\), using
that region's datum. At a vertex with active
set \(v\), only subsets \(S\subset v\) survive
the switch monomials. The coefficient of
\(F_a\) is
\(\sum_{a\subset S\subset v}(-1)^{|S\setminus a|}\).
It is zero for \(a\subsetneq v\) and one for
\(a=v\). Thus the exact datum for the nonzero
comparison is the one on \(Y_v\)'s determining
set. If the complete read lies inside the
cutoff plateau of a common region \(Y_v\),
and its own sequence satisfies the direct
top-inset criterion above, the full-switch
response estimate applies on those reads
even when \(v\) is a proper subset of the
ambient cubes. The assignment of the
region's own core and datum is needed before
using that criterion. In particular the raw
masked lift applied to a datum supported on
a missing row can vanish on the read while
the named all-one lift has that row as its
linear constraint.

The outer and removed tails can be bounded before
comparing the vertex fields. For an outer \(X\)
away from the window, the
window derivative has the distance gain of the
one-run estimate. For \(X\) meeting
\(\square^\sim\) and leaving \(\square^{\sim2}\),
\(d_j(X)\ge \sigma_j^{-1}-C\), so use part of the
old-domain decay instead. Each of the two removed
directions already has this gain on its full carrier
in \eqref{eq:local-formation-WF11}. After the window
count, the factor \(e^{-c/\sigma_j}\) in these
estimates leaves the localization reserve.
With age \(n=k+1-j\),
\(\rho_j\le C(n+1)^{1/2}\) and
\(\sigma_j^{-1}=L^{n-1}\). Thus the two one-run
bounds themselves give a factor \(L^{-\alpha j}\)
together with the fixed exponential age weight
when \(L^{n-1}\ge A(j+1)\), with \(A\) fixed
sufficiently large. Use the actual outer vertex
and removed segment comparisons at the remaining
ages.

%
%
The first removed segment is affine in the two
shift coordinates of its common response chart.
On a matched complete carrier write it as
\(A_t=A_0+t\Delta A,\ J_t=J_0+t\Delta J\),
where the current slot includes the fixed base
current. To compare the segment at the two
cutoffs, put the two endpoints in the same
frames before interpolating. Let \(h\) be the
fine spacing in these units and let \(T\) be
the relative potential transfer. Its second
derivative in the potential slots is \(O(h)\):
the holomorphic map on the dimensionless
stencil is evaluated at \(hA_t\) and its
logarithm is divided by \(Lh\).
Consequently the difference between \(T(A_t)\)
and the affine interpolation of the transferred
endpoints is bounded by
\[
 C h\,t(1-t)\|\Delta A\|^2 .
\]
For the neighboring-stencil estimate take local
units in which the reference scale is one.
Compare the stencils by parallel transport.
The base link arguments of the stencil differ
by \(O(h^2\kappa)\), where \(\kappa\) is the
bound on the base curvature on the intervening
plaquette strips. At an interpolated potential
the difference of the angular arguments is
bounded by \(Ch^2\sup_t\|\nabla_U A_t\|\).
The second derivatives differ by \(Ch\) times
the sum of these two argument differences.
The two direction slots also vary, each by
at most \(Ch\|\nabla_U\Delta A\|\).
After division by the mesh, the first difference
of the affine defect is bounded by \(Ct(1-t)\)
times
\[
 h\|\Delta A\|\|\nabla_U\Delta A\|
 +h^2\|\Delta A\|^2
              \bigl(\kappa+\sup_t\|\nabla_U A_t\|\bigr).
\]
All these norms use the complete neighboring
stencils in the same local units. For the current transfer
\(\mathbb P_{\mathbf U_t}J_t\), a derivative of
each path transport costs \(C h\|\Delta A\|\).
Its second derivative and the product rule
give the affine defect bounded by
\[
 C t(1-t)
 \bigl(h\|\Delta A\|\|\Delta J\|
          +h^2\|\Delta A\|^2\sup_t\|J_t\|\bigr).
\]
Together with the transferred endpoint estimates,
these defects give the comparison on this first
segment wherever the paths are admitted in the
common radius. The current used here is the
assigned affine current.

%
%
For the second removed segment use the old box
datum of \cite[(3.22)--(3.25)]{B87}. Write \(a\)
for the upper fluctuation in angular units,
\(u\) for the window parameter, and \(\chi\)
for the cutoff equal to one on
\(\square^{\sim3}\) and zero outside
\(\square^{\sim4}\). Denote the relative log
of the full average at a row \(c\) by \(Q_c\).
In a common box frame the two relative data are
\[
 D_{\rm mid}(c)=\chi(c)Q_c((t+u\zeta_\square)a),
 \qquad D_{\rm in}(c)=Q_c((t\chi+u\zeta_\square)a).
\]
Set \(R=D_{\rm mid}-D_{\rm in}\).
The row path runs from the inner term to the
preceding term with \(D_s=D_{\rm in}+sR\).
For this comparison use complete old-row
carriers, with every carrier reading
\(\zeta_\square\) in the plateau of \(\chi\)
and with \(\chi(c)=1\) on its row.
On that carrier \(R(c)=0\). On the other
carriers both data are independent of \(u\).
Thus \(R\) is independent of \(u\), and its
support is in the transition enlarged by the
row reads. At \(u=0\) the fixed-base estimate
for the full row logarithm bounds \(R(c)\)
by \(C|t|\delta_k L^{q-k}\) on a level-\(q\)
row of the old box.

The relative path multiplies the actual
unperturbed row field \(V_0\) on the left.
Use the same native node gauge for both
factors. The datum of the box solve is the
absolute log
\[
 b_s(c)=\frac1i\log(e^{iD_s(c)}V_0(c)).
\]
On the native second-removed path with
the common rebuilt box base, the inner
datum \(b_0\) is the full average of its
top-masked perturbation. Use the extra
row datum \(B'_s=b_s-b_0\)
in the old-box response with its original
truncated determining set. The product log
has \(B'_s=0\) off the support of \(R\) and
gives \(|B'_s(c)|\le C|s|\,|R(c)|\).
It is independent of \(u\), since \(D_{\rm in}\)
is independent of \(u\) wherever \(R\) is
nonzero. The admitted fluctuation margin
keeps \(\sup|B'_s|<\mathsf c_3a_0\) in
\cref{htc:inner-substitution}. Apply the local naturalizer
to this absolute row difference before
using the reduced response estimate.
Its fixed reach enlarges the source
support and leaves the separation from
the inner read. In level-\(j\)
units the increment of the potential, its
first difference and the assigned current
on the inner read is bounded by
\[
 C|st|\delta_k L^{j-k}\exp(-cM L^{k-j}) .
\]
Here the response norm is the one on the
native top read, and the distance is from
the enlarged transition to that read.
At \(j=k\) use the thinner inner collar
where this separation persists after
adding the complete transfer and
first-difference stencils.
At full old switches the potential and its
assigned current
are those of the native critical pair in the
box Landau representative.

Compare the row data before applying the
paired native box estimates. The difference
of the \(R\)'s is the difference of the
two full-average discrepancies with the
common cutoff on the first. For either mask
\(m=t+u\zeta_\square\) or
\(m=t\chi+u\zeta_\square\), first compare the
shifted fields at the one-step carrier.
Replacing the mask there by its value at
the coarse initial site costs at most
\(C_L\eta\operatorname{Lip}(m)\sup|a^\sharp|\)
in the relative log. Moving this constant
through the log costs at most
\(C_L(1+\|m\|_\infty)^2\sup|a^\sharp|^2\)
by Taylor's estimate at zero. The remaining
input difference is that of the transferred
unmasked angular shift and the flat shift,
multiplied by \(m\). The exact composition
for the plain relative log and its fixed-base
bound \(CL^q\) for \(q\) remaining stages
transfer these errors to the old row.
A difference between the rebuilt base
fields or their node frames is accounted for
in the same row comparison. Taking
differences in the analytic product log gives
\[
 |\delta B'_s|\le C|s|\bigl(
 |\delta R|+R_*(|\delta D_{\rm in}|+|\delta V_0|)\bigr),
 \qquad R_*=\max_{a=\sharp,\flat}|R^a| .
\]
This follows by writing \(B'_s\) as the
integral of the derivative of the product
log in the \(R\) direction. The common node
frames and paired full-row comparisons supply
the terms on the right. For paired unperturbed
row fields the last term vanishes by nonlinear
composition. With the native old-box
response estimate this compares the whole
segment on the admitted common top reads.
Apply the scalar real-to-complex comparison
on the admitted common lines before using
Cauchy's estimate in the window parameter.
For another continuation of this segment
one must first identify its assigned base,
row representative and current.

%
%
To use a paired response estimate at an activity,
transfer the actual evaluation pair. On an
admitted common read let \(W=\mathcal M U\),
let \(A^\sharp,A^\flat\) be the two response
potentials, and write \(J_{\rm tot}^a=J_{\rm base}^a+i^a\)
for their assigned currents. Use the paired
base currents \(J_{\rm base}^\flat=q_0J_{\rm base}^\sharp\).
The transfer of the fine pair averages
\(U'=e^{ihA^\sharp}U\) and uses
\(\mathbb P_{U'}J_{\rm tot}^\sharp\).
Apply the inverse of the node gauge \(v\) of
the straightened average to both components.
The resulting potential relative to \(W\) is
exactly \(\breve q(A^\sharp)\). At the start
of a coarse bond \(c\), the resulting current is
\[
 \operatorname{Ad}(v(c_-)^{-1})
                  (\mathbb P_{U'}J_{\rm tot}^\sharp)(c).
\]
On the full double-block stencils each path
transport and the node gauge change by at most
\(C_L h\sup|A^\sharp|\). This follows by
telescoping the bounded products and using the
regular contour mean on the admitted ball.
The current in this display therefore differs
from \(q_0J_{\rm tot}^\sharp\) by at most
\(C_Lh\sup|A^\sharp|\sup|J_{\rm tot}^\sharp|\).
No derivative of the current is used.

For the potential the local Taylor error of
\(\breve q-q_1\) is quadratic. Compare the
remainders on neighboring stencils after
parallel transport. The potential inputs then
differ by \(C_L h\sup|\nabla_U A^\sharp|\).
The transported base links differ by
\(C_L(h/\xi_y)^2\) when the normalized base
curvature is bounded in the local units.
This is the product estimate on the small
plaquette strips joining the translated paths.
The Taylor bounds with one more derivative
give the same mesh gain in the scaled
first-difference norm. Writing \(a_y\) for
the scaled first-order norm of the fine potential
on the full stencil and \(j_y\) for the scaled
supremum norm of the total fine current, the
errors beyond the paired response comparison
are bounded by
\[
   C(\eta/\xi_y)a_y^2
     \quad\text{and}\quad C(\eta/\xi_y)a_y j_y .
\]
The remaining current difference is
\(q_0 i^\sharp-i^\flat\) because the base parts
cancel. The remaining potential difference
is \(q_1 A^\sharp-A^\flat\). The relative log
between the two perturbed fields obeys the
same first-order bound, by the product and
logarithm estimates on the small chart.
The response comparison thus supplies the
evaluation mismatch on these complete reads
once its two native data and currents have
been identified. For the reduced currents
this use is on the common interior top reads.

%
%
To pass from a real evaluation mismatch to the
value of a gauge-invariant activity, fix a real
vertex and use a relative gauge for which the
potential, its first covariant difference and the
current mismatch have size at most \(2b_0\vartheta\),
where \(\vartheta\) bounds the infimum in the
mismatch norm. A zero infimum is handled by
taking the bound to zero.
Write the flat evaluation as \((\mathbf U,\mathbf J)\)
and the two mismatches as \(A,I\). Use the scalar
activity on the disc of pairs
\[
 \bigl(e^{i\eta_j w A}\mathbf U,\ \mathbf J+wI\bigr),
 \qquad |w|\le c/\vartheta.
\]
The flat evaluation must have the half field radii
and a fixed margin below \(b_0\) in the current slot.
Use \(0<b_0\le b_1\), \(b_0+b_1\le a^\sharp\) and
\(b_1^2\le\mathsf c_\natural b_0\). The outer substitution
theorem admits the disc when the potential and
its covariant difference stay below
\(b_0/[128K_{\rm red}(1+\mathsf c_\natural)]\),
and \(|\mathbf J+wI|<b_0\). The fixed choice of \(c\)
ensures both bounds. When the mismatch is
small enough that the disc radius is at least
two, Cauchy's inequality bounds the change
from \(0\) to \(1\) by \(C\vartheta\) times the
activity norm. Larger mismatches use the
individual bounds. The value at \(1\) is
the fine evaluation after transfer and gauge
change, so gauge invariance identifies its value.
The relative gauge is chosen only on real data.
For complex interpolation use the scalar
difference of the original two evaluations,
which are holomorphic on their admitted lines.

The single-lattice bound on the larger polydisc retains the
decay in the size of the carrier. Combining the smallness at
the vertices with this decay costs part of the localization
exponent, which pays for the $2^N$ terms. In the inner term
subtract the value at the flat pair and use the vanishing
of the first derivative there. The quadratic bound on the two
nearby linear-response evaluations then retains the small
radius factor in the difference. To retain an age weight
as well, the one-run estimate for this difference uses the
second renormalization at a H\"older exponent strictly
larger than the exponent in the age weight. The reserve in
the age exponent pays for interpolation.

The resulting regular commutation defect of an input formed
with $j$ fine levels has a factor
$\omega_{k+1-j}\min\{1,CL^{-\alpha j}\}$ in the step with $k+1$
fine levels, with some $\alpha>0$. The unmatched input at the
finest scale has the age weight $\omega_{k+1}$. A marked
remainder uses its squared radius ratio in place of the age
weight. To obtain the regular two-run bound, sum these defects
and join the two inverse covariances and the two auxiliary
Gaussian kernels by line segments in their admitted complex
balls. Cauchy's estimate on the available radius bounds the
change in the small-field expansion by the size of the Gaussian
and potential differences. The change in a scalar weight costs
$C|x^\sharp-x^\flat|/\mathsf t$ on a disc of radius $c\mathsf t$.
The Gaussian normalizer is compared separately.

In the covariant regular chain this gives, at equal deviations,
an inequality of the form
\[
 D_s\le C\lambda^{K-s}
              +\varepsilon\sum_{s<t<K}\omega_{t-s}D_t ,
\]
where $D_s$ is the distance between the families at depth $s$.
Choose $\lambda<1$ larger than the geometric rates of the
mesh defects and age weights, absorbing polynomial factors in
$K-s$. Choose $\varepsilon$ with
$\varepsilon\sum_{m\ge1}\omega_m\lambda^{-m}<1$. Downward
induction then gives $D_s\le C'\lambda^{K-s}$. The constant-source
coefficient is read from this distance by the same quadratic
test. The full step still needs the sourced estimate on
position-dependent arrays and the matching of the marked
histories. Those terms have no common covariant quadratic
cancellation to which the regular age estimate can be applied.

\paragraph{Sharp background tests on the matched histories.}
The approximate fluctuation test of
\cref{sec:large-field-the-tests-the-nested-regions-and-the-content-of-a}
compares the nonpivot values of
$V_k(\mathcal M^kU_{t,\square}(V_t))^{-1}$ with the identity.
The threshold is $2\delta_k$. Even when the two retained
fields agree, the regional background $U_{t,\square}$ is
computed with a different bare depth. Work on a common
partial history up to the first differing decision. One must
first compare the availability of the reads, since the test
convention assigns failure to an undefined read.

%
%
Consider a newly inserted ordinary plaquette test on a
cube \(\square'\) tested in both runs. On the common
prefix the incoming positive density tests cover the entire
footprint of the one-step average for the lifted datum on
\(\square'^{\sim4}\). To see this on either lattice, measure
lengths in units of the incoming \(k=t-1\) lattice. Put
\(s=LM_2\check R_t\) and \(W=LM\check R_t\). The four
enlargements cost at most \(4s\), and completing the
one-step averaging blocks adds at most \(L\).
A level-\(k\) density cube covering a corner of a plaquette
in this footprint has side at most \(Ls\), by the
neighboring-scale bound for \(\check R\). The total
\(4s+L+Ls<W\), since \(W\ge L^2s\). The ordinary
plaquette test region has four coarse layers to the
old large region. Even in the first step there is one
coarse layer to the hull of the failed bare tests. Thus
all the incoming density cubes just selected have their
positive test, and none is in an omitted second-class inset.

The full bare averaging blocks of a level-\(k\) plaquette
with a corner in such a density cube lie in its one-layer
enlargement. The positive density test gives bare
curvature less than \(\varepsilon_k\eta_k^2\) there.
The local small-curvature average estimate gives
\(|\partial V_k-1|\le2\varepsilon_k\) on the one-step
footprint, since the localized orbit has the prescribed
\(V_k\) there. At \(k=0\) use the bare test itself.
The ordinary averaging row prescribes \(V_t=\mathcal M V_k\)
on the top bonds inside \(\square'^{\sim4}\). Thus the
one-step bound gives \(|\partial V_t-1|<3L^2\varepsilon_k\)
there. At weak coupling
\(\varepsilon_1=3L^2\varepsilon_k\le a_1\) and
\(B_3\varepsilon_1\le\varepsilon_{\mathrm v}\le a_0\).
The real existence and orbit comparison theorem
therefore admits the lifted datum at the test radius
\(\varepsilon_{\mathrm v}\), independently at both cutoffs.

A newly inserted approximate fluctuation test is reached
only on a cube whose own plaquette test passed. It uses
that same orbit on the fourfold box. The third enlargement
lies in the top member with a margin containing the full
averaging blocks of its top sites. Earlier determining
rows have no endpoint there, so the only nodes in the
axial-gauge carrier are its top sites. On each of their
bare blocks the orbit curvature is bounded by
\(B_3\varepsilon_1\eta_t^2\) by the admission just proved.
Local averaging
bounds the level-\(k\) plaquettes in the contour blocks by
\(2L^{-2}B_3\varepsilon_1\). At the same weak-coupling floor,
\(L^{-2}B_3\varepsilon_1\le c_2^{\mathrm{elem}}\), so the
contour means exist.
For \(k=0\) use the bare curvature bound. The regular orbit
on this box therefore has an axial representative with a
unique center on the tested bonds.

These admission checks apply to ordinary tests at their
first insertion on the common prefix. A carried test
reproduced on its original reads belongs to its original
decision.

%
%
For a failed ordinary plaquette test, the
fourfold read can be exposed in the raw
incoming stock. Take the complete read
with every fine double block of its
top rows inside the raw footprint of
the failed component. In the first step
the incoming field there is raw Haar,
without earlier averaging constraints.
Keep the nonpivots fixed and use each
top row's private middle link to expose
its generated top output. The literal
row Haar kernel has mass at most \(2L^9\)
on a generated target. Distinct rows
read none of the other's pivots, so the
joint upper bound is the product of these
row factors times Haar on the top read.
The remaining factors use the true inverse
graphs and their reference majorants.

The admitted top curvature bound on this
read puts its axial logs in the ball
of radius \(CD\varepsilon_k\), for a box
of side \(D\) in the test units. The
radial argument at the ordinary threshold
\(\varepsilon_t\) then gives the factor
\(C n\vartheta\) in the top Haar bound.
The row costs grow at most exponentially
in the number of top rows of the read.
This calculation concerns the first
split and a later reference stock in which
these generated values remain free. At a
later ordinary step the selected incoming
pivots must remain free after the earlier
eliminations. A subsequent prescribed
outgoing row can constrain the exposed
top values and must be included in that
reference chart.

The marked tests have their own admission arguments.

%
At stage zero the first regularity guard uses the
untouched incoming datum, with the incoming profile.
Its positive branch enters the first prospective
guard. For that guard and the later prospective
guards, use the promoted datum bounds on the entering
support. These bounds, \(C e_q\) in each layer \(q\), use the actual
promoted averages before their delta functions are
integrated. Localizing the stage to its class component
gives the joined determining set with its prescribed
collar. A tested holonomy in the joined set either
agrees with a tested holonomy on the global datum,
is one in the lifted collar, or equals a tested
top holonomy of the datum being lifted. Hence the
joined datum is \(\bar\varepsilon\)-regular for
\(\bar\varepsilon>C\max_q e_q\). At weak coupling
take \(\bar\varepsilon\le a_1\) and
\(B_3\bar\varepsilon\le\varepsilon_{\mathrm v}\le a_0\).
The regular orbit for the localized set in force
therefore exists before the prospective guard is applied.
The inset and crossing tests that use this same orbit
also have defined reads. Passing their layer thresholds
requires their separate curvature estimates.

%
For the last-shell test take \(h=m_\flat\) and the
full fourfold test box in \(Z''_{h+1}\cap\Omega_h\).
The entering prepared datum bounds its original
level-\(h\) plaquette defects by \(C e_h\). Keep the
specified links meeting the inner cutoff. Every link
set to one belongs to \(\mathbb B_0\), so the preceding
positive main-shell test and the basic reference bound
give distance at most \(\delta'_m+C L^{h-m}\varepsilon_m\)
from one at that link. Telescoping at most four replaced
factors bounds the plaquette defects of the artificial
datum by
\[
             C e_h+4\delta'_m+C L^{h-m}\varepsilon_m .
\]
This lies in the fixed lifted-datum existence radius
at weak coupling, independently of the waiting depth.
This proves admission before the last-shell decision,
whose half-threshold is a smaller requirement.

%
On a passing last-shell cube the local test transfers
to the glued background \(U^0\) at the same
plaquettes. In the tested curvature units its
bound is
\[
 \tfrac12\varepsilon_h+
       C\varepsilon\bigl(L^{-\mathsf N}
                          +e^{-\delta_0D_2}\bigr).
\]
Here \(\mathsf N\) is the waiting length and
\(D_2\) is the clearance of the tested enlargement
inside its full fourfold box, in level-\(h\) units.
This is a one-run comparison of the two lifted
fields from the same inner datum. The band for
the last-shell decision is a function of that
retained level-\(h\) datum. In the framed integral
it stays fixed through the \(V'\) integration
at fixed datum, so its thickness must be estimated
in the retained-field law.
On the identity side of the mask, a tested
plaquette with a complete vacuum top box
around its carrier has normalized curvature
at most \(Ce\,e^{-cD}\), where \(e\) is
the admitted radius and \(D\) the distance
to that box's complement in level units.
This is the real collar comparison with
the vacuum at the same retained data.
Such plaquettes cannot attain a maximum
in the threshold band once this bound is
below \(\varepsilon_h/4\). The remaining
test depends on the actual determining
links of the retained \(h\)-datum. The
ordinary radial change holding those links
fixed gives no shell gain for this test.

%
%
The effect of a change of site gauge before masking
can be assigned to the erased rows. Write \(J\) for
the map setting the erased links to one. With gauge
action \(V_b^g=g(b_-)V_bg(b_+)^{-1}\), pull
\(J(V^g)\) back by \(g^{-1}\). Its kept links are
those of \(J(V)\), and its erased links are
\(g(b_-)^{-1}g(b_+)\). The plaquette norms of the
constrained background are invariant under this
gauge pullback. In the lift to the complete test
tower an induced link is either unchanged or
copies one changed link between level-\(h\) blocks.
If the erased-link jumps have logarithms of size
at most \(b\) and their logarithmic interpolation
lies in the common regular domain, apply the
curvature response bound along it. In the units
of the test the change at a plaquette is at most
\(Cb e^{-cD}\), where \(D\) is its tower distance
to the affected determining rows. Use this
comparison in each fixed tower before comparing
the two cutoffs. It isolates the gauge contribution
to the band on the actual read.

%
%
There is a shell estimate on a different chart,
when all kept links in the complete test read
are free and can be scaled together from the
vacuum. In a wholly kept box with a full
product Haar stock, fix a rooted axial tree.
The tree change of variables leaves independent
Haar on the \(n\) non-tree links. The bound in
\cref{lem:averaging-1a-3-1} puts their logarithms
in a ball of radius \(r\le C D e_h\) in the
maximum bond norm for a box of side \(D\).
The test is gauge invariant on this whole read,
so the tree gauge factors can be held fixed.
In a box meeting the cut use this calculation
only on a product log-Haar chart in which
scaling all kept logs preserves the identity
mask.

Along a ray write the log datum as \(tB\) with
\(\sup|B|=1\), \(0\le t\le r\). The kept logs
give the corresponding lifted row logs, and
the masked row logs are zero. Thus the vacuum
continuation of \cref{prop:pnc-critical-continuation}
admits a fixed complex disc in the radial variable,
on the separated tower with complete top test
stencils. For a tested plaquette \(p\) let
\(F_p(t)\) be its matrix curvature increment
divided by the squared bare mesh in level-\(h\)
units, in the holomorphic slice. The uniform
bound \eqref{pnc:physical-fine-bounds} bounds
this function on the disc, and \(F_p(0)=0\).
Use \(r\) at most half the disc radius.
The real norm is the norm on the actual orbit.
Cauchy's estimate gives
\[
             |tF_p'(t)-F_p(t)|\le Cr^2 .
\]
Put \(\lambda=\varepsilon_h/2\). For real
fields \(F_p\) is a quaternion, so the
operator norm is Euclidean in its four real
coordinates. Where \(|F_p|\) lies between
\(\lambda/2\) and \(3\lambda/2\), its
derivative in \(\log t\) is at least
\(\lambda/2-Cr^2\). If \(Cr^2\le\lambda/4\),
every active maximizer has slope at least
\(\lambda/4\) in this range. The band
of relative width \(\vartheta\le1/4\)
for the maximum can be crossed only
upward. It therefore occupies an interval
of length at most \(8\vartheta\) in
\(\log t\) on any ray.

Along each ray the Haar radial density is
\(t^{3n-1}w(t)\), where \(w\) is the product
of the sinc-squared Jacobians and is
decreasing on the small ball. Its cumulative
mass \(H\) satisfies \(tH'(t)\le3nH(t)\).
An interval of logarithmic length
\(8\vartheta\) therefore contributes at
most \(24n\vartheta\) times the mass of
the whole ray. Integrating the ray
directions gives the same relative
bound for the product Haar ball.
For \(n=0\) the band is empty.
To use it in the weighted history estimate,
these retained coordinates must survive
the outgoing inverse as a product stock.
The band is taken in the stored frame
before the exterior supremum. On a wholly
kept box the exterior compensator acts by
a gauge on the full datum, so the band
is the same for every admitted exterior.
For a cut box use the actual mask-preserving
pullback to check this property. The
remaining reference integral must have
its bound uniformly in these coordinates.

These conditions hold for the reference
integral of an unmarked late last-shell
drop-out before a later constraint, when
the full level-\(h\) box read, including
its endpoints and lift halo, is kept and
lies in \(\Omega_h\). In (1.145) those
links are in the \(dV_h\) Haar stock,
separate from the relative field \(V'\).
The ordinary strip pivots at fixed \(V_h\)
are already eliminated. Unresolved raw
targets have both endpoints in
\(\Omega_h^c\) in the old core and are
outside this stock. With the test stock
fixed, eliminate those remaining targets
against their own averaging constraints
and keep the old strip variables in their
original charts. The unmarked late
product envelope after
\eqref{hi:reference-cutoff-gaussian} then
has the same conditional integral bound
for every value of the test stock.
The radial restriction gives the stated
gain in that reference integral, with
any already justified event payments.
Subsequent constraints that prescribe
some of these links require the bound
on the corresponding conditional chart.

%
The same reference calculation applies to
a test box crossing the identity cut if
its full level-\(h\) read lies in \(\Omega_h\)
and all its kept links are in the retained
Haar stock. In the class geometry the inner
cut is one coarse layer from the hull of
the old core and two layers from the
outer prepared face. In level-\(h\) units
write \(A=LM\check R_{h+1}\) for the coarse
layer and \(d=LM_2\check R_h\) for the test
cube side. A fourfold box crossing the cut
has width \(9d\). Even after the endpoint
padding, \(9d+2<A\), by
\(\check R_h\le L\check R_{h+1}\),
\(M\ge L^2M_2\), and \(L\ge13\).
Thus the full padded read clears the old
core and stays in the prepared member.
Its kept links, including those crossing
the cut, are in the second Haar stock of
(1.143). The gauge-nest convention excludes
the crossing links from \(\mathbb B_0\),
so (1.145) retains them in \(dV_h\).
Use this stock after the compensating
change of (4c). At fixed exterior and
frame this change is a Haar translation
at the endpoints of every kept link.
With those translated coordinates fixed,
the last-shell test sees their masked
field directly. Read the other factors
at the inverse translation. The conditional
reference bound is uniform in those
values, so the band in these coordinates
can be restricted before the exterior
supremum.

To obtain the free logarithmic chart, form
the finite graph of the full read and
contract each connected component of its
identity edges. On the quotient choose a
rooted spanning tree from the kept edges.
The tree change of variables is triangular:
after fixing the tree links it multiplies
each non-tree link on the left and right
by fixed group elements. It leaves
independent Haar on the non-tree links.
The site gauge is constant on each
contracted component, so every identity
link stays at one.

If the artificial plaquette defects on
the read are at most \(e\), first apply the
ordinary axial estimate on the full box.
All its transformed links differ from
one by at most \(CDe\). Along any original
identity edge the two axial site gauges
therefore differ by at most \(CDe\).
Replace the gauge on each contracted
component by its value at a chosen vertex.
Paths in each component have length at
most \(CD^4\), so this changes the gauged
kept links by at most \(CD^5e\).
Choose the quotient tree by shortest paths
from the root. Its depth is at most \(CD\),
since contraction cannot increase distance.
Putting the field in that gauge costs at
most another \(CD\), by telescoping the
tree paths. Thus the non-tree logarithms
lie in the product ball of radius \(CD^6e\) when
that radius is small. Apply the radial
estimate when its square meets the
half-threshold condition. The whole
kept graph must be free in this calculation,
including the kept crossing links.

%
%
The same last-shell restriction applies on a marked
late branch whose triangular lower reference chart
leaves this full kept Haar stock. Hold its recorded
post-gauge coordinates and the inner variables fixed
while integrating the lower Lebesgue blocks. The
uniform inverse Jacobian and Gaussian bounds give the
same conditional majorant \(e^{C V^\circ}\) as for
the triangular reference integral. The bounds for
closed inner leaves remain functions of their own
stored variables. The last-shell band is Borel in the
fixed kept stock and can be inserted before the
compatible-evaluation supremum. Integrate it with
the Haar and quotient calculation above. This gives
the factor \(C n\vartheta\) in the weighted outer
integral, with the same inner masses and charged
events. The argument requires the stock in the
recorded coordinates to remain free after the lower
eliminations. For the comparison of runs, the paired
reads must place their disagreements in this band
uniformly over the admitted evaluations.

%
The final density test that classifies a restored component
uses the reference draw. On that support the restored
top field has plaquette defects bounded by \(C(M)\varepsilon_m\).
The estimate follows by changing the at most four
plaquette factors from the reference field, including
on the face where the reference agrees with the exterior.
Each fourfold density box based in the six-cell inset
is contained in the component. Apply the same lifted-datum
existence estimate on those full boxes before either
classification outcome.

%
%
On a complete fourfold test read wholly
in the fresh box, first use the normalized
reference law at fixed exterior, before
prescribing new outgoing rows or splitting
the two classes. In a fixed frame the
denominator is the centre action factor
times \((\theta^{-3/2}\sigma_0)^{n_{\mathsf G}}\)
times the admitted fluctuation normalizer.
The latter has bounded log. The centre
action costs at most
\(C p_0(g_m)^2(M\check R_m)^4\)
in the exponent, and \(n_{\mathsf G}\) is
at most \(C(M\check R_m)^4\). Hence the
inverse denominator is bounded by
\(\exp[C(\log\theta)^c]\) on this slice.

The numerator density is at most one
relative to the off-tree Haar product.
In physical coordinates the non-tree
translations preserve Haar. For the
gauge-invariant test and its top curvature
bound the choice of tree values is
immaterial, by site gauges fixing the
root. Thus one may integrate the tree
values against Haar as well when bounding
that event, and discard fresh variables
off its read. The top curvature bound
places the read in the small axial
logarithm ball. The radial estimate at
threshold \(\lambda=\varepsilon_m\)
gives Haar band mass at most
\(C n\vartheta\), where \(n\) counts the
free links of the read. The fixed-exterior
probability of this band is at most
\(C\vartheta\exp[C(\log\theta)^c]\).
For the weighted history envelope after
lower eliminations, the same radial
factor applies where the complete product
stock is still free, the remaining
reference bound is uniform with that
stock fixed, and the band is Borel in
the recorded coordinates.

The six-cell inset places the test box
inside the component. A test box entirely
outside the restored stock reads only the
retained field, and a box meeting the
restoration boundary can read fixed nonzero
links as well as the new draw.

%
%
There is a crossing case with the fixed
links separated from the tested stencils.
Fix the physical tree values first. On a
chart whose remaining free links have
product Haar reference, use the tree
frame to write the retained logs as
\(B_{\rm fix}\) and the free logs along a
ray as \(tB\), \(0\le t\le r\). The fixed
logs are held at every \(t\). If all the
fixed edges, including the physical
reference-tree edges, extend to a common
tree with compatible prescribed root gauges,
their values can be gauged to the identity.
Otherwise retain their actual \(B_{\rm fix}\).
The
fixed-coordinate part of the lifted
datum is formed by the same block copies.
Let \(F_p(t)\) be the normalized matrix
plaquette read of this datum. On the
fixed small response disc the previous
Cauchy calculation now gives
\[
 \bigl|t F_p'(t)-(F_p(t)-F_p(0))\bigr|\le Cr^2 .
\]
It is enough that
\(\sup_p|F_p(0)|+Cr^2\le\lambda/4\).
The active norm slopes in the band
are then at least \(\lambda/4\), giving
the same \(C n\vartheta\) radial bound.

To check the centre condition, interpolate
from the vacuum to \(B_{\rm fix}\) with
the free logs zero. When the complete
lifted carriers of all nonzero fixed logs
are at distance at least \(D_0\) from
the tested stencils, the weighted response
\eqref{pnc:weighted-first-response} gives
\(\sup_p|F_p(0)|\le C\|B_{\rm fix}\|e^{-c D_0}\)
in the normalized curvature units. Thus
the radial comparison also applies when
the full test box meets the old field but
its tested stencils and their padding
are separated from the nonzero fixed
data by the prescribed collar. This
uses the actual product stock in the
fixed tree frame and the admitted
interpolated data.

For a crossing chart with unscreened
fixed rows we must control the centre
curvature on its own constraints. The two class
weights are compared before removing any
implied reference cutoff or frame factor
on the first-class side.
In the quiet nest, the final test bases in
the old enlarged hull lie in the region
excluded from the earlier ordinary
plaquette-test insertion. The inherited
top plaquette bound there admits the
background, while the band of its derived
bare curvature still needs control in the
retained-field integral.

In each case compare like reads on the part of the
history where the earlier decisions agree. For a plaquette
test the numerical quantity is the curvature at the
finest scale of the local minimizer divided by \(\eta_t^2\)
on each respective lattice.
A bound on the response at the common top alone does
not control that read.

%
There is a mesh comparison on a matched input to
these plaquette tests. For the ordinary and final
density tests, take an aligned fourfold box and
the same level-\(i\) field on the entire lifted read
in a common frame. The third enlargement lies in
the top member of the separated tower. The tested
first enlargement and its complete difference
stencils are therefore inside that member. If
the top plaquette defect is at most \(e_i\), the
common comb gauge bounds the top logarithms by
\(N=CM_2\check R_i e_i\), and the block lift copies
them to the induced rows. The comparison of common
trial fields below now applies to this read: the
complete common row data coincide, the lower
unmatched rows are screened, and on the top
stencils the first differences have the same
two-mesh and H\"older bounds. Pair the literal
plaquettes of each orientation by inward rounding.
For a bare mesh \(\eta\) in level-\(i\) units this
gives \( |q_\sharp-q_\flat|\le
CN\eta^{\min(\alpha,\beta_H)}\) for the maxima.
It also applies to a last-shell box when the
complete masked top datum is common. Use the
ordinary incoming or final reference curvature
bound for \(e_i\), or the admitted masked-field
bound for the last shell. With the running
threshold bounds their ratios to the common
test threshold cost at most powers of the scale
parameters in this comparison. This gives the numerical part
of the test comparison on matched input.

On the region where the two reads are defined and the lesser
value is at most twice the threshold, we must show a relative
comparison with a positive power of the ratio of the bare
mesh to the tested scale, allowing the powers of $\log\theta$
in the history estimates. If the
two tested numbers $q_A,q_B$ differ by at most
$\vartheta\lambda$ at threshold $\lambda$, then
\[
 |\mathbf 1_{q_A<\lambda}-\mathbf 1_{q_B<\lambda}|
 \le \mathbf 1_{|q_A-\lambda|\le\vartheta\lambda} .
\]
For the approximate test $q_A$ can be the maximum over its
nonpivot bonds. The shell is then determined by at least one
bond attaining that maximum. The same argument applies to
each subsequent decision on the part of the history where
the preceding decisions agree.

On a disagreement with \(\vartheta\le1/2\), both
numerical reads are at least \(\lambda/2\). For the
tests with an energy witness obtained by forcing
the read below threshold at small energy, run that
calculation with threshold \(\lambda/2\). In the
ordinary plaquette and approximate tests this reduces
the witness constant. In the main and last-shell
tests it also requires a smaller bound on the
reference-field errors at weak coupling. The witness
can therefore be taken on either history, including
the passing one. Using it in the mass requires the
unspent share of the action or form on that same
history. A stopped root whose payment was reserved
for a later read uses that reserve only there.
This gives the corresponding ordinary-creation or
marked-event energy on the full compatible carrier.
Its use depends on which action or Gaussian share
is available there. The thickness of the near-threshold
set still has to enter its integral.

At fixed coarse datum the approximate test has a fixed
center on each tested bond. In the strip where the cutoff
equals one, the substitution gives $V_k(b)D_b^{-1}=W'X_b$
with $X_b$ fixed by the retained exterior datum and $D_b$
the cube center at the bond. Thus on this strip the read
is a radius of a translated stored coordinate. Where the
cutoff is less than one, the strip's center reserve keeps
the approximate test below $2\delta_k$. In the admitted
deep ordinary box the global radius bound and the local-to-global
offset also keep it away from $2\delta_k$. Both observations
use the margins in the ordinary step. To estimate the first
approximate failure, the runs must be compared before the
test has fixed $Q_t$ and hence $\Omega_t$. The ordinary
fluctuation chart on $\Omega_t^{\sim-1}$ is then unavailable
on the failing side.

%
%
Before the \(Q_t\) split one can use the
gauge draw itself on a tested footprint in
the gauged blocks. The approximate-test cubes
are inset one coarse layer in the gauge region.
Their two test layers and one bond reach stay
in that region by the cube-size floor. Fix the input before this
draw and the preceding plaquette decisions.
For a fine test bond choose an endpoint \(x\)
off the top lattice in a block where the
draw is made, and hold the other draws fixed.
If \(y\) is the block centre, the coordinate
\(W_x=V_k(y,x)u(x)^{-1}\) has the normalized
one-site density of the gauge-fixing factor.
Its Haar density is at most
\(C\mathsf t_k^{3/2}\), by the lower bound on
the gauge normalizer. Here
\(\mathsf t_k=g_k^{-2}\) is the coefficient
of the auxiliary gauge weight: it is \(x_k\)
in the own run and \(\theta_k\) in the reference
run. The transformation
is the identity at the top sites, so the
coarse average and the local approximate
test centre at the bond stay fixed. The
relative tested link is a fixed left and
right translate of \(W_x\) or \(W_x^{-1}\).

The Haar volume of the band about \(2\delta_k\)
is \(C\vartheta\delta_k^3\), for either side
of the threshold. Thus the conditional
probability for this endpoint is at most
\(C\vartheta\mathsf t_k^{3/2}\delta_k^3
\le C\vartheta p_0(g_k)^3\).
Sum over the bonds that can attain the
test maximum. Earlier approximate decisions
may be retained as indicators bounded by one
in this estimate. This gives a band bound
relative to the positive mass before the
gauge draw, for both outcomes of the test.
The residual half of the gauge exponent
in the ordinary-event majorant has the same
pointwise bound with the original normalizer.
Its mass is bounded uniformly above and is
at least one, since the positive exponent
cost has been halved. The same band factor
therefore applies in that residual reference
integral. Any test gain uses the assigned
Wilson and gauge share before this integral.
To apply it inside a later history mass,
the chosen endpoint coordinate and fixed
centre must still be available in its
reference integral.

For a bond from a block centre to an adjacent
site \(x\) in that block, the contour mean
is just the link and the local axial centre
is the identity on it. The test on this bond
is therefore radial in the actual \(W_x\).
Its threshold in eigenangle is
\(r_0=2\arcsin\delta_k\), inside the gauge draw
cutoff by the threshold ratio. In the
angular coordinate \(r\) the density is
proportional to
\(\sin^2 r\,e^{-\mathsf t_k(1-\cos r)}\) on this chart.
On both shells of width \(w=g_k/(C p_0(g_k))\)
at \(r_0\), its log derivative has absolute
value bounded by
\(C(\delta_k^{-1}+\mathsf t_k\delta_k)\le C'/w\).
The same bound holds for the residual half
weight. A one-sided band of width \(u\le w\)
has at most \(Cu/w\) times the mass of that
side's shell of width \(w\). Thus a relative
band about this single-bond decision has cost
\(C\vartheta p_0(g_k)^2\) in its own pass or
failure mass. Other tests that read this
same gauge variable can restrict the
comparison shells; they retain their
indicators in the actual prefix law.

%
%
On the failing side of a first ordinary
approximate decision, the enlarged failed
cube is outside \(\Omega_t\). A complete
endpoint block of the test in that enlargement
uses the old raw Haar field and the averaging
delta in the retained one-step operation.
Use the rooted tree coordinates of that block,
with the top site fixed and all incident
links of the moving sites included. They
separate the block quotient and the site
gauges with Jacobian one. At fixed quotient
the change from the site gauge to \(W_x\)
is a Haar translation and inversion. The
outgoing average is independent of these
site gauges. The centre is thus fixed
in the reference coordinates, both if
the outgoing datum is retained and if it
is eliminated against its own averaging
delta. The band is a function of the
recorded stock before the exterior-pair
supremum.

On an untouched piece with this complete
block free of incoming constraints on the
moving sites, keep the quotient fixed while
applying the residual product majorant.
Extract the assigned failure gain first.
The remaining one-site gauge factor is the
half-weight above, and restricting it to the
band gives \(C\vartheta p_0(g_k)^3\) with
the other factors in the same majorant.
The first step has no earlier averaging
constraints. At a later step the same
argument requires that the earlier constraints
and reads leave these gauge coordinates
available. A passing coordinate integrated
in \(\Lambda_t\) uses its own conditional
window.

%
There is a local estimate before the \(Q_t\) split.
Fix \(V_t\), the recorded plaquette decisions, and
the other free variables in a complete row chart
for \(\mathcal M V_k=V_t\). On a branch where the
actual middle pivots can be eliminated and \(b\)
is a free bond, its reference factor is Haar.
Reconstruct the pivots and evaluate the density
and gauge factors at those reconstructed values.
The centre \(D_b(V_t)\) is fixed. A band of relative
width \(\vartheta\) about the group-norm threshold
\(2\delta_k\) has Haar measure at most
\(C\vartheta\delta_k^3\). In \(SU(2)\) the
operator-norm distance from the identity is
\(2\sin(\phi/2)\), which has derivative bounded
away from zero at these small radii. Thus the
restricted integral in this slot is at most
\(C\vartheta\delta_k^3\) times a pointwise upper
bound for the reconstructed density in the slot,
including its pivot Jacobians. A union over the
free bonds that can attain the maximum gives the
bound on these charts. In the full pre-\(Q_t\)
history one must still control those densities
uniformly in the retained coordinates in the
order of the weighted mass integral.

%
The main-shell fluctuation test has a different shell
estimate. On the preceding positive exterior and
extendibility support, the box critical-orbit
construction gives \(V_\Lambda\) and the mixed
reference center \(V_0\) from the fixed exterior.
These are available before the main-shell decision
at each cutoff. In one retained frame use the relative
links \(V'(b)=V''(b)V_0(b)^{-1}\). First eliminate the
central delta chain from the top down. The relative
Haar translation and the normalized tree constraints
then leave the free indices \((\mathbb B_0\cap X)\setminus rT_0\).
The test reads the largest \(SU(2)\) eigenangle of
these links, with \(\phi(-1)=\pi\). Thus the
logarithm cut gives a numerical failure at the
threshold \(\delta'_m<\pi\).
On each free Haar factor the band
\(|\phi-\delta'_m|\le\vartheta\delta'_m\) has measure
at most \(C\vartheta(\delta'_m)^3\), by integrating
\(\sin^2\phi\,d\phi\). The condition is independent of
the exterior in the relative coordinates of that
run and frame. At a matched stored link and frame,
the difference of the two relative angles is at most
the group distance of the two centers by the
bi-invariant triangle inequality. The maximum obeys
the same bound. Comparing the centers between the
two cutoffs still uses the response comparison.
In the admitted passing-side product envelope
\eqref{hi:reference-cutoff-gaussian}
without earlier marks, use the product majorant
in the free \(B\)-coordinates after bounding the exterior
dependence. Replace a scaled \(B\)-ball of
radius \(\delta'_m/g_m\) by its inner annulus of
relative width \(\vartheta\). Its volume is at most
\(3\vartheta\) times the ball volume. Sum over the
free coordinates that can realize the maximum, with
the same other factors in the envelope. This gives a
relative \(C n_B\vartheta\) bound for that product
estimate, where \(n_B\) is the free-coordinate count.

%
%
On the failing side, consider a late drop-out
with no earlier marks in the core. The completed
frame operation leaves normalized Haar in each
non-tree shell slot. In its recorded frame it is
the relative link \(V'\) that is held fixed in
the exterior supremum in \eqref{hi:drop-factor}.
Changing the exterior reconstructs \(V''=V'V_0\)
at that same slot. The shell indicator is
therefore a function of the stored variable.
It can multiply the pointwise majorant before
the exterior supremum in
\eqref{hi:actual-pointwise-envelope}.
The unmarked late reference bound after
\eqref{hi:reference-cutoff-gaussian} is independent
of the shell Haar variables. Restricting that
product envelope to the band gains
\(C n_B\vartheta(\delta'_m)^3\) relative to its
full Haar integral, with the same strip factors.
Any event payments used here are the ones
already justified on that failing history.
To use the bound for disagreement, the paired
centres must place the disagreeing records in
this band uniformly over the admitted evaluations.

The failing-side gain also holds for marked late
branches with the triangular reference charts of
\eqref{hi:reference-gaussian-integral}, provided
the selected shell slots remain free Haar factors.
Write the outer reference measure as
\(dv_{\rm sh}\,d\mu_{v_{\rm sh}}^\circ\), where
\(v_{\rm sh}\) comprises the relative shell slots
and \(dv_{\rm sh}\) is product Haar. Keep
\(v_{\rm sh}\) and the inner variables fixed
during the successive integrations of the
lower Lebesgue blocks. Their inverse Jacobian
bounds and Gaussian reserves are uniform in
these group arguments. The strong majorants
of the closed leaves in
\eqref{hi:actual-pointwise-envelope} depend on
their own stored variables. Thus
\[
 \int g_\omega(v_{\rm sh},\cdot;\xi_{\rm in})\,
               d\mu_{v_{\rm sh}}^\circ \le e^{C_5 V^\circ}
\]
uniformly in \(v_{\rm sh}\) and \(\xi_{\rm in}\).
Integrate only the shell band in \(dv_{\rm sh}\).
The right side then acquires the factor
\(C n_B\vartheta(\delta'_m)^3\) before the same
sum and integration of the inner histories.
The event payments and volume weights are those
of the original history.

This applies with closed inner leaves and older
regular strip charts whose pivot eliminations
are complete. An older raw row with a free Haar
target can also be eliminated by substituting
its average for the target in the remaining
reads, starting from the highest such row.
Here the substitutions must leave the same
lower reference charts, with their bounds
uniform in the fixed shell. Fixing a shell link
can change the coarse datum of an ordinary
pivot chart. The unresolved lower raw targets
of the frozen core lie inside the inner cut.
They are separate from the native shell slots.

For prescribed targets in the first outgoing
rows of the late stop, fix the reconstructed
exterior and the old shell. Each row that reads
the fresh box has a free middle Haar link,
including a row crossing its face. The old
frame tree acts on the old field. Use the joint
inverse of the fresh rows at the prescribed
targets, with their native branch version, and
evaluate the retained weights at the reconstructed
links. The factor per row is bounded by \(2L^9+1\),
uniformly in the old shell. These factors use
the existing first-outgoing volume count.

An older restored box also gives its joint
first-outgoing density at prescribed targets in
its protected row stock. There is a two-level
case in which one such target is the fixed
middle link of the next raw row. Suppose the
complete two-level backward blocks of this
row lie in that restored box and the row is
still carried in the ordinary raw graph.
Its alternate face link belongs to the same
first-outgoing stock and is left free.
The restored cutoff bounds the plaquettes on
the backward blocks. The local averaging estimate in
\cref{prop:averaging-1a-1-6}, applied to the complete
two-level carriers, gives small curvature on the next
double block. The regular alternate
pivot then gives the joint bound at the fixed
middle link and fixed next output, with the
same other stored variables. Both inverse
densities enter the old core integral. They
are uniform in the relative shell and use
its slots only through the fixed parameters
of their charts. A prescribed output in a
different graph still requires the joint
bound for that graph with the shell fixed.

%
%
For a sharp preparatory collar test there is a
conditional shell bound in its native free exterior
stock. Fix the other stage variables, including the
top datum determining the centre. In the scaled
coordinate the threshold is \(t=\delta'_i/g_i\).
Use the width \(W=(C_{\rm sl}p_0(g_i))^{-1}\).
The surviving ordinary completion factors use the
separate strip stock. The other test on this
coordinate is the approximate collar bound at
\(2\delta'_i\), with the regularity read fixed in
this conditional integration. In the local convention
the two centres coincide. In the global convention
their scaled separation is at most \(W/4\) by the
collar response. Since \(W=o(t)\), the approximate
test passes on both one-sided shells of width \(W\).
The density slope in \eqref{aw:stage-slope} is at
most \(2/W\) there, using the waiting-window comparison
of \(p_0\). The radial Jacobian varies by a fixed
factor. Consequently restricting one of these shells
to distance \(u\le W\) from the threshold gains
\(Cu/W\) in the conditional stage mass. In angular
units a relative width \(\vartheta\) corresponds to
\(u=\vartheta\delta'_i/g_i\).

In the part of the outer collar receiving
a fresh normalized gauge draw, the same
one-site argument applies to the local sharp
test before its decision. Take a tested bond
with a noncoarse endpoint in a complete new
block and keep the other variables fixed.
The site gauge must fix the determining nodes
of the prospective centre. The relative link
then has a translated one-site density bounded
by \(C\mathsf t_i^{3/2}\). A band at angle \(\delta'_i\)
has conditional mass at most
\(C\vartheta(\delta'_i\sqrt{\mathsf t_i})^3
\le C\vartheta p_1(g_i)^3\).
On this band the same-centre approximate
test on the bond passes. This calculation
uses the untransformed local prefix, on
either side of the sharp threshold. If the
bond joins the block centre to its neighbour
and the stage uses the same axial gauge, the
test is radial in that draw. The preceding
one-sided shell argument with \(\delta'_i\)
gives the bound \(C\vartheta p_1(g_i)^2\)
relative to the single-bond outcome. For
the full maximum use the predecision union
bound over the bonds with such a draw.
Carried strip links and the other selectors
still use their own reference factors.

For a stopped component, the top datum determining
this centre can be reconstructed from the carrying
field in the compatible-evaluation supremum. The
conditional collar bound fixes that datum first. To
use it in the marked integral we need a relative
shell coordinate in the actual stored stock, or a
uniform centre-spread bound in a fixed chart.

%
%
On a global sharp failure with the full sharp mask in
the transformed free exterior stock, this band is
fixed in the stored coordinates \(W'_b\). For relative
width \(\vartheta\le1/8\), the band about \(\delta'_i\)
lies below the onset \(4\delta'_i/3\) of the cutoff
substitution. Indeed, above the onset the original
global angle is at least \(4\delta'_i/3-\delta'_i/8\),
outside the band. Below the onset it is exactly the
angle of \(W'_b\). The substitution there is a right
translation with Haar Jacobian one at fixed exterior.
The maximum band is therefore Borel in the stored
stock. On the part of the weighted envelope with
these slots still free, restrict one slot that attains
the maximum and use its Haar volume
\(C\vartheta(\delta'_i)^3\). Summing over the free
slots gives the band bound with the same conditional
majorant for the other variables. A tested bond
outside the transformed mask has to keep its own
read. In particular the local sharp failure uses
the preceding conditional calculation on its
untransformed prefix.

For the other marked histories we must show that restricting
the weighted mass of the operation to this shell gains a factor
$C\vartheta^a\exp[C(\log\theta)^c]$ over its large-field bound,
for some $a>0$.
Here $\theta$ is the frame at the finest depth held in the
histories counted by that operation. The comparison must also
sum these restricted masses over the possible first decisions
with the localization weights of the recursion. The window
estimate \eqref{aw:relative-window} uses
\cref{lem:aw-split-reads}, which fixes the other stage
coordinates and varies a literal nonpivot logarithm. To use
it on a derived background read would require a change of
variables controlling the conditional mass and preserving
the readsets of the surviving history. Even for a literal
translated coordinate the remaining tests can have distinct
centers. The full shell comparison uses the stage law after
summing the earlier gauge draws. A bound on the normalized
positive survival probability at a fixed exterior pair
does not bound the mass in \eqref{hi:weighted-mass}, where
the exterior-pair supremum and the mark weights enter before
the integration and summation. For the extendibility
test the infimum is a function of the prescribed top
exterior alone. Its shell indicator is constant on the
central fibre at fixed exterior.

%
%
On the complete exterior carriers of the
extendibility witness there is a uniform separation
at low energy. Let \(E\) be the assigned exterior
Wilson share at level \(i\). The axial extension
with the old exterior and the lifted-datum estimate
give a candidate whose tested curvature, in
level-\(i\) units, is at most
\(C(M)\sqrt{E/\theta_i}\). The infimum is no
larger. Hence \(E\le c(M)\theta_i\varepsilon_i^2\),
with \(c(M)\) chosen for the half-threshold,
puts the read at most \(\varepsilon_i\). This
leaves a margin of at least \(\varepsilon_i\)
to the test value \(2\varepsilon_i\).
Apply this on both runs when the complete carriers
are available in both. There is no disagreement
on that low-energy part. On the rest, a shell gain
requires a bound using the exterior read. The
central fibre estimate supplies no such bound.

%
%
The two infima can be compared on regular trials
near their common threshold. For a trial \(V\) let
\(q_b(V)\) be the normalized tested supremum at
cutoff \(b\), and put \(q_b^*=\inf_V q_b(V)\).
The box clearance puts the full averaging
blocks of the top plaquettes changed by the trial,
including its outer layer, in the tested top
inset. If \(q_b(V)\le5\varepsilon_i\), the local
averaging bound gives defect at most
\(2q_b(V)\) on each of those coarse plaquettes.
With the prepared exterior regularity on the
unchanged carriers, the same trial bonds are
therefore admitted at the other cutoff, at
weak coupling. Use the actual extension class
over that exterior.

%
%
%
%
On aligned class boxes the near-threshold
trial supplies a small datum on the entire
lift. A top plaquette of the box whose full
averaging carrier is inside the tested inset
has defect at most \(10\varepsilon_i\) by
the local average estimate. The remaining
top plaquettes are within a fixed top
stencil width of the outside of that inset.
Choose a density cube containing one of
their corners. It lies in the carried
positive region \(\Omega_i^{\sim4}\).
Its one-layer enlargement contains the full
read of the plaquette, in the top part of
its own fourfold tower. The clearance
between the trial box and this read leaves
that density datum unchanged. The positive
test and the same averaging bound give
defect at most \(2\varepsilon_i\) there.
The lift on the class box uses only its
top links and complete backward blocks.
In a rectangular lift the top field in a
common rooted comb gauge has logarithms bounded
by \(N=CM\check R_i\varepsilon_i\), and the
induced lower row logs have the same bound.

The test lift is on the minimal recursive
tower of the component. Use the same
physical comb on the top field in both
runs. The block lift is the identity on internal
bonds and copies a top link on each bond crossing
its corresponding block face. Each one-step averaging
path is either internal, with transport one, or crosses
that face exactly once, with transport equal to the
copied top link. Its contour transports are one.
Thus all paths in the mean give the same value, and
the one-step average of the finer lift is exactly the
coarser lift. Hence the literal
prescribed datum agrees on complete
common rows, and the reduced datum agrees
where its outward read also agrees.
The physical recursive tower has the
two-cell inset at every level above its
first lower layer. The two cutoffs
therefore have the same upper regions
on this aligned box. The reduced outward
read uses only the adjacent lower block,
so the unmatched stock stays in a fixed
number of lower layers after enlarging
it by that read.
At the trivial base, \cref{dnr:flat-identification}
identifies the Landau potential of the test lift
with the response to these actual reduced data.
Use the two-depth response estimate on
the common top layer, with the error weight
screening the unmatched lower stock.
On the full common top stencils it
gives a difference bounded by
\(CN\eta^\alpha\) between the transferred
fine potential and the coarse potential,
including one difference. The individual
first differences have a H\"older bound
of size \(CN\) there.

To compare the literal plaquettes, in
this chart expand the four links in the
mesh. The normalized curvature has the
linear curl and the quadratic commutator
of the potential, with an error
\(C\eta(N+N^2+N^3)\). At the trivial
base the per-length linear transfer is
the straight path mean: the coarse component is
the average of the transported fine component along the
\(L\)-step straight paths, with total weight one.
A coarse difference of this mean is an
average of fine differences, by
telescoping the \(L\) fine translations.
The difference from any of the nearby
fine first differences costs
\(CN\eta^{\beta_H}\) on the common
stencil. The same pairing compares the
potential values in the commutator.
Each tested plaquette has a same-orientation
partner within \(C\eta\) by rounding
inward in the common tested box. The
clearance to the tower boundary puts
these partners on complete top stencils.
Taking the maxima in both directions
therefore gives
\[
 |q_\sharp(V)-q_\flat(V)|
       \le CM\check R_i\varepsilon_i\,
                         \eta^{\min(\alpha,\beta_H)} .
\]
Here the box, the exterior data and
the top trial are the matched ones.
The unmatched bottom rows use their
one-run bounds and screening.
In the class range \(\check R_i\le i\),
so in the flat level units
\(\eta=L^{-i}\) this gives a relative
error \(C L^{-c i}\) for some \(c>0\).
If the exterior links instead differ
by at most \(b_{\rm e}\) in a common
top frame, use the same free trial
links and substitute the other exterior.
The lift copies these links or the
identity. Its induced row logs differ
by at most \(Cb_{\rm e}\), provided
both data stay in the small chart.
At a single cutoff the screened
derivative of \cref{dnr:flat-identification} gives
an additional \(Cb_{\rm e}\) in the
normalized curvature on the tested
top stencils. This term uses the
actual exterior matching estimate.

Suppose the smaller infimum is at most
\(4\varepsilon_i\), and the comparison of the
two bare lifts on these common trials gives
\(|q_\sharp(V)-q_\flat(V)|\le d\), with
\(d<\varepsilon_i/2\). Use a trial arbitrarily
close to the smaller infimum. Its value bounds
the other infimum above by the first plus \(d\).
Trials approaching that other infimum also lie
below \(5\varepsilon_i\), so reversing the
argument gives
\(|q_\sharp^*-q_\flat^*|\le d\).
This reduces the pointwise transfer to the
uniform local comparison on the near-threshold
trial class.

The infimum also has a radial band estimate
on a free exterior logarithmic chart. To see
this on the aligned class box, fix a rooted
tree of the exterior graph. The exterior of
the inner rectangle has paths of length
\(O(D)\) to the root along the outer faces,
where \(D=CM\check R_i\) is the box side.
Choose the tree with those path lengths.
Include the outer crossing links of the test
read as leaves. At an aligned rectangular
face each outside site on such a link has just
one neighbor inside the box.
Extend it into the inner rectangle by axial
columns from a face. The added tree links
are free trial links. In the extended tree
gauge all paths to the root still have
length \(O(D)\), and this gauge keeps the
prescribed exterior in its original tree
chart.

A near-minimizing trial in the threshold band
has top plaquette defects at most
\(C\varepsilon_i\), by the preceding trial
and density bounds. Its ordinary full-box
axial gauge has links within \(CD\varepsilon_i\)
of one. Each link in the fixed extended
tree gauge is a loop of length \(O(D)\)
in that field. Thus the exterior and free
trial logs in the fixed chart are bounded
by \(r=CD^2\varepsilon_i\).

Write \(\lambda=2\varepsilon_i\).
For a combined vector \(B\) of exterior
and trial logs with \(\sup|B|\le r\),
the vacuum continuation of the test lift
is holomorphic on the fixed datum ball.
For each normalized matrix plaquette read
\(F_p\), its Taylor series at the vacuum
gives, uniformly for \(0\le u\le1\),
\[
    |F_p(uB)-uF_p(B)|
       \le C u(1-u)r^2 .
\]
Indeed the constant term is zero and
the linear term cancels. Each term of
degree \(m\ge2\) has factor
\(|u^m-u|\le (m-1)u(1-u)\), and Cauchy's
bound sums the resulting series.

On an exterior ray let \(t_1<t_2\) both
lie in the band of relative width
\(\vartheta\le1/4\). Take a trial
arbitrarily close to the infimum at
\(t_2\) and multiply every log, exterior
and interior, by \(u=t_1/t_2\). The tree
links remain fixed and the scaled trial
lies in the admitted class at \(t_1\).
The previous inequality and passage to
the infimum give
\[
 q^*(t_1)\le u q^*(t_2)+Cu(1-u)r^2 .
\]
The condition \(Cr^2\le\lambda/4\) holds at
weak coupling: \(D\) is bounded by a fixed power of
\(\log g_i^{-2}\), while
\(\varepsilon_i=g_i p_0(g_i)\) and \(p_0\) is a fixed
power of that logarithm. Hence
\(r^2/\lambda\le C D^4\varepsilon_i\to0\) as \(g_i\to0\).
The band bounds then give \(1-u\le C\vartheta\)
and \(u\ge c>0\).
The band's logarithmic radial diameter is
therefore \(O(\vartheta)\). The Haar
radial estimate gives a relative bound
\(C n_{\rm e}\vartheta\) on a product ball
of \(n_{\rm e}\) exterior non-tree links.

At the native extendibility stop the top
exterior is fixed in the infimum defining \(q^*\).
The product bound enters a
history integral only where those exterior
links form a free stock in a containing
operation and the uniform reference
bound is taken with that stock fixed.
A prescribed outgoing row can have its
whole double block in the exterior.
In that case its constraint belongs to
the same stock, and the weighted shell
requires the bound on the constrained
coordinates.

%
%
There is an integrated exterior stock when this
extendibility stop is immediately below the floor
\(h=i+1\) of a containing operation that reaches
its shell or last-shell stop. Take an old
component whose tested box and stencil pad are
disjoint from the other marked charts. Include
both its tested exterior and its previously
generated targets in \(Z''_i\) when listing the
ordinary \(h\)-rows that read it. With the full
row carriers in the eight-cell raw buffer, every
target on this list belongs to the integrated
inner field in (1.145). The old prepared targets
in \(Z''_i\) and the old exterior links retain
their separate factors in the joint integral.
Integrate the listed \(h\)-targets against their
own Haar deltas, substituting their row values
in every remaining factor. The other \(h\)-rows
do not read the listed old stock. In the absence
of a later prescribed target on this list, this
leaves product Haar on the old exterior, with the
old inner graph evaluated on its own targets.

Apply the pointwise reference majorant of the
containing operation before taking the exterior
supremum. On its admitted chart the majorant is
uniform in the free old group links after the
other reference factors have been bounded.
The exterior tree change preserves their Haar
factors. The band depends only on these stored
exterior coordinates at level \(i\). It therefore
costs \(C n_{\rm e}\vartheta\) relative to the
product-ball majorant for this chart. With
identical top exteriors the relative tolerance
above is \(O(D\eta^{\min(\alpha,\beta_H)})\).
Since \(n_{\rm e}\le CD^4\) and \(\check R_i\le i\),
the resulting \(CD^5\eta^{\min(\alpha,\beta_H)}\)
is still geometric in \(i\).

%
%
In the same immediate-floor geometry, take
an old second-class component of level \(i\)
with a final density test. Use the complete
fourfold read and its row carriers in the
raw buffer and include its old exterior and
restored links in the list of reads. If all
affected \(h\)-targets are integrated, the
same Haar delta elimination leaves product
Haar on the old exterior of the test. A read
wholly outside the old proposal box uses
only this stock. The negative test of the
old second-class component belongs to the
containing core's payment. Keep its actual
failed cube for that payment before bounding
the remaining reference factors.

For a read crossing the old proposal
boundary, use the physical off-tree links
in the recorded frame. The old reference
centre is determined by the old exterior
stock. At fixed old exterior the relative
non-tree Haar factors give the physical
ones by right translation. After bounding
the reference density, restore the interior
site gauges for this invariant band event,
fixing the exterior sites. This replaces
the interior tree constraint by normalized
Haar integration. Together with the old
exterior factors, Fubini gives a product
Haar majorant on the complete physical
test read. The complete top curvature
bound and the radial estimate at
\(\lambda=\varepsilon_i\) give a factor
\(C n\vartheta\) relative to its product
majorant, with \(n\) the free-link count.
The reference normalization is bounded
uniformly on the admitted slice; use that
bound once if it has not already been
included in the pointwise envelope.
The band here is a function of the
recorded physical links and can be
imposed before the compatible-exterior
supremum.

%
%
Some of these estimates also hold with prescribed
\(h\)-targets. Consider rows crossing a face of
the old trial box \(\Lambda^{\rm prop}\).
Its faces are aligned with the \(h\)-blocks.
The middle link of such a row crosses the face
and is among the free trial links in (1.142).
It is thus outside the exterior read of the
infimum. The trial box includes a full \(M\)-layer
beyond \(Z''_i\), so this link is also separate
from the old prepared targets. It is an original
Haar link in the ordinary raw average.
Solve the prescribed row in this link, keeping
its two endpoint blocks and all weights that read
the solved link in the joint integral. Distinct
rows have private middle links. The row kernel
has mass at most \(2L^9+1\) per prescribed
target, uniformly with the exterior stock fixed.
Each such row still reads the old exterior in
its outside endpoint block.

There is another choice of link for rows that
cross the outer face of the old rectangular
component. The middle link can now be in the
test read, since the lift uses links meeting
that component. Orient the row outward and
choose an axial link on the part of the straight
path strictly in the outside endpoint block.
Take it beyond the test pad and before or at
the outside coarse centre. The choice is
available for large \(L\) and for the complete
face carriers in the raw buffer, disjoint from
other marked charts. This link also has original
Haar measure.

The carried positive density cubes in the outer
collar supply the regularity for this second
choice. At the face their tested enlargements
cover the two endpoint blocks, including the
full averaging reads of the top plaquettes there.
Thus the top defects on the double block are
at most \(C\varepsilon_i\), also on a failing
extendibility branch. Choose a rooted tree of
the double block avoiding the chosen link.
The small-plaquette estimate puts every admitted
field into a log ball of radius \(C_L\varepsilon_i\)
in this gauge, which depends only on the other
links. At the vacuum the derivative of the
logarithm of the row in the chosen link is
\(m\mathbf 1\), with \(m\ge L^{-3}\).
Indeed, the \(L\) starts on the straight axis in
the contour formula each contribute the link
once in the positive direction. The off-axis
starts can read it only on the reversed contour
in the outside block, also in the positive
direction. Regularity of (1.28) gives derivative
distance at most \(m/2\) from \(m\mathbf 1\) on
a small convex log ball. There is at most one
inverse there, with Haar density bounded by
\(C_L\).

An outside endpoint block of a face-crossing
row has only one neighboring block in the
aligned rectangle. Distinct such rows therefore
use different outside blocks, and none reads
the chosen link of another. First use the
private middle-link kernels for the inner-face
rows, conditional on all the outer-face links.
The outer-face rows do not read those middle
links. Then use the outer-face inverses, with
all retained factors evaluated at the solved
links. The resulting bound has the same
\(C n_{\rm e}\vartheta\) from the exterior
stock, with the product of the row constants
in the reference bound. Those constants are
charged to the number of prescribed rows.
Other affected targets must be integrated or
have their own inverse estimate with the test
stock held fixed. In particular a row wholly
in the exterior can constrain the quotient
links. Its site-gauge variables give another
case in which the band estimate applies.

For the final density test there is also
an outside-pivot construction relative to
its complete fourfold read \(S\). Take a
rectangle \(S\) aligned with the next
averaging blocks. A prescribed row crossing
its face has a straight-axis link strictly
in the outside endpoint block, beyond \(S\)
and before the outer centre. Use that link
as its pivot when it remains in the raw
Haar stock. The six-cell inset leaves the
complete double block inside the component.
The physical field on the final reference
support has top plaquette defects at most
\(C(M)\varepsilon_i\) there. The tree
avoiding this link therefore gives the same
small regular inverse ball as in the
outer-face argument, with constants depending
on \(L,M\). Different face-crossing rows use
different outside blocks. Solving them
consumes no link of the test read.
Other rows reading these pivots must have
their targets integrated or be solved
in a compatible inverse order. A row
whose whole double block lies inside
\(S\) has no outside pivot of this kind.

%
In the rooted exterior tree coordinates write
the actual exterior field as
\(V(b)=g_{b_-}^{-1}\widehat V(b)g_{b_+}\),
where \(\widehat V\) is the tree-gauged field.
The change separates product Haar on the
non-tree links and on the nonroot site gauges.
The extendibility infimum depends only on
\(\widehat V\), since a site gauge extends
to the trial fields over the inner box.
Suppose the remaining prescribed rows have
their full double blocks in the exterior stock
and form a forest of coarse bonds. Choose a
root for each tree, using the fixed root of the
exterior gauge if the tree contains it. Take
the nonroot coarse sites to have their full
fine link stars in the free stock. By covariance
each prescribed value is
\[
 W(c)=g_{c_-}^{-1}\mathcal M(\widehat V)(c)g_{c_+}.
\]
For fixed quotient links and fixed root gauges
solve the rows in the child gauges, starting
at the roots. The conditional Haar factor is
one for each solved row. Substitute the gauges
in all remaining reads before applying the
reference bound. This keeps the full quotient
Haar integral for the band. The previous
face-row inverses, if needed, can be taken
first with these gauges fixed, since the
forest rows read none of their pivots.
This forest calculation also applies to
the final density read on a physical
stock with these free site gauges. A
reference tree delta still present before
classification must be included in the
joint gauge calculation with the prescribed
rows. Its removal after a successful class
test does not supply the gauge variables
for the preceding decision. On an already
ungauged stock the plaquette norms are
gauge invariant, so at fixed quotient
links the values prescribed on a forest
change the solved site gauges and leave
the test value unchanged. The children
must have their complete link stars in
the free stock.
If the prescribed graph has cycles, solving
a spanning forest leaves a constraint on the
ordered product of the quotient row values
around each fundamental cycle. That product
must equal the product of the prescribed
values conjugated by the root gauge. The
remaining estimate is on these constrained
quotient coordinates.

%
%
Some of the remaining cycle constraints can
be treated together. Use the same forest and
give each nonforest row its private fine
middle link. The exterior tree can avoid all
these links. A transverse plaquette detour
replaces each forbidden link in an exterior
path by a nonmiddle crossing and two
within-block links. The root paths still
have length \(O(D)\). Thus the near-minimizing
trial estimate in this tree again gives
the log radius \(r=CD^2\varepsilon_i\).

Let \(x\) denote the non-tree exterior logs
apart from the selected pivots. For a coarse
vertex \(v\), let \(A_v(x)\) be the ordered
product of the quotient row values from its
forest root to \(v\), and \(T_v\) the product
of the prescribed values on that path.
If the root gauge is \(g_0\), the solved
gauge at \(v\) is \(A_v(x)^{-1}g_0T_v\).
The remaining equation for a row \(c\) from
\(v\) to \(w\) is therefore
\[
 \mathcal M(\widehat V)(c)
      =A_v(x)^{-1} Z_c A_w(x),\qquad
 Z_c=g_0 T_v W(c)T_w^{-1}g_0^{-1}.
\]
The forest paths read none of the selected
pivots. Each remaining row reads only
its own pivot from this stock. Thus the
inverse in all these links is diagonal even
when the forest paths overlap. At the
vacuum the single-row pivot derivative is
\(L^{-3}\mathbf1\) in the oriented log
coordinates by (1.31). The regular one-row
chart gives a unique inverse \(p_c=P_c(x)\)
with conditional Haar density bounded by
\(C_L\) per row on the near-band chart.
In taking the radial integral keep the
prescribed values and root gauges fixed.

Write \(\ell_*\) for one plus the greatest
sum of the two forest path lengths used
by any such row. The geometry gives
\(\ell_*\le CD^4\). On the near-band
chart \(|Z_c-\mathbf1|\le C_L\ell_*r\).
On a complex radial disc of radius
\(c_L/(\ell_*r)\), the forest products
stay in a fixed small ball by the
product estimate. The target of each
one-row inverse stays in its regular
ball at weak coupling. Cauchy's estimate
on this disc gives, for \(0\le u\le1\),
\[
 |P_c(ux)-uP_c(x)|
 \le C_L(1-u)\bigl(|Z_c-\mathbf1|+\ell_*^2r^2\bigr).
\]
Indeed \(P_c(0)\) has size \(C_L|Z_c-\mathbf1|\).
The linear part cancels, and the terms of
degree at least two cost
\(C_L(1-u)\ell_*^2r^2\).

To compare two radii on this graph, scale
a near-minimizing trial at the larger radius,
then correct all its selected exterior pivots
to their values at the smaller radius.
The corrections are to the non-tree links
in the fixed chart. Their lifted data and
their reduced data use the full local carriers
of the corrected links. For pivots whose
target products satisfy
\(|Z_c-\mathbf1|\le c_L\varepsilon_i\)
use the unweighted normalized response
bound. For the others suppose the complete
carriers are in the middle of the two-cell
outer density collar, a distance
\(d_0\ge cM\check R_i\) from all tested
stencils. In \eqref{pnc:weighted-first-response}
use the weight given by distance to the
union of their carriers. This gives a
normalized curvature response bounded by
\(Ce^{-c d_0}\) times the supremum of their
corrections.
Including the curvature scaling error of
the trial itself, the extra term in the
radial inequality is bounded by
\[
 C_L(1-u)\bigl[
       c_L\varepsilon_i+\ell_*^2r^2
       +e^{-c d_0}(\ell_*r+\ell_*^2r^2)\bigr].
\]
The terms containing \(r^2\) are
\(o(\varepsilon_i)\) by the logarithmic
class size. The screened term linear in
\(r\) is also \(o(\varepsilon_i)\), since
\(\ell_*D^2e^{-c d_0}\) tends to zero.
Choose \(c_L\) small so the total coefficient
of \(1-u\) is below \(\lambda/4\).
The earlier logarithmic radial diameter
bound applies to \(x\). The product of the
one-row inverse bounds is charged to the
number of prescribed rows, and the band
still contributes \(C n_{\rm e}\vartheta\).
For a nonscreened pivot at a general target
product one must control the curvature
produced by its offset on the constrained
exterior chart.

For a final test it suffices to screen
the pivots from the actual tested
stencils. With \(r=CD^2\varepsilon_i\),
the required bound on the linear offset
is
\[
 C_L e^{-c d_0}\ell_*D^2\le c_0 .
\]
Since \(\ell_*\le CD^4\), a sufficiently
large multiple of \(\log(2+D)\) for
\(d_0\) suffices. The outer part of the
fourfold test box can provide this
clearance to its tested first enlargement.
Use the complete lifted pivot carriers
in measuring that distance.

There is a sufficient rule for choosing
the forest in the screened case. Declare
a prescribed edge near when the complete
lifted carrier of its available private
pivot is within \(d_0\) of the tested
stencils. If those edges form a forest,
extend it to a spanning forest of the
prescribed graph. Every remaining chord
then has a far pivot, so its cycle-product
offset enters only through the screened
carriers above. The forest rows are solved
in the site gauges. The criterion uses
the whole pivot carrier, including its
double block and induced outward rows.
A cycle among the near edges leaves a
near chord under every choice of forest,
and its invariant target product remains
in the centre comparison.

%
%
%
To compare the two cutoffs on this chart,
take the same prescribed graph, forest and
exterior tree at the common top level.
Use the same remaining quotient links in
the two reference integrals. The row maps
at levels \(i\) and \(h=i+1\) have the same
double blocks at matched depths. The
forest products \(A_v(x)\) therefore agree.
Different prescribed forest values change
only the solved site gauges. The test
infimum is invariant under those gauges.
For the remaining rows the only prescribed
values entering the quotient graph are the
root-conjugated cycle products \(Z_c\).
If these agree, so do the solved quotient
links. This also treats prescribed forests,
for which there are no cycle products.
If the two prescribed fields are related
by coarse site gauges \(u_v\), their path
products satisfy
\(T_v^\flat=u_0T_v^\sharp u_v^{-1}\).
Choose the root of the exterior tree at
a fine site outside the prescribed graph.
Pair the free forest root gauges by
\(g_0^\flat=g_0^\sharp u_0^{-1}\).
This preserves their Haar measures and
gives equal \(Z_c\) on all remaining rows.
For the test read alone the two prescribed
forests may even differ, provided their
site gauges are free on the same exterior
stock. Each forest removes only its own
gauge coordinates, so the test uses the
same quotient field in the two integrals.
The same observation applies to additional
branches attached to a common cyclic
part at one vertex each. Such branches
leave its cycle products unchanged.
Use the reference bound of each graph
for its other factors.

For unequal products compare the two
one-row inverses at fixed \(x\).
The targets differ by
\(A_v(x)^{-1}(Z_c^\sharp-Z_c^\flat)A_w(x)\).
On the common regular ball the difference
of the solved log links is at most
\(C_L|Z_c^\sharp-Z_c^\flat|\). At fixed
root gauges the difference of the products
is bounded by the sum of the prescribed
row differences along the two forest
paths and the remaining row. This follows
by telescoping products of real unitary
matrices in a common root frame.
The regular Haar inverse density for
such a row is the same function of \(x\)
and \(Z_c\) at both cutoffs. Its logarithm
is Lipschitz in the solved log link on the
fixed regular ball, by the determinant
formula and the local inverse bound.
Thus the absolute log ratio of these row
factors is at most
\(C_L\sum_c|Z_c^\sharp-Z_c^\flat|\).
In particular these factors agree for
equal products.

Write \(b_{\rm near}\) for the maximum
product difference on the unscreened rows
and \(b_{\rm far}\) for the maximum on the
screened rows. The other exterior links
agree on the quotient graph. Transfer a
near-minimizing trial from the passing side
by keeping its free interior links and
replacing the exterior pivots by those of
the other graph. First compare the two
lifts at the original trial, then replace
the pivots in the other lift. The pointwise
comparison at the original trial and the
unweighted and weighted response bounds
on the two sets of pivot carriers give
an increase of at most
\[
 d=C D\varepsilon_i\,\eta^{\min(\alpha,\beta_H)}
     +C_L\bigl(b_{\rm near}+e^{-c d_0} b_{\rm far}\bigr).
\]
Here the changed data must stay in the
common regular ball. If one test passes
and the other fails, the failing infimum
is at most the passing infimum plus \(d\).
Consequently both infima are within \(d\)
of \(\lambda\). It suffices to use the
band of the reference graph on the
remaining quotient stock. When \(d/\lambda
\le1/4\) and the radial bound applies to
this graph, the factor is
\(C n_{\rm e}d/\lambda\) in the common
reference majorant. The inverse row
densities remain in that majorant. A
geometric bound in the bare mesh uses
a mesh-small product discrepancy in
the normalized width \(d/\lambda\).
The reference graph's prescribed cycle
products must be fixed at this stage,
as in the fixed-output convention for the
compatible evaluation supremum. The other
graph can vary there if the comparison by
\(d\) is uniform. The reference band is
then independent of that evaluation.
Integrate it at fixed reference root
gauges, and then integrate those gauges
against Haar.

Replacing the sharp tests by Lipschitz functions
allows pointwise
telescoping, with a separate estimate for the modified
one-run operation. The marked remainders and hence the
tuned coordinates in that argument belong to the modified
recursion.

\subsubsection*{Calibration maps and their admitted cocycle}

For the same bare integral let $C_a(\beta)$ and
$C_b(\beta)$ be two actual terminal-coupling
coordinates, on domains where the branches are
holomorphic and invertible.  Define
\begin{equation}
 H_{ab}=C_b\circ C_a^{-1}.
 \label{eq:weighted-matched-domains-WM22}
\end{equation}
On an admitted common domain with the same bare
branch, one has
\begin{equation}
 H_{bc}\circ H_{ab}=H_{ac},\qquad
 H_{ab}'(x)={C_b'(\beta)\over C_a'(\beta)},
            \qquad \beta=C_a^{-1}(x).
 \label{eq:weighted-matched-domains-WM23}
\end{equation}
These follow by composition and the chain rule.
Every intermediate inverse must be the admitted
branch.  Geometric overlap of two complex
discs does not establish that fact.

Moreover
$|H_{ab}(x)-x|=|C_b(\beta)-C_a(\beta)|$.
This is the actual calibration discrepancy to
estimate.  Equality after a change of physical
units does not bound it.  If the two
presentations are known to represent the same
bare integral on an admitted real open overlap,
their holomorphic identities extend from that
overlap by the identity theorem.  The real
admission is an input, as in the flow section.

\subsubsection*{Typed closure is a family of finite-grade estimates}

A nonscalar insertion must come with its own
gauge-covariant source core, counter-insertions,
and tensor readout.  At finite grade, label
every source monomial and all cross terms
created by the actual finite map.  If several
directions share one potential or zone reserve,
their joint radii obey \eqref{eq:weighted-matched-domains-WM15}, with the
multi-direction cost following it.  A scalar
$\alpha^2$ gain is not copied into another
channel without its own field-degree and
counter-insertion calculation.

Parity proves \eqref{eq:weighted-matched-domains-WM16} only for the actually
admitted unmarked family.  Its one-loop
coefficient in \eqref{eq:weighted-matched-domains-WM17}, the marked corrections,
and the finite-grade closure still have to
be computed for that channel.  For the
optional anisotropy Ward route an estimate
of size $K^{-c_*}(\log K)^A$ multiplied by
$g_0^{-1}\asymp K^{1/2}$ tends to zero
from this bound if $c_*>1/2$.  At equality
the logarithmic factor must itself tend to
zero; a mere bounded-power estimate does
not settle it.  A fixed-torus estimate
does not give a growing-volume Ward limit
unless its volume factors are controlled
at that rate.

The outstanding weighted task is therefore
the actual finite-map extension with the
weights in \eqref{eq:weighted-matched-domains-WM3}, not just ordinary source
holomorphy.  The outstanding zone task is
the complete incidence and original-weight
reserve in \eqref{eq:weighted-matched-domains-WM13}--\eqref{eq:weighted-matched-domains-WM15}.  The matched task
is the actual common-domain generator
estimate in \eqref{eq:weighted-matched-domains-WM20}--\eqref{eq:weighted-matched-domains-WM21}, with its
summable or direct two-cutoff error.
Typed closure requires these assertions
at the stated finite grade.  These are
separate optional routes; the estimates they require
are not additional hypotheses
of the ordinary scalar angular argument.


% END MODULE chapters/weighted_matched_domains.tex
% BEGIN MODULE chapters/notation_index_v2.tex
\section{Index of notation and conventions}
\label{ch:notation-index}

The table gives the principal notation and links to its definitions.

{\small
\setlength{\tabcolsep}{4pt}
\renewcommand{\arraystretch}{1.08}
\begin{longtable}{@{}>{\raggedright\arraybackslash}p{.22\linewidth}
>{\raggedright\arraybackslash}p{.47\linewidth}
>{\raggedright\arraybackslash}p{.25\linewidth}@{}}
\hline
\textbf{Notation}&\textbf{Meaning}&\textbf{Definition}\\
\hline
\endfirsthead
\hline
\textbf{Notation}&\textbf{Meaning}&\textbf{Definition}\\
\hline
\endhead
\hline
\endfoot

$\mu_e,q(e),\nu_G$&Zero-source conditional law, clean conditional probability, and exterior-preserving clean resampling.&\eqref{phc:disintegration}; \eqref{phc:clean-probability}; \cref{prop:phc-resampling}\\[4pt]
$d_S,b_n(D)$&Weighted exceptional mixed moment and its connected-error polynomial.&\eqref{phc:weighted-defect}; \eqref{phc:stability-polynomial}\\[4pt]
$A_B$&Mixed source derivative of the complete effective action.&\eqref{esj:source-jet}\\[4pt]
$\alpha,L(A)$&Complete residual oscillation and Gaussian quadratic coefficient size.&\eqref{fwg:oscillation}; \eqref{fwg:quadratic-size}\\[4pt]
$\mathsf D_\eta,\mathcal B_{\omega,\chi}$&Lattice displacement vector field and its exact shell flux.&\eqref{wd:displacement}; \eqref{wd:flux-definition}\\[4pt]

\multicolumn{3}{@{}l}{\textit{Lattices, fields, and action}}\\[3pt]
$G,\mathfrak g,N,d_G$&$G=\operatorname{SU}(2)$, $N=2$, and $d_G=N^2-1=3$; Hermitian traceless matrices with $\langle X,Y\rangle=\tfrac12\operatorname{tr}(XY)$.&\hyperref[notdef:gauge-convention]{gauge convention}; \cref{ch:free-curvature-exact}\\[4pt]
$L,M$&Odd blocking integer $L>1$ and localization cube side $M$. The rooted-source construction takes $M=L^{m_0}$ and compatible tori $\ell=2L^m$.&\hyperref[notdef:localization-cubes]{localization cubes}; \eqref{eq:graded-source-bound-GS1}\\[4pt]
$K,j,s,a_s,\eta_j$&Cutoff depth $K$, level $j$, remaining depth $s=K-j$, physical spacing $a_s=L^{-s}$, and bare spacing $\eta_j=L^{-j}$ in level-$j$ units.&\hyperref[notdef:level-depth]{level and depth}; \hyperref[notdef:gauge-convention]{gauge convention}\\[4pt]
$\Lambda_{K,\ell},\mathbb T^{(j)}$&Finite torus lattice $(L^{-K}\mathbb Z/\ell\mathbb Z)^4$ and the lattice retained after $j$ blockings.&\hyperref[notdef:lattice-units]{lattice units}; \hyperref[notdef:gauge-convention]{gauge convention}\\[4pt]
$U_b,U_p,\partial U$&Bond variable, ordered plaquette transport, and its alternate notation. Reversal gives $U_{\bar b}=U_b^{-1}$.&\hyperref[notdef:bond-fields]{bond fields}; \hyperref[notdef:gauge-convention]{gauge convention}\\[4pt]
$a_p,A(w,U),\beta$&Plaquette action $1-\tfrac12\operatorname{Re}\operatorname{tr}U_p$, weighted action $\sum_pw(x(p))a_p$, and bare inverse coupling $\beta=g_0^{-2}$.&\hyperref[notdef:weighted-action]{action and sources}; \hyperref[notdef:first-passage]{first passage}\\[4pt]
$\mathcal O(f),S_n^T$&Smeared action density $\sum_pf(x(p))a_p$ and its joint cumulants. No spacing prefactor occurs in $\mathcal O$.&\hyperref[notdef:observables]{observables}; \hyperref[notdef:weighted-action]{action and sources}\\[4pt]
\multicolumn{3}{@{}l}{\textit{Averages and propagators}}\\[3pt]
$\mathcal M,Q_k,Q_k^\sharp$&Contour block average, linearization of its plain iterate, and straightened linearization after removal of the coarse gauge factor.&\eqref{eq:rg-objects-R3}; \hyperref[notdef:covariant-averages]{iterated averages}\\[4pt]
$Q_k^{\mathrm{st}},Q_k^{\prime}$&Straight covariant bond average and site-function average.&\hyperref[notdef:covariant-averages]{iterated averages}\\[4pt]
$\mathrm{d}_0,\mathrm{d}_1,\mathrm{d}_\eta$&Unscaled scalar difference, plaquette coboundary, and physical difference; in the scalar block convention $\mathrm{d}_0=\eta \mathrm{d}_\eta$.&\hyperref[notdef:block-inner-products]{block inner products}; \hyperref[notdef:curvature-kernels]{curvature kernels}\\[4pt]
$Q,P,\mathcal A,G^{\prime}$&Scalar arithmetic mean, block-constant projection $P=Q^*Q$, operator $\mathcal A=\mathrm{d}_\eta^*\mathrm{d}_\eta+P$, and inverse $G^{\prime}=\mathcal A^{-1}$.&\hyperref[notdef:block-mean]{block mean}; \hyperref[notdef:scalar-gap]{scalar operator}; \hyperref[notdef:scalar-projection]{scalar projection}\\[4pt]
$R,\mathbf R,R^{(j)}$&Scalar orthogonal projection onto $\ker(QG^{\prime})$, large-field operation, and local remainder family, respectively.&\hyperref[notdef:scalar-projection]{scalar projection}; \cref{candidate:large_field}; \hyperref[notdef:activity-families]{activity families}\\[4pt]
$P_{\mathrm{ex}},K,\mathsf K$&Projection onto exact lattice two-forms and continuum curvature covariance at unit coupling. The physical lattice kernel is $a^{-4}P_{\mathrm{ex}}(p,q)$.&\hyperref[notdef:curvature-kernels]{curvature kernels}; \cref{ch:free-curvature-exact}\\[4pt]
$U_j(V),\mathsf G,\mathsf H$&Constrained regular background, constrained inverse, and response operator, with the gauges and domains specified in the background hypotheses.&\hyperref[notdef:background-orbit]{background orbit}; \hyperref[notdef:propagators]{background operators}\\[4pt]
$\Phi_\Lambda(B),r_B$\newline$Z_y$&Potential/current background on the datum ball $\sup_c|B(c)|<r_B$ and local output $Z_y=(\xi_yH,\xi_y^2\nabla H,\xi_y^3J)$.&\eqref{eq:pointed-window-remainder-PW1}; \eqref{eq:pointed-window-remainder-PW2}\\[4pt]
\multicolumn{3}{@{}l}{\textit{Couplings and sources}}\\[3pt]
$x_j,g_j,b_j$&Running inverse coupling $x_j=x_{j-1}-b_j$, $x_0=\beta$, and coupling $g_j=x_j^{-1/2}$.&\hyperref[notdef:flow-recursion]{flow recursion}; \hyperref[notdef:frames]{frames and tests}\\[4pt]
$c_0,b_j^0$&Leading drift $c_0$ and designated one-loop increment $b_j^0$. The numerical flow theorem assumes an interval-sum law for these increments.&\hyperref[notdef:flow-recursion]{flow recursion}; \eqref{eq:comparison-capsule-CP2}\\[4pt]
$\theta_j,\bar g_j,x_j^\natural$&Auxiliary frame, its coupling $\bar g_j=\theta_j^{-1/2}$, and inverse coupling computed with auxiliary choices frozen at $\theta_j=g^{-2}+c_0(K-j)$.&\hyperref[notdef:frames]{frames and tests}; \hyperref[notdef:frozen-coupling]{frozen coupling}\\[4pt]
$K(\beta,g)$\newline$\mathcal W_K(g)$&First entry into $g^{-2}<x_K\le g^{-2}+B_\sharp$ and the set of bare couplings with this stopping index.&\hyperref[notdef:first-passage]{first passage}\\[4pt]
$w,c,z,f_i$&Source triple with function $w=c+\sum_i z_if_i$ on each fixed finite real test family.&\hyperref[notdef:source-triples]{source triple}\\[4pt]
$\|\cdot\|_{\sharp,s},\|\cdot\|_{\flat,s}$\newline$\mathcal W_r^{(s)}$&Test norm $\max(\|f\|_\infty,c_ML^{-s}\operatorname{Lip}f)$, source-triple norm $|c|+\sum_i|z_i|\|f_i\|_{\sharp,s}$, and its open radius-$r$ ball.&\hyperref[notdef:source-norms]{source norms}\\[4pt]
$h_y^{(j)},\varphi_j,\zeta_j$&Anchor partition, fixed coupling profile, and running source shift defined by the exact counterterm identity.&\hyperref[notdef:source-profiles]{source profiles}\\[4pt]
$\varrho,\rho$\newline$\alpha_{0,j},\alpha_{1,j}$&Source or coupling-disc radii in the stated estimate and complex field radii. Their meanings are local to that estimate.&\hyperref[notdef:source-triples]{source triple}; \hyperref[notdef:frames]{frames and tests}; \eqref{eq:rg-objects-R4}\\[4pt]
$p_0,\check R_j$\newline$\varepsilon_j,\delta_j$&Deposit scale, localization/waiting scale, and curvature and fluctuation thresholds formed from the fixed auxiliary frame.&\hyperref[notdef:frames]{frames and tests}; \eqref{eq:rg-objects-R19}\\[4pt]
\multicolumn{3}{@{}l}{\textit{Local families and histories}}\\[3pt]
$E^{(j)},R^{(j)}$\newline$B^{(j)}$&Regular activities, remainder activities, and boundary activities; pointed families also carry a base point.&\hyperref[notdef:activity-families]{activity families}; \eqref{eq:rg-objects-R13}\\[4pt]
$\mathsf b(E),e_j,\mathcal E$&Scalar marginal extractor, local source energy, and its sum. The vacuum constant $E$ has a different role.&\eqref{ph:extractor}; \hyperref[notdef:local-source-energy]{local source energies}\\[4pt]
$d_j(X),\kappa,\kappa_{\mathcal D}$&Localization tree length, decay exponent, and domain-counting exponent. Summability requires decay greater than the counting exponent.&\hyperref[notdef:tree-length]{tree length}; \eqref{eq:graded-source-bound-GS1}\\[4pt]
$\Omega_j,\Lambda_j,Z_j$&Regular region, buffered inner region, and complement $Z_j=\Lambda_j^c$ in the specified history geometry.&\hyperref[notdef:history-regions]{history regions}\\[4pt]
$h,\omega$&Full history and local history inside a stored operation.&\hyperref[notdef:histories]{histories}; \hyperref[notdef:standing-history]{stored operations}\\[4pt]
$\Xi_\omega,\mu_\omega$\newline$\Phi_\omega,\Psi_\omega$&Stored integration space, positive reference measure, exponent, and excess source exponent. The integration domains and evaluation fibers are part of the construction.&\eqref{eq:rg-objects-R24}; \eqref{eq:flow-RF7c}; \hyperref[notdef:excess-source]{excess source exponent}\\[4pt]
$\Theta^{(k)},q_a,\Sigma_k$&History deposit/volume mark, age-survival rate $q_a=e^{-c_a}$, and integrated aged load.&\eqref{eq:rg-objects-R28}; \hyperref[notdef:age-survival]{survival rate}; \hyperref[notdef:aged-load]{aged load}\\[4pt]
$\mathfrak C_k,\mathcal U_j^c$&Class of complete presentations and complex pair domain. The current in a pair is an independent coordinate.&\hyperref[notdef:presentation-class]{presentations}; \eqref{eq:rg-objects-R4}\\[4pt]
$\mathcal G_z,\mathcal K_z,A_z$&Complete pointed total, its response-composed Hessian, and midpoint tent action for the same member.&\eqref{ph:pointed-total}; \eqref{ph:actual-kernel}; \eqref{ph:tent-action}\\[4pt]
$\rho,\sigma,a,a_{\rm out}$ (window)&Field-radius ratio $\rho=a_{\rm out}/a\ge1$, scale ratio $\sigma$, member radius, and output radius.&\eqref{eq:pointed-window-remainder-PW5}\\[4pt]
$\tau_S(w),N$ (history sum)&Complex label measures, positive at zero, and the site count in the ordinary stability criterion.&\eqref{oh:exactness}\\[4pt]
$E(z),E_{\ell,K},A_K$ (transfer)&Common cell normalizer, finite-torus vacuum constant, and companion vacuum density used in the finite-reference route.&\eqref{eq:two-point-extraction-partial-ET5}; \eqref{eq:two-point-extraction-partial-ET8}; \eqref{eq:two-point-extraction-partial-ET9}\\[4pt]
\multicolumn{3}{@{}l}{\textit{Limits and reconstruction}}\\[3pt]
$\mathfrak T_{\mathbb E}(\mathsf d;\vartheta)$\newline$\mathfrak T^{(4)}_{\mathbb E}(\mathsf d;\vartheta)$&Positive test pairs with quantitative mass and Lipschitz bounds, and quadruples whose pairs belong to this class, with small supports around collinear centres.&\hyperref[notdef:positive-test-pairs]{positive test pairs}; \hyperref[notdef:four-point-tests]{four-point tests}\\[4pt]

$\mathfrak S,\sigma$&Separated compact test tuples, including the empty tuple, and their moment functional, extended to finite linear combinations.&\hyperref[notdef:separated-tests]{separated tests}; \hyperref[notdef:moment-functional]{moment functional}\\[4pt]
$\mathsf d(X)$&Minimum distance between distinct test supports; support-containing balls do not define this distance.&\hyperref[notdef:separated-tests]{separated tests}\\[4pt]
$\sigma_S,\sigma^0$&Canonical centered functional: the partition sum over blocks of size at least two, with empty value one and singleton value zero.&\hyperref[notdef:centered-functional]{centered functional}; \eqref{cc:centered-functional}\\[4pt]
$\mathcal S_{\ne},|\varphi|_k$&Schwartz tests flat with all derivatives on every coincidence diagonal, and the weighted derivative seminorm, written $p_k(\varphi)$ in the extension proof.&\hyperref[notdef:diagonal-flat-space]{diagonal-flat space}; \eqref{eq:quantitative-flat-extension-QF2}\\[4pt]
$m(X),X^+$&Minimum cubical tree-edge count and the fixed-layer source-dependence enlargement in the rooted geometry.&\eqref{eq:graded-source-bound-GS1}; \eqref{eq:graded-source-bound-GS2}\\[4pt]
$k,B(R)$&Factorial exponent and support-radius factor in the graded all-state estimate; one may take $B(R)=C(R+1)^4$.&\cref{thm:graded-rooted-source}\\[4pt]
$k_{\rm E},\kappa_{\rm E}$&Diagonal-flat extension exponents; the quantitative theorem permits $k_{\rm E}=20$ and $\kappa_{\rm E}=k+32$.&\cref{thm:quantitative-flat-extension}\\[4pt]
$\mathsf p,\wp$&Uniform insertion exponent and analytic plate-density exponent $\wp=\mathsf p+5$.&\cref{thm:uniform-one-field-insertion}; \cref{thm:plate-density-from-insertion}\\[4pt]
$G_{\varphi,\psi},\mu_\varphi$\newline$\Gamma_\varphi,L_s$&Transversely smeared planar kernel, positive spectral measure, Laplace norm, and spectral boost.&\eqref{eq:angular-two-point-elimination-TU2}; \eqref{eq:angular-two-point-elimination-TU5}; \eqref{eq:angular-two-point-elimination-TU6}; \eqref{eq:angular-two-point-elimination-TU25}\\[4pt]
$\vartheta_*,u(d),c_m,b_n$&Fixed positive-shape threshold, ultraviolet error function, angular coefficients, and largest nonconstant angular frequency.&\eqref{eq:angular-two-point-elimination-TU9}; \eqref{eq:angular-two-point-elimination-TU10}; \eqref{eq:angular-harmonic-closure-HC1}; \hyperref[notdef:harmonic-degree]{harmonic degree}\\[4pt]
$\mathcal G^{\rm tor}$&Limits of periodically lifted torus measures at fixed spacing and coupling.&\hyperref[notdef:torus-limit-states]{torus-limit states}; \eqref{cc:specification}\\[4pt]
$\mathcal G_j,\mathcal E_j,D_j$&Translation-invariant DLR class, its ergodic states, and the decomposition law of the centering state on those states.&\hyperref[notdef:invariant-dlr]{invariant DLR states}; \hyperref[notdef:ergodic-decomposition]{ergodic decomposition}\\[4pt]
$\mathcal I,\operatorname{Pl}_\nu$\newline$\mathcal N_{\nu,r}$&Invariant sigma-field, conditional-expectation plateau pairing, and reflected size in the lattice plateau estimate.&\hyperref[notdef:invariant-plateau]{plateau quantities}; \cref{thm:lattice-plateau}\\[4pt]
$\mathsf X,\lambda,\sigma_\xi$&Compact support of the phase-coordinate law, that probability law, and its continuous component evaluations.&\hyperref[notdef:phase-coordinates]{phase coordinates}\\[4pt]
$\theta,\mathcal H,H,P$\newline$\Omega,P_0$&Time reflection, Hilbert completion of its positive form, time and momentum generators, empty-tuple vector, and projection onto $\ker H$.&\eqref{eq:euclidean-hilbert-bridge-HB1}; \eqref{eq:euclidean-hilbert-bridge-HB3}; \cref{ch:phase-vacuum}\\[4pt]
$\mathfrak U,\mathfrak B,W_4$&Ward defect, shell insertion, and signed coordinate permutation group, with the specified covariance and flux conventions.&\hyperref[notdef:ward-insertion]{Ward insertion}; \eqref{ph:W4}\\[4pt]
\end{longtable}
}

\subsection*{Conventions}

Lengths are physical unless a level or local lattice unit is stated.
In particular, $a_{K-j}$ and $\eta_j$ express different unit conventions.
An adjoint always uses the stated inner products on both spaces.

Constants may depend on fixed geometric and analytic parameters.
The stated parameter dependence and uniformity are part of each estimate.
A radius used for coefficient extraction is the radius in that estimate.

Concatenation of test tuples is used only when the resulting supports
remain separated. Centering changes the first cumulant and leaves higher
cumulants unchanged. Two-state moments retain the stated centering state.

The letters $K,R,P,\rho,N$ have several local uses, distinguished above
by their arguments and spaces. The window exponent is a local H\"older
exponent; it is independent of the field analyticity radii.

In the usual Wilson coupling convention, $\beta=2N/g_{\mathrm{std}}^2$.
All logarithms are natural. Near-identity group logarithms use the
branches specified in the defining chart. A pair current is independent
of its field coordinate unless the pair is explicitly restricted to
a background graph. Each use of a construction retains its stated
existence, domain, and uniformity hypotheses.

% END MODULE chapters/notation_index_v2.tex

% BEGIN LAYERED EXPANSION analytic_notation_index.tex
\subsection*{Analytic operators and domains}

The following symbols refer to the constructions
in the analytic estimates. Their transposes
use the physical measures unless a counting
normalization is displayed.

{\small
\setlength{\tabcolsep}{4pt}
\renewcommand{\arraystretch}{1.08}
\begin{longtable}{@{}>{\raggedright\arraybackslash}p{.20\linewidth}
>{\raggedright\arraybackslash}p{.45\linewidth}
>{\raggedright\arraybackslash}p{.28\linewidth}@{}}
\hline
\textbf{Notation}&\textbf{Meaning}&\textbf{Definition}\\
\hline
\endfirsthead
\hline
\textbf{Notation}&\textbf{Meaning}&\textbf{Definition}\\
\hline
\endhead
\hline
\endfoot
\(Y,Y^\flat,Z\)
&Balanced first-order site space, its source space
under pairing with \(\ell^2u\), and the site
second-difference energy space.
&\eqref{lcu:site-norms}.\\[4pt]
\(\mathcal A_\alpha,\mathcal F_p\)
&Strong field space and normalized finite-\(p\)
force space. A cell force has norm
\(\xi^{3-4/p}\|f\|_p\).
&\eqref{psr:strong-norm}; \eqref{scf:force-cell-norm}.\\[4pt]
\(\widehat F,\widehat M\)
&Site columns and their positive Gram matrix
in counting top coordinates; the column
injection is \(q^\sharp\Xi^{-4}\).
&\eqref{ncw:normalized-site-columns}.\\[4pt]
\(\Lambda,S,\Pi_U\)
&Linear gauge restoration, its complement
\(S=1-D\Lambda\), and the selected finite
gauge chain with derivative \(S\).
&\eqref{cng:restoration-objects};
\eqref{cng:finite-gauge-map}.\\[4pt]
\(G_F,G_0,G,G_1\)
&Positive field primitive, projected zero-current
inverse, first current inverse, and second
current inverse.
&\eqref{ncw:zero-current-Woodbury};
\eqref{ncw:first-inverse-series};
\eqref{ncw:second-inverse-series}.\\[4pt]
\(G_0^a,G^a,H^a\);
\newline\(\omega^a,G_1^a,H_1^a\)
&Weighted canonical inverses and lifts.
The multiplier is \((H^a)^\sharp J\),
formed before the second inverse.
Their auxiliary field output has two
finite-\(p\) differences.
&\eqref{wcg:G0-definition};
\eqref{wcg:first-definition};
\eqref{wcg:omega-definition};
\eqref{wcg:second-definition}.\\[4pt]
\(H,H_1,\omega_*\)
&First and second row lifts, and the multiplier
readout \(H^\sharp J\). The readout uses
the first lift.
&\eqref{ncw:first-lift};
\eqref{ncw:second-row-lift};
\eqref{ncw:bare-multiplier-split}.\\[4pt]
\(K_4,G_4,M_4\)
&Fourth-order site precision, its inverse,
and the counting coarse matrix \(qG_4q^*\).
They give a new strong family for \(G_sR\).
&\eqref{bhr:fourth-order-operator};
\eqref{bhr:coarse-M4};
\eqref{bhr:GsR-identity}.\\[4pt]
\(\mathcal S_\alpha,Y^{2,\alpha}\)
&Fourth-order source class
\(\Delta C^\alpha_2+D^\sharp L^\infty_3+L^\infty_4\)
and its strong site output space.
&\eqref{bhr:source-class};
\eqref{dsc:site-gauge-norm}.\\[4pt]
\(\mathcal L_{U,J},A(b),C(B)\)
&Selected auxiliary action, its constrained
potential, and its finite gauge completion.
The potential has \(D_bA(0)=H_1\);
the literal completion has
\(D_BC(0)=SH_1Db_U(0)\).
&\eqref{cng:auxiliary-action};
\eqref{cng:auxiliary-system};
\eqref{cng:reduced-first-jets}.\\[4pt]
\(\Gamma_P,K_\tau(P)\)
&Mass-subtracted coarse form and its outer
plain-constraint precision, including the
multiplier times the second chart derivative.
&\eqref{neg:off-current-outer-completion}.\\[4pt]
\(\mathcal W_p,\mathcal L_{0,a}\)
&Field norm with two physical \(L^p\)
differences, and the full weighted
trial saddle. The principal perturbation
acts from \(\mathcal W_p\) to \(\mathcal F_p\).
&\eqref{wca:Wp-norms};
\eqref{wca:trial-joint-operator};
\eqref{wca:principal-Neumann-bound}.\\[4pt]
\(M_{\rm c},M_{\rm G},M_{\rm h}\)
&Carrier cube size, Gaussian parameter cube
size, and a prescribed history cube size.
Their separation requirements are stated
at their respective uses.
&\eqref{ncw:family-norm};
\eqref{acg:reblocked-parameter-cost}.\\[4pt]
\(D_W(s),R_W(s)\)
&Local Gaussian covariance and its root.
At complete parameters they give the
actual matrix marginal.
&\eqref{acg:local-Gaussian-covariance};
\eqref{acg:actual-marginal-identity}.\\[4pt]
\(p_j^{\rm lit},\mathsf p_j\)
&Literal intermediate mean of the physical
critical orbit. Its raw bond amplitude is
bounded by \(C/L\) times the output amplitude.
&\eqref{hfm:uniform-amplitude-coarse-contraction};
\eqref{lcf:literal-centre}.\\[4pt]
\(e_{s,z},T_{i,k}\)
&Original fine-plaquette action pointing and
complete physical transport with exact
hat refinement.
&\eqref{hit:actual-classical-pointed-density};
\eqref{lcf:transport-definition}.\\[4pt]
\(\mathfrak H_{r,\kappa,\gamma_{\rm per}}\)
&Compatible pointed interaction space:
fixed analytic radius, carrier decay and
finite-period tail bound.
&\eqref{lcf:compatible-family-norm}.\\[4pt]
\(\mathbf F_j(x,h),N_j\)
&Vacuum-subtracted formation map and physical
Gaussian logarithm in the literal regular
step. \(N_j\) is independent of \(x,h\).
&\eqref{lcf:formation-map};
\eqref{lcf:Gaussian-stage}.\\[4pt]
\(\mathcal B_{j,e}\)
&Complete rooted sphere family:
the moving-centre logarithmic weight
jet minus the fixed-kernel jet.
Its six curvature readouts vanish.
&\eqref{fcw:pointed-boundary-source};
\cref{cor:fcw-zero-curvature-readout}.\\[4pt]
\(\widetilde\beta_a^{\mu\nu}(z)\)
&Weighted massive Wilson coefficient in one
plane, with vector regulator \(z\mathsf A\)
before congruence.
&\eqref{amc:determinant-density};
\eqref{amc:weighted-logarithm}.\\[4pt]
\(x_j,h_j,\|\cdot\|_{j,*}\)
&Tuned coupling, zero-readout regular profile,
and horizon-adapted profile norm.
&\eqref{hrt:exact-recursion};
\eqref{hrt:adapted-norm}.\\[4pt]
\(P_{j\ell},U_{ji},B^+,K^+\)
&Backward scalar and forward profile
propagators of the mixed-end linear
response, with positive majorizing kernels.
&\eqref{rts:propagators};
\eqref{rts:positive-kernels}.\\[4pt]
\(T_N,O_N,\nu_{N,t}\)
&Bare tuned scalar insertion, its
nonconstant conditional coarse image,
and the capped regular terminal law.
&\eqref{rsi:transferred-observable};
\eqref{rsi:infrared-measure}.\\[4pt]
\(\delta\alpha_N,\psi_N\)
&Bare tuning with its value at the
reference endpoint subtracted, and
its local inverse in the bare scalar
source. Neither subtracts a random field.
&\cref{lem:rsi-finite-dimensional-extraction,cor:rsi-source-ratio-subsequence}.\\[4pt]
\end{longtable}
}

% END LAYERED EXPANSION analytic_notation_index.tex

% BEGIN MODULE chapters/argument_index.tex
\subsection{Proof references}
\label{app:proof-references}

The table locates the component arguments. Each result is used with the
hypotheses stated at its cited location.

{\small
\setlength{\tabcolsep}{4pt}
\renewcommand{\arraystretch}{1.08}
\begin{longtable}{@{}>{\raggedright\arraybackslash}p{.23\linewidth}
>{\raggedright\arraybackslash}p{.72\linewidth}@{}}
\hline
\textbf{Proof task}&\textbf{Arguments and precise references}\\
\hline
\endfirsthead
\hline
\textbf{Proof task}&\textbf{Arguments and precise references}\\
\hline
\endhead
\hline
\endfoot

Projected Landau inverse
&Local gauge regularity in \cref{ch:bg-local-coulomb}; the all-mesh projected gap and actual-slice saddle inverse in \cref{ch:bg-single-scale-landau}; curvature-only periodic localization in \cref{ch:bg-single-scale-localization}.\\[4pt]
Normalized interacting extraction
&Positive history aggregation, weighted exceptional moments and exterior cumulants in \cref{ch:positive-history-comparison}; mixed source derivatives in \cref{ch:effective-source-jets}; the complete-oscillation Gaussian criterion in \cref{ch:finite-window-gaussian-law}.\\[4pt]
Four-point and Ward calculations
&Quantitative free-kernel signs and shape comparison in \cref{ch:four-point-kernel-recovered}; exact Haar divergence and contact terms in \cref{ch:exact-displacement-ward}; off-shell invariant classification and the weighted smooth Wilson expansion in \cref{ch:classical-invariants-wilson}.\\[4pt]

Averaging and the flat reference
&Pivot coordinates, the contour-gauge gap, bounded lifts, and exponential
kernels: \cref{ch:block-average-pivot-chart,ch:flat-contour-precision-gap,ch:flat-whitney-comparison,ch:flat-exponential-kernels}.\\[5pt]
First constrained background
&Nonflat first-step precision, the regular critical orbit, and the
relative-current extension:
\cref{ch:first-step-nonflat-precision,ch:first-step-background-expanded,ch:explicit-relative-current-extension}.\\[5pt]
Multiscale backgrounds
&Pairings, core domains, hierarchy identities, nonlinear responses,
and local H\"older estimates:
\cref{ch:domain-core-ledger,ch:multilevel-backgrounds,ch:nonlinear-response-reductions,ch:discrete-holder-estimate}.
Layered precision is treated in \cref{ch:layered-current-precision}.\\[5pt]
Local interactions
&The member's quadratic counterterm is given by \eqref{ph:extractor}
and \eqref{ph:own-counterterm}. The nonlinear window estimate is
\cref{thm:pointed-window-remainder}; insertion and the source bootstrap
are \cref{prop:finite-inserted-potential,prop:linear-source-bootstrap}.\\[5pt]
Histories and stability
&History counting and mass estimates: \cref{ch:history-mass-program}.
The ordinary stability upper bound and positive floor are
\eqref{oh:upper} and \eqref{oh:floor}, under their respective hypotheses.\\[5pt]
Flow and coupling comparison
&First-passage arithmetic: \cref{ch:flow-arithmetic}.
Joint-source comparison and the terminal coupling limit:
\cref{lem:polydisc-dimension-free-v7,thm:terminal-coupling-limit-v7}.
Weighted comparison domains: \cref{ch:weighted-matched-domains}.\\[5pt]
States and continuum limits
&Periodic-state mirror positivity: \cref{thm:amrp-all-periodic}.
Common state normalizers and factorial moment bounds:
\cref{thm:all-state-source-transfer,thm:graded-rooted-source}.
Subsequence compactness and flat-test extension:
\cref{thm:continuum-compactness,thm:quantitative-flat-extension}.\\[5pt]
Free and interacting two-point functions
&The free-kernel estimate is \eqref{eq:lcc-main}.
Conditional covariance and Gaussian cutoff calculations are in
\cref{ch:clean-conditional-law,ch:gaussian-cutoff-clean-law};
the extraction and common-normalizer transfer estimates are in
\cref{ch:two-point-extraction-tests}.\\[5pt]
Insertions and rotations
&Uniform insertion, plate density, the planar kernel, and harmonic
elimination:
\cref{thm:uniform-one-field-insertion,thm:plate-density-from-insertion,ch:angular-kernel-construction,thm:angular-two-point-elimination,thm:all-orders-angular-closure}.\\[5pt]
Fields and operator domains
&Energy-band fields: \cref{ch:h3-energy-band-fields}.
The matched scalar reconstruction is
\cref{thm:scalar-reconstruction-on-ordered-sector}; its representation
domains are in \cref{ch:reconstruction-representation-domains}.\\[5pt]
Phases and spectra
&The conditional probability kernel is \cref{thm:compact-conditional-law}.
Phase coordinates, the vacuum space, and the scalar direct integral are
treated in \cref{ch:phase-interface,ch:phase-vacuum,ch:scalar-phase-hilbert-bridge}.\\[5pt]
Gauge observables
&The cutoff orbit coordinates and continuum extension criteria are in
\cref{ch:su2-wilson-orbit-coordinates,ch:gauge-observable-extension-criteria}.
The scalar observable information is characterized in
\cref{ch:scalar-observable-information}.\\[5pt]
\end{longtable}
}

% END MODULE chapters/argument_index.tex

% BEGIN LAYERED EXPANSION analytic_proof_index.tex
\subsection{References for the analytic constructions}
\label{app:analytic-proof-references}

These references distinguish the completed
operators from their local coefficient
families, and the regular step from its
complementary history operations.
Each construction retains its stated
reference, geometry and parameter domain.

{\small
\setlength{\tabcolsep}{4pt}
\renewcommand{\arraystretch}{1.08}
\begin{longtable}{@{}>{\raggedright\arraybackslash}p{.27\linewidth}
>{\raggedright\arraybackslash}p{.68\linewidth}@{}}
\hline
\textbf{Construction}&\textbf{Proof and scope}\\
\hline
\endfirsthead
\hline
\textbf{Construction}&\textbf{Proof and scope}\\
\hline
\endhead
\hline
\endfoot
Strong response at an independent current
&The actual Wilson stencil, finite-\(p\)
bootstrap and current adjoints are in
\cref{ch:proper-strong-response,ch:actual-current-completion}.
The completed strong estimates do not
require \(DJ\).\\[4pt]
Local positive inverses
&\Cref{lpw:positive-local-walks} uses genuine
Dirichlet compressions, slow cutoffs and
full geometry carriers. The positive coarse
inverse is \cref{prop:cci-positive-coarse-inverse}.\\[4pt]
Named current families
&\Cref{prop:ncw-named-families} constructs
\(G_0,G,H,\omega_*,G_1,H_1\) in that order,
with normalized physical row transposes.
It retains the bare current readout separately.\\[4pt]
Strong restoration coefficients
&\Cref{prop:bhr-full-H2-gap,prop:bhr-exact-restoration,prop:bhr-strong-restoration-family}
prove the fourth-order realization, its
exact \(G_sR\) identity and its new
\(C^{2,\alpha}\) coefficient family.\\[4pt]
Strong saddle coefficients
&\Cref{prop:scf-strong-current-families}
uses the finite-\(p\) force classes and
the strong restoration. Its current
ceiling may shrink with the fixed
H\"older exponent and carrier rate.\\[4pt]
Nonlinear current germ
&\Cref{prop:cng-auxiliary-germ,prop:cng-physical-matching,prop:cng-current-update}
give the selected auxiliary action,
its physical matching and its current
update. The \(H_1\), \(SH_1\) and literal
datum jets are distinct.\\[4pt]
All-degree nonlinear coefficients
&\Cref{prop:ngc-datum-connected-family,prop:ngc-literal-connected-family,prop:ngc-joint-connected-family}
give the convergent connected tree series
for that selected germ, including joint
base and datum variables. Local-data
suprema precede all carrier sums.\\[4pt]
Weighted connected germ
&\Cref{prop:wcg-weighted-linear-family,prop:wcg-connected-germ,prop:wcg-joint-connected-family}
construct the weighted principal and
current inversions before the actual
nonlinear tree series.
\Cref{prop:wcg-WCA-physical-restriction}
identifies its admitted physical
curvature-atlas branch.\\[4pt]
Actual hierarchical row density
&\Cref{prop:hfm-fixed-leg-density}
computes the one-step second jet and
the full chain sum. Uniformity in depth
does not make the nonlinear radius
independent of \(L\).\\[4pt]
Homogeneous critical fields
&\Cref{prop:hb-actual-critical-orbit}
constructs the physical orbit. Flat
reference descent is
\cref{prop:frd-flat-atlas-descent};
\cref{prop:rtc-nonstationary-trial}
also allows a current-regular trial.\\[4pt]
Curvature-only homogeneous data
&\Cref{prop:cta-uniform-trial,prop:cta-curvature-critical-atlas}
construct the trial from a smooth
quaternionic bundle, then correct its
actual mean and current. The atlas
allows large coarse holonomies and
nontrivial bundle topology.\\[4pt]
Weighted physical atlas
&\Cref{lem:wca-Wp-joint-inverse,prop:wca-weighted-curvature-atlas,prop:wca-outer-Gaussian-gap,prop:wca-complex-outer-sector}
retain two \(L^p\) differences before
perturbing the principal part. The
physical gap includes both constraint
curvatures and the row-mass subtraction.\\[4pt]
Outer nonflat Gaussian
&\Cref{neg:outer-Gaussian-gap,acg:canonical-Gaussian-admission}
prove the real physical gap and the
canonical local Gaussian extension
near that reference, retaining the
mass subtraction and chart curvature.\\[4pt]
Cutoff covariance
&\Cref{prop:rcs-cutoff-transfer-covariance}
uses the middle crossing pivots and
real colour-ball cutoff. It proves
covariance of the actual regular
integral before pointing or filing.\\[4pt]
Physical irrelevance
&\Cref{prop:hit-physical-irrelevant-transport}
uses the literal intermediate field,
complete Ward cancellation, block
interpolation and explicit finite-period
tails. Individual local pieces need
not be invariant.\\[4pt]
Fixed-radius formation
&\Cref{prop:lcf-buffered-formation,cor:lcf-common-radius-transport,cor:lcf-exact-profile-step}
admit the old interaction on its own
ball, prove the coupling and profile
slopes, and close one regular step
on the common radius.\\[4pt]
Fixed-kernel stress response
&\Cref{prop:fcw-frozen-weight-jet,cor:fcw-zero-curvature-readout}
compare the moving weight jet with
the fixed kernel by a normalized
rooted sphere family.
\Cref{prop:rds-regular-stress-decay}
uses the corrected profile forcing
in its coupling-weighted memory sum.\\[4pt]
Actual scalar one-loop limit
&The bare Wilson calculation is
\cref{mlc:main}; physical gauge change
is \cref{pgg:physical-RF12}.
The regulator and Taylor-vertex
comparisons are
\cref{ch:mrc-actual-comparison,ch:cac-classical-convergence,ch:gaussian-polarization-lipschitz}.\\[4pt]
Diagonal one-loop response
&\Cref{amc:weighted-massive-law,prop:ach-actual-anisotropic-background,prop:ach-Taylor-convergence,prop:agc-actual-anisotropic-increment}
retain the actual weight parameter,
its principal-part perturbation, the
weighted regulator and its differentiated
Gaussian polarization.\\[4pt]
Regular tuning
&\Cref{prop:hrt-regular-trajectory}
solves the mixed end conditions.
\Cref{hrt:endpoint-slope,hrt:higher-endpoint-jets}
differentiate one frozen presentation,
with constants uniform in the horizon.\\[4pt]
Bounded initial irrelevant profiles
&\Cref{prop:rts-bounded-UV-tuning,prop:rts-exponential-universality}
retune the same finite regular problem.
\Cref{rts:Green-envelope,rts:mixed-jets}
give the mixed-end inverse and all
fixed-order initial-profile jets.\\[4pt]
Scalar normalization and nonvanishing
&\Cref{prop:rsi-scalar-insertion,prop:rsi-regular-nontriviality}
retain the numerical vacuum derivative
before centering and prove a positive
variance in the capped regular law.
\Cref{prop:rsi-fixed-coarse-subsequence}
extracts a nondegenerate coarse
subsequence at fixed volume and cutoff.\\[4pt]
\end{longtable}
}

% END LAYERED EXPANSION analytic_proof_index.tex


% BEGIN MODULE chapters/bibliography.tex
\begin{thebibliography}{OSeiler78}\footnotesize

\bibitem[B82a]{B82a} T. Ba\l{}aban, $(\mathrm{Higgs})_{2,3}$ quantum fields in a finite volume. I. A lower bound, \emph{Communications in Mathematical Physics} \textbf{85} (1982), 603--626. \href{https://doi.org/10.1007/BF01403506}{doi:10.1007/BF01403506}.

\bibitem[B82b]{B82b} T. Ba\l{}aban, $(\mathrm{Higgs})_{2,3}$ quantum fields in a finite volume. II. An upper bound, \emph{Communications in Mathematical Physics} \textbf{86} (1982), 555--594. \href{https://doi.org/10.1007/BF01214890}{doi:10.1007/BF01214890}.

\bibitem[B83a]{B83a} T. Ba\l{}aban, $(\mathrm{Higgs})_{2,3}$ quantum fields in a finite volume. III. Renormalization, \emph{Communications in Mathematical Physics} \textbf{88} (1983), 411--445. \href{https://doi.org/10.1007/BF01213217}{doi:10.1007/BF01213217}.

\bibitem[B83b]{B83b} T. Ba\l{}aban, Regularity and decay of lattice Green's functions, \emph{Communications in Mathematical Physics} \textbf{89} (1983), 571--597. \href{https://doi.org/10.1007/BF01214744}{doi:10.1007/BF01214744}.

\bibitem[B84a]{B84a} T. Ba\l{}aban, Propagators and renormalization transformations for lattice gauge theories. I, \emph{Communications in Mathematical Physics} \textbf{95} (1984), 17--40. \href{https://doi.org/10.1007/BF01215753}{doi:10.1007/BF01215753}.

\bibitem[B84b]{B84b} T. Ba\l{}aban, Propagators and renormalization transformations for lattice gauge theories. II, \emph{Communications in Mathematical Physics} \textbf{96} (1984), 223--250. \href{https://doi.org/10.1007/BF01240221}{doi:10.1007/BF01240221}.

\bibitem[B85a]{B85a} T. Ba\l{}aban, Averaging operations for lattice gauge theories, \emph{Communications in Mathematical Physics} \textbf{98} (1985), 17--51. \href{https://doi.org/10.1007/BF01211042}{doi:10.1007/BF01211042}.

\bibitem[B85b]{B85b} T. Ba\l{}aban, Spaces of regular gauge field configurations on a lattice and gauge fixing conditions, \emph{Communications in Mathematical Physics} \textbf{99} (1985), 75--102. \href{https://doi.org/10.1007/BF01466594}{doi:10.1007/BF01466594}.

\bibitem[B85c]{B85c} T. Ba\l{}aban, Propagators for lattice gauge theories in a background field, \emph{Communications in Mathematical Physics} \textbf{99} (1985), 389--434. \href{https://doi.org/10.1007/BF01240355}{doi:10.1007/BF01240355}.

\bibitem[B85d]{B85d} T. Ba\l{}aban, Ultraviolet stability of three-dimensional lattice pure gauge field theories, \emph{Communications in Mathematical Physics} \textbf{102} (1985), 255--275. \href{https://doi.org/10.1007/BF01229380}{doi:10.1007/BF01229380}.

\bibitem[B85e]{B85e} T. Ba\l{}aban, The variational problem and background fields in renormalization group method for lattice gauge theories, \emph{Communications in Mathematical Physics} \textbf{102} (1985), 277--309. \href{https://doi.org/10.1007/BF01229381}{doi:10.1007/BF01229381}.

\bibitem[B87]{B87} T. Ba\l{}aban, Renormalization group approach to lattice gauge field theories. I. Generation of effective actions in a small field approximation and a coupling constant renormalization in four dimensions, \emph{Communications in Mathematical Physics} \textbf{109} (1987), 249--301. \href{https://doi.org/10.1007/BF01215223}{doi:10.1007/BF01215223}.

\bibitem[B88a]{B88a} T. Ba\l{}aban, Renormalization group approach to lattice gauge field theories. II. Cluster expansions, \emph{Communications in Mathematical Physics} \textbf{116} (1988), 1--22. \href{https://doi.org/10.1007/BF01239022}{doi:10.1007/BF01239022}.

\bibitem[B88b]{B88b} T. Ba\l{}aban, Convergent renormalization expansions for lattice gauge theories, \emph{Communications in Mathematical Physics} \textbf{119} (1988), 243--285. \href{https://doi.org/10.1007/BF01217741}{doi:10.1007/BF01217741}.

\bibitem[B89a]{B89a} T. Ba\l{}aban, Large field renormalization. I. The basic step of the $\mathbb R$ operation, \emph{Communications in Mathematical Physics} \textbf{122} (1989), 175--202. \href{https://doi.org/10.1007/BF01257412}{doi:10.1007/BF01257412}.

\bibitem[B89b]{B89b} T. Ba\l{}aban, Large field renormalization. II. Localization, exponentiation, and bounds for the $R$ operation, \emph{Communications in Mathematical Physics} \textbf{122} (1989), 355--392. \href{https://doi.org/10.1007/BF01238433}{doi:10.1007/BF01238433}.

\bibitem[CZ52]{CZ52} A. P. Calder\'on and A. Zygmund,
On the existence of certain singular integrals,
\emph{Acta Mathematica} \textbf{88} (1952), 85--139.
\href{https://doi.org/10.1007/BF02392130}{doi:10.1007/BF02392130}.

\bibitem[Edm00]{Edmundo2000}
M.~J. Edmundo, \emph{An introduction to o-minimal structures},
arXiv:math/0012051v2 (2000).
\url{https://arxiv.org/abs/math/0012051}.

\bibitem[FILS78]{FILS78} J. Fr\"ohlich, R. Israel, E. H. Lieb and B. Simon, Phase transitions and reflection positivity. I. General theory and long range lattice models, \emph{Communications in Mathematical Physics} \textbf{62} (1978), 1--34. \href{https://doi.org/10.1007/BF01940327}{doi:10.1007/BF01940327}.

\bibitem[JW]{JW} A. Jaffe and E. Witten, Quantum Yang--Mills theory, in \emph{The Millennium Prize Problems} (J. Carlson, A. Jaffe and A. Wiles, eds.), Clay Mathematics Institute, Cambridge, MA, and American Mathematical Society, Providence, RI, 2006, 129--152. \url{https://www.claymath.org/wp-content/uploads/2022/06/yangmills.pdf}.

\bibitem[Kallenberg02]{Kallenberg2002}
O. Kallenberg, \emph{Foundations of Modern Probability}, second edition,
Probability and Its Applications, Springer, New York, 2002.
\href{https://doi.org/10.1007/978-1-4757-4015-8}{doi:10.1007/978-1-4757-4015-8}.

\bibitem[Kur98]{Kurdyka1998}
K.~Kurdyka, \emph{On gradients of functions definable in o-minimal structures},
Ann. Inst. Fourier \textbf{48} (1998), no.~3, 769--783.
\href{https://doi.org/10.5802/aif.1638}{doi:10.5802/aif.1638}.

\bibitem[DLMF]{DLMF-pivot}
NIST Digital Library of Mathematical Functions, formula 4.22.3.
\url{https://dlmf.nist.gov/4.22.E3}.

\bibitem[Hat02]{Hatcher-pivot}
A.~Hatcher, \emph{Algebraic topology}, Cambridge University Press,
Cambridge, 2002.

\bibitem[KP86]{KP86} R. Koteck\'y and D. Preiss, Cluster expansion for abstract polymer models, \emph{Communications in Mathematical Physics} \textbf{103} (1986), 491--498. \href{https://doi.org/10.1007/BF01211762}{doi:10.1007/BF01211762}.

\bibitem[OS73]{OS73} K. Osterwalder and R. Schrader, Axioms for Euclidean Green's functions, \emph{Communications in Mathematical Physics} \textbf{31} (1973), 83--112. \href{https://doi.org/10.1007/BF01645738}{doi:10.1007/BF01645738}.

\bibitem[OS75]{OS75} K. Osterwalder and R. Schrader, Axioms for Euclidean Green's functions II, \emph{Communications in Mathematical Physics} \textbf{42} (1975), 281--305. \href{https://doi.org/10.1007/BF01608978}{doi:10.1007/BF01608978}.

\bibitem[OSeiler78]{OSeiler78} K. Osterwalder and E. Seiler, Gauge field theories on a lattice, \emph{Annals of Physics} \textbf{110} (1978), 440--471. \href{https://doi.org/10.1016/0003-4916(78)90039-8}{doi:10.1016/0003-4916(78)90039-8}.

\bibitem[Wilson74]{Wilson74} K. G. Wilson, Confinement of quarks, \emph{Physical Review D} \textbf{10} (1974), 2445--2459. \href{https://doi.org/10.1103/PhysRevD.10.2445}{doi:10.1103/PhysRevD.10.2445}.

\end{thebibliography}

% END MODULE chapters/bibliography.tex

\end{document}
